% Figure 1 (F1 of docs/FIGURE_PLAN.md): the design of the study, in TikZ.
% % Figure 1 (F1 of docs/FIGURE_PLAN.md): the design of the study, in TikZ.
% % Figure 1 (F1 of docs/FIGURE_PLAN.md): the design of the study, in TikZ.
% % Figure 1 (F1 of docs/FIGURE_PLAN.md): the design of the study, in TikZ.
% \input{fig_f1} from main.tex; this file holds one figure environment and nothing else.
%
% WHAT main.tex MUST LOAD (this file loads nothing itself):
%   \usepackage{tikz}
%   \usetikzlibrary{arrows.meta}          % the only TikZ library used (arrow tips)
%   \definecolor{c188}{HTML}{1F77B4}      % the configuration colours of experiments/FIGS/
%   \definecolor{c162}{HTML}{FF7F0E}      % mlst_style.py (matplotlib's tab10 in order, as in
%   \definecolor{c43}{HTML}{2CA02C}       % Sophon's figures); the variants are drawn in their
%   \definecolor{c17}{HTML}{D62728}       % label set's colour with a dotted border, as the
%   \definecolor{crand}{HTML}{17BECF}     % plotted figures draw variants dotted
%   \definecolor{cflav}{HTML}{9467BD}
%   \input{results_generated}             % the counts below are its macros
% Visual language: Sophon's Fig. 1 (arXiv:2405.12972): the document's serif font, light fills,
% rounded boxes, light grey panels behind the stages (PI 2026-11-02). The text is set in the
% document's font at the printed sizes 8.5-10 pt (MLST 8-12 pt), never scaled.
% Text width: 15 cm (MLST full width). The picture is 15.0 cm wide and is drawn at its final
% size, never scaled, so the font sizes hold.
% Optional: \def\FoneDropXbx{} before \input{fig_f1} removes X->bb/X->bs from panel (c)
% (PI 2026-11-02: if the samples are not staged by 12 Nov, the X->bb/X->bs section goes).
%
% Every class box is drawn from the label maps (checked by hand, 2026-11-02):
%   (a) configs/labelmaps/rung_label_maps.v1.csv, columns L188, L162, R42_Q1, R16_Q1. Rows: the ten
%       decays X->QQ/gg in the order bb bc bq | cc cs cq | ss sq qq gg (the 43-class b / c / light
%       groups are then contiguous), then ee mm | tauh-taue tauh-taum, then the 27 QCD classes.
%   (b) column: configs/labelmaps/rand_label_map.v2.csv, RAND2_p1 on the same rows: groups
%       0 0 0 0 0 9 9 9 0 3 | 11 11 | 15 4 | QCD, lettered A..F in order of first appearance.
%       strips: configs/labelmaps/flavour_pair_map.v2.csv, the 22 classes of orbit X_YY_QQQQ.
%       Order with a b quark: bbbb bbcc bbss bbqq qqqb bbcq ccbs ccbq qqbc qqbs bcsq;
%       without: cccc ccss ccqq ssss ssqq qqqq qqqc qqqs ccsq sscq qqcs.
%       FLAV_F0: all 22 in group 9. FLAV_F1: the 11 with a b in group 11, the rest 9.
%       FLAV_F1R: group 11 = bbss bbqq qqqb bbcq ccbq qqbc (with b, positions 2-5, 7, 8 from 0)
%       and ccss ccqq ssss qqqc sscq (without, positions 1-3, 6, 9); filled square = group 11.
% Geometry (cm): row height 0.35 from y = -1.18, gaps of 0.14 between the quark/gluon, lepton
% and QCD blocks, the QCD row 0.70; boxes inset 0.025 top and bottom; columns 1.45 wide with
% left edges 1.75, 3.28, 4.81, 6.34 (hierarchy) and 8.25 (random partition).
\providecommand{\Kfull}{\VocabSizeLoneeighteight}
\providecommand{\Kres}{\VocabSizeLonesixtwo}
\providecommand{\Kmid}{\VocabSizeRfourtwoQone}
\providecommand{\Kcoarse}{\VocabSizeRonesixQone}
\begin{figure}[tbp]
\centering
\ifdim\linewidth<14.99cm\typeout{fig_f1.tex: WARNING: the drawing is 15 cm wide but \string\linewidth\space
  is \the\linewidth; set a 15 cm text width (it is not scaled, to keep its 8.5 pt text)}\fi
\def\FoneText{\rmfamily\fontsize{8.5}{9.5}\selectfont}%
\def\FoneMain{\rmfamily\fontsize{9}{10.5}\selectfont}%
\def\FoneBold{\rmfamily\bfseries\fontsize{9}{10.5}\selectfont}%
\def\FoneTitle{\rmfamily\bfseries\fontsize{10}{12}\selectfont}%
\ifdefined\FoneDropXbx\def\FoneXbx{}\else\def\FoneXbx{\\$X\to b\bar b$, $X\to bs$}\fi
\begin{tikzpicture}[
  x=1cm, y=1cm,
  font=\FoneMain,
  ttl/.style={anchor=base west, font=\FoneTitle, inner sep=0pt},
  cls/.style={draw=black!55, line width=0.45pt, rounded corners=1.2pt},
  hdr/.style={draw=#1!70!black, fill=#1!22, line width=0.6pt, rounded corners=2.5pt, align=center,
              font=\FoneText, inner sep=0pt},
  inbox/.style={align=center, font=\FoneText, inner sep=0pt},
  letter/.style={font=\FoneBold, inner sep=0pt},
  rowlab/.style={anchor=east, inner sep=0pt},
  txt/.style={align=left, font=\FoneText, inner sep=0pt},
  sq/.style={draw=black!65, line width=0.4pt},
  chip/.style={draw=#1!70!black, fill=#1!22, line width=0.6pt, rounded corners=3pt,
               align=center, font=\FoneText, inner sep=0pt, anchor=north west},
  stage/.style={draw=black!45, fill=black!4, line width=0.6pt, rounded corners=3pt, align=left,
                font=\FoneText, inner sep=0.1cm},
  flow/.style={-{Latex[length=1.8mm, width=1.4mm]}, line width=0.7pt, black!75},
]
% ================================================================ (a) the class hierarchy
\node[ttl] at (0,0) {(a) Class hierarchy};
% faint bands on every other row, drawn first so the boxes sit on top
\foreach \rowtop in {-1.53,-2.23,-2.93,-3.63,-4.33,-5.17,-5.87}
  \fill[black!4] (0,\rowtop) rectangle (9.70,\rowtop-0.35);
% column headers
\node[rowlab, align=right, font=\FoneText] at (1.65,-0.68) {JetClass-II\\class};
\node[hdr=c188, minimum width=1.45cm, minimum height=0.80cm, anchor=north west] at (1.75,-0.28) {\Kfull\\classes};
\node[hdr=c162, minimum width=1.45cm, minimum height=0.80cm, anchor=north west] at (3.28,-0.28) {\Kres\\classes};
\node[hdr=c43,  minimum width=1.45cm, minimum height=0.80cm, anchor=north west] at (4.81,-0.28) {\Kmid\\classes};
\node[hdr=c17,  minimum width=1.45cm, minimum height=0.80cm, anchor=north west] at (6.34,-0.28) {\Kcoarse\\classes};
% row labels
\node[rowlab] at (1.65,-1.355) {$X\to b\bar b$};
\node[rowlab] at (1.65,-1.705) {$X\to bc$};
\node[rowlab] at (1.65,-2.055) {$X\to bq$};
\node[rowlab] at (1.65,-2.405) {$X\to c\bar c$};
\node[rowlab] at (1.65,-2.755) {$X\to cs$};
\node[rowlab] at (1.65,-3.105) {$X\to cq$};
\node[rowlab] at (1.65,-3.455) {$X\to s\bar s$};
\node[rowlab] at (1.65,-3.805) {$X\to sq$};
\node[rowlab] at (1.65,-4.155) {$X\to q\bar q$};
\node[rowlab] at (1.65,-4.505) {$X\to gg$};
\node[rowlab] at (1.65,-4.995) {$X\to ee$};
\node[rowlab] at (1.65,-5.345) {$X\to\mu\mu$};
\node[rowlab] at (1.65,-5.695) {$X\to\tau_h\tau_e$};
\node[rowlab] at (1.65,-6.045) {$X\to\tau_h\tau_\mu$};
\node[rowlab] at (1.65,-6.710) {QCD jets};
% 188 and 162 classes: every resonant decay is its own class
\foreach \rowtop in {-1.18,-1.53,-1.88,-2.23,-2.58,-2.93,-3.28,-3.63,-3.98,-4.33,-4.82,-5.17,-5.52,-5.87} {
  \draw[cls, fill=c188!15] (1.75,\rowtop-0.025) rectangle (3.20,\rowtop-0.325);
  \draw[cls, fill=c162!15] (3.28,\rowtop-0.025) rectangle (4.73,\rowtop-0.325);
}
\draw[cls, fill=c188!15] (1.75,-6.385) rectangle (3.20,-7.035);
\node[inbox] at (2.475,-6.71) {\LitJetClassTwoNQcdClasses\\classes};
\draw[cls, fill=c162!15] (3.28,-6.385) rectangle (4.73,-7.035);
\node[inbox] at (4.005,-6.71) {1 class};
% 43 classes: decay topology x {b, else c, else light quark or gluon}; e and mu merged
\draw[cls, fill=c43!15] (4.81,-1.205) rectangle (6.26,-2.205);
\node[inbox] at (5.535,-1.705) {with $b$};
\draw[cls, fill=c43!15] (4.81,-2.255) rectangle (6.26,-3.255);
\node[inbox] at (5.535,-2.755) {with $c$,\\no $b$};
\draw[cls, fill=c43!15] (4.81,-3.305) rectangle (6.26,-4.655);
\node[inbox] at (5.535,-3.980) {$u$, $d$, $s$\\or $g$};
\draw[cls, fill=c43!15] (4.81,-4.845) rectangle (6.26,-5.495);
\node[inbox] at (5.535,-5.170) {$\ell\ell$};
\draw[cls, fill=c43!15] (4.81,-5.545) rectangle (6.26,-6.195);
\node[inbox] at (5.535,-5.870) {$\tau_h\tau_\ell$};
\draw[cls, fill=c43!15] (4.81,-6.385) rectangle (6.26,-7.035);
\node[inbox] at (5.535,-6.710) {1 class};
% 17 classes: number and type of prongs, no flavor
\draw[cls, fill=c17!15] (6.34,-1.205) rectangle (7.79,-4.655);
\node[inbox] at (7.065,-2.930) {two\\quarks\\or gluons};
\draw[cls, fill=c17!15] (6.34,-4.845) rectangle (7.79,-5.495);
\node[inbox] at (7.065,-5.170) {$\ell\ell$};
\draw[cls, fill=c17!15] (6.34,-5.545) rectangle (7.79,-6.195);
\node[inbox] at (7.065,-5.870) {$\tau_h\tau_\ell$};
\draw[cls, fill=c17!15] (6.34,-6.385) rectangle (7.79,-7.035);
\node[inbox] at (7.065,-6.710) {1 class};
% what each merge removes
\node[anchor=north west, font=\FoneBold, inner sep=0pt] at (0,-7.27) {What each merge removes};
\node[txt, anchor=north west, text width=6.04cm] (m1) at (1.75,-7.65)
  {the $b$, $c$ and $s$ content of the QCD jets};
\node[txt, anchor=north west, text width=6.04cm] (m2) at ([yshift=-0.10cm]m1.south west)
  {the number of $b$ and $c$ quarks, light-quark flavor, quark against gluon, and lepton
   flavor ($e$ against $\mu$)};
\node[txt, anchor=north west, text width=6.04cm] (m3) at ([yshift=-0.10cm]m2.south west)
  {quark flavor ($b$, $c$ or light); the number and type of prongs remain};
\node[txt, anchor=north east] at ([xshift=-0.10cm]m1.north west) {\Kfull\,$\to$\,\Kres};
\node[txt, anchor=north east] at ([xshift=-0.10cm]m2.north west) {\Kres\,$\to$\,\Kmid};
\node[txt, anchor=north east] at ([xshift=-0.10cm]m3.north west) {\Kmid\,$\to$\,\Kcoarse};
% ================================================================ (b) the controls
\node[ttl] at (8.25,0) {(b) Controls, \Kcoarse{} classes each};
% the first random partition on the rows of (a); one letter = one class
\node[hdr=crand, minimum width=1.45cm, minimum height=0.80cm, anchor=north west] at (8.25,-0.28) {random\\partition};
\draw[cls, fill=crand!15] (8.25,-1.205) rectangle (9.70,-2.905);
\node[letter] at (8.975,-2.055) {A};
\draw[cls, fill=crand!15] (8.25,-2.955) rectangle (9.70,-3.955);
\node[letter] at (8.975,-3.455) {B};
\draw[cls, fill=crand!15] (8.25,-4.005) rectangle (9.70,-4.305);
\node[letter] at (8.975,-4.155) {A};
\draw[cls, fill=crand!15] (8.25,-4.355) rectangle (9.70,-4.655);
\node[letter] at (8.975,-4.505) {C};
\draw[cls, fill=crand!15] (8.25,-4.845) rectangle (9.70,-5.495);
\node[letter] at (8.975,-5.170) {D};
\draw[cls, fill=crand!15] (8.25,-5.545) rectangle (9.70,-5.845);
\node[letter] at (8.975,-5.695) {E};
\draw[cls, fill=crand!15] (8.25,-5.895) rectangle (9.70,-6.195);
\node[letter] at (8.975,-6.045) {F};
\draw[cls, fill=crand!15] (8.25,-6.385) rectangle (9.70,-7.035);
\node[inbox] at (8.975,-6.710) {1 class};
% how the random partitions are built
\node[hdr=crand, anchor=north west, text width=4.75cm, inner sep=0.1cm, align=left, font=\FoneBold]
  (rh) at (10.05,-0.28) {Random partitions (\RandVtwoNPartitions)};
\node[txt, anchor=north west, text width=4.95cm] (rb) at ([yshift=-0.12cm]rh.south west)
  {\textbullet\ groups of resonant classes drawn at random, plus one class of QCD jets\\[2pt]
   \textbullet\ group shares of training jets within \RandVtwoShareTol{} of the \Kcoarse-class groups\\[2pt]
   \textbullet\ best of \RandVtwoTrials{} draws: the fewest merges shared with the \Kcoarse-class level\\[2pt]
   \textbullet\ left: the first partition; boxes with the same letter are one class};
% how the random partitions with controlled flavor boundaries are built
\node[hdr=cflav, anchor=north west, text width=4.75cm, inner sep=0.1cm, align=left, font=\FoneBold]
  (fh) at ([yshift=-0.22cm]rb.south west) {Random partitions with controlled flavor boundaries};
\node[txt, anchor=north west, text width=4.95cm] (fn) at ([yshift=-0.12cm]fh.south west)
  {whole families of decays that differ only in quark flavor; two of them split one family,
   the \RandFlavSplitOrbitClasses{} decays \RandFlavSplitOrbit:};
\begin{scope}[shift={(fn.south west)}]
  % brackets over the decays with and without a b quark (squares at a 0.2 cm pitch)
  \node[txt, anchor=south] at (1.075,-0.45) {with $b$};
  \node[txt, anchor=south] at (3.425,-0.45) {without $b$};
  \draw[line width=0.5pt] (0,-0.54) -- (0,-0.47) -- (2.15,-0.47) -- (2.15,-0.54);
  \draw[line width=0.5pt] (2.35,-0.54) -- (2.35,-0.47) -- (4.50,-0.47) -- (4.50,-0.54);
  % no flavor boundary: all 22 in one class
  \node[txt, anchor=north west] at (0,-0.62) {no flavor boundary};
  \foreach \i in {0,...,10} {
    \draw[sq, fill=white] (0.2*\i,-1.10) rectangle ++(0.15,0.15);
    \draw[sq, fill=white] (2.35+0.2*\i,-1.10) rectangle ++(0.15,0.15);
  }
  % b-quark boundary: the decays with a b quark moved to another class
  \node[txt, anchor=north west] at (0,-1.22) {$b$-quark boundary};
  \foreach \i in {0,...,10} {
    \draw[sq, fill=cflav!70!black] (0.2*\i,-1.70) rectangle ++(0.15,0.15);
    \draw[sq, fill=white] (2.35+0.2*\i,-1.70) rectangle ++(0.15,0.15);
  }
  % random flavor boundary: the same number of decays, drawn at random
  \node[txt, anchor=north west] at (0,-1.82) {random flavor boundary};
  \foreach \i in {0,1,6,9,10}
    \draw[sq, fill=white] (0.2*\i,-2.30) rectangle ++(0.15,0.15);
  \foreach \i in {2,3,4,5,7,8}
    \draw[sq, fill=cflav!70!black] (0.2*\i,-2.30) rectangle ++(0.15,0.15);
  \foreach \i in {0,4,5,7,8,10}
    \draw[sq, fill=white] (2.35+0.2*\i,-2.30) rectangle ++(0.15,0.15);
  \foreach \i in {1,2,3,6,9}
    \draw[sq, fill=cflav!70!black] (2.35+0.2*\i,-2.30) rectangle ++(0.15,0.15);
  % key
  \draw[sq, fill=cflav!70!black] (0,-2.62) rectangle ++(0.15,0.15);
  \node[txt, anchor=west] (fkey) at (0.25,-2.545) {moved to another class};
\end{scope}
% ================================================================ (c) the pipeline
% (c) starts below the lower of the two columns above, whatever their text heights
\pgfpointanchor{m3}{south}\pgfgetlastxy{\FoneXa}{\FoneYa}%
\pgfpointanchor{fkey}{south}\pgfgetlastxy{\FoneXb}{\FoneYb}%
\pgfmathsetlengthmacro{\FoneYc}{min(\FoneYa,\FoneYb)-0.45cm}%
\coordinate (ctop) at (0pt,\FoneYc);
\begin{scope}[shift={(ctop)}]
  \node[ttl, anchor=north west] at (0,0) {(c) Pretraining and evaluation};
  % what every model of one seed shares
  \node[stage, text width=2.15cm] (seed) at (1.175,-1.96)
    {\textbf{Same for every model of one seed:} the training jets, their order and the
     initial weights, except the output layer};
  % one pretraining per label set: the configuration colours of every later figure
  \node[draw=black!45, fill=black!4, line width=0.6pt, rounded corners=3pt, minimum width=5.15cm,
        minimum height=3.08cm, anchor=north west, inner sep=0pt] (pre) at (2.81,-0.42) {};
  \node[anchor=north west, font=\FoneBold, inner sep=0pt] at (2.96,-0.52) {Pretraining on JetClass-II};
  \node[chip=c188, minimum width=1.13cm, minimum height=0.42cm] at (2.96,-0.95) {\Kfull};
  \node[chip=c162, minimum width=1.13cm, minimum height=0.42cm] at (4.20,-0.95) {\Kres};
  \node[chip=c43,  minimum width=1.13cm, minimum height=0.42cm] at (5.44,-0.95) {\Kmid};
  \node[chip=c17,  minimum width=1.13cm, minimum height=0.42cm] at (6.68,-0.95) {\Kcoarse};
  \node[chip=crand, minimum width=1.70cm, minimum height=0.70cm] at (2.96,-1.47) {random\\partitions};
  \node[chip=cflav, minimum width=3.04cm, minimum height=0.70cm] at (4.77,-1.47) {controlled flavor\\boundaries};
  \node[chip=c162, densely dotted, minimum width=1.13cm, minimum height=0.42cm] at (2.96,-2.32) {\Kres};
  \node[chip=c17,  densely dotted, minimum width=1.13cm, minimum height=0.42cm] at (4.20,-2.32) {\Kcoarse};
  \node[txt, anchor=west] at (5.44,-2.53) {with auxiliary\\mass regression};
  \node[chip=c188, densely dotted, minimum width=1.13cm, minimum height=0.42cm] at (2.96,-2.94) {\Kfull};
  \node[txt, anchor=west] at (4.20,-3.15) {kinematic inputs only};
  % evaluation
  \node[stage, text width=2.4cm] (fz) at (9.68,-0.95) {frozen features:\\linear probe,\\MLP probe};
  \node[stage, text width=2.4cm] (ft) at (9.68,-2.05) {fine-tuning\\of all weights};
  \node[stage, text width=2.4cm, draw=black, fill=white, line width=1.1pt] (sc) at (9.68,-3.15) {trained from\\scratch};
  \node[stage, text width=3.4cm, minimum height=3.2cm] (task) at (13.20,-2.00)
    {\textbf{Downstream tasks}\\JetClass-II: all classes\\JetClass-II: class pairs\\\quad(both class-balanced)\\JetClass\\
     top tagging\\quark/gluon, Pythia\\quark/gluon, Herwig\\jet mass\FoneXbx};
  \draw[flow] (seed.east) -- (seed.east -| pre.west);
  \draw[flow] (fz.west -| pre.east) -- (fz.west);
  \draw[flow] (ft.west -| pre.east) -- (ft.west);
  \draw[flow] (fz.east) -- (fz.east -| task.west);
  \draw[flow] (ft.east) -- (ft.east -| task.west);
  \draw[flow] (sc.east) -- (sc.east -| task.west);
\end{scope}
\end{tikzpicture}
\caption{Design of the study.
(a)~The class hierarchy on the decays of a heavy resonance~$X$ and on the jets from quantum
chromodynamics (QCD); $q$ is a $u$ or $d$ quark, $\ell$ an electron or a muon, and $\tau_h$,
$\tau_e$, $\tau_\mu$ a $\tau$ lepton decaying to hadrons, an electron or a muon. Each box is one
class of the label set above its column, except the QCD box of the \Kfull-class column (Sophon's
label set~\cite{sophon}), which stands for \LitJetClassTwoNQcdClasses{} classes. Each level merges
classes of the level to its left; the lines below name what each merge removes.
(b)~The controls: the first random partition, where boxes with the same letter are one class, and
the random partitions with controlled flavor boundaries on the decays \RandFlavSplitOrbit{} (one
square per decay; filled squares are moved to a second class).
(c)~Models of one seed share the training jets, their order and the initial weights, except the
output layer, the mass regression target, and the first layer and initial weights of the model with
kinematic inputs only (the particle four-momenta). The hierarchy levels have \MlstNRunsWord{} seeds
each, the models with the auxiliary mass regression \MlstMassNRunsWord{}, the random partitions with
controlled flavor boundaries \RandFlavRuns{} and each random partition \RandVtwoRuns{}. Frozen
features, the output of the backbone (the network below the classification head) with its weights held
fixed, feed linear probes (logistic regressions) and multilayer perceptron (MLP) probes; fine-tuning
trains all weights. Colors identify the label sets and controls in every figure; the plots give the three partitions with
controlled flavor boundaries three shades; dotted borders mark the two
variants, which the other figures draw as dotted lines in the color of their label set.}
\label{fig:design}
\end{figure}
 from main.tex; this file holds one figure environment and nothing else.
%
% WHAT main.tex MUST LOAD (this file loads nothing itself):
%   \usepackage{tikz}
%   \usetikzlibrary{arrows.meta}          % the only TikZ library used (arrow tips)
%   \definecolor{c188}{HTML}{1F77B4}      % the configuration colours of experiments/FIGS/
%   \definecolor{c162}{HTML}{FF7F0E}      % mlst_style.py (matplotlib's tab10 in order, as in
%   \definecolor{c43}{HTML}{2CA02C}       % Sophon's figures); the variants are drawn in their
%   \definecolor{c17}{HTML}{D62728}       % label set's colour with a dotted border, as the
%   \definecolor{crand}{HTML}{17BECF}     % plotted figures draw variants dotted
%   \definecolor{cflav}{HTML}{9467BD}
%   % GENERATED by experiments/FIGS/make_tables.py -- do not edit.
% Every line carries: source file :: path inside that file :: first 16 hex of its
% sha256. paper/journal/provenance.json carries the full hash beside the value.
% Regenerate with:  python3 experiments/FIGS/make_tables.py
% Check in CI with: python3 experiments/FIGS/make_tables.py --check
%
% Macro names spell digits as words (162 -> Onesixtwo) because \newcommand takes
% letters only. A number whose source file does not exist yet is NOT emitted --
% see the missing list below, and leave a \pending{} in the text for it.
%
% MISSING INPUTS -- no macro was emitted for these:
%   community benchmarks, bench v3 readout -- experiments/FIGS/data/bench_v3_metrics*/bench_metrics_last.json
%   v2 pretraining rerun: data order and dropout fixed by seed and epoch, identical initial trunk weights, each epoch reading one of five disjoint random parts of the rows of nearly every file, fixed validation sample, checkpoint at the first maximum within the last ten epochs -- experiments/FIGS/data/v2/pretraining/
%   v2 frozen probes on the rerun models, at the first maximum within the last ten epochs and on the model whose weights are the average of the last ten epochs, BatchNorm statistics recomputed, with untrained-network and pooled-embedding reference rows -- experiments/FIGS/data/v2/probe_ladder/
%   v2 paired ratios with run and test-sample errors on the rerun models -- experiments/FIGS/data/v2/paired_errors/
%   v2 the 64- and 30-class vocabularies -- experiments/FIGS/data/v2/levels_64_30/
%   v2 the leave-one-family-out vocabularies -- experiments/FIGS/data/v2/leave_one_family_out/
%   v2 label-tree recovery on the rerun models -- experiments/FIGS/data/v2/label_recovery/
%   v2 the rerun's random partitions and flavour pair -- experiments/FIGS/data/v2/random_partitions/
%   v2 JetClass-II fine-tuning on files held out from pretraining -- experiments/FIGS/data/v2/finetune/
%   v2 fine-tuning from random initialisation and from the self-supervised model -- experiments/FIGS/data/v2/finetune_references/
%   v2 top tagging and quark/gluon tagging (Pythia and Herwig) fine-tuned from every rerun model -- experiments/FIGS/data/v2/benchmarks/
%   v2 the second and third self-supervised pretraining runs -- experiments/FIGS/data/v2/self_supervised/
%   v2 the mass-output models with the loss weight matched in loss share -- experiments/FIGS/data/v2/mass_lambda_matched/
%   v2 anomaly detection on the rerun models, per-run values -- experiments/FIGS/data/v2/anomaly/
%   v2 the open-data fits on the rerun models -- experiments/FIGS/data/v2/real_data/
%   MassLossShareVtwoOnesixtwoMass -- experiments/FIGS/data/v2/mass_lambda_matched/loss_share.json :: shares.162+mass
%   MassLossShareVtwoOnesevenMass -- experiments/FIGS/data/v2/mass_lambda_matched/loss_share.json :: shares.17+mass
%   MassLossShareVtwoOnesevenMassMatched -- experiments/FIGS/data/v2/mass_lambda_matched/loss_share.json :: shares.17+mass_matched
%   DesignVtwoSecondGpu, DesignVtwoSecondGpuRuns -- the grid's job specs put every run index on one product
%   per-arm parameters and MACs -- flops_by_arm.json (see docs/RECORD.md 2.1)
%   fine-tuning read-outs -- experiments/FIGS/data/v2/ft_mlst/: jc2, jc2_qcd
%   top-tagging recipe check -- ft_mlst/topcheck_bench_metrics.json
%   X->bb and X->bs fits -- experiments/FIGS/data/v2/xbx/
%   X->bs from scratch -- experiments/FIGS/data/v2/xbx_scratch/
%   JetClass accuracy bootstrap -- experiments/FIGS/data/v2/budget/jc1_acc_boot_N10000.json
%   MlstFtJcOneErrCoarseOverFullBootEFour -- waits on the JetClass bootstrap at 1e4 (after the equal-stopping read-out)
%   MlstLpBbQgZeroAccDropFull -- waits on pair probes with the quark/gluon inputs zeroed (open item O10, not run)
%   MlstLpBbQgZeroAccDropCoarse -- waits on pair probes with the quark/gluon inputs zeroed (open item O10, not run)
%   XbsRthirtyFullOverScratch -- waits on X->bs fits and from-scratch runs
%   XbsRthirtyFullOverScratchPr -- waits on X->bs fits and from-scratch runs
%   XbsRthirtyCoarseOverFull -- waits on X->bs fits
%   XbsRthirtyCoarseOverFullPr -- waits on X->bs fits
%   XbsRthirtyNodeOverTransferFull -- waits on X->bs fits
%   XbsRthirtyNodeOverTransferFullPr -- waits on X->bs fits
%   XbbRthirtyDirectOverTransferFull -- waits on X->bb fits
%   XbbRthirtyDirectOverTransferFullPr -- waits on X->bb fits
%   XbbRthirtyTransferCoarseOverFull -- waits on X->bb fits
%   XbbRthirtyTransferCoarseOverFullPr -- waits on X->bb fits
%
% SKIPPED TABLES:
%   T4 compute/parameters: no machine-readable source in the repository; docs/RECORD.md 2.1 names /data/results/mtx/flops/flops_by_arm.json, which is not committed

\newcommand{\PendingPretrain}{\pending{pending v2: pretraining rerun: data order and dropout fixed by seed and epoch, identical initial trunk weights, each epoch reading one of five disjoint random parts of the rows of nearly every file, fixed validation sample, checkpoint at the first maximum within the last ten epochs}}  % experiments/FIGS/make_tables.py :: V2_PENDING['Pretrain'] :: sha256 10e6ec08356a4a2b :: a marker, not a number: red in the draft until the input exists
\newcommand{\PendingProbes}{\pending{pending v2: frozen probes on the rerun models, at the first maximum within the last ten epochs and on the model whose weights are the average of the last ten epochs, BatchNorm statistics recomputed, with untrained-network and pooled-embedding reference rows}}  % experiments/FIGS/make_tables.py :: V2_PENDING['Probes'] :: sha256 10e6ec08356a4a2b :: a marker, not a number: red in the draft until the input exists
\newcommand{\PendingPaired}{\pending{pending v2: paired ratios with run and test-sample errors on the rerun models}}  % experiments/FIGS/make_tables.py :: V2_PENDING['Paired'] :: sha256 10e6ec08356a4a2b :: a marker, not a number: red in the draft until the input exists
\newcommand{\PendingLevels}{\pending{pending v2: the 64- and 30-class vocabularies}}  % experiments/FIGS/make_tables.py :: V2_PENDING['Levels'] :: sha256 10e6ec08356a4a2b :: a marker, not a number: red in the draft until the input exists
\newcommand{\PendingLofo}{\pending{pending v2: the leave-one-family-out vocabularies}}  % experiments/FIGS/make_tables.py :: V2_PENDING['Lofo'] :: sha256 10e6ec08356a4a2b :: a marker, not a number: red in the draft until the input exists
\newcommand{\PendingRecovery}{\pending{pending v2: label-tree recovery on the rerun models}}  % experiments/FIGS/make_tables.py :: V2_PENDING['Recovery'] :: sha256 10e6ec08356a4a2b :: a marker, not a number: red in the draft until the input exists
\newcommand{\PendingRandom}{\pending{pending v2: the rerun's random partitions and flavour pair}}  % experiments/FIGS/make_tables.py :: V2_PENDING['Random'] :: sha256 10e6ec08356a4a2b :: a marker, not a number: red in the draft until the input exists
\newcommand{\PendingFtHeldout}{\pending{pending v2: JetClass-II fine-tuning on files held out from pretraining}}  % experiments/FIGS/make_tables.py :: V2_PENDING['FtHeldout'] :: sha256 10e6ec08356a4a2b :: a marker, not a number: red in the draft until the input exists
\newcommand{\PendingFtRefs}{\pending{pending v2: fine-tuning from random initialisation and from the self-supervised model}}  % experiments/FIGS/make_tables.py :: V2_PENDING['FtRefs'] :: sha256 10e6ec08356a4a2b :: a marker, not a number: red in the draft until the input exists
\newcommand{\PendingBench}{\pending{pending v2: top tagging and quark/gluon tagging (Pythia and Herwig) fine-tuned from every rerun model}}  % experiments/FIGS/make_tables.py :: V2_PENDING['Bench'] :: sha256 10e6ec08356a4a2b :: a marker, not a number: red in the draft until the input exists
\newcommand{\PendingSsl}{\pending{pending v2: the second and third self-supervised pretraining runs}}  % experiments/FIGS/make_tables.py :: V2_PENDING['Ssl'] :: sha256 10e6ec08356a4a2b :: a marker, not a number: red in the draft until the input exists
\newcommand{\PendingMassLambda}{\pending{pending v2: the mass-output models with the loss weight matched in loss share}}  % experiments/FIGS/make_tables.py :: V2_PENDING['MassLambda'] :: sha256 10e6ec08356a4a2b :: a marker, not a number: red in the draft until the input exists
\newcommand{\PendingAnomaly}{\pending{pending v2: anomaly detection on the rerun models, per-run values}}  % experiments/FIGS/make_tables.py :: V2_PENDING['Anomaly'] :: sha256 10e6ec08356a4a2b :: a marker, not a number: red in the draft until the input exists
\newcommand{\PendingRealData}{\pending{pending v2: the open-data fits on the rerun models}}  % experiments/FIGS/make_tables.py :: V2_PENDING['RealData'] :: sha256 10e6ec08356a4a2b :: a marker, not a number: red in the draft until the input exists
\newcommand{\PendingConclusion}{\pending{pending v2: written once every input above is in}}  % experiments/FIGS/make_tables.py :: V2_PENDING['Conclusion'] :: sha256 10e6ec08356a4a2b :: a marker, not a number: red in the draft until the input exists
\newcommand{\AojChainDataset}{178}  % experiments/FIGS/data/aoj_checks_v2/selection_chain_v2.json :: n_dataset :: sha256 729d83085aba25f8
\newcommand{\AojChainStaged}{12{,}589{,}467}  % experiments/FIGS/data/aoj_checks_v2/selection_chain_v2.json :: n_staged :: sha256 729d83085aba25f8
\newcommand{\AojChainRhoWindow}{9{,}441{,}209}  % experiments/FIGS/data/aoj_checks_v2/selection_chain_v2.json :: n_rho_window :: sha256 729d83085aba25f8
\newcommand{\AojNJetsFit}{4{,}194{,}846}  % experiments/FIGS/data/aoj_checks_v2/selection_chain_v2.json :: n_fit :: sha256 729d83085aba25f8
\newcommand{\AojChainPtBins}{8}  % experiments/FIGS/data/aoj_checks_v2/selection_chain_v2.json :: n_pt_bins :: sha256 729d83085aba25f8
\newcommand{\AojChainEtaMax}{2.40}  % experiments/FIGS/data/aoj_checks_v2/selection_chain_v2.json :: abs_eta_max :: sha256 729d83085aba25f8
\newcommand{\MassLossShareVtwoOnesixtwoMass}{\pending{pending: realised mass-loss share in the rerun, 162 classes}}  % experiments/FIGS/make_tables.py :: SLOTS[MassLossShareVtwoOnesixtwoMass] :: sha256 10e6ec08356a4a2b :: a marker, not a number: experiments/FIGS/data/v2/mass_lambda_matched/loss_share.json :: shares.162+mass does not exist yet
\newcommand{\MassLossShareVtwoOnesevenMass}{\pending{pending: realised mass-loss share in the rerun, 17 classes}}  % experiments/FIGS/make_tables.py :: SLOTS[MassLossShareVtwoOnesevenMass] :: sha256 10e6ec08356a4a2b :: a marker, not a number: experiments/FIGS/data/v2/mass_lambda_matched/loss_share.json :: shares.17+mass does not exist yet
\newcommand{\MassLossShareVtwoOnesevenMassMatched}{\pending{pending: realised mass-loss share in the rerun, 17 classes, matched weight}}  % experiments/FIGS/make_tables.py :: SLOTS[MassLossShareVtwoOnesevenMassMatched] :: sha256 10e6ec08356a4a2b :: a marker, not a number: experiments/FIGS/data/v2/mass_lambda_matched/loss_share.json :: shares.17+mass_matched does not exist yet
\newcommand{\FtCoverageEThreeNClasses}{160}  % experiments/FT/data/jc2_v2_manifest.json :: class_coverage['train_N1000_s1.parquet'].n_classes_present :: sha256 4bd68655919b14e9
\newcommand{\FtCoverageEFourNClasses}{183}  % experiments/FT/data/jc2_v2_manifest.json :: class_coverage['train_N10000_s1.parquet'].n_classes_present :: sha256 4bd68655919b14e9
\newcommand{\FtCoverageEThreeMedian}{4}  % experiments/FT/data/jc2_v2_manifest.json :: class_coverage['train_N1000_s1.parquet'].median_per_present_class :: sha256 4bd68655919b14e9
\newcommand{\FtCoverageEThreeOneJet}{23}  % experiments/FT/data/jc2_v2_manifest.json :: class_coverage['train_N1000_s1.parquet'].n_classes_with_one_jet :: sha256 4bd68655919b14e9
\newcommand{\FtCoverageVoneEThreeNClasses}{161}  % experiments/FT/data/jc2_v1_coverage.json :: class_coverage['train_N1000_s1.parquet'].n_classes_present :: sha256 82652f57f708dc23
\newcommand{\FtCoverageVoneEThreeMedian}{5}  % experiments/FT/data/jc2_v1_coverage.json :: class_coverage['train_N1000_s1.parquet'].median_per_present_class :: sha256 82652f57f708dc23
\newcommand{\FtCoverageVoneEThreeOneJet}{25}  % experiments/FT/data/jc2_v1_coverage.json :: class_coverage['train_N1000_s1.parquet'].n_classes_with_one_jet :: sha256 82652f57f708dc23
\newcommand{\ProbeNJets}{2{,}000{,}000}  % experiments/FIGS/data/probe_ladder_v2_mlp2/analysis/seed_level_results.json :: provenance.n_jets_total :: sha256 53e1ec81343ac2c1
\newcommand{\ProbeNSeeds}{5}  % experiments/FIGS/data/probe_ladder_v2_mlp2/analysis/seed_level_results.json :: seeds_used (length) :: sha256 53e1ec81343ac2c1
\newcommand{\ProbeNSeedsWord}{five}  % experiments/FIGS/data/probe_ladder_v2_mlp2/analysis/seed_level_results.json :: seeds_used (length) :: sha256 53e1ec81343ac2c1 :: pretraining runs per vocabulary, as a word for the prose
\newcommand{\ProbeRowAlign}{\texttt{a9449fd420c473f6}}  % experiments/FIGS/data/probe_ladder_v2_mlp2/analysis/seed_level_results.json :: provenance.row_alignment_sha256 (first 16) :: sha256 53e1ec81343ac2c1
\newcommand{\ProbeEpsS}{90\%}  % experiments/FIGS/data/probe_ladder_v2_mlp2/analysis/seed_level_results.json :: levels.bvc_4prong.linear[0].headline_eps_s :: sha256 53e1ec81343ac2c1
\newcommand{\ProbeAucBvcFourprongLinearOneeighteight}{0.99320}  % experiments/FIGS/data/probe_ladder_v2_mlp2/analysis/seed_level_results.json :: levels.bvc_4prong.linear[0].mean_auc :: sha256 53e1ec81343ac2c1 :: 5 seeds, 0 saturated at AUC=1
\newcommand{\ProbeAucSdBvcFourprongLinearOneeighteight}{0.00027}  % experiments/FIGS/data/probe_ladder_v2_mlp2/analysis/seed_level_results.json :: table[bvc_4prong,linear,188].auc (SD over seeds, ddof=1) :: sha256 53e1ec81343ac2c1
\newcommand{\ProbeLogOneMinusAucBvcFourprongLinearOneeighteight}{\ensuremath{-}4.9915}  % experiments/FIGS/data/probe_ladder_v2_mlp2/analysis/seed_level_results.json :: levels.bvc_4prong.linear[0].mean :: sha256 53e1ec81343ac2c1 :: natural log of 1-AUC, lower is better
\newcommand{\ProbeLogOneMinusAucSdBvcFourprongLinearOneeighteight}{0.0389}  % experiments/FIGS/data/probe_ladder_v2_mlp2/analysis/seed_level_results.json :: levels.bvc_4prong.linear[0].seed_sd :: sha256 53e1ec81343ac2c1
\newcommand{\ProbeOmaBvcFourprongLinearOneeighteight}{\ensuremath{(6.80 \pm 0.27)\times10^{-3}}}  % experiments/FIGS/data/probe_ladder_v2_mlp2/analysis/seed_level_results.json :: table[bvc_4prong,linear,188].auc :: sha256 53e1ec81343ac2c1 :: 1-AUC, mean +- SD over 5 seeds
\newcommand{\ProbeBkgLeftBvcFourprongLinearOneeighteight}{14.0}  % experiments/FIGS/data/probe_ladder_v2_mlp2/analysis/seed_level_results.json :: levels.bvc_4prong.linear[0].rejection_points['0.90'].mean_n_bkg_pass :: sha256 53e1ec81343ac2c1 :: background jets passing the cut, mean over seeds
\newcommand{\ProbeRejBvcFourprongLinearOneeighteight}{227\,\ensuremath{\pm}\,23}  % experiments/FIGS/data/probe_ladder_v2_mlp2/analysis/seed_level_results.json :: table[bvc_4prong,linear,188].rejection_points['0.90'].rejection :: sha256 53e1ec81343ac2c1 :: at 90% signal efficiency, mean +- SD over seeds; 0 of 5 seeds at the cap
\newcommand{\ProbeAucBvcFourprongLinearOnesixtwo}{0.99311}  % experiments/FIGS/data/probe_ladder_v2_mlp2/analysis/seed_level_results.json :: levels.bvc_4prong.linear[1].mean_auc :: sha256 53e1ec81343ac2c1 :: 5 seeds, 0 saturated at AUC=1
\newcommand{\ProbeAucSdBvcFourprongLinearOnesixtwo}{0.00018}  % experiments/FIGS/data/probe_ladder_v2_mlp2/analysis/seed_level_results.json :: table[bvc_4prong,linear,162].auc (SD over seeds, ddof=1) :: sha256 53e1ec81343ac2c1
\newcommand{\ProbeLogOneMinusAucBvcFourprongLinearOnesixtwo}{\ensuremath{-}4.9777}  % experiments/FIGS/data/probe_ladder_v2_mlp2/analysis/seed_level_results.json :: levels.bvc_4prong.linear[1].mean :: sha256 53e1ec81343ac2c1 :: natural log of 1-AUC, lower is better
\newcommand{\ProbeLogOneMinusAucSdBvcFourprongLinearOnesixtwo}{0.0264}  % experiments/FIGS/data/probe_ladder_v2_mlp2/analysis/seed_level_results.json :: levels.bvc_4prong.linear[1].seed_sd :: sha256 53e1ec81343ac2c1
\newcommand{\ProbeOmaBvcFourprongLinearOnesixtwo}{\ensuremath{(6.89 \pm 0.18)\times10^{-3}}}  % experiments/FIGS/data/probe_ladder_v2_mlp2/analysis/seed_level_results.json :: table[bvc_4prong,linear,162].auc :: sha256 53e1ec81343ac2c1 :: 1-AUC, mean +- SD over 5 seeds
\newcommand{\ProbeBkgLeftBvcFourprongLinearOnesixtwo}{14.4}  % experiments/FIGS/data/probe_ladder_v2_mlp2/analysis/seed_level_results.json :: levels.bvc_4prong.linear[1].rejection_points['0.90'].mean_n_bkg_pass :: sha256 53e1ec81343ac2c1 :: background jets passing the cut, mean over seeds
\newcommand{\ProbeRejBvcFourprongLinearOnesixtwo}{220\,\ensuremath{\pm}\,40}  % experiments/FIGS/data/probe_ladder_v2_mlp2/analysis/seed_level_results.json :: table[bvc_4prong,linear,162].rejection_points['0.90'].rejection :: sha256 53e1ec81343ac2c1 :: at 90% signal efficiency, mean +- SD over seeds; 0 of 5 seeds at the cap
\newcommand{\ProbeAucBvcFourprongLinearFourthree}{0.99250}  % experiments/FIGS/data/probe_ladder_v2_mlp2/analysis/seed_level_results.json :: levels.bvc_4prong.linear[2].mean_auc :: sha256 53e1ec81343ac2c1 :: 5 seeds, 0 saturated at AUC=1
\newcommand{\ProbeAucSdBvcFourprongLinearFourthree}{0.00049}  % experiments/FIGS/data/probe_ladder_v2_mlp2/analysis/seed_level_results.json :: table[bvc_4prong,linear,43].auc (SD over seeds, ddof=1) :: sha256 53e1ec81343ac2c1
\newcommand{\ProbeLogOneMinusAucBvcFourprongLinearFourthree}{\ensuremath{-}4.8946}  % experiments/FIGS/data/probe_ladder_v2_mlp2/analysis/seed_level_results.json :: levels.bvc_4prong.linear[2].mean :: sha256 53e1ec81343ac2c1 :: natural log of 1-AUC, lower is better
\newcommand{\ProbeLogOneMinusAucSdBvcFourprongLinearFourthree}{0.0646}  % experiments/FIGS/data/probe_ladder_v2_mlp2/analysis/seed_level_results.json :: levels.bvc_4prong.linear[2].seed_sd :: sha256 53e1ec81343ac2c1
\newcommand{\ProbeOmaBvcFourprongLinearFourthree}{\ensuremath{(7.5 \pm 0.5)\times10^{-3}}}  % experiments/FIGS/data/probe_ladder_v2_mlp2/analysis/seed_level_results.json :: table[bvc_4prong,linear,43].auc :: sha256 53e1ec81343ac2c1 :: 1-AUC, mean +- SD over 5 seeds
\newcommand{\ProbeBkgLeftBvcFourprongLinearFourthree}{16.6}  % experiments/FIGS/data/probe_ladder_v2_mlp2/analysis/seed_level_results.json :: levels.bvc_4prong.linear[2].rejection_points['0.90'].mean_n_bkg_pass :: sha256 53e1ec81343ac2c1 :: background jets passing the cut, mean over seeds
\newcommand{\ProbeRejBvcFourprongLinearFourthree}{192\,\ensuremath{\pm}\,20}  % experiments/FIGS/data/probe_ladder_v2_mlp2/analysis/seed_level_results.json :: table[bvc_4prong,linear,43].rejection_points['0.90'].rejection :: sha256 53e1ec81343ac2c1 :: at 90% signal efficiency, mean +- SD over seeds; 0 of 5 seeds at the cap
\newcommand{\ProbeAucBvcFourprongLinearOneseven}{0.96936}  % experiments/FIGS/data/probe_ladder_v2_mlp2/analysis/seed_level_results.json :: levels.bvc_4prong.linear[3].mean_auc :: sha256 53e1ec81343ac2c1 :: 5 seeds, 0 saturated at AUC=1
\newcommand{\ProbeAucSdBvcFourprongLinearOneseven}{0.00514}  % experiments/FIGS/data/probe_ladder_v2_mlp2/analysis/seed_level_results.json :: table[bvc_4prong,linear,17].auc (SD over seeds, ddof=1) :: sha256 53e1ec81343ac2c1
\newcommand{\ProbeLogOneMinusAucBvcFourprongLinearOneseven}{\ensuremath{-}3.4967}  % experiments/FIGS/data/probe_ladder_v2_mlp2/analysis/seed_level_results.json :: levels.bvc_4prong.linear[3].mean :: sha256 53e1ec81343ac2c1 :: natural log of 1-AUC, lower is better
\newcommand{\ProbeLogOneMinusAucSdBvcFourprongLinearOneseven}{0.1667}  % experiments/FIGS/data/probe_ladder_v2_mlp2/analysis/seed_level_results.json :: levels.bvc_4prong.linear[3].seed_sd :: sha256 53e1ec81343ac2c1
\newcommand{\ProbeOmaBvcFourprongLinearOneseven}{\ensuremath{(3.1 \pm 0.5)\times10^{-2}}}  % experiments/FIGS/data/probe_ladder_v2_mlp2/analysis/seed_level_results.json :: table[bvc_4prong,linear,17].auc :: sha256 53e1ec81343ac2c1 :: 1-AUC, mean +- SD over 5 seeds
\newcommand{\ProbeBkgLeftBvcFourprongLinearOneseven}{137.4}  % experiments/FIGS/data/probe_ladder_v2_mlp2/analysis/seed_level_results.json :: levels.bvc_4prong.linear[3].rejection_points['0.90'].mean_n_bkg_pass :: sha256 53e1ec81343ac2c1 :: background jets passing the cut, mean over seeds
\newcommand{\ProbeRejBvcFourprongLinearOneseven}{24\,\ensuremath{\pm}\,6}  % experiments/FIGS/data/probe_ladder_v2_mlp2/analysis/seed_level_results.json :: table[bvc_4prong,linear,17].rejection_points['0.90'].rejection :: sha256 53e1ec81343ac2c1 :: at 90% signal efficiency, mean +- SD over seeds; 0 of 5 seeds at the cap
\newcommand{\ProbeOmaRatioBvcFourprongLinearOnesixtwoOverOneeighteight}{1.01}  % experiments/FIGS/data/probe_ladder_v2_mlp2/analysis/seed_level_results.json :: table[bvc_4prong,linear,162/188].auc :: sha256 53e1ec81343ac2c1 :: ratio of the seed means of 1-AUC, coarser over finer
\newcommand{\ProbeOmaRatioBvcFourprongLinearFourthreeOverOneeighteight}{1.10}  % experiments/FIGS/data/probe_ladder_v2_mlp2/analysis/seed_level_results.json :: table[bvc_4prong,linear,43/188].auc :: sha256 53e1ec81343ac2c1 :: ratio of the seed means of 1-AUC, coarser over finer
\newcommand{\ProbeOmaRatioBvcFourprongLinearOnesevenOverOneeighteight}{4.5}  % experiments/FIGS/data/probe_ladder_v2_mlp2/analysis/seed_level_results.json :: table[bvc_4prong,linear,17/188].auc :: sha256 53e1ec81343ac2c1 :: ratio of the seed means of 1-AUC, coarser over finer
\newcommand{\ProbeOmaRatioBvcFourprongLinearFourthreeOverOnesixtwo}{1.09}  % experiments/FIGS/data/probe_ladder_v2_mlp2/analysis/seed_level_results.json :: table[bvc_4prong,linear,43/162].auc :: sha256 53e1ec81343ac2c1 :: ratio of the seed means of 1-AUC, coarser over finer
\newcommand{\ProbeOmaRatioBvcFourprongLinearOnesevenOverOnesixtwo}{4.4}  % experiments/FIGS/data/probe_ladder_v2_mlp2/analysis/seed_level_results.json :: table[bvc_4prong,linear,17/162].auc :: sha256 53e1ec81343ac2c1 :: ratio of the seed means of 1-AUC, coarser over finer
\newcommand{\ProbeOmaRatioBvcFourprongLinearOnesevenOverFourthree}{4.1}  % experiments/FIGS/data/probe_ladder_v2_mlp2/analysis/seed_level_results.json :: table[bvc_4prong,linear,17/43].auc :: sha256 53e1ec81343ac2c1 :: ratio of the seed means of 1-AUC, coarser over finer
\newcommand{\ProbeAucBvcFourprongMlpOneeighteight}{0.99316}  % experiments/FIGS/data/probe_ladder_v2_mlp2/analysis/seed_level_results.json :: levels.bvc_4prong.mlp[0].mean_auc :: sha256 53e1ec81343ac2c1 :: 5 seeds, 0 saturated at AUC=1
\newcommand{\ProbeAucSdBvcFourprongMlpOneeighteight}{0.00020}  % experiments/FIGS/data/probe_ladder_v2_mlp2/analysis/seed_level_results.json :: table[bvc_4prong,mlp,188].auc (SD over seeds, ddof=1) :: sha256 53e1ec81343ac2c1
\newcommand{\ProbeLogOneMinusAucBvcFourprongMlpOneeighteight}{\ensuremath{-}4.9857}  % experiments/FIGS/data/probe_ladder_v2_mlp2/analysis/seed_level_results.json :: levels.bvc_4prong.mlp[0].mean :: sha256 53e1ec81343ac2c1 :: natural log of 1-AUC, lower is better
\newcommand{\ProbeLogOneMinusAucSdBvcFourprongMlpOneeighteight}{0.0290}  % experiments/FIGS/data/probe_ladder_v2_mlp2/analysis/seed_level_results.json :: levels.bvc_4prong.mlp[0].seed_sd :: sha256 53e1ec81343ac2c1
\newcommand{\ProbeOmaBvcFourprongMlpOneeighteight}{\ensuremath{(6.84 \pm 0.20)\times10^{-3}}}  % experiments/FIGS/data/probe_ladder_v2_mlp2/analysis/seed_level_results.json :: table[bvc_4prong,mlp,188].auc :: sha256 53e1ec81343ac2c1 :: 1-AUC, mean +- SD over 5 seeds
\newcommand{\ProbeBkgLeftBvcFourprongMlpOneeighteight}{15.2}  % experiments/FIGS/data/probe_ladder_v2_mlp2/analysis/seed_level_results.json :: levels.bvc_4prong.mlp[0].rejection_points['0.90'].mean_n_bkg_pass :: sha256 53e1ec81343ac2c1 :: background jets passing the cut, mean over seeds
\newcommand{\ProbeRejBvcFourprongMlpOneeighteight}{212\,\ensuremath{\pm}\,34}  % experiments/FIGS/data/probe_ladder_v2_mlp2/analysis/seed_level_results.json :: table[bvc_4prong,mlp,188].rejection_points['0.90'].rejection :: sha256 53e1ec81343ac2c1 :: at 90% signal efficiency, mean +- SD over seeds; 0 of 5 seeds at the cap
\newcommand{\ProbeAucBvcFourprongMlpOnesixtwo}{0.99318}  % experiments/FIGS/data/probe_ladder_v2_mlp2/analysis/seed_level_results.json :: levels.bvc_4prong.mlp[1].mean_auc :: sha256 53e1ec81343ac2c1 :: 5 seeds, 0 saturated at AUC=1
\newcommand{\ProbeAucSdBvcFourprongMlpOnesixtwo}{0.00020}  % experiments/FIGS/data/probe_ladder_v2_mlp2/analysis/seed_level_results.json :: table[bvc_4prong,mlp,162].auc (SD over seeds, ddof=1) :: sha256 53e1ec81343ac2c1
\newcommand{\ProbeLogOneMinusAucBvcFourprongMlpOnesixtwo}{\ensuremath{-}4.9886}  % experiments/FIGS/data/probe_ladder_v2_mlp2/analysis/seed_level_results.json :: levels.bvc_4prong.mlp[1].mean :: sha256 53e1ec81343ac2c1 :: natural log of 1-AUC, lower is better
\newcommand{\ProbeLogOneMinusAucSdBvcFourprongMlpOnesixtwo}{0.0298}  % experiments/FIGS/data/probe_ladder_v2_mlp2/analysis/seed_level_results.json :: levels.bvc_4prong.mlp[1].seed_sd :: sha256 53e1ec81343ac2c1
\newcommand{\ProbeOmaBvcFourprongMlpOnesixtwo}{\ensuremath{(6.82 \pm 0.20)\times10^{-3}}}  % experiments/FIGS/data/probe_ladder_v2_mlp2/analysis/seed_level_results.json :: table[bvc_4prong,mlp,162].auc :: sha256 53e1ec81343ac2c1 :: 1-AUC, mean +- SD over 5 seeds
\newcommand{\ProbeBkgLeftBvcFourprongMlpOnesixtwo}{16.6}  % experiments/FIGS/data/probe_ladder_v2_mlp2/analysis/seed_level_results.json :: levels.bvc_4prong.mlp[1].rejection_points['0.90'].mean_n_bkg_pass :: sha256 53e1ec81343ac2c1 :: background jets passing the cut, mean over seeds
\newcommand{\ProbeRejBvcFourprongMlpOnesixtwo}{200\,\ensuremath{\pm}\,40}  % experiments/FIGS/data/probe_ladder_v2_mlp2/analysis/seed_level_results.json :: table[bvc_4prong,mlp,162].rejection_points['0.90'].rejection :: sha256 53e1ec81343ac2c1 :: at 90% signal efficiency, mean +- SD over seeds; 0 of 5 seeds at the cap
\newcommand{\ProbeAucBvcFourprongMlpFourthree}{0.99232}  % experiments/FIGS/data/probe_ladder_v2_mlp2/analysis/seed_level_results.json :: levels.bvc_4prong.mlp[2].mean_auc :: sha256 53e1ec81343ac2c1 :: 5 seeds, 0 saturated at AUC=1
\newcommand{\ProbeAucSdBvcFourprongMlpFourthree}{0.00058}  % experiments/FIGS/data/probe_ladder_v2_mlp2/analysis/seed_level_results.json :: table[bvc_4prong,mlp,43].auc (SD over seeds, ddof=1) :: sha256 53e1ec81343ac2c1
\newcommand{\ProbeLogOneMinusAucBvcFourprongMlpFourthree}{\ensuremath{-}4.8719}  % experiments/FIGS/data/probe_ladder_v2_mlp2/analysis/seed_level_results.json :: levels.bvc_4prong.mlp[2].mean :: sha256 53e1ec81343ac2c1 :: natural log of 1-AUC, lower is better
\newcommand{\ProbeLogOneMinusAucSdBvcFourprongMlpFourthree}{0.0734}  % experiments/FIGS/data/probe_ladder_v2_mlp2/analysis/seed_level_results.json :: levels.bvc_4prong.mlp[2].seed_sd :: sha256 53e1ec81343ac2c1
\newcommand{\ProbeOmaBvcFourprongMlpFourthree}{\ensuremath{(7.7 \pm 0.6)\times10^{-3}}}  % experiments/FIGS/data/probe_ladder_v2_mlp2/analysis/seed_level_results.json :: table[bvc_4prong,mlp,43].auc :: sha256 53e1ec81343ac2c1 :: 1-AUC, mean +- SD over 5 seeds
\newcommand{\ProbeBkgLeftBvcFourprongMlpFourthree}{17.4}  % experiments/FIGS/data/probe_ladder_v2_mlp2/analysis/seed_level_results.json :: levels.bvc_4prong.mlp[2].rejection_points['0.90'].mean_n_bkg_pass :: sha256 53e1ec81343ac2c1 :: background jets passing the cut, mean over seeds
\newcommand{\ProbeRejBvcFourprongMlpFourthree}{184\,\ensuremath{\pm}\,23}  % experiments/FIGS/data/probe_ladder_v2_mlp2/analysis/seed_level_results.json :: table[bvc_4prong,mlp,43].rejection_points['0.90'].rejection :: sha256 53e1ec81343ac2c1 :: at 90% signal efficiency, mean +- SD over seeds; 0 of 5 seeds at the cap
\newcommand{\ProbeAucBvcFourprongMlpOneseven}{0.97766}  % experiments/FIGS/data/probe_ladder_v2_mlp2/analysis/seed_level_results.json :: levels.bvc_4prong.mlp[3].mean_auc :: sha256 53e1ec81343ac2c1 :: 5 seeds, 0 saturated at AUC=1
\newcommand{\ProbeAucSdBvcFourprongMlpOneseven}{0.00310}  % experiments/FIGS/data/probe_ladder_v2_mlp2/analysis/seed_level_results.json :: table[bvc_4prong,mlp,17].auc (SD over seeds, ddof=1) :: sha256 53e1ec81343ac2c1
\newcommand{\ProbeLogOneMinusAucBvcFourprongMlpOneseven}{\ensuremath{-}3.8087}  % experiments/FIGS/data/probe_ladder_v2_mlp2/analysis/seed_level_results.json :: levels.bvc_4prong.mlp[3].mean :: sha256 53e1ec81343ac2c1 :: natural log of 1-AUC, lower is better
\newcommand{\ProbeLogOneMinusAucSdBvcFourprongMlpOneseven}{0.1366}  % experiments/FIGS/data/probe_ladder_v2_mlp2/analysis/seed_level_results.json :: levels.bvc_4prong.mlp[3].seed_sd :: sha256 53e1ec81343ac2c1
\newcommand{\ProbeOmaBvcFourprongMlpOneseven}{\ensuremath{(2.23 \pm 0.31)\times10^{-2}}}  % experiments/FIGS/data/probe_ladder_v2_mlp2/analysis/seed_level_results.json :: table[bvc_4prong,mlp,17].auc :: sha256 53e1ec81343ac2c1 :: 1-AUC, mean +- SD over 5 seeds
\newcommand{\ProbeBkgLeftBvcFourprongMlpOneseven}{99.0}  % experiments/FIGS/data/probe_ladder_v2_mlp2/analysis/seed_level_results.json :: levels.bvc_4prong.mlp[3].rejection_points['0.90'].mean_n_bkg_pass :: sha256 53e1ec81343ac2c1 :: background jets passing the cut, mean over seeds
\newcommand{\ProbeRejBvcFourprongMlpOneseven}{33\,\ensuremath{\pm}\,5}  % experiments/FIGS/data/probe_ladder_v2_mlp2/analysis/seed_level_results.json :: table[bvc_4prong,mlp,17].rejection_points['0.90'].rejection :: sha256 53e1ec81343ac2c1 :: at 90% signal efficiency, mean +- SD over seeds; 0 of 5 seeds at the cap
\newcommand{\ProbeOmaRatioBvcFourprongMlpOnesixtwoOverOneeighteight}{1.00}  % experiments/FIGS/data/probe_ladder_v2_mlp2/analysis/seed_level_results.json :: table[bvc_4prong,mlp,162/188].auc :: sha256 53e1ec81343ac2c1 :: ratio of the seed means of 1-AUC, coarser over finer
\newcommand{\ProbeOmaRatioBvcFourprongMlpFourthreeOverOneeighteight}{1.12}  % experiments/FIGS/data/probe_ladder_v2_mlp2/analysis/seed_level_results.json :: table[bvc_4prong,mlp,43/188].auc :: sha256 53e1ec81343ac2c1 :: ratio of the seed means of 1-AUC, coarser over finer
\newcommand{\ProbeOmaRatioBvcFourprongMlpOnesevenOverOneeighteight}{3.3}  % experiments/FIGS/data/probe_ladder_v2_mlp2/analysis/seed_level_results.json :: table[bvc_4prong,mlp,17/188].auc :: sha256 53e1ec81343ac2c1 :: ratio of the seed means of 1-AUC, coarser over finer
\newcommand{\ProbeOmaRatioBvcFourprongMlpFourthreeOverOnesixtwo}{1.13}  % experiments/FIGS/data/probe_ladder_v2_mlp2/analysis/seed_level_results.json :: table[bvc_4prong,mlp,43/162].auc :: sha256 53e1ec81343ac2c1 :: ratio of the seed means of 1-AUC, coarser over finer
\newcommand{\ProbeOmaRatioBvcFourprongMlpOnesevenOverOnesixtwo}{3.3}  % experiments/FIGS/data/probe_ladder_v2_mlp2/analysis/seed_level_results.json :: table[bvc_4prong,mlp,17/162].auc :: sha256 53e1ec81343ac2c1 :: ratio of the seed means of 1-AUC, coarser over finer
\newcommand{\ProbeOmaRatioBvcFourprongMlpOnesevenOverFourthree}{2.9}  % experiments/FIGS/data/probe_ladder_v2_mlp2/analysis/seed_level_results.json :: table[bvc_4prong,mlp,17/43].auc :: sha256 53e1ec81343ac2c1 :: ratio of the seed means of 1-AUC, coarser over finer
\newcommand{\ProbeAucBvcQcdLinearOneeighteight}{0.98519}  % experiments/FIGS/data/probe_ladder_v2_mlp2/analysis/seed_level_results.json :: levels.bvc_qcd.linear[0].mean_auc :: sha256 53e1ec81343ac2c1 :: 5 seeds, 0 saturated at AUC=1
\newcommand{\ProbeAucSdBvcQcdLinearOneeighteight}{0.00061}  % experiments/FIGS/data/probe_ladder_v2_mlp2/analysis/seed_level_results.json :: table[bvc_qcd,linear,188].auc (SD over seeds, ddof=1) :: sha256 53e1ec81343ac2c1
\newcommand{\ProbeLogOneMinusAucBvcQcdLinearOneeighteight}{\ensuremath{-}4.2133}  % experiments/FIGS/data/probe_ladder_v2_mlp2/analysis/seed_level_results.json :: levels.bvc_qcd.linear[0].mean :: sha256 53e1ec81343ac2c1 :: natural log of 1-AUC, lower is better
\newcommand{\ProbeLogOneMinusAucSdBvcQcdLinearOneeighteight}{0.0422}  % experiments/FIGS/data/probe_ladder_v2_mlp2/analysis/seed_level_results.json :: levels.bvc_qcd.linear[0].seed_sd :: sha256 53e1ec81343ac2c1
\newcommand{\ProbeOmaBvcQcdLinearOneeighteight}{\ensuremath{(1.48 \pm 0.06)\times10^{-2}}}  % experiments/FIGS/data/probe_ladder_v2_mlp2/analysis/seed_level_results.json :: table[bvc_qcd,linear,188].auc :: sha256 53e1ec81343ac2c1 :: 1-AUC, mean +- SD over 5 seeds
\newcommand{\ProbeBkgLeftBvcQcdLinearOneeighteight}{96.4}  % experiments/FIGS/data/probe_ladder_v2_mlp2/analysis/seed_level_results.json :: levels.bvc_qcd.linear[0].rejection_points['0.90'].mean_n_bkg_pass :: sha256 53e1ec81343ac2c1 :: background jets passing the cut, mean over seeds
\newcommand{\ProbeRejBvcQcdLinearOneeighteight}{56\,\ensuremath{\pm}\,5}  % experiments/FIGS/data/probe_ladder_v2_mlp2/analysis/seed_level_results.json :: table[bvc_qcd,linear,188].rejection_points['0.90'].rejection :: sha256 53e1ec81343ac2c1 :: at 90% signal efficiency, mean +- SD over seeds; 0 of 5 seeds at the cap
\newcommand{\ProbeAucBvcQcdLinearOnesixtwo}{0.98305}  % experiments/FIGS/data/probe_ladder_v2_mlp2/analysis/seed_level_results.json :: levels.bvc_qcd.linear[1].mean_auc :: sha256 53e1ec81343ac2c1 :: 5 seeds, 0 saturated at AUC=1
\newcommand{\ProbeAucSdBvcQcdLinearOnesixtwo}{0.00024}  % experiments/FIGS/data/probe_ladder_v2_mlp2/analysis/seed_level_results.json :: table[bvc_qcd,linear,162].auc (SD over seeds, ddof=1) :: sha256 53e1ec81343ac2c1
\newcommand{\ProbeLogOneMinusAucBvcQcdLinearOnesixtwo}{\ensuremath{-}4.0774}  % experiments/FIGS/data/probe_ladder_v2_mlp2/analysis/seed_level_results.json :: levels.bvc_qcd.linear[1].mean :: sha256 53e1ec81343ac2c1 :: natural log of 1-AUC, lower is better
\newcommand{\ProbeLogOneMinusAucSdBvcQcdLinearOnesixtwo}{0.0142}  % experiments/FIGS/data/probe_ladder_v2_mlp2/analysis/seed_level_results.json :: levels.bvc_qcd.linear[1].seed_sd :: sha256 53e1ec81343ac2c1
\newcommand{\ProbeOmaBvcQcdLinearOnesixtwo}{\ensuremath{(1.695 \pm 0.024)\times10^{-2}}}  % experiments/FIGS/data/probe_ladder_v2_mlp2/analysis/seed_level_results.json :: table[bvc_qcd,linear,162].auc :: sha256 53e1ec81343ac2c1 :: 1-AUC, mean +- SD over 5 seeds
\newcommand{\ProbeBkgLeftBvcQcdLinearOnesixtwo}{121.4}  % experiments/FIGS/data/probe_ladder_v2_mlp2/analysis/seed_level_results.json :: levels.bvc_qcd.linear[1].rejection_points['0.90'].mean_n_bkg_pass :: sha256 53e1ec81343ac2c1 :: background jets passing the cut, mean over seeds
\newcommand{\ProbeRejBvcQcdLinearOnesixtwo}{44.3\,\ensuremath{\pm}\,2.1}  % experiments/FIGS/data/probe_ladder_v2_mlp2/analysis/seed_level_results.json :: table[bvc_qcd,linear,162].rejection_points['0.90'].rejection :: sha256 53e1ec81343ac2c1 :: at 90% signal efficiency, mean +- SD over seeds; 0 of 5 seeds at the cap
\newcommand{\ProbeAucBvcQcdLinearFourthree}{0.98281}  % experiments/FIGS/data/probe_ladder_v2_mlp2/analysis/seed_level_results.json :: levels.bvc_qcd.linear[2].mean_auc :: sha256 53e1ec81343ac2c1 :: 5 seeds, 0 saturated at AUC=1
\newcommand{\ProbeAucSdBvcQcdLinearFourthree}{0.00064}  % experiments/FIGS/data/probe_ladder_v2_mlp2/analysis/seed_level_results.json :: table[bvc_qcd,linear,43].auc (SD over seeds, ddof=1) :: sha256 53e1ec81343ac2c1
\newcommand{\ProbeLogOneMinusAucBvcQcdLinearFourthree}{\ensuremath{-}4.0638}  % experiments/FIGS/data/probe_ladder_v2_mlp2/analysis/seed_level_results.json :: levels.bvc_qcd.linear[2].mean :: sha256 53e1ec81343ac2c1 :: natural log of 1-AUC, lower is better
\newcommand{\ProbeLogOneMinusAucSdBvcQcdLinearFourthree}{0.0370}  % experiments/FIGS/data/probe_ladder_v2_mlp2/analysis/seed_level_results.json :: levels.bvc_qcd.linear[2].seed_sd :: sha256 53e1ec81343ac2c1
\newcommand{\ProbeOmaBvcQcdLinearFourthree}{\ensuremath{(1.72 \pm 0.06)\times10^{-2}}}  % experiments/FIGS/data/probe_ladder_v2_mlp2/analysis/seed_level_results.json :: table[bvc_qcd,linear,43].auc :: sha256 53e1ec81343ac2c1 :: 1-AUC, mean +- SD over 5 seeds
\newcommand{\ProbeBkgLeftBvcQcdLinearFourthree}{126.8}  % experiments/FIGS/data/probe_ladder_v2_mlp2/analysis/seed_level_results.json :: levels.bvc_qcd.linear[2].rejection_points['0.90'].mean_n_bkg_pass :: sha256 53e1ec81343ac2c1 :: background jets passing the cut, mean over seeds
\newcommand{\ProbeRejBvcQcdLinearFourthree}{43\,\ensuremath{\pm}\,4}  % experiments/FIGS/data/probe_ladder_v2_mlp2/analysis/seed_level_results.json :: table[bvc_qcd,linear,43].rejection_points['0.90'].rejection :: sha256 53e1ec81343ac2c1 :: at 90% signal efficiency, mean +- SD over seeds; 0 of 5 seeds at the cap
\newcommand{\ProbeAucBvcQcdLinearOneseven}{0.95504}  % experiments/FIGS/data/probe_ladder_v2_mlp2/analysis/seed_level_results.json :: levels.bvc_qcd.linear[3].mean_auc :: sha256 53e1ec81343ac2c1 :: 5 seeds, 0 saturated at AUC=1
\newcommand{\ProbeAucSdBvcQcdLinearOneseven}{0.00527}  % experiments/FIGS/data/probe_ladder_v2_mlp2/analysis/seed_level_results.json :: table[bvc_qcd,linear,17].auc (SD over seeds, ddof=1) :: sha256 53e1ec81343ac2c1
\newcommand{\ProbeLogOneMinusAucBvcQcdLinearOneseven}{\ensuremath{-}3.1076}  % experiments/FIGS/data/probe_ladder_v2_mlp2/analysis/seed_level_results.json :: levels.bvc_qcd.linear[3].mean :: sha256 53e1ec81343ac2c1 :: natural log of 1-AUC, lower is better
\newcommand{\ProbeLogOneMinusAucSdBvcQcdLinearOneseven}{0.1198}  % experiments/FIGS/data/probe_ladder_v2_mlp2/analysis/seed_level_results.json :: levels.bvc_qcd.linear[3].seed_sd :: sha256 53e1ec81343ac2c1
\newcommand{\ProbeOmaBvcQcdLinearOneseven}{\ensuremath{(4.5 \pm 0.5)\times10^{-2}}}  % experiments/FIGS/data/probe_ladder_v2_mlp2/analysis/seed_level_results.json :: table[bvc_qcd,linear,17].auc :: sha256 53e1ec81343ac2c1 :: 1-AUC, mean +- SD over 5 seeds
\newcommand{\ProbeBkgLeftBvcQcdLinearOneseven}{544.6}  % experiments/FIGS/data/probe_ladder_v2_mlp2/analysis/seed_level_results.json :: levels.bvc_qcd.linear[3].rejection_points['0.90'].mean_n_bkg_pass :: sha256 53e1ec81343ac2c1 :: background jets passing the cut, mean over seeds
\newcommand{\ProbeRejBvcQcdLinearOneseven}{10.1\,\ensuremath{\pm}\,1.9}  % experiments/FIGS/data/probe_ladder_v2_mlp2/analysis/seed_level_results.json :: table[bvc_qcd,linear,17].rejection_points['0.90'].rejection :: sha256 53e1ec81343ac2c1 :: at 90% signal efficiency, mean +- SD over seeds; 0 of 5 seeds at the cap
\newcommand{\ProbeOmaRatioBvcQcdLinearOnesixtwoOverOneeighteight}{1.14}  % experiments/FIGS/data/probe_ladder_v2_mlp2/analysis/seed_level_results.json :: table[bvc_qcd,linear,162/188].auc :: sha256 53e1ec81343ac2c1 :: ratio of the seed means of 1-AUC, coarser over finer
\newcommand{\ProbeOmaRatioBvcQcdLinearFourthreeOverOneeighteight}{1.16}  % experiments/FIGS/data/probe_ladder_v2_mlp2/analysis/seed_level_results.json :: table[bvc_qcd,linear,43/188].auc :: sha256 53e1ec81343ac2c1 :: ratio of the seed means of 1-AUC, coarser over finer
\newcommand{\ProbeOmaRatioBvcQcdLinearOnesevenOverOneeighteight}{3.0}  % experiments/FIGS/data/probe_ladder_v2_mlp2/analysis/seed_level_results.json :: table[bvc_qcd,linear,17/188].auc :: sha256 53e1ec81343ac2c1 :: ratio of the seed means of 1-AUC, coarser over finer
\newcommand{\ProbeOmaRatioBvcQcdLinearFourthreeOverOnesixtwo}{1.01}  % experiments/FIGS/data/probe_ladder_v2_mlp2/analysis/seed_level_results.json :: table[bvc_qcd,linear,43/162].auc :: sha256 53e1ec81343ac2c1 :: ratio of the seed means of 1-AUC, coarser over finer
\newcommand{\ProbeOmaRatioBvcQcdLinearOnesevenOverOnesixtwo}{2.7}  % experiments/FIGS/data/probe_ladder_v2_mlp2/analysis/seed_level_results.json :: table[bvc_qcd,linear,17/162].auc :: sha256 53e1ec81343ac2c1 :: ratio of the seed means of 1-AUC, coarser over finer
\newcommand{\ProbeOmaRatioBvcQcdLinearOnesevenOverFourthree}{2.6}  % experiments/FIGS/data/probe_ladder_v2_mlp2/analysis/seed_level_results.json :: table[bvc_qcd,linear,17/43].auc :: sha256 53e1ec81343ac2c1 :: ratio of the seed means of 1-AUC, coarser over finer
\newcommand{\ProbeAucBvcQcdMlpOneeighteight}{0.98527}  % experiments/FIGS/data/probe_ladder_v2_mlp2/analysis/seed_level_results.json :: levels.bvc_qcd.mlp[0].mean_auc :: sha256 53e1ec81343ac2c1 :: 5 seeds, 0 saturated at AUC=1
\newcommand{\ProbeAucSdBvcQcdMlpOneeighteight}{0.00058}  % experiments/FIGS/data/probe_ladder_v2_mlp2/analysis/seed_level_results.json :: table[bvc_qcd,mlp,188].auc (SD over seeds, ddof=1) :: sha256 53e1ec81343ac2c1
\newcommand{\ProbeLogOneMinusAucBvcQcdMlpOneeighteight}{\ensuremath{-}4.2188}  % experiments/FIGS/data/probe_ladder_v2_mlp2/analysis/seed_level_results.json :: levels.bvc_qcd.mlp[0].mean :: sha256 53e1ec81343ac2c1 :: natural log of 1-AUC, lower is better
\newcommand{\ProbeLogOneMinusAucSdBvcQcdMlpOneeighteight}{0.0400}  % experiments/FIGS/data/probe_ladder_v2_mlp2/analysis/seed_level_results.json :: levels.bvc_qcd.mlp[0].seed_sd :: sha256 53e1ec81343ac2c1
\newcommand{\ProbeOmaBvcQcdMlpOneeighteight}{\ensuremath{(1.47 \pm 0.06)\times10^{-2}}}  % experiments/FIGS/data/probe_ladder_v2_mlp2/analysis/seed_level_results.json :: table[bvc_qcd,mlp,188].auc :: sha256 53e1ec81343ac2c1 :: 1-AUC, mean +- SD over 5 seeds
\newcommand{\ProbeBkgLeftBvcQcdMlpOneeighteight}{94.6}  % experiments/FIGS/data/probe_ladder_v2_mlp2/analysis/seed_level_results.json :: levels.bvc_qcd.mlp[0].rejection_points['0.90'].mean_n_bkg_pass :: sha256 53e1ec81343ac2c1 :: background jets passing the cut, mean over seeds
\newcommand{\ProbeRejBvcQcdMlpOneeighteight}{57\,\ensuremath{\pm}\,6}  % experiments/FIGS/data/probe_ladder_v2_mlp2/analysis/seed_level_results.json :: table[bvc_qcd,mlp,188].rejection_points['0.90'].rejection :: sha256 53e1ec81343ac2c1 :: at 90% signal efficiency, mean +- SD over seeds; 0 of 5 seeds at the cap
\newcommand{\ProbeAucBvcQcdMlpOnesixtwo}{0.98361}  % experiments/FIGS/data/probe_ladder_v2_mlp2/analysis/seed_level_results.json :: levels.bvc_qcd.mlp[1].mean_auc :: sha256 53e1ec81343ac2c1 :: 5 seeds, 0 saturated at AUC=1
\newcommand{\ProbeAucSdBvcQcdMlpOnesixtwo}{0.00039}  % experiments/FIGS/data/probe_ladder_v2_mlp2/analysis/seed_level_results.json :: table[bvc_qcd,mlp,162].auc (SD over seeds, ddof=1) :: sha256 53e1ec81343ac2c1
\newcommand{\ProbeLogOneMinusAucBvcQcdMlpOnesixtwo}{\ensuremath{-}4.1113}  % experiments/FIGS/data/probe_ladder_v2_mlp2/analysis/seed_level_results.json :: levels.bvc_qcd.mlp[1].mean :: sha256 53e1ec81343ac2c1 :: natural log of 1-AUC, lower is better
\newcommand{\ProbeLogOneMinusAucSdBvcQcdMlpOnesixtwo}{0.0241}  % experiments/FIGS/data/probe_ladder_v2_mlp2/analysis/seed_level_results.json :: levels.bvc_qcd.mlp[1].seed_sd :: sha256 53e1ec81343ac2c1
\newcommand{\ProbeOmaBvcQcdMlpOnesixtwo}{\ensuremath{(1.64 \pm 0.04)\times10^{-2}}}  % experiments/FIGS/data/probe_ladder_v2_mlp2/analysis/seed_level_results.json :: table[bvc_qcd,mlp,162].auc :: sha256 53e1ec81343ac2c1 :: 1-AUC, mean +- SD over 5 seeds
\newcommand{\ProbeBkgLeftBvcQcdMlpOnesixtwo}{119.6}  % experiments/FIGS/data/probe_ladder_v2_mlp2/analysis/seed_level_results.json :: levels.bvc_qcd.mlp[1].rejection_points['0.90'].mean_n_bkg_pass :: sha256 53e1ec81343ac2c1 :: background jets passing the cut, mean over seeds
\newcommand{\ProbeRejBvcQcdMlpOnesixtwo}{45.1\,\ensuremath{\pm}\,3.2}  % experiments/FIGS/data/probe_ladder_v2_mlp2/analysis/seed_level_results.json :: table[bvc_qcd,mlp,162].rejection_points['0.90'].rejection :: sha256 53e1ec81343ac2c1 :: at 90% signal efficiency, mean +- SD over seeds; 0 of 5 seeds at the cap
\newcommand{\ProbeAucBvcQcdMlpFourthree}{0.98341}  % experiments/FIGS/data/probe_ladder_v2_mlp2/analysis/seed_level_results.json :: levels.bvc_qcd.mlp[2].mean_auc :: sha256 53e1ec81343ac2c1 :: 5 seeds, 0 saturated at AUC=1
\newcommand{\ProbeAucSdBvcQcdMlpFourthree}{0.00052}  % experiments/FIGS/data/probe_ladder_v2_mlp2/analysis/seed_level_results.json :: table[bvc_qcd,mlp,43].auc (SD over seeds, ddof=1) :: sha256 53e1ec81343ac2c1
\newcommand{\ProbeLogOneMinusAucBvcQcdMlpFourthree}{\ensuremath{-}4.0995}  % experiments/FIGS/data/probe_ladder_v2_mlp2/analysis/seed_level_results.json :: levels.bvc_qcd.mlp[2].mean :: sha256 53e1ec81343ac2c1 :: natural log of 1-AUC, lower is better
\newcommand{\ProbeLogOneMinusAucSdBvcQcdMlpFourthree}{0.0318}  % experiments/FIGS/data/probe_ladder_v2_mlp2/analysis/seed_level_results.json :: levels.bvc_qcd.mlp[2].seed_sd :: sha256 53e1ec81343ac2c1
\newcommand{\ProbeOmaBvcQcdMlpFourthree}{\ensuremath{(1.66 \pm 0.05)\times10^{-2}}}  % experiments/FIGS/data/probe_ladder_v2_mlp2/analysis/seed_level_results.json :: table[bvc_qcd,mlp,43].auc :: sha256 53e1ec81343ac2c1 :: 1-AUC, mean +- SD over 5 seeds
\newcommand{\ProbeBkgLeftBvcQcdMlpFourthree}{119.2}  % experiments/FIGS/data/probe_ladder_v2_mlp2/analysis/seed_level_results.json :: levels.bvc_qcd.mlp[2].rejection_points['0.90'].mean_n_bkg_pass :: sha256 53e1ec81343ac2c1 :: background jets passing the cut, mean over seeds
\newcommand{\ProbeRejBvcQcdMlpFourthree}{45.1\,\ensuremath{\pm}\,1.3}  % experiments/FIGS/data/probe_ladder_v2_mlp2/analysis/seed_level_results.json :: table[bvc_qcd,mlp,43].rejection_points['0.90'].rejection :: sha256 53e1ec81343ac2c1 :: at 90% signal efficiency, mean +- SD over seeds; 0 of 5 seeds at the cap
\newcommand{\ProbeAucBvcQcdMlpOneseven}{0.96226}  % experiments/FIGS/data/probe_ladder_v2_mlp2/analysis/seed_level_results.json :: levels.bvc_qcd.mlp[3].mean_auc :: sha256 53e1ec81343ac2c1 :: 5 seeds, 0 saturated at AUC=1
\newcommand{\ProbeAucSdBvcQcdMlpOneseven}{0.00373}  % experiments/FIGS/data/probe_ladder_v2_mlp2/analysis/seed_level_results.json :: table[bvc_qcd,mlp,17].auc (SD over seeds, ddof=1) :: sha256 53e1ec81343ac2c1
\newcommand{\ProbeLogOneMinusAucBvcQcdMlpOneseven}{\ensuremath{-}3.2811}  % experiments/FIGS/data/probe_ladder_v2_mlp2/analysis/seed_level_results.json :: levels.bvc_qcd.mlp[3].mean :: sha256 53e1ec81343ac2c1 :: natural log of 1-AUC, lower is better
\newcommand{\ProbeLogOneMinusAucSdBvcQcdMlpOneseven}{0.1010}  % experiments/FIGS/data/probe_ladder_v2_mlp2/analysis/seed_level_results.json :: levels.bvc_qcd.mlp[3].seed_sd :: sha256 53e1ec81343ac2c1
\newcommand{\ProbeOmaBvcQcdMlpOneseven}{\ensuremath{(3.8 \pm 0.4)\times10^{-2}}}  % experiments/FIGS/data/probe_ladder_v2_mlp2/analysis/seed_level_results.json :: table[bvc_qcd,mlp,17].auc :: sha256 53e1ec81343ac2c1 :: 1-AUC, mean +- SD over 5 seeds
\newcommand{\ProbeBkgLeftBvcQcdMlpOneseven}{458.8}  % experiments/FIGS/data/probe_ladder_v2_mlp2/analysis/seed_level_results.json :: levels.bvc_qcd.mlp[3].rejection_points['0.90'].mean_n_bkg_pass :: sha256 53e1ec81343ac2c1 :: background jets passing the cut, mean over seeds
\newcommand{\ProbeRejBvcQcdMlpOneseven}{11.8\,\ensuremath{\pm}\,1.5}  % experiments/FIGS/data/probe_ladder_v2_mlp2/analysis/seed_level_results.json :: table[bvc_qcd,mlp,17].rejection_points['0.90'].rejection :: sha256 53e1ec81343ac2c1 :: at 90% signal efficiency, mean +- SD over seeds; 0 of 5 seeds at the cap
\newcommand{\ProbeOmaRatioBvcQcdMlpOnesixtwoOverOneeighteight}{1.11}  % experiments/FIGS/data/probe_ladder_v2_mlp2/analysis/seed_level_results.json :: table[bvc_qcd,mlp,162/188].auc :: sha256 53e1ec81343ac2c1 :: ratio of the seed means of 1-AUC, coarser over finer
\newcommand{\ProbeOmaRatioBvcQcdMlpFourthreeOverOneeighteight}{1.13}  % experiments/FIGS/data/probe_ladder_v2_mlp2/analysis/seed_level_results.json :: table[bvc_qcd,mlp,43/188].auc :: sha256 53e1ec81343ac2c1 :: ratio of the seed means of 1-AUC, coarser over finer
\newcommand{\ProbeOmaRatioBvcQcdMlpOnesevenOverOneeighteight}{2.6}  % experiments/FIGS/data/probe_ladder_v2_mlp2/analysis/seed_level_results.json :: table[bvc_qcd,mlp,17/188].auc :: sha256 53e1ec81343ac2c1 :: ratio of the seed means of 1-AUC, coarser over finer
\newcommand{\ProbeOmaRatioBvcQcdMlpFourthreeOverOnesixtwo}{1.01}  % experiments/FIGS/data/probe_ladder_v2_mlp2/analysis/seed_level_results.json :: table[bvc_qcd,mlp,43/162].auc :: sha256 53e1ec81343ac2c1 :: ratio of the seed means of 1-AUC, coarser over finer
\newcommand{\ProbeOmaRatioBvcQcdMlpOnesevenOverOnesixtwo}{2.3}  % experiments/FIGS/data/probe_ladder_v2_mlp2/analysis/seed_level_results.json :: table[bvc_qcd,mlp,17/162].auc :: sha256 53e1ec81343ac2c1 :: ratio of the seed means of 1-AUC, coarser over finer
\newcommand{\ProbeOmaRatioBvcQcdMlpOnesevenOverFourthree}{2.3}  % experiments/FIGS/data/probe_ladder_v2_mlp2/analysis/seed_level_results.json :: table[bvc_qcd,mlp,17/43].auc :: sha256 53e1ec81343ac2c1 :: ratio of the seed means of 1-AUC, coarser over finer
\newcommand{\ProbeAucBvcResonantLinearOneeighteight}{0.99768}  % experiments/FIGS/data/probe_ladder_v2_mlp2/analysis/seed_level_results.json :: levels.bvc_resonant.linear[0].mean_auc :: sha256 53e1ec81343ac2c1 :: 5 seeds, 0 saturated at AUC=1
\newcommand{\ProbeAucSdBvcResonantLinearOneeighteight}{0.00007}  % experiments/FIGS/data/probe_ladder_v2_mlp2/analysis/seed_level_results.json :: table[bvc_resonant,linear,188].auc (SD over seeds, ddof=1) :: sha256 53e1ec81343ac2c1
\newcommand{\ProbeLogOneMinusAucBvcResonantLinearOneeighteight}{\ensuremath{-}6.0683}  % experiments/FIGS/data/probe_ladder_v2_mlp2/analysis/seed_level_results.json :: levels.bvc_resonant.linear[0].mean :: sha256 53e1ec81343ac2c1 :: natural log of 1-AUC, lower is better
\newcommand{\ProbeLogOneMinusAucSdBvcResonantLinearOneeighteight}{0.0317}  % experiments/FIGS/data/probe_ladder_v2_mlp2/analysis/seed_level_results.json :: levels.bvc_resonant.linear[0].seed_sd :: sha256 53e1ec81343ac2c1
\newcommand{\ProbeOmaBvcResonantLinearOneeighteight}{\ensuremath{(2.32 \pm 0.07)\times10^{-3}}}  % experiments/FIGS/data/probe_ladder_v2_mlp2/analysis/seed_level_results.json :: table[bvc_resonant,linear,188].auc :: sha256 53e1ec81343ac2c1 :: 1-AUC, mean +- SD over 5 seeds
\newcommand{\ProbeBkgLeftBvcResonantLinearOneeighteight}{9.0}  % experiments/FIGS/data/probe_ladder_v2_mlp2/analysis/seed_level_results.json :: levels.bvc_resonant.linear[0].rejection_points['0.90'].mean_n_bkg_pass :: sha256 53e1ec81343ac2c1 :: background jets passing the cut, mean over seeds
\newcommand{\ProbeRejBvcResonantLinearOneeighteight}{1{,}380\,\ensuremath{\pm}\,350}  % experiments/FIGS/data/probe_ladder_v2_mlp2/analysis/seed_level_results.json :: table[bvc_resonant,linear,188].rejection_points['0.90'].rejection :: sha256 53e1ec81343ac2c1 :: at 90% signal efficiency, mean +- SD over seeds; 0 of 5 seeds at the cap
\newcommand{\ProbeAucBvcResonantLinearOnesixtwo}{0.99764}  % experiments/FIGS/data/probe_ladder_v2_mlp2/analysis/seed_level_results.json :: levels.bvc_resonant.linear[1].mean_auc :: sha256 53e1ec81343ac2c1 :: 5 seeds, 0 saturated at AUC=1
\newcommand{\ProbeAucSdBvcResonantLinearOnesixtwo}{0.00006}  % experiments/FIGS/data/probe_ladder_v2_mlp2/analysis/seed_level_results.json :: table[bvc_resonant,linear,162].auc (SD over seeds, ddof=1) :: sha256 53e1ec81343ac2c1
\newcommand{\ProbeLogOneMinusAucBvcResonantLinearOnesixtwo}{\ensuremath{-}6.0493}  % experiments/FIGS/data/probe_ladder_v2_mlp2/analysis/seed_level_results.json :: levels.bvc_resonant.linear[1].mean :: sha256 53e1ec81343ac2c1 :: natural log of 1-AUC, lower is better
\newcommand{\ProbeLogOneMinusAucSdBvcResonantLinearOnesixtwo}{0.0261}  % experiments/FIGS/data/probe_ladder_v2_mlp2/analysis/seed_level_results.json :: levels.bvc_resonant.linear[1].seed_sd :: sha256 53e1ec81343ac2c1
\newcommand{\ProbeOmaBvcResonantLinearOnesixtwo}{\ensuremath{(2.36 \pm 0.06)\times10^{-3}}}  % experiments/FIGS/data/probe_ladder_v2_mlp2/analysis/seed_level_results.json :: table[bvc_resonant,linear,162].auc :: sha256 53e1ec81343ac2c1 :: 1-AUC, mean +- SD over 5 seeds
\newcommand{\ProbeBkgLeftBvcResonantLinearOnesixtwo}{9.4}  % experiments/FIGS/data/probe_ladder_v2_mlp2/analysis/seed_level_results.json :: levels.bvc_resonant.linear[1].rejection_points['0.90'].mean_n_bkg_pass :: sha256 53e1ec81343ac2c1 :: background jets passing the cut, mean over seeds
\newcommand{\ProbeRejBvcResonantLinearOnesixtwo}{1{,}290\,\ensuremath{\pm}\,240}  % experiments/FIGS/data/probe_ladder_v2_mlp2/analysis/seed_level_results.json :: table[bvc_resonant,linear,162].rejection_points['0.90'].rejection :: sha256 53e1ec81343ac2c1 :: at 90% signal efficiency, mean +- SD over seeds; 0 of 5 seeds at the cap
\newcommand{\ProbeAucBvcResonantLinearFourthree}{0.99748}  % experiments/FIGS/data/probe_ladder_v2_mlp2/analysis/seed_level_results.json :: levels.bvc_resonant.linear[2].mean_auc :: sha256 53e1ec81343ac2c1 :: 5 seeds, 0 saturated at AUC=1
\newcommand{\ProbeAucSdBvcResonantLinearFourthree}{0.00007}  % experiments/FIGS/data/probe_ladder_v2_mlp2/analysis/seed_level_results.json :: table[bvc_resonant,linear,43].auc (SD over seeds, ddof=1) :: sha256 53e1ec81343ac2c1
\newcommand{\ProbeLogOneMinusAucBvcResonantLinearFourthree}{\ensuremath{-}5.9834}  % experiments/FIGS/data/probe_ladder_v2_mlp2/analysis/seed_level_results.json :: levels.bvc_resonant.linear[2].mean :: sha256 53e1ec81343ac2c1 :: natural log of 1-AUC, lower is better
\newcommand{\ProbeLogOneMinusAucSdBvcResonantLinearFourthree}{0.0264}  % experiments/FIGS/data/probe_ladder_v2_mlp2/analysis/seed_level_results.json :: levels.bvc_resonant.linear[2].seed_sd :: sha256 53e1ec81343ac2c1
\newcommand{\ProbeOmaBvcResonantLinearFourthree}{\ensuremath{(2.52 \pm 0.07)\times10^{-3}}}  % experiments/FIGS/data/probe_ladder_v2_mlp2/analysis/seed_level_results.json :: table[bvc_resonant,linear,43].auc :: sha256 53e1ec81343ac2c1 :: 1-AUC, mean +- SD over 5 seeds
\newcommand{\ProbeBkgLeftBvcResonantLinearFourthree}{11.0}  % experiments/FIGS/data/probe_ladder_v2_mlp2/analysis/seed_level_results.json :: levels.bvc_resonant.linear[2].rejection_points['0.90'].mean_n_bkg_pass :: sha256 53e1ec81343ac2c1 :: background jets passing the cut, mean over seeds
\newcommand{\ProbeRejBvcResonantLinearFourthree}{1{,}130\,\ensuremath{\pm}\,260}  % experiments/FIGS/data/probe_ladder_v2_mlp2/analysis/seed_level_results.json :: table[bvc_resonant,linear,43].rejection_points['0.90'].rejection :: sha256 53e1ec81343ac2c1 :: at 90% signal efficiency, mean +- SD over seeds; 0 of 5 seeds at the cap
\newcommand{\ProbeAucBvcResonantLinearOneseven}{0.99067}  % experiments/FIGS/data/probe_ladder_v2_mlp2/analysis/seed_level_results.json :: levels.bvc_resonant.linear[3].mean_auc :: sha256 53e1ec81343ac2c1 :: 5 seeds, 0 saturated at AUC=1
\newcommand{\ProbeAucSdBvcResonantLinearOneseven}{0.00147}  % experiments/FIGS/data/probe_ladder_v2_mlp2/analysis/seed_level_results.json :: table[bvc_resonant,linear,17].auc (SD over seeds, ddof=1) :: sha256 53e1ec81343ac2c1
\newcommand{\ProbeLogOneMinusAucBvcResonantLinearOneseven}{\ensuremath{-}4.6848}  % experiments/FIGS/data/probe_ladder_v2_mlp2/analysis/seed_level_results.json :: levels.bvc_resonant.linear[3].mean :: sha256 53e1ec81343ac2c1 :: natural log of 1-AUC, lower is better
\newcommand{\ProbeLogOneMinusAucSdBvcResonantLinearOneseven}{0.1619}  % experiments/FIGS/data/probe_ladder_v2_mlp2/analysis/seed_level_results.json :: levels.bvc_resonant.linear[3].seed_sd :: sha256 53e1ec81343ac2c1
\newcommand{\ProbeOmaBvcResonantLinearOneseven}{\ensuremath{(9.3 \pm 1.5)\times10^{-3}}}  % experiments/FIGS/data/probe_ladder_v2_mlp2/analysis/seed_level_results.json :: table[bvc_resonant,linear,17].auc :: sha256 53e1ec81343ac2c1 :: 1-AUC, mean +- SD over 5 seeds
\newcommand{\ProbeBkgLeftBvcResonantLinearOneseven}{90.8}  % experiments/FIGS/data/probe_ladder_v2_mlp2/analysis/seed_level_results.json :: levels.bvc_resonant.linear[3].rejection_points['0.90'].mean_n_bkg_pass :: sha256 53e1ec81343ac2c1 :: background jets passing the cut, mean over seeds
\newcommand{\ProbeRejBvcResonantLinearOneseven}{134\,\ensuremath{\pm}\,21}  % experiments/FIGS/data/probe_ladder_v2_mlp2/analysis/seed_level_results.json :: table[bvc_resonant,linear,17].rejection_points['0.90'].rejection :: sha256 53e1ec81343ac2c1 :: at 90% signal efficiency, mean +- SD over seeds; 0 of 5 seeds at the cap
\newcommand{\ProbeOmaRatioBvcResonantLinearOnesixtwoOverOneeighteight}{1.02}  % experiments/FIGS/data/probe_ladder_v2_mlp2/analysis/seed_level_results.json :: table[bvc_resonant,linear,162/188].auc :: sha256 53e1ec81343ac2c1 :: ratio of the seed means of 1-AUC, coarser over finer
\newcommand{\ProbeOmaRatioBvcResonantLinearFourthreeOverOneeighteight}{1.09}  % experiments/FIGS/data/probe_ladder_v2_mlp2/analysis/seed_level_results.json :: table[bvc_resonant,linear,43/188].auc :: sha256 53e1ec81343ac2c1 :: ratio of the seed means of 1-AUC, coarser over finer
\newcommand{\ProbeOmaRatioBvcResonantLinearOnesevenOverOneeighteight}{4.0}  % experiments/FIGS/data/probe_ladder_v2_mlp2/analysis/seed_level_results.json :: table[bvc_resonant,linear,17/188].auc :: sha256 53e1ec81343ac2c1 :: ratio of the seed means of 1-AUC, coarser over finer
\newcommand{\ProbeOmaRatioBvcResonantLinearFourthreeOverOnesixtwo}{1.07}  % experiments/FIGS/data/probe_ladder_v2_mlp2/analysis/seed_level_results.json :: table[bvc_resonant,linear,43/162].auc :: sha256 53e1ec81343ac2c1 :: ratio of the seed means of 1-AUC, coarser over finer
\newcommand{\ProbeOmaRatioBvcResonantLinearOnesevenOverOnesixtwo}{4.0}  % experiments/FIGS/data/probe_ladder_v2_mlp2/analysis/seed_level_results.json :: table[bvc_resonant,linear,17/162].auc :: sha256 53e1ec81343ac2c1 :: ratio of the seed means of 1-AUC, coarser over finer
\newcommand{\ProbeOmaRatioBvcResonantLinearOnesevenOverFourthree}{3.7}  % experiments/FIGS/data/probe_ladder_v2_mlp2/analysis/seed_level_results.json :: table[bvc_resonant,linear,17/43].auc :: sha256 53e1ec81343ac2c1 :: ratio of the seed means of 1-AUC, coarser over finer
\newcommand{\ProbeAucBvcResonantMlpOneeighteight}{0.99778}  % experiments/FIGS/data/probe_ladder_v2_mlp2/analysis/seed_level_results.json :: levels.bvc_resonant.mlp[0].mean_auc :: sha256 53e1ec81343ac2c1 :: 5 seeds, 0 saturated at AUC=1
\newcommand{\ProbeAucSdBvcResonantMlpOneeighteight}{0.00009}  % experiments/FIGS/data/probe_ladder_v2_mlp2/analysis/seed_level_results.json :: table[bvc_resonant,mlp,188].auc (SD over seeds, ddof=1) :: sha256 53e1ec81343ac2c1
\newcommand{\ProbeLogOneMinusAucBvcResonantMlpOneeighteight}{\ensuremath{-}6.1102}  % experiments/FIGS/data/probe_ladder_v2_mlp2/analysis/seed_level_results.json :: levels.bvc_resonant.mlp[0].mean :: sha256 53e1ec81343ac2c1 :: natural log of 1-AUC, lower is better
\newcommand{\ProbeLogOneMinusAucSdBvcResonantMlpOneeighteight}{0.0383}  % experiments/FIGS/data/probe_ladder_v2_mlp2/analysis/seed_level_results.json :: levels.bvc_resonant.mlp[0].seed_sd :: sha256 53e1ec81343ac2c1
\newcommand{\ProbeOmaBvcResonantMlpOneeighteight}{\ensuremath{(2.22 \pm 0.09)\times10^{-3}}}  % experiments/FIGS/data/probe_ladder_v2_mlp2/analysis/seed_level_results.json :: table[bvc_resonant,mlp,188].auc :: sha256 53e1ec81343ac2c1 :: 1-AUC, mean +- SD over 5 seeds
\newcommand{\ProbeBkgLeftBvcResonantMlpOneeighteight}{9.8}  % experiments/FIGS/data/probe_ladder_v2_mlp2/analysis/seed_level_results.json :: levels.bvc_resonant.mlp[0].rejection_points['0.90'].mean_n_bkg_pass :: sha256 53e1ec81343ac2c1 :: background jets passing the cut, mean over seeds
\newcommand{\ProbeRejBvcResonantMlpOneeighteight}{1{,}230\,\ensuremath{\pm}\,180}  % experiments/FIGS/data/probe_ladder_v2_mlp2/analysis/seed_level_results.json :: table[bvc_resonant,mlp,188].rejection_points['0.90'].rejection :: sha256 53e1ec81343ac2c1 :: at 90% signal efficiency, mean +- SD over seeds; 0 of 5 seeds at the cap
\newcommand{\ProbeAucBvcResonantMlpOnesixtwo}{0.99768}  % experiments/FIGS/data/probe_ladder_v2_mlp2/analysis/seed_level_results.json :: levels.bvc_resonant.mlp[1].mean_auc :: sha256 53e1ec81343ac2c1 :: 5 seeds, 0 saturated at AUC=1
\newcommand{\ProbeAucSdBvcResonantMlpOnesixtwo}{0.00005}  % experiments/FIGS/data/probe_ladder_v2_mlp2/analysis/seed_level_results.json :: table[bvc_resonant,mlp,162].auc (SD over seeds, ddof=1) :: sha256 53e1ec81343ac2c1
\newcommand{\ProbeLogOneMinusAucBvcResonantMlpOnesixtwo}{\ensuremath{-}6.0667}  % experiments/FIGS/data/probe_ladder_v2_mlp2/analysis/seed_level_results.json :: levels.bvc_resonant.mlp[1].mean :: sha256 53e1ec81343ac2c1 :: natural log of 1-AUC, lower is better
\newcommand{\ProbeLogOneMinusAucSdBvcResonantMlpOnesixtwo}{0.0220}  % experiments/FIGS/data/probe_ladder_v2_mlp2/analysis/seed_level_results.json :: levels.bvc_resonant.mlp[1].seed_sd :: sha256 53e1ec81343ac2c1
\newcommand{\ProbeOmaBvcResonantMlpOnesixtwo}{\ensuremath{(2.32 \pm 0.05)\times10^{-3}}}  % experiments/FIGS/data/probe_ladder_v2_mlp2/analysis/seed_level_results.json :: table[bvc_resonant,mlp,162].auc :: sha256 53e1ec81343ac2c1 :: 1-AUC, mean +- SD over 5 seeds
\newcommand{\ProbeBkgLeftBvcResonantMlpOnesixtwo}{10.4}  % experiments/FIGS/data/probe_ladder_v2_mlp2/analysis/seed_level_results.json :: levels.bvc_resonant.mlp[1].rejection_points['0.90'].mean_n_bkg_pass :: sha256 53e1ec81343ac2c1 :: background jets passing the cut, mean over seeds
\newcommand{\ProbeRejBvcResonantMlpOnesixtwo}{1{,}170\,\ensuremath{\pm}\,220}  % experiments/FIGS/data/probe_ladder_v2_mlp2/analysis/seed_level_results.json :: table[bvc_resonant,mlp,162].rejection_points['0.90'].rejection :: sha256 53e1ec81343ac2c1 :: at 90% signal efficiency, mean +- SD over seeds; 0 of 5 seeds at the cap
\newcommand{\ProbeAucBvcResonantMlpFourthree}{0.99747}  % experiments/FIGS/data/probe_ladder_v2_mlp2/analysis/seed_level_results.json :: levels.bvc_resonant.mlp[2].mean_auc :: sha256 53e1ec81343ac2c1 :: 5 seeds, 0 saturated at AUC=1
\newcommand{\ProbeAucSdBvcResonantMlpFourthree}{0.00004}  % experiments/FIGS/data/probe_ladder_v2_mlp2/analysis/seed_level_results.json :: table[bvc_resonant,mlp,43].auc (SD over seeds, ddof=1) :: sha256 53e1ec81343ac2c1
\newcommand{\ProbeLogOneMinusAucBvcResonantMlpFourthree}{\ensuremath{-}5.9792}  % experiments/FIGS/data/probe_ladder_v2_mlp2/analysis/seed_level_results.json :: levels.bvc_resonant.mlp[2].mean :: sha256 53e1ec81343ac2c1 :: natural log of 1-AUC, lower is better
\newcommand{\ProbeLogOneMinusAucSdBvcResonantMlpFourthree}{0.0164}  % experiments/FIGS/data/probe_ladder_v2_mlp2/analysis/seed_level_results.json :: levels.bvc_resonant.mlp[2].seed_sd :: sha256 53e1ec81343ac2c1
\newcommand{\ProbeOmaBvcResonantMlpFourthree}{\ensuremath{(2.53 \pm 0.04)\times10^{-3}}}  % experiments/FIGS/data/probe_ladder_v2_mlp2/analysis/seed_level_results.json :: table[bvc_resonant,mlp,43].auc :: sha256 53e1ec81343ac2c1 :: 1-AUC, mean +- SD over 5 seeds
\newcommand{\ProbeBkgLeftBvcResonantMlpFourthree}{11.8}  % experiments/FIGS/data/probe_ladder_v2_mlp2/analysis/seed_level_results.json :: levels.bvc_resonant.mlp[2].rejection_points['0.90'].mean_n_bkg_pass :: sha256 53e1ec81343ac2c1 :: background jets passing the cut, mean over seeds
\newcommand{\ProbeRejBvcResonantMlpFourthree}{1{,}030\,\ensuremath{\pm}\,170}  % experiments/FIGS/data/probe_ladder_v2_mlp2/analysis/seed_level_results.json :: table[bvc_resonant,mlp,43].rejection_points['0.90'].rejection :: sha256 53e1ec81343ac2c1 :: at 90% signal efficiency, mean +- SD over seeds; 0 of 5 seeds at the cap
\newcommand{\ProbeAucBvcResonantMlpOneseven}{0.99274}  % experiments/FIGS/data/probe_ladder_v2_mlp2/analysis/seed_level_results.json :: levels.bvc_resonant.mlp[3].mean_auc :: sha256 53e1ec81343ac2c1 :: 5 seeds, 0 saturated at AUC=1
\newcommand{\ProbeAucSdBvcResonantMlpOneseven}{0.00089}  % experiments/FIGS/data/probe_ladder_v2_mlp2/analysis/seed_level_results.json :: table[bvc_resonant,mlp,17].auc (SD over seeds, ddof=1) :: sha256 53e1ec81343ac2c1
\newcommand{\ProbeLogOneMinusAucBvcResonantMlpOneseven}{\ensuremath{-}4.9309}  % experiments/FIGS/data/probe_ladder_v2_mlp2/analysis/seed_level_results.json :: levels.bvc_resonant.mlp[3].mean :: sha256 53e1ec81343ac2c1 :: natural log of 1-AUC, lower is better
\newcommand{\ProbeLogOneMinusAucSdBvcResonantMlpOneseven}{0.1224}  % experiments/FIGS/data/probe_ladder_v2_mlp2/analysis/seed_level_results.json :: levels.bvc_resonant.mlp[3].seed_sd :: sha256 53e1ec81343ac2c1
\newcommand{\ProbeOmaBvcResonantMlpOneseven}{\ensuremath{(7.3 \pm 0.9)\times10^{-3}}}  % experiments/FIGS/data/probe_ladder_v2_mlp2/analysis/seed_level_results.json :: table[bvc_resonant,mlp,17].auc :: sha256 53e1ec81343ac2c1 :: 1-AUC, mean +- SD over 5 seeds
\newcommand{\ProbeBkgLeftBvcResonantMlpOneseven}{62.6}  % experiments/FIGS/data/probe_ladder_v2_mlp2/analysis/seed_level_results.json :: levels.bvc_resonant.mlp[3].rejection_points['0.90'].mean_n_bkg_pass :: sha256 53e1ec81343ac2c1 :: background jets passing the cut, mean over seeds
\newcommand{\ProbeRejBvcResonantMlpOneseven}{192\,\ensuremath{\pm}\,21}  % experiments/FIGS/data/probe_ladder_v2_mlp2/analysis/seed_level_results.json :: table[bvc_resonant,mlp,17].rejection_points['0.90'].rejection :: sha256 53e1ec81343ac2c1 :: at 90% signal efficiency, mean +- SD over seeds; 0 of 5 seeds at the cap
\newcommand{\ProbeOmaRatioBvcResonantMlpOnesixtwoOverOneeighteight}{1.04}  % experiments/FIGS/data/probe_ladder_v2_mlp2/analysis/seed_level_results.json :: table[bvc_resonant,mlp,162/188].auc :: sha256 53e1ec81343ac2c1 :: ratio of the seed means of 1-AUC, coarser over finer
\newcommand{\ProbeOmaRatioBvcResonantMlpFourthreeOverOneeighteight}{1.14}  % experiments/FIGS/data/probe_ladder_v2_mlp2/analysis/seed_level_results.json :: table[bvc_resonant,mlp,43/188].auc :: sha256 53e1ec81343ac2c1 :: ratio of the seed means of 1-AUC, coarser over finer
\newcommand{\ProbeOmaRatioBvcResonantMlpOnesevenOverOneeighteight}{3.3}  % experiments/FIGS/data/probe_ladder_v2_mlp2/analysis/seed_level_results.json :: table[bvc_resonant,mlp,17/188].auc :: sha256 53e1ec81343ac2c1 :: ratio of the seed means of 1-AUC, coarser over finer
\newcommand{\ProbeOmaRatioBvcResonantMlpFourthreeOverOnesixtwo}{1.09}  % experiments/FIGS/data/probe_ladder_v2_mlp2/analysis/seed_level_results.json :: table[bvc_resonant,mlp,43/162].auc :: sha256 53e1ec81343ac2c1 :: ratio of the seed means of 1-AUC, coarser over finer
\newcommand{\ProbeOmaRatioBvcResonantMlpOnesevenOverOnesixtwo}{3.1}  % experiments/FIGS/data/probe_ladder_v2_mlp2/analysis/seed_level_results.json :: table[bvc_resonant,mlp,17/162].auc :: sha256 53e1ec81343ac2c1 :: ratio of the seed means of 1-AUC, coarser over finer
\newcommand{\ProbeOmaRatioBvcResonantMlpOnesevenOverFourthree}{2.9}  % experiments/FIGS/data/probe_ladder_v2_mlp2/analysis/seed_level_results.json :: table[bvc_resonant,mlp,17/43].auc :: sha256 53e1ec81343ac2c1 :: ratio of the seed means of 1-AUC, coarser over finer
\newcommand{\ProbeAucEeVsMmLinearOneeighteight}{\ensuremath{1^{\dagger}}}  % experiments/FIGS/data/probe_ladder_v2_mlp2/analysis/seed_level_results.json :: levels.ee_vs_mm.linear[0].mean_auc :: sha256 53e1ec81343ac2c1 :: 5 seeds, 5 saturated at AUC=1
\newcommand{\ProbeAucSdEeVsMmLinearOneeighteight}{0.00000}  % experiments/FIGS/data/probe_ladder_v2_mlp2/analysis/seed_level_results.json :: table[ee_vs_mm,linear,188].auc (SD over seeds, ddof=1) :: sha256 53e1ec81343ac2c1
\newcommand{\ProbeLogOneMinusAucEeVsMmLinearOneeighteight}{\ensuremath{-}16.9579}  % experiments/FIGS/data/probe_ladder_v2_mlp2/analysis/seed_level_results.json :: levels.ee_vs_mm.linear[0].mean :: sha256 53e1ec81343ac2c1 :: natural log of 1-AUC, lower is better
\newcommand{\ProbeLogOneMinusAucSdEeVsMmLinearOneeighteight}{0.0000}  % experiments/FIGS/data/probe_ladder_v2_mlp2/analysis/seed_level_results.json :: levels.ee_vs_mm.linear[0].seed_sd :: sha256 53e1ec81343ac2c1
\newcommand{\ProbeOmaEeVsMmLinearOneeighteight}{---}  % experiments/FIGS/data/probe_ladder_v2_mlp2/analysis/seed_level_results.json :: table[ee_vs_mm,linear,188].auc :: sha256 53e1ec81343ac2c1 :: not quoted: the AUC reached 1 in a seed, so 1-AUC is only a bound
\newcommand{\ProbeBkgLeftEeVsMmLinearOneeighteight}{0.0}  % experiments/FIGS/data/probe_ladder_v2_mlp2/analysis/seed_level_results.json :: levels.ee_vs_mm.linear[0].rejection_points['0.90'].mean_n_bkg_pass :: sha256 53e1ec81343ac2c1 :: background jets passing the cut, mean over seeds
\newcommand{\ProbeRejEeVsMmLinearOneeighteight}{\ensuremath{>}1{,}571}  % experiments/FIGS/data/probe_ladder_v2_mlp2/analysis/seed_level_results.json :: table[ee_vs_mm,linear,188].rejection_points['0.90'].rejection :: sha256 53e1ec81343ac2c1 :: at 90% signal efficiency, mean +- SD over seeds; 5 of 5 seeds at the cap
\newcommand{\ProbeAucEeVsMmLinearOnesixtwo}{\ensuremath{1^{\dagger}}}  % experiments/FIGS/data/probe_ladder_v2_mlp2/analysis/seed_level_results.json :: levels.ee_vs_mm.linear[1].mean_auc :: sha256 53e1ec81343ac2c1 :: 5 seeds, 5 saturated at AUC=1
\newcommand{\ProbeAucSdEeVsMmLinearOnesixtwo}{0.00000}  % experiments/FIGS/data/probe_ladder_v2_mlp2/analysis/seed_level_results.json :: table[ee_vs_mm,linear,162].auc (SD over seeds, ddof=1) :: sha256 53e1ec81343ac2c1
\newcommand{\ProbeLogOneMinusAucEeVsMmLinearOnesixtwo}{\ensuremath{-}16.9579}  % experiments/FIGS/data/probe_ladder_v2_mlp2/analysis/seed_level_results.json :: levels.ee_vs_mm.linear[1].mean :: sha256 53e1ec81343ac2c1 :: natural log of 1-AUC, lower is better
\newcommand{\ProbeLogOneMinusAucSdEeVsMmLinearOnesixtwo}{0.0000}  % experiments/FIGS/data/probe_ladder_v2_mlp2/analysis/seed_level_results.json :: levels.ee_vs_mm.linear[1].seed_sd :: sha256 53e1ec81343ac2c1
\newcommand{\ProbeOmaEeVsMmLinearOnesixtwo}{---}  % experiments/FIGS/data/probe_ladder_v2_mlp2/analysis/seed_level_results.json :: table[ee_vs_mm,linear,162].auc :: sha256 53e1ec81343ac2c1 :: not quoted: the AUC reached 1 in a seed, so 1-AUC is only a bound
\newcommand{\ProbeBkgLeftEeVsMmLinearOnesixtwo}{0.0}  % experiments/FIGS/data/probe_ladder_v2_mlp2/analysis/seed_level_results.json :: levels.ee_vs_mm.linear[1].rejection_points['0.90'].mean_n_bkg_pass :: sha256 53e1ec81343ac2c1 :: background jets passing the cut, mean over seeds
\newcommand{\ProbeRejEeVsMmLinearOnesixtwo}{\ensuremath{>}1{,}571}  % experiments/FIGS/data/probe_ladder_v2_mlp2/analysis/seed_level_results.json :: table[ee_vs_mm,linear,162].rejection_points['0.90'].rejection :: sha256 53e1ec81343ac2c1 :: at 90% signal efficiency, mean +- SD over seeds; 5 of 5 seeds at the cap
\newcommand{\ProbeAucEeVsMmLinearFourthree}{0.97519}  % experiments/FIGS/data/probe_ladder_v2_mlp2/analysis/seed_level_results.json :: levels.ee_vs_mm.linear[2].mean_auc :: sha256 53e1ec81343ac2c1 :: 5 seeds, 0 saturated at AUC=1
\newcommand{\ProbeAucSdEeVsMmLinearFourthree}{0.02073}  % experiments/FIGS/data/probe_ladder_v2_mlp2/analysis/seed_level_results.json :: table[ee_vs_mm,linear,43].auc (SD over seeds, ddof=1) :: sha256 53e1ec81343ac2c1
\newcommand{\ProbeLogOneMinusAucEeVsMmLinearFourthree}{\ensuremath{-}4.0326}  % experiments/FIGS/data/probe_ladder_v2_mlp2/analysis/seed_level_results.json :: levels.ee_vs_mm.linear[2].mean :: sha256 53e1ec81343ac2c1 :: natural log of 1-AUC, lower is better
\newcommand{\ProbeLogOneMinusAucSdEeVsMmLinearFourthree}{0.9467}  % experiments/FIGS/data/probe_ladder_v2_mlp2/analysis/seed_level_results.json :: levels.ee_vs_mm.linear[2].seed_sd :: sha256 53e1ec81343ac2c1
\newcommand{\ProbeOmaEeVsMmLinearFourthree}{\ensuremath{(2.5 \pm 2.1)\times10^{-2}}}  % experiments/FIGS/data/probe_ladder_v2_mlp2/analysis/seed_level_results.json :: table[ee_vs_mm,linear,43].auc :: sha256 53e1ec81343ac2c1 :: 1-AUC, mean +- SD over 5 seeds
\newcommand{\ProbeBkgLeftEeVsMmLinearFourthree}{155.8}  % experiments/FIGS/data/probe_ladder_v2_mlp2/analysis/seed_level_results.json :: levels.ee_vs_mm.linear[2].rejection_points['0.90'].mean_n_bkg_pass :: sha256 53e1ec81343ac2c1 :: background jets passing the cut, mean over seeds
\newcommand{\ProbeRejEeVsMmLinearFourthree}{80\,\ensuremath{\pm}\,70}  % experiments/FIGS/data/probe_ladder_v2_mlp2/analysis/seed_level_results.json :: table[ee_vs_mm,linear,43].rejection_points['0.90'].rejection :: sha256 53e1ec81343ac2c1 :: at 90% signal efficiency, mean +- SD over seeds; 0 of 5 seeds at the cap
\newcommand{\ProbeAucEeVsMmLinearOneseven}{0.95832}  % experiments/FIGS/data/probe_ladder_v2_mlp2/analysis/seed_level_results.json :: levels.ee_vs_mm.linear[3].mean_auc :: sha256 53e1ec81343ac2c1 :: 5 seeds, 0 saturated at AUC=1
\newcommand{\ProbeAucSdEeVsMmLinearOneseven}{0.00970}  % experiments/FIGS/data/probe_ladder_v2_mlp2/analysis/seed_level_results.json :: table[ee_vs_mm,linear,17].auc (SD over seeds, ddof=1) :: sha256 53e1ec81343ac2c1
\newcommand{\ProbeLogOneMinusAucEeVsMmLinearOneseven}{\ensuremath{-}3.2013}  % experiments/FIGS/data/probe_ladder_v2_mlp2/analysis/seed_level_results.json :: levels.ee_vs_mm.linear[3].mean :: sha256 53e1ec81343ac2c1 :: natural log of 1-AUC, lower is better
\newcommand{\ProbeLogOneMinusAucSdEeVsMmLinearOneseven}{0.2488}  % experiments/FIGS/data/probe_ladder_v2_mlp2/analysis/seed_level_results.json :: levels.ee_vs_mm.linear[3].seed_sd :: sha256 53e1ec81343ac2c1
\newcommand{\ProbeOmaEeVsMmLinearOneseven}{\ensuremath{(4.2 \pm 1.0)\times10^{-2}}}  % experiments/FIGS/data/probe_ladder_v2_mlp2/analysis/seed_level_results.json :: table[ee_vs_mm,linear,17].auc :: sha256 53e1ec81343ac2c1 :: 1-AUC, mean +- SD over 5 seeds
\newcommand{\ProbeBkgLeftEeVsMmLinearOneseven}{289.6}  % experiments/FIGS/data/probe_ladder_v2_mlp2/analysis/seed_level_results.json :: levels.ee_vs_mm.linear[3].rejection_points['0.90'].mean_n_bkg_pass :: sha256 53e1ec81343ac2c1 :: background jets passing the cut, mean over seeds
\newcommand{\ProbeRejEeVsMmLinearOneseven}{17\,\ensuremath{\pm}\,6}  % experiments/FIGS/data/probe_ladder_v2_mlp2/analysis/seed_level_results.json :: table[ee_vs_mm,linear,17].rejection_points['0.90'].rejection :: sha256 53e1ec81343ac2c1 :: at 90% signal efficiency, mean +- SD over seeds; 0 of 5 seeds at the cap
\newcommand{\ProbeOmaRatioEeVsMmLinearOnesevenOverFourthree}{1.68}  % experiments/FIGS/data/probe_ladder_v2_mlp2/analysis/seed_level_results.json :: table[ee_vs_mm,linear,17/43].auc :: sha256 53e1ec81343ac2c1 :: ratio of the seed means of 1-AUC, coarser over finer
\newcommand{\ProbeAucEeVsMmMlpOneeighteight}{\ensuremath{1^{\dagger}}}  % experiments/FIGS/data/probe_ladder_v2_mlp2/analysis/seed_level_results.json :: levels.ee_vs_mm.mlp[0].mean_auc :: sha256 53e1ec81343ac2c1 :: 5 seeds, 5 saturated at AUC=1
\newcommand{\ProbeAucSdEeVsMmMlpOneeighteight}{0.00000}  % experiments/FIGS/data/probe_ladder_v2_mlp2/analysis/seed_level_results.json :: table[ee_vs_mm,mlp,188].auc (SD over seeds, ddof=1) :: sha256 53e1ec81343ac2c1
\newcommand{\ProbeLogOneMinusAucEeVsMmMlpOneeighteight}{\ensuremath{-}16.9579}  % experiments/FIGS/data/probe_ladder_v2_mlp2/analysis/seed_level_results.json :: levels.ee_vs_mm.mlp[0].mean :: sha256 53e1ec81343ac2c1 :: natural log of 1-AUC, lower is better
\newcommand{\ProbeLogOneMinusAucSdEeVsMmMlpOneeighteight}{0.0000}  % experiments/FIGS/data/probe_ladder_v2_mlp2/analysis/seed_level_results.json :: levels.ee_vs_mm.mlp[0].seed_sd :: sha256 53e1ec81343ac2c1
\newcommand{\ProbeOmaEeVsMmMlpOneeighteight}{---}  % experiments/FIGS/data/probe_ladder_v2_mlp2/analysis/seed_level_results.json :: table[ee_vs_mm,mlp,188].auc :: sha256 53e1ec81343ac2c1 :: not quoted: the AUC reached 1 in a seed, so 1-AUC is only a bound
\newcommand{\ProbeBkgLeftEeVsMmMlpOneeighteight}{0.0}  % experiments/FIGS/data/probe_ladder_v2_mlp2/analysis/seed_level_results.json :: levels.ee_vs_mm.mlp[0].rejection_points['0.90'].mean_n_bkg_pass :: sha256 53e1ec81343ac2c1 :: background jets passing the cut, mean over seeds
\newcommand{\ProbeRejEeVsMmMlpOneeighteight}{\ensuremath{>}1{,}571}  % experiments/FIGS/data/probe_ladder_v2_mlp2/analysis/seed_level_results.json :: table[ee_vs_mm,mlp,188].rejection_points['0.90'].rejection :: sha256 53e1ec81343ac2c1 :: at 90% signal efficiency, mean +- SD over seeds; 5 of 5 seeds at the cap
\newcommand{\ProbeAucEeVsMmMlpOnesixtwo}{\ensuremath{1^{\dagger}}}  % experiments/FIGS/data/probe_ladder_v2_mlp2/analysis/seed_level_results.json :: levels.ee_vs_mm.mlp[1].mean_auc :: sha256 53e1ec81343ac2c1 :: 5 seeds, 5 saturated at AUC=1
\newcommand{\ProbeAucSdEeVsMmMlpOnesixtwo}{0.00000}  % experiments/FIGS/data/probe_ladder_v2_mlp2/analysis/seed_level_results.json :: table[ee_vs_mm,mlp,162].auc (SD over seeds, ddof=1) :: sha256 53e1ec81343ac2c1
\newcommand{\ProbeLogOneMinusAucEeVsMmMlpOnesixtwo}{\ensuremath{-}16.9579}  % experiments/FIGS/data/probe_ladder_v2_mlp2/analysis/seed_level_results.json :: levels.ee_vs_mm.mlp[1].mean :: sha256 53e1ec81343ac2c1 :: natural log of 1-AUC, lower is better
\newcommand{\ProbeLogOneMinusAucSdEeVsMmMlpOnesixtwo}{0.0000}  % experiments/FIGS/data/probe_ladder_v2_mlp2/analysis/seed_level_results.json :: levels.ee_vs_mm.mlp[1].seed_sd :: sha256 53e1ec81343ac2c1
\newcommand{\ProbeOmaEeVsMmMlpOnesixtwo}{---}  % experiments/FIGS/data/probe_ladder_v2_mlp2/analysis/seed_level_results.json :: table[ee_vs_mm,mlp,162].auc :: sha256 53e1ec81343ac2c1 :: not quoted: the AUC reached 1 in a seed, so 1-AUC is only a bound
\newcommand{\ProbeBkgLeftEeVsMmMlpOnesixtwo}{0.0}  % experiments/FIGS/data/probe_ladder_v2_mlp2/analysis/seed_level_results.json :: levels.ee_vs_mm.mlp[1].rejection_points['0.90'].mean_n_bkg_pass :: sha256 53e1ec81343ac2c1 :: background jets passing the cut, mean over seeds
\newcommand{\ProbeRejEeVsMmMlpOnesixtwo}{\ensuremath{>}1{,}571}  % experiments/FIGS/data/probe_ladder_v2_mlp2/analysis/seed_level_results.json :: table[ee_vs_mm,mlp,162].rejection_points['0.90'].rejection :: sha256 53e1ec81343ac2c1 :: at 90% signal efficiency, mean +- SD over seeds; 5 of 5 seeds at the cap
\newcommand{\ProbeAucEeVsMmMlpFourthree}{0.99657}  % experiments/FIGS/data/probe_ladder_v2_mlp2/analysis/seed_level_results.json :: levels.ee_vs_mm.mlp[2].mean_auc :: sha256 53e1ec81343ac2c1 :: 5 seeds, 0 saturated at AUC=1
\newcommand{\ProbeAucSdEeVsMmMlpFourthree}{0.00300}  % experiments/FIGS/data/probe_ladder_v2_mlp2/analysis/seed_level_results.json :: table[ee_vs_mm,mlp,43].auc (SD over seeds, ddof=1) :: sha256 53e1ec81343ac2c1
\newcommand{\ProbeLogOneMinusAucEeVsMmMlpFourthree}{\ensuremath{-}6.0371}  % experiments/FIGS/data/probe_ladder_v2_mlp2/analysis/seed_level_results.json :: levels.ee_vs_mm.mlp[2].mean :: sha256 53e1ec81343ac2c1 :: natural log of 1-AUC, lower is better
\newcommand{\ProbeLogOneMinusAucSdEeVsMmMlpFourthree}{0.9962}  % experiments/FIGS/data/probe_ladder_v2_mlp2/analysis/seed_level_results.json :: levels.ee_vs_mm.mlp[2].seed_sd :: sha256 53e1ec81343ac2c1
\newcommand{\ProbeOmaEeVsMmMlpFourthree}{\ensuremath{(3.4 \pm 3.0)\times10^{-3}}}  % experiments/FIGS/data/probe_ladder_v2_mlp2/analysis/seed_level_results.json :: table[ee_vs_mm,mlp,43].auc :: sha256 53e1ec81343ac2c1 :: 1-AUC, mean +- SD over 5 seeds
\newcommand{\ProbeBkgLeftEeVsMmMlpFourthree}{33.4}  % experiments/FIGS/data/probe_ladder_v2_mlp2/analysis/seed_level_results.json :: levels.ee_vs_mm.mlp[2].rejection_points['0.90'].mean_n_bkg_pass :: sha256 53e1ec81343ac2c1 :: background jets passing the cut, mean over seeds
\newcommand{\ProbeRejEeVsMmMlpFourthree}{400\,\ensuremath{\pm}\,500}  % experiments/FIGS/data/probe_ladder_v2_mlp2/analysis/seed_level_results.json :: table[ee_vs_mm,mlp,43].rejection_points['0.90'].rejection :: sha256 53e1ec81343ac2c1 :: at 90% signal efficiency, mean +- SD over seeds; 0 of 5 seeds at the cap
\newcommand{\ProbeAucEeVsMmMlpOneseven}{0.99338}  % experiments/FIGS/data/probe_ladder_v2_mlp2/analysis/seed_level_results.json :: levels.ee_vs_mm.mlp[3].mean_auc :: sha256 53e1ec81343ac2c1 :: 5 seeds, 0 saturated at AUC=1
\newcommand{\ProbeAucSdEeVsMmMlpOneseven}{0.00065}  % experiments/FIGS/data/probe_ladder_v2_mlp2/analysis/seed_level_results.json :: table[ee_vs_mm,mlp,17].auc (SD over seeds, ddof=1) :: sha256 53e1ec81343ac2c1
\newcommand{\ProbeLogOneMinusAucEeVsMmMlpOneseven}{\ensuremath{-}5.0211}  % experiments/FIGS/data/probe_ladder_v2_mlp2/analysis/seed_level_results.json :: levels.ee_vs_mm.mlp[3].mean :: sha256 53e1ec81343ac2c1 :: natural log of 1-AUC, lower is better
\newcommand{\ProbeLogOneMinusAucSdEeVsMmMlpOneseven}{0.0997}  % experiments/FIGS/data/probe_ladder_v2_mlp2/analysis/seed_level_results.json :: levels.ee_vs_mm.mlp[3].seed_sd :: sha256 53e1ec81343ac2c1
\newcommand{\ProbeOmaEeVsMmMlpOneseven}{\ensuremath{(6.6 \pm 0.7)\times10^{-3}}}  % experiments/FIGS/data/probe_ladder_v2_mlp2/analysis/seed_level_results.json :: table[ee_vs_mm,mlp,17].auc :: sha256 53e1ec81343ac2c1 :: 1-AUC, mean +- SD over 5 seeds
\newcommand{\ProbeBkgLeftEeVsMmMlpOneseven}{59.2}  % experiments/FIGS/data/probe_ladder_v2_mlp2/analysis/seed_level_results.json :: levels.ee_vs_mm.mlp[3].rejection_points['0.90'].mean_n_bkg_pass :: sha256 53e1ec81343ac2c1 :: background jets passing the cut, mean over seeds
\newcommand{\ProbeRejEeVsMmMlpOneseven}{81\,\ensuremath{\pm}\,12}  % experiments/FIGS/data/probe_ladder_v2_mlp2/analysis/seed_level_results.json :: table[ee_vs_mm,mlp,17].rejection_points['0.90'].rejection :: sha256 53e1ec81343ac2c1 :: at 90% signal efficiency, mean +- SD over seeds; 0 of 5 seeds at the cap
\newcommand{\ProbeOmaRatioEeVsMmMlpOnesevenOverFourthree}{1.93}  % experiments/FIGS/data/probe_ladder_v2_mlp2/analysis/seed_level_results.json :: table[ee_vs_mm,mlp,17/43].auc :: sha256 53e1ec81343ac2c1 :: ratio of the seed means of 1-AUC, coarser over finer
\newcommand{\ProbeAucRetainedTopologyLinearOneeighteight}{0.99413}  % experiments/FIGS/data/probe_ladder_v2_mlp2/analysis/seed_level_results.json :: levels.retained_topology.linear[0].mean_auc :: sha256 53e1ec81343ac2c1 :: 5 seeds, 0 saturated at AUC=1
\newcommand{\ProbeAucSdRetainedTopologyLinearOneeighteight}{0.00011}  % experiments/FIGS/data/probe_ladder_v2_mlp2/analysis/seed_level_results.json :: table[retained_topology,linear,188].auc (SD over seeds, ddof=1) :: sha256 53e1ec81343ac2c1
\newcommand{\ProbeLogOneMinusAucRetainedTopologyLinearOneeighteight}{\ensuremath{-}5.1388}  % experiments/FIGS/data/probe_ladder_v2_mlp2/analysis/seed_level_results.json :: levels.retained_topology.linear[0].mean :: sha256 53e1ec81343ac2c1 :: natural log of 1-AUC, lower is better
\newcommand{\ProbeLogOneMinusAucSdRetainedTopologyLinearOneeighteight}{0.0187}  % experiments/FIGS/data/probe_ladder_v2_mlp2/analysis/seed_level_results.json :: levels.retained_topology.linear[0].seed_sd :: sha256 53e1ec81343ac2c1
\newcommand{\ProbeOmaRetainedTopologyLinearOneeighteight}{\ensuremath{(5.87 \pm 0.11)\times10^{-3}}}  % experiments/FIGS/data/probe_ladder_v2_mlp2/analysis/seed_level_results.json :: table[retained_topology,linear,188].auc :: sha256 53e1ec81343ac2c1 :: 1-AUC, mean +- SD over 5 seeds
\newcommand{\ProbeBkgLeftRetainedTopologyLinearOneeighteight}{15.4}  % experiments/FIGS/data/probe_ladder_v2_mlp2/analysis/seed_level_results.json :: levels.retained_topology.linear[0].rejection_points['0.90'].mean_n_bkg_pass :: sha256 53e1ec81343ac2c1 :: background jets passing the cut, mean over seeds
\newcommand{\ProbeRejRetainedTopologyLinearOneeighteight}{71\,\ensuremath{\pm}\,7}  % experiments/FIGS/data/probe_ladder_v2_mlp2/analysis/seed_level_results.json :: table[retained_topology,linear,188].rejection_points['0.90'].rejection :: sha256 53e1ec81343ac2c1 :: at 90% signal efficiency, mean +- SD over seeds; 0 of 5 seeds at the cap
\newcommand{\ProbeAucRetainedTopologyLinearOnesixtwo}{0.99413}  % experiments/FIGS/data/probe_ladder_v2_mlp2/analysis/seed_level_results.json :: levels.retained_topology.linear[1].mean_auc :: sha256 53e1ec81343ac2c1 :: 5 seeds, 0 saturated at AUC=1
\newcommand{\ProbeAucSdRetainedTopologyLinearOnesixtwo}{0.00030}  % experiments/FIGS/data/probe_ladder_v2_mlp2/analysis/seed_level_results.json :: table[retained_topology,linear,162].auc (SD over seeds, ddof=1) :: sha256 53e1ec81343ac2c1
\newcommand{\ProbeLogOneMinusAucRetainedTopologyLinearOnesixtwo}{\ensuremath{-}5.1395}  % experiments/FIGS/data/probe_ladder_v2_mlp2/analysis/seed_level_results.json :: levels.retained_topology.linear[1].mean :: sha256 53e1ec81343ac2c1 :: natural log of 1-AUC, lower is better
\newcommand{\ProbeLogOneMinusAucSdRetainedTopologyLinearOnesixtwo}{0.0510}  % experiments/FIGS/data/probe_ladder_v2_mlp2/analysis/seed_level_results.json :: levels.retained_topology.linear[1].seed_sd :: sha256 53e1ec81343ac2c1
\newcommand{\ProbeOmaRetainedTopologyLinearOnesixtwo}{\ensuremath{(5.87 \pm 0.30)\times10^{-3}}}  % experiments/FIGS/data/probe_ladder_v2_mlp2/analysis/seed_level_results.json :: table[retained_topology,linear,162].auc :: sha256 53e1ec81343ac2c1 :: 1-AUC, mean +- SD over 5 seeds
\newcommand{\ProbeBkgLeftRetainedTopologyLinearOnesixtwo}{14.0}  % experiments/FIGS/data/probe_ladder_v2_mlp2/analysis/seed_level_results.json :: levels.retained_topology.linear[1].rejection_points['0.90'].mean_n_bkg_pass :: sha256 53e1ec81343ac2c1 :: background jets passing the cut, mean over seeds
\newcommand{\ProbeRejRetainedTopologyLinearOnesixtwo}{79\,\ensuremath{\pm}\,13}  % experiments/FIGS/data/probe_ladder_v2_mlp2/analysis/seed_level_results.json :: table[retained_topology,linear,162].rejection_points['0.90'].rejection :: sha256 53e1ec81343ac2c1 :: at 90% signal efficiency, mean +- SD over seeds; 0 of 5 seeds at the cap
\newcommand{\ProbeAucRetainedTopologyLinearFourthree}{0.99274}  % experiments/FIGS/data/probe_ladder_v2_mlp2/analysis/seed_level_results.json :: levels.retained_topology.linear[2].mean_auc :: sha256 53e1ec81343ac2c1 :: 5 seeds, 0 saturated at AUC=1
\newcommand{\ProbeAucSdRetainedTopologyLinearFourthree}{0.00012}  % experiments/FIGS/data/probe_ladder_v2_mlp2/analysis/seed_level_results.json :: table[retained_topology,linear,43].auc (SD over seeds, ddof=1) :: sha256 53e1ec81343ac2c1
\newcommand{\ProbeLogOneMinusAucRetainedTopologyLinearFourthree}{\ensuremath{-}4.9253}  % experiments/FIGS/data/probe_ladder_v2_mlp2/analysis/seed_level_results.json :: levels.retained_topology.linear[2].mean :: sha256 53e1ec81343ac2c1 :: natural log of 1-AUC, lower is better
\newcommand{\ProbeLogOneMinusAucSdRetainedTopologyLinearFourthree}{0.0169}  % experiments/FIGS/data/probe_ladder_v2_mlp2/analysis/seed_level_results.json :: levels.retained_topology.linear[2].seed_sd :: sha256 53e1ec81343ac2c1
\newcommand{\ProbeOmaRetainedTopologyLinearFourthree}{\ensuremath{(7.26 \pm 0.12)\times10^{-3}}}  % experiments/FIGS/data/probe_ladder_v2_mlp2/analysis/seed_level_results.json :: table[retained_topology,linear,43].auc :: sha256 53e1ec81343ac2c1 :: 1-AUC, mean +- SD over 5 seeds
\newcommand{\ProbeBkgLeftRetainedTopologyLinearFourthree}{16.4}  % experiments/FIGS/data/probe_ladder_v2_mlp2/analysis/seed_level_results.json :: levels.retained_topology.linear[2].rejection_points['0.90'].mean_n_bkg_pass :: sha256 53e1ec81343ac2c1 :: background jets passing the cut, mean over seeds
\newcommand{\ProbeRejRetainedTopologyLinearFourthree}{66\,\ensuremath{\pm}\,5}  % experiments/FIGS/data/probe_ladder_v2_mlp2/analysis/seed_level_results.json :: table[retained_topology,linear,43].rejection_points['0.90'].rejection :: sha256 53e1ec81343ac2c1 :: at 90% signal efficiency, mean +- SD over seeds; 0 of 5 seeds at the cap
\newcommand{\ProbeAucRetainedTopologyLinearOneseven}{0.99096}  % experiments/FIGS/data/probe_ladder_v2_mlp2/analysis/seed_level_results.json :: levels.retained_topology.linear[3].mean_auc :: sha256 53e1ec81343ac2c1 :: 5 seeds, 0 saturated at AUC=1
\newcommand{\ProbeAucSdRetainedTopologyLinearOneseven}{0.00092}  % experiments/FIGS/data/probe_ladder_v2_mlp2/analysis/seed_level_results.json :: table[retained_topology,linear,17].auc (SD over seeds, ddof=1) :: sha256 53e1ec81343ac2c1
\newcommand{\ProbeLogOneMinusAucRetainedTopologyLinearOneseven}{\ensuremath{-}4.7095}  % experiments/FIGS/data/probe_ladder_v2_mlp2/analysis/seed_level_results.json :: levels.retained_topology.linear[3].mean :: sha256 53e1ec81343ac2c1 :: natural log of 1-AUC, lower is better
\newcommand{\ProbeLogOneMinusAucSdRetainedTopologyLinearOneseven}{0.0979}  % experiments/FIGS/data/probe_ladder_v2_mlp2/analysis/seed_level_results.json :: levels.retained_topology.linear[3].seed_sd :: sha256 53e1ec81343ac2c1
\newcommand{\ProbeOmaRetainedTopologyLinearOneseven}{\ensuremath{(9.0 \pm 0.9)\times10^{-3}}}  % experiments/FIGS/data/probe_ladder_v2_mlp2/analysis/seed_level_results.json :: table[retained_topology,linear,17].auc :: sha256 53e1ec81343ac2c1 :: 1-AUC, mean +- SD over 5 seeds
\newcommand{\ProbeBkgLeftRetainedTopologyLinearOneseven}{21.8}  % experiments/FIGS/data/probe_ladder_v2_mlp2/analysis/seed_level_results.json :: levels.retained_topology.linear[3].rejection_points['0.90'].mean_n_bkg_pass :: sha256 53e1ec81343ac2c1 :: background jets passing the cut, mean over seeds
\newcommand{\ProbeRejRetainedTopologyLinearOneseven}{51\,\ensuremath{\pm}\,11}  % experiments/FIGS/data/probe_ladder_v2_mlp2/analysis/seed_level_results.json :: table[retained_topology,linear,17].rejection_points['0.90'].rejection :: sha256 53e1ec81343ac2c1 :: at 90% signal efficiency, mean +- SD over seeds; 0 of 5 seeds at the cap
\newcommand{\ProbeOmaRatioRetainedTopologyLinearOnesixtwoOverOneeighteight}{1.00}  % experiments/FIGS/data/probe_ladder_v2_mlp2/analysis/seed_level_results.json :: table[retained_topology,linear,162/188].auc :: sha256 53e1ec81343ac2c1 :: ratio of the seed means of 1-AUC, coarser over finer
\newcommand{\ProbeOmaRatioRetainedTopologyLinearFourthreeOverOneeighteight}{1.24}  % experiments/FIGS/data/probe_ladder_v2_mlp2/analysis/seed_level_results.json :: table[retained_topology,linear,43/188].auc :: sha256 53e1ec81343ac2c1 :: ratio of the seed means of 1-AUC, coarser over finer
\newcommand{\ProbeOmaRatioRetainedTopologyLinearOnesevenOverOneeighteight}{1.54}  % experiments/FIGS/data/probe_ladder_v2_mlp2/analysis/seed_level_results.json :: table[retained_topology,linear,17/188].auc :: sha256 53e1ec81343ac2c1 :: ratio of the seed means of 1-AUC, coarser over finer
\newcommand{\ProbeOmaRatioRetainedTopologyLinearFourthreeOverOnesixtwo}{1.24}  % experiments/FIGS/data/probe_ladder_v2_mlp2/analysis/seed_level_results.json :: table[retained_topology,linear,43/162].auc :: sha256 53e1ec81343ac2c1 :: ratio of the seed means of 1-AUC, coarser over finer
\newcommand{\ProbeOmaRatioRetainedTopologyLinearOnesevenOverOnesixtwo}{1.54}  % experiments/FIGS/data/probe_ladder_v2_mlp2/analysis/seed_level_results.json :: table[retained_topology,linear,17/162].auc :: sha256 53e1ec81343ac2c1 :: ratio of the seed means of 1-AUC, coarser over finer
\newcommand{\ProbeOmaRatioRetainedTopologyLinearOnesevenOverFourthree}{1.25}  % experiments/FIGS/data/probe_ladder_v2_mlp2/analysis/seed_level_results.json :: table[retained_topology,linear,17/43].auc :: sha256 53e1ec81343ac2c1 :: ratio of the seed means of 1-AUC, coarser over finer
\newcommand{\ProbeAucRetainedTopologyMlpOneeighteight}{0.99394}  % experiments/FIGS/data/probe_ladder_v2_mlp2/analysis/seed_level_results.json :: levels.retained_topology.mlp[0].mean_auc :: sha256 53e1ec81343ac2c1 :: 5 seeds, 0 saturated at AUC=1
\newcommand{\ProbeAucSdRetainedTopologyMlpOneeighteight}{0.00032}  % experiments/FIGS/data/probe_ladder_v2_mlp2/analysis/seed_level_results.json :: table[retained_topology,mlp,188].auc (SD over seeds, ddof=1) :: sha256 53e1ec81343ac2c1
\newcommand{\ProbeLogOneMinusAucRetainedTopologyMlpOneeighteight}{\ensuremath{-}5.1078}  % experiments/FIGS/data/probe_ladder_v2_mlp2/analysis/seed_level_results.json :: levels.retained_topology.mlp[0].mean :: sha256 53e1ec81343ac2c1 :: natural log of 1-AUC, lower is better
\newcommand{\ProbeLogOneMinusAucSdRetainedTopologyMlpOneeighteight}{0.0522}  % experiments/FIGS/data/probe_ladder_v2_mlp2/analysis/seed_level_results.json :: levels.retained_topology.mlp[0].seed_sd :: sha256 53e1ec81343ac2c1
\newcommand{\ProbeOmaRetainedTopologyMlpOneeighteight}{\ensuremath{(6.06 \pm 0.32)\times10^{-3}}}  % experiments/FIGS/data/probe_ladder_v2_mlp2/analysis/seed_level_results.json :: table[retained_topology,mlp,188].auc :: sha256 53e1ec81343ac2c1 :: 1-AUC, mean +- SD over 5 seeds
\newcommand{\ProbeBkgLeftRetainedTopologyMlpOneeighteight}{15.8}  % experiments/FIGS/data/probe_ladder_v2_mlp2/analysis/seed_level_results.json :: levels.retained_topology.mlp[0].rejection_points['0.90'].mean_n_bkg_pass :: sha256 53e1ec81343ac2c1 :: background jets passing the cut, mean over seeds
\newcommand{\ProbeRejRetainedTopologyMlpOneeighteight}{69\,\ensuremath{\pm}\,9}  % experiments/FIGS/data/probe_ladder_v2_mlp2/analysis/seed_level_results.json :: table[retained_topology,mlp,188].rejection_points['0.90'].rejection :: sha256 53e1ec81343ac2c1 :: at 90% signal efficiency, mean +- SD over seeds; 0 of 5 seeds at the cap
\newcommand{\ProbeAucRetainedTopologyMlpOnesixtwo}{0.99387}  % experiments/FIGS/data/probe_ladder_v2_mlp2/analysis/seed_level_results.json :: levels.retained_topology.mlp[1].mean_auc :: sha256 53e1ec81343ac2c1 :: 5 seeds, 0 saturated at AUC=1
\newcommand{\ProbeAucSdRetainedTopologyMlpOnesixtwo}{0.00048}  % experiments/FIGS/data/probe_ladder_v2_mlp2/analysis/seed_level_results.json :: table[retained_topology,mlp,162].auc (SD over seeds, ddof=1) :: sha256 53e1ec81343ac2c1
\newcommand{\ProbeLogOneMinusAucRetainedTopologyMlpOnesixtwo}{\ensuremath{-}5.0964}  % experiments/FIGS/data/probe_ladder_v2_mlp2/analysis/seed_level_results.json :: levels.retained_topology.mlp[1].mean :: sha256 53e1ec81343ac2c1 :: natural log of 1-AUC, lower is better
\newcommand{\ProbeLogOneMinusAucSdRetainedTopologyMlpOnesixtwo}{0.0773}  % experiments/FIGS/data/probe_ladder_v2_mlp2/analysis/seed_level_results.json :: levels.retained_topology.mlp[1].seed_sd :: sha256 53e1ec81343ac2c1
\newcommand{\ProbeOmaRetainedTopologyMlpOnesixtwo}{\ensuremath{(6.1 \pm 0.5)\times10^{-3}}}  % experiments/FIGS/data/probe_ladder_v2_mlp2/analysis/seed_level_results.json :: table[retained_topology,mlp,162].auc :: sha256 53e1ec81343ac2c1 :: 1-AUC, mean +- SD over 5 seeds
\newcommand{\ProbeBkgLeftRetainedTopologyMlpOnesixtwo}{16.6}  % experiments/FIGS/data/probe_ladder_v2_mlp2/analysis/seed_level_results.json :: levels.retained_topology.mlp[1].rejection_points['0.90'].mean_n_bkg_pass :: sha256 53e1ec81343ac2c1 :: background jets passing the cut, mean over seeds
\newcommand{\ProbeRejRetainedTopologyMlpOnesixtwo}{67\,\ensuremath{\pm}\,14}  % experiments/FIGS/data/probe_ladder_v2_mlp2/analysis/seed_level_results.json :: table[retained_topology,mlp,162].rejection_points['0.90'].rejection :: sha256 53e1ec81343ac2c1 :: at 90% signal efficiency, mean +- SD over seeds; 0 of 5 seeds at the cap
\newcommand{\ProbeAucRetainedTopologyMlpFourthree}{0.99284}  % experiments/FIGS/data/probe_ladder_v2_mlp2/analysis/seed_level_results.json :: levels.retained_topology.mlp[2].mean_auc :: sha256 53e1ec81343ac2c1 :: 5 seeds, 0 saturated at AUC=1
\newcommand{\ProbeAucSdRetainedTopologyMlpFourthree}{0.00040}  % experiments/FIGS/data/probe_ladder_v2_mlp2/analysis/seed_level_results.json :: table[retained_topology,mlp,43].auc (SD over seeds, ddof=1) :: sha256 53e1ec81343ac2c1
\newcommand{\ProbeLogOneMinusAucRetainedTopologyMlpFourthree}{\ensuremath{-}4.9400}  % experiments/FIGS/data/probe_ladder_v2_mlp2/analysis/seed_level_results.json :: levels.retained_topology.mlp[2].mean :: sha256 53e1ec81343ac2c1 :: natural log of 1-AUC, lower is better
\newcommand{\ProbeLogOneMinusAucSdRetainedTopologyMlpFourthree}{0.0570}  % experiments/FIGS/data/probe_ladder_v2_mlp2/analysis/seed_level_results.json :: levels.retained_topology.mlp[2].seed_sd :: sha256 53e1ec81343ac2c1
\newcommand{\ProbeOmaRetainedTopologyMlpFourthree}{\ensuremath{(7.2 \pm 0.4)\times10^{-3}}}  % experiments/FIGS/data/probe_ladder_v2_mlp2/analysis/seed_level_results.json :: table[retained_topology,mlp,43].auc :: sha256 53e1ec81343ac2c1 :: 1-AUC, mean +- SD over 5 seeds
\newcommand{\ProbeBkgLeftRetainedTopologyMlpFourthree}{17.8}  % experiments/FIGS/data/probe_ladder_v2_mlp2/analysis/seed_level_results.json :: levels.retained_topology.mlp[2].rejection_points['0.90'].mean_n_bkg_pass :: sha256 53e1ec81343ac2c1 :: background jets passing the cut, mean over seeds
\newcommand{\ProbeRejRetainedTopologyMlpFourthree}{61\,\ensuremath{\pm}\,7}  % experiments/FIGS/data/probe_ladder_v2_mlp2/analysis/seed_level_results.json :: table[retained_topology,mlp,43].rejection_points['0.90'].rejection :: sha256 53e1ec81343ac2c1 :: at 90% signal efficiency, mean +- SD over seeds; 0 of 5 seeds at the cap
\newcommand{\ProbeAucRetainedTopologyMlpOneseven}{0.99081}  % experiments/FIGS/data/probe_ladder_v2_mlp2/analysis/seed_level_results.json :: levels.retained_topology.mlp[3].mean_auc :: sha256 53e1ec81343ac2c1 :: 5 seeds, 0 saturated at AUC=1
\newcommand{\ProbeAucSdRetainedTopologyMlpOneseven}{0.00062}  % experiments/FIGS/data/probe_ladder_v2_mlp2/analysis/seed_level_results.json :: table[retained_topology,mlp,17].auc (SD over seeds, ddof=1) :: sha256 53e1ec81343ac2c1
\newcommand{\ProbeLogOneMinusAucRetainedTopologyMlpOneseven}{\ensuremath{-}4.6913}  % experiments/FIGS/data/probe_ladder_v2_mlp2/analysis/seed_level_results.json :: levels.retained_topology.mlp[3].mean :: sha256 53e1ec81343ac2c1 :: natural log of 1-AUC, lower is better
\newcommand{\ProbeLogOneMinusAucSdRetainedTopologyMlpOneseven}{0.0658}  % experiments/FIGS/data/probe_ladder_v2_mlp2/analysis/seed_level_results.json :: levels.retained_topology.mlp[3].seed_sd :: sha256 53e1ec81343ac2c1
\newcommand{\ProbeOmaRetainedTopologyMlpOneseven}{\ensuremath{(9.2 \pm 0.6)\times10^{-3}}}  % experiments/FIGS/data/probe_ladder_v2_mlp2/analysis/seed_level_results.json :: table[retained_topology,mlp,17].auc :: sha256 53e1ec81343ac2c1 :: 1-AUC, mean +- SD over 5 seeds
\newcommand{\ProbeBkgLeftRetainedTopologyMlpOneseven}{21.2}  % experiments/FIGS/data/probe_ladder_v2_mlp2/analysis/seed_level_results.json :: levels.retained_topology.mlp[3].rejection_points['0.90'].mean_n_bkg_pass :: sha256 53e1ec81343ac2c1 :: background jets passing the cut, mean over seeds
\newcommand{\ProbeRejRetainedTopologyMlpOneseven}{51\,\ensuremath{\pm}\,4}  % experiments/FIGS/data/probe_ladder_v2_mlp2/analysis/seed_level_results.json :: table[retained_topology,mlp,17].rejection_points['0.90'].rejection :: sha256 53e1ec81343ac2c1 :: at 90% signal efficiency, mean +- SD over seeds; 0 of 5 seeds at the cap
\newcommand{\ProbeOmaRatioRetainedTopologyMlpOnesixtwoOverOneeighteight}{1.01}  % experiments/FIGS/data/probe_ladder_v2_mlp2/analysis/seed_level_results.json :: table[retained_topology,mlp,162/188].auc :: sha256 53e1ec81343ac2c1 :: ratio of the seed means of 1-AUC, coarser over finer
\newcommand{\ProbeOmaRatioRetainedTopologyMlpFourthreeOverOneeighteight}{1.18}  % experiments/FIGS/data/probe_ladder_v2_mlp2/analysis/seed_level_results.json :: table[retained_topology,mlp,43/188].auc :: sha256 53e1ec81343ac2c1 :: ratio of the seed means of 1-AUC, coarser over finer
\newcommand{\ProbeOmaRatioRetainedTopologyMlpOnesevenOverOneeighteight}{1.52}  % experiments/FIGS/data/probe_ladder_v2_mlp2/analysis/seed_level_results.json :: table[retained_topology,mlp,17/188].auc :: sha256 53e1ec81343ac2c1 :: ratio of the seed means of 1-AUC, coarser over finer
\newcommand{\ProbeOmaRatioRetainedTopologyMlpFourthreeOverOnesixtwo}{1.17}  % experiments/FIGS/data/probe_ladder_v2_mlp2/analysis/seed_level_results.json :: table[retained_topology,mlp,43/162].auc :: sha256 53e1ec81343ac2c1 :: ratio of the seed means of 1-AUC, coarser over finer
\newcommand{\ProbeOmaRatioRetainedTopologyMlpOnesevenOverOnesixtwo}{1.50}  % experiments/FIGS/data/probe_ladder_v2_mlp2/analysis/seed_level_results.json :: table[retained_topology,mlp,17/162].auc :: sha256 53e1ec81343ac2c1 :: ratio of the seed means of 1-AUC, coarser over finer
\newcommand{\ProbeOmaRatioRetainedTopologyMlpOnesevenOverFourthree}{1.28}  % experiments/FIGS/data/probe_ladder_v2_mlp2/analysis/seed_level_results.json :: table[retained_topology,mlp,17/43].auc :: sha256 53e1ec81343ac2c1 :: ratio of the seed means of 1-AUC, coarser over finer
\newcommand{\ProbeAucVisibleContentLinearOneeighteight}{0.99917}  % experiments/FIGS/data/probe_ladder_v2_mlp2/analysis/seed_level_results.json :: levels.visible_content.linear[0].mean_auc :: sha256 53e1ec81343ac2c1 :: 5 seeds, 0 saturated at AUC=1
\newcommand{\ProbeAucSdVisibleContentLinearOneeighteight}{0.00003}  % experiments/FIGS/data/probe_ladder_v2_mlp2/analysis/seed_level_results.json :: table[visible_content,linear,188].auc (SD over seeds, ddof=1) :: sha256 53e1ec81343ac2c1
\newcommand{\ProbeLogOneMinusAucVisibleContentLinearOneeighteight}{\ensuremath{-}7.0999}  % experiments/FIGS/data/probe_ladder_v2_mlp2/analysis/seed_level_results.json :: levels.visible_content.linear[0].mean :: sha256 53e1ec81343ac2c1 :: natural log of 1-AUC, lower is better
\newcommand{\ProbeLogOneMinusAucSdVisibleContentLinearOneeighteight}{0.0358}  % experiments/FIGS/data/probe_ladder_v2_mlp2/analysis/seed_level_results.json :: levels.visible_content.linear[0].seed_sd :: sha256 53e1ec81343ac2c1
\newcommand{\ProbeOmaVisibleContentLinearOneeighteight}{\ensuremath{(8.26 \pm 0.30)\times10^{-4}}}  % experiments/FIGS/data/probe_ladder_v2_mlp2/analysis/seed_level_results.json :: table[visible_content,linear,188].auc :: sha256 53e1ec81343ac2c1 :: 1-AUC, mean +- SD over 5 seeds
\newcommand{\ProbeBkgLeftVisibleContentLinearOneeighteight}{0.4}  % experiments/FIGS/data/probe_ladder_v2_mlp2/analysis/seed_level_results.json :: levels.visible_content.linear[0].rejection_points['0.90'].mean_n_bkg_pass :: sha256 53e1ec81343ac2c1 :: background jets passing the cut, mean over seeds
\newcommand{\ProbeRejVisibleContentLinearOneeighteight}{\ensuremath{>}538}  % experiments/FIGS/data/probe_ladder_v2_mlp2/analysis/seed_level_results.json :: table[visible_content,linear,188].rejection_points['0.90'].rejection :: sha256 53e1ec81343ac2c1 :: at 90% signal efficiency, mean +- SD over seeds; 5 of 5 seeds at the cap
\newcommand{\ProbeAucVisibleContentLinearOnesixtwo}{0.99921}  % experiments/FIGS/data/probe_ladder_v2_mlp2/analysis/seed_level_results.json :: levels.visible_content.linear[1].mean_auc :: sha256 53e1ec81343ac2c1 :: 5 seeds, 0 saturated at AUC=1
\newcommand{\ProbeAucSdVisibleContentLinearOnesixtwo}{0.00004}  % experiments/FIGS/data/probe_ladder_v2_mlp2/analysis/seed_level_results.json :: table[visible_content,linear,162].auc (SD over seeds, ddof=1) :: sha256 53e1ec81343ac2c1
\newcommand{\ProbeLogOneMinusAucVisibleContentLinearOnesixtwo}{\ensuremath{-}7.1480}  % experiments/FIGS/data/probe_ladder_v2_mlp2/analysis/seed_level_results.json :: levels.visible_content.linear[1].mean :: sha256 53e1ec81343ac2c1 :: natural log of 1-AUC, lower is better
\newcommand{\ProbeLogOneMinusAucSdVisibleContentLinearOnesixtwo}{0.0449}  % experiments/FIGS/data/probe_ladder_v2_mlp2/analysis/seed_level_results.json :: levels.visible_content.linear[1].seed_sd :: sha256 53e1ec81343ac2c1
\newcommand{\ProbeOmaVisibleContentLinearOnesixtwo}{\ensuremath{(7.87 \pm 0.35)\times10^{-4}}}  % experiments/FIGS/data/probe_ladder_v2_mlp2/analysis/seed_level_results.json :: table[visible_content,linear,162].auc :: sha256 53e1ec81343ac2c1 :: 1-AUC, mean +- SD over 5 seeds
\newcommand{\ProbeBkgLeftVisibleContentLinearOnesixtwo}{0.6}  % experiments/FIGS/data/probe_ladder_v2_mlp2/analysis/seed_level_results.json :: levels.visible_content.linear[1].rejection_points['0.90'].mean_n_bkg_pass :: sha256 53e1ec81343ac2c1 :: background jets passing the cut, mean over seeds
\newcommand{\ProbeRejVisibleContentLinearOnesixtwo}{\ensuremath{>}538}  % experiments/FIGS/data/probe_ladder_v2_mlp2/analysis/seed_level_results.json :: table[visible_content,linear,162].rejection_points['0.90'].rejection :: sha256 53e1ec81343ac2c1 :: at 90% signal efficiency, mean +- SD over seeds; 5 of 5 seeds at the cap
\newcommand{\ProbeAucVisibleContentLinearFourthree}{0.99918}  % experiments/FIGS/data/probe_ladder_v2_mlp2/analysis/seed_level_results.json :: levels.visible_content.linear[2].mean_auc :: sha256 53e1ec81343ac2c1 :: 5 seeds, 0 saturated at AUC=1
\newcommand{\ProbeAucSdVisibleContentLinearFourthree}{0.00009}  % experiments/FIGS/data/probe_ladder_v2_mlp2/analysis/seed_level_results.json :: table[visible_content,linear,43].auc (SD over seeds, ddof=1) :: sha256 53e1ec81343ac2c1
\newcommand{\ProbeLogOneMinusAucVisibleContentLinearFourthree}{\ensuremath{-}7.1124}  % experiments/FIGS/data/probe_ladder_v2_mlp2/analysis/seed_level_results.json :: levels.visible_content.linear[2].mean :: sha256 53e1ec81343ac2c1 :: natural log of 1-AUC, lower is better
\newcommand{\ProbeLogOneMinusAucSdVisibleContentLinearFourthree}{0.1137}  % experiments/FIGS/data/probe_ladder_v2_mlp2/analysis/seed_level_results.json :: levels.visible_content.linear[2].seed_sd :: sha256 53e1ec81343ac2c1
\newcommand{\ProbeOmaVisibleContentLinearFourthree}{\ensuremath{(8.2 \pm 1.0)\times10^{-4}}}  % experiments/FIGS/data/probe_ladder_v2_mlp2/analysis/seed_level_results.json :: table[visible_content,linear,43].auc :: sha256 53e1ec81343ac2c1 :: 1-AUC, mean +- SD over 5 seeds
\newcommand{\ProbeBkgLeftVisibleContentLinearFourthree}{0.4}  % experiments/FIGS/data/probe_ladder_v2_mlp2/analysis/seed_level_results.json :: levels.visible_content.linear[2].rejection_points['0.90'].mean_n_bkg_pass :: sha256 53e1ec81343ac2c1 :: background jets passing the cut, mean over seeds
\newcommand{\ProbeRejVisibleContentLinearFourthree}{\ensuremath{>}538}  % experiments/FIGS/data/probe_ladder_v2_mlp2/analysis/seed_level_results.json :: table[visible_content,linear,43].rejection_points['0.90'].rejection :: sha256 53e1ec81343ac2c1 :: at 90% signal efficiency, mean +- SD over seeds; 5 of 5 seeds at the cap
\newcommand{\ProbeAucVisibleContentLinearOneseven}{0.99768}  % experiments/FIGS/data/probe_ladder_v2_mlp2/analysis/seed_level_results.json :: levels.visible_content.linear[3].mean_auc :: sha256 53e1ec81343ac2c1 :: 5 seeds, 0 saturated at AUC=1
\newcommand{\ProbeAucSdVisibleContentLinearOneseven}{0.00020}  % experiments/FIGS/data/probe_ladder_v2_mlp2/analysis/seed_level_results.json :: table[visible_content,linear,17].auc (SD over seeds, ddof=1) :: sha256 53e1ec81343ac2c1
\newcommand{\ProbeLogOneMinusAucVisibleContentLinearOneseven}{\ensuremath{-}6.0684}  % experiments/FIGS/data/probe_ladder_v2_mlp2/analysis/seed_level_results.json :: levels.visible_content.linear[3].mean :: sha256 53e1ec81343ac2c1 :: natural log of 1-AUC, lower is better
\newcommand{\ProbeLogOneMinusAucSdVisibleContentLinearOneseven}{0.0889}  % experiments/FIGS/data/probe_ladder_v2_mlp2/analysis/seed_level_results.json :: levels.visible_content.linear[3].seed_sd :: sha256 53e1ec81343ac2c1
\newcommand{\ProbeOmaVisibleContentLinearOneseven}{\ensuremath{(2.32 \pm 0.20)\times10^{-3}}}  % experiments/FIGS/data/probe_ladder_v2_mlp2/analysis/seed_level_results.json :: table[visible_content,linear,17].auc :: sha256 53e1ec81343ac2c1 :: 1-AUC, mean +- SD over 5 seeds
\newcommand{\ProbeBkgLeftVisibleContentLinearOneseven}{3.4}  % experiments/FIGS/data/probe_ladder_v2_mlp2/analysis/seed_level_results.json :: levels.visible_content.linear[3].rejection_points['0.90'].mean_n_bkg_pass :: sha256 53e1ec81343ac2c1 :: background jets passing the cut, mean over seeds
\newcommand{\ProbeRejVisibleContentLinearOneseven}{900\,\ensuremath{\pm}\,400}  % experiments/FIGS/data/probe_ladder_v2_mlp2/analysis/seed_level_results.json :: table[visible_content,linear,17].rejection_points['0.90'].rejection :: sha256 53e1ec81343ac2c1 :: at 90% signal efficiency, mean +- SD over seeds; 0 of 5 seeds at the cap
\newcommand{\ProbeOmaRatioVisibleContentLinearOnesixtwoOverOneeighteight}{0.95}  % experiments/FIGS/data/probe_ladder_v2_mlp2/analysis/seed_level_results.json :: table[visible_content,linear,162/188].auc :: sha256 53e1ec81343ac2c1 :: ratio of the seed means of 1-AUC, coarser over finer
\newcommand{\ProbeOmaRatioVisibleContentLinearFourthreeOverOneeighteight}{0.99}  % experiments/FIGS/data/probe_ladder_v2_mlp2/analysis/seed_level_results.json :: table[visible_content,linear,43/188].auc :: sha256 53e1ec81343ac2c1 :: ratio of the seed means of 1-AUC, coarser over finer
\newcommand{\ProbeOmaRatioVisibleContentLinearOnesevenOverOneeighteight}{2.8}  % experiments/FIGS/data/probe_ladder_v2_mlp2/analysis/seed_level_results.json :: table[visible_content,linear,17/188].auc :: sha256 53e1ec81343ac2c1 :: ratio of the seed means of 1-AUC, coarser over finer
\newcommand{\ProbeOmaRatioVisibleContentLinearFourthreeOverOnesixtwo}{1.04}  % experiments/FIGS/data/probe_ladder_v2_mlp2/analysis/seed_level_results.json :: table[visible_content,linear,43/162].auc :: sha256 53e1ec81343ac2c1 :: ratio of the seed means of 1-AUC, coarser over finer
\newcommand{\ProbeOmaRatioVisibleContentLinearOnesevenOverOnesixtwo}{3.0}  % experiments/FIGS/data/probe_ladder_v2_mlp2/analysis/seed_level_results.json :: table[visible_content,linear,17/162].auc :: sha256 53e1ec81343ac2c1 :: ratio of the seed means of 1-AUC, coarser over finer
\newcommand{\ProbeOmaRatioVisibleContentLinearOnesevenOverFourthree}{2.8}  % experiments/FIGS/data/probe_ladder_v2_mlp2/analysis/seed_level_results.json :: table[visible_content,linear,17/43].auc :: sha256 53e1ec81343ac2c1 :: ratio of the seed means of 1-AUC, coarser over finer
\newcommand{\ProbeAucVisibleContentMlpOneeighteight}{0.99911}  % experiments/FIGS/data/probe_ladder_v2_mlp2/analysis/seed_level_results.json :: levels.visible_content.mlp[0].mean_auc :: sha256 53e1ec81343ac2c1 :: 5 seeds, 0 saturated at AUC=1
\newcommand{\ProbeAucSdVisibleContentMlpOneeighteight}{0.00005}  % experiments/FIGS/data/probe_ladder_v2_mlp2/analysis/seed_level_results.json :: table[visible_content,mlp,188].auc (SD over seeds, ddof=1) :: sha256 53e1ec81343ac2c1
\newcommand{\ProbeLogOneMinusAucVisibleContentMlpOneeighteight}{\ensuremath{-}7.0202}  % experiments/FIGS/data/probe_ladder_v2_mlp2/analysis/seed_level_results.json :: levels.visible_content.mlp[0].mean :: sha256 53e1ec81343ac2c1 :: natural log of 1-AUC, lower is better
\newcommand{\ProbeLogOneMinusAucSdVisibleContentMlpOneeighteight}{0.0570}  % experiments/FIGS/data/probe_ladder_v2_mlp2/analysis/seed_level_results.json :: levels.visible_content.mlp[0].seed_sd :: sha256 53e1ec81343ac2c1
\newcommand{\ProbeOmaVisibleContentMlpOneeighteight}{\ensuremath{(8.9 \pm 0.5)\times10^{-4}}}  % experiments/FIGS/data/probe_ladder_v2_mlp2/analysis/seed_level_results.json :: table[visible_content,mlp,188].auc :: sha256 53e1ec81343ac2c1 :: 1-AUC, mean +- SD over 5 seeds
\newcommand{\ProbeBkgLeftVisibleContentMlpOneeighteight}{0.4}  % experiments/FIGS/data/probe_ladder_v2_mlp2/analysis/seed_level_results.json :: levels.visible_content.mlp[0].rejection_points['0.90'].mean_n_bkg_pass :: sha256 53e1ec81343ac2c1 :: background jets passing the cut, mean over seeds
\newcommand{\ProbeRejVisibleContentMlpOneeighteight}{\ensuremath{>}538}  % experiments/FIGS/data/probe_ladder_v2_mlp2/analysis/seed_level_results.json :: table[visible_content,mlp,188].rejection_points['0.90'].rejection :: sha256 53e1ec81343ac2c1 :: at 90% signal efficiency, mean +- SD over seeds; 5 of 5 seeds at the cap
\newcommand{\ProbeAucVisibleContentMlpOnesixtwo}{0.99918}  % experiments/FIGS/data/probe_ladder_v2_mlp2/analysis/seed_level_results.json :: levels.visible_content.mlp[1].mean_auc :: sha256 53e1ec81343ac2c1 :: 5 seeds, 0 saturated at AUC=1
\newcommand{\ProbeAucSdVisibleContentMlpOnesixtwo}{0.00003}  % experiments/FIGS/data/probe_ladder_v2_mlp2/analysis/seed_level_results.json :: table[visible_content,mlp,162].auc (SD over seeds, ddof=1) :: sha256 53e1ec81343ac2c1
\newcommand{\ProbeLogOneMinusAucVisibleContentMlpOnesixtwo}{\ensuremath{-}7.1010}  % experiments/FIGS/data/probe_ladder_v2_mlp2/analysis/seed_level_results.json :: levels.visible_content.mlp[1].mean :: sha256 53e1ec81343ac2c1 :: natural log of 1-AUC, lower is better
\newcommand{\ProbeLogOneMinusAucSdVisibleContentMlpOnesixtwo}{0.0371}  % experiments/FIGS/data/probe_ladder_v2_mlp2/analysis/seed_level_results.json :: levels.visible_content.mlp[1].seed_sd :: sha256 53e1ec81343ac2c1
\newcommand{\ProbeOmaVisibleContentMlpOnesixtwo}{\ensuremath{(8.25 \pm 0.30)\times10^{-4}}}  % experiments/FIGS/data/probe_ladder_v2_mlp2/analysis/seed_level_results.json :: table[visible_content,mlp,162].auc :: sha256 53e1ec81343ac2c1 :: 1-AUC, mean +- SD over 5 seeds
\newcommand{\ProbeBkgLeftVisibleContentMlpOnesixtwo}{0.2}  % experiments/FIGS/data/probe_ladder_v2_mlp2/analysis/seed_level_results.json :: levels.visible_content.mlp[1].rejection_points['0.90'].mean_n_bkg_pass :: sha256 53e1ec81343ac2c1 :: background jets passing the cut, mean over seeds
\newcommand{\ProbeRejVisibleContentMlpOnesixtwo}{\ensuremath{>}538}  % experiments/FIGS/data/probe_ladder_v2_mlp2/analysis/seed_level_results.json :: table[visible_content,mlp,162].rejection_points['0.90'].rejection :: sha256 53e1ec81343ac2c1 :: at 90% signal efficiency, mean +- SD over seeds; 5 of 5 seeds at the cap
\newcommand{\ProbeAucVisibleContentMlpFourthree}{0.99917}  % experiments/FIGS/data/probe_ladder_v2_mlp2/analysis/seed_level_results.json :: levels.visible_content.mlp[2].mean_auc :: sha256 53e1ec81343ac2c1 :: 5 seeds, 0 saturated at AUC=1
\newcommand{\ProbeAucSdVisibleContentMlpFourthree}{0.00006}  % experiments/FIGS/data/probe_ladder_v2_mlp2/analysis/seed_level_results.json :: table[visible_content,mlp,43].auc (SD over seeds, ddof=1) :: sha256 53e1ec81343ac2c1
\newcommand{\ProbeLogOneMinusAucVisibleContentMlpFourthree}{\ensuremath{-}7.0957}  % experiments/FIGS/data/probe_ladder_v2_mlp2/analysis/seed_level_results.json :: levels.visible_content.mlp[2].mean :: sha256 53e1ec81343ac2c1 :: natural log of 1-AUC, lower is better
\newcommand{\ProbeLogOneMinusAucSdVisibleContentMlpFourthree}{0.0773}  % experiments/FIGS/data/probe_ladder_v2_mlp2/analysis/seed_level_results.json :: levels.visible_content.mlp[2].seed_sd :: sha256 53e1ec81343ac2c1
\newcommand{\ProbeOmaVisibleContentMlpFourthree}{\ensuremath{(8.3 \pm 0.6)\times10^{-4}}}  % experiments/FIGS/data/probe_ladder_v2_mlp2/analysis/seed_level_results.json :: table[visible_content,mlp,43].auc :: sha256 53e1ec81343ac2c1 :: 1-AUC, mean +- SD over 5 seeds
\newcommand{\ProbeBkgLeftVisibleContentMlpFourthree}{0.4}  % experiments/FIGS/data/probe_ladder_v2_mlp2/analysis/seed_level_results.json :: levels.visible_content.mlp[2].rejection_points['0.90'].mean_n_bkg_pass :: sha256 53e1ec81343ac2c1 :: background jets passing the cut, mean over seeds
\newcommand{\ProbeRejVisibleContentMlpFourthree}{\ensuremath{>}538}  % experiments/FIGS/data/probe_ladder_v2_mlp2/analysis/seed_level_results.json :: table[visible_content,mlp,43].rejection_points['0.90'].rejection :: sha256 53e1ec81343ac2c1 :: at 90% signal efficiency, mean +- SD over seeds; 5 of 5 seeds at the cap
\newcommand{\ProbeAucVisibleContentMlpOneseven}{0.99793}  % experiments/FIGS/data/probe_ladder_v2_mlp2/analysis/seed_level_results.json :: levels.visible_content.mlp[3].mean_auc :: sha256 53e1ec81343ac2c1 :: 5 seeds, 0 saturated at AUC=1
\newcommand{\ProbeAucSdVisibleContentMlpOneseven}{0.00020}  % experiments/FIGS/data/probe_ladder_v2_mlp2/analysis/seed_level_results.json :: table[visible_content,mlp,17].auc (SD over seeds, ddof=1) :: sha256 53e1ec81343ac2c1
\newcommand{\ProbeLogOneMinusAucVisibleContentMlpOneseven}{\ensuremath{-}6.1839}  % experiments/FIGS/data/probe_ladder_v2_mlp2/analysis/seed_level_results.json :: levels.visible_content.mlp[3].mean :: sha256 53e1ec81343ac2c1 :: natural log of 1-AUC, lower is better
\newcommand{\ProbeLogOneMinusAucSdVisibleContentMlpOneseven}{0.1014}  % experiments/FIGS/data/probe_ladder_v2_mlp2/analysis/seed_level_results.json :: levels.visible_content.mlp[3].seed_sd :: sha256 53e1ec81343ac2c1
\newcommand{\ProbeOmaVisibleContentMlpOneseven}{\ensuremath{(2.07 \pm 0.20)\times10^{-3}}}  % experiments/FIGS/data/probe_ladder_v2_mlp2/analysis/seed_level_results.json :: table[visible_content,mlp,17].auc :: sha256 53e1ec81343ac2c1 :: 1-AUC, mean +- SD over 5 seeds
\newcommand{\ProbeBkgLeftVisibleContentMlpOneseven}{3.4}  % experiments/FIGS/data/probe_ladder_v2_mlp2/analysis/seed_level_results.json :: levels.visible_content.mlp[3].rejection_points['0.90'].mean_n_bkg_pass :: sha256 53e1ec81343ac2c1 :: background jets passing the cut, mean over seeds
\newcommand{\ProbeRejVisibleContentMlpOneseven}{\ensuremath{\geq}638\ensuremath{^{\ast}}}  % experiments/FIGS/data/probe_ladder_v2_mlp2/analysis/seed_level_results.json :: table[visible_content,mlp,17].rejection_points['0.90'].rejection :: sha256 53e1ec81343ac2c1 :: at 90% signal efficiency, mean +- SD over seeds; 1 of 5 seeds at the cap
\newcommand{\ProbeOmaRatioVisibleContentMlpOnesixtwoOverOneeighteight}{0.92}  % experiments/FIGS/data/probe_ladder_v2_mlp2/analysis/seed_level_results.json :: table[visible_content,mlp,162/188].auc :: sha256 53e1ec81343ac2c1 :: ratio of the seed means of 1-AUC, coarser over finer
\newcommand{\ProbeOmaRatioVisibleContentMlpFourthreeOverOneeighteight}{0.93}  % experiments/FIGS/data/probe_ladder_v2_mlp2/analysis/seed_level_results.json :: table[visible_content,mlp,43/188].auc :: sha256 53e1ec81343ac2c1 :: ratio of the seed means of 1-AUC, coarser over finer
\newcommand{\ProbeOmaRatioVisibleContentMlpOnesevenOverOneeighteight}{2.3}  % experiments/FIGS/data/probe_ladder_v2_mlp2/analysis/seed_level_results.json :: table[visible_content,mlp,17/188].auc :: sha256 53e1ec81343ac2c1 :: ratio of the seed means of 1-AUC, coarser over finer
\newcommand{\ProbeOmaRatioVisibleContentMlpFourthreeOverOnesixtwo}{1.01}  % experiments/FIGS/data/probe_ladder_v2_mlp2/analysis/seed_level_results.json :: table[visible_content,mlp,43/162].auc :: sha256 53e1ec81343ac2c1 :: ratio of the seed means of 1-AUC, coarser over finer
\newcommand{\ProbeOmaRatioVisibleContentMlpOnesevenOverOnesixtwo}{2.5}  % experiments/FIGS/data/probe_ladder_v2_mlp2/analysis/seed_level_results.json :: table[visible_content,mlp,17/162].auc :: sha256 53e1ec81343ac2c1 :: ratio of the seed means of 1-AUC, coarser over finer
\newcommand{\ProbeOmaRatioVisibleContentMlpOnesevenOverFourthree}{2.5}  % experiments/FIGS/data/probe_ladder_v2_mlp2/analysis/seed_level_results.json :: table[visible_content,mlp,17/43].auc :: sha256 53e1ec81343ac2c1 :: ratio of the seed means of 1-AUC, coarser over finer
\newcommand{\VocabSizeLoneeighteight}{188}  % configs/labelmaps/rung_label_maps.v1.csv :: column L188 (distinct group count) :: sha256 412f9aa3b58db412
\newcommand{\VocabSizeLonesixtwo}{162}  % configs/labelmaps/rung_label_maps.v1.csv :: column L162 (distinct group count) :: sha256 412f9aa3b58db412
\newcommand{\VocabSizeRsixthreeQone}{64}  % configs/labelmaps/rung_label_maps.v1.csv :: column R63_Q1 (distinct group count) :: sha256 412f9aa3b58db412
\newcommand{\VocabSizeRfourtwoQone}{43}  % configs/labelmaps/rung_label_maps.v1.csv :: column R42_Q1 (distinct group count) :: sha256 412f9aa3b58db412
\newcommand{\VocabSizeRtwonineQone}{30}  % configs/labelmaps/rung_label_maps.v1.csv :: column R29_Q1 (distinct group count) :: sha256 412f9aa3b58db412
\newcommand{\VocabSizeRonesixQone}{17}  % configs/labelmaps/rung_label_maps.v1.csv :: column R16_Q1 (distinct group count) :: sha256 412f9aa3b58db412
\newcommand{\VocabSizeRthreeVIS}{4}  % configs/labelmaps/rung_label_maps.v1.csv :: column R3_VIS (distinct group count) :: sha256 412f9aa3b58db412
\newcommand{\VocabSizeRoneQone}{2}  % configs/labelmaps/rung_label_maps.v1.csv :: column R1_Q1 (distinct group count) :: sha256 412f9aa3b58db412
\newcommand{\ProbeSplitTrain}{60\%}  % experiments/EVAL/probe.py :: SPLIT_FRACTIONS :: sha256 485db0966e660588 :: share of the probe sample; the same jets for every model
\newcommand{\ProbeSplitVal}{20\%}  % experiments/EVAL/probe.py :: SPLIT_FRACTIONS :: sha256 485db0966e660588 :: share of the probe sample; the same jets for every model
\newcommand{\ProbeSplitTest}{20\%}  % experiments/EVAL/probe.py :: SPLIT_FRACTIONS :: sha256 485db0966e660588 :: share of the probe sample; the same jets for every model
\newcommand{\ProbeMlpHidden}{256}  % experiments/EVAL/probe.py :: _fit_mlp :: sha256 485db0966e660588 :: hidden width of the MLP probe
\newcommand{\ProbeMlpInits}{3}  % experiments/EVAL/probe.py :: MLP_SEEDS :: sha256 485db0966e660588 :: MLP initialisations whose scores are averaged per cell
\newcommand{\ProbeMlpLrDrop}{10}  % experiments/EVAL/probe.py :: MLP_SCHEDULE['lr_factor'] :: sha256 485db0966e660588 :: factor the MLP's learning rate drops by on a plateau
\newcommand{\ProbeMlpPlateaus}{three}  % experiments/EVAL/probe.py :: MLP_SCHEDULE: stop_patience / (lr_patience + 1) :: sha256 485db0966e660588 :: patience windows without gain after which the MLP stops
\newcommand{\ProbeCMin}{\ensuremath{10^{-2}}}  % experiments/EVAL/probe.py :: C_GRID :: sha256 485db0966e660588 :: end of the logistic-regression inverse-regularisation grid
\newcommand{\ProbeCMax}{\ensuremath{10^{2}}}  % experiments/EVAL/probe.py :: C_GRID :: sha256 485db0966e660588 :: end of the logistic-regression inverse-regularisation grid
\newcommand{\ProbeNBkgTestBvcFourprong}{3{,}158}  % experiments/FIGS/data/probe_ladder_v2_mlp2/s1.json :: tasks.bvc_4prong.n_background_test :: sha256 f7e1de0e43be89b8 :: test-sample background jets
\newcommand{\ProbeNBkgTestBvcQcd}{5{,}371}  % experiments/FIGS/data/probe_ladder_v2_mlp2/s1.json :: tasks.bvc_qcd.n_background_test :: sha256 f7e1de0e43be89b8 :: test-sample background jets
\newcommand{\ProbeNBkgTestBvcResonant}{11{,}876}  % experiments/FIGS/data/probe_ladder_v2_mlp2/s1.json :: tasks.bvc_resonant.n_background_test :: sha256 f7e1de0e43be89b8 :: test-sample background jets
\newcommand{\ProbeNBkgTestEeVsMm}{4{,}708}  % experiments/FIGS/data/probe_ladder_v2_mlp2/s1.json :: tasks.ee_vs_mm.n_background_test :: sha256 f7e1de0e43be89b8 :: test-sample background jets
\newcommand{\ProbeNBkgTestRetainedTopology}{1{,}078}  % experiments/FIGS/data/probe_ladder_v2_mlp2/s1.json :: tasks.retained_topology.n_background_test :: sha256 f7e1de0e43be89b8 :: test-sample background jets
\newcommand{\ProbeNBkgTestVisibleContent}{2{,}554}  % experiments/FIGS/data/probe_ladder_v2_mlp2/s1.json :: tasks.visible_content.n_background_test :: sha256 f7e1de0e43be89b8 :: test-sample background jets
\newcommand{\UseDiesAtSophonEqfour}{43}  % configs/labelmaps/usecase_survival.v1.json :: sophon_eq4.dies_at :: sha256 122275213c349593 :: first vocabulary size at which it is not constructible
\newcommand{\UseLastAliveSophonEqfour}{64}  % configs/labelmaps/usecase_survival.v1.json :: sophon_eq4.last_rung_alive :: sha256 122275213c349593
\newcommand{\UseDiesAtSophonEqsixAone}{64}  % configs/labelmaps/usecase_survival.v1.json :: sophon_eq6_a1.dies_at :: sha256 122275213c349593 :: first vocabulary size at which it is not constructible
\newcommand{\UseLastAliveSophonEqsixAone}{162}  % configs/labelmaps/usecase_survival.v1.json :: sophon_eq6_a1.last_rung_alive :: sha256 122275213c349593
\newcommand{\UseDiesAtSophonEqsixAthree}{64}  % configs/labelmaps/usecase_survival.v1.json :: sophon_eq6_a3.dies_at :: sha256 122275213c349593 :: first vocabulary size at which it is not constructible
\newcommand{\UseLastAliveSophonEqsixAthree}{162}  % configs/labelmaps/usecase_survival.v1.json :: sophon_eq6_a3.last_rung_alive :: sha256 122275213c349593
\newcommand{\UseDiesAtSophonEqsevenNine}{64}  % configs/labelmaps/usecase_survival.v1.json :: sophon_eq7_9.dies_at :: sha256 122275213c349593 :: first vocabulary size at which it is not constructible
\newcommand{\UseLastAliveSophonEqsevenNine}{162}  % configs/labelmaps/usecase_survival.v1.json :: sophon_eq7_9.last_rung_alive :: sha256 122275213c349593
\newcommand{\UseDiesAtVcbEqone}{64}  % configs/labelmaps/usecase_survival.v1.json :: vcb_eq1.dies_at :: sha256 122275213c349593 :: first vocabulary size at which it is not constructible
\newcommand{\UseLastAliveVcbEqone}{162}  % configs/labelmaps/usecase_survival.v1.json :: vcb_eq1.last_rung_alive :: sha256 122275213c349593
\newcommand{\UseDiesAtGnthreexQcdSubclass}{162}  % configs/labelmaps/usecase_survival.v1.json :: gn3x_qcd_subclass.dies_at :: sha256 122275213c349593 :: first vocabulary size at which it is not constructible
\newcommand{\UseLastAliveGnthreexQcdSubclass}{188}  % configs/labelmaps/usecase_survival.v1.json :: gn3x_qcd_subclass.last_rung_alive :: sha256 122275213c349593
\newcommand{\UseWlikeLastAlive}{162}  % configs/labelmaps/rung_label_maps.v1.csv :: groups at each level vs {label_X_cs, label_X_qq} and that plus the QCD classes :: sha256 412f9aa3b58db412 :: last level at which Sophon's W-like score g_W(2) is constructible
\newcommand{\UseWlikeDiesAt}{64}  % configs/labelmaps/rung_label_maps.v1.csv :: groups at each level vs {label_X_cs, label_X_qq} and that plus the QCD classes :: sha256 412f9aa3b58db412 :: first level at which it is not
\newcommand{\DesignEpochs}{80}  % experiments/MTX/k8s/job-mtx-l188-s1-raunav.yaml :: --num-epochs :: sha256 96fef4b6337aaa27
\newcommand{\DesignSamplesPerEpoch}{10{,}240{,}000}  % experiments/MTX/k8s/job-mtx-l188-s1-raunav.yaml :: --samples-per-epoch :: sha256 96fef4b6337aaa27
\newcommand{\DesignExamplesSeen}{\ensuremath{8.192\times10^{8}}}  % experiments/MTX/k8s/job-mtx-l188-s1-raunav.yaml :: --num-epochs x --samples-per-epoch :: sha256 96fef4b6337aaa27 :: training examples drawn over the run
\newcommand{\DesignPassesSelected}{7.45}  % experiments/FIGS/data/v2/jc2_train/counts.json :: n_selected :: sha256 b0b5a8b63aa808c0 :: --num-epochs x --samples-per-epoch of job-mtx-l188-s1-raunav.yaml / n_selected: passes over the selected training jets
\newcommand{\DesignBatchSize}{512}  % experiments/MTX/k8s/job-mtx-l188-s1-raunav.yaml :: --batch-size :: sha256 96fef4b6337aaa27
\newcommand{\DesignStartLr}{\ensuremath{5\times10^{-4}}}  % experiments/MTX/k8s/job-mtx-l188-s1-raunav.yaml :: --start-lr :: sha256 96fef4b6337aaa27
\newcommand{\DesignHeadHidden}{512}  % experiments/MTX/k8s/job-mtx-l188-s1-raunav.yaml :: -o fc_params (hidden width) :: sha256 96fef4b6337aaa27
\newcommand{\DesignParticleBlocks}{8}  % experiments/E1/ParT_sophon_arch_10c.py :: num_layers=(\d+) :: sha256 5793faa7d55a8b0d
\newcommand{\DesignClassBlocks}{2}  % experiments/E1/ParT_sophon_arch_10c.py :: num_cls_layers=(\d+) :: sha256 5793faa7d55a8b0d
\newcommand{\DesignHeads}{8}  % experiments/E1/ParT_sophon_arch_10c.py :: num_heads=(\d+) :: sha256 5793faa7d55a8b0d
\newcommand{\DesignEmbedDim}{128}  % experiments/E1/ParT_sophon_arch_10c.py :: \bembed_dims=\[\d+, \d+, (\d+)\] :: sha256 5793faa7d55a8b0d
\newcommand{\DesignMpmMaskRate}{40\%}  % experiments/MTX/mpm.py :: DEFAULT_MASK_RATE :: sha256 49f0c8c6666d0211 :: share of each jet's particles hidden from the encoder
\newcommand{\DesignPtMin}{200}  % configs/arms/L188.yaml :: selection :: sha256 d5a188dc80885515 :: GeV
\newcommand{\DesignPtMax}{2{,}500}  % configs/arms/L188.yaml :: selection :: sha256 d5a188dc80885515 :: GeV
\newcommand{\DesignMsdMin}{20}  % configs/arms/L188.yaml :: selection :: sha256 d5a188dc80885515 :: GeV
\newcommand{\DesignMsdMax}{500}  % configs/arms/L188.yaml :: selection :: sha256 d5a188dc80885515 :: GeV
\newcommand{\LitJetClassTwoNJetsMillion}{139}  % experiments/FIGS/data/literature_facts.json :: facts[0].value :: sha256 713dba2a5e6e1dd1 :: arXiv:2405.12972v1 App. A, p. 11
\newcommand{\LitJetClassTwoNResonantClasses}{161}  % experiments/FIGS/data/literature_facts.json :: facts[1].value :: sha256 713dba2a5e6e1dd1 :: arXiv:2405.12972v1 App. A, p. 11
\newcommand{\LitJetClassTwoNQcdClasses}{27}  % experiments/FIGS/data/literature_facts.json :: facts[2].value :: sha256 713dba2a5e6e1dd1 :: arXiv:2405.12972v1 App. A, p. 11
\newcommand{\LitJetClassTwoMaxResonanceMass}{500}  % experiments/FIGS/data/literature_facts.json :: facts[3].value :: sha256 713dba2a5e6e1dd1 :: arXiv:2405.12972v1 Sec. II, p. 2
\newcommand{\LitMadgraphVersion}{2.9.18}  % experiments/FIGS/data/literature_facts.json :: facts[4].value :: sha256 713dba2a5e6e1dd1 :: arXiv:2405.12972v1 App. A, p. 11
\newcommand{\LitPythiaVersion}{8.3}  % experiments/FIGS/data/literature_facts.json :: facts[5].value :: sha256 713dba2a5e6e1dd1 :: arXiv:2405.12972v1 App. A, p. 11
\newcommand{\LitDelphesVersion}{3}  % experiments/FIGS/data/literature_facts.json :: facts[6].value :: sha256 713dba2a5e6e1dd1 :: arXiv:2405.12972v1 App. A, p. 11
\newcommand{\LitJetClassTwoPileupMean}{50}  % experiments/FIGS/data/literature_facts.json :: facts[7].value :: sha256 713dba2a5e6e1dd1 :: arXiv:2405.12972v1 Sec. II, p. 2
\newcommand{\LitJetClassTwoJetRadius}{0.8}  % experiments/FIGS/data/literature_facts.json :: facts[8].value :: sha256 713dba2a5e6e1dd1 :: arXiv:2405.12972v1 Sec. II, p. 2
\newcommand{\LitAtlasUnifiedEffDropPercent}{10}  % experiments/FIGS/data/literature_facts.json :: facts[9].value :: sha256 713dba2a5e6e1dd1 :: ATL-PHYS-PUB-2026-013 Sec. 6.5, p. 22
\newcommand{\LitSophonPaperParticleBlocks}{6}  % experiments/FIGS/data/literature_facts.json :: facts[10].value :: sha256 713dba2a5e6e1dd1 :: arXiv:2405.12972v1 App. B.1, p. 12
\newcommand{\LitSophonSamplesPerEpochMillion}{10}  % experiments/FIGS/data/literature_facts.json :: facts[11].value :: sha256 713dba2a5e6e1dd1 :: arXiv:2405.12972v1 App. B.1, p. 12
\newcommand{\LitSophonSubsetSamplesPerEpochMillion}{2.5}  % experiments/FIGS/data/literature_facts.json :: facts[12].value :: sha256 713dba2a5e6e1dd1 :: arXiv:2405.12972v1 App. B.2, p. 12
\newcommand{\LitGlopartNClasses}{314}  % experiments/FIGS/data/literature_facts.json :: facts[13].value :: sha256 713dba2a5e6e1dd1 :: C. Li, PhD thesis, CERN-THESIS-2024-281 Sec. 11.2.1, p. 188
\newcommand{\LitPartFtEpochs}{20}  % experiments/FIGS/data/literature_facts.json :: facts[14].value :: sha256 713dba2a5e6e1dd1 :: arXiv:2202.03772 Sec. 5.2, p. 8
\newcommand{\LitMehtaNCategories}{eight}  % experiments/FIGS/data/literature_facts.json :: facts[15].value :: sha256 713dba2a5e6e1dd1 :: arXiv:2605.05556v2 Abstract
\newcommand{\LitHongPretrainRuns}{2}  % experiments/FIGS/data/literature_facts.json :: facts[16].value :: sha256 713dba2a5e6e1dd1 :: arXiv:2410.23129 App. A.1.1, p. 18
\newcommand{\LitSophonLookaheadK}{6}  % experiments/FIGS/data/literature_facts.json :: facts[17].value :: sha256 713dba2a5e6e1dd1 :: arXiv:2405.12972v1 App. B.1, p. 12
\newcommand{\LitSophonLookaheadAlpha}{0.5}  % experiments/FIGS/data/literature_facts.json :: facts[18].value :: sha256 713dba2a5e6e1dd1 :: arXiv:2405.12972v1 App. B.1, p. 12
\newcommand{\LitSophonRadamBetaOne}{0.95}  % experiments/FIGS/data/literature_facts.json :: facts[19].value :: sha256 713dba2a5e6e1dd1 :: arXiv:2405.12972v1 App. B.1, p. 12
\newcommand{\LitSophonRadamBetaTwo}{0.999}  % experiments/FIGS/data/literature_facts.json :: facts[20].value :: sha256 713dba2a5e6e1dd1 :: arXiv:2405.12972v1 App. B.1, p. 12
\newcommand{\LitSophonRadamEps}{\ensuremath{10^{-5}}}  % experiments/FIGS/data/literature_facts.json :: facts[21].value :: sha256 713dba2a5e6e1dd1 :: arXiv:2405.12972v1 App. B.1, p. 12
\newcommand{\LitGnTwoAuxRejectionDropPercent}{30}  % experiments/FIGS/data/literature_facts.json :: facts[22].value :: sha256 713dba2a5e6e1dd1 :: arXiv:2505.19689v2 (Nature Commun. 17 (2026) 541) p. 9
\newcommand{\LitCmsPartNClasses}{37}  % experiments/FIGS/data/literature_facts.json :: facts[23].value :: sha256 713dba2a5e6e1dd1 :: arXiv:2604.09809v1 (CMS-JME-25-001) Sec. 5.2, p. 7
\newcommand{\LitCmsPartMassLambda}{0.05}  % experiments/FIGS/data/literature_facts.json :: facts[24].value :: sha256 713dba2a5e6e1dd1 :: arXiv:2604.09809v1 (CMS-JME-25-001) Sec. 5.2, p. 7
\newcommand{\LitGlopartVthreeNClasses}{374}  % experiments/FIGS/data/literature_facts.json :: facts[25].value :: sha256 713dba2a5e6e1dd1 :: CMS-DP-2026-104 (CMS public page) Architecture and outputs
\newcommand{\LitHuhCoarseClasses}{127}  % experiments/FIGS/data/literature_facts.json :: facts[26].value :: sha256 713dba2a5e6e1dd1 :: arXiv:1608.08614v2 Sec. 5.3, p. 5
\newcommand{\LitHuhKnnGapPercent}{15}  % experiments/FIGS/data/literature_facts.json :: facts[27].value :: sha256 713dba2a5e6e1dd1 :: arXiv:1608.08614v2 Sec. 5.3, p. 5
\newcommand{\LitRsthreelRuns}{3}  % experiments/FIGS/data/literature_facts.json :: facts[28].value :: sha256 713dba2a5e6e1dd1 :: arXiv:2403.07066v2 (Phys. Rev. D 111 (2025) 032010) Sec. IV, p. 5
\newcommand{\LitSophonParamsMillion}{2.3}  % experiments/FIGS/data/literature_facts.json :: facts[29].value :: sha256 713dba2a5e6e1dd1 :: arXiv:2405.12972v1 App. B.1, p. 12
\newcommand{\LitGlopartVoneNClasses}{37}  % experiments/FIGS/data/literature_facts.json :: facts[30].value :: sha256 713dba2a5e6e1dd1 :: CMS-DP-2026-104 (CMS public page) Model generations
\newcommand{\LitSophonXMass}{200}  % experiments/FIGS/data/literature_facts.json :: facts[31].value :: sha256 713dba2a5e6e1dd1 :: arXiv:2405.12972v1 Sec. III, p. 3
\newcommand{\LitSophonXLumi}{40}  % experiments/FIGS/data/literature_facts.json :: facts[32].value :: sha256 713dba2a5e6e1dd1 :: arXiv:2405.12972v1 Sec. III, Fig. 2
\newcommand{\LitSophonXWindowLo}{150}  % experiments/FIGS/data/literature_facts.json :: facts[33].value :: sha256 713dba2a5e6e1dd1 :: arXiv:2405.12972v1 Sec. III, Fig. 2
\newcommand{\LitSophonXWindowHi}{230}  % experiments/FIGS/data/literature_facts.json :: facts[34].value :: sha256 713dba2a5e6e1dd1 :: arXiv:2405.12972v1 Sec. III, Fig. 2
\newcommand{\LitSophonAblationNClasses}{42}  % experiments/FIGS/data/literature_facts.json :: facts[35].value :: sha256 713dba2a5e6e1dd1 :: arXiv:2405.12972v1 Sec. III
\newcommand{\LitSophonTransferSamplesMillion}{2.5}  % experiments/FIGS/data/literature_facts.json :: facts[36].value :: sha256 713dba2a5e6e1dd1 :: arXiv:2405.12972v1 App. B.5, p. 13
\newcommand{\LitSophonTransferBatch}{1024}  % experiments/FIGS/data/literature_facts.json :: facts[37].value :: sha256 713dba2a5e6e1dd1 :: arXiv:2405.12972v1 App. B.5, p. 13
\newcommand{\LitSophonScratchEpochs}{50}  % experiments/FIGS/data/literature_facts.json :: facts[38].value :: sha256 713dba2a5e6e1dd1 :: arXiv:2405.12972v1 App. B.3
\newcommand{\LitHuhFineClasses}{1000}  % experiments/FIGS/data/literature_facts.json :: facts[39].value :: sha256 713dba2a5e6e1dd1 :: arXiv:1608.08614v2 Sec. 1, p. 1
\newcommand{\LitHuhPascalDropMap}{2.8}  % experiments/FIGS/data/literature_facts.json :: facts[40].value :: sha256 713dba2a5e6e1dd1 :: arXiv:1608.08614v2 Sec. 1, p. 1
\newcommand{\LitHongSuperclasses}{11}  % experiments/FIGS/data/literature_facts.json :: facts[41].value :: sha256 713dba2a5e6e1dd1 :: arXiv:2410.23129 Sec. 6.2
\newcommand{\LitHongDeclineClasses}{\ensuremath{10^{3}}}  % experiments/FIGS/data/literature_facts.json :: facts[42].value :: sha256 713dba2a5e6e1dd1 :: arXiv:2410.23129 Sec. 6.2
\newcommand{\LitObjectivesFullJets}{\ensuremath{1.2\times10^{6}}}  % experiments/FIGS/data/literature_facts.json :: facts[43].value :: sha256 713dba2a5e6e1dd1 :: arXiv:2606.14870 Sec. IV A
\newcommand{\LitObjectivesFewJets}{\ensuremath{10^{4}}}  % experiments/FIGS/data/literature_facts.json :: facts[44].value :: sha256 713dba2a5e6e1dd1 :: arXiv:2606.14870 Sec. IV A
\newcommand{\LitMpmBenefitJets}{10k}  % experiments/FIGS/data/literature_facts.json :: facts[45].value :: sha256 713dba2a5e6e1dd1 :: arXiv:2401.13537 Sec. V
\newcommand{\LitPartTopRejScratch}{1602}  % experiments/FIGS/data/literature_facts.json :: facts[46].value :: sha256 713dba2a5e6e1dd1 :: arXiv:2202.03772 Table 5
\newcommand{\LitPartTopRejFt}{2766}  % experiments/FIGS/data/literature_facts.json :: facts[47].value :: sha256 713dba2a5e6e1dd1 :: arXiv:2202.03772 Table 5
\newcommand{\LitMahajanLinearGapPercent}{7}  % experiments/FIGS/data/literature_facts.json :: facts[48].value :: sha256 713dba2a5e6e1dd1 :: arXiv:1805.00932 Sec. 3.1.2
\newcommand{\LitBukchinCoarse}{33.83}  % experiments/FIGS/data/literature_facts.json :: facts[49].value :: sha256 713dba2a5e6e1dd1 :: arXiv:2012.03515 Table 2, LIVING-17 all-way
\newcommand{\LitBukchinFine}{58.95}  % experiments/FIGS/data/literature_facts.json :: facts[50].value :: sha256 713dba2a5e6e1dd1 :: arXiv:2012.03515 Table 2, LIVING-17 all-way
\newcommand{\LitKornblithRho}{\ensuremath{-0.93}}  % experiments/FIGS/data/literature_facts.json :: facts[51].value :: sha256 713dba2a5e6e1dd1 :: arXiv:2010.16402 Sec. 3.4
\newcommand{\LitPartTopSigEffPercent}{30}  % experiments/FIGS/data/literature_facts.json :: facts[52].value :: sha256 713dba2a5e6e1dd1 :: arXiv:2202.03772 Table 5 (column Rej30%)
\newcommand{\LitPartTopKinFeatures}{7}  % experiments/FIGS/data/literature_facts.json :: facts[53].value :: sha256 713dba2a5e6e1dd1 :: arXiv:2202.03772 Sec. 5.2
\newcommand{\LitBukchinAncor}{45.14}  % experiments/FIGS/data/literature_facts.json :: facts[54].value :: sha256 713dba2a5e6e1dd1 :: arXiv:2012.03515 Table 2, LIVING-17 all-way
\newcommand{\LitMahajanFtGapPp}{0.1}  % experiments/FIGS/data/literature_facts.json :: facts[55].value :: sha256 713dba2a5e6e1dd1 :: arXiv:1805.00932 Fig. 1, ImageNet-9k target
\newcommand{\LitSophonTransferLr}{5 × 10−4}  % experiments/FIGS/data/literature_facts.json :: facts[56].value :: sha256 713dba2a5e6e1dd1 :: arXiv:2405.12972v1 App. B.5
\newcommand{\LitQgJetRadius}{0.4}  % experiments/FIGS/data/literature_facts.json :: facts[57].value :: sha256 713dba2a5e6e1dd1 :: arXiv:1810.05165 Sec. 4.1
\newcommand{\LitQgPtMin}{500}  % experiments/FIGS/data/literature_facts.json :: facts[58].value :: sha256 713dba2a5e6e1dd1 :: arXiv:1810.05165 Sec. 4.1
\newcommand{\LitQgPtMax}{550}  % experiments/FIGS/data/literature_facts.json :: facts[59].value :: sha256 713dba2a5e6e1dd1 :: arXiv:1810.05165 Sec. 4.1
\newcommand{\LitPartTopAucScratch}{0.9858}  % experiments/FIGS/data/literature_facts.json :: facts[60].value :: sha256 713dba2a5e6e1dd1 :: arXiv:2202.03772 Table 5
\newcommand{\LitPartTopAucFt}{0.9877}  % experiments/FIGS/data/literature_facts.json :: facts[61].value :: sha256 713dba2a5e6e1dd1 :: arXiv:2202.03772 Table 5
\newcommand{\LitPartQgAucScratch}{0.9121}  % experiments/FIGS/data/literature_facts.json :: facts[62].value :: sha256 713dba2a5e6e1dd1 :: arXiv:2202.03772 Table 6 (exp)
\newcommand{\LitPartQgRejScratch}{101.2}  % experiments/FIGS/data/literature_facts.json :: facts[63].value :: sha256 713dba2a5e6e1dd1 :: arXiv:2202.03772 Table 6 (exp)
\newcommand{\LitPartQgAucFt}{0.9151}  % experiments/FIGS/data/literature_facts.json :: facts[64].value :: sha256 713dba2a5e6e1dd1 :: arXiv:2202.03772 Table 6 (exp)
\newcommand{\LitPartQgRejFt}{107.9}  % experiments/FIGS/data/literature_facts.json :: facts[65].value :: sha256 713dba2a5e6e1dd1 :: arXiv:2202.03772 Table 6 (exp)
\newcommand{\LitPartJcTestMillion}{20}  % experiments/FIGS/data/literature_facts.json :: facts[66].value :: sha256 713dba2a5e6e1dd1 :: arXiv:2202.03772 Sec. 3
\newcommand{\LitPartJcTwoMAcc}{0.836}  % experiments/FIGS/data/literature_facts.json :: facts[67].value :: sha256 713dba2a5e6e1dd1 :: arXiv:2202.03772 Table 3
\newcommand{\LitPartSmallSetIters}{100 k}  % experiments/FIGS/data/literature_facts.json :: facts[68].value :: sha256 713dba2a5e6e1dd1 :: arXiv:2202.03772 Sec. 5.1 (training dataset size)
\newcommand{\LitPartSmallSetPercent}{2}  % experiments/FIGS/data/literature_facts.json :: facts[69].value :: sha256 713dba2a5e6e1dd1 :: arXiv:2202.03772 Sec. 5.1 (training dataset size)
\newcommand{\LitSophonEtaMax}{2.5}  % experiments/FIGS/data/literature_facts.json :: facts[70].value :: sha256 713dba2a5e6e1dd1 :: arXiv:2405.12972 App. A
\newcommand{\LitBHadronLifetimePs}{1.5}  % experiments/FIGS/data/literature_facts.json :: facts[71].value :: sha256 713dba2a5e6e1dd1 :: arXiv:1712.07158 Sec. 4.1
\newcommand{\LitCHadronLifetimePs}{1}  % experiments/FIGS/data/literature_facts.json :: facts[72].value :: sha256 713dba2a5e6e1dd1 :: arXiv:1712.07158 Sec. 4.1
\newcommand{\LitOmnilearnPidDropProb}{0.2}  % experiments/FIGS/data/literature_facts.json :: facts[73].value :: sha256 713dba2a5e6e1dd1 :: arXiv:2502.14652 Sec. II.A
\newcommand{\LitMahajanFtNineKLarge}{48.5}  % experiments/FIGS/data/literature_facts.json :: facts[74].value :: sha256 713dba2a5e6e1dd1 :: arXiv:1805.00932 Fig. 1 (IG-1B-17k, IN-9k, full fine-tuning)
\newcommand{\LitMahajanFtNineKSmall}{48.4}  % experiments/FIGS/data/literature_facts.json :: facts[75].value :: sha256 713dba2a5e6e1dd1 :: arXiv:1805.00932 Fig. 1 (IG-940M-1.5k, IN-9k, full fine-tuning)
\newcommand{\LitMahajanFtOneKSmall}{84.2}  % experiments/FIGS/data/literature_facts.json :: facts[76].value :: sha256 713dba2a5e6e1dd1 :: arXiv:1805.00932 Fig. 1 (IG-940M-1.5k, IN-1k, full fine-tuning)
\newcommand{\LitMahajanFtOneKLarge}{83.6}  % experiments/FIGS/data/literature_facts.json :: facts[77].value :: sha256 713dba2a5e6e1dd1 :: arXiv:1805.00932 Fig. 1 (IG-1B-17k, IN-1k, full fine-tuning)
\newcommand{\LitChenGainFifteen}{0.16}  % experiments/FIGS/data/literature_facts.json :: facts[78].value :: sha256 713dba2a5e6e1dd1 :: arXiv:1901.07012 Table VIII (15 coarse classes, CIFAR-100)
\newcommand{\LitChenGainTwenty}{0.01}  % experiments/FIGS/data/literature_facts.json :: facts[79].value :: sha256 713dba2a5e6e1dd1 :: arXiv:1901.07012 Table VIII (20 coarse classes, CIFAR-100)
\newcommand{\LitChenCoarseFifteen}{15}  % experiments/FIGS/data/literature_facts.json :: facts[80].value :: sha256 713dba2a5e6e1dd1 :: arXiv:1901.07012 Table VIII
\newcommand{\LitChenCoarseTwenty}{20}  % experiments/FIGS/data/literature_facts.json :: facts[81].value :: sha256 713dba2a5e6e1dd1 :: arXiv:1901.07012 Table VIII
\newcommand{\LitChenRandomMaxGain}{2.78}  % experiments/FIGS/data/literature_facts.json :: facts[82].value :: sha256 713dba2a5e6e1dd1 :: arXiv:1901.07012 Sec. VI-A, Table VI
\newcommand{\LitGlopartMassLambda}{5}  % experiments/FIGS/data/literature_facts.json :: facts[83].value :: sha256 713dba2a5e6e1dd1 :: CERN-THESIS-2024-281 Eq. (11.7), Sec. 11.2.1
\newcommand{\LitGlopartAddedAucBase}{0.9555}  % experiments/FIGS/data/literature_facts.json :: facts[84].value :: sha256 713dba2a5e6e1dd1 :: CERN-THESIS-2024-281 Sec. 11.4.1, Fig. 11.19 (H->WW 4q vs QCD)
\newcommand{\LitGlopartAddedAucTop}{0.9593}  % experiments/FIGS/data/literature_facts.json :: facts[85].value :: sha256 713dba2a5e6e1dd1 :: CERN-THESIS-2024-281 Sec. 11.4.1, Fig. 11.19 (H->WW 4q vs QCD)
\newcommand{\LitCmsSdZcut}{0.1}  % experiments/FIGS/data/literature_facts.json :: facts[86].value :: sha256 713dba2a5e6e1dd1 :: arXiv:2004.08262 Sec. 6.1
\newcommand{\LitCmsSdBeta}{0}  % experiments/FIGS/data/literature_facts.json :: facts[87].value :: sha256 713dba2a5e6e1dd1 :: arXiv:2004.08262 Sec. 6.1
\newcommand{\LitQgNJetsMillion}{2}  % experiments/FIGS/data/literature_facts.json :: facts[88].value :: sha256 713dba2a5e6e1dd1 :: arXiv:2202.03772 Sec. 5.2
\newcommand{\LitQgTrainMillion}{1.6}  % experiments/FIGS/data/literature_facts.json :: facts[89].value :: sha256 713dba2a5e6e1dd1 :: arXiv:2202.03772 Sec. 5.2
\newcommand{\LitQgValMillion}{0.2}  % experiments/FIGS/data/literature_facts.json :: facts[90].value :: sha256 713dba2a5e6e1dd1 :: arXiv:2202.03772 Sec. 5.2
\newcommand{\LitQgRapidityMax}{1.7}  % experiments/FIGS/data/literature_facts.json :: facts[91].value :: sha256 713dba2a5e6e1dd1 :: Zenodo 10.5281/zenodo.3164691 Zenodo record 3164691, description
\newcommand{\LitHerwigVersion}{7.1.4}  % experiments/FIGS/data/literature_facts.json :: facts[92].value :: sha256 713dba2a5e6e1dd1 :: Zenodo 10.5281/zenodo.3066475 record description
\newcommand{\LitObjectivesFewestJets}{\ensuremath{10^{3}}}  % experiments/FIGS/data/literature_facts.json :: facts[93].value :: sha256 713dba2a5e6e1dd1 :: arXiv:2606.14870 Table XVII (Micro model, Npt = 10^8: Classifier 136 vs MPM 218 at 10^3, 331 vs 539 at 10^4)
\newcommand{\LitPartTopNTrainMillion}{1.2}  % experiments/FIGS/data/literature_facts.json :: facts[94].value :: sha256 713dba2a5e6e1dd1 :: arXiv:2202.03772 Sec. 5.2
\newcommand{\LitLandscapeParticlenetEnsMean}{1412}  % experiments/FIGS/data/literature_facts.json :: facts[95].value :: sha256 713dba2a5e6e1dd1 :: arXiv:1902.09914 Table 1 (ParticleNet, ensemble mean of 1/eps_B at eps_S=0.3)
\newcommand{\LitPartFtHeadReinits}{9}  % experiments/FIGS/data/literature_facts.json :: facts[96].value :: sha256 713dba2a5e6e1dd1 :: arXiv:2202.03772 Sec. 5.2
\newcommand{\LitSophonBlocksCode}{8}  % experiments/FIGS/data/literature_facts.json :: facts[97].value :: sha256 713dba2a5e6e1dd1 :: github.com/jet-universe/sophon@9dd6dd6 networks/example_ParticleTransformer_sophon.py line 67
\newcommand{\LitGrasHerwigVersion}{2.7.1}  % experiments/FIGS/data/literature_facts.json :: facts[98].value :: sha256 713dba2a5e6e1dd1 :: arXiv:1704.03878 Sec. 2 (list of generators)
\newcommand{\DesignMassLambda}{5.0}  % experiments/MTX/k8s/job-mtx-l162_mass-s1-raunav.yaml :: --mass-lambda :: sha256 8882459fe5520795 :: identical in all ten mass-output pretraining specs
\newcommand{\FtRecipeNRuns}{264}  % experiments/FIGS/data/ft_recipes/recipes_w2b_bench_v2_v3.json :: runs (count), reported cells :: sha256 7dc8c791343e1b6f
\newcommand{\FtRecipeLr}{\ensuremath{1\times10^{-4}}}  % experiments/FIGS/data/ft_recipes/recipes_w2b_bench_v2_v3.json :: runs[*].lr, pretrained starts :: sha256 7dc8c791343e1b6f
\newcommand{\FtRecipeHeadMult}{50}  % experiments/FIGS/data/ft_recipes/recipes_w2b_bench_v2_v3.json :: runs[*].head_lr_mult, pretrained starts :: sha256 7dc8c791343e1b6f
\newcommand{\FtRecipeLrHead}{\ensuremath{5\times10^{-3}}}  % experiments/FIGS/data/ft_recipes/recipes_w2b_bench_v2_v3.json :: lr x head_lr_mult :: sha256 7dc8c791343e1b6f
\newcommand{\FtRecipeWeightDecay}{0.01}  % experiments/FIGS/data/ft_recipes/recipes_w2b_bench_v2_v3.json :: runs[*].weight_decay :: sha256 7dc8c791343e1b6f
\newcommand{\FtRecipeBatch}{512}  % experiments/FIGS/data/ft_recipes/recipes_w2b_bench_v2_v3.json :: runs[*].batch_size :: sha256 7dc8c791343e1b6f
\newcommand{\FtRecipeEpochsEThree}{50}  % experiments/FIGS/data/ft_recipes/recipes_w2b_bench_v2_v3.json :: runs[leg 1,2; n_train 1000].epochs :: sha256 7dc8c791343e1b6f
\newcommand{\FtRecipeStepsEThree}{19}  % experiments/FT/k8s/job-ft-legs-w2-raunav.yaml :: job commands: samples per epoch at n_train 1000 // batch size :: sha256 f0962eddb451b7b0 :: optimizer steps per epoch, as weaver floors them
\newcommand{\FtRecipePassesEThree}{10}  % experiments/FT/k8s/job-ft-legs-w2-raunav.yaml :: job commands: samples per epoch at n_train 1000 // n_train :: sha256 f0962eddb451b7b0 :: passes over the training subset per epoch
\newcommand{\FtRecipeEpochsEFour}{50}  % experiments/FIGS/data/ft_recipes/recipes_w2b_bench_v2_v3.json :: runs[leg 1,2; n_train 10000].epochs :: sha256 7dc8c791343e1b6f
\newcommand{\FtRecipeStepsEFour}{19}  % experiments/FT/k8s/job-ft-legs-w2-raunav.yaml :: job commands: samples per epoch at n_train 10000 // batch size :: sha256 f0962eddb451b7b0 :: optimizer steps per epoch, as weaver floors them
\newcommand{\FtRecipePassesEFour}{one}  % experiments/FT/k8s/job-ft-legs-w2-raunav.yaml :: job commands: samples per epoch at n_train 10000 // n_train :: sha256 f0962eddb451b7b0 :: passes over the training subset per epoch
\newcommand{\FtRecipeEpochsEFive}{30}  % experiments/FIGS/data/ft_recipes/recipes_w2b_bench_v2_v3.json :: runs[leg 1,2; n_train 100000].epochs :: sha256 7dc8c791343e1b6f
\newcommand{\FtRecipeStepsEFive}{195}  % experiments/FT/k8s/job-ft-legs-w2-raunav.yaml :: job commands: samples per epoch at n_train 100000 // batch size :: sha256 f0962eddb451b7b0 :: optimizer steps per epoch, as weaver floors them
\newcommand{\FtRecipePassesEFive}{one}  % experiments/FT/k8s/job-ft-legs-w2-raunav.yaml :: job commands: samples per epoch at n_train 100000 // n_train :: sha256 f0962eddb451b7b0 :: passes over the training subset per epoch
\newcommand{\FtRecipeEpochsESix}{10}  % experiments/FIGS/data/ft_recipes/recipes_w2b_bench_v2_v3.json :: runs[leg 1,2; n_train 1000000].epochs :: sha256 7dc8c791343e1b6f
\newcommand{\FtRecipeStepsESix}{1{,}953}  % experiments/FT/k8s/job-ft-legs-w2-raunav.yaml :: job commands: samples per epoch at n_train 1000000 // batch size :: sha256 f0962eddb451b7b0 :: optimizer steps per epoch, as weaver floors them
\newcommand{\FtRecipePassesESix}{one}  % experiments/FT/k8s/job-ft-legs-w2-raunav.yaml :: job commands: samples per epoch at n_train 1000000 // n_train :: sha256 f0962eddb451b7b0 :: passes over the training subset per epoch
\newcommand{\FtRecipeValJets}{19{,}968}  % experiments/FT/k8s/job-ft-legs-w2-raunav.yaml :: --samples-per-epoch-val 20000 // batch x batch :: sha256 f0962eddb451b7b0 :: validation jets per epoch, floored to whole batches
\newcommand{\RecoveryAccOneeighteightByOneeighteight}{0.4123\,\ensuremath{\pm}\,0.0011}  % experiments/FIGS/data/label_recovery_curve_v1err/analysis/label_recovery_curve_summary.json :: table[rung=L188,level=188,probe=linear,n_train=max].accuracy :: sha256 1250d1f00537e0bd :: balanced accuracy, mean +- SD over 5 runs
\newcommand{\RecoveryAccOneeighteightByOnesixtwo}{0.4090\,\ensuremath{\pm}\,0.0012}  % experiments/FIGS/data/label_recovery_curve_v1err/analysis/label_recovery_curve_summary.json :: table[rung=L188,level=162,probe=linear,n_train=max].accuracy :: sha256 1250d1f00537e0bd :: balanced accuracy, mean +- SD over 5 runs
\newcommand{\RecoveryAccOneeighteightByFourthree}{0.356\,\ensuremath{\pm}\,0.008}  % experiments/FIGS/data/label_recovery_curve_v1err/analysis/label_recovery_curve_summary.json :: table[rung=L188,level=43,probe=linear,n_train=max].accuracy :: sha256 1250d1f00537e0bd :: balanced accuracy, mean +- SD over 5 runs
\newcommand{\RecoveryAccOneeighteightByOneseven}{0.306\,\ensuremath{\pm}\,0.008}  % experiments/FIGS/data/label_recovery_curve_v1err/analysis/label_recovery_curve_summary.json :: table[rung=L188,level=17,probe=linear,n_train=max].accuracy :: sha256 1250d1f00537e0bd :: balanced accuracy, mean +- SD over 5 runs
\newcommand{\RecoveryAccOnesixtwoByOneeighteight}{0.4623\,\ensuremath{\pm}\,0.0006}  % experiments/FIGS/data/label_recovery_curve_v1err/analysis/label_recovery_curve_summary.json :: table[rung=L162,level=188,probe=linear,n_train=max].accuracy :: sha256 1250d1f00537e0bd :: balanced accuracy, mean +- SD over 5 runs
\newcommand{\RecoveryAccOnesixtwoByOnesixtwo}{0.4622\,\ensuremath{\pm}\,0.0007}  % experiments/FIGS/data/label_recovery_curve_v1err/analysis/label_recovery_curve_summary.json :: table[rung=L162,level=162,probe=linear,n_train=max].accuracy :: sha256 1250d1f00537e0bd :: balanced accuracy, mean +- SD over 5 runs
\newcommand{\RecoveryAccOnesixtwoByFourthree}{0.402\,\ensuremath{\pm}\,0.009}  % experiments/FIGS/data/label_recovery_curve_v1err/analysis/label_recovery_curve_summary.json :: table[rung=L162,level=43,probe=linear,n_train=max].accuracy :: sha256 1250d1f00537e0bd :: balanced accuracy, mean +- SD over 5 runs
\newcommand{\RecoveryAccOnesixtwoByOneseven}{0.348\,\ensuremath{\pm}\,0.010}  % experiments/FIGS/data/label_recovery_curve_v1err/analysis/label_recovery_curve_summary.json :: table[rung=L162,level=17,probe=linear,n_train=max].accuracy :: sha256 1250d1f00537e0bd :: balanced accuracy, mean +- SD over 5 runs
\newcommand{\RecoveryAccSixfourByOneeighteight}{0.5979\,\ensuremath{\pm}\,0.0008}  % experiments/FIGS/data/label_recovery_curve_v1err/analysis/label_recovery_curve_summary.json :: table[rung=R63_Q1,level=188,probe=linear,n_train=max].accuracy :: sha256 1250d1f00537e0bd :: balanced accuracy, mean +- SD over 5 runs
\newcommand{\RecoveryAccSixfourByOnesixtwo}{0.5973\,\ensuremath{\pm}\,0.0013}  % experiments/FIGS/data/label_recovery_curve_v1err/analysis/label_recovery_curve_summary.json :: table[rung=R63_Q1,level=162,probe=linear,n_train=max].accuracy :: sha256 1250d1f00537e0bd :: balanced accuracy, mean +- SD over 5 runs
\newcommand{\RecoveryAccSixfourByFourthree}{0.5881\,\ensuremath{\pm}\,0.0007}  % experiments/FIGS/data/label_recovery_curve_v1err/analysis/label_recovery_curve_summary.json :: table[rung=R63_Q1,level=43,probe=linear,n_train=max].accuracy :: sha256 1250d1f00537e0bd :: balanced accuracy, mean +- SD over 5 runs
\newcommand{\RecoveryAccSixfourByOneseven}{0.531\,\ensuremath{\pm}\,0.006}  % experiments/FIGS/data/label_recovery_curve_v1err/analysis/label_recovery_curve_summary.json :: table[rung=R63_Q1,level=17,probe=linear,n_train=max].accuracy :: sha256 1250d1f00537e0bd :: balanced accuracy, mean +- SD over 5 runs
\newcommand{\RecoveryAccFourthreeByOneeighteight}{0.6770\,\ensuremath{\pm}\,0.0006}  % experiments/FIGS/data/label_recovery_curve_v1err/analysis/label_recovery_curve_summary.json :: table[rung=R42_Q1,level=188,probe=linear,n_train=max].accuracy :: sha256 1250d1f00537e0bd :: balanced accuracy, mean +- SD over 5 runs
\newcommand{\RecoveryAccFourthreeByOnesixtwo}{0.6766\,\ensuremath{\pm}\,0.0008}  % experiments/FIGS/data/label_recovery_curve_v1err/analysis/label_recovery_curve_summary.json :: table[rung=R42_Q1,level=162,probe=linear,n_train=max].accuracy :: sha256 1250d1f00537e0bd :: balanced accuracy, mean +- SD over 5 runs
\newcommand{\RecoveryAccFourthreeByFourthree}{0.6751\,\ensuremath{\pm}\,0.0011}  % experiments/FIGS/data/label_recovery_curve_v1err/analysis/label_recovery_curve_summary.json :: table[rung=R42_Q1,level=43,probe=linear,n_train=max].accuracy :: sha256 1250d1f00537e0bd :: balanced accuracy, mean +- SD over 5 runs
\newcommand{\RecoveryAccFourthreeByOneseven}{0.633\,\ensuremath{\pm}\,0.006}  % experiments/FIGS/data/label_recovery_curve_v1err/analysis/label_recovery_curve_summary.json :: table[rung=R42_Q1,level=17,probe=linear,n_train=max].accuracy :: sha256 1250d1f00537e0bd :: balanced accuracy, mean +- SD over 5 runs
\newcommand{\RecoveryAccThreezeroByOneeighteight}{0.7001\,\ensuremath{\pm}\,0.0005}  % experiments/FIGS/data/label_recovery_curve_v1err/analysis/label_recovery_curve_summary.json :: table[rung=R29_Q1,level=188,probe=linear,n_train=max].accuracy :: sha256 1250d1f00537e0bd :: balanced accuracy, mean +- SD over 5 runs
\newcommand{\RecoveryAccThreezeroByOnesixtwo}{0.7004\,\ensuremath{\pm}\,0.0007}  % experiments/FIGS/data/label_recovery_curve_v1err/analysis/label_recovery_curve_summary.json :: table[rung=R29_Q1,level=162,probe=linear,n_train=max].accuracy :: sha256 1250d1f00537e0bd :: balanced accuracy, mean +- SD over 5 runs
\newcommand{\RecoveryAccThreezeroByFourthree}{0.6998\,\ensuremath{\pm}\,0.0008}  % experiments/FIGS/data/label_recovery_curve_v1err/analysis/label_recovery_curve_summary.json :: table[rung=R29_Q1,level=43,probe=linear,n_train=max].accuracy :: sha256 1250d1f00537e0bd :: balanced accuracy, mean +- SD over 5 runs
\newcommand{\RecoveryAccThreezeroByOneseven}{0.6778\,\ensuremath{\pm}\,0.0030}  % experiments/FIGS/data/label_recovery_curve_v1err/analysis/label_recovery_curve_summary.json :: table[rung=R29_Q1,level=17,probe=linear,n_train=max].accuracy :: sha256 1250d1f00537e0bd :: balanced accuracy, mean +- SD over 5 runs
\newcommand{\RecoveryAccOnesevenByOneeighteight}{0.7675\,\ensuremath{\pm}\,0.0010}  % experiments/FIGS/data/label_recovery_curve_v1err/analysis/label_recovery_curve_summary.json :: table[rung=R16_Q1,level=188,probe=linear,n_train=max].accuracy :: sha256 1250d1f00537e0bd :: balanced accuracy, mean +- SD over 5 runs
\newcommand{\RecoveryAccOnesevenByOnesixtwo}{0.7681\,\ensuremath{\pm}\,0.0009}  % experiments/FIGS/data/label_recovery_curve_v1err/analysis/label_recovery_curve_summary.json :: table[rung=R16_Q1,level=162,probe=linear,n_train=max].accuracy :: sha256 1250d1f00537e0bd :: balanced accuracy, mean +- SD over 5 runs
\newcommand{\RecoveryAccOnesevenByFourthree}{0.7680\,\ensuremath{\pm}\,0.0010}  % experiments/FIGS/data/label_recovery_curve_v1err/analysis/label_recovery_curve_summary.json :: table[rung=R16_Q1,level=43,probe=linear,n_train=max].accuracy :: sha256 1250d1f00537e0bd :: balanced accuracy, mean +- SD over 5 runs
\newcommand{\RecoveryAccOnesevenByOneseven}{0.7681\,\ensuremath{\pm}\,0.0003}  % experiments/FIGS/data/label_recovery_curve_v1err/analysis/label_recovery_curve_summary.json :: table[rung=R16_Q1,level=17,probe=linear,n_train=max].accuracy :: sha256 1250d1f00537e0bd :: balanced accuracy, mean +- SD over 5 runs
\newcommand{\RecoveryAccFourByOneeighteight}{0.7526\,\ensuremath{\pm}\,0.0002}  % experiments/FIGS/data/label_recovery_curve_v1err/analysis/label_recovery_curve_summary.json :: table[rung=R3_VIS,level=188,probe=linear,n_train=max].accuracy :: sha256 1250d1f00537e0bd :: balanced accuracy, mean +- SD over 5 runs
\newcommand{\RecoveryAccFourByOnesixtwo}{0.7538\,\ensuremath{\pm}\,0.0017}  % experiments/FIGS/data/label_recovery_curve_v1err/analysis/label_recovery_curve_summary.json :: table[rung=R3_VIS,level=162,probe=linear,n_train=max].accuracy :: sha256 1250d1f00537e0bd :: balanced accuracy, mean +- SD over 5 runs
\newcommand{\RecoveryAccFourByFourthree}{0.7557\,\ensuremath{\pm}\,0.0015}  % experiments/FIGS/data/label_recovery_curve_v1err/analysis/label_recovery_curve_summary.json :: table[rung=R3_VIS,level=43,probe=linear,n_train=max].accuracy :: sha256 1250d1f00537e0bd :: balanced accuracy, mean +- SD over 5 runs
\newcommand{\RecoveryAccFourByOneseven}{0.7586\,\ensuremath{\pm}\,0.0008}  % experiments/FIGS/data/label_recovery_curve_v1err/analysis/label_recovery_curve_summary.json :: table[rung=R3_VIS,level=17,probe=linear,n_train=max].accuracy :: sha256 1250d1f00537e0bd :: balanced accuracy, mean +- SD over 5 runs
\newcommand{\RecoveryAccTwoByOneeighteight}{0.8688\,\ensuremath{\pm}\,0.0005}  % experiments/FIGS/data/label_recovery_curve_v1err/analysis/label_recovery_curve_summary.json :: table[rung=R1_Q1,level=188,probe=linear,n_train=max].accuracy :: sha256 1250d1f00537e0bd :: balanced accuracy, mean +- SD over 5 runs
\newcommand{\RecoveryAccTwoByOnesixtwo}{0.8703\,\ensuremath{\pm}\,0.0008}  % experiments/FIGS/data/label_recovery_curve_v1err/analysis/label_recovery_curve_summary.json :: table[rung=R1_Q1,level=162,probe=linear,n_train=max].accuracy :: sha256 1250d1f00537e0bd :: balanced accuracy, mean +- SD over 5 runs
\newcommand{\RecoveryAccTwoByFourthree}{0.8709\,\ensuremath{\pm}\,0.0011}  % experiments/FIGS/data/label_recovery_curve_v1err/analysis/label_recovery_curve_summary.json :: table[rung=R1_Q1,level=43,probe=linear,n_train=max].accuracy :: sha256 1250d1f00537e0bd :: balanced accuracy, mean +- SD over 5 runs
\newcommand{\RecoveryAccTwoByOneseven}{0.8713\,\ensuremath{\pm}\,0.0002}  % experiments/FIGS/data/label_recovery_curve_v1err/analysis/label_recovery_curve_summary.json :: table[rung=R1_Q1,level=17,probe=linear,n_train=max].accuracy :: sha256 1250d1f00537e0bd :: balanced accuracy, mean +- SD over 5 runs
\newcommand{\RecoveryOneeighteightOnesixtwoMaxAbs}{0.0033}  % experiments/FIGS/data/label_recovery_curve_v1err/analysis/label_recovery_curve_summary.json :: table[rung=*,level=188/162,probe=linear,n_train=max].accuracy :: sha256 1250d1f00537e0bd :: max |difference of the run means| over the tree
\newcommand{\RecoveryPairedOnesevenLevelOnesevenMinusOneeighteight}{+0.0006~[\ensuremath{-}0.0004, +0.0015]}  % experiments/FIGS/data/label_recovery_curve_v1err/analysis/label_recovery_curve_summary.json :: table[rung=R16_Q1,level=17 minus 188,probe=linear,n_train=max].accuracy :: sha256 1250d1f00537e0bd :: paired by run, mean difference [95% Student-t interval]
\newcommand{\RecoveryPairedOnesevenLevelOnesevenMinusOnesixtwo}{+0.00004~[\ensuremath{-}0.00123, +0.00131]}  % experiments/FIGS/data/label_recovery_curve_v1err/analysis/label_recovery_curve_summary.json :: table[rung=R16_Q1,level=17 minus 162,probe=linear,n_train=max].accuracy :: sha256 1250d1f00537e0bd :: paired by run, mean difference [95% Student-t interval]
\newcommand{\RecoveryPairedFourLevelOnesevenMinusOneeighteight}{+0.0060~[+0.0051, +0.0069]}  % experiments/FIGS/data/label_recovery_curve_v1err/analysis/label_recovery_curve_summary.json :: table[rung=R3_VIS,level=17 minus 188,probe=linear,n_train=max].accuracy :: sha256 1250d1f00537e0bd :: paired by run, mean difference [95% Student-t interval]
\newcommand{\RecoveryPairedFourLevelOnesevenMinusOnesixtwo}{+0.0048~[+0.0028, +0.0068]}  % experiments/FIGS/data/label_recovery_curve_v1err/analysis/label_recovery_curve_summary.json :: table[rung=R3_VIS,level=17 minus 162,probe=linear,n_train=max].accuracy :: sha256 1250d1f00537e0bd :: paired by run, mean difference [95% Student-t interval]
\newcommand{\RecoveryPairedTwoLevelOnesevenMinusOneeighteight}{+0.0025~[+0.0019, +0.0032]}  % experiments/FIGS/data/label_recovery_curve_v1err/analysis/label_recovery_curve_summary.json :: table[rung=R1_Q1,level=17 minus 188,probe=linear,n_train=max].accuracy :: sha256 1250d1f00537e0bd :: paired by run, mean difference [95% Student-t interval]
\newcommand{\RecoveryPairedTwoLevelOnesevenMinusOnesixtwo}{+0.00101~[+0.00003, +0.00198]}  % experiments/FIGS/data/label_recovery_curve_v1err/analysis/label_recovery_curve_summary.json :: table[rung=R1_Q1,level=17 minus 162,probe=linear,n_train=max].accuracy :: sha256 1250d1f00537e0bd :: paired by run, mean difference [95% Student-t interval]
\newcommand{\RecoveryMlpGainOneseven}{+0.0099~[+0.0078, +0.0120]}  % experiments/FIGS/data/label_recovery_curve_v1err/analysis/label_recovery_curve_summary.json :: mlp_minus_linear_at_largest[rung=L188,level=17].per_run :: sha256 1250d1f00537e0bd :: nonlinear minus linear probe, 188 native labels from the 17-class models
\newcommand{\RecoveryLastGainOneseven}{0.0125\,\ensuremath{\pm}\,0.0021}  % experiments/FIGS/data/label_recovery_curve_v1err/analysis/label_recovery_curve_summary.json :: data_gain_last_step[rung=L188,level=17] :: sha256 1250d1f00537e0bd :: balanced accuracy gained from 443000 to 1400000 training jets
\newcommand{\RecoveryLastGainOneeighteight}{0.0147\,\ensuremath{\pm}\,0.0014}  % experiments/FIGS/data/label_recovery_curve_v1err/analysis/label_recovery_curve_summary.json :: data_gain_last_step[rung=L188,level=188] :: sha256 1250d1f00537e0bd :: balanced accuracy gained from 443000 to 1400000 training jets
\newcommand{\RecoveryNTrain}{1{,}400{,}000}  % experiments/FIGS/data/label_recovery_curve_v1err/analysis/label_recovery_curve_summary.json :: sizes (max) :: sha256 1250d1f00537e0bd :: jets the linear probe is fit on at the end of the curve
\newcommand{\RecoveryNTrainMin}{14{,}000}  % experiments/FIGS/data/label_recovery_curve_v1err/analysis/label_recovery_curve_summary.json :: sizes (min) :: sha256 1250d1f00537e0bd
\newcommand{\RecoveryNTrainPrev}{443{,}000}  % experiments/FIGS/data/label_recovery_curve_v1err/analysis/label_recovery_curve_summary.json :: sizes (second largest) :: sha256 1250d1f00537e0bd
\newcommand{\RecoveryNTest}{400{,}000}  % experiments/FIGS/data/label_recovery_curve_v1err/l188-s1/label_recovery_curve.json :: n_test :: sha256 65651e4c06c7766f :: jets the probes are scored on, recorded in all 20 per-run files
\newcommand{\RecoveryMlpHidden}{512}  % experiments/FIGS/data/label_recovery_curve_v1err/l188-s1/label_recovery_curve.json :: mlp.hidden :: sha256 65651e4c06c7766f :: hidden units of the nonlinear probe
\newcommand{\RandNDraws}{three}  % experiments/FIGS/data/probe_ladder_randcontrol_mlp2/analysis/c4_random_control.json :: table[*].draw :: sha256 c675e609ff4fe3c9 :: random partitions drawn
\newcommand{\RandNDrawsCap}{Three}  % experiments/FIGS/data/probe_ladder_randcontrol_mlp2/analysis/c4_random_control.json :: table[*].draw :: sha256 c675e609ff4fe3c9 :: random partitions drawn, capitalised to open a sentence
\newcommand{\RandOmaBvcFourprongDrawOneLinear}{\ensuremath{1.05\times10^{-2}}}  % experiments/FIGS/data/probe_ladder_randcontrol_mlp2/analysis/c4_random_control.json :: table[0].control_log1m_auc :: sha256 c675e609ff4fe3c9 :: 1-AUC = exp of it, one run
\newcommand{\RandOmaBvcFourprongDrawTwoLinear}{\ensuremath{1.11\times10^{-2}}}  % experiments/FIGS/data/probe_ladder_randcontrol_mlp2/analysis/c4_random_control.json :: table[4].control_log1m_auc :: sha256 c675e609ff4fe3c9 :: 1-AUC = exp of it, one run
\newcommand{\RandOmaBvcFourprongDrawThreeLinear}{\ensuremath{9.98\times10^{-3}}}  % experiments/FIGS/data/probe_ladder_randcontrol_mlp2/analysis/c4_random_control.json :: table[8].control_log1m_auc :: sha256 c675e609ff4fe3c9 :: 1-AUC = exp of it, one run
\newcommand{\RandOmaMeanBvcFourprongLinear}{\ensuremath{(1.05 \pm 0.06)\times10^{-2}}}  % experiments/FIGS/data/probe_ladder_randcontrol_mlp2/analysis/c4_random_control.json :: table[task=bvc_4prong,probe=linear].control_log1m_auc :: sha256 c675e609ff4fe3c9 :: 1-AUC, mean +- SD over 3 draws
\newcommand{\RandOmaBvcFourprongDrawOneMlp}{\ensuremath{1.05\times10^{-2}}}  % experiments/FIGS/data/probe_ladder_randcontrol_mlp2/analysis/c4_random_control.json :: table[1].control_log1m_auc :: sha256 c675e609ff4fe3c9 :: 1-AUC = exp of it, one run
\newcommand{\RandOmaBvcFourprongDrawTwoMlp}{\ensuremath{1.04\times10^{-2}}}  % experiments/FIGS/data/probe_ladder_randcontrol_mlp2/analysis/c4_random_control.json :: table[5].control_log1m_auc :: sha256 c675e609ff4fe3c9 :: 1-AUC = exp of it, one run
\newcommand{\RandOmaBvcFourprongDrawThreeMlp}{\ensuremath{9.28\times10^{-3}}}  % experiments/FIGS/data/probe_ladder_randcontrol_mlp2/analysis/c4_random_control.json :: table[9].control_log1m_auc :: sha256 c675e609ff4fe3c9 :: 1-AUC = exp of it, one run
\newcommand{\RandOmaMeanBvcFourprongMlp}{\ensuremath{(1.01 \pm 0.07)\times10^{-2}}}  % experiments/FIGS/data/probe_ladder_randcontrol_mlp2/analysis/c4_random_control.json :: table[task=bvc_4prong,probe=mlp].control_log1m_auc :: sha256 c675e609ff4fe3c9 :: 1-AUC, mean +- SD over 3 draws
\newcommand{\RandOmaVisibleContentDrawOneLinear}{\ensuremath{1.25\times10^{-3}}}  % experiments/FIGS/data/probe_ladder_randcontrol_mlp2/analysis/c4_random_control.json :: table[2].control_log1m_auc :: sha256 c675e609ff4fe3c9 :: 1-AUC = exp of it, one run
\newcommand{\RandOmaVisibleContentDrawTwoLinear}{\ensuremath{9.89\times10^{-4}}}  % experiments/FIGS/data/probe_ladder_randcontrol_mlp2/analysis/c4_random_control.json :: table[6].control_log1m_auc :: sha256 c675e609ff4fe3c9 :: 1-AUC = exp of it, one run
\newcommand{\RandOmaVisibleContentDrawThreeLinear}{\ensuremath{1.08\times10^{-3}}}  % experiments/FIGS/data/probe_ladder_randcontrol_mlp2/analysis/c4_random_control.json :: table[10].control_log1m_auc :: sha256 c675e609ff4fe3c9 :: 1-AUC = exp of it, one run
\newcommand{\RandOmaMeanVisibleContentLinear}{\ensuremath{(1.11 \pm 0.13)\times10^{-3}}}  % experiments/FIGS/data/probe_ladder_randcontrol_mlp2/analysis/c4_random_control.json :: table[task=visible_content,probe=linear].control_log1m_auc :: sha256 c675e609ff4fe3c9 :: 1-AUC, mean +- SD over 3 draws
\newcommand{\RandOmaVisibleContentDrawOneMlp}{\ensuremath{1.21\times10^{-3}}}  % experiments/FIGS/data/probe_ladder_randcontrol_mlp2/analysis/c4_random_control.json :: table[3].control_log1m_auc :: sha256 c675e609ff4fe3c9 :: 1-AUC = exp of it, one run
\newcommand{\RandOmaVisibleContentDrawTwoMlp}{\ensuremath{1.07\times10^{-3}}}  % experiments/FIGS/data/probe_ladder_randcontrol_mlp2/analysis/c4_random_control.json :: table[7].control_log1m_auc :: sha256 c675e609ff4fe3c9 :: 1-AUC = exp of it, one run
\newcommand{\RandOmaVisibleContentDrawThreeMlp}{\ensuremath{1.13\times10^{-3}}}  % experiments/FIGS/data/probe_ladder_randcontrol_mlp2/analysis/c4_random_control.json :: table[11].control_log1m_auc :: sha256 c675e609ff4fe3c9 :: 1-AUC = exp of it, one run
\newcommand{\RandOmaMeanVisibleContentMlp}{\ensuremath{(1.14 \pm 0.07)\times10^{-3}}}  % experiments/FIGS/data/probe_ladder_randcontrol_mlp2/analysis/c4_random_control.json :: table[task=visible_content,probe=mlp].control_log1m_auc :: sha256 c675e609ff4fe3c9 :: 1-AUC, mean +- SD over 3 draws
\newcommand{\RandCostFactorBvcFourprongLinear}{1.49}  % experiments/FIGS/data/probe_ladder_randcontrol_mlp2/analysis/c4_random_control.json :: grouping_cost_post_hoc.tasks.bvc_4prong.linear.factor_control :: sha256 c675e609ff4fe3c9 :: POST HOC: random grouping vs finer models, in 1-AUC
\newcommand{\RandSemCostFactorBvcFourprongLinear}{4.42}  % experiments/FIGS/data/probe_ladder_randcontrol_mlp2/analysis/c4_random_control.json :: grouping_cost_post_hoc.tasks.bvc_4prong.linear.factor_semantic :: sha256 c675e609ff4fe3c9 :: POST HOC: semantic 17-class vs finer models
\newcommand{\RandCostFactorBvcFourprongMlp}{1.41}  % experiments/FIGS/data/probe_ladder_randcontrol_mlp2/analysis/c4_random_control.json :: grouping_cost_post_hoc.tasks.bvc_4prong.mlp.factor_control :: sha256 c675e609ff4fe3c9 :: POST HOC: random grouping vs finer models, in 1-AUC
\newcommand{\RandSemCostFactorBvcFourprongMlp}{3.26}  % experiments/FIGS/data/probe_ladder_randcontrol_mlp2/analysis/c4_random_control.json :: grouping_cost_post_hoc.tasks.bvc_4prong.mlp.factor_semantic :: sha256 c675e609ff4fe3c9 :: POST HOC: semantic 17-class vs finer models
\newcommand{\RandCostFactorVisibleContentLinear}{1.35}  % experiments/FIGS/data/probe_ladder_randcontrol_mlp2/analysis/c4_random_control.json :: grouping_cost_post_hoc.tasks.visible_content.linear.factor_control :: sha256 c675e609ff4fe3c9 :: POST HOC: random grouping vs finer models, in 1-AUC
\newcommand{\RandSemCostFactorVisibleContentLinear}{2.82}  % experiments/FIGS/data/probe_ladder_randcontrol_mlp2/analysis/c4_random_control.json :: grouping_cost_post_hoc.tasks.visible_content.linear.factor_semantic :: sha256 c675e609ff4fe3c9 :: POST HOC: semantic 17-class vs finer models
\newcommand{\RandCostFactorVisibleContentMlp}{1.33}  % experiments/FIGS/data/probe_ladder_randcontrol_mlp2/analysis/c4_random_control.json :: grouping_cost_post_hoc.tasks.visible_content.mlp.factor_control :: sha256 c675e609ff4fe3c9 :: POST HOC: random grouping vs finer models, in 1-AUC
\newcommand{\RandSemCostFactorVisibleContentMlp}{2.37}  % experiments/FIGS/data/probe_ladder_randcontrol_mlp2/analysis/c4_random_control.json :: grouping_cost_post_hoc.tasks.visible_content.mlp.factor_semantic :: sha256 c675e609ff4fe3c9 :: POST HOC: semantic 17-class vs finer models
\newcommand{\MassResNTest}{398{,}625}  % experiments/FIGS/data/mass_resolution/s1.json :: centering_detail.split[2] :: sha256 88a57bbad270d51c :: test jets sigma_eff is computed on, the same in every per-seed file and arm
\newcommand{\MassResNClasses}{184}  % experiments/FIGS/data/mass_resolution/analysis_holm/s7_mass_resolution.json :: provenance.n_classes_used :: sha256 e03cde45ece9f385 :: native classes with enough training jets to centre
\newcommand{\MassResSigmaEffRidgeOneeighteight}{0.3393\,\ensuremath{\pm}\,0.0019}  % experiments/FIGS/data/mass_resolution/analysis_holm/s7_mass_resolution.json :: table[cell=188,probe=ridge].sigma_eff :: sha256 e03cde45ece9f385 :: mean +- SD over 5 seeds
\newcommand{\MassResSigmaEffRidgeOnesixtwo}{0.3302\,\ensuremath{\pm}\,0.0029}  % experiments/FIGS/data/mass_resolution/analysis_holm/s7_mass_resolution.json :: table[cell=162,probe=ridge].sigma_eff :: sha256 e03cde45ece9f385 :: mean +- SD over 5 seeds
\newcommand{\MassResSigmaEffRidgeFourthree}{0.3228\,\ensuremath{\pm}\,0.0019}  % experiments/FIGS/data/mass_resolution/analysis_holm/s7_mass_resolution.json :: table[cell=43,probe=ridge].sigma_eff :: sha256 e03cde45ece9f385 :: mean +- SD over 5 seeds
\newcommand{\MassResSigmaEffRidgeOneseven}{0.300\,\ensuremath{\pm}\,0.010}  % experiments/FIGS/data/mass_resolution/analysis_holm/s7_mass_resolution.json :: table[cell=17,probe=ridge].sigma_eff :: sha256 e03cde45ece9f385 :: mean +- SD over 5 seeds
\newcommand{\MassResSigmaEffRidgeOnesixtwoMass}{0.2238\,\ensuremath{\pm}\,0.0019}  % experiments/FIGS/data/mass_resolution/analysis_holm/s7_mass_resolution.json :: table[cell=162+mass,probe=ridge].sigma_eff :: sha256 e03cde45ece9f385 :: mean +- SD over 5 seeds
\newcommand{\MassResSigmaEffRidgeOnesevenMass}{0.2170\,\ensuremath{\pm}\,0.0034}  % experiments/FIGS/data/mass_resolution/analysis_holm/s7_mass_resolution.json :: table[cell=17+mass,probe=ridge].sigma_eff :: sha256 e03cde45ece9f385 :: mean +- SD over 5 seeds
\newcommand{\MassResGainRidgeOnesixtwo}{\ensuremath{-}0.1063}  % experiments/FIGS/data/mass_resolution/analysis_holm/s7_mass_resolution.json :: table[cell=162+mass/162,probe=ridge].sigma_eff :: sha256 e03cde45ece9f385 :: sigma_eff with the mass output minus without, difference of the seed means; negative is better
\newcommand{\MassResGainRidgeOneseven}{\ensuremath{-}0.0831}  % experiments/FIGS/data/mass_resolution/analysis_holm/s7_mass_resolution.json :: table[cell=17+mass/17,probe=ridge].sigma_eff :: sha256 e03cde45ece9f385 :: sigma_eff with the mass output minus without, difference of the seed means; negative is better
\newcommand{\MassResDidRidge}{\ensuremath{-}0.0233}  % experiments/FIGS/data/mass_resolution/analysis_holm/s7_mass_resolution.json :: table[cell=162+mass/162/17+mass/17,probe=ridge].sigma_eff :: sha256 e03cde45ece9f385 :: gain at 162 classes minus gain at 17
\newcommand{\MassResSigmaEffMlpOneeighteight}{0.2463\,\ensuremath{\pm}\,0.0022}  % experiments/FIGS/data/mass_resolution/analysis_holm/s7_mass_resolution.json :: table[cell=188,probe=mlp].sigma_eff :: sha256 e03cde45ece9f385 :: mean +- SD over 5 seeds
\newcommand{\MassResSigmaEffMlpOnesixtwo}{0.2438\,\ensuremath{\pm}\,0.0018}  % experiments/FIGS/data/mass_resolution/analysis_holm/s7_mass_resolution.json :: table[cell=162,probe=mlp].sigma_eff :: sha256 e03cde45ece9f385 :: mean +- SD over 5 seeds
\newcommand{\MassResSigmaEffMlpFourthree}{0.2407\,\ensuremath{\pm}\,0.0009}  % experiments/FIGS/data/mass_resolution/analysis_holm/s7_mass_resolution.json :: table[cell=43,probe=mlp].sigma_eff :: sha256 e03cde45ece9f385 :: mean +- SD over 5 seeds
\newcommand{\MassResSigmaEffMlpOneseven}{0.2358\,\ensuremath{\pm}\,0.0025}  % experiments/FIGS/data/mass_resolution/analysis_holm/s7_mass_resolution.json :: table[cell=17,probe=mlp].sigma_eff :: sha256 e03cde45ece9f385 :: mean +- SD over 5 seeds
\newcommand{\MassResSigmaEffMlpOnesixtwoMass}{0.1943\,\ensuremath{\pm}\,0.0001}  % experiments/FIGS/data/mass_resolution/analysis_holm/s7_mass_resolution.json :: table[cell=162+mass,probe=mlp].sigma_eff :: sha256 e03cde45ece9f385 :: mean +- SD over 5 seeds
\newcommand{\MassResSigmaEffMlpOnesevenMass}{0.1949\,\ensuremath{\pm}\,0.0006}  % experiments/FIGS/data/mass_resolution/analysis_holm/s7_mass_resolution.json :: table[cell=17+mass,probe=mlp].sigma_eff :: sha256 e03cde45ece9f385 :: mean +- SD over 5 seeds
\newcommand{\MassResGainMlpOnesixtwo}{\ensuremath{-}0.0496}  % experiments/FIGS/data/mass_resolution/analysis_holm/s7_mass_resolution.json :: table[cell=162+mass/162,probe=mlp].sigma_eff :: sha256 e03cde45ece9f385 :: sigma_eff with the mass output minus without, difference of the seed means; negative is better
\newcommand{\MassResGainMlpOneseven}{\ensuremath{-}0.0409}  % experiments/FIGS/data/mass_resolution/analysis_holm/s7_mass_resolution.json :: table[cell=17+mass/17,probe=mlp].sigma_eff :: sha256 e03cde45ece9f385 :: sigma_eff with the mass output minus without, difference of the seed means; negative is better
\newcommand{\MassResDidMlp}{\ensuremath{-}0.0087}  % experiments/FIGS/data/mass_resolution/analysis_holm/s7_mass_resolution.json :: table[cell=162+mass/162/17+mass/17,probe=mlp].sigma_eff :: sha256 e03cde45ece9f385 :: gain at 162 classes minus gain at 17
\newcommand{\MassResTargetSigmaEff}{0.2986}  % experiments/FIGS/data/mass_resolution/analysis_holm/s7_mass_resolution.json :: table[*].target_sigma_eff :: sha256 e03cde45ece9f385 :: sigma_eff of the class-centred target itself (a model that predicts the class mean)
\newcommand{\MassResRidgeRtwoMin}{20\%}  % experiments/FIGS/data/mass_resolution/s1.json :: arms.{l188,l162,r42q1,r16q1}-s*.ridge.val_r2 :: sha256 88a57bbad270d51c :: validation R^2 of the ridge probe, smallest of 20 models without the mass output, all 5 per-seed files
\newcommand{\MassResRidgeRtwoMax}{30\%}  % experiments/FIGS/data/mass_resolution/s1.json :: arms.{l188,l162,r42q1,r16q1}-s*.ridge.val_r2 :: sha256 88a57bbad270d51c :: validation R^2 of the ridge probe, largest of 20 models without the mass output, all 5 per-seed files
\newcommand{\MassResSigmaEffOnesixtwo}{0.3302}  % experiments/FIGS/data/mass_resolution/analysis_holm/s7_mass_resolution.json :: table[cell=162,probe=ridge].sigma_eff :: sha256 e03cde45ece9f385 :: mean over 5 seeds
\newcommand{\MassResSdOnesixtwo}{0.626}  % experiments/FIGS/data/mass_resolution/analysis_holm/s7_mass_resolution.json :: table[cell=162,probe=ridge].sd :: sha256 e03cde45ece9f385 :: mean over 5 seeds
\newcommand{\MassResTailFractionOnesixtwo}{0.062}  % experiments/FIGS/data/mass_resolution/analysis_holm/s7_mass_resolution.json :: table[cell=162,probe=ridge].tail_fraction :: sha256 e03cde45ece9f385 :: mean over 5 seeds
\newcommand{\MassResSigmaEffOnesixtwoMass}{0.2238}  % experiments/FIGS/data/mass_resolution/analysis_holm/s7_mass_resolution.json :: table[cell=162+mass,probe=ridge].sigma_eff :: sha256 e03cde45ece9f385 :: mean over 5 seeds
\newcommand{\MassResSdOnesixtwoMass}{0.494}  % experiments/FIGS/data/mass_resolution/analysis_holm/s7_mass_resolution.json :: table[cell=162+mass,probe=ridge].sd :: sha256 e03cde45ece9f385 :: mean over 5 seeds
\newcommand{\MassResTailFractionOnesixtwoMass}{0.034}  % experiments/FIGS/data/mass_resolution/analysis_holm/s7_mass_resolution.json :: table[cell=162+mass,probe=ridge].tail_fraction :: sha256 e03cde45ece9f385 :: mean over 5 seeds
\newcommand{\MassResSigmaEffOneseven}{0.3001}  % experiments/FIGS/data/mass_resolution/analysis_holm/s7_mass_resolution.json :: table[cell=17,probe=ridge].sigma_eff :: sha256 e03cde45ece9f385 :: mean over 5 seeds
\newcommand{\MassResSdOneseven}{0.601}  % experiments/FIGS/data/mass_resolution/analysis_holm/s7_mass_resolution.json :: table[cell=17,probe=ridge].sd :: sha256 e03cde45ece9f385 :: mean over 5 seeds
\newcommand{\MassResTailFractionOneseven}{0.056}  % experiments/FIGS/data/mass_resolution/analysis_holm/s7_mass_resolution.json :: table[cell=17,probe=ridge].tail_fraction :: sha256 e03cde45ece9f385 :: mean over 5 seeds
\newcommand{\MassResSigmaEffOnesevenMass}{0.2170}  % experiments/FIGS/data/mass_resolution/analysis_holm/s7_mass_resolution.json :: table[cell=17+mass,probe=ridge].sigma_eff :: sha256 e03cde45ece9f385 :: mean over 5 seeds
\newcommand{\MassResSdOnesevenMass}{0.490}  % experiments/FIGS/data/mass_resolution/analysis_holm/s7_mass_resolution.json :: table[cell=17+mass,probe=ridge].sd :: sha256 e03cde45ece9f385 :: mean over 5 seeds
\newcommand{\MassResTailFractionOnesevenMass}{0.033}  % experiments/FIGS/data/mass_resolution/analysis_holm/s7_mass_resolution.json :: table[cell=17+mass,probe=ridge].tail_fraction :: sha256 e03cde45ece9f385 :: mean over 5 seeds
\newcommand{\MassResSigmaEffOneeighteight}{0.3393}  % experiments/FIGS/data/mass_resolution/analysis_holm/s7_mass_resolution.json :: table[cell=188,probe=ridge].sigma_eff :: sha256 e03cde45ece9f385 :: mean over 5 seeds
\newcommand{\MassResSdOneeighteight}{0.631}  % experiments/FIGS/data/mass_resolution/analysis_holm/s7_mass_resolution.json :: table[cell=188,probe=ridge].sd :: sha256 e03cde45ece9f385 :: mean over 5 seeds
\newcommand{\MassResTailFractionOneeighteight}{0.064}  % experiments/FIGS/data/mass_resolution/analysis_holm/s7_mass_resolution.json :: table[cell=188,probe=ridge].tail_fraction :: sha256 e03cde45ece9f385 :: mean over 5 seeds
\newcommand{\MassResSigmaEffFourthree}{0.3228}  % experiments/FIGS/data/mass_resolution/analysis_holm/s7_mass_resolution.json :: table[cell=43,probe=ridge].sigma_eff :: sha256 e03cde45ece9f385 :: mean over 5 seeds
\newcommand{\MassResSdFourthree}{0.618}  % experiments/FIGS/data/mass_resolution/analysis_holm/s7_mass_resolution.json :: table[cell=43,probe=ridge].sd :: sha256 e03cde45ece9f385 :: mean over 5 seeds
\newcommand{\MassResTailFractionFourthree}{0.060}  % experiments/FIGS/data/mass_resolution/analysis_holm/s7_mass_resolution.json :: table[cell=43,probe=ridge].tail_fraction :: sha256 e03cde45ece9f385 :: mean over 5 seeds
\newcommand{\AojEffPercent}{1\%}  % experiments/FIGS/data/aoj_full_v1/fit_v6/results.json :: eff :: sha256 2416fc21f93fcfa6 :: the score cut's pass fraction outside the mass windows
\newcommand{\AojPtMin}{500}  % experiments/AOJ/peak_fit.py :: PT_RANGE[0] :: sha256 cad06229a74532be :: GeV
\newcommand{\AojPtMax}{2{,}500}  % experiments/AOJ/peak_fit.py :: PT_RANGE[1] :: sha256 cad06229a74532be :: GeV
\newcommand{\AojRhoMin}{\ensuremath{-}5.5}  % experiments/AOJ/peak_fit.py :: RHO_RANGE[0] :: sha256 cad06229a74532be
\newcommand{\AojRhoMax}{\ensuremath{-}2.0}  % experiments/AOJ/peak_fit.py :: RHO_RANGE[1] :: sha256 cad06229a74532be
\newcommand{\AojWindowWLo}{65}  % experiments/AOJ/peak_fit.py :: PEAKS[W].window[0] :: sha256 cad06229a74532be :: GeV
\newcommand{\AojWindowWHi}{105}  % experiments/AOJ/peak_fit.py :: PEAKS[W].window[1] :: sha256 cad06229a74532be :: GeV
\newcommand{\AojWindowTopLo}{140}  % experiments/AOJ/peak_fit.py :: PEAKS[top].window[0] :: sha256 cad06229a74532be :: GeV
\newcommand{\AojWindowTopHi}{220}  % experiments/AOJ/peak_fit.py :: PEAKS[top].window[1] :: sha256 cad06229a74532be :: GeV
\newcommand{\AojFitLo}{105}  % experiments/AOJ/peak_fit.py :: PEAKS[top].fit_range[0] :: sha256 cad06229a74532be :: GeV
\newcommand{\AojFitHi}{300}  % experiments/AOJ/peak_fit.py :: PEAKS[top].fit_range[1] :: sha256 cad06229a74532be :: GeV
\newcommand{\AojMassBin}{5}  % experiments/AOJ/peak_fit.py :: MASS_BIN :: sha256 cad06229a74532be :: GeV
\newcommand{\AojMapOrderRho}{3}  % experiments/AOJ/peak_fit.py :: MAP_ORDER[0] :: sha256 cad06229a74532be :: polynomial order in rho
\newcommand{\AojMapOrderPt}{2}  % experiments/AOJ/peak_fit.py :: MAP_ORDER[1] :: sha256 cad06229a74532be :: polynomial order in ln pT
\newcommand{\AojNRunsWord}{five}  % experiments/FIGS/data/aoj_full_v1/analysis_v6/aoj_top.json :: per_label_set.label_sets.*.signal_yield.n :: sha256 5dfef902a5fe75ba :: pretraining runs per vocabulary
\newcommand{\AojNJetsRhoWindow}{9{,}441{,}209}  % experiments/FIGS/data/aoj_full_v1/fit_v6/results.json :: n_jets :: sha256 2416fc21f93fcfa6 :: staged jets inside the rho window, all files; not the jets in the mass fit
\newcommand{\AojYieldOneeighteight}{2{,}400\,\ensuremath{\pm}\,500}  % experiments/FIGS/data/aoj_full_v1/analysis_v6/aoj_top.json :: per_label_set.label_sets.188.signal_yield :: sha256 5dfef902a5fe75ba :: mean +- SD over 5 runs, shared shape
\newcommand{\AojStatErrOneeighteight}{211}  % experiments/FIGS/data/aoj_full_v1/analysis_v6/aoj_top.json :: per_label_set.label_sets.188.median_stat_err :: sha256 5dfef902a5fe75ba :: median per-fit statistical error
\newcommand{\AojYieldOwnShapeOneeighteight}{2{,}400\,\ensuremath{\pm}\,600}  % experiments/FIGS/data/aoj_full_v1/analysis_v6/aoj_top.json :: per_label_set.label_sets.188.by_shape.own_shape :: sha256 5dfef902a5fe75ba :: every fit with its own floated shape
\newcommand{\AojSpreadOverStatOneeighteight}{2.4}  % experiments/FIGS/data/aoj_full_v1/analysis_v6/aoj_top.json :: per_label_set.label_sets.188.signal_yield.sd / median_stat_err :: sha256 5dfef902a5fe75ba :: spread between runs over the median per-fit statistical error
\newcommand{\AojYieldOnesixtwo}{2{,}400\,\ensuremath{\pm}\,900}  % experiments/FIGS/data/aoj_full_v1/analysis_v6/aoj_top.json :: per_label_set.label_sets.162.signal_yield :: sha256 5dfef902a5fe75ba :: mean +- SD over 5 runs, shared shape
\newcommand{\AojStatErrOnesixtwo}{199}  % experiments/FIGS/data/aoj_full_v1/analysis_v6/aoj_top.json :: per_label_set.label_sets.162.median_stat_err :: sha256 5dfef902a5fe75ba :: median per-fit statistical error
\newcommand{\AojYieldOwnShapeOnesixtwo}{2{,}200\,\ensuremath{\pm}\,800}  % experiments/FIGS/data/aoj_full_v1/analysis_v6/aoj_top.json :: per_label_set.label_sets.162.by_shape.own_shape :: sha256 5dfef902a5fe75ba :: every fit with its own floated shape
\newcommand{\AojSpreadOverStatOnesixtwo}{4.5}  % experiments/FIGS/data/aoj_full_v1/analysis_v6/aoj_top.json :: per_label_set.label_sets.162.signal_yield.sd / median_stat_err :: sha256 5dfef902a5fe75ba :: spread between runs over the median per-fit statistical error
\newcommand{\AojYieldFourthree}{2{,}400\,\ensuremath{\pm}\,800}  % experiments/FIGS/data/aoj_full_v1/analysis_v6/aoj_top.json :: per_label_set.label_sets.43.signal_yield :: sha256 5dfef902a5fe75ba :: mean +- SD over 5 runs, shared shape
\newcommand{\AojStatErrFourthree}{184}  % experiments/FIGS/data/aoj_full_v1/analysis_v6/aoj_top.json :: per_label_set.label_sets.43.median_stat_err :: sha256 5dfef902a5fe75ba :: median per-fit statistical error
\newcommand{\AojYieldOwnShapeFourthree}{2{,}300\,\ensuremath{\pm}\,800}  % experiments/FIGS/data/aoj_full_v1/analysis_v6/aoj_top.json :: per_label_set.label_sets.43.by_shape.own_shape :: sha256 5dfef902a5fe75ba :: every fit with its own floated shape
\newcommand{\AojSpreadOverStatFourthree}{4.2}  % experiments/FIGS/data/aoj_full_v1/analysis_v6/aoj_top.json :: per_label_set.label_sets.43.signal_yield.sd / median_stat_err :: sha256 5dfef902a5fe75ba :: spread between runs over the median per-fit statistical error
\newcommand{\AojYieldOneseven}{2{,}350\,\ensuremath{\pm}\,240}  % experiments/FIGS/data/aoj_full_v1/analysis_v6/aoj_top.json :: per_label_set.label_sets.17.signal_yield :: sha256 5dfef902a5fe75ba :: mean +- SD over 5 runs, shared shape
\newcommand{\AojStatErrOneseven}{201}  % experiments/FIGS/data/aoj_full_v1/analysis_v6/aoj_top.json :: per_label_set.label_sets.17.median_stat_err :: sha256 5dfef902a5fe75ba :: median per-fit statistical error
\newcommand{\AojYieldOwnShapeOneseven}{2{,}480\,\ensuremath{\pm}\,290}  % experiments/FIGS/data/aoj_full_v1/analysis_v6/aoj_top.json :: per_label_set.label_sets.17.by_shape.own_shape :: sha256 5dfef902a5fe75ba :: every fit with its own floated shape
\newcommand{\AojSpreadOverStatOneseven}{1.2}  % experiments/FIGS/data/aoj_full_v1/analysis_v6/aoj_top.json :: per_label_set.label_sets.17.signal_yield.sd / median_stat_err :: sha256 5dfef902a5fe75ba :: spread between runs over the median per-fit statistical error
\newcommand{\AojYieldOnesixtwoMass}{2{,}420\,\ensuremath{\pm}\,240}  % experiments/FIGS/data/aoj_full_v1/analysis_v6/aoj_top.json :: per_label_set.label_sets.162+mass.signal_yield :: sha256 5dfef902a5fe75ba :: mean +- SD over 5 runs, shared shape
\newcommand{\AojStatErrOnesixtwoMass}{179}  % experiments/FIGS/data/aoj_full_v1/analysis_v6/aoj_top.json :: per_label_set.label_sets.162+mass.median_stat_err :: sha256 5dfef902a5fe75ba :: median per-fit statistical error
\newcommand{\AojYieldOwnShapeOnesixtwoMass}{2{,}350\,\ensuremath{\pm}\,320}  % experiments/FIGS/data/aoj_full_v1/analysis_v6/aoj_top.json :: per_label_set.label_sets.162+mass.by_shape.own_shape :: sha256 5dfef902a5fe75ba :: every fit with its own floated shape
\newcommand{\AojSpreadOverStatOnesixtwoMass}{1.3}  % experiments/FIGS/data/aoj_full_v1/analysis_v6/aoj_top.json :: per_label_set.label_sets.162+mass.signal_yield.sd / median_stat_err :: sha256 5dfef902a5fe75ba :: spread between runs over the median per-fit statistical error
\newcommand{\AojYieldOnesevenMass}{2{,}500\,\ensuremath{\pm}\,600}  % experiments/FIGS/data/aoj_full_v1/analysis_v6/aoj_top.json :: per_label_set.label_sets.17+mass.signal_yield :: sha256 5dfef902a5fe75ba :: mean +- SD over 5 runs, shared shape
\newcommand{\AojStatErrOnesevenMass}{200}  % experiments/FIGS/data/aoj_full_v1/analysis_v6/aoj_top.json :: per_label_set.label_sets.17+mass.median_stat_err :: sha256 5dfef902a5fe75ba :: median per-fit statistical error
\newcommand{\AojYieldOwnShapeOnesevenMass}{2{,}900\,\ensuremath{\pm}\,500}  % experiments/FIGS/data/aoj_full_v1/analysis_v6/aoj_top.json :: per_label_set.label_sets.17+mass.by_shape.own_shape :: sha256 5dfef902a5fe75ba :: every fit with its own floated shape
\newcommand{\AojSpreadOverStatOnesevenMass}{2.9}  % experiments/FIGS/data/aoj_full_v1/analysis_v6/aoj_top.json :: per_label_set.label_sets.17+mass.signal_yield.sd / median_stat_err :: sha256 5dfef902a5fe75ba :: spread between runs over the median per-fit statistical error
\newcommand{\AojOwnShapeShiftMin}{\ensuremath{-}7\%}  % experiments/FIGS/data/aoj_full_v1/analysis_v6/aoj_top.json :: per_label_set.label_sets.*.by_shape.own_shape.mean / signal_yield.mean - 1 :: sha256 5dfef902a5fe75ba :: the vocabulary means with each model's own shape, relative to the shared shape
\newcommand{\AojOwnShapeShiftMax}{+13\%}  % experiments/FIGS/data/aoj_full_v1/analysis_v6/aoj_top.json :: per_label_set.label_sets.*.by_shape.own_shape.mean / signal_yield.mean - 1 :: sha256 5dfef902a5fe75ba :: the vocabulary means with each model's own shape, relative to the shared shape
\newcommand{\AojPooledMean}{182.8}  % experiments/FIGS/data/aoj_full_v1/analysis_v6/aoj_top.json :: per_label_set.shapes.pooled :: sha256 5dfef902a5fe75ba :: GeV
\newcommand{\AojPooledWidth}{11.7}  % experiments/FIGS/data/aoj_full_v1/analysis_v6/aoj_top.json :: per_label_set.shapes.pooled :: sha256 5dfef902a5fe75ba :: GeV
\newcommand{\AojPoolSize}{30}  % experiments/FIGS/data/aoj_full_v1/analysis_v6/aoj_top.json :: per_label_set.shapes.pool (count) :: sha256 5dfef902a5fe75ba :: pretrained models the shared shape is fitted over
\newcommand{\AojNModels}{30}  % experiments/FIGS/data/aoj_full_v1/fit_v6/results.json :: models (pretrained) :: sha256 2416fc21f93fcfa6 :: fits
\newcommand{\AojOwnMeanMin}{181.6}  % experiments/FIGS/data/aoj_full_v1/fit_v6/results.json :: models.*.top.own_shape_fit.mean (min) :: sha256 2416fc21f93fcfa6 :: GeV
\newcommand{\AojOwnMeanMax}{184.0}  % experiments/FIGS/data/aoj_full_v1/fit_v6/results.json :: models.*.top.own_shape_fit.mean (max) :: sha256 2416fc21f93fcfa6 :: GeV
\newcommand{\AojOwnWidthMin}{10.0}  % experiments/FIGS/data/aoj_full_v1/fit_v6/results.json :: models.*.top.own_shape_fit.width (min) :: sha256 2416fc21f93fcfa6 :: GeV
\newcommand{\AojOwnWidthMax}{13.8}  % experiments/FIGS/data/aoj_full_v1/fit_v6/results.json :: models.*.top.own_shape_fit.width (max) :: sha256 2416fc21f93fcfa6 :: GeV
\newcommand{\AojShapeStepMean}{0.7}  % experiments/FIGS/data/aoj_full_v1/fit_v6/results.json :: pooled_shape.systematic_step.mean :: sha256 2416fc21f93fcfa6 :: GeV: the SD of the models' own floated values, one shift for every model together
\newcommand{\AojShapeSystMeanMax}{1\%}  % experiments/FIGS/data/aoj_full_v1/fit_v6/results.json :: models.*.top.shape_systematic.pooled_mean_(down|up) / signal_yield - 1 (largest |.|) :: sha256 2416fc21f93fcfa6 :: the largest yield change of any model under that shift
\newcommand{\AojShapeStepWidth}{1.0}  % experiments/FIGS/data/aoj_full_v1/fit_v6/results.json :: pooled_shape.systematic_step.width :: sha256 2416fc21f93fcfa6 :: GeV: the SD of the models' own floated values, one shift for every model together
\newcommand{\AojShapeSystWidthMax}{12\%}  % experiments/FIGS/data/aoj_full_v1/fit_v6/results.json :: models.*.top.shape_systematic.pooled_width_(down|up) / signal_yield - 1 (largest |.|) :: sha256 2416fc21f93fcfa6 :: the largest yield change of any model under that shift
\newcommand{\AojLeakEpsHi}{0.6}  % experiments/FIGS/data/aoj_full_v1/fit_v6/results.json :: models.*.top.leak_systematic.eps_ref (max) :: sha256 2416fc21f93fcfa6 :: the reference's assumed efficiency for data tops, upper variation
\newcommand{\AojLeakEpsLo}{0.4}  % experiments/FIGS/data/aoj_full_v1/fit_v6/results.json :: models.*.top.leak_systematic.eps_ref (min) :: sha256 2416fc21f93fcfa6 :: the same, lower variation
\newcommand{\AojLeakShiftMin}{+97}  % experiments/FIGS/data/aoj_full_v1/fit_v6/results.json :: models.*.top.leak_systematic.shift['0.6'] (min) :: sha256 2416fc21f93fcfa6 :: tops added to a yield
\newcommand{\AojLeakShiftMax}{+271}  % experiments/FIGS/data/aoj_full_v1/fit_v6/results.json :: models.*.top.leak_systematic.shift['0.4'] (max) :: sha256 2416fc21f93fcfa6 :: tops added to a yield
\newcommand{\AojFitPMin}{0.010}  % experiments/FIGS/data/aoj_full_v1/fit_v6/results.json :: models.*.top.asymptotic_p (min) :: sha256 2416fc21f93fcfa6 :: goodness of fit of the working-point fit, pretrained models
\newcommand{\AojFitPMax}{0.974}  % experiments/FIGS/data/aoj_full_v1/fit_v6/results.json :: models.*.top.asymptotic_p (max) :: sha256 2416fc21f93fcfa6
\newcommand{\AojWeakestYield}{885}  % experiments/FIGS/data/aoj_full_v1/fit_v6/results.json :: models.l162-s2.top.signal_yield :: sha256 2416fc21f93fcfa6 :: the smallest yield of any pretrained model
\newcommand{\AojWeakestYieldErr}{211}  % experiments/FIGS/data/aoj_full_v1/fit_v6/results.json :: models.l162-s2.top.signal_yield_err :: sha256 2416fc21f93fcfa6
\newcommand{\AojWeakestFitP}{0.94}  % experiments/FIGS/data/aoj_full_v1/fit_v6/results.json :: models.l162-s2.top.asymptotic_p :: sha256 2416fc21f93fcfa6
\newcommand{\AojSidebandZMin}{\ensuremath{-}0.43}  % experiments/FIGS/data/aoj_full_v1/fit_v6/histograms.npz :: experiments/AOJ/sideband_residuals.py: z (min over the pretrained models) :: sha256 2bd1e3d5b32d4a12 :: passing jets minus the fit, summed from the fit's lower edge to the top window
\newcommand{\AojSidebandZMax}{+0.58}  % experiments/FIGS/data/aoj_full_v1/fit_v6/histograms.npz :: experiments/AOJ/sideband_residuals.py: z (max over the pretrained models) :: sha256 2bd1e3d5b32d4a12
\newcommand{\AojWeakestSidebandZ}{0.00}  % experiments/FIGS/data/aoj_full_v1/fit_v6/histograms.npz :: experiments/AOJ/sideband_residuals.py: z of l162-s2 :: sha256 2bd1e3d5b32d4a12 :: the smallest-yield model
\newcommand{\AojSidebandMaxPull}{2.7}  % experiments/FIGS/data/aoj_full_v1/fit_v6/histograms.npz :: experiments/AOJ/sideband_residuals.py: max_abs_pull (max over the pretrained models) :: sha256 2bd1e3d5b32d4a12 :: largest single-bin pull below the top window, any pretrained model
\newcommand{\AojSharedBinLo}{125}  % experiments/FIGS/data/aoj_full_v1/fit_v6/histograms.npz :: the bin holding the most models' largest pull :: sha256 2bd1e3d5b32d4a12 :: GeV
\newcommand{\AojSharedBinHi}{130}  % experiments/FIGS/data/aoj_full_v1/fit_v6/histograms.npz :: its upper edge :: sha256 2bd1e3d5b32d4a12 :: GeV
\newcommand{\AojSharedBinNModels}{18}  % experiments/FIGS/data/aoj_full_v1/fit_v6/histograms.npz :: models whose largest pull is in that bin :: sha256 2bd1e3d5b32d4a12
\newcommand{\AojSharedBinMeanPull}{+1.4}  % experiments/FIGS/data/aoj_full_v1/fit_v6/histograms.npz :: experiments/AOJ/sideband_residuals.py mean_pulls: that bin :: sha256 2bd1e3d5b32d4a12 :: mean pull over the pretrained models
\newcommand{\AojSharedBinRank}{third}  % experiments/FIGS/data/aoj_full_v1/fit_v6/histograms.npz :: experiments/AOJ/sideband_residuals.py mean_pulls: rank of that bin's |mean pull| :: sha256 2bd1e3d5b32d4a12
\newcommand{\AojFitNBins}{39}  % experiments/FIGS/data/aoj_full_v1/fit_v6/histograms.npz :: mass bins of the fit range :: sha256 2bd1e3d5b32d4a12
\newcommand{\AojPairedYieldOnesixtwoOverOneeighteight}{0.94~[0.45, 1.95]}  % experiments/FIGS/data/aoj_full_v1/analysis_v6/aoj_top.json :: per_label_set.label_sets.{162,188}.signal_yields :: sha256 5dfef902a5fe75ba :: yield ratio paired by run, geometric mean [95% Student-t interval]
\newcommand{\AojPairedYieldFourthreeOverOneeighteight}{0.97~[0.81, 1.16]}  % experiments/FIGS/data/aoj_full_v1/analysis_v6/aoj_top.json :: per_label_set.label_sets.{43,188}.signal_yields :: sha256 5dfef902a5fe75ba :: yield ratio paired by run, geometric mean [95% Student-t interval]
\newcommand{\AojPairedYieldOnesevenOverOneeighteight}{1.01~[0.75, 1.35]}  % experiments/FIGS/data/aoj_full_v1/analysis_v6/aoj_top.json :: per_label_set.label_sets.{17,188}.signal_yields :: sha256 5dfef902a5fe75ba :: yield ratio paired by run, geometric mean [95% Student-t interval]
\newcommand{\AojPairedYieldOnesevenOverOnesixtwo}{1.07~[0.51, 2.24]}  % experiments/FIGS/data/aoj_full_v1/analysis_v6/aoj_top.json :: per_label_set.label_sets.{17,162}.signal_yields :: sha256 5dfef902a5fe75ba :: yield ratio paired by run, geometric mean [95% Student-t interval]
\newcommand{\AojPairedYieldOnesixtwoMassOverOnesixtwo}{1.10~[0.57, 2.14]}  % experiments/FIGS/data/aoj_full_v1/analysis_v6/aoj_top.json :: per_label_set.label_sets.{162+mass,162}.signal_yields :: sha256 5dfef902a5fe75ba :: yield ratio paired by run, geometric mean [95% Student-t interval]
\newcommand{\AojPairedYieldOnesevenMassOverOneseven}{1.05~[0.74, 1.49]}  % experiments/FIGS/data/aoj_full_v1/analysis_v6/aoj_top.json :: per_label_set.label_sets.{17+mass,17}.signal_yields :: sha256 5dfef902a5fe75ba :: yield ratio paired by run, geometric mean [95% Student-t interval]
\newcommand{\AojLeakPairedShiftMax}{2\%}  % experiments/FIGS/data/aoj_full_v1/fit_v6/results.json :: models.*.top.leak_systematic.signal_yield against signal_yield, paired ratios :: sha256 2416fc21f93fcfa6 :: largest change of a paired yield ratio when the fail-region tops are scaled up
\newcommand{\AojRefEffPercent}{1.38\%}  % experiments/FIGS/data/aoj_full_v1/fit_v6/results.json :: reference.top.data_efficiency :: sha256 2416fc21f93fcfa6 :: fraction of all data jets passing its cut
\newcommand{\AojRefFitP}{\ensuremath{2.56\times10^{-11}}}  % experiments/FIGS/data/aoj_full_v1/fit_v6/results.json :: reference.top.asymptotic_p :: sha256 2416fc21f93fcfa6 :: goodness of fit of the reference's own top fit
\newcommand{\AojRefValidationP}{\ensuremath{p \leq 0.005} (0 of 200 toys reach the observed deviance)}  % experiments/FIGS/data/aoj_full_v1/fit_v6/results.json :: reference.top.validation.toy_p, n_toys :: sha256 2416fc21f93fcfa6 :: toy p-value of the background validation, (k + 1) / (n + 1)
\newcommand{\AojRefBandSignalZ}{2.5}  % experiments/FIGS/data/aoj_full_v1/fit_v6/results.json :: reference.top.validation.band_signal_z :: sha256 2416fc21f93fcfa6 :: significance of a peak of the score's own shape in the validation band
\newcommand{\AojModelEffMin}{1.01\%}  % experiments/FIGS/data/aoj_full_v1/fit_v6/results.json :: models.*.top.data_efficiency (min) :: sha256 2416fc21f93fcfa6 :: pretrained models
\newcommand{\AojModelEffMax}{1.05\%}  % experiments/FIGS/data/aoj_full_v1/fit_v6/results.json :: models.*.top.data_efficiency (max) :: sha256 2416fc21f93fcfa6 :: pretrained models
\newcommand{\AojModelsValidationPMin}{0.11}  % experiments/FIGS/data/aoj_full_v1/fit_v6/results.json :: models.*.top.validation.toy_p (min) :: sha256 2416fc21f93fcfa6 :: pretrained models
\newcommand{\AojRefYield}{29{,}143}  % experiments/FIGS/data/aoj_full_v1/fit_v6/results.json :: reference.top.signal_yield :: sha256 2416fc21f93fcfa6
\newcommand{\AojRefYieldErr}{865}  % experiments/FIGS/data/aoj_full_v1/fit_v6/results.json :: reference.top.signal_yield_err :: sha256 2416fc21f93fcfa6
\newcommand{\AojRefMass}{181.0}  % experiments/FIGS/data/aoj_full_v1/fit_v6/results.json :: reference.top.floated_mean :: sha256 2416fc21f93fcfa6 :: GeV
\newcommand{\AojRefWidth}{17.8}  % experiments/FIGS/data/aoj_full_v1/fit_v6/results.json :: reference.top.floated_width :: sha256 2416fc21f93fcfa6 :: GeV
\newcommand{\AojPublicYield}{1{,}813}  % experiments/FIGS/data/aoj_full_v1/fit_v6/results.json :: models.sophon-public.top.signal_yield :: sha256 2416fc21f93fcfa6
\newcommand{\AojPublicYieldErr}{193}  % experiments/FIGS/data/aoj_full_v1/fit_v6/results.json :: models.sophon-public.top.signal_yield_err :: sha256 2416fc21f93fcfa6
\newcommand{\AojPublicMass}{182.8}  % experiments/FIGS/data/aoj_full_v1/fit_v6/results.json :: models.sophon-public.top.floated_mean :: sha256 2416fc21f93fcfa6 :: GeV
\newcommand{\AojPublicWidth}{11.7}  % experiments/FIGS/data/aoj_full_v1/fit_v6/results.json :: models.sophon-public.top.floated_width :: sha256 2416fc21f93fcfa6 :: GeV
\newcommand{\AojInjSizeMin}{1{,}000}  % experiments/FIGS/data/aoj_injection_v1/summary.json :: groups[region=top,mode=ensemble].size (min) :: sha256 d9b66466ae24b584 :: injected tops
\newcommand{\AojInjSizeMax}{4{,}000}  % experiments/FIGS/data/aoj_injection_v1/summary.json :: groups[region=top,mode=ensemble].size (max) :: sha256 d9b66466ae24b584
\newcommand{\AojInjFloatPullSdMax}{2.0}  % experiments/FIGS/data/aoj_injection_v1/summary.json :: groups[region=top,mode=bootstrap,variant=float].pull.sd (max) :: sha256 d9b66466ae24b584 :: pull width with each model's own floated shape
\newcommand{\AojInjEnsembleRecoveryMin}{0.983}  % experiments/FIGS/data/aoj_injection_v1/summary.json :: groups[region=top,mode=ensemble].ratio (min) :: sha256 d9b66466ae24b584 :: the whole procedure, shape re-derived per toy set
\newcommand{\AojInjEnsembleRecoveryMax}{0.998}  % experiments/FIGS/data/aoj_injection_v1/summary.json :: groups[region=top,mode=ensemble].ratio (max) :: sha256 d9b66466ae24b584
\newcommand{\AojInjEnsemblePullSdMin}{1.02}  % experiments/FIGS/data/aoj_injection_v1/summary.json :: groups[region=top,mode=ensemble].pull (min) :: sha256 d9b66466ae24b584 :: the whole procedure, shape re-derived per toy set
\newcommand{\AojInjEnsemblePullSdMax}{1.09}  % experiments/FIGS/data/aoj_injection_v1/summary.json :: groups[region=top,mode=ensemble].pull (max) :: sha256 d9b66466ae24b584
\newcommand{\AojClosureNShards}{10}  % experiments/FIGS/data/aoj_checks_v2/closure_all_shards.json :: n_shards :: sha256 6a6018351d7048d2 :: staged shards of the open data
\newcommand{\AojClosurePtLo}{500}  % experiments/FIGS/data/aoj_checks_v2/closure_all_shards.json :: pt_window[0] :: sha256 6a6018351d7048d2 :: GeV
\newcommand{\AojClosurePtHi}{700}  % experiments/FIGS/data/aoj_checks_v2/closure_all_shards.json :: pt_window[1] :: sha256 6a6018351d7048d2 :: GeV
\newcommand{\AojClosureDzeroIqr}{8.7}  % experiments/FIGS/data/aoj_checks_v2/closure_all_shards.json :: domain_shift.part_d0.iqr_ratio :: sha256 6a6018351d7048d2 :: open data over JetClass-II QCD, pooled over the shards
\newcommand{\AojClosureDzIqr}{10.6}  % experiments/FIGS/data/aoj_checks_v2/closure_all_shards.json :: domain_shift.part_dz.iqr_ratio :: sha256 6a6018351d7048d2 :: open data over JetClass-II QCD, pooled over the shards
\newcommand{\AojClosureNeutral}{0.65}  % experiments/FIGS/data/aoj_checks_v2/closure_all_shards.json :: features.part_isNeutralHadron.ratio.per_shard (mean) :: sha256 6a6018351d7048d2 :: open data over JetClass-II QCD
\newcommand{\AojClosureCharged}{1.22}  % experiments/FIGS/data/aoj_checks_v2/closure_all_shards.json :: features.part_isChargedHadron.ratio.per_shard (mean) :: sha256 6a6018351d7048d2 :: open data over JetClass-II QCD
\newcommand{\AojClosureNPartAoj}{61}  % experiments/FIGS/data/aoj_checks_v2/closure_all_shards.json :: features.n_particles.aoj.per_shard (mean of the shard medians) :: sha256 6a6018351d7048d2 :: median
\newcommand{\AojClosureNPartRef}{48}  % experiments/FIGS/data/aoj_checks_v2/closure_all_shards.json :: features.n_particles.reference.per_shard (mean) :: sha256 6a6018351d7048d2 :: median
\newcommand{\RandMergingDrawsBvcFourprong}{three}  % configs/labelmaps/rand_label_map.v1.csv :: RAND_d* of label_X_YY_bbqq and label_X_YY_ccqq :: sha256 f395a25f52e3cfad :: first-grid partitions that merge the task's pair
\newcommand{\RandSplittingDrawsBvcFourprong}{one and two}  % configs/labelmaps/rand_label_map.v1.csv :: RAND_d* of label_X_YY_bbqq and label_X_YY_ccqq :: sha256 f395a25f52e3cfad :: first-grid partitions that keep the pair apart
\newcommand{\RandMergingDrawsVisibleContent}{one and three}  % configs/labelmaps/rand_label_map.v1.csv :: RAND_d* of label_X_YY_bbqq and label_X_YY_cqtauhv :: sha256 f395a25f52e3cfad :: first-grid partitions that merge the task's pair
\newcommand{\RandSplittingDrawsVisibleContent}{two}  % configs/labelmaps/rand_label_map.v1.csv :: RAND_d* of label_X_YY_bbqq and label_X_YY_cqtauhv :: sha256 f395a25f52e3cfad :: first-grid partitions that keep the pair apart
\newcommand{\RandVtwoNPartitions}{five}  % configs/arms/v2_grid.json :: arms[RAND2_*] (count) :: sha256 d8a31b1486992302 :: random partitions in the rerun
\newcommand{\RandVtwoRuns}{one}  % configs/arms/v2_paper_runs.txt :: runs of RAND2_* :: sha256 b04562b14ea8cbc7 :: runs per partition
\newcommand{\RandFlavRuns}{two}  % configs/arms/v2_paper_runs.txt :: runs of FLAV_* :: sha256 b04562b14ea8cbc7 :: runs per flavour-pair model
\newcommand{\RandVtwoNPairs}{seven}  % configs/labelmaps/rand_v2_selection.json :: pairs (count) :: sha256 98cf52d49c7d2f02 :: probe pairs the balance rule covers
\newcommand{\RandVtwoMergedMin}{two}  % configs/labelmaps/rand_v2_selection.json :: merged_in :: sha256 98cf52d49c7d2f02 :: partitions merging each probe pair, checked on the map
\newcommand{\RandVtwoMergedMax}{three}  % configs/labelmaps/rand_v2_selection.json :: merged_in :: sha256 98cf52d49c7d2f02 :: partitions merging each probe pair, checked on the map
\newcommand{\RandVtwoShareTol}{5\%}  % configs/labelmaps/rand_v2_selection.json :: realised_tolerance :: sha256 98cf52d49c7d2f02 :: largest relative deviation of a realised group share from its target (rule 3)
\newcommand{\RandVtwoTrials}{32}  % scripts/build_rand_control.py :: TRIALS :: sha256 5bc91d77a3f2e81f :: share-matched assignments drawn per random partition; the one sharing the fewest merged pairs with the 17-class label set is kept
\newcommand{\RandVtwoPoolDraws}{6{,}600}  % configs/labelmaps/rand_v2_selection.json :: pool.draws :: sha256 98cf52d49c7d2f02 :: candidate partitions drawn (seeds) before the realised-share rule
\newcommand{\RandVtwoPoolPass}{27}  % configs/labelmaps/rand_v2_selection.json :: rule3_seeds (count) :: sha256 98cf52d49c7d2f02 :: candidate partitions meeting the realised-share rule (rule 3)
\newcommand{\RandFlavSplitOrbit}{\ensuremath{X\to YY\to QQQQ}}  % configs/labelmaps/flavour_pair.v2.json :: split_orbit :: sha256 3b79ce5e44c1e2f4 :: the decays F1 cuts, Q any quark
\newcommand{\RandFlavSplitOrbitClasses}{22}  % configs/labelmaps/flavour_pair_map.v2.csv :: classes whose orbit is split_orbit (count) :: sha256 9aab235099ceac43 :: classes of the split orbit
\newcommand{\RandFlavCutClasses}{11}  % configs/labelmaps/flavour_pair.v2.json :: split_classes_moved (count) :: sha256 3b79ce5e44c1e2f4 :: classes of the split orbit that F1 moves: those with a b quark
\newcommand{\RandFlavOrbitsMoved}{three}  % configs/labelmaps/flavour_pair.v2.json :: orbits_moved_B_to_A (count) :: sha256 3b79ce5e44c1e2f4 :: whole orbits F1 moves the other way
\newcommand{\RandFlavOrbitsMovedList}{\ensuremath{X\to YY\to QQ\mu\mu}, \ensuremath{X\to YY\to QQ\tau_h\tau_h} and \ensuremath{X\to \mu\mu}}  % configs/labelmaps/flavour_pair.v2.json :: orbits_moved_B_to_A :: sha256 3b79ce5e44c1e2f4 :: the orbits F1 moves the other way, Q any quark
\newcommand{\RandFlavOrbitClassesMoved}{9}  % configs/labelmaps/flavour_pair_map.v2.csv :: classes whose orbit is in orbits_moved_B_to_A :: sha256 9aab235099ceac43 :: classes in those orbits
\newcommand{\RandFlavFonerCutWithB}{6}  % configs/labelmaps/flavour_pair_map.v2.csv :: FLAV_F1R != FLAV_F0 within split_orbit, with a b quark (count) :: sha256 9aab235099ceac43 :: classes with a b quark among those the F1r cut moves
\newcommand{\RandFlavBcPairs}{11}  % configs/labelmaps/flavour_pair.v2.json :: b_c_pairs_in_split_orbit (count) :: sha256 3b79ce5e44c1e2f4 :: pairs of the split orbit's classes that differ only by exchanging b and c
\newcommand{\RandFlavFoneBcSplit}{8}  % configs/labelmaps/flavour_pair_map.v2.csv :: b_c_pairs_in_split_orbit split by FLAV_F1 (count) :: sha256 9aab235099ceac43
\newcommand{\RandFlavFonerBcSplit}{none}  % configs/labelmaps/flavour_pair_map.v2.csv :: b_c_pairs_in_split_orbit split by FLAV_F1R (count) :: sha256 9aab235099ceac43 :: PRESPEC A14: none
\newcommand{\RandFlavFonerShareTol}{1\%}  % configs/labelmaps/flavour_pair.v2.json :: F1R.realised_tolerance_vs_F1 :: sha256 3b79ce5e44c1e2f4 :: largest allowed relative difference of a group's realised share from F1's
\newcommand{\RandFlavFonerShareDev}{0.77\%}  % configs/labelmaps/realised_native_shares.v2.json :: share summed per FLAV_F1R and FLAV_F1 group of flavour_pair_map.v2.csv (largest relative difference) :: sha256 5728c12e88fa3683 :: realised shares from the final loader's dry run
\newcommand{\MlstNRunsWord}{three}  % configs/arms/v2_paper_runs.txt :: runs of L188/L162/R42_Q1/R16_Q1 :: sha256 b04562b14ea8cbc7 :: pretraining runs per configuration
\newcommand{\MlstMassNRunsWord}{two}  % configs/arms/v2_paper_runs.txt :: runs of L162_MASS/R16_Q1_MASS :: sha256 b04562b14ea8cbc7 :: pretraining runs per configuration
\newcommand{\LpCMin}{\ensuremath{1\times10^{-3}}}  % experiments/EVAL/linear_probe_v2.py :: C_GRID :: sha256 4c576b44812c18a2 :: smallest inverse regularisation
\newcommand{\LpCMax}{100}  % experiments/EVAL/linear_probe_v2.py :: C_GRID :: sha256 4c576b44812c18a2 :: largest inverse regularisation
\newcommand{\LpNTestJcii}{2{,}000{,}000}  % experiments/EVAL/linear_probe_v2.py :: TEST_FIRST.jc2 :: sha256 4c576b44812c18a2 :: JetClass-II test jets
\newcommand{\LpNTestAucJcii}{500{,}000}  % experiments/EVAL/linear_probe_v2.py :: TEST_FIRST.jc2 / AUC_STRIDE :: sha256 4c576b44812c18a2 :: JetClass-II test jets of the macro AUC
\newcommand{\LpConvergedGrad}{\ensuremath{1\times10^{-5}}}  % experiments/EVAL/linear_probe_v2.py :: CONVERGED_GRAD :: sha256 4c576b44812c18a2 :: largest gradient component a converged probe ends below
\newcommand{\FillNFeatures}{17}  % configs/finetune/TopReference.yaml :: inputs.pf_features.vars (count) :: sha256 501a6ffe5d0a8051 :: per-particle input features of the network
\newcommand{\FillTopNZero}{10}  % configs/finetune/TopReference.yaml :: inputs.pf_features.vars, part_zero_* (count) :: sha256 501a6ffe5d0a8051 :: features the data set lacks, set to zero
\newcommand{\FillTopNPresent}{7}  % configs/finetune/TopReference.yaml :: inputs.pf_features.vars, not part_zero_* (count) :: sha256 501a6ffe5d0a8051 :: features the data set provides
\newcommand{\FillQgNZero}{4}  % configs/finetune/EnergyFlowQG.yaml :: inputs.pf_features.vars, part_zero_* (count) :: sha256 db42d6471e8b865c :: features the data set lacks, set to zero
\newcommand{\FillQgNPresent}{13}  % configs/finetune/EnergyFlowQG.yaml :: inputs.pf_features.vars, not part_zero_* (count) :: sha256 db42d6471e8b865c :: features the data set provides
\newcommand{\DesignVtwoDataFraction}{20\%}  % experiments/MTX/k8s/v2/grid/job-mtx2-l188-s1-raunav.yaml :: --data-fraction :: sha256 d28ce02267a6b87a :: share of every file read per epoch
\newcommand{\DesignVtwoNValFiles}{25}  % experiments/MTX/k8s/v2/grid/job-mtx2-l188-s1-raunav.yaml :: --data-val (files, brace ranges expanded) :: sha256 d28ce02267a6b87a :: validation files of the fixed sample
\newcommand{\DesignWavgFirstEpoch}{70}  % experiments/MTX/pretrain_v2.py :: --num-epochs - WAVG_EPOCHS :: sha256 95da5a6e029855ee :: first epoch of the weight average (epochs count from 0)
\newcommand{\DesignWavgLastEpoch}{79}  % experiments/MTX/k8s/v2/grid/job-mtx2-l188-s1-raunav.yaml :: --num-epochs - 1 :: sha256 d28ce02267a6b87a :: last epoch
\newcommand{\DesignBnJets}{200{,}000}  % experiments/MTX/pretrain_v2.py :: BN_JETS :: sha256 95da5a6e029855ee :: training jets the BatchNorm statistics are recomputed on
\newcommand{\DesignVtwoNWindows}{five}  % experiments/MTX/k8s/v2/grid/job-mtx2-l188-s1-raunav.yaml :: 1 / --data-fraction :: sha256 d28ce02267a6b87a :: disjoint windows of each file's rows, one read per epoch
\newcommand{\DesignLofoNWindows}{three}  % experiments/MTX/k8s/v2/cut/grid/job-mtx2-l188lofo4p-s1-raunav.yaml :: --data-windows :: sha256 44a9573929efb7ac :: windows of the models trained without one family
\newcommand{\DesignVtwoPlanUsed}{97\%}  % experiments/FIGS/data/v2_loader/loader_dryrun/dryrun_s176_seed1.json :: epochs[*].(max_fetch_id + 1) / args.data_split_num, mean :: sha256 03b6586b26ad7a2f :: share of an epoch's planned reading used before its jets are reached, 20 epochs
\newcommand{\DesignLofoPlanUsed}{76\%}  % experiments/FIGS/data/v2_loader/loader_dryrun/dryrun_lofo4p_seed1.json :: epochs[*].(max_fetch_id + 1) / fetches_per_pass, mean :: sha256 2a1315e611d13ff5 :: the same, leaving the family out
\newcommand{\DesignLofoRowFactor}{1.29}  % experiments/FIGS/data/v2_loader/loader_dryrun/dryrun_lofo4p_seed1.json :: (mean plan used / --data-windows) over (mean plan used / k) of the full models' dry run :: sha256 2a1315e611d13ff5 :: how much more often a row of a file is read per epoch when the family is left out
\newcommand{\DesignDistinctJetsCycle}{\ensuremath{3.3\times10^{7}}}  % experiments/FIGS/data/v2_loader/loader_dryrun/dryrun_s176_seed1.json :: sum of epochs[c*k:(c+1)*k].distinct_jets, mean over 4 cycles :: sha256 03b6586b26ad7a2f :: distinct jets the full models train on in one cycle of windows
\newcommand{\DesignLofoFamilyShare}{22.4\%}  % experiments/FIGS/data/v2_loader/loader_dryrun/dryrun_s176_seed1.json :: epochs[*].native_counts over summary.labels_absent_every_epoch of the family-out run :: sha256 03b6586b26ad7a2f :: the left-out family's share of the full models' stream
\newcommand{\DesignLofoNClasses}{27}  % experiments/FIGS/data/v2_loader/loader_dryrun/dryrun_lofo4p_seed1.json :: summary.labels_absent_every_epoch (count) :: sha256 2a1315e611d13ff5 :: native classes of the left-out family
\newcommand{\DesignLofoExposureFactor}{1.29}  % experiments/FIGS/data/v2_loader/loader_dryrun/dryrun_s176_seed1.json :: 1 / (1 - family share) :: sha256 03b6586b26ad7a2f :: how much more often every other class is seen
\newcommand{\DesignLofoQcdShareFull}{14.4\%}  % experiments/FIGS/data/v2_loader/loader_dryrun/dryrun_s176_seed1.json :: summary.mean_qcd_share :: sha256 03b6586b26ad7a2f :: QCD share of the realised stream
\newcommand{\DesignLofoQcdShareLofo}{18.6\%}  % experiments/FIGS/data/v2_loader/loader_dryrun/dryrun_lofo4p_seed1.json :: summary.mean_qcd_share :: sha256 2a1315e611d13ff5 :: QCD share of the realised stream
\newcommand{\DesignLrFlatPercent}{70\%}  % experiments/MTX/pretrain_v2.py :: make_optimizer: 1 - the decay share :: sha256 95da5a6e029855ee :: share of the epochs at the initial rate
\newcommand{\DesignLrDecayFirstEpoch}{56}  % experiments/MTX/pretrain_v2.py :: num_epochs - int(num_epochs * decay share) :: sha256 95da5a6e029855ee :: first epoch at a reduced rate (epochs from 0)
\newcommand{\DesignLrEndPercent}{1\%}  % experiments/MTX/pretrain_v2.py :: make_optimizer: gamma ** n_decay :: sha256 95da5a6e029855ee :: the last epoch's rate over the initial rate
\newcommand{\DesignEndLr}{\ensuremath{5\times10^{-6}}}  % experiments/MTX/k8s/v2/grid/job-mtx2-l188-s1-raunav.yaml :: --start-lr x the final factor :: sha256 d28ce02267a6b87a :: the rate of the last epoch
\newcommand{\DesignWeaverVersion}{0.4.17}  % experiments/MTX/pretrain_v2.py :: the weaver release whose loop it reproduces :: sha256 95da5a6e029855ee
\newcommand{\DesignMeiFactor}{1.1}  % src/stats/mde.py :: MEI_LOG = ln(factor) :: sha256 626468dcc37940bb :: the smallest effect of interest, as a factor
\newcommand{\DesignVtwoGpu}{RTX 3090}  % experiments/MTX/k8s/v2/grid/job-mtx2-flavf0-s1-raunav.yaml :: nodeAffinity nvidia.com/gpu.product :: sha256 960cad1691314980 :: the product of run index 1, every configuration
\newcommand{\DesignVtwoSecondGpu}{\pending{pending: the GPU type of each run index}}  % experiments/FIGS/make_tables.py :: emit_v2_recipe[DesignVtwoSecondGpu] :: sha256 10e6ec08356a4a2b :: a marker, not a number: the grid's job specs put every run index on one product
\newcommand{\DesignVtwoSecondGpuRuns}{\pending{pending: the GPU type of each run index}}  % experiments/FIGS/make_tables.py :: emit_v2_recipe[DesignVtwoSecondGpuRuns] :: sha256 10e6ec08356a4a2b :: a marker, not a number: the grid's job specs put every run index on one product
\newcommand{\MassNRunsWord}{five}  % experiments/FIGS/data/probe_ladder_mass2x2_mlp2/s1.json :: per-seed files (count) :: sha256 906e778d7043fe67 :: pretraining runs of each mass-output configuration
\newcommand{\MassOmaLinearOnesixtwo}{\ensuremath{(2.36 \pm 0.06)\times10^{-3}}}  % experiments/FIGS/data/probe_ladder_v2_mlp2/analysis/seed_level_results.json :: table[bvc_resonant,linear,162].auc :: sha256 53e1ec81343ac2c1 :: 1-AUC, mean +- SD over 5 seeds
\newcommand{\MassOmaLinearOnesixtwoMass}{\ensuremath{(2.62 \pm 0.08)\times10^{-3}}}  % experiments/FIGS/data/probe_ladder_mass2x2_mlp2/s1.json :: tasks.bvc_resonant.arms.l162mass-s*.linear.auc :: sha256 906e778d7043fe67 :: 1-AUC, mean +- SD over the 5 per-seed files
\newcommand{\MassOmaRatioLinearOnesixtwo}{1.11}  % experiments/FIGS/data/probe_ladder_mass2x2_mlp2/s1.json :: tasks.bvc_resonant.arms.l162mass-s*.linear.auc over table[bvc_resonant,linear,162].auc :: sha256 906e778d7043fe67 :: seed mean of 1-AUC with the mass output over without
\newcommand{\MassOmaLinearOneseven}{\ensuremath{(9.3 \pm 1.5)\times10^{-3}}}  % experiments/FIGS/data/probe_ladder_v2_mlp2/analysis/seed_level_results.json :: table[bvc_resonant,linear,17].auc :: sha256 53e1ec81343ac2c1 :: 1-AUC, mean +- SD over 5 seeds
\newcommand{\MassOmaLinearOnesevenMass}{\ensuremath{(3.0 \pm 1.2)\times10^{-2}}}  % experiments/FIGS/data/probe_ladder_mass2x2_mlp2/s1.json :: tasks.bvc_resonant.arms.r16q1mass-s*.linear.auc :: sha256 906e778d7043fe67 :: 1-AUC, mean +- SD over the 5 per-seed files
\newcommand{\MassOmaRatioLinearOneseven}{3.2}  % experiments/FIGS/data/probe_ladder_mass2x2_mlp2/s1.json :: tasks.bvc_resonant.arms.r16q1mass-s*.linear.auc over table[bvc_resonant,linear,17].auc :: sha256 906e778d7043fe67 :: seed mean of 1-AUC with the mass output over without
\newcommand{\MassOmaMlpOnesixtwo}{\ensuremath{(2.32 \pm 0.05)\times10^{-3}}}  % experiments/FIGS/data/probe_ladder_v2_mlp2/analysis/seed_level_results.json :: table[bvc_resonant,mlp,162].auc :: sha256 53e1ec81343ac2c1 :: 1-AUC, mean +- SD over 5 seeds
\newcommand{\MassOmaMlpOnesixtwoMass}{\ensuremath{(2.52 \pm 0.06)\times10^{-3}}}  % experiments/FIGS/data/probe_ladder_mass2x2_mlp2/s1.json :: tasks.bvc_resonant.arms.l162mass-s*.mlp.auc :: sha256 906e778d7043fe67 :: 1-AUC, mean +- SD over the 5 per-seed files
\newcommand{\MassOmaRatioMlpOnesixtwo}{1.09}  % experiments/FIGS/data/probe_ladder_mass2x2_mlp2/s1.json :: tasks.bvc_resonant.arms.l162mass-s*.mlp.auc over table[bvc_resonant,mlp,162].auc :: sha256 906e778d7043fe67 :: seed mean of 1-AUC with the mass output over without
\newcommand{\MassOmaMlpOneseven}{\ensuremath{(7.3 \pm 0.9)\times10^{-3}}}  % experiments/FIGS/data/probe_ladder_v2_mlp2/analysis/seed_level_results.json :: table[bvc_resonant,mlp,17].auc :: sha256 53e1ec81343ac2c1 :: 1-AUC, mean +- SD over 5 seeds
\newcommand{\MassOmaMlpOnesevenMass}{\ensuremath{(1.5 \pm 0.4)\times10^{-2}}}  % experiments/FIGS/data/probe_ladder_mass2x2_mlp2/s1.json :: tasks.bvc_resonant.arms.r16q1mass-s*.mlp.auc :: sha256 906e778d7043fe67 :: 1-AUC, mean +- SD over the 5 per-seed files
\newcommand{\MassOmaRatioMlpOneseven}{2.0}  % experiments/FIGS/data/probe_ladder_mass2x2_mlp2/s1.json :: tasks.bvc_resonant.arms.r16q1mass-s*.mlp.auc over table[bvc_resonant,mlp,17].auc :: sha256 906e778d7043fe67 :: seed mean of 1-AUC with the mass output over without
\newcommand{\FtAucJciiEThreeOneeighteight}{0.902\,\ensuremath{\pm}\,0.004}  % experiments/FIGS/data/w2b_leg1_metrics_v2.json :: cells.{l188-s1,l188-s2,l188-s3,l188-s4,l188-s5}.N1000.s1.macro_auc_ovr :: sha256 b56bf70be9982044 :: mean +- SD over the row's 5 pretrained models
\newcommand{\FtAucJciiEFourOneeighteight}{0.9723\,\ensuremath{\pm}\,0.0004}  % experiments/FIGS/data/w2b_leg1_metrics_v2.json :: cells.{l188-s1,l188-s2,l188-s3,l188-s4,l188-s5}.N10000.s1.macro_auc_ovr :: sha256 b56bf70be9982044 :: mean +- SD over the row's 5 pretrained models
\newcommand{\FtAucJciiEFiveOneeighteight}{0.9793\,\ensuremath{\pm}\,0.0002}  % experiments/FIGS/data/w2b_leg1_metrics_v2.json :: cells.{l188-s1,l188-s2,l188-s3,l188-s4,l188-s5}.N100000.s1.macro_auc_ovr :: sha256 b56bf70be9982044 :: mean +- SD over the row's 5 pretrained models
\newcommand{\FtAucJciiESixOneeighteight}{0.9809\,\ensuremath{\pm}\,\ensuremath{<}0.0001}  % experiments/FIGS/data/w2b_leg1_metrics_v2.json :: cells.{l188-s1,l188-s2,l188-s3,l188-s4,l188-s5}.N1000000.s1.macro_auc_ovr :: sha256 b56bf70be9982044 :: mean +- SD over the row's 5 pretrained models
\newcommand{\FtAucJciiEThreeOnesixtwo}{0.906\,\ensuremath{\pm}\,0.007}  % experiments/FIGS/data/w2b_leg1_metrics_v2.json :: cells.{l162-s1b,l162-s2,l162-s3,l162-s4,l162-s5}.N1000.s1.macro_auc_ovr :: sha256 b56bf70be9982044 :: mean +- SD over the row's 5 pretrained models
\newcommand{\FtAucJciiEFourOnesixtwo}{0.9724\,\ensuremath{\pm}\,0.0005}  % experiments/FIGS/data/w2b_leg1_metrics_v2.json :: cells.{l162-s1b,l162-s2,l162-s3,l162-s4,l162-s5}.N10000.s1.macro_auc_ovr :: sha256 b56bf70be9982044 :: mean +- SD over the row's 5 pretrained models
\newcommand{\FtAucJciiEFiveOnesixtwo}{0.9794\,\ensuremath{\pm}\,0.0002}  % experiments/FIGS/data/w2b_leg1_metrics_v2.json :: cells.{l162-s1b,l162-s2,l162-s3,l162-s4,l162-s5}.N100000.s1.macro_auc_ovr :: sha256 b56bf70be9982044 :: mean +- SD over the row's 5 pretrained models
\newcommand{\FtAucJciiESixOnesixtwo}{0.9808\,\ensuremath{\pm}\,0.0001}  % experiments/FIGS/data/w2b_leg1_metrics_v2.json :: cells.{l162-s1b,l162-s2,l162-s3,l162-s4,l162-s5}.N1000000.s1.macro_auc_ovr :: sha256 b56bf70be9982044 :: mean +- SD over the row's 5 pretrained models
\newcommand{\FtAucJciiEThreeFourthree}{0.918\,\ensuremath{\pm}\,0.005}  % experiments/FIGS/data/w2b_leg1_metrics_v2.json :: cells.{r42q1-s1,r42q1-s2,r42q1-s3,r42q1-s4,r42q1-s5}.N1000.s1.macro_auc_ovr :: sha256 b56bf70be9982044 :: mean +- SD over the row's 5 pretrained models
\newcommand{\FtAucJciiEFourFourthree}{0.9667\,\ensuremath{\pm}\,0.0005}  % experiments/FIGS/data/w2b_leg1_metrics_v2.json :: cells.{r42q1-s1,r42q1-s2,r42q1-s3,r42q1-s4,r42q1-s5}.N10000.s1.macro_auc_ovr :: sha256 b56bf70be9982044 :: mean +- SD over the row's 5 pretrained models
\newcommand{\FtAucJciiEFiveFourthree}{0.9775\,\ensuremath{\pm}\,0.0003}  % experiments/FIGS/data/w2b_leg1_metrics_v2.json :: cells.{r42q1-s1,r42q1-s2,r42q1-s3,r42q1-s4,r42q1-s5}.N100000.s1.macro_auc_ovr :: sha256 b56bf70be9982044 :: mean +- SD over the row's 5 pretrained models
\newcommand{\FtAucJciiESixFourthree}{0.9800\,\ensuremath{\pm}\,0.0001}  % experiments/FIGS/data/w2b_leg1_metrics_v2.json :: cells.{r42q1-s1,r42q1-s2,r42q1-s3,r42q1-s4,r42q1-s5}.N1000000.s1.macro_auc_ovr :: sha256 b56bf70be9982044 :: mean +- SD over the row's 5 pretrained models
\newcommand{\FtAucJciiEThreeOneseven}{0.905\,\ensuremath{\pm}\,0.008}  % experiments/FIGS/data/w2b_leg1_metrics_v2.json :: cells.{r16q1-s1,r16q1-s2,r16q1-s3,r16q1-s4,r16q1-s5}.N1000.s1.macro_auc_ovr :: sha256 b56bf70be9982044 :: mean +- SD over the row's 5 pretrained models
\newcommand{\FtAucJciiEFourOneseven}{0.9635\,\ensuremath{\pm}\,0.0005}  % experiments/FIGS/data/w2b_leg1_metrics_v2.json :: cells.{r16q1-s1,r16q1-s2,r16q1-s3,r16q1-s4,r16q1-s5}.N10000.s1.macro_auc_ovr :: sha256 b56bf70be9982044 :: mean +- SD over the row's 5 pretrained models
\newcommand{\FtAucJciiEFiveOneseven}{0.9748\,\ensuremath{\pm}\,0.0003}  % experiments/FIGS/data/w2b_leg1_metrics_v2.json :: cells.{r16q1-s1,r16q1-s2,r16q1-s3,r16q1-s4,r16q1-s5}.N100000.s1.macro_auc_ovr :: sha256 b56bf70be9982044 :: mean +- SD over the row's 5 pretrained models
\newcommand{\FtAucJciiESixOneseven}{0.9781\,\ensuremath{\pm}\,0.0001}  % experiments/FIGS/data/w2b_leg1_metrics_v2.json :: cells.{r16q1-s1,r16q1-s2,r16q1-s3,r16q1-s4,r16q1-s5}.N1000000.s1.macro_auc_ovr :: sha256 b56bf70be9982044 :: mean +- SD over the row's 5 pretrained models
\newcommand{\FtAucJciiEThreeOnesixtwoMass}{0.921\,\ensuremath{\pm}\,0.005}  % experiments/FIGS/data/w2b_leg1_metrics_v2.json :: cells.{l162mass-s1,l162mass-s2,l162mass-s3,l162mass-s4,l162mass-s5}.N1000.s1.macro_auc_ovr :: sha256 b56bf70be9982044 :: mean +- SD over the row's 5 pretrained models
\newcommand{\FtAucJciiEFourOnesixtwoMass}{0.9724\,\ensuremath{\pm}\,0.0013}  % experiments/FIGS/data/w2b_leg1_metrics_v2.json :: cells.{l162mass-s1,l162mass-s2,l162mass-s3,l162mass-s4,l162mass-s5}.N10000.s1.macro_auc_ovr :: sha256 b56bf70be9982044 :: mean +- SD over the row's 5 pretrained models
\newcommand{\FtAucJciiEFiveOnesixtwoMass}{0.9790\,\ensuremath{\pm}\,0.0001}  % experiments/FIGS/data/w2b_leg1_metrics_v2.json :: cells.{l162mass-s1,l162mass-s2,l162mass-s3,l162mass-s4,l162mass-s5}.N100000.s1.macro_auc_ovr :: sha256 b56bf70be9982044 :: mean +- SD over the row's 5 pretrained models
\newcommand{\FtAucJciiESixOnesixtwoMass}{0.9805\,\ensuremath{\pm}\,0.0001}  % experiments/FIGS/data/w2b_leg1_metrics_v2.json :: cells.{l162mass-s1,l162mass-s2,l162mass-s3,l162mass-s4,l162mass-s5}.N1000000.s1.macro_auc_ovr :: sha256 b56bf70be9982044 :: mean +- SD over the row's 5 pretrained models
\newcommand{\FtAucJciiEThreeOnesevenMass}{0.904\,\ensuremath{\pm}\,0.007}  % experiments/FIGS/data/w2b_leg1_metrics_v2.json :: cells.{r16q1mass-s1,r16q1mass-s2,r16q1mass-s3,r16q1mass-s4,r16q1mass-s5}.N1000.s1.macro_auc_ovr :: sha256 b56bf70be9982044 :: mean +- SD over the row's 5 pretrained models
\newcommand{\FtAucJciiEFourOnesevenMass}{0.9639\,\ensuremath{\pm}\,0.0007}  % experiments/FIGS/data/w2b_leg1_metrics_v2.json :: cells.{r16q1mass-s1,r16q1mass-s2,r16q1mass-s3,r16q1mass-s4,r16q1mass-s5}.N10000.s1.macro_auc_ovr :: sha256 b56bf70be9982044 :: mean +- SD over the row's 5 pretrained models
\newcommand{\FtAucJciiEFiveOnesevenMass}{0.9740\,\ensuremath{\pm}\,0.0004}  % experiments/FIGS/data/w2b_leg1_metrics_v2.json :: cells.{r16q1mass-s1,r16q1mass-s2,r16q1mass-s3,r16q1mass-s4,r16q1mass-s5}.N100000.s1.macro_auc_ovr :: sha256 b56bf70be9982044 :: mean +- SD over the row's 5 pretrained models
\newcommand{\FtAucJciiESixOnesevenMass}{0.9774\,\ensuremath{\pm}\,0.0001}  % experiments/FIGS/data/w2b_leg1_metrics_v2.json :: cells.{r16q1mass-s1,r16q1mass-s2,r16q1mass-s3,r16q1mass-s4,r16q1mass-s5}.N1000000.s1.macro_auc_ovr :: sha256 b56bf70be9982044 :: mean +- SD over the row's 5 pretrained models
\newcommand{\FtAucJciiEThreeRandDrawOne}{0.921}  % experiments/FIGS/data/w2b_leg1_metrics_v2.json :: cells.{rand-d1-s1b}.N1000.s1.macro_auc_ovr :: sha256 b56bf70be9982044 :: one run
\newcommand{\FtAucJciiEFourRandDrawOne}{0.9675}  % experiments/FIGS/data/w2b_leg1_metrics_v2.json :: cells.{rand-d1-s1b}.N10000.s1.macro_auc_ovr :: sha256 b56bf70be9982044 :: one run
\newcommand{\FtAucJciiEFiveRandDrawOne}{0.97617}  % experiments/FIGS/data/w2b_leg1_metrics_v2.json :: cells.{rand-d1-s1b}.N100000.s1.macro_auc_ovr :: sha256 b56bf70be9982044 :: one run
\newcommand{\FtAucJciiESixRandDrawOne}{0.97941}  % experiments/FIGS/data/w2b_leg1_metrics_v2.json :: cells.{rand-d1-s1b}.N1000000.s1.macro_auc_ovr :: sha256 b56bf70be9982044 :: one run
\newcommand{\FtAucJciiEThreeRandDrawTwo}{0.914}  % experiments/FIGS/data/w2b_leg1_metrics_v2.json :: cells.{rand-d2-s2}.N1000.s1.macro_auc_ovr :: sha256 b56bf70be9982044 :: one run
\newcommand{\FtAucJciiEFourRandDrawTwo}{0.9675}  % experiments/FIGS/data/w2b_leg1_metrics_v2.json :: cells.{rand-d2-s2}.N10000.s1.macro_auc_ovr :: sha256 b56bf70be9982044 :: one run
\newcommand{\FtAucJciiEFiveRandDrawTwo}{0.97668}  % experiments/FIGS/data/w2b_leg1_metrics_v2.json :: cells.{rand-d2-s2}.N100000.s1.macro_auc_ovr :: sha256 b56bf70be9982044 :: one run
\newcommand{\FtAucJciiESixRandDrawTwo}{0.97943}  % experiments/FIGS/data/w2b_leg1_metrics_v2.json :: cells.{rand-d2-s2}.N1000000.s1.macro_auc_ovr :: sha256 b56bf70be9982044 :: one run
\newcommand{\FtAucJciiEThreeRandDrawThree}{0.918}  % experiments/FIGS/data/w2b_leg1_metrics_v2.json :: cells.{rand-d3-s3}.N1000.s1.macro_auc_ovr :: sha256 b56bf70be9982044 :: one run
\newcommand{\FtAucJciiEFourRandDrawThree}{0.9686}  % experiments/FIGS/data/w2b_leg1_metrics_v2.json :: cells.{rand-d3-s3}.N10000.s1.macro_auc_ovr :: sha256 b56bf70be9982044 :: one run
\newcommand{\FtAucJciiEFiveRandDrawThree}{0.97605}  % experiments/FIGS/data/w2b_leg1_metrics_v2.json :: cells.{rand-d3-s3}.N100000.s1.macro_auc_ovr :: sha256 b56bf70be9982044 :: one run
\newcommand{\FtAucJciiESixRandDrawThree}{0.97950}  % experiments/FIGS/data/w2b_leg1_metrics_v2.json :: cells.{rand-d3-s3}.N1000000.s1.macro_auc_ovr :: sha256 b56bf70be9982044 :: one run
\newcommand{\FtAccJciiEThreeOneeighteight}{0.445\,\ensuremath{\pm}\,0.004}  % experiments/FIGS/data/w2b_leg1_metrics_v2.json :: cells.{l188-s1,l188-s2,l188-s3,l188-s4,l188-s5}.N1000.s1.accuracy :: sha256 b56bf70be9982044 :: mean +- SD over the row's 5 pretrained models
\newcommand{\FtAccJciiEFourOneeighteight}{0.5199\,\ensuremath{\pm}\,0.0018}  % experiments/FIGS/data/w2b_leg1_metrics_v2.json :: cells.{l188-s1,l188-s2,l188-s3,l188-s4,l188-s5}.N10000.s1.accuracy :: sha256 b56bf70be9982044 :: mean +- SD over the row's 5 pretrained models
\newcommand{\FtAccJciiEFiveOneeighteight}{0.5379\,\ensuremath{\pm}\,0.0010}  % experiments/FIGS/data/w2b_leg1_metrics_v2.json :: cells.{l188-s1,l188-s2,l188-s3,l188-s4,l188-s5}.N100000.s1.accuracy :: sha256 b56bf70be9982044 :: mean +- SD over the row's 5 pretrained models
\newcommand{\FtAccJciiESixOneeighteight}{0.5443\,\ensuremath{\pm}\,0.0003}  % experiments/FIGS/data/w2b_leg1_metrics_v2.json :: cells.{l188-s1,l188-s2,l188-s3,l188-s4,l188-s5}.N1000000.s1.accuracy :: sha256 b56bf70be9982044 :: mean +- SD over the row's 5 pretrained models
\newcommand{\FtAccJciiEThreeOnesixtwo}{0.448\,\ensuremath{\pm}\,0.006}  % experiments/FIGS/data/w2b_leg1_metrics_v2.json :: cells.{l162-s1b,l162-s2,l162-s3,l162-s4,l162-s5}.N1000.s1.accuracy :: sha256 b56bf70be9982044 :: mean +- SD over the row's 5 pretrained models
\newcommand{\FtAccJciiEFourOnesixtwo}{0.5212\,\ensuremath{\pm}\,0.0013}  % experiments/FIGS/data/w2b_leg1_metrics_v2.json :: cells.{l162-s1b,l162-s2,l162-s3,l162-s4,l162-s5}.N10000.s1.accuracy :: sha256 b56bf70be9982044 :: mean +- SD over the row's 5 pretrained models
\newcommand{\FtAccJciiEFiveOnesixtwo}{0.5387\,\ensuremath{\pm}\,0.0008}  % experiments/FIGS/data/w2b_leg1_metrics_v2.json :: cells.{l162-s1b,l162-s2,l162-s3,l162-s4,l162-s5}.N100000.s1.accuracy :: sha256 b56bf70be9982044 :: mean +- SD over the row's 5 pretrained models
\newcommand{\FtAccJciiESixOnesixtwo}{0.5440\,\ensuremath{\pm}\,0.0006}  % experiments/FIGS/data/w2b_leg1_metrics_v2.json :: cells.{l162-s1b,l162-s2,l162-s3,l162-s4,l162-s5}.N1000000.s1.accuracy :: sha256 b56bf70be9982044 :: mean +- SD over the row's 5 pretrained models
\newcommand{\FtAccJciiEThreeFourthree}{0.380\,\ensuremath{\pm}\,0.005}  % experiments/FIGS/data/w2b_leg1_metrics_v2.json :: cells.{r42q1-s1,r42q1-s2,r42q1-s3,r42q1-s4,r42q1-s5}.N1000.s1.accuracy :: sha256 b56bf70be9982044 :: mean +- SD over the row's 5 pretrained models
\newcommand{\FtAccJciiEFourFourthree}{0.481\,\ensuremath{\pm}\,0.004}  % experiments/FIGS/data/w2b_leg1_metrics_v2.json :: cells.{r42q1-s1,r42q1-s2,r42q1-s3,r42q1-s4,r42q1-s5}.N10000.s1.accuracy :: sha256 b56bf70be9982044 :: mean +- SD over the row's 5 pretrained models
\newcommand{\FtAccJciiEFiveFourthree}{0.5219\,\ensuremath{\pm}\,0.0017}  % experiments/FIGS/data/w2b_leg1_metrics_v2.json :: cells.{r42q1-s1,r42q1-s2,r42q1-s3,r42q1-s4,r42q1-s5}.N100000.s1.accuracy :: sha256 b56bf70be9982044 :: mean +- SD over the row's 5 pretrained models
\newcommand{\FtAccJciiESixFourthree}{0.5362\,\ensuremath{\pm}\,0.0007}  % experiments/FIGS/data/w2b_leg1_metrics_v2.json :: cells.{r42q1-s1,r42q1-s2,r42q1-s3,r42q1-s4,r42q1-s5}.N1000000.s1.accuracy :: sha256 b56bf70be9982044 :: mean +- SD over the row's 5 pretrained models
\newcommand{\FtAccJciiEThreeOneseven}{0.332\,\ensuremath{\pm}\,0.004}  % experiments/FIGS/data/w2b_leg1_metrics_v2.json :: cells.{r16q1-s1,r16q1-s2,r16q1-s3,r16q1-s4,r16q1-s5}.N1000.s1.accuracy :: sha256 b56bf70be9982044 :: mean +- SD over the row's 5 pretrained models
\newcommand{\FtAccJciiEFourOneseven}{0.447\,\ensuremath{\pm}\,0.013}  % experiments/FIGS/data/w2b_leg1_metrics_v2.json :: cells.{r16q1-s1,r16q1-s2,r16q1-s3,r16q1-s4,r16q1-s5}.N10000.s1.accuracy :: sha256 b56bf70be9982044 :: mean +- SD over the row's 5 pretrained models
\newcommand{\FtAccJciiEFiveOneseven}{0.5078\,\ensuremath{\pm}\,0.0011}  % experiments/FIGS/data/w2b_leg1_metrics_v2.json :: cells.{r16q1-s1,r16q1-s2,r16q1-s3,r16q1-s4,r16q1-s5}.N100000.s1.accuracy :: sha256 b56bf70be9982044 :: mean +- SD over the row's 5 pretrained models
\newcommand{\FtAccJciiESixOneseven}{0.5227\,\ensuremath{\pm}\,0.0010}  % experiments/FIGS/data/w2b_leg1_metrics_v2.json :: cells.{r16q1-s1,r16q1-s2,r16q1-s3,r16q1-s4,r16q1-s5}.N1000000.s1.accuracy :: sha256 b56bf70be9982044 :: mean +- SD over the row's 5 pretrained models
\newcommand{\FtAccJciiEThreeOnesixtwoMass}{0.451\,\ensuremath{\pm}\,0.007}  % experiments/FIGS/data/w2b_leg1_metrics_v2.json :: cells.{l162mass-s1,l162mass-s2,l162mass-s3,l162mass-s4,l162mass-s5}.N1000.s1.accuracy :: sha256 b56bf70be9982044 :: mean +- SD over the row's 5 pretrained models
\newcommand{\FtAccJciiEFourOnesixtwoMass}{0.515\,\ensuremath{\pm}\,0.004}  % experiments/FIGS/data/w2b_leg1_metrics_v2.json :: cells.{l162mass-s1,l162mass-s2,l162mass-s3,l162mass-s4,l162mass-s5}.N10000.s1.accuracy :: sha256 b56bf70be9982044 :: mean +- SD over the row's 5 pretrained models
\newcommand{\FtAccJciiEFiveOnesixtwoMass}{0.5357\,\ensuremath{\pm}\,0.0007}  % experiments/FIGS/data/w2b_leg1_metrics_v2.json :: cells.{l162mass-s1,l162mass-s2,l162mass-s3,l162mass-s4,l162mass-s5}.N100000.s1.accuracy :: sha256 b56bf70be9982044 :: mean +- SD over the row's 5 pretrained models
\newcommand{\FtAccJciiESixOnesixtwoMass}{0.5409\,\ensuremath{\pm}\,0.0010}  % experiments/FIGS/data/w2b_leg1_metrics_v2.json :: cells.{l162mass-s1,l162mass-s2,l162mass-s3,l162mass-s4,l162mass-s5}.N1000000.s1.accuracy :: sha256 b56bf70be9982044 :: mean +- SD over the row's 5 pretrained models
\newcommand{\FtAccJciiEThreeOnesevenMass}{0.331\,\ensuremath{\pm}\,0.005}  % experiments/FIGS/data/w2b_leg1_metrics_v2.json :: cells.{r16q1mass-s1,r16q1mass-s2,r16q1mass-s3,r16q1mass-s4,r16q1mass-s5}.N1000.s1.accuracy :: sha256 b56bf70be9982044 :: mean +- SD over the row's 5 pretrained models
\newcommand{\FtAccJciiEFourOnesevenMass}{0.4507\,\ensuremath{\pm}\,0.0034}  % experiments/FIGS/data/w2b_leg1_metrics_v2.json :: cells.{r16q1mass-s1,r16q1mass-s2,r16q1mass-s3,r16q1mass-s4,r16q1mass-s5}.N10000.s1.accuracy :: sha256 b56bf70be9982044 :: mean +- SD over the row's 5 pretrained models
\newcommand{\FtAccJciiEFiveOnesevenMass}{0.4993\,\ensuremath{\pm}\,0.0022}  % experiments/FIGS/data/w2b_leg1_metrics_v2.json :: cells.{r16q1mass-s1,r16q1mass-s2,r16q1mass-s3,r16q1mass-s4,r16q1mass-s5}.N100000.s1.accuracy :: sha256 b56bf70be9982044 :: mean +- SD over the row's 5 pretrained models
\newcommand{\FtAccJciiESixOnesevenMass}{0.5166\,\ensuremath{\pm}\,0.0013}  % experiments/FIGS/data/w2b_leg1_metrics_v2.json :: cells.{r16q1mass-s1,r16q1mass-s2,r16q1mass-s3,r16q1mass-s4,r16q1mass-s5}.N1000000.s1.accuracy :: sha256 b56bf70be9982044 :: mean +- SD over the row's 5 pretrained models
\newcommand{\FtAccJciiEThreeRandDrawOne}{0.424}  % experiments/FIGS/data/w2b_leg1_metrics_v2.json :: cells.{rand-d1-s1b}.N1000.s1.accuracy :: sha256 b56bf70be9982044 :: one run
\newcommand{\FtAccJciiEFourRandDrawOne}{0.5070}  % experiments/FIGS/data/w2b_leg1_metrics_v2.json :: cells.{rand-d1-s1b}.N10000.s1.accuracy :: sha256 b56bf70be9982044 :: one run
\newcommand{\FtAccJciiEFiveRandDrawOne}{0.5226}  % experiments/FIGS/data/w2b_leg1_metrics_v2.json :: cells.{rand-d1-s1b}.N100000.s1.accuracy :: sha256 b56bf70be9982044 :: one run
\newcommand{\FtAccJciiESixRandDrawOne}{0.53354}  % experiments/FIGS/data/w2b_leg1_metrics_v2.json :: cells.{rand-d1-s1b}.N1000000.s1.accuracy :: sha256 b56bf70be9982044 :: one run
\newcommand{\FtAccJciiEThreeRandDrawTwo}{0.420}  % experiments/FIGS/data/w2b_leg1_metrics_v2.json :: cells.{rand-d2-s2}.N1000.s1.accuracy :: sha256 b56bf70be9982044 :: one run
\newcommand{\FtAccJciiEFourRandDrawTwo}{0.5074}  % experiments/FIGS/data/w2b_leg1_metrics_v2.json :: cells.{rand-d2-s2}.N10000.s1.accuracy :: sha256 b56bf70be9982044 :: one run
\newcommand{\FtAccJciiEFiveRandDrawTwo}{0.5243}  % experiments/FIGS/data/w2b_leg1_metrics_v2.json :: cells.{rand-d2-s2}.N100000.s1.accuracy :: sha256 b56bf70be9982044 :: one run
\newcommand{\FtAccJciiESixRandDrawTwo}{0.53441}  % experiments/FIGS/data/w2b_leg1_metrics_v2.json :: cells.{rand-d2-s2}.N1000000.s1.accuracy :: sha256 b56bf70be9982044 :: one run
\newcommand{\FtAccJciiEThreeRandDrawThree}{0.423}  % experiments/FIGS/data/w2b_leg1_metrics_v2.json :: cells.{rand-d3-s3}.N1000.s1.accuracy :: sha256 b56bf70be9982044 :: one run
\newcommand{\FtAccJciiEFourRandDrawThree}{0.5114}  % experiments/FIGS/data/w2b_leg1_metrics_v2.json :: cells.{rand-d3-s3}.N10000.s1.accuracy :: sha256 b56bf70be9982044 :: one run
\newcommand{\FtAccJciiEFiveRandDrawThree}{0.5245}  % experiments/FIGS/data/w2b_leg1_metrics_v2.json :: cells.{rand-d3-s3}.N100000.s1.accuracy :: sha256 b56bf70be9982044 :: one run
\newcommand{\FtAccJciiESixRandDrawThree}{0.53473}  % experiments/FIGS/data/w2b_leg1_metrics_v2.json :: cells.{rand-d3-s3}.N1000000.s1.accuracy :: sha256 b56bf70be9982044 :: one run
\newcommand{\FtOmaRatioJciiEThreeOnesixtwoOverOneeighteight}{0.96}  % experiments/FIGS/data/w2b_leg1_metrics_v2.json :: cells.*.N1000.s1.macro_auc_ovr :: sha256 b56bf70be9982044 :: ratio of the seed means of 1 - macro AUC, this row over the reference
\newcommand{\FtOmaRatioJciiEThreeFourthreeOverOneeighteight}{0.84}  % experiments/FIGS/data/w2b_leg1_metrics_v2.json :: cells.*.N1000.s1.macro_auc_ovr :: sha256 b56bf70be9982044 :: ratio of the seed means of 1 - macro AUC, this row over the reference
\newcommand{\FtOmaRatioJciiEThreeOnesevenOverOneeighteight}{0.97}  % experiments/FIGS/data/w2b_leg1_metrics_v2.json :: cells.*.N1000.s1.macro_auc_ovr :: sha256 b56bf70be9982044 :: ratio of the seed means of 1 - macro AUC, this row over the reference
\newcommand{\FtOmaRatioJciiEThreeOnesixtwoMassOverOneeighteight}{0.81}  % experiments/FIGS/data/w2b_leg1_metrics_v2.json :: cells.*.N1000.s1.macro_auc_ovr :: sha256 b56bf70be9982044 :: ratio of the seed means of 1 - macro AUC, this row over the reference
\newcommand{\FtOmaRatioJciiEThreeOnesevenMassOverOneeighteight}{0.98}  % experiments/FIGS/data/w2b_leg1_metrics_v2.json :: cells.*.N1000.s1.macro_auc_ovr :: sha256 b56bf70be9982044 :: ratio of the seed means of 1 - macro AUC, this row over the reference
\newcommand{\FtOmaRatioJciiEThreeRandDrawOneOverOneeighteight}{0.80}  % experiments/FIGS/data/w2b_leg1_metrics_v2.json :: cells.*.N1000.s1.macro_auc_ovr :: sha256 b56bf70be9982044 :: ratio of the seed means of 1 - macro AUC, this row over the reference
\newcommand{\FtOmaRatioJciiEThreeRandDrawTwoOverOneeighteight}{0.87}  % experiments/FIGS/data/w2b_leg1_metrics_v2.json :: cells.*.N1000.s1.macro_auc_ovr :: sha256 b56bf70be9982044 :: ratio of the seed means of 1 - macro AUC, this row over the reference
\newcommand{\FtOmaRatioJciiEThreeRandDrawThreeOverOneeighteight}{0.83}  % experiments/FIGS/data/w2b_leg1_metrics_v2.json :: cells.*.N1000.s1.macro_auc_ovr :: sha256 b56bf70be9982044 :: ratio of the seed means of 1 - macro AUC, this row over the reference
\newcommand{\FtOmaRatioJciiEThreeOneeighteightOverOnesixtwo}{1.05}  % experiments/FIGS/data/w2b_leg1_metrics_v2.json :: cells.*.N1000.s1.macro_auc_ovr :: sha256 b56bf70be9982044 :: ratio of the seed means of 1 - macro AUC, this row over the reference
\newcommand{\FtOmaRatioJciiEThreeFourthreeOverOnesixtwo}{0.88}  % experiments/FIGS/data/w2b_leg1_metrics_v2.json :: cells.*.N1000.s1.macro_auc_ovr :: sha256 b56bf70be9982044 :: ratio of the seed means of 1 - macro AUC, this row over the reference
\newcommand{\FtOmaRatioJciiEThreeOnesevenOverOnesixtwo}{1.02}  % experiments/FIGS/data/w2b_leg1_metrics_v2.json :: cells.*.N1000.s1.macro_auc_ovr :: sha256 b56bf70be9982044 :: ratio of the seed means of 1 - macro AUC, this row over the reference
\newcommand{\FtOmaRatioJciiEThreeOnesixtwoMassOverOnesixtwo}{0.85}  % experiments/FIGS/data/w2b_leg1_metrics_v2.json :: cells.*.N1000.s1.macro_auc_ovr :: sha256 b56bf70be9982044 :: ratio of the seed means of 1 - macro AUC, this row over the reference
\newcommand{\FtOmaRatioJciiEThreeOnesevenMassOverOnesixtwo}{1.02}  % experiments/FIGS/data/w2b_leg1_metrics_v2.json :: cells.*.N1000.s1.macro_auc_ovr :: sha256 b56bf70be9982044 :: ratio of the seed means of 1 - macro AUC, this row over the reference
\newcommand{\FtOmaRatioJciiEThreeRandDrawOneOverOnesixtwo}{0.84}  % experiments/FIGS/data/w2b_leg1_metrics_v2.json :: cells.*.N1000.s1.macro_auc_ovr :: sha256 b56bf70be9982044 :: ratio of the seed means of 1 - macro AUC, this row over the reference
\newcommand{\FtOmaRatioJciiEThreeRandDrawTwoOverOnesixtwo}{0.91}  % experiments/FIGS/data/w2b_leg1_metrics_v2.json :: cells.*.N1000.s1.macro_auc_ovr :: sha256 b56bf70be9982044 :: ratio of the seed means of 1 - macro AUC, this row over the reference
\newcommand{\FtOmaRatioJciiEThreeRandDrawThreeOverOnesixtwo}{0.87}  % experiments/FIGS/data/w2b_leg1_metrics_v2.json :: cells.*.N1000.s1.macro_auc_ovr :: sha256 b56bf70be9982044 :: ratio of the seed means of 1 - macro AUC, this row over the reference
\newcommand{\FtOmaRatioJciiEFourOnesixtwoOverOneeighteight}{0.99}  % experiments/FIGS/data/w2b_leg1_metrics_v2.json :: cells.*.N10000.s1.macro_auc_ovr :: sha256 b56bf70be9982044 :: ratio of the seed means of 1 - macro AUC, this row over the reference
\newcommand{\FtOmaRatioJciiEFourFourthreeOverOneeighteight}{1.20}  % experiments/FIGS/data/w2b_leg1_metrics_v2.json :: cells.*.N10000.s1.macro_auc_ovr :: sha256 b56bf70be9982044 :: ratio of the seed means of 1 - macro AUC, this row over the reference
\newcommand{\FtOmaRatioJciiEFourOnesevenOverOneeighteight}{1.32}  % experiments/FIGS/data/w2b_leg1_metrics_v2.json :: cells.*.N10000.s1.macro_auc_ovr :: sha256 b56bf70be9982044 :: ratio of the seed means of 1 - macro AUC, this row over the reference
\newcommand{\FtOmaRatioJciiEFourOnesixtwoMassOverOneeighteight}{0.99}  % experiments/FIGS/data/w2b_leg1_metrics_v2.json :: cells.*.N10000.s1.macro_auc_ovr :: sha256 b56bf70be9982044 :: ratio of the seed means of 1 - macro AUC, this row over the reference
\newcommand{\FtOmaRatioJciiEFourOnesevenMassOverOneeighteight}{1.30}  % experiments/FIGS/data/w2b_leg1_metrics_v2.json :: cells.*.N10000.s1.macro_auc_ovr :: sha256 b56bf70be9982044 :: ratio of the seed means of 1 - macro AUC, this row over the reference
\newcommand{\FtOmaRatioJciiEFourRandDrawOneOverOneeighteight}{1.17}  % experiments/FIGS/data/w2b_leg1_metrics_v2.json :: cells.*.N10000.s1.macro_auc_ovr :: sha256 b56bf70be9982044 :: ratio of the seed means of 1 - macro AUC, this row over the reference
\newcommand{\FtOmaRatioJciiEFourRandDrawTwoOverOneeighteight}{1.17}  % experiments/FIGS/data/w2b_leg1_metrics_v2.json :: cells.*.N10000.s1.macro_auc_ovr :: sha256 b56bf70be9982044 :: ratio of the seed means of 1 - macro AUC, this row over the reference
\newcommand{\FtOmaRatioJciiEFourRandDrawThreeOverOneeighteight}{1.13}  % experiments/FIGS/data/w2b_leg1_metrics_v2.json :: cells.*.N10000.s1.macro_auc_ovr :: sha256 b56bf70be9982044 :: ratio of the seed means of 1 - macro AUC, this row over the reference
\newcommand{\FtOmaRatioJciiEFourOneeighteightOverOnesixtwo}{1.01}  % experiments/FIGS/data/w2b_leg1_metrics_v2.json :: cells.*.N10000.s1.macro_auc_ovr :: sha256 b56bf70be9982044 :: ratio of the seed means of 1 - macro AUC, this row over the reference
\newcommand{\FtOmaRatioJciiEFourFourthreeOverOnesixtwo}{1.21}  % experiments/FIGS/data/w2b_leg1_metrics_v2.json :: cells.*.N10000.s1.macro_auc_ovr :: sha256 b56bf70be9982044 :: ratio of the seed means of 1 - macro AUC, this row over the reference
\newcommand{\FtOmaRatioJciiEFourOnesevenOverOnesixtwo}{1.32}  % experiments/FIGS/data/w2b_leg1_metrics_v2.json :: cells.*.N10000.s1.macro_auc_ovr :: sha256 b56bf70be9982044 :: ratio of the seed means of 1 - macro AUC, this row over the reference
\newcommand{\FtOmaRatioJciiEFourOnesixtwoMassOverOnesixtwo}{1.00}  % experiments/FIGS/data/w2b_leg1_metrics_v2.json :: cells.*.N10000.s1.macro_auc_ovr :: sha256 b56bf70be9982044 :: ratio of the seed means of 1 - macro AUC, this row over the reference
\newcommand{\FtOmaRatioJciiEFourOnesevenMassOverOnesixtwo}{1.31}  % experiments/FIGS/data/w2b_leg1_metrics_v2.json :: cells.*.N10000.s1.macro_auc_ovr :: sha256 b56bf70be9982044 :: ratio of the seed means of 1 - macro AUC, this row over the reference
\newcommand{\FtOmaRatioJciiEFourRandDrawOneOverOnesixtwo}{1.18}  % experiments/FIGS/data/w2b_leg1_metrics_v2.json :: cells.*.N10000.s1.macro_auc_ovr :: sha256 b56bf70be9982044 :: ratio of the seed means of 1 - macro AUC, this row over the reference
\newcommand{\FtOmaRatioJciiEFourRandDrawTwoOverOnesixtwo}{1.18}  % experiments/FIGS/data/w2b_leg1_metrics_v2.json :: cells.*.N10000.s1.macro_auc_ovr :: sha256 b56bf70be9982044 :: ratio of the seed means of 1 - macro AUC, this row over the reference
\newcommand{\FtOmaRatioJciiEFourRandDrawThreeOverOnesixtwo}{1.14}  % experiments/FIGS/data/w2b_leg1_metrics_v2.json :: cells.*.N10000.s1.macro_auc_ovr :: sha256 b56bf70be9982044 :: ratio of the seed means of 1 - macro AUC, this row over the reference
\newcommand{\FtOmaRatioJciiEFiveOnesixtwoOverOneeighteight}{1.00}  % experiments/FIGS/data/w2b_leg1_metrics_v2.json :: cells.*.N100000.s1.macro_auc_ovr :: sha256 b56bf70be9982044 :: ratio of the seed means of 1 - macro AUC, this row over the reference
\newcommand{\FtOmaRatioJciiEFiveFourthreeOverOneeighteight}{1.09}  % experiments/FIGS/data/w2b_leg1_metrics_v2.json :: cells.*.N100000.s1.macro_auc_ovr :: sha256 b56bf70be9982044 :: ratio of the seed means of 1 - macro AUC, this row over the reference
\newcommand{\FtOmaRatioJciiEFiveOnesevenOverOneeighteight}{1.22}  % experiments/FIGS/data/w2b_leg1_metrics_v2.json :: cells.*.N100000.s1.macro_auc_ovr :: sha256 b56bf70be9982044 :: ratio of the seed means of 1 - macro AUC, this row over the reference
\newcommand{\FtOmaRatioJciiEFiveOnesixtwoMassOverOneeighteight}{1.01}  % experiments/FIGS/data/w2b_leg1_metrics_v2.json :: cells.*.N100000.s1.macro_auc_ovr :: sha256 b56bf70be9982044 :: ratio of the seed means of 1 - macro AUC, this row over the reference
\newcommand{\FtOmaRatioJciiEFiveOnesevenMassOverOneeighteight}{1.26}  % experiments/FIGS/data/w2b_leg1_metrics_v2.json :: cells.*.N100000.s1.macro_auc_ovr :: sha256 b56bf70be9982044 :: ratio of the seed means of 1 - macro AUC, this row over the reference
\newcommand{\FtOmaRatioJciiEFiveRandDrawOneOverOneeighteight}{1.15}  % experiments/FIGS/data/w2b_leg1_metrics_v2.json :: cells.*.N100000.s1.macro_auc_ovr :: sha256 b56bf70be9982044 :: ratio of the seed means of 1 - macro AUC, this row over the reference
\newcommand{\FtOmaRatioJciiEFiveRandDrawTwoOverOneeighteight}{1.12}  % experiments/FIGS/data/w2b_leg1_metrics_v2.json :: cells.*.N100000.s1.macro_auc_ovr :: sha256 b56bf70be9982044 :: ratio of the seed means of 1 - macro AUC, this row over the reference
\newcommand{\FtOmaRatioJciiEFiveRandDrawThreeOverOneeighteight}{1.15}  % experiments/FIGS/data/w2b_leg1_metrics_v2.json :: cells.*.N100000.s1.macro_auc_ovr :: sha256 b56bf70be9982044 :: ratio of the seed means of 1 - macro AUC, this row over the reference
\newcommand{\FtOmaRatioJciiEFiveOneeighteightOverOnesixtwo}{1.00}  % experiments/FIGS/data/w2b_leg1_metrics_v2.json :: cells.*.N100000.s1.macro_auc_ovr :: sha256 b56bf70be9982044 :: ratio of the seed means of 1 - macro AUC, this row over the reference
\newcommand{\FtOmaRatioJciiEFiveFourthreeOverOnesixtwo}{1.09}  % experiments/FIGS/data/w2b_leg1_metrics_v2.json :: cells.*.N100000.s1.macro_auc_ovr :: sha256 b56bf70be9982044 :: ratio of the seed means of 1 - macro AUC, this row over the reference
\newcommand{\FtOmaRatioJciiEFiveOnesevenOverOnesixtwo}{1.22}  % experiments/FIGS/data/w2b_leg1_metrics_v2.json :: cells.*.N100000.s1.macro_auc_ovr :: sha256 b56bf70be9982044 :: ratio of the seed means of 1 - macro AUC, this row over the reference
\newcommand{\FtOmaRatioJciiEFiveOnesixtwoMassOverOnesixtwo}{1.02}  % experiments/FIGS/data/w2b_leg1_metrics_v2.json :: cells.*.N100000.s1.macro_auc_ovr :: sha256 b56bf70be9982044 :: ratio of the seed means of 1 - macro AUC, this row over the reference
\newcommand{\FtOmaRatioJciiEFiveOnesevenMassOverOnesixtwo}{1.26}  % experiments/FIGS/data/w2b_leg1_metrics_v2.json :: cells.*.N100000.s1.macro_auc_ovr :: sha256 b56bf70be9982044 :: ratio of the seed means of 1 - macro AUC, this row over the reference
\newcommand{\FtOmaRatioJciiEFiveRandDrawOneOverOnesixtwo}{1.15}  % experiments/FIGS/data/w2b_leg1_metrics_v2.json :: cells.*.N100000.s1.macro_auc_ovr :: sha256 b56bf70be9982044 :: ratio of the seed means of 1 - macro AUC, this row over the reference
\newcommand{\FtOmaRatioJciiEFiveRandDrawTwoOverOnesixtwo}{1.13}  % experiments/FIGS/data/w2b_leg1_metrics_v2.json :: cells.*.N100000.s1.macro_auc_ovr :: sha256 b56bf70be9982044 :: ratio of the seed means of 1 - macro AUC, this row over the reference
\newcommand{\FtOmaRatioJciiEFiveRandDrawThreeOverOnesixtwo}{1.16}  % experiments/FIGS/data/w2b_leg1_metrics_v2.json :: cells.*.N100000.s1.macro_auc_ovr :: sha256 b56bf70be9982044 :: ratio of the seed means of 1 - macro AUC, this row over the reference
\newcommand{\FtOmaRatioJciiESixOnesixtwoOverOneeighteight}{1.00}  % experiments/FIGS/data/w2b_leg1_metrics_v2.json :: cells.*.N1000000.s1.macro_auc_ovr :: sha256 b56bf70be9982044 :: ratio of the seed means of 1 - macro AUC, this row over the reference
\newcommand{\FtOmaRatioJciiESixFourthreeOverOneeighteight}{1.05}  % experiments/FIGS/data/w2b_leg1_metrics_v2.json :: cells.*.N1000000.s1.macro_auc_ovr :: sha256 b56bf70be9982044 :: ratio of the seed means of 1 - macro AUC, this row over the reference
\newcommand{\FtOmaRatioJciiESixOnesevenOverOneeighteight}{1.14}  % experiments/FIGS/data/w2b_leg1_metrics_v2.json :: cells.*.N1000000.s1.macro_auc_ovr :: sha256 b56bf70be9982044 :: ratio of the seed means of 1 - macro AUC, this row over the reference
\newcommand{\FtOmaRatioJciiESixOnesixtwoMassOverOneeighteight}{1.02}  % experiments/FIGS/data/w2b_leg1_metrics_v2.json :: cells.*.N1000000.s1.macro_auc_ovr :: sha256 b56bf70be9982044 :: ratio of the seed means of 1 - macro AUC, this row over the reference
\newcommand{\FtOmaRatioJciiESixOnesevenMassOverOneeighteight}{1.18}  % experiments/FIGS/data/w2b_leg1_metrics_v2.json :: cells.*.N1000000.s1.macro_auc_ovr :: sha256 b56bf70be9982044 :: ratio of the seed means of 1 - macro AUC, this row over the reference
\newcommand{\FtOmaRatioJciiESixRandDrawOneOverOneeighteight}{1.08}  % experiments/FIGS/data/w2b_leg1_metrics_v2.json :: cells.*.N1000000.s1.macro_auc_ovr :: sha256 b56bf70be9982044 :: ratio of the seed means of 1 - macro AUC, this row over the reference
\newcommand{\FtOmaRatioJciiESixRandDrawTwoOverOneeighteight}{1.07}  % experiments/FIGS/data/w2b_leg1_metrics_v2.json :: cells.*.N1000000.s1.macro_auc_ovr :: sha256 b56bf70be9982044 :: ratio of the seed means of 1 - macro AUC, this row over the reference
\newcommand{\FtOmaRatioJciiESixRandDrawThreeOverOneeighteight}{1.07}  % experiments/FIGS/data/w2b_leg1_metrics_v2.json :: cells.*.N1000000.s1.macro_auc_ovr :: sha256 b56bf70be9982044 :: ratio of the seed means of 1 - macro AUC, this row over the reference
\newcommand{\FtOmaRatioJciiESixOneeighteightOverOnesixtwo}{1.00}  % experiments/FIGS/data/w2b_leg1_metrics_v2.json :: cells.*.N1000000.s1.macro_auc_ovr :: sha256 b56bf70be9982044 :: ratio of the seed means of 1 - macro AUC, this row over the reference
\newcommand{\FtOmaRatioJciiESixFourthreeOverOnesixtwo}{1.04}  % experiments/FIGS/data/w2b_leg1_metrics_v2.json :: cells.*.N1000000.s1.macro_auc_ovr :: sha256 b56bf70be9982044 :: ratio of the seed means of 1 - macro AUC, this row over the reference
\newcommand{\FtOmaRatioJciiESixOnesevenOverOnesixtwo}{1.14}  % experiments/FIGS/data/w2b_leg1_metrics_v2.json :: cells.*.N1000000.s1.macro_auc_ovr :: sha256 b56bf70be9982044 :: ratio of the seed means of 1 - macro AUC, this row over the reference
\newcommand{\FtOmaRatioJciiESixOnesixtwoMassOverOnesixtwo}{1.02}  % experiments/FIGS/data/w2b_leg1_metrics_v2.json :: cells.*.N1000000.s1.macro_auc_ovr :: sha256 b56bf70be9982044 :: ratio of the seed means of 1 - macro AUC, this row over the reference
\newcommand{\FtOmaRatioJciiESixOnesevenMassOverOnesixtwo}{1.18}  % experiments/FIGS/data/w2b_leg1_metrics_v2.json :: cells.*.N1000000.s1.macro_auc_ovr :: sha256 b56bf70be9982044 :: ratio of the seed means of 1 - macro AUC, this row over the reference
\newcommand{\FtOmaRatioJciiESixRandDrawOneOverOnesixtwo}{1.07}  % experiments/FIGS/data/w2b_leg1_metrics_v2.json :: cells.*.N1000000.s1.macro_auc_ovr :: sha256 b56bf70be9982044 :: ratio of the seed means of 1 - macro AUC, this row over the reference
\newcommand{\FtOmaRatioJciiESixRandDrawTwoOverOnesixtwo}{1.07}  % experiments/FIGS/data/w2b_leg1_metrics_v2.json :: cells.*.N1000000.s1.macro_auc_ovr :: sha256 b56bf70be9982044 :: ratio of the seed means of 1 - macro AUC, this row over the reference
\newcommand{\FtOmaRatioJciiESixRandDrawThreeOverOnesixtwo}{1.07}  % experiments/FIGS/data/w2b_leg1_metrics_v2.json :: cells.*.N1000000.s1.macro_auc_ovr :: sha256 b56bf70be9982044 :: ratio of the seed means of 1 - macro AUC, this row over the reference
\newcommand{\PairedFtAccJciiEThreeOnesevenMinusOneeighteight}{\ensuremath{-}0.1132~[\ensuremath{-}0.1214, \ensuremath{-}0.1051]}  % experiments/FIGS/data/w2b_leg1_metrics_v2.json :: cells.{r16q1-s1,r16q1-s2,r16q1-s3,r16q1-s4,r16q1-s5,l188-s1,l188-s2,l188-s3,l188-s4,l188-s5}.N1000.s1.accuracy :: sha256 b56bf70be9982044 :: accuracy, coarse minus fine, paired by run [95% Student-t interval]
\newcommand{\PairedFtAccJciiEFourOnesevenMinusOneeighteight}{\ensuremath{-}0.0730~[\ensuremath{-}0.0902, \ensuremath{-}0.0558]}  % experiments/FIGS/data/w2b_leg1_metrics_v2.json :: cells.{r16q1-s1,r16q1-s2,r16q1-s3,r16q1-s4,r16q1-s5,l188-s1,l188-s2,l188-s3,l188-s4,l188-s5}.N10000.s1.accuracy :: sha256 b56bf70be9982044 :: accuracy, coarse minus fine, paired by run [95% Student-t interval]
\newcommand{\PairedFtAccJciiEFiveOnesevenMinusOneeighteight}{\ensuremath{-}0.0301~[\ensuremath{-}0.0319, \ensuremath{-}0.0284]}  % experiments/FIGS/data/w2b_leg1_metrics_v2.json :: cells.{r16q1-s1,r16q1-s2,r16q1-s3,r16q1-s4,r16q1-s5,l188-s1,l188-s2,l188-s3,l188-s4,l188-s5}.N100000.s1.accuracy :: sha256 b56bf70be9982044 :: accuracy, coarse minus fine, paired by run [95% Student-t interval]
\newcommand{\PairedFtAccJciiESixOnesevenMinusOneeighteight}{\ensuremath{-}0.0216~[\ensuremath{-}0.0232, \ensuremath{-}0.0199]}  % experiments/FIGS/data/w2b_leg1_metrics_v2.json :: cells.{r16q1-s1,r16q1-s2,r16q1-s3,r16q1-s4,r16q1-s5,l188-s1,l188-s2,l188-s3,l188-s4,l188-s5}.N1000000.s1.accuracy :: sha256 b56bf70be9982044 :: accuracy, coarse minus fine, paired by run [95% Student-t interval]
\newcommand{\FtNClassesJcii}{162}  % experiments/FIGS/data/w2b_leg1_metrics_v2.json :: cells.*.*.s1.n_classes_present :: sha256 b56bf70be9982044 :: classes of the fine-tuning task
\newcommand{\FtNTestAucJcii}{125{,}000}  % experiments/FIGS/data/w2b_leg1_metrics_v2.json :: cells.*.*.s1.n_jets_auc :: sha256 b56bf70be9982044 :: test jets the macro AUC is computed on
\newcommand{\FtNTestAccJcii}{500{,}000}  % experiments/FIGS/data/w2b_leg1_metrics_v2.json :: cells.*.*.s1.n_jets :: sha256 b56bf70be9982044 :: test jets the accuracy is computed on
\newcommand{\FtAucJcEThreeOneeighteight}{0.950\,\ensuremath{\pm}\,0.004}  % experiments/FIGS/data/w2b_leg2_metrics_v2.json :: cells.{l188-s1,l188-s2,l188-s3,l188-s4,l188-s5}.N1000.s1.macro_auc_ovr :: sha256 9cd059d3e2f51798 :: mean +- SD over the row's 5 pretrained models
\newcommand{\FtAucJcEFourOneeighteight}{0.9765\,\ensuremath{\pm}\,0.0008}  % experiments/FIGS/data/w2b_leg2_metrics_v2.json :: cells.{l188-s1,l188-s2,l188-s3,l188-s4,l188-s5}.N10000.s1.macro_auc_ovr :: sha256 9cd059d3e2f51798 :: mean +- SD over the row's 5 pretrained models
\newcommand{\FtAucJcEFiveOneeighteight}{0.9833\,\ensuremath{\pm}\,0.0002}  % experiments/FIGS/data/w2b_leg2_metrics_v2.json :: cells.{l188-s1,l188-s2,l188-s3,l188-s4,l188-s5}.N100000.s1.macro_auc_ovr :: sha256 9cd059d3e2f51798 :: mean +- SD over the row's 5 pretrained models
\newcommand{\FtAucJcESixOneeighteight}{0.9858\,\ensuremath{\pm}\,0.0001}  % experiments/FIGS/data/w2b_leg2_metrics_v2.json :: cells.{l188-s1,l188-s2,l188-s3,l188-s4,l188-s5}.N1000000.s1.macro_auc_ovr :: sha256 9cd059d3e2f51798 :: mean +- SD over the row's 5 pretrained models
\newcommand{\FtAucJcEThreeOnesixtwo}{0.9526\,\ensuremath{\pm}\,0.0020}  % experiments/FIGS/data/w2b_leg2_metrics_v2.json :: cells.{l162-s1b,l162-s2,l162-s3,l162-s4,l162-s5}.N1000.s1.macro_auc_ovr :: sha256 9cd059d3e2f51798 :: mean +- SD over the row's 5 pretrained models
\newcommand{\FtAucJcEFourOnesixtwo}{0.9770\,\ensuremath{\pm}\,0.0007}  % experiments/FIGS/data/w2b_leg2_metrics_v2.json :: cells.{l162-s1b,l162-s2,l162-s3,l162-s4,l162-s5}.N10000.s1.macro_auc_ovr :: sha256 9cd059d3e2f51798 :: mean +- SD over the row's 5 pretrained models
\newcommand{\FtAucJcEFiveOnesixtwo}{0.9836\,\ensuremath{\pm}\,0.0001}  % experiments/FIGS/data/w2b_leg2_metrics_v2.json :: cells.{l162-s1b,l162-s2,l162-s3,l162-s4,l162-s5}.N100000.s1.macro_auc_ovr :: sha256 9cd059d3e2f51798 :: mean +- SD over the row's 5 pretrained models
\newcommand{\FtAucJcESixOnesixtwo}{0.9858\,\ensuremath{\pm}\,0.0001}  % experiments/FIGS/data/w2b_leg2_metrics_v2.json :: cells.{l162-s1b,l162-s2,l162-s3,l162-s4,l162-s5}.N1000000.s1.macro_auc_ovr :: sha256 9cd059d3e2f51798 :: mean +- SD over the row's 5 pretrained models
\newcommand{\FtAucJcEThreeFourthree}{0.950\,\ensuremath{\pm}\,0.004}  % experiments/FIGS/data/w2b_leg2_metrics_v2.json :: cells.{r42q1-s1,r42q1-s2,r42q1-s3,r42q1-s4,r42q1-s5}.N1000.s1.macro_auc_ovr :: sha256 9cd059d3e2f51798 :: mean +- SD over the row's 5 pretrained models
\newcommand{\FtAucJcEFourFourthree}{0.9769\,\ensuremath{\pm}\,0.0005}  % experiments/FIGS/data/w2b_leg2_metrics_v2.json :: cells.{r42q1-s1,r42q1-s2,r42q1-s3,r42q1-s4,r42q1-s5}.N10000.s1.macro_auc_ovr :: sha256 9cd059d3e2f51798 :: mean +- SD over the row's 5 pretrained models
\newcommand{\FtAucJcEFiveFourthree}{0.9831\,\ensuremath{\pm}\,0.0002}  % experiments/FIGS/data/w2b_leg2_metrics_v2.json :: cells.{r42q1-s1,r42q1-s2,r42q1-s3,r42q1-s4,r42q1-s5}.N100000.s1.macro_auc_ovr :: sha256 9cd059d3e2f51798 :: mean +- SD over the row's 5 pretrained models
\newcommand{\FtAucJcESixFourthree}{0.9855\,\ensuremath{\pm}\,0.0001}  % experiments/FIGS/data/w2b_leg2_metrics_v2.json :: cells.{r42q1-s1,r42q1-s2,r42q1-s3,r42q1-s4,r42q1-s5}.N1000000.s1.macro_auc_ovr :: sha256 9cd059d3e2f51798 :: mean +- SD over the row's 5 pretrained models
\newcommand{\FtAucJcEThreeOneseven}{0.9470\,\ensuremath{\pm}\,0.0024}  % experiments/FIGS/data/w2b_leg2_metrics_v2.json :: cells.{r16q1-s1,r16q1-s2,r16q1-s3,r16q1-s4,r16q1-s5}.N1000.s1.macro_auc_ovr :: sha256 9cd059d3e2f51798 :: mean +- SD over the row's 5 pretrained models
\newcommand{\FtAucJcEFourOneseven}{0.9749\,\ensuremath{\pm}\,0.0006}  % experiments/FIGS/data/w2b_leg2_metrics_v2.json :: cells.{r16q1-s1,r16q1-s2,r16q1-s3,r16q1-s4,r16q1-s5}.N10000.s1.macro_auc_ovr :: sha256 9cd059d3e2f51798 :: mean +- SD over the row's 5 pretrained models
\newcommand{\FtAucJcEFiveOneseven}{0.9821\,\ensuremath{\pm}\,\ensuremath{<}0.0001}  % experiments/FIGS/data/w2b_leg2_metrics_v2.json :: cells.{r16q1-s1,r16q1-s2,r16q1-s3,r16q1-s4,r16q1-s5}.N100000.s1.macro_auc_ovr :: sha256 9cd059d3e2f51798 :: mean +- SD over the row's 5 pretrained models
\newcommand{\FtAucJcESixOneseven}{0.9846\,\ensuremath{\pm}\,0.0001}  % experiments/FIGS/data/w2b_leg2_metrics_v2.json :: cells.{r16q1-s1,r16q1-s2,r16q1-s3,r16q1-s4,r16q1-s5}.N1000000.s1.macro_auc_ovr :: sha256 9cd059d3e2f51798 :: mean +- SD over the row's 5 pretrained models
\newcommand{\FtAucJcEThreeOnesixtwoMass}{0.9569\,\ensuremath{\pm}\,0.0016}  % experiments/FIGS/data/w2b_leg2_metrics_v2.json :: cells.{l162mass-s1,l162mass-s2,l162mass-s3,l162mass-s4,l162mass-s5}.N1000.s1.macro_auc_ovr :: sha256 9cd059d3e2f51798 :: mean +- SD over the row's 5 pretrained models
\newcommand{\FtAucJcEFourOnesixtwoMass}{0.9768\,\ensuremath{\pm}\,0.0005}  % experiments/FIGS/data/w2b_leg2_metrics_v2.json :: cells.{l162mass-s1,l162mass-s2,l162mass-s3,l162mass-s4,l162mass-s5}.N10000.s1.macro_auc_ovr :: sha256 9cd059d3e2f51798 :: mean +- SD over the row's 5 pretrained models
\newcommand{\FtAucJcEFiveOnesixtwoMass}{0.9834\,\ensuremath{\pm}\,0.0001}  % experiments/FIGS/data/w2b_leg2_metrics_v2.json :: cells.{l162mass-s1,l162mass-s2,l162mass-s3,l162mass-s4,l162mass-s5}.N100000.s1.macro_auc_ovr :: sha256 9cd059d3e2f51798 :: mean +- SD over the row's 5 pretrained models
\newcommand{\FtAucJcESixOnesixtwoMass}{0.9856\,\ensuremath{\pm}\,0.0001}  % experiments/FIGS/data/w2b_leg2_metrics_v2.json :: cells.{l162mass-s1,l162mass-s2,l162mass-s3,l162mass-s4,l162mass-s5}.N1000000.s1.macro_auc_ovr :: sha256 9cd059d3e2f51798 :: mean +- SD over the row's 5 pretrained models
\newcommand{\FtAucJcEThreeOnesevenMass}{0.9460\,\ensuremath{\pm}\,0.0029}  % experiments/FIGS/data/w2b_leg2_metrics_v2.json :: cells.{r16q1mass-s1,r16q1mass-s2,r16q1mass-s3,r16q1mass-s4,r16q1mass-s5}.N1000.s1.macro_auc_ovr :: sha256 9cd059d3e2f51798 :: mean +- SD over the row's 5 pretrained models
\newcommand{\FtAucJcEFourOnesevenMass}{0.9737\,\ensuremath{\pm}\,0.0008}  % experiments/FIGS/data/w2b_leg2_metrics_v2.json :: cells.{r16q1mass-s1,r16q1mass-s2,r16q1mass-s3,r16q1mass-s4,r16q1mass-s5}.N10000.s1.macro_auc_ovr :: sha256 9cd059d3e2f51798 :: mean +- SD over the row's 5 pretrained models
\newcommand{\FtAucJcEFiveOnesevenMass}{0.9817\,\ensuremath{\pm}\,0.0002}  % experiments/FIGS/data/w2b_leg2_metrics_v2.json :: cells.{r16q1mass-s1,r16q1mass-s2,r16q1mass-s3,r16q1mass-s4,r16q1mass-s5}.N100000.s1.macro_auc_ovr :: sha256 9cd059d3e2f51798 :: mean +- SD over the row's 5 pretrained models
\newcommand{\FtAucJcESixOnesevenMass}{0.9842\,\ensuremath{\pm}\,0.0002}  % experiments/FIGS/data/w2b_leg2_metrics_v2.json :: cells.{r16q1mass-s1,r16q1mass-s2,r16q1mass-s3,r16q1mass-s4,r16q1mass-s5}.N1000000.s1.macro_auc_ovr :: sha256 9cd059d3e2f51798 :: mean +- SD over the row's 5 pretrained models
\newcommand{\FtAucJcEThreeRandDrawOne}{0.9477}  % experiments/FIGS/data/w2b_leg2_metrics_v2.json :: cells.{rand-d1-s1b}.N1000.s1.macro_auc_ovr :: sha256 9cd059d3e2f51798 :: one run
\newcommand{\FtAucJcEFourRandDrawOne}{0.9745}  % experiments/FIGS/data/w2b_leg2_metrics_v2.json :: cells.{rand-d1-s1b}.N10000.s1.macro_auc_ovr :: sha256 9cd059d3e2f51798 :: one run
\newcommand{\FtAucJcEFiveRandDrawOne}{0.98265}  % experiments/FIGS/data/w2b_leg2_metrics_v2.json :: cells.{rand-d1-s1b}.N100000.s1.macro_auc_ovr :: sha256 9cd059d3e2f51798 :: one run
\newcommand{\FtAucJcESixRandDrawOne}{0.98513}  % experiments/FIGS/data/w2b_leg2_metrics_v2.json :: cells.{rand-d1-s1b}.N1000000.s1.macro_auc_ovr :: sha256 9cd059d3e2f51798 :: one run
\newcommand{\FtAucJcEThreeRandDrawTwo}{0.9501}  % experiments/FIGS/data/w2b_leg2_metrics_v2.json :: cells.{rand-d2-s2}.N1000.s1.macro_auc_ovr :: sha256 9cd059d3e2f51798 :: one run
\newcommand{\FtAucJcEFourRandDrawTwo}{0.9756}  % experiments/FIGS/data/w2b_leg2_metrics_v2.json :: cells.{rand-d2-s2}.N10000.s1.macro_auc_ovr :: sha256 9cd059d3e2f51798 :: one run
\newcommand{\FtAucJcEFiveRandDrawTwo}{0.98278}  % experiments/FIGS/data/w2b_leg2_metrics_v2.json :: cells.{rand-d2-s2}.N100000.s1.macro_auc_ovr :: sha256 9cd059d3e2f51798 :: one run
\newcommand{\FtAucJcESixRandDrawTwo}{0.98512}  % experiments/FIGS/data/w2b_leg2_metrics_v2.json :: cells.{rand-d2-s2}.N1000000.s1.macro_auc_ovr :: sha256 9cd059d3e2f51798 :: one run
\newcommand{\FtAucJcEThreeRandDrawThree}{0.9470}  % experiments/FIGS/data/w2b_leg2_metrics_v2.json :: cells.{rand-d3-s3}.N1000.s1.macro_auc_ovr :: sha256 9cd059d3e2f51798 :: one run
\newcommand{\FtAucJcEFourRandDrawThree}{0.9743}  % experiments/FIGS/data/w2b_leg2_metrics_v2.json :: cells.{rand-d3-s3}.N10000.s1.macro_auc_ovr :: sha256 9cd059d3e2f51798 :: one run
\newcommand{\FtAucJcEFiveRandDrawThree}{0.98253}  % experiments/FIGS/data/w2b_leg2_metrics_v2.json :: cells.{rand-d3-s3}.N100000.s1.macro_auc_ovr :: sha256 9cd059d3e2f51798 :: one run
\newcommand{\FtAucJcESixRandDrawThree}{0.98506}  % experiments/FIGS/data/w2b_leg2_metrics_v2.json :: cells.{rand-d3-s3}.N1000000.s1.macro_auc_ovr :: sha256 9cd059d3e2f51798 :: one run
\newcommand{\FtAccJcEThreeOneeighteight}{0.7183\,\ensuremath{\pm}\,0.0032}  % experiments/FIGS/data/w2b_leg2_metrics_v2.json :: cells.{l188-s1,l188-s2,l188-s3,l188-s4,l188-s5}.N1000.s1.accuracy :: sha256 9cd059d3e2f51798 :: mean +- SD over the row's 5 pretrained models
\newcommand{\FtAccJcEFourOneeighteight}{0.8039\,\ensuremath{\pm}\,0.0030}  % experiments/FIGS/data/w2b_leg2_metrics_v2.json :: cells.{l188-s1,l188-s2,l188-s3,l188-s4,l188-s5}.N10000.s1.accuracy :: sha256 9cd059d3e2f51798 :: mean +- SD over the row's 5 pretrained models
\newcommand{\FtAccJcEFiveOneeighteight}{0.8351\,\ensuremath{\pm}\,0.0008}  % experiments/FIGS/data/w2b_leg2_metrics_v2.json :: cells.{l188-s1,l188-s2,l188-s3,l188-s4,l188-s5}.N100000.s1.accuracy :: sha256 9cd059d3e2f51798 :: mean +- SD over the row's 5 pretrained models
\newcommand{\FtAccJcESixOneeighteight}{0.8486\,\ensuremath{\pm}\,0.0005}  % experiments/FIGS/data/w2b_leg2_metrics_v2.json :: cells.{l188-s1,l188-s2,l188-s3,l188-s4,l188-s5}.N1000000.s1.accuracy :: sha256 9cd059d3e2f51798 :: mean +- SD over the row's 5 pretrained models
\newcommand{\FtAccJcEThreeOnesixtwo}{0.7268\,\ensuremath{\pm}\,0.0017}  % experiments/FIGS/data/w2b_leg2_metrics_v2.json :: cells.{l162-s1b,l162-s2,l162-s3,l162-s4,l162-s5}.N1000.s1.accuracy :: sha256 9cd059d3e2f51798 :: mean +- SD over the row's 5 pretrained models
\newcommand{\FtAccJcEFourOnesixtwo}{0.8069\,\ensuremath{\pm}\,0.0025}  % experiments/FIGS/data/w2b_leg2_metrics_v2.json :: cells.{l162-s1b,l162-s2,l162-s3,l162-s4,l162-s5}.N10000.s1.accuracy :: sha256 9cd059d3e2f51798 :: mean +- SD over the row's 5 pretrained models
\newcommand{\FtAccJcEFiveOnesixtwo}{0.8363\,\ensuremath{\pm}\,0.0008}  % experiments/FIGS/data/w2b_leg2_metrics_v2.json :: cells.{l162-s1b,l162-s2,l162-s3,l162-s4,l162-s5}.N100000.s1.accuracy :: sha256 9cd059d3e2f51798 :: mean +- SD over the row's 5 pretrained models
\newcommand{\FtAccJcESixOnesixtwo}{0.8485\,\ensuremath{\pm}\,0.0005}  % experiments/FIGS/data/w2b_leg2_metrics_v2.json :: cells.{l162-s1b,l162-s2,l162-s3,l162-s4,l162-s5}.N1000000.s1.accuracy :: sha256 9cd059d3e2f51798 :: mean +- SD over the row's 5 pretrained models
\newcommand{\FtAccJcEThreeFourthree}{0.722\,\ensuremath{\pm}\,0.007}  % experiments/FIGS/data/w2b_leg2_metrics_v2.json :: cells.{r42q1-s1,r42q1-s2,r42q1-s3,r42q1-s4,r42q1-s5}.N1000.s1.accuracy :: sha256 9cd059d3e2f51798 :: mean +- SD over the row's 5 pretrained models
\newcommand{\FtAccJcEFourFourthree}{0.8050\,\ensuremath{\pm}\,0.0026}  % experiments/FIGS/data/w2b_leg2_metrics_v2.json :: cells.{r42q1-s1,r42q1-s2,r42q1-s3,r42q1-s4,r42q1-s5}.N10000.s1.accuracy :: sha256 9cd059d3e2f51798 :: mean +- SD over the row's 5 pretrained models
\newcommand{\FtAccJcEFiveFourthree}{0.8336\,\ensuremath{\pm}\,0.0008}  % experiments/FIGS/data/w2b_leg2_metrics_v2.json :: cells.{r42q1-s1,r42q1-s2,r42q1-s3,r42q1-s4,r42q1-s5}.N100000.s1.accuracy :: sha256 9cd059d3e2f51798 :: mean +- SD over the row's 5 pretrained models
\newcommand{\FtAccJcESixFourthree}{0.8468\,\ensuremath{\pm}\,0.0003}  % experiments/FIGS/data/w2b_leg2_metrics_v2.json :: cells.{r42q1-s1,r42q1-s2,r42q1-s3,r42q1-s4,r42q1-s5}.N1000000.s1.accuracy :: sha256 9cd059d3e2f51798 :: mean +- SD over the row's 5 pretrained models
\newcommand{\FtAccJcEThreeOneseven}{0.714\,\ensuremath{\pm}\,0.005}  % experiments/FIGS/data/w2b_leg2_metrics_v2.json :: cells.{r16q1-s1,r16q1-s2,r16q1-s3,r16q1-s4,r16q1-s5}.N1000.s1.accuracy :: sha256 9cd059d3e2f51798 :: mean +- SD over the row's 5 pretrained models
\newcommand{\FtAccJcEFourOneseven}{0.7964\,\ensuremath{\pm}\,0.0019}  % experiments/FIGS/data/w2b_leg2_metrics_v2.json :: cells.{r16q1-s1,r16q1-s2,r16q1-s3,r16q1-s4,r16q1-s5}.N10000.s1.accuracy :: sha256 9cd059d3e2f51798 :: mean +- SD over the row's 5 pretrained models
\newcommand{\FtAccJcEFiveOneseven}{0.8286\,\ensuremath{\pm}\,0.0006}  % experiments/FIGS/data/w2b_leg2_metrics_v2.json :: cells.{r16q1-s1,r16q1-s2,r16q1-s3,r16q1-s4,r16q1-s5}.N100000.s1.accuracy :: sha256 9cd059d3e2f51798 :: mean +- SD over the row's 5 pretrained models
\newcommand{\FtAccJcESixOneseven}{0.8415\,\ensuremath{\pm}\,0.0007}  % experiments/FIGS/data/w2b_leg2_metrics_v2.json :: cells.{r16q1-s1,r16q1-s2,r16q1-s3,r16q1-s4,r16q1-s5}.N1000000.s1.accuracy :: sha256 9cd059d3e2f51798 :: mean +- SD over the row's 5 pretrained models
\newcommand{\FtAccJcEThreeOnesixtwoMass}{0.728\,\ensuremath{\pm}\,0.004}  % experiments/FIGS/data/w2b_leg2_metrics_v2.json :: cells.{l162mass-s1,l162mass-s2,l162mass-s3,l162mass-s4,l162mass-s5}.N1000.s1.accuracy :: sha256 9cd059d3e2f51798 :: mean +- SD over the row's 5 pretrained models
\newcommand{\FtAccJcEFourOnesixtwoMass}{0.8043\,\ensuremath{\pm}\,0.0020}  % experiments/FIGS/data/w2b_leg2_metrics_v2.json :: cells.{l162mass-s1,l162mass-s2,l162mass-s3,l162mass-s4,l162mass-s5}.N10000.s1.accuracy :: sha256 9cd059d3e2f51798 :: mean +- SD over the row's 5 pretrained models
\newcommand{\FtAccJcEFiveOnesixtwoMass}{0.8349\,\ensuremath{\pm}\,0.0005}  % experiments/FIGS/data/w2b_leg2_metrics_v2.json :: cells.{l162mass-s1,l162mass-s2,l162mass-s3,l162mass-s4,l162mass-s5}.N100000.s1.accuracy :: sha256 9cd059d3e2f51798 :: mean +- SD over the row's 5 pretrained models
\newcommand{\FtAccJcESixOnesixtwoMass}{0.8475\,\ensuremath{\pm}\,0.0009}  % experiments/FIGS/data/w2b_leg2_metrics_v2.json :: cells.{l162mass-s1,l162mass-s2,l162mass-s3,l162mass-s4,l162mass-s5}.N1000000.s1.accuracy :: sha256 9cd059d3e2f51798 :: mean +- SD over the row's 5 pretrained models
\newcommand{\FtAccJcEThreeOnesevenMass}{0.702\,\ensuremath{\pm}\,0.005}  % experiments/FIGS/data/w2b_leg2_metrics_v2.json :: cells.{r16q1mass-s1,r16q1mass-s2,r16q1mass-s3,r16q1mass-s4,r16q1mass-s5}.N1000.s1.accuracy :: sha256 9cd059d3e2f51798 :: mean +- SD over the row's 5 pretrained models
\newcommand{\FtAccJcEFourOnesevenMass}{0.7909\,\ensuremath{\pm}\,0.0031}  % experiments/FIGS/data/w2b_leg2_metrics_v2.json :: cells.{r16q1mass-s1,r16q1mass-s2,r16q1mass-s3,r16q1mass-s4,r16q1mass-s5}.N10000.s1.accuracy :: sha256 9cd059d3e2f51798 :: mean +- SD over the row's 5 pretrained models
\newcommand{\FtAccJcEFiveOnesevenMass}{0.8265\,\ensuremath{\pm}\,0.0007}  % experiments/FIGS/data/w2b_leg2_metrics_v2.json :: cells.{r16q1mass-s1,r16q1mass-s2,r16q1mass-s3,r16q1mass-s4,r16q1mass-s5}.N100000.s1.accuracy :: sha256 9cd059d3e2f51798 :: mean +- SD over the row's 5 pretrained models
\newcommand{\FtAccJcESixOnesevenMass}{0.8390\,\ensuremath{\pm}\,0.0011}  % experiments/FIGS/data/w2b_leg2_metrics_v2.json :: cells.{r16q1mass-s1,r16q1mass-s2,r16q1mass-s3,r16q1mass-s4,r16q1mass-s5}.N1000000.s1.accuracy :: sha256 9cd059d3e2f51798 :: mean +- SD over the row's 5 pretrained models
\newcommand{\FtAccJcEThreeRandDrawOne}{0.7049}  % experiments/FIGS/data/w2b_leg2_metrics_v2.json :: cells.{rand-d1-s1b}.N1000.s1.accuracy :: sha256 9cd059d3e2f51798 :: one run
\newcommand{\FtAccJcEFourRandDrawOne}{0.7974}  % experiments/FIGS/data/w2b_leg2_metrics_v2.json :: cells.{rand-d1-s1b}.N10000.s1.accuracy :: sha256 9cd059d3e2f51798 :: one run
\newcommand{\FtAccJcEFiveRandDrawOne}{0.8326}  % experiments/FIGS/data/w2b_leg2_metrics_v2.json :: cells.{rand-d1-s1b}.N100000.s1.accuracy :: sha256 9cd059d3e2f51798 :: one run
\newcommand{\FtAccJcESixRandDrawOne}{0.84467}  % experiments/FIGS/data/w2b_leg2_metrics_v2.json :: cells.{rand-d1-s1b}.N1000000.s1.accuracy :: sha256 9cd059d3e2f51798 :: one run
\newcommand{\FtAccJcEThreeRandDrawTwo}{0.7135}  % experiments/FIGS/data/w2b_leg2_metrics_v2.json :: cells.{rand-d2-s2}.N1000.s1.accuracy :: sha256 9cd059d3e2f51798 :: one run
\newcommand{\FtAccJcEFourRandDrawTwo}{0.7998}  % experiments/FIGS/data/w2b_leg2_metrics_v2.json :: cells.{rand-d2-s2}.N10000.s1.accuracy :: sha256 9cd059d3e2f51798 :: one run
\newcommand{\FtAccJcEFiveRandDrawTwo}{0.8332}  % experiments/FIGS/data/w2b_leg2_metrics_v2.json :: cells.{rand-d2-s2}.N100000.s1.accuracy :: sha256 9cd059d3e2f51798 :: one run
\newcommand{\FtAccJcESixRandDrawTwo}{0.84449}  % experiments/FIGS/data/w2b_leg2_metrics_v2.json :: cells.{rand-d2-s2}.N1000000.s1.accuracy :: sha256 9cd059d3e2f51798 :: one run
\newcommand{\FtAccJcEThreeRandDrawThree}{0.7090}  % experiments/FIGS/data/w2b_leg2_metrics_v2.json :: cells.{rand-d3-s3}.N1000.s1.accuracy :: sha256 9cd059d3e2f51798 :: one run
\newcommand{\FtAccJcEFourRandDrawThree}{0.7948}  % experiments/FIGS/data/w2b_leg2_metrics_v2.json :: cells.{rand-d3-s3}.N10000.s1.accuracy :: sha256 9cd059d3e2f51798 :: one run
\newcommand{\FtAccJcEFiveRandDrawThree}{0.8310}  % experiments/FIGS/data/w2b_leg2_metrics_v2.json :: cells.{rand-d3-s3}.N100000.s1.accuracy :: sha256 9cd059d3e2f51798 :: one run
\newcommand{\FtAccJcESixRandDrawThree}{0.84409}  % experiments/FIGS/data/w2b_leg2_metrics_v2.json :: cells.{rand-d3-s3}.N1000000.s1.accuracy :: sha256 9cd059d3e2f51798 :: one run
\newcommand{\FtOmaRatioJcEThreeOnesixtwoOverOneeighteight}{0.96}  % experiments/FIGS/data/w2b_leg2_metrics_v2.json :: cells.*.N1000.s1.macro_auc_ovr :: sha256 9cd059d3e2f51798 :: ratio of the seed means of 1 - macro AUC, this row over the reference
\newcommand{\FtOmaRatioJcEThreeFourthreeOverOneeighteight}{1.01}  % experiments/FIGS/data/w2b_leg2_metrics_v2.json :: cells.*.N1000.s1.macro_auc_ovr :: sha256 9cd059d3e2f51798 :: ratio of the seed means of 1 - macro AUC, this row over the reference
\newcommand{\FtOmaRatioJcEThreeOnesevenOverOneeighteight}{1.07}  % experiments/FIGS/data/w2b_leg2_metrics_v2.json :: cells.*.N1000.s1.macro_auc_ovr :: sha256 9cd059d3e2f51798 :: ratio of the seed means of 1 - macro AUC, this row over the reference
\newcommand{\FtOmaRatioJcEThreeOnesixtwoMassOverOneeighteight}{0.87}  % experiments/FIGS/data/w2b_leg2_metrics_v2.json :: cells.*.N1000.s1.macro_auc_ovr :: sha256 9cd059d3e2f51798 :: ratio of the seed means of 1 - macro AUC, this row over the reference
\newcommand{\FtOmaRatioJcEThreeOnesevenMassOverOneeighteight}{1.09}  % experiments/FIGS/data/w2b_leg2_metrics_v2.json :: cells.*.N1000.s1.macro_auc_ovr :: sha256 9cd059d3e2f51798 :: ratio of the seed means of 1 - macro AUC, this row over the reference
\newcommand{\FtOmaRatioJcEThreeRandDrawOneOverOneeighteight}{1.06}  % experiments/FIGS/data/w2b_leg2_metrics_v2.json :: cells.*.N1000.s1.macro_auc_ovr :: sha256 9cd059d3e2f51798 :: ratio of the seed means of 1 - macro AUC, this row over the reference
\newcommand{\FtOmaRatioJcEThreeRandDrawTwoOverOneeighteight}{1.01}  % experiments/FIGS/data/w2b_leg2_metrics_v2.json :: cells.*.N1000.s1.macro_auc_ovr :: sha256 9cd059d3e2f51798 :: ratio of the seed means of 1 - macro AUC, this row over the reference
\newcommand{\FtOmaRatioJcEThreeRandDrawThreeOverOneeighteight}{1.07}  % experiments/FIGS/data/w2b_leg2_metrics_v2.json :: cells.*.N1000.s1.macro_auc_ovr :: sha256 9cd059d3e2f51798 :: ratio of the seed means of 1 - macro AUC, this row over the reference
\newcommand{\FtOmaRatioJcEThreeOneeighteightOverOnesixtwo}{1.05}  % experiments/FIGS/data/w2b_leg2_metrics_v2.json :: cells.*.N1000.s1.macro_auc_ovr :: sha256 9cd059d3e2f51798 :: ratio of the seed means of 1 - macro AUC, this row over the reference
\newcommand{\FtOmaRatioJcEThreeFourthreeOverOnesixtwo}{1.05}  % experiments/FIGS/data/w2b_leg2_metrics_v2.json :: cells.*.N1000.s1.macro_auc_ovr :: sha256 9cd059d3e2f51798 :: ratio of the seed means of 1 - macro AUC, this row over the reference
\newcommand{\FtOmaRatioJcEThreeOnesevenOverOnesixtwo}{1.12}  % experiments/FIGS/data/w2b_leg2_metrics_v2.json :: cells.*.N1000.s1.macro_auc_ovr :: sha256 9cd059d3e2f51798 :: ratio of the seed means of 1 - macro AUC, this row over the reference
\newcommand{\FtOmaRatioJcEThreeOnesixtwoMassOverOnesixtwo}{0.91}  % experiments/FIGS/data/w2b_leg2_metrics_v2.json :: cells.*.N1000.s1.macro_auc_ovr :: sha256 9cd059d3e2f51798 :: ratio of the seed means of 1 - macro AUC, this row over the reference
\newcommand{\FtOmaRatioJcEThreeOnesevenMassOverOnesixtwo}{1.14}  % experiments/FIGS/data/w2b_leg2_metrics_v2.json :: cells.*.N1000.s1.macro_auc_ovr :: sha256 9cd059d3e2f51798 :: ratio of the seed means of 1 - macro AUC, this row over the reference
\newcommand{\FtOmaRatioJcEThreeRandDrawOneOverOnesixtwo}{1.10}  % experiments/FIGS/data/w2b_leg2_metrics_v2.json :: cells.*.N1000.s1.macro_auc_ovr :: sha256 9cd059d3e2f51798 :: ratio of the seed means of 1 - macro AUC, this row over the reference
\newcommand{\FtOmaRatioJcEThreeRandDrawTwoOverOnesixtwo}{1.05}  % experiments/FIGS/data/w2b_leg2_metrics_v2.json :: cells.*.N1000.s1.macro_auc_ovr :: sha256 9cd059d3e2f51798 :: ratio of the seed means of 1 - macro AUC, this row over the reference
\newcommand{\FtOmaRatioJcEThreeRandDrawThreeOverOnesixtwo}{1.12}  % experiments/FIGS/data/w2b_leg2_metrics_v2.json :: cells.*.N1000.s1.macro_auc_ovr :: sha256 9cd059d3e2f51798 :: ratio of the seed means of 1 - macro AUC, this row over the reference
\newcommand{\FtOmaRatioJcEFourOnesixtwoOverOneeighteight}{0.98}  % experiments/FIGS/data/w2b_leg2_metrics_v2.json :: cells.*.N10000.s1.macro_auc_ovr :: sha256 9cd059d3e2f51798 :: ratio of the seed means of 1 - macro AUC, this row over the reference
\newcommand{\FtOmaRatioJcEFourFourthreeOverOneeighteight}{0.98}  % experiments/FIGS/data/w2b_leg2_metrics_v2.json :: cells.*.N10000.s1.macro_auc_ovr :: sha256 9cd059d3e2f51798 :: ratio of the seed means of 1 - macro AUC, this row over the reference
\newcommand{\FtOmaRatioJcEFourOnesevenOverOneeighteight}{1.07}  % experiments/FIGS/data/w2b_leg2_metrics_v2.json :: cells.*.N10000.s1.macro_auc_ovr :: sha256 9cd059d3e2f51798 :: ratio of the seed means of 1 - macro AUC, this row over the reference
\newcommand{\FtOmaRatioJcEFourOnesixtwoMassOverOneeighteight}{0.99}  % experiments/FIGS/data/w2b_leg2_metrics_v2.json :: cells.*.N10000.s1.macro_auc_ovr :: sha256 9cd059d3e2f51798 :: ratio of the seed means of 1 - macro AUC, this row over the reference
\newcommand{\FtOmaRatioJcEFourOnesevenMassOverOneeighteight}{1.12}  % experiments/FIGS/data/w2b_leg2_metrics_v2.json :: cells.*.N10000.s1.macro_auc_ovr :: sha256 9cd059d3e2f51798 :: ratio of the seed means of 1 - macro AUC, this row over the reference
\newcommand{\FtOmaRatioJcEFourRandDrawOneOverOneeighteight}{1.08}  % experiments/FIGS/data/w2b_leg2_metrics_v2.json :: cells.*.N10000.s1.macro_auc_ovr :: sha256 9cd059d3e2f51798 :: ratio of the seed means of 1 - macro AUC, this row over the reference
\newcommand{\FtOmaRatioJcEFourRandDrawTwoOverOneeighteight}{1.04}  % experiments/FIGS/data/w2b_leg2_metrics_v2.json :: cells.*.N10000.s1.macro_auc_ovr :: sha256 9cd059d3e2f51798 :: ratio of the seed means of 1 - macro AUC, this row over the reference
\newcommand{\FtOmaRatioJcEFourRandDrawThreeOverOneeighteight}{1.09}  % experiments/FIGS/data/w2b_leg2_metrics_v2.json :: cells.*.N10000.s1.macro_auc_ovr :: sha256 9cd059d3e2f51798 :: ratio of the seed means of 1 - macro AUC, this row over the reference
\newcommand{\FtOmaRatioJcEFourOneeighteightOverOnesixtwo}{1.02}  % experiments/FIGS/data/w2b_leg2_metrics_v2.json :: cells.*.N10000.s1.macro_auc_ovr :: sha256 9cd059d3e2f51798 :: ratio of the seed means of 1 - macro AUC, this row over the reference
\newcommand{\FtOmaRatioJcEFourFourthreeOverOnesixtwo}{1.01}  % experiments/FIGS/data/w2b_leg2_metrics_v2.json :: cells.*.N10000.s1.macro_auc_ovr :: sha256 9cd059d3e2f51798 :: ratio of the seed means of 1 - macro AUC, this row over the reference
\newcommand{\FtOmaRatioJcEFourOnesevenOverOnesixtwo}{1.09}  % experiments/FIGS/data/w2b_leg2_metrics_v2.json :: cells.*.N10000.s1.macro_auc_ovr :: sha256 9cd059d3e2f51798 :: ratio of the seed means of 1 - macro AUC, this row over the reference
\newcommand{\FtOmaRatioJcEFourOnesixtwoMassOverOnesixtwo}{1.01}  % experiments/FIGS/data/w2b_leg2_metrics_v2.json :: cells.*.N10000.s1.macro_auc_ovr :: sha256 9cd059d3e2f51798 :: ratio of the seed means of 1 - macro AUC, this row over the reference
\newcommand{\FtOmaRatioJcEFourOnesevenMassOverOnesixtwo}{1.14}  % experiments/FIGS/data/w2b_leg2_metrics_v2.json :: cells.*.N10000.s1.macro_auc_ovr :: sha256 9cd059d3e2f51798 :: ratio of the seed means of 1 - macro AUC, this row over the reference
\newcommand{\FtOmaRatioJcEFourRandDrawOneOverOnesixtwo}{1.11}  % experiments/FIGS/data/w2b_leg2_metrics_v2.json :: cells.*.N10000.s1.macro_auc_ovr :: sha256 9cd059d3e2f51798 :: ratio of the seed means of 1 - macro AUC, this row over the reference
\newcommand{\FtOmaRatioJcEFourRandDrawTwoOverOnesixtwo}{1.06}  % experiments/FIGS/data/w2b_leg2_metrics_v2.json :: cells.*.N10000.s1.macro_auc_ovr :: sha256 9cd059d3e2f51798 :: ratio of the seed means of 1 - macro AUC, this row over the reference
\newcommand{\FtOmaRatioJcEFourRandDrawThreeOverOnesixtwo}{1.12}  % experiments/FIGS/data/w2b_leg2_metrics_v2.json :: cells.*.N10000.s1.macro_auc_ovr :: sha256 9cd059d3e2f51798 :: ratio of the seed means of 1 - macro AUC, this row over the reference
\newcommand{\FtOmaRatioJcEFiveOnesixtwoOverOneeighteight}{0.98}  % experiments/FIGS/data/w2b_leg2_metrics_v2.json :: cells.*.N100000.s1.macro_auc_ovr :: sha256 9cd059d3e2f51798 :: ratio of the seed means of 1 - macro AUC, this row over the reference
\newcommand{\FtOmaRatioJcEFiveFourthreeOverOneeighteight}{1.02}  % experiments/FIGS/data/w2b_leg2_metrics_v2.json :: cells.*.N100000.s1.macro_auc_ovr :: sha256 9cd059d3e2f51798 :: ratio of the seed means of 1 - macro AUC, this row over the reference
\newcommand{\FtOmaRatioJcEFiveOnesevenOverOneeighteight}{1.07}  % experiments/FIGS/data/w2b_leg2_metrics_v2.json :: cells.*.N100000.s1.macro_auc_ovr :: sha256 9cd059d3e2f51798 :: ratio of the seed means of 1 - macro AUC, this row over the reference
\newcommand{\FtOmaRatioJcEFiveOnesixtwoMassOverOneeighteight}{1.00}  % experiments/FIGS/data/w2b_leg2_metrics_v2.json :: cells.*.N100000.s1.macro_auc_ovr :: sha256 9cd059d3e2f51798 :: ratio of the seed means of 1 - macro AUC, this row over the reference
\newcommand{\FtOmaRatioJcEFiveOnesevenMassOverOneeighteight}{1.10}  % experiments/FIGS/data/w2b_leg2_metrics_v2.json :: cells.*.N100000.s1.macro_auc_ovr :: sha256 9cd059d3e2f51798 :: ratio of the seed means of 1 - macro AUC, this row over the reference
\newcommand{\FtOmaRatioJcEFiveRandDrawOneOverOneeighteight}{1.04}  % experiments/FIGS/data/w2b_leg2_metrics_v2.json :: cells.*.N100000.s1.macro_auc_ovr :: sha256 9cd059d3e2f51798 :: ratio of the seed means of 1 - macro AUC, this row over the reference
\newcommand{\FtOmaRatioJcEFiveRandDrawTwoOverOneeighteight}{1.03}  % experiments/FIGS/data/w2b_leg2_metrics_v2.json :: cells.*.N100000.s1.macro_auc_ovr :: sha256 9cd059d3e2f51798 :: ratio of the seed means of 1 - macro AUC, this row over the reference
\newcommand{\FtOmaRatioJcEFiveRandDrawThreeOverOneeighteight}{1.05}  % experiments/FIGS/data/w2b_leg2_metrics_v2.json :: cells.*.N100000.s1.macro_auc_ovr :: sha256 9cd059d3e2f51798 :: ratio of the seed means of 1 - macro AUC, this row over the reference
\newcommand{\FtOmaRatioJcEFiveOneeighteightOverOnesixtwo}{1.02}  % experiments/FIGS/data/w2b_leg2_metrics_v2.json :: cells.*.N100000.s1.macro_auc_ovr :: sha256 9cd059d3e2f51798 :: ratio of the seed means of 1 - macro AUC, this row over the reference
\newcommand{\FtOmaRatioJcEFiveFourthreeOverOnesixtwo}{1.03}  % experiments/FIGS/data/w2b_leg2_metrics_v2.json :: cells.*.N100000.s1.macro_auc_ovr :: sha256 9cd059d3e2f51798 :: ratio of the seed means of 1 - macro AUC, this row over the reference
\newcommand{\FtOmaRatioJcEFiveOnesevenOverOnesixtwo}{1.09}  % experiments/FIGS/data/w2b_leg2_metrics_v2.json :: cells.*.N100000.s1.macro_auc_ovr :: sha256 9cd059d3e2f51798 :: ratio of the seed means of 1 - macro AUC, this row over the reference
\newcommand{\FtOmaRatioJcEFiveOnesixtwoMassOverOnesixtwo}{1.02}  % experiments/FIGS/data/w2b_leg2_metrics_v2.json :: cells.*.N100000.s1.macro_auc_ovr :: sha256 9cd059d3e2f51798 :: ratio of the seed means of 1 - macro AUC, this row over the reference
\newcommand{\FtOmaRatioJcEFiveOnesevenMassOverOnesixtwo}{1.11}  % experiments/FIGS/data/w2b_leg2_metrics_v2.json :: cells.*.N100000.s1.macro_auc_ovr :: sha256 9cd059d3e2f51798 :: ratio of the seed means of 1 - macro AUC, this row over the reference
\newcommand{\FtOmaRatioJcEFiveRandDrawOneOverOnesixtwo}{1.06}  % experiments/FIGS/data/w2b_leg2_metrics_v2.json :: cells.*.N100000.s1.macro_auc_ovr :: sha256 9cd059d3e2f51798 :: ratio of the seed means of 1 - macro AUC, this row over the reference
\newcommand{\FtOmaRatioJcEFiveRandDrawTwoOverOnesixtwo}{1.05}  % experiments/FIGS/data/w2b_leg2_metrics_v2.json :: cells.*.N100000.s1.macro_auc_ovr :: sha256 9cd059d3e2f51798 :: ratio of the seed means of 1 - macro AUC, this row over the reference
\newcommand{\FtOmaRatioJcEFiveRandDrawThreeOverOnesixtwo}{1.07}  % experiments/FIGS/data/w2b_leg2_metrics_v2.json :: cells.*.N100000.s1.macro_auc_ovr :: sha256 9cd059d3e2f51798 :: ratio of the seed means of 1 - macro AUC, this row over the reference
\newcommand{\FtOmaRatioJcESixOnesixtwoOverOneeighteight}{1.00}  % experiments/FIGS/data/w2b_leg2_metrics_v2.json :: cells.*.N1000000.s1.macro_auc_ovr :: sha256 9cd059d3e2f51798 :: ratio of the seed means of 1 - macro AUC, this row over the reference
\newcommand{\FtOmaRatioJcESixFourthreeOverOneeighteight}{1.02}  % experiments/FIGS/data/w2b_leg2_metrics_v2.json :: cells.*.N1000000.s1.macro_auc_ovr :: sha256 9cd059d3e2f51798 :: ratio of the seed means of 1 - macro AUC, this row over the reference
\newcommand{\FtOmaRatioJcESixOnesevenOverOneeighteight}{1.08}  % experiments/FIGS/data/w2b_leg2_metrics_v2.json :: cells.*.N1000000.s1.macro_auc_ovr :: sha256 9cd059d3e2f51798 :: ratio of the seed means of 1 - macro AUC, this row over the reference
\newcommand{\FtOmaRatioJcESixOnesixtwoMassOverOneeighteight}{1.01}  % experiments/FIGS/data/w2b_leg2_metrics_v2.json :: cells.*.N1000000.s1.macro_auc_ovr :: sha256 9cd059d3e2f51798 :: ratio of the seed means of 1 - macro AUC, this row over the reference
\newcommand{\FtOmaRatioJcESixOnesevenMassOverOneeighteight}{1.11}  % experiments/FIGS/data/w2b_leg2_metrics_v2.json :: cells.*.N1000000.s1.macro_auc_ovr :: sha256 9cd059d3e2f51798 :: ratio of the seed means of 1 - macro AUC, this row over the reference
\newcommand{\FtOmaRatioJcESixRandDrawOneOverOneeighteight}{1.05}  % experiments/FIGS/data/w2b_leg2_metrics_v2.json :: cells.*.N1000000.s1.macro_auc_ovr :: sha256 9cd059d3e2f51798 :: ratio of the seed means of 1 - macro AUC, this row over the reference
\newcommand{\FtOmaRatioJcESixRandDrawTwoOverOneeighteight}{1.05}  % experiments/FIGS/data/w2b_leg2_metrics_v2.json :: cells.*.N1000000.s1.macro_auc_ovr :: sha256 9cd059d3e2f51798 :: ratio of the seed means of 1 - macro AUC, this row over the reference
\newcommand{\FtOmaRatioJcESixRandDrawThreeOverOneeighteight}{1.05}  % experiments/FIGS/data/w2b_leg2_metrics_v2.json :: cells.*.N1000000.s1.macro_auc_ovr :: sha256 9cd059d3e2f51798 :: ratio of the seed means of 1 - macro AUC, this row over the reference
\newcommand{\FtOmaRatioJcESixOneeighteightOverOnesixtwo}{1.00}  % experiments/FIGS/data/w2b_leg2_metrics_v2.json :: cells.*.N1000000.s1.macro_auc_ovr :: sha256 9cd059d3e2f51798 :: ratio of the seed means of 1 - macro AUC, this row over the reference
\newcommand{\FtOmaRatioJcESixFourthreeOverOnesixtwo}{1.02}  % experiments/FIGS/data/w2b_leg2_metrics_v2.json :: cells.*.N1000000.s1.macro_auc_ovr :: sha256 9cd059d3e2f51798 :: ratio of the seed means of 1 - macro AUC, this row over the reference
\newcommand{\FtOmaRatioJcESixOnesevenOverOnesixtwo}{1.08}  % experiments/FIGS/data/w2b_leg2_metrics_v2.json :: cells.*.N1000000.s1.macro_auc_ovr :: sha256 9cd059d3e2f51798 :: ratio of the seed means of 1 - macro AUC, this row over the reference
\newcommand{\FtOmaRatioJcESixOnesixtwoMassOverOnesixtwo}{1.01}  % experiments/FIGS/data/w2b_leg2_metrics_v2.json :: cells.*.N1000000.s1.macro_auc_ovr :: sha256 9cd059d3e2f51798 :: ratio of the seed means of 1 - macro AUC, this row over the reference
\newcommand{\FtOmaRatioJcESixOnesevenMassOverOnesixtwo}{1.11}  % experiments/FIGS/data/w2b_leg2_metrics_v2.json :: cells.*.N1000000.s1.macro_auc_ovr :: sha256 9cd059d3e2f51798 :: ratio of the seed means of 1 - macro AUC, this row over the reference
\newcommand{\FtOmaRatioJcESixRandDrawOneOverOnesixtwo}{1.05}  % experiments/FIGS/data/w2b_leg2_metrics_v2.json :: cells.*.N1000000.s1.macro_auc_ovr :: sha256 9cd059d3e2f51798 :: ratio of the seed means of 1 - macro AUC, this row over the reference
\newcommand{\FtOmaRatioJcESixRandDrawTwoOverOnesixtwo}{1.05}  % experiments/FIGS/data/w2b_leg2_metrics_v2.json :: cells.*.N1000000.s1.macro_auc_ovr :: sha256 9cd059d3e2f51798 :: ratio of the seed means of 1 - macro AUC, this row over the reference
\newcommand{\FtOmaRatioJcESixRandDrawThreeOverOnesixtwo}{1.05}  % experiments/FIGS/data/w2b_leg2_metrics_v2.json :: cells.*.N1000000.s1.macro_auc_ovr :: sha256 9cd059d3e2f51798 :: ratio of the seed means of 1 - macro AUC, this row over the reference
\newcommand{\PairedFtAccJcEThreeOnesevenMinusOneeighteight}{\ensuremath{-}0.0040~[\ensuremath{-}0.0144, +0.0064]}  % experiments/FIGS/data/w2b_leg2_metrics_v2.json :: cells.{r16q1-s1,r16q1-s2,r16q1-s3,r16q1-s4,r16q1-s5,l188-s1,l188-s2,l188-s3,l188-s4,l188-s5}.N1000.s1.accuracy :: sha256 9cd059d3e2f51798 :: accuracy, coarse minus fine, paired by run [95% Student-t interval]
\newcommand{\PairedFtAccJcEFourOnesevenMinusOneeighteight}{\ensuremath{-}0.0075~[\ensuremath{-}0.0118, \ensuremath{-}0.0033]}  % experiments/FIGS/data/w2b_leg2_metrics_v2.json :: cells.{r16q1-s1,r16q1-s2,r16q1-s3,r16q1-s4,r16q1-s5,l188-s1,l188-s2,l188-s3,l188-s4,l188-s5}.N10000.s1.accuracy :: sha256 9cd059d3e2f51798 :: accuracy, coarse minus fine, paired by run [95% Student-t interval]
\newcommand{\PairedFtAccJcEFiveOnesevenMinusOneeighteight}{\ensuremath{-}0.0066~[\ensuremath{-}0.0083, \ensuremath{-}0.0049]}  % experiments/FIGS/data/w2b_leg2_metrics_v2.json :: cells.{r16q1-s1,r16q1-s2,r16q1-s3,r16q1-s4,r16q1-s5,l188-s1,l188-s2,l188-s3,l188-s4,l188-s5}.N100000.s1.accuracy :: sha256 9cd059d3e2f51798 :: accuracy, coarse minus fine, paired by run [95% Student-t interval]
\newcommand{\PairedFtAccJcESixOnesevenMinusOneeighteight}{\ensuremath{-}0.0071~[\ensuremath{-}0.0080, \ensuremath{-}0.0062]}  % experiments/FIGS/data/w2b_leg2_metrics_v2.json :: cells.{r16q1-s1,r16q1-s2,r16q1-s3,r16q1-s4,r16q1-s5,l188-s1,l188-s2,l188-s3,l188-s4,l188-s5}.N1000000.s1.accuracy :: sha256 9cd059d3e2f51798 :: accuracy, coarse minus fine, paired by run [95% Student-t interval]
\newcommand{\FtNClassesJc}{10}  % experiments/FIGS/data/w2b_leg2_metrics_v2.json :: cells.*.*.s1.n_classes_present :: sha256 9cd059d3e2f51798 :: classes of the fine-tuning task
\newcommand{\FtNTestAucJc}{500{,}000}  % experiments/FIGS/data/w2b_leg2_metrics_v2.json :: cells.*.*.s1.n_jets_auc :: sha256 9cd059d3e2f51798 :: test jets the macro AUC is computed on
\newcommand{\FtNTestAccJc}{2{,}000{,}000}  % experiments/FIGS/data/w2b_leg2_metrics_v2.json :: cells.*.*.s1.n_jets :: sha256 9cd059d3e2f51798 :: test jets the accuracy is computed on
\newcommand{\FtSizeEThree}{\ensuremath{10^{3}}}  % experiments/FIGS/data/w2b_leg1_metrics_v2.json :: cells.*.N1000 :: sha256 b56bf70be9982044 :: fine-tuning training jets
\newcommand{\FtSizeEFour}{\ensuremath{10^{4}}}  % experiments/FIGS/data/w2b_leg1_metrics_v2.json :: cells.*.N10000 :: sha256 b56bf70be9982044 :: fine-tuning training jets
\newcommand{\FtSizeEFive}{\ensuremath{10^{5}}}  % experiments/FIGS/data/w2b_leg1_metrics_v2.json :: cells.*.N100000 :: sha256 b56bf70be9982044 :: fine-tuning training jets
\newcommand{\FtSizeESix}{\ensuremath{10^{6}}}  % experiments/FIGS/data/w2b_leg1_metrics_v2.json :: cells.*.N1000000 :: sha256 b56bf70be9982044 :: fine-tuning training jets
\newcommand{\FtNRunsWord}{five}  % experiments/FIGS/data/w2b_leg1_metrics_v2.json :: cells (pretrained models per vocabulary row) :: sha256 b56bf70be9982044 :: pretraining runs per vocabulary
\newcommand{\PairedProbeBcVsRestLinearOnesevenOverOnesixtwo}{1.64~[1.52, 1.78]}  % experiments/FIGS/data/paired_v1err/probes/ratios.json :: ratios[0].ratio, ratios[0].ci95 :: sha256 14dd51402305d147 :: 17 over 162, bc_vs_rest linear 1-auc: paired geometric mean [95% interval]
\newcommand{\PairedProbeBcVsRestMlpOnesevenOverOnesixtwo}{1.50~[1.39, 1.62]}  % experiments/FIGS/data/paired_v1err/probes/ratios.json :: ratios[3].ratio, ratios[3].ci95 :: sha256 14dd51402305d147 :: 17 over 162, bc_vs_rest mlp 1-auc: paired geometric mean [95% interval]
\newcommand{\PairedProbeBvcFourprongLinearOnesixtwoOverOneeighteight}{1.01~[0.95, 1.08]}  % experiments/FIGS/data/paired_v1err/probes/ratios.json :: ratios[6].ratio, ratios[6].ci95 :: sha256 14dd51402305d147 :: 162 over 188, bvc_4prong linear 1-auc: paired geometric mean [95% interval]
\newcommand{\PairedProbeBvcFourprongLinearFourthreeOverOneeighteight}{1.10~[1.01, 1.20]}  % experiments/FIGS/data/paired_v1err/probes/ratios.json :: ratios[7].ratio, ratios[7].ci95 :: sha256 14dd51402305d147 :: 43 over 188, bvc_4prong linear 1-auc: paired geometric mean [95% interval]
\newcommand{\PairedProbeBvcFourprongLinearOnesevenOverOneeighteight}{4.5~[3.5, 5.7]}  % experiments/FIGS/data/paired_v1err/probes/ratios.json :: ratios[8].ratio, ratios[8].ci95 :: sha256 14dd51402305d147 :: 17 over 188, bvc_4prong linear 1-auc: paired geometric mean [95% interval]
\newcommand{\PairedProbeBvcFourprongLinearFourthreeOverOnesixtwo}{1.09~[0.98, 1.20]}  % experiments/FIGS/data/paired_v1err/probes/ratios.json :: ratios[9].ratio, ratios[9].ci95 :: sha256 14dd51402305d147 :: 43 over 162, bvc_4prong linear 1-auc: paired geometric mean [95% interval]
\newcommand{\PairedProbeBvcFourprongLinearOnesevenOverOnesixtwo}{4.4~[3.4, 5.6]}  % experiments/FIGS/data/paired_v1err/probes/ratios.json :: ratios[10].ratio, ratios[10].ci95 :: sha256 14dd51402305d147 :: 17 over 162, bvc_4prong linear 1-auc: paired geometric mean [95% interval]
\newcommand{\PairedProbeBvcFourprongLinearOnesevenOverFourthree}{4.0~[3.1, 5.4]}  % experiments/FIGS/data/paired_v1err/probes/ratios.json :: ratios[11].ratio, ratios[11].ci95 :: sha256 14dd51402305d147 :: 17 over 43, bvc_4prong linear 1-auc: paired geometric mean [95% interval]
\newcommand{\PairedProbeBvcFourprongLinearOnesixtwoMassOverOnesixtwo}{1.05~[0.97, 1.13]}  % experiments/FIGS/data/paired_v1err/probes/ratios.json :: ratios[12].ratio, ratios[12].ci95 :: sha256 14dd51402305d147 :: 162+mass over 162, bvc_4prong linear 1-auc: paired geometric mean [95% interval]
\newcommand{\PairedProbeBvcFourprongLinearOnesevenMassOverOneseven}{2.1~[1.6, 2.7]}  % experiments/FIGS/data/paired_v1err/probes/ratios.json :: ratios[13].ratio, ratios[13].ci95 :: sha256 14dd51402305d147 :: 17+mass over 17, bvc_4prong linear 1-auc: paired geometric mean [95% interval]
\newcommand{\PairedProbeBvcFourprongLinearRandomOverOneeighteight}{1.54~[1.37, 1.74]}  % experiments/FIGS/data/paired_v1err/probes/ratios.json :: ratios[14].ratio, ratios[14].ci95 :: sha256 14dd51402305d147 :: random over 188, bvc_4prong linear 1-auc: paired geometric mean [95% interval]
\newcommand{\PairedProbeBvcFourprongLinearRandomOverOnesixtwo}{1.55~[1.36, 1.77]}  % experiments/FIGS/data/paired_v1err/probes/ratios.json :: ratios[15].ratio, ratios[15].ci95 :: sha256 14dd51402305d147 :: random over 162, bvc_4prong linear 1-auc: paired geometric mean [95% interval]
\newcommand{\PairedProbeBvcFourprongLinearRandomOverFourthree}{1.39~[1.27, 1.52]}  % experiments/FIGS/data/paired_v1err/probes/ratios.json :: ratios[16].ratio, ratios[16].ci95 :: sha256 14dd51402305d147 :: random over 43, bvc_4prong linear 1-auc: paired geometric mean [95% interval]
\newcommand{\PairedProbeBvcFourprongLinearRandomOverOneseven}{0.34~[0.20, 0.58]}  % experiments/FIGS/data/paired_v1err/probes/ratios.json :: ratios[17].ratio, ratios[17].ci95 :: sha256 14dd51402305d147 :: random over 17, bvc_4prong linear 1-auc: paired geometric mean [95% interval]
\newcommand{\PairedProbeBvcFourprongLinearRandomDrawTwoOverRandomDrawOne}{1.05~[0.94, 1.18]}  % experiments/FIGS/data/paired_v1err/probes/ratios.json :: ratios[18].ratio, ratios[18].ci95 :: sha256 14dd51402305d147 :: random draw 2 over random draw 1, bvc_4prong linear 1-auc: paired geometric mean [95% interval]
\newcommand{\PairedProbeBvcFourprongLinearRandomDrawThreeOverRandomDrawOne}{0.95~[0.85, 1.06]}  % experiments/FIGS/data/paired_v1err/probes/ratios.json :: ratios[19].ratio, ratios[19].ci95 :: sha256 14dd51402305d147 :: random draw 3 over random draw 1, bvc_4prong linear 1-auc: paired geometric mean [95% interval]
\newcommand{\PairedProbeBvcFourprongLinearRandomDrawThreeOverRandomDrawTwo}{0.90~[0.79, 1.02]}  % experiments/FIGS/data/paired_v1err/probes/ratios.json :: ratios[20].ratio, ratios[20].ci95 :: sha256 14dd51402305d147 :: random draw 3 over random draw 2, bvc_4prong linear 1-auc: paired geometric mean [95% interval]
\newcommand{\PairedProbeBvcFourprongMlpOnesixtwoOverOneeighteight}{1.00~[0.93, 1.07]}  % experiments/FIGS/data/paired_v1err/probes/ratios.json :: ratios[66].ratio, ratios[66].ci95 :: sha256 14dd51402305d147 :: 162 over 188, bvc_4prong mlp 1-auc: paired geometric mean [95% interval]
\newcommand{\PairedProbeBvcFourprongMlpFourthreeOverOneeighteight}{1.12~[1.03, 1.22]}  % experiments/FIGS/data/paired_v1err/probes/ratios.json :: ratios[67].ratio, ratios[67].ci95 :: sha256 14dd51402305d147 :: 43 over 188, bvc_4prong mlp 1-auc: paired geometric mean [95% interval]
\newcommand{\PairedProbeBvcFourprongMlpOnesevenOverOneeighteight}{3.2~[2.6, 4.1]}  % experiments/FIGS/data/paired_v1err/probes/ratios.json :: ratios[68].ratio, ratios[68].ci95 :: sha256 14dd51402305d147 :: 17 over 188, bvc_4prong mlp 1-auc: paired geometric mean [95% interval]
\newcommand{\PairedProbeBvcFourprongMlpFourthreeOverOnesixtwo}{1.12~[1.01, 1.26]}  % experiments/FIGS/data/paired_v1err/probes/ratios.json :: ratios[69].ratio, ratios[69].ci95 :: sha256 14dd51402305d147 :: 43 over 162, bvc_4prong mlp 1-auc: paired geometric mean [95% interval]
\newcommand{\PairedProbeBvcFourprongMlpOnesevenOverOnesixtwo}{3.3~[2.6, 4.0]}  % experiments/FIGS/data/paired_v1err/probes/ratios.json :: ratios[70].ratio, ratios[70].ci95 :: sha256 14dd51402305d147 :: 17 over 162, bvc_4prong mlp 1-auc: paired geometric mean [95% interval]
\newcommand{\PairedProbeBvcFourprongMlpOnesevenOverFourthree}{2.9~[2.3, 3.7]}  % experiments/FIGS/data/paired_v1err/probes/ratios.json :: ratios[71].ratio, ratios[71].ci95 :: sha256 14dd51402305d147 :: 17 over 43, bvc_4prong mlp 1-auc: paired geometric mean [95% interval]
\newcommand{\PairedProbeBvcFourprongMlpOnesixtwoMassOverOnesixtwo}{1.07~[0.99, 1.16]}  % experiments/FIGS/data/paired_v1err/probes/ratios.json :: ratios[72].ratio, ratios[72].ci95 :: sha256 14dd51402305d147 :: 162+mass over 162, bvc_4prong mlp 1-auc: paired geometric mean [95% interval]
\newcommand{\PairedProbeBvcFourprongMlpOnesevenMassOverOneseven}{1.97~[1.46, 2.67]}  % experiments/FIGS/data/paired_v1err/probes/ratios.json :: ratios[73].ratio, ratios[73].ci95 :: sha256 14dd51402305d147 :: 17+mass over 17, bvc_4prong mlp 1-auc: paired geometric mean [95% interval]
\newcommand{\PairedProbeBvcFourprongMlpRandomOverOneeighteight}{1.47~[1.30, 1.65]}  % experiments/FIGS/data/paired_v1err/probes/ratios.json :: ratios[74].ratio, ratios[74].ci95 :: sha256 14dd51402305d147 :: random over 188, bvc_4prong mlp 1-auc: paired geometric mean [95% interval]
\newcommand{\PairedProbeBvcFourprongMlpRandomOverOnesixtwo}{1.49~[1.24, 1.79]}  % experiments/FIGS/data/paired_v1err/probes/ratios.json :: ratios[75].ratio, ratios[75].ci95 :: sha256 14dd51402305d147 :: random over 162, bvc_4prong mlp 1-auc: paired geometric mean [95% interval]
\newcommand{\PairedProbeBvcFourprongMlpRandomOverFourthree}{1.28~[1.16, 1.43]}  % experiments/FIGS/data/paired_v1err/probes/ratios.json :: ratios[76].ratio, ratios[76].ci95 :: sha256 14dd51402305d147 :: random over 43, bvc_4prong mlp 1-auc: paired geometric mean [95% interval]
\newcommand{\PairedProbeBvcFourprongMlpRandomOverOneseven}{0.43~[0.28, 0.67]}  % experiments/FIGS/data/paired_v1err/probes/ratios.json :: ratios[77].ratio, ratios[77].ci95 :: sha256 14dd51402305d147 :: random over 17, bvc_4prong mlp 1-auc: paired geometric mean [95% interval]
\newcommand{\PairedProbeBvcFourprongMlpRandomDrawTwoOverRandomDrawOne}{1.00~[0.90, 1.10]}  % experiments/FIGS/data/paired_v1err/probes/ratios.json :: ratios[78].ratio, ratios[78].ci95 :: sha256 14dd51402305d147 :: random draw 2 over random draw 1, bvc_4prong mlp 1-auc: paired geometric mean [95% interval]
\newcommand{\PairedProbeBvcFourprongMlpRandomDrawThreeOverRandomDrawOne}{0.89~[0.79, 0.99]}  % experiments/FIGS/data/paired_v1err/probes/ratios.json :: ratios[79].ratio, ratios[79].ci95 :: sha256 14dd51402305d147 :: random draw 3 over random draw 1, bvc_4prong mlp 1-auc: paired geometric mean [95% interval]
\newcommand{\PairedProbeBvcFourprongMlpRandomDrawThreeOverRandomDrawTwo}{0.89~[0.79, 0.99]}  % experiments/FIGS/data/paired_v1err/probes/ratios.json :: ratios[80].ratio, ratios[80].ci95 :: sha256 14dd51402305d147 :: random draw 3 over random draw 2, bvc_4prong mlp 1-auc: paired geometric mean [95% interval]
\newcommand{\PairedProbeBvcQcdLinearOnesixtwoOverOneeighteight}{1.15~[1.07, 1.23]}  % experiments/FIGS/data/paired_v1err/probes/ratios.json :: ratios[126].ratio, ratios[126].ci95 :: sha256 14dd51402305d147 :: 162 over 188, bvc_qcd linear 1-auc: paired geometric mean [95% interval]
\newcommand{\PairedProbeBvcQcdLinearFourthreeOverOneeighteight}{1.16~[1.08, 1.25]}  % experiments/FIGS/data/paired_v1err/probes/ratios.json :: ratios[127].ratio, ratios[127].ci95 :: sha256 14dd51402305d147 :: 43 over 188, bvc_qcd linear 1-auc: paired geometric mean [95% interval]
\newcommand{\PairedProbeBvcQcdLinearOnesevenOverOneeighteight}{3.0~[2.5, 3.7]}  % experiments/FIGS/data/paired_v1err/probes/ratios.json :: ratios[128].ratio, ratios[128].ci95 :: sha256 14dd51402305d147 :: 17 over 188, bvc_qcd linear 1-auc: paired geometric mean [95% interval]
\newcommand{\PairedProbeBvcQcdLinearFourthreeOverOnesixtwo}{1.01~[0.96, 1.07]}  % experiments/FIGS/data/paired_v1err/probes/ratios.json :: ratios[129].ratio, ratios[129].ci95 :: sha256 14dd51402305d147 :: 43 over 162, bvc_qcd linear 1-auc: paired geometric mean [95% interval]
\newcommand{\PairedProbeBvcQcdLinearOnesevenOverOnesixtwo}{2.6~[2.2, 3.1]}  % experiments/FIGS/data/paired_v1err/probes/ratios.json :: ratios[130].ratio, ratios[130].ci95 :: sha256 14dd51402305d147 :: 17 over 162, bvc_qcd linear 1-auc: paired geometric mean [95% interval]
\newcommand{\PairedProbeBvcQcdLinearOnesevenOverFourthree}{2.6~[2.2, 3.0]}  % experiments/FIGS/data/paired_v1err/probes/ratios.json :: ratios[131].ratio, ratios[131].ci95 :: sha256 14dd51402305d147 :: 17 over 43, bvc_qcd linear 1-auc: paired geometric mean [95% interval]
\newcommand{\PairedProbeBvcQcdLinearOnesixtwoMassOverOnesixtwo}{1.060~[1.002, 1.123]}  % experiments/FIGS/data/paired_v1err/probes/ratios.json :: ratios[132].ratio, ratios[132].ci95 :: sha256 14dd51402305d147 :: 162+mass over 162, bvc_qcd linear 1-auc: paired geometric mean [95% interval]
\newcommand{\PairedProbeBvcQcdLinearOnesevenMassOverOneseven}{1.89~[1.52, 2.34]}  % experiments/FIGS/data/paired_v1err/probes/ratios.json :: ratios[133].ratio, ratios[133].ci95 :: sha256 14dd51402305d147 :: 17+mass over 17, bvc_qcd linear 1-auc: paired geometric mean [95% interval]
\newcommand{\PairedProbeBvcQcdMlpOnesixtwoOverOneeighteight}{1.11~[1.04, 1.19]}  % experiments/FIGS/data/paired_v1err/probes/ratios.json :: ratios[158].ratio, ratios[158].ci95 :: sha256 14dd51402305d147 :: 162 over 188, bvc_qcd mlp 1-auc: paired geometric mean [95% interval]
\newcommand{\PairedProbeBvcQcdMlpFourthreeOverOneeighteight}{1.13~[1.05, 1.21]}  % experiments/FIGS/data/paired_v1err/probes/ratios.json :: ratios[159].ratio, ratios[159].ci95 :: sha256 14dd51402305d147 :: 43 over 188, bvc_qcd mlp 1-auc: paired geometric mean [95% interval]
\newcommand{\PairedProbeBvcQcdMlpOnesevenOverOneeighteight}{2.6~[2.2, 3.0]}  % experiments/FIGS/data/paired_v1err/probes/ratios.json :: ratios[160].ratio, ratios[160].ci95 :: sha256 14dd51402305d147 :: 17 over 188, bvc_qcd mlp 1-auc: paired geometric mean [95% interval]
\newcommand{\PairedProbeBvcQcdMlpFourthreeOverOnesixtwo}{1.01~[0.96, 1.06]}  % experiments/FIGS/data/paired_v1err/probes/ratios.json :: ratios[161].ratio, ratios[161].ci95 :: sha256 14dd51402305d147 :: 43 over 162, bvc_qcd mlp 1-auc: paired geometric mean [95% interval]
\newcommand{\PairedProbeBvcQcdMlpOnesevenOverOnesixtwo}{2.3~[2.0, 2.6]}  % experiments/FIGS/data/paired_v1err/probes/ratios.json :: ratios[162].ratio, ratios[162].ci95 :: sha256 14dd51402305d147 :: 17 over 162, bvc_qcd mlp 1-auc: paired geometric mean [95% interval]
\newcommand{\PairedProbeBvcQcdMlpOnesevenOverFourthree}{2.3~[2.0, 2.6]}  % experiments/FIGS/data/paired_v1err/probes/ratios.json :: ratios[163].ratio, ratios[163].ci95 :: sha256 14dd51402305d147 :: 17 over 43, bvc_qcd mlp 1-auc: paired geometric mean [95% interval]
\newcommand{\PairedProbeBvcQcdMlpOnesixtwoMassOverOnesixtwo}{1.06~[1.01, 1.12]}  % experiments/FIGS/data/paired_v1err/probes/ratios.json :: ratios[164].ratio, ratios[164].ci95 :: sha256 14dd51402305d147 :: 162+mass over 162, bvc_qcd mlp 1-auc: paired geometric mean [95% interval]
\newcommand{\PairedProbeBvcQcdMlpOnesevenMassOverOneseven}{1.64~[1.33, 2.02]}  % experiments/FIGS/data/paired_v1err/probes/ratios.json :: ratios[165].ratio, ratios[165].ci95 :: sha256 14dd51402305d147 :: 17+mass over 17, bvc_qcd mlp 1-auc: paired geometric mean [95% interval]
\newcommand{\PairedProbeBvcResonantLinearOnesixtwoOverOneeighteight}{1.02~[0.96, 1.08]}  % experiments/FIGS/data/paired_v1err/probes/ratios.json :: ratios[190].ratio, ratios[190].ci95 :: sha256 14dd51402305d147 :: 162 over 188, bvc_resonant linear 1-auc: paired geometric mean [95% interval]
\newcommand{\PairedProbeBvcResonantLinearFourthreeOverOneeighteight}{1.09~[1.03, 1.15]}  % experiments/FIGS/data/paired_v1err/probes/ratios.json :: ratios[191].ratio, ratios[191].ci95 :: sha256 14dd51402305d147 :: 43 over 188, bvc_resonant linear 1-auc: paired geometric mean [95% interval]
\newcommand{\PairedProbeBvcResonantLinearOnesevenOverOneeighteight}{4.0~[3.3, 4.9]}  % experiments/FIGS/data/paired_v1err/probes/ratios.json :: ratios[192].ratio, ratios[192].ci95 :: sha256 14dd51402305d147 :: 17 over 188, bvc_resonant linear 1-auc: paired geometric mean [95% interval]
\newcommand{\PairedProbeBvcResonantLinearFourthreeOverOnesixtwo}{1.07~[1.02, 1.12]}  % experiments/FIGS/data/paired_v1err/probes/ratios.json :: ratios[193].ratio, ratios[193].ci95 :: sha256 14dd51402305d147 :: 43 over 162, bvc_resonant linear 1-auc: paired geometric mean [95% interval]
\newcommand{\PairedProbeBvcResonantLinearOnesevenOverOnesixtwo}{3.9~[3.1, 4.9]}  % experiments/FIGS/data/paired_v1err/probes/ratios.json :: ratios[194].ratio, ratios[194].ci95 :: sha256 14dd51402305d147 :: 17 over 162, bvc_resonant linear 1-auc: paired geometric mean [95% interval]
\newcommand{\PairedProbeBvcResonantLinearOnesevenOverFourthree}{3.7~[3.0, 4.5]}  % experiments/FIGS/data/paired_v1err/probes/ratios.json :: ratios[195].ratio, ratios[195].ci95 :: sha256 14dd51402305d147 :: 17 over 43, bvc_resonant linear 1-auc: paired geometric mean [95% interval]
\newcommand{\PairedProbeBvcResonantLinearOnesixtwoMassOverOnesixtwo}{1.11~[1.06, 1.16]}  % experiments/FIGS/data/paired_v1err/probes/ratios.json :: ratios[196].ratio, ratios[196].ci95 :: sha256 14dd51402305d147 :: 162+mass over 162, bvc_resonant linear 1-auc: paired geometric mean [95% interval]
\newcommand{\PairedProbeBvcResonantLinearOnesevenMassOverOneseven}{3.0~[1.9, 4.7]}  % experiments/FIGS/data/paired_v1err/probes/ratios.json :: ratios[197].ratio, ratios[197].ci95 :: sha256 14dd51402305d147 :: 17+mass over 17, bvc_resonant linear 1-auc: paired geometric mean [95% interval]
\newcommand{\PairedProbeBvcResonantMlpOnesixtwoOverOneeighteight}{1.04~[0.98, 1.12]}  % experiments/FIGS/data/paired_v1err/probes/ratios.json :: ratios[222].ratio, ratios[222].ci95 :: sha256 14dd51402305d147 :: 162 over 188, bvc_resonant mlp 1-auc: paired geometric mean [95% interval]
\newcommand{\PairedProbeBvcResonantMlpFourthreeOverOneeighteight}{1.14~[1.08, 1.21]}  % experiments/FIGS/data/paired_v1err/probes/ratios.json :: ratios[223].ratio, ratios[223].ci95 :: sha256 14dd51402305d147 :: 43 over 188, bvc_resonant mlp 1-auc: paired geometric mean [95% interval]
\newcommand{\PairedProbeBvcResonantMlpOnesevenOverOneeighteight}{3.3~[2.8, 3.8]}  % experiments/FIGS/data/paired_v1err/probes/ratios.json :: ratios[224].ratio, ratios[224].ci95 :: sha256 14dd51402305d147 :: 17 over 188, bvc_resonant mlp 1-auc: paired geometric mean [95% interval]
\newcommand{\PairedProbeBvcResonantMlpFourthreeOverOnesixtwo}{1.09~[1.04, 1.14]}  % experiments/FIGS/data/paired_v1err/probes/ratios.json :: ratios[225].ratio, ratios[225].ci95 :: sha256 14dd51402305d147 :: 43 over 162, bvc_resonant mlp 1-auc: paired geometric mean [95% interval]
\newcommand{\PairedProbeBvcResonantMlpOnesevenOverOnesixtwo}{3.1~[2.6, 3.8]}  % experiments/FIGS/data/paired_v1err/probes/ratios.json :: ratios[226].ratio, ratios[226].ci95 :: sha256 14dd51402305d147 :: 17 over 162, bvc_resonant mlp 1-auc: paired geometric mean [95% interval]
\newcommand{\PairedProbeBvcResonantMlpOnesevenOverFourthree}{2.9~[2.4, 3.4]}  % experiments/FIGS/data/paired_v1err/probes/ratios.json :: ratios[227].ratio, ratios[227].ci95 :: sha256 14dd51402305d147 :: 17 over 43, bvc_resonant mlp 1-auc: paired geometric mean [95% interval]
\newcommand{\PairedProbeBvcResonantMlpOnesixtwoMassOverOnesixtwo}{1.09~[1.04, 1.14]}  % experiments/FIGS/data/paired_v1err/probes/ratios.json :: ratios[228].ratio, ratios[228].ci95 :: sha256 14dd51402305d147 :: 162+mass over 162, bvc_resonant mlp 1-auc: paired geometric mean [95% interval]
\newcommand{\PairedProbeBvcResonantMlpOnesevenMassOverOneseven}{1.95~[1.45, 2.61]}  % experiments/FIGS/data/paired_v1err/probes/ratios.json :: ratios[229].ratio, ratios[229].ci95 :: sha256 14dd51402305d147 :: 17+mass over 17, bvc_resonant mlp 1-auc: paired geometric mean [95% interval]
\newcommand{\PairedProbeEeVsMmLinearOnesevenOverFourthree}{2.3~[0.8, 6.5]}  % experiments/FIGS/data/paired_v1err/probes/ratios.json :: ratios[259].ratio, ratios[259].ci95 :: sha256 14dd51402305d147 :: 17 over 43, ee_vs_mm linear 1-auc: paired geometric mean [95% interval]
\newcommand{\PairedProbeEeVsMmLinearOnesevenMassOverOneseven}{0.23~[0.08, 0.64]}  % experiments/FIGS/data/paired_v1err/probes/ratios.json :: ratios[261].ratio, ratios[261].ci95 :: sha256 14dd51402305d147 :: 17+mass over 17, ee_vs_mm linear 1-auc: paired geometric mean [95% interval]
\newcommand{\PairedProbeEeVsMmMlpOnesevenOverFourthree}{2.8~[0.8, 9.0]}  % experiments/FIGS/data/paired_v1err/probes/ratios.json :: ratios[291].ratio, ratios[291].ci95 :: sha256 14dd51402305d147 :: 17 over 43, ee_vs_mm mlp 1-auc: paired geometric mean [95% interval]
\newcommand{\PairedProbeEeVsMmMlpOnesevenMassOverOneseven}{0.172~[0.026, 1.125]}  % experiments/FIGS/data/paired_v1err/probes/ratios.json :: ratios[293].ratio, ratios[293].ci95 :: sha256 14dd51402305d147 :: 17+mass over 17, ee_vs_mm mlp 1-auc: paired geometric mean [95% interval]
\newcommand{\PairedProbeRetainedTopologyLinearOnesixtwoOverOneeighteight}{1.00~[0.92, 1.08]}  % experiments/FIGS/data/paired_v1err/probes/ratios.json :: ratios[318].ratio, ratios[318].ci95 :: sha256 14dd51402305d147 :: 162 over 188, retained_topology linear 1-auc: paired geometric mean [95% interval]
\newcommand{\PairedProbeRetainedTopologyLinearFourthreeOverOneeighteight}{1.24~[1.15, 1.34]}  % experiments/FIGS/data/paired_v1err/probes/ratios.json :: ratios[319].ratio, ratios[319].ci95 :: sha256 14dd51402305d147 :: 43 over 188, retained_topology linear 1-auc: paired geometric mean [95% interval]
\newcommand{\PairedProbeRetainedTopologyLinearOnesevenOverOneeighteight}{1.54~[1.35, 1.75]}  % experiments/FIGS/data/paired_v1err/probes/ratios.json :: ratios[320].ratio, ratios[320].ci95 :: sha256 14dd51402305d147 :: 17 over 188, retained_topology linear 1-auc: paired geometric mean [95% interval]
\newcommand{\PairedProbeRetainedTopologyLinearFourthreeOverOnesixtwo}{1.24~[1.14, 1.34]}  % experiments/FIGS/data/paired_v1err/probes/ratios.json :: ratios[321].ratio, ratios[321].ci95 :: sha256 14dd51402305d147 :: 43 over 162, retained_topology linear 1-auc: paired geometric mean [95% interval]
\newcommand{\PairedProbeRetainedTopologyLinearOnesevenOverOnesixtwo}{1.54~[1.33, 1.78]}  % experiments/FIGS/data/paired_v1err/probes/ratios.json :: ratios[322].ratio, ratios[322].ci95 :: sha256 14dd51402305d147 :: 17 over 162, retained_topology linear 1-auc: paired geometric mean [95% interval]
\newcommand{\PairedProbeRetainedTopologyLinearOnesevenOverFourthree}{1.24~[1.08, 1.42]}  % experiments/FIGS/data/paired_v1err/probes/ratios.json :: ratios[323].ratio, ratios[323].ci95 :: sha256 14dd51402305d147 :: 17 over 43, retained_topology linear 1-auc: paired geometric mean [95% interval]
\newcommand{\PairedProbeRetainedTopologyLinearOnesixtwoMassOverOnesixtwo}{1.06~[0.99, 1.14]}  % experiments/FIGS/data/paired_v1err/probes/ratios.json :: ratios[324].ratio, ratios[324].ci95 :: sha256 14dd51402305d147 :: 162+mass over 162, retained_topology linear 1-auc: paired geometric mean [95% interval]
\newcommand{\PairedProbeRetainedTopologyLinearOnesevenMassOverOneseven}{1.12~[0.97, 1.28]}  % experiments/FIGS/data/paired_v1err/probes/ratios.json :: ratios[325].ratio, ratios[325].ci95 :: sha256 14dd51402305d147 :: 17+mass over 17, retained_topology linear 1-auc: paired geometric mean [95% interval]
\newcommand{\PairedProbeRetainedTopologyMlpOnesixtwoOverOneeighteight}{1.01~[0.87, 1.17]}  % experiments/FIGS/data/paired_v1err/probes/ratios.json :: ratios[350].ratio, ratios[350].ci95 :: sha256 14dd51402305d147 :: 162 over 188, retained_topology mlp 1-auc: paired geometric mean [95% interval]
\newcommand{\PairedProbeRetainedTopologyMlpFourthreeOverOneeighteight}{1.18~[1.03, 1.36]}  % experiments/FIGS/data/paired_v1err/probes/ratios.json :: ratios[351].ratio, ratios[351].ci95 :: sha256 14dd51402305d147 :: 43 over 188, retained_topology mlp 1-auc: paired geometric mean [95% interval]
\newcommand{\PairedProbeRetainedTopologyMlpOnesevenOverOneeighteight}{1.52~[1.34, 1.72]}  % experiments/FIGS/data/paired_v1err/probes/ratios.json :: ratios[352].ratio, ratios[352].ci95 :: sha256 14dd51402305d147 :: 17 over 188, retained_topology mlp 1-auc: paired geometric mean [95% interval]
\newcommand{\PairedProbeRetainedTopologyMlpFourthreeOverOnesixtwo}{1.17~[1.07, 1.28]}  % experiments/FIGS/data/paired_v1err/probes/ratios.json :: ratios[353].ratio, ratios[353].ci95 :: sha256 14dd51402305d147 :: 43 over 162, retained_topology mlp 1-auc: paired geometric mean [95% interval]
\newcommand{\PairedProbeRetainedTopologyMlpOnesevenOverOnesixtwo}{1.50~[1.29, 1.75]}  % experiments/FIGS/data/paired_v1err/probes/ratios.json :: ratios[354].ratio, ratios[354].ci95 :: sha256 14dd51402305d147 :: 17 over 162, retained_topology mlp 1-auc: paired geometric mean [95% interval]
\newcommand{\PairedProbeRetainedTopologyMlpOnesevenOverFourthree}{1.28~[1.12, 1.47]}  % experiments/FIGS/data/paired_v1err/probes/ratios.json :: ratios[355].ratio, ratios[355].ci95 :: sha256 14dd51402305d147 :: 17 over 43, retained_topology mlp 1-auc: paired geometric mean [95% interval]
\newcommand{\PairedProbeRetainedTopologyMlpOnesixtwoMassOverOnesixtwo}{1.04~[0.92, 1.17]}  % experiments/FIGS/data/paired_v1err/probes/ratios.json :: ratios[356].ratio, ratios[356].ci95 :: sha256 14dd51402305d147 :: 162+mass over 162, retained_topology mlp 1-auc: paired geometric mean [95% interval]
\newcommand{\PairedProbeRetainedTopologyMlpOnesevenMassOverOneseven}{1.11~[0.99, 1.23]}  % experiments/FIGS/data/paired_v1err/probes/ratios.json :: ratios[357].ratio, ratios[357].ci95 :: sha256 14dd51402305d147 :: 17+mass over 17, retained_topology mlp 1-auc: paired geometric mean [95% interval]
\newcommand{\PairedProbeVisibleContentLinearOnesixtwoOverOneeighteight}{0.95~[0.88, 1.03]}  % experiments/FIGS/data/paired_v1err/probes/ratios.json :: ratios[382].ratio, ratios[382].ci95 :: sha256 14dd51402305d147 :: 162 over 188, visible_content linear 1-auc: paired geometric mean [95% interval]
\newcommand{\PairedProbeVisibleContentLinearFourthreeOverOneeighteight}{0.99~[0.85, 1.15]}  % experiments/FIGS/data/paired_v1err/probes/ratios.json :: ratios[383].ratio, ratios[383].ci95 :: sha256 14dd51402305d147 :: 43 over 188, visible_content linear 1-auc: paired geometric mean [95% interval]
\newcommand{\PairedProbeVisibleContentLinearOnesevenOverOneeighteight}{2.8~[2.1, 3.7]}  % experiments/FIGS/data/paired_v1err/probes/ratios.json :: ratios[384].ratio, ratios[384].ci95 :: sha256 14dd51402305d147 :: 17 over 188, visible_content linear 1-auc: paired geometric mean [95% interval]
\newcommand{\PairedProbeVisibleContentLinearFourthreeOverOnesixtwo}{1.04~[0.89, 1.20]}  % experiments/FIGS/data/paired_v1err/probes/ratios.json :: ratios[385].ratio, ratios[385].ci95 :: sha256 14dd51402305d147 :: 43 over 162, visible_content linear 1-auc: paired geometric mean [95% interval]
\newcommand{\PairedProbeVisibleContentLinearOnesevenOverOnesixtwo}{2.9~[2.3, 3.8]}  % experiments/FIGS/data/paired_v1err/probes/ratios.json :: ratios[386].ratio, ratios[386].ci95 :: sha256 14dd51402305d147 :: 17 over 162, visible_content linear 1-auc: paired geometric mean [95% interval]
\newcommand{\PairedProbeVisibleContentLinearOnesevenOverFourthree}{2.8~[2.2, 3.7]}  % experiments/FIGS/data/paired_v1err/probes/ratios.json :: ratios[387].ratio, ratios[387].ci95 :: sha256 14dd51402305d147 :: 17 over 43, visible_content linear 1-auc: paired geometric mean [95% interval]
\newcommand{\PairedProbeVisibleContentLinearOnesixtwoMassOverOnesixtwo}{1.10~[0.98, 1.23]}  % experiments/FIGS/data/paired_v1err/probes/ratios.json :: ratios[388].ratio, ratios[388].ci95 :: sha256 14dd51402305d147 :: 162+mass over 162, visible_content linear 1-auc: paired geometric mean [95% interval]
\newcommand{\PairedProbeVisibleContentLinearOnesevenMassOverOneseven}{1.44~[1.26, 1.65]}  % experiments/FIGS/data/paired_v1err/probes/ratios.json :: ratios[389].ratio, ratios[389].ci95 :: sha256 14dd51402305d147 :: 17+mass over 17, visible_content linear 1-auc: paired geometric mean [95% interval]
\newcommand{\PairedProbeVisibleContentLinearRandomOverOneeighteight}{1.31~[1.01, 1.70]}  % experiments/FIGS/data/paired_v1err/probes/ratios.json :: ratios[390].ratio, ratios[390].ci95 :: sha256 14dd51402305d147 :: random over 188, visible_content linear 1-auc: paired geometric mean [95% interval]
\newcommand{\PairedProbeVisibleContentLinearRandomOverOnesixtwo}{1.41~[1.17, 1.70]}  % experiments/FIGS/data/paired_v1err/probes/ratios.json :: ratios[391].ratio, ratios[391].ci95 :: sha256 14dd51402305d147 :: random over 162, visible_content linear 1-auc: paired geometric mean [95% interval]
\newcommand{\PairedProbeVisibleContentLinearRandomOverFourthree}{1.33~[1.12, 1.57]}  % experiments/FIGS/data/paired_v1err/probes/ratios.json :: ratios[392].ratio, ratios[392].ci95 :: sha256 14dd51402305d147 :: random over 43, visible_content linear 1-auc: paired geometric mean [95% interval]
\newcommand{\PairedProbeVisibleContentLinearRandomOverOneseven}{0.48~[0.38, 0.60]}  % experiments/FIGS/data/paired_v1err/probes/ratios.json :: ratios[393].ratio, ratios[393].ci95 :: sha256 14dd51402305d147 :: random over 17, visible_content linear 1-auc: paired geometric mean [95% interval]
\newcommand{\PairedProbeVisibleContentLinearRandomDrawTwoOverRandomDrawOne}{0.79~[0.64, 0.98]}  % experiments/FIGS/data/paired_v1err/probes/ratios.json :: ratios[394].ratio, ratios[394].ci95 :: sha256 14dd51402305d147 :: random draw 2 over random draw 1, visible_content linear 1-auc: paired geometric mean [95% interval]
\newcommand{\PairedProbeVisibleContentLinearRandomDrawThreeOverRandomDrawOne}{0.86~[0.71, 1.05]}  % experiments/FIGS/data/paired_v1err/probes/ratios.json :: ratios[395].ratio, ratios[395].ci95 :: sha256 14dd51402305d147 :: random draw 3 over random draw 1, visible_content linear 1-auc: paired geometric mean [95% interval]
\newcommand{\PairedProbeVisibleContentLinearRandomDrawThreeOverRandomDrawTwo}{1.09~[0.88, 1.35]}  % experiments/FIGS/data/paired_v1err/probes/ratios.json :: ratios[396].ratio, ratios[396].ci95 :: sha256 14dd51402305d147 :: random draw 3 over random draw 2, visible_content linear 1-auc: paired geometric mean [95% interval]
\newcommand{\PairedProbeVisibleContentMlpOnesixtwoOverOneeighteight}{0.92~[0.85, 1.01]}  % experiments/FIGS/data/paired_v1err/probes/ratios.json :: ratios[442].ratio, ratios[442].ci95 :: sha256 14dd51402305d147 :: 162 over 188, visible_content mlp 1-auc: paired geometric mean [95% interval]
\newcommand{\PairedProbeVisibleContentMlpFourthreeOverOneeighteight}{0.93~[0.80, 1.07]}  % experiments/FIGS/data/paired_v1err/probes/ratios.json :: ratios[443].ratio, ratios[443].ci95 :: sha256 14dd51402305d147 :: 43 over 188, visible_content mlp 1-auc: paired geometric mean [95% interval]
\newcommand{\PairedProbeVisibleContentMlpOnesevenOverOneeighteight}{2.3~[1.7, 3.1]}  % experiments/FIGS/data/paired_v1err/probes/ratios.json :: ratios[444].ratio, ratios[444].ci95 :: sha256 14dd51402305d147 :: 17 over 188, visible_content mlp 1-auc: paired geometric mean [95% interval]
\newcommand{\PairedProbeVisibleContentMlpFourthreeOverOnesixtwo}{1.01~[0.89, 1.14]}  % experiments/FIGS/data/paired_v1err/probes/ratios.json :: ratios[445].ratio, ratios[445].ci95 :: sha256 14dd51402305d147 :: 43 over 162, visible_content mlp 1-auc: paired geometric mean [95% interval]
\newcommand{\PairedProbeVisibleContentMlpOnesevenOverOnesixtwo}{2.5~[1.9, 3.2]}  % experiments/FIGS/data/paired_v1err/probes/ratios.json :: ratios[446].ratio, ratios[446].ci95 :: sha256 14dd51402305d147 :: 17 over 162, visible_content mlp 1-auc: paired geometric mean [95% interval]
\newcommand{\PairedProbeVisibleContentMlpOnesevenOverFourthree}{2.5~[2.0, 3.1]}  % experiments/FIGS/data/paired_v1err/probes/ratios.json :: ratios[447].ratio, ratios[447].ci95 :: sha256 14dd51402305d147 :: 17 over 43, visible_content mlp 1-auc: paired geometric mean [95% interval]
\newcommand{\PairedProbeVisibleContentMlpOnesixtwoMassOverOnesixtwo}{1.08~[0.94, 1.24]}  % experiments/FIGS/data/paired_v1err/probes/ratios.json :: ratios[448].ratio, ratios[448].ci95 :: sha256 14dd51402305d147 :: 162+mass over 162, visible_content mlp 1-auc: paired geometric mean [95% interval]
\newcommand{\PairedProbeVisibleContentMlpOnesevenMassOverOneseven}{1.41~[1.19, 1.67]}  % experiments/FIGS/data/paired_v1err/probes/ratios.json :: ratios[449].ratio, ratios[449].ci95 :: sha256 14dd51402305d147 :: 17+mass over 17, visible_content mlp 1-auc: paired geometric mean [95% interval]
\newcommand{\PairedProbeVisibleContentMlpRandomOverOneeighteight}{1.27~[1.01, 1.59]}  % experiments/FIGS/data/paired_v1err/probes/ratios.json :: ratios[450].ratio, ratios[450].ci95 :: sha256 14dd51402305d147 :: random over 188, visible_content mlp 1-auc: paired geometric mean [95% interval]
\newcommand{\PairedProbeVisibleContentMlpRandomOverOnesixtwo}{1.37~[1.12, 1.68]}  % experiments/FIGS/data/paired_v1err/probes/ratios.json :: ratios[451].ratio, ratios[451].ci95 :: sha256 14dd51402305d147 :: random over 162, visible_content mlp 1-auc: paired geometric mean [95% interval]
\newcommand{\PairedProbeVisibleContentMlpRandomOverFourthree}{1.35~[1.13, 1.61]}  % experiments/FIGS/data/paired_v1err/probes/ratios.json :: ratios[452].ratio, ratios[452].ci95 :: sha256 14dd51402305d147 :: random over 43, visible_content mlp 1-auc: paired geometric mean [95% interval]
\newcommand{\PairedProbeVisibleContentMlpRandomOverOneseven}{0.56~[0.45, 0.70]}  % experiments/FIGS/data/paired_v1err/probes/ratios.json :: ratios[453].ratio, ratios[453].ci95 :: sha256 14dd51402305d147 :: random over 17, visible_content mlp 1-auc: paired geometric mean [95% interval]
\newcommand{\PairedProbeVisibleContentMlpRandomDrawTwoOverRandomDrawOne}{0.88~[0.71, 1.09]}  % experiments/FIGS/data/paired_v1err/probes/ratios.json :: ratios[454].ratio, ratios[454].ci95 :: sha256 14dd51402305d147 :: random draw 2 over random draw 1, visible_content mlp 1-auc: paired geometric mean [95% interval]
\newcommand{\PairedProbeVisibleContentMlpRandomDrawThreeOverRandomDrawOne}{0.94~[0.75, 1.18]}  % experiments/FIGS/data/paired_v1err/probes/ratios.json :: ratios[455].ratio, ratios[455].ci95 :: sha256 14dd51402305d147 :: random draw 3 over random draw 1, visible_content mlp 1-auc: paired geometric mean [95% interval]
\newcommand{\PairedProbeVisibleContentMlpRandomDrawThreeOverRandomDrawTwo}{1.06~[0.84, 1.35]}  % experiments/FIGS/data/paired_v1err/probes/ratios.json :: ratios[456].ratio, ratios[456].ci95 :: sha256 14dd51402305d147 :: random draw 3 over random draw 2, visible_content mlp 1-auc: paired geometric mean [95% interval]
\newcommand{\PairedFtJciiEThreeOnesixtwoOverOneeighteight}{0.95~[0.87, 1.05]}  % experiments/FIGS/data/paired_v1err/finetune/ratios.json :: ratios[0].ratio, ratios[0].ci95 :: sha256 3e5cf25e859b68c2 :: 162 over 188, leg1 N1000 1-macro_auc: paired geometric mean [95% interval]
\newcommand{\PairedFtJciiEThreeFourthreeOverOneeighteight}{0.84~[0.77, 0.91]}  % experiments/FIGS/data/paired_v1err/finetune/ratios.json :: ratios[1].ratio, ratios[1].ci95 :: sha256 3e5cf25e859b68c2 :: 43 over 188, leg1 N1000 1-macro_auc: paired geometric mean [95% interval]
\newcommand{\PairedFtJciiEThreeOnesevenOverOneeighteight}{0.97~[0.86, 1.09]}  % experiments/FIGS/data/paired_v1err/finetune/ratios.json :: ratios[2].ratio, ratios[2].ci95 :: sha256 3e5cf25e859b68c2 :: 17 over 188, leg1 N1000 1-macro_auc: paired geometric mean [95% interval]
\newcommand{\PairedFtJciiEThreeFourthreeOverOnesixtwo}{0.88~[0.80, 0.97]}  % experiments/FIGS/data/paired_v1err/finetune/ratios.json :: ratios[3].ratio, ratios[3].ci95 :: sha256 3e5cf25e859b68c2 :: 43 over 162, leg1 N1000 1-macro_auc: paired geometric mean [95% interval]
\newcommand{\PairedFtJciiEThreeOnesevenOverOnesixtwo}{1.02~[0.88, 1.18]}  % experiments/FIGS/data/paired_v1err/finetune/ratios.json :: ratios[4].ratio, ratios[4].ci95 :: sha256 3e5cf25e859b68c2 :: 17 over 162, leg1 N1000 1-macro_auc: paired geometric mean [95% interval]
\newcommand{\PairedFtJciiEThreeOnesevenOverFourthree}{1.16~[1.08, 1.24]}  % experiments/FIGS/data/paired_v1err/finetune/ratios.json :: ratios[5].ratio, ratios[5].ci95 :: sha256 3e5cf25e859b68c2 :: 17 over 43, leg1 N1000 1-macro_auc: paired geometric mean [95% interval]
\newcommand{\PairedFtJciiEThreeOnesixtwoMassOverOnesixtwo}{0.85~[0.77, 0.93]}  % experiments/FIGS/data/paired_v1err/finetune/ratios.json :: ratios[6].ratio, ratios[6].ci95 :: sha256 3e5cf25e859b68c2 :: 162+mass over 162, leg1 N1000 1-macro_auc: paired geometric mean [95% interval]
\newcommand{\PairedFtJciiEThreeOnesevenMassOverOneseven}{1.00~[0.86, 1.17]}  % experiments/FIGS/data/paired_v1err/finetune/ratios.json :: ratios[7].ratio, ratios[7].ci95 :: sha256 3e5cf25e859b68c2 :: 17+mass over 17, leg1 N1000 1-macro_auc: paired geometric mean [95% interval]
\newcommand{\PairedFtJciiEThreeRandomOverOneeighteight}{0.82~[0.73, 0.93]}  % experiments/FIGS/data/paired_v1err/finetune/ratios.json :: ratios[8].ratio, ratios[8].ci95 :: sha256 3e5cf25e859b68c2 :: random over 188, leg1 N1000 1-macro_auc: paired geometric mean [95% interval]
\newcommand{\PairedFtJciiEThreeRandomOverOnesixtwo}{0.90~[0.85, 0.96]}  % experiments/FIGS/data/paired_v1err/finetune/ratios.json :: ratios[9].ratio, ratios[9].ci95 :: sha256 3e5cf25e859b68c2 :: random over 162, leg1 N1000 1-macro_auc: paired geometric mean [95% interval]
\newcommand{\PairedFtJciiEThreeRandomOverFourthree}{0.99~[0.75, 1.30]}  % experiments/FIGS/data/paired_v1err/finetune/ratios.json :: ratios[10].ratio, ratios[10].ci95 :: sha256 3e5cf25e859b68c2 :: random over 43, leg1 N1000 1-macro_auc: paired geometric mean [95% interval]
\newcommand{\PairedFtJciiEThreeRandomOverOneseven}{0.86~[0.58, 1.27]}  % experiments/FIGS/data/paired_v1err/finetune/ratios.json :: ratios[11].ratio, ratios[11].ci95 :: sha256 3e5cf25e859b68c2 :: random over 17, leg1 N1000 1-macro_auc: paired geometric mean [95% interval]
\newcommand{\PairedFtJciiEThreeRandomDrawTwoOverRandomDrawOne}{1.09~[1.07, 1.10]}  % experiments/FIGS/data/paired_v1err/finetune/ratios.json :: ratios[12].ratio, ratios[12].ci95 :: sha256 3e5cf25e859b68c2 :: random draw 2 over random draw 1, leg1 N1000 1-macro_auc: paired geometric mean [95% interval]
\newcommand{\PairedFtJciiEThreeRandomDrawThreeOverRandomDrawOne}{1.04~[1.02, 1.05]}  % experiments/FIGS/data/paired_v1err/finetune/ratios.json :: ratios[13].ratio, ratios[13].ci95 :: sha256 3e5cf25e859b68c2 :: random draw 3 over random draw 1, leg1 N1000 1-macro_auc: paired geometric mean [95% interval]
\newcommand{\PairedFtJciiEThreeRandomDrawThreeOverRandomDrawTwo}{0.95~[0.94, 0.97]}  % experiments/FIGS/data/paired_v1err/finetune/ratios.json :: ratios[14].ratio, ratios[14].ci95 :: sha256 3e5cf25e859b68c2 :: random draw 3 over random draw 2, leg1 N1000 1-macro_auc: paired geometric mean [95% interval]
\newcommand{\PairedFtJciiEFourOnesixtwoOverOneeighteight}{0.99~[0.97, 1.02]}  % experiments/FIGS/data/paired_v1err/finetune/ratios.json :: ratios[15].ratio, ratios[15].ci95 :: sha256 3e5cf25e859b68c2 :: 162 over 188, leg1 N10000 1-macro_auc: paired geometric mean [95% interval]
\newcommand{\PairedFtJciiEFourFourthreeOverOneeighteight}{1.20~[1.17, 1.23]}  % experiments/FIGS/data/paired_v1err/finetune/ratios.json :: ratios[16].ratio, ratios[16].ci95 :: sha256 3e5cf25e859b68c2 :: 43 over 188, leg1 N10000 1-macro_auc: paired geometric mean [95% interval]
\newcommand{\PairedFtJciiEFourOnesevenOverOneeighteight}{1.32~[1.29, 1.34]}  % experiments/FIGS/data/paired_v1err/finetune/ratios.json :: ratios[17].ratio, ratios[17].ci95 :: sha256 3e5cf25e859b68c2 :: 17 over 188, leg1 N10000 1-macro_auc: paired geometric mean [95% interval]
\newcommand{\PairedFtJciiEFourFourthreeOverOnesixtwo}{1.21~[1.19, 1.23]}  % experiments/FIGS/data/paired_v1err/finetune/ratios.json :: ratios[18].ratio, ratios[18].ci95 :: sha256 3e5cf25e859b68c2 :: 43 over 162, leg1 N10000 1-macro_auc: paired geometric mean [95% interval]
\newcommand{\PairedFtJciiEFourOnesevenOverOnesixtwo}{1.32~[1.29, 1.36]}  % experiments/FIGS/data/paired_v1err/finetune/ratios.json :: ratios[19].ratio, ratios[19].ci95 :: sha256 3e5cf25e859b68c2 :: 17 over 162, leg1 N10000 1-macro_auc: paired geometric mean [95% interval]
\newcommand{\PairedFtJciiEFourOnesevenOverFourthree}{1.10~[1.07, 1.13]}  % experiments/FIGS/data/paired_v1err/finetune/ratios.json :: ratios[20].ratio, ratios[20].ci95 :: sha256 3e5cf25e859b68c2 :: 17 over 43, leg1 N10000 1-macro_auc: paired geometric mean [95% interval]
\newcommand{\PairedFtJciiEFourOnesixtwoMassOverOnesixtwo}{1.00~[0.94, 1.06]}  % experiments/FIGS/data/paired_v1err/finetune/ratios.json :: ratios[21].ratio, ratios[21].ci95 :: sha256 3e5cf25e859b68c2 :: 162+mass over 162, leg1 N10000 1-macro_auc: paired geometric mean [95% interval]
\newcommand{\PairedFtJciiEFourOnesevenMassOverOneseven}{0.990~[0.976, 1.004]}  % experiments/FIGS/data/paired_v1err/finetune/ratios.json :: ratios[22].ratio, ratios[22].ci95 :: sha256 3e5cf25e859b68c2 :: 17+mass over 17, leg1 N10000 1-macro_auc: paired geometric mean [95% interval]
\newcommand{\PairedFtJciiEFourRandomOverOneeighteight}{1.15~[1.09, 1.22]}  % experiments/FIGS/data/paired_v1err/finetune/ratios.json :: ratios[23].ratio, ratios[23].ci95 :: sha256 3e5cf25e859b68c2 :: random over 188, leg1 N10000 1-macro_auc: paired geometric mean [95% interval]
\newcommand{\PairedFtJciiEFourRandomOverOnesixtwo}{1.18~[1.13, 1.22]}  % experiments/FIGS/data/paired_v1err/finetune/ratios.json :: ratios[24].ratio, ratios[24].ci95 :: sha256 3e5cf25e859b68c2 :: random over 162, leg1 N10000 1-macro_auc: paired geometric mean [95% interval]
\newcommand{\PairedFtJciiEFourRandomOverFourthree}{0.97~[0.96, 0.99]}  % experiments/FIGS/data/paired_v1err/finetune/ratios.json :: ratios[25].ratio, ratios[25].ci95 :: sha256 3e5cf25e859b68c2 :: random over 43, leg1 N10000 1-macro_auc: paired geometric mean [95% interval]
\newcommand{\PairedFtJciiEFourRandomOverOneseven}{0.88~[0.82, 0.95]}  % experiments/FIGS/data/paired_v1err/finetune/ratios.json :: ratios[26].ratio, ratios[26].ci95 :: sha256 3e5cf25e859b68c2 :: random over 17, leg1 N10000 1-macro_auc: paired geometric mean [95% interval]
\newcommand{\PairedFtJciiEFourRandomDrawTwoOverRandomDrawOne}{1.00~[0.98, 1.02]}  % experiments/FIGS/data/paired_v1err/finetune/ratios.json :: ratios[27].ratio, ratios[27].ci95 :: sha256 3e5cf25e859b68c2 :: random draw 2 over random draw 1, leg1 N10000 1-macro_auc: paired geometric mean [95% interval]
\newcommand{\PairedFtJciiEFourRandomDrawThreeOverRandomDrawOne}{0.97~[0.95, 0.98]}  % experiments/FIGS/data/paired_v1err/finetune/ratios.json :: ratios[28].ratio, ratios[28].ci95 :: sha256 3e5cf25e859b68c2 :: random draw 3 over random draw 1, leg1 N10000 1-macro_auc: paired geometric mean [95% interval]
\newcommand{\PairedFtJciiEFourRandomDrawThreeOverRandomDrawTwo}{0.97~[0.95, 0.98]}  % experiments/FIGS/data/paired_v1err/finetune/ratios.json :: ratios[29].ratio, ratios[29].ci95 :: sha256 3e5cf25e859b68c2 :: random draw 3 over random draw 2, leg1 N10000 1-macro_auc: paired geometric mean [95% interval]
\newcommand{\PairedFtJciiEFiveOnesixtwoOverOneeighteight}{1.00~[0.98, 1.01]}  % experiments/FIGS/data/paired_v1err/finetune/ratios.json :: ratios[30].ratio, ratios[30].ci95 :: sha256 3e5cf25e859b68c2 :: 162 over 188, leg1 N100000 1-macro_auc: paired geometric mean [95% interval]
\newcommand{\PairedFtJciiEFiveFourthreeOverOneeighteight}{1.09~[1.06, 1.11]}  % experiments/FIGS/data/paired_v1err/finetune/ratios.json :: ratios[31].ratio, ratios[31].ci95 :: sha256 3e5cf25e859b68c2 :: 43 over 188, leg1 N100000 1-macro_auc: paired geometric mean [95% interval]
\newcommand{\PairedFtJciiEFiveOnesevenOverOneeighteight}{1.22~[1.20, 1.23]}  % experiments/FIGS/data/paired_v1err/finetune/ratios.json :: ratios[32].ratio, ratios[32].ci95 :: sha256 3e5cf25e859b68c2 :: 17 over 188, leg1 N100000 1-macro_auc: paired geometric mean [95% interval]
\newcommand{\PairedFtJciiEFiveFourthreeOverOnesixtwo}{1.09~[1.07, 1.11]}  % experiments/FIGS/data/paired_v1err/finetune/ratios.json :: ratios[33].ratio, ratios[33].ci95 :: sha256 3e5cf25e859b68c2 :: 43 over 162, leg1 N100000 1-macro_auc: paired geometric mean [95% interval]
\newcommand{\PairedFtJciiEFiveOnesevenOverOnesixtwo}{1.22~[1.20, 1.24]}  % experiments/FIGS/data/paired_v1err/finetune/ratios.json :: ratios[34].ratio, ratios[34].ci95 :: sha256 3e5cf25e859b68c2 :: 17 over 162, leg1 N100000 1-macro_auc: paired geometric mean [95% interval]
\newcommand{\PairedFtJciiEFiveOnesevenOverFourthree}{1.12~[1.10, 1.14]}  % experiments/FIGS/data/paired_v1err/finetune/ratios.json :: ratios[35].ratio, ratios[35].ci95 :: sha256 3e5cf25e859b68c2 :: 17 over 43, leg1 N100000 1-macro_auc: paired geometric mean [95% interval]
\newcommand{\PairedFtJciiEFiveOnesixtwoMassOverOnesixtwo}{1.018~[1.002, 1.035]}  % experiments/FIGS/data/paired_v1err/finetune/ratios.json :: ratios[36].ratio, ratios[36].ci95 :: sha256 3e5cf25e859b68c2 :: 162+mass over 162, leg1 N100000 1-macro_auc: paired geometric mean [95% interval]
\newcommand{\PairedFtJciiEFiveOnesevenMassOverOneseven}{1.03~[1.01, 1.06]}  % experiments/FIGS/data/paired_v1err/finetune/ratios.json :: ratios[37].ratio, ratios[37].ci95 :: sha256 3e5cf25e859b68c2 :: 17+mass over 17, leg1 N100000 1-macro_auc: paired geometric mean [95% interval]
\newcommand{\PairedFtJciiEFiveRandomOverOneeighteight}{1.15~[1.09, 1.20]}  % experiments/FIGS/data/paired_v1err/finetune/ratios.json :: ratios[38].ratio, ratios[38].ci95 :: sha256 3e5cf25e859b68c2 :: random over 188, leg1 N100000 1-macro_auc: paired geometric mean [95% interval]
\newcommand{\PairedFtJciiEFiveRandomOverOnesixtwo}{1.15~[1.09, 1.20]}  % experiments/FIGS/data/paired_v1err/finetune/ratios.json :: ratios[39].ratio, ratios[39].ci95 :: sha256 3e5cf25e859b68c2 :: random over 162, leg1 N100000 1-macro_auc: paired geometric mean [95% interval]
\newcommand{\PairedFtJciiEFiveRandomOverFourthree}{1.05~[1.04, 1.06]}  % experiments/FIGS/data/paired_v1err/finetune/ratios.json :: ratios[40].ratio, ratios[40].ci95 :: sha256 3e5cf25e859b68c2 :: random over 43, leg1 N100000 1-macro_auc: paired geometric mean [95% interval]
\newcommand{\PairedFtJciiEFiveRandomOverOneseven}{0.95~[0.93, 0.96]}  % experiments/FIGS/data/paired_v1err/finetune/ratios.json :: ratios[41].ratio, ratios[41].ci95 :: sha256 3e5cf25e859b68c2 :: random over 17, leg1 N100000 1-macro_auc: paired geometric mean [95% interval]
\newcommand{\PairedFtJciiEFiveRandomDrawTwoOverRandomDrawOne}{0.98~[0.97, 0.99]}  % experiments/FIGS/data/paired_v1err/finetune/ratios.json :: ratios[42].ratio, ratios[42].ci95 :: sha256 3e5cf25e859b68c2 :: random draw 2 over random draw 1, leg1 N100000 1-macro_auc: paired geometric mean [95% interval]
\newcommand{\PairedFtJciiEFiveRandomDrawThreeOverRandomDrawOne}{1.00~[0.99, 1.02]}  % experiments/FIGS/data/paired_v1err/finetune/ratios.json :: ratios[43].ratio, ratios[43].ci95 :: sha256 3e5cf25e859b68c2 :: random draw 3 over random draw 1, leg1 N100000 1-macro_auc: paired geometric mean [95% interval]
\newcommand{\PairedFtJciiEFiveRandomDrawThreeOverRandomDrawTwo}{1.03~[1.01, 1.04]}  % experiments/FIGS/data/paired_v1err/finetune/ratios.json :: ratios[44].ratio, ratios[44].ci95 :: sha256 3e5cf25e859b68c2 :: random draw 3 over random draw 2, leg1 N100000 1-macro_auc: paired geometric mean [95% interval]
\newcommand{\PairedFtJciiESixOnesixtwoOverOneeighteight}{1.002~[0.998, 1.005]}  % experiments/FIGS/data/paired_v1err/finetune/ratios.json :: ratios[45].ratio, ratios[45].ci95 :: sha256 3e5cf25e859b68c2 :: 162 over 188, leg1 N1000000 1-macro_auc: paired geometric mean [95% interval]
\newcommand{\PairedFtJciiESixFourthreeOverOneeighteight}{1.045~[1.037, 1.053]}  % experiments/FIGS/data/paired_v1err/finetune/ratios.json :: ratios[46].ratio, ratios[46].ci95 :: sha256 3e5cf25e859b68c2 :: 43 over 188, leg1 N1000000 1-macro_auc: paired geometric mean [95% interval]
\newcommand{\PairedFtJciiESixOnesevenOverOneeighteight}{1.14~[1.13, 1.16]}  % experiments/FIGS/data/paired_v1err/finetune/ratios.json :: ratios[47].ratio, ratios[47].ci95 :: sha256 3e5cf25e859b68c2 :: 17 over 188, leg1 N1000000 1-macro_auc: paired geometric mean [95% interval]
\newcommand{\PairedFtJciiESixFourthreeOverOnesixtwo}{1.044~[1.036, 1.051]}  % experiments/FIGS/data/paired_v1err/finetune/ratios.json :: ratios[48].ratio, ratios[48].ci95 :: sha256 3e5cf25e859b68c2 :: 43 over 162, leg1 N1000000 1-macro_auc: paired geometric mean [95% interval]
\newcommand{\PairedFtJciiESixOnesevenOverOnesixtwo}{1.14~[1.13, 1.15]}  % experiments/FIGS/data/paired_v1err/finetune/ratios.json :: ratios[49].ratio, ratios[49].ci95 :: sha256 3e5cf25e859b68c2 :: 17 over 162, leg1 N1000000 1-macro_auc: paired geometric mean [95% interval]
\newcommand{\PairedFtJciiESixOnesevenOverFourthree}{1.09~[1.08, 1.11]}  % experiments/FIGS/data/paired_v1err/finetune/ratios.json :: ratios[50].ratio, ratios[50].ci95 :: sha256 3e5cf25e859b68c2 :: 17 over 43, leg1 N1000000 1-macro_auc: paired geometric mean [95% interval]
\newcommand{\PairedFtJciiESixOnesixtwoMassOverOnesixtwo}{1.019~[1.014, 1.023]}  % experiments/FIGS/data/paired_v1err/finetune/ratios.json :: ratios[51].ratio, ratios[51].ci95 :: sha256 3e5cf25e859b68c2 :: 162+mass over 162, leg1 N1000000 1-macro_auc: paired geometric mean [95% interval]
\newcommand{\PairedFtJciiESixOnesevenMassOverOneseven}{1.03~[1.02, 1.04]}  % experiments/FIGS/data/paired_v1err/finetune/ratios.json :: ratios[52].ratio, ratios[52].ci95 :: sha256 3e5cf25e859b68c2 :: 17+mass over 17, leg1 N1000000 1-macro_auc: paired geometric mean [95% interval]
\newcommand{\PairedFtJciiESixRandomOverOneeighteight}{1.073~[1.067, 1.080]}  % experiments/FIGS/data/paired_v1err/finetune/ratios.json :: ratios[53].ratio, ratios[53].ci95 :: sha256 3e5cf25e859b68c2 :: random over 188, leg1 N1000000 1-macro_auc: paired geometric mean [95% interval]
\newcommand{\PairedFtJciiESixRandomOverOnesixtwo}{1.072~[1.065, 1.078]}  % experiments/FIGS/data/paired_v1err/finetune/ratios.json :: ratios[54].ratio, ratios[54].ci95 :: sha256 3e5cf25e859b68c2 :: random over 162, leg1 N1000000 1-macro_auc: paired geometric mean [95% interval]
\newcommand{\PairedFtJciiESixRandomOverFourthree}{1.03~[1.02, 1.04]}  % experiments/FIGS/data/paired_v1err/finetune/ratios.json :: ratios[55].ratio, ratios[55].ci95 :: sha256 3e5cf25e859b68c2 :: random over 43, leg1 N1000000 1-macro_auc: paired geometric mean [95% interval]
\newcommand{\PairedFtJciiESixRandomOverOneseven}{0.94~[0.93, 0.95]}  % experiments/FIGS/data/paired_v1err/finetune/ratios.json :: ratios[56].ratio, ratios[56].ci95 :: sha256 3e5cf25e859b68c2 :: random over 17, leg1 N1000000 1-macro_auc: paired geometric mean [95% interval]
\newcommand{\PairedFtJciiESixRandomDrawTwoOverRandomDrawOne}{1.00~[0.99, 1.01]}  % experiments/FIGS/data/paired_v1err/finetune/ratios.json :: ratios[57].ratio, ratios[57].ci95 :: sha256 3e5cf25e859b68c2 :: random draw 2 over random draw 1, leg1 N1000000 1-macro_auc: paired geometric mean [95% interval]
\newcommand{\PairedFtJciiESixRandomDrawThreeOverRandomDrawOne}{0.996~[0.989, 1.003]}  % experiments/FIGS/data/paired_v1err/finetune/ratios.json :: ratios[58].ratio, ratios[58].ci95 :: sha256 3e5cf25e859b68c2 :: random draw 3 over random draw 1, leg1 N1000000 1-macro_auc: paired geometric mean [95% interval]
\newcommand{\PairedFtJciiESixRandomDrawThreeOverRandomDrawTwo}{0.997~[0.989, 1.004]}  % experiments/FIGS/data/paired_v1err/finetune/ratios.json :: ratios[59].ratio, ratios[59].ci95 :: sha256 3e5cf25e859b68c2 :: random draw 3 over random draw 2, leg1 N1000000 1-macro_auc: paired geometric mean [95% interval]
\newcommand{\PairedFtJcEThreeOnesixtwoOverOneeighteight}{0.96~[0.83, 1.11]}  % experiments/FIGS/data/paired_v1err/finetune/ratios.json :: ratios[60].ratio, ratios[60].ci95 :: sha256 3e5cf25e859b68c2 :: 162 over 188, leg2 N1000 1-macro_auc: paired geometric mean [95% interval]
\newcommand{\PairedFtJcEThreeFourthreeOverOneeighteight}{1.01~[0.84, 1.21]}  % experiments/FIGS/data/paired_v1err/finetune/ratios.json :: ratios[61].ratio, ratios[61].ci95 :: sha256 3e5cf25e859b68c2 :: 43 over 188, leg2 N1000 1-macro_auc: paired geometric mean [95% interval]
\newcommand{\PairedFtJcEThreeOnesevenOverOneeighteight}{1.07~[0.93, 1.23]}  % experiments/FIGS/data/paired_v1err/finetune/ratios.json :: ratios[62].ratio, ratios[62].ci95 :: sha256 3e5cf25e859b68c2 :: 17 over 188, leg2 N1000 1-macro_auc: paired geometric mean [95% interval]
\newcommand{\PairedFtJcEThreeFourthreeOverOnesixtwo}{1.05~[0.96, 1.15]}  % experiments/FIGS/data/paired_v1err/finetune/ratios.json :: ratios[63].ratio, ratios[63].ci95 :: sha256 3e5cf25e859b68c2 :: 43 over 162, leg2 N1000 1-macro_auc: paired geometric mean [95% interval]
\newcommand{\PairedFtJcEThreeOnesevenOverOnesixtwo}{1.12~[1.07, 1.17]}  % experiments/FIGS/data/paired_v1err/finetune/ratios.json :: ratios[64].ratio, ratios[64].ci95 :: sha256 3e5cf25e859b68c2 :: 17 over 162, leg2 N1000 1-macro_auc: paired geometric mean [95% interval]
\newcommand{\PairedFtJcEThreeOnesevenOverFourthree}{1.06~[0.98, 1.15]}  % experiments/FIGS/data/paired_v1err/finetune/ratios.json :: ratios[65].ratio, ratios[65].ci95 :: sha256 3e5cf25e859b68c2 :: 17 over 43, leg2 N1000 1-macro_auc: paired geometric mean [95% interval]
\newcommand{\PairedFtJcEThreeOnesixtwoMassOverOnesixtwo}{0.91~[0.88, 0.94]}  % experiments/FIGS/data/paired_v1err/finetune/ratios.json :: ratios[66].ratio, ratios[66].ci95 :: sha256 3e5cf25e859b68c2 :: 162+mass over 162, leg2 N1000 1-macro_auc: paired geometric mean [95% interval]
\newcommand{\PairedFtJcEThreeOnesevenMassOverOneseven}{1.02~[0.92, 1.13]}  % experiments/FIGS/data/paired_v1err/finetune/ratios.json :: ratios[67].ratio, ratios[67].ci95 :: sha256 3e5cf25e859b68c2 :: 17+mass over 17, leg2 N1000 1-macro_auc: paired geometric mean [95% interval]
\newcommand{\PairedFtJcEThreeRandomOverOneeighteight}{1.01~[0.78, 1.32]}  % experiments/FIGS/data/paired_v1err/finetune/ratios.json :: ratios[68].ratio, ratios[68].ci95 :: sha256 3e5cf25e859b68c2 :: random over 188, leg2 N1000 1-macro_auc: paired geometric mean [95% interval]
\newcommand{\PairedFtJcEThreeRandomOverOnesixtwo}{1.12~[1.04, 1.20]}  % experiments/FIGS/data/paired_v1err/finetune/ratios.json :: ratios[69].ratio, ratios[69].ci95 :: sha256 3e5cf25e859b68c2 :: random over 162, leg2 N1000 1-macro_auc: paired geometric mean [95% interval]
\newcommand{\PairedFtJcEThreeRandomOverFourthree}{1.06~[0.79, 1.41]}  % experiments/FIGS/data/paired_v1err/finetune/ratios.json :: ratios[70].ratio, ratios[70].ci95 :: sha256 3e5cf25e859b68c2 :: random over 43, leg2 N1000 1-macro_auc: paired geometric mean [95% interval]
\newcommand{\PairedFtJcEThreeRandomOverOneseven}{0.98~[0.84, 1.13]}  % experiments/FIGS/data/paired_v1err/finetune/ratios.json :: ratios[71].ratio, ratios[71].ci95 :: sha256 3e5cf25e859b68c2 :: random over 17, leg2 N1000 1-macro_auc: paired geometric mean [95% interval]
\newcommand{\PairedFtJcEThreeRandomDrawTwoOverRandomDrawOne}{0.954~[0.949, 0.958]}  % experiments/FIGS/data/paired_v1err/finetune/ratios.json :: ratios[72].ratio, ratios[72].ci95 :: sha256 3e5cf25e859b68c2 :: random draw 2 over random draw 1, leg2 N1000 1-macro_auc: paired geometric mean [95% interval]
\newcommand{\PairedFtJcEThreeRandomDrawThreeOverRandomDrawOne}{1.013~[1.009, 1.017]}  % experiments/FIGS/data/paired_v1err/finetune/ratios.json :: ratios[73].ratio, ratios[73].ci95 :: sha256 3e5cf25e859b68c2 :: random draw 3 over random draw 1, leg2 N1000 1-macro_auc: paired geometric mean [95% interval]
\newcommand{\PairedFtJcEThreeRandomDrawThreeOverRandomDrawTwo}{1.062~[1.057, 1.067]}  % experiments/FIGS/data/paired_v1err/finetune/ratios.json :: ratios[74].ratio, ratios[74].ci95 :: sha256 3e5cf25e859b68c2 :: random draw 3 over random draw 2, leg2 N1000 1-macro_auc: paired geometric mean [95% interval]
\newcommand{\PairedFtJcEFourOnesixtwoOverOneeighteight}{0.98~[0.93, 1.02]}  % experiments/FIGS/data/paired_v1err/finetune/ratios.json :: ratios[75].ratio, ratios[75].ci95 :: sha256 3e5cf25e859b68c2 :: 162 over 188, leg2 N10000 1-macro_auc: paired geometric mean [95% interval]
\newcommand{\PairedFtJcEFourFourthreeOverOneeighteight}{0.98~[0.93, 1.04]}  % experiments/FIGS/data/paired_v1err/finetune/ratios.json :: ratios[76].ratio, ratios[76].ci95 :: sha256 3e5cf25e859b68c2 :: 43 over 188, leg2 N10000 1-macro_auc: paired geometric mean [95% interval]
\newcommand{\PairedFtJcEFourOnesevenOverOneeighteight}{1.07~[1.02, 1.11]}  % experiments/FIGS/data/paired_v1err/finetune/ratios.json :: ratios[77].ratio, ratios[77].ci95 :: sha256 3e5cf25e859b68c2 :: 17 over 188, leg2 N10000 1-macro_auc: paired geometric mean [95% interval]
\newcommand{\PairedFtJcEFourFourthreeOverOnesixtwo}{1.01~[0.97, 1.04]}  % experiments/FIGS/data/paired_v1err/finetune/ratios.json :: ratios[78].ratio, ratios[78].ci95 :: sha256 3e5cf25e859b68c2 :: 43 over 162, leg2 N10000 1-macro_auc: paired geometric mean [95% interval]
\newcommand{\PairedFtJcEFourOnesevenOverOnesixtwo}{1.09~[1.06, 1.12]}  % experiments/FIGS/data/paired_v1err/finetune/ratios.json :: ratios[79].ratio, ratios[79].ci95 :: sha256 3e5cf25e859b68c2 :: 17 over 162, leg2 N10000 1-macro_auc: paired geometric mean [95% interval]
\newcommand{\PairedFtJcEFourOnesevenOverFourthree}{1.08~[1.07, 1.10]}  % experiments/FIGS/data/paired_v1err/finetune/ratios.json :: ratios[80].ratio, ratios[80].ci95 :: sha256 3e5cf25e859b68c2 :: 17 over 43, leg2 N10000 1-macro_auc: paired geometric mean [95% interval]
\newcommand{\PairedFtJcEFourOnesixtwoMassOverOnesixtwo}{1.01~[0.98, 1.04]}  % experiments/FIGS/data/paired_v1err/finetune/ratios.json :: ratios[81].ratio, ratios[81].ci95 :: sha256 3e5cf25e859b68c2 :: 162+mass over 162, leg2 N10000 1-macro_auc: paired geometric mean [95% interval]
\newcommand{\PairedFtJcEFourOnesevenMassOverOneseven}{1.05~[1.02, 1.08]}  % experiments/FIGS/data/paired_v1err/finetune/ratios.json :: ratios[82].ratio, ratios[82].ci95 :: sha256 3e5cf25e859b68c2 :: 17+mass over 17, leg2 N10000 1-macro_auc: paired geometric mean [95% interval]
\newcommand{\PairedFtJcEFourRandomOverOneeighteight}{1.08~[1.03, 1.14]}  % experiments/FIGS/data/paired_v1err/finetune/ratios.json :: ratios[83].ratio, ratios[83].ci95 :: sha256 3e5cf25e859b68c2 :: random over 188, leg2 N10000 1-macro_auc: paired geometric mean [95% interval]
\newcommand{\PairedFtJcEFourRandomOverOnesixtwo}{1.11~[0.99, 1.23]}  % experiments/FIGS/data/paired_v1err/finetune/ratios.json :: ratios[84].ratio, ratios[84].ci95 :: sha256 3e5cf25e859b68c2 :: random over 162, leg2 N10000 1-macro_auc: paired geometric mean [95% interval]
\newcommand{\PairedFtJcEFourRandomOverFourthree}{1.10~[0.99, 1.21]}  % experiments/FIGS/data/paired_v1err/finetune/ratios.json :: ratios[85].ratio, ratios[85].ci95 :: sha256 3e5cf25e859b68c2 :: random over 43, leg2 N10000 1-macro_auc: paired geometric mean [95% interval]
\newcommand{\PairedFtJcEFourRandomOverOneseven}{1.02~[0.94, 1.10]}  % experiments/FIGS/data/paired_v1err/finetune/ratios.json :: ratios[86].ratio, ratios[86].ci95 :: sha256 3e5cf25e859b68c2 :: random over 17, leg2 N10000 1-macro_auc: paired geometric mean [95% interval]
\newcommand{\PairedFtJcEFourRandomDrawTwoOverRandomDrawOne}{0.960~[0.955, 0.965]}  % experiments/FIGS/data/paired_v1err/finetune/ratios.json :: ratios[87].ratio, ratios[87].ci95 :: sha256 3e5cf25e859b68c2 :: random draw 2 over random draw 1, leg2 N10000 1-macro_auc: paired geometric mean [95% interval]
\newcommand{\PairedFtJcEFourRandomDrawThreeOverRandomDrawOne}{1.008~[1.004, 1.012]}  % experiments/FIGS/data/paired_v1err/finetune/ratios.json :: ratios[88].ratio, ratios[88].ci95 :: sha256 3e5cf25e859b68c2 :: random draw 3 over random draw 1, leg2 N10000 1-macro_auc: paired geometric mean [95% interval]
\newcommand{\PairedFtJcEFourRandomDrawThreeOverRandomDrawTwo}{1.050~[1.045, 1.055]}  % experiments/FIGS/data/paired_v1err/finetune/ratios.json :: ratios[89].ratio, ratios[89].ci95 :: sha256 3e5cf25e859b68c2 :: random draw 3 over random draw 2, leg2 N10000 1-macro_auc: paired geometric mean [95% interval]
\newcommand{\PairedFtJcEFiveOnesixtwoOverOneeighteight}{0.983~[0.968, 0.998]}  % experiments/FIGS/data/paired_v1err/finetune/ratios.json :: ratios[90].ratio, ratios[90].ci95 :: sha256 3e5cf25e859b68c2 :: 162 over 188, leg2 N100000 1-macro_auc: paired geometric mean [95% interval]
\newcommand{\PairedFtJcEFiveFourthreeOverOneeighteight}{1.015~[1.002, 1.029]}  % experiments/FIGS/data/paired_v1err/finetune/ratios.json :: ratios[91].ratio, ratios[91].ci95 :: sha256 3e5cf25e859b68c2 :: 43 over 188, leg2 N100000 1-macro_auc: paired geometric mean [95% interval]
\newcommand{\PairedFtJcEFiveOnesevenOverOneeighteight}{1.07~[1.06, 1.09]}  % experiments/FIGS/data/paired_v1err/finetune/ratios.json :: ratios[92].ratio, ratios[92].ci95 :: sha256 3e5cf25e859b68c2 :: 17 over 188, leg2 N100000 1-macro_auc: paired geometric mean [95% interval]
\newcommand{\PairedFtJcEFiveFourthreeOverOnesixtwo}{1.03~[1.02, 1.05]}  % experiments/FIGS/data/paired_v1err/finetune/ratios.json :: ratios[93].ratio, ratios[93].ci95 :: sha256 3e5cf25e859b68c2 :: 43 over 162, leg2 N100000 1-macro_auc: paired geometric mean [95% interval]
\newcommand{\PairedFtJcEFiveOnesevenOverOnesixtwo}{1.091~[1.088, 1.094]}  % experiments/FIGS/data/paired_v1err/finetune/ratios.json :: ratios[94].ratio, ratios[94].ci95 :: sha256 3e5cf25e859b68c2 :: 17 over 162, leg2 N100000 1-macro_auc: paired geometric mean [95% interval]
\newcommand{\PairedFtJcEFiveOnesevenOverFourthree}{1.06~[1.04, 1.07]}  % experiments/FIGS/data/paired_v1err/finetune/ratios.json :: ratios[95].ratio, ratios[95].ci95 :: sha256 3e5cf25e859b68c2 :: 17 over 43, leg2 N100000 1-macro_auc: paired geometric mean [95% interval]
\newcommand{\PairedFtJcEFiveOnesixtwoMassOverOnesixtwo}{1.016~[1.008, 1.024]}  % experiments/FIGS/data/paired_v1err/finetune/ratios.json :: ratios[96].ratio, ratios[96].ci95 :: sha256 3e5cf25e859b68c2 :: 162+mass over 162, leg2 N100000 1-macro_auc: paired geometric mean [95% interval]
\newcommand{\PairedFtJcEFiveOnesevenMassOverOneseven}{1.02~[1.01, 1.03]}  % experiments/FIGS/data/paired_v1err/finetune/ratios.json :: ratios[97].ratio, ratios[97].ci95 :: sha256 3e5cf25e859b68c2 :: 17+mass over 17, leg2 N100000 1-macro_auc: paired geometric mean [95% interval]
\newcommand{\PairedFtJcEFiveRandomOverOneeighteight}{1.044~[1.039, 1.049]}  % experiments/FIGS/data/paired_v1err/finetune/ratios.json :: ratios[98].ratio, ratios[98].ci95 :: sha256 3e5cf25e859b68c2 :: random over 188, leg2 N100000 1-macro_auc: paired geometric mean [95% interval]
\newcommand{\PairedFtJcEFiveRandomOverOnesixtwo}{1.06~[1.04, 1.09]}  % experiments/FIGS/data/paired_v1err/finetune/ratios.json :: ratios[99].ratio, ratios[99].ci95 :: sha256 3e5cf25e859b68c2 :: random over 162, leg2 N100000 1-macro_auc: paired geometric mean [95% interval]
\newcommand{\PairedFtJcEFiveRandomOverFourthree}{1.03~[1.01, 1.04]}  % experiments/FIGS/data/paired_v1err/finetune/ratios.json :: ratios[100].ratio, ratios[100].ci95 :: sha256 3e5cf25e859b68c2 :: random over 43, leg2 N100000 1-macro_auc: paired geometric mean [95% interval]
\newcommand{\PairedFtJcEFiveRandomOverOneseven}{0.97~[0.95, 0.99]}  % experiments/FIGS/data/paired_v1err/finetune/ratios.json :: ratios[101].ratio, ratios[101].ci95 :: sha256 3e5cf25e859b68c2 :: random over 17, leg2 N100000 1-macro_auc: paired geometric mean [95% interval]
\newcommand{\PairedFtJcEFiveRandomDrawTwoOverRandomDrawOne}{0.992~[0.988, 0.997]}  % experiments/FIGS/data/paired_v1err/finetune/ratios.json :: ratios[102].ratio, ratios[102].ci95 :: sha256 3e5cf25e859b68c2 :: random draw 2 over random draw 1, leg2 N100000 1-macro_auc: paired geometric mean [95% interval]
\newcommand{\PairedFtJcEFiveRandomDrawThreeOverRandomDrawOne}{1.007~[1.003, 1.011]}  % experiments/FIGS/data/paired_v1err/finetune/ratios.json :: ratios[103].ratio, ratios[103].ci95 :: sha256 3e5cf25e859b68c2 :: random draw 3 over random draw 1, leg2 N100000 1-macro_auc: paired geometric mean [95% interval]
\newcommand{\PairedFtJcEFiveRandomDrawThreeOverRandomDrawTwo}{1.015~[1.011, 1.020]}  % experiments/FIGS/data/paired_v1err/finetune/ratios.json :: ratios[104].ratio, ratios[104].ci95 :: sha256 3e5cf25e859b68c2 :: random draw 3 over random draw 2, leg2 N100000 1-macro_auc: paired geometric mean [95% interval]
\newcommand{\PairedFtJcESixOnesixtwoOverOneeighteight}{1.00~[0.99, 1.01]}  % experiments/FIGS/data/paired_v1err/finetune/ratios.json :: ratios[105].ratio, ratios[105].ci95 :: sha256 3e5cf25e859b68c2 :: 162 over 188, leg2 N1000000 1-macro_auc: paired geometric mean [95% interval]
\newcommand{\PairedFtJcESixFourthreeOverOneeighteight}{1.02~[1.01, 1.03]}  % experiments/FIGS/data/paired_v1err/finetune/ratios.json :: ratios[106].ratio, ratios[106].ci95 :: sha256 3e5cf25e859b68c2 :: 43 over 188, leg2 N1000000 1-macro_auc: paired geometric mean [95% interval]
\newcommand{\PairedFtJcESixOnesevenOverOneeighteight}{1.08~[1.07, 1.09]}  % experiments/FIGS/data/paired_v1err/finetune/ratios.json :: ratios[107].ratio, ratios[107].ci95 :: sha256 3e5cf25e859b68c2 :: 17 over 188, leg2 N1000000 1-macro_auc: paired geometric mean [95% interval]
\newcommand{\PairedFtJcESixFourthreeOverOnesixtwo}{1.02~[1.01, 1.03]}  % experiments/FIGS/data/paired_v1err/finetune/ratios.json :: ratios[108].ratio, ratios[108].ci95 :: sha256 3e5cf25e859b68c2 :: 43 over 162, leg2 N1000000 1-macro_auc: paired geometric mean [95% interval]
\newcommand{\PairedFtJcESixOnesevenOverOnesixtwo}{1.08~[1.07, 1.09]}  % experiments/FIGS/data/paired_v1err/finetune/ratios.json :: ratios[109].ratio, ratios[109].ci95 :: sha256 3e5cf25e859b68c2 :: 17 over 162, leg2 N1000000 1-macro_auc: paired geometric mean [95% interval]
\newcommand{\PairedFtJcESixOnesevenOverFourthree}{1.06~[1.05, 1.07]}  % experiments/FIGS/data/paired_v1err/finetune/ratios.json :: ratios[110].ratio, ratios[110].ci95 :: sha256 3e5cf25e859b68c2 :: 17 over 43, leg2 N1000000 1-macro_auc: paired geometric mean [95% interval]
\newcommand{\PairedFtJcESixOnesixtwoMassOverOnesixtwo}{1.012~[1.004, 1.020]}  % experiments/FIGS/data/paired_v1err/finetune/ratios.json :: ratios[111].ratio, ratios[111].ci95 :: sha256 3e5cf25e859b68c2 :: 162+mass over 162, leg2 N1000000 1-macro_auc: paired geometric mean [95% interval]
\newcommand{\PairedFtJcESixOnesevenMassOverOneseven}{1.02~[1.01, 1.04]}  % experiments/FIGS/data/paired_v1err/finetune/ratios.json :: ratios[112].ratio, ratios[112].ci95 :: sha256 3e5cf25e859b68c2 :: 17+mass over 17, leg2 N1000000 1-macro_auc: paired geometric mean [95% interval]
\newcommand{\PairedFtJcESixRandomOverOneeighteight}{1.046~[1.038, 1.053]}  % experiments/FIGS/data/paired_v1err/finetune/ratios.json :: ratios[113].ratio, ratios[113].ci95 :: sha256 3e5cf25e859b68c2 :: random over 188, leg2 N1000000 1-macro_auc: paired geometric mean [95% interval]
\newcommand{\PairedFtJcESixRandomOverOnesixtwo}{1.047~[1.038, 1.055]}  % experiments/FIGS/data/paired_v1err/finetune/ratios.json :: ratios[114].ratio, ratios[114].ci95 :: sha256 3e5cf25e859b68c2 :: random over 162, leg2 N1000000 1-macro_auc: paired geometric mean [95% interval]
\newcommand{\PairedFtJcESixRandomOverFourthree}{1.03~[1.02, 1.04]}  % experiments/FIGS/data/paired_v1err/finetune/ratios.json :: ratios[115].ratio, ratios[115].ci95 :: sha256 3e5cf25e859b68c2 :: random over 43, leg2 N1000000 1-macro_auc: paired geometric mean [95% interval]
\newcommand{\PairedFtJcESixRandomOverOneseven}{0.97~[0.96, 0.98]}  % experiments/FIGS/data/paired_v1err/finetune/ratios.json :: ratios[116].ratio, ratios[116].ci95 :: sha256 3e5cf25e859b68c2 :: random over 17, leg2 N1000000 1-macro_auc: paired geometric mean [95% interval]
\newcommand{\PairedFtJcESixRandomDrawTwoOverRandomDrawOne}{1.001~[0.998, 1.004]}  % experiments/FIGS/data/paired_v1err/finetune/ratios.json :: ratios[117].ratio, ratios[117].ci95 :: sha256 3e5cf25e859b68c2 :: random draw 2 over random draw 1, leg2 N1000000 1-macro_auc: paired geometric mean [95% interval]
\newcommand{\PairedFtJcESixRandomDrawThreeOverRandomDrawOne}{1.005~[1.002, 1.008]}  % experiments/FIGS/data/paired_v1err/finetune/ratios.json :: ratios[118].ratio, ratios[118].ci95 :: sha256 3e5cf25e859b68c2 :: random draw 3 over random draw 1, leg2 N1000000 1-macro_auc: paired geometric mean [95% interval]
\newcommand{\PairedFtJcESixRandomDrawThreeOverRandomDrawTwo}{1.004~[1.001, 1.007]}  % experiments/FIGS/data/paired_v1err/finetune/ratios.json :: ratios[119].ratio, ratios[119].ci95 :: sha256 3e5cf25e859b68c2 :: random draw 3 over random draw 2, leg2 N1000000 1-macro_auc: paired geometric mean [95% interval]
\newcommand{\PairedMassResMlpOnesixtwoOverOneeighteight}{0.990~[0.977, 1.003]}  % experiments/FIGS/data/paired_v1err/mass/ratios.json :: ratios[0].ratio, ratios[0].ci95 :: sha256 d877a7f087260f64 :: 162 over 188, resolution mlp sigma_eff: paired geometric mean [95% interval]
\newcommand{\PairedMassResMlpFourthreeOverOneeighteight}{0.98~[0.96, 0.99]}  % experiments/FIGS/data/paired_v1err/mass/ratios.json :: ratios[1].ratio, ratios[1].ci95 :: sha256 d877a7f087260f64 :: 43 over 188, resolution mlp sigma_eff: paired geometric mean [95% interval]
\newcommand{\PairedMassResMlpOnesevenOverOneeighteight}{0.96~[0.94, 0.98]}  % experiments/FIGS/data/paired_v1err/mass/ratios.json :: ratios[2].ratio, ratios[2].ci95 :: sha256 d877a7f087260f64 :: 17 over 188, resolution mlp sigma_eff: paired geometric mean [95% interval]
\newcommand{\PairedMassResMlpFourthreeOverOnesixtwo}{0.989~[0.978, 0.999]}  % experiments/FIGS/data/paired_v1err/mass/ratios.json :: ratios[3].ratio, ratios[3].ci95 :: sha256 d877a7f087260f64 :: 43 over 162, resolution mlp sigma_eff: paired geometric mean [95% interval]
\newcommand{\PairedMassResMlpOnesevenOverOnesixtwo}{0.97~[0.94, 0.99]}  % experiments/FIGS/data/paired_v1err/mass/ratios.json :: ratios[4].ratio, ratios[4].ci95 :: sha256 d877a7f087260f64 :: 17 over 162, resolution mlp sigma_eff: paired geometric mean [95% interval]
\newcommand{\PairedMassResMlpOnesevenOverFourthree}{0.979~[0.961, 0.997]}  % experiments/FIGS/data/paired_v1err/mass/ratios.json :: ratios[5].ratio, ratios[5].ci95 :: sha256 d877a7f087260f64 :: 17 over 43, resolution mlp sigma_eff: paired geometric mean [95% interval]
\newcommand{\PairedMassResMlpOnesixtwoMassOverOnesixtwo}{0.80~[0.79, 0.81]}  % experiments/FIGS/data/paired_v1err/mass/ratios.json :: ratios[6].ratio, ratios[6].ci95 :: sha256 d877a7f087260f64 :: 162+mass over 162, resolution mlp sigma_eff: paired geometric mean [95% interval]
\newcommand{\PairedMassResMlpOnesevenMassOverOneseven}{0.83~[0.82, 0.84]}  % experiments/FIGS/data/paired_v1err/mass/ratios.json :: ratios[7].ratio, ratios[7].ci95 :: sha256 d877a7f087260f64 :: 17+mass over 17, resolution mlp sigma_eff: paired geometric mean [95% interval]
\newcommand{\PairedMassResRidgeOnesixtwoOverOneeighteight}{0.97~[0.96, 0.98]}  % experiments/FIGS/data/paired_v1err/mass/ratios.json :: ratios[8].ratio, ratios[8].ci95 :: sha256 d877a7f087260f64 :: 162 over 188, resolution ridge sigma_eff: paired geometric mean [95% interval]
\newcommand{\PairedMassResRidgeFourthreeOverOneeighteight}{0.95~[0.94, 0.96]}  % experiments/FIGS/data/paired_v1err/mass/ratios.json :: ratios[9].ratio, ratios[9].ci95 :: sha256 d877a7f087260f64 :: 43 over 188, resolution ridge sigma_eff: paired geometric mean [95% interval]
\newcommand{\PairedMassResRidgeOnesevenOverOneeighteight}{0.88~[0.85, 0.92]}  % experiments/FIGS/data/paired_v1err/mass/ratios.json :: ratios[10].ratio, ratios[10].ci95 :: sha256 d877a7f087260f64 :: 17 over 188, resolution ridge sigma_eff: paired geometric mean [95% interval]
\newcommand{\PairedMassResRidgeFourthreeOverOnesixtwo}{0.98~[0.97, 0.99]}  % experiments/FIGS/data/paired_v1err/mass/ratios.json :: ratios[11].ratio, ratios[11].ci95 :: sha256 d877a7f087260f64 :: 43 over 162, resolution ridge sigma_eff: paired geometric mean [95% interval]
\newcommand{\PairedMassResRidgeOnesevenOverOnesixtwo}{0.91~[0.87, 0.95]}  % experiments/FIGS/data/paired_v1err/mass/ratios.json :: ratios[12].ratio, ratios[12].ci95 :: sha256 d877a7f087260f64 :: 17 over 162, resolution ridge sigma_eff: paired geometric mean [95% interval]
\newcommand{\PairedMassResRidgeOnesevenOverFourthree}{0.93~[0.89, 0.97]}  % experiments/FIGS/data/paired_v1err/mass/ratios.json :: ratios[13].ratio, ratios[13].ci95 :: sha256 d877a7f087260f64 :: 17 over 43, resolution ridge sigma_eff: paired geometric mean [95% interval]
\newcommand{\PairedMassResRidgeOnesixtwoMassOverOnesixtwo}{0.678~[0.672, 0.684]}  % experiments/FIGS/data/paired_v1err/mass/ratios.json :: ratios[14].ratio, ratios[14].ci95 :: sha256 d877a7f087260f64 :: 162+mass over 162, resolution ridge sigma_eff: paired geometric mean [95% interval]
\newcommand{\PairedMassResRidgeOnesevenMassOverOneseven}{0.72~[0.69, 0.76]}  % experiments/FIGS/data/paired_v1err/mass/ratios.json :: ratios[15].ratio, ratios[15].ci95 :: sha256 d877a7f087260f64 :: 17+mass over 17, resolution ridge sigma_eff: paired geometric mean [95% interval]
\newcommand{\VcbNSignal}{1{,}659}  % experiments/FIGS/data/probe_ladder_vcbwindow_mlp2/sall.json :: tasks.bc_vs_rest.n_signal_test :: sha256 5b8b5c7c28602f93 :: X->bc test jets in the window
\newcommand{\VcbNBackground}{19{,}244}  % experiments/FIGS/data/probe_ladder_vcbwindow_mlp2/sall.json :: tasks.bc_vs_rest.n_background_test :: sha256 5b8b5c7c28602f93 :: bq, cs, bqq and QCD test jets in the window
\newcommand{\VcbEpsSixty}{60}  % experiments/FIGS/data/probe_ladder_vcbwindow_mlp2/sall.json :: tasks.bc_vs_rest.eps_s :: sha256 5b8b5c7c28602f93 :: signal efficiency, percent
\newcommand{\VcbOmaLinearOnesixtwo}{\ensuremath{(3.17 \pm 0.05)\times10^{-2}}}  % experiments/FIGS/data/probe_ladder_vcbwindow_mlp2/sall.json :: tasks.bc_vs_rest.arms.l162-s*.linear.log1m_auc :: sha256 5b8b5c7c28602f93 :: 1-AUC = exp of it, mean +- SD over 5 seeds
\newcommand{\VcbLogOmaLinearOnesixtwo}{\ensuremath{-}3.451}  % experiments/FIGS/data/probe_ladder_vcbwindow_mlp2/sall.json :: tasks.bc_vs_rest.arms.l162-s*.linear.log1m_auc :: sha256 5b8b5c7c28602f93 :: mean over seeds
\newcommand{\VcbRejLinearSixtyOnesixtwo}{207\,\ensuremath{\pm}\,19}  % experiments/FIGS/data/probe_ladder_vcbwindow_mlp2/sall.json :: tasks.bc_vs_rest.arms.l162-s*.linear.rejection_at['0.60'].rejection :: sha256 5b8b5c7c28602f93 :: background rejection at this signal efficiency, mean +- SD over seeds
\newcommand{\VcbBkgLeftLinearSixtyOnesixtwo}{93.4}  % experiments/FIGS/data/probe_ladder_vcbwindow_mlp2/sall.json :: tasks.bc_vs_rest.arms.l162-s*.linear.rejection_at['0.60'].n_bkg_pass :: sha256 5b8b5c7c28602f93 :: background jets passing the cut, mean over seeds
\newcommand{\VcbOmaLinearOneseven}{\ensuremath{(5.21 \pm 0.23)\times10^{-2}}}  % experiments/FIGS/data/probe_ladder_vcbwindow_mlp2/sall.json :: tasks.bc_vs_rest.arms.r16q1-s*.linear.log1m_auc :: sha256 5b8b5c7c28602f93 :: 1-AUC = exp of it, mean +- SD over 5 seeds
\newcommand{\VcbLogOmaLinearOneseven}{\ensuremath{-}2.955}  % experiments/FIGS/data/probe_ladder_vcbwindow_mlp2/sall.json :: tasks.bc_vs_rest.arms.r16q1-s*.linear.log1m_auc :: sha256 5b8b5c7c28602f93 :: mean over seeds
\newcommand{\VcbRejLinearSixtyOneseven}{67\,\ensuremath{\pm}\,14}  % experiments/FIGS/data/probe_ladder_vcbwindow_mlp2/sall.json :: tasks.bc_vs_rest.arms.r16q1-s*.linear.rejection_at['0.60'].rejection :: sha256 5b8b5c7c28602f93 :: background rejection at this signal efficiency, mean +- SD over seeds
\newcommand{\VcbBkgLeftLinearSixtyOneseven}{295.0}  % experiments/FIGS/data/probe_ladder_vcbwindow_mlp2/sall.json :: tasks.bc_vs_rest.arms.r16q1-s*.linear.rejection_at['0.60'].n_bkg_pass :: sha256 5b8b5c7c28602f93 :: background jets passing the cut, mean over seeds
\newcommand{\VcbOmaRatioLinear}{1.64}  % experiments/FIGS/data/probe_ladder_vcbwindow_mlp2/sall.json :: tasks.bc_vs_rest.arms.*.linear.log1m_auc :: sha256 5b8b5c7c28602f93 :: seed mean of 1-AUC, 17 classes over 162
\newcommand{\VcbFactorLinear}{1.64}  % experiments/FIGS/data/probe_ladder_vcbwindow_mlp2/sall.json :: tasks.bc_vs_rest.arms.*.linear.log1m_auc :: sha256 5b8b5c7c28602f93 :: exp of the difference of the seed means of log(1-AUC), 17 classes over 162
\newcommand{\VcbOmaMlpOnesixtwo}{\ensuremath{(3.102 \pm 0.023)\times10^{-2}}}  % experiments/FIGS/data/probe_ladder_vcbwindow_mlp2/sall.json :: tasks.bc_vs_rest.arms.l162-s*.mlp.log1m_auc :: sha256 5b8b5c7c28602f93 :: 1-AUC = exp of it, mean +- SD over 5 seeds
\newcommand{\VcbLogOmaMlpOnesixtwo}{\ensuremath{-}3.473}  % experiments/FIGS/data/probe_ladder_vcbwindow_mlp2/sall.json :: tasks.bc_vs_rest.arms.l162-s*.mlp.log1m_auc :: sha256 5b8b5c7c28602f93 :: mean over seeds
\newcommand{\VcbRejMlpSixtyOnesixtwo}{230\,\ensuremath{\pm}\,31}  % experiments/FIGS/data/probe_ladder_vcbwindow_mlp2/sall.json :: tasks.bc_vs_rest.arms.l162-s*.mlp.rejection_at['0.60'].rejection :: sha256 5b8b5c7c28602f93 :: background rejection at this signal efficiency, mean +- SD over seeds
\newcommand{\VcbBkgLeftMlpSixtyOnesixtwo}{84.8}  % experiments/FIGS/data/probe_ladder_vcbwindow_mlp2/sall.json :: tasks.bc_vs_rest.arms.l162-s*.mlp.rejection_at['0.60'].n_bkg_pass :: sha256 5b8b5c7c28602f93 :: background jets passing the cut, mean over seeds
\newcommand{\VcbOmaMlpOneseven}{\ensuremath{(4.66 \pm 0.19)\times10^{-2}}}  % experiments/FIGS/data/probe_ladder_vcbwindow_mlp2/sall.json :: tasks.bc_vs_rest.arms.r16q1-s*.mlp.log1m_auc :: sha256 5b8b5c7c28602f93 :: 1-AUC = exp of it, mean +- SD over 5 seeds
\newcommand{\VcbLogOmaMlpOneseven}{\ensuremath{-}3.066}  % experiments/FIGS/data/probe_ladder_vcbwindow_mlp2/sall.json :: tasks.bc_vs_rest.arms.r16q1-s*.mlp.log1m_auc :: sha256 5b8b5c7c28602f93 :: mean over seeds
\newcommand{\VcbRejMlpSixtyOneseven}{78\,\ensuremath{\pm}\,13}  % experiments/FIGS/data/probe_ladder_vcbwindow_mlp2/sall.json :: tasks.bc_vs_rest.arms.r16q1-s*.mlp.rejection_at['0.60'].rejection :: sha256 5b8b5c7c28602f93 :: background rejection at this signal efficiency, mean +- SD over seeds
\newcommand{\VcbBkgLeftMlpSixtyOneseven}{250.8}  % experiments/FIGS/data/probe_ladder_vcbwindow_mlp2/sall.json :: tasks.bc_vs_rest.arms.r16q1-s*.mlp.rejection_at['0.60'].n_bkg_pass :: sha256 5b8b5c7c28602f93 :: background jets passing the cut, mean over seeds
\newcommand{\VcbOmaRatioMlp}{1.50}  % experiments/FIGS/data/probe_ladder_vcbwindow_mlp2/sall.json :: tasks.bc_vs_rest.arms.*.mlp.log1m_auc :: sha256 5b8b5c7c28602f93 :: seed mean of 1-AUC, 17 classes over 162
\newcommand{\VcbFactorMlp}{1.50}  % experiments/FIGS/data/probe_ladder_vcbwindow_mlp2/sall.json :: tasks.bc_vs_rest.arms.*.mlp.log1m_auc :: sha256 5b8b5c7c28602f93 :: exp of the difference of the seed means of log(1-AUC), 17 classes over 162
\newcommand{\AccLegOneOnesixtwoSonebEFour}{0.32273}  % experiments/FIGS/data/leg1_metrics.json :: summary.l162-s1b.N10000.accuracy_mean :: sha256 3f564e960b2c86a7 :: wave 1, superseded; 3 fine-tuning seeds
\newcommand{\AccSdLegOneOnesixtwoSonebEFour}{0.01894}  % experiments/FIGS/data/leg1_metrics.json :: summary.l162-s1b.N10000.accuracy_sd :: sha256 3f564e960b2c86a7 :: wave 1, superseded
\newcommand{\AccLegOneOnesixtwoSonebEFive}{0.53506}  % experiments/FIGS/data/leg1_metrics.json :: summary.l162-s1b.N100000.accuracy_mean :: sha256 3f564e960b2c86a7 :: wave 1, superseded; 3 fine-tuning seeds
\newcommand{\AccSdLegOneOnesixtwoSonebEFive}{0.00044}  % experiments/FIGS/data/leg1_metrics.json :: summary.l162-s1b.N100000.accuracy_sd :: sha256 3f564e960b2c86a7 :: wave 1, superseded
\newcommand{\AccLegOneOnesixtwoSonebESix}{0.54337}  % experiments/FIGS/data/leg1_metrics.json :: summary.l162-s1b.N1000000.accuracy_mean :: sha256 3f564e960b2c86a7 :: wave 1, superseded; 3 fine-tuning seeds
\newcommand{\AccSdLegOneOnesixtwoSonebESix}{0.00074}  % experiments/FIGS/data/leg1_metrics.json :: summary.l162-s1b.N1000000.accuracy_sd :: sha256 3f564e960b2c86a7 :: wave 1, superseded
\newcommand{\AccLegOneOnesevenStwoEFour}{0.22684}  % experiments/FIGS/data/leg1_metrics.json :: summary.r16q1-s2.N10000.accuracy_mean :: sha256 3f564e960b2c86a7 :: wave 1, superseded; 3 fine-tuning seeds
\newcommand{\AccSdLegOneOnesevenStwoEFour}{0.00805}  % experiments/FIGS/data/leg1_metrics.json :: summary.r16q1-s2.N10000.accuracy_sd :: sha256 3f564e960b2c86a7 :: wave 1, superseded
\newcommand{\AccLegOneOnesevenStwoEFive}{0.49650}  % experiments/FIGS/data/leg1_metrics.json :: summary.r16q1-s2.N100000.accuracy_mean :: sha256 3f564e960b2c86a7 :: wave 1, superseded; 3 fine-tuning seeds
\newcommand{\AccSdLegOneOnesevenStwoEFive}{0.00110}  % experiments/FIGS/data/leg1_metrics.json :: summary.r16q1-s2.N100000.accuracy_sd :: sha256 3f564e960b2c86a7 :: wave 1, superseded
\newcommand{\AccLegOneOnesevenStwoESix}{0.51969}  % experiments/FIGS/data/leg1_metrics.json :: summary.r16q1-s2.N1000000.accuracy_mean :: sha256 3f564e960b2c86a7 :: wave 1, superseded; 3 fine-tuning seeds
\newcommand{\AccSdLegOneOnesevenStwoESix}{0.00072}  % experiments/FIGS/data/leg1_metrics.json :: summary.r16q1-s2.N1000000.accuracy_sd :: sha256 3f564e960b2c86a7 :: wave 1, superseded
\newcommand{\AccLegOneOnesevenSthreeEFour}{0.24203}  % experiments/FIGS/data/leg1_metrics.json :: summary.r16q1-s3.N10000.accuracy_mean :: sha256 3f564e960b2c86a7 :: wave 1, superseded; 3 fine-tuning seeds
\newcommand{\AccSdLegOneOnesevenSthreeEFour}{0.00213}  % experiments/FIGS/data/leg1_metrics.json :: summary.r16q1-s3.N10000.accuracy_sd :: sha256 3f564e960b2c86a7 :: wave 1, superseded
\newcommand{\AccLegOneOnesevenSthreeEFive}{0.49313}  % experiments/FIGS/data/leg1_metrics.json :: summary.r16q1-s3.N100000.accuracy_mean :: sha256 3f564e960b2c86a7 :: wave 1, superseded; 3 fine-tuning seeds
\newcommand{\AccSdLegOneOnesevenSthreeEFive}{0.00294}  % experiments/FIGS/data/leg1_metrics.json :: summary.r16q1-s3.N100000.accuracy_sd :: sha256 3f564e960b2c86a7 :: wave 1, superseded
\newcommand{\AccLegOneOnesevenSthreeESix}{0.51804}  % experiments/FIGS/data/leg1_metrics.json :: summary.r16q1-s3.N1000000.accuracy_mean :: sha256 3f564e960b2c86a7 :: wave 1, superseded; 3 fine-tuning seeds
\newcommand{\AccSdLegOneOnesevenSthreeESix}{0.00049}  % experiments/FIGS/data/leg1_metrics.json :: summary.r16q1-s3.N1000000.accuracy_sd :: sha256 3f564e960b2c86a7 :: wave 1, superseded
\newcommand{\AccLegOneOnesevenSfourEFour}{0.24487}  % experiments/FIGS/data/leg1_metrics.json :: summary.r16q1-s4.N10000.accuracy_mean :: sha256 3f564e960b2c86a7 :: wave 1, superseded; 3 fine-tuning seeds
\newcommand{\AccSdLegOneOnesevenSfourEFour}{0.00645}  % experiments/FIGS/data/leg1_metrics.json :: summary.r16q1-s4.N10000.accuracy_sd :: sha256 3f564e960b2c86a7 :: wave 1, superseded
\newcommand{\AccLegOneOnesevenSfourEFive}{0.49510}  % experiments/FIGS/data/leg1_metrics.json :: summary.r16q1-s4.N100000.accuracy_mean :: sha256 3f564e960b2c86a7 :: wave 1, superseded; 3 fine-tuning seeds
\newcommand{\AccSdLegOneOnesevenSfourEFive}{0.00133}  % experiments/FIGS/data/leg1_metrics.json :: summary.r16q1-s4.N100000.accuracy_sd :: sha256 3f564e960b2c86a7 :: wave 1, superseded
\newcommand{\AccLegOneOnesevenSfourESix}{0.51872}  % experiments/FIGS/data/leg1_metrics.json :: summary.r16q1-s4.N1000000.accuracy_mean :: sha256 3f564e960b2c86a7 :: wave 1, superseded; 3 fine-tuning seeds
\newcommand{\AccSdLegOneOnesevenSfourESix}{0.00049}  % experiments/FIGS/data/leg1_metrics.json :: summary.r16q1-s4.N1000000.accuracy_sd :: sha256 3f564e960b2c86a7 :: wave 1, superseded
\newcommand{\AccLegOneScratchEFour}{0.31369}  % experiments/FIGS/data/leg1_metrics.json :: summary.scratch.N10000.accuracy_mean :: sha256 3f564e960b2c86a7 :: wave 1, superseded; 3 fine-tuning seeds
\newcommand{\AccSdLegOneScratchEFour}{0.00712}  % experiments/FIGS/data/leg1_metrics.json :: summary.scratch.N10000.accuracy_sd :: sha256 3f564e960b2c86a7 :: wave 1, superseded
\newcommand{\AccLegOneScratchEFive}{0.42015}  % experiments/FIGS/data/leg1_metrics.json :: summary.scratch.N100000.accuracy_mean :: sha256 3f564e960b2c86a7 :: wave 1, superseded; 3 fine-tuning seeds
\newcommand{\AccSdLegOneScratchEFive}{0.00267}  % experiments/FIGS/data/leg1_metrics.json :: summary.scratch.N100000.accuracy_sd :: sha256 3f564e960b2c86a7 :: wave 1, superseded
\newcommand{\AccLegOneScratchESix}{0.48621}  % experiments/FIGS/data/leg1_metrics.json :: summary.scratch.N1000000.accuracy_mean :: sha256 3f564e960b2c86a7 :: wave 1, superseded; 3 fine-tuning seeds
\newcommand{\AccSdLegOneScratchESix}{0.00100}  % experiments/FIGS/data/leg1_metrics.json :: summary.scratch.N1000000.accuracy_sd :: sha256 3f564e960b2c86a7 :: wave 1, superseded
\newcommand{\AccLegOneSophonPublicEFour}{0.21613}  % experiments/FIGS/data/leg1_metrics.json :: summary.sophon-public.N10000.accuracy_mean :: sha256 3f564e960b2c86a7 :: wave 1, superseded; 3 fine-tuning seeds
\newcommand{\AccSdLegOneSophonPublicEFour}{0.00734}  % experiments/FIGS/data/leg1_metrics.json :: summary.sophon-public.N10000.accuracy_sd :: sha256 3f564e960b2c86a7 :: wave 1, superseded
\newcommand{\AccLegOneSophonPublicEFive}{0.53187}  % experiments/FIGS/data/leg1_metrics.json :: summary.sophon-public.N100000.accuracy_mean :: sha256 3f564e960b2c86a7 :: wave 1, superseded; 3 fine-tuning seeds
\newcommand{\AccSdLegOneSophonPublicEFive}{0.00231}  % experiments/FIGS/data/leg1_metrics.json :: summary.sophon-public.N100000.accuracy_sd :: sha256 3f564e960b2c86a7 :: wave 1, superseded
\newcommand{\AccLegOneSophonPublicESix}{0.54273}  % experiments/FIGS/data/leg1_metrics.json :: summary.sophon-public.N1000000.accuracy_mean :: sha256 3f564e960b2c86a7 :: wave 1, superseded; 3 fine-tuning seeds
\newcommand{\AccSdLegOneSophonPublicESix}{0.00064}  % experiments/FIGS/data/leg1_metrics.json :: summary.sophon-public.N1000000.accuracy_sd :: sha256 3f564e960b2c86a7 :: wave 1, superseded
\newcommand{\AccLegTwoOnesixtwoSonebEFour}{0.78606}  % experiments/FIGS/data/leg2_metrics.json :: summary.l162-s1b.N10000.accuracy_mean :: sha256 701c58c3240274bb :: wave 1, superseded; 3 fine-tuning seeds
\newcommand{\AccSdLegTwoOnesixtwoSonebEFour}{0.00111}  % experiments/FIGS/data/leg2_metrics.json :: summary.l162-s1b.N10000.accuracy_sd :: sha256 701c58c3240274bb :: wave 1, superseded
\newcommand{\AccLegTwoOnesixtwoSonebEFive}{0.83598}  % experiments/FIGS/data/leg2_metrics.json :: summary.l162-s1b.N100000.accuracy_mean :: sha256 701c58c3240274bb :: wave 1, superseded; 3 fine-tuning seeds
\newcommand{\AccSdLegTwoOnesixtwoSonebEFive}{0.00037}  % experiments/FIGS/data/leg2_metrics.json :: summary.l162-s1b.N100000.accuracy_sd :: sha256 701c58c3240274bb :: wave 1, superseded
\newcommand{\AccLegTwoOnesixtwoSonebESix}{0.84704}  % experiments/FIGS/data/leg2_metrics.json :: summary.l162-s1b.N1000000.accuracy_mean :: sha256 701c58c3240274bb :: wave 1, superseded; 3 fine-tuning seeds
\newcommand{\AccSdLegTwoOnesixtwoSonebESix}{0.00027}  % experiments/FIGS/data/leg2_metrics.json :: summary.l162-s1b.N1000000.accuracy_sd :: sha256 701c58c3240274bb :: wave 1, superseded
\newcommand{\AccLegTwoOnesevenStwoEFour}{0.72156}  % experiments/FIGS/data/leg2_metrics.json :: summary.r16q1-s2.N10000.accuracy_mean :: sha256 701c58c3240274bb :: wave 1, superseded; 3 fine-tuning seeds
\newcommand{\AccSdLegTwoOnesevenStwoEFour}{0.01049}  % experiments/FIGS/data/leg2_metrics.json :: summary.r16q1-s2.N10000.accuracy_sd :: sha256 701c58c3240274bb :: wave 1, superseded
\newcommand{\AccLegTwoOnesevenStwoEFive}{0.82614}  % experiments/FIGS/data/leg2_metrics.json :: summary.r16q1-s2.N100000.accuracy_mean :: sha256 701c58c3240274bb :: wave 1, superseded; 3 fine-tuning seeds
\newcommand{\AccSdLegTwoOnesevenStwoEFive}{0.00058}  % experiments/FIGS/data/leg2_metrics.json :: summary.r16q1-s2.N100000.accuracy_sd :: sha256 701c58c3240274bb :: wave 1, superseded
\newcommand{\AccLegTwoOnesevenStwoESix}{0.83943}  % experiments/FIGS/data/leg2_metrics.json :: summary.r16q1-s2.N1000000.accuracy_mean :: sha256 701c58c3240274bb :: wave 1, superseded; 3 fine-tuning seeds
\newcommand{\AccSdLegTwoOnesevenStwoESix}{0.00052}  % experiments/FIGS/data/leg2_metrics.json :: summary.r16q1-s2.N1000000.accuracy_sd :: sha256 701c58c3240274bb :: wave 1, superseded
\newcommand{\AccLegTwoOnesevenSthreeEFour}{0.72103}  % experiments/FIGS/data/leg2_metrics.json :: summary.r16q1-s3.N10000.accuracy_mean :: sha256 701c58c3240274bb :: wave 1, superseded; 3 fine-tuning seeds
\newcommand{\AccSdLegTwoOnesevenSthreeEFour}{0.01651}  % experiments/FIGS/data/leg2_metrics.json :: summary.r16q1-s3.N10000.accuracy_sd :: sha256 701c58c3240274bb :: wave 1, superseded
\newcommand{\AccLegTwoOnesevenSthreeEFive}{0.82629}  % experiments/FIGS/data/leg2_metrics.json :: summary.r16q1-s3.N100000.accuracy_mean :: sha256 701c58c3240274bb :: wave 1, superseded; 3 fine-tuning seeds
\newcommand{\AccSdLegTwoOnesevenSthreeEFive}{0.00035}  % experiments/FIGS/data/leg2_metrics.json :: summary.r16q1-s3.N100000.accuracy_sd :: sha256 701c58c3240274bb :: wave 1, superseded
\newcommand{\AccLegTwoOnesevenSthreeESix}{0.83944}  % experiments/FIGS/data/leg2_metrics.json :: summary.r16q1-s3.N1000000.accuracy_mean :: sha256 701c58c3240274bb :: wave 1, superseded; 3 fine-tuning seeds
\newcommand{\AccSdLegTwoOnesevenSthreeESix}{0.00019}  % experiments/FIGS/data/leg2_metrics.json :: summary.r16q1-s3.N1000000.accuracy_sd :: sha256 701c58c3240274bb :: wave 1, superseded
\newcommand{\AccLegTwoOnesevenSfourEFour}{0.72953}  % experiments/FIGS/data/leg2_metrics.json :: summary.r16q1-s4.N10000.accuracy_mean :: sha256 701c58c3240274bb :: wave 1, superseded; 3 fine-tuning seeds
\newcommand{\AccSdLegTwoOnesevenSfourEFour}{0.01069}  % experiments/FIGS/data/leg2_metrics.json :: summary.r16q1-s4.N10000.accuracy_sd :: sha256 701c58c3240274bb :: wave 1, superseded
\newcommand{\AccLegTwoOnesevenSfourEFive}{0.82564}  % experiments/FIGS/data/leg2_metrics.json :: summary.r16q1-s4.N100000.accuracy_mean :: sha256 701c58c3240274bb :: wave 1, superseded; 3 fine-tuning seeds
\newcommand{\AccSdLegTwoOnesevenSfourEFive}{0.00041}  % experiments/FIGS/data/leg2_metrics.json :: summary.r16q1-s4.N100000.accuracy_sd :: sha256 701c58c3240274bb :: wave 1, superseded
\newcommand{\AccLegTwoOnesevenSfourESix}{0.83909}  % experiments/FIGS/data/leg2_metrics.json :: summary.r16q1-s4.N1000000.accuracy_mean :: sha256 701c58c3240274bb :: wave 1, superseded; 3 fine-tuning seeds
\newcommand{\AccSdLegTwoOnesevenSfourESix}{0.00037}  % experiments/FIGS/data/leg2_metrics.json :: summary.r16q1-s4.N1000000.accuracy_sd :: sha256 701c58c3240274bb :: wave 1, superseded
\newcommand{\AccLegTwoRefEonearmsSoneRef}{0.81770}  % experiments/FIGS/data/leg2_metrics.json :: summary.ref_e1arms-s1.ref.accuracy_mean :: sha256 701c58c3240274bb :: wave 1, superseded; 1 fine-tuning seeds
\newcommand{\AccSdLegTwoRefEonearmsSoneRef}{---}  % experiments/FIGS/data/leg2_metrics.json :: summary.ref_e1arms-s1.ref.accuracy_sd :: sha256 701c58c3240274bb :: wave 1, superseded
\newcommand{\AccLegTwoRefEonearmsStwoRef}{0.81602}  % experiments/FIGS/data/leg2_metrics.json :: summary.ref_e1arms-s2.ref.accuracy_mean :: sha256 701c58c3240274bb :: wave 1, superseded; 1 fine-tuning seeds
\newcommand{\AccSdLegTwoRefEonearmsStwoRef}{---}  % experiments/FIGS/data/leg2_metrics.json :: summary.ref_e1arms-s2.ref.accuracy_sd :: sha256 701c58c3240274bb :: wave 1, superseded
\newcommand{\AccLegTwoRefEonearmsSthreeRef}{0.81625}  % experiments/FIGS/data/leg2_metrics.json :: summary.ref_e1arms-s3.ref.accuracy_mean :: sha256 701c58c3240274bb :: wave 1, superseded; 1 fine-tuning seeds
\newcommand{\AccSdLegTwoRefEonearmsSthreeRef}{---}  % experiments/FIGS/data/leg2_metrics.json :: summary.ref_e1arms-s3.ref.accuracy_sd :: sha256 701c58c3240274bb :: wave 1, superseded
\newcommand{\AccLegTwoScratchEFour}{0.68230}  % experiments/FIGS/data/leg2_metrics.json :: summary.scratch.N10000.accuracy_mean :: sha256 701c58c3240274bb :: wave 1, superseded; 3 fine-tuning seeds
\newcommand{\AccSdLegTwoScratchEFour}{0.00282}  % experiments/FIGS/data/leg2_metrics.json :: summary.scratch.N10000.accuracy_sd :: sha256 701c58c3240274bb :: wave 1, superseded
\newcommand{\AccLegTwoScratchEFive}{0.77769}  % experiments/FIGS/data/leg2_metrics.json :: summary.scratch.N100000.accuracy_mean :: sha256 701c58c3240274bb :: wave 1, superseded; 3 fine-tuning seeds
\newcommand{\AccSdLegTwoScratchEFive}{0.00208}  % experiments/FIGS/data/leg2_metrics.json :: summary.scratch.N100000.accuracy_sd :: sha256 701c58c3240274bb :: wave 1, superseded
\newcommand{\AccLegTwoScratchESix}{0.82303}  % experiments/FIGS/data/leg2_metrics.json :: summary.scratch.N1000000.accuracy_mean :: sha256 701c58c3240274bb :: wave 1, superseded; 3 fine-tuning seeds
\newcommand{\AccSdLegTwoScratchESix}{0.00134}  % experiments/FIGS/data/leg2_metrics.json :: summary.scratch.N1000000.accuracy_sd :: sha256 701c58c3240274bb :: wave 1, superseded
\newcommand{\AccLegTwoSophonPublicEFour}{0.72739}  % experiments/FIGS/data/leg2_metrics.json :: summary.sophon-public.N10000.accuracy_mean :: sha256 701c58c3240274bb :: wave 1, superseded; 3 fine-tuning seeds
\newcommand{\AccSdLegTwoSophonPublicEFour}{0.00470}  % experiments/FIGS/data/leg2_metrics.json :: summary.sophon-public.N10000.accuracy_sd :: sha256 701c58c3240274bb :: wave 1, superseded
\newcommand{\AccLegTwoSophonPublicEFive}{0.83492}  % experiments/FIGS/data/leg2_metrics.json :: summary.sophon-public.N100000.accuracy_mean :: sha256 701c58c3240274bb :: wave 1, superseded; 3 fine-tuning seeds
\newcommand{\AccSdLegTwoSophonPublicEFive}{0.00049}  % experiments/FIGS/data/leg2_metrics.json :: summary.sophon-public.N100000.accuracy_sd :: sha256 701c58c3240274bb :: wave 1, superseded
\newcommand{\AccLegTwoSophonPublicESix}{0.84632}  % experiments/FIGS/data/leg2_metrics.json :: summary.sophon-public.N1000000.accuracy_mean :: sha256 701c58c3240274bb :: wave 1, superseded; 3 fine-tuning seeds
\newcommand{\AccSdLegTwoSophonPublicESix}{0.00031}  % experiments/FIGS/data/leg2_metrics.json :: summary.sophon-public.N1000000.accuracy_sd :: sha256 701c58c3240274bb :: wave 1, superseded
\newcommand{\FtBudgetSteps}{20{,}000}  % scripts/build_ft_budget_jobs.py :: EPOCHS * STEPS :: sha256 f2ed9faa8e3bd622 :: the largest number of optimizer steps below the full training set
\newcommand{\FtLongSteps}{80{,}000}  % scripts/build_ft_budget_jobs.py :: LONG_EPOCHS * STEPS :: sha256 f2ed9faa8e3bd622 :: the largest number of optimizer steps of the long budget at LONG_N jets (JetClass, JetClass-II)
\newcommand{\DesignNTrainSelected}{109{,}938{,}036}  % experiments/FIGS/data/v2/jc2_train/counts.json :: n_selected :: sha256 b0b5a8b63aa808c0 :: selected jets of the pretrain_train files (experiments/FT/count_jc2_train.py)
\newcommand{\FtBudgetEval}{500}  % scripts/build_ft_budget_jobs.py :: STEPS :: sha256 f2ed9faa8e3bd622 :: steps between validation passes at and above SMALL_N training jets
\newcommand{\FtBudgetEvalSmall}{50}  % scripts/build_ft_budget_jobs.py :: SMALL_STEPS :: sha256 f2ed9faa8e3bd622 :: steps between validation passes below SMALL_N training jets
\newcommand{\FtBudgetEvalSmallBelow}{\ensuremath{10^{5}}}  % scripts/build_ft_budget_jobs.py :: SMALL_N :: sha256 f2ed9faa8e3bd622 :: training jets below which the validation runs every SMALL_STEPS steps
\newcommand{\FtBudgetPatience}{10}  % scripts/build_ft_budget_jobs.py :: PATIENCE :: sha256 f2ed9faa8e3bd622 :: validation passes without improvement before training stops
\newcommand{\FtHeadMult}{50}  % scripts/build_ft_budget_jobs.py :: HEAD_MULT :: sha256 f2ed9faa8e3bd622 :: the head's learning-rate multiplier
\newcommand{\FtHeadCap}{\ensuremath{5\times10^{-3}}}  % scripts/build_ft_budget_jobs.py :: HEAD_CAP :: sha256 f2ed9faa8e3bd622 :: the head's largest learning rate
\newcommand{\FtBudgetBatch}{512}  % scripts/build_ft_budget_jobs.py :: BATCH :: sha256 f2ed9faa8e3bd622
\newcommand{\FtBudgetWeightDecay}{0.01}  % scripts/build_ft_budget_jobs.py :: --optimizer-option weight_decay (every cell) :: sha256 f2ed9faa8e3bd622
\newcommand{\FtSweepPre}{\ensuremath{10^{-4}}, \ensuremath{3\times10^{-4}} and \ensuremath{10^{-3}}}  % scripts/build_ft_budget_jobs.py :: SWEEP_PRE :: sha256 f2ed9faa8e3bd622 :: backbone rates of the sweep
\newcommand{\FtSweepScratch}{\ensuremath{3\times10^{-4}}, \ensuremath{10^{-3}} and \ensuremath{3\times10^{-3}}}  % scripts/build_ft_budget_jobs.py :: SWEEP_SCRATCH :: sha256 f2ed9faa8e3bd622 :: from-scratch rates of the sweep
\newcommand{\FtSweepN}{10{,}000}  % scripts/build_ft_budget_jobs.py :: SWEEP_N :: sha256 f2ed9faa8e3bd622 :: training jets of the sweep
\newcommand{\FtValJets}{200{,}000}  % scripts/build_ft_budget_jobs.py :: VAL_N :: sha256 f2ed9faa8e3bd622 :: the fixed validation set
\newcommand{\FtValJetsBal}{20{,}480}  % scripts/build_ft_budget_jobs.py :: VAL_SIZE[jc2b] :: sha256 f2ed9faa8e3bd622 :: the balanced JetClass-II validation set
\newcommand{\LpVfourCMax}{\ensuremath{10^{4}}}  % experiments/EVAL/probe_bal.py :: C_GRID_V4 :: sha256 a7804216e3179d9d :: largest inverse regularisation of the v4 probes
\newcommand{\JcBalTestPerClass}{5{,}000}  % experiments/FT/stage_jc2_balanced.py :: TEST_PER_CLASS :: sha256 0260341d5ed75fec
\newcommand{\JcBalValPerClass}{109}  % experiments/FT/stage_jc2_balanced.py :: VAL_PER_CLASS :: sha256 0260341d5ed75fec
\newcommand{\JcBalPairTrainPerSide}{\ensuremath{10^{5}}}  % experiments/FT/stage_jc2_balanced.py :: PAIR_TRAIN_PER_SIDE :: sha256 0260341d5ed75fec
\newcommand{\JcBalPairTestPerSide}{\ensuremath{2\times10^{5}}}  % experiments/FT/stage_jc2_balanced.py :: PAIR_TEST_PER_SIDE :: sha256 0260341d5ed75fec
\newcommand{\JcBalNShort}{11}  % experiments/FIGS/data/v2/jc2bal/census.json :: counts.ft_test :: sha256 a4379ee0817254cf :: classes with fewer than 5000 selected test jets
\newcommand{\JcBalRarestTest}{224}  % experiments/FIGS/data/v2/jc2bal/census.json :: counts.ft_test (minimum) :: sha256 a4379ee0817254cf :: label_QCD_bcs
\newcommand{\JcBalShortTrainESix}{11}  % experiments/FIGS/data/v2/jc2bal/census.json :: counts.ft_heldout :: sha256 a4379ee0817254cf :: classes with fewer held-out jets than the 1e6 quota
\newcommand{\JcBalNTest}{903{,}922}  % experiments/FIGS/data/v2/jc2bal/manifest.json :: outputs[test.parquet].rows :: sha256 9c7775be6ff56ac3 :: jets of the balanced test set
\newcommand{\FtLrPre}{\ensuremath{10^{-3}}}  % experiments/FIGS/data/v2/budget/lrsweep_metrics.json :: argmax of the summed validation metric, pretrained :: sha256 7cdc148054ca902c
\newcommand{\FtLrScratchTop}{\ensuremath{10^{-3}}}  % experiments/FIGS/data/v2/budget/lrsweep_metrics.json :: argmax of the validation metric on top, from scratch :: sha256 7cdc148054ca902c :: agrees with LR of the launched job-ft-budget-bench, -jc1, -jc2b, -scratch-full and -long-scratch-full from-scratch cells
\newcommand{\FtLrScratchQg}{\ensuremath{3\times10^{-3}}}  % experiments/FIGS/data/v2/budget/lrsweep_metrics.json :: argmax of the validation metric on quark/gluon, from scratch :: sha256 7cdc148054ca902c :: agrees with LR of the launched job-ft-budget-bench and -scratch-full from-scratch cells
\newcommand{\MlstNParamsFt}{2{,}209{,}270}  % experiments/FIGS/mlst_v4.py :: our_params(2, 17) :: sha256 fa536b4866ffd651 :: trainable parameters of our fine-tuned network, two outputs (top, quark/gluon)
\newcommand{\MlstNParamsFtKin}{2{,}207{,}950}  % experiments/FIGS/mlst_v4.py :: our_params(2, 7) :: sha256 fa536b4866ffd651 :: trainable parameters of our fine-tuned network, two outputs, the 7 kinematic inputs
\newcommand{\MlstNParamsJc}{2{,}213{,}374}  % experiments/FIGS/mlst_v4.py :: our_params(10, 17) :: sha256 fa536b4866ffd651 :: trainable parameters of our fine-tuned network, the ten JetClass classes
\newcommand{\MlstProbeRunsDone}{27}  % experiments/FIGS/data/v2/index.json :: runs with a linear-probe fit (count) :: sha256 c6c4138ee5c86423 :: of the paper's runs
\newcommand{\MlstPaperRuns}{27}  % configs/arms/v2_paper_runs.txt :: runs (count) :: sha256 b04562b14ea8cbc7
\newcommand{\MlstSdRelUnc}{50\%}  % configs/arms/v2_paper_runs.txt :: relative standard error of a standard deviation from the runs of a hierarchy level, 1/sqrt(2(n-1)) :: sha256 b04562b14ea8cbc7
\newcommand{\MlstCmpJcTwoResOverFull}{1.026\,\ensuremath{\pm}\,0.005}  % experiments/FIGS/data/v2/index.json :: linear probe jc2_188 N=1000000: (1 - macro_auc_ovr) of L162 over L188, run by run :: sha256 c6c4138ee5c86423
\newcommand{\MlstCmpJcTwoResOverFullPr}{2.6\% more errors}  % experiments/FIGS/data/v2/index.json :: linear probe jc2_188 N=1000000: (1 - macro_auc_ovr) of L162 over L188, run by run :: sha256 c6c4138ee5c86423 :: prose twin: two significant figures, direction in words
\newcommand{\MlstCmpJcTwoMidOverFull}{1.227\,\ensuremath{\pm}\,0.009}  % experiments/FIGS/data/v2/index.json :: linear probe jc2_188 N=1000000: (1 - macro_auc_ovr) of R42_Q1 over L188, run by run :: sha256 c6c4138ee5c86423
\newcommand{\MlstCmpJcTwoMidOverFullPr}{23\% more errors}  % experiments/FIGS/data/v2/index.json :: linear probe jc2_188 N=1000000: (1 - macro_auc_ovr) of R42_Q1 over L188, run by run :: sha256 c6c4138ee5c86423 :: prose twin: two significant figures, direction in words
\newcommand{\MlstCmpJcTwoCoarseOverFull}{1.557\,\ensuremath{\pm}\,0.019}  % experiments/FIGS/data/v2/index.json :: linear probe jc2_188 N=1000000: (1 - macro_auc_ovr) of R16_Q1 over L188, run by run :: sha256 c6c4138ee5c86423
\newcommand{\MlstCmpJcTwoCoarseOverFullPr}{56\% more errors}  % experiments/FIGS/data/v2/index.json :: linear probe jc2_188 N=1000000: (1 - macro_auc_ovr) of R16_Q1 over L188, run by run :: sha256 c6c4138ee5c86423 :: prose twin: two significant figures, direction in words
\newcommand{\MlstCmpJcTwoResMassOverRes}{1.022 and 1.016}  % experiments/FIGS/data/v2/index.json :: linear probe jc2_188 N=1000000: (1 - macro_auc_ovr) of L162_MASS over L162, run by run :: sha256 c6c4138ee5c86423
\newcommand{\MlstCmpJcTwoResMassOverResPr}{equal within 3\%}  % experiments/FIGS/data/v2/index.json :: linear probe jc2_188 N=1000000: (1 - macro_auc_ovr) of L162_MASS over L162, run by run :: sha256 c6c4138ee5c86423 :: prose twin: two significant figures, direction in words
\newcommand{\MlstCmpJcTwoResMassOverResPrMax}{equal within 3\%}  % experiments/FIGS/data/v2/index.json :: linear probe jc2_188 N=1000000: (1 - macro_auc_ovr) of L162_MASS over L162, run by run :: sha256 c6c4138ee5c86423 :: prose twin of the run farther from 1
\newcommand{\MlstCmpJcTwoCoarseMassOverCoarse}{1.054 and 0.977}  % experiments/FIGS/data/v2/index.json :: linear probe jc2_188 N=1000000: (1 - macro_auc_ovr) of R16_Q1_MASS over R16_Q1, run by run :: sha256 c6c4138ee5c86423
\newcommand{\MlstCmpJcTwoCoarseMassOverCoarsePr}{equal within 6\%}  % experiments/FIGS/data/v2/index.json :: linear probe jc2_188 N=1000000: (1 - macro_auc_ovr) of R16_Q1_MASS over R16_Q1, run by run :: sha256 c6c4138ee5c86423 :: prose twin: two significant figures, direction in words
\newcommand{\MlstCmpJcTwoCoarseMassOverCoarsePrMax}{equal within 6\%}  % experiments/FIGS/data/v2/index.json :: linear probe jc2_188 N=1000000: (1 - macro_auc_ovr) of R16_Q1_MASS over R16_Q1, run by run :: sha256 c6c4138ee5c86423 :: prose twin of the run farther from 1
\newcommand{\MlstCmpJcTwoResMassOverFull}{1.053 and 1.037}  % experiments/FIGS/data/v2/index.json :: linear probe jc2_188 N=1000000: (1 - macro_auc_ovr) of L162_MASS over L188, run by run :: sha256 c6c4138ee5c86423
\newcommand{\MlstCmpJcTwoResMassOverFullPr}{at least 3.7\% more errors}  % experiments/FIGS/data/v2/index.json :: linear probe jc2_188 N=1000000: (1 - macro_auc_ovr) of L162_MASS over L188, run by run :: sha256 c6c4138ee5c86423 :: prose twin: two significant figures, direction in words
\newcommand{\MlstCmpJcTwoResMassOverFullPrMax}{up to 5.3\% more errors}  % experiments/FIGS/data/v2/index.json :: linear probe jc2_188 N=1000000: (1 - macro_auc_ovr) of L162_MASS over L188, run by run :: sha256 c6c4138ee5c86423 :: prose twin of the run farther from 1
\newcommand{\MlstCmpJcTwoCoarseMassOverFull}{1.617 and 1.528}  % experiments/FIGS/data/v2/index.json :: linear probe jc2_188 N=1000000: (1 - macro_auc_ovr) of R16_Q1_MASS over L188, run by run :: sha256 c6c4138ee5c86423
\newcommand{\MlstCmpJcTwoCoarseMassOverFullPr}{at least 53\% more errors}  % experiments/FIGS/data/v2/index.json :: linear probe jc2_188 N=1000000: (1 - macro_auc_ovr) of R16_Q1_MASS over L188, run by run :: sha256 c6c4138ee5c86423 :: prose twin: two significant figures, direction in words
\newcommand{\MlstCmpJcTwoCoarseMassOverFullPrMax}{up to 62\% more errors}  % experiments/FIGS/data/v2/index.json :: linear probe jc2_188 N=1000000: (1 - macro_auc_ovr) of R16_Q1_MASS over L188, run by run :: sha256 c6c4138ee5c86423 :: prose twin of the run farther from 1
\newcommand{\MlstCmpJcTwoFzeroOverFull}{1.424 and 1.462}  % experiments/FIGS/data/v2/index.json :: linear probe jc2_188 N=1000000: (1 - macro_auc_ovr) of FLAV_F0 over L188, run by run :: sha256 c6c4138ee5c86423
\newcommand{\MlstCmpJcTwoFzeroOverFullPr}{at least 42\% more errors}  % experiments/FIGS/data/v2/index.json :: linear probe jc2_188 N=1000000: (1 - macro_auc_ovr) of FLAV_F0 over L188, run by run :: sha256 c6c4138ee5c86423 :: prose twin: two significant figures, direction in words
\newcommand{\MlstCmpJcTwoFzeroOverFullPrMax}{up to 46\% more errors}  % experiments/FIGS/data/v2/index.json :: linear probe jc2_188 N=1000000: (1 - macro_auc_ovr) of FLAV_F0 over L188, run by run :: sha256 c6c4138ee5c86423 :: prose twin of the run farther from 1
\newcommand{\MlstCmpJcTwoFoneOverFull}{1.218 and 1.212}  % experiments/FIGS/data/v2/index.json :: linear probe jc2_188 N=1000000: (1 - macro_auc_ovr) of FLAV_F1 over L188, run by run :: sha256 c6c4138ee5c86423
\newcommand{\MlstCmpJcTwoFoneOverFullPr}{at least 21\% more errors}  % experiments/FIGS/data/v2/index.json :: linear probe jc2_188 N=1000000: (1 - macro_auc_ovr) of FLAV_F1 over L188, run by run :: sha256 c6c4138ee5c86423 :: prose twin: two significant figures, direction in words
\newcommand{\MlstCmpJcTwoFoneOverFullPrMax}{up to 22\% more errors}  % experiments/FIGS/data/v2/index.json :: linear probe jc2_188 N=1000000: (1 - macro_auc_ovr) of FLAV_F1 over L188, run by run :: sha256 c6c4138ee5c86423 :: prose twin of the run farther from 1
\newcommand{\MlstCmpJcTwoFonerOverFull}{1.206 and 1.208}  % experiments/FIGS/data/v2/index.json :: linear probe jc2_188 N=1000000: (1 - macro_auc_ovr) of FLAV_F1R over L188, run by run :: sha256 c6c4138ee5c86423
\newcommand{\MlstCmpJcTwoFonerOverFullPr}{at least 21\% more errors}  % experiments/FIGS/data/v2/index.json :: linear probe jc2_188 N=1000000: (1 - macro_auc_ovr) of FLAV_F1R over L188, run by run :: sha256 c6c4138ee5c86423 :: prose twin: two significant figures, direction in words
\newcommand{\MlstCmpJcTwoFonerOverFullPrMax}{up to 21\% more errors}  % experiments/FIGS/data/v2/index.json :: linear probe jc2_188 N=1000000: (1 - macro_auc_ovr) of FLAV_F1R over L188, run by run :: sha256 c6c4138ee5c86423 :: prose twin of the run farther from 1
\newcommand{\MlstCmpJcTwoFzeroOverCoarse}{0.9275 and 0.9345}  % experiments/FIGS/data/v2/index.json :: linear probe jc2_188 N=1000000: (1 - macro_auc_ovr) of FLAV_F0 over R16_Q1, run by run :: sha256 c6c4138ee5c86423
\newcommand{\MlstCmpJcTwoFzeroOverCoarsePr}{at least 6.5\% fewer errors}  % experiments/FIGS/data/v2/index.json :: linear probe jc2_188 N=1000000: (1 - macro_auc_ovr) of FLAV_F0 over R16_Q1, run by run :: sha256 c6c4138ee5c86423 :: prose twin: two significant figures, direction in words
\newcommand{\MlstCmpJcTwoFzeroOverCoarsePrMax}{up to 7.2\% fewer errors}  % experiments/FIGS/data/v2/index.json :: linear probe jc2_188 N=1000000: (1 - macro_auc_ovr) of FLAV_F0 over R16_Q1, run by run :: sha256 c6c4138ee5c86423 :: prose twin of the run farther from 1
\newcommand{\MlstCmpJcTwoFoneOverCoarse}{0.794 and 0.775}  % experiments/FIGS/data/v2/index.json :: linear probe jc2_188 N=1000000: (1 - macro_auc_ovr) of FLAV_F1 over R16_Q1, run by run :: sha256 c6c4138ee5c86423
\newcommand{\MlstCmpJcTwoFoneOverCoarsePr}{at least 21\% fewer errors}  % experiments/FIGS/data/v2/index.json :: linear probe jc2_188 N=1000000: (1 - macro_auc_ovr) of FLAV_F1 over R16_Q1, run by run :: sha256 c6c4138ee5c86423 :: prose twin: two significant figures, direction in words
\newcommand{\MlstCmpJcTwoFoneOverCoarsePrMax}{up to 23\% fewer errors}  % experiments/FIGS/data/v2/index.json :: linear probe jc2_188 N=1000000: (1 - macro_auc_ovr) of FLAV_F1 over R16_Q1, run by run :: sha256 c6c4138ee5c86423 :: prose twin of the run farther from 1
\newcommand{\MlstCmpJcTwoFonerOverCoarse}{0.786 and 0.772}  % experiments/FIGS/data/v2/index.json :: linear probe jc2_188 N=1000000: (1 - macro_auc_ovr) of FLAV_F1R over R16_Q1, run by run :: sha256 c6c4138ee5c86423
\newcommand{\MlstCmpJcTwoFonerOverCoarsePr}{at least 21\% fewer errors}  % experiments/FIGS/data/v2/index.json :: linear probe jc2_188 N=1000000: (1 - macro_auc_ovr) of FLAV_F1R over R16_Q1, run by run :: sha256 c6c4138ee5c86423 :: prose twin: two significant figures, direction in words
\newcommand{\MlstCmpJcTwoFonerOverCoarsePrMax}{up to 23\% fewer errors}  % experiments/FIGS/data/v2/index.json :: linear probe jc2_188 N=1000000: (1 - macro_auc_ovr) of FLAV_F1R over R16_Q1, run by run :: sha256 c6c4138ee5c86423 :: prose twin of the run farther from 1
\newcommand{\MlstCmpJcTwoFoneOverFzero}{0.856 and 0.829}  % experiments/FIGS/data/v2/index.json :: linear probe jc2_188 N=1000000: (1 - macro_auc_ovr) of FLAV_F1 over FLAV_F0, run by run :: sha256 c6c4138ee5c86423
\newcommand{\MlstCmpJcTwoFoneOverFzeroPr}{at least 14\% fewer errors}  % experiments/FIGS/data/v2/index.json :: linear probe jc2_188 N=1000000: (1 - macro_auc_ovr) of FLAV_F1 over FLAV_F0, run by run :: sha256 c6c4138ee5c86423 :: prose twin: two significant figures, direction in words
\newcommand{\MlstCmpJcTwoFoneOverFzeroPrMax}{up to 17\% fewer errors}  % experiments/FIGS/data/v2/index.json :: linear probe jc2_188 N=1000000: (1 - macro_auc_ovr) of FLAV_F1 over FLAV_F0, run by run :: sha256 c6c4138ee5c86423 :: prose twin of the run farther from 1
\newcommand{\MlstCmpJcTwoFonerOverFzero}{0.847 and 0.826}  % experiments/FIGS/data/v2/index.json :: linear probe jc2_188 N=1000000: (1 - macro_auc_ovr) of FLAV_F1R over FLAV_F0, run by run :: sha256 c6c4138ee5c86423
\newcommand{\MlstCmpJcTwoFonerOverFzeroPr}{at least 15\% fewer errors}  % experiments/FIGS/data/v2/index.json :: linear probe jc2_188 N=1000000: (1 - macro_auc_ovr) of FLAV_F1R over FLAV_F0, run by run :: sha256 c6c4138ee5c86423 :: prose twin: two significant figures, direction in words
\newcommand{\MlstCmpJcTwoFonerOverFzeroPrMax}{up to 17\% fewer errors}  % experiments/FIGS/data/v2/index.json :: linear probe jc2_188 N=1000000: (1 - macro_auc_ovr) of FLAV_F1R over FLAV_F0, run by run :: sha256 c6c4138ee5c86423 :: prose twin of the run farther from 1
\newcommand{\MlstCmpJcTwoFonerOverFone}{0.9899 and 0.9966}  % experiments/FIGS/data/v2/index.json :: linear probe jc2_188 N=1000000: (1 - macro_auc_ovr) of FLAV_F1R over FLAV_F1, run by run :: sha256 c6c4138ee5c86423
\newcommand{\MlstCmpJcTwoFonerOverFonePr}{equal within 2\%}  % experiments/FIGS/data/v2/index.json :: linear probe jc2_188 N=1000000: (1 - macro_auc_ovr) of FLAV_F1R over FLAV_F1, run by run :: sha256 c6c4138ee5c86423 :: prose twin: two significant figures, direction in words
\newcommand{\MlstCmpJcTwoFonerOverFonePrMax}{equal within 2\%}  % experiments/FIGS/data/v2/index.json :: linear probe jc2_188 N=1000000: (1 - macro_auc_ovr) of FLAV_F1R over FLAV_F1, run by run :: sha256 c6c4138ee5c86423 :: prose twin of the run farther from 1
\newcommand{\MlstCmpJcTwoMidOverRes}{1.197\,\ensuremath{\pm}\,0.013}  % experiments/FIGS/data/v2/index.json :: linear probe jc2_188 N=1000000: (1 - macro_auc_ovr) of R42_Q1 over L162, run by run :: sha256 c6c4138ee5c86423
\newcommand{\MlstCmpJcTwoMidOverResPr}{20\% more errors}  % experiments/FIGS/data/v2/index.json :: linear probe jc2_188 N=1000000: (1 - macro_auc_ovr) of R42_Q1 over L162, run by run :: sha256 c6c4138ee5c86423 :: prose twin: two significant figures, direction in words
\newcommand{\MlstCmpJcTwoCoarseOverRes}{1.518\,\ensuremath{\pm}\,0.025}  % experiments/FIGS/data/v2/index.json :: linear probe jc2_188 N=1000000: (1 - macro_auc_ovr) of R16_Q1 over L162, run by run :: sha256 c6c4138ee5c86423
\newcommand{\MlstCmpJcTwoCoarseOverResPr}{52\% more errors}  % experiments/FIGS/data/v2/index.json :: linear probe jc2_188 N=1000000: (1 - macro_auc_ovr) of R16_Q1 over L162, run by run :: sha256 c6c4138ee5c86423 :: prose twin: two significant figures, direction in words
\newcommand{\MlstCmpJcTwoCoarseOverMid}{1.269\,\ensuremath{\pm}\,0.007}  % experiments/FIGS/data/v2/index.json :: linear probe jc2_188 N=1000000: (1 - macro_auc_ovr) of R16_Q1 over R42_Q1, run by run :: sha256 c6c4138ee5c86423
\newcommand{\MlstCmpJcTwoCoarseOverMidPr}{27\% more errors}  % experiments/FIGS/data/v2/index.json :: linear probe jc2_188 N=1000000: (1 - macro_auc_ovr) of R16_Q1 over R42_Q1, run by run :: sha256 c6c4138ee5c86423 :: prose twin: two significant figures, direction in words
\newcommand{\MlstCmpJcTwoRandOverFull}{1.125\,\ensuremath{\pm}\,0.008}  % experiments/FIGS/data/v2/index.json :: linear probe jc2_188 N=1000000: (1 - macro_auc_ovr) of each random partition over the L188 run of its seed, mean +- SD over the partitions :: sha256 c6c4138ee5c86423
\newcommand{\MlstCmpJcTwoRandOverFullPr}{12\% more errors}  % experiments/FIGS/data/v2/index.json :: linear probe jc2_188 N=1000000: (1 - macro_auc_ovr) of each random partition over the L188 run of its seed, mean +- SD over the partitions :: sha256 c6c4138ee5c86423 :: prose twin: two significant figures, direction in words
\newcommand{\MlstCmpJcTwoRandOverMid}{0.918\,\ensuremath{\pm}\,0.007}  % experiments/FIGS/data/v2/index.json :: linear probe jc2_188 N=1000000: (1 - macro_auc_ovr) of each random partition over the R42_Q1 run of its seed, mean +- SD over the partitions :: sha256 c6c4138ee5c86423
\newcommand{\MlstCmpJcTwoRandOverMidPr}{8.2\% fewer errors}  % experiments/FIGS/data/v2/index.json :: linear probe jc2_188 N=1000000: (1 - macro_auc_ovr) of each random partition over the R42_Q1 run of its seed, mean +- SD over the partitions :: sha256 c6c4138ee5c86423 :: prose twin: two significant figures, direction in words
\newcommand{\MlstCmpJcTwoRandOverCoarse}{0.724\,\ensuremath{\pm}\,0.008}  % experiments/FIGS/data/v2/index.json :: linear probe jc2_188 N=1000000: (1 - macro_auc_ovr) of each random partition over the R16_Q1 run of its seed, mean +- SD over the partitions :: sha256 c6c4138ee5c86423
\newcommand{\MlstCmpJcTwoRandOverCoarsePr}{28\% fewer errors}  % experiments/FIGS/data/v2/index.json :: linear probe jc2_188 N=1000000: (1 - macro_auc_ovr) of each random partition over the R16_Q1 run of its seed, mean +- SD over the partitions :: sha256 c6c4138ee5c86423 :: prose twin: two significant figures, direction in words
\newcommand{\MlstCmpJcTwoTrainedResOverFull}{1.026\,\ensuremath{\pm}\,0.005}  % experiments/FIGS/data/v2/index.json :: linear probe jc2_188 N=1000000: (1 - macro_auc_trained) of L162 over L188, run by run :: sha256 c6c4138ee5c86423
\newcommand{\MlstCmpJcTwoTrainedResOverFullPr}{2.6\% more errors}  % experiments/FIGS/data/v2/index.json :: linear probe jc2_188 N=1000000: (1 - macro_auc_trained) of L162 over L188, run by run :: sha256 c6c4138ee5c86423 :: prose twin: two significant figures, direction in words
\newcommand{\MlstCmpJcTwoTrainedMidOverFull}{1.227\,\ensuremath{\pm}\,0.009}  % experiments/FIGS/data/v2/index.json :: linear probe jc2_188 N=1000000: (1 - macro_auc_trained) of R42_Q1 over L188, run by run :: sha256 c6c4138ee5c86423
\newcommand{\MlstCmpJcTwoTrainedMidOverFullPr}{23\% more errors}  % experiments/FIGS/data/v2/index.json :: linear probe jc2_188 N=1000000: (1 - macro_auc_trained) of R42_Q1 over L188, run by run :: sha256 c6c4138ee5c86423 :: prose twin: two significant figures, direction in words
\newcommand{\MlstCmpJcTwoTrainedCoarseOverFull}{1.557\,\ensuremath{\pm}\,0.019}  % experiments/FIGS/data/v2/index.json :: linear probe jc2_188 N=1000000: (1 - macro_auc_trained) of R16_Q1 over L188, run by run :: sha256 c6c4138ee5c86423
\newcommand{\MlstCmpJcTwoTrainedCoarseOverFullPr}{56\% more errors}  % experiments/FIGS/data/v2/index.json :: linear probe jc2_188 N=1000000: (1 - macro_auc_trained) of R16_Q1 over L188, run by run :: sha256 c6c4138ee5c86423 :: prose twin: two significant figures, direction in words
\newcommand{\MlstCmpJcTwoTrainedResMassOverRes}{1.022 and 1.016}  % experiments/FIGS/data/v2/index.json :: linear probe jc2_188 N=1000000: (1 - macro_auc_trained) of L162_MASS over L162, run by run :: sha256 c6c4138ee5c86423
\newcommand{\MlstCmpJcTwoTrainedResMassOverResPr}{equal within 3\%}  % experiments/FIGS/data/v2/index.json :: linear probe jc2_188 N=1000000: (1 - macro_auc_trained) of L162_MASS over L162, run by run :: sha256 c6c4138ee5c86423 :: prose twin: two significant figures, direction in words
\newcommand{\MlstCmpJcTwoTrainedResMassOverResPrMax}{equal within 3\%}  % experiments/FIGS/data/v2/index.json :: linear probe jc2_188 N=1000000: (1 - macro_auc_trained) of L162_MASS over L162, run by run :: sha256 c6c4138ee5c86423 :: prose twin of the run farther from 1
\newcommand{\MlstCmpJcTwoTrainedCoarseMassOverCoarse}{1.054 and 0.977}  % experiments/FIGS/data/v2/index.json :: linear probe jc2_188 N=1000000: (1 - macro_auc_trained) of R16_Q1_MASS over R16_Q1, run by run :: sha256 c6c4138ee5c86423
\newcommand{\MlstCmpJcTwoTrainedCoarseMassOverCoarsePr}{equal within 6\%}  % experiments/FIGS/data/v2/index.json :: linear probe jc2_188 N=1000000: (1 - macro_auc_trained) of R16_Q1_MASS over R16_Q1, run by run :: sha256 c6c4138ee5c86423 :: prose twin: two significant figures, direction in words
\newcommand{\MlstCmpJcTwoTrainedCoarseMassOverCoarsePrMax}{equal within 6\%}  % experiments/FIGS/data/v2/index.json :: linear probe jc2_188 N=1000000: (1 - macro_auc_trained) of R16_Q1_MASS over R16_Q1, run by run :: sha256 c6c4138ee5c86423 :: prose twin of the run farther from 1
\newcommand{\MlstCmpJcTwoTrainedResMassOverFull}{1.053 and 1.037}  % experiments/FIGS/data/v2/index.json :: linear probe jc2_188 N=1000000: (1 - macro_auc_trained) of L162_MASS over L188, run by run :: sha256 c6c4138ee5c86423
\newcommand{\MlstCmpJcTwoTrainedResMassOverFullPr}{at least 3.7\% more errors}  % experiments/FIGS/data/v2/index.json :: linear probe jc2_188 N=1000000: (1 - macro_auc_trained) of L162_MASS over L188, run by run :: sha256 c6c4138ee5c86423 :: prose twin: two significant figures, direction in words
\newcommand{\MlstCmpJcTwoTrainedResMassOverFullPrMax}{up to 5.3\% more errors}  % experiments/FIGS/data/v2/index.json :: linear probe jc2_188 N=1000000: (1 - macro_auc_trained) of L162_MASS over L188, run by run :: sha256 c6c4138ee5c86423 :: prose twin of the run farther from 1
\newcommand{\MlstCmpJcTwoTrainedCoarseMassOverFull}{1.617 and 1.528}  % experiments/FIGS/data/v2/index.json :: linear probe jc2_188 N=1000000: (1 - macro_auc_trained) of R16_Q1_MASS over L188, run by run :: sha256 c6c4138ee5c86423
\newcommand{\MlstCmpJcTwoTrainedCoarseMassOverFullPr}{at least 53\% more errors}  % experiments/FIGS/data/v2/index.json :: linear probe jc2_188 N=1000000: (1 - macro_auc_trained) of R16_Q1_MASS over L188, run by run :: sha256 c6c4138ee5c86423 :: prose twin: two significant figures, direction in words
\newcommand{\MlstCmpJcTwoTrainedCoarseMassOverFullPrMax}{up to 62\% more errors}  % experiments/FIGS/data/v2/index.json :: linear probe jc2_188 N=1000000: (1 - macro_auc_trained) of R16_Q1_MASS over L188, run by run :: sha256 c6c4138ee5c86423 :: prose twin of the run farther from 1
\newcommand{\MlstCmpJcTwoTrainedFzeroOverFull}{1.424 and 1.462}  % experiments/FIGS/data/v2/index.json :: linear probe jc2_188 N=1000000: (1 - macro_auc_trained) of FLAV_F0 over L188, run by run :: sha256 c6c4138ee5c86423
\newcommand{\MlstCmpJcTwoTrainedFzeroOverFullPr}{at least 42\% more errors}  % experiments/FIGS/data/v2/index.json :: linear probe jc2_188 N=1000000: (1 - macro_auc_trained) of FLAV_F0 over L188, run by run :: sha256 c6c4138ee5c86423 :: prose twin: two significant figures, direction in words
\newcommand{\MlstCmpJcTwoTrainedFzeroOverFullPrMax}{up to 46\% more errors}  % experiments/FIGS/data/v2/index.json :: linear probe jc2_188 N=1000000: (1 - macro_auc_trained) of FLAV_F0 over L188, run by run :: sha256 c6c4138ee5c86423 :: prose twin of the run farther from 1
\newcommand{\MlstCmpJcTwoTrainedFoneOverFull}{1.218 and 1.212}  % experiments/FIGS/data/v2/index.json :: linear probe jc2_188 N=1000000: (1 - macro_auc_trained) of FLAV_F1 over L188, run by run :: sha256 c6c4138ee5c86423
\newcommand{\MlstCmpJcTwoTrainedFoneOverFullPr}{at least 21\% more errors}  % experiments/FIGS/data/v2/index.json :: linear probe jc2_188 N=1000000: (1 - macro_auc_trained) of FLAV_F1 over L188, run by run :: sha256 c6c4138ee5c86423 :: prose twin: two significant figures, direction in words
\newcommand{\MlstCmpJcTwoTrainedFoneOverFullPrMax}{up to 22\% more errors}  % experiments/FIGS/data/v2/index.json :: linear probe jc2_188 N=1000000: (1 - macro_auc_trained) of FLAV_F1 over L188, run by run :: sha256 c6c4138ee5c86423 :: prose twin of the run farther from 1
\newcommand{\MlstCmpJcTwoTrainedFonerOverFull}{1.206 and 1.208}  % experiments/FIGS/data/v2/index.json :: linear probe jc2_188 N=1000000: (1 - macro_auc_trained) of FLAV_F1R over L188, run by run :: sha256 c6c4138ee5c86423
\newcommand{\MlstCmpJcTwoTrainedFonerOverFullPr}{at least 21\% more errors}  % experiments/FIGS/data/v2/index.json :: linear probe jc2_188 N=1000000: (1 - macro_auc_trained) of FLAV_F1R over L188, run by run :: sha256 c6c4138ee5c86423 :: prose twin: two significant figures, direction in words
\newcommand{\MlstCmpJcTwoTrainedFonerOverFullPrMax}{up to 21\% more errors}  % experiments/FIGS/data/v2/index.json :: linear probe jc2_188 N=1000000: (1 - macro_auc_trained) of FLAV_F1R over L188, run by run :: sha256 c6c4138ee5c86423 :: prose twin of the run farther from 1
\newcommand{\MlstCmpJcTwoTrainedFzeroOverCoarse}{0.9275 and 0.9345}  % experiments/FIGS/data/v2/index.json :: linear probe jc2_188 N=1000000: (1 - macro_auc_trained) of FLAV_F0 over R16_Q1, run by run :: sha256 c6c4138ee5c86423
\newcommand{\MlstCmpJcTwoTrainedFzeroOverCoarsePr}{at least 6.5\% fewer errors}  % experiments/FIGS/data/v2/index.json :: linear probe jc2_188 N=1000000: (1 - macro_auc_trained) of FLAV_F0 over R16_Q1, run by run :: sha256 c6c4138ee5c86423 :: prose twin: two significant figures, direction in words
\newcommand{\MlstCmpJcTwoTrainedFzeroOverCoarsePrMax}{up to 7.2\% fewer errors}  % experiments/FIGS/data/v2/index.json :: linear probe jc2_188 N=1000000: (1 - macro_auc_trained) of FLAV_F0 over R16_Q1, run by run :: sha256 c6c4138ee5c86423 :: prose twin of the run farther from 1
\newcommand{\MlstCmpJcTwoTrainedFoneOverCoarse}{0.794 and 0.775}  % experiments/FIGS/data/v2/index.json :: linear probe jc2_188 N=1000000: (1 - macro_auc_trained) of FLAV_F1 over R16_Q1, run by run :: sha256 c6c4138ee5c86423
\newcommand{\MlstCmpJcTwoTrainedFoneOverCoarsePr}{at least 21\% fewer errors}  % experiments/FIGS/data/v2/index.json :: linear probe jc2_188 N=1000000: (1 - macro_auc_trained) of FLAV_F1 over R16_Q1, run by run :: sha256 c6c4138ee5c86423 :: prose twin: two significant figures, direction in words
\newcommand{\MlstCmpJcTwoTrainedFoneOverCoarsePrMax}{up to 23\% fewer errors}  % experiments/FIGS/data/v2/index.json :: linear probe jc2_188 N=1000000: (1 - macro_auc_trained) of FLAV_F1 over R16_Q1, run by run :: sha256 c6c4138ee5c86423 :: prose twin of the run farther from 1
\newcommand{\MlstCmpJcTwoTrainedFonerOverCoarse}{0.786 and 0.772}  % experiments/FIGS/data/v2/index.json :: linear probe jc2_188 N=1000000: (1 - macro_auc_trained) of FLAV_F1R over R16_Q1, run by run :: sha256 c6c4138ee5c86423
\newcommand{\MlstCmpJcTwoTrainedFonerOverCoarsePr}{at least 21\% fewer errors}  % experiments/FIGS/data/v2/index.json :: linear probe jc2_188 N=1000000: (1 - macro_auc_trained) of FLAV_F1R over R16_Q1, run by run :: sha256 c6c4138ee5c86423 :: prose twin: two significant figures, direction in words
\newcommand{\MlstCmpJcTwoTrainedFonerOverCoarsePrMax}{up to 23\% fewer errors}  % experiments/FIGS/data/v2/index.json :: linear probe jc2_188 N=1000000: (1 - macro_auc_trained) of FLAV_F1R over R16_Q1, run by run :: sha256 c6c4138ee5c86423 :: prose twin of the run farther from 1
\newcommand{\MlstCmpJcTwoTrainedFoneOverFzero}{0.856 and 0.829}  % experiments/FIGS/data/v2/index.json :: linear probe jc2_188 N=1000000: (1 - macro_auc_trained) of FLAV_F1 over FLAV_F0, run by run :: sha256 c6c4138ee5c86423
\newcommand{\MlstCmpJcTwoTrainedFoneOverFzeroPr}{at least 14\% fewer errors}  % experiments/FIGS/data/v2/index.json :: linear probe jc2_188 N=1000000: (1 - macro_auc_trained) of FLAV_F1 over FLAV_F0, run by run :: sha256 c6c4138ee5c86423 :: prose twin: two significant figures, direction in words
\newcommand{\MlstCmpJcTwoTrainedFoneOverFzeroPrMax}{up to 17\% fewer errors}  % experiments/FIGS/data/v2/index.json :: linear probe jc2_188 N=1000000: (1 - macro_auc_trained) of FLAV_F1 over FLAV_F0, run by run :: sha256 c6c4138ee5c86423 :: prose twin of the run farther from 1
\newcommand{\MlstCmpJcTwoTrainedFonerOverFzero}{0.847 and 0.826}  % experiments/FIGS/data/v2/index.json :: linear probe jc2_188 N=1000000: (1 - macro_auc_trained) of FLAV_F1R over FLAV_F0, run by run :: sha256 c6c4138ee5c86423
\newcommand{\MlstCmpJcTwoTrainedFonerOverFzeroPr}{at least 15\% fewer errors}  % experiments/FIGS/data/v2/index.json :: linear probe jc2_188 N=1000000: (1 - macro_auc_trained) of FLAV_F1R over FLAV_F0, run by run :: sha256 c6c4138ee5c86423 :: prose twin: two significant figures, direction in words
\newcommand{\MlstCmpJcTwoTrainedFonerOverFzeroPrMax}{up to 17\% fewer errors}  % experiments/FIGS/data/v2/index.json :: linear probe jc2_188 N=1000000: (1 - macro_auc_trained) of FLAV_F1R over FLAV_F0, run by run :: sha256 c6c4138ee5c86423 :: prose twin of the run farther from 1
\newcommand{\MlstCmpJcTwoTrainedFonerOverFone}{0.9899 and 0.9966}  % experiments/FIGS/data/v2/index.json :: linear probe jc2_188 N=1000000: (1 - macro_auc_trained) of FLAV_F1R over FLAV_F1, run by run :: sha256 c6c4138ee5c86423
\newcommand{\MlstCmpJcTwoTrainedFonerOverFonePr}{equal within 2\%}  % experiments/FIGS/data/v2/index.json :: linear probe jc2_188 N=1000000: (1 - macro_auc_trained) of FLAV_F1R over FLAV_F1, run by run :: sha256 c6c4138ee5c86423 :: prose twin: two significant figures, direction in words
\newcommand{\MlstCmpJcTwoTrainedFonerOverFonePrMax}{equal within 2\%}  % experiments/FIGS/data/v2/index.json :: linear probe jc2_188 N=1000000: (1 - macro_auc_trained) of FLAV_F1R over FLAV_F1, run by run :: sha256 c6c4138ee5c86423 :: prose twin of the run farther from 1
\newcommand{\MlstCmpJcTwoTrainedMidOverRes}{1.197\,\ensuremath{\pm}\,0.013}  % experiments/FIGS/data/v2/index.json :: linear probe jc2_188 N=1000000: (1 - macro_auc_trained) of R42_Q1 over L162, run by run :: sha256 c6c4138ee5c86423
\newcommand{\MlstCmpJcTwoTrainedMidOverResPr}{20\% more errors}  % experiments/FIGS/data/v2/index.json :: linear probe jc2_188 N=1000000: (1 - macro_auc_trained) of R42_Q1 over L162, run by run :: sha256 c6c4138ee5c86423 :: prose twin: two significant figures, direction in words
\newcommand{\MlstCmpJcTwoTrainedCoarseOverRes}{1.518\,\ensuremath{\pm}\,0.025}  % experiments/FIGS/data/v2/index.json :: linear probe jc2_188 N=1000000: (1 - macro_auc_trained) of R16_Q1 over L162, run by run :: sha256 c6c4138ee5c86423
\newcommand{\MlstCmpJcTwoTrainedCoarseOverResPr}{52\% more errors}  % experiments/FIGS/data/v2/index.json :: linear probe jc2_188 N=1000000: (1 - macro_auc_trained) of R16_Q1 over L162, run by run :: sha256 c6c4138ee5c86423 :: prose twin: two significant figures, direction in words
\newcommand{\MlstCmpJcTwoTrainedCoarseOverMid}{1.269\,\ensuremath{\pm}\,0.007}  % experiments/FIGS/data/v2/index.json :: linear probe jc2_188 N=1000000: (1 - macro_auc_trained) of R16_Q1 over R42_Q1, run by run :: sha256 c6c4138ee5c86423
\newcommand{\MlstCmpJcTwoTrainedCoarseOverMidPr}{27\% more errors}  % experiments/FIGS/data/v2/index.json :: linear probe jc2_188 N=1000000: (1 - macro_auc_trained) of R16_Q1 over R42_Q1, run by run :: sha256 c6c4138ee5c86423 :: prose twin: two significant figures, direction in words
\newcommand{\MlstCmpJcTwoTrainedRandOverFull}{1.125\,\ensuremath{\pm}\,0.008}  % experiments/FIGS/data/v2/index.json :: linear probe jc2_188 N=1000000: (1 - macro_auc_trained) of each random partition over the L188 run of its seed, mean +- SD over the partitions :: sha256 c6c4138ee5c86423
\newcommand{\MlstCmpJcTwoTrainedRandOverFullPr}{12\% more errors}  % experiments/FIGS/data/v2/index.json :: linear probe jc2_188 N=1000000: (1 - macro_auc_trained) of each random partition over the L188 run of its seed, mean +- SD over the partitions :: sha256 c6c4138ee5c86423 :: prose twin: two significant figures, direction in words
\newcommand{\MlstCmpJcTwoTrainedRandOverMid}{0.918\,\ensuremath{\pm}\,0.007}  % experiments/FIGS/data/v2/index.json :: linear probe jc2_188 N=1000000: (1 - macro_auc_trained) of each random partition over the R42_Q1 run of its seed, mean +- SD over the partitions :: sha256 c6c4138ee5c86423
\newcommand{\MlstCmpJcTwoTrainedRandOverMidPr}{8.2\% fewer errors}  % experiments/FIGS/data/v2/index.json :: linear probe jc2_188 N=1000000: (1 - macro_auc_trained) of each random partition over the R42_Q1 run of its seed, mean +- SD over the partitions :: sha256 c6c4138ee5c86423 :: prose twin: two significant figures, direction in words
\newcommand{\MlstCmpJcTwoTrainedRandOverCoarse}{0.724\,\ensuremath{\pm}\,0.008}  % experiments/FIGS/data/v2/index.json :: linear probe jc2_188 N=1000000: (1 - macro_auc_trained) of each random partition over the R16_Q1 run of its seed, mean +- SD over the partitions :: sha256 c6c4138ee5c86423
\newcommand{\MlstCmpJcTwoTrainedRandOverCoarsePr}{28\% fewer errors}  % experiments/FIGS/data/v2/index.json :: linear probe jc2_188 N=1000000: (1 - macro_auc_trained) of each random partition over the R16_Q1 run of its seed, mean +- SD over the partitions :: sha256 c6c4138ee5c86423 :: prose twin: two significant figures, direction in words
\newcommand{\MlstCmpJcTwoResResOverFull}{1.003\,\ensuremath{\pm}\,0.005}  % experiments/FIGS/data/v2/index.json :: linear probe jc2_188 N=1000000: (1 - macro_auc_resonant) of L162 over L188, run by run :: sha256 c6c4138ee5c86423
\newcommand{\MlstCmpJcTwoResResOverFullPr}{equal within 1\%}  % experiments/FIGS/data/v2/index.json :: linear probe jc2_188 N=1000000: (1 - macro_auc_resonant) of L162 over L188, run by run :: sha256 c6c4138ee5c86423 :: prose twin: two significant figures, direction in words
\newcommand{\MlstCmpJcTwoResMidOverFull}{1.262\,\ensuremath{\pm}\,0.012}  % experiments/FIGS/data/v2/index.json :: linear probe jc2_188 N=1000000: (1 - macro_auc_resonant) of R42_Q1 over L188, run by run :: sha256 c6c4138ee5c86423
\newcommand{\MlstCmpJcTwoResMidOverFullPr}{26\% more errors}  % experiments/FIGS/data/v2/index.json :: linear probe jc2_188 N=1000000: (1 - macro_auc_resonant) of R42_Q1 over L188, run by run :: sha256 c6c4138ee5c86423 :: prose twin: two significant figures, direction in words
\newcommand{\MlstCmpJcTwoResCoarseOverFull}{1.590\,\ensuremath{\pm}\,0.032}  % experiments/FIGS/data/v2/index.json :: linear probe jc2_188 N=1000000: (1 - macro_auc_resonant) of R16_Q1 over L188, run by run :: sha256 c6c4138ee5c86423
\newcommand{\MlstCmpJcTwoResCoarseOverFullPr}{59\% more errors}  % experiments/FIGS/data/v2/index.json :: linear probe jc2_188 N=1000000: (1 - macro_auc_resonant) of R16_Q1 over L188, run by run :: sha256 c6c4138ee5c86423 :: prose twin: two significant figures, direction in words
\newcommand{\MlstCmpJcTwoResResMassOverRes}{1.025 and 1.020}  % experiments/FIGS/data/v2/index.json :: linear probe jc2_188 N=1000000: (1 - macro_auc_resonant) of L162_MASS over L162, run by run :: sha256 c6c4138ee5c86423
\newcommand{\MlstCmpJcTwoResResMassOverResPr}{equal within 3\%}  % experiments/FIGS/data/v2/index.json :: linear probe jc2_188 N=1000000: (1 - macro_auc_resonant) of L162_MASS over L162, run by run :: sha256 c6c4138ee5c86423 :: prose twin: two significant figures, direction in words
\newcommand{\MlstCmpJcTwoResResMassOverResPrMax}{equal within 3\%}  % experiments/FIGS/data/v2/index.json :: linear probe jc2_188 N=1000000: (1 - macro_auc_resonant) of L162_MASS over L162, run by run :: sha256 c6c4138ee5c86423 :: prose twin of the run farther from 1
\newcommand{\MlstCmpJcTwoResCoarseMassOverCoarse}{1.08 and 0.96}  % experiments/FIGS/data/v2/index.json :: linear probe jc2_188 N=1000000: (1 - macro_auc_resonant) of R16_Q1_MASS over R16_Q1, run by run :: sha256 c6c4138ee5c86423
\newcommand{\MlstCmpJcTwoResCoarseMassOverCoarsePr}{equal within 8\%}  % experiments/FIGS/data/v2/index.json :: linear probe jc2_188 N=1000000: (1 - macro_auc_resonant) of R16_Q1_MASS over R16_Q1, run by run :: sha256 c6c4138ee5c86423 :: prose twin: two significant figures, direction in words
\newcommand{\MlstCmpJcTwoResCoarseMassOverCoarsePrMax}{equal within 8\%}  % experiments/FIGS/data/v2/index.json :: linear probe jc2_188 N=1000000: (1 - macro_auc_resonant) of R16_Q1_MASS over R16_Q1, run by run :: sha256 c6c4138ee5c86423 :: prose twin of the run farther from 1
\newcommand{\MlstCmpJcTwoResResMassOverFull}{1.033 and 1.018}  % experiments/FIGS/data/v2/index.json :: linear probe jc2_188 N=1000000: (1 - macro_auc_resonant) of L162_MASS over L188, run by run :: sha256 c6c4138ee5c86423
\newcommand{\MlstCmpJcTwoResResMassOverFullPr}{equal within 4\%}  % experiments/FIGS/data/v2/index.json :: linear probe jc2_188 N=1000000: (1 - macro_auc_resonant) of L162_MASS over L188, run by run :: sha256 c6c4138ee5c86423 :: prose twin: two significant figures, direction in words
\newcommand{\MlstCmpJcTwoResResMassOverFullPrMax}{equal within 4\%}  % experiments/FIGS/data/v2/index.json :: linear probe jc2_188 N=1000000: (1 - macro_auc_resonant) of L162_MASS over L188, run by run :: sha256 c6c4138ee5c86423 :: prose twin of the run farther from 1
\newcommand{\MlstCmpJcTwoResCoarseMassOverFull}{1.67 and 1.55}  % experiments/FIGS/data/v2/index.json :: linear probe jc2_188 N=1000000: (1 - macro_auc_resonant) of R16_Q1_MASS over L188, run by run :: sha256 c6c4138ee5c86423
\newcommand{\MlstCmpJcTwoResCoarseMassOverFullPr}{at least 55\% more errors}  % experiments/FIGS/data/v2/index.json :: linear probe jc2_188 N=1000000: (1 - macro_auc_resonant) of R16_Q1_MASS over L188, run by run :: sha256 c6c4138ee5c86423 :: prose twin: two significant figures, direction in words
\newcommand{\MlstCmpJcTwoResCoarseMassOverFullPrMax}{up to 67\% more errors}  % experiments/FIGS/data/v2/index.json :: linear probe jc2_188 N=1000000: (1 - macro_auc_resonant) of R16_Q1_MASS over L188, run by run :: sha256 c6c4138ee5c86423 :: prose twin of the run farther from 1
\newcommand{\MlstCmpJcTwoResFzeroOverFull}{1.408 and 1.449}  % experiments/FIGS/data/v2/index.json :: linear probe jc2_188 N=1000000: (1 - macro_auc_resonant) of FLAV_F0 over L188, run by run :: sha256 c6c4138ee5c86423
\newcommand{\MlstCmpJcTwoResFzeroOverFullPr}{at least 41\% more errors}  % experiments/FIGS/data/v2/index.json :: linear probe jc2_188 N=1000000: (1 - macro_auc_resonant) of FLAV_F0 over L188, run by run :: sha256 c6c4138ee5c86423 :: prose twin: two significant figures, direction in words
\newcommand{\MlstCmpJcTwoResFzeroOverFullPrMax}{up to 45\% more errors}  % experiments/FIGS/data/v2/index.json :: linear probe jc2_188 N=1000000: (1 - macro_auc_resonant) of FLAV_F0 over L188, run by run :: sha256 c6c4138ee5c86423 :: prose twin of the run farther from 1
\newcommand{\MlstCmpJcTwoResFoneOverFull}{1.210 and 1.204}  % experiments/FIGS/data/v2/index.json :: linear probe jc2_188 N=1000000: (1 - macro_auc_resonant) of FLAV_F1 over L188, run by run :: sha256 c6c4138ee5c86423
\newcommand{\MlstCmpJcTwoResFoneOverFullPr}{at least 20\% more errors}  % experiments/FIGS/data/v2/index.json :: linear probe jc2_188 N=1000000: (1 - macro_auc_resonant) of FLAV_F1 over L188, run by run :: sha256 c6c4138ee5c86423 :: prose twin: two significant figures, direction in words
\newcommand{\MlstCmpJcTwoResFoneOverFullPrMax}{up to 21\% more errors}  % experiments/FIGS/data/v2/index.json :: linear probe jc2_188 N=1000000: (1 - macro_auc_resonant) of FLAV_F1 over L188, run by run :: sha256 c6c4138ee5c86423 :: prose twin of the run farther from 1
\newcommand{\MlstCmpJcTwoResFonerOverFull}{1.187 and 1.188}  % experiments/FIGS/data/v2/index.json :: linear probe jc2_188 N=1000000: (1 - macro_auc_resonant) of FLAV_F1R over L188, run by run :: sha256 c6c4138ee5c86423
\newcommand{\MlstCmpJcTwoResFonerOverFullPr}{at least 19\% more errors}  % experiments/FIGS/data/v2/index.json :: linear probe jc2_188 N=1000000: (1 - macro_auc_resonant) of FLAV_F1R over L188, run by run :: sha256 c6c4138ee5c86423 :: prose twin: two significant figures, direction in words
\newcommand{\MlstCmpJcTwoResFonerOverFullPrMax}{up to 19\% more errors}  % experiments/FIGS/data/v2/index.json :: linear probe jc2_188 N=1000000: (1 - macro_auc_resonant) of FLAV_F1R over L188, run by run :: sha256 c6c4138ee5c86423 :: prose twin of the run farther from 1
\newcommand{\MlstCmpJcTwoResFzeroOverCoarse}{0.9069 and 0.9029}  % experiments/FIGS/data/v2/index.json :: linear probe jc2_188 N=1000000: (1 - macro_auc_resonant) of FLAV_F0 over R16_Q1, run by run :: sha256 c6c4138ee5c86423
\newcommand{\MlstCmpJcTwoResFzeroOverCoarsePr}{at least 9.3\% fewer errors}  % experiments/FIGS/data/v2/index.json :: linear probe jc2_188 N=1000000: (1 - macro_auc_resonant) of FLAV_F0 over R16_Q1, run by run :: sha256 c6c4138ee5c86423 :: prose twin: two significant figures, direction in words
\newcommand{\MlstCmpJcTwoResFzeroOverCoarsePrMax}{up to 9.7\% fewer errors}  % experiments/FIGS/data/v2/index.json :: linear probe jc2_188 N=1000000: (1 - macro_auc_resonant) of FLAV_F0 over R16_Q1, run by run :: sha256 c6c4138ee5c86423 :: prose twin of the run farther from 1
\newcommand{\MlstCmpJcTwoResFoneOverCoarse}{0.779 and 0.750}  % experiments/FIGS/data/v2/index.json :: linear probe jc2_188 N=1000000: (1 - macro_auc_resonant) of FLAV_F1 over R16_Q1, run by run :: sha256 c6c4138ee5c86423
\newcommand{\MlstCmpJcTwoResFoneOverCoarsePr}{at least 22\% fewer errors}  % experiments/FIGS/data/v2/index.json :: linear probe jc2_188 N=1000000: (1 - macro_auc_resonant) of FLAV_F1 over R16_Q1, run by run :: sha256 c6c4138ee5c86423 :: prose twin: two significant figures, direction in words
\newcommand{\MlstCmpJcTwoResFoneOverCoarsePrMax}{up to 25\% fewer errors}  % experiments/FIGS/data/v2/index.json :: linear probe jc2_188 N=1000000: (1 - macro_auc_resonant) of FLAV_F1 over R16_Q1, run by run :: sha256 c6c4138ee5c86423 :: prose twin of the run farther from 1
\newcommand{\MlstCmpJcTwoResFonerOverCoarse}{0.764 and 0.740}  % experiments/FIGS/data/v2/index.json :: linear probe jc2_188 N=1000000: (1 - macro_auc_resonant) of FLAV_F1R over R16_Q1, run by run :: sha256 c6c4138ee5c86423
\newcommand{\MlstCmpJcTwoResFonerOverCoarsePr}{at least 24\% fewer errors}  % experiments/FIGS/data/v2/index.json :: linear probe jc2_188 N=1000000: (1 - macro_auc_resonant) of FLAV_F1R over R16_Q1, run by run :: sha256 c6c4138ee5c86423 :: prose twin: two significant figures, direction in words
\newcommand{\MlstCmpJcTwoResFonerOverCoarsePrMax}{up to 26\% fewer errors}  % experiments/FIGS/data/v2/index.json :: linear probe jc2_188 N=1000000: (1 - macro_auc_resonant) of FLAV_F1R over R16_Q1, run by run :: sha256 c6c4138ee5c86423 :: prose twin of the run farther from 1
\newcommand{\MlstCmpJcTwoResFoneOverFzero}{0.859 and 0.831}  % experiments/FIGS/data/v2/index.json :: linear probe jc2_188 N=1000000: (1 - macro_auc_resonant) of FLAV_F1 over FLAV_F0, run by run :: sha256 c6c4138ee5c86423
\newcommand{\MlstCmpJcTwoResFoneOverFzeroPr}{at least 14\% fewer errors}  % experiments/FIGS/data/v2/index.json :: linear probe jc2_188 N=1000000: (1 - macro_auc_resonant) of FLAV_F1 over FLAV_F0, run by run :: sha256 c6c4138ee5c86423 :: prose twin: two significant figures, direction in words
\newcommand{\MlstCmpJcTwoResFoneOverFzeroPrMax}{up to 17\% fewer errors}  % experiments/FIGS/data/v2/index.json :: linear probe jc2_188 N=1000000: (1 - macro_auc_resonant) of FLAV_F1 over FLAV_F0, run by run :: sha256 c6c4138ee5c86423 :: prose twin of the run farther from 1
\newcommand{\MlstCmpJcTwoResFonerOverFzero}{0.843 and 0.820}  % experiments/FIGS/data/v2/index.json :: linear probe jc2_188 N=1000000: (1 - macro_auc_resonant) of FLAV_F1R over FLAV_F0, run by run :: sha256 c6c4138ee5c86423
\newcommand{\MlstCmpJcTwoResFonerOverFzeroPr}{at least 16\% fewer errors}  % experiments/FIGS/data/v2/index.json :: linear probe jc2_188 N=1000000: (1 - macro_auc_resonant) of FLAV_F1R over FLAV_F0, run by run :: sha256 c6c4138ee5c86423 :: prose twin: two significant figures, direction in words
\newcommand{\MlstCmpJcTwoResFonerOverFzeroPrMax}{up to 18\% fewer errors}  % experiments/FIGS/data/v2/index.json :: linear probe jc2_188 N=1000000: (1 - macro_auc_resonant) of FLAV_F1R over FLAV_F0, run by run :: sha256 c6c4138ee5c86423 :: prose twin of the run farther from 1
\newcommand{\MlstCmpJcTwoResFonerOverFone}{0.9806 and 0.9867}  % experiments/FIGS/data/v2/index.json :: linear probe jc2_188 N=1000000: (1 - macro_auc_resonant) of FLAV_F1R over FLAV_F1, run by run :: sha256 c6c4138ee5c86423
\newcommand{\MlstCmpJcTwoResFonerOverFonePr}{equal within 2\%}  % experiments/FIGS/data/v2/index.json :: linear probe jc2_188 N=1000000: (1 - macro_auc_resonant) of FLAV_F1R over FLAV_F1, run by run :: sha256 c6c4138ee5c86423 :: prose twin: two significant figures, direction in words
\newcommand{\MlstCmpJcTwoResFonerOverFonePrMax}{equal within 2\%}  % experiments/FIGS/data/v2/index.json :: linear probe jc2_188 N=1000000: (1 - macro_auc_resonant) of FLAV_F1R over FLAV_F1, run by run :: sha256 c6c4138ee5c86423 :: prose twin of the run farther from 1
\newcommand{\MlstCmpJcTwoResMidOverRes}{1.258\,\ensuremath{\pm}\,0.017}  % experiments/FIGS/data/v2/index.json :: linear probe jc2_188 N=1000000: (1 - macro_auc_resonant) of R42_Q1 over L162, run by run :: sha256 c6c4138ee5c86423
\newcommand{\MlstCmpJcTwoResMidOverResPr}{26\% more errors}  % experiments/FIGS/data/v2/index.json :: linear probe jc2_188 N=1000000: (1 - macro_auc_resonant) of R42_Q1 over L162, run by run :: sha256 c6c4138ee5c86423 :: prose twin: two significant figures, direction in words
\newcommand{\MlstCmpJcTwoResCoarseOverRes}{1.58\,\ensuremath{\pm}\,0.04}  % experiments/FIGS/data/v2/index.json :: linear probe jc2_188 N=1000000: (1 - macro_auc_resonant) of R16_Q1 over L162, run by run :: sha256 c6c4138ee5c86423
\newcommand{\MlstCmpJcTwoResCoarseOverResPr}{58\% more errors}  % experiments/FIGS/data/v2/index.json :: linear probe jc2_188 N=1000000: (1 - macro_auc_resonant) of R16_Q1 over L162, run by run :: sha256 c6c4138ee5c86423 :: prose twin: two significant figures, direction in words
\newcommand{\MlstCmpJcTwoResCoarseOverMid}{1.260\,\ensuremath{\pm}\,0.015}  % experiments/FIGS/data/v2/index.json :: linear probe jc2_188 N=1000000: (1 - macro_auc_resonant) of R16_Q1 over R42_Q1, run by run :: sha256 c6c4138ee5c86423
\newcommand{\MlstCmpJcTwoResCoarseOverMidPr}{26\% more errors}  % experiments/FIGS/data/v2/index.json :: linear probe jc2_188 N=1000000: (1 - macro_auc_resonant) of R16_Q1 over R42_Q1, run by run :: sha256 c6c4138ee5c86423 :: prose twin: two significant figures, direction in words
\newcommand{\MlstCmpJcTwoResRandOverFull}{1.098\,\ensuremath{\pm}\,0.008}  % experiments/FIGS/data/v2/index.json :: linear probe jc2_188 N=1000000: (1 - macro_auc_resonant) of each random partition over the L188 run of its seed, mean +- SD over the partitions :: sha256 c6c4138ee5c86423
\newcommand{\MlstCmpJcTwoResRandOverFullPr}{9.8\% more errors}  % experiments/FIGS/data/v2/index.json :: linear probe jc2_188 N=1000000: (1 - macro_auc_resonant) of each random partition over the L188 run of its seed, mean +- SD over the partitions :: sha256 c6c4138ee5c86423 :: prose twin: two significant figures, direction in words
\newcommand{\MlstCmpJcTwoResRandOverMid}{0.870\,\ensuremath{\pm}\,0.008}  % experiments/FIGS/data/v2/index.json :: linear probe jc2_188 N=1000000: (1 - macro_auc_resonant) of each random partition over the R42_Q1 run of its seed, mean +- SD over the partitions :: sha256 c6c4138ee5c86423
\newcommand{\MlstCmpJcTwoResRandOverMidPr}{13\% fewer errors}  % experiments/FIGS/data/v2/index.json :: linear probe jc2_188 N=1000000: (1 - macro_auc_resonant) of each random partition over the R42_Q1 run of its seed, mean +- SD over the partitions :: sha256 c6c4138ee5c86423 :: prose twin: two significant figures, direction in words
\newcommand{\MlstCmpJcTwoResRandOverCoarse}{0.693\,\ensuremath{\pm}\,0.011}  % experiments/FIGS/data/v2/index.json :: linear probe jc2_188 N=1000000: (1 - macro_auc_resonant) of each random partition over the R16_Q1 run of its seed, mean +- SD over the partitions :: sha256 c6c4138ee5c86423
\newcommand{\MlstCmpJcTwoResRandOverCoarsePr}{31\% fewer errors}  % experiments/FIGS/data/v2/index.json :: linear probe jc2_188 N=1000000: (1 - macro_auc_resonant) of each random partition over the R16_Q1 run of its seed, mean +- SD over the partitions :: sha256 c6c4138ee5c86423 :: prose twin: two significant figures, direction in words
\newcommand{\MlstCmpJcTwoQcdResOverFull}{1.087\,\ensuremath{\pm}\,0.006}  % experiments/FIGS/data/v2/index.json :: linear probe jc2_188 N=1000000: (1 - macro_auc_qcd) of L162 over L188, run by run :: sha256 c6c4138ee5c86423
\newcommand{\MlstCmpJcTwoQcdResOverFullPr}{8.7\% more errors}  % experiments/FIGS/data/v2/index.json :: linear probe jc2_188 N=1000000: (1 - macro_auc_qcd) of L162 over L188, run by run :: sha256 c6c4138ee5c86423 :: prose twin: two significant figures, direction in words
\newcommand{\MlstCmpJcTwoQcdMidOverFull}{1.131\,\ensuremath{\pm}\,0.007}  % experiments/FIGS/data/v2/index.json :: linear probe jc2_188 N=1000000: (1 - macro_auc_qcd) of R42_Q1 over L188, run by run :: sha256 c6c4138ee5c86423
\newcommand{\MlstCmpJcTwoQcdMidOverFullPr}{13\% more errors}  % experiments/FIGS/data/v2/index.json :: linear probe jc2_188 N=1000000: (1 - macro_auc_qcd) of R42_Q1 over L188, run by run :: sha256 c6c4138ee5c86423 :: prose twin: two significant figures, direction in words
\newcommand{\MlstCmpJcTwoQcdCoarseOverFull}{1.466\,\ensuremath{\pm}\,0.017}  % experiments/FIGS/data/v2/index.json :: linear probe jc2_188 N=1000000: (1 - macro_auc_qcd) of R16_Q1 over L188, run by run :: sha256 c6c4138ee5c86423
\newcommand{\MlstCmpJcTwoQcdCoarseOverFullPr}{47\% more errors}  % experiments/FIGS/data/v2/index.json :: linear probe jc2_188 N=1000000: (1 - macro_auc_qcd) of R16_Q1 over L188, run by run :: sha256 c6c4138ee5c86423 :: prose twin: two significant figures, direction in words
\newcommand{\MlstCmpJcTwoQcdResMassOverRes}{1.012 and 1.006}  % experiments/FIGS/data/v2/index.json :: linear probe jc2_188 N=1000000: (1 - macro_auc_qcd) of L162_MASS over L162, run by run :: sha256 c6c4138ee5c86423
\newcommand{\MlstCmpJcTwoQcdResMassOverResPr}{equal within 2\%}  % experiments/FIGS/data/v2/index.json :: linear probe jc2_188 N=1000000: (1 - macro_auc_qcd) of L162_MASS over L162, run by run :: sha256 c6c4138ee5c86423 :: prose twin: two significant figures, direction in words
\newcommand{\MlstCmpJcTwoQcdResMassOverResPrMax}{equal within 2\%}  % experiments/FIGS/data/v2/index.json :: linear probe jc2_188 N=1000000: (1 - macro_auc_qcd) of L162_MASS over L162, run by run :: sha256 c6c4138ee5c86423 :: prose twin of the run farther from 1
\newcommand{\MlstCmpJcTwoQcdCoarseMassOverCoarse}{0.991 and 1.014}  % experiments/FIGS/data/v2/index.json :: linear probe jc2_188 N=1000000: (1 - macro_auc_qcd) of R16_Q1_MASS over R16_Q1, run by run :: sha256 c6c4138ee5c86423
\newcommand{\MlstCmpJcTwoQcdCoarseMassOverCoarsePr}{equal within 2\%}  % experiments/FIGS/data/v2/index.json :: linear probe jc2_188 N=1000000: (1 - macro_auc_qcd) of R16_Q1_MASS over R16_Q1, run by run :: sha256 c6c4138ee5c86423 :: prose twin: two significant figures, direction in words
\newcommand{\MlstCmpJcTwoQcdCoarseMassOverCoarsePrMax}{equal within 2\%}  % experiments/FIGS/data/v2/index.json :: linear probe jc2_188 N=1000000: (1 - macro_auc_qcd) of R16_Q1_MASS over R16_Q1, run by run :: sha256 c6c4138ee5c86423 :: prose twin of the run farther from 1
\newcommand{\MlstCmpJcTwoQcdResMassOverFull}{1.107 and 1.092}  % experiments/FIGS/data/v2/index.json :: linear probe jc2_188 N=1000000: (1 - macro_auc_qcd) of L162_MASS over L188, run by run :: sha256 c6c4138ee5c86423
\newcommand{\MlstCmpJcTwoQcdResMassOverFullPr}{at least 9.2\% more errors}  % experiments/FIGS/data/v2/index.json :: linear probe jc2_188 N=1000000: (1 - macro_auc_qcd) of L162_MASS over L188, run by run :: sha256 c6c4138ee5c86423 :: prose twin: two significant figures, direction in words
\newcommand{\MlstCmpJcTwoQcdResMassOverFullPrMax}{up to 11\% more errors}  % experiments/FIGS/data/v2/index.json :: linear probe jc2_188 N=1000000: (1 - macro_auc_qcd) of L162_MASS over L188, run by run :: sha256 c6c4138ee5c86423 :: prose twin of the run farther from 1
\newcommand{\MlstCmpJcTwoQcdCoarseMassOverFull}{1.472 and 1.474}  % experiments/FIGS/data/v2/index.json :: linear probe jc2_188 N=1000000: (1 - macro_auc_qcd) of R16_Q1_MASS over L188, run by run :: sha256 c6c4138ee5c86423
\newcommand{\MlstCmpJcTwoQcdCoarseMassOverFullPr}{at least 47\% more errors}  % experiments/FIGS/data/v2/index.json :: linear probe jc2_188 N=1000000: (1 - macro_auc_qcd) of R16_Q1_MASS over L188, run by run :: sha256 c6c4138ee5c86423 :: prose twin: two significant figures, direction in words
\newcommand{\MlstCmpJcTwoQcdCoarseMassOverFullPrMax}{up to 47\% more errors}  % experiments/FIGS/data/v2/index.json :: linear probe jc2_188 N=1000000: (1 - macro_auc_qcd) of R16_Q1_MASS over L188, run by run :: sha256 c6c4138ee5c86423 :: prose twin of the run farther from 1
\newcommand{\MlstCmpJcTwoQcdFzeroOverFull}{1.466 and 1.497}  % experiments/FIGS/data/v2/index.json :: linear probe jc2_188 N=1000000: (1 - macro_auc_qcd) of FLAV_F0 over L188, run by run :: sha256 c6c4138ee5c86423
\newcommand{\MlstCmpJcTwoQcdFzeroOverFullPr}{at least 47\% more errors}  % experiments/FIGS/data/v2/index.json :: linear probe jc2_188 N=1000000: (1 - macro_auc_qcd) of FLAV_F0 over L188, run by run :: sha256 c6c4138ee5c86423 :: prose twin: two significant figures, direction in words
\newcommand{\MlstCmpJcTwoQcdFzeroOverFullPrMax}{up to 50\% more errors}  % experiments/FIGS/data/v2/index.json :: linear probe jc2_188 N=1000000: (1 - macro_auc_qcd) of FLAV_F0 over L188, run by run :: sha256 c6c4138ee5c86423 :: prose twin of the run farther from 1
\newcommand{\MlstCmpJcTwoQcdFoneOverFull}{1.241 and 1.235}  % experiments/FIGS/data/v2/index.json :: linear probe jc2_188 N=1000000: (1 - macro_auc_qcd) of FLAV_F1 over L188, run by run :: sha256 c6c4138ee5c86423
\newcommand{\MlstCmpJcTwoQcdFoneOverFullPr}{at least 23\% more errors}  % experiments/FIGS/data/v2/index.json :: linear probe jc2_188 N=1000000: (1 - macro_auc_qcd) of FLAV_F1 over L188, run by run :: sha256 c6c4138ee5c86423 :: prose twin: two significant figures, direction in words
\newcommand{\MlstCmpJcTwoQcdFoneOverFullPrMax}{up to 24\% more errors}  % experiments/FIGS/data/v2/index.json :: linear probe jc2_188 N=1000000: (1 - macro_auc_qcd) of FLAV_F1 over L188, run by run :: sha256 c6c4138ee5c86423 :: prose twin of the run farther from 1
\newcommand{\MlstCmpJcTwoQcdFonerOverFull}{1.260 and 1.263}  % experiments/FIGS/data/v2/index.json :: linear probe jc2_188 N=1000000: (1 - macro_auc_qcd) of FLAV_F1R over L188, run by run :: sha256 c6c4138ee5c86423
\newcommand{\MlstCmpJcTwoQcdFonerOverFullPr}{at least 26\% more errors}  % experiments/FIGS/data/v2/index.json :: linear probe jc2_188 N=1000000: (1 - macro_auc_qcd) of FLAV_F1R over L188, run by run :: sha256 c6c4138ee5c86423 :: prose twin: two significant figures, direction in words
\newcommand{\MlstCmpJcTwoQcdFonerOverFullPrMax}{up to 26\% more errors}  % experiments/FIGS/data/v2/index.json :: linear probe jc2_188 N=1000000: (1 - macro_auc_qcd) of FLAV_F1R over L188, run by run :: sha256 c6c4138ee5c86423 :: prose twin of the run farther from 1
\newcommand{\MlstCmpJcTwoQcdFzeroOverCoarse}{0.987 and 1.031}  % experiments/FIGS/data/v2/index.json :: linear probe jc2_188 N=1000000: (1 - macro_auc_qcd) of FLAV_F0 over R16_Q1, run by run :: sha256 c6c4138ee5c86423
\newcommand{\MlstCmpJcTwoQcdFzeroOverCoarsePr}{equal within 4\%}  % experiments/FIGS/data/v2/index.json :: linear probe jc2_188 N=1000000: (1 - macro_auc_qcd) of FLAV_F0 over R16_Q1, run by run :: sha256 c6c4138ee5c86423 :: prose twin: two significant figures, direction in words
\newcommand{\MlstCmpJcTwoQcdFzeroOverCoarsePrMax}{equal within 4\%}  % experiments/FIGS/data/v2/index.json :: linear probe jc2_188 N=1000000: (1 - macro_auc_qcd) of FLAV_F0 over R16_Q1, run by run :: sha256 c6c4138ee5c86423 :: prose twin of the run farther from 1
\newcommand{\MlstCmpJcTwoQcdFoneOverCoarse}{0.836 and 0.850}  % experiments/FIGS/data/v2/index.json :: linear probe jc2_188 N=1000000: (1 - macro_auc_qcd) of FLAV_F1 over R16_Q1, run by run :: sha256 c6c4138ee5c86423
\newcommand{\MlstCmpJcTwoQcdFoneOverCoarsePr}{at least 15\% fewer errors}  % experiments/FIGS/data/v2/index.json :: linear probe jc2_188 N=1000000: (1 - macro_auc_qcd) of FLAV_F1 over R16_Q1, run by run :: sha256 c6c4138ee5c86423 :: prose twin: two significant figures, direction in words
\newcommand{\MlstCmpJcTwoQcdFoneOverCoarsePrMax}{up to 16\% fewer errors}  % experiments/FIGS/data/v2/index.json :: linear probe jc2_188 N=1000000: (1 - macro_auc_qcd) of FLAV_F1 over R16_Q1, run by run :: sha256 c6c4138ee5c86423 :: prose twin of the run farther from 1
\newcommand{\MlstCmpJcTwoQcdFonerOverCoarse}{0.848 and 0.870}  % experiments/FIGS/data/v2/index.json :: linear probe jc2_188 N=1000000: (1 - macro_auc_qcd) of FLAV_F1R over R16_Q1, run by run :: sha256 c6c4138ee5c86423
\newcommand{\MlstCmpJcTwoQcdFonerOverCoarsePr}{at least 13\% fewer errors}  % experiments/FIGS/data/v2/index.json :: linear probe jc2_188 N=1000000: (1 - macro_auc_qcd) of FLAV_F1R over R16_Q1, run by run :: sha256 c6c4138ee5c86423 :: prose twin: two significant figures, direction in words
\newcommand{\MlstCmpJcTwoQcdFonerOverCoarsePrMax}{up to 15\% fewer errors}  % experiments/FIGS/data/v2/index.json :: linear probe jc2_188 N=1000000: (1 - macro_auc_qcd) of FLAV_F1R over R16_Q1, run by run :: sha256 c6c4138ee5c86423 :: prose twin of the run farther from 1
\newcommand{\MlstCmpJcTwoQcdFoneOverFzero}{0.847 and 0.825}  % experiments/FIGS/data/v2/index.json :: linear probe jc2_188 N=1000000: (1 - macro_auc_qcd) of FLAV_F1 over FLAV_F0, run by run :: sha256 c6c4138ee5c86423
\newcommand{\MlstCmpJcTwoQcdFoneOverFzeroPr}{at least 15\% fewer errors}  % experiments/FIGS/data/v2/index.json :: linear probe jc2_188 N=1000000: (1 - macro_auc_qcd) of FLAV_F1 over FLAV_F0, run by run :: sha256 c6c4138ee5c86423 :: prose twin: two significant figures, direction in words
\newcommand{\MlstCmpJcTwoQcdFoneOverFzeroPrMax}{up to 18\% fewer errors}  % experiments/FIGS/data/v2/index.json :: linear probe jc2_188 N=1000000: (1 - macro_auc_qcd) of FLAV_F1 over FLAV_F0, run by run :: sha256 c6c4138ee5c86423 :: prose twin of the run farther from 1
\newcommand{\MlstCmpJcTwoQcdFonerOverFzero}{0.859 and 0.844}  % experiments/FIGS/data/v2/index.json :: linear probe jc2_188 N=1000000: (1 - macro_auc_qcd) of FLAV_F1R over FLAV_F0, run by run :: sha256 c6c4138ee5c86423
\newcommand{\MlstCmpJcTwoQcdFonerOverFzeroPr}{at least 14\% fewer errors}  % experiments/FIGS/data/v2/index.json :: linear probe jc2_188 N=1000000: (1 - macro_auc_qcd) of FLAV_F1R over FLAV_F0, run by run :: sha256 c6c4138ee5c86423 :: prose twin: two significant figures, direction in words
\newcommand{\MlstCmpJcTwoQcdFonerOverFzeroPrMax}{up to 16\% fewer errors}  % experiments/FIGS/data/v2/index.json :: linear probe jc2_188 N=1000000: (1 - macro_auc_qcd) of FLAV_F1R over FLAV_F0, run by run :: sha256 c6c4138ee5c86423 :: prose twin of the run farther from 1
\newcommand{\MlstCmpJcTwoQcdFonerOverFone}{1.015 and 1.023}  % experiments/FIGS/data/v2/index.json :: linear probe jc2_188 N=1000000: (1 - macro_auc_qcd) of FLAV_F1R over FLAV_F1, run by run :: sha256 c6c4138ee5c86423
\newcommand{\MlstCmpJcTwoQcdFonerOverFonePr}{equal within 3\%}  % experiments/FIGS/data/v2/index.json :: linear probe jc2_188 N=1000000: (1 - macro_auc_qcd) of FLAV_F1R over FLAV_F1, run by run :: sha256 c6c4138ee5c86423 :: prose twin: two significant figures, direction in words
\newcommand{\MlstCmpJcTwoQcdFonerOverFonePrMax}{equal within 3\%}  % experiments/FIGS/data/v2/index.json :: linear probe jc2_188 N=1000000: (1 - macro_auc_qcd) of FLAV_F1R over FLAV_F1, run by run :: sha256 c6c4138ee5c86423 :: prose twin of the run farther from 1
\newcommand{\MlstCmpJcTwoQcdMidOverRes}{1.040\,\ensuremath{\pm}\,0.011}  % experiments/FIGS/data/v2/index.json :: linear probe jc2_188 N=1000000: (1 - macro_auc_qcd) of R42_Q1 over L162, run by run :: sha256 c6c4138ee5c86423
\newcommand{\MlstCmpJcTwoQcdMidOverResPr}{4.0\% more errors}  % experiments/FIGS/data/v2/index.json :: linear probe jc2_188 N=1000000: (1 - macro_auc_qcd) of R42_Q1 over L162, run by run :: sha256 c6c4138ee5c86423 :: prose twin: two significant figures, direction in words
\newcommand{\MlstCmpJcTwoQcdCoarseOverRes}{1.348\,\ensuremath{\pm}\,0.010}  % experiments/FIGS/data/v2/index.json :: linear probe jc2_188 N=1000000: (1 - macro_auc_qcd) of R16_Q1 over L162, run by run :: sha256 c6c4138ee5c86423
\newcommand{\MlstCmpJcTwoQcdCoarseOverResPr}{35\% more errors}  % experiments/FIGS/data/v2/index.json :: linear probe jc2_188 N=1000000: (1 - macro_auc_qcd) of R16_Q1 over L162, run by run :: sha256 c6c4138ee5c86423 :: prose twin: two significant figures, direction in words
\newcommand{\MlstCmpJcTwoQcdCoarseOverMid}{1.296\,\ensuremath{\pm}\,0.016}  % experiments/FIGS/data/v2/index.json :: linear probe jc2_188 N=1000000: (1 - macro_auc_qcd) of R16_Q1 over R42_Q1, run by run :: sha256 c6c4138ee5c86423
\newcommand{\MlstCmpJcTwoQcdCoarseOverMidPr}{30\% more errors}  % experiments/FIGS/data/v2/index.json :: linear probe jc2_188 N=1000000: (1 - macro_auc_qcd) of R16_Q1 over R42_Q1, run by run :: sha256 c6c4138ee5c86423 :: prose twin: two significant figures, direction in words
\newcommand{\MlstCmpJcTwoQcdRandOverFull}{1.198\,\ensuremath{\pm}\,0.011}  % experiments/FIGS/data/v2/index.json :: linear probe jc2_188 N=1000000: (1 - macro_auc_qcd) of each random partition over the L188 run of its seed, mean +- SD over the partitions :: sha256 c6c4138ee5c86423
\newcommand{\MlstCmpJcTwoQcdRandOverFullPr}{20\% more errors}  % experiments/FIGS/data/v2/index.json :: linear probe jc2_188 N=1000000: (1 - macro_auc_qcd) of each random partition over the L188 run of its seed, mean +- SD over the partitions :: sha256 c6c4138ee5c86423 :: prose twin: two significant figures, direction in words
\newcommand{\MlstCmpJcTwoQcdRandOverMid}{1.064\,\ensuremath{\pm}\,0.011}  % experiments/FIGS/data/v2/index.json :: linear probe jc2_188 N=1000000: (1 - macro_auc_qcd) of each random partition over the R42_Q1 run of its seed, mean +- SD over the partitions :: sha256 c6c4138ee5c86423
\newcommand{\MlstCmpJcTwoQcdRandOverMidPr}{6.4\% more errors}  % experiments/FIGS/data/v2/index.json :: linear probe jc2_188 N=1000000: (1 - macro_auc_qcd) of each random partition over the R42_Q1 run of its seed, mean +- SD over the partitions :: sha256 c6c4138ee5c86423 :: prose twin: two significant figures, direction in words
\newcommand{\MlstCmpJcTwoQcdRandOverCoarse}{0.818\,\ensuremath{\pm}\,0.013}  % experiments/FIGS/data/v2/index.json :: linear probe jc2_188 N=1000000: (1 - macro_auc_qcd) of each random partition over the R16_Q1 run of its seed, mean +- SD over the partitions :: sha256 c6c4138ee5c86423
\newcommand{\MlstCmpJcTwoQcdRandOverCoarsePr}{18\% fewer errors}  % experiments/FIGS/data/v2/index.json :: linear probe jc2_188 N=1000000: (1 - macro_auc_qcd) of each random partition over the R16_Q1 run of its seed, mean +- SD over the partitions :: sha256 c6c4138ee5c86423 :: prose twin: two significant figures, direction in words
\newcommand{\MlstCmpBcResOverFull}{0.998\,\ensuremath{\pm}\,0.007}  % experiments/FIGS/data/v2/index.json :: binary probe bc_vs_bq_cs: (1 - AUC) of L162 over L188, run by run :: sha256 c6c4138ee5c86423
\newcommand{\MlstCmpBcResOverFullPr}{equal within 1\%}  % experiments/FIGS/data/v2/index.json :: binary probe bc_vs_bq_cs: (1 - AUC) of L162 over L188, run by run :: sha256 c6c4138ee5c86423 :: prose twin: two significant figures, direction in words
\newcommand{\MlstCmpBcMidOverFull}{1.151\,\ensuremath{\pm}\,0.029}  % experiments/FIGS/data/v2/index.json :: binary probe bc_vs_bq_cs: (1 - AUC) of R42_Q1 over L188, run by run :: sha256 c6c4138ee5c86423
\newcommand{\MlstCmpBcMidOverFullPr}{15\% more errors}  % experiments/FIGS/data/v2/index.json :: binary probe bc_vs_bq_cs: (1 - AUC) of R42_Q1 over L188, run by run :: sha256 c6c4138ee5c86423 :: prose twin: two significant figures, direction in words
\newcommand{\MlstCmpBcCoarseOverFull}{1.95\,\ensuremath{\pm}\,0.06}  % experiments/FIGS/data/v2/index.json :: binary probe bc_vs_bq_cs: (1 - AUC) of R16_Q1 over L188, run by run :: sha256 c6c4138ee5c86423
\newcommand{\MlstCmpBcCoarseOverFullPr}{95\% more errors}  % experiments/FIGS/data/v2/index.json :: binary probe bc_vs_bq_cs: (1 - AUC) of R16_Q1 over L188, run by run :: sha256 c6c4138ee5c86423 :: prose twin: two significant figures, direction in words
\newcommand{\MlstCmpBcFoneOverFzero}{0.759 and 0.698}  % experiments/FIGS/data/v2/index.json :: binary probe bc_vs_bq_cs: (1 - AUC) of FLAV_F1 over FLAV_F0, run by run :: sha256 c6c4138ee5c86423
\newcommand{\MlstCmpBcFoneOverFzeroPr}{at least 24\% fewer errors}  % experiments/FIGS/data/v2/index.json :: binary probe bc_vs_bq_cs: (1 - AUC) of FLAV_F1 over FLAV_F0, run by run :: sha256 c6c4138ee5c86423 :: prose twin: two significant figures, direction in words
\newcommand{\MlstCmpBcFoneOverFzeroPrMax}{up to 30\% fewer errors}  % experiments/FIGS/data/v2/index.json :: binary probe bc_vs_bq_cs: (1 - AUC) of FLAV_F1 over FLAV_F0, run by run :: sha256 c6c4138ee5c86423 :: prose twin of the run farther from 1
\newcommand{\MlstCmpBcFonerOverFone}{1.008 and 0.990}  % experiments/FIGS/data/v2/index.json :: binary probe bc_vs_bq_cs: (1 - AUC) of FLAV_F1R over FLAV_F1, run by run :: sha256 c6c4138ee5c86423
\newcommand{\MlstCmpBcFonerOverFonePr}{equal within 1\%}  % experiments/FIGS/data/v2/index.json :: binary probe bc_vs_bq_cs: (1 - AUC) of FLAV_F1R over FLAV_F1, run by run :: sha256 c6c4138ee5c86423 :: prose twin: two significant figures, direction in words
\newcommand{\MlstCmpBcFonerOverFonePrMax}{equal within 1\%}  % experiments/FIGS/data/v2/index.json :: binary probe bc_vs_bq_cs: (1 - AUC) of FLAV_F1R over FLAV_F1, run by run :: sha256 c6c4138ee5c86423 :: prose twin of the run farther from 1
\newcommand{\MlstCmpBcFonerOverFzero}{0.765 and 0.691}  % experiments/FIGS/data/v2/index.json :: binary probe bc_vs_bq_cs: (1 - AUC) of FLAV_F1R over FLAV_F0, run by run :: sha256 c6c4138ee5c86423
\newcommand{\MlstCmpBcFonerOverFzeroPr}{at least 23\% fewer errors}  % experiments/FIGS/data/v2/index.json :: binary probe bc_vs_bq_cs: (1 - AUC) of FLAV_F1R over FLAV_F0, run by run :: sha256 c6c4138ee5c86423 :: prose twin: two significant figures, direction in words
\newcommand{\MlstCmpBcFonerOverFzeroPrMax}{up to 31\% fewer errors}  % experiments/FIGS/data/v2/index.json :: binary probe bc_vs_bq_cs: (1 - AUC) of FLAV_F1R over FLAV_F0, run by run :: sha256 c6c4138ee5c86423 :: prose twin of the run farther from 1
\newcommand{\MlstCmpBcFzeroOverCoarse}{0.99 and 1.13}  % experiments/FIGS/data/v2/index.json :: binary probe bc_vs_bq_cs: (1 - AUC) of FLAV_F0 over R16_Q1, run by run :: sha256 c6c4138ee5c86423
\newcommand{\MlstCmpBcFzeroOverCoarsePr}{equal within 20\%}  % experiments/FIGS/data/v2/index.json :: binary probe bc_vs_bq_cs: (1 - AUC) of FLAV_F0 over R16_Q1, run by run :: sha256 c6c4138ee5c86423 :: prose twin: two significant figures, direction in words
\newcommand{\MlstCmpBcFzeroOverCoarsePrMax}{equal within 20\%}  % experiments/FIGS/data/v2/index.json :: binary probe bc_vs_bq_cs: (1 - AUC) of FLAV_F0 over R16_Q1, run by run :: sha256 c6c4138ee5c86423 :: prose twin of the run farther from 1
\newcommand{\MlstCmpBcFoneOverCoarse}{0.752 and 0.786}  % experiments/FIGS/data/v2/index.json :: binary probe bc_vs_bq_cs: (1 - AUC) of FLAV_F1 over R16_Q1, run by run :: sha256 c6c4138ee5c86423
\newcommand{\MlstCmpBcFoneOverCoarsePr}{at least 21\% fewer errors}  % experiments/FIGS/data/v2/index.json :: binary probe bc_vs_bq_cs: (1 - AUC) of FLAV_F1 over R16_Q1, run by run :: sha256 c6c4138ee5c86423 :: prose twin: two significant figures, direction in words
\newcommand{\MlstCmpBcFoneOverCoarsePrMax}{up to 25\% fewer errors}  % experiments/FIGS/data/v2/index.json :: binary probe bc_vs_bq_cs: (1 - AUC) of FLAV_F1 over R16_Q1, run by run :: sha256 c6c4138ee5c86423 :: prose twin of the run farther from 1
\newcommand{\MlstCmpBcFonerOverCoarse}{0.758 and 0.778}  % experiments/FIGS/data/v2/index.json :: binary probe bc_vs_bq_cs: (1 - AUC) of FLAV_F1R over R16_Q1, run by run :: sha256 c6c4138ee5c86423
\newcommand{\MlstCmpBcFonerOverCoarsePr}{at least 22\% fewer errors}  % experiments/FIGS/data/v2/index.json :: binary probe bc_vs_bq_cs: (1 - AUC) of FLAV_F1R over R16_Q1, run by run :: sha256 c6c4138ee5c86423 :: prose twin: two significant figures, direction in words
\newcommand{\MlstCmpBcFonerOverCoarsePrMax}{up to 24\% fewer errors}  % experiments/FIGS/data/v2/index.json :: binary probe bc_vs_bq_cs: (1 - AUC) of FLAV_F1R over R16_Q1, run by run :: sha256 c6c4138ee5c86423 :: prose twin of the run farther from 1
\newcommand{\MlstCmpBcFzeroOverFull}{2.00 and 2.19}  % experiments/FIGS/data/v2/index.json :: binary probe bc_vs_bq_cs: (1 - AUC) of FLAV_F0 over L188, run by run :: sha256 c6c4138ee5c86423
\newcommand{\MlstCmpBcFzeroOverFullPr}{at least 2.0 times the errors}  % experiments/FIGS/data/v2/index.json :: binary probe bc_vs_bq_cs: (1 - AUC) of FLAV_F0 over L188, run by run :: sha256 c6c4138ee5c86423 :: prose twin: two significant figures, direction in words
\newcommand{\MlstCmpBcFzeroOverFullPrMax}{up to 2.2 times the errors}  % experiments/FIGS/data/v2/index.json :: binary probe bc_vs_bq_cs: (1 - AUC) of FLAV_F0 over L188, run by run :: sha256 c6c4138ee5c86423 :: prose twin of the run farther from 1
\newcommand{\MlstCmpBcFoneOverFull}{1.518 and 1.530}  % experiments/FIGS/data/v2/index.json :: binary probe bc_vs_bq_cs: (1 - AUC) of FLAV_F1 over L188, run by run :: sha256 c6c4138ee5c86423
\newcommand{\MlstCmpBcFoneOverFullPr}{at least 52\% more errors}  % experiments/FIGS/data/v2/index.json :: binary probe bc_vs_bq_cs: (1 - AUC) of FLAV_F1 over L188, run by run :: sha256 c6c4138ee5c86423 :: prose twin: two significant figures, direction in words
\newcommand{\MlstCmpBcFoneOverFullPrMax}{up to 53\% more errors}  % experiments/FIGS/data/v2/index.json :: binary probe bc_vs_bq_cs: (1 - AUC) of FLAV_F1 over L188, run by run :: sha256 c6c4138ee5c86423 :: prose twin of the run farther from 1
\newcommand{\MlstCmpBcFonerOverFull}{1.530 and 1.515}  % experiments/FIGS/data/v2/index.json :: binary probe bc_vs_bq_cs: (1 - AUC) of FLAV_F1R over L188, run by run :: sha256 c6c4138ee5c86423
\newcommand{\MlstCmpBcFonerOverFullPr}{at least 51\% more errors}  % experiments/FIGS/data/v2/index.json :: binary probe bc_vs_bq_cs: (1 - AUC) of FLAV_F1R over L188, run by run :: sha256 c6c4138ee5c86423 :: prose twin: two significant figures, direction in words
\newcommand{\MlstCmpBcFonerOverFullPrMax}{up to 53\% more errors}  % experiments/FIGS/data/v2/index.json :: binary probe bc_vs_bq_cs: (1 - AUC) of FLAV_F1R over L188, run by run :: sha256 c6c4138ee5c86423 :: prose twin of the run farther from 1
\newcommand{\MlstCmpBcRandOverFull}{1.15\,\ensuremath{\pm}\,0.07}  % experiments/FIGS/data/v2/index.json :: binary probe bc_vs_bq_cs: (1 - AUC) of each random partition over the 188-class run of its seed :: sha256 c6c4138ee5c86423
\newcommand{\MlstCmpBcRandOverFullPr}{15\% more errors}  % experiments/FIGS/data/v2/index.json :: binary probe bc_vs_bq_cs: (1 - AUC) of each random partition over the 188-class run of its seed :: sha256 c6c4138ee5c86423 :: prose twin: two significant figures, direction in words
\newcommand{\MlstPairRandMergedBcN}{2}  % experiments/FIGS/data/v2/index.json :: random partitions that merge bc_vs_bq_cs (count) :: sha256 c6c4138ee5c86423
\newcommand{\MlstPairRandApartBcN}{2}  % experiments/FIGS/data/v2/index.json :: random partitions that keep apart bc_vs_bq_cs (count) :: sha256 c6c4138ee5c86423
\newcommand{\MlstPairRandMergedOverApartBc}{1.12}  % experiments/FIGS/data/v2/index.json :: binary probe bc_vs_bq_cs: mean (1 - AUC) of the partitions that merge it over that of those that keep it apart (one run each) :: sha256 c6c4138ee5c86423
\newcommand{\MlstPairRandMergedOverApartBcPr}{12\% more errors}  % experiments/FIGS/data/v2/index.json :: binary probe bc_vs_bq_cs: mean (1 - AUC) of the partitions that merge it over that of those that keep it apart (one run each) :: sha256 c6c4138ee5c86423 :: prose twin: two significant figures, direction in words
\newcommand{\MlstCmpBbResOverFull}{1.015\,\ensuremath{\pm}\,0.027}  % experiments/FIGS/data/v2/index.json :: binary probe bb_vs_cc: (1 - AUC) of L162 over L188, run by run :: sha256 c6c4138ee5c86423
\newcommand{\MlstCmpBbResOverFullPr}{equal within 3\%}  % experiments/FIGS/data/v2/index.json :: binary probe bb_vs_cc: (1 - AUC) of L162 over L188, run by run :: sha256 c6c4138ee5c86423 :: prose twin: two significant figures, direction in words
\newcommand{\MlstCmpBbMidOverFull}{1.102\,\ensuremath{\pm}\,0.013}  % experiments/FIGS/data/v2/index.json :: binary probe bb_vs_cc: (1 - AUC) of R42_Q1 over L188, run by run :: sha256 c6c4138ee5c86423
\newcommand{\MlstCmpBbMidOverFullPr}{10\% more errors}  % experiments/FIGS/data/v2/index.json :: binary probe bb_vs_cc: (1 - AUC) of R42_Q1 over L188, run by run :: sha256 c6c4138ee5c86423 :: prose twin: two significant figures, direction in words
\newcommand{\MlstCmpBbCoarseOverFull}{4.22\,\ensuremath{\pm}\,0.29}  % experiments/FIGS/data/v2/index.json :: binary probe bb_vs_cc: (1 - AUC) of R16_Q1 over L188, run by run :: sha256 c6c4138ee5c86423
\newcommand{\MlstCmpBbCoarseOverFullPr}{4.2 times the errors}  % experiments/FIGS/data/v2/index.json :: binary probe bb_vs_cc: (1 - AUC) of R16_Q1 over L188, run by run :: sha256 c6c4138ee5c86423 :: prose twin: two significant figures, direction in words
\newcommand{\MlstCmpBbFoneOverFzero}{0.31 and 0.13}  % experiments/FIGS/data/v2/index.json :: binary probe bb_vs_cc: (1 - AUC) of FLAV_F1 over FLAV_F0, run by run :: sha256 c6c4138ee5c86423
\newcommand{\MlstCmpBbFoneOverFzeroPr}{at most 0.31 times the errors}  % experiments/FIGS/data/v2/index.json :: binary probe bb_vs_cc: (1 - AUC) of FLAV_F1 over FLAV_F0, run by run :: sha256 c6c4138ee5c86423 :: prose twin: two significant figures, direction in words
\newcommand{\MlstCmpBbFoneOverFzeroPrMax}{down to 0.13 times the errors}  % experiments/FIGS/data/v2/index.json :: binary probe bb_vs_cc: (1 - AUC) of FLAV_F1 over FLAV_F0, run by run :: sha256 c6c4138ee5c86423 :: prose twin of the run farther from 1
\newcommand{\MlstCmpBbFonerOverFone}{1.40 and 1.86}  % experiments/FIGS/data/v2/index.json :: binary probe bb_vs_cc: (1 - AUC) of FLAV_F1R over FLAV_F1, run by run :: sha256 c6c4138ee5c86423
\newcommand{\MlstCmpBbFonerOverFonePr}{at least 40\% more errors}  % experiments/FIGS/data/v2/index.json :: binary probe bb_vs_cc: (1 - AUC) of FLAV_F1R over FLAV_F1, run by run :: sha256 c6c4138ee5c86423 :: prose twin: two significant figures, direction in words
\newcommand{\MlstCmpBbFonerOverFonePrMax}{up to 86\% more errors}  % experiments/FIGS/data/v2/index.json :: binary probe bb_vs_cc: (1 - AUC) of FLAV_F1R over FLAV_F1, run by run :: sha256 c6c4138ee5c86423 :: prose twin of the run farther from 1
\newcommand{\MlstCmpBbFonerOverFzero}{0.44 and 0.24}  % experiments/FIGS/data/v2/index.json :: binary probe bb_vs_cc: (1 - AUC) of FLAV_F1R over FLAV_F0, run by run :: sha256 c6c4138ee5c86423
\newcommand{\MlstCmpBbFonerOverFzeroPr}{at most 0.44 times the errors}  % experiments/FIGS/data/v2/index.json :: binary probe bb_vs_cc: (1 - AUC) of FLAV_F1R over FLAV_F0, run by run :: sha256 c6c4138ee5c86423 :: prose twin: two significant figures, direction in words
\newcommand{\MlstCmpBbFonerOverFzeroPrMax}{down to 0.24 times the errors}  % experiments/FIGS/data/v2/index.json :: binary probe bb_vs_cc: (1 - AUC) of FLAV_F1R over FLAV_F0, run by run :: sha256 c6c4138ee5c86423 :: prose twin of the run farther from 1
\newcommand{\MlstCmpBbFzeroOverCoarse}{1.3 and 3.2}  % experiments/FIGS/data/v2/index.json :: binary probe bb_vs_cc: (1 - AUC) of FLAV_F0 over R16_Q1, run by run :: sha256 c6c4138ee5c86423
\newcommand{\MlstCmpBbFzeroOverCoarsePr}{at least 30\% more errors}  % experiments/FIGS/data/v2/index.json :: binary probe bb_vs_cc: (1 - AUC) of FLAV_F0 over R16_Q1, run by run :: sha256 c6c4138ee5c86423 :: prose twin: two significant figures, direction in words
\newcommand{\MlstCmpBbFzeroOverCoarsePrMax}{up to 3.2 times the errors}  % experiments/FIGS/data/v2/index.json :: binary probe bb_vs_cc: (1 - AUC) of FLAV_F0 over R16_Q1, run by run :: sha256 c6c4138ee5c86423 :: prose twin of the run farther from 1
\newcommand{\MlstCmpBbFoneOverCoarse}{0.4077 and 0.4126}  % experiments/FIGS/data/v2/index.json :: binary probe bb_vs_cc: (1 - AUC) of FLAV_F1 over R16_Q1, run by run :: sha256 c6c4138ee5c86423
\newcommand{\MlstCmpBbFoneOverCoarsePr}{at most 0.41 times the errors}  % experiments/FIGS/data/v2/index.json :: binary probe bb_vs_cc: (1 - AUC) of FLAV_F1 over R16_Q1, run by run :: sha256 c6c4138ee5c86423 :: prose twin: two significant figures, direction in words
\newcommand{\MlstCmpBbFoneOverCoarsePrMax}{down to 0.41 times the errors}  % experiments/FIGS/data/v2/index.json :: binary probe bb_vs_cc: (1 - AUC) of FLAV_F1 over R16_Q1, run by run :: sha256 c6c4138ee5c86423 :: prose twin of the run farther from 1
\newcommand{\MlstCmpBbFonerOverCoarse}{0.57 and 0.77}  % experiments/FIGS/data/v2/index.json :: binary probe bb_vs_cc: (1 - AUC) of FLAV_F1R over R16_Q1, run by run :: sha256 c6c4138ee5c86423
\newcommand{\MlstCmpBbFonerOverCoarsePr}{at least 23\% fewer errors}  % experiments/FIGS/data/v2/index.json :: binary probe bb_vs_cc: (1 - AUC) of FLAV_F1R over R16_Q1, run by run :: sha256 c6c4138ee5c86423 :: prose twin: two significant figures, direction in words
\newcommand{\MlstCmpBbFonerOverCoarsePrMax}{up to 43\% fewer errors}  % experiments/FIGS/data/v2/index.json :: binary probe bb_vs_cc: (1 - AUC) of FLAV_F1R over R16_Q1, run by run :: sha256 c6c4138ee5c86423 :: prose twin of the run farther from 1
\newcommand{\MlstCmpBbFzeroOverFull}{5.9 and 13.4}  % experiments/FIGS/data/v2/index.json :: binary probe bb_vs_cc: (1 - AUC) of FLAV_F0 over L188, run by run :: sha256 c6c4138ee5c86423
\newcommand{\MlstCmpBbFzeroOverFullPr}{at least 5.9 times the errors}  % experiments/FIGS/data/v2/index.json :: binary probe bb_vs_cc: (1 - AUC) of FLAV_F0 over L188, run by run :: sha256 c6c4138ee5c86423 :: prose twin: two significant figures, direction in words
\newcommand{\MlstCmpBbFzeroOverFullPrMax}{up to 13 times the errors}  % experiments/FIGS/data/v2/index.json :: binary probe bb_vs_cc: (1 - AUC) of FLAV_F0 over L188, run by run :: sha256 c6c4138ee5c86423 :: prose twin of the run farther from 1
\newcommand{\MlstCmpBbFoneOverFull}{1.837 and 1.748}  % experiments/FIGS/data/v2/index.json :: binary probe bb_vs_cc: (1 - AUC) of FLAV_F1 over L188, run by run :: sha256 c6c4138ee5c86423
\newcommand{\MlstCmpBbFoneOverFullPr}{at least 75\% more errors}  % experiments/FIGS/data/v2/index.json :: binary probe bb_vs_cc: (1 - AUC) of FLAV_F1 over L188, run by run :: sha256 c6c4138ee5c86423 :: prose twin: two significant figures, direction in words
\newcommand{\MlstCmpBbFoneOverFullPrMax}{up to 84\% more errors}  % experiments/FIGS/data/v2/index.json :: binary probe bb_vs_cc: (1 - AUC) of FLAV_F1 over L188, run by run :: sha256 c6c4138ee5c86423 :: prose twin of the run farther from 1
\newcommand{\MlstCmpBbFonerOverFull}{2.58 and 3.25}  % experiments/FIGS/data/v2/index.json :: binary probe bb_vs_cc: (1 - AUC) of FLAV_F1R over L188, run by run :: sha256 c6c4138ee5c86423
\newcommand{\MlstCmpBbFonerOverFullPr}{at least 2.6 times the errors}  % experiments/FIGS/data/v2/index.json :: binary probe bb_vs_cc: (1 - AUC) of FLAV_F1R over L188, run by run :: sha256 c6c4138ee5c86423 :: prose twin: two significant figures, direction in words
\newcommand{\MlstCmpBbFonerOverFullPrMax}{up to 3.2 times the errors}  % experiments/FIGS/data/v2/index.json :: binary probe bb_vs_cc: (1 - AUC) of FLAV_F1R over L188, run by run :: sha256 c6c4138ee5c86423 :: prose twin of the run farther from 1
\newcommand{\MlstCmpBbRandOverFull}{1.53\,\ensuremath{\pm}\,0.11}  % experiments/FIGS/data/v2/index.json :: binary probe bb_vs_cc: (1 - AUC) of each random partition over the 188-class run of its seed :: sha256 c6c4138ee5c86423
\newcommand{\MlstCmpBbRandOverFullPr}{53\% more errors}  % experiments/FIGS/data/v2/index.json :: binary probe bb_vs_cc: (1 - AUC) of each random partition over the 188-class run of its seed :: sha256 c6c4138ee5c86423 :: prose twin: two significant figures, direction in words
\newcommand{\MlstPairRandMergedBbN}{3}  % experiments/FIGS/data/v2/index.json :: random partitions that merge bb_vs_cc (count) :: sha256 c6c4138ee5c86423
\newcommand{\MlstPairRandApartBbN}{2}  % experiments/FIGS/data/v2/index.json :: random partitions that keep apart bb_vs_cc (count) :: sha256 c6c4138ee5c86423
\newcommand{\MlstPairRandMergedOverApartBb}{1.07}  % experiments/FIGS/data/v2/index.json :: binary probe bb_vs_cc: mean (1 - AUC) of the partitions that merge it over that of those that keep it apart (one run each) :: sha256 c6c4138ee5c86423
\newcommand{\MlstPairRandMergedOverApartBbPr}{7.3\% more errors}  % experiments/FIGS/data/v2/index.json :: binary probe bb_vs_cc: mean (1 - AUC) of the partitions that merge it over that of those that keep it apart (one run each) :: sha256 c6c4138ee5c86423 :: prose twin: two significant figures, direction in words
\newcommand{\MlstCmpBbqqResOverFull}{1.005\,\ensuremath{\pm}\,0.032}  % experiments/FIGS/data/v2/index.json :: binary probe bbqq_vs_ccqq: (1 - AUC) of L162 over L188, run by run :: sha256 c6c4138ee5c86423
\newcommand{\MlstCmpBbqqResOverFullPr}{equal within 4\%}  % experiments/FIGS/data/v2/index.json :: binary probe bbqq_vs_ccqq: (1 - AUC) of L162 over L188, run by run :: sha256 c6c4138ee5c86423 :: prose twin: two significant figures, direction in words
\newcommand{\MlstCmpBbqqMidOverFull}{1.117\,\ensuremath{\pm}\,0.008}  % experiments/FIGS/data/v2/index.json :: binary probe bbqq_vs_ccqq: (1 - AUC) of R42_Q1 over L188, run by run :: sha256 c6c4138ee5c86423
\newcommand{\MlstCmpBbqqMidOverFullPr}{12\% more errors}  % experiments/FIGS/data/v2/index.json :: binary probe bbqq_vs_ccqq: (1 - AUC) of R42_Q1 over L188, run by run :: sha256 c6c4138ee5c86423 :: prose twin: two significant figures, direction in words
\newcommand{\MlstCmpBbqqCoarseOverFull}{4.7\,\ensuremath{\pm}\,0.5}  % experiments/FIGS/data/v2/index.json :: binary probe bbqq_vs_ccqq: (1 - AUC) of R16_Q1 over L188, run by run :: sha256 c6c4138ee5c86423
\newcommand{\MlstCmpBbqqCoarseOverFullPr}{4.7 times the errors}  % experiments/FIGS/data/v2/index.json :: binary probe bbqq_vs_ccqq: (1 - AUC) of R16_Q1 over L188, run by run :: sha256 c6c4138ee5c86423 :: prose twin: two significant figures, direction in words
\newcommand{\MlstCmpBbqqFoneOverFzero}{0.236 and 0.154}  % experiments/FIGS/data/v2/index.json :: binary probe bbqq_vs_ccqq: (1 - AUC) of FLAV_F1 over FLAV_F0, run by run :: sha256 c6c4138ee5c86423
\newcommand{\MlstCmpBbqqFoneOverFzeroPr}{at most 0.24 times the errors}  % experiments/FIGS/data/v2/index.json :: binary probe bbqq_vs_ccqq: (1 - AUC) of FLAV_F1 over FLAV_F0, run by run :: sha256 c6c4138ee5c86423 :: prose twin: two significant figures, direction in words
\newcommand{\MlstCmpBbqqFoneOverFzeroPrMax}{down to 0.15 times the errors}  % experiments/FIGS/data/v2/index.json :: binary probe bbqq_vs_ccqq: (1 - AUC) of FLAV_F1 over FLAV_F0, run by run :: sha256 c6c4138ee5c86423 :: prose twin of the run farther from 1
\newcommand{\MlstCmpBbqqFonerOverFone}{1.505 and 1.577}  % experiments/FIGS/data/v2/index.json :: binary probe bbqq_vs_ccqq: (1 - AUC) of FLAV_F1R over FLAV_F1, run by run :: sha256 c6c4138ee5c86423
\newcommand{\MlstCmpBbqqFonerOverFonePr}{at least 50\% more errors}  % experiments/FIGS/data/v2/index.json :: binary probe bbqq_vs_ccqq: (1 - AUC) of FLAV_F1R over FLAV_F1, run by run :: sha256 c6c4138ee5c86423 :: prose twin: two significant figures, direction in words
\newcommand{\MlstCmpBbqqFonerOverFonePrMax}{up to 58\% more errors}  % experiments/FIGS/data/v2/index.json :: binary probe bbqq_vs_ccqq: (1 - AUC) of FLAV_F1R over FLAV_F1, run by run :: sha256 c6c4138ee5c86423 :: prose twin of the run farther from 1
\newcommand{\MlstCmpBbqqFonerOverFzero}{0.36 and 0.24}  % experiments/FIGS/data/v2/index.json :: binary probe bbqq_vs_ccqq: (1 - AUC) of FLAV_F1R over FLAV_F0, run by run :: sha256 c6c4138ee5c86423
\newcommand{\MlstCmpBbqqFonerOverFzeroPr}{at most 0.36 times the errors}  % experiments/FIGS/data/v2/index.json :: binary probe bbqq_vs_ccqq: (1 - AUC) of FLAV_F1R over FLAV_F0, run by run :: sha256 c6c4138ee5c86423 :: prose twin: two significant figures, direction in words
\newcommand{\MlstCmpBbqqFonerOverFzeroPrMax}{down to 0.24 times the errors}  % experiments/FIGS/data/v2/index.json :: binary probe bbqq_vs_ccqq: (1 - AUC) of FLAV_F1R over FLAV_F0, run by run :: sha256 c6c4138ee5c86423 :: prose twin of the run farther from 1
\newcommand{\MlstCmpBbqqFzeroOverCoarse}{1.51 and 1.79}  % experiments/FIGS/data/v2/index.json :: binary probe bbqq_vs_ccqq: (1 - AUC) of FLAV_F0 over R16_Q1, run by run :: sha256 c6c4138ee5c86423
\newcommand{\MlstCmpBbqqFzeroOverCoarsePr}{at least 51\% more errors}  % experiments/FIGS/data/v2/index.json :: binary probe bbqq_vs_ccqq: (1 - AUC) of FLAV_F0 over R16_Q1, run by run :: sha256 c6c4138ee5c86423 :: prose twin: two significant figures, direction in words
\newcommand{\MlstCmpBbqqFzeroOverCoarsePrMax}{up to 79\% more errors}  % experiments/FIGS/data/v2/index.json :: binary probe bbqq_vs_ccqq: (1 - AUC) of FLAV_F0 over R16_Q1, run by run :: sha256 c6c4138ee5c86423 :: prose twin of the run farther from 1
\newcommand{\MlstCmpBbqqFoneOverCoarse}{0.356 and 0.276}  % experiments/FIGS/data/v2/index.json :: binary probe bbqq_vs_ccqq: (1 - AUC) of FLAV_F1 over R16_Q1, run by run :: sha256 c6c4138ee5c86423
\newcommand{\MlstCmpBbqqFoneOverCoarsePr}{at most 0.36 times the errors}  % experiments/FIGS/data/v2/index.json :: binary probe bbqq_vs_ccqq: (1 - AUC) of FLAV_F1 over R16_Q1, run by run :: sha256 c6c4138ee5c86423 :: prose twin: two significant figures, direction in words
\newcommand{\MlstCmpBbqqFoneOverCoarsePrMax}{down to 0.28 times the errors}  % experiments/FIGS/data/v2/index.json :: binary probe bbqq_vs_ccqq: (1 - AUC) of FLAV_F1 over R16_Q1, run by run :: sha256 c6c4138ee5c86423 :: prose twin of the run farther from 1
\newcommand{\MlstCmpBbqqFonerOverCoarse}{0.54 and 0.44}  % experiments/FIGS/data/v2/index.json :: binary probe bbqq_vs_ccqq: (1 - AUC) of FLAV_F1R over R16_Q1, run by run :: sha256 c6c4138ee5c86423
\newcommand{\MlstCmpBbqqFonerOverCoarsePr}{at least 46\% fewer errors}  % experiments/FIGS/data/v2/index.json :: binary probe bbqq_vs_ccqq: (1 - AUC) of FLAV_F1R over R16_Q1, run by run :: sha256 c6c4138ee5c86423 :: prose twin: two significant figures, direction in words
\newcommand{\MlstCmpBbqqFonerOverCoarsePrMax}{down to 0.44 times the errors}  % experiments/FIGS/data/v2/index.json :: binary probe bbqq_vs_ccqq: (1 - AUC) of FLAV_F1R over R16_Q1, run by run :: sha256 c6c4138ee5c86423 :: prose twin of the run farther from 1
\newcommand{\MlstCmpBbqqFzeroOverFull}{6.4 and 9.4}  % experiments/FIGS/data/v2/index.json :: binary probe bbqq_vs_ccqq: (1 - AUC) of FLAV_F0 over L188, run by run :: sha256 c6c4138ee5c86423
\newcommand{\MlstCmpBbqqFzeroOverFullPr}{at least 6.4 times the errors}  % experiments/FIGS/data/v2/index.json :: binary probe bbqq_vs_ccqq: (1 - AUC) of FLAV_F0 over L188, run by run :: sha256 c6c4138ee5c86423 :: prose twin: two significant figures, direction in words
\newcommand{\MlstCmpBbqqFzeroOverFullPrMax}{up to 9.4 times the errors}  % experiments/FIGS/data/v2/index.json :: binary probe bbqq_vs_ccqq: (1 - AUC) of FLAV_F0 over L188, run by run :: sha256 c6c4138ee5c86423 :: prose twin of the run farther from 1
\newcommand{\MlstCmpBbqqFoneOverFull}{1.510 and 1.447}  % experiments/FIGS/data/v2/index.json :: binary probe bbqq_vs_ccqq: (1 - AUC) of FLAV_F1 over L188, run by run :: sha256 c6c4138ee5c86423
\newcommand{\MlstCmpBbqqFoneOverFullPr}{at least 45\% more errors}  % experiments/FIGS/data/v2/index.json :: binary probe bbqq_vs_ccqq: (1 - AUC) of FLAV_F1 over L188, run by run :: sha256 c6c4138ee5c86423 :: prose twin: two significant figures, direction in words
\newcommand{\MlstCmpBbqqFoneOverFullPrMax}{up to 51\% more errors}  % experiments/FIGS/data/v2/index.json :: binary probe bbqq_vs_ccqq: (1 - AUC) of FLAV_F1 over L188, run by run :: sha256 c6c4138ee5c86423 :: prose twin of the run farther from 1
\newcommand{\MlstCmpBbqqFonerOverFull}{2.272 and 2.281}  % experiments/FIGS/data/v2/index.json :: binary probe bbqq_vs_ccqq: (1 - AUC) of FLAV_F1R over L188, run by run :: sha256 c6c4138ee5c86423
\newcommand{\MlstCmpBbqqFonerOverFullPr}{at least 2.3 times the errors}  % experiments/FIGS/data/v2/index.json :: binary probe bbqq_vs_ccqq: (1 - AUC) of FLAV_F1R over L188, run by run :: sha256 c6c4138ee5c86423 :: prose twin: two significant figures, direction in words
\newcommand{\MlstCmpBbqqFonerOverFullPrMax}{up to 2.3 times the errors}  % experiments/FIGS/data/v2/index.json :: binary probe bbqq_vs_ccqq: (1 - AUC) of FLAV_F1R over L188, run by run :: sha256 c6c4138ee5c86423 :: prose twin of the run farther from 1
\newcommand{\MlstCmpBbqqRandOverFull}{1.42\,\ensuremath{\pm}\,0.06}  % experiments/FIGS/data/v2/index.json :: binary probe bbqq_vs_ccqq: (1 - AUC) of each random partition over the 188-class run of its seed :: sha256 c6c4138ee5c86423
\newcommand{\MlstCmpBbqqRandOverFullPr}{42\% more errors}  % experiments/FIGS/data/v2/index.json :: binary probe bbqq_vs_ccqq: (1 - AUC) of each random partition over the 188-class run of its seed :: sha256 c6c4138ee5c86423 :: prose twin: two significant figures, direction in words
\newcommand{\MlstPairRandMergedBbqqN}{2}  % experiments/FIGS/data/v2/index.json :: random partitions that merge bbqq_vs_ccqq (count) :: sha256 c6c4138ee5c86423
\newcommand{\MlstPairRandApartBbqqN}{3}  % experiments/FIGS/data/v2/index.json :: random partitions that keep apart bbqq_vs_ccqq (count) :: sha256 c6c4138ee5c86423
\newcommand{\MlstPairRandMergedOverApartBbqq}{0.996}  % experiments/FIGS/data/v2/index.json :: binary probe bbqq_vs_ccqq: mean (1 - AUC) of the partitions that merge it over that of those that keep it apart (one run each) :: sha256 c6c4138ee5c86423
\newcommand{\MlstPairRandMergedOverApartBbqqPr}{0.44\% fewer errors}  % experiments/FIGS/data/v2/index.json :: binary probe bbqq_vs_ccqq: mean (1 - AUC) of the partitions that merge it over that of those that keep it apart (one run each) :: sha256 c6c4138ee5c86423 :: prose twin: two significant figures, direction in words
\newcommand{\MlstCmpEemmResOverFull}{--}  % experiments/FIGS/data/v2/index.json :: binary probe ee_vs_mm: (1 - AUC) of L162 over L188, run by run :: sha256 c6c4138ee5c86423 :: the reference's 1 - AUC is 0
\newcommand{\MlstCmpEemmResOverFullPr}{undefined, since the reference separates the test jets completely}  % experiments/FIGS/data/v2/index.json :: binary probe ee_vs_mm: (1 - AUC) of L162 over L188, run by run :: sha256 c6c4138ee5c86423 :: the reference's 1 - AUC is 0
\newcommand{\MlstCmpEemmMidOverFull}{--}  % experiments/FIGS/data/v2/index.json :: binary probe ee_vs_mm: (1 - AUC) of R42_Q1 over L188, run by run :: sha256 c6c4138ee5c86423 :: the reference's 1 - AUC is 0
\newcommand{\MlstCmpEemmMidOverFullPr}{undefined, since the reference separates the test jets completely}  % experiments/FIGS/data/v2/index.json :: binary probe ee_vs_mm: (1 - AUC) of R42_Q1 over L188, run by run :: sha256 c6c4138ee5c86423 :: the reference's 1 - AUC is 0
\newcommand{\MlstCmpEemmCoarseOverFull}{--}  % experiments/FIGS/data/v2/index.json :: binary probe ee_vs_mm: (1 - AUC) of R16_Q1 over L188, run by run :: sha256 c6c4138ee5c86423 :: the reference's 1 - AUC is 0
\newcommand{\MlstCmpEemmCoarseOverFullPr}{undefined, since the reference separates the test jets completely}  % experiments/FIGS/data/v2/index.json :: binary probe ee_vs_mm: (1 - AUC) of R16_Q1 over L188, run by run :: sha256 c6c4138ee5c86423 :: the reference's 1 - AUC is 0
\newcommand{\MlstCmpEemmFoneOverFzero}{--}  % experiments/FIGS/data/v2/index.json :: binary probe ee_vs_mm: (1 - AUC) of FLAV_F1 over FLAV_F0, run by run :: sha256 c6c4138ee5c86423 :: the reference's 1 - AUC is 0
\newcommand{\MlstCmpEemmFoneOverFzeroPr}{undefined, since the reference separates the test jets completely}  % experiments/FIGS/data/v2/index.json :: binary probe ee_vs_mm: (1 - AUC) of FLAV_F1 over FLAV_F0, run by run :: sha256 c6c4138ee5c86423 :: the reference's 1 - AUC is 0
\newcommand{\MlstCmpEemmFoneOverFzeroPrMax}{undefined, since the reference separates the test jets completely}  % experiments/FIGS/data/v2/index.json :: binary probe ee_vs_mm: (1 - AUC) of FLAV_F1 over FLAV_F0, run by run :: sha256 c6c4138ee5c86423 :: the reference's 1 - AUC is 0
\newcommand{\MlstCmpEemmFonerOverFone}{--}  % experiments/FIGS/data/v2/index.json :: binary probe ee_vs_mm: (1 - AUC) of FLAV_F1R over FLAV_F1, run by run :: sha256 c6c4138ee5c86423 :: the reference's 1 - AUC is 0
\newcommand{\MlstCmpEemmFonerOverFonePr}{undefined, since the reference separates the test jets completely}  % experiments/FIGS/data/v2/index.json :: binary probe ee_vs_mm: (1 - AUC) of FLAV_F1R over FLAV_F1, run by run :: sha256 c6c4138ee5c86423 :: the reference's 1 - AUC is 0
\newcommand{\MlstCmpEemmFonerOverFonePrMax}{undefined, since the reference separates the test jets completely}  % experiments/FIGS/data/v2/index.json :: binary probe ee_vs_mm: (1 - AUC) of FLAV_F1R over FLAV_F1, run by run :: sha256 c6c4138ee5c86423 :: the reference's 1 - AUC is 0
\newcommand{\MlstCmpEemmFonerOverFzero}{--}  % experiments/FIGS/data/v2/index.json :: binary probe ee_vs_mm: (1 - AUC) of FLAV_F1R over FLAV_F0, run by run :: sha256 c6c4138ee5c86423 :: the reference's 1 - AUC is 0
\newcommand{\MlstCmpEemmFonerOverFzeroPr}{undefined, since the reference separates the test jets completely}  % experiments/FIGS/data/v2/index.json :: binary probe ee_vs_mm: (1 - AUC) of FLAV_F1R over FLAV_F0, run by run :: sha256 c6c4138ee5c86423 :: the reference's 1 - AUC is 0
\newcommand{\MlstCmpEemmFonerOverFzeroPrMax}{undefined, since the reference separates the test jets completely}  % experiments/FIGS/data/v2/index.json :: binary probe ee_vs_mm: (1 - AUC) of FLAV_F1R over FLAV_F0, run by run :: sha256 c6c4138ee5c86423 :: the reference's 1 - AUC is 0
\newcommand{\MlstCmpEemmFzeroOverCoarse}{0.0000022 and 0.0000003}  % experiments/FIGS/data/v2/index.json :: binary probe ee_vs_mm: (1 - AUC) of FLAV_F0 over R16_Q1, run by run :: sha256 c6c4138ee5c86423
\newcommand{\MlstCmpEemmFzeroOverCoarsePr}{at most 0.0000022 times the errors}  % experiments/FIGS/data/v2/index.json :: binary probe ee_vs_mm: (1 - AUC) of FLAV_F0 over R16_Q1, run by run :: sha256 c6c4138ee5c86423 :: prose twin: two significant figures, direction in words
\newcommand{\MlstCmpEemmFzeroOverCoarsePrMax}{down to 0.00000034 times the errors}  % experiments/FIGS/data/v2/index.json :: binary probe ee_vs_mm: (1 - AUC) of FLAV_F0 over R16_Q1, run by run :: sha256 c6c4138ee5c86423 :: prose twin of the run farther from 1
\newcommand{\MlstCmpEemmFoneOverCoarse}{0.000000 and 0.000038}  % experiments/FIGS/data/v2/index.json :: binary probe ee_vs_mm: (1 - AUC) of FLAV_F1 over R16_Q1, run by run :: sha256 c6c4138ee5c86423
\newcommand{\MlstCmpEemmFoneOverCoarsePr}{at most 0.000038 times the errors}  % experiments/FIGS/data/v2/index.json :: binary probe ee_vs_mm: (1 - AUC) of FLAV_F1 over R16_Q1, run by run :: sha256 c6c4138ee5c86423 :: prose twin: two significant figures, direction in words
\newcommand{\MlstCmpEemmFoneOverCoarsePrMax}{down to 0.000000042 times the errors}  % experiments/FIGS/data/v2/index.json :: binary probe ee_vs_mm: (1 - AUC) of FLAV_F1 over R16_Q1, run by run :: sha256 c6c4138ee5c86423 :: prose twin of the run farther from 1
\newcommand{\MlstCmpEemmFonerOverCoarse}{0.00000005 and 0.00000017}  % experiments/FIGS/data/v2/index.json :: binary probe ee_vs_mm: (1 - AUC) of FLAV_F1R over R16_Q1, run by run :: sha256 c6c4138ee5c86423
\newcommand{\MlstCmpEemmFonerOverCoarsePr}{at most 0.00000017 times the errors}  % experiments/FIGS/data/v2/index.json :: binary probe ee_vs_mm: (1 - AUC) of FLAV_F1R over R16_Q1, run by run :: sha256 c6c4138ee5c86423 :: prose twin: two significant figures, direction in words
\newcommand{\MlstCmpEemmFonerOverCoarsePrMax}{down to 0.000000055 times the errors}  % experiments/FIGS/data/v2/index.json :: binary probe ee_vs_mm: (1 - AUC) of FLAV_F1R over R16_Q1, run by run :: sha256 c6c4138ee5c86423 :: prose twin of the run farther from 1
\newcommand{\MlstCmpEemmFzeroOverFull}{--}  % experiments/FIGS/data/v2/index.json :: binary probe ee_vs_mm: (1 - AUC) of FLAV_F0 over L188, run by run :: sha256 c6c4138ee5c86423 :: the reference's 1 - AUC is 0
\newcommand{\MlstCmpEemmFzeroOverFullPr}{undefined, since the reference separates the test jets completely}  % experiments/FIGS/data/v2/index.json :: binary probe ee_vs_mm: (1 - AUC) of FLAV_F0 over L188, run by run :: sha256 c6c4138ee5c86423 :: the reference's 1 - AUC is 0
\newcommand{\MlstCmpEemmFzeroOverFullPrMax}{undefined, since the reference separates the test jets completely}  % experiments/FIGS/data/v2/index.json :: binary probe ee_vs_mm: (1 - AUC) of FLAV_F0 over L188, run by run :: sha256 c6c4138ee5c86423 :: the reference's 1 - AUC is 0
\newcommand{\MlstCmpEemmFoneOverFull}{--}  % experiments/FIGS/data/v2/index.json :: binary probe ee_vs_mm: (1 - AUC) of FLAV_F1 over L188, run by run :: sha256 c6c4138ee5c86423 :: the reference's 1 - AUC is 0
\newcommand{\MlstCmpEemmFoneOverFullPr}{undefined, since the reference separates the test jets completely}  % experiments/FIGS/data/v2/index.json :: binary probe ee_vs_mm: (1 - AUC) of FLAV_F1 over L188, run by run :: sha256 c6c4138ee5c86423 :: the reference's 1 - AUC is 0
\newcommand{\MlstCmpEemmFoneOverFullPrMax}{undefined, since the reference separates the test jets completely}  % experiments/FIGS/data/v2/index.json :: binary probe ee_vs_mm: (1 - AUC) of FLAV_F1 over L188, run by run :: sha256 c6c4138ee5c86423 :: the reference's 1 - AUC is 0
\newcommand{\MlstCmpEemmFonerOverFull}{--}  % experiments/FIGS/data/v2/index.json :: binary probe ee_vs_mm: (1 - AUC) of FLAV_F1R over L188, run by run :: sha256 c6c4138ee5c86423 :: the reference's 1 - AUC is 0
\newcommand{\MlstCmpEemmFonerOverFullPr}{undefined, since the reference separates the test jets completely}  % experiments/FIGS/data/v2/index.json :: binary probe ee_vs_mm: (1 - AUC) of FLAV_F1R over L188, run by run :: sha256 c6c4138ee5c86423 :: the reference's 1 - AUC is 0
\newcommand{\MlstCmpEemmFonerOverFullPrMax}{undefined, since the reference separates the test jets completely}  % experiments/FIGS/data/v2/index.json :: binary probe ee_vs_mm: (1 - AUC) of FLAV_F1R over L188, run by run :: sha256 c6c4138ee5c86423 :: the reference's 1 - AUC is 0
\newcommand{\MlstCmpEemmRandOverFull}{--}  % experiments/FIGS/data/v2/index.json :: binary probe ee_vs_mm: (1 - AUC) of each random partition over the 188-class run of its seed :: sha256 c6c4138ee5c86423 :: the reference's 1 - AUC is 0
\newcommand{\MlstCmpEemmRandOverFullPr}{undefined, since the reference separates the test jets completely}  % experiments/FIGS/data/v2/index.json :: binary probe ee_vs_mm: (1 - AUC) of each random partition over the 188-class run of its seed :: sha256 c6c4138ee5c86423 :: the reference's 1 - AUC is 0
\newcommand{\MlstPairRandMergedEemmN}{3}  % experiments/FIGS/data/v2/index.json :: random partitions that merge ee_vs_mm (count) :: sha256 c6c4138ee5c86423
\newcommand{\MlstPairRandApartEemmN}{2}  % experiments/FIGS/data/v2/index.json :: random partitions that keep apart ee_vs_mm (count) :: sha256 c6c4138ee5c86423
\newcommand{\MlstPairRandMergedOverApartEemm}{75.3}  % experiments/FIGS/data/v2/index.json :: binary probe ee_vs_mm: mean (1 - AUC) of the partitions that merge it over that of those that keep it apart (one run each) :: sha256 c6c4138ee5c86423
\newcommand{\MlstPairRandMergedOverApartEemmPr}{75 times the errors}  % experiments/FIGS/data/v2/index.json :: binary probe ee_vs_mm: mean (1 - AUC) of the partitions that merge it over that of those that keep it apart (one run each) :: sha256 c6c4138ee5c86423 :: prose twin: two significant figures, direction in words
\newcommand{\MlstCmpQcdBbResOverFull}{1.15\,\ensuremath{\pm}\,0.04}  % experiments/FIGS/data/v2/index.json :: binary probe qcd_bb_vs_cc: (1 - AUC) of L162 over L188, run by run :: sha256 c6c4138ee5c86423
\newcommand{\MlstCmpQcdBbResOverFullPr}{15\% more errors}  % experiments/FIGS/data/v2/index.json :: binary probe qcd_bb_vs_cc: (1 - AUC) of L162 over L188, run by run :: sha256 c6c4138ee5c86423 :: prose twin: two significant figures, direction in words
\newcommand{\MlstCmpQcdBbMidOverFull}{1.163\,\ensuremath{\pm}\,0.016}  % experiments/FIGS/data/v2/index.json :: binary probe qcd_bb_vs_cc: (1 - AUC) of R42_Q1 over L188, run by run :: sha256 c6c4138ee5c86423
\newcommand{\MlstCmpQcdBbMidOverFullPr}{16\% more errors}  % experiments/FIGS/data/v2/index.json :: binary probe qcd_bb_vs_cc: (1 - AUC) of R42_Q1 over L188, run by run :: sha256 c6c4138ee5c86423 :: prose twin: two significant figures, direction in words
\newcommand{\MlstCmpQcdBbCoarseOverFull}{3.294\,\ensuremath{\pm}\,0.033}  % experiments/FIGS/data/v2/index.json :: binary probe qcd_bb_vs_cc: (1 - AUC) of R16_Q1 over L188, run by run :: sha256 c6c4138ee5c86423
\newcommand{\MlstCmpQcdBbCoarseOverFullPr}{3.3 times the errors}  % experiments/FIGS/data/v2/index.json :: binary probe qcd_bb_vs_cc: (1 - AUC) of R16_Q1 over L188, run by run :: sha256 c6c4138ee5c86423 :: prose twin: two significant figures, direction in words
\newcommand{\MlstCmpQcdBbFoneOverFzero}{0.368 and 0.294}  % experiments/FIGS/data/v2/index.json :: binary probe qcd_bb_vs_cc: (1 - AUC) of FLAV_F1 over FLAV_F0, run by run :: sha256 c6c4138ee5c86423
\newcommand{\MlstCmpQcdBbFoneOverFzeroPr}{at most 0.37 times the errors}  % experiments/FIGS/data/v2/index.json :: binary probe qcd_bb_vs_cc: (1 - AUC) of FLAV_F1 over FLAV_F0, run by run :: sha256 c6c4138ee5c86423 :: prose twin: two significant figures, direction in words
\newcommand{\MlstCmpQcdBbFoneOverFzeroPrMax}{down to 0.29 times the errors}  % experiments/FIGS/data/v2/index.json :: binary probe qcd_bb_vs_cc: (1 - AUC) of FLAV_F1 over FLAV_F0, run by run :: sha256 c6c4138ee5c86423 :: prose twin of the run farther from 1
\newcommand{\MlstCmpQcdBbFonerOverFone}{1.361 and 1.429}  % experiments/FIGS/data/v2/index.json :: binary probe qcd_bb_vs_cc: (1 - AUC) of FLAV_F1R over FLAV_F1, run by run :: sha256 c6c4138ee5c86423
\newcommand{\MlstCmpQcdBbFonerOverFonePr}{at least 36\% more errors}  % experiments/FIGS/data/v2/index.json :: binary probe qcd_bb_vs_cc: (1 - AUC) of FLAV_F1R over FLAV_F1, run by run :: sha256 c6c4138ee5c86423 :: prose twin: two significant figures, direction in words
\newcommand{\MlstCmpQcdBbFonerOverFonePrMax}{up to 43\% more errors}  % experiments/FIGS/data/v2/index.json :: binary probe qcd_bb_vs_cc: (1 - AUC) of FLAV_F1R over FLAV_F1, run by run :: sha256 c6c4138ee5c86423 :: prose twin of the run farther from 1
\newcommand{\MlstCmpQcdBbFonerOverFzero}{0.501 and 0.420}  % experiments/FIGS/data/v2/index.json :: binary probe qcd_bb_vs_cc: (1 - AUC) of FLAV_F1R over FLAV_F0, run by run :: sha256 c6c4138ee5c86423
\newcommand{\MlstCmpQcdBbFonerOverFzeroPr}{at least 50\% fewer errors}  % experiments/FIGS/data/v2/index.json :: binary probe qcd_bb_vs_cc: (1 - AUC) of FLAV_F1R over FLAV_F0, run by run :: sha256 c6c4138ee5c86423 :: prose twin: two significant figures, direction in words
\newcommand{\MlstCmpQcdBbFonerOverFzeroPrMax}{down to 0.42 times the errors}  % experiments/FIGS/data/v2/index.json :: binary probe qcd_bb_vs_cc: (1 - AUC) of FLAV_F1R over FLAV_F0, run by run :: sha256 c6c4138ee5c86423 :: prose twin of the run farther from 1
\newcommand{\MlstCmpQcdBbFzeroOverCoarse}{1.28 and 1.58}  % experiments/FIGS/data/v2/index.json :: binary probe qcd_bb_vs_cc: (1 - AUC) of FLAV_F0 over R16_Q1, run by run :: sha256 c6c4138ee5c86423
\newcommand{\MlstCmpQcdBbFzeroOverCoarsePr}{at least 28\% more errors}  % experiments/FIGS/data/v2/index.json :: binary probe qcd_bb_vs_cc: (1 - AUC) of FLAV_F0 over R16_Q1, run by run :: sha256 c6c4138ee5c86423 :: prose twin: two significant figures, direction in words
\newcommand{\MlstCmpQcdBbFzeroOverCoarsePrMax}{up to 58\% more errors}  % experiments/FIGS/data/v2/index.json :: binary probe qcd_bb_vs_cc: (1 - AUC) of FLAV_F0 over R16_Q1, run by run :: sha256 c6c4138ee5c86423 :: prose twin of the run farther from 1
\newcommand{\MlstCmpQcdBbFoneOverCoarse}{0.4692 and 0.4648}  % experiments/FIGS/data/v2/index.json :: binary probe qcd_bb_vs_cc: (1 - AUC) of FLAV_F1 over R16_Q1, run by run :: sha256 c6c4138ee5c86423
\newcommand{\MlstCmpQcdBbFoneOverCoarsePr}{at most 0.47 times the errors}  % experiments/FIGS/data/v2/index.json :: binary probe qcd_bb_vs_cc: (1 - AUC) of FLAV_F1 over R16_Q1, run by run :: sha256 c6c4138ee5c86423 :: prose twin: two significant figures, direction in words
\newcommand{\MlstCmpQcdBbFoneOverCoarsePrMax}{down to 0.46 times the errors}  % experiments/FIGS/data/v2/index.json :: binary probe qcd_bb_vs_cc: (1 - AUC) of FLAV_F1 over R16_Q1, run by run :: sha256 c6c4138ee5c86423 :: prose twin of the run farther from 1
\newcommand{\MlstCmpQcdBbFonerOverCoarse}{0.639 and 0.664}  % experiments/FIGS/data/v2/index.json :: binary probe qcd_bb_vs_cc: (1 - AUC) of FLAV_F1R over R16_Q1, run by run :: sha256 c6c4138ee5c86423
\newcommand{\MlstCmpQcdBbFonerOverCoarsePr}{at least 34\% fewer errors}  % experiments/FIGS/data/v2/index.json :: binary probe qcd_bb_vs_cc: (1 - AUC) of FLAV_F1R over R16_Q1, run by run :: sha256 c6c4138ee5c86423 :: prose twin: two significant figures, direction in words
\newcommand{\MlstCmpQcdBbFonerOverCoarsePrMax}{up to 36\% fewer errors}  % experiments/FIGS/data/v2/index.json :: binary probe qcd_bb_vs_cc: (1 - AUC) of FLAV_F1R over R16_Q1, run by run :: sha256 c6c4138ee5c86423 :: prose twin of the run farther from 1
\newcommand{\MlstCmpQcdBbFzeroOverFull}{4.24 and 5.21}  % experiments/FIGS/data/v2/index.json :: binary probe qcd_bb_vs_cc: (1 - AUC) of FLAV_F0 over L188, run by run :: sha256 c6c4138ee5c86423
\newcommand{\MlstCmpQcdBbFzeroOverFullPr}{at least 4.2 times the errors}  % experiments/FIGS/data/v2/index.json :: binary probe qcd_bb_vs_cc: (1 - AUC) of FLAV_F0 over L188, run by run :: sha256 c6c4138ee5c86423 :: prose twin: two significant figures, direction in words
\newcommand{\MlstCmpQcdBbFzeroOverFullPrMax}{up to 5.2 times the errors}  % experiments/FIGS/data/v2/index.json :: binary probe qcd_bb_vs_cc: (1 - AUC) of FLAV_F0 over L188, run by run :: sha256 c6c4138ee5c86423 :: prose twin of the run farther from 1
\newcommand{\MlstCmpQcdBbFoneOverFull}{1.561 and 1.532}  % experiments/FIGS/data/v2/index.json :: binary probe qcd_bb_vs_cc: (1 - AUC) of FLAV_F1 over L188, run by run :: sha256 c6c4138ee5c86423
\newcommand{\MlstCmpQcdBbFoneOverFullPr}{at least 53\% more errors}  % experiments/FIGS/data/v2/index.json :: binary probe qcd_bb_vs_cc: (1 - AUC) of FLAV_F1 over L188, run by run :: sha256 c6c4138ee5c86423 :: prose twin: two significant figures, direction in words
\newcommand{\MlstCmpQcdBbFoneOverFullPrMax}{up to 56\% more errors}  % experiments/FIGS/data/v2/index.json :: binary probe qcd_bb_vs_cc: (1 - AUC) of FLAV_F1 over L188, run by run :: sha256 c6c4138ee5c86423 :: prose twin of the run farther from 1
\newcommand{\MlstCmpQcdBbFonerOverFull}{2.124 and 2.189}  % experiments/FIGS/data/v2/index.json :: binary probe qcd_bb_vs_cc: (1 - AUC) of FLAV_F1R over L188, run by run :: sha256 c6c4138ee5c86423
\newcommand{\MlstCmpQcdBbFonerOverFullPr}{at least 2.1 times the errors}  % experiments/FIGS/data/v2/index.json :: binary probe qcd_bb_vs_cc: (1 - AUC) of FLAV_F1R over L188, run by run :: sha256 c6c4138ee5c86423 :: prose twin: two significant figures, direction in words
\newcommand{\MlstCmpQcdBbFonerOverFullPrMax}{up to 2.2 times the errors}  % experiments/FIGS/data/v2/index.json :: binary probe qcd_bb_vs_cc: (1 - AUC) of FLAV_F1R over L188, run by run :: sha256 c6c4138ee5c86423 :: prose twin of the run farther from 1
\newcommand{\MlstCmpQcdBResOverFull}{1.085\,\ensuremath{\pm}\,0.004}  % experiments/FIGS/data/v2/index.json :: binary probe qcd_b_vs_light: (1 - AUC) of L162 over L188, run by run :: sha256 c6c4138ee5c86423
\newcommand{\MlstCmpQcdBResOverFullPr}{8.5\% more errors}  % experiments/FIGS/data/v2/index.json :: binary probe qcd_b_vs_light: (1 - AUC) of L162 over L188, run by run :: sha256 c6c4138ee5c86423 :: prose twin: two significant figures, direction in words
\newcommand{\MlstCmpQcdBMidOverFull}{1.081\,\ensuremath{\pm}\,0.008}  % experiments/FIGS/data/v2/index.json :: binary probe qcd_b_vs_light: (1 - AUC) of R42_Q1 over L188, run by run :: sha256 c6c4138ee5c86423
\newcommand{\MlstCmpQcdBMidOverFullPr}{8.1\% more errors}  % experiments/FIGS/data/v2/index.json :: binary probe qcd_b_vs_light: (1 - AUC) of R42_Q1 over L188, run by run :: sha256 c6c4138ee5c86423 :: prose twin: two significant figures, direction in words
\newcommand{\MlstCmpQcdBCoarseOverFull}{1.54\,\ensuremath{\pm}\,0.07}  % experiments/FIGS/data/v2/index.json :: binary probe qcd_b_vs_light: (1 - AUC) of R16_Q1 over L188, run by run :: sha256 c6c4138ee5c86423
\newcommand{\MlstCmpQcdBCoarseOverFullPr}{54\% more errors}  % experiments/FIGS/data/v2/index.json :: binary probe qcd_b_vs_light: (1 - AUC) of R16_Q1 over L188, run by run :: sha256 c6c4138ee5c86423 :: prose twin: two significant figures, direction in words
\newcommand{\MlstCmpQcdBFoneOverFzero}{0.778 and 0.835}  % experiments/FIGS/data/v2/index.json :: binary probe qcd_b_vs_light: (1 - AUC) of FLAV_F1 over FLAV_F0, run by run :: sha256 c6c4138ee5c86423
\newcommand{\MlstCmpQcdBFoneOverFzeroPr}{at least 17\% fewer errors}  % experiments/FIGS/data/v2/index.json :: binary probe qcd_b_vs_light: (1 - AUC) of FLAV_F1 over FLAV_F0, run by run :: sha256 c6c4138ee5c86423 :: prose twin: two significant figures, direction in words
\newcommand{\MlstCmpQcdBFoneOverFzeroPrMax}{up to 22\% fewer errors}  % experiments/FIGS/data/v2/index.json :: binary probe qcd_b_vs_light: (1 - AUC) of FLAV_F1 over FLAV_F0, run by run :: sha256 c6c4138ee5c86423 :: prose twin of the run farther from 1
\newcommand{\MlstCmpQcdBFonerOverFone}{1.054 and 1.061}  % experiments/FIGS/data/v2/index.json :: binary probe qcd_b_vs_light: (1 - AUC) of FLAV_F1R over FLAV_F1, run by run :: sha256 c6c4138ee5c86423
\newcommand{\MlstCmpQcdBFonerOverFonePr}{equal within 7\%}  % experiments/FIGS/data/v2/index.json :: binary probe qcd_b_vs_light: (1 - AUC) of FLAV_F1R over FLAV_F1, run by run :: sha256 c6c4138ee5c86423 :: prose twin: two significant figures, direction in words
\newcommand{\MlstCmpQcdBFonerOverFonePrMax}{equal within 7\%}  % experiments/FIGS/data/v2/index.json :: binary probe qcd_b_vs_light: (1 - AUC) of FLAV_F1R over FLAV_F1, run by run :: sha256 c6c4138ee5c86423 :: prose twin of the run farther from 1
\newcommand{\MlstCmpQcdBFonerOverFzero}{0.820 and 0.885}  % experiments/FIGS/data/v2/index.json :: binary probe qcd_b_vs_light: (1 - AUC) of FLAV_F1R over FLAV_F0, run by run :: sha256 c6c4138ee5c86423
\newcommand{\MlstCmpQcdBFonerOverFzeroPr}{at least 11\% fewer errors}  % experiments/FIGS/data/v2/index.json :: binary probe qcd_b_vs_light: (1 - AUC) of FLAV_F1R over FLAV_F0, run by run :: sha256 c6c4138ee5c86423 :: prose twin: two significant figures, direction in words
\newcommand{\MlstCmpQcdBFonerOverFzeroPrMax}{up to 18\% fewer errors}  % experiments/FIGS/data/v2/index.json :: binary probe qcd_b_vs_light: (1 - AUC) of FLAV_F1R over FLAV_F0, run by run :: sha256 c6c4138ee5c86423 :: prose twin of the run farther from 1
\newcommand{\MlstCmpQcdBFzeroOverCoarse}{0.962 and 0.981}  % experiments/FIGS/data/v2/index.json :: binary probe qcd_b_vs_light: (1 - AUC) of FLAV_F0 over R16_Q1, run by run :: sha256 c6c4138ee5c86423
\newcommand{\MlstCmpQcdBFzeroOverCoarsePr}{equal within 4\%}  % experiments/FIGS/data/v2/index.json :: binary probe qcd_b_vs_light: (1 - AUC) of FLAV_F0 over R16_Q1, run by run :: sha256 c6c4138ee5c86423 :: prose twin: two significant figures, direction in words
\newcommand{\MlstCmpQcdBFzeroOverCoarsePrMax}{equal within 4\%}  % experiments/FIGS/data/v2/index.json :: binary probe qcd_b_vs_light: (1 - AUC) of FLAV_F0 over R16_Q1, run by run :: sha256 c6c4138ee5c86423 :: prose twin of the run farther from 1
\newcommand{\MlstCmpQcdBFoneOverCoarse}{0.749 and 0.819}  % experiments/FIGS/data/v2/index.json :: binary probe qcd_b_vs_light: (1 - AUC) of FLAV_F1 over R16_Q1, run by run :: sha256 c6c4138ee5c86423
\newcommand{\MlstCmpQcdBFoneOverCoarsePr}{at least 18\% fewer errors}  % experiments/FIGS/data/v2/index.json :: binary probe qcd_b_vs_light: (1 - AUC) of FLAV_F1 over R16_Q1, run by run :: sha256 c6c4138ee5c86423 :: prose twin: two significant figures, direction in words
\newcommand{\MlstCmpQcdBFoneOverCoarsePrMax}{up to 25\% fewer errors}  % experiments/FIGS/data/v2/index.json :: binary probe qcd_b_vs_light: (1 - AUC) of FLAV_F1 over R16_Q1, run by run :: sha256 c6c4138ee5c86423 :: prose twin of the run farther from 1
\newcommand{\MlstCmpQcdBFonerOverCoarse}{0.789 and 0.869}  % experiments/FIGS/data/v2/index.json :: binary probe qcd_b_vs_light: (1 - AUC) of FLAV_F1R over R16_Q1, run by run :: sha256 c6c4138ee5c86423
\newcommand{\MlstCmpQcdBFonerOverCoarsePr}{at least 13\% fewer errors}  % experiments/FIGS/data/v2/index.json :: binary probe qcd_b_vs_light: (1 - AUC) of FLAV_F1R over R16_Q1, run by run :: sha256 c6c4138ee5c86423 :: prose twin: two significant figures, direction in words
\newcommand{\MlstCmpQcdBFonerOverCoarsePrMax}{up to 21\% fewer errors}  % experiments/FIGS/data/v2/index.json :: binary probe qcd_b_vs_light: (1 - AUC) of FLAV_F1R over R16_Q1, run by run :: sha256 c6c4138ee5c86423 :: prose twin of the run farther from 1
\newcommand{\MlstCmpQcdBFzeroOverFull}{1.56 and 1.45}  % experiments/FIGS/data/v2/index.json :: binary probe qcd_b_vs_light: (1 - AUC) of FLAV_F0 over L188, run by run :: sha256 c6c4138ee5c86423
\newcommand{\MlstCmpQcdBFzeroOverFullPr}{at least 45\% more errors}  % experiments/FIGS/data/v2/index.json :: binary probe qcd_b_vs_light: (1 - AUC) of FLAV_F0 over L188, run by run :: sha256 c6c4138ee5c86423 :: prose twin: two significant figures, direction in words
\newcommand{\MlstCmpQcdBFzeroOverFullPrMax}{up to 56\% more errors}  % experiments/FIGS/data/v2/index.json :: binary probe qcd_b_vs_light: (1 - AUC) of FLAV_F0 over L188, run by run :: sha256 c6c4138ee5c86423 :: prose twin of the run farther from 1
\newcommand{\MlstCmpQcdBFoneOverFull}{1.212 and 1.213}  % experiments/FIGS/data/v2/index.json :: binary probe qcd_b_vs_light: (1 - AUC) of FLAV_F1 over L188, run by run :: sha256 c6c4138ee5c86423
\newcommand{\MlstCmpQcdBFoneOverFullPr}{at least 21\% more errors}  % experiments/FIGS/data/v2/index.json :: binary probe qcd_b_vs_light: (1 - AUC) of FLAV_F1 over L188, run by run :: sha256 c6c4138ee5c86423 :: prose twin: two significant figures, direction in words
\newcommand{\MlstCmpQcdBFoneOverFullPrMax}{up to 21\% more errors}  % experiments/FIGS/data/v2/index.json :: binary probe qcd_b_vs_light: (1 - AUC) of FLAV_F1 over L188, run by run :: sha256 c6c4138ee5c86423 :: prose twin of the run farther from 1
\newcommand{\MlstCmpQcdBFonerOverFull}{1.277 and 1.286}  % experiments/FIGS/data/v2/index.json :: binary probe qcd_b_vs_light: (1 - AUC) of FLAV_F1R over L188, run by run :: sha256 c6c4138ee5c86423
\newcommand{\MlstCmpQcdBFonerOverFullPr}{at least 28\% more errors}  % experiments/FIGS/data/v2/index.json :: binary probe qcd_b_vs_light: (1 - AUC) of FLAV_F1R over L188, run by run :: sha256 c6c4138ee5c86423 :: prose twin: two significant figures, direction in words
\newcommand{\MlstCmpQcdBFonerOverFullPrMax}{up to 29\% more errors}  % experiments/FIGS/data/v2/index.json :: binary probe qcd_b_vs_light: (1 - AUC) of FLAV_F1R over L188, run by run :: sha256 c6c4138ee5c86423 :: prose twin of the run farther from 1
\newcommand{\MlstCmpQcdCResOverFull}{1.061\,\ensuremath{\pm}\,0.005}  % experiments/FIGS/data/v2/index.json :: binary probe qcd_c_vs_light: (1 - AUC) of L162 over L188, run by run :: sha256 c6c4138ee5c86423
\newcommand{\MlstCmpQcdCResOverFullPr}{6.1\% more errors}  % experiments/FIGS/data/v2/index.json :: binary probe qcd_c_vs_light: (1 - AUC) of L162 over L188, run by run :: sha256 c6c4138ee5c86423 :: prose twin: two significant figures, direction in words
\newcommand{\MlstCmpQcdCMidOverFull}{1.072\,\ensuremath{\pm}\,0.008}  % experiments/FIGS/data/v2/index.json :: binary probe qcd_c_vs_light: (1 - AUC) of R42_Q1 over L188, run by run :: sha256 c6c4138ee5c86423
\newcommand{\MlstCmpQcdCMidOverFullPr}{7.2\% more errors}  % experiments/FIGS/data/v2/index.json :: binary probe qcd_c_vs_light: (1 - AUC) of R42_Q1 over L188, run by run :: sha256 c6c4138ee5c86423 :: prose twin: two significant figures, direction in words
\newcommand{\MlstCmpQcdCCoarseOverFull}{1.407\,\ensuremath{\pm}\,0.021}  % experiments/FIGS/data/v2/index.json :: binary probe qcd_c_vs_light: (1 - AUC) of R16_Q1 over L188, run by run :: sha256 c6c4138ee5c86423
\newcommand{\MlstCmpQcdCCoarseOverFullPr}{41\% more errors}  % experiments/FIGS/data/v2/index.json :: binary probe qcd_c_vs_light: (1 - AUC) of R16_Q1 over L188, run by run :: sha256 c6c4138ee5c86423 :: prose twin: two significant figures, direction in words
\newcommand{\MlstCmpQcdCFoneOverFzero}{0.915 and 0.944}  % experiments/FIGS/data/v2/index.json :: binary probe qcd_c_vs_light: (1 - AUC) of FLAV_F1 over FLAV_F0, run by run :: sha256 c6c4138ee5c86423
\newcommand{\MlstCmpQcdCFoneOverFzeroPr}{at least 5.6\% fewer errors}  % experiments/FIGS/data/v2/index.json :: binary probe qcd_c_vs_light: (1 - AUC) of FLAV_F1 over FLAV_F0, run by run :: sha256 c6c4138ee5c86423 :: prose twin: two significant figures, direction in words
\newcommand{\MlstCmpQcdCFoneOverFzeroPrMax}{up to 8.5\% fewer errors}  % experiments/FIGS/data/v2/index.json :: binary probe qcd_c_vs_light: (1 - AUC) of FLAV_F1 over FLAV_F0, run by run :: sha256 c6c4138ee5c86423 :: prose twin of the run farther from 1
\newcommand{\MlstCmpQcdCFonerOverFone}{0.9535 and 0.9520}  % experiments/FIGS/data/v2/index.json :: binary probe qcd_c_vs_light: (1 - AUC) of FLAV_F1R over FLAV_F1, run by run :: sha256 c6c4138ee5c86423
\newcommand{\MlstCmpQcdCFonerOverFonePr}{at least 4.6\% fewer errors}  % experiments/FIGS/data/v2/index.json :: binary probe qcd_c_vs_light: (1 - AUC) of FLAV_F1R over FLAV_F1, run by run :: sha256 c6c4138ee5c86423 :: prose twin: two significant figures, direction in words
\newcommand{\MlstCmpQcdCFonerOverFonePrMax}{up to 4.8\% fewer errors}  % experiments/FIGS/data/v2/index.json :: binary probe qcd_c_vs_light: (1 - AUC) of FLAV_F1R over FLAV_F1, run by run :: sha256 c6c4138ee5c86423 :: prose twin of the run farther from 1
\newcommand{\MlstCmpQcdCFonerOverFzero}{0.872 and 0.899}  % experiments/FIGS/data/v2/index.json :: binary probe qcd_c_vs_light: (1 - AUC) of FLAV_F1R over FLAV_F0, run by run :: sha256 c6c4138ee5c86423
\newcommand{\MlstCmpQcdCFonerOverFzeroPr}{at least 10\% fewer errors}  % experiments/FIGS/data/v2/index.json :: binary probe qcd_c_vs_light: (1 - AUC) of FLAV_F1R over FLAV_F0, run by run :: sha256 c6c4138ee5c86423 :: prose twin: two significant figures, direction in words
\newcommand{\MlstCmpQcdCFonerOverFzeroPrMax}{up to 13\% fewer errors}  % experiments/FIGS/data/v2/index.json :: binary probe qcd_c_vs_light: (1 - AUC) of FLAV_F1R over FLAV_F0, run by run :: sha256 c6c4138ee5c86423 :: prose twin of the run farther from 1
\newcommand{\MlstCmpQcdCFzeroOverCoarse}{0.952 and 0.934}  % experiments/FIGS/data/v2/index.json :: binary probe qcd_c_vs_light: (1 - AUC) of FLAV_F0 over R16_Q1, run by run :: sha256 c6c4138ee5c86423
\newcommand{\MlstCmpQcdCFzeroOverCoarsePr}{at least 4.8\% fewer errors}  % experiments/FIGS/data/v2/index.json :: binary probe qcd_c_vs_light: (1 - AUC) of FLAV_F0 over R16_Q1, run by run :: sha256 c6c4138ee5c86423 :: prose twin: two significant figures, direction in words
\newcommand{\MlstCmpQcdCFzeroOverCoarsePrMax}{up to 6.6\% fewer errors}  % experiments/FIGS/data/v2/index.json :: binary probe qcd_c_vs_light: (1 - AUC) of FLAV_F0 over R16_Q1, run by run :: sha256 c6c4138ee5c86423 :: prose twin of the run farther from 1
\newcommand{\MlstCmpQcdCFoneOverCoarse}{0.871 and 0.882}  % experiments/FIGS/data/v2/index.json :: binary probe qcd_c_vs_light: (1 - AUC) of FLAV_F1 over R16_Q1, run by run :: sha256 c6c4138ee5c86423
\newcommand{\MlstCmpQcdCFoneOverCoarsePr}{at least 12\% fewer errors}  % experiments/FIGS/data/v2/index.json :: binary probe qcd_c_vs_light: (1 - AUC) of FLAV_F1 over R16_Q1, run by run :: sha256 c6c4138ee5c86423 :: prose twin: two significant figures, direction in words
\newcommand{\MlstCmpQcdCFoneOverCoarsePrMax}{up to 13\% fewer errors}  % experiments/FIGS/data/v2/index.json :: binary probe qcd_c_vs_light: (1 - AUC) of FLAV_F1 over R16_Q1, run by run :: sha256 c6c4138ee5c86423 :: prose twin of the run farther from 1
\newcommand{\MlstCmpQcdCFonerOverCoarse}{0.8301 and 0.8396}  % experiments/FIGS/data/v2/index.json :: binary probe qcd_c_vs_light: (1 - AUC) of FLAV_F1R over R16_Q1, run by run :: sha256 c6c4138ee5c86423
\newcommand{\MlstCmpQcdCFonerOverCoarsePr}{at least 16\% fewer errors}  % experiments/FIGS/data/v2/index.json :: binary probe qcd_c_vs_light: (1 - AUC) of FLAV_F1R over R16_Q1, run by run :: sha256 c6c4138ee5c86423 :: prose twin: two significant figures, direction in words
\newcommand{\MlstCmpQcdCFonerOverCoarsePrMax}{up to 17\% fewer errors}  % experiments/FIGS/data/v2/index.json :: binary probe qcd_c_vs_light: (1 - AUC) of FLAV_F1R over R16_Q1, run by run :: sha256 c6c4138ee5c86423 :: prose twin of the run farther from 1
\newcommand{\MlstCmpQcdCFzeroOverFull}{1.352 and 1.293}  % experiments/FIGS/data/v2/index.json :: binary probe qcd_c_vs_light: (1 - AUC) of FLAV_F0 over L188, run by run :: sha256 c6c4138ee5c86423
\newcommand{\MlstCmpQcdCFzeroOverFullPr}{at least 29\% more errors}  % experiments/FIGS/data/v2/index.json :: binary probe qcd_c_vs_light: (1 - AUC) of FLAV_F0 over L188, run by run :: sha256 c6c4138ee5c86423 :: prose twin: two significant figures, direction in words
\newcommand{\MlstCmpQcdCFzeroOverFullPrMax}{up to 35\% more errors}  % experiments/FIGS/data/v2/index.json :: binary probe qcd_c_vs_light: (1 - AUC) of FLAV_F0 over L188, run by run :: sha256 c6c4138ee5c86423 :: prose twin of the run farther from 1
\newcommand{\MlstCmpQcdCFoneOverFull}{1.237 and 1.220}  % experiments/FIGS/data/v2/index.json :: binary probe qcd_c_vs_light: (1 - AUC) of FLAV_F1 over L188, run by run :: sha256 c6c4138ee5c86423
\newcommand{\MlstCmpQcdCFoneOverFullPr}{at least 22\% more errors}  % experiments/FIGS/data/v2/index.json :: binary probe qcd_c_vs_light: (1 - AUC) of FLAV_F1 over L188, run by run :: sha256 c6c4138ee5c86423 :: prose twin: two significant figures, direction in words
\newcommand{\MlstCmpQcdCFoneOverFullPrMax}{up to 24\% more errors}  % experiments/FIGS/data/v2/index.json :: binary probe qcd_c_vs_light: (1 - AUC) of FLAV_F1 over L188, run by run :: sha256 c6c4138ee5c86423 :: prose twin of the run farther from 1
\newcommand{\MlstCmpQcdCFonerOverFull}{1.179 and 1.162}  % experiments/FIGS/data/v2/index.json :: binary probe qcd_c_vs_light: (1 - AUC) of FLAV_F1R over L188, run by run :: sha256 c6c4138ee5c86423
\newcommand{\MlstCmpQcdCFonerOverFullPr}{at least 16\% more errors}  % experiments/FIGS/data/v2/index.json :: binary probe qcd_c_vs_light: (1 - AUC) of FLAV_F1R over L188, run by run :: sha256 c6c4138ee5c86423 :: prose twin: two significant figures, direction in words
\newcommand{\MlstCmpQcdCFonerOverFullPrMax}{up to 18\% more errors}  % experiments/FIGS/data/v2/index.json :: binary probe qcd_c_vs_light: (1 - AUC) of FLAV_F1R over L188, run by run :: sha256 c6c4138ee5c86423 :: prose twin of the run farther from 1
\newcommand{\MlstLpEemmAucFull}{1.0000\,\ensuremath{\pm}\,\ensuremath{<}0.0001}  % experiments/FIGS/data/v2/index.json :: linear probe ee_vs_mm (fits_v4_sib): AUC, L188 :: sha256 c6c4138ee5c86423
\newcommand{\MlstLpEemmAucRes}{1.0000\,\ensuremath{\pm}\,\ensuremath{<}0.0001}  % experiments/FIGS/data/v2/index.json :: linear probe ee_vs_mm (fits_v4_sib): AUC, L162 :: sha256 c6c4138ee5c86423
\newcommand{\MlstLpEemmAucMid}{0.9610\,\ensuremath{\pm}\,0.0030}  % experiments/FIGS/data/v2/index.json :: linear probe ee_vs_mm (fits_v4_sib): AUC, R42_Q1 :: sha256 c6c4138ee5c86423
\newcommand{\MlstLpEemmAucCoarse}{0.960\,\ensuremath{\pm}\,0.020}  % experiments/FIGS/data/v2/index.json :: linear probe ee_vs_mm (fits_v4_sib): AUC, R16_Q1 :: sha256 c6c4138ee5c86423
\newcommand{\MlstLpEemmAucFzero}{1.0000 and 1.0000}  % experiments/FIGS/data/v2/index.json :: linear probe ee_vs_mm (fits_v4_sib): AUC, FLAV_F0 :: sha256 c6c4138ee5c86423
\newcommand{\MlstLpEemmAucFone}{1.0000 and 1.0000}  % experiments/FIGS/data/v2/index.json :: linear probe ee_vs_mm (fits_v4_sib): AUC, FLAV_F1 :: sha256 c6c4138ee5c86423
\newcommand{\MlstLpEemmAucFoner}{1.0000 and 1.0000}  % experiments/FIGS/data/v2/index.json :: linear probe ee_vs_mm (fits_v4_sib): AUC, FLAV_F1R :: sha256 c6c4138ee5c86423
\newcommand{\MlstLpEemmAucInit}{0.9995\,\ensuremath{\pm}\,0.0003}  % experiments/FIGS/data/v2/index.json :: linear probe ee_vs_mm (fits_v4_sib): AUC, INIT :: sha256 c6c4138ee5c86423
\newcommand{\MlstLpEemmAucInitBn}{1.0000\,\ensuremath{\pm}\,\ensuremath{<}0.0001}  % experiments/FIGS/data/v2/index.json :: linear probe ee_vs_mm (fits_v4_sib): AUC, INIT_BN :: sha256 c6c4138ee5c86423
\newcommand{\MlstLpEemmAucRandMerged}{1.0000\,\ensuremath{\pm}\,\ensuremath{<}0.0001}  % experiments/FIGS/data/v2/index.json :: linear probe ee_vs_mm (fits_v4_sib): AUC, the random partitions that merge the pair :: sha256 c6c4138ee5c86423
\newcommand{\MlstPairEemmNTrain}{9{,}571}  % experiments/FIGS/data/v2/index.json :: fits_v4_sib ee_vs_mm: n_train :: sha256 c6c4138ee5c86423
\newcommand{\MlstPairEemmNVal}{1{,}068}  % experiments/FIGS/data/v2/index.json :: fits_v4_sib ee_vs_mm: n_val :: sha256 c6c4138ee5c86423
\newcommand{\MlstPairEemmNTestPerSide}{128{,}833--135{,}245}  % experiments/FIGS/data/v2/index.json :: fits_v4_sib ee_vs_mm: n_test_signal and n_test_background :: sha256 c6c4138ee5c86423
\newcommand{\MlstLpJcTwoAucFull}{0.9801\,\ensuremath{\pm}\,\ensuremath{<}0.0001}  % experiments/FIGS/data/v2/index.json :: linear probe JetClass-II N=1000000: AUC averaged over the classes, L188 :: sha256 c6c4138ee5c86423
\newcommand{\MlstLpJcTwoAccFull}{0.4348\,\ensuremath{\pm}\,0.0003}  % experiments/FIGS/data/v2/index.json :: linear probe JetClass-II N=1000000: accuracy, L188 :: sha256 c6c4138ee5c86423
\newcommand{\MlstLpJcTwoAucRes}{0.9796\,\ensuremath{\pm}\,0.0001}  % experiments/FIGS/data/v2/index.json :: linear probe JetClass-II N=1000000: AUC averaged over the classes, L162 :: sha256 c6c4138ee5c86423
\newcommand{\MlstLpJcTwoAccRes}{0.4328\,\ensuremath{\pm}\,0.0008}  % experiments/FIGS/data/v2/index.json :: linear probe JetClass-II N=1000000: accuracy, L162 :: sha256 c6c4138ee5c86423
\newcommand{\MlstLpJcTwoAucMid}{0.9756\,\ensuremath{\pm}\,0.0002}  % experiments/FIGS/data/v2/index.json :: linear probe JetClass-II N=1000000: AUC averaged over the classes, R42_Q1 :: sha256 c6c4138ee5c86423
\newcommand{\MlstLpJcTwoAccMid}{0.361\,\ensuremath{\pm}\,0.005}  % experiments/FIGS/data/v2/index.json :: linear probe JetClass-II N=1000000: accuracy, R42_Q1 :: sha256 c6c4138ee5c86423
\newcommand{\MlstLpJcTwoAucCoarse}{0.9691\,\ensuremath{\pm}\,0.0004}  % experiments/FIGS/data/v2/index.json :: linear probe JetClass-II N=1000000: AUC averaged over the classes, R16_Q1 :: sha256 c6c4138ee5c86423
\newcommand{\MlstLpJcTwoAccCoarse}{0.334\,\ensuremath{\pm}\,0.007}  % experiments/FIGS/data/v2/index.json :: linear probe JetClass-II N=1000000: accuracy, R16_Q1 :: sha256 c6c4138ee5c86423
\newcommand{\MlstLpJcTwoAucResMass}{0.9791 and 0.9794}  % experiments/FIGS/data/v2/index.json :: linear probe JetClass-II N=1000000: AUC averaged over the classes, L162_MASS :: sha256 c6c4138ee5c86423
\newcommand{\MlstLpJcTwoAccResMass}{0.4271 and 0.4305}  % experiments/FIGS/data/v2/index.json :: linear probe JetClass-II N=1000000: accuracy, L162_MASS :: sha256 c6c4138ee5c86423
\newcommand{\MlstLpJcTwoAucCoarseMass}{0.9679 and 0.9696}  % experiments/FIGS/data/v2/index.json :: linear probe JetClass-II N=1000000: AUC averaged over the classes, R16_Q1_MASS :: sha256 c6c4138ee5c86423
\newcommand{\MlstLpJcTwoAccCoarseMass}{0.312 and 0.350}  % experiments/FIGS/data/v2/index.json :: linear probe JetClass-II N=1000000: accuracy, R16_Q1_MASS :: sha256 c6c4138ee5c86423
\newcommand{\MlstLpJcTwoAucFzero}{0.9717 and 0.9710}  % experiments/FIGS/data/v2/index.json :: linear probe JetClass-II N=1000000: AUC averaged over the classes, FLAV_F0 :: sha256 c6c4138ee5c86423
\newcommand{\MlstLpJcTwoAccFzero}{0.3769 and 0.3723}  % experiments/FIGS/data/v2/index.json :: linear probe JetClass-II N=1000000: accuracy, FLAV_F0 :: sha256 c6c4138ee5c86423
\newcommand{\MlstLpJcTwoAucFone}{0.9758 and 0.9759}  % experiments/FIGS/data/v2/index.json :: linear probe JetClass-II N=1000000: AUC averaged over the classes, FLAV_F1 :: sha256 c6c4138ee5c86423
\newcommand{\MlstLpJcTwoAccFone}{0.3945 and 0.3963}  % experiments/FIGS/data/v2/index.json :: linear probe JetClass-II N=1000000: accuracy, FLAV_F1 :: sha256 c6c4138ee5c86423
\newcommand{\MlstLpJcTwoAucFoner}{0.9760 and 0.9760}  % experiments/FIGS/data/v2/index.json :: linear probe JetClass-II N=1000000: AUC averaged over the classes, FLAV_F1R :: sha256 c6c4138ee5c86423
\newcommand{\MlstLpJcTwoAccFoner}{0.4071 and 0.4089}  % experiments/FIGS/data/v2/index.json :: linear probe JetClass-II N=1000000: accuracy, FLAV_F1R :: sha256 c6c4138ee5c86423
\newcommand{\MlstLpJcTwoAucInit}{0.8121\,\ensuremath{\pm}\,0.0019}  % experiments/FIGS/data/v2/index.json :: linear probe JetClass-II N=1000000: AUC averaged over the classes, INIT :: sha256 c6c4138ee5c86423
\newcommand{\MlstLpJcTwoAccInit}{0.0724\,\ensuremath{\pm}\,0.0015}  % experiments/FIGS/data/v2/index.json :: linear probe JetClass-II N=1000000: accuracy, INIT :: sha256 c6c4138ee5c86423
\newcommand{\MlstLpJcTwoAucInitBn}{0.866\,\ensuremath{\pm}\,0.004}  % experiments/FIGS/data/v2/index.json :: linear probe JetClass-II N=1000000: AUC averaged over the classes, INIT_BN :: sha256 c6c4138ee5c86423
\newcommand{\MlstLpJcTwoAccInitBn}{0.1017\,\ensuremath{\pm}\,0.0017}  % experiments/FIGS/data/v2/index.json :: linear probe JetClass-II N=1000000: accuracy, INIT_BN :: sha256 c6c4138ee5c86423
\newcommand{\MlstLpJcTwoAucRand}{0.9777\,\ensuremath{\pm}\,0.0002}  % experiments/FIGS/data/v2/index.json :: linear probe JetClass-II N=1000000: AUC averaged over the classes, the random partitions :: sha256 c6c4138ee5c86423
\newcommand{\MlstLpJcOneAccInitBn}{0.475\,\ensuremath{\pm}\,0.005}  % experiments/FIGS/data/v2/index.json :: linear probe JetClass N=1000000: accuracy, INIT_BN :: sha256 c6c4138ee5c86423
\newcommand{\MlstBnMaxDev}{0.2\%}  % experiments/FIGS/data/v2/index.json :: largest |r(best70_bn) / r(best70) - 1| of a JetClass-II ratio of 1 - AUC to the 188-class run, per cent :: sha256 c6c4138ee5c86423
\newcommand{\MlstBnMaxDevPair}{0.7\%}  % experiments/FIGS/data/v2/index.json :: largest |r(best70_bn) / r(best70) - 1| of a ratio of 1 - AUC on a class pair to the 188-class run, levels of the hierarchy, mass models and controlled partitions, per cent :: sha256 c6c4138ee5c86423 :: pairs with a check: bc_vs_bq_cs, bb_vs_cc, bbqq_vs_ccqq, qcd_bb_vs_cc, qcd_b_vs_light, qcd_c_vs_light; without: ee_vs_mm
\newcommand{\MlstGpuTypeBound}{1.4\%}  % experiments/FIGS/data/v2/index.json :: largest |label-set ratio of 1 - AUC of index 3 / mean of the other indices - 1|, JetClass-II largest set, per cent :: sha256 c6c4138ee5c86423 :: mtx-l188-s3 trained on NVIDIA GeForce RTX 3090; largest for R16_Q1; label sets L162 0.09%, R42_Q1 0.70%, R16_Q1 1.37%
\newcommand{\MlstGpuOtherSeedsDiff}{1.9\%}  % experiments/FIGS/data/v2/index.json :: largest (max - min) / mean of a label-set ratio of 1 - AUC over the indices other than 3, JetClass-II largest set, per cent :: sha256 c6c4138ee5c86423 :: label sets L162 0.90%, R42_Q1 1.14%, R16_Q1 1.93%
\newcommand{\MlstPerClassResNDoubled}{0}  % experiments/FIGS/data/v2/index.json :: JetClass-II classes whose mean ratio of 1 - AUC (L162 over 188) is at least 2 (count) :: sha256 c6c4138ee5c86423
\newcommand{\MlstPerClassResSeedSdMedian}{0.011}  % experiments/FIGS/data/v2/index.json :: median over classes of the SD over runs of the per-class ratio, L162 over 188 :: sha256 c6c4138ee5c86423
\newcommand{\MlstPerClassResTopClass}{QCD, \ensuremath{ccss}}  % experiments/FIGS/data/v2/index.json :: class with the largest mean ratio, L162 :: sha256 c6c4138ee5c86423
\newcommand{\MlstPerClassResTopRatio}{1.129\,\ensuremath{\pm}\,0.010}  % experiments/FIGS/data/v2/index.json :: per-class ratio of 1 - AUC of that class, L162 over 188, run by run :: sha256 c6c4138ee5c86423
\newcommand{\MlstPerClassResTopRatioPr}{13\% more errors}  % experiments/FIGS/data/v2/index.json :: per-class ratio of 1 - AUC of that class, L162 over 188, run by run :: sha256 c6c4138ee5c86423 :: prose twin: two significant figures, direction in words
\newcommand{\MlstPerClassResLeptonN}{0}  % experiments/FIGS/data/v2/index.json :: resonant classes sharing their L162 group with a decay that differs only in lepton flavor (count) :: sha256 c6c4138ee5c86423
\newcommand{\MlstPerClassResLeptonMedian}{--}  % experiments/FIGS/data/v2/index.json :: median over those classes of the mean per-class ratio, L162 over 188 :: sha256 c6c4138ee5c86423
\newcommand{\MlstPerClassResLeptonMedianPr}{\pending{no run yet}}  % experiments/FIGS/data/v2/index.json :: median over those classes of the mean per-class ratio, L162 over 188 :: sha256 c6c4138ee5c86423 :: prose twin: two significant figures, direction in words
\newcommand{\MlstPerClassResLeptonShare}{--}  % experiments/FIGS/data/v2/index.json :: their share of the summed ln ratio over the resonant classes, L162 :: sha256 c6c4138ee5c86423
\newcommand{\MlstPerClassResQuarkN}{0}  % experiments/FIGS/data/v2/index.json :: resonant classes sharing their L162 group with a decay that differs only in quark flavor (count) :: sha256 c6c4138ee5c86423
\newcommand{\MlstPerClassResQuarkMedian}{--}  % experiments/FIGS/data/v2/index.json :: median over those classes of the mean per-class ratio, L162 over 188 :: sha256 c6c4138ee5c86423
\newcommand{\MlstPerClassResQuarkMedianPr}{\pending{no run yet}}  % experiments/FIGS/data/v2/index.json :: median over those classes of the mean per-class ratio, L162 over 188 :: sha256 c6c4138ee5c86423 :: prose twin: two significant figures, direction in words
\newcommand{\MlstPerClassResQuarkShare}{--}  % experiments/FIGS/data/v2/index.json :: their share of the summed ln ratio over the resonant classes, L162 :: sha256 c6c4138ee5c86423
\newcommand{\MlstPerClassResQuarkLeptonN}{0}  % experiments/FIGS/data/v2/index.json :: resonant classes counted in both the quark- and the lepton-flavor set, L162 (count) :: sha256 c6c4138ee5c86423
\newcommand{\MlstPerClassMidNDoubled}{40}  % experiments/FIGS/data/v2/index.json :: JetClass-II classes whose mean ratio of 1 - AUC (R42_Q1 over 188) is at least 2 (count) :: sha256 c6c4138ee5c86423
\newcommand{\MlstPerClassMidSeedSdMedian}{0.023}  % experiments/FIGS/data/v2/index.json :: median over classes of the SD over runs of the per-class ratio, R42_Q1 over 188 :: sha256 c6c4138ee5c86423
\newcommand{\MlstPerClassMidTopClass}{\ensuremath{X\to YY\to bb\mu\mu}}  % experiments/FIGS/data/v2/index.json :: class with the largest mean ratio, R42_Q1 :: sha256 c6c4138ee5c86423
\newcommand{\MlstPerClassMidTopRatio}{4.06\,\ensuremath{\pm}\,0.14}  % experiments/FIGS/data/v2/index.json :: per-class ratio of 1 - AUC of that class, R42_Q1 over 188, run by run :: sha256 c6c4138ee5c86423
\newcommand{\MlstPerClassMidTopRatioPr}{4.1 times the errors}  % experiments/FIGS/data/v2/index.json :: per-class ratio of 1 - AUC of that class, R42_Q1 over 188, run by run :: sha256 c6c4138ee5c86423 :: prose twin: two significant figures, direction in words
\newcommand{\MlstPerClassMidLeptonN}{78}  % experiments/FIGS/data/v2/index.json :: resonant classes sharing their R42_Q1 group with a decay that differs only in lepton flavor (count) :: sha256 c6c4138ee5c86423
\newcommand{\MlstPerClassMidLeptonMedian}{2.02}  % experiments/FIGS/data/v2/index.json :: median over those classes of the mean per-class ratio, R42_Q1 over 188 :: sha256 c6c4138ee5c86423
\newcommand{\MlstPerClassMidLeptonMedianPr}{2.0 times the errors}  % experiments/FIGS/data/v2/index.json :: median over those classes of the mean per-class ratio, R42_Q1 over 188 :: sha256 c6c4138ee5c86423 :: prose twin: two significant figures, direction in words
\newcommand{\MlstPerClassMidLeptonShare}{90\%}  % experiments/FIGS/data/v2/index.json :: their share of the summed ln ratio over the resonant classes, R42_Q1 :: sha256 c6c4138ee5c86423
\newcommand{\MlstPerClassMidQuarkN}{111}  % experiments/FIGS/data/v2/index.json :: resonant classes sharing their R42_Q1 group with a decay that differs only in quark flavor (count) :: sha256 c6c4138ee5c86423
\newcommand{\MlstPerClassMidQuarkMedian}{1.14}  % experiments/FIGS/data/v2/index.json :: median over those classes of the mean per-class ratio, R42_Q1 over 188 :: sha256 c6c4138ee5c86423
\newcommand{\MlstPerClassMidQuarkMedianPr}{14\% more errors}  % experiments/FIGS/data/v2/index.json :: median over those classes of the mean per-class ratio, R42_Q1 over 188 :: sha256 c6c4138ee5c86423 :: prose twin: two significant figures, direction in words
\newcommand{\MlstPerClassMidQuarkShare}{53\%}  % experiments/FIGS/data/v2/index.json :: their share of the summed ln ratio over the resonant classes, R42_Q1 :: sha256 c6c4138ee5c86423
\newcommand{\MlstPerClassMidQuarkLeptonN}{44}  % experiments/FIGS/data/v2/index.json :: resonant classes counted in both the quark- and the lepton-flavor set, R42_Q1 (count) :: sha256 c6c4138ee5c86423
\newcommand{\MlstPerClassCoarseNDoubled}{60}  % experiments/FIGS/data/v2/index.json :: JetClass-II classes whose mean ratio of 1 - AUC (R16_Q1 over 188) is at least 2 (count) :: sha256 c6c4138ee5c86423
\newcommand{\MlstPerClassCoarseSeedSdMedian}{0.055}  % experiments/FIGS/data/v2/index.json :: median over classes of the SD over runs of the per-class ratio, R16_Q1 over 188 :: sha256 c6c4138ee5c86423
\newcommand{\MlstPerClassCoarseTopClass}{\ensuremath{X\to YY\to bbee}}  % experiments/FIGS/data/v2/index.json :: class with the largest mean ratio, R16_Q1 :: sha256 c6c4138ee5c86423
\newcommand{\MlstPerClassCoarseTopRatio}{5.2\,\ensuremath{\pm}\,0.4}  % experiments/FIGS/data/v2/index.json :: per-class ratio of 1 - AUC of that class, R16_Q1 over 188, run by run :: sha256 c6c4138ee5c86423
\newcommand{\MlstPerClassCoarseTopRatioPr}{5.2 times the errors}  % experiments/FIGS/data/v2/index.json :: per-class ratio of 1 - AUC of that class, R16_Q1 over 188, run by run :: sha256 c6c4138ee5c86423 :: prose twin: two significant figures, direction in words
\newcommand{\MlstPerClassCoarseLeptonN}{78}  % experiments/FIGS/data/v2/index.json :: resonant classes sharing their R16_Q1 group with a decay that differs only in lepton flavor (count) :: sha256 c6c4138ee5c86423
\newcommand{\MlstPerClassCoarseLeptonMedian}{2.27}  % experiments/FIGS/data/v2/index.json :: median over those classes of the mean per-class ratio, R16_Q1 over 188 :: sha256 c6c4138ee5c86423
\newcommand{\MlstPerClassCoarseLeptonMedianPr}{2.3 times the errors}  % experiments/FIGS/data/v2/index.json :: median over those classes of the mean per-class ratio, R16_Q1 over 188 :: sha256 c6c4138ee5c86423 :: prose twin: two significant figures, direction in words
\newcommand{\MlstPerClassCoarseLeptonShare}{68\%}  % experiments/FIGS/data/v2/index.json :: their share of the summed ln ratio over the resonant classes, R16_Q1 :: sha256 c6c4138ee5c86423
\newcommand{\MlstPerClassCoarseQuarkN}{141}  % experiments/FIGS/data/v2/index.json :: resonant classes sharing their R16_Q1 group with a decay that differs only in quark flavor (count) :: sha256 c6c4138ee5c86423
\newcommand{\MlstPerClassCoarseQuarkMedian}{1.65}  % experiments/FIGS/data/v2/index.json :: median over those classes of the mean per-class ratio, R16_Q1 over 188 :: sha256 c6c4138ee5c86423
\newcommand{\MlstPerClassCoarseQuarkMedianPr}{65\% more errors}  % experiments/FIGS/data/v2/index.json :: median over those classes of the mean per-class ratio, R16_Q1 over 188 :: sha256 c6c4138ee5c86423 :: prose twin: two significant figures, direction in words
\newcommand{\MlstPerClassCoarseQuarkShare}{87\%}  % experiments/FIGS/data/v2/index.json :: their share of the summed ln ratio over the resonant classes, R16_Q1 :: sha256 c6c4138ee5c86423
\newcommand{\MlstPerClassCoarseQuarkLeptonN}{64}  % experiments/FIGS/data/v2/index.json :: resonant classes counted in both the quark- and the lepton-flavor set, R16_Q1 (count) :: sha256 c6c4138ee5c86423
\newcommand{\MlstPerClassRandFlavMergedOverApart}{1.010\,\ensuremath{\pm}\,0.008}  % experiments/FIGS/data/v2/index.json :: random partitions, linear probe: median per-class ratio of the classes merged with a decay that differs only in quark flavor over that of the classes kept apart from all of them, mean +- SD over the partitions :: sha256 c6c4138ee5c86423
\newcommand{\MlstPerClassRandFlavMergedOverApartPr}{equal within 1\%}  % experiments/FIGS/data/v2/index.json :: random partitions, linear probe: median per-class ratio of the classes merged with a decay that differs only in quark flavor over that of the classes kept apart from all of them, mean +- SD over the partitions :: sha256 c6c4138ee5c86423 :: prose twin: two significant figures, direction in words
\newcommand{\MlstPerClassRandLeptonMergedOverApart}{0.997\,\ensuremath{\pm}\,0.020}  % experiments/FIGS/data/v2/index.json :: random partitions, linear probe: median per-class ratio of the classes merged with a decay that differs only in lepton flavor over that of the classes kept apart from all of them, mean +- SD over the partitions :: sha256 c6c4138ee5c86423
\newcommand{\MlstPerClassRandLeptonMergedOverApartPr}{equal within 3\%}  % experiments/FIGS/data/v2/index.json :: random partitions, linear probe: median per-class ratio of the classes merged with a decay that differs only in lepton flavor over that of the classes kept apart from all of them, mean +- SD over the partitions :: sha256 c6c4138ee5c86423 :: prose twin: two significant figures, direction in words
\newcommand{\MlstMlpNatPerClassRandFlavMergedOverApart}{1.003\,\ensuremath{\pm}\,0.007}  % experiments/FIGS/data/v2/index.json :: random partitions, nonlinear probe, natural composition: median per-class ratio of the classes merged with a decay that differs only in quark flavor over that of the classes kept apart from all of them, mean +- SD over the partitions :: sha256 c6c4138ee5c86423
\newcommand{\MlstMlpNatPerClassRandFlavMergedOverApartPr}{equal within 1\%}  % experiments/FIGS/data/v2/index.json :: random partitions, nonlinear probe, natural composition: median per-class ratio of the classes merged with a decay that differs only in quark flavor over that of the classes kept apart from all of them, mean +- SD over the partitions :: sha256 c6c4138ee5c86423 :: prose twin: two significant figures, direction in words
\newcommand{\MlstMlpNatPerClassRandLeptonMergedOverApart}{1.01\,\ensuremath{\pm}\,0.07}  % experiments/FIGS/data/v2/index.json :: random partitions, nonlinear probe, natural composition: median per-class ratio of the classes merged with a decay that differs only in lepton flavor over that of the classes kept apart from all of them, mean +- SD over the partitions :: sha256 c6c4138ee5c86423
\newcommand{\MlstMlpNatPerClassRandLeptonMergedOverApartPr}{equal within 8\%}  % experiments/FIGS/data/v2/index.json :: random partitions, nonlinear probe, natural composition: median per-class ratio of the classes merged with a decay that differs only in lepton flavor over that of the classes kept apart from all of them, mean +- SD over the partitions :: sha256 c6c4138ee5c86423 :: prose twin: two significant figures, direction in words
\newcommand{\MlstPerClassAucDiffResSeedSdMedian}{0.015}  % experiments/FIGS/data/v2/index.json :: median over classes of the SD over indices of 100 (AUC_c of L162 - AUC_c of 188), pp :: sha256 c6c4138ee5c86423
\newcommand{\MlstPerClassAucDiffResMin}{\ensuremath{-}0.73}  % experiments/FIGS/data/v2/index.json :: most negative mean over indices of 100 (AUC_c of L162 - AUC_c of 188), pp :: sha256 c6c4138ee5c86423 :: label_QCD_bcss
\newcommand{\MlstPerClassAucDiffResMinClass}{QCD, \ensuremath{bcss}}  % experiments/FIGS/data/v2/index.json :: its class, L162 :: sha256 c6c4138ee5c86423
\newcommand{\MlstPerClassAucDiffMidSeedSdMedian}{0.032}  % experiments/FIGS/data/v2/index.json :: median over classes of the SD over indices of 100 (AUC_c of R42_Q1 - AUC_c of 188), pp :: sha256 c6c4138ee5c86423
\newcommand{\MlstPerClassAucDiffMidMin}{\ensuremath{-}1.4}  % experiments/FIGS/data/v2/index.json :: most negative mean over indices of 100 (AUC_c of R42_Q1 - AUC_c of 188), pp :: sha256 c6c4138ee5c86423 :: label_X_YY_cqtauev
\newcommand{\MlstPerClassAucDiffMidMinClass}{\ensuremath{X\to YY\to cq\tau_e\nu}}  % experiments/FIGS/data/v2/index.json :: its class, R42_Q1 :: sha256 c6c4138ee5c86423
\newcommand{\MlstPerClassAucDiffCoarseSeedSdMedian}{0.077}  % experiments/FIGS/data/v2/index.json :: median over classes of the SD over indices of 100 (AUC_c of R16_Q1 - AUC_c of 188), pp :: sha256 c6c4138ee5c86423
\newcommand{\MlstPerClassAucDiffCoarseMin}{\ensuremath{-}3.1}  % experiments/FIGS/data/v2/index.json :: most negative mean over indices of 100 (AUC_c of R16_Q1 - AUC_c of 188), pp :: sha256 c6c4138ee5c86423 :: label_QCD_bcss
\newcommand{\MlstPerClassAucDiffCoarseMinClass}{QCD, \ensuremath{bcss}}  % experiments/FIGS/data/v2/index.json :: its class, R16_Q1 :: sha256 c6c4138ee5c86423
\newcommand{\MlstRocJcTwoHurtList}{\ensuremath{X\to YY\to ccq}; \ensuremath{X\to YY\to ssb}; \ensuremath{X\to YY\to ccss}; QCD, \ensuremath{bcc}; QCD, \ensuremath{bccss}; QCD, \ensuremath{bss}; QCD, \ensuremath{ccs} and QCD, \ensuremath{ccss}}  % experiments/FIGS/data/v2/index.json :: jc2_roc_classes on the class-balanced probes; candidates with at least 1000 test jets :: sha256 c6c4138ee5c86423 :: label_X_YY_ccq, label_X_YY_ssb, label_X_YY_ccss, label_QCD_bcc, label_QCD_bccss, label_QCD_bss, label_QCD_ccs, label_QCD_ccss
\newcommand{\MlstRocJcTwoStableList}{\ensuremath{X\to gg}; \ensuremath{X\to YY\to ggg}; \ensuremath{X\to YY\to g\tau_h\tau_h} and QCD, light}  % experiments/FIGS/data/v2/index.json :: jc2_roc_classes on the class-balanced probes; candidates with at least 1000 test jets :: sha256 c6c4138ee5c86423 :: label_X_gg, label_X_YY_ggg, label_X_YY_gtauhtauh, label_QCD_light
\newcommand{\MlstRocJcTwoHurtMinDrop}{2.1}  % experiments/FIGS/data/v2/index.json :: jc2_roc_classes on the class-balanced probes; candidates with at least 1000 test jets :: sha256 c6c4138ee5c86423 :: smallest AUC loss at 17 classes of the hurt, pp
\newcommand{\MlstRocJcTwoHurtMaxDrop}{2.8}  % experiments/FIGS/data/v2/index.json :: jc2_roc_classes on the class-balanced probes; candidates with at least 1000 test jets :: sha256 c6c4138ee5c86423 :: largest AUC loss at 17 classes of the hurt, pp
\newcommand{\MlstRocJcTwoStableMaxChange}{15}  % experiments/FIGS/data/v2/index.json :: jc2_roc_classes on the class-balanced probes; candidates with at least 1000 test jets :: sha256 c6c4138ee5c86423 :: largest change of 1 - AUC at any label set of the stable classes, per cent of the 188-class value
\newcommand{\MlstJcBalNThin}{4}  % experiments/FIGS/data/v2/jc2bal/census.json :: counts.ft_test :: sha256 a4379ee0817254cf :: classes with fewer than 1000 class-balanced test jets (F3 hatching)
\newcommand{\MlstMinBkgPass}{20}  % experiments/FIGS/mlst_v4.py :: JC1_MIN_BKG :: sha256 fa536b4866ffd651 :: fewest background jets behind a drawn ROC point or a quoted rejection
\newcommand{\MlstMassMinJets}{100}  % experiments/EVAL/atlas_resolution.py :: MIN_JETS :: sha256 b82cc219f79f0f14 :: jets below which a mass bin is null
\newcommand{\MlstRocCutBbFull}{0.79}  % experiments/FIGS/data/v2/index.json :: curves_v2/pairs curve bb_vs_cc: largest cut.eps_s over the L188 runs ('not reached' when a run's cut is null or leaves fewer than 2 grid points) :: sha256 c6c4138ee5c86423
\newcommand{\MlstRocCutBbRes}{0.77}  % experiments/FIGS/data/v2/index.json :: curves_v2/pairs curve bb_vs_cc: largest cut.eps_s over the L162 runs ('not reached' when a run's cut is null or leaves fewer than 2 grid points) :: sha256 c6c4138ee5c86423
\newcommand{\MlstRocCutBbMid}{0.79}  % experiments/FIGS/data/v2/index.json :: curves_v2/pairs curve bb_vs_cc: largest cut.eps_s over the R42_Q1 runs ('not reached' when a run's cut is null or leaves fewer than 2 grid points) :: sha256 c6c4138ee5c86423
\newcommand{\MlstRocCutBbCoarse}{0.13}  % experiments/FIGS/data/v2/index.json :: curves_v2/pairs curve bb_vs_cc: largest cut.eps_s over the R16_Q1 runs ('not reached' when a run's cut is null or leaves fewer than 2 grid points) :: sha256 c6c4138ee5c86423
\newcommand{\MlstRocCutBbFzero}{0.03}  % experiments/FIGS/data/v2/index.json :: curves_v2/pairs curve bb_vs_cc: largest cut.eps_s over the FLAV_F0 runs ('not reached' when a run's cut is null or leaves fewer than 2 grid points) :: sha256 c6c4138ee5c86423
\newcommand{\MlstRocCutBbFone}{0.62}  % experiments/FIGS/data/v2/index.json :: curves_v2/pairs curve bb_vs_cc: largest cut.eps_s over the FLAV_F1 runs ('not reached' when a run's cut is null or leaves fewer than 2 grid points) :: sha256 c6c4138ee5c86423
\newcommand{\MlstRocCutBbFoner}{0.51}  % experiments/FIGS/data/v2/index.json :: curves_v2/pairs curve bb_vs_cc: largest cut.eps_s over the FLAV_F1R runs ('not reached' when a run's cut is null or leaves fewer than 2 grid points) :: sha256 c6c4138ee5c86423
\newcommand{\MlstRocNBkgBb}{200{,}000}  % experiments/FIGS/data/v2/index.json :: curves_v2/pairs curve bb_vs_cc: n_background :: sha256 c6c4138ee5c86423
\newcommand{\MlstRocRejMaxBb}{10{,}000}  % experiments/FIGS/data/v2/index.json :: curves_v2/pairs curve bb_vs_cc: n_background / 20 :: sha256 c6c4138ee5c86423
\newcommand{\MlstRocCutBbqqFull}{0.67}  % experiments/FIGS/data/v2/index.json :: curves_v2/pairs curve bbqq_vs_ccqq: largest cut.eps_s over the L188 runs ('not reached' when a run's cut is null or leaves fewer than 2 grid points) :: sha256 c6c4138ee5c86423
\newcommand{\MlstRocCutBbqqRes}{0.66}  % experiments/FIGS/data/v2/index.json :: curves_v2/pairs curve bbqq_vs_ccqq: largest cut.eps_s over the L162 runs ('not reached' when a run's cut is null or leaves fewer than 2 grid points) :: sha256 c6c4138ee5c86423
\newcommand{\MlstRocCutBbqqMid}{0.62}  % experiments/FIGS/data/v2/index.json :: curves_v2/pairs curve bbqq_vs_ccqq: largest cut.eps_s over the R42_Q1 runs ('not reached' when a run's cut is null or leaves fewer than 2 grid points) :: sha256 c6c4138ee5c86423
\newcommand{\MlstRocCutBbqqCoarse}{0.06}  % experiments/FIGS/data/v2/index.json :: curves_v2/pairs curve bbqq_vs_ccqq: largest cut.eps_s over the R16_Q1 runs ('not reached' when a run's cut is null or leaves fewer than 2 grid points) :: sha256 c6c4138ee5c86423
\newcommand{\MlstRocCutBbqqFzero}{0.06}  % experiments/FIGS/data/v2/index.json :: curves_v2/pairs curve bbqq_vs_ccqq: largest cut.eps_s over the FLAV_F0 runs ('not reached' when a run's cut is null or leaves fewer than 2 grid points) :: sha256 c6c4138ee5c86423
\newcommand{\MlstRocCutBbqqFone}{0.56}  % experiments/FIGS/data/v2/index.json :: curves_v2/pairs curve bbqq_vs_ccqq: largest cut.eps_s over the FLAV_F1 runs ('not reached' when a run's cut is null or leaves fewer than 2 grid points) :: sha256 c6c4138ee5c86423
\newcommand{\MlstRocCutBbqqFoner}{0.32}  % experiments/FIGS/data/v2/index.json :: curves_v2/pairs curve bbqq_vs_ccqq: largest cut.eps_s over the FLAV_F1R runs ('not reached' when a run's cut is null or leaves fewer than 2 grid points) :: sha256 c6c4138ee5c86423
\newcommand{\MlstRocNBkgBbqq}{200{,}000}  % experiments/FIGS/data/v2/index.json :: curves_v2/pairs curve bbqq_vs_ccqq: n_background :: sha256 c6c4138ee5c86423
\newcommand{\MlstRocRejMaxBbqq}{10{,}000}  % experiments/FIGS/data/v2/index.json :: curves_v2/pairs curve bbqq_vs_ccqq: n_background / 20 :: sha256 c6c4138ee5c86423
\newcommand{\MlstRocCutBcFull}{0.24}  % experiments/FIGS/data/v2/index.json :: curves_v2/pairs curve bc_vs_bq_cs: largest cut.eps_s over the L188 runs ('not reached' when a run's cut is null or leaves fewer than 2 grid points) :: sha256 c6c4138ee5c86423
\newcommand{\MlstRocCutBcRes}{0.25}  % experiments/FIGS/data/v2/index.json :: curves_v2/pairs curve bc_vs_bq_cs: largest cut.eps_s over the L162 runs ('not reached' when a run's cut is null or leaves fewer than 2 grid points) :: sha256 c6c4138ee5c86423
\newcommand{\MlstRocCutBcMid}{0.14}  % experiments/FIGS/data/v2/index.json :: curves_v2/pairs curve bc_vs_bq_cs: largest cut.eps_s over the R42_Q1 runs ('not reached' when a run's cut is null or leaves fewer than 2 grid points) :: sha256 c6c4138ee5c86423
\newcommand{\MlstRocCutBcCoarse}{0.00}  % experiments/FIGS/data/v2/index.json :: curves_v2/pairs curve bc_vs_bq_cs: largest cut.eps_s over the R16_Q1 runs ('not reached' when a run's cut is null or leaves fewer than 2 grid points) :: sha256 c6c4138ee5c86423
\newcommand{\MlstRocCutBcFzero}{0.00}  % experiments/FIGS/data/v2/index.json :: curves_v2/pairs curve bc_vs_bq_cs: largest cut.eps_s over the FLAV_F0 runs ('not reached' when a run's cut is null or leaves fewer than 2 grid points) :: sha256 c6c4138ee5c86423
\newcommand{\MlstRocCutBcFone}{0.01}  % experiments/FIGS/data/v2/index.json :: curves_v2/pairs curve bc_vs_bq_cs: largest cut.eps_s over the FLAV_F1 runs ('not reached' when a run's cut is null or leaves fewer than 2 grid points) :: sha256 c6c4138ee5c86423
\newcommand{\MlstRocCutBcFoner}{0.01}  % experiments/FIGS/data/v2/index.json :: curves_v2/pairs curve bc_vs_bq_cs: largest cut.eps_s over the FLAV_F1R runs ('not reached' when a run's cut is null or leaves fewer than 2 grid points) :: sha256 c6c4138ee5c86423
\newcommand{\MlstRocNBkgBc}{200{,}000}  % experiments/FIGS/data/v2/index.json :: curves_v2/pairs curve bc_vs_bq_cs: n_background :: sha256 c6c4138ee5c86423
\newcommand{\MlstRocRejMaxBc}{10{,}000}  % experiments/FIGS/data/v2/index.json :: curves_v2/pairs curve bc_vs_bq_cs: n_background / 20 :: sha256 c6c4138ee5c86423
\newcommand{\MlstRocCutEemmFull}{not reached}  % experiments/FIGS/data/v2/index.json :: curves_v2/pairs curve ee_vs_mm: largest cut.eps_s over the L188 runs ('not reached' when a run's cut is null or leaves fewer than 2 grid points) :: sha256 c6c4138ee5c86423
\newcommand{\MlstRocCutEemmRes}{not reached}  % experiments/FIGS/data/v2/index.json :: curves_v2/pairs curve ee_vs_mm: largest cut.eps_s over the L162 runs ('not reached' when a run's cut is null or leaves fewer than 2 grid points) :: sha256 c6c4138ee5c86423
\newcommand{\MlstRocCutEemmMid}{0.00}  % experiments/FIGS/data/v2/index.json :: curves_v2/pairs curve ee_vs_mm: largest cut.eps_s over the R42_Q1 runs ('not reached' when a run's cut is null or leaves fewer than 2 grid points) :: sha256 c6c4138ee5c86423
\newcommand{\MlstRocCutEemmCoarse}{0.00}  % experiments/FIGS/data/v2/index.json :: curves_v2/pairs curve ee_vs_mm: largest cut.eps_s over the R16_Q1 runs ('not reached' when a run's cut is null or leaves fewer than 2 grid points) :: sha256 c6c4138ee5c86423
\newcommand{\MlstRocCutEemmFzero}{not reached}  % experiments/FIGS/data/v2/index.json :: curves_v2/pairs curve ee_vs_mm: largest cut.eps_s over the FLAV_F0 runs ('not reached' when a run's cut is null or leaves fewer than 2 grid points) :: sha256 c6c4138ee5c86423
\newcommand{\MlstRocCutEemmFone}{not reached}  % experiments/FIGS/data/v2/index.json :: curves_v2/pairs curve ee_vs_mm: largest cut.eps_s over the FLAV_F1 runs ('not reached' when a run's cut is null or leaves fewer than 2 grid points) :: sha256 c6c4138ee5c86423
\newcommand{\MlstRocCutEemmFoner}{not reached}  % experiments/FIGS/data/v2/index.json :: curves_v2/pairs curve ee_vs_mm: largest cut.eps_s over the FLAV_F1R runs ('not reached' when a run's cut is null or leaves fewer than 2 grid points) :: sha256 c6c4138ee5c86423
\newcommand{\MlstRocNBkgEemm}{128{,}833}  % experiments/FIGS/data/v2/index.json :: curves_v2/pairs curve ee_vs_mm: n_background :: sha256 c6c4138ee5c86423
\newcommand{\MlstRocRejMaxEemm}{6{,}442}  % experiments/FIGS/data/v2/index.json :: curves_v2/pairs curve ee_vs_mm: n_background / 20 :: sha256 c6c4138ee5c86423
\newcommand{\MlstRocNoCurveEemm}{the \Kfull- and \Kres-class models and the three random partitions with controlled flavor boundaries}  % experiments/FIGS/data/v2/index.json :: curves_v2/pairs curve ee_vs_mm: configurations of F4 with a run whose cut is null or leaves fewer than 2 grid points (F4 draws no curve for them) :: sha256 c6c4138ee5c86423 :: L188, L162, FLAV_F0, FLAV_F1, FLAV_F1R
\newcommand{\MlstRocCutJcOneHbbFull}{0.40}  % experiments/FIGS/data/v2/index.json :: curves_v2/jc1_ft curve label_Hbb: largest cut.eps_s over the 188-class runs ('not reached' when a run's cut is null or leaves fewer than 2 grid points) :: sha256 c6c4138ee5c86423
\newcommand{\MlstRocCutJcOneHccFull}{0.37}  % experiments/FIGS/data/v2/index.json :: curves_v2/jc1_ft curve label_Hcc: largest cut.eps_s over the 188-class runs ('not reached' when a run's cut is null or leaves fewer than 2 grid points) :: sha256 c6c4138ee5c86423
\newcommand{\MlstRocCutJcOneHggFull}{0.04}  % experiments/FIGS/data/v2/index.json :: curves_v2/jc1_ft curve label_Hgg: largest cut.eps_s over the 188-class runs ('not reached' when a run's cut is null or leaves fewer than 2 grid points) :: sha256 c6c4138ee5c86423
\newcommand{\MlstRocCutJcOneHfourqFull}{0.26}  % experiments/FIGS/data/v2/index.json :: curves_v2/jc1_ft curve label_H4q: largest cut.eps_s over the 188-class runs ('not reached' when a run's cut is null or leaves fewer than 2 grid points) :: sha256 c6c4138ee5c86423
\newcommand{\MlstRocCutJcOneHqqlFull}{0.98}  % experiments/FIGS/data/v2/index.json :: curves_v2/jc1_ft curve label_Hqql: largest cut.eps_s over the 188-class runs ('not reached' when a run's cut is null or leaves fewer than 2 grid points) :: sha256 c6c4138ee5c86423
\newcommand{\MlstRocCutJcOneTbqqFull}{0.53}  % experiments/FIGS/data/v2/index.json :: curves_v2/jc1_ft curve label_Tbqq: largest cut.eps_s over the 188-class runs ('not reached' when a run's cut is null or leaves fewer than 2 grid points) :: sha256 c6c4138ee5c86423
\newcommand{\MlstRocCutJcOneTblFull}{not reached}  % experiments/FIGS/data/v2/index.json :: curves_v2/jc1_ft curve label_Tbl: largest cut.eps_s over the 188-class runs ('not reached' when a run's cut is null or leaves fewer than 2 grid points) :: sha256 c6c4138ee5c86423
\newcommand{\MlstRocCutJcOneWqqFull}{0.15}  % experiments/FIGS/data/v2/index.json :: curves_v2/jc1_ft curve label_Wqq: largest cut.eps_s over the 188-class runs ('not reached' when a run's cut is null or leaves fewer than 2 grid points) :: sha256 c6c4138ee5c86423
\newcommand{\MlstRocCutJcOneZqqFull}{0.12}  % experiments/FIGS/data/v2/index.json :: curves_v2/jc1_ft curve label_Zqq: largest cut.eps_s over the 188-class runs ('not reached' when a run's cut is null or leaves fewer than 2 grid points) :: sha256 c6c4138ee5c86423
\newcommand{\MlstFtTopRfiftyCoarseOverFull}{0.932\,\ensuremath{\pm}\,0.018}  % experiments/FIGS/data/v2/index.json :: curves_v2/top_ft curve top: rejection at eps_S = 0.5 of R16_Q1 over L188, run by run :: sha256 c6c4138ee5c86423
\newcommand{\MlstFtTopRfiftyCoarseOverFullPr}{6.8\% less rejection}  % experiments/FIGS/data/v2/index.json :: curves_v2/top_ft curve top: rejection at eps_S = 0.5 of R16_Q1 over L188, run by run :: sha256 c6c4138ee5c86423 :: prose twin: two significant figures, direction in words
\newcommand{\MlstProbeUnconverged}{0}  % experiments/FIGS/data/v2/index.json :: linear probes of pretrained models that did not converge (count) :: sha256 c6c4138ee5c86423
\newcommand{\MlstProbeUnconvergedInit}{0}  % experiments/FIGS/data/v2/index.json :: linear probes of the randomly initialized networks that did not converge (count) :: sha256 c6c4138ee5c86423
\newcommand{\MlstProbeFits}{895}  % experiments/FIGS/data/v2/index.json :: selected linear-probe fits of the reported checkpoints (count) :: sha256 c6c4138ee5c86423
\newcommand{\MlstProbeEdgeC}{16}  % experiments/FIGS/data/v2/index.json :: of these, fits whose selected C is the largest of the grid (count) :: sha256 c6c4138ee5c86423
\newcommand{\MlstMlpUnconverged}{0}  % experiments/FIGS/data/v2/index.json :: nonlinear probes that did not reach their plateau within the epoch cap (count) :: sha256 c6c4138ee5c86423
\newcommand{\MlstMlpCmpJcTwoResOverFull}{1.019\,\ensuremath{\pm}\,0.007}  % experiments/FIGS/data/v2/index.json :: nonlinear probe jc2_188: (1 - AUC) of L162 over L188, run by run :: sha256 c6c4138ee5c86423
\newcommand{\MlstMlpCmpJcTwoResOverFullPr}{1.9\% more errors}  % experiments/FIGS/data/v2/index.json :: nonlinear probe jc2_188: (1 - AUC) of L162 over L188, run by run :: sha256 c6c4138ee5c86423 :: prose twin: two significant figures, direction in words
\newcommand{\MlstMlpCmpJcTwoMidOverFull}{1.170\,\ensuremath{\pm}\,0.006}  % experiments/FIGS/data/v2/index.json :: nonlinear probe jc2_188: (1 - AUC) of R42_Q1 over L188, run by run :: sha256 c6c4138ee5c86423
\newcommand{\MlstMlpCmpJcTwoMidOverFullPr}{17\% more errors}  % experiments/FIGS/data/v2/index.json :: nonlinear probe jc2_188: (1 - AUC) of R42_Q1 over L188, run by run :: sha256 c6c4138ee5c86423 :: prose twin: two significant figures, direction in words
\newcommand{\MlstMlpCmpJcTwoCoarseOverFull}{1.388\,\ensuremath{\pm}\,0.021}  % experiments/FIGS/data/v2/index.json :: nonlinear probe jc2_188: (1 - AUC) of R16_Q1 over L188, run by run :: sha256 c6c4138ee5c86423
\newcommand{\MlstMlpCmpJcTwoCoarseOverFullPr}{39\% more errors}  % experiments/FIGS/data/v2/index.json :: nonlinear probe jc2_188: (1 - AUC) of R16_Q1 over L188, run by run :: sha256 c6c4138ee5c86423 :: prose twin: two significant figures, direction in words
\newcommand{\MlstMlpCmpJcTwoResMassOverRes}{\pending{no run yet}}  % experiments/FIGS/data/v2/index.json :: nonlinear probe jc2_188: (1 - AUC) of L162_MASS over L162, run by run :: sha256 c6c4138ee5c86423
\newcommand{\MlstMlpCmpJcTwoResMassOverResPr}{\pending{no run yet}}  % experiments/FIGS/data/v2/index.json :: nonlinear probe jc2_188: (1 - AUC) of L162_MASS over L162, run by run :: sha256 c6c4138ee5c86423 :: prose twin: two significant figures, direction in words
\newcommand{\MlstMlpCmpJcTwoResMassOverResPrMax}{\pending{no run yet}}  % experiments/FIGS/data/v2/index.json :: nonlinear probe jc2_188: (1 - AUC) of L162_MASS over L162, run by run :: sha256 c6c4138ee5c86423 :: prose twin of the run farther from 1
\newcommand{\MlstMlpCmpJcTwoCoarseMassOverCoarse}{\pending{no run yet}}  % experiments/FIGS/data/v2/index.json :: nonlinear probe jc2_188: (1 - AUC) of R16_Q1_MASS over R16_Q1, run by run :: sha256 c6c4138ee5c86423
\newcommand{\MlstMlpCmpJcTwoCoarseMassOverCoarsePr}{\pending{no run yet}}  % experiments/FIGS/data/v2/index.json :: nonlinear probe jc2_188: (1 - AUC) of R16_Q1_MASS over R16_Q1, run by run :: sha256 c6c4138ee5c86423 :: prose twin: two significant figures, direction in words
\newcommand{\MlstMlpCmpJcTwoCoarseMassOverCoarsePrMax}{\pending{no run yet}}  % experiments/FIGS/data/v2/index.json :: nonlinear probe jc2_188: (1 - AUC) of R16_Q1_MASS over R16_Q1, run by run :: sha256 c6c4138ee5c86423 :: prose twin of the run farther from 1
\newcommand{\MlstMlpCmpJcTwoFzeroOverFull}{1.294 and 1.310}  % experiments/FIGS/data/v2/index.json :: nonlinear probe jc2_188: (1 - AUC) of FLAV_F0 over L188, run by run :: sha256 c6c4138ee5c86423
\newcommand{\MlstMlpCmpJcTwoFzeroOverFullPr}{at least 29\% more errors}  % experiments/FIGS/data/v2/index.json :: nonlinear probe jc2_188: (1 - AUC) of FLAV_F0 over L188, run by run :: sha256 c6c4138ee5c86423 :: prose twin: two significant figures, direction in words
\newcommand{\MlstMlpCmpJcTwoFzeroOverFullPrMax}{up to 31\% more errors}  % experiments/FIGS/data/v2/index.json :: nonlinear probe jc2_188: (1 - AUC) of FLAV_F0 over L188, run by run :: sha256 c6c4138ee5c86423 :: prose twin of the run farther from 1
\newcommand{\MlstMlpCmpJcTwoFoneOverFzero}{0.903 and 0.885}  % experiments/FIGS/data/v2/index.json :: nonlinear probe jc2_188: (1 - AUC) of FLAV_F1 over FLAV_F0, run by run :: sha256 c6c4138ee5c86423
\newcommand{\MlstMlpCmpJcTwoFoneOverFzeroPr}{at least 9.7\% fewer errors}  % experiments/FIGS/data/v2/index.json :: nonlinear probe jc2_188: (1 - AUC) of FLAV_F1 over FLAV_F0, run by run :: sha256 c6c4138ee5c86423 :: prose twin: two significant figures, direction in words
\newcommand{\MlstMlpCmpJcTwoFoneOverFzeroPrMax}{up to 12\% fewer errors}  % experiments/FIGS/data/v2/index.json :: nonlinear probe jc2_188: (1 - AUC) of FLAV_F1 over FLAV_F0, run by run :: sha256 c6c4138ee5c86423 :: prose twin of the run farther from 1
\newcommand{\MlstMlpCmpJcTwoFonerOverFone}{0.9916 and 0.9918}  % experiments/FIGS/data/v2/index.json :: nonlinear probe jc2_188: (1 - AUC) of FLAV_F1R over FLAV_F1, run by run :: sha256 c6c4138ee5c86423
\newcommand{\MlstMlpCmpJcTwoFonerOverFonePr}{equal within 1\%}  % experiments/FIGS/data/v2/index.json :: nonlinear probe jc2_188: (1 - AUC) of FLAV_F1R over FLAV_F1, run by run :: sha256 c6c4138ee5c86423 :: prose twin: two significant figures, direction in words
\newcommand{\MlstMlpCmpJcTwoFonerOverFonePrMax}{equal within 1\%}  % experiments/FIGS/data/v2/index.json :: nonlinear probe jc2_188: (1 - AUC) of FLAV_F1R over FLAV_F1, run by run :: sha256 c6c4138ee5c86423 :: prose twin of the run farther from 1
\newcommand{\MlstMlpCmpJcTwoFonerOverFzero}{0.895 and 0.877}  % experiments/FIGS/data/v2/index.json :: nonlinear probe jc2_188: (1 - AUC) of FLAV_F1R over FLAV_F0, run by run :: sha256 c6c4138ee5c86423
\newcommand{\MlstMlpCmpJcTwoFonerOverFzeroPr}{at least 10\% fewer errors}  % experiments/FIGS/data/v2/index.json :: nonlinear probe jc2_188: (1 - AUC) of FLAV_F1R over FLAV_F0, run by run :: sha256 c6c4138ee5c86423 :: prose twin: two significant figures, direction in words
\newcommand{\MlstMlpCmpJcTwoFonerOverFzeroPrMax}{up to 12\% fewer errors}  % experiments/FIGS/data/v2/index.json :: nonlinear probe jc2_188: (1 - AUC) of FLAV_F1R over FLAV_F0, run by run :: sha256 c6c4138ee5c86423 :: prose twin of the run farther from 1
\newcommand{\MlstMlpCmpJcTwoFzeroOverCoarse}{0.9461 and 0.9456}  % experiments/FIGS/data/v2/index.json :: nonlinear probe jc2_188: (1 - AUC) of FLAV_F0 over R16_Q1, run by run :: sha256 c6c4138ee5c86423
\newcommand{\MlstMlpCmpJcTwoFzeroOverCoarsePr}{at least 5.4\% fewer errors}  % experiments/FIGS/data/v2/index.json :: nonlinear probe jc2_188: (1 - AUC) of FLAV_F0 over R16_Q1, run by run :: sha256 c6c4138ee5c86423 :: prose twin: two significant figures, direction in words
\newcommand{\MlstMlpCmpJcTwoFzeroOverCoarsePrMax}{up to 5.4\% fewer errors}  % experiments/FIGS/data/v2/index.json :: nonlinear probe jc2_188: (1 - AUC) of FLAV_F0 over R16_Q1, run by run :: sha256 c6c4138ee5c86423 :: prose twin of the run farther from 1
\newcommand{\MlstMlpCmpJcTwoRandOverFull}{\pending{no run yet}}  % experiments/FIGS/data/v2/index.json :: nonlinear probe jc2_188: (1 - AUC) of each random partition over the 188-class run of its seed :: sha256 c6c4138ee5c86423
\newcommand{\MlstMlpCmpJcTwoRandOverFullPr}{\pending{no run yet}}  % experiments/FIGS/data/v2/index.json :: nonlinear probe jc2_188: (1 - AUC) of each random partition over the 188-class run of its seed :: sha256 c6c4138ee5c86423 :: prose twin: two significant figures, direction in words
\newcommand{\MlstMlpCmpJcTwoResResOverFull}{1.001\,\ensuremath{\pm}\,0.007}  % experiments/FIGS/data/v2/index.json :: nonlinear probe jc2_188_res: (1 - AUC) of L162 over L188, run by run :: sha256 c6c4138ee5c86423
\newcommand{\MlstMlpCmpJcTwoResResOverFullPr}{equal within 1\%}  % experiments/FIGS/data/v2/index.json :: nonlinear probe jc2_188_res: (1 - AUC) of L162 over L188, run by run :: sha256 c6c4138ee5c86423 :: prose twin: two significant figures, direction in words
\newcommand{\MlstMlpCmpJcTwoResMidOverFull}{1.192\,\ensuremath{\pm}\,0.008}  % experiments/FIGS/data/v2/index.json :: nonlinear probe jc2_188_res: (1 - AUC) of R42_Q1 over L188, run by run :: sha256 c6c4138ee5c86423
\newcommand{\MlstMlpCmpJcTwoResMidOverFullPr}{19\% more errors}  % experiments/FIGS/data/v2/index.json :: nonlinear probe jc2_188_res: (1 - AUC) of R42_Q1 over L188, run by run :: sha256 c6c4138ee5c86423 :: prose twin: two significant figures, direction in words
\newcommand{\MlstMlpCmpJcTwoResCoarseOverFull}{1.408\,\ensuremath{\pm}\,0.026}  % experiments/FIGS/data/v2/index.json :: nonlinear probe jc2_188_res: (1 - AUC) of R16_Q1 over L188, run by run :: sha256 c6c4138ee5c86423
\newcommand{\MlstMlpCmpJcTwoResCoarseOverFullPr}{41\% more errors}  % experiments/FIGS/data/v2/index.json :: nonlinear probe jc2_188_res: (1 - AUC) of R16_Q1 over L188, run by run :: sha256 c6c4138ee5c86423 :: prose twin: two significant figures, direction in words
\newcommand{\MlstMlpCmpJcTwoResResMassOverRes}{\pending{no run yet}}  % experiments/FIGS/data/v2/index.json :: nonlinear probe jc2_188_res: (1 - AUC) of L162_MASS over L162, run by run :: sha256 c6c4138ee5c86423
\newcommand{\MlstMlpCmpJcTwoResResMassOverResPr}{\pending{no run yet}}  % experiments/FIGS/data/v2/index.json :: nonlinear probe jc2_188_res: (1 - AUC) of L162_MASS over L162, run by run :: sha256 c6c4138ee5c86423 :: prose twin: two significant figures, direction in words
\newcommand{\MlstMlpCmpJcTwoResResMassOverResPrMax}{\pending{no run yet}}  % experiments/FIGS/data/v2/index.json :: nonlinear probe jc2_188_res: (1 - AUC) of L162_MASS over L162, run by run :: sha256 c6c4138ee5c86423 :: prose twin of the run farther from 1
\newcommand{\MlstMlpCmpJcTwoResCoarseMassOverCoarse}{\pending{no run yet}}  % experiments/FIGS/data/v2/index.json :: nonlinear probe jc2_188_res: (1 - AUC) of R16_Q1_MASS over R16_Q1, run by run :: sha256 c6c4138ee5c86423
\newcommand{\MlstMlpCmpJcTwoResCoarseMassOverCoarsePr}{\pending{no run yet}}  % experiments/FIGS/data/v2/index.json :: nonlinear probe jc2_188_res: (1 - AUC) of R16_Q1_MASS over R16_Q1, run by run :: sha256 c6c4138ee5c86423 :: prose twin: two significant figures, direction in words
\newcommand{\MlstMlpCmpJcTwoResCoarseMassOverCoarsePrMax}{\pending{no run yet}}  % experiments/FIGS/data/v2/index.json :: nonlinear probe jc2_188_res: (1 - AUC) of R16_Q1_MASS over R16_Q1, run by run :: sha256 c6c4138ee5c86423 :: prose twin of the run farther from 1
\newcommand{\MlstMlpCmpJcTwoResFzeroOverFull}{1.285 and 1.300}  % experiments/FIGS/data/v2/index.json :: nonlinear probe jc2_188_res: (1 - AUC) of FLAV_F0 over L188, run by run :: sha256 c6c4138ee5c86423
\newcommand{\MlstMlpCmpJcTwoResFzeroOverFullPr}{at least 29\% more errors}  % experiments/FIGS/data/v2/index.json :: nonlinear probe jc2_188_res: (1 - AUC) of FLAV_F0 over L188, run by run :: sha256 c6c4138ee5c86423 :: prose twin: two significant figures, direction in words
\newcommand{\MlstMlpCmpJcTwoResFzeroOverFullPrMax}{up to 30\% more errors}  % experiments/FIGS/data/v2/index.json :: nonlinear probe jc2_188_res: (1 - AUC) of FLAV_F0 over L188, run by run :: sha256 c6c4138ee5c86423 :: prose twin of the run farther from 1
\newcommand{\MlstMlpCmpJcTwoResFoneOverFzero}{0.909 and 0.889}  % experiments/FIGS/data/v2/index.json :: nonlinear probe jc2_188_res: (1 - AUC) of FLAV_F1 over FLAV_F0, run by run :: sha256 c6c4138ee5c86423
\newcommand{\MlstMlpCmpJcTwoResFoneOverFzeroPr}{at least 9.1\% fewer errors}  % experiments/FIGS/data/v2/index.json :: nonlinear probe jc2_188_res: (1 - AUC) of FLAV_F1 over FLAV_F0, run by run :: sha256 c6c4138ee5c86423 :: prose twin: two significant figures, direction in words
\newcommand{\MlstMlpCmpJcTwoResFoneOverFzeroPrMax}{up to 11\% fewer errors}  % experiments/FIGS/data/v2/index.json :: nonlinear probe jc2_188_res: (1 - AUC) of FLAV_F1 over FLAV_F0, run by run :: sha256 c6c4138ee5c86423 :: prose twin of the run farther from 1
\newcommand{\MlstMlpCmpJcTwoResFonerOverFone}{0.9793 and 0.9825}  % experiments/FIGS/data/v2/index.json :: nonlinear probe jc2_188_res: (1 - AUC) of FLAV_F1R over FLAV_F1, run by run :: sha256 c6c4138ee5c86423
\newcommand{\MlstMlpCmpJcTwoResFonerOverFonePr}{equal within 3\%}  % experiments/FIGS/data/v2/index.json :: nonlinear probe jc2_188_res: (1 - AUC) of FLAV_F1R over FLAV_F1, run by run :: sha256 c6c4138ee5c86423 :: prose twin: two significant figures, direction in words
\newcommand{\MlstMlpCmpJcTwoResFonerOverFonePrMax}{equal within 3\%}  % experiments/FIGS/data/v2/index.json :: nonlinear probe jc2_188_res: (1 - AUC) of FLAV_F1R over FLAV_F1, run by run :: sha256 c6c4138ee5c86423 :: prose twin of the run farther from 1
\newcommand{\MlstMlpCmpJcTwoResFonerOverFzero}{0.890 and 0.874}  % experiments/FIGS/data/v2/index.json :: nonlinear probe jc2_188_res: (1 - AUC) of FLAV_F1R over FLAV_F0, run by run :: sha256 c6c4138ee5c86423
\newcommand{\MlstMlpCmpJcTwoResFonerOverFzeroPr}{at least 11\% fewer errors}  % experiments/FIGS/data/v2/index.json :: nonlinear probe jc2_188_res: (1 - AUC) of FLAV_F1R over FLAV_F0, run by run :: sha256 c6c4138ee5c86423 :: prose twin: two significant figures, direction in words
\newcommand{\MlstMlpCmpJcTwoResFonerOverFzeroPrMax}{up to 13\% fewer errors}  % experiments/FIGS/data/v2/index.json :: nonlinear probe jc2_188_res: (1 - AUC) of FLAV_F1R over FLAV_F0, run by run :: sha256 c6c4138ee5c86423 :: prose twin of the run farther from 1
\newcommand{\MlstMlpCmpJcTwoResFzeroOverCoarse}{0.9299 and 0.9233}  % experiments/FIGS/data/v2/index.json :: nonlinear probe jc2_188_res: (1 - AUC) of FLAV_F0 over R16_Q1, run by run :: sha256 c6c4138ee5c86423
\newcommand{\MlstMlpCmpJcTwoResFzeroOverCoarsePr}{at least 7.0\% fewer errors}  % experiments/FIGS/data/v2/index.json :: nonlinear probe jc2_188_res: (1 - AUC) of FLAV_F0 over R16_Q1, run by run :: sha256 c6c4138ee5c86423 :: prose twin: two significant figures, direction in words
\newcommand{\MlstMlpCmpJcTwoResFzeroOverCoarsePrMax}{up to 7.7\% fewer errors}  % experiments/FIGS/data/v2/index.json :: nonlinear probe jc2_188_res: (1 - AUC) of FLAV_F0 over R16_Q1, run by run :: sha256 c6c4138ee5c86423 :: prose twin of the run farther from 1
\newcommand{\MlstMlpCmpJcTwoResRandOverFull}{\pending{no run yet}}  % experiments/FIGS/data/v2/index.json :: nonlinear probe jc2_188_res: (1 - AUC) of each random partition over the 188-class run of its seed :: sha256 c6c4138ee5c86423
\newcommand{\MlstMlpCmpJcTwoResRandOverFullPr}{\pending{no run yet}}  % experiments/FIGS/data/v2/index.json :: nonlinear probe jc2_188_res: (1 - AUC) of each random partition over the 188-class run of its seed :: sha256 c6c4138ee5c86423 :: prose twin: two significant figures, direction in words
\newcommand{\MlstMlpCmpJcTwoQcdResOverFull}{1.070\,\ensuremath{\pm}\,0.006}  % experiments/FIGS/data/v2/index.json :: nonlinear probe jc2_188_qcd: (1 - AUC) of L162 over L188, run by run :: sha256 c6c4138ee5c86423
\newcommand{\MlstMlpCmpJcTwoQcdResOverFullPr}{7.0\% more errors}  % experiments/FIGS/data/v2/index.json :: nonlinear probe jc2_188_qcd: (1 - AUC) of L162 over L188, run by run :: sha256 c6c4138ee5c86423 :: prose twin: two significant figures, direction in words
\newcommand{\MlstMlpCmpJcTwoQcdMidOverFull}{1.107\,\ensuremath{\pm}\,0.002}  % experiments/FIGS/data/v2/index.json :: nonlinear probe jc2_188_qcd: (1 - AUC) of R42_Q1 over L188, run by run :: sha256 c6c4138ee5c86423
\newcommand{\MlstMlpCmpJcTwoQcdMidOverFullPr}{11\% more errors}  % experiments/FIGS/data/v2/index.json :: nonlinear probe jc2_188_qcd: (1 - AUC) of R42_Q1 over L188, run by run :: sha256 c6c4138ee5c86423 :: prose twin: two significant figures, direction in words
\newcommand{\MlstMlpCmpJcTwoQcdCoarseOverFull}{1.333\,\ensuremath{\pm}\,0.012}  % experiments/FIGS/data/v2/index.json :: nonlinear probe jc2_188_qcd: (1 - AUC) of R16_Q1 over L188, run by run :: sha256 c6c4138ee5c86423
\newcommand{\MlstMlpCmpJcTwoQcdCoarseOverFullPr}{33\% more errors}  % experiments/FIGS/data/v2/index.json :: nonlinear probe jc2_188_qcd: (1 - AUC) of R16_Q1 over L188, run by run :: sha256 c6c4138ee5c86423 :: prose twin: two significant figures, direction in words
\newcommand{\MlstMlpCmpJcTwoQcdResMassOverRes}{\pending{no run yet}}  % experiments/FIGS/data/v2/index.json :: nonlinear probe jc2_188_qcd: (1 - AUC) of L162_MASS over L162, run by run :: sha256 c6c4138ee5c86423
\newcommand{\MlstMlpCmpJcTwoQcdResMassOverResPr}{\pending{no run yet}}  % experiments/FIGS/data/v2/index.json :: nonlinear probe jc2_188_qcd: (1 - AUC) of L162_MASS over L162, run by run :: sha256 c6c4138ee5c86423 :: prose twin: two significant figures, direction in words
\newcommand{\MlstMlpCmpJcTwoQcdResMassOverResPrMax}{\pending{no run yet}}  % experiments/FIGS/data/v2/index.json :: nonlinear probe jc2_188_qcd: (1 - AUC) of L162_MASS over L162, run by run :: sha256 c6c4138ee5c86423 :: prose twin of the run farther from 1
\newcommand{\MlstMlpCmpJcTwoQcdCoarseMassOverCoarse}{\pending{no run yet}}  % experiments/FIGS/data/v2/index.json :: nonlinear probe jc2_188_qcd: (1 - AUC) of R16_Q1_MASS over R16_Q1, run by run :: sha256 c6c4138ee5c86423
\newcommand{\MlstMlpCmpJcTwoQcdCoarseMassOverCoarsePr}{\pending{no run yet}}  % experiments/FIGS/data/v2/index.json :: nonlinear probe jc2_188_qcd: (1 - AUC) of R16_Q1_MASS over R16_Q1, run by run :: sha256 c6c4138ee5c86423 :: prose twin: two significant figures, direction in words
\newcommand{\MlstMlpCmpJcTwoQcdCoarseMassOverCoarsePrMax}{\pending{no run yet}}  % experiments/FIGS/data/v2/index.json :: nonlinear probe jc2_188_qcd: (1 - AUC) of R16_Q1_MASS over R16_Q1, run by run :: sha256 c6c4138ee5c86423 :: prose twin of the run farther from 1
\newcommand{\MlstMlpCmpJcTwoQcdFzeroOverFull}{1.319 and 1.339}  % experiments/FIGS/data/v2/index.json :: nonlinear probe jc2_188_qcd: (1 - AUC) of FLAV_F0 over L188, run by run :: sha256 c6c4138ee5c86423
\newcommand{\MlstMlpCmpJcTwoQcdFzeroOverFullPr}{at least 32\% more errors}  % experiments/FIGS/data/v2/index.json :: nonlinear probe jc2_188_qcd: (1 - AUC) of FLAV_F0 over L188, run by run :: sha256 c6c4138ee5c86423 :: prose twin: two significant figures, direction in words
\newcommand{\MlstMlpCmpJcTwoQcdFzeroOverFullPrMax}{up to 34\% more errors}  % experiments/FIGS/data/v2/index.json :: nonlinear probe jc2_188_qcd: (1 - AUC) of FLAV_F0 over L188, run by run :: sha256 c6c4138ee5c86423 :: prose twin of the run farther from 1
\newcommand{\MlstMlpCmpJcTwoQcdFoneOverFzero}{0.886 and 0.872}  % experiments/FIGS/data/v2/index.json :: nonlinear probe jc2_188_qcd: (1 - AUC) of FLAV_F1 over FLAV_F0, run by run :: sha256 c6c4138ee5c86423
\newcommand{\MlstMlpCmpJcTwoQcdFoneOverFzeroPr}{at least 11\% fewer errors}  % experiments/FIGS/data/v2/index.json :: nonlinear probe jc2_188_qcd: (1 - AUC) of FLAV_F1 over FLAV_F0, run by run :: sha256 c6c4138ee5c86423 :: prose twin: two significant figures, direction in words
\newcommand{\MlstMlpCmpJcTwoQcdFoneOverFzeroPrMax}{up to 13\% fewer errors}  % experiments/FIGS/data/v2/index.json :: nonlinear probe jc2_188_qcd: (1 - AUC) of FLAV_F1 over FLAV_F0, run by run :: sha256 c6c4138ee5c86423 :: prose twin of the run farther from 1
\newcommand{\MlstMlpCmpJcTwoQcdFonerOverFone}{1.026 and 1.018}  % experiments/FIGS/data/v2/index.json :: nonlinear probe jc2_188_qcd: (1 - AUC) of FLAV_F1R over FLAV_F1, run by run :: sha256 c6c4138ee5c86423
\newcommand{\MlstMlpCmpJcTwoQcdFonerOverFonePr}{at least 1.8\% more errors}  % experiments/FIGS/data/v2/index.json :: nonlinear probe jc2_188_qcd: (1 - AUC) of FLAV_F1R over FLAV_F1, run by run :: sha256 c6c4138ee5c86423 :: prose twin: two significant figures, direction in words
\newcommand{\MlstMlpCmpJcTwoQcdFonerOverFonePrMax}{up to 2.6\% more errors}  % experiments/FIGS/data/v2/index.json :: nonlinear probe jc2_188_qcd: (1 - AUC) of FLAV_F1R over FLAV_F1, run by run :: sha256 c6c4138ee5c86423 :: prose twin of the run farther from 1
\newcommand{\MlstMlpCmpJcTwoQcdFonerOverFzero}{0.909 and 0.887}  % experiments/FIGS/data/v2/index.json :: nonlinear probe jc2_188_qcd: (1 - AUC) of FLAV_F1R over FLAV_F0, run by run :: sha256 c6c4138ee5c86423
\newcommand{\MlstMlpCmpJcTwoQcdFonerOverFzeroPr}{at least 9.1\% fewer errors}  % experiments/FIGS/data/v2/index.json :: nonlinear probe jc2_188_qcd: (1 - AUC) of FLAV_F1R over FLAV_F0, run by run :: sha256 c6c4138ee5c86423 :: prose twin: two significant figures, direction in words
\newcommand{\MlstMlpCmpJcTwoQcdFonerOverFzeroPrMax}{up to 11\% fewer errors}  % experiments/FIGS/data/v2/index.json :: nonlinear probe jc2_188_qcd: (1 - AUC) of FLAV_F1R over FLAV_F0, run by run :: sha256 c6c4138ee5c86423 :: prose twin of the run farther from 1
\newcommand{\MlstMlpCmpJcTwoQcdFzeroOverCoarse}{0.993 and 1.012}  % experiments/FIGS/data/v2/index.json :: nonlinear probe jc2_188_qcd: (1 - AUC) of FLAV_F0 over R16_Q1, run by run :: sha256 c6c4138ee5c86423
\newcommand{\MlstMlpCmpJcTwoQcdFzeroOverCoarsePr}{equal within 2\%}  % experiments/FIGS/data/v2/index.json :: nonlinear probe jc2_188_qcd: (1 - AUC) of FLAV_F0 over R16_Q1, run by run :: sha256 c6c4138ee5c86423 :: prose twin: two significant figures, direction in words
\newcommand{\MlstMlpCmpJcTwoQcdFzeroOverCoarsePrMax}{equal within 2\%}  % experiments/FIGS/data/v2/index.json :: nonlinear probe jc2_188_qcd: (1 - AUC) of FLAV_F0 over R16_Q1, run by run :: sha256 c6c4138ee5c86423 :: prose twin of the run farther from 1
\newcommand{\MlstMlpCmpJcTwoQcdRandOverFull}{\pending{no run yet}}  % experiments/FIGS/data/v2/index.json :: nonlinear probe jc2_188_qcd: (1 - AUC) of each random partition over the 188-class run of its seed :: sha256 c6c4138ee5c86423
\newcommand{\MlstMlpCmpJcTwoQcdRandOverFullPr}{\pending{no run yet}}  % experiments/FIGS/data/v2/index.json :: nonlinear probe jc2_188_qcd: (1 - AUC) of each random partition over the 188-class run of its seed :: sha256 c6c4138ee5c86423 :: prose twin: two significant figures, direction in words
\newcommand{\MlstMlpCmpBcResOverFull}{1.002\,\ensuremath{\pm}\,0.006}  % experiments/FIGS/data/v2/index.json :: nonlinear probe bc_vs_bq_cs: (1 - AUC) of L162 over L188, run by run :: sha256 c6c4138ee5c86423
\newcommand{\MlstMlpCmpBcResOverFullPr}{equal within 1\%}  % experiments/FIGS/data/v2/index.json :: nonlinear probe bc_vs_bq_cs: (1 - AUC) of L162 over L188, run by run :: sha256 c6c4138ee5c86423 :: prose twin: two significant figures, direction in words
\newcommand{\MlstMlpCmpBcMidOverFull}{1.104\,\ensuremath{\pm}\,0.017}  % experiments/FIGS/data/v2/index.json :: nonlinear probe bc_vs_bq_cs: (1 - AUC) of R42_Q1 over L188, run by run :: sha256 c6c4138ee5c86423
\newcommand{\MlstMlpCmpBcMidOverFullPr}{10\% more errors}  % experiments/FIGS/data/v2/index.json :: nonlinear probe bc_vs_bq_cs: (1 - AUC) of R42_Q1 over L188, run by run :: sha256 c6c4138ee5c86423 :: prose twin: two significant figures, direction in words
\newcommand{\MlstMlpCmpBcCoarseOverFull}{1.624\,\ensuremath{\pm}\,0.022}  % experiments/FIGS/data/v2/index.json :: nonlinear probe bc_vs_bq_cs: (1 - AUC) of R16_Q1 over L188, run by run :: sha256 c6c4138ee5c86423
\newcommand{\MlstMlpCmpBcCoarseOverFullPr}{62\% more errors}  % experiments/FIGS/data/v2/index.json :: nonlinear probe bc_vs_bq_cs: (1 - AUC) of R16_Q1 over L188, run by run :: sha256 c6c4138ee5c86423 :: prose twin: two significant figures, direction in words
\newcommand{\MlstMlpCmpBcResMassOverRes}{\pending{no run yet}}  % experiments/FIGS/data/v2/index.json :: nonlinear probe bc_vs_bq_cs: (1 - AUC) of L162_MASS over L162, run by run :: sha256 c6c4138ee5c86423
\newcommand{\MlstMlpCmpBcResMassOverResPr}{\pending{no run yet}}  % experiments/FIGS/data/v2/index.json :: nonlinear probe bc_vs_bq_cs: (1 - AUC) of L162_MASS over L162, run by run :: sha256 c6c4138ee5c86423 :: prose twin: two significant figures, direction in words
\newcommand{\MlstMlpCmpBcResMassOverResPrMax}{\pending{no run yet}}  % experiments/FIGS/data/v2/index.json :: nonlinear probe bc_vs_bq_cs: (1 - AUC) of L162_MASS over L162, run by run :: sha256 c6c4138ee5c86423 :: prose twin of the run farther from 1
\newcommand{\MlstMlpCmpBcCoarseMassOverCoarse}{\pending{no run yet}}  % experiments/FIGS/data/v2/index.json :: nonlinear probe bc_vs_bq_cs: (1 - AUC) of R16_Q1_MASS over R16_Q1, run by run :: sha256 c6c4138ee5c86423
\newcommand{\MlstMlpCmpBcCoarseMassOverCoarsePr}{\pending{no run yet}}  % experiments/FIGS/data/v2/index.json :: nonlinear probe bc_vs_bq_cs: (1 - AUC) of R16_Q1_MASS over R16_Q1, run by run :: sha256 c6c4138ee5c86423 :: prose twin: two significant figures, direction in words
\newcommand{\MlstMlpCmpBcCoarseMassOverCoarsePrMax}{\pending{no run yet}}  % experiments/FIGS/data/v2/index.json :: nonlinear probe bc_vs_bq_cs: (1 - AUC) of R16_Q1_MASS over R16_Q1, run by run :: sha256 c6c4138ee5c86423 :: prose twin of the run farther from 1
\newcommand{\MlstMlpCmpBcFzeroOverFull}{1.633 and 1.679}  % experiments/FIGS/data/v2/index.json :: nonlinear probe bc_vs_bq_cs: (1 - AUC) of FLAV_F0 over L188, run by run :: sha256 c6c4138ee5c86423
\newcommand{\MlstMlpCmpBcFzeroOverFullPr}{at least 63\% more errors}  % experiments/FIGS/data/v2/index.json :: nonlinear probe bc_vs_bq_cs: (1 - AUC) of FLAV_F0 over L188, run by run :: sha256 c6c4138ee5c86423 :: prose twin: two significant figures, direction in words
\newcommand{\MlstMlpCmpBcFzeroOverFullPrMax}{up to 68\% more errors}  % experiments/FIGS/data/v2/index.json :: nonlinear probe bc_vs_bq_cs: (1 - AUC) of FLAV_F0 over L188, run by run :: sha256 c6c4138ee5c86423 :: prose twin of the run farther from 1
\newcommand{\MlstMlpCmpBcFoneOverFzero}{0.848 and 0.810}  % experiments/FIGS/data/v2/index.json :: nonlinear probe bc_vs_bq_cs: (1 - AUC) of FLAV_F1 over FLAV_F0, run by run :: sha256 c6c4138ee5c86423
\newcommand{\MlstMlpCmpBcFoneOverFzeroPr}{at least 15\% fewer errors}  % experiments/FIGS/data/v2/index.json :: nonlinear probe bc_vs_bq_cs: (1 - AUC) of FLAV_F1 over FLAV_F0, run by run :: sha256 c6c4138ee5c86423 :: prose twin: two significant figures, direction in words
\newcommand{\MlstMlpCmpBcFoneOverFzeroPrMax}{up to 19\% fewer errors}  % experiments/FIGS/data/v2/index.json :: nonlinear probe bc_vs_bq_cs: (1 - AUC) of FLAV_F1 over FLAV_F0, run by run :: sha256 c6c4138ee5c86423 :: prose twin of the run farther from 1
\newcommand{\MlstMlpCmpBcFonerOverFone}{0.9812 and 0.9815}  % experiments/FIGS/data/v2/index.json :: nonlinear probe bc_vs_bq_cs: (1 - AUC) of FLAV_F1R over FLAV_F1, run by run :: sha256 c6c4138ee5c86423
\newcommand{\MlstMlpCmpBcFonerOverFonePr}{equal within 2\%}  % experiments/FIGS/data/v2/index.json :: nonlinear probe bc_vs_bq_cs: (1 - AUC) of FLAV_F1R over FLAV_F1, run by run :: sha256 c6c4138ee5c86423 :: prose twin: two significant figures, direction in words
\newcommand{\MlstMlpCmpBcFonerOverFonePrMax}{equal within 2\%}  % experiments/FIGS/data/v2/index.json :: nonlinear probe bc_vs_bq_cs: (1 - AUC) of FLAV_F1R over FLAV_F1, run by run :: sha256 c6c4138ee5c86423 :: prose twin of the run farther from 1
\newcommand{\MlstMlpCmpBcFonerOverFzero}{0.833 and 0.795}  % experiments/FIGS/data/v2/index.json :: nonlinear probe bc_vs_bq_cs: (1 - AUC) of FLAV_F1R over FLAV_F0, run by run :: sha256 c6c4138ee5c86423
\newcommand{\MlstMlpCmpBcFonerOverFzeroPr}{at least 17\% fewer errors}  % experiments/FIGS/data/v2/index.json :: nonlinear probe bc_vs_bq_cs: (1 - AUC) of FLAV_F1R over FLAV_F0, run by run :: sha256 c6c4138ee5c86423 :: prose twin: two significant figures, direction in words
\newcommand{\MlstMlpCmpBcFonerOverFzeroPrMax}{up to 21\% fewer errors}  % experiments/FIGS/data/v2/index.json :: nonlinear probe bc_vs_bq_cs: (1 - AUC) of FLAV_F1R over FLAV_F0, run by run :: sha256 c6c4138ee5c86423 :: prose twin of the run farther from 1
\newcommand{\MlstMlpCmpBcFzeroOverCoarse}{0.994 and 1.030}  % experiments/FIGS/data/v2/index.json :: nonlinear probe bc_vs_bq_cs: (1 - AUC) of FLAV_F0 over R16_Q1, run by run :: sha256 c6c4138ee5c86423
\newcommand{\MlstMlpCmpBcFzeroOverCoarsePr}{equal within 3\%}  % experiments/FIGS/data/v2/index.json :: nonlinear probe bc_vs_bq_cs: (1 - AUC) of FLAV_F0 over R16_Q1, run by run :: sha256 c6c4138ee5c86423 :: prose twin: two significant figures, direction in words
\newcommand{\MlstMlpCmpBcFzeroOverCoarsePrMax}{equal within 3\%}  % experiments/FIGS/data/v2/index.json :: nonlinear probe bc_vs_bq_cs: (1 - AUC) of FLAV_F0 over R16_Q1, run by run :: sha256 c6c4138ee5c86423 :: prose twin of the run farther from 1
\newcommand{\MlstMlpCmpBcRandOverFull}{\pending{no run yet}}  % experiments/FIGS/data/v2/index.json :: nonlinear probe bc_vs_bq_cs: (1 - AUC) of each random partition over the 188-class run of its seed :: sha256 c6c4138ee5c86423
\newcommand{\MlstMlpCmpBcRandOverFullPr}{\pending{no run yet}}  % experiments/FIGS/data/v2/index.json :: nonlinear probe bc_vs_bq_cs: (1 - AUC) of each random partition over the 188-class run of its seed :: sha256 c6c4138ee5c86423 :: prose twin: two significant figures, direction in words
\newcommand{\MlstMlpCmpBbResOverFull}{1.023\,\ensuremath{\pm}\,0.021}  % experiments/FIGS/data/v2/index.json :: nonlinear probe bb_vs_cc: (1 - AUC) of L162 over L188, run by run :: sha256 c6c4138ee5c86423
\newcommand{\MlstMlpCmpBbResOverFullPr}{equal within 3\%}  % experiments/FIGS/data/v2/index.json :: nonlinear probe bb_vs_cc: (1 - AUC) of L162 over L188, run by run :: sha256 c6c4138ee5c86423 :: prose twin: two significant figures, direction in words
\newcommand{\MlstMlpCmpBbMidOverFull}{1.125\,\ensuremath{\pm}\,0.011}  % experiments/FIGS/data/v2/index.json :: nonlinear probe bb_vs_cc: (1 - AUC) of R42_Q1 over L188, run by run :: sha256 c6c4138ee5c86423
\newcommand{\MlstMlpCmpBbMidOverFullPr}{13\% more errors}  % experiments/FIGS/data/v2/index.json :: nonlinear probe bb_vs_cc: (1 - AUC) of R42_Q1 over L188, run by run :: sha256 c6c4138ee5c86423 :: prose twin: two significant figures, direction in words
\newcommand{\MlstMlpCmpBbCoarseOverFull}{3.04\,\ensuremath{\pm}\,0.09}  % experiments/FIGS/data/v2/index.json :: nonlinear probe bb_vs_cc: (1 - AUC) of R16_Q1 over L188, run by run :: sha256 c6c4138ee5c86423
\newcommand{\MlstMlpCmpBbCoarseOverFullPr}{3.0 times the errors}  % experiments/FIGS/data/v2/index.json :: nonlinear probe bb_vs_cc: (1 - AUC) of R16_Q1 over L188, run by run :: sha256 c6c4138ee5c86423 :: prose twin: two significant figures, direction in words
\newcommand{\MlstMlpCmpBbResMassOverRes}{\pending{no run yet}}  % experiments/FIGS/data/v2/index.json :: nonlinear probe bb_vs_cc: (1 - AUC) of L162_MASS over L162, run by run :: sha256 c6c4138ee5c86423
\newcommand{\MlstMlpCmpBbResMassOverResPr}{\pending{no run yet}}  % experiments/FIGS/data/v2/index.json :: nonlinear probe bb_vs_cc: (1 - AUC) of L162_MASS over L162, run by run :: sha256 c6c4138ee5c86423 :: prose twin: two significant figures, direction in words
\newcommand{\MlstMlpCmpBbResMassOverResPrMax}{\pending{no run yet}}  % experiments/FIGS/data/v2/index.json :: nonlinear probe bb_vs_cc: (1 - AUC) of L162_MASS over L162, run by run :: sha256 c6c4138ee5c86423 :: prose twin of the run farther from 1
\newcommand{\MlstMlpCmpBbCoarseMassOverCoarse}{\pending{no run yet}}  % experiments/FIGS/data/v2/index.json :: nonlinear probe bb_vs_cc: (1 - AUC) of R16_Q1_MASS over R16_Q1, run by run :: sha256 c6c4138ee5c86423
\newcommand{\MlstMlpCmpBbCoarseMassOverCoarsePr}{\pending{no run yet}}  % experiments/FIGS/data/v2/index.json :: nonlinear probe bb_vs_cc: (1 - AUC) of R16_Q1_MASS over R16_Q1, run by run :: sha256 c6c4138ee5c86423 :: prose twin: two significant figures, direction in words
\newcommand{\MlstMlpCmpBbCoarseMassOverCoarsePrMax}{\pending{no run yet}}  % experiments/FIGS/data/v2/index.json :: nonlinear probe bb_vs_cc: (1 - AUC) of R16_Q1_MASS over R16_Q1, run by run :: sha256 c6c4138ee5c86423 :: prose twin of the run farther from 1
\newcommand{\MlstMlpCmpBbFzeroOverFull}{3.3 and 4.3}  % experiments/FIGS/data/v2/index.json :: nonlinear probe bb_vs_cc: (1 - AUC) of FLAV_F0 over L188, run by run :: sha256 c6c4138ee5c86423
\newcommand{\MlstMlpCmpBbFzeroOverFullPr}{at least 3.3 times the errors}  % experiments/FIGS/data/v2/index.json :: nonlinear probe bb_vs_cc: (1 - AUC) of FLAV_F0 over L188, run by run :: sha256 c6c4138ee5c86423 :: prose twin: two significant figures, direction in words
\newcommand{\MlstMlpCmpBbFzeroOverFullPrMax}{up to 4.3 times the errors}  % experiments/FIGS/data/v2/index.json :: nonlinear probe bb_vs_cc: (1 - AUC) of FLAV_F0 over L188, run by run :: sha256 c6c4138ee5c86423 :: prose twin of the run farther from 1
\newcommand{\MlstMlpCmpBbFoneOverFzero}{0.52 and 0.38}  % experiments/FIGS/data/v2/index.json :: nonlinear probe bb_vs_cc: (1 - AUC) of FLAV_F1 over FLAV_F0, run by run :: sha256 c6c4138ee5c86423
\newcommand{\MlstMlpCmpBbFoneOverFzeroPr}{at least 48\% fewer errors}  % experiments/FIGS/data/v2/index.json :: nonlinear probe bb_vs_cc: (1 - AUC) of FLAV_F1 over FLAV_F0, run by run :: sha256 c6c4138ee5c86423 :: prose twin: two significant figures, direction in words
\newcommand{\MlstMlpCmpBbFoneOverFzeroPrMax}{down to 0.38 times the errors}  % experiments/FIGS/data/v2/index.json :: nonlinear probe bb_vs_cc: (1 - AUC) of FLAV_F1 over FLAV_F0, run by run :: sha256 c6c4138ee5c86423 :: prose twin of the run farther from 1
\newcommand{\MlstMlpCmpBbFonerOverFone}{1.27 and 1.45}  % experiments/FIGS/data/v2/index.json :: nonlinear probe bb_vs_cc: (1 - AUC) of FLAV_F1R over FLAV_F1, run by run :: sha256 c6c4138ee5c86423
\newcommand{\MlstMlpCmpBbFonerOverFonePr}{at least 27\% more errors}  % experiments/FIGS/data/v2/index.json :: nonlinear probe bb_vs_cc: (1 - AUC) of FLAV_F1R over FLAV_F1, run by run :: sha256 c6c4138ee5c86423 :: prose twin: two significant figures, direction in words
\newcommand{\MlstMlpCmpBbFonerOverFonePrMax}{up to 45\% more errors}  % experiments/FIGS/data/v2/index.json :: nonlinear probe bb_vs_cc: (1 - AUC) of FLAV_F1R over FLAV_F1, run by run :: sha256 c6c4138ee5c86423 :: prose twin of the run farther from 1
\newcommand{\MlstMlpCmpBbFonerOverFzero}{0.66 and 0.55}  % experiments/FIGS/data/v2/index.json :: nonlinear probe bb_vs_cc: (1 - AUC) of FLAV_F1R over FLAV_F0, run by run :: sha256 c6c4138ee5c86423
\newcommand{\MlstMlpCmpBbFonerOverFzeroPr}{at least 34\% fewer errors}  % experiments/FIGS/data/v2/index.json :: nonlinear probe bb_vs_cc: (1 - AUC) of FLAV_F1R over FLAV_F0, run by run :: sha256 c6c4138ee5c86423 :: prose twin: two significant figures, direction in words
\newcommand{\MlstMlpCmpBbFonerOverFzeroPrMax}{up to 45\% fewer errors}  % experiments/FIGS/data/v2/index.json :: nonlinear probe bb_vs_cc: (1 - AUC) of FLAV_F1R over FLAV_F0, run by run :: sha256 c6c4138ee5c86423 :: prose twin of the run farther from 1
\newcommand{\MlstMlpCmpBbFzeroOverCoarse}{1.07 and 1.41}  % experiments/FIGS/data/v2/index.json :: nonlinear probe bb_vs_cc: (1 - AUC) of FLAV_F0 over R16_Q1, run by run :: sha256 c6c4138ee5c86423
\newcommand{\MlstMlpCmpBbFzeroOverCoarsePr}{at least 7.4\% more errors}  % experiments/FIGS/data/v2/index.json :: nonlinear probe bb_vs_cc: (1 - AUC) of FLAV_F0 over R16_Q1, run by run :: sha256 c6c4138ee5c86423 :: prose twin: two significant figures, direction in words
\newcommand{\MlstMlpCmpBbFzeroOverCoarsePrMax}{up to 41\% more errors}  % experiments/FIGS/data/v2/index.json :: nonlinear probe bb_vs_cc: (1 - AUC) of FLAV_F0 over R16_Q1, run by run :: sha256 c6c4138ee5c86423 :: prose twin of the run farther from 1
\newcommand{\MlstMlpCmpBbRandOverFull}{\pending{no run yet}}  % experiments/FIGS/data/v2/index.json :: nonlinear probe bb_vs_cc: (1 - AUC) of each random partition over the 188-class run of its seed :: sha256 c6c4138ee5c86423
\newcommand{\MlstMlpCmpBbRandOverFullPr}{\pending{no run yet}}  % experiments/FIGS/data/v2/index.json :: nonlinear probe bb_vs_cc: (1 - AUC) of each random partition over the 188-class run of its seed :: sha256 c6c4138ee5c86423 :: prose twin: two significant figures, direction in words
\newcommand{\MlstMlpCmpBbqqResOverFull}{1.011\,\ensuremath{\pm}\,0.030}  % experiments/FIGS/data/v2/index.json :: nonlinear probe bbqq_vs_ccqq: (1 - AUC) of L162 over L188, run by run :: sha256 c6c4138ee5c86423
\newcommand{\MlstMlpCmpBbqqResOverFullPr}{equal within 3\%}  % experiments/FIGS/data/v2/index.json :: nonlinear probe bbqq_vs_ccqq: (1 - AUC) of L162 over L188, run by run :: sha256 c6c4138ee5c86423 :: prose twin: two significant figures, direction in words
\newcommand{\MlstMlpCmpBbqqMidOverFull}{1.112\,\ensuremath{\pm}\,0.006}  % experiments/FIGS/data/v2/index.json :: nonlinear probe bbqq_vs_ccqq: (1 - AUC) of R42_Q1 over L188, run by run :: sha256 c6c4138ee5c86423
\newcommand{\MlstMlpCmpBbqqMidOverFullPr}{11\% more errors}  % experiments/FIGS/data/v2/index.json :: nonlinear probe bbqq_vs_ccqq: (1 - AUC) of R42_Q1 over L188, run by run :: sha256 c6c4138ee5c86423 :: prose twin: two significant figures, direction in words
\newcommand{\MlstMlpCmpBbqqCoarseOverFull}{2.93\,\ensuremath{\pm}\,0.17}  % experiments/FIGS/data/v2/index.json :: nonlinear probe bbqq_vs_ccqq: (1 - AUC) of R16_Q1 over L188, run by run :: sha256 c6c4138ee5c86423
\newcommand{\MlstMlpCmpBbqqCoarseOverFullPr}{2.9 times the errors}  % experiments/FIGS/data/v2/index.json :: nonlinear probe bbqq_vs_ccqq: (1 - AUC) of R16_Q1 over L188, run by run :: sha256 c6c4138ee5c86423 :: prose twin: two significant figures, direction in words
\newcommand{\MlstMlpCmpBbqqResMassOverRes}{\pending{no run yet}}  % experiments/FIGS/data/v2/index.json :: nonlinear probe bbqq_vs_ccqq: (1 - AUC) of L162_MASS over L162, run by run :: sha256 c6c4138ee5c86423
\newcommand{\MlstMlpCmpBbqqResMassOverResPr}{\pending{no run yet}}  % experiments/FIGS/data/v2/index.json :: nonlinear probe bbqq_vs_ccqq: (1 - AUC) of L162_MASS over L162, run by run :: sha256 c6c4138ee5c86423 :: prose twin: two significant figures, direction in words
\newcommand{\MlstMlpCmpBbqqResMassOverResPrMax}{\pending{no run yet}}  % experiments/FIGS/data/v2/index.json :: nonlinear probe bbqq_vs_ccqq: (1 - AUC) of L162_MASS over L162, run by run :: sha256 c6c4138ee5c86423 :: prose twin of the run farther from 1
\newcommand{\MlstMlpCmpBbqqCoarseMassOverCoarse}{\pending{no run yet}}  % experiments/FIGS/data/v2/index.json :: nonlinear probe bbqq_vs_ccqq: (1 - AUC) of R16_Q1_MASS over R16_Q1, run by run :: sha256 c6c4138ee5c86423
\newcommand{\MlstMlpCmpBbqqCoarseMassOverCoarsePr}{\pending{no run yet}}  % experiments/FIGS/data/v2/index.json :: nonlinear probe bbqq_vs_ccqq: (1 - AUC) of R16_Q1_MASS over R16_Q1, run by run :: sha256 c6c4138ee5c86423 :: prose twin: two significant figures, direction in words
\newcommand{\MlstMlpCmpBbqqCoarseMassOverCoarsePrMax}{\pending{no run yet}}  % experiments/FIGS/data/v2/index.json :: nonlinear probe bbqq_vs_ccqq: (1 - AUC) of R16_Q1_MASS over R16_Q1, run by run :: sha256 c6c4138ee5c86423 :: prose twin of the run farther from 1
\newcommand{\MlstMlpCmpBbqqFzeroOverFull}{3.20 and 3.78}  % experiments/FIGS/data/v2/index.json :: nonlinear probe bbqq_vs_ccqq: (1 - AUC) of FLAV_F0 over L188, run by run :: sha256 c6c4138ee5c86423
\newcommand{\MlstMlpCmpBbqqFzeroOverFullPr}{at least 3.2 times the errors}  % experiments/FIGS/data/v2/index.json :: nonlinear probe bbqq_vs_ccqq: (1 - AUC) of FLAV_F0 over L188, run by run :: sha256 c6c4138ee5c86423 :: prose twin: two significant figures, direction in words
\newcommand{\MlstMlpCmpBbqqFzeroOverFullPrMax}{up to 3.8 times the errors}  % experiments/FIGS/data/v2/index.json :: nonlinear probe bbqq_vs_ccqq: (1 - AUC) of FLAV_F0 over L188, run by run :: sha256 c6c4138ee5c86423 :: prose twin of the run farther from 1
\newcommand{\MlstMlpCmpBbqqFoneOverFzero}{0.450 and 0.363}  % experiments/FIGS/data/v2/index.json :: nonlinear probe bbqq_vs_ccqq: (1 - AUC) of FLAV_F1 over FLAV_F0, run by run :: sha256 c6c4138ee5c86423
\newcommand{\MlstMlpCmpBbqqFoneOverFzeroPr}{at most 0.45 times the errors}  % experiments/FIGS/data/v2/index.json :: nonlinear probe bbqq_vs_ccqq: (1 - AUC) of FLAV_F1 over FLAV_F0, run by run :: sha256 c6c4138ee5c86423 :: prose twin: two significant figures, direction in words
\newcommand{\MlstMlpCmpBbqqFoneOverFzeroPrMax}{down to 0.36 times the errors}  % experiments/FIGS/data/v2/index.json :: nonlinear probe bbqq_vs_ccqq: (1 - AUC) of FLAV_F1 over FLAV_F0, run by run :: sha256 c6c4138ee5c86423 :: prose twin of the run farther from 1
\newcommand{\MlstMlpCmpBbqqFonerOverFone}{1.311 and 1.343}  % experiments/FIGS/data/v2/index.json :: nonlinear probe bbqq_vs_ccqq: (1 - AUC) of FLAV_F1R over FLAV_F1, run by run :: sha256 c6c4138ee5c86423
\newcommand{\MlstMlpCmpBbqqFonerOverFonePr}{at least 31\% more errors}  % experiments/FIGS/data/v2/index.json :: nonlinear probe bbqq_vs_ccqq: (1 - AUC) of FLAV_F1R over FLAV_F1, run by run :: sha256 c6c4138ee5c86423 :: prose twin: two significant figures, direction in words
\newcommand{\MlstMlpCmpBbqqFonerOverFonePrMax}{up to 34\% more errors}  % experiments/FIGS/data/v2/index.json :: nonlinear probe bbqq_vs_ccqq: (1 - AUC) of FLAV_F1R over FLAV_F1, run by run :: sha256 c6c4138ee5c86423 :: prose twin of the run farther from 1
\newcommand{\MlstMlpCmpBbqqFonerOverFzero}{0.59 and 0.49}  % experiments/FIGS/data/v2/index.json :: nonlinear probe bbqq_vs_ccqq: (1 - AUC) of FLAV_F1R over FLAV_F0, run by run :: sha256 c6c4138ee5c86423
\newcommand{\MlstMlpCmpBbqqFonerOverFzeroPr}{at least 41\% fewer errors}  % experiments/FIGS/data/v2/index.json :: nonlinear probe bbqq_vs_ccqq: (1 - AUC) of FLAV_F1R over FLAV_F0, run by run :: sha256 c6c4138ee5c86423 :: prose twin: two significant figures, direction in words
\newcommand{\MlstMlpCmpBbqqFonerOverFzeroPrMax}{down to 0.49 times the errors}  % experiments/FIGS/data/v2/index.json :: nonlinear probe bbqq_vs_ccqq: (1 - AUC) of FLAV_F1R over FLAV_F0, run by run :: sha256 c6c4138ee5c86423 :: prose twin of the run farther from 1
\newcommand{\MlstMlpCmpBbqqFzeroOverCoarse}{1.170 and 1.236}  % experiments/FIGS/data/v2/index.json :: nonlinear probe bbqq_vs_ccqq: (1 - AUC) of FLAV_F0 over R16_Q1, run by run :: sha256 c6c4138ee5c86423
\newcommand{\MlstMlpCmpBbqqFzeroOverCoarsePr}{at least 17\% more errors}  % experiments/FIGS/data/v2/index.json :: nonlinear probe bbqq_vs_ccqq: (1 - AUC) of FLAV_F0 over R16_Q1, run by run :: sha256 c6c4138ee5c86423 :: prose twin: two significant figures, direction in words
\newcommand{\MlstMlpCmpBbqqFzeroOverCoarsePrMax}{up to 24\% more errors}  % experiments/FIGS/data/v2/index.json :: nonlinear probe bbqq_vs_ccqq: (1 - AUC) of FLAV_F0 over R16_Q1, run by run :: sha256 c6c4138ee5c86423 :: prose twin of the run farther from 1
\newcommand{\MlstMlpCmpBbqqRandOverFull}{\pending{no run yet}}  % experiments/FIGS/data/v2/index.json :: nonlinear probe bbqq_vs_ccqq: (1 - AUC) of each random partition over the 188-class run of its seed :: sha256 c6c4138ee5c86423
\newcommand{\MlstMlpCmpBbqqRandOverFullPr}{\pending{no run yet}}  % experiments/FIGS/data/v2/index.json :: nonlinear probe bbqq_vs_ccqq: (1 - AUC) of each random partition over the 188-class run of its seed :: sha256 c6c4138ee5c86423 :: prose twin: two significant figures, direction in words
\newcommand{\MlstMlpCmpQcdBbResOverFull}{1.085\,\ensuremath{\pm}\,0.030}  % experiments/FIGS/data/v2/index.json :: nonlinear probe qcd_bb_vs_cc: (1 - AUC) of L162 over L188, run by run :: sha256 c6c4138ee5c86423
\newcommand{\MlstMlpCmpQcdBbResOverFullPr}{8.5\% more errors}  % experiments/FIGS/data/v2/index.json :: nonlinear probe qcd_bb_vs_cc: (1 - AUC) of L162 over L188, run by run :: sha256 c6c4138ee5c86423 :: prose twin: two significant figures, direction in words
\newcommand{\MlstMlpCmpQcdBbMidOverFull}{1.121\,\ensuremath{\pm}\,0.017}  % experiments/FIGS/data/v2/index.json :: nonlinear probe qcd_bb_vs_cc: (1 - AUC) of R42_Q1 over L188, run by run :: sha256 c6c4138ee5c86423
\newcommand{\MlstMlpCmpQcdBbMidOverFullPr}{12\% more errors}  % experiments/FIGS/data/v2/index.json :: nonlinear probe qcd_bb_vs_cc: (1 - AUC) of R42_Q1 over L188, run by run :: sha256 c6c4138ee5c86423 :: prose twin: two significant figures, direction in words
\newcommand{\MlstMlpCmpQcdBbCoarseOverFull}{2.31\,\ensuremath{\pm}\,0.04}  % experiments/FIGS/data/v2/index.json :: nonlinear probe qcd_bb_vs_cc: (1 - AUC) of R16_Q1 over L188, run by run :: sha256 c6c4138ee5c86423
\newcommand{\MlstMlpCmpQcdBbCoarseOverFullPr}{2.3 times the errors}  % experiments/FIGS/data/v2/index.json :: nonlinear probe qcd_bb_vs_cc: (1 - AUC) of R16_Q1 over L188, run by run :: sha256 c6c4138ee5c86423 :: prose twin: two significant figures, direction in words
\newcommand{\MlstMlpCmpQcdBbResMassOverRes}{\pending{no run yet}}  % experiments/FIGS/data/v2/index.json :: nonlinear probe qcd_bb_vs_cc: (1 - AUC) of L162_MASS over L162, run by run :: sha256 c6c4138ee5c86423
\newcommand{\MlstMlpCmpQcdBbResMassOverResPr}{\pending{no run yet}}  % experiments/FIGS/data/v2/index.json :: nonlinear probe qcd_bb_vs_cc: (1 - AUC) of L162_MASS over L162, run by run :: sha256 c6c4138ee5c86423 :: prose twin: two significant figures, direction in words
\newcommand{\MlstMlpCmpQcdBbResMassOverResPrMax}{\pending{no run yet}}  % experiments/FIGS/data/v2/index.json :: nonlinear probe qcd_bb_vs_cc: (1 - AUC) of L162_MASS over L162, run by run :: sha256 c6c4138ee5c86423 :: prose twin of the run farther from 1
\newcommand{\MlstMlpCmpQcdBbCoarseMassOverCoarse}{\pending{no run yet}}  % experiments/FIGS/data/v2/index.json :: nonlinear probe qcd_bb_vs_cc: (1 - AUC) of R16_Q1_MASS over R16_Q1, run by run :: sha256 c6c4138ee5c86423
\newcommand{\MlstMlpCmpQcdBbCoarseMassOverCoarsePr}{\pending{no run yet}}  % experiments/FIGS/data/v2/index.json :: nonlinear probe qcd_bb_vs_cc: (1 - AUC) of R16_Q1_MASS over R16_Q1, run by run :: sha256 c6c4138ee5c86423 :: prose twin: two significant figures, direction in words
\newcommand{\MlstMlpCmpQcdBbCoarseMassOverCoarsePrMax}{\pending{no run yet}}  % experiments/FIGS/data/v2/index.json :: nonlinear probe qcd_bb_vs_cc: (1 - AUC) of R16_Q1_MASS over R16_Q1, run by run :: sha256 c6c4138ee5c86423 :: prose twin of the run farther from 1
\newcommand{\MlstMlpCmpQcdBbFzeroOverFull}{2.50 and 2.72}  % experiments/FIGS/data/v2/index.json :: nonlinear probe qcd_bb_vs_cc: (1 - AUC) of FLAV_F0 over L188, run by run :: sha256 c6c4138ee5c86423
\newcommand{\MlstMlpCmpQcdBbFzeroOverFullPr}{at least 2.5 times the errors}  % experiments/FIGS/data/v2/index.json :: nonlinear probe qcd_bb_vs_cc: (1 - AUC) of FLAV_F0 over L188, run by run :: sha256 c6c4138ee5c86423 :: prose twin: two significant figures, direction in words
\newcommand{\MlstMlpCmpQcdBbFzeroOverFullPrMax}{up to 2.7 times the errors}  % experiments/FIGS/data/v2/index.json :: nonlinear probe qcd_bb_vs_cc: (1 - AUC) of FLAV_F0 over L188, run by run :: sha256 c6c4138ee5c86423 :: prose twin of the run farther from 1
\newcommand{\MlstMlpCmpQcdBbFoneOverFzero}{0.573 and 0.505}  % experiments/FIGS/data/v2/index.json :: nonlinear probe qcd_bb_vs_cc: (1 - AUC) of FLAV_F1 over FLAV_F0, run by run :: sha256 c6c4138ee5c86423
\newcommand{\MlstMlpCmpQcdBbFoneOverFzeroPr}{at least 43\% fewer errors}  % experiments/FIGS/data/v2/index.json :: nonlinear probe qcd_bb_vs_cc: (1 - AUC) of FLAV_F1 over FLAV_F0, run by run :: sha256 c6c4138ee5c86423 :: prose twin: two significant figures, direction in words
\newcommand{\MlstMlpCmpQcdBbFoneOverFzeroPrMax}{up to 49\% fewer errors}  % experiments/FIGS/data/v2/index.json :: nonlinear probe qcd_bb_vs_cc: (1 - AUC) of FLAV_F1 over FLAV_F0, run by run :: sha256 c6c4138ee5c86423 :: prose twin of the run farther from 1
\newcommand{\MlstMlpCmpQcdBbFonerOverFone}{1.219 and 1.258}  % experiments/FIGS/data/v2/index.json :: nonlinear probe qcd_bb_vs_cc: (1 - AUC) of FLAV_F1R over FLAV_F1, run by run :: sha256 c6c4138ee5c86423
\newcommand{\MlstMlpCmpQcdBbFonerOverFonePr}{at least 22\% more errors}  % experiments/FIGS/data/v2/index.json :: nonlinear probe qcd_bb_vs_cc: (1 - AUC) of FLAV_F1R over FLAV_F1, run by run :: sha256 c6c4138ee5c86423 :: prose twin: two significant figures, direction in words
\newcommand{\MlstMlpCmpQcdBbFonerOverFonePrMax}{up to 26\% more errors}  % experiments/FIGS/data/v2/index.json :: nonlinear probe qcd_bb_vs_cc: (1 - AUC) of FLAV_F1R over FLAV_F1, run by run :: sha256 c6c4138ee5c86423 :: prose twin of the run farther from 1
\newcommand{\MlstMlpCmpQcdBbFonerOverFzero}{0.699 and 0.635}  % experiments/FIGS/data/v2/index.json :: nonlinear probe qcd_bb_vs_cc: (1 - AUC) of FLAV_F1R over FLAV_F0, run by run :: sha256 c6c4138ee5c86423
\newcommand{\MlstMlpCmpQcdBbFonerOverFzeroPr}{at least 30\% fewer errors}  % experiments/FIGS/data/v2/index.json :: nonlinear probe qcd_bb_vs_cc: (1 - AUC) of FLAV_F1R over FLAV_F0, run by run :: sha256 c6c4138ee5c86423 :: prose twin: two significant figures, direction in words
\newcommand{\MlstMlpCmpQcdBbFonerOverFzeroPrMax}{up to 36\% fewer errors}  % experiments/FIGS/data/v2/index.json :: nonlinear probe qcd_bb_vs_cc: (1 - AUC) of FLAV_F1R over FLAV_F0, run by run :: sha256 c6c4138ee5c86423 :: prose twin of the run farther from 1
\newcommand{\MlstMlpCmpQcdBbFzeroOverCoarse}{1.103 and 1.167}  % experiments/FIGS/data/v2/index.json :: nonlinear probe qcd_bb_vs_cc: (1 - AUC) of FLAV_F0 over R16_Q1, run by run :: sha256 c6c4138ee5c86423
\newcommand{\MlstMlpCmpQcdBbFzeroOverCoarsePr}{at least 10\% more errors}  % experiments/FIGS/data/v2/index.json :: nonlinear probe qcd_bb_vs_cc: (1 - AUC) of FLAV_F0 over R16_Q1, run by run :: sha256 c6c4138ee5c86423 :: prose twin: two significant figures, direction in words
\newcommand{\MlstMlpCmpQcdBbFzeroOverCoarsePrMax}{up to 17\% more errors}  % experiments/FIGS/data/v2/index.json :: nonlinear probe qcd_bb_vs_cc: (1 - AUC) of FLAV_F0 over R16_Q1, run by run :: sha256 c6c4138ee5c86423 :: prose twin of the run farther from 1
\newcommand{\MlstMlpCmpQcdBbRandOverFull}{\pending{no run yet}}  % experiments/FIGS/data/v2/index.json :: nonlinear probe qcd_bb_vs_cc: (1 - AUC) of each random partition over the 188-class run of its seed :: sha256 c6c4138ee5c86423
\newcommand{\MlstMlpCmpQcdBbRandOverFullPr}{\pending{no run yet}}  % experiments/FIGS/data/v2/index.json :: nonlinear probe qcd_bb_vs_cc: (1 - AUC) of each random partition over the 188-class run of its seed :: sha256 c6c4138ee5c86423 :: prose twin: two significant figures, direction in words
\newcommand{\MlstMlpCmpQcdBResOverFull}{1.050\,\ensuremath{\pm}\,0.004}  % experiments/FIGS/data/v2/index.json :: nonlinear probe qcd_b_vs_light: (1 - AUC) of L162 over L188, run by run :: sha256 c6c4138ee5c86423
\newcommand{\MlstMlpCmpQcdBResOverFullPr}{5.0\% more errors}  % experiments/FIGS/data/v2/index.json :: nonlinear probe qcd_b_vs_light: (1 - AUC) of L162 over L188, run by run :: sha256 c6c4138ee5c86423 :: prose twin: two significant figures, direction in words
\newcommand{\MlstMlpCmpQcdBMidOverFull}{1.055\,\ensuremath{\pm}\,0.005}  % experiments/FIGS/data/v2/index.json :: nonlinear probe qcd_b_vs_light: (1 - AUC) of R42_Q1 over L188, run by run :: sha256 c6c4138ee5c86423
\newcommand{\MlstMlpCmpQcdBMidOverFullPr}{5.5\% more errors}  % experiments/FIGS/data/v2/index.json :: nonlinear probe qcd_b_vs_light: (1 - AUC) of R42_Q1 over L188, run by run :: sha256 c6c4138ee5c86423 :: prose twin: two significant figures, direction in words
\newcommand{\MlstMlpCmpQcdBCoarseOverFull}{1.337\,\ensuremath{\pm}\,0.021}  % experiments/FIGS/data/v2/index.json :: nonlinear probe qcd_b_vs_light: (1 - AUC) of R16_Q1 over L188, run by run :: sha256 c6c4138ee5c86423
\newcommand{\MlstMlpCmpQcdBCoarseOverFullPr}{34\% more errors}  % experiments/FIGS/data/v2/index.json :: nonlinear probe qcd_b_vs_light: (1 - AUC) of R16_Q1 over L188, run by run :: sha256 c6c4138ee5c86423 :: prose twin: two significant figures, direction in words
\newcommand{\MlstMlpCmpQcdBResMassOverRes}{\pending{no run yet}}  % experiments/FIGS/data/v2/index.json :: nonlinear probe qcd_b_vs_light: (1 - AUC) of L162_MASS over L162, run by run :: sha256 c6c4138ee5c86423
\newcommand{\MlstMlpCmpQcdBResMassOverResPr}{\pending{no run yet}}  % experiments/FIGS/data/v2/index.json :: nonlinear probe qcd_b_vs_light: (1 - AUC) of L162_MASS over L162, run by run :: sha256 c6c4138ee5c86423 :: prose twin: two significant figures, direction in words
\newcommand{\MlstMlpCmpQcdBResMassOverResPrMax}{\pending{no run yet}}  % experiments/FIGS/data/v2/index.json :: nonlinear probe qcd_b_vs_light: (1 - AUC) of L162_MASS over L162, run by run :: sha256 c6c4138ee5c86423 :: prose twin of the run farther from 1
\newcommand{\MlstMlpCmpQcdBCoarseMassOverCoarse}{\pending{no run yet}}  % experiments/FIGS/data/v2/index.json :: nonlinear probe qcd_b_vs_light: (1 - AUC) of R16_Q1_MASS over R16_Q1, run by run :: sha256 c6c4138ee5c86423
\newcommand{\MlstMlpCmpQcdBCoarseMassOverCoarsePr}{\pending{no run yet}}  % experiments/FIGS/data/v2/index.json :: nonlinear probe qcd_b_vs_light: (1 - AUC) of R16_Q1_MASS over R16_Q1, run by run :: sha256 c6c4138ee5c86423 :: prose twin: two significant figures, direction in words
\newcommand{\MlstMlpCmpQcdBCoarseMassOverCoarsePrMax}{\pending{no run yet}}  % experiments/FIGS/data/v2/index.json :: nonlinear probe qcd_b_vs_light: (1 - AUC) of R16_Q1_MASS over R16_Q1, run by run :: sha256 c6c4138ee5c86423 :: prose twin of the run farther from 1
\newcommand{\MlstMlpCmpQcdBFzeroOverFull}{1.315 and 1.288}  % experiments/FIGS/data/v2/index.json :: nonlinear probe qcd_b_vs_light: (1 - AUC) of FLAV_F0 over L188, run by run :: sha256 c6c4138ee5c86423
\newcommand{\MlstMlpCmpQcdBFzeroOverFullPr}{at least 29\% more errors}  % experiments/FIGS/data/v2/index.json :: nonlinear probe qcd_b_vs_light: (1 - AUC) of FLAV_F0 over L188, run by run :: sha256 c6c4138ee5c86423 :: prose twin: two significant figures, direction in words
\newcommand{\MlstMlpCmpQcdBFzeroOverFullPrMax}{up to 32\% more errors}  % experiments/FIGS/data/v2/index.json :: nonlinear probe qcd_b_vs_light: (1 - AUC) of FLAV_F0 over L188, run by run :: sha256 c6c4138ee5c86423 :: prose twin of the run farther from 1
\newcommand{\MlstMlpCmpQcdBFoneOverFzero}{0.8804 and 0.8849}  % experiments/FIGS/data/v2/index.json :: nonlinear probe qcd_b_vs_light: (1 - AUC) of FLAV_F1 over FLAV_F0, run by run :: sha256 c6c4138ee5c86423
\newcommand{\MlstMlpCmpQcdBFoneOverFzeroPr}{at least 12\% fewer errors}  % experiments/FIGS/data/v2/index.json :: nonlinear probe qcd_b_vs_light: (1 - AUC) of FLAV_F1 over FLAV_F0, run by run :: sha256 c6c4138ee5c86423 :: prose twin: two significant figures, direction in words
\newcommand{\MlstMlpCmpQcdBFoneOverFzeroPrMax}{up to 12\% fewer errors}  % experiments/FIGS/data/v2/index.json :: nonlinear probe qcd_b_vs_light: (1 - AUC) of FLAV_F1 over FLAV_F0, run by run :: sha256 c6c4138ee5c86423 :: prose twin of the run farther from 1
\newcommand{\MlstMlpCmpQcdBFonerOverFone}{1.037 and 1.042}  % experiments/FIGS/data/v2/index.json :: nonlinear probe qcd_b_vs_light: (1 - AUC) of FLAV_F1R over FLAV_F1, run by run :: sha256 c6c4138ee5c86423
\newcommand{\MlstMlpCmpQcdBFonerOverFonePr}{at least 3.7\% more errors}  % experiments/FIGS/data/v2/index.json :: nonlinear probe qcd_b_vs_light: (1 - AUC) of FLAV_F1R over FLAV_F1, run by run :: sha256 c6c4138ee5c86423 :: prose twin: two significant figures, direction in words
\newcommand{\MlstMlpCmpQcdBFonerOverFonePrMax}{up to 4.2\% more errors}  % experiments/FIGS/data/v2/index.json :: nonlinear probe qcd_b_vs_light: (1 - AUC) of FLAV_F1R over FLAV_F1, run by run :: sha256 c6c4138ee5c86423 :: prose twin of the run farther from 1
\newcommand{\MlstMlpCmpQcdBFonerOverFzero}{0.9130 and 0.9217}  % experiments/FIGS/data/v2/index.json :: nonlinear probe qcd_b_vs_light: (1 - AUC) of FLAV_F1R over FLAV_F0, run by run :: sha256 c6c4138ee5c86423
\newcommand{\MlstMlpCmpQcdBFonerOverFzeroPr}{at least 7.8\% fewer errors}  % experiments/FIGS/data/v2/index.json :: nonlinear probe qcd_b_vs_light: (1 - AUC) of FLAV_F1R over FLAV_F0, run by run :: sha256 c6c4138ee5c86423 :: prose twin: two significant figures, direction in words
\newcommand{\MlstMlpCmpQcdBFonerOverFzeroPrMax}{up to 8.7\% fewer errors}  % experiments/FIGS/data/v2/index.json :: nonlinear probe qcd_b_vs_light: (1 - AUC) of FLAV_F1R over FLAV_F0, run by run :: sha256 c6c4138ee5c86423 :: prose twin of the run farther from 1
\newcommand{\MlstMlpCmpQcdBFzeroOverCoarse}{0.968 and 0.978}  % experiments/FIGS/data/v2/index.json :: nonlinear probe qcd_b_vs_light: (1 - AUC) of FLAV_F0 over R16_Q1, run by run :: sha256 c6c4138ee5c86423
\newcommand{\MlstMlpCmpQcdBFzeroOverCoarsePr}{at least 2.2\% fewer errors}  % experiments/FIGS/data/v2/index.json :: nonlinear probe qcd_b_vs_light: (1 - AUC) of FLAV_F0 over R16_Q1, run by run :: sha256 c6c4138ee5c86423 :: prose twin: two significant figures, direction in words
\newcommand{\MlstMlpCmpQcdBFzeroOverCoarsePrMax}{up to 3.2\% fewer errors}  % experiments/FIGS/data/v2/index.json :: nonlinear probe qcd_b_vs_light: (1 - AUC) of FLAV_F0 over R16_Q1, run by run :: sha256 c6c4138ee5c86423 :: prose twin of the run farther from 1
\newcommand{\MlstMlpCmpQcdBRandOverFull}{\pending{no run yet}}  % experiments/FIGS/data/v2/index.json :: nonlinear probe qcd_b_vs_light: (1 - AUC) of each random partition over the 188-class run of its seed :: sha256 c6c4138ee5c86423
\newcommand{\MlstMlpCmpQcdBRandOverFullPr}{\pending{no run yet}}  % experiments/FIGS/data/v2/index.json :: nonlinear probe qcd_b_vs_light: (1 - AUC) of each random partition over the 188-class run of its seed :: sha256 c6c4138ee5c86423 :: prose twin: two significant figures, direction in words
\newcommand{\MlstMlpCmpQcdCResOverFull}{1.039\,\ensuremath{\pm}\,0.005}  % experiments/FIGS/data/v2/index.json :: nonlinear probe qcd_c_vs_light: (1 - AUC) of L162 over L188, run by run :: sha256 c6c4138ee5c86423
\newcommand{\MlstMlpCmpQcdCResOverFullPr}{3.9\% more errors}  % experiments/FIGS/data/v2/index.json :: nonlinear probe qcd_c_vs_light: (1 - AUC) of L162 over L188, run by run :: sha256 c6c4138ee5c86423 :: prose twin: two significant figures, direction in words
\newcommand{\MlstMlpCmpQcdCMidOverFull}{1.055\,\ensuremath{\pm}\,0.006}  % experiments/FIGS/data/v2/index.json :: nonlinear probe qcd_c_vs_light: (1 - AUC) of R42_Q1 over L188, run by run :: sha256 c6c4138ee5c86423
\newcommand{\MlstMlpCmpQcdCMidOverFullPr}{5.5\% more errors}  % experiments/FIGS/data/v2/index.json :: nonlinear probe qcd_c_vs_light: (1 - AUC) of R42_Q1 over L188, run by run :: sha256 c6c4138ee5c86423 :: prose twin: two significant figures, direction in words
\newcommand{\MlstMlpCmpQcdCCoarseOverFull}{1.274\,\ensuremath{\pm}\,0.008}  % experiments/FIGS/data/v2/index.json :: nonlinear probe qcd_c_vs_light: (1 - AUC) of R16_Q1 over L188, run by run :: sha256 c6c4138ee5c86423
\newcommand{\MlstMlpCmpQcdCCoarseOverFullPr}{27\% more errors}  % experiments/FIGS/data/v2/index.json :: nonlinear probe qcd_c_vs_light: (1 - AUC) of R16_Q1 over L188, run by run :: sha256 c6c4138ee5c86423 :: prose twin: two significant figures, direction in words
\newcommand{\MlstMlpCmpQcdCResMassOverRes}{\pending{no run yet}}  % experiments/FIGS/data/v2/index.json :: nonlinear probe qcd_c_vs_light: (1 - AUC) of L162_MASS over L162, run by run :: sha256 c6c4138ee5c86423
\newcommand{\MlstMlpCmpQcdCResMassOverResPr}{\pending{no run yet}}  % experiments/FIGS/data/v2/index.json :: nonlinear probe qcd_c_vs_light: (1 - AUC) of L162_MASS over L162, run by run :: sha256 c6c4138ee5c86423 :: prose twin: two significant figures, direction in words
\newcommand{\MlstMlpCmpQcdCResMassOverResPrMax}{\pending{no run yet}}  % experiments/FIGS/data/v2/index.json :: nonlinear probe qcd_c_vs_light: (1 - AUC) of L162_MASS over L162, run by run :: sha256 c6c4138ee5c86423 :: prose twin of the run farther from 1
\newcommand{\MlstMlpCmpQcdCCoarseMassOverCoarse}{\pending{no run yet}}  % experiments/FIGS/data/v2/index.json :: nonlinear probe qcd_c_vs_light: (1 - AUC) of R16_Q1_MASS over R16_Q1, run by run :: sha256 c6c4138ee5c86423
\newcommand{\MlstMlpCmpQcdCCoarseMassOverCoarsePr}{\pending{no run yet}}  % experiments/FIGS/data/v2/index.json :: nonlinear probe qcd_c_vs_light: (1 - AUC) of R16_Q1_MASS over R16_Q1, run by run :: sha256 c6c4138ee5c86423 :: prose twin: two significant figures, direction in words
\newcommand{\MlstMlpCmpQcdCCoarseMassOverCoarsePrMax}{\pending{no run yet}}  % experiments/FIGS/data/v2/index.json :: nonlinear probe qcd_c_vs_light: (1 - AUC) of R16_Q1_MASS over R16_Q1, run by run :: sha256 c6c4138ee5c86423 :: prose twin of the run farther from 1
\newcommand{\MlstMlpCmpQcdCFzeroOverFull}{1.234 and 1.211}  % experiments/FIGS/data/v2/index.json :: nonlinear probe qcd_c_vs_light: (1 - AUC) of FLAV_F0 over L188, run by run :: sha256 c6c4138ee5c86423
\newcommand{\MlstMlpCmpQcdCFzeroOverFullPr}{at least 21\% more errors}  % experiments/FIGS/data/v2/index.json :: nonlinear probe qcd_c_vs_light: (1 - AUC) of FLAV_F0 over L188, run by run :: sha256 c6c4138ee5c86423 :: prose twin: two significant figures, direction in words
\newcommand{\MlstMlpCmpQcdCFzeroOverFullPrMax}{up to 23\% more errors}  % experiments/FIGS/data/v2/index.json :: nonlinear probe qcd_c_vs_light: (1 - AUC) of FLAV_F0 over L188, run by run :: sha256 c6c4138ee5c86423 :: prose twin of the run farther from 1
\newcommand{\MlstMlpCmpQcdCFoneOverFzero}{0.9395 and 0.9472}  % experiments/FIGS/data/v2/index.json :: nonlinear probe qcd_c_vs_light: (1 - AUC) of FLAV_F1 over FLAV_F0, run by run :: sha256 c6c4138ee5c86423
\newcommand{\MlstMlpCmpQcdCFoneOverFzeroPr}{at least 5.3\% fewer errors}  % experiments/FIGS/data/v2/index.json :: nonlinear probe qcd_c_vs_light: (1 - AUC) of FLAV_F1 over FLAV_F0, run by run :: sha256 c6c4138ee5c86423 :: prose twin: two significant figures, direction in words
\newcommand{\MlstMlpCmpQcdCFoneOverFzeroPrMax}{up to 6.0\% fewer errors}  % experiments/FIGS/data/v2/index.json :: nonlinear probe qcd_c_vs_light: (1 - AUC) of FLAV_F1 over FLAV_F0, run by run :: sha256 c6c4138ee5c86423 :: prose twin of the run farther from 1
\newcommand{\MlstMlpCmpQcdCFonerOverFone}{0.9846 and 0.9767}  % experiments/FIGS/data/v2/index.json :: nonlinear probe qcd_c_vs_light: (1 - AUC) of FLAV_F1R over FLAV_F1, run by run :: sha256 c6c4138ee5c86423
\newcommand{\MlstMlpCmpQcdCFonerOverFonePr}{at least 1.5\% fewer errors}  % experiments/FIGS/data/v2/index.json :: nonlinear probe qcd_c_vs_light: (1 - AUC) of FLAV_F1R over FLAV_F1, run by run :: sha256 c6c4138ee5c86423 :: prose twin: two significant figures, direction in words
\newcommand{\MlstMlpCmpQcdCFonerOverFonePrMax}{up to 2.3\% fewer errors}  % experiments/FIGS/data/v2/index.json :: nonlinear probe qcd_c_vs_light: (1 - AUC) of FLAV_F1R over FLAV_F1, run by run :: sha256 c6c4138ee5c86423 :: prose twin of the run farther from 1
\newcommand{\MlstMlpCmpQcdCFonerOverFzero}{0.9251 and 0.9251}  % experiments/FIGS/data/v2/index.json :: nonlinear probe qcd_c_vs_light: (1 - AUC) of FLAV_F1R over FLAV_F0, run by run :: sha256 c6c4138ee5c86423
\newcommand{\MlstMlpCmpQcdCFonerOverFzeroPr}{at least 7.5\% fewer errors}  % experiments/FIGS/data/v2/index.json :: nonlinear probe qcd_c_vs_light: (1 - AUC) of FLAV_F1R over FLAV_F0, run by run :: sha256 c6c4138ee5c86423 :: prose twin: two significant figures, direction in words
\newcommand{\MlstMlpCmpQcdCFonerOverFzeroPrMax}{up to 7.5\% fewer errors}  % experiments/FIGS/data/v2/index.json :: nonlinear probe qcd_c_vs_light: (1 - AUC) of FLAV_F1R over FLAV_F0, run by run :: sha256 c6c4138ee5c86423 :: prose twin of the run farther from 1
\newcommand{\MlstMlpCmpQcdCFzeroOverCoarse}{0.968 and 0.956}  % experiments/FIGS/data/v2/index.json :: nonlinear probe qcd_c_vs_light: (1 - AUC) of FLAV_F0 over R16_Q1, run by run :: sha256 c6c4138ee5c86423
\newcommand{\MlstMlpCmpQcdCFzeroOverCoarsePr}{at least 3.2\% fewer errors}  % experiments/FIGS/data/v2/index.json :: nonlinear probe qcd_c_vs_light: (1 - AUC) of FLAV_F0 over R16_Q1, run by run :: sha256 c6c4138ee5c86423 :: prose twin: two significant figures, direction in words
\newcommand{\MlstMlpCmpQcdCFzeroOverCoarsePrMax}{up to 4.4\% fewer errors}  % experiments/FIGS/data/v2/index.json :: nonlinear probe qcd_c_vs_light: (1 - AUC) of FLAV_F0 over R16_Q1, run by run :: sha256 c6c4138ee5c86423 :: prose twin of the run farther from 1
\newcommand{\MlstMlpCmpQcdCRandOverFull}{\pending{no run yet}}  % experiments/FIGS/data/v2/index.json :: nonlinear probe qcd_c_vs_light: (1 - AUC) of each random partition over the 188-class run of its seed :: sha256 c6c4138ee5c86423
\newcommand{\MlstMlpCmpQcdCRandOverFullPr}{\pending{no run yet}}  % experiments/FIGS/data/v2/index.json :: nonlinear probe qcd_c_vs_light: (1 - AUC) of each random partition over the 188-class run of its seed :: sha256 c6c4138ee5c86423 :: prose twin: two significant figures, direction in words
\newcommand{\MlstMlpTopRthirtyFullOverInitBn}{1.00\,\ensuremath{\pm}\,0.31}  % experiments/FIGS/data/v2/index.json :: nonlinear probe top, missing inputs zero: R30 of L188 over the BatchNorm-recomputed random network of the same index :: sha256 c6c4138ee5c86423
\newcommand{\MlstMlpTopRthirtyFullOverInitBnPr}{equal within 40\%}  % experiments/FIGS/data/v2/index.json :: nonlinear probe top, missing inputs zero: R30 of L188 over the BatchNorm-recomputed random network of the same index :: sha256 c6c4138ee5c86423 :: prose twin: two significant figures, direction in words
\newcommand{\MlstMlpTopRthirtyResOverInitBn}{0.94\,\ensuremath{\pm}\,0.26}  % experiments/FIGS/data/v2/index.json :: nonlinear probe top, missing inputs zero: R30 of L162 over the BatchNorm-recomputed random network of the same index :: sha256 c6c4138ee5c86423
\newcommand{\MlstMlpTopRthirtyResOverInitBnPr}{equal within 30\%}  % experiments/FIGS/data/v2/index.json :: nonlinear probe top, missing inputs zero: R30 of L162 over the BatchNorm-recomputed random network of the same index :: sha256 c6c4138ee5c86423 :: prose twin: two significant figures, direction in words
\newcommand{\MlstMlpTopRthirtyMidOverInitBn}{0.84\,\ensuremath{\pm}\,0.17}  % experiments/FIGS/data/v2/index.json :: nonlinear probe top, missing inputs zero: R30 of R42_Q1 over the BatchNorm-recomputed random network of the same index :: sha256 c6c4138ee5c86423
\newcommand{\MlstMlpTopRthirtyMidOverInitBnPr}{equal within 20\%}  % experiments/FIGS/data/v2/index.json :: nonlinear probe top, missing inputs zero: R30 of R42_Q1 over the BatchNorm-recomputed random network of the same index :: sha256 c6c4138ee5c86423 :: prose twin: two significant figures, direction in words
\newcommand{\MlstMlpTopRthirtyCoarseOverInitBn}{0.75\,\ensuremath{\pm}\,0.05}  % experiments/FIGS/data/v2/index.json :: nonlinear probe top, missing inputs zero: R30 of R16_Q1 over the BatchNorm-recomputed random network of the same index :: sha256 c6c4138ee5c86423
\newcommand{\MlstMlpTopRthirtyCoarseOverInitBnPr}{25\% less rejection}  % experiments/FIGS/data/v2/index.json :: nonlinear probe top, missing inputs zero: R30 of R16_Q1 over the BatchNorm-recomputed random network of the same index :: sha256 c6c4138ee5c86423 :: prose twin: two significant figures, direction in words
\newcommand{\MlstMlpQgRthirtyFullOverInitBn}{1.083\,\ensuremath{\pm}\,0.026}  % experiments/FIGS/data/v2/index.json :: nonlinear probe qg, missing inputs zero: R30 of L188 over the BatchNorm-recomputed random network of the same index :: sha256 c6c4138ee5c86423
\newcommand{\MlstMlpQgRthirtyFullOverInitBnPr}{8.3\% more rejection}  % experiments/FIGS/data/v2/index.json :: nonlinear probe qg, missing inputs zero: R30 of L188 over the BatchNorm-recomputed random network of the same index :: sha256 c6c4138ee5c86423 :: prose twin: two significant figures, direction in words
\newcommand{\MlstMlpQgRthirtyResOverInitBn}{1.010\,\ensuremath{\pm}\,0.016}  % experiments/FIGS/data/v2/index.json :: nonlinear probe qg, missing inputs zero: R30 of L162 over the BatchNorm-recomputed random network of the same index :: sha256 c6c4138ee5c86423
\newcommand{\MlstMlpQgRthirtyResOverInitBnPr}{equal within 2\%}  % experiments/FIGS/data/v2/index.json :: nonlinear probe qg, missing inputs zero: R30 of L162 over the BatchNorm-recomputed random network of the same index :: sha256 c6c4138ee5c86423 :: prose twin: two significant figures, direction in words
\newcommand{\MlstMlpQgRthirtyMidOverInitBn}{1.06\,\ensuremath{\pm}\,0.05}  % experiments/FIGS/data/v2/index.json :: nonlinear probe qg, missing inputs zero: R30 of R42_Q1 over the BatchNorm-recomputed random network of the same index :: sha256 c6c4138ee5c86423
\newcommand{\MlstMlpQgRthirtyMidOverInitBnPr}{equal within 7\%}  % experiments/FIGS/data/v2/index.json :: nonlinear probe qg, missing inputs zero: R30 of R42_Q1 over the BatchNorm-recomputed random network of the same index :: sha256 c6c4138ee5c86423 :: prose twin: two significant figures, direction in words
\newcommand{\MlstMlpQgRthirtyCoarseOverInitBn}{1.04\,\ensuremath{\pm}\,0.04}  % experiments/FIGS/data/v2/index.json :: nonlinear probe qg, missing inputs zero: R30 of R16_Q1 over the BatchNorm-recomputed random network of the same index :: sha256 c6c4138ee5c86423
\newcommand{\MlstMlpQgRthirtyCoarseOverInitBnPr}{equal within 5\%}  % experiments/FIGS/data/v2/index.json :: nonlinear probe qg, missing inputs zero: R30 of R16_Q1 over the BatchNorm-recomputed random network of the same index :: sha256 c6c4138ee5c86423 :: prose twin: two significant figures, direction in words
\newcommand{\MlstMlpQgHerwigRthirtyFullOverInitBn}{0.971\,\ensuremath{\pm}\,0.034}  % experiments/FIGS/data/v2/index.json :: nonlinear probe qg_herwig, missing inputs zero: R30 of L188 over the BatchNorm-recomputed random network of the same index :: sha256 c6c4138ee5c86423
\newcommand{\MlstMlpQgHerwigRthirtyFullOverInitBnPr}{equal within 4\%}  % experiments/FIGS/data/v2/index.json :: nonlinear probe qg_herwig, missing inputs zero: R30 of L188 over the BatchNorm-recomputed random network of the same index :: sha256 c6c4138ee5c86423 :: prose twin: two significant figures, direction in words
\newcommand{\MlstMlpQgHerwigRthirtyResOverInitBn}{0.94\,\ensuremath{\pm}\,0.05}  % experiments/FIGS/data/v2/index.json :: nonlinear probe qg_herwig, missing inputs zero: R30 of L162 over the BatchNorm-recomputed random network of the same index :: sha256 c6c4138ee5c86423
\newcommand{\MlstMlpQgHerwigRthirtyResOverInitBnPr}{equal within 6\%}  % experiments/FIGS/data/v2/index.json :: nonlinear probe qg_herwig, missing inputs zero: R30 of L162 over the BatchNorm-recomputed random network of the same index :: sha256 c6c4138ee5c86423 :: prose twin: two significant figures, direction in words
\newcommand{\MlstMlpQgHerwigRthirtyMidOverInitBn}{0.945\,\ensuremath{\pm}\,0.030}  % experiments/FIGS/data/v2/index.json :: nonlinear probe qg_herwig, missing inputs zero: R30 of R42_Q1 over the BatchNorm-recomputed random network of the same index :: sha256 c6c4138ee5c86423
\newcommand{\MlstMlpQgHerwigRthirtyMidOverInitBnPr}{equal within 6\%}  % experiments/FIGS/data/v2/index.json :: nonlinear probe qg_herwig, missing inputs zero: R30 of R42_Q1 over the BatchNorm-recomputed random network of the same index :: sha256 c6c4138ee5c86423 :: prose twin: two significant figures, direction in words
\newcommand{\MlstMlpQgHerwigRthirtyCoarseOverInitBn}{0.96\,\ensuremath{\pm}\,0.04}  % experiments/FIGS/data/v2/index.json :: nonlinear probe qg_herwig, missing inputs zero: R30 of R16_Q1 over the BatchNorm-recomputed random network of the same index :: sha256 c6c4138ee5c86423
\newcommand{\MlstMlpQgHerwigRthirtyCoarseOverInitBnPr}{equal within 5\%}  % experiments/FIGS/data/v2/index.json :: nonlinear probe qg_herwig, missing inputs zero: R30 of R16_Q1 over the BatchNorm-recomputed random network of the same index :: sha256 c6c4138ee5c86423 :: prose twin: two significant figures, direction in words
\newcommand{\MlstMlpNatJcTwoResMassOverRes}{1.016 and 1.009}  % experiments/FIGS/data/v2/index.json :: nonlinear probe, natural-composition JetClass-II (fits_v3_mlp2): (1 - AUC) of L162_MASS over L162, run by run :: sha256 c6c4138ee5c86423
\newcommand{\MlstMlpNatJcTwoResMassOverResPr}{equal within 2\%}  % experiments/FIGS/data/v2/index.json :: nonlinear probe, natural-composition JetClass-II (fits_v3_mlp2): (1 - AUC) of L162_MASS over L162, run by run :: sha256 c6c4138ee5c86423 :: prose twin: two significant figures, direction in words
\newcommand{\MlstMlpNatJcTwoResMassOverResPrMax}{equal within 2\%}  % experiments/FIGS/data/v2/index.json :: nonlinear probe, natural-composition JetClass-II (fits_v3_mlp2): (1 - AUC) of L162_MASS over L162, run by run :: sha256 c6c4138ee5c86423 :: prose twin of the run farther from 1
\newcommand{\MlstMlpNatJcTwoCoarseMassOverCoarse}{1.026 and 1.019}  % experiments/FIGS/data/v2/index.json :: nonlinear probe, natural-composition JetClass-II (fits_v3_mlp2): (1 - AUC) of R16_Q1_MASS over R16_Q1, run by run :: sha256 c6c4138ee5c86423
\newcommand{\MlstMlpNatJcTwoCoarseMassOverCoarsePr}{equal within 3\%}  % experiments/FIGS/data/v2/index.json :: nonlinear probe, natural-composition JetClass-II (fits_v3_mlp2): (1 - AUC) of R16_Q1_MASS over R16_Q1, run by run :: sha256 c6c4138ee5c86423 :: prose twin: two significant figures, direction in words
\newcommand{\MlstMlpNatJcTwoCoarseMassOverCoarsePrMax}{equal within 3\%}  % experiments/FIGS/data/v2/index.json :: nonlinear probe, natural-composition JetClass-II (fits_v3_mlp2): (1 - AUC) of R16_Q1_MASS over R16_Q1, run by run :: sha256 c6c4138ee5c86423 :: prose twin of the run farther from 1
\newcommand{\LpSubstitutionTop}{random draws}  % experiments/FIGS/data/v2/index.json :: linprobe/fits_v4_fill: lowest summed validation cross-entropy over every model, largest training set :: sha256 c6c4138ee5c86423 :: zero 8.8428; mean 8.99394; draws 7.00838
\newcommand{\LpSubstitutionQg}{the training mean}  % experiments/FIGS/data/v2/index.json :: linprobe/fits_v4_fill: lowest summed validation cross-entropy over every model, largest training set :: sha256 c6c4138ee5c86423 :: zero 13.6224; mean 13.1964; draws 13.3584
\newcommand{\MlstLpTopRthirtyFullOverInitBn}{4.67\,\ensuremath{\pm}\,0.23}  % experiments/FIGS/data/v2/index.json :: linear probe top, largest set, chosen substitution: R30 of L188 over the BatchNorm-recomputed random network of the same index :: sha256 c6c4138ee5c86423
\newcommand{\MlstLpTopRthirtyFullOverInitBnPr}{4.7 times the rejection}  % experiments/FIGS/data/v2/index.json :: linear probe top, largest set, chosen substitution: R30 of L188 over the BatchNorm-recomputed random network of the same index :: sha256 c6c4138ee5c86423 :: prose twin: two significant figures, direction in words
\newcommand{\MlstLpTopRthirtyFullOverInitBnZero}{1.04\,\ensuremath{\pm}\,0.05}  % experiments/FIGS/data/v2/index.json :: linear probe top, largest set, zero: R30 of L188 over the BatchNorm-recomputed random network of the same index :: sha256 c6c4138ee5c86423
\newcommand{\MlstLpTopRthirtyFullOverInitBnZeroPr}{equal within 5\%}  % experiments/FIGS/data/v2/index.json :: linear probe top, largest set, zero: R30 of L188 over the BatchNorm-recomputed random network of the same index :: sha256 c6c4138ee5c86423 :: prose twin: two significant figures, direction in words
\newcommand{\MlstLpTopRthirtyResOverInitBn}{4.08\,\ensuremath{\pm}\,0.28}  % experiments/FIGS/data/v2/index.json :: linear probe top, largest set, chosen substitution: R30 of L162 over the BatchNorm-recomputed random network of the same index :: sha256 c6c4138ee5c86423
\newcommand{\MlstLpTopRthirtyResOverInitBnPr}{4.1 times the rejection}  % experiments/FIGS/data/v2/index.json :: linear probe top, largest set, chosen substitution: R30 of L162 over the BatchNorm-recomputed random network of the same index :: sha256 c6c4138ee5c86423 :: prose twin: two significant figures, direction in words
\newcommand{\MlstLpTopRthirtyResOverInitBnZero}{1.33\,\ensuremath{\pm}\,0.30}  % experiments/FIGS/data/v2/index.json :: linear probe top, largest set, zero: R30 of L162 over the BatchNorm-recomputed random network of the same index :: sha256 c6c4138ee5c86423
\newcommand{\MlstLpTopRthirtyResOverInitBnZeroPr}{equal within 40\%}  % experiments/FIGS/data/v2/index.json :: linear probe top, largest set, zero: R30 of L162 over the BatchNorm-recomputed random network of the same index :: sha256 c6c4138ee5c86423 :: prose twin: two significant figures, direction in words
\newcommand{\MlstLpTopRthirtyMidOverInitBn}{4.76\,\ensuremath{\pm}\,0.15}  % experiments/FIGS/data/v2/index.json :: linear probe top, largest set, chosen substitution: R30 of R42_Q1 over the BatchNorm-recomputed random network of the same index :: sha256 c6c4138ee5c86423
\newcommand{\MlstLpTopRthirtyMidOverInitBnPr}{4.8 times the rejection}  % experiments/FIGS/data/v2/index.json :: linear probe top, largest set, chosen substitution: R30 of R42_Q1 over the BatchNorm-recomputed random network of the same index :: sha256 c6c4138ee5c86423 :: prose twin: two significant figures, direction in words
\newcommand{\MlstLpTopRthirtyMidOverInitBnZero}{1.01\,\ensuremath{\pm}\,0.28}  % experiments/FIGS/data/v2/index.json :: linear probe top, largest set, zero: R30 of R42_Q1 over the BatchNorm-recomputed random network of the same index :: sha256 c6c4138ee5c86423
\newcommand{\MlstLpTopRthirtyMidOverInitBnZeroPr}{equal within 30\%}  % experiments/FIGS/data/v2/index.json :: linear probe top, largest set, zero: R30 of R42_Q1 over the BatchNorm-recomputed random network of the same index :: sha256 c6c4138ee5c86423 :: prose twin: two significant figures, direction in words
\newcommand{\MlstLpTopRthirtyCoarseOverInitBn}{5.24\,\ensuremath{\pm}\,0.06}  % experiments/FIGS/data/v2/index.json :: linear probe top, largest set, chosen substitution: R30 of R16_Q1 over the BatchNorm-recomputed random network of the same index :: sha256 c6c4138ee5c86423
\newcommand{\MlstLpTopRthirtyCoarseOverInitBnPr}{5.2 times the rejection}  % experiments/FIGS/data/v2/index.json :: linear probe top, largest set, chosen substitution: R30 of R16_Q1 over the BatchNorm-recomputed random network of the same index :: sha256 c6c4138ee5c86423 :: prose twin: two significant figures, direction in words
\newcommand{\MlstLpTopRthirtyCoarseOverInitBnZero}{0.97\,\ensuremath{\pm}\,0.16}  % experiments/FIGS/data/v2/index.json :: linear probe top, largest set, zero: R30 of R16_Q1 over the BatchNorm-recomputed random network of the same index :: sha256 c6c4138ee5c86423
\newcommand{\MlstLpTopRthirtyCoarseOverInitBnZeroPr}{equal within 20\%}  % experiments/FIGS/data/v2/index.json :: linear probe top, largest set, zero: R30 of R16_Q1 over the BatchNorm-recomputed random network of the same index :: sha256 c6c4138ee5c86423 :: prose twin: two significant figures, direction in words
\newcommand{\MlstLpTopRthirtyResMassOverInitBn}{5.71 and 6.25}  % experiments/FIGS/data/v2/index.json :: linear probe top, largest set, chosen substitution: R30 of L162_MASS over the BatchNorm-recomputed random network of the same index :: sha256 c6c4138ee5c86423
\newcommand{\MlstLpTopRthirtyResMassOverInitBnPr}{at least 5.7 times the rejection}  % experiments/FIGS/data/v2/index.json :: linear probe top, largest set, chosen substitution: R30 of L162_MASS over the BatchNorm-recomputed random network of the same index :: sha256 c6c4138ee5c86423 :: prose twin: two significant figures, direction in words
\newcommand{\MlstLpTopRthirtyResMassOverInitBnPrMax}{up to 6.3 times the rejection}  % experiments/FIGS/data/v2/index.json :: linear probe top, largest set, chosen substitution: R30 of L162_MASS over the BatchNorm-recomputed random network of the same index :: sha256 c6c4138ee5c86423 :: prose twin of the run farther from 1
\newcommand{\MlstLpTopRthirtyResMassOverInitBnZero}{1.683 and 1.678}  % experiments/FIGS/data/v2/index.json :: linear probe top, largest set, zero: R30 of L162_MASS over the BatchNorm-recomputed random network of the same index :: sha256 c6c4138ee5c86423
\newcommand{\MlstLpTopRthirtyResMassOverInitBnZeroPr}{at least 68\% more rejection}  % experiments/FIGS/data/v2/index.json :: linear probe top, largest set, zero: R30 of L162_MASS over the BatchNorm-recomputed random network of the same index :: sha256 c6c4138ee5c86423 :: prose twin: two significant figures, direction in words
\newcommand{\MlstLpTopRthirtyResMassOverInitBnZeroPrMax}{up to 68\% more rejection}  % experiments/FIGS/data/v2/index.json :: linear probe top, largest set, zero: R30 of L162_MASS over the BatchNorm-recomputed random network of the same index :: sha256 c6c4138ee5c86423 :: prose twin of the run farther from 1
\newcommand{\MlstLpTopRthirtyCoarseMassOverInitBn}{6.73 and 7.72}  % experiments/FIGS/data/v2/index.json :: linear probe top, largest set, chosen substitution: R30 of R16_Q1_MASS over the BatchNorm-recomputed random network of the same index :: sha256 c6c4138ee5c86423
\newcommand{\MlstLpTopRthirtyCoarseMassOverInitBnPr}{at least 6.7 times the rejection}  % experiments/FIGS/data/v2/index.json :: linear probe top, largest set, chosen substitution: R30 of R16_Q1_MASS over the BatchNorm-recomputed random network of the same index :: sha256 c6c4138ee5c86423 :: prose twin: two significant figures, direction in words
\newcommand{\MlstLpTopRthirtyCoarseMassOverInitBnPrMax}{up to 7.7 times the rejection}  % experiments/FIGS/data/v2/index.json :: linear probe top, largest set, chosen substitution: R30 of R16_Q1_MASS over the BatchNorm-recomputed random network of the same index :: sha256 c6c4138ee5c86423 :: prose twin of the run farther from 1
\newcommand{\MlstLpTopRthirtyCoarseMassOverInitBnZero}{2.3 and 1.2}  % experiments/FIGS/data/v2/index.json :: linear probe top, largest set, zero: R30 of R16_Q1_MASS over the BatchNorm-recomputed random network of the same index :: sha256 c6c4138ee5c86423
\newcommand{\MlstLpTopRthirtyCoarseMassOverInitBnZeroPr}{equal within 200\%}  % experiments/FIGS/data/v2/index.json :: linear probe top, largest set, zero: R30 of R16_Q1_MASS over the BatchNorm-recomputed random network of the same index :: sha256 c6c4138ee5c86423 :: prose twin: two significant figures, direction in words
\newcommand{\MlstLpTopRthirtyCoarseMassOverInitBnZeroPrMax}{equal within 200\%}  % experiments/FIGS/data/v2/index.json :: linear probe top, largest set, zero: R30 of R16_Q1_MASS over the BatchNorm-recomputed random network of the same index :: sha256 c6c4138ee5c86423 :: prose twin of the run farther from 1
\newcommand{\MlstLpTopRthirtyResOverFull}{0.874\,\ensuremath{\pm}\,0.018}  % experiments/FIGS/data/v2/index.json :: linear probe top: R30 of L162 over the 188-class run :: sha256 c6c4138ee5c86423
\newcommand{\MlstLpTopRthirtyResOverFullPr}{13\% less rejection}  % experiments/FIGS/data/v2/index.json :: linear probe top: R30 of L162 over the 188-class run :: sha256 c6c4138ee5c86423 :: prose twin: two significant figures, direction in words
\newcommand{\MlstLpTopRthirtyResOverFullZero}{1.29\,\ensuremath{\pm}\,0.31}  % experiments/FIGS/data/v2/index.json :: linear probe top: R30 of L162 over the 188-class run :: sha256 c6c4138ee5c86423
\newcommand{\MlstLpTopRthirtyResOverFullZeroPr}{equal within 40\%}  % experiments/FIGS/data/v2/index.json :: linear probe top: R30 of L162 over the 188-class run :: sha256 c6c4138ee5c86423 :: prose twin: two significant figures, direction in words
\newcommand{\MlstLpTopRthirtyMidOverFull}{1.02\,\ensuremath{\pm}\,0.08}  % experiments/FIGS/data/v2/index.json :: linear probe top: R30 of R42_Q1 over the 188-class run :: sha256 c6c4138ee5c86423
\newcommand{\MlstLpTopRthirtyMidOverFullPr}{equal within 9\%}  % experiments/FIGS/data/v2/index.json :: linear probe top: R30 of R42_Q1 over the 188-class run :: sha256 c6c4138ee5c86423 :: prose twin: two significant figures, direction in words
\newcommand{\MlstLpTopRthirtyMidOverFullZero}{0.98\,\ensuremath{\pm}\,0.28}  % experiments/FIGS/data/v2/index.json :: linear probe top: R30 of R42_Q1 over the 188-class run :: sha256 c6c4138ee5c86423
\newcommand{\MlstLpTopRthirtyMidOverFullZeroPr}{equal within 30\%}  % experiments/FIGS/data/v2/index.json :: linear probe top: R30 of R42_Q1 over the 188-class run :: sha256 c6c4138ee5c86423 :: prose twin: two significant figures, direction in words
\newcommand{\MlstLpTopRthirtyCoarseOverFull}{1.12\,\ensuremath{\pm}\,0.06}  % experiments/FIGS/data/v2/index.json :: linear probe top: R30 of R16_Q1 over the 188-class run :: sha256 c6c4138ee5c86423
\newcommand{\MlstLpTopRthirtyCoarseOverFullPr}{12\% more rejection}  % experiments/FIGS/data/v2/index.json :: linear probe top: R30 of R16_Q1 over the 188-class run :: sha256 c6c4138ee5c86423 :: prose twin: two significant figures, direction in words
\newcommand{\MlstLpTopRthirtyCoarseOverFullZero}{0.93\,\ensuremath{\pm}\,0.15}  % experiments/FIGS/data/v2/index.json :: linear probe top: R30 of R16_Q1 over the 188-class run :: sha256 c6c4138ee5c86423
\newcommand{\MlstLpTopRthirtyCoarseOverFullZeroPr}{equal within 20\%}  % experiments/FIGS/data/v2/index.json :: linear probe top: R30 of R16_Q1 over the 188-class run :: sha256 c6c4138ee5c86423 :: prose twin: two significant figures, direction in words
\newcommand{\MlstLpTopRthirtyMeanOverZeroFull}{1.25\,\ensuremath{\pm}\,0.27}  % experiments/FIGS/data/v2/index.json :: linear probe top: R30 with mean over R30 with zero, L188 :: sha256 c6c4138ee5c86423
\newcommand{\MlstLpTopRthirtyMeanOverZeroFullPr}{equal within 30\%}  % experiments/FIGS/data/v2/index.json :: linear probe top: R30 with mean over R30 with zero, L188 :: sha256 c6c4138ee5c86423 :: prose twin: two significant figures, direction in words
\newcommand{\MlstLpTopRthirtyDrawOverZeroFull}{3.9\,\ensuremath{\pm}\,0.4}  % experiments/FIGS/data/v2/index.json :: linear probe top: R30 with draws over R30 with zero, L188 :: sha256 c6c4138ee5c86423
\newcommand{\MlstLpTopRthirtyDrawOverZeroFullPr}{3.9 times the rejection}  % experiments/FIGS/data/v2/index.json :: linear probe top: R30 with draws over R30 with zero, L188 :: sha256 c6c4138ee5c86423 :: prose twin: two significant figures, direction in words
\newcommand{\MlstLpTopRthirtyMeanOverZeroInitBn}{1.08\,\ensuremath{\pm}\,0.20}  % experiments/FIGS/data/v2/index.json :: linear probe top: R30 with mean over R30 with zero, INIT_BN :: sha256 c6c4138ee5c86423
\newcommand{\MlstLpTopRthirtyMeanOverZeroInitBnPr}{equal within 30\%}  % experiments/FIGS/data/v2/index.json :: linear probe top: R30 with mean over R30 with zero, INIT_BN :: sha256 c6c4138ee5c86423 :: prose twin: two significant figures, direction in words
\newcommand{\MlstLpTopRthirtyDrawOverZeroInitBn}{0.86\,\ensuremath{\pm}\,0.14}  % experiments/FIGS/data/v2/index.json :: linear probe top: R30 with draws over R30 with zero, INIT_BN :: sha256 c6c4138ee5c86423
\newcommand{\MlstLpTopRthirtyDrawOverZeroInitBnPr}{equal within 20\%}  % experiments/FIGS/data/v2/index.json :: linear probe top: R30 with draws over R30 with zero, INIT_BN :: sha256 c6c4138ee5c86423 :: prose twin: two significant figures, direction in words
\newcommand{\MlstLpTopRthirtyKinOverFull}{3.88}  % experiments/FIGS/data/v2/index.json :: linear probe top: R30 of the kinematic-only model over the mean of the 188-class runs (zero) :: sha256 c6c4138ee5c86423
\newcommand{\MlstLpTopRthirtyKinOverFullPr}{3.9 times the rejection}  % experiments/FIGS/data/v2/index.json :: linear probe top: R30 of the kinematic-only model over the mean of the 188-class runs (zero) :: sha256 c6c4138ee5c86423 :: prose twin: two significant figures, direction in words
\newcommand{\MlstLpTopRthirtyKin}{231}  % experiments/FIGS/data/v2/index.json :: linear probe top, largest set: R30 of the kinematic-only model (one run) :: sha256 c6c4138ee5c86423
\newcommand{\MlstLpTopRthirtyFullZero}{60\,\ensuremath{\pm}\,12}  % experiments/FIGS/data/v2/index.json :: linear probe top, largest set: R30 with zero, L188 :: sha256 c6c4138ee5c86423
\newcommand{\MlstLpTopRthirtyFullMean}{73\,\ensuremath{\pm}\,15}  % experiments/FIGS/data/v2/index.json :: linear probe top, largest set: R30 with mean, L188 :: sha256 c6c4138ee5c86423
\newcommand{\MlstLpTopRthirtyFullDraws}{228\,\ensuremath{\pm}\,25}  % experiments/FIGS/data/v2/index.json :: linear probe top, largest set: R30 with draws, L188 :: sha256 c6c4138ee5c86423
\newcommand{\MlstLpTopRthirtyInitBnZero}{58\,\ensuremath{\pm}\,13}  % experiments/FIGS/data/v2/index.json :: linear probe top, largest set: R30 with zero, INIT_BN :: sha256 c6c4138ee5c86423
\newcommand{\MlstLpTopRthirtyInitBnMean}{60.4\,\ensuremath{\pm}\,2.4}  % experiments/FIGS/data/v2/index.json :: linear probe top, largest set: R30 with mean, INIT_BN :: sha256 c6c4138ee5c86423
\newcommand{\MlstLpTopRthirtyInitBnDraws}{48.6\,\ensuremath{\pm}\,3.0}  % experiments/FIGS/data/v2/index.json :: linear probe top, largest set: R30 with draws, INIT_BN :: sha256 c6c4138ee5c86423
\newcommand{\MlstLpQgRthirtyFullOverInitBn}{1.39\,\ensuremath{\pm}\,0.06}  % experiments/FIGS/data/v2/index.json :: linear probe qg, largest set, chosen substitution: R30 of L188 over the BatchNorm-recomputed random network of the same index :: sha256 c6c4138ee5c86423
\newcommand{\MlstLpQgRthirtyFullOverInitBnPr}{39\% more rejection}  % experiments/FIGS/data/v2/index.json :: linear probe qg, largest set, chosen substitution: R30 of L188 over the BatchNorm-recomputed random network of the same index :: sha256 c6c4138ee5c86423 :: prose twin: two significant figures, direction in words
\newcommand{\MlstLpQgRthirtyFullOverInitBnZero}{1.21\,\ensuremath{\pm}\,0.15}  % experiments/FIGS/data/v2/index.json :: linear probe qg, largest set, zero: R30 of L188 over the BatchNorm-recomputed random network of the same index :: sha256 c6c4138ee5c86423
\newcommand{\MlstLpQgRthirtyFullOverInitBnZeroPr}{equal within 30\%}  % experiments/FIGS/data/v2/index.json :: linear probe qg, largest set, zero: R30 of L188 over the BatchNorm-recomputed random network of the same index :: sha256 c6c4138ee5c86423 :: prose twin: two significant figures, direction in words
\newcommand{\MlstLpQgRthirtyResOverInitBn}{1.360\,\ensuremath{\pm}\,0.025}  % experiments/FIGS/data/v2/index.json :: linear probe qg, largest set, chosen substitution: R30 of L162 over the BatchNorm-recomputed random network of the same index :: sha256 c6c4138ee5c86423
\newcommand{\MlstLpQgRthirtyResOverInitBnPr}{36\% more rejection}  % experiments/FIGS/data/v2/index.json :: linear probe qg, largest set, chosen substitution: R30 of L162 over the BatchNorm-recomputed random network of the same index :: sha256 c6c4138ee5c86423 :: prose twin: two significant figures, direction in words
\newcommand{\MlstLpQgRthirtyResOverInitBnZero}{1.07\,\ensuremath{\pm}\,0.13}  % experiments/FIGS/data/v2/index.json :: linear probe qg, largest set, zero: R30 of L162 over the BatchNorm-recomputed random network of the same index :: sha256 c6c4138ee5c86423
\newcommand{\MlstLpQgRthirtyResOverInitBnZeroPr}{equal within 20\%}  % experiments/FIGS/data/v2/index.json :: linear probe qg, largest set, zero: R30 of L162 over the BatchNorm-recomputed random network of the same index :: sha256 c6c4138ee5c86423 :: prose twin: two significant figures, direction in words
\newcommand{\MlstLpQgRthirtyMidOverInitBn}{1.357\,\ensuremath{\pm}\,0.033}  % experiments/FIGS/data/v2/index.json :: linear probe qg, largest set, chosen substitution: R30 of R42_Q1 over the BatchNorm-recomputed random network of the same index :: sha256 c6c4138ee5c86423
\newcommand{\MlstLpQgRthirtyMidOverInitBnPr}{36\% more rejection}  % experiments/FIGS/data/v2/index.json :: linear probe qg, largest set, chosen substitution: R30 of R42_Q1 over the BatchNorm-recomputed random network of the same index :: sha256 c6c4138ee5c86423 :: prose twin: two significant figures, direction in words
\newcommand{\MlstLpQgRthirtyMidOverInitBnZero}{1.12\,\ensuremath{\pm}\,0.11}  % experiments/FIGS/data/v2/index.json :: linear probe qg, largest set, zero: R30 of R42_Q1 over the BatchNorm-recomputed random network of the same index :: sha256 c6c4138ee5c86423
\newcommand{\MlstLpQgRthirtyMidOverInitBnZeroPr}{equal within 20\%}  % experiments/FIGS/data/v2/index.json :: linear probe qg, largest set, zero: R30 of R42_Q1 over the BatchNorm-recomputed random network of the same index :: sha256 c6c4138ee5c86423 :: prose twin: two significant figures, direction in words
\newcommand{\MlstLpQgRthirtyCoarseOverInitBn}{1.42\,\ensuremath{\pm}\,0.04}  % experiments/FIGS/data/v2/index.json :: linear probe qg, largest set, chosen substitution: R30 of R16_Q1 over the BatchNorm-recomputed random network of the same index :: sha256 c6c4138ee5c86423
\newcommand{\MlstLpQgRthirtyCoarseOverInitBnPr}{42\% more rejection}  % experiments/FIGS/data/v2/index.json :: linear probe qg, largest set, chosen substitution: R30 of R16_Q1 over the BatchNorm-recomputed random network of the same index :: sha256 c6c4138ee5c86423 :: prose twin: two significant figures, direction in words
\newcommand{\MlstLpQgRthirtyCoarseOverInitBnZero}{1.15\,\ensuremath{\pm}\,0.17}  % experiments/FIGS/data/v2/index.json :: linear probe qg, largest set, zero: R30 of R16_Q1 over the BatchNorm-recomputed random network of the same index :: sha256 c6c4138ee5c86423
\newcommand{\MlstLpQgRthirtyCoarseOverInitBnZeroPr}{equal within 20\%}  % experiments/FIGS/data/v2/index.json :: linear probe qg, largest set, zero: R30 of R16_Q1 over the BatchNorm-recomputed random network of the same index :: sha256 c6c4138ee5c86423 :: prose twin: two significant figures, direction in words
\newcommand{\MlstLpQgRthirtyResMassOverInitBn}{1.430 and 1.364}  % experiments/FIGS/data/v2/index.json :: linear probe qg, largest set, chosen substitution: R30 of L162_MASS over the BatchNorm-recomputed random network of the same index :: sha256 c6c4138ee5c86423
\newcommand{\MlstLpQgRthirtyResMassOverInitBnPr}{at least 36\% more rejection}  % experiments/FIGS/data/v2/index.json :: linear probe qg, largest set, chosen substitution: R30 of L162_MASS over the BatchNorm-recomputed random network of the same index :: sha256 c6c4138ee5c86423 :: prose twin: two significant figures, direction in words
\newcommand{\MlstLpQgRthirtyResMassOverInitBnPrMax}{up to 43\% more rejection}  % experiments/FIGS/data/v2/index.json :: linear probe qg, largest set, chosen substitution: R30 of L162_MASS over the BatchNorm-recomputed random network of the same index :: sha256 c6c4138ee5c86423 :: prose twin of the run farther from 1
\newcommand{\MlstLpQgRthirtyResMassOverInitBnZero}{1.213 and 1.175}  % experiments/FIGS/data/v2/index.json :: linear probe qg, largest set, zero: R30 of L162_MASS over the BatchNorm-recomputed random network of the same index :: sha256 c6c4138ee5c86423
\newcommand{\MlstLpQgRthirtyResMassOverInitBnZeroPr}{equal within 30\%}  % experiments/FIGS/data/v2/index.json :: linear probe qg, largest set, zero: R30 of L162_MASS over the BatchNorm-recomputed random network of the same index :: sha256 c6c4138ee5c86423 :: prose twin: two significant figures, direction in words
\newcommand{\MlstLpQgRthirtyResMassOverInitBnZeroPrMax}{equal within 30\%}  % experiments/FIGS/data/v2/index.json :: linear probe qg, largest set, zero: R30 of L162_MASS over the BatchNorm-recomputed random network of the same index :: sha256 c6c4138ee5c86423 :: prose twin of the run farther from 1
\newcommand{\MlstLpQgRthirtyCoarseMassOverInitBn}{1.324 and 1.394}  % experiments/FIGS/data/v2/index.json :: linear probe qg, largest set, chosen substitution: R30 of R16_Q1_MASS over the BatchNorm-recomputed random network of the same index :: sha256 c6c4138ee5c86423
\newcommand{\MlstLpQgRthirtyCoarseMassOverInitBnPr}{at least 32\% more rejection}  % experiments/FIGS/data/v2/index.json :: linear probe qg, largest set, chosen substitution: R30 of R16_Q1_MASS over the BatchNorm-recomputed random network of the same index :: sha256 c6c4138ee5c86423 :: prose twin: two significant figures, direction in words
\newcommand{\MlstLpQgRthirtyCoarseMassOverInitBnPrMax}{up to 39\% more rejection}  % experiments/FIGS/data/v2/index.json :: linear probe qg, largest set, chosen substitution: R30 of R16_Q1_MASS over the BatchNorm-recomputed random network of the same index :: sha256 c6c4138ee5c86423 :: prose twin of the run farther from 1
\newcommand{\MlstLpQgRthirtyCoarseMassOverInitBnZero}{1.224 and 1.261}  % experiments/FIGS/data/v2/index.json :: linear probe qg, largest set, zero: R30 of R16_Q1_MASS over the BatchNorm-recomputed random network of the same index :: sha256 c6c4138ee5c86423
\newcommand{\MlstLpQgRthirtyCoarseMassOverInitBnZeroPr}{equal within 30\%}  % experiments/FIGS/data/v2/index.json :: linear probe qg, largest set, zero: R30 of R16_Q1_MASS over the BatchNorm-recomputed random network of the same index :: sha256 c6c4138ee5c86423 :: prose twin: two significant figures, direction in words
\newcommand{\MlstLpQgRthirtyCoarseMassOverInitBnZeroPrMax}{equal within 30\%}  % experiments/FIGS/data/v2/index.json :: linear probe qg, largest set, zero: R30 of R16_Q1_MASS over the BatchNorm-recomputed random network of the same index :: sha256 c6c4138ee5c86423 :: prose twin of the run farther from 1
\newcommand{\MlstLpQgRthirtyResOverFull}{0.977\,\ensuremath{\pm}\,0.028}  % experiments/FIGS/data/v2/index.json :: linear probe qg: R30 of L162 over the 188-class run :: sha256 c6c4138ee5c86423
\newcommand{\MlstLpQgRthirtyResOverFullPr}{equal within 3\%}  % experiments/FIGS/data/v2/index.json :: linear probe qg: R30 of L162 over the 188-class run :: sha256 c6c4138ee5c86423 :: prose twin: two significant figures, direction in words
\newcommand{\MlstLpQgRthirtyResOverFullZero}{0.89\,\ensuremath{\pm}\,0.20}  % experiments/FIGS/data/v2/index.json :: linear probe qg: R30 of L162 over the 188-class run :: sha256 c6c4138ee5c86423
\newcommand{\MlstLpQgRthirtyResOverFullZeroPr}{equal within 20\%}  % experiments/FIGS/data/v2/index.json :: linear probe qg: R30 of L162 over the 188-class run :: sha256 c6c4138ee5c86423 :: prose twin: two significant figures, direction in words
\newcommand{\MlstLpQgRthirtyMidOverFull}{0.974\,\ensuremath{\pm}\,0.031}  % experiments/FIGS/data/v2/index.json :: linear probe qg: R30 of R42_Q1 over the 188-class run :: sha256 c6c4138ee5c86423
\newcommand{\MlstLpQgRthirtyMidOverFullPr}{equal within 4\%}  % experiments/FIGS/data/v2/index.json :: linear probe qg: R30 of R42_Q1 over the 188-class run :: sha256 c6c4138ee5c86423 :: prose twin: two significant figures, direction in words
\newcommand{\MlstLpQgRthirtyMidOverFullZero}{0.94\,\ensuremath{\pm}\,0.18}  % experiments/FIGS/data/v2/index.json :: linear probe qg: R30 of R42_Q1 over the 188-class run :: sha256 c6c4138ee5c86423
\newcommand{\MlstLpQgRthirtyMidOverFullZeroPr}{equal within 20\%}  % experiments/FIGS/data/v2/index.json :: linear probe qg: R30 of R42_Q1 over the 188-class run :: sha256 c6c4138ee5c86423 :: prose twin: two significant figures, direction in words
\newcommand{\MlstLpQgRthirtyCoarseOverFull}{1.02\,\ensuremath{\pm}\,0.06}  % experiments/FIGS/data/v2/index.json :: linear probe qg: R30 of R16_Q1 over the 188-class run :: sha256 c6c4138ee5c86423
\newcommand{\MlstLpQgRthirtyCoarseOverFullPr}{equal within 6\%}  % experiments/FIGS/data/v2/index.json :: linear probe qg: R30 of R16_Q1 over the 188-class run :: sha256 c6c4138ee5c86423 :: prose twin: two significant figures, direction in words
\newcommand{\MlstLpQgRthirtyCoarseOverFullZero}{0.95\,\ensuremath{\pm}\,0.09}  % experiments/FIGS/data/v2/index.json :: linear probe qg: R30 of R16_Q1 over the 188-class run :: sha256 c6c4138ee5c86423
\newcommand{\MlstLpQgRthirtyCoarseOverFullZeroPr}{equal within 9\%}  % experiments/FIGS/data/v2/index.json :: linear probe qg: R30 of R16_Q1 over the 188-class run :: sha256 c6c4138ee5c86423 :: prose twin: two significant figures, direction in words
\newcommand{\MlstLpQgRthirtyMeanOverZeroFull}{1.12\,\ensuremath{\pm}\,0.17}  % experiments/FIGS/data/v2/index.json :: linear probe qg: R30 with mean over R30 with zero, L188 :: sha256 c6c4138ee5c86423
\newcommand{\MlstLpQgRthirtyMeanOverZeroFullPr}{equal within 20\%}  % experiments/FIGS/data/v2/index.json :: linear probe qg: R30 with mean over R30 with zero, L188 :: sha256 c6c4138ee5c86423 :: prose twin: two significant figures, direction in words
\newcommand{\MlstLpQgRthirtyDrawOverZeroFull}{1.05\,\ensuremath{\pm}\,0.12}  % experiments/FIGS/data/v2/index.json :: linear probe qg: R30 with draws over R30 with zero, L188 :: sha256 c6c4138ee5c86423
\newcommand{\MlstLpQgRthirtyDrawOverZeroFullPr}{equal within 20\%}  % experiments/FIGS/data/v2/index.json :: linear probe qg: R30 with draws over R30 with zero, L188 :: sha256 c6c4138ee5c86423 :: prose twin: two significant figures, direction in words
\newcommand{\MlstLpQgRthirtyMeanOverZeroInitBn}{0.964\,\ensuremath{\pm}\,0.030}  % experiments/FIGS/data/v2/index.json :: linear probe qg: R30 with mean over R30 with zero, INIT_BN :: sha256 c6c4138ee5c86423
\newcommand{\MlstLpQgRthirtyMeanOverZeroInitBnPr}{equal within 4\%}  % experiments/FIGS/data/v2/index.json :: linear probe qg: R30 with mean over R30 with zero, INIT_BN :: sha256 c6c4138ee5c86423 :: prose twin: two significant figures, direction in words
\newcommand{\MlstLpQgRthirtyDrawOverZeroInitBn}{0.97\,\ensuremath{\pm}\,0.04}  % experiments/FIGS/data/v2/index.json :: linear probe qg: R30 with draws over R30 with zero, INIT_BN :: sha256 c6c4138ee5c86423
\newcommand{\MlstLpQgRthirtyDrawOverZeroInitBnPr}{equal within 4\%}  % experiments/FIGS/data/v2/index.json :: linear probe qg: R30 with draws over R30 with zero, INIT_BN :: sha256 c6c4138ee5c86423 :: prose twin: two significant figures, direction in words
\newcommand{\MlstLpQgRthirtyFullZero}{66\,\ensuremath{\pm}\,7}  % experiments/FIGS/data/v2/index.json :: linear probe qg, largest set: R30 with zero, L188 :: sha256 c6c4138ee5c86423
\newcommand{\MlstLpQgRthirtyFullMean}{72.7\,\ensuremath{\pm}\,2.4}  % experiments/FIGS/data/v2/index.json :: linear probe qg, largest set: R30 with mean, L188 :: sha256 c6c4138ee5c86423
\newcommand{\MlstLpQgRthirtyFullDraws}{68.1\,\ensuremath{\pm}\,1.8}  % experiments/FIGS/data/v2/index.json :: linear probe qg, largest set: R30 with draws, L188 :: sha256 c6c4138ee5c86423
\newcommand{\MlstLpQgRthirtyInitBnZero}{54.2\,\ensuremath{\pm}\,2.4}  % experiments/FIGS/data/v2/index.json :: linear probe qg, largest set: R30 with zero, INIT_BN :: sha256 c6c4138ee5c86423
\newcommand{\MlstLpQgRthirtyInitBnMean}{52.2\,\ensuremath{\pm}\,0.8}  % experiments/FIGS/data/v2/index.json :: linear probe qg, largest set: R30 with mean, INIT_BN :: sha256 c6c4138ee5c86423
\newcommand{\MlstLpQgRthirtyInitBnDraws}{52.4\,\ensuremath{\pm}\,0.4}  % experiments/FIGS/data/v2/index.json :: linear probe qg, largest set: R30 with draws, INIT_BN :: sha256 c6c4138ee5c86423
\newcommand{\MlstLpQgHerwigRthirtyFullOverInitBn}{1.14\,\ensuremath{\pm}\,0.10}  % experiments/FIGS/data/v2/index.json :: linear probe qg_herwig, largest set, chosen substitution: R30 of L188 over the BatchNorm-recomputed random network of the same index :: sha256 c6c4138ee5c86423
\newcommand{\MlstLpQgHerwigRthirtyFullOverInitBnPr}{equal within 20\%}  % experiments/FIGS/data/v2/index.json :: linear probe qg_herwig, largest set, chosen substitution: R30 of L188 over the BatchNorm-recomputed random network of the same index :: sha256 c6c4138ee5c86423 :: prose twin: two significant figures, direction in words
\newcommand{\MlstLpQgHerwigRthirtyFullOverInitBnZero}{1.054\,\ensuremath{\pm}\,0.028}  % experiments/FIGS/data/v2/index.json :: linear probe qg_herwig, largest set, zero: R30 of L188 over the BatchNorm-recomputed random network of the same index :: sha256 c6c4138ee5c86423
\newcommand{\MlstLpQgHerwigRthirtyFullOverInitBnZeroPr}{equal within 6\%}  % experiments/FIGS/data/v2/index.json :: linear probe qg_herwig, largest set, zero: R30 of L188 over the BatchNorm-recomputed random network of the same index :: sha256 c6c4138ee5c86423 :: prose twin: two significant figures, direction in words
\newcommand{\MlstLpQgHerwigRthirtyResOverInitBn}{1.17\,\ensuremath{\pm}\,0.04}  % experiments/FIGS/data/v2/index.json :: linear probe qg_herwig, largest set, chosen substitution: R30 of L162 over the BatchNorm-recomputed random network of the same index :: sha256 c6c4138ee5c86423
\newcommand{\MlstLpQgHerwigRthirtyResOverInitBnPr}{17\% more rejection}  % experiments/FIGS/data/v2/index.json :: linear probe qg_herwig, largest set, chosen substitution: R30 of L162 over the BatchNorm-recomputed random network of the same index :: sha256 c6c4138ee5c86423 :: prose twin: two significant figures, direction in words
\newcommand{\MlstLpQgHerwigRthirtyResOverInitBnZero}{0.98\,\ensuremath{\pm}\,0.10}  % experiments/FIGS/data/v2/index.json :: linear probe qg_herwig, largest set, zero: R30 of L162 over the BatchNorm-recomputed random network of the same index :: sha256 c6c4138ee5c86423
\newcommand{\MlstLpQgHerwigRthirtyResOverInitBnZeroPr}{equal within 10\%}  % experiments/FIGS/data/v2/index.json :: linear probe qg_herwig, largest set, zero: R30 of L162 over the BatchNorm-recomputed random network of the same index :: sha256 c6c4138ee5c86423 :: prose twin: two significant figures, direction in words
\newcommand{\MlstLpQgHerwigRthirtyMidOverInitBn}{1.12\,\ensuremath{\pm}\,0.07}  % experiments/FIGS/data/v2/index.json :: linear probe qg_herwig, largest set, chosen substitution: R30 of R42_Q1 over the BatchNorm-recomputed random network of the same index :: sha256 c6c4138ee5c86423
\newcommand{\MlstLpQgHerwigRthirtyMidOverInitBnPr}{equal within 20\%}  % experiments/FIGS/data/v2/index.json :: linear probe qg_herwig, largest set, chosen substitution: R30 of R42_Q1 over the BatchNorm-recomputed random network of the same index :: sha256 c6c4138ee5c86423 :: prose twin: two significant figures, direction in words
\newcommand{\MlstLpQgHerwigRthirtyMidOverInitBnZero}{0.99\,\ensuremath{\pm}\,0.07}  % experiments/FIGS/data/v2/index.json :: linear probe qg_herwig, largest set, zero: R30 of R42_Q1 over the BatchNorm-recomputed random network of the same index :: sha256 c6c4138ee5c86423
\newcommand{\MlstLpQgHerwigRthirtyMidOverInitBnZeroPr}{equal within 7\%}  % experiments/FIGS/data/v2/index.json :: linear probe qg_herwig, largest set, zero: R30 of R42_Q1 over the BatchNorm-recomputed random network of the same index :: sha256 c6c4138ee5c86423 :: prose twin: two significant figures, direction in words
\newcommand{\MlstLpQgHerwigRthirtyCoarseOverInitBn}{1.154\,\ensuremath{\pm}\,0.015}  % experiments/FIGS/data/v2/index.json :: linear probe qg_herwig, largest set, chosen substitution: R30 of R16_Q1 over the BatchNorm-recomputed random network of the same index :: sha256 c6c4138ee5c86423
\newcommand{\MlstLpQgHerwigRthirtyCoarseOverInitBnPr}{15\% more rejection}  % experiments/FIGS/data/v2/index.json :: linear probe qg_herwig, largest set, chosen substitution: R30 of R16_Q1 over the BatchNorm-recomputed random network of the same index :: sha256 c6c4138ee5c86423 :: prose twin: two significant figures, direction in words
\newcommand{\MlstLpQgHerwigRthirtyCoarseOverInitBnZero}{1.07\,\ensuremath{\pm}\,0.05}  % experiments/FIGS/data/v2/index.json :: linear probe qg_herwig, largest set, zero: R30 of R16_Q1 over the BatchNorm-recomputed random network of the same index :: sha256 c6c4138ee5c86423
\newcommand{\MlstLpQgHerwigRthirtyCoarseOverInitBnZeroPr}{equal within 8\%}  % experiments/FIGS/data/v2/index.json :: linear probe qg_herwig, largest set, zero: R30 of R16_Q1 over the BatchNorm-recomputed random network of the same index :: sha256 c6c4138ee5c86423 :: prose twin: two significant figures, direction in words
\newcommand{\MlstLpQgHerwigRthirtyResMassOverInitBn}{1.152 and 1.138}  % experiments/FIGS/data/v2/index.json :: linear probe qg_herwig, largest set, chosen substitution: R30 of L162_MASS over the BatchNorm-recomputed random network of the same index :: sha256 c6c4138ee5c86423
\newcommand{\MlstLpQgHerwigRthirtyResMassOverInitBnPr}{equal within 20\%}  % experiments/FIGS/data/v2/index.json :: linear probe qg_herwig, largest set, chosen substitution: R30 of L162_MASS over the BatchNorm-recomputed random network of the same index :: sha256 c6c4138ee5c86423 :: prose twin: two significant figures, direction in words
\newcommand{\MlstLpQgHerwigRthirtyResMassOverInitBnPrMax}{equal within 20\%}  % experiments/FIGS/data/v2/index.json :: linear probe qg_herwig, largest set, chosen substitution: R30 of L162_MASS over the BatchNorm-recomputed random network of the same index :: sha256 c6c4138ee5c86423 :: prose twin of the run farther from 1
\newcommand{\MlstLpQgHerwigRthirtyResMassOverInitBnZero}{1.13 and 1.03}  % experiments/FIGS/data/v2/index.json :: linear probe qg_herwig, largest set, zero: R30 of L162_MASS over the BatchNorm-recomputed random network of the same index :: sha256 c6c4138ee5c86423
\newcommand{\MlstLpQgHerwigRthirtyResMassOverInitBnZeroPr}{equal within 20\%}  % experiments/FIGS/data/v2/index.json :: linear probe qg_herwig, largest set, zero: R30 of L162_MASS over the BatchNorm-recomputed random network of the same index :: sha256 c6c4138ee5c86423 :: prose twin: two significant figures, direction in words
\newcommand{\MlstLpQgHerwigRthirtyResMassOverInitBnZeroPrMax}{equal within 20\%}  % experiments/FIGS/data/v2/index.json :: linear probe qg_herwig, largest set, zero: R30 of L162_MASS over the BatchNorm-recomputed random network of the same index :: sha256 c6c4138ee5c86423 :: prose twin of the run farther from 1
\newcommand{\MlstLpQgHerwigRthirtyCoarseMassOverInitBn}{1.146 and 1.159}  % experiments/FIGS/data/v2/index.json :: linear probe qg_herwig, largest set, chosen substitution: R30 of R16_Q1_MASS over the BatchNorm-recomputed random network of the same index :: sha256 c6c4138ee5c86423
\newcommand{\MlstLpQgHerwigRthirtyCoarseMassOverInitBnPr}{equal within 20\%}  % experiments/FIGS/data/v2/index.json :: linear probe qg_herwig, largest set, chosen substitution: R30 of R16_Q1_MASS over the BatchNorm-recomputed random network of the same index :: sha256 c6c4138ee5c86423 :: prose twin: two significant figures, direction in words
\newcommand{\MlstLpQgHerwigRthirtyCoarseMassOverInitBnPrMax}{equal within 20\%}  % experiments/FIGS/data/v2/index.json :: linear probe qg_herwig, largest set, chosen substitution: R30 of R16_Q1_MASS over the BatchNorm-recomputed random network of the same index :: sha256 c6c4138ee5c86423 :: prose twin of the run farther from 1
\newcommand{\MlstLpQgHerwigRthirtyCoarseMassOverInitBnZero}{1.070 and 1.100}  % experiments/FIGS/data/v2/index.json :: linear probe qg_herwig, largest set, zero: R30 of R16_Q1_MASS over the BatchNorm-recomputed random network of the same index :: sha256 c6c4138ee5c86423
\newcommand{\MlstLpQgHerwigRthirtyCoarseMassOverInitBnZeroPr}{equal within 20\%}  % experiments/FIGS/data/v2/index.json :: linear probe qg_herwig, largest set, zero: R30 of R16_Q1_MASS over the BatchNorm-recomputed random network of the same index :: sha256 c6c4138ee5c86423 :: prose twin: two significant figures, direction in words
\newcommand{\MlstLpQgHerwigRthirtyCoarseMassOverInitBnZeroPrMax}{equal within 20\%}  % experiments/FIGS/data/v2/index.json :: linear probe qg_herwig, largest set, zero: R30 of R16_Q1_MASS over the BatchNorm-recomputed random network of the same index :: sha256 c6c4138ee5c86423 :: prose twin of the run farther from 1
\newcommand{\MlstLpQgHerwigRthirtyResOverFull}{1.03\,\ensuremath{\pm}\,0.11}  % experiments/FIGS/data/v2/index.json :: linear probe qg_herwig: R30 of L162 over the 188-class run :: sha256 c6c4138ee5c86423
\newcommand{\MlstLpQgHerwigRthirtyResOverFullPr}{equal within 20\%}  % experiments/FIGS/data/v2/index.json :: linear probe qg_herwig: R30 of L162 over the 188-class run :: sha256 c6c4138ee5c86423 :: prose twin: two significant figures, direction in words
\newcommand{\MlstLpQgHerwigRthirtyResOverFullZero}{0.93\,\ensuremath{\pm}\,0.11}  % experiments/FIGS/data/v2/index.json :: linear probe qg_herwig: R30 of L162 over the 188-class run :: sha256 c6c4138ee5c86423
\newcommand{\MlstLpQgHerwigRthirtyResOverFullZeroPr}{equal within 20\%}  % experiments/FIGS/data/v2/index.json :: linear probe qg_herwig: R30 of L162 over the 188-class run :: sha256 c6c4138ee5c86423 :: prose twin: two significant figures, direction in words
\newcommand{\MlstLpQgHerwigRthirtyMidOverFull}{0.98\,\ensuremath{\pm}\,0.06}  % experiments/FIGS/data/v2/index.json :: linear probe qg_herwig: R30 of R42_Q1 over the 188-class run :: sha256 c6c4138ee5c86423
\newcommand{\MlstLpQgHerwigRthirtyMidOverFullPr}{equal within 6\%}  % experiments/FIGS/data/v2/index.json :: linear probe qg_herwig: R30 of R42_Q1 over the 188-class run :: sha256 c6c4138ee5c86423 :: prose twin: two significant figures, direction in words
\newcommand{\MlstLpQgHerwigRthirtyMidOverFullZero}{0.94\,\ensuremath{\pm}\,0.09}  % experiments/FIGS/data/v2/index.json :: linear probe qg_herwig: R30 of R42_Q1 over the 188-class run :: sha256 c6c4138ee5c86423
\newcommand{\MlstLpQgHerwigRthirtyMidOverFullZeroPr}{equal within 9\%}  % experiments/FIGS/data/v2/index.json :: linear probe qg_herwig: R30 of R42_Q1 over the 188-class run :: sha256 c6c4138ee5c86423 :: prose twin: two significant figures, direction in words
\newcommand{\MlstLpQgHerwigRthirtyCoarseOverFull}{1.02\,\ensuremath{\pm}\,0.09}  % experiments/FIGS/data/v2/index.json :: linear probe qg_herwig: R30 of R16_Q1 over the 188-class run :: sha256 c6c4138ee5c86423
\newcommand{\MlstLpQgHerwigRthirtyCoarseOverFullPr}{equal within 10\%}  % experiments/FIGS/data/v2/index.json :: linear probe qg_herwig: R30 of R16_Q1 over the 188-class run :: sha256 c6c4138ee5c86423 :: prose twin: two significant figures, direction in words
\newcommand{\MlstLpQgHerwigRthirtyCoarseOverFullZero}{1.018\,\ensuremath{\pm}\,0.024}  % experiments/FIGS/data/v2/index.json :: linear probe qg_herwig: R30 of R16_Q1 over the 188-class run :: sha256 c6c4138ee5c86423
\newcommand{\MlstLpQgHerwigRthirtyCoarseOverFullZeroPr}{equal within 3\%}  % experiments/FIGS/data/v2/index.json :: linear probe qg_herwig: R30 of R16_Q1 over the 188-class run :: sha256 c6c4138ee5c86423 :: prose twin: two significant figures, direction in words
\newcommand{\MlstLpQgHerwigRthirtyMeanOverZeroFull}{1.07\,\ensuremath{\pm}\,0.11}  % experiments/FIGS/data/v2/index.json :: linear probe qg_herwig: R30 with mean over R30 with zero, L188 :: sha256 c6c4138ee5c86423
\newcommand{\MlstLpQgHerwigRthirtyMeanOverZeroFullPr}{equal within 20\%}  % experiments/FIGS/data/v2/index.json :: linear probe qg_herwig: R30 with mean over R30 with zero, L188 :: sha256 c6c4138ee5c86423 :: prose twin: two significant figures, direction in words
\newcommand{\MlstLpQgHerwigRthirtyDrawOverZeroFull}{1.07\,\ensuremath{\pm}\,0.06}  % experiments/FIGS/data/v2/index.json :: linear probe qg_herwig: R30 with draws over R30 with zero, L188 :: sha256 c6c4138ee5c86423
\newcommand{\MlstLpQgHerwigRthirtyDrawOverZeroFullPr}{equal within 7\%}  % experiments/FIGS/data/v2/index.json :: linear probe qg_herwig: R30 with draws over R30 with zero, L188 :: sha256 c6c4138ee5c86423 :: prose twin: two significant figures, direction in words
\newcommand{\MlstLpQgHerwigRthirtyMeanOverZeroInitBn}{0.988\,\ensuremath{\pm}\,0.019}  % experiments/FIGS/data/v2/index.json :: linear probe qg_herwig: R30 with mean over R30 with zero, INIT_BN :: sha256 c6c4138ee5c86423
\newcommand{\MlstLpQgHerwigRthirtyMeanOverZeroInitBnPr}{equal within 2\%}  % experiments/FIGS/data/v2/index.json :: linear probe qg_herwig: R30 with mean over R30 with zero, INIT_BN :: sha256 c6c4138ee5c86423 :: prose twin: two significant figures, direction in words
\newcommand{\MlstLpQgHerwigRthirtyDrawOverZeroInitBn}{0.983\,\ensuremath{\pm}\,0.032}  % experiments/FIGS/data/v2/index.json :: linear probe qg_herwig: R30 with draws over R30 with zero, INIT_BN :: sha256 c6c4138ee5c86423
\newcommand{\MlstLpQgHerwigRthirtyDrawOverZeroInitBnPr}{equal within 4\%}  % experiments/FIGS/data/v2/index.json :: linear probe qg_herwig: R30 with draws over R30 with zero, INIT_BN :: sha256 c6c4138ee5c86423 :: prose twin: two significant figures, direction in words
\newcommand{\MlstLpQgHerwigRthirtyFullZero}{32.8\,\ensuremath{\pm}\,1.0}  % experiments/FIGS/data/v2/index.json :: linear probe qg_herwig, largest set: R30 with zero, L188 :: sha256 c6c4138ee5c86423
\newcommand{\MlstLpQgHerwigRthirtyFullMean}{35.0\,\ensuremath{\pm}\,2.6}  % experiments/FIGS/data/v2/index.json :: linear probe qg_herwig, largest set: R30 with mean, L188 :: sha256 c6c4138ee5c86423
\newcommand{\MlstLpQgHerwigRthirtyFullDraws}{35.0\,\ensuremath{\pm}\,0.8}  % experiments/FIGS/data/v2/index.json :: linear probe qg_herwig, largest set: R30 with draws, L188 :: sha256 c6c4138ee5c86423
\newcommand{\MlstLpQgHerwigRthirtyInitBnZero}{31.1\,\ensuremath{\pm}\,0.5}  % experiments/FIGS/data/v2/index.json :: linear probe qg_herwig, largest set: R30 with zero, INIT_BN :: sha256 c6c4138ee5c86423
\newcommand{\MlstLpQgHerwigRthirtyInitBnMean}{30.7\,\ensuremath{\pm}\,0.5}  % experiments/FIGS/data/v2/index.json :: linear probe qg_herwig, largest set: R30 with mean, INIT_BN :: sha256 c6c4138ee5c86423
\newcommand{\MlstLpQgHerwigRthirtyInitBnDraws}{30.6\,\ensuremath{\pm}\,0.6}  % experiments/FIGS/data/v2/index.json :: linear probe qg_herwig, largest set: R30 with draws, INIT_BN :: sha256 c6c4138ee5c86423
\newcommand{\MlstLpJcTwoTopZeroAccDropFull}{88.9\,\ensuremath{\pm}\,1.0\%}  % experiments/FIGS/data/v2/index.json :: linear probe JetClass-II, largest set: relative drop in accuracy (%) with the top data's missing slots set to zero, L188 :: sha256 c6c4138ee5c86423
\newcommand{\MlstLpJcTwoTopZeroAccDropFullPr}{89\%}  % experiments/FIGS/data/v2/index.json :: linear probe JetClass-II, largest set: relative drop in accuracy (%) with the top data's missing slots set to zero, L188 :: sha256 c6c4138ee5c86423 :: prose twin, two significant figures
\newcommand{\MlstLpJcTwoTopMeanAccDropFull}{87.9\,\ensuremath{\pm}\,1.1\%}  % experiments/FIGS/data/v2/index.json :: linear probe JetClass-II, largest set: relative drop in accuracy (%) with the top data's missing slots set to mean, L188 :: sha256 c6c4138ee5c86423
\newcommand{\MlstLpJcTwoTopMeanAccDropFullPr}{88\%}  % experiments/FIGS/data/v2/index.json :: linear probe JetClass-II, largest set: relative drop in accuracy (%) with the top data's missing slots set to mean, L188 :: sha256 c6c4138ee5c86423 :: prose twin, two significant figures
\newcommand{\MlstLpJcTwoTopDrawAccDropFull}{85.98\,\ensuremath{\pm}\,0.07\%}  % experiments/FIGS/data/v2/index.json :: linear probe JetClass-II, largest set: relative drop in accuracy (%) with the top data's missing slots set to draws, L188 :: sha256 c6c4138ee5c86423
\newcommand{\MlstLpJcTwoTopDrawAccDropFullPr}{86\%}  % experiments/FIGS/data/v2/index.json :: linear probe JetClass-II, largest set: relative drop in accuracy (%) with the top data's missing slots set to draws, L188 :: sha256 c6c4138ee5c86423 :: prose twin, two significant figures
\newcommand{\MlstLpJcTwoTopZeroAccDropRes}{89.0\,\ensuremath{\pm}\,0.9\%}  % experiments/FIGS/data/v2/index.json :: linear probe JetClass-II, largest set: relative drop in accuracy (%) with the top data's missing slots set to zero, L162 :: sha256 c6c4138ee5c86423
\newcommand{\MlstLpJcTwoTopZeroAccDropResPr}{89\%}  % experiments/FIGS/data/v2/index.json :: linear probe JetClass-II, largest set: relative drop in accuracy (%) with the top data's missing slots set to zero, L162 :: sha256 c6c4138ee5c86423 :: prose twin, two significant figures
\newcommand{\MlstLpJcTwoTopMeanAccDropRes}{87.8\,\ensuremath{\pm}\,0.7\%}  % experiments/FIGS/data/v2/index.json :: linear probe JetClass-II, largest set: relative drop in accuracy (%) with the top data's missing slots set to mean, L162 :: sha256 c6c4138ee5c86423
\newcommand{\MlstLpJcTwoTopMeanAccDropResPr}{88\%}  % experiments/FIGS/data/v2/index.json :: linear probe JetClass-II, largest set: relative drop in accuracy (%) with the top data's missing slots set to mean, L162 :: sha256 c6c4138ee5c86423 :: prose twin, two significant figures
\newcommand{\MlstLpJcTwoTopDrawAccDropRes}{85.89\,\ensuremath{\pm}\,0.11\%}  % experiments/FIGS/data/v2/index.json :: linear probe JetClass-II, largest set: relative drop in accuracy (%) with the top data's missing slots set to draws, L162 :: sha256 c6c4138ee5c86423
\newcommand{\MlstLpJcTwoTopDrawAccDropResPr}{86\%}  % experiments/FIGS/data/v2/index.json :: linear probe JetClass-II, largest set: relative drop in accuracy (%) with the top data's missing slots set to draws, L162 :: sha256 c6c4138ee5c86423 :: prose twin, two significant figures
\newcommand{\MlstLpJcTwoTopZeroAccDropMid}{87.81\,\ensuremath{\pm}\,0.28\%}  % experiments/FIGS/data/v2/index.json :: linear probe JetClass-II, largest set: relative drop in accuracy (%) with the top data's missing slots set to zero, R42_Q1 :: sha256 c6c4138ee5c86423
\newcommand{\MlstLpJcTwoTopZeroAccDropMidPr}{88\%}  % experiments/FIGS/data/v2/index.json :: linear probe JetClass-II, largest set: relative drop in accuracy (%) with the top data's missing slots set to zero, R42_Q1 :: sha256 c6c4138ee5c86423 :: prose twin, two significant figures
\newcommand{\MlstLpJcTwoTopMeanAccDropMid}{86.7\,\ensuremath{\pm}\,0.9\%}  % experiments/FIGS/data/v2/index.json :: linear probe JetClass-II, largest set: relative drop in accuracy (%) with the top data's missing slots set to mean, R42_Q1 :: sha256 c6c4138ee5c86423
\newcommand{\MlstLpJcTwoTopMeanAccDropMidPr}{87\%}  % experiments/FIGS/data/v2/index.json :: linear probe JetClass-II, largest set: relative drop in accuracy (%) with the top data's missing slots set to mean, R42_Q1 :: sha256 c6c4138ee5c86423 :: prose twin, two significant figures
\newcommand{\MlstLpJcTwoTopDrawAccDropMid}{83.17\,\ensuremath{\pm}\,0.29\%}  % experiments/FIGS/data/v2/index.json :: linear probe JetClass-II, largest set: relative drop in accuracy (%) with the top data's missing slots set to draws, R42_Q1 :: sha256 c6c4138ee5c86423
\newcommand{\MlstLpJcTwoTopDrawAccDropMidPr}{83\%}  % experiments/FIGS/data/v2/index.json :: linear probe JetClass-II, largest set: relative drop in accuracy (%) with the top data's missing slots set to draws, R42_Q1 :: sha256 c6c4138ee5c86423 :: prose twin, two significant figures
\newcommand{\MlstLpJcTwoTopZeroAccDropCoarse}{86.7\,\ensuremath{\pm}\,0.7\%}  % experiments/FIGS/data/v2/index.json :: linear probe JetClass-II, largest set: relative drop in accuracy (%) with the top data's missing slots set to zero, R16_Q1 :: sha256 c6c4138ee5c86423
\newcommand{\MlstLpJcTwoTopZeroAccDropCoarsePr}{87\%}  % experiments/FIGS/data/v2/index.json :: linear probe JetClass-II, largest set: relative drop in accuracy (%) with the top data's missing slots set to zero, R16_Q1 :: sha256 c6c4138ee5c86423 :: prose twin, two significant figures
\newcommand{\MlstLpJcTwoTopMeanAccDropCoarse}{86.1\,\ensuremath{\pm}\,0.8\%}  % experiments/FIGS/data/v2/index.json :: linear probe JetClass-II, largest set: relative drop in accuracy (%) with the top data's missing slots set to mean, R16_Q1 :: sha256 c6c4138ee5c86423
\newcommand{\MlstLpJcTwoTopMeanAccDropCoarsePr}{86\%}  % experiments/FIGS/data/v2/index.json :: linear probe JetClass-II, largest set: relative drop in accuracy (%) with the top data's missing slots set to mean, R16_Q1 :: sha256 c6c4138ee5c86423 :: prose twin, two significant figures
\newcommand{\MlstLpJcTwoTopDrawAccDropCoarse}{81.4\,\ensuremath{\pm}\,0.5\%}  % experiments/FIGS/data/v2/index.json :: linear probe JetClass-II, largest set: relative drop in accuracy (%) with the top data's missing slots set to draws, R16_Q1 :: sha256 c6c4138ee5c86423
\newcommand{\MlstLpJcTwoTopDrawAccDropCoarsePr}{81\%}  % experiments/FIGS/data/v2/index.json :: linear probe JetClass-II, largest set: relative drop in accuracy (%) with the top data's missing slots set to draws, R16_Q1 :: sha256 c6c4138ee5c86423 :: prose twin, two significant figures
\newcommand{\MlstLpJcTwoTopZeroAccDropResMass}{89.2, 88.3\%}  % experiments/FIGS/data/v2/index.json :: linear probe JetClass-II, largest set: relative drop in accuracy (%) with the top data's missing slots set to zero, L162_MASS :: sha256 c6c4138ee5c86423
\newcommand{\MlstLpJcTwoTopZeroAccDropResMassPr}{88--89\%}  % experiments/FIGS/data/v2/index.json :: linear probe JetClass-II, largest set: relative drop in accuracy (%) with the top data's missing slots set to zero, L162_MASS :: sha256 c6c4138ee5c86423 :: prose twin, two significant figures
\newcommand{\MlstLpJcTwoTopMeanAccDropResMass}{88.5, 87.8\%}  % experiments/FIGS/data/v2/index.json :: linear probe JetClass-II, largest set: relative drop in accuracy (%) with the top data's missing slots set to mean, L162_MASS :: sha256 c6c4138ee5c86423
\newcommand{\MlstLpJcTwoTopMeanAccDropResMassPr}{88--88\%}  % experiments/FIGS/data/v2/index.json :: linear probe JetClass-II, largest set: relative drop in accuracy (%) with the top data's missing slots set to mean, L162_MASS :: sha256 c6c4138ee5c86423 :: prose twin, two significant figures
\newcommand{\MlstLpJcTwoTopDrawAccDropResMass}{85.4, 85.4\%}  % experiments/FIGS/data/v2/index.json :: linear probe JetClass-II, largest set: relative drop in accuracy (%) with the top data's missing slots set to draws, L162_MASS :: sha256 c6c4138ee5c86423
\newcommand{\MlstLpJcTwoTopDrawAccDropResMassPr}{85--85\%}  % experiments/FIGS/data/v2/index.json :: linear probe JetClass-II, largest set: relative drop in accuracy (%) with the top data's missing slots set to draws, L162_MASS :: sha256 c6c4138ee5c86423 :: prose twin, two significant figures
\newcommand{\MlstLpJcTwoTopZeroAccDropCoarseMass}{84.8, 87.4\%}  % experiments/FIGS/data/v2/index.json :: linear probe JetClass-II, largest set: relative drop in accuracy (%) with the top data's missing slots set to zero, R16_Q1_MASS :: sha256 c6c4138ee5c86423
\newcommand{\MlstLpJcTwoTopZeroAccDropCoarseMassPr}{85--87\%}  % experiments/FIGS/data/v2/index.json :: linear probe JetClass-II, largest set: relative drop in accuracy (%) with the top data's missing slots set to zero, R16_Q1_MASS :: sha256 c6c4138ee5c86423 :: prose twin, two significant figures
\newcommand{\MlstLpJcTwoTopMeanAccDropCoarseMass}{85.9, 86.9\%}  % experiments/FIGS/data/v2/index.json :: linear probe JetClass-II, largest set: relative drop in accuracy (%) with the top data's missing slots set to mean, R16_Q1_MASS :: sha256 c6c4138ee5c86423
\newcommand{\MlstLpJcTwoTopMeanAccDropCoarseMassPr}{86--87\%}  % experiments/FIGS/data/v2/index.json :: linear probe JetClass-II, largest set: relative drop in accuracy (%) with the top data's missing slots set to mean, R16_Q1_MASS :: sha256 c6c4138ee5c86423 :: prose twin, two significant figures
\newcommand{\MlstLpJcTwoTopDrawAccDropCoarseMass}{79.1, 81.7\%}  % experiments/FIGS/data/v2/index.json :: linear probe JetClass-II, largest set: relative drop in accuracy (%) with the top data's missing slots set to draws, R16_Q1_MASS :: sha256 c6c4138ee5c86423
\newcommand{\MlstLpJcTwoTopDrawAccDropCoarseMassPr}{79--82\%}  % experiments/FIGS/data/v2/index.json :: linear probe JetClass-II, largest set: relative drop in accuracy (%) with the top data's missing slots set to draws, R16_Q1_MASS :: sha256 c6c4138ee5c86423 :: prose twin, two significant figures
\newcommand{\MlstLpJcTwoTopZeroAccDropFzero}{88.9, 88.5\%}  % experiments/FIGS/data/v2/index.json :: linear probe JetClass-II, largest set: relative drop in accuracy (%) with the top data's missing slots set to zero, FLAV_F0 :: sha256 c6c4138ee5c86423
\newcommand{\MlstLpJcTwoTopZeroAccDropFzeroPr}{88--89\%}  % experiments/FIGS/data/v2/index.json :: linear probe JetClass-II, largest set: relative drop in accuracy (%) with the top data's missing slots set to zero, FLAV_F0 :: sha256 c6c4138ee5c86423 :: prose twin, two significant figures
\newcommand{\MlstLpJcTwoTopMeanAccDropFzero}{88.3, 87.0\%}  % experiments/FIGS/data/v2/index.json :: linear probe JetClass-II, largest set: relative drop in accuracy (%) with the top data's missing slots set to mean, FLAV_F0 :: sha256 c6c4138ee5c86423
\newcommand{\MlstLpJcTwoTopMeanAccDropFzeroPr}{87--88\%}  % experiments/FIGS/data/v2/index.json :: linear probe JetClass-II, largest set: relative drop in accuracy (%) with the top data's missing slots set to mean, FLAV_F0 :: sha256 c6c4138ee5c86423 :: prose twin, two significant figures
\newcommand{\MlstLpJcTwoTopDrawAccDropFzero}{83.3, 83.9\%}  % experiments/FIGS/data/v2/index.json :: linear probe JetClass-II, largest set: relative drop in accuracy (%) with the top data's missing slots set to draws, FLAV_F0 :: sha256 c6c4138ee5c86423
\newcommand{\MlstLpJcTwoTopDrawAccDropFzeroPr}{83--84\%}  % experiments/FIGS/data/v2/index.json :: linear probe JetClass-II, largest set: relative drop in accuracy (%) with the top data's missing slots set to draws, FLAV_F0 :: sha256 c6c4138ee5c86423 :: prose twin, two significant figures
\newcommand{\MlstLpJcTwoTopZeroAccDropFone}{87.8, 87.6\%}  % experiments/FIGS/data/v2/index.json :: linear probe JetClass-II, largest set: relative drop in accuracy (%) with the top data's missing slots set to zero, FLAV_F1 :: sha256 c6c4138ee5c86423
\newcommand{\MlstLpJcTwoTopZeroAccDropFonePr}{88--88\%}  % experiments/FIGS/data/v2/index.json :: linear probe JetClass-II, largest set: relative drop in accuracy (%) with the top data's missing slots set to zero, FLAV_F1 :: sha256 c6c4138ee5c86423 :: prose twin, two significant figures
\newcommand{\MlstLpJcTwoTopMeanAccDropFone}{88.9, 86.9\%}  % experiments/FIGS/data/v2/index.json :: linear probe JetClass-II, largest set: relative drop in accuracy (%) with the top data's missing slots set to mean, FLAV_F1 :: sha256 c6c4138ee5c86423
\newcommand{\MlstLpJcTwoTopMeanAccDropFonePr}{87--89\%}  % experiments/FIGS/data/v2/index.json :: linear probe JetClass-II, largest set: relative drop in accuracy (%) with the top data's missing slots set to mean, FLAV_F1 :: sha256 c6c4138ee5c86423 :: prose twin, two significant figures
\newcommand{\MlstLpJcTwoTopDrawAccDropFone}{84.4, 85.2\%}  % experiments/FIGS/data/v2/index.json :: linear probe JetClass-II, largest set: relative drop in accuracy (%) with the top data's missing slots set to draws, FLAV_F1 :: sha256 c6c4138ee5c86423
\newcommand{\MlstLpJcTwoTopDrawAccDropFonePr}{84--85\%}  % experiments/FIGS/data/v2/index.json :: linear probe JetClass-II, largest set: relative drop in accuracy (%) with the top data's missing slots set to draws, FLAV_F1 :: sha256 c6c4138ee5c86423 :: prose twin, two significant figures
\newcommand{\MlstLpJcTwoTopZeroAccDropFoner}{89.6, 90.3\%}  % experiments/FIGS/data/v2/index.json :: linear probe JetClass-II, largest set: relative drop in accuracy (%) with the top data's missing slots set to zero, FLAV_F1R :: sha256 c6c4138ee5c86423
\newcommand{\MlstLpJcTwoTopZeroAccDropFonerPr}{90--90\%}  % experiments/FIGS/data/v2/index.json :: linear probe JetClass-II, largest set: relative drop in accuracy (%) with the top data's missing slots set to zero, FLAV_F1R :: sha256 c6c4138ee5c86423 :: prose twin, two significant figures
\newcommand{\MlstLpJcTwoTopMeanAccDropFoner}{89.4, 89.7\%}  % experiments/FIGS/data/v2/index.json :: linear probe JetClass-II, largest set: relative drop in accuracy (%) with the top data's missing slots set to mean, FLAV_F1R :: sha256 c6c4138ee5c86423
\newcommand{\MlstLpJcTwoTopMeanAccDropFonerPr}{89--90\%}  % experiments/FIGS/data/v2/index.json :: linear probe JetClass-II, largest set: relative drop in accuracy (%) with the top data's missing slots set to mean, FLAV_F1R :: sha256 c6c4138ee5c86423 :: prose twin, two significant figures
\newcommand{\MlstLpJcTwoTopDrawAccDropFoner}{84.9, 85.6\%}  % experiments/FIGS/data/v2/index.json :: linear probe JetClass-II, largest set: relative drop in accuracy (%) with the top data's missing slots set to draws, FLAV_F1R :: sha256 c6c4138ee5c86423
\newcommand{\MlstLpJcTwoTopDrawAccDropFonerPr}{85--86\%}  % experiments/FIGS/data/v2/index.json :: linear probe JetClass-II, largest set: relative drop in accuracy (%) with the top data's missing slots set to draws, FLAV_F1R :: sha256 c6c4138ee5c86423 :: prose twin, two significant figures
\newcommand{\MlstLpJcTwoTopZeroAccDropInitBn}{56.4\,\ensuremath{\pm}\,1.4\%}  % experiments/FIGS/data/v2/index.json :: linear probe JetClass-II, largest set: relative drop in accuracy (%) with the top data's missing slots set to zero, INIT_BN :: sha256 c6c4138ee5c86423
\newcommand{\MlstLpJcTwoTopZeroAccDropInitBnPr}{56\%}  % experiments/FIGS/data/v2/index.json :: linear probe JetClass-II, largest set: relative drop in accuracy (%) with the top data's missing slots set to zero, INIT_BN :: sha256 c6c4138ee5c86423 :: prose twin, two significant figures
\newcommand{\MlstLpJcTwoTopMeanAccDropInitBn}{59.8\,\ensuremath{\pm}\,0.8\%}  % experiments/FIGS/data/v2/index.json :: linear probe JetClass-II, largest set: relative drop in accuracy (%) with the top data's missing slots set to mean, INIT_BN :: sha256 c6c4138ee5c86423
\newcommand{\MlstLpJcTwoTopMeanAccDropInitBnPr}{60\%}  % experiments/FIGS/data/v2/index.json :: linear probe JetClass-II, largest set: relative drop in accuracy (%) with the top data's missing slots set to mean, INIT_BN :: sha256 c6c4138ee5c86423 :: prose twin, two significant figures
\newcommand{\MlstLpJcTwoTopDrawAccDropInitBn}{59.8\,\ensuremath{\pm}\,1.2\%}  % experiments/FIGS/data/v2/index.json :: linear probe JetClass-II, largest set: relative drop in accuracy (%) with the top data's missing slots set to draws, INIT_BN :: sha256 c6c4138ee5c86423
\newcommand{\MlstLpJcTwoTopDrawAccDropInitBnPr}{60\%}  % experiments/FIGS/data/v2/index.json :: linear probe JetClass-II, largest set: relative drop in accuracy (%) with the top data's missing slots set to draws, INIT_BN :: sha256 c6c4138ee5c86423 :: prose twin, two significant figures
\newcommand{\MlstLpJcTwoQgZeroAccDropFull}{61.2\,\ensuremath{\pm}\,3.0\%}  % experiments/FIGS/data/v2/index.json :: linear probe JetClass-II, largest set: relative drop in accuracy (%) with the qg data's missing slots set to zero, L188 :: sha256 c6c4138ee5c86423
\newcommand{\MlstLpJcTwoQgZeroAccDropFullPr}{61\%}  % experiments/FIGS/data/v2/index.json :: linear probe JetClass-II, largest set: relative drop in accuracy (%) with the qg data's missing slots set to zero, L188 :: sha256 c6c4138ee5c86423 :: prose twin, two significant figures
\newcommand{\MlstLpJcTwoQgMeanAccDropFull}{48.8\,\ensuremath{\pm}\,1.1\%}  % experiments/FIGS/data/v2/index.json :: linear probe JetClass-II, largest set: relative drop in accuracy (%) with the qg data's missing slots set to mean, L188 :: sha256 c6c4138ee5c86423
\newcommand{\MlstLpJcTwoQgMeanAccDropFullPr}{49\%}  % experiments/FIGS/data/v2/index.json :: linear probe JetClass-II, largest set: relative drop in accuracy (%) with the qg data's missing slots set to mean, L188 :: sha256 c6c4138ee5c86423 :: prose twin, two significant figures
\newcommand{\MlstLpJcTwoQgDrawAccDropFull}{54.30\,\ensuremath{\pm}\,0.21\%}  % experiments/FIGS/data/v2/index.json :: linear probe JetClass-II, largest set: relative drop in accuracy (%) with the qg data's missing slots set to draws, L188 :: sha256 c6c4138ee5c86423
\newcommand{\MlstLpJcTwoQgDrawAccDropFullPr}{54\%}  % experiments/FIGS/data/v2/index.json :: linear probe JetClass-II, largest set: relative drop in accuracy (%) with the qg data's missing slots set to draws, L188 :: sha256 c6c4138ee5c86423 :: prose twin, two significant figures
\newcommand{\MlstLpJcTwoQgZeroAccDropRes}{64\,\ensuremath{\pm}\,5\%}  % experiments/FIGS/data/v2/index.json :: linear probe JetClass-II, largest set: relative drop in accuracy (%) with the qg data's missing slots set to zero, L162 :: sha256 c6c4138ee5c86423
\newcommand{\MlstLpJcTwoQgZeroAccDropResPr}{64\%}  % experiments/FIGS/data/v2/index.json :: linear probe JetClass-II, largest set: relative drop in accuracy (%) with the qg data's missing slots set to zero, L162 :: sha256 c6c4138ee5c86423 :: prose twin, two significant figures
\newcommand{\MlstLpJcTwoQgMeanAccDropRes}{49.5\,\ensuremath{\pm}\,0.5\%}  % experiments/FIGS/data/v2/index.json :: linear probe JetClass-II, largest set: relative drop in accuracy (%) with the qg data's missing slots set to mean, L162 :: sha256 c6c4138ee5c86423
\newcommand{\MlstLpJcTwoQgMeanAccDropResPr}{49\%}  % experiments/FIGS/data/v2/index.json :: linear probe JetClass-II, largest set: relative drop in accuracy (%) with the qg data's missing slots set to mean, L162 :: sha256 c6c4138ee5c86423 :: prose twin, two significant figures
\newcommand{\MlstLpJcTwoQgDrawAccDropRes}{54.6\,\ensuremath{\pm}\,0.6\%}  % experiments/FIGS/data/v2/index.json :: linear probe JetClass-II, largest set: relative drop in accuracy (%) with the qg data's missing slots set to draws, L162 :: sha256 c6c4138ee5c86423
\newcommand{\MlstLpJcTwoQgDrawAccDropResPr}{55\%}  % experiments/FIGS/data/v2/index.json :: linear probe JetClass-II, largest set: relative drop in accuracy (%) with the qg data's missing slots set to draws, L162 :: sha256 c6c4138ee5c86423 :: prose twin, two significant figures
\newcommand{\MlstLpJcTwoQgZeroAccDropMid}{68\,\ensuremath{\pm}\,5\%}  % experiments/FIGS/data/v2/index.json :: linear probe JetClass-II, largest set: relative drop in accuracy (%) with the qg data's missing slots set to zero, R42_Q1 :: sha256 c6c4138ee5c86423
\newcommand{\MlstLpJcTwoQgZeroAccDropMidPr}{68\%}  % experiments/FIGS/data/v2/index.json :: linear probe JetClass-II, largest set: relative drop in accuracy (%) with the qg data's missing slots set to zero, R42_Q1 :: sha256 c6c4138ee5c86423 :: prose twin, two significant figures
\newcommand{\MlstLpJcTwoQgMeanAccDropMid}{47.6\,\ensuremath{\pm}\,1.6\%}  % experiments/FIGS/data/v2/index.json :: linear probe JetClass-II, largest set: relative drop in accuracy (%) with the qg data's missing slots set to mean, R42_Q1 :: sha256 c6c4138ee5c86423
\newcommand{\MlstLpJcTwoQgMeanAccDropMidPr}{48\%}  % experiments/FIGS/data/v2/index.json :: linear probe JetClass-II, largest set: relative drop in accuracy (%) with the qg data's missing slots set to mean, R42_Q1 :: sha256 c6c4138ee5c86423 :: prose twin, two significant figures
\newcommand{\MlstLpJcTwoQgDrawAccDropMid}{53.5\,\ensuremath{\pm}\,1.4\%}  % experiments/FIGS/data/v2/index.json :: linear probe JetClass-II, largest set: relative drop in accuracy (%) with the qg data's missing slots set to draws, R42_Q1 :: sha256 c6c4138ee5c86423
\newcommand{\MlstLpJcTwoQgDrawAccDropMidPr}{54\%}  % experiments/FIGS/data/v2/index.json :: linear probe JetClass-II, largest set: relative drop in accuracy (%) with the qg data's missing slots set to draws, R42_Q1 :: sha256 c6c4138ee5c86423 :: prose twin, two significant figures
\newcommand{\MlstLpJcTwoQgZeroAccDropCoarse}{70.9\,\ensuremath{\pm}\,2.3\%}  % experiments/FIGS/data/v2/index.json :: linear probe JetClass-II, largest set: relative drop in accuracy (%) with the qg data's missing slots set to zero, R16_Q1 :: sha256 c6c4138ee5c86423
\newcommand{\MlstLpJcTwoQgZeroAccDropCoarsePr}{71\%}  % experiments/FIGS/data/v2/index.json :: linear probe JetClass-II, largest set: relative drop in accuracy (%) with the qg data's missing slots set to zero, R16_Q1 :: sha256 c6c4138ee5c86423 :: prose twin, two significant figures
\newcommand{\MlstLpJcTwoQgMeanAccDropCoarse}{50.0\,\ensuremath{\pm}\,1.0\%}  % experiments/FIGS/data/v2/index.json :: linear probe JetClass-II, largest set: relative drop in accuracy (%) with the qg data's missing slots set to mean, R16_Q1 :: sha256 c6c4138ee5c86423
\newcommand{\MlstLpJcTwoQgMeanAccDropCoarsePr}{50\%}  % experiments/FIGS/data/v2/index.json :: linear probe JetClass-II, largest set: relative drop in accuracy (%) with the qg data's missing slots set to mean, R16_Q1 :: sha256 c6c4138ee5c86423 :: prose twin, two significant figures
\newcommand{\MlstLpJcTwoQgDrawAccDropCoarse}{51.7\,\ensuremath{\pm}\,0.4\%}  % experiments/FIGS/data/v2/index.json :: linear probe JetClass-II, largest set: relative drop in accuracy (%) with the qg data's missing slots set to draws, R16_Q1 :: sha256 c6c4138ee5c86423
\newcommand{\MlstLpJcTwoQgDrawAccDropCoarsePr}{52\%}  % experiments/FIGS/data/v2/index.json :: linear probe JetClass-II, largest set: relative drop in accuracy (%) with the qg data's missing slots set to draws, R16_Q1 :: sha256 c6c4138ee5c86423 :: prose twin, two significant figures
\newcommand{\MlstLpJcTwoQgZeroAccDropResMass}{61.8, 67.9\%}  % experiments/FIGS/data/v2/index.json :: linear probe JetClass-II, largest set: relative drop in accuracy (%) with the qg data's missing slots set to zero, L162_MASS :: sha256 c6c4138ee5c86423
\newcommand{\MlstLpJcTwoQgZeroAccDropResMassPr}{62--68\%}  % experiments/FIGS/data/v2/index.json :: linear probe JetClass-II, largest set: relative drop in accuracy (%) with the qg data's missing slots set to zero, L162_MASS :: sha256 c6c4138ee5c86423 :: prose twin, two significant figures
\newcommand{\MlstLpJcTwoQgMeanAccDropResMass}{49.7, 52.2\%}  % experiments/FIGS/data/v2/index.json :: linear probe JetClass-II, largest set: relative drop in accuracy (%) with the qg data's missing slots set to mean, L162_MASS :: sha256 c6c4138ee5c86423
\newcommand{\MlstLpJcTwoQgMeanAccDropResMassPr}{50--52\%}  % experiments/FIGS/data/v2/index.json :: linear probe JetClass-II, largest set: relative drop in accuracy (%) with the qg data's missing slots set to mean, L162_MASS :: sha256 c6c4138ee5c86423 :: prose twin, two significant figures
\newcommand{\MlstLpJcTwoQgDrawAccDropResMass}{54.3, 55.1\%}  % experiments/FIGS/data/v2/index.json :: linear probe JetClass-II, largest set: relative drop in accuracy (%) with the qg data's missing slots set to draws, L162_MASS :: sha256 c6c4138ee5c86423
\newcommand{\MlstLpJcTwoQgDrawAccDropResMassPr}{54--55\%}  % experiments/FIGS/data/v2/index.json :: linear probe JetClass-II, largest set: relative drop in accuracy (%) with the qg data's missing slots set to draws, L162_MASS :: sha256 c6c4138ee5c86423 :: prose twin, two significant figures
\newcommand{\MlstLpJcTwoQgZeroAccDropCoarseMass}{66.7, 65.4\%}  % experiments/FIGS/data/v2/index.json :: linear probe JetClass-II, largest set: relative drop in accuracy (%) with the qg data's missing slots set to zero, R16_Q1_MASS :: sha256 c6c4138ee5c86423
\newcommand{\MlstLpJcTwoQgZeroAccDropCoarseMassPr}{65--67\%}  % experiments/FIGS/data/v2/index.json :: linear probe JetClass-II, largest set: relative drop in accuracy (%) with the qg data's missing slots set to zero, R16_Q1_MASS :: sha256 c6c4138ee5c86423 :: prose twin, two significant figures
\newcommand{\MlstLpJcTwoQgMeanAccDropCoarseMass}{49.0, 49.6\%}  % experiments/FIGS/data/v2/index.json :: linear probe JetClass-II, largest set: relative drop in accuracy (%) with the qg data's missing slots set to mean, R16_Q1_MASS :: sha256 c6c4138ee5c86423
\newcommand{\MlstLpJcTwoQgMeanAccDropCoarseMassPr}{49--50\%}  % experiments/FIGS/data/v2/index.json :: linear probe JetClass-II, largest set: relative drop in accuracy (%) with the qg data's missing slots set to mean, R16_Q1_MASS :: sha256 c6c4138ee5c86423 :: prose twin, two significant figures
\newcommand{\MlstLpJcTwoQgDrawAccDropCoarseMass}{50.8, 51.1\%}  % experiments/FIGS/data/v2/index.json :: linear probe JetClass-II, largest set: relative drop in accuracy (%) with the qg data's missing slots set to draws, R16_Q1_MASS :: sha256 c6c4138ee5c86423
\newcommand{\MlstLpJcTwoQgDrawAccDropCoarseMassPr}{51--51\%}  % experiments/FIGS/data/v2/index.json :: linear probe JetClass-II, largest set: relative drop in accuracy (%) with the qg data's missing slots set to draws, R16_Q1_MASS :: sha256 c6c4138ee5c86423 :: prose twin, two significant figures
\newcommand{\MlstLpJcTwoQgZeroAccDropFzero}{62.7, 62.2\%}  % experiments/FIGS/data/v2/index.json :: linear probe JetClass-II, largest set: relative drop in accuracy (%) with the qg data's missing slots set to zero, FLAV_F0 :: sha256 c6c4138ee5c86423
\newcommand{\MlstLpJcTwoQgZeroAccDropFzeroPr}{62--63\%}  % experiments/FIGS/data/v2/index.json :: linear probe JetClass-II, largest set: relative drop in accuracy (%) with the qg data's missing slots set to zero, FLAV_F0 :: sha256 c6c4138ee5c86423 :: prose twin, two significant figures
\newcommand{\MlstLpJcTwoQgMeanAccDropFzero}{51.9, 51.3\%}  % experiments/FIGS/data/v2/index.json :: linear probe JetClass-II, largest set: relative drop in accuracy (%) with the qg data's missing slots set to mean, FLAV_F0 :: sha256 c6c4138ee5c86423
\newcommand{\MlstLpJcTwoQgMeanAccDropFzeroPr}{51--52\%}  % experiments/FIGS/data/v2/index.json :: linear probe JetClass-II, largest set: relative drop in accuracy (%) with the qg data's missing slots set to mean, FLAV_F0 :: sha256 c6c4138ee5c86423 :: prose twin, two significant figures
\newcommand{\MlstLpJcTwoQgDrawAccDropFzero}{53.9, 53.4\%}  % experiments/FIGS/data/v2/index.json :: linear probe JetClass-II, largest set: relative drop in accuracy (%) with the qg data's missing slots set to draws, FLAV_F0 :: sha256 c6c4138ee5c86423
\newcommand{\MlstLpJcTwoQgDrawAccDropFzeroPr}{53--54\%}  % experiments/FIGS/data/v2/index.json :: linear probe JetClass-II, largest set: relative drop in accuracy (%) with the qg data's missing slots set to draws, FLAV_F0 :: sha256 c6c4138ee5c86423 :: prose twin, two significant figures
\newcommand{\MlstLpJcTwoQgZeroAccDropFone}{66.2, 62.8\%}  % experiments/FIGS/data/v2/index.json :: linear probe JetClass-II, largest set: relative drop in accuracy (%) with the qg data's missing slots set to zero, FLAV_F1 :: sha256 c6c4138ee5c86423
\newcommand{\MlstLpJcTwoQgZeroAccDropFonePr}{63--66\%}  % experiments/FIGS/data/v2/index.json :: linear probe JetClass-II, largest set: relative drop in accuracy (%) with the qg data's missing slots set to zero, FLAV_F1 :: sha256 c6c4138ee5c86423 :: prose twin, two significant figures
\newcommand{\MlstLpJcTwoQgMeanAccDropFone}{52.7, 52.8\%}  % experiments/FIGS/data/v2/index.json :: linear probe JetClass-II, largest set: relative drop in accuracy (%) with the qg data's missing slots set to mean, FLAV_F1 :: sha256 c6c4138ee5c86423
\newcommand{\MlstLpJcTwoQgMeanAccDropFonePr}{53--53\%}  % experiments/FIGS/data/v2/index.json :: linear probe JetClass-II, largest set: relative drop in accuracy (%) with the qg data's missing slots set to mean, FLAV_F1 :: sha256 c6c4138ee5c86423 :: prose twin, two significant figures
\newcommand{\MlstLpJcTwoQgDrawAccDropFone}{55.6, 56.5\%}  % experiments/FIGS/data/v2/index.json :: linear probe JetClass-II, largest set: relative drop in accuracy (%) with the qg data's missing slots set to draws, FLAV_F1 :: sha256 c6c4138ee5c86423
\newcommand{\MlstLpJcTwoQgDrawAccDropFonePr}{56--56\%}  % experiments/FIGS/data/v2/index.json :: linear probe JetClass-II, largest set: relative drop in accuracy (%) with the qg data's missing slots set to draws, FLAV_F1 :: sha256 c6c4138ee5c86423 :: prose twin, two significant figures
\newcommand{\MlstLpJcTwoQgZeroAccDropFoner}{67.2, 69.0\%}  % experiments/FIGS/data/v2/index.json :: linear probe JetClass-II, largest set: relative drop in accuracy (%) with the qg data's missing slots set to zero, FLAV_F1R :: sha256 c6c4138ee5c86423
\newcommand{\MlstLpJcTwoQgZeroAccDropFonerPr}{67--69\%}  % experiments/FIGS/data/v2/index.json :: linear probe JetClass-II, largest set: relative drop in accuracy (%) with the qg data's missing slots set to zero, FLAV_F1R :: sha256 c6c4138ee5c86423 :: prose twin, two significant figures
\newcommand{\MlstLpJcTwoQgMeanAccDropFoner}{52.9, 51.8\%}  % experiments/FIGS/data/v2/index.json :: linear probe JetClass-II, largest set: relative drop in accuracy (%) with the qg data's missing slots set to mean, FLAV_F1R :: sha256 c6c4138ee5c86423
\newcommand{\MlstLpJcTwoQgMeanAccDropFonerPr}{52--53\%}  % experiments/FIGS/data/v2/index.json :: linear probe JetClass-II, largest set: relative drop in accuracy (%) with the qg data's missing slots set to mean, FLAV_F1R :: sha256 c6c4138ee5c86423 :: prose twin, two significant figures
\newcommand{\MlstLpJcTwoQgDrawAccDropFoner}{57.2, 55.4\%}  % experiments/FIGS/data/v2/index.json :: linear probe JetClass-II, largest set: relative drop in accuracy (%) with the qg data's missing slots set to draws, FLAV_F1R :: sha256 c6c4138ee5c86423
\newcommand{\MlstLpJcTwoQgDrawAccDropFonerPr}{55--57\%}  % experiments/FIGS/data/v2/index.json :: linear probe JetClass-II, largest set: relative drop in accuracy (%) with the qg data's missing slots set to draws, FLAV_F1R :: sha256 c6c4138ee5c86423 :: prose twin, two significant figures
\newcommand{\MlstLpJcTwoQgZeroAccDropInitBn}{8.7\,\ensuremath{\pm}\,2.6\%}  % experiments/FIGS/data/v2/index.json :: linear probe JetClass-II, largest set: relative drop in accuracy (%) with the qg data's missing slots set to zero, INIT_BN :: sha256 c6c4138ee5c86423
\newcommand{\MlstLpJcTwoQgZeroAccDropInitBnPr}{8.7\%}  % experiments/FIGS/data/v2/index.json :: linear probe JetClass-II, largest set: relative drop in accuracy (%) with the qg data's missing slots set to zero, INIT_BN :: sha256 c6c4138ee5c86423 :: prose twin, two significant figures
\newcommand{\MlstLpJcTwoQgMeanAccDropInitBn}{10.4\,\ensuremath{\pm}\,1.4\%}  % experiments/FIGS/data/v2/index.json :: linear probe JetClass-II, largest set: relative drop in accuracy (%) with the qg data's missing slots set to mean, INIT_BN :: sha256 c6c4138ee5c86423
\newcommand{\MlstLpJcTwoQgMeanAccDropInitBnPr}{10\%}  % experiments/FIGS/data/v2/index.json :: linear probe JetClass-II, largest set: relative drop in accuracy (%) with the qg data's missing slots set to mean, INIT_BN :: sha256 c6c4138ee5c86423 :: prose twin, two significant figures
\newcommand{\MlstLpJcTwoQgDrawAccDropInitBn}{11.4\,\ensuremath{\pm}\,2.1\%}  % experiments/FIGS/data/v2/index.json :: linear probe JetClass-II, largest set: relative drop in accuracy (%) with the qg data's missing slots set to draws, INIT_BN :: sha256 c6c4138ee5c86423
\newcommand{\MlstLpJcTwoQgDrawAccDropInitBnPr}{11\%}  % experiments/FIGS/data/v2/index.json :: linear probe JetClass-II, largest set: relative drop in accuracy (%) with the qg data's missing slots set to draws, INIT_BN :: sha256 c6c4138ee5c86423 :: prose twin, two significant figures
\newcommand{\MlstMassLambda}{5}  % configs/arms/v2_grid.json :: arms[L162_MASS, R16_Q1_MASS].mass_lambda :: sha256 d8a31b1486992302 :: weight of the mass-regression term in the pretraining loss, both mass models
\newcommand{\MlstMassLossShareResMass}{13.78 and 13.87}  % experiments/FIGS/data/v2/index.json :: share of the mass-regression term in the training loss of L162_MASS, mean over epochs, percent :: sha256 c6c4138ee5c86423
\newcommand{\MlstMassLossShareCoarseMass}{28.65 and 28.67}  % experiments/FIGS/data/v2/index.json :: share of the mass-regression term in the training loss of R16_Q1_MASS, mean over epochs, percent :: sha256 c6c4138ee5c86423
\newcommand{\MlstMassUncorrected}{0.245}  % experiments/FIGS/data/v2/index.json :: head/atlas_res: groomed mass, all jets, ATLAS resolution :: sha256 c6c4138ee5c86423
\newcommand{\MlstMassUncorrectedResponse}{0.929}  % experiments/FIGS/data/v2/index.json :: head/atlas_res: groomed mass, all jets, median response :: sha256 c6c4138ee5c86423
\newcommand{\MlstMassHeadResMass}{0.1471 and 0.1463}  % experiments/FIGS/data/v2/index.json :: regression output of L162_MASS, all jets, ATLAS resolution :: sha256 c6c4138ee5c86423
\newcommand{\MlstMassHeadResMassResponse}{0.995 and 1.003}  % experiments/FIGS/data/v2/index.json :: regression output of L162_MASS, all jets, median response :: sha256 c6c4138ee5c86423
\newcommand{\MlstMassHeadResMassOverUncorr}{0.6009 and 0.5976}  % experiments/FIGS/data/v2/index.json :: resolution of the regression output of L162_MASS over the groomed mass's, all jets :: sha256 c6c4138ee5c86423
\newcommand{\MlstMassHeadResMassOverUncorrPr}{at most 0.60 times the resolution}  % experiments/FIGS/data/v2/index.json :: resolution of the regression output of L162_MASS over the groomed mass's, all jets :: sha256 c6c4138ee5c86423 :: prose twin: two significant figures, direction in words
\newcommand{\MlstMassHeadResMassOverUncorrPrMax}{down to 0.60 times the resolution}  % experiments/FIGS/data/v2/index.json :: resolution of the regression output of L162_MASS over the groomed mass's, all jets :: sha256 c6c4138ee5c86423 :: prose twin of the run farther from 1
\newcommand{\MlstMassHeadResMassOverMlpRes}{0.6645 and 0.6691}  % experiments/FIGS/data/v2/index.json :: resolution of the regression output of L162_MASS over the frozen-feature MLP of L162, run by run, all jets :: sha256 c6c4138ee5c86423
\newcommand{\MlstMassHeadResMassOverMlpResPr}{at most 0.67 times the resolution}  % experiments/FIGS/data/v2/index.json :: resolution of the regression output of L162_MASS over the frozen-feature MLP of L162, run by run, all jets :: sha256 c6c4138ee5c86423 :: prose twin: two significant figures, direction in words
\newcommand{\MlstMassHeadResMassOverMlpResPrMax}{down to 0.66 times the resolution}  % experiments/FIGS/data/v2/index.json :: resolution of the regression output of L162_MASS over the frozen-feature MLP of L162, run by run, all jets :: sha256 c6c4138ee5c86423 :: prose twin of the run farther from 1
\newcommand{\MlstMassHeadCoarseMass}{0.1450 and 0.1461}  % experiments/FIGS/data/v2/index.json :: regression output of R16_Q1_MASS, all jets, ATLAS resolution :: sha256 c6c4138ee5c86423
\newcommand{\MlstMassHeadCoarseMassResponse}{1.005 and 1.005}  % experiments/FIGS/data/v2/index.json :: regression output of R16_Q1_MASS, all jets, median response :: sha256 c6c4138ee5c86423
\newcommand{\MlstMassHeadCoarseMassOverUncorr}{0.5921 and 0.5966}  % experiments/FIGS/data/v2/index.json :: resolution of the regression output of R16_Q1_MASS over the groomed mass's, all jets :: sha256 c6c4138ee5c86423
\newcommand{\MlstMassHeadCoarseMassOverUncorrPr}{at most 0.60 times the resolution}  % experiments/FIGS/data/v2/index.json :: resolution of the regression output of R16_Q1_MASS over the groomed mass's, all jets :: sha256 c6c4138ee5c86423 :: prose twin: two significant figures, direction in words
\newcommand{\MlstMassHeadCoarseMassOverUncorrPrMax}{down to 0.59 times the resolution}  % experiments/FIGS/data/v2/index.json :: resolution of the regression output of R16_Q1_MASS over the groomed mass's, all jets :: sha256 c6c4138ee5c86423 :: prose twin of the run farther from 1
\newcommand{\MlstMassHeadCoarseMassOverMlpCoarse}{0.703 and 0.690}  % experiments/FIGS/data/v2/index.json :: resolution of the regression output of R16_Q1_MASS over the frozen-feature MLP of R16_Q1, run by run, all jets :: sha256 c6c4138ee5c86423
\newcommand{\MlstMassHeadCoarseMassOverMlpCoarsePr}{at most 0.70 times the resolution}  % experiments/FIGS/data/v2/index.json :: resolution of the regression output of R16_Q1_MASS over the frozen-feature MLP of R16_Q1, run by run, all jets :: sha256 c6c4138ee5c86423 :: prose twin: two significant figures, direction in words
\newcommand{\MlstMassHeadCoarseMassOverMlpCoarsePrMax}{down to 0.69 times the resolution}  % experiments/FIGS/data/v2/index.json :: resolution of the regression output of R16_Q1_MASS over the frozen-feature MLP of R16_Q1, run by run, all jets :: sha256 c6c4138ee5c86423 :: prose twin of the run farther from 1
\newcommand{\MlstMassMlpInit}{0.2421\,\ensuremath{\pm}\,0.0018}  % experiments/FIGS/data/v2/index.json :: frozen-feature MLP of INIT, all jets, ATLAS resolution :: sha256 c6c4138ee5c86423
\newcommand{\MlstMassMlpFull}{0.2200\,\ensuremath{\pm}\,0.0014}  % experiments/FIGS/data/v2/index.json :: frozen-feature MLP of L188, all jets, ATLAS resolution :: sha256 c6c4138ee5c86423
\newcommand{\MlstMassMlpRes}{0.2195\,\ensuremath{\pm}\,0.0017}  % experiments/FIGS/data/v2/index.json :: frozen-feature MLP of L162, all jets, ATLAS resolution :: sha256 c6c4138ee5c86423
\newcommand{\MlstMassMlpMid}{0.2158\,\ensuremath{\pm}\,0.0029}  % experiments/FIGS/data/v2/index.json :: frozen-feature MLP of R42_Q1, all jets, ATLAS resolution :: sha256 c6c4138ee5c86423
\newcommand{\MlstMassMlpCoarse}{0.2098\,\ensuremath{\pm}\,0.0031}  % experiments/FIGS/data/v2/index.json :: frozen-feature MLP of R16_Q1, all jets, ATLAS resolution :: sha256 c6c4138ee5c86423
\newcommand{\MlstMassMlpResMass}{0.1591 and 0.1586}  % experiments/FIGS/data/v2/index.json :: frozen-feature MLP of L162_MASS, all jets, ATLAS resolution :: sha256 c6c4138ee5c86423
\newcommand{\MlstMassMlpCoarseMass}{0.1583 and 0.1588}  % experiments/FIGS/data/v2/index.json :: frozen-feature MLP of R16_Q1_MASS, all jets, ATLAS resolution :: sha256 c6c4138ee5c86423
\newcommand{\MlstCmpMlpResOverFull}{0.998\,\ensuremath{\pm}\,0.006}  % experiments/FIGS/data/v2/index.json :: frozen-feature MLP resolution of L162 over L188, run by run, all jets :: sha256 c6c4138ee5c86423
\newcommand{\MlstCmpMlpResOverFullPr}{equal within 1\%}  % experiments/FIGS/data/v2/index.json :: frozen-feature MLP resolution of L162 over L188, run by run, all jets :: sha256 c6c4138ee5c86423 :: prose twin: two significant figures, direction in words
\newcommand{\MlstCmpMlpMidOverFull}{0.981\,\ensuremath{\pm}\,0.017}  % experiments/FIGS/data/v2/index.json :: frozen-feature MLP resolution of R42_Q1 over L188, run by run, all jets :: sha256 c6c4138ee5c86423
\newcommand{\MlstCmpMlpMidOverFullPr}{equal within 2\%}  % experiments/FIGS/data/v2/index.json :: frozen-feature MLP resolution of R42_Q1 over L188, run by run, all jets :: sha256 c6c4138ee5c86423 :: prose twin: two significant figures, direction in words
\newcommand{\MlstCmpMlpCoarseOverFull}{0.954\,\ensuremath{\pm}\,0.018}  % experiments/FIGS/data/v2/index.json :: frozen-feature MLP resolution of R16_Q1 over L188, run by run, all jets :: sha256 c6c4138ee5c86423
\newcommand{\MlstCmpMlpCoarseOverFullPr}{0.95 times the resolution}  % experiments/FIGS/data/v2/index.json :: frozen-feature MLP resolution of R16_Q1 over L188, run by run, all jets :: sha256 c6c4138ee5c86423 :: prose twin: two significant figures, direction in words
\newcommand{\MlstCmpMlpResMassOverRes}{0.7186 and 0.7253}  % experiments/FIGS/data/v2/index.json :: frozen-feature MLP resolution of L162_MASS over L162, run by run, all jets :: sha256 c6c4138ee5c86423
\newcommand{\MlstCmpMlpResMassOverResPr}{at most 0.73 times the resolution}  % experiments/FIGS/data/v2/index.json :: frozen-feature MLP resolution of L162_MASS over L162, run by run, all jets :: sha256 c6c4138ee5c86423 :: prose twin: two significant figures, direction in words
\newcommand{\MlstCmpMlpResMassOverResPrMax}{down to 0.72 times the resolution}  % experiments/FIGS/data/v2/index.json :: frozen-feature MLP resolution of L162_MASS over L162, run by run, all jets :: sha256 c6c4138ee5c86423 :: prose twin of the run farther from 1
\newcommand{\MlstCmpMlpCoarseMassOverCoarse}{0.768 and 0.750}  % experiments/FIGS/data/v2/index.json :: frozen-feature MLP resolution of R16_Q1_MASS over R16_Q1, run by run, all jets :: sha256 c6c4138ee5c86423
\newcommand{\MlstCmpMlpCoarseMassOverCoarsePr}{at most 0.77 times the resolution}  % experiments/FIGS/data/v2/index.json :: frozen-feature MLP resolution of R16_Q1_MASS over R16_Q1, run by run, all jets :: sha256 c6c4138ee5c86423 :: prose twin: two significant figures, direction in words
\newcommand{\MlstCmpMlpCoarseMassOverCoarsePrMax}{down to 0.75 times the resolution}  % experiments/FIGS/data/v2/index.json :: frozen-feature MLP resolution of R16_Q1_MASS over R16_Q1, run by run, all jets :: sha256 c6c4138ee5c86423 :: prose twin of the run farther from 1
\newcommand{\MlstMassUncorrectedResonant}{0.213}  % experiments/FIGS/data/v2/index.json :: head/atlas_res: groomed mass, resonant jets, ATLAS resolution :: sha256 c6c4138ee5c86423
\newcommand{\MlstMassUncorrectedResponseResonant}{0.929}  % experiments/FIGS/data/v2/index.json :: head/atlas_res: groomed mass, resonant jets, median response :: sha256 c6c4138ee5c86423
\newcommand{\MlstMassHeadResMassResonant}{0.1332 and 0.1320}  % experiments/FIGS/data/v2/index.json :: regression output of L162_MASS, resonant jets, ATLAS resolution :: sha256 c6c4138ee5c86423
\newcommand{\MlstMassHeadResMassResponseResonant}{0.996 and 1.004}  % experiments/FIGS/data/v2/index.json :: regression output of L162_MASS, resonant jets, median response :: sha256 c6c4138ee5c86423
\newcommand{\MlstMassHeadResMassOverUncorrResonant}{0.6259 and 0.6205}  % experiments/FIGS/data/v2/index.json :: resolution of the regression output of L162_MASS over the groomed mass's, resonant jets :: sha256 c6c4138ee5c86423
\newcommand{\MlstMassHeadResMassOverUncorrResonantPr}{at most 0.63 times the resolution}  % experiments/FIGS/data/v2/index.json :: resolution of the regression output of L162_MASS over the groomed mass's, resonant jets :: sha256 c6c4138ee5c86423 :: prose twin: two significant figures, direction in words
\newcommand{\MlstMassHeadResMassOverUncorrResonantPrMax}{down to 0.62 times the resolution}  % experiments/FIGS/data/v2/index.json :: resolution of the regression output of L162_MASS over the groomed mass's, resonant jets :: sha256 c6c4138ee5c86423 :: prose twin of the run farther from 1
\newcommand{\MlstMassHeadResMassOverMlpResResonant}{0.6620 and 0.6662}  % experiments/FIGS/data/v2/index.json :: resolution of the regression output of L162_MASS over the frozen-feature MLP of L162, run by run, resonant jets :: sha256 c6c4138ee5c86423
\newcommand{\MlstMassHeadResMassOverMlpResResonantPr}{at most 0.67 times the resolution}  % experiments/FIGS/data/v2/index.json :: resolution of the regression output of L162_MASS over the frozen-feature MLP of L162, run by run, resonant jets :: sha256 c6c4138ee5c86423 :: prose twin: two significant figures, direction in words
\newcommand{\MlstMassHeadResMassOverMlpResResonantPrMax}{down to 0.66 times the resolution}  % experiments/FIGS/data/v2/index.json :: resolution of the regression output of L162_MASS over the frozen-feature MLP of L162, run by run, resonant jets :: sha256 c6c4138ee5c86423 :: prose twin of the run farther from 1
\newcommand{\MlstMassHeadCoarseMassResonant}{0.1312 and 0.1320}  % experiments/FIGS/data/v2/index.json :: regression output of R16_Q1_MASS, resonant jets, ATLAS resolution :: sha256 c6c4138ee5c86423
\newcommand{\MlstMassHeadCoarseMassResponseResonant}{1.006 and 1.006}  % experiments/FIGS/data/v2/index.json :: regression output of R16_Q1_MASS, resonant jets, median response :: sha256 c6c4138ee5c86423
\newcommand{\MlstMassHeadCoarseMassOverUncorrResonant}{0.6165 and 0.6206}  % experiments/FIGS/data/v2/index.json :: resolution of the regression output of R16_Q1_MASS over the groomed mass's, resonant jets :: sha256 c6c4138ee5c86423
\newcommand{\MlstMassHeadCoarseMassOverUncorrResonantPr}{at most 0.62 times the resolution}  % experiments/FIGS/data/v2/index.json :: resolution of the regression output of R16_Q1_MASS over the groomed mass's, resonant jets :: sha256 c6c4138ee5c86423 :: prose twin: two significant figures, direction in words
\newcommand{\MlstMassHeadCoarseMassOverUncorrResonantPrMax}{down to 0.62 times the resolution}  % experiments/FIGS/data/v2/index.json :: resolution of the regression output of R16_Q1_MASS over the groomed mass's, resonant jets :: sha256 c6c4138ee5c86423 :: prose twin of the run farther from 1
\newcommand{\MlstMassHeadCoarseMassOverMlpCoarseResonant}{0.707 and 0.688}  % experiments/FIGS/data/v2/index.json :: resolution of the regression output of R16_Q1_MASS over the frozen-feature MLP of R16_Q1, run by run, resonant jets :: sha256 c6c4138ee5c86423
\newcommand{\MlstMassHeadCoarseMassOverMlpCoarseResonantPr}{at most 0.71 times the resolution}  % experiments/FIGS/data/v2/index.json :: resolution of the regression output of R16_Q1_MASS over the frozen-feature MLP of R16_Q1, run by run, resonant jets :: sha256 c6c4138ee5c86423 :: prose twin: two significant figures, direction in words
\newcommand{\MlstMassHeadCoarseMassOverMlpCoarseResonantPrMax}{down to 0.69 times the resolution}  % experiments/FIGS/data/v2/index.json :: resolution of the regression output of R16_Q1_MASS over the frozen-feature MLP of R16_Q1, run by run, resonant jets :: sha256 c6c4138ee5c86423 :: prose twin of the run farther from 1
\newcommand{\MlstMassMlpInitResonant}{0.2214\,\ensuremath{\pm}\,0.0017}  % experiments/FIGS/data/v2/index.json :: frozen-feature MLP of INIT, resonant jets, ATLAS resolution :: sha256 c6c4138ee5c86423
\newcommand{\MlstMassMlpFullResonant}{0.1990\,\ensuremath{\pm}\,0.0016}  % experiments/FIGS/data/v2/index.json :: frozen-feature MLP of L188, resonant jets, ATLAS resolution :: sha256 c6c4138ee5c86423
\newcommand{\MlstMassMlpResResonant}{0.1989\,\ensuremath{\pm}\,0.0020}  % experiments/FIGS/data/v2/index.json :: frozen-feature MLP of L162, resonant jets, ATLAS resolution :: sha256 c6c4138ee5c86423
\newcommand{\MlstMassMlpMidResonant}{0.1949\,\ensuremath{\pm}\,0.0023}  % experiments/FIGS/data/v2/index.json :: frozen-feature MLP of R42_Q1, resonant jets, ATLAS resolution :: sha256 c6c4138ee5c86423
\newcommand{\MlstMassMlpCoarseResonant}{0.1894\,\ensuremath{\pm}\,0.0035}  % experiments/FIGS/data/v2/index.json :: frozen-feature MLP of R16_Q1, resonant jets, ATLAS resolution :: sha256 c6c4138ee5c86423
\newcommand{\MlstMassMlpResMassResonant}{0.1436 and 0.1430}  % experiments/FIGS/data/v2/index.json :: frozen-feature MLP of L162_MASS, resonant jets, ATLAS resolution :: sha256 c6c4138ee5c86423
\newcommand{\MlstMassMlpCoarseMassResonant}{0.1420 and 0.1431}  % experiments/FIGS/data/v2/index.json :: frozen-feature MLP of R16_Q1_MASS, resonant jets, ATLAS resolution :: sha256 c6c4138ee5c86423
\newcommand{\MlstCmpMlpResOverFullResonant}{1.000\,\ensuremath{\pm}\,0.005}  % experiments/FIGS/data/v2/index.json :: frozen-feature MLP resolution of L162 over L188, run by run, resonant jets :: sha256 c6c4138ee5c86423
\newcommand{\MlstCmpMlpResOverFullResonantPr}{equal within 1\%}  % experiments/FIGS/data/v2/index.json :: frozen-feature MLP resolution of L162 over L188, run by run, resonant jets :: sha256 c6c4138ee5c86423 :: prose twin: two significant figures, direction in words
\newcommand{\MlstCmpMlpMidOverFullResonant}{0.980\,\ensuremath{\pm}\,0.018}  % experiments/FIGS/data/v2/index.json :: frozen-feature MLP resolution of R42_Q1 over L188, run by run, resonant jets :: sha256 c6c4138ee5c86423
\newcommand{\MlstCmpMlpMidOverFullResonantPr}{equal within 3\%}  % experiments/FIGS/data/v2/index.json :: frozen-feature MLP resolution of R42_Q1 over L188, run by run, resonant jets :: sha256 c6c4138ee5c86423 :: prose twin: two significant figures, direction in words
\newcommand{\MlstCmpMlpCoarseOverFullResonant}{0.952\,\ensuremath{\pm}\,0.023}  % experiments/FIGS/data/v2/index.json :: frozen-feature MLP resolution of R16_Q1 over L188, run by run, resonant jets :: sha256 c6c4138ee5c86423
\newcommand{\MlstCmpMlpCoarseOverFullResonantPr}{0.95 times the resolution}  % experiments/FIGS/data/v2/index.json :: frozen-feature MLP resolution of R16_Q1 over L188, run by run, resonant jets :: sha256 c6c4138ee5c86423 :: prose twin: two significant figures, direction in words
\newcommand{\MlstCmpMlpResMassOverResResonant}{0.7137 and 0.7219}  % experiments/FIGS/data/v2/index.json :: frozen-feature MLP resolution of L162_MASS over L162, run by run, resonant jets :: sha256 c6c4138ee5c86423
\newcommand{\MlstCmpMlpResMassOverResResonantPr}{at most 0.72 times the resolution}  % experiments/FIGS/data/v2/index.json :: frozen-feature MLP resolution of L162_MASS over L162, run by run, resonant jets :: sha256 c6c4138ee5c86423 :: prose twin: two significant figures, direction in words
\newcommand{\MlstCmpMlpResMassOverResResonantPrMax}{down to 0.71 times the resolution}  % experiments/FIGS/data/v2/index.json :: frozen-feature MLP resolution of L162_MASS over L162, run by run, resonant jets :: sha256 c6c4138ee5c86423 :: prose twin of the run farther from 1
\newcommand{\MlstCmpMlpCoarseMassOverCoarseResonant}{0.766 and 0.745}  % experiments/FIGS/data/v2/index.json :: frozen-feature MLP resolution of R16_Q1_MASS over R16_Q1, run by run, resonant jets :: sha256 c6c4138ee5c86423
\newcommand{\MlstCmpMlpCoarseMassOverCoarseResonantPr}{at most 0.77 times the resolution}  % experiments/FIGS/data/v2/index.json :: frozen-feature MLP resolution of R16_Q1_MASS over R16_Q1, run by run, resonant jets :: sha256 c6c4138ee5c86423 :: prose twin: two significant figures, direction in words
\newcommand{\MlstCmpMlpCoarseMassOverCoarseResonantPrMax}{down to 0.75 times the resolution}  % experiments/FIGS/data/v2/index.json :: frozen-feature MLP resolution of R16_Q1_MASS over R16_Q1, run by run, resonant jets :: sha256 c6c4138ee5c86423 :: prose twin of the run farther from 1
\newcommand{\MlstMassUncorrectedQcd}{1.52}  % experiments/FIGS/data/v2/index.json :: head/atlas_res: groomed mass, qcd jets, ATLAS resolution :: sha256 c6c4138ee5c86423
\newcommand{\MlstMassUncorrectedResponseQcd}{0.931}  % experiments/FIGS/data/v2/index.json :: head/atlas_res: groomed mass, qcd jets, median response :: sha256 c6c4138ee5c86423
\newcommand{\MlstMassHeadResMassQcd}{0.2263 and 0.2256}  % experiments/FIGS/data/v2/index.json :: regression output of L162_MASS, qcd jets, ATLAS resolution :: sha256 c6c4138ee5c86423
\newcommand{\MlstMassHeadResMassResponseQcd}{0.9897 and 0.9996}  % experiments/FIGS/data/v2/index.json :: regression output of L162_MASS, qcd jets, median response :: sha256 c6c4138ee5c86423
\newcommand{\MlstMassHeadResMassOverUncorrQcd}{0.1485 and 0.1480}  % experiments/FIGS/data/v2/index.json :: resolution of the regression output of L162_MASS over the groomed mass's, qcd jets :: sha256 c6c4138ee5c86423
\newcommand{\MlstMassHeadResMassOverUncorrQcdPr}{at most 0.15 times the resolution}  % experiments/FIGS/data/v2/index.json :: resolution of the regression output of L162_MASS over the groomed mass's, qcd jets :: sha256 c6c4138ee5c86423 :: prose twin: two significant figures, direction in words
\newcommand{\MlstMassHeadResMassOverUncorrQcdPrMax}{down to 0.15 times the resolution}  % experiments/FIGS/data/v2/index.json :: resolution of the regression output of L162_MASS over the groomed mass's, qcd jets :: sha256 c6c4138ee5c86423 :: prose twin of the run farther from 1
\newcommand{\MlstMassHeadResMassOverMlpResQcd}{0.6542 and 0.6632}  % experiments/FIGS/data/v2/index.json :: resolution of the regression output of L162_MASS over the frozen-feature MLP of L162, run by run, qcd jets :: sha256 c6c4138ee5c86423
\newcommand{\MlstMassHeadResMassOverMlpResQcdPr}{at most 0.66 times the resolution}  % experiments/FIGS/data/v2/index.json :: resolution of the regression output of L162_MASS over the frozen-feature MLP of L162, run by run, qcd jets :: sha256 c6c4138ee5c86423 :: prose twin: two significant figures, direction in words
\newcommand{\MlstMassHeadResMassOverMlpResQcdPrMax}{down to 0.65 times the resolution}  % experiments/FIGS/data/v2/index.json :: resolution of the regression output of L162_MASS over the frozen-feature MLP of L162, run by run, qcd jets :: sha256 c6c4138ee5c86423 :: prose twin of the run farther from 1
\newcommand{\MlstMassHeadCoarseMassQcd}{0.2239 and 0.2240}  % experiments/FIGS/data/v2/index.json :: regression output of R16_Q1_MASS, qcd jets, ATLAS resolution :: sha256 c6c4138ee5c86423
\newcommand{\MlstMassHeadCoarseMassResponseQcd}{1.001 and 1.000}  % experiments/FIGS/data/v2/index.json :: regression output of R16_Q1_MASS, qcd jets, median response :: sha256 c6c4138ee5c86423
\newcommand{\MlstMassHeadCoarseMassOverUncorrQcd}{0.1469 and 0.1469}  % experiments/FIGS/data/v2/index.json :: resolution of the regression output of R16_Q1_MASS over the groomed mass's, qcd jets :: sha256 c6c4138ee5c86423
\newcommand{\MlstMassHeadCoarseMassOverUncorrQcdPr}{at most 0.15 times the resolution}  % experiments/FIGS/data/v2/index.json :: resolution of the regression output of R16_Q1_MASS over the groomed mass's, qcd jets :: sha256 c6c4138ee5c86423 :: prose twin: two significant figures, direction in words
\newcommand{\MlstMassHeadCoarseMassOverUncorrQcdPrMax}{down to 0.15 times the resolution}  % experiments/FIGS/data/v2/index.json :: resolution of the regression output of R16_Q1_MASS over the groomed mass's, qcd jets :: sha256 c6c4138ee5c86423 :: prose twin of the run farther from 1
\newcommand{\MlstMassHeadCoarseMassOverMlpCoarseQcd}{0.679 and 0.698}  % experiments/FIGS/data/v2/index.json :: resolution of the regression output of R16_Q1_MASS over the frozen-feature MLP of R16_Q1, run by run, qcd jets :: sha256 c6c4138ee5c86423
\newcommand{\MlstMassHeadCoarseMassOverMlpCoarseQcdPr}{at most 0.70 times the resolution}  % experiments/FIGS/data/v2/index.json :: resolution of the regression output of R16_Q1_MASS over the frozen-feature MLP of R16_Q1, run by run, qcd jets :: sha256 c6c4138ee5c86423 :: prose twin: two significant figures, direction in words
\newcommand{\MlstMassHeadCoarseMassOverMlpCoarseQcdPrMax}{down to 0.68 times the resolution}  % experiments/FIGS/data/v2/index.json :: resolution of the regression output of R16_Q1_MASS over the frozen-feature MLP of R16_Q1, run by run, qcd jets :: sha256 c6c4138ee5c86423 :: prose twin of the run farther from 1
\newcommand{\MlstMassMlpInitQcd}{0.502\,\ensuremath{\pm}\,0.014}  % experiments/FIGS/data/v2/index.json :: frozen-feature MLP of INIT, qcd jets, ATLAS resolution :: sha256 c6c4138ee5c86423
\newcommand{\MlstMassMlpFullQcd}{0.3487\,\ensuremath{\pm}\,0.0029}  % experiments/FIGS/data/v2/index.json :: frozen-feature MLP of L188, qcd jets, ATLAS resolution :: sha256 c6c4138ee5c86423
\newcommand{\MlstMassMlpResQcd}{0.3425\,\ensuremath{\pm}\,0.0031}  % experiments/FIGS/data/v2/index.json :: frozen-feature MLP of L162, qcd jets, ATLAS resolution :: sha256 c6c4138ee5c86423
\newcommand{\MlstMassMlpMidQcd}{0.338\,\ensuremath{\pm}\,0.007}  % experiments/FIGS/data/v2/index.json :: frozen-feature MLP of R42_Q1, qcd jets, ATLAS resolution :: sha256 c6c4138ee5c86423
\newcommand{\MlstMassMlpCoarseQcd}{0.327\,\ensuremath{\pm}\,0.006}  % experiments/FIGS/data/v2/index.json :: frozen-feature MLP of R16_Q1, qcd jets, ATLAS resolution :: sha256 c6c4138ee5c86423
\newcommand{\MlstMassMlpResMassQcd}{0.2473 and 0.2452}  % experiments/FIGS/data/v2/index.json :: frozen-feature MLP of L162_MASS, qcd jets, ATLAS resolution :: sha256 c6c4138ee5c86423
\newcommand{\MlstMassMlpCoarseMassQcd}{0.2465 and 0.2463}  % experiments/FIGS/data/v2/index.json :: frozen-feature MLP of R16_Q1_MASS, qcd jets, ATLAS resolution :: sha256 c6c4138ee5c86423
\newcommand{\MlstCmpMlpResOverFullQcd}{0.982\,\ensuremath{\pm}\,0.017}  % experiments/FIGS/data/v2/index.json :: frozen-feature MLP resolution of L162 over L188, run by run, qcd jets :: sha256 c6c4138ee5c86423
\newcommand{\MlstCmpMlpResOverFullQcdPr}{equal within 2\%}  % experiments/FIGS/data/v2/index.json :: frozen-feature MLP resolution of L162 over L188, run by run, qcd jets :: sha256 c6c4138ee5c86423 :: prose twin: two significant figures, direction in words
\newcommand{\MlstCmpMlpMidOverFullQcd}{0.969\,\ensuremath{\pm}\,0.013}  % experiments/FIGS/data/v2/index.json :: frozen-feature MLP resolution of R42_Q1 over L188, run by run, qcd jets :: sha256 c6c4138ee5c86423
\newcommand{\MlstCmpMlpMidOverFullQcdPr}{0.97 times the resolution}  % experiments/FIGS/data/v2/index.json :: frozen-feature MLP resolution of R42_Q1 over L188, run by run, qcd jets :: sha256 c6c4138ee5c86423 :: prose twin: two significant figures, direction in words
\newcommand{\MlstCmpMlpCoarseOverFullQcd}{0.939\,\ensuremath{\pm}\,0.023}  % experiments/FIGS/data/v2/index.json :: frozen-feature MLP resolution of R16_Q1 over L188, run by run, qcd jets :: sha256 c6c4138ee5c86423
\newcommand{\MlstCmpMlpCoarseOverFullQcdPr}{0.94 times the resolution}  % experiments/FIGS/data/v2/index.json :: frozen-feature MLP resolution of R16_Q1 over L188, run by run, qcd jets :: sha256 c6c4138ee5c86423 :: prose twin: two significant figures, direction in words
\newcommand{\MlstCmpMlpResMassOverResQcd}{0.7147 and 0.7209}  % experiments/FIGS/data/v2/index.json :: frozen-feature MLP resolution of L162_MASS over L162, run by run, qcd jets :: sha256 c6c4138ee5c86423
\newcommand{\MlstCmpMlpResMassOverResQcdPr}{at most 0.72 times the resolution}  % experiments/FIGS/data/v2/index.json :: frozen-feature MLP resolution of L162_MASS over L162, run by run, qcd jets :: sha256 c6c4138ee5c86423 :: prose twin: two significant figures, direction in words
\newcommand{\MlstCmpMlpResMassOverResQcdPrMax}{down to 0.71 times the resolution}  % experiments/FIGS/data/v2/index.json :: frozen-feature MLP resolution of L162_MASS over L162, run by run, qcd jets :: sha256 c6c4138ee5c86423 :: prose twin of the run farther from 1
\newcommand{\MlstCmpMlpCoarseMassOverCoarseQcd}{0.747 and 0.768}  % experiments/FIGS/data/v2/index.json :: frozen-feature MLP resolution of R16_Q1_MASS over R16_Q1, run by run, qcd jets :: sha256 c6c4138ee5c86423
\newcommand{\MlstCmpMlpCoarseMassOverCoarseQcdPr}{at most 0.77 times the resolution}  % experiments/FIGS/data/v2/index.json :: frozen-feature MLP resolution of R16_Q1_MASS over R16_Q1, run by run, qcd jets :: sha256 c6c4138ee5c86423 :: prose twin: two significant figures, direction in words
\newcommand{\MlstCmpMlpCoarseMassOverCoarseQcdPrMax}{down to 0.75 times the resolution}  % experiments/FIGS/data/v2/index.json :: frozen-feature MLP resolution of R16_Q1_MASS over R16_Q1, run by run, qcd jets :: sha256 c6c4138ee5c86423 :: prose twin of the run farther from 1
\newcommand{\MlstMassRespAxesShareQcd}{76\%}  % experiments/FIGS/data/v2/index.json :: head/atlas_res: share of QCD jets whose groomed-mass response lies inside F9a(b)'s x axis :: sha256 c6c4138ee5c86423
\newcommand{\MlstSculptNPassTen}{41{,}028}  % experiments/FIGS/data/v2/index.json :: QCD test jets passing the cut at QCD efficiency 0.1, every run :: sha256 c6c4138ee5c86423
\newcommand{\MlstSculptFloorTen}{0.00944}  % experiments/FIGS/data/v2/index.json :: Jensen-Shannon distance expected for a random selection of as many QCD jets as pass at 0.1 :: sha256 c6c4138ee5c86423
\newcommand{\MlstSculptFullTen}{0.0524\,\ensuremath{\pm}\,0.0034}  % experiments/FIGS/data/v2/index.json :: Jensen-Shannon distance of the QCD m_SD distribution at QCD efficiency 0.1, L188 :: sha256 c6c4138ee5c86423
\newcommand{\MlstSculptFullOverFloorTen}{5.6\,\ensuremath{\pm}\,0.4}  % experiments/FIGS/data/v2/index.json :: distance at 0.1 over the random-selection floor, L188 :: sha256 c6c4138ee5c86423
\newcommand{\MlstSculptFullOverFloorTenPr}{5.6 times the distance}  % experiments/FIGS/data/v2/index.json :: distance at 0.1 over the random-selection floor, L188 :: sha256 c6c4138ee5c86423 :: prose twin: two significant figures, direction in words
\newcommand{\MlstSculptResTen}{0.0532\,\ensuremath{\pm}\,0.0009}  % experiments/FIGS/data/v2/index.json :: Jensen-Shannon distance of the QCD m_SD distribution at QCD efficiency 0.1, L162 :: sha256 c6c4138ee5c86423
\newcommand{\MlstSculptResOverFloorTen}{5.64\,\ensuremath{\pm}\,0.10}  % experiments/FIGS/data/v2/index.json :: distance at 0.1 over the random-selection floor, L162 :: sha256 c6c4138ee5c86423
\newcommand{\MlstSculptResOverFloorTenPr}{5.6 times the distance}  % experiments/FIGS/data/v2/index.json :: distance at 0.1 over the random-selection floor, L162 :: sha256 c6c4138ee5c86423 :: prose twin: two significant figures, direction in words
\newcommand{\MlstSculptMidTen}{0.0560\,\ensuremath{\pm}\,0.0027}  % experiments/FIGS/data/v2/index.json :: Jensen-Shannon distance of the QCD m_SD distribution at QCD efficiency 0.1, R42_Q1 :: sha256 c6c4138ee5c86423
\newcommand{\MlstSculptMidOverFloorTen}{5.94\,\ensuremath{\pm}\,0.28}  % experiments/FIGS/data/v2/index.json :: distance at 0.1 over the random-selection floor, R42_Q1 :: sha256 c6c4138ee5c86423
\newcommand{\MlstSculptMidOverFloorTenPr}{5.9 times the distance}  % experiments/FIGS/data/v2/index.json :: distance at 0.1 over the random-selection floor, R42_Q1 :: sha256 c6c4138ee5c86423 :: prose twin: two significant figures, direction in words
\newcommand{\MlstSculptCoarseTen}{0.0543\,\ensuremath{\pm}\,0.0010}  % experiments/FIGS/data/v2/index.json :: Jensen-Shannon distance of the QCD m_SD distribution at QCD efficiency 0.1, R16_Q1 :: sha256 c6c4138ee5c86423
\newcommand{\MlstSculptCoarseOverFloorTen}{5.76\,\ensuremath{\pm}\,0.10}  % experiments/FIGS/data/v2/index.json :: distance at 0.1 over the random-selection floor, R16_Q1 :: sha256 c6c4138ee5c86423
\newcommand{\MlstSculptCoarseOverFloorTenPr}{5.8 times the distance}  % experiments/FIGS/data/v2/index.json :: distance at 0.1 over the random-selection floor, R16_Q1 :: sha256 c6c4138ee5c86423 :: prose twin: two significant figures, direction in words
\newcommand{\MlstSculptResMassTen}{0.0422 and 0.0464}  % experiments/FIGS/data/v2/index.json :: Jensen-Shannon distance of the QCD m_SD distribution at QCD efficiency 0.1, L162_MASS :: sha256 c6c4138ee5c86423
\newcommand{\MlstSculptResMassOverFloorTen}{4.47 and 4.91}  % experiments/FIGS/data/v2/index.json :: distance at 0.1 over the random-selection floor, L162_MASS :: sha256 c6c4138ee5c86423
\newcommand{\MlstSculptResMassOverFloorTenPr}{at least 4.5 times the distance}  % experiments/FIGS/data/v2/index.json :: distance at 0.1 over the random-selection floor, L162_MASS :: sha256 c6c4138ee5c86423 :: prose twin: two significant figures, direction in words
\newcommand{\MlstSculptResMassOverFloorTenPrMax}{up to 4.9 times the distance}  % experiments/FIGS/data/v2/index.json :: distance at 0.1 over the random-selection floor, L162_MASS :: sha256 c6c4138ee5c86423 :: prose twin of the run farther from 1
\newcommand{\MlstSculptCoarseMassTen}{0.0522 and 0.0423}  % experiments/FIGS/data/v2/index.json :: Jensen-Shannon distance of the QCD m_SD distribution at QCD efficiency 0.1, R16_Q1_MASS :: sha256 c6c4138ee5c86423
\newcommand{\MlstSculptCoarseMassOverFloorTen}{5.5 and 4.5}  % experiments/FIGS/data/v2/index.json :: distance at 0.1 over the random-selection floor, R16_Q1_MASS :: sha256 c6c4138ee5c86423
\newcommand{\MlstSculptCoarseMassOverFloorTenPr}{at least 4.5 times the distance}  % experiments/FIGS/data/v2/index.json :: distance at 0.1 over the random-selection floor, R16_Q1_MASS :: sha256 c6c4138ee5c86423 :: prose twin: two significant figures, direction in words
\newcommand{\MlstSculptCoarseMassOverFloorTenPrMax}{up to 5.5 times the distance}  % experiments/FIGS/data/v2/index.json :: distance at 0.1 over the random-selection floor, R16_Q1_MASS :: sha256 c6c4138ee5c86423 :: prose twin of the run farther from 1
\newcommand{\MlstSculptResMassOverResTen}{0.789 and 0.860}  % experiments/FIGS/data/v2/index.json :: distance at QCD efficiency 0.1: L162_MASS over L162, run by run :: sha256 c6c4138ee5c86423
\newcommand{\MlstSculptResMassOverResTenPr}{at most 0.86 times the distance}  % experiments/FIGS/data/v2/index.json :: distance at QCD efficiency 0.1: L162_MASS over L162, run by run :: sha256 c6c4138ee5c86423 :: prose twin: two significant figures, direction in words
\newcommand{\MlstSculptResMassOverResTenPrMax}{down to 0.79 times the distance}  % experiments/FIGS/data/v2/index.json :: distance at QCD efficiency 0.1: L162_MASS over L162, run by run :: sha256 c6c4138ee5c86423 :: prose twin of the run farther from 1
\newcommand{\MlstSculptCoarseMassOverCoarseTen}{0.94 and 0.79}  % experiments/FIGS/data/v2/index.json :: distance at QCD efficiency 0.1: R16_Q1_MASS over R16_Q1, run by run :: sha256 c6c4138ee5c86423
\newcommand{\MlstSculptCoarseMassOverCoarseTenPr}{equal within 30\%}  % experiments/FIGS/data/v2/index.json :: distance at QCD efficiency 0.1: R16_Q1_MASS over R16_Q1, run by run :: sha256 c6c4138ee5c86423 :: prose twin: two significant figures, direction in words
\newcommand{\MlstSculptCoarseMassOverCoarseTenPrMax}{equal within 30\%}  % experiments/FIGS/data/v2/index.json :: distance at QCD efficiency 0.1: R16_Q1_MASS over R16_Q1, run by run :: sha256 c6c4138ee5c86423 :: prose twin of the run farther from 1
\newcommand{\MlstSculptCoarseOverFullTen}{1.04\,\ensuremath{\pm}\,0.05}  % experiments/FIGS/data/v2/index.json :: distance at QCD efficiency 0.1: R16_Q1 over L188, run by run :: sha256 c6c4138ee5c86423
\newcommand{\MlstSculptCoarseOverFullTenPr}{equal within 6\%}  % experiments/FIGS/data/v2/index.json :: distance at QCD efficiency 0.1: R16_Q1 over L188, run by run :: sha256 c6c4138ee5c86423 :: prose twin: two significant figures, direction in words
\newcommand{\MlstSculptResOverFullTen}{1.02\,\ensuremath{\pm}\,0.06}  % experiments/FIGS/data/v2/index.json :: distance at QCD efficiency 0.1: L162 over L188, run by run :: sha256 c6c4138ee5c86423
\newcommand{\MlstSculptResOverFullTenPr}{equal within 6\%}  % experiments/FIGS/data/v2/index.json :: distance at QCD efficiency 0.1: L162 over L188, run by run :: sha256 c6c4138ee5c86423 :: prose twin: two significant figures, direction in words
\newcommand{\MlstSculptMidOverFullTen}{1.069\,\ensuremath{\pm}\,0.020}  % experiments/FIGS/data/v2/index.json :: distance at QCD efficiency 0.1: R42_Q1 over L188, run by run :: sha256 c6c4138ee5c86423
\newcommand{\MlstSculptMidOverFullTenPr}{1.1 times the distance}  % experiments/FIGS/data/v2/index.json :: distance at QCD efficiency 0.1: R42_Q1 over L188, run by run :: sha256 c6c4138ee5c86423 :: prose twin: two significant figures, direction in words
\newcommand{\MlstSculptCoarseOverMidTen}{0.971\,\ensuremath{\pm}\,0.035}  % experiments/FIGS/data/v2/index.json :: distance at QCD efficiency 0.1: R16_Q1 over R42_Q1, run by run :: sha256 c6c4138ee5c86423
\newcommand{\MlstSculptCoarseOverMidTenPr}{equal within 4\%}  % experiments/FIGS/data/v2/index.json :: distance at QCD efficiency 0.1: R16_Q1 over R42_Q1, run by run :: sha256 c6c4138ee5c86423 :: prose twin: two significant figures, direction in words
\newcommand{\MlstSculptNPassOne}{4{,}103}  % experiments/FIGS/data/v2/index.json :: QCD test jets passing the cut at QCD efficiency 0.01, every run :: sha256 c6c4138ee5c86423
\newcommand{\MlstSculptFloorOne}{0.0313}  % experiments/FIGS/data/v2/index.json :: Jensen-Shannon distance expected for a random selection of as many QCD jets as pass at 0.01 :: sha256 c6c4138ee5c86423
\newcommand{\MlstSculptFullOne}{0.0439\,\ensuremath{\pm}\,0.0024}  % experiments/FIGS/data/v2/index.json :: Jensen-Shannon distance of the QCD m_SD distribution at QCD efficiency 0.01, L188 :: sha256 c6c4138ee5c86423
\newcommand{\MlstSculptFullOverFloorOne}{1.40\,\ensuremath{\pm}\,0.08}  % experiments/FIGS/data/v2/index.json :: distance at 0.01 over the random-selection floor, L188 :: sha256 c6c4138ee5c86423
\newcommand{\MlstSculptFullOverFloorOnePr}{1.4 times the distance}  % experiments/FIGS/data/v2/index.json :: distance at 0.01 over the random-selection floor, L188 :: sha256 c6c4138ee5c86423 :: prose twin: two significant figures, direction in words
\newcommand{\MlstSculptResOne}{0.0429\,\ensuremath{\pm}\,0.0029}  % experiments/FIGS/data/v2/index.json :: Jensen-Shannon distance of the QCD m_SD distribution at QCD efficiency 0.01, L162 :: sha256 c6c4138ee5c86423
\newcommand{\MlstSculptResOverFloorOne}{1.37\,\ensuremath{\pm}\,0.09}  % experiments/FIGS/data/v2/index.json :: distance at 0.01 over the random-selection floor, L162 :: sha256 c6c4138ee5c86423
\newcommand{\MlstSculptResOverFloorOnePr}{1.4 times the distance}  % experiments/FIGS/data/v2/index.json :: distance at 0.01 over the random-selection floor, L162 :: sha256 c6c4138ee5c86423 :: prose twin: two significant figures, direction in words
\newcommand{\MlstSculptMidOne}{0.051\,\ensuremath{\pm}\,0.007}  % experiments/FIGS/data/v2/index.json :: Jensen-Shannon distance of the QCD m_SD distribution at QCD efficiency 0.01, R42_Q1 :: sha256 c6c4138ee5c86423
\newcommand{\MlstSculptMidOverFloorOne}{1.62\,\ensuremath{\pm}\,0.23}  % experiments/FIGS/data/v2/index.json :: distance at 0.01 over the random-selection floor, R42_Q1 :: sha256 c6c4138ee5c86423
\newcommand{\MlstSculptMidOverFloorOnePr}{1.6 times the distance}  % experiments/FIGS/data/v2/index.json :: distance at 0.01 over the random-selection floor, R42_Q1 :: sha256 c6c4138ee5c86423 :: prose twin: two significant figures, direction in words
\newcommand{\MlstSculptCoarseOne}{0.0561\,\ensuremath{\pm}\,0.0017}  % experiments/FIGS/data/v2/index.json :: Jensen-Shannon distance of the QCD m_SD distribution at QCD efficiency 0.01, R16_Q1 :: sha256 c6c4138ee5c86423
\newcommand{\MlstSculptCoarseOverFloorOne}{1.79\,\ensuremath{\pm}\,0.06}  % experiments/FIGS/data/v2/index.json :: distance at 0.01 over the random-selection floor, R16_Q1 :: sha256 c6c4138ee5c86423
\newcommand{\MlstSculptCoarseOverFloorOnePr}{1.8 times the distance}  % experiments/FIGS/data/v2/index.json :: distance at 0.01 over the random-selection floor, R16_Q1 :: sha256 c6c4138ee5c86423 :: prose twin: two significant figures, direction in words
\newcommand{\MlstSculptResMassOne}{0.0424 and 0.0442}  % experiments/FIGS/data/v2/index.json :: Jensen-Shannon distance of the QCD m_SD distribution at QCD efficiency 0.01, L162_MASS :: sha256 c6c4138ee5c86423
\newcommand{\MlstSculptResMassOverFloorOne}{1.356 and 1.411}  % experiments/FIGS/data/v2/index.json :: distance at 0.01 over the random-selection floor, L162_MASS :: sha256 c6c4138ee5c86423
\newcommand{\MlstSculptResMassOverFloorOnePr}{at least 1.4 times the distance}  % experiments/FIGS/data/v2/index.json :: distance at 0.01 over the random-selection floor, L162_MASS :: sha256 c6c4138ee5c86423 :: prose twin: two significant figures, direction in words
\newcommand{\MlstSculptResMassOverFloorOnePrMax}{up to 1.4 times the distance}  % experiments/FIGS/data/v2/index.json :: distance at 0.01 over the random-selection floor, L162_MASS :: sha256 c6c4138ee5c86423 :: prose twin of the run farther from 1
\newcommand{\MlstSculptCoarseMassOne}{0.0405 and 0.0452}  % experiments/FIGS/data/v2/index.json :: Jensen-Shannon distance of the QCD m_SD distribution at QCD efficiency 0.01, R16_Q1_MASS :: sha256 c6c4138ee5c86423
\newcommand{\MlstSculptCoarseMassOverFloorOne}{1.29 and 1.44}  % experiments/FIGS/data/v2/index.json :: distance at 0.01 over the random-selection floor, R16_Q1_MASS :: sha256 c6c4138ee5c86423
\newcommand{\MlstSculptCoarseMassOverFloorOnePr}{at least 1.3 times the distance}  % experiments/FIGS/data/v2/index.json :: distance at 0.01 over the random-selection floor, R16_Q1_MASS :: sha256 c6c4138ee5c86423 :: prose twin: two significant figures, direction in words
\newcommand{\MlstSculptCoarseMassOverFloorOnePrMax}{up to 1.4 times the distance}  % experiments/FIGS/data/v2/index.json :: distance at 0.01 over the random-selection floor, R16_Q1_MASS :: sha256 c6c4138ee5c86423 :: prose twin of the run farther from 1
\newcommand{\MlstSculptResMassOverResOne}{0.962 and 0.983}  % experiments/FIGS/data/v2/index.json :: distance at QCD efficiency 0.01: L162_MASS over L162, run by run :: sha256 c6c4138ee5c86423
\newcommand{\MlstSculptResMassOverResOnePr}{equal within 4\%}  % experiments/FIGS/data/v2/index.json :: distance at QCD efficiency 0.01: L162_MASS over L162, run by run :: sha256 c6c4138ee5c86423 :: prose twin: two significant figures, direction in words
\newcommand{\MlstSculptResMassOverResOnePrMax}{equal within 4\%}  % experiments/FIGS/data/v2/index.json :: distance at QCD efficiency 0.01: L162_MASS over L162, run by run :: sha256 c6c4138ee5c86423 :: prose twin of the run farther from 1
\newcommand{\MlstSculptCoarseMassOverCoarseOne}{0.70 and 0.82}  % experiments/FIGS/data/v2/index.json :: distance at QCD efficiency 0.01: R16_Q1_MASS over R16_Q1, run by run :: sha256 c6c4138ee5c86423
\newcommand{\MlstSculptCoarseMassOverCoarseOnePr}{at most 0.82 times the distance}  % experiments/FIGS/data/v2/index.json :: distance at QCD efficiency 0.01: R16_Q1_MASS over R16_Q1, run by run :: sha256 c6c4138ee5c86423 :: prose twin: two significant figures, direction in words
\newcommand{\MlstSculptCoarseMassOverCoarseOnePrMax}{down to 0.70 times the distance}  % experiments/FIGS/data/v2/index.json :: distance at QCD efficiency 0.01: R16_Q1_MASS over R16_Q1, run by run :: sha256 c6c4138ee5c86423 :: prose twin of the run farther from 1
\newcommand{\MlstSculptCoarseOverFullOne}{1.28\,\ensuremath{\pm}\,0.04}  % experiments/FIGS/data/v2/index.json :: distance at QCD efficiency 0.01: R16_Q1 over L188, run by run :: sha256 c6c4138ee5c86423
\newcommand{\MlstSculptCoarseOverFullOnePr}{1.3 times the distance}  % experiments/FIGS/data/v2/index.json :: distance at QCD efficiency 0.01: R16_Q1 over L188, run by run :: sha256 c6c4138ee5c86423 :: prose twin: two significant figures, direction in words
\newcommand{\MlstSculptResOverFullOne}{0.98\,\ensuremath{\pm}\,0.05}  % experiments/FIGS/data/v2/index.json :: distance at QCD efficiency 0.01: L162 over L188, run by run :: sha256 c6c4138ee5c86423
\newcommand{\MlstSculptResOverFullOnePr}{equal within 5\%}  % experiments/FIGS/data/v2/index.json :: distance at QCD efficiency 0.01: L162 over L188, run by run :: sha256 c6c4138ee5c86423 :: prose twin: two significant figures, direction in words
\newcommand{\MlstSculptMidOverFullOne}{1.15\,\ensuremath{\pm}\,0.12}  % experiments/FIGS/data/v2/index.json :: distance at QCD efficiency 0.01: R42_Q1 over L188, run by run :: sha256 c6c4138ee5c86423
\newcommand{\MlstSculptMidOverFullOnePr}{equal within 20\%}  % experiments/FIGS/data/v2/index.json :: distance at QCD efficiency 0.01: R42_Q1 over L188, run by run :: sha256 c6c4138ee5c86423 :: prose twin: two significant figures, direction in words
\newcommand{\MlstSculptCoarseOverMidOne}{1.13\,\ensuremath{\pm}\,0.16}  % experiments/FIGS/data/v2/index.json :: distance at QCD efficiency 0.01: R16_Q1 over R42_Q1, run by run :: sha256 c6c4138ee5c86423
\newcommand{\MlstSculptCoarseOverMidOnePr}{equal within 20\%}  % experiments/FIGS/data/v2/index.json :: distance at QCD efficiency 0.01: R16_Q1 over R42_Q1, run by run :: sha256 c6c4138ee5c86423 :: prose twin: two significant figures, direction in words
\newcommand{\MlstSculptDiffEffLo}{0.52\%}  % experiments/FIGS/data/v2/index.json :: lowest stored QCD efficiency of the unbroken stretch around the tighter working point at which the 17/188 distance ratio exceeds 1 by twice its SD :: sha256 c6c4138ee5c86423 :: passing grid points: 0.001, 0.005245, 0.007937, 0.012011
\newcommand{\MlstSculptDiffEffHi}{1.2\%}  % experiments/FIGS/data/v2/index.json :: highest efficiency of that stretch :: sha256 c6c4138ee5c86423
\newcommand{\MlstSculptLowestEff}{0.10\%}  % experiments/FIGS/data/v2/index.json :: lowest stored QCD efficiency of the sculpting curves :: sha256 c6c4138ee5c86423
\newcommand{\MlstSculptLowestNPass}{410}  % experiments/FIGS/data/v2/index.json :: QCD test jets passing at the lowest stored efficiency :: sha256 c6c4138ee5c86423
\newcommand{\MlstSculptLowestMaxFloorSd}{1.7}  % experiments/FIGS/data/v2/index.json :: largest distance of a configuration's mean from the random-selection floor at the lowest stored efficiency, in SDs of the floor over random selections :: sha256 c6c4138ee5c86423
\newcommand{\MlstSculptNConfigs}{6}  % experiments/FIGS/data/v2/index.json :: configurations with a sculpting curve (count) :: sha256 c6c4138ee5c86423
\newcommand{\LpJciiTestQcdPercent}{20.5\%}  % experiments/FIGS/data/v2/jc2_test/counts.json :: QCD share of the jets, per cent :: sha256 3574ba1f6f2a59f0
\newcommand{\LpJciiTestTwoProngPercent}{38.7\%}  % experiments/FIGS/data/v2/jc2_test/counts.json :: two-prong share of the jets, per cent :: sha256 3574ba1f6f2a59f0
\newcommand{\LpJciiTrainQcdPercent}{13.5\%}  % experiments/FT/data/jc2_v2_manifest.json :: QCD share of the jets, per cent :: sha256 4bd68655919b14e9
\newcommand{\LpJciiTrainTwoProngPercent}{15.4\%}  % experiments/FT/data/jc2_v2_manifest.json :: two-prong share of the jets, per cent :: sha256 4bd68655919b14e9
\newcommand{\LpJciiTestMinAucJets}{10}  % experiments/FIGS/data/v2/jc2_test/counts.json :: fewest jets of a class among the macro-AUC rows :: sha256 3574ba1f6f2a59f0
\newcommand{\LpJciiTrainMinJets}{11}  % experiments/FT/data/jc2_v2_manifest.json :: fewest jets of a class in the 10^6 training subset :: sha256 4bd68655919b14e9
\newcommand{\LpNTestJci}{2{,}000{,}000}  % experiments/FIGS/data/v2/index.json :: jc1 test jets the probes score :: sha256 c6c4138ee5c86423
\newcommand{\LpNTrainMaxJci}{1{,}000{,}000}  % experiments/FIGS/data/v2/index.json :: jc1 largest training subset :: sha256 c6c4138ee5c86423
\newcommand{\LpNTestTop}{404{,}000}  % experiments/FIGS/data/v2/index.json :: top test jets the probes score :: sha256 c6c4138ee5c86423
\newcommand{\LpNTrainMaxTop}{1{,}200{,}000}  % experiments/FIGS/data/v2/index.json :: top largest training subset :: sha256 c6c4138ee5c86423
\newcommand{\LpNTestQg}{200{,}000}  % experiments/FIGS/data/v2/index.json :: qg test jets the probes score :: sha256 c6c4138ee5c86423
\newcommand{\LpNTrainMaxQg}{1{,}600{,}000}  % experiments/FIGS/data/v2/index.json :: qg largest training subset :: sha256 c6c4138ee5c86423
\newcommand{\LpNTestAucJci}{2{,}000{,}000}  % experiments/FIGS/data/v2/index.json :: JetClass test jets of the macro AUC :: sha256 c6c4138ee5c86423
\newcommand{\LpPairNValMin}{1{,}068}  % experiments/FIGS/data/v2/index.json :: validation jets of the pair probe with the fewest :: sha256 c6c4138ee5c86423
\newcommand{\LpPairNValMax}{20{,}195}  % experiments/FIGS/data/v2/index.json :: validation jets of the pair probe with the most :: sha256 c6c4138ee5c86423
\newcommand{\LpNValJci}{200{,}000}  % experiments/FIGS/data/v2/index.json :: JetClass probe validation jets :: sha256 c6c4138ee5c86423
\newcommand{\SculptEffHigh}{10\%}  % experiments/EVAL/mass_sculpting.py :: WORKING_POINTS :: sha256 e85d125c85dac054 :: QCD efficiency of a sculpting cut
\newcommand{\SculptEffLow}{1\%}  % experiments/EVAL/mass_sculpting.py :: WORKING_POINTS :: sha256 e85d125c85dac054 :: QCD efficiency of a sculpting cut
\newcommand{\LpMlpHidden}{256}  % experiments/EVAL/probe.py :: fit_mlp hidden layer :: sha256 485db0966e660588 :: units of the MLP probe's hidden layer
\newcommand{\LpMlpSeedsWord}{three}  % experiments/EVAL/probe.py :: MLP_SEEDS :: sha256 485db0966e660588 :: seeds of the MLP probe whose scores are averaged
\newcommand{\MassMinTrainPerClass}{50}  % experiments/EVAL/mass_resolution.py :: MIN_TRAIN_PER_CLASS :: sha256 26acb0e315094bd1 :: fewest training jets of a class the mass regression keeps
\newcommand{\MlstFtLr}{\ensuremath{10^{-4}}}  % experiments/FIGS/data/v2/ft_mlst/manifests.json :: cells[best70].lr :: sha256 a689f21e894f48ae
\newcommand{\MlstFtHeadMult}{50}  % experiments/FIGS/data/v2/ft_mlst/manifests.json :: cells[best70].head_lr_mult :: sha256 a689f21e894f48ae
\newcommand{\MlstFtWeightDecay}{0.01}  % experiments/FIGS/data/v2/ft_mlst/manifests.json :: cells[best70].weight_decay :: sha256 a689f21e894f48ae
\newcommand{\MlstFtBatch}{512}  % experiments/FIGS/data/v2/ft_mlst/manifests.json :: cells[best70].batch_size :: sha256 a689f21e894f48ae
\newcommand{\MlstFtBenchEpochs}{20}  % experiments/FIGS/data/v2/ft_mlst/manifests.json :: cells[best70, top/qg].epochs :: sha256 a689f21e894f48ae
\newcommand{\MlstFtEpochsEThree}{50}  % experiments/FIGS/data/v2/ft_mlst/manifests.json :: cells[best70, leg 1/2, n_train 1000].epochs :: sha256 a689f21e894f48ae
\newcommand{\MlstFtEpochsEFour}{50}  % experiments/FIGS/data/v2/ft_mlst/manifests.json :: cells[best70, leg 1/2, n_train 10000].epochs :: sha256 a689f21e894f48ae
\newcommand{\MlstFtEpochsEFive}{30}  % experiments/FIGS/data/v2/ft_mlst/manifests.json :: cells[best70, leg 1/2, n_train 100000].epochs :: sha256 a689f21e894f48ae
\newcommand{\MlstFtEpochsESix}{10}  % experiments/FIGS/data/v2/ft_mlst/manifests.json :: cells[best70, leg 1/2, n_train 1000000].epochs :: sha256 a689f21e894f48ae
\newcommand{\MlstFtHeadDecayShare}{30\%}  % experiments/FT/k8s/job-ft-v2-legs-best70-l188-s3-raunav.yaml :: no --lr-scheduler: weaver 0.4.17 flat+decay, the head's last epochs :: sha256 19f59aa54a385c11
\newcommand{\MlstFtHeadDecayTo}{1\%}  % experiments/FT/k8s/job-ft-v2-legs-best70-l188-s3-raunav.yaml :: no --lr-scheduler: weaver 0.4.17 flat+decay, final factor :: sha256 19f59aa54a385c11
\newcommand{\MlstFtHeadDropout}{0.1}  % experiments/FT/k8s/job-ft-v2-legs-best70-l188-s3-raunav.yaml :: fc_params dropout of the fine-tuning head :: sha256 19f59aa54a385c11
\newcommand{\MlstFtMinEpochJets}{10{,}000}  % experiments/FIGS/data/v2/ft_mlst/manifests.json :: cells[best70].samples_per_epoch (smallest) :: sha256 a689f21e894f48ae
\newcommand{\MlstFtValJcii}{20{,}480}  % experiments/FIGS/data/v2/ft_mlst/manifests.json :: cells[best70, leg 1].val_jets :: sha256 a689f21e894f48ae
\newcommand{\MlstFtValSmall}{20{,}000}  % experiments/FT/k8s/job-ft-v2-legs-best70-l188-s3-raunav.yaml :: JetClass --samples-per-epoch-val; top and q/g at N <= 1e4 :: sha256 19f59aa54a385c11
\newcommand{\MlstFtValLarge}{200{,}000}  % experiments/FT/k8s/job-ft-v2-bench-best70-l188-s1-top-raunav.yaml :: top and q/g --samples-per-epoch-val at N > 1e4 :: sha256 1e38d4f542119ea3
\newcommand{\MlstFtGpuTypes}{V100}  % experiments/FIGS/data/v2/ft_mlst/manifests.json :: cells[best70].gpu_device_name :: sha256 a689f21e894f48ae
\newcommand{\MlstFtCellsDone}{445}  % experiments/FIGS/data/v2/ft_mlst/manifests.json :: DONE paper-run fine-tuning cells (count) :: sha256 a689f21e894f48ae
\newcommand{\MlstFtRefLr}{\ensuremath{10^{-3}}}  % experiments/FIGS/data/v2/ft_mlst/manifests.json :: cells[ref_v100].lr :: sha256 a689f21e894f48ae
\newcommand{\MlstFtRefSeedsWord}{three}  % experiments/FIGS/data/v2/ft_mlst/manifests.json :: cells[ref_v100].ft_seed (count) :: sha256 a689f21e894f48ae
\newcommand{\MlstFtRefGpuTypes}{V100}  % experiments/FIGS/data/v2/ft_mlst/manifests.json :: cells[ref_v100].gpu_device_name :: sha256 a689f21e894f48ae
\newcommand{\MlstFtRefWeightDecayBench}{0.01}  % experiments/FIGS/data/v2/ft_mlst/manifests.json :: cells[ref_v100, top/qg].weight_decay :: sha256 a689f21e894f48ae
\newcommand{\MlstBootResamples}{200}  % experiments/FIGS/data/v2/index.json :: bootstrap resamples of the test jets :: sha256 c6c4138ee5c86423
\newcommand{\MlstBootTestMaxProbe}{0.21}  % experiments/FIGS/data/v2/index.json :: largest test-sample SD of the ln ratio of 1 - AUC on JetClass-II, linear probe, label sets over the 188-class model, per cent :: sha256 c6c4138ee5c86423
\newcommand{\MlstBootTestMaxFt}{0.15}  % experiments/FIGS/data/v2/index.json :: largest test-sample SD of the ln ratio of 1 - AUC on JetClass-II, fine-tuning, label sets over the 188-class model, per cent :: sha256 c6c4138ee5c86423
\newcommand{\MlstBootTestMaxRthirtyProbe}{3.6}  % experiments/FIGS/data/v2/index.json :: largest test-sample SD of the ln R30 ratio on top and quark/gluon, linear probe (counting), label sets over the 188-class model, per cent :: sha256 c6c4138ee5c86423
\newcommand{\MlstBootTestMaxRthirtyFt}{13}  % experiments/FIGS/data/v2/index.json :: largest test-sample SD of the ln R30 ratio on top and quark/gluon, fine-tuning (counting), label sets over the 188-class model, per cent :: sha256 c6c4138ee5c86423
\newcommand{\MlstFtCmpJcTwoResOverFullEFour}{0.989\,\ensuremath{\pm}\,0.011}  % experiments/FIGS/data/v2/index.json :: fine-tuning jc2 N=10000: (1 - macro_auc_ovr) of L162 over L188, run by run :: sha256 c6c4138ee5c86423
\newcommand{\MlstFtCmpJcTwoResOverFullEFourPr}{equal within 2\%}  % experiments/FIGS/data/v2/index.json :: fine-tuning jc2 N=10000: (1 - macro_auc_ovr) of L162 over L188, run by run :: sha256 c6c4138ee5c86423 :: prose twin: two significant figures, direction in words
\newcommand{\MlstFtCmpJcTwoMidOverFullEFour}{1.191\,\ensuremath{\pm}\,0.030}  % experiments/FIGS/data/v2/index.json :: fine-tuning jc2 N=10000: (1 - macro_auc_ovr) of R42_Q1 over L188, run by run :: sha256 c6c4138ee5c86423
\newcommand{\MlstFtCmpJcTwoMidOverFullEFourPr}{19\% more errors}  % experiments/FIGS/data/v2/index.json :: fine-tuning jc2 N=10000: (1 - macro_auc_ovr) of R42_Q1 over L188, run by run :: sha256 c6c4138ee5c86423 :: prose twin: two significant figures, direction in words
\newcommand{\MlstFtCmpJcTwoCoarseOverFullEFour}{1.353\,\ensuremath{\pm}\,0.003}  % experiments/FIGS/data/v2/index.json :: fine-tuning jc2 N=10000: (1 - macro_auc_ovr) of R16_Q1 over L188, run by run :: sha256 c6c4138ee5c86423
\newcommand{\MlstFtCmpJcTwoCoarseOverFullEFourPr}{35\% more errors}  % experiments/FIGS/data/v2/index.json :: fine-tuning jc2 N=10000: (1 - macro_auc_ovr) of R16_Q1 over L188, run by run :: sha256 c6c4138ee5c86423 :: prose twin: two significant figures, direction in words
\newcommand{\MlstFtCmpJcTwoResMassOverResEFour}{1.016 and 1.005}  % experiments/FIGS/data/v2/index.json :: fine-tuning jc2 N=10000: (1 - macro_auc_ovr) of L162_MASS over L162, run by run :: sha256 c6c4138ee5c86423
\newcommand{\MlstFtCmpJcTwoResMassOverResEFourPr}{equal within 2\%}  % experiments/FIGS/data/v2/index.json :: fine-tuning jc2 N=10000: (1 - macro_auc_ovr) of L162_MASS over L162, run by run :: sha256 c6c4138ee5c86423 :: prose twin: two significant figures, direction in words
\newcommand{\MlstFtCmpJcTwoResMassOverResEFourPrMax}{equal within 2\%}  % experiments/FIGS/data/v2/index.json :: fine-tuning jc2 N=10000: (1 - macro_auc_ovr) of L162_MASS over L162, run by run :: sha256 c6c4138ee5c86423 :: prose twin of the run farther from 1
\newcommand{\MlstFtCmpJcTwoCoarseMassOverCoarseEFour}{0.989 and 1.000}  % experiments/FIGS/data/v2/index.json :: fine-tuning jc2 N=10000: (1 - macro_auc_ovr) of R16_Q1_MASS over R16_Q1, run by run :: sha256 c6c4138ee5c86423
\newcommand{\MlstFtCmpJcTwoCoarseMassOverCoarseEFourPr}{equal within 2\%}  % experiments/FIGS/data/v2/index.json :: fine-tuning jc2 N=10000: (1 - macro_auc_ovr) of R16_Q1_MASS over R16_Q1, run by run :: sha256 c6c4138ee5c86423 :: prose twin: two significant figures, direction in words
\newcommand{\MlstFtCmpJcTwoCoarseMassOverCoarseEFourPrMax}{equal within 2\%}  % experiments/FIGS/data/v2/index.json :: fine-tuning jc2 N=10000: (1 - macro_auc_ovr) of R16_Q1_MASS over R16_Q1, run by run :: sha256 c6c4138ee5c86423 :: prose twin of the run farther from 1
\newcommand{\MlstFtCmpJcTwoResOverFullEFive}{1.011\,\ensuremath{\pm}\,0.005}  % experiments/FIGS/data/v2/index.json :: fine-tuning jc2 N=100000: (1 - macro_auc_ovr) of L162 over L188, run by run :: sha256 c6c4138ee5c86423
\newcommand{\MlstFtCmpJcTwoResOverFullEFivePr}{1.1\% more errors}  % experiments/FIGS/data/v2/index.json :: fine-tuning jc2 N=100000: (1 - macro_auc_ovr) of L162 over L188, run by run :: sha256 c6c4138ee5c86423 :: prose twin: two significant figures, direction in words
\newcommand{\MlstFtCmpJcTwoMidOverFullEFive}{1.064\,\ensuremath{\pm}\,0.009}  % experiments/FIGS/data/v2/index.json :: fine-tuning jc2 N=100000: (1 - macro_auc_ovr) of R42_Q1 over L188, run by run :: sha256 c6c4138ee5c86423
\newcommand{\MlstFtCmpJcTwoMidOverFullEFivePr}{6.4\% more errors}  % experiments/FIGS/data/v2/index.json :: fine-tuning jc2 N=100000: (1 - macro_auc_ovr) of R42_Q1 over L188, run by run :: sha256 c6c4138ee5c86423 :: prose twin: two significant figures, direction in words
\newcommand{\MlstFtCmpJcTwoCoarseOverFullEFive}{1.157\,\ensuremath{\pm}\,0.001}  % experiments/FIGS/data/v2/index.json :: fine-tuning jc2 N=100000: (1 - macro_auc_ovr) of R16_Q1 over L188, run by run :: sha256 c6c4138ee5c86423
\newcommand{\MlstFtCmpJcTwoCoarseOverFullEFivePr}{16\% more errors}  % experiments/FIGS/data/v2/index.json :: fine-tuning jc2 N=100000: (1 - macro_auc_ovr) of R16_Q1 over L188, run by run :: sha256 c6c4138ee5c86423 :: prose twin: two significant figures, direction in words
\newcommand{\MlstFtCmpJcTwoResMassOverResEFive}{1.027 and 1.021}  % experiments/FIGS/data/v2/index.json :: fine-tuning jc2 N=100000: (1 - macro_auc_ovr) of L162_MASS over L162, run by run :: sha256 c6c4138ee5c86423
\newcommand{\MlstFtCmpJcTwoResMassOverResEFivePr}{at least 2.1\% more errors}  % experiments/FIGS/data/v2/index.json :: fine-tuning jc2 N=100000: (1 - macro_auc_ovr) of L162_MASS over L162, run by run :: sha256 c6c4138ee5c86423 :: prose twin: two significant figures, direction in words
\newcommand{\MlstFtCmpJcTwoResMassOverResEFivePrMax}{up to 2.7\% more errors}  % experiments/FIGS/data/v2/index.json :: fine-tuning jc2 N=100000: (1 - macro_auc_ovr) of L162_MASS over L162, run by run :: sha256 c6c4138ee5c86423 :: prose twin of the run farther from 1
\newcommand{\MlstFtCmpJcTwoCoarseMassOverCoarseEFive}{1.029 and 1.022}  % experiments/FIGS/data/v2/index.json :: fine-tuning jc2 N=100000: (1 - macro_auc_ovr) of R16_Q1_MASS over R16_Q1, run by run :: sha256 c6c4138ee5c86423
\newcommand{\MlstFtCmpJcTwoCoarseMassOverCoarseEFivePr}{at least 2.2\% more errors}  % experiments/FIGS/data/v2/index.json :: fine-tuning jc2 N=100000: (1 - macro_auc_ovr) of R16_Q1_MASS over R16_Q1, run by run :: sha256 c6c4138ee5c86423 :: prose twin: two significant figures, direction in words
\newcommand{\MlstFtCmpJcTwoCoarseMassOverCoarseEFivePrMax}{up to 2.9\% more errors}  % experiments/FIGS/data/v2/index.json :: fine-tuning jc2 N=100000: (1 - macro_auc_ovr) of R16_Q1_MASS over R16_Q1, run by run :: sha256 c6c4138ee5c86423 :: prose twin of the run farther from 1
\newcommand{\MlstFtCmpJcTwoResOverFullESix}{1.014\,\ensuremath{\pm}\,0.011}  % experiments/FIGS/data/v2/index.json :: fine-tuning jc2 N=1000000: (1 - macro_auc_ovr) of L162 over L188, run by run :: sha256 c6c4138ee5c86423
\newcommand{\MlstFtCmpJcTwoResOverFullESixPr}{equal within 2\%}  % experiments/FIGS/data/v2/index.json :: fine-tuning jc2 N=1000000: (1 - macro_auc_ovr) of L162 over L188, run by run :: sha256 c6c4138ee5c86423 :: prose twin: two significant figures, direction in words
\newcommand{\MlstFtCmpJcTwoMidOverFullESix}{1.043\,\ensuremath{\pm}\,0.002}  % experiments/FIGS/data/v2/index.json :: fine-tuning jc2 N=1000000: (1 - macro_auc_ovr) of R42_Q1 over L188, run by run :: sha256 c6c4138ee5c86423
\newcommand{\MlstFtCmpJcTwoMidOverFullESixPr}{4.3\% more errors}  % experiments/FIGS/data/v2/index.json :: fine-tuning jc2 N=1000000: (1 - macro_auc_ovr) of R42_Q1 over L188, run by run :: sha256 c6c4138ee5c86423 :: prose twin: two significant figures, direction in words
\newcommand{\MlstFtCmpJcTwoCoarseOverFullESix}{1.123\,\ensuremath{\pm}\,0.003}  % experiments/FIGS/data/v2/index.json :: fine-tuning jc2 N=1000000: (1 - macro_auc_ovr) of R16_Q1 over L188, run by run :: sha256 c6c4138ee5c86423
\newcommand{\MlstFtCmpJcTwoCoarseOverFullESixPr}{12\% more errors}  % experiments/FIGS/data/v2/index.json :: fine-tuning jc2 N=1000000: (1 - macro_auc_ovr) of R16_Q1 over L188, run by run :: sha256 c6c4138ee5c86423 :: prose twin: two significant figures, direction in words
\newcommand{\MlstFtCmpJcTwoResMassOverResESix}{1.009 and 1.016}  % experiments/FIGS/data/v2/index.json :: fine-tuning jc2 N=1000000: (1 - macro_auc_ovr) of L162_MASS over L162, run by run :: sha256 c6c4138ee5c86423
\newcommand{\MlstFtCmpJcTwoResMassOverResESixPr}{equal within 2\%}  % experiments/FIGS/data/v2/index.json :: fine-tuning jc2 N=1000000: (1 - macro_auc_ovr) of L162_MASS over L162, run by run :: sha256 c6c4138ee5c86423 :: prose twin: two significant figures, direction in words
\newcommand{\MlstFtCmpJcTwoResMassOverResESixPrMax}{equal within 2\%}  % experiments/FIGS/data/v2/index.json :: fine-tuning jc2 N=1000000: (1 - macro_auc_ovr) of L162_MASS over L162, run by run :: sha256 c6c4138ee5c86423 :: prose twin of the run farther from 1
\newcommand{\MlstFtCmpJcTwoCoarseMassOverCoarseESix}{1.023 and 1.017}  % experiments/FIGS/data/v2/index.json :: fine-tuning jc2 N=1000000: (1 - macro_auc_ovr) of R16_Q1_MASS over R16_Q1, run by run :: sha256 c6c4138ee5c86423
\newcommand{\MlstFtCmpJcTwoCoarseMassOverCoarseESixPr}{at least 1.7\% more errors}  % experiments/FIGS/data/v2/index.json :: fine-tuning jc2 N=1000000: (1 - macro_auc_ovr) of R16_Q1_MASS over R16_Q1, run by run :: sha256 c6c4138ee5c86423 :: prose twin: two significant figures, direction in words
\newcommand{\MlstFtCmpJcTwoCoarseMassOverCoarseESixPrMax}{up to 2.3\% more errors}  % experiments/FIGS/data/v2/index.json :: fine-tuning jc2 N=1000000: (1 - macro_auc_ovr) of R16_Q1_MASS over R16_Q1, run by run :: sha256 c6c4138ee5c86423 :: prose twin of the run farther from 1
\newcommand{\MlstFtCmpJcTwoQcdResOverFullEFour}{1.038\,\ensuremath{\pm}\,0.005}  % experiments/FIGS/data/v2/index.json :: fine-tuning jc2_qcd N=10000: (1 - macro_auc_qcd) of L162 over L188, run by run :: sha256 c6c4138ee5c86423
\newcommand{\MlstFtCmpJcTwoQcdResOverFullEFourPr}{3.8\% more errors}  % experiments/FIGS/data/v2/index.json :: fine-tuning jc2_qcd N=10000: (1 - macro_auc_qcd) of L162 over L188, run by run :: sha256 c6c4138ee5c86423 :: prose twin: two significant figures, direction in words
\newcommand{\MlstFtCmpJcTwoQcdMidOverFullEFour}{1.164\,\ensuremath{\pm}\,0.018}  % experiments/FIGS/data/v2/index.json :: fine-tuning jc2_qcd N=10000: (1 - macro_auc_qcd) of R42_Q1 over L188, run by run :: sha256 c6c4138ee5c86423
\newcommand{\MlstFtCmpJcTwoQcdMidOverFullEFourPr}{16\% more errors}  % experiments/FIGS/data/v2/index.json :: fine-tuning jc2_qcd N=10000: (1 - macro_auc_qcd) of R42_Q1 over L188, run by run :: sha256 c6c4138ee5c86423 :: prose twin: two significant figures, direction in words
\newcommand{\MlstFtCmpJcTwoQcdCoarseOverFullEFour}{1.258\,\ensuremath{\pm}\,0.025}  % experiments/FIGS/data/v2/index.json :: fine-tuning jc2_qcd N=10000: (1 - macro_auc_qcd) of R16_Q1 over L188, run by run :: sha256 c6c4138ee5c86423
\newcommand{\MlstFtCmpJcTwoQcdCoarseOverFullEFourPr}{26\% more errors}  % experiments/FIGS/data/v2/index.json :: fine-tuning jc2_qcd N=10000: (1 - macro_auc_qcd) of R16_Q1 over L188, run by run :: sha256 c6c4138ee5c86423 :: prose twin: two significant figures, direction in words
\newcommand{\MlstFtCmpJcTwoQcdResMassOverResEFour}{1.027 and 0.997}  % experiments/FIGS/data/v2/index.json :: fine-tuning jc2_qcd N=10000: (1 - macro_auc_qcd) of L162_MASS over L162, run by run :: sha256 c6c4138ee5c86423
\newcommand{\MlstFtCmpJcTwoQcdResMassOverResEFourPr}{equal within 3\%}  % experiments/FIGS/data/v2/index.json :: fine-tuning jc2_qcd N=10000: (1 - macro_auc_qcd) of L162_MASS over L162, run by run :: sha256 c6c4138ee5c86423 :: prose twin: two significant figures, direction in words
\newcommand{\MlstFtCmpJcTwoQcdResMassOverResEFourPrMax}{equal within 3\%}  % experiments/FIGS/data/v2/index.json :: fine-tuning jc2_qcd N=10000: (1 - macro_auc_qcd) of L162_MASS over L162, run by run :: sha256 c6c4138ee5c86423 :: prose twin of the run farther from 1
\newcommand{\MlstFtCmpJcTwoQcdCoarseMassOverCoarseEFour}{1.028 and 1.048}  % experiments/FIGS/data/v2/index.json :: fine-tuning jc2_qcd N=10000: (1 - macro_auc_qcd) of R16_Q1_MASS over R16_Q1, run by run :: sha256 c6c4138ee5c86423
\newcommand{\MlstFtCmpJcTwoQcdCoarseMassOverCoarseEFourPr}{equal within 5\%}  % experiments/FIGS/data/v2/index.json :: fine-tuning jc2_qcd N=10000: (1 - macro_auc_qcd) of R16_Q1_MASS over R16_Q1, run by run :: sha256 c6c4138ee5c86423 :: prose twin: two significant figures, direction in words
\newcommand{\MlstFtCmpJcTwoQcdCoarseMassOverCoarseEFourPrMax}{equal within 5\%}  % experiments/FIGS/data/v2/index.json :: fine-tuning jc2_qcd N=10000: (1 - macro_auc_qcd) of R16_Q1_MASS over R16_Q1, run by run :: sha256 c6c4138ee5c86423 :: prose twin of the run farther from 1
\newcommand{\MlstFtCmpJcTwoQcdResOverFullEFive}{1.060\,\ensuremath{\pm}\,0.002}  % experiments/FIGS/data/v2/index.json :: fine-tuning jc2_qcd N=100000: (1 - macro_auc_qcd) of L162 over L188, run by run :: sha256 c6c4138ee5c86423
\newcommand{\MlstFtCmpJcTwoQcdResOverFullEFivePr}{6.0\% more errors}  % experiments/FIGS/data/v2/index.json :: fine-tuning jc2_qcd N=100000: (1 - macro_auc_qcd) of L162 over L188, run by run :: sha256 c6c4138ee5c86423 :: prose twin: two significant figures, direction in words
\newcommand{\MlstFtCmpJcTwoQcdMidOverFullEFive}{1.062\,\ensuremath{\pm}\,0.010}  % experiments/FIGS/data/v2/index.json :: fine-tuning jc2_qcd N=100000: (1 - macro_auc_qcd) of R42_Q1 over L188, run by run :: sha256 c6c4138ee5c86423
\newcommand{\MlstFtCmpJcTwoQcdMidOverFullEFivePr}{6.2\% more errors}  % experiments/FIGS/data/v2/index.json :: fine-tuning jc2_qcd N=100000: (1 - macro_auc_qcd) of R42_Q1 over L188, run by run :: sha256 c6c4138ee5c86423 :: prose twin: two significant figures, direction in words
\newcommand{\MlstFtCmpJcTwoQcdCoarseOverFullEFive}{1.137\,\ensuremath{\pm}\,0.001}  % experiments/FIGS/data/v2/index.json :: fine-tuning jc2_qcd N=100000: (1 - macro_auc_qcd) of R16_Q1 over L188, run by run :: sha256 c6c4138ee5c86423
\newcommand{\MlstFtCmpJcTwoQcdCoarseOverFullEFivePr}{14\% more errors}  % experiments/FIGS/data/v2/index.json :: fine-tuning jc2_qcd N=100000: (1 - macro_auc_qcd) of R16_Q1 over L188, run by run :: sha256 c6c4138ee5c86423 :: prose twin: two significant figures, direction in words
\newcommand{\MlstFtCmpJcTwoQcdResMassOverResEFive}{1.022 and 1.009}  % experiments/FIGS/data/v2/index.json :: fine-tuning jc2_qcd N=100000: (1 - macro_auc_qcd) of L162_MASS over L162, run by run :: sha256 c6c4138ee5c86423
\newcommand{\MlstFtCmpJcTwoQcdResMassOverResEFivePr}{equal within 3\%}  % experiments/FIGS/data/v2/index.json :: fine-tuning jc2_qcd N=100000: (1 - macro_auc_qcd) of L162_MASS over L162, run by run :: sha256 c6c4138ee5c86423 :: prose twin: two significant figures, direction in words
\newcommand{\MlstFtCmpJcTwoQcdResMassOverResEFivePrMax}{equal within 3\%}  % experiments/FIGS/data/v2/index.json :: fine-tuning jc2_qcd N=100000: (1 - macro_auc_qcd) of L162_MASS over L162, run by run :: sha256 c6c4138ee5c86423 :: prose twin of the run farther from 1
\newcommand{\MlstFtCmpJcTwoQcdCoarseMassOverCoarseEFive}{1.028 and 1.027}  % experiments/FIGS/data/v2/index.json :: fine-tuning jc2_qcd N=100000: (1 - macro_auc_qcd) of R16_Q1_MASS over R16_Q1, run by run :: sha256 c6c4138ee5c86423
\newcommand{\MlstFtCmpJcTwoQcdCoarseMassOverCoarseEFivePr}{at least 2.7\% more errors}  % experiments/FIGS/data/v2/index.json :: fine-tuning jc2_qcd N=100000: (1 - macro_auc_qcd) of R16_Q1_MASS over R16_Q1, run by run :: sha256 c6c4138ee5c86423 :: prose twin: two significant figures, direction in words
\newcommand{\MlstFtCmpJcTwoQcdCoarseMassOverCoarseEFivePrMax}{up to 2.8\% more errors}  % experiments/FIGS/data/v2/index.json :: fine-tuning jc2_qcd N=100000: (1 - macro_auc_qcd) of R16_Q1_MASS over R16_Q1, run by run :: sha256 c6c4138ee5c86423 :: prose twin of the run farther from 1
\newcommand{\MlstFtCmpJcTwoQcdResOverFullESix}{1.035\,\ensuremath{\pm}\,0.010}  % experiments/FIGS/data/v2/index.json :: fine-tuning jc2_qcd N=1000000: (1 - macro_auc_qcd) of L162 over L188, run by run :: sha256 c6c4138ee5c86423
\newcommand{\MlstFtCmpJcTwoQcdResOverFullESixPr}{3.5\% more errors}  % experiments/FIGS/data/v2/index.json :: fine-tuning jc2_qcd N=1000000: (1 - macro_auc_qcd) of L162 over L188, run by run :: sha256 c6c4138ee5c86423 :: prose twin: two significant figures, direction in words
\newcommand{\MlstFtCmpJcTwoQcdMidOverFullESix}{1.050\,\ensuremath{\pm}\,0.004}  % experiments/FIGS/data/v2/index.json :: fine-tuning jc2_qcd N=1000000: (1 - macro_auc_qcd) of R42_Q1 over L188, run by run :: sha256 c6c4138ee5c86423
\newcommand{\MlstFtCmpJcTwoQcdMidOverFullESixPr}{5.0\% more errors}  % experiments/FIGS/data/v2/index.json :: fine-tuning jc2_qcd N=1000000: (1 - macro_auc_qcd) of R42_Q1 over L188, run by run :: sha256 c6c4138ee5c86423 :: prose twin: two significant figures, direction in words
\newcommand{\MlstFtCmpJcTwoQcdCoarseOverFullESix}{1.122\,\ensuremath{\pm}\,0.002}  % experiments/FIGS/data/v2/index.json :: fine-tuning jc2_qcd N=1000000: (1 - macro_auc_qcd) of R16_Q1 over L188, run by run :: sha256 c6c4138ee5c86423
\newcommand{\MlstFtCmpJcTwoQcdCoarseOverFullESixPr}{12\% more errors}  % experiments/FIGS/data/v2/index.json :: fine-tuning jc2_qcd N=1000000: (1 - macro_auc_qcd) of R16_Q1 over L188, run by run :: sha256 c6c4138ee5c86423 :: prose twin: two significant figures, direction in words
\newcommand{\MlstFtCmpJcTwoQcdResMassOverResESix}{1.011 and 1.025}  % experiments/FIGS/data/v2/index.json :: fine-tuning jc2_qcd N=1000000: (1 - macro_auc_qcd) of L162_MASS over L162, run by run :: sha256 c6c4138ee5c86423
\newcommand{\MlstFtCmpJcTwoQcdResMassOverResESixPr}{equal within 3\%}  % experiments/FIGS/data/v2/index.json :: fine-tuning jc2_qcd N=1000000: (1 - macro_auc_qcd) of L162_MASS over L162, run by run :: sha256 c6c4138ee5c86423 :: prose twin: two significant figures, direction in words
\newcommand{\MlstFtCmpJcTwoQcdResMassOverResESixPrMax}{equal within 3\%}  % experiments/FIGS/data/v2/index.json :: fine-tuning jc2_qcd N=1000000: (1 - macro_auc_qcd) of L162_MASS over L162, run by run :: sha256 c6c4138ee5c86423 :: prose twin of the run farther from 1
\newcommand{\MlstFtCmpJcTwoQcdCoarseMassOverCoarseESix}{1.015 and 1.015}  % experiments/FIGS/data/v2/index.json :: fine-tuning jc2_qcd N=1000000: (1 - macro_auc_qcd) of R16_Q1_MASS over R16_Q1, run by run :: sha256 c6c4138ee5c86423
\newcommand{\MlstFtCmpJcTwoQcdCoarseMassOverCoarseESixPr}{at least 1.5\% more errors}  % experiments/FIGS/data/v2/index.json :: fine-tuning jc2_qcd N=1000000: (1 - macro_auc_qcd) of R16_Q1_MASS over R16_Q1, run by run :: sha256 c6c4138ee5c86423 :: prose twin: two significant figures, direction in words
\newcommand{\MlstFtCmpJcTwoQcdCoarseMassOverCoarseESixPrMax}{up to 1.5\% more errors}  % experiments/FIGS/data/v2/index.json :: fine-tuning jc2_qcd N=1000000: (1 - macro_auc_qcd) of R16_Q1_MASS over R16_Q1, run by run :: sha256 c6c4138ee5c86423 :: prose twin of the run farther from 1
\newcommand{\MlstFtJcTwoAccFullEFour}{0.3914\,\ensuremath{\pm}\,0.0011}  % experiments/FIGS/data/v2/index.json :: fine-tuning jc2 N=10000: accuracy, L188 :: sha256 c6c4138ee5c86423
\newcommand{\MlstFtJcTwoAccResEFour}{0.3934\,\ensuremath{\pm}\,0.0018}  % experiments/FIGS/data/v2/index.json :: fine-tuning jc2 N=10000: accuracy, L162 :: sha256 c6c4138ee5c86423
\newcommand{\MlstFtJcTwoAccMidEFour}{0.3563\,\ensuremath{\pm}\,0.0025}  % experiments/FIGS/data/v2/index.json :: fine-tuning jc2 N=10000: accuracy, R42_Q1 :: sha256 c6c4138ee5c86423
\newcommand{\MlstFtJcTwoAccCoarseEFour}{0.3219\,\ensuremath{\pm}\,0.0011}  % experiments/FIGS/data/v2/index.json :: fine-tuning jc2 N=10000: accuracy, R16_Q1 :: sha256 c6c4138ee5c86423
\newcommand{\MlstFtJcTwoAccResMassEFour}{0.381 and 0.392}  % experiments/FIGS/data/v2/index.json :: fine-tuning jc2 N=10000: accuracy, L162_MASS :: sha256 c6c4138ee5c86423
\newcommand{\MlstFtJcTwoAccCoarseMassEFour}{0.3156 and 0.3187}  % experiments/FIGS/data/v2/index.json :: fine-tuning jc2 N=10000: accuracy, R16_Q1_MASS :: sha256 c6c4138ee5c86423
\newcommand{\MlstFtJcTwoAccScratchEFour}{0.1492\,\ensuremath{\pm}\,0.0011}  % experiments/FIGS/data/v2/index.json :: fine-tuning jc2 N=10000: accuracy from scratch :: sha256 c6c4138ee5c86423
\newcommand{\MlstFtJcTwoAccFullMinusScratchEFour}{0.2423\,\ensuremath{\pm}\,0.0011}  % experiments/FIGS/data/v2/index.json :: fine-tuning jc2 N=10000: accuracy of each 188-class run minus the from-scratch mean (absolute) :: sha256 c6c4138ee5c86423
\newcommand{\MlstFtJcTwoAccFullEFive}{0.4164\,\ensuremath{\pm}\,0.0015}  % experiments/FIGS/data/v2/index.json :: fine-tuning jc2 N=100000: accuracy, L188 :: sha256 c6c4138ee5c86423
\newcommand{\MlstFtJcTwoAccResEFive}{0.4161\,\ensuremath{\pm}\,0.0015}  % experiments/FIGS/data/v2/index.json :: fine-tuning jc2 N=100000: accuracy, L162 :: sha256 c6c4138ee5c86423
\newcommand{\MlstFtJcTwoAccMidEFive}{0.4007\,\ensuremath{\pm}\,0.0013}  % experiments/FIGS/data/v2/index.json :: fine-tuning jc2 N=100000: accuracy, R42_Q1 :: sha256 c6c4138ee5c86423
\newcommand{\MlstFtJcTwoAccCoarseEFive}{0.3833\,\ensuremath{\pm}\,0.0020}  % experiments/FIGS/data/v2/index.json :: fine-tuning jc2 N=100000: accuracy, R16_Q1 :: sha256 c6c4138ee5c86423
\newcommand{\MlstFtJcTwoAccResMassEFive}{0.4093 and 0.4115}  % experiments/FIGS/data/v2/index.json :: fine-tuning jc2 N=100000: accuracy, L162_MASS :: sha256 c6c4138ee5c86423
\newcommand{\MlstFtJcTwoAccCoarseMassEFive}{0.3758 and 0.3781}  % experiments/FIGS/data/v2/index.json :: fine-tuning jc2 N=100000: accuracy, R16_Q1_MASS :: sha256 c6c4138ee5c86423
\newcommand{\MlstFtJcTwoAccScratchEFive}{0.292\,\ensuremath{\pm}\,0.004}  % experiments/FIGS/data/v2/index.json :: fine-tuning jc2 N=100000: accuracy from scratch :: sha256 c6c4138ee5c86423
\newcommand{\MlstFtJcTwoAccFullMinusScratchEFive}{0.1243\,\ensuremath{\pm}\,0.0015}  % experiments/FIGS/data/v2/index.json :: fine-tuning jc2 N=100000: accuracy of each 188-class run minus the from-scratch mean (absolute) :: sha256 c6c4138ee5c86423
\newcommand{\MlstFtJcTwoAccFullESix}{0.4263\,\ensuremath{\pm}\,0.0005}  % experiments/FIGS/data/v2/index.json :: fine-tuning jc2 N=1000000: accuracy, L188 :: sha256 c6c4138ee5c86423
\newcommand{\MlstFtJcTwoAccResESix}{0.4249\,\ensuremath{\pm}\,0.0021}  % experiments/FIGS/data/v2/index.json :: fine-tuning jc2 N=1000000: accuracy, L162 :: sha256 c6c4138ee5c86423
\newcommand{\MlstFtJcTwoAccMidESix}{0.4166\,\ensuremath{\pm}\,0.0011}  % experiments/FIGS/data/v2/index.json :: fine-tuning jc2 N=1000000: accuracy, R42_Q1 :: sha256 c6c4138ee5c86423
\newcommand{\MlstFtJcTwoAccCoarseESix}{0.4054\,\ensuremath{\pm}\,0.0020}  % experiments/FIGS/data/v2/index.json :: fine-tuning jc2 N=1000000: accuracy, R16_Q1 :: sha256 c6c4138ee5c86423
\newcommand{\MlstFtJcTwoAccResMassESix}{0.4198 and 0.4213}  % experiments/FIGS/data/v2/index.json :: fine-tuning jc2 N=1000000: accuracy, L162_MASS :: sha256 c6c4138ee5c86423
\newcommand{\MlstFtJcTwoAccCoarseMassESix}{0.4012 and 0.4000}  % experiments/FIGS/data/v2/index.json :: fine-tuning jc2 N=1000000: accuracy, R16_Q1_MASS :: sha256 c6c4138ee5c86423
\newcommand{\MlstFtJcTwoAccScratchESix}{0.3677\,\ensuremath{\pm}\,0.0012}  % experiments/FIGS/data/v2/index.json :: fine-tuning jc2 N=1000000: accuracy from scratch :: sha256 c6c4138ee5c86423
\newcommand{\MlstFtJcTwoAccFullMinusScratchESix}{0.0585\,\ensuremath{\pm}\,0.0005}  % experiments/FIGS/data/v2/index.json :: fine-tuning jc2 N=1000000: accuracy of each 188-class run minus the from-scratch mean (absolute) :: sha256 c6c4138ee5c86423
\newcommand{\MlstFtJcTwoAucFull}{0.9797\,\ensuremath{\pm}\,0.0001}  % experiments/FIGS/data/v2/index.json :: fine-tuning class-balanced JetClass-II N=1000000: AUC averaged over the classes, Full :: sha256 c6c4138ee5c86423
\newcommand{\MlstFtJcTwoAucScratch}{0.9734\,\ensuremath{\pm}\,0.0003}  % experiments/FIGS/data/v2/index.json :: fine-tuning class-balanced JetClass-II N=1000000: AUC averaged over the classes, Scratch :: sha256 c6c4138ee5c86423
\newcommand{\MlstFtJcOneAccFullEFour}{0.8117\,\ensuremath{\pm}\,0.0031}  % experiments/FIGS/data/v2/index.json :: fine-tuning jc1 N=10000: accuracy, L188 :: sha256 c6c4138ee5c86423
\newcommand{\MlstFtJcOneAccResEFour}{0.8117\,\ensuremath{\pm}\,0.0013}  % experiments/FIGS/data/v2/index.json :: fine-tuning jc1 N=10000: accuracy, L162 :: sha256 c6c4138ee5c86423
\newcommand{\MlstFtJcOneAccMidEFour}{0.8086\,\ensuremath{\pm}\,0.0031}  % experiments/FIGS/data/v2/index.json :: fine-tuning jc1 N=10000: accuracy, R42_Q1 :: sha256 c6c4138ee5c86423
\newcommand{\MlstFtJcOneAccCoarseEFour}{0.7982\,\ensuremath{\pm}\,0.0023}  % experiments/FIGS/data/v2/index.json :: fine-tuning jc1 N=10000: accuracy, R16_Q1 :: sha256 c6c4138ee5c86423
\newcommand{\MlstFtJcOneAccResMassEFour}{0.8107 and 0.8098}  % experiments/FIGS/data/v2/index.json :: fine-tuning jc1 N=10000: accuracy, L162_MASS :: sha256 c6c4138ee5c86423
\newcommand{\MlstFtJcOneAccCoarseMassEFour}{0.7944 and 0.7990}  % experiments/FIGS/data/v2/index.json :: fine-tuning jc1 N=10000: accuracy, R16_Q1_MASS :: sha256 c6c4138ee5c86423
\newcommand{\MlstFtJcOneAccScratchEFour}{0.677\,\ensuremath{\pm}\,0.004}  % experiments/FIGS/data/v2/index.json :: fine-tuning jc1 N=10000: accuracy from scratch :: sha256 c6c4138ee5c86423
\newcommand{\MlstFtJcOneAccFullMinusScratchEFour}{0.1343\,\ensuremath{\pm}\,0.0031}  % experiments/FIGS/data/v2/index.json :: fine-tuning jc1 N=10000: accuracy of each 188-class run minus the from-scratch mean (absolute) :: sha256 c6c4138ee5c86423
\newcommand{\MlstFtJcOneAccFullEFive}{0.8343\,\ensuremath{\pm}\,0.0035}  % experiments/FIGS/data/v2/index.json :: fine-tuning jc1 N=100000: accuracy, L188 :: sha256 c6c4138ee5c86423
\newcommand{\MlstFtJcOneAccResEFive}{0.8328\,\ensuremath{\pm}\,0.0005}  % experiments/FIGS/data/v2/index.json :: fine-tuning jc1 N=100000: accuracy, L162 :: sha256 c6c4138ee5c86423
\newcommand{\MlstFtJcOneAccMidEFive}{0.8324\,\ensuremath{\pm}\,0.0004}  % experiments/FIGS/data/v2/index.json :: fine-tuning jc1 N=100000: accuracy, R42_Q1 :: sha256 c6c4138ee5c86423
\newcommand{\MlstFtJcOneAccCoarseEFive}{0.8288\,\ensuremath{\pm}\,0.0005}  % experiments/FIGS/data/v2/index.json :: fine-tuning jc1 N=100000: accuracy, R16_Q1 :: sha256 c6c4138ee5c86423
\newcommand{\MlstFtJcOneAccResMassEFive}{0.8326 and 0.8348}  % experiments/FIGS/data/v2/index.json :: fine-tuning jc1 N=100000: accuracy, L162_MASS :: sha256 c6c4138ee5c86423
\newcommand{\MlstFtJcOneAccCoarseMassEFive}{0.8253 and 0.8262}  % experiments/FIGS/data/v2/index.json :: fine-tuning jc1 N=100000: accuracy, R16_Q1_MASS :: sha256 c6c4138ee5c86423
\newcommand{\MlstFtJcOneAccScratchEFive}{0.777\,\ensuremath{\pm}\,0.005}  % experiments/FIGS/data/v2/index.json :: fine-tuning jc1 N=100000: accuracy from scratch :: sha256 c6c4138ee5c86423
\newcommand{\MlstFtJcOneAccFullMinusScratchEFive}{0.0571\,\ensuremath{\pm}\,0.0035}  % experiments/FIGS/data/v2/index.json :: fine-tuning jc1 N=100000: accuracy of each 188-class run minus the from-scratch mean (absolute) :: sha256 c6c4138ee5c86423
\newcommand{\MlstFtJcOneAccFullESix}{0.8468\,\ensuremath{\pm}\,0.0006}  % experiments/FIGS/data/v2/index.json :: fine-tuning jc1 N=1000000: accuracy, L188 :: sha256 c6c4138ee5c86423
\newcommand{\MlstFtJcOneAccResESix}{0.8475\,\ensuremath{\pm}\,0.0009}  % experiments/FIGS/data/v2/index.json :: fine-tuning jc1 N=1000000: accuracy, L162 :: sha256 c6c4138ee5c86423
\newcommand{\MlstFtJcOneAccMidESix}{0.8459\,\ensuremath{\pm}\,0.0002}  % experiments/FIGS/data/v2/index.json :: fine-tuning jc1 N=1000000: accuracy, R42_Q1 :: sha256 c6c4138ee5c86423
\newcommand{\MlstFtJcOneAccCoarseESix}{0.8422\,\ensuremath{\pm}\,0.0005}  % experiments/FIGS/data/v2/index.json :: fine-tuning jc1 N=1000000: accuracy, R16_Q1 :: sha256 c6c4138ee5c86423
\newcommand{\MlstFtJcOneAccResMassESix}{0.8459 and 0.8471}  % experiments/FIGS/data/v2/index.json :: fine-tuning jc1 N=1000000: accuracy, L162_MASS :: sha256 c6c4138ee5c86423
\newcommand{\MlstFtJcOneAccCoarseMassESix}{0.8400 and 0.8398}  % experiments/FIGS/data/v2/index.json :: fine-tuning jc1 N=1000000: accuracy, R16_Q1_MASS :: sha256 c6c4138ee5c86423
\newcommand{\MlstFtJcOneAccScratchESix}{0.8267\,\ensuremath{\pm}\,0.0011}  % experiments/FIGS/data/v2/index.json :: fine-tuning jc1 N=1000000: accuracy from scratch :: sha256 c6c4138ee5c86423
\newcommand{\MlstFtJcOneAccFullMinusScratchESix}{0.0201\,\ensuremath{\pm}\,0.0006}  % experiments/FIGS/data/v2/index.json :: fine-tuning jc1 N=1000000: accuracy of each 188-class run minus the from-scratch mean (absolute) :: sha256 c6c4138ee5c86423
\newcommand{\MlstFtTopRthirtyScratchEThree}{49\,\ensuremath{\pm}\,12}  % experiments/FIGS/data/v2/index.json :: fine-tuning top N=1000: R30 from scratch, every seed :: sha256 c6c4138ee5c86423
\newcommand{\MlstFtTopRthirtyFullEThree}{72\,\ensuremath{\pm}\,19}  % experiments/FIGS/data/v2/index.json :: fine-tuning top N=1000: R30, L188 :: sha256 c6c4138ee5c86423
\newcommand{\MlstFtTopRthirtyFullOverScratchEThree}{1.5\,\ensuremath{\pm}\,0.4}  % experiments/FIGS/data/v2/index.json :: fine-tuning top N=1000: R30 of each L188 run over the from-scratch mean, SD over runs and standard error of that mean in quadrature :: sha256 c6c4138ee5c86423
\newcommand{\MlstFtTopRthirtyFullOverScratchEThreePr}{equal within 50\%}  % experiments/FIGS/data/v2/index.json :: fine-tuning top N=1000: R30 of each L188 run over the from-scratch mean, SD over runs and standard error of that mean in quadrature :: sha256 c6c4138ee5c86423 :: prose twin: two significant figures, direction in words
\newcommand{\MlstFtTopRthirtyResEThree}{135\,\ensuremath{\pm}\,26}  % experiments/FIGS/data/v2/index.json :: fine-tuning top N=1000: R30, L162 :: sha256 c6c4138ee5c86423
\newcommand{\MlstFtTopRthirtyResOverScratchEThree}{2.8\,\ensuremath{\pm}\,0.7}  % experiments/FIGS/data/v2/index.json :: fine-tuning top N=1000: R30 of each L162 run over the from-scratch mean, SD over runs and standard error of that mean in quadrature :: sha256 c6c4138ee5c86423
\newcommand{\MlstFtTopRthirtyResOverScratchEThreePr}{2.8 times the rejection}  % experiments/FIGS/data/v2/index.json :: fine-tuning top N=1000: R30 of each L162 run over the from-scratch mean, SD over runs and standard error of that mean in quadrature :: sha256 c6c4138ee5c86423 :: prose twin: two significant figures, direction in words
\newcommand{\MlstFtTopRthirtyMidEThree}{104\,\ensuremath{\pm}\,20}  % experiments/FIGS/data/v2/index.json :: fine-tuning top N=1000: R30, R42_Q1 :: sha256 c6c4138ee5c86423
\newcommand{\MlstFtTopRthirtyMidOverScratchEThree}{2.1\,\ensuremath{\pm}\,0.5}  % experiments/FIGS/data/v2/index.json :: fine-tuning top N=1000: R30 of each R42_Q1 run over the from-scratch mean, SD over runs and standard error of that mean in quadrature :: sha256 c6c4138ee5c86423
\newcommand{\MlstFtTopRthirtyMidOverScratchEThreePr}{2.1 times the rejection}  % experiments/FIGS/data/v2/index.json :: fine-tuning top N=1000: R30 of each R42_Q1 run over the from-scratch mean, SD over runs and standard error of that mean in quadrature :: sha256 c6c4138ee5c86423 :: prose twin: two significant figures, direction in words
\newcommand{\MlstFtTopRthirtyCoarseEThree}{83\,\ensuremath{\pm}\,25}  % experiments/FIGS/data/v2/index.json :: fine-tuning top N=1000: R30, R16_Q1 :: sha256 c6c4138ee5c86423
\newcommand{\MlstFtTopRthirtyCoarseOverScratchEThree}{1.7\,\ensuremath{\pm}\,0.6}  % experiments/FIGS/data/v2/index.json :: fine-tuning top N=1000: R30 of each R16_Q1 run over the from-scratch mean, SD over runs and standard error of that mean in quadrature :: sha256 c6c4138ee5c86423
\newcommand{\MlstFtTopRthirtyCoarseOverScratchEThreePr}{equal within 80\%}  % experiments/FIGS/data/v2/index.json :: fine-tuning top N=1000: R30 of each R16_Q1 run over the from-scratch mean, SD over runs and standard error of that mean in quadrature :: sha256 c6c4138ee5c86423 :: prose twin: two significant figures, direction in words
\newcommand{\MlstFtTopRthirtyResMassEThree}{105.9 and 112.7}  % experiments/FIGS/data/v2/index.json :: fine-tuning top N=1000: R30, L162_MASS :: sha256 c6c4138ee5c86423
\newcommand{\MlstFtTopRthirtyResMassOverScratchEThree}{2.18 and 2.32}  % experiments/FIGS/data/v2/index.json :: fine-tuning top N=1000: R30 of each L162_MASS run over the from-scratch mean, SD over runs and standard error of that mean in quadrature :: sha256 c6c4138ee5c86423
\newcommand{\MlstFtTopRthirtyResMassOverScratchEThreePr}{at least 2.2 times the rejection}  % experiments/FIGS/data/v2/index.json :: fine-tuning top N=1000: R30 of each L162_MASS run over the from-scratch mean, SD over runs and standard error of that mean in quadrature :: sha256 c6c4138ee5c86423 :: prose twin: two significant figures, direction in words
\newcommand{\MlstFtTopRthirtyResMassOverScratchEThreePrMax}{up to 2.3 times the rejection}  % experiments/FIGS/data/v2/index.json :: fine-tuning top N=1000: R30 of each L162_MASS run over the from-scratch mean, SD over runs and standard error of that mean in quadrature :: sha256 c6c4138ee5c86423 :: prose twin of the run farther from 1
\newcommand{\MlstFtTopRthirtyCoarseMassEThree}{59.1 and 63.9}  % experiments/FIGS/data/v2/index.json :: fine-tuning top N=1000: R30, R16_Q1_MASS :: sha256 c6c4138ee5c86423
\newcommand{\MlstFtTopRthirtyCoarseMassOverScratchEThree}{1.216 and 1.314}  % experiments/FIGS/data/v2/index.json :: fine-tuning top N=1000: R30 of each R16_Q1_MASS run over the from-scratch mean, SD over runs and standard error of that mean in quadrature :: sha256 c6c4138ee5c86423
\newcommand{\MlstFtTopRthirtyCoarseMassOverScratchEThreePr}{at least 22\% more rejection}  % experiments/FIGS/data/v2/index.json :: fine-tuning top N=1000: R30 of each R16_Q1_MASS run over the from-scratch mean, SD over runs and standard error of that mean in quadrature :: sha256 c6c4138ee5c86423 :: prose twin: two significant figures, direction in words
\newcommand{\MlstFtTopRthirtyCoarseMassOverScratchEThreePrMax}{up to 31\% more rejection}  % experiments/FIGS/data/v2/index.json :: fine-tuning top N=1000: R30 of each R16_Q1_MASS run over the from-scratch mean, SD over runs and standard error of that mean in quadrature :: sha256 c6c4138ee5c86423 :: prose twin of the run farther from 1
\newcommand{\MlstFtTopRthirtyResOverFullEThree}{2.0\,\ensuremath{\pm}\,0.8}  % experiments/FIGS/data/v2/index.json :: fine-tuning top N=1000: R30 of L162 over L188, run by run :: sha256 c6c4138ee5c86423
\newcommand{\MlstFtTopRthirtyResOverFullEThreePr}{equal within 100\%}  % experiments/FIGS/data/v2/index.json :: fine-tuning top N=1000: R30 of L162 over L188, run by run :: sha256 c6c4138ee5c86423 :: prose twin: two significant figures, direction in words
\newcommand{\MlstFtTopRthirtyMidOverFullEThree}{1.46\,\ensuremath{\pm}\,0.12}  % experiments/FIGS/data/v2/index.json :: fine-tuning top N=1000: R30 of R42_Q1 over L188, run by run :: sha256 c6c4138ee5c86423
\newcommand{\MlstFtTopRthirtyMidOverFullEThreePr}{46\% more rejection}  % experiments/FIGS/data/v2/index.json :: fine-tuning top N=1000: R30 of R42_Q1 over L188, run by run :: sha256 c6c4138ee5c86423 :: prose twin: two significant figures, direction in words
\newcommand{\MlstFtTopRthirtyCoarseOverFullEThree}{1.3\,\ensuremath{\pm}\,0.7}  % experiments/FIGS/data/v2/index.json :: fine-tuning top N=1000: R30 of R16_Q1 over L188, run by run :: sha256 c6c4138ee5c86423
\newcommand{\MlstFtTopRthirtyCoarseOverFullEThreePr}{equal within 80\%}  % experiments/FIGS/data/v2/index.json :: fine-tuning top N=1000: R30 of R16_Q1 over L188, run by run :: sha256 c6c4138ee5c86423 :: prose twin: two significant figures, direction in words
\newcommand{\MlstFtTopRthirtyResMassOverResEThree}{0.73 and 1.07}  % experiments/FIGS/data/v2/index.json :: fine-tuning top N=1000: R30 of L162_MASS over L162, run by run :: sha256 c6c4138ee5c86423
\newcommand{\MlstFtTopRthirtyResMassOverResEThreePr}{equal within 30\%}  % experiments/FIGS/data/v2/index.json :: fine-tuning top N=1000: R30 of L162_MASS over L162, run by run :: sha256 c6c4138ee5c86423 :: prose twin: two significant figures, direction in words
\newcommand{\MlstFtTopRthirtyResMassOverResEThreePrMax}{equal within 30\%}  % experiments/FIGS/data/v2/index.json :: fine-tuning top N=1000: R30 of L162_MASS over L162, run by run :: sha256 c6c4138ee5c86423 :: prose twin of the run farther from 1
\newcommand{\MlstFtTopRthirtyCoarseMassOverCoarseEThree}{0.56 and 0.75}  % experiments/FIGS/data/v2/index.json :: fine-tuning top N=1000: R30 of R16_Q1_MASS over R16_Q1, run by run :: sha256 c6c4138ee5c86423
\newcommand{\MlstFtTopRthirtyCoarseMassOverCoarseEThreePr}{equal within 50\%}  % experiments/FIGS/data/v2/index.json :: fine-tuning top N=1000: R30 of R16_Q1_MASS over R16_Q1, run by run :: sha256 c6c4138ee5c86423 :: prose twin: two significant figures, direction in words
\newcommand{\MlstFtTopRthirtyCoarseMassOverCoarseEThreePrMax}{equal within 50\%}  % experiments/FIGS/data/v2/index.json :: fine-tuning top N=1000: R30 of R16_Q1_MASS over R16_Q1, run by run :: sha256 c6c4138ee5c86423 :: prose twin of the run farther from 1
\newcommand{\MlstFtTopRthirtyRandOverFullEThree}{1.3\,\ensuremath{\pm}\,0.5}  % experiments/FIGS/data/v2/index.json :: fine-tuning top N=1000: R30 of each random partition over the L188 run of its seed :: sha256 c6c4138ee5c86423
\newcommand{\MlstFtTopRthirtyRandOverFullEThreePr}{equal within 50\%}  % experiments/FIGS/data/v2/index.json :: fine-tuning top N=1000: R30 of each random partition over the L188 run of its seed :: sha256 c6c4138ee5c86423 :: prose twin: two significant figures, direction in words
\newcommand{\MlstFtTopRthirtyRandOverCoarseEThree}{0.87\,\ensuremath{\pm}\,0.12}  % experiments/FIGS/data/v2/index.json :: fine-tuning top N=1000: R30 of each random partition over the R16_Q1 run of its seed :: sha256 c6c4138ee5c86423
\newcommand{\MlstFtTopRthirtyRandOverCoarseEThreePr}{equal within 20\%}  % experiments/FIGS/data/v2/index.json :: fine-tuning top N=1000: R30 of each random partition over the R16_Q1 run of its seed :: sha256 c6c4138ee5c86423 :: prose twin: two significant figures, direction in words
\newcommand{\MlstFtTopRthirtyScratchEFour}{340\,\ensuremath{\pm}\,50}  % experiments/FIGS/data/v2/index.json :: fine-tuning top N=10000: R30 from scratch, every seed :: sha256 c6c4138ee5c86423
\newcommand{\MlstFtTopRthirtyFullEFour}{420\,\ensuremath{\pm}\,9}  % experiments/FIGS/data/v2/index.json :: fine-tuning top N=10000: R30, L188 :: sha256 c6c4138ee5c86423
\newcommand{\MlstFtTopRthirtyFullOverScratchEFour}{1.24\,\ensuremath{\pm}\,0.11}  % experiments/FIGS/data/v2/index.json :: fine-tuning top N=10000: R30 of each L188 run over the from-scratch mean, SD over runs and standard error of that mean in quadrature :: sha256 c6c4138ee5c86423
\newcommand{\MlstFtTopRthirtyFullOverScratchEFourPr}{24\% more rejection}  % experiments/FIGS/data/v2/index.json :: fine-tuning top N=10000: R30 of each L188 run over the from-scratch mean, SD over runs and standard error of that mean in quadrature :: sha256 c6c4138ee5c86423 :: prose twin: two significant figures, direction in words
\newcommand{\MlstFtTopRthirtyResEFour}{430\,\ensuremath{\pm}\,60}  % experiments/FIGS/data/v2/index.json :: fine-tuning top N=10000: R30, L162 :: sha256 c6c4138ee5c86423
\newcommand{\MlstFtTopRthirtyResOverScratchEFour}{1.27\,\ensuremath{\pm}\,0.21}  % experiments/FIGS/data/v2/index.json :: fine-tuning top N=10000: R30 of each L162 run over the from-scratch mean, SD over runs and standard error of that mean in quadrature :: sha256 c6c4138ee5c86423
\newcommand{\MlstFtTopRthirtyResOverScratchEFourPr}{equal within 30\%}  % experiments/FIGS/data/v2/index.json :: fine-tuning top N=10000: R30 of each L162 run over the from-scratch mean, SD over runs and standard error of that mean in quadrature :: sha256 c6c4138ee5c86423 :: prose twin: two significant figures, direction in words
\newcommand{\MlstFtTopRthirtyMidEFour}{455\,\ensuremath{\pm}\,30}  % experiments/FIGS/data/v2/index.json :: fine-tuning top N=10000: R30, R42_Q1 :: sha256 c6c4138ee5c86423
\newcommand{\MlstFtTopRthirtyMidOverScratchEFour}{1.34\,\ensuremath{\pm}\,0.14}  % experiments/FIGS/data/v2/index.json :: fine-tuning top N=10000: R30 of each R42_Q1 run over the from-scratch mean, SD over runs and standard error of that mean in quadrature :: sha256 c6c4138ee5c86423
\newcommand{\MlstFtTopRthirtyMidOverScratchEFourPr}{34\% more rejection}  % experiments/FIGS/data/v2/index.json :: fine-tuning top N=10000: R30 of each R42_Q1 run over the from-scratch mean, SD over runs and standard error of that mean in quadrature :: sha256 c6c4138ee5c86423 :: prose twin: two significant figures, direction in words
\newcommand{\MlstFtTopRthirtyCoarseEFour}{350\,\ensuremath{\pm}\,40}  % experiments/FIGS/data/v2/index.json :: fine-tuning top N=10000: R30, R16_Q1 :: sha256 c6c4138ee5c86423
\newcommand{\MlstFtTopRthirtyCoarseOverScratchEFour}{1.04\,\ensuremath{\pm}\,0.14}  % experiments/FIGS/data/v2/index.json :: fine-tuning top N=10000: R30 of each R16_Q1 run over the from-scratch mean, SD over runs and standard error of that mean in quadrature :: sha256 c6c4138ee5c86423
\newcommand{\MlstFtTopRthirtyCoarseOverScratchEFourPr}{equal within 20\%}  % experiments/FIGS/data/v2/index.json :: fine-tuning top N=10000: R30 of each R16_Q1 run over the from-scratch mean, SD over runs and standard error of that mean in quadrature :: sha256 c6c4138ee5c86423 :: prose twin: two significant figures, direction in words
\newcommand{\MlstFtTopRthirtyResMassEFour}{373.2 and 376.7}  % experiments/FIGS/data/v2/index.json :: fine-tuning top N=10000: R30, L162_MASS :: sha256 c6c4138ee5c86423
\newcommand{\MlstFtTopRthirtyResMassOverScratchEFour}{1.101 and 1.111}  % experiments/FIGS/data/v2/index.json :: fine-tuning top N=10000: R30 of each L162_MASS run over the from-scratch mean, SD over runs and standard error of that mean in quadrature :: sha256 c6c4138ee5c86423
\newcommand{\MlstFtTopRthirtyResMassOverScratchEFourPr}{at least 10\% more rejection}  % experiments/FIGS/data/v2/index.json :: fine-tuning top N=10000: R30 of each L162_MASS run over the from-scratch mean, SD over runs and standard error of that mean in quadrature :: sha256 c6c4138ee5c86423 :: prose twin: two significant figures, direction in words
\newcommand{\MlstFtTopRthirtyResMassOverScratchEFourPrMax}{up to 11\% more rejection}  % experiments/FIGS/data/v2/index.json :: fine-tuning top N=10000: R30 of each L162_MASS run over the from-scratch mean, SD over runs and standard error of that mean in quadrature :: sha256 c6c4138ee5c86423 :: prose twin of the run farther from 1
\newcommand{\MlstFtTopRthirtyCoarseMassEFour}{365.8 and 359.9}  % experiments/FIGS/data/v2/index.json :: fine-tuning top N=10000: R30, R16_Q1_MASS :: sha256 c6c4138ee5c86423
\newcommand{\MlstFtTopRthirtyCoarseMassOverScratchEFour}{1.079 and 1.062}  % experiments/FIGS/data/v2/index.json :: fine-tuning top N=10000: R30 of each R16_Q1_MASS run over the from-scratch mean, SD over runs and standard error of that mean in quadrature :: sha256 c6c4138ee5c86423
\newcommand{\MlstFtTopRthirtyCoarseMassOverScratchEFourPr}{at least 6.2\% more rejection}  % experiments/FIGS/data/v2/index.json :: fine-tuning top N=10000: R30 of each R16_Q1_MASS run over the from-scratch mean, SD over runs and standard error of that mean in quadrature :: sha256 c6c4138ee5c86423 :: prose twin: two significant figures, direction in words
\newcommand{\MlstFtTopRthirtyCoarseMassOverScratchEFourPrMax}{up to 7.9\% more rejection}  % experiments/FIGS/data/v2/index.json :: fine-tuning top N=10000: R30 of each R16_Q1_MASS run over the from-scratch mean, SD over runs and standard error of that mean in quadrature :: sha256 c6c4138ee5c86423 :: prose twin of the run farther from 1
\newcommand{\MlstFtTopRthirtyResOverFullEFour}{1.02\,\ensuremath{\pm}\,0.12}  % experiments/FIGS/data/v2/index.json :: fine-tuning top N=10000: R30 of L162 over L188, run by run :: sha256 c6c4138ee5c86423
\newcommand{\MlstFtTopRthirtyResOverFullEFourPr}{equal within 20\%}  % experiments/FIGS/data/v2/index.json :: fine-tuning top N=10000: R30 of L162 over L188, run by run :: sha256 c6c4138ee5c86423 :: prose twin: two significant figures, direction in words
\newcommand{\MlstFtTopRthirtyMidOverFullEFour}{1.08\,\ensuremath{\pm}\,0.08}  % experiments/FIGS/data/v2/index.json :: fine-tuning top N=10000: R30 of R42_Q1 over L188, run by run :: sha256 c6c4138ee5c86423
\newcommand{\MlstFtTopRthirtyMidOverFullEFourPr}{equal within 9\%}  % experiments/FIGS/data/v2/index.json :: fine-tuning top N=10000: R30 of R42_Q1 over L188, run by run :: sha256 c6c4138ee5c86423 :: prose twin: two significant figures, direction in words
\newcommand{\MlstFtTopRthirtyCoarseOverFullEFour}{0.84\,\ensuremath{\pm}\,0.07}  % experiments/FIGS/data/v2/index.json :: fine-tuning top N=10000: R30 of R16_Q1 over L188, run by run :: sha256 c6c4138ee5c86423
\newcommand{\MlstFtTopRthirtyCoarseOverFullEFourPr}{16\% less rejection}  % experiments/FIGS/data/v2/index.json :: fine-tuning top N=10000: R30 of R16_Q1 over L188, run by run :: sha256 c6c4138ee5c86423 :: prose twin: two significant figures, direction in words
\newcommand{\MlstFtTopRthirtyResMassOverResEFour}{1.00 and 0.76}  % experiments/FIGS/data/v2/index.json :: fine-tuning top N=10000: R30 of L162_MASS over L162, run by run :: sha256 c6c4138ee5c86423
\newcommand{\MlstFtTopRthirtyResMassOverResEFourPr}{equal within 30\%}  % experiments/FIGS/data/v2/index.json :: fine-tuning top N=10000: R30 of L162_MASS over L162, run by run :: sha256 c6c4138ee5c86423 :: prose twin: two significant figures, direction in words
\newcommand{\MlstFtTopRthirtyResMassOverResEFourPrMax}{equal within 30\%}  % experiments/FIGS/data/v2/index.json :: fine-tuning top N=10000: R30 of L162_MASS over L162, run by run :: sha256 c6c4138ee5c86423 :: prose twin of the run farther from 1
\newcommand{\MlstFtTopRthirtyCoarseMassOverCoarseEFour}{1.13 and 0.92}  % experiments/FIGS/data/v2/index.json :: fine-tuning top N=10000: R30 of R16_Q1_MASS over R16_Q1, run by run :: sha256 c6c4138ee5c86423
\newcommand{\MlstFtTopRthirtyCoarseMassOverCoarseEFourPr}{equal within 20\%}  % experiments/FIGS/data/v2/index.json :: fine-tuning top N=10000: R30 of R16_Q1_MASS over R16_Q1, run by run :: sha256 c6c4138ee5c86423 :: prose twin: two significant figures, direction in words
\newcommand{\MlstFtTopRthirtyCoarseMassOverCoarseEFourPrMax}{equal within 20\%}  % experiments/FIGS/data/v2/index.json :: fine-tuning top N=10000: R30 of R16_Q1_MASS over R16_Q1, run by run :: sha256 c6c4138ee5c86423 :: prose twin of the run farther from 1
\newcommand{\MlstFtTopRthirtyRandOverFullEFour}{0.83\,\ensuremath{\pm}\,0.10}  % experiments/FIGS/data/v2/index.json :: fine-tuning top N=10000: R30 of each random partition over the L188 run of its seed :: sha256 c6c4138ee5c86423
\newcommand{\MlstFtTopRthirtyRandOverFullEFourPr}{equal within 20\%}  % experiments/FIGS/data/v2/index.json :: fine-tuning top N=10000: R30 of each random partition over the L188 run of its seed :: sha256 c6c4138ee5c86423 :: prose twin: two significant figures, direction in words
\newcommand{\MlstFtTopRthirtyRandOverCoarseEFour}{0.97\,\ensuremath{\pm}\,0.15}  % experiments/FIGS/data/v2/index.json :: fine-tuning top N=10000: R30 of each random partition over the R16_Q1 run of its seed :: sha256 c6c4138ee5c86423
\newcommand{\MlstFtTopRthirtyRandOverCoarseEFourPr}{equal within 20\%}  % experiments/FIGS/data/v2/index.json :: fine-tuning top N=10000: R30 of each random partition over the R16_Q1 run of its seed :: sha256 c6c4138ee5c86423 :: prose twin: two significant figures, direction in words
\newcommand{\MlstFtTopRthirtyScratchEFive}{560\,\ensuremath{\pm}\,90}  % experiments/FIGS/data/v2/index.json :: fine-tuning top N=100000: R30 from scratch, every seed :: sha256 c6c4138ee5c86423
\newcommand{\MlstFtTopRthirtyFullEFive}{850\,\ensuremath{\pm}\,90}  % experiments/FIGS/data/v2/index.json :: fine-tuning top N=100000: R30, L188 :: sha256 c6c4138ee5c86423
\newcommand{\MlstFtTopRthirtyFullOverScratchEFive}{1.50\,\ensuremath{\pm}\,0.21}  % experiments/FIGS/data/v2/index.json :: fine-tuning top N=100000: R30 of each L188 run over the from-scratch mean, SD over runs and standard error of that mean in quadrature :: sha256 c6c4138ee5c86423
\newcommand{\MlstFtTopRthirtyFullOverScratchEFivePr}{50\% more rejection}  % experiments/FIGS/data/v2/index.json :: fine-tuning top N=100000: R30 of each L188 run over the from-scratch mean, SD over runs and standard error of that mean in quadrature :: sha256 c6c4138ee5c86423 :: prose twin: two significant figures, direction in words
\newcommand{\MlstFtTopRthirtyResEFive}{810\,\ensuremath{\pm}\,60}  % experiments/FIGS/data/v2/index.json :: fine-tuning top N=100000: R30, L162 :: sha256 c6c4138ee5c86423
\newcommand{\MlstFtTopRthirtyResOverScratchEFive}{1.44\,\ensuremath{\pm}\,0.18}  % experiments/FIGS/data/v2/index.json :: fine-tuning top N=100000: R30 of each L162 run over the from-scratch mean, SD over runs and standard error of that mean in quadrature :: sha256 c6c4138ee5c86423
\newcommand{\MlstFtTopRthirtyResOverScratchEFivePr}{44\% more rejection}  % experiments/FIGS/data/v2/index.json :: fine-tuning top N=100000: R30 of each L162 run over the from-scratch mean, SD over runs and standard error of that mean in quadrature :: sha256 c6c4138ee5c86423 :: prose twin: two significant figures, direction in words
\newcommand{\MlstFtTopRthirtyMidEFive}{850\,\ensuremath{\pm}\,80}  % experiments/FIGS/data/v2/index.json :: fine-tuning top N=100000: R30, R42_Q1 :: sha256 c6c4138ee5c86423
\newcommand{\MlstFtTopRthirtyMidOverScratchEFive}{1.51\,\ensuremath{\pm}\,0.20}  % experiments/FIGS/data/v2/index.json :: fine-tuning top N=100000: R30 of each R42_Q1 run over the from-scratch mean, SD over runs and standard error of that mean in quadrature :: sha256 c6c4138ee5c86423
\newcommand{\MlstFtTopRthirtyMidOverScratchEFivePr}{51\% more rejection}  % experiments/FIGS/data/v2/index.json :: fine-tuning top N=100000: R30 of each R42_Q1 run over the from-scratch mean, SD over runs and standard error of that mean in quadrature :: sha256 c6c4138ee5c86423 :: prose twin: two significant figures, direction in words
\newcommand{\MlstFtTopRthirtyCoarseEFive}{830\,\ensuremath{\pm}\,60}  % experiments/FIGS/data/v2/index.json :: fine-tuning top N=100000: R30, R16_Q1 :: sha256 c6c4138ee5c86423
\newcommand{\MlstFtTopRthirtyCoarseOverScratchEFive}{1.47\,\ensuremath{\pm}\,0.17}  % experiments/FIGS/data/v2/index.json :: fine-tuning top N=100000: R30 of each R16_Q1 run over the from-scratch mean, SD over runs and standard error of that mean in quadrature :: sha256 c6c4138ee5c86423
\newcommand{\MlstFtTopRthirtyCoarseOverScratchEFivePr}{47\% more rejection}  % experiments/FIGS/data/v2/index.json :: fine-tuning top N=100000: R30 of each R16_Q1 run over the from-scratch mean, SD over runs and standard error of that mean in quadrature :: sha256 c6c4138ee5c86423 :: prose twin: two significant figures, direction in words
\newcommand{\MlstFtTopRthirtyResMassEFive}{848 and 848}  % experiments/FIGS/data/v2/index.json :: fine-tuning top N=100000: R30, L162_MASS :: sha256 c6c4138ee5c86423
\newcommand{\MlstFtTopRthirtyResMassOverScratchEFive}{1.50 and 1.50}  % experiments/FIGS/data/v2/index.json :: fine-tuning top N=100000: R30 of each L162_MASS run over the from-scratch mean, SD over runs and standard error of that mean in quadrature :: sha256 c6c4138ee5c86423
\newcommand{\MlstFtTopRthirtyResMassOverScratchEFivePr}{at least 50\% more rejection}  % experiments/FIGS/data/v2/index.json :: fine-tuning top N=100000: R30 of each L162_MASS run over the from-scratch mean, SD over runs and standard error of that mean in quadrature :: sha256 c6c4138ee5c86423 :: prose twin: two significant figures, direction in words
\newcommand{\MlstFtTopRthirtyResMassOverScratchEFivePrMax}{up to 50\% more rejection}  % experiments/FIGS/data/v2/index.json :: fine-tuning top N=100000: R30 of each L162_MASS run over the from-scratch mean, SD over runs and standard error of that mean in quadrature :: sha256 c6c4138ee5c86423 :: prose twin of the run farther from 1
\newcommand{\MlstFtTopRthirtyCoarseMassEFive}{856 and 774}  % experiments/FIGS/data/v2/index.json :: fine-tuning top N=100000: R30, R16_Q1_MASS :: sha256 c6c4138ee5c86423
\newcommand{\MlstFtTopRthirtyCoarseMassOverScratchEFive}{1.51 and 1.37}  % experiments/FIGS/data/v2/index.json :: fine-tuning top N=100000: R30 of each R16_Q1_MASS run over the from-scratch mean, SD over runs and standard error of that mean in quadrature :: sha256 c6c4138ee5c86423
\newcommand{\MlstFtTopRthirtyCoarseMassOverScratchEFivePr}{at least 37\% more rejection}  % experiments/FIGS/data/v2/index.json :: fine-tuning top N=100000: R30 of each R16_Q1_MASS run over the from-scratch mean, SD over runs and standard error of that mean in quadrature :: sha256 c6c4138ee5c86423 :: prose twin: two significant figures, direction in words
\newcommand{\MlstFtTopRthirtyCoarseMassOverScratchEFivePrMax}{up to 51\% more rejection}  % experiments/FIGS/data/v2/index.json :: fine-tuning top N=100000: R30 of each R16_Q1_MASS run over the from-scratch mean, SD over runs and standard error of that mean in quadrature :: sha256 c6c4138ee5c86423 :: prose twin of the run farther from 1
\newcommand{\MlstFtTopRthirtyResOverFullEFive}{0.96\,\ensuremath{\pm}\,0.05}  % experiments/FIGS/data/v2/index.json :: fine-tuning top N=100000: R30 of L162 over L188, run by run :: sha256 c6c4138ee5c86423
\newcommand{\MlstFtTopRthirtyResOverFullEFivePr}{equal within 6\%}  % experiments/FIGS/data/v2/index.json :: fine-tuning top N=100000: R30 of L162 over L188, run by run :: sha256 c6c4138ee5c86423 :: prose twin: two significant figures, direction in words
\newcommand{\MlstFtTopRthirtyMidOverFullEFive}{1.02\,\ensuremath{\pm}\,0.18}  % experiments/FIGS/data/v2/index.json :: fine-tuning top N=100000: R30 of R42_Q1 over L188, run by run :: sha256 c6c4138ee5c86423
\newcommand{\MlstFtTopRthirtyMidOverFullEFivePr}{equal within 20\%}  % experiments/FIGS/data/v2/index.json :: fine-tuning top N=100000: R30 of R42_Q1 over L188, run by run :: sha256 c6c4138ee5c86423 :: prose twin: two significant figures, direction in words
\newcommand{\MlstFtTopRthirtyCoarseOverFullEFive}{0.99\,\ensuremath{\pm}\,0.17}  % experiments/FIGS/data/v2/index.json :: fine-tuning top N=100000: R30 of R16_Q1 over L188, run by run :: sha256 c6c4138ee5c86423
\newcommand{\MlstFtTopRthirtyCoarseOverFullEFivePr}{equal within 20\%}  % experiments/FIGS/data/v2/index.json :: fine-tuning top N=100000: R30 of R16_Q1 over L188, run by run :: sha256 c6c4138ee5c86423 :: prose twin: two significant figures, direction in words
\newcommand{\MlstFtTopRthirtyResMassOverResEFive}{1.00 and 1.15}  % experiments/FIGS/data/v2/index.json :: fine-tuning top N=100000: R30 of L162_MASS over L162, run by run :: sha256 c6c4138ee5c86423
\newcommand{\MlstFtTopRthirtyResMassOverResEFivePr}{equal within 20\%}  % experiments/FIGS/data/v2/index.json :: fine-tuning top N=100000: R30 of L162_MASS over L162, run by run :: sha256 c6c4138ee5c86423 :: prose twin: two significant figures, direction in words
\newcommand{\MlstFtTopRthirtyResMassOverResEFivePrMax}{equal within 20\%}  % experiments/FIGS/data/v2/index.json :: fine-tuning top N=100000: R30 of L162_MASS over L162, run by run :: sha256 c6c4138ee5c86423 :: prose twin of the run farther from 1
\newcommand{\MlstFtTopRthirtyCoarseMassOverCoarseEFive}{1.08 and 0.87}  % experiments/FIGS/data/v2/index.json :: fine-tuning top N=100000: R30 of R16_Q1_MASS over R16_Q1, run by run :: sha256 c6c4138ee5c86423
\newcommand{\MlstFtTopRthirtyCoarseMassOverCoarseEFivePr}{equal within 20\%}  % experiments/FIGS/data/v2/index.json :: fine-tuning top N=100000: R30 of R16_Q1_MASS over R16_Q1, run by run :: sha256 c6c4138ee5c86423 :: prose twin: two significant figures, direction in words
\newcommand{\MlstFtTopRthirtyCoarseMassOverCoarseEFivePrMax}{equal within 20\%}  % experiments/FIGS/data/v2/index.json :: fine-tuning top N=100000: R30 of R16_Q1_MASS over R16_Q1, run by run :: sha256 c6c4138ee5c86423 :: prose twin of the run farther from 1
\newcommand{\MlstFtTopRthirtyRandOverFullEFive}{1.01\,\ensuremath{\pm}\,0.22}  % experiments/FIGS/data/v2/index.json :: fine-tuning top N=100000: R30 of each random partition over the L188 run of its seed :: sha256 c6c4138ee5c86423
\newcommand{\MlstFtTopRthirtyRandOverFullEFivePr}{equal within 30\%}  % experiments/FIGS/data/v2/index.json :: fine-tuning top N=100000: R30 of each random partition over the L188 run of its seed :: sha256 c6c4138ee5c86423 :: prose twin: two significant figures, direction in words
\newcommand{\MlstFtTopRthirtyRandOverCoarseEFive}{0.97\,\ensuremath{\pm}\,0.10}  % experiments/FIGS/data/v2/index.json :: fine-tuning top N=100000: R30 of each random partition over the R16_Q1 run of its seed :: sha256 c6c4138ee5c86423
\newcommand{\MlstFtTopRthirtyRandOverCoarseEFivePr}{equal within 10\%}  % experiments/FIGS/data/v2/index.json :: fine-tuning top N=100000: R30 of each random partition over the R16_Q1 run of its seed :: sha256 c6c4138ee5c86423 :: prose twin: two significant figures, direction in words
\newcommand{\MlstFtTopRthirtyScratch}{1{,}412\,\ensuremath{\pm}\,20}  % experiments/FIGS/data/v2/index.json :: fine-tuning top N=1200000: R30 from scratch, every seed :: sha256 c6c4138ee5c86423
\newcommand{\MlstFtTopRthirtyFull}{1{,}800\,\ensuremath{\pm}\,60}  % experiments/FIGS/data/v2/index.json :: fine-tuning top N=1200000: R30, L188 :: sha256 c6c4138ee5c86423
\newcommand{\MlstFtTopRthirtyFullOverScratch}{1.27\,\ensuremath{\pm}\,0.04}  % experiments/FIGS/data/v2/index.json :: fine-tuning top N=1200000: R30 of each L188 run over the from-scratch mean, SD over runs and standard error of that mean in quadrature :: sha256 c6c4138ee5c86423
\newcommand{\MlstFtTopRthirtyFullOverScratchPr}{27\% more rejection}  % experiments/FIGS/data/v2/index.json :: fine-tuning top N=1200000: R30 of each L188 run over the from-scratch mean, SD over runs and standard error of that mean in quadrature :: sha256 c6c4138ee5c86423 :: prose twin: two significant figures, direction in words
\newcommand{\MlstFtTopRthirtyRes}{1{,}790\,\ensuremath{\pm}\,60}  % experiments/FIGS/data/v2/index.json :: fine-tuning top N=1200000: R30, L162 :: sha256 c6c4138ee5c86423
\newcommand{\MlstFtTopRthirtyResOverScratch}{1.27\,\ensuremath{\pm}\,0.04}  % experiments/FIGS/data/v2/index.json :: fine-tuning top N=1200000: R30 of each L162 run over the from-scratch mean, SD over runs and standard error of that mean in quadrature :: sha256 c6c4138ee5c86423
\newcommand{\MlstFtTopRthirtyResOverScratchPr}{27\% more rejection}  % experiments/FIGS/data/v2/index.json :: fine-tuning top N=1200000: R30 of each L162 run over the from-scratch mean, SD over runs and standard error of that mean in quadrature :: sha256 c6c4138ee5c86423 :: prose twin: two significant figures, direction in words
\newcommand{\MlstFtTopRthirtyMid}{1{,}780\,\ensuremath{\pm}\,150}  % experiments/FIGS/data/v2/index.json :: fine-tuning top N=1200000: R30, R42_Q1 :: sha256 c6c4138ee5c86423
\newcommand{\MlstFtTopRthirtyMidOverScratch}{1.26\,\ensuremath{\pm}\,0.11}  % experiments/FIGS/data/v2/index.json :: fine-tuning top N=1200000: R30 of each R42_Q1 run over the from-scratch mean, SD over runs and standard error of that mean in quadrature :: sha256 c6c4138ee5c86423
\newcommand{\MlstFtTopRthirtyMidOverScratchPr}{26\% more rejection}  % experiments/FIGS/data/v2/index.json :: fine-tuning top N=1200000: R30 of each R42_Q1 run over the from-scratch mean, SD over runs and standard error of that mean in quadrature :: sha256 c6c4138ee5c86423 :: prose twin: two significant figures, direction in words
\newcommand{\MlstFtTopRthirtyCoarse}{1{,}740\,\ensuremath{\pm}\,140}  % experiments/FIGS/data/v2/index.json :: fine-tuning top N=1200000: R30, R16_Q1 :: sha256 c6c4138ee5c86423
\newcommand{\MlstFtTopRthirtyCoarseOverScratch}{1.23\,\ensuremath{\pm}\,0.10}  % experiments/FIGS/data/v2/index.json :: fine-tuning top N=1200000: R30 of each R16_Q1 run over the from-scratch mean, SD over runs and standard error of that mean in quadrature :: sha256 c6c4138ee5c86423
\newcommand{\MlstFtTopRthirtyCoarseOverScratchPr}{23\% more rejection}  % experiments/FIGS/data/v2/index.json :: fine-tuning top N=1200000: R30 of each R16_Q1 run over the from-scratch mean, SD over runs and standard error of that mean in quadrature :: sha256 c6c4138ee5c86423 :: prose twin: two significant figures, direction in words
\newcommand{\MlstFtTopRthirtyResMass}{1{,}941 and 1{,}590}  % experiments/FIGS/data/v2/index.json :: fine-tuning top N=1200000: R30, L162_MASS :: sha256 c6c4138ee5c86423
\newcommand{\MlstFtTopRthirtyResMassOverScratch}{1.37 and 1.13}  % experiments/FIGS/data/v2/index.json :: fine-tuning top N=1200000: R30 of each L162_MASS run over the from-scratch mean, SD over runs and standard error of that mean in quadrature :: sha256 c6c4138ee5c86423
\newcommand{\MlstFtTopRthirtyResMassOverScratchPr}{at least 13\% more rejection}  % experiments/FIGS/data/v2/index.json :: fine-tuning top N=1200000: R30 of each L162_MASS run over the from-scratch mean, SD over runs and standard error of that mean in quadrature :: sha256 c6c4138ee5c86423 :: prose twin: two significant figures, direction in words
\newcommand{\MlstFtTopRthirtyResMassOverScratchPrMax}{up to 37\% more rejection}  % experiments/FIGS/data/v2/index.json :: fine-tuning top N=1200000: R30 of each L162_MASS run over the from-scratch mean, SD over runs and standard error of that mean in quadrature :: sha256 c6c4138ee5c86423 :: prose twin of the run farther from 1
\newcommand{\MlstFtTopRthirtyCoarseMass}{1{,}683 and 1{,}628}  % experiments/FIGS/data/v2/index.json :: fine-tuning top N=1200000: R30, R16_Q1_MASS :: sha256 c6c4138ee5c86423
\newcommand{\MlstFtTopRthirtyCoarseMassOverScratch}{1.192 and 1.153}  % experiments/FIGS/data/v2/index.json :: fine-tuning top N=1200000: R30 of each R16_Q1_MASS run over the from-scratch mean, SD over runs and standard error of that mean in quadrature :: sha256 c6c4138ee5c86423
\newcommand{\MlstFtTopRthirtyCoarseMassOverScratchPr}{at least 15\% more rejection}  % experiments/FIGS/data/v2/index.json :: fine-tuning top N=1200000: R30 of each R16_Q1_MASS run over the from-scratch mean, SD over runs and standard error of that mean in quadrature :: sha256 c6c4138ee5c86423 :: prose twin: two significant figures, direction in words
\newcommand{\MlstFtTopRthirtyCoarseMassOverScratchPrMax}{up to 19\% more rejection}  % experiments/FIGS/data/v2/index.json :: fine-tuning top N=1200000: R30 of each R16_Q1_MASS run over the from-scratch mean, SD over runs and standard error of that mean in quadrature :: sha256 c6c4138ee5c86423 :: prose twin of the run farther from 1
\newcommand{\MlstFtTopRthirtyResOverFull}{0.997\,\ensuremath{\pm}\,0.032}  % experiments/FIGS/data/v2/index.json :: fine-tuning top N=1200000: R30 of L162 over L188, run by run :: sha256 c6c4138ee5c86423
\newcommand{\MlstFtTopRthirtyResOverFullPr}{equal within 4\%}  % experiments/FIGS/data/v2/index.json :: fine-tuning top N=1200000: R30 of L162 over L188, run by run :: sha256 c6c4138ee5c86423 :: prose twin: two significant figures, direction in words
\newcommand{\MlstFtTopRthirtyMidOverFull}{0.99\,\ensuremath{\pm}\,0.07}  % experiments/FIGS/data/v2/index.json :: fine-tuning top N=1200000: R30 of R42_Q1 over L188, run by run :: sha256 c6c4138ee5c86423
\newcommand{\MlstFtTopRthirtyMidOverFullPr}{equal within 8\%}  % experiments/FIGS/data/v2/index.json :: fine-tuning top N=1200000: R30 of R42_Q1 over L188, run by run :: sha256 c6c4138ee5c86423 :: prose twin: two significant figures, direction in words
\newcommand{\MlstFtTopRthirtyCoarseOverFull}{0.97\,\ensuremath{\pm}\,0.10}  % experiments/FIGS/data/v2/index.json :: fine-tuning top N=1200000: R30 of R16_Q1 over L188, run by run :: sha256 c6c4138ee5c86423
\newcommand{\MlstFtTopRthirtyCoarseOverFullPr}{equal within 10\%}  % experiments/FIGS/data/v2/index.json :: fine-tuning top N=1200000: R30 of R16_Q1 over L188, run by run :: sha256 c6c4138ee5c86423 :: prose twin: two significant figures, direction in words
\newcommand{\MlstFtTopRthirtyResMassOverRes}{1.12 and 0.89}  % experiments/FIGS/data/v2/index.json :: fine-tuning top N=1200000: R30 of L162_MASS over L162, run by run :: sha256 c6c4138ee5c86423
\newcommand{\MlstFtTopRthirtyResMassOverResPr}{equal within 20\%}  % experiments/FIGS/data/v2/index.json :: fine-tuning top N=1200000: R30 of L162_MASS over L162, run by run :: sha256 c6c4138ee5c86423 :: prose twin: two significant figures, direction in words
\newcommand{\MlstFtTopRthirtyResMassOverResPrMax}{equal within 20\%}  % experiments/FIGS/data/v2/index.json :: fine-tuning top N=1200000: R30 of L162_MASS over L162, run by run :: sha256 c6c4138ee5c86423 :: prose twin of the run farther from 1
\newcommand{\MlstFtTopRthirtyCoarseMassOverCoarse}{0.900 and 0.919}  % experiments/FIGS/data/v2/index.json :: fine-tuning top N=1200000: R30 of R16_Q1_MASS over R16_Q1, run by run :: sha256 c6c4138ee5c86423
\newcommand{\MlstFtTopRthirtyCoarseMassOverCoarsePr}{equal within 10\%}  % experiments/FIGS/data/v2/index.json :: fine-tuning top N=1200000: R30 of R16_Q1_MASS over R16_Q1, run by run :: sha256 c6c4138ee5c86423 :: prose twin: two significant figures, direction in words
\newcommand{\MlstFtTopRthirtyCoarseMassOverCoarsePrMax}{equal within 10\%}  % experiments/FIGS/data/v2/index.json :: fine-tuning top N=1200000: R30 of R16_Q1_MASS over R16_Q1, run by run :: sha256 c6c4138ee5c86423 :: prose twin of the run farther from 1
\newcommand{\MlstFtTopRthirtyRandOverFull}{0.97\,\ensuremath{\pm}\,0.04}  % experiments/FIGS/data/v2/index.json :: fine-tuning top N=1200000: R30 of each random partition over the L188 run of its seed :: sha256 c6c4138ee5c86423
\newcommand{\MlstFtTopRthirtyRandOverFullPr}{equal within 4\%}  % experiments/FIGS/data/v2/index.json :: fine-tuning top N=1200000: R30 of each random partition over the L188 run of its seed :: sha256 c6c4138ee5c86423 :: prose twin: two significant figures, direction in words
\newcommand{\MlstFtTopRthirtyRandOverCoarse}{0.97\,\ensuremath{\pm}\,0.09}  % experiments/FIGS/data/v2/index.json :: fine-tuning top N=1200000: R30 of each random partition over the R16_Q1 run of its seed :: sha256 c6c4138ee5c86423
\newcommand{\MlstFtTopRthirtyRandOverCoarsePr}{equal within 9\%}  % experiments/FIGS/data/v2/index.json :: fine-tuning top N=1200000: R30 of each random partition over the R16_Q1 run of its seed :: sha256 c6c4138ee5c86423 :: prose twin: two significant figures, direction in words
\newcommand{\MlstFtTopRthirtyFullOverPart}{0.650\,\ensuremath{\pm}\,0.020}  % experiments/FIGS/data/v2/index.json :: fine-tuning top N=1200000: R30 of each 188-class run over ParT-f.t.'s published value :: sha256 c6c4138ee5c86423
\newcommand{\MlstFtTopRthirtyFullOverPartPr}{35\% less rejection}  % experiments/FIGS/data/v2/index.json :: fine-tuning top N=1200000: R30 of each 188-class run over ParT-f.t.'s published value :: sha256 c6c4138ee5c86423 :: prose twin: two significant figures, direction in words
\newcommand{\MlstFtTopAucFull}{0.9861\,\ensuremath{\pm}\,0.0001}  % experiments/FIGS/data/v2/index.json :: fine-tuning top N=1200000: AUC of the 188-class runs :: sha256 c6c4138ee5c86423
\newcommand{\MlstFtTopAucScratch}{0.9852\,\ensuremath{\pm}\,0.0001}  % experiments/FIGS/data/v2/index.json :: fine-tuning top N=1200000: AUC from scratch :: sha256 c6c4138ee5c86423
\newcommand{\MlstFtTopRthirtyFullRelSd}{3\%}  % experiments/FIGS/data/v2/index.json :: fine-tuning top, full set: relative SD of R30 over the 188-class runs :: sha256 c6c4138ee5c86423
\newcommand{\MlstFtTopRthirtyKin}{1{,}800}  % experiments/FIGS/data/v2/index.json :: fine-tuning top, full set: R30 of the kinematic-only model (one run) :: sha256 c6c4138ee5c86423
\newcommand{\MlstFtTopRthirtyKinOverFull}{1.00}  % experiments/FIGS/data/v2/index.json :: fine-tuning top, full set: R30 of the kinematic-only model over the mean of the 188-class runs :: sha256 c6c4138ee5c86423
\newcommand{\MlstFtTopRthirtyKinOverFullPr}{0.23\% more rejection}  % experiments/FIGS/data/v2/index.json :: fine-tuning top, full set: R30 of the kinematic-only model over the mean of the 188-class runs :: sha256 c6c4138ee5c86423 :: prose twin: two significant figures, direction in words
\newcommand{\MlstFtTopDataEquivJets}{1{,}200}  % experiments/FIGS/data/v2/index.json :: fine-tuning top: from-scratch training jets whose R30 equals the 188-class mean at 1000 jets (interpolated linearly in R30 against log10 N) :: sha256 c6c4138ee5c86423
\newcommand{\MlstFtTopDataEquivJetsEFour}{23{,}000}  % experiments/FIGS/data/v2/index.json :: fine-tuning top: from-scratch training jets whose R30 equals the 188-class mean at 10000 jets (interpolated linearly in R30 against log10 N) :: sha256 c6c4138ee5c86423
\newcommand{\MlstFtQgRthirtyScratchEThree}{49.7\,\ensuremath{\pm}\,2.1}  % experiments/FIGS/data/v2/index.json :: fine-tuning qg N=1000: R30 from scratch, every seed :: sha256 c6c4138ee5c86423
\newcommand{\MlstFtQgRthirtyFullEThree}{47\,\ensuremath{\pm}\,10}  % experiments/FIGS/data/v2/index.json :: fine-tuning qg N=1000: R30, L188 :: sha256 c6c4138ee5c86423
\newcommand{\MlstFtQgRthirtyFullOverScratchEThree}{0.95\,\ensuremath{\pm}\,0.19}  % experiments/FIGS/data/v2/index.json :: fine-tuning qg N=1000: R30 of each L188 run over the from-scratch mean, SD over runs and standard error of that mean in quadrature :: sha256 c6c4138ee5c86423
\newcommand{\MlstFtQgRthirtyFullOverScratchEThreePr}{equal within 20\%}  % experiments/FIGS/data/v2/index.json :: fine-tuning qg N=1000: R30 of each L188 run over the from-scratch mean, SD over runs and standard error of that mean in quadrature :: sha256 c6c4138ee5c86423 :: prose twin: two significant figures, direction in words
\newcommand{\MlstFtQgRthirtyResEThree}{52.4\,\ensuremath{\pm}\,3.1}  % experiments/FIGS/data/v2/index.json :: fine-tuning qg N=1000: R30, L162 :: sha256 c6c4138ee5c86423
\newcommand{\MlstFtQgRthirtyResOverScratchEThree}{1.06\,\ensuremath{\pm}\,0.07}  % experiments/FIGS/data/v2/index.json :: fine-tuning qg N=1000: R30 of each L162 run over the from-scratch mean, SD over runs and standard error of that mean in quadrature :: sha256 c6c4138ee5c86423
\newcommand{\MlstFtQgRthirtyResOverScratchEThreePr}{equal within 7\%}  % experiments/FIGS/data/v2/index.json :: fine-tuning qg N=1000: R30 of each L162 run over the from-scratch mean, SD over runs and standard error of that mean in quadrature :: sha256 c6c4138ee5c86423 :: prose twin: two significant figures, direction in words
\newcommand{\MlstFtQgRthirtyMidEThree}{54\,\ensuremath{\pm}\,6}  % experiments/FIGS/data/v2/index.json :: fine-tuning qg N=1000: R30, R42_Q1 :: sha256 c6c4138ee5c86423
\newcommand{\MlstFtQgRthirtyMidOverScratchEThree}{1.09\,\ensuremath{\pm}\,0.12}  % experiments/FIGS/data/v2/index.json :: fine-tuning qg N=1000: R30 of each R42_Q1 run over the from-scratch mean, SD over runs and standard error of that mean in quadrature :: sha256 c6c4138ee5c86423
\newcommand{\MlstFtQgRthirtyMidOverScratchEThreePr}{equal within 20\%}  % experiments/FIGS/data/v2/index.json :: fine-tuning qg N=1000: R30 of each R42_Q1 run over the from-scratch mean, SD over runs and standard error of that mean in quadrature :: sha256 c6c4138ee5c86423 :: prose twin: two significant figures, direction in words
\newcommand{\MlstFtQgRthirtyCoarseEThree}{55\,\ensuremath{\pm}\,5}  % experiments/FIGS/data/v2/index.json :: fine-tuning qg N=1000: R30, R16_Q1 :: sha256 c6c4138ee5c86423
\newcommand{\MlstFtQgRthirtyCoarseOverScratchEThree}{1.12\,\ensuremath{\pm}\,0.10}  % experiments/FIGS/data/v2/index.json :: fine-tuning qg N=1000: R30 of each R16_Q1 run over the from-scratch mean, SD over runs and standard error of that mean in quadrature :: sha256 c6c4138ee5c86423
\newcommand{\MlstFtQgRthirtyCoarseOverScratchEThreePr}{equal within 20\%}  % experiments/FIGS/data/v2/index.json :: fine-tuning qg N=1000: R30 of each R16_Q1 run over the from-scratch mean, SD over runs and standard error of that mean in quadrature :: sha256 c6c4138ee5c86423 :: prose twin: two significant figures, direction in words
\newcommand{\MlstFtQgRthirtyResMassEThree}{59.3 and 57.1}  % experiments/FIGS/data/v2/index.json :: fine-tuning qg N=1000: R30, L162_MASS :: sha256 c6c4138ee5c86423
\newcommand{\MlstFtQgRthirtyResMassOverScratchEThree}{1.193 and 1.149}  % experiments/FIGS/data/v2/index.json :: fine-tuning qg N=1000: R30 of each L162_MASS run over the from-scratch mean, SD over runs and standard error of that mean in quadrature :: sha256 c6c4138ee5c86423
\newcommand{\MlstFtQgRthirtyResMassOverScratchEThreePr}{at least 15\% more rejection}  % experiments/FIGS/data/v2/index.json :: fine-tuning qg N=1000: R30 of each L162_MASS run over the from-scratch mean, SD over runs and standard error of that mean in quadrature :: sha256 c6c4138ee5c86423 :: prose twin: two significant figures, direction in words
\newcommand{\MlstFtQgRthirtyResMassOverScratchEThreePrMax}{up to 19\% more rejection}  % experiments/FIGS/data/v2/index.json :: fine-tuning qg N=1000: R30 of each L162_MASS run over the from-scratch mean, SD over runs and standard error of that mean in quadrature :: sha256 c6c4138ee5c86423 :: prose twin of the run farther from 1
\newcommand{\MlstFtQgRthirtyCoarseMassEThree}{54.2 and 57.7}  % experiments/FIGS/data/v2/index.json :: fine-tuning qg N=1000: R30, R16_Q1_MASS :: sha256 c6c4138ee5c86423
\newcommand{\MlstFtQgRthirtyCoarseMassOverScratchEThree}{1.091 and 1.161}  % experiments/FIGS/data/v2/index.json :: fine-tuning qg N=1000: R30 of each R16_Q1_MASS run over the from-scratch mean, SD over runs and standard error of that mean in quadrature :: sha256 c6c4138ee5c86423
\newcommand{\MlstFtQgRthirtyCoarseMassOverScratchEThreePr}{at least 9.1\% more rejection}  % experiments/FIGS/data/v2/index.json :: fine-tuning qg N=1000: R30 of each R16_Q1_MASS run over the from-scratch mean, SD over runs and standard error of that mean in quadrature :: sha256 c6c4138ee5c86423 :: prose twin: two significant figures, direction in words
\newcommand{\MlstFtQgRthirtyCoarseMassOverScratchEThreePrMax}{up to 16\% more rejection}  % experiments/FIGS/data/v2/index.json :: fine-tuning qg N=1000: R30 of each R16_Q1_MASS run over the from-scratch mean, SD over runs and standard error of that mean in quadrature :: sha256 c6c4138ee5c86423 :: prose twin of the run farther from 1
\newcommand{\MlstFtQgRthirtyResOverFullEThree}{1.14\,\ensuremath{\pm}\,0.17}  % experiments/FIGS/data/v2/index.json :: fine-tuning qg N=1000: R30 of L162 over L188, run by run :: sha256 c6c4138ee5c86423
\newcommand{\MlstFtQgRthirtyResOverFullEThreePr}{equal within 20\%}  % experiments/FIGS/data/v2/index.json :: fine-tuning qg N=1000: R30 of L162 over L188, run by run :: sha256 c6c4138ee5c86423 :: prose twin: two significant figures, direction in words
\newcommand{\MlstFtQgRthirtyMidOverFullEThree}{1.17\,\ensuremath{\pm}\,0.14}  % experiments/FIGS/data/v2/index.json :: fine-tuning qg N=1000: R30 of R42_Q1 over L188, run by run :: sha256 c6c4138ee5c86423
\newcommand{\MlstFtQgRthirtyMidOverFullEThreePr}{equal within 20\%}  % experiments/FIGS/data/v2/index.json :: fine-tuning qg N=1000: R30 of R42_Q1 over L188, run by run :: sha256 c6c4138ee5c86423 :: prose twin: two significant figures, direction in words
\newcommand{\MlstFtQgRthirtyCoarseOverFullEThree}{1.21\,\ensuremath{\pm}\,0.27}  % experiments/FIGS/data/v2/index.json :: fine-tuning qg N=1000: R30 of R16_Q1 over L188, run by run :: sha256 c6c4138ee5c86423
\newcommand{\MlstFtQgRthirtyCoarseOverFullEThreePr}{equal within 30\%}  % experiments/FIGS/data/v2/index.json :: fine-tuning qg N=1000: R30 of R16_Q1 over L188, run by run :: sha256 c6c4138ee5c86423 :: prose twin: two significant figures, direction in words
\newcommand{\MlstFtQgRthirtyResMassOverResEThree}{1.108 and 1.042}  % experiments/FIGS/data/v2/index.json :: fine-tuning qg N=1000: R30 of L162_MASS over L162, run by run :: sha256 c6c4138ee5c86423
\newcommand{\MlstFtQgRthirtyResMassOverResEThreePr}{equal within 20\%}  % experiments/FIGS/data/v2/index.json :: fine-tuning qg N=1000: R30 of L162_MASS over L162, run by run :: sha256 c6c4138ee5c86423 :: prose twin: two significant figures, direction in words
\newcommand{\MlstFtQgRthirtyResMassOverResEThreePrMax}{equal within 20\%}  % experiments/FIGS/data/v2/index.json :: fine-tuning qg N=1000: R30 of L162_MASS over L162, run by run :: sha256 c6c4138ee5c86423 :: prose twin of the run farther from 1
\newcommand{\MlstFtQgRthirtyCoarseMassOverCoarseEThree}{0.89 and 1.11}  % experiments/FIGS/data/v2/index.json :: fine-tuning qg N=1000: R30 of R16_Q1_MASS over R16_Q1, run by run :: sha256 c6c4138ee5c86423
\newcommand{\MlstFtQgRthirtyCoarseMassOverCoarseEThreePr}{equal within 20\%}  % experiments/FIGS/data/v2/index.json :: fine-tuning qg N=1000: R30 of R16_Q1_MASS over R16_Q1, run by run :: sha256 c6c4138ee5c86423 :: prose twin: two significant figures, direction in words
\newcommand{\MlstFtQgRthirtyCoarseMassOverCoarseEThreePrMax}{equal within 20\%}  % experiments/FIGS/data/v2/index.json :: fine-tuning qg N=1000: R30 of R16_Q1_MASS over R16_Q1, run by run :: sha256 c6c4138ee5c86423 :: prose twin of the run farther from 1
\newcommand{\MlstFtQgRthirtyRandOverFullEThree}{0.91\,\ensuremath{\pm}\,0.07}  % experiments/FIGS/data/v2/index.json :: fine-tuning qg N=1000: R30 of each random partition over the L188 run of its seed :: sha256 c6c4138ee5c86423
\newcommand{\MlstFtQgRthirtyRandOverFullEThreePr}{equal within 10\%}  % experiments/FIGS/data/v2/index.json :: fine-tuning qg N=1000: R30 of each random partition over the L188 run of its seed :: sha256 c6c4138ee5c86423 :: prose twin: two significant figures, direction in words
\newcommand{\MlstFtQgRthirtyRandOverCoarseEThree}{0.87\,\ensuremath{\pm}\,0.14}  % experiments/FIGS/data/v2/index.json :: fine-tuning qg N=1000: R30 of each random partition over the R16_Q1 run of its seed :: sha256 c6c4138ee5c86423
\newcommand{\MlstFtQgRthirtyRandOverCoarseEThreePr}{equal within 20\%}  % experiments/FIGS/data/v2/index.json :: fine-tuning qg N=1000: R30 of each random partition over the R16_Q1 run of its seed :: sha256 c6c4138ee5c86423 :: prose twin: two significant figures, direction in words
\newcommand{\MlstFtQgRthirtyScratchEFour}{66\,\ensuremath{\pm}\,9}  % experiments/FIGS/data/v2/index.json :: fine-tuning qg N=10000: R30 from scratch, every seed :: sha256 c6c4138ee5c86423
\newcommand{\MlstFtQgRthirtyFullEFour}{75\,\ensuremath{\pm}\,7}  % experiments/FIGS/data/v2/index.json :: fine-tuning qg N=10000: R30, L188 :: sha256 c6c4138ee5c86423
\newcommand{\MlstFtQgRthirtyFullOverScratchEFour}{1.14\,\ensuremath{\pm}\,0.14}  % experiments/FIGS/data/v2/index.json :: fine-tuning qg N=10000: R30 of each L188 run over the from-scratch mean, SD over runs and standard error of that mean in quadrature :: sha256 c6c4138ee5c86423
\newcommand{\MlstFtQgRthirtyFullOverScratchEFourPr}{equal within 20\%}  % experiments/FIGS/data/v2/index.json :: fine-tuning qg N=10000: R30 of each L188 run over the from-scratch mean, SD over runs and standard error of that mean in quadrature :: sha256 c6c4138ee5c86423 :: prose twin: two significant figures, direction in words
\newcommand{\MlstFtQgRthirtyResEFour}{76.4\,\ensuremath{\pm}\,0.6}  % experiments/FIGS/data/v2/index.json :: fine-tuning qg N=10000: R30, L162 :: sha256 c6c4138ee5c86423
\newcommand{\MlstFtQgRthirtyResOverScratchEFour}{1.15\,\ensuremath{\pm}\,0.09}  % experiments/FIGS/data/v2/index.json :: fine-tuning qg N=10000: R30 of each L162 run over the from-scratch mean, SD over runs and standard error of that mean in quadrature :: sha256 c6c4138ee5c86423
\newcommand{\MlstFtQgRthirtyResOverScratchEFourPr}{15\% more rejection}  % experiments/FIGS/data/v2/index.json :: fine-tuning qg N=10000: R30 of each L162 run over the from-scratch mean, SD over runs and standard error of that mean in quadrature :: sha256 c6c4138ee5c86423 :: prose twin: two significant figures, direction in words
\newcommand{\MlstFtQgRthirtyMidEFour}{74\,\ensuremath{\pm}\,5}  % experiments/FIGS/data/v2/index.json :: fine-tuning qg N=10000: R30, R42_Q1 :: sha256 c6c4138ee5c86423
\newcommand{\MlstFtQgRthirtyMidOverScratchEFour}{1.12\,\ensuremath{\pm}\,0.11}  % experiments/FIGS/data/v2/index.json :: fine-tuning qg N=10000: R30 of each R42_Q1 run over the from-scratch mean, SD over runs and standard error of that mean in quadrature :: sha256 c6c4138ee5c86423
\newcommand{\MlstFtQgRthirtyMidOverScratchEFourPr}{equal within 20\%}  % experiments/FIGS/data/v2/index.json :: fine-tuning qg N=10000: R30 of each R42_Q1 run over the from-scratch mean, SD over runs and standard error of that mean in quadrature :: sha256 c6c4138ee5c86423 :: prose twin: two significant figures, direction in words
\newcommand{\MlstFtQgRthirtyCoarseEFour}{80.3\,\ensuremath{\pm}\,1.1}  % experiments/FIGS/data/v2/index.json :: fine-tuning qg N=10000: R30, R16_Q1 :: sha256 c6c4138ee5c86423
\newcommand{\MlstFtQgRthirtyCoarseOverScratchEFour}{1.21\,\ensuremath{\pm}\,0.10}  % experiments/FIGS/data/v2/index.json :: fine-tuning qg N=10000: R30 of each R16_Q1 run over the from-scratch mean, SD over runs and standard error of that mean in quadrature :: sha256 c6c4138ee5c86423
\newcommand{\MlstFtQgRthirtyCoarseOverScratchEFourPr}{21\% more rejection}  % experiments/FIGS/data/v2/index.json :: fine-tuning qg N=10000: R30 of each R16_Q1 run over the from-scratch mean, SD over runs and standard error of that mean in quadrature :: sha256 c6c4138ee5c86423 :: prose twin: two significant figures, direction in words
\newcommand{\MlstFtQgRthirtyResMassEFour}{80.1 and 72.6}  % experiments/FIGS/data/v2/index.json :: fine-tuning qg N=10000: R30, L162_MASS :: sha256 c6c4138ee5c86423
\newcommand{\MlstFtQgRthirtyResMassOverScratchEFour}{1.21 and 1.09}  % experiments/FIGS/data/v2/index.json :: fine-tuning qg N=10000: R30 of each L162_MASS run over the from-scratch mean, SD over runs and standard error of that mean in quadrature :: sha256 c6c4138ee5c86423
\newcommand{\MlstFtQgRthirtyResMassOverScratchEFourPr}{at least 9.3\% more rejection}  % experiments/FIGS/data/v2/index.json :: fine-tuning qg N=10000: R30 of each L162_MASS run over the from-scratch mean, SD over runs and standard error of that mean in quadrature :: sha256 c6c4138ee5c86423 :: prose twin: two significant figures, direction in words
\newcommand{\MlstFtQgRthirtyResMassOverScratchEFourPrMax}{up to 21\% more rejection}  % experiments/FIGS/data/v2/index.json :: fine-tuning qg N=10000: R30 of each L162_MASS run over the from-scratch mean, SD over runs and standard error of that mean in quadrature :: sha256 c6c4138ee5c86423 :: prose twin of the run farther from 1
\newcommand{\MlstFtQgRthirtyCoarseMassEFour}{79.5 and 77.9}  % experiments/FIGS/data/v2/index.json :: fine-tuning qg N=10000: R30, R16_Q1_MASS :: sha256 c6c4138ee5c86423
\newcommand{\MlstFtQgRthirtyCoarseMassOverScratchEFour}{1.198 and 1.173}  % experiments/FIGS/data/v2/index.json :: fine-tuning qg N=10000: R30 of each R16_Q1_MASS run over the from-scratch mean, SD over runs and standard error of that mean in quadrature :: sha256 c6c4138ee5c86423
\newcommand{\MlstFtQgRthirtyCoarseMassOverScratchEFourPr}{at least 17\% more rejection}  % experiments/FIGS/data/v2/index.json :: fine-tuning qg N=10000: R30 of each R16_Q1_MASS run over the from-scratch mean, SD over runs and standard error of that mean in quadrature :: sha256 c6c4138ee5c86423 :: prose twin: two significant figures, direction in words
\newcommand{\MlstFtQgRthirtyCoarseMassOverScratchEFourPrMax}{up to 20\% more rejection}  % experiments/FIGS/data/v2/index.json :: fine-tuning qg N=10000: R30 of each R16_Q1_MASS run over the from-scratch mean, SD over runs and standard error of that mean in quadrature :: sha256 c6c4138ee5c86423 :: prose twin of the run farther from 1
\newcommand{\MlstFtQgRthirtyResOverFullEFour}{1.02\,\ensuremath{\pm}\,0.10}  % experiments/FIGS/data/v2/index.json :: fine-tuning qg N=10000: R30 of L162 over L188, run by run :: sha256 c6c4138ee5c86423
\newcommand{\MlstFtQgRthirtyResOverFullEFourPr}{equal within 10\%}  % experiments/FIGS/data/v2/index.json :: fine-tuning qg N=10000: R30 of L162 over L188, run by run :: sha256 c6c4138ee5c86423 :: prose twin: two significant figures, direction in words
\newcommand{\MlstFtQgRthirtyMidOverFullEFour}{0.99\,\ensuremath{\pm}\,0.14}  % experiments/FIGS/data/v2/index.json :: fine-tuning qg N=10000: R30 of R42_Q1 over L188, run by run :: sha256 c6c4138ee5c86423
\newcommand{\MlstFtQgRthirtyMidOverFullEFourPr}{equal within 20\%}  % experiments/FIGS/data/v2/index.json :: fine-tuning qg N=10000: R30 of R42_Q1 over L188, run by run :: sha256 c6c4138ee5c86423 :: prose twin: two significant figures, direction in words
\newcommand{\MlstFtQgRthirtyCoarseOverFullEFour}{1.07\,\ensuremath{\pm}\,0.13}  % experiments/FIGS/data/v2/index.json :: fine-tuning qg N=10000: R30 of R16_Q1 over L188, run by run :: sha256 c6c4138ee5c86423
\newcommand{\MlstFtQgRthirtyCoarseOverFullEFourPr}{equal within 20\%}  % experiments/FIGS/data/v2/index.json :: fine-tuning qg N=10000: R30 of R16_Q1 over L188, run by run :: sha256 c6c4138ee5c86423 :: prose twin: two significant figures, direction in words
\newcommand{\MlstFtQgRthirtyResMassOverResEFour}{1.06 and 0.95}  % experiments/FIGS/data/v2/index.json :: fine-tuning qg N=10000: R30 of L162_MASS over L162, run by run :: sha256 c6c4138ee5c86423
\newcommand{\MlstFtQgRthirtyResMassOverResEFourPr}{equal within 6\%}  % experiments/FIGS/data/v2/index.json :: fine-tuning qg N=10000: R30 of L162_MASS over L162, run by run :: sha256 c6c4138ee5c86423 :: prose twin: two significant figures, direction in words
\newcommand{\MlstFtQgRthirtyResMassOverResEFourPrMax}{equal within 6\%}  % experiments/FIGS/data/v2/index.json :: fine-tuning qg N=10000: R30 of L162_MASS over L162, run by run :: sha256 c6c4138ee5c86423 :: prose twin of the run farther from 1
\newcommand{\MlstFtQgRthirtyCoarseMassOverCoarseEFour}{0.9769 and 0.9696}  % experiments/FIGS/data/v2/index.json :: fine-tuning qg N=10000: R30 of R16_Q1_MASS over R16_Q1, run by run :: sha256 c6c4138ee5c86423
\newcommand{\MlstFtQgRthirtyCoarseMassOverCoarseEFourPr}{equal within 4\%}  % experiments/FIGS/data/v2/index.json :: fine-tuning qg N=10000: R30 of R16_Q1_MASS over R16_Q1, run by run :: sha256 c6c4138ee5c86423 :: prose twin: two significant figures, direction in words
\newcommand{\MlstFtQgRthirtyCoarseMassOverCoarseEFourPrMax}{equal within 4\%}  % experiments/FIGS/data/v2/index.json :: fine-tuning qg N=10000: R30 of R16_Q1_MASS over R16_Q1, run by run :: sha256 c6c4138ee5c86423 :: prose twin of the run farther from 1
\newcommand{\MlstFtQgRthirtyRandOverFullEFour}{1.03\,\ensuremath{\pm}\,0.18}  % experiments/FIGS/data/v2/index.json :: fine-tuning qg N=10000: R30 of each random partition over the L188 run of its seed :: sha256 c6c4138ee5c86423
\newcommand{\MlstFtQgRthirtyRandOverFullEFourPr}{equal within 20\%}  % experiments/FIGS/data/v2/index.json :: fine-tuning qg N=10000: R30 of each random partition over the L188 run of its seed :: sha256 c6c4138ee5c86423 :: prose twin: two significant figures, direction in words
\newcommand{\MlstFtQgRthirtyRandOverCoarseEFour}{0.93\,\ensuremath{\pm}\,0.12}  % experiments/FIGS/data/v2/index.json :: fine-tuning qg N=10000: R30 of each random partition over the R16_Q1 run of its seed :: sha256 c6c4138ee5c86423
\newcommand{\MlstFtQgRthirtyRandOverCoarseEFourPr}{equal within 20\%}  % experiments/FIGS/data/v2/index.json :: fine-tuning qg N=10000: R30 of each random partition over the R16_Q1 run of its seed :: sha256 c6c4138ee5c86423 :: prose twin: two significant figures, direction in words
\newcommand{\MlstFtQgRthirtyScratchEFive}{88.4\,\ensuremath{\pm}\,0.8}  % experiments/FIGS/data/v2/index.json :: fine-tuning qg N=100000: R30 from scratch, every seed :: sha256 c6c4138ee5c86423
\newcommand{\MlstFtQgRthirtyFullEFive}{99.3\,\ensuremath{\pm}\,1.4}  % experiments/FIGS/data/v2/index.json :: fine-tuning qg N=100000: R30, L188 :: sha256 c6c4138ee5c86423
\newcommand{\MlstFtQgRthirtyFullOverScratchEFive}{1.124\,\ensuremath{\pm}\,0.017}  % experiments/FIGS/data/v2/index.json :: fine-tuning qg N=100000: R30 of each L188 run over the from-scratch mean, SD over runs and standard error of that mean in quadrature :: sha256 c6c4138ee5c86423
\newcommand{\MlstFtQgRthirtyFullOverScratchEFivePr}{12\% more rejection}  % experiments/FIGS/data/v2/index.json :: fine-tuning qg N=100000: R30 of each L188 run over the from-scratch mean, SD over runs and standard error of that mean in quadrature :: sha256 c6c4138ee5c86423 :: prose twin: two significant figures, direction in words
\newcommand{\MlstFtQgRthirtyResEFive}{97.9\,\ensuremath{\pm}\,3.3}  % experiments/FIGS/data/v2/index.json :: fine-tuning qg N=100000: R30, L162 :: sha256 c6c4138ee5c86423
\newcommand{\MlstFtQgRthirtyResOverScratchEFive}{1.11\,\ensuremath{\pm}\,0.04}  % experiments/FIGS/data/v2/index.json :: fine-tuning qg N=100000: R30 of each L162 run over the from-scratch mean, SD over runs and standard error of that mean in quadrature :: sha256 c6c4138ee5c86423
\newcommand{\MlstFtQgRthirtyResOverScratchEFivePr}{11\% more rejection}  % experiments/FIGS/data/v2/index.json :: fine-tuning qg N=100000: R30 of each L162 run over the from-scratch mean, SD over runs and standard error of that mean in quadrature :: sha256 c6c4138ee5c86423 :: prose twin: two significant figures, direction in words
\newcommand{\MlstFtQgRthirtyMidEFive}{96.1\,\ensuremath{\pm}\,1.8}  % experiments/FIGS/data/v2/index.json :: fine-tuning qg N=100000: R30, R42_Q1 :: sha256 c6c4138ee5c86423
\newcommand{\MlstFtQgRthirtyMidOverScratchEFive}{1.088\,\ensuremath{\pm}\,0.021}  % experiments/FIGS/data/v2/index.json :: fine-tuning qg N=100000: R30 of each R42_Q1 run over the from-scratch mean, SD over runs and standard error of that mean in quadrature :: sha256 c6c4138ee5c86423
\newcommand{\MlstFtQgRthirtyMidOverScratchEFivePr}{8.8\% more rejection}  % experiments/FIGS/data/v2/index.json :: fine-tuning qg N=100000: R30 of each R42_Q1 run over the from-scratch mean, SD over runs and standard error of that mean in quadrature :: sha256 c6c4138ee5c86423 :: prose twin: two significant figures, direction in words
\newcommand{\MlstFtQgRthirtyCoarseEFive}{96.4\,\ensuremath{\pm}\,1.0}  % experiments/FIGS/data/v2/index.json :: fine-tuning qg N=100000: R30, R16_Q1 :: sha256 c6c4138ee5c86423
\newcommand{\MlstFtQgRthirtyCoarseOverScratchEFive}{1.091\,\ensuremath{\pm}\,0.013}  % experiments/FIGS/data/v2/index.json :: fine-tuning qg N=100000: R30 of each R16_Q1 run over the from-scratch mean, SD over runs and standard error of that mean in quadrature :: sha256 c6c4138ee5c86423
\newcommand{\MlstFtQgRthirtyCoarseOverScratchEFivePr}{9.1\% more rejection}  % experiments/FIGS/data/v2/index.json :: fine-tuning qg N=100000: R30 of each R16_Q1 run over the from-scratch mean, SD over runs and standard error of that mean in quadrature :: sha256 c6c4138ee5c86423 :: prose twin: two significant figures, direction in words
\newcommand{\MlstFtQgRthirtyResMassEFive}{95.6 and 97.6}  % experiments/FIGS/data/v2/index.json :: fine-tuning qg N=100000: R30, L162_MASS :: sha256 c6c4138ee5c86423
\newcommand{\MlstFtQgRthirtyResMassOverScratchEFive}{1.082 and 1.104}  % experiments/FIGS/data/v2/index.json :: fine-tuning qg N=100000: R30 of each L162_MASS run over the from-scratch mean, SD over runs and standard error of that mean in quadrature :: sha256 c6c4138ee5c86423
\newcommand{\MlstFtQgRthirtyResMassOverScratchEFivePr}{at least 8.2\% more rejection}  % experiments/FIGS/data/v2/index.json :: fine-tuning qg N=100000: R30 of each L162_MASS run over the from-scratch mean, SD over runs and standard error of that mean in quadrature :: sha256 c6c4138ee5c86423 :: prose twin: two significant figures, direction in words
\newcommand{\MlstFtQgRthirtyResMassOverScratchEFivePrMax}{up to 10\% more rejection}  % experiments/FIGS/data/v2/index.json :: fine-tuning qg N=100000: R30 of each L162_MASS run over the from-scratch mean, SD over runs and standard error of that mean in quadrature :: sha256 c6c4138ee5c86423 :: prose twin of the run farther from 1
\newcommand{\MlstFtQgRthirtyCoarseMassEFive}{97.85 and 97.94}  % experiments/FIGS/data/v2/index.json :: fine-tuning qg N=100000: R30, R16_Q1_MASS :: sha256 c6c4138ee5c86423
\newcommand{\MlstFtQgRthirtyCoarseMassOverScratchEFive}{1.107 and 1.108}  % experiments/FIGS/data/v2/index.json :: fine-tuning qg N=100000: R30 of each R16_Q1_MASS run over the from-scratch mean, SD over runs and standard error of that mean in quadrature :: sha256 c6c4138ee5c86423
\newcommand{\MlstFtQgRthirtyCoarseMassOverScratchEFivePr}{at least 11\% more rejection}  % experiments/FIGS/data/v2/index.json :: fine-tuning qg N=100000: R30 of each R16_Q1_MASS run over the from-scratch mean, SD over runs and standard error of that mean in quadrature :: sha256 c6c4138ee5c86423 :: prose twin: two significant figures, direction in words
\newcommand{\MlstFtQgRthirtyCoarseMassOverScratchEFivePrMax}{up to 11\% more rejection}  % experiments/FIGS/data/v2/index.json :: fine-tuning qg N=100000: R30 of each R16_Q1_MASS run over the from-scratch mean, SD over runs and standard error of that mean in quadrature :: sha256 c6c4138ee5c86423 :: prose twin of the run farther from 1
\newcommand{\MlstFtQgRthirtyResOverFullEFive}{0.985\,\ensuremath{\pm}\,0.020}  % experiments/FIGS/data/v2/index.json :: fine-tuning qg N=100000: R30 of L162 over L188, run by run :: sha256 c6c4138ee5c86423
\newcommand{\MlstFtQgRthirtyResOverFullEFivePr}{equal within 2\%}  % experiments/FIGS/data/v2/index.json :: fine-tuning qg N=100000: R30 of L162 over L188, run by run :: sha256 c6c4138ee5c86423 :: prose twin: two significant figures, direction in words
\newcommand{\MlstFtQgRthirtyMidOverFullEFive}{0.969\,\ensuremath{\pm}\,0.030}  % experiments/FIGS/data/v2/index.json :: fine-tuning qg N=100000: R30 of R42_Q1 over L188, run by run :: sha256 c6c4138ee5c86423
\newcommand{\MlstFtQgRthirtyMidOverFullEFivePr}{equal within 4\%}  % experiments/FIGS/data/v2/index.json :: fine-tuning qg N=100000: R30 of R42_Q1 over L188, run by run :: sha256 c6c4138ee5c86423 :: prose twin: two significant figures, direction in words
\newcommand{\MlstFtQgRthirtyCoarseOverFullEFive}{0.972\,\ensuremath{\pm}\,0.020}  % experiments/FIGS/data/v2/index.json :: fine-tuning qg N=100000: R30 of R16_Q1 over L188, run by run :: sha256 c6c4138ee5c86423
\newcommand{\MlstFtQgRthirtyCoarseOverFullEFivePr}{equal within 3\%}  % experiments/FIGS/data/v2/index.json :: fine-tuning qg N=100000: R30 of R16_Q1 over L188, run by run :: sha256 c6c4138ee5c86423 :: prose twin: two significant figures, direction in words
\newcommand{\MlstFtQgRthirtyResMassOverResEFive}{0.942 and 1.027}  % experiments/FIGS/data/v2/index.json :: fine-tuning qg N=100000: R30 of L162_MASS over L162, run by run :: sha256 c6c4138ee5c86423
\newcommand{\MlstFtQgRthirtyResMassOverResEFivePr}{equal within 6\%}  % experiments/FIGS/data/v2/index.json :: fine-tuning qg N=100000: R30 of L162_MASS over L162, run by run :: sha256 c6c4138ee5c86423 :: prose twin: two significant figures, direction in words
\newcommand{\MlstFtQgRthirtyResMassOverResEFivePrMax}{equal within 6\%}  % experiments/FIGS/data/v2/index.json :: fine-tuning qg N=100000: R30 of L162_MASS over L162, run by run :: sha256 c6c4138ee5c86423 :: prose twin of the run farther from 1
\newcommand{\MlstFtQgRthirtyCoarseMassOverCoarseEFive}{1.018 and 1.004}  % experiments/FIGS/data/v2/index.json :: fine-tuning qg N=100000: R30 of R16_Q1_MASS over R16_Q1, run by run :: sha256 c6c4138ee5c86423
\newcommand{\MlstFtQgRthirtyCoarseMassOverCoarseEFivePr}{equal within 2\%}  % experiments/FIGS/data/v2/index.json :: fine-tuning qg N=100000: R30 of R16_Q1_MASS over R16_Q1, run by run :: sha256 c6c4138ee5c86423 :: prose twin: two significant figures, direction in words
\newcommand{\MlstFtQgRthirtyCoarseMassOverCoarseEFivePrMax}{equal within 2\%}  % experiments/FIGS/data/v2/index.json :: fine-tuning qg N=100000: R30 of R16_Q1_MASS over R16_Q1, run by run :: sha256 c6c4138ee5c86423 :: prose twin of the run farther from 1
\newcommand{\MlstFtQgRthirtyRandOverFullEFive}{0.986\,\ensuremath{\pm}\,0.025}  % experiments/FIGS/data/v2/index.json :: fine-tuning qg N=100000: R30 of each random partition over the L188 run of its seed :: sha256 c6c4138ee5c86423
\newcommand{\MlstFtQgRthirtyRandOverFullEFivePr}{equal within 3\%}  % experiments/FIGS/data/v2/index.json :: fine-tuning qg N=100000: R30 of each random partition over the L188 run of its seed :: sha256 c6c4138ee5c86423 :: prose twin: two significant figures, direction in words
\newcommand{\MlstFtQgRthirtyRandOverCoarseEFive}{1.009\,\ensuremath{\pm}\,0.020}  % experiments/FIGS/data/v2/index.json :: fine-tuning qg N=100000: R30 of each random partition over the R16_Q1 run of its seed :: sha256 c6c4138ee5c86423
\newcommand{\MlstFtQgRthirtyRandOverCoarseEFivePr}{equal within 3\%}  % experiments/FIGS/data/v2/index.json :: fine-tuning qg N=100000: R30 of each random partition over the R16_Q1 run of its seed :: sha256 c6c4138ee5c86423 :: prose twin: two significant figures, direction in words
\newcommand{\MlstFtQgRthirtyScratch}{104.0\,\ensuremath{\pm}\,1.0}  % experiments/FIGS/data/v2/index.json :: fine-tuning qg N=1600000: R30 from scratch, every seed :: sha256 c6c4138ee5c86423
\newcommand{\MlstFtQgRthirtyFull}{107.3\,\ensuremath{\pm}\,2.2}  % experiments/FIGS/data/v2/index.json :: fine-tuning qg N=1600000: R30, L188 :: sha256 c6c4138ee5c86423
\newcommand{\MlstFtQgRthirtyFullOverScratch}{1.031\,\ensuremath{\pm}\,0.022}  % experiments/FIGS/data/v2/index.json :: fine-tuning qg N=1600000: R30 of each L188 run over the from-scratch mean, SD over runs and standard error of that mean in quadrature :: sha256 c6c4138ee5c86423
\newcommand{\MlstFtQgRthirtyFullOverScratchPr}{equal within 4\%}  % experiments/FIGS/data/v2/index.json :: fine-tuning qg N=1600000: R30 of each L188 run over the from-scratch mean, SD over runs and standard error of that mean in quadrature :: sha256 c6c4138ee5c86423 :: prose twin: two significant figures, direction in words
\newcommand{\MlstFtQgRthirtyRes}{108.0\,\ensuremath{\pm}\,1.4}  % experiments/FIGS/data/v2/index.json :: fine-tuning qg N=1600000: R30, L162 :: sha256 c6c4138ee5c86423
\newcommand{\MlstFtQgRthirtyResOverScratch}{1.038\,\ensuremath{\pm}\,0.015}  % experiments/FIGS/data/v2/index.json :: fine-tuning qg N=1600000: R30 of each L162 run over the from-scratch mean, SD over runs and standard error of that mean in quadrature :: sha256 c6c4138ee5c86423
\newcommand{\MlstFtQgRthirtyResOverScratchPr}{3.8\% more rejection}  % experiments/FIGS/data/v2/index.json :: fine-tuning qg N=1600000: R30 of each L162 run over the from-scratch mean, SD over runs and standard error of that mean in quadrature :: sha256 c6c4138ee5c86423 :: prose twin: two significant figures, direction in words
\newcommand{\MlstFtQgRthirtyMid}{108.3\,\ensuremath{\pm}\,1.0}  % experiments/FIGS/data/v2/index.json :: fine-tuning qg N=1600000: R30, R42_Q1 :: sha256 c6c4138ee5c86423
\newcommand{\MlstFtQgRthirtyMidOverScratch}{1.041\,\ensuremath{\pm}\,0.012}  % experiments/FIGS/data/v2/index.json :: fine-tuning qg N=1600000: R30 of each R42_Q1 run over the from-scratch mean, SD over runs and standard error of that mean in quadrature :: sha256 c6c4138ee5c86423
\newcommand{\MlstFtQgRthirtyMidOverScratchPr}{4.1\% more rejection}  % experiments/FIGS/data/v2/index.json :: fine-tuning qg N=1600000: R30 of each R42_Q1 run over the from-scratch mean, SD over runs and standard error of that mean in quadrature :: sha256 c6c4138ee5c86423 :: prose twin: two significant figures, direction in words
\newcommand{\MlstFtQgRthirtyCoarse}{106.5\,\ensuremath{\pm}\,0.8}  % experiments/FIGS/data/v2/index.json :: fine-tuning qg N=1600000: R30, R16_Q1 :: sha256 c6c4138ee5c86423
\newcommand{\MlstFtQgRthirtyCoarseOverScratch}{1.024\,\ensuremath{\pm}\,0.010}  % experiments/FIGS/data/v2/index.json :: fine-tuning qg N=1600000: R30 of each R16_Q1 run over the from-scratch mean, SD over runs and standard error of that mean in quadrature :: sha256 c6c4138ee5c86423
\newcommand{\MlstFtQgRthirtyCoarseOverScratchPr}{2.4\% more rejection}  % experiments/FIGS/data/v2/index.json :: fine-tuning qg N=1600000: R30 of each R16_Q1 run over the from-scratch mean, SD over runs and standard error of that mean in quadrature :: sha256 c6c4138ee5c86423 :: prose twin: two significant figures, direction in words
\newcommand{\MlstFtQgRthirtyResMass}{107.0 and 106.6}  % experiments/FIGS/data/v2/index.json :: fine-tuning qg N=1600000: R30, L162_MASS :: sha256 c6c4138ee5c86423
\newcommand{\MlstFtQgRthirtyResMassOverScratch}{1.028 and 1.025}  % experiments/FIGS/data/v2/index.json :: fine-tuning qg N=1600000: R30 of each L162_MASS run over the from-scratch mean, SD over runs and standard error of that mean in quadrature :: sha256 c6c4138ee5c86423
\newcommand{\MlstFtQgRthirtyResMassOverScratchPr}{at least 2.5\% more rejection}  % experiments/FIGS/data/v2/index.json :: fine-tuning qg N=1600000: R30 of each L162_MASS run over the from-scratch mean, SD over runs and standard error of that mean in quadrature :: sha256 c6c4138ee5c86423 :: prose twin: two significant figures, direction in words
\newcommand{\MlstFtQgRthirtyResMassOverScratchPrMax}{up to 2.8\% more rejection}  % experiments/FIGS/data/v2/index.json :: fine-tuning qg N=1600000: R30 of each L162_MASS run over the from-scratch mean, SD over runs and standard error of that mean in quadrature :: sha256 c6c4138ee5c86423 :: prose twin of the run farther from 1
\newcommand{\MlstFtQgRthirtyCoarseMass}{104.9 and 105.6}  % experiments/FIGS/data/v2/index.json :: fine-tuning qg N=1600000: R30, R16_Q1_MASS :: sha256 c6c4138ee5c86423
\newcommand{\MlstFtQgRthirtyCoarseMassOverScratch}{1.009 and 1.015}  % experiments/FIGS/data/v2/index.json :: fine-tuning qg N=1600000: R30 of each R16_Q1_MASS run over the from-scratch mean, SD over runs and standard error of that mean in quadrature :: sha256 c6c4138ee5c86423
\newcommand{\MlstFtQgRthirtyCoarseMassOverScratchPr}{at least 0.87\% more rejection}  % experiments/FIGS/data/v2/index.json :: fine-tuning qg N=1600000: R30 of each R16_Q1_MASS run over the from-scratch mean, SD over runs and standard error of that mean in quadrature :: sha256 c6c4138ee5c86423 :: prose twin: two significant figures, direction in words
\newcommand{\MlstFtQgRthirtyCoarseMassOverScratchPrMax}{up to 1.5\% more rejection}  % experiments/FIGS/data/v2/index.json :: fine-tuning qg N=1600000: R30 of each R16_Q1_MASS run over the from-scratch mean, SD over runs and standard error of that mean in quadrature :: sha256 c6c4138ee5c86423 :: prose twin of the run farther from 1
\newcommand{\MlstFtQgRthirtyResOverFull}{1.007\,\ensuremath{\pm}\,0.026}  % experiments/FIGS/data/v2/index.json :: fine-tuning qg N=1600000: R30 of L162 over L188, run by run :: sha256 c6c4138ee5c86423
\newcommand{\MlstFtQgRthirtyResOverFullPr}{equal within 3\%}  % experiments/FIGS/data/v2/index.json :: fine-tuning qg N=1600000: R30 of L162 over L188, run by run :: sha256 c6c4138ee5c86423 :: prose twin: two significant figures, direction in words
\newcommand{\MlstFtQgRthirtyMidOverFull}{1.010\,\ensuremath{\pm}\,0.029}  % experiments/FIGS/data/v2/index.json :: fine-tuning qg N=1600000: R30 of R42_Q1 over L188, run by run :: sha256 c6c4138ee5c86423
\newcommand{\MlstFtQgRthirtyMidOverFullPr}{equal within 3\%}  % experiments/FIGS/data/v2/index.json :: fine-tuning qg N=1600000: R30 of R42_Q1 over L188, run by run :: sha256 c6c4138ee5c86423 :: prose twin: two significant figures, direction in words
\newcommand{\MlstFtQgRthirtyCoarseOverFull}{0.993\,\ensuremath{\pm}\,0.026}  % experiments/FIGS/data/v2/index.json :: fine-tuning qg N=1600000: R30 of R16_Q1 over L188, run by run :: sha256 c6c4138ee5c86423
\newcommand{\MlstFtQgRthirtyCoarseOverFullPr}{equal within 3\%}  % experiments/FIGS/data/v2/index.json :: fine-tuning qg N=1600000: R30 of R16_Q1 over L188, run by run :: sha256 c6c4138ee5c86423 :: prose twin: two significant figures, direction in words
\newcommand{\MlstFtQgRthirtyResMassOverRes}{0.976 and 0.999}  % experiments/FIGS/data/v2/index.json :: fine-tuning qg N=1600000: R30 of L162_MASS over L162, run by run :: sha256 c6c4138ee5c86423
\newcommand{\MlstFtQgRthirtyResMassOverResPr}{equal within 3\%}  % experiments/FIGS/data/v2/index.json :: fine-tuning qg N=1600000: R30 of L162_MASS over L162, run by run :: sha256 c6c4138ee5c86423 :: prose twin: two significant figures, direction in words
\newcommand{\MlstFtQgRthirtyResMassOverResPrMax}{equal within 3\%}  % experiments/FIGS/data/v2/index.json :: fine-tuning qg N=1600000: R30 of L162_MASS over L162, run by run :: sha256 c6c4138ee5c86423 :: prose twin of the run farther from 1
\newcommand{\MlstFtQgRthirtyCoarseMassOverCoarse}{0.9916 and 0.9926}  % experiments/FIGS/data/v2/index.json :: fine-tuning qg N=1600000: R30 of R16_Q1_MASS over R16_Q1, run by run :: sha256 c6c4138ee5c86423
\newcommand{\MlstFtQgRthirtyCoarseMassOverCoarsePr}{equal within 1\%}  % experiments/FIGS/data/v2/index.json :: fine-tuning qg N=1600000: R30 of R16_Q1_MASS over R16_Q1, run by run :: sha256 c6c4138ee5c86423 :: prose twin: two significant figures, direction in words
\newcommand{\MlstFtQgRthirtyCoarseMassOverCoarsePrMax}{equal within 1\%}  % experiments/FIGS/data/v2/index.json :: fine-tuning qg N=1600000: R30 of R16_Q1_MASS over R16_Q1, run by run :: sha256 c6c4138ee5c86423 :: prose twin of the run farther from 1
\newcommand{\MlstFtQgRthirtyRandOverFull}{0.985\,\ensuremath{\pm}\,0.011}  % experiments/FIGS/data/v2/index.json :: fine-tuning qg N=1600000: R30 of each random partition over the L188 run of its seed :: sha256 c6c4138ee5c86423
\newcommand{\MlstFtQgRthirtyRandOverFullPr}{equal within 2\%}  % experiments/FIGS/data/v2/index.json :: fine-tuning qg N=1600000: R30 of each random partition over the L188 run of its seed :: sha256 c6c4138ee5c86423 :: prose twin: two significant figures, direction in words
\newcommand{\MlstFtQgRthirtyRandOverCoarse}{1.008\,\ensuremath{\pm}\,0.009}  % experiments/FIGS/data/v2/index.json :: fine-tuning qg N=1600000: R30 of each random partition over the R16_Q1 run of its seed :: sha256 c6c4138ee5c86423
\newcommand{\MlstFtQgRthirtyRandOverCoarsePr}{equal within 1\%}  % experiments/FIGS/data/v2/index.json :: fine-tuning qg N=1600000: R30 of each random partition over the R16_Q1 run of its seed :: sha256 c6c4138ee5c86423 :: prose twin: two significant figures, direction in words
\newcommand{\MlstFtQgRthirtyFullOverPart}{0.994\,\ensuremath{\pm}\,0.020}  % experiments/FIGS/data/v2/index.json :: fine-tuning qg N=1600000: R30 of each 188-class run over ParT-f.t.'s published value :: sha256 c6c4138ee5c86423
\newcommand{\MlstFtQgRthirtyFullOverPartPr}{equal within 3\%}  % experiments/FIGS/data/v2/index.json :: fine-tuning qg N=1600000: R30 of each 188-class run over ParT-f.t.'s published value :: sha256 c6c4138ee5c86423 :: prose twin: two significant figures, direction in words
\newcommand{\MlstFtQgAucFull}{0.9148\,\ensuremath{\pm}\,0.0001}  % experiments/FIGS/data/v2/index.json :: fine-tuning qg N=1600000: AUC of the 188-class runs :: sha256 c6c4138ee5c86423
\newcommand{\MlstFtQgAucScratch}{0.9121\,\ensuremath{\pm}\,0.0002}  % experiments/FIGS/data/v2/index.json :: fine-tuning qg N=1600000: AUC from scratch :: sha256 c6c4138ee5c86423
\newcommand{\MlstFtQgDataEquivJets}{\ensuremath{<}1{,}000}  % experiments/FIGS/data/v2/index.json :: fine-tuning qg: from-scratch training jets whose R30 equals the 188-class mean at 1000 jets (interpolated linearly in R30 against log10 N) :: sha256 c6c4138ee5c86423
\newcommand{\MlstFtQgDataEquivJetsEFour}{26{,}000}  % experiments/FIGS/data/v2/index.json :: fine-tuning qg: from-scratch training jets whose R30 equals the 188-class mean at 10000 jets (interpolated linearly in R30 against log10 N) :: sha256 c6c4138ee5c86423
\newcommand{\MlstFtQgHerwigRthirtyScratchEThree}{27.0\,\ensuremath{\pm}\,1.3}  % experiments/FIGS/data/v2/index.json :: fine-tuning qg_herwig N=1000: R30 from scratch, every seed :: sha256 c6c4138ee5c86423
\newcommand{\MlstFtQgHerwigRthirtyFullEThree}{21\,\ensuremath{\pm}\,6}  % experiments/FIGS/data/v2/index.json :: fine-tuning qg_herwig N=1000: R30, L188 :: sha256 c6c4138ee5c86423
\newcommand{\MlstFtQgHerwigRthirtyFullOverScratchEThree}{0.77\,\ensuremath{\pm}\,0.22}  % experiments/FIGS/data/v2/index.json :: fine-tuning qg_herwig N=1000: R30 of each L188 run over the from-scratch mean, SD over runs and standard error of that mean in quadrature :: sha256 c6c4138ee5c86423
\newcommand{\MlstFtQgHerwigRthirtyFullOverScratchEThreePr}{equal within 30\%}  % experiments/FIGS/data/v2/index.json :: fine-tuning qg_herwig N=1000: R30 of each L188 run over the from-scratch mean, SD over runs and standard error of that mean in quadrature :: sha256 c6c4138ee5c86423 :: prose twin: two significant figures, direction in words
\newcommand{\MlstFtQgHerwigRthirtyResEThree}{26.2\,\ensuremath{\pm}\,0.9}  % experiments/FIGS/data/v2/index.json :: fine-tuning qg_herwig N=1000: R30, L162 :: sha256 c6c4138ee5c86423
\newcommand{\MlstFtQgHerwigRthirtyResOverScratchEThree}{0.97\,\ensuremath{\pm}\,0.04}  % experiments/FIGS/data/v2/index.json :: fine-tuning qg_herwig N=1000: R30 of each L162 run over the from-scratch mean, SD over runs and standard error of that mean in quadrature :: sha256 c6c4138ee5c86423
\newcommand{\MlstFtQgHerwigRthirtyResOverScratchEThreePr}{equal within 4\%}  % experiments/FIGS/data/v2/index.json :: fine-tuning qg_herwig N=1000: R30 of each L162 run over the from-scratch mean, SD over runs and standard error of that mean in quadrature :: sha256 c6c4138ee5c86423 :: prose twin: two significant figures, direction in words
\newcommand{\MlstFtQgHerwigRthirtyMidEThree}{26.9\,\ensuremath{\pm}\,2.7}  % experiments/FIGS/data/v2/index.json :: fine-tuning qg_herwig N=1000: R30, R42_Q1 :: sha256 c6c4138ee5c86423
\newcommand{\MlstFtQgHerwigRthirtyMidOverScratchEThree}{1.00\,\ensuremath{\pm}\,0.10}  % experiments/FIGS/data/v2/index.json :: fine-tuning qg_herwig N=1000: R30 of each R42_Q1 run over the from-scratch mean, SD over runs and standard error of that mean in quadrature :: sha256 c6c4138ee5c86423
\newcommand{\MlstFtQgHerwigRthirtyMidOverScratchEThreePr}{equal within 10\%}  % experiments/FIGS/data/v2/index.json :: fine-tuning qg_herwig N=1000: R30 of each R42_Q1 run over the from-scratch mean, SD over runs and standard error of that mean in quadrature :: sha256 c6c4138ee5c86423 :: prose twin: two significant figures, direction in words
\newcommand{\MlstFtQgHerwigRthirtyCoarseEThree}{25.2\,\ensuremath{\pm}\,2.5}  % experiments/FIGS/data/v2/index.json :: fine-tuning qg_herwig N=1000: R30, R16_Q1 :: sha256 c6c4138ee5c86423
\newcommand{\MlstFtQgHerwigRthirtyCoarseOverScratchEThree}{0.94\,\ensuremath{\pm}\,0.10}  % experiments/FIGS/data/v2/index.json :: fine-tuning qg_herwig N=1000: R30 of each R16_Q1 run over the from-scratch mean, SD over runs and standard error of that mean in quadrature :: sha256 c6c4138ee5c86423
\newcommand{\MlstFtQgHerwigRthirtyCoarseOverScratchEThreePr}{equal within 10\%}  % experiments/FIGS/data/v2/index.json :: fine-tuning qg_herwig N=1000: R30 of each R16_Q1 run over the from-scratch mean, SD over runs and standard error of that mean in quadrature :: sha256 c6c4138ee5c86423 :: prose twin: two significant figures, direction in words
\newcommand{\MlstFtQgHerwigRthirtyResMassEThree}{28.0 and 29.0}  % experiments/FIGS/data/v2/index.json :: fine-tuning qg_herwig N=1000: R30, L162_MASS :: sha256 c6c4138ee5c86423
\newcommand{\MlstFtQgHerwigRthirtyResMassOverScratchEThree}{1.038 and 1.076}  % experiments/FIGS/data/v2/index.json :: fine-tuning qg_herwig N=1000: R30 of each L162_MASS run over the from-scratch mean, SD over runs and standard error of that mean in quadrature :: sha256 c6c4138ee5c86423
\newcommand{\MlstFtQgHerwigRthirtyResMassOverScratchEThreePr}{at least 3.8\% more rejection}  % experiments/FIGS/data/v2/index.json :: fine-tuning qg_herwig N=1000: R30 of each L162_MASS run over the from-scratch mean, SD over runs and standard error of that mean in quadrature :: sha256 c6c4138ee5c86423 :: prose twin: two significant figures, direction in words
\newcommand{\MlstFtQgHerwigRthirtyResMassOverScratchEThreePrMax}{up to 7.6\% more rejection}  % experiments/FIGS/data/v2/index.json :: fine-tuning qg_herwig N=1000: R30 of each L162_MASS run over the from-scratch mean, SD over runs and standard error of that mean in quadrature :: sha256 c6c4138ee5c86423 :: prose twin of the run farther from 1
\newcommand{\MlstFtQgHerwigRthirtyCoarseMassEThree}{25.8 and 27.4}  % experiments/FIGS/data/v2/index.json :: fine-tuning qg_herwig N=1000: R30, R16_Q1_MASS :: sha256 c6c4138ee5c86423
\newcommand{\MlstFtQgHerwigRthirtyCoarseMassOverScratchEThree}{0.955 and 1.015}  % experiments/FIGS/data/v2/index.json :: fine-tuning qg_herwig N=1000: R30 of each R16_Q1_MASS run over the from-scratch mean, SD over runs and standard error of that mean in quadrature :: sha256 c6c4138ee5c86423
\newcommand{\MlstFtQgHerwigRthirtyCoarseMassOverScratchEThreePr}{equal within 5\%}  % experiments/FIGS/data/v2/index.json :: fine-tuning qg_herwig N=1000: R30 of each R16_Q1_MASS run over the from-scratch mean, SD over runs and standard error of that mean in quadrature :: sha256 c6c4138ee5c86423 :: prose twin: two significant figures, direction in words
\newcommand{\MlstFtQgHerwigRthirtyCoarseMassOverScratchEThreePrMax}{equal within 5\%}  % experiments/FIGS/data/v2/index.json :: fine-tuning qg_herwig N=1000: R30 of each R16_Q1_MASS run over the from-scratch mean, SD over runs and standard error of that mean in quadrature :: sha256 c6c4138ee5c86423 :: prose twin of the run farther from 1
\newcommand{\MlstFtQgHerwigRthirtyResOverFullEThree}{1.33\,\ensuremath{\pm}\,0.33}  % experiments/FIGS/data/v2/index.json :: fine-tuning qg_herwig N=1000: R30 of L162 over L188, run by run :: sha256 c6c4138ee5c86423
\newcommand{\MlstFtQgHerwigRthirtyResOverFullEThreePr}{equal within 40\%}  % experiments/FIGS/data/v2/index.json :: fine-tuning qg_herwig N=1000: R30 of L162 over L188, run by run :: sha256 c6c4138ee5c86423 :: prose twin: two significant figures, direction in words
\newcommand{\MlstFtQgHerwigRthirtyMidOverFullEThree}{1.36\,\ensuremath{\pm}\,0.31}  % experiments/FIGS/data/v2/index.json :: fine-tuning qg_herwig N=1000: R30 of R42_Q1 over L188, run by run :: sha256 c6c4138ee5c86423
\newcommand{\MlstFtQgHerwigRthirtyMidOverFullEThreePr}{equal within 40\%}  % experiments/FIGS/data/v2/index.json :: fine-tuning qg_herwig N=1000: R30 of R42_Q1 over L188, run by run :: sha256 c6c4138ee5c86423 :: prose twin: two significant figures, direction in words
\newcommand{\MlstFtQgHerwigRthirtyCoarseOverFullEThree}{1.3\,\ensuremath{\pm}\,0.4}  % experiments/FIGS/data/v2/index.json :: fine-tuning qg_herwig N=1000: R30 of R16_Q1 over L188, run by run :: sha256 c6c4138ee5c86423
\newcommand{\MlstFtQgHerwigRthirtyCoarseOverFullEThreePr}{equal within 50\%}  % experiments/FIGS/data/v2/index.json :: fine-tuning qg_herwig N=1000: R30 of R16_Q1 over L188, run by run :: sha256 c6c4138ee5c86423 :: prose twin: two significant figures, direction in words
\newcommand{\MlstFtQgHerwigRthirtyResMassOverResEThree}{1.062 and 1.073}  % experiments/FIGS/data/v2/index.json :: fine-tuning qg_herwig N=1000: R30 of L162_MASS over L162, run by run :: sha256 c6c4138ee5c86423
\newcommand{\MlstFtQgHerwigRthirtyResMassOverResEThreePr}{equal within 8\%}  % experiments/FIGS/data/v2/index.json :: fine-tuning qg_herwig N=1000: R30 of L162_MASS over L162, run by run :: sha256 c6c4138ee5c86423 :: prose twin: two significant figures, direction in words
\newcommand{\MlstFtQgHerwigRthirtyResMassOverResEThreePrMax}{equal within 8\%}  % experiments/FIGS/data/v2/index.json :: fine-tuning qg_herwig N=1000: R30 of L162_MASS over L162, run by run :: sha256 c6c4138ee5c86423 :: prose twin of the run farther from 1
\newcommand{\MlstFtQgHerwigRthirtyCoarseMassOverCoarseEThree}{0.93 and 1.20}  % experiments/FIGS/data/v2/index.json :: fine-tuning qg_herwig N=1000: R30 of R16_Q1_MASS over R16_Q1, run by run :: sha256 c6c4138ee5c86423
\newcommand{\MlstFtQgHerwigRthirtyCoarseMassOverCoarseEThreePr}{equal within 30\%}  % experiments/FIGS/data/v2/index.json :: fine-tuning qg_herwig N=1000: R30 of R16_Q1_MASS over R16_Q1, run by run :: sha256 c6c4138ee5c86423 :: prose twin: two significant figures, direction in words
\newcommand{\MlstFtQgHerwigRthirtyCoarseMassOverCoarseEThreePrMax}{equal within 30\%}  % experiments/FIGS/data/v2/index.json :: fine-tuning qg_herwig N=1000: R30 of R16_Q1_MASS over R16_Q1, run by run :: sha256 c6c4138ee5c86423 :: prose twin of the run farther from 1
\newcommand{\MlstFtQgHerwigRthirtyRandOverFullEThree}{0.99\,\ensuremath{\pm}\,0.20}  % experiments/FIGS/data/v2/index.json :: fine-tuning qg_herwig N=1000: R30 of each random partition over the L188 run of its seed :: sha256 c6c4138ee5c86423
\newcommand{\MlstFtQgHerwigRthirtyRandOverFullEThreePr}{equal within 30\%}  % experiments/FIGS/data/v2/index.json :: fine-tuning qg_herwig N=1000: R30 of each random partition over the L188 run of its seed :: sha256 c6c4138ee5c86423 :: prose twin: two significant figures, direction in words
\newcommand{\MlstFtQgHerwigRthirtyRandOverCoarseEThree}{0.95\,\ensuremath{\pm}\,0.13}  % experiments/FIGS/data/v2/index.json :: fine-tuning qg_herwig N=1000: R30 of each random partition over the R16_Q1 run of its seed :: sha256 c6c4138ee5c86423
\newcommand{\MlstFtQgHerwigRthirtyRandOverCoarseEThreePr}{equal within 20\%}  % experiments/FIGS/data/v2/index.json :: fine-tuning qg_herwig N=1000: R30 of each random partition over the R16_Q1 run of its seed :: sha256 c6c4138ee5c86423 :: prose twin: two significant figures, direction in words
\newcommand{\MlstFtQgHerwigRthirtyScratchEFour}{31\,\ensuremath{\pm}\,7}  % experiments/FIGS/data/v2/index.json :: fine-tuning qg_herwig N=10000: R30 from scratch, every seed :: sha256 c6c4138ee5c86423
\newcommand{\MlstFtQgHerwigRthirtyFullEFour}{27.9\,\ensuremath{\pm}\,2.8}  % experiments/FIGS/data/v2/index.json :: fine-tuning qg_herwig N=10000: R30, L188 :: sha256 c6c4138ee5c86423
\newcommand{\MlstFtQgHerwigRthirtyFullOverScratchEFour}{0.89\,\ensuremath{\pm}\,0.14}  % experiments/FIGS/data/v2/index.json :: fine-tuning qg_herwig N=10000: R30 of each L188 run over the from-scratch mean, SD over runs and standard error of that mean in quadrature :: sha256 c6c4138ee5c86423
\newcommand{\MlstFtQgHerwigRthirtyFullOverScratchEFourPr}{equal within 20\%}  % experiments/FIGS/data/v2/index.json :: fine-tuning qg_herwig N=10000: R30 of each L188 run over the from-scratch mean, SD over runs and standard error of that mean in quadrature :: sha256 c6c4138ee5c86423 :: prose twin: two significant figures, direction in words
\newcommand{\MlstFtQgHerwigRthirtyResEFour}{27.4\,\ensuremath{\pm}\,0.5}  % experiments/FIGS/data/v2/index.json :: fine-tuning qg_herwig N=10000: R30, L162 :: sha256 c6c4138ee5c86423
\newcommand{\MlstFtQgHerwigRthirtyResOverScratchEFour}{0.87\,\ensuremath{\pm}\,0.11}  % experiments/FIGS/data/v2/index.json :: fine-tuning qg_herwig N=10000: R30 of each L162 run over the from-scratch mean, SD over runs and standard error of that mean in quadrature :: sha256 c6c4138ee5c86423
\newcommand{\MlstFtQgHerwigRthirtyResOverScratchEFourPr}{13\% less rejection}  % experiments/FIGS/data/v2/index.json :: fine-tuning qg_herwig N=10000: R30 of each L162 run over the from-scratch mean, SD over runs and standard error of that mean in quadrature :: sha256 c6c4138ee5c86423 :: prose twin: two significant figures, direction in words
\newcommand{\MlstFtQgHerwigRthirtyMidEFour}{28.4\,\ensuremath{\pm}\,2.1}  % experiments/FIGS/data/v2/index.json :: fine-tuning qg_herwig N=10000: R30, R42_Q1 :: sha256 c6c4138ee5c86423
\newcommand{\MlstFtQgHerwigRthirtyMidOverScratchEFour}{0.91\,\ensuremath{\pm}\,0.13}  % experiments/FIGS/data/v2/index.json :: fine-tuning qg_herwig N=10000: R30 of each R42_Q1 run over the from-scratch mean, SD over runs and standard error of that mean in quadrature :: sha256 c6c4138ee5c86423
\newcommand{\MlstFtQgHerwigRthirtyMidOverScratchEFourPr}{equal within 10\%}  % experiments/FIGS/data/v2/index.json :: fine-tuning qg_herwig N=10000: R30 of each R42_Q1 run over the from-scratch mean, SD over runs and standard error of that mean in quadrature :: sha256 c6c4138ee5c86423 :: prose twin: two significant figures, direction in words
\newcommand{\MlstFtQgHerwigRthirtyCoarseEFour}{32.7\,\ensuremath{\pm}\,1.8}  % experiments/FIGS/data/v2/index.json :: fine-tuning qg_herwig N=10000: R30, R16_Q1 :: sha256 c6c4138ee5c86423
\newcommand{\MlstFtQgHerwigRthirtyCoarseOverScratchEFour}{1.05\,\ensuremath{\pm}\,0.14}  % experiments/FIGS/data/v2/index.json :: fine-tuning qg_herwig N=10000: R30 of each R16_Q1 run over the from-scratch mean, SD over runs and standard error of that mean in quadrature :: sha256 c6c4138ee5c86423
\newcommand{\MlstFtQgHerwigRthirtyCoarseOverScratchEFourPr}{equal within 6\%}  % experiments/FIGS/data/v2/index.json :: fine-tuning qg_herwig N=10000: R30 of each R16_Q1 run over the from-scratch mean, SD over runs and standard error of that mean in quadrature :: sha256 c6c4138ee5c86423 :: prose twin: two significant figures, direction in words
\newcommand{\MlstFtQgHerwigRthirtyResMassEFour}{34.3 and 28.5}  % experiments/FIGS/data/v2/index.json :: fine-tuning qg_herwig N=10000: R30, L162_MASS :: sha256 c6c4138ee5c86423
\newcommand{\MlstFtQgHerwigRthirtyResMassOverScratchEFour}{1.09 and 0.91}  % experiments/FIGS/data/v2/index.json :: fine-tuning qg_herwig N=10000: R30 of each L162_MASS run over the from-scratch mean, SD over runs and standard error of that mean in quadrature :: sha256 c6c4138ee5c86423
\newcommand{\MlstFtQgHerwigRthirtyResMassOverScratchEFourPr}{equal within 10\%}  % experiments/FIGS/data/v2/index.json :: fine-tuning qg_herwig N=10000: R30 of each L162_MASS run over the from-scratch mean, SD over runs and standard error of that mean in quadrature :: sha256 c6c4138ee5c86423 :: prose twin: two significant figures, direction in words
\newcommand{\MlstFtQgHerwigRthirtyResMassOverScratchEFourPrMax}{equal within 10\%}  % experiments/FIGS/data/v2/index.json :: fine-tuning qg_herwig N=10000: R30 of each L162_MASS run over the from-scratch mean, SD over runs and standard error of that mean in quadrature :: sha256 c6c4138ee5c86423 :: prose twin of the run farther from 1
\newcommand{\MlstFtQgHerwigRthirtyCoarseMassEFour}{34.4 and 31.2}  % experiments/FIGS/data/v2/index.json :: fine-tuning qg_herwig N=10000: R30, R16_Q1_MASS :: sha256 c6c4138ee5c86423
\newcommand{\MlstFtQgHerwigRthirtyCoarseMassOverScratchEFour}{1.10 and 1.00}  % experiments/FIGS/data/v2/index.json :: fine-tuning qg_herwig N=10000: R30 of each R16_Q1_MASS run over the from-scratch mean, SD over runs and standard error of that mean in quadrature :: sha256 c6c4138ee5c86423
\newcommand{\MlstFtQgHerwigRthirtyCoarseMassOverScratchEFourPr}{equal within 20\%}  % experiments/FIGS/data/v2/index.json :: fine-tuning qg_herwig N=10000: R30 of each R16_Q1_MASS run over the from-scratch mean, SD over runs and standard error of that mean in quadrature :: sha256 c6c4138ee5c86423 :: prose twin: two significant figures, direction in words
\newcommand{\MlstFtQgHerwigRthirtyCoarseMassOverScratchEFourPrMax}{equal within 20\%}  % experiments/FIGS/data/v2/index.json :: fine-tuning qg_herwig N=10000: R30 of each R16_Q1_MASS run over the from-scratch mean, SD over runs and standard error of that mean in quadrature :: sha256 c6c4138ee5c86423 :: prose twin of the run farther from 1
\newcommand{\MlstFtQgHerwigRthirtyResOverFullEFour}{0.99\,\ensuremath{\pm}\,0.12}  % experiments/FIGS/data/v2/index.json :: fine-tuning qg_herwig N=10000: R30 of L162 over L188, run by run :: sha256 c6c4138ee5c86423
\newcommand{\MlstFtQgHerwigRthirtyResOverFullEFourPr}{equal within 20\%}  % experiments/FIGS/data/v2/index.json :: fine-tuning qg_herwig N=10000: R30 of L162 over L188, run by run :: sha256 c6c4138ee5c86423 :: prose twin: two significant figures, direction in words
\newcommand{\MlstFtQgHerwigRthirtyMidOverFullEFour}{1.03\,\ensuremath{\pm}\,0.18}  % experiments/FIGS/data/v2/index.json :: fine-tuning qg_herwig N=10000: R30 of R42_Q1 over L188, run by run :: sha256 c6c4138ee5c86423
\newcommand{\MlstFtQgHerwigRthirtyMidOverFullEFourPr}{equal within 20\%}  % experiments/FIGS/data/v2/index.json :: fine-tuning qg_herwig N=10000: R30 of R42_Q1 over L188, run by run :: sha256 c6c4138ee5c86423 :: prose twin: two significant figures, direction in words
\newcommand{\MlstFtQgHerwigRthirtyCoarseOverFullEFour}{1.18\,\ensuremath{\pm}\,0.08}  % experiments/FIGS/data/v2/index.json :: fine-tuning qg_herwig N=10000: R30 of R16_Q1 over L188, run by run :: sha256 c6c4138ee5c86423
\newcommand{\MlstFtQgHerwigRthirtyCoarseOverFullEFourPr}{18\% more rejection}  % experiments/FIGS/data/v2/index.json :: fine-tuning qg_herwig N=10000: R30 of R16_Q1 over L188, run by run :: sha256 c6c4138ee5c86423 :: prose twin: two significant figures, direction in words
\newcommand{\MlstFtQgHerwigRthirtyResMassOverResEFour}{1.23 and 1.03}  % experiments/FIGS/data/v2/index.json :: fine-tuning qg_herwig N=10000: R30 of L162_MASS over L162, run by run :: sha256 c6c4138ee5c86423
\newcommand{\MlstFtQgHerwigRthirtyResMassOverResEFourPr}{equal within 30\%}  % experiments/FIGS/data/v2/index.json :: fine-tuning qg_herwig N=10000: R30 of L162_MASS over L162, run by run :: sha256 c6c4138ee5c86423 :: prose twin: two significant figures, direction in words
\newcommand{\MlstFtQgHerwigRthirtyResMassOverResEFourPrMax}{equal within 30\%}  % experiments/FIGS/data/v2/index.json :: fine-tuning qg_herwig N=10000: R30 of L162_MASS over L162, run by run :: sha256 c6c4138ee5c86423 :: prose twin of the run farther from 1
\newcommand{\MlstFtQgHerwigRthirtyCoarseMassOverCoarseEFour}{1.12 and 0.91}  % experiments/FIGS/data/v2/index.json :: fine-tuning qg_herwig N=10000: R30 of R16_Q1_MASS over R16_Q1, run by run :: sha256 c6c4138ee5c86423
\newcommand{\MlstFtQgHerwigRthirtyCoarseMassOverCoarseEFourPr}{equal within 20\%}  % experiments/FIGS/data/v2/index.json :: fine-tuning qg_herwig N=10000: R30 of R16_Q1_MASS over R16_Q1, run by run :: sha256 c6c4138ee5c86423 :: prose twin: two significant figures, direction in words
\newcommand{\MlstFtQgHerwigRthirtyCoarseMassOverCoarseEFourPrMax}{equal within 20\%}  % experiments/FIGS/data/v2/index.json :: fine-tuning qg_herwig N=10000: R30 of R16_Q1_MASS over R16_Q1, run by run :: sha256 c6c4138ee5c86423 :: prose twin of the run farther from 1
\newcommand{\MlstFtQgHerwigRthirtyRandOverFullEFour}{1.13\,\ensuremath{\pm}\,0.18}  % experiments/FIGS/data/v2/index.json :: fine-tuning qg_herwig N=10000: R30 of each random partition over the L188 run of its seed :: sha256 c6c4138ee5c86423
\newcommand{\MlstFtQgHerwigRthirtyRandOverFullEFourPr}{equal within 20\%}  % experiments/FIGS/data/v2/index.json :: fine-tuning qg_herwig N=10000: R30 of each random partition over the L188 run of its seed :: sha256 c6c4138ee5c86423 :: prose twin: two significant figures, direction in words
\newcommand{\MlstFtQgHerwigRthirtyRandOverCoarseEFour}{0.93\,\ensuremath{\pm}\,0.14}  % experiments/FIGS/data/v2/index.json :: fine-tuning qg_herwig N=10000: R30 of each random partition over the R16_Q1 run of its seed :: sha256 c6c4138ee5c86423
\newcommand{\MlstFtQgHerwigRthirtyRandOverCoarseEFourPr}{equal within 20\%}  % experiments/FIGS/data/v2/index.json :: fine-tuning qg_herwig N=10000: R30 of each random partition over the R16_Q1 run of its seed :: sha256 c6c4138ee5c86423 :: prose twin: two significant figures, direction in words
\newcommand{\MlstFtQgHerwigRthirtyScratchEFive}{36.7\,\ensuremath{\pm}\,0.9}  % experiments/FIGS/data/v2/index.json :: fine-tuning qg_herwig N=100000: R30 from scratch, every seed :: sha256 c6c4138ee5c86423
\newcommand{\MlstFtQgHerwigRthirtyFullEFive}{34.3\,\ensuremath{\pm}\,1.0}  % experiments/FIGS/data/v2/index.json :: fine-tuning qg_herwig N=100000: R30, L188 :: sha256 c6c4138ee5c86423
\newcommand{\MlstFtQgHerwigRthirtyFullOverScratchEFive}{0.935\,\ensuremath{\pm}\,0.031}  % experiments/FIGS/data/v2/index.json :: fine-tuning qg_herwig N=100000: R30 of each L188 run over the from-scratch mean, SD over runs and standard error of that mean in quadrature :: sha256 c6c4138ee5c86423
\newcommand{\MlstFtQgHerwigRthirtyFullOverScratchEFivePr}{6.5\% less rejection}  % experiments/FIGS/data/v2/index.json :: fine-tuning qg_herwig N=100000: R30 of each L188 run over the from-scratch mean, SD over runs and standard error of that mean in quadrature :: sha256 c6c4138ee5c86423 :: prose twin: two significant figures, direction in words
\newcommand{\MlstFtQgHerwigRthirtyResEFive}{35.5\,\ensuremath{\pm}\,0.5}  % experiments/FIGS/data/v2/index.json :: fine-tuning qg_herwig N=100000: R30, L162 :: sha256 c6c4138ee5c86423
\newcommand{\MlstFtQgHerwigRthirtyResOverScratchEFive}{0.966\,\ensuremath{\pm}\,0.019}  % experiments/FIGS/data/v2/index.json :: fine-tuning qg_herwig N=100000: R30 of each L162 run over the from-scratch mean, SD over runs and standard error of that mean in quadrature :: sha256 c6c4138ee5c86423
\newcommand{\MlstFtQgHerwigRthirtyResOverScratchEFivePr}{3.4\% less rejection}  % experiments/FIGS/data/v2/index.json :: fine-tuning qg_herwig N=100000: R30 of each L162 run over the from-scratch mean, SD over runs and standard error of that mean in quadrature :: sha256 c6c4138ee5c86423 :: prose twin: two significant figures, direction in words
\newcommand{\MlstFtQgHerwigRthirtyMidEFive}{35.5\,\ensuremath{\pm}\,2.1}  % experiments/FIGS/data/v2/index.json :: fine-tuning qg_herwig N=100000: R30, R42_Q1 :: sha256 c6c4138ee5c86423
\newcommand{\MlstFtQgHerwigRthirtyMidOverScratchEFive}{0.97\,\ensuremath{\pm}\,0.06}  % experiments/FIGS/data/v2/index.json :: fine-tuning qg_herwig N=100000: R30 of each R42_Q1 run over the from-scratch mean, SD over runs and standard error of that mean in quadrature :: sha256 c6c4138ee5c86423
\newcommand{\MlstFtQgHerwigRthirtyMidOverScratchEFivePr}{equal within 6\%}  % experiments/FIGS/data/v2/index.json :: fine-tuning qg_herwig N=100000: R30 of each R42_Q1 run over the from-scratch mean, SD over runs and standard error of that mean in quadrature :: sha256 c6c4138ee5c86423 :: prose twin: two significant figures, direction in words
\newcommand{\MlstFtQgHerwigRthirtyCoarseEFive}{34.9\,\ensuremath{\pm}\,1.1}  % experiments/FIGS/data/v2/index.json :: fine-tuning qg_herwig N=100000: R30, R16_Q1 :: sha256 c6c4138ee5c86423
\newcommand{\MlstFtQgHerwigRthirtyCoarseOverScratchEFive}{0.950\,\ensuremath{\pm}\,0.033}  % experiments/FIGS/data/v2/index.json :: fine-tuning qg_herwig N=100000: R30 of each R16_Q1 run over the from-scratch mean, SD over runs and standard error of that mean in quadrature :: sha256 c6c4138ee5c86423
\newcommand{\MlstFtQgHerwigRthirtyCoarseOverScratchEFivePr}{equal within 6\%}  % experiments/FIGS/data/v2/index.json :: fine-tuning qg_herwig N=100000: R30 of each R16_Q1 run over the from-scratch mean, SD over runs and standard error of that mean in quadrature :: sha256 c6c4138ee5c86423 :: prose twin: two significant figures, direction in words
\newcommand{\MlstFtQgHerwigRthirtyResMassEFive}{34.0 and 36.5}  % experiments/FIGS/data/v2/index.json :: fine-tuning qg_herwig N=100000: R30, L162_MASS :: sha256 c6c4138ee5c86423
\newcommand{\MlstFtQgHerwigRthirtyResMassOverScratchEFive}{0.925 and 0.994}  % experiments/FIGS/data/v2/index.json :: fine-tuning qg_herwig N=100000: R30 of each L162_MASS run over the from-scratch mean, SD over runs and standard error of that mean in quadrature :: sha256 c6c4138ee5c86423
\newcommand{\MlstFtQgHerwigRthirtyResMassOverScratchEFivePr}{at least 0.63\% less rejection}  % experiments/FIGS/data/v2/index.json :: fine-tuning qg_herwig N=100000: R30 of each L162_MASS run over the from-scratch mean, SD over runs and standard error of that mean in quadrature :: sha256 c6c4138ee5c86423 :: prose twin: two significant figures, direction in words
\newcommand{\MlstFtQgHerwigRthirtyResMassOverScratchEFivePrMax}{up to 7.5\% less rejection}  % experiments/FIGS/data/v2/index.json :: fine-tuning qg_herwig N=100000: R30 of each L162_MASS run over the from-scratch mean, SD over runs and standard error of that mean in quadrature :: sha256 c6c4138ee5c86423 :: prose twin of the run farther from 1
\newcommand{\MlstFtQgHerwigRthirtyCoarseMassEFive}{35.9 and 39.0}  % experiments/FIGS/data/v2/index.json :: fine-tuning qg_herwig N=100000: R30, R16_Q1_MASS :: sha256 c6c4138ee5c86423
\newcommand{\MlstFtQgHerwigRthirtyCoarseMassOverScratchEFive}{0.979 and 1.064}  % experiments/FIGS/data/v2/index.json :: fine-tuning qg_herwig N=100000: R30 of each R16_Q1_MASS run over the from-scratch mean, SD over runs and standard error of that mean in quadrature :: sha256 c6c4138ee5c86423
\newcommand{\MlstFtQgHerwigRthirtyCoarseMassOverScratchEFivePr}{equal within 7\%}  % experiments/FIGS/data/v2/index.json :: fine-tuning qg_herwig N=100000: R30 of each R16_Q1_MASS run over the from-scratch mean, SD over runs and standard error of that mean in quadrature :: sha256 c6c4138ee5c86423 :: prose twin: two significant figures, direction in words
\newcommand{\MlstFtQgHerwigRthirtyCoarseMassOverScratchEFivePrMax}{equal within 7\%}  % experiments/FIGS/data/v2/index.json :: fine-tuning qg_herwig N=100000: R30 of each R16_Q1_MASS run over the from-scratch mean, SD over runs and standard error of that mean in quadrature :: sha256 c6c4138ee5c86423 :: prose twin of the run farther from 1
\newcommand{\MlstFtQgHerwigRthirtyResOverFullEFive}{1.03\,\ensuremath{\pm}\,0.04}  % experiments/FIGS/data/v2/index.json :: fine-tuning qg_herwig N=100000: R30 of L162 over L188, run by run :: sha256 c6c4138ee5c86423
\newcommand{\MlstFtQgHerwigRthirtyResOverFullEFivePr}{equal within 5\%}  % experiments/FIGS/data/v2/index.json :: fine-tuning qg_herwig N=100000: R30 of L162 over L188, run by run :: sha256 c6c4138ee5c86423 :: prose twin: two significant figures, direction in words
\newcommand{\MlstFtQgHerwigRthirtyMidOverFullEFive}{1.03\,\ensuremath{\pm}\,0.05}  % experiments/FIGS/data/v2/index.json :: fine-tuning qg_herwig N=100000: R30 of R42_Q1 over L188, run by run :: sha256 c6c4138ee5c86423
\newcommand{\MlstFtQgHerwigRthirtyMidOverFullEFivePr}{equal within 6\%}  % experiments/FIGS/data/v2/index.json :: fine-tuning qg_herwig N=100000: R30 of R42_Q1 over L188, run by run :: sha256 c6c4138ee5c86423 :: prose twin: two significant figures, direction in words
\newcommand{\MlstFtQgHerwigRthirtyCoarseOverFullEFive}{1.02\,\ensuremath{\pm}\,0.06}  % experiments/FIGS/data/v2/index.json :: fine-tuning qg_herwig N=100000: R30 of R16_Q1 over L188, run by run :: sha256 c6c4138ee5c86423
\newcommand{\MlstFtQgHerwigRthirtyCoarseOverFullEFivePr}{equal within 7\%}  % experiments/FIGS/data/v2/index.json :: fine-tuning qg_herwig N=100000: R30 of R16_Q1 over L188, run by run :: sha256 c6c4138ee5c86423 :: prose twin: two significant figures, direction in words
\newcommand{\MlstFtQgHerwigRthirtyResMassOverResEFive}{0.971 and 1.028}  % experiments/FIGS/data/v2/index.json :: fine-tuning qg_herwig N=100000: R30 of L162_MASS over L162, run by run :: sha256 c6c4138ee5c86423
\newcommand{\MlstFtQgHerwigRthirtyResMassOverResEFivePr}{equal within 3\%}  % experiments/FIGS/data/v2/index.json :: fine-tuning qg_herwig N=100000: R30 of L162_MASS over L162, run by run :: sha256 c6c4138ee5c86423 :: prose twin: two significant figures, direction in words
\newcommand{\MlstFtQgHerwigRthirtyResMassOverResEFivePrMax}{equal within 3\%}  % experiments/FIGS/data/v2/index.json :: fine-tuning qg_herwig N=100000: R30 of L162_MASS over L162, run by run :: sha256 c6c4138ee5c86423 :: prose twin of the run farther from 1
\newcommand{\MlstFtQgHerwigRthirtyCoarseMassOverCoarseEFive}{1.064 and 1.121}  % experiments/FIGS/data/v2/index.json :: fine-tuning qg_herwig N=100000: R30 of R16_Q1_MASS over R16_Q1, run by run :: sha256 c6c4138ee5c86423
\newcommand{\MlstFtQgHerwigRthirtyCoarseMassOverCoarseEFivePr}{equal within 20\%}  % experiments/FIGS/data/v2/index.json :: fine-tuning qg_herwig N=100000: R30 of R16_Q1_MASS over R16_Q1, run by run :: sha256 c6c4138ee5c86423 :: prose twin: two significant figures, direction in words
\newcommand{\MlstFtQgHerwigRthirtyCoarseMassOverCoarseEFivePrMax}{equal within 20\%}  % experiments/FIGS/data/v2/index.json :: fine-tuning qg_herwig N=100000: R30 of R16_Q1_MASS over R16_Q1, run by run :: sha256 c6c4138ee5c86423 :: prose twin of the run farther from 1
\newcommand{\MlstFtQgHerwigRthirtyRandOverFullEFive}{1.01\,\ensuremath{\pm}\,0.05}  % experiments/FIGS/data/v2/index.json :: fine-tuning qg_herwig N=100000: R30 of each random partition over the L188 run of its seed :: sha256 c6c4138ee5c86423
\newcommand{\MlstFtQgHerwigRthirtyRandOverFullEFivePr}{equal within 5\%}  % experiments/FIGS/data/v2/index.json :: fine-tuning qg_herwig N=100000: R30 of each random partition over the L188 run of its seed :: sha256 c6c4138ee5c86423 :: prose twin: two significant figures, direction in words
\newcommand{\MlstFtQgHerwigRthirtyRandOverCoarseEFive}{1.009\,\ensuremath{\pm}\,0.024}  % experiments/FIGS/data/v2/index.json :: fine-tuning qg_herwig N=100000: R30 of each random partition over the R16_Q1 run of its seed :: sha256 c6c4138ee5c86423
\newcommand{\MlstFtQgHerwigRthirtyRandOverCoarseEFivePr}{equal within 3\%}  % experiments/FIGS/data/v2/index.json :: fine-tuning qg_herwig N=100000: R30 of each random partition over the R16_Q1 run of its seed :: sha256 c6c4138ee5c86423 :: prose twin: two significant figures, direction in words
\newcommand{\MlstFtQgHerwigRthirtyScratch}{35.5\,\ensuremath{\pm}\,1.2}  % experiments/FIGS/data/v2/index.json :: fine-tuning qg_herwig N=1600000: R30 from scratch, every seed :: sha256 c6c4138ee5c86423
\newcommand{\MlstFtQgHerwigRthirtyFull}{35.2\,\ensuremath{\pm}\,2.1}  % experiments/FIGS/data/v2/index.json :: fine-tuning qg_herwig N=1600000: R30, L188 :: sha256 c6c4138ee5c86423
\newcommand{\MlstFtQgHerwigRthirtyFullOverScratch}{0.99\,\ensuremath{\pm}\,0.06}  % experiments/FIGS/data/v2/index.json :: fine-tuning qg_herwig N=1600000: R30 of each L188 run over the from-scratch mean, SD over runs and standard error of that mean in quadrature :: sha256 c6c4138ee5c86423
\newcommand{\MlstFtQgHerwigRthirtyFullOverScratchPr}{equal within 6\%}  % experiments/FIGS/data/v2/index.json :: fine-tuning qg_herwig N=1600000: R30 of each L188 run over the from-scratch mean, SD over runs and standard error of that mean in quadrature :: sha256 c6c4138ee5c86423 :: prose twin: two significant figures, direction in words
\newcommand{\MlstFtQgHerwigRthirtyRes}{34.6\,\ensuremath{\pm}\,2.0}  % experiments/FIGS/data/v2/index.json :: fine-tuning qg_herwig N=1600000: R30, L162 :: sha256 c6c4138ee5c86423
\newcommand{\MlstFtQgHerwigRthirtyResOverScratch}{0.97\,\ensuremath{\pm}\,0.06}  % experiments/FIGS/data/v2/index.json :: fine-tuning qg_herwig N=1600000: R30 of each L162 run over the from-scratch mean, SD over runs and standard error of that mean in quadrature :: sha256 c6c4138ee5c86423
\newcommand{\MlstFtQgHerwigRthirtyResOverScratchPr}{equal within 6\%}  % experiments/FIGS/data/v2/index.json :: fine-tuning qg_herwig N=1600000: R30 of each L162 run over the from-scratch mean, SD over runs and standard error of that mean in quadrature :: sha256 c6c4138ee5c86423 :: prose twin: two significant figures, direction in words
\newcommand{\MlstFtQgHerwigRthirtyMid}{36.1\,\ensuremath{\pm}\,2.7}  % experiments/FIGS/data/v2/index.json :: fine-tuning qg_herwig N=1600000: R30, R42_Q1 :: sha256 c6c4138ee5c86423
\newcommand{\MlstFtQgHerwigRthirtyMidOverScratch}{1.02\,\ensuremath{\pm}\,0.08}  % experiments/FIGS/data/v2/index.json :: fine-tuning qg_herwig N=1600000: R30 of each R42_Q1 run over the from-scratch mean, SD over runs and standard error of that mean in quadrature :: sha256 c6c4138ee5c86423
\newcommand{\MlstFtQgHerwigRthirtyMidOverScratchPr}{equal within 8\%}  % experiments/FIGS/data/v2/index.json :: fine-tuning qg_herwig N=1600000: R30 of each R42_Q1 run over the from-scratch mean, SD over runs and standard error of that mean in quadrature :: sha256 c6c4138ee5c86423 :: prose twin: two significant figures, direction in words
\newcommand{\MlstFtQgHerwigRthirtyCoarse}{36.9\,\ensuremath{\pm}\,0.9}  % experiments/FIGS/data/v2/index.json :: fine-tuning qg_herwig N=1600000: R30, R16_Q1 :: sha256 c6c4138ee5c86423
\newcommand{\MlstFtQgHerwigRthirtyCoarseOverScratch}{1.040\,\ensuremath{\pm}\,0.033}  % experiments/FIGS/data/v2/index.json :: fine-tuning qg_herwig N=1600000: R30 of each R16_Q1 run over the from-scratch mean, SD over runs and standard error of that mean in quadrature :: sha256 c6c4138ee5c86423
\newcommand{\MlstFtQgHerwigRthirtyCoarseOverScratchPr}{equal within 4\%}  % experiments/FIGS/data/v2/index.json :: fine-tuning qg_herwig N=1600000: R30 of each R16_Q1 run over the from-scratch mean, SD over runs and standard error of that mean in quadrature :: sha256 c6c4138ee5c86423 :: prose twin: two significant figures, direction in words
\newcommand{\MlstFtQgHerwigRthirtyResMass}{37.8 and 34.4}  % experiments/FIGS/data/v2/index.json :: fine-tuning qg_herwig N=1600000: R30, L162_MASS :: sha256 c6c4138ee5c86423
\newcommand{\MlstFtQgHerwigRthirtyResMassOverScratch}{1.065 and 0.970}  % experiments/FIGS/data/v2/index.json :: fine-tuning qg_herwig N=1600000: R30 of each L162_MASS run over the from-scratch mean, SD over runs and standard error of that mean in quadrature :: sha256 c6c4138ee5c86423
\newcommand{\MlstFtQgHerwigRthirtyResMassOverScratchPr}{equal within 7\%}  % experiments/FIGS/data/v2/index.json :: fine-tuning qg_herwig N=1600000: R30 of each L162_MASS run over the from-scratch mean, SD over runs and standard error of that mean in quadrature :: sha256 c6c4138ee5c86423 :: prose twin: two significant figures, direction in words
\newcommand{\MlstFtQgHerwigRthirtyResMassOverScratchPrMax}{equal within 7\%}  % experiments/FIGS/data/v2/index.json :: fine-tuning qg_herwig N=1600000: R30 of each L162_MASS run over the from-scratch mean, SD over runs and standard error of that mean in quadrature :: sha256 c6c4138ee5c86423 :: prose twin of the run farther from 1
\newcommand{\MlstFtQgHerwigRthirtyCoarseMass}{37.7 and 32.9}  % experiments/FIGS/data/v2/index.json :: fine-tuning qg_herwig N=1600000: R30, R16_Q1_MASS :: sha256 c6c4138ee5c86423
\newcommand{\MlstFtQgHerwigRthirtyCoarseMassOverScratch}{1.06 and 0.93}  % experiments/FIGS/data/v2/index.json :: fine-tuning qg_herwig N=1600000: R30 of each R16_Q1_MASS run over the from-scratch mean, SD over runs and standard error of that mean in quadrature :: sha256 c6c4138ee5c86423
\newcommand{\MlstFtQgHerwigRthirtyCoarseMassOverScratchPr}{equal within 8\%}  % experiments/FIGS/data/v2/index.json :: fine-tuning qg_herwig N=1600000: R30 of each R16_Q1_MASS run over the from-scratch mean, SD over runs and standard error of that mean in quadrature :: sha256 c6c4138ee5c86423 :: prose twin: two significant figures, direction in words
\newcommand{\MlstFtQgHerwigRthirtyCoarseMassOverScratchPrMax}{equal within 8\%}  % experiments/FIGS/data/v2/index.json :: fine-tuning qg_herwig N=1600000: R30 of each R16_Q1_MASS run over the from-scratch mean, SD over runs and standard error of that mean in quadrature :: sha256 c6c4138ee5c86423 :: prose twin of the run farther from 1
\newcommand{\MlstFtQgHerwigRthirtyResOverFull}{0.9832\,\ensuremath{\pm}\,0.0022}  % experiments/FIGS/data/v2/index.json :: fine-tuning qg_herwig N=1600000: R30 of L162 over L188, run by run :: sha256 c6c4138ee5c86423
\newcommand{\MlstFtQgHerwigRthirtyResOverFullPr}{1.7\% less rejection}  % experiments/FIGS/data/v2/index.json :: fine-tuning qg_herwig N=1600000: R30 of L162 over L188, run by run :: sha256 c6c4138ee5c86423 :: prose twin: two significant figures, direction in words
\newcommand{\MlstFtQgHerwigRthirtyMidOverFull}{1.03\,\ensuremath{\pm}\,0.05}  % experiments/FIGS/data/v2/index.json :: fine-tuning qg_herwig N=1600000: R30 of R42_Q1 over L188, run by run :: sha256 c6c4138ee5c86423
\newcommand{\MlstFtQgHerwigRthirtyMidOverFullPr}{equal within 6\%}  % experiments/FIGS/data/v2/index.json :: fine-tuning qg_herwig N=1600000: R30 of R42_Q1 over L188, run by run :: sha256 c6c4138ee5c86423 :: prose twin: two significant figures, direction in words
\newcommand{\MlstFtQgHerwigRthirtyCoarseOverFull}{1.05\,\ensuremath{\pm}\,0.06}  % experiments/FIGS/data/v2/index.json :: fine-tuning qg_herwig N=1600000: R30 of R16_Q1 over L188, run by run :: sha256 c6c4138ee5c86423
\newcommand{\MlstFtQgHerwigRthirtyCoarseOverFullPr}{equal within 6\%}  % experiments/FIGS/data/v2/index.json :: fine-tuning qg_herwig N=1600000: R30 of R16_Q1 over L188, run by run :: sha256 c6c4138ee5c86423 :: prose twin: two significant figures, direction in words
\newcommand{\MlstFtQgHerwigRthirtyResMassOverRes}{1.14 and 0.93}  % experiments/FIGS/data/v2/index.json :: fine-tuning qg_herwig N=1600000: R30 of L162_MASS over L162, run by run :: sha256 c6c4138ee5c86423
\newcommand{\MlstFtQgHerwigRthirtyResMassOverResPr}{equal within 20\%}  % experiments/FIGS/data/v2/index.json :: fine-tuning qg_herwig N=1600000: R30 of L162_MASS over L162, run by run :: sha256 c6c4138ee5c86423 :: prose twin: two significant figures, direction in words
\newcommand{\MlstFtQgHerwigRthirtyResMassOverResPrMax}{equal within 20\%}  % experiments/FIGS/data/v2/index.json :: fine-tuning qg_herwig N=1600000: R30 of L162_MASS over L162, run by run :: sha256 c6c4138ee5c86423 :: prose twin of the run farther from 1
\newcommand{\MlstFtQgHerwigRthirtyCoarseMassOverCoarse}{1.01 and 0.88}  % experiments/FIGS/data/v2/index.json :: fine-tuning qg_herwig N=1600000: R30 of R16_Q1_MASS over R16_Q1, run by run :: sha256 c6c4138ee5c86423
\newcommand{\MlstFtQgHerwigRthirtyCoarseMassOverCoarsePr}{equal within 20\%}  % experiments/FIGS/data/v2/index.json :: fine-tuning qg_herwig N=1600000: R30 of R16_Q1_MASS over R16_Q1, run by run :: sha256 c6c4138ee5c86423 :: prose twin: two significant figures, direction in words
\newcommand{\MlstFtQgHerwigRthirtyCoarseMassOverCoarsePrMax}{equal within 20\%}  % experiments/FIGS/data/v2/index.json :: fine-tuning qg_herwig N=1600000: R30 of R16_Q1_MASS over R16_Q1, run by run :: sha256 c6c4138ee5c86423 :: prose twin of the run farther from 1
\newcommand{\MlstFtQgHerwigRthirtyRandOverFull}{0.98\,\ensuremath{\pm}\,0.04}  % experiments/FIGS/data/v2/index.json :: fine-tuning qg_herwig N=1600000: R30 of each random partition over the L188 run of its seed :: sha256 c6c4138ee5c86423
\newcommand{\MlstFtQgHerwigRthirtyRandOverFullPr}{equal within 5\%}  % experiments/FIGS/data/v2/index.json :: fine-tuning qg_herwig N=1600000: R30 of each random partition over the L188 run of its seed :: sha256 c6c4138ee5c86423 :: prose twin: two significant figures, direction in words
\newcommand{\MlstFtQgHerwigRthirtyRandOverCoarse}{0.94\,\ensuremath{\pm}\,0.04}  % experiments/FIGS/data/v2/index.json :: fine-tuning qg_herwig N=1600000: R30 of each random partition over the R16_Q1 run of its seed :: sha256 c6c4138ee5c86423
\newcommand{\MlstFtQgHerwigRthirtyRandOverCoarsePr}{equal within 6\%}  % experiments/FIGS/data/v2/index.json :: fine-tuning qg_herwig N=1600000: R30 of each random partition over the R16_Q1 run of its seed :: sha256 c6c4138ee5c86423 :: prose twin: two significant figures, direction in words
\newcommand{\MlstFtQgRthirtyHerwigOverPythiaFull}{0.328\,\ensuremath{\pm}\,0.015}  % experiments/FIGS/data/v2/index.json :: fine-tuning q/g N=1600000: R30 on Herwig over R30 on Pythia, each run of L188 :: sha256 c6c4138ee5c86423
\newcommand{\MlstFtQgRthirtyHerwigOverPythiaFullPr}{0.33 times the rejection}  % experiments/FIGS/data/v2/index.json :: fine-tuning q/g N=1600000: R30 on Herwig over R30 on Pythia, each run of L188 :: sha256 c6c4138ee5c86423 :: prose twin: two significant figures, direction in words
\newcommand{\MlstFtQgRthirtyHerwigOverPythiaRes}{0.320\,\ensuremath{\pm}\,0.023}  % experiments/FIGS/data/v2/index.json :: fine-tuning q/g N=1600000: R30 on Herwig over R30 on Pythia, each run of L162 :: sha256 c6c4138ee5c86423
\newcommand{\MlstFtQgRthirtyHerwigOverPythiaResPr}{0.32 times the rejection}  % experiments/FIGS/data/v2/index.json :: fine-tuning q/g N=1600000: R30 on Herwig over R30 on Pythia, each run of L162 :: sha256 c6c4138ee5c86423 :: prose twin: two significant figures, direction in words
\newcommand{\MlstFtQgRthirtyHerwigOverPythiaMid}{0.334\,\ensuremath{\pm}\,0.023}  % experiments/FIGS/data/v2/index.json :: fine-tuning q/g N=1600000: R30 on Herwig over R30 on Pythia, each run of R42_Q1 :: sha256 c6c4138ee5c86423
\newcommand{\MlstFtQgRthirtyHerwigOverPythiaMidPr}{0.33 times the rejection}  % experiments/FIGS/data/v2/index.json :: fine-tuning q/g N=1600000: R30 on Herwig over R30 on Pythia, each run of R42_Q1 :: sha256 c6c4138ee5c86423 :: prose twin: two significant figures, direction in words
\newcommand{\MlstFtQgRthirtyHerwigOverPythiaCoarse}{0.346\,\ensuremath{\pm}\,0.011}  % experiments/FIGS/data/v2/index.json :: fine-tuning q/g N=1600000: R30 on Herwig over R30 on Pythia, each run of R16_Q1 :: sha256 c6c4138ee5c86423
\newcommand{\MlstFtQgRthirtyHerwigOverPythiaCoarsePr}{0.35 times the rejection}  % experiments/FIGS/data/v2/index.json :: fine-tuning q/g N=1600000: R30 on Herwig over R30 on Pythia, each run of R16_Q1 :: sha256 c6c4138ee5c86423 :: prose twin: two significant figures, direction in words
\newcommand{\MlstFtQgRthirtyHerwigOverPythiaResMass}{0.353 and 0.323}  % experiments/FIGS/data/v2/index.json :: fine-tuning q/g N=1600000: R30 on Herwig over R30 on Pythia, each run of L162_MASS :: sha256 c6c4138ee5c86423
\newcommand{\MlstFtQgRthirtyHerwigOverPythiaResMassPr}{at most 0.35 times the rejection}  % experiments/FIGS/data/v2/index.json :: fine-tuning q/g N=1600000: R30 on Herwig over R30 on Pythia, each run of L162_MASS :: sha256 c6c4138ee5c86423 :: prose twin: two significant figures, direction in words
\newcommand{\MlstFtQgRthirtyHerwigOverPythiaResMassPrMax}{down to 0.32 times the rejection}  % experiments/FIGS/data/v2/index.json :: fine-tuning q/g N=1600000: R30 on Herwig over R30 on Pythia, each run of L162_MASS :: sha256 c6c4138ee5c86423 :: prose twin of the run farther from 1
\newcommand{\MlstFtQgRthirtyHerwigOverPythiaCoarseMass}{0.359 and 0.311}  % experiments/FIGS/data/v2/index.json :: fine-tuning q/g N=1600000: R30 on Herwig over R30 on Pythia, each run of R16_Q1_MASS :: sha256 c6c4138ee5c86423
\newcommand{\MlstFtQgRthirtyHerwigOverPythiaCoarseMassPr}{at most 0.36 times the rejection}  % experiments/FIGS/data/v2/index.json :: fine-tuning q/g N=1600000: R30 on Herwig over R30 on Pythia, each run of R16_Q1_MASS :: sha256 c6c4138ee5c86423 :: prose twin: two significant figures, direction in words
\newcommand{\MlstFtQgRthirtyHerwigOverPythiaCoarseMassPrMax}{down to 0.31 times the rejection}  % experiments/FIGS/data/v2/index.json :: fine-tuning q/g N=1600000: R30 on Herwig over R30 on Pythia, each run of R16_Q1_MASS :: sha256 c6c4138ee5c86423 :: prose twin of the run farther from 1
\newcommand{\MlstFtQgRthirtyHerwigOverPythiaScratch}{0.341\,\ensuremath{\pm}\,0.009}  % experiments/FIGS/data/v2/index.json :: fine-tuning q/g N=1600000: R30 on Herwig over R30 on Pythia, each from-scratch seed :: sha256 c6c4138ee5c86423
\newcommand{\MlstFtQgRthirtyHerwigOverPythiaScratchPr}{0.34 times the rejection}  % experiments/FIGS/data/v2/index.json :: fine-tuning q/g N=1600000: R30 on Herwig over R30 on Pythia, each from-scratch seed :: sha256 c6c4138ee5c86423 :: prose twin: two significant figures, direction in words
\newcommand{\MlstTopCheckFullLong}{\pending{recipe check}}  % experiments/FIGS/data/v2/index.json :: the recipe check is not read out yet :: sha256 c6c4138ee5c86423
\newcommand{\MlstTopCheckScratchLong}{\pending{recipe check}}  % experiments/FIGS/data/v2/index.json :: the recipe check is not read out yet :: sha256 c6c4138ee5c86423
\newcommand{\MlstTopCheckFullShort}{\pending{recipe check}}  % experiments/FIGS/data/v2/index.json :: the recipe check is not read out yet :: sha256 c6c4138ee5c86423
\newcommand{\MlstTopCheckScratchShort}{\pending{recipe check}}  % experiments/FIGS/data/v2/index.json :: the recipe check is not read out yet :: sha256 c6c4138ee5c86423
\newcommand{\MlstLpJcOneGapFracEFour}{56.3\,\ensuremath{\pm}\,1.1\%}  % experiments/FIGS/data/v2/index.json :: jc1 N=10000: (probe - scratch) / (fine-tuned - scratch) accuracy, 188-class runs, per cent :: sha256 c6c4138ee5c86423
\newcommand{\MlstLpJcOneGapFracEFourPr}{56\%}  % experiments/FIGS/data/v2/index.json :: jc1 N=10000: (probe - scratch) / (fine-tuned - scratch) accuracy, 188-class runs, per cent :: sha256 c6c4138ee5c86423 :: prose twin, two significant figures
\newcommand{\MlstLpJcOneGapFracEFive}{\ensuremath{-}15\,\ensuremath{\pm}\,5\%}  % experiments/FIGS/data/v2/index.json :: jc1 N=100000: (probe - scratch) / (fine-tuned - scratch) accuracy, 188-class runs, per cent :: sha256 c6c4138ee5c86423
\newcommand{\MlstLpJcOneGapFracEFivePr}{\ensuremath{-}15\%}  % experiments/FIGS/data/v2/index.json :: jc1 N=100000: (probe - scratch) / (fine-tuned - scratch) accuracy, 188-class runs, per cent :: sha256 c6c4138ee5c86423 :: prose twin, two significant figures
\newcommand{\MlstLpJcTwoGapFracEFour}{96.2\,\ensuremath{\pm}\,0.6\%}  % experiments/FIGS/data/v2/index.json :: jc2_188 N=10000: (probe - scratch) / (fine-tuned - scratch) accuracy, 188-class runs, per cent :: sha256 c6c4138ee5c86423
\newcommand{\MlstLpJcTwoGapFracEFourPr}{96\%}  % experiments/FIGS/data/v2/index.json :: jc2_188 N=10000: (probe - scratch) / (fine-tuned - scratch) accuracy, 188-class runs, per cent :: sha256 c6c4138ee5c86423 :: prose twin, two significant figures
\newcommand{\MlstLpJcTwoGapFracEFive}{102.2\,\ensuremath{\pm}\,1.0\%}  % experiments/FIGS/data/v2/index.json :: jc2_188 N=100000: (probe - scratch) / (fine-tuned - scratch) accuracy, 188-class runs, per cent :: sha256 c6c4138ee5c86423
\newcommand{\MlstLpJcTwoGapFracEFivePr}{100\%}  % experiments/FIGS/data/v2/index.json :: jc2_188 N=100000: (probe - scratch) / (fine-tuned - scratch) accuracy, 188-class runs, per cent :: sha256 c6c4138ee5c86423 :: prose twin, two significant figures
\newcommand{\XbxMx}{\pending{samples not in}}  % experiments/XBS/xbx_report.py :: pending :: sha256 cb6dffe774ff5b1e
\newcommand{\XbxWindowLo}{\pending{samples not in}}  % experiments/XBS/xbx_report.py :: pending :: sha256 cb6dffe774ff5b1e
\newcommand{\XbxWindowHi}{\pending{samples not in}}  % experiments/XBS/xbx_report.py :: pending :: sha256 cb6dffe774ff5b1e
\newcommand{\XbbNTrainPerClass}{\pending{samples not in}}  % experiments/XBS/xbx_report.py :: pending :: sha256 cb6dffe774ff5b1e
\newcommand{\XbbRepetition}{\pending{samples not in}}  % experiments/XBS/xbx_report.py :: pending :: sha256 cb6dffe774ff5b1e
\newcommand{\XbbNEvalSignal}{\pending{samples not in}}  % experiments/XBS/xbx_report.py :: pending :: sha256 cb6dffe774ff5b1e
\newcommand{\XbbNEvalQcd}{\pending{samples not in}}  % experiments/XBS/xbx_report.py :: pending :: sha256 cb6dffe774ff5b1e
\newcommand{\XbsNTrainPerClass}{\pending{samples not in}}  % experiments/XBS/xbx_report.py :: pending :: sha256 cb6dffe774ff5b1e
\newcommand{\XbsRepetition}{\pending{samples not in}}  % experiments/XBS/xbx_report.py :: pending :: sha256 cb6dffe774ff5b1e
\newcommand{\XbsNEvalSignal}{\pending{samples not in}}  % experiments/XBS/xbx_report.py :: pending :: sha256 cb6dffe774ff5b1e
\newcommand{\XbsNEvalQcd}{\pending{samples not in}}  % experiments/XBS/xbx_report.py :: pending :: sha256 cb6dffe774ff5b1e
\newcommand{\MlstPartreproRthirtyPart}{1{,}711}  % experiments/FIGS/data/v2/partrepro/part.json :: r30 :: sha256 0d44b997af9e6448 :: ParT reproduction on top, part: R30 (one run)
\newcommand{\MlstPartreproRthirtyHead}{1{,}577}  % experiments/FIGS/data/v2/partrepro/head.json :: r30 :: sha256 37b82ba0a0fe9299 :: ParT reproduction on top, head: R30 (one run)
\newcommand{\MlstPartreproRthirtyRescale}{1{,}224}  % experiments/FIGS/data/v2/partrepro/rescale.json :: r30 :: sha256 0a9f1401047d07ff :: ParT reproduction on top, rescale: R30 (one run)
\newcommand{\MlstPartreproRthirtyHeadOverPart}{0.922}  % experiments/FIGS/data/v2/partrepro/head.json :: r30 / part r30 :: sha256 37b82ba0a0fe9299
\newcommand{\MlstPartreproRthirtyHeadOverPartPr}{7.8\% less rejection}  % experiments/FIGS/data/v2/partrepro/head.json :: r30 / part r30 :: sha256 37b82ba0a0fe9299 :: prose twin: two significant figures, direction in words
\newcommand{\MlstPartreproRthirtyRescaleOverPart}{0.715}  % experiments/FIGS/data/v2/partrepro/rescale.json :: r30 / part r30 :: sha256 0a9f1401047d07ff
\newcommand{\MlstPartreproRthirtyRescaleOverPartPr}{28\% less rejection}  % experiments/FIGS/data/v2/partrepro/rescale.json :: r30 / part r30 :: sha256 0a9f1401047d07ff :: prose twin: two significant figures, direction in words
\newcommand{\MlstPartreproSliceDropRatio}{\pending{jet-pT slice test not run}}  % experiments/FIGS/data/v2/index.json :: open item O10: R30 in narrow pT slices of partrepro part and rescale :: sha256 c6c4138ee5c86423
\newcommand{\MlstInputBnZeroSigmaMax}{0.96}  % experiments/FIGS/data/v2/head/input_bn_offsets.json :: offsets :: sha256 d6b821509fccecff :: largest |offset| over the five type flags
\newcommand{\MlstFtJcOneNSignal}{nine}  % experiments/FIGS/data/v2/curves/full_budget_long_leg2.json :: JC1_COLS (count) :: sha256 588e935996d1006a :: JetClass signal classes
\newcommand{\MlstFtJcOneRejHbbResOverFull}{0.92\,\ensuremath{\pm}\,0.05}  % experiments/FIGS/data/v2/curves/full_budget_long_leg2.json :: cells.<run>.N1000000.s1.jc1.per_class.label_Hbb: rejection of L162 over L188, run by run :: sha256 588e935996d1006a
\newcommand{\MlstFtJcOneRejHbbResOverFullPr}{equal within 9\%}  % experiments/FIGS/data/v2/curves/full_budget_long_leg2.json :: cells.<run>.N1000000.s1.jc1.per_class.label_Hbb: rejection of L162 over L188, run by run :: sha256 588e935996d1006a :: prose twin: two significant figures, direction in words
\newcommand{\MlstFtJcOneRejHbbMidOverFull}{0.98\,\ensuremath{\pm}\,0.12}  % experiments/FIGS/data/v2/curves/full_budget_long_leg2.json :: cells.<run>.N1000000.s1.jc1.per_class.label_Hbb: rejection of R42_Q1 over L188, run by run :: sha256 588e935996d1006a
\newcommand{\MlstFtJcOneRejHbbMidOverFullPr}{equal within 20\%}  % experiments/FIGS/data/v2/curves/full_budget_long_leg2.json :: cells.<run>.N1000000.s1.jc1.per_class.label_Hbb: rejection of R42_Q1 over L188, run by run :: sha256 588e935996d1006a :: prose twin: two significant figures, direction in words
\newcommand{\MlstFtJcOneRejHbbCoarseOverFull}{1.01\,\ensuremath{\pm}\,0.07}  % experiments/FIGS/data/v2/curves/full_budget_long_leg2.json :: cells.<run>.N1000000.s1.jc1.per_class.label_Hbb: rejection of R16_Q1 over L188, run by run :: sha256 588e935996d1006a
\newcommand{\MlstFtJcOneRejHbbCoarseOverFullPr}{equal within 7\%}  % experiments/FIGS/data/v2/curves/full_budget_long_leg2.json :: cells.<run>.N1000000.s1.jc1.per_class.label_Hbb: rejection of R16_Q1 over L188, run by run :: sha256 588e935996d1006a :: prose twin: two significant figures, direction in words
\newcommand{\MlstFtJcOneRejHbbResMassOverRes}{1.065 and 1.100}  % experiments/FIGS/data/v2/curves/full_budget_long_leg2.json :: cells.<run>.N1000000.s1.jc1.per_class.label_Hbb: rejection of L162_MASS over L162, run by run :: sha256 588e935996d1006a
\newcommand{\MlstFtJcOneRejHbbResMassOverResPr}{equal within 10\%}  % experiments/FIGS/data/v2/curves/full_budget_long_leg2.json :: cells.<run>.N1000000.s1.jc1.per_class.label_Hbb: rejection of L162_MASS over L162, run by run :: sha256 588e935996d1006a :: prose twin: two significant figures, direction in words
\newcommand{\MlstFtJcOneRejHbbResMassOverResPrMax}{equal within 10\%}  % experiments/FIGS/data/v2/curves/full_budget_long_leg2.json :: cells.<run>.N1000000.s1.jc1.per_class.label_Hbb: rejection of L162_MASS over L162, run by run :: sha256 588e935996d1006a :: prose twin of the run farther from 1
\newcommand{\MlstFtJcOneRejHbbCoarseMassOverCoarse}{0.97 and 0.86}  % experiments/FIGS/data/v2/curves/full_budget_long_leg2.json :: cells.<run>.N1000000.s1.jc1.per_class.label_Hbb: rejection of R16_Q1_MASS over R16_Q1, run by run :: sha256 588e935996d1006a
\newcommand{\MlstFtJcOneRejHbbCoarseMassOverCoarsePr}{equal within 20\%}  % experiments/FIGS/data/v2/curves/full_budget_long_leg2.json :: cells.<run>.N1000000.s1.jc1.per_class.label_Hbb: rejection of R16_Q1_MASS over R16_Q1, run by run :: sha256 588e935996d1006a :: prose twin: two significant figures, direction in words
\newcommand{\MlstFtJcOneRejHbbCoarseMassOverCoarsePrMax}{equal within 20\%}  % experiments/FIGS/data/v2/curves/full_budget_long_leg2.json :: cells.<run>.N1000000.s1.jc1.per_class.label_Hbb: rejection of R16_Q1_MASS over R16_Q1, run by run :: sha256 588e935996d1006a :: prose twin of the run farther from 1
\newcommand{\MlstFtJcOneRejHbbRandOverFull}{0.92\,\ensuremath{\pm}\,0.11}  % experiments/FIGS/data/v2/curves/full_budget_long_leg2.json :: per_class.label_Hbb: rejection of each random partition over the 188-class run of its seed :: sha256 588e935996d1006a
\newcommand{\MlstFtJcOneRejHbbRandOverFullPr}{equal within 20\%}  % experiments/FIGS/data/v2/curves/full_budget_long_leg2.json :: per_class.label_Hbb: rejection of each random partition over the 188-class run of its seed :: sha256 588e935996d1006a :: prose twin: two significant figures, direction in words
\newcommand{\MlstFtJcOneRejHbbRandOverFullRange}{0.84--1.1}  % experiments/FIGS/data/v2/curves/full_budget_long_leg2.json :: per_class.label_Hbb: smallest and largest of the partitions' ratios :: sha256 588e935996d1006a
\newcommand{\MlstFtJcOneRejHccResOverFull}{1.07\,\ensuremath{\pm}\,0.12}  % experiments/FIGS/data/v2/curves/full_budget_long_leg2.json :: cells.<run>.N1000000.s1.jc1.per_class.label_Hcc: rejection of L162 over L188, run by run :: sha256 588e935996d1006a
\newcommand{\MlstFtJcOneRejHccResOverFullPr}{equal within 20\%}  % experiments/FIGS/data/v2/curves/full_budget_long_leg2.json :: cells.<run>.N1000000.s1.jc1.per_class.label_Hcc: rejection of L162 over L188, run by run :: sha256 588e935996d1006a :: prose twin: two significant figures, direction in words
\newcommand{\MlstFtJcOneRejHccMidOverFull}{0.99\,\ensuremath{\pm}\,0.10}  % experiments/FIGS/data/v2/curves/full_budget_long_leg2.json :: cells.<run>.N1000000.s1.jc1.per_class.label_Hcc: rejection of R42_Q1 over L188, run by run :: sha256 588e935996d1006a
\newcommand{\MlstFtJcOneRejHccMidOverFullPr}{equal within 20\%}  % experiments/FIGS/data/v2/curves/full_budget_long_leg2.json :: cells.<run>.N1000000.s1.jc1.per_class.label_Hcc: rejection of R42_Q1 over L188, run by run :: sha256 588e935996d1006a :: prose twin: two significant figures, direction in words
\newcommand{\MlstFtJcOneRejHccCoarseOverFull}{0.91\,\ensuremath{\pm}\,0.07}  % experiments/FIGS/data/v2/curves/full_budget_long_leg2.json :: cells.<run>.N1000000.s1.jc1.per_class.label_Hcc: rejection of R16_Q1 over L188, run by run :: sha256 588e935996d1006a
\newcommand{\MlstFtJcOneRejHccCoarseOverFullPr}{equal within 10\%}  % experiments/FIGS/data/v2/curves/full_budget_long_leg2.json :: cells.<run>.N1000000.s1.jc1.per_class.label_Hcc: rejection of R16_Q1 over L188, run by run :: sha256 588e935996d1006a :: prose twin: two significant figures, direction in words
\newcommand{\MlstFtJcOneRejHccResMassOverRes}{0.886 and 0.816}  % experiments/FIGS/data/v2/curves/full_budget_long_leg2.json :: cells.<run>.N1000000.s1.jc1.per_class.label_Hcc: rejection of L162_MASS over L162, run by run :: sha256 588e935996d1006a
\newcommand{\MlstFtJcOneRejHccResMassOverResPr}{equal within 20\%}  % experiments/FIGS/data/v2/curves/full_budget_long_leg2.json :: cells.<run>.N1000000.s1.jc1.per_class.label_Hcc: rejection of L162_MASS over L162, run by run :: sha256 588e935996d1006a :: prose twin: two significant figures, direction in words
\newcommand{\MlstFtJcOneRejHccResMassOverResPrMax}{equal within 20\%}  % experiments/FIGS/data/v2/curves/full_budget_long_leg2.json :: cells.<run>.N1000000.s1.jc1.per_class.label_Hcc: rejection of L162_MASS over L162, run by run :: sha256 588e935996d1006a :: prose twin of the run farther from 1
\newcommand{\MlstFtJcOneRejHccCoarseMassOverCoarse}{0.98 and 0.85}  % experiments/FIGS/data/v2/curves/full_budget_long_leg2.json :: cells.<run>.N1000000.s1.jc1.per_class.label_Hcc: rejection of R16_Q1_MASS over R16_Q1, run by run :: sha256 588e935996d1006a
\newcommand{\MlstFtJcOneRejHccCoarseMassOverCoarsePr}{equal within 20\%}  % experiments/FIGS/data/v2/curves/full_budget_long_leg2.json :: cells.<run>.N1000000.s1.jc1.per_class.label_Hcc: rejection of R16_Q1_MASS over R16_Q1, run by run :: sha256 588e935996d1006a :: prose twin: two significant figures, direction in words
\newcommand{\MlstFtJcOneRejHccCoarseMassOverCoarsePrMax}{equal within 20\%}  % experiments/FIGS/data/v2/curves/full_budget_long_leg2.json :: cells.<run>.N1000000.s1.jc1.per_class.label_Hcc: rejection of R16_Q1_MASS over R16_Q1, run by run :: sha256 588e935996d1006a :: prose twin of the run farther from 1
\newcommand{\MlstFtJcOneRejHccRandOverFull}{1.03\,\ensuremath{\pm}\,0.07}  % experiments/FIGS/data/v2/curves/full_budget_long_leg2.json :: per_class.label_Hcc: rejection of each random partition over the 188-class run of its seed :: sha256 588e935996d1006a
\newcommand{\MlstFtJcOneRejHccRandOverFullPr}{equal within 8\%}  % experiments/FIGS/data/v2/curves/full_budget_long_leg2.json :: per_class.label_Hcc: rejection of each random partition over the 188-class run of its seed :: sha256 588e935996d1006a :: prose twin: two significant figures, direction in words
\newcommand{\MlstFtJcOneRejHccRandOverFullRange}{0.94--1.1}  % experiments/FIGS/data/v2/curves/full_budget_long_leg2.json :: per_class.label_Hcc: smallest and largest of the partitions' ratios :: sha256 588e935996d1006a
\newcommand{\MlstFtJcOneRejHggResOverFull}{0.98\,\ensuremath{\pm}\,0.05}  % experiments/FIGS/data/v2/curves/full_budget_long_leg2.json :: cells.<run>.N1000000.s1.jc1.per_class.label_Hgg: rejection of L162 over L188, run by run :: sha256 588e935996d1006a
\newcommand{\MlstFtJcOneRejHggResOverFullPr}{equal within 6\%}  % experiments/FIGS/data/v2/curves/full_budget_long_leg2.json :: cells.<run>.N1000000.s1.jc1.per_class.label_Hgg: rejection of L162 over L188, run by run :: sha256 588e935996d1006a :: prose twin: two significant figures, direction in words
\newcommand{\MlstFtJcOneRejHggMidOverFull}{1.000\,\ensuremath{\pm}\,0.025}  % experiments/FIGS/data/v2/curves/full_budget_long_leg2.json :: cells.<run>.N1000000.s1.jc1.per_class.label_Hgg: rejection of R42_Q1 over L188, run by run :: sha256 588e935996d1006a
\newcommand{\MlstFtJcOneRejHggMidOverFullPr}{equal within 3\%}  % experiments/FIGS/data/v2/curves/full_budget_long_leg2.json :: cells.<run>.N1000000.s1.jc1.per_class.label_Hgg: rejection of R42_Q1 over L188, run by run :: sha256 588e935996d1006a :: prose twin: two significant figures, direction in words
\newcommand{\MlstFtJcOneRejHggCoarseOverFull}{0.97\,\ensuremath{\pm}\,0.09}  % experiments/FIGS/data/v2/curves/full_budget_long_leg2.json :: cells.<run>.N1000000.s1.jc1.per_class.label_Hgg: rejection of R16_Q1 over L188, run by run :: sha256 588e935996d1006a
\newcommand{\MlstFtJcOneRejHggCoarseOverFullPr}{equal within 10\%}  % experiments/FIGS/data/v2/curves/full_budget_long_leg2.json :: cells.<run>.N1000000.s1.jc1.per_class.label_Hgg: rejection of R16_Q1 over L188, run by run :: sha256 588e935996d1006a :: prose twin: two significant figures, direction in words
\newcommand{\MlstFtJcOneRejHggResMassOverRes}{1.046 and 1.008}  % experiments/FIGS/data/v2/curves/full_budget_long_leg2.json :: cells.<run>.N1000000.s1.jc1.per_class.label_Hgg: rejection of L162_MASS over L162, run by run :: sha256 588e935996d1006a
\newcommand{\MlstFtJcOneRejHggResMassOverResPr}{equal within 5\%}  % experiments/FIGS/data/v2/curves/full_budget_long_leg2.json :: cells.<run>.N1000000.s1.jc1.per_class.label_Hgg: rejection of L162_MASS over L162, run by run :: sha256 588e935996d1006a :: prose twin: two significant figures, direction in words
\newcommand{\MlstFtJcOneRejHggResMassOverResPrMax}{equal within 5\%}  % experiments/FIGS/data/v2/curves/full_budget_long_leg2.json :: cells.<run>.N1000000.s1.jc1.per_class.label_Hgg: rejection of L162_MASS over L162, run by run :: sha256 588e935996d1006a :: prose twin of the run farther from 1
\newcommand{\MlstFtJcOneRejHggCoarseMassOverCoarse}{1.12 and 1.01}  % experiments/FIGS/data/v2/curves/full_budget_long_leg2.json :: cells.<run>.N1000000.s1.jc1.per_class.label_Hgg: rejection of R16_Q1_MASS over R16_Q1, run by run :: sha256 588e935996d1006a
\newcommand{\MlstFtJcOneRejHggCoarseMassOverCoarsePr}{equal within 20\%}  % experiments/FIGS/data/v2/curves/full_budget_long_leg2.json :: cells.<run>.N1000000.s1.jc1.per_class.label_Hgg: rejection of R16_Q1_MASS over R16_Q1, run by run :: sha256 588e935996d1006a :: prose twin: two significant figures, direction in words
\newcommand{\MlstFtJcOneRejHggCoarseMassOverCoarsePrMax}{equal within 20\%}  % experiments/FIGS/data/v2/curves/full_budget_long_leg2.json :: cells.<run>.N1000000.s1.jc1.per_class.label_Hgg: rejection of R16_Q1_MASS over R16_Q1, run by run :: sha256 588e935996d1006a :: prose twin of the run farther from 1
\newcommand{\MlstFtJcOneRejHggRandOverFull}{0.968\,\ensuremath{\pm}\,0.016}  % experiments/FIGS/data/v2/curves/full_budget_long_leg2.json :: per_class.label_Hgg: rejection of each random partition over the 188-class run of its seed :: sha256 588e935996d1006a
\newcommand{\MlstFtJcOneRejHggRandOverFullPr}{equal within 4\%}  % experiments/FIGS/data/v2/curves/full_budget_long_leg2.json :: per_class.label_Hgg: rejection of each random partition over the 188-class run of its seed :: sha256 588e935996d1006a :: prose twin: two significant figures, direction in words
\newcommand{\MlstFtJcOneRejHggRandOverFullRange}{0.94--0.98}  % experiments/FIGS/data/v2/curves/full_budget_long_leg2.json :: per_class.label_Hgg: smallest and largest of the partitions' ratios :: sha256 588e935996d1006a
\newcommand{\MlstFtJcOneRejHfourqResOverFull}{1.05\,\ensuremath{\pm}\,0.06}  % experiments/FIGS/data/v2/curves/full_budget_long_leg2.json :: cells.<run>.N1000000.s1.jc1.per_class.label_H4q: rejection of L162 over L188, run by run :: sha256 588e935996d1006a
\newcommand{\MlstFtJcOneRejHfourqResOverFullPr}{equal within 7\%}  % experiments/FIGS/data/v2/curves/full_budget_long_leg2.json :: cells.<run>.N1000000.s1.jc1.per_class.label_H4q: rejection of L162 over L188, run by run :: sha256 588e935996d1006a :: prose twin: two significant figures, direction in words
\newcommand{\MlstFtJcOneRejHfourqMidOverFull}{0.95\,\ensuremath{\pm}\,0.04}  % experiments/FIGS/data/v2/curves/full_budget_long_leg2.json :: cells.<run>.N1000000.s1.jc1.per_class.label_H4q: rejection of R42_Q1 over L188, run by run :: sha256 588e935996d1006a
\newcommand{\MlstFtJcOneRejHfourqMidOverFullPr}{equal within 6\%}  % experiments/FIGS/data/v2/curves/full_budget_long_leg2.json :: cells.<run>.N1000000.s1.jc1.per_class.label_H4q: rejection of R42_Q1 over L188, run by run :: sha256 588e935996d1006a :: prose twin: two significant figures, direction in words
\newcommand{\MlstFtJcOneRejHfourqCoarseOverFull}{0.962\,\ensuremath{\pm}\,0.025}  % experiments/FIGS/data/v2/curves/full_budget_long_leg2.json :: cells.<run>.N1000000.s1.jc1.per_class.label_H4q: rejection of R16_Q1 over L188, run by run :: sha256 588e935996d1006a
\newcommand{\MlstFtJcOneRejHfourqCoarseOverFullPr}{equal within 4\%}  % experiments/FIGS/data/v2/curves/full_budget_long_leg2.json :: cells.<run>.N1000000.s1.jc1.per_class.label_H4q: rejection of R16_Q1 over L188, run by run :: sha256 588e935996d1006a :: prose twin: two significant figures, direction in words
\newcommand{\MlstFtJcOneRejHfourqResMassOverRes}{0.912 and 0.966}  % experiments/FIGS/data/v2/curves/full_budget_long_leg2.json :: cells.<run>.N1000000.s1.jc1.per_class.label_H4q: rejection of L162_MASS over L162, run by run :: sha256 588e935996d1006a
\newcommand{\MlstFtJcOneRejHfourqResMassOverResPr}{equal within 9\%}  % experiments/FIGS/data/v2/curves/full_budget_long_leg2.json :: cells.<run>.N1000000.s1.jc1.per_class.label_H4q: rejection of L162_MASS over L162, run by run :: sha256 588e935996d1006a :: prose twin: two significant figures, direction in words
\newcommand{\MlstFtJcOneRejHfourqResMassOverResPrMax}{equal within 9\%}  % experiments/FIGS/data/v2/curves/full_budget_long_leg2.json :: cells.<run>.N1000000.s1.jc1.per_class.label_H4q: rejection of L162_MASS over L162, run by run :: sha256 588e935996d1006a :: prose twin of the run farther from 1
\newcommand{\MlstFtJcOneRejHfourqCoarseMassOverCoarse}{0.937 and 0.888}  % experiments/FIGS/data/v2/curves/full_budget_long_leg2.json :: cells.<run>.N1000000.s1.jc1.per_class.label_H4q: rejection of R16_Q1_MASS over R16_Q1, run by run :: sha256 588e935996d1006a
\newcommand{\MlstFtJcOneRejHfourqCoarseMassOverCoarsePr}{equal within 20\%}  % experiments/FIGS/data/v2/curves/full_budget_long_leg2.json :: cells.<run>.N1000000.s1.jc1.per_class.label_H4q: rejection of R16_Q1_MASS over R16_Q1, run by run :: sha256 588e935996d1006a :: prose twin: two significant figures, direction in words
\newcommand{\MlstFtJcOneRejHfourqCoarseMassOverCoarsePrMax}{equal within 20\%}  % experiments/FIGS/data/v2/curves/full_budget_long_leg2.json :: cells.<run>.N1000000.s1.jc1.per_class.label_H4q: rejection of R16_Q1_MASS over R16_Q1, run by run :: sha256 588e935996d1006a :: prose twin of the run farther from 1
\newcommand{\MlstFtJcOneRejHfourqRandOverFull}{0.93\,\ensuremath{\pm}\,0.04}  % experiments/FIGS/data/v2/curves/full_budget_long_leg2.json :: per_class.label_H4q: rejection of each random partition over the 188-class run of its seed :: sha256 588e935996d1006a
\newcommand{\MlstFtJcOneRejHfourqRandOverFullPr}{equal within 8\%}  % experiments/FIGS/data/v2/curves/full_budget_long_leg2.json :: per_class.label_H4q: rejection of each random partition over the 188-class run of its seed :: sha256 588e935996d1006a :: prose twin: two significant figures, direction in words
\newcommand{\MlstFtJcOneRejHfourqRandOverFullRange}{0.86--0.96}  % experiments/FIGS/data/v2/curves/full_budget_long_leg2.json :: per_class.label_H4q: smallest and largest of the partitions' ratios :: sha256 588e935996d1006a
\newcommand{\MlstFtJcOneRejHqqlResOverFull}{1.004\,\ensuremath{\pm}\,0.018}  % experiments/FIGS/data/v2/curves/full_budget_long_leg2.json :: cells.<run>.N1000000.s1.jc1.per_class.label_Hqql: rejection of L162 over L188, run by run :: sha256 588e935996d1006a
\newcommand{\MlstFtJcOneRejHqqlResOverFullPr}{equal within 2\%}  % experiments/FIGS/data/v2/curves/full_budget_long_leg2.json :: cells.<run>.N1000000.s1.jc1.per_class.label_Hqql: rejection of L162 over L188, run by run :: sha256 588e935996d1006a :: prose twin: two significant figures, direction in words
\newcommand{\MlstFtJcOneRejHqqlMidOverFull}{1.04\,\ensuremath{\pm}\,0.18}  % experiments/FIGS/data/v2/curves/full_budget_long_leg2.json :: cells.<run>.N1000000.s1.jc1.per_class.label_Hqql: rejection of R42_Q1 over L188, run by run :: sha256 588e935996d1006a
\newcommand{\MlstFtJcOneRejHqqlMidOverFullPr}{equal within 20\%}  % experiments/FIGS/data/v2/curves/full_budget_long_leg2.json :: cells.<run>.N1000000.s1.jc1.per_class.label_Hqql: rejection of R42_Q1 over L188, run by run :: sha256 588e935996d1006a :: prose twin: two significant figures, direction in words
\newcommand{\MlstFtJcOneRejHqqlCoarseOverFull}{0.93\,\ensuremath{\pm}\,0.16}  % experiments/FIGS/data/v2/curves/full_budget_long_leg2.json :: cells.<run>.N1000000.s1.jc1.per_class.label_Hqql: rejection of R16_Q1 over L188, run by run :: sha256 588e935996d1006a
\newcommand{\MlstFtJcOneRejHqqlCoarseOverFullPr}{equal within 20\%}  % experiments/FIGS/data/v2/curves/full_budget_long_leg2.json :: cells.<run>.N1000000.s1.jc1.per_class.label_Hqql: rejection of R16_Q1 over L188, run by run :: sha256 588e935996d1006a :: prose twin: two significant figures, direction in words
\newcommand{\MlstFtJcOneRejHqqlResMassOverRes}{0.986 and 1.000}  % experiments/FIGS/data/v2/curves/full_budget_long_leg2.json :: cells.<run>.N1000000.s1.jc1.per_class.label_Hqql: rejection of L162_MASS over L162, run by run :: sha256 588e935996d1006a
\newcommand{\MlstFtJcOneRejHqqlResMassOverResPr}{equal within 2\%}  % experiments/FIGS/data/v2/curves/full_budget_long_leg2.json :: cells.<run>.N1000000.s1.jc1.per_class.label_Hqql: rejection of L162_MASS over L162, run by run :: sha256 588e935996d1006a :: prose twin: two significant figures, direction in words
\newcommand{\MlstFtJcOneRejHqqlResMassOverResPrMax}{equal within 2\%}  % experiments/FIGS/data/v2/curves/full_budget_long_leg2.json :: cells.<run>.N1000000.s1.jc1.per_class.label_Hqql: rejection of L162_MASS over L162, run by run :: sha256 588e935996d1006a :: prose twin of the run farther from 1
\newcommand{\MlstFtJcOneRejHqqlCoarseMassOverCoarse}{0.97 and 1.09}  % experiments/FIGS/data/v2/curves/full_budget_long_leg2.json :: cells.<run>.N1000000.s1.jc1.per_class.label_Hqql: rejection of R16_Q1_MASS over R16_Q1, run by run :: sha256 588e935996d1006a
\newcommand{\MlstFtJcOneRejHqqlCoarseMassOverCoarsePr}{equal within 9\%}  % experiments/FIGS/data/v2/curves/full_budget_long_leg2.json :: cells.<run>.N1000000.s1.jc1.per_class.label_Hqql: rejection of R16_Q1_MASS over R16_Q1, run by run :: sha256 588e935996d1006a :: prose twin: two significant figures, direction in words
\newcommand{\MlstFtJcOneRejHqqlCoarseMassOverCoarsePrMax}{equal within 9\%}  % experiments/FIGS/data/v2/curves/full_budget_long_leg2.json :: cells.<run>.N1000000.s1.jc1.per_class.label_Hqql: rejection of R16_Q1_MASS over R16_Q1, run by run :: sha256 588e935996d1006a :: prose twin of the run farther from 1
\newcommand{\MlstFtJcOneRejHqqlRandOverFull}{0.94\,\ensuremath{\pm}\,0.12}  % experiments/FIGS/data/v2/curves/full_budget_long_leg2.json :: per_class.label_Hqql: rejection of each random partition over the 188-class run of its seed :: sha256 588e935996d1006a
\newcommand{\MlstFtJcOneRejHqqlRandOverFullPr}{equal within 20\%}  % experiments/FIGS/data/v2/curves/full_budget_long_leg2.json :: per_class.label_Hqql: rejection of each random partition over the 188-class run of its seed :: sha256 588e935996d1006a :: prose twin: two significant figures, direction in words
\newcommand{\MlstFtJcOneRejHqqlRandOverFullRange}{0.79--1.1}  % experiments/FIGS/data/v2/curves/full_budget_long_leg2.json :: per_class.label_Hqql: smallest and largest of the partitions' ratios :: sha256 588e935996d1006a
\newcommand{\MlstFtJcOneRejTbqqResOverFull}{1.1\,\ensuremath{\pm}\,0.4}  % experiments/FIGS/data/v2/curves/full_budget_long_leg2.json :: cells.<run>.N1000000.s1.jc1.per_class.label_Tbqq: rejection of L162 over L188, run by run :: sha256 588e935996d1006a
\newcommand{\MlstFtJcOneRejTbqqResOverFullPr}{equal within 50\%}  % experiments/FIGS/data/v2/curves/full_budget_long_leg2.json :: cells.<run>.N1000000.s1.jc1.per_class.label_Tbqq: rejection of L162 over L188, run by run :: sha256 588e935996d1006a :: prose twin: two significant figures, direction in words
\newcommand{\MlstFtJcOneRejTbqqMidOverFull}{1.10\,\ensuremath{\pm}\,0.05}  % experiments/FIGS/data/v2/curves/full_budget_long_leg2.json :: cells.<run>.N1000000.s1.jc1.per_class.label_Tbqq: rejection of R42_Q1 over L188, run by run :: sha256 588e935996d1006a
\newcommand{\MlstFtJcOneRejTbqqMidOverFullPr}{equal within 10\%}  % experiments/FIGS/data/v2/curves/full_budget_long_leg2.json :: cells.<run>.N1000000.s1.jc1.per_class.label_Tbqq: rejection of R42_Q1 over L188, run by run :: sha256 588e935996d1006a :: prose twin: two significant figures, direction in words
\newcommand{\MlstFtJcOneRejTbqqCoarseOverFull}{0.85\,\ensuremath{\pm}\,0.32}  % experiments/FIGS/data/v2/curves/full_budget_long_leg2.json :: cells.<run>.N1000000.s1.jc1.per_class.label_Tbqq: rejection of R16_Q1 over L188, run by run :: sha256 588e935996d1006a
\newcommand{\MlstFtJcOneRejTbqqCoarseOverFullPr}{equal within 40\%}  % experiments/FIGS/data/v2/curves/full_budget_long_leg2.json :: cells.<run>.N1000000.s1.jc1.per_class.label_Tbqq: rejection of R16_Q1 over L188, run by run :: sha256 588e935996d1006a :: prose twin: two significant figures, direction in words
\newcommand{\MlstFtJcOneRejTbqqResMassOverRes}{0.89 and 1.21}  % experiments/FIGS/data/v2/curves/full_budget_long_leg2.json :: cells.<run>.N1000000.s1.jc1.per_class.label_Tbqq: rejection of L162_MASS over L162, run by run :: sha256 588e935996d1006a
\newcommand{\MlstFtJcOneRejTbqqResMassOverResPr}{equal within 30\%}  % experiments/FIGS/data/v2/curves/full_budget_long_leg2.json :: cells.<run>.N1000000.s1.jc1.per_class.label_Tbqq: rejection of L162_MASS over L162, run by run :: sha256 588e935996d1006a :: prose twin: two significant figures, direction in words
\newcommand{\MlstFtJcOneRejTbqqResMassOverResPrMax}{equal within 30\%}  % experiments/FIGS/data/v2/curves/full_budget_long_leg2.json :: cells.<run>.N1000000.s1.jc1.per_class.label_Tbqq: rejection of L162_MASS over L162, run by run :: sha256 588e935996d1006a :: prose twin of the run farther from 1
\newcommand{\MlstFtJcOneRejTbqqCoarseMassOverCoarse}{1.00 and 1.79}  % experiments/FIGS/data/v2/curves/full_budget_long_leg2.json :: cells.<run>.N1000000.s1.jc1.per_class.label_Tbqq: rejection of R16_Q1_MASS over R16_Q1, run by run :: sha256 588e935996d1006a
\newcommand{\MlstFtJcOneRejTbqqCoarseMassOverCoarsePr}{equal within 80\%}  % experiments/FIGS/data/v2/curves/full_budget_long_leg2.json :: cells.<run>.N1000000.s1.jc1.per_class.label_Tbqq: rejection of R16_Q1_MASS over R16_Q1, run by run :: sha256 588e935996d1006a :: prose twin: two significant figures, direction in words
\newcommand{\MlstFtJcOneRejTbqqCoarseMassOverCoarsePrMax}{equal within 80\%}  % experiments/FIGS/data/v2/curves/full_budget_long_leg2.json :: cells.<run>.N1000000.s1.jc1.per_class.label_Tbqq: rejection of R16_Q1_MASS over R16_Q1, run by run :: sha256 588e935996d1006a :: prose twin of the run farther from 1
\newcommand{\MlstFtJcOneRejTbqqRandOverFull}{0.79\,\ensuremath{\pm}\,0.15}  % experiments/FIGS/data/v2/curves/full_budget_long_leg2.json :: per_class.label_Tbqq: rejection of each random partition over the 188-class run of its seed :: sha256 588e935996d1006a
\newcommand{\MlstFtJcOneRejTbqqRandOverFullPr}{equal within 30\%}  % experiments/FIGS/data/v2/curves/full_budget_long_leg2.json :: per_class.label_Tbqq: rejection of each random partition over the 188-class run of its seed :: sha256 588e935996d1006a :: prose twin: two significant figures, direction in words
\newcommand{\MlstFtJcOneRejTbqqRandOverFullRange}{0.64--1.0}  % experiments/FIGS/data/v2/curves/full_budget_long_leg2.json :: per_class.label_Tbqq: smallest and largest of the partitions' ratios :: sha256 588e935996d1006a
\newcommand{\MlstFtJcOneRejTblResOverFull}{0.85\,\ensuremath{\pm}\,0.17}  % experiments/FIGS/data/v2/curves/full_budget_long_leg2.json :: cells.<run>.N1000000.s1.jc1.per_class.label_Tbl: rejection of L162 over L188, run by run :: sha256 588e935996d1006a
\newcommand{\MlstFtJcOneRejTblResOverFullPr}{equal within 20\%}  % experiments/FIGS/data/v2/curves/full_budget_long_leg2.json :: cells.<run>.N1000000.s1.jc1.per_class.label_Tbl: rejection of L162 over L188, run by run :: sha256 588e935996d1006a :: prose twin: two significant figures, direction in words
\newcommand{\MlstFtJcOneRejTblMidOverFull}{0.91\,\ensuremath{\pm}\,0.15}  % experiments/FIGS/data/v2/curves/full_budget_long_leg2.json :: cells.<run>.N1000000.s1.jc1.per_class.label_Tbl: rejection of R42_Q1 over L188, run by run :: sha256 588e935996d1006a
\newcommand{\MlstFtJcOneRejTblMidOverFullPr}{equal within 20\%}  % experiments/FIGS/data/v2/curves/full_budget_long_leg2.json :: cells.<run>.N1000000.s1.jc1.per_class.label_Tbl: rejection of R42_Q1 over L188, run by run :: sha256 588e935996d1006a :: prose twin: two significant figures, direction in words
\newcommand{\MlstFtJcOneRejTblCoarseOverFull}{0.94\,\ensuremath{\pm}\,0.28}  % experiments/FIGS/data/v2/curves/full_budget_long_leg2.json :: cells.<run>.N1000000.s1.jc1.per_class.label_Tbl: rejection of R16_Q1 over L188, run by run :: sha256 588e935996d1006a
\newcommand{\MlstFtJcOneRejTblCoarseOverFullPr}{equal within 30\%}  % experiments/FIGS/data/v2/curves/full_budget_long_leg2.json :: cells.<run>.N1000000.s1.jc1.per_class.label_Tbl: rejection of R16_Q1 over L188, run by run :: sha256 588e935996d1006a :: prose twin: two significant figures, direction in words
\newcommand{\MlstFtJcOneRejTblResMassOverRes}{1.04 and 0.93}  % experiments/FIGS/data/v2/curves/full_budget_long_leg2.json :: cells.<run>.N1000000.s1.jc1.per_class.label_Tbl: rejection of L162_MASS over L162, run by run :: sha256 588e935996d1006a
\newcommand{\MlstFtJcOneRejTblResMassOverResPr}{equal within 7\%}  % experiments/FIGS/data/v2/curves/full_budget_long_leg2.json :: cells.<run>.N1000000.s1.jc1.per_class.label_Tbl: rejection of L162_MASS over L162, run by run :: sha256 588e935996d1006a :: prose twin: two significant figures, direction in words
\newcommand{\MlstFtJcOneRejTblResMassOverResPrMax}{equal within 7\%}  % experiments/FIGS/data/v2/curves/full_budget_long_leg2.json :: cells.<run>.N1000000.s1.jc1.per_class.label_Tbl: rejection of L162_MASS over L162, run by run :: sha256 588e935996d1006a :: prose twin of the run farther from 1
\newcommand{\MlstFtJcOneRejTblCoarseMassOverCoarse}{0.84 and 1.22}  % experiments/FIGS/data/v2/curves/full_budget_long_leg2.json :: cells.<run>.N1000000.s1.jc1.per_class.label_Tbl: rejection of R16_Q1_MASS over R16_Q1, run by run :: sha256 588e935996d1006a
\newcommand{\MlstFtJcOneRejTblCoarseMassOverCoarsePr}{equal within 30\%}  % experiments/FIGS/data/v2/curves/full_budget_long_leg2.json :: cells.<run>.N1000000.s1.jc1.per_class.label_Tbl: rejection of R16_Q1_MASS over R16_Q1, run by run :: sha256 588e935996d1006a :: prose twin: two significant figures, direction in words
\newcommand{\MlstFtJcOneRejTblCoarseMassOverCoarsePrMax}{equal within 30\%}  % experiments/FIGS/data/v2/curves/full_budget_long_leg2.json :: cells.<run>.N1000000.s1.jc1.per_class.label_Tbl: rejection of R16_Q1_MASS over R16_Q1, run by run :: sha256 588e935996d1006a :: prose twin of the run farther from 1
\newcommand{\MlstFtJcOneRejTblRandOverFull}{0.88\,\ensuremath{\pm}\,0.26}  % experiments/FIGS/data/v2/curves/full_budget_long_leg2.json :: per_class.label_Tbl: rejection of each random partition over the 188-class run of its seed :: sha256 588e935996d1006a
\newcommand{\MlstFtJcOneRejTblRandOverFullPr}{equal within 30\%}  % experiments/FIGS/data/v2/curves/full_budget_long_leg2.json :: per_class.label_Tbl: rejection of each random partition over the 188-class run of its seed :: sha256 588e935996d1006a :: prose twin: two significant figures, direction in words
\newcommand{\MlstFtJcOneRejTblRandOverFullRange}{0.61--1.2}  % experiments/FIGS/data/v2/curves/full_budget_long_leg2.json :: per_class.label_Tbl: smallest and largest of the partitions' ratios :: sha256 588e935996d1006a
\newcommand{\MlstFtJcOneRejWqqResOverFull}{1.039\,\ensuremath{\pm}\,0.035}  % experiments/FIGS/data/v2/curves/full_budget_long_leg2.json :: cells.<run>.N1000000.s1.jc1.per_class.label_Wqq: rejection of L162 over L188, run by run :: sha256 588e935996d1006a
\newcommand{\MlstFtJcOneRejWqqResOverFullPr}{equal within 4\%}  % experiments/FIGS/data/v2/curves/full_budget_long_leg2.json :: cells.<run>.N1000000.s1.jc1.per_class.label_Wqq: rejection of L162 over L188, run by run :: sha256 588e935996d1006a :: prose twin: two significant figures, direction in words
\newcommand{\MlstFtJcOneRejWqqMidOverFull}{1.005\,\ensuremath{\pm}\,0.032}  % experiments/FIGS/data/v2/curves/full_budget_long_leg2.json :: cells.<run>.N1000000.s1.jc1.per_class.label_Wqq: rejection of R42_Q1 over L188, run by run :: sha256 588e935996d1006a
\newcommand{\MlstFtJcOneRejWqqMidOverFullPr}{equal within 4\%}  % experiments/FIGS/data/v2/curves/full_budget_long_leg2.json :: cells.<run>.N1000000.s1.jc1.per_class.label_Wqq: rejection of R42_Q1 over L188, run by run :: sha256 588e935996d1006a :: prose twin: two significant figures, direction in words
\newcommand{\MlstFtJcOneRejWqqCoarseOverFull}{1.01\,\ensuremath{\pm}\,0.07}  % experiments/FIGS/data/v2/curves/full_budget_long_leg2.json :: cells.<run>.N1000000.s1.jc1.per_class.label_Wqq: rejection of R16_Q1 over L188, run by run :: sha256 588e935996d1006a
\newcommand{\MlstFtJcOneRejWqqCoarseOverFullPr}{equal within 8\%}  % experiments/FIGS/data/v2/curves/full_budget_long_leg2.json :: cells.<run>.N1000000.s1.jc1.per_class.label_Wqq: rejection of R16_Q1 over L188, run by run :: sha256 588e935996d1006a :: prose twin: two significant figures, direction in words
\newcommand{\MlstFtJcOneRejWqqResMassOverRes}{0.996 and 0.964}  % experiments/FIGS/data/v2/curves/full_budget_long_leg2.json :: cells.<run>.N1000000.s1.jc1.per_class.label_Wqq: rejection of L162_MASS over L162, run by run :: sha256 588e935996d1006a
\newcommand{\MlstFtJcOneRejWqqResMassOverResPr}{equal within 4\%}  % experiments/FIGS/data/v2/curves/full_budget_long_leg2.json :: cells.<run>.N1000000.s1.jc1.per_class.label_Wqq: rejection of L162_MASS over L162, run by run :: sha256 588e935996d1006a :: prose twin: two significant figures, direction in words
\newcommand{\MlstFtJcOneRejWqqResMassOverResPrMax}{equal within 4\%}  % experiments/FIGS/data/v2/curves/full_budget_long_leg2.json :: cells.<run>.N1000000.s1.jc1.per_class.label_Wqq: rejection of L162_MASS over L162, run by run :: sha256 588e935996d1006a :: prose twin of the run farther from 1
\newcommand{\MlstFtJcOneRejWqqCoarseMassOverCoarse}{1.025 and 0.926}  % experiments/FIGS/data/v2/curves/full_budget_long_leg2.json :: cells.<run>.N1000000.s1.jc1.per_class.label_Wqq: rejection of R16_Q1_MASS over R16_Q1, run by run :: sha256 588e935996d1006a
\newcommand{\MlstFtJcOneRejWqqCoarseMassOverCoarsePr}{equal within 8\%}  % experiments/FIGS/data/v2/curves/full_budget_long_leg2.json :: cells.<run>.N1000000.s1.jc1.per_class.label_Wqq: rejection of R16_Q1_MASS over R16_Q1, run by run :: sha256 588e935996d1006a :: prose twin: two significant figures, direction in words
\newcommand{\MlstFtJcOneRejWqqCoarseMassOverCoarsePrMax}{equal within 8\%}  % experiments/FIGS/data/v2/curves/full_budget_long_leg2.json :: cells.<run>.N1000000.s1.jc1.per_class.label_Wqq: rejection of R16_Q1_MASS over R16_Q1, run by run :: sha256 588e935996d1006a :: prose twin of the run farther from 1
\newcommand{\MlstFtJcOneRejWqqRandOverFull}{1.02\,\ensuremath{\pm}\,0.07}  % experiments/FIGS/data/v2/curves/full_budget_long_leg2.json :: per_class.label_Wqq: rejection of each random partition over the 188-class run of its seed :: sha256 588e935996d1006a
\newcommand{\MlstFtJcOneRejWqqRandOverFullPr}{equal within 7\%}  % experiments/FIGS/data/v2/curves/full_budget_long_leg2.json :: per_class.label_Wqq: rejection of each random partition over the 188-class run of its seed :: sha256 588e935996d1006a :: prose twin: two significant figures, direction in words
\newcommand{\MlstFtJcOneRejWqqRandOverFullRange}{0.93--1.1}  % experiments/FIGS/data/v2/curves/full_budget_long_leg2.json :: per_class.label_Wqq: smallest and largest of the partitions' ratios :: sha256 588e935996d1006a
\newcommand{\MlstFtJcOneRejZqqResOverFull}{0.987\,\ensuremath{\pm}\,0.019}  % experiments/FIGS/data/v2/curves/full_budget_long_leg2.json :: cells.<run>.N1000000.s1.jc1.per_class.label_Zqq: rejection of L162 over L188, run by run :: sha256 588e935996d1006a
\newcommand{\MlstFtJcOneRejZqqResOverFullPr}{equal within 2\%}  % experiments/FIGS/data/v2/curves/full_budget_long_leg2.json :: cells.<run>.N1000000.s1.jc1.per_class.label_Zqq: rejection of L162 over L188, run by run :: sha256 588e935996d1006a :: prose twin: two significant figures, direction in words
\newcommand{\MlstFtJcOneRejZqqMidOverFull}{0.984\,\ensuremath{\pm}\,0.010}  % experiments/FIGS/data/v2/curves/full_budget_long_leg2.json :: cells.<run>.N1000000.s1.jc1.per_class.label_Zqq: rejection of R42_Q1 over L188, run by run :: sha256 588e935996d1006a
\newcommand{\MlstFtJcOneRejZqqMidOverFullPr}{equal within 2\%}  % experiments/FIGS/data/v2/curves/full_budget_long_leg2.json :: cells.<run>.N1000000.s1.jc1.per_class.label_Zqq: rejection of R42_Q1 over L188, run by run :: sha256 588e935996d1006a :: prose twin: two significant figures, direction in words
\newcommand{\MlstFtJcOneRejZqqCoarseOverFull}{0.91\,\ensuremath{\pm}\,0.05}  % experiments/FIGS/data/v2/curves/full_budget_long_leg2.json :: cells.<run>.N1000000.s1.jc1.per_class.label_Zqq: rejection of R16_Q1 over L188, run by run :: sha256 588e935996d1006a
\newcommand{\MlstFtJcOneRejZqqCoarseOverFullPr}{equal within 9\%}  % experiments/FIGS/data/v2/curves/full_budget_long_leg2.json :: cells.<run>.N1000000.s1.jc1.per_class.label_Zqq: rejection of R16_Q1 over L188, run by run :: sha256 588e935996d1006a :: prose twin: two significant figures, direction in words
\newcommand{\MlstFtJcOneRejZqqResMassOverRes}{0.9803 and 0.9820}  % experiments/FIGS/data/v2/curves/full_budget_long_leg2.json :: cells.<run>.N1000000.s1.jc1.per_class.label_Zqq: rejection of L162_MASS over L162, run by run :: sha256 588e935996d1006a
\newcommand{\MlstFtJcOneRejZqqResMassOverResPr}{equal within 2\%}  % experiments/FIGS/data/v2/curves/full_budget_long_leg2.json :: cells.<run>.N1000000.s1.jc1.per_class.label_Zqq: rejection of L162_MASS over L162, run by run :: sha256 588e935996d1006a :: prose twin: two significant figures, direction in words
\newcommand{\MlstFtJcOneRejZqqResMassOverResPrMax}{equal within 2\%}  % experiments/FIGS/data/v2/curves/full_budget_long_leg2.json :: cells.<run>.N1000000.s1.jc1.per_class.label_Zqq: rejection of L162_MASS over L162, run by run :: sha256 588e935996d1006a :: prose twin of the run farther from 1
\newcommand{\MlstFtJcOneRejZqqCoarseMassOverCoarse}{1.06 and 0.94}  % experiments/FIGS/data/v2/curves/full_budget_long_leg2.json :: cells.<run>.N1000000.s1.jc1.per_class.label_Zqq: rejection of R16_Q1_MASS over R16_Q1, run by run :: sha256 588e935996d1006a
\newcommand{\MlstFtJcOneRejZqqCoarseMassOverCoarsePr}{equal within 7\%}  % experiments/FIGS/data/v2/curves/full_budget_long_leg2.json :: cells.<run>.N1000000.s1.jc1.per_class.label_Zqq: rejection of R16_Q1_MASS over R16_Q1, run by run :: sha256 588e935996d1006a :: prose twin: two significant figures, direction in words
\newcommand{\MlstFtJcOneRejZqqCoarseMassOverCoarsePrMax}{equal within 7\%}  % experiments/FIGS/data/v2/curves/full_budget_long_leg2.json :: cells.<run>.N1000000.s1.jc1.per_class.label_Zqq: rejection of R16_Q1_MASS over R16_Q1, run by run :: sha256 588e935996d1006a :: prose twin of the run farther from 1
\newcommand{\MlstFtJcOneRejZqqRandOverFull}{0.952\,\ensuremath{\pm}\,0.028}  % experiments/FIGS/data/v2/curves/full_budget_long_leg2.json :: per_class.label_Zqq: rejection of each random partition over the 188-class run of its seed :: sha256 588e935996d1006a
\newcommand{\MlstFtJcOneRejZqqRandOverFullPr}{equal within 5\%}  % experiments/FIGS/data/v2/curves/full_budget_long_leg2.json :: per_class.label_Zqq: rejection of each random partition over the 188-class run of its seed :: sha256 588e935996d1006a :: prose twin: two significant figures, direction in words
\newcommand{\MlstFtJcOneRejZqqRandOverFullRange}{0.93--0.99}  % experiments/FIGS/data/v2/curves/full_budget_long_leg2.json :: per_class.label_Zqq: smallest and largest of the partitions' ratios :: sha256 588e935996d1006a
\newcommand{\MlstFtJcOneResNRejDrop}{zero}  % experiments/FIGS/data/v2/curves/full_budget_long_leg2.json :: per_class rejections :: sha256 588e935996d1006a :: classes whose paired rejection ratio to the 188-class model departs from 1 by more than twice its SD and its counting uncertainty, columns with fewer than JC1_MIN_BKG passing jets excluded, ratio below 1 (count)
\newcommand{\MlstFtJcOneResNRejDiff}{zero}  % experiments/FIGS/data/v2/curves/full_budget_long_leg2.json :: per_class rejections :: sha256 588e935996d1006a :: classes whose paired rejection ratio to the 188-class model departs from 1 by more than twice its SD and its counting uncertainty, columns with fewer than JC1_MIN_BKG passing jets excluded (count)
\newcommand{\MlstFtJcOneResRejDropClasses}{none}  % experiments/FIGS/data/v2/curves/full_budget_long_leg2.json :: per_class rejections :: sha256 588e935996d1006a :: the classes of ...NRejDrop, in ParT's column order
\newcommand{\MlstFtJcOneMidNRejDrop}{zero}  % experiments/FIGS/data/v2/curves/full_budget_long_leg2.json :: per_class rejections :: sha256 588e935996d1006a :: classes whose paired rejection ratio to the 188-class model departs from 1 by more than twice its SD and its counting uncertainty, columns with fewer than JC1_MIN_BKG passing jets excluded, ratio below 1 (count)
\newcommand{\MlstFtJcOneMidNRejDiff}{zero}  % experiments/FIGS/data/v2/curves/full_budget_long_leg2.json :: per_class rejections :: sha256 588e935996d1006a :: classes whose paired rejection ratio to the 188-class model departs from 1 by more than twice its SD and its counting uncertainty, columns with fewer than JC1_MIN_BKG passing jets excluded (count)
\newcommand{\MlstFtJcOneMidRejDropClasses}{none}  % experiments/FIGS/data/v2/curves/full_budget_long_leg2.json :: per_class rejections :: sha256 588e935996d1006a :: the classes of ...NRejDrop, in ParT's column order
\newcommand{\MlstFtJcOneCoarseNRejDrop}{zero}  % experiments/FIGS/data/v2/curves/full_budget_long_leg2.json :: per_class rejections :: sha256 588e935996d1006a :: classes whose paired rejection ratio to the 188-class model departs from 1 by more than twice its SD and its counting uncertainty, columns with fewer than JC1_MIN_BKG passing jets excluded, ratio below 1 (count)
\newcommand{\MlstFtJcOneCoarseNRejDiff}{zero}  % experiments/FIGS/data/v2/curves/full_budget_long_leg2.json :: per_class rejections :: sha256 588e935996d1006a :: classes whose paired rejection ratio to the 188-class model departs from 1 by more than twice its SD and its counting uncertainty, columns with fewer than JC1_MIN_BKG passing jets excluded (count)
\newcommand{\MlstFtJcOneCoarseRejDropClasses}{none}  % experiments/FIGS/data/v2/curves/full_budget_long_leg2.json :: per_class rejections :: sha256 588e935996d1006a :: the classes of ...NRejDrop, in ParT's column order
\newcommand{\MlstFtJcOneAccRandMinusCoarse}{+0.002}  % experiments/FIGS/data/v2/curves/full_budget_long_leg2.json :: cells.*.N1000000.s1.jc1.accuracy :: sha256 588e935996d1006a :: mean over the random partitions minus the mean over the 17-class runs
\newcommand{\MlstFtJcOneAccFullBoot}{0.847}  % experiments/FIGS/data/v2/budget/jc1_acc_boot_N1000000.json :: diffs.r16q1-l188.accuracy_ref (mean) :: sha256 7cccb7b0ff900fb8
\newcommand{\MlstFtJcOneAccCoarseBoot}{0.842}  % experiments/FIGS/data/v2/budget/jc1_acc_boot_N1000000.json :: diffs.r16q1-l188.accuracy_config (mean) :: sha256 7cccb7b0ff900fb8
\newcommand{\MlstFtJcOneErrResOverFullBoot}{consistent}  % experiments/FIGS/data/v2/budget/jc1_acc_boot_N1000000.json :: diffs.l162-l188.error_ratio :: sha256 7cccb7b0ff900fb8 :: delta and sigma = hypot(run SD, bootstrap SD) in per cent; 'consistent' when |delta| <= 3 sigma
\newcommand{\MlstFtJcOneErrResOverFullZ}{0.6}  % experiments/FIGS/data/v2/budget/jc1_acc_boot_N1000000.json :: diffs.l162-l188.error_ratio.z :: sha256 7cccb7b0ff900fb8 :: |delta| / sigma
\newcommand{\MlstFtJcOneErrMidOverFullBoot}{consistent}  % experiments/FIGS/data/v2/budget/jc1_acc_boot_N1000000.json :: diffs.r42q1-l188.error_ratio :: sha256 7cccb7b0ff900fb8 :: delta and sigma = hypot(run SD, bootstrap SD) in per cent; 'consistent' when |delta| <= 3 sigma
\newcommand{\MlstFtJcOneErrMidOverFullZ}{1.4}  % experiments/FIGS/data/v2/budget/jc1_acc_boot_N1000000.json :: diffs.r42q1-l188.error_ratio.z :: sha256 7cccb7b0ff900fb8 :: |delta| / sigma
\newcommand{\MlstFtJcOneErrCoarseOverFullBoot}{(3.0\,\ensuremath{\pm}\,0.6)\% more errors}  % experiments/FIGS/data/v2/budget/jc1_acc_boot_N1000000.json :: diffs.r16q1-l188.error_ratio :: sha256 7cccb7b0ff900fb8 :: delta and sigma = hypot(run SD, bootstrap SD) in per cent; 'consistent' when |delta| <= 3 sigma
\newcommand{\MlstFtJcOneErrCoarseOverFullZ}{4.8}  % experiments/FIGS/data/v2/budget/jc1_acc_boot_N1000000.json :: diffs.r16q1-l188.error_ratio.z :: sha256 7cccb7b0ff900fb8 :: |delta| / sigma
\newcommand{\MlstFtJcOneBootReplicas}{200}  % experiments/FIGS/data/v2/budget/jc1_acc_boot_N1000000.json :: replicas :: sha256 7cccb7b0ff900fb8
\newcommand{\MlstFtJcOneBootJets}{2{,}000{,}000}  % experiments/FIGS/data/v2/budget/jc1_acc_boot_N1000000.json :: n_test_jets :: sha256 7cccb7b0ff900fb8
\newcommand{\ComputeTotalGpuHours}{3{,}435}  % experiments/FIGS/data/v2/index.json :: GPU-hours of every reported pretraining run, training and validation passes :: sha256 c6c4138ee5c86423
\newcommand{\MlstComputeTotalGpuHours}{3{,}510}  % experiments/FIGS/data/v2/index.json :: GPU-hours of every pretraining run the paper reports, the kinematic-only model included :: sha256 c6c4138ee5c86423
\newcommand{\ComputeMainGpu}{V100}  % experiments/FIGS/data/v2/index.json :: GPU type with the most hours over the four label sets :: sha256 c6c4138ee5c86423
\newcommand{\ComputeMainGpuShareMin}{56\%}  % experiments/FIGS/data/v2/index.json :: lowest share of a label set's GPU-hours on that type :: sha256 c6c4138ee5c86423
\newcommand{\ComputeMainGpuShareMax}{71\%}  % experiments/FIGS/data/v2/index.json :: highest share of a label set's GPU-hours on that type :: sha256 c6c4138ee5c86423
\newcommand{\ComputeRunsDone}{27}  % experiments/FIGS/data/v2/index.json :: paper runs with a compute record (count) :: sha256 c6c4138ee5c86423
\newcommand{\ComputeGpuTypes}{L4, RTX 3090, RTX 4090 and V100}  % experiments/FIGS/data/v2/index.json :: GPU types any finished paper run trained on :: sha256 c6c4138ee5c86423
\newcommand{\MlstComputeGpuTypes}{L4, RTX 3090, RTX 4090 and V100}  % experiments/FIGS/data/v2/index.json :: GPU types of the pretraining runs the paper reports :: sha256 c6c4138ee5c86423
\newcommand{\MlstFtJcOneErrCoarseOverFullBootEFour}{\pending{the JetClass bootstrap at 1e4 (after the equal-stopping read-out)}}  % experiments/FIGS/data/v2/index.json :: not computed yet :: sha256 c6c4138ee5c86423 :: waits on the JetClass bootstrap at 1e4 (after the equal-stopping read-out)
\newcommand{\MlstLpBbQgZeroAccDropFull}{\pending{pair probes with the quark/gluon inputs zeroed (open item O10, not run)}}  % experiments/FIGS/data/v2/index.json :: not computed yet :: sha256 c6c4138ee5c86423 :: waits on pair probes with the quark/gluon inputs zeroed (open item O10, not run)
\newcommand{\MlstLpBbQgZeroAccDropCoarse}{\pending{pair probes with the quark/gluon inputs zeroed (open item O10, not run)}}  % experiments/FIGS/data/v2/index.json :: not computed yet :: sha256 c6c4138ee5c86423 :: waits on pair probes with the quark/gluon inputs zeroed (open item O10, not run)
\newcommand{\XbsRthirtyFullOverScratch}{\pending{X->bs fits and from-scratch runs}}  % experiments/FIGS/data/v2/index.json :: not computed yet :: sha256 c6c4138ee5c86423 :: waits on X->bs fits and from-scratch runs
\newcommand{\XbsRthirtyFullOverScratchPr}{\pending{X->bs fits and from-scratch runs}}  % experiments/FIGS/data/v2/index.json :: not computed yet :: sha256 c6c4138ee5c86423 :: waits on X->bs fits and from-scratch runs
\newcommand{\XbsRthirtyCoarseOverFull}{\pending{X->bs fits}}  % experiments/FIGS/data/v2/index.json :: not computed yet :: sha256 c6c4138ee5c86423 :: waits on X->bs fits
\newcommand{\XbsRthirtyCoarseOverFullPr}{\pending{X->bs fits}}  % experiments/FIGS/data/v2/index.json :: not computed yet :: sha256 c6c4138ee5c86423 :: waits on X->bs fits
\newcommand{\XbsRthirtyNodeOverTransferFull}{\pending{X->bs fits}}  % experiments/FIGS/data/v2/index.json :: not computed yet :: sha256 c6c4138ee5c86423 :: waits on X->bs fits
\newcommand{\XbsRthirtyNodeOverTransferFullPr}{\pending{X->bs fits}}  % experiments/FIGS/data/v2/index.json :: not computed yet :: sha256 c6c4138ee5c86423 :: waits on X->bs fits
\newcommand{\XbbRthirtyDirectOverTransferFull}{\pending{X->bb fits}}  % experiments/FIGS/data/v2/index.json :: not computed yet :: sha256 c6c4138ee5c86423 :: waits on X->bb fits
\newcommand{\XbbRthirtyDirectOverTransferFullPr}{\pending{X->bb fits}}  % experiments/FIGS/data/v2/index.json :: not computed yet :: sha256 c6c4138ee5c86423 :: waits on X->bb fits
\newcommand{\XbbRthirtyTransferCoarseOverFull}{\pending{X->bb fits}}  % experiments/FIGS/data/v2/index.json :: not computed yet :: sha256 c6c4138ee5c86423 :: waits on X->bb fits
\newcommand{\XbbRthirtyTransferCoarseOverFullPr}{\pending{X->bb fits}}  % experiments/FIGS/data/v2/index.json :: not computed yet :: sha256 c6c4138ee5c86423 :: waits on X->bb fits
             % the counts below are its macros
% Visual language: Sophon's Fig. 1 (arXiv:2405.12972): the document's serif font, light fills,
% rounded boxes, light grey panels behind the stages (PI 2026-11-02). The text is set in the
% document's font at the printed sizes 8.5-10 pt (MLST 8-12 pt), never scaled.
% Text width: 15 cm (MLST full width). The picture is 15.0 cm wide and is drawn at its final
% size, never scaled, so the font sizes hold.
% Optional: \def\FoneDropXbx{} before % Figure 1 (F1 of docs/FIGURE_PLAN.md): the design of the study, in TikZ.
% \input{fig_f1} from main.tex; this file holds one figure environment and nothing else.
%
% WHAT main.tex MUST LOAD (this file loads nothing itself):
%   \usepackage{tikz}
%   \usetikzlibrary{arrows.meta}          % the only TikZ library used (arrow tips)
%   \definecolor{c188}{HTML}{1F77B4}      % the configuration colours of experiments/FIGS/
%   \definecolor{c162}{HTML}{FF7F0E}      % mlst_style.py (matplotlib's tab10 in order, as in
%   \definecolor{c43}{HTML}{2CA02C}       % Sophon's figures); the variants are drawn in their
%   \definecolor{c17}{HTML}{D62728}       % label set's colour with a dotted border, as the
%   \definecolor{crand}{HTML}{17BECF}     % plotted figures draw variants dotted
%   \definecolor{cflav}{HTML}{9467BD}
%   \input{results_generated}             % the counts below are its macros
% Visual language: Sophon's Fig. 1 (arXiv:2405.12972): the document's serif font, light fills,
% rounded boxes, light grey panels behind the stages (PI 2026-11-02). The text is set in the
% document's font at the printed sizes 8.5-10 pt (MLST 8-12 pt), never scaled.
% Text width: 15 cm (MLST full width). The picture is 15.0 cm wide and is drawn at its final
% size, never scaled, so the font sizes hold.
% Optional: \def\FoneDropXbx{} before \input{fig_f1} removes X->bb/X->bs from panel (c)
% (PI 2026-11-02: if the samples are not staged by 12 Nov, the X->bb/X->bs section goes).
%
% Every class box is drawn from the label maps (checked by hand, 2026-11-02):
%   (a) configs/labelmaps/rung_label_maps.v1.csv, columns L188, L162, R42_Q1, R16_Q1. Rows: the ten
%       decays X->QQ/gg in the order bb bc bq | cc cs cq | ss sq qq gg (the 43-class b / c / light
%       groups are then contiguous), then ee mm | tauh-taue tauh-taum, then the 27 QCD classes.
%   (b) column: configs/labelmaps/rand_label_map.v2.csv, RAND2_p1 on the same rows: groups
%       0 0 0 0 0 9 9 9 0 3 | 11 11 | 15 4 | QCD, lettered A..F in order of first appearance.
%       strips: configs/labelmaps/flavour_pair_map.v2.csv, the 22 classes of orbit X_YY_QQQQ.
%       Order with a b quark: bbbb bbcc bbss bbqq qqqb bbcq ccbs ccbq qqbc qqbs bcsq;
%       without: cccc ccss ccqq ssss ssqq qqqq qqqc qqqs ccsq sscq qqcs.
%       FLAV_F0: all 22 in group 9. FLAV_F1: the 11 with a b in group 11, the rest 9.
%       FLAV_F1R: group 11 = bbss bbqq qqqb bbcq ccbq qqbc (with b, positions 2-5, 7, 8 from 0)
%       and ccss ccqq ssss qqqc sscq (without, positions 1-3, 6, 9); filled square = group 11.
% Geometry (cm): row height 0.35 from y = -1.18, gaps of 0.14 between the quark/gluon, lepton
% and QCD blocks, the QCD row 0.70; boxes inset 0.025 top and bottom; columns 1.45 wide with
% left edges 1.75, 3.28, 4.81, 6.34 (hierarchy) and 8.25 (random partition).
\providecommand{\Kfull}{\VocabSizeLoneeighteight}
\providecommand{\Kres}{\VocabSizeLonesixtwo}
\providecommand{\Kmid}{\VocabSizeRfourtwoQone}
\providecommand{\Kcoarse}{\VocabSizeRonesixQone}
\begin{figure}[tbp]
\centering
\ifdim\linewidth<14.99cm\typeout{fig_f1.tex: WARNING: the drawing is 15 cm wide but \string\linewidth\space
  is \the\linewidth; set a 15 cm text width (it is not scaled, to keep its 8.5 pt text)}\fi
\def\FoneText{\rmfamily\fontsize{8.5}{9.5}\selectfont}%
\def\FoneMain{\rmfamily\fontsize{9}{10.5}\selectfont}%
\def\FoneBold{\rmfamily\bfseries\fontsize{9}{10.5}\selectfont}%
\def\FoneTitle{\rmfamily\bfseries\fontsize{10}{12}\selectfont}%
\ifdefined\FoneDropXbx\def\FoneXbx{}\else\def\FoneXbx{\\$X\to b\bar b$, $X\to bs$}\fi
\begin{tikzpicture}[
  x=1cm, y=1cm,
  font=\FoneMain,
  ttl/.style={anchor=base west, font=\FoneTitle, inner sep=0pt},
  cls/.style={draw=black!55, line width=0.45pt, rounded corners=1.2pt},
  hdr/.style={draw=#1!70!black, fill=#1!22, line width=0.6pt, rounded corners=2.5pt, align=center,
              font=\FoneText, inner sep=0pt},
  inbox/.style={align=center, font=\FoneText, inner sep=0pt},
  letter/.style={font=\FoneBold, inner sep=0pt},
  rowlab/.style={anchor=east, inner sep=0pt},
  txt/.style={align=left, font=\FoneText, inner sep=0pt},
  sq/.style={draw=black!65, line width=0.4pt},
  chip/.style={draw=#1!70!black, fill=#1!22, line width=0.6pt, rounded corners=3pt,
               align=center, font=\FoneText, inner sep=0pt, anchor=north west},
  stage/.style={draw=black!45, fill=black!4, line width=0.6pt, rounded corners=3pt, align=left,
                font=\FoneText, inner sep=0.1cm},
  flow/.style={-{Latex[length=1.8mm, width=1.4mm]}, line width=0.7pt, black!75},
]
% ================================================================ (a) the class hierarchy
\node[ttl] at (0,0) {(a) Class hierarchy};
% faint bands on every other row, drawn first so the boxes sit on top
\foreach \rowtop in {-1.53,-2.23,-2.93,-3.63,-4.33,-5.17,-5.87}
  \fill[black!4] (0,\rowtop) rectangle (9.70,\rowtop-0.35);
% column headers
\node[rowlab, align=right, font=\FoneText] at (1.65,-0.68) {JetClass-II\\class};
\node[hdr=c188, minimum width=1.45cm, minimum height=0.80cm, anchor=north west] at (1.75,-0.28) {\Kfull\\classes};
\node[hdr=c162, minimum width=1.45cm, minimum height=0.80cm, anchor=north west] at (3.28,-0.28) {\Kres\\classes};
\node[hdr=c43,  minimum width=1.45cm, minimum height=0.80cm, anchor=north west] at (4.81,-0.28) {\Kmid\\classes};
\node[hdr=c17,  minimum width=1.45cm, minimum height=0.80cm, anchor=north west] at (6.34,-0.28) {\Kcoarse\\classes};
% row labels
\node[rowlab] at (1.65,-1.355) {$X\to b\bar b$};
\node[rowlab] at (1.65,-1.705) {$X\to bc$};
\node[rowlab] at (1.65,-2.055) {$X\to bq$};
\node[rowlab] at (1.65,-2.405) {$X\to c\bar c$};
\node[rowlab] at (1.65,-2.755) {$X\to cs$};
\node[rowlab] at (1.65,-3.105) {$X\to cq$};
\node[rowlab] at (1.65,-3.455) {$X\to s\bar s$};
\node[rowlab] at (1.65,-3.805) {$X\to sq$};
\node[rowlab] at (1.65,-4.155) {$X\to q\bar q$};
\node[rowlab] at (1.65,-4.505) {$X\to gg$};
\node[rowlab] at (1.65,-4.995) {$X\to ee$};
\node[rowlab] at (1.65,-5.345) {$X\to\mu\mu$};
\node[rowlab] at (1.65,-5.695) {$X\to\tau_h\tau_e$};
\node[rowlab] at (1.65,-6.045) {$X\to\tau_h\tau_\mu$};
\node[rowlab] at (1.65,-6.710) {QCD jets};
% 188 and 162 classes: every resonant decay is its own class
\foreach \rowtop in {-1.18,-1.53,-1.88,-2.23,-2.58,-2.93,-3.28,-3.63,-3.98,-4.33,-4.82,-5.17,-5.52,-5.87} {
  \draw[cls, fill=c188!15] (1.75,\rowtop-0.025) rectangle (3.20,\rowtop-0.325);
  \draw[cls, fill=c162!15] (3.28,\rowtop-0.025) rectangle (4.73,\rowtop-0.325);
}
\draw[cls, fill=c188!15] (1.75,-6.385) rectangle (3.20,-7.035);
\node[inbox] at (2.475,-6.71) {\LitJetClassTwoNQcdClasses\\classes};
\draw[cls, fill=c162!15] (3.28,-6.385) rectangle (4.73,-7.035);
\node[inbox] at (4.005,-6.71) {1 class};
% 43 classes: decay topology x {b, else c, else light quark or gluon}; e and mu merged
\draw[cls, fill=c43!15] (4.81,-1.205) rectangle (6.26,-2.205);
\node[inbox] at (5.535,-1.705) {with $b$};
\draw[cls, fill=c43!15] (4.81,-2.255) rectangle (6.26,-3.255);
\node[inbox] at (5.535,-2.755) {with $c$,\\no $b$};
\draw[cls, fill=c43!15] (4.81,-3.305) rectangle (6.26,-4.655);
\node[inbox] at (5.535,-3.980) {$u$, $d$, $s$\\or $g$};
\draw[cls, fill=c43!15] (4.81,-4.845) rectangle (6.26,-5.495);
\node[inbox] at (5.535,-5.170) {$\ell\ell$};
\draw[cls, fill=c43!15] (4.81,-5.545) rectangle (6.26,-6.195);
\node[inbox] at (5.535,-5.870) {$\tau_h\tau_\ell$};
\draw[cls, fill=c43!15] (4.81,-6.385) rectangle (6.26,-7.035);
\node[inbox] at (5.535,-6.710) {1 class};
% 17 classes: number and type of prongs, no flavor
\draw[cls, fill=c17!15] (6.34,-1.205) rectangle (7.79,-4.655);
\node[inbox] at (7.065,-2.930) {two\\quarks\\or gluons};
\draw[cls, fill=c17!15] (6.34,-4.845) rectangle (7.79,-5.495);
\node[inbox] at (7.065,-5.170) {$\ell\ell$};
\draw[cls, fill=c17!15] (6.34,-5.545) rectangle (7.79,-6.195);
\node[inbox] at (7.065,-5.870) {$\tau_h\tau_\ell$};
\draw[cls, fill=c17!15] (6.34,-6.385) rectangle (7.79,-7.035);
\node[inbox] at (7.065,-6.710) {1 class};
% what each merge removes
\node[anchor=north west, font=\FoneBold, inner sep=0pt] at (0,-7.27) {What each merge removes};
\node[txt, anchor=north west, text width=6.04cm] (m1) at (1.75,-7.65)
  {the $b$, $c$ and $s$ content of the QCD jets};
\node[txt, anchor=north west, text width=6.04cm] (m2) at ([yshift=-0.10cm]m1.south west)
  {the number of $b$ and $c$ quarks, light-quark flavor, quark against gluon, and lepton
   flavor ($e$ against $\mu$)};
\node[txt, anchor=north west, text width=6.04cm] (m3) at ([yshift=-0.10cm]m2.south west)
  {quark flavor ($b$, $c$ or light); the number and type of prongs remain};
\node[txt, anchor=north east] at ([xshift=-0.10cm]m1.north west) {\Kfull\,$\to$\,\Kres};
\node[txt, anchor=north east] at ([xshift=-0.10cm]m2.north west) {\Kres\,$\to$\,\Kmid};
\node[txt, anchor=north east] at ([xshift=-0.10cm]m3.north west) {\Kmid\,$\to$\,\Kcoarse};
% ================================================================ (b) the controls
\node[ttl] at (8.25,0) {(b) Controls, \Kcoarse{} classes each};
% the first random partition on the rows of (a); one letter = one class
\node[hdr=crand, minimum width=1.45cm, minimum height=0.80cm, anchor=north west] at (8.25,-0.28) {random\\partition};
\draw[cls, fill=crand!15] (8.25,-1.205) rectangle (9.70,-2.905);
\node[letter] at (8.975,-2.055) {A};
\draw[cls, fill=crand!15] (8.25,-2.955) rectangle (9.70,-3.955);
\node[letter] at (8.975,-3.455) {B};
\draw[cls, fill=crand!15] (8.25,-4.005) rectangle (9.70,-4.305);
\node[letter] at (8.975,-4.155) {A};
\draw[cls, fill=crand!15] (8.25,-4.355) rectangle (9.70,-4.655);
\node[letter] at (8.975,-4.505) {C};
\draw[cls, fill=crand!15] (8.25,-4.845) rectangle (9.70,-5.495);
\node[letter] at (8.975,-5.170) {D};
\draw[cls, fill=crand!15] (8.25,-5.545) rectangle (9.70,-5.845);
\node[letter] at (8.975,-5.695) {E};
\draw[cls, fill=crand!15] (8.25,-5.895) rectangle (9.70,-6.195);
\node[letter] at (8.975,-6.045) {F};
\draw[cls, fill=crand!15] (8.25,-6.385) rectangle (9.70,-7.035);
\node[inbox] at (8.975,-6.710) {1 class};
% how the random partitions are built
\node[hdr=crand, anchor=north west, text width=4.75cm, inner sep=0.1cm, align=left, font=\FoneBold]
  (rh) at (10.05,-0.28) {Random partitions (\RandVtwoNPartitions)};
\node[txt, anchor=north west, text width=4.95cm] (rb) at ([yshift=-0.12cm]rh.south west)
  {\textbullet\ groups of resonant classes drawn at random, plus one class of QCD jets\\[2pt]
   \textbullet\ group shares of training jets within \RandVtwoShareTol{} of the \Kcoarse-class groups\\[2pt]
   \textbullet\ best of \RandVtwoTrials{} draws: the fewest merges shared with the \Kcoarse-class level\\[2pt]
   \textbullet\ left: the first partition; boxes with the same letter are one class};
% how the random partitions with controlled flavor boundaries are built
\node[hdr=cflav, anchor=north west, text width=4.75cm, inner sep=0.1cm, align=left, font=\FoneBold]
  (fh) at ([yshift=-0.22cm]rb.south west) {Random partitions with controlled flavor boundaries};
\node[txt, anchor=north west, text width=4.95cm] (fn) at ([yshift=-0.12cm]fh.south west)
  {whole families of decays that differ only in quark flavor; two of them split one family,
   the \RandFlavSplitOrbitClasses{} decays \RandFlavSplitOrbit:};
\begin{scope}[shift={(fn.south west)}]
  % brackets over the decays with and without a b quark (squares at a 0.2 cm pitch)
  \node[txt, anchor=south] at (1.075,-0.45) {with $b$};
  \node[txt, anchor=south] at (3.425,-0.45) {without $b$};
  \draw[line width=0.5pt] (0,-0.54) -- (0,-0.47) -- (2.15,-0.47) -- (2.15,-0.54);
  \draw[line width=0.5pt] (2.35,-0.54) -- (2.35,-0.47) -- (4.50,-0.47) -- (4.50,-0.54);
  % no flavor boundary: all 22 in one class
  \node[txt, anchor=north west] at (0,-0.62) {no flavor boundary};
  \foreach \i in {0,...,10} {
    \draw[sq, fill=white] (0.2*\i,-1.10) rectangle ++(0.15,0.15);
    \draw[sq, fill=white] (2.35+0.2*\i,-1.10) rectangle ++(0.15,0.15);
  }
  % b-quark boundary: the decays with a b quark moved to another class
  \node[txt, anchor=north west] at (0,-1.22) {$b$-quark boundary};
  \foreach \i in {0,...,10} {
    \draw[sq, fill=cflav!70!black] (0.2*\i,-1.70) rectangle ++(0.15,0.15);
    \draw[sq, fill=white] (2.35+0.2*\i,-1.70) rectangle ++(0.15,0.15);
  }
  % random flavor boundary: the same number of decays, drawn at random
  \node[txt, anchor=north west] at (0,-1.82) {random flavor boundary};
  \foreach \i in {0,1,6,9,10}
    \draw[sq, fill=white] (0.2*\i,-2.30) rectangle ++(0.15,0.15);
  \foreach \i in {2,3,4,5,7,8}
    \draw[sq, fill=cflav!70!black] (0.2*\i,-2.30) rectangle ++(0.15,0.15);
  \foreach \i in {0,4,5,7,8,10}
    \draw[sq, fill=white] (2.35+0.2*\i,-2.30) rectangle ++(0.15,0.15);
  \foreach \i in {1,2,3,6,9}
    \draw[sq, fill=cflav!70!black] (2.35+0.2*\i,-2.30) rectangle ++(0.15,0.15);
  % key
  \draw[sq, fill=cflav!70!black] (0,-2.62) rectangle ++(0.15,0.15);
  \node[txt, anchor=west] (fkey) at (0.25,-2.545) {moved to another class};
\end{scope}
% ================================================================ (c) the pipeline
% (c) starts below the lower of the two columns above, whatever their text heights
\pgfpointanchor{m3}{south}\pgfgetlastxy{\FoneXa}{\FoneYa}%
\pgfpointanchor{fkey}{south}\pgfgetlastxy{\FoneXb}{\FoneYb}%
\pgfmathsetlengthmacro{\FoneYc}{min(\FoneYa,\FoneYb)-0.45cm}%
\coordinate (ctop) at (0pt,\FoneYc);
\begin{scope}[shift={(ctop)}]
  \node[ttl, anchor=north west] at (0,0) {(c) Pretraining and evaluation};
  % what every model of one seed shares
  \node[stage, text width=2.15cm] (seed) at (1.175,-1.96)
    {\textbf{Same for every model of one seed:} the training jets, their order and the
     initial weights, except the output layer};
  % one pretraining per label set: the configuration colours of every later figure
  \node[draw=black!45, fill=black!4, line width=0.6pt, rounded corners=3pt, minimum width=5.15cm,
        minimum height=3.08cm, anchor=north west, inner sep=0pt] (pre) at (2.81,-0.42) {};
  \node[anchor=north west, font=\FoneBold, inner sep=0pt] at (2.96,-0.52) {Pretraining on JetClass-II};
  \node[chip=c188, minimum width=1.13cm, minimum height=0.42cm] at (2.96,-0.95) {\Kfull};
  \node[chip=c162, minimum width=1.13cm, minimum height=0.42cm] at (4.20,-0.95) {\Kres};
  \node[chip=c43,  minimum width=1.13cm, minimum height=0.42cm] at (5.44,-0.95) {\Kmid};
  \node[chip=c17,  minimum width=1.13cm, minimum height=0.42cm] at (6.68,-0.95) {\Kcoarse};
  \node[chip=crand, minimum width=1.70cm, minimum height=0.70cm] at (2.96,-1.47) {random\\partitions};
  \node[chip=cflav, minimum width=3.04cm, minimum height=0.70cm] at (4.77,-1.47) {controlled flavor\\boundaries};
  \node[chip=c162, densely dotted, minimum width=1.13cm, minimum height=0.42cm] at (2.96,-2.32) {\Kres};
  \node[chip=c17,  densely dotted, minimum width=1.13cm, minimum height=0.42cm] at (4.20,-2.32) {\Kcoarse};
  \node[txt, anchor=west] at (5.44,-2.53) {with auxiliary\\mass regression};
  \node[chip=c188, densely dotted, minimum width=1.13cm, minimum height=0.42cm] at (2.96,-2.94) {\Kfull};
  \node[txt, anchor=west] at (4.20,-3.15) {kinematic inputs only};
  % evaluation
  \node[stage, text width=2.4cm] (fz) at (9.68,-0.95) {frozen features:\\linear probe,\\MLP probe};
  \node[stage, text width=2.4cm] (ft) at (9.68,-2.05) {fine-tuning\\of all weights};
  \node[stage, text width=2.4cm, draw=black, fill=white, line width=1.1pt] (sc) at (9.68,-3.15) {trained from\\scratch};
  \node[stage, text width=3.4cm, minimum height=3.2cm] (task) at (13.20,-2.00)
    {\textbf{Downstream tasks}\\JetClass-II: all classes\\JetClass-II: class pairs\\\quad(both class-balanced)\\JetClass\\
     top tagging\\quark/gluon, Pythia\\quark/gluon, Herwig\\jet mass\FoneXbx};
  \draw[flow] (seed.east) -- (seed.east -| pre.west);
  \draw[flow] (fz.west -| pre.east) -- (fz.west);
  \draw[flow] (ft.west -| pre.east) -- (ft.west);
  \draw[flow] (fz.east) -- (fz.east -| task.west);
  \draw[flow] (ft.east) -- (ft.east -| task.west);
  \draw[flow] (sc.east) -- (sc.east -| task.west);
\end{scope}
\end{tikzpicture}
\caption{Design of the study.
(a)~The class hierarchy on the decays of a heavy resonance~$X$ and on the jets from quantum
chromodynamics (QCD); $q$ is a $u$ or $d$ quark, $\ell$ an electron or a muon, and $\tau_h$,
$\tau_e$, $\tau_\mu$ a $\tau$ lepton decaying to hadrons, an electron or a muon. Each box is one
class of the label set above its column, except the QCD box of the \Kfull-class column (Sophon's
label set~\cite{sophon}), which stands for \LitJetClassTwoNQcdClasses{} classes. Each level merges
classes of the level to its left; the lines below name what each merge removes.
(b)~The controls: the first random partition, where boxes with the same letter are one class, and
the random partitions with controlled flavor boundaries on the decays \RandFlavSplitOrbit{} (one
square per decay; filled squares are moved to a second class).
(c)~Models of one seed share the training jets, their order and the initial weights, except the
output layer, the mass regression target, and the first layer and initial weights of the model with
kinematic inputs only (the particle four-momenta). The hierarchy levels have \MlstNRunsWord{} seeds
each, the models with the auxiliary mass regression \MlstMassNRunsWord{}, the random partitions with
controlled flavor boundaries \RandFlavRuns{} and each random partition \RandVtwoRuns{}. Frozen
features, the output of the backbone (the network below the classification head) with its weights held
fixed, feed linear probes (logistic regressions) and multilayer perceptron (MLP) probes; fine-tuning
trains all weights. Colors identify the label sets and controls in every figure; the plots give the three partitions with
controlled flavor boundaries three shades; dotted borders mark the two
variants, which the other figures draw as dotted lines in the color of their label set.}
\label{fig:design}
\end{figure}
 removes X->bb/X->bs from panel (c)
% (PI 2026-11-02: if the samples are not staged by 12 Nov, the X->bb/X->bs section goes).
%
% Every class box is drawn from the label maps (checked by hand, 2026-11-02):
%   (a) configs/labelmaps/rung_label_maps.v1.csv, columns L188, L162, R42_Q1, R16_Q1. Rows: the ten
%       decays X->QQ/gg in the order bb bc bq | cc cs cq | ss sq qq gg (the 43-class b / c / light
%       groups are then contiguous), then ee mm | tauh-taue tauh-taum, then the 27 QCD classes.
%   (b) column: configs/labelmaps/rand_label_map.v2.csv, RAND2_p1 on the same rows: groups
%       0 0 0 0 0 9 9 9 0 3 | 11 11 | 15 4 | QCD, lettered A..F in order of first appearance.
%       strips: configs/labelmaps/flavour_pair_map.v2.csv, the 22 classes of orbit X_YY_QQQQ.
%       Order with a b quark: bbbb bbcc bbss bbqq qqqb bbcq ccbs ccbq qqbc qqbs bcsq;
%       without: cccc ccss ccqq ssss ssqq qqqq qqqc qqqs ccsq sscq qqcs.
%       FLAV_F0: all 22 in group 9. FLAV_F1: the 11 with a b in group 11, the rest 9.
%       FLAV_F1R: group 11 = bbss bbqq qqqb bbcq ccbq qqbc (with b, positions 2-5, 7, 8 from 0)
%       and ccss ccqq ssss qqqc sscq (without, positions 1-3, 6, 9); filled square = group 11.
% Geometry (cm): row height 0.35 from y = -1.18, gaps of 0.14 between the quark/gluon, lepton
% and QCD blocks, the QCD row 0.70; boxes inset 0.025 top and bottom; columns 1.45 wide with
% left edges 1.75, 3.28, 4.81, 6.34 (hierarchy) and 8.25 (random partition).
\providecommand{\Kfull}{\VocabSizeLoneeighteight}
\providecommand{\Kres}{\VocabSizeLonesixtwo}
\providecommand{\Kmid}{\VocabSizeRfourtwoQone}
\providecommand{\Kcoarse}{\VocabSizeRonesixQone}
\begin{figure}[tbp]
\centering
\ifdim\linewidth<14.99cm\typeout{fig_f1.tex: WARNING: the drawing is 15 cm wide but \string\linewidth\space
  is \the\linewidth; set a 15 cm text width (it is not scaled, to keep its 8.5 pt text)}\fi
\def\FoneText{\rmfamily\fontsize{8.5}{9.5}\selectfont}%
\def\FoneMain{\rmfamily\fontsize{9}{10.5}\selectfont}%
\def\FoneBold{\rmfamily\bfseries\fontsize{9}{10.5}\selectfont}%
\def\FoneTitle{\rmfamily\bfseries\fontsize{10}{12}\selectfont}%
\ifdefined\FoneDropXbx\def\FoneXbx{}\else\def\FoneXbx{\\$X\to b\bar b$, $X\to bs$}\fi
\begin{tikzpicture}[
  x=1cm, y=1cm,
  font=\FoneMain,
  ttl/.style={anchor=base west, font=\FoneTitle, inner sep=0pt},
  cls/.style={draw=black!55, line width=0.45pt, rounded corners=1.2pt},
  hdr/.style={draw=#1!70!black, fill=#1!22, line width=0.6pt, rounded corners=2.5pt, align=center,
              font=\FoneText, inner sep=0pt},
  inbox/.style={align=center, font=\FoneText, inner sep=0pt},
  letter/.style={font=\FoneBold, inner sep=0pt},
  rowlab/.style={anchor=east, inner sep=0pt},
  txt/.style={align=left, font=\FoneText, inner sep=0pt},
  sq/.style={draw=black!65, line width=0.4pt},
  chip/.style={draw=#1!70!black, fill=#1!22, line width=0.6pt, rounded corners=3pt,
               align=center, font=\FoneText, inner sep=0pt, anchor=north west},
  stage/.style={draw=black!45, fill=black!4, line width=0.6pt, rounded corners=3pt, align=left,
                font=\FoneText, inner sep=0.1cm},
  flow/.style={-{Latex[length=1.8mm, width=1.4mm]}, line width=0.7pt, black!75},
]
% ================================================================ (a) the class hierarchy
\node[ttl] at (0,0) {(a) Class hierarchy};
% faint bands on every other row, drawn first so the boxes sit on top
\foreach \rowtop in {-1.53,-2.23,-2.93,-3.63,-4.33,-5.17,-5.87}
  \fill[black!4] (0,\rowtop) rectangle (9.70,\rowtop-0.35);
% column headers
\node[rowlab, align=right, font=\FoneText] at (1.65,-0.68) {JetClass-II\\class};
\node[hdr=c188, minimum width=1.45cm, minimum height=0.80cm, anchor=north west] at (1.75,-0.28) {\Kfull\\classes};
\node[hdr=c162, minimum width=1.45cm, minimum height=0.80cm, anchor=north west] at (3.28,-0.28) {\Kres\\classes};
\node[hdr=c43,  minimum width=1.45cm, minimum height=0.80cm, anchor=north west] at (4.81,-0.28) {\Kmid\\classes};
\node[hdr=c17,  minimum width=1.45cm, minimum height=0.80cm, anchor=north west] at (6.34,-0.28) {\Kcoarse\\classes};
% row labels
\node[rowlab] at (1.65,-1.355) {$X\to b\bar b$};
\node[rowlab] at (1.65,-1.705) {$X\to bc$};
\node[rowlab] at (1.65,-2.055) {$X\to bq$};
\node[rowlab] at (1.65,-2.405) {$X\to c\bar c$};
\node[rowlab] at (1.65,-2.755) {$X\to cs$};
\node[rowlab] at (1.65,-3.105) {$X\to cq$};
\node[rowlab] at (1.65,-3.455) {$X\to s\bar s$};
\node[rowlab] at (1.65,-3.805) {$X\to sq$};
\node[rowlab] at (1.65,-4.155) {$X\to q\bar q$};
\node[rowlab] at (1.65,-4.505) {$X\to gg$};
\node[rowlab] at (1.65,-4.995) {$X\to ee$};
\node[rowlab] at (1.65,-5.345) {$X\to\mu\mu$};
\node[rowlab] at (1.65,-5.695) {$X\to\tau_h\tau_e$};
\node[rowlab] at (1.65,-6.045) {$X\to\tau_h\tau_\mu$};
\node[rowlab] at (1.65,-6.710) {QCD jets};
% 188 and 162 classes: every resonant decay is its own class
\foreach \rowtop in {-1.18,-1.53,-1.88,-2.23,-2.58,-2.93,-3.28,-3.63,-3.98,-4.33,-4.82,-5.17,-5.52,-5.87} {
  \draw[cls, fill=c188!15] (1.75,\rowtop-0.025) rectangle (3.20,\rowtop-0.325);
  \draw[cls, fill=c162!15] (3.28,\rowtop-0.025) rectangle (4.73,\rowtop-0.325);
}
\draw[cls, fill=c188!15] (1.75,-6.385) rectangle (3.20,-7.035);
\node[inbox] at (2.475,-6.71) {\LitJetClassTwoNQcdClasses\\classes};
\draw[cls, fill=c162!15] (3.28,-6.385) rectangle (4.73,-7.035);
\node[inbox] at (4.005,-6.71) {1 class};
% 43 classes: decay topology x {b, else c, else light quark or gluon}; e and mu merged
\draw[cls, fill=c43!15] (4.81,-1.205) rectangle (6.26,-2.205);
\node[inbox] at (5.535,-1.705) {with $b$};
\draw[cls, fill=c43!15] (4.81,-2.255) rectangle (6.26,-3.255);
\node[inbox] at (5.535,-2.755) {with $c$,\\no $b$};
\draw[cls, fill=c43!15] (4.81,-3.305) rectangle (6.26,-4.655);
\node[inbox] at (5.535,-3.980) {$u$, $d$, $s$\\or $g$};
\draw[cls, fill=c43!15] (4.81,-4.845) rectangle (6.26,-5.495);
\node[inbox] at (5.535,-5.170) {$\ell\ell$};
\draw[cls, fill=c43!15] (4.81,-5.545) rectangle (6.26,-6.195);
\node[inbox] at (5.535,-5.870) {$\tau_h\tau_\ell$};
\draw[cls, fill=c43!15] (4.81,-6.385) rectangle (6.26,-7.035);
\node[inbox] at (5.535,-6.710) {1 class};
% 17 classes: number and type of prongs, no flavor
\draw[cls, fill=c17!15] (6.34,-1.205) rectangle (7.79,-4.655);
\node[inbox] at (7.065,-2.930) {two\\quarks\\or gluons};
\draw[cls, fill=c17!15] (6.34,-4.845) rectangle (7.79,-5.495);
\node[inbox] at (7.065,-5.170) {$\ell\ell$};
\draw[cls, fill=c17!15] (6.34,-5.545) rectangle (7.79,-6.195);
\node[inbox] at (7.065,-5.870) {$\tau_h\tau_\ell$};
\draw[cls, fill=c17!15] (6.34,-6.385) rectangle (7.79,-7.035);
\node[inbox] at (7.065,-6.710) {1 class};
% what each merge removes
\node[anchor=north west, font=\FoneBold, inner sep=0pt] at (0,-7.27) {What each merge removes};
\node[txt, anchor=north west, text width=6.04cm] (m1) at (1.75,-7.65)
  {the $b$, $c$ and $s$ content of the QCD jets};
\node[txt, anchor=north west, text width=6.04cm] (m2) at ([yshift=-0.10cm]m1.south west)
  {the number of $b$ and $c$ quarks, light-quark flavor, quark against gluon, and lepton
   flavor ($e$ against $\mu$)};
\node[txt, anchor=north west, text width=6.04cm] (m3) at ([yshift=-0.10cm]m2.south west)
  {quark flavor ($b$, $c$ or light); the number and type of prongs remain};
\node[txt, anchor=north east] at ([xshift=-0.10cm]m1.north west) {\Kfull\,$\to$\,\Kres};
\node[txt, anchor=north east] at ([xshift=-0.10cm]m2.north west) {\Kres\,$\to$\,\Kmid};
\node[txt, anchor=north east] at ([xshift=-0.10cm]m3.north west) {\Kmid\,$\to$\,\Kcoarse};
% ================================================================ (b) the controls
\node[ttl] at (8.25,0) {(b) Controls, \Kcoarse{} classes each};
% the first random partition on the rows of (a); one letter = one class
\node[hdr=crand, minimum width=1.45cm, minimum height=0.80cm, anchor=north west] at (8.25,-0.28) {random\\partition};
\draw[cls, fill=crand!15] (8.25,-1.205) rectangle (9.70,-2.905);
\node[letter] at (8.975,-2.055) {A};
\draw[cls, fill=crand!15] (8.25,-2.955) rectangle (9.70,-3.955);
\node[letter] at (8.975,-3.455) {B};
\draw[cls, fill=crand!15] (8.25,-4.005) rectangle (9.70,-4.305);
\node[letter] at (8.975,-4.155) {A};
\draw[cls, fill=crand!15] (8.25,-4.355) rectangle (9.70,-4.655);
\node[letter] at (8.975,-4.505) {C};
\draw[cls, fill=crand!15] (8.25,-4.845) rectangle (9.70,-5.495);
\node[letter] at (8.975,-5.170) {D};
\draw[cls, fill=crand!15] (8.25,-5.545) rectangle (9.70,-5.845);
\node[letter] at (8.975,-5.695) {E};
\draw[cls, fill=crand!15] (8.25,-5.895) rectangle (9.70,-6.195);
\node[letter] at (8.975,-6.045) {F};
\draw[cls, fill=crand!15] (8.25,-6.385) rectangle (9.70,-7.035);
\node[inbox] at (8.975,-6.710) {1 class};
% how the random partitions are built
\node[hdr=crand, anchor=north west, text width=4.75cm, inner sep=0.1cm, align=left, font=\FoneBold]
  (rh) at (10.05,-0.28) {Random partitions (\RandVtwoNPartitions)};
\node[txt, anchor=north west, text width=4.95cm] (rb) at ([yshift=-0.12cm]rh.south west)
  {\textbullet\ groups of resonant classes drawn at random, plus one class of QCD jets\\[2pt]
   \textbullet\ group shares of training jets within \RandVtwoShareTol{} of the \Kcoarse-class groups\\[2pt]
   \textbullet\ best of \RandVtwoTrials{} draws: the fewest merges shared with the \Kcoarse-class level\\[2pt]
   \textbullet\ left: the first partition; boxes with the same letter are one class};
% how the random partitions with controlled flavor boundaries are built
\node[hdr=cflav, anchor=north west, text width=4.75cm, inner sep=0.1cm, align=left, font=\FoneBold]
  (fh) at ([yshift=-0.22cm]rb.south west) {Random partitions with controlled flavor boundaries};
\node[txt, anchor=north west, text width=4.95cm] (fn) at ([yshift=-0.12cm]fh.south west)
  {whole families of decays that differ only in quark flavor; two of them split one family,
   the \RandFlavSplitOrbitClasses{} decays \RandFlavSplitOrbit:};
\begin{scope}[shift={(fn.south west)}]
  % brackets over the decays with and without a b quark (squares at a 0.2 cm pitch)
  \node[txt, anchor=south] at (1.075,-0.45) {with $b$};
  \node[txt, anchor=south] at (3.425,-0.45) {without $b$};
  \draw[line width=0.5pt] (0,-0.54) -- (0,-0.47) -- (2.15,-0.47) -- (2.15,-0.54);
  \draw[line width=0.5pt] (2.35,-0.54) -- (2.35,-0.47) -- (4.50,-0.47) -- (4.50,-0.54);
  % no flavor boundary: all 22 in one class
  \node[txt, anchor=north west] at (0,-0.62) {no flavor boundary};
  \foreach \i in {0,...,10} {
    \draw[sq, fill=white] (0.2*\i,-1.10) rectangle ++(0.15,0.15);
    \draw[sq, fill=white] (2.35+0.2*\i,-1.10) rectangle ++(0.15,0.15);
  }
  % b-quark boundary: the decays with a b quark moved to another class
  \node[txt, anchor=north west] at (0,-1.22) {$b$-quark boundary};
  \foreach \i in {0,...,10} {
    \draw[sq, fill=cflav!70!black] (0.2*\i,-1.70) rectangle ++(0.15,0.15);
    \draw[sq, fill=white] (2.35+0.2*\i,-1.70) rectangle ++(0.15,0.15);
  }
  % random flavor boundary: the same number of decays, drawn at random
  \node[txt, anchor=north west] at (0,-1.82) {random flavor boundary};
  \foreach \i in {0,1,6,9,10}
    \draw[sq, fill=white] (0.2*\i,-2.30) rectangle ++(0.15,0.15);
  \foreach \i in {2,3,4,5,7,8}
    \draw[sq, fill=cflav!70!black] (0.2*\i,-2.30) rectangle ++(0.15,0.15);
  \foreach \i in {0,4,5,7,8,10}
    \draw[sq, fill=white] (2.35+0.2*\i,-2.30) rectangle ++(0.15,0.15);
  \foreach \i in {1,2,3,6,9}
    \draw[sq, fill=cflav!70!black] (2.35+0.2*\i,-2.30) rectangle ++(0.15,0.15);
  % key
  \draw[sq, fill=cflav!70!black] (0,-2.62) rectangle ++(0.15,0.15);
  \node[txt, anchor=west] (fkey) at (0.25,-2.545) {moved to another class};
\end{scope}
% ================================================================ (c) the pipeline
% (c) starts below the lower of the two columns above, whatever their text heights
\pgfpointanchor{m3}{south}\pgfgetlastxy{\FoneXa}{\FoneYa}%
\pgfpointanchor{fkey}{south}\pgfgetlastxy{\FoneXb}{\FoneYb}%
\pgfmathsetlengthmacro{\FoneYc}{min(\FoneYa,\FoneYb)-0.45cm}%
\coordinate (ctop) at (0pt,\FoneYc);
\begin{scope}[shift={(ctop)}]
  \node[ttl, anchor=north west] at (0,0) {(c) Pretraining and evaluation};
  % what every model of one seed shares
  \node[stage, text width=2.15cm] (seed) at (1.175,-1.96)
    {\textbf{Same for every model of one seed:} the training jets, their order and the
     initial weights, except the output layer};
  % one pretraining per label set: the configuration colours of every later figure
  \node[draw=black!45, fill=black!4, line width=0.6pt, rounded corners=3pt, minimum width=5.15cm,
        minimum height=3.08cm, anchor=north west, inner sep=0pt] (pre) at (2.81,-0.42) {};
  \node[anchor=north west, font=\FoneBold, inner sep=0pt] at (2.96,-0.52) {Pretraining on JetClass-II};
  \node[chip=c188, minimum width=1.13cm, minimum height=0.42cm] at (2.96,-0.95) {\Kfull};
  \node[chip=c162, minimum width=1.13cm, minimum height=0.42cm] at (4.20,-0.95) {\Kres};
  \node[chip=c43,  minimum width=1.13cm, minimum height=0.42cm] at (5.44,-0.95) {\Kmid};
  \node[chip=c17,  minimum width=1.13cm, minimum height=0.42cm] at (6.68,-0.95) {\Kcoarse};
  \node[chip=crand, minimum width=1.70cm, minimum height=0.70cm] at (2.96,-1.47) {random\\partitions};
  \node[chip=cflav, minimum width=3.04cm, minimum height=0.70cm] at (4.77,-1.47) {controlled flavor\\boundaries};
  \node[chip=c162, densely dotted, minimum width=1.13cm, minimum height=0.42cm] at (2.96,-2.32) {\Kres};
  \node[chip=c17,  densely dotted, minimum width=1.13cm, minimum height=0.42cm] at (4.20,-2.32) {\Kcoarse};
  \node[txt, anchor=west] at (5.44,-2.53) {with auxiliary\\mass regression};
  \node[chip=c188, densely dotted, minimum width=1.13cm, minimum height=0.42cm] at (2.96,-2.94) {\Kfull};
  \node[txt, anchor=west] at (4.20,-3.15) {kinematic inputs only};
  % evaluation
  \node[stage, text width=2.4cm] (fz) at (9.68,-0.95) {frozen features:\\linear probe,\\MLP probe};
  \node[stage, text width=2.4cm] (ft) at (9.68,-2.05) {fine-tuning\\of all weights};
  \node[stage, text width=2.4cm, draw=black, fill=white, line width=1.1pt] (sc) at (9.68,-3.15) {trained from\\scratch};
  \node[stage, text width=3.4cm, minimum height=3.2cm] (task) at (13.20,-2.00)
    {\textbf{Downstream tasks}\\JetClass-II: all classes\\JetClass-II: class pairs\\\quad(both class-balanced)\\JetClass\\
     top tagging\\quark/gluon, Pythia\\quark/gluon, Herwig\\jet mass\FoneXbx};
  \draw[flow] (seed.east) -- (seed.east -| pre.west);
  \draw[flow] (fz.west -| pre.east) -- (fz.west);
  \draw[flow] (ft.west -| pre.east) -- (ft.west);
  \draw[flow] (fz.east) -- (fz.east -| task.west);
  \draw[flow] (ft.east) -- (ft.east -| task.west);
  \draw[flow] (sc.east) -- (sc.east -| task.west);
\end{scope}
\end{tikzpicture}
\caption{Design of the study.
(a)~The class hierarchy on the decays of a heavy resonance~$X$ and on the jets from quantum
chromodynamics (QCD); $q$ is a $u$ or $d$ quark, $\ell$ an electron or a muon, and $\tau_h$,
$\tau_e$, $\tau_\mu$ a $\tau$ lepton decaying to hadrons, an electron or a muon. Each box is one
class of the label set above its column, except the QCD box of the \Kfull-class column (Sophon's
label set~\cite{sophon}), which stands for \LitJetClassTwoNQcdClasses{} classes. Each level merges
classes of the level to its left; the lines below name what each merge removes.
(b)~The controls: the first random partition, where boxes with the same letter are one class, and
the random partitions with controlled flavor boundaries on the decays \RandFlavSplitOrbit{} (one
square per decay; filled squares are moved to a second class).
(c)~Models of one seed share the training jets, their order and the initial weights, except the
output layer, the mass regression target, and the first layer and initial weights of the model with
kinematic inputs only (the particle four-momenta). The hierarchy levels have \MlstNRunsWord{} seeds
each, the models with the auxiliary mass regression \MlstMassNRunsWord{}, the random partitions with
controlled flavor boundaries \RandFlavRuns{} and each random partition \RandVtwoRuns{}. Frozen
features, the output of the backbone (the network below the classification head) with its weights held
fixed, feed linear probes (logistic regressions) and multilayer perceptron (MLP) probes; fine-tuning
trains all weights. Colors identify the label sets and controls in every figure; the plots give the three partitions with
controlled flavor boundaries three shades; dotted borders mark the two
variants, which the other figures draw as dotted lines in the color of their label set.}
\label{fig:design}
\end{figure}
 from main.tex; this file holds one figure environment and nothing else.
%
% WHAT main.tex MUST LOAD (this file loads nothing itself):
%   \usepackage{tikz}
%   \usetikzlibrary{arrows.meta}          % the only TikZ library used (arrow tips)
%   \definecolor{c188}{HTML}{1F77B4}      % the configuration colours of experiments/FIGS/
%   \definecolor{c162}{HTML}{FF7F0E}      % mlst_style.py (matplotlib's tab10 in order, as in
%   \definecolor{c43}{HTML}{2CA02C}       % Sophon's figures); the variants are drawn in their
%   \definecolor{c17}{HTML}{D62728}       % label set's colour with a dotted border, as the
%   \definecolor{crand}{HTML}{17BECF}     % plotted figures draw variants dotted
%   \definecolor{cflav}{HTML}{9467BD}
%   % GENERATED by experiments/FIGS/make_tables.py -- do not edit.
% Every line carries: source file :: path inside that file :: first 16 hex of its
% sha256. paper/journal/provenance.json carries the full hash beside the value.
% Regenerate with:  python3 experiments/FIGS/make_tables.py
% Check in CI with: python3 experiments/FIGS/make_tables.py --check
%
% Macro names spell digits as words (162 -> Onesixtwo) because \newcommand takes
% letters only. A number whose source file does not exist yet is NOT emitted --
% see the missing list below, and leave a \pending{} in the text for it.
%
% MISSING INPUTS -- no macro was emitted for these:
%   community benchmarks, bench v3 readout -- experiments/FIGS/data/bench_v3_metrics*/bench_metrics_last.json
%   v2 pretraining rerun: data order and dropout fixed by seed and epoch, identical initial trunk weights, each epoch reading one of five disjoint random parts of the rows of nearly every file, fixed validation sample, checkpoint at the first maximum within the last ten epochs -- experiments/FIGS/data/v2/pretraining/
%   v2 frozen probes on the rerun models, at the first maximum within the last ten epochs and on the model whose weights are the average of the last ten epochs, BatchNorm statistics recomputed, with untrained-network and pooled-embedding reference rows -- experiments/FIGS/data/v2/probe_ladder/
%   v2 paired ratios with run and test-sample errors on the rerun models -- experiments/FIGS/data/v2/paired_errors/
%   v2 the 64- and 30-class vocabularies -- experiments/FIGS/data/v2/levels_64_30/
%   v2 the leave-one-family-out vocabularies -- experiments/FIGS/data/v2/leave_one_family_out/
%   v2 label-tree recovery on the rerun models -- experiments/FIGS/data/v2/label_recovery/
%   v2 the rerun's random partitions and flavour pair -- experiments/FIGS/data/v2/random_partitions/
%   v2 JetClass-II fine-tuning on files held out from pretraining -- experiments/FIGS/data/v2/finetune/
%   v2 fine-tuning from random initialisation and from the self-supervised model -- experiments/FIGS/data/v2/finetune_references/
%   v2 top tagging and quark/gluon tagging (Pythia and Herwig) fine-tuned from every rerun model -- experiments/FIGS/data/v2/benchmarks/
%   v2 the second and third self-supervised pretraining runs -- experiments/FIGS/data/v2/self_supervised/
%   v2 the mass-output models with the loss weight matched in loss share -- experiments/FIGS/data/v2/mass_lambda_matched/
%   v2 anomaly detection on the rerun models, per-run values -- experiments/FIGS/data/v2/anomaly/
%   v2 the open-data fits on the rerun models -- experiments/FIGS/data/v2/real_data/
%   MassLossShareVtwoOnesixtwoMass -- experiments/FIGS/data/v2/mass_lambda_matched/loss_share.json :: shares.162+mass
%   MassLossShareVtwoOnesevenMass -- experiments/FIGS/data/v2/mass_lambda_matched/loss_share.json :: shares.17+mass
%   MassLossShareVtwoOnesevenMassMatched -- experiments/FIGS/data/v2/mass_lambda_matched/loss_share.json :: shares.17+mass_matched
%   DesignVtwoSecondGpu, DesignVtwoSecondGpuRuns -- the grid's job specs put every run index on one product
%   per-arm parameters and MACs -- flops_by_arm.json (see docs/RECORD.md 2.1)
%   fine-tuning read-outs -- experiments/FIGS/data/v2/ft_mlst/: jc2, jc2_qcd
%   top-tagging recipe check -- ft_mlst/topcheck_bench_metrics.json
%   X->bb and X->bs fits -- experiments/FIGS/data/v2/xbx/
%   X->bs from scratch -- experiments/FIGS/data/v2/xbx_scratch/
%   JetClass accuracy bootstrap -- experiments/FIGS/data/v2/budget/jc1_acc_boot_N10000.json
%   MlstFtJcOneErrCoarseOverFullBootEFour -- waits on the JetClass bootstrap at 1e4 (after the equal-stopping read-out)
%   MlstLpBbQgZeroAccDropFull -- waits on pair probes with the quark/gluon inputs zeroed (open item O10, not run)
%   MlstLpBbQgZeroAccDropCoarse -- waits on pair probes with the quark/gluon inputs zeroed (open item O10, not run)
%   XbsRthirtyFullOverScratch -- waits on X->bs fits and from-scratch runs
%   XbsRthirtyFullOverScratchPr -- waits on X->bs fits and from-scratch runs
%   XbsRthirtyCoarseOverFull -- waits on X->bs fits
%   XbsRthirtyCoarseOverFullPr -- waits on X->bs fits
%   XbsRthirtyNodeOverTransferFull -- waits on X->bs fits
%   XbsRthirtyNodeOverTransferFullPr -- waits on X->bs fits
%   XbbRthirtyDirectOverTransferFull -- waits on X->bb fits
%   XbbRthirtyDirectOverTransferFullPr -- waits on X->bb fits
%   XbbRthirtyTransferCoarseOverFull -- waits on X->bb fits
%   XbbRthirtyTransferCoarseOverFullPr -- waits on X->bb fits
%
% SKIPPED TABLES:
%   T4 compute/parameters: no machine-readable source in the repository; docs/RECORD.md 2.1 names /data/results/mtx/flops/flops_by_arm.json, which is not committed

\newcommand{\PendingPretrain}{\pending{pending v2: pretraining rerun: data order and dropout fixed by seed and epoch, identical initial trunk weights, each epoch reading one of five disjoint random parts of the rows of nearly every file, fixed validation sample, checkpoint at the first maximum within the last ten epochs}}  % experiments/FIGS/make_tables.py :: V2_PENDING['Pretrain'] :: sha256 10e6ec08356a4a2b :: a marker, not a number: red in the draft until the input exists
\newcommand{\PendingProbes}{\pending{pending v2: frozen probes on the rerun models, at the first maximum within the last ten epochs and on the model whose weights are the average of the last ten epochs, BatchNorm statistics recomputed, with untrained-network and pooled-embedding reference rows}}  % experiments/FIGS/make_tables.py :: V2_PENDING['Probes'] :: sha256 10e6ec08356a4a2b :: a marker, not a number: red in the draft until the input exists
\newcommand{\PendingPaired}{\pending{pending v2: paired ratios with run and test-sample errors on the rerun models}}  % experiments/FIGS/make_tables.py :: V2_PENDING['Paired'] :: sha256 10e6ec08356a4a2b :: a marker, not a number: red in the draft until the input exists
\newcommand{\PendingLevels}{\pending{pending v2: the 64- and 30-class vocabularies}}  % experiments/FIGS/make_tables.py :: V2_PENDING['Levels'] :: sha256 10e6ec08356a4a2b :: a marker, not a number: red in the draft until the input exists
\newcommand{\PendingLofo}{\pending{pending v2: the leave-one-family-out vocabularies}}  % experiments/FIGS/make_tables.py :: V2_PENDING['Lofo'] :: sha256 10e6ec08356a4a2b :: a marker, not a number: red in the draft until the input exists
\newcommand{\PendingRecovery}{\pending{pending v2: label-tree recovery on the rerun models}}  % experiments/FIGS/make_tables.py :: V2_PENDING['Recovery'] :: sha256 10e6ec08356a4a2b :: a marker, not a number: red in the draft until the input exists
\newcommand{\PendingRandom}{\pending{pending v2: the rerun's random partitions and flavour pair}}  % experiments/FIGS/make_tables.py :: V2_PENDING['Random'] :: sha256 10e6ec08356a4a2b :: a marker, not a number: red in the draft until the input exists
\newcommand{\PendingFtHeldout}{\pending{pending v2: JetClass-II fine-tuning on files held out from pretraining}}  % experiments/FIGS/make_tables.py :: V2_PENDING['FtHeldout'] :: sha256 10e6ec08356a4a2b :: a marker, not a number: red in the draft until the input exists
\newcommand{\PendingFtRefs}{\pending{pending v2: fine-tuning from random initialisation and from the self-supervised model}}  % experiments/FIGS/make_tables.py :: V2_PENDING['FtRefs'] :: sha256 10e6ec08356a4a2b :: a marker, not a number: red in the draft until the input exists
\newcommand{\PendingBench}{\pending{pending v2: top tagging and quark/gluon tagging (Pythia and Herwig) fine-tuned from every rerun model}}  % experiments/FIGS/make_tables.py :: V2_PENDING['Bench'] :: sha256 10e6ec08356a4a2b :: a marker, not a number: red in the draft until the input exists
\newcommand{\PendingSsl}{\pending{pending v2: the second and third self-supervised pretraining runs}}  % experiments/FIGS/make_tables.py :: V2_PENDING['Ssl'] :: sha256 10e6ec08356a4a2b :: a marker, not a number: red in the draft until the input exists
\newcommand{\PendingMassLambda}{\pending{pending v2: the mass-output models with the loss weight matched in loss share}}  % experiments/FIGS/make_tables.py :: V2_PENDING['MassLambda'] :: sha256 10e6ec08356a4a2b :: a marker, not a number: red in the draft until the input exists
\newcommand{\PendingAnomaly}{\pending{pending v2: anomaly detection on the rerun models, per-run values}}  % experiments/FIGS/make_tables.py :: V2_PENDING['Anomaly'] :: sha256 10e6ec08356a4a2b :: a marker, not a number: red in the draft until the input exists
\newcommand{\PendingRealData}{\pending{pending v2: the open-data fits on the rerun models}}  % experiments/FIGS/make_tables.py :: V2_PENDING['RealData'] :: sha256 10e6ec08356a4a2b :: a marker, not a number: red in the draft until the input exists
\newcommand{\PendingConclusion}{\pending{pending v2: written once every input above is in}}  % experiments/FIGS/make_tables.py :: V2_PENDING['Conclusion'] :: sha256 10e6ec08356a4a2b :: a marker, not a number: red in the draft until the input exists
\newcommand{\AojChainDataset}{178}  % experiments/FIGS/data/aoj_checks_v2/selection_chain_v2.json :: n_dataset :: sha256 729d83085aba25f8
\newcommand{\AojChainStaged}{12{,}589{,}467}  % experiments/FIGS/data/aoj_checks_v2/selection_chain_v2.json :: n_staged :: sha256 729d83085aba25f8
\newcommand{\AojChainRhoWindow}{9{,}441{,}209}  % experiments/FIGS/data/aoj_checks_v2/selection_chain_v2.json :: n_rho_window :: sha256 729d83085aba25f8
\newcommand{\AojNJetsFit}{4{,}194{,}846}  % experiments/FIGS/data/aoj_checks_v2/selection_chain_v2.json :: n_fit :: sha256 729d83085aba25f8
\newcommand{\AojChainPtBins}{8}  % experiments/FIGS/data/aoj_checks_v2/selection_chain_v2.json :: n_pt_bins :: sha256 729d83085aba25f8
\newcommand{\AojChainEtaMax}{2.40}  % experiments/FIGS/data/aoj_checks_v2/selection_chain_v2.json :: abs_eta_max :: sha256 729d83085aba25f8
\newcommand{\MassLossShareVtwoOnesixtwoMass}{\pending{pending: realised mass-loss share in the rerun, 162 classes}}  % experiments/FIGS/make_tables.py :: SLOTS[MassLossShareVtwoOnesixtwoMass] :: sha256 10e6ec08356a4a2b :: a marker, not a number: experiments/FIGS/data/v2/mass_lambda_matched/loss_share.json :: shares.162+mass does not exist yet
\newcommand{\MassLossShareVtwoOnesevenMass}{\pending{pending: realised mass-loss share in the rerun, 17 classes}}  % experiments/FIGS/make_tables.py :: SLOTS[MassLossShareVtwoOnesevenMass] :: sha256 10e6ec08356a4a2b :: a marker, not a number: experiments/FIGS/data/v2/mass_lambda_matched/loss_share.json :: shares.17+mass does not exist yet
\newcommand{\MassLossShareVtwoOnesevenMassMatched}{\pending{pending: realised mass-loss share in the rerun, 17 classes, matched weight}}  % experiments/FIGS/make_tables.py :: SLOTS[MassLossShareVtwoOnesevenMassMatched] :: sha256 10e6ec08356a4a2b :: a marker, not a number: experiments/FIGS/data/v2/mass_lambda_matched/loss_share.json :: shares.17+mass_matched does not exist yet
\newcommand{\FtCoverageEThreeNClasses}{160}  % experiments/FT/data/jc2_v2_manifest.json :: class_coverage['train_N1000_s1.parquet'].n_classes_present :: sha256 4bd68655919b14e9
\newcommand{\FtCoverageEFourNClasses}{183}  % experiments/FT/data/jc2_v2_manifest.json :: class_coverage['train_N10000_s1.parquet'].n_classes_present :: sha256 4bd68655919b14e9
\newcommand{\FtCoverageEThreeMedian}{4}  % experiments/FT/data/jc2_v2_manifest.json :: class_coverage['train_N1000_s1.parquet'].median_per_present_class :: sha256 4bd68655919b14e9
\newcommand{\FtCoverageEThreeOneJet}{23}  % experiments/FT/data/jc2_v2_manifest.json :: class_coverage['train_N1000_s1.parquet'].n_classes_with_one_jet :: sha256 4bd68655919b14e9
\newcommand{\FtCoverageVoneEThreeNClasses}{161}  % experiments/FT/data/jc2_v1_coverage.json :: class_coverage['train_N1000_s1.parquet'].n_classes_present :: sha256 82652f57f708dc23
\newcommand{\FtCoverageVoneEThreeMedian}{5}  % experiments/FT/data/jc2_v1_coverage.json :: class_coverage['train_N1000_s1.parquet'].median_per_present_class :: sha256 82652f57f708dc23
\newcommand{\FtCoverageVoneEThreeOneJet}{25}  % experiments/FT/data/jc2_v1_coverage.json :: class_coverage['train_N1000_s1.parquet'].n_classes_with_one_jet :: sha256 82652f57f708dc23
\newcommand{\ProbeNJets}{2{,}000{,}000}  % experiments/FIGS/data/probe_ladder_v2_mlp2/analysis/seed_level_results.json :: provenance.n_jets_total :: sha256 53e1ec81343ac2c1
\newcommand{\ProbeNSeeds}{5}  % experiments/FIGS/data/probe_ladder_v2_mlp2/analysis/seed_level_results.json :: seeds_used (length) :: sha256 53e1ec81343ac2c1
\newcommand{\ProbeNSeedsWord}{five}  % experiments/FIGS/data/probe_ladder_v2_mlp2/analysis/seed_level_results.json :: seeds_used (length) :: sha256 53e1ec81343ac2c1 :: pretraining runs per vocabulary, as a word for the prose
\newcommand{\ProbeRowAlign}{\texttt{a9449fd420c473f6}}  % experiments/FIGS/data/probe_ladder_v2_mlp2/analysis/seed_level_results.json :: provenance.row_alignment_sha256 (first 16) :: sha256 53e1ec81343ac2c1
\newcommand{\ProbeEpsS}{90\%}  % experiments/FIGS/data/probe_ladder_v2_mlp2/analysis/seed_level_results.json :: levels.bvc_4prong.linear[0].headline_eps_s :: sha256 53e1ec81343ac2c1
\newcommand{\ProbeAucBvcFourprongLinearOneeighteight}{0.99320}  % experiments/FIGS/data/probe_ladder_v2_mlp2/analysis/seed_level_results.json :: levels.bvc_4prong.linear[0].mean_auc :: sha256 53e1ec81343ac2c1 :: 5 seeds, 0 saturated at AUC=1
\newcommand{\ProbeAucSdBvcFourprongLinearOneeighteight}{0.00027}  % experiments/FIGS/data/probe_ladder_v2_mlp2/analysis/seed_level_results.json :: table[bvc_4prong,linear,188].auc (SD over seeds, ddof=1) :: sha256 53e1ec81343ac2c1
\newcommand{\ProbeLogOneMinusAucBvcFourprongLinearOneeighteight}{\ensuremath{-}4.9915}  % experiments/FIGS/data/probe_ladder_v2_mlp2/analysis/seed_level_results.json :: levels.bvc_4prong.linear[0].mean :: sha256 53e1ec81343ac2c1 :: natural log of 1-AUC, lower is better
\newcommand{\ProbeLogOneMinusAucSdBvcFourprongLinearOneeighteight}{0.0389}  % experiments/FIGS/data/probe_ladder_v2_mlp2/analysis/seed_level_results.json :: levels.bvc_4prong.linear[0].seed_sd :: sha256 53e1ec81343ac2c1
\newcommand{\ProbeOmaBvcFourprongLinearOneeighteight}{\ensuremath{(6.80 \pm 0.27)\times10^{-3}}}  % experiments/FIGS/data/probe_ladder_v2_mlp2/analysis/seed_level_results.json :: table[bvc_4prong,linear,188].auc :: sha256 53e1ec81343ac2c1 :: 1-AUC, mean +- SD over 5 seeds
\newcommand{\ProbeBkgLeftBvcFourprongLinearOneeighteight}{14.0}  % experiments/FIGS/data/probe_ladder_v2_mlp2/analysis/seed_level_results.json :: levels.bvc_4prong.linear[0].rejection_points['0.90'].mean_n_bkg_pass :: sha256 53e1ec81343ac2c1 :: background jets passing the cut, mean over seeds
\newcommand{\ProbeRejBvcFourprongLinearOneeighteight}{227\,\ensuremath{\pm}\,23}  % experiments/FIGS/data/probe_ladder_v2_mlp2/analysis/seed_level_results.json :: table[bvc_4prong,linear,188].rejection_points['0.90'].rejection :: sha256 53e1ec81343ac2c1 :: at 90% signal efficiency, mean +- SD over seeds; 0 of 5 seeds at the cap
\newcommand{\ProbeAucBvcFourprongLinearOnesixtwo}{0.99311}  % experiments/FIGS/data/probe_ladder_v2_mlp2/analysis/seed_level_results.json :: levels.bvc_4prong.linear[1].mean_auc :: sha256 53e1ec81343ac2c1 :: 5 seeds, 0 saturated at AUC=1
\newcommand{\ProbeAucSdBvcFourprongLinearOnesixtwo}{0.00018}  % experiments/FIGS/data/probe_ladder_v2_mlp2/analysis/seed_level_results.json :: table[bvc_4prong,linear,162].auc (SD over seeds, ddof=1) :: sha256 53e1ec81343ac2c1
\newcommand{\ProbeLogOneMinusAucBvcFourprongLinearOnesixtwo}{\ensuremath{-}4.9777}  % experiments/FIGS/data/probe_ladder_v2_mlp2/analysis/seed_level_results.json :: levels.bvc_4prong.linear[1].mean :: sha256 53e1ec81343ac2c1 :: natural log of 1-AUC, lower is better
\newcommand{\ProbeLogOneMinusAucSdBvcFourprongLinearOnesixtwo}{0.0264}  % experiments/FIGS/data/probe_ladder_v2_mlp2/analysis/seed_level_results.json :: levels.bvc_4prong.linear[1].seed_sd :: sha256 53e1ec81343ac2c1
\newcommand{\ProbeOmaBvcFourprongLinearOnesixtwo}{\ensuremath{(6.89 \pm 0.18)\times10^{-3}}}  % experiments/FIGS/data/probe_ladder_v2_mlp2/analysis/seed_level_results.json :: table[bvc_4prong,linear,162].auc :: sha256 53e1ec81343ac2c1 :: 1-AUC, mean +- SD over 5 seeds
\newcommand{\ProbeBkgLeftBvcFourprongLinearOnesixtwo}{14.4}  % experiments/FIGS/data/probe_ladder_v2_mlp2/analysis/seed_level_results.json :: levels.bvc_4prong.linear[1].rejection_points['0.90'].mean_n_bkg_pass :: sha256 53e1ec81343ac2c1 :: background jets passing the cut, mean over seeds
\newcommand{\ProbeRejBvcFourprongLinearOnesixtwo}{220\,\ensuremath{\pm}\,40}  % experiments/FIGS/data/probe_ladder_v2_mlp2/analysis/seed_level_results.json :: table[bvc_4prong,linear,162].rejection_points['0.90'].rejection :: sha256 53e1ec81343ac2c1 :: at 90% signal efficiency, mean +- SD over seeds; 0 of 5 seeds at the cap
\newcommand{\ProbeAucBvcFourprongLinearFourthree}{0.99250}  % experiments/FIGS/data/probe_ladder_v2_mlp2/analysis/seed_level_results.json :: levels.bvc_4prong.linear[2].mean_auc :: sha256 53e1ec81343ac2c1 :: 5 seeds, 0 saturated at AUC=1
\newcommand{\ProbeAucSdBvcFourprongLinearFourthree}{0.00049}  % experiments/FIGS/data/probe_ladder_v2_mlp2/analysis/seed_level_results.json :: table[bvc_4prong,linear,43].auc (SD over seeds, ddof=1) :: sha256 53e1ec81343ac2c1
\newcommand{\ProbeLogOneMinusAucBvcFourprongLinearFourthree}{\ensuremath{-}4.8946}  % experiments/FIGS/data/probe_ladder_v2_mlp2/analysis/seed_level_results.json :: levels.bvc_4prong.linear[2].mean :: sha256 53e1ec81343ac2c1 :: natural log of 1-AUC, lower is better
\newcommand{\ProbeLogOneMinusAucSdBvcFourprongLinearFourthree}{0.0646}  % experiments/FIGS/data/probe_ladder_v2_mlp2/analysis/seed_level_results.json :: levels.bvc_4prong.linear[2].seed_sd :: sha256 53e1ec81343ac2c1
\newcommand{\ProbeOmaBvcFourprongLinearFourthree}{\ensuremath{(7.5 \pm 0.5)\times10^{-3}}}  % experiments/FIGS/data/probe_ladder_v2_mlp2/analysis/seed_level_results.json :: table[bvc_4prong,linear,43].auc :: sha256 53e1ec81343ac2c1 :: 1-AUC, mean +- SD over 5 seeds
\newcommand{\ProbeBkgLeftBvcFourprongLinearFourthree}{16.6}  % experiments/FIGS/data/probe_ladder_v2_mlp2/analysis/seed_level_results.json :: levels.bvc_4prong.linear[2].rejection_points['0.90'].mean_n_bkg_pass :: sha256 53e1ec81343ac2c1 :: background jets passing the cut, mean over seeds
\newcommand{\ProbeRejBvcFourprongLinearFourthree}{192\,\ensuremath{\pm}\,20}  % experiments/FIGS/data/probe_ladder_v2_mlp2/analysis/seed_level_results.json :: table[bvc_4prong,linear,43].rejection_points['0.90'].rejection :: sha256 53e1ec81343ac2c1 :: at 90% signal efficiency, mean +- SD over seeds; 0 of 5 seeds at the cap
\newcommand{\ProbeAucBvcFourprongLinearOneseven}{0.96936}  % experiments/FIGS/data/probe_ladder_v2_mlp2/analysis/seed_level_results.json :: levels.bvc_4prong.linear[3].mean_auc :: sha256 53e1ec81343ac2c1 :: 5 seeds, 0 saturated at AUC=1
\newcommand{\ProbeAucSdBvcFourprongLinearOneseven}{0.00514}  % experiments/FIGS/data/probe_ladder_v2_mlp2/analysis/seed_level_results.json :: table[bvc_4prong,linear,17].auc (SD over seeds, ddof=1) :: sha256 53e1ec81343ac2c1
\newcommand{\ProbeLogOneMinusAucBvcFourprongLinearOneseven}{\ensuremath{-}3.4967}  % experiments/FIGS/data/probe_ladder_v2_mlp2/analysis/seed_level_results.json :: levels.bvc_4prong.linear[3].mean :: sha256 53e1ec81343ac2c1 :: natural log of 1-AUC, lower is better
\newcommand{\ProbeLogOneMinusAucSdBvcFourprongLinearOneseven}{0.1667}  % experiments/FIGS/data/probe_ladder_v2_mlp2/analysis/seed_level_results.json :: levels.bvc_4prong.linear[3].seed_sd :: sha256 53e1ec81343ac2c1
\newcommand{\ProbeOmaBvcFourprongLinearOneseven}{\ensuremath{(3.1 \pm 0.5)\times10^{-2}}}  % experiments/FIGS/data/probe_ladder_v2_mlp2/analysis/seed_level_results.json :: table[bvc_4prong,linear,17].auc :: sha256 53e1ec81343ac2c1 :: 1-AUC, mean +- SD over 5 seeds
\newcommand{\ProbeBkgLeftBvcFourprongLinearOneseven}{137.4}  % experiments/FIGS/data/probe_ladder_v2_mlp2/analysis/seed_level_results.json :: levels.bvc_4prong.linear[3].rejection_points['0.90'].mean_n_bkg_pass :: sha256 53e1ec81343ac2c1 :: background jets passing the cut, mean over seeds
\newcommand{\ProbeRejBvcFourprongLinearOneseven}{24\,\ensuremath{\pm}\,6}  % experiments/FIGS/data/probe_ladder_v2_mlp2/analysis/seed_level_results.json :: table[bvc_4prong,linear,17].rejection_points['0.90'].rejection :: sha256 53e1ec81343ac2c1 :: at 90% signal efficiency, mean +- SD over seeds; 0 of 5 seeds at the cap
\newcommand{\ProbeOmaRatioBvcFourprongLinearOnesixtwoOverOneeighteight}{1.01}  % experiments/FIGS/data/probe_ladder_v2_mlp2/analysis/seed_level_results.json :: table[bvc_4prong,linear,162/188].auc :: sha256 53e1ec81343ac2c1 :: ratio of the seed means of 1-AUC, coarser over finer
\newcommand{\ProbeOmaRatioBvcFourprongLinearFourthreeOverOneeighteight}{1.10}  % experiments/FIGS/data/probe_ladder_v2_mlp2/analysis/seed_level_results.json :: table[bvc_4prong,linear,43/188].auc :: sha256 53e1ec81343ac2c1 :: ratio of the seed means of 1-AUC, coarser over finer
\newcommand{\ProbeOmaRatioBvcFourprongLinearOnesevenOverOneeighteight}{4.5}  % experiments/FIGS/data/probe_ladder_v2_mlp2/analysis/seed_level_results.json :: table[bvc_4prong,linear,17/188].auc :: sha256 53e1ec81343ac2c1 :: ratio of the seed means of 1-AUC, coarser over finer
\newcommand{\ProbeOmaRatioBvcFourprongLinearFourthreeOverOnesixtwo}{1.09}  % experiments/FIGS/data/probe_ladder_v2_mlp2/analysis/seed_level_results.json :: table[bvc_4prong,linear,43/162].auc :: sha256 53e1ec81343ac2c1 :: ratio of the seed means of 1-AUC, coarser over finer
\newcommand{\ProbeOmaRatioBvcFourprongLinearOnesevenOverOnesixtwo}{4.4}  % experiments/FIGS/data/probe_ladder_v2_mlp2/analysis/seed_level_results.json :: table[bvc_4prong,linear,17/162].auc :: sha256 53e1ec81343ac2c1 :: ratio of the seed means of 1-AUC, coarser over finer
\newcommand{\ProbeOmaRatioBvcFourprongLinearOnesevenOverFourthree}{4.1}  % experiments/FIGS/data/probe_ladder_v2_mlp2/analysis/seed_level_results.json :: table[bvc_4prong,linear,17/43].auc :: sha256 53e1ec81343ac2c1 :: ratio of the seed means of 1-AUC, coarser over finer
\newcommand{\ProbeAucBvcFourprongMlpOneeighteight}{0.99316}  % experiments/FIGS/data/probe_ladder_v2_mlp2/analysis/seed_level_results.json :: levels.bvc_4prong.mlp[0].mean_auc :: sha256 53e1ec81343ac2c1 :: 5 seeds, 0 saturated at AUC=1
\newcommand{\ProbeAucSdBvcFourprongMlpOneeighteight}{0.00020}  % experiments/FIGS/data/probe_ladder_v2_mlp2/analysis/seed_level_results.json :: table[bvc_4prong,mlp,188].auc (SD over seeds, ddof=1) :: sha256 53e1ec81343ac2c1
\newcommand{\ProbeLogOneMinusAucBvcFourprongMlpOneeighteight}{\ensuremath{-}4.9857}  % experiments/FIGS/data/probe_ladder_v2_mlp2/analysis/seed_level_results.json :: levels.bvc_4prong.mlp[0].mean :: sha256 53e1ec81343ac2c1 :: natural log of 1-AUC, lower is better
\newcommand{\ProbeLogOneMinusAucSdBvcFourprongMlpOneeighteight}{0.0290}  % experiments/FIGS/data/probe_ladder_v2_mlp2/analysis/seed_level_results.json :: levels.bvc_4prong.mlp[0].seed_sd :: sha256 53e1ec81343ac2c1
\newcommand{\ProbeOmaBvcFourprongMlpOneeighteight}{\ensuremath{(6.84 \pm 0.20)\times10^{-3}}}  % experiments/FIGS/data/probe_ladder_v2_mlp2/analysis/seed_level_results.json :: table[bvc_4prong,mlp,188].auc :: sha256 53e1ec81343ac2c1 :: 1-AUC, mean +- SD over 5 seeds
\newcommand{\ProbeBkgLeftBvcFourprongMlpOneeighteight}{15.2}  % experiments/FIGS/data/probe_ladder_v2_mlp2/analysis/seed_level_results.json :: levels.bvc_4prong.mlp[0].rejection_points['0.90'].mean_n_bkg_pass :: sha256 53e1ec81343ac2c1 :: background jets passing the cut, mean over seeds
\newcommand{\ProbeRejBvcFourprongMlpOneeighteight}{212\,\ensuremath{\pm}\,34}  % experiments/FIGS/data/probe_ladder_v2_mlp2/analysis/seed_level_results.json :: table[bvc_4prong,mlp,188].rejection_points['0.90'].rejection :: sha256 53e1ec81343ac2c1 :: at 90% signal efficiency, mean +- SD over seeds; 0 of 5 seeds at the cap
\newcommand{\ProbeAucBvcFourprongMlpOnesixtwo}{0.99318}  % experiments/FIGS/data/probe_ladder_v2_mlp2/analysis/seed_level_results.json :: levels.bvc_4prong.mlp[1].mean_auc :: sha256 53e1ec81343ac2c1 :: 5 seeds, 0 saturated at AUC=1
\newcommand{\ProbeAucSdBvcFourprongMlpOnesixtwo}{0.00020}  % experiments/FIGS/data/probe_ladder_v2_mlp2/analysis/seed_level_results.json :: table[bvc_4prong,mlp,162].auc (SD over seeds, ddof=1) :: sha256 53e1ec81343ac2c1
\newcommand{\ProbeLogOneMinusAucBvcFourprongMlpOnesixtwo}{\ensuremath{-}4.9886}  % experiments/FIGS/data/probe_ladder_v2_mlp2/analysis/seed_level_results.json :: levels.bvc_4prong.mlp[1].mean :: sha256 53e1ec81343ac2c1 :: natural log of 1-AUC, lower is better
\newcommand{\ProbeLogOneMinusAucSdBvcFourprongMlpOnesixtwo}{0.0298}  % experiments/FIGS/data/probe_ladder_v2_mlp2/analysis/seed_level_results.json :: levels.bvc_4prong.mlp[1].seed_sd :: sha256 53e1ec81343ac2c1
\newcommand{\ProbeOmaBvcFourprongMlpOnesixtwo}{\ensuremath{(6.82 \pm 0.20)\times10^{-3}}}  % experiments/FIGS/data/probe_ladder_v2_mlp2/analysis/seed_level_results.json :: table[bvc_4prong,mlp,162].auc :: sha256 53e1ec81343ac2c1 :: 1-AUC, mean +- SD over 5 seeds
\newcommand{\ProbeBkgLeftBvcFourprongMlpOnesixtwo}{16.6}  % experiments/FIGS/data/probe_ladder_v2_mlp2/analysis/seed_level_results.json :: levels.bvc_4prong.mlp[1].rejection_points['0.90'].mean_n_bkg_pass :: sha256 53e1ec81343ac2c1 :: background jets passing the cut, mean over seeds
\newcommand{\ProbeRejBvcFourprongMlpOnesixtwo}{200\,\ensuremath{\pm}\,40}  % experiments/FIGS/data/probe_ladder_v2_mlp2/analysis/seed_level_results.json :: table[bvc_4prong,mlp,162].rejection_points['0.90'].rejection :: sha256 53e1ec81343ac2c1 :: at 90% signal efficiency, mean +- SD over seeds; 0 of 5 seeds at the cap
\newcommand{\ProbeAucBvcFourprongMlpFourthree}{0.99232}  % experiments/FIGS/data/probe_ladder_v2_mlp2/analysis/seed_level_results.json :: levels.bvc_4prong.mlp[2].mean_auc :: sha256 53e1ec81343ac2c1 :: 5 seeds, 0 saturated at AUC=1
\newcommand{\ProbeAucSdBvcFourprongMlpFourthree}{0.00058}  % experiments/FIGS/data/probe_ladder_v2_mlp2/analysis/seed_level_results.json :: table[bvc_4prong,mlp,43].auc (SD over seeds, ddof=1) :: sha256 53e1ec81343ac2c1
\newcommand{\ProbeLogOneMinusAucBvcFourprongMlpFourthree}{\ensuremath{-}4.8719}  % experiments/FIGS/data/probe_ladder_v2_mlp2/analysis/seed_level_results.json :: levels.bvc_4prong.mlp[2].mean :: sha256 53e1ec81343ac2c1 :: natural log of 1-AUC, lower is better
\newcommand{\ProbeLogOneMinusAucSdBvcFourprongMlpFourthree}{0.0734}  % experiments/FIGS/data/probe_ladder_v2_mlp2/analysis/seed_level_results.json :: levels.bvc_4prong.mlp[2].seed_sd :: sha256 53e1ec81343ac2c1
\newcommand{\ProbeOmaBvcFourprongMlpFourthree}{\ensuremath{(7.7 \pm 0.6)\times10^{-3}}}  % experiments/FIGS/data/probe_ladder_v2_mlp2/analysis/seed_level_results.json :: table[bvc_4prong,mlp,43].auc :: sha256 53e1ec81343ac2c1 :: 1-AUC, mean +- SD over 5 seeds
\newcommand{\ProbeBkgLeftBvcFourprongMlpFourthree}{17.4}  % experiments/FIGS/data/probe_ladder_v2_mlp2/analysis/seed_level_results.json :: levels.bvc_4prong.mlp[2].rejection_points['0.90'].mean_n_bkg_pass :: sha256 53e1ec81343ac2c1 :: background jets passing the cut, mean over seeds
\newcommand{\ProbeRejBvcFourprongMlpFourthree}{184\,\ensuremath{\pm}\,23}  % experiments/FIGS/data/probe_ladder_v2_mlp2/analysis/seed_level_results.json :: table[bvc_4prong,mlp,43].rejection_points['0.90'].rejection :: sha256 53e1ec81343ac2c1 :: at 90% signal efficiency, mean +- SD over seeds; 0 of 5 seeds at the cap
\newcommand{\ProbeAucBvcFourprongMlpOneseven}{0.97766}  % experiments/FIGS/data/probe_ladder_v2_mlp2/analysis/seed_level_results.json :: levels.bvc_4prong.mlp[3].mean_auc :: sha256 53e1ec81343ac2c1 :: 5 seeds, 0 saturated at AUC=1
\newcommand{\ProbeAucSdBvcFourprongMlpOneseven}{0.00310}  % experiments/FIGS/data/probe_ladder_v2_mlp2/analysis/seed_level_results.json :: table[bvc_4prong,mlp,17].auc (SD over seeds, ddof=1) :: sha256 53e1ec81343ac2c1
\newcommand{\ProbeLogOneMinusAucBvcFourprongMlpOneseven}{\ensuremath{-}3.8087}  % experiments/FIGS/data/probe_ladder_v2_mlp2/analysis/seed_level_results.json :: levels.bvc_4prong.mlp[3].mean :: sha256 53e1ec81343ac2c1 :: natural log of 1-AUC, lower is better
\newcommand{\ProbeLogOneMinusAucSdBvcFourprongMlpOneseven}{0.1366}  % experiments/FIGS/data/probe_ladder_v2_mlp2/analysis/seed_level_results.json :: levels.bvc_4prong.mlp[3].seed_sd :: sha256 53e1ec81343ac2c1
\newcommand{\ProbeOmaBvcFourprongMlpOneseven}{\ensuremath{(2.23 \pm 0.31)\times10^{-2}}}  % experiments/FIGS/data/probe_ladder_v2_mlp2/analysis/seed_level_results.json :: table[bvc_4prong,mlp,17].auc :: sha256 53e1ec81343ac2c1 :: 1-AUC, mean +- SD over 5 seeds
\newcommand{\ProbeBkgLeftBvcFourprongMlpOneseven}{99.0}  % experiments/FIGS/data/probe_ladder_v2_mlp2/analysis/seed_level_results.json :: levels.bvc_4prong.mlp[3].rejection_points['0.90'].mean_n_bkg_pass :: sha256 53e1ec81343ac2c1 :: background jets passing the cut, mean over seeds
\newcommand{\ProbeRejBvcFourprongMlpOneseven}{33\,\ensuremath{\pm}\,5}  % experiments/FIGS/data/probe_ladder_v2_mlp2/analysis/seed_level_results.json :: table[bvc_4prong,mlp,17].rejection_points['0.90'].rejection :: sha256 53e1ec81343ac2c1 :: at 90% signal efficiency, mean +- SD over seeds; 0 of 5 seeds at the cap
\newcommand{\ProbeOmaRatioBvcFourprongMlpOnesixtwoOverOneeighteight}{1.00}  % experiments/FIGS/data/probe_ladder_v2_mlp2/analysis/seed_level_results.json :: table[bvc_4prong,mlp,162/188].auc :: sha256 53e1ec81343ac2c1 :: ratio of the seed means of 1-AUC, coarser over finer
\newcommand{\ProbeOmaRatioBvcFourprongMlpFourthreeOverOneeighteight}{1.12}  % experiments/FIGS/data/probe_ladder_v2_mlp2/analysis/seed_level_results.json :: table[bvc_4prong,mlp,43/188].auc :: sha256 53e1ec81343ac2c1 :: ratio of the seed means of 1-AUC, coarser over finer
\newcommand{\ProbeOmaRatioBvcFourprongMlpOnesevenOverOneeighteight}{3.3}  % experiments/FIGS/data/probe_ladder_v2_mlp2/analysis/seed_level_results.json :: table[bvc_4prong,mlp,17/188].auc :: sha256 53e1ec81343ac2c1 :: ratio of the seed means of 1-AUC, coarser over finer
\newcommand{\ProbeOmaRatioBvcFourprongMlpFourthreeOverOnesixtwo}{1.13}  % experiments/FIGS/data/probe_ladder_v2_mlp2/analysis/seed_level_results.json :: table[bvc_4prong,mlp,43/162].auc :: sha256 53e1ec81343ac2c1 :: ratio of the seed means of 1-AUC, coarser over finer
\newcommand{\ProbeOmaRatioBvcFourprongMlpOnesevenOverOnesixtwo}{3.3}  % experiments/FIGS/data/probe_ladder_v2_mlp2/analysis/seed_level_results.json :: table[bvc_4prong,mlp,17/162].auc :: sha256 53e1ec81343ac2c1 :: ratio of the seed means of 1-AUC, coarser over finer
\newcommand{\ProbeOmaRatioBvcFourprongMlpOnesevenOverFourthree}{2.9}  % experiments/FIGS/data/probe_ladder_v2_mlp2/analysis/seed_level_results.json :: table[bvc_4prong,mlp,17/43].auc :: sha256 53e1ec81343ac2c1 :: ratio of the seed means of 1-AUC, coarser over finer
\newcommand{\ProbeAucBvcQcdLinearOneeighteight}{0.98519}  % experiments/FIGS/data/probe_ladder_v2_mlp2/analysis/seed_level_results.json :: levels.bvc_qcd.linear[0].mean_auc :: sha256 53e1ec81343ac2c1 :: 5 seeds, 0 saturated at AUC=1
\newcommand{\ProbeAucSdBvcQcdLinearOneeighteight}{0.00061}  % experiments/FIGS/data/probe_ladder_v2_mlp2/analysis/seed_level_results.json :: table[bvc_qcd,linear,188].auc (SD over seeds, ddof=1) :: sha256 53e1ec81343ac2c1
\newcommand{\ProbeLogOneMinusAucBvcQcdLinearOneeighteight}{\ensuremath{-}4.2133}  % experiments/FIGS/data/probe_ladder_v2_mlp2/analysis/seed_level_results.json :: levels.bvc_qcd.linear[0].mean :: sha256 53e1ec81343ac2c1 :: natural log of 1-AUC, lower is better
\newcommand{\ProbeLogOneMinusAucSdBvcQcdLinearOneeighteight}{0.0422}  % experiments/FIGS/data/probe_ladder_v2_mlp2/analysis/seed_level_results.json :: levels.bvc_qcd.linear[0].seed_sd :: sha256 53e1ec81343ac2c1
\newcommand{\ProbeOmaBvcQcdLinearOneeighteight}{\ensuremath{(1.48 \pm 0.06)\times10^{-2}}}  % experiments/FIGS/data/probe_ladder_v2_mlp2/analysis/seed_level_results.json :: table[bvc_qcd,linear,188].auc :: sha256 53e1ec81343ac2c1 :: 1-AUC, mean +- SD over 5 seeds
\newcommand{\ProbeBkgLeftBvcQcdLinearOneeighteight}{96.4}  % experiments/FIGS/data/probe_ladder_v2_mlp2/analysis/seed_level_results.json :: levels.bvc_qcd.linear[0].rejection_points['0.90'].mean_n_bkg_pass :: sha256 53e1ec81343ac2c1 :: background jets passing the cut, mean over seeds
\newcommand{\ProbeRejBvcQcdLinearOneeighteight}{56\,\ensuremath{\pm}\,5}  % experiments/FIGS/data/probe_ladder_v2_mlp2/analysis/seed_level_results.json :: table[bvc_qcd,linear,188].rejection_points['0.90'].rejection :: sha256 53e1ec81343ac2c1 :: at 90% signal efficiency, mean +- SD over seeds; 0 of 5 seeds at the cap
\newcommand{\ProbeAucBvcQcdLinearOnesixtwo}{0.98305}  % experiments/FIGS/data/probe_ladder_v2_mlp2/analysis/seed_level_results.json :: levels.bvc_qcd.linear[1].mean_auc :: sha256 53e1ec81343ac2c1 :: 5 seeds, 0 saturated at AUC=1
\newcommand{\ProbeAucSdBvcQcdLinearOnesixtwo}{0.00024}  % experiments/FIGS/data/probe_ladder_v2_mlp2/analysis/seed_level_results.json :: table[bvc_qcd,linear,162].auc (SD over seeds, ddof=1) :: sha256 53e1ec81343ac2c1
\newcommand{\ProbeLogOneMinusAucBvcQcdLinearOnesixtwo}{\ensuremath{-}4.0774}  % experiments/FIGS/data/probe_ladder_v2_mlp2/analysis/seed_level_results.json :: levels.bvc_qcd.linear[1].mean :: sha256 53e1ec81343ac2c1 :: natural log of 1-AUC, lower is better
\newcommand{\ProbeLogOneMinusAucSdBvcQcdLinearOnesixtwo}{0.0142}  % experiments/FIGS/data/probe_ladder_v2_mlp2/analysis/seed_level_results.json :: levels.bvc_qcd.linear[1].seed_sd :: sha256 53e1ec81343ac2c1
\newcommand{\ProbeOmaBvcQcdLinearOnesixtwo}{\ensuremath{(1.695 \pm 0.024)\times10^{-2}}}  % experiments/FIGS/data/probe_ladder_v2_mlp2/analysis/seed_level_results.json :: table[bvc_qcd,linear,162].auc :: sha256 53e1ec81343ac2c1 :: 1-AUC, mean +- SD over 5 seeds
\newcommand{\ProbeBkgLeftBvcQcdLinearOnesixtwo}{121.4}  % experiments/FIGS/data/probe_ladder_v2_mlp2/analysis/seed_level_results.json :: levels.bvc_qcd.linear[1].rejection_points['0.90'].mean_n_bkg_pass :: sha256 53e1ec81343ac2c1 :: background jets passing the cut, mean over seeds
\newcommand{\ProbeRejBvcQcdLinearOnesixtwo}{44.3\,\ensuremath{\pm}\,2.1}  % experiments/FIGS/data/probe_ladder_v2_mlp2/analysis/seed_level_results.json :: table[bvc_qcd,linear,162].rejection_points['0.90'].rejection :: sha256 53e1ec81343ac2c1 :: at 90% signal efficiency, mean +- SD over seeds; 0 of 5 seeds at the cap
\newcommand{\ProbeAucBvcQcdLinearFourthree}{0.98281}  % experiments/FIGS/data/probe_ladder_v2_mlp2/analysis/seed_level_results.json :: levels.bvc_qcd.linear[2].mean_auc :: sha256 53e1ec81343ac2c1 :: 5 seeds, 0 saturated at AUC=1
\newcommand{\ProbeAucSdBvcQcdLinearFourthree}{0.00064}  % experiments/FIGS/data/probe_ladder_v2_mlp2/analysis/seed_level_results.json :: table[bvc_qcd,linear,43].auc (SD over seeds, ddof=1) :: sha256 53e1ec81343ac2c1
\newcommand{\ProbeLogOneMinusAucBvcQcdLinearFourthree}{\ensuremath{-}4.0638}  % experiments/FIGS/data/probe_ladder_v2_mlp2/analysis/seed_level_results.json :: levels.bvc_qcd.linear[2].mean :: sha256 53e1ec81343ac2c1 :: natural log of 1-AUC, lower is better
\newcommand{\ProbeLogOneMinusAucSdBvcQcdLinearFourthree}{0.0370}  % experiments/FIGS/data/probe_ladder_v2_mlp2/analysis/seed_level_results.json :: levels.bvc_qcd.linear[2].seed_sd :: sha256 53e1ec81343ac2c1
\newcommand{\ProbeOmaBvcQcdLinearFourthree}{\ensuremath{(1.72 \pm 0.06)\times10^{-2}}}  % experiments/FIGS/data/probe_ladder_v2_mlp2/analysis/seed_level_results.json :: table[bvc_qcd,linear,43].auc :: sha256 53e1ec81343ac2c1 :: 1-AUC, mean +- SD over 5 seeds
\newcommand{\ProbeBkgLeftBvcQcdLinearFourthree}{126.8}  % experiments/FIGS/data/probe_ladder_v2_mlp2/analysis/seed_level_results.json :: levels.bvc_qcd.linear[2].rejection_points['0.90'].mean_n_bkg_pass :: sha256 53e1ec81343ac2c1 :: background jets passing the cut, mean over seeds
\newcommand{\ProbeRejBvcQcdLinearFourthree}{43\,\ensuremath{\pm}\,4}  % experiments/FIGS/data/probe_ladder_v2_mlp2/analysis/seed_level_results.json :: table[bvc_qcd,linear,43].rejection_points['0.90'].rejection :: sha256 53e1ec81343ac2c1 :: at 90% signal efficiency, mean +- SD over seeds; 0 of 5 seeds at the cap
\newcommand{\ProbeAucBvcQcdLinearOneseven}{0.95504}  % experiments/FIGS/data/probe_ladder_v2_mlp2/analysis/seed_level_results.json :: levels.bvc_qcd.linear[3].mean_auc :: sha256 53e1ec81343ac2c1 :: 5 seeds, 0 saturated at AUC=1
\newcommand{\ProbeAucSdBvcQcdLinearOneseven}{0.00527}  % experiments/FIGS/data/probe_ladder_v2_mlp2/analysis/seed_level_results.json :: table[bvc_qcd,linear,17].auc (SD over seeds, ddof=1) :: sha256 53e1ec81343ac2c1
\newcommand{\ProbeLogOneMinusAucBvcQcdLinearOneseven}{\ensuremath{-}3.1076}  % experiments/FIGS/data/probe_ladder_v2_mlp2/analysis/seed_level_results.json :: levels.bvc_qcd.linear[3].mean :: sha256 53e1ec81343ac2c1 :: natural log of 1-AUC, lower is better
\newcommand{\ProbeLogOneMinusAucSdBvcQcdLinearOneseven}{0.1198}  % experiments/FIGS/data/probe_ladder_v2_mlp2/analysis/seed_level_results.json :: levels.bvc_qcd.linear[3].seed_sd :: sha256 53e1ec81343ac2c1
\newcommand{\ProbeOmaBvcQcdLinearOneseven}{\ensuremath{(4.5 \pm 0.5)\times10^{-2}}}  % experiments/FIGS/data/probe_ladder_v2_mlp2/analysis/seed_level_results.json :: table[bvc_qcd,linear,17].auc :: sha256 53e1ec81343ac2c1 :: 1-AUC, mean +- SD over 5 seeds
\newcommand{\ProbeBkgLeftBvcQcdLinearOneseven}{544.6}  % experiments/FIGS/data/probe_ladder_v2_mlp2/analysis/seed_level_results.json :: levels.bvc_qcd.linear[3].rejection_points['0.90'].mean_n_bkg_pass :: sha256 53e1ec81343ac2c1 :: background jets passing the cut, mean over seeds
\newcommand{\ProbeRejBvcQcdLinearOneseven}{10.1\,\ensuremath{\pm}\,1.9}  % experiments/FIGS/data/probe_ladder_v2_mlp2/analysis/seed_level_results.json :: table[bvc_qcd,linear,17].rejection_points['0.90'].rejection :: sha256 53e1ec81343ac2c1 :: at 90% signal efficiency, mean +- SD over seeds; 0 of 5 seeds at the cap
\newcommand{\ProbeOmaRatioBvcQcdLinearOnesixtwoOverOneeighteight}{1.14}  % experiments/FIGS/data/probe_ladder_v2_mlp2/analysis/seed_level_results.json :: table[bvc_qcd,linear,162/188].auc :: sha256 53e1ec81343ac2c1 :: ratio of the seed means of 1-AUC, coarser over finer
\newcommand{\ProbeOmaRatioBvcQcdLinearFourthreeOverOneeighteight}{1.16}  % experiments/FIGS/data/probe_ladder_v2_mlp2/analysis/seed_level_results.json :: table[bvc_qcd,linear,43/188].auc :: sha256 53e1ec81343ac2c1 :: ratio of the seed means of 1-AUC, coarser over finer
\newcommand{\ProbeOmaRatioBvcQcdLinearOnesevenOverOneeighteight}{3.0}  % experiments/FIGS/data/probe_ladder_v2_mlp2/analysis/seed_level_results.json :: table[bvc_qcd,linear,17/188].auc :: sha256 53e1ec81343ac2c1 :: ratio of the seed means of 1-AUC, coarser over finer
\newcommand{\ProbeOmaRatioBvcQcdLinearFourthreeOverOnesixtwo}{1.01}  % experiments/FIGS/data/probe_ladder_v2_mlp2/analysis/seed_level_results.json :: table[bvc_qcd,linear,43/162].auc :: sha256 53e1ec81343ac2c1 :: ratio of the seed means of 1-AUC, coarser over finer
\newcommand{\ProbeOmaRatioBvcQcdLinearOnesevenOverOnesixtwo}{2.7}  % experiments/FIGS/data/probe_ladder_v2_mlp2/analysis/seed_level_results.json :: table[bvc_qcd,linear,17/162].auc :: sha256 53e1ec81343ac2c1 :: ratio of the seed means of 1-AUC, coarser over finer
\newcommand{\ProbeOmaRatioBvcQcdLinearOnesevenOverFourthree}{2.6}  % experiments/FIGS/data/probe_ladder_v2_mlp2/analysis/seed_level_results.json :: table[bvc_qcd,linear,17/43].auc :: sha256 53e1ec81343ac2c1 :: ratio of the seed means of 1-AUC, coarser over finer
\newcommand{\ProbeAucBvcQcdMlpOneeighteight}{0.98527}  % experiments/FIGS/data/probe_ladder_v2_mlp2/analysis/seed_level_results.json :: levels.bvc_qcd.mlp[0].mean_auc :: sha256 53e1ec81343ac2c1 :: 5 seeds, 0 saturated at AUC=1
\newcommand{\ProbeAucSdBvcQcdMlpOneeighteight}{0.00058}  % experiments/FIGS/data/probe_ladder_v2_mlp2/analysis/seed_level_results.json :: table[bvc_qcd,mlp,188].auc (SD over seeds, ddof=1) :: sha256 53e1ec81343ac2c1
\newcommand{\ProbeLogOneMinusAucBvcQcdMlpOneeighteight}{\ensuremath{-}4.2188}  % experiments/FIGS/data/probe_ladder_v2_mlp2/analysis/seed_level_results.json :: levels.bvc_qcd.mlp[0].mean :: sha256 53e1ec81343ac2c1 :: natural log of 1-AUC, lower is better
\newcommand{\ProbeLogOneMinusAucSdBvcQcdMlpOneeighteight}{0.0400}  % experiments/FIGS/data/probe_ladder_v2_mlp2/analysis/seed_level_results.json :: levels.bvc_qcd.mlp[0].seed_sd :: sha256 53e1ec81343ac2c1
\newcommand{\ProbeOmaBvcQcdMlpOneeighteight}{\ensuremath{(1.47 \pm 0.06)\times10^{-2}}}  % experiments/FIGS/data/probe_ladder_v2_mlp2/analysis/seed_level_results.json :: table[bvc_qcd,mlp,188].auc :: sha256 53e1ec81343ac2c1 :: 1-AUC, mean +- SD over 5 seeds
\newcommand{\ProbeBkgLeftBvcQcdMlpOneeighteight}{94.6}  % experiments/FIGS/data/probe_ladder_v2_mlp2/analysis/seed_level_results.json :: levels.bvc_qcd.mlp[0].rejection_points['0.90'].mean_n_bkg_pass :: sha256 53e1ec81343ac2c1 :: background jets passing the cut, mean over seeds
\newcommand{\ProbeRejBvcQcdMlpOneeighteight}{57\,\ensuremath{\pm}\,6}  % experiments/FIGS/data/probe_ladder_v2_mlp2/analysis/seed_level_results.json :: table[bvc_qcd,mlp,188].rejection_points['0.90'].rejection :: sha256 53e1ec81343ac2c1 :: at 90% signal efficiency, mean +- SD over seeds; 0 of 5 seeds at the cap
\newcommand{\ProbeAucBvcQcdMlpOnesixtwo}{0.98361}  % experiments/FIGS/data/probe_ladder_v2_mlp2/analysis/seed_level_results.json :: levels.bvc_qcd.mlp[1].mean_auc :: sha256 53e1ec81343ac2c1 :: 5 seeds, 0 saturated at AUC=1
\newcommand{\ProbeAucSdBvcQcdMlpOnesixtwo}{0.00039}  % experiments/FIGS/data/probe_ladder_v2_mlp2/analysis/seed_level_results.json :: table[bvc_qcd,mlp,162].auc (SD over seeds, ddof=1) :: sha256 53e1ec81343ac2c1
\newcommand{\ProbeLogOneMinusAucBvcQcdMlpOnesixtwo}{\ensuremath{-}4.1113}  % experiments/FIGS/data/probe_ladder_v2_mlp2/analysis/seed_level_results.json :: levels.bvc_qcd.mlp[1].mean :: sha256 53e1ec81343ac2c1 :: natural log of 1-AUC, lower is better
\newcommand{\ProbeLogOneMinusAucSdBvcQcdMlpOnesixtwo}{0.0241}  % experiments/FIGS/data/probe_ladder_v2_mlp2/analysis/seed_level_results.json :: levels.bvc_qcd.mlp[1].seed_sd :: sha256 53e1ec81343ac2c1
\newcommand{\ProbeOmaBvcQcdMlpOnesixtwo}{\ensuremath{(1.64 \pm 0.04)\times10^{-2}}}  % experiments/FIGS/data/probe_ladder_v2_mlp2/analysis/seed_level_results.json :: table[bvc_qcd,mlp,162].auc :: sha256 53e1ec81343ac2c1 :: 1-AUC, mean +- SD over 5 seeds
\newcommand{\ProbeBkgLeftBvcQcdMlpOnesixtwo}{119.6}  % experiments/FIGS/data/probe_ladder_v2_mlp2/analysis/seed_level_results.json :: levels.bvc_qcd.mlp[1].rejection_points['0.90'].mean_n_bkg_pass :: sha256 53e1ec81343ac2c1 :: background jets passing the cut, mean over seeds
\newcommand{\ProbeRejBvcQcdMlpOnesixtwo}{45.1\,\ensuremath{\pm}\,3.2}  % experiments/FIGS/data/probe_ladder_v2_mlp2/analysis/seed_level_results.json :: table[bvc_qcd,mlp,162].rejection_points['0.90'].rejection :: sha256 53e1ec81343ac2c1 :: at 90% signal efficiency, mean +- SD over seeds; 0 of 5 seeds at the cap
\newcommand{\ProbeAucBvcQcdMlpFourthree}{0.98341}  % experiments/FIGS/data/probe_ladder_v2_mlp2/analysis/seed_level_results.json :: levels.bvc_qcd.mlp[2].mean_auc :: sha256 53e1ec81343ac2c1 :: 5 seeds, 0 saturated at AUC=1
\newcommand{\ProbeAucSdBvcQcdMlpFourthree}{0.00052}  % experiments/FIGS/data/probe_ladder_v2_mlp2/analysis/seed_level_results.json :: table[bvc_qcd,mlp,43].auc (SD over seeds, ddof=1) :: sha256 53e1ec81343ac2c1
\newcommand{\ProbeLogOneMinusAucBvcQcdMlpFourthree}{\ensuremath{-}4.0995}  % experiments/FIGS/data/probe_ladder_v2_mlp2/analysis/seed_level_results.json :: levels.bvc_qcd.mlp[2].mean :: sha256 53e1ec81343ac2c1 :: natural log of 1-AUC, lower is better
\newcommand{\ProbeLogOneMinusAucSdBvcQcdMlpFourthree}{0.0318}  % experiments/FIGS/data/probe_ladder_v2_mlp2/analysis/seed_level_results.json :: levels.bvc_qcd.mlp[2].seed_sd :: sha256 53e1ec81343ac2c1
\newcommand{\ProbeOmaBvcQcdMlpFourthree}{\ensuremath{(1.66 \pm 0.05)\times10^{-2}}}  % experiments/FIGS/data/probe_ladder_v2_mlp2/analysis/seed_level_results.json :: table[bvc_qcd,mlp,43].auc :: sha256 53e1ec81343ac2c1 :: 1-AUC, mean +- SD over 5 seeds
\newcommand{\ProbeBkgLeftBvcQcdMlpFourthree}{119.2}  % experiments/FIGS/data/probe_ladder_v2_mlp2/analysis/seed_level_results.json :: levels.bvc_qcd.mlp[2].rejection_points['0.90'].mean_n_bkg_pass :: sha256 53e1ec81343ac2c1 :: background jets passing the cut, mean over seeds
\newcommand{\ProbeRejBvcQcdMlpFourthree}{45.1\,\ensuremath{\pm}\,1.3}  % experiments/FIGS/data/probe_ladder_v2_mlp2/analysis/seed_level_results.json :: table[bvc_qcd,mlp,43].rejection_points['0.90'].rejection :: sha256 53e1ec81343ac2c1 :: at 90% signal efficiency, mean +- SD over seeds; 0 of 5 seeds at the cap
\newcommand{\ProbeAucBvcQcdMlpOneseven}{0.96226}  % experiments/FIGS/data/probe_ladder_v2_mlp2/analysis/seed_level_results.json :: levels.bvc_qcd.mlp[3].mean_auc :: sha256 53e1ec81343ac2c1 :: 5 seeds, 0 saturated at AUC=1
\newcommand{\ProbeAucSdBvcQcdMlpOneseven}{0.00373}  % experiments/FIGS/data/probe_ladder_v2_mlp2/analysis/seed_level_results.json :: table[bvc_qcd,mlp,17].auc (SD over seeds, ddof=1) :: sha256 53e1ec81343ac2c1
\newcommand{\ProbeLogOneMinusAucBvcQcdMlpOneseven}{\ensuremath{-}3.2811}  % experiments/FIGS/data/probe_ladder_v2_mlp2/analysis/seed_level_results.json :: levels.bvc_qcd.mlp[3].mean :: sha256 53e1ec81343ac2c1 :: natural log of 1-AUC, lower is better
\newcommand{\ProbeLogOneMinusAucSdBvcQcdMlpOneseven}{0.1010}  % experiments/FIGS/data/probe_ladder_v2_mlp2/analysis/seed_level_results.json :: levels.bvc_qcd.mlp[3].seed_sd :: sha256 53e1ec81343ac2c1
\newcommand{\ProbeOmaBvcQcdMlpOneseven}{\ensuremath{(3.8 \pm 0.4)\times10^{-2}}}  % experiments/FIGS/data/probe_ladder_v2_mlp2/analysis/seed_level_results.json :: table[bvc_qcd,mlp,17].auc :: sha256 53e1ec81343ac2c1 :: 1-AUC, mean +- SD over 5 seeds
\newcommand{\ProbeBkgLeftBvcQcdMlpOneseven}{458.8}  % experiments/FIGS/data/probe_ladder_v2_mlp2/analysis/seed_level_results.json :: levels.bvc_qcd.mlp[3].rejection_points['0.90'].mean_n_bkg_pass :: sha256 53e1ec81343ac2c1 :: background jets passing the cut, mean over seeds
\newcommand{\ProbeRejBvcQcdMlpOneseven}{11.8\,\ensuremath{\pm}\,1.5}  % experiments/FIGS/data/probe_ladder_v2_mlp2/analysis/seed_level_results.json :: table[bvc_qcd,mlp,17].rejection_points['0.90'].rejection :: sha256 53e1ec81343ac2c1 :: at 90% signal efficiency, mean +- SD over seeds; 0 of 5 seeds at the cap
\newcommand{\ProbeOmaRatioBvcQcdMlpOnesixtwoOverOneeighteight}{1.11}  % experiments/FIGS/data/probe_ladder_v2_mlp2/analysis/seed_level_results.json :: table[bvc_qcd,mlp,162/188].auc :: sha256 53e1ec81343ac2c1 :: ratio of the seed means of 1-AUC, coarser over finer
\newcommand{\ProbeOmaRatioBvcQcdMlpFourthreeOverOneeighteight}{1.13}  % experiments/FIGS/data/probe_ladder_v2_mlp2/analysis/seed_level_results.json :: table[bvc_qcd,mlp,43/188].auc :: sha256 53e1ec81343ac2c1 :: ratio of the seed means of 1-AUC, coarser over finer
\newcommand{\ProbeOmaRatioBvcQcdMlpOnesevenOverOneeighteight}{2.6}  % experiments/FIGS/data/probe_ladder_v2_mlp2/analysis/seed_level_results.json :: table[bvc_qcd,mlp,17/188].auc :: sha256 53e1ec81343ac2c1 :: ratio of the seed means of 1-AUC, coarser over finer
\newcommand{\ProbeOmaRatioBvcQcdMlpFourthreeOverOnesixtwo}{1.01}  % experiments/FIGS/data/probe_ladder_v2_mlp2/analysis/seed_level_results.json :: table[bvc_qcd,mlp,43/162].auc :: sha256 53e1ec81343ac2c1 :: ratio of the seed means of 1-AUC, coarser over finer
\newcommand{\ProbeOmaRatioBvcQcdMlpOnesevenOverOnesixtwo}{2.3}  % experiments/FIGS/data/probe_ladder_v2_mlp2/analysis/seed_level_results.json :: table[bvc_qcd,mlp,17/162].auc :: sha256 53e1ec81343ac2c1 :: ratio of the seed means of 1-AUC, coarser over finer
\newcommand{\ProbeOmaRatioBvcQcdMlpOnesevenOverFourthree}{2.3}  % experiments/FIGS/data/probe_ladder_v2_mlp2/analysis/seed_level_results.json :: table[bvc_qcd,mlp,17/43].auc :: sha256 53e1ec81343ac2c1 :: ratio of the seed means of 1-AUC, coarser over finer
\newcommand{\ProbeAucBvcResonantLinearOneeighteight}{0.99768}  % experiments/FIGS/data/probe_ladder_v2_mlp2/analysis/seed_level_results.json :: levels.bvc_resonant.linear[0].mean_auc :: sha256 53e1ec81343ac2c1 :: 5 seeds, 0 saturated at AUC=1
\newcommand{\ProbeAucSdBvcResonantLinearOneeighteight}{0.00007}  % experiments/FIGS/data/probe_ladder_v2_mlp2/analysis/seed_level_results.json :: table[bvc_resonant,linear,188].auc (SD over seeds, ddof=1) :: sha256 53e1ec81343ac2c1
\newcommand{\ProbeLogOneMinusAucBvcResonantLinearOneeighteight}{\ensuremath{-}6.0683}  % experiments/FIGS/data/probe_ladder_v2_mlp2/analysis/seed_level_results.json :: levels.bvc_resonant.linear[0].mean :: sha256 53e1ec81343ac2c1 :: natural log of 1-AUC, lower is better
\newcommand{\ProbeLogOneMinusAucSdBvcResonantLinearOneeighteight}{0.0317}  % experiments/FIGS/data/probe_ladder_v2_mlp2/analysis/seed_level_results.json :: levels.bvc_resonant.linear[0].seed_sd :: sha256 53e1ec81343ac2c1
\newcommand{\ProbeOmaBvcResonantLinearOneeighteight}{\ensuremath{(2.32 \pm 0.07)\times10^{-3}}}  % experiments/FIGS/data/probe_ladder_v2_mlp2/analysis/seed_level_results.json :: table[bvc_resonant,linear,188].auc :: sha256 53e1ec81343ac2c1 :: 1-AUC, mean +- SD over 5 seeds
\newcommand{\ProbeBkgLeftBvcResonantLinearOneeighteight}{9.0}  % experiments/FIGS/data/probe_ladder_v2_mlp2/analysis/seed_level_results.json :: levels.bvc_resonant.linear[0].rejection_points['0.90'].mean_n_bkg_pass :: sha256 53e1ec81343ac2c1 :: background jets passing the cut, mean over seeds
\newcommand{\ProbeRejBvcResonantLinearOneeighteight}{1{,}380\,\ensuremath{\pm}\,350}  % experiments/FIGS/data/probe_ladder_v2_mlp2/analysis/seed_level_results.json :: table[bvc_resonant,linear,188].rejection_points['0.90'].rejection :: sha256 53e1ec81343ac2c1 :: at 90% signal efficiency, mean +- SD over seeds; 0 of 5 seeds at the cap
\newcommand{\ProbeAucBvcResonantLinearOnesixtwo}{0.99764}  % experiments/FIGS/data/probe_ladder_v2_mlp2/analysis/seed_level_results.json :: levels.bvc_resonant.linear[1].mean_auc :: sha256 53e1ec81343ac2c1 :: 5 seeds, 0 saturated at AUC=1
\newcommand{\ProbeAucSdBvcResonantLinearOnesixtwo}{0.00006}  % experiments/FIGS/data/probe_ladder_v2_mlp2/analysis/seed_level_results.json :: table[bvc_resonant,linear,162].auc (SD over seeds, ddof=1) :: sha256 53e1ec81343ac2c1
\newcommand{\ProbeLogOneMinusAucBvcResonantLinearOnesixtwo}{\ensuremath{-}6.0493}  % experiments/FIGS/data/probe_ladder_v2_mlp2/analysis/seed_level_results.json :: levels.bvc_resonant.linear[1].mean :: sha256 53e1ec81343ac2c1 :: natural log of 1-AUC, lower is better
\newcommand{\ProbeLogOneMinusAucSdBvcResonantLinearOnesixtwo}{0.0261}  % experiments/FIGS/data/probe_ladder_v2_mlp2/analysis/seed_level_results.json :: levels.bvc_resonant.linear[1].seed_sd :: sha256 53e1ec81343ac2c1
\newcommand{\ProbeOmaBvcResonantLinearOnesixtwo}{\ensuremath{(2.36 \pm 0.06)\times10^{-3}}}  % experiments/FIGS/data/probe_ladder_v2_mlp2/analysis/seed_level_results.json :: table[bvc_resonant,linear,162].auc :: sha256 53e1ec81343ac2c1 :: 1-AUC, mean +- SD over 5 seeds
\newcommand{\ProbeBkgLeftBvcResonantLinearOnesixtwo}{9.4}  % experiments/FIGS/data/probe_ladder_v2_mlp2/analysis/seed_level_results.json :: levels.bvc_resonant.linear[1].rejection_points['0.90'].mean_n_bkg_pass :: sha256 53e1ec81343ac2c1 :: background jets passing the cut, mean over seeds
\newcommand{\ProbeRejBvcResonantLinearOnesixtwo}{1{,}290\,\ensuremath{\pm}\,240}  % experiments/FIGS/data/probe_ladder_v2_mlp2/analysis/seed_level_results.json :: table[bvc_resonant,linear,162].rejection_points['0.90'].rejection :: sha256 53e1ec81343ac2c1 :: at 90% signal efficiency, mean +- SD over seeds; 0 of 5 seeds at the cap
\newcommand{\ProbeAucBvcResonantLinearFourthree}{0.99748}  % experiments/FIGS/data/probe_ladder_v2_mlp2/analysis/seed_level_results.json :: levels.bvc_resonant.linear[2].mean_auc :: sha256 53e1ec81343ac2c1 :: 5 seeds, 0 saturated at AUC=1
\newcommand{\ProbeAucSdBvcResonantLinearFourthree}{0.00007}  % experiments/FIGS/data/probe_ladder_v2_mlp2/analysis/seed_level_results.json :: table[bvc_resonant,linear,43].auc (SD over seeds, ddof=1) :: sha256 53e1ec81343ac2c1
\newcommand{\ProbeLogOneMinusAucBvcResonantLinearFourthree}{\ensuremath{-}5.9834}  % experiments/FIGS/data/probe_ladder_v2_mlp2/analysis/seed_level_results.json :: levels.bvc_resonant.linear[2].mean :: sha256 53e1ec81343ac2c1 :: natural log of 1-AUC, lower is better
\newcommand{\ProbeLogOneMinusAucSdBvcResonantLinearFourthree}{0.0264}  % experiments/FIGS/data/probe_ladder_v2_mlp2/analysis/seed_level_results.json :: levels.bvc_resonant.linear[2].seed_sd :: sha256 53e1ec81343ac2c1
\newcommand{\ProbeOmaBvcResonantLinearFourthree}{\ensuremath{(2.52 \pm 0.07)\times10^{-3}}}  % experiments/FIGS/data/probe_ladder_v2_mlp2/analysis/seed_level_results.json :: table[bvc_resonant,linear,43].auc :: sha256 53e1ec81343ac2c1 :: 1-AUC, mean +- SD over 5 seeds
\newcommand{\ProbeBkgLeftBvcResonantLinearFourthree}{11.0}  % experiments/FIGS/data/probe_ladder_v2_mlp2/analysis/seed_level_results.json :: levels.bvc_resonant.linear[2].rejection_points['0.90'].mean_n_bkg_pass :: sha256 53e1ec81343ac2c1 :: background jets passing the cut, mean over seeds
\newcommand{\ProbeRejBvcResonantLinearFourthree}{1{,}130\,\ensuremath{\pm}\,260}  % experiments/FIGS/data/probe_ladder_v2_mlp2/analysis/seed_level_results.json :: table[bvc_resonant,linear,43].rejection_points['0.90'].rejection :: sha256 53e1ec81343ac2c1 :: at 90% signal efficiency, mean +- SD over seeds; 0 of 5 seeds at the cap
\newcommand{\ProbeAucBvcResonantLinearOneseven}{0.99067}  % experiments/FIGS/data/probe_ladder_v2_mlp2/analysis/seed_level_results.json :: levels.bvc_resonant.linear[3].mean_auc :: sha256 53e1ec81343ac2c1 :: 5 seeds, 0 saturated at AUC=1
\newcommand{\ProbeAucSdBvcResonantLinearOneseven}{0.00147}  % experiments/FIGS/data/probe_ladder_v2_mlp2/analysis/seed_level_results.json :: table[bvc_resonant,linear,17].auc (SD over seeds, ddof=1) :: sha256 53e1ec81343ac2c1
\newcommand{\ProbeLogOneMinusAucBvcResonantLinearOneseven}{\ensuremath{-}4.6848}  % experiments/FIGS/data/probe_ladder_v2_mlp2/analysis/seed_level_results.json :: levels.bvc_resonant.linear[3].mean :: sha256 53e1ec81343ac2c1 :: natural log of 1-AUC, lower is better
\newcommand{\ProbeLogOneMinusAucSdBvcResonantLinearOneseven}{0.1619}  % experiments/FIGS/data/probe_ladder_v2_mlp2/analysis/seed_level_results.json :: levels.bvc_resonant.linear[3].seed_sd :: sha256 53e1ec81343ac2c1
\newcommand{\ProbeOmaBvcResonantLinearOneseven}{\ensuremath{(9.3 \pm 1.5)\times10^{-3}}}  % experiments/FIGS/data/probe_ladder_v2_mlp2/analysis/seed_level_results.json :: table[bvc_resonant,linear,17].auc :: sha256 53e1ec81343ac2c1 :: 1-AUC, mean +- SD over 5 seeds
\newcommand{\ProbeBkgLeftBvcResonantLinearOneseven}{90.8}  % experiments/FIGS/data/probe_ladder_v2_mlp2/analysis/seed_level_results.json :: levels.bvc_resonant.linear[3].rejection_points['0.90'].mean_n_bkg_pass :: sha256 53e1ec81343ac2c1 :: background jets passing the cut, mean over seeds
\newcommand{\ProbeRejBvcResonantLinearOneseven}{134\,\ensuremath{\pm}\,21}  % experiments/FIGS/data/probe_ladder_v2_mlp2/analysis/seed_level_results.json :: table[bvc_resonant,linear,17].rejection_points['0.90'].rejection :: sha256 53e1ec81343ac2c1 :: at 90% signal efficiency, mean +- SD over seeds; 0 of 5 seeds at the cap
\newcommand{\ProbeOmaRatioBvcResonantLinearOnesixtwoOverOneeighteight}{1.02}  % experiments/FIGS/data/probe_ladder_v2_mlp2/analysis/seed_level_results.json :: table[bvc_resonant,linear,162/188].auc :: sha256 53e1ec81343ac2c1 :: ratio of the seed means of 1-AUC, coarser over finer
\newcommand{\ProbeOmaRatioBvcResonantLinearFourthreeOverOneeighteight}{1.09}  % experiments/FIGS/data/probe_ladder_v2_mlp2/analysis/seed_level_results.json :: table[bvc_resonant,linear,43/188].auc :: sha256 53e1ec81343ac2c1 :: ratio of the seed means of 1-AUC, coarser over finer
\newcommand{\ProbeOmaRatioBvcResonantLinearOnesevenOverOneeighteight}{4.0}  % experiments/FIGS/data/probe_ladder_v2_mlp2/analysis/seed_level_results.json :: table[bvc_resonant,linear,17/188].auc :: sha256 53e1ec81343ac2c1 :: ratio of the seed means of 1-AUC, coarser over finer
\newcommand{\ProbeOmaRatioBvcResonantLinearFourthreeOverOnesixtwo}{1.07}  % experiments/FIGS/data/probe_ladder_v2_mlp2/analysis/seed_level_results.json :: table[bvc_resonant,linear,43/162].auc :: sha256 53e1ec81343ac2c1 :: ratio of the seed means of 1-AUC, coarser over finer
\newcommand{\ProbeOmaRatioBvcResonantLinearOnesevenOverOnesixtwo}{4.0}  % experiments/FIGS/data/probe_ladder_v2_mlp2/analysis/seed_level_results.json :: table[bvc_resonant,linear,17/162].auc :: sha256 53e1ec81343ac2c1 :: ratio of the seed means of 1-AUC, coarser over finer
\newcommand{\ProbeOmaRatioBvcResonantLinearOnesevenOverFourthree}{3.7}  % experiments/FIGS/data/probe_ladder_v2_mlp2/analysis/seed_level_results.json :: table[bvc_resonant,linear,17/43].auc :: sha256 53e1ec81343ac2c1 :: ratio of the seed means of 1-AUC, coarser over finer
\newcommand{\ProbeAucBvcResonantMlpOneeighteight}{0.99778}  % experiments/FIGS/data/probe_ladder_v2_mlp2/analysis/seed_level_results.json :: levels.bvc_resonant.mlp[0].mean_auc :: sha256 53e1ec81343ac2c1 :: 5 seeds, 0 saturated at AUC=1
\newcommand{\ProbeAucSdBvcResonantMlpOneeighteight}{0.00009}  % experiments/FIGS/data/probe_ladder_v2_mlp2/analysis/seed_level_results.json :: table[bvc_resonant,mlp,188].auc (SD over seeds, ddof=1) :: sha256 53e1ec81343ac2c1
\newcommand{\ProbeLogOneMinusAucBvcResonantMlpOneeighteight}{\ensuremath{-}6.1102}  % experiments/FIGS/data/probe_ladder_v2_mlp2/analysis/seed_level_results.json :: levels.bvc_resonant.mlp[0].mean :: sha256 53e1ec81343ac2c1 :: natural log of 1-AUC, lower is better
\newcommand{\ProbeLogOneMinusAucSdBvcResonantMlpOneeighteight}{0.0383}  % experiments/FIGS/data/probe_ladder_v2_mlp2/analysis/seed_level_results.json :: levels.bvc_resonant.mlp[0].seed_sd :: sha256 53e1ec81343ac2c1
\newcommand{\ProbeOmaBvcResonantMlpOneeighteight}{\ensuremath{(2.22 \pm 0.09)\times10^{-3}}}  % experiments/FIGS/data/probe_ladder_v2_mlp2/analysis/seed_level_results.json :: table[bvc_resonant,mlp,188].auc :: sha256 53e1ec81343ac2c1 :: 1-AUC, mean +- SD over 5 seeds
\newcommand{\ProbeBkgLeftBvcResonantMlpOneeighteight}{9.8}  % experiments/FIGS/data/probe_ladder_v2_mlp2/analysis/seed_level_results.json :: levels.bvc_resonant.mlp[0].rejection_points['0.90'].mean_n_bkg_pass :: sha256 53e1ec81343ac2c1 :: background jets passing the cut, mean over seeds
\newcommand{\ProbeRejBvcResonantMlpOneeighteight}{1{,}230\,\ensuremath{\pm}\,180}  % experiments/FIGS/data/probe_ladder_v2_mlp2/analysis/seed_level_results.json :: table[bvc_resonant,mlp,188].rejection_points['0.90'].rejection :: sha256 53e1ec81343ac2c1 :: at 90% signal efficiency, mean +- SD over seeds; 0 of 5 seeds at the cap
\newcommand{\ProbeAucBvcResonantMlpOnesixtwo}{0.99768}  % experiments/FIGS/data/probe_ladder_v2_mlp2/analysis/seed_level_results.json :: levels.bvc_resonant.mlp[1].mean_auc :: sha256 53e1ec81343ac2c1 :: 5 seeds, 0 saturated at AUC=1
\newcommand{\ProbeAucSdBvcResonantMlpOnesixtwo}{0.00005}  % experiments/FIGS/data/probe_ladder_v2_mlp2/analysis/seed_level_results.json :: table[bvc_resonant,mlp,162].auc (SD over seeds, ddof=1) :: sha256 53e1ec81343ac2c1
\newcommand{\ProbeLogOneMinusAucBvcResonantMlpOnesixtwo}{\ensuremath{-}6.0667}  % experiments/FIGS/data/probe_ladder_v2_mlp2/analysis/seed_level_results.json :: levels.bvc_resonant.mlp[1].mean :: sha256 53e1ec81343ac2c1 :: natural log of 1-AUC, lower is better
\newcommand{\ProbeLogOneMinusAucSdBvcResonantMlpOnesixtwo}{0.0220}  % experiments/FIGS/data/probe_ladder_v2_mlp2/analysis/seed_level_results.json :: levels.bvc_resonant.mlp[1].seed_sd :: sha256 53e1ec81343ac2c1
\newcommand{\ProbeOmaBvcResonantMlpOnesixtwo}{\ensuremath{(2.32 \pm 0.05)\times10^{-3}}}  % experiments/FIGS/data/probe_ladder_v2_mlp2/analysis/seed_level_results.json :: table[bvc_resonant,mlp,162].auc :: sha256 53e1ec81343ac2c1 :: 1-AUC, mean +- SD over 5 seeds
\newcommand{\ProbeBkgLeftBvcResonantMlpOnesixtwo}{10.4}  % experiments/FIGS/data/probe_ladder_v2_mlp2/analysis/seed_level_results.json :: levels.bvc_resonant.mlp[1].rejection_points['0.90'].mean_n_bkg_pass :: sha256 53e1ec81343ac2c1 :: background jets passing the cut, mean over seeds
\newcommand{\ProbeRejBvcResonantMlpOnesixtwo}{1{,}170\,\ensuremath{\pm}\,220}  % experiments/FIGS/data/probe_ladder_v2_mlp2/analysis/seed_level_results.json :: table[bvc_resonant,mlp,162].rejection_points['0.90'].rejection :: sha256 53e1ec81343ac2c1 :: at 90% signal efficiency, mean +- SD over seeds; 0 of 5 seeds at the cap
\newcommand{\ProbeAucBvcResonantMlpFourthree}{0.99747}  % experiments/FIGS/data/probe_ladder_v2_mlp2/analysis/seed_level_results.json :: levels.bvc_resonant.mlp[2].mean_auc :: sha256 53e1ec81343ac2c1 :: 5 seeds, 0 saturated at AUC=1
\newcommand{\ProbeAucSdBvcResonantMlpFourthree}{0.00004}  % experiments/FIGS/data/probe_ladder_v2_mlp2/analysis/seed_level_results.json :: table[bvc_resonant,mlp,43].auc (SD over seeds, ddof=1) :: sha256 53e1ec81343ac2c1
\newcommand{\ProbeLogOneMinusAucBvcResonantMlpFourthree}{\ensuremath{-}5.9792}  % experiments/FIGS/data/probe_ladder_v2_mlp2/analysis/seed_level_results.json :: levels.bvc_resonant.mlp[2].mean :: sha256 53e1ec81343ac2c1 :: natural log of 1-AUC, lower is better
\newcommand{\ProbeLogOneMinusAucSdBvcResonantMlpFourthree}{0.0164}  % experiments/FIGS/data/probe_ladder_v2_mlp2/analysis/seed_level_results.json :: levels.bvc_resonant.mlp[2].seed_sd :: sha256 53e1ec81343ac2c1
\newcommand{\ProbeOmaBvcResonantMlpFourthree}{\ensuremath{(2.53 \pm 0.04)\times10^{-3}}}  % experiments/FIGS/data/probe_ladder_v2_mlp2/analysis/seed_level_results.json :: table[bvc_resonant,mlp,43].auc :: sha256 53e1ec81343ac2c1 :: 1-AUC, mean +- SD over 5 seeds
\newcommand{\ProbeBkgLeftBvcResonantMlpFourthree}{11.8}  % experiments/FIGS/data/probe_ladder_v2_mlp2/analysis/seed_level_results.json :: levels.bvc_resonant.mlp[2].rejection_points['0.90'].mean_n_bkg_pass :: sha256 53e1ec81343ac2c1 :: background jets passing the cut, mean over seeds
\newcommand{\ProbeRejBvcResonantMlpFourthree}{1{,}030\,\ensuremath{\pm}\,170}  % experiments/FIGS/data/probe_ladder_v2_mlp2/analysis/seed_level_results.json :: table[bvc_resonant,mlp,43].rejection_points['0.90'].rejection :: sha256 53e1ec81343ac2c1 :: at 90% signal efficiency, mean +- SD over seeds; 0 of 5 seeds at the cap
\newcommand{\ProbeAucBvcResonantMlpOneseven}{0.99274}  % experiments/FIGS/data/probe_ladder_v2_mlp2/analysis/seed_level_results.json :: levels.bvc_resonant.mlp[3].mean_auc :: sha256 53e1ec81343ac2c1 :: 5 seeds, 0 saturated at AUC=1
\newcommand{\ProbeAucSdBvcResonantMlpOneseven}{0.00089}  % experiments/FIGS/data/probe_ladder_v2_mlp2/analysis/seed_level_results.json :: table[bvc_resonant,mlp,17].auc (SD over seeds, ddof=1) :: sha256 53e1ec81343ac2c1
\newcommand{\ProbeLogOneMinusAucBvcResonantMlpOneseven}{\ensuremath{-}4.9309}  % experiments/FIGS/data/probe_ladder_v2_mlp2/analysis/seed_level_results.json :: levels.bvc_resonant.mlp[3].mean :: sha256 53e1ec81343ac2c1 :: natural log of 1-AUC, lower is better
\newcommand{\ProbeLogOneMinusAucSdBvcResonantMlpOneseven}{0.1224}  % experiments/FIGS/data/probe_ladder_v2_mlp2/analysis/seed_level_results.json :: levels.bvc_resonant.mlp[3].seed_sd :: sha256 53e1ec81343ac2c1
\newcommand{\ProbeOmaBvcResonantMlpOneseven}{\ensuremath{(7.3 \pm 0.9)\times10^{-3}}}  % experiments/FIGS/data/probe_ladder_v2_mlp2/analysis/seed_level_results.json :: table[bvc_resonant,mlp,17].auc :: sha256 53e1ec81343ac2c1 :: 1-AUC, mean +- SD over 5 seeds
\newcommand{\ProbeBkgLeftBvcResonantMlpOneseven}{62.6}  % experiments/FIGS/data/probe_ladder_v2_mlp2/analysis/seed_level_results.json :: levels.bvc_resonant.mlp[3].rejection_points['0.90'].mean_n_bkg_pass :: sha256 53e1ec81343ac2c1 :: background jets passing the cut, mean over seeds
\newcommand{\ProbeRejBvcResonantMlpOneseven}{192\,\ensuremath{\pm}\,21}  % experiments/FIGS/data/probe_ladder_v2_mlp2/analysis/seed_level_results.json :: table[bvc_resonant,mlp,17].rejection_points['0.90'].rejection :: sha256 53e1ec81343ac2c1 :: at 90% signal efficiency, mean +- SD over seeds; 0 of 5 seeds at the cap
\newcommand{\ProbeOmaRatioBvcResonantMlpOnesixtwoOverOneeighteight}{1.04}  % experiments/FIGS/data/probe_ladder_v2_mlp2/analysis/seed_level_results.json :: table[bvc_resonant,mlp,162/188].auc :: sha256 53e1ec81343ac2c1 :: ratio of the seed means of 1-AUC, coarser over finer
\newcommand{\ProbeOmaRatioBvcResonantMlpFourthreeOverOneeighteight}{1.14}  % experiments/FIGS/data/probe_ladder_v2_mlp2/analysis/seed_level_results.json :: table[bvc_resonant,mlp,43/188].auc :: sha256 53e1ec81343ac2c1 :: ratio of the seed means of 1-AUC, coarser over finer
\newcommand{\ProbeOmaRatioBvcResonantMlpOnesevenOverOneeighteight}{3.3}  % experiments/FIGS/data/probe_ladder_v2_mlp2/analysis/seed_level_results.json :: table[bvc_resonant,mlp,17/188].auc :: sha256 53e1ec81343ac2c1 :: ratio of the seed means of 1-AUC, coarser over finer
\newcommand{\ProbeOmaRatioBvcResonantMlpFourthreeOverOnesixtwo}{1.09}  % experiments/FIGS/data/probe_ladder_v2_mlp2/analysis/seed_level_results.json :: table[bvc_resonant,mlp,43/162].auc :: sha256 53e1ec81343ac2c1 :: ratio of the seed means of 1-AUC, coarser over finer
\newcommand{\ProbeOmaRatioBvcResonantMlpOnesevenOverOnesixtwo}{3.1}  % experiments/FIGS/data/probe_ladder_v2_mlp2/analysis/seed_level_results.json :: table[bvc_resonant,mlp,17/162].auc :: sha256 53e1ec81343ac2c1 :: ratio of the seed means of 1-AUC, coarser over finer
\newcommand{\ProbeOmaRatioBvcResonantMlpOnesevenOverFourthree}{2.9}  % experiments/FIGS/data/probe_ladder_v2_mlp2/analysis/seed_level_results.json :: table[bvc_resonant,mlp,17/43].auc :: sha256 53e1ec81343ac2c1 :: ratio of the seed means of 1-AUC, coarser over finer
\newcommand{\ProbeAucEeVsMmLinearOneeighteight}{\ensuremath{1^{\dagger}}}  % experiments/FIGS/data/probe_ladder_v2_mlp2/analysis/seed_level_results.json :: levels.ee_vs_mm.linear[0].mean_auc :: sha256 53e1ec81343ac2c1 :: 5 seeds, 5 saturated at AUC=1
\newcommand{\ProbeAucSdEeVsMmLinearOneeighteight}{0.00000}  % experiments/FIGS/data/probe_ladder_v2_mlp2/analysis/seed_level_results.json :: table[ee_vs_mm,linear,188].auc (SD over seeds, ddof=1) :: sha256 53e1ec81343ac2c1
\newcommand{\ProbeLogOneMinusAucEeVsMmLinearOneeighteight}{\ensuremath{-}16.9579}  % experiments/FIGS/data/probe_ladder_v2_mlp2/analysis/seed_level_results.json :: levels.ee_vs_mm.linear[0].mean :: sha256 53e1ec81343ac2c1 :: natural log of 1-AUC, lower is better
\newcommand{\ProbeLogOneMinusAucSdEeVsMmLinearOneeighteight}{0.0000}  % experiments/FIGS/data/probe_ladder_v2_mlp2/analysis/seed_level_results.json :: levels.ee_vs_mm.linear[0].seed_sd :: sha256 53e1ec81343ac2c1
\newcommand{\ProbeOmaEeVsMmLinearOneeighteight}{---}  % experiments/FIGS/data/probe_ladder_v2_mlp2/analysis/seed_level_results.json :: table[ee_vs_mm,linear,188].auc :: sha256 53e1ec81343ac2c1 :: not quoted: the AUC reached 1 in a seed, so 1-AUC is only a bound
\newcommand{\ProbeBkgLeftEeVsMmLinearOneeighteight}{0.0}  % experiments/FIGS/data/probe_ladder_v2_mlp2/analysis/seed_level_results.json :: levels.ee_vs_mm.linear[0].rejection_points['0.90'].mean_n_bkg_pass :: sha256 53e1ec81343ac2c1 :: background jets passing the cut, mean over seeds
\newcommand{\ProbeRejEeVsMmLinearOneeighteight}{\ensuremath{>}1{,}571}  % experiments/FIGS/data/probe_ladder_v2_mlp2/analysis/seed_level_results.json :: table[ee_vs_mm,linear,188].rejection_points['0.90'].rejection :: sha256 53e1ec81343ac2c1 :: at 90% signal efficiency, mean +- SD over seeds; 5 of 5 seeds at the cap
\newcommand{\ProbeAucEeVsMmLinearOnesixtwo}{\ensuremath{1^{\dagger}}}  % experiments/FIGS/data/probe_ladder_v2_mlp2/analysis/seed_level_results.json :: levels.ee_vs_mm.linear[1].mean_auc :: sha256 53e1ec81343ac2c1 :: 5 seeds, 5 saturated at AUC=1
\newcommand{\ProbeAucSdEeVsMmLinearOnesixtwo}{0.00000}  % experiments/FIGS/data/probe_ladder_v2_mlp2/analysis/seed_level_results.json :: table[ee_vs_mm,linear,162].auc (SD over seeds, ddof=1) :: sha256 53e1ec81343ac2c1
\newcommand{\ProbeLogOneMinusAucEeVsMmLinearOnesixtwo}{\ensuremath{-}16.9579}  % experiments/FIGS/data/probe_ladder_v2_mlp2/analysis/seed_level_results.json :: levels.ee_vs_mm.linear[1].mean :: sha256 53e1ec81343ac2c1 :: natural log of 1-AUC, lower is better
\newcommand{\ProbeLogOneMinusAucSdEeVsMmLinearOnesixtwo}{0.0000}  % experiments/FIGS/data/probe_ladder_v2_mlp2/analysis/seed_level_results.json :: levels.ee_vs_mm.linear[1].seed_sd :: sha256 53e1ec81343ac2c1
\newcommand{\ProbeOmaEeVsMmLinearOnesixtwo}{---}  % experiments/FIGS/data/probe_ladder_v2_mlp2/analysis/seed_level_results.json :: table[ee_vs_mm,linear,162].auc :: sha256 53e1ec81343ac2c1 :: not quoted: the AUC reached 1 in a seed, so 1-AUC is only a bound
\newcommand{\ProbeBkgLeftEeVsMmLinearOnesixtwo}{0.0}  % experiments/FIGS/data/probe_ladder_v2_mlp2/analysis/seed_level_results.json :: levels.ee_vs_mm.linear[1].rejection_points['0.90'].mean_n_bkg_pass :: sha256 53e1ec81343ac2c1 :: background jets passing the cut, mean over seeds
\newcommand{\ProbeRejEeVsMmLinearOnesixtwo}{\ensuremath{>}1{,}571}  % experiments/FIGS/data/probe_ladder_v2_mlp2/analysis/seed_level_results.json :: table[ee_vs_mm,linear,162].rejection_points['0.90'].rejection :: sha256 53e1ec81343ac2c1 :: at 90% signal efficiency, mean +- SD over seeds; 5 of 5 seeds at the cap
\newcommand{\ProbeAucEeVsMmLinearFourthree}{0.97519}  % experiments/FIGS/data/probe_ladder_v2_mlp2/analysis/seed_level_results.json :: levels.ee_vs_mm.linear[2].mean_auc :: sha256 53e1ec81343ac2c1 :: 5 seeds, 0 saturated at AUC=1
\newcommand{\ProbeAucSdEeVsMmLinearFourthree}{0.02073}  % experiments/FIGS/data/probe_ladder_v2_mlp2/analysis/seed_level_results.json :: table[ee_vs_mm,linear,43].auc (SD over seeds, ddof=1) :: sha256 53e1ec81343ac2c1
\newcommand{\ProbeLogOneMinusAucEeVsMmLinearFourthree}{\ensuremath{-}4.0326}  % experiments/FIGS/data/probe_ladder_v2_mlp2/analysis/seed_level_results.json :: levels.ee_vs_mm.linear[2].mean :: sha256 53e1ec81343ac2c1 :: natural log of 1-AUC, lower is better
\newcommand{\ProbeLogOneMinusAucSdEeVsMmLinearFourthree}{0.9467}  % experiments/FIGS/data/probe_ladder_v2_mlp2/analysis/seed_level_results.json :: levels.ee_vs_mm.linear[2].seed_sd :: sha256 53e1ec81343ac2c1
\newcommand{\ProbeOmaEeVsMmLinearFourthree}{\ensuremath{(2.5 \pm 2.1)\times10^{-2}}}  % experiments/FIGS/data/probe_ladder_v2_mlp2/analysis/seed_level_results.json :: table[ee_vs_mm,linear,43].auc :: sha256 53e1ec81343ac2c1 :: 1-AUC, mean +- SD over 5 seeds
\newcommand{\ProbeBkgLeftEeVsMmLinearFourthree}{155.8}  % experiments/FIGS/data/probe_ladder_v2_mlp2/analysis/seed_level_results.json :: levels.ee_vs_mm.linear[2].rejection_points['0.90'].mean_n_bkg_pass :: sha256 53e1ec81343ac2c1 :: background jets passing the cut, mean over seeds
\newcommand{\ProbeRejEeVsMmLinearFourthree}{80\,\ensuremath{\pm}\,70}  % experiments/FIGS/data/probe_ladder_v2_mlp2/analysis/seed_level_results.json :: table[ee_vs_mm,linear,43].rejection_points['0.90'].rejection :: sha256 53e1ec81343ac2c1 :: at 90% signal efficiency, mean +- SD over seeds; 0 of 5 seeds at the cap
\newcommand{\ProbeAucEeVsMmLinearOneseven}{0.95832}  % experiments/FIGS/data/probe_ladder_v2_mlp2/analysis/seed_level_results.json :: levels.ee_vs_mm.linear[3].mean_auc :: sha256 53e1ec81343ac2c1 :: 5 seeds, 0 saturated at AUC=1
\newcommand{\ProbeAucSdEeVsMmLinearOneseven}{0.00970}  % experiments/FIGS/data/probe_ladder_v2_mlp2/analysis/seed_level_results.json :: table[ee_vs_mm,linear,17].auc (SD over seeds, ddof=1) :: sha256 53e1ec81343ac2c1
\newcommand{\ProbeLogOneMinusAucEeVsMmLinearOneseven}{\ensuremath{-}3.2013}  % experiments/FIGS/data/probe_ladder_v2_mlp2/analysis/seed_level_results.json :: levels.ee_vs_mm.linear[3].mean :: sha256 53e1ec81343ac2c1 :: natural log of 1-AUC, lower is better
\newcommand{\ProbeLogOneMinusAucSdEeVsMmLinearOneseven}{0.2488}  % experiments/FIGS/data/probe_ladder_v2_mlp2/analysis/seed_level_results.json :: levels.ee_vs_mm.linear[3].seed_sd :: sha256 53e1ec81343ac2c1
\newcommand{\ProbeOmaEeVsMmLinearOneseven}{\ensuremath{(4.2 \pm 1.0)\times10^{-2}}}  % experiments/FIGS/data/probe_ladder_v2_mlp2/analysis/seed_level_results.json :: table[ee_vs_mm,linear,17].auc :: sha256 53e1ec81343ac2c1 :: 1-AUC, mean +- SD over 5 seeds
\newcommand{\ProbeBkgLeftEeVsMmLinearOneseven}{289.6}  % experiments/FIGS/data/probe_ladder_v2_mlp2/analysis/seed_level_results.json :: levels.ee_vs_mm.linear[3].rejection_points['0.90'].mean_n_bkg_pass :: sha256 53e1ec81343ac2c1 :: background jets passing the cut, mean over seeds
\newcommand{\ProbeRejEeVsMmLinearOneseven}{17\,\ensuremath{\pm}\,6}  % experiments/FIGS/data/probe_ladder_v2_mlp2/analysis/seed_level_results.json :: table[ee_vs_mm,linear,17].rejection_points['0.90'].rejection :: sha256 53e1ec81343ac2c1 :: at 90% signal efficiency, mean +- SD over seeds; 0 of 5 seeds at the cap
\newcommand{\ProbeOmaRatioEeVsMmLinearOnesevenOverFourthree}{1.68}  % experiments/FIGS/data/probe_ladder_v2_mlp2/analysis/seed_level_results.json :: table[ee_vs_mm,linear,17/43].auc :: sha256 53e1ec81343ac2c1 :: ratio of the seed means of 1-AUC, coarser over finer
\newcommand{\ProbeAucEeVsMmMlpOneeighteight}{\ensuremath{1^{\dagger}}}  % experiments/FIGS/data/probe_ladder_v2_mlp2/analysis/seed_level_results.json :: levels.ee_vs_mm.mlp[0].mean_auc :: sha256 53e1ec81343ac2c1 :: 5 seeds, 5 saturated at AUC=1
\newcommand{\ProbeAucSdEeVsMmMlpOneeighteight}{0.00000}  % experiments/FIGS/data/probe_ladder_v2_mlp2/analysis/seed_level_results.json :: table[ee_vs_mm,mlp,188].auc (SD over seeds, ddof=1) :: sha256 53e1ec81343ac2c1
\newcommand{\ProbeLogOneMinusAucEeVsMmMlpOneeighteight}{\ensuremath{-}16.9579}  % experiments/FIGS/data/probe_ladder_v2_mlp2/analysis/seed_level_results.json :: levels.ee_vs_mm.mlp[0].mean :: sha256 53e1ec81343ac2c1 :: natural log of 1-AUC, lower is better
\newcommand{\ProbeLogOneMinusAucSdEeVsMmMlpOneeighteight}{0.0000}  % experiments/FIGS/data/probe_ladder_v2_mlp2/analysis/seed_level_results.json :: levels.ee_vs_mm.mlp[0].seed_sd :: sha256 53e1ec81343ac2c1
\newcommand{\ProbeOmaEeVsMmMlpOneeighteight}{---}  % experiments/FIGS/data/probe_ladder_v2_mlp2/analysis/seed_level_results.json :: table[ee_vs_mm,mlp,188].auc :: sha256 53e1ec81343ac2c1 :: not quoted: the AUC reached 1 in a seed, so 1-AUC is only a bound
\newcommand{\ProbeBkgLeftEeVsMmMlpOneeighteight}{0.0}  % experiments/FIGS/data/probe_ladder_v2_mlp2/analysis/seed_level_results.json :: levels.ee_vs_mm.mlp[0].rejection_points['0.90'].mean_n_bkg_pass :: sha256 53e1ec81343ac2c1 :: background jets passing the cut, mean over seeds
\newcommand{\ProbeRejEeVsMmMlpOneeighteight}{\ensuremath{>}1{,}571}  % experiments/FIGS/data/probe_ladder_v2_mlp2/analysis/seed_level_results.json :: table[ee_vs_mm,mlp,188].rejection_points['0.90'].rejection :: sha256 53e1ec81343ac2c1 :: at 90% signal efficiency, mean +- SD over seeds; 5 of 5 seeds at the cap
\newcommand{\ProbeAucEeVsMmMlpOnesixtwo}{\ensuremath{1^{\dagger}}}  % experiments/FIGS/data/probe_ladder_v2_mlp2/analysis/seed_level_results.json :: levels.ee_vs_mm.mlp[1].mean_auc :: sha256 53e1ec81343ac2c1 :: 5 seeds, 5 saturated at AUC=1
\newcommand{\ProbeAucSdEeVsMmMlpOnesixtwo}{0.00000}  % experiments/FIGS/data/probe_ladder_v2_mlp2/analysis/seed_level_results.json :: table[ee_vs_mm,mlp,162].auc (SD over seeds, ddof=1) :: sha256 53e1ec81343ac2c1
\newcommand{\ProbeLogOneMinusAucEeVsMmMlpOnesixtwo}{\ensuremath{-}16.9579}  % experiments/FIGS/data/probe_ladder_v2_mlp2/analysis/seed_level_results.json :: levels.ee_vs_mm.mlp[1].mean :: sha256 53e1ec81343ac2c1 :: natural log of 1-AUC, lower is better
\newcommand{\ProbeLogOneMinusAucSdEeVsMmMlpOnesixtwo}{0.0000}  % experiments/FIGS/data/probe_ladder_v2_mlp2/analysis/seed_level_results.json :: levels.ee_vs_mm.mlp[1].seed_sd :: sha256 53e1ec81343ac2c1
\newcommand{\ProbeOmaEeVsMmMlpOnesixtwo}{---}  % experiments/FIGS/data/probe_ladder_v2_mlp2/analysis/seed_level_results.json :: table[ee_vs_mm,mlp,162].auc :: sha256 53e1ec81343ac2c1 :: not quoted: the AUC reached 1 in a seed, so 1-AUC is only a bound
\newcommand{\ProbeBkgLeftEeVsMmMlpOnesixtwo}{0.0}  % experiments/FIGS/data/probe_ladder_v2_mlp2/analysis/seed_level_results.json :: levels.ee_vs_mm.mlp[1].rejection_points['0.90'].mean_n_bkg_pass :: sha256 53e1ec81343ac2c1 :: background jets passing the cut, mean over seeds
\newcommand{\ProbeRejEeVsMmMlpOnesixtwo}{\ensuremath{>}1{,}571}  % experiments/FIGS/data/probe_ladder_v2_mlp2/analysis/seed_level_results.json :: table[ee_vs_mm,mlp,162].rejection_points['0.90'].rejection :: sha256 53e1ec81343ac2c1 :: at 90% signal efficiency, mean +- SD over seeds; 5 of 5 seeds at the cap
\newcommand{\ProbeAucEeVsMmMlpFourthree}{0.99657}  % experiments/FIGS/data/probe_ladder_v2_mlp2/analysis/seed_level_results.json :: levels.ee_vs_mm.mlp[2].mean_auc :: sha256 53e1ec81343ac2c1 :: 5 seeds, 0 saturated at AUC=1
\newcommand{\ProbeAucSdEeVsMmMlpFourthree}{0.00300}  % experiments/FIGS/data/probe_ladder_v2_mlp2/analysis/seed_level_results.json :: table[ee_vs_mm,mlp,43].auc (SD over seeds, ddof=1) :: sha256 53e1ec81343ac2c1
\newcommand{\ProbeLogOneMinusAucEeVsMmMlpFourthree}{\ensuremath{-}6.0371}  % experiments/FIGS/data/probe_ladder_v2_mlp2/analysis/seed_level_results.json :: levels.ee_vs_mm.mlp[2].mean :: sha256 53e1ec81343ac2c1 :: natural log of 1-AUC, lower is better
\newcommand{\ProbeLogOneMinusAucSdEeVsMmMlpFourthree}{0.9962}  % experiments/FIGS/data/probe_ladder_v2_mlp2/analysis/seed_level_results.json :: levels.ee_vs_mm.mlp[2].seed_sd :: sha256 53e1ec81343ac2c1
\newcommand{\ProbeOmaEeVsMmMlpFourthree}{\ensuremath{(3.4 \pm 3.0)\times10^{-3}}}  % experiments/FIGS/data/probe_ladder_v2_mlp2/analysis/seed_level_results.json :: table[ee_vs_mm,mlp,43].auc :: sha256 53e1ec81343ac2c1 :: 1-AUC, mean +- SD over 5 seeds
\newcommand{\ProbeBkgLeftEeVsMmMlpFourthree}{33.4}  % experiments/FIGS/data/probe_ladder_v2_mlp2/analysis/seed_level_results.json :: levels.ee_vs_mm.mlp[2].rejection_points['0.90'].mean_n_bkg_pass :: sha256 53e1ec81343ac2c1 :: background jets passing the cut, mean over seeds
\newcommand{\ProbeRejEeVsMmMlpFourthree}{400\,\ensuremath{\pm}\,500}  % experiments/FIGS/data/probe_ladder_v2_mlp2/analysis/seed_level_results.json :: table[ee_vs_mm,mlp,43].rejection_points['0.90'].rejection :: sha256 53e1ec81343ac2c1 :: at 90% signal efficiency, mean +- SD over seeds; 0 of 5 seeds at the cap
\newcommand{\ProbeAucEeVsMmMlpOneseven}{0.99338}  % experiments/FIGS/data/probe_ladder_v2_mlp2/analysis/seed_level_results.json :: levels.ee_vs_mm.mlp[3].mean_auc :: sha256 53e1ec81343ac2c1 :: 5 seeds, 0 saturated at AUC=1
\newcommand{\ProbeAucSdEeVsMmMlpOneseven}{0.00065}  % experiments/FIGS/data/probe_ladder_v2_mlp2/analysis/seed_level_results.json :: table[ee_vs_mm,mlp,17].auc (SD over seeds, ddof=1) :: sha256 53e1ec81343ac2c1
\newcommand{\ProbeLogOneMinusAucEeVsMmMlpOneseven}{\ensuremath{-}5.0211}  % experiments/FIGS/data/probe_ladder_v2_mlp2/analysis/seed_level_results.json :: levels.ee_vs_mm.mlp[3].mean :: sha256 53e1ec81343ac2c1 :: natural log of 1-AUC, lower is better
\newcommand{\ProbeLogOneMinusAucSdEeVsMmMlpOneseven}{0.0997}  % experiments/FIGS/data/probe_ladder_v2_mlp2/analysis/seed_level_results.json :: levels.ee_vs_mm.mlp[3].seed_sd :: sha256 53e1ec81343ac2c1
\newcommand{\ProbeOmaEeVsMmMlpOneseven}{\ensuremath{(6.6 \pm 0.7)\times10^{-3}}}  % experiments/FIGS/data/probe_ladder_v2_mlp2/analysis/seed_level_results.json :: table[ee_vs_mm,mlp,17].auc :: sha256 53e1ec81343ac2c1 :: 1-AUC, mean +- SD over 5 seeds
\newcommand{\ProbeBkgLeftEeVsMmMlpOneseven}{59.2}  % experiments/FIGS/data/probe_ladder_v2_mlp2/analysis/seed_level_results.json :: levels.ee_vs_mm.mlp[3].rejection_points['0.90'].mean_n_bkg_pass :: sha256 53e1ec81343ac2c1 :: background jets passing the cut, mean over seeds
\newcommand{\ProbeRejEeVsMmMlpOneseven}{81\,\ensuremath{\pm}\,12}  % experiments/FIGS/data/probe_ladder_v2_mlp2/analysis/seed_level_results.json :: table[ee_vs_mm,mlp,17].rejection_points['0.90'].rejection :: sha256 53e1ec81343ac2c1 :: at 90% signal efficiency, mean +- SD over seeds; 0 of 5 seeds at the cap
\newcommand{\ProbeOmaRatioEeVsMmMlpOnesevenOverFourthree}{1.93}  % experiments/FIGS/data/probe_ladder_v2_mlp2/analysis/seed_level_results.json :: table[ee_vs_mm,mlp,17/43].auc :: sha256 53e1ec81343ac2c1 :: ratio of the seed means of 1-AUC, coarser over finer
\newcommand{\ProbeAucRetainedTopologyLinearOneeighteight}{0.99413}  % experiments/FIGS/data/probe_ladder_v2_mlp2/analysis/seed_level_results.json :: levels.retained_topology.linear[0].mean_auc :: sha256 53e1ec81343ac2c1 :: 5 seeds, 0 saturated at AUC=1
\newcommand{\ProbeAucSdRetainedTopologyLinearOneeighteight}{0.00011}  % experiments/FIGS/data/probe_ladder_v2_mlp2/analysis/seed_level_results.json :: table[retained_topology,linear,188].auc (SD over seeds, ddof=1) :: sha256 53e1ec81343ac2c1
\newcommand{\ProbeLogOneMinusAucRetainedTopologyLinearOneeighteight}{\ensuremath{-}5.1388}  % experiments/FIGS/data/probe_ladder_v2_mlp2/analysis/seed_level_results.json :: levels.retained_topology.linear[0].mean :: sha256 53e1ec81343ac2c1 :: natural log of 1-AUC, lower is better
\newcommand{\ProbeLogOneMinusAucSdRetainedTopologyLinearOneeighteight}{0.0187}  % experiments/FIGS/data/probe_ladder_v2_mlp2/analysis/seed_level_results.json :: levels.retained_topology.linear[0].seed_sd :: sha256 53e1ec81343ac2c1
\newcommand{\ProbeOmaRetainedTopologyLinearOneeighteight}{\ensuremath{(5.87 \pm 0.11)\times10^{-3}}}  % experiments/FIGS/data/probe_ladder_v2_mlp2/analysis/seed_level_results.json :: table[retained_topology,linear,188].auc :: sha256 53e1ec81343ac2c1 :: 1-AUC, mean +- SD over 5 seeds
\newcommand{\ProbeBkgLeftRetainedTopologyLinearOneeighteight}{15.4}  % experiments/FIGS/data/probe_ladder_v2_mlp2/analysis/seed_level_results.json :: levels.retained_topology.linear[0].rejection_points['0.90'].mean_n_bkg_pass :: sha256 53e1ec81343ac2c1 :: background jets passing the cut, mean over seeds
\newcommand{\ProbeRejRetainedTopologyLinearOneeighteight}{71\,\ensuremath{\pm}\,7}  % experiments/FIGS/data/probe_ladder_v2_mlp2/analysis/seed_level_results.json :: table[retained_topology,linear,188].rejection_points['0.90'].rejection :: sha256 53e1ec81343ac2c1 :: at 90% signal efficiency, mean +- SD over seeds; 0 of 5 seeds at the cap
\newcommand{\ProbeAucRetainedTopologyLinearOnesixtwo}{0.99413}  % experiments/FIGS/data/probe_ladder_v2_mlp2/analysis/seed_level_results.json :: levels.retained_topology.linear[1].mean_auc :: sha256 53e1ec81343ac2c1 :: 5 seeds, 0 saturated at AUC=1
\newcommand{\ProbeAucSdRetainedTopologyLinearOnesixtwo}{0.00030}  % experiments/FIGS/data/probe_ladder_v2_mlp2/analysis/seed_level_results.json :: table[retained_topology,linear,162].auc (SD over seeds, ddof=1) :: sha256 53e1ec81343ac2c1
\newcommand{\ProbeLogOneMinusAucRetainedTopologyLinearOnesixtwo}{\ensuremath{-}5.1395}  % experiments/FIGS/data/probe_ladder_v2_mlp2/analysis/seed_level_results.json :: levels.retained_topology.linear[1].mean :: sha256 53e1ec81343ac2c1 :: natural log of 1-AUC, lower is better
\newcommand{\ProbeLogOneMinusAucSdRetainedTopologyLinearOnesixtwo}{0.0510}  % experiments/FIGS/data/probe_ladder_v2_mlp2/analysis/seed_level_results.json :: levels.retained_topology.linear[1].seed_sd :: sha256 53e1ec81343ac2c1
\newcommand{\ProbeOmaRetainedTopologyLinearOnesixtwo}{\ensuremath{(5.87 \pm 0.30)\times10^{-3}}}  % experiments/FIGS/data/probe_ladder_v2_mlp2/analysis/seed_level_results.json :: table[retained_topology,linear,162].auc :: sha256 53e1ec81343ac2c1 :: 1-AUC, mean +- SD over 5 seeds
\newcommand{\ProbeBkgLeftRetainedTopologyLinearOnesixtwo}{14.0}  % experiments/FIGS/data/probe_ladder_v2_mlp2/analysis/seed_level_results.json :: levels.retained_topology.linear[1].rejection_points['0.90'].mean_n_bkg_pass :: sha256 53e1ec81343ac2c1 :: background jets passing the cut, mean over seeds
\newcommand{\ProbeRejRetainedTopologyLinearOnesixtwo}{79\,\ensuremath{\pm}\,13}  % experiments/FIGS/data/probe_ladder_v2_mlp2/analysis/seed_level_results.json :: table[retained_topology,linear,162].rejection_points['0.90'].rejection :: sha256 53e1ec81343ac2c1 :: at 90% signal efficiency, mean +- SD over seeds; 0 of 5 seeds at the cap
\newcommand{\ProbeAucRetainedTopologyLinearFourthree}{0.99274}  % experiments/FIGS/data/probe_ladder_v2_mlp2/analysis/seed_level_results.json :: levels.retained_topology.linear[2].mean_auc :: sha256 53e1ec81343ac2c1 :: 5 seeds, 0 saturated at AUC=1
\newcommand{\ProbeAucSdRetainedTopologyLinearFourthree}{0.00012}  % experiments/FIGS/data/probe_ladder_v2_mlp2/analysis/seed_level_results.json :: table[retained_topology,linear,43].auc (SD over seeds, ddof=1) :: sha256 53e1ec81343ac2c1
\newcommand{\ProbeLogOneMinusAucRetainedTopologyLinearFourthree}{\ensuremath{-}4.9253}  % experiments/FIGS/data/probe_ladder_v2_mlp2/analysis/seed_level_results.json :: levels.retained_topology.linear[2].mean :: sha256 53e1ec81343ac2c1 :: natural log of 1-AUC, lower is better
\newcommand{\ProbeLogOneMinusAucSdRetainedTopologyLinearFourthree}{0.0169}  % experiments/FIGS/data/probe_ladder_v2_mlp2/analysis/seed_level_results.json :: levels.retained_topology.linear[2].seed_sd :: sha256 53e1ec81343ac2c1
\newcommand{\ProbeOmaRetainedTopologyLinearFourthree}{\ensuremath{(7.26 \pm 0.12)\times10^{-3}}}  % experiments/FIGS/data/probe_ladder_v2_mlp2/analysis/seed_level_results.json :: table[retained_topology,linear,43].auc :: sha256 53e1ec81343ac2c1 :: 1-AUC, mean +- SD over 5 seeds
\newcommand{\ProbeBkgLeftRetainedTopologyLinearFourthree}{16.4}  % experiments/FIGS/data/probe_ladder_v2_mlp2/analysis/seed_level_results.json :: levels.retained_topology.linear[2].rejection_points['0.90'].mean_n_bkg_pass :: sha256 53e1ec81343ac2c1 :: background jets passing the cut, mean over seeds
\newcommand{\ProbeRejRetainedTopologyLinearFourthree}{66\,\ensuremath{\pm}\,5}  % experiments/FIGS/data/probe_ladder_v2_mlp2/analysis/seed_level_results.json :: table[retained_topology,linear,43].rejection_points['0.90'].rejection :: sha256 53e1ec81343ac2c1 :: at 90% signal efficiency, mean +- SD over seeds; 0 of 5 seeds at the cap
\newcommand{\ProbeAucRetainedTopologyLinearOneseven}{0.99096}  % experiments/FIGS/data/probe_ladder_v2_mlp2/analysis/seed_level_results.json :: levels.retained_topology.linear[3].mean_auc :: sha256 53e1ec81343ac2c1 :: 5 seeds, 0 saturated at AUC=1
\newcommand{\ProbeAucSdRetainedTopologyLinearOneseven}{0.00092}  % experiments/FIGS/data/probe_ladder_v2_mlp2/analysis/seed_level_results.json :: table[retained_topology,linear,17].auc (SD over seeds, ddof=1) :: sha256 53e1ec81343ac2c1
\newcommand{\ProbeLogOneMinusAucRetainedTopologyLinearOneseven}{\ensuremath{-}4.7095}  % experiments/FIGS/data/probe_ladder_v2_mlp2/analysis/seed_level_results.json :: levels.retained_topology.linear[3].mean :: sha256 53e1ec81343ac2c1 :: natural log of 1-AUC, lower is better
\newcommand{\ProbeLogOneMinusAucSdRetainedTopologyLinearOneseven}{0.0979}  % experiments/FIGS/data/probe_ladder_v2_mlp2/analysis/seed_level_results.json :: levels.retained_topology.linear[3].seed_sd :: sha256 53e1ec81343ac2c1
\newcommand{\ProbeOmaRetainedTopologyLinearOneseven}{\ensuremath{(9.0 \pm 0.9)\times10^{-3}}}  % experiments/FIGS/data/probe_ladder_v2_mlp2/analysis/seed_level_results.json :: table[retained_topology,linear,17].auc :: sha256 53e1ec81343ac2c1 :: 1-AUC, mean +- SD over 5 seeds
\newcommand{\ProbeBkgLeftRetainedTopologyLinearOneseven}{21.8}  % experiments/FIGS/data/probe_ladder_v2_mlp2/analysis/seed_level_results.json :: levels.retained_topology.linear[3].rejection_points['0.90'].mean_n_bkg_pass :: sha256 53e1ec81343ac2c1 :: background jets passing the cut, mean over seeds
\newcommand{\ProbeRejRetainedTopologyLinearOneseven}{51\,\ensuremath{\pm}\,11}  % experiments/FIGS/data/probe_ladder_v2_mlp2/analysis/seed_level_results.json :: table[retained_topology,linear,17].rejection_points['0.90'].rejection :: sha256 53e1ec81343ac2c1 :: at 90% signal efficiency, mean +- SD over seeds; 0 of 5 seeds at the cap
\newcommand{\ProbeOmaRatioRetainedTopologyLinearOnesixtwoOverOneeighteight}{1.00}  % experiments/FIGS/data/probe_ladder_v2_mlp2/analysis/seed_level_results.json :: table[retained_topology,linear,162/188].auc :: sha256 53e1ec81343ac2c1 :: ratio of the seed means of 1-AUC, coarser over finer
\newcommand{\ProbeOmaRatioRetainedTopologyLinearFourthreeOverOneeighteight}{1.24}  % experiments/FIGS/data/probe_ladder_v2_mlp2/analysis/seed_level_results.json :: table[retained_topology,linear,43/188].auc :: sha256 53e1ec81343ac2c1 :: ratio of the seed means of 1-AUC, coarser over finer
\newcommand{\ProbeOmaRatioRetainedTopologyLinearOnesevenOverOneeighteight}{1.54}  % experiments/FIGS/data/probe_ladder_v2_mlp2/analysis/seed_level_results.json :: table[retained_topology,linear,17/188].auc :: sha256 53e1ec81343ac2c1 :: ratio of the seed means of 1-AUC, coarser over finer
\newcommand{\ProbeOmaRatioRetainedTopologyLinearFourthreeOverOnesixtwo}{1.24}  % experiments/FIGS/data/probe_ladder_v2_mlp2/analysis/seed_level_results.json :: table[retained_topology,linear,43/162].auc :: sha256 53e1ec81343ac2c1 :: ratio of the seed means of 1-AUC, coarser over finer
\newcommand{\ProbeOmaRatioRetainedTopologyLinearOnesevenOverOnesixtwo}{1.54}  % experiments/FIGS/data/probe_ladder_v2_mlp2/analysis/seed_level_results.json :: table[retained_topology,linear,17/162].auc :: sha256 53e1ec81343ac2c1 :: ratio of the seed means of 1-AUC, coarser over finer
\newcommand{\ProbeOmaRatioRetainedTopologyLinearOnesevenOverFourthree}{1.25}  % experiments/FIGS/data/probe_ladder_v2_mlp2/analysis/seed_level_results.json :: table[retained_topology,linear,17/43].auc :: sha256 53e1ec81343ac2c1 :: ratio of the seed means of 1-AUC, coarser over finer
\newcommand{\ProbeAucRetainedTopologyMlpOneeighteight}{0.99394}  % experiments/FIGS/data/probe_ladder_v2_mlp2/analysis/seed_level_results.json :: levels.retained_topology.mlp[0].mean_auc :: sha256 53e1ec81343ac2c1 :: 5 seeds, 0 saturated at AUC=1
\newcommand{\ProbeAucSdRetainedTopologyMlpOneeighteight}{0.00032}  % experiments/FIGS/data/probe_ladder_v2_mlp2/analysis/seed_level_results.json :: table[retained_topology,mlp,188].auc (SD over seeds, ddof=1) :: sha256 53e1ec81343ac2c1
\newcommand{\ProbeLogOneMinusAucRetainedTopologyMlpOneeighteight}{\ensuremath{-}5.1078}  % experiments/FIGS/data/probe_ladder_v2_mlp2/analysis/seed_level_results.json :: levels.retained_topology.mlp[0].mean :: sha256 53e1ec81343ac2c1 :: natural log of 1-AUC, lower is better
\newcommand{\ProbeLogOneMinusAucSdRetainedTopologyMlpOneeighteight}{0.0522}  % experiments/FIGS/data/probe_ladder_v2_mlp2/analysis/seed_level_results.json :: levels.retained_topology.mlp[0].seed_sd :: sha256 53e1ec81343ac2c1
\newcommand{\ProbeOmaRetainedTopologyMlpOneeighteight}{\ensuremath{(6.06 \pm 0.32)\times10^{-3}}}  % experiments/FIGS/data/probe_ladder_v2_mlp2/analysis/seed_level_results.json :: table[retained_topology,mlp,188].auc :: sha256 53e1ec81343ac2c1 :: 1-AUC, mean +- SD over 5 seeds
\newcommand{\ProbeBkgLeftRetainedTopologyMlpOneeighteight}{15.8}  % experiments/FIGS/data/probe_ladder_v2_mlp2/analysis/seed_level_results.json :: levels.retained_topology.mlp[0].rejection_points['0.90'].mean_n_bkg_pass :: sha256 53e1ec81343ac2c1 :: background jets passing the cut, mean over seeds
\newcommand{\ProbeRejRetainedTopologyMlpOneeighteight}{69\,\ensuremath{\pm}\,9}  % experiments/FIGS/data/probe_ladder_v2_mlp2/analysis/seed_level_results.json :: table[retained_topology,mlp,188].rejection_points['0.90'].rejection :: sha256 53e1ec81343ac2c1 :: at 90% signal efficiency, mean +- SD over seeds; 0 of 5 seeds at the cap
\newcommand{\ProbeAucRetainedTopologyMlpOnesixtwo}{0.99387}  % experiments/FIGS/data/probe_ladder_v2_mlp2/analysis/seed_level_results.json :: levels.retained_topology.mlp[1].mean_auc :: sha256 53e1ec81343ac2c1 :: 5 seeds, 0 saturated at AUC=1
\newcommand{\ProbeAucSdRetainedTopologyMlpOnesixtwo}{0.00048}  % experiments/FIGS/data/probe_ladder_v2_mlp2/analysis/seed_level_results.json :: table[retained_topology,mlp,162].auc (SD over seeds, ddof=1) :: sha256 53e1ec81343ac2c1
\newcommand{\ProbeLogOneMinusAucRetainedTopologyMlpOnesixtwo}{\ensuremath{-}5.0964}  % experiments/FIGS/data/probe_ladder_v2_mlp2/analysis/seed_level_results.json :: levels.retained_topology.mlp[1].mean :: sha256 53e1ec81343ac2c1 :: natural log of 1-AUC, lower is better
\newcommand{\ProbeLogOneMinusAucSdRetainedTopologyMlpOnesixtwo}{0.0773}  % experiments/FIGS/data/probe_ladder_v2_mlp2/analysis/seed_level_results.json :: levels.retained_topology.mlp[1].seed_sd :: sha256 53e1ec81343ac2c1
\newcommand{\ProbeOmaRetainedTopologyMlpOnesixtwo}{\ensuremath{(6.1 \pm 0.5)\times10^{-3}}}  % experiments/FIGS/data/probe_ladder_v2_mlp2/analysis/seed_level_results.json :: table[retained_topology,mlp,162].auc :: sha256 53e1ec81343ac2c1 :: 1-AUC, mean +- SD over 5 seeds
\newcommand{\ProbeBkgLeftRetainedTopologyMlpOnesixtwo}{16.6}  % experiments/FIGS/data/probe_ladder_v2_mlp2/analysis/seed_level_results.json :: levels.retained_topology.mlp[1].rejection_points['0.90'].mean_n_bkg_pass :: sha256 53e1ec81343ac2c1 :: background jets passing the cut, mean over seeds
\newcommand{\ProbeRejRetainedTopologyMlpOnesixtwo}{67\,\ensuremath{\pm}\,14}  % experiments/FIGS/data/probe_ladder_v2_mlp2/analysis/seed_level_results.json :: table[retained_topology,mlp,162].rejection_points['0.90'].rejection :: sha256 53e1ec81343ac2c1 :: at 90% signal efficiency, mean +- SD over seeds; 0 of 5 seeds at the cap
\newcommand{\ProbeAucRetainedTopologyMlpFourthree}{0.99284}  % experiments/FIGS/data/probe_ladder_v2_mlp2/analysis/seed_level_results.json :: levels.retained_topology.mlp[2].mean_auc :: sha256 53e1ec81343ac2c1 :: 5 seeds, 0 saturated at AUC=1
\newcommand{\ProbeAucSdRetainedTopologyMlpFourthree}{0.00040}  % experiments/FIGS/data/probe_ladder_v2_mlp2/analysis/seed_level_results.json :: table[retained_topology,mlp,43].auc (SD over seeds, ddof=1) :: sha256 53e1ec81343ac2c1
\newcommand{\ProbeLogOneMinusAucRetainedTopologyMlpFourthree}{\ensuremath{-}4.9400}  % experiments/FIGS/data/probe_ladder_v2_mlp2/analysis/seed_level_results.json :: levels.retained_topology.mlp[2].mean :: sha256 53e1ec81343ac2c1 :: natural log of 1-AUC, lower is better
\newcommand{\ProbeLogOneMinusAucSdRetainedTopologyMlpFourthree}{0.0570}  % experiments/FIGS/data/probe_ladder_v2_mlp2/analysis/seed_level_results.json :: levels.retained_topology.mlp[2].seed_sd :: sha256 53e1ec81343ac2c1
\newcommand{\ProbeOmaRetainedTopologyMlpFourthree}{\ensuremath{(7.2 \pm 0.4)\times10^{-3}}}  % experiments/FIGS/data/probe_ladder_v2_mlp2/analysis/seed_level_results.json :: table[retained_topology,mlp,43].auc :: sha256 53e1ec81343ac2c1 :: 1-AUC, mean +- SD over 5 seeds
\newcommand{\ProbeBkgLeftRetainedTopologyMlpFourthree}{17.8}  % experiments/FIGS/data/probe_ladder_v2_mlp2/analysis/seed_level_results.json :: levels.retained_topology.mlp[2].rejection_points['0.90'].mean_n_bkg_pass :: sha256 53e1ec81343ac2c1 :: background jets passing the cut, mean over seeds
\newcommand{\ProbeRejRetainedTopologyMlpFourthree}{61\,\ensuremath{\pm}\,7}  % experiments/FIGS/data/probe_ladder_v2_mlp2/analysis/seed_level_results.json :: table[retained_topology,mlp,43].rejection_points['0.90'].rejection :: sha256 53e1ec81343ac2c1 :: at 90% signal efficiency, mean +- SD over seeds; 0 of 5 seeds at the cap
\newcommand{\ProbeAucRetainedTopologyMlpOneseven}{0.99081}  % experiments/FIGS/data/probe_ladder_v2_mlp2/analysis/seed_level_results.json :: levels.retained_topology.mlp[3].mean_auc :: sha256 53e1ec81343ac2c1 :: 5 seeds, 0 saturated at AUC=1
\newcommand{\ProbeAucSdRetainedTopologyMlpOneseven}{0.00062}  % experiments/FIGS/data/probe_ladder_v2_mlp2/analysis/seed_level_results.json :: table[retained_topology,mlp,17].auc (SD over seeds, ddof=1) :: sha256 53e1ec81343ac2c1
\newcommand{\ProbeLogOneMinusAucRetainedTopologyMlpOneseven}{\ensuremath{-}4.6913}  % experiments/FIGS/data/probe_ladder_v2_mlp2/analysis/seed_level_results.json :: levels.retained_topology.mlp[3].mean :: sha256 53e1ec81343ac2c1 :: natural log of 1-AUC, lower is better
\newcommand{\ProbeLogOneMinusAucSdRetainedTopologyMlpOneseven}{0.0658}  % experiments/FIGS/data/probe_ladder_v2_mlp2/analysis/seed_level_results.json :: levels.retained_topology.mlp[3].seed_sd :: sha256 53e1ec81343ac2c1
\newcommand{\ProbeOmaRetainedTopologyMlpOneseven}{\ensuremath{(9.2 \pm 0.6)\times10^{-3}}}  % experiments/FIGS/data/probe_ladder_v2_mlp2/analysis/seed_level_results.json :: table[retained_topology,mlp,17].auc :: sha256 53e1ec81343ac2c1 :: 1-AUC, mean +- SD over 5 seeds
\newcommand{\ProbeBkgLeftRetainedTopologyMlpOneseven}{21.2}  % experiments/FIGS/data/probe_ladder_v2_mlp2/analysis/seed_level_results.json :: levels.retained_topology.mlp[3].rejection_points['0.90'].mean_n_bkg_pass :: sha256 53e1ec81343ac2c1 :: background jets passing the cut, mean over seeds
\newcommand{\ProbeRejRetainedTopologyMlpOneseven}{51\,\ensuremath{\pm}\,4}  % experiments/FIGS/data/probe_ladder_v2_mlp2/analysis/seed_level_results.json :: table[retained_topology,mlp,17].rejection_points['0.90'].rejection :: sha256 53e1ec81343ac2c1 :: at 90% signal efficiency, mean +- SD over seeds; 0 of 5 seeds at the cap
\newcommand{\ProbeOmaRatioRetainedTopologyMlpOnesixtwoOverOneeighteight}{1.01}  % experiments/FIGS/data/probe_ladder_v2_mlp2/analysis/seed_level_results.json :: table[retained_topology,mlp,162/188].auc :: sha256 53e1ec81343ac2c1 :: ratio of the seed means of 1-AUC, coarser over finer
\newcommand{\ProbeOmaRatioRetainedTopologyMlpFourthreeOverOneeighteight}{1.18}  % experiments/FIGS/data/probe_ladder_v2_mlp2/analysis/seed_level_results.json :: table[retained_topology,mlp,43/188].auc :: sha256 53e1ec81343ac2c1 :: ratio of the seed means of 1-AUC, coarser over finer
\newcommand{\ProbeOmaRatioRetainedTopologyMlpOnesevenOverOneeighteight}{1.52}  % experiments/FIGS/data/probe_ladder_v2_mlp2/analysis/seed_level_results.json :: table[retained_topology,mlp,17/188].auc :: sha256 53e1ec81343ac2c1 :: ratio of the seed means of 1-AUC, coarser over finer
\newcommand{\ProbeOmaRatioRetainedTopologyMlpFourthreeOverOnesixtwo}{1.17}  % experiments/FIGS/data/probe_ladder_v2_mlp2/analysis/seed_level_results.json :: table[retained_topology,mlp,43/162].auc :: sha256 53e1ec81343ac2c1 :: ratio of the seed means of 1-AUC, coarser over finer
\newcommand{\ProbeOmaRatioRetainedTopologyMlpOnesevenOverOnesixtwo}{1.50}  % experiments/FIGS/data/probe_ladder_v2_mlp2/analysis/seed_level_results.json :: table[retained_topology,mlp,17/162].auc :: sha256 53e1ec81343ac2c1 :: ratio of the seed means of 1-AUC, coarser over finer
\newcommand{\ProbeOmaRatioRetainedTopologyMlpOnesevenOverFourthree}{1.28}  % experiments/FIGS/data/probe_ladder_v2_mlp2/analysis/seed_level_results.json :: table[retained_topology,mlp,17/43].auc :: sha256 53e1ec81343ac2c1 :: ratio of the seed means of 1-AUC, coarser over finer
\newcommand{\ProbeAucVisibleContentLinearOneeighteight}{0.99917}  % experiments/FIGS/data/probe_ladder_v2_mlp2/analysis/seed_level_results.json :: levels.visible_content.linear[0].mean_auc :: sha256 53e1ec81343ac2c1 :: 5 seeds, 0 saturated at AUC=1
\newcommand{\ProbeAucSdVisibleContentLinearOneeighteight}{0.00003}  % experiments/FIGS/data/probe_ladder_v2_mlp2/analysis/seed_level_results.json :: table[visible_content,linear,188].auc (SD over seeds, ddof=1) :: sha256 53e1ec81343ac2c1
\newcommand{\ProbeLogOneMinusAucVisibleContentLinearOneeighteight}{\ensuremath{-}7.0999}  % experiments/FIGS/data/probe_ladder_v2_mlp2/analysis/seed_level_results.json :: levels.visible_content.linear[0].mean :: sha256 53e1ec81343ac2c1 :: natural log of 1-AUC, lower is better
\newcommand{\ProbeLogOneMinusAucSdVisibleContentLinearOneeighteight}{0.0358}  % experiments/FIGS/data/probe_ladder_v2_mlp2/analysis/seed_level_results.json :: levels.visible_content.linear[0].seed_sd :: sha256 53e1ec81343ac2c1
\newcommand{\ProbeOmaVisibleContentLinearOneeighteight}{\ensuremath{(8.26 \pm 0.30)\times10^{-4}}}  % experiments/FIGS/data/probe_ladder_v2_mlp2/analysis/seed_level_results.json :: table[visible_content,linear,188].auc :: sha256 53e1ec81343ac2c1 :: 1-AUC, mean +- SD over 5 seeds
\newcommand{\ProbeBkgLeftVisibleContentLinearOneeighteight}{0.4}  % experiments/FIGS/data/probe_ladder_v2_mlp2/analysis/seed_level_results.json :: levels.visible_content.linear[0].rejection_points['0.90'].mean_n_bkg_pass :: sha256 53e1ec81343ac2c1 :: background jets passing the cut, mean over seeds
\newcommand{\ProbeRejVisibleContentLinearOneeighteight}{\ensuremath{>}538}  % experiments/FIGS/data/probe_ladder_v2_mlp2/analysis/seed_level_results.json :: table[visible_content,linear,188].rejection_points['0.90'].rejection :: sha256 53e1ec81343ac2c1 :: at 90% signal efficiency, mean +- SD over seeds; 5 of 5 seeds at the cap
\newcommand{\ProbeAucVisibleContentLinearOnesixtwo}{0.99921}  % experiments/FIGS/data/probe_ladder_v2_mlp2/analysis/seed_level_results.json :: levels.visible_content.linear[1].mean_auc :: sha256 53e1ec81343ac2c1 :: 5 seeds, 0 saturated at AUC=1
\newcommand{\ProbeAucSdVisibleContentLinearOnesixtwo}{0.00004}  % experiments/FIGS/data/probe_ladder_v2_mlp2/analysis/seed_level_results.json :: table[visible_content,linear,162].auc (SD over seeds, ddof=1) :: sha256 53e1ec81343ac2c1
\newcommand{\ProbeLogOneMinusAucVisibleContentLinearOnesixtwo}{\ensuremath{-}7.1480}  % experiments/FIGS/data/probe_ladder_v2_mlp2/analysis/seed_level_results.json :: levels.visible_content.linear[1].mean :: sha256 53e1ec81343ac2c1 :: natural log of 1-AUC, lower is better
\newcommand{\ProbeLogOneMinusAucSdVisibleContentLinearOnesixtwo}{0.0449}  % experiments/FIGS/data/probe_ladder_v2_mlp2/analysis/seed_level_results.json :: levels.visible_content.linear[1].seed_sd :: sha256 53e1ec81343ac2c1
\newcommand{\ProbeOmaVisibleContentLinearOnesixtwo}{\ensuremath{(7.87 \pm 0.35)\times10^{-4}}}  % experiments/FIGS/data/probe_ladder_v2_mlp2/analysis/seed_level_results.json :: table[visible_content,linear,162].auc :: sha256 53e1ec81343ac2c1 :: 1-AUC, mean +- SD over 5 seeds
\newcommand{\ProbeBkgLeftVisibleContentLinearOnesixtwo}{0.6}  % experiments/FIGS/data/probe_ladder_v2_mlp2/analysis/seed_level_results.json :: levels.visible_content.linear[1].rejection_points['0.90'].mean_n_bkg_pass :: sha256 53e1ec81343ac2c1 :: background jets passing the cut, mean over seeds
\newcommand{\ProbeRejVisibleContentLinearOnesixtwo}{\ensuremath{>}538}  % experiments/FIGS/data/probe_ladder_v2_mlp2/analysis/seed_level_results.json :: table[visible_content,linear,162].rejection_points['0.90'].rejection :: sha256 53e1ec81343ac2c1 :: at 90% signal efficiency, mean +- SD over seeds; 5 of 5 seeds at the cap
\newcommand{\ProbeAucVisibleContentLinearFourthree}{0.99918}  % experiments/FIGS/data/probe_ladder_v2_mlp2/analysis/seed_level_results.json :: levels.visible_content.linear[2].mean_auc :: sha256 53e1ec81343ac2c1 :: 5 seeds, 0 saturated at AUC=1
\newcommand{\ProbeAucSdVisibleContentLinearFourthree}{0.00009}  % experiments/FIGS/data/probe_ladder_v2_mlp2/analysis/seed_level_results.json :: table[visible_content,linear,43].auc (SD over seeds, ddof=1) :: sha256 53e1ec81343ac2c1
\newcommand{\ProbeLogOneMinusAucVisibleContentLinearFourthree}{\ensuremath{-}7.1124}  % experiments/FIGS/data/probe_ladder_v2_mlp2/analysis/seed_level_results.json :: levels.visible_content.linear[2].mean :: sha256 53e1ec81343ac2c1 :: natural log of 1-AUC, lower is better
\newcommand{\ProbeLogOneMinusAucSdVisibleContentLinearFourthree}{0.1137}  % experiments/FIGS/data/probe_ladder_v2_mlp2/analysis/seed_level_results.json :: levels.visible_content.linear[2].seed_sd :: sha256 53e1ec81343ac2c1
\newcommand{\ProbeOmaVisibleContentLinearFourthree}{\ensuremath{(8.2 \pm 1.0)\times10^{-4}}}  % experiments/FIGS/data/probe_ladder_v2_mlp2/analysis/seed_level_results.json :: table[visible_content,linear,43].auc :: sha256 53e1ec81343ac2c1 :: 1-AUC, mean +- SD over 5 seeds
\newcommand{\ProbeBkgLeftVisibleContentLinearFourthree}{0.4}  % experiments/FIGS/data/probe_ladder_v2_mlp2/analysis/seed_level_results.json :: levels.visible_content.linear[2].rejection_points['0.90'].mean_n_bkg_pass :: sha256 53e1ec81343ac2c1 :: background jets passing the cut, mean over seeds
\newcommand{\ProbeRejVisibleContentLinearFourthree}{\ensuremath{>}538}  % experiments/FIGS/data/probe_ladder_v2_mlp2/analysis/seed_level_results.json :: table[visible_content,linear,43].rejection_points['0.90'].rejection :: sha256 53e1ec81343ac2c1 :: at 90% signal efficiency, mean +- SD over seeds; 5 of 5 seeds at the cap
\newcommand{\ProbeAucVisibleContentLinearOneseven}{0.99768}  % experiments/FIGS/data/probe_ladder_v2_mlp2/analysis/seed_level_results.json :: levels.visible_content.linear[3].mean_auc :: sha256 53e1ec81343ac2c1 :: 5 seeds, 0 saturated at AUC=1
\newcommand{\ProbeAucSdVisibleContentLinearOneseven}{0.00020}  % experiments/FIGS/data/probe_ladder_v2_mlp2/analysis/seed_level_results.json :: table[visible_content,linear,17].auc (SD over seeds, ddof=1) :: sha256 53e1ec81343ac2c1
\newcommand{\ProbeLogOneMinusAucVisibleContentLinearOneseven}{\ensuremath{-}6.0684}  % experiments/FIGS/data/probe_ladder_v2_mlp2/analysis/seed_level_results.json :: levels.visible_content.linear[3].mean :: sha256 53e1ec81343ac2c1 :: natural log of 1-AUC, lower is better
\newcommand{\ProbeLogOneMinusAucSdVisibleContentLinearOneseven}{0.0889}  % experiments/FIGS/data/probe_ladder_v2_mlp2/analysis/seed_level_results.json :: levels.visible_content.linear[3].seed_sd :: sha256 53e1ec81343ac2c1
\newcommand{\ProbeOmaVisibleContentLinearOneseven}{\ensuremath{(2.32 \pm 0.20)\times10^{-3}}}  % experiments/FIGS/data/probe_ladder_v2_mlp2/analysis/seed_level_results.json :: table[visible_content,linear,17].auc :: sha256 53e1ec81343ac2c1 :: 1-AUC, mean +- SD over 5 seeds
\newcommand{\ProbeBkgLeftVisibleContentLinearOneseven}{3.4}  % experiments/FIGS/data/probe_ladder_v2_mlp2/analysis/seed_level_results.json :: levels.visible_content.linear[3].rejection_points['0.90'].mean_n_bkg_pass :: sha256 53e1ec81343ac2c1 :: background jets passing the cut, mean over seeds
\newcommand{\ProbeRejVisibleContentLinearOneseven}{900\,\ensuremath{\pm}\,400}  % experiments/FIGS/data/probe_ladder_v2_mlp2/analysis/seed_level_results.json :: table[visible_content,linear,17].rejection_points['0.90'].rejection :: sha256 53e1ec81343ac2c1 :: at 90% signal efficiency, mean +- SD over seeds; 0 of 5 seeds at the cap
\newcommand{\ProbeOmaRatioVisibleContentLinearOnesixtwoOverOneeighteight}{0.95}  % experiments/FIGS/data/probe_ladder_v2_mlp2/analysis/seed_level_results.json :: table[visible_content,linear,162/188].auc :: sha256 53e1ec81343ac2c1 :: ratio of the seed means of 1-AUC, coarser over finer
\newcommand{\ProbeOmaRatioVisibleContentLinearFourthreeOverOneeighteight}{0.99}  % experiments/FIGS/data/probe_ladder_v2_mlp2/analysis/seed_level_results.json :: table[visible_content,linear,43/188].auc :: sha256 53e1ec81343ac2c1 :: ratio of the seed means of 1-AUC, coarser over finer
\newcommand{\ProbeOmaRatioVisibleContentLinearOnesevenOverOneeighteight}{2.8}  % experiments/FIGS/data/probe_ladder_v2_mlp2/analysis/seed_level_results.json :: table[visible_content,linear,17/188].auc :: sha256 53e1ec81343ac2c1 :: ratio of the seed means of 1-AUC, coarser over finer
\newcommand{\ProbeOmaRatioVisibleContentLinearFourthreeOverOnesixtwo}{1.04}  % experiments/FIGS/data/probe_ladder_v2_mlp2/analysis/seed_level_results.json :: table[visible_content,linear,43/162].auc :: sha256 53e1ec81343ac2c1 :: ratio of the seed means of 1-AUC, coarser over finer
\newcommand{\ProbeOmaRatioVisibleContentLinearOnesevenOverOnesixtwo}{3.0}  % experiments/FIGS/data/probe_ladder_v2_mlp2/analysis/seed_level_results.json :: table[visible_content,linear,17/162].auc :: sha256 53e1ec81343ac2c1 :: ratio of the seed means of 1-AUC, coarser over finer
\newcommand{\ProbeOmaRatioVisibleContentLinearOnesevenOverFourthree}{2.8}  % experiments/FIGS/data/probe_ladder_v2_mlp2/analysis/seed_level_results.json :: table[visible_content,linear,17/43].auc :: sha256 53e1ec81343ac2c1 :: ratio of the seed means of 1-AUC, coarser over finer
\newcommand{\ProbeAucVisibleContentMlpOneeighteight}{0.99911}  % experiments/FIGS/data/probe_ladder_v2_mlp2/analysis/seed_level_results.json :: levels.visible_content.mlp[0].mean_auc :: sha256 53e1ec81343ac2c1 :: 5 seeds, 0 saturated at AUC=1
\newcommand{\ProbeAucSdVisibleContentMlpOneeighteight}{0.00005}  % experiments/FIGS/data/probe_ladder_v2_mlp2/analysis/seed_level_results.json :: table[visible_content,mlp,188].auc (SD over seeds, ddof=1) :: sha256 53e1ec81343ac2c1
\newcommand{\ProbeLogOneMinusAucVisibleContentMlpOneeighteight}{\ensuremath{-}7.0202}  % experiments/FIGS/data/probe_ladder_v2_mlp2/analysis/seed_level_results.json :: levels.visible_content.mlp[0].mean :: sha256 53e1ec81343ac2c1 :: natural log of 1-AUC, lower is better
\newcommand{\ProbeLogOneMinusAucSdVisibleContentMlpOneeighteight}{0.0570}  % experiments/FIGS/data/probe_ladder_v2_mlp2/analysis/seed_level_results.json :: levels.visible_content.mlp[0].seed_sd :: sha256 53e1ec81343ac2c1
\newcommand{\ProbeOmaVisibleContentMlpOneeighteight}{\ensuremath{(8.9 \pm 0.5)\times10^{-4}}}  % experiments/FIGS/data/probe_ladder_v2_mlp2/analysis/seed_level_results.json :: table[visible_content,mlp,188].auc :: sha256 53e1ec81343ac2c1 :: 1-AUC, mean +- SD over 5 seeds
\newcommand{\ProbeBkgLeftVisibleContentMlpOneeighteight}{0.4}  % experiments/FIGS/data/probe_ladder_v2_mlp2/analysis/seed_level_results.json :: levels.visible_content.mlp[0].rejection_points['0.90'].mean_n_bkg_pass :: sha256 53e1ec81343ac2c1 :: background jets passing the cut, mean over seeds
\newcommand{\ProbeRejVisibleContentMlpOneeighteight}{\ensuremath{>}538}  % experiments/FIGS/data/probe_ladder_v2_mlp2/analysis/seed_level_results.json :: table[visible_content,mlp,188].rejection_points['0.90'].rejection :: sha256 53e1ec81343ac2c1 :: at 90% signal efficiency, mean +- SD over seeds; 5 of 5 seeds at the cap
\newcommand{\ProbeAucVisibleContentMlpOnesixtwo}{0.99918}  % experiments/FIGS/data/probe_ladder_v2_mlp2/analysis/seed_level_results.json :: levels.visible_content.mlp[1].mean_auc :: sha256 53e1ec81343ac2c1 :: 5 seeds, 0 saturated at AUC=1
\newcommand{\ProbeAucSdVisibleContentMlpOnesixtwo}{0.00003}  % experiments/FIGS/data/probe_ladder_v2_mlp2/analysis/seed_level_results.json :: table[visible_content,mlp,162].auc (SD over seeds, ddof=1) :: sha256 53e1ec81343ac2c1
\newcommand{\ProbeLogOneMinusAucVisibleContentMlpOnesixtwo}{\ensuremath{-}7.1010}  % experiments/FIGS/data/probe_ladder_v2_mlp2/analysis/seed_level_results.json :: levels.visible_content.mlp[1].mean :: sha256 53e1ec81343ac2c1 :: natural log of 1-AUC, lower is better
\newcommand{\ProbeLogOneMinusAucSdVisibleContentMlpOnesixtwo}{0.0371}  % experiments/FIGS/data/probe_ladder_v2_mlp2/analysis/seed_level_results.json :: levels.visible_content.mlp[1].seed_sd :: sha256 53e1ec81343ac2c1
\newcommand{\ProbeOmaVisibleContentMlpOnesixtwo}{\ensuremath{(8.25 \pm 0.30)\times10^{-4}}}  % experiments/FIGS/data/probe_ladder_v2_mlp2/analysis/seed_level_results.json :: table[visible_content,mlp,162].auc :: sha256 53e1ec81343ac2c1 :: 1-AUC, mean +- SD over 5 seeds
\newcommand{\ProbeBkgLeftVisibleContentMlpOnesixtwo}{0.2}  % experiments/FIGS/data/probe_ladder_v2_mlp2/analysis/seed_level_results.json :: levels.visible_content.mlp[1].rejection_points['0.90'].mean_n_bkg_pass :: sha256 53e1ec81343ac2c1 :: background jets passing the cut, mean over seeds
\newcommand{\ProbeRejVisibleContentMlpOnesixtwo}{\ensuremath{>}538}  % experiments/FIGS/data/probe_ladder_v2_mlp2/analysis/seed_level_results.json :: table[visible_content,mlp,162].rejection_points['0.90'].rejection :: sha256 53e1ec81343ac2c1 :: at 90% signal efficiency, mean +- SD over seeds; 5 of 5 seeds at the cap
\newcommand{\ProbeAucVisibleContentMlpFourthree}{0.99917}  % experiments/FIGS/data/probe_ladder_v2_mlp2/analysis/seed_level_results.json :: levels.visible_content.mlp[2].mean_auc :: sha256 53e1ec81343ac2c1 :: 5 seeds, 0 saturated at AUC=1
\newcommand{\ProbeAucSdVisibleContentMlpFourthree}{0.00006}  % experiments/FIGS/data/probe_ladder_v2_mlp2/analysis/seed_level_results.json :: table[visible_content,mlp,43].auc (SD over seeds, ddof=1) :: sha256 53e1ec81343ac2c1
\newcommand{\ProbeLogOneMinusAucVisibleContentMlpFourthree}{\ensuremath{-}7.0957}  % experiments/FIGS/data/probe_ladder_v2_mlp2/analysis/seed_level_results.json :: levels.visible_content.mlp[2].mean :: sha256 53e1ec81343ac2c1 :: natural log of 1-AUC, lower is better
\newcommand{\ProbeLogOneMinusAucSdVisibleContentMlpFourthree}{0.0773}  % experiments/FIGS/data/probe_ladder_v2_mlp2/analysis/seed_level_results.json :: levels.visible_content.mlp[2].seed_sd :: sha256 53e1ec81343ac2c1
\newcommand{\ProbeOmaVisibleContentMlpFourthree}{\ensuremath{(8.3 \pm 0.6)\times10^{-4}}}  % experiments/FIGS/data/probe_ladder_v2_mlp2/analysis/seed_level_results.json :: table[visible_content,mlp,43].auc :: sha256 53e1ec81343ac2c1 :: 1-AUC, mean +- SD over 5 seeds
\newcommand{\ProbeBkgLeftVisibleContentMlpFourthree}{0.4}  % experiments/FIGS/data/probe_ladder_v2_mlp2/analysis/seed_level_results.json :: levels.visible_content.mlp[2].rejection_points['0.90'].mean_n_bkg_pass :: sha256 53e1ec81343ac2c1 :: background jets passing the cut, mean over seeds
\newcommand{\ProbeRejVisibleContentMlpFourthree}{\ensuremath{>}538}  % experiments/FIGS/data/probe_ladder_v2_mlp2/analysis/seed_level_results.json :: table[visible_content,mlp,43].rejection_points['0.90'].rejection :: sha256 53e1ec81343ac2c1 :: at 90% signal efficiency, mean +- SD over seeds; 5 of 5 seeds at the cap
\newcommand{\ProbeAucVisibleContentMlpOneseven}{0.99793}  % experiments/FIGS/data/probe_ladder_v2_mlp2/analysis/seed_level_results.json :: levels.visible_content.mlp[3].mean_auc :: sha256 53e1ec81343ac2c1 :: 5 seeds, 0 saturated at AUC=1
\newcommand{\ProbeAucSdVisibleContentMlpOneseven}{0.00020}  % experiments/FIGS/data/probe_ladder_v2_mlp2/analysis/seed_level_results.json :: table[visible_content,mlp,17].auc (SD over seeds, ddof=1) :: sha256 53e1ec81343ac2c1
\newcommand{\ProbeLogOneMinusAucVisibleContentMlpOneseven}{\ensuremath{-}6.1839}  % experiments/FIGS/data/probe_ladder_v2_mlp2/analysis/seed_level_results.json :: levels.visible_content.mlp[3].mean :: sha256 53e1ec81343ac2c1 :: natural log of 1-AUC, lower is better
\newcommand{\ProbeLogOneMinusAucSdVisibleContentMlpOneseven}{0.1014}  % experiments/FIGS/data/probe_ladder_v2_mlp2/analysis/seed_level_results.json :: levels.visible_content.mlp[3].seed_sd :: sha256 53e1ec81343ac2c1
\newcommand{\ProbeOmaVisibleContentMlpOneseven}{\ensuremath{(2.07 \pm 0.20)\times10^{-3}}}  % experiments/FIGS/data/probe_ladder_v2_mlp2/analysis/seed_level_results.json :: table[visible_content,mlp,17].auc :: sha256 53e1ec81343ac2c1 :: 1-AUC, mean +- SD over 5 seeds
\newcommand{\ProbeBkgLeftVisibleContentMlpOneseven}{3.4}  % experiments/FIGS/data/probe_ladder_v2_mlp2/analysis/seed_level_results.json :: levels.visible_content.mlp[3].rejection_points['0.90'].mean_n_bkg_pass :: sha256 53e1ec81343ac2c1 :: background jets passing the cut, mean over seeds
\newcommand{\ProbeRejVisibleContentMlpOneseven}{\ensuremath{\geq}638\ensuremath{^{\ast}}}  % experiments/FIGS/data/probe_ladder_v2_mlp2/analysis/seed_level_results.json :: table[visible_content,mlp,17].rejection_points['0.90'].rejection :: sha256 53e1ec81343ac2c1 :: at 90% signal efficiency, mean +- SD over seeds; 1 of 5 seeds at the cap
\newcommand{\ProbeOmaRatioVisibleContentMlpOnesixtwoOverOneeighteight}{0.92}  % experiments/FIGS/data/probe_ladder_v2_mlp2/analysis/seed_level_results.json :: table[visible_content,mlp,162/188].auc :: sha256 53e1ec81343ac2c1 :: ratio of the seed means of 1-AUC, coarser over finer
\newcommand{\ProbeOmaRatioVisibleContentMlpFourthreeOverOneeighteight}{0.93}  % experiments/FIGS/data/probe_ladder_v2_mlp2/analysis/seed_level_results.json :: table[visible_content,mlp,43/188].auc :: sha256 53e1ec81343ac2c1 :: ratio of the seed means of 1-AUC, coarser over finer
\newcommand{\ProbeOmaRatioVisibleContentMlpOnesevenOverOneeighteight}{2.3}  % experiments/FIGS/data/probe_ladder_v2_mlp2/analysis/seed_level_results.json :: table[visible_content,mlp,17/188].auc :: sha256 53e1ec81343ac2c1 :: ratio of the seed means of 1-AUC, coarser over finer
\newcommand{\ProbeOmaRatioVisibleContentMlpFourthreeOverOnesixtwo}{1.01}  % experiments/FIGS/data/probe_ladder_v2_mlp2/analysis/seed_level_results.json :: table[visible_content,mlp,43/162].auc :: sha256 53e1ec81343ac2c1 :: ratio of the seed means of 1-AUC, coarser over finer
\newcommand{\ProbeOmaRatioVisibleContentMlpOnesevenOverOnesixtwo}{2.5}  % experiments/FIGS/data/probe_ladder_v2_mlp2/analysis/seed_level_results.json :: table[visible_content,mlp,17/162].auc :: sha256 53e1ec81343ac2c1 :: ratio of the seed means of 1-AUC, coarser over finer
\newcommand{\ProbeOmaRatioVisibleContentMlpOnesevenOverFourthree}{2.5}  % experiments/FIGS/data/probe_ladder_v2_mlp2/analysis/seed_level_results.json :: table[visible_content,mlp,17/43].auc :: sha256 53e1ec81343ac2c1 :: ratio of the seed means of 1-AUC, coarser over finer
\newcommand{\VocabSizeLoneeighteight}{188}  % configs/labelmaps/rung_label_maps.v1.csv :: column L188 (distinct group count) :: sha256 412f9aa3b58db412
\newcommand{\VocabSizeLonesixtwo}{162}  % configs/labelmaps/rung_label_maps.v1.csv :: column L162 (distinct group count) :: sha256 412f9aa3b58db412
\newcommand{\VocabSizeRsixthreeQone}{64}  % configs/labelmaps/rung_label_maps.v1.csv :: column R63_Q1 (distinct group count) :: sha256 412f9aa3b58db412
\newcommand{\VocabSizeRfourtwoQone}{43}  % configs/labelmaps/rung_label_maps.v1.csv :: column R42_Q1 (distinct group count) :: sha256 412f9aa3b58db412
\newcommand{\VocabSizeRtwonineQone}{30}  % configs/labelmaps/rung_label_maps.v1.csv :: column R29_Q1 (distinct group count) :: sha256 412f9aa3b58db412
\newcommand{\VocabSizeRonesixQone}{17}  % configs/labelmaps/rung_label_maps.v1.csv :: column R16_Q1 (distinct group count) :: sha256 412f9aa3b58db412
\newcommand{\VocabSizeRthreeVIS}{4}  % configs/labelmaps/rung_label_maps.v1.csv :: column R3_VIS (distinct group count) :: sha256 412f9aa3b58db412
\newcommand{\VocabSizeRoneQone}{2}  % configs/labelmaps/rung_label_maps.v1.csv :: column R1_Q1 (distinct group count) :: sha256 412f9aa3b58db412
\newcommand{\ProbeSplitTrain}{60\%}  % experiments/EVAL/probe.py :: SPLIT_FRACTIONS :: sha256 485db0966e660588 :: share of the probe sample; the same jets for every model
\newcommand{\ProbeSplitVal}{20\%}  % experiments/EVAL/probe.py :: SPLIT_FRACTIONS :: sha256 485db0966e660588 :: share of the probe sample; the same jets for every model
\newcommand{\ProbeSplitTest}{20\%}  % experiments/EVAL/probe.py :: SPLIT_FRACTIONS :: sha256 485db0966e660588 :: share of the probe sample; the same jets for every model
\newcommand{\ProbeMlpHidden}{256}  % experiments/EVAL/probe.py :: _fit_mlp :: sha256 485db0966e660588 :: hidden width of the MLP probe
\newcommand{\ProbeMlpInits}{3}  % experiments/EVAL/probe.py :: MLP_SEEDS :: sha256 485db0966e660588 :: MLP initialisations whose scores are averaged per cell
\newcommand{\ProbeMlpLrDrop}{10}  % experiments/EVAL/probe.py :: MLP_SCHEDULE['lr_factor'] :: sha256 485db0966e660588 :: factor the MLP's learning rate drops by on a plateau
\newcommand{\ProbeMlpPlateaus}{three}  % experiments/EVAL/probe.py :: MLP_SCHEDULE: stop_patience / (lr_patience + 1) :: sha256 485db0966e660588 :: patience windows without gain after which the MLP stops
\newcommand{\ProbeCMin}{\ensuremath{10^{-2}}}  % experiments/EVAL/probe.py :: C_GRID :: sha256 485db0966e660588 :: end of the logistic-regression inverse-regularisation grid
\newcommand{\ProbeCMax}{\ensuremath{10^{2}}}  % experiments/EVAL/probe.py :: C_GRID :: sha256 485db0966e660588 :: end of the logistic-regression inverse-regularisation grid
\newcommand{\ProbeNBkgTestBvcFourprong}{3{,}158}  % experiments/FIGS/data/probe_ladder_v2_mlp2/s1.json :: tasks.bvc_4prong.n_background_test :: sha256 f7e1de0e43be89b8 :: test-sample background jets
\newcommand{\ProbeNBkgTestBvcQcd}{5{,}371}  % experiments/FIGS/data/probe_ladder_v2_mlp2/s1.json :: tasks.bvc_qcd.n_background_test :: sha256 f7e1de0e43be89b8 :: test-sample background jets
\newcommand{\ProbeNBkgTestBvcResonant}{11{,}876}  % experiments/FIGS/data/probe_ladder_v2_mlp2/s1.json :: tasks.bvc_resonant.n_background_test :: sha256 f7e1de0e43be89b8 :: test-sample background jets
\newcommand{\ProbeNBkgTestEeVsMm}{4{,}708}  % experiments/FIGS/data/probe_ladder_v2_mlp2/s1.json :: tasks.ee_vs_mm.n_background_test :: sha256 f7e1de0e43be89b8 :: test-sample background jets
\newcommand{\ProbeNBkgTestRetainedTopology}{1{,}078}  % experiments/FIGS/data/probe_ladder_v2_mlp2/s1.json :: tasks.retained_topology.n_background_test :: sha256 f7e1de0e43be89b8 :: test-sample background jets
\newcommand{\ProbeNBkgTestVisibleContent}{2{,}554}  % experiments/FIGS/data/probe_ladder_v2_mlp2/s1.json :: tasks.visible_content.n_background_test :: sha256 f7e1de0e43be89b8 :: test-sample background jets
\newcommand{\UseDiesAtSophonEqfour}{43}  % configs/labelmaps/usecase_survival.v1.json :: sophon_eq4.dies_at :: sha256 122275213c349593 :: first vocabulary size at which it is not constructible
\newcommand{\UseLastAliveSophonEqfour}{64}  % configs/labelmaps/usecase_survival.v1.json :: sophon_eq4.last_rung_alive :: sha256 122275213c349593
\newcommand{\UseDiesAtSophonEqsixAone}{64}  % configs/labelmaps/usecase_survival.v1.json :: sophon_eq6_a1.dies_at :: sha256 122275213c349593 :: first vocabulary size at which it is not constructible
\newcommand{\UseLastAliveSophonEqsixAone}{162}  % configs/labelmaps/usecase_survival.v1.json :: sophon_eq6_a1.last_rung_alive :: sha256 122275213c349593
\newcommand{\UseDiesAtSophonEqsixAthree}{64}  % configs/labelmaps/usecase_survival.v1.json :: sophon_eq6_a3.dies_at :: sha256 122275213c349593 :: first vocabulary size at which it is not constructible
\newcommand{\UseLastAliveSophonEqsixAthree}{162}  % configs/labelmaps/usecase_survival.v1.json :: sophon_eq6_a3.last_rung_alive :: sha256 122275213c349593
\newcommand{\UseDiesAtSophonEqsevenNine}{64}  % configs/labelmaps/usecase_survival.v1.json :: sophon_eq7_9.dies_at :: sha256 122275213c349593 :: first vocabulary size at which it is not constructible
\newcommand{\UseLastAliveSophonEqsevenNine}{162}  % configs/labelmaps/usecase_survival.v1.json :: sophon_eq7_9.last_rung_alive :: sha256 122275213c349593
\newcommand{\UseDiesAtVcbEqone}{64}  % configs/labelmaps/usecase_survival.v1.json :: vcb_eq1.dies_at :: sha256 122275213c349593 :: first vocabulary size at which it is not constructible
\newcommand{\UseLastAliveVcbEqone}{162}  % configs/labelmaps/usecase_survival.v1.json :: vcb_eq1.last_rung_alive :: sha256 122275213c349593
\newcommand{\UseDiesAtGnthreexQcdSubclass}{162}  % configs/labelmaps/usecase_survival.v1.json :: gn3x_qcd_subclass.dies_at :: sha256 122275213c349593 :: first vocabulary size at which it is not constructible
\newcommand{\UseLastAliveGnthreexQcdSubclass}{188}  % configs/labelmaps/usecase_survival.v1.json :: gn3x_qcd_subclass.last_rung_alive :: sha256 122275213c349593
\newcommand{\UseWlikeLastAlive}{162}  % configs/labelmaps/rung_label_maps.v1.csv :: groups at each level vs {label_X_cs, label_X_qq} and that plus the QCD classes :: sha256 412f9aa3b58db412 :: last level at which Sophon's W-like score g_W(2) is constructible
\newcommand{\UseWlikeDiesAt}{64}  % configs/labelmaps/rung_label_maps.v1.csv :: groups at each level vs {label_X_cs, label_X_qq} and that plus the QCD classes :: sha256 412f9aa3b58db412 :: first level at which it is not
\newcommand{\DesignEpochs}{80}  % experiments/MTX/k8s/job-mtx-l188-s1-raunav.yaml :: --num-epochs :: sha256 96fef4b6337aaa27
\newcommand{\DesignSamplesPerEpoch}{10{,}240{,}000}  % experiments/MTX/k8s/job-mtx-l188-s1-raunav.yaml :: --samples-per-epoch :: sha256 96fef4b6337aaa27
\newcommand{\DesignExamplesSeen}{\ensuremath{8.192\times10^{8}}}  % experiments/MTX/k8s/job-mtx-l188-s1-raunav.yaml :: --num-epochs x --samples-per-epoch :: sha256 96fef4b6337aaa27 :: training examples drawn over the run
\newcommand{\DesignPassesSelected}{7.45}  % experiments/FIGS/data/v2/jc2_train/counts.json :: n_selected :: sha256 b0b5a8b63aa808c0 :: --num-epochs x --samples-per-epoch of job-mtx-l188-s1-raunav.yaml / n_selected: passes over the selected training jets
\newcommand{\DesignBatchSize}{512}  % experiments/MTX/k8s/job-mtx-l188-s1-raunav.yaml :: --batch-size :: sha256 96fef4b6337aaa27
\newcommand{\DesignStartLr}{\ensuremath{5\times10^{-4}}}  % experiments/MTX/k8s/job-mtx-l188-s1-raunav.yaml :: --start-lr :: sha256 96fef4b6337aaa27
\newcommand{\DesignHeadHidden}{512}  % experiments/MTX/k8s/job-mtx-l188-s1-raunav.yaml :: -o fc_params (hidden width) :: sha256 96fef4b6337aaa27
\newcommand{\DesignParticleBlocks}{8}  % experiments/E1/ParT_sophon_arch_10c.py :: num_layers=(\d+) :: sha256 5793faa7d55a8b0d
\newcommand{\DesignClassBlocks}{2}  % experiments/E1/ParT_sophon_arch_10c.py :: num_cls_layers=(\d+) :: sha256 5793faa7d55a8b0d
\newcommand{\DesignHeads}{8}  % experiments/E1/ParT_sophon_arch_10c.py :: num_heads=(\d+) :: sha256 5793faa7d55a8b0d
\newcommand{\DesignEmbedDim}{128}  % experiments/E1/ParT_sophon_arch_10c.py :: \bembed_dims=\[\d+, \d+, (\d+)\] :: sha256 5793faa7d55a8b0d
\newcommand{\DesignMpmMaskRate}{40\%}  % experiments/MTX/mpm.py :: DEFAULT_MASK_RATE :: sha256 49f0c8c6666d0211 :: share of each jet's particles hidden from the encoder
\newcommand{\DesignPtMin}{200}  % configs/arms/L188.yaml :: selection :: sha256 d5a188dc80885515 :: GeV
\newcommand{\DesignPtMax}{2{,}500}  % configs/arms/L188.yaml :: selection :: sha256 d5a188dc80885515 :: GeV
\newcommand{\DesignMsdMin}{20}  % configs/arms/L188.yaml :: selection :: sha256 d5a188dc80885515 :: GeV
\newcommand{\DesignMsdMax}{500}  % configs/arms/L188.yaml :: selection :: sha256 d5a188dc80885515 :: GeV
\newcommand{\LitJetClassTwoNJetsMillion}{139}  % experiments/FIGS/data/literature_facts.json :: facts[0].value :: sha256 713dba2a5e6e1dd1 :: arXiv:2405.12972v1 App. A, p. 11
\newcommand{\LitJetClassTwoNResonantClasses}{161}  % experiments/FIGS/data/literature_facts.json :: facts[1].value :: sha256 713dba2a5e6e1dd1 :: arXiv:2405.12972v1 App. A, p. 11
\newcommand{\LitJetClassTwoNQcdClasses}{27}  % experiments/FIGS/data/literature_facts.json :: facts[2].value :: sha256 713dba2a5e6e1dd1 :: arXiv:2405.12972v1 App. A, p. 11
\newcommand{\LitJetClassTwoMaxResonanceMass}{500}  % experiments/FIGS/data/literature_facts.json :: facts[3].value :: sha256 713dba2a5e6e1dd1 :: arXiv:2405.12972v1 Sec. II, p. 2
\newcommand{\LitMadgraphVersion}{2.9.18}  % experiments/FIGS/data/literature_facts.json :: facts[4].value :: sha256 713dba2a5e6e1dd1 :: arXiv:2405.12972v1 App. A, p. 11
\newcommand{\LitPythiaVersion}{8.3}  % experiments/FIGS/data/literature_facts.json :: facts[5].value :: sha256 713dba2a5e6e1dd1 :: arXiv:2405.12972v1 App. A, p. 11
\newcommand{\LitDelphesVersion}{3}  % experiments/FIGS/data/literature_facts.json :: facts[6].value :: sha256 713dba2a5e6e1dd1 :: arXiv:2405.12972v1 App. A, p. 11
\newcommand{\LitJetClassTwoPileupMean}{50}  % experiments/FIGS/data/literature_facts.json :: facts[7].value :: sha256 713dba2a5e6e1dd1 :: arXiv:2405.12972v1 Sec. II, p. 2
\newcommand{\LitJetClassTwoJetRadius}{0.8}  % experiments/FIGS/data/literature_facts.json :: facts[8].value :: sha256 713dba2a5e6e1dd1 :: arXiv:2405.12972v1 Sec. II, p. 2
\newcommand{\LitAtlasUnifiedEffDropPercent}{10}  % experiments/FIGS/data/literature_facts.json :: facts[9].value :: sha256 713dba2a5e6e1dd1 :: ATL-PHYS-PUB-2026-013 Sec. 6.5, p. 22
\newcommand{\LitSophonPaperParticleBlocks}{6}  % experiments/FIGS/data/literature_facts.json :: facts[10].value :: sha256 713dba2a5e6e1dd1 :: arXiv:2405.12972v1 App. B.1, p. 12
\newcommand{\LitSophonSamplesPerEpochMillion}{10}  % experiments/FIGS/data/literature_facts.json :: facts[11].value :: sha256 713dba2a5e6e1dd1 :: arXiv:2405.12972v1 App. B.1, p. 12
\newcommand{\LitSophonSubsetSamplesPerEpochMillion}{2.5}  % experiments/FIGS/data/literature_facts.json :: facts[12].value :: sha256 713dba2a5e6e1dd1 :: arXiv:2405.12972v1 App. B.2, p. 12
\newcommand{\LitGlopartNClasses}{314}  % experiments/FIGS/data/literature_facts.json :: facts[13].value :: sha256 713dba2a5e6e1dd1 :: C. Li, PhD thesis, CERN-THESIS-2024-281 Sec. 11.2.1, p. 188
\newcommand{\LitPartFtEpochs}{20}  % experiments/FIGS/data/literature_facts.json :: facts[14].value :: sha256 713dba2a5e6e1dd1 :: arXiv:2202.03772 Sec. 5.2, p. 8
\newcommand{\LitMehtaNCategories}{eight}  % experiments/FIGS/data/literature_facts.json :: facts[15].value :: sha256 713dba2a5e6e1dd1 :: arXiv:2605.05556v2 Abstract
\newcommand{\LitHongPretrainRuns}{2}  % experiments/FIGS/data/literature_facts.json :: facts[16].value :: sha256 713dba2a5e6e1dd1 :: arXiv:2410.23129 App. A.1.1, p. 18
\newcommand{\LitSophonLookaheadK}{6}  % experiments/FIGS/data/literature_facts.json :: facts[17].value :: sha256 713dba2a5e6e1dd1 :: arXiv:2405.12972v1 App. B.1, p. 12
\newcommand{\LitSophonLookaheadAlpha}{0.5}  % experiments/FIGS/data/literature_facts.json :: facts[18].value :: sha256 713dba2a5e6e1dd1 :: arXiv:2405.12972v1 App. B.1, p. 12
\newcommand{\LitSophonRadamBetaOne}{0.95}  % experiments/FIGS/data/literature_facts.json :: facts[19].value :: sha256 713dba2a5e6e1dd1 :: arXiv:2405.12972v1 App. B.1, p. 12
\newcommand{\LitSophonRadamBetaTwo}{0.999}  % experiments/FIGS/data/literature_facts.json :: facts[20].value :: sha256 713dba2a5e6e1dd1 :: arXiv:2405.12972v1 App. B.1, p. 12
\newcommand{\LitSophonRadamEps}{\ensuremath{10^{-5}}}  % experiments/FIGS/data/literature_facts.json :: facts[21].value :: sha256 713dba2a5e6e1dd1 :: arXiv:2405.12972v1 App. B.1, p. 12
\newcommand{\LitGnTwoAuxRejectionDropPercent}{30}  % experiments/FIGS/data/literature_facts.json :: facts[22].value :: sha256 713dba2a5e6e1dd1 :: arXiv:2505.19689v2 (Nature Commun. 17 (2026) 541) p. 9
\newcommand{\LitCmsPartNClasses}{37}  % experiments/FIGS/data/literature_facts.json :: facts[23].value :: sha256 713dba2a5e6e1dd1 :: arXiv:2604.09809v1 (CMS-JME-25-001) Sec. 5.2, p. 7
\newcommand{\LitCmsPartMassLambda}{0.05}  % experiments/FIGS/data/literature_facts.json :: facts[24].value :: sha256 713dba2a5e6e1dd1 :: arXiv:2604.09809v1 (CMS-JME-25-001) Sec. 5.2, p. 7
\newcommand{\LitGlopartVthreeNClasses}{374}  % experiments/FIGS/data/literature_facts.json :: facts[25].value :: sha256 713dba2a5e6e1dd1 :: CMS-DP-2026-104 (CMS public page) Architecture and outputs
\newcommand{\LitHuhCoarseClasses}{127}  % experiments/FIGS/data/literature_facts.json :: facts[26].value :: sha256 713dba2a5e6e1dd1 :: arXiv:1608.08614v2 Sec. 5.3, p. 5
\newcommand{\LitHuhKnnGapPercent}{15}  % experiments/FIGS/data/literature_facts.json :: facts[27].value :: sha256 713dba2a5e6e1dd1 :: arXiv:1608.08614v2 Sec. 5.3, p. 5
\newcommand{\LitRsthreelRuns}{3}  % experiments/FIGS/data/literature_facts.json :: facts[28].value :: sha256 713dba2a5e6e1dd1 :: arXiv:2403.07066v2 (Phys. Rev. D 111 (2025) 032010) Sec. IV, p. 5
\newcommand{\LitSophonParamsMillion}{2.3}  % experiments/FIGS/data/literature_facts.json :: facts[29].value :: sha256 713dba2a5e6e1dd1 :: arXiv:2405.12972v1 App. B.1, p. 12
\newcommand{\LitGlopartVoneNClasses}{37}  % experiments/FIGS/data/literature_facts.json :: facts[30].value :: sha256 713dba2a5e6e1dd1 :: CMS-DP-2026-104 (CMS public page) Model generations
\newcommand{\LitSophonXMass}{200}  % experiments/FIGS/data/literature_facts.json :: facts[31].value :: sha256 713dba2a5e6e1dd1 :: arXiv:2405.12972v1 Sec. III, p. 3
\newcommand{\LitSophonXLumi}{40}  % experiments/FIGS/data/literature_facts.json :: facts[32].value :: sha256 713dba2a5e6e1dd1 :: arXiv:2405.12972v1 Sec. III, Fig. 2
\newcommand{\LitSophonXWindowLo}{150}  % experiments/FIGS/data/literature_facts.json :: facts[33].value :: sha256 713dba2a5e6e1dd1 :: arXiv:2405.12972v1 Sec. III, Fig. 2
\newcommand{\LitSophonXWindowHi}{230}  % experiments/FIGS/data/literature_facts.json :: facts[34].value :: sha256 713dba2a5e6e1dd1 :: arXiv:2405.12972v1 Sec. III, Fig. 2
\newcommand{\LitSophonAblationNClasses}{42}  % experiments/FIGS/data/literature_facts.json :: facts[35].value :: sha256 713dba2a5e6e1dd1 :: arXiv:2405.12972v1 Sec. III
\newcommand{\LitSophonTransferSamplesMillion}{2.5}  % experiments/FIGS/data/literature_facts.json :: facts[36].value :: sha256 713dba2a5e6e1dd1 :: arXiv:2405.12972v1 App. B.5, p. 13
\newcommand{\LitSophonTransferBatch}{1024}  % experiments/FIGS/data/literature_facts.json :: facts[37].value :: sha256 713dba2a5e6e1dd1 :: arXiv:2405.12972v1 App. B.5, p. 13
\newcommand{\LitSophonScratchEpochs}{50}  % experiments/FIGS/data/literature_facts.json :: facts[38].value :: sha256 713dba2a5e6e1dd1 :: arXiv:2405.12972v1 App. B.3
\newcommand{\LitHuhFineClasses}{1000}  % experiments/FIGS/data/literature_facts.json :: facts[39].value :: sha256 713dba2a5e6e1dd1 :: arXiv:1608.08614v2 Sec. 1, p. 1
\newcommand{\LitHuhPascalDropMap}{2.8}  % experiments/FIGS/data/literature_facts.json :: facts[40].value :: sha256 713dba2a5e6e1dd1 :: arXiv:1608.08614v2 Sec. 1, p. 1
\newcommand{\LitHongSuperclasses}{11}  % experiments/FIGS/data/literature_facts.json :: facts[41].value :: sha256 713dba2a5e6e1dd1 :: arXiv:2410.23129 Sec. 6.2
\newcommand{\LitHongDeclineClasses}{\ensuremath{10^{3}}}  % experiments/FIGS/data/literature_facts.json :: facts[42].value :: sha256 713dba2a5e6e1dd1 :: arXiv:2410.23129 Sec. 6.2
\newcommand{\LitObjectivesFullJets}{\ensuremath{1.2\times10^{6}}}  % experiments/FIGS/data/literature_facts.json :: facts[43].value :: sha256 713dba2a5e6e1dd1 :: arXiv:2606.14870 Sec. IV A
\newcommand{\LitObjectivesFewJets}{\ensuremath{10^{4}}}  % experiments/FIGS/data/literature_facts.json :: facts[44].value :: sha256 713dba2a5e6e1dd1 :: arXiv:2606.14870 Sec. IV A
\newcommand{\LitMpmBenefitJets}{10k}  % experiments/FIGS/data/literature_facts.json :: facts[45].value :: sha256 713dba2a5e6e1dd1 :: arXiv:2401.13537 Sec. V
\newcommand{\LitPartTopRejScratch}{1602}  % experiments/FIGS/data/literature_facts.json :: facts[46].value :: sha256 713dba2a5e6e1dd1 :: arXiv:2202.03772 Table 5
\newcommand{\LitPartTopRejFt}{2766}  % experiments/FIGS/data/literature_facts.json :: facts[47].value :: sha256 713dba2a5e6e1dd1 :: arXiv:2202.03772 Table 5
\newcommand{\LitMahajanLinearGapPercent}{7}  % experiments/FIGS/data/literature_facts.json :: facts[48].value :: sha256 713dba2a5e6e1dd1 :: arXiv:1805.00932 Sec. 3.1.2
\newcommand{\LitBukchinCoarse}{33.83}  % experiments/FIGS/data/literature_facts.json :: facts[49].value :: sha256 713dba2a5e6e1dd1 :: arXiv:2012.03515 Table 2, LIVING-17 all-way
\newcommand{\LitBukchinFine}{58.95}  % experiments/FIGS/data/literature_facts.json :: facts[50].value :: sha256 713dba2a5e6e1dd1 :: arXiv:2012.03515 Table 2, LIVING-17 all-way
\newcommand{\LitKornblithRho}{\ensuremath{-0.93}}  % experiments/FIGS/data/literature_facts.json :: facts[51].value :: sha256 713dba2a5e6e1dd1 :: arXiv:2010.16402 Sec. 3.4
\newcommand{\LitPartTopSigEffPercent}{30}  % experiments/FIGS/data/literature_facts.json :: facts[52].value :: sha256 713dba2a5e6e1dd1 :: arXiv:2202.03772 Table 5 (column Rej30%)
\newcommand{\LitPartTopKinFeatures}{7}  % experiments/FIGS/data/literature_facts.json :: facts[53].value :: sha256 713dba2a5e6e1dd1 :: arXiv:2202.03772 Sec. 5.2
\newcommand{\LitBukchinAncor}{45.14}  % experiments/FIGS/data/literature_facts.json :: facts[54].value :: sha256 713dba2a5e6e1dd1 :: arXiv:2012.03515 Table 2, LIVING-17 all-way
\newcommand{\LitMahajanFtGapPp}{0.1}  % experiments/FIGS/data/literature_facts.json :: facts[55].value :: sha256 713dba2a5e6e1dd1 :: arXiv:1805.00932 Fig. 1, ImageNet-9k target
\newcommand{\LitSophonTransferLr}{5 × 10−4}  % experiments/FIGS/data/literature_facts.json :: facts[56].value :: sha256 713dba2a5e6e1dd1 :: arXiv:2405.12972v1 App. B.5
\newcommand{\LitQgJetRadius}{0.4}  % experiments/FIGS/data/literature_facts.json :: facts[57].value :: sha256 713dba2a5e6e1dd1 :: arXiv:1810.05165 Sec. 4.1
\newcommand{\LitQgPtMin}{500}  % experiments/FIGS/data/literature_facts.json :: facts[58].value :: sha256 713dba2a5e6e1dd1 :: arXiv:1810.05165 Sec. 4.1
\newcommand{\LitQgPtMax}{550}  % experiments/FIGS/data/literature_facts.json :: facts[59].value :: sha256 713dba2a5e6e1dd1 :: arXiv:1810.05165 Sec. 4.1
\newcommand{\LitPartTopAucScratch}{0.9858}  % experiments/FIGS/data/literature_facts.json :: facts[60].value :: sha256 713dba2a5e6e1dd1 :: arXiv:2202.03772 Table 5
\newcommand{\LitPartTopAucFt}{0.9877}  % experiments/FIGS/data/literature_facts.json :: facts[61].value :: sha256 713dba2a5e6e1dd1 :: arXiv:2202.03772 Table 5
\newcommand{\LitPartQgAucScratch}{0.9121}  % experiments/FIGS/data/literature_facts.json :: facts[62].value :: sha256 713dba2a5e6e1dd1 :: arXiv:2202.03772 Table 6 (exp)
\newcommand{\LitPartQgRejScratch}{101.2}  % experiments/FIGS/data/literature_facts.json :: facts[63].value :: sha256 713dba2a5e6e1dd1 :: arXiv:2202.03772 Table 6 (exp)
\newcommand{\LitPartQgAucFt}{0.9151}  % experiments/FIGS/data/literature_facts.json :: facts[64].value :: sha256 713dba2a5e6e1dd1 :: arXiv:2202.03772 Table 6 (exp)
\newcommand{\LitPartQgRejFt}{107.9}  % experiments/FIGS/data/literature_facts.json :: facts[65].value :: sha256 713dba2a5e6e1dd1 :: arXiv:2202.03772 Table 6 (exp)
\newcommand{\LitPartJcTestMillion}{20}  % experiments/FIGS/data/literature_facts.json :: facts[66].value :: sha256 713dba2a5e6e1dd1 :: arXiv:2202.03772 Sec. 3
\newcommand{\LitPartJcTwoMAcc}{0.836}  % experiments/FIGS/data/literature_facts.json :: facts[67].value :: sha256 713dba2a5e6e1dd1 :: arXiv:2202.03772 Table 3
\newcommand{\LitPartSmallSetIters}{100 k}  % experiments/FIGS/data/literature_facts.json :: facts[68].value :: sha256 713dba2a5e6e1dd1 :: arXiv:2202.03772 Sec. 5.1 (training dataset size)
\newcommand{\LitPartSmallSetPercent}{2}  % experiments/FIGS/data/literature_facts.json :: facts[69].value :: sha256 713dba2a5e6e1dd1 :: arXiv:2202.03772 Sec. 5.1 (training dataset size)
\newcommand{\LitSophonEtaMax}{2.5}  % experiments/FIGS/data/literature_facts.json :: facts[70].value :: sha256 713dba2a5e6e1dd1 :: arXiv:2405.12972 App. A
\newcommand{\LitBHadronLifetimePs}{1.5}  % experiments/FIGS/data/literature_facts.json :: facts[71].value :: sha256 713dba2a5e6e1dd1 :: arXiv:1712.07158 Sec. 4.1
\newcommand{\LitCHadronLifetimePs}{1}  % experiments/FIGS/data/literature_facts.json :: facts[72].value :: sha256 713dba2a5e6e1dd1 :: arXiv:1712.07158 Sec. 4.1
\newcommand{\LitOmnilearnPidDropProb}{0.2}  % experiments/FIGS/data/literature_facts.json :: facts[73].value :: sha256 713dba2a5e6e1dd1 :: arXiv:2502.14652 Sec. II.A
\newcommand{\LitMahajanFtNineKLarge}{48.5}  % experiments/FIGS/data/literature_facts.json :: facts[74].value :: sha256 713dba2a5e6e1dd1 :: arXiv:1805.00932 Fig. 1 (IG-1B-17k, IN-9k, full fine-tuning)
\newcommand{\LitMahajanFtNineKSmall}{48.4}  % experiments/FIGS/data/literature_facts.json :: facts[75].value :: sha256 713dba2a5e6e1dd1 :: arXiv:1805.00932 Fig. 1 (IG-940M-1.5k, IN-9k, full fine-tuning)
\newcommand{\LitMahajanFtOneKSmall}{84.2}  % experiments/FIGS/data/literature_facts.json :: facts[76].value :: sha256 713dba2a5e6e1dd1 :: arXiv:1805.00932 Fig. 1 (IG-940M-1.5k, IN-1k, full fine-tuning)
\newcommand{\LitMahajanFtOneKLarge}{83.6}  % experiments/FIGS/data/literature_facts.json :: facts[77].value :: sha256 713dba2a5e6e1dd1 :: arXiv:1805.00932 Fig. 1 (IG-1B-17k, IN-1k, full fine-tuning)
\newcommand{\LitChenGainFifteen}{0.16}  % experiments/FIGS/data/literature_facts.json :: facts[78].value :: sha256 713dba2a5e6e1dd1 :: arXiv:1901.07012 Table VIII (15 coarse classes, CIFAR-100)
\newcommand{\LitChenGainTwenty}{0.01}  % experiments/FIGS/data/literature_facts.json :: facts[79].value :: sha256 713dba2a5e6e1dd1 :: arXiv:1901.07012 Table VIII (20 coarse classes, CIFAR-100)
\newcommand{\LitChenCoarseFifteen}{15}  % experiments/FIGS/data/literature_facts.json :: facts[80].value :: sha256 713dba2a5e6e1dd1 :: arXiv:1901.07012 Table VIII
\newcommand{\LitChenCoarseTwenty}{20}  % experiments/FIGS/data/literature_facts.json :: facts[81].value :: sha256 713dba2a5e6e1dd1 :: arXiv:1901.07012 Table VIII
\newcommand{\LitChenRandomMaxGain}{2.78}  % experiments/FIGS/data/literature_facts.json :: facts[82].value :: sha256 713dba2a5e6e1dd1 :: arXiv:1901.07012 Sec. VI-A, Table VI
\newcommand{\LitGlopartMassLambda}{5}  % experiments/FIGS/data/literature_facts.json :: facts[83].value :: sha256 713dba2a5e6e1dd1 :: CERN-THESIS-2024-281 Eq. (11.7), Sec. 11.2.1
\newcommand{\LitGlopartAddedAucBase}{0.9555}  % experiments/FIGS/data/literature_facts.json :: facts[84].value :: sha256 713dba2a5e6e1dd1 :: CERN-THESIS-2024-281 Sec. 11.4.1, Fig. 11.19 (H->WW 4q vs QCD)
\newcommand{\LitGlopartAddedAucTop}{0.9593}  % experiments/FIGS/data/literature_facts.json :: facts[85].value :: sha256 713dba2a5e6e1dd1 :: CERN-THESIS-2024-281 Sec. 11.4.1, Fig. 11.19 (H->WW 4q vs QCD)
\newcommand{\LitCmsSdZcut}{0.1}  % experiments/FIGS/data/literature_facts.json :: facts[86].value :: sha256 713dba2a5e6e1dd1 :: arXiv:2004.08262 Sec. 6.1
\newcommand{\LitCmsSdBeta}{0}  % experiments/FIGS/data/literature_facts.json :: facts[87].value :: sha256 713dba2a5e6e1dd1 :: arXiv:2004.08262 Sec. 6.1
\newcommand{\LitQgNJetsMillion}{2}  % experiments/FIGS/data/literature_facts.json :: facts[88].value :: sha256 713dba2a5e6e1dd1 :: arXiv:2202.03772 Sec. 5.2
\newcommand{\LitQgTrainMillion}{1.6}  % experiments/FIGS/data/literature_facts.json :: facts[89].value :: sha256 713dba2a5e6e1dd1 :: arXiv:2202.03772 Sec. 5.2
\newcommand{\LitQgValMillion}{0.2}  % experiments/FIGS/data/literature_facts.json :: facts[90].value :: sha256 713dba2a5e6e1dd1 :: arXiv:2202.03772 Sec. 5.2
\newcommand{\LitQgRapidityMax}{1.7}  % experiments/FIGS/data/literature_facts.json :: facts[91].value :: sha256 713dba2a5e6e1dd1 :: Zenodo 10.5281/zenodo.3164691 Zenodo record 3164691, description
\newcommand{\LitHerwigVersion}{7.1.4}  % experiments/FIGS/data/literature_facts.json :: facts[92].value :: sha256 713dba2a5e6e1dd1 :: Zenodo 10.5281/zenodo.3066475 record description
\newcommand{\LitObjectivesFewestJets}{\ensuremath{10^{3}}}  % experiments/FIGS/data/literature_facts.json :: facts[93].value :: sha256 713dba2a5e6e1dd1 :: arXiv:2606.14870 Table XVII (Micro model, Npt = 10^8: Classifier 136 vs MPM 218 at 10^3, 331 vs 539 at 10^4)
\newcommand{\LitPartTopNTrainMillion}{1.2}  % experiments/FIGS/data/literature_facts.json :: facts[94].value :: sha256 713dba2a5e6e1dd1 :: arXiv:2202.03772 Sec. 5.2
\newcommand{\LitLandscapeParticlenetEnsMean}{1412}  % experiments/FIGS/data/literature_facts.json :: facts[95].value :: sha256 713dba2a5e6e1dd1 :: arXiv:1902.09914 Table 1 (ParticleNet, ensemble mean of 1/eps_B at eps_S=0.3)
\newcommand{\LitPartFtHeadReinits}{9}  % experiments/FIGS/data/literature_facts.json :: facts[96].value :: sha256 713dba2a5e6e1dd1 :: arXiv:2202.03772 Sec. 5.2
\newcommand{\LitSophonBlocksCode}{8}  % experiments/FIGS/data/literature_facts.json :: facts[97].value :: sha256 713dba2a5e6e1dd1 :: github.com/jet-universe/sophon@9dd6dd6 networks/example_ParticleTransformer_sophon.py line 67
\newcommand{\LitGrasHerwigVersion}{2.7.1}  % experiments/FIGS/data/literature_facts.json :: facts[98].value :: sha256 713dba2a5e6e1dd1 :: arXiv:1704.03878 Sec. 2 (list of generators)
\newcommand{\DesignMassLambda}{5.0}  % experiments/MTX/k8s/job-mtx-l162_mass-s1-raunav.yaml :: --mass-lambda :: sha256 8882459fe5520795 :: identical in all ten mass-output pretraining specs
\newcommand{\FtRecipeNRuns}{264}  % experiments/FIGS/data/ft_recipes/recipes_w2b_bench_v2_v3.json :: runs (count), reported cells :: sha256 7dc8c791343e1b6f
\newcommand{\FtRecipeLr}{\ensuremath{1\times10^{-4}}}  % experiments/FIGS/data/ft_recipes/recipes_w2b_bench_v2_v3.json :: runs[*].lr, pretrained starts :: sha256 7dc8c791343e1b6f
\newcommand{\FtRecipeHeadMult}{50}  % experiments/FIGS/data/ft_recipes/recipes_w2b_bench_v2_v3.json :: runs[*].head_lr_mult, pretrained starts :: sha256 7dc8c791343e1b6f
\newcommand{\FtRecipeLrHead}{\ensuremath{5\times10^{-3}}}  % experiments/FIGS/data/ft_recipes/recipes_w2b_bench_v2_v3.json :: lr x head_lr_mult :: sha256 7dc8c791343e1b6f
\newcommand{\FtRecipeWeightDecay}{0.01}  % experiments/FIGS/data/ft_recipes/recipes_w2b_bench_v2_v3.json :: runs[*].weight_decay :: sha256 7dc8c791343e1b6f
\newcommand{\FtRecipeBatch}{512}  % experiments/FIGS/data/ft_recipes/recipes_w2b_bench_v2_v3.json :: runs[*].batch_size :: sha256 7dc8c791343e1b6f
\newcommand{\FtRecipeEpochsEThree}{50}  % experiments/FIGS/data/ft_recipes/recipes_w2b_bench_v2_v3.json :: runs[leg 1,2; n_train 1000].epochs :: sha256 7dc8c791343e1b6f
\newcommand{\FtRecipeStepsEThree}{19}  % experiments/FT/k8s/job-ft-legs-w2-raunav.yaml :: job commands: samples per epoch at n_train 1000 // batch size :: sha256 f0962eddb451b7b0 :: optimizer steps per epoch, as weaver floors them
\newcommand{\FtRecipePassesEThree}{10}  % experiments/FT/k8s/job-ft-legs-w2-raunav.yaml :: job commands: samples per epoch at n_train 1000 // n_train :: sha256 f0962eddb451b7b0 :: passes over the training subset per epoch
\newcommand{\FtRecipeEpochsEFour}{50}  % experiments/FIGS/data/ft_recipes/recipes_w2b_bench_v2_v3.json :: runs[leg 1,2; n_train 10000].epochs :: sha256 7dc8c791343e1b6f
\newcommand{\FtRecipeStepsEFour}{19}  % experiments/FT/k8s/job-ft-legs-w2-raunav.yaml :: job commands: samples per epoch at n_train 10000 // batch size :: sha256 f0962eddb451b7b0 :: optimizer steps per epoch, as weaver floors them
\newcommand{\FtRecipePassesEFour}{one}  % experiments/FT/k8s/job-ft-legs-w2-raunav.yaml :: job commands: samples per epoch at n_train 10000 // n_train :: sha256 f0962eddb451b7b0 :: passes over the training subset per epoch
\newcommand{\FtRecipeEpochsEFive}{30}  % experiments/FIGS/data/ft_recipes/recipes_w2b_bench_v2_v3.json :: runs[leg 1,2; n_train 100000].epochs :: sha256 7dc8c791343e1b6f
\newcommand{\FtRecipeStepsEFive}{195}  % experiments/FT/k8s/job-ft-legs-w2-raunav.yaml :: job commands: samples per epoch at n_train 100000 // batch size :: sha256 f0962eddb451b7b0 :: optimizer steps per epoch, as weaver floors them
\newcommand{\FtRecipePassesEFive}{one}  % experiments/FT/k8s/job-ft-legs-w2-raunav.yaml :: job commands: samples per epoch at n_train 100000 // n_train :: sha256 f0962eddb451b7b0 :: passes over the training subset per epoch
\newcommand{\FtRecipeEpochsESix}{10}  % experiments/FIGS/data/ft_recipes/recipes_w2b_bench_v2_v3.json :: runs[leg 1,2; n_train 1000000].epochs :: sha256 7dc8c791343e1b6f
\newcommand{\FtRecipeStepsESix}{1{,}953}  % experiments/FT/k8s/job-ft-legs-w2-raunav.yaml :: job commands: samples per epoch at n_train 1000000 // batch size :: sha256 f0962eddb451b7b0 :: optimizer steps per epoch, as weaver floors them
\newcommand{\FtRecipePassesESix}{one}  % experiments/FT/k8s/job-ft-legs-w2-raunav.yaml :: job commands: samples per epoch at n_train 1000000 // n_train :: sha256 f0962eddb451b7b0 :: passes over the training subset per epoch
\newcommand{\FtRecipeValJets}{19{,}968}  % experiments/FT/k8s/job-ft-legs-w2-raunav.yaml :: --samples-per-epoch-val 20000 // batch x batch :: sha256 f0962eddb451b7b0 :: validation jets per epoch, floored to whole batches
\newcommand{\RecoveryAccOneeighteightByOneeighteight}{0.4123\,\ensuremath{\pm}\,0.0011}  % experiments/FIGS/data/label_recovery_curve_v1err/analysis/label_recovery_curve_summary.json :: table[rung=L188,level=188,probe=linear,n_train=max].accuracy :: sha256 1250d1f00537e0bd :: balanced accuracy, mean +- SD over 5 runs
\newcommand{\RecoveryAccOneeighteightByOnesixtwo}{0.4090\,\ensuremath{\pm}\,0.0012}  % experiments/FIGS/data/label_recovery_curve_v1err/analysis/label_recovery_curve_summary.json :: table[rung=L188,level=162,probe=linear,n_train=max].accuracy :: sha256 1250d1f00537e0bd :: balanced accuracy, mean +- SD over 5 runs
\newcommand{\RecoveryAccOneeighteightByFourthree}{0.356\,\ensuremath{\pm}\,0.008}  % experiments/FIGS/data/label_recovery_curve_v1err/analysis/label_recovery_curve_summary.json :: table[rung=L188,level=43,probe=linear,n_train=max].accuracy :: sha256 1250d1f00537e0bd :: balanced accuracy, mean +- SD over 5 runs
\newcommand{\RecoveryAccOneeighteightByOneseven}{0.306\,\ensuremath{\pm}\,0.008}  % experiments/FIGS/data/label_recovery_curve_v1err/analysis/label_recovery_curve_summary.json :: table[rung=L188,level=17,probe=linear,n_train=max].accuracy :: sha256 1250d1f00537e0bd :: balanced accuracy, mean +- SD over 5 runs
\newcommand{\RecoveryAccOnesixtwoByOneeighteight}{0.4623\,\ensuremath{\pm}\,0.0006}  % experiments/FIGS/data/label_recovery_curve_v1err/analysis/label_recovery_curve_summary.json :: table[rung=L162,level=188,probe=linear,n_train=max].accuracy :: sha256 1250d1f00537e0bd :: balanced accuracy, mean +- SD over 5 runs
\newcommand{\RecoveryAccOnesixtwoByOnesixtwo}{0.4622\,\ensuremath{\pm}\,0.0007}  % experiments/FIGS/data/label_recovery_curve_v1err/analysis/label_recovery_curve_summary.json :: table[rung=L162,level=162,probe=linear,n_train=max].accuracy :: sha256 1250d1f00537e0bd :: balanced accuracy, mean +- SD over 5 runs
\newcommand{\RecoveryAccOnesixtwoByFourthree}{0.402\,\ensuremath{\pm}\,0.009}  % experiments/FIGS/data/label_recovery_curve_v1err/analysis/label_recovery_curve_summary.json :: table[rung=L162,level=43,probe=linear,n_train=max].accuracy :: sha256 1250d1f00537e0bd :: balanced accuracy, mean +- SD over 5 runs
\newcommand{\RecoveryAccOnesixtwoByOneseven}{0.348\,\ensuremath{\pm}\,0.010}  % experiments/FIGS/data/label_recovery_curve_v1err/analysis/label_recovery_curve_summary.json :: table[rung=L162,level=17,probe=linear,n_train=max].accuracy :: sha256 1250d1f00537e0bd :: balanced accuracy, mean +- SD over 5 runs
\newcommand{\RecoveryAccSixfourByOneeighteight}{0.5979\,\ensuremath{\pm}\,0.0008}  % experiments/FIGS/data/label_recovery_curve_v1err/analysis/label_recovery_curve_summary.json :: table[rung=R63_Q1,level=188,probe=linear,n_train=max].accuracy :: sha256 1250d1f00537e0bd :: balanced accuracy, mean +- SD over 5 runs
\newcommand{\RecoveryAccSixfourByOnesixtwo}{0.5973\,\ensuremath{\pm}\,0.0013}  % experiments/FIGS/data/label_recovery_curve_v1err/analysis/label_recovery_curve_summary.json :: table[rung=R63_Q1,level=162,probe=linear,n_train=max].accuracy :: sha256 1250d1f00537e0bd :: balanced accuracy, mean +- SD over 5 runs
\newcommand{\RecoveryAccSixfourByFourthree}{0.5881\,\ensuremath{\pm}\,0.0007}  % experiments/FIGS/data/label_recovery_curve_v1err/analysis/label_recovery_curve_summary.json :: table[rung=R63_Q1,level=43,probe=linear,n_train=max].accuracy :: sha256 1250d1f00537e0bd :: balanced accuracy, mean +- SD over 5 runs
\newcommand{\RecoveryAccSixfourByOneseven}{0.531\,\ensuremath{\pm}\,0.006}  % experiments/FIGS/data/label_recovery_curve_v1err/analysis/label_recovery_curve_summary.json :: table[rung=R63_Q1,level=17,probe=linear,n_train=max].accuracy :: sha256 1250d1f00537e0bd :: balanced accuracy, mean +- SD over 5 runs
\newcommand{\RecoveryAccFourthreeByOneeighteight}{0.6770\,\ensuremath{\pm}\,0.0006}  % experiments/FIGS/data/label_recovery_curve_v1err/analysis/label_recovery_curve_summary.json :: table[rung=R42_Q1,level=188,probe=linear,n_train=max].accuracy :: sha256 1250d1f00537e0bd :: balanced accuracy, mean +- SD over 5 runs
\newcommand{\RecoveryAccFourthreeByOnesixtwo}{0.6766\,\ensuremath{\pm}\,0.0008}  % experiments/FIGS/data/label_recovery_curve_v1err/analysis/label_recovery_curve_summary.json :: table[rung=R42_Q1,level=162,probe=linear,n_train=max].accuracy :: sha256 1250d1f00537e0bd :: balanced accuracy, mean +- SD over 5 runs
\newcommand{\RecoveryAccFourthreeByFourthree}{0.6751\,\ensuremath{\pm}\,0.0011}  % experiments/FIGS/data/label_recovery_curve_v1err/analysis/label_recovery_curve_summary.json :: table[rung=R42_Q1,level=43,probe=linear,n_train=max].accuracy :: sha256 1250d1f00537e0bd :: balanced accuracy, mean +- SD over 5 runs
\newcommand{\RecoveryAccFourthreeByOneseven}{0.633\,\ensuremath{\pm}\,0.006}  % experiments/FIGS/data/label_recovery_curve_v1err/analysis/label_recovery_curve_summary.json :: table[rung=R42_Q1,level=17,probe=linear,n_train=max].accuracy :: sha256 1250d1f00537e0bd :: balanced accuracy, mean +- SD over 5 runs
\newcommand{\RecoveryAccThreezeroByOneeighteight}{0.7001\,\ensuremath{\pm}\,0.0005}  % experiments/FIGS/data/label_recovery_curve_v1err/analysis/label_recovery_curve_summary.json :: table[rung=R29_Q1,level=188,probe=linear,n_train=max].accuracy :: sha256 1250d1f00537e0bd :: balanced accuracy, mean +- SD over 5 runs
\newcommand{\RecoveryAccThreezeroByOnesixtwo}{0.7004\,\ensuremath{\pm}\,0.0007}  % experiments/FIGS/data/label_recovery_curve_v1err/analysis/label_recovery_curve_summary.json :: table[rung=R29_Q1,level=162,probe=linear,n_train=max].accuracy :: sha256 1250d1f00537e0bd :: balanced accuracy, mean +- SD over 5 runs
\newcommand{\RecoveryAccThreezeroByFourthree}{0.6998\,\ensuremath{\pm}\,0.0008}  % experiments/FIGS/data/label_recovery_curve_v1err/analysis/label_recovery_curve_summary.json :: table[rung=R29_Q1,level=43,probe=linear,n_train=max].accuracy :: sha256 1250d1f00537e0bd :: balanced accuracy, mean +- SD over 5 runs
\newcommand{\RecoveryAccThreezeroByOneseven}{0.6778\,\ensuremath{\pm}\,0.0030}  % experiments/FIGS/data/label_recovery_curve_v1err/analysis/label_recovery_curve_summary.json :: table[rung=R29_Q1,level=17,probe=linear,n_train=max].accuracy :: sha256 1250d1f00537e0bd :: balanced accuracy, mean +- SD over 5 runs
\newcommand{\RecoveryAccOnesevenByOneeighteight}{0.7675\,\ensuremath{\pm}\,0.0010}  % experiments/FIGS/data/label_recovery_curve_v1err/analysis/label_recovery_curve_summary.json :: table[rung=R16_Q1,level=188,probe=linear,n_train=max].accuracy :: sha256 1250d1f00537e0bd :: balanced accuracy, mean +- SD over 5 runs
\newcommand{\RecoveryAccOnesevenByOnesixtwo}{0.7681\,\ensuremath{\pm}\,0.0009}  % experiments/FIGS/data/label_recovery_curve_v1err/analysis/label_recovery_curve_summary.json :: table[rung=R16_Q1,level=162,probe=linear,n_train=max].accuracy :: sha256 1250d1f00537e0bd :: balanced accuracy, mean +- SD over 5 runs
\newcommand{\RecoveryAccOnesevenByFourthree}{0.7680\,\ensuremath{\pm}\,0.0010}  % experiments/FIGS/data/label_recovery_curve_v1err/analysis/label_recovery_curve_summary.json :: table[rung=R16_Q1,level=43,probe=linear,n_train=max].accuracy :: sha256 1250d1f00537e0bd :: balanced accuracy, mean +- SD over 5 runs
\newcommand{\RecoveryAccOnesevenByOneseven}{0.7681\,\ensuremath{\pm}\,0.0003}  % experiments/FIGS/data/label_recovery_curve_v1err/analysis/label_recovery_curve_summary.json :: table[rung=R16_Q1,level=17,probe=linear,n_train=max].accuracy :: sha256 1250d1f00537e0bd :: balanced accuracy, mean +- SD over 5 runs
\newcommand{\RecoveryAccFourByOneeighteight}{0.7526\,\ensuremath{\pm}\,0.0002}  % experiments/FIGS/data/label_recovery_curve_v1err/analysis/label_recovery_curve_summary.json :: table[rung=R3_VIS,level=188,probe=linear,n_train=max].accuracy :: sha256 1250d1f00537e0bd :: balanced accuracy, mean +- SD over 5 runs
\newcommand{\RecoveryAccFourByOnesixtwo}{0.7538\,\ensuremath{\pm}\,0.0017}  % experiments/FIGS/data/label_recovery_curve_v1err/analysis/label_recovery_curve_summary.json :: table[rung=R3_VIS,level=162,probe=linear,n_train=max].accuracy :: sha256 1250d1f00537e0bd :: balanced accuracy, mean +- SD over 5 runs
\newcommand{\RecoveryAccFourByFourthree}{0.7557\,\ensuremath{\pm}\,0.0015}  % experiments/FIGS/data/label_recovery_curve_v1err/analysis/label_recovery_curve_summary.json :: table[rung=R3_VIS,level=43,probe=linear,n_train=max].accuracy :: sha256 1250d1f00537e0bd :: balanced accuracy, mean +- SD over 5 runs
\newcommand{\RecoveryAccFourByOneseven}{0.7586\,\ensuremath{\pm}\,0.0008}  % experiments/FIGS/data/label_recovery_curve_v1err/analysis/label_recovery_curve_summary.json :: table[rung=R3_VIS,level=17,probe=linear,n_train=max].accuracy :: sha256 1250d1f00537e0bd :: balanced accuracy, mean +- SD over 5 runs
\newcommand{\RecoveryAccTwoByOneeighteight}{0.8688\,\ensuremath{\pm}\,0.0005}  % experiments/FIGS/data/label_recovery_curve_v1err/analysis/label_recovery_curve_summary.json :: table[rung=R1_Q1,level=188,probe=linear,n_train=max].accuracy :: sha256 1250d1f00537e0bd :: balanced accuracy, mean +- SD over 5 runs
\newcommand{\RecoveryAccTwoByOnesixtwo}{0.8703\,\ensuremath{\pm}\,0.0008}  % experiments/FIGS/data/label_recovery_curve_v1err/analysis/label_recovery_curve_summary.json :: table[rung=R1_Q1,level=162,probe=linear,n_train=max].accuracy :: sha256 1250d1f00537e0bd :: balanced accuracy, mean +- SD over 5 runs
\newcommand{\RecoveryAccTwoByFourthree}{0.8709\,\ensuremath{\pm}\,0.0011}  % experiments/FIGS/data/label_recovery_curve_v1err/analysis/label_recovery_curve_summary.json :: table[rung=R1_Q1,level=43,probe=linear,n_train=max].accuracy :: sha256 1250d1f00537e0bd :: balanced accuracy, mean +- SD over 5 runs
\newcommand{\RecoveryAccTwoByOneseven}{0.8713\,\ensuremath{\pm}\,0.0002}  % experiments/FIGS/data/label_recovery_curve_v1err/analysis/label_recovery_curve_summary.json :: table[rung=R1_Q1,level=17,probe=linear,n_train=max].accuracy :: sha256 1250d1f00537e0bd :: balanced accuracy, mean +- SD over 5 runs
\newcommand{\RecoveryOneeighteightOnesixtwoMaxAbs}{0.0033}  % experiments/FIGS/data/label_recovery_curve_v1err/analysis/label_recovery_curve_summary.json :: table[rung=*,level=188/162,probe=linear,n_train=max].accuracy :: sha256 1250d1f00537e0bd :: max |difference of the run means| over the tree
\newcommand{\RecoveryPairedOnesevenLevelOnesevenMinusOneeighteight}{+0.0006~[\ensuremath{-}0.0004, +0.0015]}  % experiments/FIGS/data/label_recovery_curve_v1err/analysis/label_recovery_curve_summary.json :: table[rung=R16_Q1,level=17 minus 188,probe=linear,n_train=max].accuracy :: sha256 1250d1f00537e0bd :: paired by run, mean difference [95% Student-t interval]
\newcommand{\RecoveryPairedOnesevenLevelOnesevenMinusOnesixtwo}{+0.00004~[\ensuremath{-}0.00123, +0.00131]}  % experiments/FIGS/data/label_recovery_curve_v1err/analysis/label_recovery_curve_summary.json :: table[rung=R16_Q1,level=17 minus 162,probe=linear,n_train=max].accuracy :: sha256 1250d1f00537e0bd :: paired by run, mean difference [95% Student-t interval]
\newcommand{\RecoveryPairedFourLevelOnesevenMinusOneeighteight}{+0.0060~[+0.0051, +0.0069]}  % experiments/FIGS/data/label_recovery_curve_v1err/analysis/label_recovery_curve_summary.json :: table[rung=R3_VIS,level=17 minus 188,probe=linear,n_train=max].accuracy :: sha256 1250d1f00537e0bd :: paired by run, mean difference [95% Student-t interval]
\newcommand{\RecoveryPairedFourLevelOnesevenMinusOnesixtwo}{+0.0048~[+0.0028, +0.0068]}  % experiments/FIGS/data/label_recovery_curve_v1err/analysis/label_recovery_curve_summary.json :: table[rung=R3_VIS,level=17 minus 162,probe=linear,n_train=max].accuracy :: sha256 1250d1f00537e0bd :: paired by run, mean difference [95% Student-t interval]
\newcommand{\RecoveryPairedTwoLevelOnesevenMinusOneeighteight}{+0.0025~[+0.0019, +0.0032]}  % experiments/FIGS/data/label_recovery_curve_v1err/analysis/label_recovery_curve_summary.json :: table[rung=R1_Q1,level=17 minus 188,probe=linear,n_train=max].accuracy :: sha256 1250d1f00537e0bd :: paired by run, mean difference [95% Student-t interval]
\newcommand{\RecoveryPairedTwoLevelOnesevenMinusOnesixtwo}{+0.00101~[+0.00003, +0.00198]}  % experiments/FIGS/data/label_recovery_curve_v1err/analysis/label_recovery_curve_summary.json :: table[rung=R1_Q1,level=17 minus 162,probe=linear,n_train=max].accuracy :: sha256 1250d1f00537e0bd :: paired by run, mean difference [95% Student-t interval]
\newcommand{\RecoveryMlpGainOneseven}{+0.0099~[+0.0078, +0.0120]}  % experiments/FIGS/data/label_recovery_curve_v1err/analysis/label_recovery_curve_summary.json :: mlp_minus_linear_at_largest[rung=L188,level=17].per_run :: sha256 1250d1f00537e0bd :: nonlinear minus linear probe, 188 native labels from the 17-class models
\newcommand{\RecoveryLastGainOneseven}{0.0125\,\ensuremath{\pm}\,0.0021}  % experiments/FIGS/data/label_recovery_curve_v1err/analysis/label_recovery_curve_summary.json :: data_gain_last_step[rung=L188,level=17] :: sha256 1250d1f00537e0bd :: balanced accuracy gained from 443000 to 1400000 training jets
\newcommand{\RecoveryLastGainOneeighteight}{0.0147\,\ensuremath{\pm}\,0.0014}  % experiments/FIGS/data/label_recovery_curve_v1err/analysis/label_recovery_curve_summary.json :: data_gain_last_step[rung=L188,level=188] :: sha256 1250d1f00537e0bd :: balanced accuracy gained from 443000 to 1400000 training jets
\newcommand{\RecoveryNTrain}{1{,}400{,}000}  % experiments/FIGS/data/label_recovery_curve_v1err/analysis/label_recovery_curve_summary.json :: sizes (max) :: sha256 1250d1f00537e0bd :: jets the linear probe is fit on at the end of the curve
\newcommand{\RecoveryNTrainMin}{14{,}000}  % experiments/FIGS/data/label_recovery_curve_v1err/analysis/label_recovery_curve_summary.json :: sizes (min) :: sha256 1250d1f00537e0bd
\newcommand{\RecoveryNTrainPrev}{443{,}000}  % experiments/FIGS/data/label_recovery_curve_v1err/analysis/label_recovery_curve_summary.json :: sizes (second largest) :: sha256 1250d1f00537e0bd
\newcommand{\RecoveryNTest}{400{,}000}  % experiments/FIGS/data/label_recovery_curve_v1err/l188-s1/label_recovery_curve.json :: n_test :: sha256 65651e4c06c7766f :: jets the probes are scored on, recorded in all 20 per-run files
\newcommand{\RecoveryMlpHidden}{512}  % experiments/FIGS/data/label_recovery_curve_v1err/l188-s1/label_recovery_curve.json :: mlp.hidden :: sha256 65651e4c06c7766f :: hidden units of the nonlinear probe
\newcommand{\RandNDraws}{three}  % experiments/FIGS/data/probe_ladder_randcontrol_mlp2/analysis/c4_random_control.json :: table[*].draw :: sha256 c675e609ff4fe3c9 :: random partitions drawn
\newcommand{\RandNDrawsCap}{Three}  % experiments/FIGS/data/probe_ladder_randcontrol_mlp2/analysis/c4_random_control.json :: table[*].draw :: sha256 c675e609ff4fe3c9 :: random partitions drawn, capitalised to open a sentence
\newcommand{\RandOmaBvcFourprongDrawOneLinear}{\ensuremath{1.05\times10^{-2}}}  % experiments/FIGS/data/probe_ladder_randcontrol_mlp2/analysis/c4_random_control.json :: table[0].control_log1m_auc :: sha256 c675e609ff4fe3c9 :: 1-AUC = exp of it, one run
\newcommand{\RandOmaBvcFourprongDrawTwoLinear}{\ensuremath{1.11\times10^{-2}}}  % experiments/FIGS/data/probe_ladder_randcontrol_mlp2/analysis/c4_random_control.json :: table[4].control_log1m_auc :: sha256 c675e609ff4fe3c9 :: 1-AUC = exp of it, one run
\newcommand{\RandOmaBvcFourprongDrawThreeLinear}{\ensuremath{9.98\times10^{-3}}}  % experiments/FIGS/data/probe_ladder_randcontrol_mlp2/analysis/c4_random_control.json :: table[8].control_log1m_auc :: sha256 c675e609ff4fe3c9 :: 1-AUC = exp of it, one run
\newcommand{\RandOmaMeanBvcFourprongLinear}{\ensuremath{(1.05 \pm 0.06)\times10^{-2}}}  % experiments/FIGS/data/probe_ladder_randcontrol_mlp2/analysis/c4_random_control.json :: table[task=bvc_4prong,probe=linear].control_log1m_auc :: sha256 c675e609ff4fe3c9 :: 1-AUC, mean +- SD over 3 draws
\newcommand{\RandOmaBvcFourprongDrawOneMlp}{\ensuremath{1.05\times10^{-2}}}  % experiments/FIGS/data/probe_ladder_randcontrol_mlp2/analysis/c4_random_control.json :: table[1].control_log1m_auc :: sha256 c675e609ff4fe3c9 :: 1-AUC = exp of it, one run
\newcommand{\RandOmaBvcFourprongDrawTwoMlp}{\ensuremath{1.04\times10^{-2}}}  % experiments/FIGS/data/probe_ladder_randcontrol_mlp2/analysis/c4_random_control.json :: table[5].control_log1m_auc :: sha256 c675e609ff4fe3c9 :: 1-AUC = exp of it, one run
\newcommand{\RandOmaBvcFourprongDrawThreeMlp}{\ensuremath{9.28\times10^{-3}}}  % experiments/FIGS/data/probe_ladder_randcontrol_mlp2/analysis/c4_random_control.json :: table[9].control_log1m_auc :: sha256 c675e609ff4fe3c9 :: 1-AUC = exp of it, one run
\newcommand{\RandOmaMeanBvcFourprongMlp}{\ensuremath{(1.01 \pm 0.07)\times10^{-2}}}  % experiments/FIGS/data/probe_ladder_randcontrol_mlp2/analysis/c4_random_control.json :: table[task=bvc_4prong,probe=mlp].control_log1m_auc :: sha256 c675e609ff4fe3c9 :: 1-AUC, mean +- SD over 3 draws
\newcommand{\RandOmaVisibleContentDrawOneLinear}{\ensuremath{1.25\times10^{-3}}}  % experiments/FIGS/data/probe_ladder_randcontrol_mlp2/analysis/c4_random_control.json :: table[2].control_log1m_auc :: sha256 c675e609ff4fe3c9 :: 1-AUC = exp of it, one run
\newcommand{\RandOmaVisibleContentDrawTwoLinear}{\ensuremath{9.89\times10^{-4}}}  % experiments/FIGS/data/probe_ladder_randcontrol_mlp2/analysis/c4_random_control.json :: table[6].control_log1m_auc :: sha256 c675e609ff4fe3c9 :: 1-AUC = exp of it, one run
\newcommand{\RandOmaVisibleContentDrawThreeLinear}{\ensuremath{1.08\times10^{-3}}}  % experiments/FIGS/data/probe_ladder_randcontrol_mlp2/analysis/c4_random_control.json :: table[10].control_log1m_auc :: sha256 c675e609ff4fe3c9 :: 1-AUC = exp of it, one run
\newcommand{\RandOmaMeanVisibleContentLinear}{\ensuremath{(1.11 \pm 0.13)\times10^{-3}}}  % experiments/FIGS/data/probe_ladder_randcontrol_mlp2/analysis/c4_random_control.json :: table[task=visible_content,probe=linear].control_log1m_auc :: sha256 c675e609ff4fe3c9 :: 1-AUC, mean +- SD over 3 draws
\newcommand{\RandOmaVisibleContentDrawOneMlp}{\ensuremath{1.21\times10^{-3}}}  % experiments/FIGS/data/probe_ladder_randcontrol_mlp2/analysis/c4_random_control.json :: table[3].control_log1m_auc :: sha256 c675e609ff4fe3c9 :: 1-AUC = exp of it, one run
\newcommand{\RandOmaVisibleContentDrawTwoMlp}{\ensuremath{1.07\times10^{-3}}}  % experiments/FIGS/data/probe_ladder_randcontrol_mlp2/analysis/c4_random_control.json :: table[7].control_log1m_auc :: sha256 c675e609ff4fe3c9 :: 1-AUC = exp of it, one run
\newcommand{\RandOmaVisibleContentDrawThreeMlp}{\ensuremath{1.13\times10^{-3}}}  % experiments/FIGS/data/probe_ladder_randcontrol_mlp2/analysis/c4_random_control.json :: table[11].control_log1m_auc :: sha256 c675e609ff4fe3c9 :: 1-AUC = exp of it, one run
\newcommand{\RandOmaMeanVisibleContentMlp}{\ensuremath{(1.14 \pm 0.07)\times10^{-3}}}  % experiments/FIGS/data/probe_ladder_randcontrol_mlp2/analysis/c4_random_control.json :: table[task=visible_content,probe=mlp].control_log1m_auc :: sha256 c675e609ff4fe3c9 :: 1-AUC, mean +- SD over 3 draws
\newcommand{\RandCostFactorBvcFourprongLinear}{1.49}  % experiments/FIGS/data/probe_ladder_randcontrol_mlp2/analysis/c4_random_control.json :: grouping_cost_post_hoc.tasks.bvc_4prong.linear.factor_control :: sha256 c675e609ff4fe3c9 :: POST HOC: random grouping vs finer models, in 1-AUC
\newcommand{\RandSemCostFactorBvcFourprongLinear}{4.42}  % experiments/FIGS/data/probe_ladder_randcontrol_mlp2/analysis/c4_random_control.json :: grouping_cost_post_hoc.tasks.bvc_4prong.linear.factor_semantic :: sha256 c675e609ff4fe3c9 :: POST HOC: semantic 17-class vs finer models
\newcommand{\RandCostFactorBvcFourprongMlp}{1.41}  % experiments/FIGS/data/probe_ladder_randcontrol_mlp2/analysis/c4_random_control.json :: grouping_cost_post_hoc.tasks.bvc_4prong.mlp.factor_control :: sha256 c675e609ff4fe3c9 :: POST HOC: random grouping vs finer models, in 1-AUC
\newcommand{\RandSemCostFactorBvcFourprongMlp}{3.26}  % experiments/FIGS/data/probe_ladder_randcontrol_mlp2/analysis/c4_random_control.json :: grouping_cost_post_hoc.tasks.bvc_4prong.mlp.factor_semantic :: sha256 c675e609ff4fe3c9 :: POST HOC: semantic 17-class vs finer models
\newcommand{\RandCostFactorVisibleContentLinear}{1.35}  % experiments/FIGS/data/probe_ladder_randcontrol_mlp2/analysis/c4_random_control.json :: grouping_cost_post_hoc.tasks.visible_content.linear.factor_control :: sha256 c675e609ff4fe3c9 :: POST HOC: random grouping vs finer models, in 1-AUC
\newcommand{\RandSemCostFactorVisibleContentLinear}{2.82}  % experiments/FIGS/data/probe_ladder_randcontrol_mlp2/analysis/c4_random_control.json :: grouping_cost_post_hoc.tasks.visible_content.linear.factor_semantic :: sha256 c675e609ff4fe3c9 :: POST HOC: semantic 17-class vs finer models
\newcommand{\RandCostFactorVisibleContentMlp}{1.33}  % experiments/FIGS/data/probe_ladder_randcontrol_mlp2/analysis/c4_random_control.json :: grouping_cost_post_hoc.tasks.visible_content.mlp.factor_control :: sha256 c675e609ff4fe3c9 :: POST HOC: random grouping vs finer models, in 1-AUC
\newcommand{\RandSemCostFactorVisibleContentMlp}{2.37}  % experiments/FIGS/data/probe_ladder_randcontrol_mlp2/analysis/c4_random_control.json :: grouping_cost_post_hoc.tasks.visible_content.mlp.factor_semantic :: sha256 c675e609ff4fe3c9 :: POST HOC: semantic 17-class vs finer models
\newcommand{\MassResNTest}{398{,}625}  % experiments/FIGS/data/mass_resolution/s1.json :: centering_detail.split[2] :: sha256 88a57bbad270d51c :: test jets sigma_eff is computed on, the same in every per-seed file and arm
\newcommand{\MassResNClasses}{184}  % experiments/FIGS/data/mass_resolution/analysis_holm/s7_mass_resolution.json :: provenance.n_classes_used :: sha256 e03cde45ece9f385 :: native classes with enough training jets to centre
\newcommand{\MassResSigmaEffRidgeOneeighteight}{0.3393\,\ensuremath{\pm}\,0.0019}  % experiments/FIGS/data/mass_resolution/analysis_holm/s7_mass_resolution.json :: table[cell=188,probe=ridge].sigma_eff :: sha256 e03cde45ece9f385 :: mean +- SD over 5 seeds
\newcommand{\MassResSigmaEffRidgeOnesixtwo}{0.3302\,\ensuremath{\pm}\,0.0029}  % experiments/FIGS/data/mass_resolution/analysis_holm/s7_mass_resolution.json :: table[cell=162,probe=ridge].sigma_eff :: sha256 e03cde45ece9f385 :: mean +- SD over 5 seeds
\newcommand{\MassResSigmaEffRidgeFourthree}{0.3228\,\ensuremath{\pm}\,0.0019}  % experiments/FIGS/data/mass_resolution/analysis_holm/s7_mass_resolution.json :: table[cell=43,probe=ridge].sigma_eff :: sha256 e03cde45ece9f385 :: mean +- SD over 5 seeds
\newcommand{\MassResSigmaEffRidgeOneseven}{0.300\,\ensuremath{\pm}\,0.010}  % experiments/FIGS/data/mass_resolution/analysis_holm/s7_mass_resolution.json :: table[cell=17,probe=ridge].sigma_eff :: sha256 e03cde45ece9f385 :: mean +- SD over 5 seeds
\newcommand{\MassResSigmaEffRidgeOnesixtwoMass}{0.2238\,\ensuremath{\pm}\,0.0019}  % experiments/FIGS/data/mass_resolution/analysis_holm/s7_mass_resolution.json :: table[cell=162+mass,probe=ridge].sigma_eff :: sha256 e03cde45ece9f385 :: mean +- SD over 5 seeds
\newcommand{\MassResSigmaEffRidgeOnesevenMass}{0.2170\,\ensuremath{\pm}\,0.0034}  % experiments/FIGS/data/mass_resolution/analysis_holm/s7_mass_resolution.json :: table[cell=17+mass,probe=ridge].sigma_eff :: sha256 e03cde45ece9f385 :: mean +- SD over 5 seeds
\newcommand{\MassResGainRidgeOnesixtwo}{\ensuremath{-}0.1063}  % experiments/FIGS/data/mass_resolution/analysis_holm/s7_mass_resolution.json :: table[cell=162+mass/162,probe=ridge].sigma_eff :: sha256 e03cde45ece9f385 :: sigma_eff with the mass output minus without, difference of the seed means; negative is better
\newcommand{\MassResGainRidgeOneseven}{\ensuremath{-}0.0831}  % experiments/FIGS/data/mass_resolution/analysis_holm/s7_mass_resolution.json :: table[cell=17+mass/17,probe=ridge].sigma_eff :: sha256 e03cde45ece9f385 :: sigma_eff with the mass output minus without, difference of the seed means; negative is better
\newcommand{\MassResDidRidge}{\ensuremath{-}0.0233}  % experiments/FIGS/data/mass_resolution/analysis_holm/s7_mass_resolution.json :: table[cell=162+mass/162/17+mass/17,probe=ridge].sigma_eff :: sha256 e03cde45ece9f385 :: gain at 162 classes minus gain at 17
\newcommand{\MassResSigmaEffMlpOneeighteight}{0.2463\,\ensuremath{\pm}\,0.0022}  % experiments/FIGS/data/mass_resolution/analysis_holm/s7_mass_resolution.json :: table[cell=188,probe=mlp].sigma_eff :: sha256 e03cde45ece9f385 :: mean +- SD over 5 seeds
\newcommand{\MassResSigmaEffMlpOnesixtwo}{0.2438\,\ensuremath{\pm}\,0.0018}  % experiments/FIGS/data/mass_resolution/analysis_holm/s7_mass_resolution.json :: table[cell=162,probe=mlp].sigma_eff :: sha256 e03cde45ece9f385 :: mean +- SD over 5 seeds
\newcommand{\MassResSigmaEffMlpFourthree}{0.2407\,\ensuremath{\pm}\,0.0009}  % experiments/FIGS/data/mass_resolution/analysis_holm/s7_mass_resolution.json :: table[cell=43,probe=mlp].sigma_eff :: sha256 e03cde45ece9f385 :: mean +- SD over 5 seeds
\newcommand{\MassResSigmaEffMlpOneseven}{0.2358\,\ensuremath{\pm}\,0.0025}  % experiments/FIGS/data/mass_resolution/analysis_holm/s7_mass_resolution.json :: table[cell=17,probe=mlp].sigma_eff :: sha256 e03cde45ece9f385 :: mean +- SD over 5 seeds
\newcommand{\MassResSigmaEffMlpOnesixtwoMass}{0.1943\,\ensuremath{\pm}\,0.0001}  % experiments/FIGS/data/mass_resolution/analysis_holm/s7_mass_resolution.json :: table[cell=162+mass,probe=mlp].sigma_eff :: sha256 e03cde45ece9f385 :: mean +- SD over 5 seeds
\newcommand{\MassResSigmaEffMlpOnesevenMass}{0.1949\,\ensuremath{\pm}\,0.0006}  % experiments/FIGS/data/mass_resolution/analysis_holm/s7_mass_resolution.json :: table[cell=17+mass,probe=mlp].sigma_eff :: sha256 e03cde45ece9f385 :: mean +- SD over 5 seeds
\newcommand{\MassResGainMlpOnesixtwo}{\ensuremath{-}0.0496}  % experiments/FIGS/data/mass_resolution/analysis_holm/s7_mass_resolution.json :: table[cell=162+mass/162,probe=mlp].sigma_eff :: sha256 e03cde45ece9f385 :: sigma_eff with the mass output minus without, difference of the seed means; negative is better
\newcommand{\MassResGainMlpOneseven}{\ensuremath{-}0.0409}  % experiments/FIGS/data/mass_resolution/analysis_holm/s7_mass_resolution.json :: table[cell=17+mass/17,probe=mlp].sigma_eff :: sha256 e03cde45ece9f385 :: sigma_eff with the mass output minus without, difference of the seed means; negative is better
\newcommand{\MassResDidMlp}{\ensuremath{-}0.0087}  % experiments/FIGS/data/mass_resolution/analysis_holm/s7_mass_resolution.json :: table[cell=162+mass/162/17+mass/17,probe=mlp].sigma_eff :: sha256 e03cde45ece9f385 :: gain at 162 classes minus gain at 17
\newcommand{\MassResTargetSigmaEff}{0.2986}  % experiments/FIGS/data/mass_resolution/analysis_holm/s7_mass_resolution.json :: table[*].target_sigma_eff :: sha256 e03cde45ece9f385 :: sigma_eff of the class-centred target itself (a model that predicts the class mean)
\newcommand{\MassResRidgeRtwoMin}{20\%}  % experiments/FIGS/data/mass_resolution/s1.json :: arms.{l188,l162,r42q1,r16q1}-s*.ridge.val_r2 :: sha256 88a57bbad270d51c :: validation R^2 of the ridge probe, smallest of 20 models without the mass output, all 5 per-seed files
\newcommand{\MassResRidgeRtwoMax}{30\%}  % experiments/FIGS/data/mass_resolution/s1.json :: arms.{l188,l162,r42q1,r16q1}-s*.ridge.val_r2 :: sha256 88a57bbad270d51c :: validation R^2 of the ridge probe, largest of 20 models without the mass output, all 5 per-seed files
\newcommand{\MassResSigmaEffOnesixtwo}{0.3302}  % experiments/FIGS/data/mass_resolution/analysis_holm/s7_mass_resolution.json :: table[cell=162,probe=ridge].sigma_eff :: sha256 e03cde45ece9f385 :: mean over 5 seeds
\newcommand{\MassResSdOnesixtwo}{0.626}  % experiments/FIGS/data/mass_resolution/analysis_holm/s7_mass_resolution.json :: table[cell=162,probe=ridge].sd :: sha256 e03cde45ece9f385 :: mean over 5 seeds
\newcommand{\MassResTailFractionOnesixtwo}{0.062}  % experiments/FIGS/data/mass_resolution/analysis_holm/s7_mass_resolution.json :: table[cell=162,probe=ridge].tail_fraction :: sha256 e03cde45ece9f385 :: mean over 5 seeds
\newcommand{\MassResSigmaEffOnesixtwoMass}{0.2238}  % experiments/FIGS/data/mass_resolution/analysis_holm/s7_mass_resolution.json :: table[cell=162+mass,probe=ridge].sigma_eff :: sha256 e03cde45ece9f385 :: mean over 5 seeds
\newcommand{\MassResSdOnesixtwoMass}{0.494}  % experiments/FIGS/data/mass_resolution/analysis_holm/s7_mass_resolution.json :: table[cell=162+mass,probe=ridge].sd :: sha256 e03cde45ece9f385 :: mean over 5 seeds
\newcommand{\MassResTailFractionOnesixtwoMass}{0.034}  % experiments/FIGS/data/mass_resolution/analysis_holm/s7_mass_resolution.json :: table[cell=162+mass,probe=ridge].tail_fraction :: sha256 e03cde45ece9f385 :: mean over 5 seeds
\newcommand{\MassResSigmaEffOneseven}{0.3001}  % experiments/FIGS/data/mass_resolution/analysis_holm/s7_mass_resolution.json :: table[cell=17,probe=ridge].sigma_eff :: sha256 e03cde45ece9f385 :: mean over 5 seeds
\newcommand{\MassResSdOneseven}{0.601}  % experiments/FIGS/data/mass_resolution/analysis_holm/s7_mass_resolution.json :: table[cell=17,probe=ridge].sd :: sha256 e03cde45ece9f385 :: mean over 5 seeds
\newcommand{\MassResTailFractionOneseven}{0.056}  % experiments/FIGS/data/mass_resolution/analysis_holm/s7_mass_resolution.json :: table[cell=17,probe=ridge].tail_fraction :: sha256 e03cde45ece9f385 :: mean over 5 seeds
\newcommand{\MassResSigmaEffOnesevenMass}{0.2170}  % experiments/FIGS/data/mass_resolution/analysis_holm/s7_mass_resolution.json :: table[cell=17+mass,probe=ridge].sigma_eff :: sha256 e03cde45ece9f385 :: mean over 5 seeds
\newcommand{\MassResSdOnesevenMass}{0.490}  % experiments/FIGS/data/mass_resolution/analysis_holm/s7_mass_resolution.json :: table[cell=17+mass,probe=ridge].sd :: sha256 e03cde45ece9f385 :: mean over 5 seeds
\newcommand{\MassResTailFractionOnesevenMass}{0.033}  % experiments/FIGS/data/mass_resolution/analysis_holm/s7_mass_resolution.json :: table[cell=17+mass,probe=ridge].tail_fraction :: sha256 e03cde45ece9f385 :: mean over 5 seeds
\newcommand{\MassResSigmaEffOneeighteight}{0.3393}  % experiments/FIGS/data/mass_resolution/analysis_holm/s7_mass_resolution.json :: table[cell=188,probe=ridge].sigma_eff :: sha256 e03cde45ece9f385 :: mean over 5 seeds
\newcommand{\MassResSdOneeighteight}{0.631}  % experiments/FIGS/data/mass_resolution/analysis_holm/s7_mass_resolution.json :: table[cell=188,probe=ridge].sd :: sha256 e03cde45ece9f385 :: mean over 5 seeds
\newcommand{\MassResTailFractionOneeighteight}{0.064}  % experiments/FIGS/data/mass_resolution/analysis_holm/s7_mass_resolution.json :: table[cell=188,probe=ridge].tail_fraction :: sha256 e03cde45ece9f385 :: mean over 5 seeds
\newcommand{\MassResSigmaEffFourthree}{0.3228}  % experiments/FIGS/data/mass_resolution/analysis_holm/s7_mass_resolution.json :: table[cell=43,probe=ridge].sigma_eff :: sha256 e03cde45ece9f385 :: mean over 5 seeds
\newcommand{\MassResSdFourthree}{0.618}  % experiments/FIGS/data/mass_resolution/analysis_holm/s7_mass_resolution.json :: table[cell=43,probe=ridge].sd :: sha256 e03cde45ece9f385 :: mean over 5 seeds
\newcommand{\MassResTailFractionFourthree}{0.060}  % experiments/FIGS/data/mass_resolution/analysis_holm/s7_mass_resolution.json :: table[cell=43,probe=ridge].tail_fraction :: sha256 e03cde45ece9f385 :: mean over 5 seeds
\newcommand{\AojEffPercent}{1\%}  % experiments/FIGS/data/aoj_full_v1/fit_v6/results.json :: eff :: sha256 2416fc21f93fcfa6 :: the score cut's pass fraction outside the mass windows
\newcommand{\AojPtMin}{500}  % experiments/AOJ/peak_fit.py :: PT_RANGE[0] :: sha256 cad06229a74532be :: GeV
\newcommand{\AojPtMax}{2{,}500}  % experiments/AOJ/peak_fit.py :: PT_RANGE[1] :: sha256 cad06229a74532be :: GeV
\newcommand{\AojRhoMin}{\ensuremath{-}5.5}  % experiments/AOJ/peak_fit.py :: RHO_RANGE[0] :: sha256 cad06229a74532be
\newcommand{\AojRhoMax}{\ensuremath{-}2.0}  % experiments/AOJ/peak_fit.py :: RHO_RANGE[1] :: sha256 cad06229a74532be
\newcommand{\AojWindowWLo}{65}  % experiments/AOJ/peak_fit.py :: PEAKS[W].window[0] :: sha256 cad06229a74532be :: GeV
\newcommand{\AojWindowWHi}{105}  % experiments/AOJ/peak_fit.py :: PEAKS[W].window[1] :: sha256 cad06229a74532be :: GeV
\newcommand{\AojWindowTopLo}{140}  % experiments/AOJ/peak_fit.py :: PEAKS[top].window[0] :: sha256 cad06229a74532be :: GeV
\newcommand{\AojWindowTopHi}{220}  % experiments/AOJ/peak_fit.py :: PEAKS[top].window[1] :: sha256 cad06229a74532be :: GeV
\newcommand{\AojFitLo}{105}  % experiments/AOJ/peak_fit.py :: PEAKS[top].fit_range[0] :: sha256 cad06229a74532be :: GeV
\newcommand{\AojFitHi}{300}  % experiments/AOJ/peak_fit.py :: PEAKS[top].fit_range[1] :: sha256 cad06229a74532be :: GeV
\newcommand{\AojMassBin}{5}  % experiments/AOJ/peak_fit.py :: MASS_BIN :: sha256 cad06229a74532be :: GeV
\newcommand{\AojMapOrderRho}{3}  % experiments/AOJ/peak_fit.py :: MAP_ORDER[0] :: sha256 cad06229a74532be :: polynomial order in rho
\newcommand{\AojMapOrderPt}{2}  % experiments/AOJ/peak_fit.py :: MAP_ORDER[1] :: sha256 cad06229a74532be :: polynomial order in ln pT
\newcommand{\AojNRunsWord}{five}  % experiments/FIGS/data/aoj_full_v1/analysis_v6/aoj_top.json :: per_label_set.label_sets.*.signal_yield.n :: sha256 5dfef902a5fe75ba :: pretraining runs per vocabulary
\newcommand{\AojNJetsRhoWindow}{9{,}441{,}209}  % experiments/FIGS/data/aoj_full_v1/fit_v6/results.json :: n_jets :: sha256 2416fc21f93fcfa6 :: staged jets inside the rho window, all files; not the jets in the mass fit
\newcommand{\AojYieldOneeighteight}{2{,}400\,\ensuremath{\pm}\,500}  % experiments/FIGS/data/aoj_full_v1/analysis_v6/aoj_top.json :: per_label_set.label_sets.188.signal_yield :: sha256 5dfef902a5fe75ba :: mean +- SD over 5 runs, shared shape
\newcommand{\AojStatErrOneeighteight}{211}  % experiments/FIGS/data/aoj_full_v1/analysis_v6/aoj_top.json :: per_label_set.label_sets.188.median_stat_err :: sha256 5dfef902a5fe75ba :: median per-fit statistical error
\newcommand{\AojYieldOwnShapeOneeighteight}{2{,}400\,\ensuremath{\pm}\,600}  % experiments/FIGS/data/aoj_full_v1/analysis_v6/aoj_top.json :: per_label_set.label_sets.188.by_shape.own_shape :: sha256 5dfef902a5fe75ba :: every fit with its own floated shape
\newcommand{\AojSpreadOverStatOneeighteight}{2.4}  % experiments/FIGS/data/aoj_full_v1/analysis_v6/aoj_top.json :: per_label_set.label_sets.188.signal_yield.sd / median_stat_err :: sha256 5dfef902a5fe75ba :: spread between runs over the median per-fit statistical error
\newcommand{\AojYieldOnesixtwo}{2{,}400\,\ensuremath{\pm}\,900}  % experiments/FIGS/data/aoj_full_v1/analysis_v6/aoj_top.json :: per_label_set.label_sets.162.signal_yield :: sha256 5dfef902a5fe75ba :: mean +- SD over 5 runs, shared shape
\newcommand{\AojStatErrOnesixtwo}{199}  % experiments/FIGS/data/aoj_full_v1/analysis_v6/aoj_top.json :: per_label_set.label_sets.162.median_stat_err :: sha256 5dfef902a5fe75ba :: median per-fit statistical error
\newcommand{\AojYieldOwnShapeOnesixtwo}{2{,}200\,\ensuremath{\pm}\,800}  % experiments/FIGS/data/aoj_full_v1/analysis_v6/aoj_top.json :: per_label_set.label_sets.162.by_shape.own_shape :: sha256 5dfef902a5fe75ba :: every fit with its own floated shape
\newcommand{\AojSpreadOverStatOnesixtwo}{4.5}  % experiments/FIGS/data/aoj_full_v1/analysis_v6/aoj_top.json :: per_label_set.label_sets.162.signal_yield.sd / median_stat_err :: sha256 5dfef902a5fe75ba :: spread between runs over the median per-fit statistical error
\newcommand{\AojYieldFourthree}{2{,}400\,\ensuremath{\pm}\,800}  % experiments/FIGS/data/aoj_full_v1/analysis_v6/aoj_top.json :: per_label_set.label_sets.43.signal_yield :: sha256 5dfef902a5fe75ba :: mean +- SD over 5 runs, shared shape
\newcommand{\AojStatErrFourthree}{184}  % experiments/FIGS/data/aoj_full_v1/analysis_v6/aoj_top.json :: per_label_set.label_sets.43.median_stat_err :: sha256 5dfef902a5fe75ba :: median per-fit statistical error
\newcommand{\AojYieldOwnShapeFourthree}{2{,}300\,\ensuremath{\pm}\,800}  % experiments/FIGS/data/aoj_full_v1/analysis_v6/aoj_top.json :: per_label_set.label_sets.43.by_shape.own_shape :: sha256 5dfef902a5fe75ba :: every fit with its own floated shape
\newcommand{\AojSpreadOverStatFourthree}{4.2}  % experiments/FIGS/data/aoj_full_v1/analysis_v6/aoj_top.json :: per_label_set.label_sets.43.signal_yield.sd / median_stat_err :: sha256 5dfef902a5fe75ba :: spread between runs over the median per-fit statistical error
\newcommand{\AojYieldOneseven}{2{,}350\,\ensuremath{\pm}\,240}  % experiments/FIGS/data/aoj_full_v1/analysis_v6/aoj_top.json :: per_label_set.label_sets.17.signal_yield :: sha256 5dfef902a5fe75ba :: mean +- SD over 5 runs, shared shape
\newcommand{\AojStatErrOneseven}{201}  % experiments/FIGS/data/aoj_full_v1/analysis_v6/aoj_top.json :: per_label_set.label_sets.17.median_stat_err :: sha256 5dfef902a5fe75ba :: median per-fit statistical error
\newcommand{\AojYieldOwnShapeOneseven}{2{,}480\,\ensuremath{\pm}\,290}  % experiments/FIGS/data/aoj_full_v1/analysis_v6/aoj_top.json :: per_label_set.label_sets.17.by_shape.own_shape :: sha256 5dfef902a5fe75ba :: every fit with its own floated shape
\newcommand{\AojSpreadOverStatOneseven}{1.2}  % experiments/FIGS/data/aoj_full_v1/analysis_v6/aoj_top.json :: per_label_set.label_sets.17.signal_yield.sd / median_stat_err :: sha256 5dfef902a5fe75ba :: spread between runs over the median per-fit statistical error
\newcommand{\AojYieldOnesixtwoMass}{2{,}420\,\ensuremath{\pm}\,240}  % experiments/FIGS/data/aoj_full_v1/analysis_v6/aoj_top.json :: per_label_set.label_sets.162+mass.signal_yield :: sha256 5dfef902a5fe75ba :: mean +- SD over 5 runs, shared shape
\newcommand{\AojStatErrOnesixtwoMass}{179}  % experiments/FIGS/data/aoj_full_v1/analysis_v6/aoj_top.json :: per_label_set.label_sets.162+mass.median_stat_err :: sha256 5dfef902a5fe75ba :: median per-fit statistical error
\newcommand{\AojYieldOwnShapeOnesixtwoMass}{2{,}350\,\ensuremath{\pm}\,320}  % experiments/FIGS/data/aoj_full_v1/analysis_v6/aoj_top.json :: per_label_set.label_sets.162+mass.by_shape.own_shape :: sha256 5dfef902a5fe75ba :: every fit with its own floated shape
\newcommand{\AojSpreadOverStatOnesixtwoMass}{1.3}  % experiments/FIGS/data/aoj_full_v1/analysis_v6/aoj_top.json :: per_label_set.label_sets.162+mass.signal_yield.sd / median_stat_err :: sha256 5dfef902a5fe75ba :: spread between runs over the median per-fit statistical error
\newcommand{\AojYieldOnesevenMass}{2{,}500\,\ensuremath{\pm}\,600}  % experiments/FIGS/data/aoj_full_v1/analysis_v6/aoj_top.json :: per_label_set.label_sets.17+mass.signal_yield :: sha256 5dfef902a5fe75ba :: mean +- SD over 5 runs, shared shape
\newcommand{\AojStatErrOnesevenMass}{200}  % experiments/FIGS/data/aoj_full_v1/analysis_v6/aoj_top.json :: per_label_set.label_sets.17+mass.median_stat_err :: sha256 5dfef902a5fe75ba :: median per-fit statistical error
\newcommand{\AojYieldOwnShapeOnesevenMass}{2{,}900\,\ensuremath{\pm}\,500}  % experiments/FIGS/data/aoj_full_v1/analysis_v6/aoj_top.json :: per_label_set.label_sets.17+mass.by_shape.own_shape :: sha256 5dfef902a5fe75ba :: every fit with its own floated shape
\newcommand{\AojSpreadOverStatOnesevenMass}{2.9}  % experiments/FIGS/data/aoj_full_v1/analysis_v6/aoj_top.json :: per_label_set.label_sets.17+mass.signal_yield.sd / median_stat_err :: sha256 5dfef902a5fe75ba :: spread between runs over the median per-fit statistical error
\newcommand{\AojOwnShapeShiftMin}{\ensuremath{-}7\%}  % experiments/FIGS/data/aoj_full_v1/analysis_v6/aoj_top.json :: per_label_set.label_sets.*.by_shape.own_shape.mean / signal_yield.mean - 1 :: sha256 5dfef902a5fe75ba :: the vocabulary means with each model's own shape, relative to the shared shape
\newcommand{\AojOwnShapeShiftMax}{+13\%}  % experiments/FIGS/data/aoj_full_v1/analysis_v6/aoj_top.json :: per_label_set.label_sets.*.by_shape.own_shape.mean / signal_yield.mean - 1 :: sha256 5dfef902a5fe75ba :: the vocabulary means with each model's own shape, relative to the shared shape
\newcommand{\AojPooledMean}{182.8}  % experiments/FIGS/data/aoj_full_v1/analysis_v6/aoj_top.json :: per_label_set.shapes.pooled :: sha256 5dfef902a5fe75ba :: GeV
\newcommand{\AojPooledWidth}{11.7}  % experiments/FIGS/data/aoj_full_v1/analysis_v6/aoj_top.json :: per_label_set.shapes.pooled :: sha256 5dfef902a5fe75ba :: GeV
\newcommand{\AojPoolSize}{30}  % experiments/FIGS/data/aoj_full_v1/analysis_v6/aoj_top.json :: per_label_set.shapes.pool (count) :: sha256 5dfef902a5fe75ba :: pretrained models the shared shape is fitted over
\newcommand{\AojNModels}{30}  % experiments/FIGS/data/aoj_full_v1/fit_v6/results.json :: models (pretrained) :: sha256 2416fc21f93fcfa6 :: fits
\newcommand{\AojOwnMeanMin}{181.6}  % experiments/FIGS/data/aoj_full_v1/fit_v6/results.json :: models.*.top.own_shape_fit.mean (min) :: sha256 2416fc21f93fcfa6 :: GeV
\newcommand{\AojOwnMeanMax}{184.0}  % experiments/FIGS/data/aoj_full_v1/fit_v6/results.json :: models.*.top.own_shape_fit.mean (max) :: sha256 2416fc21f93fcfa6 :: GeV
\newcommand{\AojOwnWidthMin}{10.0}  % experiments/FIGS/data/aoj_full_v1/fit_v6/results.json :: models.*.top.own_shape_fit.width (min) :: sha256 2416fc21f93fcfa6 :: GeV
\newcommand{\AojOwnWidthMax}{13.8}  % experiments/FIGS/data/aoj_full_v1/fit_v6/results.json :: models.*.top.own_shape_fit.width (max) :: sha256 2416fc21f93fcfa6 :: GeV
\newcommand{\AojShapeStepMean}{0.7}  % experiments/FIGS/data/aoj_full_v1/fit_v6/results.json :: pooled_shape.systematic_step.mean :: sha256 2416fc21f93fcfa6 :: GeV: the SD of the models' own floated values, one shift for every model together
\newcommand{\AojShapeSystMeanMax}{1\%}  % experiments/FIGS/data/aoj_full_v1/fit_v6/results.json :: models.*.top.shape_systematic.pooled_mean_(down|up) / signal_yield - 1 (largest |.|) :: sha256 2416fc21f93fcfa6 :: the largest yield change of any model under that shift
\newcommand{\AojShapeStepWidth}{1.0}  % experiments/FIGS/data/aoj_full_v1/fit_v6/results.json :: pooled_shape.systematic_step.width :: sha256 2416fc21f93fcfa6 :: GeV: the SD of the models' own floated values, one shift for every model together
\newcommand{\AojShapeSystWidthMax}{12\%}  % experiments/FIGS/data/aoj_full_v1/fit_v6/results.json :: models.*.top.shape_systematic.pooled_width_(down|up) / signal_yield - 1 (largest |.|) :: sha256 2416fc21f93fcfa6 :: the largest yield change of any model under that shift
\newcommand{\AojLeakEpsHi}{0.6}  % experiments/FIGS/data/aoj_full_v1/fit_v6/results.json :: models.*.top.leak_systematic.eps_ref (max) :: sha256 2416fc21f93fcfa6 :: the reference's assumed efficiency for data tops, upper variation
\newcommand{\AojLeakEpsLo}{0.4}  % experiments/FIGS/data/aoj_full_v1/fit_v6/results.json :: models.*.top.leak_systematic.eps_ref (min) :: sha256 2416fc21f93fcfa6 :: the same, lower variation
\newcommand{\AojLeakShiftMin}{+97}  % experiments/FIGS/data/aoj_full_v1/fit_v6/results.json :: models.*.top.leak_systematic.shift['0.6'] (min) :: sha256 2416fc21f93fcfa6 :: tops added to a yield
\newcommand{\AojLeakShiftMax}{+271}  % experiments/FIGS/data/aoj_full_v1/fit_v6/results.json :: models.*.top.leak_systematic.shift['0.4'] (max) :: sha256 2416fc21f93fcfa6 :: tops added to a yield
\newcommand{\AojFitPMin}{0.010}  % experiments/FIGS/data/aoj_full_v1/fit_v6/results.json :: models.*.top.asymptotic_p (min) :: sha256 2416fc21f93fcfa6 :: goodness of fit of the working-point fit, pretrained models
\newcommand{\AojFitPMax}{0.974}  % experiments/FIGS/data/aoj_full_v1/fit_v6/results.json :: models.*.top.asymptotic_p (max) :: sha256 2416fc21f93fcfa6
\newcommand{\AojWeakestYield}{885}  % experiments/FIGS/data/aoj_full_v1/fit_v6/results.json :: models.l162-s2.top.signal_yield :: sha256 2416fc21f93fcfa6 :: the smallest yield of any pretrained model
\newcommand{\AojWeakestYieldErr}{211}  % experiments/FIGS/data/aoj_full_v1/fit_v6/results.json :: models.l162-s2.top.signal_yield_err :: sha256 2416fc21f93fcfa6
\newcommand{\AojWeakestFitP}{0.94}  % experiments/FIGS/data/aoj_full_v1/fit_v6/results.json :: models.l162-s2.top.asymptotic_p :: sha256 2416fc21f93fcfa6
\newcommand{\AojSidebandZMin}{\ensuremath{-}0.43}  % experiments/FIGS/data/aoj_full_v1/fit_v6/histograms.npz :: experiments/AOJ/sideband_residuals.py: z (min over the pretrained models) :: sha256 2bd1e3d5b32d4a12 :: passing jets minus the fit, summed from the fit's lower edge to the top window
\newcommand{\AojSidebandZMax}{+0.58}  % experiments/FIGS/data/aoj_full_v1/fit_v6/histograms.npz :: experiments/AOJ/sideband_residuals.py: z (max over the pretrained models) :: sha256 2bd1e3d5b32d4a12
\newcommand{\AojWeakestSidebandZ}{0.00}  % experiments/FIGS/data/aoj_full_v1/fit_v6/histograms.npz :: experiments/AOJ/sideband_residuals.py: z of l162-s2 :: sha256 2bd1e3d5b32d4a12 :: the smallest-yield model
\newcommand{\AojSidebandMaxPull}{2.7}  % experiments/FIGS/data/aoj_full_v1/fit_v6/histograms.npz :: experiments/AOJ/sideband_residuals.py: max_abs_pull (max over the pretrained models) :: sha256 2bd1e3d5b32d4a12 :: largest single-bin pull below the top window, any pretrained model
\newcommand{\AojSharedBinLo}{125}  % experiments/FIGS/data/aoj_full_v1/fit_v6/histograms.npz :: the bin holding the most models' largest pull :: sha256 2bd1e3d5b32d4a12 :: GeV
\newcommand{\AojSharedBinHi}{130}  % experiments/FIGS/data/aoj_full_v1/fit_v6/histograms.npz :: its upper edge :: sha256 2bd1e3d5b32d4a12 :: GeV
\newcommand{\AojSharedBinNModels}{18}  % experiments/FIGS/data/aoj_full_v1/fit_v6/histograms.npz :: models whose largest pull is in that bin :: sha256 2bd1e3d5b32d4a12
\newcommand{\AojSharedBinMeanPull}{+1.4}  % experiments/FIGS/data/aoj_full_v1/fit_v6/histograms.npz :: experiments/AOJ/sideband_residuals.py mean_pulls: that bin :: sha256 2bd1e3d5b32d4a12 :: mean pull over the pretrained models
\newcommand{\AojSharedBinRank}{third}  % experiments/FIGS/data/aoj_full_v1/fit_v6/histograms.npz :: experiments/AOJ/sideband_residuals.py mean_pulls: rank of that bin's |mean pull| :: sha256 2bd1e3d5b32d4a12
\newcommand{\AojFitNBins}{39}  % experiments/FIGS/data/aoj_full_v1/fit_v6/histograms.npz :: mass bins of the fit range :: sha256 2bd1e3d5b32d4a12
\newcommand{\AojPairedYieldOnesixtwoOverOneeighteight}{0.94~[0.45, 1.95]}  % experiments/FIGS/data/aoj_full_v1/analysis_v6/aoj_top.json :: per_label_set.label_sets.{162,188}.signal_yields :: sha256 5dfef902a5fe75ba :: yield ratio paired by run, geometric mean [95% Student-t interval]
\newcommand{\AojPairedYieldFourthreeOverOneeighteight}{0.97~[0.81, 1.16]}  % experiments/FIGS/data/aoj_full_v1/analysis_v6/aoj_top.json :: per_label_set.label_sets.{43,188}.signal_yields :: sha256 5dfef902a5fe75ba :: yield ratio paired by run, geometric mean [95% Student-t interval]
\newcommand{\AojPairedYieldOnesevenOverOneeighteight}{1.01~[0.75, 1.35]}  % experiments/FIGS/data/aoj_full_v1/analysis_v6/aoj_top.json :: per_label_set.label_sets.{17,188}.signal_yields :: sha256 5dfef902a5fe75ba :: yield ratio paired by run, geometric mean [95% Student-t interval]
\newcommand{\AojPairedYieldOnesevenOverOnesixtwo}{1.07~[0.51, 2.24]}  % experiments/FIGS/data/aoj_full_v1/analysis_v6/aoj_top.json :: per_label_set.label_sets.{17,162}.signal_yields :: sha256 5dfef902a5fe75ba :: yield ratio paired by run, geometric mean [95% Student-t interval]
\newcommand{\AojPairedYieldOnesixtwoMassOverOnesixtwo}{1.10~[0.57, 2.14]}  % experiments/FIGS/data/aoj_full_v1/analysis_v6/aoj_top.json :: per_label_set.label_sets.{162+mass,162}.signal_yields :: sha256 5dfef902a5fe75ba :: yield ratio paired by run, geometric mean [95% Student-t interval]
\newcommand{\AojPairedYieldOnesevenMassOverOneseven}{1.05~[0.74, 1.49]}  % experiments/FIGS/data/aoj_full_v1/analysis_v6/aoj_top.json :: per_label_set.label_sets.{17+mass,17}.signal_yields :: sha256 5dfef902a5fe75ba :: yield ratio paired by run, geometric mean [95% Student-t interval]
\newcommand{\AojLeakPairedShiftMax}{2\%}  % experiments/FIGS/data/aoj_full_v1/fit_v6/results.json :: models.*.top.leak_systematic.signal_yield against signal_yield, paired ratios :: sha256 2416fc21f93fcfa6 :: largest change of a paired yield ratio when the fail-region tops are scaled up
\newcommand{\AojRefEffPercent}{1.38\%}  % experiments/FIGS/data/aoj_full_v1/fit_v6/results.json :: reference.top.data_efficiency :: sha256 2416fc21f93fcfa6 :: fraction of all data jets passing its cut
\newcommand{\AojRefFitP}{\ensuremath{2.56\times10^{-11}}}  % experiments/FIGS/data/aoj_full_v1/fit_v6/results.json :: reference.top.asymptotic_p :: sha256 2416fc21f93fcfa6 :: goodness of fit of the reference's own top fit
\newcommand{\AojRefValidationP}{\ensuremath{p \leq 0.005} (0 of 200 toys reach the observed deviance)}  % experiments/FIGS/data/aoj_full_v1/fit_v6/results.json :: reference.top.validation.toy_p, n_toys :: sha256 2416fc21f93fcfa6 :: toy p-value of the background validation, (k + 1) / (n + 1)
\newcommand{\AojRefBandSignalZ}{2.5}  % experiments/FIGS/data/aoj_full_v1/fit_v6/results.json :: reference.top.validation.band_signal_z :: sha256 2416fc21f93fcfa6 :: significance of a peak of the score's own shape in the validation band
\newcommand{\AojModelEffMin}{1.01\%}  % experiments/FIGS/data/aoj_full_v1/fit_v6/results.json :: models.*.top.data_efficiency (min) :: sha256 2416fc21f93fcfa6 :: pretrained models
\newcommand{\AojModelEffMax}{1.05\%}  % experiments/FIGS/data/aoj_full_v1/fit_v6/results.json :: models.*.top.data_efficiency (max) :: sha256 2416fc21f93fcfa6 :: pretrained models
\newcommand{\AojModelsValidationPMin}{0.11}  % experiments/FIGS/data/aoj_full_v1/fit_v6/results.json :: models.*.top.validation.toy_p (min) :: sha256 2416fc21f93fcfa6 :: pretrained models
\newcommand{\AojRefYield}{29{,}143}  % experiments/FIGS/data/aoj_full_v1/fit_v6/results.json :: reference.top.signal_yield :: sha256 2416fc21f93fcfa6
\newcommand{\AojRefYieldErr}{865}  % experiments/FIGS/data/aoj_full_v1/fit_v6/results.json :: reference.top.signal_yield_err :: sha256 2416fc21f93fcfa6
\newcommand{\AojRefMass}{181.0}  % experiments/FIGS/data/aoj_full_v1/fit_v6/results.json :: reference.top.floated_mean :: sha256 2416fc21f93fcfa6 :: GeV
\newcommand{\AojRefWidth}{17.8}  % experiments/FIGS/data/aoj_full_v1/fit_v6/results.json :: reference.top.floated_width :: sha256 2416fc21f93fcfa6 :: GeV
\newcommand{\AojPublicYield}{1{,}813}  % experiments/FIGS/data/aoj_full_v1/fit_v6/results.json :: models.sophon-public.top.signal_yield :: sha256 2416fc21f93fcfa6
\newcommand{\AojPublicYieldErr}{193}  % experiments/FIGS/data/aoj_full_v1/fit_v6/results.json :: models.sophon-public.top.signal_yield_err :: sha256 2416fc21f93fcfa6
\newcommand{\AojPublicMass}{182.8}  % experiments/FIGS/data/aoj_full_v1/fit_v6/results.json :: models.sophon-public.top.floated_mean :: sha256 2416fc21f93fcfa6 :: GeV
\newcommand{\AojPublicWidth}{11.7}  % experiments/FIGS/data/aoj_full_v1/fit_v6/results.json :: models.sophon-public.top.floated_width :: sha256 2416fc21f93fcfa6 :: GeV
\newcommand{\AojInjSizeMin}{1{,}000}  % experiments/FIGS/data/aoj_injection_v1/summary.json :: groups[region=top,mode=ensemble].size (min) :: sha256 d9b66466ae24b584 :: injected tops
\newcommand{\AojInjSizeMax}{4{,}000}  % experiments/FIGS/data/aoj_injection_v1/summary.json :: groups[region=top,mode=ensemble].size (max) :: sha256 d9b66466ae24b584
\newcommand{\AojInjFloatPullSdMax}{2.0}  % experiments/FIGS/data/aoj_injection_v1/summary.json :: groups[region=top,mode=bootstrap,variant=float].pull.sd (max) :: sha256 d9b66466ae24b584 :: pull width with each model's own floated shape
\newcommand{\AojInjEnsembleRecoveryMin}{0.983}  % experiments/FIGS/data/aoj_injection_v1/summary.json :: groups[region=top,mode=ensemble].ratio (min) :: sha256 d9b66466ae24b584 :: the whole procedure, shape re-derived per toy set
\newcommand{\AojInjEnsembleRecoveryMax}{0.998}  % experiments/FIGS/data/aoj_injection_v1/summary.json :: groups[region=top,mode=ensemble].ratio (max) :: sha256 d9b66466ae24b584
\newcommand{\AojInjEnsemblePullSdMin}{1.02}  % experiments/FIGS/data/aoj_injection_v1/summary.json :: groups[region=top,mode=ensemble].pull (min) :: sha256 d9b66466ae24b584 :: the whole procedure, shape re-derived per toy set
\newcommand{\AojInjEnsemblePullSdMax}{1.09}  % experiments/FIGS/data/aoj_injection_v1/summary.json :: groups[region=top,mode=ensemble].pull (max) :: sha256 d9b66466ae24b584
\newcommand{\AojClosureNShards}{10}  % experiments/FIGS/data/aoj_checks_v2/closure_all_shards.json :: n_shards :: sha256 6a6018351d7048d2 :: staged shards of the open data
\newcommand{\AojClosurePtLo}{500}  % experiments/FIGS/data/aoj_checks_v2/closure_all_shards.json :: pt_window[0] :: sha256 6a6018351d7048d2 :: GeV
\newcommand{\AojClosurePtHi}{700}  % experiments/FIGS/data/aoj_checks_v2/closure_all_shards.json :: pt_window[1] :: sha256 6a6018351d7048d2 :: GeV
\newcommand{\AojClosureDzeroIqr}{8.7}  % experiments/FIGS/data/aoj_checks_v2/closure_all_shards.json :: domain_shift.part_d0.iqr_ratio :: sha256 6a6018351d7048d2 :: open data over JetClass-II QCD, pooled over the shards
\newcommand{\AojClosureDzIqr}{10.6}  % experiments/FIGS/data/aoj_checks_v2/closure_all_shards.json :: domain_shift.part_dz.iqr_ratio :: sha256 6a6018351d7048d2 :: open data over JetClass-II QCD, pooled over the shards
\newcommand{\AojClosureNeutral}{0.65}  % experiments/FIGS/data/aoj_checks_v2/closure_all_shards.json :: features.part_isNeutralHadron.ratio.per_shard (mean) :: sha256 6a6018351d7048d2 :: open data over JetClass-II QCD
\newcommand{\AojClosureCharged}{1.22}  % experiments/FIGS/data/aoj_checks_v2/closure_all_shards.json :: features.part_isChargedHadron.ratio.per_shard (mean) :: sha256 6a6018351d7048d2 :: open data over JetClass-II QCD
\newcommand{\AojClosureNPartAoj}{61}  % experiments/FIGS/data/aoj_checks_v2/closure_all_shards.json :: features.n_particles.aoj.per_shard (mean of the shard medians) :: sha256 6a6018351d7048d2 :: median
\newcommand{\AojClosureNPartRef}{48}  % experiments/FIGS/data/aoj_checks_v2/closure_all_shards.json :: features.n_particles.reference.per_shard (mean) :: sha256 6a6018351d7048d2 :: median
\newcommand{\RandMergingDrawsBvcFourprong}{three}  % configs/labelmaps/rand_label_map.v1.csv :: RAND_d* of label_X_YY_bbqq and label_X_YY_ccqq :: sha256 f395a25f52e3cfad :: first-grid partitions that merge the task's pair
\newcommand{\RandSplittingDrawsBvcFourprong}{one and two}  % configs/labelmaps/rand_label_map.v1.csv :: RAND_d* of label_X_YY_bbqq and label_X_YY_ccqq :: sha256 f395a25f52e3cfad :: first-grid partitions that keep the pair apart
\newcommand{\RandMergingDrawsVisibleContent}{one and three}  % configs/labelmaps/rand_label_map.v1.csv :: RAND_d* of label_X_YY_bbqq and label_X_YY_cqtauhv :: sha256 f395a25f52e3cfad :: first-grid partitions that merge the task's pair
\newcommand{\RandSplittingDrawsVisibleContent}{two}  % configs/labelmaps/rand_label_map.v1.csv :: RAND_d* of label_X_YY_bbqq and label_X_YY_cqtauhv :: sha256 f395a25f52e3cfad :: first-grid partitions that keep the pair apart
\newcommand{\RandVtwoNPartitions}{five}  % configs/arms/v2_grid.json :: arms[RAND2_*] (count) :: sha256 d8a31b1486992302 :: random partitions in the rerun
\newcommand{\RandVtwoRuns}{one}  % configs/arms/v2_paper_runs.txt :: runs of RAND2_* :: sha256 b04562b14ea8cbc7 :: runs per partition
\newcommand{\RandFlavRuns}{two}  % configs/arms/v2_paper_runs.txt :: runs of FLAV_* :: sha256 b04562b14ea8cbc7 :: runs per flavour-pair model
\newcommand{\RandVtwoNPairs}{seven}  % configs/labelmaps/rand_v2_selection.json :: pairs (count) :: sha256 98cf52d49c7d2f02 :: probe pairs the balance rule covers
\newcommand{\RandVtwoMergedMin}{two}  % configs/labelmaps/rand_v2_selection.json :: merged_in :: sha256 98cf52d49c7d2f02 :: partitions merging each probe pair, checked on the map
\newcommand{\RandVtwoMergedMax}{three}  % configs/labelmaps/rand_v2_selection.json :: merged_in :: sha256 98cf52d49c7d2f02 :: partitions merging each probe pair, checked on the map
\newcommand{\RandVtwoShareTol}{5\%}  % configs/labelmaps/rand_v2_selection.json :: realised_tolerance :: sha256 98cf52d49c7d2f02 :: largest relative deviation of a realised group share from its target (rule 3)
\newcommand{\RandVtwoTrials}{32}  % scripts/build_rand_control.py :: TRIALS :: sha256 5bc91d77a3f2e81f :: share-matched assignments drawn per random partition; the one sharing the fewest merged pairs with the 17-class label set is kept
\newcommand{\RandVtwoPoolDraws}{6{,}600}  % configs/labelmaps/rand_v2_selection.json :: pool.draws :: sha256 98cf52d49c7d2f02 :: candidate partitions drawn (seeds) before the realised-share rule
\newcommand{\RandVtwoPoolPass}{27}  % configs/labelmaps/rand_v2_selection.json :: rule3_seeds (count) :: sha256 98cf52d49c7d2f02 :: candidate partitions meeting the realised-share rule (rule 3)
\newcommand{\RandFlavSplitOrbit}{\ensuremath{X\to YY\to QQQQ}}  % configs/labelmaps/flavour_pair.v2.json :: split_orbit :: sha256 3b79ce5e44c1e2f4 :: the decays F1 cuts, Q any quark
\newcommand{\RandFlavSplitOrbitClasses}{22}  % configs/labelmaps/flavour_pair_map.v2.csv :: classes whose orbit is split_orbit (count) :: sha256 9aab235099ceac43 :: classes of the split orbit
\newcommand{\RandFlavCutClasses}{11}  % configs/labelmaps/flavour_pair.v2.json :: split_classes_moved (count) :: sha256 3b79ce5e44c1e2f4 :: classes of the split orbit that F1 moves: those with a b quark
\newcommand{\RandFlavOrbitsMoved}{three}  % configs/labelmaps/flavour_pair.v2.json :: orbits_moved_B_to_A (count) :: sha256 3b79ce5e44c1e2f4 :: whole orbits F1 moves the other way
\newcommand{\RandFlavOrbitsMovedList}{\ensuremath{X\to YY\to QQ\mu\mu}, \ensuremath{X\to YY\to QQ\tau_h\tau_h} and \ensuremath{X\to \mu\mu}}  % configs/labelmaps/flavour_pair.v2.json :: orbits_moved_B_to_A :: sha256 3b79ce5e44c1e2f4 :: the orbits F1 moves the other way, Q any quark
\newcommand{\RandFlavOrbitClassesMoved}{9}  % configs/labelmaps/flavour_pair_map.v2.csv :: classes whose orbit is in orbits_moved_B_to_A :: sha256 9aab235099ceac43 :: classes in those orbits
\newcommand{\RandFlavFonerCutWithB}{6}  % configs/labelmaps/flavour_pair_map.v2.csv :: FLAV_F1R != FLAV_F0 within split_orbit, with a b quark (count) :: sha256 9aab235099ceac43 :: classes with a b quark among those the F1r cut moves
\newcommand{\RandFlavBcPairs}{11}  % configs/labelmaps/flavour_pair.v2.json :: b_c_pairs_in_split_orbit (count) :: sha256 3b79ce5e44c1e2f4 :: pairs of the split orbit's classes that differ only by exchanging b and c
\newcommand{\RandFlavFoneBcSplit}{8}  % configs/labelmaps/flavour_pair_map.v2.csv :: b_c_pairs_in_split_orbit split by FLAV_F1 (count) :: sha256 9aab235099ceac43
\newcommand{\RandFlavFonerBcSplit}{none}  % configs/labelmaps/flavour_pair_map.v2.csv :: b_c_pairs_in_split_orbit split by FLAV_F1R (count) :: sha256 9aab235099ceac43 :: PRESPEC A14: none
\newcommand{\RandFlavFonerShareTol}{1\%}  % configs/labelmaps/flavour_pair.v2.json :: F1R.realised_tolerance_vs_F1 :: sha256 3b79ce5e44c1e2f4 :: largest allowed relative difference of a group's realised share from F1's
\newcommand{\RandFlavFonerShareDev}{0.77\%}  % configs/labelmaps/realised_native_shares.v2.json :: share summed per FLAV_F1R and FLAV_F1 group of flavour_pair_map.v2.csv (largest relative difference) :: sha256 5728c12e88fa3683 :: realised shares from the final loader's dry run
\newcommand{\MlstNRunsWord}{three}  % configs/arms/v2_paper_runs.txt :: runs of L188/L162/R42_Q1/R16_Q1 :: sha256 b04562b14ea8cbc7 :: pretraining runs per configuration
\newcommand{\MlstMassNRunsWord}{two}  % configs/arms/v2_paper_runs.txt :: runs of L162_MASS/R16_Q1_MASS :: sha256 b04562b14ea8cbc7 :: pretraining runs per configuration
\newcommand{\LpCMin}{\ensuremath{1\times10^{-3}}}  % experiments/EVAL/linear_probe_v2.py :: C_GRID :: sha256 4c576b44812c18a2 :: smallest inverse regularisation
\newcommand{\LpCMax}{100}  % experiments/EVAL/linear_probe_v2.py :: C_GRID :: sha256 4c576b44812c18a2 :: largest inverse regularisation
\newcommand{\LpNTestJcii}{2{,}000{,}000}  % experiments/EVAL/linear_probe_v2.py :: TEST_FIRST.jc2 :: sha256 4c576b44812c18a2 :: JetClass-II test jets
\newcommand{\LpNTestAucJcii}{500{,}000}  % experiments/EVAL/linear_probe_v2.py :: TEST_FIRST.jc2 / AUC_STRIDE :: sha256 4c576b44812c18a2 :: JetClass-II test jets of the macro AUC
\newcommand{\LpConvergedGrad}{\ensuremath{1\times10^{-5}}}  % experiments/EVAL/linear_probe_v2.py :: CONVERGED_GRAD :: sha256 4c576b44812c18a2 :: largest gradient component a converged probe ends below
\newcommand{\FillNFeatures}{17}  % configs/finetune/TopReference.yaml :: inputs.pf_features.vars (count) :: sha256 501a6ffe5d0a8051 :: per-particle input features of the network
\newcommand{\FillTopNZero}{10}  % configs/finetune/TopReference.yaml :: inputs.pf_features.vars, part_zero_* (count) :: sha256 501a6ffe5d0a8051 :: features the data set lacks, set to zero
\newcommand{\FillTopNPresent}{7}  % configs/finetune/TopReference.yaml :: inputs.pf_features.vars, not part_zero_* (count) :: sha256 501a6ffe5d0a8051 :: features the data set provides
\newcommand{\FillQgNZero}{4}  % configs/finetune/EnergyFlowQG.yaml :: inputs.pf_features.vars, part_zero_* (count) :: sha256 db42d6471e8b865c :: features the data set lacks, set to zero
\newcommand{\FillQgNPresent}{13}  % configs/finetune/EnergyFlowQG.yaml :: inputs.pf_features.vars, not part_zero_* (count) :: sha256 db42d6471e8b865c :: features the data set provides
\newcommand{\DesignVtwoDataFraction}{20\%}  % experiments/MTX/k8s/v2/grid/job-mtx2-l188-s1-raunav.yaml :: --data-fraction :: sha256 d28ce02267a6b87a :: share of every file read per epoch
\newcommand{\DesignVtwoNValFiles}{25}  % experiments/MTX/k8s/v2/grid/job-mtx2-l188-s1-raunav.yaml :: --data-val (files, brace ranges expanded) :: sha256 d28ce02267a6b87a :: validation files of the fixed sample
\newcommand{\DesignWavgFirstEpoch}{70}  % experiments/MTX/pretrain_v2.py :: --num-epochs - WAVG_EPOCHS :: sha256 95da5a6e029855ee :: first epoch of the weight average (epochs count from 0)
\newcommand{\DesignWavgLastEpoch}{79}  % experiments/MTX/k8s/v2/grid/job-mtx2-l188-s1-raunav.yaml :: --num-epochs - 1 :: sha256 d28ce02267a6b87a :: last epoch
\newcommand{\DesignBnJets}{200{,}000}  % experiments/MTX/pretrain_v2.py :: BN_JETS :: sha256 95da5a6e029855ee :: training jets the BatchNorm statistics are recomputed on
\newcommand{\DesignVtwoNWindows}{five}  % experiments/MTX/k8s/v2/grid/job-mtx2-l188-s1-raunav.yaml :: 1 / --data-fraction :: sha256 d28ce02267a6b87a :: disjoint windows of each file's rows, one read per epoch
\newcommand{\DesignLofoNWindows}{three}  % experiments/MTX/k8s/v2/cut/grid/job-mtx2-l188lofo4p-s1-raunav.yaml :: --data-windows :: sha256 44a9573929efb7ac :: windows of the models trained without one family
\newcommand{\DesignVtwoPlanUsed}{97\%}  % experiments/FIGS/data/v2_loader/loader_dryrun/dryrun_s176_seed1.json :: epochs[*].(max_fetch_id + 1) / args.data_split_num, mean :: sha256 03b6586b26ad7a2f :: share of an epoch's planned reading used before its jets are reached, 20 epochs
\newcommand{\DesignLofoPlanUsed}{76\%}  % experiments/FIGS/data/v2_loader/loader_dryrun/dryrun_lofo4p_seed1.json :: epochs[*].(max_fetch_id + 1) / fetches_per_pass, mean :: sha256 2a1315e611d13ff5 :: the same, leaving the family out
\newcommand{\DesignLofoRowFactor}{1.29}  % experiments/FIGS/data/v2_loader/loader_dryrun/dryrun_lofo4p_seed1.json :: (mean plan used / --data-windows) over (mean plan used / k) of the full models' dry run :: sha256 2a1315e611d13ff5 :: how much more often a row of a file is read per epoch when the family is left out
\newcommand{\DesignDistinctJetsCycle}{\ensuremath{3.3\times10^{7}}}  % experiments/FIGS/data/v2_loader/loader_dryrun/dryrun_s176_seed1.json :: sum of epochs[c*k:(c+1)*k].distinct_jets, mean over 4 cycles :: sha256 03b6586b26ad7a2f :: distinct jets the full models train on in one cycle of windows
\newcommand{\DesignLofoFamilyShare}{22.4\%}  % experiments/FIGS/data/v2_loader/loader_dryrun/dryrun_s176_seed1.json :: epochs[*].native_counts over summary.labels_absent_every_epoch of the family-out run :: sha256 03b6586b26ad7a2f :: the left-out family's share of the full models' stream
\newcommand{\DesignLofoNClasses}{27}  % experiments/FIGS/data/v2_loader/loader_dryrun/dryrun_lofo4p_seed1.json :: summary.labels_absent_every_epoch (count) :: sha256 2a1315e611d13ff5 :: native classes of the left-out family
\newcommand{\DesignLofoExposureFactor}{1.29}  % experiments/FIGS/data/v2_loader/loader_dryrun/dryrun_s176_seed1.json :: 1 / (1 - family share) :: sha256 03b6586b26ad7a2f :: how much more often every other class is seen
\newcommand{\DesignLofoQcdShareFull}{14.4\%}  % experiments/FIGS/data/v2_loader/loader_dryrun/dryrun_s176_seed1.json :: summary.mean_qcd_share :: sha256 03b6586b26ad7a2f :: QCD share of the realised stream
\newcommand{\DesignLofoQcdShareLofo}{18.6\%}  % experiments/FIGS/data/v2_loader/loader_dryrun/dryrun_lofo4p_seed1.json :: summary.mean_qcd_share :: sha256 2a1315e611d13ff5 :: QCD share of the realised stream
\newcommand{\DesignLrFlatPercent}{70\%}  % experiments/MTX/pretrain_v2.py :: make_optimizer: 1 - the decay share :: sha256 95da5a6e029855ee :: share of the epochs at the initial rate
\newcommand{\DesignLrDecayFirstEpoch}{56}  % experiments/MTX/pretrain_v2.py :: num_epochs - int(num_epochs * decay share) :: sha256 95da5a6e029855ee :: first epoch at a reduced rate (epochs from 0)
\newcommand{\DesignLrEndPercent}{1\%}  % experiments/MTX/pretrain_v2.py :: make_optimizer: gamma ** n_decay :: sha256 95da5a6e029855ee :: the last epoch's rate over the initial rate
\newcommand{\DesignEndLr}{\ensuremath{5\times10^{-6}}}  % experiments/MTX/k8s/v2/grid/job-mtx2-l188-s1-raunav.yaml :: --start-lr x the final factor :: sha256 d28ce02267a6b87a :: the rate of the last epoch
\newcommand{\DesignWeaverVersion}{0.4.17}  % experiments/MTX/pretrain_v2.py :: the weaver release whose loop it reproduces :: sha256 95da5a6e029855ee
\newcommand{\DesignMeiFactor}{1.1}  % src/stats/mde.py :: MEI_LOG = ln(factor) :: sha256 626468dcc37940bb :: the smallest effect of interest, as a factor
\newcommand{\DesignVtwoGpu}{RTX 3090}  % experiments/MTX/k8s/v2/grid/job-mtx2-flavf0-s1-raunav.yaml :: nodeAffinity nvidia.com/gpu.product :: sha256 960cad1691314980 :: the product of run index 1, every configuration
\newcommand{\DesignVtwoSecondGpu}{\pending{pending: the GPU type of each run index}}  % experiments/FIGS/make_tables.py :: emit_v2_recipe[DesignVtwoSecondGpu] :: sha256 10e6ec08356a4a2b :: a marker, not a number: the grid's job specs put every run index on one product
\newcommand{\DesignVtwoSecondGpuRuns}{\pending{pending: the GPU type of each run index}}  % experiments/FIGS/make_tables.py :: emit_v2_recipe[DesignVtwoSecondGpuRuns] :: sha256 10e6ec08356a4a2b :: a marker, not a number: the grid's job specs put every run index on one product
\newcommand{\MassNRunsWord}{five}  % experiments/FIGS/data/probe_ladder_mass2x2_mlp2/s1.json :: per-seed files (count) :: sha256 906e778d7043fe67 :: pretraining runs of each mass-output configuration
\newcommand{\MassOmaLinearOnesixtwo}{\ensuremath{(2.36 \pm 0.06)\times10^{-3}}}  % experiments/FIGS/data/probe_ladder_v2_mlp2/analysis/seed_level_results.json :: table[bvc_resonant,linear,162].auc :: sha256 53e1ec81343ac2c1 :: 1-AUC, mean +- SD over 5 seeds
\newcommand{\MassOmaLinearOnesixtwoMass}{\ensuremath{(2.62 \pm 0.08)\times10^{-3}}}  % experiments/FIGS/data/probe_ladder_mass2x2_mlp2/s1.json :: tasks.bvc_resonant.arms.l162mass-s*.linear.auc :: sha256 906e778d7043fe67 :: 1-AUC, mean +- SD over the 5 per-seed files
\newcommand{\MassOmaRatioLinearOnesixtwo}{1.11}  % experiments/FIGS/data/probe_ladder_mass2x2_mlp2/s1.json :: tasks.bvc_resonant.arms.l162mass-s*.linear.auc over table[bvc_resonant,linear,162].auc :: sha256 906e778d7043fe67 :: seed mean of 1-AUC with the mass output over without
\newcommand{\MassOmaLinearOneseven}{\ensuremath{(9.3 \pm 1.5)\times10^{-3}}}  % experiments/FIGS/data/probe_ladder_v2_mlp2/analysis/seed_level_results.json :: table[bvc_resonant,linear,17].auc :: sha256 53e1ec81343ac2c1 :: 1-AUC, mean +- SD over 5 seeds
\newcommand{\MassOmaLinearOnesevenMass}{\ensuremath{(3.0 \pm 1.2)\times10^{-2}}}  % experiments/FIGS/data/probe_ladder_mass2x2_mlp2/s1.json :: tasks.bvc_resonant.arms.r16q1mass-s*.linear.auc :: sha256 906e778d7043fe67 :: 1-AUC, mean +- SD over the 5 per-seed files
\newcommand{\MassOmaRatioLinearOneseven}{3.2}  % experiments/FIGS/data/probe_ladder_mass2x2_mlp2/s1.json :: tasks.bvc_resonant.arms.r16q1mass-s*.linear.auc over table[bvc_resonant,linear,17].auc :: sha256 906e778d7043fe67 :: seed mean of 1-AUC with the mass output over without
\newcommand{\MassOmaMlpOnesixtwo}{\ensuremath{(2.32 \pm 0.05)\times10^{-3}}}  % experiments/FIGS/data/probe_ladder_v2_mlp2/analysis/seed_level_results.json :: table[bvc_resonant,mlp,162].auc :: sha256 53e1ec81343ac2c1 :: 1-AUC, mean +- SD over 5 seeds
\newcommand{\MassOmaMlpOnesixtwoMass}{\ensuremath{(2.52 \pm 0.06)\times10^{-3}}}  % experiments/FIGS/data/probe_ladder_mass2x2_mlp2/s1.json :: tasks.bvc_resonant.arms.l162mass-s*.mlp.auc :: sha256 906e778d7043fe67 :: 1-AUC, mean +- SD over the 5 per-seed files
\newcommand{\MassOmaRatioMlpOnesixtwo}{1.09}  % experiments/FIGS/data/probe_ladder_mass2x2_mlp2/s1.json :: tasks.bvc_resonant.arms.l162mass-s*.mlp.auc over table[bvc_resonant,mlp,162].auc :: sha256 906e778d7043fe67 :: seed mean of 1-AUC with the mass output over without
\newcommand{\MassOmaMlpOneseven}{\ensuremath{(7.3 \pm 0.9)\times10^{-3}}}  % experiments/FIGS/data/probe_ladder_v2_mlp2/analysis/seed_level_results.json :: table[bvc_resonant,mlp,17].auc :: sha256 53e1ec81343ac2c1 :: 1-AUC, mean +- SD over 5 seeds
\newcommand{\MassOmaMlpOnesevenMass}{\ensuremath{(1.5 \pm 0.4)\times10^{-2}}}  % experiments/FIGS/data/probe_ladder_mass2x2_mlp2/s1.json :: tasks.bvc_resonant.arms.r16q1mass-s*.mlp.auc :: sha256 906e778d7043fe67 :: 1-AUC, mean +- SD over the 5 per-seed files
\newcommand{\MassOmaRatioMlpOneseven}{2.0}  % experiments/FIGS/data/probe_ladder_mass2x2_mlp2/s1.json :: tasks.bvc_resonant.arms.r16q1mass-s*.mlp.auc over table[bvc_resonant,mlp,17].auc :: sha256 906e778d7043fe67 :: seed mean of 1-AUC with the mass output over without
\newcommand{\FtAucJciiEThreeOneeighteight}{0.902\,\ensuremath{\pm}\,0.004}  % experiments/FIGS/data/w2b_leg1_metrics_v2.json :: cells.{l188-s1,l188-s2,l188-s3,l188-s4,l188-s5}.N1000.s1.macro_auc_ovr :: sha256 b56bf70be9982044 :: mean +- SD over the row's 5 pretrained models
\newcommand{\FtAucJciiEFourOneeighteight}{0.9723\,\ensuremath{\pm}\,0.0004}  % experiments/FIGS/data/w2b_leg1_metrics_v2.json :: cells.{l188-s1,l188-s2,l188-s3,l188-s4,l188-s5}.N10000.s1.macro_auc_ovr :: sha256 b56bf70be9982044 :: mean +- SD over the row's 5 pretrained models
\newcommand{\FtAucJciiEFiveOneeighteight}{0.9793\,\ensuremath{\pm}\,0.0002}  % experiments/FIGS/data/w2b_leg1_metrics_v2.json :: cells.{l188-s1,l188-s2,l188-s3,l188-s4,l188-s5}.N100000.s1.macro_auc_ovr :: sha256 b56bf70be9982044 :: mean +- SD over the row's 5 pretrained models
\newcommand{\FtAucJciiESixOneeighteight}{0.9809\,\ensuremath{\pm}\,\ensuremath{<}0.0001}  % experiments/FIGS/data/w2b_leg1_metrics_v2.json :: cells.{l188-s1,l188-s2,l188-s3,l188-s4,l188-s5}.N1000000.s1.macro_auc_ovr :: sha256 b56bf70be9982044 :: mean +- SD over the row's 5 pretrained models
\newcommand{\FtAucJciiEThreeOnesixtwo}{0.906\,\ensuremath{\pm}\,0.007}  % experiments/FIGS/data/w2b_leg1_metrics_v2.json :: cells.{l162-s1b,l162-s2,l162-s3,l162-s4,l162-s5}.N1000.s1.macro_auc_ovr :: sha256 b56bf70be9982044 :: mean +- SD over the row's 5 pretrained models
\newcommand{\FtAucJciiEFourOnesixtwo}{0.9724\,\ensuremath{\pm}\,0.0005}  % experiments/FIGS/data/w2b_leg1_metrics_v2.json :: cells.{l162-s1b,l162-s2,l162-s3,l162-s4,l162-s5}.N10000.s1.macro_auc_ovr :: sha256 b56bf70be9982044 :: mean +- SD over the row's 5 pretrained models
\newcommand{\FtAucJciiEFiveOnesixtwo}{0.9794\,\ensuremath{\pm}\,0.0002}  % experiments/FIGS/data/w2b_leg1_metrics_v2.json :: cells.{l162-s1b,l162-s2,l162-s3,l162-s4,l162-s5}.N100000.s1.macro_auc_ovr :: sha256 b56bf70be9982044 :: mean +- SD over the row's 5 pretrained models
\newcommand{\FtAucJciiESixOnesixtwo}{0.9808\,\ensuremath{\pm}\,0.0001}  % experiments/FIGS/data/w2b_leg1_metrics_v2.json :: cells.{l162-s1b,l162-s2,l162-s3,l162-s4,l162-s5}.N1000000.s1.macro_auc_ovr :: sha256 b56bf70be9982044 :: mean +- SD over the row's 5 pretrained models
\newcommand{\FtAucJciiEThreeFourthree}{0.918\,\ensuremath{\pm}\,0.005}  % experiments/FIGS/data/w2b_leg1_metrics_v2.json :: cells.{r42q1-s1,r42q1-s2,r42q1-s3,r42q1-s4,r42q1-s5}.N1000.s1.macro_auc_ovr :: sha256 b56bf70be9982044 :: mean +- SD over the row's 5 pretrained models
\newcommand{\FtAucJciiEFourFourthree}{0.9667\,\ensuremath{\pm}\,0.0005}  % experiments/FIGS/data/w2b_leg1_metrics_v2.json :: cells.{r42q1-s1,r42q1-s2,r42q1-s3,r42q1-s4,r42q1-s5}.N10000.s1.macro_auc_ovr :: sha256 b56bf70be9982044 :: mean +- SD over the row's 5 pretrained models
\newcommand{\FtAucJciiEFiveFourthree}{0.9775\,\ensuremath{\pm}\,0.0003}  % experiments/FIGS/data/w2b_leg1_metrics_v2.json :: cells.{r42q1-s1,r42q1-s2,r42q1-s3,r42q1-s4,r42q1-s5}.N100000.s1.macro_auc_ovr :: sha256 b56bf70be9982044 :: mean +- SD over the row's 5 pretrained models
\newcommand{\FtAucJciiESixFourthree}{0.9800\,\ensuremath{\pm}\,0.0001}  % experiments/FIGS/data/w2b_leg1_metrics_v2.json :: cells.{r42q1-s1,r42q1-s2,r42q1-s3,r42q1-s4,r42q1-s5}.N1000000.s1.macro_auc_ovr :: sha256 b56bf70be9982044 :: mean +- SD over the row's 5 pretrained models
\newcommand{\FtAucJciiEThreeOneseven}{0.905\,\ensuremath{\pm}\,0.008}  % experiments/FIGS/data/w2b_leg1_metrics_v2.json :: cells.{r16q1-s1,r16q1-s2,r16q1-s3,r16q1-s4,r16q1-s5}.N1000.s1.macro_auc_ovr :: sha256 b56bf70be9982044 :: mean +- SD over the row's 5 pretrained models
\newcommand{\FtAucJciiEFourOneseven}{0.9635\,\ensuremath{\pm}\,0.0005}  % experiments/FIGS/data/w2b_leg1_metrics_v2.json :: cells.{r16q1-s1,r16q1-s2,r16q1-s3,r16q1-s4,r16q1-s5}.N10000.s1.macro_auc_ovr :: sha256 b56bf70be9982044 :: mean +- SD over the row's 5 pretrained models
\newcommand{\FtAucJciiEFiveOneseven}{0.9748\,\ensuremath{\pm}\,0.0003}  % experiments/FIGS/data/w2b_leg1_metrics_v2.json :: cells.{r16q1-s1,r16q1-s2,r16q1-s3,r16q1-s4,r16q1-s5}.N100000.s1.macro_auc_ovr :: sha256 b56bf70be9982044 :: mean +- SD over the row's 5 pretrained models
\newcommand{\FtAucJciiESixOneseven}{0.9781\,\ensuremath{\pm}\,0.0001}  % experiments/FIGS/data/w2b_leg1_metrics_v2.json :: cells.{r16q1-s1,r16q1-s2,r16q1-s3,r16q1-s4,r16q1-s5}.N1000000.s1.macro_auc_ovr :: sha256 b56bf70be9982044 :: mean +- SD over the row's 5 pretrained models
\newcommand{\FtAucJciiEThreeOnesixtwoMass}{0.921\,\ensuremath{\pm}\,0.005}  % experiments/FIGS/data/w2b_leg1_metrics_v2.json :: cells.{l162mass-s1,l162mass-s2,l162mass-s3,l162mass-s4,l162mass-s5}.N1000.s1.macro_auc_ovr :: sha256 b56bf70be9982044 :: mean +- SD over the row's 5 pretrained models
\newcommand{\FtAucJciiEFourOnesixtwoMass}{0.9724\,\ensuremath{\pm}\,0.0013}  % experiments/FIGS/data/w2b_leg1_metrics_v2.json :: cells.{l162mass-s1,l162mass-s2,l162mass-s3,l162mass-s4,l162mass-s5}.N10000.s1.macro_auc_ovr :: sha256 b56bf70be9982044 :: mean +- SD over the row's 5 pretrained models
\newcommand{\FtAucJciiEFiveOnesixtwoMass}{0.9790\,\ensuremath{\pm}\,0.0001}  % experiments/FIGS/data/w2b_leg1_metrics_v2.json :: cells.{l162mass-s1,l162mass-s2,l162mass-s3,l162mass-s4,l162mass-s5}.N100000.s1.macro_auc_ovr :: sha256 b56bf70be9982044 :: mean +- SD over the row's 5 pretrained models
\newcommand{\FtAucJciiESixOnesixtwoMass}{0.9805\,\ensuremath{\pm}\,0.0001}  % experiments/FIGS/data/w2b_leg1_metrics_v2.json :: cells.{l162mass-s1,l162mass-s2,l162mass-s3,l162mass-s4,l162mass-s5}.N1000000.s1.macro_auc_ovr :: sha256 b56bf70be9982044 :: mean +- SD over the row's 5 pretrained models
\newcommand{\FtAucJciiEThreeOnesevenMass}{0.904\,\ensuremath{\pm}\,0.007}  % experiments/FIGS/data/w2b_leg1_metrics_v2.json :: cells.{r16q1mass-s1,r16q1mass-s2,r16q1mass-s3,r16q1mass-s4,r16q1mass-s5}.N1000.s1.macro_auc_ovr :: sha256 b56bf70be9982044 :: mean +- SD over the row's 5 pretrained models
\newcommand{\FtAucJciiEFourOnesevenMass}{0.9639\,\ensuremath{\pm}\,0.0007}  % experiments/FIGS/data/w2b_leg1_metrics_v2.json :: cells.{r16q1mass-s1,r16q1mass-s2,r16q1mass-s3,r16q1mass-s4,r16q1mass-s5}.N10000.s1.macro_auc_ovr :: sha256 b56bf70be9982044 :: mean +- SD over the row's 5 pretrained models
\newcommand{\FtAucJciiEFiveOnesevenMass}{0.9740\,\ensuremath{\pm}\,0.0004}  % experiments/FIGS/data/w2b_leg1_metrics_v2.json :: cells.{r16q1mass-s1,r16q1mass-s2,r16q1mass-s3,r16q1mass-s4,r16q1mass-s5}.N100000.s1.macro_auc_ovr :: sha256 b56bf70be9982044 :: mean +- SD over the row's 5 pretrained models
\newcommand{\FtAucJciiESixOnesevenMass}{0.9774\,\ensuremath{\pm}\,0.0001}  % experiments/FIGS/data/w2b_leg1_metrics_v2.json :: cells.{r16q1mass-s1,r16q1mass-s2,r16q1mass-s3,r16q1mass-s4,r16q1mass-s5}.N1000000.s1.macro_auc_ovr :: sha256 b56bf70be9982044 :: mean +- SD over the row's 5 pretrained models
\newcommand{\FtAucJciiEThreeRandDrawOne}{0.921}  % experiments/FIGS/data/w2b_leg1_metrics_v2.json :: cells.{rand-d1-s1b}.N1000.s1.macro_auc_ovr :: sha256 b56bf70be9982044 :: one run
\newcommand{\FtAucJciiEFourRandDrawOne}{0.9675}  % experiments/FIGS/data/w2b_leg1_metrics_v2.json :: cells.{rand-d1-s1b}.N10000.s1.macro_auc_ovr :: sha256 b56bf70be9982044 :: one run
\newcommand{\FtAucJciiEFiveRandDrawOne}{0.97617}  % experiments/FIGS/data/w2b_leg1_metrics_v2.json :: cells.{rand-d1-s1b}.N100000.s1.macro_auc_ovr :: sha256 b56bf70be9982044 :: one run
\newcommand{\FtAucJciiESixRandDrawOne}{0.97941}  % experiments/FIGS/data/w2b_leg1_metrics_v2.json :: cells.{rand-d1-s1b}.N1000000.s1.macro_auc_ovr :: sha256 b56bf70be9982044 :: one run
\newcommand{\FtAucJciiEThreeRandDrawTwo}{0.914}  % experiments/FIGS/data/w2b_leg1_metrics_v2.json :: cells.{rand-d2-s2}.N1000.s1.macro_auc_ovr :: sha256 b56bf70be9982044 :: one run
\newcommand{\FtAucJciiEFourRandDrawTwo}{0.9675}  % experiments/FIGS/data/w2b_leg1_metrics_v2.json :: cells.{rand-d2-s2}.N10000.s1.macro_auc_ovr :: sha256 b56bf70be9982044 :: one run
\newcommand{\FtAucJciiEFiveRandDrawTwo}{0.97668}  % experiments/FIGS/data/w2b_leg1_metrics_v2.json :: cells.{rand-d2-s2}.N100000.s1.macro_auc_ovr :: sha256 b56bf70be9982044 :: one run
\newcommand{\FtAucJciiESixRandDrawTwo}{0.97943}  % experiments/FIGS/data/w2b_leg1_metrics_v2.json :: cells.{rand-d2-s2}.N1000000.s1.macro_auc_ovr :: sha256 b56bf70be9982044 :: one run
\newcommand{\FtAucJciiEThreeRandDrawThree}{0.918}  % experiments/FIGS/data/w2b_leg1_metrics_v2.json :: cells.{rand-d3-s3}.N1000.s1.macro_auc_ovr :: sha256 b56bf70be9982044 :: one run
\newcommand{\FtAucJciiEFourRandDrawThree}{0.9686}  % experiments/FIGS/data/w2b_leg1_metrics_v2.json :: cells.{rand-d3-s3}.N10000.s1.macro_auc_ovr :: sha256 b56bf70be9982044 :: one run
\newcommand{\FtAucJciiEFiveRandDrawThree}{0.97605}  % experiments/FIGS/data/w2b_leg1_metrics_v2.json :: cells.{rand-d3-s3}.N100000.s1.macro_auc_ovr :: sha256 b56bf70be9982044 :: one run
\newcommand{\FtAucJciiESixRandDrawThree}{0.97950}  % experiments/FIGS/data/w2b_leg1_metrics_v2.json :: cells.{rand-d3-s3}.N1000000.s1.macro_auc_ovr :: sha256 b56bf70be9982044 :: one run
\newcommand{\FtAccJciiEThreeOneeighteight}{0.445\,\ensuremath{\pm}\,0.004}  % experiments/FIGS/data/w2b_leg1_metrics_v2.json :: cells.{l188-s1,l188-s2,l188-s3,l188-s4,l188-s5}.N1000.s1.accuracy :: sha256 b56bf70be9982044 :: mean +- SD over the row's 5 pretrained models
\newcommand{\FtAccJciiEFourOneeighteight}{0.5199\,\ensuremath{\pm}\,0.0018}  % experiments/FIGS/data/w2b_leg1_metrics_v2.json :: cells.{l188-s1,l188-s2,l188-s3,l188-s4,l188-s5}.N10000.s1.accuracy :: sha256 b56bf70be9982044 :: mean +- SD over the row's 5 pretrained models
\newcommand{\FtAccJciiEFiveOneeighteight}{0.5379\,\ensuremath{\pm}\,0.0010}  % experiments/FIGS/data/w2b_leg1_metrics_v2.json :: cells.{l188-s1,l188-s2,l188-s3,l188-s4,l188-s5}.N100000.s1.accuracy :: sha256 b56bf70be9982044 :: mean +- SD over the row's 5 pretrained models
\newcommand{\FtAccJciiESixOneeighteight}{0.5443\,\ensuremath{\pm}\,0.0003}  % experiments/FIGS/data/w2b_leg1_metrics_v2.json :: cells.{l188-s1,l188-s2,l188-s3,l188-s4,l188-s5}.N1000000.s1.accuracy :: sha256 b56bf70be9982044 :: mean +- SD over the row's 5 pretrained models
\newcommand{\FtAccJciiEThreeOnesixtwo}{0.448\,\ensuremath{\pm}\,0.006}  % experiments/FIGS/data/w2b_leg1_metrics_v2.json :: cells.{l162-s1b,l162-s2,l162-s3,l162-s4,l162-s5}.N1000.s1.accuracy :: sha256 b56bf70be9982044 :: mean +- SD over the row's 5 pretrained models
\newcommand{\FtAccJciiEFourOnesixtwo}{0.5212\,\ensuremath{\pm}\,0.0013}  % experiments/FIGS/data/w2b_leg1_metrics_v2.json :: cells.{l162-s1b,l162-s2,l162-s3,l162-s4,l162-s5}.N10000.s1.accuracy :: sha256 b56bf70be9982044 :: mean +- SD over the row's 5 pretrained models
\newcommand{\FtAccJciiEFiveOnesixtwo}{0.5387\,\ensuremath{\pm}\,0.0008}  % experiments/FIGS/data/w2b_leg1_metrics_v2.json :: cells.{l162-s1b,l162-s2,l162-s3,l162-s4,l162-s5}.N100000.s1.accuracy :: sha256 b56bf70be9982044 :: mean +- SD over the row's 5 pretrained models
\newcommand{\FtAccJciiESixOnesixtwo}{0.5440\,\ensuremath{\pm}\,0.0006}  % experiments/FIGS/data/w2b_leg1_metrics_v2.json :: cells.{l162-s1b,l162-s2,l162-s3,l162-s4,l162-s5}.N1000000.s1.accuracy :: sha256 b56bf70be9982044 :: mean +- SD over the row's 5 pretrained models
\newcommand{\FtAccJciiEThreeFourthree}{0.380\,\ensuremath{\pm}\,0.005}  % experiments/FIGS/data/w2b_leg1_metrics_v2.json :: cells.{r42q1-s1,r42q1-s2,r42q1-s3,r42q1-s4,r42q1-s5}.N1000.s1.accuracy :: sha256 b56bf70be9982044 :: mean +- SD over the row's 5 pretrained models
\newcommand{\FtAccJciiEFourFourthree}{0.481\,\ensuremath{\pm}\,0.004}  % experiments/FIGS/data/w2b_leg1_metrics_v2.json :: cells.{r42q1-s1,r42q1-s2,r42q1-s3,r42q1-s4,r42q1-s5}.N10000.s1.accuracy :: sha256 b56bf70be9982044 :: mean +- SD over the row's 5 pretrained models
\newcommand{\FtAccJciiEFiveFourthree}{0.5219\,\ensuremath{\pm}\,0.0017}  % experiments/FIGS/data/w2b_leg1_metrics_v2.json :: cells.{r42q1-s1,r42q1-s2,r42q1-s3,r42q1-s4,r42q1-s5}.N100000.s1.accuracy :: sha256 b56bf70be9982044 :: mean +- SD over the row's 5 pretrained models
\newcommand{\FtAccJciiESixFourthree}{0.5362\,\ensuremath{\pm}\,0.0007}  % experiments/FIGS/data/w2b_leg1_metrics_v2.json :: cells.{r42q1-s1,r42q1-s2,r42q1-s3,r42q1-s4,r42q1-s5}.N1000000.s1.accuracy :: sha256 b56bf70be9982044 :: mean +- SD over the row's 5 pretrained models
\newcommand{\FtAccJciiEThreeOneseven}{0.332\,\ensuremath{\pm}\,0.004}  % experiments/FIGS/data/w2b_leg1_metrics_v2.json :: cells.{r16q1-s1,r16q1-s2,r16q1-s3,r16q1-s4,r16q1-s5}.N1000.s1.accuracy :: sha256 b56bf70be9982044 :: mean +- SD over the row's 5 pretrained models
\newcommand{\FtAccJciiEFourOneseven}{0.447\,\ensuremath{\pm}\,0.013}  % experiments/FIGS/data/w2b_leg1_metrics_v2.json :: cells.{r16q1-s1,r16q1-s2,r16q1-s3,r16q1-s4,r16q1-s5}.N10000.s1.accuracy :: sha256 b56bf70be9982044 :: mean +- SD over the row's 5 pretrained models
\newcommand{\FtAccJciiEFiveOneseven}{0.5078\,\ensuremath{\pm}\,0.0011}  % experiments/FIGS/data/w2b_leg1_metrics_v2.json :: cells.{r16q1-s1,r16q1-s2,r16q1-s3,r16q1-s4,r16q1-s5}.N100000.s1.accuracy :: sha256 b56bf70be9982044 :: mean +- SD over the row's 5 pretrained models
\newcommand{\FtAccJciiESixOneseven}{0.5227\,\ensuremath{\pm}\,0.0010}  % experiments/FIGS/data/w2b_leg1_metrics_v2.json :: cells.{r16q1-s1,r16q1-s2,r16q1-s3,r16q1-s4,r16q1-s5}.N1000000.s1.accuracy :: sha256 b56bf70be9982044 :: mean +- SD over the row's 5 pretrained models
\newcommand{\FtAccJciiEThreeOnesixtwoMass}{0.451\,\ensuremath{\pm}\,0.007}  % experiments/FIGS/data/w2b_leg1_metrics_v2.json :: cells.{l162mass-s1,l162mass-s2,l162mass-s3,l162mass-s4,l162mass-s5}.N1000.s1.accuracy :: sha256 b56bf70be9982044 :: mean +- SD over the row's 5 pretrained models
\newcommand{\FtAccJciiEFourOnesixtwoMass}{0.515\,\ensuremath{\pm}\,0.004}  % experiments/FIGS/data/w2b_leg1_metrics_v2.json :: cells.{l162mass-s1,l162mass-s2,l162mass-s3,l162mass-s4,l162mass-s5}.N10000.s1.accuracy :: sha256 b56bf70be9982044 :: mean +- SD over the row's 5 pretrained models
\newcommand{\FtAccJciiEFiveOnesixtwoMass}{0.5357\,\ensuremath{\pm}\,0.0007}  % experiments/FIGS/data/w2b_leg1_metrics_v2.json :: cells.{l162mass-s1,l162mass-s2,l162mass-s3,l162mass-s4,l162mass-s5}.N100000.s1.accuracy :: sha256 b56bf70be9982044 :: mean +- SD over the row's 5 pretrained models
\newcommand{\FtAccJciiESixOnesixtwoMass}{0.5409\,\ensuremath{\pm}\,0.0010}  % experiments/FIGS/data/w2b_leg1_metrics_v2.json :: cells.{l162mass-s1,l162mass-s2,l162mass-s3,l162mass-s4,l162mass-s5}.N1000000.s1.accuracy :: sha256 b56bf70be9982044 :: mean +- SD over the row's 5 pretrained models
\newcommand{\FtAccJciiEThreeOnesevenMass}{0.331\,\ensuremath{\pm}\,0.005}  % experiments/FIGS/data/w2b_leg1_metrics_v2.json :: cells.{r16q1mass-s1,r16q1mass-s2,r16q1mass-s3,r16q1mass-s4,r16q1mass-s5}.N1000.s1.accuracy :: sha256 b56bf70be9982044 :: mean +- SD over the row's 5 pretrained models
\newcommand{\FtAccJciiEFourOnesevenMass}{0.4507\,\ensuremath{\pm}\,0.0034}  % experiments/FIGS/data/w2b_leg1_metrics_v2.json :: cells.{r16q1mass-s1,r16q1mass-s2,r16q1mass-s3,r16q1mass-s4,r16q1mass-s5}.N10000.s1.accuracy :: sha256 b56bf70be9982044 :: mean +- SD over the row's 5 pretrained models
\newcommand{\FtAccJciiEFiveOnesevenMass}{0.4993\,\ensuremath{\pm}\,0.0022}  % experiments/FIGS/data/w2b_leg1_metrics_v2.json :: cells.{r16q1mass-s1,r16q1mass-s2,r16q1mass-s3,r16q1mass-s4,r16q1mass-s5}.N100000.s1.accuracy :: sha256 b56bf70be9982044 :: mean +- SD over the row's 5 pretrained models
\newcommand{\FtAccJciiESixOnesevenMass}{0.5166\,\ensuremath{\pm}\,0.0013}  % experiments/FIGS/data/w2b_leg1_metrics_v2.json :: cells.{r16q1mass-s1,r16q1mass-s2,r16q1mass-s3,r16q1mass-s4,r16q1mass-s5}.N1000000.s1.accuracy :: sha256 b56bf70be9982044 :: mean +- SD over the row's 5 pretrained models
\newcommand{\FtAccJciiEThreeRandDrawOne}{0.424}  % experiments/FIGS/data/w2b_leg1_metrics_v2.json :: cells.{rand-d1-s1b}.N1000.s1.accuracy :: sha256 b56bf70be9982044 :: one run
\newcommand{\FtAccJciiEFourRandDrawOne}{0.5070}  % experiments/FIGS/data/w2b_leg1_metrics_v2.json :: cells.{rand-d1-s1b}.N10000.s1.accuracy :: sha256 b56bf70be9982044 :: one run
\newcommand{\FtAccJciiEFiveRandDrawOne}{0.5226}  % experiments/FIGS/data/w2b_leg1_metrics_v2.json :: cells.{rand-d1-s1b}.N100000.s1.accuracy :: sha256 b56bf70be9982044 :: one run
\newcommand{\FtAccJciiESixRandDrawOne}{0.53354}  % experiments/FIGS/data/w2b_leg1_metrics_v2.json :: cells.{rand-d1-s1b}.N1000000.s1.accuracy :: sha256 b56bf70be9982044 :: one run
\newcommand{\FtAccJciiEThreeRandDrawTwo}{0.420}  % experiments/FIGS/data/w2b_leg1_metrics_v2.json :: cells.{rand-d2-s2}.N1000.s1.accuracy :: sha256 b56bf70be9982044 :: one run
\newcommand{\FtAccJciiEFourRandDrawTwo}{0.5074}  % experiments/FIGS/data/w2b_leg1_metrics_v2.json :: cells.{rand-d2-s2}.N10000.s1.accuracy :: sha256 b56bf70be9982044 :: one run
\newcommand{\FtAccJciiEFiveRandDrawTwo}{0.5243}  % experiments/FIGS/data/w2b_leg1_metrics_v2.json :: cells.{rand-d2-s2}.N100000.s1.accuracy :: sha256 b56bf70be9982044 :: one run
\newcommand{\FtAccJciiESixRandDrawTwo}{0.53441}  % experiments/FIGS/data/w2b_leg1_metrics_v2.json :: cells.{rand-d2-s2}.N1000000.s1.accuracy :: sha256 b56bf70be9982044 :: one run
\newcommand{\FtAccJciiEThreeRandDrawThree}{0.423}  % experiments/FIGS/data/w2b_leg1_metrics_v2.json :: cells.{rand-d3-s3}.N1000.s1.accuracy :: sha256 b56bf70be9982044 :: one run
\newcommand{\FtAccJciiEFourRandDrawThree}{0.5114}  % experiments/FIGS/data/w2b_leg1_metrics_v2.json :: cells.{rand-d3-s3}.N10000.s1.accuracy :: sha256 b56bf70be9982044 :: one run
\newcommand{\FtAccJciiEFiveRandDrawThree}{0.5245}  % experiments/FIGS/data/w2b_leg1_metrics_v2.json :: cells.{rand-d3-s3}.N100000.s1.accuracy :: sha256 b56bf70be9982044 :: one run
\newcommand{\FtAccJciiESixRandDrawThree}{0.53473}  % experiments/FIGS/data/w2b_leg1_metrics_v2.json :: cells.{rand-d3-s3}.N1000000.s1.accuracy :: sha256 b56bf70be9982044 :: one run
\newcommand{\FtOmaRatioJciiEThreeOnesixtwoOverOneeighteight}{0.96}  % experiments/FIGS/data/w2b_leg1_metrics_v2.json :: cells.*.N1000.s1.macro_auc_ovr :: sha256 b56bf70be9982044 :: ratio of the seed means of 1 - macro AUC, this row over the reference
\newcommand{\FtOmaRatioJciiEThreeFourthreeOverOneeighteight}{0.84}  % experiments/FIGS/data/w2b_leg1_metrics_v2.json :: cells.*.N1000.s1.macro_auc_ovr :: sha256 b56bf70be9982044 :: ratio of the seed means of 1 - macro AUC, this row over the reference
\newcommand{\FtOmaRatioJciiEThreeOnesevenOverOneeighteight}{0.97}  % experiments/FIGS/data/w2b_leg1_metrics_v2.json :: cells.*.N1000.s1.macro_auc_ovr :: sha256 b56bf70be9982044 :: ratio of the seed means of 1 - macro AUC, this row over the reference
\newcommand{\FtOmaRatioJciiEThreeOnesixtwoMassOverOneeighteight}{0.81}  % experiments/FIGS/data/w2b_leg1_metrics_v2.json :: cells.*.N1000.s1.macro_auc_ovr :: sha256 b56bf70be9982044 :: ratio of the seed means of 1 - macro AUC, this row over the reference
\newcommand{\FtOmaRatioJciiEThreeOnesevenMassOverOneeighteight}{0.98}  % experiments/FIGS/data/w2b_leg1_metrics_v2.json :: cells.*.N1000.s1.macro_auc_ovr :: sha256 b56bf70be9982044 :: ratio of the seed means of 1 - macro AUC, this row over the reference
\newcommand{\FtOmaRatioJciiEThreeRandDrawOneOverOneeighteight}{0.80}  % experiments/FIGS/data/w2b_leg1_metrics_v2.json :: cells.*.N1000.s1.macro_auc_ovr :: sha256 b56bf70be9982044 :: ratio of the seed means of 1 - macro AUC, this row over the reference
\newcommand{\FtOmaRatioJciiEThreeRandDrawTwoOverOneeighteight}{0.87}  % experiments/FIGS/data/w2b_leg1_metrics_v2.json :: cells.*.N1000.s1.macro_auc_ovr :: sha256 b56bf70be9982044 :: ratio of the seed means of 1 - macro AUC, this row over the reference
\newcommand{\FtOmaRatioJciiEThreeRandDrawThreeOverOneeighteight}{0.83}  % experiments/FIGS/data/w2b_leg1_metrics_v2.json :: cells.*.N1000.s1.macro_auc_ovr :: sha256 b56bf70be9982044 :: ratio of the seed means of 1 - macro AUC, this row over the reference
\newcommand{\FtOmaRatioJciiEThreeOneeighteightOverOnesixtwo}{1.05}  % experiments/FIGS/data/w2b_leg1_metrics_v2.json :: cells.*.N1000.s1.macro_auc_ovr :: sha256 b56bf70be9982044 :: ratio of the seed means of 1 - macro AUC, this row over the reference
\newcommand{\FtOmaRatioJciiEThreeFourthreeOverOnesixtwo}{0.88}  % experiments/FIGS/data/w2b_leg1_metrics_v2.json :: cells.*.N1000.s1.macro_auc_ovr :: sha256 b56bf70be9982044 :: ratio of the seed means of 1 - macro AUC, this row over the reference
\newcommand{\FtOmaRatioJciiEThreeOnesevenOverOnesixtwo}{1.02}  % experiments/FIGS/data/w2b_leg1_metrics_v2.json :: cells.*.N1000.s1.macro_auc_ovr :: sha256 b56bf70be9982044 :: ratio of the seed means of 1 - macro AUC, this row over the reference
\newcommand{\FtOmaRatioJciiEThreeOnesixtwoMassOverOnesixtwo}{0.85}  % experiments/FIGS/data/w2b_leg1_metrics_v2.json :: cells.*.N1000.s1.macro_auc_ovr :: sha256 b56bf70be9982044 :: ratio of the seed means of 1 - macro AUC, this row over the reference
\newcommand{\FtOmaRatioJciiEThreeOnesevenMassOverOnesixtwo}{1.02}  % experiments/FIGS/data/w2b_leg1_metrics_v2.json :: cells.*.N1000.s1.macro_auc_ovr :: sha256 b56bf70be9982044 :: ratio of the seed means of 1 - macro AUC, this row over the reference
\newcommand{\FtOmaRatioJciiEThreeRandDrawOneOverOnesixtwo}{0.84}  % experiments/FIGS/data/w2b_leg1_metrics_v2.json :: cells.*.N1000.s1.macro_auc_ovr :: sha256 b56bf70be9982044 :: ratio of the seed means of 1 - macro AUC, this row over the reference
\newcommand{\FtOmaRatioJciiEThreeRandDrawTwoOverOnesixtwo}{0.91}  % experiments/FIGS/data/w2b_leg1_metrics_v2.json :: cells.*.N1000.s1.macro_auc_ovr :: sha256 b56bf70be9982044 :: ratio of the seed means of 1 - macro AUC, this row over the reference
\newcommand{\FtOmaRatioJciiEThreeRandDrawThreeOverOnesixtwo}{0.87}  % experiments/FIGS/data/w2b_leg1_metrics_v2.json :: cells.*.N1000.s1.macro_auc_ovr :: sha256 b56bf70be9982044 :: ratio of the seed means of 1 - macro AUC, this row over the reference
\newcommand{\FtOmaRatioJciiEFourOnesixtwoOverOneeighteight}{0.99}  % experiments/FIGS/data/w2b_leg1_metrics_v2.json :: cells.*.N10000.s1.macro_auc_ovr :: sha256 b56bf70be9982044 :: ratio of the seed means of 1 - macro AUC, this row over the reference
\newcommand{\FtOmaRatioJciiEFourFourthreeOverOneeighteight}{1.20}  % experiments/FIGS/data/w2b_leg1_metrics_v2.json :: cells.*.N10000.s1.macro_auc_ovr :: sha256 b56bf70be9982044 :: ratio of the seed means of 1 - macro AUC, this row over the reference
\newcommand{\FtOmaRatioJciiEFourOnesevenOverOneeighteight}{1.32}  % experiments/FIGS/data/w2b_leg1_metrics_v2.json :: cells.*.N10000.s1.macro_auc_ovr :: sha256 b56bf70be9982044 :: ratio of the seed means of 1 - macro AUC, this row over the reference
\newcommand{\FtOmaRatioJciiEFourOnesixtwoMassOverOneeighteight}{0.99}  % experiments/FIGS/data/w2b_leg1_metrics_v2.json :: cells.*.N10000.s1.macro_auc_ovr :: sha256 b56bf70be9982044 :: ratio of the seed means of 1 - macro AUC, this row over the reference
\newcommand{\FtOmaRatioJciiEFourOnesevenMassOverOneeighteight}{1.30}  % experiments/FIGS/data/w2b_leg1_metrics_v2.json :: cells.*.N10000.s1.macro_auc_ovr :: sha256 b56bf70be9982044 :: ratio of the seed means of 1 - macro AUC, this row over the reference
\newcommand{\FtOmaRatioJciiEFourRandDrawOneOverOneeighteight}{1.17}  % experiments/FIGS/data/w2b_leg1_metrics_v2.json :: cells.*.N10000.s1.macro_auc_ovr :: sha256 b56bf70be9982044 :: ratio of the seed means of 1 - macro AUC, this row over the reference
\newcommand{\FtOmaRatioJciiEFourRandDrawTwoOverOneeighteight}{1.17}  % experiments/FIGS/data/w2b_leg1_metrics_v2.json :: cells.*.N10000.s1.macro_auc_ovr :: sha256 b56bf70be9982044 :: ratio of the seed means of 1 - macro AUC, this row over the reference
\newcommand{\FtOmaRatioJciiEFourRandDrawThreeOverOneeighteight}{1.13}  % experiments/FIGS/data/w2b_leg1_metrics_v2.json :: cells.*.N10000.s1.macro_auc_ovr :: sha256 b56bf70be9982044 :: ratio of the seed means of 1 - macro AUC, this row over the reference
\newcommand{\FtOmaRatioJciiEFourOneeighteightOverOnesixtwo}{1.01}  % experiments/FIGS/data/w2b_leg1_metrics_v2.json :: cells.*.N10000.s1.macro_auc_ovr :: sha256 b56bf70be9982044 :: ratio of the seed means of 1 - macro AUC, this row over the reference
\newcommand{\FtOmaRatioJciiEFourFourthreeOverOnesixtwo}{1.21}  % experiments/FIGS/data/w2b_leg1_metrics_v2.json :: cells.*.N10000.s1.macro_auc_ovr :: sha256 b56bf70be9982044 :: ratio of the seed means of 1 - macro AUC, this row over the reference
\newcommand{\FtOmaRatioJciiEFourOnesevenOverOnesixtwo}{1.32}  % experiments/FIGS/data/w2b_leg1_metrics_v2.json :: cells.*.N10000.s1.macro_auc_ovr :: sha256 b56bf70be9982044 :: ratio of the seed means of 1 - macro AUC, this row over the reference
\newcommand{\FtOmaRatioJciiEFourOnesixtwoMassOverOnesixtwo}{1.00}  % experiments/FIGS/data/w2b_leg1_metrics_v2.json :: cells.*.N10000.s1.macro_auc_ovr :: sha256 b56bf70be9982044 :: ratio of the seed means of 1 - macro AUC, this row over the reference
\newcommand{\FtOmaRatioJciiEFourOnesevenMassOverOnesixtwo}{1.31}  % experiments/FIGS/data/w2b_leg1_metrics_v2.json :: cells.*.N10000.s1.macro_auc_ovr :: sha256 b56bf70be9982044 :: ratio of the seed means of 1 - macro AUC, this row over the reference
\newcommand{\FtOmaRatioJciiEFourRandDrawOneOverOnesixtwo}{1.18}  % experiments/FIGS/data/w2b_leg1_metrics_v2.json :: cells.*.N10000.s1.macro_auc_ovr :: sha256 b56bf70be9982044 :: ratio of the seed means of 1 - macro AUC, this row over the reference
\newcommand{\FtOmaRatioJciiEFourRandDrawTwoOverOnesixtwo}{1.18}  % experiments/FIGS/data/w2b_leg1_metrics_v2.json :: cells.*.N10000.s1.macro_auc_ovr :: sha256 b56bf70be9982044 :: ratio of the seed means of 1 - macro AUC, this row over the reference
\newcommand{\FtOmaRatioJciiEFourRandDrawThreeOverOnesixtwo}{1.14}  % experiments/FIGS/data/w2b_leg1_metrics_v2.json :: cells.*.N10000.s1.macro_auc_ovr :: sha256 b56bf70be9982044 :: ratio of the seed means of 1 - macro AUC, this row over the reference
\newcommand{\FtOmaRatioJciiEFiveOnesixtwoOverOneeighteight}{1.00}  % experiments/FIGS/data/w2b_leg1_metrics_v2.json :: cells.*.N100000.s1.macro_auc_ovr :: sha256 b56bf70be9982044 :: ratio of the seed means of 1 - macro AUC, this row over the reference
\newcommand{\FtOmaRatioJciiEFiveFourthreeOverOneeighteight}{1.09}  % experiments/FIGS/data/w2b_leg1_metrics_v2.json :: cells.*.N100000.s1.macro_auc_ovr :: sha256 b56bf70be9982044 :: ratio of the seed means of 1 - macro AUC, this row over the reference
\newcommand{\FtOmaRatioJciiEFiveOnesevenOverOneeighteight}{1.22}  % experiments/FIGS/data/w2b_leg1_metrics_v2.json :: cells.*.N100000.s1.macro_auc_ovr :: sha256 b56bf70be9982044 :: ratio of the seed means of 1 - macro AUC, this row over the reference
\newcommand{\FtOmaRatioJciiEFiveOnesixtwoMassOverOneeighteight}{1.01}  % experiments/FIGS/data/w2b_leg1_metrics_v2.json :: cells.*.N100000.s1.macro_auc_ovr :: sha256 b56bf70be9982044 :: ratio of the seed means of 1 - macro AUC, this row over the reference
\newcommand{\FtOmaRatioJciiEFiveOnesevenMassOverOneeighteight}{1.26}  % experiments/FIGS/data/w2b_leg1_metrics_v2.json :: cells.*.N100000.s1.macro_auc_ovr :: sha256 b56bf70be9982044 :: ratio of the seed means of 1 - macro AUC, this row over the reference
\newcommand{\FtOmaRatioJciiEFiveRandDrawOneOverOneeighteight}{1.15}  % experiments/FIGS/data/w2b_leg1_metrics_v2.json :: cells.*.N100000.s1.macro_auc_ovr :: sha256 b56bf70be9982044 :: ratio of the seed means of 1 - macro AUC, this row over the reference
\newcommand{\FtOmaRatioJciiEFiveRandDrawTwoOverOneeighteight}{1.12}  % experiments/FIGS/data/w2b_leg1_metrics_v2.json :: cells.*.N100000.s1.macro_auc_ovr :: sha256 b56bf70be9982044 :: ratio of the seed means of 1 - macro AUC, this row over the reference
\newcommand{\FtOmaRatioJciiEFiveRandDrawThreeOverOneeighteight}{1.15}  % experiments/FIGS/data/w2b_leg1_metrics_v2.json :: cells.*.N100000.s1.macro_auc_ovr :: sha256 b56bf70be9982044 :: ratio of the seed means of 1 - macro AUC, this row over the reference
\newcommand{\FtOmaRatioJciiEFiveOneeighteightOverOnesixtwo}{1.00}  % experiments/FIGS/data/w2b_leg1_metrics_v2.json :: cells.*.N100000.s1.macro_auc_ovr :: sha256 b56bf70be9982044 :: ratio of the seed means of 1 - macro AUC, this row over the reference
\newcommand{\FtOmaRatioJciiEFiveFourthreeOverOnesixtwo}{1.09}  % experiments/FIGS/data/w2b_leg1_metrics_v2.json :: cells.*.N100000.s1.macro_auc_ovr :: sha256 b56bf70be9982044 :: ratio of the seed means of 1 - macro AUC, this row over the reference
\newcommand{\FtOmaRatioJciiEFiveOnesevenOverOnesixtwo}{1.22}  % experiments/FIGS/data/w2b_leg1_metrics_v2.json :: cells.*.N100000.s1.macro_auc_ovr :: sha256 b56bf70be9982044 :: ratio of the seed means of 1 - macro AUC, this row over the reference
\newcommand{\FtOmaRatioJciiEFiveOnesixtwoMassOverOnesixtwo}{1.02}  % experiments/FIGS/data/w2b_leg1_metrics_v2.json :: cells.*.N100000.s1.macro_auc_ovr :: sha256 b56bf70be9982044 :: ratio of the seed means of 1 - macro AUC, this row over the reference
\newcommand{\FtOmaRatioJciiEFiveOnesevenMassOverOnesixtwo}{1.26}  % experiments/FIGS/data/w2b_leg1_metrics_v2.json :: cells.*.N100000.s1.macro_auc_ovr :: sha256 b56bf70be9982044 :: ratio of the seed means of 1 - macro AUC, this row over the reference
\newcommand{\FtOmaRatioJciiEFiveRandDrawOneOverOnesixtwo}{1.15}  % experiments/FIGS/data/w2b_leg1_metrics_v2.json :: cells.*.N100000.s1.macro_auc_ovr :: sha256 b56bf70be9982044 :: ratio of the seed means of 1 - macro AUC, this row over the reference
\newcommand{\FtOmaRatioJciiEFiveRandDrawTwoOverOnesixtwo}{1.13}  % experiments/FIGS/data/w2b_leg1_metrics_v2.json :: cells.*.N100000.s1.macro_auc_ovr :: sha256 b56bf70be9982044 :: ratio of the seed means of 1 - macro AUC, this row over the reference
\newcommand{\FtOmaRatioJciiEFiveRandDrawThreeOverOnesixtwo}{1.16}  % experiments/FIGS/data/w2b_leg1_metrics_v2.json :: cells.*.N100000.s1.macro_auc_ovr :: sha256 b56bf70be9982044 :: ratio of the seed means of 1 - macro AUC, this row over the reference
\newcommand{\FtOmaRatioJciiESixOnesixtwoOverOneeighteight}{1.00}  % experiments/FIGS/data/w2b_leg1_metrics_v2.json :: cells.*.N1000000.s1.macro_auc_ovr :: sha256 b56bf70be9982044 :: ratio of the seed means of 1 - macro AUC, this row over the reference
\newcommand{\FtOmaRatioJciiESixFourthreeOverOneeighteight}{1.05}  % experiments/FIGS/data/w2b_leg1_metrics_v2.json :: cells.*.N1000000.s1.macro_auc_ovr :: sha256 b56bf70be9982044 :: ratio of the seed means of 1 - macro AUC, this row over the reference
\newcommand{\FtOmaRatioJciiESixOnesevenOverOneeighteight}{1.14}  % experiments/FIGS/data/w2b_leg1_metrics_v2.json :: cells.*.N1000000.s1.macro_auc_ovr :: sha256 b56bf70be9982044 :: ratio of the seed means of 1 - macro AUC, this row over the reference
\newcommand{\FtOmaRatioJciiESixOnesixtwoMassOverOneeighteight}{1.02}  % experiments/FIGS/data/w2b_leg1_metrics_v2.json :: cells.*.N1000000.s1.macro_auc_ovr :: sha256 b56bf70be9982044 :: ratio of the seed means of 1 - macro AUC, this row over the reference
\newcommand{\FtOmaRatioJciiESixOnesevenMassOverOneeighteight}{1.18}  % experiments/FIGS/data/w2b_leg1_metrics_v2.json :: cells.*.N1000000.s1.macro_auc_ovr :: sha256 b56bf70be9982044 :: ratio of the seed means of 1 - macro AUC, this row over the reference
\newcommand{\FtOmaRatioJciiESixRandDrawOneOverOneeighteight}{1.08}  % experiments/FIGS/data/w2b_leg1_metrics_v2.json :: cells.*.N1000000.s1.macro_auc_ovr :: sha256 b56bf70be9982044 :: ratio of the seed means of 1 - macro AUC, this row over the reference
\newcommand{\FtOmaRatioJciiESixRandDrawTwoOverOneeighteight}{1.07}  % experiments/FIGS/data/w2b_leg1_metrics_v2.json :: cells.*.N1000000.s1.macro_auc_ovr :: sha256 b56bf70be9982044 :: ratio of the seed means of 1 - macro AUC, this row over the reference
\newcommand{\FtOmaRatioJciiESixRandDrawThreeOverOneeighteight}{1.07}  % experiments/FIGS/data/w2b_leg1_metrics_v2.json :: cells.*.N1000000.s1.macro_auc_ovr :: sha256 b56bf70be9982044 :: ratio of the seed means of 1 - macro AUC, this row over the reference
\newcommand{\FtOmaRatioJciiESixOneeighteightOverOnesixtwo}{1.00}  % experiments/FIGS/data/w2b_leg1_metrics_v2.json :: cells.*.N1000000.s1.macro_auc_ovr :: sha256 b56bf70be9982044 :: ratio of the seed means of 1 - macro AUC, this row over the reference
\newcommand{\FtOmaRatioJciiESixFourthreeOverOnesixtwo}{1.04}  % experiments/FIGS/data/w2b_leg1_metrics_v2.json :: cells.*.N1000000.s1.macro_auc_ovr :: sha256 b56bf70be9982044 :: ratio of the seed means of 1 - macro AUC, this row over the reference
\newcommand{\FtOmaRatioJciiESixOnesevenOverOnesixtwo}{1.14}  % experiments/FIGS/data/w2b_leg1_metrics_v2.json :: cells.*.N1000000.s1.macro_auc_ovr :: sha256 b56bf70be9982044 :: ratio of the seed means of 1 - macro AUC, this row over the reference
\newcommand{\FtOmaRatioJciiESixOnesixtwoMassOverOnesixtwo}{1.02}  % experiments/FIGS/data/w2b_leg1_metrics_v2.json :: cells.*.N1000000.s1.macro_auc_ovr :: sha256 b56bf70be9982044 :: ratio of the seed means of 1 - macro AUC, this row over the reference
\newcommand{\FtOmaRatioJciiESixOnesevenMassOverOnesixtwo}{1.18}  % experiments/FIGS/data/w2b_leg1_metrics_v2.json :: cells.*.N1000000.s1.macro_auc_ovr :: sha256 b56bf70be9982044 :: ratio of the seed means of 1 - macro AUC, this row over the reference
\newcommand{\FtOmaRatioJciiESixRandDrawOneOverOnesixtwo}{1.07}  % experiments/FIGS/data/w2b_leg1_metrics_v2.json :: cells.*.N1000000.s1.macro_auc_ovr :: sha256 b56bf70be9982044 :: ratio of the seed means of 1 - macro AUC, this row over the reference
\newcommand{\FtOmaRatioJciiESixRandDrawTwoOverOnesixtwo}{1.07}  % experiments/FIGS/data/w2b_leg1_metrics_v2.json :: cells.*.N1000000.s1.macro_auc_ovr :: sha256 b56bf70be9982044 :: ratio of the seed means of 1 - macro AUC, this row over the reference
\newcommand{\FtOmaRatioJciiESixRandDrawThreeOverOnesixtwo}{1.07}  % experiments/FIGS/data/w2b_leg1_metrics_v2.json :: cells.*.N1000000.s1.macro_auc_ovr :: sha256 b56bf70be9982044 :: ratio of the seed means of 1 - macro AUC, this row over the reference
\newcommand{\PairedFtAccJciiEThreeOnesevenMinusOneeighteight}{\ensuremath{-}0.1132~[\ensuremath{-}0.1214, \ensuremath{-}0.1051]}  % experiments/FIGS/data/w2b_leg1_metrics_v2.json :: cells.{r16q1-s1,r16q1-s2,r16q1-s3,r16q1-s4,r16q1-s5,l188-s1,l188-s2,l188-s3,l188-s4,l188-s5}.N1000.s1.accuracy :: sha256 b56bf70be9982044 :: accuracy, coarse minus fine, paired by run [95% Student-t interval]
\newcommand{\PairedFtAccJciiEFourOnesevenMinusOneeighteight}{\ensuremath{-}0.0730~[\ensuremath{-}0.0902, \ensuremath{-}0.0558]}  % experiments/FIGS/data/w2b_leg1_metrics_v2.json :: cells.{r16q1-s1,r16q1-s2,r16q1-s3,r16q1-s4,r16q1-s5,l188-s1,l188-s2,l188-s3,l188-s4,l188-s5}.N10000.s1.accuracy :: sha256 b56bf70be9982044 :: accuracy, coarse minus fine, paired by run [95% Student-t interval]
\newcommand{\PairedFtAccJciiEFiveOnesevenMinusOneeighteight}{\ensuremath{-}0.0301~[\ensuremath{-}0.0319, \ensuremath{-}0.0284]}  % experiments/FIGS/data/w2b_leg1_metrics_v2.json :: cells.{r16q1-s1,r16q1-s2,r16q1-s3,r16q1-s4,r16q1-s5,l188-s1,l188-s2,l188-s3,l188-s4,l188-s5}.N100000.s1.accuracy :: sha256 b56bf70be9982044 :: accuracy, coarse minus fine, paired by run [95% Student-t interval]
\newcommand{\PairedFtAccJciiESixOnesevenMinusOneeighteight}{\ensuremath{-}0.0216~[\ensuremath{-}0.0232, \ensuremath{-}0.0199]}  % experiments/FIGS/data/w2b_leg1_metrics_v2.json :: cells.{r16q1-s1,r16q1-s2,r16q1-s3,r16q1-s4,r16q1-s5,l188-s1,l188-s2,l188-s3,l188-s4,l188-s5}.N1000000.s1.accuracy :: sha256 b56bf70be9982044 :: accuracy, coarse minus fine, paired by run [95% Student-t interval]
\newcommand{\FtNClassesJcii}{162}  % experiments/FIGS/data/w2b_leg1_metrics_v2.json :: cells.*.*.s1.n_classes_present :: sha256 b56bf70be9982044 :: classes of the fine-tuning task
\newcommand{\FtNTestAucJcii}{125{,}000}  % experiments/FIGS/data/w2b_leg1_metrics_v2.json :: cells.*.*.s1.n_jets_auc :: sha256 b56bf70be9982044 :: test jets the macro AUC is computed on
\newcommand{\FtNTestAccJcii}{500{,}000}  % experiments/FIGS/data/w2b_leg1_metrics_v2.json :: cells.*.*.s1.n_jets :: sha256 b56bf70be9982044 :: test jets the accuracy is computed on
\newcommand{\FtAucJcEThreeOneeighteight}{0.950\,\ensuremath{\pm}\,0.004}  % experiments/FIGS/data/w2b_leg2_metrics_v2.json :: cells.{l188-s1,l188-s2,l188-s3,l188-s4,l188-s5}.N1000.s1.macro_auc_ovr :: sha256 9cd059d3e2f51798 :: mean +- SD over the row's 5 pretrained models
\newcommand{\FtAucJcEFourOneeighteight}{0.9765\,\ensuremath{\pm}\,0.0008}  % experiments/FIGS/data/w2b_leg2_metrics_v2.json :: cells.{l188-s1,l188-s2,l188-s3,l188-s4,l188-s5}.N10000.s1.macro_auc_ovr :: sha256 9cd059d3e2f51798 :: mean +- SD over the row's 5 pretrained models
\newcommand{\FtAucJcEFiveOneeighteight}{0.9833\,\ensuremath{\pm}\,0.0002}  % experiments/FIGS/data/w2b_leg2_metrics_v2.json :: cells.{l188-s1,l188-s2,l188-s3,l188-s4,l188-s5}.N100000.s1.macro_auc_ovr :: sha256 9cd059d3e2f51798 :: mean +- SD over the row's 5 pretrained models
\newcommand{\FtAucJcESixOneeighteight}{0.9858\,\ensuremath{\pm}\,0.0001}  % experiments/FIGS/data/w2b_leg2_metrics_v2.json :: cells.{l188-s1,l188-s2,l188-s3,l188-s4,l188-s5}.N1000000.s1.macro_auc_ovr :: sha256 9cd059d3e2f51798 :: mean +- SD over the row's 5 pretrained models
\newcommand{\FtAucJcEThreeOnesixtwo}{0.9526\,\ensuremath{\pm}\,0.0020}  % experiments/FIGS/data/w2b_leg2_metrics_v2.json :: cells.{l162-s1b,l162-s2,l162-s3,l162-s4,l162-s5}.N1000.s1.macro_auc_ovr :: sha256 9cd059d3e2f51798 :: mean +- SD over the row's 5 pretrained models
\newcommand{\FtAucJcEFourOnesixtwo}{0.9770\,\ensuremath{\pm}\,0.0007}  % experiments/FIGS/data/w2b_leg2_metrics_v2.json :: cells.{l162-s1b,l162-s2,l162-s3,l162-s4,l162-s5}.N10000.s1.macro_auc_ovr :: sha256 9cd059d3e2f51798 :: mean +- SD over the row's 5 pretrained models
\newcommand{\FtAucJcEFiveOnesixtwo}{0.9836\,\ensuremath{\pm}\,0.0001}  % experiments/FIGS/data/w2b_leg2_metrics_v2.json :: cells.{l162-s1b,l162-s2,l162-s3,l162-s4,l162-s5}.N100000.s1.macro_auc_ovr :: sha256 9cd059d3e2f51798 :: mean +- SD over the row's 5 pretrained models
\newcommand{\FtAucJcESixOnesixtwo}{0.9858\,\ensuremath{\pm}\,0.0001}  % experiments/FIGS/data/w2b_leg2_metrics_v2.json :: cells.{l162-s1b,l162-s2,l162-s3,l162-s4,l162-s5}.N1000000.s1.macro_auc_ovr :: sha256 9cd059d3e2f51798 :: mean +- SD over the row's 5 pretrained models
\newcommand{\FtAucJcEThreeFourthree}{0.950\,\ensuremath{\pm}\,0.004}  % experiments/FIGS/data/w2b_leg2_metrics_v2.json :: cells.{r42q1-s1,r42q1-s2,r42q1-s3,r42q1-s4,r42q1-s5}.N1000.s1.macro_auc_ovr :: sha256 9cd059d3e2f51798 :: mean +- SD over the row's 5 pretrained models
\newcommand{\FtAucJcEFourFourthree}{0.9769\,\ensuremath{\pm}\,0.0005}  % experiments/FIGS/data/w2b_leg2_metrics_v2.json :: cells.{r42q1-s1,r42q1-s2,r42q1-s3,r42q1-s4,r42q1-s5}.N10000.s1.macro_auc_ovr :: sha256 9cd059d3e2f51798 :: mean +- SD over the row's 5 pretrained models
\newcommand{\FtAucJcEFiveFourthree}{0.9831\,\ensuremath{\pm}\,0.0002}  % experiments/FIGS/data/w2b_leg2_metrics_v2.json :: cells.{r42q1-s1,r42q1-s2,r42q1-s3,r42q1-s4,r42q1-s5}.N100000.s1.macro_auc_ovr :: sha256 9cd059d3e2f51798 :: mean +- SD over the row's 5 pretrained models
\newcommand{\FtAucJcESixFourthree}{0.9855\,\ensuremath{\pm}\,0.0001}  % experiments/FIGS/data/w2b_leg2_metrics_v2.json :: cells.{r42q1-s1,r42q1-s2,r42q1-s3,r42q1-s4,r42q1-s5}.N1000000.s1.macro_auc_ovr :: sha256 9cd059d3e2f51798 :: mean +- SD over the row's 5 pretrained models
\newcommand{\FtAucJcEThreeOneseven}{0.9470\,\ensuremath{\pm}\,0.0024}  % experiments/FIGS/data/w2b_leg2_metrics_v2.json :: cells.{r16q1-s1,r16q1-s2,r16q1-s3,r16q1-s4,r16q1-s5}.N1000.s1.macro_auc_ovr :: sha256 9cd059d3e2f51798 :: mean +- SD over the row's 5 pretrained models
\newcommand{\FtAucJcEFourOneseven}{0.9749\,\ensuremath{\pm}\,0.0006}  % experiments/FIGS/data/w2b_leg2_metrics_v2.json :: cells.{r16q1-s1,r16q1-s2,r16q1-s3,r16q1-s4,r16q1-s5}.N10000.s1.macro_auc_ovr :: sha256 9cd059d3e2f51798 :: mean +- SD over the row's 5 pretrained models
\newcommand{\FtAucJcEFiveOneseven}{0.9821\,\ensuremath{\pm}\,\ensuremath{<}0.0001}  % experiments/FIGS/data/w2b_leg2_metrics_v2.json :: cells.{r16q1-s1,r16q1-s2,r16q1-s3,r16q1-s4,r16q1-s5}.N100000.s1.macro_auc_ovr :: sha256 9cd059d3e2f51798 :: mean +- SD over the row's 5 pretrained models
\newcommand{\FtAucJcESixOneseven}{0.9846\,\ensuremath{\pm}\,0.0001}  % experiments/FIGS/data/w2b_leg2_metrics_v2.json :: cells.{r16q1-s1,r16q1-s2,r16q1-s3,r16q1-s4,r16q1-s5}.N1000000.s1.macro_auc_ovr :: sha256 9cd059d3e2f51798 :: mean +- SD over the row's 5 pretrained models
\newcommand{\FtAucJcEThreeOnesixtwoMass}{0.9569\,\ensuremath{\pm}\,0.0016}  % experiments/FIGS/data/w2b_leg2_metrics_v2.json :: cells.{l162mass-s1,l162mass-s2,l162mass-s3,l162mass-s4,l162mass-s5}.N1000.s1.macro_auc_ovr :: sha256 9cd059d3e2f51798 :: mean +- SD over the row's 5 pretrained models
\newcommand{\FtAucJcEFourOnesixtwoMass}{0.9768\,\ensuremath{\pm}\,0.0005}  % experiments/FIGS/data/w2b_leg2_metrics_v2.json :: cells.{l162mass-s1,l162mass-s2,l162mass-s3,l162mass-s4,l162mass-s5}.N10000.s1.macro_auc_ovr :: sha256 9cd059d3e2f51798 :: mean +- SD over the row's 5 pretrained models
\newcommand{\FtAucJcEFiveOnesixtwoMass}{0.9834\,\ensuremath{\pm}\,0.0001}  % experiments/FIGS/data/w2b_leg2_metrics_v2.json :: cells.{l162mass-s1,l162mass-s2,l162mass-s3,l162mass-s4,l162mass-s5}.N100000.s1.macro_auc_ovr :: sha256 9cd059d3e2f51798 :: mean +- SD over the row's 5 pretrained models
\newcommand{\FtAucJcESixOnesixtwoMass}{0.9856\,\ensuremath{\pm}\,0.0001}  % experiments/FIGS/data/w2b_leg2_metrics_v2.json :: cells.{l162mass-s1,l162mass-s2,l162mass-s3,l162mass-s4,l162mass-s5}.N1000000.s1.macro_auc_ovr :: sha256 9cd059d3e2f51798 :: mean +- SD over the row's 5 pretrained models
\newcommand{\FtAucJcEThreeOnesevenMass}{0.9460\,\ensuremath{\pm}\,0.0029}  % experiments/FIGS/data/w2b_leg2_metrics_v2.json :: cells.{r16q1mass-s1,r16q1mass-s2,r16q1mass-s3,r16q1mass-s4,r16q1mass-s5}.N1000.s1.macro_auc_ovr :: sha256 9cd059d3e2f51798 :: mean +- SD over the row's 5 pretrained models
\newcommand{\FtAucJcEFourOnesevenMass}{0.9737\,\ensuremath{\pm}\,0.0008}  % experiments/FIGS/data/w2b_leg2_metrics_v2.json :: cells.{r16q1mass-s1,r16q1mass-s2,r16q1mass-s3,r16q1mass-s4,r16q1mass-s5}.N10000.s1.macro_auc_ovr :: sha256 9cd059d3e2f51798 :: mean +- SD over the row's 5 pretrained models
\newcommand{\FtAucJcEFiveOnesevenMass}{0.9817\,\ensuremath{\pm}\,0.0002}  % experiments/FIGS/data/w2b_leg2_metrics_v2.json :: cells.{r16q1mass-s1,r16q1mass-s2,r16q1mass-s3,r16q1mass-s4,r16q1mass-s5}.N100000.s1.macro_auc_ovr :: sha256 9cd059d3e2f51798 :: mean +- SD over the row's 5 pretrained models
\newcommand{\FtAucJcESixOnesevenMass}{0.9842\,\ensuremath{\pm}\,0.0002}  % experiments/FIGS/data/w2b_leg2_metrics_v2.json :: cells.{r16q1mass-s1,r16q1mass-s2,r16q1mass-s3,r16q1mass-s4,r16q1mass-s5}.N1000000.s1.macro_auc_ovr :: sha256 9cd059d3e2f51798 :: mean +- SD over the row's 5 pretrained models
\newcommand{\FtAucJcEThreeRandDrawOne}{0.9477}  % experiments/FIGS/data/w2b_leg2_metrics_v2.json :: cells.{rand-d1-s1b}.N1000.s1.macro_auc_ovr :: sha256 9cd059d3e2f51798 :: one run
\newcommand{\FtAucJcEFourRandDrawOne}{0.9745}  % experiments/FIGS/data/w2b_leg2_metrics_v2.json :: cells.{rand-d1-s1b}.N10000.s1.macro_auc_ovr :: sha256 9cd059d3e2f51798 :: one run
\newcommand{\FtAucJcEFiveRandDrawOne}{0.98265}  % experiments/FIGS/data/w2b_leg2_metrics_v2.json :: cells.{rand-d1-s1b}.N100000.s1.macro_auc_ovr :: sha256 9cd059d3e2f51798 :: one run
\newcommand{\FtAucJcESixRandDrawOne}{0.98513}  % experiments/FIGS/data/w2b_leg2_metrics_v2.json :: cells.{rand-d1-s1b}.N1000000.s1.macro_auc_ovr :: sha256 9cd059d3e2f51798 :: one run
\newcommand{\FtAucJcEThreeRandDrawTwo}{0.9501}  % experiments/FIGS/data/w2b_leg2_metrics_v2.json :: cells.{rand-d2-s2}.N1000.s1.macro_auc_ovr :: sha256 9cd059d3e2f51798 :: one run
\newcommand{\FtAucJcEFourRandDrawTwo}{0.9756}  % experiments/FIGS/data/w2b_leg2_metrics_v2.json :: cells.{rand-d2-s2}.N10000.s1.macro_auc_ovr :: sha256 9cd059d3e2f51798 :: one run
\newcommand{\FtAucJcEFiveRandDrawTwo}{0.98278}  % experiments/FIGS/data/w2b_leg2_metrics_v2.json :: cells.{rand-d2-s2}.N100000.s1.macro_auc_ovr :: sha256 9cd059d3e2f51798 :: one run
\newcommand{\FtAucJcESixRandDrawTwo}{0.98512}  % experiments/FIGS/data/w2b_leg2_metrics_v2.json :: cells.{rand-d2-s2}.N1000000.s1.macro_auc_ovr :: sha256 9cd059d3e2f51798 :: one run
\newcommand{\FtAucJcEThreeRandDrawThree}{0.9470}  % experiments/FIGS/data/w2b_leg2_metrics_v2.json :: cells.{rand-d3-s3}.N1000.s1.macro_auc_ovr :: sha256 9cd059d3e2f51798 :: one run
\newcommand{\FtAucJcEFourRandDrawThree}{0.9743}  % experiments/FIGS/data/w2b_leg2_metrics_v2.json :: cells.{rand-d3-s3}.N10000.s1.macro_auc_ovr :: sha256 9cd059d3e2f51798 :: one run
\newcommand{\FtAucJcEFiveRandDrawThree}{0.98253}  % experiments/FIGS/data/w2b_leg2_metrics_v2.json :: cells.{rand-d3-s3}.N100000.s1.macro_auc_ovr :: sha256 9cd059d3e2f51798 :: one run
\newcommand{\FtAucJcESixRandDrawThree}{0.98506}  % experiments/FIGS/data/w2b_leg2_metrics_v2.json :: cells.{rand-d3-s3}.N1000000.s1.macro_auc_ovr :: sha256 9cd059d3e2f51798 :: one run
\newcommand{\FtAccJcEThreeOneeighteight}{0.7183\,\ensuremath{\pm}\,0.0032}  % experiments/FIGS/data/w2b_leg2_metrics_v2.json :: cells.{l188-s1,l188-s2,l188-s3,l188-s4,l188-s5}.N1000.s1.accuracy :: sha256 9cd059d3e2f51798 :: mean +- SD over the row's 5 pretrained models
\newcommand{\FtAccJcEFourOneeighteight}{0.8039\,\ensuremath{\pm}\,0.0030}  % experiments/FIGS/data/w2b_leg2_metrics_v2.json :: cells.{l188-s1,l188-s2,l188-s3,l188-s4,l188-s5}.N10000.s1.accuracy :: sha256 9cd059d3e2f51798 :: mean +- SD over the row's 5 pretrained models
\newcommand{\FtAccJcEFiveOneeighteight}{0.8351\,\ensuremath{\pm}\,0.0008}  % experiments/FIGS/data/w2b_leg2_metrics_v2.json :: cells.{l188-s1,l188-s2,l188-s3,l188-s4,l188-s5}.N100000.s1.accuracy :: sha256 9cd059d3e2f51798 :: mean +- SD over the row's 5 pretrained models
\newcommand{\FtAccJcESixOneeighteight}{0.8486\,\ensuremath{\pm}\,0.0005}  % experiments/FIGS/data/w2b_leg2_metrics_v2.json :: cells.{l188-s1,l188-s2,l188-s3,l188-s4,l188-s5}.N1000000.s1.accuracy :: sha256 9cd059d3e2f51798 :: mean +- SD over the row's 5 pretrained models
\newcommand{\FtAccJcEThreeOnesixtwo}{0.7268\,\ensuremath{\pm}\,0.0017}  % experiments/FIGS/data/w2b_leg2_metrics_v2.json :: cells.{l162-s1b,l162-s2,l162-s3,l162-s4,l162-s5}.N1000.s1.accuracy :: sha256 9cd059d3e2f51798 :: mean +- SD over the row's 5 pretrained models
\newcommand{\FtAccJcEFourOnesixtwo}{0.8069\,\ensuremath{\pm}\,0.0025}  % experiments/FIGS/data/w2b_leg2_metrics_v2.json :: cells.{l162-s1b,l162-s2,l162-s3,l162-s4,l162-s5}.N10000.s1.accuracy :: sha256 9cd059d3e2f51798 :: mean +- SD over the row's 5 pretrained models
\newcommand{\FtAccJcEFiveOnesixtwo}{0.8363\,\ensuremath{\pm}\,0.0008}  % experiments/FIGS/data/w2b_leg2_metrics_v2.json :: cells.{l162-s1b,l162-s2,l162-s3,l162-s4,l162-s5}.N100000.s1.accuracy :: sha256 9cd059d3e2f51798 :: mean +- SD over the row's 5 pretrained models
\newcommand{\FtAccJcESixOnesixtwo}{0.8485\,\ensuremath{\pm}\,0.0005}  % experiments/FIGS/data/w2b_leg2_metrics_v2.json :: cells.{l162-s1b,l162-s2,l162-s3,l162-s4,l162-s5}.N1000000.s1.accuracy :: sha256 9cd059d3e2f51798 :: mean +- SD over the row's 5 pretrained models
\newcommand{\FtAccJcEThreeFourthree}{0.722\,\ensuremath{\pm}\,0.007}  % experiments/FIGS/data/w2b_leg2_metrics_v2.json :: cells.{r42q1-s1,r42q1-s2,r42q1-s3,r42q1-s4,r42q1-s5}.N1000.s1.accuracy :: sha256 9cd059d3e2f51798 :: mean +- SD over the row's 5 pretrained models
\newcommand{\FtAccJcEFourFourthree}{0.8050\,\ensuremath{\pm}\,0.0026}  % experiments/FIGS/data/w2b_leg2_metrics_v2.json :: cells.{r42q1-s1,r42q1-s2,r42q1-s3,r42q1-s4,r42q1-s5}.N10000.s1.accuracy :: sha256 9cd059d3e2f51798 :: mean +- SD over the row's 5 pretrained models
\newcommand{\FtAccJcEFiveFourthree}{0.8336\,\ensuremath{\pm}\,0.0008}  % experiments/FIGS/data/w2b_leg2_metrics_v2.json :: cells.{r42q1-s1,r42q1-s2,r42q1-s3,r42q1-s4,r42q1-s5}.N100000.s1.accuracy :: sha256 9cd059d3e2f51798 :: mean +- SD over the row's 5 pretrained models
\newcommand{\FtAccJcESixFourthree}{0.8468\,\ensuremath{\pm}\,0.0003}  % experiments/FIGS/data/w2b_leg2_metrics_v2.json :: cells.{r42q1-s1,r42q1-s2,r42q1-s3,r42q1-s4,r42q1-s5}.N1000000.s1.accuracy :: sha256 9cd059d3e2f51798 :: mean +- SD over the row's 5 pretrained models
\newcommand{\FtAccJcEThreeOneseven}{0.714\,\ensuremath{\pm}\,0.005}  % experiments/FIGS/data/w2b_leg2_metrics_v2.json :: cells.{r16q1-s1,r16q1-s2,r16q1-s3,r16q1-s4,r16q1-s5}.N1000.s1.accuracy :: sha256 9cd059d3e2f51798 :: mean +- SD over the row's 5 pretrained models
\newcommand{\FtAccJcEFourOneseven}{0.7964\,\ensuremath{\pm}\,0.0019}  % experiments/FIGS/data/w2b_leg2_metrics_v2.json :: cells.{r16q1-s1,r16q1-s2,r16q1-s3,r16q1-s4,r16q1-s5}.N10000.s1.accuracy :: sha256 9cd059d3e2f51798 :: mean +- SD over the row's 5 pretrained models
\newcommand{\FtAccJcEFiveOneseven}{0.8286\,\ensuremath{\pm}\,0.0006}  % experiments/FIGS/data/w2b_leg2_metrics_v2.json :: cells.{r16q1-s1,r16q1-s2,r16q1-s3,r16q1-s4,r16q1-s5}.N100000.s1.accuracy :: sha256 9cd059d3e2f51798 :: mean +- SD over the row's 5 pretrained models
\newcommand{\FtAccJcESixOneseven}{0.8415\,\ensuremath{\pm}\,0.0007}  % experiments/FIGS/data/w2b_leg2_metrics_v2.json :: cells.{r16q1-s1,r16q1-s2,r16q1-s3,r16q1-s4,r16q1-s5}.N1000000.s1.accuracy :: sha256 9cd059d3e2f51798 :: mean +- SD over the row's 5 pretrained models
\newcommand{\FtAccJcEThreeOnesixtwoMass}{0.728\,\ensuremath{\pm}\,0.004}  % experiments/FIGS/data/w2b_leg2_metrics_v2.json :: cells.{l162mass-s1,l162mass-s2,l162mass-s3,l162mass-s4,l162mass-s5}.N1000.s1.accuracy :: sha256 9cd059d3e2f51798 :: mean +- SD over the row's 5 pretrained models
\newcommand{\FtAccJcEFourOnesixtwoMass}{0.8043\,\ensuremath{\pm}\,0.0020}  % experiments/FIGS/data/w2b_leg2_metrics_v2.json :: cells.{l162mass-s1,l162mass-s2,l162mass-s3,l162mass-s4,l162mass-s5}.N10000.s1.accuracy :: sha256 9cd059d3e2f51798 :: mean +- SD over the row's 5 pretrained models
\newcommand{\FtAccJcEFiveOnesixtwoMass}{0.8349\,\ensuremath{\pm}\,0.0005}  % experiments/FIGS/data/w2b_leg2_metrics_v2.json :: cells.{l162mass-s1,l162mass-s2,l162mass-s3,l162mass-s4,l162mass-s5}.N100000.s1.accuracy :: sha256 9cd059d3e2f51798 :: mean +- SD over the row's 5 pretrained models
\newcommand{\FtAccJcESixOnesixtwoMass}{0.8475\,\ensuremath{\pm}\,0.0009}  % experiments/FIGS/data/w2b_leg2_metrics_v2.json :: cells.{l162mass-s1,l162mass-s2,l162mass-s3,l162mass-s4,l162mass-s5}.N1000000.s1.accuracy :: sha256 9cd059d3e2f51798 :: mean +- SD over the row's 5 pretrained models
\newcommand{\FtAccJcEThreeOnesevenMass}{0.702\,\ensuremath{\pm}\,0.005}  % experiments/FIGS/data/w2b_leg2_metrics_v2.json :: cells.{r16q1mass-s1,r16q1mass-s2,r16q1mass-s3,r16q1mass-s4,r16q1mass-s5}.N1000.s1.accuracy :: sha256 9cd059d3e2f51798 :: mean +- SD over the row's 5 pretrained models
\newcommand{\FtAccJcEFourOnesevenMass}{0.7909\,\ensuremath{\pm}\,0.0031}  % experiments/FIGS/data/w2b_leg2_metrics_v2.json :: cells.{r16q1mass-s1,r16q1mass-s2,r16q1mass-s3,r16q1mass-s4,r16q1mass-s5}.N10000.s1.accuracy :: sha256 9cd059d3e2f51798 :: mean +- SD over the row's 5 pretrained models
\newcommand{\FtAccJcEFiveOnesevenMass}{0.8265\,\ensuremath{\pm}\,0.0007}  % experiments/FIGS/data/w2b_leg2_metrics_v2.json :: cells.{r16q1mass-s1,r16q1mass-s2,r16q1mass-s3,r16q1mass-s4,r16q1mass-s5}.N100000.s1.accuracy :: sha256 9cd059d3e2f51798 :: mean +- SD over the row's 5 pretrained models
\newcommand{\FtAccJcESixOnesevenMass}{0.8390\,\ensuremath{\pm}\,0.0011}  % experiments/FIGS/data/w2b_leg2_metrics_v2.json :: cells.{r16q1mass-s1,r16q1mass-s2,r16q1mass-s3,r16q1mass-s4,r16q1mass-s5}.N1000000.s1.accuracy :: sha256 9cd059d3e2f51798 :: mean +- SD over the row's 5 pretrained models
\newcommand{\FtAccJcEThreeRandDrawOne}{0.7049}  % experiments/FIGS/data/w2b_leg2_metrics_v2.json :: cells.{rand-d1-s1b}.N1000.s1.accuracy :: sha256 9cd059d3e2f51798 :: one run
\newcommand{\FtAccJcEFourRandDrawOne}{0.7974}  % experiments/FIGS/data/w2b_leg2_metrics_v2.json :: cells.{rand-d1-s1b}.N10000.s1.accuracy :: sha256 9cd059d3e2f51798 :: one run
\newcommand{\FtAccJcEFiveRandDrawOne}{0.8326}  % experiments/FIGS/data/w2b_leg2_metrics_v2.json :: cells.{rand-d1-s1b}.N100000.s1.accuracy :: sha256 9cd059d3e2f51798 :: one run
\newcommand{\FtAccJcESixRandDrawOne}{0.84467}  % experiments/FIGS/data/w2b_leg2_metrics_v2.json :: cells.{rand-d1-s1b}.N1000000.s1.accuracy :: sha256 9cd059d3e2f51798 :: one run
\newcommand{\FtAccJcEThreeRandDrawTwo}{0.7135}  % experiments/FIGS/data/w2b_leg2_metrics_v2.json :: cells.{rand-d2-s2}.N1000.s1.accuracy :: sha256 9cd059d3e2f51798 :: one run
\newcommand{\FtAccJcEFourRandDrawTwo}{0.7998}  % experiments/FIGS/data/w2b_leg2_metrics_v2.json :: cells.{rand-d2-s2}.N10000.s1.accuracy :: sha256 9cd059d3e2f51798 :: one run
\newcommand{\FtAccJcEFiveRandDrawTwo}{0.8332}  % experiments/FIGS/data/w2b_leg2_metrics_v2.json :: cells.{rand-d2-s2}.N100000.s1.accuracy :: sha256 9cd059d3e2f51798 :: one run
\newcommand{\FtAccJcESixRandDrawTwo}{0.84449}  % experiments/FIGS/data/w2b_leg2_metrics_v2.json :: cells.{rand-d2-s2}.N1000000.s1.accuracy :: sha256 9cd059d3e2f51798 :: one run
\newcommand{\FtAccJcEThreeRandDrawThree}{0.7090}  % experiments/FIGS/data/w2b_leg2_metrics_v2.json :: cells.{rand-d3-s3}.N1000.s1.accuracy :: sha256 9cd059d3e2f51798 :: one run
\newcommand{\FtAccJcEFourRandDrawThree}{0.7948}  % experiments/FIGS/data/w2b_leg2_metrics_v2.json :: cells.{rand-d3-s3}.N10000.s1.accuracy :: sha256 9cd059d3e2f51798 :: one run
\newcommand{\FtAccJcEFiveRandDrawThree}{0.8310}  % experiments/FIGS/data/w2b_leg2_metrics_v2.json :: cells.{rand-d3-s3}.N100000.s1.accuracy :: sha256 9cd059d3e2f51798 :: one run
\newcommand{\FtAccJcESixRandDrawThree}{0.84409}  % experiments/FIGS/data/w2b_leg2_metrics_v2.json :: cells.{rand-d3-s3}.N1000000.s1.accuracy :: sha256 9cd059d3e2f51798 :: one run
\newcommand{\FtOmaRatioJcEThreeOnesixtwoOverOneeighteight}{0.96}  % experiments/FIGS/data/w2b_leg2_metrics_v2.json :: cells.*.N1000.s1.macro_auc_ovr :: sha256 9cd059d3e2f51798 :: ratio of the seed means of 1 - macro AUC, this row over the reference
\newcommand{\FtOmaRatioJcEThreeFourthreeOverOneeighteight}{1.01}  % experiments/FIGS/data/w2b_leg2_metrics_v2.json :: cells.*.N1000.s1.macro_auc_ovr :: sha256 9cd059d3e2f51798 :: ratio of the seed means of 1 - macro AUC, this row over the reference
\newcommand{\FtOmaRatioJcEThreeOnesevenOverOneeighteight}{1.07}  % experiments/FIGS/data/w2b_leg2_metrics_v2.json :: cells.*.N1000.s1.macro_auc_ovr :: sha256 9cd059d3e2f51798 :: ratio of the seed means of 1 - macro AUC, this row over the reference
\newcommand{\FtOmaRatioJcEThreeOnesixtwoMassOverOneeighteight}{0.87}  % experiments/FIGS/data/w2b_leg2_metrics_v2.json :: cells.*.N1000.s1.macro_auc_ovr :: sha256 9cd059d3e2f51798 :: ratio of the seed means of 1 - macro AUC, this row over the reference
\newcommand{\FtOmaRatioJcEThreeOnesevenMassOverOneeighteight}{1.09}  % experiments/FIGS/data/w2b_leg2_metrics_v2.json :: cells.*.N1000.s1.macro_auc_ovr :: sha256 9cd059d3e2f51798 :: ratio of the seed means of 1 - macro AUC, this row over the reference
\newcommand{\FtOmaRatioJcEThreeRandDrawOneOverOneeighteight}{1.06}  % experiments/FIGS/data/w2b_leg2_metrics_v2.json :: cells.*.N1000.s1.macro_auc_ovr :: sha256 9cd059d3e2f51798 :: ratio of the seed means of 1 - macro AUC, this row over the reference
\newcommand{\FtOmaRatioJcEThreeRandDrawTwoOverOneeighteight}{1.01}  % experiments/FIGS/data/w2b_leg2_metrics_v2.json :: cells.*.N1000.s1.macro_auc_ovr :: sha256 9cd059d3e2f51798 :: ratio of the seed means of 1 - macro AUC, this row over the reference
\newcommand{\FtOmaRatioJcEThreeRandDrawThreeOverOneeighteight}{1.07}  % experiments/FIGS/data/w2b_leg2_metrics_v2.json :: cells.*.N1000.s1.macro_auc_ovr :: sha256 9cd059d3e2f51798 :: ratio of the seed means of 1 - macro AUC, this row over the reference
\newcommand{\FtOmaRatioJcEThreeOneeighteightOverOnesixtwo}{1.05}  % experiments/FIGS/data/w2b_leg2_metrics_v2.json :: cells.*.N1000.s1.macro_auc_ovr :: sha256 9cd059d3e2f51798 :: ratio of the seed means of 1 - macro AUC, this row over the reference
\newcommand{\FtOmaRatioJcEThreeFourthreeOverOnesixtwo}{1.05}  % experiments/FIGS/data/w2b_leg2_metrics_v2.json :: cells.*.N1000.s1.macro_auc_ovr :: sha256 9cd059d3e2f51798 :: ratio of the seed means of 1 - macro AUC, this row over the reference
\newcommand{\FtOmaRatioJcEThreeOnesevenOverOnesixtwo}{1.12}  % experiments/FIGS/data/w2b_leg2_metrics_v2.json :: cells.*.N1000.s1.macro_auc_ovr :: sha256 9cd059d3e2f51798 :: ratio of the seed means of 1 - macro AUC, this row over the reference
\newcommand{\FtOmaRatioJcEThreeOnesixtwoMassOverOnesixtwo}{0.91}  % experiments/FIGS/data/w2b_leg2_metrics_v2.json :: cells.*.N1000.s1.macro_auc_ovr :: sha256 9cd059d3e2f51798 :: ratio of the seed means of 1 - macro AUC, this row over the reference
\newcommand{\FtOmaRatioJcEThreeOnesevenMassOverOnesixtwo}{1.14}  % experiments/FIGS/data/w2b_leg2_metrics_v2.json :: cells.*.N1000.s1.macro_auc_ovr :: sha256 9cd059d3e2f51798 :: ratio of the seed means of 1 - macro AUC, this row over the reference
\newcommand{\FtOmaRatioJcEThreeRandDrawOneOverOnesixtwo}{1.10}  % experiments/FIGS/data/w2b_leg2_metrics_v2.json :: cells.*.N1000.s1.macro_auc_ovr :: sha256 9cd059d3e2f51798 :: ratio of the seed means of 1 - macro AUC, this row over the reference
\newcommand{\FtOmaRatioJcEThreeRandDrawTwoOverOnesixtwo}{1.05}  % experiments/FIGS/data/w2b_leg2_metrics_v2.json :: cells.*.N1000.s1.macro_auc_ovr :: sha256 9cd059d3e2f51798 :: ratio of the seed means of 1 - macro AUC, this row over the reference
\newcommand{\FtOmaRatioJcEThreeRandDrawThreeOverOnesixtwo}{1.12}  % experiments/FIGS/data/w2b_leg2_metrics_v2.json :: cells.*.N1000.s1.macro_auc_ovr :: sha256 9cd059d3e2f51798 :: ratio of the seed means of 1 - macro AUC, this row over the reference
\newcommand{\FtOmaRatioJcEFourOnesixtwoOverOneeighteight}{0.98}  % experiments/FIGS/data/w2b_leg2_metrics_v2.json :: cells.*.N10000.s1.macro_auc_ovr :: sha256 9cd059d3e2f51798 :: ratio of the seed means of 1 - macro AUC, this row over the reference
\newcommand{\FtOmaRatioJcEFourFourthreeOverOneeighteight}{0.98}  % experiments/FIGS/data/w2b_leg2_metrics_v2.json :: cells.*.N10000.s1.macro_auc_ovr :: sha256 9cd059d3e2f51798 :: ratio of the seed means of 1 - macro AUC, this row over the reference
\newcommand{\FtOmaRatioJcEFourOnesevenOverOneeighteight}{1.07}  % experiments/FIGS/data/w2b_leg2_metrics_v2.json :: cells.*.N10000.s1.macro_auc_ovr :: sha256 9cd059d3e2f51798 :: ratio of the seed means of 1 - macro AUC, this row over the reference
\newcommand{\FtOmaRatioJcEFourOnesixtwoMassOverOneeighteight}{0.99}  % experiments/FIGS/data/w2b_leg2_metrics_v2.json :: cells.*.N10000.s1.macro_auc_ovr :: sha256 9cd059d3e2f51798 :: ratio of the seed means of 1 - macro AUC, this row over the reference
\newcommand{\FtOmaRatioJcEFourOnesevenMassOverOneeighteight}{1.12}  % experiments/FIGS/data/w2b_leg2_metrics_v2.json :: cells.*.N10000.s1.macro_auc_ovr :: sha256 9cd059d3e2f51798 :: ratio of the seed means of 1 - macro AUC, this row over the reference
\newcommand{\FtOmaRatioJcEFourRandDrawOneOverOneeighteight}{1.08}  % experiments/FIGS/data/w2b_leg2_metrics_v2.json :: cells.*.N10000.s1.macro_auc_ovr :: sha256 9cd059d3e2f51798 :: ratio of the seed means of 1 - macro AUC, this row over the reference
\newcommand{\FtOmaRatioJcEFourRandDrawTwoOverOneeighteight}{1.04}  % experiments/FIGS/data/w2b_leg2_metrics_v2.json :: cells.*.N10000.s1.macro_auc_ovr :: sha256 9cd059d3e2f51798 :: ratio of the seed means of 1 - macro AUC, this row over the reference
\newcommand{\FtOmaRatioJcEFourRandDrawThreeOverOneeighteight}{1.09}  % experiments/FIGS/data/w2b_leg2_metrics_v2.json :: cells.*.N10000.s1.macro_auc_ovr :: sha256 9cd059d3e2f51798 :: ratio of the seed means of 1 - macro AUC, this row over the reference
\newcommand{\FtOmaRatioJcEFourOneeighteightOverOnesixtwo}{1.02}  % experiments/FIGS/data/w2b_leg2_metrics_v2.json :: cells.*.N10000.s1.macro_auc_ovr :: sha256 9cd059d3e2f51798 :: ratio of the seed means of 1 - macro AUC, this row over the reference
\newcommand{\FtOmaRatioJcEFourFourthreeOverOnesixtwo}{1.01}  % experiments/FIGS/data/w2b_leg2_metrics_v2.json :: cells.*.N10000.s1.macro_auc_ovr :: sha256 9cd059d3e2f51798 :: ratio of the seed means of 1 - macro AUC, this row over the reference
\newcommand{\FtOmaRatioJcEFourOnesevenOverOnesixtwo}{1.09}  % experiments/FIGS/data/w2b_leg2_metrics_v2.json :: cells.*.N10000.s1.macro_auc_ovr :: sha256 9cd059d3e2f51798 :: ratio of the seed means of 1 - macro AUC, this row over the reference
\newcommand{\FtOmaRatioJcEFourOnesixtwoMassOverOnesixtwo}{1.01}  % experiments/FIGS/data/w2b_leg2_metrics_v2.json :: cells.*.N10000.s1.macro_auc_ovr :: sha256 9cd059d3e2f51798 :: ratio of the seed means of 1 - macro AUC, this row over the reference
\newcommand{\FtOmaRatioJcEFourOnesevenMassOverOnesixtwo}{1.14}  % experiments/FIGS/data/w2b_leg2_metrics_v2.json :: cells.*.N10000.s1.macro_auc_ovr :: sha256 9cd059d3e2f51798 :: ratio of the seed means of 1 - macro AUC, this row over the reference
\newcommand{\FtOmaRatioJcEFourRandDrawOneOverOnesixtwo}{1.11}  % experiments/FIGS/data/w2b_leg2_metrics_v2.json :: cells.*.N10000.s1.macro_auc_ovr :: sha256 9cd059d3e2f51798 :: ratio of the seed means of 1 - macro AUC, this row over the reference
\newcommand{\FtOmaRatioJcEFourRandDrawTwoOverOnesixtwo}{1.06}  % experiments/FIGS/data/w2b_leg2_metrics_v2.json :: cells.*.N10000.s1.macro_auc_ovr :: sha256 9cd059d3e2f51798 :: ratio of the seed means of 1 - macro AUC, this row over the reference
\newcommand{\FtOmaRatioJcEFourRandDrawThreeOverOnesixtwo}{1.12}  % experiments/FIGS/data/w2b_leg2_metrics_v2.json :: cells.*.N10000.s1.macro_auc_ovr :: sha256 9cd059d3e2f51798 :: ratio of the seed means of 1 - macro AUC, this row over the reference
\newcommand{\FtOmaRatioJcEFiveOnesixtwoOverOneeighteight}{0.98}  % experiments/FIGS/data/w2b_leg2_metrics_v2.json :: cells.*.N100000.s1.macro_auc_ovr :: sha256 9cd059d3e2f51798 :: ratio of the seed means of 1 - macro AUC, this row over the reference
\newcommand{\FtOmaRatioJcEFiveFourthreeOverOneeighteight}{1.02}  % experiments/FIGS/data/w2b_leg2_metrics_v2.json :: cells.*.N100000.s1.macro_auc_ovr :: sha256 9cd059d3e2f51798 :: ratio of the seed means of 1 - macro AUC, this row over the reference
\newcommand{\FtOmaRatioJcEFiveOnesevenOverOneeighteight}{1.07}  % experiments/FIGS/data/w2b_leg2_metrics_v2.json :: cells.*.N100000.s1.macro_auc_ovr :: sha256 9cd059d3e2f51798 :: ratio of the seed means of 1 - macro AUC, this row over the reference
\newcommand{\FtOmaRatioJcEFiveOnesixtwoMassOverOneeighteight}{1.00}  % experiments/FIGS/data/w2b_leg2_metrics_v2.json :: cells.*.N100000.s1.macro_auc_ovr :: sha256 9cd059d3e2f51798 :: ratio of the seed means of 1 - macro AUC, this row over the reference
\newcommand{\FtOmaRatioJcEFiveOnesevenMassOverOneeighteight}{1.10}  % experiments/FIGS/data/w2b_leg2_metrics_v2.json :: cells.*.N100000.s1.macro_auc_ovr :: sha256 9cd059d3e2f51798 :: ratio of the seed means of 1 - macro AUC, this row over the reference
\newcommand{\FtOmaRatioJcEFiveRandDrawOneOverOneeighteight}{1.04}  % experiments/FIGS/data/w2b_leg2_metrics_v2.json :: cells.*.N100000.s1.macro_auc_ovr :: sha256 9cd059d3e2f51798 :: ratio of the seed means of 1 - macro AUC, this row over the reference
\newcommand{\FtOmaRatioJcEFiveRandDrawTwoOverOneeighteight}{1.03}  % experiments/FIGS/data/w2b_leg2_metrics_v2.json :: cells.*.N100000.s1.macro_auc_ovr :: sha256 9cd059d3e2f51798 :: ratio of the seed means of 1 - macro AUC, this row over the reference
\newcommand{\FtOmaRatioJcEFiveRandDrawThreeOverOneeighteight}{1.05}  % experiments/FIGS/data/w2b_leg2_metrics_v2.json :: cells.*.N100000.s1.macro_auc_ovr :: sha256 9cd059d3e2f51798 :: ratio of the seed means of 1 - macro AUC, this row over the reference
\newcommand{\FtOmaRatioJcEFiveOneeighteightOverOnesixtwo}{1.02}  % experiments/FIGS/data/w2b_leg2_metrics_v2.json :: cells.*.N100000.s1.macro_auc_ovr :: sha256 9cd059d3e2f51798 :: ratio of the seed means of 1 - macro AUC, this row over the reference
\newcommand{\FtOmaRatioJcEFiveFourthreeOverOnesixtwo}{1.03}  % experiments/FIGS/data/w2b_leg2_metrics_v2.json :: cells.*.N100000.s1.macro_auc_ovr :: sha256 9cd059d3e2f51798 :: ratio of the seed means of 1 - macro AUC, this row over the reference
\newcommand{\FtOmaRatioJcEFiveOnesevenOverOnesixtwo}{1.09}  % experiments/FIGS/data/w2b_leg2_metrics_v2.json :: cells.*.N100000.s1.macro_auc_ovr :: sha256 9cd059d3e2f51798 :: ratio of the seed means of 1 - macro AUC, this row over the reference
\newcommand{\FtOmaRatioJcEFiveOnesixtwoMassOverOnesixtwo}{1.02}  % experiments/FIGS/data/w2b_leg2_metrics_v2.json :: cells.*.N100000.s1.macro_auc_ovr :: sha256 9cd059d3e2f51798 :: ratio of the seed means of 1 - macro AUC, this row over the reference
\newcommand{\FtOmaRatioJcEFiveOnesevenMassOverOnesixtwo}{1.11}  % experiments/FIGS/data/w2b_leg2_metrics_v2.json :: cells.*.N100000.s1.macro_auc_ovr :: sha256 9cd059d3e2f51798 :: ratio of the seed means of 1 - macro AUC, this row over the reference
\newcommand{\FtOmaRatioJcEFiveRandDrawOneOverOnesixtwo}{1.06}  % experiments/FIGS/data/w2b_leg2_metrics_v2.json :: cells.*.N100000.s1.macro_auc_ovr :: sha256 9cd059d3e2f51798 :: ratio of the seed means of 1 - macro AUC, this row over the reference
\newcommand{\FtOmaRatioJcEFiveRandDrawTwoOverOnesixtwo}{1.05}  % experiments/FIGS/data/w2b_leg2_metrics_v2.json :: cells.*.N100000.s1.macro_auc_ovr :: sha256 9cd059d3e2f51798 :: ratio of the seed means of 1 - macro AUC, this row over the reference
\newcommand{\FtOmaRatioJcEFiveRandDrawThreeOverOnesixtwo}{1.07}  % experiments/FIGS/data/w2b_leg2_metrics_v2.json :: cells.*.N100000.s1.macro_auc_ovr :: sha256 9cd059d3e2f51798 :: ratio of the seed means of 1 - macro AUC, this row over the reference
\newcommand{\FtOmaRatioJcESixOnesixtwoOverOneeighteight}{1.00}  % experiments/FIGS/data/w2b_leg2_metrics_v2.json :: cells.*.N1000000.s1.macro_auc_ovr :: sha256 9cd059d3e2f51798 :: ratio of the seed means of 1 - macro AUC, this row over the reference
\newcommand{\FtOmaRatioJcESixFourthreeOverOneeighteight}{1.02}  % experiments/FIGS/data/w2b_leg2_metrics_v2.json :: cells.*.N1000000.s1.macro_auc_ovr :: sha256 9cd059d3e2f51798 :: ratio of the seed means of 1 - macro AUC, this row over the reference
\newcommand{\FtOmaRatioJcESixOnesevenOverOneeighteight}{1.08}  % experiments/FIGS/data/w2b_leg2_metrics_v2.json :: cells.*.N1000000.s1.macro_auc_ovr :: sha256 9cd059d3e2f51798 :: ratio of the seed means of 1 - macro AUC, this row over the reference
\newcommand{\FtOmaRatioJcESixOnesixtwoMassOverOneeighteight}{1.01}  % experiments/FIGS/data/w2b_leg2_metrics_v2.json :: cells.*.N1000000.s1.macro_auc_ovr :: sha256 9cd059d3e2f51798 :: ratio of the seed means of 1 - macro AUC, this row over the reference
\newcommand{\FtOmaRatioJcESixOnesevenMassOverOneeighteight}{1.11}  % experiments/FIGS/data/w2b_leg2_metrics_v2.json :: cells.*.N1000000.s1.macro_auc_ovr :: sha256 9cd059d3e2f51798 :: ratio of the seed means of 1 - macro AUC, this row over the reference
\newcommand{\FtOmaRatioJcESixRandDrawOneOverOneeighteight}{1.05}  % experiments/FIGS/data/w2b_leg2_metrics_v2.json :: cells.*.N1000000.s1.macro_auc_ovr :: sha256 9cd059d3e2f51798 :: ratio of the seed means of 1 - macro AUC, this row over the reference
\newcommand{\FtOmaRatioJcESixRandDrawTwoOverOneeighteight}{1.05}  % experiments/FIGS/data/w2b_leg2_metrics_v2.json :: cells.*.N1000000.s1.macro_auc_ovr :: sha256 9cd059d3e2f51798 :: ratio of the seed means of 1 - macro AUC, this row over the reference
\newcommand{\FtOmaRatioJcESixRandDrawThreeOverOneeighteight}{1.05}  % experiments/FIGS/data/w2b_leg2_metrics_v2.json :: cells.*.N1000000.s1.macro_auc_ovr :: sha256 9cd059d3e2f51798 :: ratio of the seed means of 1 - macro AUC, this row over the reference
\newcommand{\FtOmaRatioJcESixOneeighteightOverOnesixtwo}{1.00}  % experiments/FIGS/data/w2b_leg2_metrics_v2.json :: cells.*.N1000000.s1.macro_auc_ovr :: sha256 9cd059d3e2f51798 :: ratio of the seed means of 1 - macro AUC, this row over the reference
\newcommand{\FtOmaRatioJcESixFourthreeOverOnesixtwo}{1.02}  % experiments/FIGS/data/w2b_leg2_metrics_v2.json :: cells.*.N1000000.s1.macro_auc_ovr :: sha256 9cd059d3e2f51798 :: ratio of the seed means of 1 - macro AUC, this row over the reference
\newcommand{\FtOmaRatioJcESixOnesevenOverOnesixtwo}{1.08}  % experiments/FIGS/data/w2b_leg2_metrics_v2.json :: cells.*.N1000000.s1.macro_auc_ovr :: sha256 9cd059d3e2f51798 :: ratio of the seed means of 1 - macro AUC, this row over the reference
\newcommand{\FtOmaRatioJcESixOnesixtwoMassOverOnesixtwo}{1.01}  % experiments/FIGS/data/w2b_leg2_metrics_v2.json :: cells.*.N1000000.s1.macro_auc_ovr :: sha256 9cd059d3e2f51798 :: ratio of the seed means of 1 - macro AUC, this row over the reference
\newcommand{\FtOmaRatioJcESixOnesevenMassOverOnesixtwo}{1.11}  % experiments/FIGS/data/w2b_leg2_metrics_v2.json :: cells.*.N1000000.s1.macro_auc_ovr :: sha256 9cd059d3e2f51798 :: ratio of the seed means of 1 - macro AUC, this row over the reference
\newcommand{\FtOmaRatioJcESixRandDrawOneOverOnesixtwo}{1.05}  % experiments/FIGS/data/w2b_leg2_metrics_v2.json :: cells.*.N1000000.s1.macro_auc_ovr :: sha256 9cd059d3e2f51798 :: ratio of the seed means of 1 - macro AUC, this row over the reference
\newcommand{\FtOmaRatioJcESixRandDrawTwoOverOnesixtwo}{1.05}  % experiments/FIGS/data/w2b_leg2_metrics_v2.json :: cells.*.N1000000.s1.macro_auc_ovr :: sha256 9cd059d3e2f51798 :: ratio of the seed means of 1 - macro AUC, this row over the reference
\newcommand{\FtOmaRatioJcESixRandDrawThreeOverOnesixtwo}{1.05}  % experiments/FIGS/data/w2b_leg2_metrics_v2.json :: cells.*.N1000000.s1.macro_auc_ovr :: sha256 9cd059d3e2f51798 :: ratio of the seed means of 1 - macro AUC, this row over the reference
\newcommand{\PairedFtAccJcEThreeOnesevenMinusOneeighteight}{\ensuremath{-}0.0040~[\ensuremath{-}0.0144, +0.0064]}  % experiments/FIGS/data/w2b_leg2_metrics_v2.json :: cells.{r16q1-s1,r16q1-s2,r16q1-s3,r16q1-s4,r16q1-s5,l188-s1,l188-s2,l188-s3,l188-s4,l188-s5}.N1000.s1.accuracy :: sha256 9cd059d3e2f51798 :: accuracy, coarse minus fine, paired by run [95% Student-t interval]
\newcommand{\PairedFtAccJcEFourOnesevenMinusOneeighteight}{\ensuremath{-}0.0075~[\ensuremath{-}0.0118, \ensuremath{-}0.0033]}  % experiments/FIGS/data/w2b_leg2_metrics_v2.json :: cells.{r16q1-s1,r16q1-s2,r16q1-s3,r16q1-s4,r16q1-s5,l188-s1,l188-s2,l188-s3,l188-s4,l188-s5}.N10000.s1.accuracy :: sha256 9cd059d3e2f51798 :: accuracy, coarse minus fine, paired by run [95% Student-t interval]
\newcommand{\PairedFtAccJcEFiveOnesevenMinusOneeighteight}{\ensuremath{-}0.0066~[\ensuremath{-}0.0083, \ensuremath{-}0.0049]}  % experiments/FIGS/data/w2b_leg2_metrics_v2.json :: cells.{r16q1-s1,r16q1-s2,r16q1-s3,r16q1-s4,r16q1-s5,l188-s1,l188-s2,l188-s3,l188-s4,l188-s5}.N100000.s1.accuracy :: sha256 9cd059d3e2f51798 :: accuracy, coarse minus fine, paired by run [95% Student-t interval]
\newcommand{\PairedFtAccJcESixOnesevenMinusOneeighteight}{\ensuremath{-}0.0071~[\ensuremath{-}0.0080, \ensuremath{-}0.0062]}  % experiments/FIGS/data/w2b_leg2_metrics_v2.json :: cells.{r16q1-s1,r16q1-s2,r16q1-s3,r16q1-s4,r16q1-s5,l188-s1,l188-s2,l188-s3,l188-s4,l188-s5}.N1000000.s1.accuracy :: sha256 9cd059d3e2f51798 :: accuracy, coarse minus fine, paired by run [95% Student-t interval]
\newcommand{\FtNClassesJc}{10}  % experiments/FIGS/data/w2b_leg2_metrics_v2.json :: cells.*.*.s1.n_classes_present :: sha256 9cd059d3e2f51798 :: classes of the fine-tuning task
\newcommand{\FtNTestAucJc}{500{,}000}  % experiments/FIGS/data/w2b_leg2_metrics_v2.json :: cells.*.*.s1.n_jets_auc :: sha256 9cd059d3e2f51798 :: test jets the macro AUC is computed on
\newcommand{\FtNTestAccJc}{2{,}000{,}000}  % experiments/FIGS/data/w2b_leg2_metrics_v2.json :: cells.*.*.s1.n_jets :: sha256 9cd059d3e2f51798 :: test jets the accuracy is computed on
\newcommand{\FtSizeEThree}{\ensuremath{10^{3}}}  % experiments/FIGS/data/w2b_leg1_metrics_v2.json :: cells.*.N1000 :: sha256 b56bf70be9982044 :: fine-tuning training jets
\newcommand{\FtSizeEFour}{\ensuremath{10^{4}}}  % experiments/FIGS/data/w2b_leg1_metrics_v2.json :: cells.*.N10000 :: sha256 b56bf70be9982044 :: fine-tuning training jets
\newcommand{\FtSizeEFive}{\ensuremath{10^{5}}}  % experiments/FIGS/data/w2b_leg1_metrics_v2.json :: cells.*.N100000 :: sha256 b56bf70be9982044 :: fine-tuning training jets
\newcommand{\FtSizeESix}{\ensuremath{10^{6}}}  % experiments/FIGS/data/w2b_leg1_metrics_v2.json :: cells.*.N1000000 :: sha256 b56bf70be9982044 :: fine-tuning training jets
\newcommand{\FtNRunsWord}{five}  % experiments/FIGS/data/w2b_leg1_metrics_v2.json :: cells (pretrained models per vocabulary row) :: sha256 b56bf70be9982044 :: pretraining runs per vocabulary
\newcommand{\PairedProbeBcVsRestLinearOnesevenOverOnesixtwo}{1.64~[1.52, 1.78]}  % experiments/FIGS/data/paired_v1err/probes/ratios.json :: ratios[0].ratio, ratios[0].ci95 :: sha256 14dd51402305d147 :: 17 over 162, bc_vs_rest linear 1-auc: paired geometric mean [95% interval]
\newcommand{\PairedProbeBcVsRestMlpOnesevenOverOnesixtwo}{1.50~[1.39, 1.62]}  % experiments/FIGS/data/paired_v1err/probes/ratios.json :: ratios[3].ratio, ratios[3].ci95 :: sha256 14dd51402305d147 :: 17 over 162, bc_vs_rest mlp 1-auc: paired geometric mean [95% interval]
\newcommand{\PairedProbeBvcFourprongLinearOnesixtwoOverOneeighteight}{1.01~[0.95, 1.08]}  % experiments/FIGS/data/paired_v1err/probes/ratios.json :: ratios[6].ratio, ratios[6].ci95 :: sha256 14dd51402305d147 :: 162 over 188, bvc_4prong linear 1-auc: paired geometric mean [95% interval]
\newcommand{\PairedProbeBvcFourprongLinearFourthreeOverOneeighteight}{1.10~[1.01, 1.20]}  % experiments/FIGS/data/paired_v1err/probes/ratios.json :: ratios[7].ratio, ratios[7].ci95 :: sha256 14dd51402305d147 :: 43 over 188, bvc_4prong linear 1-auc: paired geometric mean [95% interval]
\newcommand{\PairedProbeBvcFourprongLinearOnesevenOverOneeighteight}{4.5~[3.5, 5.7]}  % experiments/FIGS/data/paired_v1err/probes/ratios.json :: ratios[8].ratio, ratios[8].ci95 :: sha256 14dd51402305d147 :: 17 over 188, bvc_4prong linear 1-auc: paired geometric mean [95% interval]
\newcommand{\PairedProbeBvcFourprongLinearFourthreeOverOnesixtwo}{1.09~[0.98, 1.20]}  % experiments/FIGS/data/paired_v1err/probes/ratios.json :: ratios[9].ratio, ratios[9].ci95 :: sha256 14dd51402305d147 :: 43 over 162, bvc_4prong linear 1-auc: paired geometric mean [95% interval]
\newcommand{\PairedProbeBvcFourprongLinearOnesevenOverOnesixtwo}{4.4~[3.4, 5.6]}  % experiments/FIGS/data/paired_v1err/probes/ratios.json :: ratios[10].ratio, ratios[10].ci95 :: sha256 14dd51402305d147 :: 17 over 162, bvc_4prong linear 1-auc: paired geometric mean [95% interval]
\newcommand{\PairedProbeBvcFourprongLinearOnesevenOverFourthree}{4.0~[3.1, 5.4]}  % experiments/FIGS/data/paired_v1err/probes/ratios.json :: ratios[11].ratio, ratios[11].ci95 :: sha256 14dd51402305d147 :: 17 over 43, bvc_4prong linear 1-auc: paired geometric mean [95% interval]
\newcommand{\PairedProbeBvcFourprongLinearOnesixtwoMassOverOnesixtwo}{1.05~[0.97, 1.13]}  % experiments/FIGS/data/paired_v1err/probes/ratios.json :: ratios[12].ratio, ratios[12].ci95 :: sha256 14dd51402305d147 :: 162+mass over 162, bvc_4prong linear 1-auc: paired geometric mean [95% interval]
\newcommand{\PairedProbeBvcFourprongLinearOnesevenMassOverOneseven}{2.1~[1.6, 2.7]}  % experiments/FIGS/data/paired_v1err/probes/ratios.json :: ratios[13].ratio, ratios[13].ci95 :: sha256 14dd51402305d147 :: 17+mass over 17, bvc_4prong linear 1-auc: paired geometric mean [95% interval]
\newcommand{\PairedProbeBvcFourprongLinearRandomOverOneeighteight}{1.54~[1.37, 1.74]}  % experiments/FIGS/data/paired_v1err/probes/ratios.json :: ratios[14].ratio, ratios[14].ci95 :: sha256 14dd51402305d147 :: random over 188, bvc_4prong linear 1-auc: paired geometric mean [95% interval]
\newcommand{\PairedProbeBvcFourprongLinearRandomOverOnesixtwo}{1.55~[1.36, 1.77]}  % experiments/FIGS/data/paired_v1err/probes/ratios.json :: ratios[15].ratio, ratios[15].ci95 :: sha256 14dd51402305d147 :: random over 162, bvc_4prong linear 1-auc: paired geometric mean [95% interval]
\newcommand{\PairedProbeBvcFourprongLinearRandomOverFourthree}{1.39~[1.27, 1.52]}  % experiments/FIGS/data/paired_v1err/probes/ratios.json :: ratios[16].ratio, ratios[16].ci95 :: sha256 14dd51402305d147 :: random over 43, bvc_4prong linear 1-auc: paired geometric mean [95% interval]
\newcommand{\PairedProbeBvcFourprongLinearRandomOverOneseven}{0.34~[0.20, 0.58]}  % experiments/FIGS/data/paired_v1err/probes/ratios.json :: ratios[17].ratio, ratios[17].ci95 :: sha256 14dd51402305d147 :: random over 17, bvc_4prong linear 1-auc: paired geometric mean [95% interval]
\newcommand{\PairedProbeBvcFourprongLinearRandomDrawTwoOverRandomDrawOne}{1.05~[0.94, 1.18]}  % experiments/FIGS/data/paired_v1err/probes/ratios.json :: ratios[18].ratio, ratios[18].ci95 :: sha256 14dd51402305d147 :: random draw 2 over random draw 1, bvc_4prong linear 1-auc: paired geometric mean [95% interval]
\newcommand{\PairedProbeBvcFourprongLinearRandomDrawThreeOverRandomDrawOne}{0.95~[0.85, 1.06]}  % experiments/FIGS/data/paired_v1err/probes/ratios.json :: ratios[19].ratio, ratios[19].ci95 :: sha256 14dd51402305d147 :: random draw 3 over random draw 1, bvc_4prong linear 1-auc: paired geometric mean [95% interval]
\newcommand{\PairedProbeBvcFourprongLinearRandomDrawThreeOverRandomDrawTwo}{0.90~[0.79, 1.02]}  % experiments/FIGS/data/paired_v1err/probes/ratios.json :: ratios[20].ratio, ratios[20].ci95 :: sha256 14dd51402305d147 :: random draw 3 over random draw 2, bvc_4prong linear 1-auc: paired geometric mean [95% interval]
\newcommand{\PairedProbeBvcFourprongMlpOnesixtwoOverOneeighteight}{1.00~[0.93, 1.07]}  % experiments/FIGS/data/paired_v1err/probes/ratios.json :: ratios[66].ratio, ratios[66].ci95 :: sha256 14dd51402305d147 :: 162 over 188, bvc_4prong mlp 1-auc: paired geometric mean [95% interval]
\newcommand{\PairedProbeBvcFourprongMlpFourthreeOverOneeighteight}{1.12~[1.03, 1.22]}  % experiments/FIGS/data/paired_v1err/probes/ratios.json :: ratios[67].ratio, ratios[67].ci95 :: sha256 14dd51402305d147 :: 43 over 188, bvc_4prong mlp 1-auc: paired geometric mean [95% interval]
\newcommand{\PairedProbeBvcFourprongMlpOnesevenOverOneeighteight}{3.2~[2.6, 4.1]}  % experiments/FIGS/data/paired_v1err/probes/ratios.json :: ratios[68].ratio, ratios[68].ci95 :: sha256 14dd51402305d147 :: 17 over 188, bvc_4prong mlp 1-auc: paired geometric mean [95% interval]
\newcommand{\PairedProbeBvcFourprongMlpFourthreeOverOnesixtwo}{1.12~[1.01, 1.26]}  % experiments/FIGS/data/paired_v1err/probes/ratios.json :: ratios[69].ratio, ratios[69].ci95 :: sha256 14dd51402305d147 :: 43 over 162, bvc_4prong mlp 1-auc: paired geometric mean [95% interval]
\newcommand{\PairedProbeBvcFourprongMlpOnesevenOverOnesixtwo}{3.3~[2.6, 4.0]}  % experiments/FIGS/data/paired_v1err/probes/ratios.json :: ratios[70].ratio, ratios[70].ci95 :: sha256 14dd51402305d147 :: 17 over 162, bvc_4prong mlp 1-auc: paired geometric mean [95% interval]
\newcommand{\PairedProbeBvcFourprongMlpOnesevenOverFourthree}{2.9~[2.3, 3.7]}  % experiments/FIGS/data/paired_v1err/probes/ratios.json :: ratios[71].ratio, ratios[71].ci95 :: sha256 14dd51402305d147 :: 17 over 43, bvc_4prong mlp 1-auc: paired geometric mean [95% interval]
\newcommand{\PairedProbeBvcFourprongMlpOnesixtwoMassOverOnesixtwo}{1.07~[0.99, 1.16]}  % experiments/FIGS/data/paired_v1err/probes/ratios.json :: ratios[72].ratio, ratios[72].ci95 :: sha256 14dd51402305d147 :: 162+mass over 162, bvc_4prong mlp 1-auc: paired geometric mean [95% interval]
\newcommand{\PairedProbeBvcFourprongMlpOnesevenMassOverOneseven}{1.97~[1.46, 2.67]}  % experiments/FIGS/data/paired_v1err/probes/ratios.json :: ratios[73].ratio, ratios[73].ci95 :: sha256 14dd51402305d147 :: 17+mass over 17, bvc_4prong mlp 1-auc: paired geometric mean [95% interval]
\newcommand{\PairedProbeBvcFourprongMlpRandomOverOneeighteight}{1.47~[1.30, 1.65]}  % experiments/FIGS/data/paired_v1err/probes/ratios.json :: ratios[74].ratio, ratios[74].ci95 :: sha256 14dd51402305d147 :: random over 188, bvc_4prong mlp 1-auc: paired geometric mean [95% interval]
\newcommand{\PairedProbeBvcFourprongMlpRandomOverOnesixtwo}{1.49~[1.24, 1.79]}  % experiments/FIGS/data/paired_v1err/probes/ratios.json :: ratios[75].ratio, ratios[75].ci95 :: sha256 14dd51402305d147 :: random over 162, bvc_4prong mlp 1-auc: paired geometric mean [95% interval]
\newcommand{\PairedProbeBvcFourprongMlpRandomOverFourthree}{1.28~[1.16, 1.43]}  % experiments/FIGS/data/paired_v1err/probes/ratios.json :: ratios[76].ratio, ratios[76].ci95 :: sha256 14dd51402305d147 :: random over 43, bvc_4prong mlp 1-auc: paired geometric mean [95% interval]
\newcommand{\PairedProbeBvcFourprongMlpRandomOverOneseven}{0.43~[0.28, 0.67]}  % experiments/FIGS/data/paired_v1err/probes/ratios.json :: ratios[77].ratio, ratios[77].ci95 :: sha256 14dd51402305d147 :: random over 17, bvc_4prong mlp 1-auc: paired geometric mean [95% interval]
\newcommand{\PairedProbeBvcFourprongMlpRandomDrawTwoOverRandomDrawOne}{1.00~[0.90, 1.10]}  % experiments/FIGS/data/paired_v1err/probes/ratios.json :: ratios[78].ratio, ratios[78].ci95 :: sha256 14dd51402305d147 :: random draw 2 over random draw 1, bvc_4prong mlp 1-auc: paired geometric mean [95% interval]
\newcommand{\PairedProbeBvcFourprongMlpRandomDrawThreeOverRandomDrawOne}{0.89~[0.79, 0.99]}  % experiments/FIGS/data/paired_v1err/probes/ratios.json :: ratios[79].ratio, ratios[79].ci95 :: sha256 14dd51402305d147 :: random draw 3 over random draw 1, bvc_4prong mlp 1-auc: paired geometric mean [95% interval]
\newcommand{\PairedProbeBvcFourprongMlpRandomDrawThreeOverRandomDrawTwo}{0.89~[0.79, 0.99]}  % experiments/FIGS/data/paired_v1err/probes/ratios.json :: ratios[80].ratio, ratios[80].ci95 :: sha256 14dd51402305d147 :: random draw 3 over random draw 2, bvc_4prong mlp 1-auc: paired geometric mean [95% interval]
\newcommand{\PairedProbeBvcQcdLinearOnesixtwoOverOneeighteight}{1.15~[1.07, 1.23]}  % experiments/FIGS/data/paired_v1err/probes/ratios.json :: ratios[126].ratio, ratios[126].ci95 :: sha256 14dd51402305d147 :: 162 over 188, bvc_qcd linear 1-auc: paired geometric mean [95% interval]
\newcommand{\PairedProbeBvcQcdLinearFourthreeOverOneeighteight}{1.16~[1.08, 1.25]}  % experiments/FIGS/data/paired_v1err/probes/ratios.json :: ratios[127].ratio, ratios[127].ci95 :: sha256 14dd51402305d147 :: 43 over 188, bvc_qcd linear 1-auc: paired geometric mean [95% interval]
\newcommand{\PairedProbeBvcQcdLinearOnesevenOverOneeighteight}{3.0~[2.5, 3.7]}  % experiments/FIGS/data/paired_v1err/probes/ratios.json :: ratios[128].ratio, ratios[128].ci95 :: sha256 14dd51402305d147 :: 17 over 188, bvc_qcd linear 1-auc: paired geometric mean [95% interval]
\newcommand{\PairedProbeBvcQcdLinearFourthreeOverOnesixtwo}{1.01~[0.96, 1.07]}  % experiments/FIGS/data/paired_v1err/probes/ratios.json :: ratios[129].ratio, ratios[129].ci95 :: sha256 14dd51402305d147 :: 43 over 162, bvc_qcd linear 1-auc: paired geometric mean [95% interval]
\newcommand{\PairedProbeBvcQcdLinearOnesevenOverOnesixtwo}{2.6~[2.2, 3.1]}  % experiments/FIGS/data/paired_v1err/probes/ratios.json :: ratios[130].ratio, ratios[130].ci95 :: sha256 14dd51402305d147 :: 17 over 162, bvc_qcd linear 1-auc: paired geometric mean [95% interval]
\newcommand{\PairedProbeBvcQcdLinearOnesevenOverFourthree}{2.6~[2.2, 3.0]}  % experiments/FIGS/data/paired_v1err/probes/ratios.json :: ratios[131].ratio, ratios[131].ci95 :: sha256 14dd51402305d147 :: 17 over 43, bvc_qcd linear 1-auc: paired geometric mean [95% interval]
\newcommand{\PairedProbeBvcQcdLinearOnesixtwoMassOverOnesixtwo}{1.060~[1.002, 1.123]}  % experiments/FIGS/data/paired_v1err/probes/ratios.json :: ratios[132].ratio, ratios[132].ci95 :: sha256 14dd51402305d147 :: 162+mass over 162, bvc_qcd linear 1-auc: paired geometric mean [95% interval]
\newcommand{\PairedProbeBvcQcdLinearOnesevenMassOverOneseven}{1.89~[1.52, 2.34]}  % experiments/FIGS/data/paired_v1err/probes/ratios.json :: ratios[133].ratio, ratios[133].ci95 :: sha256 14dd51402305d147 :: 17+mass over 17, bvc_qcd linear 1-auc: paired geometric mean [95% interval]
\newcommand{\PairedProbeBvcQcdMlpOnesixtwoOverOneeighteight}{1.11~[1.04, 1.19]}  % experiments/FIGS/data/paired_v1err/probes/ratios.json :: ratios[158].ratio, ratios[158].ci95 :: sha256 14dd51402305d147 :: 162 over 188, bvc_qcd mlp 1-auc: paired geometric mean [95% interval]
\newcommand{\PairedProbeBvcQcdMlpFourthreeOverOneeighteight}{1.13~[1.05, 1.21]}  % experiments/FIGS/data/paired_v1err/probes/ratios.json :: ratios[159].ratio, ratios[159].ci95 :: sha256 14dd51402305d147 :: 43 over 188, bvc_qcd mlp 1-auc: paired geometric mean [95% interval]
\newcommand{\PairedProbeBvcQcdMlpOnesevenOverOneeighteight}{2.6~[2.2, 3.0]}  % experiments/FIGS/data/paired_v1err/probes/ratios.json :: ratios[160].ratio, ratios[160].ci95 :: sha256 14dd51402305d147 :: 17 over 188, bvc_qcd mlp 1-auc: paired geometric mean [95% interval]
\newcommand{\PairedProbeBvcQcdMlpFourthreeOverOnesixtwo}{1.01~[0.96, 1.06]}  % experiments/FIGS/data/paired_v1err/probes/ratios.json :: ratios[161].ratio, ratios[161].ci95 :: sha256 14dd51402305d147 :: 43 over 162, bvc_qcd mlp 1-auc: paired geometric mean [95% interval]
\newcommand{\PairedProbeBvcQcdMlpOnesevenOverOnesixtwo}{2.3~[2.0, 2.6]}  % experiments/FIGS/data/paired_v1err/probes/ratios.json :: ratios[162].ratio, ratios[162].ci95 :: sha256 14dd51402305d147 :: 17 over 162, bvc_qcd mlp 1-auc: paired geometric mean [95% interval]
\newcommand{\PairedProbeBvcQcdMlpOnesevenOverFourthree}{2.3~[2.0, 2.6]}  % experiments/FIGS/data/paired_v1err/probes/ratios.json :: ratios[163].ratio, ratios[163].ci95 :: sha256 14dd51402305d147 :: 17 over 43, bvc_qcd mlp 1-auc: paired geometric mean [95% interval]
\newcommand{\PairedProbeBvcQcdMlpOnesixtwoMassOverOnesixtwo}{1.06~[1.01, 1.12]}  % experiments/FIGS/data/paired_v1err/probes/ratios.json :: ratios[164].ratio, ratios[164].ci95 :: sha256 14dd51402305d147 :: 162+mass over 162, bvc_qcd mlp 1-auc: paired geometric mean [95% interval]
\newcommand{\PairedProbeBvcQcdMlpOnesevenMassOverOneseven}{1.64~[1.33, 2.02]}  % experiments/FIGS/data/paired_v1err/probes/ratios.json :: ratios[165].ratio, ratios[165].ci95 :: sha256 14dd51402305d147 :: 17+mass over 17, bvc_qcd mlp 1-auc: paired geometric mean [95% interval]
\newcommand{\PairedProbeBvcResonantLinearOnesixtwoOverOneeighteight}{1.02~[0.96, 1.08]}  % experiments/FIGS/data/paired_v1err/probes/ratios.json :: ratios[190].ratio, ratios[190].ci95 :: sha256 14dd51402305d147 :: 162 over 188, bvc_resonant linear 1-auc: paired geometric mean [95% interval]
\newcommand{\PairedProbeBvcResonantLinearFourthreeOverOneeighteight}{1.09~[1.03, 1.15]}  % experiments/FIGS/data/paired_v1err/probes/ratios.json :: ratios[191].ratio, ratios[191].ci95 :: sha256 14dd51402305d147 :: 43 over 188, bvc_resonant linear 1-auc: paired geometric mean [95% interval]
\newcommand{\PairedProbeBvcResonantLinearOnesevenOverOneeighteight}{4.0~[3.3, 4.9]}  % experiments/FIGS/data/paired_v1err/probes/ratios.json :: ratios[192].ratio, ratios[192].ci95 :: sha256 14dd51402305d147 :: 17 over 188, bvc_resonant linear 1-auc: paired geometric mean [95% interval]
\newcommand{\PairedProbeBvcResonantLinearFourthreeOverOnesixtwo}{1.07~[1.02, 1.12]}  % experiments/FIGS/data/paired_v1err/probes/ratios.json :: ratios[193].ratio, ratios[193].ci95 :: sha256 14dd51402305d147 :: 43 over 162, bvc_resonant linear 1-auc: paired geometric mean [95% interval]
\newcommand{\PairedProbeBvcResonantLinearOnesevenOverOnesixtwo}{3.9~[3.1, 4.9]}  % experiments/FIGS/data/paired_v1err/probes/ratios.json :: ratios[194].ratio, ratios[194].ci95 :: sha256 14dd51402305d147 :: 17 over 162, bvc_resonant linear 1-auc: paired geometric mean [95% interval]
\newcommand{\PairedProbeBvcResonantLinearOnesevenOverFourthree}{3.7~[3.0, 4.5]}  % experiments/FIGS/data/paired_v1err/probes/ratios.json :: ratios[195].ratio, ratios[195].ci95 :: sha256 14dd51402305d147 :: 17 over 43, bvc_resonant linear 1-auc: paired geometric mean [95% interval]
\newcommand{\PairedProbeBvcResonantLinearOnesixtwoMassOverOnesixtwo}{1.11~[1.06, 1.16]}  % experiments/FIGS/data/paired_v1err/probes/ratios.json :: ratios[196].ratio, ratios[196].ci95 :: sha256 14dd51402305d147 :: 162+mass over 162, bvc_resonant linear 1-auc: paired geometric mean [95% interval]
\newcommand{\PairedProbeBvcResonantLinearOnesevenMassOverOneseven}{3.0~[1.9, 4.7]}  % experiments/FIGS/data/paired_v1err/probes/ratios.json :: ratios[197].ratio, ratios[197].ci95 :: sha256 14dd51402305d147 :: 17+mass over 17, bvc_resonant linear 1-auc: paired geometric mean [95% interval]
\newcommand{\PairedProbeBvcResonantMlpOnesixtwoOverOneeighteight}{1.04~[0.98, 1.12]}  % experiments/FIGS/data/paired_v1err/probes/ratios.json :: ratios[222].ratio, ratios[222].ci95 :: sha256 14dd51402305d147 :: 162 over 188, bvc_resonant mlp 1-auc: paired geometric mean [95% interval]
\newcommand{\PairedProbeBvcResonantMlpFourthreeOverOneeighteight}{1.14~[1.08, 1.21]}  % experiments/FIGS/data/paired_v1err/probes/ratios.json :: ratios[223].ratio, ratios[223].ci95 :: sha256 14dd51402305d147 :: 43 over 188, bvc_resonant mlp 1-auc: paired geometric mean [95% interval]
\newcommand{\PairedProbeBvcResonantMlpOnesevenOverOneeighteight}{3.3~[2.8, 3.8]}  % experiments/FIGS/data/paired_v1err/probes/ratios.json :: ratios[224].ratio, ratios[224].ci95 :: sha256 14dd51402305d147 :: 17 over 188, bvc_resonant mlp 1-auc: paired geometric mean [95% interval]
\newcommand{\PairedProbeBvcResonantMlpFourthreeOverOnesixtwo}{1.09~[1.04, 1.14]}  % experiments/FIGS/data/paired_v1err/probes/ratios.json :: ratios[225].ratio, ratios[225].ci95 :: sha256 14dd51402305d147 :: 43 over 162, bvc_resonant mlp 1-auc: paired geometric mean [95% interval]
\newcommand{\PairedProbeBvcResonantMlpOnesevenOverOnesixtwo}{3.1~[2.6, 3.8]}  % experiments/FIGS/data/paired_v1err/probes/ratios.json :: ratios[226].ratio, ratios[226].ci95 :: sha256 14dd51402305d147 :: 17 over 162, bvc_resonant mlp 1-auc: paired geometric mean [95% interval]
\newcommand{\PairedProbeBvcResonantMlpOnesevenOverFourthree}{2.9~[2.4, 3.4]}  % experiments/FIGS/data/paired_v1err/probes/ratios.json :: ratios[227].ratio, ratios[227].ci95 :: sha256 14dd51402305d147 :: 17 over 43, bvc_resonant mlp 1-auc: paired geometric mean [95% interval]
\newcommand{\PairedProbeBvcResonantMlpOnesixtwoMassOverOnesixtwo}{1.09~[1.04, 1.14]}  % experiments/FIGS/data/paired_v1err/probes/ratios.json :: ratios[228].ratio, ratios[228].ci95 :: sha256 14dd51402305d147 :: 162+mass over 162, bvc_resonant mlp 1-auc: paired geometric mean [95% interval]
\newcommand{\PairedProbeBvcResonantMlpOnesevenMassOverOneseven}{1.95~[1.45, 2.61]}  % experiments/FIGS/data/paired_v1err/probes/ratios.json :: ratios[229].ratio, ratios[229].ci95 :: sha256 14dd51402305d147 :: 17+mass over 17, bvc_resonant mlp 1-auc: paired geometric mean [95% interval]
\newcommand{\PairedProbeEeVsMmLinearOnesevenOverFourthree}{2.3~[0.8, 6.5]}  % experiments/FIGS/data/paired_v1err/probes/ratios.json :: ratios[259].ratio, ratios[259].ci95 :: sha256 14dd51402305d147 :: 17 over 43, ee_vs_mm linear 1-auc: paired geometric mean [95% interval]
\newcommand{\PairedProbeEeVsMmLinearOnesevenMassOverOneseven}{0.23~[0.08, 0.64]}  % experiments/FIGS/data/paired_v1err/probes/ratios.json :: ratios[261].ratio, ratios[261].ci95 :: sha256 14dd51402305d147 :: 17+mass over 17, ee_vs_mm linear 1-auc: paired geometric mean [95% interval]
\newcommand{\PairedProbeEeVsMmMlpOnesevenOverFourthree}{2.8~[0.8, 9.0]}  % experiments/FIGS/data/paired_v1err/probes/ratios.json :: ratios[291].ratio, ratios[291].ci95 :: sha256 14dd51402305d147 :: 17 over 43, ee_vs_mm mlp 1-auc: paired geometric mean [95% interval]
\newcommand{\PairedProbeEeVsMmMlpOnesevenMassOverOneseven}{0.172~[0.026, 1.125]}  % experiments/FIGS/data/paired_v1err/probes/ratios.json :: ratios[293].ratio, ratios[293].ci95 :: sha256 14dd51402305d147 :: 17+mass over 17, ee_vs_mm mlp 1-auc: paired geometric mean [95% interval]
\newcommand{\PairedProbeRetainedTopologyLinearOnesixtwoOverOneeighteight}{1.00~[0.92, 1.08]}  % experiments/FIGS/data/paired_v1err/probes/ratios.json :: ratios[318].ratio, ratios[318].ci95 :: sha256 14dd51402305d147 :: 162 over 188, retained_topology linear 1-auc: paired geometric mean [95% interval]
\newcommand{\PairedProbeRetainedTopologyLinearFourthreeOverOneeighteight}{1.24~[1.15, 1.34]}  % experiments/FIGS/data/paired_v1err/probes/ratios.json :: ratios[319].ratio, ratios[319].ci95 :: sha256 14dd51402305d147 :: 43 over 188, retained_topology linear 1-auc: paired geometric mean [95% interval]
\newcommand{\PairedProbeRetainedTopologyLinearOnesevenOverOneeighteight}{1.54~[1.35, 1.75]}  % experiments/FIGS/data/paired_v1err/probes/ratios.json :: ratios[320].ratio, ratios[320].ci95 :: sha256 14dd51402305d147 :: 17 over 188, retained_topology linear 1-auc: paired geometric mean [95% interval]
\newcommand{\PairedProbeRetainedTopologyLinearFourthreeOverOnesixtwo}{1.24~[1.14, 1.34]}  % experiments/FIGS/data/paired_v1err/probes/ratios.json :: ratios[321].ratio, ratios[321].ci95 :: sha256 14dd51402305d147 :: 43 over 162, retained_topology linear 1-auc: paired geometric mean [95% interval]
\newcommand{\PairedProbeRetainedTopologyLinearOnesevenOverOnesixtwo}{1.54~[1.33, 1.78]}  % experiments/FIGS/data/paired_v1err/probes/ratios.json :: ratios[322].ratio, ratios[322].ci95 :: sha256 14dd51402305d147 :: 17 over 162, retained_topology linear 1-auc: paired geometric mean [95% interval]
\newcommand{\PairedProbeRetainedTopologyLinearOnesevenOverFourthree}{1.24~[1.08, 1.42]}  % experiments/FIGS/data/paired_v1err/probes/ratios.json :: ratios[323].ratio, ratios[323].ci95 :: sha256 14dd51402305d147 :: 17 over 43, retained_topology linear 1-auc: paired geometric mean [95% interval]
\newcommand{\PairedProbeRetainedTopologyLinearOnesixtwoMassOverOnesixtwo}{1.06~[0.99, 1.14]}  % experiments/FIGS/data/paired_v1err/probes/ratios.json :: ratios[324].ratio, ratios[324].ci95 :: sha256 14dd51402305d147 :: 162+mass over 162, retained_topology linear 1-auc: paired geometric mean [95% interval]
\newcommand{\PairedProbeRetainedTopologyLinearOnesevenMassOverOneseven}{1.12~[0.97, 1.28]}  % experiments/FIGS/data/paired_v1err/probes/ratios.json :: ratios[325].ratio, ratios[325].ci95 :: sha256 14dd51402305d147 :: 17+mass over 17, retained_topology linear 1-auc: paired geometric mean [95% interval]
\newcommand{\PairedProbeRetainedTopologyMlpOnesixtwoOverOneeighteight}{1.01~[0.87, 1.17]}  % experiments/FIGS/data/paired_v1err/probes/ratios.json :: ratios[350].ratio, ratios[350].ci95 :: sha256 14dd51402305d147 :: 162 over 188, retained_topology mlp 1-auc: paired geometric mean [95% interval]
\newcommand{\PairedProbeRetainedTopologyMlpFourthreeOverOneeighteight}{1.18~[1.03, 1.36]}  % experiments/FIGS/data/paired_v1err/probes/ratios.json :: ratios[351].ratio, ratios[351].ci95 :: sha256 14dd51402305d147 :: 43 over 188, retained_topology mlp 1-auc: paired geometric mean [95% interval]
\newcommand{\PairedProbeRetainedTopologyMlpOnesevenOverOneeighteight}{1.52~[1.34, 1.72]}  % experiments/FIGS/data/paired_v1err/probes/ratios.json :: ratios[352].ratio, ratios[352].ci95 :: sha256 14dd51402305d147 :: 17 over 188, retained_topology mlp 1-auc: paired geometric mean [95% interval]
\newcommand{\PairedProbeRetainedTopologyMlpFourthreeOverOnesixtwo}{1.17~[1.07, 1.28]}  % experiments/FIGS/data/paired_v1err/probes/ratios.json :: ratios[353].ratio, ratios[353].ci95 :: sha256 14dd51402305d147 :: 43 over 162, retained_topology mlp 1-auc: paired geometric mean [95% interval]
\newcommand{\PairedProbeRetainedTopologyMlpOnesevenOverOnesixtwo}{1.50~[1.29, 1.75]}  % experiments/FIGS/data/paired_v1err/probes/ratios.json :: ratios[354].ratio, ratios[354].ci95 :: sha256 14dd51402305d147 :: 17 over 162, retained_topology mlp 1-auc: paired geometric mean [95% interval]
\newcommand{\PairedProbeRetainedTopologyMlpOnesevenOverFourthree}{1.28~[1.12, 1.47]}  % experiments/FIGS/data/paired_v1err/probes/ratios.json :: ratios[355].ratio, ratios[355].ci95 :: sha256 14dd51402305d147 :: 17 over 43, retained_topology mlp 1-auc: paired geometric mean [95% interval]
\newcommand{\PairedProbeRetainedTopologyMlpOnesixtwoMassOverOnesixtwo}{1.04~[0.92, 1.17]}  % experiments/FIGS/data/paired_v1err/probes/ratios.json :: ratios[356].ratio, ratios[356].ci95 :: sha256 14dd51402305d147 :: 162+mass over 162, retained_topology mlp 1-auc: paired geometric mean [95% interval]
\newcommand{\PairedProbeRetainedTopologyMlpOnesevenMassOverOneseven}{1.11~[0.99, 1.23]}  % experiments/FIGS/data/paired_v1err/probes/ratios.json :: ratios[357].ratio, ratios[357].ci95 :: sha256 14dd51402305d147 :: 17+mass over 17, retained_topology mlp 1-auc: paired geometric mean [95% interval]
\newcommand{\PairedProbeVisibleContentLinearOnesixtwoOverOneeighteight}{0.95~[0.88, 1.03]}  % experiments/FIGS/data/paired_v1err/probes/ratios.json :: ratios[382].ratio, ratios[382].ci95 :: sha256 14dd51402305d147 :: 162 over 188, visible_content linear 1-auc: paired geometric mean [95% interval]
\newcommand{\PairedProbeVisibleContentLinearFourthreeOverOneeighteight}{0.99~[0.85, 1.15]}  % experiments/FIGS/data/paired_v1err/probes/ratios.json :: ratios[383].ratio, ratios[383].ci95 :: sha256 14dd51402305d147 :: 43 over 188, visible_content linear 1-auc: paired geometric mean [95% interval]
\newcommand{\PairedProbeVisibleContentLinearOnesevenOverOneeighteight}{2.8~[2.1, 3.7]}  % experiments/FIGS/data/paired_v1err/probes/ratios.json :: ratios[384].ratio, ratios[384].ci95 :: sha256 14dd51402305d147 :: 17 over 188, visible_content linear 1-auc: paired geometric mean [95% interval]
\newcommand{\PairedProbeVisibleContentLinearFourthreeOverOnesixtwo}{1.04~[0.89, 1.20]}  % experiments/FIGS/data/paired_v1err/probes/ratios.json :: ratios[385].ratio, ratios[385].ci95 :: sha256 14dd51402305d147 :: 43 over 162, visible_content linear 1-auc: paired geometric mean [95% interval]
\newcommand{\PairedProbeVisibleContentLinearOnesevenOverOnesixtwo}{2.9~[2.3, 3.8]}  % experiments/FIGS/data/paired_v1err/probes/ratios.json :: ratios[386].ratio, ratios[386].ci95 :: sha256 14dd51402305d147 :: 17 over 162, visible_content linear 1-auc: paired geometric mean [95% interval]
\newcommand{\PairedProbeVisibleContentLinearOnesevenOverFourthree}{2.8~[2.2, 3.7]}  % experiments/FIGS/data/paired_v1err/probes/ratios.json :: ratios[387].ratio, ratios[387].ci95 :: sha256 14dd51402305d147 :: 17 over 43, visible_content linear 1-auc: paired geometric mean [95% interval]
\newcommand{\PairedProbeVisibleContentLinearOnesixtwoMassOverOnesixtwo}{1.10~[0.98, 1.23]}  % experiments/FIGS/data/paired_v1err/probes/ratios.json :: ratios[388].ratio, ratios[388].ci95 :: sha256 14dd51402305d147 :: 162+mass over 162, visible_content linear 1-auc: paired geometric mean [95% interval]
\newcommand{\PairedProbeVisibleContentLinearOnesevenMassOverOneseven}{1.44~[1.26, 1.65]}  % experiments/FIGS/data/paired_v1err/probes/ratios.json :: ratios[389].ratio, ratios[389].ci95 :: sha256 14dd51402305d147 :: 17+mass over 17, visible_content linear 1-auc: paired geometric mean [95% interval]
\newcommand{\PairedProbeVisibleContentLinearRandomOverOneeighteight}{1.31~[1.01, 1.70]}  % experiments/FIGS/data/paired_v1err/probes/ratios.json :: ratios[390].ratio, ratios[390].ci95 :: sha256 14dd51402305d147 :: random over 188, visible_content linear 1-auc: paired geometric mean [95% interval]
\newcommand{\PairedProbeVisibleContentLinearRandomOverOnesixtwo}{1.41~[1.17, 1.70]}  % experiments/FIGS/data/paired_v1err/probes/ratios.json :: ratios[391].ratio, ratios[391].ci95 :: sha256 14dd51402305d147 :: random over 162, visible_content linear 1-auc: paired geometric mean [95% interval]
\newcommand{\PairedProbeVisibleContentLinearRandomOverFourthree}{1.33~[1.12, 1.57]}  % experiments/FIGS/data/paired_v1err/probes/ratios.json :: ratios[392].ratio, ratios[392].ci95 :: sha256 14dd51402305d147 :: random over 43, visible_content linear 1-auc: paired geometric mean [95% interval]
\newcommand{\PairedProbeVisibleContentLinearRandomOverOneseven}{0.48~[0.38, 0.60]}  % experiments/FIGS/data/paired_v1err/probes/ratios.json :: ratios[393].ratio, ratios[393].ci95 :: sha256 14dd51402305d147 :: random over 17, visible_content linear 1-auc: paired geometric mean [95% interval]
\newcommand{\PairedProbeVisibleContentLinearRandomDrawTwoOverRandomDrawOne}{0.79~[0.64, 0.98]}  % experiments/FIGS/data/paired_v1err/probes/ratios.json :: ratios[394].ratio, ratios[394].ci95 :: sha256 14dd51402305d147 :: random draw 2 over random draw 1, visible_content linear 1-auc: paired geometric mean [95% interval]
\newcommand{\PairedProbeVisibleContentLinearRandomDrawThreeOverRandomDrawOne}{0.86~[0.71, 1.05]}  % experiments/FIGS/data/paired_v1err/probes/ratios.json :: ratios[395].ratio, ratios[395].ci95 :: sha256 14dd51402305d147 :: random draw 3 over random draw 1, visible_content linear 1-auc: paired geometric mean [95% interval]
\newcommand{\PairedProbeVisibleContentLinearRandomDrawThreeOverRandomDrawTwo}{1.09~[0.88, 1.35]}  % experiments/FIGS/data/paired_v1err/probes/ratios.json :: ratios[396].ratio, ratios[396].ci95 :: sha256 14dd51402305d147 :: random draw 3 over random draw 2, visible_content linear 1-auc: paired geometric mean [95% interval]
\newcommand{\PairedProbeVisibleContentMlpOnesixtwoOverOneeighteight}{0.92~[0.85, 1.01]}  % experiments/FIGS/data/paired_v1err/probes/ratios.json :: ratios[442].ratio, ratios[442].ci95 :: sha256 14dd51402305d147 :: 162 over 188, visible_content mlp 1-auc: paired geometric mean [95% interval]
\newcommand{\PairedProbeVisibleContentMlpFourthreeOverOneeighteight}{0.93~[0.80, 1.07]}  % experiments/FIGS/data/paired_v1err/probes/ratios.json :: ratios[443].ratio, ratios[443].ci95 :: sha256 14dd51402305d147 :: 43 over 188, visible_content mlp 1-auc: paired geometric mean [95% interval]
\newcommand{\PairedProbeVisibleContentMlpOnesevenOverOneeighteight}{2.3~[1.7, 3.1]}  % experiments/FIGS/data/paired_v1err/probes/ratios.json :: ratios[444].ratio, ratios[444].ci95 :: sha256 14dd51402305d147 :: 17 over 188, visible_content mlp 1-auc: paired geometric mean [95% interval]
\newcommand{\PairedProbeVisibleContentMlpFourthreeOverOnesixtwo}{1.01~[0.89, 1.14]}  % experiments/FIGS/data/paired_v1err/probes/ratios.json :: ratios[445].ratio, ratios[445].ci95 :: sha256 14dd51402305d147 :: 43 over 162, visible_content mlp 1-auc: paired geometric mean [95% interval]
\newcommand{\PairedProbeVisibleContentMlpOnesevenOverOnesixtwo}{2.5~[1.9, 3.2]}  % experiments/FIGS/data/paired_v1err/probes/ratios.json :: ratios[446].ratio, ratios[446].ci95 :: sha256 14dd51402305d147 :: 17 over 162, visible_content mlp 1-auc: paired geometric mean [95% interval]
\newcommand{\PairedProbeVisibleContentMlpOnesevenOverFourthree}{2.5~[2.0, 3.1]}  % experiments/FIGS/data/paired_v1err/probes/ratios.json :: ratios[447].ratio, ratios[447].ci95 :: sha256 14dd51402305d147 :: 17 over 43, visible_content mlp 1-auc: paired geometric mean [95% interval]
\newcommand{\PairedProbeVisibleContentMlpOnesixtwoMassOverOnesixtwo}{1.08~[0.94, 1.24]}  % experiments/FIGS/data/paired_v1err/probes/ratios.json :: ratios[448].ratio, ratios[448].ci95 :: sha256 14dd51402305d147 :: 162+mass over 162, visible_content mlp 1-auc: paired geometric mean [95% interval]
\newcommand{\PairedProbeVisibleContentMlpOnesevenMassOverOneseven}{1.41~[1.19, 1.67]}  % experiments/FIGS/data/paired_v1err/probes/ratios.json :: ratios[449].ratio, ratios[449].ci95 :: sha256 14dd51402305d147 :: 17+mass over 17, visible_content mlp 1-auc: paired geometric mean [95% interval]
\newcommand{\PairedProbeVisibleContentMlpRandomOverOneeighteight}{1.27~[1.01, 1.59]}  % experiments/FIGS/data/paired_v1err/probes/ratios.json :: ratios[450].ratio, ratios[450].ci95 :: sha256 14dd51402305d147 :: random over 188, visible_content mlp 1-auc: paired geometric mean [95% interval]
\newcommand{\PairedProbeVisibleContentMlpRandomOverOnesixtwo}{1.37~[1.12, 1.68]}  % experiments/FIGS/data/paired_v1err/probes/ratios.json :: ratios[451].ratio, ratios[451].ci95 :: sha256 14dd51402305d147 :: random over 162, visible_content mlp 1-auc: paired geometric mean [95% interval]
\newcommand{\PairedProbeVisibleContentMlpRandomOverFourthree}{1.35~[1.13, 1.61]}  % experiments/FIGS/data/paired_v1err/probes/ratios.json :: ratios[452].ratio, ratios[452].ci95 :: sha256 14dd51402305d147 :: random over 43, visible_content mlp 1-auc: paired geometric mean [95% interval]
\newcommand{\PairedProbeVisibleContentMlpRandomOverOneseven}{0.56~[0.45, 0.70]}  % experiments/FIGS/data/paired_v1err/probes/ratios.json :: ratios[453].ratio, ratios[453].ci95 :: sha256 14dd51402305d147 :: random over 17, visible_content mlp 1-auc: paired geometric mean [95% interval]
\newcommand{\PairedProbeVisibleContentMlpRandomDrawTwoOverRandomDrawOne}{0.88~[0.71, 1.09]}  % experiments/FIGS/data/paired_v1err/probes/ratios.json :: ratios[454].ratio, ratios[454].ci95 :: sha256 14dd51402305d147 :: random draw 2 over random draw 1, visible_content mlp 1-auc: paired geometric mean [95% interval]
\newcommand{\PairedProbeVisibleContentMlpRandomDrawThreeOverRandomDrawOne}{0.94~[0.75, 1.18]}  % experiments/FIGS/data/paired_v1err/probes/ratios.json :: ratios[455].ratio, ratios[455].ci95 :: sha256 14dd51402305d147 :: random draw 3 over random draw 1, visible_content mlp 1-auc: paired geometric mean [95% interval]
\newcommand{\PairedProbeVisibleContentMlpRandomDrawThreeOverRandomDrawTwo}{1.06~[0.84, 1.35]}  % experiments/FIGS/data/paired_v1err/probes/ratios.json :: ratios[456].ratio, ratios[456].ci95 :: sha256 14dd51402305d147 :: random draw 3 over random draw 2, visible_content mlp 1-auc: paired geometric mean [95% interval]
\newcommand{\PairedFtJciiEThreeOnesixtwoOverOneeighteight}{0.95~[0.87, 1.05]}  % experiments/FIGS/data/paired_v1err/finetune/ratios.json :: ratios[0].ratio, ratios[0].ci95 :: sha256 3e5cf25e859b68c2 :: 162 over 188, leg1 N1000 1-macro_auc: paired geometric mean [95% interval]
\newcommand{\PairedFtJciiEThreeFourthreeOverOneeighteight}{0.84~[0.77, 0.91]}  % experiments/FIGS/data/paired_v1err/finetune/ratios.json :: ratios[1].ratio, ratios[1].ci95 :: sha256 3e5cf25e859b68c2 :: 43 over 188, leg1 N1000 1-macro_auc: paired geometric mean [95% interval]
\newcommand{\PairedFtJciiEThreeOnesevenOverOneeighteight}{0.97~[0.86, 1.09]}  % experiments/FIGS/data/paired_v1err/finetune/ratios.json :: ratios[2].ratio, ratios[2].ci95 :: sha256 3e5cf25e859b68c2 :: 17 over 188, leg1 N1000 1-macro_auc: paired geometric mean [95% interval]
\newcommand{\PairedFtJciiEThreeFourthreeOverOnesixtwo}{0.88~[0.80, 0.97]}  % experiments/FIGS/data/paired_v1err/finetune/ratios.json :: ratios[3].ratio, ratios[3].ci95 :: sha256 3e5cf25e859b68c2 :: 43 over 162, leg1 N1000 1-macro_auc: paired geometric mean [95% interval]
\newcommand{\PairedFtJciiEThreeOnesevenOverOnesixtwo}{1.02~[0.88, 1.18]}  % experiments/FIGS/data/paired_v1err/finetune/ratios.json :: ratios[4].ratio, ratios[4].ci95 :: sha256 3e5cf25e859b68c2 :: 17 over 162, leg1 N1000 1-macro_auc: paired geometric mean [95% interval]
\newcommand{\PairedFtJciiEThreeOnesevenOverFourthree}{1.16~[1.08, 1.24]}  % experiments/FIGS/data/paired_v1err/finetune/ratios.json :: ratios[5].ratio, ratios[5].ci95 :: sha256 3e5cf25e859b68c2 :: 17 over 43, leg1 N1000 1-macro_auc: paired geometric mean [95% interval]
\newcommand{\PairedFtJciiEThreeOnesixtwoMassOverOnesixtwo}{0.85~[0.77, 0.93]}  % experiments/FIGS/data/paired_v1err/finetune/ratios.json :: ratios[6].ratio, ratios[6].ci95 :: sha256 3e5cf25e859b68c2 :: 162+mass over 162, leg1 N1000 1-macro_auc: paired geometric mean [95% interval]
\newcommand{\PairedFtJciiEThreeOnesevenMassOverOneseven}{1.00~[0.86, 1.17]}  % experiments/FIGS/data/paired_v1err/finetune/ratios.json :: ratios[7].ratio, ratios[7].ci95 :: sha256 3e5cf25e859b68c2 :: 17+mass over 17, leg1 N1000 1-macro_auc: paired geometric mean [95% interval]
\newcommand{\PairedFtJciiEThreeRandomOverOneeighteight}{0.82~[0.73, 0.93]}  % experiments/FIGS/data/paired_v1err/finetune/ratios.json :: ratios[8].ratio, ratios[8].ci95 :: sha256 3e5cf25e859b68c2 :: random over 188, leg1 N1000 1-macro_auc: paired geometric mean [95% interval]
\newcommand{\PairedFtJciiEThreeRandomOverOnesixtwo}{0.90~[0.85, 0.96]}  % experiments/FIGS/data/paired_v1err/finetune/ratios.json :: ratios[9].ratio, ratios[9].ci95 :: sha256 3e5cf25e859b68c2 :: random over 162, leg1 N1000 1-macro_auc: paired geometric mean [95% interval]
\newcommand{\PairedFtJciiEThreeRandomOverFourthree}{0.99~[0.75, 1.30]}  % experiments/FIGS/data/paired_v1err/finetune/ratios.json :: ratios[10].ratio, ratios[10].ci95 :: sha256 3e5cf25e859b68c2 :: random over 43, leg1 N1000 1-macro_auc: paired geometric mean [95% interval]
\newcommand{\PairedFtJciiEThreeRandomOverOneseven}{0.86~[0.58, 1.27]}  % experiments/FIGS/data/paired_v1err/finetune/ratios.json :: ratios[11].ratio, ratios[11].ci95 :: sha256 3e5cf25e859b68c2 :: random over 17, leg1 N1000 1-macro_auc: paired geometric mean [95% interval]
\newcommand{\PairedFtJciiEThreeRandomDrawTwoOverRandomDrawOne}{1.09~[1.07, 1.10]}  % experiments/FIGS/data/paired_v1err/finetune/ratios.json :: ratios[12].ratio, ratios[12].ci95 :: sha256 3e5cf25e859b68c2 :: random draw 2 over random draw 1, leg1 N1000 1-macro_auc: paired geometric mean [95% interval]
\newcommand{\PairedFtJciiEThreeRandomDrawThreeOverRandomDrawOne}{1.04~[1.02, 1.05]}  % experiments/FIGS/data/paired_v1err/finetune/ratios.json :: ratios[13].ratio, ratios[13].ci95 :: sha256 3e5cf25e859b68c2 :: random draw 3 over random draw 1, leg1 N1000 1-macro_auc: paired geometric mean [95% interval]
\newcommand{\PairedFtJciiEThreeRandomDrawThreeOverRandomDrawTwo}{0.95~[0.94, 0.97]}  % experiments/FIGS/data/paired_v1err/finetune/ratios.json :: ratios[14].ratio, ratios[14].ci95 :: sha256 3e5cf25e859b68c2 :: random draw 3 over random draw 2, leg1 N1000 1-macro_auc: paired geometric mean [95% interval]
\newcommand{\PairedFtJciiEFourOnesixtwoOverOneeighteight}{0.99~[0.97, 1.02]}  % experiments/FIGS/data/paired_v1err/finetune/ratios.json :: ratios[15].ratio, ratios[15].ci95 :: sha256 3e5cf25e859b68c2 :: 162 over 188, leg1 N10000 1-macro_auc: paired geometric mean [95% interval]
\newcommand{\PairedFtJciiEFourFourthreeOverOneeighteight}{1.20~[1.17, 1.23]}  % experiments/FIGS/data/paired_v1err/finetune/ratios.json :: ratios[16].ratio, ratios[16].ci95 :: sha256 3e5cf25e859b68c2 :: 43 over 188, leg1 N10000 1-macro_auc: paired geometric mean [95% interval]
\newcommand{\PairedFtJciiEFourOnesevenOverOneeighteight}{1.32~[1.29, 1.34]}  % experiments/FIGS/data/paired_v1err/finetune/ratios.json :: ratios[17].ratio, ratios[17].ci95 :: sha256 3e5cf25e859b68c2 :: 17 over 188, leg1 N10000 1-macro_auc: paired geometric mean [95% interval]
\newcommand{\PairedFtJciiEFourFourthreeOverOnesixtwo}{1.21~[1.19, 1.23]}  % experiments/FIGS/data/paired_v1err/finetune/ratios.json :: ratios[18].ratio, ratios[18].ci95 :: sha256 3e5cf25e859b68c2 :: 43 over 162, leg1 N10000 1-macro_auc: paired geometric mean [95% interval]
\newcommand{\PairedFtJciiEFourOnesevenOverOnesixtwo}{1.32~[1.29, 1.36]}  % experiments/FIGS/data/paired_v1err/finetune/ratios.json :: ratios[19].ratio, ratios[19].ci95 :: sha256 3e5cf25e859b68c2 :: 17 over 162, leg1 N10000 1-macro_auc: paired geometric mean [95% interval]
\newcommand{\PairedFtJciiEFourOnesevenOverFourthree}{1.10~[1.07, 1.13]}  % experiments/FIGS/data/paired_v1err/finetune/ratios.json :: ratios[20].ratio, ratios[20].ci95 :: sha256 3e5cf25e859b68c2 :: 17 over 43, leg1 N10000 1-macro_auc: paired geometric mean [95% interval]
\newcommand{\PairedFtJciiEFourOnesixtwoMassOverOnesixtwo}{1.00~[0.94, 1.06]}  % experiments/FIGS/data/paired_v1err/finetune/ratios.json :: ratios[21].ratio, ratios[21].ci95 :: sha256 3e5cf25e859b68c2 :: 162+mass over 162, leg1 N10000 1-macro_auc: paired geometric mean [95% interval]
\newcommand{\PairedFtJciiEFourOnesevenMassOverOneseven}{0.990~[0.976, 1.004]}  % experiments/FIGS/data/paired_v1err/finetune/ratios.json :: ratios[22].ratio, ratios[22].ci95 :: sha256 3e5cf25e859b68c2 :: 17+mass over 17, leg1 N10000 1-macro_auc: paired geometric mean [95% interval]
\newcommand{\PairedFtJciiEFourRandomOverOneeighteight}{1.15~[1.09, 1.22]}  % experiments/FIGS/data/paired_v1err/finetune/ratios.json :: ratios[23].ratio, ratios[23].ci95 :: sha256 3e5cf25e859b68c2 :: random over 188, leg1 N10000 1-macro_auc: paired geometric mean [95% interval]
\newcommand{\PairedFtJciiEFourRandomOverOnesixtwo}{1.18~[1.13, 1.22]}  % experiments/FIGS/data/paired_v1err/finetune/ratios.json :: ratios[24].ratio, ratios[24].ci95 :: sha256 3e5cf25e859b68c2 :: random over 162, leg1 N10000 1-macro_auc: paired geometric mean [95% interval]
\newcommand{\PairedFtJciiEFourRandomOverFourthree}{0.97~[0.96, 0.99]}  % experiments/FIGS/data/paired_v1err/finetune/ratios.json :: ratios[25].ratio, ratios[25].ci95 :: sha256 3e5cf25e859b68c2 :: random over 43, leg1 N10000 1-macro_auc: paired geometric mean [95% interval]
\newcommand{\PairedFtJciiEFourRandomOverOneseven}{0.88~[0.82, 0.95]}  % experiments/FIGS/data/paired_v1err/finetune/ratios.json :: ratios[26].ratio, ratios[26].ci95 :: sha256 3e5cf25e859b68c2 :: random over 17, leg1 N10000 1-macro_auc: paired geometric mean [95% interval]
\newcommand{\PairedFtJciiEFourRandomDrawTwoOverRandomDrawOne}{1.00~[0.98, 1.02]}  % experiments/FIGS/data/paired_v1err/finetune/ratios.json :: ratios[27].ratio, ratios[27].ci95 :: sha256 3e5cf25e859b68c2 :: random draw 2 over random draw 1, leg1 N10000 1-macro_auc: paired geometric mean [95% interval]
\newcommand{\PairedFtJciiEFourRandomDrawThreeOverRandomDrawOne}{0.97~[0.95, 0.98]}  % experiments/FIGS/data/paired_v1err/finetune/ratios.json :: ratios[28].ratio, ratios[28].ci95 :: sha256 3e5cf25e859b68c2 :: random draw 3 over random draw 1, leg1 N10000 1-macro_auc: paired geometric mean [95% interval]
\newcommand{\PairedFtJciiEFourRandomDrawThreeOverRandomDrawTwo}{0.97~[0.95, 0.98]}  % experiments/FIGS/data/paired_v1err/finetune/ratios.json :: ratios[29].ratio, ratios[29].ci95 :: sha256 3e5cf25e859b68c2 :: random draw 3 over random draw 2, leg1 N10000 1-macro_auc: paired geometric mean [95% interval]
\newcommand{\PairedFtJciiEFiveOnesixtwoOverOneeighteight}{1.00~[0.98, 1.01]}  % experiments/FIGS/data/paired_v1err/finetune/ratios.json :: ratios[30].ratio, ratios[30].ci95 :: sha256 3e5cf25e859b68c2 :: 162 over 188, leg1 N100000 1-macro_auc: paired geometric mean [95% interval]
\newcommand{\PairedFtJciiEFiveFourthreeOverOneeighteight}{1.09~[1.06, 1.11]}  % experiments/FIGS/data/paired_v1err/finetune/ratios.json :: ratios[31].ratio, ratios[31].ci95 :: sha256 3e5cf25e859b68c2 :: 43 over 188, leg1 N100000 1-macro_auc: paired geometric mean [95% interval]
\newcommand{\PairedFtJciiEFiveOnesevenOverOneeighteight}{1.22~[1.20, 1.23]}  % experiments/FIGS/data/paired_v1err/finetune/ratios.json :: ratios[32].ratio, ratios[32].ci95 :: sha256 3e5cf25e859b68c2 :: 17 over 188, leg1 N100000 1-macro_auc: paired geometric mean [95% interval]
\newcommand{\PairedFtJciiEFiveFourthreeOverOnesixtwo}{1.09~[1.07, 1.11]}  % experiments/FIGS/data/paired_v1err/finetune/ratios.json :: ratios[33].ratio, ratios[33].ci95 :: sha256 3e5cf25e859b68c2 :: 43 over 162, leg1 N100000 1-macro_auc: paired geometric mean [95% interval]
\newcommand{\PairedFtJciiEFiveOnesevenOverOnesixtwo}{1.22~[1.20, 1.24]}  % experiments/FIGS/data/paired_v1err/finetune/ratios.json :: ratios[34].ratio, ratios[34].ci95 :: sha256 3e5cf25e859b68c2 :: 17 over 162, leg1 N100000 1-macro_auc: paired geometric mean [95% interval]
\newcommand{\PairedFtJciiEFiveOnesevenOverFourthree}{1.12~[1.10, 1.14]}  % experiments/FIGS/data/paired_v1err/finetune/ratios.json :: ratios[35].ratio, ratios[35].ci95 :: sha256 3e5cf25e859b68c2 :: 17 over 43, leg1 N100000 1-macro_auc: paired geometric mean [95% interval]
\newcommand{\PairedFtJciiEFiveOnesixtwoMassOverOnesixtwo}{1.018~[1.002, 1.035]}  % experiments/FIGS/data/paired_v1err/finetune/ratios.json :: ratios[36].ratio, ratios[36].ci95 :: sha256 3e5cf25e859b68c2 :: 162+mass over 162, leg1 N100000 1-macro_auc: paired geometric mean [95% interval]
\newcommand{\PairedFtJciiEFiveOnesevenMassOverOneseven}{1.03~[1.01, 1.06]}  % experiments/FIGS/data/paired_v1err/finetune/ratios.json :: ratios[37].ratio, ratios[37].ci95 :: sha256 3e5cf25e859b68c2 :: 17+mass over 17, leg1 N100000 1-macro_auc: paired geometric mean [95% interval]
\newcommand{\PairedFtJciiEFiveRandomOverOneeighteight}{1.15~[1.09, 1.20]}  % experiments/FIGS/data/paired_v1err/finetune/ratios.json :: ratios[38].ratio, ratios[38].ci95 :: sha256 3e5cf25e859b68c2 :: random over 188, leg1 N100000 1-macro_auc: paired geometric mean [95% interval]
\newcommand{\PairedFtJciiEFiveRandomOverOnesixtwo}{1.15~[1.09, 1.20]}  % experiments/FIGS/data/paired_v1err/finetune/ratios.json :: ratios[39].ratio, ratios[39].ci95 :: sha256 3e5cf25e859b68c2 :: random over 162, leg1 N100000 1-macro_auc: paired geometric mean [95% interval]
\newcommand{\PairedFtJciiEFiveRandomOverFourthree}{1.05~[1.04, 1.06]}  % experiments/FIGS/data/paired_v1err/finetune/ratios.json :: ratios[40].ratio, ratios[40].ci95 :: sha256 3e5cf25e859b68c2 :: random over 43, leg1 N100000 1-macro_auc: paired geometric mean [95% interval]
\newcommand{\PairedFtJciiEFiveRandomOverOneseven}{0.95~[0.93, 0.96]}  % experiments/FIGS/data/paired_v1err/finetune/ratios.json :: ratios[41].ratio, ratios[41].ci95 :: sha256 3e5cf25e859b68c2 :: random over 17, leg1 N100000 1-macro_auc: paired geometric mean [95% interval]
\newcommand{\PairedFtJciiEFiveRandomDrawTwoOverRandomDrawOne}{0.98~[0.97, 0.99]}  % experiments/FIGS/data/paired_v1err/finetune/ratios.json :: ratios[42].ratio, ratios[42].ci95 :: sha256 3e5cf25e859b68c2 :: random draw 2 over random draw 1, leg1 N100000 1-macro_auc: paired geometric mean [95% interval]
\newcommand{\PairedFtJciiEFiveRandomDrawThreeOverRandomDrawOne}{1.00~[0.99, 1.02]}  % experiments/FIGS/data/paired_v1err/finetune/ratios.json :: ratios[43].ratio, ratios[43].ci95 :: sha256 3e5cf25e859b68c2 :: random draw 3 over random draw 1, leg1 N100000 1-macro_auc: paired geometric mean [95% interval]
\newcommand{\PairedFtJciiEFiveRandomDrawThreeOverRandomDrawTwo}{1.03~[1.01, 1.04]}  % experiments/FIGS/data/paired_v1err/finetune/ratios.json :: ratios[44].ratio, ratios[44].ci95 :: sha256 3e5cf25e859b68c2 :: random draw 3 over random draw 2, leg1 N100000 1-macro_auc: paired geometric mean [95% interval]
\newcommand{\PairedFtJciiESixOnesixtwoOverOneeighteight}{1.002~[0.998, 1.005]}  % experiments/FIGS/data/paired_v1err/finetune/ratios.json :: ratios[45].ratio, ratios[45].ci95 :: sha256 3e5cf25e859b68c2 :: 162 over 188, leg1 N1000000 1-macro_auc: paired geometric mean [95% interval]
\newcommand{\PairedFtJciiESixFourthreeOverOneeighteight}{1.045~[1.037, 1.053]}  % experiments/FIGS/data/paired_v1err/finetune/ratios.json :: ratios[46].ratio, ratios[46].ci95 :: sha256 3e5cf25e859b68c2 :: 43 over 188, leg1 N1000000 1-macro_auc: paired geometric mean [95% interval]
\newcommand{\PairedFtJciiESixOnesevenOverOneeighteight}{1.14~[1.13, 1.16]}  % experiments/FIGS/data/paired_v1err/finetune/ratios.json :: ratios[47].ratio, ratios[47].ci95 :: sha256 3e5cf25e859b68c2 :: 17 over 188, leg1 N1000000 1-macro_auc: paired geometric mean [95% interval]
\newcommand{\PairedFtJciiESixFourthreeOverOnesixtwo}{1.044~[1.036, 1.051]}  % experiments/FIGS/data/paired_v1err/finetune/ratios.json :: ratios[48].ratio, ratios[48].ci95 :: sha256 3e5cf25e859b68c2 :: 43 over 162, leg1 N1000000 1-macro_auc: paired geometric mean [95% interval]
\newcommand{\PairedFtJciiESixOnesevenOverOnesixtwo}{1.14~[1.13, 1.15]}  % experiments/FIGS/data/paired_v1err/finetune/ratios.json :: ratios[49].ratio, ratios[49].ci95 :: sha256 3e5cf25e859b68c2 :: 17 over 162, leg1 N1000000 1-macro_auc: paired geometric mean [95% interval]
\newcommand{\PairedFtJciiESixOnesevenOverFourthree}{1.09~[1.08, 1.11]}  % experiments/FIGS/data/paired_v1err/finetune/ratios.json :: ratios[50].ratio, ratios[50].ci95 :: sha256 3e5cf25e859b68c2 :: 17 over 43, leg1 N1000000 1-macro_auc: paired geometric mean [95% interval]
\newcommand{\PairedFtJciiESixOnesixtwoMassOverOnesixtwo}{1.019~[1.014, 1.023]}  % experiments/FIGS/data/paired_v1err/finetune/ratios.json :: ratios[51].ratio, ratios[51].ci95 :: sha256 3e5cf25e859b68c2 :: 162+mass over 162, leg1 N1000000 1-macro_auc: paired geometric mean [95% interval]
\newcommand{\PairedFtJciiESixOnesevenMassOverOneseven}{1.03~[1.02, 1.04]}  % experiments/FIGS/data/paired_v1err/finetune/ratios.json :: ratios[52].ratio, ratios[52].ci95 :: sha256 3e5cf25e859b68c2 :: 17+mass over 17, leg1 N1000000 1-macro_auc: paired geometric mean [95% interval]
\newcommand{\PairedFtJciiESixRandomOverOneeighteight}{1.073~[1.067, 1.080]}  % experiments/FIGS/data/paired_v1err/finetune/ratios.json :: ratios[53].ratio, ratios[53].ci95 :: sha256 3e5cf25e859b68c2 :: random over 188, leg1 N1000000 1-macro_auc: paired geometric mean [95% interval]
\newcommand{\PairedFtJciiESixRandomOverOnesixtwo}{1.072~[1.065, 1.078]}  % experiments/FIGS/data/paired_v1err/finetune/ratios.json :: ratios[54].ratio, ratios[54].ci95 :: sha256 3e5cf25e859b68c2 :: random over 162, leg1 N1000000 1-macro_auc: paired geometric mean [95% interval]
\newcommand{\PairedFtJciiESixRandomOverFourthree}{1.03~[1.02, 1.04]}  % experiments/FIGS/data/paired_v1err/finetune/ratios.json :: ratios[55].ratio, ratios[55].ci95 :: sha256 3e5cf25e859b68c2 :: random over 43, leg1 N1000000 1-macro_auc: paired geometric mean [95% interval]
\newcommand{\PairedFtJciiESixRandomOverOneseven}{0.94~[0.93, 0.95]}  % experiments/FIGS/data/paired_v1err/finetune/ratios.json :: ratios[56].ratio, ratios[56].ci95 :: sha256 3e5cf25e859b68c2 :: random over 17, leg1 N1000000 1-macro_auc: paired geometric mean [95% interval]
\newcommand{\PairedFtJciiESixRandomDrawTwoOverRandomDrawOne}{1.00~[0.99, 1.01]}  % experiments/FIGS/data/paired_v1err/finetune/ratios.json :: ratios[57].ratio, ratios[57].ci95 :: sha256 3e5cf25e859b68c2 :: random draw 2 over random draw 1, leg1 N1000000 1-macro_auc: paired geometric mean [95% interval]
\newcommand{\PairedFtJciiESixRandomDrawThreeOverRandomDrawOne}{0.996~[0.989, 1.003]}  % experiments/FIGS/data/paired_v1err/finetune/ratios.json :: ratios[58].ratio, ratios[58].ci95 :: sha256 3e5cf25e859b68c2 :: random draw 3 over random draw 1, leg1 N1000000 1-macro_auc: paired geometric mean [95% interval]
\newcommand{\PairedFtJciiESixRandomDrawThreeOverRandomDrawTwo}{0.997~[0.989, 1.004]}  % experiments/FIGS/data/paired_v1err/finetune/ratios.json :: ratios[59].ratio, ratios[59].ci95 :: sha256 3e5cf25e859b68c2 :: random draw 3 over random draw 2, leg1 N1000000 1-macro_auc: paired geometric mean [95% interval]
\newcommand{\PairedFtJcEThreeOnesixtwoOverOneeighteight}{0.96~[0.83, 1.11]}  % experiments/FIGS/data/paired_v1err/finetune/ratios.json :: ratios[60].ratio, ratios[60].ci95 :: sha256 3e5cf25e859b68c2 :: 162 over 188, leg2 N1000 1-macro_auc: paired geometric mean [95% interval]
\newcommand{\PairedFtJcEThreeFourthreeOverOneeighteight}{1.01~[0.84, 1.21]}  % experiments/FIGS/data/paired_v1err/finetune/ratios.json :: ratios[61].ratio, ratios[61].ci95 :: sha256 3e5cf25e859b68c2 :: 43 over 188, leg2 N1000 1-macro_auc: paired geometric mean [95% interval]
\newcommand{\PairedFtJcEThreeOnesevenOverOneeighteight}{1.07~[0.93, 1.23]}  % experiments/FIGS/data/paired_v1err/finetune/ratios.json :: ratios[62].ratio, ratios[62].ci95 :: sha256 3e5cf25e859b68c2 :: 17 over 188, leg2 N1000 1-macro_auc: paired geometric mean [95% interval]
\newcommand{\PairedFtJcEThreeFourthreeOverOnesixtwo}{1.05~[0.96, 1.15]}  % experiments/FIGS/data/paired_v1err/finetune/ratios.json :: ratios[63].ratio, ratios[63].ci95 :: sha256 3e5cf25e859b68c2 :: 43 over 162, leg2 N1000 1-macro_auc: paired geometric mean [95% interval]
\newcommand{\PairedFtJcEThreeOnesevenOverOnesixtwo}{1.12~[1.07, 1.17]}  % experiments/FIGS/data/paired_v1err/finetune/ratios.json :: ratios[64].ratio, ratios[64].ci95 :: sha256 3e5cf25e859b68c2 :: 17 over 162, leg2 N1000 1-macro_auc: paired geometric mean [95% interval]
\newcommand{\PairedFtJcEThreeOnesevenOverFourthree}{1.06~[0.98, 1.15]}  % experiments/FIGS/data/paired_v1err/finetune/ratios.json :: ratios[65].ratio, ratios[65].ci95 :: sha256 3e5cf25e859b68c2 :: 17 over 43, leg2 N1000 1-macro_auc: paired geometric mean [95% interval]
\newcommand{\PairedFtJcEThreeOnesixtwoMassOverOnesixtwo}{0.91~[0.88, 0.94]}  % experiments/FIGS/data/paired_v1err/finetune/ratios.json :: ratios[66].ratio, ratios[66].ci95 :: sha256 3e5cf25e859b68c2 :: 162+mass over 162, leg2 N1000 1-macro_auc: paired geometric mean [95% interval]
\newcommand{\PairedFtJcEThreeOnesevenMassOverOneseven}{1.02~[0.92, 1.13]}  % experiments/FIGS/data/paired_v1err/finetune/ratios.json :: ratios[67].ratio, ratios[67].ci95 :: sha256 3e5cf25e859b68c2 :: 17+mass over 17, leg2 N1000 1-macro_auc: paired geometric mean [95% interval]
\newcommand{\PairedFtJcEThreeRandomOverOneeighteight}{1.01~[0.78, 1.32]}  % experiments/FIGS/data/paired_v1err/finetune/ratios.json :: ratios[68].ratio, ratios[68].ci95 :: sha256 3e5cf25e859b68c2 :: random over 188, leg2 N1000 1-macro_auc: paired geometric mean [95% interval]
\newcommand{\PairedFtJcEThreeRandomOverOnesixtwo}{1.12~[1.04, 1.20]}  % experiments/FIGS/data/paired_v1err/finetune/ratios.json :: ratios[69].ratio, ratios[69].ci95 :: sha256 3e5cf25e859b68c2 :: random over 162, leg2 N1000 1-macro_auc: paired geometric mean [95% interval]
\newcommand{\PairedFtJcEThreeRandomOverFourthree}{1.06~[0.79, 1.41]}  % experiments/FIGS/data/paired_v1err/finetune/ratios.json :: ratios[70].ratio, ratios[70].ci95 :: sha256 3e5cf25e859b68c2 :: random over 43, leg2 N1000 1-macro_auc: paired geometric mean [95% interval]
\newcommand{\PairedFtJcEThreeRandomOverOneseven}{0.98~[0.84, 1.13]}  % experiments/FIGS/data/paired_v1err/finetune/ratios.json :: ratios[71].ratio, ratios[71].ci95 :: sha256 3e5cf25e859b68c2 :: random over 17, leg2 N1000 1-macro_auc: paired geometric mean [95% interval]
\newcommand{\PairedFtJcEThreeRandomDrawTwoOverRandomDrawOne}{0.954~[0.949, 0.958]}  % experiments/FIGS/data/paired_v1err/finetune/ratios.json :: ratios[72].ratio, ratios[72].ci95 :: sha256 3e5cf25e859b68c2 :: random draw 2 over random draw 1, leg2 N1000 1-macro_auc: paired geometric mean [95% interval]
\newcommand{\PairedFtJcEThreeRandomDrawThreeOverRandomDrawOne}{1.013~[1.009, 1.017]}  % experiments/FIGS/data/paired_v1err/finetune/ratios.json :: ratios[73].ratio, ratios[73].ci95 :: sha256 3e5cf25e859b68c2 :: random draw 3 over random draw 1, leg2 N1000 1-macro_auc: paired geometric mean [95% interval]
\newcommand{\PairedFtJcEThreeRandomDrawThreeOverRandomDrawTwo}{1.062~[1.057, 1.067]}  % experiments/FIGS/data/paired_v1err/finetune/ratios.json :: ratios[74].ratio, ratios[74].ci95 :: sha256 3e5cf25e859b68c2 :: random draw 3 over random draw 2, leg2 N1000 1-macro_auc: paired geometric mean [95% interval]
\newcommand{\PairedFtJcEFourOnesixtwoOverOneeighteight}{0.98~[0.93, 1.02]}  % experiments/FIGS/data/paired_v1err/finetune/ratios.json :: ratios[75].ratio, ratios[75].ci95 :: sha256 3e5cf25e859b68c2 :: 162 over 188, leg2 N10000 1-macro_auc: paired geometric mean [95% interval]
\newcommand{\PairedFtJcEFourFourthreeOverOneeighteight}{0.98~[0.93, 1.04]}  % experiments/FIGS/data/paired_v1err/finetune/ratios.json :: ratios[76].ratio, ratios[76].ci95 :: sha256 3e5cf25e859b68c2 :: 43 over 188, leg2 N10000 1-macro_auc: paired geometric mean [95% interval]
\newcommand{\PairedFtJcEFourOnesevenOverOneeighteight}{1.07~[1.02, 1.11]}  % experiments/FIGS/data/paired_v1err/finetune/ratios.json :: ratios[77].ratio, ratios[77].ci95 :: sha256 3e5cf25e859b68c2 :: 17 over 188, leg2 N10000 1-macro_auc: paired geometric mean [95% interval]
\newcommand{\PairedFtJcEFourFourthreeOverOnesixtwo}{1.01~[0.97, 1.04]}  % experiments/FIGS/data/paired_v1err/finetune/ratios.json :: ratios[78].ratio, ratios[78].ci95 :: sha256 3e5cf25e859b68c2 :: 43 over 162, leg2 N10000 1-macro_auc: paired geometric mean [95% interval]
\newcommand{\PairedFtJcEFourOnesevenOverOnesixtwo}{1.09~[1.06, 1.12]}  % experiments/FIGS/data/paired_v1err/finetune/ratios.json :: ratios[79].ratio, ratios[79].ci95 :: sha256 3e5cf25e859b68c2 :: 17 over 162, leg2 N10000 1-macro_auc: paired geometric mean [95% interval]
\newcommand{\PairedFtJcEFourOnesevenOverFourthree}{1.08~[1.07, 1.10]}  % experiments/FIGS/data/paired_v1err/finetune/ratios.json :: ratios[80].ratio, ratios[80].ci95 :: sha256 3e5cf25e859b68c2 :: 17 over 43, leg2 N10000 1-macro_auc: paired geometric mean [95% interval]
\newcommand{\PairedFtJcEFourOnesixtwoMassOverOnesixtwo}{1.01~[0.98, 1.04]}  % experiments/FIGS/data/paired_v1err/finetune/ratios.json :: ratios[81].ratio, ratios[81].ci95 :: sha256 3e5cf25e859b68c2 :: 162+mass over 162, leg2 N10000 1-macro_auc: paired geometric mean [95% interval]
\newcommand{\PairedFtJcEFourOnesevenMassOverOneseven}{1.05~[1.02, 1.08]}  % experiments/FIGS/data/paired_v1err/finetune/ratios.json :: ratios[82].ratio, ratios[82].ci95 :: sha256 3e5cf25e859b68c2 :: 17+mass over 17, leg2 N10000 1-macro_auc: paired geometric mean [95% interval]
\newcommand{\PairedFtJcEFourRandomOverOneeighteight}{1.08~[1.03, 1.14]}  % experiments/FIGS/data/paired_v1err/finetune/ratios.json :: ratios[83].ratio, ratios[83].ci95 :: sha256 3e5cf25e859b68c2 :: random over 188, leg2 N10000 1-macro_auc: paired geometric mean [95% interval]
\newcommand{\PairedFtJcEFourRandomOverOnesixtwo}{1.11~[0.99, 1.23]}  % experiments/FIGS/data/paired_v1err/finetune/ratios.json :: ratios[84].ratio, ratios[84].ci95 :: sha256 3e5cf25e859b68c2 :: random over 162, leg2 N10000 1-macro_auc: paired geometric mean [95% interval]
\newcommand{\PairedFtJcEFourRandomOverFourthree}{1.10~[0.99, 1.21]}  % experiments/FIGS/data/paired_v1err/finetune/ratios.json :: ratios[85].ratio, ratios[85].ci95 :: sha256 3e5cf25e859b68c2 :: random over 43, leg2 N10000 1-macro_auc: paired geometric mean [95% interval]
\newcommand{\PairedFtJcEFourRandomOverOneseven}{1.02~[0.94, 1.10]}  % experiments/FIGS/data/paired_v1err/finetune/ratios.json :: ratios[86].ratio, ratios[86].ci95 :: sha256 3e5cf25e859b68c2 :: random over 17, leg2 N10000 1-macro_auc: paired geometric mean [95% interval]
\newcommand{\PairedFtJcEFourRandomDrawTwoOverRandomDrawOne}{0.960~[0.955, 0.965]}  % experiments/FIGS/data/paired_v1err/finetune/ratios.json :: ratios[87].ratio, ratios[87].ci95 :: sha256 3e5cf25e859b68c2 :: random draw 2 over random draw 1, leg2 N10000 1-macro_auc: paired geometric mean [95% interval]
\newcommand{\PairedFtJcEFourRandomDrawThreeOverRandomDrawOne}{1.008~[1.004, 1.012]}  % experiments/FIGS/data/paired_v1err/finetune/ratios.json :: ratios[88].ratio, ratios[88].ci95 :: sha256 3e5cf25e859b68c2 :: random draw 3 over random draw 1, leg2 N10000 1-macro_auc: paired geometric mean [95% interval]
\newcommand{\PairedFtJcEFourRandomDrawThreeOverRandomDrawTwo}{1.050~[1.045, 1.055]}  % experiments/FIGS/data/paired_v1err/finetune/ratios.json :: ratios[89].ratio, ratios[89].ci95 :: sha256 3e5cf25e859b68c2 :: random draw 3 over random draw 2, leg2 N10000 1-macro_auc: paired geometric mean [95% interval]
\newcommand{\PairedFtJcEFiveOnesixtwoOverOneeighteight}{0.983~[0.968, 0.998]}  % experiments/FIGS/data/paired_v1err/finetune/ratios.json :: ratios[90].ratio, ratios[90].ci95 :: sha256 3e5cf25e859b68c2 :: 162 over 188, leg2 N100000 1-macro_auc: paired geometric mean [95% interval]
\newcommand{\PairedFtJcEFiveFourthreeOverOneeighteight}{1.015~[1.002, 1.029]}  % experiments/FIGS/data/paired_v1err/finetune/ratios.json :: ratios[91].ratio, ratios[91].ci95 :: sha256 3e5cf25e859b68c2 :: 43 over 188, leg2 N100000 1-macro_auc: paired geometric mean [95% interval]
\newcommand{\PairedFtJcEFiveOnesevenOverOneeighteight}{1.07~[1.06, 1.09]}  % experiments/FIGS/data/paired_v1err/finetune/ratios.json :: ratios[92].ratio, ratios[92].ci95 :: sha256 3e5cf25e859b68c2 :: 17 over 188, leg2 N100000 1-macro_auc: paired geometric mean [95% interval]
\newcommand{\PairedFtJcEFiveFourthreeOverOnesixtwo}{1.03~[1.02, 1.05]}  % experiments/FIGS/data/paired_v1err/finetune/ratios.json :: ratios[93].ratio, ratios[93].ci95 :: sha256 3e5cf25e859b68c2 :: 43 over 162, leg2 N100000 1-macro_auc: paired geometric mean [95% interval]
\newcommand{\PairedFtJcEFiveOnesevenOverOnesixtwo}{1.091~[1.088, 1.094]}  % experiments/FIGS/data/paired_v1err/finetune/ratios.json :: ratios[94].ratio, ratios[94].ci95 :: sha256 3e5cf25e859b68c2 :: 17 over 162, leg2 N100000 1-macro_auc: paired geometric mean [95% interval]
\newcommand{\PairedFtJcEFiveOnesevenOverFourthree}{1.06~[1.04, 1.07]}  % experiments/FIGS/data/paired_v1err/finetune/ratios.json :: ratios[95].ratio, ratios[95].ci95 :: sha256 3e5cf25e859b68c2 :: 17 over 43, leg2 N100000 1-macro_auc: paired geometric mean [95% interval]
\newcommand{\PairedFtJcEFiveOnesixtwoMassOverOnesixtwo}{1.016~[1.008, 1.024]}  % experiments/FIGS/data/paired_v1err/finetune/ratios.json :: ratios[96].ratio, ratios[96].ci95 :: sha256 3e5cf25e859b68c2 :: 162+mass over 162, leg2 N100000 1-macro_auc: paired geometric mean [95% interval]
\newcommand{\PairedFtJcEFiveOnesevenMassOverOneseven}{1.02~[1.01, 1.03]}  % experiments/FIGS/data/paired_v1err/finetune/ratios.json :: ratios[97].ratio, ratios[97].ci95 :: sha256 3e5cf25e859b68c2 :: 17+mass over 17, leg2 N100000 1-macro_auc: paired geometric mean [95% interval]
\newcommand{\PairedFtJcEFiveRandomOverOneeighteight}{1.044~[1.039, 1.049]}  % experiments/FIGS/data/paired_v1err/finetune/ratios.json :: ratios[98].ratio, ratios[98].ci95 :: sha256 3e5cf25e859b68c2 :: random over 188, leg2 N100000 1-macro_auc: paired geometric mean [95% interval]
\newcommand{\PairedFtJcEFiveRandomOverOnesixtwo}{1.06~[1.04, 1.09]}  % experiments/FIGS/data/paired_v1err/finetune/ratios.json :: ratios[99].ratio, ratios[99].ci95 :: sha256 3e5cf25e859b68c2 :: random over 162, leg2 N100000 1-macro_auc: paired geometric mean [95% interval]
\newcommand{\PairedFtJcEFiveRandomOverFourthree}{1.03~[1.01, 1.04]}  % experiments/FIGS/data/paired_v1err/finetune/ratios.json :: ratios[100].ratio, ratios[100].ci95 :: sha256 3e5cf25e859b68c2 :: random over 43, leg2 N100000 1-macro_auc: paired geometric mean [95% interval]
\newcommand{\PairedFtJcEFiveRandomOverOneseven}{0.97~[0.95, 0.99]}  % experiments/FIGS/data/paired_v1err/finetune/ratios.json :: ratios[101].ratio, ratios[101].ci95 :: sha256 3e5cf25e859b68c2 :: random over 17, leg2 N100000 1-macro_auc: paired geometric mean [95% interval]
\newcommand{\PairedFtJcEFiveRandomDrawTwoOverRandomDrawOne}{0.992~[0.988, 0.997]}  % experiments/FIGS/data/paired_v1err/finetune/ratios.json :: ratios[102].ratio, ratios[102].ci95 :: sha256 3e5cf25e859b68c2 :: random draw 2 over random draw 1, leg2 N100000 1-macro_auc: paired geometric mean [95% interval]
\newcommand{\PairedFtJcEFiveRandomDrawThreeOverRandomDrawOne}{1.007~[1.003, 1.011]}  % experiments/FIGS/data/paired_v1err/finetune/ratios.json :: ratios[103].ratio, ratios[103].ci95 :: sha256 3e5cf25e859b68c2 :: random draw 3 over random draw 1, leg2 N100000 1-macro_auc: paired geometric mean [95% interval]
\newcommand{\PairedFtJcEFiveRandomDrawThreeOverRandomDrawTwo}{1.015~[1.011, 1.020]}  % experiments/FIGS/data/paired_v1err/finetune/ratios.json :: ratios[104].ratio, ratios[104].ci95 :: sha256 3e5cf25e859b68c2 :: random draw 3 over random draw 2, leg2 N100000 1-macro_auc: paired geometric mean [95% interval]
\newcommand{\PairedFtJcESixOnesixtwoOverOneeighteight}{1.00~[0.99, 1.01]}  % experiments/FIGS/data/paired_v1err/finetune/ratios.json :: ratios[105].ratio, ratios[105].ci95 :: sha256 3e5cf25e859b68c2 :: 162 over 188, leg2 N1000000 1-macro_auc: paired geometric mean [95% interval]
\newcommand{\PairedFtJcESixFourthreeOverOneeighteight}{1.02~[1.01, 1.03]}  % experiments/FIGS/data/paired_v1err/finetune/ratios.json :: ratios[106].ratio, ratios[106].ci95 :: sha256 3e5cf25e859b68c2 :: 43 over 188, leg2 N1000000 1-macro_auc: paired geometric mean [95% interval]
\newcommand{\PairedFtJcESixOnesevenOverOneeighteight}{1.08~[1.07, 1.09]}  % experiments/FIGS/data/paired_v1err/finetune/ratios.json :: ratios[107].ratio, ratios[107].ci95 :: sha256 3e5cf25e859b68c2 :: 17 over 188, leg2 N1000000 1-macro_auc: paired geometric mean [95% interval]
\newcommand{\PairedFtJcESixFourthreeOverOnesixtwo}{1.02~[1.01, 1.03]}  % experiments/FIGS/data/paired_v1err/finetune/ratios.json :: ratios[108].ratio, ratios[108].ci95 :: sha256 3e5cf25e859b68c2 :: 43 over 162, leg2 N1000000 1-macro_auc: paired geometric mean [95% interval]
\newcommand{\PairedFtJcESixOnesevenOverOnesixtwo}{1.08~[1.07, 1.09]}  % experiments/FIGS/data/paired_v1err/finetune/ratios.json :: ratios[109].ratio, ratios[109].ci95 :: sha256 3e5cf25e859b68c2 :: 17 over 162, leg2 N1000000 1-macro_auc: paired geometric mean [95% interval]
\newcommand{\PairedFtJcESixOnesevenOverFourthree}{1.06~[1.05, 1.07]}  % experiments/FIGS/data/paired_v1err/finetune/ratios.json :: ratios[110].ratio, ratios[110].ci95 :: sha256 3e5cf25e859b68c2 :: 17 over 43, leg2 N1000000 1-macro_auc: paired geometric mean [95% interval]
\newcommand{\PairedFtJcESixOnesixtwoMassOverOnesixtwo}{1.012~[1.004, 1.020]}  % experiments/FIGS/data/paired_v1err/finetune/ratios.json :: ratios[111].ratio, ratios[111].ci95 :: sha256 3e5cf25e859b68c2 :: 162+mass over 162, leg2 N1000000 1-macro_auc: paired geometric mean [95% interval]
\newcommand{\PairedFtJcESixOnesevenMassOverOneseven}{1.02~[1.01, 1.04]}  % experiments/FIGS/data/paired_v1err/finetune/ratios.json :: ratios[112].ratio, ratios[112].ci95 :: sha256 3e5cf25e859b68c2 :: 17+mass over 17, leg2 N1000000 1-macro_auc: paired geometric mean [95% interval]
\newcommand{\PairedFtJcESixRandomOverOneeighteight}{1.046~[1.038, 1.053]}  % experiments/FIGS/data/paired_v1err/finetune/ratios.json :: ratios[113].ratio, ratios[113].ci95 :: sha256 3e5cf25e859b68c2 :: random over 188, leg2 N1000000 1-macro_auc: paired geometric mean [95% interval]
\newcommand{\PairedFtJcESixRandomOverOnesixtwo}{1.047~[1.038, 1.055]}  % experiments/FIGS/data/paired_v1err/finetune/ratios.json :: ratios[114].ratio, ratios[114].ci95 :: sha256 3e5cf25e859b68c2 :: random over 162, leg2 N1000000 1-macro_auc: paired geometric mean [95% interval]
\newcommand{\PairedFtJcESixRandomOverFourthree}{1.03~[1.02, 1.04]}  % experiments/FIGS/data/paired_v1err/finetune/ratios.json :: ratios[115].ratio, ratios[115].ci95 :: sha256 3e5cf25e859b68c2 :: random over 43, leg2 N1000000 1-macro_auc: paired geometric mean [95% interval]
\newcommand{\PairedFtJcESixRandomOverOneseven}{0.97~[0.96, 0.98]}  % experiments/FIGS/data/paired_v1err/finetune/ratios.json :: ratios[116].ratio, ratios[116].ci95 :: sha256 3e5cf25e859b68c2 :: random over 17, leg2 N1000000 1-macro_auc: paired geometric mean [95% interval]
\newcommand{\PairedFtJcESixRandomDrawTwoOverRandomDrawOne}{1.001~[0.998, 1.004]}  % experiments/FIGS/data/paired_v1err/finetune/ratios.json :: ratios[117].ratio, ratios[117].ci95 :: sha256 3e5cf25e859b68c2 :: random draw 2 over random draw 1, leg2 N1000000 1-macro_auc: paired geometric mean [95% interval]
\newcommand{\PairedFtJcESixRandomDrawThreeOverRandomDrawOne}{1.005~[1.002, 1.008]}  % experiments/FIGS/data/paired_v1err/finetune/ratios.json :: ratios[118].ratio, ratios[118].ci95 :: sha256 3e5cf25e859b68c2 :: random draw 3 over random draw 1, leg2 N1000000 1-macro_auc: paired geometric mean [95% interval]
\newcommand{\PairedFtJcESixRandomDrawThreeOverRandomDrawTwo}{1.004~[1.001, 1.007]}  % experiments/FIGS/data/paired_v1err/finetune/ratios.json :: ratios[119].ratio, ratios[119].ci95 :: sha256 3e5cf25e859b68c2 :: random draw 3 over random draw 2, leg2 N1000000 1-macro_auc: paired geometric mean [95% interval]
\newcommand{\PairedMassResMlpOnesixtwoOverOneeighteight}{0.990~[0.977, 1.003]}  % experiments/FIGS/data/paired_v1err/mass/ratios.json :: ratios[0].ratio, ratios[0].ci95 :: sha256 d877a7f087260f64 :: 162 over 188, resolution mlp sigma_eff: paired geometric mean [95% interval]
\newcommand{\PairedMassResMlpFourthreeOverOneeighteight}{0.98~[0.96, 0.99]}  % experiments/FIGS/data/paired_v1err/mass/ratios.json :: ratios[1].ratio, ratios[1].ci95 :: sha256 d877a7f087260f64 :: 43 over 188, resolution mlp sigma_eff: paired geometric mean [95% interval]
\newcommand{\PairedMassResMlpOnesevenOverOneeighteight}{0.96~[0.94, 0.98]}  % experiments/FIGS/data/paired_v1err/mass/ratios.json :: ratios[2].ratio, ratios[2].ci95 :: sha256 d877a7f087260f64 :: 17 over 188, resolution mlp sigma_eff: paired geometric mean [95% interval]
\newcommand{\PairedMassResMlpFourthreeOverOnesixtwo}{0.989~[0.978, 0.999]}  % experiments/FIGS/data/paired_v1err/mass/ratios.json :: ratios[3].ratio, ratios[3].ci95 :: sha256 d877a7f087260f64 :: 43 over 162, resolution mlp sigma_eff: paired geometric mean [95% interval]
\newcommand{\PairedMassResMlpOnesevenOverOnesixtwo}{0.97~[0.94, 0.99]}  % experiments/FIGS/data/paired_v1err/mass/ratios.json :: ratios[4].ratio, ratios[4].ci95 :: sha256 d877a7f087260f64 :: 17 over 162, resolution mlp sigma_eff: paired geometric mean [95% interval]
\newcommand{\PairedMassResMlpOnesevenOverFourthree}{0.979~[0.961, 0.997]}  % experiments/FIGS/data/paired_v1err/mass/ratios.json :: ratios[5].ratio, ratios[5].ci95 :: sha256 d877a7f087260f64 :: 17 over 43, resolution mlp sigma_eff: paired geometric mean [95% interval]
\newcommand{\PairedMassResMlpOnesixtwoMassOverOnesixtwo}{0.80~[0.79, 0.81]}  % experiments/FIGS/data/paired_v1err/mass/ratios.json :: ratios[6].ratio, ratios[6].ci95 :: sha256 d877a7f087260f64 :: 162+mass over 162, resolution mlp sigma_eff: paired geometric mean [95% interval]
\newcommand{\PairedMassResMlpOnesevenMassOverOneseven}{0.83~[0.82, 0.84]}  % experiments/FIGS/data/paired_v1err/mass/ratios.json :: ratios[7].ratio, ratios[7].ci95 :: sha256 d877a7f087260f64 :: 17+mass over 17, resolution mlp sigma_eff: paired geometric mean [95% interval]
\newcommand{\PairedMassResRidgeOnesixtwoOverOneeighteight}{0.97~[0.96, 0.98]}  % experiments/FIGS/data/paired_v1err/mass/ratios.json :: ratios[8].ratio, ratios[8].ci95 :: sha256 d877a7f087260f64 :: 162 over 188, resolution ridge sigma_eff: paired geometric mean [95% interval]
\newcommand{\PairedMassResRidgeFourthreeOverOneeighteight}{0.95~[0.94, 0.96]}  % experiments/FIGS/data/paired_v1err/mass/ratios.json :: ratios[9].ratio, ratios[9].ci95 :: sha256 d877a7f087260f64 :: 43 over 188, resolution ridge sigma_eff: paired geometric mean [95% interval]
\newcommand{\PairedMassResRidgeOnesevenOverOneeighteight}{0.88~[0.85, 0.92]}  % experiments/FIGS/data/paired_v1err/mass/ratios.json :: ratios[10].ratio, ratios[10].ci95 :: sha256 d877a7f087260f64 :: 17 over 188, resolution ridge sigma_eff: paired geometric mean [95% interval]
\newcommand{\PairedMassResRidgeFourthreeOverOnesixtwo}{0.98~[0.97, 0.99]}  % experiments/FIGS/data/paired_v1err/mass/ratios.json :: ratios[11].ratio, ratios[11].ci95 :: sha256 d877a7f087260f64 :: 43 over 162, resolution ridge sigma_eff: paired geometric mean [95% interval]
\newcommand{\PairedMassResRidgeOnesevenOverOnesixtwo}{0.91~[0.87, 0.95]}  % experiments/FIGS/data/paired_v1err/mass/ratios.json :: ratios[12].ratio, ratios[12].ci95 :: sha256 d877a7f087260f64 :: 17 over 162, resolution ridge sigma_eff: paired geometric mean [95% interval]
\newcommand{\PairedMassResRidgeOnesevenOverFourthree}{0.93~[0.89, 0.97]}  % experiments/FIGS/data/paired_v1err/mass/ratios.json :: ratios[13].ratio, ratios[13].ci95 :: sha256 d877a7f087260f64 :: 17 over 43, resolution ridge sigma_eff: paired geometric mean [95% interval]
\newcommand{\PairedMassResRidgeOnesixtwoMassOverOnesixtwo}{0.678~[0.672, 0.684]}  % experiments/FIGS/data/paired_v1err/mass/ratios.json :: ratios[14].ratio, ratios[14].ci95 :: sha256 d877a7f087260f64 :: 162+mass over 162, resolution ridge sigma_eff: paired geometric mean [95% interval]
\newcommand{\PairedMassResRidgeOnesevenMassOverOneseven}{0.72~[0.69, 0.76]}  % experiments/FIGS/data/paired_v1err/mass/ratios.json :: ratios[15].ratio, ratios[15].ci95 :: sha256 d877a7f087260f64 :: 17+mass over 17, resolution ridge sigma_eff: paired geometric mean [95% interval]
\newcommand{\VcbNSignal}{1{,}659}  % experiments/FIGS/data/probe_ladder_vcbwindow_mlp2/sall.json :: tasks.bc_vs_rest.n_signal_test :: sha256 5b8b5c7c28602f93 :: X->bc test jets in the window
\newcommand{\VcbNBackground}{19{,}244}  % experiments/FIGS/data/probe_ladder_vcbwindow_mlp2/sall.json :: tasks.bc_vs_rest.n_background_test :: sha256 5b8b5c7c28602f93 :: bq, cs, bqq and QCD test jets in the window
\newcommand{\VcbEpsSixty}{60}  % experiments/FIGS/data/probe_ladder_vcbwindow_mlp2/sall.json :: tasks.bc_vs_rest.eps_s :: sha256 5b8b5c7c28602f93 :: signal efficiency, percent
\newcommand{\VcbOmaLinearOnesixtwo}{\ensuremath{(3.17 \pm 0.05)\times10^{-2}}}  % experiments/FIGS/data/probe_ladder_vcbwindow_mlp2/sall.json :: tasks.bc_vs_rest.arms.l162-s*.linear.log1m_auc :: sha256 5b8b5c7c28602f93 :: 1-AUC = exp of it, mean +- SD over 5 seeds
\newcommand{\VcbLogOmaLinearOnesixtwo}{\ensuremath{-}3.451}  % experiments/FIGS/data/probe_ladder_vcbwindow_mlp2/sall.json :: tasks.bc_vs_rest.arms.l162-s*.linear.log1m_auc :: sha256 5b8b5c7c28602f93 :: mean over seeds
\newcommand{\VcbRejLinearSixtyOnesixtwo}{207\,\ensuremath{\pm}\,19}  % experiments/FIGS/data/probe_ladder_vcbwindow_mlp2/sall.json :: tasks.bc_vs_rest.arms.l162-s*.linear.rejection_at['0.60'].rejection :: sha256 5b8b5c7c28602f93 :: background rejection at this signal efficiency, mean +- SD over seeds
\newcommand{\VcbBkgLeftLinearSixtyOnesixtwo}{93.4}  % experiments/FIGS/data/probe_ladder_vcbwindow_mlp2/sall.json :: tasks.bc_vs_rest.arms.l162-s*.linear.rejection_at['0.60'].n_bkg_pass :: sha256 5b8b5c7c28602f93 :: background jets passing the cut, mean over seeds
\newcommand{\VcbOmaLinearOneseven}{\ensuremath{(5.21 \pm 0.23)\times10^{-2}}}  % experiments/FIGS/data/probe_ladder_vcbwindow_mlp2/sall.json :: tasks.bc_vs_rest.arms.r16q1-s*.linear.log1m_auc :: sha256 5b8b5c7c28602f93 :: 1-AUC = exp of it, mean +- SD over 5 seeds
\newcommand{\VcbLogOmaLinearOneseven}{\ensuremath{-}2.955}  % experiments/FIGS/data/probe_ladder_vcbwindow_mlp2/sall.json :: tasks.bc_vs_rest.arms.r16q1-s*.linear.log1m_auc :: sha256 5b8b5c7c28602f93 :: mean over seeds
\newcommand{\VcbRejLinearSixtyOneseven}{67\,\ensuremath{\pm}\,14}  % experiments/FIGS/data/probe_ladder_vcbwindow_mlp2/sall.json :: tasks.bc_vs_rest.arms.r16q1-s*.linear.rejection_at['0.60'].rejection :: sha256 5b8b5c7c28602f93 :: background rejection at this signal efficiency, mean +- SD over seeds
\newcommand{\VcbBkgLeftLinearSixtyOneseven}{295.0}  % experiments/FIGS/data/probe_ladder_vcbwindow_mlp2/sall.json :: tasks.bc_vs_rest.arms.r16q1-s*.linear.rejection_at['0.60'].n_bkg_pass :: sha256 5b8b5c7c28602f93 :: background jets passing the cut, mean over seeds
\newcommand{\VcbOmaRatioLinear}{1.64}  % experiments/FIGS/data/probe_ladder_vcbwindow_mlp2/sall.json :: tasks.bc_vs_rest.arms.*.linear.log1m_auc :: sha256 5b8b5c7c28602f93 :: seed mean of 1-AUC, 17 classes over 162
\newcommand{\VcbFactorLinear}{1.64}  % experiments/FIGS/data/probe_ladder_vcbwindow_mlp2/sall.json :: tasks.bc_vs_rest.arms.*.linear.log1m_auc :: sha256 5b8b5c7c28602f93 :: exp of the difference of the seed means of log(1-AUC), 17 classes over 162
\newcommand{\VcbOmaMlpOnesixtwo}{\ensuremath{(3.102 \pm 0.023)\times10^{-2}}}  % experiments/FIGS/data/probe_ladder_vcbwindow_mlp2/sall.json :: tasks.bc_vs_rest.arms.l162-s*.mlp.log1m_auc :: sha256 5b8b5c7c28602f93 :: 1-AUC = exp of it, mean +- SD over 5 seeds
\newcommand{\VcbLogOmaMlpOnesixtwo}{\ensuremath{-}3.473}  % experiments/FIGS/data/probe_ladder_vcbwindow_mlp2/sall.json :: tasks.bc_vs_rest.arms.l162-s*.mlp.log1m_auc :: sha256 5b8b5c7c28602f93 :: mean over seeds
\newcommand{\VcbRejMlpSixtyOnesixtwo}{230\,\ensuremath{\pm}\,31}  % experiments/FIGS/data/probe_ladder_vcbwindow_mlp2/sall.json :: tasks.bc_vs_rest.arms.l162-s*.mlp.rejection_at['0.60'].rejection :: sha256 5b8b5c7c28602f93 :: background rejection at this signal efficiency, mean +- SD over seeds
\newcommand{\VcbBkgLeftMlpSixtyOnesixtwo}{84.8}  % experiments/FIGS/data/probe_ladder_vcbwindow_mlp2/sall.json :: tasks.bc_vs_rest.arms.l162-s*.mlp.rejection_at['0.60'].n_bkg_pass :: sha256 5b8b5c7c28602f93 :: background jets passing the cut, mean over seeds
\newcommand{\VcbOmaMlpOneseven}{\ensuremath{(4.66 \pm 0.19)\times10^{-2}}}  % experiments/FIGS/data/probe_ladder_vcbwindow_mlp2/sall.json :: tasks.bc_vs_rest.arms.r16q1-s*.mlp.log1m_auc :: sha256 5b8b5c7c28602f93 :: 1-AUC = exp of it, mean +- SD over 5 seeds
\newcommand{\VcbLogOmaMlpOneseven}{\ensuremath{-}3.066}  % experiments/FIGS/data/probe_ladder_vcbwindow_mlp2/sall.json :: tasks.bc_vs_rest.arms.r16q1-s*.mlp.log1m_auc :: sha256 5b8b5c7c28602f93 :: mean over seeds
\newcommand{\VcbRejMlpSixtyOneseven}{78\,\ensuremath{\pm}\,13}  % experiments/FIGS/data/probe_ladder_vcbwindow_mlp2/sall.json :: tasks.bc_vs_rest.arms.r16q1-s*.mlp.rejection_at['0.60'].rejection :: sha256 5b8b5c7c28602f93 :: background rejection at this signal efficiency, mean +- SD over seeds
\newcommand{\VcbBkgLeftMlpSixtyOneseven}{250.8}  % experiments/FIGS/data/probe_ladder_vcbwindow_mlp2/sall.json :: tasks.bc_vs_rest.arms.r16q1-s*.mlp.rejection_at['0.60'].n_bkg_pass :: sha256 5b8b5c7c28602f93 :: background jets passing the cut, mean over seeds
\newcommand{\VcbOmaRatioMlp}{1.50}  % experiments/FIGS/data/probe_ladder_vcbwindow_mlp2/sall.json :: tasks.bc_vs_rest.arms.*.mlp.log1m_auc :: sha256 5b8b5c7c28602f93 :: seed mean of 1-AUC, 17 classes over 162
\newcommand{\VcbFactorMlp}{1.50}  % experiments/FIGS/data/probe_ladder_vcbwindow_mlp2/sall.json :: tasks.bc_vs_rest.arms.*.mlp.log1m_auc :: sha256 5b8b5c7c28602f93 :: exp of the difference of the seed means of log(1-AUC), 17 classes over 162
\newcommand{\AccLegOneOnesixtwoSonebEFour}{0.32273}  % experiments/FIGS/data/leg1_metrics.json :: summary.l162-s1b.N10000.accuracy_mean :: sha256 3f564e960b2c86a7 :: wave 1, superseded; 3 fine-tuning seeds
\newcommand{\AccSdLegOneOnesixtwoSonebEFour}{0.01894}  % experiments/FIGS/data/leg1_metrics.json :: summary.l162-s1b.N10000.accuracy_sd :: sha256 3f564e960b2c86a7 :: wave 1, superseded
\newcommand{\AccLegOneOnesixtwoSonebEFive}{0.53506}  % experiments/FIGS/data/leg1_metrics.json :: summary.l162-s1b.N100000.accuracy_mean :: sha256 3f564e960b2c86a7 :: wave 1, superseded; 3 fine-tuning seeds
\newcommand{\AccSdLegOneOnesixtwoSonebEFive}{0.00044}  % experiments/FIGS/data/leg1_metrics.json :: summary.l162-s1b.N100000.accuracy_sd :: sha256 3f564e960b2c86a7 :: wave 1, superseded
\newcommand{\AccLegOneOnesixtwoSonebESix}{0.54337}  % experiments/FIGS/data/leg1_metrics.json :: summary.l162-s1b.N1000000.accuracy_mean :: sha256 3f564e960b2c86a7 :: wave 1, superseded; 3 fine-tuning seeds
\newcommand{\AccSdLegOneOnesixtwoSonebESix}{0.00074}  % experiments/FIGS/data/leg1_metrics.json :: summary.l162-s1b.N1000000.accuracy_sd :: sha256 3f564e960b2c86a7 :: wave 1, superseded
\newcommand{\AccLegOneOnesevenStwoEFour}{0.22684}  % experiments/FIGS/data/leg1_metrics.json :: summary.r16q1-s2.N10000.accuracy_mean :: sha256 3f564e960b2c86a7 :: wave 1, superseded; 3 fine-tuning seeds
\newcommand{\AccSdLegOneOnesevenStwoEFour}{0.00805}  % experiments/FIGS/data/leg1_metrics.json :: summary.r16q1-s2.N10000.accuracy_sd :: sha256 3f564e960b2c86a7 :: wave 1, superseded
\newcommand{\AccLegOneOnesevenStwoEFive}{0.49650}  % experiments/FIGS/data/leg1_metrics.json :: summary.r16q1-s2.N100000.accuracy_mean :: sha256 3f564e960b2c86a7 :: wave 1, superseded; 3 fine-tuning seeds
\newcommand{\AccSdLegOneOnesevenStwoEFive}{0.00110}  % experiments/FIGS/data/leg1_metrics.json :: summary.r16q1-s2.N100000.accuracy_sd :: sha256 3f564e960b2c86a7 :: wave 1, superseded
\newcommand{\AccLegOneOnesevenStwoESix}{0.51969}  % experiments/FIGS/data/leg1_metrics.json :: summary.r16q1-s2.N1000000.accuracy_mean :: sha256 3f564e960b2c86a7 :: wave 1, superseded; 3 fine-tuning seeds
\newcommand{\AccSdLegOneOnesevenStwoESix}{0.00072}  % experiments/FIGS/data/leg1_metrics.json :: summary.r16q1-s2.N1000000.accuracy_sd :: sha256 3f564e960b2c86a7 :: wave 1, superseded
\newcommand{\AccLegOneOnesevenSthreeEFour}{0.24203}  % experiments/FIGS/data/leg1_metrics.json :: summary.r16q1-s3.N10000.accuracy_mean :: sha256 3f564e960b2c86a7 :: wave 1, superseded; 3 fine-tuning seeds
\newcommand{\AccSdLegOneOnesevenSthreeEFour}{0.00213}  % experiments/FIGS/data/leg1_metrics.json :: summary.r16q1-s3.N10000.accuracy_sd :: sha256 3f564e960b2c86a7 :: wave 1, superseded
\newcommand{\AccLegOneOnesevenSthreeEFive}{0.49313}  % experiments/FIGS/data/leg1_metrics.json :: summary.r16q1-s3.N100000.accuracy_mean :: sha256 3f564e960b2c86a7 :: wave 1, superseded; 3 fine-tuning seeds
\newcommand{\AccSdLegOneOnesevenSthreeEFive}{0.00294}  % experiments/FIGS/data/leg1_metrics.json :: summary.r16q1-s3.N100000.accuracy_sd :: sha256 3f564e960b2c86a7 :: wave 1, superseded
\newcommand{\AccLegOneOnesevenSthreeESix}{0.51804}  % experiments/FIGS/data/leg1_metrics.json :: summary.r16q1-s3.N1000000.accuracy_mean :: sha256 3f564e960b2c86a7 :: wave 1, superseded; 3 fine-tuning seeds
\newcommand{\AccSdLegOneOnesevenSthreeESix}{0.00049}  % experiments/FIGS/data/leg1_metrics.json :: summary.r16q1-s3.N1000000.accuracy_sd :: sha256 3f564e960b2c86a7 :: wave 1, superseded
\newcommand{\AccLegOneOnesevenSfourEFour}{0.24487}  % experiments/FIGS/data/leg1_metrics.json :: summary.r16q1-s4.N10000.accuracy_mean :: sha256 3f564e960b2c86a7 :: wave 1, superseded; 3 fine-tuning seeds
\newcommand{\AccSdLegOneOnesevenSfourEFour}{0.00645}  % experiments/FIGS/data/leg1_metrics.json :: summary.r16q1-s4.N10000.accuracy_sd :: sha256 3f564e960b2c86a7 :: wave 1, superseded
\newcommand{\AccLegOneOnesevenSfourEFive}{0.49510}  % experiments/FIGS/data/leg1_metrics.json :: summary.r16q1-s4.N100000.accuracy_mean :: sha256 3f564e960b2c86a7 :: wave 1, superseded; 3 fine-tuning seeds
\newcommand{\AccSdLegOneOnesevenSfourEFive}{0.00133}  % experiments/FIGS/data/leg1_metrics.json :: summary.r16q1-s4.N100000.accuracy_sd :: sha256 3f564e960b2c86a7 :: wave 1, superseded
\newcommand{\AccLegOneOnesevenSfourESix}{0.51872}  % experiments/FIGS/data/leg1_metrics.json :: summary.r16q1-s4.N1000000.accuracy_mean :: sha256 3f564e960b2c86a7 :: wave 1, superseded; 3 fine-tuning seeds
\newcommand{\AccSdLegOneOnesevenSfourESix}{0.00049}  % experiments/FIGS/data/leg1_metrics.json :: summary.r16q1-s4.N1000000.accuracy_sd :: sha256 3f564e960b2c86a7 :: wave 1, superseded
\newcommand{\AccLegOneScratchEFour}{0.31369}  % experiments/FIGS/data/leg1_metrics.json :: summary.scratch.N10000.accuracy_mean :: sha256 3f564e960b2c86a7 :: wave 1, superseded; 3 fine-tuning seeds
\newcommand{\AccSdLegOneScratchEFour}{0.00712}  % experiments/FIGS/data/leg1_metrics.json :: summary.scratch.N10000.accuracy_sd :: sha256 3f564e960b2c86a7 :: wave 1, superseded
\newcommand{\AccLegOneScratchEFive}{0.42015}  % experiments/FIGS/data/leg1_metrics.json :: summary.scratch.N100000.accuracy_mean :: sha256 3f564e960b2c86a7 :: wave 1, superseded; 3 fine-tuning seeds
\newcommand{\AccSdLegOneScratchEFive}{0.00267}  % experiments/FIGS/data/leg1_metrics.json :: summary.scratch.N100000.accuracy_sd :: sha256 3f564e960b2c86a7 :: wave 1, superseded
\newcommand{\AccLegOneScratchESix}{0.48621}  % experiments/FIGS/data/leg1_metrics.json :: summary.scratch.N1000000.accuracy_mean :: sha256 3f564e960b2c86a7 :: wave 1, superseded; 3 fine-tuning seeds
\newcommand{\AccSdLegOneScratchESix}{0.00100}  % experiments/FIGS/data/leg1_metrics.json :: summary.scratch.N1000000.accuracy_sd :: sha256 3f564e960b2c86a7 :: wave 1, superseded
\newcommand{\AccLegOneSophonPublicEFour}{0.21613}  % experiments/FIGS/data/leg1_metrics.json :: summary.sophon-public.N10000.accuracy_mean :: sha256 3f564e960b2c86a7 :: wave 1, superseded; 3 fine-tuning seeds
\newcommand{\AccSdLegOneSophonPublicEFour}{0.00734}  % experiments/FIGS/data/leg1_metrics.json :: summary.sophon-public.N10000.accuracy_sd :: sha256 3f564e960b2c86a7 :: wave 1, superseded
\newcommand{\AccLegOneSophonPublicEFive}{0.53187}  % experiments/FIGS/data/leg1_metrics.json :: summary.sophon-public.N100000.accuracy_mean :: sha256 3f564e960b2c86a7 :: wave 1, superseded; 3 fine-tuning seeds
\newcommand{\AccSdLegOneSophonPublicEFive}{0.00231}  % experiments/FIGS/data/leg1_metrics.json :: summary.sophon-public.N100000.accuracy_sd :: sha256 3f564e960b2c86a7 :: wave 1, superseded
\newcommand{\AccLegOneSophonPublicESix}{0.54273}  % experiments/FIGS/data/leg1_metrics.json :: summary.sophon-public.N1000000.accuracy_mean :: sha256 3f564e960b2c86a7 :: wave 1, superseded; 3 fine-tuning seeds
\newcommand{\AccSdLegOneSophonPublicESix}{0.00064}  % experiments/FIGS/data/leg1_metrics.json :: summary.sophon-public.N1000000.accuracy_sd :: sha256 3f564e960b2c86a7 :: wave 1, superseded
\newcommand{\AccLegTwoOnesixtwoSonebEFour}{0.78606}  % experiments/FIGS/data/leg2_metrics.json :: summary.l162-s1b.N10000.accuracy_mean :: sha256 701c58c3240274bb :: wave 1, superseded; 3 fine-tuning seeds
\newcommand{\AccSdLegTwoOnesixtwoSonebEFour}{0.00111}  % experiments/FIGS/data/leg2_metrics.json :: summary.l162-s1b.N10000.accuracy_sd :: sha256 701c58c3240274bb :: wave 1, superseded
\newcommand{\AccLegTwoOnesixtwoSonebEFive}{0.83598}  % experiments/FIGS/data/leg2_metrics.json :: summary.l162-s1b.N100000.accuracy_mean :: sha256 701c58c3240274bb :: wave 1, superseded; 3 fine-tuning seeds
\newcommand{\AccSdLegTwoOnesixtwoSonebEFive}{0.00037}  % experiments/FIGS/data/leg2_metrics.json :: summary.l162-s1b.N100000.accuracy_sd :: sha256 701c58c3240274bb :: wave 1, superseded
\newcommand{\AccLegTwoOnesixtwoSonebESix}{0.84704}  % experiments/FIGS/data/leg2_metrics.json :: summary.l162-s1b.N1000000.accuracy_mean :: sha256 701c58c3240274bb :: wave 1, superseded; 3 fine-tuning seeds
\newcommand{\AccSdLegTwoOnesixtwoSonebESix}{0.00027}  % experiments/FIGS/data/leg2_metrics.json :: summary.l162-s1b.N1000000.accuracy_sd :: sha256 701c58c3240274bb :: wave 1, superseded
\newcommand{\AccLegTwoOnesevenStwoEFour}{0.72156}  % experiments/FIGS/data/leg2_metrics.json :: summary.r16q1-s2.N10000.accuracy_mean :: sha256 701c58c3240274bb :: wave 1, superseded; 3 fine-tuning seeds
\newcommand{\AccSdLegTwoOnesevenStwoEFour}{0.01049}  % experiments/FIGS/data/leg2_metrics.json :: summary.r16q1-s2.N10000.accuracy_sd :: sha256 701c58c3240274bb :: wave 1, superseded
\newcommand{\AccLegTwoOnesevenStwoEFive}{0.82614}  % experiments/FIGS/data/leg2_metrics.json :: summary.r16q1-s2.N100000.accuracy_mean :: sha256 701c58c3240274bb :: wave 1, superseded; 3 fine-tuning seeds
\newcommand{\AccSdLegTwoOnesevenStwoEFive}{0.00058}  % experiments/FIGS/data/leg2_metrics.json :: summary.r16q1-s2.N100000.accuracy_sd :: sha256 701c58c3240274bb :: wave 1, superseded
\newcommand{\AccLegTwoOnesevenStwoESix}{0.83943}  % experiments/FIGS/data/leg2_metrics.json :: summary.r16q1-s2.N1000000.accuracy_mean :: sha256 701c58c3240274bb :: wave 1, superseded; 3 fine-tuning seeds
\newcommand{\AccSdLegTwoOnesevenStwoESix}{0.00052}  % experiments/FIGS/data/leg2_metrics.json :: summary.r16q1-s2.N1000000.accuracy_sd :: sha256 701c58c3240274bb :: wave 1, superseded
\newcommand{\AccLegTwoOnesevenSthreeEFour}{0.72103}  % experiments/FIGS/data/leg2_metrics.json :: summary.r16q1-s3.N10000.accuracy_mean :: sha256 701c58c3240274bb :: wave 1, superseded; 3 fine-tuning seeds
\newcommand{\AccSdLegTwoOnesevenSthreeEFour}{0.01651}  % experiments/FIGS/data/leg2_metrics.json :: summary.r16q1-s3.N10000.accuracy_sd :: sha256 701c58c3240274bb :: wave 1, superseded
\newcommand{\AccLegTwoOnesevenSthreeEFive}{0.82629}  % experiments/FIGS/data/leg2_metrics.json :: summary.r16q1-s3.N100000.accuracy_mean :: sha256 701c58c3240274bb :: wave 1, superseded; 3 fine-tuning seeds
\newcommand{\AccSdLegTwoOnesevenSthreeEFive}{0.00035}  % experiments/FIGS/data/leg2_metrics.json :: summary.r16q1-s3.N100000.accuracy_sd :: sha256 701c58c3240274bb :: wave 1, superseded
\newcommand{\AccLegTwoOnesevenSthreeESix}{0.83944}  % experiments/FIGS/data/leg2_metrics.json :: summary.r16q1-s3.N1000000.accuracy_mean :: sha256 701c58c3240274bb :: wave 1, superseded; 3 fine-tuning seeds
\newcommand{\AccSdLegTwoOnesevenSthreeESix}{0.00019}  % experiments/FIGS/data/leg2_metrics.json :: summary.r16q1-s3.N1000000.accuracy_sd :: sha256 701c58c3240274bb :: wave 1, superseded
\newcommand{\AccLegTwoOnesevenSfourEFour}{0.72953}  % experiments/FIGS/data/leg2_metrics.json :: summary.r16q1-s4.N10000.accuracy_mean :: sha256 701c58c3240274bb :: wave 1, superseded; 3 fine-tuning seeds
\newcommand{\AccSdLegTwoOnesevenSfourEFour}{0.01069}  % experiments/FIGS/data/leg2_metrics.json :: summary.r16q1-s4.N10000.accuracy_sd :: sha256 701c58c3240274bb :: wave 1, superseded
\newcommand{\AccLegTwoOnesevenSfourEFive}{0.82564}  % experiments/FIGS/data/leg2_metrics.json :: summary.r16q1-s4.N100000.accuracy_mean :: sha256 701c58c3240274bb :: wave 1, superseded; 3 fine-tuning seeds
\newcommand{\AccSdLegTwoOnesevenSfourEFive}{0.00041}  % experiments/FIGS/data/leg2_metrics.json :: summary.r16q1-s4.N100000.accuracy_sd :: sha256 701c58c3240274bb :: wave 1, superseded
\newcommand{\AccLegTwoOnesevenSfourESix}{0.83909}  % experiments/FIGS/data/leg2_metrics.json :: summary.r16q1-s4.N1000000.accuracy_mean :: sha256 701c58c3240274bb :: wave 1, superseded; 3 fine-tuning seeds
\newcommand{\AccSdLegTwoOnesevenSfourESix}{0.00037}  % experiments/FIGS/data/leg2_metrics.json :: summary.r16q1-s4.N1000000.accuracy_sd :: sha256 701c58c3240274bb :: wave 1, superseded
\newcommand{\AccLegTwoRefEonearmsSoneRef}{0.81770}  % experiments/FIGS/data/leg2_metrics.json :: summary.ref_e1arms-s1.ref.accuracy_mean :: sha256 701c58c3240274bb :: wave 1, superseded; 1 fine-tuning seeds
\newcommand{\AccSdLegTwoRefEonearmsSoneRef}{---}  % experiments/FIGS/data/leg2_metrics.json :: summary.ref_e1arms-s1.ref.accuracy_sd :: sha256 701c58c3240274bb :: wave 1, superseded
\newcommand{\AccLegTwoRefEonearmsStwoRef}{0.81602}  % experiments/FIGS/data/leg2_metrics.json :: summary.ref_e1arms-s2.ref.accuracy_mean :: sha256 701c58c3240274bb :: wave 1, superseded; 1 fine-tuning seeds
\newcommand{\AccSdLegTwoRefEonearmsStwoRef}{---}  % experiments/FIGS/data/leg2_metrics.json :: summary.ref_e1arms-s2.ref.accuracy_sd :: sha256 701c58c3240274bb :: wave 1, superseded
\newcommand{\AccLegTwoRefEonearmsSthreeRef}{0.81625}  % experiments/FIGS/data/leg2_metrics.json :: summary.ref_e1arms-s3.ref.accuracy_mean :: sha256 701c58c3240274bb :: wave 1, superseded; 1 fine-tuning seeds
\newcommand{\AccSdLegTwoRefEonearmsSthreeRef}{---}  % experiments/FIGS/data/leg2_metrics.json :: summary.ref_e1arms-s3.ref.accuracy_sd :: sha256 701c58c3240274bb :: wave 1, superseded
\newcommand{\AccLegTwoScratchEFour}{0.68230}  % experiments/FIGS/data/leg2_metrics.json :: summary.scratch.N10000.accuracy_mean :: sha256 701c58c3240274bb :: wave 1, superseded; 3 fine-tuning seeds
\newcommand{\AccSdLegTwoScratchEFour}{0.00282}  % experiments/FIGS/data/leg2_metrics.json :: summary.scratch.N10000.accuracy_sd :: sha256 701c58c3240274bb :: wave 1, superseded
\newcommand{\AccLegTwoScratchEFive}{0.77769}  % experiments/FIGS/data/leg2_metrics.json :: summary.scratch.N100000.accuracy_mean :: sha256 701c58c3240274bb :: wave 1, superseded; 3 fine-tuning seeds
\newcommand{\AccSdLegTwoScratchEFive}{0.00208}  % experiments/FIGS/data/leg2_metrics.json :: summary.scratch.N100000.accuracy_sd :: sha256 701c58c3240274bb :: wave 1, superseded
\newcommand{\AccLegTwoScratchESix}{0.82303}  % experiments/FIGS/data/leg2_metrics.json :: summary.scratch.N1000000.accuracy_mean :: sha256 701c58c3240274bb :: wave 1, superseded; 3 fine-tuning seeds
\newcommand{\AccSdLegTwoScratchESix}{0.00134}  % experiments/FIGS/data/leg2_metrics.json :: summary.scratch.N1000000.accuracy_sd :: sha256 701c58c3240274bb :: wave 1, superseded
\newcommand{\AccLegTwoSophonPublicEFour}{0.72739}  % experiments/FIGS/data/leg2_metrics.json :: summary.sophon-public.N10000.accuracy_mean :: sha256 701c58c3240274bb :: wave 1, superseded; 3 fine-tuning seeds
\newcommand{\AccSdLegTwoSophonPublicEFour}{0.00470}  % experiments/FIGS/data/leg2_metrics.json :: summary.sophon-public.N10000.accuracy_sd :: sha256 701c58c3240274bb :: wave 1, superseded
\newcommand{\AccLegTwoSophonPublicEFive}{0.83492}  % experiments/FIGS/data/leg2_metrics.json :: summary.sophon-public.N100000.accuracy_mean :: sha256 701c58c3240274bb :: wave 1, superseded; 3 fine-tuning seeds
\newcommand{\AccSdLegTwoSophonPublicEFive}{0.00049}  % experiments/FIGS/data/leg2_metrics.json :: summary.sophon-public.N100000.accuracy_sd :: sha256 701c58c3240274bb :: wave 1, superseded
\newcommand{\AccLegTwoSophonPublicESix}{0.84632}  % experiments/FIGS/data/leg2_metrics.json :: summary.sophon-public.N1000000.accuracy_mean :: sha256 701c58c3240274bb :: wave 1, superseded; 3 fine-tuning seeds
\newcommand{\AccSdLegTwoSophonPublicESix}{0.00031}  % experiments/FIGS/data/leg2_metrics.json :: summary.sophon-public.N1000000.accuracy_sd :: sha256 701c58c3240274bb :: wave 1, superseded
\newcommand{\FtBudgetSteps}{20{,}000}  % scripts/build_ft_budget_jobs.py :: EPOCHS * STEPS :: sha256 f2ed9faa8e3bd622 :: the largest number of optimizer steps below the full training set
\newcommand{\FtLongSteps}{80{,}000}  % scripts/build_ft_budget_jobs.py :: LONG_EPOCHS * STEPS :: sha256 f2ed9faa8e3bd622 :: the largest number of optimizer steps of the long budget at LONG_N jets (JetClass, JetClass-II)
\newcommand{\DesignNTrainSelected}{109{,}938{,}036}  % experiments/FIGS/data/v2/jc2_train/counts.json :: n_selected :: sha256 b0b5a8b63aa808c0 :: selected jets of the pretrain_train files (experiments/FT/count_jc2_train.py)
\newcommand{\FtBudgetEval}{500}  % scripts/build_ft_budget_jobs.py :: STEPS :: sha256 f2ed9faa8e3bd622 :: steps between validation passes at and above SMALL_N training jets
\newcommand{\FtBudgetEvalSmall}{50}  % scripts/build_ft_budget_jobs.py :: SMALL_STEPS :: sha256 f2ed9faa8e3bd622 :: steps between validation passes below SMALL_N training jets
\newcommand{\FtBudgetEvalSmallBelow}{\ensuremath{10^{5}}}  % scripts/build_ft_budget_jobs.py :: SMALL_N :: sha256 f2ed9faa8e3bd622 :: training jets below which the validation runs every SMALL_STEPS steps
\newcommand{\FtBudgetPatience}{10}  % scripts/build_ft_budget_jobs.py :: PATIENCE :: sha256 f2ed9faa8e3bd622 :: validation passes without improvement before training stops
\newcommand{\FtHeadMult}{50}  % scripts/build_ft_budget_jobs.py :: HEAD_MULT :: sha256 f2ed9faa8e3bd622 :: the head's learning-rate multiplier
\newcommand{\FtHeadCap}{\ensuremath{5\times10^{-3}}}  % scripts/build_ft_budget_jobs.py :: HEAD_CAP :: sha256 f2ed9faa8e3bd622 :: the head's largest learning rate
\newcommand{\FtBudgetBatch}{512}  % scripts/build_ft_budget_jobs.py :: BATCH :: sha256 f2ed9faa8e3bd622
\newcommand{\FtBudgetWeightDecay}{0.01}  % scripts/build_ft_budget_jobs.py :: --optimizer-option weight_decay (every cell) :: sha256 f2ed9faa8e3bd622
\newcommand{\FtSweepPre}{\ensuremath{10^{-4}}, \ensuremath{3\times10^{-4}} and \ensuremath{10^{-3}}}  % scripts/build_ft_budget_jobs.py :: SWEEP_PRE :: sha256 f2ed9faa8e3bd622 :: backbone rates of the sweep
\newcommand{\FtSweepScratch}{\ensuremath{3\times10^{-4}}, \ensuremath{10^{-3}} and \ensuremath{3\times10^{-3}}}  % scripts/build_ft_budget_jobs.py :: SWEEP_SCRATCH :: sha256 f2ed9faa8e3bd622 :: from-scratch rates of the sweep
\newcommand{\FtSweepN}{10{,}000}  % scripts/build_ft_budget_jobs.py :: SWEEP_N :: sha256 f2ed9faa8e3bd622 :: training jets of the sweep
\newcommand{\FtValJets}{200{,}000}  % scripts/build_ft_budget_jobs.py :: VAL_N :: sha256 f2ed9faa8e3bd622 :: the fixed validation set
\newcommand{\FtValJetsBal}{20{,}480}  % scripts/build_ft_budget_jobs.py :: VAL_SIZE[jc2b] :: sha256 f2ed9faa8e3bd622 :: the balanced JetClass-II validation set
\newcommand{\LpVfourCMax}{\ensuremath{10^{4}}}  % experiments/EVAL/probe_bal.py :: C_GRID_V4 :: sha256 a7804216e3179d9d :: largest inverse regularisation of the v4 probes
\newcommand{\JcBalTestPerClass}{5{,}000}  % experiments/FT/stage_jc2_balanced.py :: TEST_PER_CLASS :: sha256 0260341d5ed75fec
\newcommand{\JcBalValPerClass}{109}  % experiments/FT/stage_jc2_balanced.py :: VAL_PER_CLASS :: sha256 0260341d5ed75fec
\newcommand{\JcBalPairTrainPerSide}{\ensuremath{10^{5}}}  % experiments/FT/stage_jc2_balanced.py :: PAIR_TRAIN_PER_SIDE :: sha256 0260341d5ed75fec
\newcommand{\JcBalPairTestPerSide}{\ensuremath{2\times10^{5}}}  % experiments/FT/stage_jc2_balanced.py :: PAIR_TEST_PER_SIDE :: sha256 0260341d5ed75fec
\newcommand{\JcBalNShort}{11}  % experiments/FIGS/data/v2/jc2bal/census.json :: counts.ft_test :: sha256 a4379ee0817254cf :: classes with fewer than 5000 selected test jets
\newcommand{\JcBalRarestTest}{224}  % experiments/FIGS/data/v2/jc2bal/census.json :: counts.ft_test (minimum) :: sha256 a4379ee0817254cf :: label_QCD_bcs
\newcommand{\JcBalShortTrainESix}{11}  % experiments/FIGS/data/v2/jc2bal/census.json :: counts.ft_heldout :: sha256 a4379ee0817254cf :: classes with fewer held-out jets than the 1e6 quota
\newcommand{\JcBalNTest}{903{,}922}  % experiments/FIGS/data/v2/jc2bal/manifest.json :: outputs[test.parquet].rows :: sha256 9c7775be6ff56ac3 :: jets of the balanced test set
\newcommand{\FtLrPre}{\ensuremath{10^{-3}}}  % experiments/FIGS/data/v2/budget/lrsweep_metrics.json :: argmax of the summed validation metric, pretrained :: sha256 7cdc148054ca902c
\newcommand{\FtLrScratchTop}{\ensuremath{10^{-3}}}  % experiments/FIGS/data/v2/budget/lrsweep_metrics.json :: argmax of the validation metric on top, from scratch :: sha256 7cdc148054ca902c :: agrees with LR of the launched job-ft-budget-bench, -jc1, -jc2b, -scratch-full and -long-scratch-full from-scratch cells
\newcommand{\FtLrScratchQg}{\ensuremath{3\times10^{-3}}}  % experiments/FIGS/data/v2/budget/lrsweep_metrics.json :: argmax of the validation metric on quark/gluon, from scratch :: sha256 7cdc148054ca902c :: agrees with LR of the launched job-ft-budget-bench and -scratch-full from-scratch cells
\newcommand{\MlstNParamsFt}{2{,}209{,}270}  % experiments/FIGS/mlst_v4.py :: our_params(2, 17) :: sha256 fa536b4866ffd651 :: trainable parameters of our fine-tuned network, two outputs (top, quark/gluon)
\newcommand{\MlstNParamsFtKin}{2{,}207{,}950}  % experiments/FIGS/mlst_v4.py :: our_params(2, 7) :: sha256 fa536b4866ffd651 :: trainable parameters of our fine-tuned network, two outputs, the 7 kinematic inputs
\newcommand{\MlstNParamsJc}{2{,}213{,}374}  % experiments/FIGS/mlst_v4.py :: our_params(10, 17) :: sha256 fa536b4866ffd651 :: trainable parameters of our fine-tuned network, the ten JetClass classes
\newcommand{\MlstProbeRunsDone}{27}  % experiments/FIGS/data/v2/index.json :: runs with a linear-probe fit (count) :: sha256 c6c4138ee5c86423 :: of the paper's runs
\newcommand{\MlstPaperRuns}{27}  % configs/arms/v2_paper_runs.txt :: runs (count) :: sha256 b04562b14ea8cbc7
\newcommand{\MlstSdRelUnc}{50\%}  % configs/arms/v2_paper_runs.txt :: relative standard error of a standard deviation from the runs of a hierarchy level, 1/sqrt(2(n-1)) :: sha256 b04562b14ea8cbc7
\newcommand{\MlstCmpJcTwoResOverFull}{1.026\,\ensuremath{\pm}\,0.005}  % experiments/FIGS/data/v2/index.json :: linear probe jc2_188 N=1000000: (1 - macro_auc_ovr) of L162 over L188, run by run :: sha256 c6c4138ee5c86423
\newcommand{\MlstCmpJcTwoResOverFullPr}{2.6\% more errors}  % experiments/FIGS/data/v2/index.json :: linear probe jc2_188 N=1000000: (1 - macro_auc_ovr) of L162 over L188, run by run :: sha256 c6c4138ee5c86423 :: prose twin: two significant figures, direction in words
\newcommand{\MlstCmpJcTwoMidOverFull}{1.227\,\ensuremath{\pm}\,0.009}  % experiments/FIGS/data/v2/index.json :: linear probe jc2_188 N=1000000: (1 - macro_auc_ovr) of R42_Q1 over L188, run by run :: sha256 c6c4138ee5c86423
\newcommand{\MlstCmpJcTwoMidOverFullPr}{23\% more errors}  % experiments/FIGS/data/v2/index.json :: linear probe jc2_188 N=1000000: (1 - macro_auc_ovr) of R42_Q1 over L188, run by run :: sha256 c6c4138ee5c86423 :: prose twin: two significant figures, direction in words
\newcommand{\MlstCmpJcTwoCoarseOverFull}{1.557\,\ensuremath{\pm}\,0.019}  % experiments/FIGS/data/v2/index.json :: linear probe jc2_188 N=1000000: (1 - macro_auc_ovr) of R16_Q1 over L188, run by run :: sha256 c6c4138ee5c86423
\newcommand{\MlstCmpJcTwoCoarseOverFullPr}{56\% more errors}  % experiments/FIGS/data/v2/index.json :: linear probe jc2_188 N=1000000: (1 - macro_auc_ovr) of R16_Q1 over L188, run by run :: sha256 c6c4138ee5c86423 :: prose twin: two significant figures, direction in words
\newcommand{\MlstCmpJcTwoResMassOverRes}{1.022 and 1.016}  % experiments/FIGS/data/v2/index.json :: linear probe jc2_188 N=1000000: (1 - macro_auc_ovr) of L162_MASS over L162, run by run :: sha256 c6c4138ee5c86423
\newcommand{\MlstCmpJcTwoResMassOverResPr}{equal within 3\%}  % experiments/FIGS/data/v2/index.json :: linear probe jc2_188 N=1000000: (1 - macro_auc_ovr) of L162_MASS over L162, run by run :: sha256 c6c4138ee5c86423 :: prose twin: two significant figures, direction in words
\newcommand{\MlstCmpJcTwoResMassOverResPrMax}{equal within 3\%}  % experiments/FIGS/data/v2/index.json :: linear probe jc2_188 N=1000000: (1 - macro_auc_ovr) of L162_MASS over L162, run by run :: sha256 c6c4138ee5c86423 :: prose twin of the run farther from 1
\newcommand{\MlstCmpJcTwoCoarseMassOverCoarse}{1.054 and 0.977}  % experiments/FIGS/data/v2/index.json :: linear probe jc2_188 N=1000000: (1 - macro_auc_ovr) of R16_Q1_MASS over R16_Q1, run by run :: sha256 c6c4138ee5c86423
\newcommand{\MlstCmpJcTwoCoarseMassOverCoarsePr}{equal within 6\%}  % experiments/FIGS/data/v2/index.json :: linear probe jc2_188 N=1000000: (1 - macro_auc_ovr) of R16_Q1_MASS over R16_Q1, run by run :: sha256 c6c4138ee5c86423 :: prose twin: two significant figures, direction in words
\newcommand{\MlstCmpJcTwoCoarseMassOverCoarsePrMax}{equal within 6\%}  % experiments/FIGS/data/v2/index.json :: linear probe jc2_188 N=1000000: (1 - macro_auc_ovr) of R16_Q1_MASS over R16_Q1, run by run :: sha256 c6c4138ee5c86423 :: prose twin of the run farther from 1
\newcommand{\MlstCmpJcTwoResMassOverFull}{1.053 and 1.037}  % experiments/FIGS/data/v2/index.json :: linear probe jc2_188 N=1000000: (1 - macro_auc_ovr) of L162_MASS over L188, run by run :: sha256 c6c4138ee5c86423
\newcommand{\MlstCmpJcTwoResMassOverFullPr}{at least 3.7\% more errors}  % experiments/FIGS/data/v2/index.json :: linear probe jc2_188 N=1000000: (1 - macro_auc_ovr) of L162_MASS over L188, run by run :: sha256 c6c4138ee5c86423 :: prose twin: two significant figures, direction in words
\newcommand{\MlstCmpJcTwoResMassOverFullPrMax}{up to 5.3\% more errors}  % experiments/FIGS/data/v2/index.json :: linear probe jc2_188 N=1000000: (1 - macro_auc_ovr) of L162_MASS over L188, run by run :: sha256 c6c4138ee5c86423 :: prose twin of the run farther from 1
\newcommand{\MlstCmpJcTwoCoarseMassOverFull}{1.617 and 1.528}  % experiments/FIGS/data/v2/index.json :: linear probe jc2_188 N=1000000: (1 - macro_auc_ovr) of R16_Q1_MASS over L188, run by run :: sha256 c6c4138ee5c86423
\newcommand{\MlstCmpJcTwoCoarseMassOverFullPr}{at least 53\% more errors}  % experiments/FIGS/data/v2/index.json :: linear probe jc2_188 N=1000000: (1 - macro_auc_ovr) of R16_Q1_MASS over L188, run by run :: sha256 c6c4138ee5c86423 :: prose twin: two significant figures, direction in words
\newcommand{\MlstCmpJcTwoCoarseMassOverFullPrMax}{up to 62\% more errors}  % experiments/FIGS/data/v2/index.json :: linear probe jc2_188 N=1000000: (1 - macro_auc_ovr) of R16_Q1_MASS over L188, run by run :: sha256 c6c4138ee5c86423 :: prose twin of the run farther from 1
\newcommand{\MlstCmpJcTwoFzeroOverFull}{1.424 and 1.462}  % experiments/FIGS/data/v2/index.json :: linear probe jc2_188 N=1000000: (1 - macro_auc_ovr) of FLAV_F0 over L188, run by run :: sha256 c6c4138ee5c86423
\newcommand{\MlstCmpJcTwoFzeroOverFullPr}{at least 42\% more errors}  % experiments/FIGS/data/v2/index.json :: linear probe jc2_188 N=1000000: (1 - macro_auc_ovr) of FLAV_F0 over L188, run by run :: sha256 c6c4138ee5c86423 :: prose twin: two significant figures, direction in words
\newcommand{\MlstCmpJcTwoFzeroOverFullPrMax}{up to 46\% more errors}  % experiments/FIGS/data/v2/index.json :: linear probe jc2_188 N=1000000: (1 - macro_auc_ovr) of FLAV_F0 over L188, run by run :: sha256 c6c4138ee5c86423 :: prose twin of the run farther from 1
\newcommand{\MlstCmpJcTwoFoneOverFull}{1.218 and 1.212}  % experiments/FIGS/data/v2/index.json :: linear probe jc2_188 N=1000000: (1 - macro_auc_ovr) of FLAV_F1 over L188, run by run :: sha256 c6c4138ee5c86423
\newcommand{\MlstCmpJcTwoFoneOverFullPr}{at least 21\% more errors}  % experiments/FIGS/data/v2/index.json :: linear probe jc2_188 N=1000000: (1 - macro_auc_ovr) of FLAV_F1 over L188, run by run :: sha256 c6c4138ee5c86423 :: prose twin: two significant figures, direction in words
\newcommand{\MlstCmpJcTwoFoneOverFullPrMax}{up to 22\% more errors}  % experiments/FIGS/data/v2/index.json :: linear probe jc2_188 N=1000000: (1 - macro_auc_ovr) of FLAV_F1 over L188, run by run :: sha256 c6c4138ee5c86423 :: prose twin of the run farther from 1
\newcommand{\MlstCmpJcTwoFonerOverFull}{1.206 and 1.208}  % experiments/FIGS/data/v2/index.json :: linear probe jc2_188 N=1000000: (1 - macro_auc_ovr) of FLAV_F1R over L188, run by run :: sha256 c6c4138ee5c86423
\newcommand{\MlstCmpJcTwoFonerOverFullPr}{at least 21\% more errors}  % experiments/FIGS/data/v2/index.json :: linear probe jc2_188 N=1000000: (1 - macro_auc_ovr) of FLAV_F1R over L188, run by run :: sha256 c6c4138ee5c86423 :: prose twin: two significant figures, direction in words
\newcommand{\MlstCmpJcTwoFonerOverFullPrMax}{up to 21\% more errors}  % experiments/FIGS/data/v2/index.json :: linear probe jc2_188 N=1000000: (1 - macro_auc_ovr) of FLAV_F1R over L188, run by run :: sha256 c6c4138ee5c86423 :: prose twin of the run farther from 1
\newcommand{\MlstCmpJcTwoFzeroOverCoarse}{0.9275 and 0.9345}  % experiments/FIGS/data/v2/index.json :: linear probe jc2_188 N=1000000: (1 - macro_auc_ovr) of FLAV_F0 over R16_Q1, run by run :: sha256 c6c4138ee5c86423
\newcommand{\MlstCmpJcTwoFzeroOverCoarsePr}{at least 6.5\% fewer errors}  % experiments/FIGS/data/v2/index.json :: linear probe jc2_188 N=1000000: (1 - macro_auc_ovr) of FLAV_F0 over R16_Q1, run by run :: sha256 c6c4138ee5c86423 :: prose twin: two significant figures, direction in words
\newcommand{\MlstCmpJcTwoFzeroOverCoarsePrMax}{up to 7.2\% fewer errors}  % experiments/FIGS/data/v2/index.json :: linear probe jc2_188 N=1000000: (1 - macro_auc_ovr) of FLAV_F0 over R16_Q1, run by run :: sha256 c6c4138ee5c86423 :: prose twin of the run farther from 1
\newcommand{\MlstCmpJcTwoFoneOverCoarse}{0.794 and 0.775}  % experiments/FIGS/data/v2/index.json :: linear probe jc2_188 N=1000000: (1 - macro_auc_ovr) of FLAV_F1 over R16_Q1, run by run :: sha256 c6c4138ee5c86423
\newcommand{\MlstCmpJcTwoFoneOverCoarsePr}{at least 21\% fewer errors}  % experiments/FIGS/data/v2/index.json :: linear probe jc2_188 N=1000000: (1 - macro_auc_ovr) of FLAV_F1 over R16_Q1, run by run :: sha256 c6c4138ee5c86423 :: prose twin: two significant figures, direction in words
\newcommand{\MlstCmpJcTwoFoneOverCoarsePrMax}{up to 23\% fewer errors}  % experiments/FIGS/data/v2/index.json :: linear probe jc2_188 N=1000000: (1 - macro_auc_ovr) of FLAV_F1 over R16_Q1, run by run :: sha256 c6c4138ee5c86423 :: prose twin of the run farther from 1
\newcommand{\MlstCmpJcTwoFonerOverCoarse}{0.786 and 0.772}  % experiments/FIGS/data/v2/index.json :: linear probe jc2_188 N=1000000: (1 - macro_auc_ovr) of FLAV_F1R over R16_Q1, run by run :: sha256 c6c4138ee5c86423
\newcommand{\MlstCmpJcTwoFonerOverCoarsePr}{at least 21\% fewer errors}  % experiments/FIGS/data/v2/index.json :: linear probe jc2_188 N=1000000: (1 - macro_auc_ovr) of FLAV_F1R over R16_Q1, run by run :: sha256 c6c4138ee5c86423 :: prose twin: two significant figures, direction in words
\newcommand{\MlstCmpJcTwoFonerOverCoarsePrMax}{up to 23\% fewer errors}  % experiments/FIGS/data/v2/index.json :: linear probe jc2_188 N=1000000: (1 - macro_auc_ovr) of FLAV_F1R over R16_Q1, run by run :: sha256 c6c4138ee5c86423 :: prose twin of the run farther from 1
\newcommand{\MlstCmpJcTwoFoneOverFzero}{0.856 and 0.829}  % experiments/FIGS/data/v2/index.json :: linear probe jc2_188 N=1000000: (1 - macro_auc_ovr) of FLAV_F1 over FLAV_F0, run by run :: sha256 c6c4138ee5c86423
\newcommand{\MlstCmpJcTwoFoneOverFzeroPr}{at least 14\% fewer errors}  % experiments/FIGS/data/v2/index.json :: linear probe jc2_188 N=1000000: (1 - macro_auc_ovr) of FLAV_F1 over FLAV_F0, run by run :: sha256 c6c4138ee5c86423 :: prose twin: two significant figures, direction in words
\newcommand{\MlstCmpJcTwoFoneOverFzeroPrMax}{up to 17\% fewer errors}  % experiments/FIGS/data/v2/index.json :: linear probe jc2_188 N=1000000: (1 - macro_auc_ovr) of FLAV_F1 over FLAV_F0, run by run :: sha256 c6c4138ee5c86423 :: prose twin of the run farther from 1
\newcommand{\MlstCmpJcTwoFonerOverFzero}{0.847 and 0.826}  % experiments/FIGS/data/v2/index.json :: linear probe jc2_188 N=1000000: (1 - macro_auc_ovr) of FLAV_F1R over FLAV_F0, run by run :: sha256 c6c4138ee5c86423
\newcommand{\MlstCmpJcTwoFonerOverFzeroPr}{at least 15\% fewer errors}  % experiments/FIGS/data/v2/index.json :: linear probe jc2_188 N=1000000: (1 - macro_auc_ovr) of FLAV_F1R over FLAV_F0, run by run :: sha256 c6c4138ee5c86423 :: prose twin: two significant figures, direction in words
\newcommand{\MlstCmpJcTwoFonerOverFzeroPrMax}{up to 17\% fewer errors}  % experiments/FIGS/data/v2/index.json :: linear probe jc2_188 N=1000000: (1 - macro_auc_ovr) of FLAV_F1R over FLAV_F0, run by run :: sha256 c6c4138ee5c86423 :: prose twin of the run farther from 1
\newcommand{\MlstCmpJcTwoFonerOverFone}{0.9899 and 0.9966}  % experiments/FIGS/data/v2/index.json :: linear probe jc2_188 N=1000000: (1 - macro_auc_ovr) of FLAV_F1R over FLAV_F1, run by run :: sha256 c6c4138ee5c86423
\newcommand{\MlstCmpJcTwoFonerOverFonePr}{equal within 2\%}  % experiments/FIGS/data/v2/index.json :: linear probe jc2_188 N=1000000: (1 - macro_auc_ovr) of FLAV_F1R over FLAV_F1, run by run :: sha256 c6c4138ee5c86423 :: prose twin: two significant figures, direction in words
\newcommand{\MlstCmpJcTwoFonerOverFonePrMax}{equal within 2\%}  % experiments/FIGS/data/v2/index.json :: linear probe jc2_188 N=1000000: (1 - macro_auc_ovr) of FLAV_F1R over FLAV_F1, run by run :: sha256 c6c4138ee5c86423 :: prose twin of the run farther from 1
\newcommand{\MlstCmpJcTwoMidOverRes}{1.197\,\ensuremath{\pm}\,0.013}  % experiments/FIGS/data/v2/index.json :: linear probe jc2_188 N=1000000: (1 - macro_auc_ovr) of R42_Q1 over L162, run by run :: sha256 c6c4138ee5c86423
\newcommand{\MlstCmpJcTwoMidOverResPr}{20\% more errors}  % experiments/FIGS/data/v2/index.json :: linear probe jc2_188 N=1000000: (1 - macro_auc_ovr) of R42_Q1 over L162, run by run :: sha256 c6c4138ee5c86423 :: prose twin: two significant figures, direction in words
\newcommand{\MlstCmpJcTwoCoarseOverRes}{1.518\,\ensuremath{\pm}\,0.025}  % experiments/FIGS/data/v2/index.json :: linear probe jc2_188 N=1000000: (1 - macro_auc_ovr) of R16_Q1 over L162, run by run :: sha256 c6c4138ee5c86423
\newcommand{\MlstCmpJcTwoCoarseOverResPr}{52\% more errors}  % experiments/FIGS/data/v2/index.json :: linear probe jc2_188 N=1000000: (1 - macro_auc_ovr) of R16_Q1 over L162, run by run :: sha256 c6c4138ee5c86423 :: prose twin: two significant figures, direction in words
\newcommand{\MlstCmpJcTwoCoarseOverMid}{1.269\,\ensuremath{\pm}\,0.007}  % experiments/FIGS/data/v2/index.json :: linear probe jc2_188 N=1000000: (1 - macro_auc_ovr) of R16_Q1 over R42_Q1, run by run :: sha256 c6c4138ee5c86423
\newcommand{\MlstCmpJcTwoCoarseOverMidPr}{27\% more errors}  % experiments/FIGS/data/v2/index.json :: linear probe jc2_188 N=1000000: (1 - macro_auc_ovr) of R16_Q1 over R42_Q1, run by run :: sha256 c6c4138ee5c86423 :: prose twin: two significant figures, direction in words
\newcommand{\MlstCmpJcTwoRandOverFull}{1.125\,\ensuremath{\pm}\,0.008}  % experiments/FIGS/data/v2/index.json :: linear probe jc2_188 N=1000000: (1 - macro_auc_ovr) of each random partition over the L188 run of its seed, mean +- SD over the partitions :: sha256 c6c4138ee5c86423
\newcommand{\MlstCmpJcTwoRandOverFullPr}{12\% more errors}  % experiments/FIGS/data/v2/index.json :: linear probe jc2_188 N=1000000: (1 - macro_auc_ovr) of each random partition over the L188 run of its seed, mean +- SD over the partitions :: sha256 c6c4138ee5c86423 :: prose twin: two significant figures, direction in words
\newcommand{\MlstCmpJcTwoRandOverMid}{0.918\,\ensuremath{\pm}\,0.007}  % experiments/FIGS/data/v2/index.json :: linear probe jc2_188 N=1000000: (1 - macro_auc_ovr) of each random partition over the R42_Q1 run of its seed, mean +- SD over the partitions :: sha256 c6c4138ee5c86423
\newcommand{\MlstCmpJcTwoRandOverMidPr}{8.2\% fewer errors}  % experiments/FIGS/data/v2/index.json :: linear probe jc2_188 N=1000000: (1 - macro_auc_ovr) of each random partition over the R42_Q1 run of its seed, mean +- SD over the partitions :: sha256 c6c4138ee5c86423 :: prose twin: two significant figures, direction in words
\newcommand{\MlstCmpJcTwoRandOverCoarse}{0.724\,\ensuremath{\pm}\,0.008}  % experiments/FIGS/data/v2/index.json :: linear probe jc2_188 N=1000000: (1 - macro_auc_ovr) of each random partition over the R16_Q1 run of its seed, mean +- SD over the partitions :: sha256 c6c4138ee5c86423
\newcommand{\MlstCmpJcTwoRandOverCoarsePr}{28\% fewer errors}  % experiments/FIGS/data/v2/index.json :: linear probe jc2_188 N=1000000: (1 - macro_auc_ovr) of each random partition over the R16_Q1 run of its seed, mean +- SD over the partitions :: sha256 c6c4138ee5c86423 :: prose twin: two significant figures, direction in words
\newcommand{\MlstCmpJcTwoTrainedResOverFull}{1.026\,\ensuremath{\pm}\,0.005}  % experiments/FIGS/data/v2/index.json :: linear probe jc2_188 N=1000000: (1 - macro_auc_trained) of L162 over L188, run by run :: sha256 c6c4138ee5c86423
\newcommand{\MlstCmpJcTwoTrainedResOverFullPr}{2.6\% more errors}  % experiments/FIGS/data/v2/index.json :: linear probe jc2_188 N=1000000: (1 - macro_auc_trained) of L162 over L188, run by run :: sha256 c6c4138ee5c86423 :: prose twin: two significant figures, direction in words
\newcommand{\MlstCmpJcTwoTrainedMidOverFull}{1.227\,\ensuremath{\pm}\,0.009}  % experiments/FIGS/data/v2/index.json :: linear probe jc2_188 N=1000000: (1 - macro_auc_trained) of R42_Q1 over L188, run by run :: sha256 c6c4138ee5c86423
\newcommand{\MlstCmpJcTwoTrainedMidOverFullPr}{23\% more errors}  % experiments/FIGS/data/v2/index.json :: linear probe jc2_188 N=1000000: (1 - macro_auc_trained) of R42_Q1 over L188, run by run :: sha256 c6c4138ee5c86423 :: prose twin: two significant figures, direction in words
\newcommand{\MlstCmpJcTwoTrainedCoarseOverFull}{1.557\,\ensuremath{\pm}\,0.019}  % experiments/FIGS/data/v2/index.json :: linear probe jc2_188 N=1000000: (1 - macro_auc_trained) of R16_Q1 over L188, run by run :: sha256 c6c4138ee5c86423
\newcommand{\MlstCmpJcTwoTrainedCoarseOverFullPr}{56\% more errors}  % experiments/FIGS/data/v2/index.json :: linear probe jc2_188 N=1000000: (1 - macro_auc_trained) of R16_Q1 over L188, run by run :: sha256 c6c4138ee5c86423 :: prose twin: two significant figures, direction in words
\newcommand{\MlstCmpJcTwoTrainedResMassOverRes}{1.022 and 1.016}  % experiments/FIGS/data/v2/index.json :: linear probe jc2_188 N=1000000: (1 - macro_auc_trained) of L162_MASS over L162, run by run :: sha256 c6c4138ee5c86423
\newcommand{\MlstCmpJcTwoTrainedResMassOverResPr}{equal within 3\%}  % experiments/FIGS/data/v2/index.json :: linear probe jc2_188 N=1000000: (1 - macro_auc_trained) of L162_MASS over L162, run by run :: sha256 c6c4138ee5c86423 :: prose twin: two significant figures, direction in words
\newcommand{\MlstCmpJcTwoTrainedResMassOverResPrMax}{equal within 3\%}  % experiments/FIGS/data/v2/index.json :: linear probe jc2_188 N=1000000: (1 - macro_auc_trained) of L162_MASS over L162, run by run :: sha256 c6c4138ee5c86423 :: prose twin of the run farther from 1
\newcommand{\MlstCmpJcTwoTrainedCoarseMassOverCoarse}{1.054 and 0.977}  % experiments/FIGS/data/v2/index.json :: linear probe jc2_188 N=1000000: (1 - macro_auc_trained) of R16_Q1_MASS over R16_Q1, run by run :: sha256 c6c4138ee5c86423
\newcommand{\MlstCmpJcTwoTrainedCoarseMassOverCoarsePr}{equal within 6\%}  % experiments/FIGS/data/v2/index.json :: linear probe jc2_188 N=1000000: (1 - macro_auc_trained) of R16_Q1_MASS over R16_Q1, run by run :: sha256 c6c4138ee5c86423 :: prose twin: two significant figures, direction in words
\newcommand{\MlstCmpJcTwoTrainedCoarseMassOverCoarsePrMax}{equal within 6\%}  % experiments/FIGS/data/v2/index.json :: linear probe jc2_188 N=1000000: (1 - macro_auc_trained) of R16_Q1_MASS over R16_Q1, run by run :: sha256 c6c4138ee5c86423 :: prose twin of the run farther from 1
\newcommand{\MlstCmpJcTwoTrainedResMassOverFull}{1.053 and 1.037}  % experiments/FIGS/data/v2/index.json :: linear probe jc2_188 N=1000000: (1 - macro_auc_trained) of L162_MASS over L188, run by run :: sha256 c6c4138ee5c86423
\newcommand{\MlstCmpJcTwoTrainedResMassOverFullPr}{at least 3.7\% more errors}  % experiments/FIGS/data/v2/index.json :: linear probe jc2_188 N=1000000: (1 - macro_auc_trained) of L162_MASS over L188, run by run :: sha256 c6c4138ee5c86423 :: prose twin: two significant figures, direction in words
\newcommand{\MlstCmpJcTwoTrainedResMassOverFullPrMax}{up to 5.3\% more errors}  % experiments/FIGS/data/v2/index.json :: linear probe jc2_188 N=1000000: (1 - macro_auc_trained) of L162_MASS over L188, run by run :: sha256 c6c4138ee5c86423 :: prose twin of the run farther from 1
\newcommand{\MlstCmpJcTwoTrainedCoarseMassOverFull}{1.617 and 1.528}  % experiments/FIGS/data/v2/index.json :: linear probe jc2_188 N=1000000: (1 - macro_auc_trained) of R16_Q1_MASS over L188, run by run :: sha256 c6c4138ee5c86423
\newcommand{\MlstCmpJcTwoTrainedCoarseMassOverFullPr}{at least 53\% more errors}  % experiments/FIGS/data/v2/index.json :: linear probe jc2_188 N=1000000: (1 - macro_auc_trained) of R16_Q1_MASS over L188, run by run :: sha256 c6c4138ee5c86423 :: prose twin: two significant figures, direction in words
\newcommand{\MlstCmpJcTwoTrainedCoarseMassOverFullPrMax}{up to 62\% more errors}  % experiments/FIGS/data/v2/index.json :: linear probe jc2_188 N=1000000: (1 - macro_auc_trained) of R16_Q1_MASS over L188, run by run :: sha256 c6c4138ee5c86423 :: prose twin of the run farther from 1
\newcommand{\MlstCmpJcTwoTrainedFzeroOverFull}{1.424 and 1.462}  % experiments/FIGS/data/v2/index.json :: linear probe jc2_188 N=1000000: (1 - macro_auc_trained) of FLAV_F0 over L188, run by run :: sha256 c6c4138ee5c86423
\newcommand{\MlstCmpJcTwoTrainedFzeroOverFullPr}{at least 42\% more errors}  % experiments/FIGS/data/v2/index.json :: linear probe jc2_188 N=1000000: (1 - macro_auc_trained) of FLAV_F0 over L188, run by run :: sha256 c6c4138ee5c86423 :: prose twin: two significant figures, direction in words
\newcommand{\MlstCmpJcTwoTrainedFzeroOverFullPrMax}{up to 46\% more errors}  % experiments/FIGS/data/v2/index.json :: linear probe jc2_188 N=1000000: (1 - macro_auc_trained) of FLAV_F0 over L188, run by run :: sha256 c6c4138ee5c86423 :: prose twin of the run farther from 1
\newcommand{\MlstCmpJcTwoTrainedFoneOverFull}{1.218 and 1.212}  % experiments/FIGS/data/v2/index.json :: linear probe jc2_188 N=1000000: (1 - macro_auc_trained) of FLAV_F1 over L188, run by run :: sha256 c6c4138ee5c86423
\newcommand{\MlstCmpJcTwoTrainedFoneOverFullPr}{at least 21\% more errors}  % experiments/FIGS/data/v2/index.json :: linear probe jc2_188 N=1000000: (1 - macro_auc_trained) of FLAV_F1 over L188, run by run :: sha256 c6c4138ee5c86423 :: prose twin: two significant figures, direction in words
\newcommand{\MlstCmpJcTwoTrainedFoneOverFullPrMax}{up to 22\% more errors}  % experiments/FIGS/data/v2/index.json :: linear probe jc2_188 N=1000000: (1 - macro_auc_trained) of FLAV_F1 over L188, run by run :: sha256 c6c4138ee5c86423 :: prose twin of the run farther from 1
\newcommand{\MlstCmpJcTwoTrainedFonerOverFull}{1.206 and 1.208}  % experiments/FIGS/data/v2/index.json :: linear probe jc2_188 N=1000000: (1 - macro_auc_trained) of FLAV_F1R over L188, run by run :: sha256 c6c4138ee5c86423
\newcommand{\MlstCmpJcTwoTrainedFonerOverFullPr}{at least 21\% more errors}  % experiments/FIGS/data/v2/index.json :: linear probe jc2_188 N=1000000: (1 - macro_auc_trained) of FLAV_F1R over L188, run by run :: sha256 c6c4138ee5c86423 :: prose twin: two significant figures, direction in words
\newcommand{\MlstCmpJcTwoTrainedFonerOverFullPrMax}{up to 21\% more errors}  % experiments/FIGS/data/v2/index.json :: linear probe jc2_188 N=1000000: (1 - macro_auc_trained) of FLAV_F1R over L188, run by run :: sha256 c6c4138ee5c86423 :: prose twin of the run farther from 1
\newcommand{\MlstCmpJcTwoTrainedFzeroOverCoarse}{0.9275 and 0.9345}  % experiments/FIGS/data/v2/index.json :: linear probe jc2_188 N=1000000: (1 - macro_auc_trained) of FLAV_F0 over R16_Q1, run by run :: sha256 c6c4138ee5c86423
\newcommand{\MlstCmpJcTwoTrainedFzeroOverCoarsePr}{at least 6.5\% fewer errors}  % experiments/FIGS/data/v2/index.json :: linear probe jc2_188 N=1000000: (1 - macro_auc_trained) of FLAV_F0 over R16_Q1, run by run :: sha256 c6c4138ee5c86423 :: prose twin: two significant figures, direction in words
\newcommand{\MlstCmpJcTwoTrainedFzeroOverCoarsePrMax}{up to 7.2\% fewer errors}  % experiments/FIGS/data/v2/index.json :: linear probe jc2_188 N=1000000: (1 - macro_auc_trained) of FLAV_F0 over R16_Q1, run by run :: sha256 c6c4138ee5c86423 :: prose twin of the run farther from 1
\newcommand{\MlstCmpJcTwoTrainedFoneOverCoarse}{0.794 and 0.775}  % experiments/FIGS/data/v2/index.json :: linear probe jc2_188 N=1000000: (1 - macro_auc_trained) of FLAV_F1 over R16_Q1, run by run :: sha256 c6c4138ee5c86423
\newcommand{\MlstCmpJcTwoTrainedFoneOverCoarsePr}{at least 21\% fewer errors}  % experiments/FIGS/data/v2/index.json :: linear probe jc2_188 N=1000000: (1 - macro_auc_trained) of FLAV_F1 over R16_Q1, run by run :: sha256 c6c4138ee5c86423 :: prose twin: two significant figures, direction in words
\newcommand{\MlstCmpJcTwoTrainedFoneOverCoarsePrMax}{up to 23\% fewer errors}  % experiments/FIGS/data/v2/index.json :: linear probe jc2_188 N=1000000: (1 - macro_auc_trained) of FLAV_F1 over R16_Q1, run by run :: sha256 c6c4138ee5c86423 :: prose twin of the run farther from 1
\newcommand{\MlstCmpJcTwoTrainedFonerOverCoarse}{0.786 and 0.772}  % experiments/FIGS/data/v2/index.json :: linear probe jc2_188 N=1000000: (1 - macro_auc_trained) of FLAV_F1R over R16_Q1, run by run :: sha256 c6c4138ee5c86423
\newcommand{\MlstCmpJcTwoTrainedFonerOverCoarsePr}{at least 21\% fewer errors}  % experiments/FIGS/data/v2/index.json :: linear probe jc2_188 N=1000000: (1 - macro_auc_trained) of FLAV_F1R over R16_Q1, run by run :: sha256 c6c4138ee5c86423 :: prose twin: two significant figures, direction in words
\newcommand{\MlstCmpJcTwoTrainedFonerOverCoarsePrMax}{up to 23\% fewer errors}  % experiments/FIGS/data/v2/index.json :: linear probe jc2_188 N=1000000: (1 - macro_auc_trained) of FLAV_F1R over R16_Q1, run by run :: sha256 c6c4138ee5c86423 :: prose twin of the run farther from 1
\newcommand{\MlstCmpJcTwoTrainedFoneOverFzero}{0.856 and 0.829}  % experiments/FIGS/data/v2/index.json :: linear probe jc2_188 N=1000000: (1 - macro_auc_trained) of FLAV_F1 over FLAV_F0, run by run :: sha256 c6c4138ee5c86423
\newcommand{\MlstCmpJcTwoTrainedFoneOverFzeroPr}{at least 14\% fewer errors}  % experiments/FIGS/data/v2/index.json :: linear probe jc2_188 N=1000000: (1 - macro_auc_trained) of FLAV_F1 over FLAV_F0, run by run :: sha256 c6c4138ee5c86423 :: prose twin: two significant figures, direction in words
\newcommand{\MlstCmpJcTwoTrainedFoneOverFzeroPrMax}{up to 17\% fewer errors}  % experiments/FIGS/data/v2/index.json :: linear probe jc2_188 N=1000000: (1 - macro_auc_trained) of FLAV_F1 over FLAV_F0, run by run :: sha256 c6c4138ee5c86423 :: prose twin of the run farther from 1
\newcommand{\MlstCmpJcTwoTrainedFonerOverFzero}{0.847 and 0.826}  % experiments/FIGS/data/v2/index.json :: linear probe jc2_188 N=1000000: (1 - macro_auc_trained) of FLAV_F1R over FLAV_F0, run by run :: sha256 c6c4138ee5c86423
\newcommand{\MlstCmpJcTwoTrainedFonerOverFzeroPr}{at least 15\% fewer errors}  % experiments/FIGS/data/v2/index.json :: linear probe jc2_188 N=1000000: (1 - macro_auc_trained) of FLAV_F1R over FLAV_F0, run by run :: sha256 c6c4138ee5c86423 :: prose twin: two significant figures, direction in words
\newcommand{\MlstCmpJcTwoTrainedFonerOverFzeroPrMax}{up to 17\% fewer errors}  % experiments/FIGS/data/v2/index.json :: linear probe jc2_188 N=1000000: (1 - macro_auc_trained) of FLAV_F1R over FLAV_F0, run by run :: sha256 c6c4138ee5c86423 :: prose twin of the run farther from 1
\newcommand{\MlstCmpJcTwoTrainedFonerOverFone}{0.9899 and 0.9966}  % experiments/FIGS/data/v2/index.json :: linear probe jc2_188 N=1000000: (1 - macro_auc_trained) of FLAV_F1R over FLAV_F1, run by run :: sha256 c6c4138ee5c86423
\newcommand{\MlstCmpJcTwoTrainedFonerOverFonePr}{equal within 2\%}  % experiments/FIGS/data/v2/index.json :: linear probe jc2_188 N=1000000: (1 - macro_auc_trained) of FLAV_F1R over FLAV_F1, run by run :: sha256 c6c4138ee5c86423 :: prose twin: two significant figures, direction in words
\newcommand{\MlstCmpJcTwoTrainedFonerOverFonePrMax}{equal within 2\%}  % experiments/FIGS/data/v2/index.json :: linear probe jc2_188 N=1000000: (1 - macro_auc_trained) of FLAV_F1R over FLAV_F1, run by run :: sha256 c6c4138ee5c86423 :: prose twin of the run farther from 1
\newcommand{\MlstCmpJcTwoTrainedMidOverRes}{1.197\,\ensuremath{\pm}\,0.013}  % experiments/FIGS/data/v2/index.json :: linear probe jc2_188 N=1000000: (1 - macro_auc_trained) of R42_Q1 over L162, run by run :: sha256 c6c4138ee5c86423
\newcommand{\MlstCmpJcTwoTrainedMidOverResPr}{20\% more errors}  % experiments/FIGS/data/v2/index.json :: linear probe jc2_188 N=1000000: (1 - macro_auc_trained) of R42_Q1 over L162, run by run :: sha256 c6c4138ee5c86423 :: prose twin: two significant figures, direction in words
\newcommand{\MlstCmpJcTwoTrainedCoarseOverRes}{1.518\,\ensuremath{\pm}\,0.025}  % experiments/FIGS/data/v2/index.json :: linear probe jc2_188 N=1000000: (1 - macro_auc_trained) of R16_Q1 over L162, run by run :: sha256 c6c4138ee5c86423
\newcommand{\MlstCmpJcTwoTrainedCoarseOverResPr}{52\% more errors}  % experiments/FIGS/data/v2/index.json :: linear probe jc2_188 N=1000000: (1 - macro_auc_trained) of R16_Q1 over L162, run by run :: sha256 c6c4138ee5c86423 :: prose twin: two significant figures, direction in words
\newcommand{\MlstCmpJcTwoTrainedCoarseOverMid}{1.269\,\ensuremath{\pm}\,0.007}  % experiments/FIGS/data/v2/index.json :: linear probe jc2_188 N=1000000: (1 - macro_auc_trained) of R16_Q1 over R42_Q1, run by run :: sha256 c6c4138ee5c86423
\newcommand{\MlstCmpJcTwoTrainedCoarseOverMidPr}{27\% more errors}  % experiments/FIGS/data/v2/index.json :: linear probe jc2_188 N=1000000: (1 - macro_auc_trained) of R16_Q1 over R42_Q1, run by run :: sha256 c6c4138ee5c86423 :: prose twin: two significant figures, direction in words
\newcommand{\MlstCmpJcTwoTrainedRandOverFull}{1.125\,\ensuremath{\pm}\,0.008}  % experiments/FIGS/data/v2/index.json :: linear probe jc2_188 N=1000000: (1 - macro_auc_trained) of each random partition over the L188 run of its seed, mean +- SD over the partitions :: sha256 c6c4138ee5c86423
\newcommand{\MlstCmpJcTwoTrainedRandOverFullPr}{12\% more errors}  % experiments/FIGS/data/v2/index.json :: linear probe jc2_188 N=1000000: (1 - macro_auc_trained) of each random partition over the L188 run of its seed, mean +- SD over the partitions :: sha256 c6c4138ee5c86423 :: prose twin: two significant figures, direction in words
\newcommand{\MlstCmpJcTwoTrainedRandOverMid}{0.918\,\ensuremath{\pm}\,0.007}  % experiments/FIGS/data/v2/index.json :: linear probe jc2_188 N=1000000: (1 - macro_auc_trained) of each random partition over the R42_Q1 run of its seed, mean +- SD over the partitions :: sha256 c6c4138ee5c86423
\newcommand{\MlstCmpJcTwoTrainedRandOverMidPr}{8.2\% fewer errors}  % experiments/FIGS/data/v2/index.json :: linear probe jc2_188 N=1000000: (1 - macro_auc_trained) of each random partition over the R42_Q1 run of its seed, mean +- SD over the partitions :: sha256 c6c4138ee5c86423 :: prose twin: two significant figures, direction in words
\newcommand{\MlstCmpJcTwoTrainedRandOverCoarse}{0.724\,\ensuremath{\pm}\,0.008}  % experiments/FIGS/data/v2/index.json :: linear probe jc2_188 N=1000000: (1 - macro_auc_trained) of each random partition over the R16_Q1 run of its seed, mean +- SD over the partitions :: sha256 c6c4138ee5c86423
\newcommand{\MlstCmpJcTwoTrainedRandOverCoarsePr}{28\% fewer errors}  % experiments/FIGS/data/v2/index.json :: linear probe jc2_188 N=1000000: (1 - macro_auc_trained) of each random partition over the R16_Q1 run of its seed, mean +- SD over the partitions :: sha256 c6c4138ee5c86423 :: prose twin: two significant figures, direction in words
\newcommand{\MlstCmpJcTwoResResOverFull}{1.003\,\ensuremath{\pm}\,0.005}  % experiments/FIGS/data/v2/index.json :: linear probe jc2_188 N=1000000: (1 - macro_auc_resonant) of L162 over L188, run by run :: sha256 c6c4138ee5c86423
\newcommand{\MlstCmpJcTwoResResOverFullPr}{equal within 1\%}  % experiments/FIGS/data/v2/index.json :: linear probe jc2_188 N=1000000: (1 - macro_auc_resonant) of L162 over L188, run by run :: sha256 c6c4138ee5c86423 :: prose twin: two significant figures, direction in words
\newcommand{\MlstCmpJcTwoResMidOverFull}{1.262\,\ensuremath{\pm}\,0.012}  % experiments/FIGS/data/v2/index.json :: linear probe jc2_188 N=1000000: (1 - macro_auc_resonant) of R42_Q1 over L188, run by run :: sha256 c6c4138ee5c86423
\newcommand{\MlstCmpJcTwoResMidOverFullPr}{26\% more errors}  % experiments/FIGS/data/v2/index.json :: linear probe jc2_188 N=1000000: (1 - macro_auc_resonant) of R42_Q1 over L188, run by run :: sha256 c6c4138ee5c86423 :: prose twin: two significant figures, direction in words
\newcommand{\MlstCmpJcTwoResCoarseOverFull}{1.590\,\ensuremath{\pm}\,0.032}  % experiments/FIGS/data/v2/index.json :: linear probe jc2_188 N=1000000: (1 - macro_auc_resonant) of R16_Q1 over L188, run by run :: sha256 c6c4138ee5c86423
\newcommand{\MlstCmpJcTwoResCoarseOverFullPr}{59\% more errors}  % experiments/FIGS/data/v2/index.json :: linear probe jc2_188 N=1000000: (1 - macro_auc_resonant) of R16_Q1 over L188, run by run :: sha256 c6c4138ee5c86423 :: prose twin: two significant figures, direction in words
\newcommand{\MlstCmpJcTwoResResMassOverRes}{1.025 and 1.020}  % experiments/FIGS/data/v2/index.json :: linear probe jc2_188 N=1000000: (1 - macro_auc_resonant) of L162_MASS over L162, run by run :: sha256 c6c4138ee5c86423
\newcommand{\MlstCmpJcTwoResResMassOverResPr}{equal within 3\%}  % experiments/FIGS/data/v2/index.json :: linear probe jc2_188 N=1000000: (1 - macro_auc_resonant) of L162_MASS over L162, run by run :: sha256 c6c4138ee5c86423 :: prose twin: two significant figures, direction in words
\newcommand{\MlstCmpJcTwoResResMassOverResPrMax}{equal within 3\%}  % experiments/FIGS/data/v2/index.json :: linear probe jc2_188 N=1000000: (1 - macro_auc_resonant) of L162_MASS over L162, run by run :: sha256 c6c4138ee5c86423 :: prose twin of the run farther from 1
\newcommand{\MlstCmpJcTwoResCoarseMassOverCoarse}{1.08 and 0.96}  % experiments/FIGS/data/v2/index.json :: linear probe jc2_188 N=1000000: (1 - macro_auc_resonant) of R16_Q1_MASS over R16_Q1, run by run :: sha256 c6c4138ee5c86423
\newcommand{\MlstCmpJcTwoResCoarseMassOverCoarsePr}{equal within 8\%}  % experiments/FIGS/data/v2/index.json :: linear probe jc2_188 N=1000000: (1 - macro_auc_resonant) of R16_Q1_MASS over R16_Q1, run by run :: sha256 c6c4138ee5c86423 :: prose twin: two significant figures, direction in words
\newcommand{\MlstCmpJcTwoResCoarseMassOverCoarsePrMax}{equal within 8\%}  % experiments/FIGS/data/v2/index.json :: linear probe jc2_188 N=1000000: (1 - macro_auc_resonant) of R16_Q1_MASS over R16_Q1, run by run :: sha256 c6c4138ee5c86423 :: prose twin of the run farther from 1
\newcommand{\MlstCmpJcTwoResResMassOverFull}{1.033 and 1.018}  % experiments/FIGS/data/v2/index.json :: linear probe jc2_188 N=1000000: (1 - macro_auc_resonant) of L162_MASS over L188, run by run :: sha256 c6c4138ee5c86423
\newcommand{\MlstCmpJcTwoResResMassOverFullPr}{equal within 4\%}  % experiments/FIGS/data/v2/index.json :: linear probe jc2_188 N=1000000: (1 - macro_auc_resonant) of L162_MASS over L188, run by run :: sha256 c6c4138ee5c86423 :: prose twin: two significant figures, direction in words
\newcommand{\MlstCmpJcTwoResResMassOverFullPrMax}{equal within 4\%}  % experiments/FIGS/data/v2/index.json :: linear probe jc2_188 N=1000000: (1 - macro_auc_resonant) of L162_MASS over L188, run by run :: sha256 c6c4138ee5c86423 :: prose twin of the run farther from 1
\newcommand{\MlstCmpJcTwoResCoarseMassOverFull}{1.67 and 1.55}  % experiments/FIGS/data/v2/index.json :: linear probe jc2_188 N=1000000: (1 - macro_auc_resonant) of R16_Q1_MASS over L188, run by run :: sha256 c6c4138ee5c86423
\newcommand{\MlstCmpJcTwoResCoarseMassOverFullPr}{at least 55\% more errors}  % experiments/FIGS/data/v2/index.json :: linear probe jc2_188 N=1000000: (1 - macro_auc_resonant) of R16_Q1_MASS over L188, run by run :: sha256 c6c4138ee5c86423 :: prose twin: two significant figures, direction in words
\newcommand{\MlstCmpJcTwoResCoarseMassOverFullPrMax}{up to 67\% more errors}  % experiments/FIGS/data/v2/index.json :: linear probe jc2_188 N=1000000: (1 - macro_auc_resonant) of R16_Q1_MASS over L188, run by run :: sha256 c6c4138ee5c86423 :: prose twin of the run farther from 1
\newcommand{\MlstCmpJcTwoResFzeroOverFull}{1.408 and 1.449}  % experiments/FIGS/data/v2/index.json :: linear probe jc2_188 N=1000000: (1 - macro_auc_resonant) of FLAV_F0 over L188, run by run :: sha256 c6c4138ee5c86423
\newcommand{\MlstCmpJcTwoResFzeroOverFullPr}{at least 41\% more errors}  % experiments/FIGS/data/v2/index.json :: linear probe jc2_188 N=1000000: (1 - macro_auc_resonant) of FLAV_F0 over L188, run by run :: sha256 c6c4138ee5c86423 :: prose twin: two significant figures, direction in words
\newcommand{\MlstCmpJcTwoResFzeroOverFullPrMax}{up to 45\% more errors}  % experiments/FIGS/data/v2/index.json :: linear probe jc2_188 N=1000000: (1 - macro_auc_resonant) of FLAV_F0 over L188, run by run :: sha256 c6c4138ee5c86423 :: prose twin of the run farther from 1
\newcommand{\MlstCmpJcTwoResFoneOverFull}{1.210 and 1.204}  % experiments/FIGS/data/v2/index.json :: linear probe jc2_188 N=1000000: (1 - macro_auc_resonant) of FLAV_F1 over L188, run by run :: sha256 c6c4138ee5c86423
\newcommand{\MlstCmpJcTwoResFoneOverFullPr}{at least 20\% more errors}  % experiments/FIGS/data/v2/index.json :: linear probe jc2_188 N=1000000: (1 - macro_auc_resonant) of FLAV_F1 over L188, run by run :: sha256 c6c4138ee5c86423 :: prose twin: two significant figures, direction in words
\newcommand{\MlstCmpJcTwoResFoneOverFullPrMax}{up to 21\% more errors}  % experiments/FIGS/data/v2/index.json :: linear probe jc2_188 N=1000000: (1 - macro_auc_resonant) of FLAV_F1 over L188, run by run :: sha256 c6c4138ee5c86423 :: prose twin of the run farther from 1
\newcommand{\MlstCmpJcTwoResFonerOverFull}{1.187 and 1.188}  % experiments/FIGS/data/v2/index.json :: linear probe jc2_188 N=1000000: (1 - macro_auc_resonant) of FLAV_F1R over L188, run by run :: sha256 c6c4138ee5c86423
\newcommand{\MlstCmpJcTwoResFonerOverFullPr}{at least 19\% more errors}  % experiments/FIGS/data/v2/index.json :: linear probe jc2_188 N=1000000: (1 - macro_auc_resonant) of FLAV_F1R over L188, run by run :: sha256 c6c4138ee5c86423 :: prose twin: two significant figures, direction in words
\newcommand{\MlstCmpJcTwoResFonerOverFullPrMax}{up to 19\% more errors}  % experiments/FIGS/data/v2/index.json :: linear probe jc2_188 N=1000000: (1 - macro_auc_resonant) of FLAV_F1R over L188, run by run :: sha256 c6c4138ee5c86423 :: prose twin of the run farther from 1
\newcommand{\MlstCmpJcTwoResFzeroOverCoarse}{0.9069 and 0.9029}  % experiments/FIGS/data/v2/index.json :: linear probe jc2_188 N=1000000: (1 - macro_auc_resonant) of FLAV_F0 over R16_Q1, run by run :: sha256 c6c4138ee5c86423
\newcommand{\MlstCmpJcTwoResFzeroOverCoarsePr}{at least 9.3\% fewer errors}  % experiments/FIGS/data/v2/index.json :: linear probe jc2_188 N=1000000: (1 - macro_auc_resonant) of FLAV_F0 over R16_Q1, run by run :: sha256 c6c4138ee5c86423 :: prose twin: two significant figures, direction in words
\newcommand{\MlstCmpJcTwoResFzeroOverCoarsePrMax}{up to 9.7\% fewer errors}  % experiments/FIGS/data/v2/index.json :: linear probe jc2_188 N=1000000: (1 - macro_auc_resonant) of FLAV_F0 over R16_Q1, run by run :: sha256 c6c4138ee5c86423 :: prose twin of the run farther from 1
\newcommand{\MlstCmpJcTwoResFoneOverCoarse}{0.779 and 0.750}  % experiments/FIGS/data/v2/index.json :: linear probe jc2_188 N=1000000: (1 - macro_auc_resonant) of FLAV_F1 over R16_Q1, run by run :: sha256 c6c4138ee5c86423
\newcommand{\MlstCmpJcTwoResFoneOverCoarsePr}{at least 22\% fewer errors}  % experiments/FIGS/data/v2/index.json :: linear probe jc2_188 N=1000000: (1 - macro_auc_resonant) of FLAV_F1 over R16_Q1, run by run :: sha256 c6c4138ee5c86423 :: prose twin: two significant figures, direction in words
\newcommand{\MlstCmpJcTwoResFoneOverCoarsePrMax}{up to 25\% fewer errors}  % experiments/FIGS/data/v2/index.json :: linear probe jc2_188 N=1000000: (1 - macro_auc_resonant) of FLAV_F1 over R16_Q1, run by run :: sha256 c6c4138ee5c86423 :: prose twin of the run farther from 1
\newcommand{\MlstCmpJcTwoResFonerOverCoarse}{0.764 and 0.740}  % experiments/FIGS/data/v2/index.json :: linear probe jc2_188 N=1000000: (1 - macro_auc_resonant) of FLAV_F1R over R16_Q1, run by run :: sha256 c6c4138ee5c86423
\newcommand{\MlstCmpJcTwoResFonerOverCoarsePr}{at least 24\% fewer errors}  % experiments/FIGS/data/v2/index.json :: linear probe jc2_188 N=1000000: (1 - macro_auc_resonant) of FLAV_F1R over R16_Q1, run by run :: sha256 c6c4138ee5c86423 :: prose twin: two significant figures, direction in words
\newcommand{\MlstCmpJcTwoResFonerOverCoarsePrMax}{up to 26\% fewer errors}  % experiments/FIGS/data/v2/index.json :: linear probe jc2_188 N=1000000: (1 - macro_auc_resonant) of FLAV_F1R over R16_Q1, run by run :: sha256 c6c4138ee5c86423 :: prose twin of the run farther from 1
\newcommand{\MlstCmpJcTwoResFoneOverFzero}{0.859 and 0.831}  % experiments/FIGS/data/v2/index.json :: linear probe jc2_188 N=1000000: (1 - macro_auc_resonant) of FLAV_F1 over FLAV_F0, run by run :: sha256 c6c4138ee5c86423
\newcommand{\MlstCmpJcTwoResFoneOverFzeroPr}{at least 14\% fewer errors}  % experiments/FIGS/data/v2/index.json :: linear probe jc2_188 N=1000000: (1 - macro_auc_resonant) of FLAV_F1 over FLAV_F0, run by run :: sha256 c6c4138ee5c86423 :: prose twin: two significant figures, direction in words
\newcommand{\MlstCmpJcTwoResFoneOverFzeroPrMax}{up to 17\% fewer errors}  % experiments/FIGS/data/v2/index.json :: linear probe jc2_188 N=1000000: (1 - macro_auc_resonant) of FLAV_F1 over FLAV_F0, run by run :: sha256 c6c4138ee5c86423 :: prose twin of the run farther from 1
\newcommand{\MlstCmpJcTwoResFonerOverFzero}{0.843 and 0.820}  % experiments/FIGS/data/v2/index.json :: linear probe jc2_188 N=1000000: (1 - macro_auc_resonant) of FLAV_F1R over FLAV_F0, run by run :: sha256 c6c4138ee5c86423
\newcommand{\MlstCmpJcTwoResFonerOverFzeroPr}{at least 16\% fewer errors}  % experiments/FIGS/data/v2/index.json :: linear probe jc2_188 N=1000000: (1 - macro_auc_resonant) of FLAV_F1R over FLAV_F0, run by run :: sha256 c6c4138ee5c86423 :: prose twin: two significant figures, direction in words
\newcommand{\MlstCmpJcTwoResFonerOverFzeroPrMax}{up to 18\% fewer errors}  % experiments/FIGS/data/v2/index.json :: linear probe jc2_188 N=1000000: (1 - macro_auc_resonant) of FLAV_F1R over FLAV_F0, run by run :: sha256 c6c4138ee5c86423 :: prose twin of the run farther from 1
\newcommand{\MlstCmpJcTwoResFonerOverFone}{0.9806 and 0.9867}  % experiments/FIGS/data/v2/index.json :: linear probe jc2_188 N=1000000: (1 - macro_auc_resonant) of FLAV_F1R over FLAV_F1, run by run :: sha256 c6c4138ee5c86423
\newcommand{\MlstCmpJcTwoResFonerOverFonePr}{equal within 2\%}  % experiments/FIGS/data/v2/index.json :: linear probe jc2_188 N=1000000: (1 - macro_auc_resonant) of FLAV_F1R over FLAV_F1, run by run :: sha256 c6c4138ee5c86423 :: prose twin: two significant figures, direction in words
\newcommand{\MlstCmpJcTwoResFonerOverFonePrMax}{equal within 2\%}  % experiments/FIGS/data/v2/index.json :: linear probe jc2_188 N=1000000: (1 - macro_auc_resonant) of FLAV_F1R over FLAV_F1, run by run :: sha256 c6c4138ee5c86423 :: prose twin of the run farther from 1
\newcommand{\MlstCmpJcTwoResMidOverRes}{1.258\,\ensuremath{\pm}\,0.017}  % experiments/FIGS/data/v2/index.json :: linear probe jc2_188 N=1000000: (1 - macro_auc_resonant) of R42_Q1 over L162, run by run :: sha256 c6c4138ee5c86423
\newcommand{\MlstCmpJcTwoResMidOverResPr}{26\% more errors}  % experiments/FIGS/data/v2/index.json :: linear probe jc2_188 N=1000000: (1 - macro_auc_resonant) of R42_Q1 over L162, run by run :: sha256 c6c4138ee5c86423 :: prose twin: two significant figures, direction in words
\newcommand{\MlstCmpJcTwoResCoarseOverRes}{1.58\,\ensuremath{\pm}\,0.04}  % experiments/FIGS/data/v2/index.json :: linear probe jc2_188 N=1000000: (1 - macro_auc_resonant) of R16_Q1 over L162, run by run :: sha256 c6c4138ee5c86423
\newcommand{\MlstCmpJcTwoResCoarseOverResPr}{58\% more errors}  % experiments/FIGS/data/v2/index.json :: linear probe jc2_188 N=1000000: (1 - macro_auc_resonant) of R16_Q1 over L162, run by run :: sha256 c6c4138ee5c86423 :: prose twin: two significant figures, direction in words
\newcommand{\MlstCmpJcTwoResCoarseOverMid}{1.260\,\ensuremath{\pm}\,0.015}  % experiments/FIGS/data/v2/index.json :: linear probe jc2_188 N=1000000: (1 - macro_auc_resonant) of R16_Q1 over R42_Q1, run by run :: sha256 c6c4138ee5c86423
\newcommand{\MlstCmpJcTwoResCoarseOverMidPr}{26\% more errors}  % experiments/FIGS/data/v2/index.json :: linear probe jc2_188 N=1000000: (1 - macro_auc_resonant) of R16_Q1 over R42_Q1, run by run :: sha256 c6c4138ee5c86423 :: prose twin: two significant figures, direction in words
\newcommand{\MlstCmpJcTwoResRandOverFull}{1.098\,\ensuremath{\pm}\,0.008}  % experiments/FIGS/data/v2/index.json :: linear probe jc2_188 N=1000000: (1 - macro_auc_resonant) of each random partition over the L188 run of its seed, mean +- SD over the partitions :: sha256 c6c4138ee5c86423
\newcommand{\MlstCmpJcTwoResRandOverFullPr}{9.8\% more errors}  % experiments/FIGS/data/v2/index.json :: linear probe jc2_188 N=1000000: (1 - macro_auc_resonant) of each random partition over the L188 run of its seed, mean +- SD over the partitions :: sha256 c6c4138ee5c86423 :: prose twin: two significant figures, direction in words
\newcommand{\MlstCmpJcTwoResRandOverMid}{0.870\,\ensuremath{\pm}\,0.008}  % experiments/FIGS/data/v2/index.json :: linear probe jc2_188 N=1000000: (1 - macro_auc_resonant) of each random partition over the R42_Q1 run of its seed, mean +- SD over the partitions :: sha256 c6c4138ee5c86423
\newcommand{\MlstCmpJcTwoResRandOverMidPr}{13\% fewer errors}  % experiments/FIGS/data/v2/index.json :: linear probe jc2_188 N=1000000: (1 - macro_auc_resonant) of each random partition over the R42_Q1 run of its seed, mean +- SD over the partitions :: sha256 c6c4138ee5c86423 :: prose twin: two significant figures, direction in words
\newcommand{\MlstCmpJcTwoResRandOverCoarse}{0.693\,\ensuremath{\pm}\,0.011}  % experiments/FIGS/data/v2/index.json :: linear probe jc2_188 N=1000000: (1 - macro_auc_resonant) of each random partition over the R16_Q1 run of its seed, mean +- SD over the partitions :: sha256 c6c4138ee5c86423
\newcommand{\MlstCmpJcTwoResRandOverCoarsePr}{31\% fewer errors}  % experiments/FIGS/data/v2/index.json :: linear probe jc2_188 N=1000000: (1 - macro_auc_resonant) of each random partition over the R16_Q1 run of its seed, mean +- SD over the partitions :: sha256 c6c4138ee5c86423 :: prose twin: two significant figures, direction in words
\newcommand{\MlstCmpJcTwoQcdResOverFull}{1.087\,\ensuremath{\pm}\,0.006}  % experiments/FIGS/data/v2/index.json :: linear probe jc2_188 N=1000000: (1 - macro_auc_qcd) of L162 over L188, run by run :: sha256 c6c4138ee5c86423
\newcommand{\MlstCmpJcTwoQcdResOverFullPr}{8.7\% more errors}  % experiments/FIGS/data/v2/index.json :: linear probe jc2_188 N=1000000: (1 - macro_auc_qcd) of L162 over L188, run by run :: sha256 c6c4138ee5c86423 :: prose twin: two significant figures, direction in words
\newcommand{\MlstCmpJcTwoQcdMidOverFull}{1.131\,\ensuremath{\pm}\,0.007}  % experiments/FIGS/data/v2/index.json :: linear probe jc2_188 N=1000000: (1 - macro_auc_qcd) of R42_Q1 over L188, run by run :: sha256 c6c4138ee5c86423
\newcommand{\MlstCmpJcTwoQcdMidOverFullPr}{13\% more errors}  % experiments/FIGS/data/v2/index.json :: linear probe jc2_188 N=1000000: (1 - macro_auc_qcd) of R42_Q1 over L188, run by run :: sha256 c6c4138ee5c86423 :: prose twin: two significant figures, direction in words
\newcommand{\MlstCmpJcTwoQcdCoarseOverFull}{1.466\,\ensuremath{\pm}\,0.017}  % experiments/FIGS/data/v2/index.json :: linear probe jc2_188 N=1000000: (1 - macro_auc_qcd) of R16_Q1 over L188, run by run :: sha256 c6c4138ee5c86423
\newcommand{\MlstCmpJcTwoQcdCoarseOverFullPr}{47\% more errors}  % experiments/FIGS/data/v2/index.json :: linear probe jc2_188 N=1000000: (1 - macro_auc_qcd) of R16_Q1 over L188, run by run :: sha256 c6c4138ee5c86423 :: prose twin: two significant figures, direction in words
\newcommand{\MlstCmpJcTwoQcdResMassOverRes}{1.012 and 1.006}  % experiments/FIGS/data/v2/index.json :: linear probe jc2_188 N=1000000: (1 - macro_auc_qcd) of L162_MASS over L162, run by run :: sha256 c6c4138ee5c86423
\newcommand{\MlstCmpJcTwoQcdResMassOverResPr}{equal within 2\%}  % experiments/FIGS/data/v2/index.json :: linear probe jc2_188 N=1000000: (1 - macro_auc_qcd) of L162_MASS over L162, run by run :: sha256 c6c4138ee5c86423 :: prose twin: two significant figures, direction in words
\newcommand{\MlstCmpJcTwoQcdResMassOverResPrMax}{equal within 2\%}  % experiments/FIGS/data/v2/index.json :: linear probe jc2_188 N=1000000: (1 - macro_auc_qcd) of L162_MASS over L162, run by run :: sha256 c6c4138ee5c86423 :: prose twin of the run farther from 1
\newcommand{\MlstCmpJcTwoQcdCoarseMassOverCoarse}{0.991 and 1.014}  % experiments/FIGS/data/v2/index.json :: linear probe jc2_188 N=1000000: (1 - macro_auc_qcd) of R16_Q1_MASS over R16_Q1, run by run :: sha256 c6c4138ee5c86423
\newcommand{\MlstCmpJcTwoQcdCoarseMassOverCoarsePr}{equal within 2\%}  % experiments/FIGS/data/v2/index.json :: linear probe jc2_188 N=1000000: (1 - macro_auc_qcd) of R16_Q1_MASS over R16_Q1, run by run :: sha256 c6c4138ee5c86423 :: prose twin: two significant figures, direction in words
\newcommand{\MlstCmpJcTwoQcdCoarseMassOverCoarsePrMax}{equal within 2\%}  % experiments/FIGS/data/v2/index.json :: linear probe jc2_188 N=1000000: (1 - macro_auc_qcd) of R16_Q1_MASS over R16_Q1, run by run :: sha256 c6c4138ee5c86423 :: prose twin of the run farther from 1
\newcommand{\MlstCmpJcTwoQcdResMassOverFull}{1.107 and 1.092}  % experiments/FIGS/data/v2/index.json :: linear probe jc2_188 N=1000000: (1 - macro_auc_qcd) of L162_MASS over L188, run by run :: sha256 c6c4138ee5c86423
\newcommand{\MlstCmpJcTwoQcdResMassOverFullPr}{at least 9.2\% more errors}  % experiments/FIGS/data/v2/index.json :: linear probe jc2_188 N=1000000: (1 - macro_auc_qcd) of L162_MASS over L188, run by run :: sha256 c6c4138ee5c86423 :: prose twin: two significant figures, direction in words
\newcommand{\MlstCmpJcTwoQcdResMassOverFullPrMax}{up to 11\% more errors}  % experiments/FIGS/data/v2/index.json :: linear probe jc2_188 N=1000000: (1 - macro_auc_qcd) of L162_MASS over L188, run by run :: sha256 c6c4138ee5c86423 :: prose twin of the run farther from 1
\newcommand{\MlstCmpJcTwoQcdCoarseMassOverFull}{1.472 and 1.474}  % experiments/FIGS/data/v2/index.json :: linear probe jc2_188 N=1000000: (1 - macro_auc_qcd) of R16_Q1_MASS over L188, run by run :: sha256 c6c4138ee5c86423
\newcommand{\MlstCmpJcTwoQcdCoarseMassOverFullPr}{at least 47\% more errors}  % experiments/FIGS/data/v2/index.json :: linear probe jc2_188 N=1000000: (1 - macro_auc_qcd) of R16_Q1_MASS over L188, run by run :: sha256 c6c4138ee5c86423 :: prose twin: two significant figures, direction in words
\newcommand{\MlstCmpJcTwoQcdCoarseMassOverFullPrMax}{up to 47\% more errors}  % experiments/FIGS/data/v2/index.json :: linear probe jc2_188 N=1000000: (1 - macro_auc_qcd) of R16_Q1_MASS over L188, run by run :: sha256 c6c4138ee5c86423 :: prose twin of the run farther from 1
\newcommand{\MlstCmpJcTwoQcdFzeroOverFull}{1.466 and 1.497}  % experiments/FIGS/data/v2/index.json :: linear probe jc2_188 N=1000000: (1 - macro_auc_qcd) of FLAV_F0 over L188, run by run :: sha256 c6c4138ee5c86423
\newcommand{\MlstCmpJcTwoQcdFzeroOverFullPr}{at least 47\% more errors}  % experiments/FIGS/data/v2/index.json :: linear probe jc2_188 N=1000000: (1 - macro_auc_qcd) of FLAV_F0 over L188, run by run :: sha256 c6c4138ee5c86423 :: prose twin: two significant figures, direction in words
\newcommand{\MlstCmpJcTwoQcdFzeroOverFullPrMax}{up to 50\% more errors}  % experiments/FIGS/data/v2/index.json :: linear probe jc2_188 N=1000000: (1 - macro_auc_qcd) of FLAV_F0 over L188, run by run :: sha256 c6c4138ee5c86423 :: prose twin of the run farther from 1
\newcommand{\MlstCmpJcTwoQcdFoneOverFull}{1.241 and 1.235}  % experiments/FIGS/data/v2/index.json :: linear probe jc2_188 N=1000000: (1 - macro_auc_qcd) of FLAV_F1 over L188, run by run :: sha256 c6c4138ee5c86423
\newcommand{\MlstCmpJcTwoQcdFoneOverFullPr}{at least 23\% more errors}  % experiments/FIGS/data/v2/index.json :: linear probe jc2_188 N=1000000: (1 - macro_auc_qcd) of FLAV_F1 over L188, run by run :: sha256 c6c4138ee5c86423 :: prose twin: two significant figures, direction in words
\newcommand{\MlstCmpJcTwoQcdFoneOverFullPrMax}{up to 24\% more errors}  % experiments/FIGS/data/v2/index.json :: linear probe jc2_188 N=1000000: (1 - macro_auc_qcd) of FLAV_F1 over L188, run by run :: sha256 c6c4138ee5c86423 :: prose twin of the run farther from 1
\newcommand{\MlstCmpJcTwoQcdFonerOverFull}{1.260 and 1.263}  % experiments/FIGS/data/v2/index.json :: linear probe jc2_188 N=1000000: (1 - macro_auc_qcd) of FLAV_F1R over L188, run by run :: sha256 c6c4138ee5c86423
\newcommand{\MlstCmpJcTwoQcdFonerOverFullPr}{at least 26\% more errors}  % experiments/FIGS/data/v2/index.json :: linear probe jc2_188 N=1000000: (1 - macro_auc_qcd) of FLAV_F1R over L188, run by run :: sha256 c6c4138ee5c86423 :: prose twin: two significant figures, direction in words
\newcommand{\MlstCmpJcTwoQcdFonerOverFullPrMax}{up to 26\% more errors}  % experiments/FIGS/data/v2/index.json :: linear probe jc2_188 N=1000000: (1 - macro_auc_qcd) of FLAV_F1R over L188, run by run :: sha256 c6c4138ee5c86423 :: prose twin of the run farther from 1
\newcommand{\MlstCmpJcTwoQcdFzeroOverCoarse}{0.987 and 1.031}  % experiments/FIGS/data/v2/index.json :: linear probe jc2_188 N=1000000: (1 - macro_auc_qcd) of FLAV_F0 over R16_Q1, run by run :: sha256 c6c4138ee5c86423
\newcommand{\MlstCmpJcTwoQcdFzeroOverCoarsePr}{equal within 4\%}  % experiments/FIGS/data/v2/index.json :: linear probe jc2_188 N=1000000: (1 - macro_auc_qcd) of FLAV_F0 over R16_Q1, run by run :: sha256 c6c4138ee5c86423 :: prose twin: two significant figures, direction in words
\newcommand{\MlstCmpJcTwoQcdFzeroOverCoarsePrMax}{equal within 4\%}  % experiments/FIGS/data/v2/index.json :: linear probe jc2_188 N=1000000: (1 - macro_auc_qcd) of FLAV_F0 over R16_Q1, run by run :: sha256 c6c4138ee5c86423 :: prose twin of the run farther from 1
\newcommand{\MlstCmpJcTwoQcdFoneOverCoarse}{0.836 and 0.850}  % experiments/FIGS/data/v2/index.json :: linear probe jc2_188 N=1000000: (1 - macro_auc_qcd) of FLAV_F1 over R16_Q1, run by run :: sha256 c6c4138ee5c86423
\newcommand{\MlstCmpJcTwoQcdFoneOverCoarsePr}{at least 15\% fewer errors}  % experiments/FIGS/data/v2/index.json :: linear probe jc2_188 N=1000000: (1 - macro_auc_qcd) of FLAV_F1 over R16_Q1, run by run :: sha256 c6c4138ee5c86423 :: prose twin: two significant figures, direction in words
\newcommand{\MlstCmpJcTwoQcdFoneOverCoarsePrMax}{up to 16\% fewer errors}  % experiments/FIGS/data/v2/index.json :: linear probe jc2_188 N=1000000: (1 - macro_auc_qcd) of FLAV_F1 over R16_Q1, run by run :: sha256 c6c4138ee5c86423 :: prose twin of the run farther from 1
\newcommand{\MlstCmpJcTwoQcdFonerOverCoarse}{0.848 and 0.870}  % experiments/FIGS/data/v2/index.json :: linear probe jc2_188 N=1000000: (1 - macro_auc_qcd) of FLAV_F1R over R16_Q1, run by run :: sha256 c6c4138ee5c86423
\newcommand{\MlstCmpJcTwoQcdFonerOverCoarsePr}{at least 13\% fewer errors}  % experiments/FIGS/data/v2/index.json :: linear probe jc2_188 N=1000000: (1 - macro_auc_qcd) of FLAV_F1R over R16_Q1, run by run :: sha256 c6c4138ee5c86423 :: prose twin: two significant figures, direction in words
\newcommand{\MlstCmpJcTwoQcdFonerOverCoarsePrMax}{up to 15\% fewer errors}  % experiments/FIGS/data/v2/index.json :: linear probe jc2_188 N=1000000: (1 - macro_auc_qcd) of FLAV_F1R over R16_Q1, run by run :: sha256 c6c4138ee5c86423 :: prose twin of the run farther from 1
\newcommand{\MlstCmpJcTwoQcdFoneOverFzero}{0.847 and 0.825}  % experiments/FIGS/data/v2/index.json :: linear probe jc2_188 N=1000000: (1 - macro_auc_qcd) of FLAV_F1 over FLAV_F0, run by run :: sha256 c6c4138ee5c86423
\newcommand{\MlstCmpJcTwoQcdFoneOverFzeroPr}{at least 15\% fewer errors}  % experiments/FIGS/data/v2/index.json :: linear probe jc2_188 N=1000000: (1 - macro_auc_qcd) of FLAV_F1 over FLAV_F0, run by run :: sha256 c6c4138ee5c86423 :: prose twin: two significant figures, direction in words
\newcommand{\MlstCmpJcTwoQcdFoneOverFzeroPrMax}{up to 18\% fewer errors}  % experiments/FIGS/data/v2/index.json :: linear probe jc2_188 N=1000000: (1 - macro_auc_qcd) of FLAV_F1 over FLAV_F0, run by run :: sha256 c6c4138ee5c86423 :: prose twin of the run farther from 1
\newcommand{\MlstCmpJcTwoQcdFonerOverFzero}{0.859 and 0.844}  % experiments/FIGS/data/v2/index.json :: linear probe jc2_188 N=1000000: (1 - macro_auc_qcd) of FLAV_F1R over FLAV_F0, run by run :: sha256 c6c4138ee5c86423
\newcommand{\MlstCmpJcTwoQcdFonerOverFzeroPr}{at least 14\% fewer errors}  % experiments/FIGS/data/v2/index.json :: linear probe jc2_188 N=1000000: (1 - macro_auc_qcd) of FLAV_F1R over FLAV_F0, run by run :: sha256 c6c4138ee5c86423 :: prose twin: two significant figures, direction in words
\newcommand{\MlstCmpJcTwoQcdFonerOverFzeroPrMax}{up to 16\% fewer errors}  % experiments/FIGS/data/v2/index.json :: linear probe jc2_188 N=1000000: (1 - macro_auc_qcd) of FLAV_F1R over FLAV_F0, run by run :: sha256 c6c4138ee5c86423 :: prose twin of the run farther from 1
\newcommand{\MlstCmpJcTwoQcdFonerOverFone}{1.015 and 1.023}  % experiments/FIGS/data/v2/index.json :: linear probe jc2_188 N=1000000: (1 - macro_auc_qcd) of FLAV_F1R over FLAV_F1, run by run :: sha256 c6c4138ee5c86423
\newcommand{\MlstCmpJcTwoQcdFonerOverFonePr}{equal within 3\%}  % experiments/FIGS/data/v2/index.json :: linear probe jc2_188 N=1000000: (1 - macro_auc_qcd) of FLAV_F1R over FLAV_F1, run by run :: sha256 c6c4138ee5c86423 :: prose twin: two significant figures, direction in words
\newcommand{\MlstCmpJcTwoQcdFonerOverFonePrMax}{equal within 3\%}  % experiments/FIGS/data/v2/index.json :: linear probe jc2_188 N=1000000: (1 - macro_auc_qcd) of FLAV_F1R over FLAV_F1, run by run :: sha256 c6c4138ee5c86423 :: prose twin of the run farther from 1
\newcommand{\MlstCmpJcTwoQcdMidOverRes}{1.040\,\ensuremath{\pm}\,0.011}  % experiments/FIGS/data/v2/index.json :: linear probe jc2_188 N=1000000: (1 - macro_auc_qcd) of R42_Q1 over L162, run by run :: sha256 c6c4138ee5c86423
\newcommand{\MlstCmpJcTwoQcdMidOverResPr}{4.0\% more errors}  % experiments/FIGS/data/v2/index.json :: linear probe jc2_188 N=1000000: (1 - macro_auc_qcd) of R42_Q1 over L162, run by run :: sha256 c6c4138ee5c86423 :: prose twin: two significant figures, direction in words
\newcommand{\MlstCmpJcTwoQcdCoarseOverRes}{1.348\,\ensuremath{\pm}\,0.010}  % experiments/FIGS/data/v2/index.json :: linear probe jc2_188 N=1000000: (1 - macro_auc_qcd) of R16_Q1 over L162, run by run :: sha256 c6c4138ee5c86423
\newcommand{\MlstCmpJcTwoQcdCoarseOverResPr}{35\% more errors}  % experiments/FIGS/data/v2/index.json :: linear probe jc2_188 N=1000000: (1 - macro_auc_qcd) of R16_Q1 over L162, run by run :: sha256 c6c4138ee5c86423 :: prose twin: two significant figures, direction in words
\newcommand{\MlstCmpJcTwoQcdCoarseOverMid}{1.296\,\ensuremath{\pm}\,0.016}  % experiments/FIGS/data/v2/index.json :: linear probe jc2_188 N=1000000: (1 - macro_auc_qcd) of R16_Q1 over R42_Q1, run by run :: sha256 c6c4138ee5c86423
\newcommand{\MlstCmpJcTwoQcdCoarseOverMidPr}{30\% more errors}  % experiments/FIGS/data/v2/index.json :: linear probe jc2_188 N=1000000: (1 - macro_auc_qcd) of R16_Q1 over R42_Q1, run by run :: sha256 c6c4138ee5c86423 :: prose twin: two significant figures, direction in words
\newcommand{\MlstCmpJcTwoQcdRandOverFull}{1.198\,\ensuremath{\pm}\,0.011}  % experiments/FIGS/data/v2/index.json :: linear probe jc2_188 N=1000000: (1 - macro_auc_qcd) of each random partition over the L188 run of its seed, mean +- SD over the partitions :: sha256 c6c4138ee5c86423
\newcommand{\MlstCmpJcTwoQcdRandOverFullPr}{20\% more errors}  % experiments/FIGS/data/v2/index.json :: linear probe jc2_188 N=1000000: (1 - macro_auc_qcd) of each random partition over the L188 run of its seed, mean +- SD over the partitions :: sha256 c6c4138ee5c86423 :: prose twin: two significant figures, direction in words
\newcommand{\MlstCmpJcTwoQcdRandOverMid}{1.064\,\ensuremath{\pm}\,0.011}  % experiments/FIGS/data/v2/index.json :: linear probe jc2_188 N=1000000: (1 - macro_auc_qcd) of each random partition over the R42_Q1 run of its seed, mean +- SD over the partitions :: sha256 c6c4138ee5c86423
\newcommand{\MlstCmpJcTwoQcdRandOverMidPr}{6.4\% more errors}  % experiments/FIGS/data/v2/index.json :: linear probe jc2_188 N=1000000: (1 - macro_auc_qcd) of each random partition over the R42_Q1 run of its seed, mean +- SD over the partitions :: sha256 c6c4138ee5c86423 :: prose twin: two significant figures, direction in words
\newcommand{\MlstCmpJcTwoQcdRandOverCoarse}{0.818\,\ensuremath{\pm}\,0.013}  % experiments/FIGS/data/v2/index.json :: linear probe jc2_188 N=1000000: (1 - macro_auc_qcd) of each random partition over the R16_Q1 run of its seed, mean +- SD over the partitions :: sha256 c6c4138ee5c86423
\newcommand{\MlstCmpJcTwoQcdRandOverCoarsePr}{18\% fewer errors}  % experiments/FIGS/data/v2/index.json :: linear probe jc2_188 N=1000000: (1 - macro_auc_qcd) of each random partition over the R16_Q1 run of its seed, mean +- SD over the partitions :: sha256 c6c4138ee5c86423 :: prose twin: two significant figures, direction in words
\newcommand{\MlstCmpBcResOverFull}{0.998\,\ensuremath{\pm}\,0.007}  % experiments/FIGS/data/v2/index.json :: binary probe bc_vs_bq_cs: (1 - AUC) of L162 over L188, run by run :: sha256 c6c4138ee5c86423
\newcommand{\MlstCmpBcResOverFullPr}{equal within 1\%}  % experiments/FIGS/data/v2/index.json :: binary probe bc_vs_bq_cs: (1 - AUC) of L162 over L188, run by run :: sha256 c6c4138ee5c86423 :: prose twin: two significant figures, direction in words
\newcommand{\MlstCmpBcMidOverFull}{1.151\,\ensuremath{\pm}\,0.029}  % experiments/FIGS/data/v2/index.json :: binary probe bc_vs_bq_cs: (1 - AUC) of R42_Q1 over L188, run by run :: sha256 c6c4138ee5c86423
\newcommand{\MlstCmpBcMidOverFullPr}{15\% more errors}  % experiments/FIGS/data/v2/index.json :: binary probe bc_vs_bq_cs: (1 - AUC) of R42_Q1 over L188, run by run :: sha256 c6c4138ee5c86423 :: prose twin: two significant figures, direction in words
\newcommand{\MlstCmpBcCoarseOverFull}{1.95\,\ensuremath{\pm}\,0.06}  % experiments/FIGS/data/v2/index.json :: binary probe bc_vs_bq_cs: (1 - AUC) of R16_Q1 over L188, run by run :: sha256 c6c4138ee5c86423
\newcommand{\MlstCmpBcCoarseOverFullPr}{95\% more errors}  % experiments/FIGS/data/v2/index.json :: binary probe bc_vs_bq_cs: (1 - AUC) of R16_Q1 over L188, run by run :: sha256 c6c4138ee5c86423 :: prose twin: two significant figures, direction in words
\newcommand{\MlstCmpBcFoneOverFzero}{0.759 and 0.698}  % experiments/FIGS/data/v2/index.json :: binary probe bc_vs_bq_cs: (1 - AUC) of FLAV_F1 over FLAV_F0, run by run :: sha256 c6c4138ee5c86423
\newcommand{\MlstCmpBcFoneOverFzeroPr}{at least 24\% fewer errors}  % experiments/FIGS/data/v2/index.json :: binary probe bc_vs_bq_cs: (1 - AUC) of FLAV_F1 over FLAV_F0, run by run :: sha256 c6c4138ee5c86423 :: prose twin: two significant figures, direction in words
\newcommand{\MlstCmpBcFoneOverFzeroPrMax}{up to 30\% fewer errors}  % experiments/FIGS/data/v2/index.json :: binary probe bc_vs_bq_cs: (1 - AUC) of FLAV_F1 over FLAV_F0, run by run :: sha256 c6c4138ee5c86423 :: prose twin of the run farther from 1
\newcommand{\MlstCmpBcFonerOverFone}{1.008 and 0.990}  % experiments/FIGS/data/v2/index.json :: binary probe bc_vs_bq_cs: (1 - AUC) of FLAV_F1R over FLAV_F1, run by run :: sha256 c6c4138ee5c86423
\newcommand{\MlstCmpBcFonerOverFonePr}{equal within 1\%}  % experiments/FIGS/data/v2/index.json :: binary probe bc_vs_bq_cs: (1 - AUC) of FLAV_F1R over FLAV_F1, run by run :: sha256 c6c4138ee5c86423 :: prose twin: two significant figures, direction in words
\newcommand{\MlstCmpBcFonerOverFonePrMax}{equal within 1\%}  % experiments/FIGS/data/v2/index.json :: binary probe bc_vs_bq_cs: (1 - AUC) of FLAV_F1R over FLAV_F1, run by run :: sha256 c6c4138ee5c86423 :: prose twin of the run farther from 1
\newcommand{\MlstCmpBcFonerOverFzero}{0.765 and 0.691}  % experiments/FIGS/data/v2/index.json :: binary probe bc_vs_bq_cs: (1 - AUC) of FLAV_F1R over FLAV_F0, run by run :: sha256 c6c4138ee5c86423
\newcommand{\MlstCmpBcFonerOverFzeroPr}{at least 23\% fewer errors}  % experiments/FIGS/data/v2/index.json :: binary probe bc_vs_bq_cs: (1 - AUC) of FLAV_F1R over FLAV_F0, run by run :: sha256 c6c4138ee5c86423 :: prose twin: two significant figures, direction in words
\newcommand{\MlstCmpBcFonerOverFzeroPrMax}{up to 31\% fewer errors}  % experiments/FIGS/data/v2/index.json :: binary probe bc_vs_bq_cs: (1 - AUC) of FLAV_F1R over FLAV_F0, run by run :: sha256 c6c4138ee5c86423 :: prose twin of the run farther from 1
\newcommand{\MlstCmpBcFzeroOverCoarse}{0.99 and 1.13}  % experiments/FIGS/data/v2/index.json :: binary probe bc_vs_bq_cs: (1 - AUC) of FLAV_F0 over R16_Q1, run by run :: sha256 c6c4138ee5c86423
\newcommand{\MlstCmpBcFzeroOverCoarsePr}{equal within 20\%}  % experiments/FIGS/data/v2/index.json :: binary probe bc_vs_bq_cs: (1 - AUC) of FLAV_F0 over R16_Q1, run by run :: sha256 c6c4138ee5c86423 :: prose twin: two significant figures, direction in words
\newcommand{\MlstCmpBcFzeroOverCoarsePrMax}{equal within 20\%}  % experiments/FIGS/data/v2/index.json :: binary probe bc_vs_bq_cs: (1 - AUC) of FLAV_F0 over R16_Q1, run by run :: sha256 c6c4138ee5c86423 :: prose twin of the run farther from 1
\newcommand{\MlstCmpBcFoneOverCoarse}{0.752 and 0.786}  % experiments/FIGS/data/v2/index.json :: binary probe bc_vs_bq_cs: (1 - AUC) of FLAV_F1 over R16_Q1, run by run :: sha256 c6c4138ee5c86423
\newcommand{\MlstCmpBcFoneOverCoarsePr}{at least 21\% fewer errors}  % experiments/FIGS/data/v2/index.json :: binary probe bc_vs_bq_cs: (1 - AUC) of FLAV_F1 over R16_Q1, run by run :: sha256 c6c4138ee5c86423 :: prose twin: two significant figures, direction in words
\newcommand{\MlstCmpBcFoneOverCoarsePrMax}{up to 25\% fewer errors}  % experiments/FIGS/data/v2/index.json :: binary probe bc_vs_bq_cs: (1 - AUC) of FLAV_F1 over R16_Q1, run by run :: sha256 c6c4138ee5c86423 :: prose twin of the run farther from 1
\newcommand{\MlstCmpBcFonerOverCoarse}{0.758 and 0.778}  % experiments/FIGS/data/v2/index.json :: binary probe bc_vs_bq_cs: (1 - AUC) of FLAV_F1R over R16_Q1, run by run :: sha256 c6c4138ee5c86423
\newcommand{\MlstCmpBcFonerOverCoarsePr}{at least 22\% fewer errors}  % experiments/FIGS/data/v2/index.json :: binary probe bc_vs_bq_cs: (1 - AUC) of FLAV_F1R over R16_Q1, run by run :: sha256 c6c4138ee5c86423 :: prose twin: two significant figures, direction in words
\newcommand{\MlstCmpBcFonerOverCoarsePrMax}{up to 24\% fewer errors}  % experiments/FIGS/data/v2/index.json :: binary probe bc_vs_bq_cs: (1 - AUC) of FLAV_F1R over R16_Q1, run by run :: sha256 c6c4138ee5c86423 :: prose twin of the run farther from 1
\newcommand{\MlstCmpBcFzeroOverFull}{2.00 and 2.19}  % experiments/FIGS/data/v2/index.json :: binary probe bc_vs_bq_cs: (1 - AUC) of FLAV_F0 over L188, run by run :: sha256 c6c4138ee5c86423
\newcommand{\MlstCmpBcFzeroOverFullPr}{at least 2.0 times the errors}  % experiments/FIGS/data/v2/index.json :: binary probe bc_vs_bq_cs: (1 - AUC) of FLAV_F0 over L188, run by run :: sha256 c6c4138ee5c86423 :: prose twin: two significant figures, direction in words
\newcommand{\MlstCmpBcFzeroOverFullPrMax}{up to 2.2 times the errors}  % experiments/FIGS/data/v2/index.json :: binary probe bc_vs_bq_cs: (1 - AUC) of FLAV_F0 over L188, run by run :: sha256 c6c4138ee5c86423 :: prose twin of the run farther from 1
\newcommand{\MlstCmpBcFoneOverFull}{1.518 and 1.530}  % experiments/FIGS/data/v2/index.json :: binary probe bc_vs_bq_cs: (1 - AUC) of FLAV_F1 over L188, run by run :: sha256 c6c4138ee5c86423
\newcommand{\MlstCmpBcFoneOverFullPr}{at least 52\% more errors}  % experiments/FIGS/data/v2/index.json :: binary probe bc_vs_bq_cs: (1 - AUC) of FLAV_F1 over L188, run by run :: sha256 c6c4138ee5c86423 :: prose twin: two significant figures, direction in words
\newcommand{\MlstCmpBcFoneOverFullPrMax}{up to 53\% more errors}  % experiments/FIGS/data/v2/index.json :: binary probe bc_vs_bq_cs: (1 - AUC) of FLAV_F1 over L188, run by run :: sha256 c6c4138ee5c86423 :: prose twin of the run farther from 1
\newcommand{\MlstCmpBcFonerOverFull}{1.530 and 1.515}  % experiments/FIGS/data/v2/index.json :: binary probe bc_vs_bq_cs: (1 - AUC) of FLAV_F1R over L188, run by run :: sha256 c6c4138ee5c86423
\newcommand{\MlstCmpBcFonerOverFullPr}{at least 51\% more errors}  % experiments/FIGS/data/v2/index.json :: binary probe bc_vs_bq_cs: (1 - AUC) of FLAV_F1R over L188, run by run :: sha256 c6c4138ee5c86423 :: prose twin: two significant figures, direction in words
\newcommand{\MlstCmpBcFonerOverFullPrMax}{up to 53\% more errors}  % experiments/FIGS/data/v2/index.json :: binary probe bc_vs_bq_cs: (1 - AUC) of FLAV_F1R over L188, run by run :: sha256 c6c4138ee5c86423 :: prose twin of the run farther from 1
\newcommand{\MlstCmpBcRandOverFull}{1.15\,\ensuremath{\pm}\,0.07}  % experiments/FIGS/data/v2/index.json :: binary probe bc_vs_bq_cs: (1 - AUC) of each random partition over the 188-class run of its seed :: sha256 c6c4138ee5c86423
\newcommand{\MlstCmpBcRandOverFullPr}{15\% more errors}  % experiments/FIGS/data/v2/index.json :: binary probe bc_vs_bq_cs: (1 - AUC) of each random partition over the 188-class run of its seed :: sha256 c6c4138ee5c86423 :: prose twin: two significant figures, direction in words
\newcommand{\MlstPairRandMergedBcN}{2}  % experiments/FIGS/data/v2/index.json :: random partitions that merge bc_vs_bq_cs (count) :: sha256 c6c4138ee5c86423
\newcommand{\MlstPairRandApartBcN}{2}  % experiments/FIGS/data/v2/index.json :: random partitions that keep apart bc_vs_bq_cs (count) :: sha256 c6c4138ee5c86423
\newcommand{\MlstPairRandMergedOverApartBc}{1.12}  % experiments/FIGS/data/v2/index.json :: binary probe bc_vs_bq_cs: mean (1 - AUC) of the partitions that merge it over that of those that keep it apart (one run each) :: sha256 c6c4138ee5c86423
\newcommand{\MlstPairRandMergedOverApartBcPr}{12\% more errors}  % experiments/FIGS/data/v2/index.json :: binary probe bc_vs_bq_cs: mean (1 - AUC) of the partitions that merge it over that of those that keep it apart (one run each) :: sha256 c6c4138ee5c86423 :: prose twin: two significant figures, direction in words
\newcommand{\MlstCmpBbResOverFull}{1.015\,\ensuremath{\pm}\,0.027}  % experiments/FIGS/data/v2/index.json :: binary probe bb_vs_cc: (1 - AUC) of L162 over L188, run by run :: sha256 c6c4138ee5c86423
\newcommand{\MlstCmpBbResOverFullPr}{equal within 3\%}  % experiments/FIGS/data/v2/index.json :: binary probe bb_vs_cc: (1 - AUC) of L162 over L188, run by run :: sha256 c6c4138ee5c86423 :: prose twin: two significant figures, direction in words
\newcommand{\MlstCmpBbMidOverFull}{1.102\,\ensuremath{\pm}\,0.013}  % experiments/FIGS/data/v2/index.json :: binary probe bb_vs_cc: (1 - AUC) of R42_Q1 over L188, run by run :: sha256 c6c4138ee5c86423
\newcommand{\MlstCmpBbMidOverFullPr}{10\% more errors}  % experiments/FIGS/data/v2/index.json :: binary probe bb_vs_cc: (1 - AUC) of R42_Q1 over L188, run by run :: sha256 c6c4138ee5c86423 :: prose twin: two significant figures, direction in words
\newcommand{\MlstCmpBbCoarseOverFull}{4.22\,\ensuremath{\pm}\,0.29}  % experiments/FIGS/data/v2/index.json :: binary probe bb_vs_cc: (1 - AUC) of R16_Q1 over L188, run by run :: sha256 c6c4138ee5c86423
\newcommand{\MlstCmpBbCoarseOverFullPr}{4.2 times the errors}  % experiments/FIGS/data/v2/index.json :: binary probe bb_vs_cc: (1 - AUC) of R16_Q1 over L188, run by run :: sha256 c6c4138ee5c86423 :: prose twin: two significant figures, direction in words
\newcommand{\MlstCmpBbFoneOverFzero}{0.31 and 0.13}  % experiments/FIGS/data/v2/index.json :: binary probe bb_vs_cc: (1 - AUC) of FLAV_F1 over FLAV_F0, run by run :: sha256 c6c4138ee5c86423
\newcommand{\MlstCmpBbFoneOverFzeroPr}{at most 0.31 times the errors}  % experiments/FIGS/data/v2/index.json :: binary probe bb_vs_cc: (1 - AUC) of FLAV_F1 over FLAV_F0, run by run :: sha256 c6c4138ee5c86423 :: prose twin: two significant figures, direction in words
\newcommand{\MlstCmpBbFoneOverFzeroPrMax}{down to 0.13 times the errors}  % experiments/FIGS/data/v2/index.json :: binary probe bb_vs_cc: (1 - AUC) of FLAV_F1 over FLAV_F0, run by run :: sha256 c6c4138ee5c86423 :: prose twin of the run farther from 1
\newcommand{\MlstCmpBbFonerOverFone}{1.40 and 1.86}  % experiments/FIGS/data/v2/index.json :: binary probe bb_vs_cc: (1 - AUC) of FLAV_F1R over FLAV_F1, run by run :: sha256 c6c4138ee5c86423
\newcommand{\MlstCmpBbFonerOverFonePr}{at least 40\% more errors}  % experiments/FIGS/data/v2/index.json :: binary probe bb_vs_cc: (1 - AUC) of FLAV_F1R over FLAV_F1, run by run :: sha256 c6c4138ee5c86423 :: prose twin: two significant figures, direction in words
\newcommand{\MlstCmpBbFonerOverFonePrMax}{up to 86\% more errors}  % experiments/FIGS/data/v2/index.json :: binary probe bb_vs_cc: (1 - AUC) of FLAV_F1R over FLAV_F1, run by run :: sha256 c6c4138ee5c86423 :: prose twin of the run farther from 1
\newcommand{\MlstCmpBbFonerOverFzero}{0.44 and 0.24}  % experiments/FIGS/data/v2/index.json :: binary probe bb_vs_cc: (1 - AUC) of FLAV_F1R over FLAV_F0, run by run :: sha256 c6c4138ee5c86423
\newcommand{\MlstCmpBbFonerOverFzeroPr}{at most 0.44 times the errors}  % experiments/FIGS/data/v2/index.json :: binary probe bb_vs_cc: (1 - AUC) of FLAV_F1R over FLAV_F0, run by run :: sha256 c6c4138ee5c86423 :: prose twin: two significant figures, direction in words
\newcommand{\MlstCmpBbFonerOverFzeroPrMax}{down to 0.24 times the errors}  % experiments/FIGS/data/v2/index.json :: binary probe bb_vs_cc: (1 - AUC) of FLAV_F1R over FLAV_F0, run by run :: sha256 c6c4138ee5c86423 :: prose twin of the run farther from 1
\newcommand{\MlstCmpBbFzeroOverCoarse}{1.3 and 3.2}  % experiments/FIGS/data/v2/index.json :: binary probe bb_vs_cc: (1 - AUC) of FLAV_F0 over R16_Q1, run by run :: sha256 c6c4138ee5c86423
\newcommand{\MlstCmpBbFzeroOverCoarsePr}{at least 30\% more errors}  % experiments/FIGS/data/v2/index.json :: binary probe bb_vs_cc: (1 - AUC) of FLAV_F0 over R16_Q1, run by run :: sha256 c6c4138ee5c86423 :: prose twin: two significant figures, direction in words
\newcommand{\MlstCmpBbFzeroOverCoarsePrMax}{up to 3.2 times the errors}  % experiments/FIGS/data/v2/index.json :: binary probe bb_vs_cc: (1 - AUC) of FLAV_F0 over R16_Q1, run by run :: sha256 c6c4138ee5c86423 :: prose twin of the run farther from 1
\newcommand{\MlstCmpBbFoneOverCoarse}{0.4077 and 0.4126}  % experiments/FIGS/data/v2/index.json :: binary probe bb_vs_cc: (1 - AUC) of FLAV_F1 over R16_Q1, run by run :: sha256 c6c4138ee5c86423
\newcommand{\MlstCmpBbFoneOverCoarsePr}{at most 0.41 times the errors}  % experiments/FIGS/data/v2/index.json :: binary probe bb_vs_cc: (1 - AUC) of FLAV_F1 over R16_Q1, run by run :: sha256 c6c4138ee5c86423 :: prose twin: two significant figures, direction in words
\newcommand{\MlstCmpBbFoneOverCoarsePrMax}{down to 0.41 times the errors}  % experiments/FIGS/data/v2/index.json :: binary probe bb_vs_cc: (1 - AUC) of FLAV_F1 over R16_Q1, run by run :: sha256 c6c4138ee5c86423 :: prose twin of the run farther from 1
\newcommand{\MlstCmpBbFonerOverCoarse}{0.57 and 0.77}  % experiments/FIGS/data/v2/index.json :: binary probe bb_vs_cc: (1 - AUC) of FLAV_F1R over R16_Q1, run by run :: sha256 c6c4138ee5c86423
\newcommand{\MlstCmpBbFonerOverCoarsePr}{at least 23\% fewer errors}  % experiments/FIGS/data/v2/index.json :: binary probe bb_vs_cc: (1 - AUC) of FLAV_F1R over R16_Q1, run by run :: sha256 c6c4138ee5c86423 :: prose twin: two significant figures, direction in words
\newcommand{\MlstCmpBbFonerOverCoarsePrMax}{up to 43\% fewer errors}  % experiments/FIGS/data/v2/index.json :: binary probe bb_vs_cc: (1 - AUC) of FLAV_F1R over R16_Q1, run by run :: sha256 c6c4138ee5c86423 :: prose twin of the run farther from 1
\newcommand{\MlstCmpBbFzeroOverFull}{5.9 and 13.4}  % experiments/FIGS/data/v2/index.json :: binary probe bb_vs_cc: (1 - AUC) of FLAV_F0 over L188, run by run :: sha256 c6c4138ee5c86423
\newcommand{\MlstCmpBbFzeroOverFullPr}{at least 5.9 times the errors}  % experiments/FIGS/data/v2/index.json :: binary probe bb_vs_cc: (1 - AUC) of FLAV_F0 over L188, run by run :: sha256 c6c4138ee5c86423 :: prose twin: two significant figures, direction in words
\newcommand{\MlstCmpBbFzeroOverFullPrMax}{up to 13 times the errors}  % experiments/FIGS/data/v2/index.json :: binary probe bb_vs_cc: (1 - AUC) of FLAV_F0 over L188, run by run :: sha256 c6c4138ee5c86423 :: prose twin of the run farther from 1
\newcommand{\MlstCmpBbFoneOverFull}{1.837 and 1.748}  % experiments/FIGS/data/v2/index.json :: binary probe bb_vs_cc: (1 - AUC) of FLAV_F1 over L188, run by run :: sha256 c6c4138ee5c86423
\newcommand{\MlstCmpBbFoneOverFullPr}{at least 75\% more errors}  % experiments/FIGS/data/v2/index.json :: binary probe bb_vs_cc: (1 - AUC) of FLAV_F1 over L188, run by run :: sha256 c6c4138ee5c86423 :: prose twin: two significant figures, direction in words
\newcommand{\MlstCmpBbFoneOverFullPrMax}{up to 84\% more errors}  % experiments/FIGS/data/v2/index.json :: binary probe bb_vs_cc: (1 - AUC) of FLAV_F1 over L188, run by run :: sha256 c6c4138ee5c86423 :: prose twin of the run farther from 1
\newcommand{\MlstCmpBbFonerOverFull}{2.58 and 3.25}  % experiments/FIGS/data/v2/index.json :: binary probe bb_vs_cc: (1 - AUC) of FLAV_F1R over L188, run by run :: sha256 c6c4138ee5c86423
\newcommand{\MlstCmpBbFonerOverFullPr}{at least 2.6 times the errors}  % experiments/FIGS/data/v2/index.json :: binary probe bb_vs_cc: (1 - AUC) of FLAV_F1R over L188, run by run :: sha256 c6c4138ee5c86423 :: prose twin: two significant figures, direction in words
\newcommand{\MlstCmpBbFonerOverFullPrMax}{up to 3.2 times the errors}  % experiments/FIGS/data/v2/index.json :: binary probe bb_vs_cc: (1 - AUC) of FLAV_F1R over L188, run by run :: sha256 c6c4138ee5c86423 :: prose twin of the run farther from 1
\newcommand{\MlstCmpBbRandOverFull}{1.53\,\ensuremath{\pm}\,0.11}  % experiments/FIGS/data/v2/index.json :: binary probe bb_vs_cc: (1 - AUC) of each random partition over the 188-class run of its seed :: sha256 c6c4138ee5c86423
\newcommand{\MlstCmpBbRandOverFullPr}{53\% more errors}  % experiments/FIGS/data/v2/index.json :: binary probe bb_vs_cc: (1 - AUC) of each random partition over the 188-class run of its seed :: sha256 c6c4138ee5c86423 :: prose twin: two significant figures, direction in words
\newcommand{\MlstPairRandMergedBbN}{3}  % experiments/FIGS/data/v2/index.json :: random partitions that merge bb_vs_cc (count) :: sha256 c6c4138ee5c86423
\newcommand{\MlstPairRandApartBbN}{2}  % experiments/FIGS/data/v2/index.json :: random partitions that keep apart bb_vs_cc (count) :: sha256 c6c4138ee5c86423
\newcommand{\MlstPairRandMergedOverApartBb}{1.07}  % experiments/FIGS/data/v2/index.json :: binary probe bb_vs_cc: mean (1 - AUC) of the partitions that merge it over that of those that keep it apart (one run each) :: sha256 c6c4138ee5c86423
\newcommand{\MlstPairRandMergedOverApartBbPr}{7.3\% more errors}  % experiments/FIGS/data/v2/index.json :: binary probe bb_vs_cc: mean (1 - AUC) of the partitions that merge it over that of those that keep it apart (one run each) :: sha256 c6c4138ee5c86423 :: prose twin: two significant figures, direction in words
\newcommand{\MlstCmpBbqqResOverFull}{1.005\,\ensuremath{\pm}\,0.032}  % experiments/FIGS/data/v2/index.json :: binary probe bbqq_vs_ccqq: (1 - AUC) of L162 over L188, run by run :: sha256 c6c4138ee5c86423
\newcommand{\MlstCmpBbqqResOverFullPr}{equal within 4\%}  % experiments/FIGS/data/v2/index.json :: binary probe bbqq_vs_ccqq: (1 - AUC) of L162 over L188, run by run :: sha256 c6c4138ee5c86423 :: prose twin: two significant figures, direction in words
\newcommand{\MlstCmpBbqqMidOverFull}{1.117\,\ensuremath{\pm}\,0.008}  % experiments/FIGS/data/v2/index.json :: binary probe bbqq_vs_ccqq: (1 - AUC) of R42_Q1 over L188, run by run :: sha256 c6c4138ee5c86423
\newcommand{\MlstCmpBbqqMidOverFullPr}{12\% more errors}  % experiments/FIGS/data/v2/index.json :: binary probe bbqq_vs_ccqq: (1 - AUC) of R42_Q1 over L188, run by run :: sha256 c6c4138ee5c86423 :: prose twin: two significant figures, direction in words
\newcommand{\MlstCmpBbqqCoarseOverFull}{4.7\,\ensuremath{\pm}\,0.5}  % experiments/FIGS/data/v2/index.json :: binary probe bbqq_vs_ccqq: (1 - AUC) of R16_Q1 over L188, run by run :: sha256 c6c4138ee5c86423
\newcommand{\MlstCmpBbqqCoarseOverFullPr}{4.7 times the errors}  % experiments/FIGS/data/v2/index.json :: binary probe bbqq_vs_ccqq: (1 - AUC) of R16_Q1 over L188, run by run :: sha256 c6c4138ee5c86423 :: prose twin: two significant figures, direction in words
\newcommand{\MlstCmpBbqqFoneOverFzero}{0.236 and 0.154}  % experiments/FIGS/data/v2/index.json :: binary probe bbqq_vs_ccqq: (1 - AUC) of FLAV_F1 over FLAV_F0, run by run :: sha256 c6c4138ee5c86423
\newcommand{\MlstCmpBbqqFoneOverFzeroPr}{at most 0.24 times the errors}  % experiments/FIGS/data/v2/index.json :: binary probe bbqq_vs_ccqq: (1 - AUC) of FLAV_F1 over FLAV_F0, run by run :: sha256 c6c4138ee5c86423 :: prose twin: two significant figures, direction in words
\newcommand{\MlstCmpBbqqFoneOverFzeroPrMax}{down to 0.15 times the errors}  % experiments/FIGS/data/v2/index.json :: binary probe bbqq_vs_ccqq: (1 - AUC) of FLAV_F1 over FLAV_F0, run by run :: sha256 c6c4138ee5c86423 :: prose twin of the run farther from 1
\newcommand{\MlstCmpBbqqFonerOverFone}{1.505 and 1.577}  % experiments/FIGS/data/v2/index.json :: binary probe bbqq_vs_ccqq: (1 - AUC) of FLAV_F1R over FLAV_F1, run by run :: sha256 c6c4138ee5c86423
\newcommand{\MlstCmpBbqqFonerOverFonePr}{at least 50\% more errors}  % experiments/FIGS/data/v2/index.json :: binary probe bbqq_vs_ccqq: (1 - AUC) of FLAV_F1R over FLAV_F1, run by run :: sha256 c6c4138ee5c86423 :: prose twin: two significant figures, direction in words
\newcommand{\MlstCmpBbqqFonerOverFonePrMax}{up to 58\% more errors}  % experiments/FIGS/data/v2/index.json :: binary probe bbqq_vs_ccqq: (1 - AUC) of FLAV_F1R over FLAV_F1, run by run :: sha256 c6c4138ee5c86423 :: prose twin of the run farther from 1
\newcommand{\MlstCmpBbqqFonerOverFzero}{0.36 and 0.24}  % experiments/FIGS/data/v2/index.json :: binary probe bbqq_vs_ccqq: (1 - AUC) of FLAV_F1R over FLAV_F0, run by run :: sha256 c6c4138ee5c86423
\newcommand{\MlstCmpBbqqFonerOverFzeroPr}{at most 0.36 times the errors}  % experiments/FIGS/data/v2/index.json :: binary probe bbqq_vs_ccqq: (1 - AUC) of FLAV_F1R over FLAV_F0, run by run :: sha256 c6c4138ee5c86423 :: prose twin: two significant figures, direction in words
\newcommand{\MlstCmpBbqqFonerOverFzeroPrMax}{down to 0.24 times the errors}  % experiments/FIGS/data/v2/index.json :: binary probe bbqq_vs_ccqq: (1 - AUC) of FLAV_F1R over FLAV_F0, run by run :: sha256 c6c4138ee5c86423 :: prose twin of the run farther from 1
\newcommand{\MlstCmpBbqqFzeroOverCoarse}{1.51 and 1.79}  % experiments/FIGS/data/v2/index.json :: binary probe bbqq_vs_ccqq: (1 - AUC) of FLAV_F0 over R16_Q1, run by run :: sha256 c6c4138ee5c86423
\newcommand{\MlstCmpBbqqFzeroOverCoarsePr}{at least 51\% more errors}  % experiments/FIGS/data/v2/index.json :: binary probe bbqq_vs_ccqq: (1 - AUC) of FLAV_F0 over R16_Q1, run by run :: sha256 c6c4138ee5c86423 :: prose twin: two significant figures, direction in words
\newcommand{\MlstCmpBbqqFzeroOverCoarsePrMax}{up to 79\% more errors}  % experiments/FIGS/data/v2/index.json :: binary probe bbqq_vs_ccqq: (1 - AUC) of FLAV_F0 over R16_Q1, run by run :: sha256 c6c4138ee5c86423 :: prose twin of the run farther from 1
\newcommand{\MlstCmpBbqqFoneOverCoarse}{0.356 and 0.276}  % experiments/FIGS/data/v2/index.json :: binary probe bbqq_vs_ccqq: (1 - AUC) of FLAV_F1 over R16_Q1, run by run :: sha256 c6c4138ee5c86423
\newcommand{\MlstCmpBbqqFoneOverCoarsePr}{at most 0.36 times the errors}  % experiments/FIGS/data/v2/index.json :: binary probe bbqq_vs_ccqq: (1 - AUC) of FLAV_F1 over R16_Q1, run by run :: sha256 c6c4138ee5c86423 :: prose twin: two significant figures, direction in words
\newcommand{\MlstCmpBbqqFoneOverCoarsePrMax}{down to 0.28 times the errors}  % experiments/FIGS/data/v2/index.json :: binary probe bbqq_vs_ccqq: (1 - AUC) of FLAV_F1 over R16_Q1, run by run :: sha256 c6c4138ee5c86423 :: prose twin of the run farther from 1
\newcommand{\MlstCmpBbqqFonerOverCoarse}{0.54 and 0.44}  % experiments/FIGS/data/v2/index.json :: binary probe bbqq_vs_ccqq: (1 - AUC) of FLAV_F1R over R16_Q1, run by run :: sha256 c6c4138ee5c86423
\newcommand{\MlstCmpBbqqFonerOverCoarsePr}{at least 46\% fewer errors}  % experiments/FIGS/data/v2/index.json :: binary probe bbqq_vs_ccqq: (1 - AUC) of FLAV_F1R over R16_Q1, run by run :: sha256 c6c4138ee5c86423 :: prose twin: two significant figures, direction in words
\newcommand{\MlstCmpBbqqFonerOverCoarsePrMax}{down to 0.44 times the errors}  % experiments/FIGS/data/v2/index.json :: binary probe bbqq_vs_ccqq: (1 - AUC) of FLAV_F1R over R16_Q1, run by run :: sha256 c6c4138ee5c86423 :: prose twin of the run farther from 1
\newcommand{\MlstCmpBbqqFzeroOverFull}{6.4 and 9.4}  % experiments/FIGS/data/v2/index.json :: binary probe bbqq_vs_ccqq: (1 - AUC) of FLAV_F0 over L188, run by run :: sha256 c6c4138ee5c86423
\newcommand{\MlstCmpBbqqFzeroOverFullPr}{at least 6.4 times the errors}  % experiments/FIGS/data/v2/index.json :: binary probe bbqq_vs_ccqq: (1 - AUC) of FLAV_F0 over L188, run by run :: sha256 c6c4138ee5c86423 :: prose twin: two significant figures, direction in words
\newcommand{\MlstCmpBbqqFzeroOverFullPrMax}{up to 9.4 times the errors}  % experiments/FIGS/data/v2/index.json :: binary probe bbqq_vs_ccqq: (1 - AUC) of FLAV_F0 over L188, run by run :: sha256 c6c4138ee5c86423 :: prose twin of the run farther from 1
\newcommand{\MlstCmpBbqqFoneOverFull}{1.510 and 1.447}  % experiments/FIGS/data/v2/index.json :: binary probe bbqq_vs_ccqq: (1 - AUC) of FLAV_F1 over L188, run by run :: sha256 c6c4138ee5c86423
\newcommand{\MlstCmpBbqqFoneOverFullPr}{at least 45\% more errors}  % experiments/FIGS/data/v2/index.json :: binary probe bbqq_vs_ccqq: (1 - AUC) of FLAV_F1 over L188, run by run :: sha256 c6c4138ee5c86423 :: prose twin: two significant figures, direction in words
\newcommand{\MlstCmpBbqqFoneOverFullPrMax}{up to 51\% more errors}  % experiments/FIGS/data/v2/index.json :: binary probe bbqq_vs_ccqq: (1 - AUC) of FLAV_F1 over L188, run by run :: sha256 c6c4138ee5c86423 :: prose twin of the run farther from 1
\newcommand{\MlstCmpBbqqFonerOverFull}{2.272 and 2.281}  % experiments/FIGS/data/v2/index.json :: binary probe bbqq_vs_ccqq: (1 - AUC) of FLAV_F1R over L188, run by run :: sha256 c6c4138ee5c86423
\newcommand{\MlstCmpBbqqFonerOverFullPr}{at least 2.3 times the errors}  % experiments/FIGS/data/v2/index.json :: binary probe bbqq_vs_ccqq: (1 - AUC) of FLAV_F1R over L188, run by run :: sha256 c6c4138ee5c86423 :: prose twin: two significant figures, direction in words
\newcommand{\MlstCmpBbqqFonerOverFullPrMax}{up to 2.3 times the errors}  % experiments/FIGS/data/v2/index.json :: binary probe bbqq_vs_ccqq: (1 - AUC) of FLAV_F1R over L188, run by run :: sha256 c6c4138ee5c86423 :: prose twin of the run farther from 1
\newcommand{\MlstCmpBbqqRandOverFull}{1.42\,\ensuremath{\pm}\,0.06}  % experiments/FIGS/data/v2/index.json :: binary probe bbqq_vs_ccqq: (1 - AUC) of each random partition over the 188-class run of its seed :: sha256 c6c4138ee5c86423
\newcommand{\MlstCmpBbqqRandOverFullPr}{42\% more errors}  % experiments/FIGS/data/v2/index.json :: binary probe bbqq_vs_ccqq: (1 - AUC) of each random partition over the 188-class run of its seed :: sha256 c6c4138ee5c86423 :: prose twin: two significant figures, direction in words
\newcommand{\MlstPairRandMergedBbqqN}{2}  % experiments/FIGS/data/v2/index.json :: random partitions that merge bbqq_vs_ccqq (count) :: sha256 c6c4138ee5c86423
\newcommand{\MlstPairRandApartBbqqN}{3}  % experiments/FIGS/data/v2/index.json :: random partitions that keep apart bbqq_vs_ccqq (count) :: sha256 c6c4138ee5c86423
\newcommand{\MlstPairRandMergedOverApartBbqq}{0.996}  % experiments/FIGS/data/v2/index.json :: binary probe bbqq_vs_ccqq: mean (1 - AUC) of the partitions that merge it over that of those that keep it apart (one run each) :: sha256 c6c4138ee5c86423
\newcommand{\MlstPairRandMergedOverApartBbqqPr}{0.44\% fewer errors}  % experiments/FIGS/data/v2/index.json :: binary probe bbqq_vs_ccqq: mean (1 - AUC) of the partitions that merge it over that of those that keep it apart (one run each) :: sha256 c6c4138ee5c86423 :: prose twin: two significant figures, direction in words
\newcommand{\MlstCmpEemmResOverFull}{--}  % experiments/FIGS/data/v2/index.json :: binary probe ee_vs_mm: (1 - AUC) of L162 over L188, run by run :: sha256 c6c4138ee5c86423 :: the reference's 1 - AUC is 0
\newcommand{\MlstCmpEemmResOverFullPr}{undefined, since the reference separates the test jets completely}  % experiments/FIGS/data/v2/index.json :: binary probe ee_vs_mm: (1 - AUC) of L162 over L188, run by run :: sha256 c6c4138ee5c86423 :: the reference's 1 - AUC is 0
\newcommand{\MlstCmpEemmMidOverFull}{--}  % experiments/FIGS/data/v2/index.json :: binary probe ee_vs_mm: (1 - AUC) of R42_Q1 over L188, run by run :: sha256 c6c4138ee5c86423 :: the reference's 1 - AUC is 0
\newcommand{\MlstCmpEemmMidOverFullPr}{undefined, since the reference separates the test jets completely}  % experiments/FIGS/data/v2/index.json :: binary probe ee_vs_mm: (1 - AUC) of R42_Q1 over L188, run by run :: sha256 c6c4138ee5c86423 :: the reference's 1 - AUC is 0
\newcommand{\MlstCmpEemmCoarseOverFull}{--}  % experiments/FIGS/data/v2/index.json :: binary probe ee_vs_mm: (1 - AUC) of R16_Q1 over L188, run by run :: sha256 c6c4138ee5c86423 :: the reference's 1 - AUC is 0
\newcommand{\MlstCmpEemmCoarseOverFullPr}{undefined, since the reference separates the test jets completely}  % experiments/FIGS/data/v2/index.json :: binary probe ee_vs_mm: (1 - AUC) of R16_Q1 over L188, run by run :: sha256 c6c4138ee5c86423 :: the reference's 1 - AUC is 0
\newcommand{\MlstCmpEemmFoneOverFzero}{--}  % experiments/FIGS/data/v2/index.json :: binary probe ee_vs_mm: (1 - AUC) of FLAV_F1 over FLAV_F0, run by run :: sha256 c6c4138ee5c86423 :: the reference's 1 - AUC is 0
\newcommand{\MlstCmpEemmFoneOverFzeroPr}{undefined, since the reference separates the test jets completely}  % experiments/FIGS/data/v2/index.json :: binary probe ee_vs_mm: (1 - AUC) of FLAV_F1 over FLAV_F0, run by run :: sha256 c6c4138ee5c86423 :: the reference's 1 - AUC is 0
\newcommand{\MlstCmpEemmFoneOverFzeroPrMax}{undefined, since the reference separates the test jets completely}  % experiments/FIGS/data/v2/index.json :: binary probe ee_vs_mm: (1 - AUC) of FLAV_F1 over FLAV_F0, run by run :: sha256 c6c4138ee5c86423 :: the reference's 1 - AUC is 0
\newcommand{\MlstCmpEemmFonerOverFone}{--}  % experiments/FIGS/data/v2/index.json :: binary probe ee_vs_mm: (1 - AUC) of FLAV_F1R over FLAV_F1, run by run :: sha256 c6c4138ee5c86423 :: the reference's 1 - AUC is 0
\newcommand{\MlstCmpEemmFonerOverFonePr}{undefined, since the reference separates the test jets completely}  % experiments/FIGS/data/v2/index.json :: binary probe ee_vs_mm: (1 - AUC) of FLAV_F1R over FLAV_F1, run by run :: sha256 c6c4138ee5c86423 :: the reference's 1 - AUC is 0
\newcommand{\MlstCmpEemmFonerOverFonePrMax}{undefined, since the reference separates the test jets completely}  % experiments/FIGS/data/v2/index.json :: binary probe ee_vs_mm: (1 - AUC) of FLAV_F1R over FLAV_F1, run by run :: sha256 c6c4138ee5c86423 :: the reference's 1 - AUC is 0
\newcommand{\MlstCmpEemmFonerOverFzero}{--}  % experiments/FIGS/data/v2/index.json :: binary probe ee_vs_mm: (1 - AUC) of FLAV_F1R over FLAV_F0, run by run :: sha256 c6c4138ee5c86423 :: the reference's 1 - AUC is 0
\newcommand{\MlstCmpEemmFonerOverFzeroPr}{undefined, since the reference separates the test jets completely}  % experiments/FIGS/data/v2/index.json :: binary probe ee_vs_mm: (1 - AUC) of FLAV_F1R over FLAV_F0, run by run :: sha256 c6c4138ee5c86423 :: the reference's 1 - AUC is 0
\newcommand{\MlstCmpEemmFonerOverFzeroPrMax}{undefined, since the reference separates the test jets completely}  % experiments/FIGS/data/v2/index.json :: binary probe ee_vs_mm: (1 - AUC) of FLAV_F1R over FLAV_F0, run by run :: sha256 c6c4138ee5c86423 :: the reference's 1 - AUC is 0
\newcommand{\MlstCmpEemmFzeroOverCoarse}{0.0000022 and 0.0000003}  % experiments/FIGS/data/v2/index.json :: binary probe ee_vs_mm: (1 - AUC) of FLAV_F0 over R16_Q1, run by run :: sha256 c6c4138ee5c86423
\newcommand{\MlstCmpEemmFzeroOverCoarsePr}{at most 0.0000022 times the errors}  % experiments/FIGS/data/v2/index.json :: binary probe ee_vs_mm: (1 - AUC) of FLAV_F0 over R16_Q1, run by run :: sha256 c6c4138ee5c86423 :: prose twin: two significant figures, direction in words
\newcommand{\MlstCmpEemmFzeroOverCoarsePrMax}{down to 0.00000034 times the errors}  % experiments/FIGS/data/v2/index.json :: binary probe ee_vs_mm: (1 - AUC) of FLAV_F0 over R16_Q1, run by run :: sha256 c6c4138ee5c86423 :: prose twin of the run farther from 1
\newcommand{\MlstCmpEemmFoneOverCoarse}{0.000000 and 0.000038}  % experiments/FIGS/data/v2/index.json :: binary probe ee_vs_mm: (1 - AUC) of FLAV_F1 over R16_Q1, run by run :: sha256 c6c4138ee5c86423
\newcommand{\MlstCmpEemmFoneOverCoarsePr}{at most 0.000038 times the errors}  % experiments/FIGS/data/v2/index.json :: binary probe ee_vs_mm: (1 - AUC) of FLAV_F1 over R16_Q1, run by run :: sha256 c6c4138ee5c86423 :: prose twin: two significant figures, direction in words
\newcommand{\MlstCmpEemmFoneOverCoarsePrMax}{down to 0.000000042 times the errors}  % experiments/FIGS/data/v2/index.json :: binary probe ee_vs_mm: (1 - AUC) of FLAV_F1 over R16_Q1, run by run :: sha256 c6c4138ee5c86423 :: prose twin of the run farther from 1
\newcommand{\MlstCmpEemmFonerOverCoarse}{0.00000005 and 0.00000017}  % experiments/FIGS/data/v2/index.json :: binary probe ee_vs_mm: (1 - AUC) of FLAV_F1R over R16_Q1, run by run :: sha256 c6c4138ee5c86423
\newcommand{\MlstCmpEemmFonerOverCoarsePr}{at most 0.00000017 times the errors}  % experiments/FIGS/data/v2/index.json :: binary probe ee_vs_mm: (1 - AUC) of FLAV_F1R over R16_Q1, run by run :: sha256 c6c4138ee5c86423 :: prose twin: two significant figures, direction in words
\newcommand{\MlstCmpEemmFonerOverCoarsePrMax}{down to 0.000000055 times the errors}  % experiments/FIGS/data/v2/index.json :: binary probe ee_vs_mm: (1 - AUC) of FLAV_F1R over R16_Q1, run by run :: sha256 c6c4138ee5c86423 :: prose twin of the run farther from 1
\newcommand{\MlstCmpEemmFzeroOverFull}{--}  % experiments/FIGS/data/v2/index.json :: binary probe ee_vs_mm: (1 - AUC) of FLAV_F0 over L188, run by run :: sha256 c6c4138ee5c86423 :: the reference's 1 - AUC is 0
\newcommand{\MlstCmpEemmFzeroOverFullPr}{undefined, since the reference separates the test jets completely}  % experiments/FIGS/data/v2/index.json :: binary probe ee_vs_mm: (1 - AUC) of FLAV_F0 over L188, run by run :: sha256 c6c4138ee5c86423 :: the reference's 1 - AUC is 0
\newcommand{\MlstCmpEemmFzeroOverFullPrMax}{undefined, since the reference separates the test jets completely}  % experiments/FIGS/data/v2/index.json :: binary probe ee_vs_mm: (1 - AUC) of FLAV_F0 over L188, run by run :: sha256 c6c4138ee5c86423 :: the reference's 1 - AUC is 0
\newcommand{\MlstCmpEemmFoneOverFull}{--}  % experiments/FIGS/data/v2/index.json :: binary probe ee_vs_mm: (1 - AUC) of FLAV_F1 over L188, run by run :: sha256 c6c4138ee5c86423 :: the reference's 1 - AUC is 0
\newcommand{\MlstCmpEemmFoneOverFullPr}{undefined, since the reference separates the test jets completely}  % experiments/FIGS/data/v2/index.json :: binary probe ee_vs_mm: (1 - AUC) of FLAV_F1 over L188, run by run :: sha256 c6c4138ee5c86423 :: the reference's 1 - AUC is 0
\newcommand{\MlstCmpEemmFoneOverFullPrMax}{undefined, since the reference separates the test jets completely}  % experiments/FIGS/data/v2/index.json :: binary probe ee_vs_mm: (1 - AUC) of FLAV_F1 over L188, run by run :: sha256 c6c4138ee5c86423 :: the reference's 1 - AUC is 0
\newcommand{\MlstCmpEemmFonerOverFull}{--}  % experiments/FIGS/data/v2/index.json :: binary probe ee_vs_mm: (1 - AUC) of FLAV_F1R over L188, run by run :: sha256 c6c4138ee5c86423 :: the reference's 1 - AUC is 0
\newcommand{\MlstCmpEemmFonerOverFullPr}{undefined, since the reference separates the test jets completely}  % experiments/FIGS/data/v2/index.json :: binary probe ee_vs_mm: (1 - AUC) of FLAV_F1R over L188, run by run :: sha256 c6c4138ee5c86423 :: the reference's 1 - AUC is 0
\newcommand{\MlstCmpEemmFonerOverFullPrMax}{undefined, since the reference separates the test jets completely}  % experiments/FIGS/data/v2/index.json :: binary probe ee_vs_mm: (1 - AUC) of FLAV_F1R over L188, run by run :: sha256 c6c4138ee5c86423 :: the reference's 1 - AUC is 0
\newcommand{\MlstCmpEemmRandOverFull}{--}  % experiments/FIGS/data/v2/index.json :: binary probe ee_vs_mm: (1 - AUC) of each random partition over the 188-class run of its seed :: sha256 c6c4138ee5c86423 :: the reference's 1 - AUC is 0
\newcommand{\MlstCmpEemmRandOverFullPr}{undefined, since the reference separates the test jets completely}  % experiments/FIGS/data/v2/index.json :: binary probe ee_vs_mm: (1 - AUC) of each random partition over the 188-class run of its seed :: sha256 c6c4138ee5c86423 :: the reference's 1 - AUC is 0
\newcommand{\MlstPairRandMergedEemmN}{3}  % experiments/FIGS/data/v2/index.json :: random partitions that merge ee_vs_mm (count) :: sha256 c6c4138ee5c86423
\newcommand{\MlstPairRandApartEemmN}{2}  % experiments/FIGS/data/v2/index.json :: random partitions that keep apart ee_vs_mm (count) :: sha256 c6c4138ee5c86423
\newcommand{\MlstPairRandMergedOverApartEemm}{75.3}  % experiments/FIGS/data/v2/index.json :: binary probe ee_vs_mm: mean (1 - AUC) of the partitions that merge it over that of those that keep it apart (one run each) :: sha256 c6c4138ee5c86423
\newcommand{\MlstPairRandMergedOverApartEemmPr}{75 times the errors}  % experiments/FIGS/data/v2/index.json :: binary probe ee_vs_mm: mean (1 - AUC) of the partitions that merge it over that of those that keep it apart (one run each) :: sha256 c6c4138ee5c86423 :: prose twin: two significant figures, direction in words
\newcommand{\MlstCmpQcdBbResOverFull}{1.15\,\ensuremath{\pm}\,0.04}  % experiments/FIGS/data/v2/index.json :: binary probe qcd_bb_vs_cc: (1 - AUC) of L162 over L188, run by run :: sha256 c6c4138ee5c86423
\newcommand{\MlstCmpQcdBbResOverFullPr}{15\% more errors}  % experiments/FIGS/data/v2/index.json :: binary probe qcd_bb_vs_cc: (1 - AUC) of L162 over L188, run by run :: sha256 c6c4138ee5c86423 :: prose twin: two significant figures, direction in words
\newcommand{\MlstCmpQcdBbMidOverFull}{1.163\,\ensuremath{\pm}\,0.016}  % experiments/FIGS/data/v2/index.json :: binary probe qcd_bb_vs_cc: (1 - AUC) of R42_Q1 over L188, run by run :: sha256 c6c4138ee5c86423
\newcommand{\MlstCmpQcdBbMidOverFullPr}{16\% more errors}  % experiments/FIGS/data/v2/index.json :: binary probe qcd_bb_vs_cc: (1 - AUC) of R42_Q1 over L188, run by run :: sha256 c6c4138ee5c86423 :: prose twin: two significant figures, direction in words
\newcommand{\MlstCmpQcdBbCoarseOverFull}{3.294\,\ensuremath{\pm}\,0.033}  % experiments/FIGS/data/v2/index.json :: binary probe qcd_bb_vs_cc: (1 - AUC) of R16_Q1 over L188, run by run :: sha256 c6c4138ee5c86423
\newcommand{\MlstCmpQcdBbCoarseOverFullPr}{3.3 times the errors}  % experiments/FIGS/data/v2/index.json :: binary probe qcd_bb_vs_cc: (1 - AUC) of R16_Q1 over L188, run by run :: sha256 c6c4138ee5c86423 :: prose twin: two significant figures, direction in words
\newcommand{\MlstCmpQcdBbFoneOverFzero}{0.368 and 0.294}  % experiments/FIGS/data/v2/index.json :: binary probe qcd_bb_vs_cc: (1 - AUC) of FLAV_F1 over FLAV_F0, run by run :: sha256 c6c4138ee5c86423
\newcommand{\MlstCmpQcdBbFoneOverFzeroPr}{at most 0.37 times the errors}  % experiments/FIGS/data/v2/index.json :: binary probe qcd_bb_vs_cc: (1 - AUC) of FLAV_F1 over FLAV_F0, run by run :: sha256 c6c4138ee5c86423 :: prose twin: two significant figures, direction in words
\newcommand{\MlstCmpQcdBbFoneOverFzeroPrMax}{down to 0.29 times the errors}  % experiments/FIGS/data/v2/index.json :: binary probe qcd_bb_vs_cc: (1 - AUC) of FLAV_F1 over FLAV_F0, run by run :: sha256 c6c4138ee5c86423 :: prose twin of the run farther from 1
\newcommand{\MlstCmpQcdBbFonerOverFone}{1.361 and 1.429}  % experiments/FIGS/data/v2/index.json :: binary probe qcd_bb_vs_cc: (1 - AUC) of FLAV_F1R over FLAV_F1, run by run :: sha256 c6c4138ee5c86423
\newcommand{\MlstCmpQcdBbFonerOverFonePr}{at least 36\% more errors}  % experiments/FIGS/data/v2/index.json :: binary probe qcd_bb_vs_cc: (1 - AUC) of FLAV_F1R over FLAV_F1, run by run :: sha256 c6c4138ee5c86423 :: prose twin: two significant figures, direction in words
\newcommand{\MlstCmpQcdBbFonerOverFonePrMax}{up to 43\% more errors}  % experiments/FIGS/data/v2/index.json :: binary probe qcd_bb_vs_cc: (1 - AUC) of FLAV_F1R over FLAV_F1, run by run :: sha256 c6c4138ee5c86423 :: prose twin of the run farther from 1
\newcommand{\MlstCmpQcdBbFonerOverFzero}{0.501 and 0.420}  % experiments/FIGS/data/v2/index.json :: binary probe qcd_bb_vs_cc: (1 - AUC) of FLAV_F1R over FLAV_F0, run by run :: sha256 c6c4138ee5c86423
\newcommand{\MlstCmpQcdBbFonerOverFzeroPr}{at least 50\% fewer errors}  % experiments/FIGS/data/v2/index.json :: binary probe qcd_bb_vs_cc: (1 - AUC) of FLAV_F1R over FLAV_F0, run by run :: sha256 c6c4138ee5c86423 :: prose twin: two significant figures, direction in words
\newcommand{\MlstCmpQcdBbFonerOverFzeroPrMax}{down to 0.42 times the errors}  % experiments/FIGS/data/v2/index.json :: binary probe qcd_bb_vs_cc: (1 - AUC) of FLAV_F1R over FLAV_F0, run by run :: sha256 c6c4138ee5c86423 :: prose twin of the run farther from 1
\newcommand{\MlstCmpQcdBbFzeroOverCoarse}{1.28 and 1.58}  % experiments/FIGS/data/v2/index.json :: binary probe qcd_bb_vs_cc: (1 - AUC) of FLAV_F0 over R16_Q1, run by run :: sha256 c6c4138ee5c86423
\newcommand{\MlstCmpQcdBbFzeroOverCoarsePr}{at least 28\% more errors}  % experiments/FIGS/data/v2/index.json :: binary probe qcd_bb_vs_cc: (1 - AUC) of FLAV_F0 over R16_Q1, run by run :: sha256 c6c4138ee5c86423 :: prose twin: two significant figures, direction in words
\newcommand{\MlstCmpQcdBbFzeroOverCoarsePrMax}{up to 58\% more errors}  % experiments/FIGS/data/v2/index.json :: binary probe qcd_bb_vs_cc: (1 - AUC) of FLAV_F0 over R16_Q1, run by run :: sha256 c6c4138ee5c86423 :: prose twin of the run farther from 1
\newcommand{\MlstCmpQcdBbFoneOverCoarse}{0.4692 and 0.4648}  % experiments/FIGS/data/v2/index.json :: binary probe qcd_bb_vs_cc: (1 - AUC) of FLAV_F1 over R16_Q1, run by run :: sha256 c6c4138ee5c86423
\newcommand{\MlstCmpQcdBbFoneOverCoarsePr}{at most 0.47 times the errors}  % experiments/FIGS/data/v2/index.json :: binary probe qcd_bb_vs_cc: (1 - AUC) of FLAV_F1 over R16_Q1, run by run :: sha256 c6c4138ee5c86423 :: prose twin: two significant figures, direction in words
\newcommand{\MlstCmpQcdBbFoneOverCoarsePrMax}{down to 0.46 times the errors}  % experiments/FIGS/data/v2/index.json :: binary probe qcd_bb_vs_cc: (1 - AUC) of FLAV_F1 over R16_Q1, run by run :: sha256 c6c4138ee5c86423 :: prose twin of the run farther from 1
\newcommand{\MlstCmpQcdBbFonerOverCoarse}{0.639 and 0.664}  % experiments/FIGS/data/v2/index.json :: binary probe qcd_bb_vs_cc: (1 - AUC) of FLAV_F1R over R16_Q1, run by run :: sha256 c6c4138ee5c86423
\newcommand{\MlstCmpQcdBbFonerOverCoarsePr}{at least 34\% fewer errors}  % experiments/FIGS/data/v2/index.json :: binary probe qcd_bb_vs_cc: (1 - AUC) of FLAV_F1R over R16_Q1, run by run :: sha256 c6c4138ee5c86423 :: prose twin: two significant figures, direction in words
\newcommand{\MlstCmpQcdBbFonerOverCoarsePrMax}{up to 36\% fewer errors}  % experiments/FIGS/data/v2/index.json :: binary probe qcd_bb_vs_cc: (1 - AUC) of FLAV_F1R over R16_Q1, run by run :: sha256 c6c4138ee5c86423 :: prose twin of the run farther from 1
\newcommand{\MlstCmpQcdBbFzeroOverFull}{4.24 and 5.21}  % experiments/FIGS/data/v2/index.json :: binary probe qcd_bb_vs_cc: (1 - AUC) of FLAV_F0 over L188, run by run :: sha256 c6c4138ee5c86423
\newcommand{\MlstCmpQcdBbFzeroOverFullPr}{at least 4.2 times the errors}  % experiments/FIGS/data/v2/index.json :: binary probe qcd_bb_vs_cc: (1 - AUC) of FLAV_F0 over L188, run by run :: sha256 c6c4138ee5c86423 :: prose twin: two significant figures, direction in words
\newcommand{\MlstCmpQcdBbFzeroOverFullPrMax}{up to 5.2 times the errors}  % experiments/FIGS/data/v2/index.json :: binary probe qcd_bb_vs_cc: (1 - AUC) of FLAV_F0 over L188, run by run :: sha256 c6c4138ee5c86423 :: prose twin of the run farther from 1
\newcommand{\MlstCmpQcdBbFoneOverFull}{1.561 and 1.532}  % experiments/FIGS/data/v2/index.json :: binary probe qcd_bb_vs_cc: (1 - AUC) of FLAV_F1 over L188, run by run :: sha256 c6c4138ee5c86423
\newcommand{\MlstCmpQcdBbFoneOverFullPr}{at least 53\% more errors}  % experiments/FIGS/data/v2/index.json :: binary probe qcd_bb_vs_cc: (1 - AUC) of FLAV_F1 over L188, run by run :: sha256 c6c4138ee5c86423 :: prose twin: two significant figures, direction in words
\newcommand{\MlstCmpQcdBbFoneOverFullPrMax}{up to 56\% more errors}  % experiments/FIGS/data/v2/index.json :: binary probe qcd_bb_vs_cc: (1 - AUC) of FLAV_F1 over L188, run by run :: sha256 c6c4138ee5c86423 :: prose twin of the run farther from 1
\newcommand{\MlstCmpQcdBbFonerOverFull}{2.124 and 2.189}  % experiments/FIGS/data/v2/index.json :: binary probe qcd_bb_vs_cc: (1 - AUC) of FLAV_F1R over L188, run by run :: sha256 c6c4138ee5c86423
\newcommand{\MlstCmpQcdBbFonerOverFullPr}{at least 2.1 times the errors}  % experiments/FIGS/data/v2/index.json :: binary probe qcd_bb_vs_cc: (1 - AUC) of FLAV_F1R over L188, run by run :: sha256 c6c4138ee5c86423 :: prose twin: two significant figures, direction in words
\newcommand{\MlstCmpQcdBbFonerOverFullPrMax}{up to 2.2 times the errors}  % experiments/FIGS/data/v2/index.json :: binary probe qcd_bb_vs_cc: (1 - AUC) of FLAV_F1R over L188, run by run :: sha256 c6c4138ee5c86423 :: prose twin of the run farther from 1
\newcommand{\MlstCmpQcdBResOverFull}{1.085\,\ensuremath{\pm}\,0.004}  % experiments/FIGS/data/v2/index.json :: binary probe qcd_b_vs_light: (1 - AUC) of L162 over L188, run by run :: sha256 c6c4138ee5c86423
\newcommand{\MlstCmpQcdBResOverFullPr}{8.5\% more errors}  % experiments/FIGS/data/v2/index.json :: binary probe qcd_b_vs_light: (1 - AUC) of L162 over L188, run by run :: sha256 c6c4138ee5c86423 :: prose twin: two significant figures, direction in words
\newcommand{\MlstCmpQcdBMidOverFull}{1.081\,\ensuremath{\pm}\,0.008}  % experiments/FIGS/data/v2/index.json :: binary probe qcd_b_vs_light: (1 - AUC) of R42_Q1 over L188, run by run :: sha256 c6c4138ee5c86423
\newcommand{\MlstCmpQcdBMidOverFullPr}{8.1\% more errors}  % experiments/FIGS/data/v2/index.json :: binary probe qcd_b_vs_light: (1 - AUC) of R42_Q1 over L188, run by run :: sha256 c6c4138ee5c86423 :: prose twin: two significant figures, direction in words
\newcommand{\MlstCmpQcdBCoarseOverFull}{1.54\,\ensuremath{\pm}\,0.07}  % experiments/FIGS/data/v2/index.json :: binary probe qcd_b_vs_light: (1 - AUC) of R16_Q1 over L188, run by run :: sha256 c6c4138ee5c86423
\newcommand{\MlstCmpQcdBCoarseOverFullPr}{54\% more errors}  % experiments/FIGS/data/v2/index.json :: binary probe qcd_b_vs_light: (1 - AUC) of R16_Q1 over L188, run by run :: sha256 c6c4138ee5c86423 :: prose twin: two significant figures, direction in words
\newcommand{\MlstCmpQcdBFoneOverFzero}{0.778 and 0.835}  % experiments/FIGS/data/v2/index.json :: binary probe qcd_b_vs_light: (1 - AUC) of FLAV_F1 over FLAV_F0, run by run :: sha256 c6c4138ee5c86423
\newcommand{\MlstCmpQcdBFoneOverFzeroPr}{at least 17\% fewer errors}  % experiments/FIGS/data/v2/index.json :: binary probe qcd_b_vs_light: (1 - AUC) of FLAV_F1 over FLAV_F0, run by run :: sha256 c6c4138ee5c86423 :: prose twin: two significant figures, direction in words
\newcommand{\MlstCmpQcdBFoneOverFzeroPrMax}{up to 22\% fewer errors}  % experiments/FIGS/data/v2/index.json :: binary probe qcd_b_vs_light: (1 - AUC) of FLAV_F1 over FLAV_F0, run by run :: sha256 c6c4138ee5c86423 :: prose twin of the run farther from 1
\newcommand{\MlstCmpQcdBFonerOverFone}{1.054 and 1.061}  % experiments/FIGS/data/v2/index.json :: binary probe qcd_b_vs_light: (1 - AUC) of FLAV_F1R over FLAV_F1, run by run :: sha256 c6c4138ee5c86423
\newcommand{\MlstCmpQcdBFonerOverFonePr}{equal within 7\%}  % experiments/FIGS/data/v2/index.json :: binary probe qcd_b_vs_light: (1 - AUC) of FLAV_F1R over FLAV_F1, run by run :: sha256 c6c4138ee5c86423 :: prose twin: two significant figures, direction in words
\newcommand{\MlstCmpQcdBFonerOverFonePrMax}{equal within 7\%}  % experiments/FIGS/data/v2/index.json :: binary probe qcd_b_vs_light: (1 - AUC) of FLAV_F1R over FLAV_F1, run by run :: sha256 c6c4138ee5c86423 :: prose twin of the run farther from 1
\newcommand{\MlstCmpQcdBFonerOverFzero}{0.820 and 0.885}  % experiments/FIGS/data/v2/index.json :: binary probe qcd_b_vs_light: (1 - AUC) of FLAV_F1R over FLAV_F0, run by run :: sha256 c6c4138ee5c86423
\newcommand{\MlstCmpQcdBFonerOverFzeroPr}{at least 11\% fewer errors}  % experiments/FIGS/data/v2/index.json :: binary probe qcd_b_vs_light: (1 - AUC) of FLAV_F1R over FLAV_F0, run by run :: sha256 c6c4138ee5c86423 :: prose twin: two significant figures, direction in words
\newcommand{\MlstCmpQcdBFonerOverFzeroPrMax}{up to 18\% fewer errors}  % experiments/FIGS/data/v2/index.json :: binary probe qcd_b_vs_light: (1 - AUC) of FLAV_F1R over FLAV_F0, run by run :: sha256 c6c4138ee5c86423 :: prose twin of the run farther from 1
\newcommand{\MlstCmpQcdBFzeroOverCoarse}{0.962 and 0.981}  % experiments/FIGS/data/v2/index.json :: binary probe qcd_b_vs_light: (1 - AUC) of FLAV_F0 over R16_Q1, run by run :: sha256 c6c4138ee5c86423
\newcommand{\MlstCmpQcdBFzeroOverCoarsePr}{equal within 4\%}  % experiments/FIGS/data/v2/index.json :: binary probe qcd_b_vs_light: (1 - AUC) of FLAV_F0 over R16_Q1, run by run :: sha256 c6c4138ee5c86423 :: prose twin: two significant figures, direction in words
\newcommand{\MlstCmpQcdBFzeroOverCoarsePrMax}{equal within 4\%}  % experiments/FIGS/data/v2/index.json :: binary probe qcd_b_vs_light: (1 - AUC) of FLAV_F0 over R16_Q1, run by run :: sha256 c6c4138ee5c86423 :: prose twin of the run farther from 1
\newcommand{\MlstCmpQcdBFoneOverCoarse}{0.749 and 0.819}  % experiments/FIGS/data/v2/index.json :: binary probe qcd_b_vs_light: (1 - AUC) of FLAV_F1 over R16_Q1, run by run :: sha256 c6c4138ee5c86423
\newcommand{\MlstCmpQcdBFoneOverCoarsePr}{at least 18\% fewer errors}  % experiments/FIGS/data/v2/index.json :: binary probe qcd_b_vs_light: (1 - AUC) of FLAV_F1 over R16_Q1, run by run :: sha256 c6c4138ee5c86423 :: prose twin: two significant figures, direction in words
\newcommand{\MlstCmpQcdBFoneOverCoarsePrMax}{up to 25\% fewer errors}  % experiments/FIGS/data/v2/index.json :: binary probe qcd_b_vs_light: (1 - AUC) of FLAV_F1 over R16_Q1, run by run :: sha256 c6c4138ee5c86423 :: prose twin of the run farther from 1
\newcommand{\MlstCmpQcdBFonerOverCoarse}{0.789 and 0.869}  % experiments/FIGS/data/v2/index.json :: binary probe qcd_b_vs_light: (1 - AUC) of FLAV_F1R over R16_Q1, run by run :: sha256 c6c4138ee5c86423
\newcommand{\MlstCmpQcdBFonerOverCoarsePr}{at least 13\% fewer errors}  % experiments/FIGS/data/v2/index.json :: binary probe qcd_b_vs_light: (1 - AUC) of FLAV_F1R over R16_Q1, run by run :: sha256 c6c4138ee5c86423 :: prose twin: two significant figures, direction in words
\newcommand{\MlstCmpQcdBFonerOverCoarsePrMax}{up to 21\% fewer errors}  % experiments/FIGS/data/v2/index.json :: binary probe qcd_b_vs_light: (1 - AUC) of FLAV_F1R over R16_Q1, run by run :: sha256 c6c4138ee5c86423 :: prose twin of the run farther from 1
\newcommand{\MlstCmpQcdBFzeroOverFull}{1.56 and 1.45}  % experiments/FIGS/data/v2/index.json :: binary probe qcd_b_vs_light: (1 - AUC) of FLAV_F0 over L188, run by run :: sha256 c6c4138ee5c86423
\newcommand{\MlstCmpQcdBFzeroOverFullPr}{at least 45\% more errors}  % experiments/FIGS/data/v2/index.json :: binary probe qcd_b_vs_light: (1 - AUC) of FLAV_F0 over L188, run by run :: sha256 c6c4138ee5c86423 :: prose twin: two significant figures, direction in words
\newcommand{\MlstCmpQcdBFzeroOverFullPrMax}{up to 56\% more errors}  % experiments/FIGS/data/v2/index.json :: binary probe qcd_b_vs_light: (1 - AUC) of FLAV_F0 over L188, run by run :: sha256 c6c4138ee5c86423 :: prose twin of the run farther from 1
\newcommand{\MlstCmpQcdBFoneOverFull}{1.212 and 1.213}  % experiments/FIGS/data/v2/index.json :: binary probe qcd_b_vs_light: (1 - AUC) of FLAV_F1 over L188, run by run :: sha256 c6c4138ee5c86423
\newcommand{\MlstCmpQcdBFoneOverFullPr}{at least 21\% more errors}  % experiments/FIGS/data/v2/index.json :: binary probe qcd_b_vs_light: (1 - AUC) of FLAV_F1 over L188, run by run :: sha256 c6c4138ee5c86423 :: prose twin: two significant figures, direction in words
\newcommand{\MlstCmpQcdBFoneOverFullPrMax}{up to 21\% more errors}  % experiments/FIGS/data/v2/index.json :: binary probe qcd_b_vs_light: (1 - AUC) of FLAV_F1 over L188, run by run :: sha256 c6c4138ee5c86423 :: prose twin of the run farther from 1
\newcommand{\MlstCmpQcdBFonerOverFull}{1.277 and 1.286}  % experiments/FIGS/data/v2/index.json :: binary probe qcd_b_vs_light: (1 - AUC) of FLAV_F1R over L188, run by run :: sha256 c6c4138ee5c86423
\newcommand{\MlstCmpQcdBFonerOverFullPr}{at least 28\% more errors}  % experiments/FIGS/data/v2/index.json :: binary probe qcd_b_vs_light: (1 - AUC) of FLAV_F1R over L188, run by run :: sha256 c6c4138ee5c86423 :: prose twin: two significant figures, direction in words
\newcommand{\MlstCmpQcdBFonerOverFullPrMax}{up to 29\% more errors}  % experiments/FIGS/data/v2/index.json :: binary probe qcd_b_vs_light: (1 - AUC) of FLAV_F1R over L188, run by run :: sha256 c6c4138ee5c86423 :: prose twin of the run farther from 1
\newcommand{\MlstCmpQcdCResOverFull}{1.061\,\ensuremath{\pm}\,0.005}  % experiments/FIGS/data/v2/index.json :: binary probe qcd_c_vs_light: (1 - AUC) of L162 over L188, run by run :: sha256 c6c4138ee5c86423
\newcommand{\MlstCmpQcdCResOverFullPr}{6.1\% more errors}  % experiments/FIGS/data/v2/index.json :: binary probe qcd_c_vs_light: (1 - AUC) of L162 over L188, run by run :: sha256 c6c4138ee5c86423 :: prose twin: two significant figures, direction in words
\newcommand{\MlstCmpQcdCMidOverFull}{1.072\,\ensuremath{\pm}\,0.008}  % experiments/FIGS/data/v2/index.json :: binary probe qcd_c_vs_light: (1 - AUC) of R42_Q1 over L188, run by run :: sha256 c6c4138ee5c86423
\newcommand{\MlstCmpQcdCMidOverFullPr}{7.2\% more errors}  % experiments/FIGS/data/v2/index.json :: binary probe qcd_c_vs_light: (1 - AUC) of R42_Q1 over L188, run by run :: sha256 c6c4138ee5c86423 :: prose twin: two significant figures, direction in words
\newcommand{\MlstCmpQcdCCoarseOverFull}{1.407\,\ensuremath{\pm}\,0.021}  % experiments/FIGS/data/v2/index.json :: binary probe qcd_c_vs_light: (1 - AUC) of R16_Q1 over L188, run by run :: sha256 c6c4138ee5c86423
\newcommand{\MlstCmpQcdCCoarseOverFullPr}{41\% more errors}  % experiments/FIGS/data/v2/index.json :: binary probe qcd_c_vs_light: (1 - AUC) of R16_Q1 over L188, run by run :: sha256 c6c4138ee5c86423 :: prose twin: two significant figures, direction in words
\newcommand{\MlstCmpQcdCFoneOverFzero}{0.915 and 0.944}  % experiments/FIGS/data/v2/index.json :: binary probe qcd_c_vs_light: (1 - AUC) of FLAV_F1 over FLAV_F0, run by run :: sha256 c6c4138ee5c86423
\newcommand{\MlstCmpQcdCFoneOverFzeroPr}{at least 5.6\% fewer errors}  % experiments/FIGS/data/v2/index.json :: binary probe qcd_c_vs_light: (1 - AUC) of FLAV_F1 over FLAV_F0, run by run :: sha256 c6c4138ee5c86423 :: prose twin: two significant figures, direction in words
\newcommand{\MlstCmpQcdCFoneOverFzeroPrMax}{up to 8.5\% fewer errors}  % experiments/FIGS/data/v2/index.json :: binary probe qcd_c_vs_light: (1 - AUC) of FLAV_F1 over FLAV_F0, run by run :: sha256 c6c4138ee5c86423 :: prose twin of the run farther from 1
\newcommand{\MlstCmpQcdCFonerOverFone}{0.9535 and 0.9520}  % experiments/FIGS/data/v2/index.json :: binary probe qcd_c_vs_light: (1 - AUC) of FLAV_F1R over FLAV_F1, run by run :: sha256 c6c4138ee5c86423
\newcommand{\MlstCmpQcdCFonerOverFonePr}{at least 4.6\% fewer errors}  % experiments/FIGS/data/v2/index.json :: binary probe qcd_c_vs_light: (1 - AUC) of FLAV_F1R over FLAV_F1, run by run :: sha256 c6c4138ee5c86423 :: prose twin: two significant figures, direction in words
\newcommand{\MlstCmpQcdCFonerOverFonePrMax}{up to 4.8\% fewer errors}  % experiments/FIGS/data/v2/index.json :: binary probe qcd_c_vs_light: (1 - AUC) of FLAV_F1R over FLAV_F1, run by run :: sha256 c6c4138ee5c86423 :: prose twin of the run farther from 1
\newcommand{\MlstCmpQcdCFonerOverFzero}{0.872 and 0.899}  % experiments/FIGS/data/v2/index.json :: binary probe qcd_c_vs_light: (1 - AUC) of FLAV_F1R over FLAV_F0, run by run :: sha256 c6c4138ee5c86423
\newcommand{\MlstCmpQcdCFonerOverFzeroPr}{at least 10\% fewer errors}  % experiments/FIGS/data/v2/index.json :: binary probe qcd_c_vs_light: (1 - AUC) of FLAV_F1R over FLAV_F0, run by run :: sha256 c6c4138ee5c86423 :: prose twin: two significant figures, direction in words
\newcommand{\MlstCmpQcdCFonerOverFzeroPrMax}{up to 13\% fewer errors}  % experiments/FIGS/data/v2/index.json :: binary probe qcd_c_vs_light: (1 - AUC) of FLAV_F1R over FLAV_F0, run by run :: sha256 c6c4138ee5c86423 :: prose twin of the run farther from 1
\newcommand{\MlstCmpQcdCFzeroOverCoarse}{0.952 and 0.934}  % experiments/FIGS/data/v2/index.json :: binary probe qcd_c_vs_light: (1 - AUC) of FLAV_F0 over R16_Q1, run by run :: sha256 c6c4138ee5c86423
\newcommand{\MlstCmpQcdCFzeroOverCoarsePr}{at least 4.8\% fewer errors}  % experiments/FIGS/data/v2/index.json :: binary probe qcd_c_vs_light: (1 - AUC) of FLAV_F0 over R16_Q1, run by run :: sha256 c6c4138ee5c86423 :: prose twin: two significant figures, direction in words
\newcommand{\MlstCmpQcdCFzeroOverCoarsePrMax}{up to 6.6\% fewer errors}  % experiments/FIGS/data/v2/index.json :: binary probe qcd_c_vs_light: (1 - AUC) of FLAV_F0 over R16_Q1, run by run :: sha256 c6c4138ee5c86423 :: prose twin of the run farther from 1
\newcommand{\MlstCmpQcdCFoneOverCoarse}{0.871 and 0.882}  % experiments/FIGS/data/v2/index.json :: binary probe qcd_c_vs_light: (1 - AUC) of FLAV_F1 over R16_Q1, run by run :: sha256 c6c4138ee5c86423
\newcommand{\MlstCmpQcdCFoneOverCoarsePr}{at least 12\% fewer errors}  % experiments/FIGS/data/v2/index.json :: binary probe qcd_c_vs_light: (1 - AUC) of FLAV_F1 over R16_Q1, run by run :: sha256 c6c4138ee5c86423 :: prose twin: two significant figures, direction in words
\newcommand{\MlstCmpQcdCFoneOverCoarsePrMax}{up to 13\% fewer errors}  % experiments/FIGS/data/v2/index.json :: binary probe qcd_c_vs_light: (1 - AUC) of FLAV_F1 over R16_Q1, run by run :: sha256 c6c4138ee5c86423 :: prose twin of the run farther from 1
\newcommand{\MlstCmpQcdCFonerOverCoarse}{0.8301 and 0.8396}  % experiments/FIGS/data/v2/index.json :: binary probe qcd_c_vs_light: (1 - AUC) of FLAV_F1R over R16_Q1, run by run :: sha256 c6c4138ee5c86423
\newcommand{\MlstCmpQcdCFonerOverCoarsePr}{at least 16\% fewer errors}  % experiments/FIGS/data/v2/index.json :: binary probe qcd_c_vs_light: (1 - AUC) of FLAV_F1R over R16_Q1, run by run :: sha256 c6c4138ee5c86423 :: prose twin: two significant figures, direction in words
\newcommand{\MlstCmpQcdCFonerOverCoarsePrMax}{up to 17\% fewer errors}  % experiments/FIGS/data/v2/index.json :: binary probe qcd_c_vs_light: (1 - AUC) of FLAV_F1R over R16_Q1, run by run :: sha256 c6c4138ee5c86423 :: prose twin of the run farther from 1
\newcommand{\MlstCmpQcdCFzeroOverFull}{1.352 and 1.293}  % experiments/FIGS/data/v2/index.json :: binary probe qcd_c_vs_light: (1 - AUC) of FLAV_F0 over L188, run by run :: sha256 c6c4138ee5c86423
\newcommand{\MlstCmpQcdCFzeroOverFullPr}{at least 29\% more errors}  % experiments/FIGS/data/v2/index.json :: binary probe qcd_c_vs_light: (1 - AUC) of FLAV_F0 over L188, run by run :: sha256 c6c4138ee5c86423 :: prose twin: two significant figures, direction in words
\newcommand{\MlstCmpQcdCFzeroOverFullPrMax}{up to 35\% more errors}  % experiments/FIGS/data/v2/index.json :: binary probe qcd_c_vs_light: (1 - AUC) of FLAV_F0 over L188, run by run :: sha256 c6c4138ee5c86423 :: prose twin of the run farther from 1
\newcommand{\MlstCmpQcdCFoneOverFull}{1.237 and 1.220}  % experiments/FIGS/data/v2/index.json :: binary probe qcd_c_vs_light: (1 - AUC) of FLAV_F1 over L188, run by run :: sha256 c6c4138ee5c86423
\newcommand{\MlstCmpQcdCFoneOverFullPr}{at least 22\% more errors}  % experiments/FIGS/data/v2/index.json :: binary probe qcd_c_vs_light: (1 - AUC) of FLAV_F1 over L188, run by run :: sha256 c6c4138ee5c86423 :: prose twin: two significant figures, direction in words
\newcommand{\MlstCmpQcdCFoneOverFullPrMax}{up to 24\% more errors}  % experiments/FIGS/data/v2/index.json :: binary probe qcd_c_vs_light: (1 - AUC) of FLAV_F1 over L188, run by run :: sha256 c6c4138ee5c86423 :: prose twin of the run farther from 1
\newcommand{\MlstCmpQcdCFonerOverFull}{1.179 and 1.162}  % experiments/FIGS/data/v2/index.json :: binary probe qcd_c_vs_light: (1 - AUC) of FLAV_F1R over L188, run by run :: sha256 c6c4138ee5c86423
\newcommand{\MlstCmpQcdCFonerOverFullPr}{at least 16\% more errors}  % experiments/FIGS/data/v2/index.json :: binary probe qcd_c_vs_light: (1 - AUC) of FLAV_F1R over L188, run by run :: sha256 c6c4138ee5c86423 :: prose twin: two significant figures, direction in words
\newcommand{\MlstCmpQcdCFonerOverFullPrMax}{up to 18\% more errors}  % experiments/FIGS/data/v2/index.json :: binary probe qcd_c_vs_light: (1 - AUC) of FLAV_F1R over L188, run by run :: sha256 c6c4138ee5c86423 :: prose twin of the run farther from 1
\newcommand{\MlstLpEemmAucFull}{1.0000\,\ensuremath{\pm}\,\ensuremath{<}0.0001}  % experiments/FIGS/data/v2/index.json :: linear probe ee_vs_mm (fits_v4_sib): AUC, L188 :: sha256 c6c4138ee5c86423
\newcommand{\MlstLpEemmAucRes}{1.0000\,\ensuremath{\pm}\,\ensuremath{<}0.0001}  % experiments/FIGS/data/v2/index.json :: linear probe ee_vs_mm (fits_v4_sib): AUC, L162 :: sha256 c6c4138ee5c86423
\newcommand{\MlstLpEemmAucMid}{0.9610\,\ensuremath{\pm}\,0.0030}  % experiments/FIGS/data/v2/index.json :: linear probe ee_vs_mm (fits_v4_sib): AUC, R42_Q1 :: sha256 c6c4138ee5c86423
\newcommand{\MlstLpEemmAucCoarse}{0.960\,\ensuremath{\pm}\,0.020}  % experiments/FIGS/data/v2/index.json :: linear probe ee_vs_mm (fits_v4_sib): AUC, R16_Q1 :: sha256 c6c4138ee5c86423
\newcommand{\MlstLpEemmAucFzero}{1.0000 and 1.0000}  % experiments/FIGS/data/v2/index.json :: linear probe ee_vs_mm (fits_v4_sib): AUC, FLAV_F0 :: sha256 c6c4138ee5c86423
\newcommand{\MlstLpEemmAucFone}{1.0000 and 1.0000}  % experiments/FIGS/data/v2/index.json :: linear probe ee_vs_mm (fits_v4_sib): AUC, FLAV_F1 :: sha256 c6c4138ee5c86423
\newcommand{\MlstLpEemmAucFoner}{1.0000 and 1.0000}  % experiments/FIGS/data/v2/index.json :: linear probe ee_vs_mm (fits_v4_sib): AUC, FLAV_F1R :: sha256 c6c4138ee5c86423
\newcommand{\MlstLpEemmAucInit}{0.9995\,\ensuremath{\pm}\,0.0003}  % experiments/FIGS/data/v2/index.json :: linear probe ee_vs_mm (fits_v4_sib): AUC, INIT :: sha256 c6c4138ee5c86423
\newcommand{\MlstLpEemmAucInitBn}{1.0000\,\ensuremath{\pm}\,\ensuremath{<}0.0001}  % experiments/FIGS/data/v2/index.json :: linear probe ee_vs_mm (fits_v4_sib): AUC, INIT_BN :: sha256 c6c4138ee5c86423
\newcommand{\MlstLpEemmAucRandMerged}{1.0000\,\ensuremath{\pm}\,\ensuremath{<}0.0001}  % experiments/FIGS/data/v2/index.json :: linear probe ee_vs_mm (fits_v4_sib): AUC, the random partitions that merge the pair :: sha256 c6c4138ee5c86423
\newcommand{\MlstPairEemmNTrain}{9{,}571}  % experiments/FIGS/data/v2/index.json :: fits_v4_sib ee_vs_mm: n_train :: sha256 c6c4138ee5c86423
\newcommand{\MlstPairEemmNVal}{1{,}068}  % experiments/FIGS/data/v2/index.json :: fits_v4_sib ee_vs_mm: n_val :: sha256 c6c4138ee5c86423
\newcommand{\MlstPairEemmNTestPerSide}{128{,}833--135{,}245}  % experiments/FIGS/data/v2/index.json :: fits_v4_sib ee_vs_mm: n_test_signal and n_test_background :: sha256 c6c4138ee5c86423
\newcommand{\MlstLpJcTwoAucFull}{0.9801\,\ensuremath{\pm}\,\ensuremath{<}0.0001}  % experiments/FIGS/data/v2/index.json :: linear probe JetClass-II N=1000000: AUC averaged over the classes, L188 :: sha256 c6c4138ee5c86423
\newcommand{\MlstLpJcTwoAccFull}{0.4348\,\ensuremath{\pm}\,0.0003}  % experiments/FIGS/data/v2/index.json :: linear probe JetClass-II N=1000000: accuracy, L188 :: sha256 c6c4138ee5c86423
\newcommand{\MlstLpJcTwoAucRes}{0.9796\,\ensuremath{\pm}\,0.0001}  % experiments/FIGS/data/v2/index.json :: linear probe JetClass-II N=1000000: AUC averaged over the classes, L162 :: sha256 c6c4138ee5c86423
\newcommand{\MlstLpJcTwoAccRes}{0.4328\,\ensuremath{\pm}\,0.0008}  % experiments/FIGS/data/v2/index.json :: linear probe JetClass-II N=1000000: accuracy, L162 :: sha256 c6c4138ee5c86423
\newcommand{\MlstLpJcTwoAucMid}{0.9756\,\ensuremath{\pm}\,0.0002}  % experiments/FIGS/data/v2/index.json :: linear probe JetClass-II N=1000000: AUC averaged over the classes, R42_Q1 :: sha256 c6c4138ee5c86423
\newcommand{\MlstLpJcTwoAccMid}{0.361\,\ensuremath{\pm}\,0.005}  % experiments/FIGS/data/v2/index.json :: linear probe JetClass-II N=1000000: accuracy, R42_Q1 :: sha256 c6c4138ee5c86423
\newcommand{\MlstLpJcTwoAucCoarse}{0.9691\,\ensuremath{\pm}\,0.0004}  % experiments/FIGS/data/v2/index.json :: linear probe JetClass-II N=1000000: AUC averaged over the classes, R16_Q1 :: sha256 c6c4138ee5c86423
\newcommand{\MlstLpJcTwoAccCoarse}{0.334\,\ensuremath{\pm}\,0.007}  % experiments/FIGS/data/v2/index.json :: linear probe JetClass-II N=1000000: accuracy, R16_Q1 :: sha256 c6c4138ee5c86423
\newcommand{\MlstLpJcTwoAucResMass}{0.9791 and 0.9794}  % experiments/FIGS/data/v2/index.json :: linear probe JetClass-II N=1000000: AUC averaged over the classes, L162_MASS :: sha256 c6c4138ee5c86423
\newcommand{\MlstLpJcTwoAccResMass}{0.4271 and 0.4305}  % experiments/FIGS/data/v2/index.json :: linear probe JetClass-II N=1000000: accuracy, L162_MASS :: sha256 c6c4138ee5c86423
\newcommand{\MlstLpJcTwoAucCoarseMass}{0.9679 and 0.9696}  % experiments/FIGS/data/v2/index.json :: linear probe JetClass-II N=1000000: AUC averaged over the classes, R16_Q1_MASS :: sha256 c6c4138ee5c86423
\newcommand{\MlstLpJcTwoAccCoarseMass}{0.312 and 0.350}  % experiments/FIGS/data/v2/index.json :: linear probe JetClass-II N=1000000: accuracy, R16_Q1_MASS :: sha256 c6c4138ee5c86423
\newcommand{\MlstLpJcTwoAucFzero}{0.9717 and 0.9710}  % experiments/FIGS/data/v2/index.json :: linear probe JetClass-II N=1000000: AUC averaged over the classes, FLAV_F0 :: sha256 c6c4138ee5c86423
\newcommand{\MlstLpJcTwoAccFzero}{0.3769 and 0.3723}  % experiments/FIGS/data/v2/index.json :: linear probe JetClass-II N=1000000: accuracy, FLAV_F0 :: sha256 c6c4138ee5c86423
\newcommand{\MlstLpJcTwoAucFone}{0.9758 and 0.9759}  % experiments/FIGS/data/v2/index.json :: linear probe JetClass-II N=1000000: AUC averaged over the classes, FLAV_F1 :: sha256 c6c4138ee5c86423
\newcommand{\MlstLpJcTwoAccFone}{0.3945 and 0.3963}  % experiments/FIGS/data/v2/index.json :: linear probe JetClass-II N=1000000: accuracy, FLAV_F1 :: sha256 c6c4138ee5c86423
\newcommand{\MlstLpJcTwoAucFoner}{0.9760 and 0.9760}  % experiments/FIGS/data/v2/index.json :: linear probe JetClass-II N=1000000: AUC averaged over the classes, FLAV_F1R :: sha256 c6c4138ee5c86423
\newcommand{\MlstLpJcTwoAccFoner}{0.4071 and 0.4089}  % experiments/FIGS/data/v2/index.json :: linear probe JetClass-II N=1000000: accuracy, FLAV_F1R :: sha256 c6c4138ee5c86423
\newcommand{\MlstLpJcTwoAucInit}{0.8121\,\ensuremath{\pm}\,0.0019}  % experiments/FIGS/data/v2/index.json :: linear probe JetClass-II N=1000000: AUC averaged over the classes, INIT :: sha256 c6c4138ee5c86423
\newcommand{\MlstLpJcTwoAccInit}{0.0724\,\ensuremath{\pm}\,0.0015}  % experiments/FIGS/data/v2/index.json :: linear probe JetClass-II N=1000000: accuracy, INIT :: sha256 c6c4138ee5c86423
\newcommand{\MlstLpJcTwoAucInitBn}{0.866\,\ensuremath{\pm}\,0.004}  % experiments/FIGS/data/v2/index.json :: linear probe JetClass-II N=1000000: AUC averaged over the classes, INIT_BN :: sha256 c6c4138ee5c86423
\newcommand{\MlstLpJcTwoAccInitBn}{0.1017\,\ensuremath{\pm}\,0.0017}  % experiments/FIGS/data/v2/index.json :: linear probe JetClass-II N=1000000: accuracy, INIT_BN :: sha256 c6c4138ee5c86423
\newcommand{\MlstLpJcTwoAucRand}{0.9777\,\ensuremath{\pm}\,0.0002}  % experiments/FIGS/data/v2/index.json :: linear probe JetClass-II N=1000000: AUC averaged over the classes, the random partitions :: sha256 c6c4138ee5c86423
\newcommand{\MlstLpJcOneAccInitBn}{0.475\,\ensuremath{\pm}\,0.005}  % experiments/FIGS/data/v2/index.json :: linear probe JetClass N=1000000: accuracy, INIT_BN :: sha256 c6c4138ee5c86423
\newcommand{\MlstBnMaxDev}{0.2\%}  % experiments/FIGS/data/v2/index.json :: largest |r(best70_bn) / r(best70) - 1| of a JetClass-II ratio of 1 - AUC to the 188-class run, per cent :: sha256 c6c4138ee5c86423
\newcommand{\MlstBnMaxDevPair}{0.7\%}  % experiments/FIGS/data/v2/index.json :: largest |r(best70_bn) / r(best70) - 1| of a ratio of 1 - AUC on a class pair to the 188-class run, levels of the hierarchy, mass models and controlled partitions, per cent :: sha256 c6c4138ee5c86423 :: pairs with a check: bc_vs_bq_cs, bb_vs_cc, bbqq_vs_ccqq, qcd_bb_vs_cc, qcd_b_vs_light, qcd_c_vs_light; without: ee_vs_mm
\newcommand{\MlstGpuTypeBound}{1.4\%}  % experiments/FIGS/data/v2/index.json :: largest |label-set ratio of 1 - AUC of index 3 / mean of the other indices - 1|, JetClass-II largest set, per cent :: sha256 c6c4138ee5c86423 :: mtx-l188-s3 trained on NVIDIA GeForce RTX 3090; largest for R16_Q1; label sets L162 0.09%, R42_Q1 0.70%, R16_Q1 1.37%
\newcommand{\MlstGpuOtherSeedsDiff}{1.9\%}  % experiments/FIGS/data/v2/index.json :: largest (max - min) / mean of a label-set ratio of 1 - AUC over the indices other than 3, JetClass-II largest set, per cent :: sha256 c6c4138ee5c86423 :: label sets L162 0.90%, R42_Q1 1.14%, R16_Q1 1.93%
\newcommand{\MlstPerClassResNDoubled}{0}  % experiments/FIGS/data/v2/index.json :: JetClass-II classes whose mean ratio of 1 - AUC (L162 over 188) is at least 2 (count) :: sha256 c6c4138ee5c86423
\newcommand{\MlstPerClassResSeedSdMedian}{0.011}  % experiments/FIGS/data/v2/index.json :: median over classes of the SD over runs of the per-class ratio, L162 over 188 :: sha256 c6c4138ee5c86423
\newcommand{\MlstPerClassResTopClass}{QCD, \ensuremath{ccss}}  % experiments/FIGS/data/v2/index.json :: class with the largest mean ratio, L162 :: sha256 c6c4138ee5c86423
\newcommand{\MlstPerClassResTopRatio}{1.129\,\ensuremath{\pm}\,0.010}  % experiments/FIGS/data/v2/index.json :: per-class ratio of 1 - AUC of that class, L162 over 188, run by run :: sha256 c6c4138ee5c86423
\newcommand{\MlstPerClassResTopRatioPr}{13\% more errors}  % experiments/FIGS/data/v2/index.json :: per-class ratio of 1 - AUC of that class, L162 over 188, run by run :: sha256 c6c4138ee5c86423 :: prose twin: two significant figures, direction in words
\newcommand{\MlstPerClassResLeptonN}{0}  % experiments/FIGS/data/v2/index.json :: resonant classes sharing their L162 group with a decay that differs only in lepton flavor (count) :: sha256 c6c4138ee5c86423
\newcommand{\MlstPerClassResLeptonMedian}{--}  % experiments/FIGS/data/v2/index.json :: median over those classes of the mean per-class ratio, L162 over 188 :: sha256 c6c4138ee5c86423
\newcommand{\MlstPerClassResLeptonMedianPr}{\pending{no run yet}}  % experiments/FIGS/data/v2/index.json :: median over those classes of the mean per-class ratio, L162 over 188 :: sha256 c6c4138ee5c86423 :: prose twin: two significant figures, direction in words
\newcommand{\MlstPerClassResLeptonShare}{--}  % experiments/FIGS/data/v2/index.json :: their share of the summed ln ratio over the resonant classes, L162 :: sha256 c6c4138ee5c86423
\newcommand{\MlstPerClassResQuarkN}{0}  % experiments/FIGS/data/v2/index.json :: resonant classes sharing their L162 group with a decay that differs only in quark flavor (count) :: sha256 c6c4138ee5c86423
\newcommand{\MlstPerClassResQuarkMedian}{--}  % experiments/FIGS/data/v2/index.json :: median over those classes of the mean per-class ratio, L162 over 188 :: sha256 c6c4138ee5c86423
\newcommand{\MlstPerClassResQuarkMedianPr}{\pending{no run yet}}  % experiments/FIGS/data/v2/index.json :: median over those classes of the mean per-class ratio, L162 over 188 :: sha256 c6c4138ee5c86423 :: prose twin: two significant figures, direction in words
\newcommand{\MlstPerClassResQuarkShare}{--}  % experiments/FIGS/data/v2/index.json :: their share of the summed ln ratio over the resonant classes, L162 :: sha256 c6c4138ee5c86423
\newcommand{\MlstPerClassResQuarkLeptonN}{0}  % experiments/FIGS/data/v2/index.json :: resonant classes counted in both the quark- and the lepton-flavor set, L162 (count) :: sha256 c6c4138ee5c86423
\newcommand{\MlstPerClassMidNDoubled}{40}  % experiments/FIGS/data/v2/index.json :: JetClass-II classes whose mean ratio of 1 - AUC (R42_Q1 over 188) is at least 2 (count) :: sha256 c6c4138ee5c86423
\newcommand{\MlstPerClassMidSeedSdMedian}{0.023}  % experiments/FIGS/data/v2/index.json :: median over classes of the SD over runs of the per-class ratio, R42_Q1 over 188 :: sha256 c6c4138ee5c86423
\newcommand{\MlstPerClassMidTopClass}{\ensuremath{X\to YY\to bb\mu\mu}}  % experiments/FIGS/data/v2/index.json :: class with the largest mean ratio, R42_Q1 :: sha256 c6c4138ee5c86423
\newcommand{\MlstPerClassMidTopRatio}{4.06\,\ensuremath{\pm}\,0.14}  % experiments/FIGS/data/v2/index.json :: per-class ratio of 1 - AUC of that class, R42_Q1 over 188, run by run :: sha256 c6c4138ee5c86423
\newcommand{\MlstPerClassMidTopRatioPr}{4.1 times the errors}  % experiments/FIGS/data/v2/index.json :: per-class ratio of 1 - AUC of that class, R42_Q1 over 188, run by run :: sha256 c6c4138ee5c86423 :: prose twin: two significant figures, direction in words
\newcommand{\MlstPerClassMidLeptonN}{78}  % experiments/FIGS/data/v2/index.json :: resonant classes sharing their R42_Q1 group with a decay that differs only in lepton flavor (count) :: sha256 c6c4138ee5c86423
\newcommand{\MlstPerClassMidLeptonMedian}{2.02}  % experiments/FIGS/data/v2/index.json :: median over those classes of the mean per-class ratio, R42_Q1 over 188 :: sha256 c6c4138ee5c86423
\newcommand{\MlstPerClassMidLeptonMedianPr}{2.0 times the errors}  % experiments/FIGS/data/v2/index.json :: median over those classes of the mean per-class ratio, R42_Q1 over 188 :: sha256 c6c4138ee5c86423 :: prose twin: two significant figures, direction in words
\newcommand{\MlstPerClassMidLeptonShare}{90\%}  % experiments/FIGS/data/v2/index.json :: their share of the summed ln ratio over the resonant classes, R42_Q1 :: sha256 c6c4138ee5c86423
\newcommand{\MlstPerClassMidQuarkN}{111}  % experiments/FIGS/data/v2/index.json :: resonant classes sharing their R42_Q1 group with a decay that differs only in quark flavor (count) :: sha256 c6c4138ee5c86423
\newcommand{\MlstPerClassMidQuarkMedian}{1.14}  % experiments/FIGS/data/v2/index.json :: median over those classes of the mean per-class ratio, R42_Q1 over 188 :: sha256 c6c4138ee5c86423
\newcommand{\MlstPerClassMidQuarkMedianPr}{14\% more errors}  % experiments/FIGS/data/v2/index.json :: median over those classes of the mean per-class ratio, R42_Q1 over 188 :: sha256 c6c4138ee5c86423 :: prose twin: two significant figures, direction in words
\newcommand{\MlstPerClassMidQuarkShare}{53\%}  % experiments/FIGS/data/v2/index.json :: their share of the summed ln ratio over the resonant classes, R42_Q1 :: sha256 c6c4138ee5c86423
\newcommand{\MlstPerClassMidQuarkLeptonN}{44}  % experiments/FIGS/data/v2/index.json :: resonant classes counted in both the quark- and the lepton-flavor set, R42_Q1 (count) :: sha256 c6c4138ee5c86423
\newcommand{\MlstPerClassCoarseNDoubled}{60}  % experiments/FIGS/data/v2/index.json :: JetClass-II classes whose mean ratio of 1 - AUC (R16_Q1 over 188) is at least 2 (count) :: sha256 c6c4138ee5c86423
\newcommand{\MlstPerClassCoarseSeedSdMedian}{0.055}  % experiments/FIGS/data/v2/index.json :: median over classes of the SD over runs of the per-class ratio, R16_Q1 over 188 :: sha256 c6c4138ee5c86423
\newcommand{\MlstPerClassCoarseTopClass}{\ensuremath{X\to YY\to bbee}}  % experiments/FIGS/data/v2/index.json :: class with the largest mean ratio, R16_Q1 :: sha256 c6c4138ee5c86423
\newcommand{\MlstPerClassCoarseTopRatio}{5.2\,\ensuremath{\pm}\,0.4}  % experiments/FIGS/data/v2/index.json :: per-class ratio of 1 - AUC of that class, R16_Q1 over 188, run by run :: sha256 c6c4138ee5c86423
\newcommand{\MlstPerClassCoarseTopRatioPr}{5.2 times the errors}  % experiments/FIGS/data/v2/index.json :: per-class ratio of 1 - AUC of that class, R16_Q1 over 188, run by run :: sha256 c6c4138ee5c86423 :: prose twin: two significant figures, direction in words
\newcommand{\MlstPerClassCoarseLeptonN}{78}  % experiments/FIGS/data/v2/index.json :: resonant classes sharing their R16_Q1 group with a decay that differs only in lepton flavor (count) :: sha256 c6c4138ee5c86423
\newcommand{\MlstPerClassCoarseLeptonMedian}{2.27}  % experiments/FIGS/data/v2/index.json :: median over those classes of the mean per-class ratio, R16_Q1 over 188 :: sha256 c6c4138ee5c86423
\newcommand{\MlstPerClassCoarseLeptonMedianPr}{2.3 times the errors}  % experiments/FIGS/data/v2/index.json :: median over those classes of the mean per-class ratio, R16_Q1 over 188 :: sha256 c6c4138ee5c86423 :: prose twin: two significant figures, direction in words
\newcommand{\MlstPerClassCoarseLeptonShare}{68\%}  % experiments/FIGS/data/v2/index.json :: their share of the summed ln ratio over the resonant classes, R16_Q1 :: sha256 c6c4138ee5c86423
\newcommand{\MlstPerClassCoarseQuarkN}{141}  % experiments/FIGS/data/v2/index.json :: resonant classes sharing their R16_Q1 group with a decay that differs only in quark flavor (count) :: sha256 c6c4138ee5c86423
\newcommand{\MlstPerClassCoarseQuarkMedian}{1.65}  % experiments/FIGS/data/v2/index.json :: median over those classes of the mean per-class ratio, R16_Q1 over 188 :: sha256 c6c4138ee5c86423
\newcommand{\MlstPerClassCoarseQuarkMedianPr}{65\% more errors}  % experiments/FIGS/data/v2/index.json :: median over those classes of the mean per-class ratio, R16_Q1 over 188 :: sha256 c6c4138ee5c86423 :: prose twin: two significant figures, direction in words
\newcommand{\MlstPerClassCoarseQuarkShare}{87\%}  % experiments/FIGS/data/v2/index.json :: their share of the summed ln ratio over the resonant classes, R16_Q1 :: sha256 c6c4138ee5c86423
\newcommand{\MlstPerClassCoarseQuarkLeptonN}{64}  % experiments/FIGS/data/v2/index.json :: resonant classes counted in both the quark- and the lepton-flavor set, R16_Q1 (count) :: sha256 c6c4138ee5c86423
\newcommand{\MlstPerClassRandFlavMergedOverApart}{1.010\,\ensuremath{\pm}\,0.008}  % experiments/FIGS/data/v2/index.json :: random partitions, linear probe: median per-class ratio of the classes merged with a decay that differs only in quark flavor over that of the classes kept apart from all of them, mean +- SD over the partitions :: sha256 c6c4138ee5c86423
\newcommand{\MlstPerClassRandFlavMergedOverApartPr}{equal within 1\%}  % experiments/FIGS/data/v2/index.json :: random partitions, linear probe: median per-class ratio of the classes merged with a decay that differs only in quark flavor over that of the classes kept apart from all of them, mean +- SD over the partitions :: sha256 c6c4138ee5c86423 :: prose twin: two significant figures, direction in words
\newcommand{\MlstPerClassRandLeptonMergedOverApart}{0.997\,\ensuremath{\pm}\,0.020}  % experiments/FIGS/data/v2/index.json :: random partitions, linear probe: median per-class ratio of the classes merged with a decay that differs only in lepton flavor over that of the classes kept apart from all of them, mean +- SD over the partitions :: sha256 c6c4138ee5c86423
\newcommand{\MlstPerClassRandLeptonMergedOverApartPr}{equal within 3\%}  % experiments/FIGS/data/v2/index.json :: random partitions, linear probe: median per-class ratio of the classes merged with a decay that differs only in lepton flavor over that of the classes kept apart from all of them, mean +- SD over the partitions :: sha256 c6c4138ee5c86423 :: prose twin: two significant figures, direction in words
\newcommand{\MlstMlpNatPerClassRandFlavMergedOverApart}{1.003\,\ensuremath{\pm}\,0.007}  % experiments/FIGS/data/v2/index.json :: random partitions, nonlinear probe, natural composition: median per-class ratio of the classes merged with a decay that differs only in quark flavor over that of the classes kept apart from all of them, mean +- SD over the partitions :: sha256 c6c4138ee5c86423
\newcommand{\MlstMlpNatPerClassRandFlavMergedOverApartPr}{equal within 1\%}  % experiments/FIGS/data/v2/index.json :: random partitions, nonlinear probe, natural composition: median per-class ratio of the classes merged with a decay that differs only in quark flavor over that of the classes kept apart from all of them, mean +- SD over the partitions :: sha256 c6c4138ee5c86423 :: prose twin: two significant figures, direction in words
\newcommand{\MlstMlpNatPerClassRandLeptonMergedOverApart}{1.01\,\ensuremath{\pm}\,0.07}  % experiments/FIGS/data/v2/index.json :: random partitions, nonlinear probe, natural composition: median per-class ratio of the classes merged with a decay that differs only in lepton flavor over that of the classes kept apart from all of them, mean +- SD over the partitions :: sha256 c6c4138ee5c86423
\newcommand{\MlstMlpNatPerClassRandLeptonMergedOverApartPr}{equal within 8\%}  % experiments/FIGS/data/v2/index.json :: random partitions, nonlinear probe, natural composition: median per-class ratio of the classes merged with a decay that differs only in lepton flavor over that of the classes kept apart from all of them, mean +- SD over the partitions :: sha256 c6c4138ee5c86423 :: prose twin: two significant figures, direction in words
\newcommand{\MlstPerClassAucDiffResSeedSdMedian}{0.015}  % experiments/FIGS/data/v2/index.json :: median over classes of the SD over indices of 100 (AUC_c of L162 - AUC_c of 188), pp :: sha256 c6c4138ee5c86423
\newcommand{\MlstPerClassAucDiffResMin}{\ensuremath{-}0.73}  % experiments/FIGS/data/v2/index.json :: most negative mean over indices of 100 (AUC_c of L162 - AUC_c of 188), pp :: sha256 c6c4138ee5c86423 :: label_QCD_bcss
\newcommand{\MlstPerClassAucDiffResMinClass}{QCD, \ensuremath{bcss}}  % experiments/FIGS/data/v2/index.json :: its class, L162 :: sha256 c6c4138ee5c86423
\newcommand{\MlstPerClassAucDiffMidSeedSdMedian}{0.032}  % experiments/FIGS/data/v2/index.json :: median over classes of the SD over indices of 100 (AUC_c of R42_Q1 - AUC_c of 188), pp :: sha256 c6c4138ee5c86423
\newcommand{\MlstPerClassAucDiffMidMin}{\ensuremath{-}1.4}  % experiments/FIGS/data/v2/index.json :: most negative mean over indices of 100 (AUC_c of R42_Q1 - AUC_c of 188), pp :: sha256 c6c4138ee5c86423 :: label_X_YY_cqtauev
\newcommand{\MlstPerClassAucDiffMidMinClass}{\ensuremath{X\to YY\to cq\tau_e\nu}}  % experiments/FIGS/data/v2/index.json :: its class, R42_Q1 :: sha256 c6c4138ee5c86423
\newcommand{\MlstPerClassAucDiffCoarseSeedSdMedian}{0.077}  % experiments/FIGS/data/v2/index.json :: median over classes of the SD over indices of 100 (AUC_c of R16_Q1 - AUC_c of 188), pp :: sha256 c6c4138ee5c86423
\newcommand{\MlstPerClassAucDiffCoarseMin}{\ensuremath{-}3.1}  % experiments/FIGS/data/v2/index.json :: most negative mean over indices of 100 (AUC_c of R16_Q1 - AUC_c of 188), pp :: sha256 c6c4138ee5c86423 :: label_QCD_bcss
\newcommand{\MlstPerClassAucDiffCoarseMinClass}{QCD, \ensuremath{bcss}}  % experiments/FIGS/data/v2/index.json :: its class, R16_Q1 :: sha256 c6c4138ee5c86423
\newcommand{\MlstRocJcTwoHurtList}{\ensuremath{X\to YY\to ccq}; \ensuremath{X\to YY\to ssb}; \ensuremath{X\to YY\to ccss}; QCD, \ensuremath{bcc}; QCD, \ensuremath{bccss}; QCD, \ensuremath{bss}; QCD, \ensuremath{ccs} and QCD, \ensuremath{ccss}}  % experiments/FIGS/data/v2/index.json :: jc2_roc_classes on the class-balanced probes; candidates with at least 1000 test jets :: sha256 c6c4138ee5c86423 :: label_X_YY_ccq, label_X_YY_ssb, label_X_YY_ccss, label_QCD_bcc, label_QCD_bccss, label_QCD_bss, label_QCD_ccs, label_QCD_ccss
\newcommand{\MlstRocJcTwoStableList}{\ensuremath{X\to gg}; \ensuremath{X\to YY\to ggg}; \ensuremath{X\to YY\to g\tau_h\tau_h} and QCD, light}  % experiments/FIGS/data/v2/index.json :: jc2_roc_classes on the class-balanced probes; candidates with at least 1000 test jets :: sha256 c6c4138ee5c86423 :: label_X_gg, label_X_YY_ggg, label_X_YY_gtauhtauh, label_QCD_light
\newcommand{\MlstRocJcTwoHurtMinDrop}{2.1}  % experiments/FIGS/data/v2/index.json :: jc2_roc_classes on the class-balanced probes; candidates with at least 1000 test jets :: sha256 c6c4138ee5c86423 :: smallest AUC loss at 17 classes of the hurt, pp
\newcommand{\MlstRocJcTwoHurtMaxDrop}{2.8}  % experiments/FIGS/data/v2/index.json :: jc2_roc_classes on the class-balanced probes; candidates with at least 1000 test jets :: sha256 c6c4138ee5c86423 :: largest AUC loss at 17 classes of the hurt, pp
\newcommand{\MlstRocJcTwoStableMaxChange}{15}  % experiments/FIGS/data/v2/index.json :: jc2_roc_classes on the class-balanced probes; candidates with at least 1000 test jets :: sha256 c6c4138ee5c86423 :: largest change of 1 - AUC at any label set of the stable classes, per cent of the 188-class value
\newcommand{\MlstJcBalNThin}{4}  % experiments/FIGS/data/v2/jc2bal/census.json :: counts.ft_test :: sha256 a4379ee0817254cf :: classes with fewer than 1000 class-balanced test jets (F3 hatching)
\newcommand{\MlstMinBkgPass}{20}  % experiments/FIGS/mlst_v4.py :: JC1_MIN_BKG :: sha256 fa536b4866ffd651 :: fewest background jets behind a drawn ROC point or a quoted rejection
\newcommand{\MlstMassMinJets}{100}  % experiments/EVAL/atlas_resolution.py :: MIN_JETS :: sha256 b82cc219f79f0f14 :: jets below which a mass bin is null
\newcommand{\MlstRocCutBbFull}{0.79}  % experiments/FIGS/data/v2/index.json :: curves_v2/pairs curve bb_vs_cc: largest cut.eps_s over the L188 runs ('not reached' when a run's cut is null or leaves fewer than 2 grid points) :: sha256 c6c4138ee5c86423
\newcommand{\MlstRocCutBbRes}{0.77}  % experiments/FIGS/data/v2/index.json :: curves_v2/pairs curve bb_vs_cc: largest cut.eps_s over the L162 runs ('not reached' when a run's cut is null or leaves fewer than 2 grid points) :: sha256 c6c4138ee5c86423
\newcommand{\MlstRocCutBbMid}{0.79}  % experiments/FIGS/data/v2/index.json :: curves_v2/pairs curve bb_vs_cc: largest cut.eps_s over the R42_Q1 runs ('not reached' when a run's cut is null or leaves fewer than 2 grid points) :: sha256 c6c4138ee5c86423
\newcommand{\MlstRocCutBbCoarse}{0.13}  % experiments/FIGS/data/v2/index.json :: curves_v2/pairs curve bb_vs_cc: largest cut.eps_s over the R16_Q1 runs ('not reached' when a run's cut is null or leaves fewer than 2 grid points) :: sha256 c6c4138ee5c86423
\newcommand{\MlstRocCutBbFzero}{0.03}  % experiments/FIGS/data/v2/index.json :: curves_v2/pairs curve bb_vs_cc: largest cut.eps_s over the FLAV_F0 runs ('not reached' when a run's cut is null or leaves fewer than 2 grid points) :: sha256 c6c4138ee5c86423
\newcommand{\MlstRocCutBbFone}{0.62}  % experiments/FIGS/data/v2/index.json :: curves_v2/pairs curve bb_vs_cc: largest cut.eps_s over the FLAV_F1 runs ('not reached' when a run's cut is null or leaves fewer than 2 grid points) :: sha256 c6c4138ee5c86423
\newcommand{\MlstRocCutBbFoner}{0.51}  % experiments/FIGS/data/v2/index.json :: curves_v2/pairs curve bb_vs_cc: largest cut.eps_s over the FLAV_F1R runs ('not reached' when a run's cut is null or leaves fewer than 2 grid points) :: sha256 c6c4138ee5c86423
\newcommand{\MlstRocNBkgBb}{200{,}000}  % experiments/FIGS/data/v2/index.json :: curves_v2/pairs curve bb_vs_cc: n_background :: sha256 c6c4138ee5c86423
\newcommand{\MlstRocRejMaxBb}{10{,}000}  % experiments/FIGS/data/v2/index.json :: curves_v2/pairs curve bb_vs_cc: n_background / 20 :: sha256 c6c4138ee5c86423
\newcommand{\MlstRocCutBbqqFull}{0.67}  % experiments/FIGS/data/v2/index.json :: curves_v2/pairs curve bbqq_vs_ccqq: largest cut.eps_s over the L188 runs ('not reached' when a run's cut is null or leaves fewer than 2 grid points) :: sha256 c6c4138ee5c86423
\newcommand{\MlstRocCutBbqqRes}{0.66}  % experiments/FIGS/data/v2/index.json :: curves_v2/pairs curve bbqq_vs_ccqq: largest cut.eps_s over the L162 runs ('not reached' when a run's cut is null or leaves fewer than 2 grid points) :: sha256 c6c4138ee5c86423
\newcommand{\MlstRocCutBbqqMid}{0.62}  % experiments/FIGS/data/v2/index.json :: curves_v2/pairs curve bbqq_vs_ccqq: largest cut.eps_s over the R42_Q1 runs ('not reached' when a run's cut is null or leaves fewer than 2 grid points) :: sha256 c6c4138ee5c86423
\newcommand{\MlstRocCutBbqqCoarse}{0.06}  % experiments/FIGS/data/v2/index.json :: curves_v2/pairs curve bbqq_vs_ccqq: largest cut.eps_s over the R16_Q1 runs ('not reached' when a run's cut is null or leaves fewer than 2 grid points) :: sha256 c6c4138ee5c86423
\newcommand{\MlstRocCutBbqqFzero}{0.06}  % experiments/FIGS/data/v2/index.json :: curves_v2/pairs curve bbqq_vs_ccqq: largest cut.eps_s over the FLAV_F0 runs ('not reached' when a run's cut is null or leaves fewer than 2 grid points) :: sha256 c6c4138ee5c86423
\newcommand{\MlstRocCutBbqqFone}{0.56}  % experiments/FIGS/data/v2/index.json :: curves_v2/pairs curve bbqq_vs_ccqq: largest cut.eps_s over the FLAV_F1 runs ('not reached' when a run's cut is null or leaves fewer than 2 grid points) :: sha256 c6c4138ee5c86423
\newcommand{\MlstRocCutBbqqFoner}{0.32}  % experiments/FIGS/data/v2/index.json :: curves_v2/pairs curve bbqq_vs_ccqq: largest cut.eps_s over the FLAV_F1R runs ('not reached' when a run's cut is null or leaves fewer than 2 grid points) :: sha256 c6c4138ee5c86423
\newcommand{\MlstRocNBkgBbqq}{200{,}000}  % experiments/FIGS/data/v2/index.json :: curves_v2/pairs curve bbqq_vs_ccqq: n_background :: sha256 c6c4138ee5c86423
\newcommand{\MlstRocRejMaxBbqq}{10{,}000}  % experiments/FIGS/data/v2/index.json :: curves_v2/pairs curve bbqq_vs_ccqq: n_background / 20 :: sha256 c6c4138ee5c86423
\newcommand{\MlstRocCutBcFull}{0.24}  % experiments/FIGS/data/v2/index.json :: curves_v2/pairs curve bc_vs_bq_cs: largest cut.eps_s over the L188 runs ('not reached' when a run's cut is null or leaves fewer than 2 grid points) :: sha256 c6c4138ee5c86423
\newcommand{\MlstRocCutBcRes}{0.25}  % experiments/FIGS/data/v2/index.json :: curves_v2/pairs curve bc_vs_bq_cs: largest cut.eps_s over the L162 runs ('not reached' when a run's cut is null or leaves fewer than 2 grid points) :: sha256 c6c4138ee5c86423
\newcommand{\MlstRocCutBcMid}{0.14}  % experiments/FIGS/data/v2/index.json :: curves_v2/pairs curve bc_vs_bq_cs: largest cut.eps_s over the R42_Q1 runs ('not reached' when a run's cut is null or leaves fewer than 2 grid points) :: sha256 c6c4138ee5c86423
\newcommand{\MlstRocCutBcCoarse}{0.00}  % experiments/FIGS/data/v2/index.json :: curves_v2/pairs curve bc_vs_bq_cs: largest cut.eps_s over the R16_Q1 runs ('not reached' when a run's cut is null or leaves fewer than 2 grid points) :: sha256 c6c4138ee5c86423
\newcommand{\MlstRocCutBcFzero}{0.00}  % experiments/FIGS/data/v2/index.json :: curves_v2/pairs curve bc_vs_bq_cs: largest cut.eps_s over the FLAV_F0 runs ('not reached' when a run's cut is null or leaves fewer than 2 grid points) :: sha256 c6c4138ee5c86423
\newcommand{\MlstRocCutBcFone}{0.01}  % experiments/FIGS/data/v2/index.json :: curves_v2/pairs curve bc_vs_bq_cs: largest cut.eps_s over the FLAV_F1 runs ('not reached' when a run's cut is null or leaves fewer than 2 grid points) :: sha256 c6c4138ee5c86423
\newcommand{\MlstRocCutBcFoner}{0.01}  % experiments/FIGS/data/v2/index.json :: curves_v2/pairs curve bc_vs_bq_cs: largest cut.eps_s over the FLAV_F1R runs ('not reached' when a run's cut is null or leaves fewer than 2 grid points) :: sha256 c6c4138ee5c86423
\newcommand{\MlstRocNBkgBc}{200{,}000}  % experiments/FIGS/data/v2/index.json :: curves_v2/pairs curve bc_vs_bq_cs: n_background :: sha256 c6c4138ee5c86423
\newcommand{\MlstRocRejMaxBc}{10{,}000}  % experiments/FIGS/data/v2/index.json :: curves_v2/pairs curve bc_vs_bq_cs: n_background / 20 :: sha256 c6c4138ee5c86423
\newcommand{\MlstRocCutEemmFull}{not reached}  % experiments/FIGS/data/v2/index.json :: curves_v2/pairs curve ee_vs_mm: largest cut.eps_s over the L188 runs ('not reached' when a run's cut is null or leaves fewer than 2 grid points) :: sha256 c6c4138ee5c86423
\newcommand{\MlstRocCutEemmRes}{not reached}  % experiments/FIGS/data/v2/index.json :: curves_v2/pairs curve ee_vs_mm: largest cut.eps_s over the L162 runs ('not reached' when a run's cut is null or leaves fewer than 2 grid points) :: sha256 c6c4138ee5c86423
\newcommand{\MlstRocCutEemmMid}{0.00}  % experiments/FIGS/data/v2/index.json :: curves_v2/pairs curve ee_vs_mm: largest cut.eps_s over the R42_Q1 runs ('not reached' when a run's cut is null or leaves fewer than 2 grid points) :: sha256 c6c4138ee5c86423
\newcommand{\MlstRocCutEemmCoarse}{0.00}  % experiments/FIGS/data/v2/index.json :: curves_v2/pairs curve ee_vs_mm: largest cut.eps_s over the R16_Q1 runs ('not reached' when a run's cut is null or leaves fewer than 2 grid points) :: sha256 c6c4138ee5c86423
\newcommand{\MlstRocCutEemmFzero}{not reached}  % experiments/FIGS/data/v2/index.json :: curves_v2/pairs curve ee_vs_mm: largest cut.eps_s over the FLAV_F0 runs ('not reached' when a run's cut is null or leaves fewer than 2 grid points) :: sha256 c6c4138ee5c86423
\newcommand{\MlstRocCutEemmFone}{not reached}  % experiments/FIGS/data/v2/index.json :: curves_v2/pairs curve ee_vs_mm: largest cut.eps_s over the FLAV_F1 runs ('not reached' when a run's cut is null or leaves fewer than 2 grid points) :: sha256 c6c4138ee5c86423
\newcommand{\MlstRocCutEemmFoner}{not reached}  % experiments/FIGS/data/v2/index.json :: curves_v2/pairs curve ee_vs_mm: largest cut.eps_s over the FLAV_F1R runs ('not reached' when a run's cut is null or leaves fewer than 2 grid points) :: sha256 c6c4138ee5c86423
\newcommand{\MlstRocNBkgEemm}{128{,}833}  % experiments/FIGS/data/v2/index.json :: curves_v2/pairs curve ee_vs_mm: n_background :: sha256 c6c4138ee5c86423
\newcommand{\MlstRocRejMaxEemm}{6{,}442}  % experiments/FIGS/data/v2/index.json :: curves_v2/pairs curve ee_vs_mm: n_background / 20 :: sha256 c6c4138ee5c86423
\newcommand{\MlstRocNoCurveEemm}{the \Kfull- and \Kres-class models and the three random partitions with controlled flavor boundaries}  % experiments/FIGS/data/v2/index.json :: curves_v2/pairs curve ee_vs_mm: configurations of F4 with a run whose cut is null or leaves fewer than 2 grid points (F4 draws no curve for them) :: sha256 c6c4138ee5c86423 :: L188, L162, FLAV_F0, FLAV_F1, FLAV_F1R
\newcommand{\MlstRocCutJcOneHbbFull}{0.40}  % experiments/FIGS/data/v2/index.json :: curves_v2/jc1_ft curve label_Hbb: largest cut.eps_s over the 188-class runs ('not reached' when a run's cut is null or leaves fewer than 2 grid points) :: sha256 c6c4138ee5c86423
\newcommand{\MlstRocCutJcOneHccFull}{0.37}  % experiments/FIGS/data/v2/index.json :: curves_v2/jc1_ft curve label_Hcc: largest cut.eps_s over the 188-class runs ('not reached' when a run's cut is null or leaves fewer than 2 grid points) :: sha256 c6c4138ee5c86423
\newcommand{\MlstRocCutJcOneHggFull}{0.04}  % experiments/FIGS/data/v2/index.json :: curves_v2/jc1_ft curve label_Hgg: largest cut.eps_s over the 188-class runs ('not reached' when a run's cut is null or leaves fewer than 2 grid points) :: sha256 c6c4138ee5c86423
\newcommand{\MlstRocCutJcOneHfourqFull}{0.26}  % experiments/FIGS/data/v2/index.json :: curves_v2/jc1_ft curve label_H4q: largest cut.eps_s over the 188-class runs ('not reached' when a run's cut is null or leaves fewer than 2 grid points) :: sha256 c6c4138ee5c86423
\newcommand{\MlstRocCutJcOneHqqlFull}{0.98}  % experiments/FIGS/data/v2/index.json :: curves_v2/jc1_ft curve label_Hqql: largest cut.eps_s over the 188-class runs ('not reached' when a run's cut is null or leaves fewer than 2 grid points) :: sha256 c6c4138ee5c86423
\newcommand{\MlstRocCutJcOneTbqqFull}{0.53}  % experiments/FIGS/data/v2/index.json :: curves_v2/jc1_ft curve label_Tbqq: largest cut.eps_s over the 188-class runs ('not reached' when a run's cut is null or leaves fewer than 2 grid points) :: sha256 c6c4138ee5c86423
\newcommand{\MlstRocCutJcOneTblFull}{not reached}  % experiments/FIGS/data/v2/index.json :: curves_v2/jc1_ft curve label_Tbl: largest cut.eps_s over the 188-class runs ('not reached' when a run's cut is null or leaves fewer than 2 grid points) :: sha256 c6c4138ee5c86423
\newcommand{\MlstRocCutJcOneWqqFull}{0.15}  % experiments/FIGS/data/v2/index.json :: curves_v2/jc1_ft curve label_Wqq: largest cut.eps_s over the 188-class runs ('not reached' when a run's cut is null or leaves fewer than 2 grid points) :: sha256 c6c4138ee5c86423
\newcommand{\MlstRocCutJcOneZqqFull}{0.12}  % experiments/FIGS/data/v2/index.json :: curves_v2/jc1_ft curve label_Zqq: largest cut.eps_s over the 188-class runs ('not reached' when a run's cut is null or leaves fewer than 2 grid points) :: sha256 c6c4138ee5c86423
\newcommand{\MlstFtTopRfiftyCoarseOverFull}{0.932\,\ensuremath{\pm}\,0.018}  % experiments/FIGS/data/v2/index.json :: curves_v2/top_ft curve top: rejection at eps_S = 0.5 of R16_Q1 over L188, run by run :: sha256 c6c4138ee5c86423
\newcommand{\MlstFtTopRfiftyCoarseOverFullPr}{6.8\% less rejection}  % experiments/FIGS/data/v2/index.json :: curves_v2/top_ft curve top: rejection at eps_S = 0.5 of R16_Q1 over L188, run by run :: sha256 c6c4138ee5c86423 :: prose twin: two significant figures, direction in words
\newcommand{\MlstProbeUnconverged}{0}  % experiments/FIGS/data/v2/index.json :: linear probes of pretrained models that did not converge (count) :: sha256 c6c4138ee5c86423
\newcommand{\MlstProbeUnconvergedInit}{0}  % experiments/FIGS/data/v2/index.json :: linear probes of the randomly initialized networks that did not converge (count) :: sha256 c6c4138ee5c86423
\newcommand{\MlstProbeFits}{895}  % experiments/FIGS/data/v2/index.json :: selected linear-probe fits of the reported checkpoints (count) :: sha256 c6c4138ee5c86423
\newcommand{\MlstProbeEdgeC}{16}  % experiments/FIGS/data/v2/index.json :: of these, fits whose selected C is the largest of the grid (count) :: sha256 c6c4138ee5c86423
\newcommand{\MlstMlpUnconverged}{0}  % experiments/FIGS/data/v2/index.json :: nonlinear probes that did not reach their plateau within the epoch cap (count) :: sha256 c6c4138ee5c86423
\newcommand{\MlstMlpCmpJcTwoResOverFull}{1.019\,\ensuremath{\pm}\,0.007}  % experiments/FIGS/data/v2/index.json :: nonlinear probe jc2_188: (1 - AUC) of L162 over L188, run by run :: sha256 c6c4138ee5c86423
\newcommand{\MlstMlpCmpJcTwoResOverFullPr}{1.9\% more errors}  % experiments/FIGS/data/v2/index.json :: nonlinear probe jc2_188: (1 - AUC) of L162 over L188, run by run :: sha256 c6c4138ee5c86423 :: prose twin: two significant figures, direction in words
\newcommand{\MlstMlpCmpJcTwoMidOverFull}{1.170\,\ensuremath{\pm}\,0.006}  % experiments/FIGS/data/v2/index.json :: nonlinear probe jc2_188: (1 - AUC) of R42_Q1 over L188, run by run :: sha256 c6c4138ee5c86423
\newcommand{\MlstMlpCmpJcTwoMidOverFullPr}{17\% more errors}  % experiments/FIGS/data/v2/index.json :: nonlinear probe jc2_188: (1 - AUC) of R42_Q1 over L188, run by run :: sha256 c6c4138ee5c86423 :: prose twin: two significant figures, direction in words
\newcommand{\MlstMlpCmpJcTwoCoarseOverFull}{1.388\,\ensuremath{\pm}\,0.021}  % experiments/FIGS/data/v2/index.json :: nonlinear probe jc2_188: (1 - AUC) of R16_Q1 over L188, run by run :: sha256 c6c4138ee5c86423
\newcommand{\MlstMlpCmpJcTwoCoarseOverFullPr}{39\% more errors}  % experiments/FIGS/data/v2/index.json :: nonlinear probe jc2_188: (1 - AUC) of R16_Q1 over L188, run by run :: sha256 c6c4138ee5c86423 :: prose twin: two significant figures, direction in words
\newcommand{\MlstMlpCmpJcTwoResMassOverRes}{\pending{no run yet}}  % experiments/FIGS/data/v2/index.json :: nonlinear probe jc2_188: (1 - AUC) of L162_MASS over L162, run by run :: sha256 c6c4138ee5c86423
\newcommand{\MlstMlpCmpJcTwoResMassOverResPr}{\pending{no run yet}}  % experiments/FIGS/data/v2/index.json :: nonlinear probe jc2_188: (1 - AUC) of L162_MASS over L162, run by run :: sha256 c6c4138ee5c86423 :: prose twin: two significant figures, direction in words
\newcommand{\MlstMlpCmpJcTwoResMassOverResPrMax}{\pending{no run yet}}  % experiments/FIGS/data/v2/index.json :: nonlinear probe jc2_188: (1 - AUC) of L162_MASS over L162, run by run :: sha256 c6c4138ee5c86423 :: prose twin of the run farther from 1
\newcommand{\MlstMlpCmpJcTwoCoarseMassOverCoarse}{\pending{no run yet}}  % experiments/FIGS/data/v2/index.json :: nonlinear probe jc2_188: (1 - AUC) of R16_Q1_MASS over R16_Q1, run by run :: sha256 c6c4138ee5c86423
\newcommand{\MlstMlpCmpJcTwoCoarseMassOverCoarsePr}{\pending{no run yet}}  % experiments/FIGS/data/v2/index.json :: nonlinear probe jc2_188: (1 - AUC) of R16_Q1_MASS over R16_Q1, run by run :: sha256 c6c4138ee5c86423 :: prose twin: two significant figures, direction in words
\newcommand{\MlstMlpCmpJcTwoCoarseMassOverCoarsePrMax}{\pending{no run yet}}  % experiments/FIGS/data/v2/index.json :: nonlinear probe jc2_188: (1 - AUC) of R16_Q1_MASS over R16_Q1, run by run :: sha256 c6c4138ee5c86423 :: prose twin of the run farther from 1
\newcommand{\MlstMlpCmpJcTwoFzeroOverFull}{1.294 and 1.310}  % experiments/FIGS/data/v2/index.json :: nonlinear probe jc2_188: (1 - AUC) of FLAV_F0 over L188, run by run :: sha256 c6c4138ee5c86423
\newcommand{\MlstMlpCmpJcTwoFzeroOverFullPr}{at least 29\% more errors}  % experiments/FIGS/data/v2/index.json :: nonlinear probe jc2_188: (1 - AUC) of FLAV_F0 over L188, run by run :: sha256 c6c4138ee5c86423 :: prose twin: two significant figures, direction in words
\newcommand{\MlstMlpCmpJcTwoFzeroOverFullPrMax}{up to 31\% more errors}  % experiments/FIGS/data/v2/index.json :: nonlinear probe jc2_188: (1 - AUC) of FLAV_F0 over L188, run by run :: sha256 c6c4138ee5c86423 :: prose twin of the run farther from 1
\newcommand{\MlstMlpCmpJcTwoFoneOverFzero}{0.903 and 0.885}  % experiments/FIGS/data/v2/index.json :: nonlinear probe jc2_188: (1 - AUC) of FLAV_F1 over FLAV_F0, run by run :: sha256 c6c4138ee5c86423
\newcommand{\MlstMlpCmpJcTwoFoneOverFzeroPr}{at least 9.7\% fewer errors}  % experiments/FIGS/data/v2/index.json :: nonlinear probe jc2_188: (1 - AUC) of FLAV_F1 over FLAV_F0, run by run :: sha256 c6c4138ee5c86423 :: prose twin: two significant figures, direction in words
\newcommand{\MlstMlpCmpJcTwoFoneOverFzeroPrMax}{up to 12\% fewer errors}  % experiments/FIGS/data/v2/index.json :: nonlinear probe jc2_188: (1 - AUC) of FLAV_F1 over FLAV_F0, run by run :: sha256 c6c4138ee5c86423 :: prose twin of the run farther from 1
\newcommand{\MlstMlpCmpJcTwoFonerOverFone}{0.9916 and 0.9918}  % experiments/FIGS/data/v2/index.json :: nonlinear probe jc2_188: (1 - AUC) of FLAV_F1R over FLAV_F1, run by run :: sha256 c6c4138ee5c86423
\newcommand{\MlstMlpCmpJcTwoFonerOverFonePr}{equal within 1\%}  % experiments/FIGS/data/v2/index.json :: nonlinear probe jc2_188: (1 - AUC) of FLAV_F1R over FLAV_F1, run by run :: sha256 c6c4138ee5c86423 :: prose twin: two significant figures, direction in words
\newcommand{\MlstMlpCmpJcTwoFonerOverFonePrMax}{equal within 1\%}  % experiments/FIGS/data/v2/index.json :: nonlinear probe jc2_188: (1 - AUC) of FLAV_F1R over FLAV_F1, run by run :: sha256 c6c4138ee5c86423 :: prose twin of the run farther from 1
\newcommand{\MlstMlpCmpJcTwoFonerOverFzero}{0.895 and 0.877}  % experiments/FIGS/data/v2/index.json :: nonlinear probe jc2_188: (1 - AUC) of FLAV_F1R over FLAV_F0, run by run :: sha256 c6c4138ee5c86423
\newcommand{\MlstMlpCmpJcTwoFonerOverFzeroPr}{at least 10\% fewer errors}  % experiments/FIGS/data/v2/index.json :: nonlinear probe jc2_188: (1 - AUC) of FLAV_F1R over FLAV_F0, run by run :: sha256 c6c4138ee5c86423 :: prose twin: two significant figures, direction in words
\newcommand{\MlstMlpCmpJcTwoFonerOverFzeroPrMax}{up to 12\% fewer errors}  % experiments/FIGS/data/v2/index.json :: nonlinear probe jc2_188: (1 - AUC) of FLAV_F1R over FLAV_F0, run by run :: sha256 c6c4138ee5c86423 :: prose twin of the run farther from 1
\newcommand{\MlstMlpCmpJcTwoFzeroOverCoarse}{0.9461 and 0.9456}  % experiments/FIGS/data/v2/index.json :: nonlinear probe jc2_188: (1 - AUC) of FLAV_F0 over R16_Q1, run by run :: sha256 c6c4138ee5c86423
\newcommand{\MlstMlpCmpJcTwoFzeroOverCoarsePr}{at least 5.4\% fewer errors}  % experiments/FIGS/data/v2/index.json :: nonlinear probe jc2_188: (1 - AUC) of FLAV_F0 over R16_Q1, run by run :: sha256 c6c4138ee5c86423 :: prose twin: two significant figures, direction in words
\newcommand{\MlstMlpCmpJcTwoFzeroOverCoarsePrMax}{up to 5.4\% fewer errors}  % experiments/FIGS/data/v2/index.json :: nonlinear probe jc2_188: (1 - AUC) of FLAV_F0 over R16_Q1, run by run :: sha256 c6c4138ee5c86423 :: prose twin of the run farther from 1
\newcommand{\MlstMlpCmpJcTwoRandOverFull}{\pending{no run yet}}  % experiments/FIGS/data/v2/index.json :: nonlinear probe jc2_188: (1 - AUC) of each random partition over the 188-class run of its seed :: sha256 c6c4138ee5c86423
\newcommand{\MlstMlpCmpJcTwoRandOverFullPr}{\pending{no run yet}}  % experiments/FIGS/data/v2/index.json :: nonlinear probe jc2_188: (1 - AUC) of each random partition over the 188-class run of its seed :: sha256 c6c4138ee5c86423 :: prose twin: two significant figures, direction in words
\newcommand{\MlstMlpCmpJcTwoResResOverFull}{1.001\,\ensuremath{\pm}\,0.007}  % experiments/FIGS/data/v2/index.json :: nonlinear probe jc2_188_res: (1 - AUC) of L162 over L188, run by run :: sha256 c6c4138ee5c86423
\newcommand{\MlstMlpCmpJcTwoResResOverFullPr}{equal within 1\%}  % experiments/FIGS/data/v2/index.json :: nonlinear probe jc2_188_res: (1 - AUC) of L162 over L188, run by run :: sha256 c6c4138ee5c86423 :: prose twin: two significant figures, direction in words
\newcommand{\MlstMlpCmpJcTwoResMidOverFull}{1.192\,\ensuremath{\pm}\,0.008}  % experiments/FIGS/data/v2/index.json :: nonlinear probe jc2_188_res: (1 - AUC) of R42_Q1 over L188, run by run :: sha256 c6c4138ee5c86423
\newcommand{\MlstMlpCmpJcTwoResMidOverFullPr}{19\% more errors}  % experiments/FIGS/data/v2/index.json :: nonlinear probe jc2_188_res: (1 - AUC) of R42_Q1 over L188, run by run :: sha256 c6c4138ee5c86423 :: prose twin: two significant figures, direction in words
\newcommand{\MlstMlpCmpJcTwoResCoarseOverFull}{1.408\,\ensuremath{\pm}\,0.026}  % experiments/FIGS/data/v2/index.json :: nonlinear probe jc2_188_res: (1 - AUC) of R16_Q1 over L188, run by run :: sha256 c6c4138ee5c86423
\newcommand{\MlstMlpCmpJcTwoResCoarseOverFullPr}{41\% more errors}  % experiments/FIGS/data/v2/index.json :: nonlinear probe jc2_188_res: (1 - AUC) of R16_Q1 over L188, run by run :: sha256 c6c4138ee5c86423 :: prose twin: two significant figures, direction in words
\newcommand{\MlstMlpCmpJcTwoResResMassOverRes}{\pending{no run yet}}  % experiments/FIGS/data/v2/index.json :: nonlinear probe jc2_188_res: (1 - AUC) of L162_MASS over L162, run by run :: sha256 c6c4138ee5c86423
\newcommand{\MlstMlpCmpJcTwoResResMassOverResPr}{\pending{no run yet}}  % experiments/FIGS/data/v2/index.json :: nonlinear probe jc2_188_res: (1 - AUC) of L162_MASS over L162, run by run :: sha256 c6c4138ee5c86423 :: prose twin: two significant figures, direction in words
\newcommand{\MlstMlpCmpJcTwoResResMassOverResPrMax}{\pending{no run yet}}  % experiments/FIGS/data/v2/index.json :: nonlinear probe jc2_188_res: (1 - AUC) of L162_MASS over L162, run by run :: sha256 c6c4138ee5c86423 :: prose twin of the run farther from 1
\newcommand{\MlstMlpCmpJcTwoResCoarseMassOverCoarse}{\pending{no run yet}}  % experiments/FIGS/data/v2/index.json :: nonlinear probe jc2_188_res: (1 - AUC) of R16_Q1_MASS over R16_Q1, run by run :: sha256 c6c4138ee5c86423
\newcommand{\MlstMlpCmpJcTwoResCoarseMassOverCoarsePr}{\pending{no run yet}}  % experiments/FIGS/data/v2/index.json :: nonlinear probe jc2_188_res: (1 - AUC) of R16_Q1_MASS over R16_Q1, run by run :: sha256 c6c4138ee5c86423 :: prose twin: two significant figures, direction in words
\newcommand{\MlstMlpCmpJcTwoResCoarseMassOverCoarsePrMax}{\pending{no run yet}}  % experiments/FIGS/data/v2/index.json :: nonlinear probe jc2_188_res: (1 - AUC) of R16_Q1_MASS over R16_Q1, run by run :: sha256 c6c4138ee5c86423 :: prose twin of the run farther from 1
\newcommand{\MlstMlpCmpJcTwoResFzeroOverFull}{1.285 and 1.300}  % experiments/FIGS/data/v2/index.json :: nonlinear probe jc2_188_res: (1 - AUC) of FLAV_F0 over L188, run by run :: sha256 c6c4138ee5c86423
\newcommand{\MlstMlpCmpJcTwoResFzeroOverFullPr}{at least 29\% more errors}  % experiments/FIGS/data/v2/index.json :: nonlinear probe jc2_188_res: (1 - AUC) of FLAV_F0 over L188, run by run :: sha256 c6c4138ee5c86423 :: prose twin: two significant figures, direction in words
\newcommand{\MlstMlpCmpJcTwoResFzeroOverFullPrMax}{up to 30\% more errors}  % experiments/FIGS/data/v2/index.json :: nonlinear probe jc2_188_res: (1 - AUC) of FLAV_F0 over L188, run by run :: sha256 c6c4138ee5c86423 :: prose twin of the run farther from 1
\newcommand{\MlstMlpCmpJcTwoResFoneOverFzero}{0.909 and 0.889}  % experiments/FIGS/data/v2/index.json :: nonlinear probe jc2_188_res: (1 - AUC) of FLAV_F1 over FLAV_F0, run by run :: sha256 c6c4138ee5c86423
\newcommand{\MlstMlpCmpJcTwoResFoneOverFzeroPr}{at least 9.1\% fewer errors}  % experiments/FIGS/data/v2/index.json :: nonlinear probe jc2_188_res: (1 - AUC) of FLAV_F1 over FLAV_F0, run by run :: sha256 c6c4138ee5c86423 :: prose twin: two significant figures, direction in words
\newcommand{\MlstMlpCmpJcTwoResFoneOverFzeroPrMax}{up to 11\% fewer errors}  % experiments/FIGS/data/v2/index.json :: nonlinear probe jc2_188_res: (1 - AUC) of FLAV_F1 over FLAV_F0, run by run :: sha256 c6c4138ee5c86423 :: prose twin of the run farther from 1
\newcommand{\MlstMlpCmpJcTwoResFonerOverFone}{0.9793 and 0.9825}  % experiments/FIGS/data/v2/index.json :: nonlinear probe jc2_188_res: (1 - AUC) of FLAV_F1R over FLAV_F1, run by run :: sha256 c6c4138ee5c86423
\newcommand{\MlstMlpCmpJcTwoResFonerOverFonePr}{equal within 3\%}  % experiments/FIGS/data/v2/index.json :: nonlinear probe jc2_188_res: (1 - AUC) of FLAV_F1R over FLAV_F1, run by run :: sha256 c6c4138ee5c86423 :: prose twin: two significant figures, direction in words
\newcommand{\MlstMlpCmpJcTwoResFonerOverFonePrMax}{equal within 3\%}  % experiments/FIGS/data/v2/index.json :: nonlinear probe jc2_188_res: (1 - AUC) of FLAV_F1R over FLAV_F1, run by run :: sha256 c6c4138ee5c86423 :: prose twin of the run farther from 1
\newcommand{\MlstMlpCmpJcTwoResFonerOverFzero}{0.890 and 0.874}  % experiments/FIGS/data/v2/index.json :: nonlinear probe jc2_188_res: (1 - AUC) of FLAV_F1R over FLAV_F0, run by run :: sha256 c6c4138ee5c86423
\newcommand{\MlstMlpCmpJcTwoResFonerOverFzeroPr}{at least 11\% fewer errors}  % experiments/FIGS/data/v2/index.json :: nonlinear probe jc2_188_res: (1 - AUC) of FLAV_F1R over FLAV_F0, run by run :: sha256 c6c4138ee5c86423 :: prose twin: two significant figures, direction in words
\newcommand{\MlstMlpCmpJcTwoResFonerOverFzeroPrMax}{up to 13\% fewer errors}  % experiments/FIGS/data/v2/index.json :: nonlinear probe jc2_188_res: (1 - AUC) of FLAV_F1R over FLAV_F0, run by run :: sha256 c6c4138ee5c86423 :: prose twin of the run farther from 1
\newcommand{\MlstMlpCmpJcTwoResFzeroOverCoarse}{0.9299 and 0.9233}  % experiments/FIGS/data/v2/index.json :: nonlinear probe jc2_188_res: (1 - AUC) of FLAV_F0 over R16_Q1, run by run :: sha256 c6c4138ee5c86423
\newcommand{\MlstMlpCmpJcTwoResFzeroOverCoarsePr}{at least 7.0\% fewer errors}  % experiments/FIGS/data/v2/index.json :: nonlinear probe jc2_188_res: (1 - AUC) of FLAV_F0 over R16_Q1, run by run :: sha256 c6c4138ee5c86423 :: prose twin: two significant figures, direction in words
\newcommand{\MlstMlpCmpJcTwoResFzeroOverCoarsePrMax}{up to 7.7\% fewer errors}  % experiments/FIGS/data/v2/index.json :: nonlinear probe jc2_188_res: (1 - AUC) of FLAV_F0 over R16_Q1, run by run :: sha256 c6c4138ee5c86423 :: prose twin of the run farther from 1
\newcommand{\MlstMlpCmpJcTwoResRandOverFull}{\pending{no run yet}}  % experiments/FIGS/data/v2/index.json :: nonlinear probe jc2_188_res: (1 - AUC) of each random partition over the 188-class run of its seed :: sha256 c6c4138ee5c86423
\newcommand{\MlstMlpCmpJcTwoResRandOverFullPr}{\pending{no run yet}}  % experiments/FIGS/data/v2/index.json :: nonlinear probe jc2_188_res: (1 - AUC) of each random partition over the 188-class run of its seed :: sha256 c6c4138ee5c86423 :: prose twin: two significant figures, direction in words
\newcommand{\MlstMlpCmpJcTwoQcdResOverFull}{1.070\,\ensuremath{\pm}\,0.006}  % experiments/FIGS/data/v2/index.json :: nonlinear probe jc2_188_qcd: (1 - AUC) of L162 over L188, run by run :: sha256 c6c4138ee5c86423
\newcommand{\MlstMlpCmpJcTwoQcdResOverFullPr}{7.0\% more errors}  % experiments/FIGS/data/v2/index.json :: nonlinear probe jc2_188_qcd: (1 - AUC) of L162 over L188, run by run :: sha256 c6c4138ee5c86423 :: prose twin: two significant figures, direction in words
\newcommand{\MlstMlpCmpJcTwoQcdMidOverFull}{1.107\,\ensuremath{\pm}\,0.002}  % experiments/FIGS/data/v2/index.json :: nonlinear probe jc2_188_qcd: (1 - AUC) of R42_Q1 over L188, run by run :: sha256 c6c4138ee5c86423
\newcommand{\MlstMlpCmpJcTwoQcdMidOverFullPr}{11\% more errors}  % experiments/FIGS/data/v2/index.json :: nonlinear probe jc2_188_qcd: (1 - AUC) of R42_Q1 over L188, run by run :: sha256 c6c4138ee5c86423 :: prose twin: two significant figures, direction in words
\newcommand{\MlstMlpCmpJcTwoQcdCoarseOverFull}{1.333\,\ensuremath{\pm}\,0.012}  % experiments/FIGS/data/v2/index.json :: nonlinear probe jc2_188_qcd: (1 - AUC) of R16_Q1 over L188, run by run :: sha256 c6c4138ee5c86423
\newcommand{\MlstMlpCmpJcTwoQcdCoarseOverFullPr}{33\% more errors}  % experiments/FIGS/data/v2/index.json :: nonlinear probe jc2_188_qcd: (1 - AUC) of R16_Q1 over L188, run by run :: sha256 c6c4138ee5c86423 :: prose twin: two significant figures, direction in words
\newcommand{\MlstMlpCmpJcTwoQcdResMassOverRes}{\pending{no run yet}}  % experiments/FIGS/data/v2/index.json :: nonlinear probe jc2_188_qcd: (1 - AUC) of L162_MASS over L162, run by run :: sha256 c6c4138ee5c86423
\newcommand{\MlstMlpCmpJcTwoQcdResMassOverResPr}{\pending{no run yet}}  % experiments/FIGS/data/v2/index.json :: nonlinear probe jc2_188_qcd: (1 - AUC) of L162_MASS over L162, run by run :: sha256 c6c4138ee5c86423 :: prose twin: two significant figures, direction in words
\newcommand{\MlstMlpCmpJcTwoQcdResMassOverResPrMax}{\pending{no run yet}}  % experiments/FIGS/data/v2/index.json :: nonlinear probe jc2_188_qcd: (1 - AUC) of L162_MASS over L162, run by run :: sha256 c6c4138ee5c86423 :: prose twin of the run farther from 1
\newcommand{\MlstMlpCmpJcTwoQcdCoarseMassOverCoarse}{\pending{no run yet}}  % experiments/FIGS/data/v2/index.json :: nonlinear probe jc2_188_qcd: (1 - AUC) of R16_Q1_MASS over R16_Q1, run by run :: sha256 c6c4138ee5c86423
\newcommand{\MlstMlpCmpJcTwoQcdCoarseMassOverCoarsePr}{\pending{no run yet}}  % experiments/FIGS/data/v2/index.json :: nonlinear probe jc2_188_qcd: (1 - AUC) of R16_Q1_MASS over R16_Q1, run by run :: sha256 c6c4138ee5c86423 :: prose twin: two significant figures, direction in words
\newcommand{\MlstMlpCmpJcTwoQcdCoarseMassOverCoarsePrMax}{\pending{no run yet}}  % experiments/FIGS/data/v2/index.json :: nonlinear probe jc2_188_qcd: (1 - AUC) of R16_Q1_MASS over R16_Q1, run by run :: sha256 c6c4138ee5c86423 :: prose twin of the run farther from 1
\newcommand{\MlstMlpCmpJcTwoQcdFzeroOverFull}{1.319 and 1.339}  % experiments/FIGS/data/v2/index.json :: nonlinear probe jc2_188_qcd: (1 - AUC) of FLAV_F0 over L188, run by run :: sha256 c6c4138ee5c86423
\newcommand{\MlstMlpCmpJcTwoQcdFzeroOverFullPr}{at least 32\% more errors}  % experiments/FIGS/data/v2/index.json :: nonlinear probe jc2_188_qcd: (1 - AUC) of FLAV_F0 over L188, run by run :: sha256 c6c4138ee5c86423 :: prose twin: two significant figures, direction in words
\newcommand{\MlstMlpCmpJcTwoQcdFzeroOverFullPrMax}{up to 34\% more errors}  % experiments/FIGS/data/v2/index.json :: nonlinear probe jc2_188_qcd: (1 - AUC) of FLAV_F0 over L188, run by run :: sha256 c6c4138ee5c86423 :: prose twin of the run farther from 1
\newcommand{\MlstMlpCmpJcTwoQcdFoneOverFzero}{0.886 and 0.872}  % experiments/FIGS/data/v2/index.json :: nonlinear probe jc2_188_qcd: (1 - AUC) of FLAV_F1 over FLAV_F0, run by run :: sha256 c6c4138ee5c86423
\newcommand{\MlstMlpCmpJcTwoQcdFoneOverFzeroPr}{at least 11\% fewer errors}  % experiments/FIGS/data/v2/index.json :: nonlinear probe jc2_188_qcd: (1 - AUC) of FLAV_F1 over FLAV_F0, run by run :: sha256 c6c4138ee5c86423 :: prose twin: two significant figures, direction in words
\newcommand{\MlstMlpCmpJcTwoQcdFoneOverFzeroPrMax}{up to 13\% fewer errors}  % experiments/FIGS/data/v2/index.json :: nonlinear probe jc2_188_qcd: (1 - AUC) of FLAV_F1 over FLAV_F0, run by run :: sha256 c6c4138ee5c86423 :: prose twin of the run farther from 1
\newcommand{\MlstMlpCmpJcTwoQcdFonerOverFone}{1.026 and 1.018}  % experiments/FIGS/data/v2/index.json :: nonlinear probe jc2_188_qcd: (1 - AUC) of FLAV_F1R over FLAV_F1, run by run :: sha256 c6c4138ee5c86423
\newcommand{\MlstMlpCmpJcTwoQcdFonerOverFonePr}{at least 1.8\% more errors}  % experiments/FIGS/data/v2/index.json :: nonlinear probe jc2_188_qcd: (1 - AUC) of FLAV_F1R over FLAV_F1, run by run :: sha256 c6c4138ee5c86423 :: prose twin: two significant figures, direction in words
\newcommand{\MlstMlpCmpJcTwoQcdFonerOverFonePrMax}{up to 2.6\% more errors}  % experiments/FIGS/data/v2/index.json :: nonlinear probe jc2_188_qcd: (1 - AUC) of FLAV_F1R over FLAV_F1, run by run :: sha256 c6c4138ee5c86423 :: prose twin of the run farther from 1
\newcommand{\MlstMlpCmpJcTwoQcdFonerOverFzero}{0.909 and 0.887}  % experiments/FIGS/data/v2/index.json :: nonlinear probe jc2_188_qcd: (1 - AUC) of FLAV_F1R over FLAV_F0, run by run :: sha256 c6c4138ee5c86423
\newcommand{\MlstMlpCmpJcTwoQcdFonerOverFzeroPr}{at least 9.1\% fewer errors}  % experiments/FIGS/data/v2/index.json :: nonlinear probe jc2_188_qcd: (1 - AUC) of FLAV_F1R over FLAV_F0, run by run :: sha256 c6c4138ee5c86423 :: prose twin: two significant figures, direction in words
\newcommand{\MlstMlpCmpJcTwoQcdFonerOverFzeroPrMax}{up to 11\% fewer errors}  % experiments/FIGS/data/v2/index.json :: nonlinear probe jc2_188_qcd: (1 - AUC) of FLAV_F1R over FLAV_F0, run by run :: sha256 c6c4138ee5c86423 :: prose twin of the run farther from 1
\newcommand{\MlstMlpCmpJcTwoQcdFzeroOverCoarse}{0.993 and 1.012}  % experiments/FIGS/data/v2/index.json :: nonlinear probe jc2_188_qcd: (1 - AUC) of FLAV_F0 over R16_Q1, run by run :: sha256 c6c4138ee5c86423
\newcommand{\MlstMlpCmpJcTwoQcdFzeroOverCoarsePr}{equal within 2\%}  % experiments/FIGS/data/v2/index.json :: nonlinear probe jc2_188_qcd: (1 - AUC) of FLAV_F0 over R16_Q1, run by run :: sha256 c6c4138ee5c86423 :: prose twin: two significant figures, direction in words
\newcommand{\MlstMlpCmpJcTwoQcdFzeroOverCoarsePrMax}{equal within 2\%}  % experiments/FIGS/data/v2/index.json :: nonlinear probe jc2_188_qcd: (1 - AUC) of FLAV_F0 over R16_Q1, run by run :: sha256 c6c4138ee5c86423 :: prose twin of the run farther from 1
\newcommand{\MlstMlpCmpJcTwoQcdRandOverFull}{\pending{no run yet}}  % experiments/FIGS/data/v2/index.json :: nonlinear probe jc2_188_qcd: (1 - AUC) of each random partition over the 188-class run of its seed :: sha256 c6c4138ee5c86423
\newcommand{\MlstMlpCmpJcTwoQcdRandOverFullPr}{\pending{no run yet}}  % experiments/FIGS/data/v2/index.json :: nonlinear probe jc2_188_qcd: (1 - AUC) of each random partition over the 188-class run of its seed :: sha256 c6c4138ee5c86423 :: prose twin: two significant figures, direction in words
\newcommand{\MlstMlpCmpBcResOverFull}{1.002\,\ensuremath{\pm}\,0.006}  % experiments/FIGS/data/v2/index.json :: nonlinear probe bc_vs_bq_cs: (1 - AUC) of L162 over L188, run by run :: sha256 c6c4138ee5c86423
\newcommand{\MlstMlpCmpBcResOverFullPr}{equal within 1\%}  % experiments/FIGS/data/v2/index.json :: nonlinear probe bc_vs_bq_cs: (1 - AUC) of L162 over L188, run by run :: sha256 c6c4138ee5c86423 :: prose twin: two significant figures, direction in words
\newcommand{\MlstMlpCmpBcMidOverFull}{1.104\,\ensuremath{\pm}\,0.017}  % experiments/FIGS/data/v2/index.json :: nonlinear probe bc_vs_bq_cs: (1 - AUC) of R42_Q1 over L188, run by run :: sha256 c6c4138ee5c86423
\newcommand{\MlstMlpCmpBcMidOverFullPr}{10\% more errors}  % experiments/FIGS/data/v2/index.json :: nonlinear probe bc_vs_bq_cs: (1 - AUC) of R42_Q1 over L188, run by run :: sha256 c6c4138ee5c86423 :: prose twin: two significant figures, direction in words
\newcommand{\MlstMlpCmpBcCoarseOverFull}{1.624\,\ensuremath{\pm}\,0.022}  % experiments/FIGS/data/v2/index.json :: nonlinear probe bc_vs_bq_cs: (1 - AUC) of R16_Q1 over L188, run by run :: sha256 c6c4138ee5c86423
\newcommand{\MlstMlpCmpBcCoarseOverFullPr}{62\% more errors}  % experiments/FIGS/data/v2/index.json :: nonlinear probe bc_vs_bq_cs: (1 - AUC) of R16_Q1 over L188, run by run :: sha256 c6c4138ee5c86423 :: prose twin: two significant figures, direction in words
\newcommand{\MlstMlpCmpBcResMassOverRes}{\pending{no run yet}}  % experiments/FIGS/data/v2/index.json :: nonlinear probe bc_vs_bq_cs: (1 - AUC) of L162_MASS over L162, run by run :: sha256 c6c4138ee5c86423
\newcommand{\MlstMlpCmpBcResMassOverResPr}{\pending{no run yet}}  % experiments/FIGS/data/v2/index.json :: nonlinear probe bc_vs_bq_cs: (1 - AUC) of L162_MASS over L162, run by run :: sha256 c6c4138ee5c86423 :: prose twin: two significant figures, direction in words
\newcommand{\MlstMlpCmpBcResMassOverResPrMax}{\pending{no run yet}}  % experiments/FIGS/data/v2/index.json :: nonlinear probe bc_vs_bq_cs: (1 - AUC) of L162_MASS over L162, run by run :: sha256 c6c4138ee5c86423 :: prose twin of the run farther from 1
\newcommand{\MlstMlpCmpBcCoarseMassOverCoarse}{\pending{no run yet}}  % experiments/FIGS/data/v2/index.json :: nonlinear probe bc_vs_bq_cs: (1 - AUC) of R16_Q1_MASS over R16_Q1, run by run :: sha256 c6c4138ee5c86423
\newcommand{\MlstMlpCmpBcCoarseMassOverCoarsePr}{\pending{no run yet}}  % experiments/FIGS/data/v2/index.json :: nonlinear probe bc_vs_bq_cs: (1 - AUC) of R16_Q1_MASS over R16_Q1, run by run :: sha256 c6c4138ee5c86423 :: prose twin: two significant figures, direction in words
\newcommand{\MlstMlpCmpBcCoarseMassOverCoarsePrMax}{\pending{no run yet}}  % experiments/FIGS/data/v2/index.json :: nonlinear probe bc_vs_bq_cs: (1 - AUC) of R16_Q1_MASS over R16_Q1, run by run :: sha256 c6c4138ee5c86423 :: prose twin of the run farther from 1
\newcommand{\MlstMlpCmpBcFzeroOverFull}{1.633 and 1.679}  % experiments/FIGS/data/v2/index.json :: nonlinear probe bc_vs_bq_cs: (1 - AUC) of FLAV_F0 over L188, run by run :: sha256 c6c4138ee5c86423
\newcommand{\MlstMlpCmpBcFzeroOverFullPr}{at least 63\% more errors}  % experiments/FIGS/data/v2/index.json :: nonlinear probe bc_vs_bq_cs: (1 - AUC) of FLAV_F0 over L188, run by run :: sha256 c6c4138ee5c86423 :: prose twin: two significant figures, direction in words
\newcommand{\MlstMlpCmpBcFzeroOverFullPrMax}{up to 68\% more errors}  % experiments/FIGS/data/v2/index.json :: nonlinear probe bc_vs_bq_cs: (1 - AUC) of FLAV_F0 over L188, run by run :: sha256 c6c4138ee5c86423 :: prose twin of the run farther from 1
\newcommand{\MlstMlpCmpBcFoneOverFzero}{0.848 and 0.810}  % experiments/FIGS/data/v2/index.json :: nonlinear probe bc_vs_bq_cs: (1 - AUC) of FLAV_F1 over FLAV_F0, run by run :: sha256 c6c4138ee5c86423
\newcommand{\MlstMlpCmpBcFoneOverFzeroPr}{at least 15\% fewer errors}  % experiments/FIGS/data/v2/index.json :: nonlinear probe bc_vs_bq_cs: (1 - AUC) of FLAV_F1 over FLAV_F0, run by run :: sha256 c6c4138ee5c86423 :: prose twin: two significant figures, direction in words
\newcommand{\MlstMlpCmpBcFoneOverFzeroPrMax}{up to 19\% fewer errors}  % experiments/FIGS/data/v2/index.json :: nonlinear probe bc_vs_bq_cs: (1 - AUC) of FLAV_F1 over FLAV_F0, run by run :: sha256 c6c4138ee5c86423 :: prose twin of the run farther from 1
\newcommand{\MlstMlpCmpBcFonerOverFone}{0.9812 and 0.9815}  % experiments/FIGS/data/v2/index.json :: nonlinear probe bc_vs_bq_cs: (1 - AUC) of FLAV_F1R over FLAV_F1, run by run :: sha256 c6c4138ee5c86423
\newcommand{\MlstMlpCmpBcFonerOverFonePr}{equal within 2\%}  % experiments/FIGS/data/v2/index.json :: nonlinear probe bc_vs_bq_cs: (1 - AUC) of FLAV_F1R over FLAV_F1, run by run :: sha256 c6c4138ee5c86423 :: prose twin: two significant figures, direction in words
\newcommand{\MlstMlpCmpBcFonerOverFonePrMax}{equal within 2\%}  % experiments/FIGS/data/v2/index.json :: nonlinear probe bc_vs_bq_cs: (1 - AUC) of FLAV_F1R over FLAV_F1, run by run :: sha256 c6c4138ee5c86423 :: prose twin of the run farther from 1
\newcommand{\MlstMlpCmpBcFonerOverFzero}{0.833 and 0.795}  % experiments/FIGS/data/v2/index.json :: nonlinear probe bc_vs_bq_cs: (1 - AUC) of FLAV_F1R over FLAV_F0, run by run :: sha256 c6c4138ee5c86423
\newcommand{\MlstMlpCmpBcFonerOverFzeroPr}{at least 17\% fewer errors}  % experiments/FIGS/data/v2/index.json :: nonlinear probe bc_vs_bq_cs: (1 - AUC) of FLAV_F1R over FLAV_F0, run by run :: sha256 c6c4138ee5c86423 :: prose twin: two significant figures, direction in words
\newcommand{\MlstMlpCmpBcFonerOverFzeroPrMax}{up to 21\% fewer errors}  % experiments/FIGS/data/v2/index.json :: nonlinear probe bc_vs_bq_cs: (1 - AUC) of FLAV_F1R over FLAV_F0, run by run :: sha256 c6c4138ee5c86423 :: prose twin of the run farther from 1
\newcommand{\MlstMlpCmpBcFzeroOverCoarse}{0.994 and 1.030}  % experiments/FIGS/data/v2/index.json :: nonlinear probe bc_vs_bq_cs: (1 - AUC) of FLAV_F0 over R16_Q1, run by run :: sha256 c6c4138ee5c86423
\newcommand{\MlstMlpCmpBcFzeroOverCoarsePr}{equal within 3\%}  % experiments/FIGS/data/v2/index.json :: nonlinear probe bc_vs_bq_cs: (1 - AUC) of FLAV_F0 over R16_Q1, run by run :: sha256 c6c4138ee5c86423 :: prose twin: two significant figures, direction in words
\newcommand{\MlstMlpCmpBcFzeroOverCoarsePrMax}{equal within 3\%}  % experiments/FIGS/data/v2/index.json :: nonlinear probe bc_vs_bq_cs: (1 - AUC) of FLAV_F0 over R16_Q1, run by run :: sha256 c6c4138ee5c86423 :: prose twin of the run farther from 1
\newcommand{\MlstMlpCmpBcRandOverFull}{\pending{no run yet}}  % experiments/FIGS/data/v2/index.json :: nonlinear probe bc_vs_bq_cs: (1 - AUC) of each random partition over the 188-class run of its seed :: sha256 c6c4138ee5c86423
\newcommand{\MlstMlpCmpBcRandOverFullPr}{\pending{no run yet}}  % experiments/FIGS/data/v2/index.json :: nonlinear probe bc_vs_bq_cs: (1 - AUC) of each random partition over the 188-class run of its seed :: sha256 c6c4138ee5c86423 :: prose twin: two significant figures, direction in words
\newcommand{\MlstMlpCmpBbResOverFull}{1.023\,\ensuremath{\pm}\,0.021}  % experiments/FIGS/data/v2/index.json :: nonlinear probe bb_vs_cc: (1 - AUC) of L162 over L188, run by run :: sha256 c6c4138ee5c86423
\newcommand{\MlstMlpCmpBbResOverFullPr}{equal within 3\%}  % experiments/FIGS/data/v2/index.json :: nonlinear probe bb_vs_cc: (1 - AUC) of L162 over L188, run by run :: sha256 c6c4138ee5c86423 :: prose twin: two significant figures, direction in words
\newcommand{\MlstMlpCmpBbMidOverFull}{1.125\,\ensuremath{\pm}\,0.011}  % experiments/FIGS/data/v2/index.json :: nonlinear probe bb_vs_cc: (1 - AUC) of R42_Q1 over L188, run by run :: sha256 c6c4138ee5c86423
\newcommand{\MlstMlpCmpBbMidOverFullPr}{13\% more errors}  % experiments/FIGS/data/v2/index.json :: nonlinear probe bb_vs_cc: (1 - AUC) of R42_Q1 over L188, run by run :: sha256 c6c4138ee5c86423 :: prose twin: two significant figures, direction in words
\newcommand{\MlstMlpCmpBbCoarseOverFull}{3.04\,\ensuremath{\pm}\,0.09}  % experiments/FIGS/data/v2/index.json :: nonlinear probe bb_vs_cc: (1 - AUC) of R16_Q1 over L188, run by run :: sha256 c6c4138ee5c86423
\newcommand{\MlstMlpCmpBbCoarseOverFullPr}{3.0 times the errors}  % experiments/FIGS/data/v2/index.json :: nonlinear probe bb_vs_cc: (1 - AUC) of R16_Q1 over L188, run by run :: sha256 c6c4138ee5c86423 :: prose twin: two significant figures, direction in words
\newcommand{\MlstMlpCmpBbResMassOverRes}{\pending{no run yet}}  % experiments/FIGS/data/v2/index.json :: nonlinear probe bb_vs_cc: (1 - AUC) of L162_MASS over L162, run by run :: sha256 c6c4138ee5c86423
\newcommand{\MlstMlpCmpBbResMassOverResPr}{\pending{no run yet}}  % experiments/FIGS/data/v2/index.json :: nonlinear probe bb_vs_cc: (1 - AUC) of L162_MASS over L162, run by run :: sha256 c6c4138ee5c86423 :: prose twin: two significant figures, direction in words
\newcommand{\MlstMlpCmpBbResMassOverResPrMax}{\pending{no run yet}}  % experiments/FIGS/data/v2/index.json :: nonlinear probe bb_vs_cc: (1 - AUC) of L162_MASS over L162, run by run :: sha256 c6c4138ee5c86423 :: prose twin of the run farther from 1
\newcommand{\MlstMlpCmpBbCoarseMassOverCoarse}{\pending{no run yet}}  % experiments/FIGS/data/v2/index.json :: nonlinear probe bb_vs_cc: (1 - AUC) of R16_Q1_MASS over R16_Q1, run by run :: sha256 c6c4138ee5c86423
\newcommand{\MlstMlpCmpBbCoarseMassOverCoarsePr}{\pending{no run yet}}  % experiments/FIGS/data/v2/index.json :: nonlinear probe bb_vs_cc: (1 - AUC) of R16_Q1_MASS over R16_Q1, run by run :: sha256 c6c4138ee5c86423 :: prose twin: two significant figures, direction in words
\newcommand{\MlstMlpCmpBbCoarseMassOverCoarsePrMax}{\pending{no run yet}}  % experiments/FIGS/data/v2/index.json :: nonlinear probe bb_vs_cc: (1 - AUC) of R16_Q1_MASS over R16_Q1, run by run :: sha256 c6c4138ee5c86423 :: prose twin of the run farther from 1
\newcommand{\MlstMlpCmpBbFzeroOverFull}{3.3 and 4.3}  % experiments/FIGS/data/v2/index.json :: nonlinear probe bb_vs_cc: (1 - AUC) of FLAV_F0 over L188, run by run :: sha256 c6c4138ee5c86423
\newcommand{\MlstMlpCmpBbFzeroOverFullPr}{at least 3.3 times the errors}  % experiments/FIGS/data/v2/index.json :: nonlinear probe bb_vs_cc: (1 - AUC) of FLAV_F0 over L188, run by run :: sha256 c6c4138ee5c86423 :: prose twin: two significant figures, direction in words
\newcommand{\MlstMlpCmpBbFzeroOverFullPrMax}{up to 4.3 times the errors}  % experiments/FIGS/data/v2/index.json :: nonlinear probe bb_vs_cc: (1 - AUC) of FLAV_F0 over L188, run by run :: sha256 c6c4138ee5c86423 :: prose twin of the run farther from 1
\newcommand{\MlstMlpCmpBbFoneOverFzero}{0.52 and 0.38}  % experiments/FIGS/data/v2/index.json :: nonlinear probe bb_vs_cc: (1 - AUC) of FLAV_F1 over FLAV_F0, run by run :: sha256 c6c4138ee5c86423
\newcommand{\MlstMlpCmpBbFoneOverFzeroPr}{at least 48\% fewer errors}  % experiments/FIGS/data/v2/index.json :: nonlinear probe bb_vs_cc: (1 - AUC) of FLAV_F1 over FLAV_F0, run by run :: sha256 c6c4138ee5c86423 :: prose twin: two significant figures, direction in words
\newcommand{\MlstMlpCmpBbFoneOverFzeroPrMax}{down to 0.38 times the errors}  % experiments/FIGS/data/v2/index.json :: nonlinear probe bb_vs_cc: (1 - AUC) of FLAV_F1 over FLAV_F0, run by run :: sha256 c6c4138ee5c86423 :: prose twin of the run farther from 1
\newcommand{\MlstMlpCmpBbFonerOverFone}{1.27 and 1.45}  % experiments/FIGS/data/v2/index.json :: nonlinear probe bb_vs_cc: (1 - AUC) of FLAV_F1R over FLAV_F1, run by run :: sha256 c6c4138ee5c86423
\newcommand{\MlstMlpCmpBbFonerOverFonePr}{at least 27\% more errors}  % experiments/FIGS/data/v2/index.json :: nonlinear probe bb_vs_cc: (1 - AUC) of FLAV_F1R over FLAV_F1, run by run :: sha256 c6c4138ee5c86423 :: prose twin: two significant figures, direction in words
\newcommand{\MlstMlpCmpBbFonerOverFonePrMax}{up to 45\% more errors}  % experiments/FIGS/data/v2/index.json :: nonlinear probe bb_vs_cc: (1 - AUC) of FLAV_F1R over FLAV_F1, run by run :: sha256 c6c4138ee5c86423 :: prose twin of the run farther from 1
\newcommand{\MlstMlpCmpBbFonerOverFzero}{0.66 and 0.55}  % experiments/FIGS/data/v2/index.json :: nonlinear probe bb_vs_cc: (1 - AUC) of FLAV_F1R over FLAV_F0, run by run :: sha256 c6c4138ee5c86423
\newcommand{\MlstMlpCmpBbFonerOverFzeroPr}{at least 34\% fewer errors}  % experiments/FIGS/data/v2/index.json :: nonlinear probe bb_vs_cc: (1 - AUC) of FLAV_F1R over FLAV_F0, run by run :: sha256 c6c4138ee5c86423 :: prose twin: two significant figures, direction in words
\newcommand{\MlstMlpCmpBbFonerOverFzeroPrMax}{up to 45\% fewer errors}  % experiments/FIGS/data/v2/index.json :: nonlinear probe bb_vs_cc: (1 - AUC) of FLAV_F1R over FLAV_F0, run by run :: sha256 c6c4138ee5c86423 :: prose twin of the run farther from 1
\newcommand{\MlstMlpCmpBbFzeroOverCoarse}{1.07 and 1.41}  % experiments/FIGS/data/v2/index.json :: nonlinear probe bb_vs_cc: (1 - AUC) of FLAV_F0 over R16_Q1, run by run :: sha256 c6c4138ee5c86423
\newcommand{\MlstMlpCmpBbFzeroOverCoarsePr}{at least 7.4\% more errors}  % experiments/FIGS/data/v2/index.json :: nonlinear probe bb_vs_cc: (1 - AUC) of FLAV_F0 over R16_Q1, run by run :: sha256 c6c4138ee5c86423 :: prose twin: two significant figures, direction in words
\newcommand{\MlstMlpCmpBbFzeroOverCoarsePrMax}{up to 41\% more errors}  % experiments/FIGS/data/v2/index.json :: nonlinear probe bb_vs_cc: (1 - AUC) of FLAV_F0 over R16_Q1, run by run :: sha256 c6c4138ee5c86423 :: prose twin of the run farther from 1
\newcommand{\MlstMlpCmpBbRandOverFull}{\pending{no run yet}}  % experiments/FIGS/data/v2/index.json :: nonlinear probe bb_vs_cc: (1 - AUC) of each random partition over the 188-class run of its seed :: sha256 c6c4138ee5c86423
\newcommand{\MlstMlpCmpBbRandOverFullPr}{\pending{no run yet}}  % experiments/FIGS/data/v2/index.json :: nonlinear probe bb_vs_cc: (1 - AUC) of each random partition over the 188-class run of its seed :: sha256 c6c4138ee5c86423 :: prose twin: two significant figures, direction in words
\newcommand{\MlstMlpCmpBbqqResOverFull}{1.011\,\ensuremath{\pm}\,0.030}  % experiments/FIGS/data/v2/index.json :: nonlinear probe bbqq_vs_ccqq: (1 - AUC) of L162 over L188, run by run :: sha256 c6c4138ee5c86423
\newcommand{\MlstMlpCmpBbqqResOverFullPr}{equal within 3\%}  % experiments/FIGS/data/v2/index.json :: nonlinear probe bbqq_vs_ccqq: (1 - AUC) of L162 over L188, run by run :: sha256 c6c4138ee5c86423 :: prose twin: two significant figures, direction in words
\newcommand{\MlstMlpCmpBbqqMidOverFull}{1.112\,\ensuremath{\pm}\,0.006}  % experiments/FIGS/data/v2/index.json :: nonlinear probe bbqq_vs_ccqq: (1 - AUC) of R42_Q1 over L188, run by run :: sha256 c6c4138ee5c86423
\newcommand{\MlstMlpCmpBbqqMidOverFullPr}{11\% more errors}  % experiments/FIGS/data/v2/index.json :: nonlinear probe bbqq_vs_ccqq: (1 - AUC) of R42_Q1 over L188, run by run :: sha256 c6c4138ee5c86423 :: prose twin: two significant figures, direction in words
\newcommand{\MlstMlpCmpBbqqCoarseOverFull}{2.93\,\ensuremath{\pm}\,0.17}  % experiments/FIGS/data/v2/index.json :: nonlinear probe bbqq_vs_ccqq: (1 - AUC) of R16_Q1 over L188, run by run :: sha256 c6c4138ee5c86423
\newcommand{\MlstMlpCmpBbqqCoarseOverFullPr}{2.9 times the errors}  % experiments/FIGS/data/v2/index.json :: nonlinear probe bbqq_vs_ccqq: (1 - AUC) of R16_Q1 over L188, run by run :: sha256 c6c4138ee5c86423 :: prose twin: two significant figures, direction in words
\newcommand{\MlstMlpCmpBbqqResMassOverRes}{\pending{no run yet}}  % experiments/FIGS/data/v2/index.json :: nonlinear probe bbqq_vs_ccqq: (1 - AUC) of L162_MASS over L162, run by run :: sha256 c6c4138ee5c86423
\newcommand{\MlstMlpCmpBbqqResMassOverResPr}{\pending{no run yet}}  % experiments/FIGS/data/v2/index.json :: nonlinear probe bbqq_vs_ccqq: (1 - AUC) of L162_MASS over L162, run by run :: sha256 c6c4138ee5c86423 :: prose twin: two significant figures, direction in words
\newcommand{\MlstMlpCmpBbqqResMassOverResPrMax}{\pending{no run yet}}  % experiments/FIGS/data/v2/index.json :: nonlinear probe bbqq_vs_ccqq: (1 - AUC) of L162_MASS over L162, run by run :: sha256 c6c4138ee5c86423 :: prose twin of the run farther from 1
\newcommand{\MlstMlpCmpBbqqCoarseMassOverCoarse}{\pending{no run yet}}  % experiments/FIGS/data/v2/index.json :: nonlinear probe bbqq_vs_ccqq: (1 - AUC) of R16_Q1_MASS over R16_Q1, run by run :: sha256 c6c4138ee5c86423
\newcommand{\MlstMlpCmpBbqqCoarseMassOverCoarsePr}{\pending{no run yet}}  % experiments/FIGS/data/v2/index.json :: nonlinear probe bbqq_vs_ccqq: (1 - AUC) of R16_Q1_MASS over R16_Q1, run by run :: sha256 c6c4138ee5c86423 :: prose twin: two significant figures, direction in words
\newcommand{\MlstMlpCmpBbqqCoarseMassOverCoarsePrMax}{\pending{no run yet}}  % experiments/FIGS/data/v2/index.json :: nonlinear probe bbqq_vs_ccqq: (1 - AUC) of R16_Q1_MASS over R16_Q1, run by run :: sha256 c6c4138ee5c86423 :: prose twin of the run farther from 1
\newcommand{\MlstMlpCmpBbqqFzeroOverFull}{3.20 and 3.78}  % experiments/FIGS/data/v2/index.json :: nonlinear probe bbqq_vs_ccqq: (1 - AUC) of FLAV_F0 over L188, run by run :: sha256 c6c4138ee5c86423
\newcommand{\MlstMlpCmpBbqqFzeroOverFullPr}{at least 3.2 times the errors}  % experiments/FIGS/data/v2/index.json :: nonlinear probe bbqq_vs_ccqq: (1 - AUC) of FLAV_F0 over L188, run by run :: sha256 c6c4138ee5c86423 :: prose twin: two significant figures, direction in words
\newcommand{\MlstMlpCmpBbqqFzeroOverFullPrMax}{up to 3.8 times the errors}  % experiments/FIGS/data/v2/index.json :: nonlinear probe bbqq_vs_ccqq: (1 - AUC) of FLAV_F0 over L188, run by run :: sha256 c6c4138ee5c86423 :: prose twin of the run farther from 1
\newcommand{\MlstMlpCmpBbqqFoneOverFzero}{0.450 and 0.363}  % experiments/FIGS/data/v2/index.json :: nonlinear probe bbqq_vs_ccqq: (1 - AUC) of FLAV_F1 over FLAV_F0, run by run :: sha256 c6c4138ee5c86423
\newcommand{\MlstMlpCmpBbqqFoneOverFzeroPr}{at most 0.45 times the errors}  % experiments/FIGS/data/v2/index.json :: nonlinear probe bbqq_vs_ccqq: (1 - AUC) of FLAV_F1 over FLAV_F0, run by run :: sha256 c6c4138ee5c86423 :: prose twin: two significant figures, direction in words
\newcommand{\MlstMlpCmpBbqqFoneOverFzeroPrMax}{down to 0.36 times the errors}  % experiments/FIGS/data/v2/index.json :: nonlinear probe bbqq_vs_ccqq: (1 - AUC) of FLAV_F1 over FLAV_F0, run by run :: sha256 c6c4138ee5c86423 :: prose twin of the run farther from 1
\newcommand{\MlstMlpCmpBbqqFonerOverFone}{1.311 and 1.343}  % experiments/FIGS/data/v2/index.json :: nonlinear probe bbqq_vs_ccqq: (1 - AUC) of FLAV_F1R over FLAV_F1, run by run :: sha256 c6c4138ee5c86423
\newcommand{\MlstMlpCmpBbqqFonerOverFonePr}{at least 31\% more errors}  % experiments/FIGS/data/v2/index.json :: nonlinear probe bbqq_vs_ccqq: (1 - AUC) of FLAV_F1R over FLAV_F1, run by run :: sha256 c6c4138ee5c86423 :: prose twin: two significant figures, direction in words
\newcommand{\MlstMlpCmpBbqqFonerOverFonePrMax}{up to 34\% more errors}  % experiments/FIGS/data/v2/index.json :: nonlinear probe bbqq_vs_ccqq: (1 - AUC) of FLAV_F1R over FLAV_F1, run by run :: sha256 c6c4138ee5c86423 :: prose twin of the run farther from 1
\newcommand{\MlstMlpCmpBbqqFonerOverFzero}{0.59 and 0.49}  % experiments/FIGS/data/v2/index.json :: nonlinear probe bbqq_vs_ccqq: (1 - AUC) of FLAV_F1R over FLAV_F0, run by run :: sha256 c6c4138ee5c86423
\newcommand{\MlstMlpCmpBbqqFonerOverFzeroPr}{at least 41\% fewer errors}  % experiments/FIGS/data/v2/index.json :: nonlinear probe bbqq_vs_ccqq: (1 - AUC) of FLAV_F1R over FLAV_F0, run by run :: sha256 c6c4138ee5c86423 :: prose twin: two significant figures, direction in words
\newcommand{\MlstMlpCmpBbqqFonerOverFzeroPrMax}{down to 0.49 times the errors}  % experiments/FIGS/data/v2/index.json :: nonlinear probe bbqq_vs_ccqq: (1 - AUC) of FLAV_F1R over FLAV_F0, run by run :: sha256 c6c4138ee5c86423 :: prose twin of the run farther from 1
\newcommand{\MlstMlpCmpBbqqFzeroOverCoarse}{1.170 and 1.236}  % experiments/FIGS/data/v2/index.json :: nonlinear probe bbqq_vs_ccqq: (1 - AUC) of FLAV_F0 over R16_Q1, run by run :: sha256 c6c4138ee5c86423
\newcommand{\MlstMlpCmpBbqqFzeroOverCoarsePr}{at least 17\% more errors}  % experiments/FIGS/data/v2/index.json :: nonlinear probe bbqq_vs_ccqq: (1 - AUC) of FLAV_F0 over R16_Q1, run by run :: sha256 c6c4138ee5c86423 :: prose twin: two significant figures, direction in words
\newcommand{\MlstMlpCmpBbqqFzeroOverCoarsePrMax}{up to 24\% more errors}  % experiments/FIGS/data/v2/index.json :: nonlinear probe bbqq_vs_ccqq: (1 - AUC) of FLAV_F0 over R16_Q1, run by run :: sha256 c6c4138ee5c86423 :: prose twin of the run farther from 1
\newcommand{\MlstMlpCmpBbqqRandOverFull}{\pending{no run yet}}  % experiments/FIGS/data/v2/index.json :: nonlinear probe bbqq_vs_ccqq: (1 - AUC) of each random partition over the 188-class run of its seed :: sha256 c6c4138ee5c86423
\newcommand{\MlstMlpCmpBbqqRandOverFullPr}{\pending{no run yet}}  % experiments/FIGS/data/v2/index.json :: nonlinear probe bbqq_vs_ccqq: (1 - AUC) of each random partition over the 188-class run of its seed :: sha256 c6c4138ee5c86423 :: prose twin: two significant figures, direction in words
\newcommand{\MlstMlpCmpQcdBbResOverFull}{1.085\,\ensuremath{\pm}\,0.030}  % experiments/FIGS/data/v2/index.json :: nonlinear probe qcd_bb_vs_cc: (1 - AUC) of L162 over L188, run by run :: sha256 c6c4138ee5c86423
\newcommand{\MlstMlpCmpQcdBbResOverFullPr}{8.5\% more errors}  % experiments/FIGS/data/v2/index.json :: nonlinear probe qcd_bb_vs_cc: (1 - AUC) of L162 over L188, run by run :: sha256 c6c4138ee5c86423 :: prose twin: two significant figures, direction in words
\newcommand{\MlstMlpCmpQcdBbMidOverFull}{1.121\,\ensuremath{\pm}\,0.017}  % experiments/FIGS/data/v2/index.json :: nonlinear probe qcd_bb_vs_cc: (1 - AUC) of R42_Q1 over L188, run by run :: sha256 c6c4138ee5c86423
\newcommand{\MlstMlpCmpQcdBbMidOverFullPr}{12\% more errors}  % experiments/FIGS/data/v2/index.json :: nonlinear probe qcd_bb_vs_cc: (1 - AUC) of R42_Q1 over L188, run by run :: sha256 c6c4138ee5c86423 :: prose twin: two significant figures, direction in words
\newcommand{\MlstMlpCmpQcdBbCoarseOverFull}{2.31\,\ensuremath{\pm}\,0.04}  % experiments/FIGS/data/v2/index.json :: nonlinear probe qcd_bb_vs_cc: (1 - AUC) of R16_Q1 over L188, run by run :: sha256 c6c4138ee5c86423
\newcommand{\MlstMlpCmpQcdBbCoarseOverFullPr}{2.3 times the errors}  % experiments/FIGS/data/v2/index.json :: nonlinear probe qcd_bb_vs_cc: (1 - AUC) of R16_Q1 over L188, run by run :: sha256 c6c4138ee5c86423 :: prose twin: two significant figures, direction in words
\newcommand{\MlstMlpCmpQcdBbResMassOverRes}{\pending{no run yet}}  % experiments/FIGS/data/v2/index.json :: nonlinear probe qcd_bb_vs_cc: (1 - AUC) of L162_MASS over L162, run by run :: sha256 c6c4138ee5c86423
\newcommand{\MlstMlpCmpQcdBbResMassOverResPr}{\pending{no run yet}}  % experiments/FIGS/data/v2/index.json :: nonlinear probe qcd_bb_vs_cc: (1 - AUC) of L162_MASS over L162, run by run :: sha256 c6c4138ee5c86423 :: prose twin: two significant figures, direction in words
\newcommand{\MlstMlpCmpQcdBbResMassOverResPrMax}{\pending{no run yet}}  % experiments/FIGS/data/v2/index.json :: nonlinear probe qcd_bb_vs_cc: (1 - AUC) of L162_MASS over L162, run by run :: sha256 c6c4138ee5c86423 :: prose twin of the run farther from 1
\newcommand{\MlstMlpCmpQcdBbCoarseMassOverCoarse}{\pending{no run yet}}  % experiments/FIGS/data/v2/index.json :: nonlinear probe qcd_bb_vs_cc: (1 - AUC) of R16_Q1_MASS over R16_Q1, run by run :: sha256 c6c4138ee5c86423
\newcommand{\MlstMlpCmpQcdBbCoarseMassOverCoarsePr}{\pending{no run yet}}  % experiments/FIGS/data/v2/index.json :: nonlinear probe qcd_bb_vs_cc: (1 - AUC) of R16_Q1_MASS over R16_Q1, run by run :: sha256 c6c4138ee5c86423 :: prose twin: two significant figures, direction in words
\newcommand{\MlstMlpCmpQcdBbCoarseMassOverCoarsePrMax}{\pending{no run yet}}  % experiments/FIGS/data/v2/index.json :: nonlinear probe qcd_bb_vs_cc: (1 - AUC) of R16_Q1_MASS over R16_Q1, run by run :: sha256 c6c4138ee5c86423 :: prose twin of the run farther from 1
\newcommand{\MlstMlpCmpQcdBbFzeroOverFull}{2.50 and 2.72}  % experiments/FIGS/data/v2/index.json :: nonlinear probe qcd_bb_vs_cc: (1 - AUC) of FLAV_F0 over L188, run by run :: sha256 c6c4138ee5c86423
\newcommand{\MlstMlpCmpQcdBbFzeroOverFullPr}{at least 2.5 times the errors}  % experiments/FIGS/data/v2/index.json :: nonlinear probe qcd_bb_vs_cc: (1 - AUC) of FLAV_F0 over L188, run by run :: sha256 c6c4138ee5c86423 :: prose twin: two significant figures, direction in words
\newcommand{\MlstMlpCmpQcdBbFzeroOverFullPrMax}{up to 2.7 times the errors}  % experiments/FIGS/data/v2/index.json :: nonlinear probe qcd_bb_vs_cc: (1 - AUC) of FLAV_F0 over L188, run by run :: sha256 c6c4138ee5c86423 :: prose twin of the run farther from 1
\newcommand{\MlstMlpCmpQcdBbFoneOverFzero}{0.573 and 0.505}  % experiments/FIGS/data/v2/index.json :: nonlinear probe qcd_bb_vs_cc: (1 - AUC) of FLAV_F1 over FLAV_F0, run by run :: sha256 c6c4138ee5c86423
\newcommand{\MlstMlpCmpQcdBbFoneOverFzeroPr}{at least 43\% fewer errors}  % experiments/FIGS/data/v2/index.json :: nonlinear probe qcd_bb_vs_cc: (1 - AUC) of FLAV_F1 over FLAV_F0, run by run :: sha256 c6c4138ee5c86423 :: prose twin: two significant figures, direction in words
\newcommand{\MlstMlpCmpQcdBbFoneOverFzeroPrMax}{up to 49\% fewer errors}  % experiments/FIGS/data/v2/index.json :: nonlinear probe qcd_bb_vs_cc: (1 - AUC) of FLAV_F1 over FLAV_F0, run by run :: sha256 c6c4138ee5c86423 :: prose twin of the run farther from 1
\newcommand{\MlstMlpCmpQcdBbFonerOverFone}{1.219 and 1.258}  % experiments/FIGS/data/v2/index.json :: nonlinear probe qcd_bb_vs_cc: (1 - AUC) of FLAV_F1R over FLAV_F1, run by run :: sha256 c6c4138ee5c86423
\newcommand{\MlstMlpCmpQcdBbFonerOverFonePr}{at least 22\% more errors}  % experiments/FIGS/data/v2/index.json :: nonlinear probe qcd_bb_vs_cc: (1 - AUC) of FLAV_F1R over FLAV_F1, run by run :: sha256 c6c4138ee5c86423 :: prose twin: two significant figures, direction in words
\newcommand{\MlstMlpCmpQcdBbFonerOverFonePrMax}{up to 26\% more errors}  % experiments/FIGS/data/v2/index.json :: nonlinear probe qcd_bb_vs_cc: (1 - AUC) of FLAV_F1R over FLAV_F1, run by run :: sha256 c6c4138ee5c86423 :: prose twin of the run farther from 1
\newcommand{\MlstMlpCmpQcdBbFonerOverFzero}{0.699 and 0.635}  % experiments/FIGS/data/v2/index.json :: nonlinear probe qcd_bb_vs_cc: (1 - AUC) of FLAV_F1R over FLAV_F0, run by run :: sha256 c6c4138ee5c86423
\newcommand{\MlstMlpCmpQcdBbFonerOverFzeroPr}{at least 30\% fewer errors}  % experiments/FIGS/data/v2/index.json :: nonlinear probe qcd_bb_vs_cc: (1 - AUC) of FLAV_F1R over FLAV_F0, run by run :: sha256 c6c4138ee5c86423 :: prose twin: two significant figures, direction in words
\newcommand{\MlstMlpCmpQcdBbFonerOverFzeroPrMax}{up to 36\% fewer errors}  % experiments/FIGS/data/v2/index.json :: nonlinear probe qcd_bb_vs_cc: (1 - AUC) of FLAV_F1R over FLAV_F0, run by run :: sha256 c6c4138ee5c86423 :: prose twin of the run farther from 1
\newcommand{\MlstMlpCmpQcdBbFzeroOverCoarse}{1.103 and 1.167}  % experiments/FIGS/data/v2/index.json :: nonlinear probe qcd_bb_vs_cc: (1 - AUC) of FLAV_F0 over R16_Q1, run by run :: sha256 c6c4138ee5c86423
\newcommand{\MlstMlpCmpQcdBbFzeroOverCoarsePr}{at least 10\% more errors}  % experiments/FIGS/data/v2/index.json :: nonlinear probe qcd_bb_vs_cc: (1 - AUC) of FLAV_F0 over R16_Q1, run by run :: sha256 c6c4138ee5c86423 :: prose twin: two significant figures, direction in words
\newcommand{\MlstMlpCmpQcdBbFzeroOverCoarsePrMax}{up to 17\% more errors}  % experiments/FIGS/data/v2/index.json :: nonlinear probe qcd_bb_vs_cc: (1 - AUC) of FLAV_F0 over R16_Q1, run by run :: sha256 c6c4138ee5c86423 :: prose twin of the run farther from 1
\newcommand{\MlstMlpCmpQcdBbRandOverFull}{\pending{no run yet}}  % experiments/FIGS/data/v2/index.json :: nonlinear probe qcd_bb_vs_cc: (1 - AUC) of each random partition over the 188-class run of its seed :: sha256 c6c4138ee5c86423
\newcommand{\MlstMlpCmpQcdBbRandOverFullPr}{\pending{no run yet}}  % experiments/FIGS/data/v2/index.json :: nonlinear probe qcd_bb_vs_cc: (1 - AUC) of each random partition over the 188-class run of its seed :: sha256 c6c4138ee5c86423 :: prose twin: two significant figures, direction in words
\newcommand{\MlstMlpCmpQcdBResOverFull}{1.050\,\ensuremath{\pm}\,0.004}  % experiments/FIGS/data/v2/index.json :: nonlinear probe qcd_b_vs_light: (1 - AUC) of L162 over L188, run by run :: sha256 c6c4138ee5c86423
\newcommand{\MlstMlpCmpQcdBResOverFullPr}{5.0\% more errors}  % experiments/FIGS/data/v2/index.json :: nonlinear probe qcd_b_vs_light: (1 - AUC) of L162 over L188, run by run :: sha256 c6c4138ee5c86423 :: prose twin: two significant figures, direction in words
\newcommand{\MlstMlpCmpQcdBMidOverFull}{1.055\,\ensuremath{\pm}\,0.005}  % experiments/FIGS/data/v2/index.json :: nonlinear probe qcd_b_vs_light: (1 - AUC) of R42_Q1 over L188, run by run :: sha256 c6c4138ee5c86423
\newcommand{\MlstMlpCmpQcdBMidOverFullPr}{5.5\% more errors}  % experiments/FIGS/data/v2/index.json :: nonlinear probe qcd_b_vs_light: (1 - AUC) of R42_Q1 over L188, run by run :: sha256 c6c4138ee5c86423 :: prose twin: two significant figures, direction in words
\newcommand{\MlstMlpCmpQcdBCoarseOverFull}{1.337\,\ensuremath{\pm}\,0.021}  % experiments/FIGS/data/v2/index.json :: nonlinear probe qcd_b_vs_light: (1 - AUC) of R16_Q1 over L188, run by run :: sha256 c6c4138ee5c86423
\newcommand{\MlstMlpCmpQcdBCoarseOverFullPr}{34\% more errors}  % experiments/FIGS/data/v2/index.json :: nonlinear probe qcd_b_vs_light: (1 - AUC) of R16_Q1 over L188, run by run :: sha256 c6c4138ee5c86423 :: prose twin: two significant figures, direction in words
\newcommand{\MlstMlpCmpQcdBResMassOverRes}{\pending{no run yet}}  % experiments/FIGS/data/v2/index.json :: nonlinear probe qcd_b_vs_light: (1 - AUC) of L162_MASS over L162, run by run :: sha256 c6c4138ee5c86423
\newcommand{\MlstMlpCmpQcdBResMassOverResPr}{\pending{no run yet}}  % experiments/FIGS/data/v2/index.json :: nonlinear probe qcd_b_vs_light: (1 - AUC) of L162_MASS over L162, run by run :: sha256 c6c4138ee5c86423 :: prose twin: two significant figures, direction in words
\newcommand{\MlstMlpCmpQcdBResMassOverResPrMax}{\pending{no run yet}}  % experiments/FIGS/data/v2/index.json :: nonlinear probe qcd_b_vs_light: (1 - AUC) of L162_MASS over L162, run by run :: sha256 c6c4138ee5c86423 :: prose twin of the run farther from 1
\newcommand{\MlstMlpCmpQcdBCoarseMassOverCoarse}{\pending{no run yet}}  % experiments/FIGS/data/v2/index.json :: nonlinear probe qcd_b_vs_light: (1 - AUC) of R16_Q1_MASS over R16_Q1, run by run :: sha256 c6c4138ee5c86423
\newcommand{\MlstMlpCmpQcdBCoarseMassOverCoarsePr}{\pending{no run yet}}  % experiments/FIGS/data/v2/index.json :: nonlinear probe qcd_b_vs_light: (1 - AUC) of R16_Q1_MASS over R16_Q1, run by run :: sha256 c6c4138ee5c86423 :: prose twin: two significant figures, direction in words
\newcommand{\MlstMlpCmpQcdBCoarseMassOverCoarsePrMax}{\pending{no run yet}}  % experiments/FIGS/data/v2/index.json :: nonlinear probe qcd_b_vs_light: (1 - AUC) of R16_Q1_MASS over R16_Q1, run by run :: sha256 c6c4138ee5c86423 :: prose twin of the run farther from 1
\newcommand{\MlstMlpCmpQcdBFzeroOverFull}{1.315 and 1.288}  % experiments/FIGS/data/v2/index.json :: nonlinear probe qcd_b_vs_light: (1 - AUC) of FLAV_F0 over L188, run by run :: sha256 c6c4138ee5c86423
\newcommand{\MlstMlpCmpQcdBFzeroOverFullPr}{at least 29\% more errors}  % experiments/FIGS/data/v2/index.json :: nonlinear probe qcd_b_vs_light: (1 - AUC) of FLAV_F0 over L188, run by run :: sha256 c6c4138ee5c86423 :: prose twin: two significant figures, direction in words
\newcommand{\MlstMlpCmpQcdBFzeroOverFullPrMax}{up to 32\% more errors}  % experiments/FIGS/data/v2/index.json :: nonlinear probe qcd_b_vs_light: (1 - AUC) of FLAV_F0 over L188, run by run :: sha256 c6c4138ee5c86423 :: prose twin of the run farther from 1
\newcommand{\MlstMlpCmpQcdBFoneOverFzero}{0.8804 and 0.8849}  % experiments/FIGS/data/v2/index.json :: nonlinear probe qcd_b_vs_light: (1 - AUC) of FLAV_F1 over FLAV_F0, run by run :: sha256 c6c4138ee5c86423
\newcommand{\MlstMlpCmpQcdBFoneOverFzeroPr}{at least 12\% fewer errors}  % experiments/FIGS/data/v2/index.json :: nonlinear probe qcd_b_vs_light: (1 - AUC) of FLAV_F1 over FLAV_F0, run by run :: sha256 c6c4138ee5c86423 :: prose twin: two significant figures, direction in words
\newcommand{\MlstMlpCmpQcdBFoneOverFzeroPrMax}{up to 12\% fewer errors}  % experiments/FIGS/data/v2/index.json :: nonlinear probe qcd_b_vs_light: (1 - AUC) of FLAV_F1 over FLAV_F0, run by run :: sha256 c6c4138ee5c86423 :: prose twin of the run farther from 1
\newcommand{\MlstMlpCmpQcdBFonerOverFone}{1.037 and 1.042}  % experiments/FIGS/data/v2/index.json :: nonlinear probe qcd_b_vs_light: (1 - AUC) of FLAV_F1R over FLAV_F1, run by run :: sha256 c6c4138ee5c86423
\newcommand{\MlstMlpCmpQcdBFonerOverFonePr}{at least 3.7\% more errors}  % experiments/FIGS/data/v2/index.json :: nonlinear probe qcd_b_vs_light: (1 - AUC) of FLAV_F1R over FLAV_F1, run by run :: sha256 c6c4138ee5c86423 :: prose twin: two significant figures, direction in words
\newcommand{\MlstMlpCmpQcdBFonerOverFonePrMax}{up to 4.2\% more errors}  % experiments/FIGS/data/v2/index.json :: nonlinear probe qcd_b_vs_light: (1 - AUC) of FLAV_F1R over FLAV_F1, run by run :: sha256 c6c4138ee5c86423 :: prose twin of the run farther from 1
\newcommand{\MlstMlpCmpQcdBFonerOverFzero}{0.9130 and 0.9217}  % experiments/FIGS/data/v2/index.json :: nonlinear probe qcd_b_vs_light: (1 - AUC) of FLAV_F1R over FLAV_F0, run by run :: sha256 c6c4138ee5c86423
\newcommand{\MlstMlpCmpQcdBFonerOverFzeroPr}{at least 7.8\% fewer errors}  % experiments/FIGS/data/v2/index.json :: nonlinear probe qcd_b_vs_light: (1 - AUC) of FLAV_F1R over FLAV_F0, run by run :: sha256 c6c4138ee5c86423 :: prose twin: two significant figures, direction in words
\newcommand{\MlstMlpCmpQcdBFonerOverFzeroPrMax}{up to 8.7\% fewer errors}  % experiments/FIGS/data/v2/index.json :: nonlinear probe qcd_b_vs_light: (1 - AUC) of FLAV_F1R over FLAV_F0, run by run :: sha256 c6c4138ee5c86423 :: prose twin of the run farther from 1
\newcommand{\MlstMlpCmpQcdBFzeroOverCoarse}{0.968 and 0.978}  % experiments/FIGS/data/v2/index.json :: nonlinear probe qcd_b_vs_light: (1 - AUC) of FLAV_F0 over R16_Q1, run by run :: sha256 c6c4138ee5c86423
\newcommand{\MlstMlpCmpQcdBFzeroOverCoarsePr}{at least 2.2\% fewer errors}  % experiments/FIGS/data/v2/index.json :: nonlinear probe qcd_b_vs_light: (1 - AUC) of FLAV_F0 over R16_Q1, run by run :: sha256 c6c4138ee5c86423 :: prose twin: two significant figures, direction in words
\newcommand{\MlstMlpCmpQcdBFzeroOverCoarsePrMax}{up to 3.2\% fewer errors}  % experiments/FIGS/data/v2/index.json :: nonlinear probe qcd_b_vs_light: (1 - AUC) of FLAV_F0 over R16_Q1, run by run :: sha256 c6c4138ee5c86423 :: prose twin of the run farther from 1
\newcommand{\MlstMlpCmpQcdBRandOverFull}{\pending{no run yet}}  % experiments/FIGS/data/v2/index.json :: nonlinear probe qcd_b_vs_light: (1 - AUC) of each random partition over the 188-class run of its seed :: sha256 c6c4138ee5c86423
\newcommand{\MlstMlpCmpQcdBRandOverFullPr}{\pending{no run yet}}  % experiments/FIGS/data/v2/index.json :: nonlinear probe qcd_b_vs_light: (1 - AUC) of each random partition over the 188-class run of its seed :: sha256 c6c4138ee5c86423 :: prose twin: two significant figures, direction in words
\newcommand{\MlstMlpCmpQcdCResOverFull}{1.039\,\ensuremath{\pm}\,0.005}  % experiments/FIGS/data/v2/index.json :: nonlinear probe qcd_c_vs_light: (1 - AUC) of L162 over L188, run by run :: sha256 c6c4138ee5c86423
\newcommand{\MlstMlpCmpQcdCResOverFullPr}{3.9\% more errors}  % experiments/FIGS/data/v2/index.json :: nonlinear probe qcd_c_vs_light: (1 - AUC) of L162 over L188, run by run :: sha256 c6c4138ee5c86423 :: prose twin: two significant figures, direction in words
\newcommand{\MlstMlpCmpQcdCMidOverFull}{1.055\,\ensuremath{\pm}\,0.006}  % experiments/FIGS/data/v2/index.json :: nonlinear probe qcd_c_vs_light: (1 - AUC) of R42_Q1 over L188, run by run :: sha256 c6c4138ee5c86423
\newcommand{\MlstMlpCmpQcdCMidOverFullPr}{5.5\% more errors}  % experiments/FIGS/data/v2/index.json :: nonlinear probe qcd_c_vs_light: (1 - AUC) of R42_Q1 over L188, run by run :: sha256 c6c4138ee5c86423 :: prose twin: two significant figures, direction in words
\newcommand{\MlstMlpCmpQcdCCoarseOverFull}{1.274\,\ensuremath{\pm}\,0.008}  % experiments/FIGS/data/v2/index.json :: nonlinear probe qcd_c_vs_light: (1 - AUC) of R16_Q1 over L188, run by run :: sha256 c6c4138ee5c86423
\newcommand{\MlstMlpCmpQcdCCoarseOverFullPr}{27\% more errors}  % experiments/FIGS/data/v2/index.json :: nonlinear probe qcd_c_vs_light: (1 - AUC) of R16_Q1 over L188, run by run :: sha256 c6c4138ee5c86423 :: prose twin: two significant figures, direction in words
\newcommand{\MlstMlpCmpQcdCResMassOverRes}{\pending{no run yet}}  % experiments/FIGS/data/v2/index.json :: nonlinear probe qcd_c_vs_light: (1 - AUC) of L162_MASS over L162, run by run :: sha256 c6c4138ee5c86423
\newcommand{\MlstMlpCmpQcdCResMassOverResPr}{\pending{no run yet}}  % experiments/FIGS/data/v2/index.json :: nonlinear probe qcd_c_vs_light: (1 - AUC) of L162_MASS over L162, run by run :: sha256 c6c4138ee5c86423 :: prose twin: two significant figures, direction in words
\newcommand{\MlstMlpCmpQcdCResMassOverResPrMax}{\pending{no run yet}}  % experiments/FIGS/data/v2/index.json :: nonlinear probe qcd_c_vs_light: (1 - AUC) of L162_MASS over L162, run by run :: sha256 c6c4138ee5c86423 :: prose twin of the run farther from 1
\newcommand{\MlstMlpCmpQcdCCoarseMassOverCoarse}{\pending{no run yet}}  % experiments/FIGS/data/v2/index.json :: nonlinear probe qcd_c_vs_light: (1 - AUC) of R16_Q1_MASS over R16_Q1, run by run :: sha256 c6c4138ee5c86423
\newcommand{\MlstMlpCmpQcdCCoarseMassOverCoarsePr}{\pending{no run yet}}  % experiments/FIGS/data/v2/index.json :: nonlinear probe qcd_c_vs_light: (1 - AUC) of R16_Q1_MASS over R16_Q1, run by run :: sha256 c6c4138ee5c86423 :: prose twin: two significant figures, direction in words
\newcommand{\MlstMlpCmpQcdCCoarseMassOverCoarsePrMax}{\pending{no run yet}}  % experiments/FIGS/data/v2/index.json :: nonlinear probe qcd_c_vs_light: (1 - AUC) of R16_Q1_MASS over R16_Q1, run by run :: sha256 c6c4138ee5c86423 :: prose twin of the run farther from 1
\newcommand{\MlstMlpCmpQcdCFzeroOverFull}{1.234 and 1.211}  % experiments/FIGS/data/v2/index.json :: nonlinear probe qcd_c_vs_light: (1 - AUC) of FLAV_F0 over L188, run by run :: sha256 c6c4138ee5c86423
\newcommand{\MlstMlpCmpQcdCFzeroOverFullPr}{at least 21\% more errors}  % experiments/FIGS/data/v2/index.json :: nonlinear probe qcd_c_vs_light: (1 - AUC) of FLAV_F0 over L188, run by run :: sha256 c6c4138ee5c86423 :: prose twin: two significant figures, direction in words
\newcommand{\MlstMlpCmpQcdCFzeroOverFullPrMax}{up to 23\% more errors}  % experiments/FIGS/data/v2/index.json :: nonlinear probe qcd_c_vs_light: (1 - AUC) of FLAV_F0 over L188, run by run :: sha256 c6c4138ee5c86423 :: prose twin of the run farther from 1
\newcommand{\MlstMlpCmpQcdCFoneOverFzero}{0.9395 and 0.9472}  % experiments/FIGS/data/v2/index.json :: nonlinear probe qcd_c_vs_light: (1 - AUC) of FLAV_F1 over FLAV_F0, run by run :: sha256 c6c4138ee5c86423
\newcommand{\MlstMlpCmpQcdCFoneOverFzeroPr}{at least 5.3\% fewer errors}  % experiments/FIGS/data/v2/index.json :: nonlinear probe qcd_c_vs_light: (1 - AUC) of FLAV_F1 over FLAV_F0, run by run :: sha256 c6c4138ee5c86423 :: prose twin: two significant figures, direction in words
\newcommand{\MlstMlpCmpQcdCFoneOverFzeroPrMax}{up to 6.0\% fewer errors}  % experiments/FIGS/data/v2/index.json :: nonlinear probe qcd_c_vs_light: (1 - AUC) of FLAV_F1 over FLAV_F0, run by run :: sha256 c6c4138ee5c86423 :: prose twin of the run farther from 1
\newcommand{\MlstMlpCmpQcdCFonerOverFone}{0.9846 and 0.9767}  % experiments/FIGS/data/v2/index.json :: nonlinear probe qcd_c_vs_light: (1 - AUC) of FLAV_F1R over FLAV_F1, run by run :: sha256 c6c4138ee5c86423
\newcommand{\MlstMlpCmpQcdCFonerOverFonePr}{at least 1.5\% fewer errors}  % experiments/FIGS/data/v2/index.json :: nonlinear probe qcd_c_vs_light: (1 - AUC) of FLAV_F1R over FLAV_F1, run by run :: sha256 c6c4138ee5c86423 :: prose twin: two significant figures, direction in words
\newcommand{\MlstMlpCmpQcdCFonerOverFonePrMax}{up to 2.3\% fewer errors}  % experiments/FIGS/data/v2/index.json :: nonlinear probe qcd_c_vs_light: (1 - AUC) of FLAV_F1R over FLAV_F1, run by run :: sha256 c6c4138ee5c86423 :: prose twin of the run farther from 1
\newcommand{\MlstMlpCmpQcdCFonerOverFzero}{0.9251 and 0.9251}  % experiments/FIGS/data/v2/index.json :: nonlinear probe qcd_c_vs_light: (1 - AUC) of FLAV_F1R over FLAV_F0, run by run :: sha256 c6c4138ee5c86423
\newcommand{\MlstMlpCmpQcdCFonerOverFzeroPr}{at least 7.5\% fewer errors}  % experiments/FIGS/data/v2/index.json :: nonlinear probe qcd_c_vs_light: (1 - AUC) of FLAV_F1R over FLAV_F0, run by run :: sha256 c6c4138ee5c86423 :: prose twin: two significant figures, direction in words
\newcommand{\MlstMlpCmpQcdCFonerOverFzeroPrMax}{up to 7.5\% fewer errors}  % experiments/FIGS/data/v2/index.json :: nonlinear probe qcd_c_vs_light: (1 - AUC) of FLAV_F1R over FLAV_F0, run by run :: sha256 c6c4138ee5c86423 :: prose twin of the run farther from 1
\newcommand{\MlstMlpCmpQcdCFzeroOverCoarse}{0.968 and 0.956}  % experiments/FIGS/data/v2/index.json :: nonlinear probe qcd_c_vs_light: (1 - AUC) of FLAV_F0 over R16_Q1, run by run :: sha256 c6c4138ee5c86423
\newcommand{\MlstMlpCmpQcdCFzeroOverCoarsePr}{at least 3.2\% fewer errors}  % experiments/FIGS/data/v2/index.json :: nonlinear probe qcd_c_vs_light: (1 - AUC) of FLAV_F0 over R16_Q1, run by run :: sha256 c6c4138ee5c86423 :: prose twin: two significant figures, direction in words
\newcommand{\MlstMlpCmpQcdCFzeroOverCoarsePrMax}{up to 4.4\% fewer errors}  % experiments/FIGS/data/v2/index.json :: nonlinear probe qcd_c_vs_light: (1 - AUC) of FLAV_F0 over R16_Q1, run by run :: sha256 c6c4138ee5c86423 :: prose twin of the run farther from 1
\newcommand{\MlstMlpCmpQcdCRandOverFull}{\pending{no run yet}}  % experiments/FIGS/data/v2/index.json :: nonlinear probe qcd_c_vs_light: (1 - AUC) of each random partition over the 188-class run of its seed :: sha256 c6c4138ee5c86423
\newcommand{\MlstMlpCmpQcdCRandOverFullPr}{\pending{no run yet}}  % experiments/FIGS/data/v2/index.json :: nonlinear probe qcd_c_vs_light: (1 - AUC) of each random partition over the 188-class run of its seed :: sha256 c6c4138ee5c86423 :: prose twin: two significant figures, direction in words
\newcommand{\MlstMlpTopRthirtyFullOverInitBn}{1.00\,\ensuremath{\pm}\,0.31}  % experiments/FIGS/data/v2/index.json :: nonlinear probe top, missing inputs zero: R30 of L188 over the BatchNorm-recomputed random network of the same index :: sha256 c6c4138ee5c86423
\newcommand{\MlstMlpTopRthirtyFullOverInitBnPr}{equal within 40\%}  % experiments/FIGS/data/v2/index.json :: nonlinear probe top, missing inputs zero: R30 of L188 over the BatchNorm-recomputed random network of the same index :: sha256 c6c4138ee5c86423 :: prose twin: two significant figures, direction in words
\newcommand{\MlstMlpTopRthirtyResOverInitBn}{0.94\,\ensuremath{\pm}\,0.26}  % experiments/FIGS/data/v2/index.json :: nonlinear probe top, missing inputs zero: R30 of L162 over the BatchNorm-recomputed random network of the same index :: sha256 c6c4138ee5c86423
\newcommand{\MlstMlpTopRthirtyResOverInitBnPr}{equal within 30\%}  % experiments/FIGS/data/v2/index.json :: nonlinear probe top, missing inputs zero: R30 of L162 over the BatchNorm-recomputed random network of the same index :: sha256 c6c4138ee5c86423 :: prose twin: two significant figures, direction in words
\newcommand{\MlstMlpTopRthirtyMidOverInitBn}{0.84\,\ensuremath{\pm}\,0.17}  % experiments/FIGS/data/v2/index.json :: nonlinear probe top, missing inputs zero: R30 of R42_Q1 over the BatchNorm-recomputed random network of the same index :: sha256 c6c4138ee5c86423
\newcommand{\MlstMlpTopRthirtyMidOverInitBnPr}{equal within 20\%}  % experiments/FIGS/data/v2/index.json :: nonlinear probe top, missing inputs zero: R30 of R42_Q1 over the BatchNorm-recomputed random network of the same index :: sha256 c6c4138ee5c86423 :: prose twin: two significant figures, direction in words
\newcommand{\MlstMlpTopRthirtyCoarseOverInitBn}{0.75\,\ensuremath{\pm}\,0.05}  % experiments/FIGS/data/v2/index.json :: nonlinear probe top, missing inputs zero: R30 of R16_Q1 over the BatchNorm-recomputed random network of the same index :: sha256 c6c4138ee5c86423
\newcommand{\MlstMlpTopRthirtyCoarseOverInitBnPr}{25\% less rejection}  % experiments/FIGS/data/v2/index.json :: nonlinear probe top, missing inputs zero: R30 of R16_Q1 over the BatchNorm-recomputed random network of the same index :: sha256 c6c4138ee5c86423 :: prose twin: two significant figures, direction in words
\newcommand{\MlstMlpQgRthirtyFullOverInitBn}{1.083\,\ensuremath{\pm}\,0.026}  % experiments/FIGS/data/v2/index.json :: nonlinear probe qg, missing inputs zero: R30 of L188 over the BatchNorm-recomputed random network of the same index :: sha256 c6c4138ee5c86423
\newcommand{\MlstMlpQgRthirtyFullOverInitBnPr}{8.3\% more rejection}  % experiments/FIGS/data/v2/index.json :: nonlinear probe qg, missing inputs zero: R30 of L188 over the BatchNorm-recomputed random network of the same index :: sha256 c6c4138ee5c86423 :: prose twin: two significant figures, direction in words
\newcommand{\MlstMlpQgRthirtyResOverInitBn}{1.010\,\ensuremath{\pm}\,0.016}  % experiments/FIGS/data/v2/index.json :: nonlinear probe qg, missing inputs zero: R30 of L162 over the BatchNorm-recomputed random network of the same index :: sha256 c6c4138ee5c86423
\newcommand{\MlstMlpQgRthirtyResOverInitBnPr}{equal within 2\%}  % experiments/FIGS/data/v2/index.json :: nonlinear probe qg, missing inputs zero: R30 of L162 over the BatchNorm-recomputed random network of the same index :: sha256 c6c4138ee5c86423 :: prose twin: two significant figures, direction in words
\newcommand{\MlstMlpQgRthirtyMidOverInitBn}{1.06\,\ensuremath{\pm}\,0.05}  % experiments/FIGS/data/v2/index.json :: nonlinear probe qg, missing inputs zero: R30 of R42_Q1 over the BatchNorm-recomputed random network of the same index :: sha256 c6c4138ee5c86423
\newcommand{\MlstMlpQgRthirtyMidOverInitBnPr}{equal within 7\%}  % experiments/FIGS/data/v2/index.json :: nonlinear probe qg, missing inputs zero: R30 of R42_Q1 over the BatchNorm-recomputed random network of the same index :: sha256 c6c4138ee5c86423 :: prose twin: two significant figures, direction in words
\newcommand{\MlstMlpQgRthirtyCoarseOverInitBn}{1.04\,\ensuremath{\pm}\,0.04}  % experiments/FIGS/data/v2/index.json :: nonlinear probe qg, missing inputs zero: R30 of R16_Q1 over the BatchNorm-recomputed random network of the same index :: sha256 c6c4138ee5c86423
\newcommand{\MlstMlpQgRthirtyCoarseOverInitBnPr}{equal within 5\%}  % experiments/FIGS/data/v2/index.json :: nonlinear probe qg, missing inputs zero: R30 of R16_Q1 over the BatchNorm-recomputed random network of the same index :: sha256 c6c4138ee5c86423 :: prose twin: two significant figures, direction in words
\newcommand{\MlstMlpQgHerwigRthirtyFullOverInitBn}{0.971\,\ensuremath{\pm}\,0.034}  % experiments/FIGS/data/v2/index.json :: nonlinear probe qg_herwig, missing inputs zero: R30 of L188 over the BatchNorm-recomputed random network of the same index :: sha256 c6c4138ee5c86423
\newcommand{\MlstMlpQgHerwigRthirtyFullOverInitBnPr}{equal within 4\%}  % experiments/FIGS/data/v2/index.json :: nonlinear probe qg_herwig, missing inputs zero: R30 of L188 over the BatchNorm-recomputed random network of the same index :: sha256 c6c4138ee5c86423 :: prose twin: two significant figures, direction in words
\newcommand{\MlstMlpQgHerwigRthirtyResOverInitBn}{0.94\,\ensuremath{\pm}\,0.05}  % experiments/FIGS/data/v2/index.json :: nonlinear probe qg_herwig, missing inputs zero: R30 of L162 over the BatchNorm-recomputed random network of the same index :: sha256 c6c4138ee5c86423
\newcommand{\MlstMlpQgHerwigRthirtyResOverInitBnPr}{equal within 6\%}  % experiments/FIGS/data/v2/index.json :: nonlinear probe qg_herwig, missing inputs zero: R30 of L162 over the BatchNorm-recomputed random network of the same index :: sha256 c6c4138ee5c86423 :: prose twin: two significant figures, direction in words
\newcommand{\MlstMlpQgHerwigRthirtyMidOverInitBn}{0.945\,\ensuremath{\pm}\,0.030}  % experiments/FIGS/data/v2/index.json :: nonlinear probe qg_herwig, missing inputs zero: R30 of R42_Q1 over the BatchNorm-recomputed random network of the same index :: sha256 c6c4138ee5c86423
\newcommand{\MlstMlpQgHerwigRthirtyMidOverInitBnPr}{equal within 6\%}  % experiments/FIGS/data/v2/index.json :: nonlinear probe qg_herwig, missing inputs zero: R30 of R42_Q1 over the BatchNorm-recomputed random network of the same index :: sha256 c6c4138ee5c86423 :: prose twin: two significant figures, direction in words
\newcommand{\MlstMlpQgHerwigRthirtyCoarseOverInitBn}{0.96\,\ensuremath{\pm}\,0.04}  % experiments/FIGS/data/v2/index.json :: nonlinear probe qg_herwig, missing inputs zero: R30 of R16_Q1 over the BatchNorm-recomputed random network of the same index :: sha256 c6c4138ee5c86423
\newcommand{\MlstMlpQgHerwigRthirtyCoarseOverInitBnPr}{equal within 5\%}  % experiments/FIGS/data/v2/index.json :: nonlinear probe qg_herwig, missing inputs zero: R30 of R16_Q1 over the BatchNorm-recomputed random network of the same index :: sha256 c6c4138ee5c86423 :: prose twin: two significant figures, direction in words
\newcommand{\MlstMlpNatJcTwoResMassOverRes}{1.016 and 1.009}  % experiments/FIGS/data/v2/index.json :: nonlinear probe, natural-composition JetClass-II (fits_v3_mlp2): (1 - AUC) of L162_MASS over L162, run by run :: sha256 c6c4138ee5c86423
\newcommand{\MlstMlpNatJcTwoResMassOverResPr}{equal within 2\%}  % experiments/FIGS/data/v2/index.json :: nonlinear probe, natural-composition JetClass-II (fits_v3_mlp2): (1 - AUC) of L162_MASS over L162, run by run :: sha256 c6c4138ee5c86423 :: prose twin: two significant figures, direction in words
\newcommand{\MlstMlpNatJcTwoResMassOverResPrMax}{equal within 2\%}  % experiments/FIGS/data/v2/index.json :: nonlinear probe, natural-composition JetClass-II (fits_v3_mlp2): (1 - AUC) of L162_MASS over L162, run by run :: sha256 c6c4138ee5c86423 :: prose twin of the run farther from 1
\newcommand{\MlstMlpNatJcTwoCoarseMassOverCoarse}{1.026 and 1.019}  % experiments/FIGS/data/v2/index.json :: nonlinear probe, natural-composition JetClass-II (fits_v3_mlp2): (1 - AUC) of R16_Q1_MASS over R16_Q1, run by run :: sha256 c6c4138ee5c86423
\newcommand{\MlstMlpNatJcTwoCoarseMassOverCoarsePr}{equal within 3\%}  % experiments/FIGS/data/v2/index.json :: nonlinear probe, natural-composition JetClass-II (fits_v3_mlp2): (1 - AUC) of R16_Q1_MASS over R16_Q1, run by run :: sha256 c6c4138ee5c86423 :: prose twin: two significant figures, direction in words
\newcommand{\MlstMlpNatJcTwoCoarseMassOverCoarsePrMax}{equal within 3\%}  % experiments/FIGS/data/v2/index.json :: nonlinear probe, natural-composition JetClass-II (fits_v3_mlp2): (1 - AUC) of R16_Q1_MASS over R16_Q1, run by run :: sha256 c6c4138ee5c86423 :: prose twin of the run farther from 1
\newcommand{\LpSubstitutionTop}{random draws}  % experiments/FIGS/data/v2/index.json :: linprobe/fits_v4_fill: lowest summed validation cross-entropy over every model, largest training set :: sha256 c6c4138ee5c86423 :: zero 8.8428; mean 8.99394; draws 7.00838
\newcommand{\LpSubstitutionQg}{the training mean}  % experiments/FIGS/data/v2/index.json :: linprobe/fits_v4_fill: lowest summed validation cross-entropy over every model, largest training set :: sha256 c6c4138ee5c86423 :: zero 13.6224; mean 13.1964; draws 13.3584
\newcommand{\MlstLpTopRthirtyFullOverInitBn}{4.67\,\ensuremath{\pm}\,0.23}  % experiments/FIGS/data/v2/index.json :: linear probe top, largest set, chosen substitution: R30 of L188 over the BatchNorm-recomputed random network of the same index :: sha256 c6c4138ee5c86423
\newcommand{\MlstLpTopRthirtyFullOverInitBnPr}{4.7 times the rejection}  % experiments/FIGS/data/v2/index.json :: linear probe top, largest set, chosen substitution: R30 of L188 over the BatchNorm-recomputed random network of the same index :: sha256 c6c4138ee5c86423 :: prose twin: two significant figures, direction in words
\newcommand{\MlstLpTopRthirtyFullOverInitBnZero}{1.04\,\ensuremath{\pm}\,0.05}  % experiments/FIGS/data/v2/index.json :: linear probe top, largest set, zero: R30 of L188 over the BatchNorm-recomputed random network of the same index :: sha256 c6c4138ee5c86423
\newcommand{\MlstLpTopRthirtyFullOverInitBnZeroPr}{equal within 5\%}  % experiments/FIGS/data/v2/index.json :: linear probe top, largest set, zero: R30 of L188 over the BatchNorm-recomputed random network of the same index :: sha256 c6c4138ee5c86423 :: prose twin: two significant figures, direction in words
\newcommand{\MlstLpTopRthirtyResOverInitBn}{4.08\,\ensuremath{\pm}\,0.28}  % experiments/FIGS/data/v2/index.json :: linear probe top, largest set, chosen substitution: R30 of L162 over the BatchNorm-recomputed random network of the same index :: sha256 c6c4138ee5c86423
\newcommand{\MlstLpTopRthirtyResOverInitBnPr}{4.1 times the rejection}  % experiments/FIGS/data/v2/index.json :: linear probe top, largest set, chosen substitution: R30 of L162 over the BatchNorm-recomputed random network of the same index :: sha256 c6c4138ee5c86423 :: prose twin: two significant figures, direction in words
\newcommand{\MlstLpTopRthirtyResOverInitBnZero}{1.33\,\ensuremath{\pm}\,0.30}  % experiments/FIGS/data/v2/index.json :: linear probe top, largest set, zero: R30 of L162 over the BatchNorm-recomputed random network of the same index :: sha256 c6c4138ee5c86423
\newcommand{\MlstLpTopRthirtyResOverInitBnZeroPr}{equal within 40\%}  % experiments/FIGS/data/v2/index.json :: linear probe top, largest set, zero: R30 of L162 over the BatchNorm-recomputed random network of the same index :: sha256 c6c4138ee5c86423 :: prose twin: two significant figures, direction in words
\newcommand{\MlstLpTopRthirtyMidOverInitBn}{4.76\,\ensuremath{\pm}\,0.15}  % experiments/FIGS/data/v2/index.json :: linear probe top, largest set, chosen substitution: R30 of R42_Q1 over the BatchNorm-recomputed random network of the same index :: sha256 c6c4138ee5c86423
\newcommand{\MlstLpTopRthirtyMidOverInitBnPr}{4.8 times the rejection}  % experiments/FIGS/data/v2/index.json :: linear probe top, largest set, chosen substitution: R30 of R42_Q1 over the BatchNorm-recomputed random network of the same index :: sha256 c6c4138ee5c86423 :: prose twin: two significant figures, direction in words
\newcommand{\MlstLpTopRthirtyMidOverInitBnZero}{1.01\,\ensuremath{\pm}\,0.28}  % experiments/FIGS/data/v2/index.json :: linear probe top, largest set, zero: R30 of R42_Q1 over the BatchNorm-recomputed random network of the same index :: sha256 c6c4138ee5c86423
\newcommand{\MlstLpTopRthirtyMidOverInitBnZeroPr}{equal within 30\%}  % experiments/FIGS/data/v2/index.json :: linear probe top, largest set, zero: R30 of R42_Q1 over the BatchNorm-recomputed random network of the same index :: sha256 c6c4138ee5c86423 :: prose twin: two significant figures, direction in words
\newcommand{\MlstLpTopRthirtyCoarseOverInitBn}{5.24\,\ensuremath{\pm}\,0.06}  % experiments/FIGS/data/v2/index.json :: linear probe top, largest set, chosen substitution: R30 of R16_Q1 over the BatchNorm-recomputed random network of the same index :: sha256 c6c4138ee5c86423
\newcommand{\MlstLpTopRthirtyCoarseOverInitBnPr}{5.2 times the rejection}  % experiments/FIGS/data/v2/index.json :: linear probe top, largest set, chosen substitution: R30 of R16_Q1 over the BatchNorm-recomputed random network of the same index :: sha256 c6c4138ee5c86423 :: prose twin: two significant figures, direction in words
\newcommand{\MlstLpTopRthirtyCoarseOverInitBnZero}{0.97\,\ensuremath{\pm}\,0.16}  % experiments/FIGS/data/v2/index.json :: linear probe top, largest set, zero: R30 of R16_Q1 over the BatchNorm-recomputed random network of the same index :: sha256 c6c4138ee5c86423
\newcommand{\MlstLpTopRthirtyCoarseOverInitBnZeroPr}{equal within 20\%}  % experiments/FIGS/data/v2/index.json :: linear probe top, largest set, zero: R30 of R16_Q1 over the BatchNorm-recomputed random network of the same index :: sha256 c6c4138ee5c86423 :: prose twin: two significant figures, direction in words
\newcommand{\MlstLpTopRthirtyResMassOverInitBn}{5.71 and 6.25}  % experiments/FIGS/data/v2/index.json :: linear probe top, largest set, chosen substitution: R30 of L162_MASS over the BatchNorm-recomputed random network of the same index :: sha256 c6c4138ee5c86423
\newcommand{\MlstLpTopRthirtyResMassOverInitBnPr}{at least 5.7 times the rejection}  % experiments/FIGS/data/v2/index.json :: linear probe top, largest set, chosen substitution: R30 of L162_MASS over the BatchNorm-recomputed random network of the same index :: sha256 c6c4138ee5c86423 :: prose twin: two significant figures, direction in words
\newcommand{\MlstLpTopRthirtyResMassOverInitBnPrMax}{up to 6.3 times the rejection}  % experiments/FIGS/data/v2/index.json :: linear probe top, largest set, chosen substitution: R30 of L162_MASS over the BatchNorm-recomputed random network of the same index :: sha256 c6c4138ee5c86423 :: prose twin of the run farther from 1
\newcommand{\MlstLpTopRthirtyResMassOverInitBnZero}{1.683 and 1.678}  % experiments/FIGS/data/v2/index.json :: linear probe top, largest set, zero: R30 of L162_MASS over the BatchNorm-recomputed random network of the same index :: sha256 c6c4138ee5c86423
\newcommand{\MlstLpTopRthirtyResMassOverInitBnZeroPr}{at least 68\% more rejection}  % experiments/FIGS/data/v2/index.json :: linear probe top, largest set, zero: R30 of L162_MASS over the BatchNorm-recomputed random network of the same index :: sha256 c6c4138ee5c86423 :: prose twin: two significant figures, direction in words
\newcommand{\MlstLpTopRthirtyResMassOverInitBnZeroPrMax}{up to 68\% more rejection}  % experiments/FIGS/data/v2/index.json :: linear probe top, largest set, zero: R30 of L162_MASS over the BatchNorm-recomputed random network of the same index :: sha256 c6c4138ee5c86423 :: prose twin of the run farther from 1
\newcommand{\MlstLpTopRthirtyCoarseMassOverInitBn}{6.73 and 7.72}  % experiments/FIGS/data/v2/index.json :: linear probe top, largest set, chosen substitution: R30 of R16_Q1_MASS over the BatchNorm-recomputed random network of the same index :: sha256 c6c4138ee5c86423
\newcommand{\MlstLpTopRthirtyCoarseMassOverInitBnPr}{at least 6.7 times the rejection}  % experiments/FIGS/data/v2/index.json :: linear probe top, largest set, chosen substitution: R30 of R16_Q1_MASS over the BatchNorm-recomputed random network of the same index :: sha256 c6c4138ee5c86423 :: prose twin: two significant figures, direction in words
\newcommand{\MlstLpTopRthirtyCoarseMassOverInitBnPrMax}{up to 7.7 times the rejection}  % experiments/FIGS/data/v2/index.json :: linear probe top, largest set, chosen substitution: R30 of R16_Q1_MASS over the BatchNorm-recomputed random network of the same index :: sha256 c6c4138ee5c86423 :: prose twin of the run farther from 1
\newcommand{\MlstLpTopRthirtyCoarseMassOverInitBnZero}{2.3 and 1.2}  % experiments/FIGS/data/v2/index.json :: linear probe top, largest set, zero: R30 of R16_Q1_MASS over the BatchNorm-recomputed random network of the same index :: sha256 c6c4138ee5c86423
\newcommand{\MlstLpTopRthirtyCoarseMassOverInitBnZeroPr}{equal within 200\%}  % experiments/FIGS/data/v2/index.json :: linear probe top, largest set, zero: R30 of R16_Q1_MASS over the BatchNorm-recomputed random network of the same index :: sha256 c6c4138ee5c86423 :: prose twin: two significant figures, direction in words
\newcommand{\MlstLpTopRthirtyCoarseMassOverInitBnZeroPrMax}{equal within 200\%}  % experiments/FIGS/data/v2/index.json :: linear probe top, largest set, zero: R30 of R16_Q1_MASS over the BatchNorm-recomputed random network of the same index :: sha256 c6c4138ee5c86423 :: prose twin of the run farther from 1
\newcommand{\MlstLpTopRthirtyResOverFull}{0.874\,\ensuremath{\pm}\,0.018}  % experiments/FIGS/data/v2/index.json :: linear probe top: R30 of L162 over the 188-class run :: sha256 c6c4138ee5c86423
\newcommand{\MlstLpTopRthirtyResOverFullPr}{13\% less rejection}  % experiments/FIGS/data/v2/index.json :: linear probe top: R30 of L162 over the 188-class run :: sha256 c6c4138ee5c86423 :: prose twin: two significant figures, direction in words
\newcommand{\MlstLpTopRthirtyResOverFullZero}{1.29\,\ensuremath{\pm}\,0.31}  % experiments/FIGS/data/v2/index.json :: linear probe top: R30 of L162 over the 188-class run :: sha256 c6c4138ee5c86423
\newcommand{\MlstLpTopRthirtyResOverFullZeroPr}{equal within 40\%}  % experiments/FIGS/data/v2/index.json :: linear probe top: R30 of L162 over the 188-class run :: sha256 c6c4138ee5c86423 :: prose twin: two significant figures, direction in words
\newcommand{\MlstLpTopRthirtyMidOverFull}{1.02\,\ensuremath{\pm}\,0.08}  % experiments/FIGS/data/v2/index.json :: linear probe top: R30 of R42_Q1 over the 188-class run :: sha256 c6c4138ee5c86423
\newcommand{\MlstLpTopRthirtyMidOverFullPr}{equal within 9\%}  % experiments/FIGS/data/v2/index.json :: linear probe top: R30 of R42_Q1 over the 188-class run :: sha256 c6c4138ee5c86423 :: prose twin: two significant figures, direction in words
\newcommand{\MlstLpTopRthirtyMidOverFullZero}{0.98\,\ensuremath{\pm}\,0.28}  % experiments/FIGS/data/v2/index.json :: linear probe top: R30 of R42_Q1 over the 188-class run :: sha256 c6c4138ee5c86423
\newcommand{\MlstLpTopRthirtyMidOverFullZeroPr}{equal within 30\%}  % experiments/FIGS/data/v2/index.json :: linear probe top: R30 of R42_Q1 over the 188-class run :: sha256 c6c4138ee5c86423 :: prose twin: two significant figures, direction in words
\newcommand{\MlstLpTopRthirtyCoarseOverFull}{1.12\,\ensuremath{\pm}\,0.06}  % experiments/FIGS/data/v2/index.json :: linear probe top: R30 of R16_Q1 over the 188-class run :: sha256 c6c4138ee5c86423
\newcommand{\MlstLpTopRthirtyCoarseOverFullPr}{12\% more rejection}  % experiments/FIGS/data/v2/index.json :: linear probe top: R30 of R16_Q1 over the 188-class run :: sha256 c6c4138ee5c86423 :: prose twin: two significant figures, direction in words
\newcommand{\MlstLpTopRthirtyCoarseOverFullZero}{0.93\,\ensuremath{\pm}\,0.15}  % experiments/FIGS/data/v2/index.json :: linear probe top: R30 of R16_Q1 over the 188-class run :: sha256 c6c4138ee5c86423
\newcommand{\MlstLpTopRthirtyCoarseOverFullZeroPr}{equal within 20\%}  % experiments/FIGS/data/v2/index.json :: linear probe top: R30 of R16_Q1 over the 188-class run :: sha256 c6c4138ee5c86423 :: prose twin: two significant figures, direction in words
\newcommand{\MlstLpTopRthirtyMeanOverZeroFull}{1.25\,\ensuremath{\pm}\,0.27}  % experiments/FIGS/data/v2/index.json :: linear probe top: R30 with mean over R30 with zero, L188 :: sha256 c6c4138ee5c86423
\newcommand{\MlstLpTopRthirtyMeanOverZeroFullPr}{equal within 30\%}  % experiments/FIGS/data/v2/index.json :: linear probe top: R30 with mean over R30 with zero, L188 :: sha256 c6c4138ee5c86423 :: prose twin: two significant figures, direction in words
\newcommand{\MlstLpTopRthirtyDrawOverZeroFull}{3.9\,\ensuremath{\pm}\,0.4}  % experiments/FIGS/data/v2/index.json :: linear probe top: R30 with draws over R30 with zero, L188 :: sha256 c6c4138ee5c86423
\newcommand{\MlstLpTopRthirtyDrawOverZeroFullPr}{3.9 times the rejection}  % experiments/FIGS/data/v2/index.json :: linear probe top: R30 with draws over R30 with zero, L188 :: sha256 c6c4138ee5c86423 :: prose twin: two significant figures, direction in words
\newcommand{\MlstLpTopRthirtyMeanOverZeroInitBn}{1.08\,\ensuremath{\pm}\,0.20}  % experiments/FIGS/data/v2/index.json :: linear probe top: R30 with mean over R30 with zero, INIT_BN :: sha256 c6c4138ee5c86423
\newcommand{\MlstLpTopRthirtyMeanOverZeroInitBnPr}{equal within 30\%}  % experiments/FIGS/data/v2/index.json :: linear probe top: R30 with mean over R30 with zero, INIT_BN :: sha256 c6c4138ee5c86423 :: prose twin: two significant figures, direction in words
\newcommand{\MlstLpTopRthirtyDrawOverZeroInitBn}{0.86\,\ensuremath{\pm}\,0.14}  % experiments/FIGS/data/v2/index.json :: linear probe top: R30 with draws over R30 with zero, INIT_BN :: sha256 c6c4138ee5c86423
\newcommand{\MlstLpTopRthirtyDrawOverZeroInitBnPr}{equal within 20\%}  % experiments/FIGS/data/v2/index.json :: linear probe top: R30 with draws over R30 with zero, INIT_BN :: sha256 c6c4138ee5c86423 :: prose twin: two significant figures, direction in words
\newcommand{\MlstLpTopRthirtyKinOverFull}{3.88}  % experiments/FIGS/data/v2/index.json :: linear probe top: R30 of the kinematic-only model over the mean of the 188-class runs (zero) :: sha256 c6c4138ee5c86423
\newcommand{\MlstLpTopRthirtyKinOverFullPr}{3.9 times the rejection}  % experiments/FIGS/data/v2/index.json :: linear probe top: R30 of the kinematic-only model over the mean of the 188-class runs (zero) :: sha256 c6c4138ee5c86423 :: prose twin: two significant figures, direction in words
\newcommand{\MlstLpTopRthirtyKin}{231}  % experiments/FIGS/data/v2/index.json :: linear probe top, largest set: R30 of the kinematic-only model (one run) :: sha256 c6c4138ee5c86423
\newcommand{\MlstLpTopRthirtyFullZero}{60\,\ensuremath{\pm}\,12}  % experiments/FIGS/data/v2/index.json :: linear probe top, largest set: R30 with zero, L188 :: sha256 c6c4138ee5c86423
\newcommand{\MlstLpTopRthirtyFullMean}{73\,\ensuremath{\pm}\,15}  % experiments/FIGS/data/v2/index.json :: linear probe top, largest set: R30 with mean, L188 :: sha256 c6c4138ee5c86423
\newcommand{\MlstLpTopRthirtyFullDraws}{228\,\ensuremath{\pm}\,25}  % experiments/FIGS/data/v2/index.json :: linear probe top, largest set: R30 with draws, L188 :: sha256 c6c4138ee5c86423
\newcommand{\MlstLpTopRthirtyInitBnZero}{58\,\ensuremath{\pm}\,13}  % experiments/FIGS/data/v2/index.json :: linear probe top, largest set: R30 with zero, INIT_BN :: sha256 c6c4138ee5c86423
\newcommand{\MlstLpTopRthirtyInitBnMean}{60.4\,\ensuremath{\pm}\,2.4}  % experiments/FIGS/data/v2/index.json :: linear probe top, largest set: R30 with mean, INIT_BN :: sha256 c6c4138ee5c86423
\newcommand{\MlstLpTopRthirtyInitBnDraws}{48.6\,\ensuremath{\pm}\,3.0}  % experiments/FIGS/data/v2/index.json :: linear probe top, largest set: R30 with draws, INIT_BN :: sha256 c6c4138ee5c86423
\newcommand{\MlstLpQgRthirtyFullOverInitBn}{1.39\,\ensuremath{\pm}\,0.06}  % experiments/FIGS/data/v2/index.json :: linear probe qg, largest set, chosen substitution: R30 of L188 over the BatchNorm-recomputed random network of the same index :: sha256 c6c4138ee5c86423
\newcommand{\MlstLpQgRthirtyFullOverInitBnPr}{39\% more rejection}  % experiments/FIGS/data/v2/index.json :: linear probe qg, largest set, chosen substitution: R30 of L188 over the BatchNorm-recomputed random network of the same index :: sha256 c6c4138ee5c86423 :: prose twin: two significant figures, direction in words
\newcommand{\MlstLpQgRthirtyFullOverInitBnZero}{1.21\,\ensuremath{\pm}\,0.15}  % experiments/FIGS/data/v2/index.json :: linear probe qg, largest set, zero: R30 of L188 over the BatchNorm-recomputed random network of the same index :: sha256 c6c4138ee5c86423
\newcommand{\MlstLpQgRthirtyFullOverInitBnZeroPr}{equal within 30\%}  % experiments/FIGS/data/v2/index.json :: linear probe qg, largest set, zero: R30 of L188 over the BatchNorm-recomputed random network of the same index :: sha256 c6c4138ee5c86423 :: prose twin: two significant figures, direction in words
\newcommand{\MlstLpQgRthirtyResOverInitBn}{1.360\,\ensuremath{\pm}\,0.025}  % experiments/FIGS/data/v2/index.json :: linear probe qg, largest set, chosen substitution: R30 of L162 over the BatchNorm-recomputed random network of the same index :: sha256 c6c4138ee5c86423
\newcommand{\MlstLpQgRthirtyResOverInitBnPr}{36\% more rejection}  % experiments/FIGS/data/v2/index.json :: linear probe qg, largest set, chosen substitution: R30 of L162 over the BatchNorm-recomputed random network of the same index :: sha256 c6c4138ee5c86423 :: prose twin: two significant figures, direction in words
\newcommand{\MlstLpQgRthirtyResOverInitBnZero}{1.07\,\ensuremath{\pm}\,0.13}  % experiments/FIGS/data/v2/index.json :: linear probe qg, largest set, zero: R30 of L162 over the BatchNorm-recomputed random network of the same index :: sha256 c6c4138ee5c86423
\newcommand{\MlstLpQgRthirtyResOverInitBnZeroPr}{equal within 20\%}  % experiments/FIGS/data/v2/index.json :: linear probe qg, largest set, zero: R30 of L162 over the BatchNorm-recomputed random network of the same index :: sha256 c6c4138ee5c86423 :: prose twin: two significant figures, direction in words
\newcommand{\MlstLpQgRthirtyMidOverInitBn}{1.357\,\ensuremath{\pm}\,0.033}  % experiments/FIGS/data/v2/index.json :: linear probe qg, largest set, chosen substitution: R30 of R42_Q1 over the BatchNorm-recomputed random network of the same index :: sha256 c6c4138ee5c86423
\newcommand{\MlstLpQgRthirtyMidOverInitBnPr}{36\% more rejection}  % experiments/FIGS/data/v2/index.json :: linear probe qg, largest set, chosen substitution: R30 of R42_Q1 over the BatchNorm-recomputed random network of the same index :: sha256 c6c4138ee5c86423 :: prose twin: two significant figures, direction in words
\newcommand{\MlstLpQgRthirtyMidOverInitBnZero}{1.12\,\ensuremath{\pm}\,0.11}  % experiments/FIGS/data/v2/index.json :: linear probe qg, largest set, zero: R30 of R42_Q1 over the BatchNorm-recomputed random network of the same index :: sha256 c6c4138ee5c86423
\newcommand{\MlstLpQgRthirtyMidOverInitBnZeroPr}{equal within 20\%}  % experiments/FIGS/data/v2/index.json :: linear probe qg, largest set, zero: R30 of R42_Q1 over the BatchNorm-recomputed random network of the same index :: sha256 c6c4138ee5c86423 :: prose twin: two significant figures, direction in words
\newcommand{\MlstLpQgRthirtyCoarseOverInitBn}{1.42\,\ensuremath{\pm}\,0.04}  % experiments/FIGS/data/v2/index.json :: linear probe qg, largest set, chosen substitution: R30 of R16_Q1 over the BatchNorm-recomputed random network of the same index :: sha256 c6c4138ee5c86423
\newcommand{\MlstLpQgRthirtyCoarseOverInitBnPr}{42\% more rejection}  % experiments/FIGS/data/v2/index.json :: linear probe qg, largest set, chosen substitution: R30 of R16_Q1 over the BatchNorm-recomputed random network of the same index :: sha256 c6c4138ee5c86423 :: prose twin: two significant figures, direction in words
\newcommand{\MlstLpQgRthirtyCoarseOverInitBnZero}{1.15\,\ensuremath{\pm}\,0.17}  % experiments/FIGS/data/v2/index.json :: linear probe qg, largest set, zero: R30 of R16_Q1 over the BatchNorm-recomputed random network of the same index :: sha256 c6c4138ee5c86423
\newcommand{\MlstLpQgRthirtyCoarseOverInitBnZeroPr}{equal within 20\%}  % experiments/FIGS/data/v2/index.json :: linear probe qg, largest set, zero: R30 of R16_Q1 over the BatchNorm-recomputed random network of the same index :: sha256 c6c4138ee5c86423 :: prose twin: two significant figures, direction in words
\newcommand{\MlstLpQgRthirtyResMassOverInitBn}{1.430 and 1.364}  % experiments/FIGS/data/v2/index.json :: linear probe qg, largest set, chosen substitution: R30 of L162_MASS over the BatchNorm-recomputed random network of the same index :: sha256 c6c4138ee5c86423
\newcommand{\MlstLpQgRthirtyResMassOverInitBnPr}{at least 36\% more rejection}  % experiments/FIGS/data/v2/index.json :: linear probe qg, largest set, chosen substitution: R30 of L162_MASS over the BatchNorm-recomputed random network of the same index :: sha256 c6c4138ee5c86423 :: prose twin: two significant figures, direction in words
\newcommand{\MlstLpQgRthirtyResMassOverInitBnPrMax}{up to 43\% more rejection}  % experiments/FIGS/data/v2/index.json :: linear probe qg, largest set, chosen substitution: R30 of L162_MASS over the BatchNorm-recomputed random network of the same index :: sha256 c6c4138ee5c86423 :: prose twin of the run farther from 1
\newcommand{\MlstLpQgRthirtyResMassOverInitBnZero}{1.213 and 1.175}  % experiments/FIGS/data/v2/index.json :: linear probe qg, largest set, zero: R30 of L162_MASS over the BatchNorm-recomputed random network of the same index :: sha256 c6c4138ee5c86423
\newcommand{\MlstLpQgRthirtyResMassOverInitBnZeroPr}{equal within 30\%}  % experiments/FIGS/data/v2/index.json :: linear probe qg, largest set, zero: R30 of L162_MASS over the BatchNorm-recomputed random network of the same index :: sha256 c6c4138ee5c86423 :: prose twin: two significant figures, direction in words
\newcommand{\MlstLpQgRthirtyResMassOverInitBnZeroPrMax}{equal within 30\%}  % experiments/FIGS/data/v2/index.json :: linear probe qg, largest set, zero: R30 of L162_MASS over the BatchNorm-recomputed random network of the same index :: sha256 c6c4138ee5c86423 :: prose twin of the run farther from 1
\newcommand{\MlstLpQgRthirtyCoarseMassOverInitBn}{1.324 and 1.394}  % experiments/FIGS/data/v2/index.json :: linear probe qg, largest set, chosen substitution: R30 of R16_Q1_MASS over the BatchNorm-recomputed random network of the same index :: sha256 c6c4138ee5c86423
\newcommand{\MlstLpQgRthirtyCoarseMassOverInitBnPr}{at least 32\% more rejection}  % experiments/FIGS/data/v2/index.json :: linear probe qg, largest set, chosen substitution: R30 of R16_Q1_MASS over the BatchNorm-recomputed random network of the same index :: sha256 c6c4138ee5c86423 :: prose twin: two significant figures, direction in words
\newcommand{\MlstLpQgRthirtyCoarseMassOverInitBnPrMax}{up to 39\% more rejection}  % experiments/FIGS/data/v2/index.json :: linear probe qg, largest set, chosen substitution: R30 of R16_Q1_MASS over the BatchNorm-recomputed random network of the same index :: sha256 c6c4138ee5c86423 :: prose twin of the run farther from 1
\newcommand{\MlstLpQgRthirtyCoarseMassOverInitBnZero}{1.224 and 1.261}  % experiments/FIGS/data/v2/index.json :: linear probe qg, largest set, zero: R30 of R16_Q1_MASS over the BatchNorm-recomputed random network of the same index :: sha256 c6c4138ee5c86423
\newcommand{\MlstLpQgRthirtyCoarseMassOverInitBnZeroPr}{equal within 30\%}  % experiments/FIGS/data/v2/index.json :: linear probe qg, largest set, zero: R30 of R16_Q1_MASS over the BatchNorm-recomputed random network of the same index :: sha256 c6c4138ee5c86423 :: prose twin: two significant figures, direction in words
\newcommand{\MlstLpQgRthirtyCoarseMassOverInitBnZeroPrMax}{equal within 30\%}  % experiments/FIGS/data/v2/index.json :: linear probe qg, largest set, zero: R30 of R16_Q1_MASS over the BatchNorm-recomputed random network of the same index :: sha256 c6c4138ee5c86423 :: prose twin of the run farther from 1
\newcommand{\MlstLpQgRthirtyResOverFull}{0.977\,\ensuremath{\pm}\,0.028}  % experiments/FIGS/data/v2/index.json :: linear probe qg: R30 of L162 over the 188-class run :: sha256 c6c4138ee5c86423
\newcommand{\MlstLpQgRthirtyResOverFullPr}{equal within 3\%}  % experiments/FIGS/data/v2/index.json :: linear probe qg: R30 of L162 over the 188-class run :: sha256 c6c4138ee5c86423 :: prose twin: two significant figures, direction in words
\newcommand{\MlstLpQgRthirtyResOverFullZero}{0.89\,\ensuremath{\pm}\,0.20}  % experiments/FIGS/data/v2/index.json :: linear probe qg: R30 of L162 over the 188-class run :: sha256 c6c4138ee5c86423
\newcommand{\MlstLpQgRthirtyResOverFullZeroPr}{equal within 20\%}  % experiments/FIGS/data/v2/index.json :: linear probe qg: R30 of L162 over the 188-class run :: sha256 c6c4138ee5c86423 :: prose twin: two significant figures, direction in words
\newcommand{\MlstLpQgRthirtyMidOverFull}{0.974\,\ensuremath{\pm}\,0.031}  % experiments/FIGS/data/v2/index.json :: linear probe qg: R30 of R42_Q1 over the 188-class run :: sha256 c6c4138ee5c86423
\newcommand{\MlstLpQgRthirtyMidOverFullPr}{equal within 4\%}  % experiments/FIGS/data/v2/index.json :: linear probe qg: R30 of R42_Q1 over the 188-class run :: sha256 c6c4138ee5c86423 :: prose twin: two significant figures, direction in words
\newcommand{\MlstLpQgRthirtyMidOverFullZero}{0.94\,\ensuremath{\pm}\,0.18}  % experiments/FIGS/data/v2/index.json :: linear probe qg: R30 of R42_Q1 over the 188-class run :: sha256 c6c4138ee5c86423
\newcommand{\MlstLpQgRthirtyMidOverFullZeroPr}{equal within 20\%}  % experiments/FIGS/data/v2/index.json :: linear probe qg: R30 of R42_Q1 over the 188-class run :: sha256 c6c4138ee5c86423 :: prose twin: two significant figures, direction in words
\newcommand{\MlstLpQgRthirtyCoarseOverFull}{1.02\,\ensuremath{\pm}\,0.06}  % experiments/FIGS/data/v2/index.json :: linear probe qg: R30 of R16_Q1 over the 188-class run :: sha256 c6c4138ee5c86423
\newcommand{\MlstLpQgRthirtyCoarseOverFullPr}{equal within 6\%}  % experiments/FIGS/data/v2/index.json :: linear probe qg: R30 of R16_Q1 over the 188-class run :: sha256 c6c4138ee5c86423 :: prose twin: two significant figures, direction in words
\newcommand{\MlstLpQgRthirtyCoarseOverFullZero}{0.95\,\ensuremath{\pm}\,0.09}  % experiments/FIGS/data/v2/index.json :: linear probe qg: R30 of R16_Q1 over the 188-class run :: sha256 c6c4138ee5c86423
\newcommand{\MlstLpQgRthirtyCoarseOverFullZeroPr}{equal within 9\%}  % experiments/FIGS/data/v2/index.json :: linear probe qg: R30 of R16_Q1 over the 188-class run :: sha256 c6c4138ee5c86423 :: prose twin: two significant figures, direction in words
\newcommand{\MlstLpQgRthirtyMeanOverZeroFull}{1.12\,\ensuremath{\pm}\,0.17}  % experiments/FIGS/data/v2/index.json :: linear probe qg: R30 with mean over R30 with zero, L188 :: sha256 c6c4138ee5c86423
\newcommand{\MlstLpQgRthirtyMeanOverZeroFullPr}{equal within 20\%}  % experiments/FIGS/data/v2/index.json :: linear probe qg: R30 with mean over R30 with zero, L188 :: sha256 c6c4138ee5c86423 :: prose twin: two significant figures, direction in words
\newcommand{\MlstLpQgRthirtyDrawOverZeroFull}{1.05\,\ensuremath{\pm}\,0.12}  % experiments/FIGS/data/v2/index.json :: linear probe qg: R30 with draws over R30 with zero, L188 :: sha256 c6c4138ee5c86423
\newcommand{\MlstLpQgRthirtyDrawOverZeroFullPr}{equal within 20\%}  % experiments/FIGS/data/v2/index.json :: linear probe qg: R30 with draws over R30 with zero, L188 :: sha256 c6c4138ee5c86423 :: prose twin: two significant figures, direction in words
\newcommand{\MlstLpQgRthirtyMeanOverZeroInitBn}{0.964\,\ensuremath{\pm}\,0.030}  % experiments/FIGS/data/v2/index.json :: linear probe qg: R30 with mean over R30 with zero, INIT_BN :: sha256 c6c4138ee5c86423
\newcommand{\MlstLpQgRthirtyMeanOverZeroInitBnPr}{equal within 4\%}  % experiments/FIGS/data/v2/index.json :: linear probe qg: R30 with mean over R30 with zero, INIT_BN :: sha256 c6c4138ee5c86423 :: prose twin: two significant figures, direction in words
\newcommand{\MlstLpQgRthirtyDrawOverZeroInitBn}{0.97\,\ensuremath{\pm}\,0.04}  % experiments/FIGS/data/v2/index.json :: linear probe qg: R30 with draws over R30 with zero, INIT_BN :: sha256 c6c4138ee5c86423
\newcommand{\MlstLpQgRthirtyDrawOverZeroInitBnPr}{equal within 4\%}  % experiments/FIGS/data/v2/index.json :: linear probe qg: R30 with draws over R30 with zero, INIT_BN :: sha256 c6c4138ee5c86423 :: prose twin: two significant figures, direction in words
\newcommand{\MlstLpQgRthirtyFullZero}{66\,\ensuremath{\pm}\,7}  % experiments/FIGS/data/v2/index.json :: linear probe qg, largest set: R30 with zero, L188 :: sha256 c6c4138ee5c86423
\newcommand{\MlstLpQgRthirtyFullMean}{72.7\,\ensuremath{\pm}\,2.4}  % experiments/FIGS/data/v2/index.json :: linear probe qg, largest set: R30 with mean, L188 :: sha256 c6c4138ee5c86423
\newcommand{\MlstLpQgRthirtyFullDraws}{68.1\,\ensuremath{\pm}\,1.8}  % experiments/FIGS/data/v2/index.json :: linear probe qg, largest set: R30 with draws, L188 :: sha256 c6c4138ee5c86423
\newcommand{\MlstLpQgRthirtyInitBnZero}{54.2\,\ensuremath{\pm}\,2.4}  % experiments/FIGS/data/v2/index.json :: linear probe qg, largest set: R30 with zero, INIT_BN :: sha256 c6c4138ee5c86423
\newcommand{\MlstLpQgRthirtyInitBnMean}{52.2\,\ensuremath{\pm}\,0.8}  % experiments/FIGS/data/v2/index.json :: linear probe qg, largest set: R30 with mean, INIT_BN :: sha256 c6c4138ee5c86423
\newcommand{\MlstLpQgRthirtyInitBnDraws}{52.4\,\ensuremath{\pm}\,0.4}  % experiments/FIGS/data/v2/index.json :: linear probe qg, largest set: R30 with draws, INIT_BN :: sha256 c6c4138ee5c86423
\newcommand{\MlstLpQgHerwigRthirtyFullOverInitBn}{1.14\,\ensuremath{\pm}\,0.10}  % experiments/FIGS/data/v2/index.json :: linear probe qg_herwig, largest set, chosen substitution: R30 of L188 over the BatchNorm-recomputed random network of the same index :: sha256 c6c4138ee5c86423
\newcommand{\MlstLpQgHerwigRthirtyFullOverInitBnPr}{equal within 20\%}  % experiments/FIGS/data/v2/index.json :: linear probe qg_herwig, largest set, chosen substitution: R30 of L188 over the BatchNorm-recomputed random network of the same index :: sha256 c6c4138ee5c86423 :: prose twin: two significant figures, direction in words
\newcommand{\MlstLpQgHerwigRthirtyFullOverInitBnZero}{1.054\,\ensuremath{\pm}\,0.028}  % experiments/FIGS/data/v2/index.json :: linear probe qg_herwig, largest set, zero: R30 of L188 over the BatchNorm-recomputed random network of the same index :: sha256 c6c4138ee5c86423
\newcommand{\MlstLpQgHerwigRthirtyFullOverInitBnZeroPr}{equal within 6\%}  % experiments/FIGS/data/v2/index.json :: linear probe qg_herwig, largest set, zero: R30 of L188 over the BatchNorm-recomputed random network of the same index :: sha256 c6c4138ee5c86423 :: prose twin: two significant figures, direction in words
\newcommand{\MlstLpQgHerwigRthirtyResOverInitBn}{1.17\,\ensuremath{\pm}\,0.04}  % experiments/FIGS/data/v2/index.json :: linear probe qg_herwig, largest set, chosen substitution: R30 of L162 over the BatchNorm-recomputed random network of the same index :: sha256 c6c4138ee5c86423
\newcommand{\MlstLpQgHerwigRthirtyResOverInitBnPr}{17\% more rejection}  % experiments/FIGS/data/v2/index.json :: linear probe qg_herwig, largest set, chosen substitution: R30 of L162 over the BatchNorm-recomputed random network of the same index :: sha256 c6c4138ee5c86423 :: prose twin: two significant figures, direction in words
\newcommand{\MlstLpQgHerwigRthirtyResOverInitBnZero}{0.98\,\ensuremath{\pm}\,0.10}  % experiments/FIGS/data/v2/index.json :: linear probe qg_herwig, largest set, zero: R30 of L162 over the BatchNorm-recomputed random network of the same index :: sha256 c6c4138ee5c86423
\newcommand{\MlstLpQgHerwigRthirtyResOverInitBnZeroPr}{equal within 10\%}  % experiments/FIGS/data/v2/index.json :: linear probe qg_herwig, largest set, zero: R30 of L162 over the BatchNorm-recomputed random network of the same index :: sha256 c6c4138ee5c86423 :: prose twin: two significant figures, direction in words
\newcommand{\MlstLpQgHerwigRthirtyMidOverInitBn}{1.12\,\ensuremath{\pm}\,0.07}  % experiments/FIGS/data/v2/index.json :: linear probe qg_herwig, largest set, chosen substitution: R30 of R42_Q1 over the BatchNorm-recomputed random network of the same index :: sha256 c6c4138ee5c86423
\newcommand{\MlstLpQgHerwigRthirtyMidOverInitBnPr}{equal within 20\%}  % experiments/FIGS/data/v2/index.json :: linear probe qg_herwig, largest set, chosen substitution: R30 of R42_Q1 over the BatchNorm-recomputed random network of the same index :: sha256 c6c4138ee5c86423 :: prose twin: two significant figures, direction in words
\newcommand{\MlstLpQgHerwigRthirtyMidOverInitBnZero}{0.99\,\ensuremath{\pm}\,0.07}  % experiments/FIGS/data/v2/index.json :: linear probe qg_herwig, largest set, zero: R30 of R42_Q1 over the BatchNorm-recomputed random network of the same index :: sha256 c6c4138ee5c86423
\newcommand{\MlstLpQgHerwigRthirtyMidOverInitBnZeroPr}{equal within 7\%}  % experiments/FIGS/data/v2/index.json :: linear probe qg_herwig, largest set, zero: R30 of R42_Q1 over the BatchNorm-recomputed random network of the same index :: sha256 c6c4138ee5c86423 :: prose twin: two significant figures, direction in words
\newcommand{\MlstLpQgHerwigRthirtyCoarseOverInitBn}{1.154\,\ensuremath{\pm}\,0.015}  % experiments/FIGS/data/v2/index.json :: linear probe qg_herwig, largest set, chosen substitution: R30 of R16_Q1 over the BatchNorm-recomputed random network of the same index :: sha256 c6c4138ee5c86423
\newcommand{\MlstLpQgHerwigRthirtyCoarseOverInitBnPr}{15\% more rejection}  % experiments/FIGS/data/v2/index.json :: linear probe qg_herwig, largest set, chosen substitution: R30 of R16_Q1 over the BatchNorm-recomputed random network of the same index :: sha256 c6c4138ee5c86423 :: prose twin: two significant figures, direction in words
\newcommand{\MlstLpQgHerwigRthirtyCoarseOverInitBnZero}{1.07\,\ensuremath{\pm}\,0.05}  % experiments/FIGS/data/v2/index.json :: linear probe qg_herwig, largest set, zero: R30 of R16_Q1 over the BatchNorm-recomputed random network of the same index :: sha256 c6c4138ee5c86423
\newcommand{\MlstLpQgHerwigRthirtyCoarseOverInitBnZeroPr}{equal within 8\%}  % experiments/FIGS/data/v2/index.json :: linear probe qg_herwig, largest set, zero: R30 of R16_Q1 over the BatchNorm-recomputed random network of the same index :: sha256 c6c4138ee5c86423 :: prose twin: two significant figures, direction in words
\newcommand{\MlstLpQgHerwigRthirtyResMassOverInitBn}{1.152 and 1.138}  % experiments/FIGS/data/v2/index.json :: linear probe qg_herwig, largest set, chosen substitution: R30 of L162_MASS over the BatchNorm-recomputed random network of the same index :: sha256 c6c4138ee5c86423
\newcommand{\MlstLpQgHerwigRthirtyResMassOverInitBnPr}{equal within 20\%}  % experiments/FIGS/data/v2/index.json :: linear probe qg_herwig, largest set, chosen substitution: R30 of L162_MASS over the BatchNorm-recomputed random network of the same index :: sha256 c6c4138ee5c86423 :: prose twin: two significant figures, direction in words
\newcommand{\MlstLpQgHerwigRthirtyResMassOverInitBnPrMax}{equal within 20\%}  % experiments/FIGS/data/v2/index.json :: linear probe qg_herwig, largest set, chosen substitution: R30 of L162_MASS over the BatchNorm-recomputed random network of the same index :: sha256 c6c4138ee5c86423 :: prose twin of the run farther from 1
\newcommand{\MlstLpQgHerwigRthirtyResMassOverInitBnZero}{1.13 and 1.03}  % experiments/FIGS/data/v2/index.json :: linear probe qg_herwig, largest set, zero: R30 of L162_MASS over the BatchNorm-recomputed random network of the same index :: sha256 c6c4138ee5c86423
\newcommand{\MlstLpQgHerwigRthirtyResMassOverInitBnZeroPr}{equal within 20\%}  % experiments/FIGS/data/v2/index.json :: linear probe qg_herwig, largest set, zero: R30 of L162_MASS over the BatchNorm-recomputed random network of the same index :: sha256 c6c4138ee5c86423 :: prose twin: two significant figures, direction in words
\newcommand{\MlstLpQgHerwigRthirtyResMassOverInitBnZeroPrMax}{equal within 20\%}  % experiments/FIGS/data/v2/index.json :: linear probe qg_herwig, largest set, zero: R30 of L162_MASS over the BatchNorm-recomputed random network of the same index :: sha256 c6c4138ee5c86423 :: prose twin of the run farther from 1
\newcommand{\MlstLpQgHerwigRthirtyCoarseMassOverInitBn}{1.146 and 1.159}  % experiments/FIGS/data/v2/index.json :: linear probe qg_herwig, largest set, chosen substitution: R30 of R16_Q1_MASS over the BatchNorm-recomputed random network of the same index :: sha256 c6c4138ee5c86423
\newcommand{\MlstLpQgHerwigRthirtyCoarseMassOverInitBnPr}{equal within 20\%}  % experiments/FIGS/data/v2/index.json :: linear probe qg_herwig, largest set, chosen substitution: R30 of R16_Q1_MASS over the BatchNorm-recomputed random network of the same index :: sha256 c6c4138ee5c86423 :: prose twin: two significant figures, direction in words
\newcommand{\MlstLpQgHerwigRthirtyCoarseMassOverInitBnPrMax}{equal within 20\%}  % experiments/FIGS/data/v2/index.json :: linear probe qg_herwig, largest set, chosen substitution: R30 of R16_Q1_MASS over the BatchNorm-recomputed random network of the same index :: sha256 c6c4138ee5c86423 :: prose twin of the run farther from 1
\newcommand{\MlstLpQgHerwigRthirtyCoarseMassOverInitBnZero}{1.070 and 1.100}  % experiments/FIGS/data/v2/index.json :: linear probe qg_herwig, largest set, zero: R30 of R16_Q1_MASS over the BatchNorm-recomputed random network of the same index :: sha256 c6c4138ee5c86423
\newcommand{\MlstLpQgHerwigRthirtyCoarseMassOverInitBnZeroPr}{equal within 20\%}  % experiments/FIGS/data/v2/index.json :: linear probe qg_herwig, largest set, zero: R30 of R16_Q1_MASS over the BatchNorm-recomputed random network of the same index :: sha256 c6c4138ee5c86423 :: prose twin: two significant figures, direction in words
\newcommand{\MlstLpQgHerwigRthirtyCoarseMassOverInitBnZeroPrMax}{equal within 20\%}  % experiments/FIGS/data/v2/index.json :: linear probe qg_herwig, largest set, zero: R30 of R16_Q1_MASS over the BatchNorm-recomputed random network of the same index :: sha256 c6c4138ee5c86423 :: prose twin of the run farther from 1
\newcommand{\MlstLpQgHerwigRthirtyResOverFull}{1.03\,\ensuremath{\pm}\,0.11}  % experiments/FIGS/data/v2/index.json :: linear probe qg_herwig: R30 of L162 over the 188-class run :: sha256 c6c4138ee5c86423
\newcommand{\MlstLpQgHerwigRthirtyResOverFullPr}{equal within 20\%}  % experiments/FIGS/data/v2/index.json :: linear probe qg_herwig: R30 of L162 over the 188-class run :: sha256 c6c4138ee5c86423 :: prose twin: two significant figures, direction in words
\newcommand{\MlstLpQgHerwigRthirtyResOverFullZero}{0.93\,\ensuremath{\pm}\,0.11}  % experiments/FIGS/data/v2/index.json :: linear probe qg_herwig: R30 of L162 over the 188-class run :: sha256 c6c4138ee5c86423
\newcommand{\MlstLpQgHerwigRthirtyResOverFullZeroPr}{equal within 20\%}  % experiments/FIGS/data/v2/index.json :: linear probe qg_herwig: R30 of L162 over the 188-class run :: sha256 c6c4138ee5c86423 :: prose twin: two significant figures, direction in words
\newcommand{\MlstLpQgHerwigRthirtyMidOverFull}{0.98\,\ensuremath{\pm}\,0.06}  % experiments/FIGS/data/v2/index.json :: linear probe qg_herwig: R30 of R42_Q1 over the 188-class run :: sha256 c6c4138ee5c86423
\newcommand{\MlstLpQgHerwigRthirtyMidOverFullPr}{equal within 6\%}  % experiments/FIGS/data/v2/index.json :: linear probe qg_herwig: R30 of R42_Q1 over the 188-class run :: sha256 c6c4138ee5c86423 :: prose twin: two significant figures, direction in words
\newcommand{\MlstLpQgHerwigRthirtyMidOverFullZero}{0.94\,\ensuremath{\pm}\,0.09}  % experiments/FIGS/data/v2/index.json :: linear probe qg_herwig: R30 of R42_Q1 over the 188-class run :: sha256 c6c4138ee5c86423
\newcommand{\MlstLpQgHerwigRthirtyMidOverFullZeroPr}{equal within 9\%}  % experiments/FIGS/data/v2/index.json :: linear probe qg_herwig: R30 of R42_Q1 over the 188-class run :: sha256 c6c4138ee5c86423 :: prose twin: two significant figures, direction in words
\newcommand{\MlstLpQgHerwigRthirtyCoarseOverFull}{1.02\,\ensuremath{\pm}\,0.09}  % experiments/FIGS/data/v2/index.json :: linear probe qg_herwig: R30 of R16_Q1 over the 188-class run :: sha256 c6c4138ee5c86423
\newcommand{\MlstLpQgHerwigRthirtyCoarseOverFullPr}{equal within 10\%}  % experiments/FIGS/data/v2/index.json :: linear probe qg_herwig: R30 of R16_Q1 over the 188-class run :: sha256 c6c4138ee5c86423 :: prose twin: two significant figures, direction in words
\newcommand{\MlstLpQgHerwigRthirtyCoarseOverFullZero}{1.018\,\ensuremath{\pm}\,0.024}  % experiments/FIGS/data/v2/index.json :: linear probe qg_herwig: R30 of R16_Q1 over the 188-class run :: sha256 c6c4138ee5c86423
\newcommand{\MlstLpQgHerwigRthirtyCoarseOverFullZeroPr}{equal within 3\%}  % experiments/FIGS/data/v2/index.json :: linear probe qg_herwig: R30 of R16_Q1 over the 188-class run :: sha256 c6c4138ee5c86423 :: prose twin: two significant figures, direction in words
\newcommand{\MlstLpQgHerwigRthirtyMeanOverZeroFull}{1.07\,\ensuremath{\pm}\,0.11}  % experiments/FIGS/data/v2/index.json :: linear probe qg_herwig: R30 with mean over R30 with zero, L188 :: sha256 c6c4138ee5c86423
\newcommand{\MlstLpQgHerwigRthirtyMeanOverZeroFullPr}{equal within 20\%}  % experiments/FIGS/data/v2/index.json :: linear probe qg_herwig: R30 with mean over R30 with zero, L188 :: sha256 c6c4138ee5c86423 :: prose twin: two significant figures, direction in words
\newcommand{\MlstLpQgHerwigRthirtyDrawOverZeroFull}{1.07\,\ensuremath{\pm}\,0.06}  % experiments/FIGS/data/v2/index.json :: linear probe qg_herwig: R30 with draws over R30 with zero, L188 :: sha256 c6c4138ee5c86423
\newcommand{\MlstLpQgHerwigRthirtyDrawOverZeroFullPr}{equal within 7\%}  % experiments/FIGS/data/v2/index.json :: linear probe qg_herwig: R30 with draws over R30 with zero, L188 :: sha256 c6c4138ee5c86423 :: prose twin: two significant figures, direction in words
\newcommand{\MlstLpQgHerwigRthirtyMeanOverZeroInitBn}{0.988\,\ensuremath{\pm}\,0.019}  % experiments/FIGS/data/v2/index.json :: linear probe qg_herwig: R30 with mean over R30 with zero, INIT_BN :: sha256 c6c4138ee5c86423
\newcommand{\MlstLpQgHerwigRthirtyMeanOverZeroInitBnPr}{equal within 2\%}  % experiments/FIGS/data/v2/index.json :: linear probe qg_herwig: R30 with mean over R30 with zero, INIT_BN :: sha256 c6c4138ee5c86423 :: prose twin: two significant figures, direction in words
\newcommand{\MlstLpQgHerwigRthirtyDrawOverZeroInitBn}{0.983\,\ensuremath{\pm}\,0.032}  % experiments/FIGS/data/v2/index.json :: linear probe qg_herwig: R30 with draws over R30 with zero, INIT_BN :: sha256 c6c4138ee5c86423
\newcommand{\MlstLpQgHerwigRthirtyDrawOverZeroInitBnPr}{equal within 4\%}  % experiments/FIGS/data/v2/index.json :: linear probe qg_herwig: R30 with draws over R30 with zero, INIT_BN :: sha256 c6c4138ee5c86423 :: prose twin: two significant figures, direction in words
\newcommand{\MlstLpQgHerwigRthirtyFullZero}{32.8\,\ensuremath{\pm}\,1.0}  % experiments/FIGS/data/v2/index.json :: linear probe qg_herwig, largest set: R30 with zero, L188 :: sha256 c6c4138ee5c86423
\newcommand{\MlstLpQgHerwigRthirtyFullMean}{35.0\,\ensuremath{\pm}\,2.6}  % experiments/FIGS/data/v2/index.json :: linear probe qg_herwig, largest set: R30 with mean, L188 :: sha256 c6c4138ee5c86423
\newcommand{\MlstLpQgHerwigRthirtyFullDraws}{35.0\,\ensuremath{\pm}\,0.8}  % experiments/FIGS/data/v2/index.json :: linear probe qg_herwig, largest set: R30 with draws, L188 :: sha256 c6c4138ee5c86423
\newcommand{\MlstLpQgHerwigRthirtyInitBnZero}{31.1\,\ensuremath{\pm}\,0.5}  % experiments/FIGS/data/v2/index.json :: linear probe qg_herwig, largest set: R30 with zero, INIT_BN :: sha256 c6c4138ee5c86423
\newcommand{\MlstLpQgHerwigRthirtyInitBnMean}{30.7\,\ensuremath{\pm}\,0.5}  % experiments/FIGS/data/v2/index.json :: linear probe qg_herwig, largest set: R30 with mean, INIT_BN :: sha256 c6c4138ee5c86423
\newcommand{\MlstLpQgHerwigRthirtyInitBnDraws}{30.6\,\ensuremath{\pm}\,0.6}  % experiments/FIGS/data/v2/index.json :: linear probe qg_herwig, largest set: R30 with draws, INIT_BN :: sha256 c6c4138ee5c86423
\newcommand{\MlstLpJcTwoTopZeroAccDropFull}{88.9\,\ensuremath{\pm}\,1.0\%}  % experiments/FIGS/data/v2/index.json :: linear probe JetClass-II, largest set: relative drop in accuracy (%) with the top data's missing slots set to zero, L188 :: sha256 c6c4138ee5c86423
\newcommand{\MlstLpJcTwoTopZeroAccDropFullPr}{89\%}  % experiments/FIGS/data/v2/index.json :: linear probe JetClass-II, largest set: relative drop in accuracy (%) with the top data's missing slots set to zero, L188 :: sha256 c6c4138ee5c86423 :: prose twin, two significant figures
\newcommand{\MlstLpJcTwoTopMeanAccDropFull}{87.9\,\ensuremath{\pm}\,1.1\%}  % experiments/FIGS/data/v2/index.json :: linear probe JetClass-II, largest set: relative drop in accuracy (%) with the top data's missing slots set to mean, L188 :: sha256 c6c4138ee5c86423
\newcommand{\MlstLpJcTwoTopMeanAccDropFullPr}{88\%}  % experiments/FIGS/data/v2/index.json :: linear probe JetClass-II, largest set: relative drop in accuracy (%) with the top data's missing slots set to mean, L188 :: sha256 c6c4138ee5c86423 :: prose twin, two significant figures
\newcommand{\MlstLpJcTwoTopDrawAccDropFull}{85.98\,\ensuremath{\pm}\,0.07\%}  % experiments/FIGS/data/v2/index.json :: linear probe JetClass-II, largest set: relative drop in accuracy (%) with the top data's missing slots set to draws, L188 :: sha256 c6c4138ee5c86423
\newcommand{\MlstLpJcTwoTopDrawAccDropFullPr}{86\%}  % experiments/FIGS/data/v2/index.json :: linear probe JetClass-II, largest set: relative drop in accuracy (%) with the top data's missing slots set to draws, L188 :: sha256 c6c4138ee5c86423 :: prose twin, two significant figures
\newcommand{\MlstLpJcTwoTopZeroAccDropRes}{89.0\,\ensuremath{\pm}\,0.9\%}  % experiments/FIGS/data/v2/index.json :: linear probe JetClass-II, largest set: relative drop in accuracy (%) with the top data's missing slots set to zero, L162 :: sha256 c6c4138ee5c86423
\newcommand{\MlstLpJcTwoTopZeroAccDropResPr}{89\%}  % experiments/FIGS/data/v2/index.json :: linear probe JetClass-II, largest set: relative drop in accuracy (%) with the top data's missing slots set to zero, L162 :: sha256 c6c4138ee5c86423 :: prose twin, two significant figures
\newcommand{\MlstLpJcTwoTopMeanAccDropRes}{87.8\,\ensuremath{\pm}\,0.7\%}  % experiments/FIGS/data/v2/index.json :: linear probe JetClass-II, largest set: relative drop in accuracy (%) with the top data's missing slots set to mean, L162 :: sha256 c6c4138ee5c86423
\newcommand{\MlstLpJcTwoTopMeanAccDropResPr}{88\%}  % experiments/FIGS/data/v2/index.json :: linear probe JetClass-II, largest set: relative drop in accuracy (%) with the top data's missing slots set to mean, L162 :: sha256 c6c4138ee5c86423 :: prose twin, two significant figures
\newcommand{\MlstLpJcTwoTopDrawAccDropRes}{85.89\,\ensuremath{\pm}\,0.11\%}  % experiments/FIGS/data/v2/index.json :: linear probe JetClass-II, largest set: relative drop in accuracy (%) with the top data's missing slots set to draws, L162 :: sha256 c6c4138ee5c86423
\newcommand{\MlstLpJcTwoTopDrawAccDropResPr}{86\%}  % experiments/FIGS/data/v2/index.json :: linear probe JetClass-II, largest set: relative drop in accuracy (%) with the top data's missing slots set to draws, L162 :: sha256 c6c4138ee5c86423 :: prose twin, two significant figures
\newcommand{\MlstLpJcTwoTopZeroAccDropMid}{87.81\,\ensuremath{\pm}\,0.28\%}  % experiments/FIGS/data/v2/index.json :: linear probe JetClass-II, largest set: relative drop in accuracy (%) with the top data's missing slots set to zero, R42_Q1 :: sha256 c6c4138ee5c86423
\newcommand{\MlstLpJcTwoTopZeroAccDropMidPr}{88\%}  % experiments/FIGS/data/v2/index.json :: linear probe JetClass-II, largest set: relative drop in accuracy (%) with the top data's missing slots set to zero, R42_Q1 :: sha256 c6c4138ee5c86423 :: prose twin, two significant figures
\newcommand{\MlstLpJcTwoTopMeanAccDropMid}{86.7\,\ensuremath{\pm}\,0.9\%}  % experiments/FIGS/data/v2/index.json :: linear probe JetClass-II, largest set: relative drop in accuracy (%) with the top data's missing slots set to mean, R42_Q1 :: sha256 c6c4138ee5c86423
\newcommand{\MlstLpJcTwoTopMeanAccDropMidPr}{87\%}  % experiments/FIGS/data/v2/index.json :: linear probe JetClass-II, largest set: relative drop in accuracy (%) with the top data's missing slots set to mean, R42_Q1 :: sha256 c6c4138ee5c86423 :: prose twin, two significant figures
\newcommand{\MlstLpJcTwoTopDrawAccDropMid}{83.17\,\ensuremath{\pm}\,0.29\%}  % experiments/FIGS/data/v2/index.json :: linear probe JetClass-II, largest set: relative drop in accuracy (%) with the top data's missing slots set to draws, R42_Q1 :: sha256 c6c4138ee5c86423
\newcommand{\MlstLpJcTwoTopDrawAccDropMidPr}{83\%}  % experiments/FIGS/data/v2/index.json :: linear probe JetClass-II, largest set: relative drop in accuracy (%) with the top data's missing slots set to draws, R42_Q1 :: sha256 c6c4138ee5c86423 :: prose twin, two significant figures
\newcommand{\MlstLpJcTwoTopZeroAccDropCoarse}{86.7\,\ensuremath{\pm}\,0.7\%}  % experiments/FIGS/data/v2/index.json :: linear probe JetClass-II, largest set: relative drop in accuracy (%) with the top data's missing slots set to zero, R16_Q1 :: sha256 c6c4138ee5c86423
\newcommand{\MlstLpJcTwoTopZeroAccDropCoarsePr}{87\%}  % experiments/FIGS/data/v2/index.json :: linear probe JetClass-II, largest set: relative drop in accuracy (%) with the top data's missing slots set to zero, R16_Q1 :: sha256 c6c4138ee5c86423 :: prose twin, two significant figures
\newcommand{\MlstLpJcTwoTopMeanAccDropCoarse}{86.1\,\ensuremath{\pm}\,0.8\%}  % experiments/FIGS/data/v2/index.json :: linear probe JetClass-II, largest set: relative drop in accuracy (%) with the top data's missing slots set to mean, R16_Q1 :: sha256 c6c4138ee5c86423
\newcommand{\MlstLpJcTwoTopMeanAccDropCoarsePr}{86\%}  % experiments/FIGS/data/v2/index.json :: linear probe JetClass-II, largest set: relative drop in accuracy (%) with the top data's missing slots set to mean, R16_Q1 :: sha256 c6c4138ee5c86423 :: prose twin, two significant figures
\newcommand{\MlstLpJcTwoTopDrawAccDropCoarse}{81.4\,\ensuremath{\pm}\,0.5\%}  % experiments/FIGS/data/v2/index.json :: linear probe JetClass-II, largest set: relative drop in accuracy (%) with the top data's missing slots set to draws, R16_Q1 :: sha256 c6c4138ee5c86423
\newcommand{\MlstLpJcTwoTopDrawAccDropCoarsePr}{81\%}  % experiments/FIGS/data/v2/index.json :: linear probe JetClass-II, largest set: relative drop in accuracy (%) with the top data's missing slots set to draws, R16_Q1 :: sha256 c6c4138ee5c86423 :: prose twin, two significant figures
\newcommand{\MlstLpJcTwoTopZeroAccDropResMass}{89.2, 88.3\%}  % experiments/FIGS/data/v2/index.json :: linear probe JetClass-II, largest set: relative drop in accuracy (%) with the top data's missing slots set to zero, L162_MASS :: sha256 c6c4138ee5c86423
\newcommand{\MlstLpJcTwoTopZeroAccDropResMassPr}{88--89\%}  % experiments/FIGS/data/v2/index.json :: linear probe JetClass-II, largest set: relative drop in accuracy (%) with the top data's missing slots set to zero, L162_MASS :: sha256 c6c4138ee5c86423 :: prose twin, two significant figures
\newcommand{\MlstLpJcTwoTopMeanAccDropResMass}{88.5, 87.8\%}  % experiments/FIGS/data/v2/index.json :: linear probe JetClass-II, largest set: relative drop in accuracy (%) with the top data's missing slots set to mean, L162_MASS :: sha256 c6c4138ee5c86423
\newcommand{\MlstLpJcTwoTopMeanAccDropResMassPr}{88--88\%}  % experiments/FIGS/data/v2/index.json :: linear probe JetClass-II, largest set: relative drop in accuracy (%) with the top data's missing slots set to mean, L162_MASS :: sha256 c6c4138ee5c86423 :: prose twin, two significant figures
\newcommand{\MlstLpJcTwoTopDrawAccDropResMass}{85.4, 85.4\%}  % experiments/FIGS/data/v2/index.json :: linear probe JetClass-II, largest set: relative drop in accuracy (%) with the top data's missing slots set to draws, L162_MASS :: sha256 c6c4138ee5c86423
\newcommand{\MlstLpJcTwoTopDrawAccDropResMassPr}{85--85\%}  % experiments/FIGS/data/v2/index.json :: linear probe JetClass-II, largest set: relative drop in accuracy (%) with the top data's missing slots set to draws, L162_MASS :: sha256 c6c4138ee5c86423 :: prose twin, two significant figures
\newcommand{\MlstLpJcTwoTopZeroAccDropCoarseMass}{84.8, 87.4\%}  % experiments/FIGS/data/v2/index.json :: linear probe JetClass-II, largest set: relative drop in accuracy (%) with the top data's missing slots set to zero, R16_Q1_MASS :: sha256 c6c4138ee5c86423
\newcommand{\MlstLpJcTwoTopZeroAccDropCoarseMassPr}{85--87\%}  % experiments/FIGS/data/v2/index.json :: linear probe JetClass-II, largest set: relative drop in accuracy (%) with the top data's missing slots set to zero, R16_Q1_MASS :: sha256 c6c4138ee5c86423 :: prose twin, two significant figures
\newcommand{\MlstLpJcTwoTopMeanAccDropCoarseMass}{85.9, 86.9\%}  % experiments/FIGS/data/v2/index.json :: linear probe JetClass-II, largest set: relative drop in accuracy (%) with the top data's missing slots set to mean, R16_Q1_MASS :: sha256 c6c4138ee5c86423
\newcommand{\MlstLpJcTwoTopMeanAccDropCoarseMassPr}{86--87\%}  % experiments/FIGS/data/v2/index.json :: linear probe JetClass-II, largest set: relative drop in accuracy (%) with the top data's missing slots set to mean, R16_Q1_MASS :: sha256 c6c4138ee5c86423 :: prose twin, two significant figures
\newcommand{\MlstLpJcTwoTopDrawAccDropCoarseMass}{79.1, 81.7\%}  % experiments/FIGS/data/v2/index.json :: linear probe JetClass-II, largest set: relative drop in accuracy (%) with the top data's missing slots set to draws, R16_Q1_MASS :: sha256 c6c4138ee5c86423
\newcommand{\MlstLpJcTwoTopDrawAccDropCoarseMassPr}{79--82\%}  % experiments/FIGS/data/v2/index.json :: linear probe JetClass-II, largest set: relative drop in accuracy (%) with the top data's missing slots set to draws, R16_Q1_MASS :: sha256 c6c4138ee5c86423 :: prose twin, two significant figures
\newcommand{\MlstLpJcTwoTopZeroAccDropFzero}{88.9, 88.5\%}  % experiments/FIGS/data/v2/index.json :: linear probe JetClass-II, largest set: relative drop in accuracy (%) with the top data's missing slots set to zero, FLAV_F0 :: sha256 c6c4138ee5c86423
\newcommand{\MlstLpJcTwoTopZeroAccDropFzeroPr}{88--89\%}  % experiments/FIGS/data/v2/index.json :: linear probe JetClass-II, largest set: relative drop in accuracy (%) with the top data's missing slots set to zero, FLAV_F0 :: sha256 c6c4138ee5c86423 :: prose twin, two significant figures
\newcommand{\MlstLpJcTwoTopMeanAccDropFzero}{88.3, 87.0\%}  % experiments/FIGS/data/v2/index.json :: linear probe JetClass-II, largest set: relative drop in accuracy (%) with the top data's missing slots set to mean, FLAV_F0 :: sha256 c6c4138ee5c86423
\newcommand{\MlstLpJcTwoTopMeanAccDropFzeroPr}{87--88\%}  % experiments/FIGS/data/v2/index.json :: linear probe JetClass-II, largest set: relative drop in accuracy (%) with the top data's missing slots set to mean, FLAV_F0 :: sha256 c6c4138ee5c86423 :: prose twin, two significant figures
\newcommand{\MlstLpJcTwoTopDrawAccDropFzero}{83.3, 83.9\%}  % experiments/FIGS/data/v2/index.json :: linear probe JetClass-II, largest set: relative drop in accuracy (%) with the top data's missing slots set to draws, FLAV_F0 :: sha256 c6c4138ee5c86423
\newcommand{\MlstLpJcTwoTopDrawAccDropFzeroPr}{83--84\%}  % experiments/FIGS/data/v2/index.json :: linear probe JetClass-II, largest set: relative drop in accuracy (%) with the top data's missing slots set to draws, FLAV_F0 :: sha256 c6c4138ee5c86423 :: prose twin, two significant figures
\newcommand{\MlstLpJcTwoTopZeroAccDropFone}{87.8, 87.6\%}  % experiments/FIGS/data/v2/index.json :: linear probe JetClass-II, largest set: relative drop in accuracy (%) with the top data's missing slots set to zero, FLAV_F1 :: sha256 c6c4138ee5c86423
\newcommand{\MlstLpJcTwoTopZeroAccDropFonePr}{88--88\%}  % experiments/FIGS/data/v2/index.json :: linear probe JetClass-II, largest set: relative drop in accuracy (%) with the top data's missing slots set to zero, FLAV_F1 :: sha256 c6c4138ee5c86423 :: prose twin, two significant figures
\newcommand{\MlstLpJcTwoTopMeanAccDropFone}{88.9, 86.9\%}  % experiments/FIGS/data/v2/index.json :: linear probe JetClass-II, largest set: relative drop in accuracy (%) with the top data's missing slots set to mean, FLAV_F1 :: sha256 c6c4138ee5c86423
\newcommand{\MlstLpJcTwoTopMeanAccDropFonePr}{87--89\%}  % experiments/FIGS/data/v2/index.json :: linear probe JetClass-II, largest set: relative drop in accuracy (%) with the top data's missing slots set to mean, FLAV_F1 :: sha256 c6c4138ee5c86423 :: prose twin, two significant figures
\newcommand{\MlstLpJcTwoTopDrawAccDropFone}{84.4, 85.2\%}  % experiments/FIGS/data/v2/index.json :: linear probe JetClass-II, largest set: relative drop in accuracy (%) with the top data's missing slots set to draws, FLAV_F1 :: sha256 c6c4138ee5c86423
\newcommand{\MlstLpJcTwoTopDrawAccDropFonePr}{84--85\%}  % experiments/FIGS/data/v2/index.json :: linear probe JetClass-II, largest set: relative drop in accuracy (%) with the top data's missing slots set to draws, FLAV_F1 :: sha256 c6c4138ee5c86423 :: prose twin, two significant figures
\newcommand{\MlstLpJcTwoTopZeroAccDropFoner}{89.6, 90.3\%}  % experiments/FIGS/data/v2/index.json :: linear probe JetClass-II, largest set: relative drop in accuracy (%) with the top data's missing slots set to zero, FLAV_F1R :: sha256 c6c4138ee5c86423
\newcommand{\MlstLpJcTwoTopZeroAccDropFonerPr}{90--90\%}  % experiments/FIGS/data/v2/index.json :: linear probe JetClass-II, largest set: relative drop in accuracy (%) with the top data's missing slots set to zero, FLAV_F1R :: sha256 c6c4138ee5c86423 :: prose twin, two significant figures
\newcommand{\MlstLpJcTwoTopMeanAccDropFoner}{89.4, 89.7\%}  % experiments/FIGS/data/v2/index.json :: linear probe JetClass-II, largest set: relative drop in accuracy (%) with the top data's missing slots set to mean, FLAV_F1R :: sha256 c6c4138ee5c86423
\newcommand{\MlstLpJcTwoTopMeanAccDropFonerPr}{89--90\%}  % experiments/FIGS/data/v2/index.json :: linear probe JetClass-II, largest set: relative drop in accuracy (%) with the top data's missing slots set to mean, FLAV_F1R :: sha256 c6c4138ee5c86423 :: prose twin, two significant figures
\newcommand{\MlstLpJcTwoTopDrawAccDropFoner}{84.9, 85.6\%}  % experiments/FIGS/data/v2/index.json :: linear probe JetClass-II, largest set: relative drop in accuracy (%) with the top data's missing slots set to draws, FLAV_F1R :: sha256 c6c4138ee5c86423
\newcommand{\MlstLpJcTwoTopDrawAccDropFonerPr}{85--86\%}  % experiments/FIGS/data/v2/index.json :: linear probe JetClass-II, largest set: relative drop in accuracy (%) with the top data's missing slots set to draws, FLAV_F1R :: sha256 c6c4138ee5c86423 :: prose twin, two significant figures
\newcommand{\MlstLpJcTwoTopZeroAccDropInitBn}{56.4\,\ensuremath{\pm}\,1.4\%}  % experiments/FIGS/data/v2/index.json :: linear probe JetClass-II, largest set: relative drop in accuracy (%) with the top data's missing slots set to zero, INIT_BN :: sha256 c6c4138ee5c86423
\newcommand{\MlstLpJcTwoTopZeroAccDropInitBnPr}{56\%}  % experiments/FIGS/data/v2/index.json :: linear probe JetClass-II, largest set: relative drop in accuracy (%) with the top data's missing slots set to zero, INIT_BN :: sha256 c6c4138ee5c86423 :: prose twin, two significant figures
\newcommand{\MlstLpJcTwoTopMeanAccDropInitBn}{59.8\,\ensuremath{\pm}\,0.8\%}  % experiments/FIGS/data/v2/index.json :: linear probe JetClass-II, largest set: relative drop in accuracy (%) with the top data's missing slots set to mean, INIT_BN :: sha256 c6c4138ee5c86423
\newcommand{\MlstLpJcTwoTopMeanAccDropInitBnPr}{60\%}  % experiments/FIGS/data/v2/index.json :: linear probe JetClass-II, largest set: relative drop in accuracy (%) with the top data's missing slots set to mean, INIT_BN :: sha256 c6c4138ee5c86423 :: prose twin, two significant figures
\newcommand{\MlstLpJcTwoTopDrawAccDropInitBn}{59.8\,\ensuremath{\pm}\,1.2\%}  % experiments/FIGS/data/v2/index.json :: linear probe JetClass-II, largest set: relative drop in accuracy (%) with the top data's missing slots set to draws, INIT_BN :: sha256 c6c4138ee5c86423
\newcommand{\MlstLpJcTwoTopDrawAccDropInitBnPr}{60\%}  % experiments/FIGS/data/v2/index.json :: linear probe JetClass-II, largest set: relative drop in accuracy (%) with the top data's missing slots set to draws, INIT_BN :: sha256 c6c4138ee5c86423 :: prose twin, two significant figures
\newcommand{\MlstLpJcTwoQgZeroAccDropFull}{61.2\,\ensuremath{\pm}\,3.0\%}  % experiments/FIGS/data/v2/index.json :: linear probe JetClass-II, largest set: relative drop in accuracy (%) with the qg data's missing slots set to zero, L188 :: sha256 c6c4138ee5c86423
\newcommand{\MlstLpJcTwoQgZeroAccDropFullPr}{61\%}  % experiments/FIGS/data/v2/index.json :: linear probe JetClass-II, largest set: relative drop in accuracy (%) with the qg data's missing slots set to zero, L188 :: sha256 c6c4138ee5c86423 :: prose twin, two significant figures
\newcommand{\MlstLpJcTwoQgMeanAccDropFull}{48.8\,\ensuremath{\pm}\,1.1\%}  % experiments/FIGS/data/v2/index.json :: linear probe JetClass-II, largest set: relative drop in accuracy (%) with the qg data's missing slots set to mean, L188 :: sha256 c6c4138ee5c86423
\newcommand{\MlstLpJcTwoQgMeanAccDropFullPr}{49\%}  % experiments/FIGS/data/v2/index.json :: linear probe JetClass-II, largest set: relative drop in accuracy (%) with the qg data's missing slots set to mean, L188 :: sha256 c6c4138ee5c86423 :: prose twin, two significant figures
\newcommand{\MlstLpJcTwoQgDrawAccDropFull}{54.30\,\ensuremath{\pm}\,0.21\%}  % experiments/FIGS/data/v2/index.json :: linear probe JetClass-II, largest set: relative drop in accuracy (%) with the qg data's missing slots set to draws, L188 :: sha256 c6c4138ee5c86423
\newcommand{\MlstLpJcTwoQgDrawAccDropFullPr}{54\%}  % experiments/FIGS/data/v2/index.json :: linear probe JetClass-II, largest set: relative drop in accuracy (%) with the qg data's missing slots set to draws, L188 :: sha256 c6c4138ee5c86423 :: prose twin, two significant figures
\newcommand{\MlstLpJcTwoQgZeroAccDropRes}{64\,\ensuremath{\pm}\,5\%}  % experiments/FIGS/data/v2/index.json :: linear probe JetClass-II, largest set: relative drop in accuracy (%) with the qg data's missing slots set to zero, L162 :: sha256 c6c4138ee5c86423
\newcommand{\MlstLpJcTwoQgZeroAccDropResPr}{64\%}  % experiments/FIGS/data/v2/index.json :: linear probe JetClass-II, largest set: relative drop in accuracy (%) with the qg data's missing slots set to zero, L162 :: sha256 c6c4138ee5c86423 :: prose twin, two significant figures
\newcommand{\MlstLpJcTwoQgMeanAccDropRes}{49.5\,\ensuremath{\pm}\,0.5\%}  % experiments/FIGS/data/v2/index.json :: linear probe JetClass-II, largest set: relative drop in accuracy (%) with the qg data's missing slots set to mean, L162 :: sha256 c6c4138ee5c86423
\newcommand{\MlstLpJcTwoQgMeanAccDropResPr}{49\%}  % experiments/FIGS/data/v2/index.json :: linear probe JetClass-II, largest set: relative drop in accuracy (%) with the qg data's missing slots set to mean, L162 :: sha256 c6c4138ee5c86423 :: prose twin, two significant figures
\newcommand{\MlstLpJcTwoQgDrawAccDropRes}{54.6\,\ensuremath{\pm}\,0.6\%}  % experiments/FIGS/data/v2/index.json :: linear probe JetClass-II, largest set: relative drop in accuracy (%) with the qg data's missing slots set to draws, L162 :: sha256 c6c4138ee5c86423
\newcommand{\MlstLpJcTwoQgDrawAccDropResPr}{55\%}  % experiments/FIGS/data/v2/index.json :: linear probe JetClass-II, largest set: relative drop in accuracy (%) with the qg data's missing slots set to draws, L162 :: sha256 c6c4138ee5c86423 :: prose twin, two significant figures
\newcommand{\MlstLpJcTwoQgZeroAccDropMid}{68\,\ensuremath{\pm}\,5\%}  % experiments/FIGS/data/v2/index.json :: linear probe JetClass-II, largest set: relative drop in accuracy (%) with the qg data's missing slots set to zero, R42_Q1 :: sha256 c6c4138ee5c86423
\newcommand{\MlstLpJcTwoQgZeroAccDropMidPr}{68\%}  % experiments/FIGS/data/v2/index.json :: linear probe JetClass-II, largest set: relative drop in accuracy (%) with the qg data's missing slots set to zero, R42_Q1 :: sha256 c6c4138ee5c86423 :: prose twin, two significant figures
\newcommand{\MlstLpJcTwoQgMeanAccDropMid}{47.6\,\ensuremath{\pm}\,1.6\%}  % experiments/FIGS/data/v2/index.json :: linear probe JetClass-II, largest set: relative drop in accuracy (%) with the qg data's missing slots set to mean, R42_Q1 :: sha256 c6c4138ee5c86423
\newcommand{\MlstLpJcTwoQgMeanAccDropMidPr}{48\%}  % experiments/FIGS/data/v2/index.json :: linear probe JetClass-II, largest set: relative drop in accuracy (%) with the qg data's missing slots set to mean, R42_Q1 :: sha256 c6c4138ee5c86423 :: prose twin, two significant figures
\newcommand{\MlstLpJcTwoQgDrawAccDropMid}{53.5\,\ensuremath{\pm}\,1.4\%}  % experiments/FIGS/data/v2/index.json :: linear probe JetClass-II, largest set: relative drop in accuracy (%) with the qg data's missing slots set to draws, R42_Q1 :: sha256 c6c4138ee5c86423
\newcommand{\MlstLpJcTwoQgDrawAccDropMidPr}{54\%}  % experiments/FIGS/data/v2/index.json :: linear probe JetClass-II, largest set: relative drop in accuracy (%) with the qg data's missing slots set to draws, R42_Q1 :: sha256 c6c4138ee5c86423 :: prose twin, two significant figures
\newcommand{\MlstLpJcTwoQgZeroAccDropCoarse}{70.9\,\ensuremath{\pm}\,2.3\%}  % experiments/FIGS/data/v2/index.json :: linear probe JetClass-II, largest set: relative drop in accuracy (%) with the qg data's missing slots set to zero, R16_Q1 :: sha256 c6c4138ee5c86423
\newcommand{\MlstLpJcTwoQgZeroAccDropCoarsePr}{71\%}  % experiments/FIGS/data/v2/index.json :: linear probe JetClass-II, largest set: relative drop in accuracy (%) with the qg data's missing slots set to zero, R16_Q1 :: sha256 c6c4138ee5c86423 :: prose twin, two significant figures
\newcommand{\MlstLpJcTwoQgMeanAccDropCoarse}{50.0\,\ensuremath{\pm}\,1.0\%}  % experiments/FIGS/data/v2/index.json :: linear probe JetClass-II, largest set: relative drop in accuracy (%) with the qg data's missing slots set to mean, R16_Q1 :: sha256 c6c4138ee5c86423
\newcommand{\MlstLpJcTwoQgMeanAccDropCoarsePr}{50\%}  % experiments/FIGS/data/v2/index.json :: linear probe JetClass-II, largest set: relative drop in accuracy (%) with the qg data's missing slots set to mean, R16_Q1 :: sha256 c6c4138ee5c86423 :: prose twin, two significant figures
\newcommand{\MlstLpJcTwoQgDrawAccDropCoarse}{51.7\,\ensuremath{\pm}\,0.4\%}  % experiments/FIGS/data/v2/index.json :: linear probe JetClass-II, largest set: relative drop in accuracy (%) with the qg data's missing slots set to draws, R16_Q1 :: sha256 c6c4138ee5c86423
\newcommand{\MlstLpJcTwoQgDrawAccDropCoarsePr}{52\%}  % experiments/FIGS/data/v2/index.json :: linear probe JetClass-II, largest set: relative drop in accuracy (%) with the qg data's missing slots set to draws, R16_Q1 :: sha256 c6c4138ee5c86423 :: prose twin, two significant figures
\newcommand{\MlstLpJcTwoQgZeroAccDropResMass}{61.8, 67.9\%}  % experiments/FIGS/data/v2/index.json :: linear probe JetClass-II, largest set: relative drop in accuracy (%) with the qg data's missing slots set to zero, L162_MASS :: sha256 c6c4138ee5c86423
\newcommand{\MlstLpJcTwoQgZeroAccDropResMassPr}{62--68\%}  % experiments/FIGS/data/v2/index.json :: linear probe JetClass-II, largest set: relative drop in accuracy (%) with the qg data's missing slots set to zero, L162_MASS :: sha256 c6c4138ee5c86423 :: prose twin, two significant figures
\newcommand{\MlstLpJcTwoQgMeanAccDropResMass}{49.7, 52.2\%}  % experiments/FIGS/data/v2/index.json :: linear probe JetClass-II, largest set: relative drop in accuracy (%) with the qg data's missing slots set to mean, L162_MASS :: sha256 c6c4138ee5c86423
\newcommand{\MlstLpJcTwoQgMeanAccDropResMassPr}{50--52\%}  % experiments/FIGS/data/v2/index.json :: linear probe JetClass-II, largest set: relative drop in accuracy (%) with the qg data's missing slots set to mean, L162_MASS :: sha256 c6c4138ee5c86423 :: prose twin, two significant figures
\newcommand{\MlstLpJcTwoQgDrawAccDropResMass}{54.3, 55.1\%}  % experiments/FIGS/data/v2/index.json :: linear probe JetClass-II, largest set: relative drop in accuracy (%) with the qg data's missing slots set to draws, L162_MASS :: sha256 c6c4138ee5c86423
\newcommand{\MlstLpJcTwoQgDrawAccDropResMassPr}{54--55\%}  % experiments/FIGS/data/v2/index.json :: linear probe JetClass-II, largest set: relative drop in accuracy (%) with the qg data's missing slots set to draws, L162_MASS :: sha256 c6c4138ee5c86423 :: prose twin, two significant figures
\newcommand{\MlstLpJcTwoQgZeroAccDropCoarseMass}{66.7, 65.4\%}  % experiments/FIGS/data/v2/index.json :: linear probe JetClass-II, largest set: relative drop in accuracy (%) with the qg data's missing slots set to zero, R16_Q1_MASS :: sha256 c6c4138ee5c86423
\newcommand{\MlstLpJcTwoQgZeroAccDropCoarseMassPr}{65--67\%}  % experiments/FIGS/data/v2/index.json :: linear probe JetClass-II, largest set: relative drop in accuracy (%) with the qg data's missing slots set to zero, R16_Q1_MASS :: sha256 c6c4138ee5c86423 :: prose twin, two significant figures
\newcommand{\MlstLpJcTwoQgMeanAccDropCoarseMass}{49.0, 49.6\%}  % experiments/FIGS/data/v2/index.json :: linear probe JetClass-II, largest set: relative drop in accuracy (%) with the qg data's missing slots set to mean, R16_Q1_MASS :: sha256 c6c4138ee5c86423
\newcommand{\MlstLpJcTwoQgMeanAccDropCoarseMassPr}{49--50\%}  % experiments/FIGS/data/v2/index.json :: linear probe JetClass-II, largest set: relative drop in accuracy (%) with the qg data's missing slots set to mean, R16_Q1_MASS :: sha256 c6c4138ee5c86423 :: prose twin, two significant figures
\newcommand{\MlstLpJcTwoQgDrawAccDropCoarseMass}{50.8, 51.1\%}  % experiments/FIGS/data/v2/index.json :: linear probe JetClass-II, largest set: relative drop in accuracy (%) with the qg data's missing slots set to draws, R16_Q1_MASS :: sha256 c6c4138ee5c86423
\newcommand{\MlstLpJcTwoQgDrawAccDropCoarseMassPr}{51--51\%}  % experiments/FIGS/data/v2/index.json :: linear probe JetClass-II, largest set: relative drop in accuracy (%) with the qg data's missing slots set to draws, R16_Q1_MASS :: sha256 c6c4138ee5c86423 :: prose twin, two significant figures
\newcommand{\MlstLpJcTwoQgZeroAccDropFzero}{62.7, 62.2\%}  % experiments/FIGS/data/v2/index.json :: linear probe JetClass-II, largest set: relative drop in accuracy (%) with the qg data's missing slots set to zero, FLAV_F0 :: sha256 c6c4138ee5c86423
\newcommand{\MlstLpJcTwoQgZeroAccDropFzeroPr}{62--63\%}  % experiments/FIGS/data/v2/index.json :: linear probe JetClass-II, largest set: relative drop in accuracy (%) with the qg data's missing slots set to zero, FLAV_F0 :: sha256 c6c4138ee5c86423 :: prose twin, two significant figures
\newcommand{\MlstLpJcTwoQgMeanAccDropFzero}{51.9, 51.3\%}  % experiments/FIGS/data/v2/index.json :: linear probe JetClass-II, largest set: relative drop in accuracy (%) with the qg data's missing slots set to mean, FLAV_F0 :: sha256 c6c4138ee5c86423
\newcommand{\MlstLpJcTwoQgMeanAccDropFzeroPr}{51--52\%}  % experiments/FIGS/data/v2/index.json :: linear probe JetClass-II, largest set: relative drop in accuracy (%) with the qg data's missing slots set to mean, FLAV_F0 :: sha256 c6c4138ee5c86423 :: prose twin, two significant figures
\newcommand{\MlstLpJcTwoQgDrawAccDropFzero}{53.9, 53.4\%}  % experiments/FIGS/data/v2/index.json :: linear probe JetClass-II, largest set: relative drop in accuracy (%) with the qg data's missing slots set to draws, FLAV_F0 :: sha256 c6c4138ee5c86423
\newcommand{\MlstLpJcTwoQgDrawAccDropFzeroPr}{53--54\%}  % experiments/FIGS/data/v2/index.json :: linear probe JetClass-II, largest set: relative drop in accuracy (%) with the qg data's missing slots set to draws, FLAV_F0 :: sha256 c6c4138ee5c86423 :: prose twin, two significant figures
\newcommand{\MlstLpJcTwoQgZeroAccDropFone}{66.2, 62.8\%}  % experiments/FIGS/data/v2/index.json :: linear probe JetClass-II, largest set: relative drop in accuracy (%) with the qg data's missing slots set to zero, FLAV_F1 :: sha256 c6c4138ee5c86423
\newcommand{\MlstLpJcTwoQgZeroAccDropFonePr}{63--66\%}  % experiments/FIGS/data/v2/index.json :: linear probe JetClass-II, largest set: relative drop in accuracy (%) with the qg data's missing slots set to zero, FLAV_F1 :: sha256 c6c4138ee5c86423 :: prose twin, two significant figures
\newcommand{\MlstLpJcTwoQgMeanAccDropFone}{52.7, 52.8\%}  % experiments/FIGS/data/v2/index.json :: linear probe JetClass-II, largest set: relative drop in accuracy (%) with the qg data's missing slots set to mean, FLAV_F1 :: sha256 c6c4138ee5c86423
\newcommand{\MlstLpJcTwoQgMeanAccDropFonePr}{53--53\%}  % experiments/FIGS/data/v2/index.json :: linear probe JetClass-II, largest set: relative drop in accuracy (%) with the qg data's missing slots set to mean, FLAV_F1 :: sha256 c6c4138ee5c86423 :: prose twin, two significant figures
\newcommand{\MlstLpJcTwoQgDrawAccDropFone}{55.6, 56.5\%}  % experiments/FIGS/data/v2/index.json :: linear probe JetClass-II, largest set: relative drop in accuracy (%) with the qg data's missing slots set to draws, FLAV_F1 :: sha256 c6c4138ee5c86423
\newcommand{\MlstLpJcTwoQgDrawAccDropFonePr}{56--56\%}  % experiments/FIGS/data/v2/index.json :: linear probe JetClass-II, largest set: relative drop in accuracy (%) with the qg data's missing slots set to draws, FLAV_F1 :: sha256 c6c4138ee5c86423 :: prose twin, two significant figures
\newcommand{\MlstLpJcTwoQgZeroAccDropFoner}{67.2, 69.0\%}  % experiments/FIGS/data/v2/index.json :: linear probe JetClass-II, largest set: relative drop in accuracy (%) with the qg data's missing slots set to zero, FLAV_F1R :: sha256 c6c4138ee5c86423
\newcommand{\MlstLpJcTwoQgZeroAccDropFonerPr}{67--69\%}  % experiments/FIGS/data/v2/index.json :: linear probe JetClass-II, largest set: relative drop in accuracy (%) with the qg data's missing slots set to zero, FLAV_F1R :: sha256 c6c4138ee5c86423 :: prose twin, two significant figures
\newcommand{\MlstLpJcTwoQgMeanAccDropFoner}{52.9, 51.8\%}  % experiments/FIGS/data/v2/index.json :: linear probe JetClass-II, largest set: relative drop in accuracy (%) with the qg data's missing slots set to mean, FLAV_F1R :: sha256 c6c4138ee5c86423
\newcommand{\MlstLpJcTwoQgMeanAccDropFonerPr}{52--53\%}  % experiments/FIGS/data/v2/index.json :: linear probe JetClass-II, largest set: relative drop in accuracy (%) with the qg data's missing slots set to mean, FLAV_F1R :: sha256 c6c4138ee5c86423 :: prose twin, two significant figures
\newcommand{\MlstLpJcTwoQgDrawAccDropFoner}{57.2, 55.4\%}  % experiments/FIGS/data/v2/index.json :: linear probe JetClass-II, largest set: relative drop in accuracy (%) with the qg data's missing slots set to draws, FLAV_F1R :: sha256 c6c4138ee5c86423
\newcommand{\MlstLpJcTwoQgDrawAccDropFonerPr}{55--57\%}  % experiments/FIGS/data/v2/index.json :: linear probe JetClass-II, largest set: relative drop in accuracy (%) with the qg data's missing slots set to draws, FLAV_F1R :: sha256 c6c4138ee5c86423 :: prose twin, two significant figures
\newcommand{\MlstLpJcTwoQgZeroAccDropInitBn}{8.7\,\ensuremath{\pm}\,2.6\%}  % experiments/FIGS/data/v2/index.json :: linear probe JetClass-II, largest set: relative drop in accuracy (%) with the qg data's missing slots set to zero, INIT_BN :: sha256 c6c4138ee5c86423
\newcommand{\MlstLpJcTwoQgZeroAccDropInitBnPr}{8.7\%}  % experiments/FIGS/data/v2/index.json :: linear probe JetClass-II, largest set: relative drop in accuracy (%) with the qg data's missing slots set to zero, INIT_BN :: sha256 c6c4138ee5c86423 :: prose twin, two significant figures
\newcommand{\MlstLpJcTwoQgMeanAccDropInitBn}{10.4\,\ensuremath{\pm}\,1.4\%}  % experiments/FIGS/data/v2/index.json :: linear probe JetClass-II, largest set: relative drop in accuracy (%) with the qg data's missing slots set to mean, INIT_BN :: sha256 c6c4138ee5c86423
\newcommand{\MlstLpJcTwoQgMeanAccDropInitBnPr}{10\%}  % experiments/FIGS/data/v2/index.json :: linear probe JetClass-II, largest set: relative drop in accuracy (%) with the qg data's missing slots set to mean, INIT_BN :: sha256 c6c4138ee5c86423 :: prose twin, two significant figures
\newcommand{\MlstLpJcTwoQgDrawAccDropInitBn}{11.4\,\ensuremath{\pm}\,2.1\%}  % experiments/FIGS/data/v2/index.json :: linear probe JetClass-II, largest set: relative drop in accuracy (%) with the qg data's missing slots set to draws, INIT_BN :: sha256 c6c4138ee5c86423
\newcommand{\MlstLpJcTwoQgDrawAccDropInitBnPr}{11\%}  % experiments/FIGS/data/v2/index.json :: linear probe JetClass-II, largest set: relative drop in accuracy (%) with the qg data's missing slots set to draws, INIT_BN :: sha256 c6c4138ee5c86423 :: prose twin, two significant figures
\newcommand{\MlstMassLambda}{5}  % configs/arms/v2_grid.json :: arms[L162_MASS, R16_Q1_MASS].mass_lambda :: sha256 d8a31b1486992302 :: weight of the mass-regression term in the pretraining loss, both mass models
\newcommand{\MlstMassLossShareResMass}{13.78 and 13.87}  % experiments/FIGS/data/v2/index.json :: share of the mass-regression term in the training loss of L162_MASS, mean over epochs, percent :: sha256 c6c4138ee5c86423
\newcommand{\MlstMassLossShareCoarseMass}{28.65 and 28.67}  % experiments/FIGS/data/v2/index.json :: share of the mass-regression term in the training loss of R16_Q1_MASS, mean over epochs, percent :: sha256 c6c4138ee5c86423
\newcommand{\MlstMassUncorrected}{0.245}  % experiments/FIGS/data/v2/index.json :: head/atlas_res: groomed mass, all jets, ATLAS resolution :: sha256 c6c4138ee5c86423
\newcommand{\MlstMassUncorrectedResponse}{0.929}  % experiments/FIGS/data/v2/index.json :: head/atlas_res: groomed mass, all jets, median response :: sha256 c6c4138ee5c86423
\newcommand{\MlstMassHeadResMass}{0.1471 and 0.1463}  % experiments/FIGS/data/v2/index.json :: regression output of L162_MASS, all jets, ATLAS resolution :: sha256 c6c4138ee5c86423
\newcommand{\MlstMassHeadResMassResponse}{0.995 and 1.003}  % experiments/FIGS/data/v2/index.json :: regression output of L162_MASS, all jets, median response :: sha256 c6c4138ee5c86423
\newcommand{\MlstMassHeadResMassOverUncorr}{0.6009 and 0.5976}  % experiments/FIGS/data/v2/index.json :: resolution of the regression output of L162_MASS over the groomed mass's, all jets :: sha256 c6c4138ee5c86423
\newcommand{\MlstMassHeadResMassOverUncorrPr}{at most 0.60 times the resolution}  % experiments/FIGS/data/v2/index.json :: resolution of the regression output of L162_MASS over the groomed mass's, all jets :: sha256 c6c4138ee5c86423 :: prose twin: two significant figures, direction in words
\newcommand{\MlstMassHeadResMassOverUncorrPrMax}{down to 0.60 times the resolution}  % experiments/FIGS/data/v2/index.json :: resolution of the regression output of L162_MASS over the groomed mass's, all jets :: sha256 c6c4138ee5c86423 :: prose twin of the run farther from 1
\newcommand{\MlstMassHeadResMassOverMlpRes}{0.6645 and 0.6691}  % experiments/FIGS/data/v2/index.json :: resolution of the regression output of L162_MASS over the frozen-feature MLP of L162, run by run, all jets :: sha256 c6c4138ee5c86423
\newcommand{\MlstMassHeadResMassOverMlpResPr}{at most 0.67 times the resolution}  % experiments/FIGS/data/v2/index.json :: resolution of the regression output of L162_MASS over the frozen-feature MLP of L162, run by run, all jets :: sha256 c6c4138ee5c86423 :: prose twin: two significant figures, direction in words
\newcommand{\MlstMassHeadResMassOverMlpResPrMax}{down to 0.66 times the resolution}  % experiments/FIGS/data/v2/index.json :: resolution of the regression output of L162_MASS over the frozen-feature MLP of L162, run by run, all jets :: sha256 c6c4138ee5c86423 :: prose twin of the run farther from 1
\newcommand{\MlstMassHeadCoarseMass}{0.1450 and 0.1461}  % experiments/FIGS/data/v2/index.json :: regression output of R16_Q1_MASS, all jets, ATLAS resolution :: sha256 c6c4138ee5c86423
\newcommand{\MlstMassHeadCoarseMassResponse}{1.005 and 1.005}  % experiments/FIGS/data/v2/index.json :: regression output of R16_Q1_MASS, all jets, median response :: sha256 c6c4138ee5c86423
\newcommand{\MlstMassHeadCoarseMassOverUncorr}{0.5921 and 0.5966}  % experiments/FIGS/data/v2/index.json :: resolution of the regression output of R16_Q1_MASS over the groomed mass's, all jets :: sha256 c6c4138ee5c86423
\newcommand{\MlstMassHeadCoarseMassOverUncorrPr}{at most 0.60 times the resolution}  % experiments/FIGS/data/v2/index.json :: resolution of the regression output of R16_Q1_MASS over the groomed mass's, all jets :: sha256 c6c4138ee5c86423 :: prose twin: two significant figures, direction in words
\newcommand{\MlstMassHeadCoarseMassOverUncorrPrMax}{down to 0.59 times the resolution}  % experiments/FIGS/data/v2/index.json :: resolution of the regression output of R16_Q1_MASS over the groomed mass's, all jets :: sha256 c6c4138ee5c86423 :: prose twin of the run farther from 1
\newcommand{\MlstMassHeadCoarseMassOverMlpCoarse}{0.703 and 0.690}  % experiments/FIGS/data/v2/index.json :: resolution of the regression output of R16_Q1_MASS over the frozen-feature MLP of R16_Q1, run by run, all jets :: sha256 c6c4138ee5c86423
\newcommand{\MlstMassHeadCoarseMassOverMlpCoarsePr}{at most 0.70 times the resolution}  % experiments/FIGS/data/v2/index.json :: resolution of the regression output of R16_Q1_MASS over the frozen-feature MLP of R16_Q1, run by run, all jets :: sha256 c6c4138ee5c86423 :: prose twin: two significant figures, direction in words
\newcommand{\MlstMassHeadCoarseMassOverMlpCoarsePrMax}{down to 0.69 times the resolution}  % experiments/FIGS/data/v2/index.json :: resolution of the regression output of R16_Q1_MASS over the frozen-feature MLP of R16_Q1, run by run, all jets :: sha256 c6c4138ee5c86423 :: prose twin of the run farther from 1
\newcommand{\MlstMassMlpInit}{0.2421\,\ensuremath{\pm}\,0.0018}  % experiments/FIGS/data/v2/index.json :: frozen-feature MLP of INIT, all jets, ATLAS resolution :: sha256 c6c4138ee5c86423
\newcommand{\MlstMassMlpFull}{0.2200\,\ensuremath{\pm}\,0.0014}  % experiments/FIGS/data/v2/index.json :: frozen-feature MLP of L188, all jets, ATLAS resolution :: sha256 c6c4138ee5c86423
\newcommand{\MlstMassMlpRes}{0.2195\,\ensuremath{\pm}\,0.0017}  % experiments/FIGS/data/v2/index.json :: frozen-feature MLP of L162, all jets, ATLAS resolution :: sha256 c6c4138ee5c86423
\newcommand{\MlstMassMlpMid}{0.2158\,\ensuremath{\pm}\,0.0029}  % experiments/FIGS/data/v2/index.json :: frozen-feature MLP of R42_Q1, all jets, ATLAS resolution :: sha256 c6c4138ee5c86423
\newcommand{\MlstMassMlpCoarse}{0.2098\,\ensuremath{\pm}\,0.0031}  % experiments/FIGS/data/v2/index.json :: frozen-feature MLP of R16_Q1, all jets, ATLAS resolution :: sha256 c6c4138ee5c86423
\newcommand{\MlstMassMlpResMass}{0.1591 and 0.1586}  % experiments/FIGS/data/v2/index.json :: frozen-feature MLP of L162_MASS, all jets, ATLAS resolution :: sha256 c6c4138ee5c86423
\newcommand{\MlstMassMlpCoarseMass}{0.1583 and 0.1588}  % experiments/FIGS/data/v2/index.json :: frozen-feature MLP of R16_Q1_MASS, all jets, ATLAS resolution :: sha256 c6c4138ee5c86423
\newcommand{\MlstCmpMlpResOverFull}{0.998\,\ensuremath{\pm}\,0.006}  % experiments/FIGS/data/v2/index.json :: frozen-feature MLP resolution of L162 over L188, run by run, all jets :: sha256 c6c4138ee5c86423
\newcommand{\MlstCmpMlpResOverFullPr}{equal within 1\%}  % experiments/FIGS/data/v2/index.json :: frozen-feature MLP resolution of L162 over L188, run by run, all jets :: sha256 c6c4138ee5c86423 :: prose twin: two significant figures, direction in words
\newcommand{\MlstCmpMlpMidOverFull}{0.981\,\ensuremath{\pm}\,0.017}  % experiments/FIGS/data/v2/index.json :: frozen-feature MLP resolution of R42_Q1 over L188, run by run, all jets :: sha256 c6c4138ee5c86423
\newcommand{\MlstCmpMlpMidOverFullPr}{equal within 2\%}  % experiments/FIGS/data/v2/index.json :: frozen-feature MLP resolution of R42_Q1 over L188, run by run, all jets :: sha256 c6c4138ee5c86423 :: prose twin: two significant figures, direction in words
\newcommand{\MlstCmpMlpCoarseOverFull}{0.954\,\ensuremath{\pm}\,0.018}  % experiments/FIGS/data/v2/index.json :: frozen-feature MLP resolution of R16_Q1 over L188, run by run, all jets :: sha256 c6c4138ee5c86423
\newcommand{\MlstCmpMlpCoarseOverFullPr}{0.95 times the resolution}  % experiments/FIGS/data/v2/index.json :: frozen-feature MLP resolution of R16_Q1 over L188, run by run, all jets :: sha256 c6c4138ee5c86423 :: prose twin: two significant figures, direction in words
\newcommand{\MlstCmpMlpResMassOverRes}{0.7186 and 0.7253}  % experiments/FIGS/data/v2/index.json :: frozen-feature MLP resolution of L162_MASS over L162, run by run, all jets :: sha256 c6c4138ee5c86423
\newcommand{\MlstCmpMlpResMassOverResPr}{at most 0.73 times the resolution}  % experiments/FIGS/data/v2/index.json :: frozen-feature MLP resolution of L162_MASS over L162, run by run, all jets :: sha256 c6c4138ee5c86423 :: prose twin: two significant figures, direction in words
\newcommand{\MlstCmpMlpResMassOverResPrMax}{down to 0.72 times the resolution}  % experiments/FIGS/data/v2/index.json :: frozen-feature MLP resolution of L162_MASS over L162, run by run, all jets :: sha256 c6c4138ee5c86423 :: prose twin of the run farther from 1
\newcommand{\MlstCmpMlpCoarseMassOverCoarse}{0.768 and 0.750}  % experiments/FIGS/data/v2/index.json :: frozen-feature MLP resolution of R16_Q1_MASS over R16_Q1, run by run, all jets :: sha256 c6c4138ee5c86423
\newcommand{\MlstCmpMlpCoarseMassOverCoarsePr}{at most 0.77 times the resolution}  % experiments/FIGS/data/v2/index.json :: frozen-feature MLP resolution of R16_Q1_MASS over R16_Q1, run by run, all jets :: sha256 c6c4138ee5c86423 :: prose twin: two significant figures, direction in words
\newcommand{\MlstCmpMlpCoarseMassOverCoarsePrMax}{down to 0.75 times the resolution}  % experiments/FIGS/data/v2/index.json :: frozen-feature MLP resolution of R16_Q1_MASS over R16_Q1, run by run, all jets :: sha256 c6c4138ee5c86423 :: prose twin of the run farther from 1
\newcommand{\MlstMassUncorrectedResonant}{0.213}  % experiments/FIGS/data/v2/index.json :: head/atlas_res: groomed mass, resonant jets, ATLAS resolution :: sha256 c6c4138ee5c86423
\newcommand{\MlstMassUncorrectedResponseResonant}{0.929}  % experiments/FIGS/data/v2/index.json :: head/atlas_res: groomed mass, resonant jets, median response :: sha256 c6c4138ee5c86423
\newcommand{\MlstMassHeadResMassResonant}{0.1332 and 0.1320}  % experiments/FIGS/data/v2/index.json :: regression output of L162_MASS, resonant jets, ATLAS resolution :: sha256 c6c4138ee5c86423
\newcommand{\MlstMassHeadResMassResponseResonant}{0.996 and 1.004}  % experiments/FIGS/data/v2/index.json :: regression output of L162_MASS, resonant jets, median response :: sha256 c6c4138ee5c86423
\newcommand{\MlstMassHeadResMassOverUncorrResonant}{0.6259 and 0.6205}  % experiments/FIGS/data/v2/index.json :: resolution of the regression output of L162_MASS over the groomed mass's, resonant jets :: sha256 c6c4138ee5c86423
\newcommand{\MlstMassHeadResMassOverUncorrResonantPr}{at most 0.63 times the resolution}  % experiments/FIGS/data/v2/index.json :: resolution of the regression output of L162_MASS over the groomed mass's, resonant jets :: sha256 c6c4138ee5c86423 :: prose twin: two significant figures, direction in words
\newcommand{\MlstMassHeadResMassOverUncorrResonantPrMax}{down to 0.62 times the resolution}  % experiments/FIGS/data/v2/index.json :: resolution of the regression output of L162_MASS over the groomed mass's, resonant jets :: sha256 c6c4138ee5c86423 :: prose twin of the run farther from 1
\newcommand{\MlstMassHeadResMassOverMlpResResonant}{0.6620 and 0.6662}  % experiments/FIGS/data/v2/index.json :: resolution of the regression output of L162_MASS over the frozen-feature MLP of L162, run by run, resonant jets :: sha256 c6c4138ee5c86423
\newcommand{\MlstMassHeadResMassOverMlpResResonantPr}{at most 0.67 times the resolution}  % experiments/FIGS/data/v2/index.json :: resolution of the regression output of L162_MASS over the frozen-feature MLP of L162, run by run, resonant jets :: sha256 c6c4138ee5c86423 :: prose twin: two significant figures, direction in words
\newcommand{\MlstMassHeadResMassOverMlpResResonantPrMax}{down to 0.66 times the resolution}  % experiments/FIGS/data/v2/index.json :: resolution of the regression output of L162_MASS over the frozen-feature MLP of L162, run by run, resonant jets :: sha256 c6c4138ee5c86423 :: prose twin of the run farther from 1
\newcommand{\MlstMassHeadCoarseMassResonant}{0.1312 and 0.1320}  % experiments/FIGS/data/v2/index.json :: regression output of R16_Q1_MASS, resonant jets, ATLAS resolution :: sha256 c6c4138ee5c86423
\newcommand{\MlstMassHeadCoarseMassResponseResonant}{1.006 and 1.006}  % experiments/FIGS/data/v2/index.json :: regression output of R16_Q1_MASS, resonant jets, median response :: sha256 c6c4138ee5c86423
\newcommand{\MlstMassHeadCoarseMassOverUncorrResonant}{0.6165 and 0.6206}  % experiments/FIGS/data/v2/index.json :: resolution of the regression output of R16_Q1_MASS over the groomed mass's, resonant jets :: sha256 c6c4138ee5c86423
\newcommand{\MlstMassHeadCoarseMassOverUncorrResonantPr}{at most 0.62 times the resolution}  % experiments/FIGS/data/v2/index.json :: resolution of the regression output of R16_Q1_MASS over the groomed mass's, resonant jets :: sha256 c6c4138ee5c86423 :: prose twin: two significant figures, direction in words
\newcommand{\MlstMassHeadCoarseMassOverUncorrResonantPrMax}{down to 0.62 times the resolution}  % experiments/FIGS/data/v2/index.json :: resolution of the regression output of R16_Q1_MASS over the groomed mass's, resonant jets :: sha256 c6c4138ee5c86423 :: prose twin of the run farther from 1
\newcommand{\MlstMassHeadCoarseMassOverMlpCoarseResonant}{0.707 and 0.688}  % experiments/FIGS/data/v2/index.json :: resolution of the regression output of R16_Q1_MASS over the frozen-feature MLP of R16_Q1, run by run, resonant jets :: sha256 c6c4138ee5c86423
\newcommand{\MlstMassHeadCoarseMassOverMlpCoarseResonantPr}{at most 0.71 times the resolution}  % experiments/FIGS/data/v2/index.json :: resolution of the regression output of R16_Q1_MASS over the frozen-feature MLP of R16_Q1, run by run, resonant jets :: sha256 c6c4138ee5c86423 :: prose twin: two significant figures, direction in words
\newcommand{\MlstMassHeadCoarseMassOverMlpCoarseResonantPrMax}{down to 0.69 times the resolution}  % experiments/FIGS/data/v2/index.json :: resolution of the regression output of R16_Q1_MASS over the frozen-feature MLP of R16_Q1, run by run, resonant jets :: sha256 c6c4138ee5c86423 :: prose twin of the run farther from 1
\newcommand{\MlstMassMlpInitResonant}{0.2214\,\ensuremath{\pm}\,0.0017}  % experiments/FIGS/data/v2/index.json :: frozen-feature MLP of INIT, resonant jets, ATLAS resolution :: sha256 c6c4138ee5c86423
\newcommand{\MlstMassMlpFullResonant}{0.1990\,\ensuremath{\pm}\,0.0016}  % experiments/FIGS/data/v2/index.json :: frozen-feature MLP of L188, resonant jets, ATLAS resolution :: sha256 c6c4138ee5c86423
\newcommand{\MlstMassMlpResResonant}{0.1989\,\ensuremath{\pm}\,0.0020}  % experiments/FIGS/data/v2/index.json :: frozen-feature MLP of L162, resonant jets, ATLAS resolution :: sha256 c6c4138ee5c86423
\newcommand{\MlstMassMlpMidResonant}{0.1949\,\ensuremath{\pm}\,0.0023}  % experiments/FIGS/data/v2/index.json :: frozen-feature MLP of R42_Q1, resonant jets, ATLAS resolution :: sha256 c6c4138ee5c86423
\newcommand{\MlstMassMlpCoarseResonant}{0.1894\,\ensuremath{\pm}\,0.0035}  % experiments/FIGS/data/v2/index.json :: frozen-feature MLP of R16_Q1, resonant jets, ATLAS resolution :: sha256 c6c4138ee5c86423
\newcommand{\MlstMassMlpResMassResonant}{0.1436 and 0.1430}  % experiments/FIGS/data/v2/index.json :: frozen-feature MLP of L162_MASS, resonant jets, ATLAS resolution :: sha256 c6c4138ee5c86423
\newcommand{\MlstMassMlpCoarseMassResonant}{0.1420 and 0.1431}  % experiments/FIGS/data/v2/index.json :: frozen-feature MLP of R16_Q1_MASS, resonant jets, ATLAS resolution :: sha256 c6c4138ee5c86423
\newcommand{\MlstCmpMlpResOverFullResonant}{1.000\,\ensuremath{\pm}\,0.005}  % experiments/FIGS/data/v2/index.json :: frozen-feature MLP resolution of L162 over L188, run by run, resonant jets :: sha256 c6c4138ee5c86423
\newcommand{\MlstCmpMlpResOverFullResonantPr}{equal within 1\%}  % experiments/FIGS/data/v2/index.json :: frozen-feature MLP resolution of L162 over L188, run by run, resonant jets :: sha256 c6c4138ee5c86423 :: prose twin: two significant figures, direction in words
\newcommand{\MlstCmpMlpMidOverFullResonant}{0.980\,\ensuremath{\pm}\,0.018}  % experiments/FIGS/data/v2/index.json :: frozen-feature MLP resolution of R42_Q1 over L188, run by run, resonant jets :: sha256 c6c4138ee5c86423
\newcommand{\MlstCmpMlpMidOverFullResonantPr}{equal within 3\%}  % experiments/FIGS/data/v2/index.json :: frozen-feature MLP resolution of R42_Q1 over L188, run by run, resonant jets :: sha256 c6c4138ee5c86423 :: prose twin: two significant figures, direction in words
\newcommand{\MlstCmpMlpCoarseOverFullResonant}{0.952\,\ensuremath{\pm}\,0.023}  % experiments/FIGS/data/v2/index.json :: frozen-feature MLP resolution of R16_Q1 over L188, run by run, resonant jets :: sha256 c6c4138ee5c86423
\newcommand{\MlstCmpMlpCoarseOverFullResonantPr}{0.95 times the resolution}  % experiments/FIGS/data/v2/index.json :: frozen-feature MLP resolution of R16_Q1 over L188, run by run, resonant jets :: sha256 c6c4138ee5c86423 :: prose twin: two significant figures, direction in words
\newcommand{\MlstCmpMlpResMassOverResResonant}{0.7137 and 0.7219}  % experiments/FIGS/data/v2/index.json :: frozen-feature MLP resolution of L162_MASS over L162, run by run, resonant jets :: sha256 c6c4138ee5c86423
\newcommand{\MlstCmpMlpResMassOverResResonantPr}{at most 0.72 times the resolution}  % experiments/FIGS/data/v2/index.json :: frozen-feature MLP resolution of L162_MASS over L162, run by run, resonant jets :: sha256 c6c4138ee5c86423 :: prose twin: two significant figures, direction in words
\newcommand{\MlstCmpMlpResMassOverResResonantPrMax}{down to 0.71 times the resolution}  % experiments/FIGS/data/v2/index.json :: frozen-feature MLP resolution of L162_MASS over L162, run by run, resonant jets :: sha256 c6c4138ee5c86423 :: prose twin of the run farther from 1
\newcommand{\MlstCmpMlpCoarseMassOverCoarseResonant}{0.766 and 0.745}  % experiments/FIGS/data/v2/index.json :: frozen-feature MLP resolution of R16_Q1_MASS over R16_Q1, run by run, resonant jets :: sha256 c6c4138ee5c86423
\newcommand{\MlstCmpMlpCoarseMassOverCoarseResonantPr}{at most 0.77 times the resolution}  % experiments/FIGS/data/v2/index.json :: frozen-feature MLP resolution of R16_Q1_MASS over R16_Q1, run by run, resonant jets :: sha256 c6c4138ee5c86423 :: prose twin: two significant figures, direction in words
\newcommand{\MlstCmpMlpCoarseMassOverCoarseResonantPrMax}{down to 0.75 times the resolution}  % experiments/FIGS/data/v2/index.json :: frozen-feature MLP resolution of R16_Q1_MASS over R16_Q1, run by run, resonant jets :: sha256 c6c4138ee5c86423 :: prose twin of the run farther from 1
\newcommand{\MlstMassUncorrectedQcd}{1.52}  % experiments/FIGS/data/v2/index.json :: head/atlas_res: groomed mass, qcd jets, ATLAS resolution :: sha256 c6c4138ee5c86423
\newcommand{\MlstMassUncorrectedResponseQcd}{0.931}  % experiments/FIGS/data/v2/index.json :: head/atlas_res: groomed mass, qcd jets, median response :: sha256 c6c4138ee5c86423
\newcommand{\MlstMassHeadResMassQcd}{0.2263 and 0.2256}  % experiments/FIGS/data/v2/index.json :: regression output of L162_MASS, qcd jets, ATLAS resolution :: sha256 c6c4138ee5c86423
\newcommand{\MlstMassHeadResMassResponseQcd}{0.9897 and 0.9996}  % experiments/FIGS/data/v2/index.json :: regression output of L162_MASS, qcd jets, median response :: sha256 c6c4138ee5c86423
\newcommand{\MlstMassHeadResMassOverUncorrQcd}{0.1485 and 0.1480}  % experiments/FIGS/data/v2/index.json :: resolution of the regression output of L162_MASS over the groomed mass's, qcd jets :: sha256 c6c4138ee5c86423
\newcommand{\MlstMassHeadResMassOverUncorrQcdPr}{at most 0.15 times the resolution}  % experiments/FIGS/data/v2/index.json :: resolution of the regression output of L162_MASS over the groomed mass's, qcd jets :: sha256 c6c4138ee5c86423 :: prose twin: two significant figures, direction in words
\newcommand{\MlstMassHeadResMassOverUncorrQcdPrMax}{down to 0.15 times the resolution}  % experiments/FIGS/data/v2/index.json :: resolution of the regression output of L162_MASS over the groomed mass's, qcd jets :: sha256 c6c4138ee5c86423 :: prose twin of the run farther from 1
\newcommand{\MlstMassHeadResMassOverMlpResQcd}{0.6542 and 0.6632}  % experiments/FIGS/data/v2/index.json :: resolution of the regression output of L162_MASS over the frozen-feature MLP of L162, run by run, qcd jets :: sha256 c6c4138ee5c86423
\newcommand{\MlstMassHeadResMassOverMlpResQcdPr}{at most 0.66 times the resolution}  % experiments/FIGS/data/v2/index.json :: resolution of the regression output of L162_MASS over the frozen-feature MLP of L162, run by run, qcd jets :: sha256 c6c4138ee5c86423 :: prose twin: two significant figures, direction in words
\newcommand{\MlstMassHeadResMassOverMlpResQcdPrMax}{down to 0.65 times the resolution}  % experiments/FIGS/data/v2/index.json :: resolution of the regression output of L162_MASS over the frozen-feature MLP of L162, run by run, qcd jets :: sha256 c6c4138ee5c86423 :: prose twin of the run farther from 1
\newcommand{\MlstMassHeadCoarseMassQcd}{0.2239 and 0.2240}  % experiments/FIGS/data/v2/index.json :: regression output of R16_Q1_MASS, qcd jets, ATLAS resolution :: sha256 c6c4138ee5c86423
\newcommand{\MlstMassHeadCoarseMassResponseQcd}{1.001 and 1.000}  % experiments/FIGS/data/v2/index.json :: regression output of R16_Q1_MASS, qcd jets, median response :: sha256 c6c4138ee5c86423
\newcommand{\MlstMassHeadCoarseMassOverUncorrQcd}{0.1469 and 0.1469}  % experiments/FIGS/data/v2/index.json :: resolution of the regression output of R16_Q1_MASS over the groomed mass's, qcd jets :: sha256 c6c4138ee5c86423
\newcommand{\MlstMassHeadCoarseMassOverUncorrQcdPr}{at most 0.15 times the resolution}  % experiments/FIGS/data/v2/index.json :: resolution of the regression output of R16_Q1_MASS over the groomed mass's, qcd jets :: sha256 c6c4138ee5c86423 :: prose twin: two significant figures, direction in words
\newcommand{\MlstMassHeadCoarseMassOverUncorrQcdPrMax}{down to 0.15 times the resolution}  % experiments/FIGS/data/v2/index.json :: resolution of the regression output of R16_Q1_MASS over the groomed mass's, qcd jets :: sha256 c6c4138ee5c86423 :: prose twin of the run farther from 1
\newcommand{\MlstMassHeadCoarseMassOverMlpCoarseQcd}{0.679 and 0.698}  % experiments/FIGS/data/v2/index.json :: resolution of the regression output of R16_Q1_MASS over the frozen-feature MLP of R16_Q1, run by run, qcd jets :: sha256 c6c4138ee5c86423
\newcommand{\MlstMassHeadCoarseMassOverMlpCoarseQcdPr}{at most 0.70 times the resolution}  % experiments/FIGS/data/v2/index.json :: resolution of the regression output of R16_Q1_MASS over the frozen-feature MLP of R16_Q1, run by run, qcd jets :: sha256 c6c4138ee5c86423 :: prose twin: two significant figures, direction in words
\newcommand{\MlstMassHeadCoarseMassOverMlpCoarseQcdPrMax}{down to 0.68 times the resolution}  % experiments/FIGS/data/v2/index.json :: resolution of the regression output of R16_Q1_MASS over the frozen-feature MLP of R16_Q1, run by run, qcd jets :: sha256 c6c4138ee5c86423 :: prose twin of the run farther from 1
\newcommand{\MlstMassMlpInitQcd}{0.502\,\ensuremath{\pm}\,0.014}  % experiments/FIGS/data/v2/index.json :: frozen-feature MLP of INIT, qcd jets, ATLAS resolution :: sha256 c6c4138ee5c86423
\newcommand{\MlstMassMlpFullQcd}{0.3487\,\ensuremath{\pm}\,0.0029}  % experiments/FIGS/data/v2/index.json :: frozen-feature MLP of L188, qcd jets, ATLAS resolution :: sha256 c6c4138ee5c86423
\newcommand{\MlstMassMlpResQcd}{0.3425\,\ensuremath{\pm}\,0.0031}  % experiments/FIGS/data/v2/index.json :: frozen-feature MLP of L162, qcd jets, ATLAS resolution :: sha256 c6c4138ee5c86423
\newcommand{\MlstMassMlpMidQcd}{0.338\,\ensuremath{\pm}\,0.007}  % experiments/FIGS/data/v2/index.json :: frozen-feature MLP of R42_Q1, qcd jets, ATLAS resolution :: sha256 c6c4138ee5c86423
\newcommand{\MlstMassMlpCoarseQcd}{0.327\,\ensuremath{\pm}\,0.006}  % experiments/FIGS/data/v2/index.json :: frozen-feature MLP of R16_Q1, qcd jets, ATLAS resolution :: sha256 c6c4138ee5c86423
\newcommand{\MlstMassMlpResMassQcd}{0.2473 and 0.2452}  % experiments/FIGS/data/v2/index.json :: frozen-feature MLP of L162_MASS, qcd jets, ATLAS resolution :: sha256 c6c4138ee5c86423
\newcommand{\MlstMassMlpCoarseMassQcd}{0.2465 and 0.2463}  % experiments/FIGS/data/v2/index.json :: frozen-feature MLP of R16_Q1_MASS, qcd jets, ATLAS resolution :: sha256 c6c4138ee5c86423
\newcommand{\MlstCmpMlpResOverFullQcd}{0.982\,\ensuremath{\pm}\,0.017}  % experiments/FIGS/data/v2/index.json :: frozen-feature MLP resolution of L162 over L188, run by run, qcd jets :: sha256 c6c4138ee5c86423
\newcommand{\MlstCmpMlpResOverFullQcdPr}{equal within 2\%}  % experiments/FIGS/data/v2/index.json :: frozen-feature MLP resolution of L162 over L188, run by run, qcd jets :: sha256 c6c4138ee5c86423 :: prose twin: two significant figures, direction in words
\newcommand{\MlstCmpMlpMidOverFullQcd}{0.969\,\ensuremath{\pm}\,0.013}  % experiments/FIGS/data/v2/index.json :: frozen-feature MLP resolution of R42_Q1 over L188, run by run, qcd jets :: sha256 c6c4138ee5c86423
\newcommand{\MlstCmpMlpMidOverFullQcdPr}{0.97 times the resolution}  % experiments/FIGS/data/v2/index.json :: frozen-feature MLP resolution of R42_Q1 over L188, run by run, qcd jets :: sha256 c6c4138ee5c86423 :: prose twin: two significant figures, direction in words
\newcommand{\MlstCmpMlpCoarseOverFullQcd}{0.939\,\ensuremath{\pm}\,0.023}  % experiments/FIGS/data/v2/index.json :: frozen-feature MLP resolution of R16_Q1 over L188, run by run, qcd jets :: sha256 c6c4138ee5c86423
\newcommand{\MlstCmpMlpCoarseOverFullQcdPr}{0.94 times the resolution}  % experiments/FIGS/data/v2/index.json :: frozen-feature MLP resolution of R16_Q1 over L188, run by run, qcd jets :: sha256 c6c4138ee5c86423 :: prose twin: two significant figures, direction in words
\newcommand{\MlstCmpMlpResMassOverResQcd}{0.7147 and 0.7209}  % experiments/FIGS/data/v2/index.json :: frozen-feature MLP resolution of L162_MASS over L162, run by run, qcd jets :: sha256 c6c4138ee5c86423
\newcommand{\MlstCmpMlpResMassOverResQcdPr}{at most 0.72 times the resolution}  % experiments/FIGS/data/v2/index.json :: frozen-feature MLP resolution of L162_MASS over L162, run by run, qcd jets :: sha256 c6c4138ee5c86423 :: prose twin: two significant figures, direction in words
\newcommand{\MlstCmpMlpResMassOverResQcdPrMax}{down to 0.71 times the resolution}  % experiments/FIGS/data/v2/index.json :: frozen-feature MLP resolution of L162_MASS over L162, run by run, qcd jets :: sha256 c6c4138ee5c86423 :: prose twin of the run farther from 1
\newcommand{\MlstCmpMlpCoarseMassOverCoarseQcd}{0.747 and 0.768}  % experiments/FIGS/data/v2/index.json :: frozen-feature MLP resolution of R16_Q1_MASS over R16_Q1, run by run, qcd jets :: sha256 c6c4138ee5c86423
\newcommand{\MlstCmpMlpCoarseMassOverCoarseQcdPr}{at most 0.77 times the resolution}  % experiments/FIGS/data/v2/index.json :: frozen-feature MLP resolution of R16_Q1_MASS over R16_Q1, run by run, qcd jets :: sha256 c6c4138ee5c86423 :: prose twin: two significant figures, direction in words
\newcommand{\MlstCmpMlpCoarseMassOverCoarseQcdPrMax}{down to 0.75 times the resolution}  % experiments/FIGS/data/v2/index.json :: frozen-feature MLP resolution of R16_Q1_MASS over R16_Q1, run by run, qcd jets :: sha256 c6c4138ee5c86423 :: prose twin of the run farther from 1
\newcommand{\MlstMassRespAxesShareQcd}{76\%}  % experiments/FIGS/data/v2/index.json :: head/atlas_res: share of QCD jets whose groomed-mass response lies inside F9a(b)'s x axis :: sha256 c6c4138ee5c86423
\newcommand{\MlstSculptNPassTen}{41{,}028}  % experiments/FIGS/data/v2/index.json :: QCD test jets passing the cut at QCD efficiency 0.1, every run :: sha256 c6c4138ee5c86423
\newcommand{\MlstSculptFloorTen}{0.00944}  % experiments/FIGS/data/v2/index.json :: Jensen-Shannon distance expected for a random selection of as many QCD jets as pass at 0.1 :: sha256 c6c4138ee5c86423
\newcommand{\MlstSculptFullTen}{0.0524\,\ensuremath{\pm}\,0.0034}  % experiments/FIGS/data/v2/index.json :: Jensen-Shannon distance of the QCD m_SD distribution at QCD efficiency 0.1, L188 :: sha256 c6c4138ee5c86423
\newcommand{\MlstSculptFullOverFloorTen}{5.6\,\ensuremath{\pm}\,0.4}  % experiments/FIGS/data/v2/index.json :: distance at 0.1 over the random-selection floor, L188 :: sha256 c6c4138ee5c86423
\newcommand{\MlstSculptFullOverFloorTenPr}{5.6 times the distance}  % experiments/FIGS/data/v2/index.json :: distance at 0.1 over the random-selection floor, L188 :: sha256 c6c4138ee5c86423 :: prose twin: two significant figures, direction in words
\newcommand{\MlstSculptResTen}{0.0532\,\ensuremath{\pm}\,0.0009}  % experiments/FIGS/data/v2/index.json :: Jensen-Shannon distance of the QCD m_SD distribution at QCD efficiency 0.1, L162 :: sha256 c6c4138ee5c86423
\newcommand{\MlstSculptResOverFloorTen}{5.64\,\ensuremath{\pm}\,0.10}  % experiments/FIGS/data/v2/index.json :: distance at 0.1 over the random-selection floor, L162 :: sha256 c6c4138ee5c86423
\newcommand{\MlstSculptResOverFloorTenPr}{5.6 times the distance}  % experiments/FIGS/data/v2/index.json :: distance at 0.1 over the random-selection floor, L162 :: sha256 c6c4138ee5c86423 :: prose twin: two significant figures, direction in words
\newcommand{\MlstSculptMidTen}{0.0560\,\ensuremath{\pm}\,0.0027}  % experiments/FIGS/data/v2/index.json :: Jensen-Shannon distance of the QCD m_SD distribution at QCD efficiency 0.1, R42_Q1 :: sha256 c6c4138ee5c86423
\newcommand{\MlstSculptMidOverFloorTen}{5.94\,\ensuremath{\pm}\,0.28}  % experiments/FIGS/data/v2/index.json :: distance at 0.1 over the random-selection floor, R42_Q1 :: sha256 c6c4138ee5c86423
\newcommand{\MlstSculptMidOverFloorTenPr}{5.9 times the distance}  % experiments/FIGS/data/v2/index.json :: distance at 0.1 over the random-selection floor, R42_Q1 :: sha256 c6c4138ee5c86423 :: prose twin: two significant figures, direction in words
\newcommand{\MlstSculptCoarseTen}{0.0543\,\ensuremath{\pm}\,0.0010}  % experiments/FIGS/data/v2/index.json :: Jensen-Shannon distance of the QCD m_SD distribution at QCD efficiency 0.1, R16_Q1 :: sha256 c6c4138ee5c86423
\newcommand{\MlstSculptCoarseOverFloorTen}{5.76\,\ensuremath{\pm}\,0.10}  % experiments/FIGS/data/v2/index.json :: distance at 0.1 over the random-selection floor, R16_Q1 :: sha256 c6c4138ee5c86423
\newcommand{\MlstSculptCoarseOverFloorTenPr}{5.8 times the distance}  % experiments/FIGS/data/v2/index.json :: distance at 0.1 over the random-selection floor, R16_Q1 :: sha256 c6c4138ee5c86423 :: prose twin: two significant figures, direction in words
\newcommand{\MlstSculptResMassTen}{0.0422 and 0.0464}  % experiments/FIGS/data/v2/index.json :: Jensen-Shannon distance of the QCD m_SD distribution at QCD efficiency 0.1, L162_MASS :: sha256 c6c4138ee5c86423
\newcommand{\MlstSculptResMassOverFloorTen}{4.47 and 4.91}  % experiments/FIGS/data/v2/index.json :: distance at 0.1 over the random-selection floor, L162_MASS :: sha256 c6c4138ee5c86423
\newcommand{\MlstSculptResMassOverFloorTenPr}{at least 4.5 times the distance}  % experiments/FIGS/data/v2/index.json :: distance at 0.1 over the random-selection floor, L162_MASS :: sha256 c6c4138ee5c86423 :: prose twin: two significant figures, direction in words
\newcommand{\MlstSculptResMassOverFloorTenPrMax}{up to 4.9 times the distance}  % experiments/FIGS/data/v2/index.json :: distance at 0.1 over the random-selection floor, L162_MASS :: sha256 c6c4138ee5c86423 :: prose twin of the run farther from 1
\newcommand{\MlstSculptCoarseMassTen}{0.0522 and 0.0423}  % experiments/FIGS/data/v2/index.json :: Jensen-Shannon distance of the QCD m_SD distribution at QCD efficiency 0.1, R16_Q1_MASS :: sha256 c6c4138ee5c86423
\newcommand{\MlstSculptCoarseMassOverFloorTen}{5.5 and 4.5}  % experiments/FIGS/data/v2/index.json :: distance at 0.1 over the random-selection floor, R16_Q1_MASS :: sha256 c6c4138ee5c86423
\newcommand{\MlstSculptCoarseMassOverFloorTenPr}{at least 4.5 times the distance}  % experiments/FIGS/data/v2/index.json :: distance at 0.1 over the random-selection floor, R16_Q1_MASS :: sha256 c6c4138ee5c86423 :: prose twin: two significant figures, direction in words
\newcommand{\MlstSculptCoarseMassOverFloorTenPrMax}{up to 5.5 times the distance}  % experiments/FIGS/data/v2/index.json :: distance at 0.1 over the random-selection floor, R16_Q1_MASS :: sha256 c6c4138ee5c86423 :: prose twin of the run farther from 1
\newcommand{\MlstSculptResMassOverResTen}{0.789 and 0.860}  % experiments/FIGS/data/v2/index.json :: distance at QCD efficiency 0.1: L162_MASS over L162, run by run :: sha256 c6c4138ee5c86423
\newcommand{\MlstSculptResMassOverResTenPr}{at most 0.86 times the distance}  % experiments/FIGS/data/v2/index.json :: distance at QCD efficiency 0.1: L162_MASS over L162, run by run :: sha256 c6c4138ee5c86423 :: prose twin: two significant figures, direction in words
\newcommand{\MlstSculptResMassOverResTenPrMax}{down to 0.79 times the distance}  % experiments/FIGS/data/v2/index.json :: distance at QCD efficiency 0.1: L162_MASS over L162, run by run :: sha256 c6c4138ee5c86423 :: prose twin of the run farther from 1
\newcommand{\MlstSculptCoarseMassOverCoarseTen}{0.94 and 0.79}  % experiments/FIGS/data/v2/index.json :: distance at QCD efficiency 0.1: R16_Q1_MASS over R16_Q1, run by run :: sha256 c6c4138ee5c86423
\newcommand{\MlstSculptCoarseMassOverCoarseTenPr}{equal within 30\%}  % experiments/FIGS/data/v2/index.json :: distance at QCD efficiency 0.1: R16_Q1_MASS over R16_Q1, run by run :: sha256 c6c4138ee5c86423 :: prose twin: two significant figures, direction in words
\newcommand{\MlstSculptCoarseMassOverCoarseTenPrMax}{equal within 30\%}  % experiments/FIGS/data/v2/index.json :: distance at QCD efficiency 0.1: R16_Q1_MASS over R16_Q1, run by run :: sha256 c6c4138ee5c86423 :: prose twin of the run farther from 1
\newcommand{\MlstSculptCoarseOverFullTen}{1.04\,\ensuremath{\pm}\,0.05}  % experiments/FIGS/data/v2/index.json :: distance at QCD efficiency 0.1: R16_Q1 over L188, run by run :: sha256 c6c4138ee5c86423
\newcommand{\MlstSculptCoarseOverFullTenPr}{equal within 6\%}  % experiments/FIGS/data/v2/index.json :: distance at QCD efficiency 0.1: R16_Q1 over L188, run by run :: sha256 c6c4138ee5c86423 :: prose twin: two significant figures, direction in words
\newcommand{\MlstSculptResOverFullTen}{1.02\,\ensuremath{\pm}\,0.06}  % experiments/FIGS/data/v2/index.json :: distance at QCD efficiency 0.1: L162 over L188, run by run :: sha256 c6c4138ee5c86423
\newcommand{\MlstSculptResOverFullTenPr}{equal within 6\%}  % experiments/FIGS/data/v2/index.json :: distance at QCD efficiency 0.1: L162 over L188, run by run :: sha256 c6c4138ee5c86423 :: prose twin: two significant figures, direction in words
\newcommand{\MlstSculptMidOverFullTen}{1.069\,\ensuremath{\pm}\,0.020}  % experiments/FIGS/data/v2/index.json :: distance at QCD efficiency 0.1: R42_Q1 over L188, run by run :: sha256 c6c4138ee5c86423
\newcommand{\MlstSculptMidOverFullTenPr}{1.1 times the distance}  % experiments/FIGS/data/v2/index.json :: distance at QCD efficiency 0.1: R42_Q1 over L188, run by run :: sha256 c6c4138ee5c86423 :: prose twin: two significant figures, direction in words
\newcommand{\MlstSculptCoarseOverMidTen}{0.971\,\ensuremath{\pm}\,0.035}  % experiments/FIGS/data/v2/index.json :: distance at QCD efficiency 0.1: R16_Q1 over R42_Q1, run by run :: sha256 c6c4138ee5c86423
\newcommand{\MlstSculptCoarseOverMidTenPr}{equal within 4\%}  % experiments/FIGS/data/v2/index.json :: distance at QCD efficiency 0.1: R16_Q1 over R42_Q1, run by run :: sha256 c6c4138ee5c86423 :: prose twin: two significant figures, direction in words
\newcommand{\MlstSculptNPassOne}{4{,}103}  % experiments/FIGS/data/v2/index.json :: QCD test jets passing the cut at QCD efficiency 0.01, every run :: sha256 c6c4138ee5c86423
\newcommand{\MlstSculptFloorOne}{0.0313}  % experiments/FIGS/data/v2/index.json :: Jensen-Shannon distance expected for a random selection of as many QCD jets as pass at 0.01 :: sha256 c6c4138ee5c86423
\newcommand{\MlstSculptFullOne}{0.0439\,\ensuremath{\pm}\,0.0024}  % experiments/FIGS/data/v2/index.json :: Jensen-Shannon distance of the QCD m_SD distribution at QCD efficiency 0.01, L188 :: sha256 c6c4138ee5c86423
\newcommand{\MlstSculptFullOverFloorOne}{1.40\,\ensuremath{\pm}\,0.08}  % experiments/FIGS/data/v2/index.json :: distance at 0.01 over the random-selection floor, L188 :: sha256 c6c4138ee5c86423
\newcommand{\MlstSculptFullOverFloorOnePr}{1.4 times the distance}  % experiments/FIGS/data/v2/index.json :: distance at 0.01 over the random-selection floor, L188 :: sha256 c6c4138ee5c86423 :: prose twin: two significant figures, direction in words
\newcommand{\MlstSculptResOne}{0.0429\,\ensuremath{\pm}\,0.0029}  % experiments/FIGS/data/v2/index.json :: Jensen-Shannon distance of the QCD m_SD distribution at QCD efficiency 0.01, L162 :: sha256 c6c4138ee5c86423
\newcommand{\MlstSculptResOverFloorOne}{1.37\,\ensuremath{\pm}\,0.09}  % experiments/FIGS/data/v2/index.json :: distance at 0.01 over the random-selection floor, L162 :: sha256 c6c4138ee5c86423
\newcommand{\MlstSculptResOverFloorOnePr}{1.4 times the distance}  % experiments/FIGS/data/v2/index.json :: distance at 0.01 over the random-selection floor, L162 :: sha256 c6c4138ee5c86423 :: prose twin: two significant figures, direction in words
\newcommand{\MlstSculptMidOne}{0.051\,\ensuremath{\pm}\,0.007}  % experiments/FIGS/data/v2/index.json :: Jensen-Shannon distance of the QCD m_SD distribution at QCD efficiency 0.01, R42_Q1 :: sha256 c6c4138ee5c86423
\newcommand{\MlstSculptMidOverFloorOne}{1.62\,\ensuremath{\pm}\,0.23}  % experiments/FIGS/data/v2/index.json :: distance at 0.01 over the random-selection floor, R42_Q1 :: sha256 c6c4138ee5c86423
\newcommand{\MlstSculptMidOverFloorOnePr}{1.6 times the distance}  % experiments/FIGS/data/v2/index.json :: distance at 0.01 over the random-selection floor, R42_Q1 :: sha256 c6c4138ee5c86423 :: prose twin: two significant figures, direction in words
\newcommand{\MlstSculptCoarseOne}{0.0561\,\ensuremath{\pm}\,0.0017}  % experiments/FIGS/data/v2/index.json :: Jensen-Shannon distance of the QCD m_SD distribution at QCD efficiency 0.01, R16_Q1 :: sha256 c6c4138ee5c86423
\newcommand{\MlstSculptCoarseOverFloorOne}{1.79\,\ensuremath{\pm}\,0.06}  % experiments/FIGS/data/v2/index.json :: distance at 0.01 over the random-selection floor, R16_Q1 :: sha256 c6c4138ee5c86423
\newcommand{\MlstSculptCoarseOverFloorOnePr}{1.8 times the distance}  % experiments/FIGS/data/v2/index.json :: distance at 0.01 over the random-selection floor, R16_Q1 :: sha256 c6c4138ee5c86423 :: prose twin: two significant figures, direction in words
\newcommand{\MlstSculptResMassOne}{0.0424 and 0.0442}  % experiments/FIGS/data/v2/index.json :: Jensen-Shannon distance of the QCD m_SD distribution at QCD efficiency 0.01, L162_MASS :: sha256 c6c4138ee5c86423
\newcommand{\MlstSculptResMassOverFloorOne}{1.356 and 1.411}  % experiments/FIGS/data/v2/index.json :: distance at 0.01 over the random-selection floor, L162_MASS :: sha256 c6c4138ee5c86423
\newcommand{\MlstSculptResMassOverFloorOnePr}{at least 1.4 times the distance}  % experiments/FIGS/data/v2/index.json :: distance at 0.01 over the random-selection floor, L162_MASS :: sha256 c6c4138ee5c86423 :: prose twin: two significant figures, direction in words
\newcommand{\MlstSculptResMassOverFloorOnePrMax}{up to 1.4 times the distance}  % experiments/FIGS/data/v2/index.json :: distance at 0.01 over the random-selection floor, L162_MASS :: sha256 c6c4138ee5c86423 :: prose twin of the run farther from 1
\newcommand{\MlstSculptCoarseMassOne}{0.0405 and 0.0452}  % experiments/FIGS/data/v2/index.json :: Jensen-Shannon distance of the QCD m_SD distribution at QCD efficiency 0.01, R16_Q1_MASS :: sha256 c6c4138ee5c86423
\newcommand{\MlstSculptCoarseMassOverFloorOne}{1.29 and 1.44}  % experiments/FIGS/data/v2/index.json :: distance at 0.01 over the random-selection floor, R16_Q1_MASS :: sha256 c6c4138ee5c86423
\newcommand{\MlstSculptCoarseMassOverFloorOnePr}{at least 1.3 times the distance}  % experiments/FIGS/data/v2/index.json :: distance at 0.01 over the random-selection floor, R16_Q1_MASS :: sha256 c6c4138ee5c86423 :: prose twin: two significant figures, direction in words
\newcommand{\MlstSculptCoarseMassOverFloorOnePrMax}{up to 1.4 times the distance}  % experiments/FIGS/data/v2/index.json :: distance at 0.01 over the random-selection floor, R16_Q1_MASS :: sha256 c6c4138ee5c86423 :: prose twin of the run farther from 1
\newcommand{\MlstSculptResMassOverResOne}{0.962 and 0.983}  % experiments/FIGS/data/v2/index.json :: distance at QCD efficiency 0.01: L162_MASS over L162, run by run :: sha256 c6c4138ee5c86423
\newcommand{\MlstSculptResMassOverResOnePr}{equal within 4\%}  % experiments/FIGS/data/v2/index.json :: distance at QCD efficiency 0.01: L162_MASS over L162, run by run :: sha256 c6c4138ee5c86423 :: prose twin: two significant figures, direction in words
\newcommand{\MlstSculptResMassOverResOnePrMax}{equal within 4\%}  % experiments/FIGS/data/v2/index.json :: distance at QCD efficiency 0.01: L162_MASS over L162, run by run :: sha256 c6c4138ee5c86423 :: prose twin of the run farther from 1
\newcommand{\MlstSculptCoarseMassOverCoarseOne}{0.70 and 0.82}  % experiments/FIGS/data/v2/index.json :: distance at QCD efficiency 0.01: R16_Q1_MASS over R16_Q1, run by run :: sha256 c6c4138ee5c86423
\newcommand{\MlstSculptCoarseMassOverCoarseOnePr}{at most 0.82 times the distance}  % experiments/FIGS/data/v2/index.json :: distance at QCD efficiency 0.01: R16_Q1_MASS over R16_Q1, run by run :: sha256 c6c4138ee5c86423 :: prose twin: two significant figures, direction in words
\newcommand{\MlstSculptCoarseMassOverCoarseOnePrMax}{down to 0.70 times the distance}  % experiments/FIGS/data/v2/index.json :: distance at QCD efficiency 0.01: R16_Q1_MASS over R16_Q1, run by run :: sha256 c6c4138ee5c86423 :: prose twin of the run farther from 1
\newcommand{\MlstSculptCoarseOverFullOne}{1.28\,\ensuremath{\pm}\,0.04}  % experiments/FIGS/data/v2/index.json :: distance at QCD efficiency 0.01: R16_Q1 over L188, run by run :: sha256 c6c4138ee5c86423
\newcommand{\MlstSculptCoarseOverFullOnePr}{1.3 times the distance}  % experiments/FIGS/data/v2/index.json :: distance at QCD efficiency 0.01: R16_Q1 over L188, run by run :: sha256 c6c4138ee5c86423 :: prose twin: two significant figures, direction in words
\newcommand{\MlstSculptResOverFullOne}{0.98\,\ensuremath{\pm}\,0.05}  % experiments/FIGS/data/v2/index.json :: distance at QCD efficiency 0.01: L162 over L188, run by run :: sha256 c6c4138ee5c86423
\newcommand{\MlstSculptResOverFullOnePr}{equal within 5\%}  % experiments/FIGS/data/v2/index.json :: distance at QCD efficiency 0.01: L162 over L188, run by run :: sha256 c6c4138ee5c86423 :: prose twin: two significant figures, direction in words
\newcommand{\MlstSculptMidOverFullOne}{1.15\,\ensuremath{\pm}\,0.12}  % experiments/FIGS/data/v2/index.json :: distance at QCD efficiency 0.01: R42_Q1 over L188, run by run :: sha256 c6c4138ee5c86423
\newcommand{\MlstSculptMidOverFullOnePr}{equal within 20\%}  % experiments/FIGS/data/v2/index.json :: distance at QCD efficiency 0.01: R42_Q1 over L188, run by run :: sha256 c6c4138ee5c86423 :: prose twin: two significant figures, direction in words
\newcommand{\MlstSculptCoarseOverMidOne}{1.13\,\ensuremath{\pm}\,0.16}  % experiments/FIGS/data/v2/index.json :: distance at QCD efficiency 0.01: R16_Q1 over R42_Q1, run by run :: sha256 c6c4138ee5c86423
\newcommand{\MlstSculptCoarseOverMidOnePr}{equal within 20\%}  % experiments/FIGS/data/v2/index.json :: distance at QCD efficiency 0.01: R16_Q1 over R42_Q1, run by run :: sha256 c6c4138ee5c86423 :: prose twin: two significant figures, direction in words
\newcommand{\MlstSculptDiffEffLo}{0.52\%}  % experiments/FIGS/data/v2/index.json :: lowest stored QCD efficiency of the unbroken stretch around the tighter working point at which the 17/188 distance ratio exceeds 1 by twice its SD :: sha256 c6c4138ee5c86423 :: passing grid points: 0.001, 0.005245, 0.007937, 0.012011
\newcommand{\MlstSculptDiffEffHi}{1.2\%}  % experiments/FIGS/data/v2/index.json :: highest efficiency of that stretch :: sha256 c6c4138ee5c86423
\newcommand{\MlstSculptLowestEff}{0.10\%}  % experiments/FIGS/data/v2/index.json :: lowest stored QCD efficiency of the sculpting curves :: sha256 c6c4138ee5c86423
\newcommand{\MlstSculptLowestNPass}{410}  % experiments/FIGS/data/v2/index.json :: QCD test jets passing at the lowest stored efficiency :: sha256 c6c4138ee5c86423
\newcommand{\MlstSculptLowestMaxFloorSd}{1.7}  % experiments/FIGS/data/v2/index.json :: largest distance of a configuration's mean from the random-selection floor at the lowest stored efficiency, in SDs of the floor over random selections :: sha256 c6c4138ee5c86423
\newcommand{\MlstSculptNConfigs}{6}  % experiments/FIGS/data/v2/index.json :: configurations with a sculpting curve (count) :: sha256 c6c4138ee5c86423
\newcommand{\LpJciiTestQcdPercent}{20.5\%}  % experiments/FIGS/data/v2/jc2_test/counts.json :: QCD share of the jets, per cent :: sha256 3574ba1f6f2a59f0
\newcommand{\LpJciiTestTwoProngPercent}{38.7\%}  % experiments/FIGS/data/v2/jc2_test/counts.json :: two-prong share of the jets, per cent :: sha256 3574ba1f6f2a59f0
\newcommand{\LpJciiTrainQcdPercent}{13.5\%}  % experiments/FT/data/jc2_v2_manifest.json :: QCD share of the jets, per cent :: sha256 4bd68655919b14e9
\newcommand{\LpJciiTrainTwoProngPercent}{15.4\%}  % experiments/FT/data/jc2_v2_manifest.json :: two-prong share of the jets, per cent :: sha256 4bd68655919b14e9
\newcommand{\LpJciiTestMinAucJets}{10}  % experiments/FIGS/data/v2/jc2_test/counts.json :: fewest jets of a class among the macro-AUC rows :: sha256 3574ba1f6f2a59f0
\newcommand{\LpJciiTrainMinJets}{11}  % experiments/FT/data/jc2_v2_manifest.json :: fewest jets of a class in the 10^6 training subset :: sha256 4bd68655919b14e9
\newcommand{\LpNTestJci}{2{,}000{,}000}  % experiments/FIGS/data/v2/index.json :: jc1 test jets the probes score :: sha256 c6c4138ee5c86423
\newcommand{\LpNTrainMaxJci}{1{,}000{,}000}  % experiments/FIGS/data/v2/index.json :: jc1 largest training subset :: sha256 c6c4138ee5c86423
\newcommand{\LpNTestTop}{404{,}000}  % experiments/FIGS/data/v2/index.json :: top test jets the probes score :: sha256 c6c4138ee5c86423
\newcommand{\LpNTrainMaxTop}{1{,}200{,}000}  % experiments/FIGS/data/v2/index.json :: top largest training subset :: sha256 c6c4138ee5c86423
\newcommand{\LpNTestQg}{200{,}000}  % experiments/FIGS/data/v2/index.json :: qg test jets the probes score :: sha256 c6c4138ee5c86423
\newcommand{\LpNTrainMaxQg}{1{,}600{,}000}  % experiments/FIGS/data/v2/index.json :: qg largest training subset :: sha256 c6c4138ee5c86423
\newcommand{\LpNTestAucJci}{2{,}000{,}000}  % experiments/FIGS/data/v2/index.json :: JetClass test jets of the macro AUC :: sha256 c6c4138ee5c86423
\newcommand{\LpPairNValMin}{1{,}068}  % experiments/FIGS/data/v2/index.json :: validation jets of the pair probe with the fewest :: sha256 c6c4138ee5c86423
\newcommand{\LpPairNValMax}{20{,}195}  % experiments/FIGS/data/v2/index.json :: validation jets of the pair probe with the most :: sha256 c6c4138ee5c86423
\newcommand{\LpNValJci}{200{,}000}  % experiments/FIGS/data/v2/index.json :: JetClass probe validation jets :: sha256 c6c4138ee5c86423
\newcommand{\SculptEffHigh}{10\%}  % experiments/EVAL/mass_sculpting.py :: WORKING_POINTS :: sha256 e85d125c85dac054 :: QCD efficiency of a sculpting cut
\newcommand{\SculptEffLow}{1\%}  % experiments/EVAL/mass_sculpting.py :: WORKING_POINTS :: sha256 e85d125c85dac054 :: QCD efficiency of a sculpting cut
\newcommand{\LpMlpHidden}{256}  % experiments/EVAL/probe.py :: fit_mlp hidden layer :: sha256 485db0966e660588 :: units of the MLP probe's hidden layer
\newcommand{\LpMlpSeedsWord}{three}  % experiments/EVAL/probe.py :: MLP_SEEDS :: sha256 485db0966e660588 :: seeds of the MLP probe whose scores are averaged
\newcommand{\MassMinTrainPerClass}{50}  % experiments/EVAL/mass_resolution.py :: MIN_TRAIN_PER_CLASS :: sha256 26acb0e315094bd1 :: fewest training jets of a class the mass regression keeps
\newcommand{\MlstFtLr}{\ensuremath{10^{-4}}}  % experiments/FIGS/data/v2/ft_mlst/manifests.json :: cells[best70].lr :: sha256 a689f21e894f48ae
\newcommand{\MlstFtHeadMult}{50}  % experiments/FIGS/data/v2/ft_mlst/manifests.json :: cells[best70].head_lr_mult :: sha256 a689f21e894f48ae
\newcommand{\MlstFtWeightDecay}{0.01}  % experiments/FIGS/data/v2/ft_mlst/manifests.json :: cells[best70].weight_decay :: sha256 a689f21e894f48ae
\newcommand{\MlstFtBatch}{512}  % experiments/FIGS/data/v2/ft_mlst/manifests.json :: cells[best70].batch_size :: sha256 a689f21e894f48ae
\newcommand{\MlstFtBenchEpochs}{20}  % experiments/FIGS/data/v2/ft_mlst/manifests.json :: cells[best70, top/qg].epochs :: sha256 a689f21e894f48ae
\newcommand{\MlstFtEpochsEThree}{50}  % experiments/FIGS/data/v2/ft_mlst/manifests.json :: cells[best70, leg 1/2, n_train 1000].epochs :: sha256 a689f21e894f48ae
\newcommand{\MlstFtEpochsEFour}{50}  % experiments/FIGS/data/v2/ft_mlst/manifests.json :: cells[best70, leg 1/2, n_train 10000].epochs :: sha256 a689f21e894f48ae
\newcommand{\MlstFtEpochsEFive}{30}  % experiments/FIGS/data/v2/ft_mlst/manifests.json :: cells[best70, leg 1/2, n_train 100000].epochs :: sha256 a689f21e894f48ae
\newcommand{\MlstFtEpochsESix}{10}  % experiments/FIGS/data/v2/ft_mlst/manifests.json :: cells[best70, leg 1/2, n_train 1000000].epochs :: sha256 a689f21e894f48ae
\newcommand{\MlstFtHeadDecayShare}{30\%}  % experiments/FT/k8s/job-ft-v2-legs-best70-l188-s3-raunav.yaml :: no --lr-scheduler: weaver 0.4.17 flat+decay, the head's last epochs :: sha256 19f59aa54a385c11
\newcommand{\MlstFtHeadDecayTo}{1\%}  % experiments/FT/k8s/job-ft-v2-legs-best70-l188-s3-raunav.yaml :: no --lr-scheduler: weaver 0.4.17 flat+decay, final factor :: sha256 19f59aa54a385c11
\newcommand{\MlstFtHeadDropout}{0.1}  % experiments/FT/k8s/job-ft-v2-legs-best70-l188-s3-raunav.yaml :: fc_params dropout of the fine-tuning head :: sha256 19f59aa54a385c11
\newcommand{\MlstFtMinEpochJets}{10{,}000}  % experiments/FIGS/data/v2/ft_mlst/manifests.json :: cells[best70].samples_per_epoch (smallest) :: sha256 a689f21e894f48ae
\newcommand{\MlstFtValJcii}{20{,}480}  % experiments/FIGS/data/v2/ft_mlst/manifests.json :: cells[best70, leg 1].val_jets :: sha256 a689f21e894f48ae
\newcommand{\MlstFtValSmall}{20{,}000}  % experiments/FT/k8s/job-ft-v2-legs-best70-l188-s3-raunav.yaml :: JetClass --samples-per-epoch-val; top and q/g at N <= 1e4 :: sha256 19f59aa54a385c11
\newcommand{\MlstFtValLarge}{200{,}000}  % experiments/FT/k8s/job-ft-v2-bench-best70-l188-s1-top-raunav.yaml :: top and q/g --samples-per-epoch-val at N > 1e4 :: sha256 1e38d4f542119ea3
\newcommand{\MlstFtGpuTypes}{V100}  % experiments/FIGS/data/v2/ft_mlst/manifests.json :: cells[best70].gpu_device_name :: sha256 a689f21e894f48ae
\newcommand{\MlstFtCellsDone}{445}  % experiments/FIGS/data/v2/ft_mlst/manifests.json :: DONE paper-run fine-tuning cells (count) :: sha256 a689f21e894f48ae
\newcommand{\MlstFtRefLr}{\ensuremath{10^{-3}}}  % experiments/FIGS/data/v2/ft_mlst/manifests.json :: cells[ref_v100].lr :: sha256 a689f21e894f48ae
\newcommand{\MlstFtRefSeedsWord}{three}  % experiments/FIGS/data/v2/ft_mlst/manifests.json :: cells[ref_v100].ft_seed (count) :: sha256 a689f21e894f48ae
\newcommand{\MlstFtRefGpuTypes}{V100}  % experiments/FIGS/data/v2/ft_mlst/manifests.json :: cells[ref_v100].gpu_device_name :: sha256 a689f21e894f48ae
\newcommand{\MlstFtRefWeightDecayBench}{0.01}  % experiments/FIGS/data/v2/ft_mlst/manifests.json :: cells[ref_v100, top/qg].weight_decay :: sha256 a689f21e894f48ae
\newcommand{\MlstBootResamples}{200}  % experiments/FIGS/data/v2/index.json :: bootstrap resamples of the test jets :: sha256 c6c4138ee5c86423
\newcommand{\MlstBootTestMaxProbe}{0.21}  % experiments/FIGS/data/v2/index.json :: largest test-sample SD of the ln ratio of 1 - AUC on JetClass-II, linear probe, label sets over the 188-class model, per cent :: sha256 c6c4138ee5c86423
\newcommand{\MlstBootTestMaxFt}{0.15}  % experiments/FIGS/data/v2/index.json :: largest test-sample SD of the ln ratio of 1 - AUC on JetClass-II, fine-tuning, label sets over the 188-class model, per cent :: sha256 c6c4138ee5c86423
\newcommand{\MlstBootTestMaxRthirtyProbe}{3.6}  % experiments/FIGS/data/v2/index.json :: largest test-sample SD of the ln R30 ratio on top and quark/gluon, linear probe (counting), label sets over the 188-class model, per cent :: sha256 c6c4138ee5c86423
\newcommand{\MlstBootTestMaxRthirtyFt}{13}  % experiments/FIGS/data/v2/index.json :: largest test-sample SD of the ln R30 ratio on top and quark/gluon, fine-tuning (counting), label sets over the 188-class model, per cent :: sha256 c6c4138ee5c86423
\newcommand{\MlstFtCmpJcTwoResOverFullEFour}{0.989\,\ensuremath{\pm}\,0.011}  % experiments/FIGS/data/v2/index.json :: fine-tuning jc2 N=10000: (1 - macro_auc_ovr) of L162 over L188, run by run :: sha256 c6c4138ee5c86423
\newcommand{\MlstFtCmpJcTwoResOverFullEFourPr}{equal within 2\%}  % experiments/FIGS/data/v2/index.json :: fine-tuning jc2 N=10000: (1 - macro_auc_ovr) of L162 over L188, run by run :: sha256 c6c4138ee5c86423 :: prose twin: two significant figures, direction in words
\newcommand{\MlstFtCmpJcTwoMidOverFullEFour}{1.191\,\ensuremath{\pm}\,0.030}  % experiments/FIGS/data/v2/index.json :: fine-tuning jc2 N=10000: (1 - macro_auc_ovr) of R42_Q1 over L188, run by run :: sha256 c6c4138ee5c86423
\newcommand{\MlstFtCmpJcTwoMidOverFullEFourPr}{19\% more errors}  % experiments/FIGS/data/v2/index.json :: fine-tuning jc2 N=10000: (1 - macro_auc_ovr) of R42_Q1 over L188, run by run :: sha256 c6c4138ee5c86423 :: prose twin: two significant figures, direction in words
\newcommand{\MlstFtCmpJcTwoCoarseOverFullEFour}{1.353\,\ensuremath{\pm}\,0.003}  % experiments/FIGS/data/v2/index.json :: fine-tuning jc2 N=10000: (1 - macro_auc_ovr) of R16_Q1 over L188, run by run :: sha256 c6c4138ee5c86423
\newcommand{\MlstFtCmpJcTwoCoarseOverFullEFourPr}{35\% more errors}  % experiments/FIGS/data/v2/index.json :: fine-tuning jc2 N=10000: (1 - macro_auc_ovr) of R16_Q1 over L188, run by run :: sha256 c6c4138ee5c86423 :: prose twin: two significant figures, direction in words
\newcommand{\MlstFtCmpJcTwoResMassOverResEFour}{1.016 and 1.005}  % experiments/FIGS/data/v2/index.json :: fine-tuning jc2 N=10000: (1 - macro_auc_ovr) of L162_MASS over L162, run by run :: sha256 c6c4138ee5c86423
\newcommand{\MlstFtCmpJcTwoResMassOverResEFourPr}{equal within 2\%}  % experiments/FIGS/data/v2/index.json :: fine-tuning jc2 N=10000: (1 - macro_auc_ovr) of L162_MASS over L162, run by run :: sha256 c6c4138ee5c86423 :: prose twin: two significant figures, direction in words
\newcommand{\MlstFtCmpJcTwoResMassOverResEFourPrMax}{equal within 2\%}  % experiments/FIGS/data/v2/index.json :: fine-tuning jc2 N=10000: (1 - macro_auc_ovr) of L162_MASS over L162, run by run :: sha256 c6c4138ee5c86423 :: prose twin of the run farther from 1
\newcommand{\MlstFtCmpJcTwoCoarseMassOverCoarseEFour}{0.989 and 1.000}  % experiments/FIGS/data/v2/index.json :: fine-tuning jc2 N=10000: (1 - macro_auc_ovr) of R16_Q1_MASS over R16_Q1, run by run :: sha256 c6c4138ee5c86423
\newcommand{\MlstFtCmpJcTwoCoarseMassOverCoarseEFourPr}{equal within 2\%}  % experiments/FIGS/data/v2/index.json :: fine-tuning jc2 N=10000: (1 - macro_auc_ovr) of R16_Q1_MASS over R16_Q1, run by run :: sha256 c6c4138ee5c86423 :: prose twin: two significant figures, direction in words
\newcommand{\MlstFtCmpJcTwoCoarseMassOverCoarseEFourPrMax}{equal within 2\%}  % experiments/FIGS/data/v2/index.json :: fine-tuning jc2 N=10000: (1 - macro_auc_ovr) of R16_Q1_MASS over R16_Q1, run by run :: sha256 c6c4138ee5c86423 :: prose twin of the run farther from 1
\newcommand{\MlstFtCmpJcTwoResOverFullEFive}{1.011\,\ensuremath{\pm}\,0.005}  % experiments/FIGS/data/v2/index.json :: fine-tuning jc2 N=100000: (1 - macro_auc_ovr) of L162 over L188, run by run :: sha256 c6c4138ee5c86423
\newcommand{\MlstFtCmpJcTwoResOverFullEFivePr}{1.1\% more errors}  % experiments/FIGS/data/v2/index.json :: fine-tuning jc2 N=100000: (1 - macro_auc_ovr) of L162 over L188, run by run :: sha256 c6c4138ee5c86423 :: prose twin: two significant figures, direction in words
\newcommand{\MlstFtCmpJcTwoMidOverFullEFive}{1.064\,\ensuremath{\pm}\,0.009}  % experiments/FIGS/data/v2/index.json :: fine-tuning jc2 N=100000: (1 - macro_auc_ovr) of R42_Q1 over L188, run by run :: sha256 c6c4138ee5c86423
\newcommand{\MlstFtCmpJcTwoMidOverFullEFivePr}{6.4\% more errors}  % experiments/FIGS/data/v2/index.json :: fine-tuning jc2 N=100000: (1 - macro_auc_ovr) of R42_Q1 over L188, run by run :: sha256 c6c4138ee5c86423 :: prose twin: two significant figures, direction in words
\newcommand{\MlstFtCmpJcTwoCoarseOverFullEFive}{1.157\,\ensuremath{\pm}\,0.001}  % experiments/FIGS/data/v2/index.json :: fine-tuning jc2 N=100000: (1 - macro_auc_ovr) of R16_Q1 over L188, run by run :: sha256 c6c4138ee5c86423
\newcommand{\MlstFtCmpJcTwoCoarseOverFullEFivePr}{16\% more errors}  % experiments/FIGS/data/v2/index.json :: fine-tuning jc2 N=100000: (1 - macro_auc_ovr) of R16_Q1 over L188, run by run :: sha256 c6c4138ee5c86423 :: prose twin: two significant figures, direction in words
\newcommand{\MlstFtCmpJcTwoResMassOverResEFive}{1.027 and 1.021}  % experiments/FIGS/data/v2/index.json :: fine-tuning jc2 N=100000: (1 - macro_auc_ovr) of L162_MASS over L162, run by run :: sha256 c6c4138ee5c86423
\newcommand{\MlstFtCmpJcTwoResMassOverResEFivePr}{at least 2.1\% more errors}  % experiments/FIGS/data/v2/index.json :: fine-tuning jc2 N=100000: (1 - macro_auc_ovr) of L162_MASS over L162, run by run :: sha256 c6c4138ee5c86423 :: prose twin: two significant figures, direction in words
\newcommand{\MlstFtCmpJcTwoResMassOverResEFivePrMax}{up to 2.7\% more errors}  % experiments/FIGS/data/v2/index.json :: fine-tuning jc2 N=100000: (1 - macro_auc_ovr) of L162_MASS over L162, run by run :: sha256 c6c4138ee5c86423 :: prose twin of the run farther from 1
\newcommand{\MlstFtCmpJcTwoCoarseMassOverCoarseEFive}{1.029 and 1.022}  % experiments/FIGS/data/v2/index.json :: fine-tuning jc2 N=100000: (1 - macro_auc_ovr) of R16_Q1_MASS over R16_Q1, run by run :: sha256 c6c4138ee5c86423
\newcommand{\MlstFtCmpJcTwoCoarseMassOverCoarseEFivePr}{at least 2.2\% more errors}  % experiments/FIGS/data/v2/index.json :: fine-tuning jc2 N=100000: (1 - macro_auc_ovr) of R16_Q1_MASS over R16_Q1, run by run :: sha256 c6c4138ee5c86423 :: prose twin: two significant figures, direction in words
\newcommand{\MlstFtCmpJcTwoCoarseMassOverCoarseEFivePrMax}{up to 2.9\% more errors}  % experiments/FIGS/data/v2/index.json :: fine-tuning jc2 N=100000: (1 - macro_auc_ovr) of R16_Q1_MASS over R16_Q1, run by run :: sha256 c6c4138ee5c86423 :: prose twin of the run farther from 1
\newcommand{\MlstFtCmpJcTwoResOverFullESix}{1.014\,\ensuremath{\pm}\,0.011}  % experiments/FIGS/data/v2/index.json :: fine-tuning jc2 N=1000000: (1 - macro_auc_ovr) of L162 over L188, run by run :: sha256 c6c4138ee5c86423
\newcommand{\MlstFtCmpJcTwoResOverFullESixPr}{equal within 2\%}  % experiments/FIGS/data/v2/index.json :: fine-tuning jc2 N=1000000: (1 - macro_auc_ovr) of L162 over L188, run by run :: sha256 c6c4138ee5c86423 :: prose twin: two significant figures, direction in words
\newcommand{\MlstFtCmpJcTwoMidOverFullESix}{1.043\,\ensuremath{\pm}\,0.002}  % experiments/FIGS/data/v2/index.json :: fine-tuning jc2 N=1000000: (1 - macro_auc_ovr) of R42_Q1 over L188, run by run :: sha256 c6c4138ee5c86423
\newcommand{\MlstFtCmpJcTwoMidOverFullESixPr}{4.3\% more errors}  % experiments/FIGS/data/v2/index.json :: fine-tuning jc2 N=1000000: (1 - macro_auc_ovr) of R42_Q1 over L188, run by run :: sha256 c6c4138ee5c86423 :: prose twin: two significant figures, direction in words
\newcommand{\MlstFtCmpJcTwoCoarseOverFullESix}{1.123\,\ensuremath{\pm}\,0.003}  % experiments/FIGS/data/v2/index.json :: fine-tuning jc2 N=1000000: (1 - macro_auc_ovr) of R16_Q1 over L188, run by run :: sha256 c6c4138ee5c86423
\newcommand{\MlstFtCmpJcTwoCoarseOverFullESixPr}{12\% more errors}  % experiments/FIGS/data/v2/index.json :: fine-tuning jc2 N=1000000: (1 - macro_auc_ovr) of R16_Q1 over L188, run by run :: sha256 c6c4138ee5c86423 :: prose twin: two significant figures, direction in words
\newcommand{\MlstFtCmpJcTwoResMassOverResESix}{1.009 and 1.016}  % experiments/FIGS/data/v2/index.json :: fine-tuning jc2 N=1000000: (1 - macro_auc_ovr) of L162_MASS over L162, run by run :: sha256 c6c4138ee5c86423
\newcommand{\MlstFtCmpJcTwoResMassOverResESixPr}{equal within 2\%}  % experiments/FIGS/data/v2/index.json :: fine-tuning jc2 N=1000000: (1 - macro_auc_ovr) of L162_MASS over L162, run by run :: sha256 c6c4138ee5c86423 :: prose twin: two significant figures, direction in words
\newcommand{\MlstFtCmpJcTwoResMassOverResESixPrMax}{equal within 2\%}  % experiments/FIGS/data/v2/index.json :: fine-tuning jc2 N=1000000: (1 - macro_auc_ovr) of L162_MASS over L162, run by run :: sha256 c6c4138ee5c86423 :: prose twin of the run farther from 1
\newcommand{\MlstFtCmpJcTwoCoarseMassOverCoarseESix}{1.023 and 1.017}  % experiments/FIGS/data/v2/index.json :: fine-tuning jc2 N=1000000: (1 - macro_auc_ovr) of R16_Q1_MASS over R16_Q1, run by run :: sha256 c6c4138ee5c86423
\newcommand{\MlstFtCmpJcTwoCoarseMassOverCoarseESixPr}{at least 1.7\% more errors}  % experiments/FIGS/data/v2/index.json :: fine-tuning jc2 N=1000000: (1 - macro_auc_ovr) of R16_Q1_MASS over R16_Q1, run by run :: sha256 c6c4138ee5c86423 :: prose twin: two significant figures, direction in words
\newcommand{\MlstFtCmpJcTwoCoarseMassOverCoarseESixPrMax}{up to 2.3\% more errors}  % experiments/FIGS/data/v2/index.json :: fine-tuning jc2 N=1000000: (1 - macro_auc_ovr) of R16_Q1_MASS over R16_Q1, run by run :: sha256 c6c4138ee5c86423 :: prose twin of the run farther from 1
\newcommand{\MlstFtCmpJcTwoQcdResOverFullEFour}{1.038\,\ensuremath{\pm}\,0.005}  % experiments/FIGS/data/v2/index.json :: fine-tuning jc2_qcd N=10000: (1 - macro_auc_qcd) of L162 over L188, run by run :: sha256 c6c4138ee5c86423
\newcommand{\MlstFtCmpJcTwoQcdResOverFullEFourPr}{3.8\% more errors}  % experiments/FIGS/data/v2/index.json :: fine-tuning jc2_qcd N=10000: (1 - macro_auc_qcd) of L162 over L188, run by run :: sha256 c6c4138ee5c86423 :: prose twin: two significant figures, direction in words
\newcommand{\MlstFtCmpJcTwoQcdMidOverFullEFour}{1.164\,\ensuremath{\pm}\,0.018}  % experiments/FIGS/data/v2/index.json :: fine-tuning jc2_qcd N=10000: (1 - macro_auc_qcd) of R42_Q1 over L188, run by run :: sha256 c6c4138ee5c86423
\newcommand{\MlstFtCmpJcTwoQcdMidOverFullEFourPr}{16\% more errors}  % experiments/FIGS/data/v2/index.json :: fine-tuning jc2_qcd N=10000: (1 - macro_auc_qcd) of R42_Q1 over L188, run by run :: sha256 c6c4138ee5c86423 :: prose twin: two significant figures, direction in words
\newcommand{\MlstFtCmpJcTwoQcdCoarseOverFullEFour}{1.258\,\ensuremath{\pm}\,0.025}  % experiments/FIGS/data/v2/index.json :: fine-tuning jc2_qcd N=10000: (1 - macro_auc_qcd) of R16_Q1 over L188, run by run :: sha256 c6c4138ee5c86423
\newcommand{\MlstFtCmpJcTwoQcdCoarseOverFullEFourPr}{26\% more errors}  % experiments/FIGS/data/v2/index.json :: fine-tuning jc2_qcd N=10000: (1 - macro_auc_qcd) of R16_Q1 over L188, run by run :: sha256 c6c4138ee5c86423 :: prose twin: two significant figures, direction in words
\newcommand{\MlstFtCmpJcTwoQcdResMassOverResEFour}{1.027 and 0.997}  % experiments/FIGS/data/v2/index.json :: fine-tuning jc2_qcd N=10000: (1 - macro_auc_qcd) of L162_MASS over L162, run by run :: sha256 c6c4138ee5c86423
\newcommand{\MlstFtCmpJcTwoQcdResMassOverResEFourPr}{equal within 3\%}  % experiments/FIGS/data/v2/index.json :: fine-tuning jc2_qcd N=10000: (1 - macro_auc_qcd) of L162_MASS over L162, run by run :: sha256 c6c4138ee5c86423 :: prose twin: two significant figures, direction in words
\newcommand{\MlstFtCmpJcTwoQcdResMassOverResEFourPrMax}{equal within 3\%}  % experiments/FIGS/data/v2/index.json :: fine-tuning jc2_qcd N=10000: (1 - macro_auc_qcd) of L162_MASS over L162, run by run :: sha256 c6c4138ee5c86423 :: prose twin of the run farther from 1
\newcommand{\MlstFtCmpJcTwoQcdCoarseMassOverCoarseEFour}{1.028 and 1.048}  % experiments/FIGS/data/v2/index.json :: fine-tuning jc2_qcd N=10000: (1 - macro_auc_qcd) of R16_Q1_MASS over R16_Q1, run by run :: sha256 c6c4138ee5c86423
\newcommand{\MlstFtCmpJcTwoQcdCoarseMassOverCoarseEFourPr}{equal within 5\%}  % experiments/FIGS/data/v2/index.json :: fine-tuning jc2_qcd N=10000: (1 - macro_auc_qcd) of R16_Q1_MASS over R16_Q1, run by run :: sha256 c6c4138ee5c86423 :: prose twin: two significant figures, direction in words
\newcommand{\MlstFtCmpJcTwoQcdCoarseMassOverCoarseEFourPrMax}{equal within 5\%}  % experiments/FIGS/data/v2/index.json :: fine-tuning jc2_qcd N=10000: (1 - macro_auc_qcd) of R16_Q1_MASS over R16_Q1, run by run :: sha256 c6c4138ee5c86423 :: prose twin of the run farther from 1
\newcommand{\MlstFtCmpJcTwoQcdResOverFullEFive}{1.060\,\ensuremath{\pm}\,0.002}  % experiments/FIGS/data/v2/index.json :: fine-tuning jc2_qcd N=100000: (1 - macro_auc_qcd) of L162 over L188, run by run :: sha256 c6c4138ee5c86423
\newcommand{\MlstFtCmpJcTwoQcdResOverFullEFivePr}{6.0\% more errors}  % experiments/FIGS/data/v2/index.json :: fine-tuning jc2_qcd N=100000: (1 - macro_auc_qcd) of L162 over L188, run by run :: sha256 c6c4138ee5c86423 :: prose twin: two significant figures, direction in words
\newcommand{\MlstFtCmpJcTwoQcdMidOverFullEFive}{1.062\,\ensuremath{\pm}\,0.010}  % experiments/FIGS/data/v2/index.json :: fine-tuning jc2_qcd N=100000: (1 - macro_auc_qcd) of R42_Q1 over L188, run by run :: sha256 c6c4138ee5c86423
\newcommand{\MlstFtCmpJcTwoQcdMidOverFullEFivePr}{6.2\% more errors}  % experiments/FIGS/data/v2/index.json :: fine-tuning jc2_qcd N=100000: (1 - macro_auc_qcd) of R42_Q1 over L188, run by run :: sha256 c6c4138ee5c86423 :: prose twin: two significant figures, direction in words
\newcommand{\MlstFtCmpJcTwoQcdCoarseOverFullEFive}{1.137\,\ensuremath{\pm}\,0.001}  % experiments/FIGS/data/v2/index.json :: fine-tuning jc2_qcd N=100000: (1 - macro_auc_qcd) of R16_Q1 over L188, run by run :: sha256 c6c4138ee5c86423
\newcommand{\MlstFtCmpJcTwoQcdCoarseOverFullEFivePr}{14\% more errors}  % experiments/FIGS/data/v2/index.json :: fine-tuning jc2_qcd N=100000: (1 - macro_auc_qcd) of R16_Q1 over L188, run by run :: sha256 c6c4138ee5c86423 :: prose twin: two significant figures, direction in words
\newcommand{\MlstFtCmpJcTwoQcdResMassOverResEFive}{1.022 and 1.009}  % experiments/FIGS/data/v2/index.json :: fine-tuning jc2_qcd N=100000: (1 - macro_auc_qcd) of L162_MASS over L162, run by run :: sha256 c6c4138ee5c86423
\newcommand{\MlstFtCmpJcTwoQcdResMassOverResEFivePr}{equal within 3\%}  % experiments/FIGS/data/v2/index.json :: fine-tuning jc2_qcd N=100000: (1 - macro_auc_qcd) of L162_MASS over L162, run by run :: sha256 c6c4138ee5c86423 :: prose twin: two significant figures, direction in words
\newcommand{\MlstFtCmpJcTwoQcdResMassOverResEFivePrMax}{equal within 3\%}  % experiments/FIGS/data/v2/index.json :: fine-tuning jc2_qcd N=100000: (1 - macro_auc_qcd) of L162_MASS over L162, run by run :: sha256 c6c4138ee5c86423 :: prose twin of the run farther from 1
\newcommand{\MlstFtCmpJcTwoQcdCoarseMassOverCoarseEFive}{1.028 and 1.027}  % experiments/FIGS/data/v2/index.json :: fine-tuning jc2_qcd N=100000: (1 - macro_auc_qcd) of R16_Q1_MASS over R16_Q1, run by run :: sha256 c6c4138ee5c86423
\newcommand{\MlstFtCmpJcTwoQcdCoarseMassOverCoarseEFivePr}{at least 2.7\% more errors}  % experiments/FIGS/data/v2/index.json :: fine-tuning jc2_qcd N=100000: (1 - macro_auc_qcd) of R16_Q1_MASS over R16_Q1, run by run :: sha256 c6c4138ee5c86423 :: prose twin: two significant figures, direction in words
\newcommand{\MlstFtCmpJcTwoQcdCoarseMassOverCoarseEFivePrMax}{up to 2.8\% more errors}  % experiments/FIGS/data/v2/index.json :: fine-tuning jc2_qcd N=100000: (1 - macro_auc_qcd) of R16_Q1_MASS over R16_Q1, run by run :: sha256 c6c4138ee5c86423 :: prose twin of the run farther from 1
\newcommand{\MlstFtCmpJcTwoQcdResOverFullESix}{1.035\,\ensuremath{\pm}\,0.010}  % experiments/FIGS/data/v2/index.json :: fine-tuning jc2_qcd N=1000000: (1 - macro_auc_qcd) of L162 over L188, run by run :: sha256 c6c4138ee5c86423
\newcommand{\MlstFtCmpJcTwoQcdResOverFullESixPr}{3.5\% more errors}  % experiments/FIGS/data/v2/index.json :: fine-tuning jc2_qcd N=1000000: (1 - macro_auc_qcd) of L162 over L188, run by run :: sha256 c6c4138ee5c86423 :: prose twin: two significant figures, direction in words
\newcommand{\MlstFtCmpJcTwoQcdMidOverFullESix}{1.050\,\ensuremath{\pm}\,0.004}  % experiments/FIGS/data/v2/index.json :: fine-tuning jc2_qcd N=1000000: (1 - macro_auc_qcd) of R42_Q1 over L188, run by run :: sha256 c6c4138ee5c86423
\newcommand{\MlstFtCmpJcTwoQcdMidOverFullESixPr}{5.0\% more errors}  % experiments/FIGS/data/v2/index.json :: fine-tuning jc2_qcd N=1000000: (1 - macro_auc_qcd) of R42_Q1 over L188, run by run :: sha256 c6c4138ee5c86423 :: prose twin: two significant figures, direction in words
\newcommand{\MlstFtCmpJcTwoQcdCoarseOverFullESix}{1.122\,\ensuremath{\pm}\,0.002}  % experiments/FIGS/data/v2/index.json :: fine-tuning jc2_qcd N=1000000: (1 - macro_auc_qcd) of R16_Q1 over L188, run by run :: sha256 c6c4138ee5c86423
\newcommand{\MlstFtCmpJcTwoQcdCoarseOverFullESixPr}{12\% more errors}  % experiments/FIGS/data/v2/index.json :: fine-tuning jc2_qcd N=1000000: (1 - macro_auc_qcd) of R16_Q1 over L188, run by run :: sha256 c6c4138ee5c86423 :: prose twin: two significant figures, direction in words
\newcommand{\MlstFtCmpJcTwoQcdResMassOverResESix}{1.011 and 1.025}  % experiments/FIGS/data/v2/index.json :: fine-tuning jc2_qcd N=1000000: (1 - macro_auc_qcd) of L162_MASS over L162, run by run :: sha256 c6c4138ee5c86423
\newcommand{\MlstFtCmpJcTwoQcdResMassOverResESixPr}{equal within 3\%}  % experiments/FIGS/data/v2/index.json :: fine-tuning jc2_qcd N=1000000: (1 - macro_auc_qcd) of L162_MASS over L162, run by run :: sha256 c6c4138ee5c86423 :: prose twin: two significant figures, direction in words
\newcommand{\MlstFtCmpJcTwoQcdResMassOverResESixPrMax}{equal within 3\%}  % experiments/FIGS/data/v2/index.json :: fine-tuning jc2_qcd N=1000000: (1 - macro_auc_qcd) of L162_MASS over L162, run by run :: sha256 c6c4138ee5c86423 :: prose twin of the run farther from 1
\newcommand{\MlstFtCmpJcTwoQcdCoarseMassOverCoarseESix}{1.015 and 1.015}  % experiments/FIGS/data/v2/index.json :: fine-tuning jc2_qcd N=1000000: (1 - macro_auc_qcd) of R16_Q1_MASS over R16_Q1, run by run :: sha256 c6c4138ee5c86423
\newcommand{\MlstFtCmpJcTwoQcdCoarseMassOverCoarseESixPr}{at least 1.5\% more errors}  % experiments/FIGS/data/v2/index.json :: fine-tuning jc2_qcd N=1000000: (1 - macro_auc_qcd) of R16_Q1_MASS over R16_Q1, run by run :: sha256 c6c4138ee5c86423 :: prose twin: two significant figures, direction in words
\newcommand{\MlstFtCmpJcTwoQcdCoarseMassOverCoarseESixPrMax}{up to 1.5\% more errors}  % experiments/FIGS/data/v2/index.json :: fine-tuning jc2_qcd N=1000000: (1 - macro_auc_qcd) of R16_Q1_MASS over R16_Q1, run by run :: sha256 c6c4138ee5c86423 :: prose twin of the run farther from 1
\newcommand{\MlstFtJcTwoAccFullEFour}{0.3914\,\ensuremath{\pm}\,0.0011}  % experiments/FIGS/data/v2/index.json :: fine-tuning jc2 N=10000: accuracy, L188 :: sha256 c6c4138ee5c86423
\newcommand{\MlstFtJcTwoAccResEFour}{0.3934\,\ensuremath{\pm}\,0.0018}  % experiments/FIGS/data/v2/index.json :: fine-tuning jc2 N=10000: accuracy, L162 :: sha256 c6c4138ee5c86423
\newcommand{\MlstFtJcTwoAccMidEFour}{0.3563\,\ensuremath{\pm}\,0.0025}  % experiments/FIGS/data/v2/index.json :: fine-tuning jc2 N=10000: accuracy, R42_Q1 :: sha256 c6c4138ee5c86423
\newcommand{\MlstFtJcTwoAccCoarseEFour}{0.3219\,\ensuremath{\pm}\,0.0011}  % experiments/FIGS/data/v2/index.json :: fine-tuning jc2 N=10000: accuracy, R16_Q1 :: sha256 c6c4138ee5c86423
\newcommand{\MlstFtJcTwoAccResMassEFour}{0.381 and 0.392}  % experiments/FIGS/data/v2/index.json :: fine-tuning jc2 N=10000: accuracy, L162_MASS :: sha256 c6c4138ee5c86423
\newcommand{\MlstFtJcTwoAccCoarseMassEFour}{0.3156 and 0.3187}  % experiments/FIGS/data/v2/index.json :: fine-tuning jc2 N=10000: accuracy, R16_Q1_MASS :: sha256 c6c4138ee5c86423
\newcommand{\MlstFtJcTwoAccScratchEFour}{0.1492\,\ensuremath{\pm}\,0.0011}  % experiments/FIGS/data/v2/index.json :: fine-tuning jc2 N=10000: accuracy from scratch :: sha256 c6c4138ee5c86423
\newcommand{\MlstFtJcTwoAccFullMinusScratchEFour}{0.2423\,\ensuremath{\pm}\,0.0011}  % experiments/FIGS/data/v2/index.json :: fine-tuning jc2 N=10000: accuracy of each 188-class run minus the from-scratch mean (absolute) :: sha256 c6c4138ee5c86423
\newcommand{\MlstFtJcTwoAccFullEFive}{0.4164\,\ensuremath{\pm}\,0.0015}  % experiments/FIGS/data/v2/index.json :: fine-tuning jc2 N=100000: accuracy, L188 :: sha256 c6c4138ee5c86423
\newcommand{\MlstFtJcTwoAccResEFive}{0.4161\,\ensuremath{\pm}\,0.0015}  % experiments/FIGS/data/v2/index.json :: fine-tuning jc2 N=100000: accuracy, L162 :: sha256 c6c4138ee5c86423
\newcommand{\MlstFtJcTwoAccMidEFive}{0.4007\,\ensuremath{\pm}\,0.0013}  % experiments/FIGS/data/v2/index.json :: fine-tuning jc2 N=100000: accuracy, R42_Q1 :: sha256 c6c4138ee5c86423
\newcommand{\MlstFtJcTwoAccCoarseEFive}{0.3833\,\ensuremath{\pm}\,0.0020}  % experiments/FIGS/data/v2/index.json :: fine-tuning jc2 N=100000: accuracy, R16_Q1 :: sha256 c6c4138ee5c86423
\newcommand{\MlstFtJcTwoAccResMassEFive}{0.4093 and 0.4115}  % experiments/FIGS/data/v2/index.json :: fine-tuning jc2 N=100000: accuracy, L162_MASS :: sha256 c6c4138ee5c86423
\newcommand{\MlstFtJcTwoAccCoarseMassEFive}{0.3758 and 0.3781}  % experiments/FIGS/data/v2/index.json :: fine-tuning jc2 N=100000: accuracy, R16_Q1_MASS :: sha256 c6c4138ee5c86423
\newcommand{\MlstFtJcTwoAccScratchEFive}{0.292\,\ensuremath{\pm}\,0.004}  % experiments/FIGS/data/v2/index.json :: fine-tuning jc2 N=100000: accuracy from scratch :: sha256 c6c4138ee5c86423
\newcommand{\MlstFtJcTwoAccFullMinusScratchEFive}{0.1243\,\ensuremath{\pm}\,0.0015}  % experiments/FIGS/data/v2/index.json :: fine-tuning jc2 N=100000: accuracy of each 188-class run minus the from-scratch mean (absolute) :: sha256 c6c4138ee5c86423
\newcommand{\MlstFtJcTwoAccFullESix}{0.4263\,\ensuremath{\pm}\,0.0005}  % experiments/FIGS/data/v2/index.json :: fine-tuning jc2 N=1000000: accuracy, L188 :: sha256 c6c4138ee5c86423
\newcommand{\MlstFtJcTwoAccResESix}{0.4249\,\ensuremath{\pm}\,0.0021}  % experiments/FIGS/data/v2/index.json :: fine-tuning jc2 N=1000000: accuracy, L162 :: sha256 c6c4138ee5c86423
\newcommand{\MlstFtJcTwoAccMidESix}{0.4166\,\ensuremath{\pm}\,0.0011}  % experiments/FIGS/data/v2/index.json :: fine-tuning jc2 N=1000000: accuracy, R42_Q1 :: sha256 c6c4138ee5c86423
\newcommand{\MlstFtJcTwoAccCoarseESix}{0.4054\,\ensuremath{\pm}\,0.0020}  % experiments/FIGS/data/v2/index.json :: fine-tuning jc2 N=1000000: accuracy, R16_Q1 :: sha256 c6c4138ee5c86423
\newcommand{\MlstFtJcTwoAccResMassESix}{0.4198 and 0.4213}  % experiments/FIGS/data/v2/index.json :: fine-tuning jc2 N=1000000: accuracy, L162_MASS :: sha256 c6c4138ee5c86423
\newcommand{\MlstFtJcTwoAccCoarseMassESix}{0.4012 and 0.4000}  % experiments/FIGS/data/v2/index.json :: fine-tuning jc2 N=1000000: accuracy, R16_Q1_MASS :: sha256 c6c4138ee5c86423
\newcommand{\MlstFtJcTwoAccScratchESix}{0.3677\,\ensuremath{\pm}\,0.0012}  % experiments/FIGS/data/v2/index.json :: fine-tuning jc2 N=1000000: accuracy from scratch :: sha256 c6c4138ee5c86423
\newcommand{\MlstFtJcTwoAccFullMinusScratchESix}{0.0585\,\ensuremath{\pm}\,0.0005}  % experiments/FIGS/data/v2/index.json :: fine-tuning jc2 N=1000000: accuracy of each 188-class run minus the from-scratch mean (absolute) :: sha256 c6c4138ee5c86423
\newcommand{\MlstFtJcTwoAucFull}{0.9797\,\ensuremath{\pm}\,0.0001}  % experiments/FIGS/data/v2/index.json :: fine-tuning class-balanced JetClass-II N=1000000: AUC averaged over the classes, Full :: sha256 c6c4138ee5c86423
\newcommand{\MlstFtJcTwoAucScratch}{0.9734\,\ensuremath{\pm}\,0.0003}  % experiments/FIGS/data/v2/index.json :: fine-tuning class-balanced JetClass-II N=1000000: AUC averaged over the classes, Scratch :: sha256 c6c4138ee5c86423
\newcommand{\MlstFtJcOneAccFullEFour}{0.8117\,\ensuremath{\pm}\,0.0031}  % experiments/FIGS/data/v2/index.json :: fine-tuning jc1 N=10000: accuracy, L188 :: sha256 c6c4138ee5c86423
\newcommand{\MlstFtJcOneAccResEFour}{0.8117\,\ensuremath{\pm}\,0.0013}  % experiments/FIGS/data/v2/index.json :: fine-tuning jc1 N=10000: accuracy, L162 :: sha256 c6c4138ee5c86423
\newcommand{\MlstFtJcOneAccMidEFour}{0.8086\,\ensuremath{\pm}\,0.0031}  % experiments/FIGS/data/v2/index.json :: fine-tuning jc1 N=10000: accuracy, R42_Q1 :: sha256 c6c4138ee5c86423
\newcommand{\MlstFtJcOneAccCoarseEFour}{0.7982\,\ensuremath{\pm}\,0.0023}  % experiments/FIGS/data/v2/index.json :: fine-tuning jc1 N=10000: accuracy, R16_Q1 :: sha256 c6c4138ee5c86423
\newcommand{\MlstFtJcOneAccResMassEFour}{0.8107 and 0.8098}  % experiments/FIGS/data/v2/index.json :: fine-tuning jc1 N=10000: accuracy, L162_MASS :: sha256 c6c4138ee5c86423
\newcommand{\MlstFtJcOneAccCoarseMassEFour}{0.7944 and 0.7990}  % experiments/FIGS/data/v2/index.json :: fine-tuning jc1 N=10000: accuracy, R16_Q1_MASS :: sha256 c6c4138ee5c86423
\newcommand{\MlstFtJcOneAccScratchEFour}{0.677\,\ensuremath{\pm}\,0.004}  % experiments/FIGS/data/v2/index.json :: fine-tuning jc1 N=10000: accuracy from scratch :: sha256 c6c4138ee5c86423
\newcommand{\MlstFtJcOneAccFullMinusScratchEFour}{0.1343\,\ensuremath{\pm}\,0.0031}  % experiments/FIGS/data/v2/index.json :: fine-tuning jc1 N=10000: accuracy of each 188-class run minus the from-scratch mean (absolute) :: sha256 c6c4138ee5c86423
\newcommand{\MlstFtJcOneAccFullEFive}{0.8343\,\ensuremath{\pm}\,0.0035}  % experiments/FIGS/data/v2/index.json :: fine-tuning jc1 N=100000: accuracy, L188 :: sha256 c6c4138ee5c86423
\newcommand{\MlstFtJcOneAccResEFive}{0.8328\,\ensuremath{\pm}\,0.0005}  % experiments/FIGS/data/v2/index.json :: fine-tuning jc1 N=100000: accuracy, L162 :: sha256 c6c4138ee5c86423
\newcommand{\MlstFtJcOneAccMidEFive}{0.8324\,\ensuremath{\pm}\,0.0004}  % experiments/FIGS/data/v2/index.json :: fine-tuning jc1 N=100000: accuracy, R42_Q1 :: sha256 c6c4138ee5c86423
\newcommand{\MlstFtJcOneAccCoarseEFive}{0.8288\,\ensuremath{\pm}\,0.0005}  % experiments/FIGS/data/v2/index.json :: fine-tuning jc1 N=100000: accuracy, R16_Q1 :: sha256 c6c4138ee5c86423
\newcommand{\MlstFtJcOneAccResMassEFive}{0.8326 and 0.8348}  % experiments/FIGS/data/v2/index.json :: fine-tuning jc1 N=100000: accuracy, L162_MASS :: sha256 c6c4138ee5c86423
\newcommand{\MlstFtJcOneAccCoarseMassEFive}{0.8253 and 0.8262}  % experiments/FIGS/data/v2/index.json :: fine-tuning jc1 N=100000: accuracy, R16_Q1_MASS :: sha256 c6c4138ee5c86423
\newcommand{\MlstFtJcOneAccScratchEFive}{0.777\,\ensuremath{\pm}\,0.005}  % experiments/FIGS/data/v2/index.json :: fine-tuning jc1 N=100000: accuracy from scratch :: sha256 c6c4138ee5c86423
\newcommand{\MlstFtJcOneAccFullMinusScratchEFive}{0.0571\,\ensuremath{\pm}\,0.0035}  % experiments/FIGS/data/v2/index.json :: fine-tuning jc1 N=100000: accuracy of each 188-class run minus the from-scratch mean (absolute) :: sha256 c6c4138ee5c86423
\newcommand{\MlstFtJcOneAccFullESix}{0.8468\,\ensuremath{\pm}\,0.0006}  % experiments/FIGS/data/v2/index.json :: fine-tuning jc1 N=1000000: accuracy, L188 :: sha256 c6c4138ee5c86423
\newcommand{\MlstFtJcOneAccResESix}{0.8475\,\ensuremath{\pm}\,0.0009}  % experiments/FIGS/data/v2/index.json :: fine-tuning jc1 N=1000000: accuracy, L162 :: sha256 c6c4138ee5c86423
\newcommand{\MlstFtJcOneAccMidESix}{0.8459\,\ensuremath{\pm}\,0.0002}  % experiments/FIGS/data/v2/index.json :: fine-tuning jc1 N=1000000: accuracy, R42_Q1 :: sha256 c6c4138ee5c86423
\newcommand{\MlstFtJcOneAccCoarseESix}{0.8422\,\ensuremath{\pm}\,0.0005}  % experiments/FIGS/data/v2/index.json :: fine-tuning jc1 N=1000000: accuracy, R16_Q1 :: sha256 c6c4138ee5c86423
\newcommand{\MlstFtJcOneAccResMassESix}{0.8459 and 0.8471}  % experiments/FIGS/data/v2/index.json :: fine-tuning jc1 N=1000000: accuracy, L162_MASS :: sha256 c6c4138ee5c86423
\newcommand{\MlstFtJcOneAccCoarseMassESix}{0.8400 and 0.8398}  % experiments/FIGS/data/v2/index.json :: fine-tuning jc1 N=1000000: accuracy, R16_Q1_MASS :: sha256 c6c4138ee5c86423
\newcommand{\MlstFtJcOneAccScratchESix}{0.8267\,\ensuremath{\pm}\,0.0011}  % experiments/FIGS/data/v2/index.json :: fine-tuning jc1 N=1000000: accuracy from scratch :: sha256 c6c4138ee5c86423
\newcommand{\MlstFtJcOneAccFullMinusScratchESix}{0.0201\,\ensuremath{\pm}\,0.0006}  % experiments/FIGS/data/v2/index.json :: fine-tuning jc1 N=1000000: accuracy of each 188-class run minus the from-scratch mean (absolute) :: sha256 c6c4138ee5c86423
\newcommand{\MlstFtTopRthirtyScratchEThree}{49\,\ensuremath{\pm}\,12}  % experiments/FIGS/data/v2/index.json :: fine-tuning top N=1000: R30 from scratch, every seed :: sha256 c6c4138ee5c86423
\newcommand{\MlstFtTopRthirtyFullEThree}{72\,\ensuremath{\pm}\,19}  % experiments/FIGS/data/v2/index.json :: fine-tuning top N=1000: R30, L188 :: sha256 c6c4138ee5c86423
\newcommand{\MlstFtTopRthirtyFullOverScratchEThree}{1.5\,\ensuremath{\pm}\,0.4}  % experiments/FIGS/data/v2/index.json :: fine-tuning top N=1000: R30 of each L188 run over the from-scratch mean, SD over runs and standard error of that mean in quadrature :: sha256 c6c4138ee5c86423
\newcommand{\MlstFtTopRthirtyFullOverScratchEThreePr}{equal within 50\%}  % experiments/FIGS/data/v2/index.json :: fine-tuning top N=1000: R30 of each L188 run over the from-scratch mean, SD over runs and standard error of that mean in quadrature :: sha256 c6c4138ee5c86423 :: prose twin: two significant figures, direction in words
\newcommand{\MlstFtTopRthirtyResEThree}{135\,\ensuremath{\pm}\,26}  % experiments/FIGS/data/v2/index.json :: fine-tuning top N=1000: R30, L162 :: sha256 c6c4138ee5c86423
\newcommand{\MlstFtTopRthirtyResOverScratchEThree}{2.8\,\ensuremath{\pm}\,0.7}  % experiments/FIGS/data/v2/index.json :: fine-tuning top N=1000: R30 of each L162 run over the from-scratch mean, SD over runs and standard error of that mean in quadrature :: sha256 c6c4138ee5c86423
\newcommand{\MlstFtTopRthirtyResOverScratchEThreePr}{2.8 times the rejection}  % experiments/FIGS/data/v2/index.json :: fine-tuning top N=1000: R30 of each L162 run over the from-scratch mean, SD over runs and standard error of that mean in quadrature :: sha256 c6c4138ee5c86423 :: prose twin: two significant figures, direction in words
\newcommand{\MlstFtTopRthirtyMidEThree}{104\,\ensuremath{\pm}\,20}  % experiments/FIGS/data/v2/index.json :: fine-tuning top N=1000: R30, R42_Q1 :: sha256 c6c4138ee5c86423
\newcommand{\MlstFtTopRthirtyMidOverScratchEThree}{2.1\,\ensuremath{\pm}\,0.5}  % experiments/FIGS/data/v2/index.json :: fine-tuning top N=1000: R30 of each R42_Q1 run over the from-scratch mean, SD over runs and standard error of that mean in quadrature :: sha256 c6c4138ee5c86423
\newcommand{\MlstFtTopRthirtyMidOverScratchEThreePr}{2.1 times the rejection}  % experiments/FIGS/data/v2/index.json :: fine-tuning top N=1000: R30 of each R42_Q1 run over the from-scratch mean, SD over runs and standard error of that mean in quadrature :: sha256 c6c4138ee5c86423 :: prose twin: two significant figures, direction in words
\newcommand{\MlstFtTopRthirtyCoarseEThree}{83\,\ensuremath{\pm}\,25}  % experiments/FIGS/data/v2/index.json :: fine-tuning top N=1000: R30, R16_Q1 :: sha256 c6c4138ee5c86423
\newcommand{\MlstFtTopRthirtyCoarseOverScratchEThree}{1.7\,\ensuremath{\pm}\,0.6}  % experiments/FIGS/data/v2/index.json :: fine-tuning top N=1000: R30 of each R16_Q1 run over the from-scratch mean, SD over runs and standard error of that mean in quadrature :: sha256 c6c4138ee5c86423
\newcommand{\MlstFtTopRthirtyCoarseOverScratchEThreePr}{equal within 80\%}  % experiments/FIGS/data/v2/index.json :: fine-tuning top N=1000: R30 of each R16_Q1 run over the from-scratch mean, SD over runs and standard error of that mean in quadrature :: sha256 c6c4138ee5c86423 :: prose twin: two significant figures, direction in words
\newcommand{\MlstFtTopRthirtyResMassEThree}{105.9 and 112.7}  % experiments/FIGS/data/v2/index.json :: fine-tuning top N=1000: R30, L162_MASS :: sha256 c6c4138ee5c86423
\newcommand{\MlstFtTopRthirtyResMassOverScratchEThree}{2.18 and 2.32}  % experiments/FIGS/data/v2/index.json :: fine-tuning top N=1000: R30 of each L162_MASS run over the from-scratch mean, SD over runs and standard error of that mean in quadrature :: sha256 c6c4138ee5c86423
\newcommand{\MlstFtTopRthirtyResMassOverScratchEThreePr}{at least 2.2 times the rejection}  % experiments/FIGS/data/v2/index.json :: fine-tuning top N=1000: R30 of each L162_MASS run over the from-scratch mean, SD over runs and standard error of that mean in quadrature :: sha256 c6c4138ee5c86423 :: prose twin: two significant figures, direction in words
\newcommand{\MlstFtTopRthirtyResMassOverScratchEThreePrMax}{up to 2.3 times the rejection}  % experiments/FIGS/data/v2/index.json :: fine-tuning top N=1000: R30 of each L162_MASS run over the from-scratch mean, SD over runs and standard error of that mean in quadrature :: sha256 c6c4138ee5c86423 :: prose twin of the run farther from 1
\newcommand{\MlstFtTopRthirtyCoarseMassEThree}{59.1 and 63.9}  % experiments/FIGS/data/v2/index.json :: fine-tuning top N=1000: R30, R16_Q1_MASS :: sha256 c6c4138ee5c86423
\newcommand{\MlstFtTopRthirtyCoarseMassOverScratchEThree}{1.216 and 1.314}  % experiments/FIGS/data/v2/index.json :: fine-tuning top N=1000: R30 of each R16_Q1_MASS run over the from-scratch mean, SD over runs and standard error of that mean in quadrature :: sha256 c6c4138ee5c86423
\newcommand{\MlstFtTopRthirtyCoarseMassOverScratchEThreePr}{at least 22\% more rejection}  % experiments/FIGS/data/v2/index.json :: fine-tuning top N=1000: R30 of each R16_Q1_MASS run over the from-scratch mean, SD over runs and standard error of that mean in quadrature :: sha256 c6c4138ee5c86423 :: prose twin: two significant figures, direction in words
\newcommand{\MlstFtTopRthirtyCoarseMassOverScratchEThreePrMax}{up to 31\% more rejection}  % experiments/FIGS/data/v2/index.json :: fine-tuning top N=1000: R30 of each R16_Q1_MASS run over the from-scratch mean, SD over runs and standard error of that mean in quadrature :: sha256 c6c4138ee5c86423 :: prose twin of the run farther from 1
\newcommand{\MlstFtTopRthirtyResOverFullEThree}{2.0\,\ensuremath{\pm}\,0.8}  % experiments/FIGS/data/v2/index.json :: fine-tuning top N=1000: R30 of L162 over L188, run by run :: sha256 c6c4138ee5c86423
\newcommand{\MlstFtTopRthirtyResOverFullEThreePr}{equal within 100\%}  % experiments/FIGS/data/v2/index.json :: fine-tuning top N=1000: R30 of L162 over L188, run by run :: sha256 c6c4138ee5c86423 :: prose twin: two significant figures, direction in words
\newcommand{\MlstFtTopRthirtyMidOverFullEThree}{1.46\,\ensuremath{\pm}\,0.12}  % experiments/FIGS/data/v2/index.json :: fine-tuning top N=1000: R30 of R42_Q1 over L188, run by run :: sha256 c6c4138ee5c86423
\newcommand{\MlstFtTopRthirtyMidOverFullEThreePr}{46\% more rejection}  % experiments/FIGS/data/v2/index.json :: fine-tuning top N=1000: R30 of R42_Q1 over L188, run by run :: sha256 c6c4138ee5c86423 :: prose twin: two significant figures, direction in words
\newcommand{\MlstFtTopRthirtyCoarseOverFullEThree}{1.3\,\ensuremath{\pm}\,0.7}  % experiments/FIGS/data/v2/index.json :: fine-tuning top N=1000: R30 of R16_Q1 over L188, run by run :: sha256 c6c4138ee5c86423
\newcommand{\MlstFtTopRthirtyCoarseOverFullEThreePr}{equal within 80\%}  % experiments/FIGS/data/v2/index.json :: fine-tuning top N=1000: R30 of R16_Q1 over L188, run by run :: sha256 c6c4138ee5c86423 :: prose twin: two significant figures, direction in words
\newcommand{\MlstFtTopRthirtyResMassOverResEThree}{0.73 and 1.07}  % experiments/FIGS/data/v2/index.json :: fine-tuning top N=1000: R30 of L162_MASS over L162, run by run :: sha256 c6c4138ee5c86423
\newcommand{\MlstFtTopRthirtyResMassOverResEThreePr}{equal within 30\%}  % experiments/FIGS/data/v2/index.json :: fine-tuning top N=1000: R30 of L162_MASS over L162, run by run :: sha256 c6c4138ee5c86423 :: prose twin: two significant figures, direction in words
\newcommand{\MlstFtTopRthirtyResMassOverResEThreePrMax}{equal within 30\%}  % experiments/FIGS/data/v2/index.json :: fine-tuning top N=1000: R30 of L162_MASS over L162, run by run :: sha256 c6c4138ee5c86423 :: prose twin of the run farther from 1
\newcommand{\MlstFtTopRthirtyCoarseMassOverCoarseEThree}{0.56 and 0.75}  % experiments/FIGS/data/v2/index.json :: fine-tuning top N=1000: R30 of R16_Q1_MASS over R16_Q1, run by run :: sha256 c6c4138ee5c86423
\newcommand{\MlstFtTopRthirtyCoarseMassOverCoarseEThreePr}{equal within 50\%}  % experiments/FIGS/data/v2/index.json :: fine-tuning top N=1000: R30 of R16_Q1_MASS over R16_Q1, run by run :: sha256 c6c4138ee5c86423 :: prose twin: two significant figures, direction in words
\newcommand{\MlstFtTopRthirtyCoarseMassOverCoarseEThreePrMax}{equal within 50\%}  % experiments/FIGS/data/v2/index.json :: fine-tuning top N=1000: R30 of R16_Q1_MASS over R16_Q1, run by run :: sha256 c6c4138ee5c86423 :: prose twin of the run farther from 1
\newcommand{\MlstFtTopRthirtyRandOverFullEThree}{1.3\,\ensuremath{\pm}\,0.5}  % experiments/FIGS/data/v2/index.json :: fine-tuning top N=1000: R30 of each random partition over the L188 run of its seed :: sha256 c6c4138ee5c86423
\newcommand{\MlstFtTopRthirtyRandOverFullEThreePr}{equal within 50\%}  % experiments/FIGS/data/v2/index.json :: fine-tuning top N=1000: R30 of each random partition over the L188 run of its seed :: sha256 c6c4138ee5c86423 :: prose twin: two significant figures, direction in words
\newcommand{\MlstFtTopRthirtyRandOverCoarseEThree}{0.87\,\ensuremath{\pm}\,0.12}  % experiments/FIGS/data/v2/index.json :: fine-tuning top N=1000: R30 of each random partition over the R16_Q1 run of its seed :: sha256 c6c4138ee5c86423
\newcommand{\MlstFtTopRthirtyRandOverCoarseEThreePr}{equal within 20\%}  % experiments/FIGS/data/v2/index.json :: fine-tuning top N=1000: R30 of each random partition over the R16_Q1 run of its seed :: sha256 c6c4138ee5c86423 :: prose twin: two significant figures, direction in words
\newcommand{\MlstFtTopRthirtyScratchEFour}{340\,\ensuremath{\pm}\,50}  % experiments/FIGS/data/v2/index.json :: fine-tuning top N=10000: R30 from scratch, every seed :: sha256 c6c4138ee5c86423
\newcommand{\MlstFtTopRthirtyFullEFour}{420\,\ensuremath{\pm}\,9}  % experiments/FIGS/data/v2/index.json :: fine-tuning top N=10000: R30, L188 :: sha256 c6c4138ee5c86423
\newcommand{\MlstFtTopRthirtyFullOverScratchEFour}{1.24\,\ensuremath{\pm}\,0.11}  % experiments/FIGS/data/v2/index.json :: fine-tuning top N=10000: R30 of each L188 run over the from-scratch mean, SD over runs and standard error of that mean in quadrature :: sha256 c6c4138ee5c86423
\newcommand{\MlstFtTopRthirtyFullOverScratchEFourPr}{24\% more rejection}  % experiments/FIGS/data/v2/index.json :: fine-tuning top N=10000: R30 of each L188 run over the from-scratch mean, SD over runs and standard error of that mean in quadrature :: sha256 c6c4138ee5c86423 :: prose twin: two significant figures, direction in words
\newcommand{\MlstFtTopRthirtyResEFour}{430\,\ensuremath{\pm}\,60}  % experiments/FIGS/data/v2/index.json :: fine-tuning top N=10000: R30, L162 :: sha256 c6c4138ee5c86423
\newcommand{\MlstFtTopRthirtyResOverScratchEFour}{1.27\,\ensuremath{\pm}\,0.21}  % experiments/FIGS/data/v2/index.json :: fine-tuning top N=10000: R30 of each L162 run over the from-scratch mean, SD over runs and standard error of that mean in quadrature :: sha256 c6c4138ee5c86423
\newcommand{\MlstFtTopRthirtyResOverScratchEFourPr}{equal within 30\%}  % experiments/FIGS/data/v2/index.json :: fine-tuning top N=10000: R30 of each L162 run over the from-scratch mean, SD over runs and standard error of that mean in quadrature :: sha256 c6c4138ee5c86423 :: prose twin: two significant figures, direction in words
\newcommand{\MlstFtTopRthirtyMidEFour}{455\,\ensuremath{\pm}\,30}  % experiments/FIGS/data/v2/index.json :: fine-tuning top N=10000: R30, R42_Q1 :: sha256 c6c4138ee5c86423
\newcommand{\MlstFtTopRthirtyMidOverScratchEFour}{1.34\,\ensuremath{\pm}\,0.14}  % experiments/FIGS/data/v2/index.json :: fine-tuning top N=10000: R30 of each R42_Q1 run over the from-scratch mean, SD over runs and standard error of that mean in quadrature :: sha256 c6c4138ee5c86423
\newcommand{\MlstFtTopRthirtyMidOverScratchEFourPr}{34\% more rejection}  % experiments/FIGS/data/v2/index.json :: fine-tuning top N=10000: R30 of each R42_Q1 run over the from-scratch mean, SD over runs and standard error of that mean in quadrature :: sha256 c6c4138ee5c86423 :: prose twin: two significant figures, direction in words
\newcommand{\MlstFtTopRthirtyCoarseEFour}{350\,\ensuremath{\pm}\,40}  % experiments/FIGS/data/v2/index.json :: fine-tuning top N=10000: R30, R16_Q1 :: sha256 c6c4138ee5c86423
\newcommand{\MlstFtTopRthirtyCoarseOverScratchEFour}{1.04\,\ensuremath{\pm}\,0.14}  % experiments/FIGS/data/v2/index.json :: fine-tuning top N=10000: R30 of each R16_Q1 run over the from-scratch mean, SD over runs and standard error of that mean in quadrature :: sha256 c6c4138ee5c86423
\newcommand{\MlstFtTopRthirtyCoarseOverScratchEFourPr}{equal within 20\%}  % experiments/FIGS/data/v2/index.json :: fine-tuning top N=10000: R30 of each R16_Q1 run over the from-scratch mean, SD over runs and standard error of that mean in quadrature :: sha256 c6c4138ee5c86423 :: prose twin: two significant figures, direction in words
\newcommand{\MlstFtTopRthirtyResMassEFour}{373.2 and 376.7}  % experiments/FIGS/data/v2/index.json :: fine-tuning top N=10000: R30, L162_MASS :: sha256 c6c4138ee5c86423
\newcommand{\MlstFtTopRthirtyResMassOverScratchEFour}{1.101 and 1.111}  % experiments/FIGS/data/v2/index.json :: fine-tuning top N=10000: R30 of each L162_MASS run over the from-scratch mean, SD over runs and standard error of that mean in quadrature :: sha256 c6c4138ee5c86423
\newcommand{\MlstFtTopRthirtyResMassOverScratchEFourPr}{at least 10\% more rejection}  % experiments/FIGS/data/v2/index.json :: fine-tuning top N=10000: R30 of each L162_MASS run over the from-scratch mean, SD over runs and standard error of that mean in quadrature :: sha256 c6c4138ee5c86423 :: prose twin: two significant figures, direction in words
\newcommand{\MlstFtTopRthirtyResMassOverScratchEFourPrMax}{up to 11\% more rejection}  % experiments/FIGS/data/v2/index.json :: fine-tuning top N=10000: R30 of each L162_MASS run over the from-scratch mean, SD over runs and standard error of that mean in quadrature :: sha256 c6c4138ee5c86423 :: prose twin of the run farther from 1
\newcommand{\MlstFtTopRthirtyCoarseMassEFour}{365.8 and 359.9}  % experiments/FIGS/data/v2/index.json :: fine-tuning top N=10000: R30, R16_Q1_MASS :: sha256 c6c4138ee5c86423
\newcommand{\MlstFtTopRthirtyCoarseMassOverScratchEFour}{1.079 and 1.062}  % experiments/FIGS/data/v2/index.json :: fine-tuning top N=10000: R30 of each R16_Q1_MASS run over the from-scratch mean, SD over runs and standard error of that mean in quadrature :: sha256 c6c4138ee5c86423
\newcommand{\MlstFtTopRthirtyCoarseMassOverScratchEFourPr}{at least 6.2\% more rejection}  % experiments/FIGS/data/v2/index.json :: fine-tuning top N=10000: R30 of each R16_Q1_MASS run over the from-scratch mean, SD over runs and standard error of that mean in quadrature :: sha256 c6c4138ee5c86423 :: prose twin: two significant figures, direction in words
\newcommand{\MlstFtTopRthirtyCoarseMassOverScratchEFourPrMax}{up to 7.9\% more rejection}  % experiments/FIGS/data/v2/index.json :: fine-tuning top N=10000: R30 of each R16_Q1_MASS run over the from-scratch mean, SD over runs and standard error of that mean in quadrature :: sha256 c6c4138ee5c86423 :: prose twin of the run farther from 1
\newcommand{\MlstFtTopRthirtyResOverFullEFour}{1.02\,\ensuremath{\pm}\,0.12}  % experiments/FIGS/data/v2/index.json :: fine-tuning top N=10000: R30 of L162 over L188, run by run :: sha256 c6c4138ee5c86423
\newcommand{\MlstFtTopRthirtyResOverFullEFourPr}{equal within 20\%}  % experiments/FIGS/data/v2/index.json :: fine-tuning top N=10000: R30 of L162 over L188, run by run :: sha256 c6c4138ee5c86423 :: prose twin: two significant figures, direction in words
\newcommand{\MlstFtTopRthirtyMidOverFullEFour}{1.08\,\ensuremath{\pm}\,0.08}  % experiments/FIGS/data/v2/index.json :: fine-tuning top N=10000: R30 of R42_Q1 over L188, run by run :: sha256 c6c4138ee5c86423
\newcommand{\MlstFtTopRthirtyMidOverFullEFourPr}{equal within 9\%}  % experiments/FIGS/data/v2/index.json :: fine-tuning top N=10000: R30 of R42_Q1 over L188, run by run :: sha256 c6c4138ee5c86423 :: prose twin: two significant figures, direction in words
\newcommand{\MlstFtTopRthirtyCoarseOverFullEFour}{0.84\,\ensuremath{\pm}\,0.07}  % experiments/FIGS/data/v2/index.json :: fine-tuning top N=10000: R30 of R16_Q1 over L188, run by run :: sha256 c6c4138ee5c86423
\newcommand{\MlstFtTopRthirtyCoarseOverFullEFourPr}{16\% less rejection}  % experiments/FIGS/data/v2/index.json :: fine-tuning top N=10000: R30 of R16_Q1 over L188, run by run :: sha256 c6c4138ee5c86423 :: prose twin: two significant figures, direction in words
\newcommand{\MlstFtTopRthirtyResMassOverResEFour}{1.00 and 0.76}  % experiments/FIGS/data/v2/index.json :: fine-tuning top N=10000: R30 of L162_MASS over L162, run by run :: sha256 c6c4138ee5c86423
\newcommand{\MlstFtTopRthirtyResMassOverResEFourPr}{equal within 30\%}  % experiments/FIGS/data/v2/index.json :: fine-tuning top N=10000: R30 of L162_MASS over L162, run by run :: sha256 c6c4138ee5c86423 :: prose twin: two significant figures, direction in words
\newcommand{\MlstFtTopRthirtyResMassOverResEFourPrMax}{equal within 30\%}  % experiments/FIGS/data/v2/index.json :: fine-tuning top N=10000: R30 of L162_MASS over L162, run by run :: sha256 c6c4138ee5c86423 :: prose twin of the run farther from 1
\newcommand{\MlstFtTopRthirtyCoarseMassOverCoarseEFour}{1.13 and 0.92}  % experiments/FIGS/data/v2/index.json :: fine-tuning top N=10000: R30 of R16_Q1_MASS over R16_Q1, run by run :: sha256 c6c4138ee5c86423
\newcommand{\MlstFtTopRthirtyCoarseMassOverCoarseEFourPr}{equal within 20\%}  % experiments/FIGS/data/v2/index.json :: fine-tuning top N=10000: R30 of R16_Q1_MASS over R16_Q1, run by run :: sha256 c6c4138ee5c86423 :: prose twin: two significant figures, direction in words
\newcommand{\MlstFtTopRthirtyCoarseMassOverCoarseEFourPrMax}{equal within 20\%}  % experiments/FIGS/data/v2/index.json :: fine-tuning top N=10000: R30 of R16_Q1_MASS over R16_Q1, run by run :: sha256 c6c4138ee5c86423 :: prose twin of the run farther from 1
\newcommand{\MlstFtTopRthirtyRandOverFullEFour}{0.83\,\ensuremath{\pm}\,0.10}  % experiments/FIGS/data/v2/index.json :: fine-tuning top N=10000: R30 of each random partition over the L188 run of its seed :: sha256 c6c4138ee5c86423
\newcommand{\MlstFtTopRthirtyRandOverFullEFourPr}{equal within 20\%}  % experiments/FIGS/data/v2/index.json :: fine-tuning top N=10000: R30 of each random partition over the L188 run of its seed :: sha256 c6c4138ee5c86423 :: prose twin: two significant figures, direction in words
\newcommand{\MlstFtTopRthirtyRandOverCoarseEFour}{0.97\,\ensuremath{\pm}\,0.15}  % experiments/FIGS/data/v2/index.json :: fine-tuning top N=10000: R30 of each random partition over the R16_Q1 run of its seed :: sha256 c6c4138ee5c86423
\newcommand{\MlstFtTopRthirtyRandOverCoarseEFourPr}{equal within 20\%}  % experiments/FIGS/data/v2/index.json :: fine-tuning top N=10000: R30 of each random partition over the R16_Q1 run of its seed :: sha256 c6c4138ee5c86423 :: prose twin: two significant figures, direction in words
\newcommand{\MlstFtTopRthirtyScratchEFive}{560\,\ensuremath{\pm}\,90}  % experiments/FIGS/data/v2/index.json :: fine-tuning top N=100000: R30 from scratch, every seed :: sha256 c6c4138ee5c86423
\newcommand{\MlstFtTopRthirtyFullEFive}{850\,\ensuremath{\pm}\,90}  % experiments/FIGS/data/v2/index.json :: fine-tuning top N=100000: R30, L188 :: sha256 c6c4138ee5c86423
\newcommand{\MlstFtTopRthirtyFullOverScratchEFive}{1.50\,\ensuremath{\pm}\,0.21}  % experiments/FIGS/data/v2/index.json :: fine-tuning top N=100000: R30 of each L188 run over the from-scratch mean, SD over runs and standard error of that mean in quadrature :: sha256 c6c4138ee5c86423
\newcommand{\MlstFtTopRthirtyFullOverScratchEFivePr}{50\% more rejection}  % experiments/FIGS/data/v2/index.json :: fine-tuning top N=100000: R30 of each L188 run over the from-scratch mean, SD over runs and standard error of that mean in quadrature :: sha256 c6c4138ee5c86423 :: prose twin: two significant figures, direction in words
\newcommand{\MlstFtTopRthirtyResEFive}{810\,\ensuremath{\pm}\,60}  % experiments/FIGS/data/v2/index.json :: fine-tuning top N=100000: R30, L162 :: sha256 c6c4138ee5c86423
\newcommand{\MlstFtTopRthirtyResOverScratchEFive}{1.44\,\ensuremath{\pm}\,0.18}  % experiments/FIGS/data/v2/index.json :: fine-tuning top N=100000: R30 of each L162 run over the from-scratch mean, SD over runs and standard error of that mean in quadrature :: sha256 c6c4138ee5c86423
\newcommand{\MlstFtTopRthirtyResOverScratchEFivePr}{44\% more rejection}  % experiments/FIGS/data/v2/index.json :: fine-tuning top N=100000: R30 of each L162 run over the from-scratch mean, SD over runs and standard error of that mean in quadrature :: sha256 c6c4138ee5c86423 :: prose twin: two significant figures, direction in words
\newcommand{\MlstFtTopRthirtyMidEFive}{850\,\ensuremath{\pm}\,80}  % experiments/FIGS/data/v2/index.json :: fine-tuning top N=100000: R30, R42_Q1 :: sha256 c6c4138ee5c86423
\newcommand{\MlstFtTopRthirtyMidOverScratchEFive}{1.51\,\ensuremath{\pm}\,0.20}  % experiments/FIGS/data/v2/index.json :: fine-tuning top N=100000: R30 of each R42_Q1 run over the from-scratch mean, SD over runs and standard error of that mean in quadrature :: sha256 c6c4138ee5c86423
\newcommand{\MlstFtTopRthirtyMidOverScratchEFivePr}{51\% more rejection}  % experiments/FIGS/data/v2/index.json :: fine-tuning top N=100000: R30 of each R42_Q1 run over the from-scratch mean, SD over runs and standard error of that mean in quadrature :: sha256 c6c4138ee5c86423 :: prose twin: two significant figures, direction in words
\newcommand{\MlstFtTopRthirtyCoarseEFive}{830\,\ensuremath{\pm}\,60}  % experiments/FIGS/data/v2/index.json :: fine-tuning top N=100000: R30, R16_Q1 :: sha256 c6c4138ee5c86423
\newcommand{\MlstFtTopRthirtyCoarseOverScratchEFive}{1.47\,\ensuremath{\pm}\,0.17}  % experiments/FIGS/data/v2/index.json :: fine-tuning top N=100000: R30 of each R16_Q1 run over the from-scratch mean, SD over runs and standard error of that mean in quadrature :: sha256 c6c4138ee5c86423
\newcommand{\MlstFtTopRthirtyCoarseOverScratchEFivePr}{47\% more rejection}  % experiments/FIGS/data/v2/index.json :: fine-tuning top N=100000: R30 of each R16_Q1 run over the from-scratch mean, SD over runs and standard error of that mean in quadrature :: sha256 c6c4138ee5c86423 :: prose twin: two significant figures, direction in words
\newcommand{\MlstFtTopRthirtyResMassEFive}{848 and 848}  % experiments/FIGS/data/v2/index.json :: fine-tuning top N=100000: R30, L162_MASS :: sha256 c6c4138ee5c86423
\newcommand{\MlstFtTopRthirtyResMassOverScratchEFive}{1.50 and 1.50}  % experiments/FIGS/data/v2/index.json :: fine-tuning top N=100000: R30 of each L162_MASS run over the from-scratch mean, SD over runs and standard error of that mean in quadrature :: sha256 c6c4138ee5c86423
\newcommand{\MlstFtTopRthirtyResMassOverScratchEFivePr}{at least 50\% more rejection}  % experiments/FIGS/data/v2/index.json :: fine-tuning top N=100000: R30 of each L162_MASS run over the from-scratch mean, SD over runs and standard error of that mean in quadrature :: sha256 c6c4138ee5c86423 :: prose twin: two significant figures, direction in words
\newcommand{\MlstFtTopRthirtyResMassOverScratchEFivePrMax}{up to 50\% more rejection}  % experiments/FIGS/data/v2/index.json :: fine-tuning top N=100000: R30 of each L162_MASS run over the from-scratch mean, SD over runs and standard error of that mean in quadrature :: sha256 c6c4138ee5c86423 :: prose twin of the run farther from 1
\newcommand{\MlstFtTopRthirtyCoarseMassEFive}{856 and 774}  % experiments/FIGS/data/v2/index.json :: fine-tuning top N=100000: R30, R16_Q1_MASS :: sha256 c6c4138ee5c86423
\newcommand{\MlstFtTopRthirtyCoarseMassOverScratchEFive}{1.51 and 1.37}  % experiments/FIGS/data/v2/index.json :: fine-tuning top N=100000: R30 of each R16_Q1_MASS run over the from-scratch mean, SD over runs and standard error of that mean in quadrature :: sha256 c6c4138ee5c86423
\newcommand{\MlstFtTopRthirtyCoarseMassOverScratchEFivePr}{at least 37\% more rejection}  % experiments/FIGS/data/v2/index.json :: fine-tuning top N=100000: R30 of each R16_Q1_MASS run over the from-scratch mean, SD over runs and standard error of that mean in quadrature :: sha256 c6c4138ee5c86423 :: prose twin: two significant figures, direction in words
\newcommand{\MlstFtTopRthirtyCoarseMassOverScratchEFivePrMax}{up to 51\% more rejection}  % experiments/FIGS/data/v2/index.json :: fine-tuning top N=100000: R30 of each R16_Q1_MASS run over the from-scratch mean, SD over runs and standard error of that mean in quadrature :: sha256 c6c4138ee5c86423 :: prose twin of the run farther from 1
\newcommand{\MlstFtTopRthirtyResOverFullEFive}{0.96\,\ensuremath{\pm}\,0.05}  % experiments/FIGS/data/v2/index.json :: fine-tuning top N=100000: R30 of L162 over L188, run by run :: sha256 c6c4138ee5c86423
\newcommand{\MlstFtTopRthirtyResOverFullEFivePr}{equal within 6\%}  % experiments/FIGS/data/v2/index.json :: fine-tuning top N=100000: R30 of L162 over L188, run by run :: sha256 c6c4138ee5c86423 :: prose twin: two significant figures, direction in words
\newcommand{\MlstFtTopRthirtyMidOverFullEFive}{1.02\,\ensuremath{\pm}\,0.18}  % experiments/FIGS/data/v2/index.json :: fine-tuning top N=100000: R30 of R42_Q1 over L188, run by run :: sha256 c6c4138ee5c86423
\newcommand{\MlstFtTopRthirtyMidOverFullEFivePr}{equal within 20\%}  % experiments/FIGS/data/v2/index.json :: fine-tuning top N=100000: R30 of R42_Q1 over L188, run by run :: sha256 c6c4138ee5c86423 :: prose twin: two significant figures, direction in words
\newcommand{\MlstFtTopRthirtyCoarseOverFullEFive}{0.99\,\ensuremath{\pm}\,0.17}  % experiments/FIGS/data/v2/index.json :: fine-tuning top N=100000: R30 of R16_Q1 over L188, run by run :: sha256 c6c4138ee5c86423
\newcommand{\MlstFtTopRthirtyCoarseOverFullEFivePr}{equal within 20\%}  % experiments/FIGS/data/v2/index.json :: fine-tuning top N=100000: R30 of R16_Q1 over L188, run by run :: sha256 c6c4138ee5c86423 :: prose twin: two significant figures, direction in words
\newcommand{\MlstFtTopRthirtyResMassOverResEFive}{1.00 and 1.15}  % experiments/FIGS/data/v2/index.json :: fine-tuning top N=100000: R30 of L162_MASS over L162, run by run :: sha256 c6c4138ee5c86423
\newcommand{\MlstFtTopRthirtyResMassOverResEFivePr}{equal within 20\%}  % experiments/FIGS/data/v2/index.json :: fine-tuning top N=100000: R30 of L162_MASS over L162, run by run :: sha256 c6c4138ee5c86423 :: prose twin: two significant figures, direction in words
\newcommand{\MlstFtTopRthirtyResMassOverResEFivePrMax}{equal within 20\%}  % experiments/FIGS/data/v2/index.json :: fine-tuning top N=100000: R30 of L162_MASS over L162, run by run :: sha256 c6c4138ee5c86423 :: prose twin of the run farther from 1
\newcommand{\MlstFtTopRthirtyCoarseMassOverCoarseEFive}{1.08 and 0.87}  % experiments/FIGS/data/v2/index.json :: fine-tuning top N=100000: R30 of R16_Q1_MASS over R16_Q1, run by run :: sha256 c6c4138ee5c86423
\newcommand{\MlstFtTopRthirtyCoarseMassOverCoarseEFivePr}{equal within 20\%}  % experiments/FIGS/data/v2/index.json :: fine-tuning top N=100000: R30 of R16_Q1_MASS over R16_Q1, run by run :: sha256 c6c4138ee5c86423 :: prose twin: two significant figures, direction in words
\newcommand{\MlstFtTopRthirtyCoarseMassOverCoarseEFivePrMax}{equal within 20\%}  % experiments/FIGS/data/v2/index.json :: fine-tuning top N=100000: R30 of R16_Q1_MASS over R16_Q1, run by run :: sha256 c6c4138ee5c86423 :: prose twin of the run farther from 1
\newcommand{\MlstFtTopRthirtyRandOverFullEFive}{1.01\,\ensuremath{\pm}\,0.22}  % experiments/FIGS/data/v2/index.json :: fine-tuning top N=100000: R30 of each random partition over the L188 run of its seed :: sha256 c6c4138ee5c86423
\newcommand{\MlstFtTopRthirtyRandOverFullEFivePr}{equal within 30\%}  % experiments/FIGS/data/v2/index.json :: fine-tuning top N=100000: R30 of each random partition over the L188 run of its seed :: sha256 c6c4138ee5c86423 :: prose twin: two significant figures, direction in words
\newcommand{\MlstFtTopRthirtyRandOverCoarseEFive}{0.97\,\ensuremath{\pm}\,0.10}  % experiments/FIGS/data/v2/index.json :: fine-tuning top N=100000: R30 of each random partition over the R16_Q1 run of its seed :: sha256 c6c4138ee5c86423
\newcommand{\MlstFtTopRthirtyRandOverCoarseEFivePr}{equal within 10\%}  % experiments/FIGS/data/v2/index.json :: fine-tuning top N=100000: R30 of each random partition over the R16_Q1 run of its seed :: sha256 c6c4138ee5c86423 :: prose twin: two significant figures, direction in words
\newcommand{\MlstFtTopRthirtyScratch}{1{,}412\,\ensuremath{\pm}\,20}  % experiments/FIGS/data/v2/index.json :: fine-tuning top N=1200000: R30 from scratch, every seed :: sha256 c6c4138ee5c86423
\newcommand{\MlstFtTopRthirtyFull}{1{,}800\,\ensuremath{\pm}\,60}  % experiments/FIGS/data/v2/index.json :: fine-tuning top N=1200000: R30, L188 :: sha256 c6c4138ee5c86423
\newcommand{\MlstFtTopRthirtyFullOverScratch}{1.27\,\ensuremath{\pm}\,0.04}  % experiments/FIGS/data/v2/index.json :: fine-tuning top N=1200000: R30 of each L188 run over the from-scratch mean, SD over runs and standard error of that mean in quadrature :: sha256 c6c4138ee5c86423
\newcommand{\MlstFtTopRthirtyFullOverScratchPr}{27\% more rejection}  % experiments/FIGS/data/v2/index.json :: fine-tuning top N=1200000: R30 of each L188 run over the from-scratch mean, SD over runs and standard error of that mean in quadrature :: sha256 c6c4138ee5c86423 :: prose twin: two significant figures, direction in words
\newcommand{\MlstFtTopRthirtyRes}{1{,}790\,\ensuremath{\pm}\,60}  % experiments/FIGS/data/v2/index.json :: fine-tuning top N=1200000: R30, L162 :: sha256 c6c4138ee5c86423
\newcommand{\MlstFtTopRthirtyResOverScratch}{1.27\,\ensuremath{\pm}\,0.04}  % experiments/FIGS/data/v2/index.json :: fine-tuning top N=1200000: R30 of each L162 run over the from-scratch mean, SD over runs and standard error of that mean in quadrature :: sha256 c6c4138ee5c86423
\newcommand{\MlstFtTopRthirtyResOverScratchPr}{27\% more rejection}  % experiments/FIGS/data/v2/index.json :: fine-tuning top N=1200000: R30 of each L162 run over the from-scratch mean, SD over runs and standard error of that mean in quadrature :: sha256 c6c4138ee5c86423 :: prose twin: two significant figures, direction in words
\newcommand{\MlstFtTopRthirtyMid}{1{,}780\,\ensuremath{\pm}\,150}  % experiments/FIGS/data/v2/index.json :: fine-tuning top N=1200000: R30, R42_Q1 :: sha256 c6c4138ee5c86423
\newcommand{\MlstFtTopRthirtyMidOverScratch}{1.26\,\ensuremath{\pm}\,0.11}  % experiments/FIGS/data/v2/index.json :: fine-tuning top N=1200000: R30 of each R42_Q1 run over the from-scratch mean, SD over runs and standard error of that mean in quadrature :: sha256 c6c4138ee5c86423
\newcommand{\MlstFtTopRthirtyMidOverScratchPr}{26\% more rejection}  % experiments/FIGS/data/v2/index.json :: fine-tuning top N=1200000: R30 of each R42_Q1 run over the from-scratch mean, SD over runs and standard error of that mean in quadrature :: sha256 c6c4138ee5c86423 :: prose twin: two significant figures, direction in words
\newcommand{\MlstFtTopRthirtyCoarse}{1{,}740\,\ensuremath{\pm}\,140}  % experiments/FIGS/data/v2/index.json :: fine-tuning top N=1200000: R30, R16_Q1 :: sha256 c6c4138ee5c86423
\newcommand{\MlstFtTopRthirtyCoarseOverScratch}{1.23\,\ensuremath{\pm}\,0.10}  % experiments/FIGS/data/v2/index.json :: fine-tuning top N=1200000: R30 of each R16_Q1 run over the from-scratch mean, SD over runs and standard error of that mean in quadrature :: sha256 c6c4138ee5c86423
\newcommand{\MlstFtTopRthirtyCoarseOverScratchPr}{23\% more rejection}  % experiments/FIGS/data/v2/index.json :: fine-tuning top N=1200000: R30 of each R16_Q1 run over the from-scratch mean, SD over runs and standard error of that mean in quadrature :: sha256 c6c4138ee5c86423 :: prose twin: two significant figures, direction in words
\newcommand{\MlstFtTopRthirtyResMass}{1{,}941 and 1{,}590}  % experiments/FIGS/data/v2/index.json :: fine-tuning top N=1200000: R30, L162_MASS :: sha256 c6c4138ee5c86423
\newcommand{\MlstFtTopRthirtyResMassOverScratch}{1.37 and 1.13}  % experiments/FIGS/data/v2/index.json :: fine-tuning top N=1200000: R30 of each L162_MASS run over the from-scratch mean, SD over runs and standard error of that mean in quadrature :: sha256 c6c4138ee5c86423
\newcommand{\MlstFtTopRthirtyResMassOverScratchPr}{at least 13\% more rejection}  % experiments/FIGS/data/v2/index.json :: fine-tuning top N=1200000: R30 of each L162_MASS run over the from-scratch mean, SD over runs and standard error of that mean in quadrature :: sha256 c6c4138ee5c86423 :: prose twin: two significant figures, direction in words
\newcommand{\MlstFtTopRthirtyResMassOverScratchPrMax}{up to 37\% more rejection}  % experiments/FIGS/data/v2/index.json :: fine-tuning top N=1200000: R30 of each L162_MASS run over the from-scratch mean, SD over runs and standard error of that mean in quadrature :: sha256 c6c4138ee5c86423 :: prose twin of the run farther from 1
\newcommand{\MlstFtTopRthirtyCoarseMass}{1{,}683 and 1{,}628}  % experiments/FIGS/data/v2/index.json :: fine-tuning top N=1200000: R30, R16_Q1_MASS :: sha256 c6c4138ee5c86423
\newcommand{\MlstFtTopRthirtyCoarseMassOverScratch}{1.192 and 1.153}  % experiments/FIGS/data/v2/index.json :: fine-tuning top N=1200000: R30 of each R16_Q1_MASS run over the from-scratch mean, SD over runs and standard error of that mean in quadrature :: sha256 c6c4138ee5c86423
\newcommand{\MlstFtTopRthirtyCoarseMassOverScratchPr}{at least 15\% more rejection}  % experiments/FIGS/data/v2/index.json :: fine-tuning top N=1200000: R30 of each R16_Q1_MASS run over the from-scratch mean, SD over runs and standard error of that mean in quadrature :: sha256 c6c4138ee5c86423 :: prose twin: two significant figures, direction in words
\newcommand{\MlstFtTopRthirtyCoarseMassOverScratchPrMax}{up to 19\% more rejection}  % experiments/FIGS/data/v2/index.json :: fine-tuning top N=1200000: R30 of each R16_Q1_MASS run over the from-scratch mean, SD over runs and standard error of that mean in quadrature :: sha256 c6c4138ee5c86423 :: prose twin of the run farther from 1
\newcommand{\MlstFtTopRthirtyResOverFull}{0.997\,\ensuremath{\pm}\,0.032}  % experiments/FIGS/data/v2/index.json :: fine-tuning top N=1200000: R30 of L162 over L188, run by run :: sha256 c6c4138ee5c86423
\newcommand{\MlstFtTopRthirtyResOverFullPr}{equal within 4\%}  % experiments/FIGS/data/v2/index.json :: fine-tuning top N=1200000: R30 of L162 over L188, run by run :: sha256 c6c4138ee5c86423 :: prose twin: two significant figures, direction in words
\newcommand{\MlstFtTopRthirtyMidOverFull}{0.99\,\ensuremath{\pm}\,0.07}  % experiments/FIGS/data/v2/index.json :: fine-tuning top N=1200000: R30 of R42_Q1 over L188, run by run :: sha256 c6c4138ee5c86423
\newcommand{\MlstFtTopRthirtyMidOverFullPr}{equal within 8\%}  % experiments/FIGS/data/v2/index.json :: fine-tuning top N=1200000: R30 of R42_Q1 over L188, run by run :: sha256 c6c4138ee5c86423 :: prose twin: two significant figures, direction in words
\newcommand{\MlstFtTopRthirtyCoarseOverFull}{0.97\,\ensuremath{\pm}\,0.10}  % experiments/FIGS/data/v2/index.json :: fine-tuning top N=1200000: R30 of R16_Q1 over L188, run by run :: sha256 c6c4138ee5c86423
\newcommand{\MlstFtTopRthirtyCoarseOverFullPr}{equal within 10\%}  % experiments/FIGS/data/v2/index.json :: fine-tuning top N=1200000: R30 of R16_Q1 over L188, run by run :: sha256 c6c4138ee5c86423 :: prose twin: two significant figures, direction in words
\newcommand{\MlstFtTopRthirtyResMassOverRes}{1.12 and 0.89}  % experiments/FIGS/data/v2/index.json :: fine-tuning top N=1200000: R30 of L162_MASS over L162, run by run :: sha256 c6c4138ee5c86423
\newcommand{\MlstFtTopRthirtyResMassOverResPr}{equal within 20\%}  % experiments/FIGS/data/v2/index.json :: fine-tuning top N=1200000: R30 of L162_MASS over L162, run by run :: sha256 c6c4138ee5c86423 :: prose twin: two significant figures, direction in words
\newcommand{\MlstFtTopRthirtyResMassOverResPrMax}{equal within 20\%}  % experiments/FIGS/data/v2/index.json :: fine-tuning top N=1200000: R30 of L162_MASS over L162, run by run :: sha256 c6c4138ee5c86423 :: prose twin of the run farther from 1
\newcommand{\MlstFtTopRthirtyCoarseMassOverCoarse}{0.900 and 0.919}  % experiments/FIGS/data/v2/index.json :: fine-tuning top N=1200000: R30 of R16_Q1_MASS over R16_Q1, run by run :: sha256 c6c4138ee5c86423
\newcommand{\MlstFtTopRthirtyCoarseMassOverCoarsePr}{equal within 10\%}  % experiments/FIGS/data/v2/index.json :: fine-tuning top N=1200000: R30 of R16_Q1_MASS over R16_Q1, run by run :: sha256 c6c4138ee5c86423 :: prose twin: two significant figures, direction in words
\newcommand{\MlstFtTopRthirtyCoarseMassOverCoarsePrMax}{equal within 10\%}  % experiments/FIGS/data/v2/index.json :: fine-tuning top N=1200000: R30 of R16_Q1_MASS over R16_Q1, run by run :: sha256 c6c4138ee5c86423 :: prose twin of the run farther from 1
\newcommand{\MlstFtTopRthirtyRandOverFull}{0.97\,\ensuremath{\pm}\,0.04}  % experiments/FIGS/data/v2/index.json :: fine-tuning top N=1200000: R30 of each random partition over the L188 run of its seed :: sha256 c6c4138ee5c86423
\newcommand{\MlstFtTopRthirtyRandOverFullPr}{equal within 4\%}  % experiments/FIGS/data/v2/index.json :: fine-tuning top N=1200000: R30 of each random partition over the L188 run of its seed :: sha256 c6c4138ee5c86423 :: prose twin: two significant figures, direction in words
\newcommand{\MlstFtTopRthirtyRandOverCoarse}{0.97\,\ensuremath{\pm}\,0.09}  % experiments/FIGS/data/v2/index.json :: fine-tuning top N=1200000: R30 of each random partition over the R16_Q1 run of its seed :: sha256 c6c4138ee5c86423
\newcommand{\MlstFtTopRthirtyRandOverCoarsePr}{equal within 9\%}  % experiments/FIGS/data/v2/index.json :: fine-tuning top N=1200000: R30 of each random partition over the R16_Q1 run of its seed :: sha256 c6c4138ee5c86423 :: prose twin: two significant figures, direction in words
\newcommand{\MlstFtTopRthirtyFullOverPart}{0.650\,\ensuremath{\pm}\,0.020}  % experiments/FIGS/data/v2/index.json :: fine-tuning top N=1200000: R30 of each 188-class run over ParT-f.t.'s published value :: sha256 c6c4138ee5c86423
\newcommand{\MlstFtTopRthirtyFullOverPartPr}{35\% less rejection}  % experiments/FIGS/data/v2/index.json :: fine-tuning top N=1200000: R30 of each 188-class run over ParT-f.t.'s published value :: sha256 c6c4138ee5c86423 :: prose twin: two significant figures, direction in words
\newcommand{\MlstFtTopAucFull}{0.9861\,\ensuremath{\pm}\,0.0001}  % experiments/FIGS/data/v2/index.json :: fine-tuning top N=1200000: AUC of the 188-class runs :: sha256 c6c4138ee5c86423
\newcommand{\MlstFtTopAucScratch}{0.9852\,\ensuremath{\pm}\,0.0001}  % experiments/FIGS/data/v2/index.json :: fine-tuning top N=1200000: AUC from scratch :: sha256 c6c4138ee5c86423
\newcommand{\MlstFtTopRthirtyFullRelSd}{3\%}  % experiments/FIGS/data/v2/index.json :: fine-tuning top, full set: relative SD of R30 over the 188-class runs :: sha256 c6c4138ee5c86423
\newcommand{\MlstFtTopRthirtyKin}{1{,}800}  % experiments/FIGS/data/v2/index.json :: fine-tuning top, full set: R30 of the kinematic-only model (one run) :: sha256 c6c4138ee5c86423
\newcommand{\MlstFtTopRthirtyKinOverFull}{1.00}  % experiments/FIGS/data/v2/index.json :: fine-tuning top, full set: R30 of the kinematic-only model over the mean of the 188-class runs :: sha256 c6c4138ee5c86423
\newcommand{\MlstFtTopRthirtyKinOverFullPr}{0.23\% more rejection}  % experiments/FIGS/data/v2/index.json :: fine-tuning top, full set: R30 of the kinematic-only model over the mean of the 188-class runs :: sha256 c6c4138ee5c86423 :: prose twin: two significant figures, direction in words
\newcommand{\MlstFtTopDataEquivJets}{1{,}200}  % experiments/FIGS/data/v2/index.json :: fine-tuning top: from-scratch training jets whose R30 equals the 188-class mean at 1000 jets (interpolated linearly in R30 against log10 N) :: sha256 c6c4138ee5c86423
\newcommand{\MlstFtTopDataEquivJetsEFour}{23{,}000}  % experiments/FIGS/data/v2/index.json :: fine-tuning top: from-scratch training jets whose R30 equals the 188-class mean at 10000 jets (interpolated linearly in R30 against log10 N) :: sha256 c6c4138ee5c86423
\newcommand{\MlstFtQgRthirtyScratchEThree}{49.7\,\ensuremath{\pm}\,2.1}  % experiments/FIGS/data/v2/index.json :: fine-tuning qg N=1000: R30 from scratch, every seed :: sha256 c6c4138ee5c86423
\newcommand{\MlstFtQgRthirtyFullEThree}{47\,\ensuremath{\pm}\,10}  % experiments/FIGS/data/v2/index.json :: fine-tuning qg N=1000: R30, L188 :: sha256 c6c4138ee5c86423
\newcommand{\MlstFtQgRthirtyFullOverScratchEThree}{0.95\,\ensuremath{\pm}\,0.19}  % experiments/FIGS/data/v2/index.json :: fine-tuning qg N=1000: R30 of each L188 run over the from-scratch mean, SD over runs and standard error of that mean in quadrature :: sha256 c6c4138ee5c86423
\newcommand{\MlstFtQgRthirtyFullOverScratchEThreePr}{equal within 20\%}  % experiments/FIGS/data/v2/index.json :: fine-tuning qg N=1000: R30 of each L188 run over the from-scratch mean, SD over runs and standard error of that mean in quadrature :: sha256 c6c4138ee5c86423 :: prose twin: two significant figures, direction in words
\newcommand{\MlstFtQgRthirtyResEThree}{52.4\,\ensuremath{\pm}\,3.1}  % experiments/FIGS/data/v2/index.json :: fine-tuning qg N=1000: R30, L162 :: sha256 c6c4138ee5c86423
\newcommand{\MlstFtQgRthirtyResOverScratchEThree}{1.06\,\ensuremath{\pm}\,0.07}  % experiments/FIGS/data/v2/index.json :: fine-tuning qg N=1000: R30 of each L162 run over the from-scratch mean, SD over runs and standard error of that mean in quadrature :: sha256 c6c4138ee5c86423
\newcommand{\MlstFtQgRthirtyResOverScratchEThreePr}{equal within 7\%}  % experiments/FIGS/data/v2/index.json :: fine-tuning qg N=1000: R30 of each L162 run over the from-scratch mean, SD over runs and standard error of that mean in quadrature :: sha256 c6c4138ee5c86423 :: prose twin: two significant figures, direction in words
\newcommand{\MlstFtQgRthirtyMidEThree}{54\,\ensuremath{\pm}\,6}  % experiments/FIGS/data/v2/index.json :: fine-tuning qg N=1000: R30, R42_Q1 :: sha256 c6c4138ee5c86423
\newcommand{\MlstFtQgRthirtyMidOverScratchEThree}{1.09\,\ensuremath{\pm}\,0.12}  % experiments/FIGS/data/v2/index.json :: fine-tuning qg N=1000: R30 of each R42_Q1 run over the from-scratch mean, SD over runs and standard error of that mean in quadrature :: sha256 c6c4138ee5c86423
\newcommand{\MlstFtQgRthirtyMidOverScratchEThreePr}{equal within 20\%}  % experiments/FIGS/data/v2/index.json :: fine-tuning qg N=1000: R30 of each R42_Q1 run over the from-scratch mean, SD over runs and standard error of that mean in quadrature :: sha256 c6c4138ee5c86423 :: prose twin: two significant figures, direction in words
\newcommand{\MlstFtQgRthirtyCoarseEThree}{55\,\ensuremath{\pm}\,5}  % experiments/FIGS/data/v2/index.json :: fine-tuning qg N=1000: R30, R16_Q1 :: sha256 c6c4138ee5c86423
\newcommand{\MlstFtQgRthirtyCoarseOverScratchEThree}{1.12\,\ensuremath{\pm}\,0.10}  % experiments/FIGS/data/v2/index.json :: fine-tuning qg N=1000: R30 of each R16_Q1 run over the from-scratch mean, SD over runs and standard error of that mean in quadrature :: sha256 c6c4138ee5c86423
\newcommand{\MlstFtQgRthirtyCoarseOverScratchEThreePr}{equal within 20\%}  % experiments/FIGS/data/v2/index.json :: fine-tuning qg N=1000: R30 of each R16_Q1 run over the from-scratch mean, SD over runs and standard error of that mean in quadrature :: sha256 c6c4138ee5c86423 :: prose twin: two significant figures, direction in words
\newcommand{\MlstFtQgRthirtyResMassEThree}{59.3 and 57.1}  % experiments/FIGS/data/v2/index.json :: fine-tuning qg N=1000: R30, L162_MASS :: sha256 c6c4138ee5c86423
\newcommand{\MlstFtQgRthirtyResMassOverScratchEThree}{1.193 and 1.149}  % experiments/FIGS/data/v2/index.json :: fine-tuning qg N=1000: R30 of each L162_MASS run over the from-scratch mean, SD over runs and standard error of that mean in quadrature :: sha256 c6c4138ee5c86423
\newcommand{\MlstFtQgRthirtyResMassOverScratchEThreePr}{at least 15\% more rejection}  % experiments/FIGS/data/v2/index.json :: fine-tuning qg N=1000: R30 of each L162_MASS run over the from-scratch mean, SD over runs and standard error of that mean in quadrature :: sha256 c6c4138ee5c86423 :: prose twin: two significant figures, direction in words
\newcommand{\MlstFtQgRthirtyResMassOverScratchEThreePrMax}{up to 19\% more rejection}  % experiments/FIGS/data/v2/index.json :: fine-tuning qg N=1000: R30 of each L162_MASS run over the from-scratch mean, SD over runs and standard error of that mean in quadrature :: sha256 c6c4138ee5c86423 :: prose twin of the run farther from 1
\newcommand{\MlstFtQgRthirtyCoarseMassEThree}{54.2 and 57.7}  % experiments/FIGS/data/v2/index.json :: fine-tuning qg N=1000: R30, R16_Q1_MASS :: sha256 c6c4138ee5c86423
\newcommand{\MlstFtQgRthirtyCoarseMassOverScratchEThree}{1.091 and 1.161}  % experiments/FIGS/data/v2/index.json :: fine-tuning qg N=1000: R30 of each R16_Q1_MASS run over the from-scratch mean, SD over runs and standard error of that mean in quadrature :: sha256 c6c4138ee5c86423
\newcommand{\MlstFtQgRthirtyCoarseMassOverScratchEThreePr}{at least 9.1\% more rejection}  % experiments/FIGS/data/v2/index.json :: fine-tuning qg N=1000: R30 of each R16_Q1_MASS run over the from-scratch mean, SD over runs and standard error of that mean in quadrature :: sha256 c6c4138ee5c86423 :: prose twin: two significant figures, direction in words
\newcommand{\MlstFtQgRthirtyCoarseMassOverScratchEThreePrMax}{up to 16\% more rejection}  % experiments/FIGS/data/v2/index.json :: fine-tuning qg N=1000: R30 of each R16_Q1_MASS run over the from-scratch mean, SD over runs and standard error of that mean in quadrature :: sha256 c6c4138ee5c86423 :: prose twin of the run farther from 1
\newcommand{\MlstFtQgRthirtyResOverFullEThree}{1.14\,\ensuremath{\pm}\,0.17}  % experiments/FIGS/data/v2/index.json :: fine-tuning qg N=1000: R30 of L162 over L188, run by run :: sha256 c6c4138ee5c86423
\newcommand{\MlstFtQgRthirtyResOverFullEThreePr}{equal within 20\%}  % experiments/FIGS/data/v2/index.json :: fine-tuning qg N=1000: R30 of L162 over L188, run by run :: sha256 c6c4138ee5c86423 :: prose twin: two significant figures, direction in words
\newcommand{\MlstFtQgRthirtyMidOverFullEThree}{1.17\,\ensuremath{\pm}\,0.14}  % experiments/FIGS/data/v2/index.json :: fine-tuning qg N=1000: R30 of R42_Q1 over L188, run by run :: sha256 c6c4138ee5c86423
\newcommand{\MlstFtQgRthirtyMidOverFullEThreePr}{equal within 20\%}  % experiments/FIGS/data/v2/index.json :: fine-tuning qg N=1000: R30 of R42_Q1 over L188, run by run :: sha256 c6c4138ee5c86423 :: prose twin: two significant figures, direction in words
\newcommand{\MlstFtQgRthirtyCoarseOverFullEThree}{1.21\,\ensuremath{\pm}\,0.27}  % experiments/FIGS/data/v2/index.json :: fine-tuning qg N=1000: R30 of R16_Q1 over L188, run by run :: sha256 c6c4138ee5c86423
\newcommand{\MlstFtQgRthirtyCoarseOverFullEThreePr}{equal within 30\%}  % experiments/FIGS/data/v2/index.json :: fine-tuning qg N=1000: R30 of R16_Q1 over L188, run by run :: sha256 c6c4138ee5c86423 :: prose twin: two significant figures, direction in words
\newcommand{\MlstFtQgRthirtyResMassOverResEThree}{1.108 and 1.042}  % experiments/FIGS/data/v2/index.json :: fine-tuning qg N=1000: R30 of L162_MASS over L162, run by run :: sha256 c6c4138ee5c86423
\newcommand{\MlstFtQgRthirtyResMassOverResEThreePr}{equal within 20\%}  % experiments/FIGS/data/v2/index.json :: fine-tuning qg N=1000: R30 of L162_MASS over L162, run by run :: sha256 c6c4138ee5c86423 :: prose twin: two significant figures, direction in words
\newcommand{\MlstFtQgRthirtyResMassOverResEThreePrMax}{equal within 20\%}  % experiments/FIGS/data/v2/index.json :: fine-tuning qg N=1000: R30 of L162_MASS over L162, run by run :: sha256 c6c4138ee5c86423 :: prose twin of the run farther from 1
\newcommand{\MlstFtQgRthirtyCoarseMassOverCoarseEThree}{0.89 and 1.11}  % experiments/FIGS/data/v2/index.json :: fine-tuning qg N=1000: R30 of R16_Q1_MASS over R16_Q1, run by run :: sha256 c6c4138ee5c86423
\newcommand{\MlstFtQgRthirtyCoarseMassOverCoarseEThreePr}{equal within 20\%}  % experiments/FIGS/data/v2/index.json :: fine-tuning qg N=1000: R30 of R16_Q1_MASS over R16_Q1, run by run :: sha256 c6c4138ee5c86423 :: prose twin: two significant figures, direction in words
\newcommand{\MlstFtQgRthirtyCoarseMassOverCoarseEThreePrMax}{equal within 20\%}  % experiments/FIGS/data/v2/index.json :: fine-tuning qg N=1000: R30 of R16_Q1_MASS over R16_Q1, run by run :: sha256 c6c4138ee5c86423 :: prose twin of the run farther from 1
\newcommand{\MlstFtQgRthirtyRandOverFullEThree}{0.91\,\ensuremath{\pm}\,0.07}  % experiments/FIGS/data/v2/index.json :: fine-tuning qg N=1000: R30 of each random partition over the L188 run of its seed :: sha256 c6c4138ee5c86423
\newcommand{\MlstFtQgRthirtyRandOverFullEThreePr}{equal within 10\%}  % experiments/FIGS/data/v2/index.json :: fine-tuning qg N=1000: R30 of each random partition over the L188 run of its seed :: sha256 c6c4138ee5c86423 :: prose twin: two significant figures, direction in words
\newcommand{\MlstFtQgRthirtyRandOverCoarseEThree}{0.87\,\ensuremath{\pm}\,0.14}  % experiments/FIGS/data/v2/index.json :: fine-tuning qg N=1000: R30 of each random partition over the R16_Q1 run of its seed :: sha256 c6c4138ee5c86423
\newcommand{\MlstFtQgRthirtyRandOverCoarseEThreePr}{equal within 20\%}  % experiments/FIGS/data/v2/index.json :: fine-tuning qg N=1000: R30 of each random partition over the R16_Q1 run of its seed :: sha256 c6c4138ee5c86423 :: prose twin: two significant figures, direction in words
\newcommand{\MlstFtQgRthirtyScratchEFour}{66\,\ensuremath{\pm}\,9}  % experiments/FIGS/data/v2/index.json :: fine-tuning qg N=10000: R30 from scratch, every seed :: sha256 c6c4138ee5c86423
\newcommand{\MlstFtQgRthirtyFullEFour}{75\,\ensuremath{\pm}\,7}  % experiments/FIGS/data/v2/index.json :: fine-tuning qg N=10000: R30, L188 :: sha256 c6c4138ee5c86423
\newcommand{\MlstFtQgRthirtyFullOverScratchEFour}{1.14\,\ensuremath{\pm}\,0.14}  % experiments/FIGS/data/v2/index.json :: fine-tuning qg N=10000: R30 of each L188 run over the from-scratch mean, SD over runs and standard error of that mean in quadrature :: sha256 c6c4138ee5c86423
\newcommand{\MlstFtQgRthirtyFullOverScratchEFourPr}{equal within 20\%}  % experiments/FIGS/data/v2/index.json :: fine-tuning qg N=10000: R30 of each L188 run over the from-scratch mean, SD over runs and standard error of that mean in quadrature :: sha256 c6c4138ee5c86423 :: prose twin: two significant figures, direction in words
\newcommand{\MlstFtQgRthirtyResEFour}{76.4\,\ensuremath{\pm}\,0.6}  % experiments/FIGS/data/v2/index.json :: fine-tuning qg N=10000: R30, L162 :: sha256 c6c4138ee5c86423
\newcommand{\MlstFtQgRthirtyResOverScratchEFour}{1.15\,\ensuremath{\pm}\,0.09}  % experiments/FIGS/data/v2/index.json :: fine-tuning qg N=10000: R30 of each L162 run over the from-scratch mean, SD over runs and standard error of that mean in quadrature :: sha256 c6c4138ee5c86423
\newcommand{\MlstFtQgRthirtyResOverScratchEFourPr}{15\% more rejection}  % experiments/FIGS/data/v2/index.json :: fine-tuning qg N=10000: R30 of each L162 run over the from-scratch mean, SD over runs and standard error of that mean in quadrature :: sha256 c6c4138ee5c86423 :: prose twin: two significant figures, direction in words
\newcommand{\MlstFtQgRthirtyMidEFour}{74\,\ensuremath{\pm}\,5}  % experiments/FIGS/data/v2/index.json :: fine-tuning qg N=10000: R30, R42_Q1 :: sha256 c6c4138ee5c86423
\newcommand{\MlstFtQgRthirtyMidOverScratchEFour}{1.12\,\ensuremath{\pm}\,0.11}  % experiments/FIGS/data/v2/index.json :: fine-tuning qg N=10000: R30 of each R42_Q1 run over the from-scratch mean, SD over runs and standard error of that mean in quadrature :: sha256 c6c4138ee5c86423
\newcommand{\MlstFtQgRthirtyMidOverScratchEFourPr}{equal within 20\%}  % experiments/FIGS/data/v2/index.json :: fine-tuning qg N=10000: R30 of each R42_Q1 run over the from-scratch mean, SD over runs and standard error of that mean in quadrature :: sha256 c6c4138ee5c86423 :: prose twin: two significant figures, direction in words
\newcommand{\MlstFtQgRthirtyCoarseEFour}{80.3\,\ensuremath{\pm}\,1.1}  % experiments/FIGS/data/v2/index.json :: fine-tuning qg N=10000: R30, R16_Q1 :: sha256 c6c4138ee5c86423
\newcommand{\MlstFtQgRthirtyCoarseOverScratchEFour}{1.21\,\ensuremath{\pm}\,0.10}  % experiments/FIGS/data/v2/index.json :: fine-tuning qg N=10000: R30 of each R16_Q1 run over the from-scratch mean, SD over runs and standard error of that mean in quadrature :: sha256 c6c4138ee5c86423
\newcommand{\MlstFtQgRthirtyCoarseOverScratchEFourPr}{21\% more rejection}  % experiments/FIGS/data/v2/index.json :: fine-tuning qg N=10000: R30 of each R16_Q1 run over the from-scratch mean, SD over runs and standard error of that mean in quadrature :: sha256 c6c4138ee5c86423 :: prose twin: two significant figures, direction in words
\newcommand{\MlstFtQgRthirtyResMassEFour}{80.1 and 72.6}  % experiments/FIGS/data/v2/index.json :: fine-tuning qg N=10000: R30, L162_MASS :: sha256 c6c4138ee5c86423
\newcommand{\MlstFtQgRthirtyResMassOverScratchEFour}{1.21 and 1.09}  % experiments/FIGS/data/v2/index.json :: fine-tuning qg N=10000: R30 of each L162_MASS run over the from-scratch mean, SD over runs and standard error of that mean in quadrature :: sha256 c6c4138ee5c86423
\newcommand{\MlstFtQgRthirtyResMassOverScratchEFourPr}{at least 9.3\% more rejection}  % experiments/FIGS/data/v2/index.json :: fine-tuning qg N=10000: R30 of each L162_MASS run over the from-scratch mean, SD over runs and standard error of that mean in quadrature :: sha256 c6c4138ee5c86423 :: prose twin: two significant figures, direction in words
\newcommand{\MlstFtQgRthirtyResMassOverScratchEFourPrMax}{up to 21\% more rejection}  % experiments/FIGS/data/v2/index.json :: fine-tuning qg N=10000: R30 of each L162_MASS run over the from-scratch mean, SD over runs and standard error of that mean in quadrature :: sha256 c6c4138ee5c86423 :: prose twin of the run farther from 1
\newcommand{\MlstFtQgRthirtyCoarseMassEFour}{79.5 and 77.9}  % experiments/FIGS/data/v2/index.json :: fine-tuning qg N=10000: R30, R16_Q1_MASS :: sha256 c6c4138ee5c86423
\newcommand{\MlstFtQgRthirtyCoarseMassOverScratchEFour}{1.198 and 1.173}  % experiments/FIGS/data/v2/index.json :: fine-tuning qg N=10000: R30 of each R16_Q1_MASS run over the from-scratch mean, SD over runs and standard error of that mean in quadrature :: sha256 c6c4138ee5c86423
\newcommand{\MlstFtQgRthirtyCoarseMassOverScratchEFourPr}{at least 17\% more rejection}  % experiments/FIGS/data/v2/index.json :: fine-tuning qg N=10000: R30 of each R16_Q1_MASS run over the from-scratch mean, SD over runs and standard error of that mean in quadrature :: sha256 c6c4138ee5c86423 :: prose twin: two significant figures, direction in words
\newcommand{\MlstFtQgRthirtyCoarseMassOverScratchEFourPrMax}{up to 20\% more rejection}  % experiments/FIGS/data/v2/index.json :: fine-tuning qg N=10000: R30 of each R16_Q1_MASS run over the from-scratch mean, SD over runs and standard error of that mean in quadrature :: sha256 c6c4138ee5c86423 :: prose twin of the run farther from 1
\newcommand{\MlstFtQgRthirtyResOverFullEFour}{1.02\,\ensuremath{\pm}\,0.10}  % experiments/FIGS/data/v2/index.json :: fine-tuning qg N=10000: R30 of L162 over L188, run by run :: sha256 c6c4138ee5c86423
\newcommand{\MlstFtQgRthirtyResOverFullEFourPr}{equal within 10\%}  % experiments/FIGS/data/v2/index.json :: fine-tuning qg N=10000: R30 of L162 over L188, run by run :: sha256 c6c4138ee5c86423 :: prose twin: two significant figures, direction in words
\newcommand{\MlstFtQgRthirtyMidOverFullEFour}{0.99\,\ensuremath{\pm}\,0.14}  % experiments/FIGS/data/v2/index.json :: fine-tuning qg N=10000: R30 of R42_Q1 over L188, run by run :: sha256 c6c4138ee5c86423
\newcommand{\MlstFtQgRthirtyMidOverFullEFourPr}{equal within 20\%}  % experiments/FIGS/data/v2/index.json :: fine-tuning qg N=10000: R30 of R42_Q1 over L188, run by run :: sha256 c6c4138ee5c86423 :: prose twin: two significant figures, direction in words
\newcommand{\MlstFtQgRthirtyCoarseOverFullEFour}{1.07\,\ensuremath{\pm}\,0.13}  % experiments/FIGS/data/v2/index.json :: fine-tuning qg N=10000: R30 of R16_Q1 over L188, run by run :: sha256 c6c4138ee5c86423
\newcommand{\MlstFtQgRthirtyCoarseOverFullEFourPr}{equal within 20\%}  % experiments/FIGS/data/v2/index.json :: fine-tuning qg N=10000: R30 of R16_Q1 over L188, run by run :: sha256 c6c4138ee5c86423 :: prose twin: two significant figures, direction in words
\newcommand{\MlstFtQgRthirtyResMassOverResEFour}{1.06 and 0.95}  % experiments/FIGS/data/v2/index.json :: fine-tuning qg N=10000: R30 of L162_MASS over L162, run by run :: sha256 c6c4138ee5c86423
\newcommand{\MlstFtQgRthirtyResMassOverResEFourPr}{equal within 6\%}  % experiments/FIGS/data/v2/index.json :: fine-tuning qg N=10000: R30 of L162_MASS over L162, run by run :: sha256 c6c4138ee5c86423 :: prose twin: two significant figures, direction in words
\newcommand{\MlstFtQgRthirtyResMassOverResEFourPrMax}{equal within 6\%}  % experiments/FIGS/data/v2/index.json :: fine-tuning qg N=10000: R30 of L162_MASS over L162, run by run :: sha256 c6c4138ee5c86423 :: prose twin of the run farther from 1
\newcommand{\MlstFtQgRthirtyCoarseMassOverCoarseEFour}{0.9769 and 0.9696}  % experiments/FIGS/data/v2/index.json :: fine-tuning qg N=10000: R30 of R16_Q1_MASS over R16_Q1, run by run :: sha256 c6c4138ee5c86423
\newcommand{\MlstFtQgRthirtyCoarseMassOverCoarseEFourPr}{equal within 4\%}  % experiments/FIGS/data/v2/index.json :: fine-tuning qg N=10000: R30 of R16_Q1_MASS over R16_Q1, run by run :: sha256 c6c4138ee5c86423 :: prose twin: two significant figures, direction in words
\newcommand{\MlstFtQgRthirtyCoarseMassOverCoarseEFourPrMax}{equal within 4\%}  % experiments/FIGS/data/v2/index.json :: fine-tuning qg N=10000: R30 of R16_Q1_MASS over R16_Q1, run by run :: sha256 c6c4138ee5c86423 :: prose twin of the run farther from 1
\newcommand{\MlstFtQgRthirtyRandOverFullEFour}{1.03\,\ensuremath{\pm}\,0.18}  % experiments/FIGS/data/v2/index.json :: fine-tuning qg N=10000: R30 of each random partition over the L188 run of its seed :: sha256 c6c4138ee5c86423
\newcommand{\MlstFtQgRthirtyRandOverFullEFourPr}{equal within 20\%}  % experiments/FIGS/data/v2/index.json :: fine-tuning qg N=10000: R30 of each random partition over the L188 run of its seed :: sha256 c6c4138ee5c86423 :: prose twin: two significant figures, direction in words
\newcommand{\MlstFtQgRthirtyRandOverCoarseEFour}{0.93\,\ensuremath{\pm}\,0.12}  % experiments/FIGS/data/v2/index.json :: fine-tuning qg N=10000: R30 of each random partition over the R16_Q1 run of its seed :: sha256 c6c4138ee5c86423
\newcommand{\MlstFtQgRthirtyRandOverCoarseEFourPr}{equal within 20\%}  % experiments/FIGS/data/v2/index.json :: fine-tuning qg N=10000: R30 of each random partition over the R16_Q1 run of its seed :: sha256 c6c4138ee5c86423 :: prose twin: two significant figures, direction in words
\newcommand{\MlstFtQgRthirtyScratchEFive}{88.4\,\ensuremath{\pm}\,0.8}  % experiments/FIGS/data/v2/index.json :: fine-tuning qg N=100000: R30 from scratch, every seed :: sha256 c6c4138ee5c86423
\newcommand{\MlstFtQgRthirtyFullEFive}{99.3\,\ensuremath{\pm}\,1.4}  % experiments/FIGS/data/v2/index.json :: fine-tuning qg N=100000: R30, L188 :: sha256 c6c4138ee5c86423
\newcommand{\MlstFtQgRthirtyFullOverScratchEFive}{1.124\,\ensuremath{\pm}\,0.017}  % experiments/FIGS/data/v2/index.json :: fine-tuning qg N=100000: R30 of each L188 run over the from-scratch mean, SD over runs and standard error of that mean in quadrature :: sha256 c6c4138ee5c86423
\newcommand{\MlstFtQgRthirtyFullOverScratchEFivePr}{12\% more rejection}  % experiments/FIGS/data/v2/index.json :: fine-tuning qg N=100000: R30 of each L188 run over the from-scratch mean, SD over runs and standard error of that mean in quadrature :: sha256 c6c4138ee5c86423 :: prose twin: two significant figures, direction in words
\newcommand{\MlstFtQgRthirtyResEFive}{97.9\,\ensuremath{\pm}\,3.3}  % experiments/FIGS/data/v2/index.json :: fine-tuning qg N=100000: R30, L162 :: sha256 c6c4138ee5c86423
\newcommand{\MlstFtQgRthirtyResOverScratchEFive}{1.11\,\ensuremath{\pm}\,0.04}  % experiments/FIGS/data/v2/index.json :: fine-tuning qg N=100000: R30 of each L162 run over the from-scratch mean, SD over runs and standard error of that mean in quadrature :: sha256 c6c4138ee5c86423
\newcommand{\MlstFtQgRthirtyResOverScratchEFivePr}{11\% more rejection}  % experiments/FIGS/data/v2/index.json :: fine-tuning qg N=100000: R30 of each L162 run over the from-scratch mean, SD over runs and standard error of that mean in quadrature :: sha256 c6c4138ee5c86423 :: prose twin: two significant figures, direction in words
\newcommand{\MlstFtQgRthirtyMidEFive}{96.1\,\ensuremath{\pm}\,1.8}  % experiments/FIGS/data/v2/index.json :: fine-tuning qg N=100000: R30, R42_Q1 :: sha256 c6c4138ee5c86423
\newcommand{\MlstFtQgRthirtyMidOverScratchEFive}{1.088\,\ensuremath{\pm}\,0.021}  % experiments/FIGS/data/v2/index.json :: fine-tuning qg N=100000: R30 of each R42_Q1 run over the from-scratch mean, SD over runs and standard error of that mean in quadrature :: sha256 c6c4138ee5c86423
\newcommand{\MlstFtQgRthirtyMidOverScratchEFivePr}{8.8\% more rejection}  % experiments/FIGS/data/v2/index.json :: fine-tuning qg N=100000: R30 of each R42_Q1 run over the from-scratch mean, SD over runs and standard error of that mean in quadrature :: sha256 c6c4138ee5c86423 :: prose twin: two significant figures, direction in words
\newcommand{\MlstFtQgRthirtyCoarseEFive}{96.4\,\ensuremath{\pm}\,1.0}  % experiments/FIGS/data/v2/index.json :: fine-tuning qg N=100000: R30, R16_Q1 :: sha256 c6c4138ee5c86423
\newcommand{\MlstFtQgRthirtyCoarseOverScratchEFive}{1.091\,\ensuremath{\pm}\,0.013}  % experiments/FIGS/data/v2/index.json :: fine-tuning qg N=100000: R30 of each R16_Q1 run over the from-scratch mean, SD over runs and standard error of that mean in quadrature :: sha256 c6c4138ee5c86423
\newcommand{\MlstFtQgRthirtyCoarseOverScratchEFivePr}{9.1\% more rejection}  % experiments/FIGS/data/v2/index.json :: fine-tuning qg N=100000: R30 of each R16_Q1 run over the from-scratch mean, SD over runs and standard error of that mean in quadrature :: sha256 c6c4138ee5c86423 :: prose twin: two significant figures, direction in words
\newcommand{\MlstFtQgRthirtyResMassEFive}{95.6 and 97.6}  % experiments/FIGS/data/v2/index.json :: fine-tuning qg N=100000: R30, L162_MASS :: sha256 c6c4138ee5c86423
\newcommand{\MlstFtQgRthirtyResMassOverScratchEFive}{1.082 and 1.104}  % experiments/FIGS/data/v2/index.json :: fine-tuning qg N=100000: R30 of each L162_MASS run over the from-scratch mean, SD over runs and standard error of that mean in quadrature :: sha256 c6c4138ee5c86423
\newcommand{\MlstFtQgRthirtyResMassOverScratchEFivePr}{at least 8.2\% more rejection}  % experiments/FIGS/data/v2/index.json :: fine-tuning qg N=100000: R30 of each L162_MASS run over the from-scratch mean, SD over runs and standard error of that mean in quadrature :: sha256 c6c4138ee5c86423 :: prose twin: two significant figures, direction in words
\newcommand{\MlstFtQgRthirtyResMassOverScratchEFivePrMax}{up to 10\% more rejection}  % experiments/FIGS/data/v2/index.json :: fine-tuning qg N=100000: R30 of each L162_MASS run over the from-scratch mean, SD over runs and standard error of that mean in quadrature :: sha256 c6c4138ee5c86423 :: prose twin of the run farther from 1
\newcommand{\MlstFtQgRthirtyCoarseMassEFive}{97.85 and 97.94}  % experiments/FIGS/data/v2/index.json :: fine-tuning qg N=100000: R30, R16_Q1_MASS :: sha256 c6c4138ee5c86423
\newcommand{\MlstFtQgRthirtyCoarseMassOverScratchEFive}{1.107 and 1.108}  % experiments/FIGS/data/v2/index.json :: fine-tuning qg N=100000: R30 of each R16_Q1_MASS run over the from-scratch mean, SD over runs and standard error of that mean in quadrature :: sha256 c6c4138ee5c86423
\newcommand{\MlstFtQgRthirtyCoarseMassOverScratchEFivePr}{at least 11\% more rejection}  % experiments/FIGS/data/v2/index.json :: fine-tuning qg N=100000: R30 of each R16_Q1_MASS run over the from-scratch mean, SD over runs and standard error of that mean in quadrature :: sha256 c6c4138ee5c86423 :: prose twin: two significant figures, direction in words
\newcommand{\MlstFtQgRthirtyCoarseMassOverScratchEFivePrMax}{up to 11\% more rejection}  % experiments/FIGS/data/v2/index.json :: fine-tuning qg N=100000: R30 of each R16_Q1_MASS run over the from-scratch mean, SD over runs and standard error of that mean in quadrature :: sha256 c6c4138ee5c86423 :: prose twin of the run farther from 1
\newcommand{\MlstFtQgRthirtyResOverFullEFive}{0.985\,\ensuremath{\pm}\,0.020}  % experiments/FIGS/data/v2/index.json :: fine-tuning qg N=100000: R30 of L162 over L188, run by run :: sha256 c6c4138ee5c86423
\newcommand{\MlstFtQgRthirtyResOverFullEFivePr}{equal within 2\%}  % experiments/FIGS/data/v2/index.json :: fine-tuning qg N=100000: R30 of L162 over L188, run by run :: sha256 c6c4138ee5c86423 :: prose twin: two significant figures, direction in words
\newcommand{\MlstFtQgRthirtyMidOverFullEFive}{0.969\,\ensuremath{\pm}\,0.030}  % experiments/FIGS/data/v2/index.json :: fine-tuning qg N=100000: R30 of R42_Q1 over L188, run by run :: sha256 c6c4138ee5c86423
\newcommand{\MlstFtQgRthirtyMidOverFullEFivePr}{equal within 4\%}  % experiments/FIGS/data/v2/index.json :: fine-tuning qg N=100000: R30 of R42_Q1 over L188, run by run :: sha256 c6c4138ee5c86423 :: prose twin: two significant figures, direction in words
\newcommand{\MlstFtQgRthirtyCoarseOverFullEFive}{0.972\,\ensuremath{\pm}\,0.020}  % experiments/FIGS/data/v2/index.json :: fine-tuning qg N=100000: R30 of R16_Q1 over L188, run by run :: sha256 c6c4138ee5c86423
\newcommand{\MlstFtQgRthirtyCoarseOverFullEFivePr}{equal within 3\%}  % experiments/FIGS/data/v2/index.json :: fine-tuning qg N=100000: R30 of R16_Q1 over L188, run by run :: sha256 c6c4138ee5c86423 :: prose twin: two significant figures, direction in words
\newcommand{\MlstFtQgRthirtyResMassOverResEFive}{0.942 and 1.027}  % experiments/FIGS/data/v2/index.json :: fine-tuning qg N=100000: R30 of L162_MASS over L162, run by run :: sha256 c6c4138ee5c86423
\newcommand{\MlstFtQgRthirtyResMassOverResEFivePr}{equal within 6\%}  % experiments/FIGS/data/v2/index.json :: fine-tuning qg N=100000: R30 of L162_MASS over L162, run by run :: sha256 c6c4138ee5c86423 :: prose twin: two significant figures, direction in words
\newcommand{\MlstFtQgRthirtyResMassOverResEFivePrMax}{equal within 6\%}  % experiments/FIGS/data/v2/index.json :: fine-tuning qg N=100000: R30 of L162_MASS over L162, run by run :: sha256 c6c4138ee5c86423 :: prose twin of the run farther from 1
\newcommand{\MlstFtQgRthirtyCoarseMassOverCoarseEFive}{1.018 and 1.004}  % experiments/FIGS/data/v2/index.json :: fine-tuning qg N=100000: R30 of R16_Q1_MASS over R16_Q1, run by run :: sha256 c6c4138ee5c86423
\newcommand{\MlstFtQgRthirtyCoarseMassOverCoarseEFivePr}{equal within 2\%}  % experiments/FIGS/data/v2/index.json :: fine-tuning qg N=100000: R30 of R16_Q1_MASS over R16_Q1, run by run :: sha256 c6c4138ee5c86423 :: prose twin: two significant figures, direction in words
\newcommand{\MlstFtQgRthirtyCoarseMassOverCoarseEFivePrMax}{equal within 2\%}  % experiments/FIGS/data/v2/index.json :: fine-tuning qg N=100000: R30 of R16_Q1_MASS over R16_Q1, run by run :: sha256 c6c4138ee5c86423 :: prose twin of the run farther from 1
\newcommand{\MlstFtQgRthirtyRandOverFullEFive}{0.986\,\ensuremath{\pm}\,0.025}  % experiments/FIGS/data/v2/index.json :: fine-tuning qg N=100000: R30 of each random partition over the L188 run of its seed :: sha256 c6c4138ee5c86423
\newcommand{\MlstFtQgRthirtyRandOverFullEFivePr}{equal within 3\%}  % experiments/FIGS/data/v2/index.json :: fine-tuning qg N=100000: R30 of each random partition over the L188 run of its seed :: sha256 c6c4138ee5c86423 :: prose twin: two significant figures, direction in words
\newcommand{\MlstFtQgRthirtyRandOverCoarseEFive}{1.009\,\ensuremath{\pm}\,0.020}  % experiments/FIGS/data/v2/index.json :: fine-tuning qg N=100000: R30 of each random partition over the R16_Q1 run of its seed :: sha256 c6c4138ee5c86423
\newcommand{\MlstFtQgRthirtyRandOverCoarseEFivePr}{equal within 3\%}  % experiments/FIGS/data/v2/index.json :: fine-tuning qg N=100000: R30 of each random partition over the R16_Q1 run of its seed :: sha256 c6c4138ee5c86423 :: prose twin: two significant figures, direction in words
\newcommand{\MlstFtQgRthirtyScratch}{104.0\,\ensuremath{\pm}\,1.0}  % experiments/FIGS/data/v2/index.json :: fine-tuning qg N=1600000: R30 from scratch, every seed :: sha256 c6c4138ee5c86423
\newcommand{\MlstFtQgRthirtyFull}{107.3\,\ensuremath{\pm}\,2.2}  % experiments/FIGS/data/v2/index.json :: fine-tuning qg N=1600000: R30, L188 :: sha256 c6c4138ee5c86423
\newcommand{\MlstFtQgRthirtyFullOverScratch}{1.031\,\ensuremath{\pm}\,0.022}  % experiments/FIGS/data/v2/index.json :: fine-tuning qg N=1600000: R30 of each L188 run over the from-scratch mean, SD over runs and standard error of that mean in quadrature :: sha256 c6c4138ee5c86423
\newcommand{\MlstFtQgRthirtyFullOverScratchPr}{equal within 4\%}  % experiments/FIGS/data/v2/index.json :: fine-tuning qg N=1600000: R30 of each L188 run over the from-scratch mean, SD over runs and standard error of that mean in quadrature :: sha256 c6c4138ee5c86423 :: prose twin: two significant figures, direction in words
\newcommand{\MlstFtQgRthirtyRes}{108.0\,\ensuremath{\pm}\,1.4}  % experiments/FIGS/data/v2/index.json :: fine-tuning qg N=1600000: R30, L162 :: sha256 c6c4138ee5c86423
\newcommand{\MlstFtQgRthirtyResOverScratch}{1.038\,\ensuremath{\pm}\,0.015}  % experiments/FIGS/data/v2/index.json :: fine-tuning qg N=1600000: R30 of each L162 run over the from-scratch mean, SD over runs and standard error of that mean in quadrature :: sha256 c6c4138ee5c86423
\newcommand{\MlstFtQgRthirtyResOverScratchPr}{3.8\% more rejection}  % experiments/FIGS/data/v2/index.json :: fine-tuning qg N=1600000: R30 of each L162 run over the from-scratch mean, SD over runs and standard error of that mean in quadrature :: sha256 c6c4138ee5c86423 :: prose twin: two significant figures, direction in words
\newcommand{\MlstFtQgRthirtyMid}{108.3\,\ensuremath{\pm}\,1.0}  % experiments/FIGS/data/v2/index.json :: fine-tuning qg N=1600000: R30, R42_Q1 :: sha256 c6c4138ee5c86423
\newcommand{\MlstFtQgRthirtyMidOverScratch}{1.041\,\ensuremath{\pm}\,0.012}  % experiments/FIGS/data/v2/index.json :: fine-tuning qg N=1600000: R30 of each R42_Q1 run over the from-scratch mean, SD over runs and standard error of that mean in quadrature :: sha256 c6c4138ee5c86423
\newcommand{\MlstFtQgRthirtyMidOverScratchPr}{4.1\% more rejection}  % experiments/FIGS/data/v2/index.json :: fine-tuning qg N=1600000: R30 of each R42_Q1 run over the from-scratch mean, SD over runs and standard error of that mean in quadrature :: sha256 c6c4138ee5c86423 :: prose twin: two significant figures, direction in words
\newcommand{\MlstFtQgRthirtyCoarse}{106.5\,\ensuremath{\pm}\,0.8}  % experiments/FIGS/data/v2/index.json :: fine-tuning qg N=1600000: R30, R16_Q1 :: sha256 c6c4138ee5c86423
\newcommand{\MlstFtQgRthirtyCoarseOverScratch}{1.024\,\ensuremath{\pm}\,0.010}  % experiments/FIGS/data/v2/index.json :: fine-tuning qg N=1600000: R30 of each R16_Q1 run over the from-scratch mean, SD over runs and standard error of that mean in quadrature :: sha256 c6c4138ee5c86423
\newcommand{\MlstFtQgRthirtyCoarseOverScratchPr}{2.4\% more rejection}  % experiments/FIGS/data/v2/index.json :: fine-tuning qg N=1600000: R30 of each R16_Q1 run over the from-scratch mean, SD over runs and standard error of that mean in quadrature :: sha256 c6c4138ee5c86423 :: prose twin: two significant figures, direction in words
\newcommand{\MlstFtQgRthirtyResMass}{107.0 and 106.6}  % experiments/FIGS/data/v2/index.json :: fine-tuning qg N=1600000: R30, L162_MASS :: sha256 c6c4138ee5c86423
\newcommand{\MlstFtQgRthirtyResMassOverScratch}{1.028 and 1.025}  % experiments/FIGS/data/v2/index.json :: fine-tuning qg N=1600000: R30 of each L162_MASS run over the from-scratch mean, SD over runs and standard error of that mean in quadrature :: sha256 c6c4138ee5c86423
\newcommand{\MlstFtQgRthirtyResMassOverScratchPr}{at least 2.5\% more rejection}  % experiments/FIGS/data/v2/index.json :: fine-tuning qg N=1600000: R30 of each L162_MASS run over the from-scratch mean, SD over runs and standard error of that mean in quadrature :: sha256 c6c4138ee5c86423 :: prose twin: two significant figures, direction in words
\newcommand{\MlstFtQgRthirtyResMassOverScratchPrMax}{up to 2.8\% more rejection}  % experiments/FIGS/data/v2/index.json :: fine-tuning qg N=1600000: R30 of each L162_MASS run over the from-scratch mean, SD over runs and standard error of that mean in quadrature :: sha256 c6c4138ee5c86423 :: prose twin of the run farther from 1
\newcommand{\MlstFtQgRthirtyCoarseMass}{104.9 and 105.6}  % experiments/FIGS/data/v2/index.json :: fine-tuning qg N=1600000: R30, R16_Q1_MASS :: sha256 c6c4138ee5c86423
\newcommand{\MlstFtQgRthirtyCoarseMassOverScratch}{1.009 and 1.015}  % experiments/FIGS/data/v2/index.json :: fine-tuning qg N=1600000: R30 of each R16_Q1_MASS run over the from-scratch mean, SD over runs and standard error of that mean in quadrature :: sha256 c6c4138ee5c86423
\newcommand{\MlstFtQgRthirtyCoarseMassOverScratchPr}{at least 0.87\% more rejection}  % experiments/FIGS/data/v2/index.json :: fine-tuning qg N=1600000: R30 of each R16_Q1_MASS run over the from-scratch mean, SD over runs and standard error of that mean in quadrature :: sha256 c6c4138ee5c86423 :: prose twin: two significant figures, direction in words
\newcommand{\MlstFtQgRthirtyCoarseMassOverScratchPrMax}{up to 1.5\% more rejection}  % experiments/FIGS/data/v2/index.json :: fine-tuning qg N=1600000: R30 of each R16_Q1_MASS run over the from-scratch mean, SD over runs and standard error of that mean in quadrature :: sha256 c6c4138ee5c86423 :: prose twin of the run farther from 1
\newcommand{\MlstFtQgRthirtyResOverFull}{1.007\,\ensuremath{\pm}\,0.026}  % experiments/FIGS/data/v2/index.json :: fine-tuning qg N=1600000: R30 of L162 over L188, run by run :: sha256 c6c4138ee5c86423
\newcommand{\MlstFtQgRthirtyResOverFullPr}{equal within 3\%}  % experiments/FIGS/data/v2/index.json :: fine-tuning qg N=1600000: R30 of L162 over L188, run by run :: sha256 c6c4138ee5c86423 :: prose twin: two significant figures, direction in words
\newcommand{\MlstFtQgRthirtyMidOverFull}{1.010\,\ensuremath{\pm}\,0.029}  % experiments/FIGS/data/v2/index.json :: fine-tuning qg N=1600000: R30 of R42_Q1 over L188, run by run :: sha256 c6c4138ee5c86423
\newcommand{\MlstFtQgRthirtyMidOverFullPr}{equal within 3\%}  % experiments/FIGS/data/v2/index.json :: fine-tuning qg N=1600000: R30 of R42_Q1 over L188, run by run :: sha256 c6c4138ee5c86423 :: prose twin: two significant figures, direction in words
\newcommand{\MlstFtQgRthirtyCoarseOverFull}{0.993\,\ensuremath{\pm}\,0.026}  % experiments/FIGS/data/v2/index.json :: fine-tuning qg N=1600000: R30 of R16_Q1 over L188, run by run :: sha256 c6c4138ee5c86423
\newcommand{\MlstFtQgRthirtyCoarseOverFullPr}{equal within 3\%}  % experiments/FIGS/data/v2/index.json :: fine-tuning qg N=1600000: R30 of R16_Q1 over L188, run by run :: sha256 c6c4138ee5c86423 :: prose twin: two significant figures, direction in words
\newcommand{\MlstFtQgRthirtyResMassOverRes}{0.976 and 0.999}  % experiments/FIGS/data/v2/index.json :: fine-tuning qg N=1600000: R30 of L162_MASS over L162, run by run :: sha256 c6c4138ee5c86423
\newcommand{\MlstFtQgRthirtyResMassOverResPr}{equal within 3\%}  % experiments/FIGS/data/v2/index.json :: fine-tuning qg N=1600000: R30 of L162_MASS over L162, run by run :: sha256 c6c4138ee5c86423 :: prose twin: two significant figures, direction in words
\newcommand{\MlstFtQgRthirtyResMassOverResPrMax}{equal within 3\%}  % experiments/FIGS/data/v2/index.json :: fine-tuning qg N=1600000: R30 of L162_MASS over L162, run by run :: sha256 c6c4138ee5c86423 :: prose twin of the run farther from 1
\newcommand{\MlstFtQgRthirtyCoarseMassOverCoarse}{0.9916 and 0.9926}  % experiments/FIGS/data/v2/index.json :: fine-tuning qg N=1600000: R30 of R16_Q1_MASS over R16_Q1, run by run :: sha256 c6c4138ee5c86423
\newcommand{\MlstFtQgRthirtyCoarseMassOverCoarsePr}{equal within 1\%}  % experiments/FIGS/data/v2/index.json :: fine-tuning qg N=1600000: R30 of R16_Q1_MASS over R16_Q1, run by run :: sha256 c6c4138ee5c86423 :: prose twin: two significant figures, direction in words
\newcommand{\MlstFtQgRthirtyCoarseMassOverCoarsePrMax}{equal within 1\%}  % experiments/FIGS/data/v2/index.json :: fine-tuning qg N=1600000: R30 of R16_Q1_MASS over R16_Q1, run by run :: sha256 c6c4138ee5c86423 :: prose twin of the run farther from 1
\newcommand{\MlstFtQgRthirtyRandOverFull}{0.985\,\ensuremath{\pm}\,0.011}  % experiments/FIGS/data/v2/index.json :: fine-tuning qg N=1600000: R30 of each random partition over the L188 run of its seed :: sha256 c6c4138ee5c86423
\newcommand{\MlstFtQgRthirtyRandOverFullPr}{equal within 2\%}  % experiments/FIGS/data/v2/index.json :: fine-tuning qg N=1600000: R30 of each random partition over the L188 run of its seed :: sha256 c6c4138ee5c86423 :: prose twin: two significant figures, direction in words
\newcommand{\MlstFtQgRthirtyRandOverCoarse}{1.008\,\ensuremath{\pm}\,0.009}  % experiments/FIGS/data/v2/index.json :: fine-tuning qg N=1600000: R30 of each random partition over the R16_Q1 run of its seed :: sha256 c6c4138ee5c86423
\newcommand{\MlstFtQgRthirtyRandOverCoarsePr}{equal within 1\%}  % experiments/FIGS/data/v2/index.json :: fine-tuning qg N=1600000: R30 of each random partition over the R16_Q1 run of its seed :: sha256 c6c4138ee5c86423 :: prose twin: two significant figures, direction in words
\newcommand{\MlstFtQgRthirtyFullOverPart}{0.994\,\ensuremath{\pm}\,0.020}  % experiments/FIGS/data/v2/index.json :: fine-tuning qg N=1600000: R30 of each 188-class run over ParT-f.t.'s published value :: sha256 c6c4138ee5c86423
\newcommand{\MlstFtQgRthirtyFullOverPartPr}{equal within 3\%}  % experiments/FIGS/data/v2/index.json :: fine-tuning qg N=1600000: R30 of each 188-class run over ParT-f.t.'s published value :: sha256 c6c4138ee5c86423 :: prose twin: two significant figures, direction in words
\newcommand{\MlstFtQgAucFull}{0.9148\,\ensuremath{\pm}\,0.0001}  % experiments/FIGS/data/v2/index.json :: fine-tuning qg N=1600000: AUC of the 188-class runs :: sha256 c6c4138ee5c86423
\newcommand{\MlstFtQgAucScratch}{0.9121\,\ensuremath{\pm}\,0.0002}  % experiments/FIGS/data/v2/index.json :: fine-tuning qg N=1600000: AUC from scratch :: sha256 c6c4138ee5c86423
\newcommand{\MlstFtQgDataEquivJets}{\ensuremath{<}1{,}000}  % experiments/FIGS/data/v2/index.json :: fine-tuning qg: from-scratch training jets whose R30 equals the 188-class mean at 1000 jets (interpolated linearly in R30 against log10 N) :: sha256 c6c4138ee5c86423
\newcommand{\MlstFtQgDataEquivJetsEFour}{26{,}000}  % experiments/FIGS/data/v2/index.json :: fine-tuning qg: from-scratch training jets whose R30 equals the 188-class mean at 10000 jets (interpolated linearly in R30 against log10 N) :: sha256 c6c4138ee5c86423
\newcommand{\MlstFtQgHerwigRthirtyScratchEThree}{27.0\,\ensuremath{\pm}\,1.3}  % experiments/FIGS/data/v2/index.json :: fine-tuning qg_herwig N=1000: R30 from scratch, every seed :: sha256 c6c4138ee5c86423
\newcommand{\MlstFtQgHerwigRthirtyFullEThree}{21\,\ensuremath{\pm}\,6}  % experiments/FIGS/data/v2/index.json :: fine-tuning qg_herwig N=1000: R30, L188 :: sha256 c6c4138ee5c86423
\newcommand{\MlstFtQgHerwigRthirtyFullOverScratchEThree}{0.77\,\ensuremath{\pm}\,0.22}  % experiments/FIGS/data/v2/index.json :: fine-tuning qg_herwig N=1000: R30 of each L188 run over the from-scratch mean, SD over runs and standard error of that mean in quadrature :: sha256 c6c4138ee5c86423
\newcommand{\MlstFtQgHerwigRthirtyFullOverScratchEThreePr}{equal within 30\%}  % experiments/FIGS/data/v2/index.json :: fine-tuning qg_herwig N=1000: R30 of each L188 run over the from-scratch mean, SD over runs and standard error of that mean in quadrature :: sha256 c6c4138ee5c86423 :: prose twin: two significant figures, direction in words
\newcommand{\MlstFtQgHerwigRthirtyResEThree}{26.2\,\ensuremath{\pm}\,0.9}  % experiments/FIGS/data/v2/index.json :: fine-tuning qg_herwig N=1000: R30, L162 :: sha256 c6c4138ee5c86423
\newcommand{\MlstFtQgHerwigRthirtyResOverScratchEThree}{0.97\,\ensuremath{\pm}\,0.04}  % experiments/FIGS/data/v2/index.json :: fine-tuning qg_herwig N=1000: R30 of each L162 run over the from-scratch mean, SD over runs and standard error of that mean in quadrature :: sha256 c6c4138ee5c86423
\newcommand{\MlstFtQgHerwigRthirtyResOverScratchEThreePr}{equal within 4\%}  % experiments/FIGS/data/v2/index.json :: fine-tuning qg_herwig N=1000: R30 of each L162 run over the from-scratch mean, SD over runs and standard error of that mean in quadrature :: sha256 c6c4138ee5c86423 :: prose twin: two significant figures, direction in words
\newcommand{\MlstFtQgHerwigRthirtyMidEThree}{26.9\,\ensuremath{\pm}\,2.7}  % experiments/FIGS/data/v2/index.json :: fine-tuning qg_herwig N=1000: R30, R42_Q1 :: sha256 c6c4138ee5c86423
\newcommand{\MlstFtQgHerwigRthirtyMidOverScratchEThree}{1.00\,\ensuremath{\pm}\,0.10}  % experiments/FIGS/data/v2/index.json :: fine-tuning qg_herwig N=1000: R30 of each R42_Q1 run over the from-scratch mean, SD over runs and standard error of that mean in quadrature :: sha256 c6c4138ee5c86423
\newcommand{\MlstFtQgHerwigRthirtyMidOverScratchEThreePr}{equal within 10\%}  % experiments/FIGS/data/v2/index.json :: fine-tuning qg_herwig N=1000: R30 of each R42_Q1 run over the from-scratch mean, SD over runs and standard error of that mean in quadrature :: sha256 c6c4138ee5c86423 :: prose twin: two significant figures, direction in words
\newcommand{\MlstFtQgHerwigRthirtyCoarseEThree}{25.2\,\ensuremath{\pm}\,2.5}  % experiments/FIGS/data/v2/index.json :: fine-tuning qg_herwig N=1000: R30, R16_Q1 :: sha256 c6c4138ee5c86423
\newcommand{\MlstFtQgHerwigRthirtyCoarseOverScratchEThree}{0.94\,\ensuremath{\pm}\,0.10}  % experiments/FIGS/data/v2/index.json :: fine-tuning qg_herwig N=1000: R30 of each R16_Q1 run over the from-scratch mean, SD over runs and standard error of that mean in quadrature :: sha256 c6c4138ee5c86423
\newcommand{\MlstFtQgHerwigRthirtyCoarseOverScratchEThreePr}{equal within 10\%}  % experiments/FIGS/data/v2/index.json :: fine-tuning qg_herwig N=1000: R30 of each R16_Q1 run over the from-scratch mean, SD over runs and standard error of that mean in quadrature :: sha256 c6c4138ee5c86423 :: prose twin: two significant figures, direction in words
\newcommand{\MlstFtQgHerwigRthirtyResMassEThree}{28.0 and 29.0}  % experiments/FIGS/data/v2/index.json :: fine-tuning qg_herwig N=1000: R30, L162_MASS :: sha256 c6c4138ee5c86423
\newcommand{\MlstFtQgHerwigRthirtyResMassOverScratchEThree}{1.038 and 1.076}  % experiments/FIGS/data/v2/index.json :: fine-tuning qg_herwig N=1000: R30 of each L162_MASS run over the from-scratch mean, SD over runs and standard error of that mean in quadrature :: sha256 c6c4138ee5c86423
\newcommand{\MlstFtQgHerwigRthirtyResMassOverScratchEThreePr}{at least 3.8\% more rejection}  % experiments/FIGS/data/v2/index.json :: fine-tuning qg_herwig N=1000: R30 of each L162_MASS run over the from-scratch mean, SD over runs and standard error of that mean in quadrature :: sha256 c6c4138ee5c86423 :: prose twin: two significant figures, direction in words
\newcommand{\MlstFtQgHerwigRthirtyResMassOverScratchEThreePrMax}{up to 7.6\% more rejection}  % experiments/FIGS/data/v2/index.json :: fine-tuning qg_herwig N=1000: R30 of each L162_MASS run over the from-scratch mean, SD over runs and standard error of that mean in quadrature :: sha256 c6c4138ee5c86423 :: prose twin of the run farther from 1
\newcommand{\MlstFtQgHerwigRthirtyCoarseMassEThree}{25.8 and 27.4}  % experiments/FIGS/data/v2/index.json :: fine-tuning qg_herwig N=1000: R30, R16_Q1_MASS :: sha256 c6c4138ee5c86423
\newcommand{\MlstFtQgHerwigRthirtyCoarseMassOverScratchEThree}{0.955 and 1.015}  % experiments/FIGS/data/v2/index.json :: fine-tuning qg_herwig N=1000: R30 of each R16_Q1_MASS run over the from-scratch mean, SD over runs and standard error of that mean in quadrature :: sha256 c6c4138ee5c86423
\newcommand{\MlstFtQgHerwigRthirtyCoarseMassOverScratchEThreePr}{equal within 5\%}  % experiments/FIGS/data/v2/index.json :: fine-tuning qg_herwig N=1000: R30 of each R16_Q1_MASS run over the from-scratch mean, SD over runs and standard error of that mean in quadrature :: sha256 c6c4138ee5c86423 :: prose twin: two significant figures, direction in words
\newcommand{\MlstFtQgHerwigRthirtyCoarseMassOverScratchEThreePrMax}{equal within 5\%}  % experiments/FIGS/data/v2/index.json :: fine-tuning qg_herwig N=1000: R30 of each R16_Q1_MASS run over the from-scratch mean, SD over runs and standard error of that mean in quadrature :: sha256 c6c4138ee5c86423 :: prose twin of the run farther from 1
\newcommand{\MlstFtQgHerwigRthirtyResOverFullEThree}{1.33\,\ensuremath{\pm}\,0.33}  % experiments/FIGS/data/v2/index.json :: fine-tuning qg_herwig N=1000: R30 of L162 over L188, run by run :: sha256 c6c4138ee5c86423
\newcommand{\MlstFtQgHerwigRthirtyResOverFullEThreePr}{equal within 40\%}  % experiments/FIGS/data/v2/index.json :: fine-tuning qg_herwig N=1000: R30 of L162 over L188, run by run :: sha256 c6c4138ee5c86423 :: prose twin: two significant figures, direction in words
\newcommand{\MlstFtQgHerwigRthirtyMidOverFullEThree}{1.36\,\ensuremath{\pm}\,0.31}  % experiments/FIGS/data/v2/index.json :: fine-tuning qg_herwig N=1000: R30 of R42_Q1 over L188, run by run :: sha256 c6c4138ee5c86423
\newcommand{\MlstFtQgHerwigRthirtyMidOverFullEThreePr}{equal within 40\%}  % experiments/FIGS/data/v2/index.json :: fine-tuning qg_herwig N=1000: R30 of R42_Q1 over L188, run by run :: sha256 c6c4138ee5c86423 :: prose twin: two significant figures, direction in words
\newcommand{\MlstFtQgHerwigRthirtyCoarseOverFullEThree}{1.3\,\ensuremath{\pm}\,0.4}  % experiments/FIGS/data/v2/index.json :: fine-tuning qg_herwig N=1000: R30 of R16_Q1 over L188, run by run :: sha256 c6c4138ee5c86423
\newcommand{\MlstFtQgHerwigRthirtyCoarseOverFullEThreePr}{equal within 50\%}  % experiments/FIGS/data/v2/index.json :: fine-tuning qg_herwig N=1000: R30 of R16_Q1 over L188, run by run :: sha256 c6c4138ee5c86423 :: prose twin: two significant figures, direction in words
\newcommand{\MlstFtQgHerwigRthirtyResMassOverResEThree}{1.062 and 1.073}  % experiments/FIGS/data/v2/index.json :: fine-tuning qg_herwig N=1000: R30 of L162_MASS over L162, run by run :: sha256 c6c4138ee5c86423
\newcommand{\MlstFtQgHerwigRthirtyResMassOverResEThreePr}{equal within 8\%}  % experiments/FIGS/data/v2/index.json :: fine-tuning qg_herwig N=1000: R30 of L162_MASS over L162, run by run :: sha256 c6c4138ee5c86423 :: prose twin: two significant figures, direction in words
\newcommand{\MlstFtQgHerwigRthirtyResMassOverResEThreePrMax}{equal within 8\%}  % experiments/FIGS/data/v2/index.json :: fine-tuning qg_herwig N=1000: R30 of L162_MASS over L162, run by run :: sha256 c6c4138ee5c86423 :: prose twin of the run farther from 1
\newcommand{\MlstFtQgHerwigRthirtyCoarseMassOverCoarseEThree}{0.93 and 1.20}  % experiments/FIGS/data/v2/index.json :: fine-tuning qg_herwig N=1000: R30 of R16_Q1_MASS over R16_Q1, run by run :: sha256 c6c4138ee5c86423
\newcommand{\MlstFtQgHerwigRthirtyCoarseMassOverCoarseEThreePr}{equal within 30\%}  % experiments/FIGS/data/v2/index.json :: fine-tuning qg_herwig N=1000: R30 of R16_Q1_MASS over R16_Q1, run by run :: sha256 c6c4138ee5c86423 :: prose twin: two significant figures, direction in words
\newcommand{\MlstFtQgHerwigRthirtyCoarseMassOverCoarseEThreePrMax}{equal within 30\%}  % experiments/FIGS/data/v2/index.json :: fine-tuning qg_herwig N=1000: R30 of R16_Q1_MASS over R16_Q1, run by run :: sha256 c6c4138ee5c86423 :: prose twin of the run farther from 1
\newcommand{\MlstFtQgHerwigRthirtyRandOverFullEThree}{0.99\,\ensuremath{\pm}\,0.20}  % experiments/FIGS/data/v2/index.json :: fine-tuning qg_herwig N=1000: R30 of each random partition over the L188 run of its seed :: sha256 c6c4138ee5c86423
\newcommand{\MlstFtQgHerwigRthirtyRandOverFullEThreePr}{equal within 30\%}  % experiments/FIGS/data/v2/index.json :: fine-tuning qg_herwig N=1000: R30 of each random partition over the L188 run of its seed :: sha256 c6c4138ee5c86423 :: prose twin: two significant figures, direction in words
\newcommand{\MlstFtQgHerwigRthirtyRandOverCoarseEThree}{0.95\,\ensuremath{\pm}\,0.13}  % experiments/FIGS/data/v2/index.json :: fine-tuning qg_herwig N=1000: R30 of each random partition over the R16_Q1 run of its seed :: sha256 c6c4138ee5c86423
\newcommand{\MlstFtQgHerwigRthirtyRandOverCoarseEThreePr}{equal within 20\%}  % experiments/FIGS/data/v2/index.json :: fine-tuning qg_herwig N=1000: R30 of each random partition over the R16_Q1 run of its seed :: sha256 c6c4138ee5c86423 :: prose twin: two significant figures, direction in words
\newcommand{\MlstFtQgHerwigRthirtyScratchEFour}{31\,\ensuremath{\pm}\,7}  % experiments/FIGS/data/v2/index.json :: fine-tuning qg_herwig N=10000: R30 from scratch, every seed :: sha256 c6c4138ee5c86423
\newcommand{\MlstFtQgHerwigRthirtyFullEFour}{27.9\,\ensuremath{\pm}\,2.8}  % experiments/FIGS/data/v2/index.json :: fine-tuning qg_herwig N=10000: R30, L188 :: sha256 c6c4138ee5c86423
\newcommand{\MlstFtQgHerwigRthirtyFullOverScratchEFour}{0.89\,\ensuremath{\pm}\,0.14}  % experiments/FIGS/data/v2/index.json :: fine-tuning qg_herwig N=10000: R30 of each L188 run over the from-scratch mean, SD over runs and standard error of that mean in quadrature :: sha256 c6c4138ee5c86423
\newcommand{\MlstFtQgHerwigRthirtyFullOverScratchEFourPr}{equal within 20\%}  % experiments/FIGS/data/v2/index.json :: fine-tuning qg_herwig N=10000: R30 of each L188 run over the from-scratch mean, SD over runs and standard error of that mean in quadrature :: sha256 c6c4138ee5c86423 :: prose twin: two significant figures, direction in words
\newcommand{\MlstFtQgHerwigRthirtyResEFour}{27.4\,\ensuremath{\pm}\,0.5}  % experiments/FIGS/data/v2/index.json :: fine-tuning qg_herwig N=10000: R30, L162 :: sha256 c6c4138ee5c86423
\newcommand{\MlstFtQgHerwigRthirtyResOverScratchEFour}{0.87\,\ensuremath{\pm}\,0.11}  % experiments/FIGS/data/v2/index.json :: fine-tuning qg_herwig N=10000: R30 of each L162 run over the from-scratch mean, SD over runs and standard error of that mean in quadrature :: sha256 c6c4138ee5c86423
\newcommand{\MlstFtQgHerwigRthirtyResOverScratchEFourPr}{13\% less rejection}  % experiments/FIGS/data/v2/index.json :: fine-tuning qg_herwig N=10000: R30 of each L162 run over the from-scratch mean, SD over runs and standard error of that mean in quadrature :: sha256 c6c4138ee5c86423 :: prose twin: two significant figures, direction in words
\newcommand{\MlstFtQgHerwigRthirtyMidEFour}{28.4\,\ensuremath{\pm}\,2.1}  % experiments/FIGS/data/v2/index.json :: fine-tuning qg_herwig N=10000: R30, R42_Q1 :: sha256 c6c4138ee5c86423
\newcommand{\MlstFtQgHerwigRthirtyMidOverScratchEFour}{0.91\,\ensuremath{\pm}\,0.13}  % experiments/FIGS/data/v2/index.json :: fine-tuning qg_herwig N=10000: R30 of each R42_Q1 run over the from-scratch mean, SD over runs and standard error of that mean in quadrature :: sha256 c6c4138ee5c86423
\newcommand{\MlstFtQgHerwigRthirtyMidOverScratchEFourPr}{equal within 10\%}  % experiments/FIGS/data/v2/index.json :: fine-tuning qg_herwig N=10000: R30 of each R42_Q1 run over the from-scratch mean, SD over runs and standard error of that mean in quadrature :: sha256 c6c4138ee5c86423 :: prose twin: two significant figures, direction in words
\newcommand{\MlstFtQgHerwigRthirtyCoarseEFour}{32.7\,\ensuremath{\pm}\,1.8}  % experiments/FIGS/data/v2/index.json :: fine-tuning qg_herwig N=10000: R30, R16_Q1 :: sha256 c6c4138ee5c86423
\newcommand{\MlstFtQgHerwigRthirtyCoarseOverScratchEFour}{1.05\,\ensuremath{\pm}\,0.14}  % experiments/FIGS/data/v2/index.json :: fine-tuning qg_herwig N=10000: R30 of each R16_Q1 run over the from-scratch mean, SD over runs and standard error of that mean in quadrature :: sha256 c6c4138ee5c86423
\newcommand{\MlstFtQgHerwigRthirtyCoarseOverScratchEFourPr}{equal within 6\%}  % experiments/FIGS/data/v2/index.json :: fine-tuning qg_herwig N=10000: R30 of each R16_Q1 run over the from-scratch mean, SD over runs and standard error of that mean in quadrature :: sha256 c6c4138ee5c86423 :: prose twin: two significant figures, direction in words
\newcommand{\MlstFtQgHerwigRthirtyResMassEFour}{34.3 and 28.5}  % experiments/FIGS/data/v2/index.json :: fine-tuning qg_herwig N=10000: R30, L162_MASS :: sha256 c6c4138ee5c86423
\newcommand{\MlstFtQgHerwigRthirtyResMassOverScratchEFour}{1.09 and 0.91}  % experiments/FIGS/data/v2/index.json :: fine-tuning qg_herwig N=10000: R30 of each L162_MASS run over the from-scratch mean, SD over runs and standard error of that mean in quadrature :: sha256 c6c4138ee5c86423
\newcommand{\MlstFtQgHerwigRthirtyResMassOverScratchEFourPr}{equal within 10\%}  % experiments/FIGS/data/v2/index.json :: fine-tuning qg_herwig N=10000: R30 of each L162_MASS run over the from-scratch mean, SD over runs and standard error of that mean in quadrature :: sha256 c6c4138ee5c86423 :: prose twin: two significant figures, direction in words
\newcommand{\MlstFtQgHerwigRthirtyResMassOverScratchEFourPrMax}{equal within 10\%}  % experiments/FIGS/data/v2/index.json :: fine-tuning qg_herwig N=10000: R30 of each L162_MASS run over the from-scratch mean, SD over runs and standard error of that mean in quadrature :: sha256 c6c4138ee5c86423 :: prose twin of the run farther from 1
\newcommand{\MlstFtQgHerwigRthirtyCoarseMassEFour}{34.4 and 31.2}  % experiments/FIGS/data/v2/index.json :: fine-tuning qg_herwig N=10000: R30, R16_Q1_MASS :: sha256 c6c4138ee5c86423
\newcommand{\MlstFtQgHerwigRthirtyCoarseMassOverScratchEFour}{1.10 and 1.00}  % experiments/FIGS/data/v2/index.json :: fine-tuning qg_herwig N=10000: R30 of each R16_Q1_MASS run over the from-scratch mean, SD over runs and standard error of that mean in quadrature :: sha256 c6c4138ee5c86423
\newcommand{\MlstFtQgHerwigRthirtyCoarseMassOverScratchEFourPr}{equal within 20\%}  % experiments/FIGS/data/v2/index.json :: fine-tuning qg_herwig N=10000: R30 of each R16_Q1_MASS run over the from-scratch mean, SD over runs and standard error of that mean in quadrature :: sha256 c6c4138ee5c86423 :: prose twin: two significant figures, direction in words
\newcommand{\MlstFtQgHerwigRthirtyCoarseMassOverScratchEFourPrMax}{equal within 20\%}  % experiments/FIGS/data/v2/index.json :: fine-tuning qg_herwig N=10000: R30 of each R16_Q1_MASS run over the from-scratch mean, SD over runs and standard error of that mean in quadrature :: sha256 c6c4138ee5c86423 :: prose twin of the run farther from 1
\newcommand{\MlstFtQgHerwigRthirtyResOverFullEFour}{0.99\,\ensuremath{\pm}\,0.12}  % experiments/FIGS/data/v2/index.json :: fine-tuning qg_herwig N=10000: R30 of L162 over L188, run by run :: sha256 c6c4138ee5c86423
\newcommand{\MlstFtQgHerwigRthirtyResOverFullEFourPr}{equal within 20\%}  % experiments/FIGS/data/v2/index.json :: fine-tuning qg_herwig N=10000: R30 of L162 over L188, run by run :: sha256 c6c4138ee5c86423 :: prose twin: two significant figures, direction in words
\newcommand{\MlstFtQgHerwigRthirtyMidOverFullEFour}{1.03\,\ensuremath{\pm}\,0.18}  % experiments/FIGS/data/v2/index.json :: fine-tuning qg_herwig N=10000: R30 of R42_Q1 over L188, run by run :: sha256 c6c4138ee5c86423
\newcommand{\MlstFtQgHerwigRthirtyMidOverFullEFourPr}{equal within 20\%}  % experiments/FIGS/data/v2/index.json :: fine-tuning qg_herwig N=10000: R30 of R42_Q1 over L188, run by run :: sha256 c6c4138ee5c86423 :: prose twin: two significant figures, direction in words
\newcommand{\MlstFtQgHerwigRthirtyCoarseOverFullEFour}{1.18\,\ensuremath{\pm}\,0.08}  % experiments/FIGS/data/v2/index.json :: fine-tuning qg_herwig N=10000: R30 of R16_Q1 over L188, run by run :: sha256 c6c4138ee5c86423
\newcommand{\MlstFtQgHerwigRthirtyCoarseOverFullEFourPr}{18\% more rejection}  % experiments/FIGS/data/v2/index.json :: fine-tuning qg_herwig N=10000: R30 of R16_Q1 over L188, run by run :: sha256 c6c4138ee5c86423 :: prose twin: two significant figures, direction in words
\newcommand{\MlstFtQgHerwigRthirtyResMassOverResEFour}{1.23 and 1.03}  % experiments/FIGS/data/v2/index.json :: fine-tuning qg_herwig N=10000: R30 of L162_MASS over L162, run by run :: sha256 c6c4138ee5c86423
\newcommand{\MlstFtQgHerwigRthirtyResMassOverResEFourPr}{equal within 30\%}  % experiments/FIGS/data/v2/index.json :: fine-tuning qg_herwig N=10000: R30 of L162_MASS over L162, run by run :: sha256 c6c4138ee5c86423 :: prose twin: two significant figures, direction in words
\newcommand{\MlstFtQgHerwigRthirtyResMassOverResEFourPrMax}{equal within 30\%}  % experiments/FIGS/data/v2/index.json :: fine-tuning qg_herwig N=10000: R30 of L162_MASS over L162, run by run :: sha256 c6c4138ee5c86423 :: prose twin of the run farther from 1
\newcommand{\MlstFtQgHerwigRthirtyCoarseMassOverCoarseEFour}{1.12 and 0.91}  % experiments/FIGS/data/v2/index.json :: fine-tuning qg_herwig N=10000: R30 of R16_Q1_MASS over R16_Q1, run by run :: sha256 c6c4138ee5c86423
\newcommand{\MlstFtQgHerwigRthirtyCoarseMassOverCoarseEFourPr}{equal within 20\%}  % experiments/FIGS/data/v2/index.json :: fine-tuning qg_herwig N=10000: R30 of R16_Q1_MASS over R16_Q1, run by run :: sha256 c6c4138ee5c86423 :: prose twin: two significant figures, direction in words
\newcommand{\MlstFtQgHerwigRthirtyCoarseMassOverCoarseEFourPrMax}{equal within 20\%}  % experiments/FIGS/data/v2/index.json :: fine-tuning qg_herwig N=10000: R30 of R16_Q1_MASS over R16_Q1, run by run :: sha256 c6c4138ee5c86423 :: prose twin of the run farther from 1
\newcommand{\MlstFtQgHerwigRthirtyRandOverFullEFour}{1.13\,\ensuremath{\pm}\,0.18}  % experiments/FIGS/data/v2/index.json :: fine-tuning qg_herwig N=10000: R30 of each random partition over the L188 run of its seed :: sha256 c6c4138ee5c86423
\newcommand{\MlstFtQgHerwigRthirtyRandOverFullEFourPr}{equal within 20\%}  % experiments/FIGS/data/v2/index.json :: fine-tuning qg_herwig N=10000: R30 of each random partition over the L188 run of its seed :: sha256 c6c4138ee5c86423 :: prose twin: two significant figures, direction in words
\newcommand{\MlstFtQgHerwigRthirtyRandOverCoarseEFour}{0.93\,\ensuremath{\pm}\,0.14}  % experiments/FIGS/data/v2/index.json :: fine-tuning qg_herwig N=10000: R30 of each random partition over the R16_Q1 run of its seed :: sha256 c6c4138ee5c86423
\newcommand{\MlstFtQgHerwigRthirtyRandOverCoarseEFourPr}{equal within 20\%}  % experiments/FIGS/data/v2/index.json :: fine-tuning qg_herwig N=10000: R30 of each random partition over the R16_Q1 run of its seed :: sha256 c6c4138ee5c86423 :: prose twin: two significant figures, direction in words
\newcommand{\MlstFtQgHerwigRthirtyScratchEFive}{36.7\,\ensuremath{\pm}\,0.9}  % experiments/FIGS/data/v2/index.json :: fine-tuning qg_herwig N=100000: R30 from scratch, every seed :: sha256 c6c4138ee5c86423
\newcommand{\MlstFtQgHerwigRthirtyFullEFive}{34.3\,\ensuremath{\pm}\,1.0}  % experiments/FIGS/data/v2/index.json :: fine-tuning qg_herwig N=100000: R30, L188 :: sha256 c6c4138ee5c86423
\newcommand{\MlstFtQgHerwigRthirtyFullOverScratchEFive}{0.935\,\ensuremath{\pm}\,0.031}  % experiments/FIGS/data/v2/index.json :: fine-tuning qg_herwig N=100000: R30 of each L188 run over the from-scratch mean, SD over runs and standard error of that mean in quadrature :: sha256 c6c4138ee5c86423
\newcommand{\MlstFtQgHerwigRthirtyFullOverScratchEFivePr}{6.5\% less rejection}  % experiments/FIGS/data/v2/index.json :: fine-tuning qg_herwig N=100000: R30 of each L188 run over the from-scratch mean, SD over runs and standard error of that mean in quadrature :: sha256 c6c4138ee5c86423 :: prose twin: two significant figures, direction in words
\newcommand{\MlstFtQgHerwigRthirtyResEFive}{35.5\,\ensuremath{\pm}\,0.5}  % experiments/FIGS/data/v2/index.json :: fine-tuning qg_herwig N=100000: R30, L162 :: sha256 c6c4138ee5c86423
\newcommand{\MlstFtQgHerwigRthirtyResOverScratchEFive}{0.966\,\ensuremath{\pm}\,0.019}  % experiments/FIGS/data/v2/index.json :: fine-tuning qg_herwig N=100000: R30 of each L162 run over the from-scratch mean, SD over runs and standard error of that mean in quadrature :: sha256 c6c4138ee5c86423
\newcommand{\MlstFtQgHerwigRthirtyResOverScratchEFivePr}{3.4\% less rejection}  % experiments/FIGS/data/v2/index.json :: fine-tuning qg_herwig N=100000: R30 of each L162 run over the from-scratch mean, SD over runs and standard error of that mean in quadrature :: sha256 c6c4138ee5c86423 :: prose twin: two significant figures, direction in words
\newcommand{\MlstFtQgHerwigRthirtyMidEFive}{35.5\,\ensuremath{\pm}\,2.1}  % experiments/FIGS/data/v2/index.json :: fine-tuning qg_herwig N=100000: R30, R42_Q1 :: sha256 c6c4138ee5c86423
\newcommand{\MlstFtQgHerwigRthirtyMidOverScratchEFive}{0.97\,\ensuremath{\pm}\,0.06}  % experiments/FIGS/data/v2/index.json :: fine-tuning qg_herwig N=100000: R30 of each R42_Q1 run over the from-scratch mean, SD over runs and standard error of that mean in quadrature :: sha256 c6c4138ee5c86423
\newcommand{\MlstFtQgHerwigRthirtyMidOverScratchEFivePr}{equal within 6\%}  % experiments/FIGS/data/v2/index.json :: fine-tuning qg_herwig N=100000: R30 of each R42_Q1 run over the from-scratch mean, SD over runs and standard error of that mean in quadrature :: sha256 c6c4138ee5c86423 :: prose twin: two significant figures, direction in words
\newcommand{\MlstFtQgHerwigRthirtyCoarseEFive}{34.9\,\ensuremath{\pm}\,1.1}  % experiments/FIGS/data/v2/index.json :: fine-tuning qg_herwig N=100000: R30, R16_Q1 :: sha256 c6c4138ee5c86423
\newcommand{\MlstFtQgHerwigRthirtyCoarseOverScratchEFive}{0.950\,\ensuremath{\pm}\,0.033}  % experiments/FIGS/data/v2/index.json :: fine-tuning qg_herwig N=100000: R30 of each R16_Q1 run over the from-scratch mean, SD over runs and standard error of that mean in quadrature :: sha256 c6c4138ee5c86423
\newcommand{\MlstFtQgHerwigRthirtyCoarseOverScratchEFivePr}{equal within 6\%}  % experiments/FIGS/data/v2/index.json :: fine-tuning qg_herwig N=100000: R30 of each R16_Q1 run over the from-scratch mean, SD over runs and standard error of that mean in quadrature :: sha256 c6c4138ee5c86423 :: prose twin: two significant figures, direction in words
\newcommand{\MlstFtQgHerwigRthirtyResMassEFive}{34.0 and 36.5}  % experiments/FIGS/data/v2/index.json :: fine-tuning qg_herwig N=100000: R30, L162_MASS :: sha256 c6c4138ee5c86423
\newcommand{\MlstFtQgHerwigRthirtyResMassOverScratchEFive}{0.925 and 0.994}  % experiments/FIGS/data/v2/index.json :: fine-tuning qg_herwig N=100000: R30 of each L162_MASS run over the from-scratch mean, SD over runs and standard error of that mean in quadrature :: sha256 c6c4138ee5c86423
\newcommand{\MlstFtQgHerwigRthirtyResMassOverScratchEFivePr}{at least 0.63\% less rejection}  % experiments/FIGS/data/v2/index.json :: fine-tuning qg_herwig N=100000: R30 of each L162_MASS run over the from-scratch mean, SD over runs and standard error of that mean in quadrature :: sha256 c6c4138ee5c86423 :: prose twin: two significant figures, direction in words
\newcommand{\MlstFtQgHerwigRthirtyResMassOverScratchEFivePrMax}{up to 7.5\% less rejection}  % experiments/FIGS/data/v2/index.json :: fine-tuning qg_herwig N=100000: R30 of each L162_MASS run over the from-scratch mean, SD over runs and standard error of that mean in quadrature :: sha256 c6c4138ee5c86423 :: prose twin of the run farther from 1
\newcommand{\MlstFtQgHerwigRthirtyCoarseMassEFive}{35.9 and 39.0}  % experiments/FIGS/data/v2/index.json :: fine-tuning qg_herwig N=100000: R30, R16_Q1_MASS :: sha256 c6c4138ee5c86423
\newcommand{\MlstFtQgHerwigRthirtyCoarseMassOverScratchEFive}{0.979 and 1.064}  % experiments/FIGS/data/v2/index.json :: fine-tuning qg_herwig N=100000: R30 of each R16_Q1_MASS run over the from-scratch mean, SD over runs and standard error of that mean in quadrature :: sha256 c6c4138ee5c86423
\newcommand{\MlstFtQgHerwigRthirtyCoarseMassOverScratchEFivePr}{equal within 7\%}  % experiments/FIGS/data/v2/index.json :: fine-tuning qg_herwig N=100000: R30 of each R16_Q1_MASS run over the from-scratch mean, SD over runs and standard error of that mean in quadrature :: sha256 c6c4138ee5c86423 :: prose twin: two significant figures, direction in words
\newcommand{\MlstFtQgHerwigRthirtyCoarseMassOverScratchEFivePrMax}{equal within 7\%}  % experiments/FIGS/data/v2/index.json :: fine-tuning qg_herwig N=100000: R30 of each R16_Q1_MASS run over the from-scratch mean, SD over runs and standard error of that mean in quadrature :: sha256 c6c4138ee5c86423 :: prose twin of the run farther from 1
\newcommand{\MlstFtQgHerwigRthirtyResOverFullEFive}{1.03\,\ensuremath{\pm}\,0.04}  % experiments/FIGS/data/v2/index.json :: fine-tuning qg_herwig N=100000: R30 of L162 over L188, run by run :: sha256 c6c4138ee5c86423
\newcommand{\MlstFtQgHerwigRthirtyResOverFullEFivePr}{equal within 5\%}  % experiments/FIGS/data/v2/index.json :: fine-tuning qg_herwig N=100000: R30 of L162 over L188, run by run :: sha256 c6c4138ee5c86423 :: prose twin: two significant figures, direction in words
\newcommand{\MlstFtQgHerwigRthirtyMidOverFullEFive}{1.03\,\ensuremath{\pm}\,0.05}  % experiments/FIGS/data/v2/index.json :: fine-tuning qg_herwig N=100000: R30 of R42_Q1 over L188, run by run :: sha256 c6c4138ee5c86423
\newcommand{\MlstFtQgHerwigRthirtyMidOverFullEFivePr}{equal within 6\%}  % experiments/FIGS/data/v2/index.json :: fine-tuning qg_herwig N=100000: R30 of R42_Q1 over L188, run by run :: sha256 c6c4138ee5c86423 :: prose twin: two significant figures, direction in words
\newcommand{\MlstFtQgHerwigRthirtyCoarseOverFullEFive}{1.02\,\ensuremath{\pm}\,0.06}  % experiments/FIGS/data/v2/index.json :: fine-tuning qg_herwig N=100000: R30 of R16_Q1 over L188, run by run :: sha256 c6c4138ee5c86423
\newcommand{\MlstFtQgHerwigRthirtyCoarseOverFullEFivePr}{equal within 7\%}  % experiments/FIGS/data/v2/index.json :: fine-tuning qg_herwig N=100000: R30 of R16_Q1 over L188, run by run :: sha256 c6c4138ee5c86423 :: prose twin: two significant figures, direction in words
\newcommand{\MlstFtQgHerwigRthirtyResMassOverResEFive}{0.971 and 1.028}  % experiments/FIGS/data/v2/index.json :: fine-tuning qg_herwig N=100000: R30 of L162_MASS over L162, run by run :: sha256 c6c4138ee5c86423
\newcommand{\MlstFtQgHerwigRthirtyResMassOverResEFivePr}{equal within 3\%}  % experiments/FIGS/data/v2/index.json :: fine-tuning qg_herwig N=100000: R30 of L162_MASS over L162, run by run :: sha256 c6c4138ee5c86423 :: prose twin: two significant figures, direction in words
\newcommand{\MlstFtQgHerwigRthirtyResMassOverResEFivePrMax}{equal within 3\%}  % experiments/FIGS/data/v2/index.json :: fine-tuning qg_herwig N=100000: R30 of L162_MASS over L162, run by run :: sha256 c6c4138ee5c86423 :: prose twin of the run farther from 1
\newcommand{\MlstFtQgHerwigRthirtyCoarseMassOverCoarseEFive}{1.064 and 1.121}  % experiments/FIGS/data/v2/index.json :: fine-tuning qg_herwig N=100000: R30 of R16_Q1_MASS over R16_Q1, run by run :: sha256 c6c4138ee5c86423
\newcommand{\MlstFtQgHerwigRthirtyCoarseMassOverCoarseEFivePr}{equal within 20\%}  % experiments/FIGS/data/v2/index.json :: fine-tuning qg_herwig N=100000: R30 of R16_Q1_MASS over R16_Q1, run by run :: sha256 c6c4138ee5c86423 :: prose twin: two significant figures, direction in words
\newcommand{\MlstFtQgHerwigRthirtyCoarseMassOverCoarseEFivePrMax}{equal within 20\%}  % experiments/FIGS/data/v2/index.json :: fine-tuning qg_herwig N=100000: R30 of R16_Q1_MASS over R16_Q1, run by run :: sha256 c6c4138ee5c86423 :: prose twin of the run farther from 1
\newcommand{\MlstFtQgHerwigRthirtyRandOverFullEFive}{1.01\,\ensuremath{\pm}\,0.05}  % experiments/FIGS/data/v2/index.json :: fine-tuning qg_herwig N=100000: R30 of each random partition over the L188 run of its seed :: sha256 c6c4138ee5c86423
\newcommand{\MlstFtQgHerwigRthirtyRandOverFullEFivePr}{equal within 5\%}  % experiments/FIGS/data/v2/index.json :: fine-tuning qg_herwig N=100000: R30 of each random partition over the L188 run of its seed :: sha256 c6c4138ee5c86423 :: prose twin: two significant figures, direction in words
\newcommand{\MlstFtQgHerwigRthirtyRandOverCoarseEFive}{1.009\,\ensuremath{\pm}\,0.024}  % experiments/FIGS/data/v2/index.json :: fine-tuning qg_herwig N=100000: R30 of each random partition over the R16_Q1 run of its seed :: sha256 c6c4138ee5c86423
\newcommand{\MlstFtQgHerwigRthirtyRandOverCoarseEFivePr}{equal within 3\%}  % experiments/FIGS/data/v2/index.json :: fine-tuning qg_herwig N=100000: R30 of each random partition over the R16_Q1 run of its seed :: sha256 c6c4138ee5c86423 :: prose twin: two significant figures, direction in words
\newcommand{\MlstFtQgHerwigRthirtyScratch}{35.5\,\ensuremath{\pm}\,1.2}  % experiments/FIGS/data/v2/index.json :: fine-tuning qg_herwig N=1600000: R30 from scratch, every seed :: sha256 c6c4138ee5c86423
\newcommand{\MlstFtQgHerwigRthirtyFull}{35.2\,\ensuremath{\pm}\,2.1}  % experiments/FIGS/data/v2/index.json :: fine-tuning qg_herwig N=1600000: R30, L188 :: sha256 c6c4138ee5c86423
\newcommand{\MlstFtQgHerwigRthirtyFullOverScratch}{0.99\,\ensuremath{\pm}\,0.06}  % experiments/FIGS/data/v2/index.json :: fine-tuning qg_herwig N=1600000: R30 of each L188 run over the from-scratch mean, SD over runs and standard error of that mean in quadrature :: sha256 c6c4138ee5c86423
\newcommand{\MlstFtQgHerwigRthirtyFullOverScratchPr}{equal within 6\%}  % experiments/FIGS/data/v2/index.json :: fine-tuning qg_herwig N=1600000: R30 of each L188 run over the from-scratch mean, SD over runs and standard error of that mean in quadrature :: sha256 c6c4138ee5c86423 :: prose twin: two significant figures, direction in words
\newcommand{\MlstFtQgHerwigRthirtyRes}{34.6\,\ensuremath{\pm}\,2.0}  % experiments/FIGS/data/v2/index.json :: fine-tuning qg_herwig N=1600000: R30, L162 :: sha256 c6c4138ee5c86423
\newcommand{\MlstFtQgHerwigRthirtyResOverScratch}{0.97\,\ensuremath{\pm}\,0.06}  % experiments/FIGS/data/v2/index.json :: fine-tuning qg_herwig N=1600000: R30 of each L162 run over the from-scratch mean, SD over runs and standard error of that mean in quadrature :: sha256 c6c4138ee5c86423
\newcommand{\MlstFtQgHerwigRthirtyResOverScratchPr}{equal within 6\%}  % experiments/FIGS/data/v2/index.json :: fine-tuning qg_herwig N=1600000: R30 of each L162 run over the from-scratch mean, SD over runs and standard error of that mean in quadrature :: sha256 c6c4138ee5c86423 :: prose twin: two significant figures, direction in words
\newcommand{\MlstFtQgHerwigRthirtyMid}{36.1\,\ensuremath{\pm}\,2.7}  % experiments/FIGS/data/v2/index.json :: fine-tuning qg_herwig N=1600000: R30, R42_Q1 :: sha256 c6c4138ee5c86423
\newcommand{\MlstFtQgHerwigRthirtyMidOverScratch}{1.02\,\ensuremath{\pm}\,0.08}  % experiments/FIGS/data/v2/index.json :: fine-tuning qg_herwig N=1600000: R30 of each R42_Q1 run over the from-scratch mean, SD over runs and standard error of that mean in quadrature :: sha256 c6c4138ee5c86423
\newcommand{\MlstFtQgHerwigRthirtyMidOverScratchPr}{equal within 8\%}  % experiments/FIGS/data/v2/index.json :: fine-tuning qg_herwig N=1600000: R30 of each R42_Q1 run over the from-scratch mean, SD over runs and standard error of that mean in quadrature :: sha256 c6c4138ee5c86423 :: prose twin: two significant figures, direction in words
\newcommand{\MlstFtQgHerwigRthirtyCoarse}{36.9\,\ensuremath{\pm}\,0.9}  % experiments/FIGS/data/v2/index.json :: fine-tuning qg_herwig N=1600000: R30, R16_Q1 :: sha256 c6c4138ee5c86423
\newcommand{\MlstFtQgHerwigRthirtyCoarseOverScratch}{1.040\,\ensuremath{\pm}\,0.033}  % experiments/FIGS/data/v2/index.json :: fine-tuning qg_herwig N=1600000: R30 of each R16_Q1 run over the from-scratch mean, SD over runs and standard error of that mean in quadrature :: sha256 c6c4138ee5c86423
\newcommand{\MlstFtQgHerwigRthirtyCoarseOverScratchPr}{equal within 4\%}  % experiments/FIGS/data/v2/index.json :: fine-tuning qg_herwig N=1600000: R30 of each R16_Q1 run over the from-scratch mean, SD over runs and standard error of that mean in quadrature :: sha256 c6c4138ee5c86423 :: prose twin: two significant figures, direction in words
\newcommand{\MlstFtQgHerwigRthirtyResMass}{37.8 and 34.4}  % experiments/FIGS/data/v2/index.json :: fine-tuning qg_herwig N=1600000: R30, L162_MASS :: sha256 c6c4138ee5c86423
\newcommand{\MlstFtQgHerwigRthirtyResMassOverScratch}{1.065 and 0.970}  % experiments/FIGS/data/v2/index.json :: fine-tuning qg_herwig N=1600000: R30 of each L162_MASS run over the from-scratch mean, SD over runs and standard error of that mean in quadrature :: sha256 c6c4138ee5c86423
\newcommand{\MlstFtQgHerwigRthirtyResMassOverScratchPr}{equal within 7\%}  % experiments/FIGS/data/v2/index.json :: fine-tuning qg_herwig N=1600000: R30 of each L162_MASS run over the from-scratch mean, SD over runs and standard error of that mean in quadrature :: sha256 c6c4138ee5c86423 :: prose twin: two significant figures, direction in words
\newcommand{\MlstFtQgHerwigRthirtyResMassOverScratchPrMax}{equal within 7\%}  % experiments/FIGS/data/v2/index.json :: fine-tuning qg_herwig N=1600000: R30 of each L162_MASS run over the from-scratch mean, SD over runs and standard error of that mean in quadrature :: sha256 c6c4138ee5c86423 :: prose twin of the run farther from 1
\newcommand{\MlstFtQgHerwigRthirtyCoarseMass}{37.7 and 32.9}  % experiments/FIGS/data/v2/index.json :: fine-tuning qg_herwig N=1600000: R30, R16_Q1_MASS :: sha256 c6c4138ee5c86423
\newcommand{\MlstFtQgHerwigRthirtyCoarseMassOverScratch}{1.06 and 0.93}  % experiments/FIGS/data/v2/index.json :: fine-tuning qg_herwig N=1600000: R30 of each R16_Q1_MASS run over the from-scratch mean, SD over runs and standard error of that mean in quadrature :: sha256 c6c4138ee5c86423
\newcommand{\MlstFtQgHerwigRthirtyCoarseMassOverScratchPr}{equal within 8\%}  % experiments/FIGS/data/v2/index.json :: fine-tuning qg_herwig N=1600000: R30 of each R16_Q1_MASS run over the from-scratch mean, SD over runs and standard error of that mean in quadrature :: sha256 c6c4138ee5c86423 :: prose twin: two significant figures, direction in words
\newcommand{\MlstFtQgHerwigRthirtyCoarseMassOverScratchPrMax}{equal within 8\%}  % experiments/FIGS/data/v2/index.json :: fine-tuning qg_herwig N=1600000: R30 of each R16_Q1_MASS run over the from-scratch mean, SD over runs and standard error of that mean in quadrature :: sha256 c6c4138ee5c86423 :: prose twin of the run farther from 1
\newcommand{\MlstFtQgHerwigRthirtyResOverFull}{0.9832\,\ensuremath{\pm}\,0.0022}  % experiments/FIGS/data/v2/index.json :: fine-tuning qg_herwig N=1600000: R30 of L162 over L188, run by run :: sha256 c6c4138ee5c86423
\newcommand{\MlstFtQgHerwigRthirtyResOverFullPr}{1.7\% less rejection}  % experiments/FIGS/data/v2/index.json :: fine-tuning qg_herwig N=1600000: R30 of L162 over L188, run by run :: sha256 c6c4138ee5c86423 :: prose twin: two significant figures, direction in words
\newcommand{\MlstFtQgHerwigRthirtyMidOverFull}{1.03\,\ensuremath{\pm}\,0.05}  % experiments/FIGS/data/v2/index.json :: fine-tuning qg_herwig N=1600000: R30 of R42_Q1 over L188, run by run :: sha256 c6c4138ee5c86423
\newcommand{\MlstFtQgHerwigRthirtyMidOverFullPr}{equal within 6\%}  % experiments/FIGS/data/v2/index.json :: fine-tuning qg_herwig N=1600000: R30 of R42_Q1 over L188, run by run :: sha256 c6c4138ee5c86423 :: prose twin: two significant figures, direction in words
\newcommand{\MlstFtQgHerwigRthirtyCoarseOverFull}{1.05\,\ensuremath{\pm}\,0.06}  % experiments/FIGS/data/v2/index.json :: fine-tuning qg_herwig N=1600000: R30 of R16_Q1 over L188, run by run :: sha256 c6c4138ee5c86423
\newcommand{\MlstFtQgHerwigRthirtyCoarseOverFullPr}{equal within 6\%}  % experiments/FIGS/data/v2/index.json :: fine-tuning qg_herwig N=1600000: R30 of R16_Q1 over L188, run by run :: sha256 c6c4138ee5c86423 :: prose twin: two significant figures, direction in words
\newcommand{\MlstFtQgHerwigRthirtyResMassOverRes}{1.14 and 0.93}  % experiments/FIGS/data/v2/index.json :: fine-tuning qg_herwig N=1600000: R30 of L162_MASS over L162, run by run :: sha256 c6c4138ee5c86423
\newcommand{\MlstFtQgHerwigRthirtyResMassOverResPr}{equal within 20\%}  % experiments/FIGS/data/v2/index.json :: fine-tuning qg_herwig N=1600000: R30 of L162_MASS over L162, run by run :: sha256 c6c4138ee5c86423 :: prose twin: two significant figures, direction in words
\newcommand{\MlstFtQgHerwigRthirtyResMassOverResPrMax}{equal within 20\%}  % experiments/FIGS/data/v2/index.json :: fine-tuning qg_herwig N=1600000: R30 of L162_MASS over L162, run by run :: sha256 c6c4138ee5c86423 :: prose twin of the run farther from 1
\newcommand{\MlstFtQgHerwigRthirtyCoarseMassOverCoarse}{1.01 and 0.88}  % experiments/FIGS/data/v2/index.json :: fine-tuning qg_herwig N=1600000: R30 of R16_Q1_MASS over R16_Q1, run by run :: sha256 c6c4138ee5c86423
\newcommand{\MlstFtQgHerwigRthirtyCoarseMassOverCoarsePr}{equal within 20\%}  % experiments/FIGS/data/v2/index.json :: fine-tuning qg_herwig N=1600000: R30 of R16_Q1_MASS over R16_Q1, run by run :: sha256 c6c4138ee5c86423 :: prose twin: two significant figures, direction in words
\newcommand{\MlstFtQgHerwigRthirtyCoarseMassOverCoarsePrMax}{equal within 20\%}  % experiments/FIGS/data/v2/index.json :: fine-tuning qg_herwig N=1600000: R30 of R16_Q1_MASS over R16_Q1, run by run :: sha256 c6c4138ee5c86423 :: prose twin of the run farther from 1
\newcommand{\MlstFtQgHerwigRthirtyRandOverFull}{0.98\,\ensuremath{\pm}\,0.04}  % experiments/FIGS/data/v2/index.json :: fine-tuning qg_herwig N=1600000: R30 of each random partition over the L188 run of its seed :: sha256 c6c4138ee5c86423
\newcommand{\MlstFtQgHerwigRthirtyRandOverFullPr}{equal within 5\%}  % experiments/FIGS/data/v2/index.json :: fine-tuning qg_herwig N=1600000: R30 of each random partition over the L188 run of its seed :: sha256 c6c4138ee5c86423 :: prose twin: two significant figures, direction in words
\newcommand{\MlstFtQgHerwigRthirtyRandOverCoarse}{0.94\,\ensuremath{\pm}\,0.04}  % experiments/FIGS/data/v2/index.json :: fine-tuning qg_herwig N=1600000: R30 of each random partition over the R16_Q1 run of its seed :: sha256 c6c4138ee5c86423
\newcommand{\MlstFtQgHerwigRthirtyRandOverCoarsePr}{equal within 6\%}  % experiments/FIGS/data/v2/index.json :: fine-tuning qg_herwig N=1600000: R30 of each random partition over the R16_Q1 run of its seed :: sha256 c6c4138ee5c86423 :: prose twin: two significant figures, direction in words
\newcommand{\MlstFtQgRthirtyHerwigOverPythiaFull}{0.328\,\ensuremath{\pm}\,0.015}  % experiments/FIGS/data/v2/index.json :: fine-tuning q/g N=1600000: R30 on Herwig over R30 on Pythia, each run of L188 :: sha256 c6c4138ee5c86423
\newcommand{\MlstFtQgRthirtyHerwigOverPythiaFullPr}{0.33 times the rejection}  % experiments/FIGS/data/v2/index.json :: fine-tuning q/g N=1600000: R30 on Herwig over R30 on Pythia, each run of L188 :: sha256 c6c4138ee5c86423 :: prose twin: two significant figures, direction in words
\newcommand{\MlstFtQgRthirtyHerwigOverPythiaRes}{0.320\,\ensuremath{\pm}\,0.023}  % experiments/FIGS/data/v2/index.json :: fine-tuning q/g N=1600000: R30 on Herwig over R30 on Pythia, each run of L162 :: sha256 c6c4138ee5c86423
\newcommand{\MlstFtQgRthirtyHerwigOverPythiaResPr}{0.32 times the rejection}  % experiments/FIGS/data/v2/index.json :: fine-tuning q/g N=1600000: R30 on Herwig over R30 on Pythia, each run of L162 :: sha256 c6c4138ee5c86423 :: prose twin: two significant figures, direction in words
\newcommand{\MlstFtQgRthirtyHerwigOverPythiaMid}{0.334\,\ensuremath{\pm}\,0.023}  % experiments/FIGS/data/v2/index.json :: fine-tuning q/g N=1600000: R30 on Herwig over R30 on Pythia, each run of R42_Q1 :: sha256 c6c4138ee5c86423
\newcommand{\MlstFtQgRthirtyHerwigOverPythiaMidPr}{0.33 times the rejection}  % experiments/FIGS/data/v2/index.json :: fine-tuning q/g N=1600000: R30 on Herwig over R30 on Pythia, each run of R42_Q1 :: sha256 c6c4138ee5c86423 :: prose twin: two significant figures, direction in words
\newcommand{\MlstFtQgRthirtyHerwigOverPythiaCoarse}{0.346\,\ensuremath{\pm}\,0.011}  % experiments/FIGS/data/v2/index.json :: fine-tuning q/g N=1600000: R30 on Herwig over R30 on Pythia, each run of R16_Q1 :: sha256 c6c4138ee5c86423
\newcommand{\MlstFtQgRthirtyHerwigOverPythiaCoarsePr}{0.35 times the rejection}  % experiments/FIGS/data/v2/index.json :: fine-tuning q/g N=1600000: R30 on Herwig over R30 on Pythia, each run of R16_Q1 :: sha256 c6c4138ee5c86423 :: prose twin: two significant figures, direction in words
\newcommand{\MlstFtQgRthirtyHerwigOverPythiaResMass}{0.353 and 0.323}  % experiments/FIGS/data/v2/index.json :: fine-tuning q/g N=1600000: R30 on Herwig over R30 on Pythia, each run of L162_MASS :: sha256 c6c4138ee5c86423
\newcommand{\MlstFtQgRthirtyHerwigOverPythiaResMassPr}{at most 0.35 times the rejection}  % experiments/FIGS/data/v2/index.json :: fine-tuning q/g N=1600000: R30 on Herwig over R30 on Pythia, each run of L162_MASS :: sha256 c6c4138ee5c86423 :: prose twin: two significant figures, direction in words
\newcommand{\MlstFtQgRthirtyHerwigOverPythiaResMassPrMax}{down to 0.32 times the rejection}  % experiments/FIGS/data/v2/index.json :: fine-tuning q/g N=1600000: R30 on Herwig over R30 on Pythia, each run of L162_MASS :: sha256 c6c4138ee5c86423 :: prose twin of the run farther from 1
\newcommand{\MlstFtQgRthirtyHerwigOverPythiaCoarseMass}{0.359 and 0.311}  % experiments/FIGS/data/v2/index.json :: fine-tuning q/g N=1600000: R30 on Herwig over R30 on Pythia, each run of R16_Q1_MASS :: sha256 c6c4138ee5c86423
\newcommand{\MlstFtQgRthirtyHerwigOverPythiaCoarseMassPr}{at most 0.36 times the rejection}  % experiments/FIGS/data/v2/index.json :: fine-tuning q/g N=1600000: R30 on Herwig over R30 on Pythia, each run of R16_Q1_MASS :: sha256 c6c4138ee5c86423 :: prose twin: two significant figures, direction in words
\newcommand{\MlstFtQgRthirtyHerwigOverPythiaCoarseMassPrMax}{down to 0.31 times the rejection}  % experiments/FIGS/data/v2/index.json :: fine-tuning q/g N=1600000: R30 on Herwig over R30 on Pythia, each run of R16_Q1_MASS :: sha256 c6c4138ee5c86423 :: prose twin of the run farther from 1
\newcommand{\MlstFtQgRthirtyHerwigOverPythiaScratch}{0.341\,\ensuremath{\pm}\,0.009}  % experiments/FIGS/data/v2/index.json :: fine-tuning q/g N=1600000: R30 on Herwig over R30 on Pythia, each from-scratch seed :: sha256 c6c4138ee5c86423
\newcommand{\MlstFtQgRthirtyHerwigOverPythiaScratchPr}{0.34 times the rejection}  % experiments/FIGS/data/v2/index.json :: fine-tuning q/g N=1600000: R30 on Herwig over R30 on Pythia, each from-scratch seed :: sha256 c6c4138ee5c86423 :: prose twin: two significant figures, direction in words
\newcommand{\MlstTopCheckFullLong}{\pending{recipe check}}  % experiments/FIGS/data/v2/index.json :: the recipe check is not read out yet :: sha256 c6c4138ee5c86423
\newcommand{\MlstTopCheckScratchLong}{\pending{recipe check}}  % experiments/FIGS/data/v2/index.json :: the recipe check is not read out yet :: sha256 c6c4138ee5c86423
\newcommand{\MlstTopCheckFullShort}{\pending{recipe check}}  % experiments/FIGS/data/v2/index.json :: the recipe check is not read out yet :: sha256 c6c4138ee5c86423
\newcommand{\MlstTopCheckScratchShort}{\pending{recipe check}}  % experiments/FIGS/data/v2/index.json :: the recipe check is not read out yet :: sha256 c6c4138ee5c86423
\newcommand{\MlstLpJcOneGapFracEFour}{56.3\,\ensuremath{\pm}\,1.1\%}  % experiments/FIGS/data/v2/index.json :: jc1 N=10000: (probe - scratch) / (fine-tuned - scratch) accuracy, 188-class runs, per cent :: sha256 c6c4138ee5c86423
\newcommand{\MlstLpJcOneGapFracEFourPr}{56\%}  % experiments/FIGS/data/v2/index.json :: jc1 N=10000: (probe - scratch) / (fine-tuned - scratch) accuracy, 188-class runs, per cent :: sha256 c6c4138ee5c86423 :: prose twin, two significant figures
\newcommand{\MlstLpJcOneGapFracEFive}{\ensuremath{-}15\,\ensuremath{\pm}\,5\%}  % experiments/FIGS/data/v2/index.json :: jc1 N=100000: (probe - scratch) / (fine-tuned - scratch) accuracy, 188-class runs, per cent :: sha256 c6c4138ee5c86423
\newcommand{\MlstLpJcOneGapFracEFivePr}{\ensuremath{-}15\%}  % experiments/FIGS/data/v2/index.json :: jc1 N=100000: (probe - scratch) / (fine-tuned - scratch) accuracy, 188-class runs, per cent :: sha256 c6c4138ee5c86423 :: prose twin, two significant figures
\newcommand{\MlstLpJcTwoGapFracEFour}{96.2\,\ensuremath{\pm}\,0.6\%}  % experiments/FIGS/data/v2/index.json :: jc2_188 N=10000: (probe - scratch) / (fine-tuned - scratch) accuracy, 188-class runs, per cent :: sha256 c6c4138ee5c86423
\newcommand{\MlstLpJcTwoGapFracEFourPr}{96\%}  % experiments/FIGS/data/v2/index.json :: jc2_188 N=10000: (probe - scratch) / (fine-tuned - scratch) accuracy, 188-class runs, per cent :: sha256 c6c4138ee5c86423 :: prose twin, two significant figures
\newcommand{\MlstLpJcTwoGapFracEFive}{102.2\,\ensuremath{\pm}\,1.0\%}  % experiments/FIGS/data/v2/index.json :: jc2_188 N=100000: (probe - scratch) / (fine-tuned - scratch) accuracy, 188-class runs, per cent :: sha256 c6c4138ee5c86423
\newcommand{\MlstLpJcTwoGapFracEFivePr}{100\%}  % experiments/FIGS/data/v2/index.json :: jc2_188 N=100000: (probe - scratch) / (fine-tuned - scratch) accuracy, 188-class runs, per cent :: sha256 c6c4138ee5c86423 :: prose twin, two significant figures
\newcommand{\XbxMx}{\pending{samples not in}}  % experiments/XBS/xbx_report.py :: pending :: sha256 cb6dffe774ff5b1e
\newcommand{\XbxWindowLo}{\pending{samples not in}}  % experiments/XBS/xbx_report.py :: pending :: sha256 cb6dffe774ff5b1e
\newcommand{\XbxWindowHi}{\pending{samples not in}}  % experiments/XBS/xbx_report.py :: pending :: sha256 cb6dffe774ff5b1e
\newcommand{\XbbNTrainPerClass}{\pending{samples not in}}  % experiments/XBS/xbx_report.py :: pending :: sha256 cb6dffe774ff5b1e
\newcommand{\XbbRepetition}{\pending{samples not in}}  % experiments/XBS/xbx_report.py :: pending :: sha256 cb6dffe774ff5b1e
\newcommand{\XbbNEvalSignal}{\pending{samples not in}}  % experiments/XBS/xbx_report.py :: pending :: sha256 cb6dffe774ff5b1e
\newcommand{\XbbNEvalQcd}{\pending{samples not in}}  % experiments/XBS/xbx_report.py :: pending :: sha256 cb6dffe774ff5b1e
\newcommand{\XbsNTrainPerClass}{\pending{samples not in}}  % experiments/XBS/xbx_report.py :: pending :: sha256 cb6dffe774ff5b1e
\newcommand{\XbsRepetition}{\pending{samples not in}}  % experiments/XBS/xbx_report.py :: pending :: sha256 cb6dffe774ff5b1e
\newcommand{\XbsNEvalSignal}{\pending{samples not in}}  % experiments/XBS/xbx_report.py :: pending :: sha256 cb6dffe774ff5b1e
\newcommand{\XbsNEvalQcd}{\pending{samples not in}}  % experiments/XBS/xbx_report.py :: pending :: sha256 cb6dffe774ff5b1e
\newcommand{\MlstPartreproRthirtyPart}{1{,}711}  % experiments/FIGS/data/v2/partrepro/part.json :: r30 :: sha256 0d44b997af9e6448 :: ParT reproduction on top, part: R30 (one run)
\newcommand{\MlstPartreproRthirtyHead}{1{,}577}  % experiments/FIGS/data/v2/partrepro/head.json :: r30 :: sha256 37b82ba0a0fe9299 :: ParT reproduction on top, head: R30 (one run)
\newcommand{\MlstPartreproRthirtyRescale}{1{,}224}  % experiments/FIGS/data/v2/partrepro/rescale.json :: r30 :: sha256 0a9f1401047d07ff :: ParT reproduction on top, rescale: R30 (one run)
\newcommand{\MlstPartreproRthirtyHeadOverPart}{0.922}  % experiments/FIGS/data/v2/partrepro/head.json :: r30 / part r30 :: sha256 37b82ba0a0fe9299
\newcommand{\MlstPartreproRthirtyHeadOverPartPr}{7.8\% less rejection}  % experiments/FIGS/data/v2/partrepro/head.json :: r30 / part r30 :: sha256 37b82ba0a0fe9299 :: prose twin: two significant figures, direction in words
\newcommand{\MlstPartreproRthirtyRescaleOverPart}{0.715}  % experiments/FIGS/data/v2/partrepro/rescale.json :: r30 / part r30 :: sha256 0a9f1401047d07ff
\newcommand{\MlstPartreproRthirtyRescaleOverPartPr}{28\% less rejection}  % experiments/FIGS/data/v2/partrepro/rescale.json :: r30 / part r30 :: sha256 0a9f1401047d07ff :: prose twin: two significant figures, direction in words
\newcommand{\MlstPartreproSliceDropRatio}{\pending{jet-pT slice test not run}}  % experiments/FIGS/data/v2/index.json :: open item O10: R30 in narrow pT slices of partrepro part and rescale :: sha256 c6c4138ee5c86423
\newcommand{\MlstInputBnZeroSigmaMax}{0.96}  % experiments/FIGS/data/v2/head/input_bn_offsets.json :: offsets :: sha256 d6b821509fccecff :: largest |offset| over the five type flags
\newcommand{\MlstFtJcOneNSignal}{nine}  % experiments/FIGS/data/v2/curves/full_budget_long_leg2.json :: JC1_COLS (count) :: sha256 588e935996d1006a :: JetClass signal classes
\newcommand{\MlstFtJcOneRejHbbResOverFull}{0.92\,\ensuremath{\pm}\,0.05}  % experiments/FIGS/data/v2/curves/full_budget_long_leg2.json :: cells.<run>.N1000000.s1.jc1.per_class.label_Hbb: rejection of L162 over L188, run by run :: sha256 588e935996d1006a
\newcommand{\MlstFtJcOneRejHbbResOverFullPr}{equal within 9\%}  % experiments/FIGS/data/v2/curves/full_budget_long_leg2.json :: cells.<run>.N1000000.s1.jc1.per_class.label_Hbb: rejection of L162 over L188, run by run :: sha256 588e935996d1006a :: prose twin: two significant figures, direction in words
\newcommand{\MlstFtJcOneRejHbbMidOverFull}{0.98\,\ensuremath{\pm}\,0.12}  % experiments/FIGS/data/v2/curves/full_budget_long_leg2.json :: cells.<run>.N1000000.s1.jc1.per_class.label_Hbb: rejection of R42_Q1 over L188, run by run :: sha256 588e935996d1006a
\newcommand{\MlstFtJcOneRejHbbMidOverFullPr}{equal within 20\%}  % experiments/FIGS/data/v2/curves/full_budget_long_leg2.json :: cells.<run>.N1000000.s1.jc1.per_class.label_Hbb: rejection of R42_Q1 over L188, run by run :: sha256 588e935996d1006a :: prose twin: two significant figures, direction in words
\newcommand{\MlstFtJcOneRejHbbCoarseOverFull}{1.01\,\ensuremath{\pm}\,0.07}  % experiments/FIGS/data/v2/curves/full_budget_long_leg2.json :: cells.<run>.N1000000.s1.jc1.per_class.label_Hbb: rejection of R16_Q1 over L188, run by run :: sha256 588e935996d1006a
\newcommand{\MlstFtJcOneRejHbbCoarseOverFullPr}{equal within 7\%}  % experiments/FIGS/data/v2/curves/full_budget_long_leg2.json :: cells.<run>.N1000000.s1.jc1.per_class.label_Hbb: rejection of R16_Q1 over L188, run by run :: sha256 588e935996d1006a :: prose twin: two significant figures, direction in words
\newcommand{\MlstFtJcOneRejHbbResMassOverRes}{1.065 and 1.100}  % experiments/FIGS/data/v2/curves/full_budget_long_leg2.json :: cells.<run>.N1000000.s1.jc1.per_class.label_Hbb: rejection of L162_MASS over L162, run by run :: sha256 588e935996d1006a
\newcommand{\MlstFtJcOneRejHbbResMassOverResPr}{equal within 10\%}  % experiments/FIGS/data/v2/curves/full_budget_long_leg2.json :: cells.<run>.N1000000.s1.jc1.per_class.label_Hbb: rejection of L162_MASS over L162, run by run :: sha256 588e935996d1006a :: prose twin: two significant figures, direction in words
\newcommand{\MlstFtJcOneRejHbbResMassOverResPrMax}{equal within 10\%}  % experiments/FIGS/data/v2/curves/full_budget_long_leg2.json :: cells.<run>.N1000000.s1.jc1.per_class.label_Hbb: rejection of L162_MASS over L162, run by run :: sha256 588e935996d1006a :: prose twin of the run farther from 1
\newcommand{\MlstFtJcOneRejHbbCoarseMassOverCoarse}{0.97 and 0.86}  % experiments/FIGS/data/v2/curves/full_budget_long_leg2.json :: cells.<run>.N1000000.s1.jc1.per_class.label_Hbb: rejection of R16_Q1_MASS over R16_Q1, run by run :: sha256 588e935996d1006a
\newcommand{\MlstFtJcOneRejHbbCoarseMassOverCoarsePr}{equal within 20\%}  % experiments/FIGS/data/v2/curves/full_budget_long_leg2.json :: cells.<run>.N1000000.s1.jc1.per_class.label_Hbb: rejection of R16_Q1_MASS over R16_Q1, run by run :: sha256 588e935996d1006a :: prose twin: two significant figures, direction in words
\newcommand{\MlstFtJcOneRejHbbCoarseMassOverCoarsePrMax}{equal within 20\%}  % experiments/FIGS/data/v2/curves/full_budget_long_leg2.json :: cells.<run>.N1000000.s1.jc1.per_class.label_Hbb: rejection of R16_Q1_MASS over R16_Q1, run by run :: sha256 588e935996d1006a :: prose twin of the run farther from 1
\newcommand{\MlstFtJcOneRejHbbRandOverFull}{0.92\,\ensuremath{\pm}\,0.11}  % experiments/FIGS/data/v2/curves/full_budget_long_leg2.json :: per_class.label_Hbb: rejection of each random partition over the 188-class run of its seed :: sha256 588e935996d1006a
\newcommand{\MlstFtJcOneRejHbbRandOverFullPr}{equal within 20\%}  % experiments/FIGS/data/v2/curves/full_budget_long_leg2.json :: per_class.label_Hbb: rejection of each random partition over the 188-class run of its seed :: sha256 588e935996d1006a :: prose twin: two significant figures, direction in words
\newcommand{\MlstFtJcOneRejHbbRandOverFullRange}{0.84--1.1}  % experiments/FIGS/data/v2/curves/full_budget_long_leg2.json :: per_class.label_Hbb: smallest and largest of the partitions' ratios :: sha256 588e935996d1006a
\newcommand{\MlstFtJcOneRejHccResOverFull}{1.07\,\ensuremath{\pm}\,0.12}  % experiments/FIGS/data/v2/curves/full_budget_long_leg2.json :: cells.<run>.N1000000.s1.jc1.per_class.label_Hcc: rejection of L162 over L188, run by run :: sha256 588e935996d1006a
\newcommand{\MlstFtJcOneRejHccResOverFullPr}{equal within 20\%}  % experiments/FIGS/data/v2/curves/full_budget_long_leg2.json :: cells.<run>.N1000000.s1.jc1.per_class.label_Hcc: rejection of L162 over L188, run by run :: sha256 588e935996d1006a :: prose twin: two significant figures, direction in words
\newcommand{\MlstFtJcOneRejHccMidOverFull}{0.99\,\ensuremath{\pm}\,0.10}  % experiments/FIGS/data/v2/curves/full_budget_long_leg2.json :: cells.<run>.N1000000.s1.jc1.per_class.label_Hcc: rejection of R42_Q1 over L188, run by run :: sha256 588e935996d1006a
\newcommand{\MlstFtJcOneRejHccMidOverFullPr}{equal within 20\%}  % experiments/FIGS/data/v2/curves/full_budget_long_leg2.json :: cells.<run>.N1000000.s1.jc1.per_class.label_Hcc: rejection of R42_Q1 over L188, run by run :: sha256 588e935996d1006a :: prose twin: two significant figures, direction in words
\newcommand{\MlstFtJcOneRejHccCoarseOverFull}{0.91\,\ensuremath{\pm}\,0.07}  % experiments/FIGS/data/v2/curves/full_budget_long_leg2.json :: cells.<run>.N1000000.s1.jc1.per_class.label_Hcc: rejection of R16_Q1 over L188, run by run :: sha256 588e935996d1006a
\newcommand{\MlstFtJcOneRejHccCoarseOverFullPr}{equal within 10\%}  % experiments/FIGS/data/v2/curves/full_budget_long_leg2.json :: cells.<run>.N1000000.s1.jc1.per_class.label_Hcc: rejection of R16_Q1 over L188, run by run :: sha256 588e935996d1006a :: prose twin: two significant figures, direction in words
\newcommand{\MlstFtJcOneRejHccResMassOverRes}{0.886 and 0.816}  % experiments/FIGS/data/v2/curves/full_budget_long_leg2.json :: cells.<run>.N1000000.s1.jc1.per_class.label_Hcc: rejection of L162_MASS over L162, run by run :: sha256 588e935996d1006a
\newcommand{\MlstFtJcOneRejHccResMassOverResPr}{equal within 20\%}  % experiments/FIGS/data/v2/curves/full_budget_long_leg2.json :: cells.<run>.N1000000.s1.jc1.per_class.label_Hcc: rejection of L162_MASS over L162, run by run :: sha256 588e935996d1006a :: prose twin: two significant figures, direction in words
\newcommand{\MlstFtJcOneRejHccResMassOverResPrMax}{equal within 20\%}  % experiments/FIGS/data/v2/curves/full_budget_long_leg2.json :: cells.<run>.N1000000.s1.jc1.per_class.label_Hcc: rejection of L162_MASS over L162, run by run :: sha256 588e935996d1006a :: prose twin of the run farther from 1
\newcommand{\MlstFtJcOneRejHccCoarseMassOverCoarse}{0.98 and 0.85}  % experiments/FIGS/data/v2/curves/full_budget_long_leg2.json :: cells.<run>.N1000000.s1.jc1.per_class.label_Hcc: rejection of R16_Q1_MASS over R16_Q1, run by run :: sha256 588e935996d1006a
\newcommand{\MlstFtJcOneRejHccCoarseMassOverCoarsePr}{equal within 20\%}  % experiments/FIGS/data/v2/curves/full_budget_long_leg2.json :: cells.<run>.N1000000.s1.jc1.per_class.label_Hcc: rejection of R16_Q1_MASS over R16_Q1, run by run :: sha256 588e935996d1006a :: prose twin: two significant figures, direction in words
\newcommand{\MlstFtJcOneRejHccCoarseMassOverCoarsePrMax}{equal within 20\%}  % experiments/FIGS/data/v2/curves/full_budget_long_leg2.json :: cells.<run>.N1000000.s1.jc1.per_class.label_Hcc: rejection of R16_Q1_MASS over R16_Q1, run by run :: sha256 588e935996d1006a :: prose twin of the run farther from 1
\newcommand{\MlstFtJcOneRejHccRandOverFull}{1.03\,\ensuremath{\pm}\,0.07}  % experiments/FIGS/data/v2/curves/full_budget_long_leg2.json :: per_class.label_Hcc: rejection of each random partition over the 188-class run of its seed :: sha256 588e935996d1006a
\newcommand{\MlstFtJcOneRejHccRandOverFullPr}{equal within 8\%}  % experiments/FIGS/data/v2/curves/full_budget_long_leg2.json :: per_class.label_Hcc: rejection of each random partition over the 188-class run of its seed :: sha256 588e935996d1006a :: prose twin: two significant figures, direction in words
\newcommand{\MlstFtJcOneRejHccRandOverFullRange}{0.94--1.1}  % experiments/FIGS/data/v2/curves/full_budget_long_leg2.json :: per_class.label_Hcc: smallest and largest of the partitions' ratios :: sha256 588e935996d1006a
\newcommand{\MlstFtJcOneRejHggResOverFull}{0.98\,\ensuremath{\pm}\,0.05}  % experiments/FIGS/data/v2/curves/full_budget_long_leg2.json :: cells.<run>.N1000000.s1.jc1.per_class.label_Hgg: rejection of L162 over L188, run by run :: sha256 588e935996d1006a
\newcommand{\MlstFtJcOneRejHggResOverFullPr}{equal within 6\%}  % experiments/FIGS/data/v2/curves/full_budget_long_leg2.json :: cells.<run>.N1000000.s1.jc1.per_class.label_Hgg: rejection of L162 over L188, run by run :: sha256 588e935996d1006a :: prose twin: two significant figures, direction in words
\newcommand{\MlstFtJcOneRejHggMidOverFull}{1.000\,\ensuremath{\pm}\,0.025}  % experiments/FIGS/data/v2/curves/full_budget_long_leg2.json :: cells.<run>.N1000000.s1.jc1.per_class.label_Hgg: rejection of R42_Q1 over L188, run by run :: sha256 588e935996d1006a
\newcommand{\MlstFtJcOneRejHggMidOverFullPr}{equal within 3\%}  % experiments/FIGS/data/v2/curves/full_budget_long_leg2.json :: cells.<run>.N1000000.s1.jc1.per_class.label_Hgg: rejection of R42_Q1 over L188, run by run :: sha256 588e935996d1006a :: prose twin: two significant figures, direction in words
\newcommand{\MlstFtJcOneRejHggCoarseOverFull}{0.97\,\ensuremath{\pm}\,0.09}  % experiments/FIGS/data/v2/curves/full_budget_long_leg2.json :: cells.<run>.N1000000.s1.jc1.per_class.label_Hgg: rejection of R16_Q1 over L188, run by run :: sha256 588e935996d1006a
\newcommand{\MlstFtJcOneRejHggCoarseOverFullPr}{equal within 10\%}  % experiments/FIGS/data/v2/curves/full_budget_long_leg2.json :: cells.<run>.N1000000.s1.jc1.per_class.label_Hgg: rejection of R16_Q1 over L188, run by run :: sha256 588e935996d1006a :: prose twin: two significant figures, direction in words
\newcommand{\MlstFtJcOneRejHggResMassOverRes}{1.046 and 1.008}  % experiments/FIGS/data/v2/curves/full_budget_long_leg2.json :: cells.<run>.N1000000.s1.jc1.per_class.label_Hgg: rejection of L162_MASS over L162, run by run :: sha256 588e935996d1006a
\newcommand{\MlstFtJcOneRejHggResMassOverResPr}{equal within 5\%}  % experiments/FIGS/data/v2/curves/full_budget_long_leg2.json :: cells.<run>.N1000000.s1.jc1.per_class.label_Hgg: rejection of L162_MASS over L162, run by run :: sha256 588e935996d1006a :: prose twin: two significant figures, direction in words
\newcommand{\MlstFtJcOneRejHggResMassOverResPrMax}{equal within 5\%}  % experiments/FIGS/data/v2/curves/full_budget_long_leg2.json :: cells.<run>.N1000000.s1.jc1.per_class.label_Hgg: rejection of L162_MASS over L162, run by run :: sha256 588e935996d1006a :: prose twin of the run farther from 1
\newcommand{\MlstFtJcOneRejHggCoarseMassOverCoarse}{1.12 and 1.01}  % experiments/FIGS/data/v2/curves/full_budget_long_leg2.json :: cells.<run>.N1000000.s1.jc1.per_class.label_Hgg: rejection of R16_Q1_MASS over R16_Q1, run by run :: sha256 588e935996d1006a
\newcommand{\MlstFtJcOneRejHggCoarseMassOverCoarsePr}{equal within 20\%}  % experiments/FIGS/data/v2/curves/full_budget_long_leg2.json :: cells.<run>.N1000000.s1.jc1.per_class.label_Hgg: rejection of R16_Q1_MASS over R16_Q1, run by run :: sha256 588e935996d1006a :: prose twin: two significant figures, direction in words
\newcommand{\MlstFtJcOneRejHggCoarseMassOverCoarsePrMax}{equal within 20\%}  % experiments/FIGS/data/v2/curves/full_budget_long_leg2.json :: cells.<run>.N1000000.s1.jc1.per_class.label_Hgg: rejection of R16_Q1_MASS over R16_Q1, run by run :: sha256 588e935996d1006a :: prose twin of the run farther from 1
\newcommand{\MlstFtJcOneRejHggRandOverFull}{0.968\,\ensuremath{\pm}\,0.016}  % experiments/FIGS/data/v2/curves/full_budget_long_leg2.json :: per_class.label_Hgg: rejection of each random partition over the 188-class run of its seed :: sha256 588e935996d1006a
\newcommand{\MlstFtJcOneRejHggRandOverFullPr}{equal within 4\%}  % experiments/FIGS/data/v2/curves/full_budget_long_leg2.json :: per_class.label_Hgg: rejection of each random partition over the 188-class run of its seed :: sha256 588e935996d1006a :: prose twin: two significant figures, direction in words
\newcommand{\MlstFtJcOneRejHggRandOverFullRange}{0.94--0.98}  % experiments/FIGS/data/v2/curves/full_budget_long_leg2.json :: per_class.label_Hgg: smallest and largest of the partitions' ratios :: sha256 588e935996d1006a
\newcommand{\MlstFtJcOneRejHfourqResOverFull}{1.05\,\ensuremath{\pm}\,0.06}  % experiments/FIGS/data/v2/curves/full_budget_long_leg2.json :: cells.<run>.N1000000.s1.jc1.per_class.label_H4q: rejection of L162 over L188, run by run :: sha256 588e935996d1006a
\newcommand{\MlstFtJcOneRejHfourqResOverFullPr}{equal within 7\%}  % experiments/FIGS/data/v2/curves/full_budget_long_leg2.json :: cells.<run>.N1000000.s1.jc1.per_class.label_H4q: rejection of L162 over L188, run by run :: sha256 588e935996d1006a :: prose twin: two significant figures, direction in words
\newcommand{\MlstFtJcOneRejHfourqMidOverFull}{0.95\,\ensuremath{\pm}\,0.04}  % experiments/FIGS/data/v2/curves/full_budget_long_leg2.json :: cells.<run>.N1000000.s1.jc1.per_class.label_H4q: rejection of R42_Q1 over L188, run by run :: sha256 588e935996d1006a
\newcommand{\MlstFtJcOneRejHfourqMidOverFullPr}{equal within 6\%}  % experiments/FIGS/data/v2/curves/full_budget_long_leg2.json :: cells.<run>.N1000000.s1.jc1.per_class.label_H4q: rejection of R42_Q1 over L188, run by run :: sha256 588e935996d1006a :: prose twin: two significant figures, direction in words
\newcommand{\MlstFtJcOneRejHfourqCoarseOverFull}{0.962\,\ensuremath{\pm}\,0.025}  % experiments/FIGS/data/v2/curves/full_budget_long_leg2.json :: cells.<run>.N1000000.s1.jc1.per_class.label_H4q: rejection of R16_Q1 over L188, run by run :: sha256 588e935996d1006a
\newcommand{\MlstFtJcOneRejHfourqCoarseOverFullPr}{equal within 4\%}  % experiments/FIGS/data/v2/curves/full_budget_long_leg2.json :: cells.<run>.N1000000.s1.jc1.per_class.label_H4q: rejection of R16_Q1 over L188, run by run :: sha256 588e935996d1006a :: prose twin: two significant figures, direction in words
\newcommand{\MlstFtJcOneRejHfourqResMassOverRes}{0.912 and 0.966}  % experiments/FIGS/data/v2/curves/full_budget_long_leg2.json :: cells.<run>.N1000000.s1.jc1.per_class.label_H4q: rejection of L162_MASS over L162, run by run :: sha256 588e935996d1006a
\newcommand{\MlstFtJcOneRejHfourqResMassOverResPr}{equal within 9\%}  % experiments/FIGS/data/v2/curves/full_budget_long_leg2.json :: cells.<run>.N1000000.s1.jc1.per_class.label_H4q: rejection of L162_MASS over L162, run by run :: sha256 588e935996d1006a :: prose twin: two significant figures, direction in words
\newcommand{\MlstFtJcOneRejHfourqResMassOverResPrMax}{equal within 9\%}  % experiments/FIGS/data/v2/curves/full_budget_long_leg2.json :: cells.<run>.N1000000.s1.jc1.per_class.label_H4q: rejection of L162_MASS over L162, run by run :: sha256 588e935996d1006a :: prose twin of the run farther from 1
\newcommand{\MlstFtJcOneRejHfourqCoarseMassOverCoarse}{0.937 and 0.888}  % experiments/FIGS/data/v2/curves/full_budget_long_leg2.json :: cells.<run>.N1000000.s1.jc1.per_class.label_H4q: rejection of R16_Q1_MASS over R16_Q1, run by run :: sha256 588e935996d1006a
\newcommand{\MlstFtJcOneRejHfourqCoarseMassOverCoarsePr}{equal within 20\%}  % experiments/FIGS/data/v2/curves/full_budget_long_leg2.json :: cells.<run>.N1000000.s1.jc1.per_class.label_H4q: rejection of R16_Q1_MASS over R16_Q1, run by run :: sha256 588e935996d1006a :: prose twin: two significant figures, direction in words
\newcommand{\MlstFtJcOneRejHfourqCoarseMassOverCoarsePrMax}{equal within 20\%}  % experiments/FIGS/data/v2/curves/full_budget_long_leg2.json :: cells.<run>.N1000000.s1.jc1.per_class.label_H4q: rejection of R16_Q1_MASS over R16_Q1, run by run :: sha256 588e935996d1006a :: prose twin of the run farther from 1
\newcommand{\MlstFtJcOneRejHfourqRandOverFull}{0.93\,\ensuremath{\pm}\,0.04}  % experiments/FIGS/data/v2/curves/full_budget_long_leg2.json :: per_class.label_H4q: rejection of each random partition over the 188-class run of its seed :: sha256 588e935996d1006a
\newcommand{\MlstFtJcOneRejHfourqRandOverFullPr}{equal within 8\%}  % experiments/FIGS/data/v2/curves/full_budget_long_leg2.json :: per_class.label_H4q: rejection of each random partition over the 188-class run of its seed :: sha256 588e935996d1006a :: prose twin: two significant figures, direction in words
\newcommand{\MlstFtJcOneRejHfourqRandOverFullRange}{0.86--0.96}  % experiments/FIGS/data/v2/curves/full_budget_long_leg2.json :: per_class.label_H4q: smallest and largest of the partitions' ratios :: sha256 588e935996d1006a
\newcommand{\MlstFtJcOneRejHqqlResOverFull}{1.004\,\ensuremath{\pm}\,0.018}  % experiments/FIGS/data/v2/curves/full_budget_long_leg2.json :: cells.<run>.N1000000.s1.jc1.per_class.label_Hqql: rejection of L162 over L188, run by run :: sha256 588e935996d1006a
\newcommand{\MlstFtJcOneRejHqqlResOverFullPr}{equal within 2\%}  % experiments/FIGS/data/v2/curves/full_budget_long_leg2.json :: cells.<run>.N1000000.s1.jc1.per_class.label_Hqql: rejection of L162 over L188, run by run :: sha256 588e935996d1006a :: prose twin: two significant figures, direction in words
\newcommand{\MlstFtJcOneRejHqqlMidOverFull}{1.04\,\ensuremath{\pm}\,0.18}  % experiments/FIGS/data/v2/curves/full_budget_long_leg2.json :: cells.<run>.N1000000.s1.jc1.per_class.label_Hqql: rejection of R42_Q1 over L188, run by run :: sha256 588e935996d1006a
\newcommand{\MlstFtJcOneRejHqqlMidOverFullPr}{equal within 20\%}  % experiments/FIGS/data/v2/curves/full_budget_long_leg2.json :: cells.<run>.N1000000.s1.jc1.per_class.label_Hqql: rejection of R42_Q1 over L188, run by run :: sha256 588e935996d1006a :: prose twin: two significant figures, direction in words
\newcommand{\MlstFtJcOneRejHqqlCoarseOverFull}{0.93\,\ensuremath{\pm}\,0.16}  % experiments/FIGS/data/v2/curves/full_budget_long_leg2.json :: cells.<run>.N1000000.s1.jc1.per_class.label_Hqql: rejection of R16_Q1 over L188, run by run :: sha256 588e935996d1006a
\newcommand{\MlstFtJcOneRejHqqlCoarseOverFullPr}{equal within 20\%}  % experiments/FIGS/data/v2/curves/full_budget_long_leg2.json :: cells.<run>.N1000000.s1.jc1.per_class.label_Hqql: rejection of R16_Q1 over L188, run by run :: sha256 588e935996d1006a :: prose twin: two significant figures, direction in words
\newcommand{\MlstFtJcOneRejHqqlResMassOverRes}{0.986 and 1.000}  % experiments/FIGS/data/v2/curves/full_budget_long_leg2.json :: cells.<run>.N1000000.s1.jc1.per_class.label_Hqql: rejection of L162_MASS over L162, run by run :: sha256 588e935996d1006a
\newcommand{\MlstFtJcOneRejHqqlResMassOverResPr}{equal within 2\%}  % experiments/FIGS/data/v2/curves/full_budget_long_leg2.json :: cells.<run>.N1000000.s1.jc1.per_class.label_Hqql: rejection of L162_MASS over L162, run by run :: sha256 588e935996d1006a :: prose twin: two significant figures, direction in words
\newcommand{\MlstFtJcOneRejHqqlResMassOverResPrMax}{equal within 2\%}  % experiments/FIGS/data/v2/curves/full_budget_long_leg2.json :: cells.<run>.N1000000.s1.jc1.per_class.label_Hqql: rejection of L162_MASS over L162, run by run :: sha256 588e935996d1006a :: prose twin of the run farther from 1
\newcommand{\MlstFtJcOneRejHqqlCoarseMassOverCoarse}{0.97 and 1.09}  % experiments/FIGS/data/v2/curves/full_budget_long_leg2.json :: cells.<run>.N1000000.s1.jc1.per_class.label_Hqql: rejection of R16_Q1_MASS over R16_Q1, run by run :: sha256 588e935996d1006a
\newcommand{\MlstFtJcOneRejHqqlCoarseMassOverCoarsePr}{equal within 9\%}  % experiments/FIGS/data/v2/curves/full_budget_long_leg2.json :: cells.<run>.N1000000.s1.jc1.per_class.label_Hqql: rejection of R16_Q1_MASS over R16_Q1, run by run :: sha256 588e935996d1006a :: prose twin: two significant figures, direction in words
\newcommand{\MlstFtJcOneRejHqqlCoarseMassOverCoarsePrMax}{equal within 9\%}  % experiments/FIGS/data/v2/curves/full_budget_long_leg2.json :: cells.<run>.N1000000.s1.jc1.per_class.label_Hqql: rejection of R16_Q1_MASS over R16_Q1, run by run :: sha256 588e935996d1006a :: prose twin of the run farther from 1
\newcommand{\MlstFtJcOneRejHqqlRandOverFull}{0.94\,\ensuremath{\pm}\,0.12}  % experiments/FIGS/data/v2/curves/full_budget_long_leg2.json :: per_class.label_Hqql: rejection of each random partition over the 188-class run of its seed :: sha256 588e935996d1006a
\newcommand{\MlstFtJcOneRejHqqlRandOverFullPr}{equal within 20\%}  % experiments/FIGS/data/v2/curves/full_budget_long_leg2.json :: per_class.label_Hqql: rejection of each random partition over the 188-class run of its seed :: sha256 588e935996d1006a :: prose twin: two significant figures, direction in words
\newcommand{\MlstFtJcOneRejHqqlRandOverFullRange}{0.79--1.1}  % experiments/FIGS/data/v2/curves/full_budget_long_leg2.json :: per_class.label_Hqql: smallest and largest of the partitions' ratios :: sha256 588e935996d1006a
\newcommand{\MlstFtJcOneRejTbqqResOverFull}{1.1\,\ensuremath{\pm}\,0.4}  % experiments/FIGS/data/v2/curves/full_budget_long_leg2.json :: cells.<run>.N1000000.s1.jc1.per_class.label_Tbqq: rejection of L162 over L188, run by run :: sha256 588e935996d1006a
\newcommand{\MlstFtJcOneRejTbqqResOverFullPr}{equal within 50\%}  % experiments/FIGS/data/v2/curves/full_budget_long_leg2.json :: cells.<run>.N1000000.s1.jc1.per_class.label_Tbqq: rejection of L162 over L188, run by run :: sha256 588e935996d1006a :: prose twin: two significant figures, direction in words
\newcommand{\MlstFtJcOneRejTbqqMidOverFull}{1.10\,\ensuremath{\pm}\,0.05}  % experiments/FIGS/data/v2/curves/full_budget_long_leg2.json :: cells.<run>.N1000000.s1.jc1.per_class.label_Tbqq: rejection of R42_Q1 over L188, run by run :: sha256 588e935996d1006a
\newcommand{\MlstFtJcOneRejTbqqMidOverFullPr}{equal within 10\%}  % experiments/FIGS/data/v2/curves/full_budget_long_leg2.json :: cells.<run>.N1000000.s1.jc1.per_class.label_Tbqq: rejection of R42_Q1 over L188, run by run :: sha256 588e935996d1006a :: prose twin: two significant figures, direction in words
\newcommand{\MlstFtJcOneRejTbqqCoarseOverFull}{0.85\,\ensuremath{\pm}\,0.32}  % experiments/FIGS/data/v2/curves/full_budget_long_leg2.json :: cells.<run>.N1000000.s1.jc1.per_class.label_Tbqq: rejection of R16_Q1 over L188, run by run :: sha256 588e935996d1006a
\newcommand{\MlstFtJcOneRejTbqqCoarseOverFullPr}{equal within 40\%}  % experiments/FIGS/data/v2/curves/full_budget_long_leg2.json :: cells.<run>.N1000000.s1.jc1.per_class.label_Tbqq: rejection of R16_Q1 over L188, run by run :: sha256 588e935996d1006a :: prose twin: two significant figures, direction in words
\newcommand{\MlstFtJcOneRejTbqqResMassOverRes}{0.89 and 1.21}  % experiments/FIGS/data/v2/curves/full_budget_long_leg2.json :: cells.<run>.N1000000.s1.jc1.per_class.label_Tbqq: rejection of L162_MASS over L162, run by run :: sha256 588e935996d1006a
\newcommand{\MlstFtJcOneRejTbqqResMassOverResPr}{equal within 30\%}  % experiments/FIGS/data/v2/curves/full_budget_long_leg2.json :: cells.<run>.N1000000.s1.jc1.per_class.label_Tbqq: rejection of L162_MASS over L162, run by run :: sha256 588e935996d1006a :: prose twin: two significant figures, direction in words
\newcommand{\MlstFtJcOneRejTbqqResMassOverResPrMax}{equal within 30\%}  % experiments/FIGS/data/v2/curves/full_budget_long_leg2.json :: cells.<run>.N1000000.s1.jc1.per_class.label_Tbqq: rejection of L162_MASS over L162, run by run :: sha256 588e935996d1006a :: prose twin of the run farther from 1
\newcommand{\MlstFtJcOneRejTbqqCoarseMassOverCoarse}{1.00 and 1.79}  % experiments/FIGS/data/v2/curves/full_budget_long_leg2.json :: cells.<run>.N1000000.s1.jc1.per_class.label_Tbqq: rejection of R16_Q1_MASS over R16_Q1, run by run :: sha256 588e935996d1006a
\newcommand{\MlstFtJcOneRejTbqqCoarseMassOverCoarsePr}{equal within 80\%}  % experiments/FIGS/data/v2/curves/full_budget_long_leg2.json :: cells.<run>.N1000000.s1.jc1.per_class.label_Tbqq: rejection of R16_Q1_MASS over R16_Q1, run by run :: sha256 588e935996d1006a :: prose twin: two significant figures, direction in words
\newcommand{\MlstFtJcOneRejTbqqCoarseMassOverCoarsePrMax}{equal within 80\%}  % experiments/FIGS/data/v2/curves/full_budget_long_leg2.json :: cells.<run>.N1000000.s1.jc1.per_class.label_Tbqq: rejection of R16_Q1_MASS over R16_Q1, run by run :: sha256 588e935996d1006a :: prose twin of the run farther from 1
\newcommand{\MlstFtJcOneRejTbqqRandOverFull}{0.79\,\ensuremath{\pm}\,0.15}  % experiments/FIGS/data/v2/curves/full_budget_long_leg2.json :: per_class.label_Tbqq: rejection of each random partition over the 188-class run of its seed :: sha256 588e935996d1006a
\newcommand{\MlstFtJcOneRejTbqqRandOverFullPr}{equal within 30\%}  % experiments/FIGS/data/v2/curves/full_budget_long_leg2.json :: per_class.label_Tbqq: rejection of each random partition over the 188-class run of its seed :: sha256 588e935996d1006a :: prose twin: two significant figures, direction in words
\newcommand{\MlstFtJcOneRejTbqqRandOverFullRange}{0.64--1.0}  % experiments/FIGS/data/v2/curves/full_budget_long_leg2.json :: per_class.label_Tbqq: smallest and largest of the partitions' ratios :: sha256 588e935996d1006a
\newcommand{\MlstFtJcOneRejTblResOverFull}{0.85\,\ensuremath{\pm}\,0.17}  % experiments/FIGS/data/v2/curves/full_budget_long_leg2.json :: cells.<run>.N1000000.s1.jc1.per_class.label_Tbl: rejection of L162 over L188, run by run :: sha256 588e935996d1006a
\newcommand{\MlstFtJcOneRejTblResOverFullPr}{equal within 20\%}  % experiments/FIGS/data/v2/curves/full_budget_long_leg2.json :: cells.<run>.N1000000.s1.jc1.per_class.label_Tbl: rejection of L162 over L188, run by run :: sha256 588e935996d1006a :: prose twin: two significant figures, direction in words
\newcommand{\MlstFtJcOneRejTblMidOverFull}{0.91\,\ensuremath{\pm}\,0.15}  % experiments/FIGS/data/v2/curves/full_budget_long_leg2.json :: cells.<run>.N1000000.s1.jc1.per_class.label_Tbl: rejection of R42_Q1 over L188, run by run :: sha256 588e935996d1006a
\newcommand{\MlstFtJcOneRejTblMidOverFullPr}{equal within 20\%}  % experiments/FIGS/data/v2/curves/full_budget_long_leg2.json :: cells.<run>.N1000000.s1.jc1.per_class.label_Tbl: rejection of R42_Q1 over L188, run by run :: sha256 588e935996d1006a :: prose twin: two significant figures, direction in words
\newcommand{\MlstFtJcOneRejTblCoarseOverFull}{0.94\,\ensuremath{\pm}\,0.28}  % experiments/FIGS/data/v2/curves/full_budget_long_leg2.json :: cells.<run>.N1000000.s1.jc1.per_class.label_Tbl: rejection of R16_Q1 over L188, run by run :: sha256 588e935996d1006a
\newcommand{\MlstFtJcOneRejTblCoarseOverFullPr}{equal within 30\%}  % experiments/FIGS/data/v2/curves/full_budget_long_leg2.json :: cells.<run>.N1000000.s1.jc1.per_class.label_Tbl: rejection of R16_Q1 over L188, run by run :: sha256 588e935996d1006a :: prose twin: two significant figures, direction in words
\newcommand{\MlstFtJcOneRejTblResMassOverRes}{1.04 and 0.93}  % experiments/FIGS/data/v2/curves/full_budget_long_leg2.json :: cells.<run>.N1000000.s1.jc1.per_class.label_Tbl: rejection of L162_MASS over L162, run by run :: sha256 588e935996d1006a
\newcommand{\MlstFtJcOneRejTblResMassOverResPr}{equal within 7\%}  % experiments/FIGS/data/v2/curves/full_budget_long_leg2.json :: cells.<run>.N1000000.s1.jc1.per_class.label_Tbl: rejection of L162_MASS over L162, run by run :: sha256 588e935996d1006a :: prose twin: two significant figures, direction in words
\newcommand{\MlstFtJcOneRejTblResMassOverResPrMax}{equal within 7\%}  % experiments/FIGS/data/v2/curves/full_budget_long_leg2.json :: cells.<run>.N1000000.s1.jc1.per_class.label_Tbl: rejection of L162_MASS over L162, run by run :: sha256 588e935996d1006a :: prose twin of the run farther from 1
\newcommand{\MlstFtJcOneRejTblCoarseMassOverCoarse}{0.84 and 1.22}  % experiments/FIGS/data/v2/curves/full_budget_long_leg2.json :: cells.<run>.N1000000.s1.jc1.per_class.label_Tbl: rejection of R16_Q1_MASS over R16_Q1, run by run :: sha256 588e935996d1006a
\newcommand{\MlstFtJcOneRejTblCoarseMassOverCoarsePr}{equal within 30\%}  % experiments/FIGS/data/v2/curves/full_budget_long_leg2.json :: cells.<run>.N1000000.s1.jc1.per_class.label_Tbl: rejection of R16_Q1_MASS over R16_Q1, run by run :: sha256 588e935996d1006a :: prose twin: two significant figures, direction in words
\newcommand{\MlstFtJcOneRejTblCoarseMassOverCoarsePrMax}{equal within 30\%}  % experiments/FIGS/data/v2/curves/full_budget_long_leg2.json :: cells.<run>.N1000000.s1.jc1.per_class.label_Tbl: rejection of R16_Q1_MASS over R16_Q1, run by run :: sha256 588e935996d1006a :: prose twin of the run farther from 1
\newcommand{\MlstFtJcOneRejTblRandOverFull}{0.88\,\ensuremath{\pm}\,0.26}  % experiments/FIGS/data/v2/curves/full_budget_long_leg2.json :: per_class.label_Tbl: rejection of each random partition over the 188-class run of its seed :: sha256 588e935996d1006a
\newcommand{\MlstFtJcOneRejTblRandOverFullPr}{equal within 30\%}  % experiments/FIGS/data/v2/curves/full_budget_long_leg2.json :: per_class.label_Tbl: rejection of each random partition over the 188-class run of its seed :: sha256 588e935996d1006a :: prose twin: two significant figures, direction in words
\newcommand{\MlstFtJcOneRejTblRandOverFullRange}{0.61--1.2}  % experiments/FIGS/data/v2/curves/full_budget_long_leg2.json :: per_class.label_Tbl: smallest and largest of the partitions' ratios :: sha256 588e935996d1006a
\newcommand{\MlstFtJcOneRejWqqResOverFull}{1.039\,\ensuremath{\pm}\,0.035}  % experiments/FIGS/data/v2/curves/full_budget_long_leg2.json :: cells.<run>.N1000000.s1.jc1.per_class.label_Wqq: rejection of L162 over L188, run by run :: sha256 588e935996d1006a
\newcommand{\MlstFtJcOneRejWqqResOverFullPr}{equal within 4\%}  % experiments/FIGS/data/v2/curves/full_budget_long_leg2.json :: cells.<run>.N1000000.s1.jc1.per_class.label_Wqq: rejection of L162 over L188, run by run :: sha256 588e935996d1006a :: prose twin: two significant figures, direction in words
\newcommand{\MlstFtJcOneRejWqqMidOverFull}{1.005\,\ensuremath{\pm}\,0.032}  % experiments/FIGS/data/v2/curves/full_budget_long_leg2.json :: cells.<run>.N1000000.s1.jc1.per_class.label_Wqq: rejection of R42_Q1 over L188, run by run :: sha256 588e935996d1006a
\newcommand{\MlstFtJcOneRejWqqMidOverFullPr}{equal within 4\%}  % experiments/FIGS/data/v2/curves/full_budget_long_leg2.json :: cells.<run>.N1000000.s1.jc1.per_class.label_Wqq: rejection of R42_Q1 over L188, run by run :: sha256 588e935996d1006a :: prose twin: two significant figures, direction in words
\newcommand{\MlstFtJcOneRejWqqCoarseOverFull}{1.01\,\ensuremath{\pm}\,0.07}  % experiments/FIGS/data/v2/curves/full_budget_long_leg2.json :: cells.<run>.N1000000.s1.jc1.per_class.label_Wqq: rejection of R16_Q1 over L188, run by run :: sha256 588e935996d1006a
\newcommand{\MlstFtJcOneRejWqqCoarseOverFullPr}{equal within 8\%}  % experiments/FIGS/data/v2/curves/full_budget_long_leg2.json :: cells.<run>.N1000000.s1.jc1.per_class.label_Wqq: rejection of R16_Q1 over L188, run by run :: sha256 588e935996d1006a :: prose twin: two significant figures, direction in words
\newcommand{\MlstFtJcOneRejWqqResMassOverRes}{0.996 and 0.964}  % experiments/FIGS/data/v2/curves/full_budget_long_leg2.json :: cells.<run>.N1000000.s1.jc1.per_class.label_Wqq: rejection of L162_MASS over L162, run by run :: sha256 588e935996d1006a
\newcommand{\MlstFtJcOneRejWqqResMassOverResPr}{equal within 4\%}  % experiments/FIGS/data/v2/curves/full_budget_long_leg2.json :: cells.<run>.N1000000.s1.jc1.per_class.label_Wqq: rejection of L162_MASS over L162, run by run :: sha256 588e935996d1006a :: prose twin: two significant figures, direction in words
\newcommand{\MlstFtJcOneRejWqqResMassOverResPrMax}{equal within 4\%}  % experiments/FIGS/data/v2/curves/full_budget_long_leg2.json :: cells.<run>.N1000000.s1.jc1.per_class.label_Wqq: rejection of L162_MASS over L162, run by run :: sha256 588e935996d1006a :: prose twin of the run farther from 1
\newcommand{\MlstFtJcOneRejWqqCoarseMassOverCoarse}{1.025 and 0.926}  % experiments/FIGS/data/v2/curves/full_budget_long_leg2.json :: cells.<run>.N1000000.s1.jc1.per_class.label_Wqq: rejection of R16_Q1_MASS over R16_Q1, run by run :: sha256 588e935996d1006a
\newcommand{\MlstFtJcOneRejWqqCoarseMassOverCoarsePr}{equal within 8\%}  % experiments/FIGS/data/v2/curves/full_budget_long_leg2.json :: cells.<run>.N1000000.s1.jc1.per_class.label_Wqq: rejection of R16_Q1_MASS over R16_Q1, run by run :: sha256 588e935996d1006a :: prose twin: two significant figures, direction in words
\newcommand{\MlstFtJcOneRejWqqCoarseMassOverCoarsePrMax}{equal within 8\%}  % experiments/FIGS/data/v2/curves/full_budget_long_leg2.json :: cells.<run>.N1000000.s1.jc1.per_class.label_Wqq: rejection of R16_Q1_MASS over R16_Q1, run by run :: sha256 588e935996d1006a :: prose twin of the run farther from 1
\newcommand{\MlstFtJcOneRejWqqRandOverFull}{1.02\,\ensuremath{\pm}\,0.07}  % experiments/FIGS/data/v2/curves/full_budget_long_leg2.json :: per_class.label_Wqq: rejection of each random partition over the 188-class run of its seed :: sha256 588e935996d1006a
\newcommand{\MlstFtJcOneRejWqqRandOverFullPr}{equal within 7\%}  % experiments/FIGS/data/v2/curves/full_budget_long_leg2.json :: per_class.label_Wqq: rejection of each random partition over the 188-class run of its seed :: sha256 588e935996d1006a :: prose twin: two significant figures, direction in words
\newcommand{\MlstFtJcOneRejWqqRandOverFullRange}{0.93--1.1}  % experiments/FIGS/data/v2/curves/full_budget_long_leg2.json :: per_class.label_Wqq: smallest and largest of the partitions' ratios :: sha256 588e935996d1006a
\newcommand{\MlstFtJcOneRejZqqResOverFull}{0.987\,\ensuremath{\pm}\,0.019}  % experiments/FIGS/data/v2/curves/full_budget_long_leg2.json :: cells.<run>.N1000000.s1.jc1.per_class.label_Zqq: rejection of L162 over L188, run by run :: sha256 588e935996d1006a
\newcommand{\MlstFtJcOneRejZqqResOverFullPr}{equal within 2\%}  % experiments/FIGS/data/v2/curves/full_budget_long_leg2.json :: cells.<run>.N1000000.s1.jc1.per_class.label_Zqq: rejection of L162 over L188, run by run :: sha256 588e935996d1006a :: prose twin: two significant figures, direction in words
\newcommand{\MlstFtJcOneRejZqqMidOverFull}{0.984\,\ensuremath{\pm}\,0.010}  % experiments/FIGS/data/v2/curves/full_budget_long_leg2.json :: cells.<run>.N1000000.s1.jc1.per_class.label_Zqq: rejection of R42_Q1 over L188, run by run :: sha256 588e935996d1006a
\newcommand{\MlstFtJcOneRejZqqMidOverFullPr}{equal within 2\%}  % experiments/FIGS/data/v2/curves/full_budget_long_leg2.json :: cells.<run>.N1000000.s1.jc1.per_class.label_Zqq: rejection of R42_Q1 over L188, run by run :: sha256 588e935996d1006a :: prose twin: two significant figures, direction in words
\newcommand{\MlstFtJcOneRejZqqCoarseOverFull}{0.91\,\ensuremath{\pm}\,0.05}  % experiments/FIGS/data/v2/curves/full_budget_long_leg2.json :: cells.<run>.N1000000.s1.jc1.per_class.label_Zqq: rejection of R16_Q1 over L188, run by run :: sha256 588e935996d1006a
\newcommand{\MlstFtJcOneRejZqqCoarseOverFullPr}{equal within 9\%}  % experiments/FIGS/data/v2/curves/full_budget_long_leg2.json :: cells.<run>.N1000000.s1.jc1.per_class.label_Zqq: rejection of R16_Q1 over L188, run by run :: sha256 588e935996d1006a :: prose twin: two significant figures, direction in words
\newcommand{\MlstFtJcOneRejZqqResMassOverRes}{0.9803 and 0.9820}  % experiments/FIGS/data/v2/curves/full_budget_long_leg2.json :: cells.<run>.N1000000.s1.jc1.per_class.label_Zqq: rejection of L162_MASS over L162, run by run :: sha256 588e935996d1006a
\newcommand{\MlstFtJcOneRejZqqResMassOverResPr}{equal within 2\%}  % experiments/FIGS/data/v2/curves/full_budget_long_leg2.json :: cells.<run>.N1000000.s1.jc1.per_class.label_Zqq: rejection of L162_MASS over L162, run by run :: sha256 588e935996d1006a :: prose twin: two significant figures, direction in words
\newcommand{\MlstFtJcOneRejZqqResMassOverResPrMax}{equal within 2\%}  % experiments/FIGS/data/v2/curves/full_budget_long_leg2.json :: cells.<run>.N1000000.s1.jc1.per_class.label_Zqq: rejection of L162_MASS over L162, run by run :: sha256 588e935996d1006a :: prose twin of the run farther from 1
\newcommand{\MlstFtJcOneRejZqqCoarseMassOverCoarse}{1.06 and 0.94}  % experiments/FIGS/data/v2/curves/full_budget_long_leg2.json :: cells.<run>.N1000000.s1.jc1.per_class.label_Zqq: rejection of R16_Q1_MASS over R16_Q1, run by run :: sha256 588e935996d1006a
\newcommand{\MlstFtJcOneRejZqqCoarseMassOverCoarsePr}{equal within 7\%}  % experiments/FIGS/data/v2/curves/full_budget_long_leg2.json :: cells.<run>.N1000000.s1.jc1.per_class.label_Zqq: rejection of R16_Q1_MASS over R16_Q1, run by run :: sha256 588e935996d1006a :: prose twin: two significant figures, direction in words
\newcommand{\MlstFtJcOneRejZqqCoarseMassOverCoarsePrMax}{equal within 7\%}  % experiments/FIGS/data/v2/curves/full_budget_long_leg2.json :: cells.<run>.N1000000.s1.jc1.per_class.label_Zqq: rejection of R16_Q1_MASS over R16_Q1, run by run :: sha256 588e935996d1006a :: prose twin of the run farther from 1
\newcommand{\MlstFtJcOneRejZqqRandOverFull}{0.952\,\ensuremath{\pm}\,0.028}  % experiments/FIGS/data/v2/curves/full_budget_long_leg2.json :: per_class.label_Zqq: rejection of each random partition over the 188-class run of its seed :: sha256 588e935996d1006a
\newcommand{\MlstFtJcOneRejZqqRandOverFullPr}{equal within 5\%}  % experiments/FIGS/data/v2/curves/full_budget_long_leg2.json :: per_class.label_Zqq: rejection of each random partition over the 188-class run of its seed :: sha256 588e935996d1006a :: prose twin: two significant figures, direction in words
\newcommand{\MlstFtJcOneRejZqqRandOverFullRange}{0.93--0.99}  % experiments/FIGS/data/v2/curves/full_budget_long_leg2.json :: per_class.label_Zqq: smallest and largest of the partitions' ratios :: sha256 588e935996d1006a
\newcommand{\MlstFtJcOneResNRejDrop}{zero}  % experiments/FIGS/data/v2/curves/full_budget_long_leg2.json :: per_class rejections :: sha256 588e935996d1006a :: classes whose paired rejection ratio to the 188-class model departs from 1 by more than twice its SD and its counting uncertainty, columns with fewer than JC1_MIN_BKG passing jets excluded, ratio below 1 (count)
\newcommand{\MlstFtJcOneResNRejDiff}{zero}  % experiments/FIGS/data/v2/curves/full_budget_long_leg2.json :: per_class rejections :: sha256 588e935996d1006a :: classes whose paired rejection ratio to the 188-class model departs from 1 by more than twice its SD and its counting uncertainty, columns with fewer than JC1_MIN_BKG passing jets excluded (count)
\newcommand{\MlstFtJcOneResRejDropClasses}{none}  % experiments/FIGS/data/v2/curves/full_budget_long_leg2.json :: per_class rejections :: sha256 588e935996d1006a :: the classes of ...NRejDrop, in ParT's column order
\newcommand{\MlstFtJcOneMidNRejDrop}{zero}  % experiments/FIGS/data/v2/curves/full_budget_long_leg2.json :: per_class rejections :: sha256 588e935996d1006a :: classes whose paired rejection ratio to the 188-class model departs from 1 by more than twice its SD and its counting uncertainty, columns with fewer than JC1_MIN_BKG passing jets excluded, ratio below 1 (count)
\newcommand{\MlstFtJcOneMidNRejDiff}{zero}  % experiments/FIGS/data/v2/curves/full_budget_long_leg2.json :: per_class rejections :: sha256 588e935996d1006a :: classes whose paired rejection ratio to the 188-class model departs from 1 by more than twice its SD and its counting uncertainty, columns with fewer than JC1_MIN_BKG passing jets excluded (count)
\newcommand{\MlstFtJcOneMidRejDropClasses}{none}  % experiments/FIGS/data/v2/curves/full_budget_long_leg2.json :: per_class rejections :: sha256 588e935996d1006a :: the classes of ...NRejDrop, in ParT's column order
\newcommand{\MlstFtJcOneCoarseNRejDrop}{zero}  % experiments/FIGS/data/v2/curves/full_budget_long_leg2.json :: per_class rejections :: sha256 588e935996d1006a :: classes whose paired rejection ratio to the 188-class model departs from 1 by more than twice its SD and its counting uncertainty, columns with fewer than JC1_MIN_BKG passing jets excluded, ratio below 1 (count)
\newcommand{\MlstFtJcOneCoarseNRejDiff}{zero}  % experiments/FIGS/data/v2/curves/full_budget_long_leg2.json :: per_class rejections :: sha256 588e935996d1006a :: classes whose paired rejection ratio to the 188-class model departs from 1 by more than twice its SD and its counting uncertainty, columns with fewer than JC1_MIN_BKG passing jets excluded (count)
\newcommand{\MlstFtJcOneCoarseRejDropClasses}{none}  % experiments/FIGS/data/v2/curves/full_budget_long_leg2.json :: per_class rejections :: sha256 588e935996d1006a :: the classes of ...NRejDrop, in ParT's column order
\newcommand{\MlstFtJcOneAccRandMinusCoarse}{+0.002}  % experiments/FIGS/data/v2/curves/full_budget_long_leg2.json :: cells.*.N1000000.s1.jc1.accuracy :: sha256 588e935996d1006a :: mean over the random partitions minus the mean over the 17-class runs
\newcommand{\MlstFtJcOneAccFullBoot}{0.847}  % experiments/FIGS/data/v2/budget/jc1_acc_boot_N1000000.json :: diffs.r16q1-l188.accuracy_ref (mean) :: sha256 7cccb7b0ff900fb8
\newcommand{\MlstFtJcOneAccCoarseBoot}{0.842}  % experiments/FIGS/data/v2/budget/jc1_acc_boot_N1000000.json :: diffs.r16q1-l188.accuracy_config (mean) :: sha256 7cccb7b0ff900fb8
\newcommand{\MlstFtJcOneErrResOverFullBoot}{consistent}  % experiments/FIGS/data/v2/budget/jc1_acc_boot_N1000000.json :: diffs.l162-l188.error_ratio :: sha256 7cccb7b0ff900fb8 :: delta and sigma = hypot(run SD, bootstrap SD) in per cent; 'consistent' when |delta| <= 3 sigma
\newcommand{\MlstFtJcOneErrResOverFullZ}{0.6}  % experiments/FIGS/data/v2/budget/jc1_acc_boot_N1000000.json :: diffs.l162-l188.error_ratio.z :: sha256 7cccb7b0ff900fb8 :: |delta| / sigma
\newcommand{\MlstFtJcOneErrMidOverFullBoot}{consistent}  % experiments/FIGS/data/v2/budget/jc1_acc_boot_N1000000.json :: diffs.r42q1-l188.error_ratio :: sha256 7cccb7b0ff900fb8 :: delta and sigma = hypot(run SD, bootstrap SD) in per cent; 'consistent' when |delta| <= 3 sigma
\newcommand{\MlstFtJcOneErrMidOverFullZ}{1.4}  % experiments/FIGS/data/v2/budget/jc1_acc_boot_N1000000.json :: diffs.r42q1-l188.error_ratio.z :: sha256 7cccb7b0ff900fb8 :: |delta| / sigma
\newcommand{\MlstFtJcOneErrCoarseOverFullBoot}{(3.0\,\ensuremath{\pm}\,0.6)\% more errors}  % experiments/FIGS/data/v2/budget/jc1_acc_boot_N1000000.json :: diffs.r16q1-l188.error_ratio :: sha256 7cccb7b0ff900fb8 :: delta and sigma = hypot(run SD, bootstrap SD) in per cent; 'consistent' when |delta| <= 3 sigma
\newcommand{\MlstFtJcOneErrCoarseOverFullZ}{4.8}  % experiments/FIGS/data/v2/budget/jc1_acc_boot_N1000000.json :: diffs.r16q1-l188.error_ratio.z :: sha256 7cccb7b0ff900fb8 :: |delta| / sigma
\newcommand{\MlstFtJcOneBootReplicas}{200}  % experiments/FIGS/data/v2/budget/jc1_acc_boot_N1000000.json :: replicas :: sha256 7cccb7b0ff900fb8
\newcommand{\MlstFtJcOneBootJets}{2{,}000{,}000}  % experiments/FIGS/data/v2/budget/jc1_acc_boot_N1000000.json :: n_test_jets :: sha256 7cccb7b0ff900fb8
\newcommand{\ComputeTotalGpuHours}{3{,}435}  % experiments/FIGS/data/v2/index.json :: GPU-hours of every reported pretraining run, training and validation passes :: sha256 c6c4138ee5c86423
\newcommand{\MlstComputeTotalGpuHours}{3{,}510}  % experiments/FIGS/data/v2/index.json :: GPU-hours of every pretraining run the paper reports, the kinematic-only model included :: sha256 c6c4138ee5c86423
\newcommand{\ComputeMainGpu}{V100}  % experiments/FIGS/data/v2/index.json :: GPU type with the most hours over the four label sets :: sha256 c6c4138ee5c86423
\newcommand{\ComputeMainGpuShareMin}{56\%}  % experiments/FIGS/data/v2/index.json :: lowest share of a label set's GPU-hours on that type :: sha256 c6c4138ee5c86423
\newcommand{\ComputeMainGpuShareMax}{71\%}  % experiments/FIGS/data/v2/index.json :: highest share of a label set's GPU-hours on that type :: sha256 c6c4138ee5c86423
\newcommand{\ComputeRunsDone}{27}  % experiments/FIGS/data/v2/index.json :: paper runs with a compute record (count) :: sha256 c6c4138ee5c86423
\newcommand{\ComputeGpuTypes}{L4, RTX 3090, RTX 4090 and V100}  % experiments/FIGS/data/v2/index.json :: GPU types any finished paper run trained on :: sha256 c6c4138ee5c86423
\newcommand{\MlstComputeGpuTypes}{L4, RTX 3090, RTX 4090 and V100}  % experiments/FIGS/data/v2/index.json :: GPU types of the pretraining runs the paper reports :: sha256 c6c4138ee5c86423
\newcommand{\MlstFtJcOneErrCoarseOverFullBootEFour}{\pending{the JetClass bootstrap at 1e4 (after the equal-stopping read-out)}}  % experiments/FIGS/data/v2/index.json :: not computed yet :: sha256 c6c4138ee5c86423 :: waits on the JetClass bootstrap at 1e4 (after the equal-stopping read-out)
\newcommand{\MlstLpBbQgZeroAccDropFull}{\pending{pair probes with the quark/gluon inputs zeroed (open item O10, not run)}}  % experiments/FIGS/data/v2/index.json :: not computed yet :: sha256 c6c4138ee5c86423 :: waits on pair probes with the quark/gluon inputs zeroed (open item O10, not run)
\newcommand{\MlstLpBbQgZeroAccDropCoarse}{\pending{pair probes with the quark/gluon inputs zeroed (open item O10, not run)}}  % experiments/FIGS/data/v2/index.json :: not computed yet :: sha256 c6c4138ee5c86423 :: waits on pair probes with the quark/gluon inputs zeroed (open item O10, not run)
\newcommand{\XbsRthirtyFullOverScratch}{\pending{X->bs fits and from-scratch runs}}  % experiments/FIGS/data/v2/index.json :: not computed yet :: sha256 c6c4138ee5c86423 :: waits on X->bs fits and from-scratch runs
\newcommand{\XbsRthirtyFullOverScratchPr}{\pending{X->bs fits and from-scratch runs}}  % experiments/FIGS/data/v2/index.json :: not computed yet :: sha256 c6c4138ee5c86423 :: waits on X->bs fits and from-scratch runs
\newcommand{\XbsRthirtyCoarseOverFull}{\pending{X->bs fits}}  % experiments/FIGS/data/v2/index.json :: not computed yet :: sha256 c6c4138ee5c86423 :: waits on X->bs fits
\newcommand{\XbsRthirtyCoarseOverFullPr}{\pending{X->bs fits}}  % experiments/FIGS/data/v2/index.json :: not computed yet :: sha256 c6c4138ee5c86423 :: waits on X->bs fits
\newcommand{\XbsRthirtyNodeOverTransferFull}{\pending{X->bs fits}}  % experiments/FIGS/data/v2/index.json :: not computed yet :: sha256 c6c4138ee5c86423 :: waits on X->bs fits
\newcommand{\XbsRthirtyNodeOverTransferFullPr}{\pending{X->bs fits}}  % experiments/FIGS/data/v2/index.json :: not computed yet :: sha256 c6c4138ee5c86423 :: waits on X->bs fits
\newcommand{\XbbRthirtyDirectOverTransferFull}{\pending{X->bb fits}}  % experiments/FIGS/data/v2/index.json :: not computed yet :: sha256 c6c4138ee5c86423 :: waits on X->bb fits
\newcommand{\XbbRthirtyDirectOverTransferFullPr}{\pending{X->bb fits}}  % experiments/FIGS/data/v2/index.json :: not computed yet :: sha256 c6c4138ee5c86423 :: waits on X->bb fits
\newcommand{\XbbRthirtyTransferCoarseOverFull}{\pending{X->bb fits}}  % experiments/FIGS/data/v2/index.json :: not computed yet :: sha256 c6c4138ee5c86423 :: waits on X->bb fits
\newcommand{\XbbRthirtyTransferCoarseOverFullPr}{\pending{X->bb fits}}  % experiments/FIGS/data/v2/index.json :: not computed yet :: sha256 c6c4138ee5c86423 :: waits on X->bb fits
             % the counts below are its macros
% Visual language: Sophon's Fig. 1 (arXiv:2405.12972): the document's serif font, light fills,
% rounded boxes, light grey panels behind the stages (PI 2026-11-02). The text is set in the
% document's font at the printed sizes 8.5-10 pt (MLST 8-12 pt), never scaled.
% Text width: 15 cm (MLST full width). The picture is 15.0 cm wide and is drawn at its final
% size, never scaled, so the font sizes hold.
% Optional: \def\FoneDropXbx{} before % Figure 1 (F1 of docs/FIGURE_PLAN.md): the design of the study, in TikZ.
% % Figure 1 (F1 of docs/FIGURE_PLAN.md): the design of the study, in TikZ.
% \input{fig_f1} from main.tex; this file holds one figure environment and nothing else.
%
% WHAT main.tex MUST LOAD (this file loads nothing itself):
%   \usepackage{tikz}
%   \usetikzlibrary{arrows.meta}          % the only TikZ library used (arrow tips)
%   \definecolor{c188}{HTML}{1F77B4}      % the configuration colours of experiments/FIGS/
%   \definecolor{c162}{HTML}{FF7F0E}      % mlst_style.py (matplotlib's tab10 in order, as in
%   \definecolor{c43}{HTML}{2CA02C}       % Sophon's figures); the variants are drawn in their
%   \definecolor{c17}{HTML}{D62728}       % label set's colour with a dotted border, as the
%   \definecolor{crand}{HTML}{17BECF}     % plotted figures draw variants dotted
%   \definecolor{cflav}{HTML}{9467BD}
%   \input{results_generated}             % the counts below are its macros
% Visual language: Sophon's Fig. 1 (arXiv:2405.12972): the document's serif font, light fills,
% rounded boxes, light grey panels behind the stages (PI 2026-11-02). The text is set in the
% document's font at the printed sizes 8.5-10 pt (MLST 8-12 pt), never scaled.
% Text width: 15 cm (MLST full width). The picture is 15.0 cm wide and is drawn at its final
% size, never scaled, so the font sizes hold.
% Optional: \def\FoneDropXbx{} before \input{fig_f1} removes X->bb/X->bs from panel (c)
% (PI 2026-11-02: if the samples are not staged by 12 Nov, the X->bb/X->bs section goes).
%
% Every class box is drawn from the label maps (checked by hand, 2026-11-02):
%   (a) configs/labelmaps/rung_label_maps.v1.csv, columns L188, L162, R42_Q1, R16_Q1. Rows: the ten
%       decays X->QQ/gg in the order bb bc bq | cc cs cq | ss sq qq gg (the 43-class b / c / light
%       groups are then contiguous), then ee mm | tauh-taue tauh-taum, then the 27 QCD classes.
%   (b) column: configs/labelmaps/rand_label_map.v2.csv, RAND2_p1 on the same rows: groups
%       0 0 0 0 0 9 9 9 0 3 | 11 11 | 15 4 | QCD, lettered A..F in order of first appearance.
%       strips: configs/labelmaps/flavour_pair_map.v2.csv, the 22 classes of orbit X_YY_QQQQ.
%       Order with a b quark: bbbb bbcc bbss bbqq qqqb bbcq ccbs ccbq qqbc qqbs bcsq;
%       without: cccc ccss ccqq ssss ssqq qqqq qqqc qqqs ccsq sscq qqcs.
%       FLAV_F0: all 22 in group 9. FLAV_F1: the 11 with a b in group 11, the rest 9.
%       FLAV_F1R: group 11 = bbss bbqq qqqb bbcq ccbq qqbc (with b, positions 2-5, 7, 8 from 0)
%       and ccss ccqq ssss qqqc sscq (without, positions 1-3, 6, 9); filled square = group 11.
% Geometry (cm): row height 0.35 from y = -1.18, gaps of 0.14 between the quark/gluon, lepton
% and QCD blocks, the QCD row 0.70; boxes inset 0.025 top and bottom; columns 1.45 wide with
% left edges 1.75, 3.28, 4.81, 6.34 (hierarchy) and 8.25 (random partition).
\providecommand{\Kfull}{\VocabSizeLoneeighteight}
\providecommand{\Kres}{\VocabSizeLonesixtwo}
\providecommand{\Kmid}{\VocabSizeRfourtwoQone}
\providecommand{\Kcoarse}{\VocabSizeRonesixQone}
\begin{figure}[tbp]
\centering
\ifdim\linewidth<14.99cm\typeout{fig_f1.tex: WARNING: the drawing is 15 cm wide but \string\linewidth\space
  is \the\linewidth; set a 15 cm text width (it is not scaled, to keep its 8.5 pt text)}\fi
\def\FoneText{\rmfamily\fontsize{8.5}{9.5}\selectfont}%
\def\FoneMain{\rmfamily\fontsize{9}{10.5}\selectfont}%
\def\FoneBold{\rmfamily\bfseries\fontsize{9}{10.5}\selectfont}%
\def\FoneTitle{\rmfamily\bfseries\fontsize{10}{12}\selectfont}%
\ifdefined\FoneDropXbx\def\FoneXbx{}\else\def\FoneXbx{\\$X\to b\bar b$, $X\to bs$}\fi
\begin{tikzpicture}[
  x=1cm, y=1cm,
  font=\FoneMain,
  ttl/.style={anchor=base west, font=\FoneTitle, inner sep=0pt},
  cls/.style={draw=black!55, line width=0.45pt, rounded corners=1.2pt},
  hdr/.style={draw=#1!70!black, fill=#1!22, line width=0.6pt, rounded corners=2.5pt, align=center,
              font=\FoneText, inner sep=0pt},
  inbox/.style={align=center, font=\FoneText, inner sep=0pt},
  letter/.style={font=\FoneBold, inner sep=0pt},
  rowlab/.style={anchor=east, inner sep=0pt},
  txt/.style={align=left, font=\FoneText, inner sep=0pt},
  sq/.style={draw=black!65, line width=0.4pt},
  chip/.style={draw=#1!70!black, fill=#1!22, line width=0.6pt, rounded corners=3pt,
               align=center, font=\FoneText, inner sep=0pt, anchor=north west},
  stage/.style={draw=black!45, fill=black!4, line width=0.6pt, rounded corners=3pt, align=left,
                font=\FoneText, inner sep=0.1cm},
  flow/.style={-{Latex[length=1.8mm, width=1.4mm]}, line width=0.7pt, black!75},
]
% ================================================================ (a) the class hierarchy
\node[ttl] at (0,0) {(a) Class hierarchy};
% faint bands on every other row, drawn first so the boxes sit on top
\foreach \rowtop in {-1.53,-2.23,-2.93,-3.63,-4.33,-5.17,-5.87}
  \fill[black!4] (0,\rowtop) rectangle (9.70,\rowtop-0.35);
% column headers
\node[rowlab, align=right, font=\FoneText] at (1.65,-0.68) {JetClass-II\\class};
\node[hdr=c188, minimum width=1.45cm, minimum height=0.80cm, anchor=north west] at (1.75,-0.28) {\Kfull\\classes};
\node[hdr=c162, minimum width=1.45cm, minimum height=0.80cm, anchor=north west] at (3.28,-0.28) {\Kres\\classes};
\node[hdr=c43,  minimum width=1.45cm, minimum height=0.80cm, anchor=north west] at (4.81,-0.28) {\Kmid\\classes};
\node[hdr=c17,  minimum width=1.45cm, minimum height=0.80cm, anchor=north west] at (6.34,-0.28) {\Kcoarse\\classes};
% row labels
\node[rowlab] at (1.65,-1.355) {$X\to b\bar b$};
\node[rowlab] at (1.65,-1.705) {$X\to bc$};
\node[rowlab] at (1.65,-2.055) {$X\to bq$};
\node[rowlab] at (1.65,-2.405) {$X\to c\bar c$};
\node[rowlab] at (1.65,-2.755) {$X\to cs$};
\node[rowlab] at (1.65,-3.105) {$X\to cq$};
\node[rowlab] at (1.65,-3.455) {$X\to s\bar s$};
\node[rowlab] at (1.65,-3.805) {$X\to sq$};
\node[rowlab] at (1.65,-4.155) {$X\to q\bar q$};
\node[rowlab] at (1.65,-4.505) {$X\to gg$};
\node[rowlab] at (1.65,-4.995) {$X\to ee$};
\node[rowlab] at (1.65,-5.345) {$X\to\mu\mu$};
\node[rowlab] at (1.65,-5.695) {$X\to\tau_h\tau_e$};
\node[rowlab] at (1.65,-6.045) {$X\to\tau_h\tau_\mu$};
\node[rowlab] at (1.65,-6.710) {QCD jets};
% 188 and 162 classes: every resonant decay is its own class
\foreach \rowtop in {-1.18,-1.53,-1.88,-2.23,-2.58,-2.93,-3.28,-3.63,-3.98,-4.33,-4.82,-5.17,-5.52,-5.87} {
  \draw[cls, fill=c188!15] (1.75,\rowtop-0.025) rectangle (3.20,\rowtop-0.325);
  \draw[cls, fill=c162!15] (3.28,\rowtop-0.025) rectangle (4.73,\rowtop-0.325);
}
\draw[cls, fill=c188!15] (1.75,-6.385) rectangle (3.20,-7.035);
\node[inbox] at (2.475,-6.71) {\LitJetClassTwoNQcdClasses\\classes};
\draw[cls, fill=c162!15] (3.28,-6.385) rectangle (4.73,-7.035);
\node[inbox] at (4.005,-6.71) {1 class};
% 43 classes: decay topology x {b, else c, else light quark or gluon}; e and mu merged
\draw[cls, fill=c43!15] (4.81,-1.205) rectangle (6.26,-2.205);
\node[inbox] at (5.535,-1.705) {with $b$};
\draw[cls, fill=c43!15] (4.81,-2.255) rectangle (6.26,-3.255);
\node[inbox] at (5.535,-2.755) {with $c$,\\no $b$};
\draw[cls, fill=c43!15] (4.81,-3.305) rectangle (6.26,-4.655);
\node[inbox] at (5.535,-3.980) {$u$, $d$, $s$\\or $g$};
\draw[cls, fill=c43!15] (4.81,-4.845) rectangle (6.26,-5.495);
\node[inbox] at (5.535,-5.170) {$\ell\ell$};
\draw[cls, fill=c43!15] (4.81,-5.545) rectangle (6.26,-6.195);
\node[inbox] at (5.535,-5.870) {$\tau_h\tau_\ell$};
\draw[cls, fill=c43!15] (4.81,-6.385) rectangle (6.26,-7.035);
\node[inbox] at (5.535,-6.710) {1 class};
% 17 classes: number and type of prongs, no flavor
\draw[cls, fill=c17!15] (6.34,-1.205) rectangle (7.79,-4.655);
\node[inbox] at (7.065,-2.930) {two\\quarks\\or gluons};
\draw[cls, fill=c17!15] (6.34,-4.845) rectangle (7.79,-5.495);
\node[inbox] at (7.065,-5.170) {$\ell\ell$};
\draw[cls, fill=c17!15] (6.34,-5.545) rectangle (7.79,-6.195);
\node[inbox] at (7.065,-5.870) {$\tau_h\tau_\ell$};
\draw[cls, fill=c17!15] (6.34,-6.385) rectangle (7.79,-7.035);
\node[inbox] at (7.065,-6.710) {1 class};
% what each merge removes
\node[anchor=north west, font=\FoneBold, inner sep=0pt] at (0,-7.27) {What each merge removes};
\node[txt, anchor=north west, text width=6.04cm] (m1) at (1.75,-7.65)
  {the $b$, $c$ and $s$ content of the QCD jets};
\node[txt, anchor=north west, text width=6.04cm] (m2) at ([yshift=-0.10cm]m1.south west)
  {the number of $b$ and $c$ quarks, light-quark flavor, quark against gluon, and lepton
   flavor ($e$ against $\mu$)};
\node[txt, anchor=north west, text width=6.04cm] (m3) at ([yshift=-0.10cm]m2.south west)
  {quark flavor ($b$, $c$ or light); the number and type of prongs remain};
\node[txt, anchor=north east] at ([xshift=-0.10cm]m1.north west) {\Kfull\,$\to$\,\Kres};
\node[txt, anchor=north east] at ([xshift=-0.10cm]m2.north west) {\Kres\,$\to$\,\Kmid};
\node[txt, anchor=north east] at ([xshift=-0.10cm]m3.north west) {\Kmid\,$\to$\,\Kcoarse};
% ================================================================ (b) the controls
\node[ttl] at (8.25,0) {(b) Controls, \Kcoarse{} classes each};
% the first random partition on the rows of (a); one letter = one class
\node[hdr=crand, minimum width=1.45cm, minimum height=0.80cm, anchor=north west] at (8.25,-0.28) {random\\partition};
\draw[cls, fill=crand!15] (8.25,-1.205) rectangle (9.70,-2.905);
\node[letter] at (8.975,-2.055) {A};
\draw[cls, fill=crand!15] (8.25,-2.955) rectangle (9.70,-3.955);
\node[letter] at (8.975,-3.455) {B};
\draw[cls, fill=crand!15] (8.25,-4.005) rectangle (9.70,-4.305);
\node[letter] at (8.975,-4.155) {A};
\draw[cls, fill=crand!15] (8.25,-4.355) rectangle (9.70,-4.655);
\node[letter] at (8.975,-4.505) {C};
\draw[cls, fill=crand!15] (8.25,-4.845) rectangle (9.70,-5.495);
\node[letter] at (8.975,-5.170) {D};
\draw[cls, fill=crand!15] (8.25,-5.545) rectangle (9.70,-5.845);
\node[letter] at (8.975,-5.695) {E};
\draw[cls, fill=crand!15] (8.25,-5.895) rectangle (9.70,-6.195);
\node[letter] at (8.975,-6.045) {F};
\draw[cls, fill=crand!15] (8.25,-6.385) rectangle (9.70,-7.035);
\node[inbox] at (8.975,-6.710) {1 class};
% how the random partitions are built
\node[hdr=crand, anchor=north west, text width=4.75cm, inner sep=0.1cm, align=left, font=\FoneBold]
  (rh) at (10.05,-0.28) {Random partitions (\RandVtwoNPartitions)};
\node[txt, anchor=north west, text width=4.95cm] (rb) at ([yshift=-0.12cm]rh.south west)
  {\textbullet\ groups of resonant classes drawn at random, plus one class of QCD jets\\[2pt]
   \textbullet\ group shares of training jets within \RandVtwoShareTol{} of the \Kcoarse-class groups\\[2pt]
   \textbullet\ best of \RandVtwoTrials{} draws: the fewest merges shared with the \Kcoarse-class level\\[2pt]
   \textbullet\ left: the first partition; boxes with the same letter are one class};
% how the random partitions with controlled flavor boundaries are built
\node[hdr=cflav, anchor=north west, text width=4.75cm, inner sep=0.1cm, align=left, font=\FoneBold]
  (fh) at ([yshift=-0.22cm]rb.south west) {Random partitions with controlled flavor boundaries};
\node[txt, anchor=north west, text width=4.95cm] (fn) at ([yshift=-0.12cm]fh.south west)
  {whole families of decays that differ only in quark flavor; two of them split one family,
   the \RandFlavSplitOrbitClasses{} decays \RandFlavSplitOrbit:};
\begin{scope}[shift={(fn.south west)}]
  % brackets over the decays with and without a b quark (squares at a 0.2 cm pitch)
  \node[txt, anchor=south] at (1.075,-0.45) {with $b$};
  \node[txt, anchor=south] at (3.425,-0.45) {without $b$};
  \draw[line width=0.5pt] (0,-0.54) -- (0,-0.47) -- (2.15,-0.47) -- (2.15,-0.54);
  \draw[line width=0.5pt] (2.35,-0.54) -- (2.35,-0.47) -- (4.50,-0.47) -- (4.50,-0.54);
  % no flavor boundary: all 22 in one class
  \node[txt, anchor=north west] at (0,-0.62) {no flavor boundary};
  \foreach \i in {0,...,10} {
    \draw[sq, fill=white] (0.2*\i,-1.10) rectangle ++(0.15,0.15);
    \draw[sq, fill=white] (2.35+0.2*\i,-1.10) rectangle ++(0.15,0.15);
  }
  % b-quark boundary: the decays with a b quark moved to another class
  \node[txt, anchor=north west] at (0,-1.22) {$b$-quark boundary};
  \foreach \i in {0,...,10} {
    \draw[sq, fill=cflav!70!black] (0.2*\i,-1.70) rectangle ++(0.15,0.15);
    \draw[sq, fill=white] (2.35+0.2*\i,-1.70) rectangle ++(0.15,0.15);
  }
  % random flavor boundary: the same number of decays, drawn at random
  \node[txt, anchor=north west] at (0,-1.82) {random flavor boundary};
  \foreach \i in {0,1,6,9,10}
    \draw[sq, fill=white] (0.2*\i,-2.30) rectangle ++(0.15,0.15);
  \foreach \i in {2,3,4,5,7,8}
    \draw[sq, fill=cflav!70!black] (0.2*\i,-2.30) rectangle ++(0.15,0.15);
  \foreach \i in {0,4,5,7,8,10}
    \draw[sq, fill=white] (2.35+0.2*\i,-2.30) rectangle ++(0.15,0.15);
  \foreach \i in {1,2,3,6,9}
    \draw[sq, fill=cflav!70!black] (2.35+0.2*\i,-2.30) rectangle ++(0.15,0.15);
  % key
  \draw[sq, fill=cflav!70!black] (0,-2.62) rectangle ++(0.15,0.15);
  \node[txt, anchor=west] (fkey) at (0.25,-2.545) {moved to another class};
\end{scope}
% ================================================================ (c) the pipeline
% (c) starts below the lower of the two columns above, whatever their text heights
\pgfpointanchor{m3}{south}\pgfgetlastxy{\FoneXa}{\FoneYa}%
\pgfpointanchor{fkey}{south}\pgfgetlastxy{\FoneXb}{\FoneYb}%
\pgfmathsetlengthmacro{\FoneYc}{min(\FoneYa,\FoneYb)-0.45cm}%
\coordinate (ctop) at (0pt,\FoneYc);
\begin{scope}[shift={(ctop)}]
  \node[ttl, anchor=north west] at (0,0) {(c) Pretraining and evaluation};
  % what every model of one seed shares
  \node[stage, text width=2.15cm] (seed) at (1.175,-1.96)
    {\textbf{Same for every model of one seed:} the training jets, their order and the
     initial weights, except the output layer};
  % one pretraining per label set: the configuration colours of every later figure
  \node[draw=black!45, fill=black!4, line width=0.6pt, rounded corners=3pt, minimum width=5.15cm,
        minimum height=3.08cm, anchor=north west, inner sep=0pt] (pre) at (2.81,-0.42) {};
  \node[anchor=north west, font=\FoneBold, inner sep=0pt] at (2.96,-0.52) {Pretraining on JetClass-II};
  \node[chip=c188, minimum width=1.13cm, minimum height=0.42cm] at (2.96,-0.95) {\Kfull};
  \node[chip=c162, minimum width=1.13cm, minimum height=0.42cm] at (4.20,-0.95) {\Kres};
  \node[chip=c43,  minimum width=1.13cm, minimum height=0.42cm] at (5.44,-0.95) {\Kmid};
  \node[chip=c17,  minimum width=1.13cm, minimum height=0.42cm] at (6.68,-0.95) {\Kcoarse};
  \node[chip=crand, minimum width=1.70cm, minimum height=0.70cm] at (2.96,-1.47) {random\\partitions};
  \node[chip=cflav, minimum width=3.04cm, minimum height=0.70cm] at (4.77,-1.47) {controlled flavor\\boundaries};
  \node[chip=c162, densely dotted, minimum width=1.13cm, minimum height=0.42cm] at (2.96,-2.32) {\Kres};
  \node[chip=c17,  densely dotted, minimum width=1.13cm, minimum height=0.42cm] at (4.20,-2.32) {\Kcoarse};
  \node[txt, anchor=west] at (5.44,-2.53) {with auxiliary\\mass regression};
  \node[chip=c188, densely dotted, minimum width=1.13cm, minimum height=0.42cm] at (2.96,-2.94) {\Kfull};
  \node[txt, anchor=west] at (4.20,-3.15) {kinematic inputs only};
  % evaluation
  \node[stage, text width=2.4cm] (fz) at (9.68,-0.95) {frozen features:\\linear probe,\\MLP probe};
  \node[stage, text width=2.4cm] (ft) at (9.68,-2.05) {fine-tuning\\of all weights};
  \node[stage, text width=2.4cm, draw=black, fill=white, line width=1.1pt] (sc) at (9.68,-3.15) {trained from\\scratch};
  \node[stage, text width=3.4cm, minimum height=3.2cm] (task) at (13.20,-2.00)
    {\textbf{Downstream tasks}\\JetClass-II: all classes\\JetClass-II: class pairs\\\quad(both class-balanced)\\JetClass\\
     top tagging\\quark/gluon, Pythia\\quark/gluon, Herwig\\jet mass\FoneXbx};
  \draw[flow] (seed.east) -- (seed.east -| pre.west);
  \draw[flow] (fz.west -| pre.east) -- (fz.west);
  \draw[flow] (ft.west -| pre.east) -- (ft.west);
  \draw[flow] (fz.east) -- (fz.east -| task.west);
  \draw[flow] (ft.east) -- (ft.east -| task.west);
  \draw[flow] (sc.east) -- (sc.east -| task.west);
\end{scope}
\end{tikzpicture}
\caption{Design of the study.
(a)~The class hierarchy on the decays of a heavy resonance~$X$ and on the jets from quantum
chromodynamics (QCD); $q$ is a $u$ or $d$ quark, $\ell$ an electron or a muon, and $\tau_h$,
$\tau_e$, $\tau_\mu$ a $\tau$ lepton decaying to hadrons, an electron or a muon. Each box is one
class of the label set above its column, except the QCD box of the \Kfull-class column (Sophon's
label set~\cite{sophon}), which stands for \LitJetClassTwoNQcdClasses{} classes. Each level merges
classes of the level to its left; the lines below name what each merge removes.
(b)~The controls: the first random partition, where boxes with the same letter are one class, and
the random partitions with controlled flavor boundaries on the decays \RandFlavSplitOrbit{} (one
square per decay; filled squares are moved to a second class).
(c)~Models of one seed share the training jets, their order and the initial weights, except the
output layer, the mass regression target, and the first layer and initial weights of the model with
kinematic inputs only (the particle four-momenta). The hierarchy levels have \MlstNRunsWord{} seeds
each, the models with the auxiliary mass regression \MlstMassNRunsWord{}, the random partitions with
controlled flavor boundaries \RandFlavRuns{} and each random partition \RandVtwoRuns{}. Frozen
features, the output of the backbone (the network below the classification head) with its weights held
fixed, feed linear probes (logistic regressions) and multilayer perceptron (MLP) probes; fine-tuning
trains all weights. Colors identify the label sets and controls in every figure; the plots give the three partitions with
controlled flavor boundaries three shades; dotted borders mark the two
variants, which the other figures draw as dotted lines in the color of their label set.}
\label{fig:design}
\end{figure}
 from main.tex; this file holds one figure environment and nothing else.
%
% WHAT main.tex MUST LOAD (this file loads nothing itself):
%   \usepackage{tikz}
%   \usetikzlibrary{arrows.meta}          % the only TikZ library used (arrow tips)
%   \definecolor{c188}{HTML}{1F77B4}      % the configuration colours of experiments/FIGS/
%   \definecolor{c162}{HTML}{FF7F0E}      % mlst_style.py (matplotlib's tab10 in order, as in
%   \definecolor{c43}{HTML}{2CA02C}       % Sophon's figures); the variants are drawn in their
%   \definecolor{c17}{HTML}{D62728}       % label set's colour with a dotted border, as the
%   \definecolor{crand}{HTML}{17BECF}     % plotted figures draw variants dotted
%   \definecolor{cflav}{HTML}{9467BD}
%   % GENERATED by experiments/FIGS/make_tables.py -- do not edit.
% Every line carries: source file :: path inside that file :: first 16 hex of its
% sha256. paper/journal/provenance.json carries the full hash beside the value.
% Regenerate with:  python3 experiments/FIGS/make_tables.py
% Check in CI with: python3 experiments/FIGS/make_tables.py --check
%
% Macro names spell digits as words (162 -> Onesixtwo) because \newcommand takes
% letters only. A number whose source file does not exist yet is NOT emitted --
% see the missing list below, and leave a \pending{} in the text for it.
%
% MISSING INPUTS -- no macro was emitted for these:
%   community benchmarks, bench v3 readout -- experiments/FIGS/data/bench_v3_metrics*/bench_metrics_last.json
%   v2 pretraining rerun: data order and dropout fixed by seed and epoch, identical initial trunk weights, each epoch reading one of five disjoint random parts of the rows of nearly every file, fixed validation sample, checkpoint at the first maximum within the last ten epochs -- experiments/FIGS/data/v2/pretraining/
%   v2 frozen probes on the rerun models, at the first maximum within the last ten epochs and on the model whose weights are the average of the last ten epochs, BatchNorm statistics recomputed, with untrained-network and pooled-embedding reference rows -- experiments/FIGS/data/v2/probe_ladder/
%   v2 paired ratios with run and test-sample errors on the rerun models -- experiments/FIGS/data/v2/paired_errors/
%   v2 the 64- and 30-class vocabularies -- experiments/FIGS/data/v2/levels_64_30/
%   v2 the leave-one-family-out vocabularies -- experiments/FIGS/data/v2/leave_one_family_out/
%   v2 label-tree recovery on the rerun models -- experiments/FIGS/data/v2/label_recovery/
%   v2 the rerun's random partitions and flavour pair -- experiments/FIGS/data/v2/random_partitions/
%   v2 JetClass-II fine-tuning on files held out from pretraining -- experiments/FIGS/data/v2/finetune/
%   v2 fine-tuning from random initialisation and from the self-supervised model -- experiments/FIGS/data/v2/finetune_references/
%   v2 top tagging and quark/gluon tagging (Pythia and Herwig) fine-tuned from every rerun model -- experiments/FIGS/data/v2/benchmarks/
%   v2 the second and third self-supervised pretraining runs -- experiments/FIGS/data/v2/self_supervised/
%   v2 the mass-output models with the loss weight matched in loss share -- experiments/FIGS/data/v2/mass_lambda_matched/
%   v2 anomaly detection on the rerun models, per-run values -- experiments/FIGS/data/v2/anomaly/
%   v2 the open-data fits on the rerun models -- experiments/FIGS/data/v2/real_data/
%   MassLossShareVtwoOnesixtwoMass -- experiments/FIGS/data/v2/mass_lambda_matched/loss_share.json :: shares.162+mass
%   MassLossShareVtwoOnesevenMass -- experiments/FIGS/data/v2/mass_lambda_matched/loss_share.json :: shares.17+mass
%   MassLossShareVtwoOnesevenMassMatched -- experiments/FIGS/data/v2/mass_lambda_matched/loss_share.json :: shares.17+mass_matched
%   DesignVtwoSecondGpu, DesignVtwoSecondGpuRuns -- the grid's job specs put every run index on one product
%   per-arm parameters and MACs -- flops_by_arm.json (see docs/RECORD.md 2.1)
%   fine-tuning read-outs -- experiments/FIGS/data/v2/ft_mlst/: jc2, jc2_qcd
%   top-tagging recipe check -- ft_mlst/topcheck_bench_metrics.json
%   X->bb and X->bs fits -- experiments/FIGS/data/v2/xbx/
%   X->bs from scratch -- experiments/FIGS/data/v2/xbx_scratch/
%   JetClass accuracy bootstrap -- experiments/FIGS/data/v2/budget/jc1_acc_boot_N10000.json
%   MlstFtJcOneErrCoarseOverFullBootEFour -- waits on the JetClass bootstrap at 1e4 (after the equal-stopping read-out)
%   MlstLpBbQgZeroAccDropFull -- waits on pair probes with the quark/gluon inputs zeroed (open item O10, not run)
%   MlstLpBbQgZeroAccDropCoarse -- waits on pair probes with the quark/gluon inputs zeroed (open item O10, not run)
%   XbsRthirtyFullOverScratch -- waits on X->bs fits and from-scratch runs
%   XbsRthirtyFullOverScratchPr -- waits on X->bs fits and from-scratch runs
%   XbsRthirtyCoarseOverFull -- waits on X->bs fits
%   XbsRthirtyCoarseOverFullPr -- waits on X->bs fits
%   XbsRthirtyNodeOverTransferFull -- waits on X->bs fits
%   XbsRthirtyNodeOverTransferFullPr -- waits on X->bs fits
%   XbbRthirtyDirectOverTransferFull -- waits on X->bb fits
%   XbbRthirtyDirectOverTransferFullPr -- waits on X->bb fits
%   XbbRthirtyTransferCoarseOverFull -- waits on X->bb fits
%   XbbRthirtyTransferCoarseOverFullPr -- waits on X->bb fits
%
% SKIPPED TABLES:
%   T4 compute/parameters: no machine-readable source in the repository; docs/RECORD.md 2.1 names /data/results/mtx/flops/flops_by_arm.json, which is not committed

\newcommand{\PendingPretrain}{\pending{pending v2: pretraining rerun: data order and dropout fixed by seed and epoch, identical initial trunk weights, each epoch reading one of five disjoint random parts of the rows of nearly every file, fixed validation sample, checkpoint at the first maximum within the last ten epochs}}  % experiments/FIGS/make_tables.py :: V2_PENDING['Pretrain'] :: sha256 10e6ec08356a4a2b :: a marker, not a number: red in the draft until the input exists
\newcommand{\PendingProbes}{\pending{pending v2: frozen probes on the rerun models, at the first maximum within the last ten epochs and on the model whose weights are the average of the last ten epochs, BatchNorm statistics recomputed, with untrained-network and pooled-embedding reference rows}}  % experiments/FIGS/make_tables.py :: V2_PENDING['Probes'] :: sha256 10e6ec08356a4a2b :: a marker, not a number: red in the draft until the input exists
\newcommand{\PendingPaired}{\pending{pending v2: paired ratios with run and test-sample errors on the rerun models}}  % experiments/FIGS/make_tables.py :: V2_PENDING['Paired'] :: sha256 10e6ec08356a4a2b :: a marker, not a number: red in the draft until the input exists
\newcommand{\PendingLevels}{\pending{pending v2: the 64- and 30-class vocabularies}}  % experiments/FIGS/make_tables.py :: V2_PENDING['Levels'] :: sha256 10e6ec08356a4a2b :: a marker, not a number: red in the draft until the input exists
\newcommand{\PendingLofo}{\pending{pending v2: the leave-one-family-out vocabularies}}  % experiments/FIGS/make_tables.py :: V2_PENDING['Lofo'] :: sha256 10e6ec08356a4a2b :: a marker, not a number: red in the draft until the input exists
\newcommand{\PendingRecovery}{\pending{pending v2: label-tree recovery on the rerun models}}  % experiments/FIGS/make_tables.py :: V2_PENDING['Recovery'] :: sha256 10e6ec08356a4a2b :: a marker, not a number: red in the draft until the input exists
\newcommand{\PendingRandom}{\pending{pending v2: the rerun's random partitions and flavour pair}}  % experiments/FIGS/make_tables.py :: V2_PENDING['Random'] :: sha256 10e6ec08356a4a2b :: a marker, not a number: red in the draft until the input exists
\newcommand{\PendingFtHeldout}{\pending{pending v2: JetClass-II fine-tuning on files held out from pretraining}}  % experiments/FIGS/make_tables.py :: V2_PENDING['FtHeldout'] :: sha256 10e6ec08356a4a2b :: a marker, not a number: red in the draft until the input exists
\newcommand{\PendingFtRefs}{\pending{pending v2: fine-tuning from random initialisation and from the self-supervised model}}  % experiments/FIGS/make_tables.py :: V2_PENDING['FtRefs'] :: sha256 10e6ec08356a4a2b :: a marker, not a number: red in the draft until the input exists
\newcommand{\PendingBench}{\pending{pending v2: top tagging and quark/gluon tagging (Pythia and Herwig) fine-tuned from every rerun model}}  % experiments/FIGS/make_tables.py :: V2_PENDING['Bench'] :: sha256 10e6ec08356a4a2b :: a marker, not a number: red in the draft until the input exists
\newcommand{\PendingSsl}{\pending{pending v2: the second and third self-supervised pretraining runs}}  % experiments/FIGS/make_tables.py :: V2_PENDING['Ssl'] :: sha256 10e6ec08356a4a2b :: a marker, not a number: red in the draft until the input exists
\newcommand{\PendingMassLambda}{\pending{pending v2: the mass-output models with the loss weight matched in loss share}}  % experiments/FIGS/make_tables.py :: V2_PENDING['MassLambda'] :: sha256 10e6ec08356a4a2b :: a marker, not a number: red in the draft until the input exists
\newcommand{\PendingAnomaly}{\pending{pending v2: anomaly detection on the rerun models, per-run values}}  % experiments/FIGS/make_tables.py :: V2_PENDING['Anomaly'] :: sha256 10e6ec08356a4a2b :: a marker, not a number: red in the draft until the input exists
\newcommand{\PendingRealData}{\pending{pending v2: the open-data fits on the rerun models}}  % experiments/FIGS/make_tables.py :: V2_PENDING['RealData'] :: sha256 10e6ec08356a4a2b :: a marker, not a number: red in the draft until the input exists
\newcommand{\PendingConclusion}{\pending{pending v2: written once every input above is in}}  % experiments/FIGS/make_tables.py :: V2_PENDING['Conclusion'] :: sha256 10e6ec08356a4a2b :: a marker, not a number: red in the draft until the input exists
\newcommand{\AojChainDataset}{178}  % experiments/FIGS/data/aoj_checks_v2/selection_chain_v2.json :: n_dataset :: sha256 729d83085aba25f8
\newcommand{\AojChainStaged}{12{,}589{,}467}  % experiments/FIGS/data/aoj_checks_v2/selection_chain_v2.json :: n_staged :: sha256 729d83085aba25f8
\newcommand{\AojChainRhoWindow}{9{,}441{,}209}  % experiments/FIGS/data/aoj_checks_v2/selection_chain_v2.json :: n_rho_window :: sha256 729d83085aba25f8
\newcommand{\AojNJetsFit}{4{,}194{,}846}  % experiments/FIGS/data/aoj_checks_v2/selection_chain_v2.json :: n_fit :: sha256 729d83085aba25f8
\newcommand{\AojChainPtBins}{8}  % experiments/FIGS/data/aoj_checks_v2/selection_chain_v2.json :: n_pt_bins :: sha256 729d83085aba25f8
\newcommand{\AojChainEtaMax}{2.40}  % experiments/FIGS/data/aoj_checks_v2/selection_chain_v2.json :: abs_eta_max :: sha256 729d83085aba25f8
\newcommand{\MassLossShareVtwoOnesixtwoMass}{\pending{pending: realised mass-loss share in the rerun, 162 classes}}  % experiments/FIGS/make_tables.py :: SLOTS[MassLossShareVtwoOnesixtwoMass] :: sha256 10e6ec08356a4a2b :: a marker, not a number: experiments/FIGS/data/v2/mass_lambda_matched/loss_share.json :: shares.162+mass does not exist yet
\newcommand{\MassLossShareVtwoOnesevenMass}{\pending{pending: realised mass-loss share in the rerun, 17 classes}}  % experiments/FIGS/make_tables.py :: SLOTS[MassLossShareVtwoOnesevenMass] :: sha256 10e6ec08356a4a2b :: a marker, not a number: experiments/FIGS/data/v2/mass_lambda_matched/loss_share.json :: shares.17+mass does not exist yet
\newcommand{\MassLossShareVtwoOnesevenMassMatched}{\pending{pending: realised mass-loss share in the rerun, 17 classes, matched weight}}  % experiments/FIGS/make_tables.py :: SLOTS[MassLossShareVtwoOnesevenMassMatched] :: sha256 10e6ec08356a4a2b :: a marker, not a number: experiments/FIGS/data/v2/mass_lambda_matched/loss_share.json :: shares.17+mass_matched does not exist yet
\newcommand{\FtCoverageEThreeNClasses}{160}  % experiments/FT/data/jc2_v2_manifest.json :: class_coverage['train_N1000_s1.parquet'].n_classes_present :: sha256 4bd68655919b14e9
\newcommand{\FtCoverageEFourNClasses}{183}  % experiments/FT/data/jc2_v2_manifest.json :: class_coverage['train_N10000_s1.parquet'].n_classes_present :: sha256 4bd68655919b14e9
\newcommand{\FtCoverageEThreeMedian}{4}  % experiments/FT/data/jc2_v2_manifest.json :: class_coverage['train_N1000_s1.parquet'].median_per_present_class :: sha256 4bd68655919b14e9
\newcommand{\FtCoverageEThreeOneJet}{23}  % experiments/FT/data/jc2_v2_manifest.json :: class_coverage['train_N1000_s1.parquet'].n_classes_with_one_jet :: sha256 4bd68655919b14e9
\newcommand{\FtCoverageVoneEThreeNClasses}{161}  % experiments/FT/data/jc2_v1_coverage.json :: class_coverage['train_N1000_s1.parquet'].n_classes_present :: sha256 82652f57f708dc23
\newcommand{\FtCoverageVoneEThreeMedian}{5}  % experiments/FT/data/jc2_v1_coverage.json :: class_coverage['train_N1000_s1.parquet'].median_per_present_class :: sha256 82652f57f708dc23
\newcommand{\FtCoverageVoneEThreeOneJet}{25}  % experiments/FT/data/jc2_v1_coverage.json :: class_coverage['train_N1000_s1.parquet'].n_classes_with_one_jet :: sha256 82652f57f708dc23
\newcommand{\ProbeNJets}{2{,}000{,}000}  % experiments/FIGS/data/probe_ladder_v2_mlp2/analysis/seed_level_results.json :: provenance.n_jets_total :: sha256 53e1ec81343ac2c1
\newcommand{\ProbeNSeeds}{5}  % experiments/FIGS/data/probe_ladder_v2_mlp2/analysis/seed_level_results.json :: seeds_used (length) :: sha256 53e1ec81343ac2c1
\newcommand{\ProbeNSeedsWord}{five}  % experiments/FIGS/data/probe_ladder_v2_mlp2/analysis/seed_level_results.json :: seeds_used (length) :: sha256 53e1ec81343ac2c1 :: pretraining runs per vocabulary, as a word for the prose
\newcommand{\ProbeRowAlign}{\texttt{a9449fd420c473f6}}  % experiments/FIGS/data/probe_ladder_v2_mlp2/analysis/seed_level_results.json :: provenance.row_alignment_sha256 (first 16) :: sha256 53e1ec81343ac2c1
\newcommand{\ProbeEpsS}{90\%}  % experiments/FIGS/data/probe_ladder_v2_mlp2/analysis/seed_level_results.json :: levels.bvc_4prong.linear[0].headline_eps_s :: sha256 53e1ec81343ac2c1
\newcommand{\ProbeAucBvcFourprongLinearOneeighteight}{0.99320}  % experiments/FIGS/data/probe_ladder_v2_mlp2/analysis/seed_level_results.json :: levels.bvc_4prong.linear[0].mean_auc :: sha256 53e1ec81343ac2c1 :: 5 seeds, 0 saturated at AUC=1
\newcommand{\ProbeAucSdBvcFourprongLinearOneeighteight}{0.00027}  % experiments/FIGS/data/probe_ladder_v2_mlp2/analysis/seed_level_results.json :: table[bvc_4prong,linear,188].auc (SD over seeds, ddof=1) :: sha256 53e1ec81343ac2c1
\newcommand{\ProbeLogOneMinusAucBvcFourprongLinearOneeighteight}{\ensuremath{-}4.9915}  % experiments/FIGS/data/probe_ladder_v2_mlp2/analysis/seed_level_results.json :: levels.bvc_4prong.linear[0].mean :: sha256 53e1ec81343ac2c1 :: natural log of 1-AUC, lower is better
\newcommand{\ProbeLogOneMinusAucSdBvcFourprongLinearOneeighteight}{0.0389}  % experiments/FIGS/data/probe_ladder_v2_mlp2/analysis/seed_level_results.json :: levels.bvc_4prong.linear[0].seed_sd :: sha256 53e1ec81343ac2c1
\newcommand{\ProbeOmaBvcFourprongLinearOneeighteight}{\ensuremath{(6.80 \pm 0.27)\times10^{-3}}}  % experiments/FIGS/data/probe_ladder_v2_mlp2/analysis/seed_level_results.json :: table[bvc_4prong,linear,188].auc :: sha256 53e1ec81343ac2c1 :: 1-AUC, mean +- SD over 5 seeds
\newcommand{\ProbeBkgLeftBvcFourprongLinearOneeighteight}{14.0}  % experiments/FIGS/data/probe_ladder_v2_mlp2/analysis/seed_level_results.json :: levels.bvc_4prong.linear[0].rejection_points['0.90'].mean_n_bkg_pass :: sha256 53e1ec81343ac2c1 :: background jets passing the cut, mean over seeds
\newcommand{\ProbeRejBvcFourprongLinearOneeighteight}{227\,\ensuremath{\pm}\,23}  % experiments/FIGS/data/probe_ladder_v2_mlp2/analysis/seed_level_results.json :: table[bvc_4prong,linear,188].rejection_points['0.90'].rejection :: sha256 53e1ec81343ac2c1 :: at 90% signal efficiency, mean +- SD over seeds; 0 of 5 seeds at the cap
\newcommand{\ProbeAucBvcFourprongLinearOnesixtwo}{0.99311}  % experiments/FIGS/data/probe_ladder_v2_mlp2/analysis/seed_level_results.json :: levels.bvc_4prong.linear[1].mean_auc :: sha256 53e1ec81343ac2c1 :: 5 seeds, 0 saturated at AUC=1
\newcommand{\ProbeAucSdBvcFourprongLinearOnesixtwo}{0.00018}  % experiments/FIGS/data/probe_ladder_v2_mlp2/analysis/seed_level_results.json :: table[bvc_4prong,linear,162].auc (SD over seeds, ddof=1) :: sha256 53e1ec81343ac2c1
\newcommand{\ProbeLogOneMinusAucBvcFourprongLinearOnesixtwo}{\ensuremath{-}4.9777}  % experiments/FIGS/data/probe_ladder_v2_mlp2/analysis/seed_level_results.json :: levels.bvc_4prong.linear[1].mean :: sha256 53e1ec81343ac2c1 :: natural log of 1-AUC, lower is better
\newcommand{\ProbeLogOneMinusAucSdBvcFourprongLinearOnesixtwo}{0.0264}  % experiments/FIGS/data/probe_ladder_v2_mlp2/analysis/seed_level_results.json :: levels.bvc_4prong.linear[1].seed_sd :: sha256 53e1ec81343ac2c1
\newcommand{\ProbeOmaBvcFourprongLinearOnesixtwo}{\ensuremath{(6.89 \pm 0.18)\times10^{-3}}}  % experiments/FIGS/data/probe_ladder_v2_mlp2/analysis/seed_level_results.json :: table[bvc_4prong,linear,162].auc :: sha256 53e1ec81343ac2c1 :: 1-AUC, mean +- SD over 5 seeds
\newcommand{\ProbeBkgLeftBvcFourprongLinearOnesixtwo}{14.4}  % experiments/FIGS/data/probe_ladder_v2_mlp2/analysis/seed_level_results.json :: levels.bvc_4prong.linear[1].rejection_points['0.90'].mean_n_bkg_pass :: sha256 53e1ec81343ac2c1 :: background jets passing the cut, mean over seeds
\newcommand{\ProbeRejBvcFourprongLinearOnesixtwo}{220\,\ensuremath{\pm}\,40}  % experiments/FIGS/data/probe_ladder_v2_mlp2/analysis/seed_level_results.json :: table[bvc_4prong,linear,162].rejection_points['0.90'].rejection :: sha256 53e1ec81343ac2c1 :: at 90% signal efficiency, mean +- SD over seeds; 0 of 5 seeds at the cap
\newcommand{\ProbeAucBvcFourprongLinearFourthree}{0.99250}  % experiments/FIGS/data/probe_ladder_v2_mlp2/analysis/seed_level_results.json :: levels.bvc_4prong.linear[2].mean_auc :: sha256 53e1ec81343ac2c1 :: 5 seeds, 0 saturated at AUC=1
\newcommand{\ProbeAucSdBvcFourprongLinearFourthree}{0.00049}  % experiments/FIGS/data/probe_ladder_v2_mlp2/analysis/seed_level_results.json :: table[bvc_4prong,linear,43].auc (SD over seeds, ddof=1) :: sha256 53e1ec81343ac2c1
\newcommand{\ProbeLogOneMinusAucBvcFourprongLinearFourthree}{\ensuremath{-}4.8946}  % experiments/FIGS/data/probe_ladder_v2_mlp2/analysis/seed_level_results.json :: levels.bvc_4prong.linear[2].mean :: sha256 53e1ec81343ac2c1 :: natural log of 1-AUC, lower is better
\newcommand{\ProbeLogOneMinusAucSdBvcFourprongLinearFourthree}{0.0646}  % experiments/FIGS/data/probe_ladder_v2_mlp2/analysis/seed_level_results.json :: levels.bvc_4prong.linear[2].seed_sd :: sha256 53e1ec81343ac2c1
\newcommand{\ProbeOmaBvcFourprongLinearFourthree}{\ensuremath{(7.5 \pm 0.5)\times10^{-3}}}  % experiments/FIGS/data/probe_ladder_v2_mlp2/analysis/seed_level_results.json :: table[bvc_4prong,linear,43].auc :: sha256 53e1ec81343ac2c1 :: 1-AUC, mean +- SD over 5 seeds
\newcommand{\ProbeBkgLeftBvcFourprongLinearFourthree}{16.6}  % experiments/FIGS/data/probe_ladder_v2_mlp2/analysis/seed_level_results.json :: levels.bvc_4prong.linear[2].rejection_points['0.90'].mean_n_bkg_pass :: sha256 53e1ec81343ac2c1 :: background jets passing the cut, mean over seeds
\newcommand{\ProbeRejBvcFourprongLinearFourthree}{192\,\ensuremath{\pm}\,20}  % experiments/FIGS/data/probe_ladder_v2_mlp2/analysis/seed_level_results.json :: table[bvc_4prong,linear,43].rejection_points['0.90'].rejection :: sha256 53e1ec81343ac2c1 :: at 90% signal efficiency, mean +- SD over seeds; 0 of 5 seeds at the cap
\newcommand{\ProbeAucBvcFourprongLinearOneseven}{0.96936}  % experiments/FIGS/data/probe_ladder_v2_mlp2/analysis/seed_level_results.json :: levels.bvc_4prong.linear[3].mean_auc :: sha256 53e1ec81343ac2c1 :: 5 seeds, 0 saturated at AUC=1
\newcommand{\ProbeAucSdBvcFourprongLinearOneseven}{0.00514}  % experiments/FIGS/data/probe_ladder_v2_mlp2/analysis/seed_level_results.json :: table[bvc_4prong,linear,17].auc (SD over seeds, ddof=1) :: sha256 53e1ec81343ac2c1
\newcommand{\ProbeLogOneMinusAucBvcFourprongLinearOneseven}{\ensuremath{-}3.4967}  % experiments/FIGS/data/probe_ladder_v2_mlp2/analysis/seed_level_results.json :: levels.bvc_4prong.linear[3].mean :: sha256 53e1ec81343ac2c1 :: natural log of 1-AUC, lower is better
\newcommand{\ProbeLogOneMinusAucSdBvcFourprongLinearOneseven}{0.1667}  % experiments/FIGS/data/probe_ladder_v2_mlp2/analysis/seed_level_results.json :: levels.bvc_4prong.linear[3].seed_sd :: sha256 53e1ec81343ac2c1
\newcommand{\ProbeOmaBvcFourprongLinearOneseven}{\ensuremath{(3.1 \pm 0.5)\times10^{-2}}}  % experiments/FIGS/data/probe_ladder_v2_mlp2/analysis/seed_level_results.json :: table[bvc_4prong,linear,17].auc :: sha256 53e1ec81343ac2c1 :: 1-AUC, mean +- SD over 5 seeds
\newcommand{\ProbeBkgLeftBvcFourprongLinearOneseven}{137.4}  % experiments/FIGS/data/probe_ladder_v2_mlp2/analysis/seed_level_results.json :: levels.bvc_4prong.linear[3].rejection_points['0.90'].mean_n_bkg_pass :: sha256 53e1ec81343ac2c1 :: background jets passing the cut, mean over seeds
\newcommand{\ProbeRejBvcFourprongLinearOneseven}{24\,\ensuremath{\pm}\,6}  % experiments/FIGS/data/probe_ladder_v2_mlp2/analysis/seed_level_results.json :: table[bvc_4prong,linear,17].rejection_points['0.90'].rejection :: sha256 53e1ec81343ac2c1 :: at 90% signal efficiency, mean +- SD over seeds; 0 of 5 seeds at the cap
\newcommand{\ProbeOmaRatioBvcFourprongLinearOnesixtwoOverOneeighteight}{1.01}  % experiments/FIGS/data/probe_ladder_v2_mlp2/analysis/seed_level_results.json :: table[bvc_4prong,linear,162/188].auc :: sha256 53e1ec81343ac2c1 :: ratio of the seed means of 1-AUC, coarser over finer
\newcommand{\ProbeOmaRatioBvcFourprongLinearFourthreeOverOneeighteight}{1.10}  % experiments/FIGS/data/probe_ladder_v2_mlp2/analysis/seed_level_results.json :: table[bvc_4prong,linear,43/188].auc :: sha256 53e1ec81343ac2c1 :: ratio of the seed means of 1-AUC, coarser over finer
\newcommand{\ProbeOmaRatioBvcFourprongLinearOnesevenOverOneeighteight}{4.5}  % experiments/FIGS/data/probe_ladder_v2_mlp2/analysis/seed_level_results.json :: table[bvc_4prong,linear,17/188].auc :: sha256 53e1ec81343ac2c1 :: ratio of the seed means of 1-AUC, coarser over finer
\newcommand{\ProbeOmaRatioBvcFourprongLinearFourthreeOverOnesixtwo}{1.09}  % experiments/FIGS/data/probe_ladder_v2_mlp2/analysis/seed_level_results.json :: table[bvc_4prong,linear,43/162].auc :: sha256 53e1ec81343ac2c1 :: ratio of the seed means of 1-AUC, coarser over finer
\newcommand{\ProbeOmaRatioBvcFourprongLinearOnesevenOverOnesixtwo}{4.4}  % experiments/FIGS/data/probe_ladder_v2_mlp2/analysis/seed_level_results.json :: table[bvc_4prong,linear,17/162].auc :: sha256 53e1ec81343ac2c1 :: ratio of the seed means of 1-AUC, coarser over finer
\newcommand{\ProbeOmaRatioBvcFourprongLinearOnesevenOverFourthree}{4.1}  % experiments/FIGS/data/probe_ladder_v2_mlp2/analysis/seed_level_results.json :: table[bvc_4prong,linear,17/43].auc :: sha256 53e1ec81343ac2c1 :: ratio of the seed means of 1-AUC, coarser over finer
\newcommand{\ProbeAucBvcFourprongMlpOneeighteight}{0.99316}  % experiments/FIGS/data/probe_ladder_v2_mlp2/analysis/seed_level_results.json :: levels.bvc_4prong.mlp[0].mean_auc :: sha256 53e1ec81343ac2c1 :: 5 seeds, 0 saturated at AUC=1
\newcommand{\ProbeAucSdBvcFourprongMlpOneeighteight}{0.00020}  % experiments/FIGS/data/probe_ladder_v2_mlp2/analysis/seed_level_results.json :: table[bvc_4prong,mlp,188].auc (SD over seeds, ddof=1) :: sha256 53e1ec81343ac2c1
\newcommand{\ProbeLogOneMinusAucBvcFourprongMlpOneeighteight}{\ensuremath{-}4.9857}  % experiments/FIGS/data/probe_ladder_v2_mlp2/analysis/seed_level_results.json :: levels.bvc_4prong.mlp[0].mean :: sha256 53e1ec81343ac2c1 :: natural log of 1-AUC, lower is better
\newcommand{\ProbeLogOneMinusAucSdBvcFourprongMlpOneeighteight}{0.0290}  % experiments/FIGS/data/probe_ladder_v2_mlp2/analysis/seed_level_results.json :: levels.bvc_4prong.mlp[0].seed_sd :: sha256 53e1ec81343ac2c1
\newcommand{\ProbeOmaBvcFourprongMlpOneeighteight}{\ensuremath{(6.84 \pm 0.20)\times10^{-3}}}  % experiments/FIGS/data/probe_ladder_v2_mlp2/analysis/seed_level_results.json :: table[bvc_4prong,mlp,188].auc :: sha256 53e1ec81343ac2c1 :: 1-AUC, mean +- SD over 5 seeds
\newcommand{\ProbeBkgLeftBvcFourprongMlpOneeighteight}{15.2}  % experiments/FIGS/data/probe_ladder_v2_mlp2/analysis/seed_level_results.json :: levels.bvc_4prong.mlp[0].rejection_points['0.90'].mean_n_bkg_pass :: sha256 53e1ec81343ac2c1 :: background jets passing the cut, mean over seeds
\newcommand{\ProbeRejBvcFourprongMlpOneeighteight}{212\,\ensuremath{\pm}\,34}  % experiments/FIGS/data/probe_ladder_v2_mlp2/analysis/seed_level_results.json :: table[bvc_4prong,mlp,188].rejection_points['0.90'].rejection :: sha256 53e1ec81343ac2c1 :: at 90% signal efficiency, mean +- SD over seeds; 0 of 5 seeds at the cap
\newcommand{\ProbeAucBvcFourprongMlpOnesixtwo}{0.99318}  % experiments/FIGS/data/probe_ladder_v2_mlp2/analysis/seed_level_results.json :: levels.bvc_4prong.mlp[1].mean_auc :: sha256 53e1ec81343ac2c1 :: 5 seeds, 0 saturated at AUC=1
\newcommand{\ProbeAucSdBvcFourprongMlpOnesixtwo}{0.00020}  % experiments/FIGS/data/probe_ladder_v2_mlp2/analysis/seed_level_results.json :: table[bvc_4prong,mlp,162].auc (SD over seeds, ddof=1) :: sha256 53e1ec81343ac2c1
\newcommand{\ProbeLogOneMinusAucBvcFourprongMlpOnesixtwo}{\ensuremath{-}4.9886}  % experiments/FIGS/data/probe_ladder_v2_mlp2/analysis/seed_level_results.json :: levels.bvc_4prong.mlp[1].mean :: sha256 53e1ec81343ac2c1 :: natural log of 1-AUC, lower is better
\newcommand{\ProbeLogOneMinusAucSdBvcFourprongMlpOnesixtwo}{0.0298}  % experiments/FIGS/data/probe_ladder_v2_mlp2/analysis/seed_level_results.json :: levels.bvc_4prong.mlp[1].seed_sd :: sha256 53e1ec81343ac2c1
\newcommand{\ProbeOmaBvcFourprongMlpOnesixtwo}{\ensuremath{(6.82 \pm 0.20)\times10^{-3}}}  % experiments/FIGS/data/probe_ladder_v2_mlp2/analysis/seed_level_results.json :: table[bvc_4prong,mlp,162].auc :: sha256 53e1ec81343ac2c1 :: 1-AUC, mean +- SD over 5 seeds
\newcommand{\ProbeBkgLeftBvcFourprongMlpOnesixtwo}{16.6}  % experiments/FIGS/data/probe_ladder_v2_mlp2/analysis/seed_level_results.json :: levels.bvc_4prong.mlp[1].rejection_points['0.90'].mean_n_bkg_pass :: sha256 53e1ec81343ac2c1 :: background jets passing the cut, mean over seeds
\newcommand{\ProbeRejBvcFourprongMlpOnesixtwo}{200\,\ensuremath{\pm}\,40}  % experiments/FIGS/data/probe_ladder_v2_mlp2/analysis/seed_level_results.json :: table[bvc_4prong,mlp,162].rejection_points['0.90'].rejection :: sha256 53e1ec81343ac2c1 :: at 90% signal efficiency, mean +- SD over seeds; 0 of 5 seeds at the cap
\newcommand{\ProbeAucBvcFourprongMlpFourthree}{0.99232}  % experiments/FIGS/data/probe_ladder_v2_mlp2/analysis/seed_level_results.json :: levels.bvc_4prong.mlp[2].mean_auc :: sha256 53e1ec81343ac2c1 :: 5 seeds, 0 saturated at AUC=1
\newcommand{\ProbeAucSdBvcFourprongMlpFourthree}{0.00058}  % experiments/FIGS/data/probe_ladder_v2_mlp2/analysis/seed_level_results.json :: table[bvc_4prong,mlp,43].auc (SD over seeds, ddof=1) :: sha256 53e1ec81343ac2c1
\newcommand{\ProbeLogOneMinusAucBvcFourprongMlpFourthree}{\ensuremath{-}4.8719}  % experiments/FIGS/data/probe_ladder_v2_mlp2/analysis/seed_level_results.json :: levels.bvc_4prong.mlp[2].mean :: sha256 53e1ec81343ac2c1 :: natural log of 1-AUC, lower is better
\newcommand{\ProbeLogOneMinusAucSdBvcFourprongMlpFourthree}{0.0734}  % experiments/FIGS/data/probe_ladder_v2_mlp2/analysis/seed_level_results.json :: levels.bvc_4prong.mlp[2].seed_sd :: sha256 53e1ec81343ac2c1
\newcommand{\ProbeOmaBvcFourprongMlpFourthree}{\ensuremath{(7.7 \pm 0.6)\times10^{-3}}}  % experiments/FIGS/data/probe_ladder_v2_mlp2/analysis/seed_level_results.json :: table[bvc_4prong,mlp,43].auc :: sha256 53e1ec81343ac2c1 :: 1-AUC, mean +- SD over 5 seeds
\newcommand{\ProbeBkgLeftBvcFourprongMlpFourthree}{17.4}  % experiments/FIGS/data/probe_ladder_v2_mlp2/analysis/seed_level_results.json :: levels.bvc_4prong.mlp[2].rejection_points['0.90'].mean_n_bkg_pass :: sha256 53e1ec81343ac2c1 :: background jets passing the cut, mean over seeds
\newcommand{\ProbeRejBvcFourprongMlpFourthree}{184\,\ensuremath{\pm}\,23}  % experiments/FIGS/data/probe_ladder_v2_mlp2/analysis/seed_level_results.json :: table[bvc_4prong,mlp,43].rejection_points['0.90'].rejection :: sha256 53e1ec81343ac2c1 :: at 90% signal efficiency, mean +- SD over seeds; 0 of 5 seeds at the cap
\newcommand{\ProbeAucBvcFourprongMlpOneseven}{0.97766}  % experiments/FIGS/data/probe_ladder_v2_mlp2/analysis/seed_level_results.json :: levels.bvc_4prong.mlp[3].mean_auc :: sha256 53e1ec81343ac2c1 :: 5 seeds, 0 saturated at AUC=1
\newcommand{\ProbeAucSdBvcFourprongMlpOneseven}{0.00310}  % experiments/FIGS/data/probe_ladder_v2_mlp2/analysis/seed_level_results.json :: table[bvc_4prong,mlp,17].auc (SD over seeds, ddof=1) :: sha256 53e1ec81343ac2c1
\newcommand{\ProbeLogOneMinusAucBvcFourprongMlpOneseven}{\ensuremath{-}3.8087}  % experiments/FIGS/data/probe_ladder_v2_mlp2/analysis/seed_level_results.json :: levels.bvc_4prong.mlp[3].mean :: sha256 53e1ec81343ac2c1 :: natural log of 1-AUC, lower is better
\newcommand{\ProbeLogOneMinusAucSdBvcFourprongMlpOneseven}{0.1366}  % experiments/FIGS/data/probe_ladder_v2_mlp2/analysis/seed_level_results.json :: levels.bvc_4prong.mlp[3].seed_sd :: sha256 53e1ec81343ac2c1
\newcommand{\ProbeOmaBvcFourprongMlpOneseven}{\ensuremath{(2.23 \pm 0.31)\times10^{-2}}}  % experiments/FIGS/data/probe_ladder_v2_mlp2/analysis/seed_level_results.json :: table[bvc_4prong,mlp,17].auc :: sha256 53e1ec81343ac2c1 :: 1-AUC, mean +- SD over 5 seeds
\newcommand{\ProbeBkgLeftBvcFourprongMlpOneseven}{99.0}  % experiments/FIGS/data/probe_ladder_v2_mlp2/analysis/seed_level_results.json :: levels.bvc_4prong.mlp[3].rejection_points['0.90'].mean_n_bkg_pass :: sha256 53e1ec81343ac2c1 :: background jets passing the cut, mean over seeds
\newcommand{\ProbeRejBvcFourprongMlpOneseven}{33\,\ensuremath{\pm}\,5}  % experiments/FIGS/data/probe_ladder_v2_mlp2/analysis/seed_level_results.json :: table[bvc_4prong,mlp,17].rejection_points['0.90'].rejection :: sha256 53e1ec81343ac2c1 :: at 90% signal efficiency, mean +- SD over seeds; 0 of 5 seeds at the cap
\newcommand{\ProbeOmaRatioBvcFourprongMlpOnesixtwoOverOneeighteight}{1.00}  % experiments/FIGS/data/probe_ladder_v2_mlp2/analysis/seed_level_results.json :: table[bvc_4prong,mlp,162/188].auc :: sha256 53e1ec81343ac2c1 :: ratio of the seed means of 1-AUC, coarser over finer
\newcommand{\ProbeOmaRatioBvcFourprongMlpFourthreeOverOneeighteight}{1.12}  % experiments/FIGS/data/probe_ladder_v2_mlp2/analysis/seed_level_results.json :: table[bvc_4prong,mlp,43/188].auc :: sha256 53e1ec81343ac2c1 :: ratio of the seed means of 1-AUC, coarser over finer
\newcommand{\ProbeOmaRatioBvcFourprongMlpOnesevenOverOneeighteight}{3.3}  % experiments/FIGS/data/probe_ladder_v2_mlp2/analysis/seed_level_results.json :: table[bvc_4prong,mlp,17/188].auc :: sha256 53e1ec81343ac2c1 :: ratio of the seed means of 1-AUC, coarser over finer
\newcommand{\ProbeOmaRatioBvcFourprongMlpFourthreeOverOnesixtwo}{1.13}  % experiments/FIGS/data/probe_ladder_v2_mlp2/analysis/seed_level_results.json :: table[bvc_4prong,mlp,43/162].auc :: sha256 53e1ec81343ac2c1 :: ratio of the seed means of 1-AUC, coarser over finer
\newcommand{\ProbeOmaRatioBvcFourprongMlpOnesevenOverOnesixtwo}{3.3}  % experiments/FIGS/data/probe_ladder_v2_mlp2/analysis/seed_level_results.json :: table[bvc_4prong,mlp,17/162].auc :: sha256 53e1ec81343ac2c1 :: ratio of the seed means of 1-AUC, coarser over finer
\newcommand{\ProbeOmaRatioBvcFourprongMlpOnesevenOverFourthree}{2.9}  % experiments/FIGS/data/probe_ladder_v2_mlp2/analysis/seed_level_results.json :: table[bvc_4prong,mlp,17/43].auc :: sha256 53e1ec81343ac2c1 :: ratio of the seed means of 1-AUC, coarser over finer
\newcommand{\ProbeAucBvcQcdLinearOneeighteight}{0.98519}  % experiments/FIGS/data/probe_ladder_v2_mlp2/analysis/seed_level_results.json :: levels.bvc_qcd.linear[0].mean_auc :: sha256 53e1ec81343ac2c1 :: 5 seeds, 0 saturated at AUC=1
\newcommand{\ProbeAucSdBvcQcdLinearOneeighteight}{0.00061}  % experiments/FIGS/data/probe_ladder_v2_mlp2/analysis/seed_level_results.json :: table[bvc_qcd,linear,188].auc (SD over seeds, ddof=1) :: sha256 53e1ec81343ac2c1
\newcommand{\ProbeLogOneMinusAucBvcQcdLinearOneeighteight}{\ensuremath{-}4.2133}  % experiments/FIGS/data/probe_ladder_v2_mlp2/analysis/seed_level_results.json :: levels.bvc_qcd.linear[0].mean :: sha256 53e1ec81343ac2c1 :: natural log of 1-AUC, lower is better
\newcommand{\ProbeLogOneMinusAucSdBvcQcdLinearOneeighteight}{0.0422}  % experiments/FIGS/data/probe_ladder_v2_mlp2/analysis/seed_level_results.json :: levels.bvc_qcd.linear[0].seed_sd :: sha256 53e1ec81343ac2c1
\newcommand{\ProbeOmaBvcQcdLinearOneeighteight}{\ensuremath{(1.48 \pm 0.06)\times10^{-2}}}  % experiments/FIGS/data/probe_ladder_v2_mlp2/analysis/seed_level_results.json :: table[bvc_qcd,linear,188].auc :: sha256 53e1ec81343ac2c1 :: 1-AUC, mean +- SD over 5 seeds
\newcommand{\ProbeBkgLeftBvcQcdLinearOneeighteight}{96.4}  % experiments/FIGS/data/probe_ladder_v2_mlp2/analysis/seed_level_results.json :: levels.bvc_qcd.linear[0].rejection_points['0.90'].mean_n_bkg_pass :: sha256 53e1ec81343ac2c1 :: background jets passing the cut, mean over seeds
\newcommand{\ProbeRejBvcQcdLinearOneeighteight}{56\,\ensuremath{\pm}\,5}  % experiments/FIGS/data/probe_ladder_v2_mlp2/analysis/seed_level_results.json :: table[bvc_qcd,linear,188].rejection_points['0.90'].rejection :: sha256 53e1ec81343ac2c1 :: at 90% signal efficiency, mean +- SD over seeds; 0 of 5 seeds at the cap
\newcommand{\ProbeAucBvcQcdLinearOnesixtwo}{0.98305}  % experiments/FIGS/data/probe_ladder_v2_mlp2/analysis/seed_level_results.json :: levels.bvc_qcd.linear[1].mean_auc :: sha256 53e1ec81343ac2c1 :: 5 seeds, 0 saturated at AUC=1
\newcommand{\ProbeAucSdBvcQcdLinearOnesixtwo}{0.00024}  % experiments/FIGS/data/probe_ladder_v2_mlp2/analysis/seed_level_results.json :: table[bvc_qcd,linear,162].auc (SD over seeds, ddof=1) :: sha256 53e1ec81343ac2c1
\newcommand{\ProbeLogOneMinusAucBvcQcdLinearOnesixtwo}{\ensuremath{-}4.0774}  % experiments/FIGS/data/probe_ladder_v2_mlp2/analysis/seed_level_results.json :: levels.bvc_qcd.linear[1].mean :: sha256 53e1ec81343ac2c1 :: natural log of 1-AUC, lower is better
\newcommand{\ProbeLogOneMinusAucSdBvcQcdLinearOnesixtwo}{0.0142}  % experiments/FIGS/data/probe_ladder_v2_mlp2/analysis/seed_level_results.json :: levels.bvc_qcd.linear[1].seed_sd :: sha256 53e1ec81343ac2c1
\newcommand{\ProbeOmaBvcQcdLinearOnesixtwo}{\ensuremath{(1.695 \pm 0.024)\times10^{-2}}}  % experiments/FIGS/data/probe_ladder_v2_mlp2/analysis/seed_level_results.json :: table[bvc_qcd,linear,162].auc :: sha256 53e1ec81343ac2c1 :: 1-AUC, mean +- SD over 5 seeds
\newcommand{\ProbeBkgLeftBvcQcdLinearOnesixtwo}{121.4}  % experiments/FIGS/data/probe_ladder_v2_mlp2/analysis/seed_level_results.json :: levels.bvc_qcd.linear[1].rejection_points['0.90'].mean_n_bkg_pass :: sha256 53e1ec81343ac2c1 :: background jets passing the cut, mean over seeds
\newcommand{\ProbeRejBvcQcdLinearOnesixtwo}{44.3\,\ensuremath{\pm}\,2.1}  % experiments/FIGS/data/probe_ladder_v2_mlp2/analysis/seed_level_results.json :: table[bvc_qcd,linear,162].rejection_points['0.90'].rejection :: sha256 53e1ec81343ac2c1 :: at 90% signal efficiency, mean +- SD over seeds; 0 of 5 seeds at the cap
\newcommand{\ProbeAucBvcQcdLinearFourthree}{0.98281}  % experiments/FIGS/data/probe_ladder_v2_mlp2/analysis/seed_level_results.json :: levels.bvc_qcd.linear[2].mean_auc :: sha256 53e1ec81343ac2c1 :: 5 seeds, 0 saturated at AUC=1
\newcommand{\ProbeAucSdBvcQcdLinearFourthree}{0.00064}  % experiments/FIGS/data/probe_ladder_v2_mlp2/analysis/seed_level_results.json :: table[bvc_qcd,linear,43].auc (SD over seeds, ddof=1) :: sha256 53e1ec81343ac2c1
\newcommand{\ProbeLogOneMinusAucBvcQcdLinearFourthree}{\ensuremath{-}4.0638}  % experiments/FIGS/data/probe_ladder_v2_mlp2/analysis/seed_level_results.json :: levels.bvc_qcd.linear[2].mean :: sha256 53e1ec81343ac2c1 :: natural log of 1-AUC, lower is better
\newcommand{\ProbeLogOneMinusAucSdBvcQcdLinearFourthree}{0.0370}  % experiments/FIGS/data/probe_ladder_v2_mlp2/analysis/seed_level_results.json :: levels.bvc_qcd.linear[2].seed_sd :: sha256 53e1ec81343ac2c1
\newcommand{\ProbeOmaBvcQcdLinearFourthree}{\ensuremath{(1.72 \pm 0.06)\times10^{-2}}}  % experiments/FIGS/data/probe_ladder_v2_mlp2/analysis/seed_level_results.json :: table[bvc_qcd,linear,43].auc :: sha256 53e1ec81343ac2c1 :: 1-AUC, mean +- SD over 5 seeds
\newcommand{\ProbeBkgLeftBvcQcdLinearFourthree}{126.8}  % experiments/FIGS/data/probe_ladder_v2_mlp2/analysis/seed_level_results.json :: levels.bvc_qcd.linear[2].rejection_points['0.90'].mean_n_bkg_pass :: sha256 53e1ec81343ac2c1 :: background jets passing the cut, mean over seeds
\newcommand{\ProbeRejBvcQcdLinearFourthree}{43\,\ensuremath{\pm}\,4}  % experiments/FIGS/data/probe_ladder_v2_mlp2/analysis/seed_level_results.json :: table[bvc_qcd,linear,43].rejection_points['0.90'].rejection :: sha256 53e1ec81343ac2c1 :: at 90% signal efficiency, mean +- SD over seeds; 0 of 5 seeds at the cap
\newcommand{\ProbeAucBvcQcdLinearOneseven}{0.95504}  % experiments/FIGS/data/probe_ladder_v2_mlp2/analysis/seed_level_results.json :: levels.bvc_qcd.linear[3].mean_auc :: sha256 53e1ec81343ac2c1 :: 5 seeds, 0 saturated at AUC=1
\newcommand{\ProbeAucSdBvcQcdLinearOneseven}{0.00527}  % experiments/FIGS/data/probe_ladder_v2_mlp2/analysis/seed_level_results.json :: table[bvc_qcd,linear,17].auc (SD over seeds, ddof=1) :: sha256 53e1ec81343ac2c1
\newcommand{\ProbeLogOneMinusAucBvcQcdLinearOneseven}{\ensuremath{-}3.1076}  % experiments/FIGS/data/probe_ladder_v2_mlp2/analysis/seed_level_results.json :: levels.bvc_qcd.linear[3].mean :: sha256 53e1ec81343ac2c1 :: natural log of 1-AUC, lower is better
\newcommand{\ProbeLogOneMinusAucSdBvcQcdLinearOneseven}{0.1198}  % experiments/FIGS/data/probe_ladder_v2_mlp2/analysis/seed_level_results.json :: levels.bvc_qcd.linear[3].seed_sd :: sha256 53e1ec81343ac2c1
\newcommand{\ProbeOmaBvcQcdLinearOneseven}{\ensuremath{(4.5 \pm 0.5)\times10^{-2}}}  % experiments/FIGS/data/probe_ladder_v2_mlp2/analysis/seed_level_results.json :: table[bvc_qcd,linear,17].auc :: sha256 53e1ec81343ac2c1 :: 1-AUC, mean +- SD over 5 seeds
\newcommand{\ProbeBkgLeftBvcQcdLinearOneseven}{544.6}  % experiments/FIGS/data/probe_ladder_v2_mlp2/analysis/seed_level_results.json :: levels.bvc_qcd.linear[3].rejection_points['0.90'].mean_n_bkg_pass :: sha256 53e1ec81343ac2c1 :: background jets passing the cut, mean over seeds
\newcommand{\ProbeRejBvcQcdLinearOneseven}{10.1\,\ensuremath{\pm}\,1.9}  % experiments/FIGS/data/probe_ladder_v2_mlp2/analysis/seed_level_results.json :: table[bvc_qcd,linear,17].rejection_points['0.90'].rejection :: sha256 53e1ec81343ac2c1 :: at 90% signal efficiency, mean +- SD over seeds; 0 of 5 seeds at the cap
\newcommand{\ProbeOmaRatioBvcQcdLinearOnesixtwoOverOneeighteight}{1.14}  % experiments/FIGS/data/probe_ladder_v2_mlp2/analysis/seed_level_results.json :: table[bvc_qcd,linear,162/188].auc :: sha256 53e1ec81343ac2c1 :: ratio of the seed means of 1-AUC, coarser over finer
\newcommand{\ProbeOmaRatioBvcQcdLinearFourthreeOverOneeighteight}{1.16}  % experiments/FIGS/data/probe_ladder_v2_mlp2/analysis/seed_level_results.json :: table[bvc_qcd,linear,43/188].auc :: sha256 53e1ec81343ac2c1 :: ratio of the seed means of 1-AUC, coarser over finer
\newcommand{\ProbeOmaRatioBvcQcdLinearOnesevenOverOneeighteight}{3.0}  % experiments/FIGS/data/probe_ladder_v2_mlp2/analysis/seed_level_results.json :: table[bvc_qcd,linear,17/188].auc :: sha256 53e1ec81343ac2c1 :: ratio of the seed means of 1-AUC, coarser over finer
\newcommand{\ProbeOmaRatioBvcQcdLinearFourthreeOverOnesixtwo}{1.01}  % experiments/FIGS/data/probe_ladder_v2_mlp2/analysis/seed_level_results.json :: table[bvc_qcd,linear,43/162].auc :: sha256 53e1ec81343ac2c1 :: ratio of the seed means of 1-AUC, coarser over finer
\newcommand{\ProbeOmaRatioBvcQcdLinearOnesevenOverOnesixtwo}{2.7}  % experiments/FIGS/data/probe_ladder_v2_mlp2/analysis/seed_level_results.json :: table[bvc_qcd,linear,17/162].auc :: sha256 53e1ec81343ac2c1 :: ratio of the seed means of 1-AUC, coarser over finer
\newcommand{\ProbeOmaRatioBvcQcdLinearOnesevenOverFourthree}{2.6}  % experiments/FIGS/data/probe_ladder_v2_mlp2/analysis/seed_level_results.json :: table[bvc_qcd,linear,17/43].auc :: sha256 53e1ec81343ac2c1 :: ratio of the seed means of 1-AUC, coarser over finer
\newcommand{\ProbeAucBvcQcdMlpOneeighteight}{0.98527}  % experiments/FIGS/data/probe_ladder_v2_mlp2/analysis/seed_level_results.json :: levels.bvc_qcd.mlp[0].mean_auc :: sha256 53e1ec81343ac2c1 :: 5 seeds, 0 saturated at AUC=1
\newcommand{\ProbeAucSdBvcQcdMlpOneeighteight}{0.00058}  % experiments/FIGS/data/probe_ladder_v2_mlp2/analysis/seed_level_results.json :: table[bvc_qcd,mlp,188].auc (SD over seeds, ddof=1) :: sha256 53e1ec81343ac2c1
\newcommand{\ProbeLogOneMinusAucBvcQcdMlpOneeighteight}{\ensuremath{-}4.2188}  % experiments/FIGS/data/probe_ladder_v2_mlp2/analysis/seed_level_results.json :: levels.bvc_qcd.mlp[0].mean :: sha256 53e1ec81343ac2c1 :: natural log of 1-AUC, lower is better
\newcommand{\ProbeLogOneMinusAucSdBvcQcdMlpOneeighteight}{0.0400}  % experiments/FIGS/data/probe_ladder_v2_mlp2/analysis/seed_level_results.json :: levels.bvc_qcd.mlp[0].seed_sd :: sha256 53e1ec81343ac2c1
\newcommand{\ProbeOmaBvcQcdMlpOneeighteight}{\ensuremath{(1.47 \pm 0.06)\times10^{-2}}}  % experiments/FIGS/data/probe_ladder_v2_mlp2/analysis/seed_level_results.json :: table[bvc_qcd,mlp,188].auc :: sha256 53e1ec81343ac2c1 :: 1-AUC, mean +- SD over 5 seeds
\newcommand{\ProbeBkgLeftBvcQcdMlpOneeighteight}{94.6}  % experiments/FIGS/data/probe_ladder_v2_mlp2/analysis/seed_level_results.json :: levels.bvc_qcd.mlp[0].rejection_points['0.90'].mean_n_bkg_pass :: sha256 53e1ec81343ac2c1 :: background jets passing the cut, mean over seeds
\newcommand{\ProbeRejBvcQcdMlpOneeighteight}{57\,\ensuremath{\pm}\,6}  % experiments/FIGS/data/probe_ladder_v2_mlp2/analysis/seed_level_results.json :: table[bvc_qcd,mlp,188].rejection_points['0.90'].rejection :: sha256 53e1ec81343ac2c1 :: at 90% signal efficiency, mean +- SD over seeds; 0 of 5 seeds at the cap
\newcommand{\ProbeAucBvcQcdMlpOnesixtwo}{0.98361}  % experiments/FIGS/data/probe_ladder_v2_mlp2/analysis/seed_level_results.json :: levels.bvc_qcd.mlp[1].mean_auc :: sha256 53e1ec81343ac2c1 :: 5 seeds, 0 saturated at AUC=1
\newcommand{\ProbeAucSdBvcQcdMlpOnesixtwo}{0.00039}  % experiments/FIGS/data/probe_ladder_v2_mlp2/analysis/seed_level_results.json :: table[bvc_qcd,mlp,162].auc (SD over seeds, ddof=1) :: sha256 53e1ec81343ac2c1
\newcommand{\ProbeLogOneMinusAucBvcQcdMlpOnesixtwo}{\ensuremath{-}4.1113}  % experiments/FIGS/data/probe_ladder_v2_mlp2/analysis/seed_level_results.json :: levels.bvc_qcd.mlp[1].mean :: sha256 53e1ec81343ac2c1 :: natural log of 1-AUC, lower is better
\newcommand{\ProbeLogOneMinusAucSdBvcQcdMlpOnesixtwo}{0.0241}  % experiments/FIGS/data/probe_ladder_v2_mlp2/analysis/seed_level_results.json :: levels.bvc_qcd.mlp[1].seed_sd :: sha256 53e1ec81343ac2c1
\newcommand{\ProbeOmaBvcQcdMlpOnesixtwo}{\ensuremath{(1.64 \pm 0.04)\times10^{-2}}}  % experiments/FIGS/data/probe_ladder_v2_mlp2/analysis/seed_level_results.json :: table[bvc_qcd,mlp,162].auc :: sha256 53e1ec81343ac2c1 :: 1-AUC, mean +- SD over 5 seeds
\newcommand{\ProbeBkgLeftBvcQcdMlpOnesixtwo}{119.6}  % experiments/FIGS/data/probe_ladder_v2_mlp2/analysis/seed_level_results.json :: levels.bvc_qcd.mlp[1].rejection_points['0.90'].mean_n_bkg_pass :: sha256 53e1ec81343ac2c1 :: background jets passing the cut, mean over seeds
\newcommand{\ProbeRejBvcQcdMlpOnesixtwo}{45.1\,\ensuremath{\pm}\,3.2}  % experiments/FIGS/data/probe_ladder_v2_mlp2/analysis/seed_level_results.json :: table[bvc_qcd,mlp,162].rejection_points['0.90'].rejection :: sha256 53e1ec81343ac2c1 :: at 90% signal efficiency, mean +- SD over seeds; 0 of 5 seeds at the cap
\newcommand{\ProbeAucBvcQcdMlpFourthree}{0.98341}  % experiments/FIGS/data/probe_ladder_v2_mlp2/analysis/seed_level_results.json :: levels.bvc_qcd.mlp[2].mean_auc :: sha256 53e1ec81343ac2c1 :: 5 seeds, 0 saturated at AUC=1
\newcommand{\ProbeAucSdBvcQcdMlpFourthree}{0.00052}  % experiments/FIGS/data/probe_ladder_v2_mlp2/analysis/seed_level_results.json :: table[bvc_qcd,mlp,43].auc (SD over seeds, ddof=1) :: sha256 53e1ec81343ac2c1
\newcommand{\ProbeLogOneMinusAucBvcQcdMlpFourthree}{\ensuremath{-}4.0995}  % experiments/FIGS/data/probe_ladder_v2_mlp2/analysis/seed_level_results.json :: levels.bvc_qcd.mlp[2].mean :: sha256 53e1ec81343ac2c1 :: natural log of 1-AUC, lower is better
\newcommand{\ProbeLogOneMinusAucSdBvcQcdMlpFourthree}{0.0318}  % experiments/FIGS/data/probe_ladder_v2_mlp2/analysis/seed_level_results.json :: levels.bvc_qcd.mlp[2].seed_sd :: sha256 53e1ec81343ac2c1
\newcommand{\ProbeOmaBvcQcdMlpFourthree}{\ensuremath{(1.66 \pm 0.05)\times10^{-2}}}  % experiments/FIGS/data/probe_ladder_v2_mlp2/analysis/seed_level_results.json :: table[bvc_qcd,mlp,43].auc :: sha256 53e1ec81343ac2c1 :: 1-AUC, mean +- SD over 5 seeds
\newcommand{\ProbeBkgLeftBvcQcdMlpFourthree}{119.2}  % experiments/FIGS/data/probe_ladder_v2_mlp2/analysis/seed_level_results.json :: levels.bvc_qcd.mlp[2].rejection_points['0.90'].mean_n_bkg_pass :: sha256 53e1ec81343ac2c1 :: background jets passing the cut, mean over seeds
\newcommand{\ProbeRejBvcQcdMlpFourthree}{45.1\,\ensuremath{\pm}\,1.3}  % experiments/FIGS/data/probe_ladder_v2_mlp2/analysis/seed_level_results.json :: table[bvc_qcd,mlp,43].rejection_points['0.90'].rejection :: sha256 53e1ec81343ac2c1 :: at 90% signal efficiency, mean +- SD over seeds; 0 of 5 seeds at the cap
\newcommand{\ProbeAucBvcQcdMlpOneseven}{0.96226}  % experiments/FIGS/data/probe_ladder_v2_mlp2/analysis/seed_level_results.json :: levels.bvc_qcd.mlp[3].mean_auc :: sha256 53e1ec81343ac2c1 :: 5 seeds, 0 saturated at AUC=1
\newcommand{\ProbeAucSdBvcQcdMlpOneseven}{0.00373}  % experiments/FIGS/data/probe_ladder_v2_mlp2/analysis/seed_level_results.json :: table[bvc_qcd,mlp,17].auc (SD over seeds, ddof=1) :: sha256 53e1ec81343ac2c1
\newcommand{\ProbeLogOneMinusAucBvcQcdMlpOneseven}{\ensuremath{-}3.2811}  % experiments/FIGS/data/probe_ladder_v2_mlp2/analysis/seed_level_results.json :: levels.bvc_qcd.mlp[3].mean :: sha256 53e1ec81343ac2c1 :: natural log of 1-AUC, lower is better
\newcommand{\ProbeLogOneMinusAucSdBvcQcdMlpOneseven}{0.1010}  % experiments/FIGS/data/probe_ladder_v2_mlp2/analysis/seed_level_results.json :: levels.bvc_qcd.mlp[3].seed_sd :: sha256 53e1ec81343ac2c1
\newcommand{\ProbeOmaBvcQcdMlpOneseven}{\ensuremath{(3.8 \pm 0.4)\times10^{-2}}}  % experiments/FIGS/data/probe_ladder_v2_mlp2/analysis/seed_level_results.json :: table[bvc_qcd,mlp,17].auc :: sha256 53e1ec81343ac2c1 :: 1-AUC, mean +- SD over 5 seeds
\newcommand{\ProbeBkgLeftBvcQcdMlpOneseven}{458.8}  % experiments/FIGS/data/probe_ladder_v2_mlp2/analysis/seed_level_results.json :: levels.bvc_qcd.mlp[3].rejection_points['0.90'].mean_n_bkg_pass :: sha256 53e1ec81343ac2c1 :: background jets passing the cut, mean over seeds
\newcommand{\ProbeRejBvcQcdMlpOneseven}{11.8\,\ensuremath{\pm}\,1.5}  % experiments/FIGS/data/probe_ladder_v2_mlp2/analysis/seed_level_results.json :: table[bvc_qcd,mlp,17].rejection_points['0.90'].rejection :: sha256 53e1ec81343ac2c1 :: at 90% signal efficiency, mean +- SD over seeds; 0 of 5 seeds at the cap
\newcommand{\ProbeOmaRatioBvcQcdMlpOnesixtwoOverOneeighteight}{1.11}  % experiments/FIGS/data/probe_ladder_v2_mlp2/analysis/seed_level_results.json :: table[bvc_qcd,mlp,162/188].auc :: sha256 53e1ec81343ac2c1 :: ratio of the seed means of 1-AUC, coarser over finer
\newcommand{\ProbeOmaRatioBvcQcdMlpFourthreeOverOneeighteight}{1.13}  % experiments/FIGS/data/probe_ladder_v2_mlp2/analysis/seed_level_results.json :: table[bvc_qcd,mlp,43/188].auc :: sha256 53e1ec81343ac2c1 :: ratio of the seed means of 1-AUC, coarser over finer
\newcommand{\ProbeOmaRatioBvcQcdMlpOnesevenOverOneeighteight}{2.6}  % experiments/FIGS/data/probe_ladder_v2_mlp2/analysis/seed_level_results.json :: table[bvc_qcd,mlp,17/188].auc :: sha256 53e1ec81343ac2c1 :: ratio of the seed means of 1-AUC, coarser over finer
\newcommand{\ProbeOmaRatioBvcQcdMlpFourthreeOverOnesixtwo}{1.01}  % experiments/FIGS/data/probe_ladder_v2_mlp2/analysis/seed_level_results.json :: table[bvc_qcd,mlp,43/162].auc :: sha256 53e1ec81343ac2c1 :: ratio of the seed means of 1-AUC, coarser over finer
\newcommand{\ProbeOmaRatioBvcQcdMlpOnesevenOverOnesixtwo}{2.3}  % experiments/FIGS/data/probe_ladder_v2_mlp2/analysis/seed_level_results.json :: table[bvc_qcd,mlp,17/162].auc :: sha256 53e1ec81343ac2c1 :: ratio of the seed means of 1-AUC, coarser over finer
\newcommand{\ProbeOmaRatioBvcQcdMlpOnesevenOverFourthree}{2.3}  % experiments/FIGS/data/probe_ladder_v2_mlp2/analysis/seed_level_results.json :: table[bvc_qcd,mlp,17/43].auc :: sha256 53e1ec81343ac2c1 :: ratio of the seed means of 1-AUC, coarser over finer
\newcommand{\ProbeAucBvcResonantLinearOneeighteight}{0.99768}  % experiments/FIGS/data/probe_ladder_v2_mlp2/analysis/seed_level_results.json :: levels.bvc_resonant.linear[0].mean_auc :: sha256 53e1ec81343ac2c1 :: 5 seeds, 0 saturated at AUC=1
\newcommand{\ProbeAucSdBvcResonantLinearOneeighteight}{0.00007}  % experiments/FIGS/data/probe_ladder_v2_mlp2/analysis/seed_level_results.json :: table[bvc_resonant,linear,188].auc (SD over seeds, ddof=1) :: sha256 53e1ec81343ac2c1
\newcommand{\ProbeLogOneMinusAucBvcResonantLinearOneeighteight}{\ensuremath{-}6.0683}  % experiments/FIGS/data/probe_ladder_v2_mlp2/analysis/seed_level_results.json :: levels.bvc_resonant.linear[0].mean :: sha256 53e1ec81343ac2c1 :: natural log of 1-AUC, lower is better
\newcommand{\ProbeLogOneMinusAucSdBvcResonantLinearOneeighteight}{0.0317}  % experiments/FIGS/data/probe_ladder_v2_mlp2/analysis/seed_level_results.json :: levels.bvc_resonant.linear[0].seed_sd :: sha256 53e1ec81343ac2c1
\newcommand{\ProbeOmaBvcResonantLinearOneeighteight}{\ensuremath{(2.32 \pm 0.07)\times10^{-3}}}  % experiments/FIGS/data/probe_ladder_v2_mlp2/analysis/seed_level_results.json :: table[bvc_resonant,linear,188].auc :: sha256 53e1ec81343ac2c1 :: 1-AUC, mean +- SD over 5 seeds
\newcommand{\ProbeBkgLeftBvcResonantLinearOneeighteight}{9.0}  % experiments/FIGS/data/probe_ladder_v2_mlp2/analysis/seed_level_results.json :: levels.bvc_resonant.linear[0].rejection_points['0.90'].mean_n_bkg_pass :: sha256 53e1ec81343ac2c1 :: background jets passing the cut, mean over seeds
\newcommand{\ProbeRejBvcResonantLinearOneeighteight}{1{,}380\,\ensuremath{\pm}\,350}  % experiments/FIGS/data/probe_ladder_v2_mlp2/analysis/seed_level_results.json :: table[bvc_resonant,linear,188].rejection_points['0.90'].rejection :: sha256 53e1ec81343ac2c1 :: at 90% signal efficiency, mean +- SD over seeds; 0 of 5 seeds at the cap
\newcommand{\ProbeAucBvcResonantLinearOnesixtwo}{0.99764}  % experiments/FIGS/data/probe_ladder_v2_mlp2/analysis/seed_level_results.json :: levels.bvc_resonant.linear[1].mean_auc :: sha256 53e1ec81343ac2c1 :: 5 seeds, 0 saturated at AUC=1
\newcommand{\ProbeAucSdBvcResonantLinearOnesixtwo}{0.00006}  % experiments/FIGS/data/probe_ladder_v2_mlp2/analysis/seed_level_results.json :: table[bvc_resonant,linear,162].auc (SD over seeds, ddof=1) :: sha256 53e1ec81343ac2c1
\newcommand{\ProbeLogOneMinusAucBvcResonantLinearOnesixtwo}{\ensuremath{-}6.0493}  % experiments/FIGS/data/probe_ladder_v2_mlp2/analysis/seed_level_results.json :: levels.bvc_resonant.linear[1].mean :: sha256 53e1ec81343ac2c1 :: natural log of 1-AUC, lower is better
\newcommand{\ProbeLogOneMinusAucSdBvcResonantLinearOnesixtwo}{0.0261}  % experiments/FIGS/data/probe_ladder_v2_mlp2/analysis/seed_level_results.json :: levels.bvc_resonant.linear[1].seed_sd :: sha256 53e1ec81343ac2c1
\newcommand{\ProbeOmaBvcResonantLinearOnesixtwo}{\ensuremath{(2.36 \pm 0.06)\times10^{-3}}}  % experiments/FIGS/data/probe_ladder_v2_mlp2/analysis/seed_level_results.json :: table[bvc_resonant,linear,162].auc :: sha256 53e1ec81343ac2c1 :: 1-AUC, mean +- SD over 5 seeds
\newcommand{\ProbeBkgLeftBvcResonantLinearOnesixtwo}{9.4}  % experiments/FIGS/data/probe_ladder_v2_mlp2/analysis/seed_level_results.json :: levels.bvc_resonant.linear[1].rejection_points['0.90'].mean_n_bkg_pass :: sha256 53e1ec81343ac2c1 :: background jets passing the cut, mean over seeds
\newcommand{\ProbeRejBvcResonantLinearOnesixtwo}{1{,}290\,\ensuremath{\pm}\,240}  % experiments/FIGS/data/probe_ladder_v2_mlp2/analysis/seed_level_results.json :: table[bvc_resonant,linear,162].rejection_points['0.90'].rejection :: sha256 53e1ec81343ac2c1 :: at 90% signal efficiency, mean +- SD over seeds; 0 of 5 seeds at the cap
\newcommand{\ProbeAucBvcResonantLinearFourthree}{0.99748}  % experiments/FIGS/data/probe_ladder_v2_mlp2/analysis/seed_level_results.json :: levels.bvc_resonant.linear[2].mean_auc :: sha256 53e1ec81343ac2c1 :: 5 seeds, 0 saturated at AUC=1
\newcommand{\ProbeAucSdBvcResonantLinearFourthree}{0.00007}  % experiments/FIGS/data/probe_ladder_v2_mlp2/analysis/seed_level_results.json :: table[bvc_resonant,linear,43].auc (SD over seeds, ddof=1) :: sha256 53e1ec81343ac2c1
\newcommand{\ProbeLogOneMinusAucBvcResonantLinearFourthree}{\ensuremath{-}5.9834}  % experiments/FIGS/data/probe_ladder_v2_mlp2/analysis/seed_level_results.json :: levels.bvc_resonant.linear[2].mean :: sha256 53e1ec81343ac2c1 :: natural log of 1-AUC, lower is better
\newcommand{\ProbeLogOneMinusAucSdBvcResonantLinearFourthree}{0.0264}  % experiments/FIGS/data/probe_ladder_v2_mlp2/analysis/seed_level_results.json :: levels.bvc_resonant.linear[2].seed_sd :: sha256 53e1ec81343ac2c1
\newcommand{\ProbeOmaBvcResonantLinearFourthree}{\ensuremath{(2.52 \pm 0.07)\times10^{-3}}}  % experiments/FIGS/data/probe_ladder_v2_mlp2/analysis/seed_level_results.json :: table[bvc_resonant,linear,43].auc :: sha256 53e1ec81343ac2c1 :: 1-AUC, mean +- SD over 5 seeds
\newcommand{\ProbeBkgLeftBvcResonantLinearFourthree}{11.0}  % experiments/FIGS/data/probe_ladder_v2_mlp2/analysis/seed_level_results.json :: levels.bvc_resonant.linear[2].rejection_points['0.90'].mean_n_bkg_pass :: sha256 53e1ec81343ac2c1 :: background jets passing the cut, mean over seeds
\newcommand{\ProbeRejBvcResonantLinearFourthree}{1{,}130\,\ensuremath{\pm}\,260}  % experiments/FIGS/data/probe_ladder_v2_mlp2/analysis/seed_level_results.json :: table[bvc_resonant,linear,43].rejection_points['0.90'].rejection :: sha256 53e1ec81343ac2c1 :: at 90% signal efficiency, mean +- SD over seeds; 0 of 5 seeds at the cap
\newcommand{\ProbeAucBvcResonantLinearOneseven}{0.99067}  % experiments/FIGS/data/probe_ladder_v2_mlp2/analysis/seed_level_results.json :: levels.bvc_resonant.linear[3].mean_auc :: sha256 53e1ec81343ac2c1 :: 5 seeds, 0 saturated at AUC=1
\newcommand{\ProbeAucSdBvcResonantLinearOneseven}{0.00147}  % experiments/FIGS/data/probe_ladder_v2_mlp2/analysis/seed_level_results.json :: table[bvc_resonant,linear,17].auc (SD over seeds, ddof=1) :: sha256 53e1ec81343ac2c1
\newcommand{\ProbeLogOneMinusAucBvcResonantLinearOneseven}{\ensuremath{-}4.6848}  % experiments/FIGS/data/probe_ladder_v2_mlp2/analysis/seed_level_results.json :: levels.bvc_resonant.linear[3].mean :: sha256 53e1ec81343ac2c1 :: natural log of 1-AUC, lower is better
\newcommand{\ProbeLogOneMinusAucSdBvcResonantLinearOneseven}{0.1619}  % experiments/FIGS/data/probe_ladder_v2_mlp2/analysis/seed_level_results.json :: levels.bvc_resonant.linear[3].seed_sd :: sha256 53e1ec81343ac2c1
\newcommand{\ProbeOmaBvcResonantLinearOneseven}{\ensuremath{(9.3 \pm 1.5)\times10^{-3}}}  % experiments/FIGS/data/probe_ladder_v2_mlp2/analysis/seed_level_results.json :: table[bvc_resonant,linear,17].auc :: sha256 53e1ec81343ac2c1 :: 1-AUC, mean +- SD over 5 seeds
\newcommand{\ProbeBkgLeftBvcResonantLinearOneseven}{90.8}  % experiments/FIGS/data/probe_ladder_v2_mlp2/analysis/seed_level_results.json :: levels.bvc_resonant.linear[3].rejection_points['0.90'].mean_n_bkg_pass :: sha256 53e1ec81343ac2c1 :: background jets passing the cut, mean over seeds
\newcommand{\ProbeRejBvcResonantLinearOneseven}{134\,\ensuremath{\pm}\,21}  % experiments/FIGS/data/probe_ladder_v2_mlp2/analysis/seed_level_results.json :: table[bvc_resonant,linear,17].rejection_points['0.90'].rejection :: sha256 53e1ec81343ac2c1 :: at 90% signal efficiency, mean +- SD over seeds; 0 of 5 seeds at the cap
\newcommand{\ProbeOmaRatioBvcResonantLinearOnesixtwoOverOneeighteight}{1.02}  % experiments/FIGS/data/probe_ladder_v2_mlp2/analysis/seed_level_results.json :: table[bvc_resonant,linear,162/188].auc :: sha256 53e1ec81343ac2c1 :: ratio of the seed means of 1-AUC, coarser over finer
\newcommand{\ProbeOmaRatioBvcResonantLinearFourthreeOverOneeighteight}{1.09}  % experiments/FIGS/data/probe_ladder_v2_mlp2/analysis/seed_level_results.json :: table[bvc_resonant,linear,43/188].auc :: sha256 53e1ec81343ac2c1 :: ratio of the seed means of 1-AUC, coarser over finer
\newcommand{\ProbeOmaRatioBvcResonantLinearOnesevenOverOneeighteight}{4.0}  % experiments/FIGS/data/probe_ladder_v2_mlp2/analysis/seed_level_results.json :: table[bvc_resonant,linear,17/188].auc :: sha256 53e1ec81343ac2c1 :: ratio of the seed means of 1-AUC, coarser over finer
\newcommand{\ProbeOmaRatioBvcResonantLinearFourthreeOverOnesixtwo}{1.07}  % experiments/FIGS/data/probe_ladder_v2_mlp2/analysis/seed_level_results.json :: table[bvc_resonant,linear,43/162].auc :: sha256 53e1ec81343ac2c1 :: ratio of the seed means of 1-AUC, coarser over finer
\newcommand{\ProbeOmaRatioBvcResonantLinearOnesevenOverOnesixtwo}{4.0}  % experiments/FIGS/data/probe_ladder_v2_mlp2/analysis/seed_level_results.json :: table[bvc_resonant,linear,17/162].auc :: sha256 53e1ec81343ac2c1 :: ratio of the seed means of 1-AUC, coarser over finer
\newcommand{\ProbeOmaRatioBvcResonantLinearOnesevenOverFourthree}{3.7}  % experiments/FIGS/data/probe_ladder_v2_mlp2/analysis/seed_level_results.json :: table[bvc_resonant,linear,17/43].auc :: sha256 53e1ec81343ac2c1 :: ratio of the seed means of 1-AUC, coarser over finer
\newcommand{\ProbeAucBvcResonantMlpOneeighteight}{0.99778}  % experiments/FIGS/data/probe_ladder_v2_mlp2/analysis/seed_level_results.json :: levels.bvc_resonant.mlp[0].mean_auc :: sha256 53e1ec81343ac2c1 :: 5 seeds, 0 saturated at AUC=1
\newcommand{\ProbeAucSdBvcResonantMlpOneeighteight}{0.00009}  % experiments/FIGS/data/probe_ladder_v2_mlp2/analysis/seed_level_results.json :: table[bvc_resonant,mlp,188].auc (SD over seeds, ddof=1) :: sha256 53e1ec81343ac2c1
\newcommand{\ProbeLogOneMinusAucBvcResonantMlpOneeighteight}{\ensuremath{-}6.1102}  % experiments/FIGS/data/probe_ladder_v2_mlp2/analysis/seed_level_results.json :: levels.bvc_resonant.mlp[0].mean :: sha256 53e1ec81343ac2c1 :: natural log of 1-AUC, lower is better
\newcommand{\ProbeLogOneMinusAucSdBvcResonantMlpOneeighteight}{0.0383}  % experiments/FIGS/data/probe_ladder_v2_mlp2/analysis/seed_level_results.json :: levels.bvc_resonant.mlp[0].seed_sd :: sha256 53e1ec81343ac2c1
\newcommand{\ProbeOmaBvcResonantMlpOneeighteight}{\ensuremath{(2.22 \pm 0.09)\times10^{-3}}}  % experiments/FIGS/data/probe_ladder_v2_mlp2/analysis/seed_level_results.json :: table[bvc_resonant,mlp,188].auc :: sha256 53e1ec81343ac2c1 :: 1-AUC, mean +- SD over 5 seeds
\newcommand{\ProbeBkgLeftBvcResonantMlpOneeighteight}{9.8}  % experiments/FIGS/data/probe_ladder_v2_mlp2/analysis/seed_level_results.json :: levels.bvc_resonant.mlp[0].rejection_points['0.90'].mean_n_bkg_pass :: sha256 53e1ec81343ac2c1 :: background jets passing the cut, mean over seeds
\newcommand{\ProbeRejBvcResonantMlpOneeighteight}{1{,}230\,\ensuremath{\pm}\,180}  % experiments/FIGS/data/probe_ladder_v2_mlp2/analysis/seed_level_results.json :: table[bvc_resonant,mlp,188].rejection_points['0.90'].rejection :: sha256 53e1ec81343ac2c1 :: at 90% signal efficiency, mean +- SD over seeds; 0 of 5 seeds at the cap
\newcommand{\ProbeAucBvcResonantMlpOnesixtwo}{0.99768}  % experiments/FIGS/data/probe_ladder_v2_mlp2/analysis/seed_level_results.json :: levels.bvc_resonant.mlp[1].mean_auc :: sha256 53e1ec81343ac2c1 :: 5 seeds, 0 saturated at AUC=1
\newcommand{\ProbeAucSdBvcResonantMlpOnesixtwo}{0.00005}  % experiments/FIGS/data/probe_ladder_v2_mlp2/analysis/seed_level_results.json :: table[bvc_resonant,mlp,162].auc (SD over seeds, ddof=1) :: sha256 53e1ec81343ac2c1
\newcommand{\ProbeLogOneMinusAucBvcResonantMlpOnesixtwo}{\ensuremath{-}6.0667}  % experiments/FIGS/data/probe_ladder_v2_mlp2/analysis/seed_level_results.json :: levels.bvc_resonant.mlp[1].mean :: sha256 53e1ec81343ac2c1 :: natural log of 1-AUC, lower is better
\newcommand{\ProbeLogOneMinusAucSdBvcResonantMlpOnesixtwo}{0.0220}  % experiments/FIGS/data/probe_ladder_v2_mlp2/analysis/seed_level_results.json :: levels.bvc_resonant.mlp[1].seed_sd :: sha256 53e1ec81343ac2c1
\newcommand{\ProbeOmaBvcResonantMlpOnesixtwo}{\ensuremath{(2.32 \pm 0.05)\times10^{-3}}}  % experiments/FIGS/data/probe_ladder_v2_mlp2/analysis/seed_level_results.json :: table[bvc_resonant,mlp,162].auc :: sha256 53e1ec81343ac2c1 :: 1-AUC, mean +- SD over 5 seeds
\newcommand{\ProbeBkgLeftBvcResonantMlpOnesixtwo}{10.4}  % experiments/FIGS/data/probe_ladder_v2_mlp2/analysis/seed_level_results.json :: levels.bvc_resonant.mlp[1].rejection_points['0.90'].mean_n_bkg_pass :: sha256 53e1ec81343ac2c1 :: background jets passing the cut, mean over seeds
\newcommand{\ProbeRejBvcResonantMlpOnesixtwo}{1{,}170\,\ensuremath{\pm}\,220}  % experiments/FIGS/data/probe_ladder_v2_mlp2/analysis/seed_level_results.json :: table[bvc_resonant,mlp,162].rejection_points['0.90'].rejection :: sha256 53e1ec81343ac2c1 :: at 90% signal efficiency, mean +- SD over seeds; 0 of 5 seeds at the cap
\newcommand{\ProbeAucBvcResonantMlpFourthree}{0.99747}  % experiments/FIGS/data/probe_ladder_v2_mlp2/analysis/seed_level_results.json :: levels.bvc_resonant.mlp[2].mean_auc :: sha256 53e1ec81343ac2c1 :: 5 seeds, 0 saturated at AUC=1
\newcommand{\ProbeAucSdBvcResonantMlpFourthree}{0.00004}  % experiments/FIGS/data/probe_ladder_v2_mlp2/analysis/seed_level_results.json :: table[bvc_resonant,mlp,43].auc (SD over seeds, ddof=1) :: sha256 53e1ec81343ac2c1
\newcommand{\ProbeLogOneMinusAucBvcResonantMlpFourthree}{\ensuremath{-}5.9792}  % experiments/FIGS/data/probe_ladder_v2_mlp2/analysis/seed_level_results.json :: levels.bvc_resonant.mlp[2].mean :: sha256 53e1ec81343ac2c1 :: natural log of 1-AUC, lower is better
\newcommand{\ProbeLogOneMinusAucSdBvcResonantMlpFourthree}{0.0164}  % experiments/FIGS/data/probe_ladder_v2_mlp2/analysis/seed_level_results.json :: levels.bvc_resonant.mlp[2].seed_sd :: sha256 53e1ec81343ac2c1
\newcommand{\ProbeOmaBvcResonantMlpFourthree}{\ensuremath{(2.53 \pm 0.04)\times10^{-3}}}  % experiments/FIGS/data/probe_ladder_v2_mlp2/analysis/seed_level_results.json :: table[bvc_resonant,mlp,43].auc :: sha256 53e1ec81343ac2c1 :: 1-AUC, mean +- SD over 5 seeds
\newcommand{\ProbeBkgLeftBvcResonantMlpFourthree}{11.8}  % experiments/FIGS/data/probe_ladder_v2_mlp2/analysis/seed_level_results.json :: levels.bvc_resonant.mlp[2].rejection_points['0.90'].mean_n_bkg_pass :: sha256 53e1ec81343ac2c1 :: background jets passing the cut, mean over seeds
\newcommand{\ProbeRejBvcResonantMlpFourthree}{1{,}030\,\ensuremath{\pm}\,170}  % experiments/FIGS/data/probe_ladder_v2_mlp2/analysis/seed_level_results.json :: table[bvc_resonant,mlp,43].rejection_points['0.90'].rejection :: sha256 53e1ec81343ac2c1 :: at 90% signal efficiency, mean +- SD over seeds; 0 of 5 seeds at the cap
\newcommand{\ProbeAucBvcResonantMlpOneseven}{0.99274}  % experiments/FIGS/data/probe_ladder_v2_mlp2/analysis/seed_level_results.json :: levels.bvc_resonant.mlp[3].mean_auc :: sha256 53e1ec81343ac2c1 :: 5 seeds, 0 saturated at AUC=1
\newcommand{\ProbeAucSdBvcResonantMlpOneseven}{0.00089}  % experiments/FIGS/data/probe_ladder_v2_mlp2/analysis/seed_level_results.json :: table[bvc_resonant,mlp,17].auc (SD over seeds, ddof=1) :: sha256 53e1ec81343ac2c1
\newcommand{\ProbeLogOneMinusAucBvcResonantMlpOneseven}{\ensuremath{-}4.9309}  % experiments/FIGS/data/probe_ladder_v2_mlp2/analysis/seed_level_results.json :: levels.bvc_resonant.mlp[3].mean :: sha256 53e1ec81343ac2c1 :: natural log of 1-AUC, lower is better
\newcommand{\ProbeLogOneMinusAucSdBvcResonantMlpOneseven}{0.1224}  % experiments/FIGS/data/probe_ladder_v2_mlp2/analysis/seed_level_results.json :: levels.bvc_resonant.mlp[3].seed_sd :: sha256 53e1ec81343ac2c1
\newcommand{\ProbeOmaBvcResonantMlpOneseven}{\ensuremath{(7.3 \pm 0.9)\times10^{-3}}}  % experiments/FIGS/data/probe_ladder_v2_mlp2/analysis/seed_level_results.json :: table[bvc_resonant,mlp,17].auc :: sha256 53e1ec81343ac2c1 :: 1-AUC, mean +- SD over 5 seeds
\newcommand{\ProbeBkgLeftBvcResonantMlpOneseven}{62.6}  % experiments/FIGS/data/probe_ladder_v2_mlp2/analysis/seed_level_results.json :: levels.bvc_resonant.mlp[3].rejection_points['0.90'].mean_n_bkg_pass :: sha256 53e1ec81343ac2c1 :: background jets passing the cut, mean over seeds
\newcommand{\ProbeRejBvcResonantMlpOneseven}{192\,\ensuremath{\pm}\,21}  % experiments/FIGS/data/probe_ladder_v2_mlp2/analysis/seed_level_results.json :: table[bvc_resonant,mlp,17].rejection_points['0.90'].rejection :: sha256 53e1ec81343ac2c1 :: at 90% signal efficiency, mean +- SD over seeds; 0 of 5 seeds at the cap
\newcommand{\ProbeOmaRatioBvcResonantMlpOnesixtwoOverOneeighteight}{1.04}  % experiments/FIGS/data/probe_ladder_v2_mlp2/analysis/seed_level_results.json :: table[bvc_resonant,mlp,162/188].auc :: sha256 53e1ec81343ac2c1 :: ratio of the seed means of 1-AUC, coarser over finer
\newcommand{\ProbeOmaRatioBvcResonantMlpFourthreeOverOneeighteight}{1.14}  % experiments/FIGS/data/probe_ladder_v2_mlp2/analysis/seed_level_results.json :: table[bvc_resonant,mlp,43/188].auc :: sha256 53e1ec81343ac2c1 :: ratio of the seed means of 1-AUC, coarser over finer
\newcommand{\ProbeOmaRatioBvcResonantMlpOnesevenOverOneeighteight}{3.3}  % experiments/FIGS/data/probe_ladder_v2_mlp2/analysis/seed_level_results.json :: table[bvc_resonant,mlp,17/188].auc :: sha256 53e1ec81343ac2c1 :: ratio of the seed means of 1-AUC, coarser over finer
\newcommand{\ProbeOmaRatioBvcResonantMlpFourthreeOverOnesixtwo}{1.09}  % experiments/FIGS/data/probe_ladder_v2_mlp2/analysis/seed_level_results.json :: table[bvc_resonant,mlp,43/162].auc :: sha256 53e1ec81343ac2c1 :: ratio of the seed means of 1-AUC, coarser over finer
\newcommand{\ProbeOmaRatioBvcResonantMlpOnesevenOverOnesixtwo}{3.1}  % experiments/FIGS/data/probe_ladder_v2_mlp2/analysis/seed_level_results.json :: table[bvc_resonant,mlp,17/162].auc :: sha256 53e1ec81343ac2c1 :: ratio of the seed means of 1-AUC, coarser over finer
\newcommand{\ProbeOmaRatioBvcResonantMlpOnesevenOverFourthree}{2.9}  % experiments/FIGS/data/probe_ladder_v2_mlp2/analysis/seed_level_results.json :: table[bvc_resonant,mlp,17/43].auc :: sha256 53e1ec81343ac2c1 :: ratio of the seed means of 1-AUC, coarser over finer
\newcommand{\ProbeAucEeVsMmLinearOneeighteight}{\ensuremath{1^{\dagger}}}  % experiments/FIGS/data/probe_ladder_v2_mlp2/analysis/seed_level_results.json :: levels.ee_vs_mm.linear[0].mean_auc :: sha256 53e1ec81343ac2c1 :: 5 seeds, 5 saturated at AUC=1
\newcommand{\ProbeAucSdEeVsMmLinearOneeighteight}{0.00000}  % experiments/FIGS/data/probe_ladder_v2_mlp2/analysis/seed_level_results.json :: table[ee_vs_mm,linear,188].auc (SD over seeds, ddof=1) :: sha256 53e1ec81343ac2c1
\newcommand{\ProbeLogOneMinusAucEeVsMmLinearOneeighteight}{\ensuremath{-}16.9579}  % experiments/FIGS/data/probe_ladder_v2_mlp2/analysis/seed_level_results.json :: levels.ee_vs_mm.linear[0].mean :: sha256 53e1ec81343ac2c1 :: natural log of 1-AUC, lower is better
\newcommand{\ProbeLogOneMinusAucSdEeVsMmLinearOneeighteight}{0.0000}  % experiments/FIGS/data/probe_ladder_v2_mlp2/analysis/seed_level_results.json :: levels.ee_vs_mm.linear[0].seed_sd :: sha256 53e1ec81343ac2c1
\newcommand{\ProbeOmaEeVsMmLinearOneeighteight}{---}  % experiments/FIGS/data/probe_ladder_v2_mlp2/analysis/seed_level_results.json :: table[ee_vs_mm,linear,188].auc :: sha256 53e1ec81343ac2c1 :: not quoted: the AUC reached 1 in a seed, so 1-AUC is only a bound
\newcommand{\ProbeBkgLeftEeVsMmLinearOneeighteight}{0.0}  % experiments/FIGS/data/probe_ladder_v2_mlp2/analysis/seed_level_results.json :: levels.ee_vs_mm.linear[0].rejection_points['0.90'].mean_n_bkg_pass :: sha256 53e1ec81343ac2c1 :: background jets passing the cut, mean over seeds
\newcommand{\ProbeRejEeVsMmLinearOneeighteight}{\ensuremath{>}1{,}571}  % experiments/FIGS/data/probe_ladder_v2_mlp2/analysis/seed_level_results.json :: table[ee_vs_mm,linear,188].rejection_points['0.90'].rejection :: sha256 53e1ec81343ac2c1 :: at 90% signal efficiency, mean +- SD over seeds; 5 of 5 seeds at the cap
\newcommand{\ProbeAucEeVsMmLinearOnesixtwo}{\ensuremath{1^{\dagger}}}  % experiments/FIGS/data/probe_ladder_v2_mlp2/analysis/seed_level_results.json :: levels.ee_vs_mm.linear[1].mean_auc :: sha256 53e1ec81343ac2c1 :: 5 seeds, 5 saturated at AUC=1
\newcommand{\ProbeAucSdEeVsMmLinearOnesixtwo}{0.00000}  % experiments/FIGS/data/probe_ladder_v2_mlp2/analysis/seed_level_results.json :: table[ee_vs_mm,linear,162].auc (SD over seeds, ddof=1) :: sha256 53e1ec81343ac2c1
\newcommand{\ProbeLogOneMinusAucEeVsMmLinearOnesixtwo}{\ensuremath{-}16.9579}  % experiments/FIGS/data/probe_ladder_v2_mlp2/analysis/seed_level_results.json :: levels.ee_vs_mm.linear[1].mean :: sha256 53e1ec81343ac2c1 :: natural log of 1-AUC, lower is better
\newcommand{\ProbeLogOneMinusAucSdEeVsMmLinearOnesixtwo}{0.0000}  % experiments/FIGS/data/probe_ladder_v2_mlp2/analysis/seed_level_results.json :: levels.ee_vs_mm.linear[1].seed_sd :: sha256 53e1ec81343ac2c1
\newcommand{\ProbeOmaEeVsMmLinearOnesixtwo}{---}  % experiments/FIGS/data/probe_ladder_v2_mlp2/analysis/seed_level_results.json :: table[ee_vs_mm,linear,162].auc :: sha256 53e1ec81343ac2c1 :: not quoted: the AUC reached 1 in a seed, so 1-AUC is only a bound
\newcommand{\ProbeBkgLeftEeVsMmLinearOnesixtwo}{0.0}  % experiments/FIGS/data/probe_ladder_v2_mlp2/analysis/seed_level_results.json :: levels.ee_vs_mm.linear[1].rejection_points['0.90'].mean_n_bkg_pass :: sha256 53e1ec81343ac2c1 :: background jets passing the cut, mean over seeds
\newcommand{\ProbeRejEeVsMmLinearOnesixtwo}{\ensuremath{>}1{,}571}  % experiments/FIGS/data/probe_ladder_v2_mlp2/analysis/seed_level_results.json :: table[ee_vs_mm,linear,162].rejection_points['0.90'].rejection :: sha256 53e1ec81343ac2c1 :: at 90% signal efficiency, mean +- SD over seeds; 5 of 5 seeds at the cap
\newcommand{\ProbeAucEeVsMmLinearFourthree}{0.97519}  % experiments/FIGS/data/probe_ladder_v2_mlp2/analysis/seed_level_results.json :: levels.ee_vs_mm.linear[2].mean_auc :: sha256 53e1ec81343ac2c1 :: 5 seeds, 0 saturated at AUC=1
\newcommand{\ProbeAucSdEeVsMmLinearFourthree}{0.02073}  % experiments/FIGS/data/probe_ladder_v2_mlp2/analysis/seed_level_results.json :: table[ee_vs_mm,linear,43].auc (SD over seeds, ddof=1) :: sha256 53e1ec81343ac2c1
\newcommand{\ProbeLogOneMinusAucEeVsMmLinearFourthree}{\ensuremath{-}4.0326}  % experiments/FIGS/data/probe_ladder_v2_mlp2/analysis/seed_level_results.json :: levels.ee_vs_mm.linear[2].mean :: sha256 53e1ec81343ac2c1 :: natural log of 1-AUC, lower is better
\newcommand{\ProbeLogOneMinusAucSdEeVsMmLinearFourthree}{0.9467}  % experiments/FIGS/data/probe_ladder_v2_mlp2/analysis/seed_level_results.json :: levels.ee_vs_mm.linear[2].seed_sd :: sha256 53e1ec81343ac2c1
\newcommand{\ProbeOmaEeVsMmLinearFourthree}{\ensuremath{(2.5 \pm 2.1)\times10^{-2}}}  % experiments/FIGS/data/probe_ladder_v2_mlp2/analysis/seed_level_results.json :: table[ee_vs_mm,linear,43].auc :: sha256 53e1ec81343ac2c1 :: 1-AUC, mean +- SD over 5 seeds
\newcommand{\ProbeBkgLeftEeVsMmLinearFourthree}{155.8}  % experiments/FIGS/data/probe_ladder_v2_mlp2/analysis/seed_level_results.json :: levels.ee_vs_mm.linear[2].rejection_points['0.90'].mean_n_bkg_pass :: sha256 53e1ec81343ac2c1 :: background jets passing the cut, mean over seeds
\newcommand{\ProbeRejEeVsMmLinearFourthree}{80\,\ensuremath{\pm}\,70}  % experiments/FIGS/data/probe_ladder_v2_mlp2/analysis/seed_level_results.json :: table[ee_vs_mm,linear,43].rejection_points['0.90'].rejection :: sha256 53e1ec81343ac2c1 :: at 90% signal efficiency, mean +- SD over seeds; 0 of 5 seeds at the cap
\newcommand{\ProbeAucEeVsMmLinearOneseven}{0.95832}  % experiments/FIGS/data/probe_ladder_v2_mlp2/analysis/seed_level_results.json :: levels.ee_vs_mm.linear[3].mean_auc :: sha256 53e1ec81343ac2c1 :: 5 seeds, 0 saturated at AUC=1
\newcommand{\ProbeAucSdEeVsMmLinearOneseven}{0.00970}  % experiments/FIGS/data/probe_ladder_v2_mlp2/analysis/seed_level_results.json :: table[ee_vs_mm,linear,17].auc (SD over seeds, ddof=1) :: sha256 53e1ec81343ac2c1
\newcommand{\ProbeLogOneMinusAucEeVsMmLinearOneseven}{\ensuremath{-}3.2013}  % experiments/FIGS/data/probe_ladder_v2_mlp2/analysis/seed_level_results.json :: levels.ee_vs_mm.linear[3].mean :: sha256 53e1ec81343ac2c1 :: natural log of 1-AUC, lower is better
\newcommand{\ProbeLogOneMinusAucSdEeVsMmLinearOneseven}{0.2488}  % experiments/FIGS/data/probe_ladder_v2_mlp2/analysis/seed_level_results.json :: levels.ee_vs_mm.linear[3].seed_sd :: sha256 53e1ec81343ac2c1
\newcommand{\ProbeOmaEeVsMmLinearOneseven}{\ensuremath{(4.2 \pm 1.0)\times10^{-2}}}  % experiments/FIGS/data/probe_ladder_v2_mlp2/analysis/seed_level_results.json :: table[ee_vs_mm,linear,17].auc :: sha256 53e1ec81343ac2c1 :: 1-AUC, mean +- SD over 5 seeds
\newcommand{\ProbeBkgLeftEeVsMmLinearOneseven}{289.6}  % experiments/FIGS/data/probe_ladder_v2_mlp2/analysis/seed_level_results.json :: levels.ee_vs_mm.linear[3].rejection_points['0.90'].mean_n_bkg_pass :: sha256 53e1ec81343ac2c1 :: background jets passing the cut, mean over seeds
\newcommand{\ProbeRejEeVsMmLinearOneseven}{17\,\ensuremath{\pm}\,6}  % experiments/FIGS/data/probe_ladder_v2_mlp2/analysis/seed_level_results.json :: table[ee_vs_mm,linear,17].rejection_points['0.90'].rejection :: sha256 53e1ec81343ac2c1 :: at 90% signal efficiency, mean +- SD over seeds; 0 of 5 seeds at the cap
\newcommand{\ProbeOmaRatioEeVsMmLinearOnesevenOverFourthree}{1.68}  % experiments/FIGS/data/probe_ladder_v2_mlp2/analysis/seed_level_results.json :: table[ee_vs_mm,linear,17/43].auc :: sha256 53e1ec81343ac2c1 :: ratio of the seed means of 1-AUC, coarser over finer
\newcommand{\ProbeAucEeVsMmMlpOneeighteight}{\ensuremath{1^{\dagger}}}  % experiments/FIGS/data/probe_ladder_v2_mlp2/analysis/seed_level_results.json :: levels.ee_vs_mm.mlp[0].mean_auc :: sha256 53e1ec81343ac2c1 :: 5 seeds, 5 saturated at AUC=1
\newcommand{\ProbeAucSdEeVsMmMlpOneeighteight}{0.00000}  % experiments/FIGS/data/probe_ladder_v2_mlp2/analysis/seed_level_results.json :: table[ee_vs_mm,mlp,188].auc (SD over seeds, ddof=1) :: sha256 53e1ec81343ac2c1
\newcommand{\ProbeLogOneMinusAucEeVsMmMlpOneeighteight}{\ensuremath{-}16.9579}  % experiments/FIGS/data/probe_ladder_v2_mlp2/analysis/seed_level_results.json :: levels.ee_vs_mm.mlp[0].mean :: sha256 53e1ec81343ac2c1 :: natural log of 1-AUC, lower is better
\newcommand{\ProbeLogOneMinusAucSdEeVsMmMlpOneeighteight}{0.0000}  % experiments/FIGS/data/probe_ladder_v2_mlp2/analysis/seed_level_results.json :: levels.ee_vs_mm.mlp[0].seed_sd :: sha256 53e1ec81343ac2c1
\newcommand{\ProbeOmaEeVsMmMlpOneeighteight}{---}  % experiments/FIGS/data/probe_ladder_v2_mlp2/analysis/seed_level_results.json :: table[ee_vs_mm,mlp,188].auc :: sha256 53e1ec81343ac2c1 :: not quoted: the AUC reached 1 in a seed, so 1-AUC is only a bound
\newcommand{\ProbeBkgLeftEeVsMmMlpOneeighteight}{0.0}  % experiments/FIGS/data/probe_ladder_v2_mlp2/analysis/seed_level_results.json :: levels.ee_vs_mm.mlp[0].rejection_points['0.90'].mean_n_bkg_pass :: sha256 53e1ec81343ac2c1 :: background jets passing the cut, mean over seeds
\newcommand{\ProbeRejEeVsMmMlpOneeighteight}{\ensuremath{>}1{,}571}  % experiments/FIGS/data/probe_ladder_v2_mlp2/analysis/seed_level_results.json :: table[ee_vs_mm,mlp,188].rejection_points['0.90'].rejection :: sha256 53e1ec81343ac2c1 :: at 90% signal efficiency, mean +- SD over seeds; 5 of 5 seeds at the cap
\newcommand{\ProbeAucEeVsMmMlpOnesixtwo}{\ensuremath{1^{\dagger}}}  % experiments/FIGS/data/probe_ladder_v2_mlp2/analysis/seed_level_results.json :: levels.ee_vs_mm.mlp[1].mean_auc :: sha256 53e1ec81343ac2c1 :: 5 seeds, 5 saturated at AUC=1
\newcommand{\ProbeAucSdEeVsMmMlpOnesixtwo}{0.00000}  % experiments/FIGS/data/probe_ladder_v2_mlp2/analysis/seed_level_results.json :: table[ee_vs_mm,mlp,162].auc (SD over seeds, ddof=1) :: sha256 53e1ec81343ac2c1
\newcommand{\ProbeLogOneMinusAucEeVsMmMlpOnesixtwo}{\ensuremath{-}16.9579}  % experiments/FIGS/data/probe_ladder_v2_mlp2/analysis/seed_level_results.json :: levels.ee_vs_mm.mlp[1].mean :: sha256 53e1ec81343ac2c1 :: natural log of 1-AUC, lower is better
\newcommand{\ProbeLogOneMinusAucSdEeVsMmMlpOnesixtwo}{0.0000}  % experiments/FIGS/data/probe_ladder_v2_mlp2/analysis/seed_level_results.json :: levels.ee_vs_mm.mlp[1].seed_sd :: sha256 53e1ec81343ac2c1
\newcommand{\ProbeOmaEeVsMmMlpOnesixtwo}{---}  % experiments/FIGS/data/probe_ladder_v2_mlp2/analysis/seed_level_results.json :: table[ee_vs_mm,mlp,162].auc :: sha256 53e1ec81343ac2c1 :: not quoted: the AUC reached 1 in a seed, so 1-AUC is only a bound
\newcommand{\ProbeBkgLeftEeVsMmMlpOnesixtwo}{0.0}  % experiments/FIGS/data/probe_ladder_v2_mlp2/analysis/seed_level_results.json :: levels.ee_vs_mm.mlp[1].rejection_points['0.90'].mean_n_bkg_pass :: sha256 53e1ec81343ac2c1 :: background jets passing the cut, mean over seeds
\newcommand{\ProbeRejEeVsMmMlpOnesixtwo}{\ensuremath{>}1{,}571}  % experiments/FIGS/data/probe_ladder_v2_mlp2/analysis/seed_level_results.json :: table[ee_vs_mm,mlp,162].rejection_points['0.90'].rejection :: sha256 53e1ec81343ac2c1 :: at 90% signal efficiency, mean +- SD over seeds; 5 of 5 seeds at the cap
\newcommand{\ProbeAucEeVsMmMlpFourthree}{0.99657}  % experiments/FIGS/data/probe_ladder_v2_mlp2/analysis/seed_level_results.json :: levels.ee_vs_mm.mlp[2].mean_auc :: sha256 53e1ec81343ac2c1 :: 5 seeds, 0 saturated at AUC=1
\newcommand{\ProbeAucSdEeVsMmMlpFourthree}{0.00300}  % experiments/FIGS/data/probe_ladder_v2_mlp2/analysis/seed_level_results.json :: table[ee_vs_mm,mlp,43].auc (SD over seeds, ddof=1) :: sha256 53e1ec81343ac2c1
\newcommand{\ProbeLogOneMinusAucEeVsMmMlpFourthree}{\ensuremath{-}6.0371}  % experiments/FIGS/data/probe_ladder_v2_mlp2/analysis/seed_level_results.json :: levels.ee_vs_mm.mlp[2].mean :: sha256 53e1ec81343ac2c1 :: natural log of 1-AUC, lower is better
\newcommand{\ProbeLogOneMinusAucSdEeVsMmMlpFourthree}{0.9962}  % experiments/FIGS/data/probe_ladder_v2_mlp2/analysis/seed_level_results.json :: levels.ee_vs_mm.mlp[2].seed_sd :: sha256 53e1ec81343ac2c1
\newcommand{\ProbeOmaEeVsMmMlpFourthree}{\ensuremath{(3.4 \pm 3.0)\times10^{-3}}}  % experiments/FIGS/data/probe_ladder_v2_mlp2/analysis/seed_level_results.json :: table[ee_vs_mm,mlp,43].auc :: sha256 53e1ec81343ac2c1 :: 1-AUC, mean +- SD over 5 seeds
\newcommand{\ProbeBkgLeftEeVsMmMlpFourthree}{33.4}  % experiments/FIGS/data/probe_ladder_v2_mlp2/analysis/seed_level_results.json :: levels.ee_vs_mm.mlp[2].rejection_points['0.90'].mean_n_bkg_pass :: sha256 53e1ec81343ac2c1 :: background jets passing the cut, mean over seeds
\newcommand{\ProbeRejEeVsMmMlpFourthree}{400\,\ensuremath{\pm}\,500}  % experiments/FIGS/data/probe_ladder_v2_mlp2/analysis/seed_level_results.json :: table[ee_vs_mm,mlp,43].rejection_points['0.90'].rejection :: sha256 53e1ec81343ac2c1 :: at 90% signal efficiency, mean +- SD over seeds; 0 of 5 seeds at the cap
\newcommand{\ProbeAucEeVsMmMlpOneseven}{0.99338}  % experiments/FIGS/data/probe_ladder_v2_mlp2/analysis/seed_level_results.json :: levels.ee_vs_mm.mlp[3].mean_auc :: sha256 53e1ec81343ac2c1 :: 5 seeds, 0 saturated at AUC=1
\newcommand{\ProbeAucSdEeVsMmMlpOneseven}{0.00065}  % experiments/FIGS/data/probe_ladder_v2_mlp2/analysis/seed_level_results.json :: table[ee_vs_mm,mlp,17].auc (SD over seeds, ddof=1) :: sha256 53e1ec81343ac2c1
\newcommand{\ProbeLogOneMinusAucEeVsMmMlpOneseven}{\ensuremath{-}5.0211}  % experiments/FIGS/data/probe_ladder_v2_mlp2/analysis/seed_level_results.json :: levels.ee_vs_mm.mlp[3].mean :: sha256 53e1ec81343ac2c1 :: natural log of 1-AUC, lower is better
\newcommand{\ProbeLogOneMinusAucSdEeVsMmMlpOneseven}{0.0997}  % experiments/FIGS/data/probe_ladder_v2_mlp2/analysis/seed_level_results.json :: levels.ee_vs_mm.mlp[3].seed_sd :: sha256 53e1ec81343ac2c1
\newcommand{\ProbeOmaEeVsMmMlpOneseven}{\ensuremath{(6.6 \pm 0.7)\times10^{-3}}}  % experiments/FIGS/data/probe_ladder_v2_mlp2/analysis/seed_level_results.json :: table[ee_vs_mm,mlp,17].auc :: sha256 53e1ec81343ac2c1 :: 1-AUC, mean +- SD over 5 seeds
\newcommand{\ProbeBkgLeftEeVsMmMlpOneseven}{59.2}  % experiments/FIGS/data/probe_ladder_v2_mlp2/analysis/seed_level_results.json :: levels.ee_vs_mm.mlp[3].rejection_points['0.90'].mean_n_bkg_pass :: sha256 53e1ec81343ac2c1 :: background jets passing the cut, mean over seeds
\newcommand{\ProbeRejEeVsMmMlpOneseven}{81\,\ensuremath{\pm}\,12}  % experiments/FIGS/data/probe_ladder_v2_mlp2/analysis/seed_level_results.json :: table[ee_vs_mm,mlp,17].rejection_points['0.90'].rejection :: sha256 53e1ec81343ac2c1 :: at 90% signal efficiency, mean +- SD over seeds; 0 of 5 seeds at the cap
\newcommand{\ProbeOmaRatioEeVsMmMlpOnesevenOverFourthree}{1.93}  % experiments/FIGS/data/probe_ladder_v2_mlp2/analysis/seed_level_results.json :: table[ee_vs_mm,mlp,17/43].auc :: sha256 53e1ec81343ac2c1 :: ratio of the seed means of 1-AUC, coarser over finer
\newcommand{\ProbeAucRetainedTopologyLinearOneeighteight}{0.99413}  % experiments/FIGS/data/probe_ladder_v2_mlp2/analysis/seed_level_results.json :: levels.retained_topology.linear[0].mean_auc :: sha256 53e1ec81343ac2c1 :: 5 seeds, 0 saturated at AUC=1
\newcommand{\ProbeAucSdRetainedTopologyLinearOneeighteight}{0.00011}  % experiments/FIGS/data/probe_ladder_v2_mlp2/analysis/seed_level_results.json :: table[retained_topology,linear,188].auc (SD over seeds, ddof=1) :: sha256 53e1ec81343ac2c1
\newcommand{\ProbeLogOneMinusAucRetainedTopologyLinearOneeighteight}{\ensuremath{-}5.1388}  % experiments/FIGS/data/probe_ladder_v2_mlp2/analysis/seed_level_results.json :: levels.retained_topology.linear[0].mean :: sha256 53e1ec81343ac2c1 :: natural log of 1-AUC, lower is better
\newcommand{\ProbeLogOneMinusAucSdRetainedTopologyLinearOneeighteight}{0.0187}  % experiments/FIGS/data/probe_ladder_v2_mlp2/analysis/seed_level_results.json :: levels.retained_topology.linear[0].seed_sd :: sha256 53e1ec81343ac2c1
\newcommand{\ProbeOmaRetainedTopologyLinearOneeighteight}{\ensuremath{(5.87 \pm 0.11)\times10^{-3}}}  % experiments/FIGS/data/probe_ladder_v2_mlp2/analysis/seed_level_results.json :: table[retained_topology,linear,188].auc :: sha256 53e1ec81343ac2c1 :: 1-AUC, mean +- SD over 5 seeds
\newcommand{\ProbeBkgLeftRetainedTopologyLinearOneeighteight}{15.4}  % experiments/FIGS/data/probe_ladder_v2_mlp2/analysis/seed_level_results.json :: levels.retained_topology.linear[0].rejection_points['0.90'].mean_n_bkg_pass :: sha256 53e1ec81343ac2c1 :: background jets passing the cut, mean over seeds
\newcommand{\ProbeRejRetainedTopologyLinearOneeighteight}{71\,\ensuremath{\pm}\,7}  % experiments/FIGS/data/probe_ladder_v2_mlp2/analysis/seed_level_results.json :: table[retained_topology,linear,188].rejection_points['0.90'].rejection :: sha256 53e1ec81343ac2c1 :: at 90% signal efficiency, mean +- SD over seeds; 0 of 5 seeds at the cap
\newcommand{\ProbeAucRetainedTopologyLinearOnesixtwo}{0.99413}  % experiments/FIGS/data/probe_ladder_v2_mlp2/analysis/seed_level_results.json :: levels.retained_topology.linear[1].mean_auc :: sha256 53e1ec81343ac2c1 :: 5 seeds, 0 saturated at AUC=1
\newcommand{\ProbeAucSdRetainedTopologyLinearOnesixtwo}{0.00030}  % experiments/FIGS/data/probe_ladder_v2_mlp2/analysis/seed_level_results.json :: table[retained_topology,linear,162].auc (SD over seeds, ddof=1) :: sha256 53e1ec81343ac2c1
\newcommand{\ProbeLogOneMinusAucRetainedTopologyLinearOnesixtwo}{\ensuremath{-}5.1395}  % experiments/FIGS/data/probe_ladder_v2_mlp2/analysis/seed_level_results.json :: levels.retained_topology.linear[1].mean :: sha256 53e1ec81343ac2c1 :: natural log of 1-AUC, lower is better
\newcommand{\ProbeLogOneMinusAucSdRetainedTopologyLinearOnesixtwo}{0.0510}  % experiments/FIGS/data/probe_ladder_v2_mlp2/analysis/seed_level_results.json :: levels.retained_topology.linear[1].seed_sd :: sha256 53e1ec81343ac2c1
\newcommand{\ProbeOmaRetainedTopologyLinearOnesixtwo}{\ensuremath{(5.87 \pm 0.30)\times10^{-3}}}  % experiments/FIGS/data/probe_ladder_v2_mlp2/analysis/seed_level_results.json :: table[retained_topology,linear,162].auc :: sha256 53e1ec81343ac2c1 :: 1-AUC, mean +- SD over 5 seeds
\newcommand{\ProbeBkgLeftRetainedTopologyLinearOnesixtwo}{14.0}  % experiments/FIGS/data/probe_ladder_v2_mlp2/analysis/seed_level_results.json :: levels.retained_topology.linear[1].rejection_points['0.90'].mean_n_bkg_pass :: sha256 53e1ec81343ac2c1 :: background jets passing the cut, mean over seeds
\newcommand{\ProbeRejRetainedTopologyLinearOnesixtwo}{79\,\ensuremath{\pm}\,13}  % experiments/FIGS/data/probe_ladder_v2_mlp2/analysis/seed_level_results.json :: table[retained_topology,linear,162].rejection_points['0.90'].rejection :: sha256 53e1ec81343ac2c1 :: at 90% signal efficiency, mean +- SD over seeds; 0 of 5 seeds at the cap
\newcommand{\ProbeAucRetainedTopologyLinearFourthree}{0.99274}  % experiments/FIGS/data/probe_ladder_v2_mlp2/analysis/seed_level_results.json :: levels.retained_topology.linear[2].mean_auc :: sha256 53e1ec81343ac2c1 :: 5 seeds, 0 saturated at AUC=1
\newcommand{\ProbeAucSdRetainedTopologyLinearFourthree}{0.00012}  % experiments/FIGS/data/probe_ladder_v2_mlp2/analysis/seed_level_results.json :: table[retained_topology,linear,43].auc (SD over seeds, ddof=1) :: sha256 53e1ec81343ac2c1
\newcommand{\ProbeLogOneMinusAucRetainedTopologyLinearFourthree}{\ensuremath{-}4.9253}  % experiments/FIGS/data/probe_ladder_v2_mlp2/analysis/seed_level_results.json :: levels.retained_topology.linear[2].mean :: sha256 53e1ec81343ac2c1 :: natural log of 1-AUC, lower is better
\newcommand{\ProbeLogOneMinusAucSdRetainedTopologyLinearFourthree}{0.0169}  % experiments/FIGS/data/probe_ladder_v2_mlp2/analysis/seed_level_results.json :: levels.retained_topology.linear[2].seed_sd :: sha256 53e1ec81343ac2c1
\newcommand{\ProbeOmaRetainedTopologyLinearFourthree}{\ensuremath{(7.26 \pm 0.12)\times10^{-3}}}  % experiments/FIGS/data/probe_ladder_v2_mlp2/analysis/seed_level_results.json :: table[retained_topology,linear,43].auc :: sha256 53e1ec81343ac2c1 :: 1-AUC, mean +- SD over 5 seeds
\newcommand{\ProbeBkgLeftRetainedTopologyLinearFourthree}{16.4}  % experiments/FIGS/data/probe_ladder_v2_mlp2/analysis/seed_level_results.json :: levels.retained_topology.linear[2].rejection_points['0.90'].mean_n_bkg_pass :: sha256 53e1ec81343ac2c1 :: background jets passing the cut, mean over seeds
\newcommand{\ProbeRejRetainedTopologyLinearFourthree}{66\,\ensuremath{\pm}\,5}  % experiments/FIGS/data/probe_ladder_v2_mlp2/analysis/seed_level_results.json :: table[retained_topology,linear,43].rejection_points['0.90'].rejection :: sha256 53e1ec81343ac2c1 :: at 90% signal efficiency, mean +- SD over seeds; 0 of 5 seeds at the cap
\newcommand{\ProbeAucRetainedTopologyLinearOneseven}{0.99096}  % experiments/FIGS/data/probe_ladder_v2_mlp2/analysis/seed_level_results.json :: levels.retained_topology.linear[3].mean_auc :: sha256 53e1ec81343ac2c1 :: 5 seeds, 0 saturated at AUC=1
\newcommand{\ProbeAucSdRetainedTopologyLinearOneseven}{0.00092}  % experiments/FIGS/data/probe_ladder_v2_mlp2/analysis/seed_level_results.json :: table[retained_topology,linear,17].auc (SD over seeds, ddof=1) :: sha256 53e1ec81343ac2c1
\newcommand{\ProbeLogOneMinusAucRetainedTopologyLinearOneseven}{\ensuremath{-}4.7095}  % experiments/FIGS/data/probe_ladder_v2_mlp2/analysis/seed_level_results.json :: levels.retained_topology.linear[3].mean :: sha256 53e1ec81343ac2c1 :: natural log of 1-AUC, lower is better
\newcommand{\ProbeLogOneMinusAucSdRetainedTopologyLinearOneseven}{0.0979}  % experiments/FIGS/data/probe_ladder_v2_mlp2/analysis/seed_level_results.json :: levels.retained_topology.linear[3].seed_sd :: sha256 53e1ec81343ac2c1
\newcommand{\ProbeOmaRetainedTopologyLinearOneseven}{\ensuremath{(9.0 \pm 0.9)\times10^{-3}}}  % experiments/FIGS/data/probe_ladder_v2_mlp2/analysis/seed_level_results.json :: table[retained_topology,linear,17].auc :: sha256 53e1ec81343ac2c1 :: 1-AUC, mean +- SD over 5 seeds
\newcommand{\ProbeBkgLeftRetainedTopologyLinearOneseven}{21.8}  % experiments/FIGS/data/probe_ladder_v2_mlp2/analysis/seed_level_results.json :: levels.retained_topology.linear[3].rejection_points['0.90'].mean_n_bkg_pass :: sha256 53e1ec81343ac2c1 :: background jets passing the cut, mean over seeds
\newcommand{\ProbeRejRetainedTopologyLinearOneseven}{51\,\ensuremath{\pm}\,11}  % experiments/FIGS/data/probe_ladder_v2_mlp2/analysis/seed_level_results.json :: table[retained_topology,linear,17].rejection_points['0.90'].rejection :: sha256 53e1ec81343ac2c1 :: at 90% signal efficiency, mean +- SD over seeds; 0 of 5 seeds at the cap
\newcommand{\ProbeOmaRatioRetainedTopologyLinearOnesixtwoOverOneeighteight}{1.00}  % experiments/FIGS/data/probe_ladder_v2_mlp2/analysis/seed_level_results.json :: table[retained_topology,linear,162/188].auc :: sha256 53e1ec81343ac2c1 :: ratio of the seed means of 1-AUC, coarser over finer
\newcommand{\ProbeOmaRatioRetainedTopologyLinearFourthreeOverOneeighteight}{1.24}  % experiments/FIGS/data/probe_ladder_v2_mlp2/analysis/seed_level_results.json :: table[retained_topology,linear,43/188].auc :: sha256 53e1ec81343ac2c1 :: ratio of the seed means of 1-AUC, coarser over finer
\newcommand{\ProbeOmaRatioRetainedTopologyLinearOnesevenOverOneeighteight}{1.54}  % experiments/FIGS/data/probe_ladder_v2_mlp2/analysis/seed_level_results.json :: table[retained_topology,linear,17/188].auc :: sha256 53e1ec81343ac2c1 :: ratio of the seed means of 1-AUC, coarser over finer
\newcommand{\ProbeOmaRatioRetainedTopologyLinearFourthreeOverOnesixtwo}{1.24}  % experiments/FIGS/data/probe_ladder_v2_mlp2/analysis/seed_level_results.json :: table[retained_topology,linear,43/162].auc :: sha256 53e1ec81343ac2c1 :: ratio of the seed means of 1-AUC, coarser over finer
\newcommand{\ProbeOmaRatioRetainedTopologyLinearOnesevenOverOnesixtwo}{1.54}  % experiments/FIGS/data/probe_ladder_v2_mlp2/analysis/seed_level_results.json :: table[retained_topology,linear,17/162].auc :: sha256 53e1ec81343ac2c1 :: ratio of the seed means of 1-AUC, coarser over finer
\newcommand{\ProbeOmaRatioRetainedTopologyLinearOnesevenOverFourthree}{1.25}  % experiments/FIGS/data/probe_ladder_v2_mlp2/analysis/seed_level_results.json :: table[retained_topology,linear,17/43].auc :: sha256 53e1ec81343ac2c1 :: ratio of the seed means of 1-AUC, coarser over finer
\newcommand{\ProbeAucRetainedTopologyMlpOneeighteight}{0.99394}  % experiments/FIGS/data/probe_ladder_v2_mlp2/analysis/seed_level_results.json :: levels.retained_topology.mlp[0].mean_auc :: sha256 53e1ec81343ac2c1 :: 5 seeds, 0 saturated at AUC=1
\newcommand{\ProbeAucSdRetainedTopologyMlpOneeighteight}{0.00032}  % experiments/FIGS/data/probe_ladder_v2_mlp2/analysis/seed_level_results.json :: table[retained_topology,mlp,188].auc (SD over seeds, ddof=1) :: sha256 53e1ec81343ac2c1
\newcommand{\ProbeLogOneMinusAucRetainedTopologyMlpOneeighteight}{\ensuremath{-}5.1078}  % experiments/FIGS/data/probe_ladder_v2_mlp2/analysis/seed_level_results.json :: levels.retained_topology.mlp[0].mean :: sha256 53e1ec81343ac2c1 :: natural log of 1-AUC, lower is better
\newcommand{\ProbeLogOneMinusAucSdRetainedTopologyMlpOneeighteight}{0.0522}  % experiments/FIGS/data/probe_ladder_v2_mlp2/analysis/seed_level_results.json :: levels.retained_topology.mlp[0].seed_sd :: sha256 53e1ec81343ac2c1
\newcommand{\ProbeOmaRetainedTopologyMlpOneeighteight}{\ensuremath{(6.06 \pm 0.32)\times10^{-3}}}  % experiments/FIGS/data/probe_ladder_v2_mlp2/analysis/seed_level_results.json :: table[retained_topology,mlp,188].auc :: sha256 53e1ec81343ac2c1 :: 1-AUC, mean +- SD over 5 seeds
\newcommand{\ProbeBkgLeftRetainedTopologyMlpOneeighteight}{15.8}  % experiments/FIGS/data/probe_ladder_v2_mlp2/analysis/seed_level_results.json :: levels.retained_topology.mlp[0].rejection_points['0.90'].mean_n_bkg_pass :: sha256 53e1ec81343ac2c1 :: background jets passing the cut, mean over seeds
\newcommand{\ProbeRejRetainedTopologyMlpOneeighteight}{69\,\ensuremath{\pm}\,9}  % experiments/FIGS/data/probe_ladder_v2_mlp2/analysis/seed_level_results.json :: table[retained_topology,mlp,188].rejection_points['0.90'].rejection :: sha256 53e1ec81343ac2c1 :: at 90% signal efficiency, mean +- SD over seeds; 0 of 5 seeds at the cap
\newcommand{\ProbeAucRetainedTopologyMlpOnesixtwo}{0.99387}  % experiments/FIGS/data/probe_ladder_v2_mlp2/analysis/seed_level_results.json :: levels.retained_topology.mlp[1].mean_auc :: sha256 53e1ec81343ac2c1 :: 5 seeds, 0 saturated at AUC=1
\newcommand{\ProbeAucSdRetainedTopologyMlpOnesixtwo}{0.00048}  % experiments/FIGS/data/probe_ladder_v2_mlp2/analysis/seed_level_results.json :: table[retained_topology,mlp,162].auc (SD over seeds, ddof=1) :: sha256 53e1ec81343ac2c1
\newcommand{\ProbeLogOneMinusAucRetainedTopologyMlpOnesixtwo}{\ensuremath{-}5.0964}  % experiments/FIGS/data/probe_ladder_v2_mlp2/analysis/seed_level_results.json :: levels.retained_topology.mlp[1].mean :: sha256 53e1ec81343ac2c1 :: natural log of 1-AUC, lower is better
\newcommand{\ProbeLogOneMinusAucSdRetainedTopologyMlpOnesixtwo}{0.0773}  % experiments/FIGS/data/probe_ladder_v2_mlp2/analysis/seed_level_results.json :: levels.retained_topology.mlp[1].seed_sd :: sha256 53e1ec81343ac2c1
\newcommand{\ProbeOmaRetainedTopologyMlpOnesixtwo}{\ensuremath{(6.1 \pm 0.5)\times10^{-3}}}  % experiments/FIGS/data/probe_ladder_v2_mlp2/analysis/seed_level_results.json :: table[retained_topology,mlp,162].auc :: sha256 53e1ec81343ac2c1 :: 1-AUC, mean +- SD over 5 seeds
\newcommand{\ProbeBkgLeftRetainedTopologyMlpOnesixtwo}{16.6}  % experiments/FIGS/data/probe_ladder_v2_mlp2/analysis/seed_level_results.json :: levels.retained_topology.mlp[1].rejection_points['0.90'].mean_n_bkg_pass :: sha256 53e1ec81343ac2c1 :: background jets passing the cut, mean over seeds
\newcommand{\ProbeRejRetainedTopologyMlpOnesixtwo}{67\,\ensuremath{\pm}\,14}  % experiments/FIGS/data/probe_ladder_v2_mlp2/analysis/seed_level_results.json :: table[retained_topology,mlp,162].rejection_points['0.90'].rejection :: sha256 53e1ec81343ac2c1 :: at 90% signal efficiency, mean +- SD over seeds; 0 of 5 seeds at the cap
\newcommand{\ProbeAucRetainedTopologyMlpFourthree}{0.99284}  % experiments/FIGS/data/probe_ladder_v2_mlp2/analysis/seed_level_results.json :: levels.retained_topology.mlp[2].mean_auc :: sha256 53e1ec81343ac2c1 :: 5 seeds, 0 saturated at AUC=1
\newcommand{\ProbeAucSdRetainedTopologyMlpFourthree}{0.00040}  % experiments/FIGS/data/probe_ladder_v2_mlp2/analysis/seed_level_results.json :: table[retained_topology,mlp,43].auc (SD over seeds, ddof=1) :: sha256 53e1ec81343ac2c1
\newcommand{\ProbeLogOneMinusAucRetainedTopologyMlpFourthree}{\ensuremath{-}4.9400}  % experiments/FIGS/data/probe_ladder_v2_mlp2/analysis/seed_level_results.json :: levels.retained_topology.mlp[2].mean :: sha256 53e1ec81343ac2c1 :: natural log of 1-AUC, lower is better
\newcommand{\ProbeLogOneMinusAucSdRetainedTopologyMlpFourthree}{0.0570}  % experiments/FIGS/data/probe_ladder_v2_mlp2/analysis/seed_level_results.json :: levels.retained_topology.mlp[2].seed_sd :: sha256 53e1ec81343ac2c1
\newcommand{\ProbeOmaRetainedTopologyMlpFourthree}{\ensuremath{(7.2 \pm 0.4)\times10^{-3}}}  % experiments/FIGS/data/probe_ladder_v2_mlp2/analysis/seed_level_results.json :: table[retained_topology,mlp,43].auc :: sha256 53e1ec81343ac2c1 :: 1-AUC, mean +- SD over 5 seeds
\newcommand{\ProbeBkgLeftRetainedTopologyMlpFourthree}{17.8}  % experiments/FIGS/data/probe_ladder_v2_mlp2/analysis/seed_level_results.json :: levels.retained_topology.mlp[2].rejection_points['0.90'].mean_n_bkg_pass :: sha256 53e1ec81343ac2c1 :: background jets passing the cut, mean over seeds
\newcommand{\ProbeRejRetainedTopologyMlpFourthree}{61\,\ensuremath{\pm}\,7}  % experiments/FIGS/data/probe_ladder_v2_mlp2/analysis/seed_level_results.json :: table[retained_topology,mlp,43].rejection_points['0.90'].rejection :: sha256 53e1ec81343ac2c1 :: at 90% signal efficiency, mean +- SD over seeds; 0 of 5 seeds at the cap
\newcommand{\ProbeAucRetainedTopologyMlpOneseven}{0.99081}  % experiments/FIGS/data/probe_ladder_v2_mlp2/analysis/seed_level_results.json :: levels.retained_topology.mlp[3].mean_auc :: sha256 53e1ec81343ac2c1 :: 5 seeds, 0 saturated at AUC=1
\newcommand{\ProbeAucSdRetainedTopologyMlpOneseven}{0.00062}  % experiments/FIGS/data/probe_ladder_v2_mlp2/analysis/seed_level_results.json :: table[retained_topology,mlp,17].auc (SD over seeds, ddof=1) :: sha256 53e1ec81343ac2c1
\newcommand{\ProbeLogOneMinusAucRetainedTopologyMlpOneseven}{\ensuremath{-}4.6913}  % experiments/FIGS/data/probe_ladder_v2_mlp2/analysis/seed_level_results.json :: levels.retained_topology.mlp[3].mean :: sha256 53e1ec81343ac2c1 :: natural log of 1-AUC, lower is better
\newcommand{\ProbeLogOneMinusAucSdRetainedTopologyMlpOneseven}{0.0658}  % experiments/FIGS/data/probe_ladder_v2_mlp2/analysis/seed_level_results.json :: levels.retained_topology.mlp[3].seed_sd :: sha256 53e1ec81343ac2c1
\newcommand{\ProbeOmaRetainedTopologyMlpOneseven}{\ensuremath{(9.2 \pm 0.6)\times10^{-3}}}  % experiments/FIGS/data/probe_ladder_v2_mlp2/analysis/seed_level_results.json :: table[retained_topology,mlp,17].auc :: sha256 53e1ec81343ac2c1 :: 1-AUC, mean +- SD over 5 seeds
\newcommand{\ProbeBkgLeftRetainedTopologyMlpOneseven}{21.2}  % experiments/FIGS/data/probe_ladder_v2_mlp2/analysis/seed_level_results.json :: levels.retained_topology.mlp[3].rejection_points['0.90'].mean_n_bkg_pass :: sha256 53e1ec81343ac2c1 :: background jets passing the cut, mean over seeds
\newcommand{\ProbeRejRetainedTopologyMlpOneseven}{51\,\ensuremath{\pm}\,4}  % experiments/FIGS/data/probe_ladder_v2_mlp2/analysis/seed_level_results.json :: table[retained_topology,mlp,17].rejection_points['0.90'].rejection :: sha256 53e1ec81343ac2c1 :: at 90% signal efficiency, mean +- SD over seeds; 0 of 5 seeds at the cap
\newcommand{\ProbeOmaRatioRetainedTopologyMlpOnesixtwoOverOneeighteight}{1.01}  % experiments/FIGS/data/probe_ladder_v2_mlp2/analysis/seed_level_results.json :: table[retained_topology,mlp,162/188].auc :: sha256 53e1ec81343ac2c1 :: ratio of the seed means of 1-AUC, coarser over finer
\newcommand{\ProbeOmaRatioRetainedTopologyMlpFourthreeOverOneeighteight}{1.18}  % experiments/FIGS/data/probe_ladder_v2_mlp2/analysis/seed_level_results.json :: table[retained_topology,mlp,43/188].auc :: sha256 53e1ec81343ac2c1 :: ratio of the seed means of 1-AUC, coarser over finer
\newcommand{\ProbeOmaRatioRetainedTopologyMlpOnesevenOverOneeighteight}{1.52}  % experiments/FIGS/data/probe_ladder_v2_mlp2/analysis/seed_level_results.json :: table[retained_topology,mlp,17/188].auc :: sha256 53e1ec81343ac2c1 :: ratio of the seed means of 1-AUC, coarser over finer
\newcommand{\ProbeOmaRatioRetainedTopologyMlpFourthreeOverOnesixtwo}{1.17}  % experiments/FIGS/data/probe_ladder_v2_mlp2/analysis/seed_level_results.json :: table[retained_topology,mlp,43/162].auc :: sha256 53e1ec81343ac2c1 :: ratio of the seed means of 1-AUC, coarser over finer
\newcommand{\ProbeOmaRatioRetainedTopologyMlpOnesevenOverOnesixtwo}{1.50}  % experiments/FIGS/data/probe_ladder_v2_mlp2/analysis/seed_level_results.json :: table[retained_topology,mlp,17/162].auc :: sha256 53e1ec81343ac2c1 :: ratio of the seed means of 1-AUC, coarser over finer
\newcommand{\ProbeOmaRatioRetainedTopologyMlpOnesevenOverFourthree}{1.28}  % experiments/FIGS/data/probe_ladder_v2_mlp2/analysis/seed_level_results.json :: table[retained_topology,mlp,17/43].auc :: sha256 53e1ec81343ac2c1 :: ratio of the seed means of 1-AUC, coarser over finer
\newcommand{\ProbeAucVisibleContentLinearOneeighteight}{0.99917}  % experiments/FIGS/data/probe_ladder_v2_mlp2/analysis/seed_level_results.json :: levels.visible_content.linear[0].mean_auc :: sha256 53e1ec81343ac2c1 :: 5 seeds, 0 saturated at AUC=1
\newcommand{\ProbeAucSdVisibleContentLinearOneeighteight}{0.00003}  % experiments/FIGS/data/probe_ladder_v2_mlp2/analysis/seed_level_results.json :: table[visible_content,linear,188].auc (SD over seeds, ddof=1) :: sha256 53e1ec81343ac2c1
\newcommand{\ProbeLogOneMinusAucVisibleContentLinearOneeighteight}{\ensuremath{-}7.0999}  % experiments/FIGS/data/probe_ladder_v2_mlp2/analysis/seed_level_results.json :: levels.visible_content.linear[0].mean :: sha256 53e1ec81343ac2c1 :: natural log of 1-AUC, lower is better
\newcommand{\ProbeLogOneMinusAucSdVisibleContentLinearOneeighteight}{0.0358}  % experiments/FIGS/data/probe_ladder_v2_mlp2/analysis/seed_level_results.json :: levels.visible_content.linear[0].seed_sd :: sha256 53e1ec81343ac2c1
\newcommand{\ProbeOmaVisibleContentLinearOneeighteight}{\ensuremath{(8.26 \pm 0.30)\times10^{-4}}}  % experiments/FIGS/data/probe_ladder_v2_mlp2/analysis/seed_level_results.json :: table[visible_content,linear,188].auc :: sha256 53e1ec81343ac2c1 :: 1-AUC, mean +- SD over 5 seeds
\newcommand{\ProbeBkgLeftVisibleContentLinearOneeighteight}{0.4}  % experiments/FIGS/data/probe_ladder_v2_mlp2/analysis/seed_level_results.json :: levels.visible_content.linear[0].rejection_points['0.90'].mean_n_bkg_pass :: sha256 53e1ec81343ac2c1 :: background jets passing the cut, mean over seeds
\newcommand{\ProbeRejVisibleContentLinearOneeighteight}{\ensuremath{>}538}  % experiments/FIGS/data/probe_ladder_v2_mlp2/analysis/seed_level_results.json :: table[visible_content,linear,188].rejection_points['0.90'].rejection :: sha256 53e1ec81343ac2c1 :: at 90% signal efficiency, mean +- SD over seeds; 5 of 5 seeds at the cap
\newcommand{\ProbeAucVisibleContentLinearOnesixtwo}{0.99921}  % experiments/FIGS/data/probe_ladder_v2_mlp2/analysis/seed_level_results.json :: levels.visible_content.linear[1].mean_auc :: sha256 53e1ec81343ac2c1 :: 5 seeds, 0 saturated at AUC=1
\newcommand{\ProbeAucSdVisibleContentLinearOnesixtwo}{0.00004}  % experiments/FIGS/data/probe_ladder_v2_mlp2/analysis/seed_level_results.json :: table[visible_content,linear,162].auc (SD over seeds, ddof=1) :: sha256 53e1ec81343ac2c1
\newcommand{\ProbeLogOneMinusAucVisibleContentLinearOnesixtwo}{\ensuremath{-}7.1480}  % experiments/FIGS/data/probe_ladder_v2_mlp2/analysis/seed_level_results.json :: levels.visible_content.linear[1].mean :: sha256 53e1ec81343ac2c1 :: natural log of 1-AUC, lower is better
\newcommand{\ProbeLogOneMinusAucSdVisibleContentLinearOnesixtwo}{0.0449}  % experiments/FIGS/data/probe_ladder_v2_mlp2/analysis/seed_level_results.json :: levels.visible_content.linear[1].seed_sd :: sha256 53e1ec81343ac2c1
\newcommand{\ProbeOmaVisibleContentLinearOnesixtwo}{\ensuremath{(7.87 \pm 0.35)\times10^{-4}}}  % experiments/FIGS/data/probe_ladder_v2_mlp2/analysis/seed_level_results.json :: table[visible_content,linear,162].auc :: sha256 53e1ec81343ac2c1 :: 1-AUC, mean +- SD over 5 seeds
\newcommand{\ProbeBkgLeftVisibleContentLinearOnesixtwo}{0.6}  % experiments/FIGS/data/probe_ladder_v2_mlp2/analysis/seed_level_results.json :: levels.visible_content.linear[1].rejection_points['0.90'].mean_n_bkg_pass :: sha256 53e1ec81343ac2c1 :: background jets passing the cut, mean over seeds
\newcommand{\ProbeRejVisibleContentLinearOnesixtwo}{\ensuremath{>}538}  % experiments/FIGS/data/probe_ladder_v2_mlp2/analysis/seed_level_results.json :: table[visible_content,linear,162].rejection_points['0.90'].rejection :: sha256 53e1ec81343ac2c1 :: at 90% signal efficiency, mean +- SD over seeds; 5 of 5 seeds at the cap
\newcommand{\ProbeAucVisibleContentLinearFourthree}{0.99918}  % experiments/FIGS/data/probe_ladder_v2_mlp2/analysis/seed_level_results.json :: levels.visible_content.linear[2].mean_auc :: sha256 53e1ec81343ac2c1 :: 5 seeds, 0 saturated at AUC=1
\newcommand{\ProbeAucSdVisibleContentLinearFourthree}{0.00009}  % experiments/FIGS/data/probe_ladder_v2_mlp2/analysis/seed_level_results.json :: table[visible_content,linear,43].auc (SD over seeds, ddof=1) :: sha256 53e1ec81343ac2c1
\newcommand{\ProbeLogOneMinusAucVisibleContentLinearFourthree}{\ensuremath{-}7.1124}  % experiments/FIGS/data/probe_ladder_v2_mlp2/analysis/seed_level_results.json :: levels.visible_content.linear[2].mean :: sha256 53e1ec81343ac2c1 :: natural log of 1-AUC, lower is better
\newcommand{\ProbeLogOneMinusAucSdVisibleContentLinearFourthree}{0.1137}  % experiments/FIGS/data/probe_ladder_v2_mlp2/analysis/seed_level_results.json :: levels.visible_content.linear[2].seed_sd :: sha256 53e1ec81343ac2c1
\newcommand{\ProbeOmaVisibleContentLinearFourthree}{\ensuremath{(8.2 \pm 1.0)\times10^{-4}}}  % experiments/FIGS/data/probe_ladder_v2_mlp2/analysis/seed_level_results.json :: table[visible_content,linear,43].auc :: sha256 53e1ec81343ac2c1 :: 1-AUC, mean +- SD over 5 seeds
\newcommand{\ProbeBkgLeftVisibleContentLinearFourthree}{0.4}  % experiments/FIGS/data/probe_ladder_v2_mlp2/analysis/seed_level_results.json :: levels.visible_content.linear[2].rejection_points['0.90'].mean_n_bkg_pass :: sha256 53e1ec81343ac2c1 :: background jets passing the cut, mean over seeds
\newcommand{\ProbeRejVisibleContentLinearFourthree}{\ensuremath{>}538}  % experiments/FIGS/data/probe_ladder_v2_mlp2/analysis/seed_level_results.json :: table[visible_content,linear,43].rejection_points['0.90'].rejection :: sha256 53e1ec81343ac2c1 :: at 90% signal efficiency, mean +- SD over seeds; 5 of 5 seeds at the cap
\newcommand{\ProbeAucVisibleContentLinearOneseven}{0.99768}  % experiments/FIGS/data/probe_ladder_v2_mlp2/analysis/seed_level_results.json :: levels.visible_content.linear[3].mean_auc :: sha256 53e1ec81343ac2c1 :: 5 seeds, 0 saturated at AUC=1
\newcommand{\ProbeAucSdVisibleContentLinearOneseven}{0.00020}  % experiments/FIGS/data/probe_ladder_v2_mlp2/analysis/seed_level_results.json :: table[visible_content,linear,17].auc (SD over seeds, ddof=1) :: sha256 53e1ec81343ac2c1
\newcommand{\ProbeLogOneMinusAucVisibleContentLinearOneseven}{\ensuremath{-}6.0684}  % experiments/FIGS/data/probe_ladder_v2_mlp2/analysis/seed_level_results.json :: levels.visible_content.linear[3].mean :: sha256 53e1ec81343ac2c1 :: natural log of 1-AUC, lower is better
\newcommand{\ProbeLogOneMinusAucSdVisibleContentLinearOneseven}{0.0889}  % experiments/FIGS/data/probe_ladder_v2_mlp2/analysis/seed_level_results.json :: levels.visible_content.linear[3].seed_sd :: sha256 53e1ec81343ac2c1
\newcommand{\ProbeOmaVisibleContentLinearOneseven}{\ensuremath{(2.32 \pm 0.20)\times10^{-3}}}  % experiments/FIGS/data/probe_ladder_v2_mlp2/analysis/seed_level_results.json :: table[visible_content,linear,17].auc :: sha256 53e1ec81343ac2c1 :: 1-AUC, mean +- SD over 5 seeds
\newcommand{\ProbeBkgLeftVisibleContentLinearOneseven}{3.4}  % experiments/FIGS/data/probe_ladder_v2_mlp2/analysis/seed_level_results.json :: levels.visible_content.linear[3].rejection_points['0.90'].mean_n_bkg_pass :: sha256 53e1ec81343ac2c1 :: background jets passing the cut, mean over seeds
\newcommand{\ProbeRejVisibleContentLinearOneseven}{900\,\ensuremath{\pm}\,400}  % experiments/FIGS/data/probe_ladder_v2_mlp2/analysis/seed_level_results.json :: table[visible_content,linear,17].rejection_points['0.90'].rejection :: sha256 53e1ec81343ac2c1 :: at 90% signal efficiency, mean +- SD over seeds; 0 of 5 seeds at the cap
\newcommand{\ProbeOmaRatioVisibleContentLinearOnesixtwoOverOneeighteight}{0.95}  % experiments/FIGS/data/probe_ladder_v2_mlp2/analysis/seed_level_results.json :: table[visible_content,linear,162/188].auc :: sha256 53e1ec81343ac2c1 :: ratio of the seed means of 1-AUC, coarser over finer
\newcommand{\ProbeOmaRatioVisibleContentLinearFourthreeOverOneeighteight}{0.99}  % experiments/FIGS/data/probe_ladder_v2_mlp2/analysis/seed_level_results.json :: table[visible_content,linear,43/188].auc :: sha256 53e1ec81343ac2c1 :: ratio of the seed means of 1-AUC, coarser over finer
\newcommand{\ProbeOmaRatioVisibleContentLinearOnesevenOverOneeighteight}{2.8}  % experiments/FIGS/data/probe_ladder_v2_mlp2/analysis/seed_level_results.json :: table[visible_content,linear,17/188].auc :: sha256 53e1ec81343ac2c1 :: ratio of the seed means of 1-AUC, coarser over finer
\newcommand{\ProbeOmaRatioVisibleContentLinearFourthreeOverOnesixtwo}{1.04}  % experiments/FIGS/data/probe_ladder_v2_mlp2/analysis/seed_level_results.json :: table[visible_content,linear,43/162].auc :: sha256 53e1ec81343ac2c1 :: ratio of the seed means of 1-AUC, coarser over finer
\newcommand{\ProbeOmaRatioVisibleContentLinearOnesevenOverOnesixtwo}{3.0}  % experiments/FIGS/data/probe_ladder_v2_mlp2/analysis/seed_level_results.json :: table[visible_content,linear,17/162].auc :: sha256 53e1ec81343ac2c1 :: ratio of the seed means of 1-AUC, coarser over finer
\newcommand{\ProbeOmaRatioVisibleContentLinearOnesevenOverFourthree}{2.8}  % experiments/FIGS/data/probe_ladder_v2_mlp2/analysis/seed_level_results.json :: table[visible_content,linear,17/43].auc :: sha256 53e1ec81343ac2c1 :: ratio of the seed means of 1-AUC, coarser over finer
\newcommand{\ProbeAucVisibleContentMlpOneeighteight}{0.99911}  % experiments/FIGS/data/probe_ladder_v2_mlp2/analysis/seed_level_results.json :: levels.visible_content.mlp[0].mean_auc :: sha256 53e1ec81343ac2c1 :: 5 seeds, 0 saturated at AUC=1
\newcommand{\ProbeAucSdVisibleContentMlpOneeighteight}{0.00005}  % experiments/FIGS/data/probe_ladder_v2_mlp2/analysis/seed_level_results.json :: table[visible_content,mlp,188].auc (SD over seeds, ddof=1) :: sha256 53e1ec81343ac2c1
\newcommand{\ProbeLogOneMinusAucVisibleContentMlpOneeighteight}{\ensuremath{-}7.0202}  % experiments/FIGS/data/probe_ladder_v2_mlp2/analysis/seed_level_results.json :: levels.visible_content.mlp[0].mean :: sha256 53e1ec81343ac2c1 :: natural log of 1-AUC, lower is better
\newcommand{\ProbeLogOneMinusAucSdVisibleContentMlpOneeighteight}{0.0570}  % experiments/FIGS/data/probe_ladder_v2_mlp2/analysis/seed_level_results.json :: levels.visible_content.mlp[0].seed_sd :: sha256 53e1ec81343ac2c1
\newcommand{\ProbeOmaVisibleContentMlpOneeighteight}{\ensuremath{(8.9 \pm 0.5)\times10^{-4}}}  % experiments/FIGS/data/probe_ladder_v2_mlp2/analysis/seed_level_results.json :: table[visible_content,mlp,188].auc :: sha256 53e1ec81343ac2c1 :: 1-AUC, mean +- SD over 5 seeds
\newcommand{\ProbeBkgLeftVisibleContentMlpOneeighteight}{0.4}  % experiments/FIGS/data/probe_ladder_v2_mlp2/analysis/seed_level_results.json :: levels.visible_content.mlp[0].rejection_points['0.90'].mean_n_bkg_pass :: sha256 53e1ec81343ac2c1 :: background jets passing the cut, mean over seeds
\newcommand{\ProbeRejVisibleContentMlpOneeighteight}{\ensuremath{>}538}  % experiments/FIGS/data/probe_ladder_v2_mlp2/analysis/seed_level_results.json :: table[visible_content,mlp,188].rejection_points['0.90'].rejection :: sha256 53e1ec81343ac2c1 :: at 90% signal efficiency, mean +- SD over seeds; 5 of 5 seeds at the cap
\newcommand{\ProbeAucVisibleContentMlpOnesixtwo}{0.99918}  % experiments/FIGS/data/probe_ladder_v2_mlp2/analysis/seed_level_results.json :: levels.visible_content.mlp[1].mean_auc :: sha256 53e1ec81343ac2c1 :: 5 seeds, 0 saturated at AUC=1
\newcommand{\ProbeAucSdVisibleContentMlpOnesixtwo}{0.00003}  % experiments/FIGS/data/probe_ladder_v2_mlp2/analysis/seed_level_results.json :: table[visible_content,mlp,162].auc (SD over seeds, ddof=1) :: sha256 53e1ec81343ac2c1
\newcommand{\ProbeLogOneMinusAucVisibleContentMlpOnesixtwo}{\ensuremath{-}7.1010}  % experiments/FIGS/data/probe_ladder_v2_mlp2/analysis/seed_level_results.json :: levels.visible_content.mlp[1].mean :: sha256 53e1ec81343ac2c1 :: natural log of 1-AUC, lower is better
\newcommand{\ProbeLogOneMinusAucSdVisibleContentMlpOnesixtwo}{0.0371}  % experiments/FIGS/data/probe_ladder_v2_mlp2/analysis/seed_level_results.json :: levels.visible_content.mlp[1].seed_sd :: sha256 53e1ec81343ac2c1
\newcommand{\ProbeOmaVisibleContentMlpOnesixtwo}{\ensuremath{(8.25 \pm 0.30)\times10^{-4}}}  % experiments/FIGS/data/probe_ladder_v2_mlp2/analysis/seed_level_results.json :: table[visible_content,mlp,162].auc :: sha256 53e1ec81343ac2c1 :: 1-AUC, mean +- SD over 5 seeds
\newcommand{\ProbeBkgLeftVisibleContentMlpOnesixtwo}{0.2}  % experiments/FIGS/data/probe_ladder_v2_mlp2/analysis/seed_level_results.json :: levels.visible_content.mlp[1].rejection_points['0.90'].mean_n_bkg_pass :: sha256 53e1ec81343ac2c1 :: background jets passing the cut, mean over seeds
\newcommand{\ProbeRejVisibleContentMlpOnesixtwo}{\ensuremath{>}538}  % experiments/FIGS/data/probe_ladder_v2_mlp2/analysis/seed_level_results.json :: table[visible_content,mlp,162].rejection_points['0.90'].rejection :: sha256 53e1ec81343ac2c1 :: at 90% signal efficiency, mean +- SD over seeds; 5 of 5 seeds at the cap
\newcommand{\ProbeAucVisibleContentMlpFourthree}{0.99917}  % experiments/FIGS/data/probe_ladder_v2_mlp2/analysis/seed_level_results.json :: levels.visible_content.mlp[2].mean_auc :: sha256 53e1ec81343ac2c1 :: 5 seeds, 0 saturated at AUC=1
\newcommand{\ProbeAucSdVisibleContentMlpFourthree}{0.00006}  % experiments/FIGS/data/probe_ladder_v2_mlp2/analysis/seed_level_results.json :: table[visible_content,mlp,43].auc (SD over seeds, ddof=1) :: sha256 53e1ec81343ac2c1
\newcommand{\ProbeLogOneMinusAucVisibleContentMlpFourthree}{\ensuremath{-}7.0957}  % experiments/FIGS/data/probe_ladder_v2_mlp2/analysis/seed_level_results.json :: levels.visible_content.mlp[2].mean :: sha256 53e1ec81343ac2c1 :: natural log of 1-AUC, lower is better
\newcommand{\ProbeLogOneMinusAucSdVisibleContentMlpFourthree}{0.0773}  % experiments/FIGS/data/probe_ladder_v2_mlp2/analysis/seed_level_results.json :: levels.visible_content.mlp[2].seed_sd :: sha256 53e1ec81343ac2c1
\newcommand{\ProbeOmaVisibleContentMlpFourthree}{\ensuremath{(8.3 \pm 0.6)\times10^{-4}}}  % experiments/FIGS/data/probe_ladder_v2_mlp2/analysis/seed_level_results.json :: table[visible_content,mlp,43].auc :: sha256 53e1ec81343ac2c1 :: 1-AUC, mean +- SD over 5 seeds
\newcommand{\ProbeBkgLeftVisibleContentMlpFourthree}{0.4}  % experiments/FIGS/data/probe_ladder_v2_mlp2/analysis/seed_level_results.json :: levels.visible_content.mlp[2].rejection_points['0.90'].mean_n_bkg_pass :: sha256 53e1ec81343ac2c1 :: background jets passing the cut, mean over seeds
\newcommand{\ProbeRejVisibleContentMlpFourthree}{\ensuremath{>}538}  % experiments/FIGS/data/probe_ladder_v2_mlp2/analysis/seed_level_results.json :: table[visible_content,mlp,43].rejection_points['0.90'].rejection :: sha256 53e1ec81343ac2c1 :: at 90% signal efficiency, mean +- SD over seeds; 5 of 5 seeds at the cap
\newcommand{\ProbeAucVisibleContentMlpOneseven}{0.99793}  % experiments/FIGS/data/probe_ladder_v2_mlp2/analysis/seed_level_results.json :: levels.visible_content.mlp[3].mean_auc :: sha256 53e1ec81343ac2c1 :: 5 seeds, 0 saturated at AUC=1
\newcommand{\ProbeAucSdVisibleContentMlpOneseven}{0.00020}  % experiments/FIGS/data/probe_ladder_v2_mlp2/analysis/seed_level_results.json :: table[visible_content,mlp,17].auc (SD over seeds, ddof=1) :: sha256 53e1ec81343ac2c1
\newcommand{\ProbeLogOneMinusAucVisibleContentMlpOneseven}{\ensuremath{-}6.1839}  % experiments/FIGS/data/probe_ladder_v2_mlp2/analysis/seed_level_results.json :: levels.visible_content.mlp[3].mean :: sha256 53e1ec81343ac2c1 :: natural log of 1-AUC, lower is better
\newcommand{\ProbeLogOneMinusAucSdVisibleContentMlpOneseven}{0.1014}  % experiments/FIGS/data/probe_ladder_v2_mlp2/analysis/seed_level_results.json :: levels.visible_content.mlp[3].seed_sd :: sha256 53e1ec81343ac2c1
\newcommand{\ProbeOmaVisibleContentMlpOneseven}{\ensuremath{(2.07 \pm 0.20)\times10^{-3}}}  % experiments/FIGS/data/probe_ladder_v2_mlp2/analysis/seed_level_results.json :: table[visible_content,mlp,17].auc :: sha256 53e1ec81343ac2c1 :: 1-AUC, mean +- SD over 5 seeds
\newcommand{\ProbeBkgLeftVisibleContentMlpOneseven}{3.4}  % experiments/FIGS/data/probe_ladder_v2_mlp2/analysis/seed_level_results.json :: levels.visible_content.mlp[3].rejection_points['0.90'].mean_n_bkg_pass :: sha256 53e1ec81343ac2c1 :: background jets passing the cut, mean over seeds
\newcommand{\ProbeRejVisibleContentMlpOneseven}{\ensuremath{\geq}638\ensuremath{^{\ast}}}  % experiments/FIGS/data/probe_ladder_v2_mlp2/analysis/seed_level_results.json :: table[visible_content,mlp,17].rejection_points['0.90'].rejection :: sha256 53e1ec81343ac2c1 :: at 90% signal efficiency, mean +- SD over seeds; 1 of 5 seeds at the cap
\newcommand{\ProbeOmaRatioVisibleContentMlpOnesixtwoOverOneeighteight}{0.92}  % experiments/FIGS/data/probe_ladder_v2_mlp2/analysis/seed_level_results.json :: table[visible_content,mlp,162/188].auc :: sha256 53e1ec81343ac2c1 :: ratio of the seed means of 1-AUC, coarser over finer
\newcommand{\ProbeOmaRatioVisibleContentMlpFourthreeOverOneeighteight}{0.93}  % experiments/FIGS/data/probe_ladder_v2_mlp2/analysis/seed_level_results.json :: table[visible_content,mlp,43/188].auc :: sha256 53e1ec81343ac2c1 :: ratio of the seed means of 1-AUC, coarser over finer
\newcommand{\ProbeOmaRatioVisibleContentMlpOnesevenOverOneeighteight}{2.3}  % experiments/FIGS/data/probe_ladder_v2_mlp2/analysis/seed_level_results.json :: table[visible_content,mlp,17/188].auc :: sha256 53e1ec81343ac2c1 :: ratio of the seed means of 1-AUC, coarser over finer
\newcommand{\ProbeOmaRatioVisibleContentMlpFourthreeOverOnesixtwo}{1.01}  % experiments/FIGS/data/probe_ladder_v2_mlp2/analysis/seed_level_results.json :: table[visible_content,mlp,43/162].auc :: sha256 53e1ec81343ac2c1 :: ratio of the seed means of 1-AUC, coarser over finer
\newcommand{\ProbeOmaRatioVisibleContentMlpOnesevenOverOnesixtwo}{2.5}  % experiments/FIGS/data/probe_ladder_v2_mlp2/analysis/seed_level_results.json :: table[visible_content,mlp,17/162].auc :: sha256 53e1ec81343ac2c1 :: ratio of the seed means of 1-AUC, coarser over finer
\newcommand{\ProbeOmaRatioVisibleContentMlpOnesevenOverFourthree}{2.5}  % experiments/FIGS/data/probe_ladder_v2_mlp2/analysis/seed_level_results.json :: table[visible_content,mlp,17/43].auc :: sha256 53e1ec81343ac2c1 :: ratio of the seed means of 1-AUC, coarser over finer
\newcommand{\VocabSizeLoneeighteight}{188}  % configs/labelmaps/rung_label_maps.v1.csv :: column L188 (distinct group count) :: sha256 412f9aa3b58db412
\newcommand{\VocabSizeLonesixtwo}{162}  % configs/labelmaps/rung_label_maps.v1.csv :: column L162 (distinct group count) :: sha256 412f9aa3b58db412
\newcommand{\VocabSizeRsixthreeQone}{64}  % configs/labelmaps/rung_label_maps.v1.csv :: column R63_Q1 (distinct group count) :: sha256 412f9aa3b58db412
\newcommand{\VocabSizeRfourtwoQone}{43}  % configs/labelmaps/rung_label_maps.v1.csv :: column R42_Q1 (distinct group count) :: sha256 412f9aa3b58db412
\newcommand{\VocabSizeRtwonineQone}{30}  % configs/labelmaps/rung_label_maps.v1.csv :: column R29_Q1 (distinct group count) :: sha256 412f9aa3b58db412
\newcommand{\VocabSizeRonesixQone}{17}  % configs/labelmaps/rung_label_maps.v1.csv :: column R16_Q1 (distinct group count) :: sha256 412f9aa3b58db412
\newcommand{\VocabSizeRthreeVIS}{4}  % configs/labelmaps/rung_label_maps.v1.csv :: column R3_VIS (distinct group count) :: sha256 412f9aa3b58db412
\newcommand{\VocabSizeRoneQone}{2}  % configs/labelmaps/rung_label_maps.v1.csv :: column R1_Q1 (distinct group count) :: sha256 412f9aa3b58db412
\newcommand{\ProbeSplitTrain}{60\%}  % experiments/EVAL/probe.py :: SPLIT_FRACTIONS :: sha256 485db0966e660588 :: share of the probe sample; the same jets for every model
\newcommand{\ProbeSplitVal}{20\%}  % experiments/EVAL/probe.py :: SPLIT_FRACTIONS :: sha256 485db0966e660588 :: share of the probe sample; the same jets for every model
\newcommand{\ProbeSplitTest}{20\%}  % experiments/EVAL/probe.py :: SPLIT_FRACTIONS :: sha256 485db0966e660588 :: share of the probe sample; the same jets for every model
\newcommand{\ProbeMlpHidden}{256}  % experiments/EVAL/probe.py :: _fit_mlp :: sha256 485db0966e660588 :: hidden width of the MLP probe
\newcommand{\ProbeMlpInits}{3}  % experiments/EVAL/probe.py :: MLP_SEEDS :: sha256 485db0966e660588 :: MLP initialisations whose scores are averaged per cell
\newcommand{\ProbeMlpLrDrop}{10}  % experiments/EVAL/probe.py :: MLP_SCHEDULE['lr_factor'] :: sha256 485db0966e660588 :: factor the MLP's learning rate drops by on a plateau
\newcommand{\ProbeMlpPlateaus}{three}  % experiments/EVAL/probe.py :: MLP_SCHEDULE: stop_patience / (lr_patience + 1) :: sha256 485db0966e660588 :: patience windows without gain after which the MLP stops
\newcommand{\ProbeCMin}{\ensuremath{10^{-2}}}  % experiments/EVAL/probe.py :: C_GRID :: sha256 485db0966e660588 :: end of the logistic-regression inverse-regularisation grid
\newcommand{\ProbeCMax}{\ensuremath{10^{2}}}  % experiments/EVAL/probe.py :: C_GRID :: sha256 485db0966e660588 :: end of the logistic-regression inverse-regularisation grid
\newcommand{\ProbeNBkgTestBvcFourprong}{3{,}158}  % experiments/FIGS/data/probe_ladder_v2_mlp2/s1.json :: tasks.bvc_4prong.n_background_test :: sha256 f7e1de0e43be89b8 :: test-sample background jets
\newcommand{\ProbeNBkgTestBvcQcd}{5{,}371}  % experiments/FIGS/data/probe_ladder_v2_mlp2/s1.json :: tasks.bvc_qcd.n_background_test :: sha256 f7e1de0e43be89b8 :: test-sample background jets
\newcommand{\ProbeNBkgTestBvcResonant}{11{,}876}  % experiments/FIGS/data/probe_ladder_v2_mlp2/s1.json :: tasks.bvc_resonant.n_background_test :: sha256 f7e1de0e43be89b8 :: test-sample background jets
\newcommand{\ProbeNBkgTestEeVsMm}{4{,}708}  % experiments/FIGS/data/probe_ladder_v2_mlp2/s1.json :: tasks.ee_vs_mm.n_background_test :: sha256 f7e1de0e43be89b8 :: test-sample background jets
\newcommand{\ProbeNBkgTestRetainedTopology}{1{,}078}  % experiments/FIGS/data/probe_ladder_v2_mlp2/s1.json :: tasks.retained_topology.n_background_test :: sha256 f7e1de0e43be89b8 :: test-sample background jets
\newcommand{\ProbeNBkgTestVisibleContent}{2{,}554}  % experiments/FIGS/data/probe_ladder_v2_mlp2/s1.json :: tasks.visible_content.n_background_test :: sha256 f7e1de0e43be89b8 :: test-sample background jets
\newcommand{\UseDiesAtSophonEqfour}{43}  % configs/labelmaps/usecase_survival.v1.json :: sophon_eq4.dies_at :: sha256 122275213c349593 :: first vocabulary size at which it is not constructible
\newcommand{\UseLastAliveSophonEqfour}{64}  % configs/labelmaps/usecase_survival.v1.json :: sophon_eq4.last_rung_alive :: sha256 122275213c349593
\newcommand{\UseDiesAtSophonEqsixAone}{64}  % configs/labelmaps/usecase_survival.v1.json :: sophon_eq6_a1.dies_at :: sha256 122275213c349593 :: first vocabulary size at which it is not constructible
\newcommand{\UseLastAliveSophonEqsixAone}{162}  % configs/labelmaps/usecase_survival.v1.json :: sophon_eq6_a1.last_rung_alive :: sha256 122275213c349593
\newcommand{\UseDiesAtSophonEqsixAthree}{64}  % configs/labelmaps/usecase_survival.v1.json :: sophon_eq6_a3.dies_at :: sha256 122275213c349593 :: first vocabulary size at which it is not constructible
\newcommand{\UseLastAliveSophonEqsixAthree}{162}  % configs/labelmaps/usecase_survival.v1.json :: sophon_eq6_a3.last_rung_alive :: sha256 122275213c349593
\newcommand{\UseDiesAtSophonEqsevenNine}{64}  % configs/labelmaps/usecase_survival.v1.json :: sophon_eq7_9.dies_at :: sha256 122275213c349593 :: first vocabulary size at which it is not constructible
\newcommand{\UseLastAliveSophonEqsevenNine}{162}  % configs/labelmaps/usecase_survival.v1.json :: sophon_eq7_9.last_rung_alive :: sha256 122275213c349593
\newcommand{\UseDiesAtVcbEqone}{64}  % configs/labelmaps/usecase_survival.v1.json :: vcb_eq1.dies_at :: sha256 122275213c349593 :: first vocabulary size at which it is not constructible
\newcommand{\UseLastAliveVcbEqone}{162}  % configs/labelmaps/usecase_survival.v1.json :: vcb_eq1.last_rung_alive :: sha256 122275213c349593
\newcommand{\UseDiesAtGnthreexQcdSubclass}{162}  % configs/labelmaps/usecase_survival.v1.json :: gn3x_qcd_subclass.dies_at :: sha256 122275213c349593 :: first vocabulary size at which it is not constructible
\newcommand{\UseLastAliveGnthreexQcdSubclass}{188}  % configs/labelmaps/usecase_survival.v1.json :: gn3x_qcd_subclass.last_rung_alive :: sha256 122275213c349593
\newcommand{\UseWlikeLastAlive}{162}  % configs/labelmaps/rung_label_maps.v1.csv :: groups at each level vs {label_X_cs, label_X_qq} and that plus the QCD classes :: sha256 412f9aa3b58db412 :: last level at which Sophon's W-like score g_W(2) is constructible
\newcommand{\UseWlikeDiesAt}{64}  % configs/labelmaps/rung_label_maps.v1.csv :: groups at each level vs {label_X_cs, label_X_qq} and that plus the QCD classes :: sha256 412f9aa3b58db412 :: first level at which it is not
\newcommand{\DesignEpochs}{80}  % experiments/MTX/k8s/job-mtx-l188-s1-raunav.yaml :: --num-epochs :: sha256 96fef4b6337aaa27
\newcommand{\DesignSamplesPerEpoch}{10{,}240{,}000}  % experiments/MTX/k8s/job-mtx-l188-s1-raunav.yaml :: --samples-per-epoch :: sha256 96fef4b6337aaa27
\newcommand{\DesignExamplesSeen}{\ensuremath{8.192\times10^{8}}}  % experiments/MTX/k8s/job-mtx-l188-s1-raunav.yaml :: --num-epochs x --samples-per-epoch :: sha256 96fef4b6337aaa27 :: training examples drawn over the run
\newcommand{\DesignPassesSelected}{7.45}  % experiments/FIGS/data/v2/jc2_train/counts.json :: n_selected :: sha256 b0b5a8b63aa808c0 :: --num-epochs x --samples-per-epoch of job-mtx-l188-s1-raunav.yaml / n_selected: passes over the selected training jets
\newcommand{\DesignBatchSize}{512}  % experiments/MTX/k8s/job-mtx-l188-s1-raunav.yaml :: --batch-size :: sha256 96fef4b6337aaa27
\newcommand{\DesignStartLr}{\ensuremath{5\times10^{-4}}}  % experiments/MTX/k8s/job-mtx-l188-s1-raunav.yaml :: --start-lr :: sha256 96fef4b6337aaa27
\newcommand{\DesignHeadHidden}{512}  % experiments/MTX/k8s/job-mtx-l188-s1-raunav.yaml :: -o fc_params (hidden width) :: sha256 96fef4b6337aaa27
\newcommand{\DesignParticleBlocks}{8}  % experiments/E1/ParT_sophon_arch_10c.py :: num_layers=(\d+) :: sha256 5793faa7d55a8b0d
\newcommand{\DesignClassBlocks}{2}  % experiments/E1/ParT_sophon_arch_10c.py :: num_cls_layers=(\d+) :: sha256 5793faa7d55a8b0d
\newcommand{\DesignHeads}{8}  % experiments/E1/ParT_sophon_arch_10c.py :: num_heads=(\d+) :: sha256 5793faa7d55a8b0d
\newcommand{\DesignEmbedDim}{128}  % experiments/E1/ParT_sophon_arch_10c.py :: \bembed_dims=\[\d+, \d+, (\d+)\] :: sha256 5793faa7d55a8b0d
\newcommand{\DesignMpmMaskRate}{40\%}  % experiments/MTX/mpm.py :: DEFAULT_MASK_RATE :: sha256 49f0c8c6666d0211 :: share of each jet's particles hidden from the encoder
\newcommand{\DesignPtMin}{200}  % configs/arms/L188.yaml :: selection :: sha256 d5a188dc80885515 :: GeV
\newcommand{\DesignPtMax}{2{,}500}  % configs/arms/L188.yaml :: selection :: sha256 d5a188dc80885515 :: GeV
\newcommand{\DesignMsdMin}{20}  % configs/arms/L188.yaml :: selection :: sha256 d5a188dc80885515 :: GeV
\newcommand{\DesignMsdMax}{500}  % configs/arms/L188.yaml :: selection :: sha256 d5a188dc80885515 :: GeV
\newcommand{\LitJetClassTwoNJetsMillion}{139}  % experiments/FIGS/data/literature_facts.json :: facts[0].value :: sha256 713dba2a5e6e1dd1 :: arXiv:2405.12972v1 App. A, p. 11
\newcommand{\LitJetClassTwoNResonantClasses}{161}  % experiments/FIGS/data/literature_facts.json :: facts[1].value :: sha256 713dba2a5e6e1dd1 :: arXiv:2405.12972v1 App. A, p. 11
\newcommand{\LitJetClassTwoNQcdClasses}{27}  % experiments/FIGS/data/literature_facts.json :: facts[2].value :: sha256 713dba2a5e6e1dd1 :: arXiv:2405.12972v1 App. A, p. 11
\newcommand{\LitJetClassTwoMaxResonanceMass}{500}  % experiments/FIGS/data/literature_facts.json :: facts[3].value :: sha256 713dba2a5e6e1dd1 :: arXiv:2405.12972v1 Sec. II, p. 2
\newcommand{\LitMadgraphVersion}{2.9.18}  % experiments/FIGS/data/literature_facts.json :: facts[4].value :: sha256 713dba2a5e6e1dd1 :: arXiv:2405.12972v1 App. A, p. 11
\newcommand{\LitPythiaVersion}{8.3}  % experiments/FIGS/data/literature_facts.json :: facts[5].value :: sha256 713dba2a5e6e1dd1 :: arXiv:2405.12972v1 App. A, p. 11
\newcommand{\LitDelphesVersion}{3}  % experiments/FIGS/data/literature_facts.json :: facts[6].value :: sha256 713dba2a5e6e1dd1 :: arXiv:2405.12972v1 App. A, p. 11
\newcommand{\LitJetClassTwoPileupMean}{50}  % experiments/FIGS/data/literature_facts.json :: facts[7].value :: sha256 713dba2a5e6e1dd1 :: arXiv:2405.12972v1 Sec. II, p. 2
\newcommand{\LitJetClassTwoJetRadius}{0.8}  % experiments/FIGS/data/literature_facts.json :: facts[8].value :: sha256 713dba2a5e6e1dd1 :: arXiv:2405.12972v1 Sec. II, p. 2
\newcommand{\LitAtlasUnifiedEffDropPercent}{10}  % experiments/FIGS/data/literature_facts.json :: facts[9].value :: sha256 713dba2a5e6e1dd1 :: ATL-PHYS-PUB-2026-013 Sec. 6.5, p. 22
\newcommand{\LitSophonPaperParticleBlocks}{6}  % experiments/FIGS/data/literature_facts.json :: facts[10].value :: sha256 713dba2a5e6e1dd1 :: arXiv:2405.12972v1 App. B.1, p. 12
\newcommand{\LitSophonSamplesPerEpochMillion}{10}  % experiments/FIGS/data/literature_facts.json :: facts[11].value :: sha256 713dba2a5e6e1dd1 :: arXiv:2405.12972v1 App. B.1, p. 12
\newcommand{\LitSophonSubsetSamplesPerEpochMillion}{2.5}  % experiments/FIGS/data/literature_facts.json :: facts[12].value :: sha256 713dba2a5e6e1dd1 :: arXiv:2405.12972v1 App. B.2, p. 12
\newcommand{\LitGlopartNClasses}{314}  % experiments/FIGS/data/literature_facts.json :: facts[13].value :: sha256 713dba2a5e6e1dd1 :: C. Li, PhD thesis, CERN-THESIS-2024-281 Sec. 11.2.1, p. 188
\newcommand{\LitPartFtEpochs}{20}  % experiments/FIGS/data/literature_facts.json :: facts[14].value :: sha256 713dba2a5e6e1dd1 :: arXiv:2202.03772 Sec. 5.2, p. 8
\newcommand{\LitMehtaNCategories}{eight}  % experiments/FIGS/data/literature_facts.json :: facts[15].value :: sha256 713dba2a5e6e1dd1 :: arXiv:2605.05556v2 Abstract
\newcommand{\LitHongPretrainRuns}{2}  % experiments/FIGS/data/literature_facts.json :: facts[16].value :: sha256 713dba2a5e6e1dd1 :: arXiv:2410.23129 App. A.1.1, p. 18
\newcommand{\LitSophonLookaheadK}{6}  % experiments/FIGS/data/literature_facts.json :: facts[17].value :: sha256 713dba2a5e6e1dd1 :: arXiv:2405.12972v1 App. B.1, p. 12
\newcommand{\LitSophonLookaheadAlpha}{0.5}  % experiments/FIGS/data/literature_facts.json :: facts[18].value :: sha256 713dba2a5e6e1dd1 :: arXiv:2405.12972v1 App. B.1, p. 12
\newcommand{\LitSophonRadamBetaOne}{0.95}  % experiments/FIGS/data/literature_facts.json :: facts[19].value :: sha256 713dba2a5e6e1dd1 :: arXiv:2405.12972v1 App. B.1, p. 12
\newcommand{\LitSophonRadamBetaTwo}{0.999}  % experiments/FIGS/data/literature_facts.json :: facts[20].value :: sha256 713dba2a5e6e1dd1 :: arXiv:2405.12972v1 App. B.1, p. 12
\newcommand{\LitSophonRadamEps}{\ensuremath{10^{-5}}}  % experiments/FIGS/data/literature_facts.json :: facts[21].value :: sha256 713dba2a5e6e1dd1 :: arXiv:2405.12972v1 App. B.1, p. 12
\newcommand{\LitGnTwoAuxRejectionDropPercent}{30}  % experiments/FIGS/data/literature_facts.json :: facts[22].value :: sha256 713dba2a5e6e1dd1 :: arXiv:2505.19689v2 (Nature Commun. 17 (2026) 541) p. 9
\newcommand{\LitCmsPartNClasses}{37}  % experiments/FIGS/data/literature_facts.json :: facts[23].value :: sha256 713dba2a5e6e1dd1 :: arXiv:2604.09809v1 (CMS-JME-25-001) Sec. 5.2, p. 7
\newcommand{\LitCmsPartMassLambda}{0.05}  % experiments/FIGS/data/literature_facts.json :: facts[24].value :: sha256 713dba2a5e6e1dd1 :: arXiv:2604.09809v1 (CMS-JME-25-001) Sec. 5.2, p. 7
\newcommand{\LitGlopartVthreeNClasses}{374}  % experiments/FIGS/data/literature_facts.json :: facts[25].value :: sha256 713dba2a5e6e1dd1 :: CMS-DP-2026-104 (CMS public page) Architecture and outputs
\newcommand{\LitHuhCoarseClasses}{127}  % experiments/FIGS/data/literature_facts.json :: facts[26].value :: sha256 713dba2a5e6e1dd1 :: arXiv:1608.08614v2 Sec. 5.3, p. 5
\newcommand{\LitHuhKnnGapPercent}{15}  % experiments/FIGS/data/literature_facts.json :: facts[27].value :: sha256 713dba2a5e6e1dd1 :: arXiv:1608.08614v2 Sec. 5.3, p. 5
\newcommand{\LitRsthreelRuns}{3}  % experiments/FIGS/data/literature_facts.json :: facts[28].value :: sha256 713dba2a5e6e1dd1 :: arXiv:2403.07066v2 (Phys. Rev. D 111 (2025) 032010) Sec. IV, p. 5
\newcommand{\LitSophonParamsMillion}{2.3}  % experiments/FIGS/data/literature_facts.json :: facts[29].value :: sha256 713dba2a5e6e1dd1 :: arXiv:2405.12972v1 App. B.1, p. 12
\newcommand{\LitGlopartVoneNClasses}{37}  % experiments/FIGS/data/literature_facts.json :: facts[30].value :: sha256 713dba2a5e6e1dd1 :: CMS-DP-2026-104 (CMS public page) Model generations
\newcommand{\LitSophonXMass}{200}  % experiments/FIGS/data/literature_facts.json :: facts[31].value :: sha256 713dba2a5e6e1dd1 :: arXiv:2405.12972v1 Sec. III, p. 3
\newcommand{\LitSophonXLumi}{40}  % experiments/FIGS/data/literature_facts.json :: facts[32].value :: sha256 713dba2a5e6e1dd1 :: arXiv:2405.12972v1 Sec. III, Fig. 2
\newcommand{\LitSophonXWindowLo}{150}  % experiments/FIGS/data/literature_facts.json :: facts[33].value :: sha256 713dba2a5e6e1dd1 :: arXiv:2405.12972v1 Sec. III, Fig. 2
\newcommand{\LitSophonXWindowHi}{230}  % experiments/FIGS/data/literature_facts.json :: facts[34].value :: sha256 713dba2a5e6e1dd1 :: arXiv:2405.12972v1 Sec. III, Fig. 2
\newcommand{\LitSophonAblationNClasses}{42}  % experiments/FIGS/data/literature_facts.json :: facts[35].value :: sha256 713dba2a5e6e1dd1 :: arXiv:2405.12972v1 Sec. III
\newcommand{\LitSophonTransferSamplesMillion}{2.5}  % experiments/FIGS/data/literature_facts.json :: facts[36].value :: sha256 713dba2a5e6e1dd1 :: arXiv:2405.12972v1 App. B.5, p. 13
\newcommand{\LitSophonTransferBatch}{1024}  % experiments/FIGS/data/literature_facts.json :: facts[37].value :: sha256 713dba2a5e6e1dd1 :: arXiv:2405.12972v1 App. B.5, p. 13
\newcommand{\LitSophonScratchEpochs}{50}  % experiments/FIGS/data/literature_facts.json :: facts[38].value :: sha256 713dba2a5e6e1dd1 :: arXiv:2405.12972v1 App. B.3
\newcommand{\LitHuhFineClasses}{1000}  % experiments/FIGS/data/literature_facts.json :: facts[39].value :: sha256 713dba2a5e6e1dd1 :: arXiv:1608.08614v2 Sec. 1, p. 1
\newcommand{\LitHuhPascalDropMap}{2.8}  % experiments/FIGS/data/literature_facts.json :: facts[40].value :: sha256 713dba2a5e6e1dd1 :: arXiv:1608.08614v2 Sec. 1, p. 1
\newcommand{\LitHongSuperclasses}{11}  % experiments/FIGS/data/literature_facts.json :: facts[41].value :: sha256 713dba2a5e6e1dd1 :: arXiv:2410.23129 Sec. 6.2
\newcommand{\LitHongDeclineClasses}{\ensuremath{10^{3}}}  % experiments/FIGS/data/literature_facts.json :: facts[42].value :: sha256 713dba2a5e6e1dd1 :: arXiv:2410.23129 Sec. 6.2
\newcommand{\LitObjectivesFullJets}{\ensuremath{1.2\times10^{6}}}  % experiments/FIGS/data/literature_facts.json :: facts[43].value :: sha256 713dba2a5e6e1dd1 :: arXiv:2606.14870 Sec. IV A
\newcommand{\LitObjectivesFewJets}{\ensuremath{10^{4}}}  % experiments/FIGS/data/literature_facts.json :: facts[44].value :: sha256 713dba2a5e6e1dd1 :: arXiv:2606.14870 Sec. IV A
\newcommand{\LitMpmBenefitJets}{10k}  % experiments/FIGS/data/literature_facts.json :: facts[45].value :: sha256 713dba2a5e6e1dd1 :: arXiv:2401.13537 Sec. V
\newcommand{\LitPartTopRejScratch}{1602}  % experiments/FIGS/data/literature_facts.json :: facts[46].value :: sha256 713dba2a5e6e1dd1 :: arXiv:2202.03772 Table 5
\newcommand{\LitPartTopRejFt}{2766}  % experiments/FIGS/data/literature_facts.json :: facts[47].value :: sha256 713dba2a5e6e1dd1 :: arXiv:2202.03772 Table 5
\newcommand{\LitMahajanLinearGapPercent}{7}  % experiments/FIGS/data/literature_facts.json :: facts[48].value :: sha256 713dba2a5e6e1dd1 :: arXiv:1805.00932 Sec. 3.1.2
\newcommand{\LitBukchinCoarse}{33.83}  % experiments/FIGS/data/literature_facts.json :: facts[49].value :: sha256 713dba2a5e6e1dd1 :: arXiv:2012.03515 Table 2, LIVING-17 all-way
\newcommand{\LitBukchinFine}{58.95}  % experiments/FIGS/data/literature_facts.json :: facts[50].value :: sha256 713dba2a5e6e1dd1 :: arXiv:2012.03515 Table 2, LIVING-17 all-way
\newcommand{\LitKornblithRho}{\ensuremath{-0.93}}  % experiments/FIGS/data/literature_facts.json :: facts[51].value :: sha256 713dba2a5e6e1dd1 :: arXiv:2010.16402 Sec. 3.4
\newcommand{\LitPartTopSigEffPercent}{30}  % experiments/FIGS/data/literature_facts.json :: facts[52].value :: sha256 713dba2a5e6e1dd1 :: arXiv:2202.03772 Table 5 (column Rej30%)
\newcommand{\LitPartTopKinFeatures}{7}  % experiments/FIGS/data/literature_facts.json :: facts[53].value :: sha256 713dba2a5e6e1dd1 :: arXiv:2202.03772 Sec. 5.2
\newcommand{\LitBukchinAncor}{45.14}  % experiments/FIGS/data/literature_facts.json :: facts[54].value :: sha256 713dba2a5e6e1dd1 :: arXiv:2012.03515 Table 2, LIVING-17 all-way
\newcommand{\LitMahajanFtGapPp}{0.1}  % experiments/FIGS/data/literature_facts.json :: facts[55].value :: sha256 713dba2a5e6e1dd1 :: arXiv:1805.00932 Fig. 1, ImageNet-9k target
\newcommand{\LitSophonTransferLr}{5 × 10−4}  % experiments/FIGS/data/literature_facts.json :: facts[56].value :: sha256 713dba2a5e6e1dd1 :: arXiv:2405.12972v1 App. B.5
\newcommand{\LitQgJetRadius}{0.4}  % experiments/FIGS/data/literature_facts.json :: facts[57].value :: sha256 713dba2a5e6e1dd1 :: arXiv:1810.05165 Sec. 4.1
\newcommand{\LitQgPtMin}{500}  % experiments/FIGS/data/literature_facts.json :: facts[58].value :: sha256 713dba2a5e6e1dd1 :: arXiv:1810.05165 Sec. 4.1
\newcommand{\LitQgPtMax}{550}  % experiments/FIGS/data/literature_facts.json :: facts[59].value :: sha256 713dba2a5e6e1dd1 :: arXiv:1810.05165 Sec. 4.1
\newcommand{\LitPartTopAucScratch}{0.9858}  % experiments/FIGS/data/literature_facts.json :: facts[60].value :: sha256 713dba2a5e6e1dd1 :: arXiv:2202.03772 Table 5
\newcommand{\LitPartTopAucFt}{0.9877}  % experiments/FIGS/data/literature_facts.json :: facts[61].value :: sha256 713dba2a5e6e1dd1 :: arXiv:2202.03772 Table 5
\newcommand{\LitPartQgAucScratch}{0.9121}  % experiments/FIGS/data/literature_facts.json :: facts[62].value :: sha256 713dba2a5e6e1dd1 :: arXiv:2202.03772 Table 6 (exp)
\newcommand{\LitPartQgRejScratch}{101.2}  % experiments/FIGS/data/literature_facts.json :: facts[63].value :: sha256 713dba2a5e6e1dd1 :: arXiv:2202.03772 Table 6 (exp)
\newcommand{\LitPartQgAucFt}{0.9151}  % experiments/FIGS/data/literature_facts.json :: facts[64].value :: sha256 713dba2a5e6e1dd1 :: arXiv:2202.03772 Table 6 (exp)
\newcommand{\LitPartQgRejFt}{107.9}  % experiments/FIGS/data/literature_facts.json :: facts[65].value :: sha256 713dba2a5e6e1dd1 :: arXiv:2202.03772 Table 6 (exp)
\newcommand{\LitPartJcTestMillion}{20}  % experiments/FIGS/data/literature_facts.json :: facts[66].value :: sha256 713dba2a5e6e1dd1 :: arXiv:2202.03772 Sec. 3
\newcommand{\LitPartJcTwoMAcc}{0.836}  % experiments/FIGS/data/literature_facts.json :: facts[67].value :: sha256 713dba2a5e6e1dd1 :: arXiv:2202.03772 Table 3
\newcommand{\LitPartSmallSetIters}{100 k}  % experiments/FIGS/data/literature_facts.json :: facts[68].value :: sha256 713dba2a5e6e1dd1 :: arXiv:2202.03772 Sec. 5.1 (training dataset size)
\newcommand{\LitPartSmallSetPercent}{2}  % experiments/FIGS/data/literature_facts.json :: facts[69].value :: sha256 713dba2a5e6e1dd1 :: arXiv:2202.03772 Sec. 5.1 (training dataset size)
\newcommand{\LitSophonEtaMax}{2.5}  % experiments/FIGS/data/literature_facts.json :: facts[70].value :: sha256 713dba2a5e6e1dd1 :: arXiv:2405.12972 App. A
\newcommand{\LitBHadronLifetimePs}{1.5}  % experiments/FIGS/data/literature_facts.json :: facts[71].value :: sha256 713dba2a5e6e1dd1 :: arXiv:1712.07158 Sec. 4.1
\newcommand{\LitCHadronLifetimePs}{1}  % experiments/FIGS/data/literature_facts.json :: facts[72].value :: sha256 713dba2a5e6e1dd1 :: arXiv:1712.07158 Sec. 4.1
\newcommand{\LitOmnilearnPidDropProb}{0.2}  % experiments/FIGS/data/literature_facts.json :: facts[73].value :: sha256 713dba2a5e6e1dd1 :: arXiv:2502.14652 Sec. II.A
\newcommand{\LitMahajanFtNineKLarge}{48.5}  % experiments/FIGS/data/literature_facts.json :: facts[74].value :: sha256 713dba2a5e6e1dd1 :: arXiv:1805.00932 Fig. 1 (IG-1B-17k, IN-9k, full fine-tuning)
\newcommand{\LitMahajanFtNineKSmall}{48.4}  % experiments/FIGS/data/literature_facts.json :: facts[75].value :: sha256 713dba2a5e6e1dd1 :: arXiv:1805.00932 Fig. 1 (IG-940M-1.5k, IN-9k, full fine-tuning)
\newcommand{\LitMahajanFtOneKSmall}{84.2}  % experiments/FIGS/data/literature_facts.json :: facts[76].value :: sha256 713dba2a5e6e1dd1 :: arXiv:1805.00932 Fig. 1 (IG-940M-1.5k, IN-1k, full fine-tuning)
\newcommand{\LitMahajanFtOneKLarge}{83.6}  % experiments/FIGS/data/literature_facts.json :: facts[77].value :: sha256 713dba2a5e6e1dd1 :: arXiv:1805.00932 Fig. 1 (IG-1B-17k, IN-1k, full fine-tuning)
\newcommand{\LitChenGainFifteen}{0.16}  % experiments/FIGS/data/literature_facts.json :: facts[78].value :: sha256 713dba2a5e6e1dd1 :: arXiv:1901.07012 Table VIII (15 coarse classes, CIFAR-100)
\newcommand{\LitChenGainTwenty}{0.01}  % experiments/FIGS/data/literature_facts.json :: facts[79].value :: sha256 713dba2a5e6e1dd1 :: arXiv:1901.07012 Table VIII (20 coarse classes, CIFAR-100)
\newcommand{\LitChenCoarseFifteen}{15}  % experiments/FIGS/data/literature_facts.json :: facts[80].value :: sha256 713dba2a5e6e1dd1 :: arXiv:1901.07012 Table VIII
\newcommand{\LitChenCoarseTwenty}{20}  % experiments/FIGS/data/literature_facts.json :: facts[81].value :: sha256 713dba2a5e6e1dd1 :: arXiv:1901.07012 Table VIII
\newcommand{\LitChenRandomMaxGain}{2.78}  % experiments/FIGS/data/literature_facts.json :: facts[82].value :: sha256 713dba2a5e6e1dd1 :: arXiv:1901.07012 Sec. VI-A, Table VI
\newcommand{\LitGlopartMassLambda}{5}  % experiments/FIGS/data/literature_facts.json :: facts[83].value :: sha256 713dba2a5e6e1dd1 :: CERN-THESIS-2024-281 Eq. (11.7), Sec. 11.2.1
\newcommand{\LitGlopartAddedAucBase}{0.9555}  % experiments/FIGS/data/literature_facts.json :: facts[84].value :: sha256 713dba2a5e6e1dd1 :: CERN-THESIS-2024-281 Sec. 11.4.1, Fig. 11.19 (H->WW 4q vs QCD)
\newcommand{\LitGlopartAddedAucTop}{0.9593}  % experiments/FIGS/data/literature_facts.json :: facts[85].value :: sha256 713dba2a5e6e1dd1 :: CERN-THESIS-2024-281 Sec. 11.4.1, Fig. 11.19 (H->WW 4q vs QCD)
\newcommand{\LitCmsSdZcut}{0.1}  % experiments/FIGS/data/literature_facts.json :: facts[86].value :: sha256 713dba2a5e6e1dd1 :: arXiv:2004.08262 Sec. 6.1
\newcommand{\LitCmsSdBeta}{0}  % experiments/FIGS/data/literature_facts.json :: facts[87].value :: sha256 713dba2a5e6e1dd1 :: arXiv:2004.08262 Sec. 6.1
\newcommand{\LitQgNJetsMillion}{2}  % experiments/FIGS/data/literature_facts.json :: facts[88].value :: sha256 713dba2a5e6e1dd1 :: arXiv:2202.03772 Sec. 5.2
\newcommand{\LitQgTrainMillion}{1.6}  % experiments/FIGS/data/literature_facts.json :: facts[89].value :: sha256 713dba2a5e6e1dd1 :: arXiv:2202.03772 Sec. 5.2
\newcommand{\LitQgValMillion}{0.2}  % experiments/FIGS/data/literature_facts.json :: facts[90].value :: sha256 713dba2a5e6e1dd1 :: arXiv:2202.03772 Sec. 5.2
\newcommand{\LitQgRapidityMax}{1.7}  % experiments/FIGS/data/literature_facts.json :: facts[91].value :: sha256 713dba2a5e6e1dd1 :: Zenodo 10.5281/zenodo.3164691 Zenodo record 3164691, description
\newcommand{\LitHerwigVersion}{7.1.4}  % experiments/FIGS/data/literature_facts.json :: facts[92].value :: sha256 713dba2a5e6e1dd1 :: Zenodo 10.5281/zenodo.3066475 record description
\newcommand{\LitObjectivesFewestJets}{\ensuremath{10^{3}}}  % experiments/FIGS/data/literature_facts.json :: facts[93].value :: sha256 713dba2a5e6e1dd1 :: arXiv:2606.14870 Table XVII (Micro model, Npt = 10^8: Classifier 136 vs MPM 218 at 10^3, 331 vs 539 at 10^4)
\newcommand{\LitPartTopNTrainMillion}{1.2}  % experiments/FIGS/data/literature_facts.json :: facts[94].value :: sha256 713dba2a5e6e1dd1 :: arXiv:2202.03772 Sec. 5.2
\newcommand{\LitLandscapeParticlenetEnsMean}{1412}  % experiments/FIGS/data/literature_facts.json :: facts[95].value :: sha256 713dba2a5e6e1dd1 :: arXiv:1902.09914 Table 1 (ParticleNet, ensemble mean of 1/eps_B at eps_S=0.3)
\newcommand{\LitPartFtHeadReinits}{9}  % experiments/FIGS/data/literature_facts.json :: facts[96].value :: sha256 713dba2a5e6e1dd1 :: arXiv:2202.03772 Sec. 5.2
\newcommand{\LitSophonBlocksCode}{8}  % experiments/FIGS/data/literature_facts.json :: facts[97].value :: sha256 713dba2a5e6e1dd1 :: github.com/jet-universe/sophon@9dd6dd6 networks/example_ParticleTransformer_sophon.py line 67
\newcommand{\LitGrasHerwigVersion}{2.7.1}  % experiments/FIGS/data/literature_facts.json :: facts[98].value :: sha256 713dba2a5e6e1dd1 :: arXiv:1704.03878 Sec. 2 (list of generators)
\newcommand{\DesignMassLambda}{5.0}  % experiments/MTX/k8s/job-mtx-l162_mass-s1-raunav.yaml :: --mass-lambda :: sha256 8882459fe5520795 :: identical in all ten mass-output pretraining specs
\newcommand{\FtRecipeNRuns}{264}  % experiments/FIGS/data/ft_recipes/recipes_w2b_bench_v2_v3.json :: runs (count), reported cells :: sha256 7dc8c791343e1b6f
\newcommand{\FtRecipeLr}{\ensuremath{1\times10^{-4}}}  % experiments/FIGS/data/ft_recipes/recipes_w2b_bench_v2_v3.json :: runs[*].lr, pretrained starts :: sha256 7dc8c791343e1b6f
\newcommand{\FtRecipeHeadMult}{50}  % experiments/FIGS/data/ft_recipes/recipes_w2b_bench_v2_v3.json :: runs[*].head_lr_mult, pretrained starts :: sha256 7dc8c791343e1b6f
\newcommand{\FtRecipeLrHead}{\ensuremath{5\times10^{-3}}}  % experiments/FIGS/data/ft_recipes/recipes_w2b_bench_v2_v3.json :: lr x head_lr_mult :: sha256 7dc8c791343e1b6f
\newcommand{\FtRecipeWeightDecay}{0.01}  % experiments/FIGS/data/ft_recipes/recipes_w2b_bench_v2_v3.json :: runs[*].weight_decay :: sha256 7dc8c791343e1b6f
\newcommand{\FtRecipeBatch}{512}  % experiments/FIGS/data/ft_recipes/recipes_w2b_bench_v2_v3.json :: runs[*].batch_size :: sha256 7dc8c791343e1b6f
\newcommand{\FtRecipeEpochsEThree}{50}  % experiments/FIGS/data/ft_recipes/recipes_w2b_bench_v2_v3.json :: runs[leg 1,2; n_train 1000].epochs :: sha256 7dc8c791343e1b6f
\newcommand{\FtRecipeStepsEThree}{19}  % experiments/FT/k8s/job-ft-legs-w2-raunav.yaml :: job commands: samples per epoch at n_train 1000 // batch size :: sha256 f0962eddb451b7b0 :: optimizer steps per epoch, as weaver floors them
\newcommand{\FtRecipePassesEThree}{10}  % experiments/FT/k8s/job-ft-legs-w2-raunav.yaml :: job commands: samples per epoch at n_train 1000 // n_train :: sha256 f0962eddb451b7b0 :: passes over the training subset per epoch
\newcommand{\FtRecipeEpochsEFour}{50}  % experiments/FIGS/data/ft_recipes/recipes_w2b_bench_v2_v3.json :: runs[leg 1,2; n_train 10000].epochs :: sha256 7dc8c791343e1b6f
\newcommand{\FtRecipeStepsEFour}{19}  % experiments/FT/k8s/job-ft-legs-w2-raunav.yaml :: job commands: samples per epoch at n_train 10000 // batch size :: sha256 f0962eddb451b7b0 :: optimizer steps per epoch, as weaver floors them
\newcommand{\FtRecipePassesEFour}{one}  % experiments/FT/k8s/job-ft-legs-w2-raunav.yaml :: job commands: samples per epoch at n_train 10000 // n_train :: sha256 f0962eddb451b7b0 :: passes over the training subset per epoch
\newcommand{\FtRecipeEpochsEFive}{30}  % experiments/FIGS/data/ft_recipes/recipes_w2b_bench_v2_v3.json :: runs[leg 1,2; n_train 100000].epochs :: sha256 7dc8c791343e1b6f
\newcommand{\FtRecipeStepsEFive}{195}  % experiments/FT/k8s/job-ft-legs-w2-raunav.yaml :: job commands: samples per epoch at n_train 100000 // batch size :: sha256 f0962eddb451b7b0 :: optimizer steps per epoch, as weaver floors them
\newcommand{\FtRecipePassesEFive}{one}  % experiments/FT/k8s/job-ft-legs-w2-raunav.yaml :: job commands: samples per epoch at n_train 100000 // n_train :: sha256 f0962eddb451b7b0 :: passes over the training subset per epoch
\newcommand{\FtRecipeEpochsESix}{10}  % experiments/FIGS/data/ft_recipes/recipes_w2b_bench_v2_v3.json :: runs[leg 1,2; n_train 1000000].epochs :: sha256 7dc8c791343e1b6f
\newcommand{\FtRecipeStepsESix}{1{,}953}  % experiments/FT/k8s/job-ft-legs-w2-raunav.yaml :: job commands: samples per epoch at n_train 1000000 // batch size :: sha256 f0962eddb451b7b0 :: optimizer steps per epoch, as weaver floors them
\newcommand{\FtRecipePassesESix}{one}  % experiments/FT/k8s/job-ft-legs-w2-raunav.yaml :: job commands: samples per epoch at n_train 1000000 // n_train :: sha256 f0962eddb451b7b0 :: passes over the training subset per epoch
\newcommand{\FtRecipeValJets}{19{,}968}  % experiments/FT/k8s/job-ft-legs-w2-raunav.yaml :: --samples-per-epoch-val 20000 // batch x batch :: sha256 f0962eddb451b7b0 :: validation jets per epoch, floored to whole batches
\newcommand{\RecoveryAccOneeighteightByOneeighteight}{0.4123\,\ensuremath{\pm}\,0.0011}  % experiments/FIGS/data/label_recovery_curve_v1err/analysis/label_recovery_curve_summary.json :: table[rung=L188,level=188,probe=linear,n_train=max].accuracy :: sha256 1250d1f00537e0bd :: balanced accuracy, mean +- SD over 5 runs
\newcommand{\RecoveryAccOneeighteightByOnesixtwo}{0.4090\,\ensuremath{\pm}\,0.0012}  % experiments/FIGS/data/label_recovery_curve_v1err/analysis/label_recovery_curve_summary.json :: table[rung=L188,level=162,probe=linear,n_train=max].accuracy :: sha256 1250d1f00537e0bd :: balanced accuracy, mean +- SD over 5 runs
\newcommand{\RecoveryAccOneeighteightByFourthree}{0.356\,\ensuremath{\pm}\,0.008}  % experiments/FIGS/data/label_recovery_curve_v1err/analysis/label_recovery_curve_summary.json :: table[rung=L188,level=43,probe=linear,n_train=max].accuracy :: sha256 1250d1f00537e0bd :: balanced accuracy, mean +- SD over 5 runs
\newcommand{\RecoveryAccOneeighteightByOneseven}{0.306\,\ensuremath{\pm}\,0.008}  % experiments/FIGS/data/label_recovery_curve_v1err/analysis/label_recovery_curve_summary.json :: table[rung=L188,level=17,probe=linear,n_train=max].accuracy :: sha256 1250d1f00537e0bd :: balanced accuracy, mean +- SD over 5 runs
\newcommand{\RecoveryAccOnesixtwoByOneeighteight}{0.4623\,\ensuremath{\pm}\,0.0006}  % experiments/FIGS/data/label_recovery_curve_v1err/analysis/label_recovery_curve_summary.json :: table[rung=L162,level=188,probe=linear,n_train=max].accuracy :: sha256 1250d1f00537e0bd :: balanced accuracy, mean +- SD over 5 runs
\newcommand{\RecoveryAccOnesixtwoByOnesixtwo}{0.4622\,\ensuremath{\pm}\,0.0007}  % experiments/FIGS/data/label_recovery_curve_v1err/analysis/label_recovery_curve_summary.json :: table[rung=L162,level=162,probe=linear,n_train=max].accuracy :: sha256 1250d1f00537e0bd :: balanced accuracy, mean +- SD over 5 runs
\newcommand{\RecoveryAccOnesixtwoByFourthree}{0.402\,\ensuremath{\pm}\,0.009}  % experiments/FIGS/data/label_recovery_curve_v1err/analysis/label_recovery_curve_summary.json :: table[rung=L162,level=43,probe=linear,n_train=max].accuracy :: sha256 1250d1f00537e0bd :: balanced accuracy, mean +- SD over 5 runs
\newcommand{\RecoveryAccOnesixtwoByOneseven}{0.348\,\ensuremath{\pm}\,0.010}  % experiments/FIGS/data/label_recovery_curve_v1err/analysis/label_recovery_curve_summary.json :: table[rung=L162,level=17,probe=linear,n_train=max].accuracy :: sha256 1250d1f00537e0bd :: balanced accuracy, mean +- SD over 5 runs
\newcommand{\RecoveryAccSixfourByOneeighteight}{0.5979\,\ensuremath{\pm}\,0.0008}  % experiments/FIGS/data/label_recovery_curve_v1err/analysis/label_recovery_curve_summary.json :: table[rung=R63_Q1,level=188,probe=linear,n_train=max].accuracy :: sha256 1250d1f00537e0bd :: balanced accuracy, mean +- SD over 5 runs
\newcommand{\RecoveryAccSixfourByOnesixtwo}{0.5973\,\ensuremath{\pm}\,0.0013}  % experiments/FIGS/data/label_recovery_curve_v1err/analysis/label_recovery_curve_summary.json :: table[rung=R63_Q1,level=162,probe=linear,n_train=max].accuracy :: sha256 1250d1f00537e0bd :: balanced accuracy, mean +- SD over 5 runs
\newcommand{\RecoveryAccSixfourByFourthree}{0.5881\,\ensuremath{\pm}\,0.0007}  % experiments/FIGS/data/label_recovery_curve_v1err/analysis/label_recovery_curve_summary.json :: table[rung=R63_Q1,level=43,probe=linear,n_train=max].accuracy :: sha256 1250d1f00537e0bd :: balanced accuracy, mean +- SD over 5 runs
\newcommand{\RecoveryAccSixfourByOneseven}{0.531\,\ensuremath{\pm}\,0.006}  % experiments/FIGS/data/label_recovery_curve_v1err/analysis/label_recovery_curve_summary.json :: table[rung=R63_Q1,level=17,probe=linear,n_train=max].accuracy :: sha256 1250d1f00537e0bd :: balanced accuracy, mean +- SD over 5 runs
\newcommand{\RecoveryAccFourthreeByOneeighteight}{0.6770\,\ensuremath{\pm}\,0.0006}  % experiments/FIGS/data/label_recovery_curve_v1err/analysis/label_recovery_curve_summary.json :: table[rung=R42_Q1,level=188,probe=linear,n_train=max].accuracy :: sha256 1250d1f00537e0bd :: balanced accuracy, mean +- SD over 5 runs
\newcommand{\RecoveryAccFourthreeByOnesixtwo}{0.6766\,\ensuremath{\pm}\,0.0008}  % experiments/FIGS/data/label_recovery_curve_v1err/analysis/label_recovery_curve_summary.json :: table[rung=R42_Q1,level=162,probe=linear,n_train=max].accuracy :: sha256 1250d1f00537e0bd :: balanced accuracy, mean +- SD over 5 runs
\newcommand{\RecoveryAccFourthreeByFourthree}{0.6751\,\ensuremath{\pm}\,0.0011}  % experiments/FIGS/data/label_recovery_curve_v1err/analysis/label_recovery_curve_summary.json :: table[rung=R42_Q1,level=43,probe=linear,n_train=max].accuracy :: sha256 1250d1f00537e0bd :: balanced accuracy, mean +- SD over 5 runs
\newcommand{\RecoveryAccFourthreeByOneseven}{0.633\,\ensuremath{\pm}\,0.006}  % experiments/FIGS/data/label_recovery_curve_v1err/analysis/label_recovery_curve_summary.json :: table[rung=R42_Q1,level=17,probe=linear,n_train=max].accuracy :: sha256 1250d1f00537e0bd :: balanced accuracy, mean +- SD over 5 runs
\newcommand{\RecoveryAccThreezeroByOneeighteight}{0.7001\,\ensuremath{\pm}\,0.0005}  % experiments/FIGS/data/label_recovery_curve_v1err/analysis/label_recovery_curve_summary.json :: table[rung=R29_Q1,level=188,probe=linear,n_train=max].accuracy :: sha256 1250d1f00537e0bd :: balanced accuracy, mean +- SD over 5 runs
\newcommand{\RecoveryAccThreezeroByOnesixtwo}{0.7004\,\ensuremath{\pm}\,0.0007}  % experiments/FIGS/data/label_recovery_curve_v1err/analysis/label_recovery_curve_summary.json :: table[rung=R29_Q1,level=162,probe=linear,n_train=max].accuracy :: sha256 1250d1f00537e0bd :: balanced accuracy, mean +- SD over 5 runs
\newcommand{\RecoveryAccThreezeroByFourthree}{0.6998\,\ensuremath{\pm}\,0.0008}  % experiments/FIGS/data/label_recovery_curve_v1err/analysis/label_recovery_curve_summary.json :: table[rung=R29_Q1,level=43,probe=linear,n_train=max].accuracy :: sha256 1250d1f00537e0bd :: balanced accuracy, mean +- SD over 5 runs
\newcommand{\RecoveryAccThreezeroByOneseven}{0.6778\,\ensuremath{\pm}\,0.0030}  % experiments/FIGS/data/label_recovery_curve_v1err/analysis/label_recovery_curve_summary.json :: table[rung=R29_Q1,level=17,probe=linear,n_train=max].accuracy :: sha256 1250d1f00537e0bd :: balanced accuracy, mean +- SD over 5 runs
\newcommand{\RecoveryAccOnesevenByOneeighteight}{0.7675\,\ensuremath{\pm}\,0.0010}  % experiments/FIGS/data/label_recovery_curve_v1err/analysis/label_recovery_curve_summary.json :: table[rung=R16_Q1,level=188,probe=linear,n_train=max].accuracy :: sha256 1250d1f00537e0bd :: balanced accuracy, mean +- SD over 5 runs
\newcommand{\RecoveryAccOnesevenByOnesixtwo}{0.7681\,\ensuremath{\pm}\,0.0009}  % experiments/FIGS/data/label_recovery_curve_v1err/analysis/label_recovery_curve_summary.json :: table[rung=R16_Q1,level=162,probe=linear,n_train=max].accuracy :: sha256 1250d1f00537e0bd :: balanced accuracy, mean +- SD over 5 runs
\newcommand{\RecoveryAccOnesevenByFourthree}{0.7680\,\ensuremath{\pm}\,0.0010}  % experiments/FIGS/data/label_recovery_curve_v1err/analysis/label_recovery_curve_summary.json :: table[rung=R16_Q1,level=43,probe=linear,n_train=max].accuracy :: sha256 1250d1f00537e0bd :: balanced accuracy, mean +- SD over 5 runs
\newcommand{\RecoveryAccOnesevenByOneseven}{0.7681\,\ensuremath{\pm}\,0.0003}  % experiments/FIGS/data/label_recovery_curve_v1err/analysis/label_recovery_curve_summary.json :: table[rung=R16_Q1,level=17,probe=linear,n_train=max].accuracy :: sha256 1250d1f00537e0bd :: balanced accuracy, mean +- SD over 5 runs
\newcommand{\RecoveryAccFourByOneeighteight}{0.7526\,\ensuremath{\pm}\,0.0002}  % experiments/FIGS/data/label_recovery_curve_v1err/analysis/label_recovery_curve_summary.json :: table[rung=R3_VIS,level=188,probe=linear,n_train=max].accuracy :: sha256 1250d1f00537e0bd :: balanced accuracy, mean +- SD over 5 runs
\newcommand{\RecoveryAccFourByOnesixtwo}{0.7538\,\ensuremath{\pm}\,0.0017}  % experiments/FIGS/data/label_recovery_curve_v1err/analysis/label_recovery_curve_summary.json :: table[rung=R3_VIS,level=162,probe=linear,n_train=max].accuracy :: sha256 1250d1f00537e0bd :: balanced accuracy, mean +- SD over 5 runs
\newcommand{\RecoveryAccFourByFourthree}{0.7557\,\ensuremath{\pm}\,0.0015}  % experiments/FIGS/data/label_recovery_curve_v1err/analysis/label_recovery_curve_summary.json :: table[rung=R3_VIS,level=43,probe=linear,n_train=max].accuracy :: sha256 1250d1f00537e0bd :: balanced accuracy, mean +- SD over 5 runs
\newcommand{\RecoveryAccFourByOneseven}{0.7586\,\ensuremath{\pm}\,0.0008}  % experiments/FIGS/data/label_recovery_curve_v1err/analysis/label_recovery_curve_summary.json :: table[rung=R3_VIS,level=17,probe=linear,n_train=max].accuracy :: sha256 1250d1f00537e0bd :: balanced accuracy, mean +- SD over 5 runs
\newcommand{\RecoveryAccTwoByOneeighteight}{0.8688\,\ensuremath{\pm}\,0.0005}  % experiments/FIGS/data/label_recovery_curve_v1err/analysis/label_recovery_curve_summary.json :: table[rung=R1_Q1,level=188,probe=linear,n_train=max].accuracy :: sha256 1250d1f00537e0bd :: balanced accuracy, mean +- SD over 5 runs
\newcommand{\RecoveryAccTwoByOnesixtwo}{0.8703\,\ensuremath{\pm}\,0.0008}  % experiments/FIGS/data/label_recovery_curve_v1err/analysis/label_recovery_curve_summary.json :: table[rung=R1_Q1,level=162,probe=linear,n_train=max].accuracy :: sha256 1250d1f00537e0bd :: balanced accuracy, mean +- SD over 5 runs
\newcommand{\RecoveryAccTwoByFourthree}{0.8709\,\ensuremath{\pm}\,0.0011}  % experiments/FIGS/data/label_recovery_curve_v1err/analysis/label_recovery_curve_summary.json :: table[rung=R1_Q1,level=43,probe=linear,n_train=max].accuracy :: sha256 1250d1f00537e0bd :: balanced accuracy, mean +- SD over 5 runs
\newcommand{\RecoveryAccTwoByOneseven}{0.8713\,\ensuremath{\pm}\,0.0002}  % experiments/FIGS/data/label_recovery_curve_v1err/analysis/label_recovery_curve_summary.json :: table[rung=R1_Q1,level=17,probe=linear,n_train=max].accuracy :: sha256 1250d1f00537e0bd :: balanced accuracy, mean +- SD over 5 runs
\newcommand{\RecoveryOneeighteightOnesixtwoMaxAbs}{0.0033}  % experiments/FIGS/data/label_recovery_curve_v1err/analysis/label_recovery_curve_summary.json :: table[rung=*,level=188/162,probe=linear,n_train=max].accuracy :: sha256 1250d1f00537e0bd :: max |difference of the run means| over the tree
\newcommand{\RecoveryPairedOnesevenLevelOnesevenMinusOneeighteight}{+0.0006~[\ensuremath{-}0.0004, +0.0015]}  % experiments/FIGS/data/label_recovery_curve_v1err/analysis/label_recovery_curve_summary.json :: table[rung=R16_Q1,level=17 minus 188,probe=linear,n_train=max].accuracy :: sha256 1250d1f00537e0bd :: paired by run, mean difference [95% Student-t interval]
\newcommand{\RecoveryPairedOnesevenLevelOnesevenMinusOnesixtwo}{+0.00004~[\ensuremath{-}0.00123, +0.00131]}  % experiments/FIGS/data/label_recovery_curve_v1err/analysis/label_recovery_curve_summary.json :: table[rung=R16_Q1,level=17 minus 162,probe=linear,n_train=max].accuracy :: sha256 1250d1f00537e0bd :: paired by run, mean difference [95% Student-t interval]
\newcommand{\RecoveryPairedFourLevelOnesevenMinusOneeighteight}{+0.0060~[+0.0051, +0.0069]}  % experiments/FIGS/data/label_recovery_curve_v1err/analysis/label_recovery_curve_summary.json :: table[rung=R3_VIS,level=17 minus 188,probe=linear,n_train=max].accuracy :: sha256 1250d1f00537e0bd :: paired by run, mean difference [95% Student-t interval]
\newcommand{\RecoveryPairedFourLevelOnesevenMinusOnesixtwo}{+0.0048~[+0.0028, +0.0068]}  % experiments/FIGS/data/label_recovery_curve_v1err/analysis/label_recovery_curve_summary.json :: table[rung=R3_VIS,level=17 minus 162,probe=linear,n_train=max].accuracy :: sha256 1250d1f00537e0bd :: paired by run, mean difference [95% Student-t interval]
\newcommand{\RecoveryPairedTwoLevelOnesevenMinusOneeighteight}{+0.0025~[+0.0019, +0.0032]}  % experiments/FIGS/data/label_recovery_curve_v1err/analysis/label_recovery_curve_summary.json :: table[rung=R1_Q1,level=17 minus 188,probe=linear,n_train=max].accuracy :: sha256 1250d1f00537e0bd :: paired by run, mean difference [95% Student-t interval]
\newcommand{\RecoveryPairedTwoLevelOnesevenMinusOnesixtwo}{+0.00101~[+0.00003, +0.00198]}  % experiments/FIGS/data/label_recovery_curve_v1err/analysis/label_recovery_curve_summary.json :: table[rung=R1_Q1,level=17 minus 162,probe=linear,n_train=max].accuracy :: sha256 1250d1f00537e0bd :: paired by run, mean difference [95% Student-t interval]
\newcommand{\RecoveryMlpGainOneseven}{+0.0099~[+0.0078, +0.0120]}  % experiments/FIGS/data/label_recovery_curve_v1err/analysis/label_recovery_curve_summary.json :: mlp_minus_linear_at_largest[rung=L188,level=17].per_run :: sha256 1250d1f00537e0bd :: nonlinear minus linear probe, 188 native labels from the 17-class models
\newcommand{\RecoveryLastGainOneseven}{0.0125\,\ensuremath{\pm}\,0.0021}  % experiments/FIGS/data/label_recovery_curve_v1err/analysis/label_recovery_curve_summary.json :: data_gain_last_step[rung=L188,level=17] :: sha256 1250d1f00537e0bd :: balanced accuracy gained from 443000 to 1400000 training jets
\newcommand{\RecoveryLastGainOneeighteight}{0.0147\,\ensuremath{\pm}\,0.0014}  % experiments/FIGS/data/label_recovery_curve_v1err/analysis/label_recovery_curve_summary.json :: data_gain_last_step[rung=L188,level=188] :: sha256 1250d1f00537e0bd :: balanced accuracy gained from 443000 to 1400000 training jets
\newcommand{\RecoveryNTrain}{1{,}400{,}000}  % experiments/FIGS/data/label_recovery_curve_v1err/analysis/label_recovery_curve_summary.json :: sizes (max) :: sha256 1250d1f00537e0bd :: jets the linear probe is fit on at the end of the curve
\newcommand{\RecoveryNTrainMin}{14{,}000}  % experiments/FIGS/data/label_recovery_curve_v1err/analysis/label_recovery_curve_summary.json :: sizes (min) :: sha256 1250d1f00537e0bd
\newcommand{\RecoveryNTrainPrev}{443{,}000}  % experiments/FIGS/data/label_recovery_curve_v1err/analysis/label_recovery_curve_summary.json :: sizes (second largest) :: sha256 1250d1f00537e0bd
\newcommand{\RecoveryNTest}{400{,}000}  % experiments/FIGS/data/label_recovery_curve_v1err/l188-s1/label_recovery_curve.json :: n_test :: sha256 65651e4c06c7766f :: jets the probes are scored on, recorded in all 20 per-run files
\newcommand{\RecoveryMlpHidden}{512}  % experiments/FIGS/data/label_recovery_curve_v1err/l188-s1/label_recovery_curve.json :: mlp.hidden :: sha256 65651e4c06c7766f :: hidden units of the nonlinear probe
\newcommand{\RandNDraws}{three}  % experiments/FIGS/data/probe_ladder_randcontrol_mlp2/analysis/c4_random_control.json :: table[*].draw :: sha256 c675e609ff4fe3c9 :: random partitions drawn
\newcommand{\RandNDrawsCap}{Three}  % experiments/FIGS/data/probe_ladder_randcontrol_mlp2/analysis/c4_random_control.json :: table[*].draw :: sha256 c675e609ff4fe3c9 :: random partitions drawn, capitalised to open a sentence
\newcommand{\RandOmaBvcFourprongDrawOneLinear}{\ensuremath{1.05\times10^{-2}}}  % experiments/FIGS/data/probe_ladder_randcontrol_mlp2/analysis/c4_random_control.json :: table[0].control_log1m_auc :: sha256 c675e609ff4fe3c9 :: 1-AUC = exp of it, one run
\newcommand{\RandOmaBvcFourprongDrawTwoLinear}{\ensuremath{1.11\times10^{-2}}}  % experiments/FIGS/data/probe_ladder_randcontrol_mlp2/analysis/c4_random_control.json :: table[4].control_log1m_auc :: sha256 c675e609ff4fe3c9 :: 1-AUC = exp of it, one run
\newcommand{\RandOmaBvcFourprongDrawThreeLinear}{\ensuremath{9.98\times10^{-3}}}  % experiments/FIGS/data/probe_ladder_randcontrol_mlp2/analysis/c4_random_control.json :: table[8].control_log1m_auc :: sha256 c675e609ff4fe3c9 :: 1-AUC = exp of it, one run
\newcommand{\RandOmaMeanBvcFourprongLinear}{\ensuremath{(1.05 \pm 0.06)\times10^{-2}}}  % experiments/FIGS/data/probe_ladder_randcontrol_mlp2/analysis/c4_random_control.json :: table[task=bvc_4prong,probe=linear].control_log1m_auc :: sha256 c675e609ff4fe3c9 :: 1-AUC, mean +- SD over 3 draws
\newcommand{\RandOmaBvcFourprongDrawOneMlp}{\ensuremath{1.05\times10^{-2}}}  % experiments/FIGS/data/probe_ladder_randcontrol_mlp2/analysis/c4_random_control.json :: table[1].control_log1m_auc :: sha256 c675e609ff4fe3c9 :: 1-AUC = exp of it, one run
\newcommand{\RandOmaBvcFourprongDrawTwoMlp}{\ensuremath{1.04\times10^{-2}}}  % experiments/FIGS/data/probe_ladder_randcontrol_mlp2/analysis/c4_random_control.json :: table[5].control_log1m_auc :: sha256 c675e609ff4fe3c9 :: 1-AUC = exp of it, one run
\newcommand{\RandOmaBvcFourprongDrawThreeMlp}{\ensuremath{9.28\times10^{-3}}}  % experiments/FIGS/data/probe_ladder_randcontrol_mlp2/analysis/c4_random_control.json :: table[9].control_log1m_auc :: sha256 c675e609ff4fe3c9 :: 1-AUC = exp of it, one run
\newcommand{\RandOmaMeanBvcFourprongMlp}{\ensuremath{(1.01 \pm 0.07)\times10^{-2}}}  % experiments/FIGS/data/probe_ladder_randcontrol_mlp2/analysis/c4_random_control.json :: table[task=bvc_4prong,probe=mlp].control_log1m_auc :: sha256 c675e609ff4fe3c9 :: 1-AUC, mean +- SD over 3 draws
\newcommand{\RandOmaVisibleContentDrawOneLinear}{\ensuremath{1.25\times10^{-3}}}  % experiments/FIGS/data/probe_ladder_randcontrol_mlp2/analysis/c4_random_control.json :: table[2].control_log1m_auc :: sha256 c675e609ff4fe3c9 :: 1-AUC = exp of it, one run
\newcommand{\RandOmaVisibleContentDrawTwoLinear}{\ensuremath{9.89\times10^{-4}}}  % experiments/FIGS/data/probe_ladder_randcontrol_mlp2/analysis/c4_random_control.json :: table[6].control_log1m_auc :: sha256 c675e609ff4fe3c9 :: 1-AUC = exp of it, one run
\newcommand{\RandOmaVisibleContentDrawThreeLinear}{\ensuremath{1.08\times10^{-3}}}  % experiments/FIGS/data/probe_ladder_randcontrol_mlp2/analysis/c4_random_control.json :: table[10].control_log1m_auc :: sha256 c675e609ff4fe3c9 :: 1-AUC = exp of it, one run
\newcommand{\RandOmaMeanVisibleContentLinear}{\ensuremath{(1.11 \pm 0.13)\times10^{-3}}}  % experiments/FIGS/data/probe_ladder_randcontrol_mlp2/analysis/c4_random_control.json :: table[task=visible_content,probe=linear].control_log1m_auc :: sha256 c675e609ff4fe3c9 :: 1-AUC, mean +- SD over 3 draws
\newcommand{\RandOmaVisibleContentDrawOneMlp}{\ensuremath{1.21\times10^{-3}}}  % experiments/FIGS/data/probe_ladder_randcontrol_mlp2/analysis/c4_random_control.json :: table[3].control_log1m_auc :: sha256 c675e609ff4fe3c9 :: 1-AUC = exp of it, one run
\newcommand{\RandOmaVisibleContentDrawTwoMlp}{\ensuremath{1.07\times10^{-3}}}  % experiments/FIGS/data/probe_ladder_randcontrol_mlp2/analysis/c4_random_control.json :: table[7].control_log1m_auc :: sha256 c675e609ff4fe3c9 :: 1-AUC = exp of it, one run
\newcommand{\RandOmaVisibleContentDrawThreeMlp}{\ensuremath{1.13\times10^{-3}}}  % experiments/FIGS/data/probe_ladder_randcontrol_mlp2/analysis/c4_random_control.json :: table[11].control_log1m_auc :: sha256 c675e609ff4fe3c9 :: 1-AUC = exp of it, one run
\newcommand{\RandOmaMeanVisibleContentMlp}{\ensuremath{(1.14 \pm 0.07)\times10^{-3}}}  % experiments/FIGS/data/probe_ladder_randcontrol_mlp2/analysis/c4_random_control.json :: table[task=visible_content,probe=mlp].control_log1m_auc :: sha256 c675e609ff4fe3c9 :: 1-AUC, mean +- SD over 3 draws
\newcommand{\RandCostFactorBvcFourprongLinear}{1.49}  % experiments/FIGS/data/probe_ladder_randcontrol_mlp2/analysis/c4_random_control.json :: grouping_cost_post_hoc.tasks.bvc_4prong.linear.factor_control :: sha256 c675e609ff4fe3c9 :: POST HOC: random grouping vs finer models, in 1-AUC
\newcommand{\RandSemCostFactorBvcFourprongLinear}{4.42}  % experiments/FIGS/data/probe_ladder_randcontrol_mlp2/analysis/c4_random_control.json :: grouping_cost_post_hoc.tasks.bvc_4prong.linear.factor_semantic :: sha256 c675e609ff4fe3c9 :: POST HOC: semantic 17-class vs finer models
\newcommand{\RandCostFactorBvcFourprongMlp}{1.41}  % experiments/FIGS/data/probe_ladder_randcontrol_mlp2/analysis/c4_random_control.json :: grouping_cost_post_hoc.tasks.bvc_4prong.mlp.factor_control :: sha256 c675e609ff4fe3c9 :: POST HOC: random grouping vs finer models, in 1-AUC
\newcommand{\RandSemCostFactorBvcFourprongMlp}{3.26}  % experiments/FIGS/data/probe_ladder_randcontrol_mlp2/analysis/c4_random_control.json :: grouping_cost_post_hoc.tasks.bvc_4prong.mlp.factor_semantic :: sha256 c675e609ff4fe3c9 :: POST HOC: semantic 17-class vs finer models
\newcommand{\RandCostFactorVisibleContentLinear}{1.35}  % experiments/FIGS/data/probe_ladder_randcontrol_mlp2/analysis/c4_random_control.json :: grouping_cost_post_hoc.tasks.visible_content.linear.factor_control :: sha256 c675e609ff4fe3c9 :: POST HOC: random grouping vs finer models, in 1-AUC
\newcommand{\RandSemCostFactorVisibleContentLinear}{2.82}  % experiments/FIGS/data/probe_ladder_randcontrol_mlp2/analysis/c4_random_control.json :: grouping_cost_post_hoc.tasks.visible_content.linear.factor_semantic :: sha256 c675e609ff4fe3c9 :: POST HOC: semantic 17-class vs finer models
\newcommand{\RandCostFactorVisibleContentMlp}{1.33}  % experiments/FIGS/data/probe_ladder_randcontrol_mlp2/analysis/c4_random_control.json :: grouping_cost_post_hoc.tasks.visible_content.mlp.factor_control :: sha256 c675e609ff4fe3c9 :: POST HOC: random grouping vs finer models, in 1-AUC
\newcommand{\RandSemCostFactorVisibleContentMlp}{2.37}  % experiments/FIGS/data/probe_ladder_randcontrol_mlp2/analysis/c4_random_control.json :: grouping_cost_post_hoc.tasks.visible_content.mlp.factor_semantic :: sha256 c675e609ff4fe3c9 :: POST HOC: semantic 17-class vs finer models
\newcommand{\MassResNTest}{398{,}625}  % experiments/FIGS/data/mass_resolution/s1.json :: centering_detail.split[2] :: sha256 88a57bbad270d51c :: test jets sigma_eff is computed on, the same in every per-seed file and arm
\newcommand{\MassResNClasses}{184}  % experiments/FIGS/data/mass_resolution/analysis_holm/s7_mass_resolution.json :: provenance.n_classes_used :: sha256 e03cde45ece9f385 :: native classes with enough training jets to centre
\newcommand{\MassResSigmaEffRidgeOneeighteight}{0.3393\,\ensuremath{\pm}\,0.0019}  % experiments/FIGS/data/mass_resolution/analysis_holm/s7_mass_resolution.json :: table[cell=188,probe=ridge].sigma_eff :: sha256 e03cde45ece9f385 :: mean +- SD over 5 seeds
\newcommand{\MassResSigmaEffRidgeOnesixtwo}{0.3302\,\ensuremath{\pm}\,0.0029}  % experiments/FIGS/data/mass_resolution/analysis_holm/s7_mass_resolution.json :: table[cell=162,probe=ridge].sigma_eff :: sha256 e03cde45ece9f385 :: mean +- SD over 5 seeds
\newcommand{\MassResSigmaEffRidgeFourthree}{0.3228\,\ensuremath{\pm}\,0.0019}  % experiments/FIGS/data/mass_resolution/analysis_holm/s7_mass_resolution.json :: table[cell=43,probe=ridge].sigma_eff :: sha256 e03cde45ece9f385 :: mean +- SD over 5 seeds
\newcommand{\MassResSigmaEffRidgeOneseven}{0.300\,\ensuremath{\pm}\,0.010}  % experiments/FIGS/data/mass_resolution/analysis_holm/s7_mass_resolution.json :: table[cell=17,probe=ridge].sigma_eff :: sha256 e03cde45ece9f385 :: mean +- SD over 5 seeds
\newcommand{\MassResSigmaEffRidgeOnesixtwoMass}{0.2238\,\ensuremath{\pm}\,0.0019}  % experiments/FIGS/data/mass_resolution/analysis_holm/s7_mass_resolution.json :: table[cell=162+mass,probe=ridge].sigma_eff :: sha256 e03cde45ece9f385 :: mean +- SD over 5 seeds
\newcommand{\MassResSigmaEffRidgeOnesevenMass}{0.2170\,\ensuremath{\pm}\,0.0034}  % experiments/FIGS/data/mass_resolution/analysis_holm/s7_mass_resolution.json :: table[cell=17+mass,probe=ridge].sigma_eff :: sha256 e03cde45ece9f385 :: mean +- SD over 5 seeds
\newcommand{\MassResGainRidgeOnesixtwo}{\ensuremath{-}0.1063}  % experiments/FIGS/data/mass_resolution/analysis_holm/s7_mass_resolution.json :: table[cell=162+mass/162,probe=ridge].sigma_eff :: sha256 e03cde45ece9f385 :: sigma_eff with the mass output minus without, difference of the seed means; negative is better
\newcommand{\MassResGainRidgeOneseven}{\ensuremath{-}0.0831}  % experiments/FIGS/data/mass_resolution/analysis_holm/s7_mass_resolution.json :: table[cell=17+mass/17,probe=ridge].sigma_eff :: sha256 e03cde45ece9f385 :: sigma_eff with the mass output minus without, difference of the seed means; negative is better
\newcommand{\MassResDidRidge}{\ensuremath{-}0.0233}  % experiments/FIGS/data/mass_resolution/analysis_holm/s7_mass_resolution.json :: table[cell=162+mass/162/17+mass/17,probe=ridge].sigma_eff :: sha256 e03cde45ece9f385 :: gain at 162 classes minus gain at 17
\newcommand{\MassResSigmaEffMlpOneeighteight}{0.2463\,\ensuremath{\pm}\,0.0022}  % experiments/FIGS/data/mass_resolution/analysis_holm/s7_mass_resolution.json :: table[cell=188,probe=mlp].sigma_eff :: sha256 e03cde45ece9f385 :: mean +- SD over 5 seeds
\newcommand{\MassResSigmaEffMlpOnesixtwo}{0.2438\,\ensuremath{\pm}\,0.0018}  % experiments/FIGS/data/mass_resolution/analysis_holm/s7_mass_resolution.json :: table[cell=162,probe=mlp].sigma_eff :: sha256 e03cde45ece9f385 :: mean +- SD over 5 seeds
\newcommand{\MassResSigmaEffMlpFourthree}{0.2407\,\ensuremath{\pm}\,0.0009}  % experiments/FIGS/data/mass_resolution/analysis_holm/s7_mass_resolution.json :: table[cell=43,probe=mlp].sigma_eff :: sha256 e03cde45ece9f385 :: mean +- SD over 5 seeds
\newcommand{\MassResSigmaEffMlpOneseven}{0.2358\,\ensuremath{\pm}\,0.0025}  % experiments/FIGS/data/mass_resolution/analysis_holm/s7_mass_resolution.json :: table[cell=17,probe=mlp].sigma_eff :: sha256 e03cde45ece9f385 :: mean +- SD over 5 seeds
\newcommand{\MassResSigmaEffMlpOnesixtwoMass}{0.1943\,\ensuremath{\pm}\,0.0001}  % experiments/FIGS/data/mass_resolution/analysis_holm/s7_mass_resolution.json :: table[cell=162+mass,probe=mlp].sigma_eff :: sha256 e03cde45ece9f385 :: mean +- SD over 5 seeds
\newcommand{\MassResSigmaEffMlpOnesevenMass}{0.1949\,\ensuremath{\pm}\,0.0006}  % experiments/FIGS/data/mass_resolution/analysis_holm/s7_mass_resolution.json :: table[cell=17+mass,probe=mlp].sigma_eff :: sha256 e03cde45ece9f385 :: mean +- SD over 5 seeds
\newcommand{\MassResGainMlpOnesixtwo}{\ensuremath{-}0.0496}  % experiments/FIGS/data/mass_resolution/analysis_holm/s7_mass_resolution.json :: table[cell=162+mass/162,probe=mlp].sigma_eff :: sha256 e03cde45ece9f385 :: sigma_eff with the mass output minus without, difference of the seed means; negative is better
\newcommand{\MassResGainMlpOneseven}{\ensuremath{-}0.0409}  % experiments/FIGS/data/mass_resolution/analysis_holm/s7_mass_resolution.json :: table[cell=17+mass/17,probe=mlp].sigma_eff :: sha256 e03cde45ece9f385 :: sigma_eff with the mass output minus without, difference of the seed means; negative is better
\newcommand{\MassResDidMlp}{\ensuremath{-}0.0087}  % experiments/FIGS/data/mass_resolution/analysis_holm/s7_mass_resolution.json :: table[cell=162+mass/162/17+mass/17,probe=mlp].sigma_eff :: sha256 e03cde45ece9f385 :: gain at 162 classes minus gain at 17
\newcommand{\MassResTargetSigmaEff}{0.2986}  % experiments/FIGS/data/mass_resolution/analysis_holm/s7_mass_resolution.json :: table[*].target_sigma_eff :: sha256 e03cde45ece9f385 :: sigma_eff of the class-centred target itself (a model that predicts the class mean)
\newcommand{\MassResRidgeRtwoMin}{20\%}  % experiments/FIGS/data/mass_resolution/s1.json :: arms.{l188,l162,r42q1,r16q1}-s*.ridge.val_r2 :: sha256 88a57bbad270d51c :: validation R^2 of the ridge probe, smallest of 20 models without the mass output, all 5 per-seed files
\newcommand{\MassResRidgeRtwoMax}{30\%}  % experiments/FIGS/data/mass_resolution/s1.json :: arms.{l188,l162,r42q1,r16q1}-s*.ridge.val_r2 :: sha256 88a57bbad270d51c :: validation R^2 of the ridge probe, largest of 20 models without the mass output, all 5 per-seed files
\newcommand{\MassResSigmaEffOnesixtwo}{0.3302}  % experiments/FIGS/data/mass_resolution/analysis_holm/s7_mass_resolution.json :: table[cell=162,probe=ridge].sigma_eff :: sha256 e03cde45ece9f385 :: mean over 5 seeds
\newcommand{\MassResSdOnesixtwo}{0.626}  % experiments/FIGS/data/mass_resolution/analysis_holm/s7_mass_resolution.json :: table[cell=162,probe=ridge].sd :: sha256 e03cde45ece9f385 :: mean over 5 seeds
\newcommand{\MassResTailFractionOnesixtwo}{0.062}  % experiments/FIGS/data/mass_resolution/analysis_holm/s7_mass_resolution.json :: table[cell=162,probe=ridge].tail_fraction :: sha256 e03cde45ece9f385 :: mean over 5 seeds
\newcommand{\MassResSigmaEffOnesixtwoMass}{0.2238}  % experiments/FIGS/data/mass_resolution/analysis_holm/s7_mass_resolution.json :: table[cell=162+mass,probe=ridge].sigma_eff :: sha256 e03cde45ece9f385 :: mean over 5 seeds
\newcommand{\MassResSdOnesixtwoMass}{0.494}  % experiments/FIGS/data/mass_resolution/analysis_holm/s7_mass_resolution.json :: table[cell=162+mass,probe=ridge].sd :: sha256 e03cde45ece9f385 :: mean over 5 seeds
\newcommand{\MassResTailFractionOnesixtwoMass}{0.034}  % experiments/FIGS/data/mass_resolution/analysis_holm/s7_mass_resolution.json :: table[cell=162+mass,probe=ridge].tail_fraction :: sha256 e03cde45ece9f385 :: mean over 5 seeds
\newcommand{\MassResSigmaEffOneseven}{0.3001}  % experiments/FIGS/data/mass_resolution/analysis_holm/s7_mass_resolution.json :: table[cell=17,probe=ridge].sigma_eff :: sha256 e03cde45ece9f385 :: mean over 5 seeds
\newcommand{\MassResSdOneseven}{0.601}  % experiments/FIGS/data/mass_resolution/analysis_holm/s7_mass_resolution.json :: table[cell=17,probe=ridge].sd :: sha256 e03cde45ece9f385 :: mean over 5 seeds
\newcommand{\MassResTailFractionOneseven}{0.056}  % experiments/FIGS/data/mass_resolution/analysis_holm/s7_mass_resolution.json :: table[cell=17,probe=ridge].tail_fraction :: sha256 e03cde45ece9f385 :: mean over 5 seeds
\newcommand{\MassResSigmaEffOnesevenMass}{0.2170}  % experiments/FIGS/data/mass_resolution/analysis_holm/s7_mass_resolution.json :: table[cell=17+mass,probe=ridge].sigma_eff :: sha256 e03cde45ece9f385 :: mean over 5 seeds
\newcommand{\MassResSdOnesevenMass}{0.490}  % experiments/FIGS/data/mass_resolution/analysis_holm/s7_mass_resolution.json :: table[cell=17+mass,probe=ridge].sd :: sha256 e03cde45ece9f385 :: mean over 5 seeds
\newcommand{\MassResTailFractionOnesevenMass}{0.033}  % experiments/FIGS/data/mass_resolution/analysis_holm/s7_mass_resolution.json :: table[cell=17+mass,probe=ridge].tail_fraction :: sha256 e03cde45ece9f385 :: mean over 5 seeds
\newcommand{\MassResSigmaEffOneeighteight}{0.3393}  % experiments/FIGS/data/mass_resolution/analysis_holm/s7_mass_resolution.json :: table[cell=188,probe=ridge].sigma_eff :: sha256 e03cde45ece9f385 :: mean over 5 seeds
\newcommand{\MassResSdOneeighteight}{0.631}  % experiments/FIGS/data/mass_resolution/analysis_holm/s7_mass_resolution.json :: table[cell=188,probe=ridge].sd :: sha256 e03cde45ece9f385 :: mean over 5 seeds
\newcommand{\MassResTailFractionOneeighteight}{0.064}  % experiments/FIGS/data/mass_resolution/analysis_holm/s7_mass_resolution.json :: table[cell=188,probe=ridge].tail_fraction :: sha256 e03cde45ece9f385 :: mean over 5 seeds
\newcommand{\MassResSigmaEffFourthree}{0.3228}  % experiments/FIGS/data/mass_resolution/analysis_holm/s7_mass_resolution.json :: table[cell=43,probe=ridge].sigma_eff :: sha256 e03cde45ece9f385 :: mean over 5 seeds
\newcommand{\MassResSdFourthree}{0.618}  % experiments/FIGS/data/mass_resolution/analysis_holm/s7_mass_resolution.json :: table[cell=43,probe=ridge].sd :: sha256 e03cde45ece9f385 :: mean over 5 seeds
\newcommand{\MassResTailFractionFourthree}{0.060}  % experiments/FIGS/data/mass_resolution/analysis_holm/s7_mass_resolution.json :: table[cell=43,probe=ridge].tail_fraction :: sha256 e03cde45ece9f385 :: mean over 5 seeds
\newcommand{\AojEffPercent}{1\%}  % experiments/FIGS/data/aoj_full_v1/fit_v6/results.json :: eff :: sha256 2416fc21f93fcfa6 :: the score cut's pass fraction outside the mass windows
\newcommand{\AojPtMin}{500}  % experiments/AOJ/peak_fit.py :: PT_RANGE[0] :: sha256 cad06229a74532be :: GeV
\newcommand{\AojPtMax}{2{,}500}  % experiments/AOJ/peak_fit.py :: PT_RANGE[1] :: sha256 cad06229a74532be :: GeV
\newcommand{\AojRhoMin}{\ensuremath{-}5.5}  % experiments/AOJ/peak_fit.py :: RHO_RANGE[0] :: sha256 cad06229a74532be
\newcommand{\AojRhoMax}{\ensuremath{-}2.0}  % experiments/AOJ/peak_fit.py :: RHO_RANGE[1] :: sha256 cad06229a74532be
\newcommand{\AojWindowWLo}{65}  % experiments/AOJ/peak_fit.py :: PEAKS[W].window[0] :: sha256 cad06229a74532be :: GeV
\newcommand{\AojWindowWHi}{105}  % experiments/AOJ/peak_fit.py :: PEAKS[W].window[1] :: sha256 cad06229a74532be :: GeV
\newcommand{\AojWindowTopLo}{140}  % experiments/AOJ/peak_fit.py :: PEAKS[top].window[0] :: sha256 cad06229a74532be :: GeV
\newcommand{\AojWindowTopHi}{220}  % experiments/AOJ/peak_fit.py :: PEAKS[top].window[1] :: sha256 cad06229a74532be :: GeV
\newcommand{\AojFitLo}{105}  % experiments/AOJ/peak_fit.py :: PEAKS[top].fit_range[0] :: sha256 cad06229a74532be :: GeV
\newcommand{\AojFitHi}{300}  % experiments/AOJ/peak_fit.py :: PEAKS[top].fit_range[1] :: sha256 cad06229a74532be :: GeV
\newcommand{\AojMassBin}{5}  % experiments/AOJ/peak_fit.py :: MASS_BIN :: sha256 cad06229a74532be :: GeV
\newcommand{\AojMapOrderRho}{3}  % experiments/AOJ/peak_fit.py :: MAP_ORDER[0] :: sha256 cad06229a74532be :: polynomial order in rho
\newcommand{\AojMapOrderPt}{2}  % experiments/AOJ/peak_fit.py :: MAP_ORDER[1] :: sha256 cad06229a74532be :: polynomial order in ln pT
\newcommand{\AojNRunsWord}{five}  % experiments/FIGS/data/aoj_full_v1/analysis_v6/aoj_top.json :: per_label_set.label_sets.*.signal_yield.n :: sha256 5dfef902a5fe75ba :: pretraining runs per vocabulary
\newcommand{\AojNJetsRhoWindow}{9{,}441{,}209}  % experiments/FIGS/data/aoj_full_v1/fit_v6/results.json :: n_jets :: sha256 2416fc21f93fcfa6 :: staged jets inside the rho window, all files; not the jets in the mass fit
\newcommand{\AojYieldOneeighteight}{2{,}400\,\ensuremath{\pm}\,500}  % experiments/FIGS/data/aoj_full_v1/analysis_v6/aoj_top.json :: per_label_set.label_sets.188.signal_yield :: sha256 5dfef902a5fe75ba :: mean +- SD over 5 runs, shared shape
\newcommand{\AojStatErrOneeighteight}{211}  % experiments/FIGS/data/aoj_full_v1/analysis_v6/aoj_top.json :: per_label_set.label_sets.188.median_stat_err :: sha256 5dfef902a5fe75ba :: median per-fit statistical error
\newcommand{\AojYieldOwnShapeOneeighteight}{2{,}400\,\ensuremath{\pm}\,600}  % experiments/FIGS/data/aoj_full_v1/analysis_v6/aoj_top.json :: per_label_set.label_sets.188.by_shape.own_shape :: sha256 5dfef902a5fe75ba :: every fit with its own floated shape
\newcommand{\AojSpreadOverStatOneeighteight}{2.4}  % experiments/FIGS/data/aoj_full_v1/analysis_v6/aoj_top.json :: per_label_set.label_sets.188.signal_yield.sd / median_stat_err :: sha256 5dfef902a5fe75ba :: spread between runs over the median per-fit statistical error
\newcommand{\AojYieldOnesixtwo}{2{,}400\,\ensuremath{\pm}\,900}  % experiments/FIGS/data/aoj_full_v1/analysis_v6/aoj_top.json :: per_label_set.label_sets.162.signal_yield :: sha256 5dfef902a5fe75ba :: mean +- SD over 5 runs, shared shape
\newcommand{\AojStatErrOnesixtwo}{199}  % experiments/FIGS/data/aoj_full_v1/analysis_v6/aoj_top.json :: per_label_set.label_sets.162.median_stat_err :: sha256 5dfef902a5fe75ba :: median per-fit statistical error
\newcommand{\AojYieldOwnShapeOnesixtwo}{2{,}200\,\ensuremath{\pm}\,800}  % experiments/FIGS/data/aoj_full_v1/analysis_v6/aoj_top.json :: per_label_set.label_sets.162.by_shape.own_shape :: sha256 5dfef902a5fe75ba :: every fit with its own floated shape
\newcommand{\AojSpreadOverStatOnesixtwo}{4.5}  % experiments/FIGS/data/aoj_full_v1/analysis_v6/aoj_top.json :: per_label_set.label_sets.162.signal_yield.sd / median_stat_err :: sha256 5dfef902a5fe75ba :: spread between runs over the median per-fit statistical error
\newcommand{\AojYieldFourthree}{2{,}400\,\ensuremath{\pm}\,800}  % experiments/FIGS/data/aoj_full_v1/analysis_v6/aoj_top.json :: per_label_set.label_sets.43.signal_yield :: sha256 5dfef902a5fe75ba :: mean +- SD over 5 runs, shared shape
\newcommand{\AojStatErrFourthree}{184}  % experiments/FIGS/data/aoj_full_v1/analysis_v6/aoj_top.json :: per_label_set.label_sets.43.median_stat_err :: sha256 5dfef902a5fe75ba :: median per-fit statistical error
\newcommand{\AojYieldOwnShapeFourthree}{2{,}300\,\ensuremath{\pm}\,800}  % experiments/FIGS/data/aoj_full_v1/analysis_v6/aoj_top.json :: per_label_set.label_sets.43.by_shape.own_shape :: sha256 5dfef902a5fe75ba :: every fit with its own floated shape
\newcommand{\AojSpreadOverStatFourthree}{4.2}  % experiments/FIGS/data/aoj_full_v1/analysis_v6/aoj_top.json :: per_label_set.label_sets.43.signal_yield.sd / median_stat_err :: sha256 5dfef902a5fe75ba :: spread between runs over the median per-fit statistical error
\newcommand{\AojYieldOneseven}{2{,}350\,\ensuremath{\pm}\,240}  % experiments/FIGS/data/aoj_full_v1/analysis_v6/aoj_top.json :: per_label_set.label_sets.17.signal_yield :: sha256 5dfef902a5fe75ba :: mean +- SD over 5 runs, shared shape
\newcommand{\AojStatErrOneseven}{201}  % experiments/FIGS/data/aoj_full_v1/analysis_v6/aoj_top.json :: per_label_set.label_sets.17.median_stat_err :: sha256 5dfef902a5fe75ba :: median per-fit statistical error
\newcommand{\AojYieldOwnShapeOneseven}{2{,}480\,\ensuremath{\pm}\,290}  % experiments/FIGS/data/aoj_full_v1/analysis_v6/aoj_top.json :: per_label_set.label_sets.17.by_shape.own_shape :: sha256 5dfef902a5fe75ba :: every fit with its own floated shape
\newcommand{\AojSpreadOverStatOneseven}{1.2}  % experiments/FIGS/data/aoj_full_v1/analysis_v6/aoj_top.json :: per_label_set.label_sets.17.signal_yield.sd / median_stat_err :: sha256 5dfef902a5fe75ba :: spread between runs over the median per-fit statistical error
\newcommand{\AojYieldOnesixtwoMass}{2{,}420\,\ensuremath{\pm}\,240}  % experiments/FIGS/data/aoj_full_v1/analysis_v6/aoj_top.json :: per_label_set.label_sets.162+mass.signal_yield :: sha256 5dfef902a5fe75ba :: mean +- SD over 5 runs, shared shape
\newcommand{\AojStatErrOnesixtwoMass}{179}  % experiments/FIGS/data/aoj_full_v1/analysis_v6/aoj_top.json :: per_label_set.label_sets.162+mass.median_stat_err :: sha256 5dfef902a5fe75ba :: median per-fit statistical error
\newcommand{\AojYieldOwnShapeOnesixtwoMass}{2{,}350\,\ensuremath{\pm}\,320}  % experiments/FIGS/data/aoj_full_v1/analysis_v6/aoj_top.json :: per_label_set.label_sets.162+mass.by_shape.own_shape :: sha256 5dfef902a5fe75ba :: every fit with its own floated shape
\newcommand{\AojSpreadOverStatOnesixtwoMass}{1.3}  % experiments/FIGS/data/aoj_full_v1/analysis_v6/aoj_top.json :: per_label_set.label_sets.162+mass.signal_yield.sd / median_stat_err :: sha256 5dfef902a5fe75ba :: spread between runs over the median per-fit statistical error
\newcommand{\AojYieldOnesevenMass}{2{,}500\,\ensuremath{\pm}\,600}  % experiments/FIGS/data/aoj_full_v1/analysis_v6/aoj_top.json :: per_label_set.label_sets.17+mass.signal_yield :: sha256 5dfef902a5fe75ba :: mean +- SD over 5 runs, shared shape
\newcommand{\AojStatErrOnesevenMass}{200}  % experiments/FIGS/data/aoj_full_v1/analysis_v6/aoj_top.json :: per_label_set.label_sets.17+mass.median_stat_err :: sha256 5dfef902a5fe75ba :: median per-fit statistical error
\newcommand{\AojYieldOwnShapeOnesevenMass}{2{,}900\,\ensuremath{\pm}\,500}  % experiments/FIGS/data/aoj_full_v1/analysis_v6/aoj_top.json :: per_label_set.label_sets.17+mass.by_shape.own_shape :: sha256 5dfef902a5fe75ba :: every fit with its own floated shape
\newcommand{\AojSpreadOverStatOnesevenMass}{2.9}  % experiments/FIGS/data/aoj_full_v1/analysis_v6/aoj_top.json :: per_label_set.label_sets.17+mass.signal_yield.sd / median_stat_err :: sha256 5dfef902a5fe75ba :: spread between runs over the median per-fit statistical error
\newcommand{\AojOwnShapeShiftMin}{\ensuremath{-}7\%}  % experiments/FIGS/data/aoj_full_v1/analysis_v6/aoj_top.json :: per_label_set.label_sets.*.by_shape.own_shape.mean / signal_yield.mean - 1 :: sha256 5dfef902a5fe75ba :: the vocabulary means with each model's own shape, relative to the shared shape
\newcommand{\AojOwnShapeShiftMax}{+13\%}  % experiments/FIGS/data/aoj_full_v1/analysis_v6/aoj_top.json :: per_label_set.label_sets.*.by_shape.own_shape.mean / signal_yield.mean - 1 :: sha256 5dfef902a5fe75ba :: the vocabulary means with each model's own shape, relative to the shared shape
\newcommand{\AojPooledMean}{182.8}  % experiments/FIGS/data/aoj_full_v1/analysis_v6/aoj_top.json :: per_label_set.shapes.pooled :: sha256 5dfef902a5fe75ba :: GeV
\newcommand{\AojPooledWidth}{11.7}  % experiments/FIGS/data/aoj_full_v1/analysis_v6/aoj_top.json :: per_label_set.shapes.pooled :: sha256 5dfef902a5fe75ba :: GeV
\newcommand{\AojPoolSize}{30}  % experiments/FIGS/data/aoj_full_v1/analysis_v6/aoj_top.json :: per_label_set.shapes.pool (count) :: sha256 5dfef902a5fe75ba :: pretrained models the shared shape is fitted over
\newcommand{\AojNModels}{30}  % experiments/FIGS/data/aoj_full_v1/fit_v6/results.json :: models (pretrained) :: sha256 2416fc21f93fcfa6 :: fits
\newcommand{\AojOwnMeanMin}{181.6}  % experiments/FIGS/data/aoj_full_v1/fit_v6/results.json :: models.*.top.own_shape_fit.mean (min) :: sha256 2416fc21f93fcfa6 :: GeV
\newcommand{\AojOwnMeanMax}{184.0}  % experiments/FIGS/data/aoj_full_v1/fit_v6/results.json :: models.*.top.own_shape_fit.mean (max) :: sha256 2416fc21f93fcfa6 :: GeV
\newcommand{\AojOwnWidthMin}{10.0}  % experiments/FIGS/data/aoj_full_v1/fit_v6/results.json :: models.*.top.own_shape_fit.width (min) :: sha256 2416fc21f93fcfa6 :: GeV
\newcommand{\AojOwnWidthMax}{13.8}  % experiments/FIGS/data/aoj_full_v1/fit_v6/results.json :: models.*.top.own_shape_fit.width (max) :: sha256 2416fc21f93fcfa6 :: GeV
\newcommand{\AojShapeStepMean}{0.7}  % experiments/FIGS/data/aoj_full_v1/fit_v6/results.json :: pooled_shape.systematic_step.mean :: sha256 2416fc21f93fcfa6 :: GeV: the SD of the models' own floated values, one shift for every model together
\newcommand{\AojShapeSystMeanMax}{1\%}  % experiments/FIGS/data/aoj_full_v1/fit_v6/results.json :: models.*.top.shape_systematic.pooled_mean_(down|up) / signal_yield - 1 (largest |.|) :: sha256 2416fc21f93fcfa6 :: the largest yield change of any model under that shift
\newcommand{\AojShapeStepWidth}{1.0}  % experiments/FIGS/data/aoj_full_v1/fit_v6/results.json :: pooled_shape.systematic_step.width :: sha256 2416fc21f93fcfa6 :: GeV: the SD of the models' own floated values, one shift for every model together
\newcommand{\AojShapeSystWidthMax}{12\%}  % experiments/FIGS/data/aoj_full_v1/fit_v6/results.json :: models.*.top.shape_systematic.pooled_width_(down|up) / signal_yield - 1 (largest |.|) :: sha256 2416fc21f93fcfa6 :: the largest yield change of any model under that shift
\newcommand{\AojLeakEpsHi}{0.6}  % experiments/FIGS/data/aoj_full_v1/fit_v6/results.json :: models.*.top.leak_systematic.eps_ref (max) :: sha256 2416fc21f93fcfa6 :: the reference's assumed efficiency for data tops, upper variation
\newcommand{\AojLeakEpsLo}{0.4}  % experiments/FIGS/data/aoj_full_v1/fit_v6/results.json :: models.*.top.leak_systematic.eps_ref (min) :: sha256 2416fc21f93fcfa6 :: the same, lower variation
\newcommand{\AojLeakShiftMin}{+97}  % experiments/FIGS/data/aoj_full_v1/fit_v6/results.json :: models.*.top.leak_systematic.shift['0.6'] (min) :: sha256 2416fc21f93fcfa6 :: tops added to a yield
\newcommand{\AojLeakShiftMax}{+271}  % experiments/FIGS/data/aoj_full_v1/fit_v6/results.json :: models.*.top.leak_systematic.shift['0.4'] (max) :: sha256 2416fc21f93fcfa6 :: tops added to a yield
\newcommand{\AojFitPMin}{0.010}  % experiments/FIGS/data/aoj_full_v1/fit_v6/results.json :: models.*.top.asymptotic_p (min) :: sha256 2416fc21f93fcfa6 :: goodness of fit of the working-point fit, pretrained models
\newcommand{\AojFitPMax}{0.974}  % experiments/FIGS/data/aoj_full_v1/fit_v6/results.json :: models.*.top.asymptotic_p (max) :: sha256 2416fc21f93fcfa6
\newcommand{\AojWeakestYield}{885}  % experiments/FIGS/data/aoj_full_v1/fit_v6/results.json :: models.l162-s2.top.signal_yield :: sha256 2416fc21f93fcfa6 :: the smallest yield of any pretrained model
\newcommand{\AojWeakestYieldErr}{211}  % experiments/FIGS/data/aoj_full_v1/fit_v6/results.json :: models.l162-s2.top.signal_yield_err :: sha256 2416fc21f93fcfa6
\newcommand{\AojWeakestFitP}{0.94}  % experiments/FIGS/data/aoj_full_v1/fit_v6/results.json :: models.l162-s2.top.asymptotic_p :: sha256 2416fc21f93fcfa6
\newcommand{\AojSidebandZMin}{\ensuremath{-}0.43}  % experiments/FIGS/data/aoj_full_v1/fit_v6/histograms.npz :: experiments/AOJ/sideband_residuals.py: z (min over the pretrained models) :: sha256 2bd1e3d5b32d4a12 :: passing jets minus the fit, summed from the fit's lower edge to the top window
\newcommand{\AojSidebandZMax}{+0.58}  % experiments/FIGS/data/aoj_full_v1/fit_v6/histograms.npz :: experiments/AOJ/sideband_residuals.py: z (max over the pretrained models) :: sha256 2bd1e3d5b32d4a12
\newcommand{\AojWeakestSidebandZ}{0.00}  % experiments/FIGS/data/aoj_full_v1/fit_v6/histograms.npz :: experiments/AOJ/sideband_residuals.py: z of l162-s2 :: sha256 2bd1e3d5b32d4a12 :: the smallest-yield model
\newcommand{\AojSidebandMaxPull}{2.7}  % experiments/FIGS/data/aoj_full_v1/fit_v6/histograms.npz :: experiments/AOJ/sideband_residuals.py: max_abs_pull (max over the pretrained models) :: sha256 2bd1e3d5b32d4a12 :: largest single-bin pull below the top window, any pretrained model
\newcommand{\AojSharedBinLo}{125}  % experiments/FIGS/data/aoj_full_v1/fit_v6/histograms.npz :: the bin holding the most models' largest pull :: sha256 2bd1e3d5b32d4a12 :: GeV
\newcommand{\AojSharedBinHi}{130}  % experiments/FIGS/data/aoj_full_v1/fit_v6/histograms.npz :: its upper edge :: sha256 2bd1e3d5b32d4a12 :: GeV
\newcommand{\AojSharedBinNModels}{18}  % experiments/FIGS/data/aoj_full_v1/fit_v6/histograms.npz :: models whose largest pull is in that bin :: sha256 2bd1e3d5b32d4a12
\newcommand{\AojSharedBinMeanPull}{+1.4}  % experiments/FIGS/data/aoj_full_v1/fit_v6/histograms.npz :: experiments/AOJ/sideband_residuals.py mean_pulls: that bin :: sha256 2bd1e3d5b32d4a12 :: mean pull over the pretrained models
\newcommand{\AojSharedBinRank}{third}  % experiments/FIGS/data/aoj_full_v1/fit_v6/histograms.npz :: experiments/AOJ/sideband_residuals.py mean_pulls: rank of that bin's |mean pull| :: sha256 2bd1e3d5b32d4a12
\newcommand{\AojFitNBins}{39}  % experiments/FIGS/data/aoj_full_v1/fit_v6/histograms.npz :: mass bins of the fit range :: sha256 2bd1e3d5b32d4a12
\newcommand{\AojPairedYieldOnesixtwoOverOneeighteight}{0.94~[0.45, 1.95]}  % experiments/FIGS/data/aoj_full_v1/analysis_v6/aoj_top.json :: per_label_set.label_sets.{162,188}.signal_yields :: sha256 5dfef902a5fe75ba :: yield ratio paired by run, geometric mean [95% Student-t interval]
\newcommand{\AojPairedYieldFourthreeOverOneeighteight}{0.97~[0.81, 1.16]}  % experiments/FIGS/data/aoj_full_v1/analysis_v6/aoj_top.json :: per_label_set.label_sets.{43,188}.signal_yields :: sha256 5dfef902a5fe75ba :: yield ratio paired by run, geometric mean [95% Student-t interval]
\newcommand{\AojPairedYieldOnesevenOverOneeighteight}{1.01~[0.75, 1.35]}  % experiments/FIGS/data/aoj_full_v1/analysis_v6/aoj_top.json :: per_label_set.label_sets.{17,188}.signal_yields :: sha256 5dfef902a5fe75ba :: yield ratio paired by run, geometric mean [95% Student-t interval]
\newcommand{\AojPairedYieldOnesevenOverOnesixtwo}{1.07~[0.51, 2.24]}  % experiments/FIGS/data/aoj_full_v1/analysis_v6/aoj_top.json :: per_label_set.label_sets.{17,162}.signal_yields :: sha256 5dfef902a5fe75ba :: yield ratio paired by run, geometric mean [95% Student-t interval]
\newcommand{\AojPairedYieldOnesixtwoMassOverOnesixtwo}{1.10~[0.57, 2.14]}  % experiments/FIGS/data/aoj_full_v1/analysis_v6/aoj_top.json :: per_label_set.label_sets.{162+mass,162}.signal_yields :: sha256 5dfef902a5fe75ba :: yield ratio paired by run, geometric mean [95% Student-t interval]
\newcommand{\AojPairedYieldOnesevenMassOverOneseven}{1.05~[0.74, 1.49]}  % experiments/FIGS/data/aoj_full_v1/analysis_v6/aoj_top.json :: per_label_set.label_sets.{17+mass,17}.signal_yields :: sha256 5dfef902a5fe75ba :: yield ratio paired by run, geometric mean [95% Student-t interval]
\newcommand{\AojLeakPairedShiftMax}{2\%}  % experiments/FIGS/data/aoj_full_v1/fit_v6/results.json :: models.*.top.leak_systematic.signal_yield against signal_yield, paired ratios :: sha256 2416fc21f93fcfa6 :: largest change of a paired yield ratio when the fail-region tops are scaled up
\newcommand{\AojRefEffPercent}{1.38\%}  % experiments/FIGS/data/aoj_full_v1/fit_v6/results.json :: reference.top.data_efficiency :: sha256 2416fc21f93fcfa6 :: fraction of all data jets passing its cut
\newcommand{\AojRefFitP}{\ensuremath{2.56\times10^{-11}}}  % experiments/FIGS/data/aoj_full_v1/fit_v6/results.json :: reference.top.asymptotic_p :: sha256 2416fc21f93fcfa6 :: goodness of fit of the reference's own top fit
\newcommand{\AojRefValidationP}{\ensuremath{p \leq 0.005} (0 of 200 toys reach the observed deviance)}  % experiments/FIGS/data/aoj_full_v1/fit_v6/results.json :: reference.top.validation.toy_p, n_toys :: sha256 2416fc21f93fcfa6 :: toy p-value of the background validation, (k + 1) / (n + 1)
\newcommand{\AojRefBandSignalZ}{2.5}  % experiments/FIGS/data/aoj_full_v1/fit_v6/results.json :: reference.top.validation.band_signal_z :: sha256 2416fc21f93fcfa6 :: significance of a peak of the score's own shape in the validation band
\newcommand{\AojModelEffMin}{1.01\%}  % experiments/FIGS/data/aoj_full_v1/fit_v6/results.json :: models.*.top.data_efficiency (min) :: sha256 2416fc21f93fcfa6 :: pretrained models
\newcommand{\AojModelEffMax}{1.05\%}  % experiments/FIGS/data/aoj_full_v1/fit_v6/results.json :: models.*.top.data_efficiency (max) :: sha256 2416fc21f93fcfa6 :: pretrained models
\newcommand{\AojModelsValidationPMin}{0.11}  % experiments/FIGS/data/aoj_full_v1/fit_v6/results.json :: models.*.top.validation.toy_p (min) :: sha256 2416fc21f93fcfa6 :: pretrained models
\newcommand{\AojRefYield}{29{,}143}  % experiments/FIGS/data/aoj_full_v1/fit_v6/results.json :: reference.top.signal_yield :: sha256 2416fc21f93fcfa6
\newcommand{\AojRefYieldErr}{865}  % experiments/FIGS/data/aoj_full_v1/fit_v6/results.json :: reference.top.signal_yield_err :: sha256 2416fc21f93fcfa6
\newcommand{\AojRefMass}{181.0}  % experiments/FIGS/data/aoj_full_v1/fit_v6/results.json :: reference.top.floated_mean :: sha256 2416fc21f93fcfa6 :: GeV
\newcommand{\AojRefWidth}{17.8}  % experiments/FIGS/data/aoj_full_v1/fit_v6/results.json :: reference.top.floated_width :: sha256 2416fc21f93fcfa6 :: GeV
\newcommand{\AojPublicYield}{1{,}813}  % experiments/FIGS/data/aoj_full_v1/fit_v6/results.json :: models.sophon-public.top.signal_yield :: sha256 2416fc21f93fcfa6
\newcommand{\AojPublicYieldErr}{193}  % experiments/FIGS/data/aoj_full_v1/fit_v6/results.json :: models.sophon-public.top.signal_yield_err :: sha256 2416fc21f93fcfa6
\newcommand{\AojPublicMass}{182.8}  % experiments/FIGS/data/aoj_full_v1/fit_v6/results.json :: models.sophon-public.top.floated_mean :: sha256 2416fc21f93fcfa6 :: GeV
\newcommand{\AojPublicWidth}{11.7}  % experiments/FIGS/data/aoj_full_v1/fit_v6/results.json :: models.sophon-public.top.floated_width :: sha256 2416fc21f93fcfa6 :: GeV
\newcommand{\AojInjSizeMin}{1{,}000}  % experiments/FIGS/data/aoj_injection_v1/summary.json :: groups[region=top,mode=ensemble].size (min) :: sha256 d9b66466ae24b584 :: injected tops
\newcommand{\AojInjSizeMax}{4{,}000}  % experiments/FIGS/data/aoj_injection_v1/summary.json :: groups[region=top,mode=ensemble].size (max) :: sha256 d9b66466ae24b584
\newcommand{\AojInjFloatPullSdMax}{2.0}  % experiments/FIGS/data/aoj_injection_v1/summary.json :: groups[region=top,mode=bootstrap,variant=float].pull.sd (max) :: sha256 d9b66466ae24b584 :: pull width with each model's own floated shape
\newcommand{\AojInjEnsembleRecoveryMin}{0.983}  % experiments/FIGS/data/aoj_injection_v1/summary.json :: groups[region=top,mode=ensemble].ratio (min) :: sha256 d9b66466ae24b584 :: the whole procedure, shape re-derived per toy set
\newcommand{\AojInjEnsembleRecoveryMax}{0.998}  % experiments/FIGS/data/aoj_injection_v1/summary.json :: groups[region=top,mode=ensemble].ratio (max) :: sha256 d9b66466ae24b584
\newcommand{\AojInjEnsemblePullSdMin}{1.02}  % experiments/FIGS/data/aoj_injection_v1/summary.json :: groups[region=top,mode=ensemble].pull (min) :: sha256 d9b66466ae24b584 :: the whole procedure, shape re-derived per toy set
\newcommand{\AojInjEnsemblePullSdMax}{1.09}  % experiments/FIGS/data/aoj_injection_v1/summary.json :: groups[region=top,mode=ensemble].pull (max) :: sha256 d9b66466ae24b584
\newcommand{\AojClosureNShards}{10}  % experiments/FIGS/data/aoj_checks_v2/closure_all_shards.json :: n_shards :: sha256 6a6018351d7048d2 :: staged shards of the open data
\newcommand{\AojClosurePtLo}{500}  % experiments/FIGS/data/aoj_checks_v2/closure_all_shards.json :: pt_window[0] :: sha256 6a6018351d7048d2 :: GeV
\newcommand{\AojClosurePtHi}{700}  % experiments/FIGS/data/aoj_checks_v2/closure_all_shards.json :: pt_window[1] :: sha256 6a6018351d7048d2 :: GeV
\newcommand{\AojClosureDzeroIqr}{8.7}  % experiments/FIGS/data/aoj_checks_v2/closure_all_shards.json :: domain_shift.part_d0.iqr_ratio :: sha256 6a6018351d7048d2 :: open data over JetClass-II QCD, pooled over the shards
\newcommand{\AojClosureDzIqr}{10.6}  % experiments/FIGS/data/aoj_checks_v2/closure_all_shards.json :: domain_shift.part_dz.iqr_ratio :: sha256 6a6018351d7048d2 :: open data over JetClass-II QCD, pooled over the shards
\newcommand{\AojClosureNeutral}{0.65}  % experiments/FIGS/data/aoj_checks_v2/closure_all_shards.json :: features.part_isNeutralHadron.ratio.per_shard (mean) :: sha256 6a6018351d7048d2 :: open data over JetClass-II QCD
\newcommand{\AojClosureCharged}{1.22}  % experiments/FIGS/data/aoj_checks_v2/closure_all_shards.json :: features.part_isChargedHadron.ratio.per_shard (mean) :: sha256 6a6018351d7048d2 :: open data over JetClass-II QCD
\newcommand{\AojClosureNPartAoj}{61}  % experiments/FIGS/data/aoj_checks_v2/closure_all_shards.json :: features.n_particles.aoj.per_shard (mean of the shard medians) :: sha256 6a6018351d7048d2 :: median
\newcommand{\AojClosureNPartRef}{48}  % experiments/FIGS/data/aoj_checks_v2/closure_all_shards.json :: features.n_particles.reference.per_shard (mean) :: sha256 6a6018351d7048d2 :: median
\newcommand{\RandMergingDrawsBvcFourprong}{three}  % configs/labelmaps/rand_label_map.v1.csv :: RAND_d* of label_X_YY_bbqq and label_X_YY_ccqq :: sha256 f395a25f52e3cfad :: first-grid partitions that merge the task's pair
\newcommand{\RandSplittingDrawsBvcFourprong}{one and two}  % configs/labelmaps/rand_label_map.v1.csv :: RAND_d* of label_X_YY_bbqq and label_X_YY_ccqq :: sha256 f395a25f52e3cfad :: first-grid partitions that keep the pair apart
\newcommand{\RandMergingDrawsVisibleContent}{one and three}  % configs/labelmaps/rand_label_map.v1.csv :: RAND_d* of label_X_YY_bbqq and label_X_YY_cqtauhv :: sha256 f395a25f52e3cfad :: first-grid partitions that merge the task's pair
\newcommand{\RandSplittingDrawsVisibleContent}{two}  % configs/labelmaps/rand_label_map.v1.csv :: RAND_d* of label_X_YY_bbqq and label_X_YY_cqtauhv :: sha256 f395a25f52e3cfad :: first-grid partitions that keep the pair apart
\newcommand{\RandVtwoNPartitions}{five}  % configs/arms/v2_grid.json :: arms[RAND2_*] (count) :: sha256 d8a31b1486992302 :: random partitions in the rerun
\newcommand{\RandVtwoRuns}{one}  % configs/arms/v2_paper_runs.txt :: runs of RAND2_* :: sha256 b04562b14ea8cbc7 :: runs per partition
\newcommand{\RandFlavRuns}{two}  % configs/arms/v2_paper_runs.txt :: runs of FLAV_* :: sha256 b04562b14ea8cbc7 :: runs per flavour-pair model
\newcommand{\RandVtwoNPairs}{seven}  % configs/labelmaps/rand_v2_selection.json :: pairs (count) :: sha256 98cf52d49c7d2f02 :: probe pairs the balance rule covers
\newcommand{\RandVtwoMergedMin}{two}  % configs/labelmaps/rand_v2_selection.json :: merged_in :: sha256 98cf52d49c7d2f02 :: partitions merging each probe pair, checked on the map
\newcommand{\RandVtwoMergedMax}{three}  % configs/labelmaps/rand_v2_selection.json :: merged_in :: sha256 98cf52d49c7d2f02 :: partitions merging each probe pair, checked on the map
\newcommand{\RandVtwoShareTol}{5\%}  % configs/labelmaps/rand_v2_selection.json :: realised_tolerance :: sha256 98cf52d49c7d2f02 :: largest relative deviation of a realised group share from its target (rule 3)
\newcommand{\RandVtwoTrials}{32}  % scripts/build_rand_control.py :: TRIALS :: sha256 5bc91d77a3f2e81f :: share-matched assignments drawn per random partition; the one sharing the fewest merged pairs with the 17-class label set is kept
\newcommand{\RandVtwoPoolDraws}{6{,}600}  % configs/labelmaps/rand_v2_selection.json :: pool.draws :: sha256 98cf52d49c7d2f02 :: candidate partitions drawn (seeds) before the realised-share rule
\newcommand{\RandVtwoPoolPass}{27}  % configs/labelmaps/rand_v2_selection.json :: rule3_seeds (count) :: sha256 98cf52d49c7d2f02 :: candidate partitions meeting the realised-share rule (rule 3)
\newcommand{\RandFlavSplitOrbit}{\ensuremath{X\to YY\to QQQQ}}  % configs/labelmaps/flavour_pair.v2.json :: split_orbit :: sha256 3b79ce5e44c1e2f4 :: the decays F1 cuts, Q any quark
\newcommand{\RandFlavSplitOrbitClasses}{22}  % configs/labelmaps/flavour_pair_map.v2.csv :: classes whose orbit is split_orbit (count) :: sha256 9aab235099ceac43 :: classes of the split orbit
\newcommand{\RandFlavCutClasses}{11}  % configs/labelmaps/flavour_pair.v2.json :: split_classes_moved (count) :: sha256 3b79ce5e44c1e2f4 :: classes of the split orbit that F1 moves: those with a b quark
\newcommand{\RandFlavOrbitsMoved}{three}  % configs/labelmaps/flavour_pair.v2.json :: orbits_moved_B_to_A (count) :: sha256 3b79ce5e44c1e2f4 :: whole orbits F1 moves the other way
\newcommand{\RandFlavOrbitsMovedList}{\ensuremath{X\to YY\to QQ\mu\mu}, \ensuremath{X\to YY\to QQ\tau_h\tau_h} and \ensuremath{X\to \mu\mu}}  % configs/labelmaps/flavour_pair.v2.json :: orbits_moved_B_to_A :: sha256 3b79ce5e44c1e2f4 :: the orbits F1 moves the other way, Q any quark
\newcommand{\RandFlavOrbitClassesMoved}{9}  % configs/labelmaps/flavour_pair_map.v2.csv :: classes whose orbit is in orbits_moved_B_to_A :: sha256 9aab235099ceac43 :: classes in those orbits
\newcommand{\RandFlavFonerCutWithB}{6}  % configs/labelmaps/flavour_pair_map.v2.csv :: FLAV_F1R != FLAV_F0 within split_orbit, with a b quark (count) :: sha256 9aab235099ceac43 :: classes with a b quark among those the F1r cut moves
\newcommand{\RandFlavBcPairs}{11}  % configs/labelmaps/flavour_pair.v2.json :: b_c_pairs_in_split_orbit (count) :: sha256 3b79ce5e44c1e2f4 :: pairs of the split orbit's classes that differ only by exchanging b and c
\newcommand{\RandFlavFoneBcSplit}{8}  % configs/labelmaps/flavour_pair_map.v2.csv :: b_c_pairs_in_split_orbit split by FLAV_F1 (count) :: sha256 9aab235099ceac43
\newcommand{\RandFlavFonerBcSplit}{none}  % configs/labelmaps/flavour_pair_map.v2.csv :: b_c_pairs_in_split_orbit split by FLAV_F1R (count) :: sha256 9aab235099ceac43 :: PRESPEC A14: none
\newcommand{\RandFlavFonerShareTol}{1\%}  % configs/labelmaps/flavour_pair.v2.json :: F1R.realised_tolerance_vs_F1 :: sha256 3b79ce5e44c1e2f4 :: largest allowed relative difference of a group's realised share from F1's
\newcommand{\RandFlavFonerShareDev}{0.77\%}  % configs/labelmaps/realised_native_shares.v2.json :: share summed per FLAV_F1R and FLAV_F1 group of flavour_pair_map.v2.csv (largest relative difference) :: sha256 5728c12e88fa3683 :: realised shares from the final loader's dry run
\newcommand{\MlstNRunsWord}{three}  % configs/arms/v2_paper_runs.txt :: runs of L188/L162/R42_Q1/R16_Q1 :: sha256 b04562b14ea8cbc7 :: pretraining runs per configuration
\newcommand{\MlstMassNRunsWord}{two}  % configs/arms/v2_paper_runs.txt :: runs of L162_MASS/R16_Q1_MASS :: sha256 b04562b14ea8cbc7 :: pretraining runs per configuration
\newcommand{\LpCMin}{\ensuremath{1\times10^{-3}}}  % experiments/EVAL/linear_probe_v2.py :: C_GRID :: sha256 4c576b44812c18a2 :: smallest inverse regularisation
\newcommand{\LpCMax}{100}  % experiments/EVAL/linear_probe_v2.py :: C_GRID :: sha256 4c576b44812c18a2 :: largest inverse regularisation
\newcommand{\LpNTestJcii}{2{,}000{,}000}  % experiments/EVAL/linear_probe_v2.py :: TEST_FIRST.jc2 :: sha256 4c576b44812c18a2 :: JetClass-II test jets
\newcommand{\LpNTestAucJcii}{500{,}000}  % experiments/EVAL/linear_probe_v2.py :: TEST_FIRST.jc2 / AUC_STRIDE :: sha256 4c576b44812c18a2 :: JetClass-II test jets of the macro AUC
\newcommand{\LpConvergedGrad}{\ensuremath{1\times10^{-5}}}  % experiments/EVAL/linear_probe_v2.py :: CONVERGED_GRAD :: sha256 4c576b44812c18a2 :: largest gradient component a converged probe ends below
\newcommand{\FillNFeatures}{17}  % configs/finetune/TopReference.yaml :: inputs.pf_features.vars (count) :: sha256 501a6ffe5d0a8051 :: per-particle input features of the network
\newcommand{\FillTopNZero}{10}  % configs/finetune/TopReference.yaml :: inputs.pf_features.vars, part_zero_* (count) :: sha256 501a6ffe5d0a8051 :: features the data set lacks, set to zero
\newcommand{\FillTopNPresent}{7}  % configs/finetune/TopReference.yaml :: inputs.pf_features.vars, not part_zero_* (count) :: sha256 501a6ffe5d0a8051 :: features the data set provides
\newcommand{\FillQgNZero}{4}  % configs/finetune/EnergyFlowQG.yaml :: inputs.pf_features.vars, part_zero_* (count) :: sha256 db42d6471e8b865c :: features the data set lacks, set to zero
\newcommand{\FillQgNPresent}{13}  % configs/finetune/EnergyFlowQG.yaml :: inputs.pf_features.vars, not part_zero_* (count) :: sha256 db42d6471e8b865c :: features the data set provides
\newcommand{\DesignVtwoDataFraction}{20\%}  % experiments/MTX/k8s/v2/grid/job-mtx2-l188-s1-raunav.yaml :: --data-fraction :: sha256 d28ce02267a6b87a :: share of every file read per epoch
\newcommand{\DesignVtwoNValFiles}{25}  % experiments/MTX/k8s/v2/grid/job-mtx2-l188-s1-raunav.yaml :: --data-val (files, brace ranges expanded) :: sha256 d28ce02267a6b87a :: validation files of the fixed sample
\newcommand{\DesignWavgFirstEpoch}{70}  % experiments/MTX/pretrain_v2.py :: --num-epochs - WAVG_EPOCHS :: sha256 95da5a6e029855ee :: first epoch of the weight average (epochs count from 0)
\newcommand{\DesignWavgLastEpoch}{79}  % experiments/MTX/k8s/v2/grid/job-mtx2-l188-s1-raunav.yaml :: --num-epochs - 1 :: sha256 d28ce02267a6b87a :: last epoch
\newcommand{\DesignBnJets}{200{,}000}  % experiments/MTX/pretrain_v2.py :: BN_JETS :: sha256 95da5a6e029855ee :: training jets the BatchNorm statistics are recomputed on
\newcommand{\DesignVtwoNWindows}{five}  % experiments/MTX/k8s/v2/grid/job-mtx2-l188-s1-raunav.yaml :: 1 / --data-fraction :: sha256 d28ce02267a6b87a :: disjoint windows of each file's rows, one read per epoch
\newcommand{\DesignLofoNWindows}{three}  % experiments/MTX/k8s/v2/cut/grid/job-mtx2-l188lofo4p-s1-raunav.yaml :: --data-windows :: sha256 44a9573929efb7ac :: windows of the models trained without one family
\newcommand{\DesignVtwoPlanUsed}{97\%}  % experiments/FIGS/data/v2_loader/loader_dryrun/dryrun_s176_seed1.json :: epochs[*].(max_fetch_id + 1) / args.data_split_num, mean :: sha256 03b6586b26ad7a2f :: share of an epoch's planned reading used before its jets are reached, 20 epochs
\newcommand{\DesignLofoPlanUsed}{76\%}  % experiments/FIGS/data/v2_loader/loader_dryrun/dryrun_lofo4p_seed1.json :: epochs[*].(max_fetch_id + 1) / fetches_per_pass, mean :: sha256 2a1315e611d13ff5 :: the same, leaving the family out
\newcommand{\DesignLofoRowFactor}{1.29}  % experiments/FIGS/data/v2_loader/loader_dryrun/dryrun_lofo4p_seed1.json :: (mean plan used / --data-windows) over (mean plan used / k) of the full models' dry run :: sha256 2a1315e611d13ff5 :: how much more often a row of a file is read per epoch when the family is left out
\newcommand{\DesignDistinctJetsCycle}{\ensuremath{3.3\times10^{7}}}  % experiments/FIGS/data/v2_loader/loader_dryrun/dryrun_s176_seed1.json :: sum of epochs[c*k:(c+1)*k].distinct_jets, mean over 4 cycles :: sha256 03b6586b26ad7a2f :: distinct jets the full models train on in one cycle of windows
\newcommand{\DesignLofoFamilyShare}{22.4\%}  % experiments/FIGS/data/v2_loader/loader_dryrun/dryrun_s176_seed1.json :: epochs[*].native_counts over summary.labels_absent_every_epoch of the family-out run :: sha256 03b6586b26ad7a2f :: the left-out family's share of the full models' stream
\newcommand{\DesignLofoNClasses}{27}  % experiments/FIGS/data/v2_loader/loader_dryrun/dryrun_lofo4p_seed1.json :: summary.labels_absent_every_epoch (count) :: sha256 2a1315e611d13ff5 :: native classes of the left-out family
\newcommand{\DesignLofoExposureFactor}{1.29}  % experiments/FIGS/data/v2_loader/loader_dryrun/dryrun_s176_seed1.json :: 1 / (1 - family share) :: sha256 03b6586b26ad7a2f :: how much more often every other class is seen
\newcommand{\DesignLofoQcdShareFull}{14.4\%}  % experiments/FIGS/data/v2_loader/loader_dryrun/dryrun_s176_seed1.json :: summary.mean_qcd_share :: sha256 03b6586b26ad7a2f :: QCD share of the realised stream
\newcommand{\DesignLofoQcdShareLofo}{18.6\%}  % experiments/FIGS/data/v2_loader/loader_dryrun/dryrun_lofo4p_seed1.json :: summary.mean_qcd_share :: sha256 2a1315e611d13ff5 :: QCD share of the realised stream
\newcommand{\DesignLrFlatPercent}{70\%}  % experiments/MTX/pretrain_v2.py :: make_optimizer: 1 - the decay share :: sha256 95da5a6e029855ee :: share of the epochs at the initial rate
\newcommand{\DesignLrDecayFirstEpoch}{56}  % experiments/MTX/pretrain_v2.py :: num_epochs - int(num_epochs * decay share) :: sha256 95da5a6e029855ee :: first epoch at a reduced rate (epochs from 0)
\newcommand{\DesignLrEndPercent}{1\%}  % experiments/MTX/pretrain_v2.py :: make_optimizer: gamma ** n_decay :: sha256 95da5a6e029855ee :: the last epoch's rate over the initial rate
\newcommand{\DesignEndLr}{\ensuremath{5\times10^{-6}}}  % experiments/MTX/k8s/v2/grid/job-mtx2-l188-s1-raunav.yaml :: --start-lr x the final factor :: sha256 d28ce02267a6b87a :: the rate of the last epoch
\newcommand{\DesignWeaverVersion}{0.4.17}  % experiments/MTX/pretrain_v2.py :: the weaver release whose loop it reproduces :: sha256 95da5a6e029855ee
\newcommand{\DesignMeiFactor}{1.1}  % src/stats/mde.py :: MEI_LOG = ln(factor) :: sha256 626468dcc37940bb :: the smallest effect of interest, as a factor
\newcommand{\DesignVtwoGpu}{RTX 3090}  % experiments/MTX/k8s/v2/grid/job-mtx2-flavf0-s1-raunav.yaml :: nodeAffinity nvidia.com/gpu.product :: sha256 960cad1691314980 :: the product of run index 1, every configuration
\newcommand{\DesignVtwoSecondGpu}{\pending{pending: the GPU type of each run index}}  % experiments/FIGS/make_tables.py :: emit_v2_recipe[DesignVtwoSecondGpu] :: sha256 10e6ec08356a4a2b :: a marker, not a number: the grid's job specs put every run index on one product
\newcommand{\DesignVtwoSecondGpuRuns}{\pending{pending: the GPU type of each run index}}  % experiments/FIGS/make_tables.py :: emit_v2_recipe[DesignVtwoSecondGpuRuns] :: sha256 10e6ec08356a4a2b :: a marker, not a number: the grid's job specs put every run index on one product
\newcommand{\MassNRunsWord}{five}  % experiments/FIGS/data/probe_ladder_mass2x2_mlp2/s1.json :: per-seed files (count) :: sha256 906e778d7043fe67 :: pretraining runs of each mass-output configuration
\newcommand{\MassOmaLinearOnesixtwo}{\ensuremath{(2.36 \pm 0.06)\times10^{-3}}}  % experiments/FIGS/data/probe_ladder_v2_mlp2/analysis/seed_level_results.json :: table[bvc_resonant,linear,162].auc :: sha256 53e1ec81343ac2c1 :: 1-AUC, mean +- SD over 5 seeds
\newcommand{\MassOmaLinearOnesixtwoMass}{\ensuremath{(2.62 \pm 0.08)\times10^{-3}}}  % experiments/FIGS/data/probe_ladder_mass2x2_mlp2/s1.json :: tasks.bvc_resonant.arms.l162mass-s*.linear.auc :: sha256 906e778d7043fe67 :: 1-AUC, mean +- SD over the 5 per-seed files
\newcommand{\MassOmaRatioLinearOnesixtwo}{1.11}  % experiments/FIGS/data/probe_ladder_mass2x2_mlp2/s1.json :: tasks.bvc_resonant.arms.l162mass-s*.linear.auc over table[bvc_resonant,linear,162].auc :: sha256 906e778d7043fe67 :: seed mean of 1-AUC with the mass output over without
\newcommand{\MassOmaLinearOneseven}{\ensuremath{(9.3 \pm 1.5)\times10^{-3}}}  % experiments/FIGS/data/probe_ladder_v2_mlp2/analysis/seed_level_results.json :: table[bvc_resonant,linear,17].auc :: sha256 53e1ec81343ac2c1 :: 1-AUC, mean +- SD over 5 seeds
\newcommand{\MassOmaLinearOnesevenMass}{\ensuremath{(3.0 \pm 1.2)\times10^{-2}}}  % experiments/FIGS/data/probe_ladder_mass2x2_mlp2/s1.json :: tasks.bvc_resonant.arms.r16q1mass-s*.linear.auc :: sha256 906e778d7043fe67 :: 1-AUC, mean +- SD over the 5 per-seed files
\newcommand{\MassOmaRatioLinearOneseven}{3.2}  % experiments/FIGS/data/probe_ladder_mass2x2_mlp2/s1.json :: tasks.bvc_resonant.arms.r16q1mass-s*.linear.auc over table[bvc_resonant,linear,17].auc :: sha256 906e778d7043fe67 :: seed mean of 1-AUC with the mass output over without
\newcommand{\MassOmaMlpOnesixtwo}{\ensuremath{(2.32 \pm 0.05)\times10^{-3}}}  % experiments/FIGS/data/probe_ladder_v2_mlp2/analysis/seed_level_results.json :: table[bvc_resonant,mlp,162].auc :: sha256 53e1ec81343ac2c1 :: 1-AUC, mean +- SD over 5 seeds
\newcommand{\MassOmaMlpOnesixtwoMass}{\ensuremath{(2.52 \pm 0.06)\times10^{-3}}}  % experiments/FIGS/data/probe_ladder_mass2x2_mlp2/s1.json :: tasks.bvc_resonant.arms.l162mass-s*.mlp.auc :: sha256 906e778d7043fe67 :: 1-AUC, mean +- SD over the 5 per-seed files
\newcommand{\MassOmaRatioMlpOnesixtwo}{1.09}  % experiments/FIGS/data/probe_ladder_mass2x2_mlp2/s1.json :: tasks.bvc_resonant.arms.l162mass-s*.mlp.auc over table[bvc_resonant,mlp,162].auc :: sha256 906e778d7043fe67 :: seed mean of 1-AUC with the mass output over without
\newcommand{\MassOmaMlpOneseven}{\ensuremath{(7.3 \pm 0.9)\times10^{-3}}}  % experiments/FIGS/data/probe_ladder_v2_mlp2/analysis/seed_level_results.json :: table[bvc_resonant,mlp,17].auc :: sha256 53e1ec81343ac2c1 :: 1-AUC, mean +- SD over 5 seeds
\newcommand{\MassOmaMlpOnesevenMass}{\ensuremath{(1.5 \pm 0.4)\times10^{-2}}}  % experiments/FIGS/data/probe_ladder_mass2x2_mlp2/s1.json :: tasks.bvc_resonant.arms.r16q1mass-s*.mlp.auc :: sha256 906e778d7043fe67 :: 1-AUC, mean +- SD over the 5 per-seed files
\newcommand{\MassOmaRatioMlpOneseven}{2.0}  % experiments/FIGS/data/probe_ladder_mass2x2_mlp2/s1.json :: tasks.bvc_resonant.arms.r16q1mass-s*.mlp.auc over table[bvc_resonant,mlp,17].auc :: sha256 906e778d7043fe67 :: seed mean of 1-AUC with the mass output over without
\newcommand{\FtAucJciiEThreeOneeighteight}{0.902\,\ensuremath{\pm}\,0.004}  % experiments/FIGS/data/w2b_leg1_metrics_v2.json :: cells.{l188-s1,l188-s2,l188-s3,l188-s4,l188-s5}.N1000.s1.macro_auc_ovr :: sha256 b56bf70be9982044 :: mean +- SD over the row's 5 pretrained models
\newcommand{\FtAucJciiEFourOneeighteight}{0.9723\,\ensuremath{\pm}\,0.0004}  % experiments/FIGS/data/w2b_leg1_metrics_v2.json :: cells.{l188-s1,l188-s2,l188-s3,l188-s4,l188-s5}.N10000.s1.macro_auc_ovr :: sha256 b56bf70be9982044 :: mean +- SD over the row's 5 pretrained models
\newcommand{\FtAucJciiEFiveOneeighteight}{0.9793\,\ensuremath{\pm}\,0.0002}  % experiments/FIGS/data/w2b_leg1_metrics_v2.json :: cells.{l188-s1,l188-s2,l188-s3,l188-s4,l188-s5}.N100000.s1.macro_auc_ovr :: sha256 b56bf70be9982044 :: mean +- SD over the row's 5 pretrained models
\newcommand{\FtAucJciiESixOneeighteight}{0.9809\,\ensuremath{\pm}\,\ensuremath{<}0.0001}  % experiments/FIGS/data/w2b_leg1_metrics_v2.json :: cells.{l188-s1,l188-s2,l188-s3,l188-s4,l188-s5}.N1000000.s1.macro_auc_ovr :: sha256 b56bf70be9982044 :: mean +- SD over the row's 5 pretrained models
\newcommand{\FtAucJciiEThreeOnesixtwo}{0.906\,\ensuremath{\pm}\,0.007}  % experiments/FIGS/data/w2b_leg1_metrics_v2.json :: cells.{l162-s1b,l162-s2,l162-s3,l162-s4,l162-s5}.N1000.s1.macro_auc_ovr :: sha256 b56bf70be9982044 :: mean +- SD over the row's 5 pretrained models
\newcommand{\FtAucJciiEFourOnesixtwo}{0.9724\,\ensuremath{\pm}\,0.0005}  % experiments/FIGS/data/w2b_leg1_metrics_v2.json :: cells.{l162-s1b,l162-s2,l162-s3,l162-s4,l162-s5}.N10000.s1.macro_auc_ovr :: sha256 b56bf70be9982044 :: mean +- SD over the row's 5 pretrained models
\newcommand{\FtAucJciiEFiveOnesixtwo}{0.9794\,\ensuremath{\pm}\,0.0002}  % experiments/FIGS/data/w2b_leg1_metrics_v2.json :: cells.{l162-s1b,l162-s2,l162-s3,l162-s4,l162-s5}.N100000.s1.macro_auc_ovr :: sha256 b56bf70be9982044 :: mean +- SD over the row's 5 pretrained models
\newcommand{\FtAucJciiESixOnesixtwo}{0.9808\,\ensuremath{\pm}\,0.0001}  % experiments/FIGS/data/w2b_leg1_metrics_v2.json :: cells.{l162-s1b,l162-s2,l162-s3,l162-s4,l162-s5}.N1000000.s1.macro_auc_ovr :: sha256 b56bf70be9982044 :: mean +- SD over the row's 5 pretrained models
\newcommand{\FtAucJciiEThreeFourthree}{0.918\,\ensuremath{\pm}\,0.005}  % experiments/FIGS/data/w2b_leg1_metrics_v2.json :: cells.{r42q1-s1,r42q1-s2,r42q1-s3,r42q1-s4,r42q1-s5}.N1000.s1.macro_auc_ovr :: sha256 b56bf70be9982044 :: mean +- SD over the row's 5 pretrained models
\newcommand{\FtAucJciiEFourFourthree}{0.9667\,\ensuremath{\pm}\,0.0005}  % experiments/FIGS/data/w2b_leg1_metrics_v2.json :: cells.{r42q1-s1,r42q1-s2,r42q1-s3,r42q1-s4,r42q1-s5}.N10000.s1.macro_auc_ovr :: sha256 b56bf70be9982044 :: mean +- SD over the row's 5 pretrained models
\newcommand{\FtAucJciiEFiveFourthree}{0.9775\,\ensuremath{\pm}\,0.0003}  % experiments/FIGS/data/w2b_leg1_metrics_v2.json :: cells.{r42q1-s1,r42q1-s2,r42q1-s3,r42q1-s4,r42q1-s5}.N100000.s1.macro_auc_ovr :: sha256 b56bf70be9982044 :: mean +- SD over the row's 5 pretrained models
\newcommand{\FtAucJciiESixFourthree}{0.9800\,\ensuremath{\pm}\,0.0001}  % experiments/FIGS/data/w2b_leg1_metrics_v2.json :: cells.{r42q1-s1,r42q1-s2,r42q1-s3,r42q1-s4,r42q1-s5}.N1000000.s1.macro_auc_ovr :: sha256 b56bf70be9982044 :: mean +- SD over the row's 5 pretrained models
\newcommand{\FtAucJciiEThreeOneseven}{0.905\,\ensuremath{\pm}\,0.008}  % experiments/FIGS/data/w2b_leg1_metrics_v2.json :: cells.{r16q1-s1,r16q1-s2,r16q1-s3,r16q1-s4,r16q1-s5}.N1000.s1.macro_auc_ovr :: sha256 b56bf70be9982044 :: mean +- SD over the row's 5 pretrained models
\newcommand{\FtAucJciiEFourOneseven}{0.9635\,\ensuremath{\pm}\,0.0005}  % experiments/FIGS/data/w2b_leg1_metrics_v2.json :: cells.{r16q1-s1,r16q1-s2,r16q1-s3,r16q1-s4,r16q1-s5}.N10000.s1.macro_auc_ovr :: sha256 b56bf70be9982044 :: mean +- SD over the row's 5 pretrained models
\newcommand{\FtAucJciiEFiveOneseven}{0.9748\,\ensuremath{\pm}\,0.0003}  % experiments/FIGS/data/w2b_leg1_metrics_v2.json :: cells.{r16q1-s1,r16q1-s2,r16q1-s3,r16q1-s4,r16q1-s5}.N100000.s1.macro_auc_ovr :: sha256 b56bf70be9982044 :: mean +- SD over the row's 5 pretrained models
\newcommand{\FtAucJciiESixOneseven}{0.9781\,\ensuremath{\pm}\,0.0001}  % experiments/FIGS/data/w2b_leg1_metrics_v2.json :: cells.{r16q1-s1,r16q1-s2,r16q1-s3,r16q1-s4,r16q1-s5}.N1000000.s1.macro_auc_ovr :: sha256 b56bf70be9982044 :: mean +- SD over the row's 5 pretrained models
\newcommand{\FtAucJciiEThreeOnesixtwoMass}{0.921\,\ensuremath{\pm}\,0.005}  % experiments/FIGS/data/w2b_leg1_metrics_v2.json :: cells.{l162mass-s1,l162mass-s2,l162mass-s3,l162mass-s4,l162mass-s5}.N1000.s1.macro_auc_ovr :: sha256 b56bf70be9982044 :: mean +- SD over the row's 5 pretrained models
\newcommand{\FtAucJciiEFourOnesixtwoMass}{0.9724\,\ensuremath{\pm}\,0.0013}  % experiments/FIGS/data/w2b_leg1_metrics_v2.json :: cells.{l162mass-s1,l162mass-s2,l162mass-s3,l162mass-s4,l162mass-s5}.N10000.s1.macro_auc_ovr :: sha256 b56bf70be9982044 :: mean +- SD over the row's 5 pretrained models
\newcommand{\FtAucJciiEFiveOnesixtwoMass}{0.9790\,\ensuremath{\pm}\,0.0001}  % experiments/FIGS/data/w2b_leg1_metrics_v2.json :: cells.{l162mass-s1,l162mass-s2,l162mass-s3,l162mass-s4,l162mass-s5}.N100000.s1.macro_auc_ovr :: sha256 b56bf70be9982044 :: mean +- SD over the row's 5 pretrained models
\newcommand{\FtAucJciiESixOnesixtwoMass}{0.9805\,\ensuremath{\pm}\,0.0001}  % experiments/FIGS/data/w2b_leg1_metrics_v2.json :: cells.{l162mass-s1,l162mass-s2,l162mass-s3,l162mass-s4,l162mass-s5}.N1000000.s1.macro_auc_ovr :: sha256 b56bf70be9982044 :: mean +- SD over the row's 5 pretrained models
\newcommand{\FtAucJciiEThreeOnesevenMass}{0.904\,\ensuremath{\pm}\,0.007}  % experiments/FIGS/data/w2b_leg1_metrics_v2.json :: cells.{r16q1mass-s1,r16q1mass-s2,r16q1mass-s3,r16q1mass-s4,r16q1mass-s5}.N1000.s1.macro_auc_ovr :: sha256 b56bf70be9982044 :: mean +- SD over the row's 5 pretrained models
\newcommand{\FtAucJciiEFourOnesevenMass}{0.9639\,\ensuremath{\pm}\,0.0007}  % experiments/FIGS/data/w2b_leg1_metrics_v2.json :: cells.{r16q1mass-s1,r16q1mass-s2,r16q1mass-s3,r16q1mass-s4,r16q1mass-s5}.N10000.s1.macro_auc_ovr :: sha256 b56bf70be9982044 :: mean +- SD over the row's 5 pretrained models
\newcommand{\FtAucJciiEFiveOnesevenMass}{0.9740\,\ensuremath{\pm}\,0.0004}  % experiments/FIGS/data/w2b_leg1_metrics_v2.json :: cells.{r16q1mass-s1,r16q1mass-s2,r16q1mass-s3,r16q1mass-s4,r16q1mass-s5}.N100000.s1.macro_auc_ovr :: sha256 b56bf70be9982044 :: mean +- SD over the row's 5 pretrained models
\newcommand{\FtAucJciiESixOnesevenMass}{0.9774\,\ensuremath{\pm}\,0.0001}  % experiments/FIGS/data/w2b_leg1_metrics_v2.json :: cells.{r16q1mass-s1,r16q1mass-s2,r16q1mass-s3,r16q1mass-s4,r16q1mass-s5}.N1000000.s1.macro_auc_ovr :: sha256 b56bf70be9982044 :: mean +- SD over the row's 5 pretrained models
\newcommand{\FtAucJciiEThreeRandDrawOne}{0.921}  % experiments/FIGS/data/w2b_leg1_metrics_v2.json :: cells.{rand-d1-s1b}.N1000.s1.macro_auc_ovr :: sha256 b56bf70be9982044 :: one run
\newcommand{\FtAucJciiEFourRandDrawOne}{0.9675}  % experiments/FIGS/data/w2b_leg1_metrics_v2.json :: cells.{rand-d1-s1b}.N10000.s1.macro_auc_ovr :: sha256 b56bf70be9982044 :: one run
\newcommand{\FtAucJciiEFiveRandDrawOne}{0.97617}  % experiments/FIGS/data/w2b_leg1_metrics_v2.json :: cells.{rand-d1-s1b}.N100000.s1.macro_auc_ovr :: sha256 b56bf70be9982044 :: one run
\newcommand{\FtAucJciiESixRandDrawOne}{0.97941}  % experiments/FIGS/data/w2b_leg1_metrics_v2.json :: cells.{rand-d1-s1b}.N1000000.s1.macro_auc_ovr :: sha256 b56bf70be9982044 :: one run
\newcommand{\FtAucJciiEThreeRandDrawTwo}{0.914}  % experiments/FIGS/data/w2b_leg1_metrics_v2.json :: cells.{rand-d2-s2}.N1000.s1.macro_auc_ovr :: sha256 b56bf70be9982044 :: one run
\newcommand{\FtAucJciiEFourRandDrawTwo}{0.9675}  % experiments/FIGS/data/w2b_leg1_metrics_v2.json :: cells.{rand-d2-s2}.N10000.s1.macro_auc_ovr :: sha256 b56bf70be9982044 :: one run
\newcommand{\FtAucJciiEFiveRandDrawTwo}{0.97668}  % experiments/FIGS/data/w2b_leg1_metrics_v2.json :: cells.{rand-d2-s2}.N100000.s1.macro_auc_ovr :: sha256 b56bf70be9982044 :: one run
\newcommand{\FtAucJciiESixRandDrawTwo}{0.97943}  % experiments/FIGS/data/w2b_leg1_metrics_v2.json :: cells.{rand-d2-s2}.N1000000.s1.macro_auc_ovr :: sha256 b56bf70be9982044 :: one run
\newcommand{\FtAucJciiEThreeRandDrawThree}{0.918}  % experiments/FIGS/data/w2b_leg1_metrics_v2.json :: cells.{rand-d3-s3}.N1000.s1.macro_auc_ovr :: sha256 b56bf70be9982044 :: one run
\newcommand{\FtAucJciiEFourRandDrawThree}{0.9686}  % experiments/FIGS/data/w2b_leg1_metrics_v2.json :: cells.{rand-d3-s3}.N10000.s1.macro_auc_ovr :: sha256 b56bf70be9982044 :: one run
\newcommand{\FtAucJciiEFiveRandDrawThree}{0.97605}  % experiments/FIGS/data/w2b_leg1_metrics_v2.json :: cells.{rand-d3-s3}.N100000.s1.macro_auc_ovr :: sha256 b56bf70be9982044 :: one run
\newcommand{\FtAucJciiESixRandDrawThree}{0.97950}  % experiments/FIGS/data/w2b_leg1_metrics_v2.json :: cells.{rand-d3-s3}.N1000000.s1.macro_auc_ovr :: sha256 b56bf70be9982044 :: one run
\newcommand{\FtAccJciiEThreeOneeighteight}{0.445\,\ensuremath{\pm}\,0.004}  % experiments/FIGS/data/w2b_leg1_metrics_v2.json :: cells.{l188-s1,l188-s2,l188-s3,l188-s4,l188-s5}.N1000.s1.accuracy :: sha256 b56bf70be9982044 :: mean +- SD over the row's 5 pretrained models
\newcommand{\FtAccJciiEFourOneeighteight}{0.5199\,\ensuremath{\pm}\,0.0018}  % experiments/FIGS/data/w2b_leg1_metrics_v2.json :: cells.{l188-s1,l188-s2,l188-s3,l188-s4,l188-s5}.N10000.s1.accuracy :: sha256 b56bf70be9982044 :: mean +- SD over the row's 5 pretrained models
\newcommand{\FtAccJciiEFiveOneeighteight}{0.5379\,\ensuremath{\pm}\,0.0010}  % experiments/FIGS/data/w2b_leg1_metrics_v2.json :: cells.{l188-s1,l188-s2,l188-s3,l188-s4,l188-s5}.N100000.s1.accuracy :: sha256 b56bf70be9982044 :: mean +- SD over the row's 5 pretrained models
\newcommand{\FtAccJciiESixOneeighteight}{0.5443\,\ensuremath{\pm}\,0.0003}  % experiments/FIGS/data/w2b_leg1_metrics_v2.json :: cells.{l188-s1,l188-s2,l188-s3,l188-s4,l188-s5}.N1000000.s1.accuracy :: sha256 b56bf70be9982044 :: mean +- SD over the row's 5 pretrained models
\newcommand{\FtAccJciiEThreeOnesixtwo}{0.448\,\ensuremath{\pm}\,0.006}  % experiments/FIGS/data/w2b_leg1_metrics_v2.json :: cells.{l162-s1b,l162-s2,l162-s3,l162-s4,l162-s5}.N1000.s1.accuracy :: sha256 b56bf70be9982044 :: mean +- SD over the row's 5 pretrained models
\newcommand{\FtAccJciiEFourOnesixtwo}{0.5212\,\ensuremath{\pm}\,0.0013}  % experiments/FIGS/data/w2b_leg1_metrics_v2.json :: cells.{l162-s1b,l162-s2,l162-s3,l162-s4,l162-s5}.N10000.s1.accuracy :: sha256 b56bf70be9982044 :: mean +- SD over the row's 5 pretrained models
\newcommand{\FtAccJciiEFiveOnesixtwo}{0.5387\,\ensuremath{\pm}\,0.0008}  % experiments/FIGS/data/w2b_leg1_metrics_v2.json :: cells.{l162-s1b,l162-s2,l162-s3,l162-s4,l162-s5}.N100000.s1.accuracy :: sha256 b56bf70be9982044 :: mean +- SD over the row's 5 pretrained models
\newcommand{\FtAccJciiESixOnesixtwo}{0.5440\,\ensuremath{\pm}\,0.0006}  % experiments/FIGS/data/w2b_leg1_metrics_v2.json :: cells.{l162-s1b,l162-s2,l162-s3,l162-s4,l162-s5}.N1000000.s1.accuracy :: sha256 b56bf70be9982044 :: mean +- SD over the row's 5 pretrained models
\newcommand{\FtAccJciiEThreeFourthree}{0.380\,\ensuremath{\pm}\,0.005}  % experiments/FIGS/data/w2b_leg1_metrics_v2.json :: cells.{r42q1-s1,r42q1-s2,r42q1-s3,r42q1-s4,r42q1-s5}.N1000.s1.accuracy :: sha256 b56bf70be9982044 :: mean +- SD over the row's 5 pretrained models
\newcommand{\FtAccJciiEFourFourthree}{0.481\,\ensuremath{\pm}\,0.004}  % experiments/FIGS/data/w2b_leg1_metrics_v2.json :: cells.{r42q1-s1,r42q1-s2,r42q1-s3,r42q1-s4,r42q1-s5}.N10000.s1.accuracy :: sha256 b56bf70be9982044 :: mean +- SD over the row's 5 pretrained models
\newcommand{\FtAccJciiEFiveFourthree}{0.5219\,\ensuremath{\pm}\,0.0017}  % experiments/FIGS/data/w2b_leg1_metrics_v2.json :: cells.{r42q1-s1,r42q1-s2,r42q1-s3,r42q1-s4,r42q1-s5}.N100000.s1.accuracy :: sha256 b56bf70be9982044 :: mean +- SD over the row's 5 pretrained models
\newcommand{\FtAccJciiESixFourthree}{0.5362\,\ensuremath{\pm}\,0.0007}  % experiments/FIGS/data/w2b_leg1_metrics_v2.json :: cells.{r42q1-s1,r42q1-s2,r42q1-s3,r42q1-s4,r42q1-s5}.N1000000.s1.accuracy :: sha256 b56bf70be9982044 :: mean +- SD over the row's 5 pretrained models
\newcommand{\FtAccJciiEThreeOneseven}{0.332\,\ensuremath{\pm}\,0.004}  % experiments/FIGS/data/w2b_leg1_metrics_v2.json :: cells.{r16q1-s1,r16q1-s2,r16q1-s3,r16q1-s4,r16q1-s5}.N1000.s1.accuracy :: sha256 b56bf70be9982044 :: mean +- SD over the row's 5 pretrained models
\newcommand{\FtAccJciiEFourOneseven}{0.447\,\ensuremath{\pm}\,0.013}  % experiments/FIGS/data/w2b_leg1_metrics_v2.json :: cells.{r16q1-s1,r16q1-s2,r16q1-s3,r16q1-s4,r16q1-s5}.N10000.s1.accuracy :: sha256 b56bf70be9982044 :: mean +- SD over the row's 5 pretrained models
\newcommand{\FtAccJciiEFiveOneseven}{0.5078\,\ensuremath{\pm}\,0.0011}  % experiments/FIGS/data/w2b_leg1_metrics_v2.json :: cells.{r16q1-s1,r16q1-s2,r16q1-s3,r16q1-s4,r16q1-s5}.N100000.s1.accuracy :: sha256 b56bf70be9982044 :: mean +- SD over the row's 5 pretrained models
\newcommand{\FtAccJciiESixOneseven}{0.5227\,\ensuremath{\pm}\,0.0010}  % experiments/FIGS/data/w2b_leg1_metrics_v2.json :: cells.{r16q1-s1,r16q1-s2,r16q1-s3,r16q1-s4,r16q1-s5}.N1000000.s1.accuracy :: sha256 b56bf70be9982044 :: mean +- SD over the row's 5 pretrained models
\newcommand{\FtAccJciiEThreeOnesixtwoMass}{0.451\,\ensuremath{\pm}\,0.007}  % experiments/FIGS/data/w2b_leg1_metrics_v2.json :: cells.{l162mass-s1,l162mass-s2,l162mass-s3,l162mass-s4,l162mass-s5}.N1000.s1.accuracy :: sha256 b56bf70be9982044 :: mean +- SD over the row's 5 pretrained models
\newcommand{\FtAccJciiEFourOnesixtwoMass}{0.515\,\ensuremath{\pm}\,0.004}  % experiments/FIGS/data/w2b_leg1_metrics_v2.json :: cells.{l162mass-s1,l162mass-s2,l162mass-s3,l162mass-s4,l162mass-s5}.N10000.s1.accuracy :: sha256 b56bf70be9982044 :: mean +- SD over the row's 5 pretrained models
\newcommand{\FtAccJciiEFiveOnesixtwoMass}{0.5357\,\ensuremath{\pm}\,0.0007}  % experiments/FIGS/data/w2b_leg1_metrics_v2.json :: cells.{l162mass-s1,l162mass-s2,l162mass-s3,l162mass-s4,l162mass-s5}.N100000.s1.accuracy :: sha256 b56bf70be9982044 :: mean +- SD over the row's 5 pretrained models
\newcommand{\FtAccJciiESixOnesixtwoMass}{0.5409\,\ensuremath{\pm}\,0.0010}  % experiments/FIGS/data/w2b_leg1_metrics_v2.json :: cells.{l162mass-s1,l162mass-s2,l162mass-s3,l162mass-s4,l162mass-s5}.N1000000.s1.accuracy :: sha256 b56bf70be9982044 :: mean +- SD over the row's 5 pretrained models
\newcommand{\FtAccJciiEThreeOnesevenMass}{0.331\,\ensuremath{\pm}\,0.005}  % experiments/FIGS/data/w2b_leg1_metrics_v2.json :: cells.{r16q1mass-s1,r16q1mass-s2,r16q1mass-s3,r16q1mass-s4,r16q1mass-s5}.N1000.s1.accuracy :: sha256 b56bf70be9982044 :: mean +- SD over the row's 5 pretrained models
\newcommand{\FtAccJciiEFourOnesevenMass}{0.4507\,\ensuremath{\pm}\,0.0034}  % experiments/FIGS/data/w2b_leg1_metrics_v2.json :: cells.{r16q1mass-s1,r16q1mass-s2,r16q1mass-s3,r16q1mass-s4,r16q1mass-s5}.N10000.s1.accuracy :: sha256 b56bf70be9982044 :: mean +- SD over the row's 5 pretrained models
\newcommand{\FtAccJciiEFiveOnesevenMass}{0.4993\,\ensuremath{\pm}\,0.0022}  % experiments/FIGS/data/w2b_leg1_metrics_v2.json :: cells.{r16q1mass-s1,r16q1mass-s2,r16q1mass-s3,r16q1mass-s4,r16q1mass-s5}.N100000.s1.accuracy :: sha256 b56bf70be9982044 :: mean +- SD over the row's 5 pretrained models
\newcommand{\FtAccJciiESixOnesevenMass}{0.5166\,\ensuremath{\pm}\,0.0013}  % experiments/FIGS/data/w2b_leg1_metrics_v2.json :: cells.{r16q1mass-s1,r16q1mass-s2,r16q1mass-s3,r16q1mass-s4,r16q1mass-s5}.N1000000.s1.accuracy :: sha256 b56bf70be9982044 :: mean +- SD over the row's 5 pretrained models
\newcommand{\FtAccJciiEThreeRandDrawOne}{0.424}  % experiments/FIGS/data/w2b_leg1_metrics_v2.json :: cells.{rand-d1-s1b}.N1000.s1.accuracy :: sha256 b56bf70be9982044 :: one run
\newcommand{\FtAccJciiEFourRandDrawOne}{0.5070}  % experiments/FIGS/data/w2b_leg1_metrics_v2.json :: cells.{rand-d1-s1b}.N10000.s1.accuracy :: sha256 b56bf70be9982044 :: one run
\newcommand{\FtAccJciiEFiveRandDrawOne}{0.5226}  % experiments/FIGS/data/w2b_leg1_metrics_v2.json :: cells.{rand-d1-s1b}.N100000.s1.accuracy :: sha256 b56bf70be9982044 :: one run
\newcommand{\FtAccJciiESixRandDrawOne}{0.53354}  % experiments/FIGS/data/w2b_leg1_metrics_v2.json :: cells.{rand-d1-s1b}.N1000000.s1.accuracy :: sha256 b56bf70be9982044 :: one run
\newcommand{\FtAccJciiEThreeRandDrawTwo}{0.420}  % experiments/FIGS/data/w2b_leg1_metrics_v2.json :: cells.{rand-d2-s2}.N1000.s1.accuracy :: sha256 b56bf70be9982044 :: one run
\newcommand{\FtAccJciiEFourRandDrawTwo}{0.5074}  % experiments/FIGS/data/w2b_leg1_metrics_v2.json :: cells.{rand-d2-s2}.N10000.s1.accuracy :: sha256 b56bf70be9982044 :: one run
\newcommand{\FtAccJciiEFiveRandDrawTwo}{0.5243}  % experiments/FIGS/data/w2b_leg1_metrics_v2.json :: cells.{rand-d2-s2}.N100000.s1.accuracy :: sha256 b56bf70be9982044 :: one run
\newcommand{\FtAccJciiESixRandDrawTwo}{0.53441}  % experiments/FIGS/data/w2b_leg1_metrics_v2.json :: cells.{rand-d2-s2}.N1000000.s1.accuracy :: sha256 b56bf70be9982044 :: one run
\newcommand{\FtAccJciiEThreeRandDrawThree}{0.423}  % experiments/FIGS/data/w2b_leg1_metrics_v2.json :: cells.{rand-d3-s3}.N1000.s1.accuracy :: sha256 b56bf70be9982044 :: one run
\newcommand{\FtAccJciiEFourRandDrawThree}{0.5114}  % experiments/FIGS/data/w2b_leg1_metrics_v2.json :: cells.{rand-d3-s3}.N10000.s1.accuracy :: sha256 b56bf70be9982044 :: one run
\newcommand{\FtAccJciiEFiveRandDrawThree}{0.5245}  % experiments/FIGS/data/w2b_leg1_metrics_v2.json :: cells.{rand-d3-s3}.N100000.s1.accuracy :: sha256 b56bf70be9982044 :: one run
\newcommand{\FtAccJciiESixRandDrawThree}{0.53473}  % experiments/FIGS/data/w2b_leg1_metrics_v2.json :: cells.{rand-d3-s3}.N1000000.s1.accuracy :: sha256 b56bf70be9982044 :: one run
\newcommand{\FtOmaRatioJciiEThreeOnesixtwoOverOneeighteight}{0.96}  % experiments/FIGS/data/w2b_leg1_metrics_v2.json :: cells.*.N1000.s1.macro_auc_ovr :: sha256 b56bf70be9982044 :: ratio of the seed means of 1 - macro AUC, this row over the reference
\newcommand{\FtOmaRatioJciiEThreeFourthreeOverOneeighteight}{0.84}  % experiments/FIGS/data/w2b_leg1_metrics_v2.json :: cells.*.N1000.s1.macro_auc_ovr :: sha256 b56bf70be9982044 :: ratio of the seed means of 1 - macro AUC, this row over the reference
\newcommand{\FtOmaRatioJciiEThreeOnesevenOverOneeighteight}{0.97}  % experiments/FIGS/data/w2b_leg1_metrics_v2.json :: cells.*.N1000.s1.macro_auc_ovr :: sha256 b56bf70be9982044 :: ratio of the seed means of 1 - macro AUC, this row over the reference
\newcommand{\FtOmaRatioJciiEThreeOnesixtwoMassOverOneeighteight}{0.81}  % experiments/FIGS/data/w2b_leg1_metrics_v2.json :: cells.*.N1000.s1.macro_auc_ovr :: sha256 b56bf70be9982044 :: ratio of the seed means of 1 - macro AUC, this row over the reference
\newcommand{\FtOmaRatioJciiEThreeOnesevenMassOverOneeighteight}{0.98}  % experiments/FIGS/data/w2b_leg1_metrics_v2.json :: cells.*.N1000.s1.macro_auc_ovr :: sha256 b56bf70be9982044 :: ratio of the seed means of 1 - macro AUC, this row over the reference
\newcommand{\FtOmaRatioJciiEThreeRandDrawOneOverOneeighteight}{0.80}  % experiments/FIGS/data/w2b_leg1_metrics_v2.json :: cells.*.N1000.s1.macro_auc_ovr :: sha256 b56bf70be9982044 :: ratio of the seed means of 1 - macro AUC, this row over the reference
\newcommand{\FtOmaRatioJciiEThreeRandDrawTwoOverOneeighteight}{0.87}  % experiments/FIGS/data/w2b_leg1_metrics_v2.json :: cells.*.N1000.s1.macro_auc_ovr :: sha256 b56bf70be9982044 :: ratio of the seed means of 1 - macro AUC, this row over the reference
\newcommand{\FtOmaRatioJciiEThreeRandDrawThreeOverOneeighteight}{0.83}  % experiments/FIGS/data/w2b_leg1_metrics_v2.json :: cells.*.N1000.s1.macro_auc_ovr :: sha256 b56bf70be9982044 :: ratio of the seed means of 1 - macro AUC, this row over the reference
\newcommand{\FtOmaRatioJciiEThreeOneeighteightOverOnesixtwo}{1.05}  % experiments/FIGS/data/w2b_leg1_metrics_v2.json :: cells.*.N1000.s1.macro_auc_ovr :: sha256 b56bf70be9982044 :: ratio of the seed means of 1 - macro AUC, this row over the reference
\newcommand{\FtOmaRatioJciiEThreeFourthreeOverOnesixtwo}{0.88}  % experiments/FIGS/data/w2b_leg1_metrics_v2.json :: cells.*.N1000.s1.macro_auc_ovr :: sha256 b56bf70be9982044 :: ratio of the seed means of 1 - macro AUC, this row over the reference
\newcommand{\FtOmaRatioJciiEThreeOnesevenOverOnesixtwo}{1.02}  % experiments/FIGS/data/w2b_leg1_metrics_v2.json :: cells.*.N1000.s1.macro_auc_ovr :: sha256 b56bf70be9982044 :: ratio of the seed means of 1 - macro AUC, this row over the reference
\newcommand{\FtOmaRatioJciiEThreeOnesixtwoMassOverOnesixtwo}{0.85}  % experiments/FIGS/data/w2b_leg1_metrics_v2.json :: cells.*.N1000.s1.macro_auc_ovr :: sha256 b56bf70be9982044 :: ratio of the seed means of 1 - macro AUC, this row over the reference
\newcommand{\FtOmaRatioJciiEThreeOnesevenMassOverOnesixtwo}{1.02}  % experiments/FIGS/data/w2b_leg1_metrics_v2.json :: cells.*.N1000.s1.macro_auc_ovr :: sha256 b56bf70be9982044 :: ratio of the seed means of 1 - macro AUC, this row over the reference
\newcommand{\FtOmaRatioJciiEThreeRandDrawOneOverOnesixtwo}{0.84}  % experiments/FIGS/data/w2b_leg1_metrics_v2.json :: cells.*.N1000.s1.macro_auc_ovr :: sha256 b56bf70be9982044 :: ratio of the seed means of 1 - macro AUC, this row over the reference
\newcommand{\FtOmaRatioJciiEThreeRandDrawTwoOverOnesixtwo}{0.91}  % experiments/FIGS/data/w2b_leg1_metrics_v2.json :: cells.*.N1000.s1.macro_auc_ovr :: sha256 b56bf70be9982044 :: ratio of the seed means of 1 - macro AUC, this row over the reference
\newcommand{\FtOmaRatioJciiEThreeRandDrawThreeOverOnesixtwo}{0.87}  % experiments/FIGS/data/w2b_leg1_metrics_v2.json :: cells.*.N1000.s1.macro_auc_ovr :: sha256 b56bf70be9982044 :: ratio of the seed means of 1 - macro AUC, this row over the reference
\newcommand{\FtOmaRatioJciiEFourOnesixtwoOverOneeighteight}{0.99}  % experiments/FIGS/data/w2b_leg1_metrics_v2.json :: cells.*.N10000.s1.macro_auc_ovr :: sha256 b56bf70be9982044 :: ratio of the seed means of 1 - macro AUC, this row over the reference
\newcommand{\FtOmaRatioJciiEFourFourthreeOverOneeighteight}{1.20}  % experiments/FIGS/data/w2b_leg1_metrics_v2.json :: cells.*.N10000.s1.macro_auc_ovr :: sha256 b56bf70be9982044 :: ratio of the seed means of 1 - macro AUC, this row over the reference
\newcommand{\FtOmaRatioJciiEFourOnesevenOverOneeighteight}{1.32}  % experiments/FIGS/data/w2b_leg1_metrics_v2.json :: cells.*.N10000.s1.macro_auc_ovr :: sha256 b56bf70be9982044 :: ratio of the seed means of 1 - macro AUC, this row over the reference
\newcommand{\FtOmaRatioJciiEFourOnesixtwoMassOverOneeighteight}{0.99}  % experiments/FIGS/data/w2b_leg1_metrics_v2.json :: cells.*.N10000.s1.macro_auc_ovr :: sha256 b56bf70be9982044 :: ratio of the seed means of 1 - macro AUC, this row over the reference
\newcommand{\FtOmaRatioJciiEFourOnesevenMassOverOneeighteight}{1.30}  % experiments/FIGS/data/w2b_leg1_metrics_v2.json :: cells.*.N10000.s1.macro_auc_ovr :: sha256 b56bf70be9982044 :: ratio of the seed means of 1 - macro AUC, this row over the reference
\newcommand{\FtOmaRatioJciiEFourRandDrawOneOverOneeighteight}{1.17}  % experiments/FIGS/data/w2b_leg1_metrics_v2.json :: cells.*.N10000.s1.macro_auc_ovr :: sha256 b56bf70be9982044 :: ratio of the seed means of 1 - macro AUC, this row over the reference
\newcommand{\FtOmaRatioJciiEFourRandDrawTwoOverOneeighteight}{1.17}  % experiments/FIGS/data/w2b_leg1_metrics_v2.json :: cells.*.N10000.s1.macro_auc_ovr :: sha256 b56bf70be9982044 :: ratio of the seed means of 1 - macro AUC, this row over the reference
\newcommand{\FtOmaRatioJciiEFourRandDrawThreeOverOneeighteight}{1.13}  % experiments/FIGS/data/w2b_leg1_metrics_v2.json :: cells.*.N10000.s1.macro_auc_ovr :: sha256 b56bf70be9982044 :: ratio of the seed means of 1 - macro AUC, this row over the reference
\newcommand{\FtOmaRatioJciiEFourOneeighteightOverOnesixtwo}{1.01}  % experiments/FIGS/data/w2b_leg1_metrics_v2.json :: cells.*.N10000.s1.macro_auc_ovr :: sha256 b56bf70be9982044 :: ratio of the seed means of 1 - macro AUC, this row over the reference
\newcommand{\FtOmaRatioJciiEFourFourthreeOverOnesixtwo}{1.21}  % experiments/FIGS/data/w2b_leg1_metrics_v2.json :: cells.*.N10000.s1.macro_auc_ovr :: sha256 b56bf70be9982044 :: ratio of the seed means of 1 - macro AUC, this row over the reference
\newcommand{\FtOmaRatioJciiEFourOnesevenOverOnesixtwo}{1.32}  % experiments/FIGS/data/w2b_leg1_metrics_v2.json :: cells.*.N10000.s1.macro_auc_ovr :: sha256 b56bf70be9982044 :: ratio of the seed means of 1 - macro AUC, this row over the reference
\newcommand{\FtOmaRatioJciiEFourOnesixtwoMassOverOnesixtwo}{1.00}  % experiments/FIGS/data/w2b_leg1_metrics_v2.json :: cells.*.N10000.s1.macro_auc_ovr :: sha256 b56bf70be9982044 :: ratio of the seed means of 1 - macro AUC, this row over the reference
\newcommand{\FtOmaRatioJciiEFourOnesevenMassOverOnesixtwo}{1.31}  % experiments/FIGS/data/w2b_leg1_metrics_v2.json :: cells.*.N10000.s1.macro_auc_ovr :: sha256 b56bf70be9982044 :: ratio of the seed means of 1 - macro AUC, this row over the reference
\newcommand{\FtOmaRatioJciiEFourRandDrawOneOverOnesixtwo}{1.18}  % experiments/FIGS/data/w2b_leg1_metrics_v2.json :: cells.*.N10000.s1.macro_auc_ovr :: sha256 b56bf70be9982044 :: ratio of the seed means of 1 - macro AUC, this row over the reference
\newcommand{\FtOmaRatioJciiEFourRandDrawTwoOverOnesixtwo}{1.18}  % experiments/FIGS/data/w2b_leg1_metrics_v2.json :: cells.*.N10000.s1.macro_auc_ovr :: sha256 b56bf70be9982044 :: ratio of the seed means of 1 - macro AUC, this row over the reference
\newcommand{\FtOmaRatioJciiEFourRandDrawThreeOverOnesixtwo}{1.14}  % experiments/FIGS/data/w2b_leg1_metrics_v2.json :: cells.*.N10000.s1.macro_auc_ovr :: sha256 b56bf70be9982044 :: ratio of the seed means of 1 - macro AUC, this row over the reference
\newcommand{\FtOmaRatioJciiEFiveOnesixtwoOverOneeighteight}{1.00}  % experiments/FIGS/data/w2b_leg1_metrics_v2.json :: cells.*.N100000.s1.macro_auc_ovr :: sha256 b56bf70be9982044 :: ratio of the seed means of 1 - macro AUC, this row over the reference
\newcommand{\FtOmaRatioJciiEFiveFourthreeOverOneeighteight}{1.09}  % experiments/FIGS/data/w2b_leg1_metrics_v2.json :: cells.*.N100000.s1.macro_auc_ovr :: sha256 b56bf70be9982044 :: ratio of the seed means of 1 - macro AUC, this row over the reference
\newcommand{\FtOmaRatioJciiEFiveOnesevenOverOneeighteight}{1.22}  % experiments/FIGS/data/w2b_leg1_metrics_v2.json :: cells.*.N100000.s1.macro_auc_ovr :: sha256 b56bf70be9982044 :: ratio of the seed means of 1 - macro AUC, this row over the reference
\newcommand{\FtOmaRatioJciiEFiveOnesixtwoMassOverOneeighteight}{1.01}  % experiments/FIGS/data/w2b_leg1_metrics_v2.json :: cells.*.N100000.s1.macro_auc_ovr :: sha256 b56bf70be9982044 :: ratio of the seed means of 1 - macro AUC, this row over the reference
\newcommand{\FtOmaRatioJciiEFiveOnesevenMassOverOneeighteight}{1.26}  % experiments/FIGS/data/w2b_leg1_metrics_v2.json :: cells.*.N100000.s1.macro_auc_ovr :: sha256 b56bf70be9982044 :: ratio of the seed means of 1 - macro AUC, this row over the reference
\newcommand{\FtOmaRatioJciiEFiveRandDrawOneOverOneeighteight}{1.15}  % experiments/FIGS/data/w2b_leg1_metrics_v2.json :: cells.*.N100000.s1.macro_auc_ovr :: sha256 b56bf70be9982044 :: ratio of the seed means of 1 - macro AUC, this row over the reference
\newcommand{\FtOmaRatioJciiEFiveRandDrawTwoOverOneeighteight}{1.12}  % experiments/FIGS/data/w2b_leg1_metrics_v2.json :: cells.*.N100000.s1.macro_auc_ovr :: sha256 b56bf70be9982044 :: ratio of the seed means of 1 - macro AUC, this row over the reference
\newcommand{\FtOmaRatioJciiEFiveRandDrawThreeOverOneeighteight}{1.15}  % experiments/FIGS/data/w2b_leg1_metrics_v2.json :: cells.*.N100000.s1.macro_auc_ovr :: sha256 b56bf70be9982044 :: ratio of the seed means of 1 - macro AUC, this row over the reference
\newcommand{\FtOmaRatioJciiEFiveOneeighteightOverOnesixtwo}{1.00}  % experiments/FIGS/data/w2b_leg1_metrics_v2.json :: cells.*.N100000.s1.macro_auc_ovr :: sha256 b56bf70be9982044 :: ratio of the seed means of 1 - macro AUC, this row over the reference
\newcommand{\FtOmaRatioJciiEFiveFourthreeOverOnesixtwo}{1.09}  % experiments/FIGS/data/w2b_leg1_metrics_v2.json :: cells.*.N100000.s1.macro_auc_ovr :: sha256 b56bf70be9982044 :: ratio of the seed means of 1 - macro AUC, this row over the reference
\newcommand{\FtOmaRatioJciiEFiveOnesevenOverOnesixtwo}{1.22}  % experiments/FIGS/data/w2b_leg1_metrics_v2.json :: cells.*.N100000.s1.macro_auc_ovr :: sha256 b56bf70be9982044 :: ratio of the seed means of 1 - macro AUC, this row over the reference
\newcommand{\FtOmaRatioJciiEFiveOnesixtwoMassOverOnesixtwo}{1.02}  % experiments/FIGS/data/w2b_leg1_metrics_v2.json :: cells.*.N100000.s1.macro_auc_ovr :: sha256 b56bf70be9982044 :: ratio of the seed means of 1 - macro AUC, this row over the reference
\newcommand{\FtOmaRatioJciiEFiveOnesevenMassOverOnesixtwo}{1.26}  % experiments/FIGS/data/w2b_leg1_metrics_v2.json :: cells.*.N100000.s1.macro_auc_ovr :: sha256 b56bf70be9982044 :: ratio of the seed means of 1 - macro AUC, this row over the reference
\newcommand{\FtOmaRatioJciiEFiveRandDrawOneOverOnesixtwo}{1.15}  % experiments/FIGS/data/w2b_leg1_metrics_v2.json :: cells.*.N100000.s1.macro_auc_ovr :: sha256 b56bf70be9982044 :: ratio of the seed means of 1 - macro AUC, this row over the reference
\newcommand{\FtOmaRatioJciiEFiveRandDrawTwoOverOnesixtwo}{1.13}  % experiments/FIGS/data/w2b_leg1_metrics_v2.json :: cells.*.N100000.s1.macro_auc_ovr :: sha256 b56bf70be9982044 :: ratio of the seed means of 1 - macro AUC, this row over the reference
\newcommand{\FtOmaRatioJciiEFiveRandDrawThreeOverOnesixtwo}{1.16}  % experiments/FIGS/data/w2b_leg1_metrics_v2.json :: cells.*.N100000.s1.macro_auc_ovr :: sha256 b56bf70be9982044 :: ratio of the seed means of 1 - macro AUC, this row over the reference
\newcommand{\FtOmaRatioJciiESixOnesixtwoOverOneeighteight}{1.00}  % experiments/FIGS/data/w2b_leg1_metrics_v2.json :: cells.*.N1000000.s1.macro_auc_ovr :: sha256 b56bf70be9982044 :: ratio of the seed means of 1 - macro AUC, this row over the reference
\newcommand{\FtOmaRatioJciiESixFourthreeOverOneeighteight}{1.05}  % experiments/FIGS/data/w2b_leg1_metrics_v2.json :: cells.*.N1000000.s1.macro_auc_ovr :: sha256 b56bf70be9982044 :: ratio of the seed means of 1 - macro AUC, this row over the reference
\newcommand{\FtOmaRatioJciiESixOnesevenOverOneeighteight}{1.14}  % experiments/FIGS/data/w2b_leg1_metrics_v2.json :: cells.*.N1000000.s1.macro_auc_ovr :: sha256 b56bf70be9982044 :: ratio of the seed means of 1 - macro AUC, this row over the reference
\newcommand{\FtOmaRatioJciiESixOnesixtwoMassOverOneeighteight}{1.02}  % experiments/FIGS/data/w2b_leg1_metrics_v2.json :: cells.*.N1000000.s1.macro_auc_ovr :: sha256 b56bf70be9982044 :: ratio of the seed means of 1 - macro AUC, this row over the reference
\newcommand{\FtOmaRatioJciiESixOnesevenMassOverOneeighteight}{1.18}  % experiments/FIGS/data/w2b_leg1_metrics_v2.json :: cells.*.N1000000.s1.macro_auc_ovr :: sha256 b56bf70be9982044 :: ratio of the seed means of 1 - macro AUC, this row over the reference
\newcommand{\FtOmaRatioJciiESixRandDrawOneOverOneeighteight}{1.08}  % experiments/FIGS/data/w2b_leg1_metrics_v2.json :: cells.*.N1000000.s1.macro_auc_ovr :: sha256 b56bf70be9982044 :: ratio of the seed means of 1 - macro AUC, this row over the reference
\newcommand{\FtOmaRatioJciiESixRandDrawTwoOverOneeighteight}{1.07}  % experiments/FIGS/data/w2b_leg1_metrics_v2.json :: cells.*.N1000000.s1.macro_auc_ovr :: sha256 b56bf70be9982044 :: ratio of the seed means of 1 - macro AUC, this row over the reference
\newcommand{\FtOmaRatioJciiESixRandDrawThreeOverOneeighteight}{1.07}  % experiments/FIGS/data/w2b_leg1_metrics_v2.json :: cells.*.N1000000.s1.macro_auc_ovr :: sha256 b56bf70be9982044 :: ratio of the seed means of 1 - macro AUC, this row over the reference
\newcommand{\FtOmaRatioJciiESixOneeighteightOverOnesixtwo}{1.00}  % experiments/FIGS/data/w2b_leg1_metrics_v2.json :: cells.*.N1000000.s1.macro_auc_ovr :: sha256 b56bf70be9982044 :: ratio of the seed means of 1 - macro AUC, this row over the reference
\newcommand{\FtOmaRatioJciiESixFourthreeOverOnesixtwo}{1.04}  % experiments/FIGS/data/w2b_leg1_metrics_v2.json :: cells.*.N1000000.s1.macro_auc_ovr :: sha256 b56bf70be9982044 :: ratio of the seed means of 1 - macro AUC, this row over the reference
\newcommand{\FtOmaRatioJciiESixOnesevenOverOnesixtwo}{1.14}  % experiments/FIGS/data/w2b_leg1_metrics_v2.json :: cells.*.N1000000.s1.macro_auc_ovr :: sha256 b56bf70be9982044 :: ratio of the seed means of 1 - macro AUC, this row over the reference
\newcommand{\FtOmaRatioJciiESixOnesixtwoMassOverOnesixtwo}{1.02}  % experiments/FIGS/data/w2b_leg1_metrics_v2.json :: cells.*.N1000000.s1.macro_auc_ovr :: sha256 b56bf70be9982044 :: ratio of the seed means of 1 - macro AUC, this row over the reference
\newcommand{\FtOmaRatioJciiESixOnesevenMassOverOnesixtwo}{1.18}  % experiments/FIGS/data/w2b_leg1_metrics_v2.json :: cells.*.N1000000.s1.macro_auc_ovr :: sha256 b56bf70be9982044 :: ratio of the seed means of 1 - macro AUC, this row over the reference
\newcommand{\FtOmaRatioJciiESixRandDrawOneOverOnesixtwo}{1.07}  % experiments/FIGS/data/w2b_leg1_metrics_v2.json :: cells.*.N1000000.s1.macro_auc_ovr :: sha256 b56bf70be9982044 :: ratio of the seed means of 1 - macro AUC, this row over the reference
\newcommand{\FtOmaRatioJciiESixRandDrawTwoOverOnesixtwo}{1.07}  % experiments/FIGS/data/w2b_leg1_metrics_v2.json :: cells.*.N1000000.s1.macro_auc_ovr :: sha256 b56bf70be9982044 :: ratio of the seed means of 1 - macro AUC, this row over the reference
\newcommand{\FtOmaRatioJciiESixRandDrawThreeOverOnesixtwo}{1.07}  % experiments/FIGS/data/w2b_leg1_metrics_v2.json :: cells.*.N1000000.s1.macro_auc_ovr :: sha256 b56bf70be9982044 :: ratio of the seed means of 1 - macro AUC, this row over the reference
\newcommand{\PairedFtAccJciiEThreeOnesevenMinusOneeighteight}{\ensuremath{-}0.1132~[\ensuremath{-}0.1214, \ensuremath{-}0.1051]}  % experiments/FIGS/data/w2b_leg1_metrics_v2.json :: cells.{r16q1-s1,r16q1-s2,r16q1-s3,r16q1-s4,r16q1-s5,l188-s1,l188-s2,l188-s3,l188-s4,l188-s5}.N1000.s1.accuracy :: sha256 b56bf70be9982044 :: accuracy, coarse minus fine, paired by run [95% Student-t interval]
\newcommand{\PairedFtAccJciiEFourOnesevenMinusOneeighteight}{\ensuremath{-}0.0730~[\ensuremath{-}0.0902, \ensuremath{-}0.0558]}  % experiments/FIGS/data/w2b_leg1_metrics_v2.json :: cells.{r16q1-s1,r16q1-s2,r16q1-s3,r16q1-s4,r16q1-s5,l188-s1,l188-s2,l188-s3,l188-s4,l188-s5}.N10000.s1.accuracy :: sha256 b56bf70be9982044 :: accuracy, coarse minus fine, paired by run [95% Student-t interval]
\newcommand{\PairedFtAccJciiEFiveOnesevenMinusOneeighteight}{\ensuremath{-}0.0301~[\ensuremath{-}0.0319, \ensuremath{-}0.0284]}  % experiments/FIGS/data/w2b_leg1_metrics_v2.json :: cells.{r16q1-s1,r16q1-s2,r16q1-s3,r16q1-s4,r16q1-s5,l188-s1,l188-s2,l188-s3,l188-s4,l188-s5}.N100000.s1.accuracy :: sha256 b56bf70be9982044 :: accuracy, coarse minus fine, paired by run [95% Student-t interval]
\newcommand{\PairedFtAccJciiESixOnesevenMinusOneeighteight}{\ensuremath{-}0.0216~[\ensuremath{-}0.0232, \ensuremath{-}0.0199]}  % experiments/FIGS/data/w2b_leg1_metrics_v2.json :: cells.{r16q1-s1,r16q1-s2,r16q1-s3,r16q1-s4,r16q1-s5,l188-s1,l188-s2,l188-s3,l188-s4,l188-s5}.N1000000.s1.accuracy :: sha256 b56bf70be9982044 :: accuracy, coarse minus fine, paired by run [95% Student-t interval]
\newcommand{\FtNClassesJcii}{162}  % experiments/FIGS/data/w2b_leg1_metrics_v2.json :: cells.*.*.s1.n_classes_present :: sha256 b56bf70be9982044 :: classes of the fine-tuning task
\newcommand{\FtNTestAucJcii}{125{,}000}  % experiments/FIGS/data/w2b_leg1_metrics_v2.json :: cells.*.*.s1.n_jets_auc :: sha256 b56bf70be9982044 :: test jets the macro AUC is computed on
\newcommand{\FtNTestAccJcii}{500{,}000}  % experiments/FIGS/data/w2b_leg1_metrics_v2.json :: cells.*.*.s1.n_jets :: sha256 b56bf70be9982044 :: test jets the accuracy is computed on
\newcommand{\FtAucJcEThreeOneeighteight}{0.950\,\ensuremath{\pm}\,0.004}  % experiments/FIGS/data/w2b_leg2_metrics_v2.json :: cells.{l188-s1,l188-s2,l188-s3,l188-s4,l188-s5}.N1000.s1.macro_auc_ovr :: sha256 9cd059d3e2f51798 :: mean +- SD over the row's 5 pretrained models
\newcommand{\FtAucJcEFourOneeighteight}{0.9765\,\ensuremath{\pm}\,0.0008}  % experiments/FIGS/data/w2b_leg2_metrics_v2.json :: cells.{l188-s1,l188-s2,l188-s3,l188-s4,l188-s5}.N10000.s1.macro_auc_ovr :: sha256 9cd059d3e2f51798 :: mean +- SD over the row's 5 pretrained models
\newcommand{\FtAucJcEFiveOneeighteight}{0.9833\,\ensuremath{\pm}\,0.0002}  % experiments/FIGS/data/w2b_leg2_metrics_v2.json :: cells.{l188-s1,l188-s2,l188-s3,l188-s4,l188-s5}.N100000.s1.macro_auc_ovr :: sha256 9cd059d3e2f51798 :: mean +- SD over the row's 5 pretrained models
\newcommand{\FtAucJcESixOneeighteight}{0.9858\,\ensuremath{\pm}\,0.0001}  % experiments/FIGS/data/w2b_leg2_metrics_v2.json :: cells.{l188-s1,l188-s2,l188-s3,l188-s4,l188-s5}.N1000000.s1.macro_auc_ovr :: sha256 9cd059d3e2f51798 :: mean +- SD over the row's 5 pretrained models
\newcommand{\FtAucJcEThreeOnesixtwo}{0.9526\,\ensuremath{\pm}\,0.0020}  % experiments/FIGS/data/w2b_leg2_metrics_v2.json :: cells.{l162-s1b,l162-s2,l162-s3,l162-s4,l162-s5}.N1000.s1.macro_auc_ovr :: sha256 9cd059d3e2f51798 :: mean +- SD over the row's 5 pretrained models
\newcommand{\FtAucJcEFourOnesixtwo}{0.9770\,\ensuremath{\pm}\,0.0007}  % experiments/FIGS/data/w2b_leg2_metrics_v2.json :: cells.{l162-s1b,l162-s2,l162-s3,l162-s4,l162-s5}.N10000.s1.macro_auc_ovr :: sha256 9cd059d3e2f51798 :: mean +- SD over the row's 5 pretrained models
\newcommand{\FtAucJcEFiveOnesixtwo}{0.9836\,\ensuremath{\pm}\,0.0001}  % experiments/FIGS/data/w2b_leg2_metrics_v2.json :: cells.{l162-s1b,l162-s2,l162-s3,l162-s4,l162-s5}.N100000.s1.macro_auc_ovr :: sha256 9cd059d3e2f51798 :: mean +- SD over the row's 5 pretrained models
\newcommand{\FtAucJcESixOnesixtwo}{0.9858\,\ensuremath{\pm}\,0.0001}  % experiments/FIGS/data/w2b_leg2_metrics_v2.json :: cells.{l162-s1b,l162-s2,l162-s3,l162-s4,l162-s5}.N1000000.s1.macro_auc_ovr :: sha256 9cd059d3e2f51798 :: mean +- SD over the row's 5 pretrained models
\newcommand{\FtAucJcEThreeFourthree}{0.950\,\ensuremath{\pm}\,0.004}  % experiments/FIGS/data/w2b_leg2_metrics_v2.json :: cells.{r42q1-s1,r42q1-s2,r42q1-s3,r42q1-s4,r42q1-s5}.N1000.s1.macro_auc_ovr :: sha256 9cd059d3e2f51798 :: mean +- SD over the row's 5 pretrained models
\newcommand{\FtAucJcEFourFourthree}{0.9769\,\ensuremath{\pm}\,0.0005}  % experiments/FIGS/data/w2b_leg2_metrics_v2.json :: cells.{r42q1-s1,r42q1-s2,r42q1-s3,r42q1-s4,r42q1-s5}.N10000.s1.macro_auc_ovr :: sha256 9cd059d3e2f51798 :: mean +- SD over the row's 5 pretrained models
\newcommand{\FtAucJcEFiveFourthree}{0.9831\,\ensuremath{\pm}\,0.0002}  % experiments/FIGS/data/w2b_leg2_metrics_v2.json :: cells.{r42q1-s1,r42q1-s2,r42q1-s3,r42q1-s4,r42q1-s5}.N100000.s1.macro_auc_ovr :: sha256 9cd059d3e2f51798 :: mean +- SD over the row's 5 pretrained models
\newcommand{\FtAucJcESixFourthree}{0.9855\,\ensuremath{\pm}\,0.0001}  % experiments/FIGS/data/w2b_leg2_metrics_v2.json :: cells.{r42q1-s1,r42q1-s2,r42q1-s3,r42q1-s4,r42q1-s5}.N1000000.s1.macro_auc_ovr :: sha256 9cd059d3e2f51798 :: mean +- SD over the row's 5 pretrained models
\newcommand{\FtAucJcEThreeOneseven}{0.9470\,\ensuremath{\pm}\,0.0024}  % experiments/FIGS/data/w2b_leg2_metrics_v2.json :: cells.{r16q1-s1,r16q1-s2,r16q1-s3,r16q1-s4,r16q1-s5}.N1000.s1.macro_auc_ovr :: sha256 9cd059d3e2f51798 :: mean +- SD over the row's 5 pretrained models
\newcommand{\FtAucJcEFourOneseven}{0.9749\,\ensuremath{\pm}\,0.0006}  % experiments/FIGS/data/w2b_leg2_metrics_v2.json :: cells.{r16q1-s1,r16q1-s2,r16q1-s3,r16q1-s4,r16q1-s5}.N10000.s1.macro_auc_ovr :: sha256 9cd059d3e2f51798 :: mean +- SD over the row's 5 pretrained models
\newcommand{\FtAucJcEFiveOneseven}{0.9821\,\ensuremath{\pm}\,\ensuremath{<}0.0001}  % experiments/FIGS/data/w2b_leg2_metrics_v2.json :: cells.{r16q1-s1,r16q1-s2,r16q1-s3,r16q1-s4,r16q1-s5}.N100000.s1.macro_auc_ovr :: sha256 9cd059d3e2f51798 :: mean +- SD over the row's 5 pretrained models
\newcommand{\FtAucJcESixOneseven}{0.9846\,\ensuremath{\pm}\,0.0001}  % experiments/FIGS/data/w2b_leg2_metrics_v2.json :: cells.{r16q1-s1,r16q1-s2,r16q1-s3,r16q1-s4,r16q1-s5}.N1000000.s1.macro_auc_ovr :: sha256 9cd059d3e2f51798 :: mean +- SD over the row's 5 pretrained models
\newcommand{\FtAucJcEThreeOnesixtwoMass}{0.9569\,\ensuremath{\pm}\,0.0016}  % experiments/FIGS/data/w2b_leg2_metrics_v2.json :: cells.{l162mass-s1,l162mass-s2,l162mass-s3,l162mass-s4,l162mass-s5}.N1000.s1.macro_auc_ovr :: sha256 9cd059d3e2f51798 :: mean +- SD over the row's 5 pretrained models
\newcommand{\FtAucJcEFourOnesixtwoMass}{0.9768\,\ensuremath{\pm}\,0.0005}  % experiments/FIGS/data/w2b_leg2_metrics_v2.json :: cells.{l162mass-s1,l162mass-s2,l162mass-s3,l162mass-s4,l162mass-s5}.N10000.s1.macro_auc_ovr :: sha256 9cd059d3e2f51798 :: mean +- SD over the row's 5 pretrained models
\newcommand{\FtAucJcEFiveOnesixtwoMass}{0.9834\,\ensuremath{\pm}\,0.0001}  % experiments/FIGS/data/w2b_leg2_metrics_v2.json :: cells.{l162mass-s1,l162mass-s2,l162mass-s3,l162mass-s4,l162mass-s5}.N100000.s1.macro_auc_ovr :: sha256 9cd059d3e2f51798 :: mean +- SD over the row's 5 pretrained models
\newcommand{\FtAucJcESixOnesixtwoMass}{0.9856\,\ensuremath{\pm}\,0.0001}  % experiments/FIGS/data/w2b_leg2_metrics_v2.json :: cells.{l162mass-s1,l162mass-s2,l162mass-s3,l162mass-s4,l162mass-s5}.N1000000.s1.macro_auc_ovr :: sha256 9cd059d3e2f51798 :: mean +- SD over the row's 5 pretrained models
\newcommand{\FtAucJcEThreeOnesevenMass}{0.9460\,\ensuremath{\pm}\,0.0029}  % experiments/FIGS/data/w2b_leg2_metrics_v2.json :: cells.{r16q1mass-s1,r16q1mass-s2,r16q1mass-s3,r16q1mass-s4,r16q1mass-s5}.N1000.s1.macro_auc_ovr :: sha256 9cd059d3e2f51798 :: mean +- SD over the row's 5 pretrained models
\newcommand{\FtAucJcEFourOnesevenMass}{0.9737\,\ensuremath{\pm}\,0.0008}  % experiments/FIGS/data/w2b_leg2_metrics_v2.json :: cells.{r16q1mass-s1,r16q1mass-s2,r16q1mass-s3,r16q1mass-s4,r16q1mass-s5}.N10000.s1.macro_auc_ovr :: sha256 9cd059d3e2f51798 :: mean +- SD over the row's 5 pretrained models
\newcommand{\FtAucJcEFiveOnesevenMass}{0.9817\,\ensuremath{\pm}\,0.0002}  % experiments/FIGS/data/w2b_leg2_metrics_v2.json :: cells.{r16q1mass-s1,r16q1mass-s2,r16q1mass-s3,r16q1mass-s4,r16q1mass-s5}.N100000.s1.macro_auc_ovr :: sha256 9cd059d3e2f51798 :: mean +- SD over the row's 5 pretrained models
\newcommand{\FtAucJcESixOnesevenMass}{0.9842\,\ensuremath{\pm}\,0.0002}  % experiments/FIGS/data/w2b_leg2_metrics_v2.json :: cells.{r16q1mass-s1,r16q1mass-s2,r16q1mass-s3,r16q1mass-s4,r16q1mass-s5}.N1000000.s1.macro_auc_ovr :: sha256 9cd059d3e2f51798 :: mean +- SD over the row's 5 pretrained models
\newcommand{\FtAucJcEThreeRandDrawOne}{0.9477}  % experiments/FIGS/data/w2b_leg2_metrics_v2.json :: cells.{rand-d1-s1b}.N1000.s1.macro_auc_ovr :: sha256 9cd059d3e2f51798 :: one run
\newcommand{\FtAucJcEFourRandDrawOne}{0.9745}  % experiments/FIGS/data/w2b_leg2_metrics_v2.json :: cells.{rand-d1-s1b}.N10000.s1.macro_auc_ovr :: sha256 9cd059d3e2f51798 :: one run
\newcommand{\FtAucJcEFiveRandDrawOne}{0.98265}  % experiments/FIGS/data/w2b_leg2_metrics_v2.json :: cells.{rand-d1-s1b}.N100000.s1.macro_auc_ovr :: sha256 9cd059d3e2f51798 :: one run
\newcommand{\FtAucJcESixRandDrawOne}{0.98513}  % experiments/FIGS/data/w2b_leg2_metrics_v2.json :: cells.{rand-d1-s1b}.N1000000.s1.macro_auc_ovr :: sha256 9cd059d3e2f51798 :: one run
\newcommand{\FtAucJcEThreeRandDrawTwo}{0.9501}  % experiments/FIGS/data/w2b_leg2_metrics_v2.json :: cells.{rand-d2-s2}.N1000.s1.macro_auc_ovr :: sha256 9cd059d3e2f51798 :: one run
\newcommand{\FtAucJcEFourRandDrawTwo}{0.9756}  % experiments/FIGS/data/w2b_leg2_metrics_v2.json :: cells.{rand-d2-s2}.N10000.s1.macro_auc_ovr :: sha256 9cd059d3e2f51798 :: one run
\newcommand{\FtAucJcEFiveRandDrawTwo}{0.98278}  % experiments/FIGS/data/w2b_leg2_metrics_v2.json :: cells.{rand-d2-s2}.N100000.s1.macro_auc_ovr :: sha256 9cd059d3e2f51798 :: one run
\newcommand{\FtAucJcESixRandDrawTwo}{0.98512}  % experiments/FIGS/data/w2b_leg2_metrics_v2.json :: cells.{rand-d2-s2}.N1000000.s1.macro_auc_ovr :: sha256 9cd059d3e2f51798 :: one run
\newcommand{\FtAucJcEThreeRandDrawThree}{0.9470}  % experiments/FIGS/data/w2b_leg2_metrics_v2.json :: cells.{rand-d3-s3}.N1000.s1.macro_auc_ovr :: sha256 9cd059d3e2f51798 :: one run
\newcommand{\FtAucJcEFourRandDrawThree}{0.9743}  % experiments/FIGS/data/w2b_leg2_metrics_v2.json :: cells.{rand-d3-s3}.N10000.s1.macro_auc_ovr :: sha256 9cd059d3e2f51798 :: one run
\newcommand{\FtAucJcEFiveRandDrawThree}{0.98253}  % experiments/FIGS/data/w2b_leg2_metrics_v2.json :: cells.{rand-d3-s3}.N100000.s1.macro_auc_ovr :: sha256 9cd059d3e2f51798 :: one run
\newcommand{\FtAucJcESixRandDrawThree}{0.98506}  % experiments/FIGS/data/w2b_leg2_metrics_v2.json :: cells.{rand-d3-s3}.N1000000.s1.macro_auc_ovr :: sha256 9cd059d3e2f51798 :: one run
\newcommand{\FtAccJcEThreeOneeighteight}{0.7183\,\ensuremath{\pm}\,0.0032}  % experiments/FIGS/data/w2b_leg2_metrics_v2.json :: cells.{l188-s1,l188-s2,l188-s3,l188-s4,l188-s5}.N1000.s1.accuracy :: sha256 9cd059d3e2f51798 :: mean +- SD over the row's 5 pretrained models
\newcommand{\FtAccJcEFourOneeighteight}{0.8039\,\ensuremath{\pm}\,0.0030}  % experiments/FIGS/data/w2b_leg2_metrics_v2.json :: cells.{l188-s1,l188-s2,l188-s3,l188-s4,l188-s5}.N10000.s1.accuracy :: sha256 9cd059d3e2f51798 :: mean +- SD over the row's 5 pretrained models
\newcommand{\FtAccJcEFiveOneeighteight}{0.8351\,\ensuremath{\pm}\,0.0008}  % experiments/FIGS/data/w2b_leg2_metrics_v2.json :: cells.{l188-s1,l188-s2,l188-s3,l188-s4,l188-s5}.N100000.s1.accuracy :: sha256 9cd059d3e2f51798 :: mean +- SD over the row's 5 pretrained models
\newcommand{\FtAccJcESixOneeighteight}{0.8486\,\ensuremath{\pm}\,0.0005}  % experiments/FIGS/data/w2b_leg2_metrics_v2.json :: cells.{l188-s1,l188-s2,l188-s3,l188-s4,l188-s5}.N1000000.s1.accuracy :: sha256 9cd059d3e2f51798 :: mean +- SD over the row's 5 pretrained models
\newcommand{\FtAccJcEThreeOnesixtwo}{0.7268\,\ensuremath{\pm}\,0.0017}  % experiments/FIGS/data/w2b_leg2_metrics_v2.json :: cells.{l162-s1b,l162-s2,l162-s3,l162-s4,l162-s5}.N1000.s1.accuracy :: sha256 9cd059d3e2f51798 :: mean +- SD over the row's 5 pretrained models
\newcommand{\FtAccJcEFourOnesixtwo}{0.8069\,\ensuremath{\pm}\,0.0025}  % experiments/FIGS/data/w2b_leg2_metrics_v2.json :: cells.{l162-s1b,l162-s2,l162-s3,l162-s4,l162-s5}.N10000.s1.accuracy :: sha256 9cd059d3e2f51798 :: mean +- SD over the row's 5 pretrained models
\newcommand{\FtAccJcEFiveOnesixtwo}{0.8363\,\ensuremath{\pm}\,0.0008}  % experiments/FIGS/data/w2b_leg2_metrics_v2.json :: cells.{l162-s1b,l162-s2,l162-s3,l162-s4,l162-s5}.N100000.s1.accuracy :: sha256 9cd059d3e2f51798 :: mean +- SD over the row's 5 pretrained models
\newcommand{\FtAccJcESixOnesixtwo}{0.8485\,\ensuremath{\pm}\,0.0005}  % experiments/FIGS/data/w2b_leg2_metrics_v2.json :: cells.{l162-s1b,l162-s2,l162-s3,l162-s4,l162-s5}.N1000000.s1.accuracy :: sha256 9cd059d3e2f51798 :: mean +- SD over the row's 5 pretrained models
\newcommand{\FtAccJcEThreeFourthree}{0.722\,\ensuremath{\pm}\,0.007}  % experiments/FIGS/data/w2b_leg2_metrics_v2.json :: cells.{r42q1-s1,r42q1-s2,r42q1-s3,r42q1-s4,r42q1-s5}.N1000.s1.accuracy :: sha256 9cd059d3e2f51798 :: mean +- SD over the row's 5 pretrained models
\newcommand{\FtAccJcEFourFourthree}{0.8050\,\ensuremath{\pm}\,0.0026}  % experiments/FIGS/data/w2b_leg2_metrics_v2.json :: cells.{r42q1-s1,r42q1-s2,r42q1-s3,r42q1-s4,r42q1-s5}.N10000.s1.accuracy :: sha256 9cd059d3e2f51798 :: mean +- SD over the row's 5 pretrained models
\newcommand{\FtAccJcEFiveFourthree}{0.8336\,\ensuremath{\pm}\,0.0008}  % experiments/FIGS/data/w2b_leg2_metrics_v2.json :: cells.{r42q1-s1,r42q1-s2,r42q1-s3,r42q1-s4,r42q1-s5}.N100000.s1.accuracy :: sha256 9cd059d3e2f51798 :: mean +- SD over the row's 5 pretrained models
\newcommand{\FtAccJcESixFourthree}{0.8468\,\ensuremath{\pm}\,0.0003}  % experiments/FIGS/data/w2b_leg2_metrics_v2.json :: cells.{r42q1-s1,r42q1-s2,r42q1-s3,r42q1-s4,r42q1-s5}.N1000000.s1.accuracy :: sha256 9cd059d3e2f51798 :: mean +- SD over the row's 5 pretrained models
\newcommand{\FtAccJcEThreeOneseven}{0.714\,\ensuremath{\pm}\,0.005}  % experiments/FIGS/data/w2b_leg2_metrics_v2.json :: cells.{r16q1-s1,r16q1-s2,r16q1-s3,r16q1-s4,r16q1-s5}.N1000.s1.accuracy :: sha256 9cd059d3e2f51798 :: mean +- SD over the row's 5 pretrained models
\newcommand{\FtAccJcEFourOneseven}{0.7964\,\ensuremath{\pm}\,0.0019}  % experiments/FIGS/data/w2b_leg2_metrics_v2.json :: cells.{r16q1-s1,r16q1-s2,r16q1-s3,r16q1-s4,r16q1-s5}.N10000.s1.accuracy :: sha256 9cd059d3e2f51798 :: mean +- SD over the row's 5 pretrained models
\newcommand{\FtAccJcEFiveOneseven}{0.8286\,\ensuremath{\pm}\,0.0006}  % experiments/FIGS/data/w2b_leg2_metrics_v2.json :: cells.{r16q1-s1,r16q1-s2,r16q1-s3,r16q1-s4,r16q1-s5}.N100000.s1.accuracy :: sha256 9cd059d3e2f51798 :: mean +- SD over the row's 5 pretrained models
\newcommand{\FtAccJcESixOneseven}{0.8415\,\ensuremath{\pm}\,0.0007}  % experiments/FIGS/data/w2b_leg2_metrics_v2.json :: cells.{r16q1-s1,r16q1-s2,r16q1-s3,r16q1-s4,r16q1-s5}.N1000000.s1.accuracy :: sha256 9cd059d3e2f51798 :: mean +- SD over the row's 5 pretrained models
\newcommand{\FtAccJcEThreeOnesixtwoMass}{0.728\,\ensuremath{\pm}\,0.004}  % experiments/FIGS/data/w2b_leg2_metrics_v2.json :: cells.{l162mass-s1,l162mass-s2,l162mass-s3,l162mass-s4,l162mass-s5}.N1000.s1.accuracy :: sha256 9cd059d3e2f51798 :: mean +- SD over the row's 5 pretrained models
\newcommand{\FtAccJcEFourOnesixtwoMass}{0.8043\,\ensuremath{\pm}\,0.0020}  % experiments/FIGS/data/w2b_leg2_metrics_v2.json :: cells.{l162mass-s1,l162mass-s2,l162mass-s3,l162mass-s4,l162mass-s5}.N10000.s1.accuracy :: sha256 9cd059d3e2f51798 :: mean +- SD over the row's 5 pretrained models
\newcommand{\FtAccJcEFiveOnesixtwoMass}{0.8349\,\ensuremath{\pm}\,0.0005}  % experiments/FIGS/data/w2b_leg2_metrics_v2.json :: cells.{l162mass-s1,l162mass-s2,l162mass-s3,l162mass-s4,l162mass-s5}.N100000.s1.accuracy :: sha256 9cd059d3e2f51798 :: mean +- SD over the row's 5 pretrained models
\newcommand{\FtAccJcESixOnesixtwoMass}{0.8475\,\ensuremath{\pm}\,0.0009}  % experiments/FIGS/data/w2b_leg2_metrics_v2.json :: cells.{l162mass-s1,l162mass-s2,l162mass-s3,l162mass-s4,l162mass-s5}.N1000000.s1.accuracy :: sha256 9cd059d3e2f51798 :: mean +- SD over the row's 5 pretrained models
\newcommand{\FtAccJcEThreeOnesevenMass}{0.702\,\ensuremath{\pm}\,0.005}  % experiments/FIGS/data/w2b_leg2_metrics_v2.json :: cells.{r16q1mass-s1,r16q1mass-s2,r16q1mass-s3,r16q1mass-s4,r16q1mass-s5}.N1000.s1.accuracy :: sha256 9cd059d3e2f51798 :: mean +- SD over the row's 5 pretrained models
\newcommand{\FtAccJcEFourOnesevenMass}{0.7909\,\ensuremath{\pm}\,0.0031}  % experiments/FIGS/data/w2b_leg2_metrics_v2.json :: cells.{r16q1mass-s1,r16q1mass-s2,r16q1mass-s3,r16q1mass-s4,r16q1mass-s5}.N10000.s1.accuracy :: sha256 9cd059d3e2f51798 :: mean +- SD over the row's 5 pretrained models
\newcommand{\FtAccJcEFiveOnesevenMass}{0.8265\,\ensuremath{\pm}\,0.0007}  % experiments/FIGS/data/w2b_leg2_metrics_v2.json :: cells.{r16q1mass-s1,r16q1mass-s2,r16q1mass-s3,r16q1mass-s4,r16q1mass-s5}.N100000.s1.accuracy :: sha256 9cd059d3e2f51798 :: mean +- SD over the row's 5 pretrained models
\newcommand{\FtAccJcESixOnesevenMass}{0.8390\,\ensuremath{\pm}\,0.0011}  % experiments/FIGS/data/w2b_leg2_metrics_v2.json :: cells.{r16q1mass-s1,r16q1mass-s2,r16q1mass-s3,r16q1mass-s4,r16q1mass-s5}.N1000000.s1.accuracy :: sha256 9cd059d3e2f51798 :: mean +- SD over the row's 5 pretrained models
\newcommand{\FtAccJcEThreeRandDrawOne}{0.7049}  % experiments/FIGS/data/w2b_leg2_metrics_v2.json :: cells.{rand-d1-s1b}.N1000.s1.accuracy :: sha256 9cd059d3e2f51798 :: one run
\newcommand{\FtAccJcEFourRandDrawOne}{0.7974}  % experiments/FIGS/data/w2b_leg2_metrics_v2.json :: cells.{rand-d1-s1b}.N10000.s1.accuracy :: sha256 9cd059d3e2f51798 :: one run
\newcommand{\FtAccJcEFiveRandDrawOne}{0.8326}  % experiments/FIGS/data/w2b_leg2_metrics_v2.json :: cells.{rand-d1-s1b}.N100000.s1.accuracy :: sha256 9cd059d3e2f51798 :: one run
\newcommand{\FtAccJcESixRandDrawOne}{0.84467}  % experiments/FIGS/data/w2b_leg2_metrics_v2.json :: cells.{rand-d1-s1b}.N1000000.s1.accuracy :: sha256 9cd059d3e2f51798 :: one run
\newcommand{\FtAccJcEThreeRandDrawTwo}{0.7135}  % experiments/FIGS/data/w2b_leg2_metrics_v2.json :: cells.{rand-d2-s2}.N1000.s1.accuracy :: sha256 9cd059d3e2f51798 :: one run
\newcommand{\FtAccJcEFourRandDrawTwo}{0.7998}  % experiments/FIGS/data/w2b_leg2_metrics_v2.json :: cells.{rand-d2-s2}.N10000.s1.accuracy :: sha256 9cd059d3e2f51798 :: one run
\newcommand{\FtAccJcEFiveRandDrawTwo}{0.8332}  % experiments/FIGS/data/w2b_leg2_metrics_v2.json :: cells.{rand-d2-s2}.N100000.s1.accuracy :: sha256 9cd059d3e2f51798 :: one run
\newcommand{\FtAccJcESixRandDrawTwo}{0.84449}  % experiments/FIGS/data/w2b_leg2_metrics_v2.json :: cells.{rand-d2-s2}.N1000000.s1.accuracy :: sha256 9cd059d3e2f51798 :: one run
\newcommand{\FtAccJcEThreeRandDrawThree}{0.7090}  % experiments/FIGS/data/w2b_leg2_metrics_v2.json :: cells.{rand-d3-s3}.N1000.s1.accuracy :: sha256 9cd059d3e2f51798 :: one run
\newcommand{\FtAccJcEFourRandDrawThree}{0.7948}  % experiments/FIGS/data/w2b_leg2_metrics_v2.json :: cells.{rand-d3-s3}.N10000.s1.accuracy :: sha256 9cd059d3e2f51798 :: one run
\newcommand{\FtAccJcEFiveRandDrawThree}{0.8310}  % experiments/FIGS/data/w2b_leg2_metrics_v2.json :: cells.{rand-d3-s3}.N100000.s1.accuracy :: sha256 9cd059d3e2f51798 :: one run
\newcommand{\FtAccJcESixRandDrawThree}{0.84409}  % experiments/FIGS/data/w2b_leg2_metrics_v2.json :: cells.{rand-d3-s3}.N1000000.s1.accuracy :: sha256 9cd059d3e2f51798 :: one run
\newcommand{\FtOmaRatioJcEThreeOnesixtwoOverOneeighteight}{0.96}  % experiments/FIGS/data/w2b_leg2_metrics_v2.json :: cells.*.N1000.s1.macro_auc_ovr :: sha256 9cd059d3e2f51798 :: ratio of the seed means of 1 - macro AUC, this row over the reference
\newcommand{\FtOmaRatioJcEThreeFourthreeOverOneeighteight}{1.01}  % experiments/FIGS/data/w2b_leg2_metrics_v2.json :: cells.*.N1000.s1.macro_auc_ovr :: sha256 9cd059d3e2f51798 :: ratio of the seed means of 1 - macro AUC, this row over the reference
\newcommand{\FtOmaRatioJcEThreeOnesevenOverOneeighteight}{1.07}  % experiments/FIGS/data/w2b_leg2_metrics_v2.json :: cells.*.N1000.s1.macro_auc_ovr :: sha256 9cd059d3e2f51798 :: ratio of the seed means of 1 - macro AUC, this row over the reference
\newcommand{\FtOmaRatioJcEThreeOnesixtwoMassOverOneeighteight}{0.87}  % experiments/FIGS/data/w2b_leg2_metrics_v2.json :: cells.*.N1000.s1.macro_auc_ovr :: sha256 9cd059d3e2f51798 :: ratio of the seed means of 1 - macro AUC, this row over the reference
\newcommand{\FtOmaRatioJcEThreeOnesevenMassOverOneeighteight}{1.09}  % experiments/FIGS/data/w2b_leg2_metrics_v2.json :: cells.*.N1000.s1.macro_auc_ovr :: sha256 9cd059d3e2f51798 :: ratio of the seed means of 1 - macro AUC, this row over the reference
\newcommand{\FtOmaRatioJcEThreeRandDrawOneOverOneeighteight}{1.06}  % experiments/FIGS/data/w2b_leg2_metrics_v2.json :: cells.*.N1000.s1.macro_auc_ovr :: sha256 9cd059d3e2f51798 :: ratio of the seed means of 1 - macro AUC, this row over the reference
\newcommand{\FtOmaRatioJcEThreeRandDrawTwoOverOneeighteight}{1.01}  % experiments/FIGS/data/w2b_leg2_metrics_v2.json :: cells.*.N1000.s1.macro_auc_ovr :: sha256 9cd059d3e2f51798 :: ratio of the seed means of 1 - macro AUC, this row over the reference
\newcommand{\FtOmaRatioJcEThreeRandDrawThreeOverOneeighteight}{1.07}  % experiments/FIGS/data/w2b_leg2_metrics_v2.json :: cells.*.N1000.s1.macro_auc_ovr :: sha256 9cd059d3e2f51798 :: ratio of the seed means of 1 - macro AUC, this row over the reference
\newcommand{\FtOmaRatioJcEThreeOneeighteightOverOnesixtwo}{1.05}  % experiments/FIGS/data/w2b_leg2_metrics_v2.json :: cells.*.N1000.s1.macro_auc_ovr :: sha256 9cd059d3e2f51798 :: ratio of the seed means of 1 - macro AUC, this row over the reference
\newcommand{\FtOmaRatioJcEThreeFourthreeOverOnesixtwo}{1.05}  % experiments/FIGS/data/w2b_leg2_metrics_v2.json :: cells.*.N1000.s1.macro_auc_ovr :: sha256 9cd059d3e2f51798 :: ratio of the seed means of 1 - macro AUC, this row over the reference
\newcommand{\FtOmaRatioJcEThreeOnesevenOverOnesixtwo}{1.12}  % experiments/FIGS/data/w2b_leg2_metrics_v2.json :: cells.*.N1000.s1.macro_auc_ovr :: sha256 9cd059d3e2f51798 :: ratio of the seed means of 1 - macro AUC, this row over the reference
\newcommand{\FtOmaRatioJcEThreeOnesixtwoMassOverOnesixtwo}{0.91}  % experiments/FIGS/data/w2b_leg2_metrics_v2.json :: cells.*.N1000.s1.macro_auc_ovr :: sha256 9cd059d3e2f51798 :: ratio of the seed means of 1 - macro AUC, this row over the reference
\newcommand{\FtOmaRatioJcEThreeOnesevenMassOverOnesixtwo}{1.14}  % experiments/FIGS/data/w2b_leg2_metrics_v2.json :: cells.*.N1000.s1.macro_auc_ovr :: sha256 9cd059d3e2f51798 :: ratio of the seed means of 1 - macro AUC, this row over the reference
\newcommand{\FtOmaRatioJcEThreeRandDrawOneOverOnesixtwo}{1.10}  % experiments/FIGS/data/w2b_leg2_metrics_v2.json :: cells.*.N1000.s1.macro_auc_ovr :: sha256 9cd059d3e2f51798 :: ratio of the seed means of 1 - macro AUC, this row over the reference
\newcommand{\FtOmaRatioJcEThreeRandDrawTwoOverOnesixtwo}{1.05}  % experiments/FIGS/data/w2b_leg2_metrics_v2.json :: cells.*.N1000.s1.macro_auc_ovr :: sha256 9cd059d3e2f51798 :: ratio of the seed means of 1 - macro AUC, this row over the reference
\newcommand{\FtOmaRatioJcEThreeRandDrawThreeOverOnesixtwo}{1.12}  % experiments/FIGS/data/w2b_leg2_metrics_v2.json :: cells.*.N1000.s1.macro_auc_ovr :: sha256 9cd059d3e2f51798 :: ratio of the seed means of 1 - macro AUC, this row over the reference
\newcommand{\FtOmaRatioJcEFourOnesixtwoOverOneeighteight}{0.98}  % experiments/FIGS/data/w2b_leg2_metrics_v2.json :: cells.*.N10000.s1.macro_auc_ovr :: sha256 9cd059d3e2f51798 :: ratio of the seed means of 1 - macro AUC, this row over the reference
\newcommand{\FtOmaRatioJcEFourFourthreeOverOneeighteight}{0.98}  % experiments/FIGS/data/w2b_leg2_metrics_v2.json :: cells.*.N10000.s1.macro_auc_ovr :: sha256 9cd059d3e2f51798 :: ratio of the seed means of 1 - macro AUC, this row over the reference
\newcommand{\FtOmaRatioJcEFourOnesevenOverOneeighteight}{1.07}  % experiments/FIGS/data/w2b_leg2_metrics_v2.json :: cells.*.N10000.s1.macro_auc_ovr :: sha256 9cd059d3e2f51798 :: ratio of the seed means of 1 - macro AUC, this row over the reference
\newcommand{\FtOmaRatioJcEFourOnesixtwoMassOverOneeighteight}{0.99}  % experiments/FIGS/data/w2b_leg2_metrics_v2.json :: cells.*.N10000.s1.macro_auc_ovr :: sha256 9cd059d3e2f51798 :: ratio of the seed means of 1 - macro AUC, this row over the reference
\newcommand{\FtOmaRatioJcEFourOnesevenMassOverOneeighteight}{1.12}  % experiments/FIGS/data/w2b_leg2_metrics_v2.json :: cells.*.N10000.s1.macro_auc_ovr :: sha256 9cd059d3e2f51798 :: ratio of the seed means of 1 - macro AUC, this row over the reference
\newcommand{\FtOmaRatioJcEFourRandDrawOneOverOneeighteight}{1.08}  % experiments/FIGS/data/w2b_leg2_metrics_v2.json :: cells.*.N10000.s1.macro_auc_ovr :: sha256 9cd059d3e2f51798 :: ratio of the seed means of 1 - macro AUC, this row over the reference
\newcommand{\FtOmaRatioJcEFourRandDrawTwoOverOneeighteight}{1.04}  % experiments/FIGS/data/w2b_leg2_metrics_v2.json :: cells.*.N10000.s1.macro_auc_ovr :: sha256 9cd059d3e2f51798 :: ratio of the seed means of 1 - macro AUC, this row over the reference
\newcommand{\FtOmaRatioJcEFourRandDrawThreeOverOneeighteight}{1.09}  % experiments/FIGS/data/w2b_leg2_metrics_v2.json :: cells.*.N10000.s1.macro_auc_ovr :: sha256 9cd059d3e2f51798 :: ratio of the seed means of 1 - macro AUC, this row over the reference
\newcommand{\FtOmaRatioJcEFourOneeighteightOverOnesixtwo}{1.02}  % experiments/FIGS/data/w2b_leg2_metrics_v2.json :: cells.*.N10000.s1.macro_auc_ovr :: sha256 9cd059d3e2f51798 :: ratio of the seed means of 1 - macro AUC, this row over the reference
\newcommand{\FtOmaRatioJcEFourFourthreeOverOnesixtwo}{1.01}  % experiments/FIGS/data/w2b_leg2_metrics_v2.json :: cells.*.N10000.s1.macro_auc_ovr :: sha256 9cd059d3e2f51798 :: ratio of the seed means of 1 - macro AUC, this row over the reference
\newcommand{\FtOmaRatioJcEFourOnesevenOverOnesixtwo}{1.09}  % experiments/FIGS/data/w2b_leg2_metrics_v2.json :: cells.*.N10000.s1.macro_auc_ovr :: sha256 9cd059d3e2f51798 :: ratio of the seed means of 1 - macro AUC, this row over the reference
\newcommand{\FtOmaRatioJcEFourOnesixtwoMassOverOnesixtwo}{1.01}  % experiments/FIGS/data/w2b_leg2_metrics_v2.json :: cells.*.N10000.s1.macro_auc_ovr :: sha256 9cd059d3e2f51798 :: ratio of the seed means of 1 - macro AUC, this row over the reference
\newcommand{\FtOmaRatioJcEFourOnesevenMassOverOnesixtwo}{1.14}  % experiments/FIGS/data/w2b_leg2_metrics_v2.json :: cells.*.N10000.s1.macro_auc_ovr :: sha256 9cd059d3e2f51798 :: ratio of the seed means of 1 - macro AUC, this row over the reference
\newcommand{\FtOmaRatioJcEFourRandDrawOneOverOnesixtwo}{1.11}  % experiments/FIGS/data/w2b_leg2_metrics_v2.json :: cells.*.N10000.s1.macro_auc_ovr :: sha256 9cd059d3e2f51798 :: ratio of the seed means of 1 - macro AUC, this row over the reference
\newcommand{\FtOmaRatioJcEFourRandDrawTwoOverOnesixtwo}{1.06}  % experiments/FIGS/data/w2b_leg2_metrics_v2.json :: cells.*.N10000.s1.macro_auc_ovr :: sha256 9cd059d3e2f51798 :: ratio of the seed means of 1 - macro AUC, this row over the reference
\newcommand{\FtOmaRatioJcEFourRandDrawThreeOverOnesixtwo}{1.12}  % experiments/FIGS/data/w2b_leg2_metrics_v2.json :: cells.*.N10000.s1.macro_auc_ovr :: sha256 9cd059d3e2f51798 :: ratio of the seed means of 1 - macro AUC, this row over the reference
\newcommand{\FtOmaRatioJcEFiveOnesixtwoOverOneeighteight}{0.98}  % experiments/FIGS/data/w2b_leg2_metrics_v2.json :: cells.*.N100000.s1.macro_auc_ovr :: sha256 9cd059d3e2f51798 :: ratio of the seed means of 1 - macro AUC, this row over the reference
\newcommand{\FtOmaRatioJcEFiveFourthreeOverOneeighteight}{1.02}  % experiments/FIGS/data/w2b_leg2_metrics_v2.json :: cells.*.N100000.s1.macro_auc_ovr :: sha256 9cd059d3e2f51798 :: ratio of the seed means of 1 - macro AUC, this row over the reference
\newcommand{\FtOmaRatioJcEFiveOnesevenOverOneeighteight}{1.07}  % experiments/FIGS/data/w2b_leg2_metrics_v2.json :: cells.*.N100000.s1.macro_auc_ovr :: sha256 9cd059d3e2f51798 :: ratio of the seed means of 1 - macro AUC, this row over the reference
\newcommand{\FtOmaRatioJcEFiveOnesixtwoMassOverOneeighteight}{1.00}  % experiments/FIGS/data/w2b_leg2_metrics_v2.json :: cells.*.N100000.s1.macro_auc_ovr :: sha256 9cd059d3e2f51798 :: ratio of the seed means of 1 - macro AUC, this row over the reference
\newcommand{\FtOmaRatioJcEFiveOnesevenMassOverOneeighteight}{1.10}  % experiments/FIGS/data/w2b_leg2_metrics_v2.json :: cells.*.N100000.s1.macro_auc_ovr :: sha256 9cd059d3e2f51798 :: ratio of the seed means of 1 - macro AUC, this row over the reference
\newcommand{\FtOmaRatioJcEFiveRandDrawOneOverOneeighteight}{1.04}  % experiments/FIGS/data/w2b_leg2_metrics_v2.json :: cells.*.N100000.s1.macro_auc_ovr :: sha256 9cd059d3e2f51798 :: ratio of the seed means of 1 - macro AUC, this row over the reference
\newcommand{\FtOmaRatioJcEFiveRandDrawTwoOverOneeighteight}{1.03}  % experiments/FIGS/data/w2b_leg2_metrics_v2.json :: cells.*.N100000.s1.macro_auc_ovr :: sha256 9cd059d3e2f51798 :: ratio of the seed means of 1 - macro AUC, this row over the reference
\newcommand{\FtOmaRatioJcEFiveRandDrawThreeOverOneeighteight}{1.05}  % experiments/FIGS/data/w2b_leg2_metrics_v2.json :: cells.*.N100000.s1.macro_auc_ovr :: sha256 9cd059d3e2f51798 :: ratio of the seed means of 1 - macro AUC, this row over the reference
\newcommand{\FtOmaRatioJcEFiveOneeighteightOverOnesixtwo}{1.02}  % experiments/FIGS/data/w2b_leg2_metrics_v2.json :: cells.*.N100000.s1.macro_auc_ovr :: sha256 9cd059d3e2f51798 :: ratio of the seed means of 1 - macro AUC, this row over the reference
\newcommand{\FtOmaRatioJcEFiveFourthreeOverOnesixtwo}{1.03}  % experiments/FIGS/data/w2b_leg2_metrics_v2.json :: cells.*.N100000.s1.macro_auc_ovr :: sha256 9cd059d3e2f51798 :: ratio of the seed means of 1 - macro AUC, this row over the reference
\newcommand{\FtOmaRatioJcEFiveOnesevenOverOnesixtwo}{1.09}  % experiments/FIGS/data/w2b_leg2_metrics_v2.json :: cells.*.N100000.s1.macro_auc_ovr :: sha256 9cd059d3e2f51798 :: ratio of the seed means of 1 - macro AUC, this row over the reference
\newcommand{\FtOmaRatioJcEFiveOnesixtwoMassOverOnesixtwo}{1.02}  % experiments/FIGS/data/w2b_leg2_metrics_v2.json :: cells.*.N100000.s1.macro_auc_ovr :: sha256 9cd059d3e2f51798 :: ratio of the seed means of 1 - macro AUC, this row over the reference
\newcommand{\FtOmaRatioJcEFiveOnesevenMassOverOnesixtwo}{1.11}  % experiments/FIGS/data/w2b_leg2_metrics_v2.json :: cells.*.N100000.s1.macro_auc_ovr :: sha256 9cd059d3e2f51798 :: ratio of the seed means of 1 - macro AUC, this row over the reference
\newcommand{\FtOmaRatioJcEFiveRandDrawOneOverOnesixtwo}{1.06}  % experiments/FIGS/data/w2b_leg2_metrics_v2.json :: cells.*.N100000.s1.macro_auc_ovr :: sha256 9cd059d3e2f51798 :: ratio of the seed means of 1 - macro AUC, this row over the reference
\newcommand{\FtOmaRatioJcEFiveRandDrawTwoOverOnesixtwo}{1.05}  % experiments/FIGS/data/w2b_leg2_metrics_v2.json :: cells.*.N100000.s1.macro_auc_ovr :: sha256 9cd059d3e2f51798 :: ratio of the seed means of 1 - macro AUC, this row over the reference
\newcommand{\FtOmaRatioJcEFiveRandDrawThreeOverOnesixtwo}{1.07}  % experiments/FIGS/data/w2b_leg2_metrics_v2.json :: cells.*.N100000.s1.macro_auc_ovr :: sha256 9cd059d3e2f51798 :: ratio of the seed means of 1 - macro AUC, this row over the reference
\newcommand{\FtOmaRatioJcESixOnesixtwoOverOneeighteight}{1.00}  % experiments/FIGS/data/w2b_leg2_metrics_v2.json :: cells.*.N1000000.s1.macro_auc_ovr :: sha256 9cd059d3e2f51798 :: ratio of the seed means of 1 - macro AUC, this row over the reference
\newcommand{\FtOmaRatioJcESixFourthreeOverOneeighteight}{1.02}  % experiments/FIGS/data/w2b_leg2_metrics_v2.json :: cells.*.N1000000.s1.macro_auc_ovr :: sha256 9cd059d3e2f51798 :: ratio of the seed means of 1 - macro AUC, this row over the reference
\newcommand{\FtOmaRatioJcESixOnesevenOverOneeighteight}{1.08}  % experiments/FIGS/data/w2b_leg2_metrics_v2.json :: cells.*.N1000000.s1.macro_auc_ovr :: sha256 9cd059d3e2f51798 :: ratio of the seed means of 1 - macro AUC, this row over the reference
\newcommand{\FtOmaRatioJcESixOnesixtwoMassOverOneeighteight}{1.01}  % experiments/FIGS/data/w2b_leg2_metrics_v2.json :: cells.*.N1000000.s1.macro_auc_ovr :: sha256 9cd059d3e2f51798 :: ratio of the seed means of 1 - macro AUC, this row over the reference
\newcommand{\FtOmaRatioJcESixOnesevenMassOverOneeighteight}{1.11}  % experiments/FIGS/data/w2b_leg2_metrics_v2.json :: cells.*.N1000000.s1.macro_auc_ovr :: sha256 9cd059d3e2f51798 :: ratio of the seed means of 1 - macro AUC, this row over the reference
\newcommand{\FtOmaRatioJcESixRandDrawOneOverOneeighteight}{1.05}  % experiments/FIGS/data/w2b_leg2_metrics_v2.json :: cells.*.N1000000.s1.macro_auc_ovr :: sha256 9cd059d3e2f51798 :: ratio of the seed means of 1 - macro AUC, this row over the reference
\newcommand{\FtOmaRatioJcESixRandDrawTwoOverOneeighteight}{1.05}  % experiments/FIGS/data/w2b_leg2_metrics_v2.json :: cells.*.N1000000.s1.macro_auc_ovr :: sha256 9cd059d3e2f51798 :: ratio of the seed means of 1 - macro AUC, this row over the reference
\newcommand{\FtOmaRatioJcESixRandDrawThreeOverOneeighteight}{1.05}  % experiments/FIGS/data/w2b_leg2_metrics_v2.json :: cells.*.N1000000.s1.macro_auc_ovr :: sha256 9cd059d3e2f51798 :: ratio of the seed means of 1 - macro AUC, this row over the reference
\newcommand{\FtOmaRatioJcESixOneeighteightOverOnesixtwo}{1.00}  % experiments/FIGS/data/w2b_leg2_metrics_v2.json :: cells.*.N1000000.s1.macro_auc_ovr :: sha256 9cd059d3e2f51798 :: ratio of the seed means of 1 - macro AUC, this row over the reference
\newcommand{\FtOmaRatioJcESixFourthreeOverOnesixtwo}{1.02}  % experiments/FIGS/data/w2b_leg2_metrics_v2.json :: cells.*.N1000000.s1.macro_auc_ovr :: sha256 9cd059d3e2f51798 :: ratio of the seed means of 1 - macro AUC, this row over the reference
\newcommand{\FtOmaRatioJcESixOnesevenOverOnesixtwo}{1.08}  % experiments/FIGS/data/w2b_leg2_metrics_v2.json :: cells.*.N1000000.s1.macro_auc_ovr :: sha256 9cd059d3e2f51798 :: ratio of the seed means of 1 - macro AUC, this row over the reference
\newcommand{\FtOmaRatioJcESixOnesixtwoMassOverOnesixtwo}{1.01}  % experiments/FIGS/data/w2b_leg2_metrics_v2.json :: cells.*.N1000000.s1.macro_auc_ovr :: sha256 9cd059d3e2f51798 :: ratio of the seed means of 1 - macro AUC, this row over the reference
\newcommand{\FtOmaRatioJcESixOnesevenMassOverOnesixtwo}{1.11}  % experiments/FIGS/data/w2b_leg2_metrics_v2.json :: cells.*.N1000000.s1.macro_auc_ovr :: sha256 9cd059d3e2f51798 :: ratio of the seed means of 1 - macro AUC, this row over the reference
\newcommand{\FtOmaRatioJcESixRandDrawOneOverOnesixtwo}{1.05}  % experiments/FIGS/data/w2b_leg2_metrics_v2.json :: cells.*.N1000000.s1.macro_auc_ovr :: sha256 9cd059d3e2f51798 :: ratio of the seed means of 1 - macro AUC, this row over the reference
\newcommand{\FtOmaRatioJcESixRandDrawTwoOverOnesixtwo}{1.05}  % experiments/FIGS/data/w2b_leg2_metrics_v2.json :: cells.*.N1000000.s1.macro_auc_ovr :: sha256 9cd059d3e2f51798 :: ratio of the seed means of 1 - macro AUC, this row over the reference
\newcommand{\FtOmaRatioJcESixRandDrawThreeOverOnesixtwo}{1.05}  % experiments/FIGS/data/w2b_leg2_metrics_v2.json :: cells.*.N1000000.s1.macro_auc_ovr :: sha256 9cd059d3e2f51798 :: ratio of the seed means of 1 - macro AUC, this row over the reference
\newcommand{\PairedFtAccJcEThreeOnesevenMinusOneeighteight}{\ensuremath{-}0.0040~[\ensuremath{-}0.0144, +0.0064]}  % experiments/FIGS/data/w2b_leg2_metrics_v2.json :: cells.{r16q1-s1,r16q1-s2,r16q1-s3,r16q1-s4,r16q1-s5,l188-s1,l188-s2,l188-s3,l188-s4,l188-s5}.N1000.s1.accuracy :: sha256 9cd059d3e2f51798 :: accuracy, coarse minus fine, paired by run [95% Student-t interval]
\newcommand{\PairedFtAccJcEFourOnesevenMinusOneeighteight}{\ensuremath{-}0.0075~[\ensuremath{-}0.0118, \ensuremath{-}0.0033]}  % experiments/FIGS/data/w2b_leg2_metrics_v2.json :: cells.{r16q1-s1,r16q1-s2,r16q1-s3,r16q1-s4,r16q1-s5,l188-s1,l188-s2,l188-s3,l188-s4,l188-s5}.N10000.s1.accuracy :: sha256 9cd059d3e2f51798 :: accuracy, coarse minus fine, paired by run [95% Student-t interval]
\newcommand{\PairedFtAccJcEFiveOnesevenMinusOneeighteight}{\ensuremath{-}0.0066~[\ensuremath{-}0.0083, \ensuremath{-}0.0049]}  % experiments/FIGS/data/w2b_leg2_metrics_v2.json :: cells.{r16q1-s1,r16q1-s2,r16q1-s3,r16q1-s4,r16q1-s5,l188-s1,l188-s2,l188-s3,l188-s4,l188-s5}.N100000.s1.accuracy :: sha256 9cd059d3e2f51798 :: accuracy, coarse minus fine, paired by run [95% Student-t interval]
\newcommand{\PairedFtAccJcESixOnesevenMinusOneeighteight}{\ensuremath{-}0.0071~[\ensuremath{-}0.0080, \ensuremath{-}0.0062]}  % experiments/FIGS/data/w2b_leg2_metrics_v2.json :: cells.{r16q1-s1,r16q1-s2,r16q1-s3,r16q1-s4,r16q1-s5,l188-s1,l188-s2,l188-s3,l188-s4,l188-s5}.N1000000.s1.accuracy :: sha256 9cd059d3e2f51798 :: accuracy, coarse minus fine, paired by run [95% Student-t interval]
\newcommand{\FtNClassesJc}{10}  % experiments/FIGS/data/w2b_leg2_metrics_v2.json :: cells.*.*.s1.n_classes_present :: sha256 9cd059d3e2f51798 :: classes of the fine-tuning task
\newcommand{\FtNTestAucJc}{500{,}000}  % experiments/FIGS/data/w2b_leg2_metrics_v2.json :: cells.*.*.s1.n_jets_auc :: sha256 9cd059d3e2f51798 :: test jets the macro AUC is computed on
\newcommand{\FtNTestAccJc}{2{,}000{,}000}  % experiments/FIGS/data/w2b_leg2_metrics_v2.json :: cells.*.*.s1.n_jets :: sha256 9cd059d3e2f51798 :: test jets the accuracy is computed on
\newcommand{\FtSizeEThree}{\ensuremath{10^{3}}}  % experiments/FIGS/data/w2b_leg1_metrics_v2.json :: cells.*.N1000 :: sha256 b56bf70be9982044 :: fine-tuning training jets
\newcommand{\FtSizeEFour}{\ensuremath{10^{4}}}  % experiments/FIGS/data/w2b_leg1_metrics_v2.json :: cells.*.N10000 :: sha256 b56bf70be9982044 :: fine-tuning training jets
\newcommand{\FtSizeEFive}{\ensuremath{10^{5}}}  % experiments/FIGS/data/w2b_leg1_metrics_v2.json :: cells.*.N100000 :: sha256 b56bf70be9982044 :: fine-tuning training jets
\newcommand{\FtSizeESix}{\ensuremath{10^{6}}}  % experiments/FIGS/data/w2b_leg1_metrics_v2.json :: cells.*.N1000000 :: sha256 b56bf70be9982044 :: fine-tuning training jets
\newcommand{\FtNRunsWord}{five}  % experiments/FIGS/data/w2b_leg1_metrics_v2.json :: cells (pretrained models per vocabulary row) :: sha256 b56bf70be9982044 :: pretraining runs per vocabulary
\newcommand{\PairedProbeBcVsRestLinearOnesevenOverOnesixtwo}{1.64~[1.52, 1.78]}  % experiments/FIGS/data/paired_v1err/probes/ratios.json :: ratios[0].ratio, ratios[0].ci95 :: sha256 14dd51402305d147 :: 17 over 162, bc_vs_rest linear 1-auc: paired geometric mean [95% interval]
\newcommand{\PairedProbeBcVsRestMlpOnesevenOverOnesixtwo}{1.50~[1.39, 1.62]}  % experiments/FIGS/data/paired_v1err/probes/ratios.json :: ratios[3].ratio, ratios[3].ci95 :: sha256 14dd51402305d147 :: 17 over 162, bc_vs_rest mlp 1-auc: paired geometric mean [95% interval]
\newcommand{\PairedProbeBvcFourprongLinearOnesixtwoOverOneeighteight}{1.01~[0.95, 1.08]}  % experiments/FIGS/data/paired_v1err/probes/ratios.json :: ratios[6].ratio, ratios[6].ci95 :: sha256 14dd51402305d147 :: 162 over 188, bvc_4prong linear 1-auc: paired geometric mean [95% interval]
\newcommand{\PairedProbeBvcFourprongLinearFourthreeOverOneeighteight}{1.10~[1.01, 1.20]}  % experiments/FIGS/data/paired_v1err/probes/ratios.json :: ratios[7].ratio, ratios[7].ci95 :: sha256 14dd51402305d147 :: 43 over 188, bvc_4prong linear 1-auc: paired geometric mean [95% interval]
\newcommand{\PairedProbeBvcFourprongLinearOnesevenOverOneeighteight}{4.5~[3.5, 5.7]}  % experiments/FIGS/data/paired_v1err/probes/ratios.json :: ratios[8].ratio, ratios[8].ci95 :: sha256 14dd51402305d147 :: 17 over 188, bvc_4prong linear 1-auc: paired geometric mean [95% interval]
\newcommand{\PairedProbeBvcFourprongLinearFourthreeOverOnesixtwo}{1.09~[0.98, 1.20]}  % experiments/FIGS/data/paired_v1err/probes/ratios.json :: ratios[9].ratio, ratios[9].ci95 :: sha256 14dd51402305d147 :: 43 over 162, bvc_4prong linear 1-auc: paired geometric mean [95% interval]
\newcommand{\PairedProbeBvcFourprongLinearOnesevenOverOnesixtwo}{4.4~[3.4, 5.6]}  % experiments/FIGS/data/paired_v1err/probes/ratios.json :: ratios[10].ratio, ratios[10].ci95 :: sha256 14dd51402305d147 :: 17 over 162, bvc_4prong linear 1-auc: paired geometric mean [95% interval]
\newcommand{\PairedProbeBvcFourprongLinearOnesevenOverFourthree}{4.0~[3.1, 5.4]}  % experiments/FIGS/data/paired_v1err/probes/ratios.json :: ratios[11].ratio, ratios[11].ci95 :: sha256 14dd51402305d147 :: 17 over 43, bvc_4prong linear 1-auc: paired geometric mean [95% interval]
\newcommand{\PairedProbeBvcFourprongLinearOnesixtwoMassOverOnesixtwo}{1.05~[0.97, 1.13]}  % experiments/FIGS/data/paired_v1err/probes/ratios.json :: ratios[12].ratio, ratios[12].ci95 :: sha256 14dd51402305d147 :: 162+mass over 162, bvc_4prong linear 1-auc: paired geometric mean [95% interval]
\newcommand{\PairedProbeBvcFourprongLinearOnesevenMassOverOneseven}{2.1~[1.6, 2.7]}  % experiments/FIGS/data/paired_v1err/probes/ratios.json :: ratios[13].ratio, ratios[13].ci95 :: sha256 14dd51402305d147 :: 17+mass over 17, bvc_4prong linear 1-auc: paired geometric mean [95% interval]
\newcommand{\PairedProbeBvcFourprongLinearRandomOverOneeighteight}{1.54~[1.37, 1.74]}  % experiments/FIGS/data/paired_v1err/probes/ratios.json :: ratios[14].ratio, ratios[14].ci95 :: sha256 14dd51402305d147 :: random over 188, bvc_4prong linear 1-auc: paired geometric mean [95% interval]
\newcommand{\PairedProbeBvcFourprongLinearRandomOverOnesixtwo}{1.55~[1.36, 1.77]}  % experiments/FIGS/data/paired_v1err/probes/ratios.json :: ratios[15].ratio, ratios[15].ci95 :: sha256 14dd51402305d147 :: random over 162, bvc_4prong linear 1-auc: paired geometric mean [95% interval]
\newcommand{\PairedProbeBvcFourprongLinearRandomOverFourthree}{1.39~[1.27, 1.52]}  % experiments/FIGS/data/paired_v1err/probes/ratios.json :: ratios[16].ratio, ratios[16].ci95 :: sha256 14dd51402305d147 :: random over 43, bvc_4prong linear 1-auc: paired geometric mean [95% interval]
\newcommand{\PairedProbeBvcFourprongLinearRandomOverOneseven}{0.34~[0.20, 0.58]}  % experiments/FIGS/data/paired_v1err/probes/ratios.json :: ratios[17].ratio, ratios[17].ci95 :: sha256 14dd51402305d147 :: random over 17, bvc_4prong linear 1-auc: paired geometric mean [95% interval]
\newcommand{\PairedProbeBvcFourprongLinearRandomDrawTwoOverRandomDrawOne}{1.05~[0.94, 1.18]}  % experiments/FIGS/data/paired_v1err/probes/ratios.json :: ratios[18].ratio, ratios[18].ci95 :: sha256 14dd51402305d147 :: random draw 2 over random draw 1, bvc_4prong linear 1-auc: paired geometric mean [95% interval]
\newcommand{\PairedProbeBvcFourprongLinearRandomDrawThreeOverRandomDrawOne}{0.95~[0.85, 1.06]}  % experiments/FIGS/data/paired_v1err/probes/ratios.json :: ratios[19].ratio, ratios[19].ci95 :: sha256 14dd51402305d147 :: random draw 3 over random draw 1, bvc_4prong linear 1-auc: paired geometric mean [95% interval]
\newcommand{\PairedProbeBvcFourprongLinearRandomDrawThreeOverRandomDrawTwo}{0.90~[0.79, 1.02]}  % experiments/FIGS/data/paired_v1err/probes/ratios.json :: ratios[20].ratio, ratios[20].ci95 :: sha256 14dd51402305d147 :: random draw 3 over random draw 2, bvc_4prong linear 1-auc: paired geometric mean [95% interval]
\newcommand{\PairedProbeBvcFourprongMlpOnesixtwoOverOneeighteight}{1.00~[0.93, 1.07]}  % experiments/FIGS/data/paired_v1err/probes/ratios.json :: ratios[66].ratio, ratios[66].ci95 :: sha256 14dd51402305d147 :: 162 over 188, bvc_4prong mlp 1-auc: paired geometric mean [95% interval]
\newcommand{\PairedProbeBvcFourprongMlpFourthreeOverOneeighteight}{1.12~[1.03, 1.22]}  % experiments/FIGS/data/paired_v1err/probes/ratios.json :: ratios[67].ratio, ratios[67].ci95 :: sha256 14dd51402305d147 :: 43 over 188, bvc_4prong mlp 1-auc: paired geometric mean [95% interval]
\newcommand{\PairedProbeBvcFourprongMlpOnesevenOverOneeighteight}{3.2~[2.6, 4.1]}  % experiments/FIGS/data/paired_v1err/probes/ratios.json :: ratios[68].ratio, ratios[68].ci95 :: sha256 14dd51402305d147 :: 17 over 188, bvc_4prong mlp 1-auc: paired geometric mean [95% interval]
\newcommand{\PairedProbeBvcFourprongMlpFourthreeOverOnesixtwo}{1.12~[1.01, 1.26]}  % experiments/FIGS/data/paired_v1err/probes/ratios.json :: ratios[69].ratio, ratios[69].ci95 :: sha256 14dd51402305d147 :: 43 over 162, bvc_4prong mlp 1-auc: paired geometric mean [95% interval]
\newcommand{\PairedProbeBvcFourprongMlpOnesevenOverOnesixtwo}{3.3~[2.6, 4.0]}  % experiments/FIGS/data/paired_v1err/probes/ratios.json :: ratios[70].ratio, ratios[70].ci95 :: sha256 14dd51402305d147 :: 17 over 162, bvc_4prong mlp 1-auc: paired geometric mean [95% interval]
\newcommand{\PairedProbeBvcFourprongMlpOnesevenOverFourthree}{2.9~[2.3, 3.7]}  % experiments/FIGS/data/paired_v1err/probes/ratios.json :: ratios[71].ratio, ratios[71].ci95 :: sha256 14dd51402305d147 :: 17 over 43, bvc_4prong mlp 1-auc: paired geometric mean [95% interval]
\newcommand{\PairedProbeBvcFourprongMlpOnesixtwoMassOverOnesixtwo}{1.07~[0.99, 1.16]}  % experiments/FIGS/data/paired_v1err/probes/ratios.json :: ratios[72].ratio, ratios[72].ci95 :: sha256 14dd51402305d147 :: 162+mass over 162, bvc_4prong mlp 1-auc: paired geometric mean [95% interval]
\newcommand{\PairedProbeBvcFourprongMlpOnesevenMassOverOneseven}{1.97~[1.46, 2.67]}  % experiments/FIGS/data/paired_v1err/probes/ratios.json :: ratios[73].ratio, ratios[73].ci95 :: sha256 14dd51402305d147 :: 17+mass over 17, bvc_4prong mlp 1-auc: paired geometric mean [95% interval]
\newcommand{\PairedProbeBvcFourprongMlpRandomOverOneeighteight}{1.47~[1.30, 1.65]}  % experiments/FIGS/data/paired_v1err/probes/ratios.json :: ratios[74].ratio, ratios[74].ci95 :: sha256 14dd51402305d147 :: random over 188, bvc_4prong mlp 1-auc: paired geometric mean [95% interval]
\newcommand{\PairedProbeBvcFourprongMlpRandomOverOnesixtwo}{1.49~[1.24, 1.79]}  % experiments/FIGS/data/paired_v1err/probes/ratios.json :: ratios[75].ratio, ratios[75].ci95 :: sha256 14dd51402305d147 :: random over 162, bvc_4prong mlp 1-auc: paired geometric mean [95% interval]
\newcommand{\PairedProbeBvcFourprongMlpRandomOverFourthree}{1.28~[1.16, 1.43]}  % experiments/FIGS/data/paired_v1err/probes/ratios.json :: ratios[76].ratio, ratios[76].ci95 :: sha256 14dd51402305d147 :: random over 43, bvc_4prong mlp 1-auc: paired geometric mean [95% interval]
\newcommand{\PairedProbeBvcFourprongMlpRandomOverOneseven}{0.43~[0.28, 0.67]}  % experiments/FIGS/data/paired_v1err/probes/ratios.json :: ratios[77].ratio, ratios[77].ci95 :: sha256 14dd51402305d147 :: random over 17, bvc_4prong mlp 1-auc: paired geometric mean [95% interval]
\newcommand{\PairedProbeBvcFourprongMlpRandomDrawTwoOverRandomDrawOne}{1.00~[0.90, 1.10]}  % experiments/FIGS/data/paired_v1err/probes/ratios.json :: ratios[78].ratio, ratios[78].ci95 :: sha256 14dd51402305d147 :: random draw 2 over random draw 1, bvc_4prong mlp 1-auc: paired geometric mean [95% interval]
\newcommand{\PairedProbeBvcFourprongMlpRandomDrawThreeOverRandomDrawOne}{0.89~[0.79, 0.99]}  % experiments/FIGS/data/paired_v1err/probes/ratios.json :: ratios[79].ratio, ratios[79].ci95 :: sha256 14dd51402305d147 :: random draw 3 over random draw 1, bvc_4prong mlp 1-auc: paired geometric mean [95% interval]
\newcommand{\PairedProbeBvcFourprongMlpRandomDrawThreeOverRandomDrawTwo}{0.89~[0.79, 0.99]}  % experiments/FIGS/data/paired_v1err/probes/ratios.json :: ratios[80].ratio, ratios[80].ci95 :: sha256 14dd51402305d147 :: random draw 3 over random draw 2, bvc_4prong mlp 1-auc: paired geometric mean [95% interval]
\newcommand{\PairedProbeBvcQcdLinearOnesixtwoOverOneeighteight}{1.15~[1.07, 1.23]}  % experiments/FIGS/data/paired_v1err/probes/ratios.json :: ratios[126].ratio, ratios[126].ci95 :: sha256 14dd51402305d147 :: 162 over 188, bvc_qcd linear 1-auc: paired geometric mean [95% interval]
\newcommand{\PairedProbeBvcQcdLinearFourthreeOverOneeighteight}{1.16~[1.08, 1.25]}  % experiments/FIGS/data/paired_v1err/probes/ratios.json :: ratios[127].ratio, ratios[127].ci95 :: sha256 14dd51402305d147 :: 43 over 188, bvc_qcd linear 1-auc: paired geometric mean [95% interval]
\newcommand{\PairedProbeBvcQcdLinearOnesevenOverOneeighteight}{3.0~[2.5, 3.7]}  % experiments/FIGS/data/paired_v1err/probes/ratios.json :: ratios[128].ratio, ratios[128].ci95 :: sha256 14dd51402305d147 :: 17 over 188, bvc_qcd linear 1-auc: paired geometric mean [95% interval]
\newcommand{\PairedProbeBvcQcdLinearFourthreeOverOnesixtwo}{1.01~[0.96, 1.07]}  % experiments/FIGS/data/paired_v1err/probes/ratios.json :: ratios[129].ratio, ratios[129].ci95 :: sha256 14dd51402305d147 :: 43 over 162, bvc_qcd linear 1-auc: paired geometric mean [95% interval]
\newcommand{\PairedProbeBvcQcdLinearOnesevenOverOnesixtwo}{2.6~[2.2, 3.1]}  % experiments/FIGS/data/paired_v1err/probes/ratios.json :: ratios[130].ratio, ratios[130].ci95 :: sha256 14dd51402305d147 :: 17 over 162, bvc_qcd linear 1-auc: paired geometric mean [95% interval]
\newcommand{\PairedProbeBvcQcdLinearOnesevenOverFourthree}{2.6~[2.2, 3.0]}  % experiments/FIGS/data/paired_v1err/probes/ratios.json :: ratios[131].ratio, ratios[131].ci95 :: sha256 14dd51402305d147 :: 17 over 43, bvc_qcd linear 1-auc: paired geometric mean [95% interval]
\newcommand{\PairedProbeBvcQcdLinearOnesixtwoMassOverOnesixtwo}{1.060~[1.002, 1.123]}  % experiments/FIGS/data/paired_v1err/probes/ratios.json :: ratios[132].ratio, ratios[132].ci95 :: sha256 14dd51402305d147 :: 162+mass over 162, bvc_qcd linear 1-auc: paired geometric mean [95% interval]
\newcommand{\PairedProbeBvcQcdLinearOnesevenMassOverOneseven}{1.89~[1.52, 2.34]}  % experiments/FIGS/data/paired_v1err/probes/ratios.json :: ratios[133].ratio, ratios[133].ci95 :: sha256 14dd51402305d147 :: 17+mass over 17, bvc_qcd linear 1-auc: paired geometric mean [95% interval]
\newcommand{\PairedProbeBvcQcdMlpOnesixtwoOverOneeighteight}{1.11~[1.04, 1.19]}  % experiments/FIGS/data/paired_v1err/probes/ratios.json :: ratios[158].ratio, ratios[158].ci95 :: sha256 14dd51402305d147 :: 162 over 188, bvc_qcd mlp 1-auc: paired geometric mean [95% interval]
\newcommand{\PairedProbeBvcQcdMlpFourthreeOverOneeighteight}{1.13~[1.05, 1.21]}  % experiments/FIGS/data/paired_v1err/probes/ratios.json :: ratios[159].ratio, ratios[159].ci95 :: sha256 14dd51402305d147 :: 43 over 188, bvc_qcd mlp 1-auc: paired geometric mean [95% interval]
\newcommand{\PairedProbeBvcQcdMlpOnesevenOverOneeighteight}{2.6~[2.2, 3.0]}  % experiments/FIGS/data/paired_v1err/probes/ratios.json :: ratios[160].ratio, ratios[160].ci95 :: sha256 14dd51402305d147 :: 17 over 188, bvc_qcd mlp 1-auc: paired geometric mean [95% interval]
\newcommand{\PairedProbeBvcQcdMlpFourthreeOverOnesixtwo}{1.01~[0.96, 1.06]}  % experiments/FIGS/data/paired_v1err/probes/ratios.json :: ratios[161].ratio, ratios[161].ci95 :: sha256 14dd51402305d147 :: 43 over 162, bvc_qcd mlp 1-auc: paired geometric mean [95% interval]
\newcommand{\PairedProbeBvcQcdMlpOnesevenOverOnesixtwo}{2.3~[2.0, 2.6]}  % experiments/FIGS/data/paired_v1err/probes/ratios.json :: ratios[162].ratio, ratios[162].ci95 :: sha256 14dd51402305d147 :: 17 over 162, bvc_qcd mlp 1-auc: paired geometric mean [95% interval]
\newcommand{\PairedProbeBvcQcdMlpOnesevenOverFourthree}{2.3~[2.0, 2.6]}  % experiments/FIGS/data/paired_v1err/probes/ratios.json :: ratios[163].ratio, ratios[163].ci95 :: sha256 14dd51402305d147 :: 17 over 43, bvc_qcd mlp 1-auc: paired geometric mean [95% interval]
\newcommand{\PairedProbeBvcQcdMlpOnesixtwoMassOverOnesixtwo}{1.06~[1.01, 1.12]}  % experiments/FIGS/data/paired_v1err/probes/ratios.json :: ratios[164].ratio, ratios[164].ci95 :: sha256 14dd51402305d147 :: 162+mass over 162, bvc_qcd mlp 1-auc: paired geometric mean [95% interval]
\newcommand{\PairedProbeBvcQcdMlpOnesevenMassOverOneseven}{1.64~[1.33, 2.02]}  % experiments/FIGS/data/paired_v1err/probes/ratios.json :: ratios[165].ratio, ratios[165].ci95 :: sha256 14dd51402305d147 :: 17+mass over 17, bvc_qcd mlp 1-auc: paired geometric mean [95% interval]
\newcommand{\PairedProbeBvcResonantLinearOnesixtwoOverOneeighteight}{1.02~[0.96, 1.08]}  % experiments/FIGS/data/paired_v1err/probes/ratios.json :: ratios[190].ratio, ratios[190].ci95 :: sha256 14dd51402305d147 :: 162 over 188, bvc_resonant linear 1-auc: paired geometric mean [95% interval]
\newcommand{\PairedProbeBvcResonantLinearFourthreeOverOneeighteight}{1.09~[1.03, 1.15]}  % experiments/FIGS/data/paired_v1err/probes/ratios.json :: ratios[191].ratio, ratios[191].ci95 :: sha256 14dd51402305d147 :: 43 over 188, bvc_resonant linear 1-auc: paired geometric mean [95% interval]
\newcommand{\PairedProbeBvcResonantLinearOnesevenOverOneeighteight}{4.0~[3.3, 4.9]}  % experiments/FIGS/data/paired_v1err/probes/ratios.json :: ratios[192].ratio, ratios[192].ci95 :: sha256 14dd51402305d147 :: 17 over 188, bvc_resonant linear 1-auc: paired geometric mean [95% interval]
\newcommand{\PairedProbeBvcResonantLinearFourthreeOverOnesixtwo}{1.07~[1.02, 1.12]}  % experiments/FIGS/data/paired_v1err/probes/ratios.json :: ratios[193].ratio, ratios[193].ci95 :: sha256 14dd51402305d147 :: 43 over 162, bvc_resonant linear 1-auc: paired geometric mean [95% interval]
\newcommand{\PairedProbeBvcResonantLinearOnesevenOverOnesixtwo}{3.9~[3.1, 4.9]}  % experiments/FIGS/data/paired_v1err/probes/ratios.json :: ratios[194].ratio, ratios[194].ci95 :: sha256 14dd51402305d147 :: 17 over 162, bvc_resonant linear 1-auc: paired geometric mean [95% interval]
\newcommand{\PairedProbeBvcResonantLinearOnesevenOverFourthree}{3.7~[3.0, 4.5]}  % experiments/FIGS/data/paired_v1err/probes/ratios.json :: ratios[195].ratio, ratios[195].ci95 :: sha256 14dd51402305d147 :: 17 over 43, bvc_resonant linear 1-auc: paired geometric mean [95% interval]
\newcommand{\PairedProbeBvcResonantLinearOnesixtwoMassOverOnesixtwo}{1.11~[1.06, 1.16]}  % experiments/FIGS/data/paired_v1err/probes/ratios.json :: ratios[196].ratio, ratios[196].ci95 :: sha256 14dd51402305d147 :: 162+mass over 162, bvc_resonant linear 1-auc: paired geometric mean [95% interval]
\newcommand{\PairedProbeBvcResonantLinearOnesevenMassOverOneseven}{3.0~[1.9, 4.7]}  % experiments/FIGS/data/paired_v1err/probes/ratios.json :: ratios[197].ratio, ratios[197].ci95 :: sha256 14dd51402305d147 :: 17+mass over 17, bvc_resonant linear 1-auc: paired geometric mean [95% interval]
\newcommand{\PairedProbeBvcResonantMlpOnesixtwoOverOneeighteight}{1.04~[0.98, 1.12]}  % experiments/FIGS/data/paired_v1err/probes/ratios.json :: ratios[222].ratio, ratios[222].ci95 :: sha256 14dd51402305d147 :: 162 over 188, bvc_resonant mlp 1-auc: paired geometric mean [95% interval]
\newcommand{\PairedProbeBvcResonantMlpFourthreeOverOneeighteight}{1.14~[1.08, 1.21]}  % experiments/FIGS/data/paired_v1err/probes/ratios.json :: ratios[223].ratio, ratios[223].ci95 :: sha256 14dd51402305d147 :: 43 over 188, bvc_resonant mlp 1-auc: paired geometric mean [95% interval]
\newcommand{\PairedProbeBvcResonantMlpOnesevenOverOneeighteight}{3.3~[2.8, 3.8]}  % experiments/FIGS/data/paired_v1err/probes/ratios.json :: ratios[224].ratio, ratios[224].ci95 :: sha256 14dd51402305d147 :: 17 over 188, bvc_resonant mlp 1-auc: paired geometric mean [95% interval]
\newcommand{\PairedProbeBvcResonantMlpFourthreeOverOnesixtwo}{1.09~[1.04, 1.14]}  % experiments/FIGS/data/paired_v1err/probes/ratios.json :: ratios[225].ratio, ratios[225].ci95 :: sha256 14dd51402305d147 :: 43 over 162, bvc_resonant mlp 1-auc: paired geometric mean [95% interval]
\newcommand{\PairedProbeBvcResonantMlpOnesevenOverOnesixtwo}{3.1~[2.6, 3.8]}  % experiments/FIGS/data/paired_v1err/probes/ratios.json :: ratios[226].ratio, ratios[226].ci95 :: sha256 14dd51402305d147 :: 17 over 162, bvc_resonant mlp 1-auc: paired geometric mean [95% interval]
\newcommand{\PairedProbeBvcResonantMlpOnesevenOverFourthree}{2.9~[2.4, 3.4]}  % experiments/FIGS/data/paired_v1err/probes/ratios.json :: ratios[227].ratio, ratios[227].ci95 :: sha256 14dd51402305d147 :: 17 over 43, bvc_resonant mlp 1-auc: paired geometric mean [95% interval]
\newcommand{\PairedProbeBvcResonantMlpOnesixtwoMassOverOnesixtwo}{1.09~[1.04, 1.14]}  % experiments/FIGS/data/paired_v1err/probes/ratios.json :: ratios[228].ratio, ratios[228].ci95 :: sha256 14dd51402305d147 :: 162+mass over 162, bvc_resonant mlp 1-auc: paired geometric mean [95% interval]
\newcommand{\PairedProbeBvcResonantMlpOnesevenMassOverOneseven}{1.95~[1.45, 2.61]}  % experiments/FIGS/data/paired_v1err/probes/ratios.json :: ratios[229].ratio, ratios[229].ci95 :: sha256 14dd51402305d147 :: 17+mass over 17, bvc_resonant mlp 1-auc: paired geometric mean [95% interval]
\newcommand{\PairedProbeEeVsMmLinearOnesevenOverFourthree}{2.3~[0.8, 6.5]}  % experiments/FIGS/data/paired_v1err/probes/ratios.json :: ratios[259].ratio, ratios[259].ci95 :: sha256 14dd51402305d147 :: 17 over 43, ee_vs_mm linear 1-auc: paired geometric mean [95% interval]
\newcommand{\PairedProbeEeVsMmLinearOnesevenMassOverOneseven}{0.23~[0.08, 0.64]}  % experiments/FIGS/data/paired_v1err/probes/ratios.json :: ratios[261].ratio, ratios[261].ci95 :: sha256 14dd51402305d147 :: 17+mass over 17, ee_vs_mm linear 1-auc: paired geometric mean [95% interval]
\newcommand{\PairedProbeEeVsMmMlpOnesevenOverFourthree}{2.8~[0.8, 9.0]}  % experiments/FIGS/data/paired_v1err/probes/ratios.json :: ratios[291].ratio, ratios[291].ci95 :: sha256 14dd51402305d147 :: 17 over 43, ee_vs_mm mlp 1-auc: paired geometric mean [95% interval]
\newcommand{\PairedProbeEeVsMmMlpOnesevenMassOverOneseven}{0.172~[0.026, 1.125]}  % experiments/FIGS/data/paired_v1err/probes/ratios.json :: ratios[293].ratio, ratios[293].ci95 :: sha256 14dd51402305d147 :: 17+mass over 17, ee_vs_mm mlp 1-auc: paired geometric mean [95% interval]
\newcommand{\PairedProbeRetainedTopologyLinearOnesixtwoOverOneeighteight}{1.00~[0.92, 1.08]}  % experiments/FIGS/data/paired_v1err/probes/ratios.json :: ratios[318].ratio, ratios[318].ci95 :: sha256 14dd51402305d147 :: 162 over 188, retained_topology linear 1-auc: paired geometric mean [95% interval]
\newcommand{\PairedProbeRetainedTopologyLinearFourthreeOverOneeighteight}{1.24~[1.15, 1.34]}  % experiments/FIGS/data/paired_v1err/probes/ratios.json :: ratios[319].ratio, ratios[319].ci95 :: sha256 14dd51402305d147 :: 43 over 188, retained_topology linear 1-auc: paired geometric mean [95% interval]
\newcommand{\PairedProbeRetainedTopologyLinearOnesevenOverOneeighteight}{1.54~[1.35, 1.75]}  % experiments/FIGS/data/paired_v1err/probes/ratios.json :: ratios[320].ratio, ratios[320].ci95 :: sha256 14dd51402305d147 :: 17 over 188, retained_topology linear 1-auc: paired geometric mean [95% interval]
\newcommand{\PairedProbeRetainedTopologyLinearFourthreeOverOnesixtwo}{1.24~[1.14, 1.34]}  % experiments/FIGS/data/paired_v1err/probes/ratios.json :: ratios[321].ratio, ratios[321].ci95 :: sha256 14dd51402305d147 :: 43 over 162, retained_topology linear 1-auc: paired geometric mean [95% interval]
\newcommand{\PairedProbeRetainedTopologyLinearOnesevenOverOnesixtwo}{1.54~[1.33, 1.78]}  % experiments/FIGS/data/paired_v1err/probes/ratios.json :: ratios[322].ratio, ratios[322].ci95 :: sha256 14dd51402305d147 :: 17 over 162, retained_topology linear 1-auc: paired geometric mean [95% interval]
\newcommand{\PairedProbeRetainedTopologyLinearOnesevenOverFourthree}{1.24~[1.08, 1.42]}  % experiments/FIGS/data/paired_v1err/probes/ratios.json :: ratios[323].ratio, ratios[323].ci95 :: sha256 14dd51402305d147 :: 17 over 43, retained_topology linear 1-auc: paired geometric mean [95% interval]
\newcommand{\PairedProbeRetainedTopologyLinearOnesixtwoMassOverOnesixtwo}{1.06~[0.99, 1.14]}  % experiments/FIGS/data/paired_v1err/probes/ratios.json :: ratios[324].ratio, ratios[324].ci95 :: sha256 14dd51402305d147 :: 162+mass over 162, retained_topology linear 1-auc: paired geometric mean [95% interval]
\newcommand{\PairedProbeRetainedTopologyLinearOnesevenMassOverOneseven}{1.12~[0.97, 1.28]}  % experiments/FIGS/data/paired_v1err/probes/ratios.json :: ratios[325].ratio, ratios[325].ci95 :: sha256 14dd51402305d147 :: 17+mass over 17, retained_topology linear 1-auc: paired geometric mean [95% interval]
\newcommand{\PairedProbeRetainedTopologyMlpOnesixtwoOverOneeighteight}{1.01~[0.87, 1.17]}  % experiments/FIGS/data/paired_v1err/probes/ratios.json :: ratios[350].ratio, ratios[350].ci95 :: sha256 14dd51402305d147 :: 162 over 188, retained_topology mlp 1-auc: paired geometric mean [95% interval]
\newcommand{\PairedProbeRetainedTopologyMlpFourthreeOverOneeighteight}{1.18~[1.03, 1.36]}  % experiments/FIGS/data/paired_v1err/probes/ratios.json :: ratios[351].ratio, ratios[351].ci95 :: sha256 14dd51402305d147 :: 43 over 188, retained_topology mlp 1-auc: paired geometric mean [95% interval]
\newcommand{\PairedProbeRetainedTopologyMlpOnesevenOverOneeighteight}{1.52~[1.34, 1.72]}  % experiments/FIGS/data/paired_v1err/probes/ratios.json :: ratios[352].ratio, ratios[352].ci95 :: sha256 14dd51402305d147 :: 17 over 188, retained_topology mlp 1-auc: paired geometric mean [95% interval]
\newcommand{\PairedProbeRetainedTopologyMlpFourthreeOverOnesixtwo}{1.17~[1.07, 1.28]}  % experiments/FIGS/data/paired_v1err/probes/ratios.json :: ratios[353].ratio, ratios[353].ci95 :: sha256 14dd51402305d147 :: 43 over 162, retained_topology mlp 1-auc: paired geometric mean [95% interval]
\newcommand{\PairedProbeRetainedTopologyMlpOnesevenOverOnesixtwo}{1.50~[1.29, 1.75]}  % experiments/FIGS/data/paired_v1err/probes/ratios.json :: ratios[354].ratio, ratios[354].ci95 :: sha256 14dd51402305d147 :: 17 over 162, retained_topology mlp 1-auc: paired geometric mean [95% interval]
\newcommand{\PairedProbeRetainedTopologyMlpOnesevenOverFourthree}{1.28~[1.12, 1.47]}  % experiments/FIGS/data/paired_v1err/probes/ratios.json :: ratios[355].ratio, ratios[355].ci95 :: sha256 14dd51402305d147 :: 17 over 43, retained_topology mlp 1-auc: paired geometric mean [95% interval]
\newcommand{\PairedProbeRetainedTopologyMlpOnesixtwoMassOverOnesixtwo}{1.04~[0.92, 1.17]}  % experiments/FIGS/data/paired_v1err/probes/ratios.json :: ratios[356].ratio, ratios[356].ci95 :: sha256 14dd51402305d147 :: 162+mass over 162, retained_topology mlp 1-auc: paired geometric mean [95% interval]
\newcommand{\PairedProbeRetainedTopologyMlpOnesevenMassOverOneseven}{1.11~[0.99, 1.23]}  % experiments/FIGS/data/paired_v1err/probes/ratios.json :: ratios[357].ratio, ratios[357].ci95 :: sha256 14dd51402305d147 :: 17+mass over 17, retained_topology mlp 1-auc: paired geometric mean [95% interval]
\newcommand{\PairedProbeVisibleContentLinearOnesixtwoOverOneeighteight}{0.95~[0.88, 1.03]}  % experiments/FIGS/data/paired_v1err/probes/ratios.json :: ratios[382].ratio, ratios[382].ci95 :: sha256 14dd51402305d147 :: 162 over 188, visible_content linear 1-auc: paired geometric mean [95% interval]
\newcommand{\PairedProbeVisibleContentLinearFourthreeOverOneeighteight}{0.99~[0.85, 1.15]}  % experiments/FIGS/data/paired_v1err/probes/ratios.json :: ratios[383].ratio, ratios[383].ci95 :: sha256 14dd51402305d147 :: 43 over 188, visible_content linear 1-auc: paired geometric mean [95% interval]
\newcommand{\PairedProbeVisibleContentLinearOnesevenOverOneeighteight}{2.8~[2.1, 3.7]}  % experiments/FIGS/data/paired_v1err/probes/ratios.json :: ratios[384].ratio, ratios[384].ci95 :: sha256 14dd51402305d147 :: 17 over 188, visible_content linear 1-auc: paired geometric mean [95% interval]
\newcommand{\PairedProbeVisibleContentLinearFourthreeOverOnesixtwo}{1.04~[0.89, 1.20]}  % experiments/FIGS/data/paired_v1err/probes/ratios.json :: ratios[385].ratio, ratios[385].ci95 :: sha256 14dd51402305d147 :: 43 over 162, visible_content linear 1-auc: paired geometric mean [95% interval]
\newcommand{\PairedProbeVisibleContentLinearOnesevenOverOnesixtwo}{2.9~[2.3, 3.8]}  % experiments/FIGS/data/paired_v1err/probes/ratios.json :: ratios[386].ratio, ratios[386].ci95 :: sha256 14dd51402305d147 :: 17 over 162, visible_content linear 1-auc: paired geometric mean [95% interval]
\newcommand{\PairedProbeVisibleContentLinearOnesevenOverFourthree}{2.8~[2.2, 3.7]}  % experiments/FIGS/data/paired_v1err/probes/ratios.json :: ratios[387].ratio, ratios[387].ci95 :: sha256 14dd51402305d147 :: 17 over 43, visible_content linear 1-auc: paired geometric mean [95% interval]
\newcommand{\PairedProbeVisibleContentLinearOnesixtwoMassOverOnesixtwo}{1.10~[0.98, 1.23]}  % experiments/FIGS/data/paired_v1err/probes/ratios.json :: ratios[388].ratio, ratios[388].ci95 :: sha256 14dd51402305d147 :: 162+mass over 162, visible_content linear 1-auc: paired geometric mean [95% interval]
\newcommand{\PairedProbeVisibleContentLinearOnesevenMassOverOneseven}{1.44~[1.26, 1.65]}  % experiments/FIGS/data/paired_v1err/probes/ratios.json :: ratios[389].ratio, ratios[389].ci95 :: sha256 14dd51402305d147 :: 17+mass over 17, visible_content linear 1-auc: paired geometric mean [95% interval]
\newcommand{\PairedProbeVisibleContentLinearRandomOverOneeighteight}{1.31~[1.01, 1.70]}  % experiments/FIGS/data/paired_v1err/probes/ratios.json :: ratios[390].ratio, ratios[390].ci95 :: sha256 14dd51402305d147 :: random over 188, visible_content linear 1-auc: paired geometric mean [95% interval]
\newcommand{\PairedProbeVisibleContentLinearRandomOverOnesixtwo}{1.41~[1.17, 1.70]}  % experiments/FIGS/data/paired_v1err/probes/ratios.json :: ratios[391].ratio, ratios[391].ci95 :: sha256 14dd51402305d147 :: random over 162, visible_content linear 1-auc: paired geometric mean [95% interval]
\newcommand{\PairedProbeVisibleContentLinearRandomOverFourthree}{1.33~[1.12, 1.57]}  % experiments/FIGS/data/paired_v1err/probes/ratios.json :: ratios[392].ratio, ratios[392].ci95 :: sha256 14dd51402305d147 :: random over 43, visible_content linear 1-auc: paired geometric mean [95% interval]
\newcommand{\PairedProbeVisibleContentLinearRandomOverOneseven}{0.48~[0.38, 0.60]}  % experiments/FIGS/data/paired_v1err/probes/ratios.json :: ratios[393].ratio, ratios[393].ci95 :: sha256 14dd51402305d147 :: random over 17, visible_content linear 1-auc: paired geometric mean [95% interval]
\newcommand{\PairedProbeVisibleContentLinearRandomDrawTwoOverRandomDrawOne}{0.79~[0.64, 0.98]}  % experiments/FIGS/data/paired_v1err/probes/ratios.json :: ratios[394].ratio, ratios[394].ci95 :: sha256 14dd51402305d147 :: random draw 2 over random draw 1, visible_content linear 1-auc: paired geometric mean [95% interval]
\newcommand{\PairedProbeVisibleContentLinearRandomDrawThreeOverRandomDrawOne}{0.86~[0.71, 1.05]}  % experiments/FIGS/data/paired_v1err/probes/ratios.json :: ratios[395].ratio, ratios[395].ci95 :: sha256 14dd51402305d147 :: random draw 3 over random draw 1, visible_content linear 1-auc: paired geometric mean [95% interval]
\newcommand{\PairedProbeVisibleContentLinearRandomDrawThreeOverRandomDrawTwo}{1.09~[0.88, 1.35]}  % experiments/FIGS/data/paired_v1err/probes/ratios.json :: ratios[396].ratio, ratios[396].ci95 :: sha256 14dd51402305d147 :: random draw 3 over random draw 2, visible_content linear 1-auc: paired geometric mean [95% interval]
\newcommand{\PairedProbeVisibleContentMlpOnesixtwoOverOneeighteight}{0.92~[0.85, 1.01]}  % experiments/FIGS/data/paired_v1err/probes/ratios.json :: ratios[442].ratio, ratios[442].ci95 :: sha256 14dd51402305d147 :: 162 over 188, visible_content mlp 1-auc: paired geometric mean [95% interval]
\newcommand{\PairedProbeVisibleContentMlpFourthreeOverOneeighteight}{0.93~[0.80, 1.07]}  % experiments/FIGS/data/paired_v1err/probes/ratios.json :: ratios[443].ratio, ratios[443].ci95 :: sha256 14dd51402305d147 :: 43 over 188, visible_content mlp 1-auc: paired geometric mean [95% interval]
\newcommand{\PairedProbeVisibleContentMlpOnesevenOverOneeighteight}{2.3~[1.7, 3.1]}  % experiments/FIGS/data/paired_v1err/probes/ratios.json :: ratios[444].ratio, ratios[444].ci95 :: sha256 14dd51402305d147 :: 17 over 188, visible_content mlp 1-auc: paired geometric mean [95% interval]
\newcommand{\PairedProbeVisibleContentMlpFourthreeOverOnesixtwo}{1.01~[0.89, 1.14]}  % experiments/FIGS/data/paired_v1err/probes/ratios.json :: ratios[445].ratio, ratios[445].ci95 :: sha256 14dd51402305d147 :: 43 over 162, visible_content mlp 1-auc: paired geometric mean [95% interval]
\newcommand{\PairedProbeVisibleContentMlpOnesevenOverOnesixtwo}{2.5~[1.9, 3.2]}  % experiments/FIGS/data/paired_v1err/probes/ratios.json :: ratios[446].ratio, ratios[446].ci95 :: sha256 14dd51402305d147 :: 17 over 162, visible_content mlp 1-auc: paired geometric mean [95% interval]
\newcommand{\PairedProbeVisibleContentMlpOnesevenOverFourthree}{2.5~[2.0, 3.1]}  % experiments/FIGS/data/paired_v1err/probes/ratios.json :: ratios[447].ratio, ratios[447].ci95 :: sha256 14dd51402305d147 :: 17 over 43, visible_content mlp 1-auc: paired geometric mean [95% interval]
\newcommand{\PairedProbeVisibleContentMlpOnesixtwoMassOverOnesixtwo}{1.08~[0.94, 1.24]}  % experiments/FIGS/data/paired_v1err/probes/ratios.json :: ratios[448].ratio, ratios[448].ci95 :: sha256 14dd51402305d147 :: 162+mass over 162, visible_content mlp 1-auc: paired geometric mean [95% interval]
\newcommand{\PairedProbeVisibleContentMlpOnesevenMassOverOneseven}{1.41~[1.19, 1.67]}  % experiments/FIGS/data/paired_v1err/probes/ratios.json :: ratios[449].ratio, ratios[449].ci95 :: sha256 14dd51402305d147 :: 17+mass over 17, visible_content mlp 1-auc: paired geometric mean [95% interval]
\newcommand{\PairedProbeVisibleContentMlpRandomOverOneeighteight}{1.27~[1.01, 1.59]}  % experiments/FIGS/data/paired_v1err/probes/ratios.json :: ratios[450].ratio, ratios[450].ci95 :: sha256 14dd51402305d147 :: random over 188, visible_content mlp 1-auc: paired geometric mean [95% interval]
\newcommand{\PairedProbeVisibleContentMlpRandomOverOnesixtwo}{1.37~[1.12, 1.68]}  % experiments/FIGS/data/paired_v1err/probes/ratios.json :: ratios[451].ratio, ratios[451].ci95 :: sha256 14dd51402305d147 :: random over 162, visible_content mlp 1-auc: paired geometric mean [95% interval]
\newcommand{\PairedProbeVisibleContentMlpRandomOverFourthree}{1.35~[1.13, 1.61]}  % experiments/FIGS/data/paired_v1err/probes/ratios.json :: ratios[452].ratio, ratios[452].ci95 :: sha256 14dd51402305d147 :: random over 43, visible_content mlp 1-auc: paired geometric mean [95% interval]
\newcommand{\PairedProbeVisibleContentMlpRandomOverOneseven}{0.56~[0.45, 0.70]}  % experiments/FIGS/data/paired_v1err/probes/ratios.json :: ratios[453].ratio, ratios[453].ci95 :: sha256 14dd51402305d147 :: random over 17, visible_content mlp 1-auc: paired geometric mean [95% interval]
\newcommand{\PairedProbeVisibleContentMlpRandomDrawTwoOverRandomDrawOne}{0.88~[0.71, 1.09]}  % experiments/FIGS/data/paired_v1err/probes/ratios.json :: ratios[454].ratio, ratios[454].ci95 :: sha256 14dd51402305d147 :: random draw 2 over random draw 1, visible_content mlp 1-auc: paired geometric mean [95% interval]
\newcommand{\PairedProbeVisibleContentMlpRandomDrawThreeOverRandomDrawOne}{0.94~[0.75, 1.18]}  % experiments/FIGS/data/paired_v1err/probes/ratios.json :: ratios[455].ratio, ratios[455].ci95 :: sha256 14dd51402305d147 :: random draw 3 over random draw 1, visible_content mlp 1-auc: paired geometric mean [95% interval]
\newcommand{\PairedProbeVisibleContentMlpRandomDrawThreeOverRandomDrawTwo}{1.06~[0.84, 1.35]}  % experiments/FIGS/data/paired_v1err/probes/ratios.json :: ratios[456].ratio, ratios[456].ci95 :: sha256 14dd51402305d147 :: random draw 3 over random draw 2, visible_content mlp 1-auc: paired geometric mean [95% interval]
\newcommand{\PairedFtJciiEThreeOnesixtwoOverOneeighteight}{0.95~[0.87, 1.05]}  % experiments/FIGS/data/paired_v1err/finetune/ratios.json :: ratios[0].ratio, ratios[0].ci95 :: sha256 3e5cf25e859b68c2 :: 162 over 188, leg1 N1000 1-macro_auc: paired geometric mean [95% interval]
\newcommand{\PairedFtJciiEThreeFourthreeOverOneeighteight}{0.84~[0.77, 0.91]}  % experiments/FIGS/data/paired_v1err/finetune/ratios.json :: ratios[1].ratio, ratios[1].ci95 :: sha256 3e5cf25e859b68c2 :: 43 over 188, leg1 N1000 1-macro_auc: paired geometric mean [95% interval]
\newcommand{\PairedFtJciiEThreeOnesevenOverOneeighteight}{0.97~[0.86, 1.09]}  % experiments/FIGS/data/paired_v1err/finetune/ratios.json :: ratios[2].ratio, ratios[2].ci95 :: sha256 3e5cf25e859b68c2 :: 17 over 188, leg1 N1000 1-macro_auc: paired geometric mean [95% interval]
\newcommand{\PairedFtJciiEThreeFourthreeOverOnesixtwo}{0.88~[0.80, 0.97]}  % experiments/FIGS/data/paired_v1err/finetune/ratios.json :: ratios[3].ratio, ratios[3].ci95 :: sha256 3e5cf25e859b68c2 :: 43 over 162, leg1 N1000 1-macro_auc: paired geometric mean [95% interval]
\newcommand{\PairedFtJciiEThreeOnesevenOverOnesixtwo}{1.02~[0.88, 1.18]}  % experiments/FIGS/data/paired_v1err/finetune/ratios.json :: ratios[4].ratio, ratios[4].ci95 :: sha256 3e5cf25e859b68c2 :: 17 over 162, leg1 N1000 1-macro_auc: paired geometric mean [95% interval]
\newcommand{\PairedFtJciiEThreeOnesevenOverFourthree}{1.16~[1.08, 1.24]}  % experiments/FIGS/data/paired_v1err/finetune/ratios.json :: ratios[5].ratio, ratios[5].ci95 :: sha256 3e5cf25e859b68c2 :: 17 over 43, leg1 N1000 1-macro_auc: paired geometric mean [95% interval]
\newcommand{\PairedFtJciiEThreeOnesixtwoMassOverOnesixtwo}{0.85~[0.77, 0.93]}  % experiments/FIGS/data/paired_v1err/finetune/ratios.json :: ratios[6].ratio, ratios[6].ci95 :: sha256 3e5cf25e859b68c2 :: 162+mass over 162, leg1 N1000 1-macro_auc: paired geometric mean [95% interval]
\newcommand{\PairedFtJciiEThreeOnesevenMassOverOneseven}{1.00~[0.86, 1.17]}  % experiments/FIGS/data/paired_v1err/finetune/ratios.json :: ratios[7].ratio, ratios[7].ci95 :: sha256 3e5cf25e859b68c2 :: 17+mass over 17, leg1 N1000 1-macro_auc: paired geometric mean [95% interval]
\newcommand{\PairedFtJciiEThreeRandomOverOneeighteight}{0.82~[0.73, 0.93]}  % experiments/FIGS/data/paired_v1err/finetune/ratios.json :: ratios[8].ratio, ratios[8].ci95 :: sha256 3e5cf25e859b68c2 :: random over 188, leg1 N1000 1-macro_auc: paired geometric mean [95% interval]
\newcommand{\PairedFtJciiEThreeRandomOverOnesixtwo}{0.90~[0.85, 0.96]}  % experiments/FIGS/data/paired_v1err/finetune/ratios.json :: ratios[9].ratio, ratios[9].ci95 :: sha256 3e5cf25e859b68c2 :: random over 162, leg1 N1000 1-macro_auc: paired geometric mean [95% interval]
\newcommand{\PairedFtJciiEThreeRandomOverFourthree}{0.99~[0.75, 1.30]}  % experiments/FIGS/data/paired_v1err/finetune/ratios.json :: ratios[10].ratio, ratios[10].ci95 :: sha256 3e5cf25e859b68c2 :: random over 43, leg1 N1000 1-macro_auc: paired geometric mean [95% interval]
\newcommand{\PairedFtJciiEThreeRandomOverOneseven}{0.86~[0.58, 1.27]}  % experiments/FIGS/data/paired_v1err/finetune/ratios.json :: ratios[11].ratio, ratios[11].ci95 :: sha256 3e5cf25e859b68c2 :: random over 17, leg1 N1000 1-macro_auc: paired geometric mean [95% interval]
\newcommand{\PairedFtJciiEThreeRandomDrawTwoOverRandomDrawOne}{1.09~[1.07, 1.10]}  % experiments/FIGS/data/paired_v1err/finetune/ratios.json :: ratios[12].ratio, ratios[12].ci95 :: sha256 3e5cf25e859b68c2 :: random draw 2 over random draw 1, leg1 N1000 1-macro_auc: paired geometric mean [95% interval]
\newcommand{\PairedFtJciiEThreeRandomDrawThreeOverRandomDrawOne}{1.04~[1.02, 1.05]}  % experiments/FIGS/data/paired_v1err/finetune/ratios.json :: ratios[13].ratio, ratios[13].ci95 :: sha256 3e5cf25e859b68c2 :: random draw 3 over random draw 1, leg1 N1000 1-macro_auc: paired geometric mean [95% interval]
\newcommand{\PairedFtJciiEThreeRandomDrawThreeOverRandomDrawTwo}{0.95~[0.94, 0.97]}  % experiments/FIGS/data/paired_v1err/finetune/ratios.json :: ratios[14].ratio, ratios[14].ci95 :: sha256 3e5cf25e859b68c2 :: random draw 3 over random draw 2, leg1 N1000 1-macro_auc: paired geometric mean [95% interval]
\newcommand{\PairedFtJciiEFourOnesixtwoOverOneeighteight}{0.99~[0.97, 1.02]}  % experiments/FIGS/data/paired_v1err/finetune/ratios.json :: ratios[15].ratio, ratios[15].ci95 :: sha256 3e5cf25e859b68c2 :: 162 over 188, leg1 N10000 1-macro_auc: paired geometric mean [95% interval]
\newcommand{\PairedFtJciiEFourFourthreeOverOneeighteight}{1.20~[1.17, 1.23]}  % experiments/FIGS/data/paired_v1err/finetune/ratios.json :: ratios[16].ratio, ratios[16].ci95 :: sha256 3e5cf25e859b68c2 :: 43 over 188, leg1 N10000 1-macro_auc: paired geometric mean [95% interval]
\newcommand{\PairedFtJciiEFourOnesevenOverOneeighteight}{1.32~[1.29, 1.34]}  % experiments/FIGS/data/paired_v1err/finetune/ratios.json :: ratios[17].ratio, ratios[17].ci95 :: sha256 3e5cf25e859b68c2 :: 17 over 188, leg1 N10000 1-macro_auc: paired geometric mean [95% interval]
\newcommand{\PairedFtJciiEFourFourthreeOverOnesixtwo}{1.21~[1.19, 1.23]}  % experiments/FIGS/data/paired_v1err/finetune/ratios.json :: ratios[18].ratio, ratios[18].ci95 :: sha256 3e5cf25e859b68c2 :: 43 over 162, leg1 N10000 1-macro_auc: paired geometric mean [95% interval]
\newcommand{\PairedFtJciiEFourOnesevenOverOnesixtwo}{1.32~[1.29, 1.36]}  % experiments/FIGS/data/paired_v1err/finetune/ratios.json :: ratios[19].ratio, ratios[19].ci95 :: sha256 3e5cf25e859b68c2 :: 17 over 162, leg1 N10000 1-macro_auc: paired geometric mean [95% interval]
\newcommand{\PairedFtJciiEFourOnesevenOverFourthree}{1.10~[1.07, 1.13]}  % experiments/FIGS/data/paired_v1err/finetune/ratios.json :: ratios[20].ratio, ratios[20].ci95 :: sha256 3e5cf25e859b68c2 :: 17 over 43, leg1 N10000 1-macro_auc: paired geometric mean [95% interval]
\newcommand{\PairedFtJciiEFourOnesixtwoMassOverOnesixtwo}{1.00~[0.94, 1.06]}  % experiments/FIGS/data/paired_v1err/finetune/ratios.json :: ratios[21].ratio, ratios[21].ci95 :: sha256 3e5cf25e859b68c2 :: 162+mass over 162, leg1 N10000 1-macro_auc: paired geometric mean [95% interval]
\newcommand{\PairedFtJciiEFourOnesevenMassOverOneseven}{0.990~[0.976, 1.004]}  % experiments/FIGS/data/paired_v1err/finetune/ratios.json :: ratios[22].ratio, ratios[22].ci95 :: sha256 3e5cf25e859b68c2 :: 17+mass over 17, leg1 N10000 1-macro_auc: paired geometric mean [95% interval]
\newcommand{\PairedFtJciiEFourRandomOverOneeighteight}{1.15~[1.09, 1.22]}  % experiments/FIGS/data/paired_v1err/finetune/ratios.json :: ratios[23].ratio, ratios[23].ci95 :: sha256 3e5cf25e859b68c2 :: random over 188, leg1 N10000 1-macro_auc: paired geometric mean [95% interval]
\newcommand{\PairedFtJciiEFourRandomOverOnesixtwo}{1.18~[1.13, 1.22]}  % experiments/FIGS/data/paired_v1err/finetune/ratios.json :: ratios[24].ratio, ratios[24].ci95 :: sha256 3e5cf25e859b68c2 :: random over 162, leg1 N10000 1-macro_auc: paired geometric mean [95% interval]
\newcommand{\PairedFtJciiEFourRandomOverFourthree}{0.97~[0.96, 0.99]}  % experiments/FIGS/data/paired_v1err/finetune/ratios.json :: ratios[25].ratio, ratios[25].ci95 :: sha256 3e5cf25e859b68c2 :: random over 43, leg1 N10000 1-macro_auc: paired geometric mean [95% interval]
\newcommand{\PairedFtJciiEFourRandomOverOneseven}{0.88~[0.82, 0.95]}  % experiments/FIGS/data/paired_v1err/finetune/ratios.json :: ratios[26].ratio, ratios[26].ci95 :: sha256 3e5cf25e859b68c2 :: random over 17, leg1 N10000 1-macro_auc: paired geometric mean [95% interval]
\newcommand{\PairedFtJciiEFourRandomDrawTwoOverRandomDrawOne}{1.00~[0.98, 1.02]}  % experiments/FIGS/data/paired_v1err/finetune/ratios.json :: ratios[27].ratio, ratios[27].ci95 :: sha256 3e5cf25e859b68c2 :: random draw 2 over random draw 1, leg1 N10000 1-macro_auc: paired geometric mean [95% interval]
\newcommand{\PairedFtJciiEFourRandomDrawThreeOverRandomDrawOne}{0.97~[0.95, 0.98]}  % experiments/FIGS/data/paired_v1err/finetune/ratios.json :: ratios[28].ratio, ratios[28].ci95 :: sha256 3e5cf25e859b68c2 :: random draw 3 over random draw 1, leg1 N10000 1-macro_auc: paired geometric mean [95% interval]
\newcommand{\PairedFtJciiEFourRandomDrawThreeOverRandomDrawTwo}{0.97~[0.95, 0.98]}  % experiments/FIGS/data/paired_v1err/finetune/ratios.json :: ratios[29].ratio, ratios[29].ci95 :: sha256 3e5cf25e859b68c2 :: random draw 3 over random draw 2, leg1 N10000 1-macro_auc: paired geometric mean [95% interval]
\newcommand{\PairedFtJciiEFiveOnesixtwoOverOneeighteight}{1.00~[0.98, 1.01]}  % experiments/FIGS/data/paired_v1err/finetune/ratios.json :: ratios[30].ratio, ratios[30].ci95 :: sha256 3e5cf25e859b68c2 :: 162 over 188, leg1 N100000 1-macro_auc: paired geometric mean [95% interval]
\newcommand{\PairedFtJciiEFiveFourthreeOverOneeighteight}{1.09~[1.06, 1.11]}  % experiments/FIGS/data/paired_v1err/finetune/ratios.json :: ratios[31].ratio, ratios[31].ci95 :: sha256 3e5cf25e859b68c2 :: 43 over 188, leg1 N100000 1-macro_auc: paired geometric mean [95% interval]
\newcommand{\PairedFtJciiEFiveOnesevenOverOneeighteight}{1.22~[1.20, 1.23]}  % experiments/FIGS/data/paired_v1err/finetune/ratios.json :: ratios[32].ratio, ratios[32].ci95 :: sha256 3e5cf25e859b68c2 :: 17 over 188, leg1 N100000 1-macro_auc: paired geometric mean [95% interval]
\newcommand{\PairedFtJciiEFiveFourthreeOverOnesixtwo}{1.09~[1.07, 1.11]}  % experiments/FIGS/data/paired_v1err/finetune/ratios.json :: ratios[33].ratio, ratios[33].ci95 :: sha256 3e5cf25e859b68c2 :: 43 over 162, leg1 N100000 1-macro_auc: paired geometric mean [95% interval]
\newcommand{\PairedFtJciiEFiveOnesevenOverOnesixtwo}{1.22~[1.20, 1.24]}  % experiments/FIGS/data/paired_v1err/finetune/ratios.json :: ratios[34].ratio, ratios[34].ci95 :: sha256 3e5cf25e859b68c2 :: 17 over 162, leg1 N100000 1-macro_auc: paired geometric mean [95% interval]
\newcommand{\PairedFtJciiEFiveOnesevenOverFourthree}{1.12~[1.10, 1.14]}  % experiments/FIGS/data/paired_v1err/finetune/ratios.json :: ratios[35].ratio, ratios[35].ci95 :: sha256 3e5cf25e859b68c2 :: 17 over 43, leg1 N100000 1-macro_auc: paired geometric mean [95% interval]
\newcommand{\PairedFtJciiEFiveOnesixtwoMassOverOnesixtwo}{1.018~[1.002, 1.035]}  % experiments/FIGS/data/paired_v1err/finetune/ratios.json :: ratios[36].ratio, ratios[36].ci95 :: sha256 3e5cf25e859b68c2 :: 162+mass over 162, leg1 N100000 1-macro_auc: paired geometric mean [95% interval]
\newcommand{\PairedFtJciiEFiveOnesevenMassOverOneseven}{1.03~[1.01, 1.06]}  % experiments/FIGS/data/paired_v1err/finetune/ratios.json :: ratios[37].ratio, ratios[37].ci95 :: sha256 3e5cf25e859b68c2 :: 17+mass over 17, leg1 N100000 1-macro_auc: paired geometric mean [95% interval]
\newcommand{\PairedFtJciiEFiveRandomOverOneeighteight}{1.15~[1.09, 1.20]}  % experiments/FIGS/data/paired_v1err/finetune/ratios.json :: ratios[38].ratio, ratios[38].ci95 :: sha256 3e5cf25e859b68c2 :: random over 188, leg1 N100000 1-macro_auc: paired geometric mean [95% interval]
\newcommand{\PairedFtJciiEFiveRandomOverOnesixtwo}{1.15~[1.09, 1.20]}  % experiments/FIGS/data/paired_v1err/finetune/ratios.json :: ratios[39].ratio, ratios[39].ci95 :: sha256 3e5cf25e859b68c2 :: random over 162, leg1 N100000 1-macro_auc: paired geometric mean [95% interval]
\newcommand{\PairedFtJciiEFiveRandomOverFourthree}{1.05~[1.04, 1.06]}  % experiments/FIGS/data/paired_v1err/finetune/ratios.json :: ratios[40].ratio, ratios[40].ci95 :: sha256 3e5cf25e859b68c2 :: random over 43, leg1 N100000 1-macro_auc: paired geometric mean [95% interval]
\newcommand{\PairedFtJciiEFiveRandomOverOneseven}{0.95~[0.93, 0.96]}  % experiments/FIGS/data/paired_v1err/finetune/ratios.json :: ratios[41].ratio, ratios[41].ci95 :: sha256 3e5cf25e859b68c2 :: random over 17, leg1 N100000 1-macro_auc: paired geometric mean [95% interval]
\newcommand{\PairedFtJciiEFiveRandomDrawTwoOverRandomDrawOne}{0.98~[0.97, 0.99]}  % experiments/FIGS/data/paired_v1err/finetune/ratios.json :: ratios[42].ratio, ratios[42].ci95 :: sha256 3e5cf25e859b68c2 :: random draw 2 over random draw 1, leg1 N100000 1-macro_auc: paired geometric mean [95% interval]
\newcommand{\PairedFtJciiEFiveRandomDrawThreeOverRandomDrawOne}{1.00~[0.99, 1.02]}  % experiments/FIGS/data/paired_v1err/finetune/ratios.json :: ratios[43].ratio, ratios[43].ci95 :: sha256 3e5cf25e859b68c2 :: random draw 3 over random draw 1, leg1 N100000 1-macro_auc: paired geometric mean [95% interval]
\newcommand{\PairedFtJciiEFiveRandomDrawThreeOverRandomDrawTwo}{1.03~[1.01, 1.04]}  % experiments/FIGS/data/paired_v1err/finetune/ratios.json :: ratios[44].ratio, ratios[44].ci95 :: sha256 3e5cf25e859b68c2 :: random draw 3 over random draw 2, leg1 N100000 1-macro_auc: paired geometric mean [95% interval]
\newcommand{\PairedFtJciiESixOnesixtwoOverOneeighteight}{1.002~[0.998, 1.005]}  % experiments/FIGS/data/paired_v1err/finetune/ratios.json :: ratios[45].ratio, ratios[45].ci95 :: sha256 3e5cf25e859b68c2 :: 162 over 188, leg1 N1000000 1-macro_auc: paired geometric mean [95% interval]
\newcommand{\PairedFtJciiESixFourthreeOverOneeighteight}{1.045~[1.037, 1.053]}  % experiments/FIGS/data/paired_v1err/finetune/ratios.json :: ratios[46].ratio, ratios[46].ci95 :: sha256 3e5cf25e859b68c2 :: 43 over 188, leg1 N1000000 1-macro_auc: paired geometric mean [95% interval]
\newcommand{\PairedFtJciiESixOnesevenOverOneeighteight}{1.14~[1.13, 1.16]}  % experiments/FIGS/data/paired_v1err/finetune/ratios.json :: ratios[47].ratio, ratios[47].ci95 :: sha256 3e5cf25e859b68c2 :: 17 over 188, leg1 N1000000 1-macro_auc: paired geometric mean [95% interval]
\newcommand{\PairedFtJciiESixFourthreeOverOnesixtwo}{1.044~[1.036, 1.051]}  % experiments/FIGS/data/paired_v1err/finetune/ratios.json :: ratios[48].ratio, ratios[48].ci95 :: sha256 3e5cf25e859b68c2 :: 43 over 162, leg1 N1000000 1-macro_auc: paired geometric mean [95% interval]
\newcommand{\PairedFtJciiESixOnesevenOverOnesixtwo}{1.14~[1.13, 1.15]}  % experiments/FIGS/data/paired_v1err/finetune/ratios.json :: ratios[49].ratio, ratios[49].ci95 :: sha256 3e5cf25e859b68c2 :: 17 over 162, leg1 N1000000 1-macro_auc: paired geometric mean [95% interval]
\newcommand{\PairedFtJciiESixOnesevenOverFourthree}{1.09~[1.08, 1.11]}  % experiments/FIGS/data/paired_v1err/finetune/ratios.json :: ratios[50].ratio, ratios[50].ci95 :: sha256 3e5cf25e859b68c2 :: 17 over 43, leg1 N1000000 1-macro_auc: paired geometric mean [95% interval]
\newcommand{\PairedFtJciiESixOnesixtwoMassOverOnesixtwo}{1.019~[1.014, 1.023]}  % experiments/FIGS/data/paired_v1err/finetune/ratios.json :: ratios[51].ratio, ratios[51].ci95 :: sha256 3e5cf25e859b68c2 :: 162+mass over 162, leg1 N1000000 1-macro_auc: paired geometric mean [95% interval]
\newcommand{\PairedFtJciiESixOnesevenMassOverOneseven}{1.03~[1.02, 1.04]}  % experiments/FIGS/data/paired_v1err/finetune/ratios.json :: ratios[52].ratio, ratios[52].ci95 :: sha256 3e5cf25e859b68c2 :: 17+mass over 17, leg1 N1000000 1-macro_auc: paired geometric mean [95% interval]
\newcommand{\PairedFtJciiESixRandomOverOneeighteight}{1.073~[1.067, 1.080]}  % experiments/FIGS/data/paired_v1err/finetune/ratios.json :: ratios[53].ratio, ratios[53].ci95 :: sha256 3e5cf25e859b68c2 :: random over 188, leg1 N1000000 1-macro_auc: paired geometric mean [95% interval]
\newcommand{\PairedFtJciiESixRandomOverOnesixtwo}{1.072~[1.065, 1.078]}  % experiments/FIGS/data/paired_v1err/finetune/ratios.json :: ratios[54].ratio, ratios[54].ci95 :: sha256 3e5cf25e859b68c2 :: random over 162, leg1 N1000000 1-macro_auc: paired geometric mean [95% interval]
\newcommand{\PairedFtJciiESixRandomOverFourthree}{1.03~[1.02, 1.04]}  % experiments/FIGS/data/paired_v1err/finetune/ratios.json :: ratios[55].ratio, ratios[55].ci95 :: sha256 3e5cf25e859b68c2 :: random over 43, leg1 N1000000 1-macro_auc: paired geometric mean [95% interval]
\newcommand{\PairedFtJciiESixRandomOverOneseven}{0.94~[0.93, 0.95]}  % experiments/FIGS/data/paired_v1err/finetune/ratios.json :: ratios[56].ratio, ratios[56].ci95 :: sha256 3e5cf25e859b68c2 :: random over 17, leg1 N1000000 1-macro_auc: paired geometric mean [95% interval]
\newcommand{\PairedFtJciiESixRandomDrawTwoOverRandomDrawOne}{1.00~[0.99, 1.01]}  % experiments/FIGS/data/paired_v1err/finetune/ratios.json :: ratios[57].ratio, ratios[57].ci95 :: sha256 3e5cf25e859b68c2 :: random draw 2 over random draw 1, leg1 N1000000 1-macro_auc: paired geometric mean [95% interval]
\newcommand{\PairedFtJciiESixRandomDrawThreeOverRandomDrawOne}{0.996~[0.989, 1.003]}  % experiments/FIGS/data/paired_v1err/finetune/ratios.json :: ratios[58].ratio, ratios[58].ci95 :: sha256 3e5cf25e859b68c2 :: random draw 3 over random draw 1, leg1 N1000000 1-macro_auc: paired geometric mean [95% interval]
\newcommand{\PairedFtJciiESixRandomDrawThreeOverRandomDrawTwo}{0.997~[0.989, 1.004]}  % experiments/FIGS/data/paired_v1err/finetune/ratios.json :: ratios[59].ratio, ratios[59].ci95 :: sha256 3e5cf25e859b68c2 :: random draw 3 over random draw 2, leg1 N1000000 1-macro_auc: paired geometric mean [95% interval]
\newcommand{\PairedFtJcEThreeOnesixtwoOverOneeighteight}{0.96~[0.83, 1.11]}  % experiments/FIGS/data/paired_v1err/finetune/ratios.json :: ratios[60].ratio, ratios[60].ci95 :: sha256 3e5cf25e859b68c2 :: 162 over 188, leg2 N1000 1-macro_auc: paired geometric mean [95% interval]
\newcommand{\PairedFtJcEThreeFourthreeOverOneeighteight}{1.01~[0.84, 1.21]}  % experiments/FIGS/data/paired_v1err/finetune/ratios.json :: ratios[61].ratio, ratios[61].ci95 :: sha256 3e5cf25e859b68c2 :: 43 over 188, leg2 N1000 1-macro_auc: paired geometric mean [95% interval]
\newcommand{\PairedFtJcEThreeOnesevenOverOneeighteight}{1.07~[0.93, 1.23]}  % experiments/FIGS/data/paired_v1err/finetune/ratios.json :: ratios[62].ratio, ratios[62].ci95 :: sha256 3e5cf25e859b68c2 :: 17 over 188, leg2 N1000 1-macro_auc: paired geometric mean [95% interval]
\newcommand{\PairedFtJcEThreeFourthreeOverOnesixtwo}{1.05~[0.96, 1.15]}  % experiments/FIGS/data/paired_v1err/finetune/ratios.json :: ratios[63].ratio, ratios[63].ci95 :: sha256 3e5cf25e859b68c2 :: 43 over 162, leg2 N1000 1-macro_auc: paired geometric mean [95% interval]
\newcommand{\PairedFtJcEThreeOnesevenOverOnesixtwo}{1.12~[1.07, 1.17]}  % experiments/FIGS/data/paired_v1err/finetune/ratios.json :: ratios[64].ratio, ratios[64].ci95 :: sha256 3e5cf25e859b68c2 :: 17 over 162, leg2 N1000 1-macro_auc: paired geometric mean [95% interval]
\newcommand{\PairedFtJcEThreeOnesevenOverFourthree}{1.06~[0.98, 1.15]}  % experiments/FIGS/data/paired_v1err/finetune/ratios.json :: ratios[65].ratio, ratios[65].ci95 :: sha256 3e5cf25e859b68c2 :: 17 over 43, leg2 N1000 1-macro_auc: paired geometric mean [95% interval]
\newcommand{\PairedFtJcEThreeOnesixtwoMassOverOnesixtwo}{0.91~[0.88, 0.94]}  % experiments/FIGS/data/paired_v1err/finetune/ratios.json :: ratios[66].ratio, ratios[66].ci95 :: sha256 3e5cf25e859b68c2 :: 162+mass over 162, leg2 N1000 1-macro_auc: paired geometric mean [95% interval]
\newcommand{\PairedFtJcEThreeOnesevenMassOverOneseven}{1.02~[0.92, 1.13]}  % experiments/FIGS/data/paired_v1err/finetune/ratios.json :: ratios[67].ratio, ratios[67].ci95 :: sha256 3e5cf25e859b68c2 :: 17+mass over 17, leg2 N1000 1-macro_auc: paired geometric mean [95% interval]
\newcommand{\PairedFtJcEThreeRandomOverOneeighteight}{1.01~[0.78, 1.32]}  % experiments/FIGS/data/paired_v1err/finetune/ratios.json :: ratios[68].ratio, ratios[68].ci95 :: sha256 3e5cf25e859b68c2 :: random over 188, leg2 N1000 1-macro_auc: paired geometric mean [95% interval]
\newcommand{\PairedFtJcEThreeRandomOverOnesixtwo}{1.12~[1.04, 1.20]}  % experiments/FIGS/data/paired_v1err/finetune/ratios.json :: ratios[69].ratio, ratios[69].ci95 :: sha256 3e5cf25e859b68c2 :: random over 162, leg2 N1000 1-macro_auc: paired geometric mean [95% interval]
\newcommand{\PairedFtJcEThreeRandomOverFourthree}{1.06~[0.79, 1.41]}  % experiments/FIGS/data/paired_v1err/finetune/ratios.json :: ratios[70].ratio, ratios[70].ci95 :: sha256 3e5cf25e859b68c2 :: random over 43, leg2 N1000 1-macro_auc: paired geometric mean [95% interval]
\newcommand{\PairedFtJcEThreeRandomOverOneseven}{0.98~[0.84, 1.13]}  % experiments/FIGS/data/paired_v1err/finetune/ratios.json :: ratios[71].ratio, ratios[71].ci95 :: sha256 3e5cf25e859b68c2 :: random over 17, leg2 N1000 1-macro_auc: paired geometric mean [95% interval]
\newcommand{\PairedFtJcEThreeRandomDrawTwoOverRandomDrawOne}{0.954~[0.949, 0.958]}  % experiments/FIGS/data/paired_v1err/finetune/ratios.json :: ratios[72].ratio, ratios[72].ci95 :: sha256 3e5cf25e859b68c2 :: random draw 2 over random draw 1, leg2 N1000 1-macro_auc: paired geometric mean [95% interval]
\newcommand{\PairedFtJcEThreeRandomDrawThreeOverRandomDrawOne}{1.013~[1.009, 1.017]}  % experiments/FIGS/data/paired_v1err/finetune/ratios.json :: ratios[73].ratio, ratios[73].ci95 :: sha256 3e5cf25e859b68c2 :: random draw 3 over random draw 1, leg2 N1000 1-macro_auc: paired geometric mean [95% interval]
\newcommand{\PairedFtJcEThreeRandomDrawThreeOverRandomDrawTwo}{1.062~[1.057, 1.067]}  % experiments/FIGS/data/paired_v1err/finetune/ratios.json :: ratios[74].ratio, ratios[74].ci95 :: sha256 3e5cf25e859b68c2 :: random draw 3 over random draw 2, leg2 N1000 1-macro_auc: paired geometric mean [95% interval]
\newcommand{\PairedFtJcEFourOnesixtwoOverOneeighteight}{0.98~[0.93, 1.02]}  % experiments/FIGS/data/paired_v1err/finetune/ratios.json :: ratios[75].ratio, ratios[75].ci95 :: sha256 3e5cf25e859b68c2 :: 162 over 188, leg2 N10000 1-macro_auc: paired geometric mean [95% interval]
\newcommand{\PairedFtJcEFourFourthreeOverOneeighteight}{0.98~[0.93, 1.04]}  % experiments/FIGS/data/paired_v1err/finetune/ratios.json :: ratios[76].ratio, ratios[76].ci95 :: sha256 3e5cf25e859b68c2 :: 43 over 188, leg2 N10000 1-macro_auc: paired geometric mean [95% interval]
\newcommand{\PairedFtJcEFourOnesevenOverOneeighteight}{1.07~[1.02, 1.11]}  % experiments/FIGS/data/paired_v1err/finetune/ratios.json :: ratios[77].ratio, ratios[77].ci95 :: sha256 3e5cf25e859b68c2 :: 17 over 188, leg2 N10000 1-macro_auc: paired geometric mean [95% interval]
\newcommand{\PairedFtJcEFourFourthreeOverOnesixtwo}{1.01~[0.97, 1.04]}  % experiments/FIGS/data/paired_v1err/finetune/ratios.json :: ratios[78].ratio, ratios[78].ci95 :: sha256 3e5cf25e859b68c2 :: 43 over 162, leg2 N10000 1-macro_auc: paired geometric mean [95% interval]
\newcommand{\PairedFtJcEFourOnesevenOverOnesixtwo}{1.09~[1.06, 1.12]}  % experiments/FIGS/data/paired_v1err/finetune/ratios.json :: ratios[79].ratio, ratios[79].ci95 :: sha256 3e5cf25e859b68c2 :: 17 over 162, leg2 N10000 1-macro_auc: paired geometric mean [95% interval]
\newcommand{\PairedFtJcEFourOnesevenOverFourthree}{1.08~[1.07, 1.10]}  % experiments/FIGS/data/paired_v1err/finetune/ratios.json :: ratios[80].ratio, ratios[80].ci95 :: sha256 3e5cf25e859b68c2 :: 17 over 43, leg2 N10000 1-macro_auc: paired geometric mean [95% interval]
\newcommand{\PairedFtJcEFourOnesixtwoMassOverOnesixtwo}{1.01~[0.98, 1.04]}  % experiments/FIGS/data/paired_v1err/finetune/ratios.json :: ratios[81].ratio, ratios[81].ci95 :: sha256 3e5cf25e859b68c2 :: 162+mass over 162, leg2 N10000 1-macro_auc: paired geometric mean [95% interval]
\newcommand{\PairedFtJcEFourOnesevenMassOverOneseven}{1.05~[1.02, 1.08]}  % experiments/FIGS/data/paired_v1err/finetune/ratios.json :: ratios[82].ratio, ratios[82].ci95 :: sha256 3e5cf25e859b68c2 :: 17+mass over 17, leg2 N10000 1-macro_auc: paired geometric mean [95% interval]
\newcommand{\PairedFtJcEFourRandomOverOneeighteight}{1.08~[1.03, 1.14]}  % experiments/FIGS/data/paired_v1err/finetune/ratios.json :: ratios[83].ratio, ratios[83].ci95 :: sha256 3e5cf25e859b68c2 :: random over 188, leg2 N10000 1-macro_auc: paired geometric mean [95% interval]
\newcommand{\PairedFtJcEFourRandomOverOnesixtwo}{1.11~[0.99, 1.23]}  % experiments/FIGS/data/paired_v1err/finetune/ratios.json :: ratios[84].ratio, ratios[84].ci95 :: sha256 3e5cf25e859b68c2 :: random over 162, leg2 N10000 1-macro_auc: paired geometric mean [95% interval]
\newcommand{\PairedFtJcEFourRandomOverFourthree}{1.10~[0.99, 1.21]}  % experiments/FIGS/data/paired_v1err/finetune/ratios.json :: ratios[85].ratio, ratios[85].ci95 :: sha256 3e5cf25e859b68c2 :: random over 43, leg2 N10000 1-macro_auc: paired geometric mean [95% interval]
\newcommand{\PairedFtJcEFourRandomOverOneseven}{1.02~[0.94, 1.10]}  % experiments/FIGS/data/paired_v1err/finetune/ratios.json :: ratios[86].ratio, ratios[86].ci95 :: sha256 3e5cf25e859b68c2 :: random over 17, leg2 N10000 1-macro_auc: paired geometric mean [95% interval]
\newcommand{\PairedFtJcEFourRandomDrawTwoOverRandomDrawOne}{0.960~[0.955, 0.965]}  % experiments/FIGS/data/paired_v1err/finetune/ratios.json :: ratios[87].ratio, ratios[87].ci95 :: sha256 3e5cf25e859b68c2 :: random draw 2 over random draw 1, leg2 N10000 1-macro_auc: paired geometric mean [95% interval]
\newcommand{\PairedFtJcEFourRandomDrawThreeOverRandomDrawOne}{1.008~[1.004, 1.012]}  % experiments/FIGS/data/paired_v1err/finetune/ratios.json :: ratios[88].ratio, ratios[88].ci95 :: sha256 3e5cf25e859b68c2 :: random draw 3 over random draw 1, leg2 N10000 1-macro_auc: paired geometric mean [95% interval]
\newcommand{\PairedFtJcEFourRandomDrawThreeOverRandomDrawTwo}{1.050~[1.045, 1.055]}  % experiments/FIGS/data/paired_v1err/finetune/ratios.json :: ratios[89].ratio, ratios[89].ci95 :: sha256 3e5cf25e859b68c2 :: random draw 3 over random draw 2, leg2 N10000 1-macro_auc: paired geometric mean [95% interval]
\newcommand{\PairedFtJcEFiveOnesixtwoOverOneeighteight}{0.983~[0.968, 0.998]}  % experiments/FIGS/data/paired_v1err/finetune/ratios.json :: ratios[90].ratio, ratios[90].ci95 :: sha256 3e5cf25e859b68c2 :: 162 over 188, leg2 N100000 1-macro_auc: paired geometric mean [95% interval]
\newcommand{\PairedFtJcEFiveFourthreeOverOneeighteight}{1.015~[1.002, 1.029]}  % experiments/FIGS/data/paired_v1err/finetune/ratios.json :: ratios[91].ratio, ratios[91].ci95 :: sha256 3e5cf25e859b68c2 :: 43 over 188, leg2 N100000 1-macro_auc: paired geometric mean [95% interval]
\newcommand{\PairedFtJcEFiveOnesevenOverOneeighteight}{1.07~[1.06, 1.09]}  % experiments/FIGS/data/paired_v1err/finetune/ratios.json :: ratios[92].ratio, ratios[92].ci95 :: sha256 3e5cf25e859b68c2 :: 17 over 188, leg2 N100000 1-macro_auc: paired geometric mean [95% interval]
\newcommand{\PairedFtJcEFiveFourthreeOverOnesixtwo}{1.03~[1.02, 1.05]}  % experiments/FIGS/data/paired_v1err/finetune/ratios.json :: ratios[93].ratio, ratios[93].ci95 :: sha256 3e5cf25e859b68c2 :: 43 over 162, leg2 N100000 1-macro_auc: paired geometric mean [95% interval]
\newcommand{\PairedFtJcEFiveOnesevenOverOnesixtwo}{1.091~[1.088, 1.094]}  % experiments/FIGS/data/paired_v1err/finetune/ratios.json :: ratios[94].ratio, ratios[94].ci95 :: sha256 3e5cf25e859b68c2 :: 17 over 162, leg2 N100000 1-macro_auc: paired geometric mean [95% interval]
\newcommand{\PairedFtJcEFiveOnesevenOverFourthree}{1.06~[1.04, 1.07]}  % experiments/FIGS/data/paired_v1err/finetune/ratios.json :: ratios[95].ratio, ratios[95].ci95 :: sha256 3e5cf25e859b68c2 :: 17 over 43, leg2 N100000 1-macro_auc: paired geometric mean [95% interval]
\newcommand{\PairedFtJcEFiveOnesixtwoMassOverOnesixtwo}{1.016~[1.008, 1.024]}  % experiments/FIGS/data/paired_v1err/finetune/ratios.json :: ratios[96].ratio, ratios[96].ci95 :: sha256 3e5cf25e859b68c2 :: 162+mass over 162, leg2 N100000 1-macro_auc: paired geometric mean [95% interval]
\newcommand{\PairedFtJcEFiveOnesevenMassOverOneseven}{1.02~[1.01, 1.03]}  % experiments/FIGS/data/paired_v1err/finetune/ratios.json :: ratios[97].ratio, ratios[97].ci95 :: sha256 3e5cf25e859b68c2 :: 17+mass over 17, leg2 N100000 1-macro_auc: paired geometric mean [95% interval]
\newcommand{\PairedFtJcEFiveRandomOverOneeighteight}{1.044~[1.039, 1.049]}  % experiments/FIGS/data/paired_v1err/finetune/ratios.json :: ratios[98].ratio, ratios[98].ci95 :: sha256 3e5cf25e859b68c2 :: random over 188, leg2 N100000 1-macro_auc: paired geometric mean [95% interval]
\newcommand{\PairedFtJcEFiveRandomOverOnesixtwo}{1.06~[1.04, 1.09]}  % experiments/FIGS/data/paired_v1err/finetune/ratios.json :: ratios[99].ratio, ratios[99].ci95 :: sha256 3e5cf25e859b68c2 :: random over 162, leg2 N100000 1-macro_auc: paired geometric mean [95% interval]
\newcommand{\PairedFtJcEFiveRandomOverFourthree}{1.03~[1.01, 1.04]}  % experiments/FIGS/data/paired_v1err/finetune/ratios.json :: ratios[100].ratio, ratios[100].ci95 :: sha256 3e5cf25e859b68c2 :: random over 43, leg2 N100000 1-macro_auc: paired geometric mean [95% interval]
\newcommand{\PairedFtJcEFiveRandomOverOneseven}{0.97~[0.95, 0.99]}  % experiments/FIGS/data/paired_v1err/finetune/ratios.json :: ratios[101].ratio, ratios[101].ci95 :: sha256 3e5cf25e859b68c2 :: random over 17, leg2 N100000 1-macro_auc: paired geometric mean [95% interval]
\newcommand{\PairedFtJcEFiveRandomDrawTwoOverRandomDrawOne}{0.992~[0.988, 0.997]}  % experiments/FIGS/data/paired_v1err/finetune/ratios.json :: ratios[102].ratio, ratios[102].ci95 :: sha256 3e5cf25e859b68c2 :: random draw 2 over random draw 1, leg2 N100000 1-macro_auc: paired geometric mean [95% interval]
\newcommand{\PairedFtJcEFiveRandomDrawThreeOverRandomDrawOne}{1.007~[1.003, 1.011]}  % experiments/FIGS/data/paired_v1err/finetune/ratios.json :: ratios[103].ratio, ratios[103].ci95 :: sha256 3e5cf25e859b68c2 :: random draw 3 over random draw 1, leg2 N100000 1-macro_auc: paired geometric mean [95% interval]
\newcommand{\PairedFtJcEFiveRandomDrawThreeOverRandomDrawTwo}{1.015~[1.011, 1.020]}  % experiments/FIGS/data/paired_v1err/finetune/ratios.json :: ratios[104].ratio, ratios[104].ci95 :: sha256 3e5cf25e859b68c2 :: random draw 3 over random draw 2, leg2 N100000 1-macro_auc: paired geometric mean [95% interval]
\newcommand{\PairedFtJcESixOnesixtwoOverOneeighteight}{1.00~[0.99, 1.01]}  % experiments/FIGS/data/paired_v1err/finetune/ratios.json :: ratios[105].ratio, ratios[105].ci95 :: sha256 3e5cf25e859b68c2 :: 162 over 188, leg2 N1000000 1-macro_auc: paired geometric mean [95% interval]
\newcommand{\PairedFtJcESixFourthreeOverOneeighteight}{1.02~[1.01, 1.03]}  % experiments/FIGS/data/paired_v1err/finetune/ratios.json :: ratios[106].ratio, ratios[106].ci95 :: sha256 3e5cf25e859b68c2 :: 43 over 188, leg2 N1000000 1-macro_auc: paired geometric mean [95% interval]
\newcommand{\PairedFtJcESixOnesevenOverOneeighteight}{1.08~[1.07, 1.09]}  % experiments/FIGS/data/paired_v1err/finetune/ratios.json :: ratios[107].ratio, ratios[107].ci95 :: sha256 3e5cf25e859b68c2 :: 17 over 188, leg2 N1000000 1-macro_auc: paired geometric mean [95% interval]
\newcommand{\PairedFtJcESixFourthreeOverOnesixtwo}{1.02~[1.01, 1.03]}  % experiments/FIGS/data/paired_v1err/finetune/ratios.json :: ratios[108].ratio, ratios[108].ci95 :: sha256 3e5cf25e859b68c2 :: 43 over 162, leg2 N1000000 1-macro_auc: paired geometric mean [95% interval]
\newcommand{\PairedFtJcESixOnesevenOverOnesixtwo}{1.08~[1.07, 1.09]}  % experiments/FIGS/data/paired_v1err/finetune/ratios.json :: ratios[109].ratio, ratios[109].ci95 :: sha256 3e5cf25e859b68c2 :: 17 over 162, leg2 N1000000 1-macro_auc: paired geometric mean [95% interval]
\newcommand{\PairedFtJcESixOnesevenOverFourthree}{1.06~[1.05, 1.07]}  % experiments/FIGS/data/paired_v1err/finetune/ratios.json :: ratios[110].ratio, ratios[110].ci95 :: sha256 3e5cf25e859b68c2 :: 17 over 43, leg2 N1000000 1-macro_auc: paired geometric mean [95% interval]
\newcommand{\PairedFtJcESixOnesixtwoMassOverOnesixtwo}{1.012~[1.004, 1.020]}  % experiments/FIGS/data/paired_v1err/finetune/ratios.json :: ratios[111].ratio, ratios[111].ci95 :: sha256 3e5cf25e859b68c2 :: 162+mass over 162, leg2 N1000000 1-macro_auc: paired geometric mean [95% interval]
\newcommand{\PairedFtJcESixOnesevenMassOverOneseven}{1.02~[1.01, 1.04]}  % experiments/FIGS/data/paired_v1err/finetune/ratios.json :: ratios[112].ratio, ratios[112].ci95 :: sha256 3e5cf25e859b68c2 :: 17+mass over 17, leg2 N1000000 1-macro_auc: paired geometric mean [95% interval]
\newcommand{\PairedFtJcESixRandomOverOneeighteight}{1.046~[1.038, 1.053]}  % experiments/FIGS/data/paired_v1err/finetune/ratios.json :: ratios[113].ratio, ratios[113].ci95 :: sha256 3e5cf25e859b68c2 :: random over 188, leg2 N1000000 1-macro_auc: paired geometric mean [95% interval]
\newcommand{\PairedFtJcESixRandomOverOnesixtwo}{1.047~[1.038, 1.055]}  % experiments/FIGS/data/paired_v1err/finetune/ratios.json :: ratios[114].ratio, ratios[114].ci95 :: sha256 3e5cf25e859b68c2 :: random over 162, leg2 N1000000 1-macro_auc: paired geometric mean [95% interval]
\newcommand{\PairedFtJcESixRandomOverFourthree}{1.03~[1.02, 1.04]}  % experiments/FIGS/data/paired_v1err/finetune/ratios.json :: ratios[115].ratio, ratios[115].ci95 :: sha256 3e5cf25e859b68c2 :: random over 43, leg2 N1000000 1-macro_auc: paired geometric mean [95% interval]
\newcommand{\PairedFtJcESixRandomOverOneseven}{0.97~[0.96, 0.98]}  % experiments/FIGS/data/paired_v1err/finetune/ratios.json :: ratios[116].ratio, ratios[116].ci95 :: sha256 3e5cf25e859b68c2 :: random over 17, leg2 N1000000 1-macro_auc: paired geometric mean [95% interval]
\newcommand{\PairedFtJcESixRandomDrawTwoOverRandomDrawOne}{1.001~[0.998, 1.004]}  % experiments/FIGS/data/paired_v1err/finetune/ratios.json :: ratios[117].ratio, ratios[117].ci95 :: sha256 3e5cf25e859b68c2 :: random draw 2 over random draw 1, leg2 N1000000 1-macro_auc: paired geometric mean [95% interval]
\newcommand{\PairedFtJcESixRandomDrawThreeOverRandomDrawOne}{1.005~[1.002, 1.008]}  % experiments/FIGS/data/paired_v1err/finetune/ratios.json :: ratios[118].ratio, ratios[118].ci95 :: sha256 3e5cf25e859b68c2 :: random draw 3 over random draw 1, leg2 N1000000 1-macro_auc: paired geometric mean [95% interval]
\newcommand{\PairedFtJcESixRandomDrawThreeOverRandomDrawTwo}{1.004~[1.001, 1.007]}  % experiments/FIGS/data/paired_v1err/finetune/ratios.json :: ratios[119].ratio, ratios[119].ci95 :: sha256 3e5cf25e859b68c2 :: random draw 3 over random draw 2, leg2 N1000000 1-macro_auc: paired geometric mean [95% interval]
\newcommand{\PairedMassResMlpOnesixtwoOverOneeighteight}{0.990~[0.977, 1.003]}  % experiments/FIGS/data/paired_v1err/mass/ratios.json :: ratios[0].ratio, ratios[0].ci95 :: sha256 d877a7f087260f64 :: 162 over 188, resolution mlp sigma_eff: paired geometric mean [95% interval]
\newcommand{\PairedMassResMlpFourthreeOverOneeighteight}{0.98~[0.96, 0.99]}  % experiments/FIGS/data/paired_v1err/mass/ratios.json :: ratios[1].ratio, ratios[1].ci95 :: sha256 d877a7f087260f64 :: 43 over 188, resolution mlp sigma_eff: paired geometric mean [95% interval]
\newcommand{\PairedMassResMlpOnesevenOverOneeighteight}{0.96~[0.94, 0.98]}  % experiments/FIGS/data/paired_v1err/mass/ratios.json :: ratios[2].ratio, ratios[2].ci95 :: sha256 d877a7f087260f64 :: 17 over 188, resolution mlp sigma_eff: paired geometric mean [95% interval]
\newcommand{\PairedMassResMlpFourthreeOverOnesixtwo}{0.989~[0.978, 0.999]}  % experiments/FIGS/data/paired_v1err/mass/ratios.json :: ratios[3].ratio, ratios[3].ci95 :: sha256 d877a7f087260f64 :: 43 over 162, resolution mlp sigma_eff: paired geometric mean [95% interval]
\newcommand{\PairedMassResMlpOnesevenOverOnesixtwo}{0.97~[0.94, 0.99]}  % experiments/FIGS/data/paired_v1err/mass/ratios.json :: ratios[4].ratio, ratios[4].ci95 :: sha256 d877a7f087260f64 :: 17 over 162, resolution mlp sigma_eff: paired geometric mean [95% interval]
\newcommand{\PairedMassResMlpOnesevenOverFourthree}{0.979~[0.961, 0.997]}  % experiments/FIGS/data/paired_v1err/mass/ratios.json :: ratios[5].ratio, ratios[5].ci95 :: sha256 d877a7f087260f64 :: 17 over 43, resolution mlp sigma_eff: paired geometric mean [95% interval]
\newcommand{\PairedMassResMlpOnesixtwoMassOverOnesixtwo}{0.80~[0.79, 0.81]}  % experiments/FIGS/data/paired_v1err/mass/ratios.json :: ratios[6].ratio, ratios[6].ci95 :: sha256 d877a7f087260f64 :: 162+mass over 162, resolution mlp sigma_eff: paired geometric mean [95% interval]
\newcommand{\PairedMassResMlpOnesevenMassOverOneseven}{0.83~[0.82, 0.84]}  % experiments/FIGS/data/paired_v1err/mass/ratios.json :: ratios[7].ratio, ratios[7].ci95 :: sha256 d877a7f087260f64 :: 17+mass over 17, resolution mlp sigma_eff: paired geometric mean [95% interval]
\newcommand{\PairedMassResRidgeOnesixtwoOverOneeighteight}{0.97~[0.96, 0.98]}  % experiments/FIGS/data/paired_v1err/mass/ratios.json :: ratios[8].ratio, ratios[8].ci95 :: sha256 d877a7f087260f64 :: 162 over 188, resolution ridge sigma_eff: paired geometric mean [95% interval]
\newcommand{\PairedMassResRidgeFourthreeOverOneeighteight}{0.95~[0.94, 0.96]}  % experiments/FIGS/data/paired_v1err/mass/ratios.json :: ratios[9].ratio, ratios[9].ci95 :: sha256 d877a7f087260f64 :: 43 over 188, resolution ridge sigma_eff: paired geometric mean [95% interval]
\newcommand{\PairedMassResRidgeOnesevenOverOneeighteight}{0.88~[0.85, 0.92]}  % experiments/FIGS/data/paired_v1err/mass/ratios.json :: ratios[10].ratio, ratios[10].ci95 :: sha256 d877a7f087260f64 :: 17 over 188, resolution ridge sigma_eff: paired geometric mean [95% interval]
\newcommand{\PairedMassResRidgeFourthreeOverOnesixtwo}{0.98~[0.97, 0.99]}  % experiments/FIGS/data/paired_v1err/mass/ratios.json :: ratios[11].ratio, ratios[11].ci95 :: sha256 d877a7f087260f64 :: 43 over 162, resolution ridge sigma_eff: paired geometric mean [95% interval]
\newcommand{\PairedMassResRidgeOnesevenOverOnesixtwo}{0.91~[0.87, 0.95]}  % experiments/FIGS/data/paired_v1err/mass/ratios.json :: ratios[12].ratio, ratios[12].ci95 :: sha256 d877a7f087260f64 :: 17 over 162, resolution ridge sigma_eff: paired geometric mean [95% interval]
\newcommand{\PairedMassResRidgeOnesevenOverFourthree}{0.93~[0.89, 0.97]}  % experiments/FIGS/data/paired_v1err/mass/ratios.json :: ratios[13].ratio, ratios[13].ci95 :: sha256 d877a7f087260f64 :: 17 over 43, resolution ridge sigma_eff: paired geometric mean [95% interval]
\newcommand{\PairedMassResRidgeOnesixtwoMassOverOnesixtwo}{0.678~[0.672, 0.684]}  % experiments/FIGS/data/paired_v1err/mass/ratios.json :: ratios[14].ratio, ratios[14].ci95 :: sha256 d877a7f087260f64 :: 162+mass over 162, resolution ridge sigma_eff: paired geometric mean [95% interval]
\newcommand{\PairedMassResRidgeOnesevenMassOverOneseven}{0.72~[0.69, 0.76]}  % experiments/FIGS/data/paired_v1err/mass/ratios.json :: ratios[15].ratio, ratios[15].ci95 :: sha256 d877a7f087260f64 :: 17+mass over 17, resolution ridge sigma_eff: paired geometric mean [95% interval]
\newcommand{\VcbNSignal}{1{,}659}  % experiments/FIGS/data/probe_ladder_vcbwindow_mlp2/sall.json :: tasks.bc_vs_rest.n_signal_test :: sha256 5b8b5c7c28602f93 :: X->bc test jets in the window
\newcommand{\VcbNBackground}{19{,}244}  % experiments/FIGS/data/probe_ladder_vcbwindow_mlp2/sall.json :: tasks.bc_vs_rest.n_background_test :: sha256 5b8b5c7c28602f93 :: bq, cs, bqq and QCD test jets in the window
\newcommand{\VcbEpsSixty}{60}  % experiments/FIGS/data/probe_ladder_vcbwindow_mlp2/sall.json :: tasks.bc_vs_rest.eps_s :: sha256 5b8b5c7c28602f93 :: signal efficiency, percent
\newcommand{\VcbOmaLinearOnesixtwo}{\ensuremath{(3.17 \pm 0.05)\times10^{-2}}}  % experiments/FIGS/data/probe_ladder_vcbwindow_mlp2/sall.json :: tasks.bc_vs_rest.arms.l162-s*.linear.log1m_auc :: sha256 5b8b5c7c28602f93 :: 1-AUC = exp of it, mean +- SD over 5 seeds
\newcommand{\VcbLogOmaLinearOnesixtwo}{\ensuremath{-}3.451}  % experiments/FIGS/data/probe_ladder_vcbwindow_mlp2/sall.json :: tasks.bc_vs_rest.arms.l162-s*.linear.log1m_auc :: sha256 5b8b5c7c28602f93 :: mean over seeds
\newcommand{\VcbRejLinearSixtyOnesixtwo}{207\,\ensuremath{\pm}\,19}  % experiments/FIGS/data/probe_ladder_vcbwindow_mlp2/sall.json :: tasks.bc_vs_rest.arms.l162-s*.linear.rejection_at['0.60'].rejection :: sha256 5b8b5c7c28602f93 :: background rejection at this signal efficiency, mean +- SD over seeds
\newcommand{\VcbBkgLeftLinearSixtyOnesixtwo}{93.4}  % experiments/FIGS/data/probe_ladder_vcbwindow_mlp2/sall.json :: tasks.bc_vs_rest.arms.l162-s*.linear.rejection_at['0.60'].n_bkg_pass :: sha256 5b8b5c7c28602f93 :: background jets passing the cut, mean over seeds
\newcommand{\VcbOmaLinearOneseven}{\ensuremath{(5.21 \pm 0.23)\times10^{-2}}}  % experiments/FIGS/data/probe_ladder_vcbwindow_mlp2/sall.json :: tasks.bc_vs_rest.arms.r16q1-s*.linear.log1m_auc :: sha256 5b8b5c7c28602f93 :: 1-AUC = exp of it, mean +- SD over 5 seeds
\newcommand{\VcbLogOmaLinearOneseven}{\ensuremath{-}2.955}  % experiments/FIGS/data/probe_ladder_vcbwindow_mlp2/sall.json :: tasks.bc_vs_rest.arms.r16q1-s*.linear.log1m_auc :: sha256 5b8b5c7c28602f93 :: mean over seeds
\newcommand{\VcbRejLinearSixtyOneseven}{67\,\ensuremath{\pm}\,14}  % experiments/FIGS/data/probe_ladder_vcbwindow_mlp2/sall.json :: tasks.bc_vs_rest.arms.r16q1-s*.linear.rejection_at['0.60'].rejection :: sha256 5b8b5c7c28602f93 :: background rejection at this signal efficiency, mean +- SD over seeds
\newcommand{\VcbBkgLeftLinearSixtyOneseven}{295.0}  % experiments/FIGS/data/probe_ladder_vcbwindow_mlp2/sall.json :: tasks.bc_vs_rest.arms.r16q1-s*.linear.rejection_at['0.60'].n_bkg_pass :: sha256 5b8b5c7c28602f93 :: background jets passing the cut, mean over seeds
\newcommand{\VcbOmaRatioLinear}{1.64}  % experiments/FIGS/data/probe_ladder_vcbwindow_mlp2/sall.json :: tasks.bc_vs_rest.arms.*.linear.log1m_auc :: sha256 5b8b5c7c28602f93 :: seed mean of 1-AUC, 17 classes over 162
\newcommand{\VcbFactorLinear}{1.64}  % experiments/FIGS/data/probe_ladder_vcbwindow_mlp2/sall.json :: tasks.bc_vs_rest.arms.*.linear.log1m_auc :: sha256 5b8b5c7c28602f93 :: exp of the difference of the seed means of log(1-AUC), 17 classes over 162
\newcommand{\VcbOmaMlpOnesixtwo}{\ensuremath{(3.102 \pm 0.023)\times10^{-2}}}  % experiments/FIGS/data/probe_ladder_vcbwindow_mlp2/sall.json :: tasks.bc_vs_rest.arms.l162-s*.mlp.log1m_auc :: sha256 5b8b5c7c28602f93 :: 1-AUC = exp of it, mean +- SD over 5 seeds
\newcommand{\VcbLogOmaMlpOnesixtwo}{\ensuremath{-}3.473}  % experiments/FIGS/data/probe_ladder_vcbwindow_mlp2/sall.json :: tasks.bc_vs_rest.arms.l162-s*.mlp.log1m_auc :: sha256 5b8b5c7c28602f93 :: mean over seeds
\newcommand{\VcbRejMlpSixtyOnesixtwo}{230\,\ensuremath{\pm}\,31}  % experiments/FIGS/data/probe_ladder_vcbwindow_mlp2/sall.json :: tasks.bc_vs_rest.arms.l162-s*.mlp.rejection_at['0.60'].rejection :: sha256 5b8b5c7c28602f93 :: background rejection at this signal efficiency, mean +- SD over seeds
\newcommand{\VcbBkgLeftMlpSixtyOnesixtwo}{84.8}  % experiments/FIGS/data/probe_ladder_vcbwindow_mlp2/sall.json :: tasks.bc_vs_rest.arms.l162-s*.mlp.rejection_at['0.60'].n_bkg_pass :: sha256 5b8b5c7c28602f93 :: background jets passing the cut, mean over seeds
\newcommand{\VcbOmaMlpOneseven}{\ensuremath{(4.66 \pm 0.19)\times10^{-2}}}  % experiments/FIGS/data/probe_ladder_vcbwindow_mlp2/sall.json :: tasks.bc_vs_rest.arms.r16q1-s*.mlp.log1m_auc :: sha256 5b8b5c7c28602f93 :: 1-AUC = exp of it, mean +- SD over 5 seeds
\newcommand{\VcbLogOmaMlpOneseven}{\ensuremath{-}3.066}  % experiments/FIGS/data/probe_ladder_vcbwindow_mlp2/sall.json :: tasks.bc_vs_rest.arms.r16q1-s*.mlp.log1m_auc :: sha256 5b8b5c7c28602f93 :: mean over seeds
\newcommand{\VcbRejMlpSixtyOneseven}{78\,\ensuremath{\pm}\,13}  % experiments/FIGS/data/probe_ladder_vcbwindow_mlp2/sall.json :: tasks.bc_vs_rest.arms.r16q1-s*.mlp.rejection_at['0.60'].rejection :: sha256 5b8b5c7c28602f93 :: background rejection at this signal efficiency, mean +- SD over seeds
\newcommand{\VcbBkgLeftMlpSixtyOneseven}{250.8}  % experiments/FIGS/data/probe_ladder_vcbwindow_mlp2/sall.json :: tasks.bc_vs_rest.arms.r16q1-s*.mlp.rejection_at['0.60'].n_bkg_pass :: sha256 5b8b5c7c28602f93 :: background jets passing the cut, mean over seeds
\newcommand{\VcbOmaRatioMlp}{1.50}  % experiments/FIGS/data/probe_ladder_vcbwindow_mlp2/sall.json :: tasks.bc_vs_rest.arms.*.mlp.log1m_auc :: sha256 5b8b5c7c28602f93 :: seed mean of 1-AUC, 17 classes over 162
\newcommand{\VcbFactorMlp}{1.50}  % experiments/FIGS/data/probe_ladder_vcbwindow_mlp2/sall.json :: tasks.bc_vs_rest.arms.*.mlp.log1m_auc :: sha256 5b8b5c7c28602f93 :: exp of the difference of the seed means of log(1-AUC), 17 classes over 162
\newcommand{\AccLegOneOnesixtwoSonebEFour}{0.32273}  % experiments/FIGS/data/leg1_metrics.json :: summary.l162-s1b.N10000.accuracy_mean :: sha256 3f564e960b2c86a7 :: wave 1, superseded; 3 fine-tuning seeds
\newcommand{\AccSdLegOneOnesixtwoSonebEFour}{0.01894}  % experiments/FIGS/data/leg1_metrics.json :: summary.l162-s1b.N10000.accuracy_sd :: sha256 3f564e960b2c86a7 :: wave 1, superseded
\newcommand{\AccLegOneOnesixtwoSonebEFive}{0.53506}  % experiments/FIGS/data/leg1_metrics.json :: summary.l162-s1b.N100000.accuracy_mean :: sha256 3f564e960b2c86a7 :: wave 1, superseded; 3 fine-tuning seeds
\newcommand{\AccSdLegOneOnesixtwoSonebEFive}{0.00044}  % experiments/FIGS/data/leg1_metrics.json :: summary.l162-s1b.N100000.accuracy_sd :: sha256 3f564e960b2c86a7 :: wave 1, superseded
\newcommand{\AccLegOneOnesixtwoSonebESix}{0.54337}  % experiments/FIGS/data/leg1_metrics.json :: summary.l162-s1b.N1000000.accuracy_mean :: sha256 3f564e960b2c86a7 :: wave 1, superseded; 3 fine-tuning seeds
\newcommand{\AccSdLegOneOnesixtwoSonebESix}{0.00074}  % experiments/FIGS/data/leg1_metrics.json :: summary.l162-s1b.N1000000.accuracy_sd :: sha256 3f564e960b2c86a7 :: wave 1, superseded
\newcommand{\AccLegOneOnesevenStwoEFour}{0.22684}  % experiments/FIGS/data/leg1_metrics.json :: summary.r16q1-s2.N10000.accuracy_mean :: sha256 3f564e960b2c86a7 :: wave 1, superseded; 3 fine-tuning seeds
\newcommand{\AccSdLegOneOnesevenStwoEFour}{0.00805}  % experiments/FIGS/data/leg1_metrics.json :: summary.r16q1-s2.N10000.accuracy_sd :: sha256 3f564e960b2c86a7 :: wave 1, superseded
\newcommand{\AccLegOneOnesevenStwoEFive}{0.49650}  % experiments/FIGS/data/leg1_metrics.json :: summary.r16q1-s2.N100000.accuracy_mean :: sha256 3f564e960b2c86a7 :: wave 1, superseded; 3 fine-tuning seeds
\newcommand{\AccSdLegOneOnesevenStwoEFive}{0.00110}  % experiments/FIGS/data/leg1_metrics.json :: summary.r16q1-s2.N100000.accuracy_sd :: sha256 3f564e960b2c86a7 :: wave 1, superseded
\newcommand{\AccLegOneOnesevenStwoESix}{0.51969}  % experiments/FIGS/data/leg1_metrics.json :: summary.r16q1-s2.N1000000.accuracy_mean :: sha256 3f564e960b2c86a7 :: wave 1, superseded; 3 fine-tuning seeds
\newcommand{\AccSdLegOneOnesevenStwoESix}{0.00072}  % experiments/FIGS/data/leg1_metrics.json :: summary.r16q1-s2.N1000000.accuracy_sd :: sha256 3f564e960b2c86a7 :: wave 1, superseded
\newcommand{\AccLegOneOnesevenSthreeEFour}{0.24203}  % experiments/FIGS/data/leg1_metrics.json :: summary.r16q1-s3.N10000.accuracy_mean :: sha256 3f564e960b2c86a7 :: wave 1, superseded; 3 fine-tuning seeds
\newcommand{\AccSdLegOneOnesevenSthreeEFour}{0.00213}  % experiments/FIGS/data/leg1_metrics.json :: summary.r16q1-s3.N10000.accuracy_sd :: sha256 3f564e960b2c86a7 :: wave 1, superseded
\newcommand{\AccLegOneOnesevenSthreeEFive}{0.49313}  % experiments/FIGS/data/leg1_metrics.json :: summary.r16q1-s3.N100000.accuracy_mean :: sha256 3f564e960b2c86a7 :: wave 1, superseded; 3 fine-tuning seeds
\newcommand{\AccSdLegOneOnesevenSthreeEFive}{0.00294}  % experiments/FIGS/data/leg1_metrics.json :: summary.r16q1-s3.N100000.accuracy_sd :: sha256 3f564e960b2c86a7 :: wave 1, superseded
\newcommand{\AccLegOneOnesevenSthreeESix}{0.51804}  % experiments/FIGS/data/leg1_metrics.json :: summary.r16q1-s3.N1000000.accuracy_mean :: sha256 3f564e960b2c86a7 :: wave 1, superseded; 3 fine-tuning seeds
\newcommand{\AccSdLegOneOnesevenSthreeESix}{0.00049}  % experiments/FIGS/data/leg1_metrics.json :: summary.r16q1-s3.N1000000.accuracy_sd :: sha256 3f564e960b2c86a7 :: wave 1, superseded
\newcommand{\AccLegOneOnesevenSfourEFour}{0.24487}  % experiments/FIGS/data/leg1_metrics.json :: summary.r16q1-s4.N10000.accuracy_mean :: sha256 3f564e960b2c86a7 :: wave 1, superseded; 3 fine-tuning seeds
\newcommand{\AccSdLegOneOnesevenSfourEFour}{0.00645}  % experiments/FIGS/data/leg1_metrics.json :: summary.r16q1-s4.N10000.accuracy_sd :: sha256 3f564e960b2c86a7 :: wave 1, superseded
\newcommand{\AccLegOneOnesevenSfourEFive}{0.49510}  % experiments/FIGS/data/leg1_metrics.json :: summary.r16q1-s4.N100000.accuracy_mean :: sha256 3f564e960b2c86a7 :: wave 1, superseded; 3 fine-tuning seeds
\newcommand{\AccSdLegOneOnesevenSfourEFive}{0.00133}  % experiments/FIGS/data/leg1_metrics.json :: summary.r16q1-s4.N100000.accuracy_sd :: sha256 3f564e960b2c86a7 :: wave 1, superseded
\newcommand{\AccLegOneOnesevenSfourESix}{0.51872}  % experiments/FIGS/data/leg1_metrics.json :: summary.r16q1-s4.N1000000.accuracy_mean :: sha256 3f564e960b2c86a7 :: wave 1, superseded; 3 fine-tuning seeds
\newcommand{\AccSdLegOneOnesevenSfourESix}{0.00049}  % experiments/FIGS/data/leg1_metrics.json :: summary.r16q1-s4.N1000000.accuracy_sd :: sha256 3f564e960b2c86a7 :: wave 1, superseded
\newcommand{\AccLegOneScratchEFour}{0.31369}  % experiments/FIGS/data/leg1_metrics.json :: summary.scratch.N10000.accuracy_mean :: sha256 3f564e960b2c86a7 :: wave 1, superseded; 3 fine-tuning seeds
\newcommand{\AccSdLegOneScratchEFour}{0.00712}  % experiments/FIGS/data/leg1_metrics.json :: summary.scratch.N10000.accuracy_sd :: sha256 3f564e960b2c86a7 :: wave 1, superseded
\newcommand{\AccLegOneScratchEFive}{0.42015}  % experiments/FIGS/data/leg1_metrics.json :: summary.scratch.N100000.accuracy_mean :: sha256 3f564e960b2c86a7 :: wave 1, superseded; 3 fine-tuning seeds
\newcommand{\AccSdLegOneScratchEFive}{0.00267}  % experiments/FIGS/data/leg1_metrics.json :: summary.scratch.N100000.accuracy_sd :: sha256 3f564e960b2c86a7 :: wave 1, superseded
\newcommand{\AccLegOneScratchESix}{0.48621}  % experiments/FIGS/data/leg1_metrics.json :: summary.scratch.N1000000.accuracy_mean :: sha256 3f564e960b2c86a7 :: wave 1, superseded; 3 fine-tuning seeds
\newcommand{\AccSdLegOneScratchESix}{0.00100}  % experiments/FIGS/data/leg1_metrics.json :: summary.scratch.N1000000.accuracy_sd :: sha256 3f564e960b2c86a7 :: wave 1, superseded
\newcommand{\AccLegOneSophonPublicEFour}{0.21613}  % experiments/FIGS/data/leg1_metrics.json :: summary.sophon-public.N10000.accuracy_mean :: sha256 3f564e960b2c86a7 :: wave 1, superseded; 3 fine-tuning seeds
\newcommand{\AccSdLegOneSophonPublicEFour}{0.00734}  % experiments/FIGS/data/leg1_metrics.json :: summary.sophon-public.N10000.accuracy_sd :: sha256 3f564e960b2c86a7 :: wave 1, superseded
\newcommand{\AccLegOneSophonPublicEFive}{0.53187}  % experiments/FIGS/data/leg1_metrics.json :: summary.sophon-public.N100000.accuracy_mean :: sha256 3f564e960b2c86a7 :: wave 1, superseded; 3 fine-tuning seeds
\newcommand{\AccSdLegOneSophonPublicEFive}{0.00231}  % experiments/FIGS/data/leg1_metrics.json :: summary.sophon-public.N100000.accuracy_sd :: sha256 3f564e960b2c86a7 :: wave 1, superseded
\newcommand{\AccLegOneSophonPublicESix}{0.54273}  % experiments/FIGS/data/leg1_metrics.json :: summary.sophon-public.N1000000.accuracy_mean :: sha256 3f564e960b2c86a7 :: wave 1, superseded; 3 fine-tuning seeds
\newcommand{\AccSdLegOneSophonPublicESix}{0.00064}  % experiments/FIGS/data/leg1_metrics.json :: summary.sophon-public.N1000000.accuracy_sd :: sha256 3f564e960b2c86a7 :: wave 1, superseded
\newcommand{\AccLegTwoOnesixtwoSonebEFour}{0.78606}  % experiments/FIGS/data/leg2_metrics.json :: summary.l162-s1b.N10000.accuracy_mean :: sha256 701c58c3240274bb :: wave 1, superseded; 3 fine-tuning seeds
\newcommand{\AccSdLegTwoOnesixtwoSonebEFour}{0.00111}  % experiments/FIGS/data/leg2_metrics.json :: summary.l162-s1b.N10000.accuracy_sd :: sha256 701c58c3240274bb :: wave 1, superseded
\newcommand{\AccLegTwoOnesixtwoSonebEFive}{0.83598}  % experiments/FIGS/data/leg2_metrics.json :: summary.l162-s1b.N100000.accuracy_mean :: sha256 701c58c3240274bb :: wave 1, superseded; 3 fine-tuning seeds
\newcommand{\AccSdLegTwoOnesixtwoSonebEFive}{0.00037}  % experiments/FIGS/data/leg2_metrics.json :: summary.l162-s1b.N100000.accuracy_sd :: sha256 701c58c3240274bb :: wave 1, superseded
\newcommand{\AccLegTwoOnesixtwoSonebESix}{0.84704}  % experiments/FIGS/data/leg2_metrics.json :: summary.l162-s1b.N1000000.accuracy_mean :: sha256 701c58c3240274bb :: wave 1, superseded; 3 fine-tuning seeds
\newcommand{\AccSdLegTwoOnesixtwoSonebESix}{0.00027}  % experiments/FIGS/data/leg2_metrics.json :: summary.l162-s1b.N1000000.accuracy_sd :: sha256 701c58c3240274bb :: wave 1, superseded
\newcommand{\AccLegTwoOnesevenStwoEFour}{0.72156}  % experiments/FIGS/data/leg2_metrics.json :: summary.r16q1-s2.N10000.accuracy_mean :: sha256 701c58c3240274bb :: wave 1, superseded; 3 fine-tuning seeds
\newcommand{\AccSdLegTwoOnesevenStwoEFour}{0.01049}  % experiments/FIGS/data/leg2_metrics.json :: summary.r16q1-s2.N10000.accuracy_sd :: sha256 701c58c3240274bb :: wave 1, superseded
\newcommand{\AccLegTwoOnesevenStwoEFive}{0.82614}  % experiments/FIGS/data/leg2_metrics.json :: summary.r16q1-s2.N100000.accuracy_mean :: sha256 701c58c3240274bb :: wave 1, superseded; 3 fine-tuning seeds
\newcommand{\AccSdLegTwoOnesevenStwoEFive}{0.00058}  % experiments/FIGS/data/leg2_metrics.json :: summary.r16q1-s2.N100000.accuracy_sd :: sha256 701c58c3240274bb :: wave 1, superseded
\newcommand{\AccLegTwoOnesevenStwoESix}{0.83943}  % experiments/FIGS/data/leg2_metrics.json :: summary.r16q1-s2.N1000000.accuracy_mean :: sha256 701c58c3240274bb :: wave 1, superseded; 3 fine-tuning seeds
\newcommand{\AccSdLegTwoOnesevenStwoESix}{0.00052}  % experiments/FIGS/data/leg2_metrics.json :: summary.r16q1-s2.N1000000.accuracy_sd :: sha256 701c58c3240274bb :: wave 1, superseded
\newcommand{\AccLegTwoOnesevenSthreeEFour}{0.72103}  % experiments/FIGS/data/leg2_metrics.json :: summary.r16q1-s3.N10000.accuracy_mean :: sha256 701c58c3240274bb :: wave 1, superseded; 3 fine-tuning seeds
\newcommand{\AccSdLegTwoOnesevenSthreeEFour}{0.01651}  % experiments/FIGS/data/leg2_metrics.json :: summary.r16q1-s3.N10000.accuracy_sd :: sha256 701c58c3240274bb :: wave 1, superseded
\newcommand{\AccLegTwoOnesevenSthreeEFive}{0.82629}  % experiments/FIGS/data/leg2_metrics.json :: summary.r16q1-s3.N100000.accuracy_mean :: sha256 701c58c3240274bb :: wave 1, superseded; 3 fine-tuning seeds
\newcommand{\AccSdLegTwoOnesevenSthreeEFive}{0.00035}  % experiments/FIGS/data/leg2_metrics.json :: summary.r16q1-s3.N100000.accuracy_sd :: sha256 701c58c3240274bb :: wave 1, superseded
\newcommand{\AccLegTwoOnesevenSthreeESix}{0.83944}  % experiments/FIGS/data/leg2_metrics.json :: summary.r16q1-s3.N1000000.accuracy_mean :: sha256 701c58c3240274bb :: wave 1, superseded; 3 fine-tuning seeds
\newcommand{\AccSdLegTwoOnesevenSthreeESix}{0.00019}  % experiments/FIGS/data/leg2_metrics.json :: summary.r16q1-s3.N1000000.accuracy_sd :: sha256 701c58c3240274bb :: wave 1, superseded
\newcommand{\AccLegTwoOnesevenSfourEFour}{0.72953}  % experiments/FIGS/data/leg2_metrics.json :: summary.r16q1-s4.N10000.accuracy_mean :: sha256 701c58c3240274bb :: wave 1, superseded; 3 fine-tuning seeds
\newcommand{\AccSdLegTwoOnesevenSfourEFour}{0.01069}  % experiments/FIGS/data/leg2_metrics.json :: summary.r16q1-s4.N10000.accuracy_sd :: sha256 701c58c3240274bb :: wave 1, superseded
\newcommand{\AccLegTwoOnesevenSfourEFive}{0.82564}  % experiments/FIGS/data/leg2_metrics.json :: summary.r16q1-s4.N100000.accuracy_mean :: sha256 701c58c3240274bb :: wave 1, superseded; 3 fine-tuning seeds
\newcommand{\AccSdLegTwoOnesevenSfourEFive}{0.00041}  % experiments/FIGS/data/leg2_metrics.json :: summary.r16q1-s4.N100000.accuracy_sd :: sha256 701c58c3240274bb :: wave 1, superseded
\newcommand{\AccLegTwoOnesevenSfourESix}{0.83909}  % experiments/FIGS/data/leg2_metrics.json :: summary.r16q1-s4.N1000000.accuracy_mean :: sha256 701c58c3240274bb :: wave 1, superseded; 3 fine-tuning seeds
\newcommand{\AccSdLegTwoOnesevenSfourESix}{0.00037}  % experiments/FIGS/data/leg2_metrics.json :: summary.r16q1-s4.N1000000.accuracy_sd :: sha256 701c58c3240274bb :: wave 1, superseded
\newcommand{\AccLegTwoRefEonearmsSoneRef}{0.81770}  % experiments/FIGS/data/leg2_metrics.json :: summary.ref_e1arms-s1.ref.accuracy_mean :: sha256 701c58c3240274bb :: wave 1, superseded; 1 fine-tuning seeds
\newcommand{\AccSdLegTwoRefEonearmsSoneRef}{---}  % experiments/FIGS/data/leg2_metrics.json :: summary.ref_e1arms-s1.ref.accuracy_sd :: sha256 701c58c3240274bb :: wave 1, superseded
\newcommand{\AccLegTwoRefEonearmsStwoRef}{0.81602}  % experiments/FIGS/data/leg2_metrics.json :: summary.ref_e1arms-s2.ref.accuracy_mean :: sha256 701c58c3240274bb :: wave 1, superseded; 1 fine-tuning seeds
\newcommand{\AccSdLegTwoRefEonearmsStwoRef}{---}  % experiments/FIGS/data/leg2_metrics.json :: summary.ref_e1arms-s2.ref.accuracy_sd :: sha256 701c58c3240274bb :: wave 1, superseded
\newcommand{\AccLegTwoRefEonearmsSthreeRef}{0.81625}  % experiments/FIGS/data/leg2_metrics.json :: summary.ref_e1arms-s3.ref.accuracy_mean :: sha256 701c58c3240274bb :: wave 1, superseded; 1 fine-tuning seeds
\newcommand{\AccSdLegTwoRefEonearmsSthreeRef}{---}  % experiments/FIGS/data/leg2_metrics.json :: summary.ref_e1arms-s3.ref.accuracy_sd :: sha256 701c58c3240274bb :: wave 1, superseded
\newcommand{\AccLegTwoScratchEFour}{0.68230}  % experiments/FIGS/data/leg2_metrics.json :: summary.scratch.N10000.accuracy_mean :: sha256 701c58c3240274bb :: wave 1, superseded; 3 fine-tuning seeds
\newcommand{\AccSdLegTwoScratchEFour}{0.00282}  % experiments/FIGS/data/leg2_metrics.json :: summary.scratch.N10000.accuracy_sd :: sha256 701c58c3240274bb :: wave 1, superseded
\newcommand{\AccLegTwoScratchEFive}{0.77769}  % experiments/FIGS/data/leg2_metrics.json :: summary.scratch.N100000.accuracy_mean :: sha256 701c58c3240274bb :: wave 1, superseded; 3 fine-tuning seeds
\newcommand{\AccSdLegTwoScratchEFive}{0.00208}  % experiments/FIGS/data/leg2_metrics.json :: summary.scratch.N100000.accuracy_sd :: sha256 701c58c3240274bb :: wave 1, superseded
\newcommand{\AccLegTwoScratchESix}{0.82303}  % experiments/FIGS/data/leg2_metrics.json :: summary.scratch.N1000000.accuracy_mean :: sha256 701c58c3240274bb :: wave 1, superseded; 3 fine-tuning seeds
\newcommand{\AccSdLegTwoScratchESix}{0.00134}  % experiments/FIGS/data/leg2_metrics.json :: summary.scratch.N1000000.accuracy_sd :: sha256 701c58c3240274bb :: wave 1, superseded
\newcommand{\AccLegTwoSophonPublicEFour}{0.72739}  % experiments/FIGS/data/leg2_metrics.json :: summary.sophon-public.N10000.accuracy_mean :: sha256 701c58c3240274bb :: wave 1, superseded; 3 fine-tuning seeds
\newcommand{\AccSdLegTwoSophonPublicEFour}{0.00470}  % experiments/FIGS/data/leg2_metrics.json :: summary.sophon-public.N10000.accuracy_sd :: sha256 701c58c3240274bb :: wave 1, superseded
\newcommand{\AccLegTwoSophonPublicEFive}{0.83492}  % experiments/FIGS/data/leg2_metrics.json :: summary.sophon-public.N100000.accuracy_mean :: sha256 701c58c3240274bb :: wave 1, superseded; 3 fine-tuning seeds
\newcommand{\AccSdLegTwoSophonPublicEFive}{0.00049}  % experiments/FIGS/data/leg2_metrics.json :: summary.sophon-public.N100000.accuracy_sd :: sha256 701c58c3240274bb :: wave 1, superseded
\newcommand{\AccLegTwoSophonPublicESix}{0.84632}  % experiments/FIGS/data/leg2_metrics.json :: summary.sophon-public.N1000000.accuracy_mean :: sha256 701c58c3240274bb :: wave 1, superseded; 3 fine-tuning seeds
\newcommand{\AccSdLegTwoSophonPublicESix}{0.00031}  % experiments/FIGS/data/leg2_metrics.json :: summary.sophon-public.N1000000.accuracy_sd :: sha256 701c58c3240274bb :: wave 1, superseded
\newcommand{\FtBudgetSteps}{20{,}000}  % scripts/build_ft_budget_jobs.py :: EPOCHS * STEPS :: sha256 f2ed9faa8e3bd622 :: the largest number of optimizer steps below the full training set
\newcommand{\FtLongSteps}{80{,}000}  % scripts/build_ft_budget_jobs.py :: LONG_EPOCHS * STEPS :: sha256 f2ed9faa8e3bd622 :: the largest number of optimizer steps of the long budget at LONG_N jets (JetClass, JetClass-II)
\newcommand{\DesignNTrainSelected}{109{,}938{,}036}  % experiments/FIGS/data/v2/jc2_train/counts.json :: n_selected :: sha256 b0b5a8b63aa808c0 :: selected jets of the pretrain_train files (experiments/FT/count_jc2_train.py)
\newcommand{\FtBudgetEval}{500}  % scripts/build_ft_budget_jobs.py :: STEPS :: sha256 f2ed9faa8e3bd622 :: steps between validation passes at and above SMALL_N training jets
\newcommand{\FtBudgetEvalSmall}{50}  % scripts/build_ft_budget_jobs.py :: SMALL_STEPS :: sha256 f2ed9faa8e3bd622 :: steps between validation passes below SMALL_N training jets
\newcommand{\FtBudgetEvalSmallBelow}{\ensuremath{10^{5}}}  % scripts/build_ft_budget_jobs.py :: SMALL_N :: sha256 f2ed9faa8e3bd622 :: training jets below which the validation runs every SMALL_STEPS steps
\newcommand{\FtBudgetPatience}{10}  % scripts/build_ft_budget_jobs.py :: PATIENCE :: sha256 f2ed9faa8e3bd622 :: validation passes without improvement before training stops
\newcommand{\FtHeadMult}{50}  % scripts/build_ft_budget_jobs.py :: HEAD_MULT :: sha256 f2ed9faa8e3bd622 :: the head's learning-rate multiplier
\newcommand{\FtHeadCap}{\ensuremath{5\times10^{-3}}}  % scripts/build_ft_budget_jobs.py :: HEAD_CAP :: sha256 f2ed9faa8e3bd622 :: the head's largest learning rate
\newcommand{\FtBudgetBatch}{512}  % scripts/build_ft_budget_jobs.py :: BATCH :: sha256 f2ed9faa8e3bd622
\newcommand{\FtBudgetWeightDecay}{0.01}  % scripts/build_ft_budget_jobs.py :: --optimizer-option weight_decay (every cell) :: sha256 f2ed9faa8e3bd622
\newcommand{\FtSweepPre}{\ensuremath{10^{-4}}, \ensuremath{3\times10^{-4}} and \ensuremath{10^{-3}}}  % scripts/build_ft_budget_jobs.py :: SWEEP_PRE :: sha256 f2ed9faa8e3bd622 :: backbone rates of the sweep
\newcommand{\FtSweepScratch}{\ensuremath{3\times10^{-4}}, \ensuremath{10^{-3}} and \ensuremath{3\times10^{-3}}}  % scripts/build_ft_budget_jobs.py :: SWEEP_SCRATCH :: sha256 f2ed9faa8e3bd622 :: from-scratch rates of the sweep
\newcommand{\FtSweepN}{10{,}000}  % scripts/build_ft_budget_jobs.py :: SWEEP_N :: sha256 f2ed9faa8e3bd622 :: training jets of the sweep
\newcommand{\FtValJets}{200{,}000}  % scripts/build_ft_budget_jobs.py :: VAL_N :: sha256 f2ed9faa8e3bd622 :: the fixed validation set
\newcommand{\FtValJetsBal}{20{,}480}  % scripts/build_ft_budget_jobs.py :: VAL_SIZE[jc2b] :: sha256 f2ed9faa8e3bd622 :: the balanced JetClass-II validation set
\newcommand{\LpVfourCMax}{\ensuremath{10^{4}}}  % experiments/EVAL/probe_bal.py :: C_GRID_V4 :: sha256 a7804216e3179d9d :: largest inverse regularisation of the v4 probes
\newcommand{\JcBalTestPerClass}{5{,}000}  % experiments/FT/stage_jc2_balanced.py :: TEST_PER_CLASS :: sha256 0260341d5ed75fec
\newcommand{\JcBalValPerClass}{109}  % experiments/FT/stage_jc2_balanced.py :: VAL_PER_CLASS :: sha256 0260341d5ed75fec
\newcommand{\JcBalPairTrainPerSide}{\ensuremath{10^{5}}}  % experiments/FT/stage_jc2_balanced.py :: PAIR_TRAIN_PER_SIDE :: sha256 0260341d5ed75fec
\newcommand{\JcBalPairTestPerSide}{\ensuremath{2\times10^{5}}}  % experiments/FT/stage_jc2_balanced.py :: PAIR_TEST_PER_SIDE :: sha256 0260341d5ed75fec
\newcommand{\JcBalNShort}{11}  % experiments/FIGS/data/v2/jc2bal/census.json :: counts.ft_test :: sha256 a4379ee0817254cf :: classes with fewer than 5000 selected test jets
\newcommand{\JcBalRarestTest}{224}  % experiments/FIGS/data/v2/jc2bal/census.json :: counts.ft_test (minimum) :: sha256 a4379ee0817254cf :: label_QCD_bcs
\newcommand{\JcBalShortTrainESix}{11}  % experiments/FIGS/data/v2/jc2bal/census.json :: counts.ft_heldout :: sha256 a4379ee0817254cf :: classes with fewer held-out jets than the 1e6 quota
\newcommand{\JcBalNTest}{903{,}922}  % experiments/FIGS/data/v2/jc2bal/manifest.json :: outputs[test.parquet].rows :: sha256 9c7775be6ff56ac3 :: jets of the balanced test set
\newcommand{\FtLrPre}{\ensuremath{10^{-3}}}  % experiments/FIGS/data/v2/budget/lrsweep_metrics.json :: argmax of the summed validation metric, pretrained :: sha256 7cdc148054ca902c
\newcommand{\FtLrScratchTop}{\ensuremath{10^{-3}}}  % experiments/FIGS/data/v2/budget/lrsweep_metrics.json :: argmax of the validation metric on top, from scratch :: sha256 7cdc148054ca902c :: agrees with LR of the launched job-ft-budget-bench, -jc1, -jc2b, -scratch-full and -long-scratch-full from-scratch cells
\newcommand{\FtLrScratchQg}{\ensuremath{3\times10^{-3}}}  % experiments/FIGS/data/v2/budget/lrsweep_metrics.json :: argmax of the validation metric on quark/gluon, from scratch :: sha256 7cdc148054ca902c :: agrees with LR of the launched job-ft-budget-bench and -scratch-full from-scratch cells
\newcommand{\MlstNParamsFt}{2{,}209{,}270}  % experiments/FIGS/mlst_v4.py :: our_params(2, 17) :: sha256 fa536b4866ffd651 :: trainable parameters of our fine-tuned network, two outputs (top, quark/gluon)
\newcommand{\MlstNParamsFtKin}{2{,}207{,}950}  % experiments/FIGS/mlst_v4.py :: our_params(2, 7) :: sha256 fa536b4866ffd651 :: trainable parameters of our fine-tuned network, two outputs, the 7 kinematic inputs
\newcommand{\MlstNParamsJc}{2{,}213{,}374}  % experiments/FIGS/mlst_v4.py :: our_params(10, 17) :: sha256 fa536b4866ffd651 :: trainable parameters of our fine-tuned network, the ten JetClass classes
\newcommand{\MlstProbeRunsDone}{27}  % experiments/FIGS/data/v2/index.json :: runs with a linear-probe fit (count) :: sha256 c6c4138ee5c86423 :: of the paper's runs
\newcommand{\MlstPaperRuns}{27}  % configs/arms/v2_paper_runs.txt :: runs (count) :: sha256 b04562b14ea8cbc7
\newcommand{\MlstSdRelUnc}{50\%}  % configs/arms/v2_paper_runs.txt :: relative standard error of a standard deviation from the runs of a hierarchy level, 1/sqrt(2(n-1)) :: sha256 b04562b14ea8cbc7
\newcommand{\MlstCmpJcTwoResOverFull}{1.026\,\ensuremath{\pm}\,0.005}  % experiments/FIGS/data/v2/index.json :: linear probe jc2_188 N=1000000: (1 - macro_auc_ovr) of L162 over L188, run by run :: sha256 c6c4138ee5c86423
\newcommand{\MlstCmpJcTwoResOverFullPr}{2.6\% more errors}  % experiments/FIGS/data/v2/index.json :: linear probe jc2_188 N=1000000: (1 - macro_auc_ovr) of L162 over L188, run by run :: sha256 c6c4138ee5c86423 :: prose twin: two significant figures, direction in words
\newcommand{\MlstCmpJcTwoMidOverFull}{1.227\,\ensuremath{\pm}\,0.009}  % experiments/FIGS/data/v2/index.json :: linear probe jc2_188 N=1000000: (1 - macro_auc_ovr) of R42_Q1 over L188, run by run :: sha256 c6c4138ee5c86423
\newcommand{\MlstCmpJcTwoMidOverFullPr}{23\% more errors}  % experiments/FIGS/data/v2/index.json :: linear probe jc2_188 N=1000000: (1 - macro_auc_ovr) of R42_Q1 over L188, run by run :: sha256 c6c4138ee5c86423 :: prose twin: two significant figures, direction in words
\newcommand{\MlstCmpJcTwoCoarseOverFull}{1.557\,\ensuremath{\pm}\,0.019}  % experiments/FIGS/data/v2/index.json :: linear probe jc2_188 N=1000000: (1 - macro_auc_ovr) of R16_Q1 over L188, run by run :: sha256 c6c4138ee5c86423
\newcommand{\MlstCmpJcTwoCoarseOverFullPr}{56\% more errors}  % experiments/FIGS/data/v2/index.json :: linear probe jc2_188 N=1000000: (1 - macro_auc_ovr) of R16_Q1 over L188, run by run :: sha256 c6c4138ee5c86423 :: prose twin: two significant figures, direction in words
\newcommand{\MlstCmpJcTwoResMassOverRes}{1.022 and 1.016}  % experiments/FIGS/data/v2/index.json :: linear probe jc2_188 N=1000000: (1 - macro_auc_ovr) of L162_MASS over L162, run by run :: sha256 c6c4138ee5c86423
\newcommand{\MlstCmpJcTwoResMassOverResPr}{equal within 3\%}  % experiments/FIGS/data/v2/index.json :: linear probe jc2_188 N=1000000: (1 - macro_auc_ovr) of L162_MASS over L162, run by run :: sha256 c6c4138ee5c86423 :: prose twin: two significant figures, direction in words
\newcommand{\MlstCmpJcTwoResMassOverResPrMax}{equal within 3\%}  % experiments/FIGS/data/v2/index.json :: linear probe jc2_188 N=1000000: (1 - macro_auc_ovr) of L162_MASS over L162, run by run :: sha256 c6c4138ee5c86423 :: prose twin of the run farther from 1
\newcommand{\MlstCmpJcTwoCoarseMassOverCoarse}{1.054 and 0.977}  % experiments/FIGS/data/v2/index.json :: linear probe jc2_188 N=1000000: (1 - macro_auc_ovr) of R16_Q1_MASS over R16_Q1, run by run :: sha256 c6c4138ee5c86423
\newcommand{\MlstCmpJcTwoCoarseMassOverCoarsePr}{equal within 6\%}  % experiments/FIGS/data/v2/index.json :: linear probe jc2_188 N=1000000: (1 - macro_auc_ovr) of R16_Q1_MASS over R16_Q1, run by run :: sha256 c6c4138ee5c86423 :: prose twin: two significant figures, direction in words
\newcommand{\MlstCmpJcTwoCoarseMassOverCoarsePrMax}{equal within 6\%}  % experiments/FIGS/data/v2/index.json :: linear probe jc2_188 N=1000000: (1 - macro_auc_ovr) of R16_Q1_MASS over R16_Q1, run by run :: sha256 c6c4138ee5c86423 :: prose twin of the run farther from 1
\newcommand{\MlstCmpJcTwoResMassOverFull}{1.053 and 1.037}  % experiments/FIGS/data/v2/index.json :: linear probe jc2_188 N=1000000: (1 - macro_auc_ovr) of L162_MASS over L188, run by run :: sha256 c6c4138ee5c86423
\newcommand{\MlstCmpJcTwoResMassOverFullPr}{at least 3.7\% more errors}  % experiments/FIGS/data/v2/index.json :: linear probe jc2_188 N=1000000: (1 - macro_auc_ovr) of L162_MASS over L188, run by run :: sha256 c6c4138ee5c86423 :: prose twin: two significant figures, direction in words
\newcommand{\MlstCmpJcTwoResMassOverFullPrMax}{up to 5.3\% more errors}  % experiments/FIGS/data/v2/index.json :: linear probe jc2_188 N=1000000: (1 - macro_auc_ovr) of L162_MASS over L188, run by run :: sha256 c6c4138ee5c86423 :: prose twin of the run farther from 1
\newcommand{\MlstCmpJcTwoCoarseMassOverFull}{1.617 and 1.528}  % experiments/FIGS/data/v2/index.json :: linear probe jc2_188 N=1000000: (1 - macro_auc_ovr) of R16_Q1_MASS over L188, run by run :: sha256 c6c4138ee5c86423
\newcommand{\MlstCmpJcTwoCoarseMassOverFullPr}{at least 53\% more errors}  % experiments/FIGS/data/v2/index.json :: linear probe jc2_188 N=1000000: (1 - macro_auc_ovr) of R16_Q1_MASS over L188, run by run :: sha256 c6c4138ee5c86423 :: prose twin: two significant figures, direction in words
\newcommand{\MlstCmpJcTwoCoarseMassOverFullPrMax}{up to 62\% more errors}  % experiments/FIGS/data/v2/index.json :: linear probe jc2_188 N=1000000: (1 - macro_auc_ovr) of R16_Q1_MASS over L188, run by run :: sha256 c6c4138ee5c86423 :: prose twin of the run farther from 1
\newcommand{\MlstCmpJcTwoFzeroOverFull}{1.424 and 1.462}  % experiments/FIGS/data/v2/index.json :: linear probe jc2_188 N=1000000: (1 - macro_auc_ovr) of FLAV_F0 over L188, run by run :: sha256 c6c4138ee5c86423
\newcommand{\MlstCmpJcTwoFzeroOverFullPr}{at least 42\% more errors}  % experiments/FIGS/data/v2/index.json :: linear probe jc2_188 N=1000000: (1 - macro_auc_ovr) of FLAV_F0 over L188, run by run :: sha256 c6c4138ee5c86423 :: prose twin: two significant figures, direction in words
\newcommand{\MlstCmpJcTwoFzeroOverFullPrMax}{up to 46\% more errors}  % experiments/FIGS/data/v2/index.json :: linear probe jc2_188 N=1000000: (1 - macro_auc_ovr) of FLAV_F0 over L188, run by run :: sha256 c6c4138ee5c86423 :: prose twin of the run farther from 1
\newcommand{\MlstCmpJcTwoFoneOverFull}{1.218 and 1.212}  % experiments/FIGS/data/v2/index.json :: linear probe jc2_188 N=1000000: (1 - macro_auc_ovr) of FLAV_F1 over L188, run by run :: sha256 c6c4138ee5c86423
\newcommand{\MlstCmpJcTwoFoneOverFullPr}{at least 21\% more errors}  % experiments/FIGS/data/v2/index.json :: linear probe jc2_188 N=1000000: (1 - macro_auc_ovr) of FLAV_F1 over L188, run by run :: sha256 c6c4138ee5c86423 :: prose twin: two significant figures, direction in words
\newcommand{\MlstCmpJcTwoFoneOverFullPrMax}{up to 22\% more errors}  % experiments/FIGS/data/v2/index.json :: linear probe jc2_188 N=1000000: (1 - macro_auc_ovr) of FLAV_F1 over L188, run by run :: sha256 c6c4138ee5c86423 :: prose twin of the run farther from 1
\newcommand{\MlstCmpJcTwoFonerOverFull}{1.206 and 1.208}  % experiments/FIGS/data/v2/index.json :: linear probe jc2_188 N=1000000: (1 - macro_auc_ovr) of FLAV_F1R over L188, run by run :: sha256 c6c4138ee5c86423
\newcommand{\MlstCmpJcTwoFonerOverFullPr}{at least 21\% more errors}  % experiments/FIGS/data/v2/index.json :: linear probe jc2_188 N=1000000: (1 - macro_auc_ovr) of FLAV_F1R over L188, run by run :: sha256 c6c4138ee5c86423 :: prose twin: two significant figures, direction in words
\newcommand{\MlstCmpJcTwoFonerOverFullPrMax}{up to 21\% more errors}  % experiments/FIGS/data/v2/index.json :: linear probe jc2_188 N=1000000: (1 - macro_auc_ovr) of FLAV_F1R over L188, run by run :: sha256 c6c4138ee5c86423 :: prose twin of the run farther from 1
\newcommand{\MlstCmpJcTwoFzeroOverCoarse}{0.9275 and 0.9345}  % experiments/FIGS/data/v2/index.json :: linear probe jc2_188 N=1000000: (1 - macro_auc_ovr) of FLAV_F0 over R16_Q1, run by run :: sha256 c6c4138ee5c86423
\newcommand{\MlstCmpJcTwoFzeroOverCoarsePr}{at least 6.5\% fewer errors}  % experiments/FIGS/data/v2/index.json :: linear probe jc2_188 N=1000000: (1 - macro_auc_ovr) of FLAV_F0 over R16_Q1, run by run :: sha256 c6c4138ee5c86423 :: prose twin: two significant figures, direction in words
\newcommand{\MlstCmpJcTwoFzeroOverCoarsePrMax}{up to 7.2\% fewer errors}  % experiments/FIGS/data/v2/index.json :: linear probe jc2_188 N=1000000: (1 - macro_auc_ovr) of FLAV_F0 over R16_Q1, run by run :: sha256 c6c4138ee5c86423 :: prose twin of the run farther from 1
\newcommand{\MlstCmpJcTwoFoneOverCoarse}{0.794 and 0.775}  % experiments/FIGS/data/v2/index.json :: linear probe jc2_188 N=1000000: (1 - macro_auc_ovr) of FLAV_F1 over R16_Q1, run by run :: sha256 c6c4138ee5c86423
\newcommand{\MlstCmpJcTwoFoneOverCoarsePr}{at least 21\% fewer errors}  % experiments/FIGS/data/v2/index.json :: linear probe jc2_188 N=1000000: (1 - macro_auc_ovr) of FLAV_F1 over R16_Q1, run by run :: sha256 c6c4138ee5c86423 :: prose twin: two significant figures, direction in words
\newcommand{\MlstCmpJcTwoFoneOverCoarsePrMax}{up to 23\% fewer errors}  % experiments/FIGS/data/v2/index.json :: linear probe jc2_188 N=1000000: (1 - macro_auc_ovr) of FLAV_F1 over R16_Q1, run by run :: sha256 c6c4138ee5c86423 :: prose twin of the run farther from 1
\newcommand{\MlstCmpJcTwoFonerOverCoarse}{0.786 and 0.772}  % experiments/FIGS/data/v2/index.json :: linear probe jc2_188 N=1000000: (1 - macro_auc_ovr) of FLAV_F1R over R16_Q1, run by run :: sha256 c6c4138ee5c86423
\newcommand{\MlstCmpJcTwoFonerOverCoarsePr}{at least 21\% fewer errors}  % experiments/FIGS/data/v2/index.json :: linear probe jc2_188 N=1000000: (1 - macro_auc_ovr) of FLAV_F1R over R16_Q1, run by run :: sha256 c6c4138ee5c86423 :: prose twin: two significant figures, direction in words
\newcommand{\MlstCmpJcTwoFonerOverCoarsePrMax}{up to 23\% fewer errors}  % experiments/FIGS/data/v2/index.json :: linear probe jc2_188 N=1000000: (1 - macro_auc_ovr) of FLAV_F1R over R16_Q1, run by run :: sha256 c6c4138ee5c86423 :: prose twin of the run farther from 1
\newcommand{\MlstCmpJcTwoFoneOverFzero}{0.856 and 0.829}  % experiments/FIGS/data/v2/index.json :: linear probe jc2_188 N=1000000: (1 - macro_auc_ovr) of FLAV_F1 over FLAV_F0, run by run :: sha256 c6c4138ee5c86423
\newcommand{\MlstCmpJcTwoFoneOverFzeroPr}{at least 14\% fewer errors}  % experiments/FIGS/data/v2/index.json :: linear probe jc2_188 N=1000000: (1 - macro_auc_ovr) of FLAV_F1 over FLAV_F0, run by run :: sha256 c6c4138ee5c86423 :: prose twin: two significant figures, direction in words
\newcommand{\MlstCmpJcTwoFoneOverFzeroPrMax}{up to 17\% fewer errors}  % experiments/FIGS/data/v2/index.json :: linear probe jc2_188 N=1000000: (1 - macro_auc_ovr) of FLAV_F1 over FLAV_F0, run by run :: sha256 c6c4138ee5c86423 :: prose twin of the run farther from 1
\newcommand{\MlstCmpJcTwoFonerOverFzero}{0.847 and 0.826}  % experiments/FIGS/data/v2/index.json :: linear probe jc2_188 N=1000000: (1 - macro_auc_ovr) of FLAV_F1R over FLAV_F0, run by run :: sha256 c6c4138ee5c86423
\newcommand{\MlstCmpJcTwoFonerOverFzeroPr}{at least 15\% fewer errors}  % experiments/FIGS/data/v2/index.json :: linear probe jc2_188 N=1000000: (1 - macro_auc_ovr) of FLAV_F1R over FLAV_F0, run by run :: sha256 c6c4138ee5c86423 :: prose twin: two significant figures, direction in words
\newcommand{\MlstCmpJcTwoFonerOverFzeroPrMax}{up to 17\% fewer errors}  % experiments/FIGS/data/v2/index.json :: linear probe jc2_188 N=1000000: (1 - macro_auc_ovr) of FLAV_F1R over FLAV_F0, run by run :: sha256 c6c4138ee5c86423 :: prose twin of the run farther from 1
\newcommand{\MlstCmpJcTwoFonerOverFone}{0.9899 and 0.9966}  % experiments/FIGS/data/v2/index.json :: linear probe jc2_188 N=1000000: (1 - macro_auc_ovr) of FLAV_F1R over FLAV_F1, run by run :: sha256 c6c4138ee5c86423
\newcommand{\MlstCmpJcTwoFonerOverFonePr}{equal within 2\%}  % experiments/FIGS/data/v2/index.json :: linear probe jc2_188 N=1000000: (1 - macro_auc_ovr) of FLAV_F1R over FLAV_F1, run by run :: sha256 c6c4138ee5c86423 :: prose twin: two significant figures, direction in words
\newcommand{\MlstCmpJcTwoFonerOverFonePrMax}{equal within 2\%}  % experiments/FIGS/data/v2/index.json :: linear probe jc2_188 N=1000000: (1 - macro_auc_ovr) of FLAV_F1R over FLAV_F1, run by run :: sha256 c6c4138ee5c86423 :: prose twin of the run farther from 1
\newcommand{\MlstCmpJcTwoMidOverRes}{1.197\,\ensuremath{\pm}\,0.013}  % experiments/FIGS/data/v2/index.json :: linear probe jc2_188 N=1000000: (1 - macro_auc_ovr) of R42_Q1 over L162, run by run :: sha256 c6c4138ee5c86423
\newcommand{\MlstCmpJcTwoMidOverResPr}{20\% more errors}  % experiments/FIGS/data/v2/index.json :: linear probe jc2_188 N=1000000: (1 - macro_auc_ovr) of R42_Q1 over L162, run by run :: sha256 c6c4138ee5c86423 :: prose twin: two significant figures, direction in words
\newcommand{\MlstCmpJcTwoCoarseOverRes}{1.518\,\ensuremath{\pm}\,0.025}  % experiments/FIGS/data/v2/index.json :: linear probe jc2_188 N=1000000: (1 - macro_auc_ovr) of R16_Q1 over L162, run by run :: sha256 c6c4138ee5c86423
\newcommand{\MlstCmpJcTwoCoarseOverResPr}{52\% more errors}  % experiments/FIGS/data/v2/index.json :: linear probe jc2_188 N=1000000: (1 - macro_auc_ovr) of R16_Q1 over L162, run by run :: sha256 c6c4138ee5c86423 :: prose twin: two significant figures, direction in words
\newcommand{\MlstCmpJcTwoCoarseOverMid}{1.269\,\ensuremath{\pm}\,0.007}  % experiments/FIGS/data/v2/index.json :: linear probe jc2_188 N=1000000: (1 - macro_auc_ovr) of R16_Q1 over R42_Q1, run by run :: sha256 c6c4138ee5c86423
\newcommand{\MlstCmpJcTwoCoarseOverMidPr}{27\% more errors}  % experiments/FIGS/data/v2/index.json :: linear probe jc2_188 N=1000000: (1 - macro_auc_ovr) of R16_Q1 over R42_Q1, run by run :: sha256 c6c4138ee5c86423 :: prose twin: two significant figures, direction in words
\newcommand{\MlstCmpJcTwoRandOverFull}{1.125\,\ensuremath{\pm}\,0.008}  % experiments/FIGS/data/v2/index.json :: linear probe jc2_188 N=1000000: (1 - macro_auc_ovr) of each random partition over the L188 run of its seed, mean +- SD over the partitions :: sha256 c6c4138ee5c86423
\newcommand{\MlstCmpJcTwoRandOverFullPr}{12\% more errors}  % experiments/FIGS/data/v2/index.json :: linear probe jc2_188 N=1000000: (1 - macro_auc_ovr) of each random partition over the L188 run of its seed, mean +- SD over the partitions :: sha256 c6c4138ee5c86423 :: prose twin: two significant figures, direction in words
\newcommand{\MlstCmpJcTwoRandOverMid}{0.918\,\ensuremath{\pm}\,0.007}  % experiments/FIGS/data/v2/index.json :: linear probe jc2_188 N=1000000: (1 - macro_auc_ovr) of each random partition over the R42_Q1 run of its seed, mean +- SD over the partitions :: sha256 c6c4138ee5c86423
\newcommand{\MlstCmpJcTwoRandOverMidPr}{8.2\% fewer errors}  % experiments/FIGS/data/v2/index.json :: linear probe jc2_188 N=1000000: (1 - macro_auc_ovr) of each random partition over the R42_Q1 run of its seed, mean +- SD over the partitions :: sha256 c6c4138ee5c86423 :: prose twin: two significant figures, direction in words
\newcommand{\MlstCmpJcTwoRandOverCoarse}{0.724\,\ensuremath{\pm}\,0.008}  % experiments/FIGS/data/v2/index.json :: linear probe jc2_188 N=1000000: (1 - macro_auc_ovr) of each random partition over the R16_Q1 run of its seed, mean +- SD over the partitions :: sha256 c6c4138ee5c86423
\newcommand{\MlstCmpJcTwoRandOverCoarsePr}{28\% fewer errors}  % experiments/FIGS/data/v2/index.json :: linear probe jc2_188 N=1000000: (1 - macro_auc_ovr) of each random partition over the R16_Q1 run of its seed, mean +- SD over the partitions :: sha256 c6c4138ee5c86423 :: prose twin: two significant figures, direction in words
\newcommand{\MlstCmpJcTwoTrainedResOverFull}{1.026\,\ensuremath{\pm}\,0.005}  % experiments/FIGS/data/v2/index.json :: linear probe jc2_188 N=1000000: (1 - macro_auc_trained) of L162 over L188, run by run :: sha256 c6c4138ee5c86423
\newcommand{\MlstCmpJcTwoTrainedResOverFullPr}{2.6\% more errors}  % experiments/FIGS/data/v2/index.json :: linear probe jc2_188 N=1000000: (1 - macro_auc_trained) of L162 over L188, run by run :: sha256 c6c4138ee5c86423 :: prose twin: two significant figures, direction in words
\newcommand{\MlstCmpJcTwoTrainedMidOverFull}{1.227\,\ensuremath{\pm}\,0.009}  % experiments/FIGS/data/v2/index.json :: linear probe jc2_188 N=1000000: (1 - macro_auc_trained) of R42_Q1 over L188, run by run :: sha256 c6c4138ee5c86423
\newcommand{\MlstCmpJcTwoTrainedMidOverFullPr}{23\% more errors}  % experiments/FIGS/data/v2/index.json :: linear probe jc2_188 N=1000000: (1 - macro_auc_trained) of R42_Q1 over L188, run by run :: sha256 c6c4138ee5c86423 :: prose twin: two significant figures, direction in words
\newcommand{\MlstCmpJcTwoTrainedCoarseOverFull}{1.557\,\ensuremath{\pm}\,0.019}  % experiments/FIGS/data/v2/index.json :: linear probe jc2_188 N=1000000: (1 - macro_auc_trained) of R16_Q1 over L188, run by run :: sha256 c6c4138ee5c86423
\newcommand{\MlstCmpJcTwoTrainedCoarseOverFullPr}{56\% more errors}  % experiments/FIGS/data/v2/index.json :: linear probe jc2_188 N=1000000: (1 - macro_auc_trained) of R16_Q1 over L188, run by run :: sha256 c6c4138ee5c86423 :: prose twin: two significant figures, direction in words
\newcommand{\MlstCmpJcTwoTrainedResMassOverRes}{1.022 and 1.016}  % experiments/FIGS/data/v2/index.json :: linear probe jc2_188 N=1000000: (1 - macro_auc_trained) of L162_MASS over L162, run by run :: sha256 c6c4138ee5c86423
\newcommand{\MlstCmpJcTwoTrainedResMassOverResPr}{equal within 3\%}  % experiments/FIGS/data/v2/index.json :: linear probe jc2_188 N=1000000: (1 - macro_auc_trained) of L162_MASS over L162, run by run :: sha256 c6c4138ee5c86423 :: prose twin: two significant figures, direction in words
\newcommand{\MlstCmpJcTwoTrainedResMassOverResPrMax}{equal within 3\%}  % experiments/FIGS/data/v2/index.json :: linear probe jc2_188 N=1000000: (1 - macro_auc_trained) of L162_MASS over L162, run by run :: sha256 c6c4138ee5c86423 :: prose twin of the run farther from 1
\newcommand{\MlstCmpJcTwoTrainedCoarseMassOverCoarse}{1.054 and 0.977}  % experiments/FIGS/data/v2/index.json :: linear probe jc2_188 N=1000000: (1 - macro_auc_trained) of R16_Q1_MASS over R16_Q1, run by run :: sha256 c6c4138ee5c86423
\newcommand{\MlstCmpJcTwoTrainedCoarseMassOverCoarsePr}{equal within 6\%}  % experiments/FIGS/data/v2/index.json :: linear probe jc2_188 N=1000000: (1 - macro_auc_trained) of R16_Q1_MASS over R16_Q1, run by run :: sha256 c6c4138ee5c86423 :: prose twin: two significant figures, direction in words
\newcommand{\MlstCmpJcTwoTrainedCoarseMassOverCoarsePrMax}{equal within 6\%}  % experiments/FIGS/data/v2/index.json :: linear probe jc2_188 N=1000000: (1 - macro_auc_trained) of R16_Q1_MASS over R16_Q1, run by run :: sha256 c6c4138ee5c86423 :: prose twin of the run farther from 1
\newcommand{\MlstCmpJcTwoTrainedResMassOverFull}{1.053 and 1.037}  % experiments/FIGS/data/v2/index.json :: linear probe jc2_188 N=1000000: (1 - macro_auc_trained) of L162_MASS over L188, run by run :: sha256 c6c4138ee5c86423
\newcommand{\MlstCmpJcTwoTrainedResMassOverFullPr}{at least 3.7\% more errors}  % experiments/FIGS/data/v2/index.json :: linear probe jc2_188 N=1000000: (1 - macro_auc_trained) of L162_MASS over L188, run by run :: sha256 c6c4138ee5c86423 :: prose twin: two significant figures, direction in words
\newcommand{\MlstCmpJcTwoTrainedResMassOverFullPrMax}{up to 5.3\% more errors}  % experiments/FIGS/data/v2/index.json :: linear probe jc2_188 N=1000000: (1 - macro_auc_trained) of L162_MASS over L188, run by run :: sha256 c6c4138ee5c86423 :: prose twin of the run farther from 1
\newcommand{\MlstCmpJcTwoTrainedCoarseMassOverFull}{1.617 and 1.528}  % experiments/FIGS/data/v2/index.json :: linear probe jc2_188 N=1000000: (1 - macro_auc_trained) of R16_Q1_MASS over L188, run by run :: sha256 c6c4138ee5c86423
\newcommand{\MlstCmpJcTwoTrainedCoarseMassOverFullPr}{at least 53\% more errors}  % experiments/FIGS/data/v2/index.json :: linear probe jc2_188 N=1000000: (1 - macro_auc_trained) of R16_Q1_MASS over L188, run by run :: sha256 c6c4138ee5c86423 :: prose twin: two significant figures, direction in words
\newcommand{\MlstCmpJcTwoTrainedCoarseMassOverFullPrMax}{up to 62\% more errors}  % experiments/FIGS/data/v2/index.json :: linear probe jc2_188 N=1000000: (1 - macro_auc_trained) of R16_Q1_MASS over L188, run by run :: sha256 c6c4138ee5c86423 :: prose twin of the run farther from 1
\newcommand{\MlstCmpJcTwoTrainedFzeroOverFull}{1.424 and 1.462}  % experiments/FIGS/data/v2/index.json :: linear probe jc2_188 N=1000000: (1 - macro_auc_trained) of FLAV_F0 over L188, run by run :: sha256 c6c4138ee5c86423
\newcommand{\MlstCmpJcTwoTrainedFzeroOverFullPr}{at least 42\% more errors}  % experiments/FIGS/data/v2/index.json :: linear probe jc2_188 N=1000000: (1 - macro_auc_trained) of FLAV_F0 over L188, run by run :: sha256 c6c4138ee5c86423 :: prose twin: two significant figures, direction in words
\newcommand{\MlstCmpJcTwoTrainedFzeroOverFullPrMax}{up to 46\% more errors}  % experiments/FIGS/data/v2/index.json :: linear probe jc2_188 N=1000000: (1 - macro_auc_trained) of FLAV_F0 over L188, run by run :: sha256 c6c4138ee5c86423 :: prose twin of the run farther from 1
\newcommand{\MlstCmpJcTwoTrainedFoneOverFull}{1.218 and 1.212}  % experiments/FIGS/data/v2/index.json :: linear probe jc2_188 N=1000000: (1 - macro_auc_trained) of FLAV_F1 over L188, run by run :: sha256 c6c4138ee5c86423
\newcommand{\MlstCmpJcTwoTrainedFoneOverFullPr}{at least 21\% more errors}  % experiments/FIGS/data/v2/index.json :: linear probe jc2_188 N=1000000: (1 - macro_auc_trained) of FLAV_F1 over L188, run by run :: sha256 c6c4138ee5c86423 :: prose twin: two significant figures, direction in words
\newcommand{\MlstCmpJcTwoTrainedFoneOverFullPrMax}{up to 22\% more errors}  % experiments/FIGS/data/v2/index.json :: linear probe jc2_188 N=1000000: (1 - macro_auc_trained) of FLAV_F1 over L188, run by run :: sha256 c6c4138ee5c86423 :: prose twin of the run farther from 1
\newcommand{\MlstCmpJcTwoTrainedFonerOverFull}{1.206 and 1.208}  % experiments/FIGS/data/v2/index.json :: linear probe jc2_188 N=1000000: (1 - macro_auc_trained) of FLAV_F1R over L188, run by run :: sha256 c6c4138ee5c86423
\newcommand{\MlstCmpJcTwoTrainedFonerOverFullPr}{at least 21\% more errors}  % experiments/FIGS/data/v2/index.json :: linear probe jc2_188 N=1000000: (1 - macro_auc_trained) of FLAV_F1R over L188, run by run :: sha256 c6c4138ee5c86423 :: prose twin: two significant figures, direction in words
\newcommand{\MlstCmpJcTwoTrainedFonerOverFullPrMax}{up to 21\% more errors}  % experiments/FIGS/data/v2/index.json :: linear probe jc2_188 N=1000000: (1 - macro_auc_trained) of FLAV_F1R over L188, run by run :: sha256 c6c4138ee5c86423 :: prose twin of the run farther from 1
\newcommand{\MlstCmpJcTwoTrainedFzeroOverCoarse}{0.9275 and 0.9345}  % experiments/FIGS/data/v2/index.json :: linear probe jc2_188 N=1000000: (1 - macro_auc_trained) of FLAV_F0 over R16_Q1, run by run :: sha256 c6c4138ee5c86423
\newcommand{\MlstCmpJcTwoTrainedFzeroOverCoarsePr}{at least 6.5\% fewer errors}  % experiments/FIGS/data/v2/index.json :: linear probe jc2_188 N=1000000: (1 - macro_auc_trained) of FLAV_F0 over R16_Q1, run by run :: sha256 c6c4138ee5c86423 :: prose twin: two significant figures, direction in words
\newcommand{\MlstCmpJcTwoTrainedFzeroOverCoarsePrMax}{up to 7.2\% fewer errors}  % experiments/FIGS/data/v2/index.json :: linear probe jc2_188 N=1000000: (1 - macro_auc_trained) of FLAV_F0 over R16_Q1, run by run :: sha256 c6c4138ee5c86423 :: prose twin of the run farther from 1
\newcommand{\MlstCmpJcTwoTrainedFoneOverCoarse}{0.794 and 0.775}  % experiments/FIGS/data/v2/index.json :: linear probe jc2_188 N=1000000: (1 - macro_auc_trained) of FLAV_F1 over R16_Q1, run by run :: sha256 c6c4138ee5c86423
\newcommand{\MlstCmpJcTwoTrainedFoneOverCoarsePr}{at least 21\% fewer errors}  % experiments/FIGS/data/v2/index.json :: linear probe jc2_188 N=1000000: (1 - macro_auc_trained) of FLAV_F1 over R16_Q1, run by run :: sha256 c6c4138ee5c86423 :: prose twin: two significant figures, direction in words
\newcommand{\MlstCmpJcTwoTrainedFoneOverCoarsePrMax}{up to 23\% fewer errors}  % experiments/FIGS/data/v2/index.json :: linear probe jc2_188 N=1000000: (1 - macro_auc_trained) of FLAV_F1 over R16_Q1, run by run :: sha256 c6c4138ee5c86423 :: prose twin of the run farther from 1
\newcommand{\MlstCmpJcTwoTrainedFonerOverCoarse}{0.786 and 0.772}  % experiments/FIGS/data/v2/index.json :: linear probe jc2_188 N=1000000: (1 - macro_auc_trained) of FLAV_F1R over R16_Q1, run by run :: sha256 c6c4138ee5c86423
\newcommand{\MlstCmpJcTwoTrainedFonerOverCoarsePr}{at least 21\% fewer errors}  % experiments/FIGS/data/v2/index.json :: linear probe jc2_188 N=1000000: (1 - macro_auc_trained) of FLAV_F1R over R16_Q1, run by run :: sha256 c6c4138ee5c86423 :: prose twin: two significant figures, direction in words
\newcommand{\MlstCmpJcTwoTrainedFonerOverCoarsePrMax}{up to 23\% fewer errors}  % experiments/FIGS/data/v2/index.json :: linear probe jc2_188 N=1000000: (1 - macro_auc_trained) of FLAV_F1R over R16_Q1, run by run :: sha256 c6c4138ee5c86423 :: prose twin of the run farther from 1
\newcommand{\MlstCmpJcTwoTrainedFoneOverFzero}{0.856 and 0.829}  % experiments/FIGS/data/v2/index.json :: linear probe jc2_188 N=1000000: (1 - macro_auc_trained) of FLAV_F1 over FLAV_F0, run by run :: sha256 c6c4138ee5c86423
\newcommand{\MlstCmpJcTwoTrainedFoneOverFzeroPr}{at least 14\% fewer errors}  % experiments/FIGS/data/v2/index.json :: linear probe jc2_188 N=1000000: (1 - macro_auc_trained) of FLAV_F1 over FLAV_F0, run by run :: sha256 c6c4138ee5c86423 :: prose twin: two significant figures, direction in words
\newcommand{\MlstCmpJcTwoTrainedFoneOverFzeroPrMax}{up to 17\% fewer errors}  % experiments/FIGS/data/v2/index.json :: linear probe jc2_188 N=1000000: (1 - macro_auc_trained) of FLAV_F1 over FLAV_F0, run by run :: sha256 c6c4138ee5c86423 :: prose twin of the run farther from 1
\newcommand{\MlstCmpJcTwoTrainedFonerOverFzero}{0.847 and 0.826}  % experiments/FIGS/data/v2/index.json :: linear probe jc2_188 N=1000000: (1 - macro_auc_trained) of FLAV_F1R over FLAV_F0, run by run :: sha256 c6c4138ee5c86423
\newcommand{\MlstCmpJcTwoTrainedFonerOverFzeroPr}{at least 15\% fewer errors}  % experiments/FIGS/data/v2/index.json :: linear probe jc2_188 N=1000000: (1 - macro_auc_trained) of FLAV_F1R over FLAV_F0, run by run :: sha256 c6c4138ee5c86423 :: prose twin: two significant figures, direction in words
\newcommand{\MlstCmpJcTwoTrainedFonerOverFzeroPrMax}{up to 17\% fewer errors}  % experiments/FIGS/data/v2/index.json :: linear probe jc2_188 N=1000000: (1 - macro_auc_trained) of FLAV_F1R over FLAV_F0, run by run :: sha256 c6c4138ee5c86423 :: prose twin of the run farther from 1
\newcommand{\MlstCmpJcTwoTrainedFonerOverFone}{0.9899 and 0.9966}  % experiments/FIGS/data/v2/index.json :: linear probe jc2_188 N=1000000: (1 - macro_auc_trained) of FLAV_F1R over FLAV_F1, run by run :: sha256 c6c4138ee5c86423
\newcommand{\MlstCmpJcTwoTrainedFonerOverFonePr}{equal within 2\%}  % experiments/FIGS/data/v2/index.json :: linear probe jc2_188 N=1000000: (1 - macro_auc_trained) of FLAV_F1R over FLAV_F1, run by run :: sha256 c6c4138ee5c86423 :: prose twin: two significant figures, direction in words
\newcommand{\MlstCmpJcTwoTrainedFonerOverFonePrMax}{equal within 2\%}  % experiments/FIGS/data/v2/index.json :: linear probe jc2_188 N=1000000: (1 - macro_auc_trained) of FLAV_F1R over FLAV_F1, run by run :: sha256 c6c4138ee5c86423 :: prose twin of the run farther from 1
\newcommand{\MlstCmpJcTwoTrainedMidOverRes}{1.197\,\ensuremath{\pm}\,0.013}  % experiments/FIGS/data/v2/index.json :: linear probe jc2_188 N=1000000: (1 - macro_auc_trained) of R42_Q1 over L162, run by run :: sha256 c6c4138ee5c86423
\newcommand{\MlstCmpJcTwoTrainedMidOverResPr}{20\% more errors}  % experiments/FIGS/data/v2/index.json :: linear probe jc2_188 N=1000000: (1 - macro_auc_trained) of R42_Q1 over L162, run by run :: sha256 c6c4138ee5c86423 :: prose twin: two significant figures, direction in words
\newcommand{\MlstCmpJcTwoTrainedCoarseOverRes}{1.518\,\ensuremath{\pm}\,0.025}  % experiments/FIGS/data/v2/index.json :: linear probe jc2_188 N=1000000: (1 - macro_auc_trained) of R16_Q1 over L162, run by run :: sha256 c6c4138ee5c86423
\newcommand{\MlstCmpJcTwoTrainedCoarseOverResPr}{52\% more errors}  % experiments/FIGS/data/v2/index.json :: linear probe jc2_188 N=1000000: (1 - macro_auc_trained) of R16_Q1 over L162, run by run :: sha256 c6c4138ee5c86423 :: prose twin: two significant figures, direction in words
\newcommand{\MlstCmpJcTwoTrainedCoarseOverMid}{1.269\,\ensuremath{\pm}\,0.007}  % experiments/FIGS/data/v2/index.json :: linear probe jc2_188 N=1000000: (1 - macro_auc_trained) of R16_Q1 over R42_Q1, run by run :: sha256 c6c4138ee5c86423
\newcommand{\MlstCmpJcTwoTrainedCoarseOverMidPr}{27\% more errors}  % experiments/FIGS/data/v2/index.json :: linear probe jc2_188 N=1000000: (1 - macro_auc_trained) of R16_Q1 over R42_Q1, run by run :: sha256 c6c4138ee5c86423 :: prose twin: two significant figures, direction in words
\newcommand{\MlstCmpJcTwoTrainedRandOverFull}{1.125\,\ensuremath{\pm}\,0.008}  % experiments/FIGS/data/v2/index.json :: linear probe jc2_188 N=1000000: (1 - macro_auc_trained) of each random partition over the L188 run of its seed, mean +- SD over the partitions :: sha256 c6c4138ee5c86423
\newcommand{\MlstCmpJcTwoTrainedRandOverFullPr}{12\% more errors}  % experiments/FIGS/data/v2/index.json :: linear probe jc2_188 N=1000000: (1 - macro_auc_trained) of each random partition over the L188 run of its seed, mean +- SD over the partitions :: sha256 c6c4138ee5c86423 :: prose twin: two significant figures, direction in words
\newcommand{\MlstCmpJcTwoTrainedRandOverMid}{0.918\,\ensuremath{\pm}\,0.007}  % experiments/FIGS/data/v2/index.json :: linear probe jc2_188 N=1000000: (1 - macro_auc_trained) of each random partition over the R42_Q1 run of its seed, mean +- SD over the partitions :: sha256 c6c4138ee5c86423
\newcommand{\MlstCmpJcTwoTrainedRandOverMidPr}{8.2\% fewer errors}  % experiments/FIGS/data/v2/index.json :: linear probe jc2_188 N=1000000: (1 - macro_auc_trained) of each random partition over the R42_Q1 run of its seed, mean +- SD over the partitions :: sha256 c6c4138ee5c86423 :: prose twin: two significant figures, direction in words
\newcommand{\MlstCmpJcTwoTrainedRandOverCoarse}{0.724\,\ensuremath{\pm}\,0.008}  % experiments/FIGS/data/v2/index.json :: linear probe jc2_188 N=1000000: (1 - macro_auc_trained) of each random partition over the R16_Q1 run of its seed, mean +- SD over the partitions :: sha256 c6c4138ee5c86423
\newcommand{\MlstCmpJcTwoTrainedRandOverCoarsePr}{28\% fewer errors}  % experiments/FIGS/data/v2/index.json :: linear probe jc2_188 N=1000000: (1 - macro_auc_trained) of each random partition over the R16_Q1 run of its seed, mean +- SD over the partitions :: sha256 c6c4138ee5c86423 :: prose twin: two significant figures, direction in words
\newcommand{\MlstCmpJcTwoResResOverFull}{1.003\,\ensuremath{\pm}\,0.005}  % experiments/FIGS/data/v2/index.json :: linear probe jc2_188 N=1000000: (1 - macro_auc_resonant) of L162 over L188, run by run :: sha256 c6c4138ee5c86423
\newcommand{\MlstCmpJcTwoResResOverFullPr}{equal within 1\%}  % experiments/FIGS/data/v2/index.json :: linear probe jc2_188 N=1000000: (1 - macro_auc_resonant) of L162 over L188, run by run :: sha256 c6c4138ee5c86423 :: prose twin: two significant figures, direction in words
\newcommand{\MlstCmpJcTwoResMidOverFull}{1.262\,\ensuremath{\pm}\,0.012}  % experiments/FIGS/data/v2/index.json :: linear probe jc2_188 N=1000000: (1 - macro_auc_resonant) of R42_Q1 over L188, run by run :: sha256 c6c4138ee5c86423
\newcommand{\MlstCmpJcTwoResMidOverFullPr}{26\% more errors}  % experiments/FIGS/data/v2/index.json :: linear probe jc2_188 N=1000000: (1 - macro_auc_resonant) of R42_Q1 over L188, run by run :: sha256 c6c4138ee5c86423 :: prose twin: two significant figures, direction in words
\newcommand{\MlstCmpJcTwoResCoarseOverFull}{1.590\,\ensuremath{\pm}\,0.032}  % experiments/FIGS/data/v2/index.json :: linear probe jc2_188 N=1000000: (1 - macro_auc_resonant) of R16_Q1 over L188, run by run :: sha256 c6c4138ee5c86423
\newcommand{\MlstCmpJcTwoResCoarseOverFullPr}{59\% more errors}  % experiments/FIGS/data/v2/index.json :: linear probe jc2_188 N=1000000: (1 - macro_auc_resonant) of R16_Q1 over L188, run by run :: sha256 c6c4138ee5c86423 :: prose twin: two significant figures, direction in words
\newcommand{\MlstCmpJcTwoResResMassOverRes}{1.025 and 1.020}  % experiments/FIGS/data/v2/index.json :: linear probe jc2_188 N=1000000: (1 - macro_auc_resonant) of L162_MASS over L162, run by run :: sha256 c6c4138ee5c86423
\newcommand{\MlstCmpJcTwoResResMassOverResPr}{equal within 3\%}  % experiments/FIGS/data/v2/index.json :: linear probe jc2_188 N=1000000: (1 - macro_auc_resonant) of L162_MASS over L162, run by run :: sha256 c6c4138ee5c86423 :: prose twin: two significant figures, direction in words
\newcommand{\MlstCmpJcTwoResResMassOverResPrMax}{equal within 3\%}  % experiments/FIGS/data/v2/index.json :: linear probe jc2_188 N=1000000: (1 - macro_auc_resonant) of L162_MASS over L162, run by run :: sha256 c6c4138ee5c86423 :: prose twin of the run farther from 1
\newcommand{\MlstCmpJcTwoResCoarseMassOverCoarse}{1.08 and 0.96}  % experiments/FIGS/data/v2/index.json :: linear probe jc2_188 N=1000000: (1 - macro_auc_resonant) of R16_Q1_MASS over R16_Q1, run by run :: sha256 c6c4138ee5c86423
\newcommand{\MlstCmpJcTwoResCoarseMassOverCoarsePr}{equal within 8\%}  % experiments/FIGS/data/v2/index.json :: linear probe jc2_188 N=1000000: (1 - macro_auc_resonant) of R16_Q1_MASS over R16_Q1, run by run :: sha256 c6c4138ee5c86423 :: prose twin: two significant figures, direction in words
\newcommand{\MlstCmpJcTwoResCoarseMassOverCoarsePrMax}{equal within 8\%}  % experiments/FIGS/data/v2/index.json :: linear probe jc2_188 N=1000000: (1 - macro_auc_resonant) of R16_Q1_MASS over R16_Q1, run by run :: sha256 c6c4138ee5c86423 :: prose twin of the run farther from 1
\newcommand{\MlstCmpJcTwoResResMassOverFull}{1.033 and 1.018}  % experiments/FIGS/data/v2/index.json :: linear probe jc2_188 N=1000000: (1 - macro_auc_resonant) of L162_MASS over L188, run by run :: sha256 c6c4138ee5c86423
\newcommand{\MlstCmpJcTwoResResMassOverFullPr}{equal within 4\%}  % experiments/FIGS/data/v2/index.json :: linear probe jc2_188 N=1000000: (1 - macro_auc_resonant) of L162_MASS over L188, run by run :: sha256 c6c4138ee5c86423 :: prose twin: two significant figures, direction in words
\newcommand{\MlstCmpJcTwoResResMassOverFullPrMax}{equal within 4\%}  % experiments/FIGS/data/v2/index.json :: linear probe jc2_188 N=1000000: (1 - macro_auc_resonant) of L162_MASS over L188, run by run :: sha256 c6c4138ee5c86423 :: prose twin of the run farther from 1
\newcommand{\MlstCmpJcTwoResCoarseMassOverFull}{1.67 and 1.55}  % experiments/FIGS/data/v2/index.json :: linear probe jc2_188 N=1000000: (1 - macro_auc_resonant) of R16_Q1_MASS over L188, run by run :: sha256 c6c4138ee5c86423
\newcommand{\MlstCmpJcTwoResCoarseMassOverFullPr}{at least 55\% more errors}  % experiments/FIGS/data/v2/index.json :: linear probe jc2_188 N=1000000: (1 - macro_auc_resonant) of R16_Q1_MASS over L188, run by run :: sha256 c6c4138ee5c86423 :: prose twin: two significant figures, direction in words
\newcommand{\MlstCmpJcTwoResCoarseMassOverFullPrMax}{up to 67\% more errors}  % experiments/FIGS/data/v2/index.json :: linear probe jc2_188 N=1000000: (1 - macro_auc_resonant) of R16_Q1_MASS over L188, run by run :: sha256 c6c4138ee5c86423 :: prose twin of the run farther from 1
\newcommand{\MlstCmpJcTwoResFzeroOverFull}{1.408 and 1.449}  % experiments/FIGS/data/v2/index.json :: linear probe jc2_188 N=1000000: (1 - macro_auc_resonant) of FLAV_F0 over L188, run by run :: sha256 c6c4138ee5c86423
\newcommand{\MlstCmpJcTwoResFzeroOverFullPr}{at least 41\% more errors}  % experiments/FIGS/data/v2/index.json :: linear probe jc2_188 N=1000000: (1 - macro_auc_resonant) of FLAV_F0 over L188, run by run :: sha256 c6c4138ee5c86423 :: prose twin: two significant figures, direction in words
\newcommand{\MlstCmpJcTwoResFzeroOverFullPrMax}{up to 45\% more errors}  % experiments/FIGS/data/v2/index.json :: linear probe jc2_188 N=1000000: (1 - macro_auc_resonant) of FLAV_F0 over L188, run by run :: sha256 c6c4138ee5c86423 :: prose twin of the run farther from 1
\newcommand{\MlstCmpJcTwoResFoneOverFull}{1.210 and 1.204}  % experiments/FIGS/data/v2/index.json :: linear probe jc2_188 N=1000000: (1 - macro_auc_resonant) of FLAV_F1 over L188, run by run :: sha256 c6c4138ee5c86423
\newcommand{\MlstCmpJcTwoResFoneOverFullPr}{at least 20\% more errors}  % experiments/FIGS/data/v2/index.json :: linear probe jc2_188 N=1000000: (1 - macro_auc_resonant) of FLAV_F1 over L188, run by run :: sha256 c6c4138ee5c86423 :: prose twin: two significant figures, direction in words
\newcommand{\MlstCmpJcTwoResFoneOverFullPrMax}{up to 21\% more errors}  % experiments/FIGS/data/v2/index.json :: linear probe jc2_188 N=1000000: (1 - macro_auc_resonant) of FLAV_F1 over L188, run by run :: sha256 c6c4138ee5c86423 :: prose twin of the run farther from 1
\newcommand{\MlstCmpJcTwoResFonerOverFull}{1.187 and 1.188}  % experiments/FIGS/data/v2/index.json :: linear probe jc2_188 N=1000000: (1 - macro_auc_resonant) of FLAV_F1R over L188, run by run :: sha256 c6c4138ee5c86423
\newcommand{\MlstCmpJcTwoResFonerOverFullPr}{at least 19\% more errors}  % experiments/FIGS/data/v2/index.json :: linear probe jc2_188 N=1000000: (1 - macro_auc_resonant) of FLAV_F1R over L188, run by run :: sha256 c6c4138ee5c86423 :: prose twin: two significant figures, direction in words
\newcommand{\MlstCmpJcTwoResFonerOverFullPrMax}{up to 19\% more errors}  % experiments/FIGS/data/v2/index.json :: linear probe jc2_188 N=1000000: (1 - macro_auc_resonant) of FLAV_F1R over L188, run by run :: sha256 c6c4138ee5c86423 :: prose twin of the run farther from 1
\newcommand{\MlstCmpJcTwoResFzeroOverCoarse}{0.9069 and 0.9029}  % experiments/FIGS/data/v2/index.json :: linear probe jc2_188 N=1000000: (1 - macro_auc_resonant) of FLAV_F0 over R16_Q1, run by run :: sha256 c6c4138ee5c86423
\newcommand{\MlstCmpJcTwoResFzeroOverCoarsePr}{at least 9.3\% fewer errors}  % experiments/FIGS/data/v2/index.json :: linear probe jc2_188 N=1000000: (1 - macro_auc_resonant) of FLAV_F0 over R16_Q1, run by run :: sha256 c6c4138ee5c86423 :: prose twin: two significant figures, direction in words
\newcommand{\MlstCmpJcTwoResFzeroOverCoarsePrMax}{up to 9.7\% fewer errors}  % experiments/FIGS/data/v2/index.json :: linear probe jc2_188 N=1000000: (1 - macro_auc_resonant) of FLAV_F0 over R16_Q1, run by run :: sha256 c6c4138ee5c86423 :: prose twin of the run farther from 1
\newcommand{\MlstCmpJcTwoResFoneOverCoarse}{0.779 and 0.750}  % experiments/FIGS/data/v2/index.json :: linear probe jc2_188 N=1000000: (1 - macro_auc_resonant) of FLAV_F1 over R16_Q1, run by run :: sha256 c6c4138ee5c86423
\newcommand{\MlstCmpJcTwoResFoneOverCoarsePr}{at least 22\% fewer errors}  % experiments/FIGS/data/v2/index.json :: linear probe jc2_188 N=1000000: (1 - macro_auc_resonant) of FLAV_F1 over R16_Q1, run by run :: sha256 c6c4138ee5c86423 :: prose twin: two significant figures, direction in words
\newcommand{\MlstCmpJcTwoResFoneOverCoarsePrMax}{up to 25\% fewer errors}  % experiments/FIGS/data/v2/index.json :: linear probe jc2_188 N=1000000: (1 - macro_auc_resonant) of FLAV_F1 over R16_Q1, run by run :: sha256 c6c4138ee5c86423 :: prose twin of the run farther from 1
\newcommand{\MlstCmpJcTwoResFonerOverCoarse}{0.764 and 0.740}  % experiments/FIGS/data/v2/index.json :: linear probe jc2_188 N=1000000: (1 - macro_auc_resonant) of FLAV_F1R over R16_Q1, run by run :: sha256 c6c4138ee5c86423
\newcommand{\MlstCmpJcTwoResFonerOverCoarsePr}{at least 24\% fewer errors}  % experiments/FIGS/data/v2/index.json :: linear probe jc2_188 N=1000000: (1 - macro_auc_resonant) of FLAV_F1R over R16_Q1, run by run :: sha256 c6c4138ee5c86423 :: prose twin: two significant figures, direction in words
\newcommand{\MlstCmpJcTwoResFonerOverCoarsePrMax}{up to 26\% fewer errors}  % experiments/FIGS/data/v2/index.json :: linear probe jc2_188 N=1000000: (1 - macro_auc_resonant) of FLAV_F1R over R16_Q1, run by run :: sha256 c6c4138ee5c86423 :: prose twin of the run farther from 1
\newcommand{\MlstCmpJcTwoResFoneOverFzero}{0.859 and 0.831}  % experiments/FIGS/data/v2/index.json :: linear probe jc2_188 N=1000000: (1 - macro_auc_resonant) of FLAV_F1 over FLAV_F0, run by run :: sha256 c6c4138ee5c86423
\newcommand{\MlstCmpJcTwoResFoneOverFzeroPr}{at least 14\% fewer errors}  % experiments/FIGS/data/v2/index.json :: linear probe jc2_188 N=1000000: (1 - macro_auc_resonant) of FLAV_F1 over FLAV_F0, run by run :: sha256 c6c4138ee5c86423 :: prose twin: two significant figures, direction in words
\newcommand{\MlstCmpJcTwoResFoneOverFzeroPrMax}{up to 17\% fewer errors}  % experiments/FIGS/data/v2/index.json :: linear probe jc2_188 N=1000000: (1 - macro_auc_resonant) of FLAV_F1 over FLAV_F0, run by run :: sha256 c6c4138ee5c86423 :: prose twin of the run farther from 1
\newcommand{\MlstCmpJcTwoResFonerOverFzero}{0.843 and 0.820}  % experiments/FIGS/data/v2/index.json :: linear probe jc2_188 N=1000000: (1 - macro_auc_resonant) of FLAV_F1R over FLAV_F0, run by run :: sha256 c6c4138ee5c86423
\newcommand{\MlstCmpJcTwoResFonerOverFzeroPr}{at least 16\% fewer errors}  % experiments/FIGS/data/v2/index.json :: linear probe jc2_188 N=1000000: (1 - macro_auc_resonant) of FLAV_F1R over FLAV_F0, run by run :: sha256 c6c4138ee5c86423 :: prose twin: two significant figures, direction in words
\newcommand{\MlstCmpJcTwoResFonerOverFzeroPrMax}{up to 18\% fewer errors}  % experiments/FIGS/data/v2/index.json :: linear probe jc2_188 N=1000000: (1 - macro_auc_resonant) of FLAV_F1R over FLAV_F0, run by run :: sha256 c6c4138ee5c86423 :: prose twin of the run farther from 1
\newcommand{\MlstCmpJcTwoResFonerOverFone}{0.9806 and 0.9867}  % experiments/FIGS/data/v2/index.json :: linear probe jc2_188 N=1000000: (1 - macro_auc_resonant) of FLAV_F1R over FLAV_F1, run by run :: sha256 c6c4138ee5c86423
\newcommand{\MlstCmpJcTwoResFonerOverFonePr}{equal within 2\%}  % experiments/FIGS/data/v2/index.json :: linear probe jc2_188 N=1000000: (1 - macro_auc_resonant) of FLAV_F1R over FLAV_F1, run by run :: sha256 c6c4138ee5c86423 :: prose twin: two significant figures, direction in words
\newcommand{\MlstCmpJcTwoResFonerOverFonePrMax}{equal within 2\%}  % experiments/FIGS/data/v2/index.json :: linear probe jc2_188 N=1000000: (1 - macro_auc_resonant) of FLAV_F1R over FLAV_F1, run by run :: sha256 c6c4138ee5c86423 :: prose twin of the run farther from 1
\newcommand{\MlstCmpJcTwoResMidOverRes}{1.258\,\ensuremath{\pm}\,0.017}  % experiments/FIGS/data/v2/index.json :: linear probe jc2_188 N=1000000: (1 - macro_auc_resonant) of R42_Q1 over L162, run by run :: sha256 c6c4138ee5c86423
\newcommand{\MlstCmpJcTwoResMidOverResPr}{26\% more errors}  % experiments/FIGS/data/v2/index.json :: linear probe jc2_188 N=1000000: (1 - macro_auc_resonant) of R42_Q1 over L162, run by run :: sha256 c6c4138ee5c86423 :: prose twin: two significant figures, direction in words
\newcommand{\MlstCmpJcTwoResCoarseOverRes}{1.58\,\ensuremath{\pm}\,0.04}  % experiments/FIGS/data/v2/index.json :: linear probe jc2_188 N=1000000: (1 - macro_auc_resonant) of R16_Q1 over L162, run by run :: sha256 c6c4138ee5c86423
\newcommand{\MlstCmpJcTwoResCoarseOverResPr}{58\% more errors}  % experiments/FIGS/data/v2/index.json :: linear probe jc2_188 N=1000000: (1 - macro_auc_resonant) of R16_Q1 over L162, run by run :: sha256 c6c4138ee5c86423 :: prose twin: two significant figures, direction in words
\newcommand{\MlstCmpJcTwoResCoarseOverMid}{1.260\,\ensuremath{\pm}\,0.015}  % experiments/FIGS/data/v2/index.json :: linear probe jc2_188 N=1000000: (1 - macro_auc_resonant) of R16_Q1 over R42_Q1, run by run :: sha256 c6c4138ee5c86423
\newcommand{\MlstCmpJcTwoResCoarseOverMidPr}{26\% more errors}  % experiments/FIGS/data/v2/index.json :: linear probe jc2_188 N=1000000: (1 - macro_auc_resonant) of R16_Q1 over R42_Q1, run by run :: sha256 c6c4138ee5c86423 :: prose twin: two significant figures, direction in words
\newcommand{\MlstCmpJcTwoResRandOverFull}{1.098\,\ensuremath{\pm}\,0.008}  % experiments/FIGS/data/v2/index.json :: linear probe jc2_188 N=1000000: (1 - macro_auc_resonant) of each random partition over the L188 run of its seed, mean +- SD over the partitions :: sha256 c6c4138ee5c86423
\newcommand{\MlstCmpJcTwoResRandOverFullPr}{9.8\% more errors}  % experiments/FIGS/data/v2/index.json :: linear probe jc2_188 N=1000000: (1 - macro_auc_resonant) of each random partition over the L188 run of its seed, mean +- SD over the partitions :: sha256 c6c4138ee5c86423 :: prose twin: two significant figures, direction in words
\newcommand{\MlstCmpJcTwoResRandOverMid}{0.870\,\ensuremath{\pm}\,0.008}  % experiments/FIGS/data/v2/index.json :: linear probe jc2_188 N=1000000: (1 - macro_auc_resonant) of each random partition over the R42_Q1 run of its seed, mean +- SD over the partitions :: sha256 c6c4138ee5c86423
\newcommand{\MlstCmpJcTwoResRandOverMidPr}{13\% fewer errors}  % experiments/FIGS/data/v2/index.json :: linear probe jc2_188 N=1000000: (1 - macro_auc_resonant) of each random partition over the R42_Q1 run of its seed, mean +- SD over the partitions :: sha256 c6c4138ee5c86423 :: prose twin: two significant figures, direction in words
\newcommand{\MlstCmpJcTwoResRandOverCoarse}{0.693\,\ensuremath{\pm}\,0.011}  % experiments/FIGS/data/v2/index.json :: linear probe jc2_188 N=1000000: (1 - macro_auc_resonant) of each random partition over the R16_Q1 run of its seed, mean +- SD over the partitions :: sha256 c6c4138ee5c86423
\newcommand{\MlstCmpJcTwoResRandOverCoarsePr}{31\% fewer errors}  % experiments/FIGS/data/v2/index.json :: linear probe jc2_188 N=1000000: (1 - macro_auc_resonant) of each random partition over the R16_Q1 run of its seed, mean +- SD over the partitions :: sha256 c6c4138ee5c86423 :: prose twin: two significant figures, direction in words
\newcommand{\MlstCmpJcTwoQcdResOverFull}{1.087\,\ensuremath{\pm}\,0.006}  % experiments/FIGS/data/v2/index.json :: linear probe jc2_188 N=1000000: (1 - macro_auc_qcd) of L162 over L188, run by run :: sha256 c6c4138ee5c86423
\newcommand{\MlstCmpJcTwoQcdResOverFullPr}{8.7\% more errors}  % experiments/FIGS/data/v2/index.json :: linear probe jc2_188 N=1000000: (1 - macro_auc_qcd) of L162 over L188, run by run :: sha256 c6c4138ee5c86423 :: prose twin: two significant figures, direction in words
\newcommand{\MlstCmpJcTwoQcdMidOverFull}{1.131\,\ensuremath{\pm}\,0.007}  % experiments/FIGS/data/v2/index.json :: linear probe jc2_188 N=1000000: (1 - macro_auc_qcd) of R42_Q1 over L188, run by run :: sha256 c6c4138ee5c86423
\newcommand{\MlstCmpJcTwoQcdMidOverFullPr}{13\% more errors}  % experiments/FIGS/data/v2/index.json :: linear probe jc2_188 N=1000000: (1 - macro_auc_qcd) of R42_Q1 over L188, run by run :: sha256 c6c4138ee5c86423 :: prose twin: two significant figures, direction in words
\newcommand{\MlstCmpJcTwoQcdCoarseOverFull}{1.466\,\ensuremath{\pm}\,0.017}  % experiments/FIGS/data/v2/index.json :: linear probe jc2_188 N=1000000: (1 - macro_auc_qcd) of R16_Q1 over L188, run by run :: sha256 c6c4138ee5c86423
\newcommand{\MlstCmpJcTwoQcdCoarseOverFullPr}{47\% more errors}  % experiments/FIGS/data/v2/index.json :: linear probe jc2_188 N=1000000: (1 - macro_auc_qcd) of R16_Q1 over L188, run by run :: sha256 c6c4138ee5c86423 :: prose twin: two significant figures, direction in words
\newcommand{\MlstCmpJcTwoQcdResMassOverRes}{1.012 and 1.006}  % experiments/FIGS/data/v2/index.json :: linear probe jc2_188 N=1000000: (1 - macro_auc_qcd) of L162_MASS over L162, run by run :: sha256 c6c4138ee5c86423
\newcommand{\MlstCmpJcTwoQcdResMassOverResPr}{equal within 2\%}  % experiments/FIGS/data/v2/index.json :: linear probe jc2_188 N=1000000: (1 - macro_auc_qcd) of L162_MASS over L162, run by run :: sha256 c6c4138ee5c86423 :: prose twin: two significant figures, direction in words
\newcommand{\MlstCmpJcTwoQcdResMassOverResPrMax}{equal within 2\%}  % experiments/FIGS/data/v2/index.json :: linear probe jc2_188 N=1000000: (1 - macro_auc_qcd) of L162_MASS over L162, run by run :: sha256 c6c4138ee5c86423 :: prose twin of the run farther from 1
\newcommand{\MlstCmpJcTwoQcdCoarseMassOverCoarse}{0.991 and 1.014}  % experiments/FIGS/data/v2/index.json :: linear probe jc2_188 N=1000000: (1 - macro_auc_qcd) of R16_Q1_MASS over R16_Q1, run by run :: sha256 c6c4138ee5c86423
\newcommand{\MlstCmpJcTwoQcdCoarseMassOverCoarsePr}{equal within 2\%}  % experiments/FIGS/data/v2/index.json :: linear probe jc2_188 N=1000000: (1 - macro_auc_qcd) of R16_Q1_MASS over R16_Q1, run by run :: sha256 c6c4138ee5c86423 :: prose twin: two significant figures, direction in words
\newcommand{\MlstCmpJcTwoQcdCoarseMassOverCoarsePrMax}{equal within 2\%}  % experiments/FIGS/data/v2/index.json :: linear probe jc2_188 N=1000000: (1 - macro_auc_qcd) of R16_Q1_MASS over R16_Q1, run by run :: sha256 c6c4138ee5c86423 :: prose twin of the run farther from 1
\newcommand{\MlstCmpJcTwoQcdResMassOverFull}{1.107 and 1.092}  % experiments/FIGS/data/v2/index.json :: linear probe jc2_188 N=1000000: (1 - macro_auc_qcd) of L162_MASS over L188, run by run :: sha256 c6c4138ee5c86423
\newcommand{\MlstCmpJcTwoQcdResMassOverFullPr}{at least 9.2\% more errors}  % experiments/FIGS/data/v2/index.json :: linear probe jc2_188 N=1000000: (1 - macro_auc_qcd) of L162_MASS over L188, run by run :: sha256 c6c4138ee5c86423 :: prose twin: two significant figures, direction in words
\newcommand{\MlstCmpJcTwoQcdResMassOverFullPrMax}{up to 11\% more errors}  % experiments/FIGS/data/v2/index.json :: linear probe jc2_188 N=1000000: (1 - macro_auc_qcd) of L162_MASS over L188, run by run :: sha256 c6c4138ee5c86423 :: prose twin of the run farther from 1
\newcommand{\MlstCmpJcTwoQcdCoarseMassOverFull}{1.472 and 1.474}  % experiments/FIGS/data/v2/index.json :: linear probe jc2_188 N=1000000: (1 - macro_auc_qcd) of R16_Q1_MASS over L188, run by run :: sha256 c6c4138ee5c86423
\newcommand{\MlstCmpJcTwoQcdCoarseMassOverFullPr}{at least 47\% more errors}  % experiments/FIGS/data/v2/index.json :: linear probe jc2_188 N=1000000: (1 - macro_auc_qcd) of R16_Q1_MASS over L188, run by run :: sha256 c6c4138ee5c86423 :: prose twin: two significant figures, direction in words
\newcommand{\MlstCmpJcTwoQcdCoarseMassOverFullPrMax}{up to 47\% more errors}  % experiments/FIGS/data/v2/index.json :: linear probe jc2_188 N=1000000: (1 - macro_auc_qcd) of R16_Q1_MASS over L188, run by run :: sha256 c6c4138ee5c86423 :: prose twin of the run farther from 1
\newcommand{\MlstCmpJcTwoQcdFzeroOverFull}{1.466 and 1.497}  % experiments/FIGS/data/v2/index.json :: linear probe jc2_188 N=1000000: (1 - macro_auc_qcd) of FLAV_F0 over L188, run by run :: sha256 c6c4138ee5c86423
\newcommand{\MlstCmpJcTwoQcdFzeroOverFullPr}{at least 47\% more errors}  % experiments/FIGS/data/v2/index.json :: linear probe jc2_188 N=1000000: (1 - macro_auc_qcd) of FLAV_F0 over L188, run by run :: sha256 c6c4138ee5c86423 :: prose twin: two significant figures, direction in words
\newcommand{\MlstCmpJcTwoQcdFzeroOverFullPrMax}{up to 50\% more errors}  % experiments/FIGS/data/v2/index.json :: linear probe jc2_188 N=1000000: (1 - macro_auc_qcd) of FLAV_F0 over L188, run by run :: sha256 c6c4138ee5c86423 :: prose twin of the run farther from 1
\newcommand{\MlstCmpJcTwoQcdFoneOverFull}{1.241 and 1.235}  % experiments/FIGS/data/v2/index.json :: linear probe jc2_188 N=1000000: (1 - macro_auc_qcd) of FLAV_F1 over L188, run by run :: sha256 c6c4138ee5c86423
\newcommand{\MlstCmpJcTwoQcdFoneOverFullPr}{at least 23\% more errors}  % experiments/FIGS/data/v2/index.json :: linear probe jc2_188 N=1000000: (1 - macro_auc_qcd) of FLAV_F1 over L188, run by run :: sha256 c6c4138ee5c86423 :: prose twin: two significant figures, direction in words
\newcommand{\MlstCmpJcTwoQcdFoneOverFullPrMax}{up to 24\% more errors}  % experiments/FIGS/data/v2/index.json :: linear probe jc2_188 N=1000000: (1 - macro_auc_qcd) of FLAV_F1 over L188, run by run :: sha256 c6c4138ee5c86423 :: prose twin of the run farther from 1
\newcommand{\MlstCmpJcTwoQcdFonerOverFull}{1.260 and 1.263}  % experiments/FIGS/data/v2/index.json :: linear probe jc2_188 N=1000000: (1 - macro_auc_qcd) of FLAV_F1R over L188, run by run :: sha256 c6c4138ee5c86423
\newcommand{\MlstCmpJcTwoQcdFonerOverFullPr}{at least 26\% more errors}  % experiments/FIGS/data/v2/index.json :: linear probe jc2_188 N=1000000: (1 - macro_auc_qcd) of FLAV_F1R over L188, run by run :: sha256 c6c4138ee5c86423 :: prose twin: two significant figures, direction in words
\newcommand{\MlstCmpJcTwoQcdFonerOverFullPrMax}{up to 26\% more errors}  % experiments/FIGS/data/v2/index.json :: linear probe jc2_188 N=1000000: (1 - macro_auc_qcd) of FLAV_F1R over L188, run by run :: sha256 c6c4138ee5c86423 :: prose twin of the run farther from 1
\newcommand{\MlstCmpJcTwoQcdFzeroOverCoarse}{0.987 and 1.031}  % experiments/FIGS/data/v2/index.json :: linear probe jc2_188 N=1000000: (1 - macro_auc_qcd) of FLAV_F0 over R16_Q1, run by run :: sha256 c6c4138ee5c86423
\newcommand{\MlstCmpJcTwoQcdFzeroOverCoarsePr}{equal within 4\%}  % experiments/FIGS/data/v2/index.json :: linear probe jc2_188 N=1000000: (1 - macro_auc_qcd) of FLAV_F0 over R16_Q1, run by run :: sha256 c6c4138ee5c86423 :: prose twin: two significant figures, direction in words
\newcommand{\MlstCmpJcTwoQcdFzeroOverCoarsePrMax}{equal within 4\%}  % experiments/FIGS/data/v2/index.json :: linear probe jc2_188 N=1000000: (1 - macro_auc_qcd) of FLAV_F0 over R16_Q1, run by run :: sha256 c6c4138ee5c86423 :: prose twin of the run farther from 1
\newcommand{\MlstCmpJcTwoQcdFoneOverCoarse}{0.836 and 0.850}  % experiments/FIGS/data/v2/index.json :: linear probe jc2_188 N=1000000: (1 - macro_auc_qcd) of FLAV_F1 over R16_Q1, run by run :: sha256 c6c4138ee5c86423
\newcommand{\MlstCmpJcTwoQcdFoneOverCoarsePr}{at least 15\% fewer errors}  % experiments/FIGS/data/v2/index.json :: linear probe jc2_188 N=1000000: (1 - macro_auc_qcd) of FLAV_F1 over R16_Q1, run by run :: sha256 c6c4138ee5c86423 :: prose twin: two significant figures, direction in words
\newcommand{\MlstCmpJcTwoQcdFoneOverCoarsePrMax}{up to 16\% fewer errors}  % experiments/FIGS/data/v2/index.json :: linear probe jc2_188 N=1000000: (1 - macro_auc_qcd) of FLAV_F1 over R16_Q1, run by run :: sha256 c6c4138ee5c86423 :: prose twin of the run farther from 1
\newcommand{\MlstCmpJcTwoQcdFonerOverCoarse}{0.848 and 0.870}  % experiments/FIGS/data/v2/index.json :: linear probe jc2_188 N=1000000: (1 - macro_auc_qcd) of FLAV_F1R over R16_Q1, run by run :: sha256 c6c4138ee5c86423
\newcommand{\MlstCmpJcTwoQcdFonerOverCoarsePr}{at least 13\% fewer errors}  % experiments/FIGS/data/v2/index.json :: linear probe jc2_188 N=1000000: (1 - macro_auc_qcd) of FLAV_F1R over R16_Q1, run by run :: sha256 c6c4138ee5c86423 :: prose twin: two significant figures, direction in words
\newcommand{\MlstCmpJcTwoQcdFonerOverCoarsePrMax}{up to 15\% fewer errors}  % experiments/FIGS/data/v2/index.json :: linear probe jc2_188 N=1000000: (1 - macro_auc_qcd) of FLAV_F1R over R16_Q1, run by run :: sha256 c6c4138ee5c86423 :: prose twin of the run farther from 1
\newcommand{\MlstCmpJcTwoQcdFoneOverFzero}{0.847 and 0.825}  % experiments/FIGS/data/v2/index.json :: linear probe jc2_188 N=1000000: (1 - macro_auc_qcd) of FLAV_F1 over FLAV_F0, run by run :: sha256 c6c4138ee5c86423
\newcommand{\MlstCmpJcTwoQcdFoneOverFzeroPr}{at least 15\% fewer errors}  % experiments/FIGS/data/v2/index.json :: linear probe jc2_188 N=1000000: (1 - macro_auc_qcd) of FLAV_F1 over FLAV_F0, run by run :: sha256 c6c4138ee5c86423 :: prose twin: two significant figures, direction in words
\newcommand{\MlstCmpJcTwoQcdFoneOverFzeroPrMax}{up to 18\% fewer errors}  % experiments/FIGS/data/v2/index.json :: linear probe jc2_188 N=1000000: (1 - macro_auc_qcd) of FLAV_F1 over FLAV_F0, run by run :: sha256 c6c4138ee5c86423 :: prose twin of the run farther from 1
\newcommand{\MlstCmpJcTwoQcdFonerOverFzero}{0.859 and 0.844}  % experiments/FIGS/data/v2/index.json :: linear probe jc2_188 N=1000000: (1 - macro_auc_qcd) of FLAV_F1R over FLAV_F0, run by run :: sha256 c6c4138ee5c86423
\newcommand{\MlstCmpJcTwoQcdFonerOverFzeroPr}{at least 14\% fewer errors}  % experiments/FIGS/data/v2/index.json :: linear probe jc2_188 N=1000000: (1 - macro_auc_qcd) of FLAV_F1R over FLAV_F0, run by run :: sha256 c6c4138ee5c86423 :: prose twin: two significant figures, direction in words
\newcommand{\MlstCmpJcTwoQcdFonerOverFzeroPrMax}{up to 16\% fewer errors}  % experiments/FIGS/data/v2/index.json :: linear probe jc2_188 N=1000000: (1 - macro_auc_qcd) of FLAV_F1R over FLAV_F0, run by run :: sha256 c6c4138ee5c86423 :: prose twin of the run farther from 1
\newcommand{\MlstCmpJcTwoQcdFonerOverFone}{1.015 and 1.023}  % experiments/FIGS/data/v2/index.json :: linear probe jc2_188 N=1000000: (1 - macro_auc_qcd) of FLAV_F1R over FLAV_F1, run by run :: sha256 c6c4138ee5c86423
\newcommand{\MlstCmpJcTwoQcdFonerOverFonePr}{equal within 3\%}  % experiments/FIGS/data/v2/index.json :: linear probe jc2_188 N=1000000: (1 - macro_auc_qcd) of FLAV_F1R over FLAV_F1, run by run :: sha256 c6c4138ee5c86423 :: prose twin: two significant figures, direction in words
\newcommand{\MlstCmpJcTwoQcdFonerOverFonePrMax}{equal within 3\%}  % experiments/FIGS/data/v2/index.json :: linear probe jc2_188 N=1000000: (1 - macro_auc_qcd) of FLAV_F1R over FLAV_F1, run by run :: sha256 c6c4138ee5c86423 :: prose twin of the run farther from 1
\newcommand{\MlstCmpJcTwoQcdMidOverRes}{1.040\,\ensuremath{\pm}\,0.011}  % experiments/FIGS/data/v2/index.json :: linear probe jc2_188 N=1000000: (1 - macro_auc_qcd) of R42_Q1 over L162, run by run :: sha256 c6c4138ee5c86423
\newcommand{\MlstCmpJcTwoQcdMidOverResPr}{4.0\% more errors}  % experiments/FIGS/data/v2/index.json :: linear probe jc2_188 N=1000000: (1 - macro_auc_qcd) of R42_Q1 over L162, run by run :: sha256 c6c4138ee5c86423 :: prose twin: two significant figures, direction in words
\newcommand{\MlstCmpJcTwoQcdCoarseOverRes}{1.348\,\ensuremath{\pm}\,0.010}  % experiments/FIGS/data/v2/index.json :: linear probe jc2_188 N=1000000: (1 - macro_auc_qcd) of R16_Q1 over L162, run by run :: sha256 c6c4138ee5c86423
\newcommand{\MlstCmpJcTwoQcdCoarseOverResPr}{35\% more errors}  % experiments/FIGS/data/v2/index.json :: linear probe jc2_188 N=1000000: (1 - macro_auc_qcd) of R16_Q1 over L162, run by run :: sha256 c6c4138ee5c86423 :: prose twin: two significant figures, direction in words
\newcommand{\MlstCmpJcTwoQcdCoarseOverMid}{1.296\,\ensuremath{\pm}\,0.016}  % experiments/FIGS/data/v2/index.json :: linear probe jc2_188 N=1000000: (1 - macro_auc_qcd) of R16_Q1 over R42_Q1, run by run :: sha256 c6c4138ee5c86423
\newcommand{\MlstCmpJcTwoQcdCoarseOverMidPr}{30\% more errors}  % experiments/FIGS/data/v2/index.json :: linear probe jc2_188 N=1000000: (1 - macro_auc_qcd) of R16_Q1 over R42_Q1, run by run :: sha256 c6c4138ee5c86423 :: prose twin: two significant figures, direction in words
\newcommand{\MlstCmpJcTwoQcdRandOverFull}{1.198\,\ensuremath{\pm}\,0.011}  % experiments/FIGS/data/v2/index.json :: linear probe jc2_188 N=1000000: (1 - macro_auc_qcd) of each random partition over the L188 run of its seed, mean +- SD over the partitions :: sha256 c6c4138ee5c86423
\newcommand{\MlstCmpJcTwoQcdRandOverFullPr}{20\% more errors}  % experiments/FIGS/data/v2/index.json :: linear probe jc2_188 N=1000000: (1 - macro_auc_qcd) of each random partition over the L188 run of its seed, mean +- SD over the partitions :: sha256 c6c4138ee5c86423 :: prose twin: two significant figures, direction in words
\newcommand{\MlstCmpJcTwoQcdRandOverMid}{1.064\,\ensuremath{\pm}\,0.011}  % experiments/FIGS/data/v2/index.json :: linear probe jc2_188 N=1000000: (1 - macro_auc_qcd) of each random partition over the R42_Q1 run of its seed, mean +- SD over the partitions :: sha256 c6c4138ee5c86423
\newcommand{\MlstCmpJcTwoQcdRandOverMidPr}{6.4\% more errors}  % experiments/FIGS/data/v2/index.json :: linear probe jc2_188 N=1000000: (1 - macro_auc_qcd) of each random partition over the R42_Q1 run of its seed, mean +- SD over the partitions :: sha256 c6c4138ee5c86423 :: prose twin: two significant figures, direction in words
\newcommand{\MlstCmpJcTwoQcdRandOverCoarse}{0.818\,\ensuremath{\pm}\,0.013}  % experiments/FIGS/data/v2/index.json :: linear probe jc2_188 N=1000000: (1 - macro_auc_qcd) of each random partition over the R16_Q1 run of its seed, mean +- SD over the partitions :: sha256 c6c4138ee5c86423
\newcommand{\MlstCmpJcTwoQcdRandOverCoarsePr}{18\% fewer errors}  % experiments/FIGS/data/v2/index.json :: linear probe jc2_188 N=1000000: (1 - macro_auc_qcd) of each random partition over the R16_Q1 run of its seed, mean +- SD over the partitions :: sha256 c6c4138ee5c86423 :: prose twin: two significant figures, direction in words
\newcommand{\MlstCmpBcResOverFull}{0.998\,\ensuremath{\pm}\,0.007}  % experiments/FIGS/data/v2/index.json :: binary probe bc_vs_bq_cs: (1 - AUC) of L162 over L188, run by run :: sha256 c6c4138ee5c86423
\newcommand{\MlstCmpBcResOverFullPr}{equal within 1\%}  % experiments/FIGS/data/v2/index.json :: binary probe bc_vs_bq_cs: (1 - AUC) of L162 over L188, run by run :: sha256 c6c4138ee5c86423 :: prose twin: two significant figures, direction in words
\newcommand{\MlstCmpBcMidOverFull}{1.151\,\ensuremath{\pm}\,0.029}  % experiments/FIGS/data/v2/index.json :: binary probe bc_vs_bq_cs: (1 - AUC) of R42_Q1 over L188, run by run :: sha256 c6c4138ee5c86423
\newcommand{\MlstCmpBcMidOverFullPr}{15\% more errors}  % experiments/FIGS/data/v2/index.json :: binary probe bc_vs_bq_cs: (1 - AUC) of R42_Q1 over L188, run by run :: sha256 c6c4138ee5c86423 :: prose twin: two significant figures, direction in words
\newcommand{\MlstCmpBcCoarseOverFull}{1.95\,\ensuremath{\pm}\,0.06}  % experiments/FIGS/data/v2/index.json :: binary probe bc_vs_bq_cs: (1 - AUC) of R16_Q1 over L188, run by run :: sha256 c6c4138ee5c86423
\newcommand{\MlstCmpBcCoarseOverFullPr}{95\% more errors}  % experiments/FIGS/data/v2/index.json :: binary probe bc_vs_bq_cs: (1 - AUC) of R16_Q1 over L188, run by run :: sha256 c6c4138ee5c86423 :: prose twin: two significant figures, direction in words
\newcommand{\MlstCmpBcFoneOverFzero}{0.759 and 0.698}  % experiments/FIGS/data/v2/index.json :: binary probe bc_vs_bq_cs: (1 - AUC) of FLAV_F1 over FLAV_F0, run by run :: sha256 c6c4138ee5c86423
\newcommand{\MlstCmpBcFoneOverFzeroPr}{at least 24\% fewer errors}  % experiments/FIGS/data/v2/index.json :: binary probe bc_vs_bq_cs: (1 - AUC) of FLAV_F1 over FLAV_F0, run by run :: sha256 c6c4138ee5c86423 :: prose twin: two significant figures, direction in words
\newcommand{\MlstCmpBcFoneOverFzeroPrMax}{up to 30\% fewer errors}  % experiments/FIGS/data/v2/index.json :: binary probe bc_vs_bq_cs: (1 - AUC) of FLAV_F1 over FLAV_F0, run by run :: sha256 c6c4138ee5c86423 :: prose twin of the run farther from 1
\newcommand{\MlstCmpBcFonerOverFone}{1.008 and 0.990}  % experiments/FIGS/data/v2/index.json :: binary probe bc_vs_bq_cs: (1 - AUC) of FLAV_F1R over FLAV_F1, run by run :: sha256 c6c4138ee5c86423
\newcommand{\MlstCmpBcFonerOverFonePr}{equal within 1\%}  % experiments/FIGS/data/v2/index.json :: binary probe bc_vs_bq_cs: (1 - AUC) of FLAV_F1R over FLAV_F1, run by run :: sha256 c6c4138ee5c86423 :: prose twin: two significant figures, direction in words
\newcommand{\MlstCmpBcFonerOverFonePrMax}{equal within 1\%}  % experiments/FIGS/data/v2/index.json :: binary probe bc_vs_bq_cs: (1 - AUC) of FLAV_F1R over FLAV_F1, run by run :: sha256 c6c4138ee5c86423 :: prose twin of the run farther from 1
\newcommand{\MlstCmpBcFonerOverFzero}{0.765 and 0.691}  % experiments/FIGS/data/v2/index.json :: binary probe bc_vs_bq_cs: (1 - AUC) of FLAV_F1R over FLAV_F0, run by run :: sha256 c6c4138ee5c86423
\newcommand{\MlstCmpBcFonerOverFzeroPr}{at least 23\% fewer errors}  % experiments/FIGS/data/v2/index.json :: binary probe bc_vs_bq_cs: (1 - AUC) of FLAV_F1R over FLAV_F0, run by run :: sha256 c6c4138ee5c86423 :: prose twin: two significant figures, direction in words
\newcommand{\MlstCmpBcFonerOverFzeroPrMax}{up to 31\% fewer errors}  % experiments/FIGS/data/v2/index.json :: binary probe bc_vs_bq_cs: (1 - AUC) of FLAV_F1R over FLAV_F0, run by run :: sha256 c6c4138ee5c86423 :: prose twin of the run farther from 1
\newcommand{\MlstCmpBcFzeroOverCoarse}{0.99 and 1.13}  % experiments/FIGS/data/v2/index.json :: binary probe bc_vs_bq_cs: (1 - AUC) of FLAV_F0 over R16_Q1, run by run :: sha256 c6c4138ee5c86423
\newcommand{\MlstCmpBcFzeroOverCoarsePr}{equal within 20\%}  % experiments/FIGS/data/v2/index.json :: binary probe bc_vs_bq_cs: (1 - AUC) of FLAV_F0 over R16_Q1, run by run :: sha256 c6c4138ee5c86423 :: prose twin: two significant figures, direction in words
\newcommand{\MlstCmpBcFzeroOverCoarsePrMax}{equal within 20\%}  % experiments/FIGS/data/v2/index.json :: binary probe bc_vs_bq_cs: (1 - AUC) of FLAV_F0 over R16_Q1, run by run :: sha256 c6c4138ee5c86423 :: prose twin of the run farther from 1
\newcommand{\MlstCmpBcFoneOverCoarse}{0.752 and 0.786}  % experiments/FIGS/data/v2/index.json :: binary probe bc_vs_bq_cs: (1 - AUC) of FLAV_F1 over R16_Q1, run by run :: sha256 c6c4138ee5c86423
\newcommand{\MlstCmpBcFoneOverCoarsePr}{at least 21\% fewer errors}  % experiments/FIGS/data/v2/index.json :: binary probe bc_vs_bq_cs: (1 - AUC) of FLAV_F1 over R16_Q1, run by run :: sha256 c6c4138ee5c86423 :: prose twin: two significant figures, direction in words
\newcommand{\MlstCmpBcFoneOverCoarsePrMax}{up to 25\% fewer errors}  % experiments/FIGS/data/v2/index.json :: binary probe bc_vs_bq_cs: (1 - AUC) of FLAV_F1 over R16_Q1, run by run :: sha256 c6c4138ee5c86423 :: prose twin of the run farther from 1
\newcommand{\MlstCmpBcFonerOverCoarse}{0.758 and 0.778}  % experiments/FIGS/data/v2/index.json :: binary probe bc_vs_bq_cs: (1 - AUC) of FLAV_F1R over R16_Q1, run by run :: sha256 c6c4138ee5c86423
\newcommand{\MlstCmpBcFonerOverCoarsePr}{at least 22\% fewer errors}  % experiments/FIGS/data/v2/index.json :: binary probe bc_vs_bq_cs: (1 - AUC) of FLAV_F1R over R16_Q1, run by run :: sha256 c6c4138ee5c86423 :: prose twin: two significant figures, direction in words
\newcommand{\MlstCmpBcFonerOverCoarsePrMax}{up to 24\% fewer errors}  % experiments/FIGS/data/v2/index.json :: binary probe bc_vs_bq_cs: (1 - AUC) of FLAV_F1R over R16_Q1, run by run :: sha256 c6c4138ee5c86423 :: prose twin of the run farther from 1
\newcommand{\MlstCmpBcFzeroOverFull}{2.00 and 2.19}  % experiments/FIGS/data/v2/index.json :: binary probe bc_vs_bq_cs: (1 - AUC) of FLAV_F0 over L188, run by run :: sha256 c6c4138ee5c86423
\newcommand{\MlstCmpBcFzeroOverFullPr}{at least 2.0 times the errors}  % experiments/FIGS/data/v2/index.json :: binary probe bc_vs_bq_cs: (1 - AUC) of FLAV_F0 over L188, run by run :: sha256 c6c4138ee5c86423 :: prose twin: two significant figures, direction in words
\newcommand{\MlstCmpBcFzeroOverFullPrMax}{up to 2.2 times the errors}  % experiments/FIGS/data/v2/index.json :: binary probe bc_vs_bq_cs: (1 - AUC) of FLAV_F0 over L188, run by run :: sha256 c6c4138ee5c86423 :: prose twin of the run farther from 1
\newcommand{\MlstCmpBcFoneOverFull}{1.518 and 1.530}  % experiments/FIGS/data/v2/index.json :: binary probe bc_vs_bq_cs: (1 - AUC) of FLAV_F1 over L188, run by run :: sha256 c6c4138ee5c86423
\newcommand{\MlstCmpBcFoneOverFullPr}{at least 52\% more errors}  % experiments/FIGS/data/v2/index.json :: binary probe bc_vs_bq_cs: (1 - AUC) of FLAV_F1 over L188, run by run :: sha256 c6c4138ee5c86423 :: prose twin: two significant figures, direction in words
\newcommand{\MlstCmpBcFoneOverFullPrMax}{up to 53\% more errors}  % experiments/FIGS/data/v2/index.json :: binary probe bc_vs_bq_cs: (1 - AUC) of FLAV_F1 over L188, run by run :: sha256 c6c4138ee5c86423 :: prose twin of the run farther from 1
\newcommand{\MlstCmpBcFonerOverFull}{1.530 and 1.515}  % experiments/FIGS/data/v2/index.json :: binary probe bc_vs_bq_cs: (1 - AUC) of FLAV_F1R over L188, run by run :: sha256 c6c4138ee5c86423
\newcommand{\MlstCmpBcFonerOverFullPr}{at least 51\% more errors}  % experiments/FIGS/data/v2/index.json :: binary probe bc_vs_bq_cs: (1 - AUC) of FLAV_F1R over L188, run by run :: sha256 c6c4138ee5c86423 :: prose twin: two significant figures, direction in words
\newcommand{\MlstCmpBcFonerOverFullPrMax}{up to 53\% more errors}  % experiments/FIGS/data/v2/index.json :: binary probe bc_vs_bq_cs: (1 - AUC) of FLAV_F1R over L188, run by run :: sha256 c6c4138ee5c86423 :: prose twin of the run farther from 1
\newcommand{\MlstCmpBcRandOverFull}{1.15\,\ensuremath{\pm}\,0.07}  % experiments/FIGS/data/v2/index.json :: binary probe bc_vs_bq_cs: (1 - AUC) of each random partition over the 188-class run of its seed :: sha256 c6c4138ee5c86423
\newcommand{\MlstCmpBcRandOverFullPr}{15\% more errors}  % experiments/FIGS/data/v2/index.json :: binary probe bc_vs_bq_cs: (1 - AUC) of each random partition over the 188-class run of its seed :: sha256 c6c4138ee5c86423 :: prose twin: two significant figures, direction in words
\newcommand{\MlstPairRandMergedBcN}{2}  % experiments/FIGS/data/v2/index.json :: random partitions that merge bc_vs_bq_cs (count) :: sha256 c6c4138ee5c86423
\newcommand{\MlstPairRandApartBcN}{2}  % experiments/FIGS/data/v2/index.json :: random partitions that keep apart bc_vs_bq_cs (count) :: sha256 c6c4138ee5c86423
\newcommand{\MlstPairRandMergedOverApartBc}{1.12}  % experiments/FIGS/data/v2/index.json :: binary probe bc_vs_bq_cs: mean (1 - AUC) of the partitions that merge it over that of those that keep it apart (one run each) :: sha256 c6c4138ee5c86423
\newcommand{\MlstPairRandMergedOverApartBcPr}{12\% more errors}  % experiments/FIGS/data/v2/index.json :: binary probe bc_vs_bq_cs: mean (1 - AUC) of the partitions that merge it over that of those that keep it apart (one run each) :: sha256 c6c4138ee5c86423 :: prose twin: two significant figures, direction in words
\newcommand{\MlstCmpBbResOverFull}{1.015\,\ensuremath{\pm}\,0.027}  % experiments/FIGS/data/v2/index.json :: binary probe bb_vs_cc: (1 - AUC) of L162 over L188, run by run :: sha256 c6c4138ee5c86423
\newcommand{\MlstCmpBbResOverFullPr}{equal within 3\%}  % experiments/FIGS/data/v2/index.json :: binary probe bb_vs_cc: (1 - AUC) of L162 over L188, run by run :: sha256 c6c4138ee5c86423 :: prose twin: two significant figures, direction in words
\newcommand{\MlstCmpBbMidOverFull}{1.102\,\ensuremath{\pm}\,0.013}  % experiments/FIGS/data/v2/index.json :: binary probe bb_vs_cc: (1 - AUC) of R42_Q1 over L188, run by run :: sha256 c6c4138ee5c86423
\newcommand{\MlstCmpBbMidOverFullPr}{10\% more errors}  % experiments/FIGS/data/v2/index.json :: binary probe bb_vs_cc: (1 - AUC) of R42_Q1 over L188, run by run :: sha256 c6c4138ee5c86423 :: prose twin: two significant figures, direction in words
\newcommand{\MlstCmpBbCoarseOverFull}{4.22\,\ensuremath{\pm}\,0.29}  % experiments/FIGS/data/v2/index.json :: binary probe bb_vs_cc: (1 - AUC) of R16_Q1 over L188, run by run :: sha256 c6c4138ee5c86423
\newcommand{\MlstCmpBbCoarseOverFullPr}{4.2 times the errors}  % experiments/FIGS/data/v2/index.json :: binary probe bb_vs_cc: (1 - AUC) of R16_Q1 over L188, run by run :: sha256 c6c4138ee5c86423 :: prose twin: two significant figures, direction in words
\newcommand{\MlstCmpBbFoneOverFzero}{0.31 and 0.13}  % experiments/FIGS/data/v2/index.json :: binary probe bb_vs_cc: (1 - AUC) of FLAV_F1 over FLAV_F0, run by run :: sha256 c6c4138ee5c86423
\newcommand{\MlstCmpBbFoneOverFzeroPr}{at most 0.31 times the errors}  % experiments/FIGS/data/v2/index.json :: binary probe bb_vs_cc: (1 - AUC) of FLAV_F1 over FLAV_F0, run by run :: sha256 c6c4138ee5c86423 :: prose twin: two significant figures, direction in words
\newcommand{\MlstCmpBbFoneOverFzeroPrMax}{down to 0.13 times the errors}  % experiments/FIGS/data/v2/index.json :: binary probe bb_vs_cc: (1 - AUC) of FLAV_F1 over FLAV_F0, run by run :: sha256 c6c4138ee5c86423 :: prose twin of the run farther from 1
\newcommand{\MlstCmpBbFonerOverFone}{1.40 and 1.86}  % experiments/FIGS/data/v2/index.json :: binary probe bb_vs_cc: (1 - AUC) of FLAV_F1R over FLAV_F1, run by run :: sha256 c6c4138ee5c86423
\newcommand{\MlstCmpBbFonerOverFonePr}{at least 40\% more errors}  % experiments/FIGS/data/v2/index.json :: binary probe bb_vs_cc: (1 - AUC) of FLAV_F1R over FLAV_F1, run by run :: sha256 c6c4138ee5c86423 :: prose twin: two significant figures, direction in words
\newcommand{\MlstCmpBbFonerOverFonePrMax}{up to 86\% more errors}  % experiments/FIGS/data/v2/index.json :: binary probe bb_vs_cc: (1 - AUC) of FLAV_F1R over FLAV_F1, run by run :: sha256 c6c4138ee5c86423 :: prose twin of the run farther from 1
\newcommand{\MlstCmpBbFonerOverFzero}{0.44 and 0.24}  % experiments/FIGS/data/v2/index.json :: binary probe bb_vs_cc: (1 - AUC) of FLAV_F1R over FLAV_F0, run by run :: sha256 c6c4138ee5c86423
\newcommand{\MlstCmpBbFonerOverFzeroPr}{at most 0.44 times the errors}  % experiments/FIGS/data/v2/index.json :: binary probe bb_vs_cc: (1 - AUC) of FLAV_F1R over FLAV_F0, run by run :: sha256 c6c4138ee5c86423 :: prose twin: two significant figures, direction in words
\newcommand{\MlstCmpBbFonerOverFzeroPrMax}{down to 0.24 times the errors}  % experiments/FIGS/data/v2/index.json :: binary probe bb_vs_cc: (1 - AUC) of FLAV_F1R over FLAV_F0, run by run :: sha256 c6c4138ee5c86423 :: prose twin of the run farther from 1
\newcommand{\MlstCmpBbFzeroOverCoarse}{1.3 and 3.2}  % experiments/FIGS/data/v2/index.json :: binary probe bb_vs_cc: (1 - AUC) of FLAV_F0 over R16_Q1, run by run :: sha256 c6c4138ee5c86423
\newcommand{\MlstCmpBbFzeroOverCoarsePr}{at least 30\% more errors}  % experiments/FIGS/data/v2/index.json :: binary probe bb_vs_cc: (1 - AUC) of FLAV_F0 over R16_Q1, run by run :: sha256 c6c4138ee5c86423 :: prose twin: two significant figures, direction in words
\newcommand{\MlstCmpBbFzeroOverCoarsePrMax}{up to 3.2 times the errors}  % experiments/FIGS/data/v2/index.json :: binary probe bb_vs_cc: (1 - AUC) of FLAV_F0 over R16_Q1, run by run :: sha256 c6c4138ee5c86423 :: prose twin of the run farther from 1
\newcommand{\MlstCmpBbFoneOverCoarse}{0.4077 and 0.4126}  % experiments/FIGS/data/v2/index.json :: binary probe bb_vs_cc: (1 - AUC) of FLAV_F1 over R16_Q1, run by run :: sha256 c6c4138ee5c86423
\newcommand{\MlstCmpBbFoneOverCoarsePr}{at most 0.41 times the errors}  % experiments/FIGS/data/v2/index.json :: binary probe bb_vs_cc: (1 - AUC) of FLAV_F1 over R16_Q1, run by run :: sha256 c6c4138ee5c86423 :: prose twin: two significant figures, direction in words
\newcommand{\MlstCmpBbFoneOverCoarsePrMax}{down to 0.41 times the errors}  % experiments/FIGS/data/v2/index.json :: binary probe bb_vs_cc: (1 - AUC) of FLAV_F1 over R16_Q1, run by run :: sha256 c6c4138ee5c86423 :: prose twin of the run farther from 1
\newcommand{\MlstCmpBbFonerOverCoarse}{0.57 and 0.77}  % experiments/FIGS/data/v2/index.json :: binary probe bb_vs_cc: (1 - AUC) of FLAV_F1R over R16_Q1, run by run :: sha256 c6c4138ee5c86423
\newcommand{\MlstCmpBbFonerOverCoarsePr}{at least 23\% fewer errors}  % experiments/FIGS/data/v2/index.json :: binary probe bb_vs_cc: (1 - AUC) of FLAV_F1R over R16_Q1, run by run :: sha256 c6c4138ee5c86423 :: prose twin: two significant figures, direction in words
\newcommand{\MlstCmpBbFonerOverCoarsePrMax}{up to 43\% fewer errors}  % experiments/FIGS/data/v2/index.json :: binary probe bb_vs_cc: (1 - AUC) of FLAV_F1R over R16_Q1, run by run :: sha256 c6c4138ee5c86423 :: prose twin of the run farther from 1
\newcommand{\MlstCmpBbFzeroOverFull}{5.9 and 13.4}  % experiments/FIGS/data/v2/index.json :: binary probe bb_vs_cc: (1 - AUC) of FLAV_F0 over L188, run by run :: sha256 c6c4138ee5c86423
\newcommand{\MlstCmpBbFzeroOverFullPr}{at least 5.9 times the errors}  % experiments/FIGS/data/v2/index.json :: binary probe bb_vs_cc: (1 - AUC) of FLAV_F0 over L188, run by run :: sha256 c6c4138ee5c86423 :: prose twin: two significant figures, direction in words
\newcommand{\MlstCmpBbFzeroOverFullPrMax}{up to 13 times the errors}  % experiments/FIGS/data/v2/index.json :: binary probe bb_vs_cc: (1 - AUC) of FLAV_F0 over L188, run by run :: sha256 c6c4138ee5c86423 :: prose twin of the run farther from 1
\newcommand{\MlstCmpBbFoneOverFull}{1.837 and 1.748}  % experiments/FIGS/data/v2/index.json :: binary probe bb_vs_cc: (1 - AUC) of FLAV_F1 over L188, run by run :: sha256 c6c4138ee5c86423
\newcommand{\MlstCmpBbFoneOverFullPr}{at least 75\% more errors}  % experiments/FIGS/data/v2/index.json :: binary probe bb_vs_cc: (1 - AUC) of FLAV_F1 over L188, run by run :: sha256 c6c4138ee5c86423 :: prose twin: two significant figures, direction in words
\newcommand{\MlstCmpBbFoneOverFullPrMax}{up to 84\% more errors}  % experiments/FIGS/data/v2/index.json :: binary probe bb_vs_cc: (1 - AUC) of FLAV_F1 over L188, run by run :: sha256 c6c4138ee5c86423 :: prose twin of the run farther from 1
\newcommand{\MlstCmpBbFonerOverFull}{2.58 and 3.25}  % experiments/FIGS/data/v2/index.json :: binary probe bb_vs_cc: (1 - AUC) of FLAV_F1R over L188, run by run :: sha256 c6c4138ee5c86423
\newcommand{\MlstCmpBbFonerOverFullPr}{at least 2.6 times the errors}  % experiments/FIGS/data/v2/index.json :: binary probe bb_vs_cc: (1 - AUC) of FLAV_F1R over L188, run by run :: sha256 c6c4138ee5c86423 :: prose twin: two significant figures, direction in words
\newcommand{\MlstCmpBbFonerOverFullPrMax}{up to 3.2 times the errors}  % experiments/FIGS/data/v2/index.json :: binary probe bb_vs_cc: (1 - AUC) of FLAV_F1R over L188, run by run :: sha256 c6c4138ee5c86423 :: prose twin of the run farther from 1
\newcommand{\MlstCmpBbRandOverFull}{1.53\,\ensuremath{\pm}\,0.11}  % experiments/FIGS/data/v2/index.json :: binary probe bb_vs_cc: (1 - AUC) of each random partition over the 188-class run of its seed :: sha256 c6c4138ee5c86423
\newcommand{\MlstCmpBbRandOverFullPr}{53\% more errors}  % experiments/FIGS/data/v2/index.json :: binary probe bb_vs_cc: (1 - AUC) of each random partition over the 188-class run of its seed :: sha256 c6c4138ee5c86423 :: prose twin: two significant figures, direction in words
\newcommand{\MlstPairRandMergedBbN}{3}  % experiments/FIGS/data/v2/index.json :: random partitions that merge bb_vs_cc (count) :: sha256 c6c4138ee5c86423
\newcommand{\MlstPairRandApartBbN}{2}  % experiments/FIGS/data/v2/index.json :: random partitions that keep apart bb_vs_cc (count) :: sha256 c6c4138ee5c86423
\newcommand{\MlstPairRandMergedOverApartBb}{1.07}  % experiments/FIGS/data/v2/index.json :: binary probe bb_vs_cc: mean (1 - AUC) of the partitions that merge it over that of those that keep it apart (one run each) :: sha256 c6c4138ee5c86423
\newcommand{\MlstPairRandMergedOverApartBbPr}{7.3\% more errors}  % experiments/FIGS/data/v2/index.json :: binary probe bb_vs_cc: mean (1 - AUC) of the partitions that merge it over that of those that keep it apart (one run each) :: sha256 c6c4138ee5c86423 :: prose twin: two significant figures, direction in words
\newcommand{\MlstCmpBbqqResOverFull}{1.005\,\ensuremath{\pm}\,0.032}  % experiments/FIGS/data/v2/index.json :: binary probe bbqq_vs_ccqq: (1 - AUC) of L162 over L188, run by run :: sha256 c6c4138ee5c86423
\newcommand{\MlstCmpBbqqResOverFullPr}{equal within 4\%}  % experiments/FIGS/data/v2/index.json :: binary probe bbqq_vs_ccqq: (1 - AUC) of L162 over L188, run by run :: sha256 c6c4138ee5c86423 :: prose twin: two significant figures, direction in words
\newcommand{\MlstCmpBbqqMidOverFull}{1.117\,\ensuremath{\pm}\,0.008}  % experiments/FIGS/data/v2/index.json :: binary probe bbqq_vs_ccqq: (1 - AUC) of R42_Q1 over L188, run by run :: sha256 c6c4138ee5c86423
\newcommand{\MlstCmpBbqqMidOverFullPr}{12\% more errors}  % experiments/FIGS/data/v2/index.json :: binary probe bbqq_vs_ccqq: (1 - AUC) of R42_Q1 over L188, run by run :: sha256 c6c4138ee5c86423 :: prose twin: two significant figures, direction in words
\newcommand{\MlstCmpBbqqCoarseOverFull}{4.7\,\ensuremath{\pm}\,0.5}  % experiments/FIGS/data/v2/index.json :: binary probe bbqq_vs_ccqq: (1 - AUC) of R16_Q1 over L188, run by run :: sha256 c6c4138ee5c86423
\newcommand{\MlstCmpBbqqCoarseOverFullPr}{4.7 times the errors}  % experiments/FIGS/data/v2/index.json :: binary probe bbqq_vs_ccqq: (1 - AUC) of R16_Q1 over L188, run by run :: sha256 c6c4138ee5c86423 :: prose twin: two significant figures, direction in words
\newcommand{\MlstCmpBbqqFoneOverFzero}{0.236 and 0.154}  % experiments/FIGS/data/v2/index.json :: binary probe bbqq_vs_ccqq: (1 - AUC) of FLAV_F1 over FLAV_F0, run by run :: sha256 c6c4138ee5c86423
\newcommand{\MlstCmpBbqqFoneOverFzeroPr}{at most 0.24 times the errors}  % experiments/FIGS/data/v2/index.json :: binary probe bbqq_vs_ccqq: (1 - AUC) of FLAV_F1 over FLAV_F0, run by run :: sha256 c6c4138ee5c86423 :: prose twin: two significant figures, direction in words
\newcommand{\MlstCmpBbqqFoneOverFzeroPrMax}{down to 0.15 times the errors}  % experiments/FIGS/data/v2/index.json :: binary probe bbqq_vs_ccqq: (1 - AUC) of FLAV_F1 over FLAV_F0, run by run :: sha256 c6c4138ee5c86423 :: prose twin of the run farther from 1
\newcommand{\MlstCmpBbqqFonerOverFone}{1.505 and 1.577}  % experiments/FIGS/data/v2/index.json :: binary probe bbqq_vs_ccqq: (1 - AUC) of FLAV_F1R over FLAV_F1, run by run :: sha256 c6c4138ee5c86423
\newcommand{\MlstCmpBbqqFonerOverFonePr}{at least 50\% more errors}  % experiments/FIGS/data/v2/index.json :: binary probe bbqq_vs_ccqq: (1 - AUC) of FLAV_F1R over FLAV_F1, run by run :: sha256 c6c4138ee5c86423 :: prose twin: two significant figures, direction in words
\newcommand{\MlstCmpBbqqFonerOverFonePrMax}{up to 58\% more errors}  % experiments/FIGS/data/v2/index.json :: binary probe bbqq_vs_ccqq: (1 - AUC) of FLAV_F1R over FLAV_F1, run by run :: sha256 c6c4138ee5c86423 :: prose twin of the run farther from 1
\newcommand{\MlstCmpBbqqFonerOverFzero}{0.36 and 0.24}  % experiments/FIGS/data/v2/index.json :: binary probe bbqq_vs_ccqq: (1 - AUC) of FLAV_F1R over FLAV_F0, run by run :: sha256 c6c4138ee5c86423
\newcommand{\MlstCmpBbqqFonerOverFzeroPr}{at most 0.36 times the errors}  % experiments/FIGS/data/v2/index.json :: binary probe bbqq_vs_ccqq: (1 - AUC) of FLAV_F1R over FLAV_F0, run by run :: sha256 c6c4138ee5c86423 :: prose twin: two significant figures, direction in words
\newcommand{\MlstCmpBbqqFonerOverFzeroPrMax}{down to 0.24 times the errors}  % experiments/FIGS/data/v2/index.json :: binary probe bbqq_vs_ccqq: (1 - AUC) of FLAV_F1R over FLAV_F0, run by run :: sha256 c6c4138ee5c86423 :: prose twin of the run farther from 1
\newcommand{\MlstCmpBbqqFzeroOverCoarse}{1.51 and 1.79}  % experiments/FIGS/data/v2/index.json :: binary probe bbqq_vs_ccqq: (1 - AUC) of FLAV_F0 over R16_Q1, run by run :: sha256 c6c4138ee5c86423
\newcommand{\MlstCmpBbqqFzeroOverCoarsePr}{at least 51\% more errors}  % experiments/FIGS/data/v2/index.json :: binary probe bbqq_vs_ccqq: (1 - AUC) of FLAV_F0 over R16_Q1, run by run :: sha256 c6c4138ee5c86423 :: prose twin: two significant figures, direction in words
\newcommand{\MlstCmpBbqqFzeroOverCoarsePrMax}{up to 79\% more errors}  % experiments/FIGS/data/v2/index.json :: binary probe bbqq_vs_ccqq: (1 - AUC) of FLAV_F0 over R16_Q1, run by run :: sha256 c6c4138ee5c86423 :: prose twin of the run farther from 1
\newcommand{\MlstCmpBbqqFoneOverCoarse}{0.356 and 0.276}  % experiments/FIGS/data/v2/index.json :: binary probe bbqq_vs_ccqq: (1 - AUC) of FLAV_F1 over R16_Q1, run by run :: sha256 c6c4138ee5c86423
\newcommand{\MlstCmpBbqqFoneOverCoarsePr}{at most 0.36 times the errors}  % experiments/FIGS/data/v2/index.json :: binary probe bbqq_vs_ccqq: (1 - AUC) of FLAV_F1 over R16_Q1, run by run :: sha256 c6c4138ee5c86423 :: prose twin: two significant figures, direction in words
\newcommand{\MlstCmpBbqqFoneOverCoarsePrMax}{down to 0.28 times the errors}  % experiments/FIGS/data/v2/index.json :: binary probe bbqq_vs_ccqq: (1 - AUC) of FLAV_F1 over R16_Q1, run by run :: sha256 c6c4138ee5c86423 :: prose twin of the run farther from 1
\newcommand{\MlstCmpBbqqFonerOverCoarse}{0.54 and 0.44}  % experiments/FIGS/data/v2/index.json :: binary probe bbqq_vs_ccqq: (1 - AUC) of FLAV_F1R over R16_Q1, run by run :: sha256 c6c4138ee5c86423
\newcommand{\MlstCmpBbqqFonerOverCoarsePr}{at least 46\% fewer errors}  % experiments/FIGS/data/v2/index.json :: binary probe bbqq_vs_ccqq: (1 - AUC) of FLAV_F1R over R16_Q1, run by run :: sha256 c6c4138ee5c86423 :: prose twin: two significant figures, direction in words
\newcommand{\MlstCmpBbqqFonerOverCoarsePrMax}{down to 0.44 times the errors}  % experiments/FIGS/data/v2/index.json :: binary probe bbqq_vs_ccqq: (1 - AUC) of FLAV_F1R over R16_Q1, run by run :: sha256 c6c4138ee5c86423 :: prose twin of the run farther from 1
\newcommand{\MlstCmpBbqqFzeroOverFull}{6.4 and 9.4}  % experiments/FIGS/data/v2/index.json :: binary probe bbqq_vs_ccqq: (1 - AUC) of FLAV_F0 over L188, run by run :: sha256 c6c4138ee5c86423
\newcommand{\MlstCmpBbqqFzeroOverFullPr}{at least 6.4 times the errors}  % experiments/FIGS/data/v2/index.json :: binary probe bbqq_vs_ccqq: (1 - AUC) of FLAV_F0 over L188, run by run :: sha256 c6c4138ee5c86423 :: prose twin: two significant figures, direction in words
\newcommand{\MlstCmpBbqqFzeroOverFullPrMax}{up to 9.4 times the errors}  % experiments/FIGS/data/v2/index.json :: binary probe bbqq_vs_ccqq: (1 - AUC) of FLAV_F0 over L188, run by run :: sha256 c6c4138ee5c86423 :: prose twin of the run farther from 1
\newcommand{\MlstCmpBbqqFoneOverFull}{1.510 and 1.447}  % experiments/FIGS/data/v2/index.json :: binary probe bbqq_vs_ccqq: (1 - AUC) of FLAV_F1 over L188, run by run :: sha256 c6c4138ee5c86423
\newcommand{\MlstCmpBbqqFoneOverFullPr}{at least 45\% more errors}  % experiments/FIGS/data/v2/index.json :: binary probe bbqq_vs_ccqq: (1 - AUC) of FLAV_F1 over L188, run by run :: sha256 c6c4138ee5c86423 :: prose twin: two significant figures, direction in words
\newcommand{\MlstCmpBbqqFoneOverFullPrMax}{up to 51\% more errors}  % experiments/FIGS/data/v2/index.json :: binary probe bbqq_vs_ccqq: (1 - AUC) of FLAV_F1 over L188, run by run :: sha256 c6c4138ee5c86423 :: prose twin of the run farther from 1
\newcommand{\MlstCmpBbqqFonerOverFull}{2.272 and 2.281}  % experiments/FIGS/data/v2/index.json :: binary probe bbqq_vs_ccqq: (1 - AUC) of FLAV_F1R over L188, run by run :: sha256 c6c4138ee5c86423
\newcommand{\MlstCmpBbqqFonerOverFullPr}{at least 2.3 times the errors}  % experiments/FIGS/data/v2/index.json :: binary probe bbqq_vs_ccqq: (1 - AUC) of FLAV_F1R over L188, run by run :: sha256 c6c4138ee5c86423 :: prose twin: two significant figures, direction in words
\newcommand{\MlstCmpBbqqFonerOverFullPrMax}{up to 2.3 times the errors}  % experiments/FIGS/data/v2/index.json :: binary probe bbqq_vs_ccqq: (1 - AUC) of FLAV_F1R over L188, run by run :: sha256 c6c4138ee5c86423 :: prose twin of the run farther from 1
\newcommand{\MlstCmpBbqqRandOverFull}{1.42\,\ensuremath{\pm}\,0.06}  % experiments/FIGS/data/v2/index.json :: binary probe bbqq_vs_ccqq: (1 - AUC) of each random partition over the 188-class run of its seed :: sha256 c6c4138ee5c86423
\newcommand{\MlstCmpBbqqRandOverFullPr}{42\% more errors}  % experiments/FIGS/data/v2/index.json :: binary probe bbqq_vs_ccqq: (1 - AUC) of each random partition over the 188-class run of its seed :: sha256 c6c4138ee5c86423 :: prose twin: two significant figures, direction in words
\newcommand{\MlstPairRandMergedBbqqN}{2}  % experiments/FIGS/data/v2/index.json :: random partitions that merge bbqq_vs_ccqq (count) :: sha256 c6c4138ee5c86423
\newcommand{\MlstPairRandApartBbqqN}{3}  % experiments/FIGS/data/v2/index.json :: random partitions that keep apart bbqq_vs_ccqq (count) :: sha256 c6c4138ee5c86423
\newcommand{\MlstPairRandMergedOverApartBbqq}{0.996}  % experiments/FIGS/data/v2/index.json :: binary probe bbqq_vs_ccqq: mean (1 - AUC) of the partitions that merge it over that of those that keep it apart (one run each) :: sha256 c6c4138ee5c86423
\newcommand{\MlstPairRandMergedOverApartBbqqPr}{0.44\% fewer errors}  % experiments/FIGS/data/v2/index.json :: binary probe bbqq_vs_ccqq: mean (1 - AUC) of the partitions that merge it over that of those that keep it apart (one run each) :: sha256 c6c4138ee5c86423 :: prose twin: two significant figures, direction in words
\newcommand{\MlstCmpEemmResOverFull}{--}  % experiments/FIGS/data/v2/index.json :: binary probe ee_vs_mm: (1 - AUC) of L162 over L188, run by run :: sha256 c6c4138ee5c86423 :: the reference's 1 - AUC is 0
\newcommand{\MlstCmpEemmResOverFullPr}{undefined, since the reference separates the test jets completely}  % experiments/FIGS/data/v2/index.json :: binary probe ee_vs_mm: (1 - AUC) of L162 over L188, run by run :: sha256 c6c4138ee5c86423 :: the reference's 1 - AUC is 0
\newcommand{\MlstCmpEemmMidOverFull}{--}  % experiments/FIGS/data/v2/index.json :: binary probe ee_vs_mm: (1 - AUC) of R42_Q1 over L188, run by run :: sha256 c6c4138ee5c86423 :: the reference's 1 - AUC is 0
\newcommand{\MlstCmpEemmMidOverFullPr}{undefined, since the reference separates the test jets completely}  % experiments/FIGS/data/v2/index.json :: binary probe ee_vs_mm: (1 - AUC) of R42_Q1 over L188, run by run :: sha256 c6c4138ee5c86423 :: the reference's 1 - AUC is 0
\newcommand{\MlstCmpEemmCoarseOverFull}{--}  % experiments/FIGS/data/v2/index.json :: binary probe ee_vs_mm: (1 - AUC) of R16_Q1 over L188, run by run :: sha256 c6c4138ee5c86423 :: the reference's 1 - AUC is 0
\newcommand{\MlstCmpEemmCoarseOverFullPr}{undefined, since the reference separates the test jets completely}  % experiments/FIGS/data/v2/index.json :: binary probe ee_vs_mm: (1 - AUC) of R16_Q1 over L188, run by run :: sha256 c6c4138ee5c86423 :: the reference's 1 - AUC is 0
\newcommand{\MlstCmpEemmFoneOverFzero}{--}  % experiments/FIGS/data/v2/index.json :: binary probe ee_vs_mm: (1 - AUC) of FLAV_F1 over FLAV_F0, run by run :: sha256 c6c4138ee5c86423 :: the reference's 1 - AUC is 0
\newcommand{\MlstCmpEemmFoneOverFzeroPr}{undefined, since the reference separates the test jets completely}  % experiments/FIGS/data/v2/index.json :: binary probe ee_vs_mm: (1 - AUC) of FLAV_F1 over FLAV_F0, run by run :: sha256 c6c4138ee5c86423 :: the reference's 1 - AUC is 0
\newcommand{\MlstCmpEemmFoneOverFzeroPrMax}{undefined, since the reference separates the test jets completely}  % experiments/FIGS/data/v2/index.json :: binary probe ee_vs_mm: (1 - AUC) of FLAV_F1 over FLAV_F0, run by run :: sha256 c6c4138ee5c86423 :: the reference's 1 - AUC is 0
\newcommand{\MlstCmpEemmFonerOverFone}{--}  % experiments/FIGS/data/v2/index.json :: binary probe ee_vs_mm: (1 - AUC) of FLAV_F1R over FLAV_F1, run by run :: sha256 c6c4138ee5c86423 :: the reference's 1 - AUC is 0
\newcommand{\MlstCmpEemmFonerOverFonePr}{undefined, since the reference separates the test jets completely}  % experiments/FIGS/data/v2/index.json :: binary probe ee_vs_mm: (1 - AUC) of FLAV_F1R over FLAV_F1, run by run :: sha256 c6c4138ee5c86423 :: the reference's 1 - AUC is 0
\newcommand{\MlstCmpEemmFonerOverFonePrMax}{undefined, since the reference separates the test jets completely}  % experiments/FIGS/data/v2/index.json :: binary probe ee_vs_mm: (1 - AUC) of FLAV_F1R over FLAV_F1, run by run :: sha256 c6c4138ee5c86423 :: the reference's 1 - AUC is 0
\newcommand{\MlstCmpEemmFonerOverFzero}{--}  % experiments/FIGS/data/v2/index.json :: binary probe ee_vs_mm: (1 - AUC) of FLAV_F1R over FLAV_F0, run by run :: sha256 c6c4138ee5c86423 :: the reference's 1 - AUC is 0
\newcommand{\MlstCmpEemmFonerOverFzeroPr}{undefined, since the reference separates the test jets completely}  % experiments/FIGS/data/v2/index.json :: binary probe ee_vs_mm: (1 - AUC) of FLAV_F1R over FLAV_F0, run by run :: sha256 c6c4138ee5c86423 :: the reference's 1 - AUC is 0
\newcommand{\MlstCmpEemmFonerOverFzeroPrMax}{undefined, since the reference separates the test jets completely}  % experiments/FIGS/data/v2/index.json :: binary probe ee_vs_mm: (1 - AUC) of FLAV_F1R over FLAV_F0, run by run :: sha256 c6c4138ee5c86423 :: the reference's 1 - AUC is 0
\newcommand{\MlstCmpEemmFzeroOverCoarse}{0.0000022 and 0.0000003}  % experiments/FIGS/data/v2/index.json :: binary probe ee_vs_mm: (1 - AUC) of FLAV_F0 over R16_Q1, run by run :: sha256 c6c4138ee5c86423
\newcommand{\MlstCmpEemmFzeroOverCoarsePr}{at most 0.0000022 times the errors}  % experiments/FIGS/data/v2/index.json :: binary probe ee_vs_mm: (1 - AUC) of FLAV_F0 over R16_Q1, run by run :: sha256 c6c4138ee5c86423 :: prose twin: two significant figures, direction in words
\newcommand{\MlstCmpEemmFzeroOverCoarsePrMax}{down to 0.00000034 times the errors}  % experiments/FIGS/data/v2/index.json :: binary probe ee_vs_mm: (1 - AUC) of FLAV_F0 over R16_Q1, run by run :: sha256 c6c4138ee5c86423 :: prose twin of the run farther from 1
\newcommand{\MlstCmpEemmFoneOverCoarse}{0.000000 and 0.000038}  % experiments/FIGS/data/v2/index.json :: binary probe ee_vs_mm: (1 - AUC) of FLAV_F1 over R16_Q1, run by run :: sha256 c6c4138ee5c86423
\newcommand{\MlstCmpEemmFoneOverCoarsePr}{at most 0.000038 times the errors}  % experiments/FIGS/data/v2/index.json :: binary probe ee_vs_mm: (1 - AUC) of FLAV_F1 over R16_Q1, run by run :: sha256 c6c4138ee5c86423 :: prose twin: two significant figures, direction in words
\newcommand{\MlstCmpEemmFoneOverCoarsePrMax}{down to 0.000000042 times the errors}  % experiments/FIGS/data/v2/index.json :: binary probe ee_vs_mm: (1 - AUC) of FLAV_F1 over R16_Q1, run by run :: sha256 c6c4138ee5c86423 :: prose twin of the run farther from 1
\newcommand{\MlstCmpEemmFonerOverCoarse}{0.00000005 and 0.00000017}  % experiments/FIGS/data/v2/index.json :: binary probe ee_vs_mm: (1 - AUC) of FLAV_F1R over R16_Q1, run by run :: sha256 c6c4138ee5c86423
\newcommand{\MlstCmpEemmFonerOverCoarsePr}{at most 0.00000017 times the errors}  % experiments/FIGS/data/v2/index.json :: binary probe ee_vs_mm: (1 - AUC) of FLAV_F1R over R16_Q1, run by run :: sha256 c6c4138ee5c86423 :: prose twin: two significant figures, direction in words
\newcommand{\MlstCmpEemmFonerOverCoarsePrMax}{down to 0.000000055 times the errors}  % experiments/FIGS/data/v2/index.json :: binary probe ee_vs_mm: (1 - AUC) of FLAV_F1R over R16_Q1, run by run :: sha256 c6c4138ee5c86423 :: prose twin of the run farther from 1
\newcommand{\MlstCmpEemmFzeroOverFull}{--}  % experiments/FIGS/data/v2/index.json :: binary probe ee_vs_mm: (1 - AUC) of FLAV_F0 over L188, run by run :: sha256 c6c4138ee5c86423 :: the reference's 1 - AUC is 0
\newcommand{\MlstCmpEemmFzeroOverFullPr}{undefined, since the reference separates the test jets completely}  % experiments/FIGS/data/v2/index.json :: binary probe ee_vs_mm: (1 - AUC) of FLAV_F0 over L188, run by run :: sha256 c6c4138ee5c86423 :: the reference's 1 - AUC is 0
\newcommand{\MlstCmpEemmFzeroOverFullPrMax}{undefined, since the reference separates the test jets completely}  % experiments/FIGS/data/v2/index.json :: binary probe ee_vs_mm: (1 - AUC) of FLAV_F0 over L188, run by run :: sha256 c6c4138ee5c86423 :: the reference's 1 - AUC is 0
\newcommand{\MlstCmpEemmFoneOverFull}{--}  % experiments/FIGS/data/v2/index.json :: binary probe ee_vs_mm: (1 - AUC) of FLAV_F1 over L188, run by run :: sha256 c6c4138ee5c86423 :: the reference's 1 - AUC is 0
\newcommand{\MlstCmpEemmFoneOverFullPr}{undefined, since the reference separates the test jets completely}  % experiments/FIGS/data/v2/index.json :: binary probe ee_vs_mm: (1 - AUC) of FLAV_F1 over L188, run by run :: sha256 c6c4138ee5c86423 :: the reference's 1 - AUC is 0
\newcommand{\MlstCmpEemmFoneOverFullPrMax}{undefined, since the reference separates the test jets completely}  % experiments/FIGS/data/v2/index.json :: binary probe ee_vs_mm: (1 - AUC) of FLAV_F1 over L188, run by run :: sha256 c6c4138ee5c86423 :: the reference's 1 - AUC is 0
\newcommand{\MlstCmpEemmFonerOverFull}{--}  % experiments/FIGS/data/v2/index.json :: binary probe ee_vs_mm: (1 - AUC) of FLAV_F1R over L188, run by run :: sha256 c6c4138ee5c86423 :: the reference's 1 - AUC is 0
\newcommand{\MlstCmpEemmFonerOverFullPr}{undefined, since the reference separates the test jets completely}  % experiments/FIGS/data/v2/index.json :: binary probe ee_vs_mm: (1 - AUC) of FLAV_F1R over L188, run by run :: sha256 c6c4138ee5c86423 :: the reference's 1 - AUC is 0
\newcommand{\MlstCmpEemmFonerOverFullPrMax}{undefined, since the reference separates the test jets completely}  % experiments/FIGS/data/v2/index.json :: binary probe ee_vs_mm: (1 - AUC) of FLAV_F1R over L188, run by run :: sha256 c6c4138ee5c86423 :: the reference's 1 - AUC is 0
\newcommand{\MlstCmpEemmRandOverFull}{--}  % experiments/FIGS/data/v2/index.json :: binary probe ee_vs_mm: (1 - AUC) of each random partition over the 188-class run of its seed :: sha256 c6c4138ee5c86423 :: the reference's 1 - AUC is 0
\newcommand{\MlstCmpEemmRandOverFullPr}{undefined, since the reference separates the test jets completely}  % experiments/FIGS/data/v2/index.json :: binary probe ee_vs_mm: (1 - AUC) of each random partition over the 188-class run of its seed :: sha256 c6c4138ee5c86423 :: the reference's 1 - AUC is 0
\newcommand{\MlstPairRandMergedEemmN}{3}  % experiments/FIGS/data/v2/index.json :: random partitions that merge ee_vs_mm (count) :: sha256 c6c4138ee5c86423
\newcommand{\MlstPairRandApartEemmN}{2}  % experiments/FIGS/data/v2/index.json :: random partitions that keep apart ee_vs_mm (count) :: sha256 c6c4138ee5c86423
\newcommand{\MlstPairRandMergedOverApartEemm}{75.3}  % experiments/FIGS/data/v2/index.json :: binary probe ee_vs_mm: mean (1 - AUC) of the partitions that merge it over that of those that keep it apart (one run each) :: sha256 c6c4138ee5c86423
\newcommand{\MlstPairRandMergedOverApartEemmPr}{75 times the errors}  % experiments/FIGS/data/v2/index.json :: binary probe ee_vs_mm: mean (1 - AUC) of the partitions that merge it over that of those that keep it apart (one run each) :: sha256 c6c4138ee5c86423 :: prose twin: two significant figures, direction in words
\newcommand{\MlstCmpQcdBbResOverFull}{1.15\,\ensuremath{\pm}\,0.04}  % experiments/FIGS/data/v2/index.json :: binary probe qcd_bb_vs_cc: (1 - AUC) of L162 over L188, run by run :: sha256 c6c4138ee5c86423
\newcommand{\MlstCmpQcdBbResOverFullPr}{15\% more errors}  % experiments/FIGS/data/v2/index.json :: binary probe qcd_bb_vs_cc: (1 - AUC) of L162 over L188, run by run :: sha256 c6c4138ee5c86423 :: prose twin: two significant figures, direction in words
\newcommand{\MlstCmpQcdBbMidOverFull}{1.163\,\ensuremath{\pm}\,0.016}  % experiments/FIGS/data/v2/index.json :: binary probe qcd_bb_vs_cc: (1 - AUC) of R42_Q1 over L188, run by run :: sha256 c6c4138ee5c86423
\newcommand{\MlstCmpQcdBbMidOverFullPr}{16\% more errors}  % experiments/FIGS/data/v2/index.json :: binary probe qcd_bb_vs_cc: (1 - AUC) of R42_Q1 over L188, run by run :: sha256 c6c4138ee5c86423 :: prose twin: two significant figures, direction in words
\newcommand{\MlstCmpQcdBbCoarseOverFull}{3.294\,\ensuremath{\pm}\,0.033}  % experiments/FIGS/data/v2/index.json :: binary probe qcd_bb_vs_cc: (1 - AUC) of R16_Q1 over L188, run by run :: sha256 c6c4138ee5c86423
\newcommand{\MlstCmpQcdBbCoarseOverFullPr}{3.3 times the errors}  % experiments/FIGS/data/v2/index.json :: binary probe qcd_bb_vs_cc: (1 - AUC) of R16_Q1 over L188, run by run :: sha256 c6c4138ee5c86423 :: prose twin: two significant figures, direction in words
\newcommand{\MlstCmpQcdBbFoneOverFzero}{0.368 and 0.294}  % experiments/FIGS/data/v2/index.json :: binary probe qcd_bb_vs_cc: (1 - AUC) of FLAV_F1 over FLAV_F0, run by run :: sha256 c6c4138ee5c86423
\newcommand{\MlstCmpQcdBbFoneOverFzeroPr}{at most 0.37 times the errors}  % experiments/FIGS/data/v2/index.json :: binary probe qcd_bb_vs_cc: (1 - AUC) of FLAV_F1 over FLAV_F0, run by run :: sha256 c6c4138ee5c86423 :: prose twin: two significant figures, direction in words
\newcommand{\MlstCmpQcdBbFoneOverFzeroPrMax}{down to 0.29 times the errors}  % experiments/FIGS/data/v2/index.json :: binary probe qcd_bb_vs_cc: (1 - AUC) of FLAV_F1 over FLAV_F0, run by run :: sha256 c6c4138ee5c86423 :: prose twin of the run farther from 1
\newcommand{\MlstCmpQcdBbFonerOverFone}{1.361 and 1.429}  % experiments/FIGS/data/v2/index.json :: binary probe qcd_bb_vs_cc: (1 - AUC) of FLAV_F1R over FLAV_F1, run by run :: sha256 c6c4138ee5c86423
\newcommand{\MlstCmpQcdBbFonerOverFonePr}{at least 36\% more errors}  % experiments/FIGS/data/v2/index.json :: binary probe qcd_bb_vs_cc: (1 - AUC) of FLAV_F1R over FLAV_F1, run by run :: sha256 c6c4138ee5c86423 :: prose twin: two significant figures, direction in words
\newcommand{\MlstCmpQcdBbFonerOverFonePrMax}{up to 43\% more errors}  % experiments/FIGS/data/v2/index.json :: binary probe qcd_bb_vs_cc: (1 - AUC) of FLAV_F1R over FLAV_F1, run by run :: sha256 c6c4138ee5c86423 :: prose twin of the run farther from 1
\newcommand{\MlstCmpQcdBbFonerOverFzero}{0.501 and 0.420}  % experiments/FIGS/data/v2/index.json :: binary probe qcd_bb_vs_cc: (1 - AUC) of FLAV_F1R over FLAV_F0, run by run :: sha256 c6c4138ee5c86423
\newcommand{\MlstCmpQcdBbFonerOverFzeroPr}{at least 50\% fewer errors}  % experiments/FIGS/data/v2/index.json :: binary probe qcd_bb_vs_cc: (1 - AUC) of FLAV_F1R over FLAV_F0, run by run :: sha256 c6c4138ee5c86423 :: prose twin: two significant figures, direction in words
\newcommand{\MlstCmpQcdBbFonerOverFzeroPrMax}{down to 0.42 times the errors}  % experiments/FIGS/data/v2/index.json :: binary probe qcd_bb_vs_cc: (1 - AUC) of FLAV_F1R over FLAV_F0, run by run :: sha256 c6c4138ee5c86423 :: prose twin of the run farther from 1
\newcommand{\MlstCmpQcdBbFzeroOverCoarse}{1.28 and 1.58}  % experiments/FIGS/data/v2/index.json :: binary probe qcd_bb_vs_cc: (1 - AUC) of FLAV_F0 over R16_Q1, run by run :: sha256 c6c4138ee5c86423
\newcommand{\MlstCmpQcdBbFzeroOverCoarsePr}{at least 28\% more errors}  % experiments/FIGS/data/v2/index.json :: binary probe qcd_bb_vs_cc: (1 - AUC) of FLAV_F0 over R16_Q1, run by run :: sha256 c6c4138ee5c86423 :: prose twin: two significant figures, direction in words
\newcommand{\MlstCmpQcdBbFzeroOverCoarsePrMax}{up to 58\% more errors}  % experiments/FIGS/data/v2/index.json :: binary probe qcd_bb_vs_cc: (1 - AUC) of FLAV_F0 over R16_Q1, run by run :: sha256 c6c4138ee5c86423 :: prose twin of the run farther from 1
\newcommand{\MlstCmpQcdBbFoneOverCoarse}{0.4692 and 0.4648}  % experiments/FIGS/data/v2/index.json :: binary probe qcd_bb_vs_cc: (1 - AUC) of FLAV_F1 over R16_Q1, run by run :: sha256 c6c4138ee5c86423
\newcommand{\MlstCmpQcdBbFoneOverCoarsePr}{at most 0.47 times the errors}  % experiments/FIGS/data/v2/index.json :: binary probe qcd_bb_vs_cc: (1 - AUC) of FLAV_F1 over R16_Q1, run by run :: sha256 c6c4138ee5c86423 :: prose twin: two significant figures, direction in words
\newcommand{\MlstCmpQcdBbFoneOverCoarsePrMax}{down to 0.46 times the errors}  % experiments/FIGS/data/v2/index.json :: binary probe qcd_bb_vs_cc: (1 - AUC) of FLAV_F1 over R16_Q1, run by run :: sha256 c6c4138ee5c86423 :: prose twin of the run farther from 1
\newcommand{\MlstCmpQcdBbFonerOverCoarse}{0.639 and 0.664}  % experiments/FIGS/data/v2/index.json :: binary probe qcd_bb_vs_cc: (1 - AUC) of FLAV_F1R over R16_Q1, run by run :: sha256 c6c4138ee5c86423
\newcommand{\MlstCmpQcdBbFonerOverCoarsePr}{at least 34\% fewer errors}  % experiments/FIGS/data/v2/index.json :: binary probe qcd_bb_vs_cc: (1 - AUC) of FLAV_F1R over R16_Q1, run by run :: sha256 c6c4138ee5c86423 :: prose twin: two significant figures, direction in words
\newcommand{\MlstCmpQcdBbFonerOverCoarsePrMax}{up to 36\% fewer errors}  % experiments/FIGS/data/v2/index.json :: binary probe qcd_bb_vs_cc: (1 - AUC) of FLAV_F1R over R16_Q1, run by run :: sha256 c6c4138ee5c86423 :: prose twin of the run farther from 1
\newcommand{\MlstCmpQcdBbFzeroOverFull}{4.24 and 5.21}  % experiments/FIGS/data/v2/index.json :: binary probe qcd_bb_vs_cc: (1 - AUC) of FLAV_F0 over L188, run by run :: sha256 c6c4138ee5c86423
\newcommand{\MlstCmpQcdBbFzeroOverFullPr}{at least 4.2 times the errors}  % experiments/FIGS/data/v2/index.json :: binary probe qcd_bb_vs_cc: (1 - AUC) of FLAV_F0 over L188, run by run :: sha256 c6c4138ee5c86423 :: prose twin: two significant figures, direction in words
\newcommand{\MlstCmpQcdBbFzeroOverFullPrMax}{up to 5.2 times the errors}  % experiments/FIGS/data/v2/index.json :: binary probe qcd_bb_vs_cc: (1 - AUC) of FLAV_F0 over L188, run by run :: sha256 c6c4138ee5c86423 :: prose twin of the run farther from 1
\newcommand{\MlstCmpQcdBbFoneOverFull}{1.561 and 1.532}  % experiments/FIGS/data/v2/index.json :: binary probe qcd_bb_vs_cc: (1 - AUC) of FLAV_F1 over L188, run by run :: sha256 c6c4138ee5c86423
\newcommand{\MlstCmpQcdBbFoneOverFullPr}{at least 53\% more errors}  % experiments/FIGS/data/v2/index.json :: binary probe qcd_bb_vs_cc: (1 - AUC) of FLAV_F1 over L188, run by run :: sha256 c6c4138ee5c86423 :: prose twin: two significant figures, direction in words
\newcommand{\MlstCmpQcdBbFoneOverFullPrMax}{up to 56\% more errors}  % experiments/FIGS/data/v2/index.json :: binary probe qcd_bb_vs_cc: (1 - AUC) of FLAV_F1 over L188, run by run :: sha256 c6c4138ee5c86423 :: prose twin of the run farther from 1
\newcommand{\MlstCmpQcdBbFonerOverFull}{2.124 and 2.189}  % experiments/FIGS/data/v2/index.json :: binary probe qcd_bb_vs_cc: (1 - AUC) of FLAV_F1R over L188, run by run :: sha256 c6c4138ee5c86423
\newcommand{\MlstCmpQcdBbFonerOverFullPr}{at least 2.1 times the errors}  % experiments/FIGS/data/v2/index.json :: binary probe qcd_bb_vs_cc: (1 - AUC) of FLAV_F1R over L188, run by run :: sha256 c6c4138ee5c86423 :: prose twin: two significant figures, direction in words
\newcommand{\MlstCmpQcdBbFonerOverFullPrMax}{up to 2.2 times the errors}  % experiments/FIGS/data/v2/index.json :: binary probe qcd_bb_vs_cc: (1 - AUC) of FLAV_F1R over L188, run by run :: sha256 c6c4138ee5c86423 :: prose twin of the run farther from 1
\newcommand{\MlstCmpQcdBResOverFull}{1.085\,\ensuremath{\pm}\,0.004}  % experiments/FIGS/data/v2/index.json :: binary probe qcd_b_vs_light: (1 - AUC) of L162 over L188, run by run :: sha256 c6c4138ee5c86423
\newcommand{\MlstCmpQcdBResOverFullPr}{8.5\% more errors}  % experiments/FIGS/data/v2/index.json :: binary probe qcd_b_vs_light: (1 - AUC) of L162 over L188, run by run :: sha256 c6c4138ee5c86423 :: prose twin: two significant figures, direction in words
\newcommand{\MlstCmpQcdBMidOverFull}{1.081\,\ensuremath{\pm}\,0.008}  % experiments/FIGS/data/v2/index.json :: binary probe qcd_b_vs_light: (1 - AUC) of R42_Q1 over L188, run by run :: sha256 c6c4138ee5c86423
\newcommand{\MlstCmpQcdBMidOverFullPr}{8.1\% more errors}  % experiments/FIGS/data/v2/index.json :: binary probe qcd_b_vs_light: (1 - AUC) of R42_Q1 over L188, run by run :: sha256 c6c4138ee5c86423 :: prose twin: two significant figures, direction in words
\newcommand{\MlstCmpQcdBCoarseOverFull}{1.54\,\ensuremath{\pm}\,0.07}  % experiments/FIGS/data/v2/index.json :: binary probe qcd_b_vs_light: (1 - AUC) of R16_Q1 over L188, run by run :: sha256 c6c4138ee5c86423
\newcommand{\MlstCmpQcdBCoarseOverFullPr}{54\% more errors}  % experiments/FIGS/data/v2/index.json :: binary probe qcd_b_vs_light: (1 - AUC) of R16_Q1 over L188, run by run :: sha256 c6c4138ee5c86423 :: prose twin: two significant figures, direction in words
\newcommand{\MlstCmpQcdBFoneOverFzero}{0.778 and 0.835}  % experiments/FIGS/data/v2/index.json :: binary probe qcd_b_vs_light: (1 - AUC) of FLAV_F1 over FLAV_F0, run by run :: sha256 c6c4138ee5c86423
\newcommand{\MlstCmpQcdBFoneOverFzeroPr}{at least 17\% fewer errors}  % experiments/FIGS/data/v2/index.json :: binary probe qcd_b_vs_light: (1 - AUC) of FLAV_F1 over FLAV_F0, run by run :: sha256 c6c4138ee5c86423 :: prose twin: two significant figures, direction in words
\newcommand{\MlstCmpQcdBFoneOverFzeroPrMax}{up to 22\% fewer errors}  % experiments/FIGS/data/v2/index.json :: binary probe qcd_b_vs_light: (1 - AUC) of FLAV_F1 over FLAV_F0, run by run :: sha256 c6c4138ee5c86423 :: prose twin of the run farther from 1
\newcommand{\MlstCmpQcdBFonerOverFone}{1.054 and 1.061}  % experiments/FIGS/data/v2/index.json :: binary probe qcd_b_vs_light: (1 - AUC) of FLAV_F1R over FLAV_F1, run by run :: sha256 c6c4138ee5c86423
\newcommand{\MlstCmpQcdBFonerOverFonePr}{equal within 7\%}  % experiments/FIGS/data/v2/index.json :: binary probe qcd_b_vs_light: (1 - AUC) of FLAV_F1R over FLAV_F1, run by run :: sha256 c6c4138ee5c86423 :: prose twin: two significant figures, direction in words
\newcommand{\MlstCmpQcdBFonerOverFonePrMax}{equal within 7\%}  % experiments/FIGS/data/v2/index.json :: binary probe qcd_b_vs_light: (1 - AUC) of FLAV_F1R over FLAV_F1, run by run :: sha256 c6c4138ee5c86423 :: prose twin of the run farther from 1
\newcommand{\MlstCmpQcdBFonerOverFzero}{0.820 and 0.885}  % experiments/FIGS/data/v2/index.json :: binary probe qcd_b_vs_light: (1 - AUC) of FLAV_F1R over FLAV_F0, run by run :: sha256 c6c4138ee5c86423
\newcommand{\MlstCmpQcdBFonerOverFzeroPr}{at least 11\% fewer errors}  % experiments/FIGS/data/v2/index.json :: binary probe qcd_b_vs_light: (1 - AUC) of FLAV_F1R over FLAV_F0, run by run :: sha256 c6c4138ee5c86423 :: prose twin: two significant figures, direction in words
\newcommand{\MlstCmpQcdBFonerOverFzeroPrMax}{up to 18\% fewer errors}  % experiments/FIGS/data/v2/index.json :: binary probe qcd_b_vs_light: (1 - AUC) of FLAV_F1R over FLAV_F0, run by run :: sha256 c6c4138ee5c86423 :: prose twin of the run farther from 1
\newcommand{\MlstCmpQcdBFzeroOverCoarse}{0.962 and 0.981}  % experiments/FIGS/data/v2/index.json :: binary probe qcd_b_vs_light: (1 - AUC) of FLAV_F0 over R16_Q1, run by run :: sha256 c6c4138ee5c86423
\newcommand{\MlstCmpQcdBFzeroOverCoarsePr}{equal within 4\%}  % experiments/FIGS/data/v2/index.json :: binary probe qcd_b_vs_light: (1 - AUC) of FLAV_F0 over R16_Q1, run by run :: sha256 c6c4138ee5c86423 :: prose twin: two significant figures, direction in words
\newcommand{\MlstCmpQcdBFzeroOverCoarsePrMax}{equal within 4\%}  % experiments/FIGS/data/v2/index.json :: binary probe qcd_b_vs_light: (1 - AUC) of FLAV_F0 over R16_Q1, run by run :: sha256 c6c4138ee5c86423 :: prose twin of the run farther from 1
\newcommand{\MlstCmpQcdBFoneOverCoarse}{0.749 and 0.819}  % experiments/FIGS/data/v2/index.json :: binary probe qcd_b_vs_light: (1 - AUC) of FLAV_F1 over R16_Q1, run by run :: sha256 c6c4138ee5c86423
\newcommand{\MlstCmpQcdBFoneOverCoarsePr}{at least 18\% fewer errors}  % experiments/FIGS/data/v2/index.json :: binary probe qcd_b_vs_light: (1 - AUC) of FLAV_F1 over R16_Q1, run by run :: sha256 c6c4138ee5c86423 :: prose twin: two significant figures, direction in words
\newcommand{\MlstCmpQcdBFoneOverCoarsePrMax}{up to 25\% fewer errors}  % experiments/FIGS/data/v2/index.json :: binary probe qcd_b_vs_light: (1 - AUC) of FLAV_F1 over R16_Q1, run by run :: sha256 c6c4138ee5c86423 :: prose twin of the run farther from 1
\newcommand{\MlstCmpQcdBFonerOverCoarse}{0.789 and 0.869}  % experiments/FIGS/data/v2/index.json :: binary probe qcd_b_vs_light: (1 - AUC) of FLAV_F1R over R16_Q1, run by run :: sha256 c6c4138ee5c86423
\newcommand{\MlstCmpQcdBFonerOverCoarsePr}{at least 13\% fewer errors}  % experiments/FIGS/data/v2/index.json :: binary probe qcd_b_vs_light: (1 - AUC) of FLAV_F1R over R16_Q1, run by run :: sha256 c6c4138ee5c86423 :: prose twin: two significant figures, direction in words
\newcommand{\MlstCmpQcdBFonerOverCoarsePrMax}{up to 21\% fewer errors}  % experiments/FIGS/data/v2/index.json :: binary probe qcd_b_vs_light: (1 - AUC) of FLAV_F1R over R16_Q1, run by run :: sha256 c6c4138ee5c86423 :: prose twin of the run farther from 1
\newcommand{\MlstCmpQcdBFzeroOverFull}{1.56 and 1.45}  % experiments/FIGS/data/v2/index.json :: binary probe qcd_b_vs_light: (1 - AUC) of FLAV_F0 over L188, run by run :: sha256 c6c4138ee5c86423
\newcommand{\MlstCmpQcdBFzeroOverFullPr}{at least 45\% more errors}  % experiments/FIGS/data/v2/index.json :: binary probe qcd_b_vs_light: (1 - AUC) of FLAV_F0 over L188, run by run :: sha256 c6c4138ee5c86423 :: prose twin: two significant figures, direction in words
\newcommand{\MlstCmpQcdBFzeroOverFullPrMax}{up to 56\% more errors}  % experiments/FIGS/data/v2/index.json :: binary probe qcd_b_vs_light: (1 - AUC) of FLAV_F0 over L188, run by run :: sha256 c6c4138ee5c86423 :: prose twin of the run farther from 1
\newcommand{\MlstCmpQcdBFoneOverFull}{1.212 and 1.213}  % experiments/FIGS/data/v2/index.json :: binary probe qcd_b_vs_light: (1 - AUC) of FLAV_F1 over L188, run by run :: sha256 c6c4138ee5c86423
\newcommand{\MlstCmpQcdBFoneOverFullPr}{at least 21\% more errors}  % experiments/FIGS/data/v2/index.json :: binary probe qcd_b_vs_light: (1 - AUC) of FLAV_F1 over L188, run by run :: sha256 c6c4138ee5c86423 :: prose twin: two significant figures, direction in words
\newcommand{\MlstCmpQcdBFoneOverFullPrMax}{up to 21\% more errors}  % experiments/FIGS/data/v2/index.json :: binary probe qcd_b_vs_light: (1 - AUC) of FLAV_F1 over L188, run by run :: sha256 c6c4138ee5c86423 :: prose twin of the run farther from 1
\newcommand{\MlstCmpQcdBFonerOverFull}{1.277 and 1.286}  % experiments/FIGS/data/v2/index.json :: binary probe qcd_b_vs_light: (1 - AUC) of FLAV_F1R over L188, run by run :: sha256 c6c4138ee5c86423
\newcommand{\MlstCmpQcdBFonerOverFullPr}{at least 28\% more errors}  % experiments/FIGS/data/v2/index.json :: binary probe qcd_b_vs_light: (1 - AUC) of FLAV_F1R over L188, run by run :: sha256 c6c4138ee5c86423 :: prose twin: two significant figures, direction in words
\newcommand{\MlstCmpQcdBFonerOverFullPrMax}{up to 29\% more errors}  % experiments/FIGS/data/v2/index.json :: binary probe qcd_b_vs_light: (1 - AUC) of FLAV_F1R over L188, run by run :: sha256 c6c4138ee5c86423 :: prose twin of the run farther from 1
\newcommand{\MlstCmpQcdCResOverFull}{1.061\,\ensuremath{\pm}\,0.005}  % experiments/FIGS/data/v2/index.json :: binary probe qcd_c_vs_light: (1 - AUC) of L162 over L188, run by run :: sha256 c6c4138ee5c86423
\newcommand{\MlstCmpQcdCResOverFullPr}{6.1\% more errors}  % experiments/FIGS/data/v2/index.json :: binary probe qcd_c_vs_light: (1 - AUC) of L162 over L188, run by run :: sha256 c6c4138ee5c86423 :: prose twin: two significant figures, direction in words
\newcommand{\MlstCmpQcdCMidOverFull}{1.072\,\ensuremath{\pm}\,0.008}  % experiments/FIGS/data/v2/index.json :: binary probe qcd_c_vs_light: (1 - AUC) of R42_Q1 over L188, run by run :: sha256 c6c4138ee5c86423
\newcommand{\MlstCmpQcdCMidOverFullPr}{7.2\% more errors}  % experiments/FIGS/data/v2/index.json :: binary probe qcd_c_vs_light: (1 - AUC) of R42_Q1 over L188, run by run :: sha256 c6c4138ee5c86423 :: prose twin: two significant figures, direction in words
\newcommand{\MlstCmpQcdCCoarseOverFull}{1.407\,\ensuremath{\pm}\,0.021}  % experiments/FIGS/data/v2/index.json :: binary probe qcd_c_vs_light: (1 - AUC) of R16_Q1 over L188, run by run :: sha256 c6c4138ee5c86423
\newcommand{\MlstCmpQcdCCoarseOverFullPr}{41\% more errors}  % experiments/FIGS/data/v2/index.json :: binary probe qcd_c_vs_light: (1 - AUC) of R16_Q1 over L188, run by run :: sha256 c6c4138ee5c86423 :: prose twin: two significant figures, direction in words
\newcommand{\MlstCmpQcdCFoneOverFzero}{0.915 and 0.944}  % experiments/FIGS/data/v2/index.json :: binary probe qcd_c_vs_light: (1 - AUC) of FLAV_F1 over FLAV_F0, run by run :: sha256 c6c4138ee5c86423
\newcommand{\MlstCmpQcdCFoneOverFzeroPr}{at least 5.6\% fewer errors}  % experiments/FIGS/data/v2/index.json :: binary probe qcd_c_vs_light: (1 - AUC) of FLAV_F1 over FLAV_F0, run by run :: sha256 c6c4138ee5c86423 :: prose twin: two significant figures, direction in words
\newcommand{\MlstCmpQcdCFoneOverFzeroPrMax}{up to 8.5\% fewer errors}  % experiments/FIGS/data/v2/index.json :: binary probe qcd_c_vs_light: (1 - AUC) of FLAV_F1 over FLAV_F0, run by run :: sha256 c6c4138ee5c86423 :: prose twin of the run farther from 1
\newcommand{\MlstCmpQcdCFonerOverFone}{0.9535 and 0.9520}  % experiments/FIGS/data/v2/index.json :: binary probe qcd_c_vs_light: (1 - AUC) of FLAV_F1R over FLAV_F1, run by run :: sha256 c6c4138ee5c86423
\newcommand{\MlstCmpQcdCFonerOverFonePr}{at least 4.6\% fewer errors}  % experiments/FIGS/data/v2/index.json :: binary probe qcd_c_vs_light: (1 - AUC) of FLAV_F1R over FLAV_F1, run by run :: sha256 c6c4138ee5c86423 :: prose twin: two significant figures, direction in words
\newcommand{\MlstCmpQcdCFonerOverFonePrMax}{up to 4.8\% fewer errors}  % experiments/FIGS/data/v2/index.json :: binary probe qcd_c_vs_light: (1 - AUC) of FLAV_F1R over FLAV_F1, run by run :: sha256 c6c4138ee5c86423 :: prose twin of the run farther from 1
\newcommand{\MlstCmpQcdCFonerOverFzero}{0.872 and 0.899}  % experiments/FIGS/data/v2/index.json :: binary probe qcd_c_vs_light: (1 - AUC) of FLAV_F1R over FLAV_F0, run by run :: sha256 c6c4138ee5c86423
\newcommand{\MlstCmpQcdCFonerOverFzeroPr}{at least 10\% fewer errors}  % experiments/FIGS/data/v2/index.json :: binary probe qcd_c_vs_light: (1 - AUC) of FLAV_F1R over FLAV_F0, run by run :: sha256 c6c4138ee5c86423 :: prose twin: two significant figures, direction in words
\newcommand{\MlstCmpQcdCFonerOverFzeroPrMax}{up to 13\% fewer errors}  % experiments/FIGS/data/v2/index.json :: binary probe qcd_c_vs_light: (1 - AUC) of FLAV_F1R over FLAV_F0, run by run :: sha256 c6c4138ee5c86423 :: prose twin of the run farther from 1
\newcommand{\MlstCmpQcdCFzeroOverCoarse}{0.952 and 0.934}  % experiments/FIGS/data/v2/index.json :: binary probe qcd_c_vs_light: (1 - AUC) of FLAV_F0 over R16_Q1, run by run :: sha256 c6c4138ee5c86423
\newcommand{\MlstCmpQcdCFzeroOverCoarsePr}{at least 4.8\% fewer errors}  % experiments/FIGS/data/v2/index.json :: binary probe qcd_c_vs_light: (1 - AUC) of FLAV_F0 over R16_Q1, run by run :: sha256 c6c4138ee5c86423 :: prose twin: two significant figures, direction in words
\newcommand{\MlstCmpQcdCFzeroOverCoarsePrMax}{up to 6.6\% fewer errors}  % experiments/FIGS/data/v2/index.json :: binary probe qcd_c_vs_light: (1 - AUC) of FLAV_F0 over R16_Q1, run by run :: sha256 c6c4138ee5c86423 :: prose twin of the run farther from 1
\newcommand{\MlstCmpQcdCFoneOverCoarse}{0.871 and 0.882}  % experiments/FIGS/data/v2/index.json :: binary probe qcd_c_vs_light: (1 - AUC) of FLAV_F1 over R16_Q1, run by run :: sha256 c6c4138ee5c86423
\newcommand{\MlstCmpQcdCFoneOverCoarsePr}{at least 12\% fewer errors}  % experiments/FIGS/data/v2/index.json :: binary probe qcd_c_vs_light: (1 - AUC) of FLAV_F1 over R16_Q1, run by run :: sha256 c6c4138ee5c86423 :: prose twin: two significant figures, direction in words
\newcommand{\MlstCmpQcdCFoneOverCoarsePrMax}{up to 13\% fewer errors}  % experiments/FIGS/data/v2/index.json :: binary probe qcd_c_vs_light: (1 - AUC) of FLAV_F1 over R16_Q1, run by run :: sha256 c6c4138ee5c86423 :: prose twin of the run farther from 1
\newcommand{\MlstCmpQcdCFonerOverCoarse}{0.8301 and 0.8396}  % experiments/FIGS/data/v2/index.json :: binary probe qcd_c_vs_light: (1 - AUC) of FLAV_F1R over R16_Q1, run by run :: sha256 c6c4138ee5c86423
\newcommand{\MlstCmpQcdCFonerOverCoarsePr}{at least 16\% fewer errors}  % experiments/FIGS/data/v2/index.json :: binary probe qcd_c_vs_light: (1 - AUC) of FLAV_F1R over R16_Q1, run by run :: sha256 c6c4138ee5c86423 :: prose twin: two significant figures, direction in words
\newcommand{\MlstCmpQcdCFonerOverCoarsePrMax}{up to 17\% fewer errors}  % experiments/FIGS/data/v2/index.json :: binary probe qcd_c_vs_light: (1 - AUC) of FLAV_F1R over R16_Q1, run by run :: sha256 c6c4138ee5c86423 :: prose twin of the run farther from 1
\newcommand{\MlstCmpQcdCFzeroOverFull}{1.352 and 1.293}  % experiments/FIGS/data/v2/index.json :: binary probe qcd_c_vs_light: (1 - AUC) of FLAV_F0 over L188, run by run :: sha256 c6c4138ee5c86423
\newcommand{\MlstCmpQcdCFzeroOverFullPr}{at least 29\% more errors}  % experiments/FIGS/data/v2/index.json :: binary probe qcd_c_vs_light: (1 - AUC) of FLAV_F0 over L188, run by run :: sha256 c6c4138ee5c86423 :: prose twin: two significant figures, direction in words
\newcommand{\MlstCmpQcdCFzeroOverFullPrMax}{up to 35\% more errors}  % experiments/FIGS/data/v2/index.json :: binary probe qcd_c_vs_light: (1 - AUC) of FLAV_F0 over L188, run by run :: sha256 c6c4138ee5c86423 :: prose twin of the run farther from 1
\newcommand{\MlstCmpQcdCFoneOverFull}{1.237 and 1.220}  % experiments/FIGS/data/v2/index.json :: binary probe qcd_c_vs_light: (1 - AUC) of FLAV_F1 over L188, run by run :: sha256 c6c4138ee5c86423
\newcommand{\MlstCmpQcdCFoneOverFullPr}{at least 22\% more errors}  % experiments/FIGS/data/v2/index.json :: binary probe qcd_c_vs_light: (1 - AUC) of FLAV_F1 over L188, run by run :: sha256 c6c4138ee5c86423 :: prose twin: two significant figures, direction in words
\newcommand{\MlstCmpQcdCFoneOverFullPrMax}{up to 24\% more errors}  % experiments/FIGS/data/v2/index.json :: binary probe qcd_c_vs_light: (1 - AUC) of FLAV_F1 over L188, run by run :: sha256 c6c4138ee5c86423 :: prose twin of the run farther from 1
\newcommand{\MlstCmpQcdCFonerOverFull}{1.179 and 1.162}  % experiments/FIGS/data/v2/index.json :: binary probe qcd_c_vs_light: (1 - AUC) of FLAV_F1R over L188, run by run :: sha256 c6c4138ee5c86423
\newcommand{\MlstCmpQcdCFonerOverFullPr}{at least 16\% more errors}  % experiments/FIGS/data/v2/index.json :: binary probe qcd_c_vs_light: (1 - AUC) of FLAV_F1R over L188, run by run :: sha256 c6c4138ee5c86423 :: prose twin: two significant figures, direction in words
\newcommand{\MlstCmpQcdCFonerOverFullPrMax}{up to 18\% more errors}  % experiments/FIGS/data/v2/index.json :: binary probe qcd_c_vs_light: (1 - AUC) of FLAV_F1R over L188, run by run :: sha256 c6c4138ee5c86423 :: prose twin of the run farther from 1
\newcommand{\MlstLpEemmAucFull}{1.0000\,\ensuremath{\pm}\,\ensuremath{<}0.0001}  % experiments/FIGS/data/v2/index.json :: linear probe ee_vs_mm (fits_v4_sib): AUC, L188 :: sha256 c6c4138ee5c86423
\newcommand{\MlstLpEemmAucRes}{1.0000\,\ensuremath{\pm}\,\ensuremath{<}0.0001}  % experiments/FIGS/data/v2/index.json :: linear probe ee_vs_mm (fits_v4_sib): AUC, L162 :: sha256 c6c4138ee5c86423
\newcommand{\MlstLpEemmAucMid}{0.9610\,\ensuremath{\pm}\,0.0030}  % experiments/FIGS/data/v2/index.json :: linear probe ee_vs_mm (fits_v4_sib): AUC, R42_Q1 :: sha256 c6c4138ee5c86423
\newcommand{\MlstLpEemmAucCoarse}{0.960\,\ensuremath{\pm}\,0.020}  % experiments/FIGS/data/v2/index.json :: linear probe ee_vs_mm (fits_v4_sib): AUC, R16_Q1 :: sha256 c6c4138ee5c86423
\newcommand{\MlstLpEemmAucFzero}{1.0000 and 1.0000}  % experiments/FIGS/data/v2/index.json :: linear probe ee_vs_mm (fits_v4_sib): AUC, FLAV_F0 :: sha256 c6c4138ee5c86423
\newcommand{\MlstLpEemmAucFone}{1.0000 and 1.0000}  % experiments/FIGS/data/v2/index.json :: linear probe ee_vs_mm (fits_v4_sib): AUC, FLAV_F1 :: sha256 c6c4138ee5c86423
\newcommand{\MlstLpEemmAucFoner}{1.0000 and 1.0000}  % experiments/FIGS/data/v2/index.json :: linear probe ee_vs_mm (fits_v4_sib): AUC, FLAV_F1R :: sha256 c6c4138ee5c86423
\newcommand{\MlstLpEemmAucInit}{0.9995\,\ensuremath{\pm}\,0.0003}  % experiments/FIGS/data/v2/index.json :: linear probe ee_vs_mm (fits_v4_sib): AUC, INIT :: sha256 c6c4138ee5c86423
\newcommand{\MlstLpEemmAucInitBn}{1.0000\,\ensuremath{\pm}\,\ensuremath{<}0.0001}  % experiments/FIGS/data/v2/index.json :: linear probe ee_vs_mm (fits_v4_sib): AUC, INIT_BN :: sha256 c6c4138ee5c86423
\newcommand{\MlstLpEemmAucRandMerged}{1.0000\,\ensuremath{\pm}\,\ensuremath{<}0.0001}  % experiments/FIGS/data/v2/index.json :: linear probe ee_vs_mm (fits_v4_sib): AUC, the random partitions that merge the pair :: sha256 c6c4138ee5c86423
\newcommand{\MlstPairEemmNTrain}{9{,}571}  % experiments/FIGS/data/v2/index.json :: fits_v4_sib ee_vs_mm: n_train :: sha256 c6c4138ee5c86423
\newcommand{\MlstPairEemmNVal}{1{,}068}  % experiments/FIGS/data/v2/index.json :: fits_v4_sib ee_vs_mm: n_val :: sha256 c6c4138ee5c86423
\newcommand{\MlstPairEemmNTestPerSide}{128{,}833--135{,}245}  % experiments/FIGS/data/v2/index.json :: fits_v4_sib ee_vs_mm: n_test_signal and n_test_background :: sha256 c6c4138ee5c86423
\newcommand{\MlstLpJcTwoAucFull}{0.9801\,\ensuremath{\pm}\,\ensuremath{<}0.0001}  % experiments/FIGS/data/v2/index.json :: linear probe JetClass-II N=1000000: AUC averaged over the classes, L188 :: sha256 c6c4138ee5c86423
\newcommand{\MlstLpJcTwoAccFull}{0.4348\,\ensuremath{\pm}\,0.0003}  % experiments/FIGS/data/v2/index.json :: linear probe JetClass-II N=1000000: accuracy, L188 :: sha256 c6c4138ee5c86423
\newcommand{\MlstLpJcTwoAucRes}{0.9796\,\ensuremath{\pm}\,0.0001}  % experiments/FIGS/data/v2/index.json :: linear probe JetClass-II N=1000000: AUC averaged over the classes, L162 :: sha256 c6c4138ee5c86423
\newcommand{\MlstLpJcTwoAccRes}{0.4328\,\ensuremath{\pm}\,0.0008}  % experiments/FIGS/data/v2/index.json :: linear probe JetClass-II N=1000000: accuracy, L162 :: sha256 c6c4138ee5c86423
\newcommand{\MlstLpJcTwoAucMid}{0.9756\,\ensuremath{\pm}\,0.0002}  % experiments/FIGS/data/v2/index.json :: linear probe JetClass-II N=1000000: AUC averaged over the classes, R42_Q1 :: sha256 c6c4138ee5c86423
\newcommand{\MlstLpJcTwoAccMid}{0.361\,\ensuremath{\pm}\,0.005}  % experiments/FIGS/data/v2/index.json :: linear probe JetClass-II N=1000000: accuracy, R42_Q1 :: sha256 c6c4138ee5c86423
\newcommand{\MlstLpJcTwoAucCoarse}{0.9691\,\ensuremath{\pm}\,0.0004}  % experiments/FIGS/data/v2/index.json :: linear probe JetClass-II N=1000000: AUC averaged over the classes, R16_Q1 :: sha256 c6c4138ee5c86423
\newcommand{\MlstLpJcTwoAccCoarse}{0.334\,\ensuremath{\pm}\,0.007}  % experiments/FIGS/data/v2/index.json :: linear probe JetClass-II N=1000000: accuracy, R16_Q1 :: sha256 c6c4138ee5c86423
\newcommand{\MlstLpJcTwoAucResMass}{0.9791 and 0.9794}  % experiments/FIGS/data/v2/index.json :: linear probe JetClass-II N=1000000: AUC averaged over the classes, L162_MASS :: sha256 c6c4138ee5c86423
\newcommand{\MlstLpJcTwoAccResMass}{0.4271 and 0.4305}  % experiments/FIGS/data/v2/index.json :: linear probe JetClass-II N=1000000: accuracy, L162_MASS :: sha256 c6c4138ee5c86423
\newcommand{\MlstLpJcTwoAucCoarseMass}{0.9679 and 0.9696}  % experiments/FIGS/data/v2/index.json :: linear probe JetClass-II N=1000000: AUC averaged over the classes, R16_Q1_MASS :: sha256 c6c4138ee5c86423
\newcommand{\MlstLpJcTwoAccCoarseMass}{0.312 and 0.350}  % experiments/FIGS/data/v2/index.json :: linear probe JetClass-II N=1000000: accuracy, R16_Q1_MASS :: sha256 c6c4138ee5c86423
\newcommand{\MlstLpJcTwoAucFzero}{0.9717 and 0.9710}  % experiments/FIGS/data/v2/index.json :: linear probe JetClass-II N=1000000: AUC averaged over the classes, FLAV_F0 :: sha256 c6c4138ee5c86423
\newcommand{\MlstLpJcTwoAccFzero}{0.3769 and 0.3723}  % experiments/FIGS/data/v2/index.json :: linear probe JetClass-II N=1000000: accuracy, FLAV_F0 :: sha256 c6c4138ee5c86423
\newcommand{\MlstLpJcTwoAucFone}{0.9758 and 0.9759}  % experiments/FIGS/data/v2/index.json :: linear probe JetClass-II N=1000000: AUC averaged over the classes, FLAV_F1 :: sha256 c6c4138ee5c86423
\newcommand{\MlstLpJcTwoAccFone}{0.3945 and 0.3963}  % experiments/FIGS/data/v2/index.json :: linear probe JetClass-II N=1000000: accuracy, FLAV_F1 :: sha256 c6c4138ee5c86423
\newcommand{\MlstLpJcTwoAucFoner}{0.9760 and 0.9760}  % experiments/FIGS/data/v2/index.json :: linear probe JetClass-II N=1000000: AUC averaged over the classes, FLAV_F1R :: sha256 c6c4138ee5c86423
\newcommand{\MlstLpJcTwoAccFoner}{0.4071 and 0.4089}  % experiments/FIGS/data/v2/index.json :: linear probe JetClass-II N=1000000: accuracy, FLAV_F1R :: sha256 c6c4138ee5c86423
\newcommand{\MlstLpJcTwoAucInit}{0.8121\,\ensuremath{\pm}\,0.0019}  % experiments/FIGS/data/v2/index.json :: linear probe JetClass-II N=1000000: AUC averaged over the classes, INIT :: sha256 c6c4138ee5c86423
\newcommand{\MlstLpJcTwoAccInit}{0.0724\,\ensuremath{\pm}\,0.0015}  % experiments/FIGS/data/v2/index.json :: linear probe JetClass-II N=1000000: accuracy, INIT :: sha256 c6c4138ee5c86423
\newcommand{\MlstLpJcTwoAucInitBn}{0.866\,\ensuremath{\pm}\,0.004}  % experiments/FIGS/data/v2/index.json :: linear probe JetClass-II N=1000000: AUC averaged over the classes, INIT_BN :: sha256 c6c4138ee5c86423
\newcommand{\MlstLpJcTwoAccInitBn}{0.1017\,\ensuremath{\pm}\,0.0017}  % experiments/FIGS/data/v2/index.json :: linear probe JetClass-II N=1000000: accuracy, INIT_BN :: sha256 c6c4138ee5c86423
\newcommand{\MlstLpJcTwoAucRand}{0.9777\,\ensuremath{\pm}\,0.0002}  % experiments/FIGS/data/v2/index.json :: linear probe JetClass-II N=1000000: AUC averaged over the classes, the random partitions :: sha256 c6c4138ee5c86423
\newcommand{\MlstLpJcOneAccInitBn}{0.475\,\ensuremath{\pm}\,0.005}  % experiments/FIGS/data/v2/index.json :: linear probe JetClass N=1000000: accuracy, INIT_BN :: sha256 c6c4138ee5c86423
\newcommand{\MlstBnMaxDev}{0.2\%}  % experiments/FIGS/data/v2/index.json :: largest |r(best70_bn) / r(best70) - 1| of a JetClass-II ratio of 1 - AUC to the 188-class run, per cent :: sha256 c6c4138ee5c86423
\newcommand{\MlstBnMaxDevPair}{0.7\%}  % experiments/FIGS/data/v2/index.json :: largest |r(best70_bn) / r(best70) - 1| of a ratio of 1 - AUC on a class pair to the 188-class run, levels of the hierarchy, mass models and controlled partitions, per cent :: sha256 c6c4138ee5c86423 :: pairs with a check: bc_vs_bq_cs, bb_vs_cc, bbqq_vs_ccqq, qcd_bb_vs_cc, qcd_b_vs_light, qcd_c_vs_light; without: ee_vs_mm
\newcommand{\MlstGpuTypeBound}{1.4\%}  % experiments/FIGS/data/v2/index.json :: largest |label-set ratio of 1 - AUC of index 3 / mean of the other indices - 1|, JetClass-II largest set, per cent :: sha256 c6c4138ee5c86423 :: mtx-l188-s3 trained on NVIDIA GeForce RTX 3090; largest for R16_Q1; label sets L162 0.09%, R42_Q1 0.70%, R16_Q1 1.37%
\newcommand{\MlstGpuOtherSeedsDiff}{1.9\%}  % experiments/FIGS/data/v2/index.json :: largest (max - min) / mean of a label-set ratio of 1 - AUC over the indices other than 3, JetClass-II largest set, per cent :: sha256 c6c4138ee5c86423 :: label sets L162 0.90%, R42_Q1 1.14%, R16_Q1 1.93%
\newcommand{\MlstPerClassResNDoubled}{0}  % experiments/FIGS/data/v2/index.json :: JetClass-II classes whose mean ratio of 1 - AUC (L162 over 188) is at least 2 (count) :: sha256 c6c4138ee5c86423
\newcommand{\MlstPerClassResSeedSdMedian}{0.011}  % experiments/FIGS/data/v2/index.json :: median over classes of the SD over runs of the per-class ratio, L162 over 188 :: sha256 c6c4138ee5c86423
\newcommand{\MlstPerClassResTopClass}{QCD, \ensuremath{ccss}}  % experiments/FIGS/data/v2/index.json :: class with the largest mean ratio, L162 :: sha256 c6c4138ee5c86423
\newcommand{\MlstPerClassResTopRatio}{1.129\,\ensuremath{\pm}\,0.010}  % experiments/FIGS/data/v2/index.json :: per-class ratio of 1 - AUC of that class, L162 over 188, run by run :: sha256 c6c4138ee5c86423
\newcommand{\MlstPerClassResTopRatioPr}{13\% more errors}  % experiments/FIGS/data/v2/index.json :: per-class ratio of 1 - AUC of that class, L162 over 188, run by run :: sha256 c6c4138ee5c86423 :: prose twin: two significant figures, direction in words
\newcommand{\MlstPerClassResLeptonN}{0}  % experiments/FIGS/data/v2/index.json :: resonant classes sharing their L162 group with a decay that differs only in lepton flavor (count) :: sha256 c6c4138ee5c86423
\newcommand{\MlstPerClassResLeptonMedian}{--}  % experiments/FIGS/data/v2/index.json :: median over those classes of the mean per-class ratio, L162 over 188 :: sha256 c6c4138ee5c86423
\newcommand{\MlstPerClassResLeptonMedianPr}{\pending{no run yet}}  % experiments/FIGS/data/v2/index.json :: median over those classes of the mean per-class ratio, L162 over 188 :: sha256 c6c4138ee5c86423 :: prose twin: two significant figures, direction in words
\newcommand{\MlstPerClassResLeptonShare}{--}  % experiments/FIGS/data/v2/index.json :: their share of the summed ln ratio over the resonant classes, L162 :: sha256 c6c4138ee5c86423
\newcommand{\MlstPerClassResQuarkN}{0}  % experiments/FIGS/data/v2/index.json :: resonant classes sharing their L162 group with a decay that differs only in quark flavor (count) :: sha256 c6c4138ee5c86423
\newcommand{\MlstPerClassResQuarkMedian}{--}  % experiments/FIGS/data/v2/index.json :: median over those classes of the mean per-class ratio, L162 over 188 :: sha256 c6c4138ee5c86423
\newcommand{\MlstPerClassResQuarkMedianPr}{\pending{no run yet}}  % experiments/FIGS/data/v2/index.json :: median over those classes of the mean per-class ratio, L162 over 188 :: sha256 c6c4138ee5c86423 :: prose twin: two significant figures, direction in words
\newcommand{\MlstPerClassResQuarkShare}{--}  % experiments/FIGS/data/v2/index.json :: their share of the summed ln ratio over the resonant classes, L162 :: sha256 c6c4138ee5c86423
\newcommand{\MlstPerClassResQuarkLeptonN}{0}  % experiments/FIGS/data/v2/index.json :: resonant classes counted in both the quark- and the lepton-flavor set, L162 (count) :: sha256 c6c4138ee5c86423
\newcommand{\MlstPerClassMidNDoubled}{40}  % experiments/FIGS/data/v2/index.json :: JetClass-II classes whose mean ratio of 1 - AUC (R42_Q1 over 188) is at least 2 (count) :: sha256 c6c4138ee5c86423
\newcommand{\MlstPerClassMidSeedSdMedian}{0.023}  % experiments/FIGS/data/v2/index.json :: median over classes of the SD over runs of the per-class ratio, R42_Q1 over 188 :: sha256 c6c4138ee5c86423
\newcommand{\MlstPerClassMidTopClass}{\ensuremath{X\to YY\to bb\mu\mu}}  % experiments/FIGS/data/v2/index.json :: class with the largest mean ratio, R42_Q1 :: sha256 c6c4138ee5c86423
\newcommand{\MlstPerClassMidTopRatio}{4.06\,\ensuremath{\pm}\,0.14}  % experiments/FIGS/data/v2/index.json :: per-class ratio of 1 - AUC of that class, R42_Q1 over 188, run by run :: sha256 c6c4138ee5c86423
\newcommand{\MlstPerClassMidTopRatioPr}{4.1 times the errors}  % experiments/FIGS/data/v2/index.json :: per-class ratio of 1 - AUC of that class, R42_Q1 over 188, run by run :: sha256 c6c4138ee5c86423 :: prose twin: two significant figures, direction in words
\newcommand{\MlstPerClassMidLeptonN}{78}  % experiments/FIGS/data/v2/index.json :: resonant classes sharing their R42_Q1 group with a decay that differs only in lepton flavor (count) :: sha256 c6c4138ee5c86423
\newcommand{\MlstPerClassMidLeptonMedian}{2.02}  % experiments/FIGS/data/v2/index.json :: median over those classes of the mean per-class ratio, R42_Q1 over 188 :: sha256 c6c4138ee5c86423
\newcommand{\MlstPerClassMidLeptonMedianPr}{2.0 times the errors}  % experiments/FIGS/data/v2/index.json :: median over those classes of the mean per-class ratio, R42_Q1 over 188 :: sha256 c6c4138ee5c86423 :: prose twin: two significant figures, direction in words
\newcommand{\MlstPerClassMidLeptonShare}{90\%}  % experiments/FIGS/data/v2/index.json :: their share of the summed ln ratio over the resonant classes, R42_Q1 :: sha256 c6c4138ee5c86423
\newcommand{\MlstPerClassMidQuarkN}{111}  % experiments/FIGS/data/v2/index.json :: resonant classes sharing their R42_Q1 group with a decay that differs only in quark flavor (count) :: sha256 c6c4138ee5c86423
\newcommand{\MlstPerClassMidQuarkMedian}{1.14}  % experiments/FIGS/data/v2/index.json :: median over those classes of the mean per-class ratio, R42_Q1 over 188 :: sha256 c6c4138ee5c86423
\newcommand{\MlstPerClassMidQuarkMedianPr}{14\% more errors}  % experiments/FIGS/data/v2/index.json :: median over those classes of the mean per-class ratio, R42_Q1 over 188 :: sha256 c6c4138ee5c86423 :: prose twin: two significant figures, direction in words
\newcommand{\MlstPerClassMidQuarkShare}{53\%}  % experiments/FIGS/data/v2/index.json :: their share of the summed ln ratio over the resonant classes, R42_Q1 :: sha256 c6c4138ee5c86423
\newcommand{\MlstPerClassMidQuarkLeptonN}{44}  % experiments/FIGS/data/v2/index.json :: resonant classes counted in both the quark- and the lepton-flavor set, R42_Q1 (count) :: sha256 c6c4138ee5c86423
\newcommand{\MlstPerClassCoarseNDoubled}{60}  % experiments/FIGS/data/v2/index.json :: JetClass-II classes whose mean ratio of 1 - AUC (R16_Q1 over 188) is at least 2 (count) :: sha256 c6c4138ee5c86423
\newcommand{\MlstPerClassCoarseSeedSdMedian}{0.055}  % experiments/FIGS/data/v2/index.json :: median over classes of the SD over runs of the per-class ratio, R16_Q1 over 188 :: sha256 c6c4138ee5c86423
\newcommand{\MlstPerClassCoarseTopClass}{\ensuremath{X\to YY\to bbee}}  % experiments/FIGS/data/v2/index.json :: class with the largest mean ratio, R16_Q1 :: sha256 c6c4138ee5c86423
\newcommand{\MlstPerClassCoarseTopRatio}{5.2\,\ensuremath{\pm}\,0.4}  % experiments/FIGS/data/v2/index.json :: per-class ratio of 1 - AUC of that class, R16_Q1 over 188, run by run :: sha256 c6c4138ee5c86423
\newcommand{\MlstPerClassCoarseTopRatioPr}{5.2 times the errors}  % experiments/FIGS/data/v2/index.json :: per-class ratio of 1 - AUC of that class, R16_Q1 over 188, run by run :: sha256 c6c4138ee5c86423 :: prose twin: two significant figures, direction in words
\newcommand{\MlstPerClassCoarseLeptonN}{78}  % experiments/FIGS/data/v2/index.json :: resonant classes sharing their R16_Q1 group with a decay that differs only in lepton flavor (count) :: sha256 c6c4138ee5c86423
\newcommand{\MlstPerClassCoarseLeptonMedian}{2.27}  % experiments/FIGS/data/v2/index.json :: median over those classes of the mean per-class ratio, R16_Q1 over 188 :: sha256 c6c4138ee5c86423
\newcommand{\MlstPerClassCoarseLeptonMedianPr}{2.3 times the errors}  % experiments/FIGS/data/v2/index.json :: median over those classes of the mean per-class ratio, R16_Q1 over 188 :: sha256 c6c4138ee5c86423 :: prose twin: two significant figures, direction in words
\newcommand{\MlstPerClassCoarseLeptonShare}{68\%}  % experiments/FIGS/data/v2/index.json :: their share of the summed ln ratio over the resonant classes, R16_Q1 :: sha256 c6c4138ee5c86423
\newcommand{\MlstPerClassCoarseQuarkN}{141}  % experiments/FIGS/data/v2/index.json :: resonant classes sharing their R16_Q1 group with a decay that differs only in quark flavor (count) :: sha256 c6c4138ee5c86423
\newcommand{\MlstPerClassCoarseQuarkMedian}{1.65}  % experiments/FIGS/data/v2/index.json :: median over those classes of the mean per-class ratio, R16_Q1 over 188 :: sha256 c6c4138ee5c86423
\newcommand{\MlstPerClassCoarseQuarkMedianPr}{65\% more errors}  % experiments/FIGS/data/v2/index.json :: median over those classes of the mean per-class ratio, R16_Q1 over 188 :: sha256 c6c4138ee5c86423 :: prose twin: two significant figures, direction in words
\newcommand{\MlstPerClassCoarseQuarkShare}{87\%}  % experiments/FIGS/data/v2/index.json :: their share of the summed ln ratio over the resonant classes, R16_Q1 :: sha256 c6c4138ee5c86423
\newcommand{\MlstPerClassCoarseQuarkLeptonN}{64}  % experiments/FIGS/data/v2/index.json :: resonant classes counted in both the quark- and the lepton-flavor set, R16_Q1 (count) :: sha256 c6c4138ee5c86423
\newcommand{\MlstPerClassRandFlavMergedOverApart}{1.010\,\ensuremath{\pm}\,0.008}  % experiments/FIGS/data/v2/index.json :: random partitions, linear probe: median per-class ratio of the classes merged with a decay that differs only in quark flavor over that of the classes kept apart from all of them, mean +- SD over the partitions :: sha256 c6c4138ee5c86423
\newcommand{\MlstPerClassRandFlavMergedOverApartPr}{equal within 1\%}  % experiments/FIGS/data/v2/index.json :: random partitions, linear probe: median per-class ratio of the classes merged with a decay that differs only in quark flavor over that of the classes kept apart from all of them, mean +- SD over the partitions :: sha256 c6c4138ee5c86423 :: prose twin: two significant figures, direction in words
\newcommand{\MlstPerClassRandLeptonMergedOverApart}{0.997\,\ensuremath{\pm}\,0.020}  % experiments/FIGS/data/v2/index.json :: random partitions, linear probe: median per-class ratio of the classes merged with a decay that differs only in lepton flavor over that of the classes kept apart from all of them, mean +- SD over the partitions :: sha256 c6c4138ee5c86423
\newcommand{\MlstPerClassRandLeptonMergedOverApartPr}{equal within 3\%}  % experiments/FIGS/data/v2/index.json :: random partitions, linear probe: median per-class ratio of the classes merged with a decay that differs only in lepton flavor over that of the classes kept apart from all of them, mean +- SD over the partitions :: sha256 c6c4138ee5c86423 :: prose twin: two significant figures, direction in words
\newcommand{\MlstMlpNatPerClassRandFlavMergedOverApart}{1.003\,\ensuremath{\pm}\,0.007}  % experiments/FIGS/data/v2/index.json :: random partitions, nonlinear probe, natural composition: median per-class ratio of the classes merged with a decay that differs only in quark flavor over that of the classes kept apart from all of them, mean +- SD over the partitions :: sha256 c6c4138ee5c86423
\newcommand{\MlstMlpNatPerClassRandFlavMergedOverApartPr}{equal within 1\%}  % experiments/FIGS/data/v2/index.json :: random partitions, nonlinear probe, natural composition: median per-class ratio of the classes merged with a decay that differs only in quark flavor over that of the classes kept apart from all of them, mean +- SD over the partitions :: sha256 c6c4138ee5c86423 :: prose twin: two significant figures, direction in words
\newcommand{\MlstMlpNatPerClassRandLeptonMergedOverApart}{1.01\,\ensuremath{\pm}\,0.07}  % experiments/FIGS/data/v2/index.json :: random partitions, nonlinear probe, natural composition: median per-class ratio of the classes merged with a decay that differs only in lepton flavor over that of the classes kept apart from all of them, mean +- SD over the partitions :: sha256 c6c4138ee5c86423
\newcommand{\MlstMlpNatPerClassRandLeptonMergedOverApartPr}{equal within 8\%}  % experiments/FIGS/data/v2/index.json :: random partitions, nonlinear probe, natural composition: median per-class ratio of the classes merged with a decay that differs only in lepton flavor over that of the classes kept apart from all of them, mean +- SD over the partitions :: sha256 c6c4138ee5c86423 :: prose twin: two significant figures, direction in words
\newcommand{\MlstPerClassAucDiffResSeedSdMedian}{0.015}  % experiments/FIGS/data/v2/index.json :: median over classes of the SD over indices of 100 (AUC_c of L162 - AUC_c of 188), pp :: sha256 c6c4138ee5c86423
\newcommand{\MlstPerClassAucDiffResMin}{\ensuremath{-}0.73}  % experiments/FIGS/data/v2/index.json :: most negative mean over indices of 100 (AUC_c of L162 - AUC_c of 188), pp :: sha256 c6c4138ee5c86423 :: label_QCD_bcss
\newcommand{\MlstPerClassAucDiffResMinClass}{QCD, \ensuremath{bcss}}  % experiments/FIGS/data/v2/index.json :: its class, L162 :: sha256 c6c4138ee5c86423
\newcommand{\MlstPerClassAucDiffMidSeedSdMedian}{0.032}  % experiments/FIGS/data/v2/index.json :: median over classes of the SD over indices of 100 (AUC_c of R42_Q1 - AUC_c of 188), pp :: sha256 c6c4138ee5c86423
\newcommand{\MlstPerClassAucDiffMidMin}{\ensuremath{-}1.4}  % experiments/FIGS/data/v2/index.json :: most negative mean over indices of 100 (AUC_c of R42_Q1 - AUC_c of 188), pp :: sha256 c6c4138ee5c86423 :: label_X_YY_cqtauev
\newcommand{\MlstPerClassAucDiffMidMinClass}{\ensuremath{X\to YY\to cq\tau_e\nu}}  % experiments/FIGS/data/v2/index.json :: its class, R42_Q1 :: sha256 c6c4138ee5c86423
\newcommand{\MlstPerClassAucDiffCoarseSeedSdMedian}{0.077}  % experiments/FIGS/data/v2/index.json :: median over classes of the SD over indices of 100 (AUC_c of R16_Q1 - AUC_c of 188), pp :: sha256 c6c4138ee5c86423
\newcommand{\MlstPerClassAucDiffCoarseMin}{\ensuremath{-}3.1}  % experiments/FIGS/data/v2/index.json :: most negative mean over indices of 100 (AUC_c of R16_Q1 - AUC_c of 188), pp :: sha256 c6c4138ee5c86423 :: label_QCD_bcss
\newcommand{\MlstPerClassAucDiffCoarseMinClass}{QCD, \ensuremath{bcss}}  % experiments/FIGS/data/v2/index.json :: its class, R16_Q1 :: sha256 c6c4138ee5c86423
\newcommand{\MlstRocJcTwoHurtList}{\ensuremath{X\to YY\to ccq}; \ensuremath{X\to YY\to ssb}; \ensuremath{X\to YY\to ccss}; QCD, \ensuremath{bcc}; QCD, \ensuremath{bccss}; QCD, \ensuremath{bss}; QCD, \ensuremath{ccs} and QCD, \ensuremath{ccss}}  % experiments/FIGS/data/v2/index.json :: jc2_roc_classes on the class-balanced probes; candidates with at least 1000 test jets :: sha256 c6c4138ee5c86423 :: label_X_YY_ccq, label_X_YY_ssb, label_X_YY_ccss, label_QCD_bcc, label_QCD_bccss, label_QCD_bss, label_QCD_ccs, label_QCD_ccss
\newcommand{\MlstRocJcTwoStableList}{\ensuremath{X\to gg}; \ensuremath{X\to YY\to ggg}; \ensuremath{X\to YY\to g\tau_h\tau_h} and QCD, light}  % experiments/FIGS/data/v2/index.json :: jc2_roc_classes on the class-balanced probes; candidates with at least 1000 test jets :: sha256 c6c4138ee5c86423 :: label_X_gg, label_X_YY_ggg, label_X_YY_gtauhtauh, label_QCD_light
\newcommand{\MlstRocJcTwoHurtMinDrop}{2.1}  % experiments/FIGS/data/v2/index.json :: jc2_roc_classes on the class-balanced probes; candidates with at least 1000 test jets :: sha256 c6c4138ee5c86423 :: smallest AUC loss at 17 classes of the hurt, pp
\newcommand{\MlstRocJcTwoHurtMaxDrop}{2.8}  % experiments/FIGS/data/v2/index.json :: jc2_roc_classes on the class-balanced probes; candidates with at least 1000 test jets :: sha256 c6c4138ee5c86423 :: largest AUC loss at 17 classes of the hurt, pp
\newcommand{\MlstRocJcTwoStableMaxChange}{15}  % experiments/FIGS/data/v2/index.json :: jc2_roc_classes on the class-balanced probes; candidates with at least 1000 test jets :: sha256 c6c4138ee5c86423 :: largest change of 1 - AUC at any label set of the stable classes, per cent of the 188-class value
\newcommand{\MlstJcBalNThin}{4}  % experiments/FIGS/data/v2/jc2bal/census.json :: counts.ft_test :: sha256 a4379ee0817254cf :: classes with fewer than 1000 class-balanced test jets (F3 hatching)
\newcommand{\MlstMinBkgPass}{20}  % experiments/FIGS/mlst_v4.py :: JC1_MIN_BKG :: sha256 fa536b4866ffd651 :: fewest background jets behind a drawn ROC point or a quoted rejection
\newcommand{\MlstMassMinJets}{100}  % experiments/EVAL/atlas_resolution.py :: MIN_JETS :: sha256 b82cc219f79f0f14 :: jets below which a mass bin is null
\newcommand{\MlstRocCutBbFull}{0.79}  % experiments/FIGS/data/v2/index.json :: curves_v2/pairs curve bb_vs_cc: largest cut.eps_s over the L188 runs ('not reached' when a run's cut is null or leaves fewer than 2 grid points) :: sha256 c6c4138ee5c86423
\newcommand{\MlstRocCutBbRes}{0.77}  % experiments/FIGS/data/v2/index.json :: curves_v2/pairs curve bb_vs_cc: largest cut.eps_s over the L162 runs ('not reached' when a run's cut is null or leaves fewer than 2 grid points) :: sha256 c6c4138ee5c86423
\newcommand{\MlstRocCutBbMid}{0.79}  % experiments/FIGS/data/v2/index.json :: curves_v2/pairs curve bb_vs_cc: largest cut.eps_s over the R42_Q1 runs ('not reached' when a run's cut is null or leaves fewer than 2 grid points) :: sha256 c6c4138ee5c86423
\newcommand{\MlstRocCutBbCoarse}{0.13}  % experiments/FIGS/data/v2/index.json :: curves_v2/pairs curve bb_vs_cc: largest cut.eps_s over the R16_Q1 runs ('not reached' when a run's cut is null or leaves fewer than 2 grid points) :: sha256 c6c4138ee5c86423
\newcommand{\MlstRocCutBbFzero}{0.03}  % experiments/FIGS/data/v2/index.json :: curves_v2/pairs curve bb_vs_cc: largest cut.eps_s over the FLAV_F0 runs ('not reached' when a run's cut is null or leaves fewer than 2 grid points) :: sha256 c6c4138ee5c86423
\newcommand{\MlstRocCutBbFone}{0.62}  % experiments/FIGS/data/v2/index.json :: curves_v2/pairs curve bb_vs_cc: largest cut.eps_s over the FLAV_F1 runs ('not reached' when a run's cut is null or leaves fewer than 2 grid points) :: sha256 c6c4138ee5c86423
\newcommand{\MlstRocCutBbFoner}{0.51}  % experiments/FIGS/data/v2/index.json :: curves_v2/pairs curve bb_vs_cc: largest cut.eps_s over the FLAV_F1R runs ('not reached' when a run's cut is null or leaves fewer than 2 grid points) :: sha256 c6c4138ee5c86423
\newcommand{\MlstRocNBkgBb}{200{,}000}  % experiments/FIGS/data/v2/index.json :: curves_v2/pairs curve bb_vs_cc: n_background :: sha256 c6c4138ee5c86423
\newcommand{\MlstRocRejMaxBb}{10{,}000}  % experiments/FIGS/data/v2/index.json :: curves_v2/pairs curve bb_vs_cc: n_background / 20 :: sha256 c6c4138ee5c86423
\newcommand{\MlstRocCutBbqqFull}{0.67}  % experiments/FIGS/data/v2/index.json :: curves_v2/pairs curve bbqq_vs_ccqq: largest cut.eps_s over the L188 runs ('not reached' when a run's cut is null or leaves fewer than 2 grid points) :: sha256 c6c4138ee5c86423
\newcommand{\MlstRocCutBbqqRes}{0.66}  % experiments/FIGS/data/v2/index.json :: curves_v2/pairs curve bbqq_vs_ccqq: largest cut.eps_s over the L162 runs ('not reached' when a run's cut is null or leaves fewer than 2 grid points) :: sha256 c6c4138ee5c86423
\newcommand{\MlstRocCutBbqqMid}{0.62}  % experiments/FIGS/data/v2/index.json :: curves_v2/pairs curve bbqq_vs_ccqq: largest cut.eps_s over the R42_Q1 runs ('not reached' when a run's cut is null or leaves fewer than 2 grid points) :: sha256 c6c4138ee5c86423
\newcommand{\MlstRocCutBbqqCoarse}{0.06}  % experiments/FIGS/data/v2/index.json :: curves_v2/pairs curve bbqq_vs_ccqq: largest cut.eps_s over the R16_Q1 runs ('not reached' when a run's cut is null or leaves fewer than 2 grid points) :: sha256 c6c4138ee5c86423
\newcommand{\MlstRocCutBbqqFzero}{0.06}  % experiments/FIGS/data/v2/index.json :: curves_v2/pairs curve bbqq_vs_ccqq: largest cut.eps_s over the FLAV_F0 runs ('not reached' when a run's cut is null or leaves fewer than 2 grid points) :: sha256 c6c4138ee5c86423
\newcommand{\MlstRocCutBbqqFone}{0.56}  % experiments/FIGS/data/v2/index.json :: curves_v2/pairs curve bbqq_vs_ccqq: largest cut.eps_s over the FLAV_F1 runs ('not reached' when a run's cut is null or leaves fewer than 2 grid points) :: sha256 c6c4138ee5c86423
\newcommand{\MlstRocCutBbqqFoner}{0.32}  % experiments/FIGS/data/v2/index.json :: curves_v2/pairs curve bbqq_vs_ccqq: largest cut.eps_s over the FLAV_F1R runs ('not reached' when a run's cut is null or leaves fewer than 2 grid points) :: sha256 c6c4138ee5c86423
\newcommand{\MlstRocNBkgBbqq}{200{,}000}  % experiments/FIGS/data/v2/index.json :: curves_v2/pairs curve bbqq_vs_ccqq: n_background :: sha256 c6c4138ee5c86423
\newcommand{\MlstRocRejMaxBbqq}{10{,}000}  % experiments/FIGS/data/v2/index.json :: curves_v2/pairs curve bbqq_vs_ccqq: n_background / 20 :: sha256 c6c4138ee5c86423
\newcommand{\MlstRocCutBcFull}{0.24}  % experiments/FIGS/data/v2/index.json :: curves_v2/pairs curve bc_vs_bq_cs: largest cut.eps_s over the L188 runs ('not reached' when a run's cut is null or leaves fewer than 2 grid points) :: sha256 c6c4138ee5c86423
\newcommand{\MlstRocCutBcRes}{0.25}  % experiments/FIGS/data/v2/index.json :: curves_v2/pairs curve bc_vs_bq_cs: largest cut.eps_s over the L162 runs ('not reached' when a run's cut is null or leaves fewer than 2 grid points) :: sha256 c6c4138ee5c86423
\newcommand{\MlstRocCutBcMid}{0.14}  % experiments/FIGS/data/v2/index.json :: curves_v2/pairs curve bc_vs_bq_cs: largest cut.eps_s over the R42_Q1 runs ('not reached' when a run's cut is null or leaves fewer than 2 grid points) :: sha256 c6c4138ee5c86423
\newcommand{\MlstRocCutBcCoarse}{0.00}  % experiments/FIGS/data/v2/index.json :: curves_v2/pairs curve bc_vs_bq_cs: largest cut.eps_s over the R16_Q1 runs ('not reached' when a run's cut is null or leaves fewer than 2 grid points) :: sha256 c6c4138ee5c86423
\newcommand{\MlstRocCutBcFzero}{0.00}  % experiments/FIGS/data/v2/index.json :: curves_v2/pairs curve bc_vs_bq_cs: largest cut.eps_s over the FLAV_F0 runs ('not reached' when a run's cut is null or leaves fewer than 2 grid points) :: sha256 c6c4138ee5c86423
\newcommand{\MlstRocCutBcFone}{0.01}  % experiments/FIGS/data/v2/index.json :: curves_v2/pairs curve bc_vs_bq_cs: largest cut.eps_s over the FLAV_F1 runs ('not reached' when a run's cut is null or leaves fewer than 2 grid points) :: sha256 c6c4138ee5c86423
\newcommand{\MlstRocCutBcFoner}{0.01}  % experiments/FIGS/data/v2/index.json :: curves_v2/pairs curve bc_vs_bq_cs: largest cut.eps_s over the FLAV_F1R runs ('not reached' when a run's cut is null or leaves fewer than 2 grid points) :: sha256 c6c4138ee5c86423
\newcommand{\MlstRocNBkgBc}{200{,}000}  % experiments/FIGS/data/v2/index.json :: curves_v2/pairs curve bc_vs_bq_cs: n_background :: sha256 c6c4138ee5c86423
\newcommand{\MlstRocRejMaxBc}{10{,}000}  % experiments/FIGS/data/v2/index.json :: curves_v2/pairs curve bc_vs_bq_cs: n_background / 20 :: sha256 c6c4138ee5c86423
\newcommand{\MlstRocCutEemmFull}{not reached}  % experiments/FIGS/data/v2/index.json :: curves_v2/pairs curve ee_vs_mm: largest cut.eps_s over the L188 runs ('not reached' when a run's cut is null or leaves fewer than 2 grid points) :: sha256 c6c4138ee5c86423
\newcommand{\MlstRocCutEemmRes}{not reached}  % experiments/FIGS/data/v2/index.json :: curves_v2/pairs curve ee_vs_mm: largest cut.eps_s over the L162 runs ('not reached' when a run's cut is null or leaves fewer than 2 grid points) :: sha256 c6c4138ee5c86423
\newcommand{\MlstRocCutEemmMid}{0.00}  % experiments/FIGS/data/v2/index.json :: curves_v2/pairs curve ee_vs_mm: largest cut.eps_s over the R42_Q1 runs ('not reached' when a run's cut is null or leaves fewer than 2 grid points) :: sha256 c6c4138ee5c86423
\newcommand{\MlstRocCutEemmCoarse}{0.00}  % experiments/FIGS/data/v2/index.json :: curves_v2/pairs curve ee_vs_mm: largest cut.eps_s over the R16_Q1 runs ('not reached' when a run's cut is null or leaves fewer than 2 grid points) :: sha256 c6c4138ee5c86423
\newcommand{\MlstRocCutEemmFzero}{not reached}  % experiments/FIGS/data/v2/index.json :: curves_v2/pairs curve ee_vs_mm: largest cut.eps_s over the FLAV_F0 runs ('not reached' when a run's cut is null or leaves fewer than 2 grid points) :: sha256 c6c4138ee5c86423
\newcommand{\MlstRocCutEemmFone}{not reached}  % experiments/FIGS/data/v2/index.json :: curves_v2/pairs curve ee_vs_mm: largest cut.eps_s over the FLAV_F1 runs ('not reached' when a run's cut is null or leaves fewer than 2 grid points) :: sha256 c6c4138ee5c86423
\newcommand{\MlstRocCutEemmFoner}{not reached}  % experiments/FIGS/data/v2/index.json :: curves_v2/pairs curve ee_vs_mm: largest cut.eps_s over the FLAV_F1R runs ('not reached' when a run's cut is null or leaves fewer than 2 grid points) :: sha256 c6c4138ee5c86423
\newcommand{\MlstRocNBkgEemm}{128{,}833}  % experiments/FIGS/data/v2/index.json :: curves_v2/pairs curve ee_vs_mm: n_background :: sha256 c6c4138ee5c86423
\newcommand{\MlstRocRejMaxEemm}{6{,}442}  % experiments/FIGS/data/v2/index.json :: curves_v2/pairs curve ee_vs_mm: n_background / 20 :: sha256 c6c4138ee5c86423
\newcommand{\MlstRocNoCurveEemm}{the \Kfull- and \Kres-class models and the three random partitions with controlled flavor boundaries}  % experiments/FIGS/data/v2/index.json :: curves_v2/pairs curve ee_vs_mm: configurations of F4 with a run whose cut is null or leaves fewer than 2 grid points (F4 draws no curve for them) :: sha256 c6c4138ee5c86423 :: L188, L162, FLAV_F0, FLAV_F1, FLAV_F1R
\newcommand{\MlstRocCutJcOneHbbFull}{0.40}  % experiments/FIGS/data/v2/index.json :: curves_v2/jc1_ft curve label_Hbb: largest cut.eps_s over the 188-class runs ('not reached' when a run's cut is null or leaves fewer than 2 grid points) :: sha256 c6c4138ee5c86423
\newcommand{\MlstRocCutJcOneHccFull}{0.37}  % experiments/FIGS/data/v2/index.json :: curves_v2/jc1_ft curve label_Hcc: largest cut.eps_s over the 188-class runs ('not reached' when a run's cut is null or leaves fewer than 2 grid points) :: sha256 c6c4138ee5c86423
\newcommand{\MlstRocCutJcOneHggFull}{0.04}  % experiments/FIGS/data/v2/index.json :: curves_v2/jc1_ft curve label_Hgg: largest cut.eps_s over the 188-class runs ('not reached' when a run's cut is null or leaves fewer than 2 grid points) :: sha256 c6c4138ee5c86423
\newcommand{\MlstRocCutJcOneHfourqFull}{0.26}  % experiments/FIGS/data/v2/index.json :: curves_v2/jc1_ft curve label_H4q: largest cut.eps_s over the 188-class runs ('not reached' when a run's cut is null or leaves fewer than 2 grid points) :: sha256 c6c4138ee5c86423
\newcommand{\MlstRocCutJcOneHqqlFull}{0.98}  % experiments/FIGS/data/v2/index.json :: curves_v2/jc1_ft curve label_Hqql: largest cut.eps_s over the 188-class runs ('not reached' when a run's cut is null or leaves fewer than 2 grid points) :: sha256 c6c4138ee5c86423
\newcommand{\MlstRocCutJcOneTbqqFull}{0.53}  % experiments/FIGS/data/v2/index.json :: curves_v2/jc1_ft curve label_Tbqq: largest cut.eps_s over the 188-class runs ('not reached' when a run's cut is null or leaves fewer than 2 grid points) :: sha256 c6c4138ee5c86423
\newcommand{\MlstRocCutJcOneTblFull}{not reached}  % experiments/FIGS/data/v2/index.json :: curves_v2/jc1_ft curve label_Tbl: largest cut.eps_s over the 188-class runs ('not reached' when a run's cut is null or leaves fewer than 2 grid points) :: sha256 c6c4138ee5c86423
\newcommand{\MlstRocCutJcOneWqqFull}{0.15}  % experiments/FIGS/data/v2/index.json :: curves_v2/jc1_ft curve label_Wqq: largest cut.eps_s over the 188-class runs ('not reached' when a run's cut is null or leaves fewer than 2 grid points) :: sha256 c6c4138ee5c86423
\newcommand{\MlstRocCutJcOneZqqFull}{0.12}  % experiments/FIGS/data/v2/index.json :: curves_v2/jc1_ft curve label_Zqq: largest cut.eps_s over the 188-class runs ('not reached' when a run's cut is null or leaves fewer than 2 grid points) :: sha256 c6c4138ee5c86423
\newcommand{\MlstFtTopRfiftyCoarseOverFull}{0.932\,\ensuremath{\pm}\,0.018}  % experiments/FIGS/data/v2/index.json :: curves_v2/top_ft curve top: rejection at eps_S = 0.5 of R16_Q1 over L188, run by run :: sha256 c6c4138ee5c86423
\newcommand{\MlstFtTopRfiftyCoarseOverFullPr}{6.8\% less rejection}  % experiments/FIGS/data/v2/index.json :: curves_v2/top_ft curve top: rejection at eps_S = 0.5 of R16_Q1 over L188, run by run :: sha256 c6c4138ee5c86423 :: prose twin: two significant figures, direction in words
\newcommand{\MlstProbeUnconverged}{0}  % experiments/FIGS/data/v2/index.json :: linear probes of pretrained models that did not converge (count) :: sha256 c6c4138ee5c86423
\newcommand{\MlstProbeUnconvergedInit}{0}  % experiments/FIGS/data/v2/index.json :: linear probes of the randomly initialized networks that did not converge (count) :: sha256 c6c4138ee5c86423
\newcommand{\MlstProbeFits}{895}  % experiments/FIGS/data/v2/index.json :: selected linear-probe fits of the reported checkpoints (count) :: sha256 c6c4138ee5c86423
\newcommand{\MlstProbeEdgeC}{16}  % experiments/FIGS/data/v2/index.json :: of these, fits whose selected C is the largest of the grid (count) :: sha256 c6c4138ee5c86423
\newcommand{\MlstMlpUnconverged}{0}  % experiments/FIGS/data/v2/index.json :: nonlinear probes that did not reach their plateau within the epoch cap (count) :: sha256 c6c4138ee5c86423
\newcommand{\MlstMlpCmpJcTwoResOverFull}{1.019\,\ensuremath{\pm}\,0.007}  % experiments/FIGS/data/v2/index.json :: nonlinear probe jc2_188: (1 - AUC) of L162 over L188, run by run :: sha256 c6c4138ee5c86423
\newcommand{\MlstMlpCmpJcTwoResOverFullPr}{1.9\% more errors}  % experiments/FIGS/data/v2/index.json :: nonlinear probe jc2_188: (1 - AUC) of L162 over L188, run by run :: sha256 c6c4138ee5c86423 :: prose twin: two significant figures, direction in words
\newcommand{\MlstMlpCmpJcTwoMidOverFull}{1.170\,\ensuremath{\pm}\,0.006}  % experiments/FIGS/data/v2/index.json :: nonlinear probe jc2_188: (1 - AUC) of R42_Q1 over L188, run by run :: sha256 c6c4138ee5c86423
\newcommand{\MlstMlpCmpJcTwoMidOverFullPr}{17\% more errors}  % experiments/FIGS/data/v2/index.json :: nonlinear probe jc2_188: (1 - AUC) of R42_Q1 over L188, run by run :: sha256 c6c4138ee5c86423 :: prose twin: two significant figures, direction in words
\newcommand{\MlstMlpCmpJcTwoCoarseOverFull}{1.388\,\ensuremath{\pm}\,0.021}  % experiments/FIGS/data/v2/index.json :: nonlinear probe jc2_188: (1 - AUC) of R16_Q1 over L188, run by run :: sha256 c6c4138ee5c86423
\newcommand{\MlstMlpCmpJcTwoCoarseOverFullPr}{39\% more errors}  % experiments/FIGS/data/v2/index.json :: nonlinear probe jc2_188: (1 - AUC) of R16_Q1 over L188, run by run :: sha256 c6c4138ee5c86423 :: prose twin: two significant figures, direction in words
\newcommand{\MlstMlpCmpJcTwoResMassOverRes}{\pending{no run yet}}  % experiments/FIGS/data/v2/index.json :: nonlinear probe jc2_188: (1 - AUC) of L162_MASS over L162, run by run :: sha256 c6c4138ee5c86423
\newcommand{\MlstMlpCmpJcTwoResMassOverResPr}{\pending{no run yet}}  % experiments/FIGS/data/v2/index.json :: nonlinear probe jc2_188: (1 - AUC) of L162_MASS over L162, run by run :: sha256 c6c4138ee5c86423 :: prose twin: two significant figures, direction in words
\newcommand{\MlstMlpCmpJcTwoResMassOverResPrMax}{\pending{no run yet}}  % experiments/FIGS/data/v2/index.json :: nonlinear probe jc2_188: (1 - AUC) of L162_MASS over L162, run by run :: sha256 c6c4138ee5c86423 :: prose twin of the run farther from 1
\newcommand{\MlstMlpCmpJcTwoCoarseMassOverCoarse}{\pending{no run yet}}  % experiments/FIGS/data/v2/index.json :: nonlinear probe jc2_188: (1 - AUC) of R16_Q1_MASS over R16_Q1, run by run :: sha256 c6c4138ee5c86423
\newcommand{\MlstMlpCmpJcTwoCoarseMassOverCoarsePr}{\pending{no run yet}}  % experiments/FIGS/data/v2/index.json :: nonlinear probe jc2_188: (1 - AUC) of R16_Q1_MASS over R16_Q1, run by run :: sha256 c6c4138ee5c86423 :: prose twin: two significant figures, direction in words
\newcommand{\MlstMlpCmpJcTwoCoarseMassOverCoarsePrMax}{\pending{no run yet}}  % experiments/FIGS/data/v2/index.json :: nonlinear probe jc2_188: (1 - AUC) of R16_Q1_MASS over R16_Q1, run by run :: sha256 c6c4138ee5c86423 :: prose twin of the run farther from 1
\newcommand{\MlstMlpCmpJcTwoFzeroOverFull}{1.294 and 1.310}  % experiments/FIGS/data/v2/index.json :: nonlinear probe jc2_188: (1 - AUC) of FLAV_F0 over L188, run by run :: sha256 c6c4138ee5c86423
\newcommand{\MlstMlpCmpJcTwoFzeroOverFullPr}{at least 29\% more errors}  % experiments/FIGS/data/v2/index.json :: nonlinear probe jc2_188: (1 - AUC) of FLAV_F0 over L188, run by run :: sha256 c6c4138ee5c86423 :: prose twin: two significant figures, direction in words
\newcommand{\MlstMlpCmpJcTwoFzeroOverFullPrMax}{up to 31\% more errors}  % experiments/FIGS/data/v2/index.json :: nonlinear probe jc2_188: (1 - AUC) of FLAV_F0 over L188, run by run :: sha256 c6c4138ee5c86423 :: prose twin of the run farther from 1
\newcommand{\MlstMlpCmpJcTwoFoneOverFzero}{0.903 and 0.885}  % experiments/FIGS/data/v2/index.json :: nonlinear probe jc2_188: (1 - AUC) of FLAV_F1 over FLAV_F0, run by run :: sha256 c6c4138ee5c86423
\newcommand{\MlstMlpCmpJcTwoFoneOverFzeroPr}{at least 9.7\% fewer errors}  % experiments/FIGS/data/v2/index.json :: nonlinear probe jc2_188: (1 - AUC) of FLAV_F1 over FLAV_F0, run by run :: sha256 c6c4138ee5c86423 :: prose twin: two significant figures, direction in words
\newcommand{\MlstMlpCmpJcTwoFoneOverFzeroPrMax}{up to 12\% fewer errors}  % experiments/FIGS/data/v2/index.json :: nonlinear probe jc2_188: (1 - AUC) of FLAV_F1 over FLAV_F0, run by run :: sha256 c6c4138ee5c86423 :: prose twin of the run farther from 1
\newcommand{\MlstMlpCmpJcTwoFonerOverFone}{0.9916 and 0.9918}  % experiments/FIGS/data/v2/index.json :: nonlinear probe jc2_188: (1 - AUC) of FLAV_F1R over FLAV_F1, run by run :: sha256 c6c4138ee5c86423
\newcommand{\MlstMlpCmpJcTwoFonerOverFonePr}{equal within 1\%}  % experiments/FIGS/data/v2/index.json :: nonlinear probe jc2_188: (1 - AUC) of FLAV_F1R over FLAV_F1, run by run :: sha256 c6c4138ee5c86423 :: prose twin: two significant figures, direction in words
\newcommand{\MlstMlpCmpJcTwoFonerOverFonePrMax}{equal within 1\%}  % experiments/FIGS/data/v2/index.json :: nonlinear probe jc2_188: (1 - AUC) of FLAV_F1R over FLAV_F1, run by run :: sha256 c6c4138ee5c86423 :: prose twin of the run farther from 1
\newcommand{\MlstMlpCmpJcTwoFonerOverFzero}{0.895 and 0.877}  % experiments/FIGS/data/v2/index.json :: nonlinear probe jc2_188: (1 - AUC) of FLAV_F1R over FLAV_F0, run by run :: sha256 c6c4138ee5c86423
\newcommand{\MlstMlpCmpJcTwoFonerOverFzeroPr}{at least 10\% fewer errors}  % experiments/FIGS/data/v2/index.json :: nonlinear probe jc2_188: (1 - AUC) of FLAV_F1R over FLAV_F0, run by run :: sha256 c6c4138ee5c86423 :: prose twin: two significant figures, direction in words
\newcommand{\MlstMlpCmpJcTwoFonerOverFzeroPrMax}{up to 12\% fewer errors}  % experiments/FIGS/data/v2/index.json :: nonlinear probe jc2_188: (1 - AUC) of FLAV_F1R over FLAV_F0, run by run :: sha256 c6c4138ee5c86423 :: prose twin of the run farther from 1
\newcommand{\MlstMlpCmpJcTwoFzeroOverCoarse}{0.9461 and 0.9456}  % experiments/FIGS/data/v2/index.json :: nonlinear probe jc2_188: (1 - AUC) of FLAV_F0 over R16_Q1, run by run :: sha256 c6c4138ee5c86423
\newcommand{\MlstMlpCmpJcTwoFzeroOverCoarsePr}{at least 5.4\% fewer errors}  % experiments/FIGS/data/v2/index.json :: nonlinear probe jc2_188: (1 - AUC) of FLAV_F0 over R16_Q1, run by run :: sha256 c6c4138ee5c86423 :: prose twin: two significant figures, direction in words
\newcommand{\MlstMlpCmpJcTwoFzeroOverCoarsePrMax}{up to 5.4\% fewer errors}  % experiments/FIGS/data/v2/index.json :: nonlinear probe jc2_188: (1 - AUC) of FLAV_F0 over R16_Q1, run by run :: sha256 c6c4138ee5c86423 :: prose twin of the run farther from 1
\newcommand{\MlstMlpCmpJcTwoRandOverFull}{\pending{no run yet}}  % experiments/FIGS/data/v2/index.json :: nonlinear probe jc2_188: (1 - AUC) of each random partition over the 188-class run of its seed :: sha256 c6c4138ee5c86423
\newcommand{\MlstMlpCmpJcTwoRandOverFullPr}{\pending{no run yet}}  % experiments/FIGS/data/v2/index.json :: nonlinear probe jc2_188: (1 - AUC) of each random partition over the 188-class run of its seed :: sha256 c6c4138ee5c86423 :: prose twin: two significant figures, direction in words
\newcommand{\MlstMlpCmpJcTwoResResOverFull}{1.001\,\ensuremath{\pm}\,0.007}  % experiments/FIGS/data/v2/index.json :: nonlinear probe jc2_188_res: (1 - AUC) of L162 over L188, run by run :: sha256 c6c4138ee5c86423
\newcommand{\MlstMlpCmpJcTwoResResOverFullPr}{equal within 1\%}  % experiments/FIGS/data/v2/index.json :: nonlinear probe jc2_188_res: (1 - AUC) of L162 over L188, run by run :: sha256 c6c4138ee5c86423 :: prose twin: two significant figures, direction in words
\newcommand{\MlstMlpCmpJcTwoResMidOverFull}{1.192\,\ensuremath{\pm}\,0.008}  % experiments/FIGS/data/v2/index.json :: nonlinear probe jc2_188_res: (1 - AUC) of R42_Q1 over L188, run by run :: sha256 c6c4138ee5c86423
\newcommand{\MlstMlpCmpJcTwoResMidOverFullPr}{19\% more errors}  % experiments/FIGS/data/v2/index.json :: nonlinear probe jc2_188_res: (1 - AUC) of R42_Q1 over L188, run by run :: sha256 c6c4138ee5c86423 :: prose twin: two significant figures, direction in words
\newcommand{\MlstMlpCmpJcTwoResCoarseOverFull}{1.408\,\ensuremath{\pm}\,0.026}  % experiments/FIGS/data/v2/index.json :: nonlinear probe jc2_188_res: (1 - AUC) of R16_Q1 over L188, run by run :: sha256 c6c4138ee5c86423
\newcommand{\MlstMlpCmpJcTwoResCoarseOverFullPr}{41\% more errors}  % experiments/FIGS/data/v2/index.json :: nonlinear probe jc2_188_res: (1 - AUC) of R16_Q1 over L188, run by run :: sha256 c6c4138ee5c86423 :: prose twin: two significant figures, direction in words
\newcommand{\MlstMlpCmpJcTwoResResMassOverRes}{\pending{no run yet}}  % experiments/FIGS/data/v2/index.json :: nonlinear probe jc2_188_res: (1 - AUC) of L162_MASS over L162, run by run :: sha256 c6c4138ee5c86423
\newcommand{\MlstMlpCmpJcTwoResResMassOverResPr}{\pending{no run yet}}  % experiments/FIGS/data/v2/index.json :: nonlinear probe jc2_188_res: (1 - AUC) of L162_MASS over L162, run by run :: sha256 c6c4138ee5c86423 :: prose twin: two significant figures, direction in words
\newcommand{\MlstMlpCmpJcTwoResResMassOverResPrMax}{\pending{no run yet}}  % experiments/FIGS/data/v2/index.json :: nonlinear probe jc2_188_res: (1 - AUC) of L162_MASS over L162, run by run :: sha256 c6c4138ee5c86423 :: prose twin of the run farther from 1
\newcommand{\MlstMlpCmpJcTwoResCoarseMassOverCoarse}{\pending{no run yet}}  % experiments/FIGS/data/v2/index.json :: nonlinear probe jc2_188_res: (1 - AUC) of R16_Q1_MASS over R16_Q1, run by run :: sha256 c6c4138ee5c86423
\newcommand{\MlstMlpCmpJcTwoResCoarseMassOverCoarsePr}{\pending{no run yet}}  % experiments/FIGS/data/v2/index.json :: nonlinear probe jc2_188_res: (1 - AUC) of R16_Q1_MASS over R16_Q1, run by run :: sha256 c6c4138ee5c86423 :: prose twin: two significant figures, direction in words
\newcommand{\MlstMlpCmpJcTwoResCoarseMassOverCoarsePrMax}{\pending{no run yet}}  % experiments/FIGS/data/v2/index.json :: nonlinear probe jc2_188_res: (1 - AUC) of R16_Q1_MASS over R16_Q1, run by run :: sha256 c6c4138ee5c86423 :: prose twin of the run farther from 1
\newcommand{\MlstMlpCmpJcTwoResFzeroOverFull}{1.285 and 1.300}  % experiments/FIGS/data/v2/index.json :: nonlinear probe jc2_188_res: (1 - AUC) of FLAV_F0 over L188, run by run :: sha256 c6c4138ee5c86423
\newcommand{\MlstMlpCmpJcTwoResFzeroOverFullPr}{at least 29\% more errors}  % experiments/FIGS/data/v2/index.json :: nonlinear probe jc2_188_res: (1 - AUC) of FLAV_F0 over L188, run by run :: sha256 c6c4138ee5c86423 :: prose twin: two significant figures, direction in words
\newcommand{\MlstMlpCmpJcTwoResFzeroOverFullPrMax}{up to 30\% more errors}  % experiments/FIGS/data/v2/index.json :: nonlinear probe jc2_188_res: (1 - AUC) of FLAV_F0 over L188, run by run :: sha256 c6c4138ee5c86423 :: prose twin of the run farther from 1
\newcommand{\MlstMlpCmpJcTwoResFoneOverFzero}{0.909 and 0.889}  % experiments/FIGS/data/v2/index.json :: nonlinear probe jc2_188_res: (1 - AUC) of FLAV_F1 over FLAV_F0, run by run :: sha256 c6c4138ee5c86423
\newcommand{\MlstMlpCmpJcTwoResFoneOverFzeroPr}{at least 9.1\% fewer errors}  % experiments/FIGS/data/v2/index.json :: nonlinear probe jc2_188_res: (1 - AUC) of FLAV_F1 over FLAV_F0, run by run :: sha256 c6c4138ee5c86423 :: prose twin: two significant figures, direction in words
\newcommand{\MlstMlpCmpJcTwoResFoneOverFzeroPrMax}{up to 11\% fewer errors}  % experiments/FIGS/data/v2/index.json :: nonlinear probe jc2_188_res: (1 - AUC) of FLAV_F1 over FLAV_F0, run by run :: sha256 c6c4138ee5c86423 :: prose twin of the run farther from 1
\newcommand{\MlstMlpCmpJcTwoResFonerOverFone}{0.9793 and 0.9825}  % experiments/FIGS/data/v2/index.json :: nonlinear probe jc2_188_res: (1 - AUC) of FLAV_F1R over FLAV_F1, run by run :: sha256 c6c4138ee5c86423
\newcommand{\MlstMlpCmpJcTwoResFonerOverFonePr}{equal within 3\%}  % experiments/FIGS/data/v2/index.json :: nonlinear probe jc2_188_res: (1 - AUC) of FLAV_F1R over FLAV_F1, run by run :: sha256 c6c4138ee5c86423 :: prose twin: two significant figures, direction in words
\newcommand{\MlstMlpCmpJcTwoResFonerOverFonePrMax}{equal within 3\%}  % experiments/FIGS/data/v2/index.json :: nonlinear probe jc2_188_res: (1 - AUC) of FLAV_F1R over FLAV_F1, run by run :: sha256 c6c4138ee5c86423 :: prose twin of the run farther from 1
\newcommand{\MlstMlpCmpJcTwoResFonerOverFzero}{0.890 and 0.874}  % experiments/FIGS/data/v2/index.json :: nonlinear probe jc2_188_res: (1 - AUC) of FLAV_F1R over FLAV_F0, run by run :: sha256 c6c4138ee5c86423
\newcommand{\MlstMlpCmpJcTwoResFonerOverFzeroPr}{at least 11\% fewer errors}  % experiments/FIGS/data/v2/index.json :: nonlinear probe jc2_188_res: (1 - AUC) of FLAV_F1R over FLAV_F0, run by run :: sha256 c6c4138ee5c86423 :: prose twin: two significant figures, direction in words
\newcommand{\MlstMlpCmpJcTwoResFonerOverFzeroPrMax}{up to 13\% fewer errors}  % experiments/FIGS/data/v2/index.json :: nonlinear probe jc2_188_res: (1 - AUC) of FLAV_F1R over FLAV_F0, run by run :: sha256 c6c4138ee5c86423 :: prose twin of the run farther from 1
\newcommand{\MlstMlpCmpJcTwoResFzeroOverCoarse}{0.9299 and 0.9233}  % experiments/FIGS/data/v2/index.json :: nonlinear probe jc2_188_res: (1 - AUC) of FLAV_F0 over R16_Q1, run by run :: sha256 c6c4138ee5c86423
\newcommand{\MlstMlpCmpJcTwoResFzeroOverCoarsePr}{at least 7.0\% fewer errors}  % experiments/FIGS/data/v2/index.json :: nonlinear probe jc2_188_res: (1 - AUC) of FLAV_F0 over R16_Q1, run by run :: sha256 c6c4138ee5c86423 :: prose twin: two significant figures, direction in words
\newcommand{\MlstMlpCmpJcTwoResFzeroOverCoarsePrMax}{up to 7.7\% fewer errors}  % experiments/FIGS/data/v2/index.json :: nonlinear probe jc2_188_res: (1 - AUC) of FLAV_F0 over R16_Q1, run by run :: sha256 c6c4138ee5c86423 :: prose twin of the run farther from 1
\newcommand{\MlstMlpCmpJcTwoResRandOverFull}{\pending{no run yet}}  % experiments/FIGS/data/v2/index.json :: nonlinear probe jc2_188_res: (1 - AUC) of each random partition over the 188-class run of its seed :: sha256 c6c4138ee5c86423
\newcommand{\MlstMlpCmpJcTwoResRandOverFullPr}{\pending{no run yet}}  % experiments/FIGS/data/v2/index.json :: nonlinear probe jc2_188_res: (1 - AUC) of each random partition over the 188-class run of its seed :: sha256 c6c4138ee5c86423 :: prose twin: two significant figures, direction in words
\newcommand{\MlstMlpCmpJcTwoQcdResOverFull}{1.070\,\ensuremath{\pm}\,0.006}  % experiments/FIGS/data/v2/index.json :: nonlinear probe jc2_188_qcd: (1 - AUC) of L162 over L188, run by run :: sha256 c6c4138ee5c86423
\newcommand{\MlstMlpCmpJcTwoQcdResOverFullPr}{7.0\% more errors}  % experiments/FIGS/data/v2/index.json :: nonlinear probe jc2_188_qcd: (1 - AUC) of L162 over L188, run by run :: sha256 c6c4138ee5c86423 :: prose twin: two significant figures, direction in words
\newcommand{\MlstMlpCmpJcTwoQcdMidOverFull}{1.107\,\ensuremath{\pm}\,0.002}  % experiments/FIGS/data/v2/index.json :: nonlinear probe jc2_188_qcd: (1 - AUC) of R42_Q1 over L188, run by run :: sha256 c6c4138ee5c86423
\newcommand{\MlstMlpCmpJcTwoQcdMidOverFullPr}{11\% more errors}  % experiments/FIGS/data/v2/index.json :: nonlinear probe jc2_188_qcd: (1 - AUC) of R42_Q1 over L188, run by run :: sha256 c6c4138ee5c86423 :: prose twin: two significant figures, direction in words
\newcommand{\MlstMlpCmpJcTwoQcdCoarseOverFull}{1.333\,\ensuremath{\pm}\,0.012}  % experiments/FIGS/data/v2/index.json :: nonlinear probe jc2_188_qcd: (1 - AUC) of R16_Q1 over L188, run by run :: sha256 c6c4138ee5c86423
\newcommand{\MlstMlpCmpJcTwoQcdCoarseOverFullPr}{33\% more errors}  % experiments/FIGS/data/v2/index.json :: nonlinear probe jc2_188_qcd: (1 - AUC) of R16_Q1 over L188, run by run :: sha256 c6c4138ee5c86423 :: prose twin: two significant figures, direction in words
\newcommand{\MlstMlpCmpJcTwoQcdResMassOverRes}{\pending{no run yet}}  % experiments/FIGS/data/v2/index.json :: nonlinear probe jc2_188_qcd: (1 - AUC) of L162_MASS over L162, run by run :: sha256 c6c4138ee5c86423
\newcommand{\MlstMlpCmpJcTwoQcdResMassOverResPr}{\pending{no run yet}}  % experiments/FIGS/data/v2/index.json :: nonlinear probe jc2_188_qcd: (1 - AUC) of L162_MASS over L162, run by run :: sha256 c6c4138ee5c86423 :: prose twin: two significant figures, direction in words
\newcommand{\MlstMlpCmpJcTwoQcdResMassOverResPrMax}{\pending{no run yet}}  % experiments/FIGS/data/v2/index.json :: nonlinear probe jc2_188_qcd: (1 - AUC) of L162_MASS over L162, run by run :: sha256 c6c4138ee5c86423 :: prose twin of the run farther from 1
\newcommand{\MlstMlpCmpJcTwoQcdCoarseMassOverCoarse}{\pending{no run yet}}  % experiments/FIGS/data/v2/index.json :: nonlinear probe jc2_188_qcd: (1 - AUC) of R16_Q1_MASS over R16_Q1, run by run :: sha256 c6c4138ee5c86423
\newcommand{\MlstMlpCmpJcTwoQcdCoarseMassOverCoarsePr}{\pending{no run yet}}  % experiments/FIGS/data/v2/index.json :: nonlinear probe jc2_188_qcd: (1 - AUC) of R16_Q1_MASS over R16_Q1, run by run :: sha256 c6c4138ee5c86423 :: prose twin: two significant figures, direction in words
\newcommand{\MlstMlpCmpJcTwoQcdCoarseMassOverCoarsePrMax}{\pending{no run yet}}  % experiments/FIGS/data/v2/index.json :: nonlinear probe jc2_188_qcd: (1 - AUC) of R16_Q1_MASS over R16_Q1, run by run :: sha256 c6c4138ee5c86423 :: prose twin of the run farther from 1
\newcommand{\MlstMlpCmpJcTwoQcdFzeroOverFull}{1.319 and 1.339}  % experiments/FIGS/data/v2/index.json :: nonlinear probe jc2_188_qcd: (1 - AUC) of FLAV_F0 over L188, run by run :: sha256 c6c4138ee5c86423
\newcommand{\MlstMlpCmpJcTwoQcdFzeroOverFullPr}{at least 32\% more errors}  % experiments/FIGS/data/v2/index.json :: nonlinear probe jc2_188_qcd: (1 - AUC) of FLAV_F0 over L188, run by run :: sha256 c6c4138ee5c86423 :: prose twin: two significant figures, direction in words
\newcommand{\MlstMlpCmpJcTwoQcdFzeroOverFullPrMax}{up to 34\% more errors}  % experiments/FIGS/data/v2/index.json :: nonlinear probe jc2_188_qcd: (1 - AUC) of FLAV_F0 over L188, run by run :: sha256 c6c4138ee5c86423 :: prose twin of the run farther from 1
\newcommand{\MlstMlpCmpJcTwoQcdFoneOverFzero}{0.886 and 0.872}  % experiments/FIGS/data/v2/index.json :: nonlinear probe jc2_188_qcd: (1 - AUC) of FLAV_F1 over FLAV_F0, run by run :: sha256 c6c4138ee5c86423
\newcommand{\MlstMlpCmpJcTwoQcdFoneOverFzeroPr}{at least 11\% fewer errors}  % experiments/FIGS/data/v2/index.json :: nonlinear probe jc2_188_qcd: (1 - AUC) of FLAV_F1 over FLAV_F0, run by run :: sha256 c6c4138ee5c86423 :: prose twin: two significant figures, direction in words
\newcommand{\MlstMlpCmpJcTwoQcdFoneOverFzeroPrMax}{up to 13\% fewer errors}  % experiments/FIGS/data/v2/index.json :: nonlinear probe jc2_188_qcd: (1 - AUC) of FLAV_F1 over FLAV_F0, run by run :: sha256 c6c4138ee5c86423 :: prose twin of the run farther from 1
\newcommand{\MlstMlpCmpJcTwoQcdFonerOverFone}{1.026 and 1.018}  % experiments/FIGS/data/v2/index.json :: nonlinear probe jc2_188_qcd: (1 - AUC) of FLAV_F1R over FLAV_F1, run by run :: sha256 c6c4138ee5c86423
\newcommand{\MlstMlpCmpJcTwoQcdFonerOverFonePr}{at least 1.8\% more errors}  % experiments/FIGS/data/v2/index.json :: nonlinear probe jc2_188_qcd: (1 - AUC) of FLAV_F1R over FLAV_F1, run by run :: sha256 c6c4138ee5c86423 :: prose twin: two significant figures, direction in words
\newcommand{\MlstMlpCmpJcTwoQcdFonerOverFonePrMax}{up to 2.6\% more errors}  % experiments/FIGS/data/v2/index.json :: nonlinear probe jc2_188_qcd: (1 - AUC) of FLAV_F1R over FLAV_F1, run by run :: sha256 c6c4138ee5c86423 :: prose twin of the run farther from 1
\newcommand{\MlstMlpCmpJcTwoQcdFonerOverFzero}{0.909 and 0.887}  % experiments/FIGS/data/v2/index.json :: nonlinear probe jc2_188_qcd: (1 - AUC) of FLAV_F1R over FLAV_F0, run by run :: sha256 c6c4138ee5c86423
\newcommand{\MlstMlpCmpJcTwoQcdFonerOverFzeroPr}{at least 9.1\% fewer errors}  % experiments/FIGS/data/v2/index.json :: nonlinear probe jc2_188_qcd: (1 - AUC) of FLAV_F1R over FLAV_F0, run by run :: sha256 c6c4138ee5c86423 :: prose twin: two significant figures, direction in words
\newcommand{\MlstMlpCmpJcTwoQcdFonerOverFzeroPrMax}{up to 11\% fewer errors}  % experiments/FIGS/data/v2/index.json :: nonlinear probe jc2_188_qcd: (1 - AUC) of FLAV_F1R over FLAV_F0, run by run :: sha256 c6c4138ee5c86423 :: prose twin of the run farther from 1
\newcommand{\MlstMlpCmpJcTwoQcdFzeroOverCoarse}{0.993 and 1.012}  % experiments/FIGS/data/v2/index.json :: nonlinear probe jc2_188_qcd: (1 - AUC) of FLAV_F0 over R16_Q1, run by run :: sha256 c6c4138ee5c86423
\newcommand{\MlstMlpCmpJcTwoQcdFzeroOverCoarsePr}{equal within 2\%}  % experiments/FIGS/data/v2/index.json :: nonlinear probe jc2_188_qcd: (1 - AUC) of FLAV_F0 over R16_Q1, run by run :: sha256 c6c4138ee5c86423 :: prose twin: two significant figures, direction in words
\newcommand{\MlstMlpCmpJcTwoQcdFzeroOverCoarsePrMax}{equal within 2\%}  % experiments/FIGS/data/v2/index.json :: nonlinear probe jc2_188_qcd: (1 - AUC) of FLAV_F0 over R16_Q1, run by run :: sha256 c6c4138ee5c86423 :: prose twin of the run farther from 1
\newcommand{\MlstMlpCmpJcTwoQcdRandOverFull}{\pending{no run yet}}  % experiments/FIGS/data/v2/index.json :: nonlinear probe jc2_188_qcd: (1 - AUC) of each random partition over the 188-class run of its seed :: sha256 c6c4138ee5c86423
\newcommand{\MlstMlpCmpJcTwoQcdRandOverFullPr}{\pending{no run yet}}  % experiments/FIGS/data/v2/index.json :: nonlinear probe jc2_188_qcd: (1 - AUC) of each random partition over the 188-class run of its seed :: sha256 c6c4138ee5c86423 :: prose twin: two significant figures, direction in words
\newcommand{\MlstMlpCmpBcResOverFull}{1.002\,\ensuremath{\pm}\,0.006}  % experiments/FIGS/data/v2/index.json :: nonlinear probe bc_vs_bq_cs: (1 - AUC) of L162 over L188, run by run :: sha256 c6c4138ee5c86423
\newcommand{\MlstMlpCmpBcResOverFullPr}{equal within 1\%}  % experiments/FIGS/data/v2/index.json :: nonlinear probe bc_vs_bq_cs: (1 - AUC) of L162 over L188, run by run :: sha256 c6c4138ee5c86423 :: prose twin: two significant figures, direction in words
\newcommand{\MlstMlpCmpBcMidOverFull}{1.104\,\ensuremath{\pm}\,0.017}  % experiments/FIGS/data/v2/index.json :: nonlinear probe bc_vs_bq_cs: (1 - AUC) of R42_Q1 over L188, run by run :: sha256 c6c4138ee5c86423
\newcommand{\MlstMlpCmpBcMidOverFullPr}{10\% more errors}  % experiments/FIGS/data/v2/index.json :: nonlinear probe bc_vs_bq_cs: (1 - AUC) of R42_Q1 over L188, run by run :: sha256 c6c4138ee5c86423 :: prose twin: two significant figures, direction in words
\newcommand{\MlstMlpCmpBcCoarseOverFull}{1.624\,\ensuremath{\pm}\,0.022}  % experiments/FIGS/data/v2/index.json :: nonlinear probe bc_vs_bq_cs: (1 - AUC) of R16_Q1 over L188, run by run :: sha256 c6c4138ee5c86423
\newcommand{\MlstMlpCmpBcCoarseOverFullPr}{62\% more errors}  % experiments/FIGS/data/v2/index.json :: nonlinear probe bc_vs_bq_cs: (1 - AUC) of R16_Q1 over L188, run by run :: sha256 c6c4138ee5c86423 :: prose twin: two significant figures, direction in words
\newcommand{\MlstMlpCmpBcResMassOverRes}{\pending{no run yet}}  % experiments/FIGS/data/v2/index.json :: nonlinear probe bc_vs_bq_cs: (1 - AUC) of L162_MASS over L162, run by run :: sha256 c6c4138ee5c86423
\newcommand{\MlstMlpCmpBcResMassOverResPr}{\pending{no run yet}}  % experiments/FIGS/data/v2/index.json :: nonlinear probe bc_vs_bq_cs: (1 - AUC) of L162_MASS over L162, run by run :: sha256 c6c4138ee5c86423 :: prose twin: two significant figures, direction in words
\newcommand{\MlstMlpCmpBcResMassOverResPrMax}{\pending{no run yet}}  % experiments/FIGS/data/v2/index.json :: nonlinear probe bc_vs_bq_cs: (1 - AUC) of L162_MASS over L162, run by run :: sha256 c6c4138ee5c86423 :: prose twin of the run farther from 1
\newcommand{\MlstMlpCmpBcCoarseMassOverCoarse}{\pending{no run yet}}  % experiments/FIGS/data/v2/index.json :: nonlinear probe bc_vs_bq_cs: (1 - AUC) of R16_Q1_MASS over R16_Q1, run by run :: sha256 c6c4138ee5c86423
\newcommand{\MlstMlpCmpBcCoarseMassOverCoarsePr}{\pending{no run yet}}  % experiments/FIGS/data/v2/index.json :: nonlinear probe bc_vs_bq_cs: (1 - AUC) of R16_Q1_MASS over R16_Q1, run by run :: sha256 c6c4138ee5c86423 :: prose twin: two significant figures, direction in words
\newcommand{\MlstMlpCmpBcCoarseMassOverCoarsePrMax}{\pending{no run yet}}  % experiments/FIGS/data/v2/index.json :: nonlinear probe bc_vs_bq_cs: (1 - AUC) of R16_Q1_MASS over R16_Q1, run by run :: sha256 c6c4138ee5c86423 :: prose twin of the run farther from 1
\newcommand{\MlstMlpCmpBcFzeroOverFull}{1.633 and 1.679}  % experiments/FIGS/data/v2/index.json :: nonlinear probe bc_vs_bq_cs: (1 - AUC) of FLAV_F0 over L188, run by run :: sha256 c6c4138ee5c86423
\newcommand{\MlstMlpCmpBcFzeroOverFullPr}{at least 63\% more errors}  % experiments/FIGS/data/v2/index.json :: nonlinear probe bc_vs_bq_cs: (1 - AUC) of FLAV_F0 over L188, run by run :: sha256 c6c4138ee5c86423 :: prose twin: two significant figures, direction in words
\newcommand{\MlstMlpCmpBcFzeroOverFullPrMax}{up to 68\% more errors}  % experiments/FIGS/data/v2/index.json :: nonlinear probe bc_vs_bq_cs: (1 - AUC) of FLAV_F0 over L188, run by run :: sha256 c6c4138ee5c86423 :: prose twin of the run farther from 1
\newcommand{\MlstMlpCmpBcFoneOverFzero}{0.848 and 0.810}  % experiments/FIGS/data/v2/index.json :: nonlinear probe bc_vs_bq_cs: (1 - AUC) of FLAV_F1 over FLAV_F0, run by run :: sha256 c6c4138ee5c86423
\newcommand{\MlstMlpCmpBcFoneOverFzeroPr}{at least 15\% fewer errors}  % experiments/FIGS/data/v2/index.json :: nonlinear probe bc_vs_bq_cs: (1 - AUC) of FLAV_F1 over FLAV_F0, run by run :: sha256 c6c4138ee5c86423 :: prose twin: two significant figures, direction in words
\newcommand{\MlstMlpCmpBcFoneOverFzeroPrMax}{up to 19\% fewer errors}  % experiments/FIGS/data/v2/index.json :: nonlinear probe bc_vs_bq_cs: (1 - AUC) of FLAV_F1 over FLAV_F0, run by run :: sha256 c6c4138ee5c86423 :: prose twin of the run farther from 1
\newcommand{\MlstMlpCmpBcFonerOverFone}{0.9812 and 0.9815}  % experiments/FIGS/data/v2/index.json :: nonlinear probe bc_vs_bq_cs: (1 - AUC) of FLAV_F1R over FLAV_F1, run by run :: sha256 c6c4138ee5c86423
\newcommand{\MlstMlpCmpBcFonerOverFonePr}{equal within 2\%}  % experiments/FIGS/data/v2/index.json :: nonlinear probe bc_vs_bq_cs: (1 - AUC) of FLAV_F1R over FLAV_F1, run by run :: sha256 c6c4138ee5c86423 :: prose twin: two significant figures, direction in words
\newcommand{\MlstMlpCmpBcFonerOverFonePrMax}{equal within 2\%}  % experiments/FIGS/data/v2/index.json :: nonlinear probe bc_vs_bq_cs: (1 - AUC) of FLAV_F1R over FLAV_F1, run by run :: sha256 c6c4138ee5c86423 :: prose twin of the run farther from 1
\newcommand{\MlstMlpCmpBcFonerOverFzero}{0.833 and 0.795}  % experiments/FIGS/data/v2/index.json :: nonlinear probe bc_vs_bq_cs: (1 - AUC) of FLAV_F1R over FLAV_F0, run by run :: sha256 c6c4138ee5c86423
\newcommand{\MlstMlpCmpBcFonerOverFzeroPr}{at least 17\% fewer errors}  % experiments/FIGS/data/v2/index.json :: nonlinear probe bc_vs_bq_cs: (1 - AUC) of FLAV_F1R over FLAV_F0, run by run :: sha256 c6c4138ee5c86423 :: prose twin: two significant figures, direction in words
\newcommand{\MlstMlpCmpBcFonerOverFzeroPrMax}{up to 21\% fewer errors}  % experiments/FIGS/data/v2/index.json :: nonlinear probe bc_vs_bq_cs: (1 - AUC) of FLAV_F1R over FLAV_F0, run by run :: sha256 c6c4138ee5c86423 :: prose twin of the run farther from 1
\newcommand{\MlstMlpCmpBcFzeroOverCoarse}{0.994 and 1.030}  % experiments/FIGS/data/v2/index.json :: nonlinear probe bc_vs_bq_cs: (1 - AUC) of FLAV_F0 over R16_Q1, run by run :: sha256 c6c4138ee5c86423
\newcommand{\MlstMlpCmpBcFzeroOverCoarsePr}{equal within 3\%}  % experiments/FIGS/data/v2/index.json :: nonlinear probe bc_vs_bq_cs: (1 - AUC) of FLAV_F0 over R16_Q1, run by run :: sha256 c6c4138ee5c86423 :: prose twin: two significant figures, direction in words
\newcommand{\MlstMlpCmpBcFzeroOverCoarsePrMax}{equal within 3\%}  % experiments/FIGS/data/v2/index.json :: nonlinear probe bc_vs_bq_cs: (1 - AUC) of FLAV_F0 over R16_Q1, run by run :: sha256 c6c4138ee5c86423 :: prose twin of the run farther from 1
\newcommand{\MlstMlpCmpBcRandOverFull}{\pending{no run yet}}  % experiments/FIGS/data/v2/index.json :: nonlinear probe bc_vs_bq_cs: (1 - AUC) of each random partition over the 188-class run of its seed :: sha256 c6c4138ee5c86423
\newcommand{\MlstMlpCmpBcRandOverFullPr}{\pending{no run yet}}  % experiments/FIGS/data/v2/index.json :: nonlinear probe bc_vs_bq_cs: (1 - AUC) of each random partition over the 188-class run of its seed :: sha256 c6c4138ee5c86423 :: prose twin: two significant figures, direction in words
\newcommand{\MlstMlpCmpBbResOverFull}{1.023\,\ensuremath{\pm}\,0.021}  % experiments/FIGS/data/v2/index.json :: nonlinear probe bb_vs_cc: (1 - AUC) of L162 over L188, run by run :: sha256 c6c4138ee5c86423
\newcommand{\MlstMlpCmpBbResOverFullPr}{equal within 3\%}  % experiments/FIGS/data/v2/index.json :: nonlinear probe bb_vs_cc: (1 - AUC) of L162 over L188, run by run :: sha256 c6c4138ee5c86423 :: prose twin: two significant figures, direction in words
\newcommand{\MlstMlpCmpBbMidOverFull}{1.125\,\ensuremath{\pm}\,0.011}  % experiments/FIGS/data/v2/index.json :: nonlinear probe bb_vs_cc: (1 - AUC) of R42_Q1 over L188, run by run :: sha256 c6c4138ee5c86423
\newcommand{\MlstMlpCmpBbMidOverFullPr}{13\% more errors}  % experiments/FIGS/data/v2/index.json :: nonlinear probe bb_vs_cc: (1 - AUC) of R42_Q1 over L188, run by run :: sha256 c6c4138ee5c86423 :: prose twin: two significant figures, direction in words
\newcommand{\MlstMlpCmpBbCoarseOverFull}{3.04\,\ensuremath{\pm}\,0.09}  % experiments/FIGS/data/v2/index.json :: nonlinear probe bb_vs_cc: (1 - AUC) of R16_Q1 over L188, run by run :: sha256 c6c4138ee5c86423
\newcommand{\MlstMlpCmpBbCoarseOverFullPr}{3.0 times the errors}  % experiments/FIGS/data/v2/index.json :: nonlinear probe bb_vs_cc: (1 - AUC) of R16_Q1 over L188, run by run :: sha256 c6c4138ee5c86423 :: prose twin: two significant figures, direction in words
\newcommand{\MlstMlpCmpBbResMassOverRes}{\pending{no run yet}}  % experiments/FIGS/data/v2/index.json :: nonlinear probe bb_vs_cc: (1 - AUC) of L162_MASS over L162, run by run :: sha256 c6c4138ee5c86423
\newcommand{\MlstMlpCmpBbResMassOverResPr}{\pending{no run yet}}  % experiments/FIGS/data/v2/index.json :: nonlinear probe bb_vs_cc: (1 - AUC) of L162_MASS over L162, run by run :: sha256 c6c4138ee5c86423 :: prose twin: two significant figures, direction in words
\newcommand{\MlstMlpCmpBbResMassOverResPrMax}{\pending{no run yet}}  % experiments/FIGS/data/v2/index.json :: nonlinear probe bb_vs_cc: (1 - AUC) of L162_MASS over L162, run by run :: sha256 c6c4138ee5c86423 :: prose twin of the run farther from 1
\newcommand{\MlstMlpCmpBbCoarseMassOverCoarse}{\pending{no run yet}}  % experiments/FIGS/data/v2/index.json :: nonlinear probe bb_vs_cc: (1 - AUC) of R16_Q1_MASS over R16_Q1, run by run :: sha256 c6c4138ee5c86423
\newcommand{\MlstMlpCmpBbCoarseMassOverCoarsePr}{\pending{no run yet}}  % experiments/FIGS/data/v2/index.json :: nonlinear probe bb_vs_cc: (1 - AUC) of R16_Q1_MASS over R16_Q1, run by run :: sha256 c6c4138ee5c86423 :: prose twin: two significant figures, direction in words
\newcommand{\MlstMlpCmpBbCoarseMassOverCoarsePrMax}{\pending{no run yet}}  % experiments/FIGS/data/v2/index.json :: nonlinear probe bb_vs_cc: (1 - AUC) of R16_Q1_MASS over R16_Q1, run by run :: sha256 c6c4138ee5c86423 :: prose twin of the run farther from 1
\newcommand{\MlstMlpCmpBbFzeroOverFull}{3.3 and 4.3}  % experiments/FIGS/data/v2/index.json :: nonlinear probe bb_vs_cc: (1 - AUC) of FLAV_F0 over L188, run by run :: sha256 c6c4138ee5c86423
\newcommand{\MlstMlpCmpBbFzeroOverFullPr}{at least 3.3 times the errors}  % experiments/FIGS/data/v2/index.json :: nonlinear probe bb_vs_cc: (1 - AUC) of FLAV_F0 over L188, run by run :: sha256 c6c4138ee5c86423 :: prose twin: two significant figures, direction in words
\newcommand{\MlstMlpCmpBbFzeroOverFullPrMax}{up to 4.3 times the errors}  % experiments/FIGS/data/v2/index.json :: nonlinear probe bb_vs_cc: (1 - AUC) of FLAV_F0 over L188, run by run :: sha256 c6c4138ee5c86423 :: prose twin of the run farther from 1
\newcommand{\MlstMlpCmpBbFoneOverFzero}{0.52 and 0.38}  % experiments/FIGS/data/v2/index.json :: nonlinear probe bb_vs_cc: (1 - AUC) of FLAV_F1 over FLAV_F0, run by run :: sha256 c6c4138ee5c86423
\newcommand{\MlstMlpCmpBbFoneOverFzeroPr}{at least 48\% fewer errors}  % experiments/FIGS/data/v2/index.json :: nonlinear probe bb_vs_cc: (1 - AUC) of FLAV_F1 over FLAV_F0, run by run :: sha256 c6c4138ee5c86423 :: prose twin: two significant figures, direction in words
\newcommand{\MlstMlpCmpBbFoneOverFzeroPrMax}{down to 0.38 times the errors}  % experiments/FIGS/data/v2/index.json :: nonlinear probe bb_vs_cc: (1 - AUC) of FLAV_F1 over FLAV_F0, run by run :: sha256 c6c4138ee5c86423 :: prose twin of the run farther from 1
\newcommand{\MlstMlpCmpBbFonerOverFone}{1.27 and 1.45}  % experiments/FIGS/data/v2/index.json :: nonlinear probe bb_vs_cc: (1 - AUC) of FLAV_F1R over FLAV_F1, run by run :: sha256 c6c4138ee5c86423
\newcommand{\MlstMlpCmpBbFonerOverFonePr}{at least 27\% more errors}  % experiments/FIGS/data/v2/index.json :: nonlinear probe bb_vs_cc: (1 - AUC) of FLAV_F1R over FLAV_F1, run by run :: sha256 c6c4138ee5c86423 :: prose twin: two significant figures, direction in words
\newcommand{\MlstMlpCmpBbFonerOverFonePrMax}{up to 45\% more errors}  % experiments/FIGS/data/v2/index.json :: nonlinear probe bb_vs_cc: (1 - AUC) of FLAV_F1R over FLAV_F1, run by run :: sha256 c6c4138ee5c86423 :: prose twin of the run farther from 1
\newcommand{\MlstMlpCmpBbFonerOverFzero}{0.66 and 0.55}  % experiments/FIGS/data/v2/index.json :: nonlinear probe bb_vs_cc: (1 - AUC) of FLAV_F1R over FLAV_F0, run by run :: sha256 c6c4138ee5c86423
\newcommand{\MlstMlpCmpBbFonerOverFzeroPr}{at least 34\% fewer errors}  % experiments/FIGS/data/v2/index.json :: nonlinear probe bb_vs_cc: (1 - AUC) of FLAV_F1R over FLAV_F0, run by run :: sha256 c6c4138ee5c86423 :: prose twin: two significant figures, direction in words
\newcommand{\MlstMlpCmpBbFonerOverFzeroPrMax}{up to 45\% fewer errors}  % experiments/FIGS/data/v2/index.json :: nonlinear probe bb_vs_cc: (1 - AUC) of FLAV_F1R over FLAV_F0, run by run :: sha256 c6c4138ee5c86423 :: prose twin of the run farther from 1
\newcommand{\MlstMlpCmpBbFzeroOverCoarse}{1.07 and 1.41}  % experiments/FIGS/data/v2/index.json :: nonlinear probe bb_vs_cc: (1 - AUC) of FLAV_F0 over R16_Q1, run by run :: sha256 c6c4138ee5c86423
\newcommand{\MlstMlpCmpBbFzeroOverCoarsePr}{at least 7.4\% more errors}  % experiments/FIGS/data/v2/index.json :: nonlinear probe bb_vs_cc: (1 - AUC) of FLAV_F0 over R16_Q1, run by run :: sha256 c6c4138ee5c86423 :: prose twin: two significant figures, direction in words
\newcommand{\MlstMlpCmpBbFzeroOverCoarsePrMax}{up to 41\% more errors}  % experiments/FIGS/data/v2/index.json :: nonlinear probe bb_vs_cc: (1 - AUC) of FLAV_F0 over R16_Q1, run by run :: sha256 c6c4138ee5c86423 :: prose twin of the run farther from 1
\newcommand{\MlstMlpCmpBbRandOverFull}{\pending{no run yet}}  % experiments/FIGS/data/v2/index.json :: nonlinear probe bb_vs_cc: (1 - AUC) of each random partition over the 188-class run of its seed :: sha256 c6c4138ee5c86423
\newcommand{\MlstMlpCmpBbRandOverFullPr}{\pending{no run yet}}  % experiments/FIGS/data/v2/index.json :: nonlinear probe bb_vs_cc: (1 - AUC) of each random partition over the 188-class run of its seed :: sha256 c6c4138ee5c86423 :: prose twin: two significant figures, direction in words
\newcommand{\MlstMlpCmpBbqqResOverFull}{1.011\,\ensuremath{\pm}\,0.030}  % experiments/FIGS/data/v2/index.json :: nonlinear probe bbqq_vs_ccqq: (1 - AUC) of L162 over L188, run by run :: sha256 c6c4138ee5c86423
\newcommand{\MlstMlpCmpBbqqResOverFullPr}{equal within 3\%}  % experiments/FIGS/data/v2/index.json :: nonlinear probe bbqq_vs_ccqq: (1 - AUC) of L162 over L188, run by run :: sha256 c6c4138ee5c86423 :: prose twin: two significant figures, direction in words
\newcommand{\MlstMlpCmpBbqqMidOverFull}{1.112\,\ensuremath{\pm}\,0.006}  % experiments/FIGS/data/v2/index.json :: nonlinear probe bbqq_vs_ccqq: (1 - AUC) of R42_Q1 over L188, run by run :: sha256 c6c4138ee5c86423
\newcommand{\MlstMlpCmpBbqqMidOverFullPr}{11\% more errors}  % experiments/FIGS/data/v2/index.json :: nonlinear probe bbqq_vs_ccqq: (1 - AUC) of R42_Q1 over L188, run by run :: sha256 c6c4138ee5c86423 :: prose twin: two significant figures, direction in words
\newcommand{\MlstMlpCmpBbqqCoarseOverFull}{2.93\,\ensuremath{\pm}\,0.17}  % experiments/FIGS/data/v2/index.json :: nonlinear probe bbqq_vs_ccqq: (1 - AUC) of R16_Q1 over L188, run by run :: sha256 c6c4138ee5c86423
\newcommand{\MlstMlpCmpBbqqCoarseOverFullPr}{2.9 times the errors}  % experiments/FIGS/data/v2/index.json :: nonlinear probe bbqq_vs_ccqq: (1 - AUC) of R16_Q1 over L188, run by run :: sha256 c6c4138ee5c86423 :: prose twin: two significant figures, direction in words
\newcommand{\MlstMlpCmpBbqqResMassOverRes}{\pending{no run yet}}  % experiments/FIGS/data/v2/index.json :: nonlinear probe bbqq_vs_ccqq: (1 - AUC) of L162_MASS over L162, run by run :: sha256 c6c4138ee5c86423
\newcommand{\MlstMlpCmpBbqqResMassOverResPr}{\pending{no run yet}}  % experiments/FIGS/data/v2/index.json :: nonlinear probe bbqq_vs_ccqq: (1 - AUC) of L162_MASS over L162, run by run :: sha256 c6c4138ee5c86423 :: prose twin: two significant figures, direction in words
\newcommand{\MlstMlpCmpBbqqResMassOverResPrMax}{\pending{no run yet}}  % experiments/FIGS/data/v2/index.json :: nonlinear probe bbqq_vs_ccqq: (1 - AUC) of L162_MASS over L162, run by run :: sha256 c6c4138ee5c86423 :: prose twin of the run farther from 1
\newcommand{\MlstMlpCmpBbqqCoarseMassOverCoarse}{\pending{no run yet}}  % experiments/FIGS/data/v2/index.json :: nonlinear probe bbqq_vs_ccqq: (1 - AUC) of R16_Q1_MASS over R16_Q1, run by run :: sha256 c6c4138ee5c86423
\newcommand{\MlstMlpCmpBbqqCoarseMassOverCoarsePr}{\pending{no run yet}}  % experiments/FIGS/data/v2/index.json :: nonlinear probe bbqq_vs_ccqq: (1 - AUC) of R16_Q1_MASS over R16_Q1, run by run :: sha256 c6c4138ee5c86423 :: prose twin: two significant figures, direction in words
\newcommand{\MlstMlpCmpBbqqCoarseMassOverCoarsePrMax}{\pending{no run yet}}  % experiments/FIGS/data/v2/index.json :: nonlinear probe bbqq_vs_ccqq: (1 - AUC) of R16_Q1_MASS over R16_Q1, run by run :: sha256 c6c4138ee5c86423 :: prose twin of the run farther from 1
\newcommand{\MlstMlpCmpBbqqFzeroOverFull}{3.20 and 3.78}  % experiments/FIGS/data/v2/index.json :: nonlinear probe bbqq_vs_ccqq: (1 - AUC) of FLAV_F0 over L188, run by run :: sha256 c6c4138ee5c86423
\newcommand{\MlstMlpCmpBbqqFzeroOverFullPr}{at least 3.2 times the errors}  % experiments/FIGS/data/v2/index.json :: nonlinear probe bbqq_vs_ccqq: (1 - AUC) of FLAV_F0 over L188, run by run :: sha256 c6c4138ee5c86423 :: prose twin: two significant figures, direction in words
\newcommand{\MlstMlpCmpBbqqFzeroOverFullPrMax}{up to 3.8 times the errors}  % experiments/FIGS/data/v2/index.json :: nonlinear probe bbqq_vs_ccqq: (1 - AUC) of FLAV_F0 over L188, run by run :: sha256 c6c4138ee5c86423 :: prose twin of the run farther from 1
\newcommand{\MlstMlpCmpBbqqFoneOverFzero}{0.450 and 0.363}  % experiments/FIGS/data/v2/index.json :: nonlinear probe bbqq_vs_ccqq: (1 - AUC) of FLAV_F1 over FLAV_F0, run by run :: sha256 c6c4138ee5c86423
\newcommand{\MlstMlpCmpBbqqFoneOverFzeroPr}{at most 0.45 times the errors}  % experiments/FIGS/data/v2/index.json :: nonlinear probe bbqq_vs_ccqq: (1 - AUC) of FLAV_F1 over FLAV_F0, run by run :: sha256 c6c4138ee5c86423 :: prose twin: two significant figures, direction in words
\newcommand{\MlstMlpCmpBbqqFoneOverFzeroPrMax}{down to 0.36 times the errors}  % experiments/FIGS/data/v2/index.json :: nonlinear probe bbqq_vs_ccqq: (1 - AUC) of FLAV_F1 over FLAV_F0, run by run :: sha256 c6c4138ee5c86423 :: prose twin of the run farther from 1
\newcommand{\MlstMlpCmpBbqqFonerOverFone}{1.311 and 1.343}  % experiments/FIGS/data/v2/index.json :: nonlinear probe bbqq_vs_ccqq: (1 - AUC) of FLAV_F1R over FLAV_F1, run by run :: sha256 c6c4138ee5c86423
\newcommand{\MlstMlpCmpBbqqFonerOverFonePr}{at least 31\% more errors}  % experiments/FIGS/data/v2/index.json :: nonlinear probe bbqq_vs_ccqq: (1 - AUC) of FLAV_F1R over FLAV_F1, run by run :: sha256 c6c4138ee5c86423 :: prose twin: two significant figures, direction in words
\newcommand{\MlstMlpCmpBbqqFonerOverFonePrMax}{up to 34\% more errors}  % experiments/FIGS/data/v2/index.json :: nonlinear probe bbqq_vs_ccqq: (1 - AUC) of FLAV_F1R over FLAV_F1, run by run :: sha256 c6c4138ee5c86423 :: prose twin of the run farther from 1
\newcommand{\MlstMlpCmpBbqqFonerOverFzero}{0.59 and 0.49}  % experiments/FIGS/data/v2/index.json :: nonlinear probe bbqq_vs_ccqq: (1 - AUC) of FLAV_F1R over FLAV_F0, run by run :: sha256 c6c4138ee5c86423
\newcommand{\MlstMlpCmpBbqqFonerOverFzeroPr}{at least 41\% fewer errors}  % experiments/FIGS/data/v2/index.json :: nonlinear probe bbqq_vs_ccqq: (1 - AUC) of FLAV_F1R over FLAV_F0, run by run :: sha256 c6c4138ee5c86423 :: prose twin: two significant figures, direction in words
\newcommand{\MlstMlpCmpBbqqFonerOverFzeroPrMax}{down to 0.49 times the errors}  % experiments/FIGS/data/v2/index.json :: nonlinear probe bbqq_vs_ccqq: (1 - AUC) of FLAV_F1R over FLAV_F0, run by run :: sha256 c6c4138ee5c86423 :: prose twin of the run farther from 1
\newcommand{\MlstMlpCmpBbqqFzeroOverCoarse}{1.170 and 1.236}  % experiments/FIGS/data/v2/index.json :: nonlinear probe bbqq_vs_ccqq: (1 - AUC) of FLAV_F0 over R16_Q1, run by run :: sha256 c6c4138ee5c86423
\newcommand{\MlstMlpCmpBbqqFzeroOverCoarsePr}{at least 17\% more errors}  % experiments/FIGS/data/v2/index.json :: nonlinear probe bbqq_vs_ccqq: (1 - AUC) of FLAV_F0 over R16_Q1, run by run :: sha256 c6c4138ee5c86423 :: prose twin: two significant figures, direction in words
\newcommand{\MlstMlpCmpBbqqFzeroOverCoarsePrMax}{up to 24\% more errors}  % experiments/FIGS/data/v2/index.json :: nonlinear probe bbqq_vs_ccqq: (1 - AUC) of FLAV_F0 over R16_Q1, run by run :: sha256 c6c4138ee5c86423 :: prose twin of the run farther from 1
\newcommand{\MlstMlpCmpBbqqRandOverFull}{\pending{no run yet}}  % experiments/FIGS/data/v2/index.json :: nonlinear probe bbqq_vs_ccqq: (1 - AUC) of each random partition over the 188-class run of its seed :: sha256 c6c4138ee5c86423
\newcommand{\MlstMlpCmpBbqqRandOverFullPr}{\pending{no run yet}}  % experiments/FIGS/data/v2/index.json :: nonlinear probe bbqq_vs_ccqq: (1 - AUC) of each random partition over the 188-class run of its seed :: sha256 c6c4138ee5c86423 :: prose twin: two significant figures, direction in words
\newcommand{\MlstMlpCmpQcdBbResOverFull}{1.085\,\ensuremath{\pm}\,0.030}  % experiments/FIGS/data/v2/index.json :: nonlinear probe qcd_bb_vs_cc: (1 - AUC) of L162 over L188, run by run :: sha256 c6c4138ee5c86423
\newcommand{\MlstMlpCmpQcdBbResOverFullPr}{8.5\% more errors}  % experiments/FIGS/data/v2/index.json :: nonlinear probe qcd_bb_vs_cc: (1 - AUC) of L162 over L188, run by run :: sha256 c6c4138ee5c86423 :: prose twin: two significant figures, direction in words
\newcommand{\MlstMlpCmpQcdBbMidOverFull}{1.121\,\ensuremath{\pm}\,0.017}  % experiments/FIGS/data/v2/index.json :: nonlinear probe qcd_bb_vs_cc: (1 - AUC) of R42_Q1 over L188, run by run :: sha256 c6c4138ee5c86423
\newcommand{\MlstMlpCmpQcdBbMidOverFullPr}{12\% more errors}  % experiments/FIGS/data/v2/index.json :: nonlinear probe qcd_bb_vs_cc: (1 - AUC) of R42_Q1 over L188, run by run :: sha256 c6c4138ee5c86423 :: prose twin: two significant figures, direction in words
\newcommand{\MlstMlpCmpQcdBbCoarseOverFull}{2.31\,\ensuremath{\pm}\,0.04}  % experiments/FIGS/data/v2/index.json :: nonlinear probe qcd_bb_vs_cc: (1 - AUC) of R16_Q1 over L188, run by run :: sha256 c6c4138ee5c86423
\newcommand{\MlstMlpCmpQcdBbCoarseOverFullPr}{2.3 times the errors}  % experiments/FIGS/data/v2/index.json :: nonlinear probe qcd_bb_vs_cc: (1 - AUC) of R16_Q1 over L188, run by run :: sha256 c6c4138ee5c86423 :: prose twin: two significant figures, direction in words
\newcommand{\MlstMlpCmpQcdBbResMassOverRes}{\pending{no run yet}}  % experiments/FIGS/data/v2/index.json :: nonlinear probe qcd_bb_vs_cc: (1 - AUC) of L162_MASS over L162, run by run :: sha256 c6c4138ee5c86423
\newcommand{\MlstMlpCmpQcdBbResMassOverResPr}{\pending{no run yet}}  % experiments/FIGS/data/v2/index.json :: nonlinear probe qcd_bb_vs_cc: (1 - AUC) of L162_MASS over L162, run by run :: sha256 c6c4138ee5c86423 :: prose twin: two significant figures, direction in words
\newcommand{\MlstMlpCmpQcdBbResMassOverResPrMax}{\pending{no run yet}}  % experiments/FIGS/data/v2/index.json :: nonlinear probe qcd_bb_vs_cc: (1 - AUC) of L162_MASS over L162, run by run :: sha256 c6c4138ee5c86423 :: prose twin of the run farther from 1
\newcommand{\MlstMlpCmpQcdBbCoarseMassOverCoarse}{\pending{no run yet}}  % experiments/FIGS/data/v2/index.json :: nonlinear probe qcd_bb_vs_cc: (1 - AUC) of R16_Q1_MASS over R16_Q1, run by run :: sha256 c6c4138ee5c86423
\newcommand{\MlstMlpCmpQcdBbCoarseMassOverCoarsePr}{\pending{no run yet}}  % experiments/FIGS/data/v2/index.json :: nonlinear probe qcd_bb_vs_cc: (1 - AUC) of R16_Q1_MASS over R16_Q1, run by run :: sha256 c6c4138ee5c86423 :: prose twin: two significant figures, direction in words
\newcommand{\MlstMlpCmpQcdBbCoarseMassOverCoarsePrMax}{\pending{no run yet}}  % experiments/FIGS/data/v2/index.json :: nonlinear probe qcd_bb_vs_cc: (1 - AUC) of R16_Q1_MASS over R16_Q1, run by run :: sha256 c6c4138ee5c86423 :: prose twin of the run farther from 1
\newcommand{\MlstMlpCmpQcdBbFzeroOverFull}{2.50 and 2.72}  % experiments/FIGS/data/v2/index.json :: nonlinear probe qcd_bb_vs_cc: (1 - AUC) of FLAV_F0 over L188, run by run :: sha256 c6c4138ee5c86423
\newcommand{\MlstMlpCmpQcdBbFzeroOverFullPr}{at least 2.5 times the errors}  % experiments/FIGS/data/v2/index.json :: nonlinear probe qcd_bb_vs_cc: (1 - AUC) of FLAV_F0 over L188, run by run :: sha256 c6c4138ee5c86423 :: prose twin: two significant figures, direction in words
\newcommand{\MlstMlpCmpQcdBbFzeroOverFullPrMax}{up to 2.7 times the errors}  % experiments/FIGS/data/v2/index.json :: nonlinear probe qcd_bb_vs_cc: (1 - AUC) of FLAV_F0 over L188, run by run :: sha256 c6c4138ee5c86423 :: prose twin of the run farther from 1
\newcommand{\MlstMlpCmpQcdBbFoneOverFzero}{0.573 and 0.505}  % experiments/FIGS/data/v2/index.json :: nonlinear probe qcd_bb_vs_cc: (1 - AUC) of FLAV_F1 over FLAV_F0, run by run :: sha256 c6c4138ee5c86423
\newcommand{\MlstMlpCmpQcdBbFoneOverFzeroPr}{at least 43\% fewer errors}  % experiments/FIGS/data/v2/index.json :: nonlinear probe qcd_bb_vs_cc: (1 - AUC) of FLAV_F1 over FLAV_F0, run by run :: sha256 c6c4138ee5c86423 :: prose twin: two significant figures, direction in words
\newcommand{\MlstMlpCmpQcdBbFoneOverFzeroPrMax}{up to 49\% fewer errors}  % experiments/FIGS/data/v2/index.json :: nonlinear probe qcd_bb_vs_cc: (1 - AUC) of FLAV_F1 over FLAV_F0, run by run :: sha256 c6c4138ee5c86423 :: prose twin of the run farther from 1
\newcommand{\MlstMlpCmpQcdBbFonerOverFone}{1.219 and 1.258}  % experiments/FIGS/data/v2/index.json :: nonlinear probe qcd_bb_vs_cc: (1 - AUC) of FLAV_F1R over FLAV_F1, run by run :: sha256 c6c4138ee5c86423
\newcommand{\MlstMlpCmpQcdBbFonerOverFonePr}{at least 22\% more errors}  % experiments/FIGS/data/v2/index.json :: nonlinear probe qcd_bb_vs_cc: (1 - AUC) of FLAV_F1R over FLAV_F1, run by run :: sha256 c6c4138ee5c86423 :: prose twin: two significant figures, direction in words
\newcommand{\MlstMlpCmpQcdBbFonerOverFonePrMax}{up to 26\% more errors}  % experiments/FIGS/data/v2/index.json :: nonlinear probe qcd_bb_vs_cc: (1 - AUC) of FLAV_F1R over FLAV_F1, run by run :: sha256 c6c4138ee5c86423 :: prose twin of the run farther from 1
\newcommand{\MlstMlpCmpQcdBbFonerOverFzero}{0.699 and 0.635}  % experiments/FIGS/data/v2/index.json :: nonlinear probe qcd_bb_vs_cc: (1 - AUC) of FLAV_F1R over FLAV_F0, run by run :: sha256 c6c4138ee5c86423
\newcommand{\MlstMlpCmpQcdBbFonerOverFzeroPr}{at least 30\% fewer errors}  % experiments/FIGS/data/v2/index.json :: nonlinear probe qcd_bb_vs_cc: (1 - AUC) of FLAV_F1R over FLAV_F0, run by run :: sha256 c6c4138ee5c86423 :: prose twin: two significant figures, direction in words
\newcommand{\MlstMlpCmpQcdBbFonerOverFzeroPrMax}{up to 36\% fewer errors}  % experiments/FIGS/data/v2/index.json :: nonlinear probe qcd_bb_vs_cc: (1 - AUC) of FLAV_F1R over FLAV_F0, run by run :: sha256 c6c4138ee5c86423 :: prose twin of the run farther from 1
\newcommand{\MlstMlpCmpQcdBbFzeroOverCoarse}{1.103 and 1.167}  % experiments/FIGS/data/v2/index.json :: nonlinear probe qcd_bb_vs_cc: (1 - AUC) of FLAV_F0 over R16_Q1, run by run :: sha256 c6c4138ee5c86423
\newcommand{\MlstMlpCmpQcdBbFzeroOverCoarsePr}{at least 10\% more errors}  % experiments/FIGS/data/v2/index.json :: nonlinear probe qcd_bb_vs_cc: (1 - AUC) of FLAV_F0 over R16_Q1, run by run :: sha256 c6c4138ee5c86423 :: prose twin: two significant figures, direction in words
\newcommand{\MlstMlpCmpQcdBbFzeroOverCoarsePrMax}{up to 17\% more errors}  % experiments/FIGS/data/v2/index.json :: nonlinear probe qcd_bb_vs_cc: (1 - AUC) of FLAV_F0 over R16_Q1, run by run :: sha256 c6c4138ee5c86423 :: prose twin of the run farther from 1
\newcommand{\MlstMlpCmpQcdBbRandOverFull}{\pending{no run yet}}  % experiments/FIGS/data/v2/index.json :: nonlinear probe qcd_bb_vs_cc: (1 - AUC) of each random partition over the 188-class run of its seed :: sha256 c6c4138ee5c86423
\newcommand{\MlstMlpCmpQcdBbRandOverFullPr}{\pending{no run yet}}  % experiments/FIGS/data/v2/index.json :: nonlinear probe qcd_bb_vs_cc: (1 - AUC) of each random partition over the 188-class run of its seed :: sha256 c6c4138ee5c86423 :: prose twin: two significant figures, direction in words
\newcommand{\MlstMlpCmpQcdBResOverFull}{1.050\,\ensuremath{\pm}\,0.004}  % experiments/FIGS/data/v2/index.json :: nonlinear probe qcd_b_vs_light: (1 - AUC) of L162 over L188, run by run :: sha256 c6c4138ee5c86423
\newcommand{\MlstMlpCmpQcdBResOverFullPr}{5.0\% more errors}  % experiments/FIGS/data/v2/index.json :: nonlinear probe qcd_b_vs_light: (1 - AUC) of L162 over L188, run by run :: sha256 c6c4138ee5c86423 :: prose twin: two significant figures, direction in words
\newcommand{\MlstMlpCmpQcdBMidOverFull}{1.055\,\ensuremath{\pm}\,0.005}  % experiments/FIGS/data/v2/index.json :: nonlinear probe qcd_b_vs_light: (1 - AUC) of R42_Q1 over L188, run by run :: sha256 c6c4138ee5c86423
\newcommand{\MlstMlpCmpQcdBMidOverFullPr}{5.5\% more errors}  % experiments/FIGS/data/v2/index.json :: nonlinear probe qcd_b_vs_light: (1 - AUC) of R42_Q1 over L188, run by run :: sha256 c6c4138ee5c86423 :: prose twin: two significant figures, direction in words
\newcommand{\MlstMlpCmpQcdBCoarseOverFull}{1.337\,\ensuremath{\pm}\,0.021}  % experiments/FIGS/data/v2/index.json :: nonlinear probe qcd_b_vs_light: (1 - AUC) of R16_Q1 over L188, run by run :: sha256 c6c4138ee5c86423
\newcommand{\MlstMlpCmpQcdBCoarseOverFullPr}{34\% more errors}  % experiments/FIGS/data/v2/index.json :: nonlinear probe qcd_b_vs_light: (1 - AUC) of R16_Q1 over L188, run by run :: sha256 c6c4138ee5c86423 :: prose twin: two significant figures, direction in words
\newcommand{\MlstMlpCmpQcdBResMassOverRes}{\pending{no run yet}}  % experiments/FIGS/data/v2/index.json :: nonlinear probe qcd_b_vs_light: (1 - AUC) of L162_MASS over L162, run by run :: sha256 c6c4138ee5c86423
\newcommand{\MlstMlpCmpQcdBResMassOverResPr}{\pending{no run yet}}  % experiments/FIGS/data/v2/index.json :: nonlinear probe qcd_b_vs_light: (1 - AUC) of L162_MASS over L162, run by run :: sha256 c6c4138ee5c86423 :: prose twin: two significant figures, direction in words
\newcommand{\MlstMlpCmpQcdBResMassOverResPrMax}{\pending{no run yet}}  % experiments/FIGS/data/v2/index.json :: nonlinear probe qcd_b_vs_light: (1 - AUC) of L162_MASS over L162, run by run :: sha256 c6c4138ee5c86423 :: prose twin of the run farther from 1
\newcommand{\MlstMlpCmpQcdBCoarseMassOverCoarse}{\pending{no run yet}}  % experiments/FIGS/data/v2/index.json :: nonlinear probe qcd_b_vs_light: (1 - AUC) of R16_Q1_MASS over R16_Q1, run by run :: sha256 c6c4138ee5c86423
\newcommand{\MlstMlpCmpQcdBCoarseMassOverCoarsePr}{\pending{no run yet}}  % experiments/FIGS/data/v2/index.json :: nonlinear probe qcd_b_vs_light: (1 - AUC) of R16_Q1_MASS over R16_Q1, run by run :: sha256 c6c4138ee5c86423 :: prose twin: two significant figures, direction in words
\newcommand{\MlstMlpCmpQcdBCoarseMassOverCoarsePrMax}{\pending{no run yet}}  % experiments/FIGS/data/v2/index.json :: nonlinear probe qcd_b_vs_light: (1 - AUC) of R16_Q1_MASS over R16_Q1, run by run :: sha256 c6c4138ee5c86423 :: prose twin of the run farther from 1
\newcommand{\MlstMlpCmpQcdBFzeroOverFull}{1.315 and 1.288}  % experiments/FIGS/data/v2/index.json :: nonlinear probe qcd_b_vs_light: (1 - AUC) of FLAV_F0 over L188, run by run :: sha256 c6c4138ee5c86423
\newcommand{\MlstMlpCmpQcdBFzeroOverFullPr}{at least 29\% more errors}  % experiments/FIGS/data/v2/index.json :: nonlinear probe qcd_b_vs_light: (1 - AUC) of FLAV_F0 over L188, run by run :: sha256 c6c4138ee5c86423 :: prose twin: two significant figures, direction in words
\newcommand{\MlstMlpCmpQcdBFzeroOverFullPrMax}{up to 32\% more errors}  % experiments/FIGS/data/v2/index.json :: nonlinear probe qcd_b_vs_light: (1 - AUC) of FLAV_F0 over L188, run by run :: sha256 c6c4138ee5c86423 :: prose twin of the run farther from 1
\newcommand{\MlstMlpCmpQcdBFoneOverFzero}{0.8804 and 0.8849}  % experiments/FIGS/data/v2/index.json :: nonlinear probe qcd_b_vs_light: (1 - AUC) of FLAV_F1 over FLAV_F0, run by run :: sha256 c6c4138ee5c86423
\newcommand{\MlstMlpCmpQcdBFoneOverFzeroPr}{at least 12\% fewer errors}  % experiments/FIGS/data/v2/index.json :: nonlinear probe qcd_b_vs_light: (1 - AUC) of FLAV_F1 over FLAV_F0, run by run :: sha256 c6c4138ee5c86423 :: prose twin: two significant figures, direction in words
\newcommand{\MlstMlpCmpQcdBFoneOverFzeroPrMax}{up to 12\% fewer errors}  % experiments/FIGS/data/v2/index.json :: nonlinear probe qcd_b_vs_light: (1 - AUC) of FLAV_F1 over FLAV_F0, run by run :: sha256 c6c4138ee5c86423 :: prose twin of the run farther from 1
\newcommand{\MlstMlpCmpQcdBFonerOverFone}{1.037 and 1.042}  % experiments/FIGS/data/v2/index.json :: nonlinear probe qcd_b_vs_light: (1 - AUC) of FLAV_F1R over FLAV_F1, run by run :: sha256 c6c4138ee5c86423
\newcommand{\MlstMlpCmpQcdBFonerOverFonePr}{at least 3.7\% more errors}  % experiments/FIGS/data/v2/index.json :: nonlinear probe qcd_b_vs_light: (1 - AUC) of FLAV_F1R over FLAV_F1, run by run :: sha256 c6c4138ee5c86423 :: prose twin: two significant figures, direction in words
\newcommand{\MlstMlpCmpQcdBFonerOverFonePrMax}{up to 4.2\% more errors}  % experiments/FIGS/data/v2/index.json :: nonlinear probe qcd_b_vs_light: (1 - AUC) of FLAV_F1R over FLAV_F1, run by run :: sha256 c6c4138ee5c86423 :: prose twin of the run farther from 1
\newcommand{\MlstMlpCmpQcdBFonerOverFzero}{0.9130 and 0.9217}  % experiments/FIGS/data/v2/index.json :: nonlinear probe qcd_b_vs_light: (1 - AUC) of FLAV_F1R over FLAV_F0, run by run :: sha256 c6c4138ee5c86423
\newcommand{\MlstMlpCmpQcdBFonerOverFzeroPr}{at least 7.8\% fewer errors}  % experiments/FIGS/data/v2/index.json :: nonlinear probe qcd_b_vs_light: (1 - AUC) of FLAV_F1R over FLAV_F0, run by run :: sha256 c6c4138ee5c86423 :: prose twin: two significant figures, direction in words
\newcommand{\MlstMlpCmpQcdBFonerOverFzeroPrMax}{up to 8.7\% fewer errors}  % experiments/FIGS/data/v2/index.json :: nonlinear probe qcd_b_vs_light: (1 - AUC) of FLAV_F1R over FLAV_F0, run by run :: sha256 c6c4138ee5c86423 :: prose twin of the run farther from 1
\newcommand{\MlstMlpCmpQcdBFzeroOverCoarse}{0.968 and 0.978}  % experiments/FIGS/data/v2/index.json :: nonlinear probe qcd_b_vs_light: (1 - AUC) of FLAV_F0 over R16_Q1, run by run :: sha256 c6c4138ee5c86423
\newcommand{\MlstMlpCmpQcdBFzeroOverCoarsePr}{at least 2.2\% fewer errors}  % experiments/FIGS/data/v2/index.json :: nonlinear probe qcd_b_vs_light: (1 - AUC) of FLAV_F0 over R16_Q1, run by run :: sha256 c6c4138ee5c86423 :: prose twin: two significant figures, direction in words
\newcommand{\MlstMlpCmpQcdBFzeroOverCoarsePrMax}{up to 3.2\% fewer errors}  % experiments/FIGS/data/v2/index.json :: nonlinear probe qcd_b_vs_light: (1 - AUC) of FLAV_F0 over R16_Q1, run by run :: sha256 c6c4138ee5c86423 :: prose twin of the run farther from 1
\newcommand{\MlstMlpCmpQcdBRandOverFull}{\pending{no run yet}}  % experiments/FIGS/data/v2/index.json :: nonlinear probe qcd_b_vs_light: (1 - AUC) of each random partition over the 188-class run of its seed :: sha256 c6c4138ee5c86423
\newcommand{\MlstMlpCmpQcdBRandOverFullPr}{\pending{no run yet}}  % experiments/FIGS/data/v2/index.json :: nonlinear probe qcd_b_vs_light: (1 - AUC) of each random partition over the 188-class run of its seed :: sha256 c6c4138ee5c86423 :: prose twin: two significant figures, direction in words
\newcommand{\MlstMlpCmpQcdCResOverFull}{1.039\,\ensuremath{\pm}\,0.005}  % experiments/FIGS/data/v2/index.json :: nonlinear probe qcd_c_vs_light: (1 - AUC) of L162 over L188, run by run :: sha256 c6c4138ee5c86423
\newcommand{\MlstMlpCmpQcdCResOverFullPr}{3.9\% more errors}  % experiments/FIGS/data/v2/index.json :: nonlinear probe qcd_c_vs_light: (1 - AUC) of L162 over L188, run by run :: sha256 c6c4138ee5c86423 :: prose twin: two significant figures, direction in words
\newcommand{\MlstMlpCmpQcdCMidOverFull}{1.055\,\ensuremath{\pm}\,0.006}  % experiments/FIGS/data/v2/index.json :: nonlinear probe qcd_c_vs_light: (1 - AUC) of R42_Q1 over L188, run by run :: sha256 c6c4138ee5c86423
\newcommand{\MlstMlpCmpQcdCMidOverFullPr}{5.5\% more errors}  % experiments/FIGS/data/v2/index.json :: nonlinear probe qcd_c_vs_light: (1 - AUC) of R42_Q1 over L188, run by run :: sha256 c6c4138ee5c86423 :: prose twin: two significant figures, direction in words
\newcommand{\MlstMlpCmpQcdCCoarseOverFull}{1.274\,\ensuremath{\pm}\,0.008}  % experiments/FIGS/data/v2/index.json :: nonlinear probe qcd_c_vs_light: (1 - AUC) of R16_Q1 over L188, run by run :: sha256 c6c4138ee5c86423
\newcommand{\MlstMlpCmpQcdCCoarseOverFullPr}{27\% more errors}  % experiments/FIGS/data/v2/index.json :: nonlinear probe qcd_c_vs_light: (1 - AUC) of R16_Q1 over L188, run by run :: sha256 c6c4138ee5c86423 :: prose twin: two significant figures, direction in words
\newcommand{\MlstMlpCmpQcdCResMassOverRes}{\pending{no run yet}}  % experiments/FIGS/data/v2/index.json :: nonlinear probe qcd_c_vs_light: (1 - AUC) of L162_MASS over L162, run by run :: sha256 c6c4138ee5c86423
\newcommand{\MlstMlpCmpQcdCResMassOverResPr}{\pending{no run yet}}  % experiments/FIGS/data/v2/index.json :: nonlinear probe qcd_c_vs_light: (1 - AUC) of L162_MASS over L162, run by run :: sha256 c6c4138ee5c86423 :: prose twin: two significant figures, direction in words
\newcommand{\MlstMlpCmpQcdCResMassOverResPrMax}{\pending{no run yet}}  % experiments/FIGS/data/v2/index.json :: nonlinear probe qcd_c_vs_light: (1 - AUC) of L162_MASS over L162, run by run :: sha256 c6c4138ee5c86423 :: prose twin of the run farther from 1
\newcommand{\MlstMlpCmpQcdCCoarseMassOverCoarse}{\pending{no run yet}}  % experiments/FIGS/data/v2/index.json :: nonlinear probe qcd_c_vs_light: (1 - AUC) of R16_Q1_MASS over R16_Q1, run by run :: sha256 c6c4138ee5c86423
\newcommand{\MlstMlpCmpQcdCCoarseMassOverCoarsePr}{\pending{no run yet}}  % experiments/FIGS/data/v2/index.json :: nonlinear probe qcd_c_vs_light: (1 - AUC) of R16_Q1_MASS over R16_Q1, run by run :: sha256 c6c4138ee5c86423 :: prose twin: two significant figures, direction in words
\newcommand{\MlstMlpCmpQcdCCoarseMassOverCoarsePrMax}{\pending{no run yet}}  % experiments/FIGS/data/v2/index.json :: nonlinear probe qcd_c_vs_light: (1 - AUC) of R16_Q1_MASS over R16_Q1, run by run :: sha256 c6c4138ee5c86423 :: prose twin of the run farther from 1
\newcommand{\MlstMlpCmpQcdCFzeroOverFull}{1.234 and 1.211}  % experiments/FIGS/data/v2/index.json :: nonlinear probe qcd_c_vs_light: (1 - AUC) of FLAV_F0 over L188, run by run :: sha256 c6c4138ee5c86423
\newcommand{\MlstMlpCmpQcdCFzeroOverFullPr}{at least 21\% more errors}  % experiments/FIGS/data/v2/index.json :: nonlinear probe qcd_c_vs_light: (1 - AUC) of FLAV_F0 over L188, run by run :: sha256 c6c4138ee5c86423 :: prose twin: two significant figures, direction in words
\newcommand{\MlstMlpCmpQcdCFzeroOverFullPrMax}{up to 23\% more errors}  % experiments/FIGS/data/v2/index.json :: nonlinear probe qcd_c_vs_light: (1 - AUC) of FLAV_F0 over L188, run by run :: sha256 c6c4138ee5c86423 :: prose twin of the run farther from 1
\newcommand{\MlstMlpCmpQcdCFoneOverFzero}{0.9395 and 0.9472}  % experiments/FIGS/data/v2/index.json :: nonlinear probe qcd_c_vs_light: (1 - AUC) of FLAV_F1 over FLAV_F0, run by run :: sha256 c6c4138ee5c86423
\newcommand{\MlstMlpCmpQcdCFoneOverFzeroPr}{at least 5.3\% fewer errors}  % experiments/FIGS/data/v2/index.json :: nonlinear probe qcd_c_vs_light: (1 - AUC) of FLAV_F1 over FLAV_F0, run by run :: sha256 c6c4138ee5c86423 :: prose twin: two significant figures, direction in words
\newcommand{\MlstMlpCmpQcdCFoneOverFzeroPrMax}{up to 6.0\% fewer errors}  % experiments/FIGS/data/v2/index.json :: nonlinear probe qcd_c_vs_light: (1 - AUC) of FLAV_F1 over FLAV_F0, run by run :: sha256 c6c4138ee5c86423 :: prose twin of the run farther from 1
\newcommand{\MlstMlpCmpQcdCFonerOverFone}{0.9846 and 0.9767}  % experiments/FIGS/data/v2/index.json :: nonlinear probe qcd_c_vs_light: (1 - AUC) of FLAV_F1R over FLAV_F1, run by run :: sha256 c6c4138ee5c86423
\newcommand{\MlstMlpCmpQcdCFonerOverFonePr}{at least 1.5\% fewer errors}  % experiments/FIGS/data/v2/index.json :: nonlinear probe qcd_c_vs_light: (1 - AUC) of FLAV_F1R over FLAV_F1, run by run :: sha256 c6c4138ee5c86423 :: prose twin: two significant figures, direction in words
\newcommand{\MlstMlpCmpQcdCFonerOverFonePrMax}{up to 2.3\% fewer errors}  % experiments/FIGS/data/v2/index.json :: nonlinear probe qcd_c_vs_light: (1 - AUC) of FLAV_F1R over FLAV_F1, run by run :: sha256 c6c4138ee5c86423 :: prose twin of the run farther from 1
\newcommand{\MlstMlpCmpQcdCFonerOverFzero}{0.9251 and 0.9251}  % experiments/FIGS/data/v2/index.json :: nonlinear probe qcd_c_vs_light: (1 - AUC) of FLAV_F1R over FLAV_F0, run by run :: sha256 c6c4138ee5c86423
\newcommand{\MlstMlpCmpQcdCFonerOverFzeroPr}{at least 7.5\% fewer errors}  % experiments/FIGS/data/v2/index.json :: nonlinear probe qcd_c_vs_light: (1 - AUC) of FLAV_F1R over FLAV_F0, run by run :: sha256 c6c4138ee5c86423 :: prose twin: two significant figures, direction in words
\newcommand{\MlstMlpCmpQcdCFonerOverFzeroPrMax}{up to 7.5\% fewer errors}  % experiments/FIGS/data/v2/index.json :: nonlinear probe qcd_c_vs_light: (1 - AUC) of FLAV_F1R over FLAV_F0, run by run :: sha256 c6c4138ee5c86423 :: prose twin of the run farther from 1
\newcommand{\MlstMlpCmpQcdCFzeroOverCoarse}{0.968 and 0.956}  % experiments/FIGS/data/v2/index.json :: nonlinear probe qcd_c_vs_light: (1 - AUC) of FLAV_F0 over R16_Q1, run by run :: sha256 c6c4138ee5c86423
\newcommand{\MlstMlpCmpQcdCFzeroOverCoarsePr}{at least 3.2\% fewer errors}  % experiments/FIGS/data/v2/index.json :: nonlinear probe qcd_c_vs_light: (1 - AUC) of FLAV_F0 over R16_Q1, run by run :: sha256 c6c4138ee5c86423 :: prose twin: two significant figures, direction in words
\newcommand{\MlstMlpCmpQcdCFzeroOverCoarsePrMax}{up to 4.4\% fewer errors}  % experiments/FIGS/data/v2/index.json :: nonlinear probe qcd_c_vs_light: (1 - AUC) of FLAV_F0 over R16_Q1, run by run :: sha256 c6c4138ee5c86423 :: prose twin of the run farther from 1
\newcommand{\MlstMlpCmpQcdCRandOverFull}{\pending{no run yet}}  % experiments/FIGS/data/v2/index.json :: nonlinear probe qcd_c_vs_light: (1 - AUC) of each random partition over the 188-class run of its seed :: sha256 c6c4138ee5c86423
\newcommand{\MlstMlpCmpQcdCRandOverFullPr}{\pending{no run yet}}  % experiments/FIGS/data/v2/index.json :: nonlinear probe qcd_c_vs_light: (1 - AUC) of each random partition over the 188-class run of its seed :: sha256 c6c4138ee5c86423 :: prose twin: two significant figures, direction in words
\newcommand{\MlstMlpTopRthirtyFullOverInitBn}{1.00\,\ensuremath{\pm}\,0.31}  % experiments/FIGS/data/v2/index.json :: nonlinear probe top, missing inputs zero: R30 of L188 over the BatchNorm-recomputed random network of the same index :: sha256 c6c4138ee5c86423
\newcommand{\MlstMlpTopRthirtyFullOverInitBnPr}{equal within 40\%}  % experiments/FIGS/data/v2/index.json :: nonlinear probe top, missing inputs zero: R30 of L188 over the BatchNorm-recomputed random network of the same index :: sha256 c6c4138ee5c86423 :: prose twin: two significant figures, direction in words
\newcommand{\MlstMlpTopRthirtyResOverInitBn}{0.94\,\ensuremath{\pm}\,0.26}  % experiments/FIGS/data/v2/index.json :: nonlinear probe top, missing inputs zero: R30 of L162 over the BatchNorm-recomputed random network of the same index :: sha256 c6c4138ee5c86423
\newcommand{\MlstMlpTopRthirtyResOverInitBnPr}{equal within 30\%}  % experiments/FIGS/data/v2/index.json :: nonlinear probe top, missing inputs zero: R30 of L162 over the BatchNorm-recomputed random network of the same index :: sha256 c6c4138ee5c86423 :: prose twin: two significant figures, direction in words
\newcommand{\MlstMlpTopRthirtyMidOverInitBn}{0.84\,\ensuremath{\pm}\,0.17}  % experiments/FIGS/data/v2/index.json :: nonlinear probe top, missing inputs zero: R30 of R42_Q1 over the BatchNorm-recomputed random network of the same index :: sha256 c6c4138ee5c86423
\newcommand{\MlstMlpTopRthirtyMidOverInitBnPr}{equal within 20\%}  % experiments/FIGS/data/v2/index.json :: nonlinear probe top, missing inputs zero: R30 of R42_Q1 over the BatchNorm-recomputed random network of the same index :: sha256 c6c4138ee5c86423 :: prose twin: two significant figures, direction in words
\newcommand{\MlstMlpTopRthirtyCoarseOverInitBn}{0.75\,\ensuremath{\pm}\,0.05}  % experiments/FIGS/data/v2/index.json :: nonlinear probe top, missing inputs zero: R30 of R16_Q1 over the BatchNorm-recomputed random network of the same index :: sha256 c6c4138ee5c86423
\newcommand{\MlstMlpTopRthirtyCoarseOverInitBnPr}{25\% less rejection}  % experiments/FIGS/data/v2/index.json :: nonlinear probe top, missing inputs zero: R30 of R16_Q1 over the BatchNorm-recomputed random network of the same index :: sha256 c6c4138ee5c86423 :: prose twin: two significant figures, direction in words
\newcommand{\MlstMlpQgRthirtyFullOverInitBn}{1.083\,\ensuremath{\pm}\,0.026}  % experiments/FIGS/data/v2/index.json :: nonlinear probe qg, missing inputs zero: R30 of L188 over the BatchNorm-recomputed random network of the same index :: sha256 c6c4138ee5c86423
\newcommand{\MlstMlpQgRthirtyFullOverInitBnPr}{8.3\% more rejection}  % experiments/FIGS/data/v2/index.json :: nonlinear probe qg, missing inputs zero: R30 of L188 over the BatchNorm-recomputed random network of the same index :: sha256 c6c4138ee5c86423 :: prose twin: two significant figures, direction in words
\newcommand{\MlstMlpQgRthirtyResOverInitBn}{1.010\,\ensuremath{\pm}\,0.016}  % experiments/FIGS/data/v2/index.json :: nonlinear probe qg, missing inputs zero: R30 of L162 over the BatchNorm-recomputed random network of the same index :: sha256 c6c4138ee5c86423
\newcommand{\MlstMlpQgRthirtyResOverInitBnPr}{equal within 2\%}  % experiments/FIGS/data/v2/index.json :: nonlinear probe qg, missing inputs zero: R30 of L162 over the BatchNorm-recomputed random network of the same index :: sha256 c6c4138ee5c86423 :: prose twin: two significant figures, direction in words
\newcommand{\MlstMlpQgRthirtyMidOverInitBn}{1.06\,\ensuremath{\pm}\,0.05}  % experiments/FIGS/data/v2/index.json :: nonlinear probe qg, missing inputs zero: R30 of R42_Q1 over the BatchNorm-recomputed random network of the same index :: sha256 c6c4138ee5c86423
\newcommand{\MlstMlpQgRthirtyMidOverInitBnPr}{equal within 7\%}  % experiments/FIGS/data/v2/index.json :: nonlinear probe qg, missing inputs zero: R30 of R42_Q1 over the BatchNorm-recomputed random network of the same index :: sha256 c6c4138ee5c86423 :: prose twin: two significant figures, direction in words
\newcommand{\MlstMlpQgRthirtyCoarseOverInitBn}{1.04\,\ensuremath{\pm}\,0.04}  % experiments/FIGS/data/v2/index.json :: nonlinear probe qg, missing inputs zero: R30 of R16_Q1 over the BatchNorm-recomputed random network of the same index :: sha256 c6c4138ee5c86423
\newcommand{\MlstMlpQgRthirtyCoarseOverInitBnPr}{equal within 5\%}  % experiments/FIGS/data/v2/index.json :: nonlinear probe qg, missing inputs zero: R30 of R16_Q1 over the BatchNorm-recomputed random network of the same index :: sha256 c6c4138ee5c86423 :: prose twin: two significant figures, direction in words
\newcommand{\MlstMlpQgHerwigRthirtyFullOverInitBn}{0.971\,\ensuremath{\pm}\,0.034}  % experiments/FIGS/data/v2/index.json :: nonlinear probe qg_herwig, missing inputs zero: R30 of L188 over the BatchNorm-recomputed random network of the same index :: sha256 c6c4138ee5c86423
\newcommand{\MlstMlpQgHerwigRthirtyFullOverInitBnPr}{equal within 4\%}  % experiments/FIGS/data/v2/index.json :: nonlinear probe qg_herwig, missing inputs zero: R30 of L188 over the BatchNorm-recomputed random network of the same index :: sha256 c6c4138ee5c86423 :: prose twin: two significant figures, direction in words
\newcommand{\MlstMlpQgHerwigRthirtyResOverInitBn}{0.94\,\ensuremath{\pm}\,0.05}  % experiments/FIGS/data/v2/index.json :: nonlinear probe qg_herwig, missing inputs zero: R30 of L162 over the BatchNorm-recomputed random network of the same index :: sha256 c6c4138ee5c86423
\newcommand{\MlstMlpQgHerwigRthirtyResOverInitBnPr}{equal within 6\%}  % experiments/FIGS/data/v2/index.json :: nonlinear probe qg_herwig, missing inputs zero: R30 of L162 over the BatchNorm-recomputed random network of the same index :: sha256 c6c4138ee5c86423 :: prose twin: two significant figures, direction in words
\newcommand{\MlstMlpQgHerwigRthirtyMidOverInitBn}{0.945\,\ensuremath{\pm}\,0.030}  % experiments/FIGS/data/v2/index.json :: nonlinear probe qg_herwig, missing inputs zero: R30 of R42_Q1 over the BatchNorm-recomputed random network of the same index :: sha256 c6c4138ee5c86423
\newcommand{\MlstMlpQgHerwigRthirtyMidOverInitBnPr}{equal within 6\%}  % experiments/FIGS/data/v2/index.json :: nonlinear probe qg_herwig, missing inputs zero: R30 of R42_Q1 over the BatchNorm-recomputed random network of the same index :: sha256 c6c4138ee5c86423 :: prose twin: two significant figures, direction in words
\newcommand{\MlstMlpQgHerwigRthirtyCoarseOverInitBn}{0.96\,\ensuremath{\pm}\,0.04}  % experiments/FIGS/data/v2/index.json :: nonlinear probe qg_herwig, missing inputs zero: R30 of R16_Q1 over the BatchNorm-recomputed random network of the same index :: sha256 c6c4138ee5c86423
\newcommand{\MlstMlpQgHerwigRthirtyCoarseOverInitBnPr}{equal within 5\%}  % experiments/FIGS/data/v2/index.json :: nonlinear probe qg_herwig, missing inputs zero: R30 of R16_Q1 over the BatchNorm-recomputed random network of the same index :: sha256 c6c4138ee5c86423 :: prose twin: two significant figures, direction in words
\newcommand{\MlstMlpNatJcTwoResMassOverRes}{1.016 and 1.009}  % experiments/FIGS/data/v2/index.json :: nonlinear probe, natural-composition JetClass-II (fits_v3_mlp2): (1 - AUC) of L162_MASS over L162, run by run :: sha256 c6c4138ee5c86423
\newcommand{\MlstMlpNatJcTwoResMassOverResPr}{equal within 2\%}  % experiments/FIGS/data/v2/index.json :: nonlinear probe, natural-composition JetClass-II (fits_v3_mlp2): (1 - AUC) of L162_MASS over L162, run by run :: sha256 c6c4138ee5c86423 :: prose twin: two significant figures, direction in words
\newcommand{\MlstMlpNatJcTwoResMassOverResPrMax}{equal within 2\%}  % experiments/FIGS/data/v2/index.json :: nonlinear probe, natural-composition JetClass-II (fits_v3_mlp2): (1 - AUC) of L162_MASS over L162, run by run :: sha256 c6c4138ee5c86423 :: prose twin of the run farther from 1
\newcommand{\MlstMlpNatJcTwoCoarseMassOverCoarse}{1.026 and 1.019}  % experiments/FIGS/data/v2/index.json :: nonlinear probe, natural-composition JetClass-II (fits_v3_mlp2): (1 - AUC) of R16_Q1_MASS over R16_Q1, run by run :: sha256 c6c4138ee5c86423
\newcommand{\MlstMlpNatJcTwoCoarseMassOverCoarsePr}{equal within 3\%}  % experiments/FIGS/data/v2/index.json :: nonlinear probe, natural-composition JetClass-II (fits_v3_mlp2): (1 - AUC) of R16_Q1_MASS over R16_Q1, run by run :: sha256 c6c4138ee5c86423 :: prose twin: two significant figures, direction in words
\newcommand{\MlstMlpNatJcTwoCoarseMassOverCoarsePrMax}{equal within 3\%}  % experiments/FIGS/data/v2/index.json :: nonlinear probe, natural-composition JetClass-II (fits_v3_mlp2): (1 - AUC) of R16_Q1_MASS over R16_Q1, run by run :: sha256 c6c4138ee5c86423 :: prose twin of the run farther from 1
\newcommand{\LpSubstitutionTop}{random draws}  % experiments/FIGS/data/v2/index.json :: linprobe/fits_v4_fill: lowest summed validation cross-entropy over every model, largest training set :: sha256 c6c4138ee5c86423 :: zero 8.8428; mean 8.99394; draws 7.00838
\newcommand{\LpSubstitutionQg}{the training mean}  % experiments/FIGS/data/v2/index.json :: linprobe/fits_v4_fill: lowest summed validation cross-entropy over every model, largest training set :: sha256 c6c4138ee5c86423 :: zero 13.6224; mean 13.1964; draws 13.3584
\newcommand{\MlstLpTopRthirtyFullOverInitBn}{4.67\,\ensuremath{\pm}\,0.23}  % experiments/FIGS/data/v2/index.json :: linear probe top, largest set, chosen substitution: R30 of L188 over the BatchNorm-recomputed random network of the same index :: sha256 c6c4138ee5c86423
\newcommand{\MlstLpTopRthirtyFullOverInitBnPr}{4.7 times the rejection}  % experiments/FIGS/data/v2/index.json :: linear probe top, largest set, chosen substitution: R30 of L188 over the BatchNorm-recomputed random network of the same index :: sha256 c6c4138ee5c86423 :: prose twin: two significant figures, direction in words
\newcommand{\MlstLpTopRthirtyFullOverInitBnZero}{1.04\,\ensuremath{\pm}\,0.05}  % experiments/FIGS/data/v2/index.json :: linear probe top, largest set, zero: R30 of L188 over the BatchNorm-recomputed random network of the same index :: sha256 c6c4138ee5c86423
\newcommand{\MlstLpTopRthirtyFullOverInitBnZeroPr}{equal within 5\%}  % experiments/FIGS/data/v2/index.json :: linear probe top, largest set, zero: R30 of L188 over the BatchNorm-recomputed random network of the same index :: sha256 c6c4138ee5c86423 :: prose twin: two significant figures, direction in words
\newcommand{\MlstLpTopRthirtyResOverInitBn}{4.08\,\ensuremath{\pm}\,0.28}  % experiments/FIGS/data/v2/index.json :: linear probe top, largest set, chosen substitution: R30 of L162 over the BatchNorm-recomputed random network of the same index :: sha256 c6c4138ee5c86423
\newcommand{\MlstLpTopRthirtyResOverInitBnPr}{4.1 times the rejection}  % experiments/FIGS/data/v2/index.json :: linear probe top, largest set, chosen substitution: R30 of L162 over the BatchNorm-recomputed random network of the same index :: sha256 c6c4138ee5c86423 :: prose twin: two significant figures, direction in words
\newcommand{\MlstLpTopRthirtyResOverInitBnZero}{1.33\,\ensuremath{\pm}\,0.30}  % experiments/FIGS/data/v2/index.json :: linear probe top, largest set, zero: R30 of L162 over the BatchNorm-recomputed random network of the same index :: sha256 c6c4138ee5c86423
\newcommand{\MlstLpTopRthirtyResOverInitBnZeroPr}{equal within 40\%}  % experiments/FIGS/data/v2/index.json :: linear probe top, largest set, zero: R30 of L162 over the BatchNorm-recomputed random network of the same index :: sha256 c6c4138ee5c86423 :: prose twin: two significant figures, direction in words
\newcommand{\MlstLpTopRthirtyMidOverInitBn}{4.76\,\ensuremath{\pm}\,0.15}  % experiments/FIGS/data/v2/index.json :: linear probe top, largest set, chosen substitution: R30 of R42_Q1 over the BatchNorm-recomputed random network of the same index :: sha256 c6c4138ee5c86423
\newcommand{\MlstLpTopRthirtyMidOverInitBnPr}{4.8 times the rejection}  % experiments/FIGS/data/v2/index.json :: linear probe top, largest set, chosen substitution: R30 of R42_Q1 over the BatchNorm-recomputed random network of the same index :: sha256 c6c4138ee5c86423 :: prose twin: two significant figures, direction in words
\newcommand{\MlstLpTopRthirtyMidOverInitBnZero}{1.01\,\ensuremath{\pm}\,0.28}  % experiments/FIGS/data/v2/index.json :: linear probe top, largest set, zero: R30 of R42_Q1 over the BatchNorm-recomputed random network of the same index :: sha256 c6c4138ee5c86423
\newcommand{\MlstLpTopRthirtyMidOverInitBnZeroPr}{equal within 30\%}  % experiments/FIGS/data/v2/index.json :: linear probe top, largest set, zero: R30 of R42_Q1 over the BatchNorm-recomputed random network of the same index :: sha256 c6c4138ee5c86423 :: prose twin: two significant figures, direction in words
\newcommand{\MlstLpTopRthirtyCoarseOverInitBn}{5.24\,\ensuremath{\pm}\,0.06}  % experiments/FIGS/data/v2/index.json :: linear probe top, largest set, chosen substitution: R30 of R16_Q1 over the BatchNorm-recomputed random network of the same index :: sha256 c6c4138ee5c86423
\newcommand{\MlstLpTopRthirtyCoarseOverInitBnPr}{5.2 times the rejection}  % experiments/FIGS/data/v2/index.json :: linear probe top, largest set, chosen substitution: R30 of R16_Q1 over the BatchNorm-recomputed random network of the same index :: sha256 c6c4138ee5c86423 :: prose twin: two significant figures, direction in words
\newcommand{\MlstLpTopRthirtyCoarseOverInitBnZero}{0.97\,\ensuremath{\pm}\,0.16}  % experiments/FIGS/data/v2/index.json :: linear probe top, largest set, zero: R30 of R16_Q1 over the BatchNorm-recomputed random network of the same index :: sha256 c6c4138ee5c86423
\newcommand{\MlstLpTopRthirtyCoarseOverInitBnZeroPr}{equal within 20\%}  % experiments/FIGS/data/v2/index.json :: linear probe top, largest set, zero: R30 of R16_Q1 over the BatchNorm-recomputed random network of the same index :: sha256 c6c4138ee5c86423 :: prose twin: two significant figures, direction in words
\newcommand{\MlstLpTopRthirtyResMassOverInitBn}{5.71 and 6.25}  % experiments/FIGS/data/v2/index.json :: linear probe top, largest set, chosen substitution: R30 of L162_MASS over the BatchNorm-recomputed random network of the same index :: sha256 c6c4138ee5c86423
\newcommand{\MlstLpTopRthirtyResMassOverInitBnPr}{at least 5.7 times the rejection}  % experiments/FIGS/data/v2/index.json :: linear probe top, largest set, chosen substitution: R30 of L162_MASS over the BatchNorm-recomputed random network of the same index :: sha256 c6c4138ee5c86423 :: prose twin: two significant figures, direction in words
\newcommand{\MlstLpTopRthirtyResMassOverInitBnPrMax}{up to 6.3 times the rejection}  % experiments/FIGS/data/v2/index.json :: linear probe top, largest set, chosen substitution: R30 of L162_MASS over the BatchNorm-recomputed random network of the same index :: sha256 c6c4138ee5c86423 :: prose twin of the run farther from 1
\newcommand{\MlstLpTopRthirtyResMassOverInitBnZero}{1.683 and 1.678}  % experiments/FIGS/data/v2/index.json :: linear probe top, largest set, zero: R30 of L162_MASS over the BatchNorm-recomputed random network of the same index :: sha256 c6c4138ee5c86423
\newcommand{\MlstLpTopRthirtyResMassOverInitBnZeroPr}{at least 68\% more rejection}  % experiments/FIGS/data/v2/index.json :: linear probe top, largest set, zero: R30 of L162_MASS over the BatchNorm-recomputed random network of the same index :: sha256 c6c4138ee5c86423 :: prose twin: two significant figures, direction in words
\newcommand{\MlstLpTopRthirtyResMassOverInitBnZeroPrMax}{up to 68\% more rejection}  % experiments/FIGS/data/v2/index.json :: linear probe top, largest set, zero: R30 of L162_MASS over the BatchNorm-recomputed random network of the same index :: sha256 c6c4138ee5c86423 :: prose twin of the run farther from 1
\newcommand{\MlstLpTopRthirtyCoarseMassOverInitBn}{6.73 and 7.72}  % experiments/FIGS/data/v2/index.json :: linear probe top, largest set, chosen substitution: R30 of R16_Q1_MASS over the BatchNorm-recomputed random network of the same index :: sha256 c6c4138ee5c86423
\newcommand{\MlstLpTopRthirtyCoarseMassOverInitBnPr}{at least 6.7 times the rejection}  % experiments/FIGS/data/v2/index.json :: linear probe top, largest set, chosen substitution: R30 of R16_Q1_MASS over the BatchNorm-recomputed random network of the same index :: sha256 c6c4138ee5c86423 :: prose twin: two significant figures, direction in words
\newcommand{\MlstLpTopRthirtyCoarseMassOverInitBnPrMax}{up to 7.7 times the rejection}  % experiments/FIGS/data/v2/index.json :: linear probe top, largest set, chosen substitution: R30 of R16_Q1_MASS over the BatchNorm-recomputed random network of the same index :: sha256 c6c4138ee5c86423 :: prose twin of the run farther from 1
\newcommand{\MlstLpTopRthirtyCoarseMassOverInitBnZero}{2.3 and 1.2}  % experiments/FIGS/data/v2/index.json :: linear probe top, largest set, zero: R30 of R16_Q1_MASS over the BatchNorm-recomputed random network of the same index :: sha256 c6c4138ee5c86423
\newcommand{\MlstLpTopRthirtyCoarseMassOverInitBnZeroPr}{equal within 200\%}  % experiments/FIGS/data/v2/index.json :: linear probe top, largest set, zero: R30 of R16_Q1_MASS over the BatchNorm-recomputed random network of the same index :: sha256 c6c4138ee5c86423 :: prose twin: two significant figures, direction in words
\newcommand{\MlstLpTopRthirtyCoarseMassOverInitBnZeroPrMax}{equal within 200\%}  % experiments/FIGS/data/v2/index.json :: linear probe top, largest set, zero: R30 of R16_Q1_MASS over the BatchNorm-recomputed random network of the same index :: sha256 c6c4138ee5c86423 :: prose twin of the run farther from 1
\newcommand{\MlstLpTopRthirtyResOverFull}{0.874\,\ensuremath{\pm}\,0.018}  % experiments/FIGS/data/v2/index.json :: linear probe top: R30 of L162 over the 188-class run :: sha256 c6c4138ee5c86423
\newcommand{\MlstLpTopRthirtyResOverFullPr}{13\% less rejection}  % experiments/FIGS/data/v2/index.json :: linear probe top: R30 of L162 over the 188-class run :: sha256 c6c4138ee5c86423 :: prose twin: two significant figures, direction in words
\newcommand{\MlstLpTopRthirtyResOverFullZero}{1.29\,\ensuremath{\pm}\,0.31}  % experiments/FIGS/data/v2/index.json :: linear probe top: R30 of L162 over the 188-class run :: sha256 c6c4138ee5c86423
\newcommand{\MlstLpTopRthirtyResOverFullZeroPr}{equal within 40\%}  % experiments/FIGS/data/v2/index.json :: linear probe top: R30 of L162 over the 188-class run :: sha256 c6c4138ee5c86423 :: prose twin: two significant figures, direction in words
\newcommand{\MlstLpTopRthirtyMidOverFull}{1.02\,\ensuremath{\pm}\,0.08}  % experiments/FIGS/data/v2/index.json :: linear probe top: R30 of R42_Q1 over the 188-class run :: sha256 c6c4138ee5c86423
\newcommand{\MlstLpTopRthirtyMidOverFullPr}{equal within 9\%}  % experiments/FIGS/data/v2/index.json :: linear probe top: R30 of R42_Q1 over the 188-class run :: sha256 c6c4138ee5c86423 :: prose twin: two significant figures, direction in words
\newcommand{\MlstLpTopRthirtyMidOverFullZero}{0.98\,\ensuremath{\pm}\,0.28}  % experiments/FIGS/data/v2/index.json :: linear probe top: R30 of R42_Q1 over the 188-class run :: sha256 c6c4138ee5c86423
\newcommand{\MlstLpTopRthirtyMidOverFullZeroPr}{equal within 30\%}  % experiments/FIGS/data/v2/index.json :: linear probe top: R30 of R42_Q1 over the 188-class run :: sha256 c6c4138ee5c86423 :: prose twin: two significant figures, direction in words
\newcommand{\MlstLpTopRthirtyCoarseOverFull}{1.12\,\ensuremath{\pm}\,0.06}  % experiments/FIGS/data/v2/index.json :: linear probe top: R30 of R16_Q1 over the 188-class run :: sha256 c6c4138ee5c86423
\newcommand{\MlstLpTopRthirtyCoarseOverFullPr}{12\% more rejection}  % experiments/FIGS/data/v2/index.json :: linear probe top: R30 of R16_Q1 over the 188-class run :: sha256 c6c4138ee5c86423 :: prose twin: two significant figures, direction in words
\newcommand{\MlstLpTopRthirtyCoarseOverFullZero}{0.93\,\ensuremath{\pm}\,0.15}  % experiments/FIGS/data/v2/index.json :: linear probe top: R30 of R16_Q1 over the 188-class run :: sha256 c6c4138ee5c86423
\newcommand{\MlstLpTopRthirtyCoarseOverFullZeroPr}{equal within 20\%}  % experiments/FIGS/data/v2/index.json :: linear probe top: R30 of R16_Q1 over the 188-class run :: sha256 c6c4138ee5c86423 :: prose twin: two significant figures, direction in words
\newcommand{\MlstLpTopRthirtyMeanOverZeroFull}{1.25\,\ensuremath{\pm}\,0.27}  % experiments/FIGS/data/v2/index.json :: linear probe top: R30 with mean over R30 with zero, L188 :: sha256 c6c4138ee5c86423
\newcommand{\MlstLpTopRthirtyMeanOverZeroFullPr}{equal within 30\%}  % experiments/FIGS/data/v2/index.json :: linear probe top: R30 with mean over R30 with zero, L188 :: sha256 c6c4138ee5c86423 :: prose twin: two significant figures, direction in words
\newcommand{\MlstLpTopRthirtyDrawOverZeroFull}{3.9\,\ensuremath{\pm}\,0.4}  % experiments/FIGS/data/v2/index.json :: linear probe top: R30 with draws over R30 with zero, L188 :: sha256 c6c4138ee5c86423
\newcommand{\MlstLpTopRthirtyDrawOverZeroFullPr}{3.9 times the rejection}  % experiments/FIGS/data/v2/index.json :: linear probe top: R30 with draws over R30 with zero, L188 :: sha256 c6c4138ee5c86423 :: prose twin: two significant figures, direction in words
\newcommand{\MlstLpTopRthirtyMeanOverZeroInitBn}{1.08\,\ensuremath{\pm}\,0.20}  % experiments/FIGS/data/v2/index.json :: linear probe top: R30 with mean over R30 with zero, INIT_BN :: sha256 c6c4138ee5c86423
\newcommand{\MlstLpTopRthirtyMeanOverZeroInitBnPr}{equal within 30\%}  % experiments/FIGS/data/v2/index.json :: linear probe top: R30 with mean over R30 with zero, INIT_BN :: sha256 c6c4138ee5c86423 :: prose twin: two significant figures, direction in words
\newcommand{\MlstLpTopRthirtyDrawOverZeroInitBn}{0.86\,\ensuremath{\pm}\,0.14}  % experiments/FIGS/data/v2/index.json :: linear probe top: R30 with draws over R30 with zero, INIT_BN :: sha256 c6c4138ee5c86423
\newcommand{\MlstLpTopRthirtyDrawOverZeroInitBnPr}{equal within 20\%}  % experiments/FIGS/data/v2/index.json :: linear probe top: R30 with draws over R30 with zero, INIT_BN :: sha256 c6c4138ee5c86423 :: prose twin: two significant figures, direction in words
\newcommand{\MlstLpTopRthirtyKinOverFull}{3.88}  % experiments/FIGS/data/v2/index.json :: linear probe top: R30 of the kinematic-only model over the mean of the 188-class runs (zero) :: sha256 c6c4138ee5c86423
\newcommand{\MlstLpTopRthirtyKinOverFullPr}{3.9 times the rejection}  % experiments/FIGS/data/v2/index.json :: linear probe top: R30 of the kinematic-only model over the mean of the 188-class runs (zero) :: sha256 c6c4138ee5c86423 :: prose twin: two significant figures, direction in words
\newcommand{\MlstLpTopRthirtyKin}{231}  % experiments/FIGS/data/v2/index.json :: linear probe top, largest set: R30 of the kinematic-only model (one run) :: sha256 c6c4138ee5c86423
\newcommand{\MlstLpTopRthirtyFullZero}{60\,\ensuremath{\pm}\,12}  % experiments/FIGS/data/v2/index.json :: linear probe top, largest set: R30 with zero, L188 :: sha256 c6c4138ee5c86423
\newcommand{\MlstLpTopRthirtyFullMean}{73\,\ensuremath{\pm}\,15}  % experiments/FIGS/data/v2/index.json :: linear probe top, largest set: R30 with mean, L188 :: sha256 c6c4138ee5c86423
\newcommand{\MlstLpTopRthirtyFullDraws}{228\,\ensuremath{\pm}\,25}  % experiments/FIGS/data/v2/index.json :: linear probe top, largest set: R30 with draws, L188 :: sha256 c6c4138ee5c86423
\newcommand{\MlstLpTopRthirtyInitBnZero}{58\,\ensuremath{\pm}\,13}  % experiments/FIGS/data/v2/index.json :: linear probe top, largest set: R30 with zero, INIT_BN :: sha256 c6c4138ee5c86423
\newcommand{\MlstLpTopRthirtyInitBnMean}{60.4\,\ensuremath{\pm}\,2.4}  % experiments/FIGS/data/v2/index.json :: linear probe top, largest set: R30 with mean, INIT_BN :: sha256 c6c4138ee5c86423
\newcommand{\MlstLpTopRthirtyInitBnDraws}{48.6\,\ensuremath{\pm}\,3.0}  % experiments/FIGS/data/v2/index.json :: linear probe top, largest set: R30 with draws, INIT_BN :: sha256 c6c4138ee5c86423
\newcommand{\MlstLpQgRthirtyFullOverInitBn}{1.39\,\ensuremath{\pm}\,0.06}  % experiments/FIGS/data/v2/index.json :: linear probe qg, largest set, chosen substitution: R30 of L188 over the BatchNorm-recomputed random network of the same index :: sha256 c6c4138ee5c86423
\newcommand{\MlstLpQgRthirtyFullOverInitBnPr}{39\% more rejection}  % experiments/FIGS/data/v2/index.json :: linear probe qg, largest set, chosen substitution: R30 of L188 over the BatchNorm-recomputed random network of the same index :: sha256 c6c4138ee5c86423 :: prose twin: two significant figures, direction in words
\newcommand{\MlstLpQgRthirtyFullOverInitBnZero}{1.21\,\ensuremath{\pm}\,0.15}  % experiments/FIGS/data/v2/index.json :: linear probe qg, largest set, zero: R30 of L188 over the BatchNorm-recomputed random network of the same index :: sha256 c6c4138ee5c86423
\newcommand{\MlstLpQgRthirtyFullOverInitBnZeroPr}{equal within 30\%}  % experiments/FIGS/data/v2/index.json :: linear probe qg, largest set, zero: R30 of L188 over the BatchNorm-recomputed random network of the same index :: sha256 c6c4138ee5c86423 :: prose twin: two significant figures, direction in words
\newcommand{\MlstLpQgRthirtyResOverInitBn}{1.360\,\ensuremath{\pm}\,0.025}  % experiments/FIGS/data/v2/index.json :: linear probe qg, largest set, chosen substitution: R30 of L162 over the BatchNorm-recomputed random network of the same index :: sha256 c6c4138ee5c86423
\newcommand{\MlstLpQgRthirtyResOverInitBnPr}{36\% more rejection}  % experiments/FIGS/data/v2/index.json :: linear probe qg, largest set, chosen substitution: R30 of L162 over the BatchNorm-recomputed random network of the same index :: sha256 c6c4138ee5c86423 :: prose twin: two significant figures, direction in words
\newcommand{\MlstLpQgRthirtyResOverInitBnZero}{1.07\,\ensuremath{\pm}\,0.13}  % experiments/FIGS/data/v2/index.json :: linear probe qg, largest set, zero: R30 of L162 over the BatchNorm-recomputed random network of the same index :: sha256 c6c4138ee5c86423
\newcommand{\MlstLpQgRthirtyResOverInitBnZeroPr}{equal within 20\%}  % experiments/FIGS/data/v2/index.json :: linear probe qg, largest set, zero: R30 of L162 over the BatchNorm-recomputed random network of the same index :: sha256 c6c4138ee5c86423 :: prose twin: two significant figures, direction in words
\newcommand{\MlstLpQgRthirtyMidOverInitBn}{1.357\,\ensuremath{\pm}\,0.033}  % experiments/FIGS/data/v2/index.json :: linear probe qg, largest set, chosen substitution: R30 of R42_Q1 over the BatchNorm-recomputed random network of the same index :: sha256 c6c4138ee5c86423
\newcommand{\MlstLpQgRthirtyMidOverInitBnPr}{36\% more rejection}  % experiments/FIGS/data/v2/index.json :: linear probe qg, largest set, chosen substitution: R30 of R42_Q1 over the BatchNorm-recomputed random network of the same index :: sha256 c6c4138ee5c86423 :: prose twin: two significant figures, direction in words
\newcommand{\MlstLpQgRthirtyMidOverInitBnZero}{1.12\,\ensuremath{\pm}\,0.11}  % experiments/FIGS/data/v2/index.json :: linear probe qg, largest set, zero: R30 of R42_Q1 over the BatchNorm-recomputed random network of the same index :: sha256 c6c4138ee5c86423
\newcommand{\MlstLpQgRthirtyMidOverInitBnZeroPr}{equal within 20\%}  % experiments/FIGS/data/v2/index.json :: linear probe qg, largest set, zero: R30 of R42_Q1 over the BatchNorm-recomputed random network of the same index :: sha256 c6c4138ee5c86423 :: prose twin: two significant figures, direction in words
\newcommand{\MlstLpQgRthirtyCoarseOverInitBn}{1.42\,\ensuremath{\pm}\,0.04}  % experiments/FIGS/data/v2/index.json :: linear probe qg, largest set, chosen substitution: R30 of R16_Q1 over the BatchNorm-recomputed random network of the same index :: sha256 c6c4138ee5c86423
\newcommand{\MlstLpQgRthirtyCoarseOverInitBnPr}{42\% more rejection}  % experiments/FIGS/data/v2/index.json :: linear probe qg, largest set, chosen substitution: R30 of R16_Q1 over the BatchNorm-recomputed random network of the same index :: sha256 c6c4138ee5c86423 :: prose twin: two significant figures, direction in words
\newcommand{\MlstLpQgRthirtyCoarseOverInitBnZero}{1.15\,\ensuremath{\pm}\,0.17}  % experiments/FIGS/data/v2/index.json :: linear probe qg, largest set, zero: R30 of R16_Q1 over the BatchNorm-recomputed random network of the same index :: sha256 c6c4138ee5c86423
\newcommand{\MlstLpQgRthirtyCoarseOverInitBnZeroPr}{equal within 20\%}  % experiments/FIGS/data/v2/index.json :: linear probe qg, largest set, zero: R30 of R16_Q1 over the BatchNorm-recomputed random network of the same index :: sha256 c6c4138ee5c86423 :: prose twin: two significant figures, direction in words
\newcommand{\MlstLpQgRthirtyResMassOverInitBn}{1.430 and 1.364}  % experiments/FIGS/data/v2/index.json :: linear probe qg, largest set, chosen substitution: R30 of L162_MASS over the BatchNorm-recomputed random network of the same index :: sha256 c6c4138ee5c86423
\newcommand{\MlstLpQgRthirtyResMassOverInitBnPr}{at least 36\% more rejection}  % experiments/FIGS/data/v2/index.json :: linear probe qg, largest set, chosen substitution: R30 of L162_MASS over the BatchNorm-recomputed random network of the same index :: sha256 c6c4138ee5c86423 :: prose twin: two significant figures, direction in words
\newcommand{\MlstLpQgRthirtyResMassOverInitBnPrMax}{up to 43\% more rejection}  % experiments/FIGS/data/v2/index.json :: linear probe qg, largest set, chosen substitution: R30 of L162_MASS over the BatchNorm-recomputed random network of the same index :: sha256 c6c4138ee5c86423 :: prose twin of the run farther from 1
\newcommand{\MlstLpQgRthirtyResMassOverInitBnZero}{1.213 and 1.175}  % experiments/FIGS/data/v2/index.json :: linear probe qg, largest set, zero: R30 of L162_MASS over the BatchNorm-recomputed random network of the same index :: sha256 c6c4138ee5c86423
\newcommand{\MlstLpQgRthirtyResMassOverInitBnZeroPr}{equal within 30\%}  % experiments/FIGS/data/v2/index.json :: linear probe qg, largest set, zero: R30 of L162_MASS over the BatchNorm-recomputed random network of the same index :: sha256 c6c4138ee5c86423 :: prose twin: two significant figures, direction in words
\newcommand{\MlstLpQgRthirtyResMassOverInitBnZeroPrMax}{equal within 30\%}  % experiments/FIGS/data/v2/index.json :: linear probe qg, largest set, zero: R30 of L162_MASS over the BatchNorm-recomputed random network of the same index :: sha256 c6c4138ee5c86423 :: prose twin of the run farther from 1
\newcommand{\MlstLpQgRthirtyCoarseMassOverInitBn}{1.324 and 1.394}  % experiments/FIGS/data/v2/index.json :: linear probe qg, largest set, chosen substitution: R30 of R16_Q1_MASS over the BatchNorm-recomputed random network of the same index :: sha256 c6c4138ee5c86423
\newcommand{\MlstLpQgRthirtyCoarseMassOverInitBnPr}{at least 32\% more rejection}  % experiments/FIGS/data/v2/index.json :: linear probe qg, largest set, chosen substitution: R30 of R16_Q1_MASS over the BatchNorm-recomputed random network of the same index :: sha256 c6c4138ee5c86423 :: prose twin: two significant figures, direction in words
\newcommand{\MlstLpQgRthirtyCoarseMassOverInitBnPrMax}{up to 39\% more rejection}  % experiments/FIGS/data/v2/index.json :: linear probe qg, largest set, chosen substitution: R30 of R16_Q1_MASS over the BatchNorm-recomputed random network of the same index :: sha256 c6c4138ee5c86423 :: prose twin of the run farther from 1
\newcommand{\MlstLpQgRthirtyCoarseMassOverInitBnZero}{1.224 and 1.261}  % experiments/FIGS/data/v2/index.json :: linear probe qg, largest set, zero: R30 of R16_Q1_MASS over the BatchNorm-recomputed random network of the same index :: sha256 c6c4138ee5c86423
\newcommand{\MlstLpQgRthirtyCoarseMassOverInitBnZeroPr}{equal within 30\%}  % experiments/FIGS/data/v2/index.json :: linear probe qg, largest set, zero: R30 of R16_Q1_MASS over the BatchNorm-recomputed random network of the same index :: sha256 c6c4138ee5c86423 :: prose twin: two significant figures, direction in words
\newcommand{\MlstLpQgRthirtyCoarseMassOverInitBnZeroPrMax}{equal within 30\%}  % experiments/FIGS/data/v2/index.json :: linear probe qg, largest set, zero: R30 of R16_Q1_MASS over the BatchNorm-recomputed random network of the same index :: sha256 c6c4138ee5c86423 :: prose twin of the run farther from 1
\newcommand{\MlstLpQgRthirtyResOverFull}{0.977\,\ensuremath{\pm}\,0.028}  % experiments/FIGS/data/v2/index.json :: linear probe qg: R30 of L162 over the 188-class run :: sha256 c6c4138ee5c86423
\newcommand{\MlstLpQgRthirtyResOverFullPr}{equal within 3\%}  % experiments/FIGS/data/v2/index.json :: linear probe qg: R30 of L162 over the 188-class run :: sha256 c6c4138ee5c86423 :: prose twin: two significant figures, direction in words
\newcommand{\MlstLpQgRthirtyResOverFullZero}{0.89\,\ensuremath{\pm}\,0.20}  % experiments/FIGS/data/v2/index.json :: linear probe qg: R30 of L162 over the 188-class run :: sha256 c6c4138ee5c86423
\newcommand{\MlstLpQgRthirtyResOverFullZeroPr}{equal within 20\%}  % experiments/FIGS/data/v2/index.json :: linear probe qg: R30 of L162 over the 188-class run :: sha256 c6c4138ee5c86423 :: prose twin: two significant figures, direction in words
\newcommand{\MlstLpQgRthirtyMidOverFull}{0.974\,\ensuremath{\pm}\,0.031}  % experiments/FIGS/data/v2/index.json :: linear probe qg: R30 of R42_Q1 over the 188-class run :: sha256 c6c4138ee5c86423
\newcommand{\MlstLpQgRthirtyMidOverFullPr}{equal within 4\%}  % experiments/FIGS/data/v2/index.json :: linear probe qg: R30 of R42_Q1 over the 188-class run :: sha256 c6c4138ee5c86423 :: prose twin: two significant figures, direction in words
\newcommand{\MlstLpQgRthirtyMidOverFullZero}{0.94\,\ensuremath{\pm}\,0.18}  % experiments/FIGS/data/v2/index.json :: linear probe qg: R30 of R42_Q1 over the 188-class run :: sha256 c6c4138ee5c86423
\newcommand{\MlstLpQgRthirtyMidOverFullZeroPr}{equal within 20\%}  % experiments/FIGS/data/v2/index.json :: linear probe qg: R30 of R42_Q1 over the 188-class run :: sha256 c6c4138ee5c86423 :: prose twin: two significant figures, direction in words
\newcommand{\MlstLpQgRthirtyCoarseOverFull}{1.02\,\ensuremath{\pm}\,0.06}  % experiments/FIGS/data/v2/index.json :: linear probe qg: R30 of R16_Q1 over the 188-class run :: sha256 c6c4138ee5c86423
\newcommand{\MlstLpQgRthirtyCoarseOverFullPr}{equal within 6\%}  % experiments/FIGS/data/v2/index.json :: linear probe qg: R30 of R16_Q1 over the 188-class run :: sha256 c6c4138ee5c86423 :: prose twin: two significant figures, direction in words
\newcommand{\MlstLpQgRthirtyCoarseOverFullZero}{0.95\,\ensuremath{\pm}\,0.09}  % experiments/FIGS/data/v2/index.json :: linear probe qg: R30 of R16_Q1 over the 188-class run :: sha256 c6c4138ee5c86423
\newcommand{\MlstLpQgRthirtyCoarseOverFullZeroPr}{equal within 9\%}  % experiments/FIGS/data/v2/index.json :: linear probe qg: R30 of R16_Q1 over the 188-class run :: sha256 c6c4138ee5c86423 :: prose twin: two significant figures, direction in words
\newcommand{\MlstLpQgRthirtyMeanOverZeroFull}{1.12\,\ensuremath{\pm}\,0.17}  % experiments/FIGS/data/v2/index.json :: linear probe qg: R30 with mean over R30 with zero, L188 :: sha256 c6c4138ee5c86423
\newcommand{\MlstLpQgRthirtyMeanOverZeroFullPr}{equal within 20\%}  % experiments/FIGS/data/v2/index.json :: linear probe qg: R30 with mean over R30 with zero, L188 :: sha256 c6c4138ee5c86423 :: prose twin: two significant figures, direction in words
\newcommand{\MlstLpQgRthirtyDrawOverZeroFull}{1.05\,\ensuremath{\pm}\,0.12}  % experiments/FIGS/data/v2/index.json :: linear probe qg: R30 with draws over R30 with zero, L188 :: sha256 c6c4138ee5c86423
\newcommand{\MlstLpQgRthirtyDrawOverZeroFullPr}{equal within 20\%}  % experiments/FIGS/data/v2/index.json :: linear probe qg: R30 with draws over R30 with zero, L188 :: sha256 c6c4138ee5c86423 :: prose twin: two significant figures, direction in words
\newcommand{\MlstLpQgRthirtyMeanOverZeroInitBn}{0.964\,\ensuremath{\pm}\,0.030}  % experiments/FIGS/data/v2/index.json :: linear probe qg: R30 with mean over R30 with zero, INIT_BN :: sha256 c6c4138ee5c86423
\newcommand{\MlstLpQgRthirtyMeanOverZeroInitBnPr}{equal within 4\%}  % experiments/FIGS/data/v2/index.json :: linear probe qg: R30 with mean over R30 with zero, INIT_BN :: sha256 c6c4138ee5c86423 :: prose twin: two significant figures, direction in words
\newcommand{\MlstLpQgRthirtyDrawOverZeroInitBn}{0.97\,\ensuremath{\pm}\,0.04}  % experiments/FIGS/data/v2/index.json :: linear probe qg: R30 with draws over R30 with zero, INIT_BN :: sha256 c6c4138ee5c86423
\newcommand{\MlstLpQgRthirtyDrawOverZeroInitBnPr}{equal within 4\%}  % experiments/FIGS/data/v2/index.json :: linear probe qg: R30 with draws over R30 with zero, INIT_BN :: sha256 c6c4138ee5c86423 :: prose twin: two significant figures, direction in words
\newcommand{\MlstLpQgRthirtyFullZero}{66\,\ensuremath{\pm}\,7}  % experiments/FIGS/data/v2/index.json :: linear probe qg, largest set: R30 with zero, L188 :: sha256 c6c4138ee5c86423
\newcommand{\MlstLpQgRthirtyFullMean}{72.7\,\ensuremath{\pm}\,2.4}  % experiments/FIGS/data/v2/index.json :: linear probe qg, largest set: R30 with mean, L188 :: sha256 c6c4138ee5c86423
\newcommand{\MlstLpQgRthirtyFullDraws}{68.1\,\ensuremath{\pm}\,1.8}  % experiments/FIGS/data/v2/index.json :: linear probe qg, largest set: R30 with draws, L188 :: sha256 c6c4138ee5c86423
\newcommand{\MlstLpQgRthirtyInitBnZero}{54.2\,\ensuremath{\pm}\,2.4}  % experiments/FIGS/data/v2/index.json :: linear probe qg, largest set: R30 with zero, INIT_BN :: sha256 c6c4138ee5c86423
\newcommand{\MlstLpQgRthirtyInitBnMean}{52.2\,\ensuremath{\pm}\,0.8}  % experiments/FIGS/data/v2/index.json :: linear probe qg, largest set: R30 with mean, INIT_BN :: sha256 c6c4138ee5c86423
\newcommand{\MlstLpQgRthirtyInitBnDraws}{52.4\,\ensuremath{\pm}\,0.4}  % experiments/FIGS/data/v2/index.json :: linear probe qg, largest set: R30 with draws, INIT_BN :: sha256 c6c4138ee5c86423
\newcommand{\MlstLpQgHerwigRthirtyFullOverInitBn}{1.14\,\ensuremath{\pm}\,0.10}  % experiments/FIGS/data/v2/index.json :: linear probe qg_herwig, largest set, chosen substitution: R30 of L188 over the BatchNorm-recomputed random network of the same index :: sha256 c6c4138ee5c86423
\newcommand{\MlstLpQgHerwigRthirtyFullOverInitBnPr}{equal within 20\%}  % experiments/FIGS/data/v2/index.json :: linear probe qg_herwig, largest set, chosen substitution: R30 of L188 over the BatchNorm-recomputed random network of the same index :: sha256 c6c4138ee5c86423 :: prose twin: two significant figures, direction in words
\newcommand{\MlstLpQgHerwigRthirtyFullOverInitBnZero}{1.054\,\ensuremath{\pm}\,0.028}  % experiments/FIGS/data/v2/index.json :: linear probe qg_herwig, largest set, zero: R30 of L188 over the BatchNorm-recomputed random network of the same index :: sha256 c6c4138ee5c86423
\newcommand{\MlstLpQgHerwigRthirtyFullOverInitBnZeroPr}{equal within 6\%}  % experiments/FIGS/data/v2/index.json :: linear probe qg_herwig, largest set, zero: R30 of L188 over the BatchNorm-recomputed random network of the same index :: sha256 c6c4138ee5c86423 :: prose twin: two significant figures, direction in words
\newcommand{\MlstLpQgHerwigRthirtyResOverInitBn}{1.17\,\ensuremath{\pm}\,0.04}  % experiments/FIGS/data/v2/index.json :: linear probe qg_herwig, largest set, chosen substitution: R30 of L162 over the BatchNorm-recomputed random network of the same index :: sha256 c6c4138ee5c86423
\newcommand{\MlstLpQgHerwigRthirtyResOverInitBnPr}{17\% more rejection}  % experiments/FIGS/data/v2/index.json :: linear probe qg_herwig, largest set, chosen substitution: R30 of L162 over the BatchNorm-recomputed random network of the same index :: sha256 c6c4138ee5c86423 :: prose twin: two significant figures, direction in words
\newcommand{\MlstLpQgHerwigRthirtyResOverInitBnZero}{0.98\,\ensuremath{\pm}\,0.10}  % experiments/FIGS/data/v2/index.json :: linear probe qg_herwig, largest set, zero: R30 of L162 over the BatchNorm-recomputed random network of the same index :: sha256 c6c4138ee5c86423
\newcommand{\MlstLpQgHerwigRthirtyResOverInitBnZeroPr}{equal within 10\%}  % experiments/FIGS/data/v2/index.json :: linear probe qg_herwig, largest set, zero: R30 of L162 over the BatchNorm-recomputed random network of the same index :: sha256 c6c4138ee5c86423 :: prose twin: two significant figures, direction in words
\newcommand{\MlstLpQgHerwigRthirtyMidOverInitBn}{1.12\,\ensuremath{\pm}\,0.07}  % experiments/FIGS/data/v2/index.json :: linear probe qg_herwig, largest set, chosen substitution: R30 of R42_Q1 over the BatchNorm-recomputed random network of the same index :: sha256 c6c4138ee5c86423
\newcommand{\MlstLpQgHerwigRthirtyMidOverInitBnPr}{equal within 20\%}  % experiments/FIGS/data/v2/index.json :: linear probe qg_herwig, largest set, chosen substitution: R30 of R42_Q1 over the BatchNorm-recomputed random network of the same index :: sha256 c6c4138ee5c86423 :: prose twin: two significant figures, direction in words
\newcommand{\MlstLpQgHerwigRthirtyMidOverInitBnZero}{0.99\,\ensuremath{\pm}\,0.07}  % experiments/FIGS/data/v2/index.json :: linear probe qg_herwig, largest set, zero: R30 of R42_Q1 over the BatchNorm-recomputed random network of the same index :: sha256 c6c4138ee5c86423
\newcommand{\MlstLpQgHerwigRthirtyMidOverInitBnZeroPr}{equal within 7\%}  % experiments/FIGS/data/v2/index.json :: linear probe qg_herwig, largest set, zero: R30 of R42_Q1 over the BatchNorm-recomputed random network of the same index :: sha256 c6c4138ee5c86423 :: prose twin: two significant figures, direction in words
\newcommand{\MlstLpQgHerwigRthirtyCoarseOverInitBn}{1.154\,\ensuremath{\pm}\,0.015}  % experiments/FIGS/data/v2/index.json :: linear probe qg_herwig, largest set, chosen substitution: R30 of R16_Q1 over the BatchNorm-recomputed random network of the same index :: sha256 c6c4138ee5c86423
\newcommand{\MlstLpQgHerwigRthirtyCoarseOverInitBnPr}{15\% more rejection}  % experiments/FIGS/data/v2/index.json :: linear probe qg_herwig, largest set, chosen substitution: R30 of R16_Q1 over the BatchNorm-recomputed random network of the same index :: sha256 c6c4138ee5c86423 :: prose twin: two significant figures, direction in words
\newcommand{\MlstLpQgHerwigRthirtyCoarseOverInitBnZero}{1.07\,\ensuremath{\pm}\,0.05}  % experiments/FIGS/data/v2/index.json :: linear probe qg_herwig, largest set, zero: R30 of R16_Q1 over the BatchNorm-recomputed random network of the same index :: sha256 c6c4138ee5c86423
\newcommand{\MlstLpQgHerwigRthirtyCoarseOverInitBnZeroPr}{equal within 8\%}  % experiments/FIGS/data/v2/index.json :: linear probe qg_herwig, largest set, zero: R30 of R16_Q1 over the BatchNorm-recomputed random network of the same index :: sha256 c6c4138ee5c86423 :: prose twin: two significant figures, direction in words
\newcommand{\MlstLpQgHerwigRthirtyResMassOverInitBn}{1.152 and 1.138}  % experiments/FIGS/data/v2/index.json :: linear probe qg_herwig, largest set, chosen substitution: R30 of L162_MASS over the BatchNorm-recomputed random network of the same index :: sha256 c6c4138ee5c86423
\newcommand{\MlstLpQgHerwigRthirtyResMassOverInitBnPr}{equal within 20\%}  % experiments/FIGS/data/v2/index.json :: linear probe qg_herwig, largest set, chosen substitution: R30 of L162_MASS over the BatchNorm-recomputed random network of the same index :: sha256 c6c4138ee5c86423 :: prose twin: two significant figures, direction in words
\newcommand{\MlstLpQgHerwigRthirtyResMassOverInitBnPrMax}{equal within 20\%}  % experiments/FIGS/data/v2/index.json :: linear probe qg_herwig, largest set, chosen substitution: R30 of L162_MASS over the BatchNorm-recomputed random network of the same index :: sha256 c6c4138ee5c86423 :: prose twin of the run farther from 1
\newcommand{\MlstLpQgHerwigRthirtyResMassOverInitBnZero}{1.13 and 1.03}  % experiments/FIGS/data/v2/index.json :: linear probe qg_herwig, largest set, zero: R30 of L162_MASS over the BatchNorm-recomputed random network of the same index :: sha256 c6c4138ee5c86423
\newcommand{\MlstLpQgHerwigRthirtyResMassOverInitBnZeroPr}{equal within 20\%}  % experiments/FIGS/data/v2/index.json :: linear probe qg_herwig, largest set, zero: R30 of L162_MASS over the BatchNorm-recomputed random network of the same index :: sha256 c6c4138ee5c86423 :: prose twin: two significant figures, direction in words
\newcommand{\MlstLpQgHerwigRthirtyResMassOverInitBnZeroPrMax}{equal within 20\%}  % experiments/FIGS/data/v2/index.json :: linear probe qg_herwig, largest set, zero: R30 of L162_MASS over the BatchNorm-recomputed random network of the same index :: sha256 c6c4138ee5c86423 :: prose twin of the run farther from 1
\newcommand{\MlstLpQgHerwigRthirtyCoarseMassOverInitBn}{1.146 and 1.159}  % experiments/FIGS/data/v2/index.json :: linear probe qg_herwig, largest set, chosen substitution: R30 of R16_Q1_MASS over the BatchNorm-recomputed random network of the same index :: sha256 c6c4138ee5c86423
\newcommand{\MlstLpQgHerwigRthirtyCoarseMassOverInitBnPr}{equal within 20\%}  % experiments/FIGS/data/v2/index.json :: linear probe qg_herwig, largest set, chosen substitution: R30 of R16_Q1_MASS over the BatchNorm-recomputed random network of the same index :: sha256 c6c4138ee5c86423 :: prose twin: two significant figures, direction in words
\newcommand{\MlstLpQgHerwigRthirtyCoarseMassOverInitBnPrMax}{equal within 20\%}  % experiments/FIGS/data/v2/index.json :: linear probe qg_herwig, largest set, chosen substitution: R30 of R16_Q1_MASS over the BatchNorm-recomputed random network of the same index :: sha256 c6c4138ee5c86423 :: prose twin of the run farther from 1
\newcommand{\MlstLpQgHerwigRthirtyCoarseMassOverInitBnZero}{1.070 and 1.100}  % experiments/FIGS/data/v2/index.json :: linear probe qg_herwig, largest set, zero: R30 of R16_Q1_MASS over the BatchNorm-recomputed random network of the same index :: sha256 c6c4138ee5c86423
\newcommand{\MlstLpQgHerwigRthirtyCoarseMassOverInitBnZeroPr}{equal within 20\%}  % experiments/FIGS/data/v2/index.json :: linear probe qg_herwig, largest set, zero: R30 of R16_Q1_MASS over the BatchNorm-recomputed random network of the same index :: sha256 c6c4138ee5c86423 :: prose twin: two significant figures, direction in words
\newcommand{\MlstLpQgHerwigRthirtyCoarseMassOverInitBnZeroPrMax}{equal within 20\%}  % experiments/FIGS/data/v2/index.json :: linear probe qg_herwig, largest set, zero: R30 of R16_Q1_MASS over the BatchNorm-recomputed random network of the same index :: sha256 c6c4138ee5c86423 :: prose twin of the run farther from 1
\newcommand{\MlstLpQgHerwigRthirtyResOverFull}{1.03\,\ensuremath{\pm}\,0.11}  % experiments/FIGS/data/v2/index.json :: linear probe qg_herwig: R30 of L162 over the 188-class run :: sha256 c6c4138ee5c86423
\newcommand{\MlstLpQgHerwigRthirtyResOverFullPr}{equal within 20\%}  % experiments/FIGS/data/v2/index.json :: linear probe qg_herwig: R30 of L162 over the 188-class run :: sha256 c6c4138ee5c86423 :: prose twin: two significant figures, direction in words
\newcommand{\MlstLpQgHerwigRthirtyResOverFullZero}{0.93\,\ensuremath{\pm}\,0.11}  % experiments/FIGS/data/v2/index.json :: linear probe qg_herwig: R30 of L162 over the 188-class run :: sha256 c6c4138ee5c86423
\newcommand{\MlstLpQgHerwigRthirtyResOverFullZeroPr}{equal within 20\%}  % experiments/FIGS/data/v2/index.json :: linear probe qg_herwig: R30 of L162 over the 188-class run :: sha256 c6c4138ee5c86423 :: prose twin: two significant figures, direction in words
\newcommand{\MlstLpQgHerwigRthirtyMidOverFull}{0.98\,\ensuremath{\pm}\,0.06}  % experiments/FIGS/data/v2/index.json :: linear probe qg_herwig: R30 of R42_Q1 over the 188-class run :: sha256 c6c4138ee5c86423
\newcommand{\MlstLpQgHerwigRthirtyMidOverFullPr}{equal within 6\%}  % experiments/FIGS/data/v2/index.json :: linear probe qg_herwig: R30 of R42_Q1 over the 188-class run :: sha256 c6c4138ee5c86423 :: prose twin: two significant figures, direction in words
\newcommand{\MlstLpQgHerwigRthirtyMidOverFullZero}{0.94\,\ensuremath{\pm}\,0.09}  % experiments/FIGS/data/v2/index.json :: linear probe qg_herwig: R30 of R42_Q1 over the 188-class run :: sha256 c6c4138ee5c86423
\newcommand{\MlstLpQgHerwigRthirtyMidOverFullZeroPr}{equal within 9\%}  % experiments/FIGS/data/v2/index.json :: linear probe qg_herwig: R30 of R42_Q1 over the 188-class run :: sha256 c6c4138ee5c86423 :: prose twin: two significant figures, direction in words
\newcommand{\MlstLpQgHerwigRthirtyCoarseOverFull}{1.02\,\ensuremath{\pm}\,0.09}  % experiments/FIGS/data/v2/index.json :: linear probe qg_herwig: R30 of R16_Q1 over the 188-class run :: sha256 c6c4138ee5c86423
\newcommand{\MlstLpQgHerwigRthirtyCoarseOverFullPr}{equal within 10\%}  % experiments/FIGS/data/v2/index.json :: linear probe qg_herwig: R30 of R16_Q1 over the 188-class run :: sha256 c6c4138ee5c86423 :: prose twin: two significant figures, direction in words
\newcommand{\MlstLpQgHerwigRthirtyCoarseOverFullZero}{1.018\,\ensuremath{\pm}\,0.024}  % experiments/FIGS/data/v2/index.json :: linear probe qg_herwig: R30 of R16_Q1 over the 188-class run :: sha256 c6c4138ee5c86423
\newcommand{\MlstLpQgHerwigRthirtyCoarseOverFullZeroPr}{equal within 3\%}  % experiments/FIGS/data/v2/index.json :: linear probe qg_herwig: R30 of R16_Q1 over the 188-class run :: sha256 c6c4138ee5c86423 :: prose twin: two significant figures, direction in words
\newcommand{\MlstLpQgHerwigRthirtyMeanOverZeroFull}{1.07\,\ensuremath{\pm}\,0.11}  % experiments/FIGS/data/v2/index.json :: linear probe qg_herwig: R30 with mean over R30 with zero, L188 :: sha256 c6c4138ee5c86423
\newcommand{\MlstLpQgHerwigRthirtyMeanOverZeroFullPr}{equal within 20\%}  % experiments/FIGS/data/v2/index.json :: linear probe qg_herwig: R30 with mean over R30 with zero, L188 :: sha256 c6c4138ee5c86423 :: prose twin: two significant figures, direction in words
\newcommand{\MlstLpQgHerwigRthirtyDrawOverZeroFull}{1.07\,\ensuremath{\pm}\,0.06}  % experiments/FIGS/data/v2/index.json :: linear probe qg_herwig: R30 with draws over R30 with zero, L188 :: sha256 c6c4138ee5c86423
\newcommand{\MlstLpQgHerwigRthirtyDrawOverZeroFullPr}{equal within 7\%}  % experiments/FIGS/data/v2/index.json :: linear probe qg_herwig: R30 with draws over R30 with zero, L188 :: sha256 c6c4138ee5c86423 :: prose twin: two significant figures, direction in words
\newcommand{\MlstLpQgHerwigRthirtyMeanOverZeroInitBn}{0.988\,\ensuremath{\pm}\,0.019}  % experiments/FIGS/data/v2/index.json :: linear probe qg_herwig: R30 with mean over R30 with zero, INIT_BN :: sha256 c6c4138ee5c86423
\newcommand{\MlstLpQgHerwigRthirtyMeanOverZeroInitBnPr}{equal within 2\%}  % experiments/FIGS/data/v2/index.json :: linear probe qg_herwig: R30 with mean over R30 with zero, INIT_BN :: sha256 c6c4138ee5c86423 :: prose twin: two significant figures, direction in words
\newcommand{\MlstLpQgHerwigRthirtyDrawOverZeroInitBn}{0.983\,\ensuremath{\pm}\,0.032}  % experiments/FIGS/data/v2/index.json :: linear probe qg_herwig: R30 with draws over R30 with zero, INIT_BN :: sha256 c6c4138ee5c86423
\newcommand{\MlstLpQgHerwigRthirtyDrawOverZeroInitBnPr}{equal within 4\%}  % experiments/FIGS/data/v2/index.json :: linear probe qg_herwig: R30 with draws over R30 with zero, INIT_BN :: sha256 c6c4138ee5c86423 :: prose twin: two significant figures, direction in words
\newcommand{\MlstLpQgHerwigRthirtyFullZero}{32.8\,\ensuremath{\pm}\,1.0}  % experiments/FIGS/data/v2/index.json :: linear probe qg_herwig, largest set: R30 with zero, L188 :: sha256 c6c4138ee5c86423
\newcommand{\MlstLpQgHerwigRthirtyFullMean}{35.0\,\ensuremath{\pm}\,2.6}  % experiments/FIGS/data/v2/index.json :: linear probe qg_herwig, largest set: R30 with mean, L188 :: sha256 c6c4138ee5c86423
\newcommand{\MlstLpQgHerwigRthirtyFullDraws}{35.0\,\ensuremath{\pm}\,0.8}  % experiments/FIGS/data/v2/index.json :: linear probe qg_herwig, largest set: R30 with draws, L188 :: sha256 c6c4138ee5c86423
\newcommand{\MlstLpQgHerwigRthirtyInitBnZero}{31.1\,\ensuremath{\pm}\,0.5}  % experiments/FIGS/data/v2/index.json :: linear probe qg_herwig, largest set: R30 with zero, INIT_BN :: sha256 c6c4138ee5c86423
\newcommand{\MlstLpQgHerwigRthirtyInitBnMean}{30.7\,\ensuremath{\pm}\,0.5}  % experiments/FIGS/data/v2/index.json :: linear probe qg_herwig, largest set: R30 with mean, INIT_BN :: sha256 c6c4138ee5c86423
\newcommand{\MlstLpQgHerwigRthirtyInitBnDraws}{30.6\,\ensuremath{\pm}\,0.6}  % experiments/FIGS/data/v2/index.json :: linear probe qg_herwig, largest set: R30 with draws, INIT_BN :: sha256 c6c4138ee5c86423
\newcommand{\MlstLpJcTwoTopZeroAccDropFull}{88.9\,\ensuremath{\pm}\,1.0\%}  % experiments/FIGS/data/v2/index.json :: linear probe JetClass-II, largest set: relative drop in accuracy (%) with the top data's missing slots set to zero, L188 :: sha256 c6c4138ee5c86423
\newcommand{\MlstLpJcTwoTopZeroAccDropFullPr}{89\%}  % experiments/FIGS/data/v2/index.json :: linear probe JetClass-II, largest set: relative drop in accuracy (%) with the top data's missing slots set to zero, L188 :: sha256 c6c4138ee5c86423 :: prose twin, two significant figures
\newcommand{\MlstLpJcTwoTopMeanAccDropFull}{87.9\,\ensuremath{\pm}\,1.1\%}  % experiments/FIGS/data/v2/index.json :: linear probe JetClass-II, largest set: relative drop in accuracy (%) with the top data's missing slots set to mean, L188 :: sha256 c6c4138ee5c86423
\newcommand{\MlstLpJcTwoTopMeanAccDropFullPr}{88\%}  % experiments/FIGS/data/v2/index.json :: linear probe JetClass-II, largest set: relative drop in accuracy (%) with the top data's missing slots set to mean, L188 :: sha256 c6c4138ee5c86423 :: prose twin, two significant figures
\newcommand{\MlstLpJcTwoTopDrawAccDropFull}{85.98\,\ensuremath{\pm}\,0.07\%}  % experiments/FIGS/data/v2/index.json :: linear probe JetClass-II, largest set: relative drop in accuracy (%) with the top data's missing slots set to draws, L188 :: sha256 c6c4138ee5c86423
\newcommand{\MlstLpJcTwoTopDrawAccDropFullPr}{86\%}  % experiments/FIGS/data/v2/index.json :: linear probe JetClass-II, largest set: relative drop in accuracy (%) with the top data's missing slots set to draws, L188 :: sha256 c6c4138ee5c86423 :: prose twin, two significant figures
\newcommand{\MlstLpJcTwoTopZeroAccDropRes}{89.0\,\ensuremath{\pm}\,0.9\%}  % experiments/FIGS/data/v2/index.json :: linear probe JetClass-II, largest set: relative drop in accuracy (%) with the top data's missing slots set to zero, L162 :: sha256 c6c4138ee5c86423
\newcommand{\MlstLpJcTwoTopZeroAccDropResPr}{89\%}  % experiments/FIGS/data/v2/index.json :: linear probe JetClass-II, largest set: relative drop in accuracy (%) with the top data's missing slots set to zero, L162 :: sha256 c6c4138ee5c86423 :: prose twin, two significant figures
\newcommand{\MlstLpJcTwoTopMeanAccDropRes}{87.8\,\ensuremath{\pm}\,0.7\%}  % experiments/FIGS/data/v2/index.json :: linear probe JetClass-II, largest set: relative drop in accuracy (%) with the top data's missing slots set to mean, L162 :: sha256 c6c4138ee5c86423
\newcommand{\MlstLpJcTwoTopMeanAccDropResPr}{88\%}  % experiments/FIGS/data/v2/index.json :: linear probe JetClass-II, largest set: relative drop in accuracy (%) with the top data's missing slots set to mean, L162 :: sha256 c6c4138ee5c86423 :: prose twin, two significant figures
\newcommand{\MlstLpJcTwoTopDrawAccDropRes}{85.89\,\ensuremath{\pm}\,0.11\%}  % experiments/FIGS/data/v2/index.json :: linear probe JetClass-II, largest set: relative drop in accuracy (%) with the top data's missing slots set to draws, L162 :: sha256 c6c4138ee5c86423
\newcommand{\MlstLpJcTwoTopDrawAccDropResPr}{86\%}  % experiments/FIGS/data/v2/index.json :: linear probe JetClass-II, largest set: relative drop in accuracy (%) with the top data's missing slots set to draws, L162 :: sha256 c6c4138ee5c86423 :: prose twin, two significant figures
\newcommand{\MlstLpJcTwoTopZeroAccDropMid}{87.81\,\ensuremath{\pm}\,0.28\%}  % experiments/FIGS/data/v2/index.json :: linear probe JetClass-II, largest set: relative drop in accuracy (%) with the top data's missing slots set to zero, R42_Q1 :: sha256 c6c4138ee5c86423
\newcommand{\MlstLpJcTwoTopZeroAccDropMidPr}{88\%}  % experiments/FIGS/data/v2/index.json :: linear probe JetClass-II, largest set: relative drop in accuracy (%) with the top data's missing slots set to zero, R42_Q1 :: sha256 c6c4138ee5c86423 :: prose twin, two significant figures
\newcommand{\MlstLpJcTwoTopMeanAccDropMid}{86.7\,\ensuremath{\pm}\,0.9\%}  % experiments/FIGS/data/v2/index.json :: linear probe JetClass-II, largest set: relative drop in accuracy (%) with the top data's missing slots set to mean, R42_Q1 :: sha256 c6c4138ee5c86423
\newcommand{\MlstLpJcTwoTopMeanAccDropMidPr}{87\%}  % experiments/FIGS/data/v2/index.json :: linear probe JetClass-II, largest set: relative drop in accuracy (%) with the top data's missing slots set to mean, R42_Q1 :: sha256 c6c4138ee5c86423 :: prose twin, two significant figures
\newcommand{\MlstLpJcTwoTopDrawAccDropMid}{83.17\,\ensuremath{\pm}\,0.29\%}  % experiments/FIGS/data/v2/index.json :: linear probe JetClass-II, largest set: relative drop in accuracy (%) with the top data's missing slots set to draws, R42_Q1 :: sha256 c6c4138ee5c86423
\newcommand{\MlstLpJcTwoTopDrawAccDropMidPr}{83\%}  % experiments/FIGS/data/v2/index.json :: linear probe JetClass-II, largest set: relative drop in accuracy (%) with the top data's missing slots set to draws, R42_Q1 :: sha256 c6c4138ee5c86423 :: prose twin, two significant figures
\newcommand{\MlstLpJcTwoTopZeroAccDropCoarse}{86.7\,\ensuremath{\pm}\,0.7\%}  % experiments/FIGS/data/v2/index.json :: linear probe JetClass-II, largest set: relative drop in accuracy (%) with the top data's missing slots set to zero, R16_Q1 :: sha256 c6c4138ee5c86423
\newcommand{\MlstLpJcTwoTopZeroAccDropCoarsePr}{87\%}  % experiments/FIGS/data/v2/index.json :: linear probe JetClass-II, largest set: relative drop in accuracy (%) with the top data's missing slots set to zero, R16_Q1 :: sha256 c6c4138ee5c86423 :: prose twin, two significant figures
\newcommand{\MlstLpJcTwoTopMeanAccDropCoarse}{86.1\,\ensuremath{\pm}\,0.8\%}  % experiments/FIGS/data/v2/index.json :: linear probe JetClass-II, largest set: relative drop in accuracy (%) with the top data's missing slots set to mean, R16_Q1 :: sha256 c6c4138ee5c86423
\newcommand{\MlstLpJcTwoTopMeanAccDropCoarsePr}{86\%}  % experiments/FIGS/data/v2/index.json :: linear probe JetClass-II, largest set: relative drop in accuracy (%) with the top data's missing slots set to mean, R16_Q1 :: sha256 c6c4138ee5c86423 :: prose twin, two significant figures
\newcommand{\MlstLpJcTwoTopDrawAccDropCoarse}{81.4\,\ensuremath{\pm}\,0.5\%}  % experiments/FIGS/data/v2/index.json :: linear probe JetClass-II, largest set: relative drop in accuracy (%) with the top data's missing slots set to draws, R16_Q1 :: sha256 c6c4138ee5c86423
\newcommand{\MlstLpJcTwoTopDrawAccDropCoarsePr}{81\%}  % experiments/FIGS/data/v2/index.json :: linear probe JetClass-II, largest set: relative drop in accuracy (%) with the top data's missing slots set to draws, R16_Q1 :: sha256 c6c4138ee5c86423 :: prose twin, two significant figures
\newcommand{\MlstLpJcTwoTopZeroAccDropResMass}{89.2, 88.3\%}  % experiments/FIGS/data/v2/index.json :: linear probe JetClass-II, largest set: relative drop in accuracy (%) with the top data's missing slots set to zero, L162_MASS :: sha256 c6c4138ee5c86423
\newcommand{\MlstLpJcTwoTopZeroAccDropResMassPr}{88--89\%}  % experiments/FIGS/data/v2/index.json :: linear probe JetClass-II, largest set: relative drop in accuracy (%) with the top data's missing slots set to zero, L162_MASS :: sha256 c6c4138ee5c86423 :: prose twin, two significant figures
\newcommand{\MlstLpJcTwoTopMeanAccDropResMass}{88.5, 87.8\%}  % experiments/FIGS/data/v2/index.json :: linear probe JetClass-II, largest set: relative drop in accuracy (%) with the top data's missing slots set to mean, L162_MASS :: sha256 c6c4138ee5c86423
\newcommand{\MlstLpJcTwoTopMeanAccDropResMassPr}{88--88\%}  % experiments/FIGS/data/v2/index.json :: linear probe JetClass-II, largest set: relative drop in accuracy (%) with the top data's missing slots set to mean, L162_MASS :: sha256 c6c4138ee5c86423 :: prose twin, two significant figures
\newcommand{\MlstLpJcTwoTopDrawAccDropResMass}{85.4, 85.4\%}  % experiments/FIGS/data/v2/index.json :: linear probe JetClass-II, largest set: relative drop in accuracy (%) with the top data's missing slots set to draws, L162_MASS :: sha256 c6c4138ee5c86423
\newcommand{\MlstLpJcTwoTopDrawAccDropResMassPr}{85--85\%}  % experiments/FIGS/data/v2/index.json :: linear probe JetClass-II, largest set: relative drop in accuracy (%) with the top data's missing slots set to draws, L162_MASS :: sha256 c6c4138ee5c86423 :: prose twin, two significant figures
\newcommand{\MlstLpJcTwoTopZeroAccDropCoarseMass}{84.8, 87.4\%}  % experiments/FIGS/data/v2/index.json :: linear probe JetClass-II, largest set: relative drop in accuracy (%) with the top data's missing slots set to zero, R16_Q1_MASS :: sha256 c6c4138ee5c86423
\newcommand{\MlstLpJcTwoTopZeroAccDropCoarseMassPr}{85--87\%}  % experiments/FIGS/data/v2/index.json :: linear probe JetClass-II, largest set: relative drop in accuracy (%) with the top data's missing slots set to zero, R16_Q1_MASS :: sha256 c6c4138ee5c86423 :: prose twin, two significant figures
\newcommand{\MlstLpJcTwoTopMeanAccDropCoarseMass}{85.9, 86.9\%}  % experiments/FIGS/data/v2/index.json :: linear probe JetClass-II, largest set: relative drop in accuracy (%) with the top data's missing slots set to mean, R16_Q1_MASS :: sha256 c6c4138ee5c86423
\newcommand{\MlstLpJcTwoTopMeanAccDropCoarseMassPr}{86--87\%}  % experiments/FIGS/data/v2/index.json :: linear probe JetClass-II, largest set: relative drop in accuracy (%) with the top data's missing slots set to mean, R16_Q1_MASS :: sha256 c6c4138ee5c86423 :: prose twin, two significant figures
\newcommand{\MlstLpJcTwoTopDrawAccDropCoarseMass}{79.1, 81.7\%}  % experiments/FIGS/data/v2/index.json :: linear probe JetClass-II, largest set: relative drop in accuracy (%) with the top data's missing slots set to draws, R16_Q1_MASS :: sha256 c6c4138ee5c86423
\newcommand{\MlstLpJcTwoTopDrawAccDropCoarseMassPr}{79--82\%}  % experiments/FIGS/data/v2/index.json :: linear probe JetClass-II, largest set: relative drop in accuracy (%) with the top data's missing slots set to draws, R16_Q1_MASS :: sha256 c6c4138ee5c86423 :: prose twin, two significant figures
\newcommand{\MlstLpJcTwoTopZeroAccDropFzero}{88.9, 88.5\%}  % experiments/FIGS/data/v2/index.json :: linear probe JetClass-II, largest set: relative drop in accuracy (%) with the top data's missing slots set to zero, FLAV_F0 :: sha256 c6c4138ee5c86423
\newcommand{\MlstLpJcTwoTopZeroAccDropFzeroPr}{88--89\%}  % experiments/FIGS/data/v2/index.json :: linear probe JetClass-II, largest set: relative drop in accuracy (%) with the top data's missing slots set to zero, FLAV_F0 :: sha256 c6c4138ee5c86423 :: prose twin, two significant figures
\newcommand{\MlstLpJcTwoTopMeanAccDropFzero}{88.3, 87.0\%}  % experiments/FIGS/data/v2/index.json :: linear probe JetClass-II, largest set: relative drop in accuracy (%) with the top data's missing slots set to mean, FLAV_F0 :: sha256 c6c4138ee5c86423
\newcommand{\MlstLpJcTwoTopMeanAccDropFzeroPr}{87--88\%}  % experiments/FIGS/data/v2/index.json :: linear probe JetClass-II, largest set: relative drop in accuracy (%) with the top data's missing slots set to mean, FLAV_F0 :: sha256 c6c4138ee5c86423 :: prose twin, two significant figures
\newcommand{\MlstLpJcTwoTopDrawAccDropFzero}{83.3, 83.9\%}  % experiments/FIGS/data/v2/index.json :: linear probe JetClass-II, largest set: relative drop in accuracy (%) with the top data's missing slots set to draws, FLAV_F0 :: sha256 c6c4138ee5c86423
\newcommand{\MlstLpJcTwoTopDrawAccDropFzeroPr}{83--84\%}  % experiments/FIGS/data/v2/index.json :: linear probe JetClass-II, largest set: relative drop in accuracy (%) with the top data's missing slots set to draws, FLAV_F0 :: sha256 c6c4138ee5c86423 :: prose twin, two significant figures
\newcommand{\MlstLpJcTwoTopZeroAccDropFone}{87.8, 87.6\%}  % experiments/FIGS/data/v2/index.json :: linear probe JetClass-II, largest set: relative drop in accuracy (%) with the top data's missing slots set to zero, FLAV_F1 :: sha256 c6c4138ee5c86423
\newcommand{\MlstLpJcTwoTopZeroAccDropFonePr}{88--88\%}  % experiments/FIGS/data/v2/index.json :: linear probe JetClass-II, largest set: relative drop in accuracy (%) with the top data's missing slots set to zero, FLAV_F1 :: sha256 c6c4138ee5c86423 :: prose twin, two significant figures
\newcommand{\MlstLpJcTwoTopMeanAccDropFone}{88.9, 86.9\%}  % experiments/FIGS/data/v2/index.json :: linear probe JetClass-II, largest set: relative drop in accuracy (%) with the top data's missing slots set to mean, FLAV_F1 :: sha256 c6c4138ee5c86423
\newcommand{\MlstLpJcTwoTopMeanAccDropFonePr}{87--89\%}  % experiments/FIGS/data/v2/index.json :: linear probe JetClass-II, largest set: relative drop in accuracy (%) with the top data's missing slots set to mean, FLAV_F1 :: sha256 c6c4138ee5c86423 :: prose twin, two significant figures
\newcommand{\MlstLpJcTwoTopDrawAccDropFone}{84.4, 85.2\%}  % experiments/FIGS/data/v2/index.json :: linear probe JetClass-II, largest set: relative drop in accuracy (%) with the top data's missing slots set to draws, FLAV_F1 :: sha256 c6c4138ee5c86423
\newcommand{\MlstLpJcTwoTopDrawAccDropFonePr}{84--85\%}  % experiments/FIGS/data/v2/index.json :: linear probe JetClass-II, largest set: relative drop in accuracy (%) with the top data's missing slots set to draws, FLAV_F1 :: sha256 c6c4138ee5c86423 :: prose twin, two significant figures
\newcommand{\MlstLpJcTwoTopZeroAccDropFoner}{89.6, 90.3\%}  % experiments/FIGS/data/v2/index.json :: linear probe JetClass-II, largest set: relative drop in accuracy (%) with the top data's missing slots set to zero, FLAV_F1R :: sha256 c6c4138ee5c86423
\newcommand{\MlstLpJcTwoTopZeroAccDropFonerPr}{90--90\%}  % experiments/FIGS/data/v2/index.json :: linear probe JetClass-II, largest set: relative drop in accuracy (%) with the top data's missing slots set to zero, FLAV_F1R :: sha256 c6c4138ee5c86423 :: prose twin, two significant figures
\newcommand{\MlstLpJcTwoTopMeanAccDropFoner}{89.4, 89.7\%}  % experiments/FIGS/data/v2/index.json :: linear probe JetClass-II, largest set: relative drop in accuracy (%) with the top data's missing slots set to mean, FLAV_F1R :: sha256 c6c4138ee5c86423
\newcommand{\MlstLpJcTwoTopMeanAccDropFonerPr}{89--90\%}  % experiments/FIGS/data/v2/index.json :: linear probe JetClass-II, largest set: relative drop in accuracy (%) with the top data's missing slots set to mean, FLAV_F1R :: sha256 c6c4138ee5c86423 :: prose twin, two significant figures
\newcommand{\MlstLpJcTwoTopDrawAccDropFoner}{84.9, 85.6\%}  % experiments/FIGS/data/v2/index.json :: linear probe JetClass-II, largest set: relative drop in accuracy (%) with the top data's missing slots set to draws, FLAV_F1R :: sha256 c6c4138ee5c86423
\newcommand{\MlstLpJcTwoTopDrawAccDropFonerPr}{85--86\%}  % experiments/FIGS/data/v2/index.json :: linear probe JetClass-II, largest set: relative drop in accuracy (%) with the top data's missing slots set to draws, FLAV_F1R :: sha256 c6c4138ee5c86423 :: prose twin, two significant figures
\newcommand{\MlstLpJcTwoTopZeroAccDropInitBn}{56.4\,\ensuremath{\pm}\,1.4\%}  % experiments/FIGS/data/v2/index.json :: linear probe JetClass-II, largest set: relative drop in accuracy (%) with the top data's missing slots set to zero, INIT_BN :: sha256 c6c4138ee5c86423
\newcommand{\MlstLpJcTwoTopZeroAccDropInitBnPr}{56\%}  % experiments/FIGS/data/v2/index.json :: linear probe JetClass-II, largest set: relative drop in accuracy (%) with the top data's missing slots set to zero, INIT_BN :: sha256 c6c4138ee5c86423 :: prose twin, two significant figures
\newcommand{\MlstLpJcTwoTopMeanAccDropInitBn}{59.8\,\ensuremath{\pm}\,0.8\%}  % experiments/FIGS/data/v2/index.json :: linear probe JetClass-II, largest set: relative drop in accuracy (%) with the top data's missing slots set to mean, INIT_BN :: sha256 c6c4138ee5c86423
\newcommand{\MlstLpJcTwoTopMeanAccDropInitBnPr}{60\%}  % experiments/FIGS/data/v2/index.json :: linear probe JetClass-II, largest set: relative drop in accuracy (%) with the top data's missing slots set to mean, INIT_BN :: sha256 c6c4138ee5c86423 :: prose twin, two significant figures
\newcommand{\MlstLpJcTwoTopDrawAccDropInitBn}{59.8\,\ensuremath{\pm}\,1.2\%}  % experiments/FIGS/data/v2/index.json :: linear probe JetClass-II, largest set: relative drop in accuracy (%) with the top data's missing slots set to draws, INIT_BN :: sha256 c6c4138ee5c86423
\newcommand{\MlstLpJcTwoTopDrawAccDropInitBnPr}{60\%}  % experiments/FIGS/data/v2/index.json :: linear probe JetClass-II, largest set: relative drop in accuracy (%) with the top data's missing slots set to draws, INIT_BN :: sha256 c6c4138ee5c86423 :: prose twin, two significant figures
\newcommand{\MlstLpJcTwoQgZeroAccDropFull}{61.2\,\ensuremath{\pm}\,3.0\%}  % experiments/FIGS/data/v2/index.json :: linear probe JetClass-II, largest set: relative drop in accuracy (%) with the qg data's missing slots set to zero, L188 :: sha256 c6c4138ee5c86423
\newcommand{\MlstLpJcTwoQgZeroAccDropFullPr}{61\%}  % experiments/FIGS/data/v2/index.json :: linear probe JetClass-II, largest set: relative drop in accuracy (%) with the qg data's missing slots set to zero, L188 :: sha256 c6c4138ee5c86423 :: prose twin, two significant figures
\newcommand{\MlstLpJcTwoQgMeanAccDropFull}{48.8\,\ensuremath{\pm}\,1.1\%}  % experiments/FIGS/data/v2/index.json :: linear probe JetClass-II, largest set: relative drop in accuracy (%) with the qg data's missing slots set to mean, L188 :: sha256 c6c4138ee5c86423
\newcommand{\MlstLpJcTwoQgMeanAccDropFullPr}{49\%}  % experiments/FIGS/data/v2/index.json :: linear probe JetClass-II, largest set: relative drop in accuracy (%) with the qg data's missing slots set to mean, L188 :: sha256 c6c4138ee5c86423 :: prose twin, two significant figures
\newcommand{\MlstLpJcTwoQgDrawAccDropFull}{54.30\,\ensuremath{\pm}\,0.21\%}  % experiments/FIGS/data/v2/index.json :: linear probe JetClass-II, largest set: relative drop in accuracy (%) with the qg data's missing slots set to draws, L188 :: sha256 c6c4138ee5c86423
\newcommand{\MlstLpJcTwoQgDrawAccDropFullPr}{54\%}  % experiments/FIGS/data/v2/index.json :: linear probe JetClass-II, largest set: relative drop in accuracy (%) with the qg data's missing slots set to draws, L188 :: sha256 c6c4138ee5c86423 :: prose twin, two significant figures
\newcommand{\MlstLpJcTwoQgZeroAccDropRes}{64\,\ensuremath{\pm}\,5\%}  % experiments/FIGS/data/v2/index.json :: linear probe JetClass-II, largest set: relative drop in accuracy (%) with the qg data's missing slots set to zero, L162 :: sha256 c6c4138ee5c86423
\newcommand{\MlstLpJcTwoQgZeroAccDropResPr}{64\%}  % experiments/FIGS/data/v2/index.json :: linear probe JetClass-II, largest set: relative drop in accuracy (%) with the qg data's missing slots set to zero, L162 :: sha256 c6c4138ee5c86423 :: prose twin, two significant figures
\newcommand{\MlstLpJcTwoQgMeanAccDropRes}{49.5\,\ensuremath{\pm}\,0.5\%}  % experiments/FIGS/data/v2/index.json :: linear probe JetClass-II, largest set: relative drop in accuracy (%) with the qg data's missing slots set to mean, L162 :: sha256 c6c4138ee5c86423
\newcommand{\MlstLpJcTwoQgMeanAccDropResPr}{49\%}  % experiments/FIGS/data/v2/index.json :: linear probe JetClass-II, largest set: relative drop in accuracy (%) with the qg data's missing slots set to mean, L162 :: sha256 c6c4138ee5c86423 :: prose twin, two significant figures
\newcommand{\MlstLpJcTwoQgDrawAccDropRes}{54.6\,\ensuremath{\pm}\,0.6\%}  % experiments/FIGS/data/v2/index.json :: linear probe JetClass-II, largest set: relative drop in accuracy (%) with the qg data's missing slots set to draws, L162 :: sha256 c6c4138ee5c86423
\newcommand{\MlstLpJcTwoQgDrawAccDropResPr}{55\%}  % experiments/FIGS/data/v2/index.json :: linear probe JetClass-II, largest set: relative drop in accuracy (%) with the qg data's missing slots set to draws, L162 :: sha256 c6c4138ee5c86423 :: prose twin, two significant figures
\newcommand{\MlstLpJcTwoQgZeroAccDropMid}{68\,\ensuremath{\pm}\,5\%}  % experiments/FIGS/data/v2/index.json :: linear probe JetClass-II, largest set: relative drop in accuracy (%) with the qg data's missing slots set to zero, R42_Q1 :: sha256 c6c4138ee5c86423
\newcommand{\MlstLpJcTwoQgZeroAccDropMidPr}{68\%}  % experiments/FIGS/data/v2/index.json :: linear probe JetClass-II, largest set: relative drop in accuracy (%) with the qg data's missing slots set to zero, R42_Q1 :: sha256 c6c4138ee5c86423 :: prose twin, two significant figures
\newcommand{\MlstLpJcTwoQgMeanAccDropMid}{47.6\,\ensuremath{\pm}\,1.6\%}  % experiments/FIGS/data/v2/index.json :: linear probe JetClass-II, largest set: relative drop in accuracy (%) with the qg data's missing slots set to mean, R42_Q1 :: sha256 c6c4138ee5c86423
\newcommand{\MlstLpJcTwoQgMeanAccDropMidPr}{48\%}  % experiments/FIGS/data/v2/index.json :: linear probe JetClass-II, largest set: relative drop in accuracy (%) with the qg data's missing slots set to mean, R42_Q1 :: sha256 c6c4138ee5c86423 :: prose twin, two significant figures
\newcommand{\MlstLpJcTwoQgDrawAccDropMid}{53.5\,\ensuremath{\pm}\,1.4\%}  % experiments/FIGS/data/v2/index.json :: linear probe JetClass-II, largest set: relative drop in accuracy (%) with the qg data's missing slots set to draws, R42_Q1 :: sha256 c6c4138ee5c86423
\newcommand{\MlstLpJcTwoQgDrawAccDropMidPr}{54\%}  % experiments/FIGS/data/v2/index.json :: linear probe JetClass-II, largest set: relative drop in accuracy (%) with the qg data's missing slots set to draws, R42_Q1 :: sha256 c6c4138ee5c86423 :: prose twin, two significant figures
\newcommand{\MlstLpJcTwoQgZeroAccDropCoarse}{70.9\,\ensuremath{\pm}\,2.3\%}  % experiments/FIGS/data/v2/index.json :: linear probe JetClass-II, largest set: relative drop in accuracy (%) with the qg data's missing slots set to zero, R16_Q1 :: sha256 c6c4138ee5c86423
\newcommand{\MlstLpJcTwoQgZeroAccDropCoarsePr}{71\%}  % experiments/FIGS/data/v2/index.json :: linear probe JetClass-II, largest set: relative drop in accuracy (%) with the qg data's missing slots set to zero, R16_Q1 :: sha256 c6c4138ee5c86423 :: prose twin, two significant figures
\newcommand{\MlstLpJcTwoQgMeanAccDropCoarse}{50.0\,\ensuremath{\pm}\,1.0\%}  % experiments/FIGS/data/v2/index.json :: linear probe JetClass-II, largest set: relative drop in accuracy (%) with the qg data's missing slots set to mean, R16_Q1 :: sha256 c6c4138ee5c86423
\newcommand{\MlstLpJcTwoQgMeanAccDropCoarsePr}{50\%}  % experiments/FIGS/data/v2/index.json :: linear probe JetClass-II, largest set: relative drop in accuracy (%) with the qg data's missing slots set to mean, R16_Q1 :: sha256 c6c4138ee5c86423 :: prose twin, two significant figures
\newcommand{\MlstLpJcTwoQgDrawAccDropCoarse}{51.7\,\ensuremath{\pm}\,0.4\%}  % experiments/FIGS/data/v2/index.json :: linear probe JetClass-II, largest set: relative drop in accuracy (%) with the qg data's missing slots set to draws, R16_Q1 :: sha256 c6c4138ee5c86423
\newcommand{\MlstLpJcTwoQgDrawAccDropCoarsePr}{52\%}  % experiments/FIGS/data/v2/index.json :: linear probe JetClass-II, largest set: relative drop in accuracy (%) with the qg data's missing slots set to draws, R16_Q1 :: sha256 c6c4138ee5c86423 :: prose twin, two significant figures
\newcommand{\MlstLpJcTwoQgZeroAccDropResMass}{61.8, 67.9\%}  % experiments/FIGS/data/v2/index.json :: linear probe JetClass-II, largest set: relative drop in accuracy (%) with the qg data's missing slots set to zero, L162_MASS :: sha256 c6c4138ee5c86423
\newcommand{\MlstLpJcTwoQgZeroAccDropResMassPr}{62--68\%}  % experiments/FIGS/data/v2/index.json :: linear probe JetClass-II, largest set: relative drop in accuracy (%) with the qg data's missing slots set to zero, L162_MASS :: sha256 c6c4138ee5c86423 :: prose twin, two significant figures
\newcommand{\MlstLpJcTwoQgMeanAccDropResMass}{49.7, 52.2\%}  % experiments/FIGS/data/v2/index.json :: linear probe JetClass-II, largest set: relative drop in accuracy (%) with the qg data's missing slots set to mean, L162_MASS :: sha256 c6c4138ee5c86423
\newcommand{\MlstLpJcTwoQgMeanAccDropResMassPr}{50--52\%}  % experiments/FIGS/data/v2/index.json :: linear probe JetClass-II, largest set: relative drop in accuracy (%) with the qg data's missing slots set to mean, L162_MASS :: sha256 c6c4138ee5c86423 :: prose twin, two significant figures
\newcommand{\MlstLpJcTwoQgDrawAccDropResMass}{54.3, 55.1\%}  % experiments/FIGS/data/v2/index.json :: linear probe JetClass-II, largest set: relative drop in accuracy (%) with the qg data's missing slots set to draws, L162_MASS :: sha256 c6c4138ee5c86423
\newcommand{\MlstLpJcTwoQgDrawAccDropResMassPr}{54--55\%}  % experiments/FIGS/data/v2/index.json :: linear probe JetClass-II, largest set: relative drop in accuracy (%) with the qg data's missing slots set to draws, L162_MASS :: sha256 c6c4138ee5c86423 :: prose twin, two significant figures
\newcommand{\MlstLpJcTwoQgZeroAccDropCoarseMass}{66.7, 65.4\%}  % experiments/FIGS/data/v2/index.json :: linear probe JetClass-II, largest set: relative drop in accuracy (%) with the qg data's missing slots set to zero, R16_Q1_MASS :: sha256 c6c4138ee5c86423
\newcommand{\MlstLpJcTwoQgZeroAccDropCoarseMassPr}{65--67\%}  % experiments/FIGS/data/v2/index.json :: linear probe JetClass-II, largest set: relative drop in accuracy (%) with the qg data's missing slots set to zero, R16_Q1_MASS :: sha256 c6c4138ee5c86423 :: prose twin, two significant figures
\newcommand{\MlstLpJcTwoQgMeanAccDropCoarseMass}{49.0, 49.6\%}  % experiments/FIGS/data/v2/index.json :: linear probe JetClass-II, largest set: relative drop in accuracy (%) with the qg data's missing slots set to mean, R16_Q1_MASS :: sha256 c6c4138ee5c86423
\newcommand{\MlstLpJcTwoQgMeanAccDropCoarseMassPr}{49--50\%}  % experiments/FIGS/data/v2/index.json :: linear probe JetClass-II, largest set: relative drop in accuracy (%) with the qg data's missing slots set to mean, R16_Q1_MASS :: sha256 c6c4138ee5c86423 :: prose twin, two significant figures
\newcommand{\MlstLpJcTwoQgDrawAccDropCoarseMass}{50.8, 51.1\%}  % experiments/FIGS/data/v2/index.json :: linear probe JetClass-II, largest set: relative drop in accuracy (%) with the qg data's missing slots set to draws, R16_Q1_MASS :: sha256 c6c4138ee5c86423
\newcommand{\MlstLpJcTwoQgDrawAccDropCoarseMassPr}{51--51\%}  % experiments/FIGS/data/v2/index.json :: linear probe JetClass-II, largest set: relative drop in accuracy (%) with the qg data's missing slots set to draws, R16_Q1_MASS :: sha256 c6c4138ee5c86423 :: prose twin, two significant figures
\newcommand{\MlstLpJcTwoQgZeroAccDropFzero}{62.7, 62.2\%}  % experiments/FIGS/data/v2/index.json :: linear probe JetClass-II, largest set: relative drop in accuracy (%) with the qg data's missing slots set to zero, FLAV_F0 :: sha256 c6c4138ee5c86423
\newcommand{\MlstLpJcTwoQgZeroAccDropFzeroPr}{62--63\%}  % experiments/FIGS/data/v2/index.json :: linear probe JetClass-II, largest set: relative drop in accuracy (%) with the qg data's missing slots set to zero, FLAV_F0 :: sha256 c6c4138ee5c86423 :: prose twin, two significant figures
\newcommand{\MlstLpJcTwoQgMeanAccDropFzero}{51.9, 51.3\%}  % experiments/FIGS/data/v2/index.json :: linear probe JetClass-II, largest set: relative drop in accuracy (%) with the qg data's missing slots set to mean, FLAV_F0 :: sha256 c6c4138ee5c86423
\newcommand{\MlstLpJcTwoQgMeanAccDropFzeroPr}{51--52\%}  % experiments/FIGS/data/v2/index.json :: linear probe JetClass-II, largest set: relative drop in accuracy (%) with the qg data's missing slots set to mean, FLAV_F0 :: sha256 c6c4138ee5c86423 :: prose twin, two significant figures
\newcommand{\MlstLpJcTwoQgDrawAccDropFzero}{53.9, 53.4\%}  % experiments/FIGS/data/v2/index.json :: linear probe JetClass-II, largest set: relative drop in accuracy (%) with the qg data's missing slots set to draws, FLAV_F0 :: sha256 c6c4138ee5c86423
\newcommand{\MlstLpJcTwoQgDrawAccDropFzeroPr}{53--54\%}  % experiments/FIGS/data/v2/index.json :: linear probe JetClass-II, largest set: relative drop in accuracy (%) with the qg data's missing slots set to draws, FLAV_F0 :: sha256 c6c4138ee5c86423 :: prose twin, two significant figures
\newcommand{\MlstLpJcTwoQgZeroAccDropFone}{66.2, 62.8\%}  % experiments/FIGS/data/v2/index.json :: linear probe JetClass-II, largest set: relative drop in accuracy (%) with the qg data's missing slots set to zero, FLAV_F1 :: sha256 c6c4138ee5c86423
\newcommand{\MlstLpJcTwoQgZeroAccDropFonePr}{63--66\%}  % experiments/FIGS/data/v2/index.json :: linear probe JetClass-II, largest set: relative drop in accuracy (%) with the qg data's missing slots set to zero, FLAV_F1 :: sha256 c6c4138ee5c86423 :: prose twin, two significant figures
\newcommand{\MlstLpJcTwoQgMeanAccDropFone}{52.7, 52.8\%}  % experiments/FIGS/data/v2/index.json :: linear probe JetClass-II, largest set: relative drop in accuracy (%) with the qg data's missing slots set to mean, FLAV_F1 :: sha256 c6c4138ee5c86423
\newcommand{\MlstLpJcTwoQgMeanAccDropFonePr}{53--53\%}  % experiments/FIGS/data/v2/index.json :: linear probe JetClass-II, largest set: relative drop in accuracy (%) with the qg data's missing slots set to mean, FLAV_F1 :: sha256 c6c4138ee5c86423 :: prose twin, two significant figures
\newcommand{\MlstLpJcTwoQgDrawAccDropFone}{55.6, 56.5\%}  % experiments/FIGS/data/v2/index.json :: linear probe JetClass-II, largest set: relative drop in accuracy (%) with the qg data's missing slots set to draws, FLAV_F1 :: sha256 c6c4138ee5c86423
\newcommand{\MlstLpJcTwoQgDrawAccDropFonePr}{56--56\%}  % experiments/FIGS/data/v2/index.json :: linear probe JetClass-II, largest set: relative drop in accuracy (%) with the qg data's missing slots set to draws, FLAV_F1 :: sha256 c6c4138ee5c86423 :: prose twin, two significant figures
\newcommand{\MlstLpJcTwoQgZeroAccDropFoner}{67.2, 69.0\%}  % experiments/FIGS/data/v2/index.json :: linear probe JetClass-II, largest set: relative drop in accuracy (%) with the qg data's missing slots set to zero, FLAV_F1R :: sha256 c6c4138ee5c86423
\newcommand{\MlstLpJcTwoQgZeroAccDropFonerPr}{67--69\%}  % experiments/FIGS/data/v2/index.json :: linear probe JetClass-II, largest set: relative drop in accuracy (%) with the qg data's missing slots set to zero, FLAV_F1R :: sha256 c6c4138ee5c86423 :: prose twin, two significant figures
\newcommand{\MlstLpJcTwoQgMeanAccDropFoner}{52.9, 51.8\%}  % experiments/FIGS/data/v2/index.json :: linear probe JetClass-II, largest set: relative drop in accuracy (%) with the qg data's missing slots set to mean, FLAV_F1R :: sha256 c6c4138ee5c86423
\newcommand{\MlstLpJcTwoQgMeanAccDropFonerPr}{52--53\%}  % experiments/FIGS/data/v2/index.json :: linear probe JetClass-II, largest set: relative drop in accuracy (%) with the qg data's missing slots set to mean, FLAV_F1R :: sha256 c6c4138ee5c86423 :: prose twin, two significant figures
\newcommand{\MlstLpJcTwoQgDrawAccDropFoner}{57.2, 55.4\%}  % experiments/FIGS/data/v2/index.json :: linear probe JetClass-II, largest set: relative drop in accuracy (%) with the qg data's missing slots set to draws, FLAV_F1R :: sha256 c6c4138ee5c86423
\newcommand{\MlstLpJcTwoQgDrawAccDropFonerPr}{55--57\%}  % experiments/FIGS/data/v2/index.json :: linear probe JetClass-II, largest set: relative drop in accuracy (%) with the qg data's missing slots set to draws, FLAV_F1R :: sha256 c6c4138ee5c86423 :: prose twin, two significant figures
\newcommand{\MlstLpJcTwoQgZeroAccDropInitBn}{8.7\,\ensuremath{\pm}\,2.6\%}  % experiments/FIGS/data/v2/index.json :: linear probe JetClass-II, largest set: relative drop in accuracy (%) with the qg data's missing slots set to zero, INIT_BN :: sha256 c6c4138ee5c86423
\newcommand{\MlstLpJcTwoQgZeroAccDropInitBnPr}{8.7\%}  % experiments/FIGS/data/v2/index.json :: linear probe JetClass-II, largest set: relative drop in accuracy (%) with the qg data's missing slots set to zero, INIT_BN :: sha256 c6c4138ee5c86423 :: prose twin, two significant figures
\newcommand{\MlstLpJcTwoQgMeanAccDropInitBn}{10.4\,\ensuremath{\pm}\,1.4\%}  % experiments/FIGS/data/v2/index.json :: linear probe JetClass-II, largest set: relative drop in accuracy (%) with the qg data's missing slots set to mean, INIT_BN :: sha256 c6c4138ee5c86423
\newcommand{\MlstLpJcTwoQgMeanAccDropInitBnPr}{10\%}  % experiments/FIGS/data/v2/index.json :: linear probe JetClass-II, largest set: relative drop in accuracy (%) with the qg data's missing slots set to mean, INIT_BN :: sha256 c6c4138ee5c86423 :: prose twin, two significant figures
\newcommand{\MlstLpJcTwoQgDrawAccDropInitBn}{11.4\,\ensuremath{\pm}\,2.1\%}  % experiments/FIGS/data/v2/index.json :: linear probe JetClass-II, largest set: relative drop in accuracy (%) with the qg data's missing slots set to draws, INIT_BN :: sha256 c6c4138ee5c86423
\newcommand{\MlstLpJcTwoQgDrawAccDropInitBnPr}{11\%}  % experiments/FIGS/data/v2/index.json :: linear probe JetClass-II, largest set: relative drop in accuracy (%) with the qg data's missing slots set to draws, INIT_BN :: sha256 c6c4138ee5c86423 :: prose twin, two significant figures
\newcommand{\MlstMassLambda}{5}  % configs/arms/v2_grid.json :: arms[L162_MASS, R16_Q1_MASS].mass_lambda :: sha256 d8a31b1486992302 :: weight of the mass-regression term in the pretraining loss, both mass models
\newcommand{\MlstMassLossShareResMass}{13.78 and 13.87}  % experiments/FIGS/data/v2/index.json :: share of the mass-regression term in the training loss of L162_MASS, mean over epochs, percent :: sha256 c6c4138ee5c86423
\newcommand{\MlstMassLossShareCoarseMass}{28.65 and 28.67}  % experiments/FIGS/data/v2/index.json :: share of the mass-regression term in the training loss of R16_Q1_MASS, mean over epochs, percent :: sha256 c6c4138ee5c86423
\newcommand{\MlstMassUncorrected}{0.245}  % experiments/FIGS/data/v2/index.json :: head/atlas_res: groomed mass, all jets, ATLAS resolution :: sha256 c6c4138ee5c86423
\newcommand{\MlstMassUncorrectedResponse}{0.929}  % experiments/FIGS/data/v2/index.json :: head/atlas_res: groomed mass, all jets, median response :: sha256 c6c4138ee5c86423
\newcommand{\MlstMassHeadResMass}{0.1471 and 0.1463}  % experiments/FIGS/data/v2/index.json :: regression output of L162_MASS, all jets, ATLAS resolution :: sha256 c6c4138ee5c86423
\newcommand{\MlstMassHeadResMassResponse}{0.995 and 1.003}  % experiments/FIGS/data/v2/index.json :: regression output of L162_MASS, all jets, median response :: sha256 c6c4138ee5c86423
\newcommand{\MlstMassHeadResMassOverUncorr}{0.6009 and 0.5976}  % experiments/FIGS/data/v2/index.json :: resolution of the regression output of L162_MASS over the groomed mass's, all jets :: sha256 c6c4138ee5c86423
\newcommand{\MlstMassHeadResMassOverUncorrPr}{at most 0.60 times the resolution}  % experiments/FIGS/data/v2/index.json :: resolution of the regression output of L162_MASS over the groomed mass's, all jets :: sha256 c6c4138ee5c86423 :: prose twin: two significant figures, direction in words
\newcommand{\MlstMassHeadResMassOverUncorrPrMax}{down to 0.60 times the resolution}  % experiments/FIGS/data/v2/index.json :: resolution of the regression output of L162_MASS over the groomed mass's, all jets :: sha256 c6c4138ee5c86423 :: prose twin of the run farther from 1
\newcommand{\MlstMassHeadResMassOverMlpRes}{0.6645 and 0.6691}  % experiments/FIGS/data/v2/index.json :: resolution of the regression output of L162_MASS over the frozen-feature MLP of L162, run by run, all jets :: sha256 c6c4138ee5c86423
\newcommand{\MlstMassHeadResMassOverMlpResPr}{at most 0.67 times the resolution}  % experiments/FIGS/data/v2/index.json :: resolution of the regression output of L162_MASS over the frozen-feature MLP of L162, run by run, all jets :: sha256 c6c4138ee5c86423 :: prose twin: two significant figures, direction in words
\newcommand{\MlstMassHeadResMassOverMlpResPrMax}{down to 0.66 times the resolution}  % experiments/FIGS/data/v2/index.json :: resolution of the regression output of L162_MASS over the frozen-feature MLP of L162, run by run, all jets :: sha256 c6c4138ee5c86423 :: prose twin of the run farther from 1
\newcommand{\MlstMassHeadCoarseMass}{0.1450 and 0.1461}  % experiments/FIGS/data/v2/index.json :: regression output of R16_Q1_MASS, all jets, ATLAS resolution :: sha256 c6c4138ee5c86423
\newcommand{\MlstMassHeadCoarseMassResponse}{1.005 and 1.005}  % experiments/FIGS/data/v2/index.json :: regression output of R16_Q1_MASS, all jets, median response :: sha256 c6c4138ee5c86423
\newcommand{\MlstMassHeadCoarseMassOverUncorr}{0.5921 and 0.5966}  % experiments/FIGS/data/v2/index.json :: resolution of the regression output of R16_Q1_MASS over the groomed mass's, all jets :: sha256 c6c4138ee5c86423
\newcommand{\MlstMassHeadCoarseMassOverUncorrPr}{at most 0.60 times the resolution}  % experiments/FIGS/data/v2/index.json :: resolution of the regression output of R16_Q1_MASS over the groomed mass's, all jets :: sha256 c6c4138ee5c86423 :: prose twin: two significant figures, direction in words
\newcommand{\MlstMassHeadCoarseMassOverUncorrPrMax}{down to 0.59 times the resolution}  % experiments/FIGS/data/v2/index.json :: resolution of the regression output of R16_Q1_MASS over the groomed mass's, all jets :: sha256 c6c4138ee5c86423 :: prose twin of the run farther from 1
\newcommand{\MlstMassHeadCoarseMassOverMlpCoarse}{0.703 and 0.690}  % experiments/FIGS/data/v2/index.json :: resolution of the regression output of R16_Q1_MASS over the frozen-feature MLP of R16_Q1, run by run, all jets :: sha256 c6c4138ee5c86423
\newcommand{\MlstMassHeadCoarseMassOverMlpCoarsePr}{at most 0.70 times the resolution}  % experiments/FIGS/data/v2/index.json :: resolution of the regression output of R16_Q1_MASS over the frozen-feature MLP of R16_Q1, run by run, all jets :: sha256 c6c4138ee5c86423 :: prose twin: two significant figures, direction in words
\newcommand{\MlstMassHeadCoarseMassOverMlpCoarsePrMax}{down to 0.69 times the resolution}  % experiments/FIGS/data/v2/index.json :: resolution of the regression output of R16_Q1_MASS over the frozen-feature MLP of R16_Q1, run by run, all jets :: sha256 c6c4138ee5c86423 :: prose twin of the run farther from 1
\newcommand{\MlstMassMlpInit}{0.2421\,\ensuremath{\pm}\,0.0018}  % experiments/FIGS/data/v2/index.json :: frozen-feature MLP of INIT, all jets, ATLAS resolution :: sha256 c6c4138ee5c86423
\newcommand{\MlstMassMlpFull}{0.2200\,\ensuremath{\pm}\,0.0014}  % experiments/FIGS/data/v2/index.json :: frozen-feature MLP of L188, all jets, ATLAS resolution :: sha256 c6c4138ee5c86423
\newcommand{\MlstMassMlpRes}{0.2195\,\ensuremath{\pm}\,0.0017}  % experiments/FIGS/data/v2/index.json :: frozen-feature MLP of L162, all jets, ATLAS resolution :: sha256 c6c4138ee5c86423
\newcommand{\MlstMassMlpMid}{0.2158\,\ensuremath{\pm}\,0.0029}  % experiments/FIGS/data/v2/index.json :: frozen-feature MLP of R42_Q1, all jets, ATLAS resolution :: sha256 c6c4138ee5c86423
\newcommand{\MlstMassMlpCoarse}{0.2098\,\ensuremath{\pm}\,0.0031}  % experiments/FIGS/data/v2/index.json :: frozen-feature MLP of R16_Q1, all jets, ATLAS resolution :: sha256 c6c4138ee5c86423
\newcommand{\MlstMassMlpResMass}{0.1591 and 0.1586}  % experiments/FIGS/data/v2/index.json :: frozen-feature MLP of L162_MASS, all jets, ATLAS resolution :: sha256 c6c4138ee5c86423
\newcommand{\MlstMassMlpCoarseMass}{0.1583 and 0.1588}  % experiments/FIGS/data/v2/index.json :: frozen-feature MLP of R16_Q1_MASS, all jets, ATLAS resolution :: sha256 c6c4138ee5c86423
\newcommand{\MlstCmpMlpResOverFull}{0.998\,\ensuremath{\pm}\,0.006}  % experiments/FIGS/data/v2/index.json :: frozen-feature MLP resolution of L162 over L188, run by run, all jets :: sha256 c6c4138ee5c86423
\newcommand{\MlstCmpMlpResOverFullPr}{equal within 1\%}  % experiments/FIGS/data/v2/index.json :: frozen-feature MLP resolution of L162 over L188, run by run, all jets :: sha256 c6c4138ee5c86423 :: prose twin: two significant figures, direction in words
\newcommand{\MlstCmpMlpMidOverFull}{0.981\,\ensuremath{\pm}\,0.017}  % experiments/FIGS/data/v2/index.json :: frozen-feature MLP resolution of R42_Q1 over L188, run by run, all jets :: sha256 c6c4138ee5c86423
\newcommand{\MlstCmpMlpMidOverFullPr}{equal within 2\%}  % experiments/FIGS/data/v2/index.json :: frozen-feature MLP resolution of R42_Q1 over L188, run by run, all jets :: sha256 c6c4138ee5c86423 :: prose twin: two significant figures, direction in words
\newcommand{\MlstCmpMlpCoarseOverFull}{0.954\,\ensuremath{\pm}\,0.018}  % experiments/FIGS/data/v2/index.json :: frozen-feature MLP resolution of R16_Q1 over L188, run by run, all jets :: sha256 c6c4138ee5c86423
\newcommand{\MlstCmpMlpCoarseOverFullPr}{0.95 times the resolution}  % experiments/FIGS/data/v2/index.json :: frozen-feature MLP resolution of R16_Q1 over L188, run by run, all jets :: sha256 c6c4138ee5c86423 :: prose twin: two significant figures, direction in words
\newcommand{\MlstCmpMlpResMassOverRes}{0.7186 and 0.7253}  % experiments/FIGS/data/v2/index.json :: frozen-feature MLP resolution of L162_MASS over L162, run by run, all jets :: sha256 c6c4138ee5c86423
\newcommand{\MlstCmpMlpResMassOverResPr}{at most 0.73 times the resolution}  % experiments/FIGS/data/v2/index.json :: frozen-feature MLP resolution of L162_MASS over L162, run by run, all jets :: sha256 c6c4138ee5c86423 :: prose twin: two significant figures, direction in words
\newcommand{\MlstCmpMlpResMassOverResPrMax}{down to 0.72 times the resolution}  % experiments/FIGS/data/v2/index.json :: frozen-feature MLP resolution of L162_MASS over L162, run by run, all jets :: sha256 c6c4138ee5c86423 :: prose twin of the run farther from 1
\newcommand{\MlstCmpMlpCoarseMassOverCoarse}{0.768 and 0.750}  % experiments/FIGS/data/v2/index.json :: frozen-feature MLP resolution of R16_Q1_MASS over R16_Q1, run by run, all jets :: sha256 c6c4138ee5c86423
\newcommand{\MlstCmpMlpCoarseMassOverCoarsePr}{at most 0.77 times the resolution}  % experiments/FIGS/data/v2/index.json :: frozen-feature MLP resolution of R16_Q1_MASS over R16_Q1, run by run, all jets :: sha256 c6c4138ee5c86423 :: prose twin: two significant figures, direction in words
\newcommand{\MlstCmpMlpCoarseMassOverCoarsePrMax}{down to 0.75 times the resolution}  % experiments/FIGS/data/v2/index.json :: frozen-feature MLP resolution of R16_Q1_MASS over R16_Q1, run by run, all jets :: sha256 c6c4138ee5c86423 :: prose twin of the run farther from 1
\newcommand{\MlstMassUncorrectedResonant}{0.213}  % experiments/FIGS/data/v2/index.json :: head/atlas_res: groomed mass, resonant jets, ATLAS resolution :: sha256 c6c4138ee5c86423
\newcommand{\MlstMassUncorrectedResponseResonant}{0.929}  % experiments/FIGS/data/v2/index.json :: head/atlas_res: groomed mass, resonant jets, median response :: sha256 c6c4138ee5c86423
\newcommand{\MlstMassHeadResMassResonant}{0.1332 and 0.1320}  % experiments/FIGS/data/v2/index.json :: regression output of L162_MASS, resonant jets, ATLAS resolution :: sha256 c6c4138ee5c86423
\newcommand{\MlstMassHeadResMassResponseResonant}{0.996 and 1.004}  % experiments/FIGS/data/v2/index.json :: regression output of L162_MASS, resonant jets, median response :: sha256 c6c4138ee5c86423
\newcommand{\MlstMassHeadResMassOverUncorrResonant}{0.6259 and 0.6205}  % experiments/FIGS/data/v2/index.json :: resolution of the regression output of L162_MASS over the groomed mass's, resonant jets :: sha256 c6c4138ee5c86423
\newcommand{\MlstMassHeadResMassOverUncorrResonantPr}{at most 0.63 times the resolution}  % experiments/FIGS/data/v2/index.json :: resolution of the regression output of L162_MASS over the groomed mass's, resonant jets :: sha256 c6c4138ee5c86423 :: prose twin: two significant figures, direction in words
\newcommand{\MlstMassHeadResMassOverUncorrResonantPrMax}{down to 0.62 times the resolution}  % experiments/FIGS/data/v2/index.json :: resolution of the regression output of L162_MASS over the groomed mass's, resonant jets :: sha256 c6c4138ee5c86423 :: prose twin of the run farther from 1
\newcommand{\MlstMassHeadResMassOverMlpResResonant}{0.6620 and 0.6662}  % experiments/FIGS/data/v2/index.json :: resolution of the regression output of L162_MASS over the frozen-feature MLP of L162, run by run, resonant jets :: sha256 c6c4138ee5c86423
\newcommand{\MlstMassHeadResMassOverMlpResResonantPr}{at most 0.67 times the resolution}  % experiments/FIGS/data/v2/index.json :: resolution of the regression output of L162_MASS over the frozen-feature MLP of L162, run by run, resonant jets :: sha256 c6c4138ee5c86423 :: prose twin: two significant figures, direction in words
\newcommand{\MlstMassHeadResMassOverMlpResResonantPrMax}{down to 0.66 times the resolution}  % experiments/FIGS/data/v2/index.json :: resolution of the regression output of L162_MASS over the frozen-feature MLP of L162, run by run, resonant jets :: sha256 c6c4138ee5c86423 :: prose twin of the run farther from 1
\newcommand{\MlstMassHeadCoarseMassResonant}{0.1312 and 0.1320}  % experiments/FIGS/data/v2/index.json :: regression output of R16_Q1_MASS, resonant jets, ATLAS resolution :: sha256 c6c4138ee5c86423
\newcommand{\MlstMassHeadCoarseMassResponseResonant}{1.006 and 1.006}  % experiments/FIGS/data/v2/index.json :: regression output of R16_Q1_MASS, resonant jets, median response :: sha256 c6c4138ee5c86423
\newcommand{\MlstMassHeadCoarseMassOverUncorrResonant}{0.6165 and 0.6206}  % experiments/FIGS/data/v2/index.json :: resolution of the regression output of R16_Q1_MASS over the groomed mass's, resonant jets :: sha256 c6c4138ee5c86423
\newcommand{\MlstMassHeadCoarseMassOverUncorrResonantPr}{at most 0.62 times the resolution}  % experiments/FIGS/data/v2/index.json :: resolution of the regression output of R16_Q1_MASS over the groomed mass's, resonant jets :: sha256 c6c4138ee5c86423 :: prose twin: two significant figures, direction in words
\newcommand{\MlstMassHeadCoarseMassOverUncorrResonantPrMax}{down to 0.62 times the resolution}  % experiments/FIGS/data/v2/index.json :: resolution of the regression output of R16_Q1_MASS over the groomed mass's, resonant jets :: sha256 c6c4138ee5c86423 :: prose twin of the run farther from 1
\newcommand{\MlstMassHeadCoarseMassOverMlpCoarseResonant}{0.707 and 0.688}  % experiments/FIGS/data/v2/index.json :: resolution of the regression output of R16_Q1_MASS over the frozen-feature MLP of R16_Q1, run by run, resonant jets :: sha256 c6c4138ee5c86423
\newcommand{\MlstMassHeadCoarseMassOverMlpCoarseResonantPr}{at most 0.71 times the resolution}  % experiments/FIGS/data/v2/index.json :: resolution of the regression output of R16_Q1_MASS over the frozen-feature MLP of R16_Q1, run by run, resonant jets :: sha256 c6c4138ee5c86423 :: prose twin: two significant figures, direction in words
\newcommand{\MlstMassHeadCoarseMassOverMlpCoarseResonantPrMax}{down to 0.69 times the resolution}  % experiments/FIGS/data/v2/index.json :: resolution of the regression output of R16_Q1_MASS over the frozen-feature MLP of R16_Q1, run by run, resonant jets :: sha256 c6c4138ee5c86423 :: prose twin of the run farther from 1
\newcommand{\MlstMassMlpInitResonant}{0.2214\,\ensuremath{\pm}\,0.0017}  % experiments/FIGS/data/v2/index.json :: frozen-feature MLP of INIT, resonant jets, ATLAS resolution :: sha256 c6c4138ee5c86423
\newcommand{\MlstMassMlpFullResonant}{0.1990\,\ensuremath{\pm}\,0.0016}  % experiments/FIGS/data/v2/index.json :: frozen-feature MLP of L188, resonant jets, ATLAS resolution :: sha256 c6c4138ee5c86423
\newcommand{\MlstMassMlpResResonant}{0.1989\,\ensuremath{\pm}\,0.0020}  % experiments/FIGS/data/v2/index.json :: frozen-feature MLP of L162, resonant jets, ATLAS resolution :: sha256 c6c4138ee5c86423
\newcommand{\MlstMassMlpMidResonant}{0.1949\,\ensuremath{\pm}\,0.0023}  % experiments/FIGS/data/v2/index.json :: frozen-feature MLP of R42_Q1, resonant jets, ATLAS resolution :: sha256 c6c4138ee5c86423
\newcommand{\MlstMassMlpCoarseResonant}{0.1894\,\ensuremath{\pm}\,0.0035}  % experiments/FIGS/data/v2/index.json :: frozen-feature MLP of R16_Q1, resonant jets, ATLAS resolution :: sha256 c6c4138ee5c86423
\newcommand{\MlstMassMlpResMassResonant}{0.1436 and 0.1430}  % experiments/FIGS/data/v2/index.json :: frozen-feature MLP of L162_MASS, resonant jets, ATLAS resolution :: sha256 c6c4138ee5c86423
\newcommand{\MlstMassMlpCoarseMassResonant}{0.1420 and 0.1431}  % experiments/FIGS/data/v2/index.json :: frozen-feature MLP of R16_Q1_MASS, resonant jets, ATLAS resolution :: sha256 c6c4138ee5c86423
\newcommand{\MlstCmpMlpResOverFullResonant}{1.000\,\ensuremath{\pm}\,0.005}  % experiments/FIGS/data/v2/index.json :: frozen-feature MLP resolution of L162 over L188, run by run, resonant jets :: sha256 c6c4138ee5c86423
\newcommand{\MlstCmpMlpResOverFullResonantPr}{equal within 1\%}  % experiments/FIGS/data/v2/index.json :: frozen-feature MLP resolution of L162 over L188, run by run, resonant jets :: sha256 c6c4138ee5c86423 :: prose twin: two significant figures, direction in words
\newcommand{\MlstCmpMlpMidOverFullResonant}{0.980\,\ensuremath{\pm}\,0.018}  % experiments/FIGS/data/v2/index.json :: frozen-feature MLP resolution of R42_Q1 over L188, run by run, resonant jets :: sha256 c6c4138ee5c86423
\newcommand{\MlstCmpMlpMidOverFullResonantPr}{equal within 3\%}  % experiments/FIGS/data/v2/index.json :: frozen-feature MLP resolution of R42_Q1 over L188, run by run, resonant jets :: sha256 c6c4138ee5c86423 :: prose twin: two significant figures, direction in words
\newcommand{\MlstCmpMlpCoarseOverFullResonant}{0.952\,\ensuremath{\pm}\,0.023}  % experiments/FIGS/data/v2/index.json :: frozen-feature MLP resolution of R16_Q1 over L188, run by run, resonant jets :: sha256 c6c4138ee5c86423
\newcommand{\MlstCmpMlpCoarseOverFullResonantPr}{0.95 times the resolution}  % experiments/FIGS/data/v2/index.json :: frozen-feature MLP resolution of R16_Q1 over L188, run by run, resonant jets :: sha256 c6c4138ee5c86423 :: prose twin: two significant figures, direction in words
\newcommand{\MlstCmpMlpResMassOverResResonant}{0.7137 and 0.7219}  % experiments/FIGS/data/v2/index.json :: frozen-feature MLP resolution of L162_MASS over L162, run by run, resonant jets :: sha256 c6c4138ee5c86423
\newcommand{\MlstCmpMlpResMassOverResResonantPr}{at most 0.72 times the resolution}  % experiments/FIGS/data/v2/index.json :: frozen-feature MLP resolution of L162_MASS over L162, run by run, resonant jets :: sha256 c6c4138ee5c86423 :: prose twin: two significant figures, direction in words
\newcommand{\MlstCmpMlpResMassOverResResonantPrMax}{down to 0.71 times the resolution}  % experiments/FIGS/data/v2/index.json :: frozen-feature MLP resolution of L162_MASS over L162, run by run, resonant jets :: sha256 c6c4138ee5c86423 :: prose twin of the run farther from 1
\newcommand{\MlstCmpMlpCoarseMassOverCoarseResonant}{0.766 and 0.745}  % experiments/FIGS/data/v2/index.json :: frozen-feature MLP resolution of R16_Q1_MASS over R16_Q1, run by run, resonant jets :: sha256 c6c4138ee5c86423
\newcommand{\MlstCmpMlpCoarseMassOverCoarseResonantPr}{at most 0.77 times the resolution}  % experiments/FIGS/data/v2/index.json :: frozen-feature MLP resolution of R16_Q1_MASS over R16_Q1, run by run, resonant jets :: sha256 c6c4138ee5c86423 :: prose twin: two significant figures, direction in words
\newcommand{\MlstCmpMlpCoarseMassOverCoarseResonantPrMax}{down to 0.75 times the resolution}  % experiments/FIGS/data/v2/index.json :: frozen-feature MLP resolution of R16_Q1_MASS over R16_Q1, run by run, resonant jets :: sha256 c6c4138ee5c86423 :: prose twin of the run farther from 1
\newcommand{\MlstMassUncorrectedQcd}{1.52}  % experiments/FIGS/data/v2/index.json :: head/atlas_res: groomed mass, qcd jets, ATLAS resolution :: sha256 c6c4138ee5c86423
\newcommand{\MlstMassUncorrectedResponseQcd}{0.931}  % experiments/FIGS/data/v2/index.json :: head/atlas_res: groomed mass, qcd jets, median response :: sha256 c6c4138ee5c86423
\newcommand{\MlstMassHeadResMassQcd}{0.2263 and 0.2256}  % experiments/FIGS/data/v2/index.json :: regression output of L162_MASS, qcd jets, ATLAS resolution :: sha256 c6c4138ee5c86423
\newcommand{\MlstMassHeadResMassResponseQcd}{0.9897 and 0.9996}  % experiments/FIGS/data/v2/index.json :: regression output of L162_MASS, qcd jets, median response :: sha256 c6c4138ee5c86423
\newcommand{\MlstMassHeadResMassOverUncorrQcd}{0.1485 and 0.1480}  % experiments/FIGS/data/v2/index.json :: resolution of the regression output of L162_MASS over the groomed mass's, qcd jets :: sha256 c6c4138ee5c86423
\newcommand{\MlstMassHeadResMassOverUncorrQcdPr}{at most 0.15 times the resolution}  % experiments/FIGS/data/v2/index.json :: resolution of the regression output of L162_MASS over the groomed mass's, qcd jets :: sha256 c6c4138ee5c86423 :: prose twin: two significant figures, direction in words
\newcommand{\MlstMassHeadResMassOverUncorrQcdPrMax}{down to 0.15 times the resolution}  % experiments/FIGS/data/v2/index.json :: resolution of the regression output of L162_MASS over the groomed mass's, qcd jets :: sha256 c6c4138ee5c86423 :: prose twin of the run farther from 1
\newcommand{\MlstMassHeadResMassOverMlpResQcd}{0.6542 and 0.6632}  % experiments/FIGS/data/v2/index.json :: resolution of the regression output of L162_MASS over the frozen-feature MLP of L162, run by run, qcd jets :: sha256 c6c4138ee5c86423
\newcommand{\MlstMassHeadResMassOverMlpResQcdPr}{at most 0.66 times the resolution}  % experiments/FIGS/data/v2/index.json :: resolution of the regression output of L162_MASS over the frozen-feature MLP of L162, run by run, qcd jets :: sha256 c6c4138ee5c86423 :: prose twin: two significant figures, direction in words
\newcommand{\MlstMassHeadResMassOverMlpResQcdPrMax}{down to 0.65 times the resolution}  % experiments/FIGS/data/v2/index.json :: resolution of the regression output of L162_MASS over the frozen-feature MLP of L162, run by run, qcd jets :: sha256 c6c4138ee5c86423 :: prose twin of the run farther from 1
\newcommand{\MlstMassHeadCoarseMassQcd}{0.2239 and 0.2240}  % experiments/FIGS/data/v2/index.json :: regression output of R16_Q1_MASS, qcd jets, ATLAS resolution :: sha256 c6c4138ee5c86423
\newcommand{\MlstMassHeadCoarseMassResponseQcd}{1.001 and 1.000}  % experiments/FIGS/data/v2/index.json :: regression output of R16_Q1_MASS, qcd jets, median response :: sha256 c6c4138ee5c86423
\newcommand{\MlstMassHeadCoarseMassOverUncorrQcd}{0.1469 and 0.1469}  % experiments/FIGS/data/v2/index.json :: resolution of the regression output of R16_Q1_MASS over the groomed mass's, qcd jets :: sha256 c6c4138ee5c86423
\newcommand{\MlstMassHeadCoarseMassOverUncorrQcdPr}{at most 0.15 times the resolution}  % experiments/FIGS/data/v2/index.json :: resolution of the regression output of R16_Q1_MASS over the groomed mass's, qcd jets :: sha256 c6c4138ee5c86423 :: prose twin: two significant figures, direction in words
\newcommand{\MlstMassHeadCoarseMassOverUncorrQcdPrMax}{down to 0.15 times the resolution}  % experiments/FIGS/data/v2/index.json :: resolution of the regression output of R16_Q1_MASS over the groomed mass's, qcd jets :: sha256 c6c4138ee5c86423 :: prose twin of the run farther from 1
\newcommand{\MlstMassHeadCoarseMassOverMlpCoarseQcd}{0.679 and 0.698}  % experiments/FIGS/data/v2/index.json :: resolution of the regression output of R16_Q1_MASS over the frozen-feature MLP of R16_Q1, run by run, qcd jets :: sha256 c6c4138ee5c86423
\newcommand{\MlstMassHeadCoarseMassOverMlpCoarseQcdPr}{at most 0.70 times the resolution}  % experiments/FIGS/data/v2/index.json :: resolution of the regression output of R16_Q1_MASS over the frozen-feature MLP of R16_Q1, run by run, qcd jets :: sha256 c6c4138ee5c86423 :: prose twin: two significant figures, direction in words
\newcommand{\MlstMassHeadCoarseMassOverMlpCoarseQcdPrMax}{down to 0.68 times the resolution}  % experiments/FIGS/data/v2/index.json :: resolution of the regression output of R16_Q1_MASS over the frozen-feature MLP of R16_Q1, run by run, qcd jets :: sha256 c6c4138ee5c86423 :: prose twin of the run farther from 1
\newcommand{\MlstMassMlpInitQcd}{0.502\,\ensuremath{\pm}\,0.014}  % experiments/FIGS/data/v2/index.json :: frozen-feature MLP of INIT, qcd jets, ATLAS resolution :: sha256 c6c4138ee5c86423
\newcommand{\MlstMassMlpFullQcd}{0.3487\,\ensuremath{\pm}\,0.0029}  % experiments/FIGS/data/v2/index.json :: frozen-feature MLP of L188, qcd jets, ATLAS resolution :: sha256 c6c4138ee5c86423
\newcommand{\MlstMassMlpResQcd}{0.3425\,\ensuremath{\pm}\,0.0031}  % experiments/FIGS/data/v2/index.json :: frozen-feature MLP of L162, qcd jets, ATLAS resolution :: sha256 c6c4138ee5c86423
\newcommand{\MlstMassMlpMidQcd}{0.338\,\ensuremath{\pm}\,0.007}  % experiments/FIGS/data/v2/index.json :: frozen-feature MLP of R42_Q1, qcd jets, ATLAS resolution :: sha256 c6c4138ee5c86423
\newcommand{\MlstMassMlpCoarseQcd}{0.327\,\ensuremath{\pm}\,0.006}  % experiments/FIGS/data/v2/index.json :: frozen-feature MLP of R16_Q1, qcd jets, ATLAS resolution :: sha256 c6c4138ee5c86423
\newcommand{\MlstMassMlpResMassQcd}{0.2473 and 0.2452}  % experiments/FIGS/data/v2/index.json :: frozen-feature MLP of L162_MASS, qcd jets, ATLAS resolution :: sha256 c6c4138ee5c86423
\newcommand{\MlstMassMlpCoarseMassQcd}{0.2465 and 0.2463}  % experiments/FIGS/data/v2/index.json :: frozen-feature MLP of R16_Q1_MASS, qcd jets, ATLAS resolution :: sha256 c6c4138ee5c86423
\newcommand{\MlstCmpMlpResOverFullQcd}{0.982\,\ensuremath{\pm}\,0.017}  % experiments/FIGS/data/v2/index.json :: frozen-feature MLP resolution of L162 over L188, run by run, qcd jets :: sha256 c6c4138ee5c86423
\newcommand{\MlstCmpMlpResOverFullQcdPr}{equal within 2\%}  % experiments/FIGS/data/v2/index.json :: frozen-feature MLP resolution of L162 over L188, run by run, qcd jets :: sha256 c6c4138ee5c86423 :: prose twin: two significant figures, direction in words
\newcommand{\MlstCmpMlpMidOverFullQcd}{0.969\,\ensuremath{\pm}\,0.013}  % experiments/FIGS/data/v2/index.json :: frozen-feature MLP resolution of R42_Q1 over L188, run by run, qcd jets :: sha256 c6c4138ee5c86423
\newcommand{\MlstCmpMlpMidOverFullQcdPr}{0.97 times the resolution}  % experiments/FIGS/data/v2/index.json :: frozen-feature MLP resolution of R42_Q1 over L188, run by run, qcd jets :: sha256 c6c4138ee5c86423 :: prose twin: two significant figures, direction in words
\newcommand{\MlstCmpMlpCoarseOverFullQcd}{0.939\,\ensuremath{\pm}\,0.023}  % experiments/FIGS/data/v2/index.json :: frozen-feature MLP resolution of R16_Q1 over L188, run by run, qcd jets :: sha256 c6c4138ee5c86423
\newcommand{\MlstCmpMlpCoarseOverFullQcdPr}{0.94 times the resolution}  % experiments/FIGS/data/v2/index.json :: frozen-feature MLP resolution of R16_Q1 over L188, run by run, qcd jets :: sha256 c6c4138ee5c86423 :: prose twin: two significant figures, direction in words
\newcommand{\MlstCmpMlpResMassOverResQcd}{0.7147 and 0.7209}  % experiments/FIGS/data/v2/index.json :: frozen-feature MLP resolution of L162_MASS over L162, run by run, qcd jets :: sha256 c6c4138ee5c86423
\newcommand{\MlstCmpMlpResMassOverResQcdPr}{at most 0.72 times the resolution}  % experiments/FIGS/data/v2/index.json :: frozen-feature MLP resolution of L162_MASS over L162, run by run, qcd jets :: sha256 c6c4138ee5c86423 :: prose twin: two significant figures, direction in words
\newcommand{\MlstCmpMlpResMassOverResQcdPrMax}{down to 0.71 times the resolution}  % experiments/FIGS/data/v2/index.json :: frozen-feature MLP resolution of L162_MASS over L162, run by run, qcd jets :: sha256 c6c4138ee5c86423 :: prose twin of the run farther from 1
\newcommand{\MlstCmpMlpCoarseMassOverCoarseQcd}{0.747 and 0.768}  % experiments/FIGS/data/v2/index.json :: frozen-feature MLP resolution of R16_Q1_MASS over R16_Q1, run by run, qcd jets :: sha256 c6c4138ee5c86423
\newcommand{\MlstCmpMlpCoarseMassOverCoarseQcdPr}{at most 0.77 times the resolution}  % experiments/FIGS/data/v2/index.json :: frozen-feature MLP resolution of R16_Q1_MASS over R16_Q1, run by run, qcd jets :: sha256 c6c4138ee5c86423 :: prose twin: two significant figures, direction in words
\newcommand{\MlstCmpMlpCoarseMassOverCoarseQcdPrMax}{down to 0.75 times the resolution}  % experiments/FIGS/data/v2/index.json :: frozen-feature MLP resolution of R16_Q1_MASS over R16_Q1, run by run, qcd jets :: sha256 c6c4138ee5c86423 :: prose twin of the run farther from 1
\newcommand{\MlstMassRespAxesShareQcd}{76\%}  % experiments/FIGS/data/v2/index.json :: head/atlas_res: share of QCD jets whose groomed-mass response lies inside F9a(b)'s x axis :: sha256 c6c4138ee5c86423
\newcommand{\MlstSculptNPassTen}{41{,}028}  % experiments/FIGS/data/v2/index.json :: QCD test jets passing the cut at QCD efficiency 0.1, every run :: sha256 c6c4138ee5c86423
\newcommand{\MlstSculptFloorTen}{0.00944}  % experiments/FIGS/data/v2/index.json :: Jensen-Shannon distance expected for a random selection of as many QCD jets as pass at 0.1 :: sha256 c6c4138ee5c86423
\newcommand{\MlstSculptFullTen}{0.0524\,\ensuremath{\pm}\,0.0034}  % experiments/FIGS/data/v2/index.json :: Jensen-Shannon distance of the QCD m_SD distribution at QCD efficiency 0.1, L188 :: sha256 c6c4138ee5c86423
\newcommand{\MlstSculptFullOverFloorTen}{5.6\,\ensuremath{\pm}\,0.4}  % experiments/FIGS/data/v2/index.json :: distance at 0.1 over the random-selection floor, L188 :: sha256 c6c4138ee5c86423
\newcommand{\MlstSculptFullOverFloorTenPr}{5.6 times the distance}  % experiments/FIGS/data/v2/index.json :: distance at 0.1 over the random-selection floor, L188 :: sha256 c6c4138ee5c86423 :: prose twin: two significant figures, direction in words
\newcommand{\MlstSculptResTen}{0.0532\,\ensuremath{\pm}\,0.0009}  % experiments/FIGS/data/v2/index.json :: Jensen-Shannon distance of the QCD m_SD distribution at QCD efficiency 0.1, L162 :: sha256 c6c4138ee5c86423
\newcommand{\MlstSculptResOverFloorTen}{5.64\,\ensuremath{\pm}\,0.10}  % experiments/FIGS/data/v2/index.json :: distance at 0.1 over the random-selection floor, L162 :: sha256 c6c4138ee5c86423
\newcommand{\MlstSculptResOverFloorTenPr}{5.6 times the distance}  % experiments/FIGS/data/v2/index.json :: distance at 0.1 over the random-selection floor, L162 :: sha256 c6c4138ee5c86423 :: prose twin: two significant figures, direction in words
\newcommand{\MlstSculptMidTen}{0.0560\,\ensuremath{\pm}\,0.0027}  % experiments/FIGS/data/v2/index.json :: Jensen-Shannon distance of the QCD m_SD distribution at QCD efficiency 0.1, R42_Q1 :: sha256 c6c4138ee5c86423
\newcommand{\MlstSculptMidOverFloorTen}{5.94\,\ensuremath{\pm}\,0.28}  % experiments/FIGS/data/v2/index.json :: distance at 0.1 over the random-selection floor, R42_Q1 :: sha256 c6c4138ee5c86423
\newcommand{\MlstSculptMidOverFloorTenPr}{5.9 times the distance}  % experiments/FIGS/data/v2/index.json :: distance at 0.1 over the random-selection floor, R42_Q1 :: sha256 c6c4138ee5c86423 :: prose twin: two significant figures, direction in words
\newcommand{\MlstSculptCoarseTen}{0.0543\,\ensuremath{\pm}\,0.0010}  % experiments/FIGS/data/v2/index.json :: Jensen-Shannon distance of the QCD m_SD distribution at QCD efficiency 0.1, R16_Q1 :: sha256 c6c4138ee5c86423
\newcommand{\MlstSculptCoarseOverFloorTen}{5.76\,\ensuremath{\pm}\,0.10}  % experiments/FIGS/data/v2/index.json :: distance at 0.1 over the random-selection floor, R16_Q1 :: sha256 c6c4138ee5c86423
\newcommand{\MlstSculptCoarseOverFloorTenPr}{5.8 times the distance}  % experiments/FIGS/data/v2/index.json :: distance at 0.1 over the random-selection floor, R16_Q1 :: sha256 c6c4138ee5c86423 :: prose twin: two significant figures, direction in words
\newcommand{\MlstSculptResMassTen}{0.0422 and 0.0464}  % experiments/FIGS/data/v2/index.json :: Jensen-Shannon distance of the QCD m_SD distribution at QCD efficiency 0.1, L162_MASS :: sha256 c6c4138ee5c86423
\newcommand{\MlstSculptResMassOverFloorTen}{4.47 and 4.91}  % experiments/FIGS/data/v2/index.json :: distance at 0.1 over the random-selection floor, L162_MASS :: sha256 c6c4138ee5c86423
\newcommand{\MlstSculptResMassOverFloorTenPr}{at least 4.5 times the distance}  % experiments/FIGS/data/v2/index.json :: distance at 0.1 over the random-selection floor, L162_MASS :: sha256 c6c4138ee5c86423 :: prose twin: two significant figures, direction in words
\newcommand{\MlstSculptResMassOverFloorTenPrMax}{up to 4.9 times the distance}  % experiments/FIGS/data/v2/index.json :: distance at 0.1 over the random-selection floor, L162_MASS :: sha256 c6c4138ee5c86423 :: prose twin of the run farther from 1
\newcommand{\MlstSculptCoarseMassTen}{0.0522 and 0.0423}  % experiments/FIGS/data/v2/index.json :: Jensen-Shannon distance of the QCD m_SD distribution at QCD efficiency 0.1, R16_Q1_MASS :: sha256 c6c4138ee5c86423
\newcommand{\MlstSculptCoarseMassOverFloorTen}{5.5 and 4.5}  % experiments/FIGS/data/v2/index.json :: distance at 0.1 over the random-selection floor, R16_Q1_MASS :: sha256 c6c4138ee5c86423
\newcommand{\MlstSculptCoarseMassOverFloorTenPr}{at least 4.5 times the distance}  % experiments/FIGS/data/v2/index.json :: distance at 0.1 over the random-selection floor, R16_Q1_MASS :: sha256 c6c4138ee5c86423 :: prose twin: two significant figures, direction in words
\newcommand{\MlstSculptCoarseMassOverFloorTenPrMax}{up to 5.5 times the distance}  % experiments/FIGS/data/v2/index.json :: distance at 0.1 over the random-selection floor, R16_Q1_MASS :: sha256 c6c4138ee5c86423 :: prose twin of the run farther from 1
\newcommand{\MlstSculptResMassOverResTen}{0.789 and 0.860}  % experiments/FIGS/data/v2/index.json :: distance at QCD efficiency 0.1: L162_MASS over L162, run by run :: sha256 c6c4138ee5c86423
\newcommand{\MlstSculptResMassOverResTenPr}{at most 0.86 times the distance}  % experiments/FIGS/data/v2/index.json :: distance at QCD efficiency 0.1: L162_MASS over L162, run by run :: sha256 c6c4138ee5c86423 :: prose twin: two significant figures, direction in words
\newcommand{\MlstSculptResMassOverResTenPrMax}{down to 0.79 times the distance}  % experiments/FIGS/data/v2/index.json :: distance at QCD efficiency 0.1: L162_MASS over L162, run by run :: sha256 c6c4138ee5c86423 :: prose twin of the run farther from 1
\newcommand{\MlstSculptCoarseMassOverCoarseTen}{0.94 and 0.79}  % experiments/FIGS/data/v2/index.json :: distance at QCD efficiency 0.1: R16_Q1_MASS over R16_Q1, run by run :: sha256 c6c4138ee5c86423
\newcommand{\MlstSculptCoarseMassOverCoarseTenPr}{equal within 30\%}  % experiments/FIGS/data/v2/index.json :: distance at QCD efficiency 0.1: R16_Q1_MASS over R16_Q1, run by run :: sha256 c6c4138ee5c86423 :: prose twin: two significant figures, direction in words
\newcommand{\MlstSculptCoarseMassOverCoarseTenPrMax}{equal within 30\%}  % experiments/FIGS/data/v2/index.json :: distance at QCD efficiency 0.1: R16_Q1_MASS over R16_Q1, run by run :: sha256 c6c4138ee5c86423 :: prose twin of the run farther from 1
\newcommand{\MlstSculptCoarseOverFullTen}{1.04\,\ensuremath{\pm}\,0.05}  % experiments/FIGS/data/v2/index.json :: distance at QCD efficiency 0.1: R16_Q1 over L188, run by run :: sha256 c6c4138ee5c86423
\newcommand{\MlstSculptCoarseOverFullTenPr}{equal within 6\%}  % experiments/FIGS/data/v2/index.json :: distance at QCD efficiency 0.1: R16_Q1 over L188, run by run :: sha256 c6c4138ee5c86423 :: prose twin: two significant figures, direction in words
\newcommand{\MlstSculptResOverFullTen}{1.02\,\ensuremath{\pm}\,0.06}  % experiments/FIGS/data/v2/index.json :: distance at QCD efficiency 0.1: L162 over L188, run by run :: sha256 c6c4138ee5c86423
\newcommand{\MlstSculptResOverFullTenPr}{equal within 6\%}  % experiments/FIGS/data/v2/index.json :: distance at QCD efficiency 0.1: L162 over L188, run by run :: sha256 c6c4138ee5c86423 :: prose twin: two significant figures, direction in words
\newcommand{\MlstSculptMidOverFullTen}{1.069\,\ensuremath{\pm}\,0.020}  % experiments/FIGS/data/v2/index.json :: distance at QCD efficiency 0.1: R42_Q1 over L188, run by run :: sha256 c6c4138ee5c86423
\newcommand{\MlstSculptMidOverFullTenPr}{1.1 times the distance}  % experiments/FIGS/data/v2/index.json :: distance at QCD efficiency 0.1: R42_Q1 over L188, run by run :: sha256 c6c4138ee5c86423 :: prose twin: two significant figures, direction in words
\newcommand{\MlstSculptCoarseOverMidTen}{0.971\,\ensuremath{\pm}\,0.035}  % experiments/FIGS/data/v2/index.json :: distance at QCD efficiency 0.1: R16_Q1 over R42_Q1, run by run :: sha256 c6c4138ee5c86423
\newcommand{\MlstSculptCoarseOverMidTenPr}{equal within 4\%}  % experiments/FIGS/data/v2/index.json :: distance at QCD efficiency 0.1: R16_Q1 over R42_Q1, run by run :: sha256 c6c4138ee5c86423 :: prose twin: two significant figures, direction in words
\newcommand{\MlstSculptNPassOne}{4{,}103}  % experiments/FIGS/data/v2/index.json :: QCD test jets passing the cut at QCD efficiency 0.01, every run :: sha256 c6c4138ee5c86423
\newcommand{\MlstSculptFloorOne}{0.0313}  % experiments/FIGS/data/v2/index.json :: Jensen-Shannon distance expected for a random selection of as many QCD jets as pass at 0.01 :: sha256 c6c4138ee5c86423
\newcommand{\MlstSculptFullOne}{0.0439\,\ensuremath{\pm}\,0.0024}  % experiments/FIGS/data/v2/index.json :: Jensen-Shannon distance of the QCD m_SD distribution at QCD efficiency 0.01, L188 :: sha256 c6c4138ee5c86423
\newcommand{\MlstSculptFullOverFloorOne}{1.40\,\ensuremath{\pm}\,0.08}  % experiments/FIGS/data/v2/index.json :: distance at 0.01 over the random-selection floor, L188 :: sha256 c6c4138ee5c86423
\newcommand{\MlstSculptFullOverFloorOnePr}{1.4 times the distance}  % experiments/FIGS/data/v2/index.json :: distance at 0.01 over the random-selection floor, L188 :: sha256 c6c4138ee5c86423 :: prose twin: two significant figures, direction in words
\newcommand{\MlstSculptResOne}{0.0429\,\ensuremath{\pm}\,0.0029}  % experiments/FIGS/data/v2/index.json :: Jensen-Shannon distance of the QCD m_SD distribution at QCD efficiency 0.01, L162 :: sha256 c6c4138ee5c86423
\newcommand{\MlstSculptResOverFloorOne}{1.37\,\ensuremath{\pm}\,0.09}  % experiments/FIGS/data/v2/index.json :: distance at 0.01 over the random-selection floor, L162 :: sha256 c6c4138ee5c86423
\newcommand{\MlstSculptResOverFloorOnePr}{1.4 times the distance}  % experiments/FIGS/data/v2/index.json :: distance at 0.01 over the random-selection floor, L162 :: sha256 c6c4138ee5c86423 :: prose twin: two significant figures, direction in words
\newcommand{\MlstSculptMidOne}{0.051\,\ensuremath{\pm}\,0.007}  % experiments/FIGS/data/v2/index.json :: Jensen-Shannon distance of the QCD m_SD distribution at QCD efficiency 0.01, R42_Q1 :: sha256 c6c4138ee5c86423
\newcommand{\MlstSculptMidOverFloorOne}{1.62\,\ensuremath{\pm}\,0.23}  % experiments/FIGS/data/v2/index.json :: distance at 0.01 over the random-selection floor, R42_Q1 :: sha256 c6c4138ee5c86423
\newcommand{\MlstSculptMidOverFloorOnePr}{1.6 times the distance}  % experiments/FIGS/data/v2/index.json :: distance at 0.01 over the random-selection floor, R42_Q1 :: sha256 c6c4138ee5c86423 :: prose twin: two significant figures, direction in words
\newcommand{\MlstSculptCoarseOne}{0.0561\,\ensuremath{\pm}\,0.0017}  % experiments/FIGS/data/v2/index.json :: Jensen-Shannon distance of the QCD m_SD distribution at QCD efficiency 0.01, R16_Q1 :: sha256 c6c4138ee5c86423
\newcommand{\MlstSculptCoarseOverFloorOne}{1.79\,\ensuremath{\pm}\,0.06}  % experiments/FIGS/data/v2/index.json :: distance at 0.01 over the random-selection floor, R16_Q1 :: sha256 c6c4138ee5c86423
\newcommand{\MlstSculptCoarseOverFloorOnePr}{1.8 times the distance}  % experiments/FIGS/data/v2/index.json :: distance at 0.01 over the random-selection floor, R16_Q1 :: sha256 c6c4138ee5c86423 :: prose twin: two significant figures, direction in words
\newcommand{\MlstSculptResMassOne}{0.0424 and 0.0442}  % experiments/FIGS/data/v2/index.json :: Jensen-Shannon distance of the QCD m_SD distribution at QCD efficiency 0.01, L162_MASS :: sha256 c6c4138ee5c86423
\newcommand{\MlstSculptResMassOverFloorOne}{1.356 and 1.411}  % experiments/FIGS/data/v2/index.json :: distance at 0.01 over the random-selection floor, L162_MASS :: sha256 c6c4138ee5c86423
\newcommand{\MlstSculptResMassOverFloorOnePr}{at least 1.4 times the distance}  % experiments/FIGS/data/v2/index.json :: distance at 0.01 over the random-selection floor, L162_MASS :: sha256 c6c4138ee5c86423 :: prose twin: two significant figures, direction in words
\newcommand{\MlstSculptResMassOverFloorOnePrMax}{up to 1.4 times the distance}  % experiments/FIGS/data/v2/index.json :: distance at 0.01 over the random-selection floor, L162_MASS :: sha256 c6c4138ee5c86423 :: prose twin of the run farther from 1
\newcommand{\MlstSculptCoarseMassOne}{0.0405 and 0.0452}  % experiments/FIGS/data/v2/index.json :: Jensen-Shannon distance of the QCD m_SD distribution at QCD efficiency 0.01, R16_Q1_MASS :: sha256 c6c4138ee5c86423
\newcommand{\MlstSculptCoarseMassOverFloorOne}{1.29 and 1.44}  % experiments/FIGS/data/v2/index.json :: distance at 0.01 over the random-selection floor, R16_Q1_MASS :: sha256 c6c4138ee5c86423
\newcommand{\MlstSculptCoarseMassOverFloorOnePr}{at least 1.3 times the distance}  % experiments/FIGS/data/v2/index.json :: distance at 0.01 over the random-selection floor, R16_Q1_MASS :: sha256 c6c4138ee5c86423 :: prose twin: two significant figures, direction in words
\newcommand{\MlstSculptCoarseMassOverFloorOnePrMax}{up to 1.4 times the distance}  % experiments/FIGS/data/v2/index.json :: distance at 0.01 over the random-selection floor, R16_Q1_MASS :: sha256 c6c4138ee5c86423 :: prose twin of the run farther from 1
\newcommand{\MlstSculptResMassOverResOne}{0.962 and 0.983}  % experiments/FIGS/data/v2/index.json :: distance at QCD efficiency 0.01: L162_MASS over L162, run by run :: sha256 c6c4138ee5c86423
\newcommand{\MlstSculptResMassOverResOnePr}{equal within 4\%}  % experiments/FIGS/data/v2/index.json :: distance at QCD efficiency 0.01: L162_MASS over L162, run by run :: sha256 c6c4138ee5c86423 :: prose twin: two significant figures, direction in words
\newcommand{\MlstSculptResMassOverResOnePrMax}{equal within 4\%}  % experiments/FIGS/data/v2/index.json :: distance at QCD efficiency 0.01: L162_MASS over L162, run by run :: sha256 c6c4138ee5c86423 :: prose twin of the run farther from 1
\newcommand{\MlstSculptCoarseMassOverCoarseOne}{0.70 and 0.82}  % experiments/FIGS/data/v2/index.json :: distance at QCD efficiency 0.01: R16_Q1_MASS over R16_Q1, run by run :: sha256 c6c4138ee5c86423
\newcommand{\MlstSculptCoarseMassOverCoarseOnePr}{at most 0.82 times the distance}  % experiments/FIGS/data/v2/index.json :: distance at QCD efficiency 0.01: R16_Q1_MASS over R16_Q1, run by run :: sha256 c6c4138ee5c86423 :: prose twin: two significant figures, direction in words
\newcommand{\MlstSculptCoarseMassOverCoarseOnePrMax}{down to 0.70 times the distance}  % experiments/FIGS/data/v2/index.json :: distance at QCD efficiency 0.01: R16_Q1_MASS over R16_Q1, run by run :: sha256 c6c4138ee5c86423 :: prose twin of the run farther from 1
\newcommand{\MlstSculptCoarseOverFullOne}{1.28\,\ensuremath{\pm}\,0.04}  % experiments/FIGS/data/v2/index.json :: distance at QCD efficiency 0.01: R16_Q1 over L188, run by run :: sha256 c6c4138ee5c86423
\newcommand{\MlstSculptCoarseOverFullOnePr}{1.3 times the distance}  % experiments/FIGS/data/v2/index.json :: distance at QCD efficiency 0.01: R16_Q1 over L188, run by run :: sha256 c6c4138ee5c86423 :: prose twin: two significant figures, direction in words
\newcommand{\MlstSculptResOverFullOne}{0.98\,\ensuremath{\pm}\,0.05}  % experiments/FIGS/data/v2/index.json :: distance at QCD efficiency 0.01: L162 over L188, run by run :: sha256 c6c4138ee5c86423
\newcommand{\MlstSculptResOverFullOnePr}{equal within 5\%}  % experiments/FIGS/data/v2/index.json :: distance at QCD efficiency 0.01: L162 over L188, run by run :: sha256 c6c4138ee5c86423 :: prose twin: two significant figures, direction in words
\newcommand{\MlstSculptMidOverFullOne}{1.15\,\ensuremath{\pm}\,0.12}  % experiments/FIGS/data/v2/index.json :: distance at QCD efficiency 0.01: R42_Q1 over L188, run by run :: sha256 c6c4138ee5c86423
\newcommand{\MlstSculptMidOverFullOnePr}{equal within 20\%}  % experiments/FIGS/data/v2/index.json :: distance at QCD efficiency 0.01: R42_Q1 over L188, run by run :: sha256 c6c4138ee5c86423 :: prose twin: two significant figures, direction in words
\newcommand{\MlstSculptCoarseOverMidOne}{1.13\,\ensuremath{\pm}\,0.16}  % experiments/FIGS/data/v2/index.json :: distance at QCD efficiency 0.01: R16_Q1 over R42_Q1, run by run :: sha256 c6c4138ee5c86423
\newcommand{\MlstSculptCoarseOverMidOnePr}{equal within 20\%}  % experiments/FIGS/data/v2/index.json :: distance at QCD efficiency 0.01: R16_Q1 over R42_Q1, run by run :: sha256 c6c4138ee5c86423 :: prose twin: two significant figures, direction in words
\newcommand{\MlstSculptDiffEffLo}{0.52\%}  % experiments/FIGS/data/v2/index.json :: lowest stored QCD efficiency of the unbroken stretch around the tighter working point at which the 17/188 distance ratio exceeds 1 by twice its SD :: sha256 c6c4138ee5c86423 :: passing grid points: 0.001, 0.005245, 0.007937, 0.012011
\newcommand{\MlstSculptDiffEffHi}{1.2\%}  % experiments/FIGS/data/v2/index.json :: highest efficiency of that stretch :: sha256 c6c4138ee5c86423
\newcommand{\MlstSculptLowestEff}{0.10\%}  % experiments/FIGS/data/v2/index.json :: lowest stored QCD efficiency of the sculpting curves :: sha256 c6c4138ee5c86423
\newcommand{\MlstSculptLowestNPass}{410}  % experiments/FIGS/data/v2/index.json :: QCD test jets passing at the lowest stored efficiency :: sha256 c6c4138ee5c86423
\newcommand{\MlstSculptLowestMaxFloorSd}{1.7}  % experiments/FIGS/data/v2/index.json :: largest distance of a configuration's mean from the random-selection floor at the lowest stored efficiency, in SDs of the floor over random selections :: sha256 c6c4138ee5c86423
\newcommand{\MlstSculptNConfigs}{6}  % experiments/FIGS/data/v2/index.json :: configurations with a sculpting curve (count) :: sha256 c6c4138ee5c86423
\newcommand{\LpJciiTestQcdPercent}{20.5\%}  % experiments/FIGS/data/v2/jc2_test/counts.json :: QCD share of the jets, per cent :: sha256 3574ba1f6f2a59f0
\newcommand{\LpJciiTestTwoProngPercent}{38.7\%}  % experiments/FIGS/data/v2/jc2_test/counts.json :: two-prong share of the jets, per cent :: sha256 3574ba1f6f2a59f0
\newcommand{\LpJciiTrainQcdPercent}{13.5\%}  % experiments/FT/data/jc2_v2_manifest.json :: QCD share of the jets, per cent :: sha256 4bd68655919b14e9
\newcommand{\LpJciiTrainTwoProngPercent}{15.4\%}  % experiments/FT/data/jc2_v2_manifest.json :: two-prong share of the jets, per cent :: sha256 4bd68655919b14e9
\newcommand{\LpJciiTestMinAucJets}{10}  % experiments/FIGS/data/v2/jc2_test/counts.json :: fewest jets of a class among the macro-AUC rows :: sha256 3574ba1f6f2a59f0
\newcommand{\LpJciiTrainMinJets}{11}  % experiments/FT/data/jc2_v2_manifest.json :: fewest jets of a class in the 10^6 training subset :: sha256 4bd68655919b14e9
\newcommand{\LpNTestJci}{2{,}000{,}000}  % experiments/FIGS/data/v2/index.json :: jc1 test jets the probes score :: sha256 c6c4138ee5c86423
\newcommand{\LpNTrainMaxJci}{1{,}000{,}000}  % experiments/FIGS/data/v2/index.json :: jc1 largest training subset :: sha256 c6c4138ee5c86423
\newcommand{\LpNTestTop}{404{,}000}  % experiments/FIGS/data/v2/index.json :: top test jets the probes score :: sha256 c6c4138ee5c86423
\newcommand{\LpNTrainMaxTop}{1{,}200{,}000}  % experiments/FIGS/data/v2/index.json :: top largest training subset :: sha256 c6c4138ee5c86423
\newcommand{\LpNTestQg}{200{,}000}  % experiments/FIGS/data/v2/index.json :: qg test jets the probes score :: sha256 c6c4138ee5c86423
\newcommand{\LpNTrainMaxQg}{1{,}600{,}000}  % experiments/FIGS/data/v2/index.json :: qg largest training subset :: sha256 c6c4138ee5c86423
\newcommand{\LpNTestAucJci}{2{,}000{,}000}  % experiments/FIGS/data/v2/index.json :: JetClass test jets of the macro AUC :: sha256 c6c4138ee5c86423
\newcommand{\LpPairNValMin}{1{,}068}  % experiments/FIGS/data/v2/index.json :: validation jets of the pair probe with the fewest :: sha256 c6c4138ee5c86423
\newcommand{\LpPairNValMax}{20{,}195}  % experiments/FIGS/data/v2/index.json :: validation jets of the pair probe with the most :: sha256 c6c4138ee5c86423
\newcommand{\LpNValJci}{200{,}000}  % experiments/FIGS/data/v2/index.json :: JetClass probe validation jets :: sha256 c6c4138ee5c86423
\newcommand{\SculptEffHigh}{10\%}  % experiments/EVAL/mass_sculpting.py :: WORKING_POINTS :: sha256 e85d125c85dac054 :: QCD efficiency of a sculpting cut
\newcommand{\SculptEffLow}{1\%}  % experiments/EVAL/mass_sculpting.py :: WORKING_POINTS :: sha256 e85d125c85dac054 :: QCD efficiency of a sculpting cut
\newcommand{\LpMlpHidden}{256}  % experiments/EVAL/probe.py :: fit_mlp hidden layer :: sha256 485db0966e660588 :: units of the MLP probe's hidden layer
\newcommand{\LpMlpSeedsWord}{three}  % experiments/EVAL/probe.py :: MLP_SEEDS :: sha256 485db0966e660588 :: seeds of the MLP probe whose scores are averaged
\newcommand{\MassMinTrainPerClass}{50}  % experiments/EVAL/mass_resolution.py :: MIN_TRAIN_PER_CLASS :: sha256 26acb0e315094bd1 :: fewest training jets of a class the mass regression keeps
\newcommand{\MlstFtLr}{\ensuremath{10^{-4}}}  % experiments/FIGS/data/v2/ft_mlst/manifests.json :: cells[best70].lr :: sha256 a689f21e894f48ae
\newcommand{\MlstFtHeadMult}{50}  % experiments/FIGS/data/v2/ft_mlst/manifests.json :: cells[best70].head_lr_mult :: sha256 a689f21e894f48ae
\newcommand{\MlstFtWeightDecay}{0.01}  % experiments/FIGS/data/v2/ft_mlst/manifests.json :: cells[best70].weight_decay :: sha256 a689f21e894f48ae
\newcommand{\MlstFtBatch}{512}  % experiments/FIGS/data/v2/ft_mlst/manifests.json :: cells[best70].batch_size :: sha256 a689f21e894f48ae
\newcommand{\MlstFtBenchEpochs}{20}  % experiments/FIGS/data/v2/ft_mlst/manifests.json :: cells[best70, top/qg].epochs :: sha256 a689f21e894f48ae
\newcommand{\MlstFtEpochsEThree}{50}  % experiments/FIGS/data/v2/ft_mlst/manifests.json :: cells[best70, leg 1/2, n_train 1000].epochs :: sha256 a689f21e894f48ae
\newcommand{\MlstFtEpochsEFour}{50}  % experiments/FIGS/data/v2/ft_mlst/manifests.json :: cells[best70, leg 1/2, n_train 10000].epochs :: sha256 a689f21e894f48ae
\newcommand{\MlstFtEpochsEFive}{30}  % experiments/FIGS/data/v2/ft_mlst/manifests.json :: cells[best70, leg 1/2, n_train 100000].epochs :: sha256 a689f21e894f48ae
\newcommand{\MlstFtEpochsESix}{10}  % experiments/FIGS/data/v2/ft_mlst/manifests.json :: cells[best70, leg 1/2, n_train 1000000].epochs :: sha256 a689f21e894f48ae
\newcommand{\MlstFtHeadDecayShare}{30\%}  % experiments/FT/k8s/job-ft-v2-legs-best70-l188-s3-raunav.yaml :: no --lr-scheduler: weaver 0.4.17 flat+decay, the head's last epochs :: sha256 19f59aa54a385c11
\newcommand{\MlstFtHeadDecayTo}{1\%}  % experiments/FT/k8s/job-ft-v2-legs-best70-l188-s3-raunav.yaml :: no --lr-scheduler: weaver 0.4.17 flat+decay, final factor :: sha256 19f59aa54a385c11
\newcommand{\MlstFtHeadDropout}{0.1}  % experiments/FT/k8s/job-ft-v2-legs-best70-l188-s3-raunav.yaml :: fc_params dropout of the fine-tuning head :: sha256 19f59aa54a385c11
\newcommand{\MlstFtMinEpochJets}{10{,}000}  % experiments/FIGS/data/v2/ft_mlst/manifests.json :: cells[best70].samples_per_epoch (smallest) :: sha256 a689f21e894f48ae
\newcommand{\MlstFtValJcii}{20{,}480}  % experiments/FIGS/data/v2/ft_mlst/manifests.json :: cells[best70, leg 1].val_jets :: sha256 a689f21e894f48ae
\newcommand{\MlstFtValSmall}{20{,}000}  % experiments/FT/k8s/job-ft-v2-legs-best70-l188-s3-raunav.yaml :: JetClass --samples-per-epoch-val; top and q/g at N <= 1e4 :: sha256 19f59aa54a385c11
\newcommand{\MlstFtValLarge}{200{,}000}  % experiments/FT/k8s/job-ft-v2-bench-best70-l188-s1-top-raunav.yaml :: top and q/g --samples-per-epoch-val at N > 1e4 :: sha256 1e38d4f542119ea3
\newcommand{\MlstFtGpuTypes}{V100}  % experiments/FIGS/data/v2/ft_mlst/manifests.json :: cells[best70].gpu_device_name :: sha256 a689f21e894f48ae
\newcommand{\MlstFtCellsDone}{445}  % experiments/FIGS/data/v2/ft_mlst/manifests.json :: DONE paper-run fine-tuning cells (count) :: sha256 a689f21e894f48ae
\newcommand{\MlstFtRefLr}{\ensuremath{10^{-3}}}  % experiments/FIGS/data/v2/ft_mlst/manifests.json :: cells[ref_v100].lr :: sha256 a689f21e894f48ae
\newcommand{\MlstFtRefSeedsWord}{three}  % experiments/FIGS/data/v2/ft_mlst/manifests.json :: cells[ref_v100].ft_seed (count) :: sha256 a689f21e894f48ae
\newcommand{\MlstFtRefGpuTypes}{V100}  % experiments/FIGS/data/v2/ft_mlst/manifests.json :: cells[ref_v100].gpu_device_name :: sha256 a689f21e894f48ae
\newcommand{\MlstFtRefWeightDecayBench}{0.01}  % experiments/FIGS/data/v2/ft_mlst/manifests.json :: cells[ref_v100, top/qg].weight_decay :: sha256 a689f21e894f48ae
\newcommand{\MlstBootResamples}{200}  % experiments/FIGS/data/v2/index.json :: bootstrap resamples of the test jets :: sha256 c6c4138ee5c86423
\newcommand{\MlstBootTestMaxProbe}{0.21}  % experiments/FIGS/data/v2/index.json :: largest test-sample SD of the ln ratio of 1 - AUC on JetClass-II, linear probe, label sets over the 188-class model, per cent :: sha256 c6c4138ee5c86423
\newcommand{\MlstBootTestMaxFt}{0.15}  % experiments/FIGS/data/v2/index.json :: largest test-sample SD of the ln ratio of 1 - AUC on JetClass-II, fine-tuning, label sets over the 188-class model, per cent :: sha256 c6c4138ee5c86423
\newcommand{\MlstBootTestMaxRthirtyProbe}{3.6}  % experiments/FIGS/data/v2/index.json :: largest test-sample SD of the ln R30 ratio on top and quark/gluon, linear probe (counting), label sets over the 188-class model, per cent :: sha256 c6c4138ee5c86423
\newcommand{\MlstBootTestMaxRthirtyFt}{13}  % experiments/FIGS/data/v2/index.json :: largest test-sample SD of the ln R30 ratio on top and quark/gluon, fine-tuning (counting), label sets over the 188-class model, per cent :: sha256 c6c4138ee5c86423
\newcommand{\MlstFtCmpJcTwoResOverFullEFour}{0.989\,\ensuremath{\pm}\,0.011}  % experiments/FIGS/data/v2/index.json :: fine-tuning jc2 N=10000: (1 - macro_auc_ovr) of L162 over L188, run by run :: sha256 c6c4138ee5c86423
\newcommand{\MlstFtCmpJcTwoResOverFullEFourPr}{equal within 2\%}  % experiments/FIGS/data/v2/index.json :: fine-tuning jc2 N=10000: (1 - macro_auc_ovr) of L162 over L188, run by run :: sha256 c6c4138ee5c86423 :: prose twin: two significant figures, direction in words
\newcommand{\MlstFtCmpJcTwoMidOverFullEFour}{1.191\,\ensuremath{\pm}\,0.030}  % experiments/FIGS/data/v2/index.json :: fine-tuning jc2 N=10000: (1 - macro_auc_ovr) of R42_Q1 over L188, run by run :: sha256 c6c4138ee5c86423
\newcommand{\MlstFtCmpJcTwoMidOverFullEFourPr}{19\% more errors}  % experiments/FIGS/data/v2/index.json :: fine-tuning jc2 N=10000: (1 - macro_auc_ovr) of R42_Q1 over L188, run by run :: sha256 c6c4138ee5c86423 :: prose twin: two significant figures, direction in words
\newcommand{\MlstFtCmpJcTwoCoarseOverFullEFour}{1.353\,\ensuremath{\pm}\,0.003}  % experiments/FIGS/data/v2/index.json :: fine-tuning jc2 N=10000: (1 - macro_auc_ovr) of R16_Q1 over L188, run by run :: sha256 c6c4138ee5c86423
\newcommand{\MlstFtCmpJcTwoCoarseOverFullEFourPr}{35\% more errors}  % experiments/FIGS/data/v2/index.json :: fine-tuning jc2 N=10000: (1 - macro_auc_ovr) of R16_Q1 over L188, run by run :: sha256 c6c4138ee5c86423 :: prose twin: two significant figures, direction in words
\newcommand{\MlstFtCmpJcTwoResMassOverResEFour}{1.016 and 1.005}  % experiments/FIGS/data/v2/index.json :: fine-tuning jc2 N=10000: (1 - macro_auc_ovr) of L162_MASS over L162, run by run :: sha256 c6c4138ee5c86423
\newcommand{\MlstFtCmpJcTwoResMassOverResEFourPr}{equal within 2\%}  % experiments/FIGS/data/v2/index.json :: fine-tuning jc2 N=10000: (1 - macro_auc_ovr) of L162_MASS over L162, run by run :: sha256 c6c4138ee5c86423 :: prose twin: two significant figures, direction in words
\newcommand{\MlstFtCmpJcTwoResMassOverResEFourPrMax}{equal within 2\%}  % experiments/FIGS/data/v2/index.json :: fine-tuning jc2 N=10000: (1 - macro_auc_ovr) of L162_MASS over L162, run by run :: sha256 c6c4138ee5c86423 :: prose twin of the run farther from 1
\newcommand{\MlstFtCmpJcTwoCoarseMassOverCoarseEFour}{0.989 and 1.000}  % experiments/FIGS/data/v2/index.json :: fine-tuning jc2 N=10000: (1 - macro_auc_ovr) of R16_Q1_MASS over R16_Q1, run by run :: sha256 c6c4138ee5c86423
\newcommand{\MlstFtCmpJcTwoCoarseMassOverCoarseEFourPr}{equal within 2\%}  % experiments/FIGS/data/v2/index.json :: fine-tuning jc2 N=10000: (1 - macro_auc_ovr) of R16_Q1_MASS over R16_Q1, run by run :: sha256 c6c4138ee5c86423 :: prose twin: two significant figures, direction in words
\newcommand{\MlstFtCmpJcTwoCoarseMassOverCoarseEFourPrMax}{equal within 2\%}  % experiments/FIGS/data/v2/index.json :: fine-tuning jc2 N=10000: (1 - macro_auc_ovr) of R16_Q1_MASS over R16_Q1, run by run :: sha256 c6c4138ee5c86423 :: prose twin of the run farther from 1
\newcommand{\MlstFtCmpJcTwoResOverFullEFive}{1.011\,\ensuremath{\pm}\,0.005}  % experiments/FIGS/data/v2/index.json :: fine-tuning jc2 N=100000: (1 - macro_auc_ovr) of L162 over L188, run by run :: sha256 c6c4138ee5c86423
\newcommand{\MlstFtCmpJcTwoResOverFullEFivePr}{1.1\% more errors}  % experiments/FIGS/data/v2/index.json :: fine-tuning jc2 N=100000: (1 - macro_auc_ovr) of L162 over L188, run by run :: sha256 c6c4138ee5c86423 :: prose twin: two significant figures, direction in words
\newcommand{\MlstFtCmpJcTwoMidOverFullEFive}{1.064\,\ensuremath{\pm}\,0.009}  % experiments/FIGS/data/v2/index.json :: fine-tuning jc2 N=100000: (1 - macro_auc_ovr) of R42_Q1 over L188, run by run :: sha256 c6c4138ee5c86423
\newcommand{\MlstFtCmpJcTwoMidOverFullEFivePr}{6.4\% more errors}  % experiments/FIGS/data/v2/index.json :: fine-tuning jc2 N=100000: (1 - macro_auc_ovr) of R42_Q1 over L188, run by run :: sha256 c6c4138ee5c86423 :: prose twin: two significant figures, direction in words
\newcommand{\MlstFtCmpJcTwoCoarseOverFullEFive}{1.157\,\ensuremath{\pm}\,0.001}  % experiments/FIGS/data/v2/index.json :: fine-tuning jc2 N=100000: (1 - macro_auc_ovr) of R16_Q1 over L188, run by run :: sha256 c6c4138ee5c86423
\newcommand{\MlstFtCmpJcTwoCoarseOverFullEFivePr}{16\% more errors}  % experiments/FIGS/data/v2/index.json :: fine-tuning jc2 N=100000: (1 - macro_auc_ovr) of R16_Q1 over L188, run by run :: sha256 c6c4138ee5c86423 :: prose twin: two significant figures, direction in words
\newcommand{\MlstFtCmpJcTwoResMassOverResEFive}{1.027 and 1.021}  % experiments/FIGS/data/v2/index.json :: fine-tuning jc2 N=100000: (1 - macro_auc_ovr) of L162_MASS over L162, run by run :: sha256 c6c4138ee5c86423
\newcommand{\MlstFtCmpJcTwoResMassOverResEFivePr}{at least 2.1\% more errors}  % experiments/FIGS/data/v2/index.json :: fine-tuning jc2 N=100000: (1 - macro_auc_ovr) of L162_MASS over L162, run by run :: sha256 c6c4138ee5c86423 :: prose twin: two significant figures, direction in words
\newcommand{\MlstFtCmpJcTwoResMassOverResEFivePrMax}{up to 2.7\% more errors}  % experiments/FIGS/data/v2/index.json :: fine-tuning jc2 N=100000: (1 - macro_auc_ovr) of L162_MASS over L162, run by run :: sha256 c6c4138ee5c86423 :: prose twin of the run farther from 1
\newcommand{\MlstFtCmpJcTwoCoarseMassOverCoarseEFive}{1.029 and 1.022}  % experiments/FIGS/data/v2/index.json :: fine-tuning jc2 N=100000: (1 - macro_auc_ovr) of R16_Q1_MASS over R16_Q1, run by run :: sha256 c6c4138ee5c86423
\newcommand{\MlstFtCmpJcTwoCoarseMassOverCoarseEFivePr}{at least 2.2\% more errors}  % experiments/FIGS/data/v2/index.json :: fine-tuning jc2 N=100000: (1 - macro_auc_ovr) of R16_Q1_MASS over R16_Q1, run by run :: sha256 c6c4138ee5c86423 :: prose twin: two significant figures, direction in words
\newcommand{\MlstFtCmpJcTwoCoarseMassOverCoarseEFivePrMax}{up to 2.9\% more errors}  % experiments/FIGS/data/v2/index.json :: fine-tuning jc2 N=100000: (1 - macro_auc_ovr) of R16_Q1_MASS over R16_Q1, run by run :: sha256 c6c4138ee5c86423 :: prose twin of the run farther from 1
\newcommand{\MlstFtCmpJcTwoResOverFullESix}{1.014\,\ensuremath{\pm}\,0.011}  % experiments/FIGS/data/v2/index.json :: fine-tuning jc2 N=1000000: (1 - macro_auc_ovr) of L162 over L188, run by run :: sha256 c6c4138ee5c86423
\newcommand{\MlstFtCmpJcTwoResOverFullESixPr}{equal within 2\%}  % experiments/FIGS/data/v2/index.json :: fine-tuning jc2 N=1000000: (1 - macro_auc_ovr) of L162 over L188, run by run :: sha256 c6c4138ee5c86423 :: prose twin: two significant figures, direction in words
\newcommand{\MlstFtCmpJcTwoMidOverFullESix}{1.043\,\ensuremath{\pm}\,0.002}  % experiments/FIGS/data/v2/index.json :: fine-tuning jc2 N=1000000: (1 - macro_auc_ovr) of R42_Q1 over L188, run by run :: sha256 c6c4138ee5c86423
\newcommand{\MlstFtCmpJcTwoMidOverFullESixPr}{4.3\% more errors}  % experiments/FIGS/data/v2/index.json :: fine-tuning jc2 N=1000000: (1 - macro_auc_ovr) of R42_Q1 over L188, run by run :: sha256 c6c4138ee5c86423 :: prose twin: two significant figures, direction in words
\newcommand{\MlstFtCmpJcTwoCoarseOverFullESix}{1.123\,\ensuremath{\pm}\,0.003}  % experiments/FIGS/data/v2/index.json :: fine-tuning jc2 N=1000000: (1 - macro_auc_ovr) of R16_Q1 over L188, run by run :: sha256 c6c4138ee5c86423
\newcommand{\MlstFtCmpJcTwoCoarseOverFullESixPr}{12\% more errors}  % experiments/FIGS/data/v2/index.json :: fine-tuning jc2 N=1000000: (1 - macro_auc_ovr) of R16_Q1 over L188, run by run :: sha256 c6c4138ee5c86423 :: prose twin: two significant figures, direction in words
\newcommand{\MlstFtCmpJcTwoResMassOverResESix}{1.009 and 1.016}  % experiments/FIGS/data/v2/index.json :: fine-tuning jc2 N=1000000: (1 - macro_auc_ovr) of L162_MASS over L162, run by run :: sha256 c6c4138ee5c86423
\newcommand{\MlstFtCmpJcTwoResMassOverResESixPr}{equal within 2\%}  % experiments/FIGS/data/v2/index.json :: fine-tuning jc2 N=1000000: (1 - macro_auc_ovr) of L162_MASS over L162, run by run :: sha256 c6c4138ee5c86423 :: prose twin: two significant figures, direction in words
\newcommand{\MlstFtCmpJcTwoResMassOverResESixPrMax}{equal within 2\%}  % experiments/FIGS/data/v2/index.json :: fine-tuning jc2 N=1000000: (1 - macro_auc_ovr) of L162_MASS over L162, run by run :: sha256 c6c4138ee5c86423 :: prose twin of the run farther from 1
\newcommand{\MlstFtCmpJcTwoCoarseMassOverCoarseESix}{1.023 and 1.017}  % experiments/FIGS/data/v2/index.json :: fine-tuning jc2 N=1000000: (1 - macro_auc_ovr) of R16_Q1_MASS over R16_Q1, run by run :: sha256 c6c4138ee5c86423
\newcommand{\MlstFtCmpJcTwoCoarseMassOverCoarseESixPr}{at least 1.7\% more errors}  % experiments/FIGS/data/v2/index.json :: fine-tuning jc2 N=1000000: (1 - macro_auc_ovr) of R16_Q1_MASS over R16_Q1, run by run :: sha256 c6c4138ee5c86423 :: prose twin: two significant figures, direction in words
\newcommand{\MlstFtCmpJcTwoCoarseMassOverCoarseESixPrMax}{up to 2.3\% more errors}  % experiments/FIGS/data/v2/index.json :: fine-tuning jc2 N=1000000: (1 - macro_auc_ovr) of R16_Q1_MASS over R16_Q1, run by run :: sha256 c6c4138ee5c86423 :: prose twin of the run farther from 1
\newcommand{\MlstFtCmpJcTwoQcdResOverFullEFour}{1.038\,\ensuremath{\pm}\,0.005}  % experiments/FIGS/data/v2/index.json :: fine-tuning jc2_qcd N=10000: (1 - macro_auc_qcd) of L162 over L188, run by run :: sha256 c6c4138ee5c86423
\newcommand{\MlstFtCmpJcTwoQcdResOverFullEFourPr}{3.8\% more errors}  % experiments/FIGS/data/v2/index.json :: fine-tuning jc2_qcd N=10000: (1 - macro_auc_qcd) of L162 over L188, run by run :: sha256 c6c4138ee5c86423 :: prose twin: two significant figures, direction in words
\newcommand{\MlstFtCmpJcTwoQcdMidOverFullEFour}{1.164\,\ensuremath{\pm}\,0.018}  % experiments/FIGS/data/v2/index.json :: fine-tuning jc2_qcd N=10000: (1 - macro_auc_qcd) of R42_Q1 over L188, run by run :: sha256 c6c4138ee5c86423
\newcommand{\MlstFtCmpJcTwoQcdMidOverFullEFourPr}{16\% more errors}  % experiments/FIGS/data/v2/index.json :: fine-tuning jc2_qcd N=10000: (1 - macro_auc_qcd) of R42_Q1 over L188, run by run :: sha256 c6c4138ee5c86423 :: prose twin: two significant figures, direction in words
\newcommand{\MlstFtCmpJcTwoQcdCoarseOverFullEFour}{1.258\,\ensuremath{\pm}\,0.025}  % experiments/FIGS/data/v2/index.json :: fine-tuning jc2_qcd N=10000: (1 - macro_auc_qcd) of R16_Q1 over L188, run by run :: sha256 c6c4138ee5c86423
\newcommand{\MlstFtCmpJcTwoQcdCoarseOverFullEFourPr}{26\% more errors}  % experiments/FIGS/data/v2/index.json :: fine-tuning jc2_qcd N=10000: (1 - macro_auc_qcd) of R16_Q1 over L188, run by run :: sha256 c6c4138ee5c86423 :: prose twin: two significant figures, direction in words
\newcommand{\MlstFtCmpJcTwoQcdResMassOverResEFour}{1.027 and 0.997}  % experiments/FIGS/data/v2/index.json :: fine-tuning jc2_qcd N=10000: (1 - macro_auc_qcd) of L162_MASS over L162, run by run :: sha256 c6c4138ee5c86423
\newcommand{\MlstFtCmpJcTwoQcdResMassOverResEFourPr}{equal within 3\%}  % experiments/FIGS/data/v2/index.json :: fine-tuning jc2_qcd N=10000: (1 - macro_auc_qcd) of L162_MASS over L162, run by run :: sha256 c6c4138ee5c86423 :: prose twin: two significant figures, direction in words
\newcommand{\MlstFtCmpJcTwoQcdResMassOverResEFourPrMax}{equal within 3\%}  % experiments/FIGS/data/v2/index.json :: fine-tuning jc2_qcd N=10000: (1 - macro_auc_qcd) of L162_MASS over L162, run by run :: sha256 c6c4138ee5c86423 :: prose twin of the run farther from 1
\newcommand{\MlstFtCmpJcTwoQcdCoarseMassOverCoarseEFour}{1.028 and 1.048}  % experiments/FIGS/data/v2/index.json :: fine-tuning jc2_qcd N=10000: (1 - macro_auc_qcd) of R16_Q1_MASS over R16_Q1, run by run :: sha256 c6c4138ee5c86423
\newcommand{\MlstFtCmpJcTwoQcdCoarseMassOverCoarseEFourPr}{equal within 5\%}  % experiments/FIGS/data/v2/index.json :: fine-tuning jc2_qcd N=10000: (1 - macro_auc_qcd) of R16_Q1_MASS over R16_Q1, run by run :: sha256 c6c4138ee5c86423 :: prose twin: two significant figures, direction in words
\newcommand{\MlstFtCmpJcTwoQcdCoarseMassOverCoarseEFourPrMax}{equal within 5\%}  % experiments/FIGS/data/v2/index.json :: fine-tuning jc2_qcd N=10000: (1 - macro_auc_qcd) of R16_Q1_MASS over R16_Q1, run by run :: sha256 c6c4138ee5c86423 :: prose twin of the run farther from 1
\newcommand{\MlstFtCmpJcTwoQcdResOverFullEFive}{1.060\,\ensuremath{\pm}\,0.002}  % experiments/FIGS/data/v2/index.json :: fine-tuning jc2_qcd N=100000: (1 - macro_auc_qcd) of L162 over L188, run by run :: sha256 c6c4138ee5c86423
\newcommand{\MlstFtCmpJcTwoQcdResOverFullEFivePr}{6.0\% more errors}  % experiments/FIGS/data/v2/index.json :: fine-tuning jc2_qcd N=100000: (1 - macro_auc_qcd) of L162 over L188, run by run :: sha256 c6c4138ee5c86423 :: prose twin: two significant figures, direction in words
\newcommand{\MlstFtCmpJcTwoQcdMidOverFullEFive}{1.062\,\ensuremath{\pm}\,0.010}  % experiments/FIGS/data/v2/index.json :: fine-tuning jc2_qcd N=100000: (1 - macro_auc_qcd) of R42_Q1 over L188, run by run :: sha256 c6c4138ee5c86423
\newcommand{\MlstFtCmpJcTwoQcdMidOverFullEFivePr}{6.2\% more errors}  % experiments/FIGS/data/v2/index.json :: fine-tuning jc2_qcd N=100000: (1 - macro_auc_qcd) of R42_Q1 over L188, run by run :: sha256 c6c4138ee5c86423 :: prose twin: two significant figures, direction in words
\newcommand{\MlstFtCmpJcTwoQcdCoarseOverFullEFive}{1.137\,\ensuremath{\pm}\,0.001}  % experiments/FIGS/data/v2/index.json :: fine-tuning jc2_qcd N=100000: (1 - macro_auc_qcd) of R16_Q1 over L188, run by run :: sha256 c6c4138ee5c86423
\newcommand{\MlstFtCmpJcTwoQcdCoarseOverFullEFivePr}{14\% more errors}  % experiments/FIGS/data/v2/index.json :: fine-tuning jc2_qcd N=100000: (1 - macro_auc_qcd) of R16_Q1 over L188, run by run :: sha256 c6c4138ee5c86423 :: prose twin: two significant figures, direction in words
\newcommand{\MlstFtCmpJcTwoQcdResMassOverResEFive}{1.022 and 1.009}  % experiments/FIGS/data/v2/index.json :: fine-tuning jc2_qcd N=100000: (1 - macro_auc_qcd) of L162_MASS over L162, run by run :: sha256 c6c4138ee5c86423
\newcommand{\MlstFtCmpJcTwoQcdResMassOverResEFivePr}{equal within 3\%}  % experiments/FIGS/data/v2/index.json :: fine-tuning jc2_qcd N=100000: (1 - macro_auc_qcd) of L162_MASS over L162, run by run :: sha256 c6c4138ee5c86423 :: prose twin: two significant figures, direction in words
\newcommand{\MlstFtCmpJcTwoQcdResMassOverResEFivePrMax}{equal within 3\%}  % experiments/FIGS/data/v2/index.json :: fine-tuning jc2_qcd N=100000: (1 - macro_auc_qcd) of L162_MASS over L162, run by run :: sha256 c6c4138ee5c86423 :: prose twin of the run farther from 1
\newcommand{\MlstFtCmpJcTwoQcdCoarseMassOverCoarseEFive}{1.028 and 1.027}  % experiments/FIGS/data/v2/index.json :: fine-tuning jc2_qcd N=100000: (1 - macro_auc_qcd) of R16_Q1_MASS over R16_Q1, run by run :: sha256 c6c4138ee5c86423
\newcommand{\MlstFtCmpJcTwoQcdCoarseMassOverCoarseEFivePr}{at least 2.7\% more errors}  % experiments/FIGS/data/v2/index.json :: fine-tuning jc2_qcd N=100000: (1 - macro_auc_qcd) of R16_Q1_MASS over R16_Q1, run by run :: sha256 c6c4138ee5c86423 :: prose twin: two significant figures, direction in words
\newcommand{\MlstFtCmpJcTwoQcdCoarseMassOverCoarseEFivePrMax}{up to 2.8\% more errors}  % experiments/FIGS/data/v2/index.json :: fine-tuning jc2_qcd N=100000: (1 - macro_auc_qcd) of R16_Q1_MASS over R16_Q1, run by run :: sha256 c6c4138ee5c86423 :: prose twin of the run farther from 1
\newcommand{\MlstFtCmpJcTwoQcdResOverFullESix}{1.035\,\ensuremath{\pm}\,0.010}  % experiments/FIGS/data/v2/index.json :: fine-tuning jc2_qcd N=1000000: (1 - macro_auc_qcd) of L162 over L188, run by run :: sha256 c6c4138ee5c86423
\newcommand{\MlstFtCmpJcTwoQcdResOverFullESixPr}{3.5\% more errors}  % experiments/FIGS/data/v2/index.json :: fine-tuning jc2_qcd N=1000000: (1 - macro_auc_qcd) of L162 over L188, run by run :: sha256 c6c4138ee5c86423 :: prose twin: two significant figures, direction in words
\newcommand{\MlstFtCmpJcTwoQcdMidOverFullESix}{1.050\,\ensuremath{\pm}\,0.004}  % experiments/FIGS/data/v2/index.json :: fine-tuning jc2_qcd N=1000000: (1 - macro_auc_qcd) of R42_Q1 over L188, run by run :: sha256 c6c4138ee5c86423
\newcommand{\MlstFtCmpJcTwoQcdMidOverFullESixPr}{5.0\% more errors}  % experiments/FIGS/data/v2/index.json :: fine-tuning jc2_qcd N=1000000: (1 - macro_auc_qcd) of R42_Q1 over L188, run by run :: sha256 c6c4138ee5c86423 :: prose twin: two significant figures, direction in words
\newcommand{\MlstFtCmpJcTwoQcdCoarseOverFullESix}{1.122\,\ensuremath{\pm}\,0.002}  % experiments/FIGS/data/v2/index.json :: fine-tuning jc2_qcd N=1000000: (1 - macro_auc_qcd) of R16_Q1 over L188, run by run :: sha256 c6c4138ee5c86423
\newcommand{\MlstFtCmpJcTwoQcdCoarseOverFullESixPr}{12\% more errors}  % experiments/FIGS/data/v2/index.json :: fine-tuning jc2_qcd N=1000000: (1 - macro_auc_qcd) of R16_Q1 over L188, run by run :: sha256 c6c4138ee5c86423 :: prose twin: two significant figures, direction in words
\newcommand{\MlstFtCmpJcTwoQcdResMassOverResESix}{1.011 and 1.025}  % experiments/FIGS/data/v2/index.json :: fine-tuning jc2_qcd N=1000000: (1 - macro_auc_qcd) of L162_MASS over L162, run by run :: sha256 c6c4138ee5c86423
\newcommand{\MlstFtCmpJcTwoQcdResMassOverResESixPr}{equal within 3\%}  % experiments/FIGS/data/v2/index.json :: fine-tuning jc2_qcd N=1000000: (1 - macro_auc_qcd) of L162_MASS over L162, run by run :: sha256 c6c4138ee5c86423 :: prose twin: two significant figures, direction in words
\newcommand{\MlstFtCmpJcTwoQcdResMassOverResESixPrMax}{equal within 3\%}  % experiments/FIGS/data/v2/index.json :: fine-tuning jc2_qcd N=1000000: (1 - macro_auc_qcd) of L162_MASS over L162, run by run :: sha256 c6c4138ee5c86423 :: prose twin of the run farther from 1
\newcommand{\MlstFtCmpJcTwoQcdCoarseMassOverCoarseESix}{1.015 and 1.015}  % experiments/FIGS/data/v2/index.json :: fine-tuning jc2_qcd N=1000000: (1 - macro_auc_qcd) of R16_Q1_MASS over R16_Q1, run by run :: sha256 c6c4138ee5c86423
\newcommand{\MlstFtCmpJcTwoQcdCoarseMassOverCoarseESixPr}{at least 1.5\% more errors}  % experiments/FIGS/data/v2/index.json :: fine-tuning jc2_qcd N=1000000: (1 - macro_auc_qcd) of R16_Q1_MASS over R16_Q1, run by run :: sha256 c6c4138ee5c86423 :: prose twin: two significant figures, direction in words
\newcommand{\MlstFtCmpJcTwoQcdCoarseMassOverCoarseESixPrMax}{up to 1.5\% more errors}  % experiments/FIGS/data/v2/index.json :: fine-tuning jc2_qcd N=1000000: (1 - macro_auc_qcd) of R16_Q1_MASS over R16_Q1, run by run :: sha256 c6c4138ee5c86423 :: prose twin of the run farther from 1
\newcommand{\MlstFtJcTwoAccFullEFour}{0.3914\,\ensuremath{\pm}\,0.0011}  % experiments/FIGS/data/v2/index.json :: fine-tuning jc2 N=10000: accuracy, L188 :: sha256 c6c4138ee5c86423
\newcommand{\MlstFtJcTwoAccResEFour}{0.3934\,\ensuremath{\pm}\,0.0018}  % experiments/FIGS/data/v2/index.json :: fine-tuning jc2 N=10000: accuracy, L162 :: sha256 c6c4138ee5c86423
\newcommand{\MlstFtJcTwoAccMidEFour}{0.3563\,\ensuremath{\pm}\,0.0025}  % experiments/FIGS/data/v2/index.json :: fine-tuning jc2 N=10000: accuracy, R42_Q1 :: sha256 c6c4138ee5c86423
\newcommand{\MlstFtJcTwoAccCoarseEFour}{0.3219\,\ensuremath{\pm}\,0.0011}  % experiments/FIGS/data/v2/index.json :: fine-tuning jc2 N=10000: accuracy, R16_Q1 :: sha256 c6c4138ee5c86423
\newcommand{\MlstFtJcTwoAccResMassEFour}{0.381 and 0.392}  % experiments/FIGS/data/v2/index.json :: fine-tuning jc2 N=10000: accuracy, L162_MASS :: sha256 c6c4138ee5c86423
\newcommand{\MlstFtJcTwoAccCoarseMassEFour}{0.3156 and 0.3187}  % experiments/FIGS/data/v2/index.json :: fine-tuning jc2 N=10000: accuracy, R16_Q1_MASS :: sha256 c6c4138ee5c86423
\newcommand{\MlstFtJcTwoAccScratchEFour}{0.1492\,\ensuremath{\pm}\,0.0011}  % experiments/FIGS/data/v2/index.json :: fine-tuning jc2 N=10000: accuracy from scratch :: sha256 c6c4138ee5c86423
\newcommand{\MlstFtJcTwoAccFullMinusScratchEFour}{0.2423\,\ensuremath{\pm}\,0.0011}  % experiments/FIGS/data/v2/index.json :: fine-tuning jc2 N=10000: accuracy of each 188-class run minus the from-scratch mean (absolute) :: sha256 c6c4138ee5c86423
\newcommand{\MlstFtJcTwoAccFullEFive}{0.4164\,\ensuremath{\pm}\,0.0015}  % experiments/FIGS/data/v2/index.json :: fine-tuning jc2 N=100000: accuracy, L188 :: sha256 c6c4138ee5c86423
\newcommand{\MlstFtJcTwoAccResEFive}{0.4161\,\ensuremath{\pm}\,0.0015}  % experiments/FIGS/data/v2/index.json :: fine-tuning jc2 N=100000: accuracy, L162 :: sha256 c6c4138ee5c86423
\newcommand{\MlstFtJcTwoAccMidEFive}{0.4007\,\ensuremath{\pm}\,0.0013}  % experiments/FIGS/data/v2/index.json :: fine-tuning jc2 N=100000: accuracy, R42_Q1 :: sha256 c6c4138ee5c86423
\newcommand{\MlstFtJcTwoAccCoarseEFive}{0.3833\,\ensuremath{\pm}\,0.0020}  % experiments/FIGS/data/v2/index.json :: fine-tuning jc2 N=100000: accuracy, R16_Q1 :: sha256 c6c4138ee5c86423
\newcommand{\MlstFtJcTwoAccResMassEFive}{0.4093 and 0.4115}  % experiments/FIGS/data/v2/index.json :: fine-tuning jc2 N=100000: accuracy, L162_MASS :: sha256 c6c4138ee5c86423
\newcommand{\MlstFtJcTwoAccCoarseMassEFive}{0.3758 and 0.3781}  % experiments/FIGS/data/v2/index.json :: fine-tuning jc2 N=100000: accuracy, R16_Q1_MASS :: sha256 c6c4138ee5c86423
\newcommand{\MlstFtJcTwoAccScratchEFive}{0.292\,\ensuremath{\pm}\,0.004}  % experiments/FIGS/data/v2/index.json :: fine-tuning jc2 N=100000: accuracy from scratch :: sha256 c6c4138ee5c86423
\newcommand{\MlstFtJcTwoAccFullMinusScratchEFive}{0.1243\,\ensuremath{\pm}\,0.0015}  % experiments/FIGS/data/v2/index.json :: fine-tuning jc2 N=100000: accuracy of each 188-class run minus the from-scratch mean (absolute) :: sha256 c6c4138ee5c86423
\newcommand{\MlstFtJcTwoAccFullESix}{0.4263\,\ensuremath{\pm}\,0.0005}  % experiments/FIGS/data/v2/index.json :: fine-tuning jc2 N=1000000: accuracy, L188 :: sha256 c6c4138ee5c86423
\newcommand{\MlstFtJcTwoAccResESix}{0.4249\,\ensuremath{\pm}\,0.0021}  % experiments/FIGS/data/v2/index.json :: fine-tuning jc2 N=1000000: accuracy, L162 :: sha256 c6c4138ee5c86423
\newcommand{\MlstFtJcTwoAccMidESix}{0.4166\,\ensuremath{\pm}\,0.0011}  % experiments/FIGS/data/v2/index.json :: fine-tuning jc2 N=1000000: accuracy, R42_Q1 :: sha256 c6c4138ee5c86423
\newcommand{\MlstFtJcTwoAccCoarseESix}{0.4054\,\ensuremath{\pm}\,0.0020}  % experiments/FIGS/data/v2/index.json :: fine-tuning jc2 N=1000000: accuracy, R16_Q1 :: sha256 c6c4138ee5c86423
\newcommand{\MlstFtJcTwoAccResMassESix}{0.4198 and 0.4213}  % experiments/FIGS/data/v2/index.json :: fine-tuning jc2 N=1000000: accuracy, L162_MASS :: sha256 c6c4138ee5c86423
\newcommand{\MlstFtJcTwoAccCoarseMassESix}{0.4012 and 0.4000}  % experiments/FIGS/data/v2/index.json :: fine-tuning jc2 N=1000000: accuracy, R16_Q1_MASS :: sha256 c6c4138ee5c86423
\newcommand{\MlstFtJcTwoAccScratchESix}{0.3677\,\ensuremath{\pm}\,0.0012}  % experiments/FIGS/data/v2/index.json :: fine-tuning jc2 N=1000000: accuracy from scratch :: sha256 c6c4138ee5c86423
\newcommand{\MlstFtJcTwoAccFullMinusScratchESix}{0.0585\,\ensuremath{\pm}\,0.0005}  % experiments/FIGS/data/v2/index.json :: fine-tuning jc2 N=1000000: accuracy of each 188-class run minus the from-scratch mean (absolute) :: sha256 c6c4138ee5c86423
\newcommand{\MlstFtJcTwoAucFull}{0.9797\,\ensuremath{\pm}\,0.0001}  % experiments/FIGS/data/v2/index.json :: fine-tuning class-balanced JetClass-II N=1000000: AUC averaged over the classes, Full :: sha256 c6c4138ee5c86423
\newcommand{\MlstFtJcTwoAucScratch}{0.9734\,\ensuremath{\pm}\,0.0003}  % experiments/FIGS/data/v2/index.json :: fine-tuning class-balanced JetClass-II N=1000000: AUC averaged over the classes, Scratch :: sha256 c6c4138ee5c86423
\newcommand{\MlstFtJcOneAccFullEFour}{0.8117\,\ensuremath{\pm}\,0.0031}  % experiments/FIGS/data/v2/index.json :: fine-tuning jc1 N=10000: accuracy, L188 :: sha256 c6c4138ee5c86423
\newcommand{\MlstFtJcOneAccResEFour}{0.8117\,\ensuremath{\pm}\,0.0013}  % experiments/FIGS/data/v2/index.json :: fine-tuning jc1 N=10000: accuracy, L162 :: sha256 c6c4138ee5c86423
\newcommand{\MlstFtJcOneAccMidEFour}{0.8086\,\ensuremath{\pm}\,0.0031}  % experiments/FIGS/data/v2/index.json :: fine-tuning jc1 N=10000: accuracy, R42_Q1 :: sha256 c6c4138ee5c86423
\newcommand{\MlstFtJcOneAccCoarseEFour}{0.7982\,\ensuremath{\pm}\,0.0023}  % experiments/FIGS/data/v2/index.json :: fine-tuning jc1 N=10000: accuracy, R16_Q1 :: sha256 c6c4138ee5c86423
\newcommand{\MlstFtJcOneAccResMassEFour}{0.8107 and 0.8098}  % experiments/FIGS/data/v2/index.json :: fine-tuning jc1 N=10000: accuracy, L162_MASS :: sha256 c6c4138ee5c86423
\newcommand{\MlstFtJcOneAccCoarseMassEFour}{0.7944 and 0.7990}  % experiments/FIGS/data/v2/index.json :: fine-tuning jc1 N=10000: accuracy, R16_Q1_MASS :: sha256 c6c4138ee5c86423
\newcommand{\MlstFtJcOneAccScratchEFour}{0.677\,\ensuremath{\pm}\,0.004}  % experiments/FIGS/data/v2/index.json :: fine-tuning jc1 N=10000: accuracy from scratch :: sha256 c6c4138ee5c86423
\newcommand{\MlstFtJcOneAccFullMinusScratchEFour}{0.1343\,\ensuremath{\pm}\,0.0031}  % experiments/FIGS/data/v2/index.json :: fine-tuning jc1 N=10000: accuracy of each 188-class run minus the from-scratch mean (absolute) :: sha256 c6c4138ee5c86423
\newcommand{\MlstFtJcOneAccFullEFive}{0.8343\,\ensuremath{\pm}\,0.0035}  % experiments/FIGS/data/v2/index.json :: fine-tuning jc1 N=100000: accuracy, L188 :: sha256 c6c4138ee5c86423
\newcommand{\MlstFtJcOneAccResEFive}{0.8328\,\ensuremath{\pm}\,0.0005}  % experiments/FIGS/data/v2/index.json :: fine-tuning jc1 N=100000: accuracy, L162 :: sha256 c6c4138ee5c86423
\newcommand{\MlstFtJcOneAccMidEFive}{0.8324\,\ensuremath{\pm}\,0.0004}  % experiments/FIGS/data/v2/index.json :: fine-tuning jc1 N=100000: accuracy, R42_Q1 :: sha256 c6c4138ee5c86423
\newcommand{\MlstFtJcOneAccCoarseEFive}{0.8288\,\ensuremath{\pm}\,0.0005}  % experiments/FIGS/data/v2/index.json :: fine-tuning jc1 N=100000: accuracy, R16_Q1 :: sha256 c6c4138ee5c86423
\newcommand{\MlstFtJcOneAccResMassEFive}{0.8326 and 0.8348}  % experiments/FIGS/data/v2/index.json :: fine-tuning jc1 N=100000: accuracy, L162_MASS :: sha256 c6c4138ee5c86423
\newcommand{\MlstFtJcOneAccCoarseMassEFive}{0.8253 and 0.8262}  % experiments/FIGS/data/v2/index.json :: fine-tuning jc1 N=100000: accuracy, R16_Q1_MASS :: sha256 c6c4138ee5c86423
\newcommand{\MlstFtJcOneAccScratchEFive}{0.777\,\ensuremath{\pm}\,0.005}  % experiments/FIGS/data/v2/index.json :: fine-tuning jc1 N=100000: accuracy from scratch :: sha256 c6c4138ee5c86423
\newcommand{\MlstFtJcOneAccFullMinusScratchEFive}{0.0571\,\ensuremath{\pm}\,0.0035}  % experiments/FIGS/data/v2/index.json :: fine-tuning jc1 N=100000: accuracy of each 188-class run minus the from-scratch mean (absolute) :: sha256 c6c4138ee5c86423
\newcommand{\MlstFtJcOneAccFullESix}{0.8468\,\ensuremath{\pm}\,0.0006}  % experiments/FIGS/data/v2/index.json :: fine-tuning jc1 N=1000000: accuracy, L188 :: sha256 c6c4138ee5c86423
\newcommand{\MlstFtJcOneAccResESix}{0.8475\,\ensuremath{\pm}\,0.0009}  % experiments/FIGS/data/v2/index.json :: fine-tuning jc1 N=1000000: accuracy, L162 :: sha256 c6c4138ee5c86423
\newcommand{\MlstFtJcOneAccMidESix}{0.8459\,\ensuremath{\pm}\,0.0002}  % experiments/FIGS/data/v2/index.json :: fine-tuning jc1 N=1000000: accuracy, R42_Q1 :: sha256 c6c4138ee5c86423
\newcommand{\MlstFtJcOneAccCoarseESix}{0.8422\,\ensuremath{\pm}\,0.0005}  % experiments/FIGS/data/v2/index.json :: fine-tuning jc1 N=1000000: accuracy, R16_Q1 :: sha256 c6c4138ee5c86423
\newcommand{\MlstFtJcOneAccResMassESix}{0.8459 and 0.8471}  % experiments/FIGS/data/v2/index.json :: fine-tuning jc1 N=1000000: accuracy, L162_MASS :: sha256 c6c4138ee5c86423
\newcommand{\MlstFtJcOneAccCoarseMassESix}{0.8400 and 0.8398}  % experiments/FIGS/data/v2/index.json :: fine-tuning jc1 N=1000000: accuracy, R16_Q1_MASS :: sha256 c6c4138ee5c86423
\newcommand{\MlstFtJcOneAccScratchESix}{0.8267\,\ensuremath{\pm}\,0.0011}  % experiments/FIGS/data/v2/index.json :: fine-tuning jc1 N=1000000: accuracy from scratch :: sha256 c6c4138ee5c86423
\newcommand{\MlstFtJcOneAccFullMinusScratchESix}{0.0201\,\ensuremath{\pm}\,0.0006}  % experiments/FIGS/data/v2/index.json :: fine-tuning jc1 N=1000000: accuracy of each 188-class run minus the from-scratch mean (absolute) :: sha256 c6c4138ee5c86423
\newcommand{\MlstFtTopRthirtyScratchEThree}{49\,\ensuremath{\pm}\,12}  % experiments/FIGS/data/v2/index.json :: fine-tuning top N=1000: R30 from scratch, every seed :: sha256 c6c4138ee5c86423
\newcommand{\MlstFtTopRthirtyFullEThree}{72\,\ensuremath{\pm}\,19}  % experiments/FIGS/data/v2/index.json :: fine-tuning top N=1000: R30, L188 :: sha256 c6c4138ee5c86423
\newcommand{\MlstFtTopRthirtyFullOverScratchEThree}{1.5\,\ensuremath{\pm}\,0.4}  % experiments/FIGS/data/v2/index.json :: fine-tuning top N=1000: R30 of each L188 run over the from-scratch mean, SD over runs and standard error of that mean in quadrature :: sha256 c6c4138ee5c86423
\newcommand{\MlstFtTopRthirtyFullOverScratchEThreePr}{equal within 50\%}  % experiments/FIGS/data/v2/index.json :: fine-tuning top N=1000: R30 of each L188 run over the from-scratch mean, SD over runs and standard error of that mean in quadrature :: sha256 c6c4138ee5c86423 :: prose twin: two significant figures, direction in words
\newcommand{\MlstFtTopRthirtyResEThree}{135\,\ensuremath{\pm}\,26}  % experiments/FIGS/data/v2/index.json :: fine-tuning top N=1000: R30, L162 :: sha256 c6c4138ee5c86423
\newcommand{\MlstFtTopRthirtyResOverScratchEThree}{2.8\,\ensuremath{\pm}\,0.7}  % experiments/FIGS/data/v2/index.json :: fine-tuning top N=1000: R30 of each L162 run over the from-scratch mean, SD over runs and standard error of that mean in quadrature :: sha256 c6c4138ee5c86423
\newcommand{\MlstFtTopRthirtyResOverScratchEThreePr}{2.8 times the rejection}  % experiments/FIGS/data/v2/index.json :: fine-tuning top N=1000: R30 of each L162 run over the from-scratch mean, SD over runs and standard error of that mean in quadrature :: sha256 c6c4138ee5c86423 :: prose twin: two significant figures, direction in words
\newcommand{\MlstFtTopRthirtyMidEThree}{104\,\ensuremath{\pm}\,20}  % experiments/FIGS/data/v2/index.json :: fine-tuning top N=1000: R30, R42_Q1 :: sha256 c6c4138ee5c86423
\newcommand{\MlstFtTopRthirtyMidOverScratchEThree}{2.1\,\ensuremath{\pm}\,0.5}  % experiments/FIGS/data/v2/index.json :: fine-tuning top N=1000: R30 of each R42_Q1 run over the from-scratch mean, SD over runs and standard error of that mean in quadrature :: sha256 c6c4138ee5c86423
\newcommand{\MlstFtTopRthirtyMidOverScratchEThreePr}{2.1 times the rejection}  % experiments/FIGS/data/v2/index.json :: fine-tuning top N=1000: R30 of each R42_Q1 run over the from-scratch mean, SD over runs and standard error of that mean in quadrature :: sha256 c6c4138ee5c86423 :: prose twin: two significant figures, direction in words
\newcommand{\MlstFtTopRthirtyCoarseEThree}{83\,\ensuremath{\pm}\,25}  % experiments/FIGS/data/v2/index.json :: fine-tuning top N=1000: R30, R16_Q1 :: sha256 c6c4138ee5c86423
\newcommand{\MlstFtTopRthirtyCoarseOverScratchEThree}{1.7\,\ensuremath{\pm}\,0.6}  % experiments/FIGS/data/v2/index.json :: fine-tuning top N=1000: R30 of each R16_Q1 run over the from-scratch mean, SD over runs and standard error of that mean in quadrature :: sha256 c6c4138ee5c86423
\newcommand{\MlstFtTopRthirtyCoarseOverScratchEThreePr}{equal within 80\%}  % experiments/FIGS/data/v2/index.json :: fine-tuning top N=1000: R30 of each R16_Q1 run over the from-scratch mean, SD over runs and standard error of that mean in quadrature :: sha256 c6c4138ee5c86423 :: prose twin: two significant figures, direction in words
\newcommand{\MlstFtTopRthirtyResMassEThree}{105.9 and 112.7}  % experiments/FIGS/data/v2/index.json :: fine-tuning top N=1000: R30, L162_MASS :: sha256 c6c4138ee5c86423
\newcommand{\MlstFtTopRthirtyResMassOverScratchEThree}{2.18 and 2.32}  % experiments/FIGS/data/v2/index.json :: fine-tuning top N=1000: R30 of each L162_MASS run over the from-scratch mean, SD over runs and standard error of that mean in quadrature :: sha256 c6c4138ee5c86423
\newcommand{\MlstFtTopRthirtyResMassOverScratchEThreePr}{at least 2.2 times the rejection}  % experiments/FIGS/data/v2/index.json :: fine-tuning top N=1000: R30 of each L162_MASS run over the from-scratch mean, SD over runs and standard error of that mean in quadrature :: sha256 c6c4138ee5c86423 :: prose twin: two significant figures, direction in words
\newcommand{\MlstFtTopRthirtyResMassOverScratchEThreePrMax}{up to 2.3 times the rejection}  % experiments/FIGS/data/v2/index.json :: fine-tuning top N=1000: R30 of each L162_MASS run over the from-scratch mean, SD over runs and standard error of that mean in quadrature :: sha256 c6c4138ee5c86423 :: prose twin of the run farther from 1
\newcommand{\MlstFtTopRthirtyCoarseMassEThree}{59.1 and 63.9}  % experiments/FIGS/data/v2/index.json :: fine-tuning top N=1000: R30, R16_Q1_MASS :: sha256 c6c4138ee5c86423
\newcommand{\MlstFtTopRthirtyCoarseMassOverScratchEThree}{1.216 and 1.314}  % experiments/FIGS/data/v2/index.json :: fine-tuning top N=1000: R30 of each R16_Q1_MASS run over the from-scratch mean, SD over runs and standard error of that mean in quadrature :: sha256 c6c4138ee5c86423
\newcommand{\MlstFtTopRthirtyCoarseMassOverScratchEThreePr}{at least 22\% more rejection}  % experiments/FIGS/data/v2/index.json :: fine-tuning top N=1000: R30 of each R16_Q1_MASS run over the from-scratch mean, SD over runs and standard error of that mean in quadrature :: sha256 c6c4138ee5c86423 :: prose twin: two significant figures, direction in words
\newcommand{\MlstFtTopRthirtyCoarseMassOverScratchEThreePrMax}{up to 31\% more rejection}  % experiments/FIGS/data/v2/index.json :: fine-tuning top N=1000: R30 of each R16_Q1_MASS run over the from-scratch mean, SD over runs and standard error of that mean in quadrature :: sha256 c6c4138ee5c86423 :: prose twin of the run farther from 1
\newcommand{\MlstFtTopRthirtyResOverFullEThree}{2.0\,\ensuremath{\pm}\,0.8}  % experiments/FIGS/data/v2/index.json :: fine-tuning top N=1000: R30 of L162 over L188, run by run :: sha256 c6c4138ee5c86423
\newcommand{\MlstFtTopRthirtyResOverFullEThreePr}{equal within 100\%}  % experiments/FIGS/data/v2/index.json :: fine-tuning top N=1000: R30 of L162 over L188, run by run :: sha256 c6c4138ee5c86423 :: prose twin: two significant figures, direction in words
\newcommand{\MlstFtTopRthirtyMidOverFullEThree}{1.46\,\ensuremath{\pm}\,0.12}  % experiments/FIGS/data/v2/index.json :: fine-tuning top N=1000: R30 of R42_Q1 over L188, run by run :: sha256 c6c4138ee5c86423
\newcommand{\MlstFtTopRthirtyMidOverFullEThreePr}{46\% more rejection}  % experiments/FIGS/data/v2/index.json :: fine-tuning top N=1000: R30 of R42_Q1 over L188, run by run :: sha256 c6c4138ee5c86423 :: prose twin: two significant figures, direction in words
\newcommand{\MlstFtTopRthirtyCoarseOverFullEThree}{1.3\,\ensuremath{\pm}\,0.7}  % experiments/FIGS/data/v2/index.json :: fine-tuning top N=1000: R30 of R16_Q1 over L188, run by run :: sha256 c6c4138ee5c86423
\newcommand{\MlstFtTopRthirtyCoarseOverFullEThreePr}{equal within 80\%}  % experiments/FIGS/data/v2/index.json :: fine-tuning top N=1000: R30 of R16_Q1 over L188, run by run :: sha256 c6c4138ee5c86423 :: prose twin: two significant figures, direction in words
\newcommand{\MlstFtTopRthirtyResMassOverResEThree}{0.73 and 1.07}  % experiments/FIGS/data/v2/index.json :: fine-tuning top N=1000: R30 of L162_MASS over L162, run by run :: sha256 c6c4138ee5c86423
\newcommand{\MlstFtTopRthirtyResMassOverResEThreePr}{equal within 30\%}  % experiments/FIGS/data/v2/index.json :: fine-tuning top N=1000: R30 of L162_MASS over L162, run by run :: sha256 c6c4138ee5c86423 :: prose twin: two significant figures, direction in words
\newcommand{\MlstFtTopRthirtyResMassOverResEThreePrMax}{equal within 30\%}  % experiments/FIGS/data/v2/index.json :: fine-tuning top N=1000: R30 of L162_MASS over L162, run by run :: sha256 c6c4138ee5c86423 :: prose twin of the run farther from 1
\newcommand{\MlstFtTopRthirtyCoarseMassOverCoarseEThree}{0.56 and 0.75}  % experiments/FIGS/data/v2/index.json :: fine-tuning top N=1000: R30 of R16_Q1_MASS over R16_Q1, run by run :: sha256 c6c4138ee5c86423
\newcommand{\MlstFtTopRthirtyCoarseMassOverCoarseEThreePr}{equal within 50\%}  % experiments/FIGS/data/v2/index.json :: fine-tuning top N=1000: R30 of R16_Q1_MASS over R16_Q1, run by run :: sha256 c6c4138ee5c86423 :: prose twin: two significant figures, direction in words
\newcommand{\MlstFtTopRthirtyCoarseMassOverCoarseEThreePrMax}{equal within 50\%}  % experiments/FIGS/data/v2/index.json :: fine-tuning top N=1000: R30 of R16_Q1_MASS over R16_Q1, run by run :: sha256 c6c4138ee5c86423 :: prose twin of the run farther from 1
\newcommand{\MlstFtTopRthirtyRandOverFullEThree}{1.3\,\ensuremath{\pm}\,0.5}  % experiments/FIGS/data/v2/index.json :: fine-tuning top N=1000: R30 of each random partition over the L188 run of its seed :: sha256 c6c4138ee5c86423
\newcommand{\MlstFtTopRthirtyRandOverFullEThreePr}{equal within 50\%}  % experiments/FIGS/data/v2/index.json :: fine-tuning top N=1000: R30 of each random partition over the L188 run of its seed :: sha256 c6c4138ee5c86423 :: prose twin: two significant figures, direction in words
\newcommand{\MlstFtTopRthirtyRandOverCoarseEThree}{0.87\,\ensuremath{\pm}\,0.12}  % experiments/FIGS/data/v2/index.json :: fine-tuning top N=1000: R30 of each random partition over the R16_Q1 run of its seed :: sha256 c6c4138ee5c86423
\newcommand{\MlstFtTopRthirtyRandOverCoarseEThreePr}{equal within 20\%}  % experiments/FIGS/data/v2/index.json :: fine-tuning top N=1000: R30 of each random partition over the R16_Q1 run of its seed :: sha256 c6c4138ee5c86423 :: prose twin: two significant figures, direction in words
\newcommand{\MlstFtTopRthirtyScratchEFour}{340\,\ensuremath{\pm}\,50}  % experiments/FIGS/data/v2/index.json :: fine-tuning top N=10000: R30 from scratch, every seed :: sha256 c6c4138ee5c86423
\newcommand{\MlstFtTopRthirtyFullEFour}{420\,\ensuremath{\pm}\,9}  % experiments/FIGS/data/v2/index.json :: fine-tuning top N=10000: R30, L188 :: sha256 c6c4138ee5c86423
\newcommand{\MlstFtTopRthirtyFullOverScratchEFour}{1.24\,\ensuremath{\pm}\,0.11}  % experiments/FIGS/data/v2/index.json :: fine-tuning top N=10000: R30 of each L188 run over the from-scratch mean, SD over runs and standard error of that mean in quadrature :: sha256 c6c4138ee5c86423
\newcommand{\MlstFtTopRthirtyFullOverScratchEFourPr}{24\% more rejection}  % experiments/FIGS/data/v2/index.json :: fine-tuning top N=10000: R30 of each L188 run over the from-scratch mean, SD over runs and standard error of that mean in quadrature :: sha256 c6c4138ee5c86423 :: prose twin: two significant figures, direction in words
\newcommand{\MlstFtTopRthirtyResEFour}{430\,\ensuremath{\pm}\,60}  % experiments/FIGS/data/v2/index.json :: fine-tuning top N=10000: R30, L162 :: sha256 c6c4138ee5c86423
\newcommand{\MlstFtTopRthirtyResOverScratchEFour}{1.27\,\ensuremath{\pm}\,0.21}  % experiments/FIGS/data/v2/index.json :: fine-tuning top N=10000: R30 of each L162 run over the from-scratch mean, SD over runs and standard error of that mean in quadrature :: sha256 c6c4138ee5c86423
\newcommand{\MlstFtTopRthirtyResOverScratchEFourPr}{equal within 30\%}  % experiments/FIGS/data/v2/index.json :: fine-tuning top N=10000: R30 of each L162 run over the from-scratch mean, SD over runs and standard error of that mean in quadrature :: sha256 c6c4138ee5c86423 :: prose twin: two significant figures, direction in words
\newcommand{\MlstFtTopRthirtyMidEFour}{455\,\ensuremath{\pm}\,30}  % experiments/FIGS/data/v2/index.json :: fine-tuning top N=10000: R30, R42_Q1 :: sha256 c6c4138ee5c86423
\newcommand{\MlstFtTopRthirtyMidOverScratchEFour}{1.34\,\ensuremath{\pm}\,0.14}  % experiments/FIGS/data/v2/index.json :: fine-tuning top N=10000: R30 of each R42_Q1 run over the from-scratch mean, SD over runs and standard error of that mean in quadrature :: sha256 c6c4138ee5c86423
\newcommand{\MlstFtTopRthirtyMidOverScratchEFourPr}{34\% more rejection}  % experiments/FIGS/data/v2/index.json :: fine-tuning top N=10000: R30 of each R42_Q1 run over the from-scratch mean, SD over runs and standard error of that mean in quadrature :: sha256 c6c4138ee5c86423 :: prose twin: two significant figures, direction in words
\newcommand{\MlstFtTopRthirtyCoarseEFour}{350\,\ensuremath{\pm}\,40}  % experiments/FIGS/data/v2/index.json :: fine-tuning top N=10000: R30, R16_Q1 :: sha256 c6c4138ee5c86423
\newcommand{\MlstFtTopRthirtyCoarseOverScratchEFour}{1.04\,\ensuremath{\pm}\,0.14}  % experiments/FIGS/data/v2/index.json :: fine-tuning top N=10000: R30 of each R16_Q1 run over the from-scratch mean, SD over runs and standard error of that mean in quadrature :: sha256 c6c4138ee5c86423
\newcommand{\MlstFtTopRthirtyCoarseOverScratchEFourPr}{equal within 20\%}  % experiments/FIGS/data/v2/index.json :: fine-tuning top N=10000: R30 of each R16_Q1 run over the from-scratch mean, SD over runs and standard error of that mean in quadrature :: sha256 c6c4138ee5c86423 :: prose twin: two significant figures, direction in words
\newcommand{\MlstFtTopRthirtyResMassEFour}{373.2 and 376.7}  % experiments/FIGS/data/v2/index.json :: fine-tuning top N=10000: R30, L162_MASS :: sha256 c6c4138ee5c86423
\newcommand{\MlstFtTopRthirtyResMassOverScratchEFour}{1.101 and 1.111}  % experiments/FIGS/data/v2/index.json :: fine-tuning top N=10000: R30 of each L162_MASS run over the from-scratch mean, SD over runs and standard error of that mean in quadrature :: sha256 c6c4138ee5c86423
\newcommand{\MlstFtTopRthirtyResMassOverScratchEFourPr}{at least 10\% more rejection}  % experiments/FIGS/data/v2/index.json :: fine-tuning top N=10000: R30 of each L162_MASS run over the from-scratch mean, SD over runs and standard error of that mean in quadrature :: sha256 c6c4138ee5c86423 :: prose twin: two significant figures, direction in words
\newcommand{\MlstFtTopRthirtyResMassOverScratchEFourPrMax}{up to 11\% more rejection}  % experiments/FIGS/data/v2/index.json :: fine-tuning top N=10000: R30 of each L162_MASS run over the from-scratch mean, SD over runs and standard error of that mean in quadrature :: sha256 c6c4138ee5c86423 :: prose twin of the run farther from 1
\newcommand{\MlstFtTopRthirtyCoarseMassEFour}{365.8 and 359.9}  % experiments/FIGS/data/v2/index.json :: fine-tuning top N=10000: R30, R16_Q1_MASS :: sha256 c6c4138ee5c86423
\newcommand{\MlstFtTopRthirtyCoarseMassOverScratchEFour}{1.079 and 1.062}  % experiments/FIGS/data/v2/index.json :: fine-tuning top N=10000: R30 of each R16_Q1_MASS run over the from-scratch mean, SD over runs and standard error of that mean in quadrature :: sha256 c6c4138ee5c86423
\newcommand{\MlstFtTopRthirtyCoarseMassOverScratchEFourPr}{at least 6.2\% more rejection}  % experiments/FIGS/data/v2/index.json :: fine-tuning top N=10000: R30 of each R16_Q1_MASS run over the from-scratch mean, SD over runs and standard error of that mean in quadrature :: sha256 c6c4138ee5c86423 :: prose twin: two significant figures, direction in words
\newcommand{\MlstFtTopRthirtyCoarseMassOverScratchEFourPrMax}{up to 7.9\% more rejection}  % experiments/FIGS/data/v2/index.json :: fine-tuning top N=10000: R30 of each R16_Q1_MASS run over the from-scratch mean, SD over runs and standard error of that mean in quadrature :: sha256 c6c4138ee5c86423 :: prose twin of the run farther from 1
\newcommand{\MlstFtTopRthirtyResOverFullEFour}{1.02\,\ensuremath{\pm}\,0.12}  % experiments/FIGS/data/v2/index.json :: fine-tuning top N=10000: R30 of L162 over L188, run by run :: sha256 c6c4138ee5c86423
\newcommand{\MlstFtTopRthirtyResOverFullEFourPr}{equal within 20\%}  % experiments/FIGS/data/v2/index.json :: fine-tuning top N=10000: R30 of L162 over L188, run by run :: sha256 c6c4138ee5c86423 :: prose twin: two significant figures, direction in words
\newcommand{\MlstFtTopRthirtyMidOverFullEFour}{1.08\,\ensuremath{\pm}\,0.08}  % experiments/FIGS/data/v2/index.json :: fine-tuning top N=10000: R30 of R42_Q1 over L188, run by run :: sha256 c6c4138ee5c86423
\newcommand{\MlstFtTopRthirtyMidOverFullEFourPr}{equal within 9\%}  % experiments/FIGS/data/v2/index.json :: fine-tuning top N=10000: R30 of R42_Q1 over L188, run by run :: sha256 c6c4138ee5c86423 :: prose twin: two significant figures, direction in words
\newcommand{\MlstFtTopRthirtyCoarseOverFullEFour}{0.84\,\ensuremath{\pm}\,0.07}  % experiments/FIGS/data/v2/index.json :: fine-tuning top N=10000: R30 of R16_Q1 over L188, run by run :: sha256 c6c4138ee5c86423
\newcommand{\MlstFtTopRthirtyCoarseOverFullEFourPr}{16\% less rejection}  % experiments/FIGS/data/v2/index.json :: fine-tuning top N=10000: R30 of R16_Q1 over L188, run by run :: sha256 c6c4138ee5c86423 :: prose twin: two significant figures, direction in words
\newcommand{\MlstFtTopRthirtyResMassOverResEFour}{1.00 and 0.76}  % experiments/FIGS/data/v2/index.json :: fine-tuning top N=10000: R30 of L162_MASS over L162, run by run :: sha256 c6c4138ee5c86423
\newcommand{\MlstFtTopRthirtyResMassOverResEFourPr}{equal within 30\%}  % experiments/FIGS/data/v2/index.json :: fine-tuning top N=10000: R30 of L162_MASS over L162, run by run :: sha256 c6c4138ee5c86423 :: prose twin: two significant figures, direction in words
\newcommand{\MlstFtTopRthirtyResMassOverResEFourPrMax}{equal within 30\%}  % experiments/FIGS/data/v2/index.json :: fine-tuning top N=10000: R30 of L162_MASS over L162, run by run :: sha256 c6c4138ee5c86423 :: prose twin of the run farther from 1
\newcommand{\MlstFtTopRthirtyCoarseMassOverCoarseEFour}{1.13 and 0.92}  % experiments/FIGS/data/v2/index.json :: fine-tuning top N=10000: R30 of R16_Q1_MASS over R16_Q1, run by run :: sha256 c6c4138ee5c86423
\newcommand{\MlstFtTopRthirtyCoarseMassOverCoarseEFourPr}{equal within 20\%}  % experiments/FIGS/data/v2/index.json :: fine-tuning top N=10000: R30 of R16_Q1_MASS over R16_Q1, run by run :: sha256 c6c4138ee5c86423 :: prose twin: two significant figures, direction in words
\newcommand{\MlstFtTopRthirtyCoarseMassOverCoarseEFourPrMax}{equal within 20\%}  % experiments/FIGS/data/v2/index.json :: fine-tuning top N=10000: R30 of R16_Q1_MASS over R16_Q1, run by run :: sha256 c6c4138ee5c86423 :: prose twin of the run farther from 1
\newcommand{\MlstFtTopRthirtyRandOverFullEFour}{0.83\,\ensuremath{\pm}\,0.10}  % experiments/FIGS/data/v2/index.json :: fine-tuning top N=10000: R30 of each random partition over the L188 run of its seed :: sha256 c6c4138ee5c86423
\newcommand{\MlstFtTopRthirtyRandOverFullEFourPr}{equal within 20\%}  % experiments/FIGS/data/v2/index.json :: fine-tuning top N=10000: R30 of each random partition over the L188 run of its seed :: sha256 c6c4138ee5c86423 :: prose twin: two significant figures, direction in words
\newcommand{\MlstFtTopRthirtyRandOverCoarseEFour}{0.97\,\ensuremath{\pm}\,0.15}  % experiments/FIGS/data/v2/index.json :: fine-tuning top N=10000: R30 of each random partition over the R16_Q1 run of its seed :: sha256 c6c4138ee5c86423
\newcommand{\MlstFtTopRthirtyRandOverCoarseEFourPr}{equal within 20\%}  % experiments/FIGS/data/v2/index.json :: fine-tuning top N=10000: R30 of each random partition over the R16_Q1 run of its seed :: sha256 c6c4138ee5c86423 :: prose twin: two significant figures, direction in words
\newcommand{\MlstFtTopRthirtyScratchEFive}{560\,\ensuremath{\pm}\,90}  % experiments/FIGS/data/v2/index.json :: fine-tuning top N=100000: R30 from scratch, every seed :: sha256 c6c4138ee5c86423
\newcommand{\MlstFtTopRthirtyFullEFive}{850\,\ensuremath{\pm}\,90}  % experiments/FIGS/data/v2/index.json :: fine-tuning top N=100000: R30, L188 :: sha256 c6c4138ee5c86423
\newcommand{\MlstFtTopRthirtyFullOverScratchEFive}{1.50\,\ensuremath{\pm}\,0.21}  % experiments/FIGS/data/v2/index.json :: fine-tuning top N=100000: R30 of each L188 run over the from-scratch mean, SD over runs and standard error of that mean in quadrature :: sha256 c6c4138ee5c86423
\newcommand{\MlstFtTopRthirtyFullOverScratchEFivePr}{50\% more rejection}  % experiments/FIGS/data/v2/index.json :: fine-tuning top N=100000: R30 of each L188 run over the from-scratch mean, SD over runs and standard error of that mean in quadrature :: sha256 c6c4138ee5c86423 :: prose twin: two significant figures, direction in words
\newcommand{\MlstFtTopRthirtyResEFive}{810\,\ensuremath{\pm}\,60}  % experiments/FIGS/data/v2/index.json :: fine-tuning top N=100000: R30, L162 :: sha256 c6c4138ee5c86423
\newcommand{\MlstFtTopRthirtyResOverScratchEFive}{1.44\,\ensuremath{\pm}\,0.18}  % experiments/FIGS/data/v2/index.json :: fine-tuning top N=100000: R30 of each L162 run over the from-scratch mean, SD over runs and standard error of that mean in quadrature :: sha256 c6c4138ee5c86423
\newcommand{\MlstFtTopRthirtyResOverScratchEFivePr}{44\% more rejection}  % experiments/FIGS/data/v2/index.json :: fine-tuning top N=100000: R30 of each L162 run over the from-scratch mean, SD over runs and standard error of that mean in quadrature :: sha256 c6c4138ee5c86423 :: prose twin: two significant figures, direction in words
\newcommand{\MlstFtTopRthirtyMidEFive}{850\,\ensuremath{\pm}\,80}  % experiments/FIGS/data/v2/index.json :: fine-tuning top N=100000: R30, R42_Q1 :: sha256 c6c4138ee5c86423
\newcommand{\MlstFtTopRthirtyMidOverScratchEFive}{1.51\,\ensuremath{\pm}\,0.20}  % experiments/FIGS/data/v2/index.json :: fine-tuning top N=100000: R30 of each R42_Q1 run over the from-scratch mean, SD over runs and standard error of that mean in quadrature :: sha256 c6c4138ee5c86423
\newcommand{\MlstFtTopRthirtyMidOverScratchEFivePr}{51\% more rejection}  % experiments/FIGS/data/v2/index.json :: fine-tuning top N=100000: R30 of each R42_Q1 run over the from-scratch mean, SD over runs and standard error of that mean in quadrature :: sha256 c6c4138ee5c86423 :: prose twin: two significant figures, direction in words
\newcommand{\MlstFtTopRthirtyCoarseEFive}{830\,\ensuremath{\pm}\,60}  % experiments/FIGS/data/v2/index.json :: fine-tuning top N=100000: R30, R16_Q1 :: sha256 c6c4138ee5c86423
\newcommand{\MlstFtTopRthirtyCoarseOverScratchEFive}{1.47\,\ensuremath{\pm}\,0.17}  % experiments/FIGS/data/v2/index.json :: fine-tuning top N=100000: R30 of each R16_Q1 run over the from-scratch mean, SD over runs and standard error of that mean in quadrature :: sha256 c6c4138ee5c86423
\newcommand{\MlstFtTopRthirtyCoarseOverScratchEFivePr}{47\% more rejection}  % experiments/FIGS/data/v2/index.json :: fine-tuning top N=100000: R30 of each R16_Q1 run over the from-scratch mean, SD over runs and standard error of that mean in quadrature :: sha256 c6c4138ee5c86423 :: prose twin: two significant figures, direction in words
\newcommand{\MlstFtTopRthirtyResMassEFive}{848 and 848}  % experiments/FIGS/data/v2/index.json :: fine-tuning top N=100000: R30, L162_MASS :: sha256 c6c4138ee5c86423
\newcommand{\MlstFtTopRthirtyResMassOverScratchEFive}{1.50 and 1.50}  % experiments/FIGS/data/v2/index.json :: fine-tuning top N=100000: R30 of each L162_MASS run over the from-scratch mean, SD over runs and standard error of that mean in quadrature :: sha256 c6c4138ee5c86423
\newcommand{\MlstFtTopRthirtyResMassOverScratchEFivePr}{at least 50\% more rejection}  % experiments/FIGS/data/v2/index.json :: fine-tuning top N=100000: R30 of each L162_MASS run over the from-scratch mean, SD over runs and standard error of that mean in quadrature :: sha256 c6c4138ee5c86423 :: prose twin: two significant figures, direction in words
\newcommand{\MlstFtTopRthirtyResMassOverScratchEFivePrMax}{up to 50\% more rejection}  % experiments/FIGS/data/v2/index.json :: fine-tuning top N=100000: R30 of each L162_MASS run over the from-scratch mean, SD over runs and standard error of that mean in quadrature :: sha256 c6c4138ee5c86423 :: prose twin of the run farther from 1
\newcommand{\MlstFtTopRthirtyCoarseMassEFive}{856 and 774}  % experiments/FIGS/data/v2/index.json :: fine-tuning top N=100000: R30, R16_Q1_MASS :: sha256 c6c4138ee5c86423
\newcommand{\MlstFtTopRthirtyCoarseMassOverScratchEFive}{1.51 and 1.37}  % experiments/FIGS/data/v2/index.json :: fine-tuning top N=100000: R30 of each R16_Q1_MASS run over the from-scratch mean, SD over runs and standard error of that mean in quadrature :: sha256 c6c4138ee5c86423
\newcommand{\MlstFtTopRthirtyCoarseMassOverScratchEFivePr}{at least 37\% more rejection}  % experiments/FIGS/data/v2/index.json :: fine-tuning top N=100000: R30 of each R16_Q1_MASS run over the from-scratch mean, SD over runs and standard error of that mean in quadrature :: sha256 c6c4138ee5c86423 :: prose twin: two significant figures, direction in words
\newcommand{\MlstFtTopRthirtyCoarseMassOverScratchEFivePrMax}{up to 51\% more rejection}  % experiments/FIGS/data/v2/index.json :: fine-tuning top N=100000: R30 of each R16_Q1_MASS run over the from-scratch mean, SD over runs and standard error of that mean in quadrature :: sha256 c6c4138ee5c86423 :: prose twin of the run farther from 1
\newcommand{\MlstFtTopRthirtyResOverFullEFive}{0.96\,\ensuremath{\pm}\,0.05}  % experiments/FIGS/data/v2/index.json :: fine-tuning top N=100000: R30 of L162 over L188, run by run :: sha256 c6c4138ee5c86423
\newcommand{\MlstFtTopRthirtyResOverFullEFivePr}{equal within 6\%}  % experiments/FIGS/data/v2/index.json :: fine-tuning top N=100000: R30 of L162 over L188, run by run :: sha256 c6c4138ee5c86423 :: prose twin: two significant figures, direction in words
\newcommand{\MlstFtTopRthirtyMidOverFullEFive}{1.02\,\ensuremath{\pm}\,0.18}  % experiments/FIGS/data/v2/index.json :: fine-tuning top N=100000: R30 of R42_Q1 over L188, run by run :: sha256 c6c4138ee5c86423
\newcommand{\MlstFtTopRthirtyMidOverFullEFivePr}{equal within 20\%}  % experiments/FIGS/data/v2/index.json :: fine-tuning top N=100000: R30 of R42_Q1 over L188, run by run :: sha256 c6c4138ee5c86423 :: prose twin: two significant figures, direction in words
\newcommand{\MlstFtTopRthirtyCoarseOverFullEFive}{0.99\,\ensuremath{\pm}\,0.17}  % experiments/FIGS/data/v2/index.json :: fine-tuning top N=100000: R30 of R16_Q1 over L188, run by run :: sha256 c6c4138ee5c86423
\newcommand{\MlstFtTopRthirtyCoarseOverFullEFivePr}{equal within 20\%}  % experiments/FIGS/data/v2/index.json :: fine-tuning top N=100000: R30 of R16_Q1 over L188, run by run :: sha256 c6c4138ee5c86423 :: prose twin: two significant figures, direction in words
\newcommand{\MlstFtTopRthirtyResMassOverResEFive}{1.00 and 1.15}  % experiments/FIGS/data/v2/index.json :: fine-tuning top N=100000: R30 of L162_MASS over L162, run by run :: sha256 c6c4138ee5c86423
\newcommand{\MlstFtTopRthirtyResMassOverResEFivePr}{equal within 20\%}  % experiments/FIGS/data/v2/index.json :: fine-tuning top N=100000: R30 of L162_MASS over L162, run by run :: sha256 c6c4138ee5c86423 :: prose twin: two significant figures, direction in words
\newcommand{\MlstFtTopRthirtyResMassOverResEFivePrMax}{equal within 20\%}  % experiments/FIGS/data/v2/index.json :: fine-tuning top N=100000: R30 of L162_MASS over L162, run by run :: sha256 c6c4138ee5c86423 :: prose twin of the run farther from 1
\newcommand{\MlstFtTopRthirtyCoarseMassOverCoarseEFive}{1.08 and 0.87}  % experiments/FIGS/data/v2/index.json :: fine-tuning top N=100000: R30 of R16_Q1_MASS over R16_Q1, run by run :: sha256 c6c4138ee5c86423
\newcommand{\MlstFtTopRthirtyCoarseMassOverCoarseEFivePr}{equal within 20\%}  % experiments/FIGS/data/v2/index.json :: fine-tuning top N=100000: R30 of R16_Q1_MASS over R16_Q1, run by run :: sha256 c6c4138ee5c86423 :: prose twin: two significant figures, direction in words
\newcommand{\MlstFtTopRthirtyCoarseMassOverCoarseEFivePrMax}{equal within 20\%}  % experiments/FIGS/data/v2/index.json :: fine-tuning top N=100000: R30 of R16_Q1_MASS over R16_Q1, run by run :: sha256 c6c4138ee5c86423 :: prose twin of the run farther from 1
\newcommand{\MlstFtTopRthirtyRandOverFullEFive}{1.01\,\ensuremath{\pm}\,0.22}  % experiments/FIGS/data/v2/index.json :: fine-tuning top N=100000: R30 of each random partition over the L188 run of its seed :: sha256 c6c4138ee5c86423
\newcommand{\MlstFtTopRthirtyRandOverFullEFivePr}{equal within 30\%}  % experiments/FIGS/data/v2/index.json :: fine-tuning top N=100000: R30 of each random partition over the L188 run of its seed :: sha256 c6c4138ee5c86423 :: prose twin: two significant figures, direction in words
\newcommand{\MlstFtTopRthirtyRandOverCoarseEFive}{0.97\,\ensuremath{\pm}\,0.10}  % experiments/FIGS/data/v2/index.json :: fine-tuning top N=100000: R30 of each random partition over the R16_Q1 run of its seed :: sha256 c6c4138ee5c86423
\newcommand{\MlstFtTopRthirtyRandOverCoarseEFivePr}{equal within 10\%}  % experiments/FIGS/data/v2/index.json :: fine-tuning top N=100000: R30 of each random partition over the R16_Q1 run of its seed :: sha256 c6c4138ee5c86423 :: prose twin: two significant figures, direction in words
\newcommand{\MlstFtTopRthirtyScratch}{1{,}412\,\ensuremath{\pm}\,20}  % experiments/FIGS/data/v2/index.json :: fine-tuning top N=1200000: R30 from scratch, every seed :: sha256 c6c4138ee5c86423
\newcommand{\MlstFtTopRthirtyFull}{1{,}800\,\ensuremath{\pm}\,60}  % experiments/FIGS/data/v2/index.json :: fine-tuning top N=1200000: R30, L188 :: sha256 c6c4138ee5c86423
\newcommand{\MlstFtTopRthirtyFullOverScratch}{1.27\,\ensuremath{\pm}\,0.04}  % experiments/FIGS/data/v2/index.json :: fine-tuning top N=1200000: R30 of each L188 run over the from-scratch mean, SD over runs and standard error of that mean in quadrature :: sha256 c6c4138ee5c86423
\newcommand{\MlstFtTopRthirtyFullOverScratchPr}{27\% more rejection}  % experiments/FIGS/data/v2/index.json :: fine-tuning top N=1200000: R30 of each L188 run over the from-scratch mean, SD over runs and standard error of that mean in quadrature :: sha256 c6c4138ee5c86423 :: prose twin: two significant figures, direction in words
\newcommand{\MlstFtTopRthirtyRes}{1{,}790\,\ensuremath{\pm}\,60}  % experiments/FIGS/data/v2/index.json :: fine-tuning top N=1200000: R30, L162 :: sha256 c6c4138ee5c86423
\newcommand{\MlstFtTopRthirtyResOverScratch}{1.27\,\ensuremath{\pm}\,0.04}  % experiments/FIGS/data/v2/index.json :: fine-tuning top N=1200000: R30 of each L162 run over the from-scratch mean, SD over runs and standard error of that mean in quadrature :: sha256 c6c4138ee5c86423
\newcommand{\MlstFtTopRthirtyResOverScratchPr}{27\% more rejection}  % experiments/FIGS/data/v2/index.json :: fine-tuning top N=1200000: R30 of each L162 run over the from-scratch mean, SD over runs and standard error of that mean in quadrature :: sha256 c6c4138ee5c86423 :: prose twin: two significant figures, direction in words
\newcommand{\MlstFtTopRthirtyMid}{1{,}780\,\ensuremath{\pm}\,150}  % experiments/FIGS/data/v2/index.json :: fine-tuning top N=1200000: R30, R42_Q1 :: sha256 c6c4138ee5c86423
\newcommand{\MlstFtTopRthirtyMidOverScratch}{1.26\,\ensuremath{\pm}\,0.11}  % experiments/FIGS/data/v2/index.json :: fine-tuning top N=1200000: R30 of each R42_Q1 run over the from-scratch mean, SD over runs and standard error of that mean in quadrature :: sha256 c6c4138ee5c86423
\newcommand{\MlstFtTopRthirtyMidOverScratchPr}{26\% more rejection}  % experiments/FIGS/data/v2/index.json :: fine-tuning top N=1200000: R30 of each R42_Q1 run over the from-scratch mean, SD over runs and standard error of that mean in quadrature :: sha256 c6c4138ee5c86423 :: prose twin: two significant figures, direction in words
\newcommand{\MlstFtTopRthirtyCoarse}{1{,}740\,\ensuremath{\pm}\,140}  % experiments/FIGS/data/v2/index.json :: fine-tuning top N=1200000: R30, R16_Q1 :: sha256 c6c4138ee5c86423
\newcommand{\MlstFtTopRthirtyCoarseOverScratch}{1.23\,\ensuremath{\pm}\,0.10}  % experiments/FIGS/data/v2/index.json :: fine-tuning top N=1200000: R30 of each R16_Q1 run over the from-scratch mean, SD over runs and standard error of that mean in quadrature :: sha256 c6c4138ee5c86423
\newcommand{\MlstFtTopRthirtyCoarseOverScratchPr}{23\% more rejection}  % experiments/FIGS/data/v2/index.json :: fine-tuning top N=1200000: R30 of each R16_Q1 run over the from-scratch mean, SD over runs and standard error of that mean in quadrature :: sha256 c6c4138ee5c86423 :: prose twin: two significant figures, direction in words
\newcommand{\MlstFtTopRthirtyResMass}{1{,}941 and 1{,}590}  % experiments/FIGS/data/v2/index.json :: fine-tuning top N=1200000: R30, L162_MASS :: sha256 c6c4138ee5c86423
\newcommand{\MlstFtTopRthirtyResMassOverScratch}{1.37 and 1.13}  % experiments/FIGS/data/v2/index.json :: fine-tuning top N=1200000: R30 of each L162_MASS run over the from-scratch mean, SD over runs and standard error of that mean in quadrature :: sha256 c6c4138ee5c86423
\newcommand{\MlstFtTopRthirtyResMassOverScratchPr}{at least 13\% more rejection}  % experiments/FIGS/data/v2/index.json :: fine-tuning top N=1200000: R30 of each L162_MASS run over the from-scratch mean, SD over runs and standard error of that mean in quadrature :: sha256 c6c4138ee5c86423 :: prose twin: two significant figures, direction in words
\newcommand{\MlstFtTopRthirtyResMassOverScratchPrMax}{up to 37\% more rejection}  % experiments/FIGS/data/v2/index.json :: fine-tuning top N=1200000: R30 of each L162_MASS run over the from-scratch mean, SD over runs and standard error of that mean in quadrature :: sha256 c6c4138ee5c86423 :: prose twin of the run farther from 1
\newcommand{\MlstFtTopRthirtyCoarseMass}{1{,}683 and 1{,}628}  % experiments/FIGS/data/v2/index.json :: fine-tuning top N=1200000: R30, R16_Q1_MASS :: sha256 c6c4138ee5c86423
\newcommand{\MlstFtTopRthirtyCoarseMassOverScratch}{1.192 and 1.153}  % experiments/FIGS/data/v2/index.json :: fine-tuning top N=1200000: R30 of each R16_Q1_MASS run over the from-scratch mean, SD over runs and standard error of that mean in quadrature :: sha256 c6c4138ee5c86423
\newcommand{\MlstFtTopRthirtyCoarseMassOverScratchPr}{at least 15\% more rejection}  % experiments/FIGS/data/v2/index.json :: fine-tuning top N=1200000: R30 of each R16_Q1_MASS run over the from-scratch mean, SD over runs and standard error of that mean in quadrature :: sha256 c6c4138ee5c86423 :: prose twin: two significant figures, direction in words
\newcommand{\MlstFtTopRthirtyCoarseMassOverScratchPrMax}{up to 19\% more rejection}  % experiments/FIGS/data/v2/index.json :: fine-tuning top N=1200000: R30 of each R16_Q1_MASS run over the from-scratch mean, SD over runs and standard error of that mean in quadrature :: sha256 c6c4138ee5c86423 :: prose twin of the run farther from 1
\newcommand{\MlstFtTopRthirtyResOverFull}{0.997\,\ensuremath{\pm}\,0.032}  % experiments/FIGS/data/v2/index.json :: fine-tuning top N=1200000: R30 of L162 over L188, run by run :: sha256 c6c4138ee5c86423
\newcommand{\MlstFtTopRthirtyResOverFullPr}{equal within 4\%}  % experiments/FIGS/data/v2/index.json :: fine-tuning top N=1200000: R30 of L162 over L188, run by run :: sha256 c6c4138ee5c86423 :: prose twin: two significant figures, direction in words
\newcommand{\MlstFtTopRthirtyMidOverFull}{0.99\,\ensuremath{\pm}\,0.07}  % experiments/FIGS/data/v2/index.json :: fine-tuning top N=1200000: R30 of R42_Q1 over L188, run by run :: sha256 c6c4138ee5c86423
\newcommand{\MlstFtTopRthirtyMidOverFullPr}{equal within 8\%}  % experiments/FIGS/data/v2/index.json :: fine-tuning top N=1200000: R30 of R42_Q1 over L188, run by run :: sha256 c6c4138ee5c86423 :: prose twin: two significant figures, direction in words
\newcommand{\MlstFtTopRthirtyCoarseOverFull}{0.97\,\ensuremath{\pm}\,0.10}  % experiments/FIGS/data/v2/index.json :: fine-tuning top N=1200000: R30 of R16_Q1 over L188, run by run :: sha256 c6c4138ee5c86423
\newcommand{\MlstFtTopRthirtyCoarseOverFullPr}{equal within 10\%}  % experiments/FIGS/data/v2/index.json :: fine-tuning top N=1200000: R30 of R16_Q1 over L188, run by run :: sha256 c6c4138ee5c86423 :: prose twin: two significant figures, direction in words
\newcommand{\MlstFtTopRthirtyResMassOverRes}{1.12 and 0.89}  % experiments/FIGS/data/v2/index.json :: fine-tuning top N=1200000: R30 of L162_MASS over L162, run by run :: sha256 c6c4138ee5c86423
\newcommand{\MlstFtTopRthirtyResMassOverResPr}{equal within 20\%}  % experiments/FIGS/data/v2/index.json :: fine-tuning top N=1200000: R30 of L162_MASS over L162, run by run :: sha256 c6c4138ee5c86423 :: prose twin: two significant figures, direction in words
\newcommand{\MlstFtTopRthirtyResMassOverResPrMax}{equal within 20\%}  % experiments/FIGS/data/v2/index.json :: fine-tuning top N=1200000: R30 of L162_MASS over L162, run by run :: sha256 c6c4138ee5c86423 :: prose twin of the run farther from 1
\newcommand{\MlstFtTopRthirtyCoarseMassOverCoarse}{0.900 and 0.919}  % experiments/FIGS/data/v2/index.json :: fine-tuning top N=1200000: R30 of R16_Q1_MASS over R16_Q1, run by run :: sha256 c6c4138ee5c86423
\newcommand{\MlstFtTopRthirtyCoarseMassOverCoarsePr}{equal within 10\%}  % experiments/FIGS/data/v2/index.json :: fine-tuning top N=1200000: R30 of R16_Q1_MASS over R16_Q1, run by run :: sha256 c6c4138ee5c86423 :: prose twin: two significant figures, direction in words
\newcommand{\MlstFtTopRthirtyCoarseMassOverCoarsePrMax}{equal within 10\%}  % experiments/FIGS/data/v2/index.json :: fine-tuning top N=1200000: R30 of R16_Q1_MASS over R16_Q1, run by run :: sha256 c6c4138ee5c86423 :: prose twin of the run farther from 1
\newcommand{\MlstFtTopRthirtyRandOverFull}{0.97\,\ensuremath{\pm}\,0.04}  % experiments/FIGS/data/v2/index.json :: fine-tuning top N=1200000: R30 of each random partition over the L188 run of its seed :: sha256 c6c4138ee5c86423
\newcommand{\MlstFtTopRthirtyRandOverFullPr}{equal within 4\%}  % experiments/FIGS/data/v2/index.json :: fine-tuning top N=1200000: R30 of each random partition over the L188 run of its seed :: sha256 c6c4138ee5c86423 :: prose twin: two significant figures, direction in words
\newcommand{\MlstFtTopRthirtyRandOverCoarse}{0.97\,\ensuremath{\pm}\,0.09}  % experiments/FIGS/data/v2/index.json :: fine-tuning top N=1200000: R30 of each random partition over the R16_Q1 run of its seed :: sha256 c6c4138ee5c86423
\newcommand{\MlstFtTopRthirtyRandOverCoarsePr}{equal within 9\%}  % experiments/FIGS/data/v2/index.json :: fine-tuning top N=1200000: R30 of each random partition over the R16_Q1 run of its seed :: sha256 c6c4138ee5c86423 :: prose twin: two significant figures, direction in words
\newcommand{\MlstFtTopRthirtyFullOverPart}{0.650\,\ensuremath{\pm}\,0.020}  % experiments/FIGS/data/v2/index.json :: fine-tuning top N=1200000: R30 of each 188-class run over ParT-f.t.'s published value :: sha256 c6c4138ee5c86423
\newcommand{\MlstFtTopRthirtyFullOverPartPr}{35\% less rejection}  % experiments/FIGS/data/v2/index.json :: fine-tuning top N=1200000: R30 of each 188-class run over ParT-f.t.'s published value :: sha256 c6c4138ee5c86423 :: prose twin: two significant figures, direction in words
\newcommand{\MlstFtTopAucFull}{0.9861\,\ensuremath{\pm}\,0.0001}  % experiments/FIGS/data/v2/index.json :: fine-tuning top N=1200000: AUC of the 188-class runs :: sha256 c6c4138ee5c86423
\newcommand{\MlstFtTopAucScratch}{0.9852\,\ensuremath{\pm}\,0.0001}  % experiments/FIGS/data/v2/index.json :: fine-tuning top N=1200000: AUC from scratch :: sha256 c6c4138ee5c86423
\newcommand{\MlstFtTopRthirtyFullRelSd}{3\%}  % experiments/FIGS/data/v2/index.json :: fine-tuning top, full set: relative SD of R30 over the 188-class runs :: sha256 c6c4138ee5c86423
\newcommand{\MlstFtTopRthirtyKin}{1{,}800}  % experiments/FIGS/data/v2/index.json :: fine-tuning top, full set: R30 of the kinematic-only model (one run) :: sha256 c6c4138ee5c86423
\newcommand{\MlstFtTopRthirtyKinOverFull}{1.00}  % experiments/FIGS/data/v2/index.json :: fine-tuning top, full set: R30 of the kinematic-only model over the mean of the 188-class runs :: sha256 c6c4138ee5c86423
\newcommand{\MlstFtTopRthirtyKinOverFullPr}{0.23\% more rejection}  % experiments/FIGS/data/v2/index.json :: fine-tuning top, full set: R30 of the kinematic-only model over the mean of the 188-class runs :: sha256 c6c4138ee5c86423 :: prose twin: two significant figures, direction in words
\newcommand{\MlstFtTopDataEquivJets}{1{,}200}  % experiments/FIGS/data/v2/index.json :: fine-tuning top: from-scratch training jets whose R30 equals the 188-class mean at 1000 jets (interpolated linearly in R30 against log10 N) :: sha256 c6c4138ee5c86423
\newcommand{\MlstFtTopDataEquivJetsEFour}{23{,}000}  % experiments/FIGS/data/v2/index.json :: fine-tuning top: from-scratch training jets whose R30 equals the 188-class mean at 10000 jets (interpolated linearly in R30 against log10 N) :: sha256 c6c4138ee5c86423
\newcommand{\MlstFtQgRthirtyScratchEThree}{49.7\,\ensuremath{\pm}\,2.1}  % experiments/FIGS/data/v2/index.json :: fine-tuning qg N=1000: R30 from scratch, every seed :: sha256 c6c4138ee5c86423
\newcommand{\MlstFtQgRthirtyFullEThree}{47\,\ensuremath{\pm}\,10}  % experiments/FIGS/data/v2/index.json :: fine-tuning qg N=1000: R30, L188 :: sha256 c6c4138ee5c86423
\newcommand{\MlstFtQgRthirtyFullOverScratchEThree}{0.95\,\ensuremath{\pm}\,0.19}  % experiments/FIGS/data/v2/index.json :: fine-tuning qg N=1000: R30 of each L188 run over the from-scratch mean, SD over runs and standard error of that mean in quadrature :: sha256 c6c4138ee5c86423
\newcommand{\MlstFtQgRthirtyFullOverScratchEThreePr}{equal within 20\%}  % experiments/FIGS/data/v2/index.json :: fine-tuning qg N=1000: R30 of each L188 run over the from-scratch mean, SD over runs and standard error of that mean in quadrature :: sha256 c6c4138ee5c86423 :: prose twin: two significant figures, direction in words
\newcommand{\MlstFtQgRthirtyResEThree}{52.4\,\ensuremath{\pm}\,3.1}  % experiments/FIGS/data/v2/index.json :: fine-tuning qg N=1000: R30, L162 :: sha256 c6c4138ee5c86423
\newcommand{\MlstFtQgRthirtyResOverScratchEThree}{1.06\,\ensuremath{\pm}\,0.07}  % experiments/FIGS/data/v2/index.json :: fine-tuning qg N=1000: R30 of each L162 run over the from-scratch mean, SD over runs and standard error of that mean in quadrature :: sha256 c6c4138ee5c86423
\newcommand{\MlstFtQgRthirtyResOverScratchEThreePr}{equal within 7\%}  % experiments/FIGS/data/v2/index.json :: fine-tuning qg N=1000: R30 of each L162 run over the from-scratch mean, SD over runs and standard error of that mean in quadrature :: sha256 c6c4138ee5c86423 :: prose twin: two significant figures, direction in words
\newcommand{\MlstFtQgRthirtyMidEThree}{54\,\ensuremath{\pm}\,6}  % experiments/FIGS/data/v2/index.json :: fine-tuning qg N=1000: R30, R42_Q1 :: sha256 c6c4138ee5c86423
\newcommand{\MlstFtQgRthirtyMidOverScratchEThree}{1.09\,\ensuremath{\pm}\,0.12}  % experiments/FIGS/data/v2/index.json :: fine-tuning qg N=1000: R30 of each R42_Q1 run over the from-scratch mean, SD over runs and standard error of that mean in quadrature :: sha256 c6c4138ee5c86423
\newcommand{\MlstFtQgRthirtyMidOverScratchEThreePr}{equal within 20\%}  % experiments/FIGS/data/v2/index.json :: fine-tuning qg N=1000: R30 of each R42_Q1 run over the from-scratch mean, SD over runs and standard error of that mean in quadrature :: sha256 c6c4138ee5c86423 :: prose twin: two significant figures, direction in words
\newcommand{\MlstFtQgRthirtyCoarseEThree}{55\,\ensuremath{\pm}\,5}  % experiments/FIGS/data/v2/index.json :: fine-tuning qg N=1000: R30, R16_Q1 :: sha256 c6c4138ee5c86423
\newcommand{\MlstFtQgRthirtyCoarseOverScratchEThree}{1.12\,\ensuremath{\pm}\,0.10}  % experiments/FIGS/data/v2/index.json :: fine-tuning qg N=1000: R30 of each R16_Q1 run over the from-scratch mean, SD over runs and standard error of that mean in quadrature :: sha256 c6c4138ee5c86423
\newcommand{\MlstFtQgRthirtyCoarseOverScratchEThreePr}{equal within 20\%}  % experiments/FIGS/data/v2/index.json :: fine-tuning qg N=1000: R30 of each R16_Q1 run over the from-scratch mean, SD over runs and standard error of that mean in quadrature :: sha256 c6c4138ee5c86423 :: prose twin: two significant figures, direction in words
\newcommand{\MlstFtQgRthirtyResMassEThree}{59.3 and 57.1}  % experiments/FIGS/data/v2/index.json :: fine-tuning qg N=1000: R30, L162_MASS :: sha256 c6c4138ee5c86423
\newcommand{\MlstFtQgRthirtyResMassOverScratchEThree}{1.193 and 1.149}  % experiments/FIGS/data/v2/index.json :: fine-tuning qg N=1000: R30 of each L162_MASS run over the from-scratch mean, SD over runs and standard error of that mean in quadrature :: sha256 c6c4138ee5c86423
\newcommand{\MlstFtQgRthirtyResMassOverScratchEThreePr}{at least 15\% more rejection}  % experiments/FIGS/data/v2/index.json :: fine-tuning qg N=1000: R30 of each L162_MASS run over the from-scratch mean, SD over runs and standard error of that mean in quadrature :: sha256 c6c4138ee5c86423 :: prose twin: two significant figures, direction in words
\newcommand{\MlstFtQgRthirtyResMassOverScratchEThreePrMax}{up to 19\% more rejection}  % experiments/FIGS/data/v2/index.json :: fine-tuning qg N=1000: R30 of each L162_MASS run over the from-scratch mean, SD over runs and standard error of that mean in quadrature :: sha256 c6c4138ee5c86423 :: prose twin of the run farther from 1
\newcommand{\MlstFtQgRthirtyCoarseMassEThree}{54.2 and 57.7}  % experiments/FIGS/data/v2/index.json :: fine-tuning qg N=1000: R30, R16_Q1_MASS :: sha256 c6c4138ee5c86423
\newcommand{\MlstFtQgRthirtyCoarseMassOverScratchEThree}{1.091 and 1.161}  % experiments/FIGS/data/v2/index.json :: fine-tuning qg N=1000: R30 of each R16_Q1_MASS run over the from-scratch mean, SD over runs and standard error of that mean in quadrature :: sha256 c6c4138ee5c86423
\newcommand{\MlstFtQgRthirtyCoarseMassOverScratchEThreePr}{at least 9.1\% more rejection}  % experiments/FIGS/data/v2/index.json :: fine-tuning qg N=1000: R30 of each R16_Q1_MASS run over the from-scratch mean, SD over runs and standard error of that mean in quadrature :: sha256 c6c4138ee5c86423 :: prose twin: two significant figures, direction in words
\newcommand{\MlstFtQgRthirtyCoarseMassOverScratchEThreePrMax}{up to 16\% more rejection}  % experiments/FIGS/data/v2/index.json :: fine-tuning qg N=1000: R30 of each R16_Q1_MASS run over the from-scratch mean, SD over runs and standard error of that mean in quadrature :: sha256 c6c4138ee5c86423 :: prose twin of the run farther from 1
\newcommand{\MlstFtQgRthirtyResOverFullEThree}{1.14\,\ensuremath{\pm}\,0.17}  % experiments/FIGS/data/v2/index.json :: fine-tuning qg N=1000: R30 of L162 over L188, run by run :: sha256 c6c4138ee5c86423
\newcommand{\MlstFtQgRthirtyResOverFullEThreePr}{equal within 20\%}  % experiments/FIGS/data/v2/index.json :: fine-tuning qg N=1000: R30 of L162 over L188, run by run :: sha256 c6c4138ee5c86423 :: prose twin: two significant figures, direction in words
\newcommand{\MlstFtQgRthirtyMidOverFullEThree}{1.17\,\ensuremath{\pm}\,0.14}  % experiments/FIGS/data/v2/index.json :: fine-tuning qg N=1000: R30 of R42_Q1 over L188, run by run :: sha256 c6c4138ee5c86423
\newcommand{\MlstFtQgRthirtyMidOverFullEThreePr}{equal within 20\%}  % experiments/FIGS/data/v2/index.json :: fine-tuning qg N=1000: R30 of R42_Q1 over L188, run by run :: sha256 c6c4138ee5c86423 :: prose twin: two significant figures, direction in words
\newcommand{\MlstFtQgRthirtyCoarseOverFullEThree}{1.21\,\ensuremath{\pm}\,0.27}  % experiments/FIGS/data/v2/index.json :: fine-tuning qg N=1000: R30 of R16_Q1 over L188, run by run :: sha256 c6c4138ee5c86423
\newcommand{\MlstFtQgRthirtyCoarseOverFullEThreePr}{equal within 30\%}  % experiments/FIGS/data/v2/index.json :: fine-tuning qg N=1000: R30 of R16_Q1 over L188, run by run :: sha256 c6c4138ee5c86423 :: prose twin: two significant figures, direction in words
\newcommand{\MlstFtQgRthirtyResMassOverResEThree}{1.108 and 1.042}  % experiments/FIGS/data/v2/index.json :: fine-tuning qg N=1000: R30 of L162_MASS over L162, run by run :: sha256 c6c4138ee5c86423
\newcommand{\MlstFtQgRthirtyResMassOverResEThreePr}{equal within 20\%}  % experiments/FIGS/data/v2/index.json :: fine-tuning qg N=1000: R30 of L162_MASS over L162, run by run :: sha256 c6c4138ee5c86423 :: prose twin: two significant figures, direction in words
\newcommand{\MlstFtQgRthirtyResMassOverResEThreePrMax}{equal within 20\%}  % experiments/FIGS/data/v2/index.json :: fine-tuning qg N=1000: R30 of L162_MASS over L162, run by run :: sha256 c6c4138ee5c86423 :: prose twin of the run farther from 1
\newcommand{\MlstFtQgRthirtyCoarseMassOverCoarseEThree}{0.89 and 1.11}  % experiments/FIGS/data/v2/index.json :: fine-tuning qg N=1000: R30 of R16_Q1_MASS over R16_Q1, run by run :: sha256 c6c4138ee5c86423
\newcommand{\MlstFtQgRthirtyCoarseMassOverCoarseEThreePr}{equal within 20\%}  % experiments/FIGS/data/v2/index.json :: fine-tuning qg N=1000: R30 of R16_Q1_MASS over R16_Q1, run by run :: sha256 c6c4138ee5c86423 :: prose twin: two significant figures, direction in words
\newcommand{\MlstFtQgRthirtyCoarseMassOverCoarseEThreePrMax}{equal within 20\%}  % experiments/FIGS/data/v2/index.json :: fine-tuning qg N=1000: R30 of R16_Q1_MASS over R16_Q1, run by run :: sha256 c6c4138ee5c86423 :: prose twin of the run farther from 1
\newcommand{\MlstFtQgRthirtyRandOverFullEThree}{0.91\,\ensuremath{\pm}\,0.07}  % experiments/FIGS/data/v2/index.json :: fine-tuning qg N=1000: R30 of each random partition over the L188 run of its seed :: sha256 c6c4138ee5c86423
\newcommand{\MlstFtQgRthirtyRandOverFullEThreePr}{equal within 10\%}  % experiments/FIGS/data/v2/index.json :: fine-tuning qg N=1000: R30 of each random partition over the L188 run of its seed :: sha256 c6c4138ee5c86423 :: prose twin: two significant figures, direction in words
\newcommand{\MlstFtQgRthirtyRandOverCoarseEThree}{0.87\,\ensuremath{\pm}\,0.14}  % experiments/FIGS/data/v2/index.json :: fine-tuning qg N=1000: R30 of each random partition over the R16_Q1 run of its seed :: sha256 c6c4138ee5c86423
\newcommand{\MlstFtQgRthirtyRandOverCoarseEThreePr}{equal within 20\%}  % experiments/FIGS/data/v2/index.json :: fine-tuning qg N=1000: R30 of each random partition over the R16_Q1 run of its seed :: sha256 c6c4138ee5c86423 :: prose twin: two significant figures, direction in words
\newcommand{\MlstFtQgRthirtyScratchEFour}{66\,\ensuremath{\pm}\,9}  % experiments/FIGS/data/v2/index.json :: fine-tuning qg N=10000: R30 from scratch, every seed :: sha256 c6c4138ee5c86423
\newcommand{\MlstFtQgRthirtyFullEFour}{75\,\ensuremath{\pm}\,7}  % experiments/FIGS/data/v2/index.json :: fine-tuning qg N=10000: R30, L188 :: sha256 c6c4138ee5c86423
\newcommand{\MlstFtQgRthirtyFullOverScratchEFour}{1.14\,\ensuremath{\pm}\,0.14}  % experiments/FIGS/data/v2/index.json :: fine-tuning qg N=10000: R30 of each L188 run over the from-scratch mean, SD over runs and standard error of that mean in quadrature :: sha256 c6c4138ee5c86423
\newcommand{\MlstFtQgRthirtyFullOverScratchEFourPr}{equal within 20\%}  % experiments/FIGS/data/v2/index.json :: fine-tuning qg N=10000: R30 of each L188 run over the from-scratch mean, SD over runs and standard error of that mean in quadrature :: sha256 c6c4138ee5c86423 :: prose twin: two significant figures, direction in words
\newcommand{\MlstFtQgRthirtyResEFour}{76.4\,\ensuremath{\pm}\,0.6}  % experiments/FIGS/data/v2/index.json :: fine-tuning qg N=10000: R30, L162 :: sha256 c6c4138ee5c86423
\newcommand{\MlstFtQgRthirtyResOverScratchEFour}{1.15\,\ensuremath{\pm}\,0.09}  % experiments/FIGS/data/v2/index.json :: fine-tuning qg N=10000: R30 of each L162 run over the from-scratch mean, SD over runs and standard error of that mean in quadrature :: sha256 c6c4138ee5c86423
\newcommand{\MlstFtQgRthirtyResOverScratchEFourPr}{15\% more rejection}  % experiments/FIGS/data/v2/index.json :: fine-tuning qg N=10000: R30 of each L162 run over the from-scratch mean, SD over runs and standard error of that mean in quadrature :: sha256 c6c4138ee5c86423 :: prose twin: two significant figures, direction in words
\newcommand{\MlstFtQgRthirtyMidEFour}{74\,\ensuremath{\pm}\,5}  % experiments/FIGS/data/v2/index.json :: fine-tuning qg N=10000: R30, R42_Q1 :: sha256 c6c4138ee5c86423
\newcommand{\MlstFtQgRthirtyMidOverScratchEFour}{1.12\,\ensuremath{\pm}\,0.11}  % experiments/FIGS/data/v2/index.json :: fine-tuning qg N=10000: R30 of each R42_Q1 run over the from-scratch mean, SD over runs and standard error of that mean in quadrature :: sha256 c6c4138ee5c86423
\newcommand{\MlstFtQgRthirtyMidOverScratchEFourPr}{equal within 20\%}  % experiments/FIGS/data/v2/index.json :: fine-tuning qg N=10000: R30 of each R42_Q1 run over the from-scratch mean, SD over runs and standard error of that mean in quadrature :: sha256 c6c4138ee5c86423 :: prose twin: two significant figures, direction in words
\newcommand{\MlstFtQgRthirtyCoarseEFour}{80.3\,\ensuremath{\pm}\,1.1}  % experiments/FIGS/data/v2/index.json :: fine-tuning qg N=10000: R30, R16_Q1 :: sha256 c6c4138ee5c86423
\newcommand{\MlstFtQgRthirtyCoarseOverScratchEFour}{1.21\,\ensuremath{\pm}\,0.10}  % experiments/FIGS/data/v2/index.json :: fine-tuning qg N=10000: R30 of each R16_Q1 run over the from-scratch mean, SD over runs and standard error of that mean in quadrature :: sha256 c6c4138ee5c86423
\newcommand{\MlstFtQgRthirtyCoarseOverScratchEFourPr}{21\% more rejection}  % experiments/FIGS/data/v2/index.json :: fine-tuning qg N=10000: R30 of each R16_Q1 run over the from-scratch mean, SD over runs and standard error of that mean in quadrature :: sha256 c6c4138ee5c86423 :: prose twin: two significant figures, direction in words
\newcommand{\MlstFtQgRthirtyResMassEFour}{80.1 and 72.6}  % experiments/FIGS/data/v2/index.json :: fine-tuning qg N=10000: R30, L162_MASS :: sha256 c6c4138ee5c86423
\newcommand{\MlstFtQgRthirtyResMassOverScratchEFour}{1.21 and 1.09}  % experiments/FIGS/data/v2/index.json :: fine-tuning qg N=10000: R30 of each L162_MASS run over the from-scratch mean, SD over runs and standard error of that mean in quadrature :: sha256 c6c4138ee5c86423
\newcommand{\MlstFtQgRthirtyResMassOverScratchEFourPr}{at least 9.3\% more rejection}  % experiments/FIGS/data/v2/index.json :: fine-tuning qg N=10000: R30 of each L162_MASS run over the from-scratch mean, SD over runs and standard error of that mean in quadrature :: sha256 c6c4138ee5c86423 :: prose twin: two significant figures, direction in words
\newcommand{\MlstFtQgRthirtyResMassOverScratchEFourPrMax}{up to 21\% more rejection}  % experiments/FIGS/data/v2/index.json :: fine-tuning qg N=10000: R30 of each L162_MASS run over the from-scratch mean, SD over runs and standard error of that mean in quadrature :: sha256 c6c4138ee5c86423 :: prose twin of the run farther from 1
\newcommand{\MlstFtQgRthirtyCoarseMassEFour}{79.5 and 77.9}  % experiments/FIGS/data/v2/index.json :: fine-tuning qg N=10000: R30, R16_Q1_MASS :: sha256 c6c4138ee5c86423
\newcommand{\MlstFtQgRthirtyCoarseMassOverScratchEFour}{1.198 and 1.173}  % experiments/FIGS/data/v2/index.json :: fine-tuning qg N=10000: R30 of each R16_Q1_MASS run over the from-scratch mean, SD over runs and standard error of that mean in quadrature :: sha256 c6c4138ee5c86423
\newcommand{\MlstFtQgRthirtyCoarseMassOverScratchEFourPr}{at least 17\% more rejection}  % experiments/FIGS/data/v2/index.json :: fine-tuning qg N=10000: R30 of each R16_Q1_MASS run over the from-scratch mean, SD over runs and standard error of that mean in quadrature :: sha256 c6c4138ee5c86423 :: prose twin: two significant figures, direction in words
\newcommand{\MlstFtQgRthirtyCoarseMassOverScratchEFourPrMax}{up to 20\% more rejection}  % experiments/FIGS/data/v2/index.json :: fine-tuning qg N=10000: R30 of each R16_Q1_MASS run over the from-scratch mean, SD over runs and standard error of that mean in quadrature :: sha256 c6c4138ee5c86423 :: prose twin of the run farther from 1
\newcommand{\MlstFtQgRthirtyResOverFullEFour}{1.02\,\ensuremath{\pm}\,0.10}  % experiments/FIGS/data/v2/index.json :: fine-tuning qg N=10000: R30 of L162 over L188, run by run :: sha256 c6c4138ee5c86423
\newcommand{\MlstFtQgRthirtyResOverFullEFourPr}{equal within 10\%}  % experiments/FIGS/data/v2/index.json :: fine-tuning qg N=10000: R30 of L162 over L188, run by run :: sha256 c6c4138ee5c86423 :: prose twin: two significant figures, direction in words
\newcommand{\MlstFtQgRthirtyMidOverFullEFour}{0.99\,\ensuremath{\pm}\,0.14}  % experiments/FIGS/data/v2/index.json :: fine-tuning qg N=10000: R30 of R42_Q1 over L188, run by run :: sha256 c6c4138ee5c86423
\newcommand{\MlstFtQgRthirtyMidOverFullEFourPr}{equal within 20\%}  % experiments/FIGS/data/v2/index.json :: fine-tuning qg N=10000: R30 of R42_Q1 over L188, run by run :: sha256 c6c4138ee5c86423 :: prose twin: two significant figures, direction in words
\newcommand{\MlstFtQgRthirtyCoarseOverFullEFour}{1.07\,\ensuremath{\pm}\,0.13}  % experiments/FIGS/data/v2/index.json :: fine-tuning qg N=10000: R30 of R16_Q1 over L188, run by run :: sha256 c6c4138ee5c86423
\newcommand{\MlstFtQgRthirtyCoarseOverFullEFourPr}{equal within 20\%}  % experiments/FIGS/data/v2/index.json :: fine-tuning qg N=10000: R30 of R16_Q1 over L188, run by run :: sha256 c6c4138ee5c86423 :: prose twin: two significant figures, direction in words
\newcommand{\MlstFtQgRthirtyResMassOverResEFour}{1.06 and 0.95}  % experiments/FIGS/data/v2/index.json :: fine-tuning qg N=10000: R30 of L162_MASS over L162, run by run :: sha256 c6c4138ee5c86423
\newcommand{\MlstFtQgRthirtyResMassOverResEFourPr}{equal within 6\%}  % experiments/FIGS/data/v2/index.json :: fine-tuning qg N=10000: R30 of L162_MASS over L162, run by run :: sha256 c6c4138ee5c86423 :: prose twin: two significant figures, direction in words
\newcommand{\MlstFtQgRthirtyResMassOverResEFourPrMax}{equal within 6\%}  % experiments/FIGS/data/v2/index.json :: fine-tuning qg N=10000: R30 of L162_MASS over L162, run by run :: sha256 c6c4138ee5c86423 :: prose twin of the run farther from 1
\newcommand{\MlstFtQgRthirtyCoarseMassOverCoarseEFour}{0.9769 and 0.9696}  % experiments/FIGS/data/v2/index.json :: fine-tuning qg N=10000: R30 of R16_Q1_MASS over R16_Q1, run by run :: sha256 c6c4138ee5c86423
\newcommand{\MlstFtQgRthirtyCoarseMassOverCoarseEFourPr}{equal within 4\%}  % experiments/FIGS/data/v2/index.json :: fine-tuning qg N=10000: R30 of R16_Q1_MASS over R16_Q1, run by run :: sha256 c6c4138ee5c86423 :: prose twin: two significant figures, direction in words
\newcommand{\MlstFtQgRthirtyCoarseMassOverCoarseEFourPrMax}{equal within 4\%}  % experiments/FIGS/data/v2/index.json :: fine-tuning qg N=10000: R30 of R16_Q1_MASS over R16_Q1, run by run :: sha256 c6c4138ee5c86423 :: prose twin of the run farther from 1
\newcommand{\MlstFtQgRthirtyRandOverFullEFour}{1.03\,\ensuremath{\pm}\,0.18}  % experiments/FIGS/data/v2/index.json :: fine-tuning qg N=10000: R30 of each random partition over the L188 run of its seed :: sha256 c6c4138ee5c86423
\newcommand{\MlstFtQgRthirtyRandOverFullEFourPr}{equal within 20\%}  % experiments/FIGS/data/v2/index.json :: fine-tuning qg N=10000: R30 of each random partition over the L188 run of its seed :: sha256 c6c4138ee5c86423 :: prose twin: two significant figures, direction in words
\newcommand{\MlstFtQgRthirtyRandOverCoarseEFour}{0.93\,\ensuremath{\pm}\,0.12}  % experiments/FIGS/data/v2/index.json :: fine-tuning qg N=10000: R30 of each random partition over the R16_Q1 run of its seed :: sha256 c6c4138ee5c86423
\newcommand{\MlstFtQgRthirtyRandOverCoarseEFourPr}{equal within 20\%}  % experiments/FIGS/data/v2/index.json :: fine-tuning qg N=10000: R30 of each random partition over the R16_Q1 run of its seed :: sha256 c6c4138ee5c86423 :: prose twin: two significant figures, direction in words
\newcommand{\MlstFtQgRthirtyScratchEFive}{88.4\,\ensuremath{\pm}\,0.8}  % experiments/FIGS/data/v2/index.json :: fine-tuning qg N=100000: R30 from scratch, every seed :: sha256 c6c4138ee5c86423
\newcommand{\MlstFtQgRthirtyFullEFive}{99.3\,\ensuremath{\pm}\,1.4}  % experiments/FIGS/data/v2/index.json :: fine-tuning qg N=100000: R30, L188 :: sha256 c6c4138ee5c86423
\newcommand{\MlstFtQgRthirtyFullOverScratchEFive}{1.124\,\ensuremath{\pm}\,0.017}  % experiments/FIGS/data/v2/index.json :: fine-tuning qg N=100000: R30 of each L188 run over the from-scratch mean, SD over runs and standard error of that mean in quadrature :: sha256 c6c4138ee5c86423
\newcommand{\MlstFtQgRthirtyFullOverScratchEFivePr}{12\% more rejection}  % experiments/FIGS/data/v2/index.json :: fine-tuning qg N=100000: R30 of each L188 run over the from-scratch mean, SD over runs and standard error of that mean in quadrature :: sha256 c6c4138ee5c86423 :: prose twin: two significant figures, direction in words
\newcommand{\MlstFtQgRthirtyResEFive}{97.9\,\ensuremath{\pm}\,3.3}  % experiments/FIGS/data/v2/index.json :: fine-tuning qg N=100000: R30, L162 :: sha256 c6c4138ee5c86423
\newcommand{\MlstFtQgRthirtyResOverScratchEFive}{1.11\,\ensuremath{\pm}\,0.04}  % experiments/FIGS/data/v2/index.json :: fine-tuning qg N=100000: R30 of each L162 run over the from-scratch mean, SD over runs and standard error of that mean in quadrature :: sha256 c6c4138ee5c86423
\newcommand{\MlstFtQgRthirtyResOverScratchEFivePr}{11\% more rejection}  % experiments/FIGS/data/v2/index.json :: fine-tuning qg N=100000: R30 of each L162 run over the from-scratch mean, SD over runs and standard error of that mean in quadrature :: sha256 c6c4138ee5c86423 :: prose twin: two significant figures, direction in words
\newcommand{\MlstFtQgRthirtyMidEFive}{96.1\,\ensuremath{\pm}\,1.8}  % experiments/FIGS/data/v2/index.json :: fine-tuning qg N=100000: R30, R42_Q1 :: sha256 c6c4138ee5c86423
\newcommand{\MlstFtQgRthirtyMidOverScratchEFive}{1.088\,\ensuremath{\pm}\,0.021}  % experiments/FIGS/data/v2/index.json :: fine-tuning qg N=100000: R30 of each R42_Q1 run over the from-scratch mean, SD over runs and standard error of that mean in quadrature :: sha256 c6c4138ee5c86423
\newcommand{\MlstFtQgRthirtyMidOverScratchEFivePr}{8.8\% more rejection}  % experiments/FIGS/data/v2/index.json :: fine-tuning qg N=100000: R30 of each R42_Q1 run over the from-scratch mean, SD over runs and standard error of that mean in quadrature :: sha256 c6c4138ee5c86423 :: prose twin: two significant figures, direction in words
\newcommand{\MlstFtQgRthirtyCoarseEFive}{96.4\,\ensuremath{\pm}\,1.0}  % experiments/FIGS/data/v2/index.json :: fine-tuning qg N=100000: R30, R16_Q1 :: sha256 c6c4138ee5c86423
\newcommand{\MlstFtQgRthirtyCoarseOverScratchEFive}{1.091\,\ensuremath{\pm}\,0.013}  % experiments/FIGS/data/v2/index.json :: fine-tuning qg N=100000: R30 of each R16_Q1 run over the from-scratch mean, SD over runs and standard error of that mean in quadrature :: sha256 c6c4138ee5c86423
\newcommand{\MlstFtQgRthirtyCoarseOverScratchEFivePr}{9.1\% more rejection}  % experiments/FIGS/data/v2/index.json :: fine-tuning qg N=100000: R30 of each R16_Q1 run over the from-scratch mean, SD over runs and standard error of that mean in quadrature :: sha256 c6c4138ee5c86423 :: prose twin: two significant figures, direction in words
\newcommand{\MlstFtQgRthirtyResMassEFive}{95.6 and 97.6}  % experiments/FIGS/data/v2/index.json :: fine-tuning qg N=100000: R30, L162_MASS :: sha256 c6c4138ee5c86423
\newcommand{\MlstFtQgRthirtyResMassOverScratchEFive}{1.082 and 1.104}  % experiments/FIGS/data/v2/index.json :: fine-tuning qg N=100000: R30 of each L162_MASS run over the from-scratch mean, SD over runs and standard error of that mean in quadrature :: sha256 c6c4138ee5c86423
\newcommand{\MlstFtQgRthirtyResMassOverScratchEFivePr}{at least 8.2\% more rejection}  % experiments/FIGS/data/v2/index.json :: fine-tuning qg N=100000: R30 of each L162_MASS run over the from-scratch mean, SD over runs and standard error of that mean in quadrature :: sha256 c6c4138ee5c86423 :: prose twin: two significant figures, direction in words
\newcommand{\MlstFtQgRthirtyResMassOverScratchEFivePrMax}{up to 10\% more rejection}  % experiments/FIGS/data/v2/index.json :: fine-tuning qg N=100000: R30 of each L162_MASS run over the from-scratch mean, SD over runs and standard error of that mean in quadrature :: sha256 c6c4138ee5c86423 :: prose twin of the run farther from 1
\newcommand{\MlstFtQgRthirtyCoarseMassEFive}{97.85 and 97.94}  % experiments/FIGS/data/v2/index.json :: fine-tuning qg N=100000: R30, R16_Q1_MASS :: sha256 c6c4138ee5c86423
\newcommand{\MlstFtQgRthirtyCoarseMassOverScratchEFive}{1.107 and 1.108}  % experiments/FIGS/data/v2/index.json :: fine-tuning qg N=100000: R30 of each R16_Q1_MASS run over the from-scratch mean, SD over runs and standard error of that mean in quadrature :: sha256 c6c4138ee5c86423
\newcommand{\MlstFtQgRthirtyCoarseMassOverScratchEFivePr}{at least 11\% more rejection}  % experiments/FIGS/data/v2/index.json :: fine-tuning qg N=100000: R30 of each R16_Q1_MASS run over the from-scratch mean, SD over runs and standard error of that mean in quadrature :: sha256 c6c4138ee5c86423 :: prose twin: two significant figures, direction in words
\newcommand{\MlstFtQgRthirtyCoarseMassOverScratchEFivePrMax}{up to 11\% more rejection}  % experiments/FIGS/data/v2/index.json :: fine-tuning qg N=100000: R30 of each R16_Q1_MASS run over the from-scratch mean, SD over runs and standard error of that mean in quadrature :: sha256 c6c4138ee5c86423 :: prose twin of the run farther from 1
\newcommand{\MlstFtQgRthirtyResOverFullEFive}{0.985\,\ensuremath{\pm}\,0.020}  % experiments/FIGS/data/v2/index.json :: fine-tuning qg N=100000: R30 of L162 over L188, run by run :: sha256 c6c4138ee5c86423
\newcommand{\MlstFtQgRthirtyResOverFullEFivePr}{equal within 2\%}  % experiments/FIGS/data/v2/index.json :: fine-tuning qg N=100000: R30 of L162 over L188, run by run :: sha256 c6c4138ee5c86423 :: prose twin: two significant figures, direction in words
\newcommand{\MlstFtQgRthirtyMidOverFullEFive}{0.969\,\ensuremath{\pm}\,0.030}  % experiments/FIGS/data/v2/index.json :: fine-tuning qg N=100000: R30 of R42_Q1 over L188, run by run :: sha256 c6c4138ee5c86423
\newcommand{\MlstFtQgRthirtyMidOverFullEFivePr}{equal within 4\%}  % experiments/FIGS/data/v2/index.json :: fine-tuning qg N=100000: R30 of R42_Q1 over L188, run by run :: sha256 c6c4138ee5c86423 :: prose twin: two significant figures, direction in words
\newcommand{\MlstFtQgRthirtyCoarseOverFullEFive}{0.972\,\ensuremath{\pm}\,0.020}  % experiments/FIGS/data/v2/index.json :: fine-tuning qg N=100000: R30 of R16_Q1 over L188, run by run :: sha256 c6c4138ee5c86423
\newcommand{\MlstFtQgRthirtyCoarseOverFullEFivePr}{equal within 3\%}  % experiments/FIGS/data/v2/index.json :: fine-tuning qg N=100000: R30 of R16_Q1 over L188, run by run :: sha256 c6c4138ee5c86423 :: prose twin: two significant figures, direction in words
\newcommand{\MlstFtQgRthirtyResMassOverResEFive}{0.942 and 1.027}  % experiments/FIGS/data/v2/index.json :: fine-tuning qg N=100000: R30 of L162_MASS over L162, run by run :: sha256 c6c4138ee5c86423
\newcommand{\MlstFtQgRthirtyResMassOverResEFivePr}{equal within 6\%}  % experiments/FIGS/data/v2/index.json :: fine-tuning qg N=100000: R30 of L162_MASS over L162, run by run :: sha256 c6c4138ee5c86423 :: prose twin: two significant figures, direction in words
\newcommand{\MlstFtQgRthirtyResMassOverResEFivePrMax}{equal within 6\%}  % experiments/FIGS/data/v2/index.json :: fine-tuning qg N=100000: R30 of L162_MASS over L162, run by run :: sha256 c6c4138ee5c86423 :: prose twin of the run farther from 1
\newcommand{\MlstFtQgRthirtyCoarseMassOverCoarseEFive}{1.018 and 1.004}  % experiments/FIGS/data/v2/index.json :: fine-tuning qg N=100000: R30 of R16_Q1_MASS over R16_Q1, run by run :: sha256 c6c4138ee5c86423
\newcommand{\MlstFtQgRthirtyCoarseMassOverCoarseEFivePr}{equal within 2\%}  % experiments/FIGS/data/v2/index.json :: fine-tuning qg N=100000: R30 of R16_Q1_MASS over R16_Q1, run by run :: sha256 c6c4138ee5c86423 :: prose twin: two significant figures, direction in words
\newcommand{\MlstFtQgRthirtyCoarseMassOverCoarseEFivePrMax}{equal within 2\%}  % experiments/FIGS/data/v2/index.json :: fine-tuning qg N=100000: R30 of R16_Q1_MASS over R16_Q1, run by run :: sha256 c6c4138ee5c86423 :: prose twin of the run farther from 1
\newcommand{\MlstFtQgRthirtyRandOverFullEFive}{0.986\,\ensuremath{\pm}\,0.025}  % experiments/FIGS/data/v2/index.json :: fine-tuning qg N=100000: R30 of each random partition over the L188 run of its seed :: sha256 c6c4138ee5c86423
\newcommand{\MlstFtQgRthirtyRandOverFullEFivePr}{equal within 3\%}  % experiments/FIGS/data/v2/index.json :: fine-tuning qg N=100000: R30 of each random partition over the L188 run of its seed :: sha256 c6c4138ee5c86423 :: prose twin: two significant figures, direction in words
\newcommand{\MlstFtQgRthirtyRandOverCoarseEFive}{1.009\,\ensuremath{\pm}\,0.020}  % experiments/FIGS/data/v2/index.json :: fine-tuning qg N=100000: R30 of each random partition over the R16_Q1 run of its seed :: sha256 c6c4138ee5c86423
\newcommand{\MlstFtQgRthirtyRandOverCoarseEFivePr}{equal within 3\%}  % experiments/FIGS/data/v2/index.json :: fine-tuning qg N=100000: R30 of each random partition over the R16_Q1 run of its seed :: sha256 c6c4138ee5c86423 :: prose twin: two significant figures, direction in words
\newcommand{\MlstFtQgRthirtyScratch}{104.0\,\ensuremath{\pm}\,1.0}  % experiments/FIGS/data/v2/index.json :: fine-tuning qg N=1600000: R30 from scratch, every seed :: sha256 c6c4138ee5c86423
\newcommand{\MlstFtQgRthirtyFull}{107.3\,\ensuremath{\pm}\,2.2}  % experiments/FIGS/data/v2/index.json :: fine-tuning qg N=1600000: R30, L188 :: sha256 c6c4138ee5c86423
\newcommand{\MlstFtQgRthirtyFullOverScratch}{1.031\,\ensuremath{\pm}\,0.022}  % experiments/FIGS/data/v2/index.json :: fine-tuning qg N=1600000: R30 of each L188 run over the from-scratch mean, SD over runs and standard error of that mean in quadrature :: sha256 c6c4138ee5c86423
\newcommand{\MlstFtQgRthirtyFullOverScratchPr}{equal within 4\%}  % experiments/FIGS/data/v2/index.json :: fine-tuning qg N=1600000: R30 of each L188 run over the from-scratch mean, SD over runs and standard error of that mean in quadrature :: sha256 c6c4138ee5c86423 :: prose twin: two significant figures, direction in words
\newcommand{\MlstFtQgRthirtyRes}{108.0\,\ensuremath{\pm}\,1.4}  % experiments/FIGS/data/v2/index.json :: fine-tuning qg N=1600000: R30, L162 :: sha256 c6c4138ee5c86423
\newcommand{\MlstFtQgRthirtyResOverScratch}{1.038\,\ensuremath{\pm}\,0.015}  % experiments/FIGS/data/v2/index.json :: fine-tuning qg N=1600000: R30 of each L162 run over the from-scratch mean, SD over runs and standard error of that mean in quadrature :: sha256 c6c4138ee5c86423
\newcommand{\MlstFtQgRthirtyResOverScratchPr}{3.8\% more rejection}  % experiments/FIGS/data/v2/index.json :: fine-tuning qg N=1600000: R30 of each L162 run over the from-scratch mean, SD over runs and standard error of that mean in quadrature :: sha256 c6c4138ee5c86423 :: prose twin: two significant figures, direction in words
\newcommand{\MlstFtQgRthirtyMid}{108.3\,\ensuremath{\pm}\,1.0}  % experiments/FIGS/data/v2/index.json :: fine-tuning qg N=1600000: R30, R42_Q1 :: sha256 c6c4138ee5c86423
\newcommand{\MlstFtQgRthirtyMidOverScratch}{1.041\,\ensuremath{\pm}\,0.012}  % experiments/FIGS/data/v2/index.json :: fine-tuning qg N=1600000: R30 of each R42_Q1 run over the from-scratch mean, SD over runs and standard error of that mean in quadrature :: sha256 c6c4138ee5c86423
\newcommand{\MlstFtQgRthirtyMidOverScratchPr}{4.1\% more rejection}  % experiments/FIGS/data/v2/index.json :: fine-tuning qg N=1600000: R30 of each R42_Q1 run over the from-scratch mean, SD over runs and standard error of that mean in quadrature :: sha256 c6c4138ee5c86423 :: prose twin: two significant figures, direction in words
\newcommand{\MlstFtQgRthirtyCoarse}{106.5\,\ensuremath{\pm}\,0.8}  % experiments/FIGS/data/v2/index.json :: fine-tuning qg N=1600000: R30, R16_Q1 :: sha256 c6c4138ee5c86423
\newcommand{\MlstFtQgRthirtyCoarseOverScratch}{1.024\,\ensuremath{\pm}\,0.010}  % experiments/FIGS/data/v2/index.json :: fine-tuning qg N=1600000: R30 of each R16_Q1 run over the from-scratch mean, SD over runs and standard error of that mean in quadrature :: sha256 c6c4138ee5c86423
\newcommand{\MlstFtQgRthirtyCoarseOverScratchPr}{2.4\% more rejection}  % experiments/FIGS/data/v2/index.json :: fine-tuning qg N=1600000: R30 of each R16_Q1 run over the from-scratch mean, SD over runs and standard error of that mean in quadrature :: sha256 c6c4138ee5c86423 :: prose twin: two significant figures, direction in words
\newcommand{\MlstFtQgRthirtyResMass}{107.0 and 106.6}  % experiments/FIGS/data/v2/index.json :: fine-tuning qg N=1600000: R30, L162_MASS :: sha256 c6c4138ee5c86423
\newcommand{\MlstFtQgRthirtyResMassOverScratch}{1.028 and 1.025}  % experiments/FIGS/data/v2/index.json :: fine-tuning qg N=1600000: R30 of each L162_MASS run over the from-scratch mean, SD over runs and standard error of that mean in quadrature :: sha256 c6c4138ee5c86423
\newcommand{\MlstFtQgRthirtyResMassOverScratchPr}{at least 2.5\% more rejection}  % experiments/FIGS/data/v2/index.json :: fine-tuning qg N=1600000: R30 of each L162_MASS run over the from-scratch mean, SD over runs and standard error of that mean in quadrature :: sha256 c6c4138ee5c86423 :: prose twin: two significant figures, direction in words
\newcommand{\MlstFtQgRthirtyResMassOverScratchPrMax}{up to 2.8\% more rejection}  % experiments/FIGS/data/v2/index.json :: fine-tuning qg N=1600000: R30 of each L162_MASS run over the from-scratch mean, SD over runs and standard error of that mean in quadrature :: sha256 c6c4138ee5c86423 :: prose twin of the run farther from 1
\newcommand{\MlstFtQgRthirtyCoarseMass}{104.9 and 105.6}  % experiments/FIGS/data/v2/index.json :: fine-tuning qg N=1600000: R30, R16_Q1_MASS :: sha256 c6c4138ee5c86423
\newcommand{\MlstFtQgRthirtyCoarseMassOverScratch}{1.009 and 1.015}  % experiments/FIGS/data/v2/index.json :: fine-tuning qg N=1600000: R30 of each R16_Q1_MASS run over the from-scratch mean, SD over runs and standard error of that mean in quadrature :: sha256 c6c4138ee5c86423
\newcommand{\MlstFtQgRthirtyCoarseMassOverScratchPr}{at least 0.87\% more rejection}  % experiments/FIGS/data/v2/index.json :: fine-tuning qg N=1600000: R30 of each R16_Q1_MASS run over the from-scratch mean, SD over runs and standard error of that mean in quadrature :: sha256 c6c4138ee5c86423 :: prose twin: two significant figures, direction in words
\newcommand{\MlstFtQgRthirtyCoarseMassOverScratchPrMax}{up to 1.5\% more rejection}  % experiments/FIGS/data/v2/index.json :: fine-tuning qg N=1600000: R30 of each R16_Q1_MASS run over the from-scratch mean, SD over runs and standard error of that mean in quadrature :: sha256 c6c4138ee5c86423 :: prose twin of the run farther from 1
\newcommand{\MlstFtQgRthirtyResOverFull}{1.007\,\ensuremath{\pm}\,0.026}  % experiments/FIGS/data/v2/index.json :: fine-tuning qg N=1600000: R30 of L162 over L188, run by run :: sha256 c6c4138ee5c86423
\newcommand{\MlstFtQgRthirtyResOverFullPr}{equal within 3\%}  % experiments/FIGS/data/v2/index.json :: fine-tuning qg N=1600000: R30 of L162 over L188, run by run :: sha256 c6c4138ee5c86423 :: prose twin: two significant figures, direction in words
\newcommand{\MlstFtQgRthirtyMidOverFull}{1.010\,\ensuremath{\pm}\,0.029}  % experiments/FIGS/data/v2/index.json :: fine-tuning qg N=1600000: R30 of R42_Q1 over L188, run by run :: sha256 c6c4138ee5c86423
\newcommand{\MlstFtQgRthirtyMidOverFullPr}{equal within 3\%}  % experiments/FIGS/data/v2/index.json :: fine-tuning qg N=1600000: R30 of R42_Q1 over L188, run by run :: sha256 c6c4138ee5c86423 :: prose twin: two significant figures, direction in words
\newcommand{\MlstFtQgRthirtyCoarseOverFull}{0.993\,\ensuremath{\pm}\,0.026}  % experiments/FIGS/data/v2/index.json :: fine-tuning qg N=1600000: R30 of R16_Q1 over L188, run by run :: sha256 c6c4138ee5c86423
\newcommand{\MlstFtQgRthirtyCoarseOverFullPr}{equal within 3\%}  % experiments/FIGS/data/v2/index.json :: fine-tuning qg N=1600000: R30 of R16_Q1 over L188, run by run :: sha256 c6c4138ee5c86423 :: prose twin: two significant figures, direction in words
\newcommand{\MlstFtQgRthirtyResMassOverRes}{0.976 and 0.999}  % experiments/FIGS/data/v2/index.json :: fine-tuning qg N=1600000: R30 of L162_MASS over L162, run by run :: sha256 c6c4138ee5c86423
\newcommand{\MlstFtQgRthirtyResMassOverResPr}{equal within 3\%}  % experiments/FIGS/data/v2/index.json :: fine-tuning qg N=1600000: R30 of L162_MASS over L162, run by run :: sha256 c6c4138ee5c86423 :: prose twin: two significant figures, direction in words
\newcommand{\MlstFtQgRthirtyResMassOverResPrMax}{equal within 3\%}  % experiments/FIGS/data/v2/index.json :: fine-tuning qg N=1600000: R30 of L162_MASS over L162, run by run :: sha256 c6c4138ee5c86423 :: prose twin of the run farther from 1
\newcommand{\MlstFtQgRthirtyCoarseMassOverCoarse}{0.9916 and 0.9926}  % experiments/FIGS/data/v2/index.json :: fine-tuning qg N=1600000: R30 of R16_Q1_MASS over R16_Q1, run by run :: sha256 c6c4138ee5c86423
\newcommand{\MlstFtQgRthirtyCoarseMassOverCoarsePr}{equal within 1\%}  % experiments/FIGS/data/v2/index.json :: fine-tuning qg N=1600000: R30 of R16_Q1_MASS over R16_Q1, run by run :: sha256 c6c4138ee5c86423 :: prose twin: two significant figures, direction in words
\newcommand{\MlstFtQgRthirtyCoarseMassOverCoarsePrMax}{equal within 1\%}  % experiments/FIGS/data/v2/index.json :: fine-tuning qg N=1600000: R30 of R16_Q1_MASS over R16_Q1, run by run :: sha256 c6c4138ee5c86423 :: prose twin of the run farther from 1
\newcommand{\MlstFtQgRthirtyRandOverFull}{0.985\,\ensuremath{\pm}\,0.011}  % experiments/FIGS/data/v2/index.json :: fine-tuning qg N=1600000: R30 of each random partition over the L188 run of its seed :: sha256 c6c4138ee5c86423
\newcommand{\MlstFtQgRthirtyRandOverFullPr}{equal within 2\%}  % experiments/FIGS/data/v2/index.json :: fine-tuning qg N=1600000: R30 of each random partition over the L188 run of its seed :: sha256 c6c4138ee5c86423 :: prose twin: two significant figures, direction in words
\newcommand{\MlstFtQgRthirtyRandOverCoarse}{1.008\,\ensuremath{\pm}\,0.009}  % experiments/FIGS/data/v2/index.json :: fine-tuning qg N=1600000: R30 of each random partition over the R16_Q1 run of its seed :: sha256 c6c4138ee5c86423
\newcommand{\MlstFtQgRthirtyRandOverCoarsePr}{equal within 1\%}  % experiments/FIGS/data/v2/index.json :: fine-tuning qg N=1600000: R30 of each random partition over the R16_Q1 run of its seed :: sha256 c6c4138ee5c86423 :: prose twin: two significant figures, direction in words
\newcommand{\MlstFtQgRthirtyFullOverPart}{0.994\,\ensuremath{\pm}\,0.020}  % experiments/FIGS/data/v2/index.json :: fine-tuning qg N=1600000: R30 of each 188-class run over ParT-f.t.'s published value :: sha256 c6c4138ee5c86423
\newcommand{\MlstFtQgRthirtyFullOverPartPr}{equal within 3\%}  % experiments/FIGS/data/v2/index.json :: fine-tuning qg N=1600000: R30 of each 188-class run over ParT-f.t.'s published value :: sha256 c6c4138ee5c86423 :: prose twin: two significant figures, direction in words
\newcommand{\MlstFtQgAucFull}{0.9148\,\ensuremath{\pm}\,0.0001}  % experiments/FIGS/data/v2/index.json :: fine-tuning qg N=1600000: AUC of the 188-class runs :: sha256 c6c4138ee5c86423
\newcommand{\MlstFtQgAucScratch}{0.9121\,\ensuremath{\pm}\,0.0002}  % experiments/FIGS/data/v2/index.json :: fine-tuning qg N=1600000: AUC from scratch :: sha256 c6c4138ee5c86423
\newcommand{\MlstFtQgDataEquivJets}{\ensuremath{<}1{,}000}  % experiments/FIGS/data/v2/index.json :: fine-tuning qg: from-scratch training jets whose R30 equals the 188-class mean at 1000 jets (interpolated linearly in R30 against log10 N) :: sha256 c6c4138ee5c86423
\newcommand{\MlstFtQgDataEquivJetsEFour}{26{,}000}  % experiments/FIGS/data/v2/index.json :: fine-tuning qg: from-scratch training jets whose R30 equals the 188-class mean at 10000 jets (interpolated linearly in R30 against log10 N) :: sha256 c6c4138ee5c86423
\newcommand{\MlstFtQgHerwigRthirtyScratchEThree}{27.0\,\ensuremath{\pm}\,1.3}  % experiments/FIGS/data/v2/index.json :: fine-tuning qg_herwig N=1000: R30 from scratch, every seed :: sha256 c6c4138ee5c86423
\newcommand{\MlstFtQgHerwigRthirtyFullEThree}{21\,\ensuremath{\pm}\,6}  % experiments/FIGS/data/v2/index.json :: fine-tuning qg_herwig N=1000: R30, L188 :: sha256 c6c4138ee5c86423
\newcommand{\MlstFtQgHerwigRthirtyFullOverScratchEThree}{0.77\,\ensuremath{\pm}\,0.22}  % experiments/FIGS/data/v2/index.json :: fine-tuning qg_herwig N=1000: R30 of each L188 run over the from-scratch mean, SD over runs and standard error of that mean in quadrature :: sha256 c6c4138ee5c86423
\newcommand{\MlstFtQgHerwigRthirtyFullOverScratchEThreePr}{equal within 30\%}  % experiments/FIGS/data/v2/index.json :: fine-tuning qg_herwig N=1000: R30 of each L188 run over the from-scratch mean, SD over runs and standard error of that mean in quadrature :: sha256 c6c4138ee5c86423 :: prose twin: two significant figures, direction in words
\newcommand{\MlstFtQgHerwigRthirtyResEThree}{26.2\,\ensuremath{\pm}\,0.9}  % experiments/FIGS/data/v2/index.json :: fine-tuning qg_herwig N=1000: R30, L162 :: sha256 c6c4138ee5c86423
\newcommand{\MlstFtQgHerwigRthirtyResOverScratchEThree}{0.97\,\ensuremath{\pm}\,0.04}  % experiments/FIGS/data/v2/index.json :: fine-tuning qg_herwig N=1000: R30 of each L162 run over the from-scratch mean, SD over runs and standard error of that mean in quadrature :: sha256 c6c4138ee5c86423
\newcommand{\MlstFtQgHerwigRthirtyResOverScratchEThreePr}{equal within 4\%}  % experiments/FIGS/data/v2/index.json :: fine-tuning qg_herwig N=1000: R30 of each L162 run over the from-scratch mean, SD over runs and standard error of that mean in quadrature :: sha256 c6c4138ee5c86423 :: prose twin: two significant figures, direction in words
\newcommand{\MlstFtQgHerwigRthirtyMidEThree}{26.9\,\ensuremath{\pm}\,2.7}  % experiments/FIGS/data/v2/index.json :: fine-tuning qg_herwig N=1000: R30, R42_Q1 :: sha256 c6c4138ee5c86423
\newcommand{\MlstFtQgHerwigRthirtyMidOverScratchEThree}{1.00\,\ensuremath{\pm}\,0.10}  % experiments/FIGS/data/v2/index.json :: fine-tuning qg_herwig N=1000: R30 of each R42_Q1 run over the from-scratch mean, SD over runs and standard error of that mean in quadrature :: sha256 c6c4138ee5c86423
\newcommand{\MlstFtQgHerwigRthirtyMidOverScratchEThreePr}{equal within 10\%}  % experiments/FIGS/data/v2/index.json :: fine-tuning qg_herwig N=1000: R30 of each R42_Q1 run over the from-scratch mean, SD over runs and standard error of that mean in quadrature :: sha256 c6c4138ee5c86423 :: prose twin: two significant figures, direction in words
\newcommand{\MlstFtQgHerwigRthirtyCoarseEThree}{25.2\,\ensuremath{\pm}\,2.5}  % experiments/FIGS/data/v2/index.json :: fine-tuning qg_herwig N=1000: R30, R16_Q1 :: sha256 c6c4138ee5c86423
\newcommand{\MlstFtQgHerwigRthirtyCoarseOverScratchEThree}{0.94\,\ensuremath{\pm}\,0.10}  % experiments/FIGS/data/v2/index.json :: fine-tuning qg_herwig N=1000: R30 of each R16_Q1 run over the from-scratch mean, SD over runs and standard error of that mean in quadrature :: sha256 c6c4138ee5c86423
\newcommand{\MlstFtQgHerwigRthirtyCoarseOverScratchEThreePr}{equal within 10\%}  % experiments/FIGS/data/v2/index.json :: fine-tuning qg_herwig N=1000: R30 of each R16_Q1 run over the from-scratch mean, SD over runs and standard error of that mean in quadrature :: sha256 c6c4138ee5c86423 :: prose twin: two significant figures, direction in words
\newcommand{\MlstFtQgHerwigRthirtyResMassEThree}{28.0 and 29.0}  % experiments/FIGS/data/v2/index.json :: fine-tuning qg_herwig N=1000: R30, L162_MASS :: sha256 c6c4138ee5c86423
\newcommand{\MlstFtQgHerwigRthirtyResMassOverScratchEThree}{1.038 and 1.076}  % experiments/FIGS/data/v2/index.json :: fine-tuning qg_herwig N=1000: R30 of each L162_MASS run over the from-scratch mean, SD over runs and standard error of that mean in quadrature :: sha256 c6c4138ee5c86423
\newcommand{\MlstFtQgHerwigRthirtyResMassOverScratchEThreePr}{at least 3.8\% more rejection}  % experiments/FIGS/data/v2/index.json :: fine-tuning qg_herwig N=1000: R30 of each L162_MASS run over the from-scratch mean, SD over runs and standard error of that mean in quadrature :: sha256 c6c4138ee5c86423 :: prose twin: two significant figures, direction in words
\newcommand{\MlstFtQgHerwigRthirtyResMassOverScratchEThreePrMax}{up to 7.6\% more rejection}  % experiments/FIGS/data/v2/index.json :: fine-tuning qg_herwig N=1000: R30 of each L162_MASS run over the from-scratch mean, SD over runs and standard error of that mean in quadrature :: sha256 c6c4138ee5c86423 :: prose twin of the run farther from 1
\newcommand{\MlstFtQgHerwigRthirtyCoarseMassEThree}{25.8 and 27.4}  % experiments/FIGS/data/v2/index.json :: fine-tuning qg_herwig N=1000: R30, R16_Q1_MASS :: sha256 c6c4138ee5c86423
\newcommand{\MlstFtQgHerwigRthirtyCoarseMassOverScratchEThree}{0.955 and 1.015}  % experiments/FIGS/data/v2/index.json :: fine-tuning qg_herwig N=1000: R30 of each R16_Q1_MASS run over the from-scratch mean, SD over runs and standard error of that mean in quadrature :: sha256 c6c4138ee5c86423
\newcommand{\MlstFtQgHerwigRthirtyCoarseMassOverScratchEThreePr}{equal within 5\%}  % experiments/FIGS/data/v2/index.json :: fine-tuning qg_herwig N=1000: R30 of each R16_Q1_MASS run over the from-scratch mean, SD over runs and standard error of that mean in quadrature :: sha256 c6c4138ee5c86423 :: prose twin: two significant figures, direction in words
\newcommand{\MlstFtQgHerwigRthirtyCoarseMassOverScratchEThreePrMax}{equal within 5\%}  % experiments/FIGS/data/v2/index.json :: fine-tuning qg_herwig N=1000: R30 of each R16_Q1_MASS run over the from-scratch mean, SD over runs and standard error of that mean in quadrature :: sha256 c6c4138ee5c86423 :: prose twin of the run farther from 1
\newcommand{\MlstFtQgHerwigRthirtyResOverFullEThree}{1.33\,\ensuremath{\pm}\,0.33}  % experiments/FIGS/data/v2/index.json :: fine-tuning qg_herwig N=1000: R30 of L162 over L188, run by run :: sha256 c6c4138ee5c86423
\newcommand{\MlstFtQgHerwigRthirtyResOverFullEThreePr}{equal within 40\%}  % experiments/FIGS/data/v2/index.json :: fine-tuning qg_herwig N=1000: R30 of L162 over L188, run by run :: sha256 c6c4138ee5c86423 :: prose twin: two significant figures, direction in words
\newcommand{\MlstFtQgHerwigRthirtyMidOverFullEThree}{1.36\,\ensuremath{\pm}\,0.31}  % experiments/FIGS/data/v2/index.json :: fine-tuning qg_herwig N=1000: R30 of R42_Q1 over L188, run by run :: sha256 c6c4138ee5c86423
\newcommand{\MlstFtQgHerwigRthirtyMidOverFullEThreePr}{equal within 40\%}  % experiments/FIGS/data/v2/index.json :: fine-tuning qg_herwig N=1000: R30 of R42_Q1 over L188, run by run :: sha256 c6c4138ee5c86423 :: prose twin: two significant figures, direction in words
\newcommand{\MlstFtQgHerwigRthirtyCoarseOverFullEThree}{1.3\,\ensuremath{\pm}\,0.4}  % experiments/FIGS/data/v2/index.json :: fine-tuning qg_herwig N=1000: R30 of R16_Q1 over L188, run by run :: sha256 c6c4138ee5c86423
\newcommand{\MlstFtQgHerwigRthirtyCoarseOverFullEThreePr}{equal within 50\%}  % experiments/FIGS/data/v2/index.json :: fine-tuning qg_herwig N=1000: R30 of R16_Q1 over L188, run by run :: sha256 c6c4138ee5c86423 :: prose twin: two significant figures, direction in words
\newcommand{\MlstFtQgHerwigRthirtyResMassOverResEThree}{1.062 and 1.073}  % experiments/FIGS/data/v2/index.json :: fine-tuning qg_herwig N=1000: R30 of L162_MASS over L162, run by run :: sha256 c6c4138ee5c86423
\newcommand{\MlstFtQgHerwigRthirtyResMassOverResEThreePr}{equal within 8\%}  % experiments/FIGS/data/v2/index.json :: fine-tuning qg_herwig N=1000: R30 of L162_MASS over L162, run by run :: sha256 c6c4138ee5c86423 :: prose twin: two significant figures, direction in words
\newcommand{\MlstFtQgHerwigRthirtyResMassOverResEThreePrMax}{equal within 8\%}  % experiments/FIGS/data/v2/index.json :: fine-tuning qg_herwig N=1000: R30 of L162_MASS over L162, run by run :: sha256 c6c4138ee5c86423 :: prose twin of the run farther from 1
\newcommand{\MlstFtQgHerwigRthirtyCoarseMassOverCoarseEThree}{0.93 and 1.20}  % experiments/FIGS/data/v2/index.json :: fine-tuning qg_herwig N=1000: R30 of R16_Q1_MASS over R16_Q1, run by run :: sha256 c6c4138ee5c86423
\newcommand{\MlstFtQgHerwigRthirtyCoarseMassOverCoarseEThreePr}{equal within 30\%}  % experiments/FIGS/data/v2/index.json :: fine-tuning qg_herwig N=1000: R30 of R16_Q1_MASS over R16_Q1, run by run :: sha256 c6c4138ee5c86423 :: prose twin: two significant figures, direction in words
\newcommand{\MlstFtQgHerwigRthirtyCoarseMassOverCoarseEThreePrMax}{equal within 30\%}  % experiments/FIGS/data/v2/index.json :: fine-tuning qg_herwig N=1000: R30 of R16_Q1_MASS over R16_Q1, run by run :: sha256 c6c4138ee5c86423 :: prose twin of the run farther from 1
\newcommand{\MlstFtQgHerwigRthirtyRandOverFullEThree}{0.99\,\ensuremath{\pm}\,0.20}  % experiments/FIGS/data/v2/index.json :: fine-tuning qg_herwig N=1000: R30 of each random partition over the L188 run of its seed :: sha256 c6c4138ee5c86423
\newcommand{\MlstFtQgHerwigRthirtyRandOverFullEThreePr}{equal within 30\%}  % experiments/FIGS/data/v2/index.json :: fine-tuning qg_herwig N=1000: R30 of each random partition over the L188 run of its seed :: sha256 c6c4138ee5c86423 :: prose twin: two significant figures, direction in words
\newcommand{\MlstFtQgHerwigRthirtyRandOverCoarseEThree}{0.95\,\ensuremath{\pm}\,0.13}  % experiments/FIGS/data/v2/index.json :: fine-tuning qg_herwig N=1000: R30 of each random partition over the R16_Q1 run of its seed :: sha256 c6c4138ee5c86423
\newcommand{\MlstFtQgHerwigRthirtyRandOverCoarseEThreePr}{equal within 20\%}  % experiments/FIGS/data/v2/index.json :: fine-tuning qg_herwig N=1000: R30 of each random partition over the R16_Q1 run of its seed :: sha256 c6c4138ee5c86423 :: prose twin: two significant figures, direction in words
\newcommand{\MlstFtQgHerwigRthirtyScratchEFour}{31\,\ensuremath{\pm}\,7}  % experiments/FIGS/data/v2/index.json :: fine-tuning qg_herwig N=10000: R30 from scratch, every seed :: sha256 c6c4138ee5c86423
\newcommand{\MlstFtQgHerwigRthirtyFullEFour}{27.9\,\ensuremath{\pm}\,2.8}  % experiments/FIGS/data/v2/index.json :: fine-tuning qg_herwig N=10000: R30, L188 :: sha256 c6c4138ee5c86423
\newcommand{\MlstFtQgHerwigRthirtyFullOverScratchEFour}{0.89\,\ensuremath{\pm}\,0.14}  % experiments/FIGS/data/v2/index.json :: fine-tuning qg_herwig N=10000: R30 of each L188 run over the from-scratch mean, SD over runs and standard error of that mean in quadrature :: sha256 c6c4138ee5c86423
\newcommand{\MlstFtQgHerwigRthirtyFullOverScratchEFourPr}{equal within 20\%}  % experiments/FIGS/data/v2/index.json :: fine-tuning qg_herwig N=10000: R30 of each L188 run over the from-scratch mean, SD over runs and standard error of that mean in quadrature :: sha256 c6c4138ee5c86423 :: prose twin: two significant figures, direction in words
\newcommand{\MlstFtQgHerwigRthirtyResEFour}{27.4\,\ensuremath{\pm}\,0.5}  % experiments/FIGS/data/v2/index.json :: fine-tuning qg_herwig N=10000: R30, L162 :: sha256 c6c4138ee5c86423
\newcommand{\MlstFtQgHerwigRthirtyResOverScratchEFour}{0.87\,\ensuremath{\pm}\,0.11}  % experiments/FIGS/data/v2/index.json :: fine-tuning qg_herwig N=10000: R30 of each L162 run over the from-scratch mean, SD over runs and standard error of that mean in quadrature :: sha256 c6c4138ee5c86423
\newcommand{\MlstFtQgHerwigRthirtyResOverScratchEFourPr}{13\% less rejection}  % experiments/FIGS/data/v2/index.json :: fine-tuning qg_herwig N=10000: R30 of each L162 run over the from-scratch mean, SD over runs and standard error of that mean in quadrature :: sha256 c6c4138ee5c86423 :: prose twin: two significant figures, direction in words
\newcommand{\MlstFtQgHerwigRthirtyMidEFour}{28.4\,\ensuremath{\pm}\,2.1}  % experiments/FIGS/data/v2/index.json :: fine-tuning qg_herwig N=10000: R30, R42_Q1 :: sha256 c6c4138ee5c86423
\newcommand{\MlstFtQgHerwigRthirtyMidOverScratchEFour}{0.91\,\ensuremath{\pm}\,0.13}  % experiments/FIGS/data/v2/index.json :: fine-tuning qg_herwig N=10000: R30 of each R42_Q1 run over the from-scratch mean, SD over runs and standard error of that mean in quadrature :: sha256 c6c4138ee5c86423
\newcommand{\MlstFtQgHerwigRthirtyMidOverScratchEFourPr}{equal within 10\%}  % experiments/FIGS/data/v2/index.json :: fine-tuning qg_herwig N=10000: R30 of each R42_Q1 run over the from-scratch mean, SD over runs and standard error of that mean in quadrature :: sha256 c6c4138ee5c86423 :: prose twin: two significant figures, direction in words
\newcommand{\MlstFtQgHerwigRthirtyCoarseEFour}{32.7\,\ensuremath{\pm}\,1.8}  % experiments/FIGS/data/v2/index.json :: fine-tuning qg_herwig N=10000: R30, R16_Q1 :: sha256 c6c4138ee5c86423
\newcommand{\MlstFtQgHerwigRthirtyCoarseOverScratchEFour}{1.05\,\ensuremath{\pm}\,0.14}  % experiments/FIGS/data/v2/index.json :: fine-tuning qg_herwig N=10000: R30 of each R16_Q1 run over the from-scratch mean, SD over runs and standard error of that mean in quadrature :: sha256 c6c4138ee5c86423
\newcommand{\MlstFtQgHerwigRthirtyCoarseOverScratchEFourPr}{equal within 6\%}  % experiments/FIGS/data/v2/index.json :: fine-tuning qg_herwig N=10000: R30 of each R16_Q1 run over the from-scratch mean, SD over runs and standard error of that mean in quadrature :: sha256 c6c4138ee5c86423 :: prose twin: two significant figures, direction in words
\newcommand{\MlstFtQgHerwigRthirtyResMassEFour}{34.3 and 28.5}  % experiments/FIGS/data/v2/index.json :: fine-tuning qg_herwig N=10000: R30, L162_MASS :: sha256 c6c4138ee5c86423
\newcommand{\MlstFtQgHerwigRthirtyResMassOverScratchEFour}{1.09 and 0.91}  % experiments/FIGS/data/v2/index.json :: fine-tuning qg_herwig N=10000: R30 of each L162_MASS run over the from-scratch mean, SD over runs and standard error of that mean in quadrature :: sha256 c6c4138ee5c86423
\newcommand{\MlstFtQgHerwigRthirtyResMassOverScratchEFourPr}{equal within 10\%}  % experiments/FIGS/data/v2/index.json :: fine-tuning qg_herwig N=10000: R30 of each L162_MASS run over the from-scratch mean, SD over runs and standard error of that mean in quadrature :: sha256 c6c4138ee5c86423 :: prose twin: two significant figures, direction in words
\newcommand{\MlstFtQgHerwigRthirtyResMassOverScratchEFourPrMax}{equal within 10\%}  % experiments/FIGS/data/v2/index.json :: fine-tuning qg_herwig N=10000: R30 of each L162_MASS run over the from-scratch mean, SD over runs and standard error of that mean in quadrature :: sha256 c6c4138ee5c86423 :: prose twin of the run farther from 1
\newcommand{\MlstFtQgHerwigRthirtyCoarseMassEFour}{34.4 and 31.2}  % experiments/FIGS/data/v2/index.json :: fine-tuning qg_herwig N=10000: R30, R16_Q1_MASS :: sha256 c6c4138ee5c86423
\newcommand{\MlstFtQgHerwigRthirtyCoarseMassOverScratchEFour}{1.10 and 1.00}  % experiments/FIGS/data/v2/index.json :: fine-tuning qg_herwig N=10000: R30 of each R16_Q1_MASS run over the from-scratch mean, SD over runs and standard error of that mean in quadrature :: sha256 c6c4138ee5c86423
\newcommand{\MlstFtQgHerwigRthirtyCoarseMassOverScratchEFourPr}{equal within 20\%}  % experiments/FIGS/data/v2/index.json :: fine-tuning qg_herwig N=10000: R30 of each R16_Q1_MASS run over the from-scratch mean, SD over runs and standard error of that mean in quadrature :: sha256 c6c4138ee5c86423 :: prose twin: two significant figures, direction in words
\newcommand{\MlstFtQgHerwigRthirtyCoarseMassOverScratchEFourPrMax}{equal within 20\%}  % experiments/FIGS/data/v2/index.json :: fine-tuning qg_herwig N=10000: R30 of each R16_Q1_MASS run over the from-scratch mean, SD over runs and standard error of that mean in quadrature :: sha256 c6c4138ee5c86423 :: prose twin of the run farther from 1
\newcommand{\MlstFtQgHerwigRthirtyResOverFullEFour}{0.99\,\ensuremath{\pm}\,0.12}  % experiments/FIGS/data/v2/index.json :: fine-tuning qg_herwig N=10000: R30 of L162 over L188, run by run :: sha256 c6c4138ee5c86423
\newcommand{\MlstFtQgHerwigRthirtyResOverFullEFourPr}{equal within 20\%}  % experiments/FIGS/data/v2/index.json :: fine-tuning qg_herwig N=10000: R30 of L162 over L188, run by run :: sha256 c6c4138ee5c86423 :: prose twin: two significant figures, direction in words
\newcommand{\MlstFtQgHerwigRthirtyMidOverFullEFour}{1.03\,\ensuremath{\pm}\,0.18}  % experiments/FIGS/data/v2/index.json :: fine-tuning qg_herwig N=10000: R30 of R42_Q1 over L188, run by run :: sha256 c6c4138ee5c86423
\newcommand{\MlstFtQgHerwigRthirtyMidOverFullEFourPr}{equal within 20\%}  % experiments/FIGS/data/v2/index.json :: fine-tuning qg_herwig N=10000: R30 of R42_Q1 over L188, run by run :: sha256 c6c4138ee5c86423 :: prose twin: two significant figures, direction in words
\newcommand{\MlstFtQgHerwigRthirtyCoarseOverFullEFour}{1.18\,\ensuremath{\pm}\,0.08}  % experiments/FIGS/data/v2/index.json :: fine-tuning qg_herwig N=10000: R30 of R16_Q1 over L188, run by run :: sha256 c6c4138ee5c86423
\newcommand{\MlstFtQgHerwigRthirtyCoarseOverFullEFourPr}{18\% more rejection}  % experiments/FIGS/data/v2/index.json :: fine-tuning qg_herwig N=10000: R30 of R16_Q1 over L188, run by run :: sha256 c6c4138ee5c86423 :: prose twin: two significant figures, direction in words
\newcommand{\MlstFtQgHerwigRthirtyResMassOverResEFour}{1.23 and 1.03}  % experiments/FIGS/data/v2/index.json :: fine-tuning qg_herwig N=10000: R30 of L162_MASS over L162, run by run :: sha256 c6c4138ee5c86423
\newcommand{\MlstFtQgHerwigRthirtyResMassOverResEFourPr}{equal within 30\%}  % experiments/FIGS/data/v2/index.json :: fine-tuning qg_herwig N=10000: R30 of L162_MASS over L162, run by run :: sha256 c6c4138ee5c86423 :: prose twin: two significant figures, direction in words
\newcommand{\MlstFtQgHerwigRthirtyResMassOverResEFourPrMax}{equal within 30\%}  % experiments/FIGS/data/v2/index.json :: fine-tuning qg_herwig N=10000: R30 of L162_MASS over L162, run by run :: sha256 c6c4138ee5c86423 :: prose twin of the run farther from 1
\newcommand{\MlstFtQgHerwigRthirtyCoarseMassOverCoarseEFour}{1.12 and 0.91}  % experiments/FIGS/data/v2/index.json :: fine-tuning qg_herwig N=10000: R30 of R16_Q1_MASS over R16_Q1, run by run :: sha256 c6c4138ee5c86423
\newcommand{\MlstFtQgHerwigRthirtyCoarseMassOverCoarseEFourPr}{equal within 20\%}  % experiments/FIGS/data/v2/index.json :: fine-tuning qg_herwig N=10000: R30 of R16_Q1_MASS over R16_Q1, run by run :: sha256 c6c4138ee5c86423 :: prose twin: two significant figures, direction in words
\newcommand{\MlstFtQgHerwigRthirtyCoarseMassOverCoarseEFourPrMax}{equal within 20\%}  % experiments/FIGS/data/v2/index.json :: fine-tuning qg_herwig N=10000: R30 of R16_Q1_MASS over R16_Q1, run by run :: sha256 c6c4138ee5c86423 :: prose twin of the run farther from 1
\newcommand{\MlstFtQgHerwigRthirtyRandOverFullEFour}{1.13\,\ensuremath{\pm}\,0.18}  % experiments/FIGS/data/v2/index.json :: fine-tuning qg_herwig N=10000: R30 of each random partition over the L188 run of its seed :: sha256 c6c4138ee5c86423
\newcommand{\MlstFtQgHerwigRthirtyRandOverFullEFourPr}{equal within 20\%}  % experiments/FIGS/data/v2/index.json :: fine-tuning qg_herwig N=10000: R30 of each random partition over the L188 run of its seed :: sha256 c6c4138ee5c86423 :: prose twin: two significant figures, direction in words
\newcommand{\MlstFtQgHerwigRthirtyRandOverCoarseEFour}{0.93\,\ensuremath{\pm}\,0.14}  % experiments/FIGS/data/v2/index.json :: fine-tuning qg_herwig N=10000: R30 of each random partition over the R16_Q1 run of its seed :: sha256 c6c4138ee5c86423
\newcommand{\MlstFtQgHerwigRthirtyRandOverCoarseEFourPr}{equal within 20\%}  % experiments/FIGS/data/v2/index.json :: fine-tuning qg_herwig N=10000: R30 of each random partition over the R16_Q1 run of its seed :: sha256 c6c4138ee5c86423 :: prose twin: two significant figures, direction in words
\newcommand{\MlstFtQgHerwigRthirtyScratchEFive}{36.7\,\ensuremath{\pm}\,0.9}  % experiments/FIGS/data/v2/index.json :: fine-tuning qg_herwig N=100000: R30 from scratch, every seed :: sha256 c6c4138ee5c86423
\newcommand{\MlstFtQgHerwigRthirtyFullEFive}{34.3\,\ensuremath{\pm}\,1.0}  % experiments/FIGS/data/v2/index.json :: fine-tuning qg_herwig N=100000: R30, L188 :: sha256 c6c4138ee5c86423
\newcommand{\MlstFtQgHerwigRthirtyFullOverScratchEFive}{0.935\,\ensuremath{\pm}\,0.031}  % experiments/FIGS/data/v2/index.json :: fine-tuning qg_herwig N=100000: R30 of each L188 run over the from-scratch mean, SD over runs and standard error of that mean in quadrature :: sha256 c6c4138ee5c86423
\newcommand{\MlstFtQgHerwigRthirtyFullOverScratchEFivePr}{6.5\% less rejection}  % experiments/FIGS/data/v2/index.json :: fine-tuning qg_herwig N=100000: R30 of each L188 run over the from-scratch mean, SD over runs and standard error of that mean in quadrature :: sha256 c6c4138ee5c86423 :: prose twin: two significant figures, direction in words
\newcommand{\MlstFtQgHerwigRthirtyResEFive}{35.5\,\ensuremath{\pm}\,0.5}  % experiments/FIGS/data/v2/index.json :: fine-tuning qg_herwig N=100000: R30, L162 :: sha256 c6c4138ee5c86423
\newcommand{\MlstFtQgHerwigRthirtyResOverScratchEFive}{0.966\,\ensuremath{\pm}\,0.019}  % experiments/FIGS/data/v2/index.json :: fine-tuning qg_herwig N=100000: R30 of each L162 run over the from-scratch mean, SD over runs and standard error of that mean in quadrature :: sha256 c6c4138ee5c86423
\newcommand{\MlstFtQgHerwigRthirtyResOverScratchEFivePr}{3.4\% less rejection}  % experiments/FIGS/data/v2/index.json :: fine-tuning qg_herwig N=100000: R30 of each L162 run over the from-scratch mean, SD over runs and standard error of that mean in quadrature :: sha256 c6c4138ee5c86423 :: prose twin: two significant figures, direction in words
\newcommand{\MlstFtQgHerwigRthirtyMidEFive}{35.5\,\ensuremath{\pm}\,2.1}  % experiments/FIGS/data/v2/index.json :: fine-tuning qg_herwig N=100000: R30, R42_Q1 :: sha256 c6c4138ee5c86423
\newcommand{\MlstFtQgHerwigRthirtyMidOverScratchEFive}{0.97\,\ensuremath{\pm}\,0.06}  % experiments/FIGS/data/v2/index.json :: fine-tuning qg_herwig N=100000: R30 of each R42_Q1 run over the from-scratch mean, SD over runs and standard error of that mean in quadrature :: sha256 c6c4138ee5c86423
\newcommand{\MlstFtQgHerwigRthirtyMidOverScratchEFivePr}{equal within 6\%}  % experiments/FIGS/data/v2/index.json :: fine-tuning qg_herwig N=100000: R30 of each R42_Q1 run over the from-scratch mean, SD over runs and standard error of that mean in quadrature :: sha256 c6c4138ee5c86423 :: prose twin: two significant figures, direction in words
\newcommand{\MlstFtQgHerwigRthirtyCoarseEFive}{34.9\,\ensuremath{\pm}\,1.1}  % experiments/FIGS/data/v2/index.json :: fine-tuning qg_herwig N=100000: R30, R16_Q1 :: sha256 c6c4138ee5c86423
\newcommand{\MlstFtQgHerwigRthirtyCoarseOverScratchEFive}{0.950\,\ensuremath{\pm}\,0.033}  % experiments/FIGS/data/v2/index.json :: fine-tuning qg_herwig N=100000: R30 of each R16_Q1 run over the from-scratch mean, SD over runs and standard error of that mean in quadrature :: sha256 c6c4138ee5c86423
\newcommand{\MlstFtQgHerwigRthirtyCoarseOverScratchEFivePr}{equal within 6\%}  % experiments/FIGS/data/v2/index.json :: fine-tuning qg_herwig N=100000: R30 of each R16_Q1 run over the from-scratch mean, SD over runs and standard error of that mean in quadrature :: sha256 c6c4138ee5c86423 :: prose twin: two significant figures, direction in words
\newcommand{\MlstFtQgHerwigRthirtyResMassEFive}{34.0 and 36.5}  % experiments/FIGS/data/v2/index.json :: fine-tuning qg_herwig N=100000: R30, L162_MASS :: sha256 c6c4138ee5c86423
\newcommand{\MlstFtQgHerwigRthirtyResMassOverScratchEFive}{0.925 and 0.994}  % experiments/FIGS/data/v2/index.json :: fine-tuning qg_herwig N=100000: R30 of each L162_MASS run over the from-scratch mean, SD over runs and standard error of that mean in quadrature :: sha256 c6c4138ee5c86423
\newcommand{\MlstFtQgHerwigRthirtyResMassOverScratchEFivePr}{at least 0.63\% less rejection}  % experiments/FIGS/data/v2/index.json :: fine-tuning qg_herwig N=100000: R30 of each L162_MASS run over the from-scratch mean, SD over runs and standard error of that mean in quadrature :: sha256 c6c4138ee5c86423 :: prose twin: two significant figures, direction in words
\newcommand{\MlstFtQgHerwigRthirtyResMassOverScratchEFivePrMax}{up to 7.5\% less rejection}  % experiments/FIGS/data/v2/index.json :: fine-tuning qg_herwig N=100000: R30 of each L162_MASS run over the from-scratch mean, SD over runs and standard error of that mean in quadrature :: sha256 c6c4138ee5c86423 :: prose twin of the run farther from 1
\newcommand{\MlstFtQgHerwigRthirtyCoarseMassEFive}{35.9 and 39.0}  % experiments/FIGS/data/v2/index.json :: fine-tuning qg_herwig N=100000: R30, R16_Q1_MASS :: sha256 c6c4138ee5c86423
\newcommand{\MlstFtQgHerwigRthirtyCoarseMassOverScratchEFive}{0.979 and 1.064}  % experiments/FIGS/data/v2/index.json :: fine-tuning qg_herwig N=100000: R30 of each R16_Q1_MASS run over the from-scratch mean, SD over runs and standard error of that mean in quadrature :: sha256 c6c4138ee5c86423
\newcommand{\MlstFtQgHerwigRthirtyCoarseMassOverScratchEFivePr}{equal within 7\%}  % experiments/FIGS/data/v2/index.json :: fine-tuning qg_herwig N=100000: R30 of each R16_Q1_MASS run over the from-scratch mean, SD over runs and standard error of that mean in quadrature :: sha256 c6c4138ee5c86423 :: prose twin: two significant figures, direction in words
\newcommand{\MlstFtQgHerwigRthirtyCoarseMassOverScratchEFivePrMax}{equal within 7\%}  % experiments/FIGS/data/v2/index.json :: fine-tuning qg_herwig N=100000: R30 of each R16_Q1_MASS run over the from-scratch mean, SD over runs and standard error of that mean in quadrature :: sha256 c6c4138ee5c86423 :: prose twin of the run farther from 1
\newcommand{\MlstFtQgHerwigRthirtyResOverFullEFive}{1.03\,\ensuremath{\pm}\,0.04}  % experiments/FIGS/data/v2/index.json :: fine-tuning qg_herwig N=100000: R30 of L162 over L188, run by run :: sha256 c6c4138ee5c86423
\newcommand{\MlstFtQgHerwigRthirtyResOverFullEFivePr}{equal within 5\%}  % experiments/FIGS/data/v2/index.json :: fine-tuning qg_herwig N=100000: R30 of L162 over L188, run by run :: sha256 c6c4138ee5c86423 :: prose twin: two significant figures, direction in words
\newcommand{\MlstFtQgHerwigRthirtyMidOverFullEFive}{1.03\,\ensuremath{\pm}\,0.05}  % experiments/FIGS/data/v2/index.json :: fine-tuning qg_herwig N=100000: R30 of R42_Q1 over L188, run by run :: sha256 c6c4138ee5c86423
\newcommand{\MlstFtQgHerwigRthirtyMidOverFullEFivePr}{equal within 6\%}  % experiments/FIGS/data/v2/index.json :: fine-tuning qg_herwig N=100000: R30 of R42_Q1 over L188, run by run :: sha256 c6c4138ee5c86423 :: prose twin: two significant figures, direction in words
\newcommand{\MlstFtQgHerwigRthirtyCoarseOverFullEFive}{1.02\,\ensuremath{\pm}\,0.06}  % experiments/FIGS/data/v2/index.json :: fine-tuning qg_herwig N=100000: R30 of R16_Q1 over L188, run by run :: sha256 c6c4138ee5c86423
\newcommand{\MlstFtQgHerwigRthirtyCoarseOverFullEFivePr}{equal within 7\%}  % experiments/FIGS/data/v2/index.json :: fine-tuning qg_herwig N=100000: R30 of R16_Q1 over L188, run by run :: sha256 c6c4138ee5c86423 :: prose twin: two significant figures, direction in words
\newcommand{\MlstFtQgHerwigRthirtyResMassOverResEFive}{0.971 and 1.028}  % experiments/FIGS/data/v2/index.json :: fine-tuning qg_herwig N=100000: R30 of L162_MASS over L162, run by run :: sha256 c6c4138ee5c86423
\newcommand{\MlstFtQgHerwigRthirtyResMassOverResEFivePr}{equal within 3\%}  % experiments/FIGS/data/v2/index.json :: fine-tuning qg_herwig N=100000: R30 of L162_MASS over L162, run by run :: sha256 c6c4138ee5c86423 :: prose twin: two significant figures, direction in words
\newcommand{\MlstFtQgHerwigRthirtyResMassOverResEFivePrMax}{equal within 3\%}  % experiments/FIGS/data/v2/index.json :: fine-tuning qg_herwig N=100000: R30 of L162_MASS over L162, run by run :: sha256 c6c4138ee5c86423 :: prose twin of the run farther from 1
\newcommand{\MlstFtQgHerwigRthirtyCoarseMassOverCoarseEFive}{1.064 and 1.121}  % experiments/FIGS/data/v2/index.json :: fine-tuning qg_herwig N=100000: R30 of R16_Q1_MASS over R16_Q1, run by run :: sha256 c6c4138ee5c86423
\newcommand{\MlstFtQgHerwigRthirtyCoarseMassOverCoarseEFivePr}{equal within 20\%}  % experiments/FIGS/data/v2/index.json :: fine-tuning qg_herwig N=100000: R30 of R16_Q1_MASS over R16_Q1, run by run :: sha256 c6c4138ee5c86423 :: prose twin: two significant figures, direction in words
\newcommand{\MlstFtQgHerwigRthirtyCoarseMassOverCoarseEFivePrMax}{equal within 20\%}  % experiments/FIGS/data/v2/index.json :: fine-tuning qg_herwig N=100000: R30 of R16_Q1_MASS over R16_Q1, run by run :: sha256 c6c4138ee5c86423 :: prose twin of the run farther from 1
\newcommand{\MlstFtQgHerwigRthirtyRandOverFullEFive}{1.01\,\ensuremath{\pm}\,0.05}  % experiments/FIGS/data/v2/index.json :: fine-tuning qg_herwig N=100000: R30 of each random partition over the L188 run of its seed :: sha256 c6c4138ee5c86423
\newcommand{\MlstFtQgHerwigRthirtyRandOverFullEFivePr}{equal within 5\%}  % experiments/FIGS/data/v2/index.json :: fine-tuning qg_herwig N=100000: R30 of each random partition over the L188 run of its seed :: sha256 c6c4138ee5c86423 :: prose twin: two significant figures, direction in words
\newcommand{\MlstFtQgHerwigRthirtyRandOverCoarseEFive}{1.009\,\ensuremath{\pm}\,0.024}  % experiments/FIGS/data/v2/index.json :: fine-tuning qg_herwig N=100000: R30 of each random partition over the R16_Q1 run of its seed :: sha256 c6c4138ee5c86423
\newcommand{\MlstFtQgHerwigRthirtyRandOverCoarseEFivePr}{equal within 3\%}  % experiments/FIGS/data/v2/index.json :: fine-tuning qg_herwig N=100000: R30 of each random partition over the R16_Q1 run of its seed :: sha256 c6c4138ee5c86423 :: prose twin: two significant figures, direction in words
\newcommand{\MlstFtQgHerwigRthirtyScratch}{35.5\,\ensuremath{\pm}\,1.2}  % experiments/FIGS/data/v2/index.json :: fine-tuning qg_herwig N=1600000: R30 from scratch, every seed :: sha256 c6c4138ee5c86423
\newcommand{\MlstFtQgHerwigRthirtyFull}{35.2\,\ensuremath{\pm}\,2.1}  % experiments/FIGS/data/v2/index.json :: fine-tuning qg_herwig N=1600000: R30, L188 :: sha256 c6c4138ee5c86423
\newcommand{\MlstFtQgHerwigRthirtyFullOverScratch}{0.99\,\ensuremath{\pm}\,0.06}  % experiments/FIGS/data/v2/index.json :: fine-tuning qg_herwig N=1600000: R30 of each L188 run over the from-scratch mean, SD over runs and standard error of that mean in quadrature :: sha256 c6c4138ee5c86423
\newcommand{\MlstFtQgHerwigRthirtyFullOverScratchPr}{equal within 6\%}  % experiments/FIGS/data/v2/index.json :: fine-tuning qg_herwig N=1600000: R30 of each L188 run over the from-scratch mean, SD over runs and standard error of that mean in quadrature :: sha256 c6c4138ee5c86423 :: prose twin: two significant figures, direction in words
\newcommand{\MlstFtQgHerwigRthirtyRes}{34.6\,\ensuremath{\pm}\,2.0}  % experiments/FIGS/data/v2/index.json :: fine-tuning qg_herwig N=1600000: R30, L162 :: sha256 c6c4138ee5c86423
\newcommand{\MlstFtQgHerwigRthirtyResOverScratch}{0.97\,\ensuremath{\pm}\,0.06}  % experiments/FIGS/data/v2/index.json :: fine-tuning qg_herwig N=1600000: R30 of each L162 run over the from-scratch mean, SD over runs and standard error of that mean in quadrature :: sha256 c6c4138ee5c86423
\newcommand{\MlstFtQgHerwigRthirtyResOverScratchPr}{equal within 6\%}  % experiments/FIGS/data/v2/index.json :: fine-tuning qg_herwig N=1600000: R30 of each L162 run over the from-scratch mean, SD over runs and standard error of that mean in quadrature :: sha256 c6c4138ee5c86423 :: prose twin: two significant figures, direction in words
\newcommand{\MlstFtQgHerwigRthirtyMid}{36.1\,\ensuremath{\pm}\,2.7}  % experiments/FIGS/data/v2/index.json :: fine-tuning qg_herwig N=1600000: R30, R42_Q1 :: sha256 c6c4138ee5c86423
\newcommand{\MlstFtQgHerwigRthirtyMidOverScratch}{1.02\,\ensuremath{\pm}\,0.08}  % experiments/FIGS/data/v2/index.json :: fine-tuning qg_herwig N=1600000: R30 of each R42_Q1 run over the from-scratch mean, SD over runs and standard error of that mean in quadrature :: sha256 c6c4138ee5c86423
\newcommand{\MlstFtQgHerwigRthirtyMidOverScratchPr}{equal within 8\%}  % experiments/FIGS/data/v2/index.json :: fine-tuning qg_herwig N=1600000: R30 of each R42_Q1 run over the from-scratch mean, SD over runs and standard error of that mean in quadrature :: sha256 c6c4138ee5c86423 :: prose twin: two significant figures, direction in words
\newcommand{\MlstFtQgHerwigRthirtyCoarse}{36.9\,\ensuremath{\pm}\,0.9}  % experiments/FIGS/data/v2/index.json :: fine-tuning qg_herwig N=1600000: R30, R16_Q1 :: sha256 c6c4138ee5c86423
\newcommand{\MlstFtQgHerwigRthirtyCoarseOverScratch}{1.040\,\ensuremath{\pm}\,0.033}  % experiments/FIGS/data/v2/index.json :: fine-tuning qg_herwig N=1600000: R30 of each R16_Q1 run over the from-scratch mean, SD over runs and standard error of that mean in quadrature :: sha256 c6c4138ee5c86423
\newcommand{\MlstFtQgHerwigRthirtyCoarseOverScratchPr}{equal within 4\%}  % experiments/FIGS/data/v2/index.json :: fine-tuning qg_herwig N=1600000: R30 of each R16_Q1 run over the from-scratch mean, SD over runs and standard error of that mean in quadrature :: sha256 c6c4138ee5c86423 :: prose twin: two significant figures, direction in words
\newcommand{\MlstFtQgHerwigRthirtyResMass}{37.8 and 34.4}  % experiments/FIGS/data/v2/index.json :: fine-tuning qg_herwig N=1600000: R30, L162_MASS :: sha256 c6c4138ee5c86423
\newcommand{\MlstFtQgHerwigRthirtyResMassOverScratch}{1.065 and 0.970}  % experiments/FIGS/data/v2/index.json :: fine-tuning qg_herwig N=1600000: R30 of each L162_MASS run over the from-scratch mean, SD over runs and standard error of that mean in quadrature :: sha256 c6c4138ee5c86423
\newcommand{\MlstFtQgHerwigRthirtyResMassOverScratchPr}{equal within 7\%}  % experiments/FIGS/data/v2/index.json :: fine-tuning qg_herwig N=1600000: R30 of each L162_MASS run over the from-scratch mean, SD over runs and standard error of that mean in quadrature :: sha256 c6c4138ee5c86423 :: prose twin: two significant figures, direction in words
\newcommand{\MlstFtQgHerwigRthirtyResMassOverScratchPrMax}{equal within 7\%}  % experiments/FIGS/data/v2/index.json :: fine-tuning qg_herwig N=1600000: R30 of each L162_MASS run over the from-scratch mean, SD over runs and standard error of that mean in quadrature :: sha256 c6c4138ee5c86423 :: prose twin of the run farther from 1
\newcommand{\MlstFtQgHerwigRthirtyCoarseMass}{37.7 and 32.9}  % experiments/FIGS/data/v2/index.json :: fine-tuning qg_herwig N=1600000: R30, R16_Q1_MASS :: sha256 c6c4138ee5c86423
\newcommand{\MlstFtQgHerwigRthirtyCoarseMassOverScratch}{1.06 and 0.93}  % experiments/FIGS/data/v2/index.json :: fine-tuning qg_herwig N=1600000: R30 of each R16_Q1_MASS run over the from-scratch mean, SD over runs and standard error of that mean in quadrature :: sha256 c6c4138ee5c86423
\newcommand{\MlstFtQgHerwigRthirtyCoarseMassOverScratchPr}{equal within 8\%}  % experiments/FIGS/data/v2/index.json :: fine-tuning qg_herwig N=1600000: R30 of each R16_Q1_MASS run over the from-scratch mean, SD over runs and standard error of that mean in quadrature :: sha256 c6c4138ee5c86423 :: prose twin: two significant figures, direction in words
\newcommand{\MlstFtQgHerwigRthirtyCoarseMassOverScratchPrMax}{equal within 8\%}  % experiments/FIGS/data/v2/index.json :: fine-tuning qg_herwig N=1600000: R30 of each R16_Q1_MASS run over the from-scratch mean, SD over runs and standard error of that mean in quadrature :: sha256 c6c4138ee5c86423 :: prose twin of the run farther from 1
\newcommand{\MlstFtQgHerwigRthirtyResOverFull}{0.9832\,\ensuremath{\pm}\,0.0022}  % experiments/FIGS/data/v2/index.json :: fine-tuning qg_herwig N=1600000: R30 of L162 over L188, run by run :: sha256 c6c4138ee5c86423
\newcommand{\MlstFtQgHerwigRthirtyResOverFullPr}{1.7\% less rejection}  % experiments/FIGS/data/v2/index.json :: fine-tuning qg_herwig N=1600000: R30 of L162 over L188, run by run :: sha256 c6c4138ee5c86423 :: prose twin: two significant figures, direction in words
\newcommand{\MlstFtQgHerwigRthirtyMidOverFull}{1.03\,\ensuremath{\pm}\,0.05}  % experiments/FIGS/data/v2/index.json :: fine-tuning qg_herwig N=1600000: R30 of R42_Q1 over L188, run by run :: sha256 c6c4138ee5c86423
\newcommand{\MlstFtQgHerwigRthirtyMidOverFullPr}{equal within 6\%}  % experiments/FIGS/data/v2/index.json :: fine-tuning qg_herwig N=1600000: R30 of R42_Q1 over L188, run by run :: sha256 c6c4138ee5c86423 :: prose twin: two significant figures, direction in words
\newcommand{\MlstFtQgHerwigRthirtyCoarseOverFull}{1.05\,\ensuremath{\pm}\,0.06}  % experiments/FIGS/data/v2/index.json :: fine-tuning qg_herwig N=1600000: R30 of R16_Q1 over L188, run by run :: sha256 c6c4138ee5c86423
\newcommand{\MlstFtQgHerwigRthirtyCoarseOverFullPr}{equal within 6\%}  % experiments/FIGS/data/v2/index.json :: fine-tuning qg_herwig N=1600000: R30 of R16_Q1 over L188, run by run :: sha256 c6c4138ee5c86423 :: prose twin: two significant figures, direction in words
\newcommand{\MlstFtQgHerwigRthirtyResMassOverRes}{1.14 and 0.93}  % experiments/FIGS/data/v2/index.json :: fine-tuning qg_herwig N=1600000: R30 of L162_MASS over L162, run by run :: sha256 c6c4138ee5c86423
\newcommand{\MlstFtQgHerwigRthirtyResMassOverResPr}{equal within 20\%}  % experiments/FIGS/data/v2/index.json :: fine-tuning qg_herwig N=1600000: R30 of L162_MASS over L162, run by run :: sha256 c6c4138ee5c86423 :: prose twin: two significant figures, direction in words
\newcommand{\MlstFtQgHerwigRthirtyResMassOverResPrMax}{equal within 20\%}  % experiments/FIGS/data/v2/index.json :: fine-tuning qg_herwig N=1600000: R30 of L162_MASS over L162, run by run :: sha256 c6c4138ee5c86423 :: prose twin of the run farther from 1
\newcommand{\MlstFtQgHerwigRthirtyCoarseMassOverCoarse}{1.01 and 0.88}  % experiments/FIGS/data/v2/index.json :: fine-tuning qg_herwig N=1600000: R30 of R16_Q1_MASS over R16_Q1, run by run :: sha256 c6c4138ee5c86423
\newcommand{\MlstFtQgHerwigRthirtyCoarseMassOverCoarsePr}{equal within 20\%}  % experiments/FIGS/data/v2/index.json :: fine-tuning qg_herwig N=1600000: R30 of R16_Q1_MASS over R16_Q1, run by run :: sha256 c6c4138ee5c86423 :: prose twin: two significant figures, direction in words
\newcommand{\MlstFtQgHerwigRthirtyCoarseMassOverCoarsePrMax}{equal within 20\%}  % experiments/FIGS/data/v2/index.json :: fine-tuning qg_herwig N=1600000: R30 of R16_Q1_MASS over R16_Q1, run by run :: sha256 c6c4138ee5c86423 :: prose twin of the run farther from 1
\newcommand{\MlstFtQgHerwigRthirtyRandOverFull}{0.98\,\ensuremath{\pm}\,0.04}  % experiments/FIGS/data/v2/index.json :: fine-tuning qg_herwig N=1600000: R30 of each random partition over the L188 run of its seed :: sha256 c6c4138ee5c86423
\newcommand{\MlstFtQgHerwigRthirtyRandOverFullPr}{equal within 5\%}  % experiments/FIGS/data/v2/index.json :: fine-tuning qg_herwig N=1600000: R30 of each random partition over the L188 run of its seed :: sha256 c6c4138ee5c86423 :: prose twin: two significant figures, direction in words
\newcommand{\MlstFtQgHerwigRthirtyRandOverCoarse}{0.94\,\ensuremath{\pm}\,0.04}  % experiments/FIGS/data/v2/index.json :: fine-tuning qg_herwig N=1600000: R30 of each random partition over the R16_Q1 run of its seed :: sha256 c6c4138ee5c86423
\newcommand{\MlstFtQgHerwigRthirtyRandOverCoarsePr}{equal within 6\%}  % experiments/FIGS/data/v2/index.json :: fine-tuning qg_herwig N=1600000: R30 of each random partition over the R16_Q1 run of its seed :: sha256 c6c4138ee5c86423 :: prose twin: two significant figures, direction in words
\newcommand{\MlstFtQgRthirtyHerwigOverPythiaFull}{0.328\,\ensuremath{\pm}\,0.015}  % experiments/FIGS/data/v2/index.json :: fine-tuning q/g N=1600000: R30 on Herwig over R30 on Pythia, each run of L188 :: sha256 c6c4138ee5c86423
\newcommand{\MlstFtQgRthirtyHerwigOverPythiaFullPr}{0.33 times the rejection}  % experiments/FIGS/data/v2/index.json :: fine-tuning q/g N=1600000: R30 on Herwig over R30 on Pythia, each run of L188 :: sha256 c6c4138ee5c86423 :: prose twin: two significant figures, direction in words
\newcommand{\MlstFtQgRthirtyHerwigOverPythiaRes}{0.320\,\ensuremath{\pm}\,0.023}  % experiments/FIGS/data/v2/index.json :: fine-tuning q/g N=1600000: R30 on Herwig over R30 on Pythia, each run of L162 :: sha256 c6c4138ee5c86423
\newcommand{\MlstFtQgRthirtyHerwigOverPythiaResPr}{0.32 times the rejection}  % experiments/FIGS/data/v2/index.json :: fine-tuning q/g N=1600000: R30 on Herwig over R30 on Pythia, each run of L162 :: sha256 c6c4138ee5c86423 :: prose twin: two significant figures, direction in words
\newcommand{\MlstFtQgRthirtyHerwigOverPythiaMid}{0.334\,\ensuremath{\pm}\,0.023}  % experiments/FIGS/data/v2/index.json :: fine-tuning q/g N=1600000: R30 on Herwig over R30 on Pythia, each run of R42_Q1 :: sha256 c6c4138ee5c86423
\newcommand{\MlstFtQgRthirtyHerwigOverPythiaMidPr}{0.33 times the rejection}  % experiments/FIGS/data/v2/index.json :: fine-tuning q/g N=1600000: R30 on Herwig over R30 on Pythia, each run of R42_Q1 :: sha256 c6c4138ee5c86423 :: prose twin: two significant figures, direction in words
\newcommand{\MlstFtQgRthirtyHerwigOverPythiaCoarse}{0.346\,\ensuremath{\pm}\,0.011}  % experiments/FIGS/data/v2/index.json :: fine-tuning q/g N=1600000: R30 on Herwig over R30 on Pythia, each run of R16_Q1 :: sha256 c6c4138ee5c86423
\newcommand{\MlstFtQgRthirtyHerwigOverPythiaCoarsePr}{0.35 times the rejection}  % experiments/FIGS/data/v2/index.json :: fine-tuning q/g N=1600000: R30 on Herwig over R30 on Pythia, each run of R16_Q1 :: sha256 c6c4138ee5c86423 :: prose twin: two significant figures, direction in words
\newcommand{\MlstFtQgRthirtyHerwigOverPythiaResMass}{0.353 and 0.323}  % experiments/FIGS/data/v2/index.json :: fine-tuning q/g N=1600000: R30 on Herwig over R30 on Pythia, each run of L162_MASS :: sha256 c6c4138ee5c86423
\newcommand{\MlstFtQgRthirtyHerwigOverPythiaResMassPr}{at most 0.35 times the rejection}  % experiments/FIGS/data/v2/index.json :: fine-tuning q/g N=1600000: R30 on Herwig over R30 on Pythia, each run of L162_MASS :: sha256 c6c4138ee5c86423 :: prose twin: two significant figures, direction in words
\newcommand{\MlstFtQgRthirtyHerwigOverPythiaResMassPrMax}{down to 0.32 times the rejection}  % experiments/FIGS/data/v2/index.json :: fine-tuning q/g N=1600000: R30 on Herwig over R30 on Pythia, each run of L162_MASS :: sha256 c6c4138ee5c86423 :: prose twin of the run farther from 1
\newcommand{\MlstFtQgRthirtyHerwigOverPythiaCoarseMass}{0.359 and 0.311}  % experiments/FIGS/data/v2/index.json :: fine-tuning q/g N=1600000: R30 on Herwig over R30 on Pythia, each run of R16_Q1_MASS :: sha256 c6c4138ee5c86423
\newcommand{\MlstFtQgRthirtyHerwigOverPythiaCoarseMassPr}{at most 0.36 times the rejection}  % experiments/FIGS/data/v2/index.json :: fine-tuning q/g N=1600000: R30 on Herwig over R30 on Pythia, each run of R16_Q1_MASS :: sha256 c6c4138ee5c86423 :: prose twin: two significant figures, direction in words
\newcommand{\MlstFtQgRthirtyHerwigOverPythiaCoarseMassPrMax}{down to 0.31 times the rejection}  % experiments/FIGS/data/v2/index.json :: fine-tuning q/g N=1600000: R30 on Herwig over R30 on Pythia, each run of R16_Q1_MASS :: sha256 c6c4138ee5c86423 :: prose twin of the run farther from 1
\newcommand{\MlstFtQgRthirtyHerwigOverPythiaScratch}{0.341\,\ensuremath{\pm}\,0.009}  % experiments/FIGS/data/v2/index.json :: fine-tuning q/g N=1600000: R30 on Herwig over R30 on Pythia, each from-scratch seed :: sha256 c6c4138ee5c86423
\newcommand{\MlstFtQgRthirtyHerwigOverPythiaScratchPr}{0.34 times the rejection}  % experiments/FIGS/data/v2/index.json :: fine-tuning q/g N=1600000: R30 on Herwig over R30 on Pythia, each from-scratch seed :: sha256 c6c4138ee5c86423 :: prose twin: two significant figures, direction in words
\newcommand{\MlstTopCheckFullLong}{\pending{recipe check}}  % experiments/FIGS/data/v2/index.json :: the recipe check is not read out yet :: sha256 c6c4138ee5c86423
\newcommand{\MlstTopCheckScratchLong}{\pending{recipe check}}  % experiments/FIGS/data/v2/index.json :: the recipe check is not read out yet :: sha256 c6c4138ee5c86423
\newcommand{\MlstTopCheckFullShort}{\pending{recipe check}}  % experiments/FIGS/data/v2/index.json :: the recipe check is not read out yet :: sha256 c6c4138ee5c86423
\newcommand{\MlstTopCheckScratchShort}{\pending{recipe check}}  % experiments/FIGS/data/v2/index.json :: the recipe check is not read out yet :: sha256 c6c4138ee5c86423
\newcommand{\MlstLpJcOneGapFracEFour}{56.3\,\ensuremath{\pm}\,1.1\%}  % experiments/FIGS/data/v2/index.json :: jc1 N=10000: (probe - scratch) / (fine-tuned - scratch) accuracy, 188-class runs, per cent :: sha256 c6c4138ee5c86423
\newcommand{\MlstLpJcOneGapFracEFourPr}{56\%}  % experiments/FIGS/data/v2/index.json :: jc1 N=10000: (probe - scratch) / (fine-tuned - scratch) accuracy, 188-class runs, per cent :: sha256 c6c4138ee5c86423 :: prose twin, two significant figures
\newcommand{\MlstLpJcOneGapFracEFive}{\ensuremath{-}15\,\ensuremath{\pm}\,5\%}  % experiments/FIGS/data/v2/index.json :: jc1 N=100000: (probe - scratch) / (fine-tuned - scratch) accuracy, 188-class runs, per cent :: sha256 c6c4138ee5c86423
\newcommand{\MlstLpJcOneGapFracEFivePr}{\ensuremath{-}15\%}  % experiments/FIGS/data/v2/index.json :: jc1 N=100000: (probe - scratch) / (fine-tuned - scratch) accuracy, 188-class runs, per cent :: sha256 c6c4138ee5c86423 :: prose twin, two significant figures
\newcommand{\MlstLpJcTwoGapFracEFour}{96.2\,\ensuremath{\pm}\,0.6\%}  % experiments/FIGS/data/v2/index.json :: jc2_188 N=10000: (probe - scratch) / (fine-tuned - scratch) accuracy, 188-class runs, per cent :: sha256 c6c4138ee5c86423
\newcommand{\MlstLpJcTwoGapFracEFourPr}{96\%}  % experiments/FIGS/data/v2/index.json :: jc2_188 N=10000: (probe - scratch) / (fine-tuned - scratch) accuracy, 188-class runs, per cent :: sha256 c6c4138ee5c86423 :: prose twin, two significant figures
\newcommand{\MlstLpJcTwoGapFracEFive}{102.2\,\ensuremath{\pm}\,1.0\%}  % experiments/FIGS/data/v2/index.json :: jc2_188 N=100000: (probe - scratch) / (fine-tuned - scratch) accuracy, 188-class runs, per cent :: sha256 c6c4138ee5c86423
\newcommand{\MlstLpJcTwoGapFracEFivePr}{100\%}  % experiments/FIGS/data/v2/index.json :: jc2_188 N=100000: (probe - scratch) / (fine-tuned - scratch) accuracy, 188-class runs, per cent :: sha256 c6c4138ee5c86423 :: prose twin, two significant figures
\newcommand{\XbxMx}{\pending{samples not in}}  % experiments/XBS/xbx_report.py :: pending :: sha256 cb6dffe774ff5b1e
\newcommand{\XbxWindowLo}{\pending{samples not in}}  % experiments/XBS/xbx_report.py :: pending :: sha256 cb6dffe774ff5b1e
\newcommand{\XbxWindowHi}{\pending{samples not in}}  % experiments/XBS/xbx_report.py :: pending :: sha256 cb6dffe774ff5b1e
\newcommand{\XbbNTrainPerClass}{\pending{samples not in}}  % experiments/XBS/xbx_report.py :: pending :: sha256 cb6dffe774ff5b1e
\newcommand{\XbbRepetition}{\pending{samples not in}}  % experiments/XBS/xbx_report.py :: pending :: sha256 cb6dffe774ff5b1e
\newcommand{\XbbNEvalSignal}{\pending{samples not in}}  % experiments/XBS/xbx_report.py :: pending :: sha256 cb6dffe774ff5b1e
\newcommand{\XbbNEvalQcd}{\pending{samples not in}}  % experiments/XBS/xbx_report.py :: pending :: sha256 cb6dffe774ff5b1e
\newcommand{\XbsNTrainPerClass}{\pending{samples not in}}  % experiments/XBS/xbx_report.py :: pending :: sha256 cb6dffe774ff5b1e
\newcommand{\XbsRepetition}{\pending{samples not in}}  % experiments/XBS/xbx_report.py :: pending :: sha256 cb6dffe774ff5b1e
\newcommand{\XbsNEvalSignal}{\pending{samples not in}}  % experiments/XBS/xbx_report.py :: pending :: sha256 cb6dffe774ff5b1e
\newcommand{\XbsNEvalQcd}{\pending{samples not in}}  % experiments/XBS/xbx_report.py :: pending :: sha256 cb6dffe774ff5b1e
\newcommand{\MlstPartreproRthirtyPart}{1{,}711}  % experiments/FIGS/data/v2/partrepro/part.json :: r30 :: sha256 0d44b997af9e6448 :: ParT reproduction on top, part: R30 (one run)
\newcommand{\MlstPartreproRthirtyHead}{1{,}577}  % experiments/FIGS/data/v2/partrepro/head.json :: r30 :: sha256 37b82ba0a0fe9299 :: ParT reproduction on top, head: R30 (one run)
\newcommand{\MlstPartreproRthirtyRescale}{1{,}224}  % experiments/FIGS/data/v2/partrepro/rescale.json :: r30 :: sha256 0a9f1401047d07ff :: ParT reproduction on top, rescale: R30 (one run)
\newcommand{\MlstPartreproRthirtyHeadOverPart}{0.922}  % experiments/FIGS/data/v2/partrepro/head.json :: r30 / part r30 :: sha256 37b82ba0a0fe9299
\newcommand{\MlstPartreproRthirtyHeadOverPartPr}{7.8\% less rejection}  % experiments/FIGS/data/v2/partrepro/head.json :: r30 / part r30 :: sha256 37b82ba0a0fe9299 :: prose twin: two significant figures, direction in words
\newcommand{\MlstPartreproRthirtyRescaleOverPart}{0.715}  % experiments/FIGS/data/v2/partrepro/rescale.json :: r30 / part r30 :: sha256 0a9f1401047d07ff
\newcommand{\MlstPartreproRthirtyRescaleOverPartPr}{28\% less rejection}  % experiments/FIGS/data/v2/partrepro/rescale.json :: r30 / part r30 :: sha256 0a9f1401047d07ff :: prose twin: two significant figures, direction in words
\newcommand{\MlstPartreproSliceDropRatio}{\pending{jet-pT slice test not run}}  % experiments/FIGS/data/v2/index.json :: open item O10: R30 in narrow pT slices of partrepro part and rescale :: sha256 c6c4138ee5c86423
\newcommand{\MlstInputBnZeroSigmaMax}{0.96}  % experiments/FIGS/data/v2/head/input_bn_offsets.json :: offsets :: sha256 d6b821509fccecff :: largest |offset| over the five type flags
\newcommand{\MlstFtJcOneNSignal}{nine}  % experiments/FIGS/data/v2/curves/full_budget_long_leg2.json :: JC1_COLS (count) :: sha256 588e935996d1006a :: JetClass signal classes
\newcommand{\MlstFtJcOneRejHbbResOverFull}{0.92\,\ensuremath{\pm}\,0.05}  % experiments/FIGS/data/v2/curves/full_budget_long_leg2.json :: cells.<run>.N1000000.s1.jc1.per_class.label_Hbb: rejection of L162 over L188, run by run :: sha256 588e935996d1006a
\newcommand{\MlstFtJcOneRejHbbResOverFullPr}{equal within 9\%}  % experiments/FIGS/data/v2/curves/full_budget_long_leg2.json :: cells.<run>.N1000000.s1.jc1.per_class.label_Hbb: rejection of L162 over L188, run by run :: sha256 588e935996d1006a :: prose twin: two significant figures, direction in words
\newcommand{\MlstFtJcOneRejHbbMidOverFull}{0.98\,\ensuremath{\pm}\,0.12}  % experiments/FIGS/data/v2/curves/full_budget_long_leg2.json :: cells.<run>.N1000000.s1.jc1.per_class.label_Hbb: rejection of R42_Q1 over L188, run by run :: sha256 588e935996d1006a
\newcommand{\MlstFtJcOneRejHbbMidOverFullPr}{equal within 20\%}  % experiments/FIGS/data/v2/curves/full_budget_long_leg2.json :: cells.<run>.N1000000.s1.jc1.per_class.label_Hbb: rejection of R42_Q1 over L188, run by run :: sha256 588e935996d1006a :: prose twin: two significant figures, direction in words
\newcommand{\MlstFtJcOneRejHbbCoarseOverFull}{1.01\,\ensuremath{\pm}\,0.07}  % experiments/FIGS/data/v2/curves/full_budget_long_leg2.json :: cells.<run>.N1000000.s1.jc1.per_class.label_Hbb: rejection of R16_Q1 over L188, run by run :: sha256 588e935996d1006a
\newcommand{\MlstFtJcOneRejHbbCoarseOverFullPr}{equal within 7\%}  % experiments/FIGS/data/v2/curves/full_budget_long_leg2.json :: cells.<run>.N1000000.s1.jc1.per_class.label_Hbb: rejection of R16_Q1 over L188, run by run :: sha256 588e935996d1006a :: prose twin: two significant figures, direction in words
\newcommand{\MlstFtJcOneRejHbbResMassOverRes}{1.065 and 1.100}  % experiments/FIGS/data/v2/curves/full_budget_long_leg2.json :: cells.<run>.N1000000.s1.jc1.per_class.label_Hbb: rejection of L162_MASS over L162, run by run :: sha256 588e935996d1006a
\newcommand{\MlstFtJcOneRejHbbResMassOverResPr}{equal within 10\%}  % experiments/FIGS/data/v2/curves/full_budget_long_leg2.json :: cells.<run>.N1000000.s1.jc1.per_class.label_Hbb: rejection of L162_MASS over L162, run by run :: sha256 588e935996d1006a :: prose twin: two significant figures, direction in words
\newcommand{\MlstFtJcOneRejHbbResMassOverResPrMax}{equal within 10\%}  % experiments/FIGS/data/v2/curves/full_budget_long_leg2.json :: cells.<run>.N1000000.s1.jc1.per_class.label_Hbb: rejection of L162_MASS over L162, run by run :: sha256 588e935996d1006a :: prose twin of the run farther from 1
\newcommand{\MlstFtJcOneRejHbbCoarseMassOverCoarse}{0.97 and 0.86}  % experiments/FIGS/data/v2/curves/full_budget_long_leg2.json :: cells.<run>.N1000000.s1.jc1.per_class.label_Hbb: rejection of R16_Q1_MASS over R16_Q1, run by run :: sha256 588e935996d1006a
\newcommand{\MlstFtJcOneRejHbbCoarseMassOverCoarsePr}{equal within 20\%}  % experiments/FIGS/data/v2/curves/full_budget_long_leg2.json :: cells.<run>.N1000000.s1.jc1.per_class.label_Hbb: rejection of R16_Q1_MASS over R16_Q1, run by run :: sha256 588e935996d1006a :: prose twin: two significant figures, direction in words
\newcommand{\MlstFtJcOneRejHbbCoarseMassOverCoarsePrMax}{equal within 20\%}  % experiments/FIGS/data/v2/curves/full_budget_long_leg2.json :: cells.<run>.N1000000.s1.jc1.per_class.label_Hbb: rejection of R16_Q1_MASS over R16_Q1, run by run :: sha256 588e935996d1006a :: prose twin of the run farther from 1
\newcommand{\MlstFtJcOneRejHbbRandOverFull}{0.92\,\ensuremath{\pm}\,0.11}  % experiments/FIGS/data/v2/curves/full_budget_long_leg2.json :: per_class.label_Hbb: rejection of each random partition over the 188-class run of its seed :: sha256 588e935996d1006a
\newcommand{\MlstFtJcOneRejHbbRandOverFullPr}{equal within 20\%}  % experiments/FIGS/data/v2/curves/full_budget_long_leg2.json :: per_class.label_Hbb: rejection of each random partition over the 188-class run of its seed :: sha256 588e935996d1006a :: prose twin: two significant figures, direction in words
\newcommand{\MlstFtJcOneRejHbbRandOverFullRange}{0.84--1.1}  % experiments/FIGS/data/v2/curves/full_budget_long_leg2.json :: per_class.label_Hbb: smallest and largest of the partitions' ratios :: sha256 588e935996d1006a
\newcommand{\MlstFtJcOneRejHccResOverFull}{1.07\,\ensuremath{\pm}\,0.12}  % experiments/FIGS/data/v2/curves/full_budget_long_leg2.json :: cells.<run>.N1000000.s1.jc1.per_class.label_Hcc: rejection of L162 over L188, run by run :: sha256 588e935996d1006a
\newcommand{\MlstFtJcOneRejHccResOverFullPr}{equal within 20\%}  % experiments/FIGS/data/v2/curves/full_budget_long_leg2.json :: cells.<run>.N1000000.s1.jc1.per_class.label_Hcc: rejection of L162 over L188, run by run :: sha256 588e935996d1006a :: prose twin: two significant figures, direction in words
\newcommand{\MlstFtJcOneRejHccMidOverFull}{0.99\,\ensuremath{\pm}\,0.10}  % experiments/FIGS/data/v2/curves/full_budget_long_leg2.json :: cells.<run>.N1000000.s1.jc1.per_class.label_Hcc: rejection of R42_Q1 over L188, run by run :: sha256 588e935996d1006a
\newcommand{\MlstFtJcOneRejHccMidOverFullPr}{equal within 20\%}  % experiments/FIGS/data/v2/curves/full_budget_long_leg2.json :: cells.<run>.N1000000.s1.jc1.per_class.label_Hcc: rejection of R42_Q1 over L188, run by run :: sha256 588e935996d1006a :: prose twin: two significant figures, direction in words
\newcommand{\MlstFtJcOneRejHccCoarseOverFull}{0.91\,\ensuremath{\pm}\,0.07}  % experiments/FIGS/data/v2/curves/full_budget_long_leg2.json :: cells.<run>.N1000000.s1.jc1.per_class.label_Hcc: rejection of R16_Q1 over L188, run by run :: sha256 588e935996d1006a
\newcommand{\MlstFtJcOneRejHccCoarseOverFullPr}{equal within 10\%}  % experiments/FIGS/data/v2/curves/full_budget_long_leg2.json :: cells.<run>.N1000000.s1.jc1.per_class.label_Hcc: rejection of R16_Q1 over L188, run by run :: sha256 588e935996d1006a :: prose twin: two significant figures, direction in words
\newcommand{\MlstFtJcOneRejHccResMassOverRes}{0.886 and 0.816}  % experiments/FIGS/data/v2/curves/full_budget_long_leg2.json :: cells.<run>.N1000000.s1.jc1.per_class.label_Hcc: rejection of L162_MASS over L162, run by run :: sha256 588e935996d1006a
\newcommand{\MlstFtJcOneRejHccResMassOverResPr}{equal within 20\%}  % experiments/FIGS/data/v2/curves/full_budget_long_leg2.json :: cells.<run>.N1000000.s1.jc1.per_class.label_Hcc: rejection of L162_MASS over L162, run by run :: sha256 588e935996d1006a :: prose twin: two significant figures, direction in words
\newcommand{\MlstFtJcOneRejHccResMassOverResPrMax}{equal within 20\%}  % experiments/FIGS/data/v2/curves/full_budget_long_leg2.json :: cells.<run>.N1000000.s1.jc1.per_class.label_Hcc: rejection of L162_MASS over L162, run by run :: sha256 588e935996d1006a :: prose twin of the run farther from 1
\newcommand{\MlstFtJcOneRejHccCoarseMassOverCoarse}{0.98 and 0.85}  % experiments/FIGS/data/v2/curves/full_budget_long_leg2.json :: cells.<run>.N1000000.s1.jc1.per_class.label_Hcc: rejection of R16_Q1_MASS over R16_Q1, run by run :: sha256 588e935996d1006a
\newcommand{\MlstFtJcOneRejHccCoarseMassOverCoarsePr}{equal within 20\%}  % experiments/FIGS/data/v2/curves/full_budget_long_leg2.json :: cells.<run>.N1000000.s1.jc1.per_class.label_Hcc: rejection of R16_Q1_MASS over R16_Q1, run by run :: sha256 588e935996d1006a :: prose twin: two significant figures, direction in words
\newcommand{\MlstFtJcOneRejHccCoarseMassOverCoarsePrMax}{equal within 20\%}  % experiments/FIGS/data/v2/curves/full_budget_long_leg2.json :: cells.<run>.N1000000.s1.jc1.per_class.label_Hcc: rejection of R16_Q1_MASS over R16_Q1, run by run :: sha256 588e935996d1006a :: prose twin of the run farther from 1
\newcommand{\MlstFtJcOneRejHccRandOverFull}{1.03\,\ensuremath{\pm}\,0.07}  % experiments/FIGS/data/v2/curves/full_budget_long_leg2.json :: per_class.label_Hcc: rejection of each random partition over the 188-class run of its seed :: sha256 588e935996d1006a
\newcommand{\MlstFtJcOneRejHccRandOverFullPr}{equal within 8\%}  % experiments/FIGS/data/v2/curves/full_budget_long_leg2.json :: per_class.label_Hcc: rejection of each random partition over the 188-class run of its seed :: sha256 588e935996d1006a :: prose twin: two significant figures, direction in words
\newcommand{\MlstFtJcOneRejHccRandOverFullRange}{0.94--1.1}  % experiments/FIGS/data/v2/curves/full_budget_long_leg2.json :: per_class.label_Hcc: smallest and largest of the partitions' ratios :: sha256 588e935996d1006a
\newcommand{\MlstFtJcOneRejHggResOverFull}{0.98\,\ensuremath{\pm}\,0.05}  % experiments/FIGS/data/v2/curves/full_budget_long_leg2.json :: cells.<run>.N1000000.s1.jc1.per_class.label_Hgg: rejection of L162 over L188, run by run :: sha256 588e935996d1006a
\newcommand{\MlstFtJcOneRejHggResOverFullPr}{equal within 6\%}  % experiments/FIGS/data/v2/curves/full_budget_long_leg2.json :: cells.<run>.N1000000.s1.jc1.per_class.label_Hgg: rejection of L162 over L188, run by run :: sha256 588e935996d1006a :: prose twin: two significant figures, direction in words
\newcommand{\MlstFtJcOneRejHggMidOverFull}{1.000\,\ensuremath{\pm}\,0.025}  % experiments/FIGS/data/v2/curves/full_budget_long_leg2.json :: cells.<run>.N1000000.s1.jc1.per_class.label_Hgg: rejection of R42_Q1 over L188, run by run :: sha256 588e935996d1006a
\newcommand{\MlstFtJcOneRejHggMidOverFullPr}{equal within 3\%}  % experiments/FIGS/data/v2/curves/full_budget_long_leg2.json :: cells.<run>.N1000000.s1.jc1.per_class.label_Hgg: rejection of R42_Q1 over L188, run by run :: sha256 588e935996d1006a :: prose twin: two significant figures, direction in words
\newcommand{\MlstFtJcOneRejHggCoarseOverFull}{0.97\,\ensuremath{\pm}\,0.09}  % experiments/FIGS/data/v2/curves/full_budget_long_leg2.json :: cells.<run>.N1000000.s1.jc1.per_class.label_Hgg: rejection of R16_Q1 over L188, run by run :: sha256 588e935996d1006a
\newcommand{\MlstFtJcOneRejHggCoarseOverFullPr}{equal within 10\%}  % experiments/FIGS/data/v2/curves/full_budget_long_leg2.json :: cells.<run>.N1000000.s1.jc1.per_class.label_Hgg: rejection of R16_Q1 over L188, run by run :: sha256 588e935996d1006a :: prose twin: two significant figures, direction in words
\newcommand{\MlstFtJcOneRejHggResMassOverRes}{1.046 and 1.008}  % experiments/FIGS/data/v2/curves/full_budget_long_leg2.json :: cells.<run>.N1000000.s1.jc1.per_class.label_Hgg: rejection of L162_MASS over L162, run by run :: sha256 588e935996d1006a
\newcommand{\MlstFtJcOneRejHggResMassOverResPr}{equal within 5\%}  % experiments/FIGS/data/v2/curves/full_budget_long_leg2.json :: cells.<run>.N1000000.s1.jc1.per_class.label_Hgg: rejection of L162_MASS over L162, run by run :: sha256 588e935996d1006a :: prose twin: two significant figures, direction in words
\newcommand{\MlstFtJcOneRejHggResMassOverResPrMax}{equal within 5\%}  % experiments/FIGS/data/v2/curves/full_budget_long_leg2.json :: cells.<run>.N1000000.s1.jc1.per_class.label_Hgg: rejection of L162_MASS over L162, run by run :: sha256 588e935996d1006a :: prose twin of the run farther from 1
\newcommand{\MlstFtJcOneRejHggCoarseMassOverCoarse}{1.12 and 1.01}  % experiments/FIGS/data/v2/curves/full_budget_long_leg2.json :: cells.<run>.N1000000.s1.jc1.per_class.label_Hgg: rejection of R16_Q1_MASS over R16_Q1, run by run :: sha256 588e935996d1006a
\newcommand{\MlstFtJcOneRejHggCoarseMassOverCoarsePr}{equal within 20\%}  % experiments/FIGS/data/v2/curves/full_budget_long_leg2.json :: cells.<run>.N1000000.s1.jc1.per_class.label_Hgg: rejection of R16_Q1_MASS over R16_Q1, run by run :: sha256 588e935996d1006a :: prose twin: two significant figures, direction in words
\newcommand{\MlstFtJcOneRejHggCoarseMassOverCoarsePrMax}{equal within 20\%}  % experiments/FIGS/data/v2/curves/full_budget_long_leg2.json :: cells.<run>.N1000000.s1.jc1.per_class.label_Hgg: rejection of R16_Q1_MASS over R16_Q1, run by run :: sha256 588e935996d1006a :: prose twin of the run farther from 1
\newcommand{\MlstFtJcOneRejHggRandOverFull}{0.968\,\ensuremath{\pm}\,0.016}  % experiments/FIGS/data/v2/curves/full_budget_long_leg2.json :: per_class.label_Hgg: rejection of each random partition over the 188-class run of its seed :: sha256 588e935996d1006a
\newcommand{\MlstFtJcOneRejHggRandOverFullPr}{equal within 4\%}  % experiments/FIGS/data/v2/curves/full_budget_long_leg2.json :: per_class.label_Hgg: rejection of each random partition over the 188-class run of its seed :: sha256 588e935996d1006a :: prose twin: two significant figures, direction in words
\newcommand{\MlstFtJcOneRejHggRandOverFullRange}{0.94--0.98}  % experiments/FIGS/data/v2/curves/full_budget_long_leg2.json :: per_class.label_Hgg: smallest and largest of the partitions' ratios :: sha256 588e935996d1006a
\newcommand{\MlstFtJcOneRejHfourqResOverFull}{1.05\,\ensuremath{\pm}\,0.06}  % experiments/FIGS/data/v2/curves/full_budget_long_leg2.json :: cells.<run>.N1000000.s1.jc1.per_class.label_H4q: rejection of L162 over L188, run by run :: sha256 588e935996d1006a
\newcommand{\MlstFtJcOneRejHfourqResOverFullPr}{equal within 7\%}  % experiments/FIGS/data/v2/curves/full_budget_long_leg2.json :: cells.<run>.N1000000.s1.jc1.per_class.label_H4q: rejection of L162 over L188, run by run :: sha256 588e935996d1006a :: prose twin: two significant figures, direction in words
\newcommand{\MlstFtJcOneRejHfourqMidOverFull}{0.95\,\ensuremath{\pm}\,0.04}  % experiments/FIGS/data/v2/curves/full_budget_long_leg2.json :: cells.<run>.N1000000.s1.jc1.per_class.label_H4q: rejection of R42_Q1 over L188, run by run :: sha256 588e935996d1006a
\newcommand{\MlstFtJcOneRejHfourqMidOverFullPr}{equal within 6\%}  % experiments/FIGS/data/v2/curves/full_budget_long_leg2.json :: cells.<run>.N1000000.s1.jc1.per_class.label_H4q: rejection of R42_Q1 over L188, run by run :: sha256 588e935996d1006a :: prose twin: two significant figures, direction in words
\newcommand{\MlstFtJcOneRejHfourqCoarseOverFull}{0.962\,\ensuremath{\pm}\,0.025}  % experiments/FIGS/data/v2/curves/full_budget_long_leg2.json :: cells.<run>.N1000000.s1.jc1.per_class.label_H4q: rejection of R16_Q1 over L188, run by run :: sha256 588e935996d1006a
\newcommand{\MlstFtJcOneRejHfourqCoarseOverFullPr}{equal within 4\%}  % experiments/FIGS/data/v2/curves/full_budget_long_leg2.json :: cells.<run>.N1000000.s1.jc1.per_class.label_H4q: rejection of R16_Q1 over L188, run by run :: sha256 588e935996d1006a :: prose twin: two significant figures, direction in words
\newcommand{\MlstFtJcOneRejHfourqResMassOverRes}{0.912 and 0.966}  % experiments/FIGS/data/v2/curves/full_budget_long_leg2.json :: cells.<run>.N1000000.s1.jc1.per_class.label_H4q: rejection of L162_MASS over L162, run by run :: sha256 588e935996d1006a
\newcommand{\MlstFtJcOneRejHfourqResMassOverResPr}{equal within 9\%}  % experiments/FIGS/data/v2/curves/full_budget_long_leg2.json :: cells.<run>.N1000000.s1.jc1.per_class.label_H4q: rejection of L162_MASS over L162, run by run :: sha256 588e935996d1006a :: prose twin: two significant figures, direction in words
\newcommand{\MlstFtJcOneRejHfourqResMassOverResPrMax}{equal within 9\%}  % experiments/FIGS/data/v2/curves/full_budget_long_leg2.json :: cells.<run>.N1000000.s1.jc1.per_class.label_H4q: rejection of L162_MASS over L162, run by run :: sha256 588e935996d1006a :: prose twin of the run farther from 1
\newcommand{\MlstFtJcOneRejHfourqCoarseMassOverCoarse}{0.937 and 0.888}  % experiments/FIGS/data/v2/curves/full_budget_long_leg2.json :: cells.<run>.N1000000.s1.jc1.per_class.label_H4q: rejection of R16_Q1_MASS over R16_Q1, run by run :: sha256 588e935996d1006a
\newcommand{\MlstFtJcOneRejHfourqCoarseMassOverCoarsePr}{equal within 20\%}  % experiments/FIGS/data/v2/curves/full_budget_long_leg2.json :: cells.<run>.N1000000.s1.jc1.per_class.label_H4q: rejection of R16_Q1_MASS over R16_Q1, run by run :: sha256 588e935996d1006a :: prose twin: two significant figures, direction in words
\newcommand{\MlstFtJcOneRejHfourqCoarseMassOverCoarsePrMax}{equal within 20\%}  % experiments/FIGS/data/v2/curves/full_budget_long_leg2.json :: cells.<run>.N1000000.s1.jc1.per_class.label_H4q: rejection of R16_Q1_MASS over R16_Q1, run by run :: sha256 588e935996d1006a :: prose twin of the run farther from 1
\newcommand{\MlstFtJcOneRejHfourqRandOverFull}{0.93\,\ensuremath{\pm}\,0.04}  % experiments/FIGS/data/v2/curves/full_budget_long_leg2.json :: per_class.label_H4q: rejection of each random partition over the 188-class run of its seed :: sha256 588e935996d1006a
\newcommand{\MlstFtJcOneRejHfourqRandOverFullPr}{equal within 8\%}  % experiments/FIGS/data/v2/curves/full_budget_long_leg2.json :: per_class.label_H4q: rejection of each random partition over the 188-class run of its seed :: sha256 588e935996d1006a :: prose twin: two significant figures, direction in words
\newcommand{\MlstFtJcOneRejHfourqRandOverFullRange}{0.86--0.96}  % experiments/FIGS/data/v2/curves/full_budget_long_leg2.json :: per_class.label_H4q: smallest and largest of the partitions' ratios :: sha256 588e935996d1006a
\newcommand{\MlstFtJcOneRejHqqlResOverFull}{1.004\,\ensuremath{\pm}\,0.018}  % experiments/FIGS/data/v2/curves/full_budget_long_leg2.json :: cells.<run>.N1000000.s1.jc1.per_class.label_Hqql: rejection of L162 over L188, run by run :: sha256 588e935996d1006a
\newcommand{\MlstFtJcOneRejHqqlResOverFullPr}{equal within 2\%}  % experiments/FIGS/data/v2/curves/full_budget_long_leg2.json :: cells.<run>.N1000000.s1.jc1.per_class.label_Hqql: rejection of L162 over L188, run by run :: sha256 588e935996d1006a :: prose twin: two significant figures, direction in words
\newcommand{\MlstFtJcOneRejHqqlMidOverFull}{1.04\,\ensuremath{\pm}\,0.18}  % experiments/FIGS/data/v2/curves/full_budget_long_leg2.json :: cells.<run>.N1000000.s1.jc1.per_class.label_Hqql: rejection of R42_Q1 over L188, run by run :: sha256 588e935996d1006a
\newcommand{\MlstFtJcOneRejHqqlMidOverFullPr}{equal within 20\%}  % experiments/FIGS/data/v2/curves/full_budget_long_leg2.json :: cells.<run>.N1000000.s1.jc1.per_class.label_Hqql: rejection of R42_Q1 over L188, run by run :: sha256 588e935996d1006a :: prose twin: two significant figures, direction in words
\newcommand{\MlstFtJcOneRejHqqlCoarseOverFull}{0.93\,\ensuremath{\pm}\,0.16}  % experiments/FIGS/data/v2/curves/full_budget_long_leg2.json :: cells.<run>.N1000000.s1.jc1.per_class.label_Hqql: rejection of R16_Q1 over L188, run by run :: sha256 588e935996d1006a
\newcommand{\MlstFtJcOneRejHqqlCoarseOverFullPr}{equal within 20\%}  % experiments/FIGS/data/v2/curves/full_budget_long_leg2.json :: cells.<run>.N1000000.s1.jc1.per_class.label_Hqql: rejection of R16_Q1 over L188, run by run :: sha256 588e935996d1006a :: prose twin: two significant figures, direction in words
\newcommand{\MlstFtJcOneRejHqqlResMassOverRes}{0.986 and 1.000}  % experiments/FIGS/data/v2/curves/full_budget_long_leg2.json :: cells.<run>.N1000000.s1.jc1.per_class.label_Hqql: rejection of L162_MASS over L162, run by run :: sha256 588e935996d1006a
\newcommand{\MlstFtJcOneRejHqqlResMassOverResPr}{equal within 2\%}  % experiments/FIGS/data/v2/curves/full_budget_long_leg2.json :: cells.<run>.N1000000.s1.jc1.per_class.label_Hqql: rejection of L162_MASS over L162, run by run :: sha256 588e935996d1006a :: prose twin: two significant figures, direction in words
\newcommand{\MlstFtJcOneRejHqqlResMassOverResPrMax}{equal within 2\%}  % experiments/FIGS/data/v2/curves/full_budget_long_leg2.json :: cells.<run>.N1000000.s1.jc1.per_class.label_Hqql: rejection of L162_MASS over L162, run by run :: sha256 588e935996d1006a :: prose twin of the run farther from 1
\newcommand{\MlstFtJcOneRejHqqlCoarseMassOverCoarse}{0.97 and 1.09}  % experiments/FIGS/data/v2/curves/full_budget_long_leg2.json :: cells.<run>.N1000000.s1.jc1.per_class.label_Hqql: rejection of R16_Q1_MASS over R16_Q1, run by run :: sha256 588e935996d1006a
\newcommand{\MlstFtJcOneRejHqqlCoarseMassOverCoarsePr}{equal within 9\%}  % experiments/FIGS/data/v2/curves/full_budget_long_leg2.json :: cells.<run>.N1000000.s1.jc1.per_class.label_Hqql: rejection of R16_Q1_MASS over R16_Q1, run by run :: sha256 588e935996d1006a :: prose twin: two significant figures, direction in words
\newcommand{\MlstFtJcOneRejHqqlCoarseMassOverCoarsePrMax}{equal within 9\%}  % experiments/FIGS/data/v2/curves/full_budget_long_leg2.json :: cells.<run>.N1000000.s1.jc1.per_class.label_Hqql: rejection of R16_Q1_MASS over R16_Q1, run by run :: sha256 588e935996d1006a :: prose twin of the run farther from 1
\newcommand{\MlstFtJcOneRejHqqlRandOverFull}{0.94\,\ensuremath{\pm}\,0.12}  % experiments/FIGS/data/v2/curves/full_budget_long_leg2.json :: per_class.label_Hqql: rejection of each random partition over the 188-class run of its seed :: sha256 588e935996d1006a
\newcommand{\MlstFtJcOneRejHqqlRandOverFullPr}{equal within 20\%}  % experiments/FIGS/data/v2/curves/full_budget_long_leg2.json :: per_class.label_Hqql: rejection of each random partition over the 188-class run of its seed :: sha256 588e935996d1006a :: prose twin: two significant figures, direction in words
\newcommand{\MlstFtJcOneRejHqqlRandOverFullRange}{0.79--1.1}  % experiments/FIGS/data/v2/curves/full_budget_long_leg2.json :: per_class.label_Hqql: smallest and largest of the partitions' ratios :: sha256 588e935996d1006a
\newcommand{\MlstFtJcOneRejTbqqResOverFull}{1.1\,\ensuremath{\pm}\,0.4}  % experiments/FIGS/data/v2/curves/full_budget_long_leg2.json :: cells.<run>.N1000000.s1.jc1.per_class.label_Tbqq: rejection of L162 over L188, run by run :: sha256 588e935996d1006a
\newcommand{\MlstFtJcOneRejTbqqResOverFullPr}{equal within 50\%}  % experiments/FIGS/data/v2/curves/full_budget_long_leg2.json :: cells.<run>.N1000000.s1.jc1.per_class.label_Tbqq: rejection of L162 over L188, run by run :: sha256 588e935996d1006a :: prose twin: two significant figures, direction in words
\newcommand{\MlstFtJcOneRejTbqqMidOverFull}{1.10\,\ensuremath{\pm}\,0.05}  % experiments/FIGS/data/v2/curves/full_budget_long_leg2.json :: cells.<run>.N1000000.s1.jc1.per_class.label_Tbqq: rejection of R42_Q1 over L188, run by run :: sha256 588e935996d1006a
\newcommand{\MlstFtJcOneRejTbqqMidOverFullPr}{equal within 10\%}  % experiments/FIGS/data/v2/curves/full_budget_long_leg2.json :: cells.<run>.N1000000.s1.jc1.per_class.label_Tbqq: rejection of R42_Q1 over L188, run by run :: sha256 588e935996d1006a :: prose twin: two significant figures, direction in words
\newcommand{\MlstFtJcOneRejTbqqCoarseOverFull}{0.85\,\ensuremath{\pm}\,0.32}  % experiments/FIGS/data/v2/curves/full_budget_long_leg2.json :: cells.<run>.N1000000.s1.jc1.per_class.label_Tbqq: rejection of R16_Q1 over L188, run by run :: sha256 588e935996d1006a
\newcommand{\MlstFtJcOneRejTbqqCoarseOverFullPr}{equal within 40\%}  % experiments/FIGS/data/v2/curves/full_budget_long_leg2.json :: cells.<run>.N1000000.s1.jc1.per_class.label_Tbqq: rejection of R16_Q1 over L188, run by run :: sha256 588e935996d1006a :: prose twin: two significant figures, direction in words
\newcommand{\MlstFtJcOneRejTbqqResMassOverRes}{0.89 and 1.21}  % experiments/FIGS/data/v2/curves/full_budget_long_leg2.json :: cells.<run>.N1000000.s1.jc1.per_class.label_Tbqq: rejection of L162_MASS over L162, run by run :: sha256 588e935996d1006a
\newcommand{\MlstFtJcOneRejTbqqResMassOverResPr}{equal within 30\%}  % experiments/FIGS/data/v2/curves/full_budget_long_leg2.json :: cells.<run>.N1000000.s1.jc1.per_class.label_Tbqq: rejection of L162_MASS over L162, run by run :: sha256 588e935996d1006a :: prose twin: two significant figures, direction in words
\newcommand{\MlstFtJcOneRejTbqqResMassOverResPrMax}{equal within 30\%}  % experiments/FIGS/data/v2/curves/full_budget_long_leg2.json :: cells.<run>.N1000000.s1.jc1.per_class.label_Tbqq: rejection of L162_MASS over L162, run by run :: sha256 588e935996d1006a :: prose twin of the run farther from 1
\newcommand{\MlstFtJcOneRejTbqqCoarseMassOverCoarse}{1.00 and 1.79}  % experiments/FIGS/data/v2/curves/full_budget_long_leg2.json :: cells.<run>.N1000000.s1.jc1.per_class.label_Tbqq: rejection of R16_Q1_MASS over R16_Q1, run by run :: sha256 588e935996d1006a
\newcommand{\MlstFtJcOneRejTbqqCoarseMassOverCoarsePr}{equal within 80\%}  % experiments/FIGS/data/v2/curves/full_budget_long_leg2.json :: cells.<run>.N1000000.s1.jc1.per_class.label_Tbqq: rejection of R16_Q1_MASS over R16_Q1, run by run :: sha256 588e935996d1006a :: prose twin: two significant figures, direction in words
\newcommand{\MlstFtJcOneRejTbqqCoarseMassOverCoarsePrMax}{equal within 80\%}  % experiments/FIGS/data/v2/curves/full_budget_long_leg2.json :: cells.<run>.N1000000.s1.jc1.per_class.label_Tbqq: rejection of R16_Q1_MASS over R16_Q1, run by run :: sha256 588e935996d1006a :: prose twin of the run farther from 1
\newcommand{\MlstFtJcOneRejTbqqRandOverFull}{0.79\,\ensuremath{\pm}\,0.15}  % experiments/FIGS/data/v2/curves/full_budget_long_leg2.json :: per_class.label_Tbqq: rejection of each random partition over the 188-class run of its seed :: sha256 588e935996d1006a
\newcommand{\MlstFtJcOneRejTbqqRandOverFullPr}{equal within 30\%}  % experiments/FIGS/data/v2/curves/full_budget_long_leg2.json :: per_class.label_Tbqq: rejection of each random partition over the 188-class run of its seed :: sha256 588e935996d1006a :: prose twin: two significant figures, direction in words
\newcommand{\MlstFtJcOneRejTbqqRandOverFullRange}{0.64--1.0}  % experiments/FIGS/data/v2/curves/full_budget_long_leg2.json :: per_class.label_Tbqq: smallest and largest of the partitions' ratios :: sha256 588e935996d1006a
\newcommand{\MlstFtJcOneRejTblResOverFull}{0.85\,\ensuremath{\pm}\,0.17}  % experiments/FIGS/data/v2/curves/full_budget_long_leg2.json :: cells.<run>.N1000000.s1.jc1.per_class.label_Tbl: rejection of L162 over L188, run by run :: sha256 588e935996d1006a
\newcommand{\MlstFtJcOneRejTblResOverFullPr}{equal within 20\%}  % experiments/FIGS/data/v2/curves/full_budget_long_leg2.json :: cells.<run>.N1000000.s1.jc1.per_class.label_Tbl: rejection of L162 over L188, run by run :: sha256 588e935996d1006a :: prose twin: two significant figures, direction in words
\newcommand{\MlstFtJcOneRejTblMidOverFull}{0.91\,\ensuremath{\pm}\,0.15}  % experiments/FIGS/data/v2/curves/full_budget_long_leg2.json :: cells.<run>.N1000000.s1.jc1.per_class.label_Tbl: rejection of R42_Q1 over L188, run by run :: sha256 588e935996d1006a
\newcommand{\MlstFtJcOneRejTblMidOverFullPr}{equal within 20\%}  % experiments/FIGS/data/v2/curves/full_budget_long_leg2.json :: cells.<run>.N1000000.s1.jc1.per_class.label_Tbl: rejection of R42_Q1 over L188, run by run :: sha256 588e935996d1006a :: prose twin: two significant figures, direction in words
\newcommand{\MlstFtJcOneRejTblCoarseOverFull}{0.94\,\ensuremath{\pm}\,0.28}  % experiments/FIGS/data/v2/curves/full_budget_long_leg2.json :: cells.<run>.N1000000.s1.jc1.per_class.label_Tbl: rejection of R16_Q1 over L188, run by run :: sha256 588e935996d1006a
\newcommand{\MlstFtJcOneRejTblCoarseOverFullPr}{equal within 30\%}  % experiments/FIGS/data/v2/curves/full_budget_long_leg2.json :: cells.<run>.N1000000.s1.jc1.per_class.label_Tbl: rejection of R16_Q1 over L188, run by run :: sha256 588e935996d1006a :: prose twin: two significant figures, direction in words
\newcommand{\MlstFtJcOneRejTblResMassOverRes}{1.04 and 0.93}  % experiments/FIGS/data/v2/curves/full_budget_long_leg2.json :: cells.<run>.N1000000.s1.jc1.per_class.label_Tbl: rejection of L162_MASS over L162, run by run :: sha256 588e935996d1006a
\newcommand{\MlstFtJcOneRejTblResMassOverResPr}{equal within 7\%}  % experiments/FIGS/data/v2/curves/full_budget_long_leg2.json :: cells.<run>.N1000000.s1.jc1.per_class.label_Tbl: rejection of L162_MASS over L162, run by run :: sha256 588e935996d1006a :: prose twin: two significant figures, direction in words
\newcommand{\MlstFtJcOneRejTblResMassOverResPrMax}{equal within 7\%}  % experiments/FIGS/data/v2/curves/full_budget_long_leg2.json :: cells.<run>.N1000000.s1.jc1.per_class.label_Tbl: rejection of L162_MASS over L162, run by run :: sha256 588e935996d1006a :: prose twin of the run farther from 1
\newcommand{\MlstFtJcOneRejTblCoarseMassOverCoarse}{0.84 and 1.22}  % experiments/FIGS/data/v2/curves/full_budget_long_leg2.json :: cells.<run>.N1000000.s1.jc1.per_class.label_Tbl: rejection of R16_Q1_MASS over R16_Q1, run by run :: sha256 588e935996d1006a
\newcommand{\MlstFtJcOneRejTblCoarseMassOverCoarsePr}{equal within 30\%}  % experiments/FIGS/data/v2/curves/full_budget_long_leg2.json :: cells.<run>.N1000000.s1.jc1.per_class.label_Tbl: rejection of R16_Q1_MASS over R16_Q1, run by run :: sha256 588e935996d1006a :: prose twin: two significant figures, direction in words
\newcommand{\MlstFtJcOneRejTblCoarseMassOverCoarsePrMax}{equal within 30\%}  % experiments/FIGS/data/v2/curves/full_budget_long_leg2.json :: cells.<run>.N1000000.s1.jc1.per_class.label_Tbl: rejection of R16_Q1_MASS over R16_Q1, run by run :: sha256 588e935996d1006a :: prose twin of the run farther from 1
\newcommand{\MlstFtJcOneRejTblRandOverFull}{0.88\,\ensuremath{\pm}\,0.26}  % experiments/FIGS/data/v2/curves/full_budget_long_leg2.json :: per_class.label_Tbl: rejection of each random partition over the 188-class run of its seed :: sha256 588e935996d1006a
\newcommand{\MlstFtJcOneRejTblRandOverFullPr}{equal within 30\%}  % experiments/FIGS/data/v2/curves/full_budget_long_leg2.json :: per_class.label_Tbl: rejection of each random partition over the 188-class run of its seed :: sha256 588e935996d1006a :: prose twin: two significant figures, direction in words
\newcommand{\MlstFtJcOneRejTblRandOverFullRange}{0.61--1.2}  % experiments/FIGS/data/v2/curves/full_budget_long_leg2.json :: per_class.label_Tbl: smallest and largest of the partitions' ratios :: sha256 588e935996d1006a
\newcommand{\MlstFtJcOneRejWqqResOverFull}{1.039\,\ensuremath{\pm}\,0.035}  % experiments/FIGS/data/v2/curves/full_budget_long_leg2.json :: cells.<run>.N1000000.s1.jc1.per_class.label_Wqq: rejection of L162 over L188, run by run :: sha256 588e935996d1006a
\newcommand{\MlstFtJcOneRejWqqResOverFullPr}{equal within 4\%}  % experiments/FIGS/data/v2/curves/full_budget_long_leg2.json :: cells.<run>.N1000000.s1.jc1.per_class.label_Wqq: rejection of L162 over L188, run by run :: sha256 588e935996d1006a :: prose twin: two significant figures, direction in words
\newcommand{\MlstFtJcOneRejWqqMidOverFull}{1.005\,\ensuremath{\pm}\,0.032}  % experiments/FIGS/data/v2/curves/full_budget_long_leg2.json :: cells.<run>.N1000000.s1.jc1.per_class.label_Wqq: rejection of R42_Q1 over L188, run by run :: sha256 588e935996d1006a
\newcommand{\MlstFtJcOneRejWqqMidOverFullPr}{equal within 4\%}  % experiments/FIGS/data/v2/curves/full_budget_long_leg2.json :: cells.<run>.N1000000.s1.jc1.per_class.label_Wqq: rejection of R42_Q1 over L188, run by run :: sha256 588e935996d1006a :: prose twin: two significant figures, direction in words
\newcommand{\MlstFtJcOneRejWqqCoarseOverFull}{1.01\,\ensuremath{\pm}\,0.07}  % experiments/FIGS/data/v2/curves/full_budget_long_leg2.json :: cells.<run>.N1000000.s1.jc1.per_class.label_Wqq: rejection of R16_Q1 over L188, run by run :: sha256 588e935996d1006a
\newcommand{\MlstFtJcOneRejWqqCoarseOverFullPr}{equal within 8\%}  % experiments/FIGS/data/v2/curves/full_budget_long_leg2.json :: cells.<run>.N1000000.s1.jc1.per_class.label_Wqq: rejection of R16_Q1 over L188, run by run :: sha256 588e935996d1006a :: prose twin: two significant figures, direction in words
\newcommand{\MlstFtJcOneRejWqqResMassOverRes}{0.996 and 0.964}  % experiments/FIGS/data/v2/curves/full_budget_long_leg2.json :: cells.<run>.N1000000.s1.jc1.per_class.label_Wqq: rejection of L162_MASS over L162, run by run :: sha256 588e935996d1006a
\newcommand{\MlstFtJcOneRejWqqResMassOverResPr}{equal within 4\%}  % experiments/FIGS/data/v2/curves/full_budget_long_leg2.json :: cells.<run>.N1000000.s1.jc1.per_class.label_Wqq: rejection of L162_MASS over L162, run by run :: sha256 588e935996d1006a :: prose twin: two significant figures, direction in words
\newcommand{\MlstFtJcOneRejWqqResMassOverResPrMax}{equal within 4\%}  % experiments/FIGS/data/v2/curves/full_budget_long_leg2.json :: cells.<run>.N1000000.s1.jc1.per_class.label_Wqq: rejection of L162_MASS over L162, run by run :: sha256 588e935996d1006a :: prose twin of the run farther from 1
\newcommand{\MlstFtJcOneRejWqqCoarseMassOverCoarse}{1.025 and 0.926}  % experiments/FIGS/data/v2/curves/full_budget_long_leg2.json :: cells.<run>.N1000000.s1.jc1.per_class.label_Wqq: rejection of R16_Q1_MASS over R16_Q1, run by run :: sha256 588e935996d1006a
\newcommand{\MlstFtJcOneRejWqqCoarseMassOverCoarsePr}{equal within 8\%}  % experiments/FIGS/data/v2/curves/full_budget_long_leg2.json :: cells.<run>.N1000000.s1.jc1.per_class.label_Wqq: rejection of R16_Q1_MASS over R16_Q1, run by run :: sha256 588e935996d1006a :: prose twin: two significant figures, direction in words
\newcommand{\MlstFtJcOneRejWqqCoarseMassOverCoarsePrMax}{equal within 8\%}  % experiments/FIGS/data/v2/curves/full_budget_long_leg2.json :: cells.<run>.N1000000.s1.jc1.per_class.label_Wqq: rejection of R16_Q1_MASS over R16_Q1, run by run :: sha256 588e935996d1006a :: prose twin of the run farther from 1
\newcommand{\MlstFtJcOneRejWqqRandOverFull}{1.02\,\ensuremath{\pm}\,0.07}  % experiments/FIGS/data/v2/curves/full_budget_long_leg2.json :: per_class.label_Wqq: rejection of each random partition over the 188-class run of its seed :: sha256 588e935996d1006a
\newcommand{\MlstFtJcOneRejWqqRandOverFullPr}{equal within 7\%}  % experiments/FIGS/data/v2/curves/full_budget_long_leg2.json :: per_class.label_Wqq: rejection of each random partition over the 188-class run of its seed :: sha256 588e935996d1006a :: prose twin: two significant figures, direction in words
\newcommand{\MlstFtJcOneRejWqqRandOverFullRange}{0.93--1.1}  % experiments/FIGS/data/v2/curves/full_budget_long_leg2.json :: per_class.label_Wqq: smallest and largest of the partitions' ratios :: sha256 588e935996d1006a
\newcommand{\MlstFtJcOneRejZqqResOverFull}{0.987\,\ensuremath{\pm}\,0.019}  % experiments/FIGS/data/v2/curves/full_budget_long_leg2.json :: cells.<run>.N1000000.s1.jc1.per_class.label_Zqq: rejection of L162 over L188, run by run :: sha256 588e935996d1006a
\newcommand{\MlstFtJcOneRejZqqResOverFullPr}{equal within 2\%}  % experiments/FIGS/data/v2/curves/full_budget_long_leg2.json :: cells.<run>.N1000000.s1.jc1.per_class.label_Zqq: rejection of L162 over L188, run by run :: sha256 588e935996d1006a :: prose twin: two significant figures, direction in words
\newcommand{\MlstFtJcOneRejZqqMidOverFull}{0.984\,\ensuremath{\pm}\,0.010}  % experiments/FIGS/data/v2/curves/full_budget_long_leg2.json :: cells.<run>.N1000000.s1.jc1.per_class.label_Zqq: rejection of R42_Q1 over L188, run by run :: sha256 588e935996d1006a
\newcommand{\MlstFtJcOneRejZqqMidOverFullPr}{equal within 2\%}  % experiments/FIGS/data/v2/curves/full_budget_long_leg2.json :: cells.<run>.N1000000.s1.jc1.per_class.label_Zqq: rejection of R42_Q1 over L188, run by run :: sha256 588e935996d1006a :: prose twin: two significant figures, direction in words
\newcommand{\MlstFtJcOneRejZqqCoarseOverFull}{0.91\,\ensuremath{\pm}\,0.05}  % experiments/FIGS/data/v2/curves/full_budget_long_leg2.json :: cells.<run>.N1000000.s1.jc1.per_class.label_Zqq: rejection of R16_Q1 over L188, run by run :: sha256 588e935996d1006a
\newcommand{\MlstFtJcOneRejZqqCoarseOverFullPr}{equal within 9\%}  % experiments/FIGS/data/v2/curves/full_budget_long_leg2.json :: cells.<run>.N1000000.s1.jc1.per_class.label_Zqq: rejection of R16_Q1 over L188, run by run :: sha256 588e935996d1006a :: prose twin: two significant figures, direction in words
\newcommand{\MlstFtJcOneRejZqqResMassOverRes}{0.9803 and 0.9820}  % experiments/FIGS/data/v2/curves/full_budget_long_leg2.json :: cells.<run>.N1000000.s1.jc1.per_class.label_Zqq: rejection of L162_MASS over L162, run by run :: sha256 588e935996d1006a
\newcommand{\MlstFtJcOneRejZqqResMassOverResPr}{equal within 2\%}  % experiments/FIGS/data/v2/curves/full_budget_long_leg2.json :: cells.<run>.N1000000.s1.jc1.per_class.label_Zqq: rejection of L162_MASS over L162, run by run :: sha256 588e935996d1006a :: prose twin: two significant figures, direction in words
\newcommand{\MlstFtJcOneRejZqqResMassOverResPrMax}{equal within 2\%}  % experiments/FIGS/data/v2/curves/full_budget_long_leg2.json :: cells.<run>.N1000000.s1.jc1.per_class.label_Zqq: rejection of L162_MASS over L162, run by run :: sha256 588e935996d1006a :: prose twin of the run farther from 1
\newcommand{\MlstFtJcOneRejZqqCoarseMassOverCoarse}{1.06 and 0.94}  % experiments/FIGS/data/v2/curves/full_budget_long_leg2.json :: cells.<run>.N1000000.s1.jc1.per_class.label_Zqq: rejection of R16_Q1_MASS over R16_Q1, run by run :: sha256 588e935996d1006a
\newcommand{\MlstFtJcOneRejZqqCoarseMassOverCoarsePr}{equal within 7\%}  % experiments/FIGS/data/v2/curves/full_budget_long_leg2.json :: cells.<run>.N1000000.s1.jc1.per_class.label_Zqq: rejection of R16_Q1_MASS over R16_Q1, run by run :: sha256 588e935996d1006a :: prose twin: two significant figures, direction in words
\newcommand{\MlstFtJcOneRejZqqCoarseMassOverCoarsePrMax}{equal within 7\%}  % experiments/FIGS/data/v2/curves/full_budget_long_leg2.json :: cells.<run>.N1000000.s1.jc1.per_class.label_Zqq: rejection of R16_Q1_MASS over R16_Q1, run by run :: sha256 588e935996d1006a :: prose twin of the run farther from 1
\newcommand{\MlstFtJcOneRejZqqRandOverFull}{0.952\,\ensuremath{\pm}\,0.028}  % experiments/FIGS/data/v2/curves/full_budget_long_leg2.json :: per_class.label_Zqq: rejection of each random partition over the 188-class run of its seed :: sha256 588e935996d1006a
\newcommand{\MlstFtJcOneRejZqqRandOverFullPr}{equal within 5\%}  % experiments/FIGS/data/v2/curves/full_budget_long_leg2.json :: per_class.label_Zqq: rejection of each random partition over the 188-class run of its seed :: sha256 588e935996d1006a :: prose twin: two significant figures, direction in words
\newcommand{\MlstFtJcOneRejZqqRandOverFullRange}{0.93--0.99}  % experiments/FIGS/data/v2/curves/full_budget_long_leg2.json :: per_class.label_Zqq: smallest and largest of the partitions' ratios :: sha256 588e935996d1006a
\newcommand{\MlstFtJcOneResNRejDrop}{zero}  % experiments/FIGS/data/v2/curves/full_budget_long_leg2.json :: per_class rejections :: sha256 588e935996d1006a :: classes whose paired rejection ratio to the 188-class model departs from 1 by more than twice its SD and its counting uncertainty, columns with fewer than JC1_MIN_BKG passing jets excluded, ratio below 1 (count)
\newcommand{\MlstFtJcOneResNRejDiff}{zero}  % experiments/FIGS/data/v2/curves/full_budget_long_leg2.json :: per_class rejections :: sha256 588e935996d1006a :: classes whose paired rejection ratio to the 188-class model departs from 1 by more than twice its SD and its counting uncertainty, columns with fewer than JC1_MIN_BKG passing jets excluded (count)
\newcommand{\MlstFtJcOneResRejDropClasses}{none}  % experiments/FIGS/data/v2/curves/full_budget_long_leg2.json :: per_class rejections :: sha256 588e935996d1006a :: the classes of ...NRejDrop, in ParT's column order
\newcommand{\MlstFtJcOneMidNRejDrop}{zero}  % experiments/FIGS/data/v2/curves/full_budget_long_leg2.json :: per_class rejections :: sha256 588e935996d1006a :: classes whose paired rejection ratio to the 188-class model departs from 1 by more than twice its SD and its counting uncertainty, columns with fewer than JC1_MIN_BKG passing jets excluded, ratio below 1 (count)
\newcommand{\MlstFtJcOneMidNRejDiff}{zero}  % experiments/FIGS/data/v2/curves/full_budget_long_leg2.json :: per_class rejections :: sha256 588e935996d1006a :: classes whose paired rejection ratio to the 188-class model departs from 1 by more than twice its SD and its counting uncertainty, columns with fewer than JC1_MIN_BKG passing jets excluded (count)
\newcommand{\MlstFtJcOneMidRejDropClasses}{none}  % experiments/FIGS/data/v2/curves/full_budget_long_leg2.json :: per_class rejections :: sha256 588e935996d1006a :: the classes of ...NRejDrop, in ParT's column order
\newcommand{\MlstFtJcOneCoarseNRejDrop}{zero}  % experiments/FIGS/data/v2/curves/full_budget_long_leg2.json :: per_class rejections :: sha256 588e935996d1006a :: classes whose paired rejection ratio to the 188-class model departs from 1 by more than twice its SD and its counting uncertainty, columns with fewer than JC1_MIN_BKG passing jets excluded, ratio below 1 (count)
\newcommand{\MlstFtJcOneCoarseNRejDiff}{zero}  % experiments/FIGS/data/v2/curves/full_budget_long_leg2.json :: per_class rejections :: sha256 588e935996d1006a :: classes whose paired rejection ratio to the 188-class model departs from 1 by more than twice its SD and its counting uncertainty, columns with fewer than JC1_MIN_BKG passing jets excluded (count)
\newcommand{\MlstFtJcOneCoarseRejDropClasses}{none}  % experiments/FIGS/data/v2/curves/full_budget_long_leg2.json :: per_class rejections :: sha256 588e935996d1006a :: the classes of ...NRejDrop, in ParT's column order
\newcommand{\MlstFtJcOneAccRandMinusCoarse}{+0.002}  % experiments/FIGS/data/v2/curves/full_budget_long_leg2.json :: cells.*.N1000000.s1.jc1.accuracy :: sha256 588e935996d1006a :: mean over the random partitions minus the mean over the 17-class runs
\newcommand{\MlstFtJcOneAccFullBoot}{0.847}  % experiments/FIGS/data/v2/budget/jc1_acc_boot_N1000000.json :: diffs.r16q1-l188.accuracy_ref (mean) :: sha256 7cccb7b0ff900fb8
\newcommand{\MlstFtJcOneAccCoarseBoot}{0.842}  % experiments/FIGS/data/v2/budget/jc1_acc_boot_N1000000.json :: diffs.r16q1-l188.accuracy_config (mean) :: sha256 7cccb7b0ff900fb8
\newcommand{\MlstFtJcOneErrResOverFullBoot}{consistent}  % experiments/FIGS/data/v2/budget/jc1_acc_boot_N1000000.json :: diffs.l162-l188.error_ratio :: sha256 7cccb7b0ff900fb8 :: delta and sigma = hypot(run SD, bootstrap SD) in per cent; 'consistent' when |delta| <= 3 sigma
\newcommand{\MlstFtJcOneErrResOverFullZ}{0.6}  % experiments/FIGS/data/v2/budget/jc1_acc_boot_N1000000.json :: diffs.l162-l188.error_ratio.z :: sha256 7cccb7b0ff900fb8 :: |delta| / sigma
\newcommand{\MlstFtJcOneErrMidOverFullBoot}{consistent}  % experiments/FIGS/data/v2/budget/jc1_acc_boot_N1000000.json :: diffs.r42q1-l188.error_ratio :: sha256 7cccb7b0ff900fb8 :: delta and sigma = hypot(run SD, bootstrap SD) in per cent; 'consistent' when |delta| <= 3 sigma
\newcommand{\MlstFtJcOneErrMidOverFullZ}{1.4}  % experiments/FIGS/data/v2/budget/jc1_acc_boot_N1000000.json :: diffs.r42q1-l188.error_ratio.z :: sha256 7cccb7b0ff900fb8 :: |delta| / sigma
\newcommand{\MlstFtJcOneErrCoarseOverFullBoot}{(3.0\,\ensuremath{\pm}\,0.6)\% more errors}  % experiments/FIGS/data/v2/budget/jc1_acc_boot_N1000000.json :: diffs.r16q1-l188.error_ratio :: sha256 7cccb7b0ff900fb8 :: delta and sigma = hypot(run SD, bootstrap SD) in per cent; 'consistent' when |delta| <= 3 sigma
\newcommand{\MlstFtJcOneErrCoarseOverFullZ}{4.8}  % experiments/FIGS/data/v2/budget/jc1_acc_boot_N1000000.json :: diffs.r16q1-l188.error_ratio.z :: sha256 7cccb7b0ff900fb8 :: |delta| / sigma
\newcommand{\MlstFtJcOneBootReplicas}{200}  % experiments/FIGS/data/v2/budget/jc1_acc_boot_N1000000.json :: replicas :: sha256 7cccb7b0ff900fb8
\newcommand{\MlstFtJcOneBootJets}{2{,}000{,}000}  % experiments/FIGS/data/v2/budget/jc1_acc_boot_N1000000.json :: n_test_jets :: sha256 7cccb7b0ff900fb8
\newcommand{\ComputeTotalGpuHours}{3{,}435}  % experiments/FIGS/data/v2/index.json :: GPU-hours of every reported pretraining run, training and validation passes :: sha256 c6c4138ee5c86423
\newcommand{\MlstComputeTotalGpuHours}{3{,}510}  % experiments/FIGS/data/v2/index.json :: GPU-hours of every pretraining run the paper reports, the kinematic-only model included :: sha256 c6c4138ee5c86423
\newcommand{\ComputeMainGpu}{V100}  % experiments/FIGS/data/v2/index.json :: GPU type with the most hours over the four label sets :: sha256 c6c4138ee5c86423
\newcommand{\ComputeMainGpuShareMin}{56\%}  % experiments/FIGS/data/v2/index.json :: lowest share of a label set's GPU-hours on that type :: sha256 c6c4138ee5c86423
\newcommand{\ComputeMainGpuShareMax}{71\%}  % experiments/FIGS/data/v2/index.json :: highest share of a label set's GPU-hours on that type :: sha256 c6c4138ee5c86423
\newcommand{\ComputeRunsDone}{27}  % experiments/FIGS/data/v2/index.json :: paper runs with a compute record (count) :: sha256 c6c4138ee5c86423
\newcommand{\ComputeGpuTypes}{L4, RTX 3090, RTX 4090 and V100}  % experiments/FIGS/data/v2/index.json :: GPU types any finished paper run trained on :: sha256 c6c4138ee5c86423
\newcommand{\MlstComputeGpuTypes}{L4, RTX 3090, RTX 4090 and V100}  % experiments/FIGS/data/v2/index.json :: GPU types of the pretraining runs the paper reports :: sha256 c6c4138ee5c86423
\newcommand{\MlstFtJcOneErrCoarseOverFullBootEFour}{\pending{the JetClass bootstrap at 1e4 (after the equal-stopping read-out)}}  % experiments/FIGS/data/v2/index.json :: not computed yet :: sha256 c6c4138ee5c86423 :: waits on the JetClass bootstrap at 1e4 (after the equal-stopping read-out)
\newcommand{\MlstLpBbQgZeroAccDropFull}{\pending{pair probes with the quark/gluon inputs zeroed (open item O10, not run)}}  % experiments/FIGS/data/v2/index.json :: not computed yet :: sha256 c6c4138ee5c86423 :: waits on pair probes with the quark/gluon inputs zeroed (open item O10, not run)
\newcommand{\MlstLpBbQgZeroAccDropCoarse}{\pending{pair probes with the quark/gluon inputs zeroed (open item O10, not run)}}  % experiments/FIGS/data/v2/index.json :: not computed yet :: sha256 c6c4138ee5c86423 :: waits on pair probes with the quark/gluon inputs zeroed (open item O10, not run)
\newcommand{\XbsRthirtyFullOverScratch}{\pending{X->bs fits and from-scratch runs}}  % experiments/FIGS/data/v2/index.json :: not computed yet :: sha256 c6c4138ee5c86423 :: waits on X->bs fits and from-scratch runs
\newcommand{\XbsRthirtyFullOverScratchPr}{\pending{X->bs fits and from-scratch runs}}  % experiments/FIGS/data/v2/index.json :: not computed yet :: sha256 c6c4138ee5c86423 :: waits on X->bs fits and from-scratch runs
\newcommand{\XbsRthirtyCoarseOverFull}{\pending{X->bs fits}}  % experiments/FIGS/data/v2/index.json :: not computed yet :: sha256 c6c4138ee5c86423 :: waits on X->bs fits
\newcommand{\XbsRthirtyCoarseOverFullPr}{\pending{X->bs fits}}  % experiments/FIGS/data/v2/index.json :: not computed yet :: sha256 c6c4138ee5c86423 :: waits on X->bs fits
\newcommand{\XbsRthirtyNodeOverTransferFull}{\pending{X->bs fits}}  % experiments/FIGS/data/v2/index.json :: not computed yet :: sha256 c6c4138ee5c86423 :: waits on X->bs fits
\newcommand{\XbsRthirtyNodeOverTransferFullPr}{\pending{X->bs fits}}  % experiments/FIGS/data/v2/index.json :: not computed yet :: sha256 c6c4138ee5c86423 :: waits on X->bs fits
\newcommand{\XbbRthirtyDirectOverTransferFull}{\pending{X->bb fits}}  % experiments/FIGS/data/v2/index.json :: not computed yet :: sha256 c6c4138ee5c86423 :: waits on X->bb fits
\newcommand{\XbbRthirtyDirectOverTransferFullPr}{\pending{X->bb fits}}  % experiments/FIGS/data/v2/index.json :: not computed yet :: sha256 c6c4138ee5c86423 :: waits on X->bb fits
\newcommand{\XbbRthirtyTransferCoarseOverFull}{\pending{X->bb fits}}  % experiments/FIGS/data/v2/index.json :: not computed yet :: sha256 c6c4138ee5c86423 :: waits on X->bb fits
\newcommand{\XbbRthirtyTransferCoarseOverFullPr}{\pending{X->bb fits}}  % experiments/FIGS/data/v2/index.json :: not computed yet :: sha256 c6c4138ee5c86423 :: waits on X->bb fits
             % the counts below are its macros
% Visual language: Sophon's Fig. 1 (arXiv:2405.12972): the document's serif font, light fills,
% rounded boxes, light grey panels behind the stages (PI 2026-11-02). The text is set in the
% document's font at the printed sizes 8.5-10 pt (MLST 8-12 pt), never scaled.
% Text width: 15 cm (MLST full width). The picture is 15.0 cm wide and is drawn at its final
% size, never scaled, so the font sizes hold.
% Optional: \def\FoneDropXbx{} before % Figure 1 (F1 of docs/FIGURE_PLAN.md): the design of the study, in TikZ.
% \input{fig_f1} from main.tex; this file holds one figure environment and nothing else.
%
% WHAT main.tex MUST LOAD (this file loads nothing itself):
%   \usepackage{tikz}
%   \usetikzlibrary{arrows.meta}          % the only TikZ library used (arrow tips)
%   \definecolor{c188}{HTML}{1F77B4}      % the configuration colours of experiments/FIGS/
%   \definecolor{c162}{HTML}{FF7F0E}      % mlst_style.py (matplotlib's tab10 in order, as in
%   \definecolor{c43}{HTML}{2CA02C}       % Sophon's figures); the variants are drawn in their
%   \definecolor{c17}{HTML}{D62728}       % label set's colour with a dotted border, as the
%   \definecolor{crand}{HTML}{17BECF}     % plotted figures draw variants dotted
%   \definecolor{cflav}{HTML}{9467BD}
%   \input{results_generated}             % the counts below are its macros
% Visual language: Sophon's Fig. 1 (arXiv:2405.12972): the document's serif font, light fills,
% rounded boxes, light grey panels behind the stages (PI 2026-11-02). The text is set in the
% document's font at the printed sizes 8.5-10 pt (MLST 8-12 pt), never scaled.
% Text width: 15 cm (MLST full width). The picture is 15.0 cm wide and is drawn at its final
% size, never scaled, so the font sizes hold.
% Optional: \def\FoneDropXbx{} before \input{fig_f1} removes X->bb/X->bs from panel (c)
% (PI 2026-11-02: if the samples are not staged by 12 Nov, the X->bb/X->bs section goes).
%
% Every class box is drawn from the label maps (checked by hand, 2026-11-02):
%   (a) configs/labelmaps/rung_label_maps.v1.csv, columns L188, L162, R42_Q1, R16_Q1. Rows: the ten
%       decays X->QQ/gg in the order bb bc bq | cc cs cq | ss sq qq gg (the 43-class b / c / light
%       groups are then contiguous), then ee mm | tauh-taue tauh-taum, then the 27 QCD classes.
%   (b) column: configs/labelmaps/rand_label_map.v2.csv, RAND2_p1 on the same rows: groups
%       0 0 0 0 0 9 9 9 0 3 | 11 11 | 15 4 | QCD, lettered A..F in order of first appearance.
%       strips: configs/labelmaps/flavour_pair_map.v2.csv, the 22 classes of orbit X_YY_QQQQ.
%       Order with a b quark: bbbb bbcc bbss bbqq qqqb bbcq ccbs ccbq qqbc qqbs bcsq;
%       without: cccc ccss ccqq ssss ssqq qqqq qqqc qqqs ccsq sscq qqcs.
%       FLAV_F0: all 22 in group 9. FLAV_F1: the 11 with a b in group 11, the rest 9.
%       FLAV_F1R: group 11 = bbss bbqq qqqb bbcq ccbq qqbc (with b, positions 2-5, 7, 8 from 0)
%       and ccss ccqq ssss qqqc sscq (without, positions 1-3, 6, 9); filled square = group 11.
% Geometry (cm): row height 0.35 from y = -1.18, gaps of 0.14 between the quark/gluon, lepton
% and QCD blocks, the QCD row 0.70; boxes inset 0.025 top and bottom; columns 1.45 wide with
% left edges 1.75, 3.28, 4.81, 6.34 (hierarchy) and 8.25 (random partition).
\providecommand{\Kfull}{\VocabSizeLoneeighteight}
\providecommand{\Kres}{\VocabSizeLonesixtwo}
\providecommand{\Kmid}{\VocabSizeRfourtwoQone}
\providecommand{\Kcoarse}{\VocabSizeRonesixQone}
\begin{figure}[tbp]
\centering
\ifdim\linewidth<14.99cm\typeout{fig_f1.tex: WARNING: the drawing is 15 cm wide but \string\linewidth\space
  is \the\linewidth; set a 15 cm text width (it is not scaled, to keep its 8.5 pt text)}\fi
\def\FoneText{\rmfamily\fontsize{8.5}{9.5}\selectfont}%
\def\FoneMain{\rmfamily\fontsize{9}{10.5}\selectfont}%
\def\FoneBold{\rmfamily\bfseries\fontsize{9}{10.5}\selectfont}%
\def\FoneTitle{\rmfamily\bfseries\fontsize{10}{12}\selectfont}%
\ifdefined\FoneDropXbx\def\FoneXbx{}\else\def\FoneXbx{\\$X\to b\bar b$, $X\to bs$}\fi
\begin{tikzpicture}[
  x=1cm, y=1cm,
  font=\FoneMain,
  ttl/.style={anchor=base west, font=\FoneTitle, inner sep=0pt},
  cls/.style={draw=black!55, line width=0.45pt, rounded corners=1.2pt},
  hdr/.style={draw=#1!70!black, fill=#1!22, line width=0.6pt, rounded corners=2.5pt, align=center,
              font=\FoneText, inner sep=0pt},
  inbox/.style={align=center, font=\FoneText, inner sep=0pt},
  letter/.style={font=\FoneBold, inner sep=0pt},
  rowlab/.style={anchor=east, inner sep=0pt},
  txt/.style={align=left, font=\FoneText, inner sep=0pt},
  sq/.style={draw=black!65, line width=0.4pt},
  chip/.style={draw=#1!70!black, fill=#1!22, line width=0.6pt, rounded corners=3pt,
               align=center, font=\FoneText, inner sep=0pt, anchor=north west},
  stage/.style={draw=black!45, fill=black!4, line width=0.6pt, rounded corners=3pt, align=left,
                font=\FoneText, inner sep=0.1cm},
  flow/.style={-{Latex[length=1.8mm, width=1.4mm]}, line width=0.7pt, black!75},
]
% ================================================================ (a) the class hierarchy
\node[ttl] at (0,0) {(a) Class hierarchy};
% faint bands on every other row, drawn first so the boxes sit on top
\foreach \rowtop in {-1.53,-2.23,-2.93,-3.63,-4.33,-5.17,-5.87}
  \fill[black!4] (0,\rowtop) rectangle (9.70,\rowtop-0.35);
% column headers
\node[rowlab, align=right, font=\FoneText] at (1.65,-0.68) {JetClass-II\\class};
\node[hdr=c188, minimum width=1.45cm, minimum height=0.80cm, anchor=north west] at (1.75,-0.28) {\Kfull\\classes};
\node[hdr=c162, minimum width=1.45cm, minimum height=0.80cm, anchor=north west] at (3.28,-0.28) {\Kres\\classes};
\node[hdr=c43,  minimum width=1.45cm, minimum height=0.80cm, anchor=north west] at (4.81,-0.28) {\Kmid\\classes};
\node[hdr=c17,  minimum width=1.45cm, minimum height=0.80cm, anchor=north west] at (6.34,-0.28) {\Kcoarse\\classes};
% row labels
\node[rowlab] at (1.65,-1.355) {$X\to b\bar b$};
\node[rowlab] at (1.65,-1.705) {$X\to bc$};
\node[rowlab] at (1.65,-2.055) {$X\to bq$};
\node[rowlab] at (1.65,-2.405) {$X\to c\bar c$};
\node[rowlab] at (1.65,-2.755) {$X\to cs$};
\node[rowlab] at (1.65,-3.105) {$X\to cq$};
\node[rowlab] at (1.65,-3.455) {$X\to s\bar s$};
\node[rowlab] at (1.65,-3.805) {$X\to sq$};
\node[rowlab] at (1.65,-4.155) {$X\to q\bar q$};
\node[rowlab] at (1.65,-4.505) {$X\to gg$};
\node[rowlab] at (1.65,-4.995) {$X\to ee$};
\node[rowlab] at (1.65,-5.345) {$X\to\mu\mu$};
\node[rowlab] at (1.65,-5.695) {$X\to\tau_h\tau_e$};
\node[rowlab] at (1.65,-6.045) {$X\to\tau_h\tau_\mu$};
\node[rowlab] at (1.65,-6.710) {QCD jets};
% 188 and 162 classes: every resonant decay is its own class
\foreach \rowtop in {-1.18,-1.53,-1.88,-2.23,-2.58,-2.93,-3.28,-3.63,-3.98,-4.33,-4.82,-5.17,-5.52,-5.87} {
  \draw[cls, fill=c188!15] (1.75,\rowtop-0.025) rectangle (3.20,\rowtop-0.325);
  \draw[cls, fill=c162!15] (3.28,\rowtop-0.025) rectangle (4.73,\rowtop-0.325);
}
\draw[cls, fill=c188!15] (1.75,-6.385) rectangle (3.20,-7.035);
\node[inbox] at (2.475,-6.71) {\LitJetClassTwoNQcdClasses\\classes};
\draw[cls, fill=c162!15] (3.28,-6.385) rectangle (4.73,-7.035);
\node[inbox] at (4.005,-6.71) {1 class};
% 43 classes: decay topology x {b, else c, else light quark or gluon}; e and mu merged
\draw[cls, fill=c43!15] (4.81,-1.205) rectangle (6.26,-2.205);
\node[inbox] at (5.535,-1.705) {with $b$};
\draw[cls, fill=c43!15] (4.81,-2.255) rectangle (6.26,-3.255);
\node[inbox] at (5.535,-2.755) {with $c$,\\no $b$};
\draw[cls, fill=c43!15] (4.81,-3.305) rectangle (6.26,-4.655);
\node[inbox] at (5.535,-3.980) {$u$, $d$, $s$\\or $g$};
\draw[cls, fill=c43!15] (4.81,-4.845) rectangle (6.26,-5.495);
\node[inbox] at (5.535,-5.170) {$\ell\ell$};
\draw[cls, fill=c43!15] (4.81,-5.545) rectangle (6.26,-6.195);
\node[inbox] at (5.535,-5.870) {$\tau_h\tau_\ell$};
\draw[cls, fill=c43!15] (4.81,-6.385) rectangle (6.26,-7.035);
\node[inbox] at (5.535,-6.710) {1 class};
% 17 classes: number and type of prongs, no flavor
\draw[cls, fill=c17!15] (6.34,-1.205) rectangle (7.79,-4.655);
\node[inbox] at (7.065,-2.930) {two\\quarks\\or gluons};
\draw[cls, fill=c17!15] (6.34,-4.845) rectangle (7.79,-5.495);
\node[inbox] at (7.065,-5.170) {$\ell\ell$};
\draw[cls, fill=c17!15] (6.34,-5.545) rectangle (7.79,-6.195);
\node[inbox] at (7.065,-5.870) {$\tau_h\tau_\ell$};
\draw[cls, fill=c17!15] (6.34,-6.385) rectangle (7.79,-7.035);
\node[inbox] at (7.065,-6.710) {1 class};
% what each merge removes
\node[anchor=north west, font=\FoneBold, inner sep=0pt] at (0,-7.27) {What each merge removes};
\node[txt, anchor=north west, text width=6.04cm] (m1) at (1.75,-7.65)
  {the $b$, $c$ and $s$ content of the QCD jets};
\node[txt, anchor=north west, text width=6.04cm] (m2) at ([yshift=-0.10cm]m1.south west)
  {the number of $b$ and $c$ quarks, light-quark flavor, quark against gluon, and lepton
   flavor ($e$ against $\mu$)};
\node[txt, anchor=north west, text width=6.04cm] (m3) at ([yshift=-0.10cm]m2.south west)
  {quark flavor ($b$, $c$ or light); the number and type of prongs remain};
\node[txt, anchor=north east] at ([xshift=-0.10cm]m1.north west) {\Kfull\,$\to$\,\Kres};
\node[txt, anchor=north east] at ([xshift=-0.10cm]m2.north west) {\Kres\,$\to$\,\Kmid};
\node[txt, anchor=north east] at ([xshift=-0.10cm]m3.north west) {\Kmid\,$\to$\,\Kcoarse};
% ================================================================ (b) the controls
\node[ttl] at (8.25,0) {(b) Controls, \Kcoarse{} classes each};
% the first random partition on the rows of (a); one letter = one class
\node[hdr=crand, minimum width=1.45cm, minimum height=0.80cm, anchor=north west] at (8.25,-0.28) {random\\partition};
\draw[cls, fill=crand!15] (8.25,-1.205) rectangle (9.70,-2.905);
\node[letter] at (8.975,-2.055) {A};
\draw[cls, fill=crand!15] (8.25,-2.955) rectangle (9.70,-3.955);
\node[letter] at (8.975,-3.455) {B};
\draw[cls, fill=crand!15] (8.25,-4.005) rectangle (9.70,-4.305);
\node[letter] at (8.975,-4.155) {A};
\draw[cls, fill=crand!15] (8.25,-4.355) rectangle (9.70,-4.655);
\node[letter] at (8.975,-4.505) {C};
\draw[cls, fill=crand!15] (8.25,-4.845) rectangle (9.70,-5.495);
\node[letter] at (8.975,-5.170) {D};
\draw[cls, fill=crand!15] (8.25,-5.545) rectangle (9.70,-5.845);
\node[letter] at (8.975,-5.695) {E};
\draw[cls, fill=crand!15] (8.25,-5.895) rectangle (9.70,-6.195);
\node[letter] at (8.975,-6.045) {F};
\draw[cls, fill=crand!15] (8.25,-6.385) rectangle (9.70,-7.035);
\node[inbox] at (8.975,-6.710) {1 class};
% how the random partitions are built
\node[hdr=crand, anchor=north west, text width=4.75cm, inner sep=0.1cm, align=left, font=\FoneBold]
  (rh) at (10.05,-0.28) {Random partitions (\RandVtwoNPartitions)};
\node[txt, anchor=north west, text width=4.95cm] (rb) at ([yshift=-0.12cm]rh.south west)
  {\textbullet\ groups of resonant classes drawn at random, plus one class of QCD jets\\[2pt]
   \textbullet\ group shares of training jets within \RandVtwoShareTol{} of the \Kcoarse-class groups\\[2pt]
   \textbullet\ best of \RandVtwoTrials{} draws: the fewest merges shared with the \Kcoarse-class level\\[2pt]
   \textbullet\ left: the first partition; boxes with the same letter are one class};
% how the random partitions with controlled flavor boundaries are built
\node[hdr=cflav, anchor=north west, text width=4.75cm, inner sep=0.1cm, align=left, font=\FoneBold]
  (fh) at ([yshift=-0.22cm]rb.south west) {Random partitions with controlled flavor boundaries};
\node[txt, anchor=north west, text width=4.95cm] (fn) at ([yshift=-0.12cm]fh.south west)
  {whole families of decays that differ only in quark flavor; two of them split one family,
   the \RandFlavSplitOrbitClasses{} decays \RandFlavSplitOrbit:};
\begin{scope}[shift={(fn.south west)}]
  % brackets over the decays with and without a b quark (squares at a 0.2 cm pitch)
  \node[txt, anchor=south] at (1.075,-0.45) {with $b$};
  \node[txt, anchor=south] at (3.425,-0.45) {without $b$};
  \draw[line width=0.5pt] (0,-0.54) -- (0,-0.47) -- (2.15,-0.47) -- (2.15,-0.54);
  \draw[line width=0.5pt] (2.35,-0.54) -- (2.35,-0.47) -- (4.50,-0.47) -- (4.50,-0.54);
  % no flavor boundary: all 22 in one class
  \node[txt, anchor=north west] at (0,-0.62) {no flavor boundary};
  \foreach \i in {0,...,10} {
    \draw[sq, fill=white] (0.2*\i,-1.10) rectangle ++(0.15,0.15);
    \draw[sq, fill=white] (2.35+0.2*\i,-1.10) rectangle ++(0.15,0.15);
  }
  % b-quark boundary: the decays with a b quark moved to another class
  \node[txt, anchor=north west] at (0,-1.22) {$b$-quark boundary};
  \foreach \i in {0,...,10} {
    \draw[sq, fill=cflav!70!black] (0.2*\i,-1.70) rectangle ++(0.15,0.15);
    \draw[sq, fill=white] (2.35+0.2*\i,-1.70) rectangle ++(0.15,0.15);
  }
  % random flavor boundary: the same number of decays, drawn at random
  \node[txt, anchor=north west] at (0,-1.82) {random flavor boundary};
  \foreach \i in {0,1,6,9,10}
    \draw[sq, fill=white] (0.2*\i,-2.30) rectangle ++(0.15,0.15);
  \foreach \i in {2,3,4,5,7,8}
    \draw[sq, fill=cflav!70!black] (0.2*\i,-2.30) rectangle ++(0.15,0.15);
  \foreach \i in {0,4,5,7,8,10}
    \draw[sq, fill=white] (2.35+0.2*\i,-2.30) rectangle ++(0.15,0.15);
  \foreach \i in {1,2,3,6,9}
    \draw[sq, fill=cflav!70!black] (2.35+0.2*\i,-2.30) rectangle ++(0.15,0.15);
  % key
  \draw[sq, fill=cflav!70!black] (0,-2.62) rectangle ++(0.15,0.15);
  \node[txt, anchor=west] (fkey) at (0.25,-2.545) {moved to another class};
\end{scope}
% ================================================================ (c) the pipeline
% (c) starts below the lower of the two columns above, whatever their text heights
\pgfpointanchor{m3}{south}\pgfgetlastxy{\FoneXa}{\FoneYa}%
\pgfpointanchor{fkey}{south}\pgfgetlastxy{\FoneXb}{\FoneYb}%
\pgfmathsetlengthmacro{\FoneYc}{min(\FoneYa,\FoneYb)-0.45cm}%
\coordinate (ctop) at (0pt,\FoneYc);
\begin{scope}[shift={(ctop)}]
  \node[ttl, anchor=north west] at (0,0) {(c) Pretraining and evaluation};
  % what every model of one seed shares
  \node[stage, text width=2.15cm] (seed) at (1.175,-1.96)
    {\textbf{Same for every model of one seed:} the training jets, their order and the
     initial weights, except the output layer};
  % one pretraining per label set: the configuration colours of every later figure
  \node[draw=black!45, fill=black!4, line width=0.6pt, rounded corners=3pt, minimum width=5.15cm,
        minimum height=3.08cm, anchor=north west, inner sep=0pt] (pre) at (2.81,-0.42) {};
  \node[anchor=north west, font=\FoneBold, inner sep=0pt] at (2.96,-0.52) {Pretraining on JetClass-II};
  \node[chip=c188, minimum width=1.13cm, minimum height=0.42cm] at (2.96,-0.95) {\Kfull};
  \node[chip=c162, minimum width=1.13cm, minimum height=0.42cm] at (4.20,-0.95) {\Kres};
  \node[chip=c43,  minimum width=1.13cm, minimum height=0.42cm] at (5.44,-0.95) {\Kmid};
  \node[chip=c17,  minimum width=1.13cm, minimum height=0.42cm] at (6.68,-0.95) {\Kcoarse};
  \node[chip=crand, minimum width=1.70cm, minimum height=0.70cm] at (2.96,-1.47) {random\\partitions};
  \node[chip=cflav, minimum width=3.04cm, minimum height=0.70cm] at (4.77,-1.47) {controlled flavor\\boundaries};
  \node[chip=c162, densely dotted, minimum width=1.13cm, minimum height=0.42cm] at (2.96,-2.32) {\Kres};
  \node[chip=c17,  densely dotted, minimum width=1.13cm, minimum height=0.42cm] at (4.20,-2.32) {\Kcoarse};
  \node[txt, anchor=west] at (5.44,-2.53) {with auxiliary\\mass regression};
  \node[chip=c188, densely dotted, minimum width=1.13cm, minimum height=0.42cm] at (2.96,-2.94) {\Kfull};
  \node[txt, anchor=west] at (4.20,-3.15) {kinematic inputs only};
  % evaluation
  \node[stage, text width=2.4cm] (fz) at (9.68,-0.95) {frozen features:\\linear probe,\\MLP probe};
  \node[stage, text width=2.4cm] (ft) at (9.68,-2.05) {fine-tuning\\of all weights};
  \node[stage, text width=2.4cm, draw=black, fill=white, line width=1.1pt] (sc) at (9.68,-3.15) {trained from\\scratch};
  \node[stage, text width=3.4cm, minimum height=3.2cm] (task) at (13.20,-2.00)
    {\textbf{Downstream tasks}\\JetClass-II: all classes\\JetClass-II: class pairs\\\quad(both class-balanced)\\JetClass\\
     top tagging\\quark/gluon, Pythia\\quark/gluon, Herwig\\jet mass\FoneXbx};
  \draw[flow] (seed.east) -- (seed.east -| pre.west);
  \draw[flow] (fz.west -| pre.east) -- (fz.west);
  \draw[flow] (ft.west -| pre.east) -- (ft.west);
  \draw[flow] (fz.east) -- (fz.east -| task.west);
  \draw[flow] (ft.east) -- (ft.east -| task.west);
  \draw[flow] (sc.east) -- (sc.east -| task.west);
\end{scope}
\end{tikzpicture}
\caption{Design of the study.
(a)~The class hierarchy on the decays of a heavy resonance~$X$ and on the jets from quantum
chromodynamics (QCD); $q$ is a $u$ or $d$ quark, $\ell$ an electron or a muon, and $\tau_h$,
$\tau_e$, $\tau_\mu$ a $\tau$ lepton decaying to hadrons, an electron or a muon. Each box is one
class of the label set above its column, except the QCD box of the \Kfull-class column (Sophon's
label set~\cite{sophon}), which stands for \LitJetClassTwoNQcdClasses{} classes. Each level merges
classes of the level to its left; the lines below name what each merge removes.
(b)~The controls: the first random partition, where boxes with the same letter are one class, and
the random partitions with controlled flavor boundaries on the decays \RandFlavSplitOrbit{} (one
square per decay; filled squares are moved to a second class).
(c)~Models of one seed share the training jets, their order and the initial weights, except the
output layer, the mass regression target, and the first layer and initial weights of the model with
kinematic inputs only (the particle four-momenta). The hierarchy levels have \MlstNRunsWord{} seeds
each, the models with the auxiliary mass regression \MlstMassNRunsWord{}, the random partitions with
controlled flavor boundaries \RandFlavRuns{} and each random partition \RandVtwoRuns{}. Frozen
features, the output of the backbone (the network below the classification head) with its weights held
fixed, feed linear probes (logistic regressions) and multilayer perceptron (MLP) probes; fine-tuning
trains all weights. Colors identify the label sets and controls in every figure; the plots give the three partitions with
controlled flavor boundaries three shades; dotted borders mark the two
variants, which the other figures draw as dotted lines in the color of their label set.}
\label{fig:design}
\end{figure}
 removes X->bb/X->bs from panel (c)
% (PI 2026-11-02: if the samples are not staged by 12 Nov, the X->bb/X->bs section goes).
%
% Every class box is drawn from the label maps (checked by hand, 2026-11-02):
%   (a) configs/labelmaps/rung_label_maps.v1.csv, columns L188, L162, R42_Q1, R16_Q1. Rows: the ten
%       decays X->QQ/gg in the order bb bc bq | cc cs cq | ss sq qq gg (the 43-class b / c / light
%       groups are then contiguous), then ee mm | tauh-taue tauh-taum, then the 27 QCD classes.
%   (b) column: configs/labelmaps/rand_label_map.v2.csv, RAND2_p1 on the same rows: groups
%       0 0 0 0 0 9 9 9 0 3 | 11 11 | 15 4 | QCD, lettered A..F in order of first appearance.
%       strips: configs/labelmaps/flavour_pair_map.v2.csv, the 22 classes of orbit X_YY_QQQQ.
%       Order with a b quark: bbbb bbcc bbss bbqq qqqb bbcq ccbs ccbq qqbc qqbs bcsq;
%       without: cccc ccss ccqq ssss ssqq qqqq qqqc qqqs ccsq sscq qqcs.
%       FLAV_F0: all 22 in group 9. FLAV_F1: the 11 with a b in group 11, the rest 9.
%       FLAV_F1R: group 11 = bbss bbqq qqqb bbcq ccbq qqbc (with b, positions 2-5, 7, 8 from 0)
%       and ccss ccqq ssss qqqc sscq (without, positions 1-3, 6, 9); filled square = group 11.
% Geometry (cm): row height 0.35 from y = -1.18, gaps of 0.14 between the quark/gluon, lepton
% and QCD blocks, the QCD row 0.70; boxes inset 0.025 top and bottom; columns 1.45 wide with
% left edges 1.75, 3.28, 4.81, 6.34 (hierarchy) and 8.25 (random partition).
\providecommand{\Kfull}{\VocabSizeLoneeighteight}
\providecommand{\Kres}{\VocabSizeLonesixtwo}
\providecommand{\Kmid}{\VocabSizeRfourtwoQone}
\providecommand{\Kcoarse}{\VocabSizeRonesixQone}
\begin{figure}[tbp]
\centering
\ifdim\linewidth<14.99cm\typeout{fig_f1.tex: WARNING: the drawing is 15 cm wide but \string\linewidth\space
  is \the\linewidth; set a 15 cm text width (it is not scaled, to keep its 8.5 pt text)}\fi
\def\FoneText{\rmfamily\fontsize{8.5}{9.5}\selectfont}%
\def\FoneMain{\rmfamily\fontsize{9}{10.5}\selectfont}%
\def\FoneBold{\rmfamily\bfseries\fontsize{9}{10.5}\selectfont}%
\def\FoneTitle{\rmfamily\bfseries\fontsize{10}{12}\selectfont}%
\ifdefined\FoneDropXbx\def\FoneXbx{}\else\def\FoneXbx{\\$X\to b\bar b$, $X\to bs$}\fi
\begin{tikzpicture}[
  x=1cm, y=1cm,
  font=\FoneMain,
  ttl/.style={anchor=base west, font=\FoneTitle, inner sep=0pt},
  cls/.style={draw=black!55, line width=0.45pt, rounded corners=1.2pt},
  hdr/.style={draw=#1!70!black, fill=#1!22, line width=0.6pt, rounded corners=2.5pt, align=center,
              font=\FoneText, inner sep=0pt},
  inbox/.style={align=center, font=\FoneText, inner sep=0pt},
  letter/.style={font=\FoneBold, inner sep=0pt},
  rowlab/.style={anchor=east, inner sep=0pt},
  txt/.style={align=left, font=\FoneText, inner sep=0pt},
  sq/.style={draw=black!65, line width=0.4pt},
  chip/.style={draw=#1!70!black, fill=#1!22, line width=0.6pt, rounded corners=3pt,
               align=center, font=\FoneText, inner sep=0pt, anchor=north west},
  stage/.style={draw=black!45, fill=black!4, line width=0.6pt, rounded corners=3pt, align=left,
                font=\FoneText, inner sep=0.1cm},
  flow/.style={-{Latex[length=1.8mm, width=1.4mm]}, line width=0.7pt, black!75},
]
% ================================================================ (a) the class hierarchy
\node[ttl] at (0,0) {(a) Class hierarchy};
% faint bands on every other row, drawn first so the boxes sit on top
\foreach \rowtop in {-1.53,-2.23,-2.93,-3.63,-4.33,-5.17,-5.87}
  \fill[black!4] (0,\rowtop) rectangle (9.70,\rowtop-0.35);
% column headers
\node[rowlab, align=right, font=\FoneText] at (1.65,-0.68) {JetClass-II\\class};
\node[hdr=c188, minimum width=1.45cm, minimum height=0.80cm, anchor=north west] at (1.75,-0.28) {\Kfull\\classes};
\node[hdr=c162, minimum width=1.45cm, minimum height=0.80cm, anchor=north west] at (3.28,-0.28) {\Kres\\classes};
\node[hdr=c43,  minimum width=1.45cm, minimum height=0.80cm, anchor=north west] at (4.81,-0.28) {\Kmid\\classes};
\node[hdr=c17,  minimum width=1.45cm, minimum height=0.80cm, anchor=north west] at (6.34,-0.28) {\Kcoarse\\classes};
% row labels
\node[rowlab] at (1.65,-1.355) {$X\to b\bar b$};
\node[rowlab] at (1.65,-1.705) {$X\to bc$};
\node[rowlab] at (1.65,-2.055) {$X\to bq$};
\node[rowlab] at (1.65,-2.405) {$X\to c\bar c$};
\node[rowlab] at (1.65,-2.755) {$X\to cs$};
\node[rowlab] at (1.65,-3.105) {$X\to cq$};
\node[rowlab] at (1.65,-3.455) {$X\to s\bar s$};
\node[rowlab] at (1.65,-3.805) {$X\to sq$};
\node[rowlab] at (1.65,-4.155) {$X\to q\bar q$};
\node[rowlab] at (1.65,-4.505) {$X\to gg$};
\node[rowlab] at (1.65,-4.995) {$X\to ee$};
\node[rowlab] at (1.65,-5.345) {$X\to\mu\mu$};
\node[rowlab] at (1.65,-5.695) {$X\to\tau_h\tau_e$};
\node[rowlab] at (1.65,-6.045) {$X\to\tau_h\tau_\mu$};
\node[rowlab] at (1.65,-6.710) {QCD jets};
% 188 and 162 classes: every resonant decay is its own class
\foreach \rowtop in {-1.18,-1.53,-1.88,-2.23,-2.58,-2.93,-3.28,-3.63,-3.98,-4.33,-4.82,-5.17,-5.52,-5.87} {
  \draw[cls, fill=c188!15] (1.75,\rowtop-0.025) rectangle (3.20,\rowtop-0.325);
  \draw[cls, fill=c162!15] (3.28,\rowtop-0.025) rectangle (4.73,\rowtop-0.325);
}
\draw[cls, fill=c188!15] (1.75,-6.385) rectangle (3.20,-7.035);
\node[inbox] at (2.475,-6.71) {\LitJetClassTwoNQcdClasses\\classes};
\draw[cls, fill=c162!15] (3.28,-6.385) rectangle (4.73,-7.035);
\node[inbox] at (4.005,-6.71) {1 class};
% 43 classes: decay topology x {b, else c, else light quark or gluon}; e and mu merged
\draw[cls, fill=c43!15] (4.81,-1.205) rectangle (6.26,-2.205);
\node[inbox] at (5.535,-1.705) {with $b$};
\draw[cls, fill=c43!15] (4.81,-2.255) rectangle (6.26,-3.255);
\node[inbox] at (5.535,-2.755) {with $c$,\\no $b$};
\draw[cls, fill=c43!15] (4.81,-3.305) rectangle (6.26,-4.655);
\node[inbox] at (5.535,-3.980) {$u$, $d$, $s$\\or $g$};
\draw[cls, fill=c43!15] (4.81,-4.845) rectangle (6.26,-5.495);
\node[inbox] at (5.535,-5.170) {$\ell\ell$};
\draw[cls, fill=c43!15] (4.81,-5.545) rectangle (6.26,-6.195);
\node[inbox] at (5.535,-5.870) {$\tau_h\tau_\ell$};
\draw[cls, fill=c43!15] (4.81,-6.385) rectangle (6.26,-7.035);
\node[inbox] at (5.535,-6.710) {1 class};
% 17 classes: number and type of prongs, no flavor
\draw[cls, fill=c17!15] (6.34,-1.205) rectangle (7.79,-4.655);
\node[inbox] at (7.065,-2.930) {two\\quarks\\or gluons};
\draw[cls, fill=c17!15] (6.34,-4.845) rectangle (7.79,-5.495);
\node[inbox] at (7.065,-5.170) {$\ell\ell$};
\draw[cls, fill=c17!15] (6.34,-5.545) rectangle (7.79,-6.195);
\node[inbox] at (7.065,-5.870) {$\tau_h\tau_\ell$};
\draw[cls, fill=c17!15] (6.34,-6.385) rectangle (7.79,-7.035);
\node[inbox] at (7.065,-6.710) {1 class};
% what each merge removes
\node[anchor=north west, font=\FoneBold, inner sep=0pt] at (0,-7.27) {What each merge removes};
\node[txt, anchor=north west, text width=6.04cm] (m1) at (1.75,-7.65)
  {the $b$, $c$ and $s$ content of the QCD jets};
\node[txt, anchor=north west, text width=6.04cm] (m2) at ([yshift=-0.10cm]m1.south west)
  {the number of $b$ and $c$ quarks, light-quark flavor, quark against gluon, and lepton
   flavor ($e$ against $\mu$)};
\node[txt, anchor=north west, text width=6.04cm] (m3) at ([yshift=-0.10cm]m2.south west)
  {quark flavor ($b$, $c$ or light); the number and type of prongs remain};
\node[txt, anchor=north east] at ([xshift=-0.10cm]m1.north west) {\Kfull\,$\to$\,\Kres};
\node[txt, anchor=north east] at ([xshift=-0.10cm]m2.north west) {\Kres\,$\to$\,\Kmid};
\node[txt, anchor=north east] at ([xshift=-0.10cm]m3.north west) {\Kmid\,$\to$\,\Kcoarse};
% ================================================================ (b) the controls
\node[ttl] at (8.25,0) {(b) Controls, \Kcoarse{} classes each};
% the first random partition on the rows of (a); one letter = one class
\node[hdr=crand, minimum width=1.45cm, minimum height=0.80cm, anchor=north west] at (8.25,-0.28) {random\\partition};
\draw[cls, fill=crand!15] (8.25,-1.205) rectangle (9.70,-2.905);
\node[letter] at (8.975,-2.055) {A};
\draw[cls, fill=crand!15] (8.25,-2.955) rectangle (9.70,-3.955);
\node[letter] at (8.975,-3.455) {B};
\draw[cls, fill=crand!15] (8.25,-4.005) rectangle (9.70,-4.305);
\node[letter] at (8.975,-4.155) {A};
\draw[cls, fill=crand!15] (8.25,-4.355) rectangle (9.70,-4.655);
\node[letter] at (8.975,-4.505) {C};
\draw[cls, fill=crand!15] (8.25,-4.845) rectangle (9.70,-5.495);
\node[letter] at (8.975,-5.170) {D};
\draw[cls, fill=crand!15] (8.25,-5.545) rectangle (9.70,-5.845);
\node[letter] at (8.975,-5.695) {E};
\draw[cls, fill=crand!15] (8.25,-5.895) rectangle (9.70,-6.195);
\node[letter] at (8.975,-6.045) {F};
\draw[cls, fill=crand!15] (8.25,-6.385) rectangle (9.70,-7.035);
\node[inbox] at (8.975,-6.710) {1 class};
% how the random partitions are built
\node[hdr=crand, anchor=north west, text width=4.75cm, inner sep=0.1cm, align=left, font=\FoneBold]
  (rh) at (10.05,-0.28) {Random partitions (\RandVtwoNPartitions)};
\node[txt, anchor=north west, text width=4.95cm] (rb) at ([yshift=-0.12cm]rh.south west)
  {\textbullet\ groups of resonant classes drawn at random, plus one class of QCD jets\\[2pt]
   \textbullet\ group shares of training jets within \RandVtwoShareTol{} of the \Kcoarse-class groups\\[2pt]
   \textbullet\ best of \RandVtwoTrials{} draws: the fewest merges shared with the \Kcoarse-class level\\[2pt]
   \textbullet\ left: the first partition; boxes with the same letter are one class};
% how the random partitions with controlled flavor boundaries are built
\node[hdr=cflav, anchor=north west, text width=4.75cm, inner sep=0.1cm, align=left, font=\FoneBold]
  (fh) at ([yshift=-0.22cm]rb.south west) {Random partitions with controlled flavor boundaries};
\node[txt, anchor=north west, text width=4.95cm] (fn) at ([yshift=-0.12cm]fh.south west)
  {whole families of decays that differ only in quark flavor; two of them split one family,
   the \RandFlavSplitOrbitClasses{} decays \RandFlavSplitOrbit:};
\begin{scope}[shift={(fn.south west)}]
  % brackets over the decays with and without a b quark (squares at a 0.2 cm pitch)
  \node[txt, anchor=south] at (1.075,-0.45) {with $b$};
  \node[txt, anchor=south] at (3.425,-0.45) {without $b$};
  \draw[line width=0.5pt] (0,-0.54) -- (0,-0.47) -- (2.15,-0.47) -- (2.15,-0.54);
  \draw[line width=0.5pt] (2.35,-0.54) -- (2.35,-0.47) -- (4.50,-0.47) -- (4.50,-0.54);
  % no flavor boundary: all 22 in one class
  \node[txt, anchor=north west] at (0,-0.62) {no flavor boundary};
  \foreach \i in {0,...,10} {
    \draw[sq, fill=white] (0.2*\i,-1.10) rectangle ++(0.15,0.15);
    \draw[sq, fill=white] (2.35+0.2*\i,-1.10) rectangle ++(0.15,0.15);
  }
  % b-quark boundary: the decays with a b quark moved to another class
  \node[txt, anchor=north west] at (0,-1.22) {$b$-quark boundary};
  \foreach \i in {0,...,10} {
    \draw[sq, fill=cflav!70!black] (0.2*\i,-1.70) rectangle ++(0.15,0.15);
    \draw[sq, fill=white] (2.35+0.2*\i,-1.70) rectangle ++(0.15,0.15);
  }
  % random flavor boundary: the same number of decays, drawn at random
  \node[txt, anchor=north west] at (0,-1.82) {random flavor boundary};
  \foreach \i in {0,1,6,9,10}
    \draw[sq, fill=white] (0.2*\i,-2.30) rectangle ++(0.15,0.15);
  \foreach \i in {2,3,4,5,7,8}
    \draw[sq, fill=cflav!70!black] (0.2*\i,-2.30) rectangle ++(0.15,0.15);
  \foreach \i in {0,4,5,7,8,10}
    \draw[sq, fill=white] (2.35+0.2*\i,-2.30) rectangle ++(0.15,0.15);
  \foreach \i in {1,2,3,6,9}
    \draw[sq, fill=cflav!70!black] (2.35+0.2*\i,-2.30) rectangle ++(0.15,0.15);
  % key
  \draw[sq, fill=cflav!70!black] (0,-2.62) rectangle ++(0.15,0.15);
  \node[txt, anchor=west] (fkey) at (0.25,-2.545) {moved to another class};
\end{scope}
% ================================================================ (c) the pipeline
% (c) starts below the lower of the two columns above, whatever their text heights
\pgfpointanchor{m3}{south}\pgfgetlastxy{\FoneXa}{\FoneYa}%
\pgfpointanchor{fkey}{south}\pgfgetlastxy{\FoneXb}{\FoneYb}%
\pgfmathsetlengthmacro{\FoneYc}{min(\FoneYa,\FoneYb)-0.45cm}%
\coordinate (ctop) at (0pt,\FoneYc);
\begin{scope}[shift={(ctop)}]
  \node[ttl, anchor=north west] at (0,0) {(c) Pretraining and evaluation};
  % what every model of one seed shares
  \node[stage, text width=2.15cm] (seed) at (1.175,-1.96)
    {\textbf{Same for every model of one seed:} the training jets, their order and the
     initial weights, except the output layer};
  % one pretraining per label set: the configuration colours of every later figure
  \node[draw=black!45, fill=black!4, line width=0.6pt, rounded corners=3pt, minimum width=5.15cm,
        minimum height=3.08cm, anchor=north west, inner sep=0pt] (pre) at (2.81,-0.42) {};
  \node[anchor=north west, font=\FoneBold, inner sep=0pt] at (2.96,-0.52) {Pretraining on JetClass-II};
  \node[chip=c188, minimum width=1.13cm, minimum height=0.42cm] at (2.96,-0.95) {\Kfull};
  \node[chip=c162, minimum width=1.13cm, minimum height=0.42cm] at (4.20,-0.95) {\Kres};
  \node[chip=c43,  minimum width=1.13cm, minimum height=0.42cm] at (5.44,-0.95) {\Kmid};
  \node[chip=c17,  minimum width=1.13cm, minimum height=0.42cm] at (6.68,-0.95) {\Kcoarse};
  \node[chip=crand, minimum width=1.70cm, minimum height=0.70cm] at (2.96,-1.47) {random\\partitions};
  \node[chip=cflav, minimum width=3.04cm, minimum height=0.70cm] at (4.77,-1.47) {controlled flavor\\boundaries};
  \node[chip=c162, densely dotted, minimum width=1.13cm, minimum height=0.42cm] at (2.96,-2.32) {\Kres};
  \node[chip=c17,  densely dotted, minimum width=1.13cm, minimum height=0.42cm] at (4.20,-2.32) {\Kcoarse};
  \node[txt, anchor=west] at (5.44,-2.53) {with auxiliary\\mass regression};
  \node[chip=c188, densely dotted, minimum width=1.13cm, minimum height=0.42cm] at (2.96,-2.94) {\Kfull};
  \node[txt, anchor=west] at (4.20,-3.15) {kinematic inputs only};
  % evaluation
  \node[stage, text width=2.4cm] (fz) at (9.68,-0.95) {frozen features:\\linear probe,\\MLP probe};
  \node[stage, text width=2.4cm] (ft) at (9.68,-2.05) {fine-tuning\\of all weights};
  \node[stage, text width=2.4cm, draw=black, fill=white, line width=1.1pt] (sc) at (9.68,-3.15) {trained from\\scratch};
  \node[stage, text width=3.4cm, minimum height=3.2cm] (task) at (13.20,-2.00)
    {\textbf{Downstream tasks}\\JetClass-II: all classes\\JetClass-II: class pairs\\\quad(both class-balanced)\\JetClass\\
     top tagging\\quark/gluon, Pythia\\quark/gluon, Herwig\\jet mass\FoneXbx};
  \draw[flow] (seed.east) -- (seed.east -| pre.west);
  \draw[flow] (fz.west -| pre.east) -- (fz.west);
  \draw[flow] (ft.west -| pre.east) -- (ft.west);
  \draw[flow] (fz.east) -- (fz.east -| task.west);
  \draw[flow] (ft.east) -- (ft.east -| task.west);
  \draw[flow] (sc.east) -- (sc.east -| task.west);
\end{scope}
\end{tikzpicture}
\caption{Design of the study.
(a)~The class hierarchy on the decays of a heavy resonance~$X$ and on the jets from quantum
chromodynamics (QCD); $q$ is a $u$ or $d$ quark, $\ell$ an electron or a muon, and $\tau_h$,
$\tau_e$, $\tau_\mu$ a $\tau$ lepton decaying to hadrons, an electron or a muon. Each box is one
class of the label set above its column, except the QCD box of the \Kfull-class column (Sophon's
label set~\cite{sophon}), which stands for \LitJetClassTwoNQcdClasses{} classes. Each level merges
classes of the level to its left; the lines below name what each merge removes.
(b)~The controls: the first random partition, where boxes with the same letter are one class, and
the random partitions with controlled flavor boundaries on the decays \RandFlavSplitOrbit{} (one
square per decay; filled squares are moved to a second class).
(c)~Models of one seed share the training jets, their order and the initial weights, except the
output layer, the mass regression target, and the first layer and initial weights of the model with
kinematic inputs only (the particle four-momenta). The hierarchy levels have \MlstNRunsWord{} seeds
each, the models with the auxiliary mass regression \MlstMassNRunsWord{}, the random partitions with
controlled flavor boundaries \RandFlavRuns{} and each random partition \RandVtwoRuns{}. Frozen
features, the output of the backbone (the network below the classification head) with its weights held
fixed, feed linear probes (logistic regressions) and multilayer perceptron (MLP) probes; fine-tuning
trains all weights. Colors identify the label sets and controls in every figure; the plots give the three partitions with
controlled flavor boundaries three shades; dotted borders mark the two
variants, which the other figures draw as dotted lines in the color of their label set.}
\label{fig:design}
\end{figure}
 removes X->bb/X->bs from panel (c)
% (PI 2026-11-02: if the samples are not staged by 12 Nov, the X->bb/X->bs section goes).
%
% Every class box is drawn from the label maps (checked by hand, 2026-11-02):
%   (a) configs/labelmaps/rung_label_maps.v1.csv, columns L188, L162, R42_Q1, R16_Q1. Rows: the ten
%       decays X->QQ/gg in the order bb bc bq | cc cs cq | ss sq qq gg (the 43-class b / c / light
%       groups are then contiguous), then ee mm | tauh-taue tauh-taum, then the 27 QCD classes.
%   (b) column: configs/labelmaps/rand_label_map.v2.csv, RAND2_p1 on the same rows: groups
%       0 0 0 0 0 9 9 9 0 3 | 11 11 | 15 4 | QCD, lettered A..F in order of first appearance.
%       strips: configs/labelmaps/flavour_pair_map.v2.csv, the 22 classes of orbit X_YY_QQQQ.
%       Order with a b quark: bbbb bbcc bbss bbqq qqqb bbcq ccbs ccbq qqbc qqbs bcsq;
%       without: cccc ccss ccqq ssss ssqq qqqq qqqc qqqs ccsq sscq qqcs.
%       FLAV_F0: all 22 in group 9. FLAV_F1: the 11 with a b in group 11, the rest 9.
%       FLAV_F1R: group 11 = bbss bbqq qqqb bbcq ccbq qqbc (with b, positions 2-5, 7, 8 from 0)
%       and ccss ccqq ssss qqqc sscq (without, positions 1-3, 6, 9); filled square = group 11.
% Geometry (cm): row height 0.35 from y = -1.18, gaps of 0.14 between the quark/gluon, lepton
% and QCD blocks, the QCD row 0.70; boxes inset 0.025 top and bottom; columns 1.45 wide with
% left edges 1.75, 3.28, 4.81, 6.34 (hierarchy) and 8.25 (random partition).
\providecommand{\Kfull}{\VocabSizeLoneeighteight}
\providecommand{\Kres}{\VocabSizeLonesixtwo}
\providecommand{\Kmid}{\VocabSizeRfourtwoQone}
\providecommand{\Kcoarse}{\VocabSizeRonesixQone}
\begin{figure}[tbp]
\centering
\ifdim\linewidth<14.99cm\typeout{fig_f1.tex: WARNING: the drawing is 15 cm wide but \string\linewidth\space
  is \the\linewidth; set a 15 cm text width (it is not scaled, to keep its 8.5 pt text)}\fi
\def\FoneText{\rmfamily\fontsize{8.5}{9.5}\selectfont}%
\def\FoneMain{\rmfamily\fontsize{9}{10.5}\selectfont}%
\def\FoneBold{\rmfamily\bfseries\fontsize{9}{10.5}\selectfont}%
\def\FoneTitle{\rmfamily\bfseries\fontsize{10}{12}\selectfont}%
\ifdefined\FoneDropXbx\def\FoneXbx{}\else\def\FoneXbx{\\$X\to b\bar b$, $X\to bs$}\fi
\begin{tikzpicture}[
  x=1cm, y=1cm,
  font=\FoneMain,
  ttl/.style={anchor=base west, font=\FoneTitle, inner sep=0pt},
  cls/.style={draw=black!55, line width=0.45pt, rounded corners=1.2pt},
  hdr/.style={draw=#1!70!black, fill=#1!22, line width=0.6pt, rounded corners=2.5pt, align=center,
              font=\FoneText, inner sep=0pt},
  inbox/.style={align=center, font=\FoneText, inner sep=0pt},
  letter/.style={font=\FoneBold, inner sep=0pt},
  rowlab/.style={anchor=east, inner sep=0pt},
  txt/.style={align=left, font=\FoneText, inner sep=0pt},
  sq/.style={draw=black!65, line width=0.4pt},
  chip/.style={draw=#1!70!black, fill=#1!22, line width=0.6pt, rounded corners=3pt,
               align=center, font=\FoneText, inner sep=0pt, anchor=north west},
  stage/.style={draw=black!45, fill=black!4, line width=0.6pt, rounded corners=3pt, align=left,
                font=\FoneText, inner sep=0.1cm},
  flow/.style={-{Latex[length=1.8mm, width=1.4mm]}, line width=0.7pt, black!75},
]
% ================================================================ (a) the class hierarchy
\node[ttl] at (0,0) {(a) Class hierarchy};
% faint bands on every other row, drawn first so the boxes sit on top
\foreach \rowtop in {-1.53,-2.23,-2.93,-3.63,-4.33,-5.17,-5.87}
  \fill[black!4] (0,\rowtop) rectangle (9.70,\rowtop-0.35);
% column headers
\node[rowlab, align=right, font=\FoneText] at (1.65,-0.68) {JetClass-II\\class};
\node[hdr=c188, minimum width=1.45cm, minimum height=0.80cm, anchor=north west] at (1.75,-0.28) {\Kfull\\classes};
\node[hdr=c162, minimum width=1.45cm, minimum height=0.80cm, anchor=north west] at (3.28,-0.28) {\Kres\\classes};
\node[hdr=c43,  minimum width=1.45cm, minimum height=0.80cm, anchor=north west] at (4.81,-0.28) {\Kmid\\classes};
\node[hdr=c17,  minimum width=1.45cm, minimum height=0.80cm, anchor=north west] at (6.34,-0.28) {\Kcoarse\\classes};
% row labels
\node[rowlab] at (1.65,-1.355) {$X\to b\bar b$};
\node[rowlab] at (1.65,-1.705) {$X\to bc$};
\node[rowlab] at (1.65,-2.055) {$X\to bq$};
\node[rowlab] at (1.65,-2.405) {$X\to c\bar c$};
\node[rowlab] at (1.65,-2.755) {$X\to cs$};
\node[rowlab] at (1.65,-3.105) {$X\to cq$};
\node[rowlab] at (1.65,-3.455) {$X\to s\bar s$};
\node[rowlab] at (1.65,-3.805) {$X\to sq$};
\node[rowlab] at (1.65,-4.155) {$X\to q\bar q$};
\node[rowlab] at (1.65,-4.505) {$X\to gg$};
\node[rowlab] at (1.65,-4.995) {$X\to ee$};
\node[rowlab] at (1.65,-5.345) {$X\to\mu\mu$};
\node[rowlab] at (1.65,-5.695) {$X\to\tau_h\tau_e$};
\node[rowlab] at (1.65,-6.045) {$X\to\tau_h\tau_\mu$};
\node[rowlab] at (1.65,-6.710) {QCD jets};
% 188 and 162 classes: every resonant decay is its own class
\foreach \rowtop in {-1.18,-1.53,-1.88,-2.23,-2.58,-2.93,-3.28,-3.63,-3.98,-4.33,-4.82,-5.17,-5.52,-5.87} {
  \draw[cls, fill=c188!15] (1.75,\rowtop-0.025) rectangle (3.20,\rowtop-0.325);
  \draw[cls, fill=c162!15] (3.28,\rowtop-0.025) rectangle (4.73,\rowtop-0.325);
}
\draw[cls, fill=c188!15] (1.75,-6.385) rectangle (3.20,-7.035);
\node[inbox] at (2.475,-6.71) {\LitJetClassTwoNQcdClasses\\classes};
\draw[cls, fill=c162!15] (3.28,-6.385) rectangle (4.73,-7.035);
\node[inbox] at (4.005,-6.71) {1 class};
% 43 classes: decay topology x {b, else c, else light quark or gluon}; e and mu merged
\draw[cls, fill=c43!15] (4.81,-1.205) rectangle (6.26,-2.205);
\node[inbox] at (5.535,-1.705) {with $b$};
\draw[cls, fill=c43!15] (4.81,-2.255) rectangle (6.26,-3.255);
\node[inbox] at (5.535,-2.755) {with $c$,\\no $b$};
\draw[cls, fill=c43!15] (4.81,-3.305) rectangle (6.26,-4.655);
\node[inbox] at (5.535,-3.980) {$u$, $d$, $s$\\or $g$};
\draw[cls, fill=c43!15] (4.81,-4.845) rectangle (6.26,-5.495);
\node[inbox] at (5.535,-5.170) {$\ell\ell$};
\draw[cls, fill=c43!15] (4.81,-5.545) rectangle (6.26,-6.195);
\node[inbox] at (5.535,-5.870) {$\tau_h\tau_\ell$};
\draw[cls, fill=c43!15] (4.81,-6.385) rectangle (6.26,-7.035);
\node[inbox] at (5.535,-6.710) {1 class};
% 17 classes: number and type of prongs, no flavor
\draw[cls, fill=c17!15] (6.34,-1.205) rectangle (7.79,-4.655);
\node[inbox] at (7.065,-2.930) {two\\quarks\\or gluons};
\draw[cls, fill=c17!15] (6.34,-4.845) rectangle (7.79,-5.495);
\node[inbox] at (7.065,-5.170) {$\ell\ell$};
\draw[cls, fill=c17!15] (6.34,-5.545) rectangle (7.79,-6.195);
\node[inbox] at (7.065,-5.870) {$\tau_h\tau_\ell$};
\draw[cls, fill=c17!15] (6.34,-6.385) rectangle (7.79,-7.035);
\node[inbox] at (7.065,-6.710) {1 class};
% what each merge removes
\node[anchor=north west, font=\FoneBold, inner sep=0pt] at (0,-7.27) {What each merge removes};
\node[txt, anchor=north west, text width=6.04cm] (m1) at (1.75,-7.65)
  {the $b$, $c$ and $s$ content of the QCD jets};
\node[txt, anchor=north west, text width=6.04cm] (m2) at ([yshift=-0.10cm]m1.south west)
  {the number of $b$ and $c$ quarks, light-quark flavor, quark against gluon, and lepton
   flavor ($e$ against $\mu$)};
\node[txt, anchor=north west, text width=6.04cm] (m3) at ([yshift=-0.10cm]m2.south west)
  {quark flavor ($b$, $c$ or light); the number and type of prongs remain};
\node[txt, anchor=north east] at ([xshift=-0.10cm]m1.north west) {\Kfull\,$\to$\,\Kres};
\node[txt, anchor=north east] at ([xshift=-0.10cm]m2.north west) {\Kres\,$\to$\,\Kmid};
\node[txt, anchor=north east] at ([xshift=-0.10cm]m3.north west) {\Kmid\,$\to$\,\Kcoarse};
% ================================================================ (b) the controls
\node[ttl] at (8.25,0) {(b) Controls, \Kcoarse{} classes each};
% the first random partition on the rows of (a); one letter = one class
\node[hdr=crand, minimum width=1.45cm, minimum height=0.80cm, anchor=north west] at (8.25,-0.28) {random\\partition};
\draw[cls, fill=crand!15] (8.25,-1.205) rectangle (9.70,-2.905);
\node[letter] at (8.975,-2.055) {A};
\draw[cls, fill=crand!15] (8.25,-2.955) rectangle (9.70,-3.955);
\node[letter] at (8.975,-3.455) {B};
\draw[cls, fill=crand!15] (8.25,-4.005) rectangle (9.70,-4.305);
\node[letter] at (8.975,-4.155) {A};
\draw[cls, fill=crand!15] (8.25,-4.355) rectangle (9.70,-4.655);
\node[letter] at (8.975,-4.505) {C};
\draw[cls, fill=crand!15] (8.25,-4.845) rectangle (9.70,-5.495);
\node[letter] at (8.975,-5.170) {D};
\draw[cls, fill=crand!15] (8.25,-5.545) rectangle (9.70,-5.845);
\node[letter] at (8.975,-5.695) {E};
\draw[cls, fill=crand!15] (8.25,-5.895) rectangle (9.70,-6.195);
\node[letter] at (8.975,-6.045) {F};
\draw[cls, fill=crand!15] (8.25,-6.385) rectangle (9.70,-7.035);
\node[inbox] at (8.975,-6.710) {1 class};
% how the random partitions are built
\node[hdr=crand, anchor=north west, text width=4.75cm, inner sep=0.1cm, align=left, font=\FoneBold]
  (rh) at (10.05,-0.28) {Random partitions (\RandVtwoNPartitions)};
\node[txt, anchor=north west, text width=4.95cm] (rb) at ([yshift=-0.12cm]rh.south west)
  {\textbullet\ groups of resonant classes drawn at random, plus one class of QCD jets\\[2pt]
   \textbullet\ group shares of training jets within \RandVtwoShareTol{} of the \Kcoarse-class groups\\[2pt]
   \textbullet\ best of \RandVtwoTrials{} draws: the fewest merges shared with the \Kcoarse-class level\\[2pt]
   \textbullet\ left: the first partition; boxes with the same letter are one class};
% how the random partitions with controlled flavor boundaries are built
\node[hdr=cflav, anchor=north west, text width=4.75cm, inner sep=0.1cm, align=left, font=\FoneBold]
  (fh) at ([yshift=-0.22cm]rb.south west) {Random partitions with controlled flavor boundaries};
\node[txt, anchor=north west, text width=4.95cm] (fn) at ([yshift=-0.12cm]fh.south west)
  {whole families of decays that differ only in quark flavor; two of them split one family,
   the \RandFlavSplitOrbitClasses{} decays \RandFlavSplitOrbit:};
\begin{scope}[shift={(fn.south west)}]
  % brackets over the decays with and without a b quark (squares at a 0.2 cm pitch)
  \node[txt, anchor=south] at (1.075,-0.45) {with $b$};
  \node[txt, anchor=south] at (3.425,-0.45) {without $b$};
  \draw[line width=0.5pt] (0,-0.54) -- (0,-0.47) -- (2.15,-0.47) -- (2.15,-0.54);
  \draw[line width=0.5pt] (2.35,-0.54) -- (2.35,-0.47) -- (4.50,-0.47) -- (4.50,-0.54);
  % no flavor boundary: all 22 in one class
  \node[txt, anchor=north west] at (0,-0.62) {no flavor boundary};
  \foreach \i in {0,...,10} {
    \draw[sq, fill=white] (0.2*\i,-1.10) rectangle ++(0.15,0.15);
    \draw[sq, fill=white] (2.35+0.2*\i,-1.10) rectangle ++(0.15,0.15);
  }
  % b-quark boundary: the decays with a b quark moved to another class
  \node[txt, anchor=north west] at (0,-1.22) {$b$-quark boundary};
  \foreach \i in {0,...,10} {
    \draw[sq, fill=cflav!70!black] (0.2*\i,-1.70) rectangle ++(0.15,0.15);
    \draw[sq, fill=white] (2.35+0.2*\i,-1.70) rectangle ++(0.15,0.15);
  }
  % random flavor boundary: the same number of decays, drawn at random
  \node[txt, anchor=north west] at (0,-1.82) {random flavor boundary};
  \foreach \i in {0,1,6,9,10}
    \draw[sq, fill=white] (0.2*\i,-2.30) rectangle ++(0.15,0.15);
  \foreach \i in {2,3,4,5,7,8}
    \draw[sq, fill=cflav!70!black] (0.2*\i,-2.30) rectangle ++(0.15,0.15);
  \foreach \i in {0,4,5,7,8,10}
    \draw[sq, fill=white] (2.35+0.2*\i,-2.30) rectangle ++(0.15,0.15);
  \foreach \i in {1,2,3,6,9}
    \draw[sq, fill=cflav!70!black] (2.35+0.2*\i,-2.30) rectangle ++(0.15,0.15);
  % key
  \draw[sq, fill=cflav!70!black] (0,-2.62) rectangle ++(0.15,0.15);
  \node[txt, anchor=west] (fkey) at (0.25,-2.545) {moved to another class};
\end{scope}
% ================================================================ (c) the pipeline
% (c) starts below the lower of the two columns above, whatever their text heights
\pgfpointanchor{m3}{south}\pgfgetlastxy{\FoneXa}{\FoneYa}%
\pgfpointanchor{fkey}{south}\pgfgetlastxy{\FoneXb}{\FoneYb}%
\pgfmathsetlengthmacro{\FoneYc}{min(\FoneYa,\FoneYb)-0.45cm}%
\coordinate (ctop) at (0pt,\FoneYc);
\begin{scope}[shift={(ctop)}]
  \node[ttl, anchor=north west] at (0,0) {(c) Pretraining and evaluation};
  % what every model of one seed shares
  \node[stage, text width=2.15cm] (seed) at (1.175,-1.96)
    {\textbf{Same for every model of one seed:} the training jets, their order and the
     initial weights, except the output layer};
  % one pretraining per label set: the configuration colours of every later figure
  \node[draw=black!45, fill=black!4, line width=0.6pt, rounded corners=3pt, minimum width=5.15cm,
        minimum height=3.08cm, anchor=north west, inner sep=0pt] (pre) at (2.81,-0.42) {};
  \node[anchor=north west, font=\FoneBold, inner sep=0pt] at (2.96,-0.52) {Pretraining on JetClass-II};
  \node[chip=c188, minimum width=1.13cm, minimum height=0.42cm] at (2.96,-0.95) {\Kfull};
  \node[chip=c162, minimum width=1.13cm, minimum height=0.42cm] at (4.20,-0.95) {\Kres};
  \node[chip=c43,  minimum width=1.13cm, minimum height=0.42cm] at (5.44,-0.95) {\Kmid};
  \node[chip=c17,  minimum width=1.13cm, minimum height=0.42cm] at (6.68,-0.95) {\Kcoarse};
  \node[chip=crand, minimum width=1.70cm, minimum height=0.70cm] at (2.96,-1.47) {random\\partitions};
  \node[chip=cflav, minimum width=3.04cm, minimum height=0.70cm] at (4.77,-1.47) {controlled flavor\\boundaries};
  \node[chip=c162, densely dotted, minimum width=1.13cm, minimum height=0.42cm] at (2.96,-2.32) {\Kres};
  \node[chip=c17,  densely dotted, minimum width=1.13cm, minimum height=0.42cm] at (4.20,-2.32) {\Kcoarse};
  \node[txt, anchor=west] at (5.44,-2.53) {with auxiliary\\mass regression};
  \node[chip=c188, densely dotted, minimum width=1.13cm, minimum height=0.42cm] at (2.96,-2.94) {\Kfull};
  \node[txt, anchor=west] at (4.20,-3.15) {kinematic inputs only};
  % evaluation
  \node[stage, text width=2.4cm] (fz) at (9.68,-0.95) {frozen features:\\linear probe,\\MLP probe};
  \node[stage, text width=2.4cm] (ft) at (9.68,-2.05) {fine-tuning\\of all weights};
  \node[stage, text width=2.4cm, draw=black, fill=white, line width=1.1pt] (sc) at (9.68,-3.15) {trained from\\scratch};
  \node[stage, text width=3.4cm, minimum height=3.2cm] (task) at (13.20,-2.00)
    {\textbf{Downstream tasks}\\JetClass-II: all classes\\JetClass-II: class pairs\\\quad(both class-balanced)\\JetClass\\
     top tagging\\quark/gluon, Pythia\\quark/gluon, Herwig\\jet mass\FoneXbx};
  \draw[flow] (seed.east) -- (seed.east -| pre.west);
  \draw[flow] (fz.west -| pre.east) -- (fz.west);
  \draw[flow] (ft.west -| pre.east) -- (ft.west);
  \draw[flow] (fz.east) -- (fz.east -| task.west);
  \draw[flow] (ft.east) -- (ft.east -| task.west);
  \draw[flow] (sc.east) -- (sc.east -| task.west);
\end{scope}
\end{tikzpicture}
\caption{Design of the study.
(a)~The class hierarchy on the decays of a heavy resonance~$X$ and on the jets from quantum
chromodynamics (QCD); $q$ is a $u$ or $d$ quark, $\ell$ an electron or a muon, and $\tau_h$,
$\tau_e$, $\tau_\mu$ a $\tau$ lepton decaying to hadrons, an electron or a muon. Each box is one
class of the label set above its column, except the QCD box of the \Kfull-class column (Sophon's
label set~\cite{sophon}), which stands for \LitJetClassTwoNQcdClasses{} classes. Each level merges
classes of the level to its left; the lines below name what each merge removes.
(b)~The controls: the first random partition, where boxes with the same letter are one class, and
the random partitions with controlled flavor boundaries on the decays \RandFlavSplitOrbit{} (one
square per decay; filled squares are moved to a second class).
(c)~Models of one seed share the training jets, their order and the initial weights, except the
output layer, the mass regression target, and the first layer and initial weights of the model with
kinematic inputs only (the particle four-momenta). The hierarchy levels have \MlstNRunsWord{} seeds
each, the models with the auxiliary mass regression \MlstMassNRunsWord{}, the random partitions with
controlled flavor boundaries \RandFlavRuns{} and each random partition \RandVtwoRuns{}. Frozen
features, the output of the backbone (the network below the classification head) with its weights held
fixed, feed linear probes (logistic regressions) and multilayer perceptron (MLP) probes; fine-tuning
trains all weights. Colors identify the label sets and controls in every figure; the plots give the three partitions with
controlled flavor boundaries three shades; dotted borders mark the two
variants, which the other figures draw as dotted lines in the color of their label set.}
\label{fig:design}
\end{figure}
 from main.tex; this file holds one figure environment and nothing else.
%
% WHAT main.tex MUST LOAD (this file loads nothing itself):
%   \usepackage{tikz}
%   \usetikzlibrary{arrows.meta}          % the only TikZ library used (arrow tips)
%   \definecolor{c188}{HTML}{1F77B4}      % the configuration colours of experiments/FIGS/
%   \definecolor{c162}{HTML}{FF7F0E}      % mlst_style.py (matplotlib's tab10 in order, as in
%   \definecolor{c43}{HTML}{2CA02C}       % Sophon's figures); the variants are drawn in their
%   \definecolor{c17}{HTML}{D62728}       % label set's colour with a dotted border, as the
%   \definecolor{crand}{HTML}{17BECF}     % plotted figures draw variants dotted
%   \definecolor{cflav}{HTML}{9467BD}
%   % GENERATED by experiments/FIGS/make_tables.py -- do not edit.
% Every line carries: source file :: path inside that file :: first 16 hex of its
% sha256. paper/journal/provenance.json carries the full hash beside the value.
% Regenerate with:  python3 experiments/FIGS/make_tables.py
% Check in CI with: python3 experiments/FIGS/make_tables.py --check
%
% Macro names spell digits as words (162 -> Onesixtwo) because \newcommand takes
% letters only. A number whose source file does not exist yet is NOT emitted --
% see the missing list below, and leave a \pending{} in the text for it.
%
% MISSING INPUTS -- no macro was emitted for these:
%   community benchmarks, bench v3 readout -- experiments/FIGS/data/bench_v3_metrics*/bench_metrics_last.json
%   v2 pretraining rerun: data order and dropout fixed by seed and epoch, identical initial trunk weights, each epoch reading one of five disjoint random parts of the rows of nearly every file, fixed validation sample, checkpoint at the first maximum within the last ten epochs -- experiments/FIGS/data/v2/pretraining/
%   v2 frozen probes on the rerun models, at the first maximum within the last ten epochs and on the model whose weights are the average of the last ten epochs, BatchNorm statistics recomputed, with untrained-network and pooled-embedding reference rows -- experiments/FIGS/data/v2/probe_ladder/
%   v2 paired ratios with run and test-sample errors on the rerun models -- experiments/FIGS/data/v2/paired_errors/
%   v2 the 64- and 30-class vocabularies -- experiments/FIGS/data/v2/levels_64_30/
%   v2 the leave-one-family-out vocabularies -- experiments/FIGS/data/v2/leave_one_family_out/
%   v2 label-tree recovery on the rerun models -- experiments/FIGS/data/v2/label_recovery/
%   v2 the rerun's random partitions and flavour pair -- experiments/FIGS/data/v2/random_partitions/
%   v2 JetClass-II fine-tuning on files held out from pretraining -- experiments/FIGS/data/v2/finetune/
%   v2 fine-tuning from random initialisation and from the self-supervised model -- experiments/FIGS/data/v2/finetune_references/
%   v2 top tagging and quark/gluon tagging (Pythia and Herwig) fine-tuned from every rerun model -- experiments/FIGS/data/v2/benchmarks/
%   v2 the second and third self-supervised pretraining runs -- experiments/FIGS/data/v2/self_supervised/
%   v2 the mass-output models with the loss weight matched in loss share -- experiments/FIGS/data/v2/mass_lambda_matched/
%   v2 anomaly detection on the rerun models, per-run values -- experiments/FIGS/data/v2/anomaly/
%   v2 the open-data fits on the rerun models -- experiments/FIGS/data/v2/real_data/
%   MassLossShareVtwoOnesixtwoMass -- experiments/FIGS/data/v2/mass_lambda_matched/loss_share.json :: shares.162+mass
%   MassLossShareVtwoOnesevenMass -- experiments/FIGS/data/v2/mass_lambda_matched/loss_share.json :: shares.17+mass
%   MassLossShareVtwoOnesevenMassMatched -- experiments/FIGS/data/v2/mass_lambda_matched/loss_share.json :: shares.17+mass_matched
%   DesignVtwoSecondGpu, DesignVtwoSecondGpuRuns -- the grid's job specs put every run index on one product
%   per-arm parameters and MACs -- flops_by_arm.json (see docs/RECORD.md 2.1)
%   fine-tuning read-outs -- experiments/FIGS/data/v2/ft_mlst/: jc2, jc2_qcd
%   top-tagging recipe check -- ft_mlst/topcheck_bench_metrics.json
%   X->bb and X->bs fits -- experiments/FIGS/data/v2/xbx/
%   X->bs from scratch -- experiments/FIGS/data/v2/xbx_scratch/
%   JetClass accuracy bootstrap -- experiments/FIGS/data/v2/budget/jc1_acc_boot_N10000.json
%   MlstFtJcOneErrCoarseOverFullBootEFour -- waits on the JetClass bootstrap at 1e4 (after the equal-stopping read-out)
%   MlstLpBbQgZeroAccDropFull -- waits on pair probes with the quark/gluon inputs zeroed (open item O10, not run)
%   MlstLpBbQgZeroAccDropCoarse -- waits on pair probes with the quark/gluon inputs zeroed (open item O10, not run)
%   XbsRthirtyFullOverScratch -- waits on X->bs fits and from-scratch runs
%   XbsRthirtyFullOverScratchPr -- waits on X->bs fits and from-scratch runs
%   XbsRthirtyCoarseOverFull -- waits on X->bs fits
%   XbsRthirtyCoarseOverFullPr -- waits on X->bs fits
%   XbsRthirtyNodeOverTransferFull -- waits on X->bs fits
%   XbsRthirtyNodeOverTransferFullPr -- waits on X->bs fits
%   XbbRthirtyDirectOverTransferFull -- waits on X->bb fits
%   XbbRthirtyDirectOverTransferFullPr -- waits on X->bb fits
%   XbbRthirtyTransferCoarseOverFull -- waits on X->bb fits
%   XbbRthirtyTransferCoarseOverFullPr -- waits on X->bb fits
%
% SKIPPED TABLES:
%   T4 compute/parameters: no machine-readable source in the repository; docs/RECORD.md 2.1 names /data/results/mtx/flops/flops_by_arm.json, which is not committed

\newcommand{\PendingPretrain}{\pending{pending v2: pretraining rerun: data order and dropout fixed by seed and epoch, identical initial trunk weights, each epoch reading one of five disjoint random parts of the rows of nearly every file, fixed validation sample, checkpoint at the first maximum within the last ten epochs}}  % experiments/FIGS/make_tables.py :: V2_PENDING['Pretrain'] :: sha256 10e6ec08356a4a2b :: a marker, not a number: red in the draft until the input exists
\newcommand{\PendingProbes}{\pending{pending v2: frozen probes on the rerun models, at the first maximum within the last ten epochs and on the model whose weights are the average of the last ten epochs, BatchNorm statistics recomputed, with untrained-network and pooled-embedding reference rows}}  % experiments/FIGS/make_tables.py :: V2_PENDING['Probes'] :: sha256 10e6ec08356a4a2b :: a marker, not a number: red in the draft until the input exists
\newcommand{\PendingPaired}{\pending{pending v2: paired ratios with run and test-sample errors on the rerun models}}  % experiments/FIGS/make_tables.py :: V2_PENDING['Paired'] :: sha256 10e6ec08356a4a2b :: a marker, not a number: red in the draft until the input exists
\newcommand{\PendingLevels}{\pending{pending v2: the 64- and 30-class vocabularies}}  % experiments/FIGS/make_tables.py :: V2_PENDING['Levels'] :: sha256 10e6ec08356a4a2b :: a marker, not a number: red in the draft until the input exists
\newcommand{\PendingLofo}{\pending{pending v2: the leave-one-family-out vocabularies}}  % experiments/FIGS/make_tables.py :: V2_PENDING['Lofo'] :: sha256 10e6ec08356a4a2b :: a marker, not a number: red in the draft until the input exists
\newcommand{\PendingRecovery}{\pending{pending v2: label-tree recovery on the rerun models}}  % experiments/FIGS/make_tables.py :: V2_PENDING['Recovery'] :: sha256 10e6ec08356a4a2b :: a marker, not a number: red in the draft until the input exists
\newcommand{\PendingRandom}{\pending{pending v2: the rerun's random partitions and flavour pair}}  % experiments/FIGS/make_tables.py :: V2_PENDING['Random'] :: sha256 10e6ec08356a4a2b :: a marker, not a number: red in the draft until the input exists
\newcommand{\PendingFtHeldout}{\pending{pending v2: JetClass-II fine-tuning on files held out from pretraining}}  % experiments/FIGS/make_tables.py :: V2_PENDING['FtHeldout'] :: sha256 10e6ec08356a4a2b :: a marker, not a number: red in the draft until the input exists
\newcommand{\PendingFtRefs}{\pending{pending v2: fine-tuning from random initialisation and from the self-supervised model}}  % experiments/FIGS/make_tables.py :: V2_PENDING['FtRefs'] :: sha256 10e6ec08356a4a2b :: a marker, not a number: red in the draft until the input exists
\newcommand{\PendingBench}{\pending{pending v2: top tagging and quark/gluon tagging (Pythia and Herwig) fine-tuned from every rerun model}}  % experiments/FIGS/make_tables.py :: V2_PENDING['Bench'] :: sha256 10e6ec08356a4a2b :: a marker, not a number: red in the draft until the input exists
\newcommand{\PendingSsl}{\pending{pending v2: the second and third self-supervised pretraining runs}}  % experiments/FIGS/make_tables.py :: V2_PENDING['Ssl'] :: sha256 10e6ec08356a4a2b :: a marker, not a number: red in the draft until the input exists
\newcommand{\PendingMassLambda}{\pending{pending v2: the mass-output models with the loss weight matched in loss share}}  % experiments/FIGS/make_tables.py :: V2_PENDING['MassLambda'] :: sha256 10e6ec08356a4a2b :: a marker, not a number: red in the draft until the input exists
\newcommand{\PendingAnomaly}{\pending{pending v2: anomaly detection on the rerun models, per-run values}}  % experiments/FIGS/make_tables.py :: V2_PENDING['Anomaly'] :: sha256 10e6ec08356a4a2b :: a marker, not a number: red in the draft until the input exists
\newcommand{\PendingRealData}{\pending{pending v2: the open-data fits on the rerun models}}  % experiments/FIGS/make_tables.py :: V2_PENDING['RealData'] :: sha256 10e6ec08356a4a2b :: a marker, not a number: red in the draft until the input exists
\newcommand{\PendingConclusion}{\pending{pending v2: written once every input above is in}}  % experiments/FIGS/make_tables.py :: V2_PENDING['Conclusion'] :: sha256 10e6ec08356a4a2b :: a marker, not a number: red in the draft until the input exists
\newcommand{\AojChainDataset}{178}  % experiments/FIGS/data/aoj_checks_v2/selection_chain_v2.json :: n_dataset :: sha256 729d83085aba25f8
\newcommand{\AojChainStaged}{12{,}589{,}467}  % experiments/FIGS/data/aoj_checks_v2/selection_chain_v2.json :: n_staged :: sha256 729d83085aba25f8
\newcommand{\AojChainRhoWindow}{9{,}441{,}209}  % experiments/FIGS/data/aoj_checks_v2/selection_chain_v2.json :: n_rho_window :: sha256 729d83085aba25f8
\newcommand{\AojNJetsFit}{4{,}194{,}846}  % experiments/FIGS/data/aoj_checks_v2/selection_chain_v2.json :: n_fit :: sha256 729d83085aba25f8
\newcommand{\AojChainPtBins}{8}  % experiments/FIGS/data/aoj_checks_v2/selection_chain_v2.json :: n_pt_bins :: sha256 729d83085aba25f8
\newcommand{\AojChainEtaMax}{2.40}  % experiments/FIGS/data/aoj_checks_v2/selection_chain_v2.json :: abs_eta_max :: sha256 729d83085aba25f8
\newcommand{\MassLossShareVtwoOnesixtwoMass}{\pending{pending: realised mass-loss share in the rerun, 162 classes}}  % experiments/FIGS/make_tables.py :: SLOTS[MassLossShareVtwoOnesixtwoMass] :: sha256 10e6ec08356a4a2b :: a marker, not a number: experiments/FIGS/data/v2/mass_lambda_matched/loss_share.json :: shares.162+mass does not exist yet
\newcommand{\MassLossShareVtwoOnesevenMass}{\pending{pending: realised mass-loss share in the rerun, 17 classes}}  % experiments/FIGS/make_tables.py :: SLOTS[MassLossShareVtwoOnesevenMass] :: sha256 10e6ec08356a4a2b :: a marker, not a number: experiments/FIGS/data/v2/mass_lambda_matched/loss_share.json :: shares.17+mass does not exist yet
\newcommand{\MassLossShareVtwoOnesevenMassMatched}{\pending{pending: realised mass-loss share in the rerun, 17 classes, matched weight}}  % experiments/FIGS/make_tables.py :: SLOTS[MassLossShareVtwoOnesevenMassMatched] :: sha256 10e6ec08356a4a2b :: a marker, not a number: experiments/FIGS/data/v2/mass_lambda_matched/loss_share.json :: shares.17+mass_matched does not exist yet
\newcommand{\FtCoverageEThreeNClasses}{160}  % experiments/FT/data/jc2_v2_manifest.json :: class_coverage['train_N1000_s1.parquet'].n_classes_present :: sha256 4bd68655919b14e9
\newcommand{\FtCoverageEFourNClasses}{183}  % experiments/FT/data/jc2_v2_manifest.json :: class_coverage['train_N10000_s1.parquet'].n_classes_present :: sha256 4bd68655919b14e9
\newcommand{\FtCoverageEThreeMedian}{4}  % experiments/FT/data/jc2_v2_manifest.json :: class_coverage['train_N1000_s1.parquet'].median_per_present_class :: sha256 4bd68655919b14e9
\newcommand{\FtCoverageEThreeOneJet}{23}  % experiments/FT/data/jc2_v2_manifest.json :: class_coverage['train_N1000_s1.parquet'].n_classes_with_one_jet :: sha256 4bd68655919b14e9
\newcommand{\FtCoverageVoneEThreeNClasses}{161}  % experiments/FT/data/jc2_v1_coverage.json :: class_coverage['train_N1000_s1.parquet'].n_classes_present :: sha256 82652f57f708dc23
\newcommand{\FtCoverageVoneEThreeMedian}{5}  % experiments/FT/data/jc2_v1_coverage.json :: class_coverage['train_N1000_s1.parquet'].median_per_present_class :: sha256 82652f57f708dc23
\newcommand{\FtCoverageVoneEThreeOneJet}{25}  % experiments/FT/data/jc2_v1_coverage.json :: class_coverage['train_N1000_s1.parquet'].n_classes_with_one_jet :: sha256 82652f57f708dc23
\newcommand{\ProbeNJets}{2{,}000{,}000}  % experiments/FIGS/data/probe_ladder_v2_mlp2/analysis/seed_level_results.json :: provenance.n_jets_total :: sha256 53e1ec81343ac2c1
\newcommand{\ProbeNSeeds}{5}  % experiments/FIGS/data/probe_ladder_v2_mlp2/analysis/seed_level_results.json :: seeds_used (length) :: sha256 53e1ec81343ac2c1
\newcommand{\ProbeNSeedsWord}{five}  % experiments/FIGS/data/probe_ladder_v2_mlp2/analysis/seed_level_results.json :: seeds_used (length) :: sha256 53e1ec81343ac2c1 :: pretraining runs per vocabulary, as a word for the prose
\newcommand{\ProbeRowAlign}{\texttt{a9449fd420c473f6}}  % experiments/FIGS/data/probe_ladder_v2_mlp2/analysis/seed_level_results.json :: provenance.row_alignment_sha256 (first 16) :: sha256 53e1ec81343ac2c1
\newcommand{\ProbeEpsS}{90\%}  % experiments/FIGS/data/probe_ladder_v2_mlp2/analysis/seed_level_results.json :: levels.bvc_4prong.linear[0].headline_eps_s :: sha256 53e1ec81343ac2c1
\newcommand{\ProbeAucBvcFourprongLinearOneeighteight}{0.99320}  % experiments/FIGS/data/probe_ladder_v2_mlp2/analysis/seed_level_results.json :: levels.bvc_4prong.linear[0].mean_auc :: sha256 53e1ec81343ac2c1 :: 5 seeds, 0 saturated at AUC=1
\newcommand{\ProbeAucSdBvcFourprongLinearOneeighteight}{0.00027}  % experiments/FIGS/data/probe_ladder_v2_mlp2/analysis/seed_level_results.json :: table[bvc_4prong,linear,188].auc (SD over seeds, ddof=1) :: sha256 53e1ec81343ac2c1
\newcommand{\ProbeLogOneMinusAucBvcFourprongLinearOneeighteight}{\ensuremath{-}4.9915}  % experiments/FIGS/data/probe_ladder_v2_mlp2/analysis/seed_level_results.json :: levels.bvc_4prong.linear[0].mean :: sha256 53e1ec81343ac2c1 :: natural log of 1-AUC, lower is better
\newcommand{\ProbeLogOneMinusAucSdBvcFourprongLinearOneeighteight}{0.0389}  % experiments/FIGS/data/probe_ladder_v2_mlp2/analysis/seed_level_results.json :: levels.bvc_4prong.linear[0].seed_sd :: sha256 53e1ec81343ac2c1
\newcommand{\ProbeOmaBvcFourprongLinearOneeighteight}{\ensuremath{(6.80 \pm 0.27)\times10^{-3}}}  % experiments/FIGS/data/probe_ladder_v2_mlp2/analysis/seed_level_results.json :: table[bvc_4prong,linear,188].auc :: sha256 53e1ec81343ac2c1 :: 1-AUC, mean +- SD over 5 seeds
\newcommand{\ProbeBkgLeftBvcFourprongLinearOneeighteight}{14.0}  % experiments/FIGS/data/probe_ladder_v2_mlp2/analysis/seed_level_results.json :: levels.bvc_4prong.linear[0].rejection_points['0.90'].mean_n_bkg_pass :: sha256 53e1ec81343ac2c1 :: background jets passing the cut, mean over seeds
\newcommand{\ProbeRejBvcFourprongLinearOneeighteight}{227\,\ensuremath{\pm}\,23}  % experiments/FIGS/data/probe_ladder_v2_mlp2/analysis/seed_level_results.json :: table[bvc_4prong,linear,188].rejection_points['0.90'].rejection :: sha256 53e1ec81343ac2c1 :: at 90% signal efficiency, mean +- SD over seeds; 0 of 5 seeds at the cap
\newcommand{\ProbeAucBvcFourprongLinearOnesixtwo}{0.99311}  % experiments/FIGS/data/probe_ladder_v2_mlp2/analysis/seed_level_results.json :: levels.bvc_4prong.linear[1].mean_auc :: sha256 53e1ec81343ac2c1 :: 5 seeds, 0 saturated at AUC=1
\newcommand{\ProbeAucSdBvcFourprongLinearOnesixtwo}{0.00018}  % experiments/FIGS/data/probe_ladder_v2_mlp2/analysis/seed_level_results.json :: table[bvc_4prong,linear,162].auc (SD over seeds, ddof=1) :: sha256 53e1ec81343ac2c1
\newcommand{\ProbeLogOneMinusAucBvcFourprongLinearOnesixtwo}{\ensuremath{-}4.9777}  % experiments/FIGS/data/probe_ladder_v2_mlp2/analysis/seed_level_results.json :: levels.bvc_4prong.linear[1].mean :: sha256 53e1ec81343ac2c1 :: natural log of 1-AUC, lower is better
\newcommand{\ProbeLogOneMinusAucSdBvcFourprongLinearOnesixtwo}{0.0264}  % experiments/FIGS/data/probe_ladder_v2_mlp2/analysis/seed_level_results.json :: levels.bvc_4prong.linear[1].seed_sd :: sha256 53e1ec81343ac2c1
\newcommand{\ProbeOmaBvcFourprongLinearOnesixtwo}{\ensuremath{(6.89 \pm 0.18)\times10^{-3}}}  % experiments/FIGS/data/probe_ladder_v2_mlp2/analysis/seed_level_results.json :: table[bvc_4prong,linear,162].auc :: sha256 53e1ec81343ac2c1 :: 1-AUC, mean +- SD over 5 seeds
\newcommand{\ProbeBkgLeftBvcFourprongLinearOnesixtwo}{14.4}  % experiments/FIGS/data/probe_ladder_v2_mlp2/analysis/seed_level_results.json :: levels.bvc_4prong.linear[1].rejection_points['0.90'].mean_n_bkg_pass :: sha256 53e1ec81343ac2c1 :: background jets passing the cut, mean over seeds
\newcommand{\ProbeRejBvcFourprongLinearOnesixtwo}{220\,\ensuremath{\pm}\,40}  % experiments/FIGS/data/probe_ladder_v2_mlp2/analysis/seed_level_results.json :: table[bvc_4prong,linear,162].rejection_points['0.90'].rejection :: sha256 53e1ec81343ac2c1 :: at 90% signal efficiency, mean +- SD over seeds; 0 of 5 seeds at the cap
\newcommand{\ProbeAucBvcFourprongLinearFourthree}{0.99250}  % experiments/FIGS/data/probe_ladder_v2_mlp2/analysis/seed_level_results.json :: levels.bvc_4prong.linear[2].mean_auc :: sha256 53e1ec81343ac2c1 :: 5 seeds, 0 saturated at AUC=1
\newcommand{\ProbeAucSdBvcFourprongLinearFourthree}{0.00049}  % experiments/FIGS/data/probe_ladder_v2_mlp2/analysis/seed_level_results.json :: table[bvc_4prong,linear,43].auc (SD over seeds, ddof=1) :: sha256 53e1ec81343ac2c1
\newcommand{\ProbeLogOneMinusAucBvcFourprongLinearFourthree}{\ensuremath{-}4.8946}  % experiments/FIGS/data/probe_ladder_v2_mlp2/analysis/seed_level_results.json :: levels.bvc_4prong.linear[2].mean :: sha256 53e1ec81343ac2c1 :: natural log of 1-AUC, lower is better
\newcommand{\ProbeLogOneMinusAucSdBvcFourprongLinearFourthree}{0.0646}  % experiments/FIGS/data/probe_ladder_v2_mlp2/analysis/seed_level_results.json :: levels.bvc_4prong.linear[2].seed_sd :: sha256 53e1ec81343ac2c1
\newcommand{\ProbeOmaBvcFourprongLinearFourthree}{\ensuremath{(7.5 \pm 0.5)\times10^{-3}}}  % experiments/FIGS/data/probe_ladder_v2_mlp2/analysis/seed_level_results.json :: table[bvc_4prong,linear,43].auc :: sha256 53e1ec81343ac2c1 :: 1-AUC, mean +- SD over 5 seeds
\newcommand{\ProbeBkgLeftBvcFourprongLinearFourthree}{16.6}  % experiments/FIGS/data/probe_ladder_v2_mlp2/analysis/seed_level_results.json :: levels.bvc_4prong.linear[2].rejection_points['0.90'].mean_n_bkg_pass :: sha256 53e1ec81343ac2c1 :: background jets passing the cut, mean over seeds
\newcommand{\ProbeRejBvcFourprongLinearFourthree}{192\,\ensuremath{\pm}\,20}  % experiments/FIGS/data/probe_ladder_v2_mlp2/analysis/seed_level_results.json :: table[bvc_4prong,linear,43].rejection_points['0.90'].rejection :: sha256 53e1ec81343ac2c1 :: at 90% signal efficiency, mean +- SD over seeds; 0 of 5 seeds at the cap
\newcommand{\ProbeAucBvcFourprongLinearOneseven}{0.96936}  % experiments/FIGS/data/probe_ladder_v2_mlp2/analysis/seed_level_results.json :: levels.bvc_4prong.linear[3].mean_auc :: sha256 53e1ec81343ac2c1 :: 5 seeds, 0 saturated at AUC=1
\newcommand{\ProbeAucSdBvcFourprongLinearOneseven}{0.00514}  % experiments/FIGS/data/probe_ladder_v2_mlp2/analysis/seed_level_results.json :: table[bvc_4prong,linear,17].auc (SD over seeds, ddof=1) :: sha256 53e1ec81343ac2c1
\newcommand{\ProbeLogOneMinusAucBvcFourprongLinearOneseven}{\ensuremath{-}3.4967}  % experiments/FIGS/data/probe_ladder_v2_mlp2/analysis/seed_level_results.json :: levels.bvc_4prong.linear[3].mean :: sha256 53e1ec81343ac2c1 :: natural log of 1-AUC, lower is better
\newcommand{\ProbeLogOneMinusAucSdBvcFourprongLinearOneseven}{0.1667}  % experiments/FIGS/data/probe_ladder_v2_mlp2/analysis/seed_level_results.json :: levels.bvc_4prong.linear[3].seed_sd :: sha256 53e1ec81343ac2c1
\newcommand{\ProbeOmaBvcFourprongLinearOneseven}{\ensuremath{(3.1 \pm 0.5)\times10^{-2}}}  % experiments/FIGS/data/probe_ladder_v2_mlp2/analysis/seed_level_results.json :: table[bvc_4prong,linear,17].auc :: sha256 53e1ec81343ac2c1 :: 1-AUC, mean +- SD over 5 seeds
\newcommand{\ProbeBkgLeftBvcFourprongLinearOneseven}{137.4}  % experiments/FIGS/data/probe_ladder_v2_mlp2/analysis/seed_level_results.json :: levels.bvc_4prong.linear[3].rejection_points['0.90'].mean_n_bkg_pass :: sha256 53e1ec81343ac2c1 :: background jets passing the cut, mean over seeds
\newcommand{\ProbeRejBvcFourprongLinearOneseven}{24\,\ensuremath{\pm}\,6}  % experiments/FIGS/data/probe_ladder_v2_mlp2/analysis/seed_level_results.json :: table[bvc_4prong,linear,17].rejection_points['0.90'].rejection :: sha256 53e1ec81343ac2c1 :: at 90% signal efficiency, mean +- SD over seeds; 0 of 5 seeds at the cap
\newcommand{\ProbeOmaRatioBvcFourprongLinearOnesixtwoOverOneeighteight}{1.01}  % experiments/FIGS/data/probe_ladder_v2_mlp2/analysis/seed_level_results.json :: table[bvc_4prong,linear,162/188].auc :: sha256 53e1ec81343ac2c1 :: ratio of the seed means of 1-AUC, coarser over finer
\newcommand{\ProbeOmaRatioBvcFourprongLinearFourthreeOverOneeighteight}{1.10}  % experiments/FIGS/data/probe_ladder_v2_mlp2/analysis/seed_level_results.json :: table[bvc_4prong,linear,43/188].auc :: sha256 53e1ec81343ac2c1 :: ratio of the seed means of 1-AUC, coarser over finer
\newcommand{\ProbeOmaRatioBvcFourprongLinearOnesevenOverOneeighteight}{4.5}  % experiments/FIGS/data/probe_ladder_v2_mlp2/analysis/seed_level_results.json :: table[bvc_4prong,linear,17/188].auc :: sha256 53e1ec81343ac2c1 :: ratio of the seed means of 1-AUC, coarser over finer
\newcommand{\ProbeOmaRatioBvcFourprongLinearFourthreeOverOnesixtwo}{1.09}  % experiments/FIGS/data/probe_ladder_v2_mlp2/analysis/seed_level_results.json :: table[bvc_4prong,linear,43/162].auc :: sha256 53e1ec81343ac2c1 :: ratio of the seed means of 1-AUC, coarser over finer
\newcommand{\ProbeOmaRatioBvcFourprongLinearOnesevenOverOnesixtwo}{4.4}  % experiments/FIGS/data/probe_ladder_v2_mlp2/analysis/seed_level_results.json :: table[bvc_4prong,linear,17/162].auc :: sha256 53e1ec81343ac2c1 :: ratio of the seed means of 1-AUC, coarser over finer
\newcommand{\ProbeOmaRatioBvcFourprongLinearOnesevenOverFourthree}{4.1}  % experiments/FIGS/data/probe_ladder_v2_mlp2/analysis/seed_level_results.json :: table[bvc_4prong,linear,17/43].auc :: sha256 53e1ec81343ac2c1 :: ratio of the seed means of 1-AUC, coarser over finer
\newcommand{\ProbeAucBvcFourprongMlpOneeighteight}{0.99316}  % experiments/FIGS/data/probe_ladder_v2_mlp2/analysis/seed_level_results.json :: levels.bvc_4prong.mlp[0].mean_auc :: sha256 53e1ec81343ac2c1 :: 5 seeds, 0 saturated at AUC=1
\newcommand{\ProbeAucSdBvcFourprongMlpOneeighteight}{0.00020}  % experiments/FIGS/data/probe_ladder_v2_mlp2/analysis/seed_level_results.json :: table[bvc_4prong,mlp,188].auc (SD over seeds, ddof=1) :: sha256 53e1ec81343ac2c1
\newcommand{\ProbeLogOneMinusAucBvcFourprongMlpOneeighteight}{\ensuremath{-}4.9857}  % experiments/FIGS/data/probe_ladder_v2_mlp2/analysis/seed_level_results.json :: levels.bvc_4prong.mlp[0].mean :: sha256 53e1ec81343ac2c1 :: natural log of 1-AUC, lower is better
\newcommand{\ProbeLogOneMinusAucSdBvcFourprongMlpOneeighteight}{0.0290}  % experiments/FIGS/data/probe_ladder_v2_mlp2/analysis/seed_level_results.json :: levels.bvc_4prong.mlp[0].seed_sd :: sha256 53e1ec81343ac2c1
\newcommand{\ProbeOmaBvcFourprongMlpOneeighteight}{\ensuremath{(6.84 \pm 0.20)\times10^{-3}}}  % experiments/FIGS/data/probe_ladder_v2_mlp2/analysis/seed_level_results.json :: table[bvc_4prong,mlp,188].auc :: sha256 53e1ec81343ac2c1 :: 1-AUC, mean +- SD over 5 seeds
\newcommand{\ProbeBkgLeftBvcFourprongMlpOneeighteight}{15.2}  % experiments/FIGS/data/probe_ladder_v2_mlp2/analysis/seed_level_results.json :: levels.bvc_4prong.mlp[0].rejection_points['0.90'].mean_n_bkg_pass :: sha256 53e1ec81343ac2c1 :: background jets passing the cut, mean over seeds
\newcommand{\ProbeRejBvcFourprongMlpOneeighteight}{212\,\ensuremath{\pm}\,34}  % experiments/FIGS/data/probe_ladder_v2_mlp2/analysis/seed_level_results.json :: table[bvc_4prong,mlp,188].rejection_points['0.90'].rejection :: sha256 53e1ec81343ac2c1 :: at 90% signal efficiency, mean +- SD over seeds; 0 of 5 seeds at the cap
\newcommand{\ProbeAucBvcFourprongMlpOnesixtwo}{0.99318}  % experiments/FIGS/data/probe_ladder_v2_mlp2/analysis/seed_level_results.json :: levels.bvc_4prong.mlp[1].mean_auc :: sha256 53e1ec81343ac2c1 :: 5 seeds, 0 saturated at AUC=1
\newcommand{\ProbeAucSdBvcFourprongMlpOnesixtwo}{0.00020}  % experiments/FIGS/data/probe_ladder_v2_mlp2/analysis/seed_level_results.json :: table[bvc_4prong,mlp,162].auc (SD over seeds, ddof=1) :: sha256 53e1ec81343ac2c1
\newcommand{\ProbeLogOneMinusAucBvcFourprongMlpOnesixtwo}{\ensuremath{-}4.9886}  % experiments/FIGS/data/probe_ladder_v2_mlp2/analysis/seed_level_results.json :: levels.bvc_4prong.mlp[1].mean :: sha256 53e1ec81343ac2c1 :: natural log of 1-AUC, lower is better
\newcommand{\ProbeLogOneMinusAucSdBvcFourprongMlpOnesixtwo}{0.0298}  % experiments/FIGS/data/probe_ladder_v2_mlp2/analysis/seed_level_results.json :: levels.bvc_4prong.mlp[1].seed_sd :: sha256 53e1ec81343ac2c1
\newcommand{\ProbeOmaBvcFourprongMlpOnesixtwo}{\ensuremath{(6.82 \pm 0.20)\times10^{-3}}}  % experiments/FIGS/data/probe_ladder_v2_mlp2/analysis/seed_level_results.json :: table[bvc_4prong,mlp,162].auc :: sha256 53e1ec81343ac2c1 :: 1-AUC, mean +- SD over 5 seeds
\newcommand{\ProbeBkgLeftBvcFourprongMlpOnesixtwo}{16.6}  % experiments/FIGS/data/probe_ladder_v2_mlp2/analysis/seed_level_results.json :: levels.bvc_4prong.mlp[1].rejection_points['0.90'].mean_n_bkg_pass :: sha256 53e1ec81343ac2c1 :: background jets passing the cut, mean over seeds
\newcommand{\ProbeRejBvcFourprongMlpOnesixtwo}{200\,\ensuremath{\pm}\,40}  % experiments/FIGS/data/probe_ladder_v2_mlp2/analysis/seed_level_results.json :: table[bvc_4prong,mlp,162].rejection_points['0.90'].rejection :: sha256 53e1ec81343ac2c1 :: at 90% signal efficiency, mean +- SD over seeds; 0 of 5 seeds at the cap
\newcommand{\ProbeAucBvcFourprongMlpFourthree}{0.99232}  % experiments/FIGS/data/probe_ladder_v2_mlp2/analysis/seed_level_results.json :: levels.bvc_4prong.mlp[2].mean_auc :: sha256 53e1ec81343ac2c1 :: 5 seeds, 0 saturated at AUC=1
\newcommand{\ProbeAucSdBvcFourprongMlpFourthree}{0.00058}  % experiments/FIGS/data/probe_ladder_v2_mlp2/analysis/seed_level_results.json :: table[bvc_4prong,mlp,43].auc (SD over seeds, ddof=1) :: sha256 53e1ec81343ac2c1
\newcommand{\ProbeLogOneMinusAucBvcFourprongMlpFourthree}{\ensuremath{-}4.8719}  % experiments/FIGS/data/probe_ladder_v2_mlp2/analysis/seed_level_results.json :: levels.bvc_4prong.mlp[2].mean :: sha256 53e1ec81343ac2c1 :: natural log of 1-AUC, lower is better
\newcommand{\ProbeLogOneMinusAucSdBvcFourprongMlpFourthree}{0.0734}  % experiments/FIGS/data/probe_ladder_v2_mlp2/analysis/seed_level_results.json :: levels.bvc_4prong.mlp[2].seed_sd :: sha256 53e1ec81343ac2c1
\newcommand{\ProbeOmaBvcFourprongMlpFourthree}{\ensuremath{(7.7 \pm 0.6)\times10^{-3}}}  % experiments/FIGS/data/probe_ladder_v2_mlp2/analysis/seed_level_results.json :: table[bvc_4prong,mlp,43].auc :: sha256 53e1ec81343ac2c1 :: 1-AUC, mean +- SD over 5 seeds
\newcommand{\ProbeBkgLeftBvcFourprongMlpFourthree}{17.4}  % experiments/FIGS/data/probe_ladder_v2_mlp2/analysis/seed_level_results.json :: levels.bvc_4prong.mlp[2].rejection_points['0.90'].mean_n_bkg_pass :: sha256 53e1ec81343ac2c1 :: background jets passing the cut, mean over seeds
\newcommand{\ProbeRejBvcFourprongMlpFourthree}{184\,\ensuremath{\pm}\,23}  % experiments/FIGS/data/probe_ladder_v2_mlp2/analysis/seed_level_results.json :: table[bvc_4prong,mlp,43].rejection_points['0.90'].rejection :: sha256 53e1ec81343ac2c1 :: at 90% signal efficiency, mean +- SD over seeds; 0 of 5 seeds at the cap
\newcommand{\ProbeAucBvcFourprongMlpOneseven}{0.97766}  % experiments/FIGS/data/probe_ladder_v2_mlp2/analysis/seed_level_results.json :: levels.bvc_4prong.mlp[3].mean_auc :: sha256 53e1ec81343ac2c1 :: 5 seeds, 0 saturated at AUC=1
\newcommand{\ProbeAucSdBvcFourprongMlpOneseven}{0.00310}  % experiments/FIGS/data/probe_ladder_v2_mlp2/analysis/seed_level_results.json :: table[bvc_4prong,mlp,17].auc (SD over seeds, ddof=1) :: sha256 53e1ec81343ac2c1
\newcommand{\ProbeLogOneMinusAucBvcFourprongMlpOneseven}{\ensuremath{-}3.8087}  % experiments/FIGS/data/probe_ladder_v2_mlp2/analysis/seed_level_results.json :: levels.bvc_4prong.mlp[3].mean :: sha256 53e1ec81343ac2c1 :: natural log of 1-AUC, lower is better
\newcommand{\ProbeLogOneMinusAucSdBvcFourprongMlpOneseven}{0.1366}  % experiments/FIGS/data/probe_ladder_v2_mlp2/analysis/seed_level_results.json :: levels.bvc_4prong.mlp[3].seed_sd :: sha256 53e1ec81343ac2c1
\newcommand{\ProbeOmaBvcFourprongMlpOneseven}{\ensuremath{(2.23 \pm 0.31)\times10^{-2}}}  % experiments/FIGS/data/probe_ladder_v2_mlp2/analysis/seed_level_results.json :: table[bvc_4prong,mlp,17].auc :: sha256 53e1ec81343ac2c1 :: 1-AUC, mean +- SD over 5 seeds
\newcommand{\ProbeBkgLeftBvcFourprongMlpOneseven}{99.0}  % experiments/FIGS/data/probe_ladder_v2_mlp2/analysis/seed_level_results.json :: levels.bvc_4prong.mlp[3].rejection_points['0.90'].mean_n_bkg_pass :: sha256 53e1ec81343ac2c1 :: background jets passing the cut, mean over seeds
\newcommand{\ProbeRejBvcFourprongMlpOneseven}{33\,\ensuremath{\pm}\,5}  % experiments/FIGS/data/probe_ladder_v2_mlp2/analysis/seed_level_results.json :: table[bvc_4prong,mlp,17].rejection_points['0.90'].rejection :: sha256 53e1ec81343ac2c1 :: at 90% signal efficiency, mean +- SD over seeds; 0 of 5 seeds at the cap
\newcommand{\ProbeOmaRatioBvcFourprongMlpOnesixtwoOverOneeighteight}{1.00}  % experiments/FIGS/data/probe_ladder_v2_mlp2/analysis/seed_level_results.json :: table[bvc_4prong,mlp,162/188].auc :: sha256 53e1ec81343ac2c1 :: ratio of the seed means of 1-AUC, coarser over finer
\newcommand{\ProbeOmaRatioBvcFourprongMlpFourthreeOverOneeighteight}{1.12}  % experiments/FIGS/data/probe_ladder_v2_mlp2/analysis/seed_level_results.json :: table[bvc_4prong,mlp,43/188].auc :: sha256 53e1ec81343ac2c1 :: ratio of the seed means of 1-AUC, coarser over finer
\newcommand{\ProbeOmaRatioBvcFourprongMlpOnesevenOverOneeighteight}{3.3}  % experiments/FIGS/data/probe_ladder_v2_mlp2/analysis/seed_level_results.json :: table[bvc_4prong,mlp,17/188].auc :: sha256 53e1ec81343ac2c1 :: ratio of the seed means of 1-AUC, coarser over finer
\newcommand{\ProbeOmaRatioBvcFourprongMlpFourthreeOverOnesixtwo}{1.13}  % experiments/FIGS/data/probe_ladder_v2_mlp2/analysis/seed_level_results.json :: table[bvc_4prong,mlp,43/162].auc :: sha256 53e1ec81343ac2c1 :: ratio of the seed means of 1-AUC, coarser over finer
\newcommand{\ProbeOmaRatioBvcFourprongMlpOnesevenOverOnesixtwo}{3.3}  % experiments/FIGS/data/probe_ladder_v2_mlp2/analysis/seed_level_results.json :: table[bvc_4prong,mlp,17/162].auc :: sha256 53e1ec81343ac2c1 :: ratio of the seed means of 1-AUC, coarser over finer
\newcommand{\ProbeOmaRatioBvcFourprongMlpOnesevenOverFourthree}{2.9}  % experiments/FIGS/data/probe_ladder_v2_mlp2/analysis/seed_level_results.json :: table[bvc_4prong,mlp,17/43].auc :: sha256 53e1ec81343ac2c1 :: ratio of the seed means of 1-AUC, coarser over finer
\newcommand{\ProbeAucBvcQcdLinearOneeighteight}{0.98519}  % experiments/FIGS/data/probe_ladder_v2_mlp2/analysis/seed_level_results.json :: levels.bvc_qcd.linear[0].mean_auc :: sha256 53e1ec81343ac2c1 :: 5 seeds, 0 saturated at AUC=1
\newcommand{\ProbeAucSdBvcQcdLinearOneeighteight}{0.00061}  % experiments/FIGS/data/probe_ladder_v2_mlp2/analysis/seed_level_results.json :: table[bvc_qcd,linear,188].auc (SD over seeds, ddof=1) :: sha256 53e1ec81343ac2c1
\newcommand{\ProbeLogOneMinusAucBvcQcdLinearOneeighteight}{\ensuremath{-}4.2133}  % experiments/FIGS/data/probe_ladder_v2_mlp2/analysis/seed_level_results.json :: levels.bvc_qcd.linear[0].mean :: sha256 53e1ec81343ac2c1 :: natural log of 1-AUC, lower is better
\newcommand{\ProbeLogOneMinusAucSdBvcQcdLinearOneeighteight}{0.0422}  % experiments/FIGS/data/probe_ladder_v2_mlp2/analysis/seed_level_results.json :: levels.bvc_qcd.linear[0].seed_sd :: sha256 53e1ec81343ac2c1
\newcommand{\ProbeOmaBvcQcdLinearOneeighteight}{\ensuremath{(1.48 \pm 0.06)\times10^{-2}}}  % experiments/FIGS/data/probe_ladder_v2_mlp2/analysis/seed_level_results.json :: table[bvc_qcd,linear,188].auc :: sha256 53e1ec81343ac2c1 :: 1-AUC, mean +- SD over 5 seeds
\newcommand{\ProbeBkgLeftBvcQcdLinearOneeighteight}{96.4}  % experiments/FIGS/data/probe_ladder_v2_mlp2/analysis/seed_level_results.json :: levels.bvc_qcd.linear[0].rejection_points['0.90'].mean_n_bkg_pass :: sha256 53e1ec81343ac2c1 :: background jets passing the cut, mean over seeds
\newcommand{\ProbeRejBvcQcdLinearOneeighteight}{56\,\ensuremath{\pm}\,5}  % experiments/FIGS/data/probe_ladder_v2_mlp2/analysis/seed_level_results.json :: table[bvc_qcd,linear,188].rejection_points['0.90'].rejection :: sha256 53e1ec81343ac2c1 :: at 90% signal efficiency, mean +- SD over seeds; 0 of 5 seeds at the cap
\newcommand{\ProbeAucBvcQcdLinearOnesixtwo}{0.98305}  % experiments/FIGS/data/probe_ladder_v2_mlp2/analysis/seed_level_results.json :: levels.bvc_qcd.linear[1].mean_auc :: sha256 53e1ec81343ac2c1 :: 5 seeds, 0 saturated at AUC=1
\newcommand{\ProbeAucSdBvcQcdLinearOnesixtwo}{0.00024}  % experiments/FIGS/data/probe_ladder_v2_mlp2/analysis/seed_level_results.json :: table[bvc_qcd,linear,162].auc (SD over seeds, ddof=1) :: sha256 53e1ec81343ac2c1
\newcommand{\ProbeLogOneMinusAucBvcQcdLinearOnesixtwo}{\ensuremath{-}4.0774}  % experiments/FIGS/data/probe_ladder_v2_mlp2/analysis/seed_level_results.json :: levels.bvc_qcd.linear[1].mean :: sha256 53e1ec81343ac2c1 :: natural log of 1-AUC, lower is better
\newcommand{\ProbeLogOneMinusAucSdBvcQcdLinearOnesixtwo}{0.0142}  % experiments/FIGS/data/probe_ladder_v2_mlp2/analysis/seed_level_results.json :: levels.bvc_qcd.linear[1].seed_sd :: sha256 53e1ec81343ac2c1
\newcommand{\ProbeOmaBvcQcdLinearOnesixtwo}{\ensuremath{(1.695 \pm 0.024)\times10^{-2}}}  % experiments/FIGS/data/probe_ladder_v2_mlp2/analysis/seed_level_results.json :: table[bvc_qcd,linear,162].auc :: sha256 53e1ec81343ac2c1 :: 1-AUC, mean +- SD over 5 seeds
\newcommand{\ProbeBkgLeftBvcQcdLinearOnesixtwo}{121.4}  % experiments/FIGS/data/probe_ladder_v2_mlp2/analysis/seed_level_results.json :: levels.bvc_qcd.linear[1].rejection_points['0.90'].mean_n_bkg_pass :: sha256 53e1ec81343ac2c1 :: background jets passing the cut, mean over seeds
\newcommand{\ProbeRejBvcQcdLinearOnesixtwo}{44.3\,\ensuremath{\pm}\,2.1}  % experiments/FIGS/data/probe_ladder_v2_mlp2/analysis/seed_level_results.json :: table[bvc_qcd,linear,162].rejection_points['0.90'].rejection :: sha256 53e1ec81343ac2c1 :: at 90% signal efficiency, mean +- SD over seeds; 0 of 5 seeds at the cap
\newcommand{\ProbeAucBvcQcdLinearFourthree}{0.98281}  % experiments/FIGS/data/probe_ladder_v2_mlp2/analysis/seed_level_results.json :: levels.bvc_qcd.linear[2].mean_auc :: sha256 53e1ec81343ac2c1 :: 5 seeds, 0 saturated at AUC=1
\newcommand{\ProbeAucSdBvcQcdLinearFourthree}{0.00064}  % experiments/FIGS/data/probe_ladder_v2_mlp2/analysis/seed_level_results.json :: table[bvc_qcd,linear,43].auc (SD over seeds, ddof=1) :: sha256 53e1ec81343ac2c1
\newcommand{\ProbeLogOneMinusAucBvcQcdLinearFourthree}{\ensuremath{-}4.0638}  % experiments/FIGS/data/probe_ladder_v2_mlp2/analysis/seed_level_results.json :: levels.bvc_qcd.linear[2].mean :: sha256 53e1ec81343ac2c1 :: natural log of 1-AUC, lower is better
\newcommand{\ProbeLogOneMinusAucSdBvcQcdLinearFourthree}{0.0370}  % experiments/FIGS/data/probe_ladder_v2_mlp2/analysis/seed_level_results.json :: levels.bvc_qcd.linear[2].seed_sd :: sha256 53e1ec81343ac2c1
\newcommand{\ProbeOmaBvcQcdLinearFourthree}{\ensuremath{(1.72 \pm 0.06)\times10^{-2}}}  % experiments/FIGS/data/probe_ladder_v2_mlp2/analysis/seed_level_results.json :: table[bvc_qcd,linear,43].auc :: sha256 53e1ec81343ac2c1 :: 1-AUC, mean +- SD over 5 seeds
\newcommand{\ProbeBkgLeftBvcQcdLinearFourthree}{126.8}  % experiments/FIGS/data/probe_ladder_v2_mlp2/analysis/seed_level_results.json :: levels.bvc_qcd.linear[2].rejection_points['0.90'].mean_n_bkg_pass :: sha256 53e1ec81343ac2c1 :: background jets passing the cut, mean over seeds
\newcommand{\ProbeRejBvcQcdLinearFourthree}{43\,\ensuremath{\pm}\,4}  % experiments/FIGS/data/probe_ladder_v2_mlp2/analysis/seed_level_results.json :: table[bvc_qcd,linear,43].rejection_points['0.90'].rejection :: sha256 53e1ec81343ac2c1 :: at 90% signal efficiency, mean +- SD over seeds; 0 of 5 seeds at the cap
\newcommand{\ProbeAucBvcQcdLinearOneseven}{0.95504}  % experiments/FIGS/data/probe_ladder_v2_mlp2/analysis/seed_level_results.json :: levels.bvc_qcd.linear[3].mean_auc :: sha256 53e1ec81343ac2c1 :: 5 seeds, 0 saturated at AUC=1
\newcommand{\ProbeAucSdBvcQcdLinearOneseven}{0.00527}  % experiments/FIGS/data/probe_ladder_v2_mlp2/analysis/seed_level_results.json :: table[bvc_qcd,linear,17].auc (SD over seeds, ddof=1) :: sha256 53e1ec81343ac2c1
\newcommand{\ProbeLogOneMinusAucBvcQcdLinearOneseven}{\ensuremath{-}3.1076}  % experiments/FIGS/data/probe_ladder_v2_mlp2/analysis/seed_level_results.json :: levels.bvc_qcd.linear[3].mean :: sha256 53e1ec81343ac2c1 :: natural log of 1-AUC, lower is better
\newcommand{\ProbeLogOneMinusAucSdBvcQcdLinearOneseven}{0.1198}  % experiments/FIGS/data/probe_ladder_v2_mlp2/analysis/seed_level_results.json :: levels.bvc_qcd.linear[3].seed_sd :: sha256 53e1ec81343ac2c1
\newcommand{\ProbeOmaBvcQcdLinearOneseven}{\ensuremath{(4.5 \pm 0.5)\times10^{-2}}}  % experiments/FIGS/data/probe_ladder_v2_mlp2/analysis/seed_level_results.json :: table[bvc_qcd,linear,17].auc :: sha256 53e1ec81343ac2c1 :: 1-AUC, mean +- SD over 5 seeds
\newcommand{\ProbeBkgLeftBvcQcdLinearOneseven}{544.6}  % experiments/FIGS/data/probe_ladder_v2_mlp2/analysis/seed_level_results.json :: levels.bvc_qcd.linear[3].rejection_points['0.90'].mean_n_bkg_pass :: sha256 53e1ec81343ac2c1 :: background jets passing the cut, mean over seeds
\newcommand{\ProbeRejBvcQcdLinearOneseven}{10.1\,\ensuremath{\pm}\,1.9}  % experiments/FIGS/data/probe_ladder_v2_mlp2/analysis/seed_level_results.json :: table[bvc_qcd,linear,17].rejection_points['0.90'].rejection :: sha256 53e1ec81343ac2c1 :: at 90% signal efficiency, mean +- SD over seeds; 0 of 5 seeds at the cap
\newcommand{\ProbeOmaRatioBvcQcdLinearOnesixtwoOverOneeighteight}{1.14}  % experiments/FIGS/data/probe_ladder_v2_mlp2/analysis/seed_level_results.json :: table[bvc_qcd,linear,162/188].auc :: sha256 53e1ec81343ac2c1 :: ratio of the seed means of 1-AUC, coarser over finer
\newcommand{\ProbeOmaRatioBvcQcdLinearFourthreeOverOneeighteight}{1.16}  % experiments/FIGS/data/probe_ladder_v2_mlp2/analysis/seed_level_results.json :: table[bvc_qcd,linear,43/188].auc :: sha256 53e1ec81343ac2c1 :: ratio of the seed means of 1-AUC, coarser over finer
\newcommand{\ProbeOmaRatioBvcQcdLinearOnesevenOverOneeighteight}{3.0}  % experiments/FIGS/data/probe_ladder_v2_mlp2/analysis/seed_level_results.json :: table[bvc_qcd,linear,17/188].auc :: sha256 53e1ec81343ac2c1 :: ratio of the seed means of 1-AUC, coarser over finer
\newcommand{\ProbeOmaRatioBvcQcdLinearFourthreeOverOnesixtwo}{1.01}  % experiments/FIGS/data/probe_ladder_v2_mlp2/analysis/seed_level_results.json :: table[bvc_qcd,linear,43/162].auc :: sha256 53e1ec81343ac2c1 :: ratio of the seed means of 1-AUC, coarser over finer
\newcommand{\ProbeOmaRatioBvcQcdLinearOnesevenOverOnesixtwo}{2.7}  % experiments/FIGS/data/probe_ladder_v2_mlp2/analysis/seed_level_results.json :: table[bvc_qcd,linear,17/162].auc :: sha256 53e1ec81343ac2c1 :: ratio of the seed means of 1-AUC, coarser over finer
\newcommand{\ProbeOmaRatioBvcQcdLinearOnesevenOverFourthree}{2.6}  % experiments/FIGS/data/probe_ladder_v2_mlp2/analysis/seed_level_results.json :: table[bvc_qcd,linear,17/43].auc :: sha256 53e1ec81343ac2c1 :: ratio of the seed means of 1-AUC, coarser over finer
\newcommand{\ProbeAucBvcQcdMlpOneeighteight}{0.98527}  % experiments/FIGS/data/probe_ladder_v2_mlp2/analysis/seed_level_results.json :: levels.bvc_qcd.mlp[0].mean_auc :: sha256 53e1ec81343ac2c1 :: 5 seeds, 0 saturated at AUC=1
\newcommand{\ProbeAucSdBvcQcdMlpOneeighteight}{0.00058}  % experiments/FIGS/data/probe_ladder_v2_mlp2/analysis/seed_level_results.json :: table[bvc_qcd,mlp,188].auc (SD over seeds, ddof=1) :: sha256 53e1ec81343ac2c1
\newcommand{\ProbeLogOneMinusAucBvcQcdMlpOneeighteight}{\ensuremath{-}4.2188}  % experiments/FIGS/data/probe_ladder_v2_mlp2/analysis/seed_level_results.json :: levels.bvc_qcd.mlp[0].mean :: sha256 53e1ec81343ac2c1 :: natural log of 1-AUC, lower is better
\newcommand{\ProbeLogOneMinusAucSdBvcQcdMlpOneeighteight}{0.0400}  % experiments/FIGS/data/probe_ladder_v2_mlp2/analysis/seed_level_results.json :: levels.bvc_qcd.mlp[0].seed_sd :: sha256 53e1ec81343ac2c1
\newcommand{\ProbeOmaBvcQcdMlpOneeighteight}{\ensuremath{(1.47 \pm 0.06)\times10^{-2}}}  % experiments/FIGS/data/probe_ladder_v2_mlp2/analysis/seed_level_results.json :: table[bvc_qcd,mlp,188].auc :: sha256 53e1ec81343ac2c1 :: 1-AUC, mean +- SD over 5 seeds
\newcommand{\ProbeBkgLeftBvcQcdMlpOneeighteight}{94.6}  % experiments/FIGS/data/probe_ladder_v2_mlp2/analysis/seed_level_results.json :: levels.bvc_qcd.mlp[0].rejection_points['0.90'].mean_n_bkg_pass :: sha256 53e1ec81343ac2c1 :: background jets passing the cut, mean over seeds
\newcommand{\ProbeRejBvcQcdMlpOneeighteight}{57\,\ensuremath{\pm}\,6}  % experiments/FIGS/data/probe_ladder_v2_mlp2/analysis/seed_level_results.json :: table[bvc_qcd,mlp,188].rejection_points['0.90'].rejection :: sha256 53e1ec81343ac2c1 :: at 90% signal efficiency, mean +- SD over seeds; 0 of 5 seeds at the cap
\newcommand{\ProbeAucBvcQcdMlpOnesixtwo}{0.98361}  % experiments/FIGS/data/probe_ladder_v2_mlp2/analysis/seed_level_results.json :: levels.bvc_qcd.mlp[1].mean_auc :: sha256 53e1ec81343ac2c1 :: 5 seeds, 0 saturated at AUC=1
\newcommand{\ProbeAucSdBvcQcdMlpOnesixtwo}{0.00039}  % experiments/FIGS/data/probe_ladder_v2_mlp2/analysis/seed_level_results.json :: table[bvc_qcd,mlp,162].auc (SD over seeds, ddof=1) :: sha256 53e1ec81343ac2c1
\newcommand{\ProbeLogOneMinusAucBvcQcdMlpOnesixtwo}{\ensuremath{-}4.1113}  % experiments/FIGS/data/probe_ladder_v2_mlp2/analysis/seed_level_results.json :: levels.bvc_qcd.mlp[1].mean :: sha256 53e1ec81343ac2c1 :: natural log of 1-AUC, lower is better
\newcommand{\ProbeLogOneMinusAucSdBvcQcdMlpOnesixtwo}{0.0241}  % experiments/FIGS/data/probe_ladder_v2_mlp2/analysis/seed_level_results.json :: levels.bvc_qcd.mlp[1].seed_sd :: sha256 53e1ec81343ac2c1
\newcommand{\ProbeOmaBvcQcdMlpOnesixtwo}{\ensuremath{(1.64 \pm 0.04)\times10^{-2}}}  % experiments/FIGS/data/probe_ladder_v2_mlp2/analysis/seed_level_results.json :: table[bvc_qcd,mlp,162].auc :: sha256 53e1ec81343ac2c1 :: 1-AUC, mean +- SD over 5 seeds
\newcommand{\ProbeBkgLeftBvcQcdMlpOnesixtwo}{119.6}  % experiments/FIGS/data/probe_ladder_v2_mlp2/analysis/seed_level_results.json :: levels.bvc_qcd.mlp[1].rejection_points['0.90'].mean_n_bkg_pass :: sha256 53e1ec81343ac2c1 :: background jets passing the cut, mean over seeds
\newcommand{\ProbeRejBvcQcdMlpOnesixtwo}{45.1\,\ensuremath{\pm}\,3.2}  % experiments/FIGS/data/probe_ladder_v2_mlp2/analysis/seed_level_results.json :: table[bvc_qcd,mlp,162].rejection_points['0.90'].rejection :: sha256 53e1ec81343ac2c1 :: at 90% signal efficiency, mean +- SD over seeds; 0 of 5 seeds at the cap
\newcommand{\ProbeAucBvcQcdMlpFourthree}{0.98341}  % experiments/FIGS/data/probe_ladder_v2_mlp2/analysis/seed_level_results.json :: levels.bvc_qcd.mlp[2].mean_auc :: sha256 53e1ec81343ac2c1 :: 5 seeds, 0 saturated at AUC=1
\newcommand{\ProbeAucSdBvcQcdMlpFourthree}{0.00052}  % experiments/FIGS/data/probe_ladder_v2_mlp2/analysis/seed_level_results.json :: table[bvc_qcd,mlp,43].auc (SD over seeds, ddof=1) :: sha256 53e1ec81343ac2c1
\newcommand{\ProbeLogOneMinusAucBvcQcdMlpFourthree}{\ensuremath{-}4.0995}  % experiments/FIGS/data/probe_ladder_v2_mlp2/analysis/seed_level_results.json :: levels.bvc_qcd.mlp[2].mean :: sha256 53e1ec81343ac2c1 :: natural log of 1-AUC, lower is better
\newcommand{\ProbeLogOneMinusAucSdBvcQcdMlpFourthree}{0.0318}  % experiments/FIGS/data/probe_ladder_v2_mlp2/analysis/seed_level_results.json :: levels.bvc_qcd.mlp[2].seed_sd :: sha256 53e1ec81343ac2c1
\newcommand{\ProbeOmaBvcQcdMlpFourthree}{\ensuremath{(1.66 \pm 0.05)\times10^{-2}}}  % experiments/FIGS/data/probe_ladder_v2_mlp2/analysis/seed_level_results.json :: table[bvc_qcd,mlp,43].auc :: sha256 53e1ec81343ac2c1 :: 1-AUC, mean +- SD over 5 seeds
\newcommand{\ProbeBkgLeftBvcQcdMlpFourthree}{119.2}  % experiments/FIGS/data/probe_ladder_v2_mlp2/analysis/seed_level_results.json :: levels.bvc_qcd.mlp[2].rejection_points['0.90'].mean_n_bkg_pass :: sha256 53e1ec81343ac2c1 :: background jets passing the cut, mean over seeds
\newcommand{\ProbeRejBvcQcdMlpFourthree}{45.1\,\ensuremath{\pm}\,1.3}  % experiments/FIGS/data/probe_ladder_v2_mlp2/analysis/seed_level_results.json :: table[bvc_qcd,mlp,43].rejection_points['0.90'].rejection :: sha256 53e1ec81343ac2c1 :: at 90% signal efficiency, mean +- SD over seeds; 0 of 5 seeds at the cap
\newcommand{\ProbeAucBvcQcdMlpOneseven}{0.96226}  % experiments/FIGS/data/probe_ladder_v2_mlp2/analysis/seed_level_results.json :: levels.bvc_qcd.mlp[3].mean_auc :: sha256 53e1ec81343ac2c1 :: 5 seeds, 0 saturated at AUC=1
\newcommand{\ProbeAucSdBvcQcdMlpOneseven}{0.00373}  % experiments/FIGS/data/probe_ladder_v2_mlp2/analysis/seed_level_results.json :: table[bvc_qcd,mlp,17].auc (SD over seeds, ddof=1) :: sha256 53e1ec81343ac2c1
\newcommand{\ProbeLogOneMinusAucBvcQcdMlpOneseven}{\ensuremath{-}3.2811}  % experiments/FIGS/data/probe_ladder_v2_mlp2/analysis/seed_level_results.json :: levels.bvc_qcd.mlp[3].mean :: sha256 53e1ec81343ac2c1 :: natural log of 1-AUC, lower is better
\newcommand{\ProbeLogOneMinusAucSdBvcQcdMlpOneseven}{0.1010}  % experiments/FIGS/data/probe_ladder_v2_mlp2/analysis/seed_level_results.json :: levels.bvc_qcd.mlp[3].seed_sd :: sha256 53e1ec81343ac2c1
\newcommand{\ProbeOmaBvcQcdMlpOneseven}{\ensuremath{(3.8 \pm 0.4)\times10^{-2}}}  % experiments/FIGS/data/probe_ladder_v2_mlp2/analysis/seed_level_results.json :: table[bvc_qcd,mlp,17].auc :: sha256 53e1ec81343ac2c1 :: 1-AUC, mean +- SD over 5 seeds
\newcommand{\ProbeBkgLeftBvcQcdMlpOneseven}{458.8}  % experiments/FIGS/data/probe_ladder_v2_mlp2/analysis/seed_level_results.json :: levels.bvc_qcd.mlp[3].rejection_points['0.90'].mean_n_bkg_pass :: sha256 53e1ec81343ac2c1 :: background jets passing the cut, mean over seeds
\newcommand{\ProbeRejBvcQcdMlpOneseven}{11.8\,\ensuremath{\pm}\,1.5}  % experiments/FIGS/data/probe_ladder_v2_mlp2/analysis/seed_level_results.json :: table[bvc_qcd,mlp,17].rejection_points['0.90'].rejection :: sha256 53e1ec81343ac2c1 :: at 90% signal efficiency, mean +- SD over seeds; 0 of 5 seeds at the cap
\newcommand{\ProbeOmaRatioBvcQcdMlpOnesixtwoOverOneeighteight}{1.11}  % experiments/FIGS/data/probe_ladder_v2_mlp2/analysis/seed_level_results.json :: table[bvc_qcd,mlp,162/188].auc :: sha256 53e1ec81343ac2c1 :: ratio of the seed means of 1-AUC, coarser over finer
\newcommand{\ProbeOmaRatioBvcQcdMlpFourthreeOverOneeighteight}{1.13}  % experiments/FIGS/data/probe_ladder_v2_mlp2/analysis/seed_level_results.json :: table[bvc_qcd,mlp,43/188].auc :: sha256 53e1ec81343ac2c1 :: ratio of the seed means of 1-AUC, coarser over finer
\newcommand{\ProbeOmaRatioBvcQcdMlpOnesevenOverOneeighteight}{2.6}  % experiments/FIGS/data/probe_ladder_v2_mlp2/analysis/seed_level_results.json :: table[bvc_qcd,mlp,17/188].auc :: sha256 53e1ec81343ac2c1 :: ratio of the seed means of 1-AUC, coarser over finer
\newcommand{\ProbeOmaRatioBvcQcdMlpFourthreeOverOnesixtwo}{1.01}  % experiments/FIGS/data/probe_ladder_v2_mlp2/analysis/seed_level_results.json :: table[bvc_qcd,mlp,43/162].auc :: sha256 53e1ec81343ac2c1 :: ratio of the seed means of 1-AUC, coarser over finer
\newcommand{\ProbeOmaRatioBvcQcdMlpOnesevenOverOnesixtwo}{2.3}  % experiments/FIGS/data/probe_ladder_v2_mlp2/analysis/seed_level_results.json :: table[bvc_qcd,mlp,17/162].auc :: sha256 53e1ec81343ac2c1 :: ratio of the seed means of 1-AUC, coarser over finer
\newcommand{\ProbeOmaRatioBvcQcdMlpOnesevenOverFourthree}{2.3}  % experiments/FIGS/data/probe_ladder_v2_mlp2/analysis/seed_level_results.json :: table[bvc_qcd,mlp,17/43].auc :: sha256 53e1ec81343ac2c1 :: ratio of the seed means of 1-AUC, coarser over finer
\newcommand{\ProbeAucBvcResonantLinearOneeighteight}{0.99768}  % experiments/FIGS/data/probe_ladder_v2_mlp2/analysis/seed_level_results.json :: levels.bvc_resonant.linear[0].mean_auc :: sha256 53e1ec81343ac2c1 :: 5 seeds, 0 saturated at AUC=1
\newcommand{\ProbeAucSdBvcResonantLinearOneeighteight}{0.00007}  % experiments/FIGS/data/probe_ladder_v2_mlp2/analysis/seed_level_results.json :: table[bvc_resonant,linear,188].auc (SD over seeds, ddof=1) :: sha256 53e1ec81343ac2c1
\newcommand{\ProbeLogOneMinusAucBvcResonantLinearOneeighteight}{\ensuremath{-}6.0683}  % experiments/FIGS/data/probe_ladder_v2_mlp2/analysis/seed_level_results.json :: levels.bvc_resonant.linear[0].mean :: sha256 53e1ec81343ac2c1 :: natural log of 1-AUC, lower is better
\newcommand{\ProbeLogOneMinusAucSdBvcResonantLinearOneeighteight}{0.0317}  % experiments/FIGS/data/probe_ladder_v2_mlp2/analysis/seed_level_results.json :: levels.bvc_resonant.linear[0].seed_sd :: sha256 53e1ec81343ac2c1
\newcommand{\ProbeOmaBvcResonantLinearOneeighteight}{\ensuremath{(2.32 \pm 0.07)\times10^{-3}}}  % experiments/FIGS/data/probe_ladder_v2_mlp2/analysis/seed_level_results.json :: table[bvc_resonant,linear,188].auc :: sha256 53e1ec81343ac2c1 :: 1-AUC, mean +- SD over 5 seeds
\newcommand{\ProbeBkgLeftBvcResonantLinearOneeighteight}{9.0}  % experiments/FIGS/data/probe_ladder_v2_mlp2/analysis/seed_level_results.json :: levels.bvc_resonant.linear[0].rejection_points['0.90'].mean_n_bkg_pass :: sha256 53e1ec81343ac2c1 :: background jets passing the cut, mean over seeds
\newcommand{\ProbeRejBvcResonantLinearOneeighteight}{1{,}380\,\ensuremath{\pm}\,350}  % experiments/FIGS/data/probe_ladder_v2_mlp2/analysis/seed_level_results.json :: table[bvc_resonant,linear,188].rejection_points['0.90'].rejection :: sha256 53e1ec81343ac2c1 :: at 90% signal efficiency, mean +- SD over seeds; 0 of 5 seeds at the cap
\newcommand{\ProbeAucBvcResonantLinearOnesixtwo}{0.99764}  % experiments/FIGS/data/probe_ladder_v2_mlp2/analysis/seed_level_results.json :: levels.bvc_resonant.linear[1].mean_auc :: sha256 53e1ec81343ac2c1 :: 5 seeds, 0 saturated at AUC=1
\newcommand{\ProbeAucSdBvcResonantLinearOnesixtwo}{0.00006}  % experiments/FIGS/data/probe_ladder_v2_mlp2/analysis/seed_level_results.json :: table[bvc_resonant,linear,162].auc (SD over seeds, ddof=1) :: sha256 53e1ec81343ac2c1
\newcommand{\ProbeLogOneMinusAucBvcResonantLinearOnesixtwo}{\ensuremath{-}6.0493}  % experiments/FIGS/data/probe_ladder_v2_mlp2/analysis/seed_level_results.json :: levels.bvc_resonant.linear[1].mean :: sha256 53e1ec81343ac2c1 :: natural log of 1-AUC, lower is better
\newcommand{\ProbeLogOneMinusAucSdBvcResonantLinearOnesixtwo}{0.0261}  % experiments/FIGS/data/probe_ladder_v2_mlp2/analysis/seed_level_results.json :: levels.bvc_resonant.linear[1].seed_sd :: sha256 53e1ec81343ac2c1
\newcommand{\ProbeOmaBvcResonantLinearOnesixtwo}{\ensuremath{(2.36 \pm 0.06)\times10^{-3}}}  % experiments/FIGS/data/probe_ladder_v2_mlp2/analysis/seed_level_results.json :: table[bvc_resonant,linear,162].auc :: sha256 53e1ec81343ac2c1 :: 1-AUC, mean +- SD over 5 seeds
\newcommand{\ProbeBkgLeftBvcResonantLinearOnesixtwo}{9.4}  % experiments/FIGS/data/probe_ladder_v2_mlp2/analysis/seed_level_results.json :: levels.bvc_resonant.linear[1].rejection_points['0.90'].mean_n_bkg_pass :: sha256 53e1ec81343ac2c1 :: background jets passing the cut, mean over seeds
\newcommand{\ProbeRejBvcResonantLinearOnesixtwo}{1{,}290\,\ensuremath{\pm}\,240}  % experiments/FIGS/data/probe_ladder_v2_mlp2/analysis/seed_level_results.json :: table[bvc_resonant,linear,162].rejection_points['0.90'].rejection :: sha256 53e1ec81343ac2c1 :: at 90% signal efficiency, mean +- SD over seeds; 0 of 5 seeds at the cap
\newcommand{\ProbeAucBvcResonantLinearFourthree}{0.99748}  % experiments/FIGS/data/probe_ladder_v2_mlp2/analysis/seed_level_results.json :: levels.bvc_resonant.linear[2].mean_auc :: sha256 53e1ec81343ac2c1 :: 5 seeds, 0 saturated at AUC=1
\newcommand{\ProbeAucSdBvcResonantLinearFourthree}{0.00007}  % experiments/FIGS/data/probe_ladder_v2_mlp2/analysis/seed_level_results.json :: table[bvc_resonant,linear,43].auc (SD over seeds, ddof=1) :: sha256 53e1ec81343ac2c1
\newcommand{\ProbeLogOneMinusAucBvcResonantLinearFourthree}{\ensuremath{-}5.9834}  % experiments/FIGS/data/probe_ladder_v2_mlp2/analysis/seed_level_results.json :: levels.bvc_resonant.linear[2].mean :: sha256 53e1ec81343ac2c1 :: natural log of 1-AUC, lower is better
\newcommand{\ProbeLogOneMinusAucSdBvcResonantLinearFourthree}{0.0264}  % experiments/FIGS/data/probe_ladder_v2_mlp2/analysis/seed_level_results.json :: levels.bvc_resonant.linear[2].seed_sd :: sha256 53e1ec81343ac2c1
\newcommand{\ProbeOmaBvcResonantLinearFourthree}{\ensuremath{(2.52 \pm 0.07)\times10^{-3}}}  % experiments/FIGS/data/probe_ladder_v2_mlp2/analysis/seed_level_results.json :: table[bvc_resonant,linear,43].auc :: sha256 53e1ec81343ac2c1 :: 1-AUC, mean +- SD over 5 seeds
\newcommand{\ProbeBkgLeftBvcResonantLinearFourthree}{11.0}  % experiments/FIGS/data/probe_ladder_v2_mlp2/analysis/seed_level_results.json :: levels.bvc_resonant.linear[2].rejection_points['0.90'].mean_n_bkg_pass :: sha256 53e1ec81343ac2c1 :: background jets passing the cut, mean over seeds
\newcommand{\ProbeRejBvcResonantLinearFourthree}{1{,}130\,\ensuremath{\pm}\,260}  % experiments/FIGS/data/probe_ladder_v2_mlp2/analysis/seed_level_results.json :: table[bvc_resonant,linear,43].rejection_points['0.90'].rejection :: sha256 53e1ec81343ac2c1 :: at 90% signal efficiency, mean +- SD over seeds; 0 of 5 seeds at the cap
\newcommand{\ProbeAucBvcResonantLinearOneseven}{0.99067}  % experiments/FIGS/data/probe_ladder_v2_mlp2/analysis/seed_level_results.json :: levels.bvc_resonant.linear[3].mean_auc :: sha256 53e1ec81343ac2c1 :: 5 seeds, 0 saturated at AUC=1
\newcommand{\ProbeAucSdBvcResonantLinearOneseven}{0.00147}  % experiments/FIGS/data/probe_ladder_v2_mlp2/analysis/seed_level_results.json :: table[bvc_resonant,linear,17].auc (SD over seeds, ddof=1) :: sha256 53e1ec81343ac2c1
\newcommand{\ProbeLogOneMinusAucBvcResonantLinearOneseven}{\ensuremath{-}4.6848}  % experiments/FIGS/data/probe_ladder_v2_mlp2/analysis/seed_level_results.json :: levels.bvc_resonant.linear[3].mean :: sha256 53e1ec81343ac2c1 :: natural log of 1-AUC, lower is better
\newcommand{\ProbeLogOneMinusAucSdBvcResonantLinearOneseven}{0.1619}  % experiments/FIGS/data/probe_ladder_v2_mlp2/analysis/seed_level_results.json :: levels.bvc_resonant.linear[3].seed_sd :: sha256 53e1ec81343ac2c1
\newcommand{\ProbeOmaBvcResonantLinearOneseven}{\ensuremath{(9.3 \pm 1.5)\times10^{-3}}}  % experiments/FIGS/data/probe_ladder_v2_mlp2/analysis/seed_level_results.json :: table[bvc_resonant,linear,17].auc :: sha256 53e1ec81343ac2c1 :: 1-AUC, mean +- SD over 5 seeds
\newcommand{\ProbeBkgLeftBvcResonantLinearOneseven}{90.8}  % experiments/FIGS/data/probe_ladder_v2_mlp2/analysis/seed_level_results.json :: levels.bvc_resonant.linear[3].rejection_points['0.90'].mean_n_bkg_pass :: sha256 53e1ec81343ac2c1 :: background jets passing the cut, mean over seeds
\newcommand{\ProbeRejBvcResonantLinearOneseven}{134\,\ensuremath{\pm}\,21}  % experiments/FIGS/data/probe_ladder_v2_mlp2/analysis/seed_level_results.json :: table[bvc_resonant,linear,17].rejection_points['0.90'].rejection :: sha256 53e1ec81343ac2c1 :: at 90% signal efficiency, mean +- SD over seeds; 0 of 5 seeds at the cap
\newcommand{\ProbeOmaRatioBvcResonantLinearOnesixtwoOverOneeighteight}{1.02}  % experiments/FIGS/data/probe_ladder_v2_mlp2/analysis/seed_level_results.json :: table[bvc_resonant,linear,162/188].auc :: sha256 53e1ec81343ac2c1 :: ratio of the seed means of 1-AUC, coarser over finer
\newcommand{\ProbeOmaRatioBvcResonantLinearFourthreeOverOneeighteight}{1.09}  % experiments/FIGS/data/probe_ladder_v2_mlp2/analysis/seed_level_results.json :: table[bvc_resonant,linear,43/188].auc :: sha256 53e1ec81343ac2c1 :: ratio of the seed means of 1-AUC, coarser over finer
\newcommand{\ProbeOmaRatioBvcResonantLinearOnesevenOverOneeighteight}{4.0}  % experiments/FIGS/data/probe_ladder_v2_mlp2/analysis/seed_level_results.json :: table[bvc_resonant,linear,17/188].auc :: sha256 53e1ec81343ac2c1 :: ratio of the seed means of 1-AUC, coarser over finer
\newcommand{\ProbeOmaRatioBvcResonantLinearFourthreeOverOnesixtwo}{1.07}  % experiments/FIGS/data/probe_ladder_v2_mlp2/analysis/seed_level_results.json :: table[bvc_resonant,linear,43/162].auc :: sha256 53e1ec81343ac2c1 :: ratio of the seed means of 1-AUC, coarser over finer
\newcommand{\ProbeOmaRatioBvcResonantLinearOnesevenOverOnesixtwo}{4.0}  % experiments/FIGS/data/probe_ladder_v2_mlp2/analysis/seed_level_results.json :: table[bvc_resonant,linear,17/162].auc :: sha256 53e1ec81343ac2c1 :: ratio of the seed means of 1-AUC, coarser over finer
\newcommand{\ProbeOmaRatioBvcResonantLinearOnesevenOverFourthree}{3.7}  % experiments/FIGS/data/probe_ladder_v2_mlp2/analysis/seed_level_results.json :: table[bvc_resonant,linear,17/43].auc :: sha256 53e1ec81343ac2c1 :: ratio of the seed means of 1-AUC, coarser over finer
\newcommand{\ProbeAucBvcResonantMlpOneeighteight}{0.99778}  % experiments/FIGS/data/probe_ladder_v2_mlp2/analysis/seed_level_results.json :: levels.bvc_resonant.mlp[0].mean_auc :: sha256 53e1ec81343ac2c1 :: 5 seeds, 0 saturated at AUC=1
\newcommand{\ProbeAucSdBvcResonantMlpOneeighteight}{0.00009}  % experiments/FIGS/data/probe_ladder_v2_mlp2/analysis/seed_level_results.json :: table[bvc_resonant,mlp,188].auc (SD over seeds, ddof=1) :: sha256 53e1ec81343ac2c1
\newcommand{\ProbeLogOneMinusAucBvcResonantMlpOneeighteight}{\ensuremath{-}6.1102}  % experiments/FIGS/data/probe_ladder_v2_mlp2/analysis/seed_level_results.json :: levels.bvc_resonant.mlp[0].mean :: sha256 53e1ec81343ac2c1 :: natural log of 1-AUC, lower is better
\newcommand{\ProbeLogOneMinusAucSdBvcResonantMlpOneeighteight}{0.0383}  % experiments/FIGS/data/probe_ladder_v2_mlp2/analysis/seed_level_results.json :: levels.bvc_resonant.mlp[0].seed_sd :: sha256 53e1ec81343ac2c1
\newcommand{\ProbeOmaBvcResonantMlpOneeighteight}{\ensuremath{(2.22 \pm 0.09)\times10^{-3}}}  % experiments/FIGS/data/probe_ladder_v2_mlp2/analysis/seed_level_results.json :: table[bvc_resonant,mlp,188].auc :: sha256 53e1ec81343ac2c1 :: 1-AUC, mean +- SD over 5 seeds
\newcommand{\ProbeBkgLeftBvcResonantMlpOneeighteight}{9.8}  % experiments/FIGS/data/probe_ladder_v2_mlp2/analysis/seed_level_results.json :: levels.bvc_resonant.mlp[0].rejection_points['0.90'].mean_n_bkg_pass :: sha256 53e1ec81343ac2c1 :: background jets passing the cut, mean over seeds
\newcommand{\ProbeRejBvcResonantMlpOneeighteight}{1{,}230\,\ensuremath{\pm}\,180}  % experiments/FIGS/data/probe_ladder_v2_mlp2/analysis/seed_level_results.json :: table[bvc_resonant,mlp,188].rejection_points['0.90'].rejection :: sha256 53e1ec81343ac2c1 :: at 90% signal efficiency, mean +- SD over seeds; 0 of 5 seeds at the cap
\newcommand{\ProbeAucBvcResonantMlpOnesixtwo}{0.99768}  % experiments/FIGS/data/probe_ladder_v2_mlp2/analysis/seed_level_results.json :: levels.bvc_resonant.mlp[1].mean_auc :: sha256 53e1ec81343ac2c1 :: 5 seeds, 0 saturated at AUC=1
\newcommand{\ProbeAucSdBvcResonantMlpOnesixtwo}{0.00005}  % experiments/FIGS/data/probe_ladder_v2_mlp2/analysis/seed_level_results.json :: table[bvc_resonant,mlp,162].auc (SD over seeds, ddof=1) :: sha256 53e1ec81343ac2c1
\newcommand{\ProbeLogOneMinusAucBvcResonantMlpOnesixtwo}{\ensuremath{-}6.0667}  % experiments/FIGS/data/probe_ladder_v2_mlp2/analysis/seed_level_results.json :: levels.bvc_resonant.mlp[1].mean :: sha256 53e1ec81343ac2c1 :: natural log of 1-AUC, lower is better
\newcommand{\ProbeLogOneMinusAucSdBvcResonantMlpOnesixtwo}{0.0220}  % experiments/FIGS/data/probe_ladder_v2_mlp2/analysis/seed_level_results.json :: levels.bvc_resonant.mlp[1].seed_sd :: sha256 53e1ec81343ac2c1
\newcommand{\ProbeOmaBvcResonantMlpOnesixtwo}{\ensuremath{(2.32 \pm 0.05)\times10^{-3}}}  % experiments/FIGS/data/probe_ladder_v2_mlp2/analysis/seed_level_results.json :: table[bvc_resonant,mlp,162].auc :: sha256 53e1ec81343ac2c1 :: 1-AUC, mean +- SD over 5 seeds
\newcommand{\ProbeBkgLeftBvcResonantMlpOnesixtwo}{10.4}  % experiments/FIGS/data/probe_ladder_v2_mlp2/analysis/seed_level_results.json :: levels.bvc_resonant.mlp[1].rejection_points['0.90'].mean_n_bkg_pass :: sha256 53e1ec81343ac2c1 :: background jets passing the cut, mean over seeds
\newcommand{\ProbeRejBvcResonantMlpOnesixtwo}{1{,}170\,\ensuremath{\pm}\,220}  % experiments/FIGS/data/probe_ladder_v2_mlp2/analysis/seed_level_results.json :: table[bvc_resonant,mlp,162].rejection_points['0.90'].rejection :: sha256 53e1ec81343ac2c1 :: at 90% signal efficiency, mean +- SD over seeds; 0 of 5 seeds at the cap
\newcommand{\ProbeAucBvcResonantMlpFourthree}{0.99747}  % experiments/FIGS/data/probe_ladder_v2_mlp2/analysis/seed_level_results.json :: levels.bvc_resonant.mlp[2].mean_auc :: sha256 53e1ec81343ac2c1 :: 5 seeds, 0 saturated at AUC=1
\newcommand{\ProbeAucSdBvcResonantMlpFourthree}{0.00004}  % experiments/FIGS/data/probe_ladder_v2_mlp2/analysis/seed_level_results.json :: table[bvc_resonant,mlp,43].auc (SD over seeds, ddof=1) :: sha256 53e1ec81343ac2c1
\newcommand{\ProbeLogOneMinusAucBvcResonantMlpFourthree}{\ensuremath{-}5.9792}  % experiments/FIGS/data/probe_ladder_v2_mlp2/analysis/seed_level_results.json :: levels.bvc_resonant.mlp[2].mean :: sha256 53e1ec81343ac2c1 :: natural log of 1-AUC, lower is better
\newcommand{\ProbeLogOneMinusAucSdBvcResonantMlpFourthree}{0.0164}  % experiments/FIGS/data/probe_ladder_v2_mlp2/analysis/seed_level_results.json :: levels.bvc_resonant.mlp[2].seed_sd :: sha256 53e1ec81343ac2c1
\newcommand{\ProbeOmaBvcResonantMlpFourthree}{\ensuremath{(2.53 \pm 0.04)\times10^{-3}}}  % experiments/FIGS/data/probe_ladder_v2_mlp2/analysis/seed_level_results.json :: table[bvc_resonant,mlp,43].auc :: sha256 53e1ec81343ac2c1 :: 1-AUC, mean +- SD over 5 seeds
\newcommand{\ProbeBkgLeftBvcResonantMlpFourthree}{11.8}  % experiments/FIGS/data/probe_ladder_v2_mlp2/analysis/seed_level_results.json :: levels.bvc_resonant.mlp[2].rejection_points['0.90'].mean_n_bkg_pass :: sha256 53e1ec81343ac2c1 :: background jets passing the cut, mean over seeds
\newcommand{\ProbeRejBvcResonantMlpFourthree}{1{,}030\,\ensuremath{\pm}\,170}  % experiments/FIGS/data/probe_ladder_v2_mlp2/analysis/seed_level_results.json :: table[bvc_resonant,mlp,43].rejection_points['0.90'].rejection :: sha256 53e1ec81343ac2c1 :: at 90% signal efficiency, mean +- SD over seeds; 0 of 5 seeds at the cap
\newcommand{\ProbeAucBvcResonantMlpOneseven}{0.99274}  % experiments/FIGS/data/probe_ladder_v2_mlp2/analysis/seed_level_results.json :: levels.bvc_resonant.mlp[3].mean_auc :: sha256 53e1ec81343ac2c1 :: 5 seeds, 0 saturated at AUC=1
\newcommand{\ProbeAucSdBvcResonantMlpOneseven}{0.00089}  % experiments/FIGS/data/probe_ladder_v2_mlp2/analysis/seed_level_results.json :: table[bvc_resonant,mlp,17].auc (SD over seeds, ddof=1) :: sha256 53e1ec81343ac2c1
\newcommand{\ProbeLogOneMinusAucBvcResonantMlpOneseven}{\ensuremath{-}4.9309}  % experiments/FIGS/data/probe_ladder_v2_mlp2/analysis/seed_level_results.json :: levels.bvc_resonant.mlp[3].mean :: sha256 53e1ec81343ac2c1 :: natural log of 1-AUC, lower is better
\newcommand{\ProbeLogOneMinusAucSdBvcResonantMlpOneseven}{0.1224}  % experiments/FIGS/data/probe_ladder_v2_mlp2/analysis/seed_level_results.json :: levels.bvc_resonant.mlp[3].seed_sd :: sha256 53e1ec81343ac2c1
\newcommand{\ProbeOmaBvcResonantMlpOneseven}{\ensuremath{(7.3 \pm 0.9)\times10^{-3}}}  % experiments/FIGS/data/probe_ladder_v2_mlp2/analysis/seed_level_results.json :: table[bvc_resonant,mlp,17].auc :: sha256 53e1ec81343ac2c1 :: 1-AUC, mean +- SD over 5 seeds
\newcommand{\ProbeBkgLeftBvcResonantMlpOneseven}{62.6}  % experiments/FIGS/data/probe_ladder_v2_mlp2/analysis/seed_level_results.json :: levels.bvc_resonant.mlp[3].rejection_points['0.90'].mean_n_bkg_pass :: sha256 53e1ec81343ac2c1 :: background jets passing the cut, mean over seeds
\newcommand{\ProbeRejBvcResonantMlpOneseven}{192\,\ensuremath{\pm}\,21}  % experiments/FIGS/data/probe_ladder_v2_mlp2/analysis/seed_level_results.json :: table[bvc_resonant,mlp,17].rejection_points['0.90'].rejection :: sha256 53e1ec81343ac2c1 :: at 90% signal efficiency, mean +- SD over seeds; 0 of 5 seeds at the cap
\newcommand{\ProbeOmaRatioBvcResonantMlpOnesixtwoOverOneeighteight}{1.04}  % experiments/FIGS/data/probe_ladder_v2_mlp2/analysis/seed_level_results.json :: table[bvc_resonant,mlp,162/188].auc :: sha256 53e1ec81343ac2c1 :: ratio of the seed means of 1-AUC, coarser over finer
\newcommand{\ProbeOmaRatioBvcResonantMlpFourthreeOverOneeighteight}{1.14}  % experiments/FIGS/data/probe_ladder_v2_mlp2/analysis/seed_level_results.json :: table[bvc_resonant,mlp,43/188].auc :: sha256 53e1ec81343ac2c1 :: ratio of the seed means of 1-AUC, coarser over finer
\newcommand{\ProbeOmaRatioBvcResonantMlpOnesevenOverOneeighteight}{3.3}  % experiments/FIGS/data/probe_ladder_v2_mlp2/analysis/seed_level_results.json :: table[bvc_resonant,mlp,17/188].auc :: sha256 53e1ec81343ac2c1 :: ratio of the seed means of 1-AUC, coarser over finer
\newcommand{\ProbeOmaRatioBvcResonantMlpFourthreeOverOnesixtwo}{1.09}  % experiments/FIGS/data/probe_ladder_v2_mlp2/analysis/seed_level_results.json :: table[bvc_resonant,mlp,43/162].auc :: sha256 53e1ec81343ac2c1 :: ratio of the seed means of 1-AUC, coarser over finer
\newcommand{\ProbeOmaRatioBvcResonantMlpOnesevenOverOnesixtwo}{3.1}  % experiments/FIGS/data/probe_ladder_v2_mlp2/analysis/seed_level_results.json :: table[bvc_resonant,mlp,17/162].auc :: sha256 53e1ec81343ac2c1 :: ratio of the seed means of 1-AUC, coarser over finer
\newcommand{\ProbeOmaRatioBvcResonantMlpOnesevenOverFourthree}{2.9}  % experiments/FIGS/data/probe_ladder_v2_mlp2/analysis/seed_level_results.json :: table[bvc_resonant,mlp,17/43].auc :: sha256 53e1ec81343ac2c1 :: ratio of the seed means of 1-AUC, coarser over finer
\newcommand{\ProbeAucEeVsMmLinearOneeighteight}{\ensuremath{1^{\dagger}}}  % experiments/FIGS/data/probe_ladder_v2_mlp2/analysis/seed_level_results.json :: levels.ee_vs_mm.linear[0].mean_auc :: sha256 53e1ec81343ac2c1 :: 5 seeds, 5 saturated at AUC=1
\newcommand{\ProbeAucSdEeVsMmLinearOneeighteight}{0.00000}  % experiments/FIGS/data/probe_ladder_v2_mlp2/analysis/seed_level_results.json :: table[ee_vs_mm,linear,188].auc (SD over seeds, ddof=1) :: sha256 53e1ec81343ac2c1
\newcommand{\ProbeLogOneMinusAucEeVsMmLinearOneeighteight}{\ensuremath{-}16.9579}  % experiments/FIGS/data/probe_ladder_v2_mlp2/analysis/seed_level_results.json :: levels.ee_vs_mm.linear[0].mean :: sha256 53e1ec81343ac2c1 :: natural log of 1-AUC, lower is better
\newcommand{\ProbeLogOneMinusAucSdEeVsMmLinearOneeighteight}{0.0000}  % experiments/FIGS/data/probe_ladder_v2_mlp2/analysis/seed_level_results.json :: levels.ee_vs_mm.linear[0].seed_sd :: sha256 53e1ec81343ac2c1
\newcommand{\ProbeOmaEeVsMmLinearOneeighteight}{---}  % experiments/FIGS/data/probe_ladder_v2_mlp2/analysis/seed_level_results.json :: table[ee_vs_mm,linear,188].auc :: sha256 53e1ec81343ac2c1 :: not quoted: the AUC reached 1 in a seed, so 1-AUC is only a bound
\newcommand{\ProbeBkgLeftEeVsMmLinearOneeighteight}{0.0}  % experiments/FIGS/data/probe_ladder_v2_mlp2/analysis/seed_level_results.json :: levels.ee_vs_mm.linear[0].rejection_points['0.90'].mean_n_bkg_pass :: sha256 53e1ec81343ac2c1 :: background jets passing the cut, mean over seeds
\newcommand{\ProbeRejEeVsMmLinearOneeighteight}{\ensuremath{>}1{,}571}  % experiments/FIGS/data/probe_ladder_v2_mlp2/analysis/seed_level_results.json :: table[ee_vs_mm,linear,188].rejection_points['0.90'].rejection :: sha256 53e1ec81343ac2c1 :: at 90% signal efficiency, mean +- SD over seeds; 5 of 5 seeds at the cap
\newcommand{\ProbeAucEeVsMmLinearOnesixtwo}{\ensuremath{1^{\dagger}}}  % experiments/FIGS/data/probe_ladder_v2_mlp2/analysis/seed_level_results.json :: levels.ee_vs_mm.linear[1].mean_auc :: sha256 53e1ec81343ac2c1 :: 5 seeds, 5 saturated at AUC=1
\newcommand{\ProbeAucSdEeVsMmLinearOnesixtwo}{0.00000}  % experiments/FIGS/data/probe_ladder_v2_mlp2/analysis/seed_level_results.json :: table[ee_vs_mm,linear,162].auc (SD over seeds, ddof=1) :: sha256 53e1ec81343ac2c1
\newcommand{\ProbeLogOneMinusAucEeVsMmLinearOnesixtwo}{\ensuremath{-}16.9579}  % experiments/FIGS/data/probe_ladder_v2_mlp2/analysis/seed_level_results.json :: levels.ee_vs_mm.linear[1].mean :: sha256 53e1ec81343ac2c1 :: natural log of 1-AUC, lower is better
\newcommand{\ProbeLogOneMinusAucSdEeVsMmLinearOnesixtwo}{0.0000}  % experiments/FIGS/data/probe_ladder_v2_mlp2/analysis/seed_level_results.json :: levels.ee_vs_mm.linear[1].seed_sd :: sha256 53e1ec81343ac2c1
\newcommand{\ProbeOmaEeVsMmLinearOnesixtwo}{---}  % experiments/FIGS/data/probe_ladder_v2_mlp2/analysis/seed_level_results.json :: table[ee_vs_mm,linear,162].auc :: sha256 53e1ec81343ac2c1 :: not quoted: the AUC reached 1 in a seed, so 1-AUC is only a bound
\newcommand{\ProbeBkgLeftEeVsMmLinearOnesixtwo}{0.0}  % experiments/FIGS/data/probe_ladder_v2_mlp2/analysis/seed_level_results.json :: levels.ee_vs_mm.linear[1].rejection_points['0.90'].mean_n_bkg_pass :: sha256 53e1ec81343ac2c1 :: background jets passing the cut, mean over seeds
\newcommand{\ProbeRejEeVsMmLinearOnesixtwo}{\ensuremath{>}1{,}571}  % experiments/FIGS/data/probe_ladder_v2_mlp2/analysis/seed_level_results.json :: table[ee_vs_mm,linear,162].rejection_points['0.90'].rejection :: sha256 53e1ec81343ac2c1 :: at 90% signal efficiency, mean +- SD over seeds; 5 of 5 seeds at the cap
\newcommand{\ProbeAucEeVsMmLinearFourthree}{0.97519}  % experiments/FIGS/data/probe_ladder_v2_mlp2/analysis/seed_level_results.json :: levels.ee_vs_mm.linear[2].mean_auc :: sha256 53e1ec81343ac2c1 :: 5 seeds, 0 saturated at AUC=1
\newcommand{\ProbeAucSdEeVsMmLinearFourthree}{0.02073}  % experiments/FIGS/data/probe_ladder_v2_mlp2/analysis/seed_level_results.json :: table[ee_vs_mm,linear,43].auc (SD over seeds, ddof=1) :: sha256 53e1ec81343ac2c1
\newcommand{\ProbeLogOneMinusAucEeVsMmLinearFourthree}{\ensuremath{-}4.0326}  % experiments/FIGS/data/probe_ladder_v2_mlp2/analysis/seed_level_results.json :: levels.ee_vs_mm.linear[2].mean :: sha256 53e1ec81343ac2c1 :: natural log of 1-AUC, lower is better
\newcommand{\ProbeLogOneMinusAucSdEeVsMmLinearFourthree}{0.9467}  % experiments/FIGS/data/probe_ladder_v2_mlp2/analysis/seed_level_results.json :: levels.ee_vs_mm.linear[2].seed_sd :: sha256 53e1ec81343ac2c1
\newcommand{\ProbeOmaEeVsMmLinearFourthree}{\ensuremath{(2.5 \pm 2.1)\times10^{-2}}}  % experiments/FIGS/data/probe_ladder_v2_mlp2/analysis/seed_level_results.json :: table[ee_vs_mm,linear,43].auc :: sha256 53e1ec81343ac2c1 :: 1-AUC, mean +- SD over 5 seeds
\newcommand{\ProbeBkgLeftEeVsMmLinearFourthree}{155.8}  % experiments/FIGS/data/probe_ladder_v2_mlp2/analysis/seed_level_results.json :: levels.ee_vs_mm.linear[2].rejection_points['0.90'].mean_n_bkg_pass :: sha256 53e1ec81343ac2c1 :: background jets passing the cut, mean over seeds
\newcommand{\ProbeRejEeVsMmLinearFourthree}{80\,\ensuremath{\pm}\,70}  % experiments/FIGS/data/probe_ladder_v2_mlp2/analysis/seed_level_results.json :: table[ee_vs_mm,linear,43].rejection_points['0.90'].rejection :: sha256 53e1ec81343ac2c1 :: at 90% signal efficiency, mean +- SD over seeds; 0 of 5 seeds at the cap
\newcommand{\ProbeAucEeVsMmLinearOneseven}{0.95832}  % experiments/FIGS/data/probe_ladder_v2_mlp2/analysis/seed_level_results.json :: levels.ee_vs_mm.linear[3].mean_auc :: sha256 53e1ec81343ac2c1 :: 5 seeds, 0 saturated at AUC=1
\newcommand{\ProbeAucSdEeVsMmLinearOneseven}{0.00970}  % experiments/FIGS/data/probe_ladder_v2_mlp2/analysis/seed_level_results.json :: table[ee_vs_mm,linear,17].auc (SD over seeds, ddof=1) :: sha256 53e1ec81343ac2c1
\newcommand{\ProbeLogOneMinusAucEeVsMmLinearOneseven}{\ensuremath{-}3.2013}  % experiments/FIGS/data/probe_ladder_v2_mlp2/analysis/seed_level_results.json :: levels.ee_vs_mm.linear[3].mean :: sha256 53e1ec81343ac2c1 :: natural log of 1-AUC, lower is better
\newcommand{\ProbeLogOneMinusAucSdEeVsMmLinearOneseven}{0.2488}  % experiments/FIGS/data/probe_ladder_v2_mlp2/analysis/seed_level_results.json :: levels.ee_vs_mm.linear[3].seed_sd :: sha256 53e1ec81343ac2c1
\newcommand{\ProbeOmaEeVsMmLinearOneseven}{\ensuremath{(4.2 \pm 1.0)\times10^{-2}}}  % experiments/FIGS/data/probe_ladder_v2_mlp2/analysis/seed_level_results.json :: table[ee_vs_mm,linear,17].auc :: sha256 53e1ec81343ac2c1 :: 1-AUC, mean +- SD over 5 seeds
\newcommand{\ProbeBkgLeftEeVsMmLinearOneseven}{289.6}  % experiments/FIGS/data/probe_ladder_v2_mlp2/analysis/seed_level_results.json :: levels.ee_vs_mm.linear[3].rejection_points['0.90'].mean_n_bkg_pass :: sha256 53e1ec81343ac2c1 :: background jets passing the cut, mean over seeds
\newcommand{\ProbeRejEeVsMmLinearOneseven}{17\,\ensuremath{\pm}\,6}  % experiments/FIGS/data/probe_ladder_v2_mlp2/analysis/seed_level_results.json :: table[ee_vs_mm,linear,17].rejection_points['0.90'].rejection :: sha256 53e1ec81343ac2c1 :: at 90% signal efficiency, mean +- SD over seeds; 0 of 5 seeds at the cap
\newcommand{\ProbeOmaRatioEeVsMmLinearOnesevenOverFourthree}{1.68}  % experiments/FIGS/data/probe_ladder_v2_mlp2/analysis/seed_level_results.json :: table[ee_vs_mm,linear,17/43].auc :: sha256 53e1ec81343ac2c1 :: ratio of the seed means of 1-AUC, coarser over finer
\newcommand{\ProbeAucEeVsMmMlpOneeighteight}{\ensuremath{1^{\dagger}}}  % experiments/FIGS/data/probe_ladder_v2_mlp2/analysis/seed_level_results.json :: levels.ee_vs_mm.mlp[0].mean_auc :: sha256 53e1ec81343ac2c1 :: 5 seeds, 5 saturated at AUC=1
\newcommand{\ProbeAucSdEeVsMmMlpOneeighteight}{0.00000}  % experiments/FIGS/data/probe_ladder_v2_mlp2/analysis/seed_level_results.json :: table[ee_vs_mm,mlp,188].auc (SD over seeds, ddof=1) :: sha256 53e1ec81343ac2c1
\newcommand{\ProbeLogOneMinusAucEeVsMmMlpOneeighteight}{\ensuremath{-}16.9579}  % experiments/FIGS/data/probe_ladder_v2_mlp2/analysis/seed_level_results.json :: levels.ee_vs_mm.mlp[0].mean :: sha256 53e1ec81343ac2c1 :: natural log of 1-AUC, lower is better
\newcommand{\ProbeLogOneMinusAucSdEeVsMmMlpOneeighteight}{0.0000}  % experiments/FIGS/data/probe_ladder_v2_mlp2/analysis/seed_level_results.json :: levels.ee_vs_mm.mlp[0].seed_sd :: sha256 53e1ec81343ac2c1
\newcommand{\ProbeOmaEeVsMmMlpOneeighteight}{---}  % experiments/FIGS/data/probe_ladder_v2_mlp2/analysis/seed_level_results.json :: table[ee_vs_mm,mlp,188].auc :: sha256 53e1ec81343ac2c1 :: not quoted: the AUC reached 1 in a seed, so 1-AUC is only a bound
\newcommand{\ProbeBkgLeftEeVsMmMlpOneeighteight}{0.0}  % experiments/FIGS/data/probe_ladder_v2_mlp2/analysis/seed_level_results.json :: levels.ee_vs_mm.mlp[0].rejection_points['0.90'].mean_n_bkg_pass :: sha256 53e1ec81343ac2c1 :: background jets passing the cut, mean over seeds
\newcommand{\ProbeRejEeVsMmMlpOneeighteight}{\ensuremath{>}1{,}571}  % experiments/FIGS/data/probe_ladder_v2_mlp2/analysis/seed_level_results.json :: table[ee_vs_mm,mlp,188].rejection_points['0.90'].rejection :: sha256 53e1ec81343ac2c1 :: at 90% signal efficiency, mean +- SD over seeds; 5 of 5 seeds at the cap
\newcommand{\ProbeAucEeVsMmMlpOnesixtwo}{\ensuremath{1^{\dagger}}}  % experiments/FIGS/data/probe_ladder_v2_mlp2/analysis/seed_level_results.json :: levels.ee_vs_mm.mlp[1].mean_auc :: sha256 53e1ec81343ac2c1 :: 5 seeds, 5 saturated at AUC=1
\newcommand{\ProbeAucSdEeVsMmMlpOnesixtwo}{0.00000}  % experiments/FIGS/data/probe_ladder_v2_mlp2/analysis/seed_level_results.json :: table[ee_vs_mm,mlp,162].auc (SD over seeds, ddof=1) :: sha256 53e1ec81343ac2c1
\newcommand{\ProbeLogOneMinusAucEeVsMmMlpOnesixtwo}{\ensuremath{-}16.9579}  % experiments/FIGS/data/probe_ladder_v2_mlp2/analysis/seed_level_results.json :: levels.ee_vs_mm.mlp[1].mean :: sha256 53e1ec81343ac2c1 :: natural log of 1-AUC, lower is better
\newcommand{\ProbeLogOneMinusAucSdEeVsMmMlpOnesixtwo}{0.0000}  % experiments/FIGS/data/probe_ladder_v2_mlp2/analysis/seed_level_results.json :: levels.ee_vs_mm.mlp[1].seed_sd :: sha256 53e1ec81343ac2c1
\newcommand{\ProbeOmaEeVsMmMlpOnesixtwo}{---}  % experiments/FIGS/data/probe_ladder_v2_mlp2/analysis/seed_level_results.json :: table[ee_vs_mm,mlp,162].auc :: sha256 53e1ec81343ac2c1 :: not quoted: the AUC reached 1 in a seed, so 1-AUC is only a bound
\newcommand{\ProbeBkgLeftEeVsMmMlpOnesixtwo}{0.0}  % experiments/FIGS/data/probe_ladder_v2_mlp2/analysis/seed_level_results.json :: levels.ee_vs_mm.mlp[1].rejection_points['0.90'].mean_n_bkg_pass :: sha256 53e1ec81343ac2c1 :: background jets passing the cut, mean over seeds
\newcommand{\ProbeRejEeVsMmMlpOnesixtwo}{\ensuremath{>}1{,}571}  % experiments/FIGS/data/probe_ladder_v2_mlp2/analysis/seed_level_results.json :: table[ee_vs_mm,mlp,162].rejection_points['0.90'].rejection :: sha256 53e1ec81343ac2c1 :: at 90% signal efficiency, mean +- SD over seeds; 5 of 5 seeds at the cap
\newcommand{\ProbeAucEeVsMmMlpFourthree}{0.99657}  % experiments/FIGS/data/probe_ladder_v2_mlp2/analysis/seed_level_results.json :: levels.ee_vs_mm.mlp[2].mean_auc :: sha256 53e1ec81343ac2c1 :: 5 seeds, 0 saturated at AUC=1
\newcommand{\ProbeAucSdEeVsMmMlpFourthree}{0.00300}  % experiments/FIGS/data/probe_ladder_v2_mlp2/analysis/seed_level_results.json :: table[ee_vs_mm,mlp,43].auc (SD over seeds, ddof=1) :: sha256 53e1ec81343ac2c1
\newcommand{\ProbeLogOneMinusAucEeVsMmMlpFourthree}{\ensuremath{-}6.0371}  % experiments/FIGS/data/probe_ladder_v2_mlp2/analysis/seed_level_results.json :: levels.ee_vs_mm.mlp[2].mean :: sha256 53e1ec81343ac2c1 :: natural log of 1-AUC, lower is better
\newcommand{\ProbeLogOneMinusAucSdEeVsMmMlpFourthree}{0.9962}  % experiments/FIGS/data/probe_ladder_v2_mlp2/analysis/seed_level_results.json :: levels.ee_vs_mm.mlp[2].seed_sd :: sha256 53e1ec81343ac2c1
\newcommand{\ProbeOmaEeVsMmMlpFourthree}{\ensuremath{(3.4 \pm 3.0)\times10^{-3}}}  % experiments/FIGS/data/probe_ladder_v2_mlp2/analysis/seed_level_results.json :: table[ee_vs_mm,mlp,43].auc :: sha256 53e1ec81343ac2c1 :: 1-AUC, mean +- SD over 5 seeds
\newcommand{\ProbeBkgLeftEeVsMmMlpFourthree}{33.4}  % experiments/FIGS/data/probe_ladder_v2_mlp2/analysis/seed_level_results.json :: levels.ee_vs_mm.mlp[2].rejection_points['0.90'].mean_n_bkg_pass :: sha256 53e1ec81343ac2c1 :: background jets passing the cut, mean over seeds
\newcommand{\ProbeRejEeVsMmMlpFourthree}{400\,\ensuremath{\pm}\,500}  % experiments/FIGS/data/probe_ladder_v2_mlp2/analysis/seed_level_results.json :: table[ee_vs_mm,mlp,43].rejection_points['0.90'].rejection :: sha256 53e1ec81343ac2c1 :: at 90% signal efficiency, mean +- SD over seeds; 0 of 5 seeds at the cap
\newcommand{\ProbeAucEeVsMmMlpOneseven}{0.99338}  % experiments/FIGS/data/probe_ladder_v2_mlp2/analysis/seed_level_results.json :: levels.ee_vs_mm.mlp[3].mean_auc :: sha256 53e1ec81343ac2c1 :: 5 seeds, 0 saturated at AUC=1
\newcommand{\ProbeAucSdEeVsMmMlpOneseven}{0.00065}  % experiments/FIGS/data/probe_ladder_v2_mlp2/analysis/seed_level_results.json :: table[ee_vs_mm,mlp,17].auc (SD over seeds, ddof=1) :: sha256 53e1ec81343ac2c1
\newcommand{\ProbeLogOneMinusAucEeVsMmMlpOneseven}{\ensuremath{-}5.0211}  % experiments/FIGS/data/probe_ladder_v2_mlp2/analysis/seed_level_results.json :: levels.ee_vs_mm.mlp[3].mean :: sha256 53e1ec81343ac2c1 :: natural log of 1-AUC, lower is better
\newcommand{\ProbeLogOneMinusAucSdEeVsMmMlpOneseven}{0.0997}  % experiments/FIGS/data/probe_ladder_v2_mlp2/analysis/seed_level_results.json :: levels.ee_vs_mm.mlp[3].seed_sd :: sha256 53e1ec81343ac2c1
\newcommand{\ProbeOmaEeVsMmMlpOneseven}{\ensuremath{(6.6 \pm 0.7)\times10^{-3}}}  % experiments/FIGS/data/probe_ladder_v2_mlp2/analysis/seed_level_results.json :: table[ee_vs_mm,mlp,17].auc :: sha256 53e1ec81343ac2c1 :: 1-AUC, mean +- SD over 5 seeds
\newcommand{\ProbeBkgLeftEeVsMmMlpOneseven}{59.2}  % experiments/FIGS/data/probe_ladder_v2_mlp2/analysis/seed_level_results.json :: levels.ee_vs_mm.mlp[3].rejection_points['0.90'].mean_n_bkg_pass :: sha256 53e1ec81343ac2c1 :: background jets passing the cut, mean over seeds
\newcommand{\ProbeRejEeVsMmMlpOneseven}{81\,\ensuremath{\pm}\,12}  % experiments/FIGS/data/probe_ladder_v2_mlp2/analysis/seed_level_results.json :: table[ee_vs_mm,mlp,17].rejection_points['0.90'].rejection :: sha256 53e1ec81343ac2c1 :: at 90% signal efficiency, mean +- SD over seeds; 0 of 5 seeds at the cap
\newcommand{\ProbeOmaRatioEeVsMmMlpOnesevenOverFourthree}{1.93}  % experiments/FIGS/data/probe_ladder_v2_mlp2/analysis/seed_level_results.json :: table[ee_vs_mm,mlp,17/43].auc :: sha256 53e1ec81343ac2c1 :: ratio of the seed means of 1-AUC, coarser over finer
\newcommand{\ProbeAucRetainedTopologyLinearOneeighteight}{0.99413}  % experiments/FIGS/data/probe_ladder_v2_mlp2/analysis/seed_level_results.json :: levels.retained_topology.linear[0].mean_auc :: sha256 53e1ec81343ac2c1 :: 5 seeds, 0 saturated at AUC=1
\newcommand{\ProbeAucSdRetainedTopologyLinearOneeighteight}{0.00011}  % experiments/FIGS/data/probe_ladder_v2_mlp2/analysis/seed_level_results.json :: table[retained_topology,linear,188].auc (SD over seeds, ddof=1) :: sha256 53e1ec81343ac2c1
\newcommand{\ProbeLogOneMinusAucRetainedTopologyLinearOneeighteight}{\ensuremath{-}5.1388}  % experiments/FIGS/data/probe_ladder_v2_mlp2/analysis/seed_level_results.json :: levels.retained_topology.linear[0].mean :: sha256 53e1ec81343ac2c1 :: natural log of 1-AUC, lower is better
\newcommand{\ProbeLogOneMinusAucSdRetainedTopologyLinearOneeighteight}{0.0187}  % experiments/FIGS/data/probe_ladder_v2_mlp2/analysis/seed_level_results.json :: levels.retained_topology.linear[0].seed_sd :: sha256 53e1ec81343ac2c1
\newcommand{\ProbeOmaRetainedTopologyLinearOneeighteight}{\ensuremath{(5.87 \pm 0.11)\times10^{-3}}}  % experiments/FIGS/data/probe_ladder_v2_mlp2/analysis/seed_level_results.json :: table[retained_topology,linear,188].auc :: sha256 53e1ec81343ac2c1 :: 1-AUC, mean +- SD over 5 seeds
\newcommand{\ProbeBkgLeftRetainedTopologyLinearOneeighteight}{15.4}  % experiments/FIGS/data/probe_ladder_v2_mlp2/analysis/seed_level_results.json :: levels.retained_topology.linear[0].rejection_points['0.90'].mean_n_bkg_pass :: sha256 53e1ec81343ac2c1 :: background jets passing the cut, mean over seeds
\newcommand{\ProbeRejRetainedTopologyLinearOneeighteight}{71\,\ensuremath{\pm}\,7}  % experiments/FIGS/data/probe_ladder_v2_mlp2/analysis/seed_level_results.json :: table[retained_topology,linear,188].rejection_points['0.90'].rejection :: sha256 53e1ec81343ac2c1 :: at 90% signal efficiency, mean +- SD over seeds; 0 of 5 seeds at the cap
\newcommand{\ProbeAucRetainedTopologyLinearOnesixtwo}{0.99413}  % experiments/FIGS/data/probe_ladder_v2_mlp2/analysis/seed_level_results.json :: levels.retained_topology.linear[1].mean_auc :: sha256 53e1ec81343ac2c1 :: 5 seeds, 0 saturated at AUC=1
\newcommand{\ProbeAucSdRetainedTopologyLinearOnesixtwo}{0.00030}  % experiments/FIGS/data/probe_ladder_v2_mlp2/analysis/seed_level_results.json :: table[retained_topology,linear,162].auc (SD over seeds, ddof=1) :: sha256 53e1ec81343ac2c1
\newcommand{\ProbeLogOneMinusAucRetainedTopologyLinearOnesixtwo}{\ensuremath{-}5.1395}  % experiments/FIGS/data/probe_ladder_v2_mlp2/analysis/seed_level_results.json :: levels.retained_topology.linear[1].mean :: sha256 53e1ec81343ac2c1 :: natural log of 1-AUC, lower is better
\newcommand{\ProbeLogOneMinusAucSdRetainedTopologyLinearOnesixtwo}{0.0510}  % experiments/FIGS/data/probe_ladder_v2_mlp2/analysis/seed_level_results.json :: levels.retained_topology.linear[1].seed_sd :: sha256 53e1ec81343ac2c1
\newcommand{\ProbeOmaRetainedTopologyLinearOnesixtwo}{\ensuremath{(5.87 \pm 0.30)\times10^{-3}}}  % experiments/FIGS/data/probe_ladder_v2_mlp2/analysis/seed_level_results.json :: table[retained_topology,linear,162].auc :: sha256 53e1ec81343ac2c1 :: 1-AUC, mean +- SD over 5 seeds
\newcommand{\ProbeBkgLeftRetainedTopologyLinearOnesixtwo}{14.0}  % experiments/FIGS/data/probe_ladder_v2_mlp2/analysis/seed_level_results.json :: levels.retained_topology.linear[1].rejection_points['0.90'].mean_n_bkg_pass :: sha256 53e1ec81343ac2c1 :: background jets passing the cut, mean over seeds
\newcommand{\ProbeRejRetainedTopologyLinearOnesixtwo}{79\,\ensuremath{\pm}\,13}  % experiments/FIGS/data/probe_ladder_v2_mlp2/analysis/seed_level_results.json :: table[retained_topology,linear,162].rejection_points['0.90'].rejection :: sha256 53e1ec81343ac2c1 :: at 90% signal efficiency, mean +- SD over seeds; 0 of 5 seeds at the cap
\newcommand{\ProbeAucRetainedTopologyLinearFourthree}{0.99274}  % experiments/FIGS/data/probe_ladder_v2_mlp2/analysis/seed_level_results.json :: levels.retained_topology.linear[2].mean_auc :: sha256 53e1ec81343ac2c1 :: 5 seeds, 0 saturated at AUC=1
\newcommand{\ProbeAucSdRetainedTopologyLinearFourthree}{0.00012}  % experiments/FIGS/data/probe_ladder_v2_mlp2/analysis/seed_level_results.json :: table[retained_topology,linear,43].auc (SD over seeds, ddof=1) :: sha256 53e1ec81343ac2c1
\newcommand{\ProbeLogOneMinusAucRetainedTopologyLinearFourthree}{\ensuremath{-}4.9253}  % experiments/FIGS/data/probe_ladder_v2_mlp2/analysis/seed_level_results.json :: levels.retained_topology.linear[2].mean :: sha256 53e1ec81343ac2c1 :: natural log of 1-AUC, lower is better
\newcommand{\ProbeLogOneMinusAucSdRetainedTopologyLinearFourthree}{0.0169}  % experiments/FIGS/data/probe_ladder_v2_mlp2/analysis/seed_level_results.json :: levels.retained_topology.linear[2].seed_sd :: sha256 53e1ec81343ac2c1
\newcommand{\ProbeOmaRetainedTopologyLinearFourthree}{\ensuremath{(7.26 \pm 0.12)\times10^{-3}}}  % experiments/FIGS/data/probe_ladder_v2_mlp2/analysis/seed_level_results.json :: table[retained_topology,linear,43].auc :: sha256 53e1ec81343ac2c1 :: 1-AUC, mean +- SD over 5 seeds
\newcommand{\ProbeBkgLeftRetainedTopologyLinearFourthree}{16.4}  % experiments/FIGS/data/probe_ladder_v2_mlp2/analysis/seed_level_results.json :: levels.retained_topology.linear[2].rejection_points['0.90'].mean_n_bkg_pass :: sha256 53e1ec81343ac2c1 :: background jets passing the cut, mean over seeds
\newcommand{\ProbeRejRetainedTopologyLinearFourthree}{66\,\ensuremath{\pm}\,5}  % experiments/FIGS/data/probe_ladder_v2_mlp2/analysis/seed_level_results.json :: table[retained_topology,linear,43].rejection_points['0.90'].rejection :: sha256 53e1ec81343ac2c1 :: at 90% signal efficiency, mean +- SD over seeds; 0 of 5 seeds at the cap
\newcommand{\ProbeAucRetainedTopologyLinearOneseven}{0.99096}  % experiments/FIGS/data/probe_ladder_v2_mlp2/analysis/seed_level_results.json :: levels.retained_topology.linear[3].mean_auc :: sha256 53e1ec81343ac2c1 :: 5 seeds, 0 saturated at AUC=1
\newcommand{\ProbeAucSdRetainedTopologyLinearOneseven}{0.00092}  % experiments/FIGS/data/probe_ladder_v2_mlp2/analysis/seed_level_results.json :: table[retained_topology,linear,17].auc (SD over seeds, ddof=1) :: sha256 53e1ec81343ac2c1
\newcommand{\ProbeLogOneMinusAucRetainedTopologyLinearOneseven}{\ensuremath{-}4.7095}  % experiments/FIGS/data/probe_ladder_v2_mlp2/analysis/seed_level_results.json :: levels.retained_topology.linear[3].mean :: sha256 53e1ec81343ac2c1 :: natural log of 1-AUC, lower is better
\newcommand{\ProbeLogOneMinusAucSdRetainedTopologyLinearOneseven}{0.0979}  % experiments/FIGS/data/probe_ladder_v2_mlp2/analysis/seed_level_results.json :: levels.retained_topology.linear[3].seed_sd :: sha256 53e1ec81343ac2c1
\newcommand{\ProbeOmaRetainedTopologyLinearOneseven}{\ensuremath{(9.0 \pm 0.9)\times10^{-3}}}  % experiments/FIGS/data/probe_ladder_v2_mlp2/analysis/seed_level_results.json :: table[retained_topology,linear,17].auc :: sha256 53e1ec81343ac2c1 :: 1-AUC, mean +- SD over 5 seeds
\newcommand{\ProbeBkgLeftRetainedTopologyLinearOneseven}{21.8}  % experiments/FIGS/data/probe_ladder_v2_mlp2/analysis/seed_level_results.json :: levels.retained_topology.linear[3].rejection_points['0.90'].mean_n_bkg_pass :: sha256 53e1ec81343ac2c1 :: background jets passing the cut, mean over seeds
\newcommand{\ProbeRejRetainedTopologyLinearOneseven}{51\,\ensuremath{\pm}\,11}  % experiments/FIGS/data/probe_ladder_v2_mlp2/analysis/seed_level_results.json :: table[retained_topology,linear,17].rejection_points['0.90'].rejection :: sha256 53e1ec81343ac2c1 :: at 90% signal efficiency, mean +- SD over seeds; 0 of 5 seeds at the cap
\newcommand{\ProbeOmaRatioRetainedTopologyLinearOnesixtwoOverOneeighteight}{1.00}  % experiments/FIGS/data/probe_ladder_v2_mlp2/analysis/seed_level_results.json :: table[retained_topology,linear,162/188].auc :: sha256 53e1ec81343ac2c1 :: ratio of the seed means of 1-AUC, coarser over finer
\newcommand{\ProbeOmaRatioRetainedTopologyLinearFourthreeOverOneeighteight}{1.24}  % experiments/FIGS/data/probe_ladder_v2_mlp2/analysis/seed_level_results.json :: table[retained_topology,linear,43/188].auc :: sha256 53e1ec81343ac2c1 :: ratio of the seed means of 1-AUC, coarser over finer
\newcommand{\ProbeOmaRatioRetainedTopologyLinearOnesevenOverOneeighteight}{1.54}  % experiments/FIGS/data/probe_ladder_v2_mlp2/analysis/seed_level_results.json :: table[retained_topology,linear,17/188].auc :: sha256 53e1ec81343ac2c1 :: ratio of the seed means of 1-AUC, coarser over finer
\newcommand{\ProbeOmaRatioRetainedTopologyLinearFourthreeOverOnesixtwo}{1.24}  % experiments/FIGS/data/probe_ladder_v2_mlp2/analysis/seed_level_results.json :: table[retained_topology,linear,43/162].auc :: sha256 53e1ec81343ac2c1 :: ratio of the seed means of 1-AUC, coarser over finer
\newcommand{\ProbeOmaRatioRetainedTopologyLinearOnesevenOverOnesixtwo}{1.54}  % experiments/FIGS/data/probe_ladder_v2_mlp2/analysis/seed_level_results.json :: table[retained_topology,linear,17/162].auc :: sha256 53e1ec81343ac2c1 :: ratio of the seed means of 1-AUC, coarser over finer
\newcommand{\ProbeOmaRatioRetainedTopologyLinearOnesevenOverFourthree}{1.25}  % experiments/FIGS/data/probe_ladder_v2_mlp2/analysis/seed_level_results.json :: table[retained_topology,linear,17/43].auc :: sha256 53e1ec81343ac2c1 :: ratio of the seed means of 1-AUC, coarser over finer
\newcommand{\ProbeAucRetainedTopologyMlpOneeighteight}{0.99394}  % experiments/FIGS/data/probe_ladder_v2_mlp2/analysis/seed_level_results.json :: levels.retained_topology.mlp[0].mean_auc :: sha256 53e1ec81343ac2c1 :: 5 seeds, 0 saturated at AUC=1
\newcommand{\ProbeAucSdRetainedTopologyMlpOneeighteight}{0.00032}  % experiments/FIGS/data/probe_ladder_v2_mlp2/analysis/seed_level_results.json :: table[retained_topology,mlp,188].auc (SD over seeds, ddof=1) :: sha256 53e1ec81343ac2c1
\newcommand{\ProbeLogOneMinusAucRetainedTopologyMlpOneeighteight}{\ensuremath{-}5.1078}  % experiments/FIGS/data/probe_ladder_v2_mlp2/analysis/seed_level_results.json :: levels.retained_topology.mlp[0].mean :: sha256 53e1ec81343ac2c1 :: natural log of 1-AUC, lower is better
\newcommand{\ProbeLogOneMinusAucSdRetainedTopologyMlpOneeighteight}{0.0522}  % experiments/FIGS/data/probe_ladder_v2_mlp2/analysis/seed_level_results.json :: levels.retained_topology.mlp[0].seed_sd :: sha256 53e1ec81343ac2c1
\newcommand{\ProbeOmaRetainedTopologyMlpOneeighteight}{\ensuremath{(6.06 \pm 0.32)\times10^{-3}}}  % experiments/FIGS/data/probe_ladder_v2_mlp2/analysis/seed_level_results.json :: table[retained_topology,mlp,188].auc :: sha256 53e1ec81343ac2c1 :: 1-AUC, mean +- SD over 5 seeds
\newcommand{\ProbeBkgLeftRetainedTopologyMlpOneeighteight}{15.8}  % experiments/FIGS/data/probe_ladder_v2_mlp2/analysis/seed_level_results.json :: levels.retained_topology.mlp[0].rejection_points['0.90'].mean_n_bkg_pass :: sha256 53e1ec81343ac2c1 :: background jets passing the cut, mean over seeds
\newcommand{\ProbeRejRetainedTopologyMlpOneeighteight}{69\,\ensuremath{\pm}\,9}  % experiments/FIGS/data/probe_ladder_v2_mlp2/analysis/seed_level_results.json :: table[retained_topology,mlp,188].rejection_points['0.90'].rejection :: sha256 53e1ec81343ac2c1 :: at 90% signal efficiency, mean +- SD over seeds; 0 of 5 seeds at the cap
\newcommand{\ProbeAucRetainedTopologyMlpOnesixtwo}{0.99387}  % experiments/FIGS/data/probe_ladder_v2_mlp2/analysis/seed_level_results.json :: levels.retained_topology.mlp[1].mean_auc :: sha256 53e1ec81343ac2c1 :: 5 seeds, 0 saturated at AUC=1
\newcommand{\ProbeAucSdRetainedTopologyMlpOnesixtwo}{0.00048}  % experiments/FIGS/data/probe_ladder_v2_mlp2/analysis/seed_level_results.json :: table[retained_topology,mlp,162].auc (SD over seeds, ddof=1) :: sha256 53e1ec81343ac2c1
\newcommand{\ProbeLogOneMinusAucRetainedTopologyMlpOnesixtwo}{\ensuremath{-}5.0964}  % experiments/FIGS/data/probe_ladder_v2_mlp2/analysis/seed_level_results.json :: levels.retained_topology.mlp[1].mean :: sha256 53e1ec81343ac2c1 :: natural log of 1-AUC, lower is better
\newcommand{\ProbeLogOneMinusAucSdRetainedTopologyMlpOnesixtwo}{0.0773}  % experiments/FIGS/data/probe_ladder_v2_mlp2/analysis/seed_level_results.json :: levels.retained_topology.mlp[1].seed_sd :: sha256 53e1ec81343ac2c1
\newcommand{\ProbeOmaRetainedTopologyMlpOnesixtwo}{\ensuremath{(6.1 \pm 0.5)\times10^{-3}}}  % experiments/FIGS/data/probe_ladder_v2_mlp2/analysis/seed_level_results.json :: table[retained_topology,mlp,162].auc :: sha256 53e1ec81343ac2c1 :: 1-AUC, mean +- SD over 5 seeds
\newcommand{\ProbeBkgLeftRetainedTopologyMlpOnesixtwo}{16.6}  % experiments/FIGS/data/probe_ladder_v2_mlp2/analysis/seed_level_results.json :: levels.retained_topology.mlp[1].rejection_points['0.90'].mean_n_bkg_pass :: sha256 53e1ec81343ac2c1 :: background jets passing the cut, mean over seeds
\newcommand{\ProbeRejRetainedTopologyMlpOnesixtwo}{67\,\ensuremath{\pm}\,14}  % experiments/FIGS/data/probe_ladder_v2_mlp2/analysis/seed_level_results.json :: table[retained_topology,mlp,162].rejection_points['0.90'].rejection :: sha256 53e1ec81343ac2c1 :: at 90% signal efficiency, mean +- SD over seeds; 0 of 5 seeds at the cap
\newcommand{\ProbeAucRetainedTopologyMlpFourthree}{0.99284}  % experiments/FIGS/data/probe_ladder_v2_mlp2/analysis/seed_level_results.json :: levels.retained_topology.mlp[2].mean_auc :: sha256 53e1ec81343ac2c1 :: 5 seeds, 0 saturated at AUC=1
\newcommand{\ProbeAucSdRetainedTopologyMlpFourthree}{0.00040}  % experiments/FIGS/data/probe_ladder_v2_mlp2/analysis/seed_level_results.json :: table[retained_topology,mlp,43].auc (SD over seeds, ddof=1) :: sha256 53e1ec81343ac2c1
\newcommand{\ProbeLogOneMinusAucRetainedTopologyMlpFourthree}{\ensuremath{-}4.9400}  % experiments/FIGS/data/probe_ladder_v2_mlp2/analysis/seed_level_results.json :: levels.retained_topology.mlp[2].mean :: sha256 53e1ec81343ac2c1 :: natural log of 1-AUC, lower is better
\newcommand{\ProbeLogOneMinusAucSdRetainedTopologyMlpFourthree}{0.0570}  % experiments/FIGS/data/probe_ladder_v2_mlp2/analysis/seed_level_results.json :: levels.retained_topology.mlp[2].seed_sd :: sha256 53e1ec81343ac2c1
\newcommand{\ProbeOmaRetainedTopologyMlpFourthree}{\ensuremath{(7.2 \pm 0.4)\times10^{-3}}}  % experiments/FIGS/data/probe_ladder_v2_mlp2/analysis/seed_level_results.json :: table[retained_topology,mlp,43].auc :: sha256 53e1ec81343ac2c1 :: 1-AUC, mean +- SD over 5 seeds
\newcommand{\ProbeBkgLeftRetainedTopologyMlpFourthree}{17.8}  % experiments/FIGS/data/probe_ladder_v2_mlp2/analysis/seed_level_results.json :: levels.retained_topology.mlp[2].rejection_points['0.90'].mean_n_bkg_pass :: sha256 53e1ec81343ac2c1 :: background jets passing the cut, mean over seeds
\newcommand{\ProbeRejRetainedTopologyMlpFourthree}{61\,\ensuremath{\pm}\,7}  % experiments/FIGS/data/probe_ladder_v2_mlp2/analysis/seed_level_results.json :: table[retained_topology,mlp,43].rejection_points['0.90'].rejection :: sha256 53e1ec81343ac2c1 :: at 90% signal efficiency, mean +- SD over seeds; 0 of 5 seeds at the cap
\newcommand{\ProbeAucRetainedTopologyMlpOneseven}{0.99081}  % experiments/FIGS/data/probe_ladder_v2_mlp2/analysis/seed_level_results.json :: levels.retained_topology.mlp[3].mean_auc :: sha256 53e1ec81343ac2c1 :: 5 seeds, 0 saturated at AUC=1
\newcommand{\ProbeAucSdRetainedTopologyMlpOneseven}{0.00062}  % experiments/FIGS/data/probe_ladder_v2_mlp2/analysis/seed_level_results.json :: table[retained_topology,mlp,17].auc (SD over seeds, ddof=1) :: sha256 53e1ec81343ac2c1
\newcommand{\ProbeLogOneMinusAucRetainedTopologyMlpOneseven}{\ensuremath{-}4.6913}  % experiments/FIGS/data/probe_ladder_v2_mlp2/analysis/seed_level_results.json :: levels.retained_topology.mlp[3].mean :: sha256 53e1ec81343ac2c1 :: natural log of 1-AUC, lower is better
\newcommand{\ProbeLogOneMinusAucSdRetainedTopologyMlpOneseven}{0.0658}  % experiments/FIGS/data/probe_ladder_v2_mlp2/analysis/seed_level_results.json :: levels.retained_topology.mlp[3].seed_sd :: sha256 53e1ec81343ac2c1
\newcommand{\ProbeOmaRetainedTopologyMlpOneseven}{\ensuremath{(9.2 \pm 0.6)\times10^{-3}}}  % experiments/FIGS/data/probe_ladder_v2_mlp2/analysis/seed_level_results.json :: table[retained_topology,mlp,17].auc :: sha256 53e1ec81343ac2c1 :: 1-AUC, mean +- SD over 5 seeds
\newcommand{\ProbeBkgLeftRetainedTopologyMlpOneseven}{21.2}  % experiments/FIGS/data/probe_ladder_v2_mlp2/analysis/seed_level_results.json :: levels.retained_topology.mlp[3].rejection_points['0.90'].mean_n_bkg_pass :: sha256 53e1ec81343ac2c1 :: background jets passing the cut, mean over seeds
\newcommand{\ProbeRejRetainedTopologyMlpOneseven}{51\,\ensuremath{\pm}\,4}  % experiments/FIGS/data/probe_ladder_v2_mlp2/analysis/seed_level_results.json :: table[retained_topology,mlp,17].rejection_points['0.90'].rejection :: sha256 53e1ec81343ac2c1 :: at 90% signal efficiency, mean +- SD over seeds; 0 of 5 seeds at the cap
\newcommand{\ProbeOmaRatioRetainedTopologyMlpOnesixtwoOverOneeighteight}{1.01}  % experiments/FIGS/data/probe_ladder_v2_mlp2/analysis/seed_level_results.json :: table[retained_topology,mlp,162/188].auc :: sha256 53e1ec81343ac2c1 :: ratio of the seed means of 1-AUC, coarser over finer
\newcommand{\ProbeOmaRatioRetainedTopologyMlpFourthreeOverOneeighteight}{1.18}  % experiments/FIGS/data/probe_ladder_v2_mlp2/analysis/seed_level_results.json :: table[retained_topology,mlp,43/188].auc :: sha256 53e1ec81343ac2c1 :: ratio of the seed means of 1-AUC, coarser over finer
\newcommand{\ProbeOmaRatioRetainedTopologyMlpOnesevenOverOneeighteight}{1.52}  % experiments/FIGS/data/probe_ladder_v2_mlp2/analysis/seed_level_results.json :: table[retained_topology,mlp,17/188].auc :: sha256 53e1ec81343ac2c1 :: ratio of the seed means of 1-AUC, coarser over finer
\newcommand{\ProbeOmaRatioRetainedTopologyMlpFourthreeOverOnesixtwo}{1.17}  % experiments/FIGS/data/probe_ladder_v2_mlp2/analysis/seed_level_results.json :: table[retained_topology,mlp,43/162].auc :: sha256 53e1ec81343ac2c1 :: ratio of the seed means of 1-AUC, coarser over finer
\newcommand{\ProbeOmaRatioRetainedTopologyMlpOnesevenOverOnesixtwo}{1.50}  % experiments/FIGS/data/probe_ladder_v2_mlp2/analysis/seed_level_results.json :: table[retained_topology,mlp,17/162].auc :: sha256 53e1ec81343ac2c1 :: ratio of the seed means of 1-AUC, coarser over finer
\newcommand{\ProbeOmaRatioRetainedTopologyMlpOnesevenOverFourthree}{1.28}  % experiments/FIGS/data/probe_ladder_v2_mlp2/analysis/seed_level_results.json :: table[retained_topology,mlp,17/43].auc :: sha256 53e1ec81343ac2c1 :: ratio of the seed means of 1-AUC, coarser over finer
\newcommand{\ProbeAucVisibleContentLinearOneeighteight}{0.99917}  % experiments/FIGS/data/probe_ladder_v2_mlp2/analysis/seed_level_results.json :: levels.visible_content.linear[0].mean_auc :: sha256 53e1ec81343ac2c1 :: 5 seeds, 0 saturated at AUC=1
\newcommand{\ProbeAucSdVisibleContentLinearOneeighteight}{0.00003}  % experiments/FIGS/data/probe_ladder_v2_mlp2/analysis/seed_level_results.json :: table[visible_content,linear,188].auc (SD over seeds, ddof=1) :: sha256 53e1ec81343ac2c1
\newcommand{\ProbeLogOneMinusAucVisibleContentLinearOneeighteight}{\ensuremath{-}7.0999}  % experiments/FIGS/data/probe_ladder_v2_mlp2/analysis/seed_level_results.json :: levels.visible_content.linear[0].mean :: sha256 53e1ec81343ac2c1 :: natural log of 1-AUC, lower is better
\newcommand{\ProbeLogOneMinusAucSdVisibleContentLinearOneeighteight}{0.0358}  % experiments/FIGS/data/probe_ladder_v2_mlp2/analysis/seed_level_results.json :: levels.visible_content.linear[0].seed_sd :: sha256 53e1ec81343ac2c1
\newcommand{\ProbeOmaVisibleContentLinearOneeighteight}{\ensuremath{(8.26 \pm 0.30)\times10^{-4}}}  % experiments/FIGS/data/probe_ladder_v2_mlp2/analysis/seed_level_results.json :: table[visible_content,linear,188].auc :: sha256 53e1ec81343ac2c1 :: 1-AUC, mean +- SD over 5 seeds
\newcommand{\ProbeBkgLeftVisibleContentLinearOneeighteight}{0.4}  % experiments/FIGS/data/probe_ladder_v2_mlp2/analysis/seed_level_results.json :: levels.visible_content.linear[0].rejection_points['0.90'].mean_n_bkg_pass :: sha256 53e1ec81343ac2c1 :: background jets passing the cut, mean over seeds
\newcommand{\ProbeRejVisibleContentLinearOneeighteight}{\ensuremath{>}538}  % experiments/FIGS/data/probe_ladder_v2_mlp2/analysis/seed_level_results.json :: table[visible_content,linear,188].rejection_points['0.90'].rejection :: sha256 53e1ec81343ac2c1 :: at 90% signal efficiency, mean +- SD over seeds; 5 of 5 seeds at the cap
\newcommand{\ProbeAucVisibleContentLinearOnesixtwo}{0.99921}  % experiments/FIGS/data/probe_ladder_v2_mlp2/analysis/seed_level_results.json :: levels.visible_content.linear[1].mean_auc :: sha256 53e1ec81343ac2c1 :: 5 seeds, 0 saturated at AUC=1
\newcommand{\ProbeAucSdVisibleContentLinearOnesixtwo}{0.00004}  % experiments/FIGS/data/probe_ladder_v2_mlp2/analysis/seed_level_results.json :: table[visible_content,linear,162].auc (SD over seeds, ddof=1) :: sha256 53e1ec81343ac2c1
\newcommand{\ProbeLogOneMinusAucVisibleContentLinearOnesixtwo}{\ensuremath{-}7.1480}  % experiments/FIGS/data/probe_ladder_v2_mlp2/analysis/seed_level_results.json :: levels.visible_content.linear[1].mean :: sha256 53e1ec81343ac2c1 :: natural log of 1-AUC, lower is better
\newcommand{\ProbeLogOneMinusAucSdVisibleContentLinearOnesixtwo}{0.0449}  % experiments/FIGS/data/probe_ladder_v2_mlp2/analysis/seed_level_results.json :: levels.visible_content.linear[1].seed_sd :: sha256 53e1ec81343ac2c1
\newcommand{\ProbeOmaVisibleContentLinearOnesixtwo}{\ensuremath{(7.87 \pm 0.35)\times10^{-4}}}  % experiments/FIGS/data/probe_ladder_v2_mlp2/analysis/seed_level_results.json :: table[visible_content,linear,162].auc :: sha256 53e1ec81343ac2c1 :: 1-AUC, mean +- SD over 5 seeds
\newcommand{\ProbeBkgLeftVisibleContentLinearOnesixtwo}{0.6}  % experiments/FIGS/data/probe_ladder_v2_mlp2/analysis/seed_level_results.json :: levels.visible_content.linear[1].rejection_points['0.90'].mean_n_bkg_pass :: sha256 53e1ec81343ac2c1 :: background jets passing the cut, mean over seeds
\newcommand{\ProbeRejVisibleContentLinearOnesixtwo}{\ensuremath{>}538}  % experiments/FIGS/data/probe_ladder_v2_mlp2/analysis/seed_level_results.json :: table[visible_content,linear,162].rejection_points['0.90'].rejection :: sha256 53e1ec81343ac2c1 :: at 90% signal efficiency, mean +- SD over seeds; 5 of 5 seeds at the cap
\newcommand{\ProbeAucVisibleContentLinearFourthree}{0.99918}  % experiments/FIGS/data/probe_ladder_v2_mlp2/analysis/seed_level_results.json :: levels.visible_content.linear[2].mean_auc :: sha256 53e1ec81343ac2c1 :: 5 seeds, 0 saturated at AUC=1
\newcommand{\ProbeAucSdVisibleContentLinearFourthree}{0.00009}  % experiments/FIGS/data/probe_ladder_v2_mlp2/analysis/seed_level_results.json :: table[visible_content,linear,43].auc (SD over seeds, ddof=1) :: sha256 53e1ec81343ac2c1
\newcommand{\ProbeLogOneMinusAucVisibleContentLinearFourthree}{\ensuremath{-}7.1124}  % experiments/FIGS/data/probe_ladder_v2_mlp2/analysis/seed_level_results.json :: levels.visible_content.linear[2].mean :: sha256 53e1ec81343ac2c1 :: natural log of 1-AUC, lower is better
\newcommand{\ProbeLogOneMinusAucSdVisibleContentLinearFourthree}{0.1137}  % experiments/FIGS/data/probe_ladder_v2_mlp2/analysis/seed_level_results.json :: levels.visible_content.linear[2].seed_sd :: sha256 53e1ec81343ac2c1
\newcommand{\ProbeOmaVisibleContentLinearFourthree}{\ensuremath{(8.2 \pm 1.0)\times10^{-4}}}  % experiments/FIGS/data/probe_ladder_v2_mlp2/analysis/seed_level_results.json :: table[visible_content,linear,43].auc :: sha256 53e1ec81343ac2c1 :: 1-AUC, mean +- SD over 5 seeds
\newcommand{\ProbeBkgLeftVisibleContentLinearFourthree}{0.4}  % experiments/FIGS/data/probe_ladder_v2_mlp2/analysis/seed_level_results.json :: levels.visible_content.linear[2].rejection_points['0.90'].mean_n_bkg_pass :: sha256 53e1ec81343ac2c1 :: background jets passing the cut, mean over seeds
\newcommand{\ProbeRejVisibleContentLinearFourthree}{\ensuremath{>}538}  % experiments/FIGS/data/probe_ladder_v2_mlp2/analysis/seed_level_results.json :: table[visible_content,linear,43].rejection_points['0.90'].rejection :: sha256 53e1ec81343ac2c1 :: at 90% signal efficiency, mean +- SD over seeds; 5 of 5 seeds at the cap
\newcommand{\ProbeAucVisibleContentLinearOneseven}{0.99768}  % experiments/FIGS/data/probe_ladder_v2_mlp2/analysis/seed_level_results.json :: levels.visible_content.linear[3].mean_auc :: sha256 53e1ec81343ac2c1 :: 5 seeds, 0 saturated at AUC=1
\newcommand{\ProbeAucSdVisibleContentLinearOneseven}{0.00020}  % experiments/FIGS/data/probe_ladder_v2_mlp2/analysis/seed_level_results.json :: table[visible_content,linear,17].auc (SD over seeds, ddof=1) :: sha256 53e1ec81343ac2c1
\newcommand{\ProbeLogOneMinusAucVisibleContentLinearOneseven}{\ensuremath{-}6.0684}  % experiments/FIGS/data/probe_ladder_v2_mlp2/analysis/seed_level_results.json :: levels.visible_content.linear[3].mean :: sha256 53e1ec81343ac2c1 :: natural log of 1-AUC, lower is better
\newcommand{\ProbeLogOneMinusAucSdVisibleContentLinearOneseven}{0.0889}  % experiments/FIGS/data/probe_ladder_v2_mlp2/analysis/seed_level_results.json :: levels.visible_content.linear[3].seed_sd :: sha256 53e1ec81343ac2c1
\newcommand{\ProbeOmaVisibleContentLinearOneseven}{\ensuremath{(2.32 \pm 0.20)\times10^{-3}}}  % experiments/FIGS/data/probe_ladder_v2_mlp2/analysis/seed_level_results.json :: table[visible_content,linear,17].auc :: sha256 53e1ec81343ac2c1 :: 1-AUC, mean +- SD over 5 seeds
\newcommand{\ProbeBkgLeftVisibleContentLinearOneseven}{3.4}  % experiments/FIGS/data/probe_ladder_v2_mlp2/analysis/seed_level_results.json :: levels.visible_content.linear[3].rejection_points['0.90'].mean_n_bkg_pass :: sha256 53e1ec81343ac2c1 :: background jets passing the cut, mean over seeds
\newcommand{\ProbeRejVisibleContentLinearOneseven}{900\,\ensuremath{\pm}\,400}  % experiments/FIGS/data/probe_ladder_v2_mlp2/analysis/seed_level_results.json :: table[visible_content,linear,17].rejection_points['0.90'].rejection :: sha256 53e1ec81343ac2c1 :: at 90% signal efficiency, mean +- SD over seeds; 0 of 5 seeds at the cap
\newcommand{\ProbeOmaRatioVisibleContentLinearOnesixtwoOverOneeighteight}{0.95}  % experiments/FIGS/data/probe_ladder_v2_mlp2/analysis/seed_level_results.json :: table[visible_content,linear,162/188].auc :: sha256 53e1ec81343ac2c1 :: ratio of the seed means of 1-AUC, coarser over finer
\newcommand{\ProbeOmaRatioVisibleContentLinearFourthreeOverOneeighteight}{0.99}  % experiments/FIGS/data/probe_ladder_v2_mlp2/analysis/seed_level_results.json :: table[visible_content,linear,43/188].auc :: sha256 53e1ec81343ac2c1 :: ratio of the seed means of 1-AUC, coarser over finer
\newcommand{\ProbeOmaRatioVisibleContentLinearOnesevenOverOneeighteight}{2.8}  % experiments/FIGS/data/probe_ladder_v2_mlp2/analysis/seed_level_results.json :: table[visible_content,linear,17/188].auc :: sha256 53e1ec81343ac2c1 :: ratio of the seed means of 1-AUC, coarser over finer
\newcommand{\ProbeOmaRatioVisibleContentLinearFourthreeOverOnesixtwo}{1.04}  % experiments/FIGS/data/probe_ladder_v2_mlp2/analysis/seed_level_results.json :: table[visible_content,linear,43/162].auc :: sha256 53e1ec81343ac2c1 :: ratio of the seed means of 1-AUC, coarser over finer
\newcommand{\ProbeOmaRatioVisibleContentLinearOnesevenOverOnesixtwo}{3.0}  % experiments/FIGS/data/probe_ladder_v2_mlp2/analysis/seed_level_results.json :: table[visible_content,linear,17/162].auc :: sha256 53e1ec81343ac2c1 :: ratio of the seed means of 1-AUC, coarser over finer
\newcommand{\ProbeOmaRatioVisibleContentLinearOnesevenOverFourthree}{2.8}  % experiments/FIGS/data/probe_ladder_v2_mlp2/analysis/seed_level_results.json :: table[visible_content,linear,17/43].auc :: sha256 53e1ec81343ac2c1 :: ratio of the seed means of 1-AUC, coarser over finer
\newcommand{\ProbeAucVisibleContentMlpOneeighteight}{0.99911}  % experiments/FIGS/data/probe_ladder_v2_mlp2/analysis/seed_level_results.json :: levels.visible_content.mlp[0].mean_auc :: sha256 53e1ec81343ac2c1 :: 5 seeds, 0 saturated at AUC=1
\newcommand{\ProbeAucSdVisibleContentMlpOneeighteight}{0.00005}  % experiments/FIGS/data/probe_ladder_v2_mlp2/analysis/seed_level_results.json :: table[visible_content,mlp,188].auc (SD over seeds, ddof=1) :: sha256 53e1ec81343ac2c1
\newcommand{\ProbeLogOneMinusAucVisibleContentMlpOneeighteight}{\ensuremath{-}7.0202}  % experiments/FIGS/data/probe_ladder_v2_mlp2/analysis/seed_level_results.json :: levels.visible_content.mlp[0].mean :: sha256 53e1ec81343ac2c1 :: natural log of 1-AUC, lower is better
\newcommand{\ProbeLogOneMinusAucSdVisibleContentMlpOneeighteight}{0.0570}  % experiments/FIGS/data/probe_ladder_v2_mlp2/analysis/seed_level_results.json :: levels.visible_content.mlp[0].seed_sd :: sha256 53e1ec81343ac2c1
\newcommand{\ProbeOmaVisibleContentMlpOneeighteight}{\ensuremath{(8.9 \pm 0.5)\times10^{-4}}}  % experiments/FIGS/data/probe_ladder_v2_mlp2/analysis/seed_level_results.json :: table[visible_content,mlp,188].auc :: sha256 53e1ec81343ac2c1 :: 1-AUC, mean +- SD over 5 seeds
\newcommand{\ProbeBkgLeftVisibleContentMlpOneeighteight}{0.4}  % experiments/FIGS/data/probe_ladder_v2_mlp2/analysis/seed_level_results.json :: levels.visible_content.mlp[0].rejection_points['0.90'].mean_n_bkg_pass :: sha256 53e1ec81343ac2c1 :: background jets passing the cut, mean over seeds
\newcommand{\ProbeRejVisibleContentMlpOneeighteight}{\ensuremath{>}538}  % experiments/FIGS/data/probe_ladder_v2_mlp2/analysis/seed_level_results.json :: table[visible_content,mlp,188].rejection_points['0.90'].rejection :: sha256 53e1ec81343ac2c1 :: at 90% signal efficiency, mean +- SD over seeds; 5 of 5 seeds at the cap
\newcommand{\ProbeAucVisibleContentMlpOnesixtwo}{0.99918}  % experiments/FIGS/data/probe_ladder_v2_mlp2/analysis/seed_level_results.json :: levels.visible_content.mlp[1].mean_auc :: sha256 53e1ec81343ac2c1 :: 5 seeds, 0 saturated at AUC=1
\newcommand{\ProbeAucSdVisibleContentMlpOnesixtwo}{0.00003}  % experiments/FIGS/data/probe_ladder_v2_mlp2/analysis/seed_level_results.json :: table[visible_content,mlp,162].auc (SD over seeds, ddof=1) :: sha256 53e1ec81343ac2c1
\newcommand{\ProbeLogOneMinusAucVisibleContentMlpOnesixtwo}{\ensuremath{-}7.1010}  % experiments/FIGS/data/probe_ladder_v2_mlp2/analysis/seed_level_results.json :: levels.visible_content.mlp[1].mean :: sha256 53e1ec81343ac2c1 :: natural log of 1-AUC, lower is better
\newcommand{\ProbeLogOneMinusAucSdVisibleContentMlpOnesixtwo}{0.0371}  % experiments/FIGS/data/probe_ladder_v2_mlp2/analysis/seed_level_results.json :: levels.visible_content.mlp[1].seed_sd :: sha256 53e1ec81343ac2c1
\newcommand{\ProbeOmaVisibleContentMlpOnesixtwo}{\ensuremath{(8.25 \pm 0.30)\times10^{-4}}}  % experiments/FIGS/data/probe_ladder_v2_mlp2/analysis/seed_level_results.json :: table[visible_content,mlp,162].auc :: sha256 53e1ec81343ac2c1 :: 1-AUC, mean +- SD over 5 seeds
\newcommand{\ProbeBkgLeftVisibleContentMlpOnesixtwo}{0.2}  % experiments/FIGS/data/probe_ladder_v2_mlp2/analysis/seed_level_results.json :: levels.visible_content.mlp[1].rejection_points['0.90'].mean_n_bkg_pass :: sha256 53e1ec81343ac2c1 :: background jets passing the cut, mean over seeds
\newcommand{\ProbeRejVisibleContentMlpOnesixtwo}{\ensuremath{>}538}  % experiments/FIGS/data/probe_ladder_v2_mlp2/analysis/seed_level_results.json :: table[visible_content,mlp,162].rejection_points['0.90'].rejection :: sha256 53e1ec81343ac2c1 :: at 90% signal efficiency, mean +- SD over seeds; 5 of 5 seeds at the cap
\newcommand{\ProbeAucVisibleContentMlpFourthree}{0.99917}  % experiments/FIGS/data/probe_ladder_v2_mlp2/analysis/seed_level_results.json :: levels.visible_content.mlp[2].mean_auc :: sha256 53e1ec81343ac2c1 :: 5 seeds, 0 saturated at AUC=1
\newcommand{\ProbeAucSdVisibleContentMlpFourthree}{0.00006}  % experiments/FIGS/data/probe_ladder_v2_mlp2/analysis/seed_level_results.json :: table[visible_content,mlp,43].auc (SD over seeds, ddof=1) :: sha256 53e1ec81343ac2c1
\newcommand{\ProbeLogOneMinusAucVisibleContentMlpFourthree}{\ensuremath{-}7.0957}  % experiments/FIGS/data/probe_ladder_v2_mlp2/analysis/seed_level_results.json :: levels.visible_content.mlp[2].mean :: sha256 53e1ec81343ac2c1 :: natural log of 1-AUC, lower is better
\newcommand{\ProbeLogOneMinusAucSdVisibleContentMlpFourthree}{0.0773}  % experiments/FIGS/data/probe_ladder_v2_mlp2/analysis/seed_level_results.json :: levels.visible_content.mlp[2].seed_sd :: sha256 53e1ec81343ac2c1
\newcommand{\ProbeOmaVisibleContentMlpFourthree}{\ensuremath{(8.3 \pm 0.6)\times10^{-4}}}  % experiments/FIGS/data/probe_ladder_v2_mlp2/analysis/seed_level_results.json :: table[visible_content,mlp,43].auc :: sha256 53e1ec81343ac2c1 :: 1-AUC, mean +- SD over 5 seeds
\newcommand{\ProbeBkgLeftVisibleContentMlpFourthree}{0.4}  % experiments/FIGS/data/probe_ladder_v2_mlp2/analysis/seed_level_results.json :: levels.visible_content.mlp[2].rejection_points['0.90'].mean_n_bkg_pass :: sha256 53e1ec81343ac2c1 :: background jets passing the cut, mean over seeds
\newcommand{\ProbeRejVisibleContentMlpFourthree}{\ensuremath{>}538}  % experiments/FIGS/data/probe_ladder_v2_mlp2/analysis/seed_level_results.json :: table[visible_content,mlp,43].rejection_points['0.90'].rejection :: sha256 53e1ec81343ac2c1 :: at 90% signal efficiency, mean +- SD over seeds; 5 of 5 seeds at the cap
\newcommand{\ProbeAucVisibleContentMlpOneseven}{0.99793}  % experiments/FIGS/data/probe_ladder_v2_mlp2/analysis/seed_level_results.json :: levels.visible_content.mlp[3].mean_auc :: sha256 53e1ec81343ac2c1 :: 5 seeds, 0 saturated at AUC=1
\newcommand{\ProbeAucSdVisibleContentMlpOneseven}{0.00020}  % experiments/FIGS/data/probe_ladder_v2_mlp2/analysis/seed_level_results.json :: table[visible_content,mlp,17].auc (SD over seeds, ddof=1) :: sha256 53e1ec81343ac2c1
\newcommand{\ProbeLogOneMinusAucVisibleContentMlpOneseven}{\ensuremath{-}6.1839}  % experiments/FIGS/data/probe_ladder_v2_mlp2/analysis/seed_level_results.json :: levels.visible_content.mlp[3].mean :: sha256 53e1ec81343ac2c1 :: natural log of 1-AUC, lower is better
\newcommand{\ProbeLogOneMinusAucSdVisibleContentMlpOneseven}{0.1014}  % experiments/FIGS/data/probe_ladder_v2_mlp2/analysis/seed_level_results.json :: levels.visible_content.mlp[3].seed_sd :: sha256 53e1ec81343ac2c1
\newcommand{\ProbeOmaVisibleContentMlpOneseven}{\ensuremath{(2.07 \pm 0.20)\times10^{-3}}}  % experiments/FIGS/data/probe_ladder_v2_mlp2/analysis/seed_level_results.json :: table[visible_content,mlp,17].auc :: sha256 53e1ec81343ac2c1 :: 1-AUC, mean +- SD over 5 seeds
\newcommand{\ProbeBkgLeftVisibleContentMlpOneseven}{3.4}  % experiments/FIGS/data/probe_ladder_v2_mlp2/analysis/seed_level_results.json :: levels.visible_content.mlp[3].rejection_points['0.90'].mean_n_bkg_pass :: sha256 53e1ec81343ac2c1 :: background jets passing the cut, mean over seeds
\newcommand{\ProbeRejVisibleContentMlpOneseven}{\ensuremath{\geq}638\ensuremath{^{\ast}}}  % experiments/FIGS/data/probe_ladder_v2_mlp2/analysis/seed_level_results.json :: table[visible_content,mlp,17].rejection_points['0.90'].rejection :: sha256 53e1ec81343ac2c1 :: at 90% signal efficiency, mean +- SD over seeds; 1 of 5 seeds at the cap
\newcommand{\ProbeOmaRatioVisibleContentMlpOnesixtwoOverOneeighteight}{0.92}  % experiments/FIGS/data/probe_ladder_v2_mlp2/analysis/seed_level_results.json :: table[visible_content,mlp,162/188].auc :: sha256 53e1ec81343ac2c1 :: ratio of the seed means of 1-AUC, coarser over finer
\newcommand{\ProbeOmaRatioVisibleContentMlpFourthreeOverOneeighteight}{0.93}  % experiments/FIGS/data/probe_ladder_v2_mlp2/analysis/seed_level_results.json :: table[visible_content,mlp,43/188].auc :: sha256 53e1ec81343ac2c1 :: ratio of the seed means of 1-AUC, coarser over finer
\newcommand{\ProbeOmaRatioVisibleContentMlpOnesevenOverOneeighteight}{2.3}  % experiments/FIGS/data/probe_ladder_v2_mlp2/analysis/seed_level_results.json :: table[visible_content,mlp,17/188].auc :: sha256 53e1ec81343ac2c1 :: ratio of the seed means of 1-AUC, coarser over finer
\newcommand{\ProbeOmaRatioVisibleContentMlpFourthreeOverOnesixtwo}{1.01}  % experiments/FIGS/data/probe_ladder_v2_mlp2/analysis/seed_level_results.json :: table[visible_content,mlp,43/162].auc :: sha256 53e1ec81343ac2c1 :: ratio of the seed means of 1-AUC, coarser over finer
\newcommand{\ProbeOmaRatioVisibleContentMlpOnesevenOverOnesixtwo}{2.5}  % experiments/FIGS/data/probe_ladder_v2_mlp2/analysis/seed_level_results.json :: table[visible_content,mlp,17/162].auc :: sha256 53e1ec81343ac2c1 :: ratio of the seed means of 1-AUC, coarser over finer
\newcommand{\ProbeOmaRatioVisibleContentMlpOnesevenOverFourthree}{2.5}  % experiments/FIGS/data/probe_ladder_v2_mlp2/analysis/seed_level_results.json :: table[visible_content,mlp,17/43].auc :: sha256 53e1ec81343ac2c1 :: ratio of the seed means of 1-AUC, coarser over finer
\newcommand{\VocabSizeLoneeighteight}{188}  % configs/labelmaps/rung_label_maps.v1.csv :: column L188 (distinct group count) :: sha256 412f9aa3b58db412
\newcommand{\VocabSizeLonesixtwo}{162}  % configs/labelmaps/rung_label_maps.v1.csv :: column L162 (distinct group count) :: sha256 412f9aa3b58db412
\newcommand{\VocabSizeRsixthreeQone}{64}  % configs/labelmaps/rung_label_maps.v1.csv :: column R63_Q1 (distinct group count) :: sha256 412f9aa3b58db412
\newcommand{\VocabSizeRfourtwoQone}{43}  % configs/labelmaps/rung_label_maps.v1.csv :: column R42_Q1 (distinct group count) :: sha256 412f9aa3b58db412
\newcommand{\VocabSizeRtwonineQone}{30}  % configs/labelmaps/rung_label_maps.v1.csv :: column R29_Q1 (distinct group count) :: sha256 412f9aa3b58db412
\newcommand{\VocabSizeRonesixQone}{17}  % configs/labelmaps/rung_label_maps.v1.csv :: column R16_Q1 (distinct group count) :: sha256 412f9aa3b58db412
\newcommand{\VocabSizeRthreeVIS}{4}  % configs/labelmaps/rung_label_maps.v1.csv :: column R3_VIS (distinct group count) :: sha256 412f9aa3b58db412
\newcommand{\VocabSizeRoneQone}{2}  % configs/labelmaps/rung_label_maps.v1.csv :: column R1_Q1 (distinct group count) :: sha256 412f9aa3b58db412
\newcommand{\ProbeSplitTrain}{60\%}  % experiments/EVAL/probe.py :: SPLIT_FRACTIONS :: sha256 485db0966e660588 :: share of the probe sample; the same jets for every model
\newcommand{\ProbeSplitVal}{20\%}  % experiments/EVAL/probe.py :: SPLIT_FRACTIONS :: sha256 485db0966e660588 :: share of the probe sample; the same jets for every model
\newcommand{\ProbeSplitTest}{20\%}  % experiments/EVAL/probe.py :: SPLIT_FRACTIONS :: sha256 485db0966e660588 :: share of the probe sample; the same jets for every model
\newcommand{\ProbeMlpHidden}{256}  % experiments/EVAL/probe.py :: _fit_mlp :: sha256 485db0966e660588 :: hidden width of the MLP probe
\newcommand{\ProbeMlpInits}{3}  % experiments/EVAL/probe.py :: MLP_SEEDS :: sha256 485db0966e660588 :: MLP initialisations whose scores are averaged per cell
\newcommand{\ProbeMlpLrDrop}{10}  % experiments/EVAL/probe.py :: MLP_SCHEDULE['lr_factor'] :: sha256 485db0966e660588 :: factor the MLP's learning rate drops by on a plateau
\newcommand{\ProbeMlpPlateaus}{three}  % experiments/EVAL/probe.py :: MLP_SCHEDULE: stop_patience / (lr_patience + 1) :: sha256 485db0966e660588 :: patience windows without gain after which the MLP stops
\newcommand{\ProbeCMin}{\ensuremath{10^{-2}}}  % experiments/EVAL/probe.py :: C_GRID :: sha256 485db0966e660588 :: end of the logistic-regression inverse-regularisation grid
\newcommand{\ProbeCMax}{\ensuremath{10^{2}}}  % experiments/EVAL/probe.py :: C_GRID :: sha256 485db0966e660588 :: end of the logistic-regression inverse-regularisation grid
\newcommand{\ProbeNBkgTestBvcFourprong}{3{,}158}  % experiments/FIGS/data/probe_ladder_v2_mlp2/s1.json :: tasks.bvc_4prong.n_background_test :: sha256 f7e1de0e43be89b8 :: test-sample background jets
\newcommand{\ProbeNBkgTestBvcQcd}{5{,}371}  % experiments/FIGS/data/probe_ladder_v2_mlp2/s1.json :: tasks.bvc_qcd.n_background_test :: sha256 f7e1de0e43be89b8 :: test-sample background jets
\newcommand{\ProbeNBkgTestBvcResonant}{11{,}876}  % experiments/FIGS/data/probe_ladder_v2_mlp2/s1.json :: tasks.bvc_resonant.n_background_test :: sha256 f7e1de0e43be89b8 :: test-sample background jets
\newcommand{\ProbeNBkgTestEeVsMm}{4{,}708}  % experiments/FIGS/data/probe_ladder_v2_mlp2/s1.json :: tasks.ee_vs_mm.n_background_test :: sha256 f7e1de0e43be89b8 :: test-sample background jets
\newcommand{\ProbeNBkgTestRetainedTopology}{1{,}078}  % experiments/FIGS/data/probe_ladder_v2_mlp2/s1.json :: tasks.retained_topology.n_background_test :: sha256 f7e1de0e43be89b8 :: test-sample background jets
\newcommand{\ProbeNBkgTestVisibleContent}{2{,}554}  % experiments/FIGS/data/probe_ladder_v2_mlp2/s1.json :: tasks.visible_content.n_background_test :: sha256 f7e1de0e43be89b8 :: test-sample background jets
\newcommand{\UseDiesAtSophonEqfour}{43}  % configs/labelmaps/usecase_survival.v1.json :: sophon_eq4.dies_at :: sha256 122275213c349593 :: first vocabulary size at which it is not constructible
\newcommand{\UseLastAliveSophonEqfour}{64}  % configs/labelmaps/usecase_survival.v1.json :: sophon_eq4.last_rung_alive :: sha256 122275213c349593
\newcommand{\UseDiesAtSophonEqsixAone}{64}  % configs/labelmaps/usecase_survival.v1.json :: sophon_eq6_a1.dies_at :: sha256 122275213c349593 :: first vocabulary size at which it is not constructible
\newcommand{\UseLastAliveSophonEqsixAone}{162}  % configs/labelmaps/usecase_survival.v1.json :: sophon_eq6_a1.last_rung_alive :: sha256 122275213c349593
\newcommand{\UseDiesAtSophonEqsixAthree}{64}  % configs/labelmaps/usecase_survival.v1.json :: sophon_eq6_a3.dies_at :: sha256 122275213c349593 :: first vocabulary size at which it is not constructible
\newcommand{\UseLastAliveSophonEqsixAthree}{162}  % configs/labelmaps/usecase_survival.v1.json :: sophon_eq6_a3.last_rung_alive :: sha256 122275213c349593
\newcommand{\UseDiesAtSophonEqsevenNine}{64}  % configs/labelmaps/usecase_survival.v1.json :: sophon_eq7_9.dies_at :: sha256 122275213c349593 :: first vocabulary size at which it is not constructible
\newcommand{\UseLastAliveSophonEqsevenNine}{162}  % configs/labelmaps/usecase_survival.v1.json :: sophon_eq7_9.last_rung_alive :: sha256 122275213c349593
\newcommand{\UseDiesAtVcbEqone}{64}  % configs/labelmaps/usecase_survival.v1.json :: vcb_eq1.dies_at :: sha256 122275213c349593 :: first vocabulary size at which it is not constructible
\newcommand{\UseLastAliveVcbEqone}{162}  % configs/labelmaps/usecase_survival.v1.json :: vcb_eq1.last_rung_alive :: sha256 122275213c349593
\newcommand{\UseDiesAtGnthreexQcdSubclass}{162}  % configs/labelmaps/usecase_survival.v1.json :: gn3x_qcd_subclass.dies_at :: sha256 122275213c349593 :: first vocabulary size at which it is not constructible
\newcommand{\UseLastAliveGnthreexQcdSubclass}{188}  % configs/labelmaps/usecase_survival.v1.json :: gn3x_qcd_subclass.last_rung_alive :: sha256 122275213c349593
\newcommand{\UseWlikeLastAlive}{162}  % configs/labelmaps/rung_label_maps.v1.csv :: groups at each level vs {label_X_cs, label_X_qq} and that plus the QCD classes :: sha256 412f9aa3b58db412 :: last level at which Sophon's W-like score g_W(2) is constructible
\newcommand{\UseWlikeDiesAt}{64}  % configs/labelmaps/rung_label_maps.v1.csv :: groups at each level vs {label_X_cs, label_X_qq} and that plus the QCD classes :: sha256 412f9aa3b58db412 :: first level at which it is not
\newcommand{\DesignEpochs}{80}  % experiments/MTX/k8s/job-mtx-l188-s1-raunav.yaml :: --num-epochs :: sha256 96fef4b6337aaa27
\newcommand{\DesignSamplesPerEpoch}{10{,}240{,}000}  % experiments/MTX/k8s/job-mtx-l188-s1-raunav.yaml :: --samples-per-epoch :: sha256 96fef4b6337aaa27
\newcommand{\DesignExamplesSeen}{\ensuremath{8.192\times10^{8}}}  % experiments/MTX/k8s/job-mtx-l188-s1-raunav.yaml :: --num-epochs x --samples-per-epoch :: sha256 96fef4b6337aaa27 :: training examples drawn over the run
\newcommand{\DesignPassesSelected}{7.45}  % experiments/FIGS/data/v2/jc2_train/counts.json :: n_selected :: sha256 b0b5a8b63aa808c0 :: --num-epochs x --samples-per-epoch of job-mtx-l188-s1-raunav.yaml / n_selected: passes over the selected training jets
\newcommand{\DesignBatchSize}{512}  % experiments/MTX/k8s/job-mtx-l188-s1-raunav.yaml :: --batch-size :: sha256 96fef4b6337aaa27
\newcommand{\DesignStartLr}{\ensuremath{5\times10^{-4}}}  % experiments/MTX/k8s/job-mtx-l188-s1-raunav.yaml :: --start-lr :: sha256 96fef4b6337aaa27
\newcommand{\DesignHeadHidden}{512}  % experiments/MTX/k8s/job-mtx-l188-s1-raunav.yaml :: -o fc_params (hidden width) :: sha256 96fef4b6337aaa27
\newcommand{\DesignParticleBlocks}{8}  % experiments/E1/ParT_sophon_arch_10c.py :: num_layers=(\d+) :: sha256 5793faa7d55a8b0d
\newcommand{\DesignClassBlocks}{2}  % experiments/E1/ParT_sophon_arch_10c.py :: num_cls_layers=(\d+) :: sha256 5793faa7d55a8b0d
\newcommand{\DesignHeads}{8}  % experiments/E1/ParT_sophon_arch_10c.py :: num_heads=(\d+) :: sha256 5793faa7d55a8b0d
\newcommand{\DesignEmbedDim}{128}  % experiments/E1/ParT_sophon_arch_10c.py :: \bembed_dims=\[\d+, \d+, (\d+)\] :: sha256 5793faa7d55a8b0d
\newcommand{\DesignMpmMaskRate}{40\%}  % experiments/MTX/mpm.py :: DEFAULT_MASK_RATE :: sha256 49f0c8c6666d0211 :: share of each jet's particles hidden from the encoder
\newcommand{\DesignPtMin}{200}  % configs/arms/L188.yaml :: selection :: sha256 d5a188dc80885515 :: GeV
\newcommand{\DesignPtMax}{2{,}500}  % configs/arms/L188.yaml :: selection :: sha256 d5a188dc80885515 :: GeV
\newcommand{\DesignMsdMin}{20}  % configs/arms/L188.yaml :: selection :: sha256 d5a188dc80885515 :: GeV
\newcommand{\DesignMsdMax}{500}  % configs/arms/L188.yaml :: selection :: sha256 d5a188dc80885515 :: GeV
\newcommand{\LitJetClassTwoNJetsMillion}{139}  % experiments/FIGS/data/literature_facts.json :: facts[0].value :: sha256 713dba2a5e6e1dd1 :: arXiv:2405.12972v1 App. A, p. 11
\newcommand{\LitJetClassTwoNResonantClasses}{161}  % experiments/FIGS/data/literature_facts.json :: facts[1].value :: sha256 713dba2a5e6e1dd1 :: arXiv:2405.12972v1 App. A, p. 11
\newcommand{\LitJetClassTwoNQcdClasses}{27}  % experiments/FIGS/data/literature_facts.json :: facts[2].value :: sha256 713dba2a5e6e1dd1 :: arXiv:2405.12972v1 App. A, p. 11
\newcommand{\LitJetClassTwoMaxResonanceMass}{500}  % experiments/FIGS/data/literature_facts.json :: facts[3].value :: sha256 713dba2a5e6e1dd1 :: arXiv:2405.12972v1 Sec. II, p. 2
\newcommand{\LitMadgraphVersion}{2.9.18}  % experiments/FIGS/data/literature_facts.json :: facts[4].value :: sha256 713dba2a5e6e1dd1 :: arXiv:2405.12972v1 App. A, p. 11
\newcommand{\LitPythiaVersion}{8.3}  % experiments/FIGS/data/literature_facts.json :: facts[5].value :: sha256 713dba2a5e6e1dd1 :: arXiv:2405.12972v1 App. A, p. 11
\newcommand{\LitDelphesVersion}{3}  % experiments/FIGS/data/literature_facts.json :: facts[6].value :: sha256 713dba2a5e6e1dd1 :: arXiv:2405.12972v1 App. A, p. 11
\newcommand{\LitJetClassTwoPileupMean}{50}  % experiments/FIGS/data/literature_facts.json :: facts[7].value :: sha256 713dba2a5e6e1dd1 :: arXiv:2405.12972v1 Sec. II, p. 2
\newcommand{\LitJetClassTwoJetRadius}{0.8}  % experiments/FIGS/data/literature_facts.json :: facts[8].value :: sha256 713dba2a5e6e1dd1 :: arXiv:2405.12972v1 Sec. II, p. 2
\newcommand{\LitAtlasUnifiedEffDropPercent}{10}  % experiments/FIGS/data/literature_facts.json :: facts[9].value :: sha256 713dba2a5e6e1dd1 :: ATL-PHYS-PUB-2026-013 Sec. 6.5, p. 22
\newcommand{\LitSophonPaperParticleBlocks}{6}  % experiments/FIGS/data/literature_facts.json :: facts[10].value :: sha256 713dba2a5e6e1dd1 :: arXiv:2405.12972v1 App. B.1, p. 12
\newcommand{\LitSophonSamplesPerEpochMillion}{10}  % experiments/FIGS/data/literature_facts.json :: facts[11].value :: sha256 713dba2a5e6e1dd1 :: arXiv:2405.12972v1 App. B.1, p. 12
\newcommand{\LitSophonSubsetSamplesPerEpochMillion}{2.5}  % experiments/FIGS/data/literature_facts.json :: facts[12].value :: sha256 713dba2a5e6e1dd1 :: arXiv:2405.12972v1 App. B.2, p. 12
\newcommand{\LitGlopartNClasses}{314}  % experiments/FIGS/data/literature_facts.json :: facts[13].value :: sha256 713dba2a5e6e1dd1 :: C. Li, PhD thesis, CERN-THESIS-2024-281 Sec. 11.2.1, p. 188
\newcommand{\LitPartFtEpochs}{20}  % experiments/FIGS/data/literature_facts.json :: facts[14].value :: sha256 713dba2a5e6e1dd1 :: arXiv:2202.03772 Sec. 5.2, p. 8
\newcommand{\LitMehtaNCategories}{eight}  % experiments/FIGS/data/literature_facts.json :: facts[15].value :: sha256 713dba2a5e6e1dd1 :: arXiv:2605.05556v2 Abstract
\newcommand{\LitHongPretrainRuns}{2}  % experiments/FIGS/data/literature_facts.json :: facts[16].value :: sha256 713dba2a5e6e1dd1 :: arXiv:2410.23129 App. A.1.1, p. 18
\newcommand{\LitSophonLookaheadK}{6}  % experiments/FIGS/data/literature_facts.json :: facts[17].value :: sha256 713dba2a5e6e1dd1 :: arXiv:2405.12972v1 App. B.1, p. 12
\newcommand{\LitSophonLookaheadAlpha}{0.5}  % experiments/FIGS/data/literature_facts.json :: facts[18].value :: sha256 713dba2a5e6e1dd1 :: arXiv:2405.12972v1 App. B.1, p. 12
\newcommand{\LitSophonRadamBetaOne}{0.95}  % experiments/FIGS/data/literature_facts.json :: facts[19].value :: sha256 713dba2a5e6e1dd1 :: arXiv:2405.12972v1 App. B.1, p. 12
\newcommand{\LitSophonRadamBetaTwo}{0.999}  % experiments/FIGS/data/literature_facts.json :: facts[20].value :: sha256 713dba2a5e6e1dd1 :: arXiv:2405.12972v1 App. B.1, p. 12
\newcommand{\LitSophonRadamEps}{\ensuremath{10^{-5}}}  % experiments/FIGS/data/literature_facts.json :: facts[21].value :: sha256 713dba2a5e6e1dd1 :: arXiv:2405.12972v1 App. B.1, p. 12
\newcommand{\LitGnTwoAuxRejectionDropPercent}{30}  % experiments/FIGS/data/literature_facts.json :: facts[22].value :: sha256 713dba2a5e6e1dd1 :: arXiv:2505.19689v2 (Nature Commun. 17 (2026) 541) p. 9
\newcommand{\LitCmsPartNClasses}{37}  % experiments/FIGS/data/literature_facts.json :: facts[23].value :: sha256 713dba2a5e6e1dd1 :: arXiv:2604.09809v1 (CMS-JME-25-001) Sec. 5.2, p. 7
\newcommand{\LitCmsPartMassLambda}{0.05}  % experiments/FIGS/data/literature_facts.json :: facts[24].value :: sha256 713dba2a5e6e1dd1 :: arXiv:2604.09809v1 (CMS-JME-25-001) Sec. 5.2, p. 7
\newcommand{\LitGlopartVthreeNClasses}{374}  % experiments/FIGS/data/literature_facts.json :: facts[25].value :: sha256 713dba2a5e6e1dd1 :: CMS-DP-2026-104 (CMS public page) Architecture and outputs
\newcommand{\LitHuhCoarseClasses}{127}  % experiments/FIGS/data/literature_facts.json :: facts[26].value :: sha256 713dba2a5e6e1dd1 :: arXiv:1608.08614v2 Sec. 5.3, p. 5
\newcommand{\LitHuhKnnGapPercent}{15}  % experiments/FIGS/data/literature_facts.json :: facts[27].value :: sha256 713dba2a5e6e1dd1 :: arXiv:1608.08614v2 Sec. 5.3, p. 5
\newcommand{\LitRsthreelRuns}{3}  % experiments/FIGS/data/literature_facts.json :: facts[28].value :: sha256 713dba2a5e6e1dd1 :: arXiv:2403.07066v2 (Phys. Rev. D 111 (2025) 032010) Sec. IV, p. 5
\newcommand{\LitSophonParamsMillion}{2.3}  % experiments/FIGS/data/literature_facts.json :: facts[29].value :: sha256 713dba2a5e6e1dd1 :: arXiv:2405.12972v1 App. B.1, p. 12
\newcommand{\LitGlopartVoneNClasses}{37}  % experiments/FIGS/data/literature_facts.json :: facts[30].value :: sha256 713dba2a5e6e1dd1 :: CMS-DP-2026-104 (CMS public page) Model generations
\newcommand{\LitSophonXMass}{200}  % experiments/FIGS/data/literature_facts.json :: facts[31].value :: sha256 713dba2a5e6e1dd1 :: arXiv:2405.12972v1 Sec. III, p. 3
\newcommand{\LitSophonXLumi}{40}  % experiments/FIGS/data/literature_facts.json :: facts[32].value :: sha256 713dba2a5e6e1dd1 :: arXiv:2405.12972v1 Sec. III, Fig. 2
\newcommand{\LitSophonXWindowLo}{150}  % experiments/FIGS/data/literature_facts.json :: facts[33].value :: sha256 713dba2a5e6e1dd1 :: arXiv:2405.12972v1 Sec. III, Fig. 2
\newcommand{\LitSophonXWindowHi}{230}  % experiments/FIGS/data/literature_facts.json :: facts[34].value :: sha256 713dba2a5e6e1dd1 :: arXiv:2405.12972v1 Sec. III, Fig. 2
\newcommand{\LitSophonAblationNClasses}{42}  % experiments/FIGS/data/literature_facts.json :: facts[35].value :: sha256 713dba2a5e6e1dd1 :: arXiv:2405.12972v1 Sec. III
\newcommand{\LitSophonTransferSamplesMillion}{2.5}  % experiments/FIGS/data/literature_facts.json :: facts[36].value :: sha256 713dba2a5e6e1dd1 :: arXiv:2405.12972v1 App. B.5, p. 13
\newcommand{\LitSophonTransferBatch}{1024}  % experiments/FIGS/data/literature_facts.json :: facts[37].value :: sha256 713dba2a5e6e1dd1 :: arXiv:2405.12972v1 App. B.5, p. 13
\newcommand{\LitSophonScratchEpochs}{50}  % experiments/FIGS/data/literature_facts.json :: facts[38].value :: sha256 713dba2a5e6e1dd1 :: arXiv:2405.12972v1 App. B.3
\newcommand{\LitHuhFineClasses}{1000}  % experiments/FIGS/data/literature_facts.json :: facts[39].value :: sha256 713dba2a5e6e1dd1 :: arXiv:1608.08614v2 Sec. 1, p. 1
\newcommand{\LitHuhPascalDropMap}{2.8}  % experiments/FIGS/data/literature_facts.json :: facts[40].value :: sha256 713dba2a5e6e1dd1 :: arXiv:1608.08614v2 Sec. 1, p. 1
\newcommand{\LitHongSuperclasses}{11}  % experiments/FIGS/data/literature_facts.json :: facts[41].value :: sha256 713dba2a5e6e1dd1 :: arXiv:2410.23129 Sec. 6.2
\newcommand{\LitHongDeclineClasses}{\ensuremath{10^{3}}}  % experiments/FIGS/data/literature_facts.json :: facts[42].value :: sha256 713dba2a5e6e1dd1 :: arXiv:2410.23129 Sec. 6.2
\newcommand{\LitObjectivesFullJets}{\ensuremath{1.2\times10^{6}}}  % experiments/FIGS/data/literature_facts.json :: facts[43].value :: sha256 713dba2a5e6e1dd1 :: arXiv:2606.14870 Sec. IV A
\newcommand{\LitObjectivesFewJets}{\ensuremath{10^{4}}}  % experiments/FIGS/data/literature_facts.json :: facts[44].value :: sha256 713dba2a5e6e1dd1 :: arXiv:2606.14870 Sec. IV A
\newcommand{\LitMpmBenefitJets}{10k}  % experiments/FIGS/data/literature_facts.json :: facts[45].value :: sha256 713dba2a5e6e1dd1 :: arXiv:2401.13537 Sec. V
\newcommand{\LitPartTopRejScratch}{1602}  % experiments/FIGS/data/literature_facts.json :: facts[46].value :: sha256 713dba2a5e6e1dd1 :: arXiv:2202.03772 Table 5
\newcommand{\LitPartTopRejFt}{2766}  % experiments/FIGS/data/literature_facts.json :: facts[47].value :: sha256 713dba2a5e6e1dd1 :: arXiv:2202.03772 Table 5
\newcommand{\LitMahajanLinearGapPercent}{7}  % experiments/FIGS/data/literature_facts.json :: facts[48].value :: sha256 713dba2a5e6e1dd1 :: arXiv:1805.00932 Sec. 3.1.2
\newcommand{\LitBukchinCoarse}{33.83}  % experiments/FIGS/data/literature_facts.json :: facts[49].value :: sha256 713dba2a5e6e1dd1 :: arXiv:2012.03515 Table 2, LIVING-17 all-way
\newcommand{\LitBukchinFine}{58.95}  % experiments/FIGS/data/literature_facts.json :: facts[50].value :: sha256 713dba2a5e6e1dd1 :: arXiv:2012.03515 Table 2, LIVING-17 all-way
\newcommand{\LitKornblithRho}{\ensuremath{-0.93}}  % experiments/FIGS/data/literature_facts.json :: facts[51].value :: sha256 713dba2a5e6e1dd1 :: arXiv:2010.16402 Sec. 3.4
\newcommand{\LitPartTopSigEffPercent}{30}  % experiments/FIGS/data/literature_facts.json :: facts[52].value :: sha256 713dba2a5e6e1dd1 :: arXiv:2202.03772 Table 5 (column Rej30%)
\newcommand{\LitPartTopKinFeatures}{7}  % experiments/FIGS/data/literature_facts.json :: facts[53].value :: sha256 713dba2a5e6e1dd1 :: arXiv:2202.03772 Sec. 5.2
\newcommand{\LitBukchinAncor}{45.14}  % experiments/FIGS/data/literature_facts.json :: facts[54].value :: sha256 713dba2a5e6e1dd1 :: arXiv:2012.03515 Table 2, LIVING-17 all-way
\newcommand{\LitMahajanFtGapPp}{0.1}  % experiments/FIGS/data/literature_facts.json :: facts[55].value :: sha256 713dba2a5e6e1dd1 :: arXiv:1805.00932 Fig. 1, ImageNet-9k target
\newcommand{\LitSophonTransferLr}{5 × 10−4}  % experiments/FIGS/data/literature_facts.json :: facts[56].value :: sha256 713dba2a5e6e1dd1 :: arXiv:2405.12972v1 App. B.5
\newcommand{\LitQgJetRadius}{0.4}  % experiments/FIGS/data/literature_facts.json :: facts[57].value :: sha256 713dba2a5e6e1dd1 :: arXiv:1810.05165 Sec. 4.1
\newcommand{\LitQgPtMin}{500}  % experiments/FIGS/data/literature_facts.json :: facts[58].value :: sha256 713dba2a5e6e1dd1 :: arXiv:1810.05165 Sec. 4.1
\newcommand{\LitQgPtMax}{550}  % experiments/FIGS/data/literature_facts.json :: facts[59].value :: sha256 713dba2a5e6e1dd1 :: arXiv:1810.05165 Sec. 4.1
\newcommand{\LitPartTopAucScratch}{0.9858}  % experiments/FIGS/data/literature_facts.json :: facts[60].value :: sha256 713dba2a5e6e1dd1 :: arXiv:2202.03772 Table 5
\newcommand{\LitPartTopAucFt}{0.9877}  % experiments/FIGS/data/literature_facts.json :: facts[61].value :: sha256 713dba2a5e6e1dd1 :: arXiv:2202.03772 Table 5
\newcommand{\LitPartQgAucScratch}{0.9121}  % experiments/FIGS/data/literature_facts.json :: facts[62].value :: sha256 713dba2a5e6e1dd1 :: arXiv:2202.03772 Table 6 (exp)
\newcommand{\LitPartQgRejScratch}{101.2}  % experiments/FIGS/data/literature_facts.json :: facts[63].value :: sha256 713dba2a5e6e1dd1 :: arXiv:2202.03772 Table 6 (exp)
\newcommand{\LitPartQgAucFt}{0.9151}  % experiments/FIGS/data/literature_facts.json :: facts[64].value :: sha256 713dba2a5e6e1dd1 :: arXiv:2202.03772 Table 6 (exp)
\newcommand{\LitPartQgRejFt}{107.9}  % experiments/FIGS/data/literature_facts.json :: facts[65].value :: sha256 713dba2a5e6e1dd1 :: arXiv:2202.03772 Table 6 (exp)
\newcommand{\LitPartJcTestMillion}{20}  % experiments/FIGS/data/literature_facts.json :: facts[66].value :: sha256 713dba2a5e6e1dd1 :: arXiv:2202.03772 Sec. 3
\newcommand{\LitPartJcTwoMAcc}{0.836}  % experiments/FIGS/data/literature_facts.json :: facts[67].value :: sha256 713dba2a5e6e1dd1 :: arXiv:2202.03772 Table 3
\newcommand{\LitPartSmallSetIters}{100 k}  % experiments/FIGS/data/literature_facts.json :: facts[68].value :: sha256 713dba2a5e6e1dd1 :: arXiv:2202.03772 Sec. 5.1 (training dataset size)
\newcommand{\LitPartSmallSetPercent}{2}  % experiments/FIGS/data/literature_facts.json :: facts[69].value :: sha256 713dba2a5e6e1dd1 :: arXiv:2202.03772 Sec. 5.1 (training dataset size)
\newcommand{\LitSophonEtaMax}{2.5}  % experiments/FIGS/data/literature_facts.json :: facts[70].value :: sha256 713dba2a5e6e1dd1 :: arXiv:2405.12972 App. A
\newcommand{\LitBHadronLifetimePs}{1.5}  % experiments/FIGS/data/literature_facts.json :: facts[71].value :: sha256 713dba2a5e6e1dd1 :: arXiv:1712.07158 Sec. 4.1
\newcommand{\LitCHadronLifetimePs}{1}  % experiments/FIGS/data/literature_facts.json :: facts[72].value :: sha256 713dba2a5e6e1dd1 :: arXiv:1712.07158 Sec. 4.1
\newcommand{\LitOmnilearnPidDropProb}{0.2}  % experiments/FIGS/data/literature_facts.json :: facts[73].value :: sha256 713dba2a5e6e1dd1 :: arXiv:2502.14652 Sec. II.A
\newcommand{\LitMahajanFtNineKLarge}{48.5}  % experiments/FIGS/data/literature_facts.json :: facts[74].value :: sha256 713dba2a5e6e1dd1 :: arXiv:1805.00932 Fig. 1 (IG-1B-17k, IN-9k, full fine-tuning)
\newcommand{\LitMahajanFtNineKSmall}{48.4}  % experiments/FIGS/data/literature_facts.json :: facts[75].value :: sha256 713dba2a5e6e1dd1 :: arXiv:1805.00932 Fig. 1 (IG-940M-1.5k, IN-9k, full fine-tuning)
\newcommand{\LitMahajanFtOneKSmall}{84.2}  % experiments/FIGS/data/literature_facts.json :: facts[76].value :: sha256 713dba2a5e6e1dd1 :: arXiv:1805.00932 Fig. 1 (IG-940M-1.5k, IN-1k, full fine-tuning)
\newcommand{\LitMahajanFtOneKLarge}{83.6}  % experiments/FIGS/data/literature_facts.json :: facts[77].value :: sha256 713dba2a5e6e1dd1 :: arXiv:1805.00932 Fig. 1 (IG-1B-17k, IN-1k, full fine-tuning)
\newcommand{\LitChenGainFifteen}{0.16}  % experiments/FIGS/data/literature_facts.json :: facts[78].value :: sha256 713dba2a5e6e1dd1 :: arXiv:1901.07012 Table VIII (15 coarse classes, CIFAR-100)
\newcommand{\LitChenGainTwenty}{0.01}  % experiments/FIGS/data/literature_facts.json :: facts[79].value :: sha256 713dba2a5e6e1dd1 :: arXiv:1901.07012 Table VIII (20 coarse classes, CIFAR-100)
\newcommand{\LitChenCoarseFifteen}{15}  % experiments/FIGS/data/literature_facts.json :: facts[80].value :: sha256 713dba2a5e6e1dd1 :: arXiv:1901.07012 Table VIII
\newcommand{\LitChenCoarseTwenty}{20}  % experiments/FIGS/data/literature_facts.json :: facts[81].value :: sha256 713dba2a5e6e1dd1 :: arXiv:1901.07012 Table VIII
\newcommand{\LitChenRandomMaxGain}{2.78}  % experiments/FIGS/data/literature_facts.json :: facts[82].value :: sha256 713dba2a5e6e1dd1 :: arXiv:1901.07012 Sec. VI-A, Table VI
\newcommand{\LitGlopartMassLambda}{5}  % experiments/FIGS/data/literature_facts.json :: facts[83].value :: sha256 713dba2a5e6e1dd1 :: CERN-THESIS-2024-281 Eq. (11.7), Sec. 11.2.1
\newcommand{\LitGlopartAddedAucBase}{0.9555}  % experiments/FIGS/data/literature_facts.json :: facts[84].value :: sha256 713dba2a5e6e1dd1 :: CERN-THESIS-2024-281 Sec. 11.4.1, Fig. 11.19 (H->WW 4q vs QCD)
\newcommand{\LitGlopartAddedAucTop}{0.9593}  % experiments/FIGS/data/literature_facts.json :: facts[85].value :: sha256 713dba2a5e6e1dd1 :: CERN-THESIS-2024-281 Sec. 11.4.1, Fig. 11.19 (H->WW 4q vs QCD)
\newcommand{\LitCmsSdZcut}{0.1}  % experiments/FIGS/data/literature_facts.json :: facts[86].value :: sha256 713dba2a5e6e1dd1 :: arXiv:2004.08262 Sec. 6.1
\newcommand{\LitCmsSdBeta}{0}  % experiments/FIGS/data/literature_facts.json :: facts[87].value :: sha256 713dba2a5e6e1dd1 :: arXiv:2004.08262 Sec. 6.1
\newcommand{\LitQgNJetsMillion}{2}  % experiments/FIGS/data/literature_facts.json :: facts[88].value :: sha256 713dba2a5e6e1dd1 :: arXiv:2202.03772 Sec. 5.2
\newcommand{\LitQgTrainMillion}{1.6}  % experiments/FIGS/data/literature_facts.json :: facts[89].value :: sha256 713dba2a5e6e1dd1 :: arXiv:2202.03772 Sec. 5.2
\newcommand{\LitQgValMillion}{0.2}  % experiments/FIGS/data/literature_facts.json :: facts[90].value :: sha256 713dba2a5e6e1dd1 :: arXiv:2202.03772 Sec. 5.2
\newcommand{\LitQgRapidityMax}{1.7}  % experiments/FIGS/data/literature_facts.json :: facts[91].value :: sha256 713dba2a5e6e1dd1 :: Zenodo 10.5281/zenodo.3164691 Zenodo record 3164691, description
\newcommand{\LitHerwigVersion}{7.1.4}  % experiments/FIGS/data/literature_facts.json :: facts[92].value :: sha256 713dba2a5e6e1dd1 :: Zenodo 10.5281/zenodo.3066475 record description
\newcommand{\LitObjectivesFewestJets}{\ensuremath{10^{3}}}  % experiments/FIGS/data/literature_facts.json :: facts[93].value :: sha256 713dba2a5e6e1dd1 :: arXiv:2606.14870 Table XVII (Micro model, Npt = 10^8: Classifier 136 vs MPM 218 at 10^3, 331 vs 539 at 10^4)
\newcommand{\LitPartTopNTrainMillion}{1.2}  % experiments/FIGS/data/literature_facts.json :: facts[94].value :: sha256 713dba2a5e6e1dd1 :: arXiv:2202.03772 Sec. 5.2
\newcommand{\LitLandscapeParticlenetEnsMean}{1412}  % experiments/FIGS/data/literature_facts.json :: facts[95].value :: sha256 713dba2a5e6e1dd1 :: arXiv:1902.09914 Table 1 (ParticleNet, ensemble mean of 1/eps_B at eps_S=0.3)
\newcommand{\LitPartFtHeadReinits}{9}  % experiments/FIGS/data/literature_facts.json :: facts[96].value :: sha256 713dba2a5e6e1dd1 :: arXiv:2202.03772 Sec. 5.2
\newcommand{\LitSophonBlocksCode}{8}  % experiments/FIGS/data/literature_facts.json :: facts[97].value :: sha256 713dba2a5e6e1dd1 :: github.com/jet-universe/sophon@9dd6dd6 networks/example_ParticleTransformer_sophon.py line 67
\newcommand{\LitGrasHerwigVersion}{2.7.1}  % experiments/FIGS/data/literature_facts.json :: facts[98].value :: sha256 713dba2a5e6e1dd1 :: arXiv:1704.03878 Sec. 2 (list of generators)
\newcommand{\DesignMassLambda}{5.0}  % experiments/MTX/k8s/job-mtx-l162_mass-s1-raunav.yaml :: --mass-lambda :: sha256 8882459fe5520795 :: identical in all ten mass-output pretraining specs
\newcommand{\FtRecipeNRuns}{264}  % experiments/FIGS/data/ft_recipes/recipes_w2b_bench_v2_v3.json :: runs (count), reported cells :: sha256 7dc8c791343e1b6f
\newcommand{\FtRecipeLr}{\ensuremath{1\times10^{-4}}}  % experiments/FIGS/data/ft_recipes/recipes_w2b_bench_v2_v3.json :: runs[*].lr, pretrained starts :: sha256 7dc8c791343e1b6f
\newcommand{\FtRecipeHeadMult}{50}  % experiments/FIGS/data/ft_recipes/recipes_w2b_bench_v2_v3.json :: runs[*].head_lr_mult, pretrained starts :: sha256 7dc8c791343e1b6f
\newcommand{\FtRecipeLrHead}{\ensuremath{5\times10^{-3}}}  % experiments/FIGS/data/ft_recipes/recipes_w2b_bench_v2_v3.json :: lr x head_lr_mult :: sha256 7dc8c791343e1b6f
\newcommand{\FtRecipeWeightDecay}{0.01}  % experiments/FIGS/data/ft_recipes/recipes_w2b_bench_v2_v3.json :: runs[*].weight_decay :: sha256 7dc8c791343e1b6f
\newcommand{\FtRecipeBatch}{512}  % experiments/FIGS/data/ft_recipes/recipes_w2b_bench_v2_v3.json :: runs[*].batch_size :: sha256 7dc8c791343e1b6f
\newcommand{\FtRecipeEpochsEThree}{50}  % experiments/FIGS/data/ft_recipes/recipes_w2b_bench_v2_v3.json :: runs[leg 1,2; n_train 1000].epochs :: sha256 7dc8c791343e1b6f
\newcommand{\FtRecipeStepsEThree}{19}  % experiments/FT/k8s/job-ft-legs-w2-raunav.yaml :: job commands: samples per epoch at n_train 1000 // batch size :: sha256 f0962eddb451b7b0 :: optimizer steps per epoch, as weaver floors them
\newcommand{\FtRecipePassesEThree}{10}  % experiments/FT/k8s/job-ft-legs-w2-raunav.yaml :: job commands: samples per epoch at n_train 1000 // n_train :: sha256 f0962eddb451b7b0 :: passes over the training subset per epoch
\newcommand{\FtRecipeEpochsEFour}{50}  % experiments/FIGS/data/ft_recipes/recipes_w2b_bench_v2_v3.json :: runs[leg 1,2; n_train 10000].epochs :: sha256 7dc8c791343e1b6f
\newcommand{\FtRecipeStepsEFour}{19}  % experiments/FT/k8s/job-ft-legs-w2-raunav.yaml :: job commands: samples per epoch at n_train 10000 // batch size :: sha256 f0962eddb451b7b0 :: optimizer steps per epoch, as weaver floors them
\newcommand{\FtRecipePassesEFour}{one}  % experiments/FT/k8s/job-ft-legs-w2-raunav.yaml :: job commands: samples per epoch at n_train 10000 // n_train :: sha256 f0962eddb451b7b0 :: passes over the training subset per epoch
\newcommand{\FtRecipeEpochsEFive}{30}  % experiments/FIGS/data/ft_recipes/recipes_w2b_bench_v2_v3.json :: runs[leg 1,2; n_train 100000].epochs :: sha256 7dc8c791343e1b6f
\newcommand{\FtRecipeStepsEFive}{195}  % experiments/FT/k8s/job-ft-legs-w2-raunav.yaml :: job commands: samples per epoch at n_train 100000 // batch size :: sha256 f0962eddb451b7b0 :: optimizer steps per epoch, as weaver floors them
\newcommand{\FtRecipePassesEFive}{one}  % experiments/FT/k8s/job-ft-legs-w2-raunav.yaml :: job commands: samples per epoch at n_train 100000 // n_train :: sha256 f0962eddb451b7b0 :: passes over the training subset per epoch
\newcommand{\FtRecipeEpochsESix}{10}  % experiments/FIGS/data/ft_recipes/recipes_w2b_bench_v2_v3.json :: runs[leg 1,2; n_train 1000000].epochs :: sha256 7dc8c791343e1b6f
\newcommand{\FtRecipeStepsESix}{1{,}953}  % experiments/FT/k8s/job-ft-legs-w2-raunav.yaml :: job commands: samples per epoch at n_train 1000000 // batch size :: sha256 f0962eddb451b7b0 :: optimizer steps per epoch, as weaver floors them
\newcommand{\FtRecipePassesESix}{one}  % experiments/FT/k8s/job-ft-legs-w2-raunav.yaml :: job commands: samples per epoch at n_train 1000000 // n_train :: sha256 f0962eddb451b7b0 :: passes over the training subset per epoch
\newcommand{\FtRecipeValJets}{19{,}968}  % experiments/FT/k8s/job-ft-legs-w2-raunav.yaml :: --samples-per-epoch-val 20000 // batch x batch :: sha256 f0962eddb451b7b0 :: validation jets per epoch, floored to whole batches
\newcommand{\RecoveryAccOneeighteightByOneeighteight}{0.4123\,\ensuremath{\pm}\,0.0011}  % experiments/FIGS/data/label_recovery_curve_v1err/analysis/label_recovery_curve_summary.json :: table[rung=L188,level=188,probe=linear,n_train=max].accuracy :: sha256 1250d1f00537e0bd :: balanced accuracy, mean +- SD over 5 runs
\newcommand{\RecoveryAccOneeighteightByOnesixtwo}{0.4090\,\ensuremath{\pm}\,0.0012}  % experiments/FIGS/data/label_recovery_curve_v1err/analysis/label_recovery_curve_summary.json :: table[rung=L188,level=162,probe=linear,n_train=max].accuracy :: sha256 1250d1f00537e0bd :: balanced accuracy, mean +- SD over 5 runs
\newcommand{\RecoveryAccOneeighteightByFourthree}{0.356\,\ensuremath{\pm}\,0.008}  % experiments/FIGS/data/label_recovery_curve_v1err/analysis/label_recovery_curve_summary.json :: table[rung=L188,level=43,probe=linear,n_train=max].accuracy :: sha256 1250d1f00537e0bd :: balanced accuracy, mean +- SD over 5 runs
\newcommand{\RecoveryAccOneeighteightByOneseven}{0.306\,\ensuremath{\pm}\,0.008}  % experiments/FIGS/data/label_recovery_curve_v1err/analysis/label_recovery_curve_summary.json :: table[rung=L188,level=17,probe=linear,n_train=max].accuracy :: sha256 1250d1f00537e0bd :: balanced accuracy, mean +- SD over 5 runs
\newcommand{\RecoveryAccOnesixtwoByOneeighteight}{0.4623\,\ensuremath{\pm}\,0.0006}  % experiments/FIGS/data/label_recovery_curve_v1err/analysis/label_recovery_curve_summary.json :: table[rung=L162,level=188,probe=linear,n_train=max].accuracy :: sha256 1250d1f00537e0bd :: balanced accuracy, mean +- SD over 5 runs
\newcommand{\RecoveryAccOnesixtwoByOnesixtwo}{0.4622\,\ensuremath{\pm}\,0.0007}  % experiments/FIGS/data/label_recovery_curve_v1err/analysis/label_recovery_curve_summary.json :: table[rung=L162,level=162,probe=linear,n_train=max].accuracy :: sha256 1250d1f00537e0bd :: balanced accuracy, mean +- SD over 5 runs
\newcommand{\RecoveryAccOnesixtwoByFourthree}{0.402\,\ensuremath{\pm}\,0.009}  % experiments/FIGS/data/label_recovery_curve_v1err/analysis/label_recovery_curve_summary.json :: table[rung=L162,level=43,probe=linear,n_train=max].accuracy :: sha256 1250d1f00537e0bd :: balanced accuracy, mean +- SD over 5 runs
\newcommand{\RecoveryAccOnesixtwoByOneseven}{0.348\,\ensuremath{\pm}\,0.010}  % experiments/FIGS/data/label_recovery_curve_v1err/analysis/label_recovery_curve_summary.json :: table[rung=L162,level=17,probe=linear,n_train=max].accuracy :: sha256 1250d1f00537e0bd :: balanced accuracy, mean +- SD over 5 runs
\newcommand{\RecoveryAccSixfourByOneeighteight}{0.5979\,\ensuremath{\pm}\,0.0008}  % experiments/FIGS/data/label_recovery_curve_v1err/analysis/label_recovery_curve_summary.json :: table[rung=R63_Q1,level=188,probe=linear,n_train=max].accuracy :: sha256 1250d1f00537e0bd :: balanced accuracy, mean +- SD over 5 runs
\newcommand{\RecoveryAccSixfourByOnesixtwo}{0.5973\,\ensuremath{\pm}\,0.0013}  % experiments/FIGS/data/label_recovery_curve_v1err/analysis/label_recovery_curve_summary.json :: table[rung=R63_Q1,level=162,probe=linear,n_train=max].accuracy :: sha256 1250d1f00537e0bd :: balanced accuracy, mean +- SD over 5 runs
\newcommand{\RecoveryAccSixfourByFourthree}{0.5881\,\ensuremath{\pm}\,0.0007}  % experiments/FIGS/data/label_recovery_curve_v1err/analysis/label_recovery_curve_summary.json :: table[rung=R63_Q1,level=43,probe=linear,n_train=max].accuracy :: sha256 1250d1f00537e0bd :: balanced accuracy, mean +- SD over 5 runs
\newcommand{\RecoveryAccSixfourByOneseven}{0.531\,\ensuremath{\pm}\,0.006}  % experiments/FIGS/data/label_recovery_curve_v1err/analysis/label_recovery_curve_summary.json :: table[rung=R63_Q1,level=17,probe=linear,n_train=max].accuracy :: sha256 1250d1f00537e0bd :: balanced accuracy, mean +- SD over 5 runs
\newcommand{\RecoveryAccFourthreeByOneeighteight}{0.6770\,\ensuremath{\pm}\,0.0006}  % experiments/FIGS/data/label_recovery_curve_v1err/analysis/label_recovery_curve_summary.json :: table[rung=R42_Q1,level=188,probe=linear,n_train=max].accuracy :: sha256 1250d1f00537e0bd :: balanced accuracy, mean +- SD over 5 runs
\newcommand{\RecoveryAccFourthreeByOnesixtwo}{0.6766\,\ensuremath{\pm}\,0.0008}  % experiments/FIGS/data/label_recovery_curve_v1err/analysis/label_recovery_curve_summary.json :: table[rung=R42_Q1,level=162,probe=linear,n_train=max].accuracy :: sha256 1250d1f00537e0bd :: balanced accuracy, mean +- SD over 5 runs
\newcommand{\RecoveryAccFourthreeByFourthree}{0.6751\,\ensuremath{\pm}\,0.0011}  % experiments/FIGS/data/label_recovery_curve_v1err/analysis/label_recovery_curve_summary.json :: table[rung=R42_Q1,level=43,probe=linear,n_train=max].accuracy :: sha256 1250d1f00537e0bd :: balanced accuracy, mean +- SD over 5 runs
\newcommand{\RecoveryAccFourthreeByOneseven}{0.633\,\ensuremath{\pm}\,0.006}  % experiments/FIGS/data/label_recovery_curve_v1err/analysis/label_recovery_curve_summary.json :: table[rung=R42_Q1,level=17,probe=linear,n_train=max].accuracy :: sha256 1250d1f00537e0bd :: balanced accuracy, mean +- SD over 5 runs
\newcommand{\RecoveryAccThreezeroByOneeighteight}{0.7001\,\ensuremath{\pm}\,0.0005}  % experiments/FIGS/data/label_recovery_curve_v1err/analysis/label_recovery_curve_summary.json :: table[rung=R29_Q1,level=188,probe=linear,n_train=max].accuracy :: sha256 1250d1f00537e0bd :: balanced accuracy, mean +- SD over 5 runs
\newcommand{\RecoveryAccThreezeroByOnesixtwo}{0.7004\,\ensuremath{\pm}\,0.0007}  % experiments/FIGS/data/label_recovery_curve_v1err/analysis/label_recovery_curve_summary.json :: table[rung=R29_Q1,level=162,probe=linear,n_train=max].accuracy :: sha256 1250d1f00537e0bd :: balanced accuracy, mean +- SD over 5 runs
\newcommand{\RecoveryAccThreezeroByFourthree}{0.6998\,\ensuremath{\pm}\,0.0008}  % experiments/FIGS/data/label_recovery_curve_v1err/analysis/label_recovery_curve_summary.json :: table[rung=R29_Q1,level=43,probe=linear,n_train=max].accuracy :: sha256 1250d1f00537e0bd :: balanced accuracy, mean +- SD over 5 runs
\newcommand{\RecoveryAccThreezeroByOneseven}{0.6778\,\ensuremath{\pm}\,0.0030}  % experiments/FIGS/data/label_recovery_curve_v1err/analysis/label_recovery_curve_summary.json :: table[rung=R29_Q1,level=17,probe=linear,n_train=max].accuracy :: sha256 1250d1f00537e0bd :: balanced accuracy, mean +- SD over 5 runs
\newcommand{\RecoveryAccOnesevenByOneeighteight}{0.7675\,\ensuremath{\pm}\,0.0010}  % experiments/FIGS/data/label_recovery_curve_v1err/analysis/label_recovery_curve_summary.json :: table[rung=R16_Q1,level=188,probe=linear,n_train=max].accuracy :: sha256 1250d1f00537e0bd :: balanced accuracy, mean +- SD over 5 runs
\newcommand{\RecoveryAccOnesevenByOnesixtwo}{0.7681\,\ensuremath{\pm}\,0.0009}  % experiments/FIGS/data/label_recovery_curve_v1err/analysis/label_recovery_curve_summary.json :: table[rung=R16_Q1,level=162,probe=linear,n_train=max].accuracy :: sha256 1250d1f00537e0bd :: balanced accuracy, mean +- SD over 5 runs
\newcommand{\RecoveryAccOnesevenByFourthree}{0.7680\,\ensuremath{\pm}\,0.0010}  % experiments/FIGS/data/label_recovery_curve_v1err/analysis/label_recovery_curve_summary.json :: table[rung=R16_Q1,level=43,probe=linear,n_train=max].accuracy :: sha256 1250d1f00537e0bd :: balanced accuracy, mean +- SD over 5 runs
\newcommand{\RecoveryAccOnesevenByOneseven}{0.7681\,\ensuremath{\pm}\,0.0003}  % experiments/FIGS/data/label_recovery_curve_v1err/analysis/label_recovery_curve_summary.json :: table[rung=R16_Q1,level=17,probe=linear,n_train=max].accuracy :: sha256 1250d1f00537e0bd :: balanced accuracy, mean +- SD over 5 runs
\newcommand{\RecoveryAccFourByOneeighteight}{0.7526\,\ensuremath{\pm}\,0.0002}  % experiments/FIGS/data/label_recovery_curve_v1err/analysis/label_recovery_curve_summary.json :: table[rung=R3_VIS,level=188,probe=linear,n_train=max].accuracy :: sha256 1250d1f00537e0bd :: balanced accuracy, mean +- SD over 5 runs
\newcommand{\RecoveryAccFourByOnesixtwo}{0.7538\,\ensuremath{\pm}\,0.0017}  % experiments/FIGS/data/label_recovery_curve_v1err/analysis/label_recovery_curve_summary.json :: table[rung=R3_VIS,level=162,probe=linear,n_train=max].accuracy :: sha256 1250d1f00537e0bd :: balanced accuracy, mean +- SD over 5 runs
\newcommand{\RecoveryAccFourByFourthree}{0.7557\,\ensuremath{\pm}\,0.0015}  % experiments/FIGS/data/label_recovery_curve_v1err/analysis/label_recovery_curve_summary.json :: table[rung=R3_VIS,level=43,probe=linear,n_train=max].accuracy :: sha256 1250d1f00537e0bd :: balanced accuracy, mean +- SD over 5 runs
\newcommand{\RecoveryAccFourByOneseven}{0.7586\,\ensuremath{\pm}\,0.0008}  % experiments/FIGS/data/label_recovery_curve_v1err/analysis/label_recovery_curve_summary.json :: table[rung=R3_VIS,level=17,probe=linear,n_train=max].accuracy :: sha256 1250d1f00537e0bd :: balanced accuracy, mean +- SD over 5 runs
\newcommand{\RecoveryAccTwoByOneeighteight}{0.8688\,\ensuremath{\pm}\,0.0005}  % experiments/FIGS/data/label_recovery_curve_v1err/analysis/label_recovery_curve_summary.json :: table[rung=R1_Q1,level=188,probe=linear,n_train=max].accuracy :: sha256 1250d1f00537e0bd :: balanced accuracy, mean +- SD over 5 runs
\newcommand{\RecoveryAccTwoByOnesixtwo}{0.8703\,\ensuremath{\pm}\,0.0008}  % experiments/FIGS/data/label_recovery_curve_v1err/analysis/label_recovery_curve_summary.json :: table[rung=R1_Q1,level=162,probe=linear,n_train=max].accuracy :: sha256 1250d1f00537e0bd :: balanced accuracy, mean +- SD over 5 runs
\newcommand{\RecoveryAccTwoByFourthree}{0.8709\,\ensuremath{\pm}\,0.0011}  % experiments/FIGS/data/label_recovery_curve_v1err/analysis/label_recovery_curve_summary.json :: table[rung=R1_Q1,level=43,probe=linear,n_train=max].accuracy :: sha256 1250d1f00537e0bd :: balanced accuracy, mean +- SD over 5 runs
\newcommand{\RecoveryAccTwoByOneseven}{0.8713\,\ensuremath{\pm}\,0.0002}  % experiments/FIGS/data/label_recovery_curve_v1err/analysis/label_recovery_curve_summary.json :: table[rung=R1_Q1,level=17,probe=linear,n_train=max].accuracy :: sha256 1250d1f00537e0bd :: balanced accuracy, mean +- SD over 5 runs
\newcommand{\RecoveryOneeighteightOnesixtwoMaxAbs}{0.0033}  % experiments/FIGS/data/label_recovery_curve_v1err/analysis/label_recovery_curve_summary.json :: table[rung=*,level=188/162,probe=linear,n_train=max].accuracy :: sha256 1250d1f00537e0bd :: max |difference of the run means| over the tree
\newcommand{\RecoveryPairedOnesevenLevelOnesevenMinusOneeighteight}{+0.0006~[\ensuremath{-}0.0004, +0.0015]}  % experiments/FIGS/data/label_recovery_curve_v1err/analysis/label_recovery_curve_summary.json :: table[rung=R16_Q1,level=17 minus 188,probe=linear,n_train=max].accuracy :: sha256 1250d1f00537e0bd :: paired by run, mean difference [95% Student-t interval]
\newcommand{\RecoveryPairedOnesevenLevelOnesevenMinusOnesixtwo}{+0.00004~[\ensuremath{-}0.00123, +0.00131]}  % experiments/FIGS/data/label_recovery_curve_v1err/analysis/label_recovery_curve_summary.json :: table[rung=R16_Q1,level=17 minus 162,probe=linear,n_train=max].accuracy :: sha256 1250d1f00537e0bd :: paired by run, mean difference [95% Student-t interval]
\newcommand{\RecoveryPairedFourLevelOnesevenMinusOneeighteight}{+0.0060~[+0.0051, +0.0069]}  % experiments/FIGS/data/label_recovery_curve_v1err/analysis/label_recovery_curve_summary.json :: table[rung=R3_VIS,level=17 minus 188,probe=linear,n_train=max].accuracy :: sha256 1250d1f00537e0bd :: paired by run, mean difference [95% Student-t interval]
\newcommand{\RecoveryPairedFourLevelOnesevenMinusOnesixtwo}{+0.0048~[+0.0028, +0.0068]}  % experiments/FIGS/data/label_recovery_curve_v1err/analysis/label_recovery_curve_summary.json :: table[rung=R3_VIS,level=17 minus 162,probe=linear,n_train=max].accuracy :: sha256 1250d1f00537e0bd :: paired by run, mean difference [95% Student-t interval]
\newcommand{\RecoveryPairedTwoLevelOnesevenMinusOneeighteight}{+0.0025~[+0.0019, +0.0032]}  % experiments/FIGS/data/label_recovery_curve_v1err/analysis/label_recovery_curve_summary.json :: table[rung=R1_Q1,level=17 minus 188,probe=linear,n_train=max].accuracy :: sha256 1250d1f00537e0bd :: paired by run, mean difference [95% Student-t interval]
\newcommand{\RecoveryPairedTwoLevelOnesevenMinusOnesixtwo}{+0.00101~[+0.00003, +0.00198]}  % experiments/FIGS/data/label_recovery_curve_v1err/analysis/label_recovery_curve_summary.json :: table[rung=R1_Q1,level=17 minus 162,probe=linear,n_train=max].accuracy :: sha256 1250d1f00537e0bd :: paired by run, mean difference [95% Student-t interval]
\newcommand{\RecoveryMlpGainOneseven}{+0.0099~[+0.0078, +0.0120]}  % experiments/FIGS/data/label_recovery_curve_v1err/analysis/label_recovery_curve_summary.json :: mlp_minus_linear_at_largest[rung=L188,level=17].per_run :: sha256 1250d1f00537e0bd :: nonlinear minus linear probe, 188 native labels from the 17-class models
\newcommand{\RecoveryLastGainOneseven}{0.0125\,\ensuremath{\pm}\,0.0021}  % experiments/FIGS/data/label_recovery_curve_v1err/analysis/label_recovery_curve_summary.json :: data_gain_last_step[rung=L188,level=17] :: sha256 1250d1f00537e0bd :: balanced accuracy gained from 443000 to 1400000 training jets
\newcommand{\RecoveryLastGainOneeighteight}{0.0147\,\ensuremath{\pm}\,0.0014}  % experiments/FIGS/data/label_recovery_curve_v1err/analysis/label_recovery_curve_summary.json :: data_gain_last_step[rung=L188,level=188] :: sha256 1250d1f00537e0bd :: balanced accuracy gained from 443000 to 1400000 training jets
\newcommand{\RecoveryNTrain}{1{,}400{,}000}  % experiments/FIGS/data/label_recovery_curve_v1err/analysis/label_recovery_curve_summary.json :: sizes (max) :: sha256 1250d1f00537e0bd :: jets the linear probe is fit on at the end of the curve
\newcommand{\RecoveryNTrainMin}{14{,}000}  % experiments/FIGS/data/label_recovery_curve_v1err/analysis/label_recovery_curve_summary.json :: sizes (min) :: sha256 1250d1f00537e0bd
\newcommand{\RecoveryNTrainPrev}{443{,}000}  % experiments/FIGS/data/label_recovery_curve_v1err/analysis/label_recovery_curve_summary.json :: sizes (second largest) :: sha256 1250d1f00537e0bd
\newcommand{\RecoveryNTest}{400{,}000}  % experiments/FIGS/data/label_recovery_curve_v1err/l188-s1/label_recovery_curve.json :: n_test :: sha256 65651e4c06c7766f :: jets the probes are scored on, recorded in all 20 per-run files
\newcommand{\RecoveryMlpHidden}{512}  % experiments/FIGS/data/label_recovery_curve_v1err/l188-s1/label_recovery_curve.json :: mlp.hidden :: sha256 65651e4c06c7766f :: hidden units of the nonlinear probe
\newcommand{\RandNDraws}{three}  % experiments/FIGS/data/probe_ladder_randcontrol_mlp2/analysis/c4_random_control.json :: table[*].draw :: sha256 c675e609ff4fe3c9 :: random partitions drawn
\newcommand{\RandNDrawsCap}{Three}  % experiments/FIGS/data/probe_ladder_randcontrol_mlp2/analysis/c4_random_control.json :: table[*].draw :: sha256 c675e609ff4fe3c9 :: random partitions drawn, capitalised to open a sentence
\newcommand{\RandOmaBvcFourprongDrawOneLinear}{\ensuremath{1.05\times10^{-2}}}  % experiments/FIGS/data/probe_ladder_randcontrol_mlp2/analysis/c4_random_control.json :: table[0].control_log1m_auc :: sha256 c675e609ff4fe3c9 :: 1-AUC = exp of it, one run
\newcommand{\RandOmaBvcFourprongDrawTwoLinear}{\ensuremath{1.11\times10^{-2}}}  % experiments/FIGS/data/probe_ladder_randcontrol_mlp2/analysis/c4_random_control.json :: table[4].control_log1m_auc :: sha256 c675e609ff4fe3c9 :: 1-AUC = exp of it, one run
\newcommand{\RandOmaBvcFourprongDrawThreeLinear}{\ensuremath{9.98\times10^{-3}}}  % experiments/FIGS/data/probe_ladder_randcontrol_mlp2/analysis/c4_random_control.json :: table[8].control_log1m_auc :: sha256 c675e609ff4fe3c9 :: 1-AUC = exp of it, one run
\newcommand{\RandOmaMeanBvcFourprongLinear}{\ensuremath{(1.05 \pm 0.06)\times10^{-2}}}  % experiments/FIGS/data/probe_ladder_randcontrol_mlp2/analysis/c4_random_control.json :: table[task=bvc_4prong,probe=linear].control_log1m_auc :: sha256 c675e609ff4fe3c9 :: 1-AUC, mean +- SD over 3 draws
\newcommand{\RandOmaBvcFourprongDrawOneMlp}{\ensuremath{1.05\times10^{-2}}}  % experiments/FIGS/data/probe_ladder_randcontrol_mlp2/analysis/c4_random_control.json :: table[1].control_log1m_auc :: sha256 c675e609ff4fe3c9 :: 1-AUC = exp of it, one run
\newcommand{\RandOmaBvcFourprongDrawTwoMlp}{\ensuremath{1.04\times10^{-2}}}  % experiments/FIGS/data/probe_ladder_randcontrol_mlp2/analysis/c4_random_control.json :: table[5].control_log1m_auc :: sha256 c675e609ff4fe3c9 :: 1-AUC = exp of it, one run
\newcommand{\RandOmaBvcFourprongDrawThreeMlp}{\ensuremath{9.28\times10^{-3}}}  % experiments/FIGS/data/probe_ladder_randcontrol_mlp2/analysis/c4_random_control.json :: table[9].control_log1m_auc :: sha256 c675e609ff4fe3c9 :: 1-AUC = exp of it, one run
\newcommand{\RandOmaMeanBvcFourprongMlp}{\ensuremath{(1.01 \pm 0.07)\times10^{-2}}}  % experiments/FIGS/data/probe_ladder_randcontrol_mlp2/analysis/c4_random_control.json :: table[task=bvc_4prong,probe=mlp].control_log1m_auc :: sha256 c675e609ff4fe3c9 :: 1-AUC, mean +- SD over 3 draws
\newcommand{\RandOmaVisibleContentDrawOneLinear}{\ensuremath{1.25\times10^{-3}}}  % experiments/FIGS/data/probe_ladder_randcontrol_mlp2/analysis/c4_random_control.json :: table[2].control_log1m_auc :: sha256 c675e609ff4fe3c9 :: 1-AUC = exp of it, one run
\newcommand{\RandOmaVisibleContentDrawTwoLinear}{\ensuremath{9.89\times10^{-4}}}  % experiments/FIGS/data/probe_ladder_randcontrol_mlp2/analysis/c4_random_control.json :: table[6].control_log1m_auc :: sha256 c675e609ff4fe3c9 :: 1-AUC = exp of it, one run
\newcommand{\RandOmaVisibleContentDrawThreeLinear}{\ensuremath{1.08\times10^{-3}}}  % experiments/FIGS/data/probe_ladder_randcontrol_mlp2/analysis/c4_random_control.json :: table[10].control_log1m_auc :: sha256 c675e609ff4fe3c9 :: 1-AUC = exp of it, one run
\newcommand{\RandOmaMeanVisibleContentLinear}{\ensuremath{(1.11 \pm 0.13)\times10^{-3}}}  % experiments/FIGS/data/probe_ladder_randcontrol_mlp2/analysis/c4_random_control.json :: table[task=visible_content,probe=linear].control_log1m_auc :: sha256 c675e609ff4fe3c9 :: 1-AUC, mean +- SD over 3 draws
\newcommand{\RandOmaVisibleContentDrawOneMlp}{\ensuremath{1.21\times10^{-3}}}  % experiments/FIGS/data/probe_ladder_randcontrol_mlp2/analysis/c4_random_control.json :: table[3].control_log1m_auc :: sha256 c675e609ff4fe3c9 :: 1-AUC = exp of it, one run
\newcommand{\RandOmaVisibleContentDrawTwoMlp}{\ensuremath{1.07\times10^{-3}}}  % experiments/FIGS/data/probe_ladder_randcontrol_mlp2/analysis/c4_random_control.json :: table[7].control_log1m_auc :: sha256 c675e609ff4fe3c9 :: 1-AUC = exp of it, one run
\newcommand{\RandOmaVisibleContentDrawThreeMlp}{\ensuremath{1.13\times10^{-3}}}  % experiments/FIGS/data/probe_ladder_randcontrol_mlp2/analysis/c4_random_control.json :: table[11].control_log1m_auc :: sha256 c675e609ff4fe3c9 :: 1-AUC = exp of it, one run
\newcommand{\RandOmaMeanVisibleContentMlp}{\ensuremath{(1.14 \pm 0.07)\times10^{-3}}}  % experiments/FIGS/data/probe_ladder_randcontrol_mlp2/analysis/c4_random_control.json :: table[task=visible_content,probe=mlp].control_log1m_auc :: sha256 c675e609ff4fe3c9 :: 1-AUC, mean +- SD over 3 draws
\newcommand{\RandCostFactorBvcFourprongLinear}{1.49}  % experiments/FIGS/data/probe_ladder_randcontrol_mlp2/analysis/c4_random_control.json :: grouping_cost_post_hoc.tasks.bvc_4prong.linear.factor_control :: sha256 c675e609ff4fe3c9 :: POST HOC: random grouping vs finer models, in 1-AUC
\newcommand{\RandSemCostFactorBvcFourprongLinear}{4.42}  % experiments/FIGS/data/probe_ladder_randcontrol_mlp2/analysis/c4_random_control.json :: grouping_cost_post_hoc.tasks.bvc_4prong.linear.factor_semantic :: sha256 c675e609ff4fe3c9 :: POST HOC: semantic 17-class vs finer models
\newcommand{\RandCostFactorBvcFourprongMlp}{1.41}  % experiments/FIGS/data/probe_ladder_randcontrol_mlp2/analysis/c4_random_control.json :: grouping_cost_post_hoc.tasks.bvc_4prong.mlp.factor_control :: sha256 c675e609ff4fe3c9 :: POST HOC: random grouping vs finer models, in 1-AUC
\newcommand{\RandSemCostFactorBvcFourprongMlp}{3.26}  % experiments/FIGS/data/probe_ladder_randcontrol_mlp2/analysis/c4_random_control.json :: grouping_cost_post_hoc.tasks.bvc_4prong.mlp.factor_semantic :: sha256 c675e609ff4fe3c9 :: POST HOC: semantic 17-class vs finer models
\newcommand{\RandCostFactorVisibleContentLinear}{1.35}  % experiments/FIGS/data/probe_ladder_randcontrol_mlp2/analysis/c4_random_control.json :: grouping_cost_post_hoc.tasks.visible_content.linear.factor_control :: sha256 c675e609ff4fe3c9 :: POST HOC: random grouping vs finer models, in 1-AUC
\newcommand{\RandSemCostFactorVisibleContentLinear}{2.82}  % experiments/FIGS/data/probe_ladder_randcontrol_mlp2/analysis/c4_random_control.json :: grouping_cost_post_hoc.tasks.visible_content.linear.factor_semantic :: sha256 c675e609ff4fe3c9 :: POST HOC: semantic 17-class vs finer models
\newcommand{\RandCostFactorVisibleContentMlp}{1.33}  % experiments/FIGS/data/probe_ladder_randcontrol_mlp2/analysis/c4_random_control.json :: grouping_cost_post_hoc.tasks.visible_content.mlp.factor_control :: sha256 c675e609ff4fe3c9 :: POST HOC: random grouping vs finer models, in 1-AUC
\newcommand{\RandSemCostFactorVisibleContentMlp}{2.37}  % experiments/FIGS/data/probe_ladder_randcontrol_mlp2/analysis/c4_random_control.json :: grouping_cost_post_hoc.tasks.visible_content.mlp.factor_semantic :: sha256 c675e609ff4fe3c9 :: POST HOC: semantic 17-class vs finer models
\newcommand{\MassResNTest}{398{,}625}  % experiments/FIGS/data/mass_resolution/s1.json :: centering_detail.split[2] :: sha256 88a57bbad270d51c :: test jets sigma_eff is computed on, the same in every per-seed file and arm
\newcommand{\MassResNClasses}{184}  % experiments/FIGS/data/mass_resolution/analysis_holm/s7_mass_resolution.json :: provenance.n_classes_used :: sha256 e03cde45ece9f385 :: native classes with enough training jets to centre
\newcommand{\MassResSigmaEffRidgeOneeighteight}{0.3393\,\ensuremath{\pm}\,0.0019}  % experiments/FIGS/data/mass_resolution/analysis_holm/s7_mass_resolution.json :: table[cell=188,probe=ridge].sigma_eff :: sha256 e03cde45ece9f385 :: mean +- SD over 5 seeds
\newcommand{\MassResSigmaEffRidgeOnesixtwo}{0.3302\,\ensuremath{\pm}\,0.0029}  % experiments/FIGS/data/mass_resolution/analysis_holm/s7_mass_resolution.json :: table[cell=162,probe=ridge].sigma_eff :: sha256 e03cde45ece9f385 :: mean +- SD over 5 seeds
\newcommand{\MassResSigmaEffRidgeFourthree}{0.3228\,\ensuremath{\pm}\,0.0019}  % experiments/FIGS/data/mass_resolution/analysis_holm/s7_mass_resolution.json :: table[cell=43,probe=ridge].sigma_eff :: sha256 e03cde45ece9f385 :: mean +- SD over 5 seeds
\newcommand{\MassResSigmaEffRidgeOneseven}{0.300\,\ensuremath{\pm}\,0.010}  % experiments/FIGS/data/mass_resolution/analysis_holm/s7_mass_resolution.json :: table[cell=17,probe=ridge].sigma_eff :: sha256 e03cde45ece9f385 :: mean +- SD over 5 seeds
\newcommand{\MassResSigmaEffRidgeOnesixtwoMass}{0.2238\,\ensuremath{\pm}\,0.0019}  % experiments/FIGS/data/mass_resolution/analysis_holm/s7_mass_resolution.json :: table[cell=162+mass,probe=ridge].sigma_eff :: sha256 e03cde45ece9f385 :: mean +- SD over 5 seeds
\newcommand{\MassResSigmaEffRidgeOnesevenMass}{0.2170\,\ensuremath{\pm}\,0.0034}  % experiments/FIGS/data/mass_resolution/analysis_holm/s7_mass_resolution.json :: table[cell=17+mass,probe=ridge].sigma_eff :: sha256 e03cde45ece9f385 :: mean +- SD over 5 seeds
\newcommand{\MassResGainRidgeOnesixtwo}{\ensuremath{-}0.1063}  % experiments/FIGS/data/mass_resolution/analysis_holm/s7_mass_resolution.json :: table[cell=162+mass/162,probe=ridge].sigma_eff :: sha256 e03cde45ece9f385 :: sigma_eff with the mass output minus without, difference of the seed means; negative is better
\newcommand{\MassResGainRidgeOneseven}{\ensuremath{-}0.0831}  % experiments/FIGS/data/mass_resolution/analysis_holm/s7_mass_resolution.json :: table[cell=17+mass/17,probe=ridge].sigma_eff :: sha256 e03cde45ece9f385 :: sigma_eff with the mass output minus without, difference of the seed means; negative is better
\newcommand{\MassResDidRidge}{\ensuremath{-}0.0233}  % experiments/FIGS/data/mass_resolution/analysis_holm/s7_mass_resolution.json :: table[cell=162+mass/162/17+mass/17,probe=ridge].sigma_eff :: sha256 e03cde45ece9f385 :: gain at 162 classes minus gain at 17
\newcommand{\MassResSigmaEffMlpOneeighteight}{0.2463\,\ensuremath{\pm}\,0.0022}  % experiments/FIGS/data/mass_resolution/analysis_holm/s7_mass_resolution.json :: table[cell=188,probe=mlp].sigma_eff :: sha256 e03cde45ece9f385 :: mean +- SD over 5 seeds
\newcommand{\MassResSigmaEffMlpOnesixtwo}{0.2438\,\ensuremath{\pm}\,0.0018}  % experiments/FIGS/data/mass_resolution/analysis_holm/s7_mass_resolution.json :: table[cell=162,probe=mlp].sigma_eff :: sha256 e03cde45ece9f385 :: mean +- SD over 5 seeds
\newcommand{\MassResSigmaEffMlpFourthree}{0.2407\,\ensuremath{\pm}\,0.0009}  % experiments/FIGS/data/mass_resolution/analysis_holm/s7_mass_resolution.json :: table[cell=43,probe=mlp].sigma_eff :: sha256 e03cde45ece9f385 :: mean +- SD over 5 seeds
\newcommand{\MassResSigmaEffMlpOneseven}{0.2358\,\ensuremath{\pm}\,0.0025}  % experiments/FIGS/data/mass_resolution/analysis_holm/s7_mass_resolution.json :: table[cell=17,probe=mlp].sigma_eff :: sha256 e03cde45ece9f385 :: mean +- SD over 5 seeds
\newcommand{\MassResSigmaEffMlpOnesixtwoMass}{0.1943\,\ensuremath{\pm}\,0.0001}  % experiments/FIGS/data/mass_resolution/analysis_holm/s7_mass_resolution.json :: table[cell=162+mass,probe=mlp].sigma_eff :: sha256 e03cde45ece9f385 :: mean +- SD over 5 seeds
\newcommand{\MassResSigmaEffMlpOnesevenMass}{0.1949\,\ensuremath{\pm}\,0.0006}  % experiments/FIGS/data/mass_resolution/analysis_holm/s7_mass_resolution.json :: table[cell=17+mass,probe=mlp].sigma_eff :: sha256 e03cde45ece9f385 :: mean +- SD over 5 seeds
\newcommand{\MassResGainMlpOnesixtwo}{\ensuremath{-}0.0496}  % experiments/FIGS/data/mass_resolution/analysis_holm/s7_mass_resolution.json :: table[cell=162+mass/162,probe=mlp].sigma_eff :: sha256 e03cde45ece9f385 :: sigma_eff with the mass output minus without, difference of the seed means; negative is better
\newcommand{\MassResGainMlpOneseven}{\ensuremath{-}0.0409}  % experiments/FIGS/data/mass_resolution/analysis_holm/s7_mass_resolution.json :: table[cell=17+mass/17,probe=mlp].sigma_eff :: sha256 e03cde45ece9f385 :: sigma_eff with the mass output minus without, difference of the seed means; negative is better
\newcommand{\MassResDidMlp}{\ensuremath{-}0.0087}  % experiments/FIGS/data/mass_resolution/analysis_holm/s7_mass_resolution.json :: table[cell=162+mass/162/17+mass/17,probe=mlp].sigma_eff :: sha256 e03cde45ece9f385 :: gain at 162 classes minus gain at 17
\newcommand{\MassResTargetSigmaEff}{0.2986}  % experiments/FIGS/data/mass_resolution/analysis_holm/s7_mass_resolution.json :: table[*].target_sigma_eff :: sha256 e03cde45ece9f385 :: sigma_eff of the class-centred target itself (a model that predicts the class mean)
\newcommand{\MassResRidgeRtwoMin}{20\%}  % experiments/FIGS/data/mass_resolution/s1.json :: arms.{l188,l162,r42q1,r16q1}-s*.ridge.val_r2 :: sha256 88a57bbad270d51c :: validation R^2 of the ridge probe, smallest of 20 models without the mass output, all 5 per-seed files
\newcommand{\MassResRidgeRtwoMax}{30\%}  % experiments/FIGS/data/mass_resolution/s1.json :: arms.{l188,l162,r42q1,r16q1}-s*.ridge.val_r2 :: sha256 88a57bbad270d51c :: validation R^2 of the ridge probe, largest of 20 models without the mass output, all 5 per-seed files
\newcommand{\MassResSigmaEffOnesixtwo}{0.3302}  % experiments/FIGS/data/mass_resolution/analysis_holm/s7_mass_resolution.json :: table[cell=162,probe=ridge].sigma_eff :: sha256 e03cde45ece9f385 :: mean over 5 seeds
\newcommand{\MassResSdOnesixtwo}{0.626}  % experiments/FIGS/data/mass_resolution/analysis_holm/s7_mass_resolution.json :: table[cell=162,probe=ridge].sd :: sha256 e03cde45ece9f385 :: mean over 5 seeds
\newcommand{\MassResTailFractionOnesixtwo}{0.062}  % experiments/FIGS/data/mass_resolution/analysis_holm/s7_mass_resolution.json :: table[cell=162,probe=ridge].tail_fraction :: sha256 e03cde45ece9f385 :: mean over 5 seeds
\newcommand{\MassResSigmaEffOnesixtwoMass}{0.2238}  % experiments/FIGS/data/mass_resolution/analysis_holm/s7_mass_resolution.json :: table[cell=162+mass,probe=ridge].sigma_eff :: sha256 e03cde45ece9f385 :: mean over 5 seeds
\newcommand{\MassResSdOnesixtwoMass}{0.494}  % experiments/FIGS/data/mass_resolution/analysis_holm/s7_mass_resolution.json :: table[cell=162+mass,probe=ridge].sd :: sha256 e03cde45ece9f385 :: mean over 5 seeds
\newcommand{\MassResTailFractionOnesixtwoMass}{0.034}  % experiments/FIGS/data/mass_resolution/analysis_holm/s7_mass_resolution.json :: table[cell=162+mass,probe=ridge].tail_fraction :: sha256 e03cde45ece9f385 :: mean over 5 seeds
\newcommand{\MassResSigmaEffOneseven}{0.3001}  % experiments/FIGS/data/mass_resolution/analysis_holm/s7_mass_resolution.json :: table[cell=17,probe=ridge].sigma_eff :: sha256 e03cde45ece9f385 :: mean over 5 seeds
\newcommand{\MassResSdOneseven}{0.601}  % experiments/FIGS/data/mass_resolution/analysis_holm/s7_mass_resolution.json :: table[cell=17,probe=ridge].sd :: sha256 e03cde45ece9f385 :: mean over 5 seeds
\newcommand{\MassResTailFractionOneseven}{0.056}  % experiments/FIGS/data/mass_resolution/analysis_holm/s7_mass_resolution.json :: table[cell=17,probe=ridge].tail_fraction :: sha256 e03cde45ece9f385 :: mean over 5 seeds
\newcommand{\MassResSigmaEffOnesevenMass}{0.2170}  % experiments/FIGS/data/mass_resolution/analysis_holm/s7_mass_resolution.json :: table[cell=17+mass,probe=ridge].sigma_eff :: sha256 e03cde45ece9f385 :: mean over 5 seeds
\newcommand{\MassResSdOnesevenMass}{0.490}  % experiments/FIGS/data/mass_resolution/analysis_holm/s7_mass_resolution.json :: table[cell=17+mass,probe=ridge].sd :: sha256 e03cde45ece9f385 :: mean over 5 seeds
\newcommand{\MassResTailFractionOnesevenMass}{0.033}  % experiments/FIGS/data/mass_resolution/analysis_holm/s7_mass_resolution.json :: table[cell=17+mass,probe=ridge].tail_fraction :: sha256 e03cde45ece9f385 :: mean over 5 seeds
\newcommand{\MassResSigmaEffOneeighteight}{0.3393}  % experiments/FIGS/data/mass_resolution/analysis_holm/s7_mass_resolution.json :: table[cell=188,probe=ridge].sigma_eff :: sha256 e03cde45ece9f385 :: mean over 5 seeds
\newcommand{\MassResSdOneeighteight}{0.631}  % experiments/FIGS/data/mass_resolution/analysis_holm/s7_mass_resolution.json :: table[cell=188,probe=ridge].sd :: sha256 e03cde45ece9f385 :: mean over 5 seeds
\newcommand{\MassResTailFractionOneeighteight}{0.064}  % experiments/FIGS/data/mass_resolution/analysis_holm/s7_mass_resolution.json :: table[cell=188,probe=ridge].tail_fraction :: sha256 e03cde45ece9f385 :: mean over 5 seeds
\newcommand{\MassResSigmaEffFourthree}{0.3228}  % experiments/FIGS/data/mass_resolution/analysis_holm/s7_mass_resolution.json :: table[cell=43,probe=ridge].sigma_eff :: sha256 e03cde45ece9f385 :: mean over 5 seeds
\newcommand{\MassResSdFourthree}{0.618}  % experiments/FIGS/data/mass_resolution/analysis_holm/s7_mass_resolution.json :: table[cell=43,probe=ridge].sd :: sha256 e03cde45ece9f385 :: mean over 5 seeds
\newcommand{\MassResTailFractionFourthree}{0.060}  % experiments/FIGS/data/mass_resolution/analysis_holm/s7_mass_resolution.json :: table[cell=43,probe=ridge].tail_fraction :: sha256 e03cde45ece9f385 :: mean over 5 seeds
\newcommand{\AojEffPercent}{1\%}  % experiments/FIGS/data/aoj_full_v1/fit_v6/results.json :: eff :: sha256 2416fc21f93fcfa6 :: the score cut's pass fraction outside the mass windows
\newcommand{\AojPtMin}{500}  % experiments/AOJ/peak_fit.py :: PT_RANGE[0] :: sha256 cad06229a74532be :: GeV
\newcommand{\AojPtMax}{2{,}500}  % experiments/AOJ/peak_fit.py :: PT_RANGE[1] :: sha256 cad06229a74532be :: GeV
\newcommand{\AojRhoMin}{\ensuremath{-}5.5}  % experiments/AOJ/peak_fit.py :: RHO_RANGE[0] :: sha256 cad06229a74532be
\newcommand{\AojRhoMax}{\ensuremath{-}2.0}  % experiments/AOJ/peak_fit.py :: RHO_RANGE[1] :: sha256 cad06229a74532be
\newcommand{\AojWindowWLo}{65}  % experiments/AOJ/peak_fit.py :: PEAKS[W].window[0] :: sha256 cad06229a74532be :: GeV
\newcommand{\AojWindowWHi}{105}  % experiments/AOJ/peak_fit.py :: PEAKS[W].window[1] :: sha256 cad06229a74532be :: GeV
\newcommand{\AojWindowTopLo}{140}  % experiments/AOJ/peak_fit.py :: PEAKS[top].window[0] :: sha256 cad06229a74532be :: GeV
\newcommand{\AojWindowTopHi}{220}  % experiments/AOJ/peak_fit.py :: PEAKS[top].window[1] :: sha256 cad06229a74532be :: GeV
\newcommand{\AojFitLo}{105}  % experiments/AOJ/peak_fit.py :: PEAKS[top].fit_range[0] :: sha256 cad06229a74532be :: GeV
\newcommand{\AojFitHi}{300}  % experiments/AOJ/peak_fit.py :: PEAKS[top].fit_range[1] :: sha256 cad06229a74532be :: GeV
\newcommand{\AojMassBin}{5}  % experiments/AOJ/peak_fit.py :: MASS_BIN :: sha256 cad06229a74532be :: GeV
\newcommand{\AojMapOrderRho}{3}  % experiments/AOJ/peak_fit.py :: MAP_ORDER[0] :: sha256 cad06229a74532be :: polynomial order in rho
\newcommand{\AojMapOrderPt}{2}  % experiments/AOJ/peak_fit.py :: MAP_ORDER[1] :: sha256 cad06229a74532be :: polynomial order in ln pT
\newcommand{\AojNRunsWord}{five}  % experiments/FIGS/data/aoj_full_v1/analysis_v6/aoj_top.json :: per_label_set.label_sets.*.signal_yield.n :: sha256 5dfef902a5fe75ba :: pretraining runs per vocabulary
\newcommand{\AojNJetsRhoWindow}{9{,}441{,}209}  % experiments/FIGS/data/aoj_full_v1/fit_v6/results.json :: n_jets :: sha256 2416fc21f93fcfa6 :: staged jets inside the rho window, all files; not the jets in the mass fit
\newcommand{\AojYieldOneeighteight}{2{,}400\,\ensuremath{\pm}\,500}  % experiments/FIGS/data/aoj_full_v1/analysis_v6/aoj_top.json :: per_label_set.label_sets.188.signal_yield :: sha256 5dfef902a5fe75ba :: mean +- SD over 5 runs, shared shape
\newcommand{\AojStatErrOneeighteight}{211}  % experiments/FIGS/data/aoj_full_v1/analysis_v6/aoj_top.json :: per_label_set.label_sets.188.median_stat_err :: sha256 5dfef902a5fe75ba :: median per-fit statistical error
\newcommand{\AojYieldOwnShapeOneeighteight}{2{,}400\,\ensuremath{\pm}\,600}  % experiments/FIGS/data/aoj_full_v1/analysis_v6/aoj_top.json :: per_label_set.label_sets.188.by_shape.own_shape :: sha256 5dfef902a5fe75ba :: every fit with its own floated shape
\newcommand{\AojSpreadOverStatOneeighteight}{2.4}  % experiments/FIGS/data/aoj_full_v1/analysis_v6/aoj_top.json :: per_label_set.label_sets.188.signal_yield.sd / median_stat_err :: sha256 5dfef902a5fe75ba :: spread between runs over the median per-fit statistical error
\newcommand{\AojYieldOnesixtwo}{2{,}400\,\ensuremath{\pm}\,900}  % experiments/FIGS/data/aoj_full_v1/analysis_v6/aoj_top.json :: per_label_set.label_sets.162.signal_yield :: sha256 5dfef902a5fe75ba :: mean +- SD over 5 runs, shared shape
\newcommand{\AojStatErrOnesixtwo}{199}  % experiments/FIGS/data/aoj_full_v1/analysis_v6/aoj_top.json :: per_label_set.label_sets.162.median_stat_err :: sha256 5dfef902a5fe75ba :: median per-fit statistical error
\newcommand{\AojYieldOwnShapeOnesixtwo}{2{,}200\,\ensuremath{\pm}\,800}  % experiments/FIGS/data/aoj_full_v1/analysis_v6/aoj_top.json :: per_label_set.label_sets.162.by_shape.own_shape :: sha256 5dfef902a5fe75ba :: every fit with its own floated shape
\newcommand{\AojSpreadOverStatOnesixtwo}{4.5}  % experiments/FIGS/data/aoj_full_v1/analysis_v6/aoj_top.json :: per_label_set.label_sets.162.signal_yield.sd / median_stat_err :: sha256 5dfef902a5fe75ba :: spread between runs over the median per-fit statistical error
\newcommand{\AojYieldFourthree}{2{,}400\,\ensuremath{\pm}\,800}  % experiments/FIGS/data/aoj_full_v1/analysis_v6/aoj_top.json :: per_label_set.label_sets.43.signal_yield :: sha256 5dfef902a5fe75ba :: mean +- SD over 5 runs, shared shape
\newcommand{\AojStatErrFourthree}{184}  % experiments/FIGS/data/aoj_full_v1/analysis_v6/aoj_top.json :: per_label_set.label_sets.43.median_stat_err :: sha256 5dfef902a5fe75ba :: median per-fit statistical error
\newcommand{\AojYieldOwnShapeFourthree}{2{,}300\,\ensuremath{\pm}\,800}  % experiments/FIGS/data/aoj_full_v1/analysis_v6/aoj_top.json :: per_label_set.label_sets.43.by_shape.own_shape :: sha256 5dfef902a5fe75ba :: every fit with its own floated shape
\newcommand{\AojSpreadOverStatFourthree}{4.2}  % experiments/FIGS/data/aoj_full_v1/analysis_v6/aoj_top.json :: per_label_set.label_sets.43.signal_yield.sd / median_stat_err :: sha256 5dfef902a5fe75ba :: spread between runs over the median per-fit statistical error
\newcommand{\AojYieldOneseven}{2{,}350\,\ensuremath{\pm}\,240}  % experiments/FIGS/data/aoj_full_v1/analysis_v6/aoj_top.json :: per_label_set.label_sets.17.signal_yield :: sha256 5dfef902a5fe75ba :: mean +- SD over 5 runs, shared shape
\newcommand{\AojStatErrOneseven}{201}  % experiments/FIGS/data/aoj_full_v1/analysis_v6/aoj_top.json :: per_label_set.label_sets.17.median_stat_err :: sha256 5dfef902a5fe75ba :: median per-fit statistical error
\newcommand{\AojYieldOwnShapeOneseven}{2{,}480\,\ensuremath{\pm}\,290}  % experiments/FIGS/data/aoj_full_v1/analysis_v6/aoj_top.json :: per_label_set.label_sets.17.by_shape.own_shape :: sha256 5dfef902a5fe75ba :: every fit with its own floated shape
\newcommand{\AojSpreadOverStatOneseven}{1.2}  % experiments/FIGS/data/aoj_full_v1/analysis_v6/aoj_top.json :: per_label_set.label_sets.17.signal_yield.sd / median_stat_err :: sha256 5dfef902a5fe75ba :: spread between runs over the median per-fit statistical error
\newcommand{\AojYieldOnesixtwoMass}{2{,}420\,\ensuremath{\pm}\,240}  % experiments/FIGS/data/aoj_full_v1/analysis_v6/aoj_top.json :: per_label_set.label_sets.162+mass.signal_yield :: sha256 5dfef902a5fe75ba :: mean +- SD over 5 runs, shared shape
\newcommand{\AojStatErrOnesixtwoMass}{179}  % experiments/FIGS/data/aoj_full_v1/analysis_v6/aoj_top.json :: per_label_set.label_sets.162+mass.median_stat_err :: sha256 5dfef902a5fe75ba :: median per-fit statistical error
\newcommand{\AojYieldOwnShapeOnesixtwoMass}{2{,}350\,\ensuremath{\pm}\,320}  % experiments/FIGS/data/aoj_full_v1/analysis_v6/aoj_top.json :: per_label_set.label_sets.162+mass.by_shape.own_shape :: sha256 5dfef902a5fe75ba :: every fit with its own floated shape
\newcommand{\AojSpreadOverStatOnesixtwoMass}{1.3}  % experiments/FIGS/data/aoj_full_v1/analysis_v6/aoj_top.json :: per_label_set.label_sets.162+mass.signal_yield.sd / median_stat_err :: sha256 5dfef902a5fe75ba :: spread between runs over the median per-fit statistical error
\newcommand{\AojYieldOnesevenMass}{2{,}500\,\ensuremath{\pm}\,600}  % experiments/FIGS/data/aoj_full_v1/analysis_v6/aoj_top.json :: per_label_set.label_sets.17+mass.signal_yield :: sha256 5dfef902a5fe75ba :: mean +- SD over 5 runs, shared shape
\newcommand{\AojStatErrOnesevenMass}{200}  % experiments/FIGS/data/aoj_full_v1/analysis_v6/aoj_top.json :: per_label_set.label_sets.17+mass.median_stat_err :: sha256 5dfef902a5fe75ba :: median per-fit statistical error
\newcommand{\AojYieldOwnShapeOnesevenMass}{2{,}900\,\ensuremath{\pm}\,500}  % experiments/FIGS/data/aoj_full_v1/analysis_v6/aoj_top.json :: per_label_set.label_sets.17+mass.by_shape.own_shape :: sha256 5dfef902a5fe75ba :: every fit with its own floated shape
\newcommand{\AojSpreadOverStatOnesevenMass}{2.9}  % experiments/FIGS/data/aoj_full_v1/analysis_v6/aoj_top.json :: per_label_set.label_sets.17+mass.signal_yield.sd / median_stat_err :: sha256 5dfef902a5fe75ba :: spread between runs over the median per-fit statistical error
\newcommand{\AojOwnShapeShiftMin}{\ensuremath{-}7\%}  % experiments/FIGS/data/aoj_full_v1/analysis_v6/aoj_top.json :: per_label_set.label_sets.*.by_shape.own_shape.mean / signal_yield.mean - 1 :: sha256 5dfef902a5fe75ba :: the vocabulary means with each model's own shape, relative to the shared shape
\newcommand{\AojOwnShapeShiftMax}{+13\%}  % experiments/FIGS/data/aoj_full_v1/analysis_v6/aoj_top.json :: per_label_set.label_sets.*.by_shape.own_shape.mean / signal_yield.mean - 1 :: sha256 5dfef902a5fe75ba :: the vocabulary means with each model's own shape, relative to the shared shape
\newcommand{\AojPooledMean}{182.8}  % experiments/FIGS/data/aoj_full_v1/analysis_v6/aoj_top.json :: per_label_set.shapes.pooled :: sha256 5dfef902a5fe75ba :: GeV
\newcommand{\AojPooledWidth}{11.7}  % experiments/FIGS/data/aoj_full_v1/analysis_v6/aoj_top.json :: per_label_set.shapes.pooled :: sha256 5dfef902a5fe75ba :: GeV
\newcommand{\AojPoolSize}{30}  % experiments/FIGS/data/aoj_full_v1/analysis_v6/aoj_top.json :: per_label_set.shapes.pool (count) :: sha256 5dfef902a5fe75ba :: pretrained models the shared shape is fitted over
\newcommand{\AojNModels}{30}  % experiments/FIGS/data/aoj_full_v1/fit_v6/results.json :: models (pretrained) :: sha256 2416fc21f93fcfa6 :: fits
\newcommand{\AojOwnMeanMin}{181.6}  % experiments/FIGS/data/aoj_full_v1/fit_v6/results.json :: models.*.top.own_shape_fit.mean (min) :: sha256 2416fc21f93fcfa6 :: GeV
\newcommand{\AojOwnMeanMax}{184.0}  % experiments/FIGS/data/aoj_full_v1/fit_v6/results.json :: models.*.top.own_shape_fit.mean (max) :: sha256 2416fc21f93fcfa6 :: GeV
\newcommand{\AojOwnWidthMin}{10.0}  % experiments/FIGS/data/aoj_full_v1/fit_v6/results.json :: models.*.top.own_shape_fit.width (min) :: sha256 2416fc21f93fcfa6 :: GeV
\newcommand{\AojOwnWidthMax}{13.8}  % experiments/FIGS/data/aoj_full_v1/fit_v6/results.json :: models.*.top.own_shape_fit.width (max) :: sha256 2416fc21f93fcfa6 :: GeV
\newcommand{\AojShapeStepMean}{0.7}  % experiments/FIGS/data/aoj_full_v1/fit_v6/results.json :: pooled_shape.systematic_step.mean :: sha256 2416fc21f93fcfa6 :: GeV: the SD of the models' own floated values, one shift for every model together
\newcommand{\AojShapeSystMeanMax}{1\%}  % experiments/FIGS/data/aoj_full_v1/fit_v6/results.json :: models.*.top.shape_systematic.pooled_mean_(down|up) / signal_yield - 1 (largest |.|) :: sha256 2416fc21f93fcfa6 :: the largest yield change of any model under that shift
\newcommand{\AojShapeStepWidth}{1.0}  % experiments/FIGS/data/aoj_full_v1/fit_v6/results.json :: pooled_shape.systematic_step.width :: sha256 2416fc21f93fcfa6 :: GeV: the SD of the models' own floated values, one shift for every model together
\newcommand{\AojShapeSystWidthMax}{12\%}  % experiments/FIGS/data/aoj_full_v1/fit_v6/results.json :: models.*.top.shape_systematic.pooled_width_(down|up) / signal_yield - 1 (largest |.|) :: sha256 2416fc21f93fcfa6 :: the largest yield change of any model under that shift
\newcommand{\AojLeakEpsHi}{0.6}  % experiments/FIGS/data/aoj_full_v1/fit_v6/results.json :: models.*.top.leak_systematic.eps_ref (max) :: sha256 2416fc21f93fcfa6 :: the reference's assumed efficiency for data tops, upper variation
\newcommand{\AojLeakEpsLo}{0.4}  % experiments/FIGS/data/aoj_full_v1/fit_v6/results.json :: models.*.top.leak_systematic.eps_ref (min) :: sha256 2416fc21f93fcfa6 :: the same, lower variation
\newcommand{\AojLeakShiftMin}{+97}  % experiments/FIGS/data/aoj_full_v1/fit_v6/results.json :: models.*.top.leak_systematic.shift['0.6'] (min) :: sha256 2416fc21f93fcfa6 :: tops added to a yield
\newcommand{\AojLeakShiftMax}{+271}  % experiments/FIGS/data/aoj_full_v1/fit_v6/results.json :: models.*.top.leak_systematic.shift['0.4'] (max) :: sha256 2416fc21f93fcfa6 :: tops added to a yield
\newcommand{\AojFitPMin}{0.010}  % experiments/FIGS/data/aoj_full_v1/fit_v6/results.json :: models.*.top.asymptotic_p (min) :: sha256 2416fc21f93fcfa6 :: goodness of fit of the working-point fit, pretrained models
\newcommand{\AojFitPMax}{0.974}  % experiments/FIGS/data/aoj_full_v1/fit_v6/results.json :: models.*.top.asymptotic_p (max) :: sha256 2416fc21f93fcfa6
\newcommand{\AojWeakestYield}{885}  % experiments/FIGS/data/aoj_full_v1/fit_v6/results.json :: models.l162-s2.top.signal_yield :: sha256 2416fc21f93fcfa6 :: the smallest yield of any pretrained model
\newcommand{\AojWeakestYieldErr}{211}  % experiments/FIGS/data/aoj_full_v1/fit_v6/results.json :: models.l162-s2.top.signal_yield_err :: sha256 2416fc21f93fcfa6
\newcommand{\AojWeakestFitP}{0.94}  % experiments/FIGS/data/aoj_full_v1/fit_v6/results.json :: models.l162-s2.top.asymptotic_p :: sha256 2416fc21f93fcfa6
\newcommand{\AojSidebandZMin}{\ensuremath{-}0.43}  % experiments/FIGS/data/aoj_full_v1/fit_v6/histograms.npz :: experiments/AOJ/sideband_residuals.py: z (min over the pretrained models) :: sha256 2bd1e3d5b32d4a12 :: passing jets minus the fit, summed from the fit's lower edge to the top window
\newcommand{\AojSidebandZMax}{+0.58}  % experiments/FIGS/data/aoj_full_v1/fit_v6/histograms.npz :: experiments/AOJ/sideband_residuals.py: z (max over the pretrained models) :: sha256 2bd1e3d5b32d4a12
\newcommand{\AojWeakestSidebandZ}{0.00}  % experiments/FIGS/data/aoj_full_v1/fit_v6/histograms.npz :: experiments/AOJ/sideband_residuals.py: z of l162-s2 :: sha256 2bd1e3d5b32d4a12 :: the smallest-yield model
\newcommand{\AojSidebandMaxPull}{2.7}  % experiments/FIGS/data/aoj_full_v1/fit_v6/histograms.npz :: experiments/AOJ/sideband_residuals.py: max_abs_pull (max over the pretrained models) :: sha256 2bd1e3d5b32d4a12 :: largest single-bin pull below the top window, any pretrained model
\newcommand{\AojSharedBinLo}{125}  % experiments/FIGS/data/aoj_full_v1/fit_v6/histograms.npz :: the bin holding the most models' largest pull :: sha256 2bd1e3d5b32d4a12 :: GeV
\newcommand{\AojSharedBinHi}{130}  % experiments/FIGS/data/aoj_full_v1/fit_v6/histograms.npz :: its upper edge :: sha256 2bd1e3d5b32d4a12 :: GeV
\newcommand{\AojSharedBinNModels}{18}  % experiments/FIGS/data/aoj_full_v1/fit_v6/histograms.npz :: models whose largest pull is in that bin :: sha256 2bd1e3d5b32d4a12
\newcommand{\AojSharedBinMeanPull}{+1.4}  % experiments/FIGS/data/aoj_full_v1/fit_v6/histograms.npz :: experiments/AOJ/sideband_residuals.py mean_pulls: that bin :: sha256 2bd1e3d5b32d4a12 :: mean pull over the pretrained models
\newcommand{\AojSharedBinRank}{third}  % experiments/FIGS/data/aoj_full_v1/fit_v6/histograms.npz :: experiments/AOJ/sideband_residuals.py mean_pulls: rank of that bin's |mean pull| :: sha256 2bd1e3d5b32d4a12
\newcommand{\AojFitNBins}{39}  % experiments/FIGS/data/aoj_full_v1/fit_v6/histograms.npz :: mass bins of the fit range :: sha256 2bd1e3d5b32d4a12
\newcommand{\AojPairedYieldOnesixtwoOverOneeighteight}{0.94~[0.45, 1.95]}  % experiments/FIGS/data/aoj_full_v1/analysis_v6/aoj_top.json :: per_label_set.label_sets.{162,188}.signal_yields :: sha256 5dfef902a5fe75ba :: yield ratio paired by run, geometric mean [95% Student-t interval]
\newcommand{\AojPairedYieldFourthreeOverOneeighteight}{0.97~[0.81, 1.16]}  % experiments/FIGS/data/aoj_full_v1/analysis_v6/aoj_top.json :: per_label_set.label_sets.{43,188}.signal_yields :: sha256 5dfef902a5fe75ba :: yield ratio paired by run, geometric mean [95% Student-t interval]
\newcommand{\AojPairedYieldOnesevenOverOneeighteight}{1.01~[0.75, 1.35]}  % experiments/FIGS/data/aoj_full_v1/analysis_v6/aoj_top.json :: per_label_set.label_sets.{17,188}.signal_yields :: sha256 5dfef902a5fe75ba :: yield ratio paired by run, geometric mean [95% Student-t interval]
\newcommand{\AojPairedYieldOnesevenOverOnesixtwo}{1.07~[0.51, 2.24]}  % experiments/FIGS/data/aoj_full_v1/analysis_v6/aoj_top.json :: per_label_set.label_sets.{17,162}.signal_yields :: sha256 5dfef902a5fe75ba :: yield ratio paired by run, geometric mean [95% Student-t interval]
\newcommand{\AojPairedYieldOnesixtwoMassOverOnesixtwo}{1.10~[0.57, 2.14]}  % experiments/FIGS/data/aoj_full_v1/analysis_v6/aoj_top.json :: per_label_set.label_sets.{162+mass,162}.signal_yields :: sha256 5dfef902a5fe75ba :: yield ratio paired by run, geometric mean [95% Student-t interval]
\newcommand{\AojPairedYieldOnesevenMassOverOneseven}{1.05~[0.74, 1.49]}  % experiments/FIGS/data/aoj_full_v1/analysis_v6/aoj_top.json :: per_label_set.label_sets.{17+mass,17}.signal_yields :: sha256 5dfef902a5fe75ba :: yield ratio paired by run, geometric mean [95% Student-t interval]
\newcommand{\AojLeakPairedShiftMax}{2\%}  % experiments/FIGS/data/aoj_full_v1/fit_v6/results.json :: models.*.top.leak_systematic.signal_yield against signal_yield, paired ratios :: sha256 2416fc21f93fcfa6 :: largest change of a paired yield ratio when the fail-region tops are scaled up
\newcommand{\AojRefEffPercent}{1.38\%}  % experiments/FIGS/data/aoj_full_v1/fit_v6/results.json :: reference.top.data_efficiency :: sha256 2416fc21f93fcfa6 :: fraction of all data jets passing its cut
\newcommand{\AojRefFitP}{\ensuremath{2.56\times10^{-11}}}  % experiments/FIGS/data/aoj_full_v1/fit_v6/results.json :: reference.top.asymptotic_p :: sha256 2416fc21f93fcfa6 :: goodness of fit of the reference's own top fit
\newcommand{\AojRefValidationP}{\ensuremath{p \leq 0.005} (0 of 200 toys reach the observed deviance)}  % experiments/FIGS/data/aoj_full_v1/fit_v6/results.json :: reference.top.validation.toy_p, n_toys :: sha256 2416fc21f93fcfa6 :: toy p-value of the background validation, (k + 1) / (n + 1)
\newcommand{\AojRefBandSignalZ}{2.5}  % experiments/FIGS/data/aoj_full_v1/fit_v6/results.json :: reference.top.validation.band_signal_z :: sha256 2416fc21f93fcfa6 :: significance of a peak of the score's own shape in the validation band
\newcommand{\AojModelEffMin}{1.01\%}  % experiments/FIGS/data/aoj_full_v1/fit_v6/results.json :: models.*.top.data_efficiency (min) :: sha256 2416fc21f93fcfa6 :: pretrained models
\newcommand{\AojModelEffMax}{1.05\%}  % experiments/FIGS/data/aoj_full_v1/fit_v6/results.json :: models.*.top.data_efficiency (max) :: sha256 2416fc21f93fcfa6 :: pretrained models
\newcommand{\AojModelsValidationPMin}{0.11}  % experiments/FIGS/data/aoj_full_v1/fit_v6/results.json :: models.*.top.validation.toy_p (min) :: sha256 2416fc21f93fcfa6 :: pretrained models
\newcommand{\AojRefYield}{29{,}143}  % experiments/FIGS/data/aoj_full_v1/fit_v6/results.json :: reference.top.signal_yield :: sha256 2416fc21f93fcfa6
\newcommand{\AojRefYieldErr}{865}  % experiments/FIGS/data/aoj_full_v1/fit_v6/results.json :: reference.top.signal_yield_err :: sha256 2416fc21f93fcfa6
\newcommand{\AojRefMass}{181.0}  % experiments/FIGS/data/aoj_full_v1/fit_v6/results.json :: reference.top.floated_mean :: sha256 2416fc21f93fcfa6 :: GeV
\newcommand{\AojRefWidth}{17.8}  % experiments/FIGS/data/aoj_full_v1/fit_v6/results.json :: reference.top.floated_width :: sha256 2416fc21f93fcfa6 :: GeV
\newcommand{\AojPublicYield}{1{,}813}  % experiments/FIGS/data/aoj_full_v1/fit_v6/results.json :: models.sophon-public.top.signal_yield :: sha256 2416fc21f93fcfa6
\newcommand{\AojPublicYieldErr}{193}  % experiments/FIGS/data/aoj_full_v1/fit_v6/results.json :: models.sophon-public.top.signal_yield_err :: sha256 2416fc21f93fcfa6
\newcommand{\AojPublicMass}{182.8}  % experiments/FIGS/data/aoj_full_v1/fit_v6/results.json :: models.sophon-public.top.floated_mean :: sha256 2416fc21f93fcfa6 :: GeV
\newcommand{\AojPublicWidth}{11.7}  % experiments/FIGS/data/aoj_full_v1/fit_v6/results.json :: models.sophon-public.top.floated_width :: sha256 2416fc21f93fcfa6 :: GeV
\newcommand{\AojInjSizeMin}{1{,}000}  % experiments/FIGS/data/aoj_injection_v1/summary.json :: groups[region=top,mode=ensemble].size (min) :: sha256 d9b66466ae24b584 :: injected tops
\newcommand{\AojInjSizeMax}{4{,}000}  % experiments/FIGS/data/aoj_injection_v1/summary.json :: groups[region=top,mode=ensemble].size (max) :: sha256 d9b66466ae24b584
\newcommand{\AojInjFloatPullSdMax}{2.0}  % experiments/FIGS/data/aoj_injection_v1/summary.json :: groups[region=top,mode=bootstrap,variant=float].pull.sd (max) :: sha256 d9b66466ae24b584 :: pull width with each model's own floated shape
\newcommand{\AojInjEnsembleRecoveryMin}{0.983}  % experiments/FIGS/data/aoj_injection_v1/summary.json :: groups[region=top,mode=ensemble].ratio (min) :: sha256 d9b66466ae24b584 :: the whole procedure, shape re-derived per toy set
\newcommand{\AojInjEnsembleRecoveryMax}{0.998}  % experiments/FIGS/data/aoj_injection_v1/summary.json :: groups[region=top,mode=ensemble].ratio (max) :: sha256 d9b66466ae24b584
\newcommand{\AojInjEnsemblePullSdMin}{1.02}  % experiments/FIGS/data/aoj_injection_v1/summary.json :: groups[region=top,mode=ensemble].pull (min) :: sha256 d9b66466ae24b584 :: the whole procedure, shape re-derived per toy set
\newcommand{\AojInjEnsemblePullSdMax}{1.09}  % experiments/FIGS/data/aoj_injection_v1/summary.json :: groups[region=top,mode=ensemble].pull (max) :: sha256 d9b66466ae24b584
\newcommand{\AojClosureNShards}{10}  % experiments/FIGS/data/aoj_checks_v2/closure_all_shards.json :: n_shards :: sha256 6a6018351d7048d2 :: staged shards of the open data
\newcommand{\AojClosurePtLo}{500}  % experiments/FIGS/data/aoj_checks_v2/closure_all_shards.json :: pt_window[0] :: sha256 6a6018351d7048d2 :: GeV
\newcommand{\AojClosurePtHi}{700}  % experiments/FIGS/data/aoj_checks_v2/closure_all_shards.json :: pt_window[1] :: sha256 6a6018351d7048d2 :: GeV
\newcommand{\AojClosureDzeroIqr}{8.7}  % experiments/FIGS/data/aoj_checks_v2/closure_all_shards.json :: domain_shift.part_d0.iqr_ratio :: sha256 6a6018351d7048d2 :: open data over JetClass-II QCD, pooled over the shards
\newcommand{\AojClosureDzIqr}{10.6}  % experiments/FIGS/data/aoj_checks_v2/closure_all_shards.json :: domain_shift.part_dz.iqr_ratio :: sha256 6a6018351d7048d2 :: open data over JetClass-II QCD, pooled over the shards
\newcommand{\AojClosureNeutral}{0.65}  % experiments/FIGS/data/aoj_checks_v2/closure_all_shards.json :: features.part_isNeutralHadron.ratio.per_shard (mean) :: sha256 6a6018351d7048d2 :: open data over JetClass-II QCD
\newcommand{\AojClosureCharged}{1.22}  % experiments/FIGS/data/aoj_checks_v2/closure_all_shards.json :: features.part_isChargedHadron.ratio.per_shard (mean) :: sha256 6a6018351d7048d2 :: open data over JetClass-II QCD
\newcommand{\AojClosureNPartAoj}{61}  % experiments/FIGS/data/aoj_checks_v2/closure_all_shards.json :: features.n_particles.aoj.per_shard (mean of the shard medians) :: sha256 6a6018351d7048d2 :: median
\newcommand{\AojClosureNPartRef}{48}  % experiments/FIGS/data/aoj_checks_v2/closure_all_shards.json :: features.n_particles.reference.per_shard (mean) :: sha256 6a6018351d7048d2 :: median
\newcommand{\RandMergingDrawsBvcFourprong}{three}  % configs/labelmaps/rand_label_map.v1.csv :: RAND_d* of label_X_YY_bbqq and label_X_YY_ccqq :: sha256 f395a25f52e3cfad :: first-grid partitions that merge the task's pair
\newcommand{\RandSplittingDrawsBvcFourprong}{one and two}  % configs/labelmaps/rand_label_map.v1.csv :: RAND_d* of label_X_YY_bbqq and label_X_YY_ccqq :: sha256 f395a25f52e3cfad :: first-grid partitions that keep the pair apart
\newcommand{\RandMergingDrawsVisibleContent}{one and three}  % configs/labelmaps/rand_label_map.v1.csv :: RAND_d* of label_X_YY_bbqq and label_X_YY_cqtauhv :: sha256 f395a25f52e3cfad :: first-grid partitions that merge the task's pair
\newcommand{\RandSplittingDrawsVisibleContent}{two}  % configs/labelmaps/rand_label_map.v1.csv :: RAND_d* of label_X_YY_bbqq and label_X_YY_cqtauhv :: sha256 f395a25f52e3cfad :: first-grid partitions that keep the pair apart
\newcommand{\RandVtwoNPartitions}{five}  % configs/arms/v2_grid.json :: arms[RAND2_*] (count) :: sha256 d8a31b1486992302 :: random partitions in the rerun
\newcommand{\RandVtwoRuns}{one}  % configs/arms/v2_paper_runs.txt :: runs of RAND2_* :: sha256 b04562b14ea8cbc7 :: runs per partition
\newcommand{\RandFlavRuns}{two}  % configs/arms/v2_paper_runs.txt :: runs of FLAV_* :: sha256 b04562b14ea8cbc7 :: runs per flavour-pair model
\newcommand{\RandVtwoNPairs}{seven}  % configs/labelmaps/rand_v2_selection.json :: pairs (count) :: sha256 98cf52d49c7d2f02 :: probe pairs the balance rule covers
\newcommand{\RandVtwoMergedMin}{two}  % configs/labelmaps/rand_v2_selection.json :: merged_in :: sha256 98cf52d49c7d2f02 :: partitions merging each probe pair, checked on the map
\newcommand{\RandVtwoMergedMax}{three}  % configs/labelmaps/rand_v2_selection.json :: merged_in :: sha256 98cf52d49c7d2f02 :: partitions merging each probe pair, checked on the map
\newcommand{\RandVtwoShareTol}{5\%}  % configs/labelmaps/rand_v2_selection.json :: realised_tolerance :: sha256 98cf52d49c7d2f02 :: largest relative deviation of a realised group share from its target (rule 3)
\newcommand{\RandVtwoTrials}{32}  % scripts/build_rand_control.py :: TRIALS :: sha256 5bc91d77a3f2e81f :: share-matched assignments drawn per random partition; the one sharing the fewest merged pairs with the 17-class label set is kept
\newcommand{\RandVtwoPoolDraws}{6{,}600}  % configs/labelmaps/rand_v2_selection.json :: pool.draws :: sha256 98cf52d49c7d2f02 :: candidate partitions drawn (seeds) before the realised-share rule
\newcommand{\RandVtwoPoolPass}{27}  % configs/labelmaps/rand_v2_selection.json :: rule3_seeds (count) :: sha256 98cf52d49c7d2f02 :: candidate partitions meeting the realised-share rule (rule 3)
\newcommand{\RandFlavSplitOrbit}{\ensuremath{X\to YY\to QQQQ}}  % configs/labelmaps/flavour_pair.v2.json :: split_orbit :: sha256 3b79ce5e44c1e2f4 :: the decays F1 cuts, Q any quark
\newcommand{\RandFlavSplitOrbitClasses}{22}  % configs/labelmaps/flavour_pair_map.v2.csv :: classes whose orbit is split_orbit (count) :: sha256 9aab235099ceac43 :: classes of the split orbit
\newcommand{\RandFlavCutClasses}{11}  % configs/labelmaps/flavour_pair.v2.json :: split_classes_moved (count) :: sha256 3b79ce5e44c1e2f4 :: classes of the split orbit that F1 moves: those with a b quark
\newcommand{\RandFlavOrbitsMoved}{three}  % configs/labelmaps/flavour_pair.v2.json :: orbits_moved_B_to_A (count) :: sha256 3b79ce5e44c1e2f4 :: whole orbits F1 moves the other way
\newcommand{\RandFlavOrbitsMovedList}{\ensuremath{X\to YY\to QQ\mu\mu}, \ensuremath{X\to YY\to QQ\tau_h\tau_h} and \ensuremath{X\to \mu\mu}}  % configs/labelmaps/flavour_pair.v2.json :: orbits_moved_B_to_A :: sha256 3b79ce5e44c1e2f4 :: the orbits F1 moves the other way, Q any quark
\newcommand{\RandFlavOrbitClassesMoved}{9}  % configs/labelmaps/flavour_pair_map.v2.csv :: classes whose orbit is in orbits_moved_B_to_A :: sha256 9aab235099ceac43 :: classes in those orbits
\newcommand{\RandFlavFonerCutWithB}{6}  % configs/labelmaps/flavour_pair_map.v2.csv :: FLAV_F1R != FLAV_F0 within split_orbit, with a b quark (count) :: sha256 9aab235099ceac43 :: classes with a b quark among those the F1r cut moves
\newcommand{\RandFlavBcPairs}{11}  % configs/labelmaps/flavour_pair.v2.json :: b_c_pairs_in_split_orbit (count) :: sha256 3b79ce5e44c1e2f4 :: pairs of the split orbit's classes that differ only by exchanging b and c
\newcommand{\RandFlavFoneBcSplit}{8}  % configs/labelmaps/flavour_pair_map.v2.csv :: b_c_pairs_in_split_orbit split by FLAV_F1 (count) :: sha256 9aab235099ceac43
\newcommand{\RandFlavFonerBcSplit}{none}  % configs/labelmaps/flavour_pair_map.v2.csv :: b_c_pairs_in_split_orbit split by FLAV_F1R (count) :: sha256 9aab235099ceac43 :: PRESPEC A14: none
\newcommand{\RandFlavFonerShareTol}{1\%}  % configs/labelmaps/flavour_pair.v2.json :: F1R.realised_tolerance_vs_F1 :: sha256 3b79ce5e44c1e2f4 :: largest allowed relative difference of a group's realised share from F1's
\newcommand{\RandFlavFonerShareDev}{0.77\%}  % configs/labelmaps/realised_native_shares.v2.json :: share summed per FLAV_F1R and FLAV_F1 group of flavour_pair_map.v2.csv (largest relative difference) :: sha256 5728c12e88fa3683 :: realised shares from the final loader's dry run
\newcommand{\MlstNRunsWord}{three}  % configs/arms/v2_paper_runs.txt :: runs of L188/L162/R42_Q1/R16_Q1 :: sha256 b04562b14ea8cbc7 :: pretraining runs per configuration
\newcommand{\MlstMassNRunsWord}{two}  % configs/arms/v2_paper_runs.txt :: runs of L162_MASS/R16_Q1_MASS :: sha256 b04562b14ea8cbc7 :: pretraining runs per configuration
\newcommand{\LpCMin}{\ensuremath{1\times10^{-3}}}  % experiments/EVAL/linear_probe_v2.py :: C_GRID :: sha256 4c576b44812c18a2 :: smallest inverse regularisation
\newcommand{\LpCMax}{100}  % experiments/EVAL/linear_probe_v2.py :: C_GRID :: sha256 4c576b44812c18a2 :: largest inverse regularisation
\newcommand{\LpNTestJcii}{2{,}000{,}000}  % experiments/EVAL/linear_probe_v2.py :: TEST_FIRST.jc2 :: sha256 4c576b44812c18a2 :: JetClass-II test jets
\newcommand{\LpNTestAucJcii}{500{,}000}  % experiments/EVAL/linear_probe_v2.py :: TEST_FIRST.jc2 / AUC_STRIDE :: sha256 4c576b44812c18a2 :: JetClass-II test jets of the macro AUC
\newcommand{\LpConvergedGrad}{\ensuremath{1\times10^{-5}}}  % experiments/EVAL/linear_probe_v2.py :: CONVERGED_GRAD :: sha256 4c576b44812c18a2 :: largest gradient component a converged probe ends below
\newcommand{\FillNFeatures}{17}  % configs/finetune/TopReference.yaml :: inputs.pf_features.vars (count) :: sha256 501a6ffe5d0a8051 :: per-particle input features of the network
\newcommand{\FillTopNZero}{10}  % configs/finetune/TopReference.yaml :: inputs.pf_features.vars, part_zero_* (count) :: sha256 501a6ffe5d0a8051 :: features the data set lacks, set to zero
\newcommand{\FillTopNPresent}{7}  % configs/finetune/TopReference.yaml :: inputs.pf_features.vars, not part_zero_* (count) :: sha256 501a6ffe5d0a8051 :: features the data set provides
\newcommand{\FillQgNZero}{4}  % configs/finetune/EnergyFlowQG.yaml :: inputs.pf_features.vars, part_zero_* (count) :: sha256 db42d6471e8b865c :: features the data set lacks, set to zero
\newcommand{\FillQgNPresent}{13}  % configs/finetune/EnergyFlowQG.yaml :: inputs.pf_features.vars, not part_zero_* (count) :: sha256 db42d6471e8b865c :: features the data set provides
\newcommand{\DesignVtwoDataFraction}{20\%}  % experiments/MTX/k8s/v2/grid/job-mtx2-l188-s1-raunav.yaml :: --data-fraction :: sha256 d28ce02267a6b87a :: share of every file read per epoch
\newcommand{\DesignVtwoNValFiles}{25}  % experiments/MTX/k8s/v2/grid/job-mtx2-l188-s1-raunav.yaml :: --data-val (files, brace ranges expanded) :: sha256 d28ce02267a6b87a :: validation files of the fixed sample
\newcommand{\DesignWavgFirstEpoch}{70}  % experiments/MTX/pretrain_v2.py :: --num-epochs - WAVG_EPOCHS :: sha256 95da5a6e029855ee :: first epoch of the weight average (epochs count from 0)
\newcommand{\DesignWavgLastEpoch}{79}  % experiments/MTX/k8s/v2/grid/job-mtx2-l188-s1-raunav.yaml :: --num-epochs - 1 :: sha256 d28ce02267a6b87a :: last epoch
\newcommand{\DesignBnJets}{200{,}000}  % experiments/MTX/pretrain_v2.py :: BN_JETS :: sha256 95da5a6e029855ee :: training jets the BatchNorm statistics are recomputed on
\newcommand{\DesignVtwoNWindows}{five}  % experiments/MTX/k8s/v2/grid/job-mtx2-l188-s1-raunav.yaml :: 1 / --data-fraction :: sha256 d28ce02267a6b87a :: disjoint windows of each file's rows, one read per epoch
\newcommand{\DesignLofoNWindows}{three}  % experiments/MTX/k8s/v2/cut/grid/job-mtx2-l188lofo4p-s1-raunav.yaml :: --data-windows :: sha256 44a9573929efb7ac :: windows of the models trained without one family
\newcommand{\DesignVtwoPlanUsed}{97\%}  % experiments/FIGS/data/v2_loader/loader_dryrun/dryrun_s176_seed1.json :: epochs[*].(max_fetch_id + 1) / args.data_split_num, mean :: sha256 03b6586b26ad7a2f :: share of an epoch's planned reading used before its jets are reached, 20 epochs
\newcommand{\DesignLofoPlanUsed}{76\%}  % experiments/FIGS/data/v2_loader/loader_dryrun/dryrun_lofo4p_seed1.json :: epochs[*].(max_fetch_id + 1) / fetches_per_pass, mean :: sha256 2a1315e611d13ff5 :: the same, leaving the family out
\newcommand{\DesignLofoRowFactor}{1.29}  % experiments/FIGS/data/v2_loader/loader_dryrun/dryrun_lofo4p_seed1.json :: (mean plan used / --data-windows) over (mean plan used / k) of the full models' dry run :: sha256 2a1315e611d13ff5 :: how much more often a row of a file is read per epoch when the family is left out
\newcommand{\DesignDistinctJetsCycle}{\ensuremath{3.3\times10^{7}}}  % experiments/FIGS/data/v2_loader/loader_dryrun/dryrun_s176_seed1.json :: sum of epochs[c*k:(c+1)*k].distinct_jets, mean over 4 cycles :: sha256 03b6586b26ad7a2f :: distinct jets the full models train on in one cycle of windows
\newcommand{\DesignLofoFamilyShare}{22.4\%}  % experiments/FIGS/data/v2_loader/loader_dryrun/dryrun_s176_seed1.json :: epochs[*].native_counts over summary.labels_absent_every_epoch of the family-out run :: sha256 03b6586b26ad7a2f :: the left-out family's share of the full models' stream
\newcommand{\DesignLofoNClasses}{27}  % experiments/FIGS/data/v2_loader/loader_dryrun/dryrun_lofo4p_seed1.json :: summary.labels_absent_every_epoch (count) :: sha256 2a1315e611d13ff5 :: native classes of the left-out family
\newcommand{\DesignLofoExposureFactor}{1.29}  % experiments/FIGS/data/v2_loader/loader_dryrun/dryrun_s176_seed1.json :: 1 / (1 - family share) :: sha256 03b6586b26ad7a2f :: how much more often every other class is seen
\newcommand{\DesignLofoQcdShareFull}{14.4\%}  % experiments/FIGS/data/v2_loader/loader_dryrun/dryrun_s176_seed1.json :: summary.mean_qcd_share :: sha256 03b6586b26ad7a2f :: QCD share of the realised stream
\newcommand{\DesignLofoQcdShareLofo}{18.6\%}  % experiments/FIGS/data/v2_loader/loader_dryrun/dryrun_lofo4p_seed1.json :: summary.mean_qcd_share :: sha256 2a1315e611d13ff5 :: QCD share of the realised stream
\newcommand{\DesignLrFlatPercent}{70\%}  % experiments/MTX/pretrain_v2.py :: make_optimizer: 1 - the decay share :: sha256 95da5a6e029855ee :: share of the epochs at the initial rate
\newcommand{\DesignLrDecayFirstEpoch}{56}  % experiments/MTX/pretrain_v2.py :: num_epochs - int(num_epochs * decay share) :: sha256 95da5a6e029855ee :: first epoch at a reduced rate (epochs from 0)
\newcommand{\DesignLrEndPercent}{1\%}  % experiments/MTX/pretrain_v2.py :: make_optimizer: gamma ** n_decay :: sha256 95da5a6e029855ee :: the last epoch's rate over the initial rate
\newcommand{\DesignEndLr}{\ensuremath{5\times10^{-6}}}  % experiments/MTX/k8s/v2/grid/job-mtx2-l188-s1-raunav.yaml :: --start-lr x the final factor :: sha256 d28ce02267a6b87a :: the rate of the last epoch
\newcommand{\DesignWeaverVersion}{0.4.17}  % experiments/MTX/pretrain_v2.py :: the weaver release whose loop it reproduces :: sha256 95da5a6e029855ee
\newcommand{\DesignMeiFactor}{1.1}  % src/stats/mde.py :: MEI_LOG = ln(factor) :: sha256 626468dcc37940bb :: the smallest effect of interest, as a factor
\newcommand{\DesignVtwoGpu}{RTX 3090}  % experiments/MTX/k8s/v2/grid/job-mtx2-flavf0-s1-raunav.yaml :: nodeAffinity nvidia.com/gpu.product :: sha256 960cad1691314980 :: the product of run index 1, every configuration
\newcommand{\DesignVtwoSecondGpu}{\pending{pending: the GPU type of each run index}}  % experiments/FIGS/make_tables.py :: emit_v2_recipe[DesignVtwoSecondGpu] :: sha256 10e6ec08356a4a2b :: a marker, not a number: the grid's job specs put every run index on one product
\newcommand{\DesignVtwoSecondGpuRuns}{\pending{pending: the GPU type of each run index}}  % experiments/FIGS/make_tables.py :: emit_v2_recipe[DesignVtwoSecondGpuRuns] :: sha256 10e6ec08356a4a2b :: a marker, not a number: the grid's job specs put every run index on one product
\newcommand{\MassNRunsWord}{five}  % experiments/FIGS/data/probe_ladder_mass2x2_mlp2/s1.json :: per-seed files (count) :: sha256 906e778d7043fe67 :: pretraining runs of each mass-output configuration
\newcommand{\MassOmaLinearOnesixtwo}{\ensuremath{(2.36 \pm 0.06)\times10^{-3}}}  % experiments/FIGS/data/probe_ladder_v2_mlp2/analysis/seed_level_results.json :: table[bvc_resonant,linear,162].auc :: sha256 53e1ec81343ac2c1 :: 1-AUC, mean +- SD over 5 seeds
\newcommand{\MassOmaLinearOnesixtwoMass}{\ensuremath{(2.62 \pm 0.08)\times10^{-3}}}  % experiments/FIGS/data/probe_ladder_mass2x2_mlp2/s1.json :: tasks.bvc_resonant.arms.l162mass-s*.linear.auc :: sha256 906e778d7043fe67 :: 1-AUC, mean +- SD over the 5 per-seed files
\newcommand{\MassOmaRatioLinearOnesixtwo}{1.11}  % experiments/FIGS/data/probe_ladder_mass2x2_mlp2/s1.json :: tasks.bvc_resonant.arms.l162mass-s*.linear.auc over table[bvc_resonant,linear,162].auc :: sha256 906e778d7043fe67 :: seed mean of 1-AUC with the mass output over without
\newcommand{\MassOmaLinearOneseven}{\ensuremath{(9.3 \pm 1.5)\times10^{-3}}}  % experiments/FIGS/data/probe_ladder_v2_mlp2/analysis/seed_level_results.json :: table[bvc_resonant,linear,17].auc :: sha256 53e1ec81343ac2c1 :: 1-AUC, mean +- SD over 5 seeds
\newcommand{\MassOmaLinearOnesevenMass}{\ensuremath{(3.0 \pm 1.2)\times10^{-2}}}  % experiments/FIGS/data/probe_ladder_mass2x2_mlp2/s1.json :: tasks.bvc_resonant.arms.r16q1mass-s*.linear.auc :: sha256 906e778d7043fe67 :: 1-AUC, mean +- SD over the 5 per-seed files
\newcommand{\MassOmaRatioLinearOneseven}{3.2}  % experiments/FIGS/data/probe_ladder_mass2x2_mlp2/s1.json :: tasks.bvc_resonant.arms.r16q1mass-s*.linear.auc over table[bvc_resonant,linear,17].auc :: sha256 906e778d7043fe67 :: seed mean of 1-AUC with the mass output over without
\newcommand{\MassOmaMlpOnesixtwo}{\ensuremath{(2.32 \pm 0.05)\times10^{-3}}}  % experiments/FIGS/data/probe_ladder_v2_mlp2/analysis/seed_level_results.json :: table[bvc_resonant,mlp,162].auc :: sha256 53e1ec81343ac2c1 :: 1-AUC, mean +- SD over 5 seeds
\newcommand{\MassOmaMlpOnesixtwoMass}{\ensuremath{(2.52 \pm 0.06)\times10^{-3}}}  % experiments/FIGS/data/probe_ladder_mass2x2_mlp2/s1.json :: tasks.bvc_resonant.arms.l162mass-s*.mlp.auc :: sha256 906e778d7043fe67 :: 1-AUC, mean +- SD over the 5 per-seed files
\newcommand{\MassOmaRatioMlpOnesixtwo}{1.09}  % experiments/FIGS/data/probe_ladder_mass2x2_mlp2/s1.json :: tasks.bvc_resonant.arms.l162mass-s*.mlp.auc over table[bvc_resonant,mlp,162].auc :: sha256 906e778d7043fe67 :: seed mean of 1-AUC with the mass output over without
\newcommand{\MassOmaMlpOneseven}{\ensuremath{(7.3 \pm 0.9)\times10^{-3}}}  % experiments/FIGS/data/probe_ladder_v2_mlp2/analysis/seed_level_results.json :: table[bvc_resonant,mlp,17].auc :: sha256 53e1ec81343ac2c1 :: 1-AUC, mean +- SD over 5 seeds
\newcommand{\MassOmaMlpOnesevenMass}{\ensuremath{(1.5 \pm 0.4)\times10^{-2}}}  % experiments/FIGS/data/probe_ladder_mass2x2_mlp2/s1.json :: tasks.bvc_resonant.arms.r16q1mass-s*.mlp.auc :: sha256 906e778d7043fe67 :: 1-AUC, mean +- SD over the 5 per-seed files
\newcommand{\MassOmaRatioMlpOneseven}{2.0}  % experiments/FIGS/data/probe_ladder_mass2x2_mlp2/s1.json :: tasks.bvc_resonant.arms.r16q1mass-s*.mlp.auc over table[bvc_resonant,mlp,17].auc :: sha256 906e778d7043fe67 :: seed mean of 1-AUC with the mass output over without
\newcommand{\FtAucJciiEThreeOneeighteight}{0.902\,\ensuremath{\pm}\,0.004}  % experiments/FIGS/data/w2b_leg1_metrics_v2.json :: cells.{l188-s1,l188-s2,l188-s3,l188-s4,l188-s5}.N1000.s1.macro_auc_ovr :: sha256 b56bf70be9982044 :: mean +- SD over the row's 5 pretrained models
\newcommand{\FtAucJciiEFourOneeighteight}{0.9723\,\ensuremath{\pm}\,0.0004}  % experiments/FIGS/data/w2b_leg1_metrics_v2.json :: cells.{l188-s1,l188-s2,l188-s3,l188-s4,l188-s5}.N10000.s1.macro_auc_ovr :: sha256 b56bf70be9982044 :: mean +- SD over the row's 5 pretrained models
\newcommand{\FtAucJciiEFiveOneeighteight}{0.9793\,\ensuremath{\pm}\,0.0002}  % experiments/FIGS/data/w2b_leg1_metrics_v2.json :: cells.{l188-s1,l188-s2,l188-s3,l188-s4,l188-s5}.N100000.s1.macro_auc_ovr :: sha256 b56bf70be9982044 :: mean +- SD over the row's 5 pretrained models
\newcommand{\FtAucJciiESixOneeighteight}{0.9809\,\ensuremath{\pm}\,\ensuremath{<}0.0001}  % experiments/FIGS/data/w2b_leg1_metrics_v2.json :: cells.{l188-s1,l188-s2,l188-s3,l188-s4,l188-s5}.N1000000.s1.macro_auc_ovr :: sha256 b56bf70be9982044 :: mean +- SD over the row's 5 pretrained models
\newcommand{\FtAucJciiEThreeOnesixtwo}{0.906\,\ensuremath{\pm}\,0.007}  % experiments/FIGS/data/w2b_leg1_metrics_v2.json :: cells.{l162-s1b,l162-s2,l162-s3,l162-s4,l162-s5}.N1000.s1.macro_auc_ovr :: sha256 b56bf70be9982044 :: mean +- SD over the row's 5 pretrained models
\newcommand{\FtAucJciiEFourOnesixtwo}{0.9724\,\ensuremath{\pm}\,0.0005}  % experiments/FIGS/data/w2b_leg1_metrics_v2.json :: cells.{l162-s1b,l162-s2,l162-s3,l162-s4,l162-s5}.N10000.s1.macro_auc_ovr :: sha256 b56bf70be9982044 :: mean +- SD over the row's 5 pretrained models
\newcommand{\FtAucJciiEFiveOnesixtwo}{0.9794\,\ensuremath{\pm}\,0.0002}  % experiments/FIGS/data/w2b_leg1_metrics_v2.json :: cells.{l162-s1b,l162-s2,l162-s3,l162-s4,l162-s5}.N100000.s1.macro_auc_ovr :: sha256 b56bf70be9982044 :: mean +- SD over the row's 5 pretrained models
\newcommand{\FtAucJciiESixOnesixtwo}{0.9808\,\ensuremath{\pm}\,0.0001}  % experiments/FIGS/data/w2b_leg1_metrics_v2.json :: cells.{l162-s1b,l162-s2,l162-s3,l162-s4,l162-s5}.N1000000.s1.macro_auc_ovr :: sha256 b56bf70be9982044 :: mean +- SD over the row's 5 pretrained models
\newcommand{\FtAucJciiEThreeFourthree}{0.918\,\ensuremath{\pm}\,0.005}  % experiments/FIGS/data/w2b_leg1_metrics_v2.json :: cells.{r42q1-s1,r42q1-s2,r42q1-s3,r42q1-s4,r42q1-s5}.N1000.s1.macro_auc_ovr :: sha256 b56bf70be9982044 :: mean +- SD over the row's 5 pretrained models
\newcommand{\FtAucJciiEFourFourthree}{0.9667\,\ensuremath{\pm}\,0.0005}  % experiments/FIGS/data/w2b_leg1_metrics_v2.json :: cells.{r42q1-s1,r42q1-s2,r42q1-s3,r42q1-s4,r42q1-s5}.N10000.s1.macro_auc_ovr :: sha256 b56bf70be9982044 :: mean +- SD over the row's 5 pretrained models
\newcommand{\FtAucJciiEFiveFourthree}{0.9775\,\ensuremath{\pm}\,0.0003}  % experiments/FIGS/data/w2b_leg1_metrics_v2.json :: cells.{r42q1-s1,r42q1-s2,r42q1-s3,r42q1-s4,r42q1-s5}.N100000.s1.macro_auc_ovr :: sha256 b56bf70be9982044 :: mean +- SD over the row's 5 pretrained models
\newcommand{\FtAucJciiESixFourthree}{0.9800\,\ensuremath{\pm}\,0.0001}  % experiments/FIGS/data/w2b_leg1_metrics_v2.json :: cells.{r42q1-s1,r42q1-s2,r42q1-s3,r42q1-s4,r42q1-s5}.N1000000.s1.macro_auc_ovr :: sha256 b56bf70be9982044 :: mean +- SD over the row's 5 pretrained models
\newcommand{\FtAucJciiEThreeOneseven}{0.905\,\ensuremath{\pm}\,0.008}  % experiments/FIGS/data/w2b_leg1_metrics_v2.json :: cells.{r16q1-s1,r16q1-s2,r16q1-s3,r16q1-s4,r16q1-s5}.N1000.s1.macro_auc_ovr :: sha256 b56bf70be9982044 :: mean +- SD over the row's 5 pretrained models
\newcommand{\FtAucJciiEFourOneseven}{0.9635\,\ensuremath{\pm}\,0.0005}  % experiments/FIGS/data/w2b_leg1_metrics_v2.json :: cells.{r16q1-s1,r16q1-s2,r16q1-s3,r16q1-s4,r16q1-s5}.N10000.s1.macro_auc_ovr :: sha256 b56bf70be9982044 :: mean +- SD over the row's 5 pretrained models
\newcommand{\FtAucJciiEFiveOneseven}{0.9748\,\ensuremath{\pm}\,0.0003}  % experiments/FIGS/data/w2b_leg1_metrics_v2.json :: cells.{r16q1-s1,r16q1-s2,r16q1-s3,r16q1-s4,r16q1-s5}.N100000.s1.macro_auc_ovr :: sha256 b56bf70be9982044 :: mean +- SD over the row's 5 pretrained models
\newcommand{\FtAucJciiESixOneseven}{0.9781\,\ensuremath{\pm}\,0.0001}  % experiments/FIGS/data/w2b_leg1_metrics_v2.json :: cells.{r16q1-s1,r16q1-s2,r16q1-s3,r16q1-s4,r16q1-s5}.N1000000.s1.macro_auc_ovr :: sha256 b56bf70be9982044 :: mean +- SD over the row's 5 pretrained models
\newcommand{\FtAucJciiEThreeOnesixtwoMass}{0.921\,\ensuremath{\pm}\,0.005}  % experiments/FIGS/data/w2b_leg1_metrics_v2.json :: cells.{l162mass-s1,l162mass-s2,l162mass-s3,l162mass-s4,l162mass-s5}.N1000.s1.macro_auc_ovr :: sha256 b56bf70be9982044 :: mean +- SD over the row's 5 pretrained models
\newcommand{\FtAucJciiEFourOnesixtwoMass}{0.9724\,\ensuremath{\pm}\,0.0013}  % experiments/FIGS/data/w2b_leg1_metrics_v2.json :: cells.{l162mass-s1,l162mass-s2,l162mass-s3,l162mass-s4,l162mass-s5}.N10000.s1.macro_auc_ovr :: sha256 b56bf70be9982044 :: mean +- SD over the row's 5 pretrained models
\newcommand{\FtAucJciiEFiveOnesixtwoMass}{0.9790\,\ensuremath{\pm}\,0.0001}  % experiments/FIGS/data/w2b_leg1_metrics_v2.json :: cells.{l162mass-s1,l162mass-s2,l162mass-s3,l162mass-s4,l162mass-s5}.N100000.s1.macro_auc_ovr :: sha256 b56bf70be9982044 :: mean +- SD over the row's 5 pretrained models
\newcommand{\FtAucJciiESixOnesixtwoMass}{0.9805\,\ensuremath{\pm}\,0.0001}  % experiments/FIGS/data/w2b_leg1_metrics_v2.json :: cells.{l162mass-s1,l162mass-s2,l162mass-s3,l162mass-s4,l162mass-s5}.N1000000.s1.macro_auc_ovr :: sha256 b56bf70be9982044 :: mean +- SD over the row's 5 pretrained models
\newcommand{\FtAucJciiEThreeOnesevenMass}{0.904\,\ensuremath{\pm}\,0.007}  % experiments/FIGS/data/w2b_leg1_metrics_v2.json :: cells.{r16q1mass-s1,r16q1mass-s2,r16q1mass-s3,r16q1mass-s4,r16q1mass-s5}.N1000.s1.macro_auc_ovr :: sha256 b56bf70be9982044 :: mean +- SD over the row's 5 pretrained models
\newcommand{\FtAucJciiEFourOnesevenMass}{0.9639\,\ensuremath{\pm}\,0.0007}  % experiments/FIGS/data/w2b_leg1_metrics_v2.json :: cells.{r16q1mass-s1,r16q1mass-s2,r16q1mass-s3,r16q1mass-s4,r16q1mass-s5}.N10000.s1.macro_auc_ovr :: sha256 b56bf70be9982044 :: mean +- SD over the row's 5 pretrained models
\newcommand{\FtAucJciiEFiveOnesevenMass}{0.9740\,\ensuremath{\pm}\,0.0004}  % experiments/FIGS/data/w2b_leg1_metrics_v2.json :: cells.{r16q1mass-s1,r16q1mass-s2,r16q1mass-s3,r16q1mass-s4,r16q1mass-s5}.N100000.s1.macro_auc_ovr :: sha256 b56bf70be9982044 :: mean +- SD over the row's 5 pretrained models
\newcommand{\FtAucJciiESixOnesevenMass}{0.9774\,\ensuremath{\pm}\,0.0001}  % experiments/FIGS/data/w2b_leg1_metrics_v2.json :: cells.{r16q1mass-s1,r16q1mass-s2,r16q1mass-s3,r16q1mass-s4,r16q1mass-s5}.N1000000.s1.macro_auc_ovr :: sha256 b56bf70be9982044 :: mean +- SD over the row's 5 pretrained models
\newcommand{\FtAucJciiEThreeRandDrawOne}{0.921}  % experiments/FIGS/data/w2b_leg1_metrics_v2.json :: cells.{rand-d1-s1b}.N1000.s1.macro_auc_ovr :: sha256 b56bf70be9982044 :: one run
\newcommand{\FtAucJciiEFourRandDrawOne}{0.9675}  % experiments/FIGS/data/w2b_leg1_metrics_v2.json :: cells.{rand-d1-s1b}.N10000.s1.macro_auc_ovr :: sha256 b56bf70be9982044 :: one run
\newcommand{\FtAucJciiEFiveRandDrawOne}{0.97617}  % experiments/FIGS/data/w2b_leg1_metrics_v2.json :: cells.{rand-d1-s1b}.N100000.s1.macro_auc_ovr :: sha256 b56bf70be9982044 :: one run
\newcommand{\FtAucJciiESixRandDrawOne}{0.97941}  % experiments/FIGS/data/w2b_leg1_metrics_v2.json :: cells.{rand-d1-s1b}.N1000000.s1.macro_auc_ovr :: sha256 b56bf70be9982044 :: one run
\newcommand{\FtAucJciiEThreeRandDrawTwo}{0.914}  % experiments/FIGS/data/w2b_leg1_metrics_v2.json :: cells.{rand-d2-s2}.N1000.s1.macro_auc_ovr :: sha256 b56bf70be9982044 :: one run
\newcommand{\FtAucJciiEFourRandDrawTwo}{0.9675}  % experiments/FIGS/data/w2b_leg1_metrics_v2.json :: cells.{rand-d2-s2}.N10000.s1.macro_auc_ovr :: sha256 b56bf70be9982044 :: one run
\newcommand{\FtAucJciiEFiveRandDrawTwo}{0.97668}  % experiments/FIGS/data/w2b_leg1_metrics_v2.json :: cells.{rand-d2-s2}.N100000.s1.macro_auc_ovr :: sha256 b56bf70be9982044 :: one run
\newcommand{\FtAucJciiESixRandDrawTwo}{0.97943}  % experiments/FIGS/data/w2b_leg1_metrics_v2.json :: cells.{rand-d2-s2}.N1000000.s1.macro_auc_ovr :: sha256 b56bf70be9982044 :: one run
\newcommand{\FtAucJciiEThreeRandDrawThree}{0.918}  % experiments/FIGS/data/w2b_leg1_metrics_v2.json :: cells.{rand-d3-s3}.N1000.s1.macro_auc_ovr :: sha256 b56bf70be9982044 :: one run
\newcommand{\FtAucJciiEFourRandDrawThree}{0.9686}  % experiments/FIGS/data/w2b_leg1_metrics_v2.json :: cells.{rand-d3-s3}.N10000.s1.macro_auc_ovr :: sha256 b56bf70be9982044 :: one run
\newcommand{\FtAucJciiEFiveRandDrawThree}{0.97605}  % experiments/FIGS/data/w2b_leg1_metrics_v2.json :: cells.{rand-d3-s3}.N100000.s1.macro_auc_ovr :: sha256 b56bf70be9982044 :: one run
\newcommand{\FtAucJciiESixRandDrawThree}{0.97950}  % experiments/FIGS/data/w2b_leg1_metrics_v2.json :: cells.{rand-d3-s3}.N1000000.s1.macro_auc_ovr :: sha256 b56bf70be9982044 :: one run
\newcommand{\FtAccJciiEThreeOneeighteight}{0.445\,\ensuremath{\pm}\,0.004}  % experiments/FIGS/data/w2b_leg1_metrics_v2.json :: cells.{l188-s1,l188-s2,l188-s3,l188-s4,l188-s5}.N1000.s1.accuracy :: sha256 b56bf70be9982044 :: mean +- SD over the row's 5 pretrained models
\newcommand{\FtAccJciiEFourOneeighteight}{0.5199\,\ensuremath{\pm}\,0.0018}  % experiments/FIGS/data/w2b_leg1_metrics_v2.json :: cells.{l188-s1,l188-s2,l188-s3,l188-s4,l188-s5}.N10000.s1.accuracy :: sha256 b56bf70be9982044 :: mean +- SD over the row's 5 pretrained models
\newcommand{\FtAccJciiEFiveOneeighteight}{0.5379\,\ensuremath{\pm}\,0.0010}  % experiments/FIGS/data/w2b_leg1_metrics_v2.json :: cells.{l188-s1,l188-s2,l188-s3,l188-s4,l188-s5}.N100000.s1.accuracy :: sha256 b56bf70be9982044 :: mean +- SD over the row's 5 pretrained models
\newcommand{\FtAccJciiESixOneeighteight}{0.5443\,\ensuremath{\pm}\,0.0003}  % experiments/FIGS/data/w2b_leg1_metrics_v2.json :: cells.{l188-s1,l188-s2,l188-s3,l188-s4,l188-s5}.N1000000.s1.accuracy :: sha256 b56bf70be9982044 :: mean +- SD over the row's 5 pretrained models
\newcommand{\FtAccJciiEThreeOnesixtwo}{0.448\,\ensuremath{\pm}\,0.006}  % experiments/FIGS/data/w2b_leg1_metrics_v2.json :: cells.{l162-s1b,l162-s2,l162-s3,l162-s4,l162-s5}.N1000.s1.accuracy :: sha256 b56bf70be9982044 :: mean +- SD over the row's 5 pretrained models
\newcommand{\FtAccJciiEFourOnesixtwo}{0.5212\,\ensuremath{\pm}\,0.0013}  % experiments/FIGS/data/w2b_leg1_metrics_v2.json :: cells.{l162-s1b,l162-s2,l162-s3,l162-s4,l162-s5}.N10000.s1.accuracy :: sha256 b56bf70be9982044 :: mean +- SD over the row's 5 pretrained models
\newcommand{\FtAccJciiEFiveOnesixtwo}{0.5387\,\ensuremath{\pm}\,0.0008}  % experiments/FIGS/data/w2b_leg1_metrics_v2.json :: cells.{l162-s1b,l162-s2,l162-s3,l162-s4,l162-s5}.N100000.s1.accuracy :: sha256 b56bf70be9982044 :: mean +- SD over the row's 5 pretrained models
\newcommand{\FtAccJciiESixOnesixtwo}{0.5440\,\ensuremath{\pm}\,0.0006}  % experiments/FIGS/data/w2b_leg1_metrics_v2.json :: cells.{l162-s1b,l162-s2,l162-s3,l162-s4,l162-s5}.N1000000.s1.accuracy :: sha256 b56bf70be9982044 :: mean +- SD over the row's 5 pretrained models
\newcommand{\FtAccJciiEThreeFourthree}{0.380\,\ensuremath{\pm}\,0.005}  % experiments/FIGS/data/w2b_leg1_metrics_v2.json :: cells.{r42q1-s1,r42q1-s2,r42q1-s3,r42q1-s4,r42q1-s5}.N1000.s1.accuracy :: sha256 b56bf70be9982044 :: mean +- SD over the row's 5 pretrained models
\newcommand{\FtAccJciiEFourFourthree}{0.481\,\ensuremath{\pm}\,0.004}  % experiments/FIGS/data/w2b_leg1_metrics_v2.json :: cells.{r42q1-s1,r42q1-s2,r42q1-s3,r42q1-s4,r42q1-s5}.N10000.s1.accuracy :: sha256 b56bf70be9982044 :: mean +- SD over the row's 5 pretrained models
\newcommand{\FtAccJciiEFiveFourthree}{0.5219\,\ensuremath{\pm}\,0.0017}  % experiments/FIGS/data/w2b_leg1_metrics_v2.json :: cells.{r42q1-s1,r42q1-s2,r42q1-s3,r42q1-s4,r42q1-s5}.N100000.s1.accuracy :: sha256 b56bf70be9982044 :: mean +- SD over the row's 5 pretrained models
\newcommand{\FtAccJciiESixFourthree}{0.5362\,\ensuremath{\pm}\,0.0007}  % experiments/FIGS/data/w2b_leg1_metrics_v2.json :: cells.{r42q1-s1,r42q1-s2,r42q1-s3,r42q1-s4,r42q1-s5}.N1000000.s1.accuracy :: sha256 b56bf70be9982044 :: mean +- SD over the row's 5 pretrained models
\newcommand{\FtAccJciiEThreeOneseven}{0.332\,\ensuremath{\pm}\,0.004}  % experiments/FIGS/data/w2b_leg1_metrics_v2.json :: cells.{r16q1-s1,r16q1-s2,r16q1-s3,r16q1-s4,r16q1-s5}.N1000.s1.accuracy :: sha256 b56bf70be9982044 :: mean +- SD over the row's 5 pretrained models
\newcommand{\FtAccJciiEFourOneseven}{0.447\,\ensuremath{\pm}\,0.013}  % experiments/FIGS/data/w2b_leg1_metrics_v2.json :: cells.{r16q1-s1,r16q1-s2,r16q1-s3,r16q1-s4,r16q1-s5}.N10000.s1.accuracy :: sha256 b56bf70be9982044 :: mean +- SD over the row's 5 pretrained models
\newcommand{\FtAccJciiEFiveOneseven}{0.5078\,\ensuremath{\pm}\,0.0011}  % experiments/FIGS/data/w2b_leg1_metrics_v2.json :: cells.{r16q1-s1,r16q1-s2,r16q1-s3,r16q1-s4,r16q1-s5}.N100000.s1.accuracy :: sha256 b56bf70be9982044 :: mean +- SD over the row's 5 pretrained models
\newcommand{\FtAccJciiESixOneseven}{0.5227\,\ensuremath{\pm}\,0.0010}  % experiments/FIGS/data/w2b_leg1_metrics_v2.json :: cells.{r16q1-s1,r16q1-s2,r16q1-s3,r16q1-s4,r16q1-s5}.N1000000.s1.accuracy :: sha256 b56bf70be9982044 :: mean +- SD over the row's 5 pretrained models
\newcommand{\FtAccJciiEThreeOnesixtwoMass}{0.451\,\ensuremath{\pm}\,0.007}  % experiments/FIGS/data/w2b_leg1_metrics_v2.json :: cells.{l162mass-s1,l162mass-s2,l162mass-s3,l162mass-s4,l162mass-s5}.N1000.s1.accuracy :: sha256 b56bf70be9982044 :: mean +- SD over the row's 5 pretrained models
\newcommand{\FtAccJciiEFourOnesixtwoMass}{0.515\,\ensuremath{\pm}\,0.004}  % experiments/FIGS/data/w2b_leg1_metrics_v2.json :: cells.{l162mass-s1,l162mass-s2,l162mass-s3,l162mass-s4,l162mass-s5}.N10000.s1.accuracy :: sha256 b56bf70be9982044 :: mean +- SD over the row's 5 pretrained models
\newcommand{\FtAccJciiEFiveOnesixtwoMass}{0.5357\,\ensuremath{\pm}\,0.0007}  % experiments/FIGS/data/w2b_leg1_metrics_v2.json :: cells.{l162mass-s1,l162mass-s2,l162mass-s3,l162mass-s4,l162mass-s5}.N100000.s1.accuracy :: sha256 b56bf70be9982044 :: mean +- SD over the row's 5 pretrained models
\newcommand{\FtAccJciiESixOnesixtwoMass}{0.5409\,\ensuremath{\pm}\,0.0010}  % experiments/FIGS/data/w2b_leg1_metrics_v2.json :: cells.{l162mass-s1,l162mass-s2,l162mass-s3,l162mass-s4,l162mass-s5}.N1000000.s1.accuracy :: sha256 b56bf70be9982044 :: mean +- SD over the row's 5 pretrained models
\newcommand{\FtAccJciiEThreeOnesevenMass}{0.331\,\ensuremath{\pm}\,0.005}  % experiments/FIGS/data/w2b_leg1_metrics_v2.json :: cells.{r16q1mass-s1,r16q1mass-s2,r16q1mass-s3,r16q1mass-s4,r16q1mass-s5}.N1000.s1.accuracy :: sha256 b56bf70be9982044 :: mean +- SD over the row's 5 pretrained models
\newcommand{\FtAccJciiEFourOnesevenMass}{0.4507\,\ensuremath{\pm}\,0.0034}  % experiments/FIGS/data/w2b_leg1_metrics_v2.json :: cells.{r16q1mass-s1,r16q1mass-s2,r16q1mass-s3,r16q1mass-s4,r16q1mass-s5}.N10000.s1.accuracy :: sha256 b56bf70be9982044 :: mean +- SD over the row's 5 pretrained models
\newcommand{\FtAccJciiEFiveOnesevenMass}{0.4993\,\ensuremath{\pm}\,0.0022}  % experiments/FIGS/data/w2b_leg1_metrics_v2.json :: cells.{r16q1mass-s1,r16q1mass-s2,r16q1mass-s3,r16q1mass-s4,r16q1mass-s5}.N100000.s1.accuracy :: sha256 b56bf70be9982044 :: mean +- SD over the row's 5 pretrained models
\newcommand{\FtAccJciiESixOnesevenMass}{0.5166\,\ensuremath{\pm}\,0.0013}  % experiments/FIGS/data/w2b_leg1_metrics_v2.json :: cells.{r16q1mass-s1,r16q1mass-s2,r16q1mass-s3,r16q1mass-s4,r16q1mass-s5}.N1000000.s1.accuracy :: sha256 b56bf70be9982044 :: mean +- SD over the row's 5 pretrained models
\newcommand{\FtAccJciiEThreeRandDrawOne}{0.424}  % experiments/FIGS/data/w2b_leg1_metrics_v2.json :: cells.{rand-d1-s1b}.N1000.s1.accuracy :: sha256 b56bf70be9982044 :: one run
\newcommand{\FtAccJciiEFourRandDrawOne}{0.5070}  % experiments/FIGS/data/w2b_leg1_metrics_v2.json :: cells.{rand-d1-s1b}.N10000.s1.accuracy :: sha256 b56bf70be9982044 :: one run
\newcommand{\FtAccJciiEFiveRandDrawOne}{0.5226}  % experiments/FIGS/data/w2b_leg1_metrics_v2.json :: cells.{rand-d1-s1b}.N100000.s1.accuracy :: sha256 b56bf70be9982044 :: one run
\newcommand{\FtAccJciiESixRandDrawOne}{0.53354}  % experiments/FIGS/data/w2b_leg1_metrics_v2.json :: cells.{rand-d1-s1b}.N1000000.s1.accuracy :: sha256 b56bf70be9982044 :: one run
\newcommand{\FtAccJciiEThreeRandDrawTwo}{0.420}  % experiments/FIGS/data/w2b_leg1_metrics_v2.json :: cells.{rand-d2-s2}.N1000.s1.accuracy :: sha256 b56bf70be9982044 :: one run
\newcommand{\FtAccJciiEFourRandDrawTwo}{0.5074}  % experiments/FIGS/data/w2b_leg1_metrics_v2.json :: cells.{rand-d2-s2}.N10000.s1.accuracy :: sha256 b56bf70be9982044 :: one run
\newcommand{\FtAccJciiEFiveRandDrawTwo}{0.5243}  % experiments/FIGS/data/w2b_leg1_metrics_v2.json :: cells.{rand-d2-s2}.N100000.s1.accuracy :: sha256 b56bf70be9982044 :: one run
\newcommand{\FtAccJciiESixRandDrawTwo}{0.53441}  % experiments/FIGS/data/w2b_leg1_metrics_v2.json :: cells.{rand-d2-s2}.N1000000.s1.accuracy :: sha256 b56bf70be9982044 :: one run
\newcommand{\FtAccJciiEThreeRandDrawThree}{0.423}  % experiments/FIGS/data/w2b_leg1_metrics_v2.json :: cells.{rand-d3-s3}.N1000.s1.accuracy :: sha256 b56bf70be9982044 :: one run
\newcommand{\FtAccJciiEFourRandDrawThree}{0.5114}  % experiments/FIGS/data/w2b_leg1_metrics_v2.json :: cells.{rand-d3-s3}.N10000.s1.accuracy :: sha256 b56bf70be9982044 :: one run
\newcommand{\FtAccJciiEFiveRandDrawThree}{0.5245}  % experiments/FIGS/data/w2b_leg1_metrics_v2.json :: cells.{rand-d3-s3}.N100000.s1.accuracy :: sha256 b56bf70be9982044 :: one run
\newcommand{\FtAccJciiESixRandDrawThree}{0.53473}  % experiments/FIGS/data/w2b_leg1_metrics_v2.json :: cells.{rand-d3-s3}.N1000000.s1.accuracy :: sha256 b56bf70be9982044 :: one run
\newcommand{\FtOmaRatioJciiEThreeOnesixtwoOverOneeighteight}{0.96}  % experiments/FIGS/data/w2b_leg1_metrics_v2.json :: cells.*.N1000.s1.macro_auc_ovr :: sha256 b56bf70be9982044 :: ratio of the seed means of 1 - macro AUC, this row over the reference
\newcommand{\FtOmaRatioJciiEThreeFourthreeOverOneeighteight}{0.84}  % experiments/FIGS/data/w2b_leg1_metrics_v2.json :: cells.*.N1000.s1.macro_auc_ovr :: sha256 b56bf70be9982044 :: ratio of the seed means of 1 - macro AUC, this row over the reference
\newcommand{\FtOmaRatioJciiEThreeOnesevenOverOneeighteight}{0.97}  % experiments/FIGS/data/w2b_leg1_metrics_v2.json :: cells.*.N1000.s1.macro_auc_ovr :: sha256 b56bf70be9982044 :: ratio of the seed means of 1 - macro AUC, this row over the reference
\newcommand{\FtOmaRatioJciiEThreeOnesixtwoMassOverOneeighteight}{0.81}  % experiments/FIGS/data/w2b_leg1_metrics_v2.json :: cells.*.N1000.s1.macro_auc_ovr :: sha256 b56bf70be9982044 :: ratio of the seed means of 1 - macro AUC, this row over the reference
\newcommand{\FtOmaRatioJciiEThreeOnesevenMassOverOneeighteight}{0.98}  % experiments/FIGS/data/w2b_leg1_metrics_v2.json :: cells.*.N1000.s1.macro_auc_ovr :: sha256 b56bf70be9982044 :: ratio of the seed means of 1 - macro AUC, this row over the reference
\newcommand{\FtOmaRatioJciiEThreeRandDrawOneOverOneeighteight}{0.80}  % experiments/FIGS/data/w2b_leg1_metrics_v2.json :: cells.*.N1000.s1.macro_auc_ovr :: sha256 b56bf70be9982044 :: ratio of the seed means of 1 - macro AUC, this row over the reference
\newcommand{\FtOmaRatioJciiEThreeRandDrawTwoOverOneeighteight}{0.87}  % experiments/FIGS/data/w2b_leg1_metrics_v2.json :: cells.*.N1000.s1.macro_auc_ovr :: sha256 b56bf70be9982044 :: ratio of the seed means of 1 - macro AUC, this row over the reference
\newcommand{\FtOmaRatioJciiEThreeRandDrawThreeOverOneeighteight}{0.83}  % experiments/FIGS/data/w2b_leg1_metrics_v2.json :: cells.*.N1000.s1.macro_auc_ovr :: sha256 b56bf70be9982044 :: ratio of the seed means of 1 - macro AUC, this row over the reference
\newcommand{\FtOmaRatioJciiEThreeOneeighteightOverOnesixtwo}{1.05}  % experiments/FIGS/data/w2b_leg1_metrics_v2.json :: cells.*.N1000.s1.macro_auc_ovr :: sha256 b56bf70be9982044 :: ratio of the seed means of 1 - macro AUC, this row over the reference
\newcommand{\FtOmaRatioJciiEThreeFourthreeOverOnesixtwo}{0.88}  % experiments/FIGS/data/w2b_leg1_metrics_v2.json :: cells.*.N1000.s1.macro_auc_ovr :: sha256 b56bf70be9982044 :: ratio of the seed means of 1 - macro AUC, this row over the reference
\newcommand{\FtOmaRatioJciiEThreeOnesevenOverOnesixtwo}{1.02}  % experiments/FIGS/data/w2b_leg1_metrics_v2.json :: cells.*.N1000.s1.macro_auc_ovr :: sha256 b56bf70be9982044 :: ratio of the seed means of 1 - macro AUC, this row over the reference
\newcommand{\FtOmaRatioJciiEThreeOnesixtwoMassOverOnesixtwo}{0.85}  % experiments/FIGS/data/w2b_leg1_metrics_v2.json :: cells.*.N1000.s1.macro_auc_ovr :: sha256 b56bf70be9982044 :: ratio of the seed means of 1 - macro AUC, this row over the reference
\newcommand{\FtOmaRatioJciiEThreeOnesevenMassOverOnesixtwo}{1.02}  % experiments/FIGS/data/w2b_leg1_metrics_v2.json :: cells.*.N1000.s1.macro_auc_ovr :: sha256 b56bf70be9982044 :: ratio of the seed means of 1 - macro AUC, this row over the reference
\newcommand{\FtOmaRatioJciiEThreeRandDrawOneOverOnesixtwo}{0.84}  % experiments/FIGS/data/w2b_leg1_metrics_v2.json :: cells.*.N1000.s1.macro_auc_ovr :: sha256 b56bf70be9982044 :: ratio of the seed means of 1 - macro AUC, this row over the reference
\newcommand{\FtOmaRatioJciiEThreeRandDrawTwoOverOnesixtwo}{0.91}  % experiments/FIGS/data/w2b_leg1_metrics_v2.json :: cells.*.N1000.s1.macro_auc_ovr :: sha256 b56bf70be9982044 :: ratio of the seed means of 1 - macro AUC, this row over the reference
\newcommand{\FtOmaRatioJciiEThreeRandDrawThreeOverOnesixtwo}{0.87}  % experiments/FIGS/data/w2b_leg1_metrics_v2.json :: cells.*.N1000.s1.macro_auc_ovr :: sha256 b56bf70be9982044 :: ratio of the seed means of 1 - macro AUC, this row over the reference
\newcommand{\FtOmaRatioJciiEFourOnesixtwoOverOneeighteight}{0.99}  % experiments/FIGS/data/w2b_leg1_metrics_v2.json :: cells.*.N10000.s1.macro_auc_ovr :: sha256 b56bf70be9982044 :: ratio of the seed means of 1 - macro AUC, this row over the reference
\newcommand{\FtOmaRatioJciiEFourFourthreeOverOneeighteight}{1.20}  % experiments/FIGS/data/w2b_leg1_metrics_v2.json :: cells.*.N10000.s1.macro_auc_ovr :: sha256 b56bf70be9982044 :: ratio of the seed means of 1 - macro AUC, this row over the reference
\newcommand{\FtOmaRatioJciiEFourOnesevenOverOneeighteight}{1.32}  % experiments/FIGS/data/w2b_leg1_metrics_v2.json :: cells.*.N10000.s1.macro_auc_ovr :: sha256 b56bf70be9982044 :: ratio of the seed means of 1 - macro AUC, this row over the reference
\newcommand{\FtOmaRatioJciiEFourOnesixtwoMassOverOneeighteight}{0.99}  % experiments/FIGS/data/w2b_leg1_metrics_v2.json :: cells.*.N10000.s1.macro_auc_ovr :: sha256 b56bf70be9982044 :: ratio of the seed means of 1 - macro AUC, this row over the reference
\newcommand{\FtOmaRatioJciiEFourOnesevenMassOverOneeighteight}{1.30}  % experiments/FIGS/data/w2b_leg1_metrics_v2.json :: cells.*.N10000.s1.macro_auc_ovr :: sha256 b56bf70be9982044 :: ratio of the seed means of 1 - macro AUC, this row over the reference
\newcommand{\FtOmaRatioJciiEFourRandDrawOneOverOneeighteight}{1.17}  % experiments/FIGS/data/w2b_leg1_metrics_v2.json :: cells.*.N10000.s1.macro_auc_ovr :: sha256 b56bf70be9982044 :: ratio of the seed means of 1 - macro AUC, this row over the reference
\newcommand{\FtOmaRatioJciiEFourRandDrawTwoOverOneeighteight}{1.17}  % experiments/FIGS/data/w2b_leg1_metrics_v2.json :: cells.*.N10000.s1.macro_auc_ovr :: sha256 b56bf70be9982044 :: ratio of the seed means of 1 - macro AUC, this row over the reference
\newcommand{\FtOmaRatioJciiEFourRandDrawThreeOverOneeighteight}{1.13}  % experiments/FIGS/data/w2b_leg1_metrics_v2.json :: cells.*.N10000.s1.macro_auc_ovr :: sha256 b56bf70be9982044 :: ratio of the seed means of 1 - macro AUC, this row over the reference
\newcommand{\FtOmaRatioJciiEFourOneeighteightOverOnesixtwo}{1.01}  % experiments/FIGS/data/w2b_leg1_metrics_v2.json :: cells.*.N10000.s1.macro_auc_ovr :: sha256 b56bf70be9982044 :: ratio of the seed means of 1 - macro AUC, this row over the reference
\newcommand{\FtOmaRatioJciiEFourFourthreeOverOnesixtwo}{1.21}  % experiments/FIGS/data/w2b_leg1_metrics_v2.json :: cells.*.N10000.s1.macro_auc_ovr :: sha256 b56bf70be9982044 :: ratio of the seed means of 1 - macro AUC, this row over the reference
\newcommand{\FtOmaRatioJciiEFourOnesevenOverOnesixtwo}{1.32}  % experiments/FIGS/data/w2b_leg1_metrics_v2.json :: cells.*.N10000.s1.macro_auc_ovr :: sha256 b56bf70be9982044 :: ratio of the seed means of 1 - macro AUC, this row over the reference
\newcommand{\FtOmaRatioJciiEFourOnesixtwoMassOverOnesixtwo}{1.00}  % experiments/FIGS/data/w2b_leg1_metrics_v2.json :: cells.*.N10000.s1.macro_auc_ovr :: sha256 b56bf70be9982044 :: ratio of the seed means of 1 - macro AUC, this row over the reference
\newcommand{\FtOmaRatioJciiEFourOnesevenMassOverOnesixtwo}{1.31}  % experiments/FIGS/data/w2b_leg1_metrics_v2.json :: cells.*.N10000.s1.macro_auc_ovr :: sha256 b56bf70be9982044 :: ratio of the seed means of 1 - macro AUC, this row over the reference
\newcommand{\FtOmaRatioJciiEFourRandDrawOneOverOnesixtwo}{1.18}  % experiments/FIGS/data/w2b_leg1_metrics_v2.json :: cells.*.N10000.s1.macro_auc_ovr :: sha256 b56bf70be9982044 :: ratio of the seed means of 1 - macro AUC, this row over the reference
\newcommand{\FtOmaRatioJciiEFourRandDrawTwoOverOnesixtwo}{1.18}  % experiments/FIGS/data/w2b_leg1_metrics_v2.json :: cells.*.N10000.s1.macro_auc_ovr :: sha256 b56bf70be9982044 :: ratio of the seed means of 1 - macro AUC, this row over the reference
\newcommand{\FtOmaRatioJciiEFourRandDrawThreeOverOnesixtwo}{1.14}  % experiments/FIGS/data/w2b_leg1_metrics_v2.json :: cells.*.N10000.s1.macro_auc_ovr :: sha256 b56bf70be9982044 :: ratio of the seed means of 1 - macro AUC, this row over the reference
\newcommand{\FtOmaRatioJciiEFiveOnesixtwoOverOneeighteight}{1.00}  % experiments/FIGS/data/w2b_leg1_metrics_v2.json :: cells.*.N100000.s1.macro_auc_ovr :: sha256 b56bf70be9982044 :: ratio of the seed means of 1 - macro AUC, this row over the reference
\newcommand{\FtOmaRatioJciiEFiveFourthreeOverOneeighteight}{1.09}  % experiments/FIGS/data/w2b_leg1_metrics_v2.json :: cells.*.N100000.s1.macro_auc_ovr :: sha256 b56bf70be9982044 :: ratio of the seed means of 1 - macro AUC, this row over the reference
\newcommand{\FtOmaRatioJciiEFiveOnesevenOverOneeighteight}{1.22}  % experiments/FIGS/data/w2b_leg1_metrics_v2.json :: cells.*.N100000.s1.macro_auc_ovr :: sha256 b56bf70be9982044 :: ratio of the seed means of 1 - macro AUC, this row over the reference
\newcommand{\FtOmaRatioJciiEFiveOnesixtwoMassOverOneeighteight}{1.01}  % experiments/FIGS/data/w2b_leg1_metrics_v2.json :: cells.*.N100000.s1.macro_auc_ovr :: sha256 b56bf70be9982044 :: ratio of the seed means of 1 - macro AUC, this row over the reference
\newcommand{\FtOmaRatioJciiEFiveOnesevenMassOverOneeighteight}{1.26}  % experiments/FIGS/data/w2b_leg1_metrics_v2.json :: cells.*.N100000.s1.macro_auc_ovr :: sha256 b56bf70be9982044 :: ratio of the seed means of 1 - macro AUC, this row over the reference
\newcommand{\FtOmaRatioJciiEFiveRandDrawOneOverOneeighteight}{1.15}  % experiments/FIGS/data/w2b_leg1_metrics_v2.json :: cells.*.N100000.s1.macro_auc_ovr :: sha256 b56bf70be9982044 :: ratio of the seed means of 1 - macro AUC, this row over the reference
\newcommand{\FtOmaRatioJciiEFiveRandDrawTwoOverOneeighteight}{1.12}  % experiments/FIGS/data/w2b_leg1_metrics_v2.json :: cells.*.N100000.s1.macro_auc_ovr :: sha256 b56bf70be9982044 :: ratio of the seed means of 1 - macro AUC, this row over the reference
\newcommand{\FtOmaRatioJciiEFiveRandDrawThreeOverOneeighteight}{1.15}  % experiments/FIGS/data/w2b_leg1_metrics_v2.json :: cells.*.N100000.s1.macro_auc_ovr :: sha256 b56bf70be9982044 :: ratio of the seed means of 1 - macro AUC, this row over the reference
\newcommand{\FtOmaRatioJciiEFiveOneeighteightOverOnesixtwo}{1.00}  % experiments/FIGS/data/w2b_leg1_metrics_v2.json :: cells.*.N100000.s1.macro_auc_ovr :: sha256 b56bf70be9982044 :: ratio of the seed means of 1 - macro AUC, this row over the reference
\newcommand{\FtOmaRatioJciiEFiveFourthreeOverOnesixtwo}{1.09}  % experiments/FIGS/data/w2b_leg1_metrics_v2.json :: cells.*.N100000.s1.macro_auc_ovr :: sha256 b56bf70be9982044 :: ratio of the seed means of 1 - macro AUC, this row over the reference
\newcommand{\FtOmaRatioJciiEFiveOnesevenOverOnesixtwo}{1.22}  % experiments/FIGS/data/w2b_leg1_metrics_v2.json :: cells.*.N100000.s1.macro_auc_ovr :: sha256 b56bf70be9982044 :: ratio of the seed means of 1 - macro AUC, this row over the reference
\newcommand{\FtOmaRatioJciiEFiveOnesixtwoMassOverOnesixtwo}{1.02}  % experiments/FIGS/data/w2b_leg1_metrics_v2.json :: cells.*.N100000.s1.macro_auc_ovr :: sha256 b56bf70be9982044 :: ratio of the seed means of 1 - macro AUC, this row over the reference
\newcommand{\FtOmaRatioJciiEFiveOnesevenMassOverOnesixtwo}{1.26}  % experiments/FIGS/data/w2b_leg1_metrics_v2.json :: cells.*.N100000.s1.macro_auc_ovr :: sha256 b56bf70be9982044 :: ratio of the seed means of 1 - macro AUC, this row over the reference
\newcommand{\FtOmaRatioJciiEFiveRandDrawOneOverOnesixtwo}{1.15}  % experiments/FIGS/data/w2b_leg1_metrics_v2.json :: cells.*.N100000.s1.macro_auc_ovr :: sha256 b56bf70be9982044 :: ratio of the seed means of 1 - macro AUC, this row over the reference
\newcommand{\FtOmaRatioJciiEFiveRandDrawTwoOverOnesixtwo}{1.13}  % experiments/FIGS/data/w2b_leg1_metrics_v2.json :: cells.*.N100000.s1.macro_auc_ovr :: sha256 b56bf70be9982044 :: ratio of the seed means of 1 - macro AUC, this row over the reference
\newcommand{\FtOmaRatioJciiEFiveRandDrawThreeOverOnesixtwo}{1.16}  % experiments/FIGS/data/w2b_leg1_metrics_v2.json :: cells.*.N100000.s1.macro_auc_ovr :: sha256 b56bf70be9982044 :: ratio of the seed means of 1 - macro AUC, this row over the reference
\newcommand{\FtOmaRatioJciiESixOnesixtwoOverOneeighteight}{1.00}  % experiments/FIGS/data/w2b_leg1_metrics_v2.json :: cells.*.N1000000.s1.macro_auc_ovr :: sha256 b56bf70be9982044 :: ratio of the seed means of 1 - macro AUC, this row over the reference
\newcommand{\FtOmaRatioJciiESixFourthreeOverOneeighteight}{1.05}  % experiments/FIGS/data/w2b_leg1_metrics_v2.json :: cells.*.N1000000.s1.macro_auc_ovr :: sha256 b56bf70be9982044 :: ratio of the seed means of 1 - macro AUC, this row over the reference
\newcommand{\FtOmaRatioJciiESixOnesevenOverOneeighteight}{1.14}  % experiments/FIGS/data/w2b_leg1_metrics_v2.json :: cells.*.N1000000.s1.macro_auc_ovr :: sha256 b56bf70be9982044 :: ratio of the seed means of 1 - macro AUC, this row over the reference
\newcommand{\FtOmaRatioJciiESixOnesixtwoMassOverOneeighteight}{1.02}  % experiments/FIGS/data/w2b_leg1_metrics_v2.json :: cells.*.N1000000.s1.macro_auc_ovr :: sha256 b56bf70be9982044 :: ratio of the seed means of 1 - macro AUC, this row over the reference
\newcommand{\FtOmaRatioJciiESixOnesevenMassOverOneeighteight}{1.18}  % experiments/FIGS/data/w2b_leg1_metrics_v2.json :: cells.*.N1000000.s1.macro_auc_ovr :: sha256 b56bf70be9982044 :: ratio of the seed means of 1 - macro AUC, this row over the reference
\newcommand{\FtOmaRatioJciiESixRandDrawOneOverOneeighteight}{1.08}  % experiments/FIGS/data/w2b_leg1_metrics_v2.json :: cells.*.N1000000.s1.macro_auc_ovr :: sha256 b56bf70be9982044 :: ratio of the seed means of 1 - macro AUC, this row over the reference
\newcommand{\FtOmaRatioJciiESixRandDrawTwoOverOneeighteight}{1.07}  % experiments/FIGS/data/w2b_leg1_metrics_v2.json :: cells.*.N1000000.s1.macro_auc_ovr :: sha256 b56bf70be9982044 :: ratio of the seed means of 1 - macro AUC, this row over the reference
\newcommand{\FtOmaRatioJciiESixRandDrawThreeOverOneeighteight}{1.07}  % experiments/FIGS/data/w2b_leg1_metrics_v2.json :: cells.*.N1000000.s1.macro_auc_ovr :: sha256 b56bf70be9982044 :: ratio of the seed means of 1 - macro AUC, this row over the reference
\newcommand{\FtOmaRatioJciiESixOneeighteightOverOnesixtwo}{1.00}  % experiments/FIGS/data/w2b_leg1_metrics_v2.json :: cells.*.N1000000.s1.macro_auc_ovr :: sha256 b56bf70be9982044 :: ratio of the seed means of 1 - macro AUC, this row over the reference
\newcommand{\FtOmaRatioJciiESixFourthreeOverOnesixtwo}{1.04}  % experiments/FIGS/data/w2b_leg1_metrics_v2.json :: cells.*.N1000000.s1.macro_auc_ovr :: sha256 b56bf70be9982044 :: ratio of the seed means of 1 - macro AUC, this row over the reference
\newcommand{\FtOmaRatioJciiESixOnesevenOverOnesixtwo}{1.14}  % experiments/FIGS/data/w2b_leg1_metrics_v2.json :: cells.*.N1000000.s1.macro_auc_ovr :: sha256 b56bf70be9982044 :: ratio of the seed means of 1 - macro AUC, this row over the reference
\newcommand{\FtOmaRatioJciiESixOnesixtwoMassOverOnesixtwo}{1.02}  % experiments/FIGS/data/w2b_leg1_metrics_v2.json :: cells.*.N1000000.s1.macro_auc_ovr :: sha256 b56bf70be9982044 :: ratio of the seed means of 1 - macro AUC, this row over the reference
\newcommand{\FtOmaRatioJciiESixOnesevenMassOverOnesixtwo}{1.18}  % experiments/FIGS/data/w2b_leg1_metrics_v2.json :: cells.*.N1000000.s1.macro_auc_ovr :: sha256 b56bf70be9982044 :: ratio of the seed means of 1 - macro AUC, this row over the reference
\newcommand{\FtOmaRatioJciiESixRandDrawOneOverOnesixtwo}{1.07}  % experiments/FIGS/data/w2b_leg1_metrics_v2.json :: cells.*.N1000000.s1.macro_auc_ovr :: sha256 b56bf70be9982044 :: ratio of the seed means of 1 - macro AUC, this row over the reference
\newcommand{\FtOmaRatioJciiESixRandDrawTwoOverOnesixtwo}{1.07}  % experiments/FIGS/data/w2b_leg1_metrics_v2.json :: cells.*.N1000000.s1.macro_auc_ovr :: sha256 b56bf70be9982044 :: ratio of the seed means of 1 - macro AUC, this row over the reference
\newcommand{\FtOmaRatioJciiESixRandDrawThreeOverOnesixtwo}{1.07}  % experiments/FIGS/data/w2b_leg1_metrics_v2.json :: cells.*.N1000000.s1.macro_auc_ovr :: sha256 b56bf70be9982044 :: ratio of the seed means of 1 - macro AUC, this row over the reference
\newcommand{\PairedFtAccJciiEThreeOnesevenMinusOneeighteight}{\ensuremath{-}0.1132~[\ensuremath{-}0.1214, \ensuremath{-}0.1051]}  % experiments/FIGS/data/w2b_leg1_metrics_v2.json :: cells.{r16q1-s1,r16q1-s2,r16q1-s3,r16q1-s4,r16q1-s5,l188-s1,l188-s2,l188-s3,l188-s4,l188-s5}.N1000.s1.accuracy :: sha256 b56bf70be9982044 :: accuracy, coarse minus fine, paired by run [95% Student-t interval]
\newcommand{\PairedFtAccJciiEFourOnesevenMinusOneeighteight}{\ensuremath{-}0.0730~[\ensuremath{-}0.0902, \ensuremath{-}0.0558]}  % experiments/FIGS/data/w2b_leg1_metrics_v2.json :: cells.{r16q1-s1,r16q1-s2,r16q1-s3,r16q1-s4,r16q1-s5,l188-s1,l188-s2,l188-s3,l188-s4,l188-s5}.N10000.s1.accuracy :: sha256 b56bf70be9982044 :: accuracy, coarse minus fine, paired by run [95% Student-t interval]
\newcommand{\PairedFtAccJciiEFiveOnesevenMinusOneeighteight}{\ensuremath{-}0.0301~[\ensuremath{-}0.0319, \ensuremath{-}0.0284]}  % experiments/FIGS/data/w2b_leg1_metrics_v2.json :: cells.{r16q1-s1,r16q1-s2,r16q1-s3,r16q1-s4,r16q1-s5,l188-s1,l188-s2,l188-s3,l188-s4,l188-s5}.N100000.s1.accuracy :: sha256 b56bf70be9982044 :: accuracy, coarse minus fine, paired by run [95% Student-t interval]
\newcommand{\PairedFtAccJciiESixOnesevenMinusOneeighteight}{\ensuremath{-}0.0216~[\ensuremath{-}0.0232, \ensuremath{-}0.0199]}  % experiments/FIGS/data/w2b_leg1_metrics_v2.json :: cells.{r16q1-s1,r16q1-s2,r16q1-s3,r16q1-s4,r16q1-s5,l188-s1,l188-s2,l188-s3,l188-s4,l188-s5}.N1000000.s1.accuracy :: sha256 b56bf70be9982044 :: accuracy, coarse minus fine, paired by run [95% Student-t interval]
\newcommand{\FtNClassesJcii}{162}  % experiments/FIGS/data/w2b_leg1_metrics_v2.json :: cells.*.*.s1.n_classes_present :: sha256 b56bf70be9982044 :: classes of the fine-tuning task
\newcommand{\FtNTestAucJcii}{125{,}000}  % experiments/FIGS/data/w2b_leg1_metrics_v2.json :: cells.*.*.s1.n_jets_auc :: sha256 b56bf70be9982044 :: test jets the macro AUC is computed on
\newcommand{\FtNTestAccJcii}{500{,}000}  % experiments/FIGS/data/w2b_leg1_metrics_v2.json :: cells.*.*.s1.n_jets :: sha256 b56bf70be9982044 :: test jets the accuracy is computed on
\newcommand{\FtAucJcEThreeOneeighteight}{0.950\,\ensuremath{\pm}\,0.004}  % experiments/FIGS/data/w2b_leg2_metrics_v2.json :: cells.{l188-s1,l188-s2,l188-s3,l188-s4,l188-s5}.N1000.s1.macro_auc_ovr :: sha256 9cd059d3e2f51798 :: mean +- SD over the row's 5 pretrained models
\newcommand{\FtAucJcEFourOneeighteight}{0.9765\,\ensuremath{\pm}\,0.0008}  % experiments/FIGS/data/w2b_leg2_metrics_v2.json :: cells.{l188-s1,l188-s2,l188-s3,l188-s4,l188-s5}.N10000.s1.macro_auc_ovr :: sha256 9cd059d3e2f51798 :: mean +- SD over the row's 5 pretrained models
\newcommand{\FtAucJcEFiveOneeighteight}{0.9833\,\ensuremath{\pm}\,0.0002}  % experiments/FIGS/data/w2b_leg2_metrics_v2.json :: cells.{l188-s1,l188-s2,l188-s3,l188-s4,l188-s5}.N100000.s1.macro_auc_ovr :: sha256 9cd059d3e2f51798 :: mean +- SD over the row's 5 pretrained models
\newcommand{\FtAucJcESixOneeighteight}{0.9858\,\ensuremath{\pm}\,0.0001}  % experiments/FIGS/data/w2b_leg2_metrics_v2.json :: cells.{l188-s1,l188-s2,l188-s3,l188-s4,l188-s5}.N1000000.s1.macro_auc_ovr :: sha256 9cd059d3e2f51798 :: mean +- SD over the row's 5 pretrained models
\newcommand{\FtAucJcEThreeOnesixtwo}{0.9526\,\ensuremath{\pm}\,0.0020}  % experiments/FIGS/data/w2b_leg2_metrics_v2.json :: cells.{l162-s1b,l162-s2,l162-s3,l162-s4,l162-s5}.N1000.s1.macro_auc_ovr :: sha256 9cd059d3e2f51798 :: mean +- SD over the row's 5 pretrained models
\newcommand{\FtAucJcEFourOnesixtwo}{0.9770\,\ensuremath{\pm}\,0.0007}  % experiments/FIGS/data/w2b_leg2_metrics_v2.json :: cells.{l162-s1b,l162-s2,l162-s3,l162-s4,l162-s5}.N10000.s1.macro_auc_ovr :: sha256 9cd059d3e2f51798 :: mean +- SD over the row's 5 pretrained models
\newcommand{\FtAucJcEFiveOnesixtwo}{0.9836\,\ensuremath{\pm}\,0.0001}  % experiments/FIGS/data/w2b_leg2_metrics_v2.json :: cells.{l162-s1b,l162-s2,l162-s3,l162-s4,l162-s5}.N100000.s1.macro_auc_ovr :: sha256 9cd059d3e2f51798 :: mean +- SD over the row's 5 pretrained models
\newcommand{\FtAucJcESixOnesixtwo}{0.9858\,\ensuremath{\pm}\,0.0001}  % experiments/FIGS/data/w2b_leg2_metrics_v2.json :: cells.{l162-s1b,l162-s2,l162-s3,l162-s4,l162-s5}.N1000000.s1.macro_auc_ovr :: sha256 9cd059d3e2f51798 :: mean +- SD over the row's 5 pretrained models
\newcommand{\FtAucJcEThreeFourthree}{0.950\,\ensuremath{\pm}\,0.004}  % experiments/FIGS/data/w2b_leg2_metrics_v2.json :: cells.{r42q1-s1,r42q1-s2,r42q1-s3,r42q1-s4,r42q1-s5}.N1000.s1.macro_auc_ovr :: sha256 9cd059d3e2f51798 :: mean +- SD over the row's 5 pretrained models
\newcommand{\FtAucJcEFourFourthree}{0.9769\,\ensuremath{\pm}\,0.0005}  % experiments/FIGS/data/w2b_leg2_metrics_v2.json :: cells.{r42q1-s1,r42q1-s2,r42q1-s3,r42q1-s4,r42q1-s5}.N10000.s1.macro_auc_ovr :: sha256 9cd059d3e2f51798 :: mean +- SD over the row's 5 pretrained models
\newcommand{\FtAucJcEFiveFourthree}{0.9831\,\ensuremath{\pm}\,0.0002}  % experiments/FIGS/data/w2b_leg2_metrics_v2.json :: cells.{r42q1-s1,r42q1-s2,r42q1-s3,r42q1-s4,r42q1-s5}.N100000.s1.macro_auc_ovr :: sha256 9cd059d3e2f51798 :: mean +- SD over the row's 5 pretrained models
\newcommand{\FtAucJcESixFourthree}{0.9855\,\ensuremath{\pm}\,0.0001}  % experiments/FIGS/data/w2b_leg2_metrics_v2.json :: cells.{r42q1-s1,r42q1-s2,r42q1-s3,r42q1-s4,r42q1-s5}.N1000000.s1.macro_auc_ovr :: sha256 9cd059d3e2f51798 :: mean +- SD over the row's 5 pretrained models
\newcommand{\FtAucJcEThreeOneseven}{0.9470\,\ensuremath{\pm}\,0.0024}  % experiments/FIGS/data/w2b_leg2_metrics_v2.json :: cells.{r16q1-s1,r16q1-s2,r16q1-s3,r16q1-s4,r16q1-s5}.N1000.s1.macro_auc_ovr :: sha256 9cd059d3e2f51798 :: mean +- SD over the row's 5 pretrained models
\newcommand{\FtAucJcEFourOneseven}{0.9749\,\ensuremath{\pm}\,0.0006}  % experiments/FIGS/data/w2b_leg2_metrics_v2.json :: cells.{r16q1-s1,r16q1-s2,r16q1-s3,r16q1-s4,r16q1-s5}.N10000.s1.macro_auc_ovr :: sha256 9cd059d3e2f51798 :: mean +- SD over the row's 5 pretrained models
\newcommand{\FtAucJcEFiveOneseven}{0.9821\,\ensuremath{\pm}\,\ensuremath{<}0.0001}  % experiments/FIGS/data/w2b_leg2_metrics_v2.json :: cells.{r16q1-s1,r16q1-s2,r16q1-s3,r16q1-s4,r16q1-s5}.N100000.s1.macro_auc_ovr :: sha256 9cd059d3e2f51798 :: mean +- SD over the row's 5 pretrained models
\newcommand{\FtAucJcESixOneseven}{0.9846\,\ensuremath{\pm}\,0.0001}  % experiments/FIGS/data/w2b_leg2_metrics_v2.json :: cells.{r16q1-s1,r16q1-s2,r16q1-s3,r16q1-s4,r16q1-s5}.N1000000.s1.macro_auc_ovr :: sha256 9cd059d3e2f51798 :: mean +- SD over the row's 5 pretrained models
\newcommand{\FtAucJcEThreeOnesixtwoMass}{0.9569\,\ensuremath{\pm}\,0.0016}  % experiments/FIGS/data/w2b_leg2_metrics_v2.json :: cells.{l162mass-s1,l162mass-s2,l162mass-s3,l162mass-s4,l162mass-s5}.N1000.s1.macro_auc_ovr :: sha256 9cd059d3e2f51798 :: mean +- SD over the row's 5 pretrained models
\newcommand{\FtAucJcEFourOnesixtwoMass}{0.9768\,\ensuremath{\pm}\,0.0005}  % experiments/FIGS/data/w2b_leg2_metrics_v2.json :: cells.{l162mass-s1,l162mass-s2,l162mass-s3,l162mass-s4,l162mass-s5}.N10000.s1.macro_auc_ovr :: sha256 9cd059d3e2f51798 :: mean +- SD over the row's 5 pretrained models
\newcommand{\FtAucJcEFiveOnesixtwoMass}{0.9834\,\ensuremath{\pm}\,0.0001}  % experiments/FIGS/data/w2b_leg2_metrics_v2.json :: cells.{l162mass-s1,l162mass-s2,l162mass-s3,l162mass-s4,l162mass-s5}.N100000.s1.macro_auc_ovr :: sha256 9cd059d3e2f51798 :: mean +- SD over the row's 5 pretrained models
\newcommand{\FtAucJcESixOnesixtwoMass}{0.9856\,\ensuremath{\pm}\,0.0001}  % experiments/FIGS/data/w2b_leg2_metrics_v2.json :: cells.{l162mass-s1,l162mass-s2,l162mass-s3,l162mass-s4,l162mass-s5}.N1000000.s1.macro_auc_ovr :: sha256 9cd059d3e2f51798 :: mean +- SD over the row's 5 pretrained models
\newcommand{\FtAucJcEThreeOnesevenMass}{0.9460\,\ensuremath{\pm}\,0.0029}  % experiments/FIGS/data/w2b_leg2_metrics_v2.json :: cells.{r16q1mass-s1,r16q1mass-s2,r16q1mass-s3,r16q1mass-s4,r16q1mass-s5}.N1000.s1.macro_auc_ovr :: sha256 9cd059d3e2f51798 :: mean +- SD over the row's 5 pretrained models
\newcommand{\FtAucJcEFourOnesevenMass}{0.9737\,\ensuremath{\pm}\,0.0008}  % experiments/FIGS/data/w2b_leg2_metrics_v2.json :: cells.{r16q1mass-s1,r16q1mass-s2,r16q1mass-s3,r16q1mass-s4,r16q1mass-s5}.N10000.s1.macro_auc_ovr :: sha256 9cd059d3e2f51798 :: mean +- SD over the row's 5 pretrained models
\newcommand{\FtAucJcEFiveOnesevenMass}{0.9817\,\ensuremath{\pm}\,0.0002}  % experiments/FIGS/data/w2b_leg2_metrics_v2.json :: cells.{r16q1mass-s1,r16q1mass-s2,r16q1mass-s3,r16q1mass-s4,r16q1mass-s5}.N100000.s1.macro_auc_ovr :: sha256 9cd059d3e2f51798 :: mean +- SD over the row's 5 pretrained models
\newcommand{\FtAucJcESixOnesevenMass}{0.9842\,\ensuremath{\pm}\,0.0002}  % experiments/FIGS/data/w2b_leg2_metrics_v2.json :: cells.{r16q1mass-s1,r16q1mass-s2,r16q1mass-s3,r16q1mass-s4,r16q1mass-s5}.N1000000.s1.macro_auc_ovr :: sha256 9cd059d3e2f51798 :: mean +- SD over the row's 5 pretrained models
\newcommand{\FtAucJcEThreeRandDrawOne}{0.9477}  % experiments/FIGS/data/w2b_leg2_metrics_v2.json :: cells.{rand-d1-s1b}.N1000.s1.macro_auc_ovr :: sha256 9cd059d3e2f51798 :: one run
\newcommand{\FtAucJcEFourRandDrawOne}{0.9745}  % experiments/FIGS/data/w2b_leg2_metrics_v2.json :: cells.{rand-d1-s1b}.N10000.s1.macro_auc_ovr :: sha256 9cd059d3e2f51798 :: one run
\newcommand{\FtAucJcEFiveRandDrawOne}{0.98265}  % experiments/FIGS/data/w2b_leg2_metrics_v2.json :: cells.{rand-d1-s1b}.N100000.s1.macro_auc_ovr :: sha256 9cd059d3e2f51798 :: one run
\newcommand{\FtAucJcESixRandDrawOne}{0.98513}  % experiments/FIGS/data/w2b_leg2_metrics_v2.json :: cells.{rand-d1-s1b}.N1000000.s1.macro_auc_ovr :: sha256 9cd059d3e2f51798 :: one run
\newcommand{\FtAucJcEThreeRandDrawTwo}{0.9501}  % experiments/FIGS/data/w2b_leg2_metrics_v2.json :: cells.{rand-d2-s2}.N1000.s1.macro_auc_ovr :: sha256 9cd059d3e2f51798 :: one run
\newcommand{\FtAucJcEFourRandDrawTwo}{0.9756}  % experiments/FIGS/data/w2b_leg2_metrics_v2.json :: cells.{rand-d2-s2}.N10000.s1.macro_auc_ovr :: sha256 9cd059d3e2f51798 :: one run
\newcommand{\FtAucJcEFiveRandDrawTwo}{0.98278}  % experiments/FIGS/data/w2b_leg2_metrics_v2.json :: cells.{rand-d2-s2}.N100000.s1.macro_auc_ovr :: sha256 9cd059d3e2f51798 :: one run
\newcommand{\FtAucJcESixRandDrawTwo}{0.98512}  % experiments/FIGS/data/w2b_leg2_metrics_v2.json :: cells.{rand-d2-s2}.N1000000.s1.macro_auc_ovr :: sha256 9cd059d3e2f51798 :: one run
\newcommand{\FtAucJcEThreeRandDrawThree}{0.9470}  % experiments/FIGS/data/w2b_leg2_metrics_v2.json :: cells.{rand-d3-s3}.N1000.s1.macro_auc_ovr :: sha256 9cd059d3e2f51798 :: one run
\newcommand{\FtAucJcEFourRandDrawThree}{0.9743}  % experiments/FIGS/data/w2b_leg2_metrics_v2.json :: cells.{rand-d3-s3}.N10000.s1.macro_auc_ovr :: sha256 9cd059d3e2f51798 :: one run
\newcommand{\FtAucJcEFiveRandDrawThree}{0.98253}  % experiments/FIGS/data/w2b_leg2_metrics_v2.json :: cells.{rand-d3-s3}.N100000.s1.macro_auc_ovr :: sha256 9cd059d3e2f51798 :: one run
\newcommand{\FtAucJcESixRandDrawThree}{0.98506}  % experiments/FIGS/data/w2b_leg2_metrics_v2.json :: cells.{rand-d3-s3}.N1000000.s1.macro_auc_ovr :: sha256 9cd059d3e2f51798 :: one run
\newcommand{\FtAccJcEThreeOneeighteight}{0.7183\,\ensuremath{\pm}\,0.0032}  % experiments/FIGS/data/w2b_leg2_metrics_v2.json :: cells.{l188-s1,l188-s2,l188-s3,l188-s4,l188-s5}.N1000.s1.accuracy :: sha256 9cd059d3e2f51798 :: mean +- SD over the row's 5 pretrained models
\newcommand{\FtAccJcEFourOneeighteight}{0.8039\,\ensuremath{\pm}\,0.0030}  % experiments/FIGS/data/w2b_leg2_metrics_v2.json :: cells.{l188-s1,l188-s2,l188-s3,l188-s4,l188-s5}.N10000.s1.accuracy :: sha256 9cd059d3e2f51798 :: mean +- SD over the row's 5 pretrained models
\newcommand{\FtAccJcEFiveOneeighteight}{0.8351\,\ensuremath{\pm}\,0.0008}  % experiments/FIGS/data/w2b_leg2_metrics_v2.json :: cells.{l188-s1,l188-s2,l188-s3,l188-s4,l188-s5}.N100000.s1.accuracy :: sha256 9cd059d3e2f51798 :: mean +- SD over the row's 5 pretrained models
\newcommand{\FtAccJcESixOneeighteight}{0.8486\,\ensuremath{\pm}\,0.0005}  % experiments/FIGS/data/w2b_leg2_metrics_v2.json :: cells.{l188-s1,l188-s2,l188-s3,l188-s4,l188-s5}.N1000000.s1.accuracy :: sha256 9cd059d3e2f51798 :: mean +- SD over the row's 5 pretrained models
\newcommand{\FtAccJcEThreeOnesixtwo}{0.7268\,\ensuremath{\pm}\,0.0017}  % experiments/FIGS/data/w2b_leg2_metrics_v2.json :: cells.{l162-s1b,l162-s2,l162-s3,l162-s4,l162-s5}.N1000.s1.accuracy :: sha256 9cd059d3e2f51798 :: mean +- SD over the row's 5 pretrained models
\newcommand{\FtAccJcEFourOnesixtwo}{0.8069\,\ensuremath{\pm}\,0.0025}  % experiments/FIGS/data/w2b_leg2_metrics_v2.json :: cells.{l162-s1b,l162-s2,l162-s3,l162-s4,l162-s5}.N10000.s1.accuracy :: sha256 9cd059d3e2f51798 :: mean +- SD over the row's 5 pretrained models
\newcommand{\FtAccJcEFiveOnesixtwo}{0.8363\,\ensuremath{\pm}\,0.0008}  % experiments/FIGS/data/w2b_leg2_metrics_v2.json :: cells.{l162-s1b,l162-s2,l162-s3,l162-s4,l162-s5}.N100000.s1.accuracy :: sha256 9cd059d3e2f51798 :: mean +- SD over the row's 5 pretrained models
\newcommand{\FtAccJcESixOnesixtwo}{0.8485\,\ensuremath{\pm}\,0.0005}  % experiments/FIGS/data/w2b_leg2_metrics_v2.json :: cells.{l162-s1b,l162-s2,l162-s3,l162-s4,l162-s5}.N1000000.s1.accuracy :: sha256 9cd059d3e2f51798 :: mean +- SD over the row's 5 pretrained models
\newcommand{\FtAccJcEThreeFourthree}{0.722\,\ensuremath{\pm}\,0.007}  % experiments/FIGS/data/w2b_leg2_metrics_v2.json :: cells.{r42q1-s1,r42q1-s2,r42q1-s3,r42q1-s4,r42q1-s5}.N1000.s1.accuracy :: sha256 9cd059d3e2f51798 :: mean +- SD over the row's 5 pretrained models
\newcommand{\FtAccJcEFourFourthree}{0.8050\,\ensuremath{\pm}\,0.0026}  % experiments/FIGS/data/w2b_leg2_metrics_v2.json :: cells.{r42q1-s1,r42q1-s2,r42q1-s3,r42q1-s4,r42q1-s5}.N10000.s1.accuracy :: sha256 9cd059d3e2f51798 :: mean +- SD over the row's 5 pretrained models
\newcommand{\FtAccJcEFiveFourthree}{0.8336\,\ensuremath{\pm}\,0.0008}  % experiments/FIGS/data/w2b_leg2_metrics_v2.json :: cells.{r42q1-s1,r42q1-s2,r42q1-s3,r42q1-s4,r42q1-s5}.N100000.s1.accuracy :: sha256 9cd059d3e2f51798 :: mean +- SD over the row's 5 pretrained models
\newcommand{\FtAccJcESixFourthree}{0.8468\,\ensuremath{\pm}\,0.0003}  % experiments/FIGS/data/w2b_leg2_metrics_v2.json :: cells.{r42q1-s1,r42q1-s2,r42q1-s3,r42q1-s4,r42q1-s5}.N1000000.s1.accuracy :: sha256 9cd059d3e2f51798 :: mean +- SD over the row's 5 pretrained models
\newcommand{\FtAccJcEThreeOneseven}{0.714\,\ensuremath{\pm}\,0.005}  % experiments/FIGS/data/w2b_leg2_metrics_v2.json :: cells.{r16q1-s1,r16q1-s2,r16q1-s3,r16q1-s4,r16q1-s5}.N1000.s1.accuracy :: sha256 9cd059d3e2f51798 :: mean +- SD over the row's 5 pretrained models
\newcommand{\FtAccJcEFourOneseven}{0.7964\,\ensuremath{\pm}\,0.0019}  % experiments/FIGS/data/w2b_leg2_metrics_v2.json :: cells.{r16q1-s1,r16q1-s2,r16q1-s3,r16q1-s4,r16q1-s5}.N10000.s1.accuracy :: sha256 9cd059d3e2f51798 :: mean +- SD over the row's 5 pretrained models
\newcommand{\FtAccJcEFiveOneseven}{0.8286\,\ensuremath{\pm}\,0.0006}  % experiments/FIGS/data/w2b_leg2_metrics_v2.json :: cells.{r16q1-s1,r16q1-s2,r16q1-s3,r16q1-s4,r16q1-s5}.N100000.s1.accuracy :: sha256 9cd059d3e2f51798 :: mean +- SD over the row's 5 pretrained models
\newcommand{\FtAccJcESixOneseven}{0.8415\,\ensuremath{\pm}\,0.0007}  % experiments/FIGS/data/w2b_leg2_metrics_v2.json :: cells.{r16q1-s1,r16q1-s2,r16q1-s3,r16q1-s4,r16q1-s5}.N1000000.s1.accuracy :: sha256 9cd059d3e2f51798 :: mean +- SD over the row's 5 pretrained models
\newcommand{\FtAccJcEThreeOnesixtwoMass}{0.728\,\ensuremath{\pm}\,0.004}  % experiments/FIGS/data/w2b_leg2_metrics_v2.json :: cells.{l162mass-s1,l162mass-s2,l162mass-s3,l162mass-s4,l162mass-s5}.N1000.s1.accuracy :: sha256 9cd059d3e2f51798 :: mean +- SD over the row's 5 pretrained models
\newcommand{\FtAccJcEFourOnesixtwoMass}{0.8043\,\ensuremath{\pm}\,0.0020}  % experiments/FIGS/data/w2b_leg2_metrics_v2.json :: cells.{l162mass-s1,l162mass-s2,l162mass-s3,l162mass-s4,l162mass-s5}.N10000.s1.accuracy :: sha256 9cd059d3e2f51798 :: mean +- SD over the row's 5 pretrained models
\newcommand{\FtAccJcEFiveOnesixtwoMass}{0.8349\,\ensuremath{\pm}\,0.0005}  % experiments/FIGS/data/w2b_leg2_metrics_v2.json :: cells.{l162mass-s1,l162mass-s2,l162mass-s3,l162mass-s4,l162mass-s5}.N100000.s1.accuracy :: sha256 9cd059d3e2f51798 :: mean +- SD over the row's 5 pretrained models
\newcommand{\FtAccJcESixOnesixtwoMass}{0.8475\,\ensuremath{\pm}\,0.0009}  % experiments/FIGS/data/w2b_leg2_metrics_v2.json :: cells.{l162mass-s1,l162mass-s2,l162mass-s3,l162mass-s4,l162mass-s5}.N1000000.s1.accuracy :: sha256 9cd059d3e2f51798 :: mean +- SD over the row's 5 pretrained models
\newcommand{\FtAccJcEThreeOnesevenMass}{0.702\,\ensuremath{\pm}\,0.005}  % experiments/FIGS/data/w2b_leg2_metrics_v2.json :: cells.{r16q1mass-s1,r16q1mass-s2,r16q1mass-s3,r16q1mass-s4,r16q1mass-s5}.N1000.s1.accuracy :: sha256 9cd059d3e2f51798 :: mean +- SD over the row's 5 pretrained models
\newcommand{\FtAccJcEFourOnesevenMass}{0.7909\,\ensuremath{\pm}\,0.0031}  % experiments/FIGS/data/w2b_leg2_metrics_v2.json :: cells.{r16q1mass-s1,r16q1mass-s2,r16q1mass-s3,r16q1mass-s4,r16q1mass-s5}.N10000.s1.accuracy :: sha256 9cd059d3e2f51798 :: mean +- SD over the row's 5 pretrained models
\newcommand{\FtAccJcEFiveOnesevenMass}{0.8265\,\ensuremath{\pm}\,0.0007}  % experiments/FIGS/data/w2b_leg2_metrics_v2.json :: cells.{r16q1mass-s1,r16q1mass-s2,r16q1mass-s3,r16q1mass-s4,r16q1mass-s5}.N100000.s1.accuracy :: sha256 9cd059d3e2f51798 :: mean +- SD over the row's 5 pretrained models
\newcommand{\FtAccJcESixOnesevenMass}{0.8390\,\ensuremath{\pm}\,0.0011}  % experiments/FIGS/data/w2b_leg2_metrics_v2.json :: cells.{r16q1mass-s1,r16q1mass-s2,r16q1mass-s3,r16q1mass-s4,r16q1mass-s5}.N1000000.s1.accuracy :: sha256 9cd059d3e2f51798 :: mean +- SD over the row's 5 pretrained models
\newcommand{\FtAccJcEThreeRandDrawOne}{0.7049}  % experiments/FIGS/data/w2b_leg2_metrics_v2.json :: cells.{rand-d1-s1b}.N1000.s1.accuracy :: sha256 9cd059d3e2f51798 :: one run
\newcommand{\FtAccJcEFourRandDrawOne}{0.7974}  % experiments/FIGS/data/w2b_leg2_metrics_v2.json :: cells.{rand-d1-s1b}.N10000.s1.accuracy :: sha256 9cd059d3e2f51798 :: one run
\newcommand{\FtAccJcEFiveRandDrawOne}{0.8326}  % experiments/FIGS/data/w2b_leg2_metrics_v2.json :: cells.{rand-d1-s1b}.N100000.s1.accuracy :: sha256 9cd059d3e2f51798 :: one run
\newcommand{\FtAccJcESixRandDrawOne}{0.84467}  % experiments/FIGS/data/w2b_leg2_metrics_v2.json :: cells.{rand-d1-s1b}.N1000000.s1.accuracy :: sha256 9cd059d3e2f51798 :: one run
\newcommand{\FtAccJcEThreeRandDrawTwo}{0.7135}  % experiments/FIGS/data/w2b_leg2_metrics_v2.json :: cells.{rand-d2-s2}.N1000.s1.accuracy :: sha256 9cd059d3e2f51798 :: one run
\newcommand{\FtAccJcEFourRandDrawTwo}{0.7998}  % experiments/FIGS/data/w2b_leg2_metrics_v2.json :: cells.{rand-d2-s2}.N10000.s1.accuracy :: sha256 9cd059d3e2f51798 :: one run
\newcommand{\FtAccJcEFiveRandDrawTwo}{0.8332}  % experiments/FIGS/data/w2b_leg2_metrics_v2.json :: cells.{rand-d2-s2}.N100000.s1.accuracy :: sha256 9cd059d3e2f51798 :: one run
\newcommand{\FtAccJcESixRandDrawTwo}{0.84449}  % experiments/FIGS/data/w2b_leg2_metrics_v2.json :: cells.{rand-d2-s2}.N1000000.s1.accuracy :: sha256 9cd059d3e2f51798 :: one run
\newcommand{\FtAccJcEThreeRandDrawThree}{0.7090}  % experiments/FIGS/data/w2b_leg2_metrics_v2.json :: cells.{rand-d3-s3}.N1000.s1.accuracy :: sha256 9cd059d3e2f51798 :: one run
\newcommand{\FtAccJcEFourRandDrawThree}{0.7948}  % experiments/FIGS/data/w2b_leg2_metrics_v2.json :: cells.{rand-d3-s3}.N10000.s1.accuracy :: sha256 9cd059d3e2f51798 :: one run
\newcommand{\FtAccJcEFiveRandDrawThree}{0.8310}  % experiments/FIGS/data/w2b_leg2_metrics_v2.json :: cells.{rand-d3-s3}.N100000.s1.accuracy :: sha256 9cd059d3e2f51798 :: one run
\newcommand{\FtAccJcESixRandDrawThree}{0.84409}  % experiments/FIGS/data/w2b_leg2_metrics_v2.json :: cells.{rand-d3-s3}.N1000000.s1.accuracy :: sha256 9cd059d3e2f51798 :: one run
\newcommand{\FtOmaRatioJcEThreeOnesixtwoOverOneeighteight}{0.96}  % experiments/FIGS/data/w2b_leg2_metrics_v2.json :: cells.*.N1000.s1.macro_auc_ovr :: sha256 9cd059d3e2f51798 :: ratio of the seed means of 1 - macro AUC, this row over the reference
\newcommand{\FtOmaRatioJcEThreeFourthreeOverOneeighteight}{1.01}  % experiments/FIGS/data/w2b_leg2_metrics_v2.json :: cells.*.N1000.s1.macro_auc_ovr :: sha256 9cd059d3e2f51798 :: ratio of the seed means of 1 - macro AUC, this row over the reference
\newcommand{\FtOmaRatioJcEThreeOnesevenOverOneeighteight}{1.07}  % experiments/FIGS/data/w2b_leg2_metrics_v2.json :: cells.*.N1000.s1.macro_auc_ovr :: sha256 9cd059d3e2f51798 :: ratio of the seed means of 1 - macro AUC, this row over the reference
\newcommand{\FtOmaRatioJcEThreeOnesixtwoMassOverOneeighteight}{0.87}  % experiments/FIGS/data/w2b_leg2_metrics_v2.json :: cells.*.N1000.s1.macro_auc_ovr :: sha256 9cd059d3e2f51798 :: ratio of the seed means of 1 - macro AUC, this row over the reference
\newcommand{\FtOmaRatioJcEThreeOnesevenMassOverOneeighteight}{1.09}  % experiments/FIGS/data/w2b_leg2_metrics_v2.json :: cells.*.N1000.s1.macro_auc_ovr :: sha256 9cd059d3e2f51798 :: ratio of the seed means of 1 - macro AUC, this row over the reference
\newcommand{\FtOmaRatioJcEThreeRandDrawOneOverOneeighteight}{1.06}  % experiments/FIGS/data/w2b_leg2_metrics_v2.json :: cells.*.N1000.s1.macro_auc_ovr :: sha256 9cd059d3e2f51798 :: ratio of the seed means of 1 - macro AUC, this row over the reference
\newcommand{\FtOmaRatioJcEThreeRandDrawTwoOverOneeighteight}{1.01}  % experiments/FIGS/data/w2b_leg2_metrics_v2.json :: cells.*.N1000.s1.macro_auc_ovr :: sha256 9cd059d3e2f51798 :: ratio of the seed means of 1 - macro AUC, this row over the reference
\newcommand{\FtOmaRatioJcEThreeRandDrawThreeOverOneeighteight}{1.07}  % experiments/FIGS/data/w2b_leg2_metrics_v2.json :: cells.*.N1000.s1.macro_auc_ovr :: sha256 9cd059d3e2f51798 :: ratio of the seed means of 1 - macro AUC, this row over the reference
\newcommand{\FtOmaRatioJcEThreeOneeighteightOverOnesixtwo}{1.05}  % experiments/FIGS/data/w2b_leg2_metrics_v2.json :: cells.*.N1000.s1.macro_auc_ovr :: sha256 9cd059d3e2f51798 :: ratio of the seed means of 1 - macro AUC, this row over the reference
\newcommand{\FtOmaRatioJcEThreeFourthreeOverOnesixtwo}{1.05}  % experiments/FIGS/data/w2b_leg2_metrics_v2.json :: cells.*.N1000.s1.macro_auc_ovr :: sha256 9cd059d3e2f51798 :: ratio of the seed means of 1 - macro AUC, this row over the reference
\newcommand{\FtOmaRatioJcEThreeOnesevenOverOnesixtwo}{1.12}  % experiments/FIGS/data/w2b_leg2_metrics_v2.json :: cells.*.N1000.s1.macro_auc_ovr :: sha256 9cd059d3e2f51798 :: ratio of the seed means of 1 - macro AUC, this row over the reference
\newcommand{\FtOmaRatioJcEThreeOnesixtwoMassOverOnesixtwo}{0.91}  % experiments/FIGS/data/w2b_leg2_metrics_v2.json :: cells.*.N1000.s1.macro_auc_ovr :: sha256 9cd059d3e2f51798 :: ratio of the seed means of 1 - macro AUC, this row over the reference
\newcommand{\FtOmaRatioJcEThreeOnesevenMassOverOnesixtwo}{1.14}  % experiments/FIGS/data/w2b_leg2_metrics_v2.json :: cells.*.N1000.s1.macro_auc_ovr :: sha256 9cd059d3e2f51798 :: ratio of the seed means of 1 - macro AUC, this row over the reference
\newcommand{\FtOmaRatioJcEThreeRandDrawOneOverOnesixtwo}{1.10}  % experiments/FIGS/data/w2b_leg2_metrics_v2.json :: cells.*.N1000.s1.macro_auc_ovr :: sha256 9cd059d3e2f51798 :: ratio of the seed means of 1 - macro AUC, this row over the reference
\newcommand{\FtOmaRatioJcEThreeRandDrawTwoOverOnesixtwo}{1.05}  % experiments/FIGS/data/w2b_leg2_metrics_v2.json :: cells.*.N1000.s1.macro_auc_ovr :: sha256 9cd059d3e2f51798 :: ratio of the seed means of 1 - macro AUC, this row over the reference
\newcommand{\FtOmaRatioJcEThreeRandDrawThreeOverOnesixtwo}{1.12}  % experiments/FIGS/data/w2b_leg2_metrics_v2.json :: cells.*.N1000.s1.macro_auc_ovr :: sha256 9cd059d3e2f51798 :: ratio of the seed means of 1 - macro AUC, this row over the reference
\newcommand{\FtOmaRatioJcEFourOnesixtwoOverOneeighteight}{0.98}  % experiments/FIGS/data/w2b_leg2_metrics_v2.json :: cells.*.N10000.s1.macro_auc_ovr :: sha256 9cd059d3e2f51798 :: ratio of the seed means of 1 - macro AUC, this row over the reference
\newcommand{\FtOmaRatioJcEFourFourthreeOverOneeighteight}{0.98}  % experiments/FIGS/data/w2b_leg2_metrics_v2.json :: cells.*.N10000.s1.macro_auc_ovr :: sha256 9cd059d3e2f51798 :: ratio of the seed means of 1 - macro AUC, this row over the reference
\newcommand{\FtOmaRatioJcEFourOnesevenOverOneeighteight}{1.07}  % experiments/FIGS/data/w2b_leg2_metrics_v2.json :: cells.*.N10000.s1.macro_auc_ovr :: sha256 9cd059d3e2f51798 :: ratio of the seed means of 1 - macro AUC, this row over the reference
\newcommand{\FtOmaRatioJcEFourOnesixtwoMassOverOneeighteight}{0.99}  % experiments/FIGS/data/w2b_leg2_metrics_v2.json :: cells.*.N10000.s1.macro_auc_ovr :: sha256 9cd059d3e2f51798 :: ratio of the seed means of 1 - macro AUC, this row over the reference
\newcommand{\FtOmaRatioJcEFourOnesevenMassOverOneeighteight}{1.12}  % experiments/FIGS/data/w2b_leg2_metrics_v2.json :: cells.*.N10000.s1.macro_auc_ovr :: sha256 9cd059d3e2f51798 :: ratio of the seed means of 1 - macro AUC, this row over the reference
\newcommand{\FtOmaRatioJcEFourRandDrawOneOverOneeighteight}{1.08}  % experiments/FIGS/data/w2b_leg2_metrics_v2.json :: cells.*.N10000.s1.macro_auc_ovr :: sha256 9cd059d3e2f51798 :: ratio of the seed means of 1 - macro AUC, this row over the reference
\newcommand{\FtOmaRatioJcEFourRandDrawTwoOverOneeighteight}{1.04}  % experiments/FIGS/data/w2b_leg2_metrics_v2.json :: cells.*.N10000.s1.macro_auc_ovr :: sha256 9cd059d3e2f51798 :: ratio of the seed means of 1 - macro AUC, this row over the reference
\newcommand{\FtOmaRatioJcEFourRandDrawThreeOverOneeighteight}{1.09}  % experiments/FIGS/data/w2b_leg2_metrics_v2.json :: cells.*.N10000.s1.macro_auc_ovr :: sha256 9cd059d3e2f51798 :: ratio of the seed means of 1 - macro AUC, this row over the reference
\newcommand{\FtOmaRatioJcEFourOneeighteightOverOnesixtwo}{1.02}  % experiments/FIGS/data/w2b_leg2_metrics_v2.json :: cells.*.N10000.s1.macro_auc_ovr :: sha256 9cd059d3e2f51798 :: ratio of the seed means of 1 - macro AUC, this row over the reference
\newcommand{\FtOmaRatioJcEFourFourthreeOverOnesixtwo}{1.01}  % experiments/FIGS/data/w2b_leg2_metrics_v2.json :: cells.*.N10000.s1.macro_auc_ovr :: sha256 9cd059d3e2f51798 :: ratio of the seed means of 1 - macro AUC, this row over the reference
\newcommand{\FtOmaRatioJcEFourOnesevenOverOnesixtwo}{1.09}  % experiments/FIGS/data/w2b_leg2_metrics_v2.json :: cells.*.N10000.s1.macro_auc_ovr :: sha256 9cd059d3e2f51798 :: ratio of the seed means of 1 - macro AUC, this row over the reference
\newcommand{\FtOmaRatioJcEFourOnesixtwoMassOverOnesixtwo}{1.01}  % experiments/FIGS/data/w2b_leg2_metrics_v2.json :: cells.*.N10000.s1.macro_auc_ovr :: sha256 9cd059d3e2f51798 :: ratio of the seed means of 1 - macro AUC, this row over the reference
\newcommand{\FtOmaRatioJcEFourOnesevenMassOverOnesixtwo}{1.14}  % experiments/FIGS/data/w2b_leg2_metrics_v2.json :: cells.*.N10000.s1.macro_auc_ovr :: sha256 9cd059d3e2f51798 :: ratio of the seed means of 1 - macro AUC, this row over the reference
\newcommand{\FtOmaRatioJcEFourRandDrawOneOverOnesixtwo}{1.11}  % experiments/FIGS/data/w2b_leg2_metrics_v2.json :: cells.*.N10000.s1.macro_auc_ovr :: sha256 9cd059d3e2f51798 :: ratio of the seed means of 1 - macro AUC, this row over the reference
\newcommand{\FtOmaRatioJcEFourRandDrawTwoOverOnesixtwo}{1.06}  % experiments/FIGS/data/w2b_leg2_metrics_v2.json :: cells.*.N10000.s1.macro_auc_ovr :: sha256 9cd059d3e2f51798 :: ratio of the seed means of 1 - macro AUC, this row over the reference
\newcommand{\FtOmaRatioJcEFourRandDrawThreeOverOnesixtwo}{1.12}  % experiments/FIGS/data/w2b_leg2_metrics_v2.json :: cells.*.N10000.s1.macro_auc_ovr :: sha256 9cd059d3e2f51798 :: ratio of the seed means of 1 - macro AUC, this row over the reference
\newcommand{\FtOmaRatioJcEFiveOnesixtwoOverOneeighteight}{0.98}  % experiments/FIGS/data/w2b_leg2_metrics_v2.json :: cells.*.N100000.s1.macro_auc_ovr :: sha256 9cd059d3e2f51798 :: ratio of the seed means of 1 - macro AUC, this row over the reference
\newcommand{\FtOmaRatioJcEFiveFourthreeOverOneeighteight}{1.02}  % experiments/FIGS/data/w2b_leg2_metrics_v2.json :: cells.*.N100000.s1.macro_auc_ovr :: sha256 9cd059d3e2f51798 :: ratio of the seed means of 1 - macro AUC, this row over the reference
\newcommand{\FtOmaRatioJcEFiveOnesevenOverOneeighteight}{1.07}  % experiments/FIGS/data/w2b_leg2_metrics_v2.json :: cells.*.N100000.s1.macro_auc_ovr :: sha256 9cd059d3e2f51798 :: ratio of the seed means of 1 - macro AUC, this row over the reference
\newcommand{\FtOmaRatioJcEFiveOnesixtwoMassOverOneeighteight}{1.00}  % experiments/FIGS/data/w2b_leg2_metrics_v2.json :: cells.*.N100000.s1.macro_auc_ovr :: sha256 9cd059d3e2f51798 :: ratio of the seed means of 1 - macro AUC, this row over the reference
\newcommand{\FtOmaRatioJcEFiveOnesevenMassOverOneeighteight}{1.10}  % experiments/FIGS/data/w2b_leg2_metrics_v2.json :: cells.*.N100000.s1.macro_auc_ovr :: sha256 9cd059d3e2f51798 :: ratio of the seed means of 1 - macro AUC, this row over the reference
\newcommand{\FtOmaRatioJcEFiveRandDrawOneOverOneeighteight}{1.04}  % experiments/FIGS/data/w2b_leg2_metrics_v2.json :: cells.*.N100000.s1.macro_auc_ovr :: sha256 9cd059d3e2f51798 :: ratio of the seed means of 1 - macro AUC, this row over the reference
\newcommand{\FtOmaRatioJcEFiveRandDrawTwoOverOneeighteight}{1.03}  % experiments/FIGS/data/w2b_leg2_metrics_v2.json :: cells.*.N100000.s1.macro_auc_ovr :: sha256 9cd059d3e2f51798 :: ratio of the seed means of 1 - macro AUC, this row over the reference
\newcommand{\FtOmaRatioJcEFiveRandDrawThreeOverOneeighteight}{1.05}  % experiments/FIGS/data/w2b_leg2_metrics_v2.json :: cells.*.N100000.s1.macro_auc_ovr :: sha256 9cd059d3e2f51798 :: ratio of the seed means of 1 - macro AUC, this row over the reference
\newcommand{\FtOmaRatioJcEFiveOneeighteightOverOnesixtwo}{1.02}  % experiments/FIGS/data/w2b_leg2_metrics_v2.json :: cells.*.N100000.s1.macro_auc_ovr :: sha256 9cd059d3e2f51798 :: ratio of the seed means of 1 - macro AUC, this row over the reference
\newcommand{\FtOmaRatioJcEFiveFourthreeOverOnesixtwo}{1.03}  % experiments/FIGS/data/w2b_leg2_metrics_v2.json :: cells.*.N100000.s1.macro_auc_ovr :: sha256 9cd059d3e2f51798 :: ratio of the seed means of 1 - macro AUC, this row over the reference
\newcommand{\FtOmaRatioJcEFiveOnesevenOverOnesixtwo}{1.09}  % experiments/FIGS/data/w2b_leg2_metrics_v2.json :: cells.*.N100000.s1.macro_auc_ovr :: sha256 9cd059d3e2f51798 :: ratio of the seed means of 1 - macro AUC, this row over the reference
\newcommand{\FtOmaRatioJcEFiveOnesixtwoMassOverOnesixtwo}{1.02}  % experiments/FIGS/data/w2b_leg2_metrics_v2.json :: cells.*.N100000.s1.macro_auc_ovr :: sha256 9cd059d3e2f51798 :: ratio of the seed means of 1 - macro AUC, this row over the reference
\newcommand{\FtOmaRatioJcEFiveOnesevenMassOverOnesixtwo}{1.11}  % experiments/FIGS/data/w2b_leg2_metrics_v2.json :: cells.*.N100000.s1.macro_auc_ovr :: sha256 9cd059d3e2f51798 :: ratio of the seed means of 1 - macro AUC, this row over the reference
\newcommand{\FtOmaRatioJcEFiveRandDrawOneOverOnesixtwo}{1.06}  % experiments/FIGS/data/w2b_leg2_metrics_v2.json :: cells.*.N100000.s1.macro_auc_ovr :: sha256 9cd059d3e2f51798 :: ratio of the seed means of 1 - macro AUC, this row over the reference
\newcommand{\FtOmaRatioJcEFiveRandDrawTwoOverOnesixtwo}{1.05}  % experiments/FIGS/data/w2b_leg2_metrics_v2.json :: cells.*.N100000.s1.macro_auc_ovr :: sha256 9cd059d3e2f51798 :: ratio of the seed means of 1 - macro AUC, this row over the reference
\newcommand{\FtOmaRatioJcEFiveRandDrawThreeOverOnesixtwo}{1.07}  % experiments/FIGS/data/w2b_leg2_metrics_v2.json :: cells.*.N100000.s1.macro_auc_ovr :: sha256 9cd059d3e2f51798 :: ratio of the seed means of 1 - macro AUC, this row over the reference
\newcommand{\FtOmaRatioJcESixOnesixtwoOverOneeighteight}{1.00}  % experiments/FIGS/data/w2b_leg2_metrics_v2.json :: cells.*.N1000000.s1.macro_auc_ovr :: sha256 9cd059d3e2f51798 :: ratio of the seed means of 1 - macro AUC, this row over the reference
\newcommand{\FtOmaRatioJcESixFourthreeOverOneeighteight}{1.02}  % experiments/FIGS/data/w2b_leg2_metrics_v2.json :: cells.*.N1000000.s1.macro_auc_ovr :: sha256 9cd059d3e2f51798 :: ratio of the seed means of 1 - macro AUC, this row over the reference
\newcommand{\FtOmaRatioJcESixOnesevenOverOneeighteight}{1.08}  % experiments/FIGS/data/w2b_leg2_metrics_v2.json :: cells.*.N1000000.s1.macro_auc_ovr :: sha256 9cd059d3e2f51798 :: ratio of the seed means of 1 - macro AUC, this row over the reference
\newcommand{\FtOmaRatioJcESixOnesixtwoMassOverOneeighteight}{1.01}  % experiments/FIGS/data/w2b_leg2_metrics_v2.json :: cells.*.N1000000.s1.macro_auc_ovr :: sha256 9cd059d3e2f51798 :: ratio of the seed means of 1 - macro AUC, this row over the reference
\newcommand{\FtOmaRatioJcESixOnesevenMassOverOneeighteight}{1.11}  % experiments/FIGS/data/w2b_leg2_metrics_v2.json :: cells.*.N1000000.s1.macro_auc_ovr :: sha256 9cd059d3e2f51798 :: ratio of the seed means of 1 - macro AUC, this row over the reference
\newcommand{\FtOmaRatioJcESixRandDrawOneOverOneeighteight}{1.05}  % experiments/FIGS/data/w2b_leg2_metrics_v2.json :: cells.*.N1000000.s1.macro_auc_ovr :: sha256 9cd059d3e2f51798 :: ratio of the seed means of 1 - macro AUC, this row over the reference
\newcommand{\FtOmaRatioJcESixRandDrawTwoOverOneeighteight}{1.05}  % experiments/FIGS/data/w2b_leg2_metrics_v2.json :: cells.*.N1000000.s1.macro_auc_ovr :: sha256 9cd059d3e2f51798 :: ratio of the seed means of 1 - macro AUC, this row over the reference
\newcommand{\FtOmaRatioJcESixRandDrawThreeOverOneeighteight}{1.05}  % experiments/FIGS/data/w2b_leg2_metrics_v2.json :: cells.*.N1000000.s1.macro_auc_ovr :: sha256 9cd059d3e2f51798 :: ratio of the seed means of 1 - macro AUC, this row over the reference
\newcommand{\FtOmaRatioJcESixOneeighteightOverOnesixtwo}{1.00}  % experiments/FIGS/data/w2b_leg2_metrics_v2.json :: cells.*.N1000000.s1.macro_auc_ovr :: sha256 9cd059d3e2f51798 :: ratio of the seed means of 1 - macro AUC, this row over the reference
\newcommand{\FtOmaRatioJcESixFourthreeOverOnesixtwo}{1.02}  % experiments/FIGS/data/w2b_leg2_metrics_v2.json :: cells.*.N1000000.s1.macro_auc_ovr :: sha256 9cd059d3e2f51798 :: ratio of the seed means of 1 - macro AUC, this row over the reference
\newcommand{\FtOmaRatioJcESixOnesevenOverOnesixtwo}{1.08}  % experiments/FIGS/data/w2b_leg2_metrics_v2.json :: cells.*.N1000000.s1.macro_auc_ovr :: sha256 9cd059d3e2f51798 :: ratio of the seed means of 1 - macro AUC, this row over the reference
\newcommand{\FtOmaRatioJcESixOnesixtwoMassOverOnesixtwo}{1.01}  % experiments/FIGS/data/w2b_leg2_metrics_v2.json :: cells.*.N1000000.s1.macro_auc_ovr :: sha256 9cd059d3e2f51798 :: ratio of the seed means of 1 - macro AUC, this row over the reference
\newcommand{\FtOmaRatioJcESixOnesevenMassOverOnesixtwo}{1.11}  % experiments/FIGS/data/w2b_leg2_metrics_v2.json :: cells.*.N1000000.s1.macro_auc_ovr :: sha256 9cd059d3e2f51798 :: ratio of the seed means of 1 - macro AUC, this row over the reference
\newcommand{\FtOmaRatioJcESixRandDrawOneOverOnesixtwo}{1.05}  % experiments/FIGS/data/w2b_leg2_metrics_v2.json :: cells.*.N1000000.s1.macro_auc_ovr :: sha256 9cd059d3e2f51798 :: ratio of the seed means of 1 - macro AUC, this row over the reference
\newcommand{\FtOmaRatioJcESixRandDrawTwoOverOnesixtwo}{1.05}  % experiments/FIGS/data/w2b_leg2_metrics_v2.json :: cells.*.N1000000.s1.macro_auc_ovr :: sha256 9cd059d3e2f51798 :: ratio of the seed means of 1 - macro AUC, this row over the reference
\newcommand{\FtOmaRatioJcESixRandDrawThreeOverOnesixtwo}{1.05}  % experiments/FIGS/data/w2b_leg2_metrics_v2.json :: cells.*.N1000000.s1.macro_auc_ovr :: sha256 9cd059d3e2f51798 :: ratio of the seed means of 1 - macro AUC, this row over the reference
\newcommand{\PairedFtAccJcEThreeOnesevenMinusOneeighteight}{\ensuremath{-}0.0040~[\ensuremath{-}0.0144, +0.0064]}  % experiments/FIGS/data/w2b_leg2_metrics_v2.json :: cells.{r16q1-s1,r16q1-s2,r16q1-s3,r16q1-s4,r16q1-s5,l188-s1,l188-s2,l188-s3,l188-s4,l188-s5}.N1000.s1.accuracy :: sha256 9cd059d3e2f51798 :: accuracy, coarse minus fine, paired by run [95% Student-t interval]
\newcommand{\PairedFtAccJcEFourOnesevenMinusOneeighteight}{\ensuremath{-}0.0075~[\ensuremath{-}0.0118, \ensuremath{-}0.0033]}  % experiments/FIGS/data/w2b_leg2_metrics_v2.json :: cells.{r16q1-s1,r16q1-s2,r16q1-s3,r16q1-s4,r16q1-s5,l188-s1,l188-s2,l188-s3,l188-s4,l188-s5}.N10000.s1.accuracy :: sha256 9cd059d3e2f51798 :: accuracy, coarse minus fine, paired by run [95% Student-t interval]
\newcommand{\PairedFtAccJcEFiveOnesevenMinusOneeighteight}{\ensuremath{-}0.0066~[\ensuremath{-}0.0083, \ensuremath{-}0.0049]}  % experiments/FIGS/data/w2b_leg2_metrics_v2.json :: cells.{r16q1-s1,r16q1-s2,r16q1-s3,r16q1-s4,r16q1-s5,l188-s1,l188-s2,l188-s3,l188-s4,l188-s5}.N100000.s1.accuracy :: sha256 9cd059d3e2f51798 :: accuracy, coarse minus fine, paired by run [95% Student-t interval]
\newcommand{\PairedFtAccJcESixOnesevenMinusOneeighteight}{\ensuremath{-}0.0071~[\ensuremath{-}0.0080, \ensuremath{-}0.0062]}  % experiments/FIGS/data/w2b_leg2_metrics_v2.json :: cells.{r16q1-s1,r16q1-s2,r16q1-s3,r16q1-s4,r16q1-s5,l188-s1,l188-s2,l188-s3,l188-s4,l188-s5}.N1000000.s1.accuracy :: sha256 9cd059d3e2f51798 :: accuracy, coarse minus fine, paired by run [95% Student-t interval]
\newcommand{\FtNClassesJc}{10}  % experiments/FIGS/data/w2b_leg2_metrics_v2.json :: cells.*.*.s1.n_classes_present :: sha256 9cd059d3e2f51798 :: classes of the fine-tuning task
\newcommand{\FtNTestAucJc}{500{,}000}  % experiments/FIGS/data/w2b_leg2_metrics_v2.json :: cells.*.*.s1.n_jets_auc :: sha256 9cd059d3e2f51798 :: test jets the macro AUC is computed on
\newcommand{\FtNTestAccJc}{2{,}000{,}000}  % experiments/FIGS/data/w2b_leg2_metrics_v2.json :: cells.*.*.s1.n_jets :: sha256 9cd059d3e2f51798 :: test jets the accuracy is computed on
\newcommand{\FtSizeEThree}{\ensuremath{10^{3}}}  % experiments/FIGS/data/w2b_leg1_metrics_v2.json :: cells.*.N1000 :: sha256 b56bf70be9982044 :: fine-tuning training jets
\newcommand{\FtSizeEFour}{\ensuremath{10^{4}}}  % experiments/FIGS/data/w2b_leg1_metrics_v2.json :: cells.*.N10000 :: sha256 b56bf70be9982044 :: fine-tuning training jets
\newcommand{\FtSizeEFive}{\ensuremath{10^{5}}}  % experiments/FIGS/data/w2b_leg1_metrics_v2.json :: cells.*.N100000 :: sha256 b56bf70be9982044 :: fine-tuning training jets
\newcommand{\FtSizeESix}{\ensuremath{10^{6}}}  % experiments/FIGS/data/w2b_leg1_metrics_v2.json :: cells.*.N1000000 :: sha256 b56bf70be9982044 :: fine-tuning training jets
\newcommand{\FtNRunsWord}{five}  % experiments/FIGS/data/w2b_leg1_metrics_v2.json :: cells (pretrained models per vocabulary row) :: sha256 b56bf70be9982044 :: pretraining runs per vocabulary
\newcommand{\PairedProbeBcVsRestLinearOnesevenOverOnesixtwo}{1.64~[1.52, 1.78]}  % experiments/FIGS/data/paired_v1err/probes/ratios.json :: ratios[0].ratio, ratios[0].ci95 :: sha256 14dd51402305d147 :: 17 over 162, bc_vs_rest linear 1-auc: paired geometric mean [95% interval]
\newcommand{\PairedProbeBcVsRestMlpOnesevenOverOnesixtwo}{1.50~[1.39, 1.62]}  % experiments/FIGS/data/paired_v1err/probes/ratios.json :: ratios[3].ratio, ratios[3].ci95 :: sha256 14dd51402305d147 :: 17 over 162, bc_vs_rest mlp 1-auc: paired geometric mean [95% interval]
\newcommand{\PairedProbeBvcFourprongLinearOnesixtwoOverOneeighteight}{1.01~[0.95, 1.08]}  % experiments/FIGS/data/paired_v1err/probes/ratios.json :: ratios[6].ratio, ratios[6].ci95 :: sha256 14dd51402305d147 :: 162 over 188, bvc_4prong linear 1-auc: paired geometric mean [95% interval]
\newcommand{\PairedProbeBvcFourprongLinearFourthreeOverOneeighteight}{1.10~[1.01, 1.20]}  % experiments/FIGS/data/paired_v1err/probes/ratios.json :: ratios[7].ratio, ratios[7].ci95 :: sha256 14dd51402305d147 :: 43 over 188, bvc_4prong linear 1-auc: paired geometric mean [95% interval]
\newcommand{\PairedProbeBvcFourprongLinearOnesevenOverOneeighteight}{4.5~[3.5, 5.7]}  % experiments/FIGS/data/paired_v1err/probes/ratios.json :: ratios[8].ratio, ratios[8].ci95 :: sha256 14dd51402305d147 :: 17 over 188, bvc_4prong linear 1-auc: paired geometric mean [95% interval]
\newcommand{\PairedProbeBvcFourprongLinearFourthreeOverOnesixtwo}{1.09~[0.98, 1.20]}  % experiments/FIGS/data/paired_v1err/probes/ratios.json :: ratios[9].ratio, ratios[9].ci95 :: sha256 14dd51402305d147 :: 43 over 162, bvc_4prong linear 1-auc: paired geometric mean [95% interval]
\newcommand{\PairedProbeBvcFourprongLinearOnesevenOverOnesixtwo}{4.4~[3.4, 5.6]}  % experiments/FIGS/data/paired_v1err/probes/ratios.json :: ratios[10].ratio, ratios[10].ci95 :: sha256 14dd51402305d147 :: 17 over 162, bvc_4prong linear 1-auc: paired geometric mean [95% interval]
\newcommand{\PairedProbeBvcFourprongLinearOnesevenOverFourthree}{4.0~[3.1, 5.4]}  % experiments/FIGS/data/paired_v1err/probes/ratios.json :: ratios[11].ratio, ratios[11].ci95 :: sha256 14dd51402305d147 :: 17 over 43, bvc_4prong linear 1-auc: paired geometric mean [95% interval]
\newcommand{\PairedProbeBvcFourprongLinearOnesixtwoMassOverOnesixtwo}{1.05~[0.97, 1.13]}  % experiments/FIGS/data/paired_v1err/probes/ratios.json :: ratios[12].ratio, ratios[12].ci95 :: sha256 14dd51402305d147 :: 162+mass over 162, bvc_4prong linear 1-auc: paired geometric mean [95% interval]
\newcommand{\PairedProbeBvcFourprongLinearOnesevenMassOverOneseven}{2.1~[1.6, 2.7]}  % experiments/FIGS/data/paired_v1err/probes/ratios.json :: ratios[13].ratio, ratios[13].ci95 :: sha256 14dd51402305d147 :: 17+mass over 17, bvc_4prong linear 1-auc: paired geometric mean [95% interval]
\newcommand{\PairedProbeBvcFourprongLinearRandomOverOneeighteight}{1.54~[1.37, 1.74]}  % experiments/FIGS/data/paired_v1err/probes/ratios.json :: ratios[14].ratio, ratios[14].ci95 :: sha256 14dd51402305d147 :: random over 188, bvc_4prong linear 1-auc: paired geometric mean [95% interval]
\newcommand{\PairedProbeBvcFourprongLinearRandomOverOnesixtwo}{1.55~[1.36, 1.77]}  % experiments/FIGS/data/paired_v1err/probes/ratios.json :: ratios[15].ratio, ratios[15].ci95 :: sha256 14dd51402305d147 :: random over 162, bvc_4prong linear 1-auc: paired geometric mean [95% interval]
\newcommand{\PairedProbeBvcFourprongLinearRandomOverFourthree}{1.39~[1.27, 1.52]}  % experiments/FIGS/data/paired_v1err/probes/ratios.json :: ratios[16].ratio, ratios[16].ci95 :: sha256 14dd51402305d147 :: random over 43, bvc_4prong linear 1-auc: paired geometric mean [95% interval]
\newcommand{\PairedProbeBvcFourprongLinearRandomOverOneseven}{0.34~[0.20, 0.58]}  % experiments/FIGS/data/paired_v1err/probes/ratios.json :: ratios[17].ratio, ratios[17].ci95 :: sha256 14dd51402305d147 :: random over 17, bvc_4prong linear 1-auc: paired geometric mean [95% interval]
\newcommand{\PairedProbeBvcFourprongLinearRandomDrawTwoOverRandomDrawOne}{1.05~[0.94, 1.18]}  % experiments/FIGS/data/paired_v1err/probes/ratios.json :: ratios[18].ratio, ratios[18].ci95 :: sha256 14dd51402305d147 :: random draw 2 over random draw 1, bvc_4prong linear 1-auc: paired geometric mean [95% interval]
\newcommand{\PairedProbeBvcFourprongLinearRandomDrawThreeOverRandomDrawOne}{0.95~[0.85, 1.06]}  % experiments/FIGS/data/paired_v1err/probes/ratios.json :: ratios[19].ratio, ratios[19].ci95 :: sha256 14dd51402305d147 :: random draw 3 over random draw 1, bvc_4prong linear 1-auc: paired geometric mean [95% interval]
\newcommand{\PairedProbeBvcFourprongLinearRandomDrawThreeOverRandomDrawTwo}{0.90~[0.79, 1.02]}  % experiments/FIGS/data/paired_v1err/probes/ratios.json :: ratios[20].ratio, ratios[20].ci95 :: sha256 14dd51402305d147 :: random draw 3 over random draw 2, bvc_4prong linear 1-auc: paired geometric mean [95% interval]
\newcommand{\PairedProbeBvcFourprongMlpOnesixtwoOverOneeighteight}{1.00~[0.93, 1.07]}  % experiments/FIGS/data/paired_v1err/probes/ratios.json :: ratios[66].ratio, ratios[66].ci95 :: sha256 14dd51402305d147 :: 162 over 188, bvc_4prong mlp 1-auc: paired geometric mean [95% interval]
\newcommand{\PairedProbeBvcFourprongMlpFourthreeOverOneeighteight}{1.12~[1.03, 1.22]}  % experiments/FIGS/data/paired_v1err/probes/ratios.json :: ratios[67].ratio, ratios[67].ci95 :: sha256 14dd51402305d147 :: 43 over 188, bvc_4prong mlp 1-auc: paired geometric mean [95% interval]
\newcommand{\PairedProbeBvcFourprongMlpOnesevenOverOneeighteight}{3.2~[2.6, 4.1]}  % experiments/FIGS/data/paired_v1err/probes/ratios.json :: ratios[68].ratio, ratios[68].ci95 :: sha256 14dd51402305d147 :: 17 over 188, bvc_4prong mlp 1-auc: paired geometric mean [95% interval]
\newcommand{\PairedProbeBvcFourprongMlpFourthreeOverOnesixtwo}{1.12~[1.01, 1.26]}  % experiments/FIGS/data/paired_v1err/probes/ratios.json :: ratios[69].ratio, ratios[69].ci95 :: sha256 14dd51402305d147 :: 43 over 162, bvc_4prong mlp 1-auc: paired geometric mean [95% interval]
\newcommand{\PairedProbeBvcFourprongMlpOnesevenOverOnesixtwo}{3.3~[2.6, 4.0]}  % experiments/FIGS/data/paired_v1err/probes/ratios.json :: ratios[70].ratio, ratios[70].ci95 :: sha256 14dd51402305d147 :: 17 over 162, bvc_4prong mlp 1-auc: paired geometric mean [95% interval]
\newcommand{\PairedProbeBvcFourprongMlpOnesevenOverFourthree}{2.9~[2.3, 3.7]}  % experiments/FIGS/data/paired_v1err/probes/ratios.json :: ratios[71].ratio, ratios[71].ci95 :: sha256 14dd51402305d147 :: 17 over 43, bvc_4prong mlp 1-auc: paired geometric mean [95% interval]
\newcommand{\PairedProbeBvcFourprongMlpOnesixtwoMassOverOnesixtwo}{1.07~[0.99, 1.16]}  % experiments/FIGS/data/paired_v1err/probes/ratios.json :: ratios[72].ratio, ratios[72].ci95 :: sha256 14dd51402305d147 :: 162+mass over 162, bvc_4prong mlp 1-auc: paired geometric mean [95% interval]
\newcommand{\PairedProbeBvcFourprongMlpOnesevenMassOverOneseven}{1.97~[1.46, 2.67]}  % experiments/FIGS/data/paired_v1err/probes/ratios.json :: ratios[73].ratio, ratios[73].ci95 :: sha256 14dd51402305d147 :: 17+mass over 17, bvc_4prong mlp 1-auc: paired geometric mean [95% interval]
\newcommand{\PairedProbeBvcFourprongMlpRandomOverOneeighteight}{1.47~[1.30, 1.65]}  % experiments/FIGS/data/paired_v1err/probes/ratios.json :: ratios[74].ratio, ratios[74].ci95 :: sha256 14dd51402305d147 :: random over 188, bvc_4prong mlp 1-auc: paired geometric mean [95% interval]
\newcommand{\PairedProbeBvcFourprongMlpRandomOverOnesixtwo}{1.49~[1.24, 1.79]}  % experiments/FIGS/data/paired_v1err/probes/ratios.json :: ratios[75].ratio, ratios[75].ci95 :: sha256 14dd51402305d147 :: random over 162, bvc_4prong mlp 1-auc: paired geometric mean [95% interval]
\newcommand{\PairedProbeBvcFourprongMlpRandomOverFourthree}{1.28~[1.16, 1.43]}  % experiments/FIGS/data/paired_v1err/probes/ratios.json :: ratios[76].ratio, ratios[76].ci95 :: sha256 14dd51402305d147 :: random over 43, bvc_4prong mlp 1-auc: paired geometric mean [95% interval]
\newcommand{\PairedProbeBvcFourprongMlpRandomOverOneseven}{0.43~[0.28, 0.67]}  % experiments/FIGS/data/paired_v1err/probes/ratios.json :: ratios[77].ratio, ratios[77].ci95 :: sha256 14dd51402305d147 :: random over 17, bvc_4prong mlp 1-auc: paired geometric mean [95% interval]
\newcommand{\PairedProbeBvcFourprongMlpRandomDrawTwoOverRandomDrawOne}{1.00~[0.90, 1.10]}  % experiments/FIGS/data/paired_v1err/probes/ratios.json :: ratios[78].ratio, ratios[78].ci95 :: sha256 14dd51402305d147 :: random draw 2 over random draw 1, bvc_4prong mlp 1-auc: paired geometric mean [95% interval]
\newcommand{\PairedProbeBvcFourprongMlpRandomDrawThreeOverRandomDrawOne}{0.89~[0.79, 0.99]}  % experiments/FIGS/data/paired_v1err/probes/ratios.json :: ratios[79].ratio, ratios[79].ci95 :: sha256 14dd51402305d147 :: random draw 3 over random draw 1, bvc_4prong mlp 1-auc: paired geometric mean [95% interval]
\newcommand{\PairedProbeBvcFourprongMlpRandomDrawThreeOverRandomDrawTwo}{0.89~[0.79, 0.99]}  % experiments/FIGS/data/paired_v1err/probes/ratios.json :: ratios[80].ratio, ratios[80].ci95 :: sha256 14dd51402305d147 :: random draw 3 over random draw 2, bvc_4prong mlp 1-auc: paired geometric mean [95% interval]
\newcommand{\PairedProbeBvcQcdLinearOnesixtwoOverOneeighteight}{1.15~[1.07, 1.23]}  % experiments/FIGS/data/paired_v1err/probes/ratios.json :: ratios[126].ratio, ratios[126].ci95 :: sha256 14dd51402305d147 :: 162 over 188, bvc_qcd linear 1-auc: paired geometric mean [95% interval]
\newcommand{\PairedProbeBvcQcdLinearFourthreeOverOneeighteight}{1.16~[1.08, 1.25]}  % experiments/FIGS/data/paired_v1err/probes/ratios.json :: ratios[127].ratio, ratios[127].ci95 :: sha256 14dd51402305d147 :: 43 over 188, bvc_qcd linear 1-auc: paired geometric mean [95% interval]
\newcommand{\PairedProbeBvcQcdLinearOnesevenOverOneeighteight}{3.0~[2.5, 3.7]}  % experiments/FIGS/data/paired_v1err/probes/ratios.json :: ratios[128].ratio, ratios[128].ci95 :: sha256 14dd51402305d147 :: 17 over 188, bvc_qcd linear 1-auc: paired geometric mean [95% interval]
\newcommand{\PairedProbeBvcQcdLinearFourthreeOverOnesixtwo}{1.01~[0.96, 1.07]}  % experiments/FIGS/data/paired_v1err/probes/ratios.json :: ratios[129].ratio, ratios[129].ci95 :: sha256 14dd51402305d147 :: 43 over 162, bvc_qcd linear 1-auc: paired geometric mean [95% interval]
\newcommand{\PairedProbeBvcQcdLinearOnesevenOverOnesixtwo}{2.6~[2.2, 3.1]}  % experiments/FIGS/data/paired_v1err/probes/ratios.json :: ratios[130].ratio, ratios[130].ci95 :: sha256 14dd51402305d147 :: 17 over 162, bvc_qcd linear 1-auc: paired geometric mean [95% interval]
\newcommand{\PairedProbeBvcQcdLinearOnesevenOverFourthree}{2.6~[2.2, 3.0]}  % experiments/FIGS/data/paired_v1err/probes/ratios.json :: ratios[131].ratio, ratios[131].ci95 :: sha256 14dd51402305d147 :: 17 over 43, bvc_qcd linear 1-auc: paired geometric mean [95% interval]
\newcommand{\PairedProbeBvcQcdLinearOnesixtwoMassOverOnesixtwo}{1.060~[1.002, 1.123]}  % experiments/FIGS/data/paired_v1err/probes/ratios.json :: ratios[132].ratio, ratios[132].ci95 :: sha256 14dd51402305d147 :: 162+mass over 162, bvc_qcd linear 1-auc: paired geometric mean [95% interval]
\newcommand{\PairedProbeBvcQcdLinearOnesevenMassOverOneseven}{1.89~[1.52, 2.34]}  % experiments/FIGS/data/paired_v1err/probes/ratios.json :: ratios[133].ratio, ratios[133].ci95 :: sha256 14dd51402305d147 :: 17+mass over 17, bvc_qcd linear 1-auc: paired geometric mean [95% interval]
\newcommand{\PairedProbeBvcQcdMlpOnesixtwoOverOneeighteight}{1.11~[1.04, 1.19]}  % experiments/FIGS/data/paired_v1err/probes/ratios.json :: ratios[158].ratio, ratios[158].ci95 :: sha256 14dd51402305d147 :: 162 over 188, bvc_qcd mlp 1-auc: paired geometric mean [95% interval]
\newcommand{\PairedProbeBvcQcdMlpFourthreeOverOneeighteight}{1.13~[1.05, 1.21]}  % experiments/FIGS/data/paired_v1err/probes/ratios.json :: ratios[159].ratio, ratios[159].ci95 :: sha256 14dd51402305d147 :: 43 over 188, bvc_qcd mlp 1-auc: paired geometric mean [95% interval]
\newcommand{\PairedProbeBvcQcdMlpOnesevenOverOneeighteight}{2.6~[2.2, 3.0]}  % experiments/FIGS/data/paired_v1err/probes/ratios.json :: ratios[160].ratio, ratios[160].ci95 :: sha256 14dd51402305d147 :: 17 over 188, bvc_qcd mlp 1-auc: paired geometric mean [95% interval]
\newcommand{\PairedProbeBvcQcdMlpFourthreeOverOnesixtwo}{1.01~[0.96, 1.06]}  % experiments/FIGS/data/paired_v1err/probes/ratios.json :: ratios[161].ratio, ratios[161].ci95 :: sha256 14dd51402305d147 :: 43 over 162, bvc_qcd mlp 1-auc: paired geometric mean [95% interval]
\newcommand{\PairedProbeBvcQcdMlpOnesevenOverOnesixtwo}{2.3~[2.0, 2.6]}  % experiments/FIGS/data/paired_v1err/probes/ratios.json :: ratios[162].ratio, ratios[162].ci95 :: sha256 14dd51402305d147 :: 17 over 162, bvc_qcd mlp 1-auc: paired geometric mean [95% interval]
\newcommand{\PairedProbeBvcQcdMlpOnesevenOverFourthree}{2.3~[2.0, 2.6]}  % experiments/FIGS/data/paired_v1err/probes/ratios.json :: ratios[163].ratio, ratios[163].ci95 :: sha256 14dd51402305d147 :: 17 over 43, bvc_qcd mlp 1-auc: paired geometric mean [95% interval]
\newcommand{\PairedProbeBvcQcdMlpOnesixtwoMassOverOnesixtwo}{1.06~[1.01, 1.12]}  % experiments/FIGS/data/paired_v1err/probes/ratios.json :: ratios[164].ratio, ratios[164].ci95 :: sha256 14dd51402305d147 :: 162+mass over 162, bvc_qcd mlp 1-auc: paired geometric mean [95% interval]
\newcommand{\PairedProbeBvcQcdMlpOnesevenMassOverOneseven}{1.64~[1.33, 2.02]}  % experiments/FIGS/data/paired_v1err/probes/ratios.json :: ratios[165].ratio, ratios[165].ci95 :: sha256 14dd51402305d147 :: 17+mass over 17, bvc_qcd mlp 1-auc: paired geometric mean [95% interval]
\newcommand{\PairedProbeBvcResonantLinearOnesixtwoOverOneeighteight}{1.02~[0.96, 1.08]}  % experiments/FIGS/data/paired_v1err/probes/ratios.json :: ratios[190].ratio, ratios[190].ci95 :: sha256 14dd51402305d147 :: 162 over 188, bvc_resonant linear 1-auc: paired geometric mean [95% interval]
\newcommand{\PairedProbeBvcResonantLinearFourthreeOverOneeighteight}{1.09~[1.03, 1.15]}  % experiments/FIGS/data/paired_v1err/probes/ratios.json :: ratios[191].ratio, ratios[191].ci95 :: sha256 14dd51402305d147 :: 43 over 188, bvc_resonant linear 1-auc: paired geometric mean [95% interval]
\newcommand{\PairedProbeBvcResonantLinearOnesevenOverOneeighteight}{4.0~[3.3, 4.9]}  % experiments/FIGS/data/paired_v1err/probes/ratios.json :: ratios[192].ratio, ratios[192].ci95 :: sha256 14dd51402305d147 :: 17 over 188, bvc_resonant linear 1-auc: paired geometric mean [95% interval]
\newcommand{\PairedProbeBvcResonantLinearFourthreeOverOnesixtwo}{1.07~[1.02, 1.12]}  % experiments/FIGS/data/paired_v1err/probes/ratios.json :: ratios[193].ratio, ratios[193].ci95 :: sha256 14dd51402305d147 :: 43 over 162, bvc_resonant linear 1-auc: paired geometric mean [95% interval]
\newcommand{\PairedProbeBvcResonantLinearOnesevenOverOnesixtwo}{3.9~[3.1, 4.9]}  % experiments/FIGS/data/paired_v1err/probes/ratios.json :: ratios[194].ratio, ratios[194].ci95 :: sha256 14dd51402305d147 :: 17 over 162, bvc_resonant linear 1-auc: paired geometric mean [95% interval]
\newcommand{\PairedProbeBvcResonantLinearOnesevenOverFourthree}{3.7~[3.0, 4.5]}  % experiments/FIGS/data/paired_v1err/probes/ratios.json :: ratios[195].ratio, ratios[195].ci95 :: sha256 14dd51402305d147 :: 17 over 43, bvc_resonant linear 1-auc: paired geometric mean [95% interval]
\newcommand{\PairedProbeBvcResonantLinearOnesixtwoMassOverOnesixtwo}{1.11~[1.06, 1.16]}  % experiments/FIGS/data/paired_v1err/probes/ratios.json :: ratios[196].ratio, ratios[196].ci95 :: sha256 14dd51402305d147 :: 162+mass over 162, bvc_resonant linear 1-auc: paired geometric mean [95% interval]
\newcommand{\PairedProbeBvcResonantLinearOnesevenMassOverOneseven}{3.0~[1.9, 4.7]}  % experiments/FIGS/data/paired_v1err/probes/ratios.json :: ratios[197].ratio, ratios[197].ci95 :: sha256 14dd51402305d147 :: 17+mass over 17, bvc_resonant linear 1-auc: paired geometric mean [95% interval]
\newcommand{\PairedProbeBvcResonantMlpOnesixtwoOverOneeighteight}{1.04~[0.98, 1.12]}  % experiments/FIGS/data/paired_v1err/probes/ratios.json :: ratios[222].ratio, ratios[222].ci95 :: sha256 14dd51402305d147 :: 162 over 188, bvc_resonant mlp 1-auc: paired geometric mean [95% interval]
\newcommand{\PairedProbeBvcResonantMlpFourthreeOverOneeighteight}{1.14~[1.08, 1.21]}  % experiments/FIGS/data/paired_v1err/probes/ratios.json :: ratios[223].ratio, ratios[223].ci95 :: sha256 14dd51402305d147 :: 43 over 188, bvc_resonant mlp 1-auc: paired geometric mean [95% interval]
\newcommand{\PairedProbeBvcResonantMlpOnesevenOverOneeighteight}{3.3~[2.8, 3.8]}  % experiments/FIGS/data/paired_v1err/probes/ratios.json :: ratios[224].ratio, ratios[224].ci95 :: sha256 14dd51402305d147 :: 17 over 188, bvc_resonant mlp 1-auc: paired geometric mean [95% interval]
\newcommand{\PairedProbeBvcResonantMlpFourthreeOverOnesixtwo}{1.09~[1.04, 1.14]}  % experiments/FIGS/data/paired_v1err/probes/ratios.json :: ratios[225].ratio, ratios[225].ci95 :: sha256 14dd51402305d147 :: 43 over 162, bvc_resonant mlp 1-auc: paired geometric mean [95% interval]
\newcommand{\PairedProbeBvcResonantMlpOnesevenOverOnesixtwo}{3.1~[2.6, 3.8]}  % experiments/FIGS/data/paired_v1err/probes/ratios.json :: ratios[226].ratio, ratios[226].ci95 :: sha256 14dd51402305d147 :: 17 over 162, bvc_resonant mlp 1-auc: paired geometric mean [95% interval]
\newcommand{\PairedProbeBvcResonantMlpOnesevenOverFourthree}{2.9~[2.4, 3.4]}  % experiments/FIGS/data/paired_v1err/probes/ratios.json :: ratios[227].ratio, ratios[227].ci95 :: sha256 14dd51402305d147 :: 17 over 43, bvc_resonant mlp 1-auc: paired geometric mean [95% interval]
\newcommand{\PairedProbeBvcResonantMlpOnesixtwoMassOverOnesixtwo}{1.09~[1.04, 1.14]}  % experiments/FIGS/data/paired_v1err/probes/ratios.json :: ratios[228].ratio, ratios[228].ci95 :: sha256 14dd51402305d147 :: 162+mass over 162, bvc_resonant mlp 1-auc: paired geometric mean [95% interval]
\newcommand{\PairedProbeBvcResonantMlpOnesevenMassOverOneseven}{1.95~[1.45, 2.61]}  % experiments/FIGS/data/paired_v1err/probes/ratios.json :: ratios[229].ratio, ratios[229].ci95 :: sha256 14dd51402305d147 :: 17+mass over 17, bvc_resonant mlp 1-auc: paired geometric mean [95% interval]
\newcommand{\PairedProbeEeVsMmLinearOnesevenOverFourthree}{2.3~[0.8, 6.5]}  % experiments/FIGS/data/paired_v1err/probes/ratios.json :: ratios[259].ratio, ratios[259].ci95 :: sha256 14dd51402305d147 :: 17 over 43, ee_vs_mm linear 1-auc: paired geometric mean [95% interval]
\newcommand{\PairedProbeEeVsMmLinearOnesevenMassOverOneseven}{0.23~[0.08, 0.64]}  % experiments/FIGS/data/paired_v1err/probes/ratios.json :: ratios[261].ratio, ratios[261].ci95 :: sha256 14dd51402305d147 :: 17+mass over 17, ee_vs_mm linear 1-auc: paired geometric mean [95% interval]
\newcommand{\PairedProbeEeVsMmMlpOnesevenOverFourthree}{2.8~[0.8, 9.0]}  % experiments/FIGS/data/paired_v1err/probes/ratios.json :: ratios[291].ratio, ratios[291].ci95 :: sha256 14dd51402305d147 :: 17 over 43, ee_vs_mm mlp 1-auc: paired geometric mean [95% interval]
\newcommand{\PairedProbeEeVsMmMlpOnesevenMassOverOneseven}{0.172~[0.026, 1.125]}  % experiments/FIGS/data/paired_v1err/probes/ratios.json :: ratios[293].ratio, ratios[293].ci95 :: sha256 14dd51402305d147 :: 17+mass over 17, ee_vs_mm mlp 1-auc: paired geometric mean [95% interval]
\newcommand{\PairedProbeRetainedTopologyLinearOnesixtwoOverOneeighteight}{1.00~[0.92, 1.08]}  % experiments/FIGS/data/paired_v1err/probes/ratios.json :: ratios[318].ratio, ratios[318].ci95 :: sha256 14dd51402305d147 :: 162 over 188, retained_topology linear 1-auc: paired geometric mean [95% interval]
\newcommand{\PairedProbeRetainedTopologyLinearFourthreeOverOneeighteight}{1.24~[1.15, 1.34]}  % experiments/FIGS/data/paired_v1err/probes/ratios.json :: ratios[319].ratio, ratios[319].ci95 :: sha256 14dd51402305d147 :: 43 over 188, retained_topology linear 1-auc: paired geometric mean [95% interval]
\newcommand{\PairedProbeRetainedTopologyLinearOnesevenOverOneeighteight}{1.54~[1.35, 1.75]}  % experiments/FIGS/data/paired_v1err/probes/ratios.json :: ratios[320].ratio, ratios[320].ci95 :: sha256 14dd51402305d147 :: 17 over 188, retained_topology linear 1-auc: paired geometric mean [95% interval]
\newcommand{\PairedProbeRetainedTopologyLinearFourthreeOverOnesixtwo}{1.24~[1.14, 1.34]}  % experiments/FIGS/data/paired_v1err/probes/ratios.json :: ratios[321].ratio, ratios[321].ci95 :: sha256 14dd51402305d147 :: 43 over 162, retained_topology linear 1-auc: paired geometric mean [95% interval]
\newcommand{\PairedProbeRetainedTopologyLinearOnesevenOverOnesixtwo}{1.54~[1.33, 1.78]}  % experiments/FIGS/data/paired_v1err/probes/ratios.json :: ratios[322].ratio, ratios[322].ci95 :: sha256 14dd51402305d147 :: 17 over 162, retained_topology linear 1-auc: paired geometric mean [95% interval]
\newcommand{\PairedProbeRetainedTopologyLinearOnesevenOverFourthree}{1.24~[1.08, 1.42]}  % experiments/FIGS/data/paired_v1err/probes/ratios.json :: ratios[323].ratio, ratios[323].ci95 :: sha256 14dd51402305d147 :: 17 over 43, retained_topology linear 1-auc: paired geometric mean [95% interval]
\newcommand{\PairedProbeRetainedTopologyLinearOnesixtwoMassOverOnesixtwo}{1.06~[0.99, 1.14]}  % experiments/FIGS/data/paired_v1err/probes/ratios.json :: ratios[324].ratio, ratios[324].ci95 :: sha256 14dd51402305d147 :: 162+mass over 162, retained_topology linear 1-auc: paired geometric mean [95% interval]
\newcommand{\PairedProbeRetainedTopologyLinearOnesevenMassOverOneseven}{1.12~[0.97, 1.28]}  % experiments/FIGS/data/paired_v1err/probes/ratios.json :: ratios[325].ratio, ratios[325].ci95 :: sha256 14dd51402305d147 :: 17+mass over 17, retained_topology linear 1-auc: paired geometric mean [95% interval]
\newcommand{\PairedProbeRetainedTopologyMlpOnesixtwoOverOneeighteight}{1.01~[0.87, 1.17]}  % experiments/FIGS/data/paired_v1err/probes/ratios.json :: ratios[350].ratio, ratios[350].ci95 :: sha256 14dd51402305d147 :: 162 over 188, retained_topology mlp 1-auc: paired geometric mean [95% interval]
\newcommand{\PairedProbeRetainedTopologyMlpFourthreeOverOneeighteight}{1.18~[1.03, 1.36]}  % experiments/FIGS/data/paired_v1err/probes/ratios.json :: ratios[351].ratio, ratios[351].ci95 :: sha256 14dd51402305d147 :: 43 over 188, retained_topology mlp 1-auc: paired geometric mean [95% interval]
\newcommand{\PairedProbeRetainedTopologyMlpOnesevenOverOneeighteight}{1.52~[1.34, 1.72]}  % experiments/FIGS/data/paired_v1err/probes/ratios.json :: ratios[352].ratio, ratios[352].ci95 :: sha256 14dd51402305d147 :: 17 over 188, retained_topology mlp 1-auc: paired geometric mean [95% interval]
\newcommand{\PairedProbeRetainedTopologyMlpFourthreeOverOnesixtwo}{1.17~[1.07, 1.28]}  % experiments/FIGS/data/paired_v1err/probes/ratios.json :: ratios[353].ratio, ratios[353].ci95 :: sha256 14dd51402305d147 :: 43 over 162, retained_topology mlp 1-auc: paired geometric mean [95% interval]
\newcommand{\PairedProbeRetainedTopologyMlpOnesevenOverOnesixtwo}{1.50~[1.29, 1.75]}  % experiments/FIGS/data/paired_v1err/probes/ratios.json :: ratios[354].ratio, ratios[354].ci95 :: sha256 14dd51402305d147 :: 17 over 162, retained_topology mlp 1-auc: paired geometric mean [95% interval]
\newcommand{\PairedProbeRetainedTopologyMlpOnesevenOverFourthree}{1.28~[1.12, 1.47]}  % experiments/FIGS/data/paired_v1err/probes/ratios.json :: ratios[355].ratio, ratios[355].ci95 :: sha256 14dd51402305d147 :: 17 over 43, retained_topology mlp 1-auc: paired geometric mean [95% interval]
\newcommand{\PairedProbeRetainedTopologyMlpOnesixtwoMassOverOnesixtwo}{1.04~[0.92, 1.17]}  % experiments/FIGS/data/paired_v1err/probes/ratios.json :: ratios[356].ratio, ratios[356].ci95 :: sha256 14dd51402305d147 :: 162+mass over 162, retained_topology mlp 1-auc: paired geometric mean [95% interval]
\newcommand{\PairedProbeRetainedTopologyMlpOnesevenMassOverOneseven}{1.11~[0.99, 1.23]}  % experiments/FIGS/data/paired_v1err/probes/ratios.json :: ratios[357].ratio, ratios[357].ci95 :: sha256 14dd51402305d147 :: 17+mass over 17, retained_topology mlp 1-auc: paired geometric mean [95% interval]
\newcommand{\PairedProbeVisibleContentLinearOnesixtwoOverOneeighteight}{0.95~[0.88, 1.03]}  % experiments/FIGS/data/paired_v1err/probes/ratios.json :: ratios[382].ratio, ratios[382].ci95 :: sha256 14dd51402305d147 :: 162 over 188, visible_content linear 1-auc: paired geometric mean [95% interval]
\newcommand{\PairedProbeVisibleContentLinearFourthreeOverOneeighteight}{0.99~[0.85, 1.15]}  % experiments/FIGS/data/paired_v1err/probes/ratios.json :: ratios[383].ratio, ratios[383].ci95 :: sha256 14dd51402305d147 :: 43 over 188, visible_content linear 1-auc: paired geometric mean [95% interval]
\newcommand{\PairedProbeVisibleContentLinearOnesevenOverOneeighteight}{2.8~[2.1, 3.7]}  % experiments/FIGS/data/paired_v1err/probes/ratios.json :: ratios[384].ratio, ratios[384].ci95 :: sha256 14dd51402305d147 :: 17 over 188, visible_content linear 1-auc: paired geometric mean [95% interval]
\newcommand{\PairedProbeVisibleContentLinearFourthreeOverOnesixtwo}{1.04~[0.89, 1.20]}  % experiments/FIGS/data/paired_v1err/probes/ratios.json :: ratios[385].ratio, ratios[385].ci95 :: sha256 14dd51402305d147 :: 43 over 162, visible_content linear 1-auc: paired geometric mean [95% interval]
\newcommand{\PairedProbeVisibleContentLinearOnesevenOverOnesixtwo}{2.9~[2.3, 3.8]}  % experiments/FIGS/data/paired_v1err/probes/ratios.json :: ratios[386].ratio, ratios[386].ci95 :: sha256 14dd51402305d147 :: 17 over 162, visible_content linear 1-auc: paired geometric mean [95% interval]
\newcommand{\PairedProbeVisibleContentLinearOnesevenOverFourthree}{2.8~[2.2, 3.7]}  % experiments/FIGS/data/paired_v1err/probes/ratios.json :: ratios[387].ratio, ratios[387].ci95 :: sha256 14dd51402305d147 :: 17 over 43, visible_content linear 1-auc: paired geometric mean [95% interval]
\newcommand{\PairedProbeVisibleContentLinearOnesixtwoMassOverOnesixtwo}{1.10~[0.98, 1.23]}  % experiments/FIGS/data/paired_v1err/probes/ratios.json :: ratios[388].ratio, ratios[388].ci95 :: sha256 14dd51402305d147 :: 162+mass over 162, visible_content linear 1-auc: paired geometric mean [95% interval]
\newcommand{\PairedProbeVisibleContentLinearOnesevenMassOverOneseven}{1.44~[1.26, 1.65]}  % experiments/FIGS/data/paired_v1err/probes/ratios.json :: ratios[389].ratio, ratios[389].ci95 :: sha256 14dd51402305d147 :: 17+mass over 17, visible_content linear 1-auc: paired geometric mean [95% interval]
\newcommand{\PairedProbeVisibleContentLinearRandomOverOneeighteight}{1.31~[1.01, 1.70]}  % experiments/FIGS/data/paired_v1err/probes/ratios.json :: ratios[390].ratio, ratios[390].ci95 :: sha256 14dd51402305d147 :: random over 188, visible_content linear 1-auc: paired geometric mean [95% interval]
\newcommand{\PairedProbeVisibleContentLinearRandomOverOnesixtwo}{1.41~[1.17, 1.70]}  % experiments/FIGS/data/paired_v1err/probes/ratios.json :: ratios[391].ratio, ratios[391].ci95 :: sha256 14dd51402305d147 :: random over 162, visible_content linear 1-auc: paired geometric mean [95% interval]
\newcommand{\PairedProbeVisibleContentLinearRandomOverFourthree}{1.33~[1.12, 1.57]}  % experiments/FIGS/data/paired_v1err/probes/ratios.json :: ratios[392].ratio, ratios[392].ci95 :: sha256 14dd51402305d147 :: random over 43, visible_content linear 1-auc: paired geometric mean [95% interval]
\newcommand{\PairedProbeVisibleContentLinearRandomOverOneseven}{0.48~[0.38, 0.60]}  % experiments/FIGS/data/paired_v1err/probes/ratios.json :: ratios[393].ratio, ratios[393].ci95 :: sha256 14dd51402305d147 :: random over 17, visible_content linear 1-auc: paired geometric mean [95% interval]
\newcommand{\PairedProbeVisibleContentLinearRandomDrawTwoOverRandomDrawOne}{0.79~[0.64, 0.98]}  % experiments/FIGS/data/paired_v1err/probes/ratios.json :: ratios[394].ratio, ratios[394].ci95 :: sha256 14dd51402305d147 :: random draw 2 over random draw 1, visible_content linear 1-auc: paired geometric mean [95% interval]
\newcommand{\PairedProbeVisibleContentLinearRandomDrawThreeOverRandomDrawOne}{0.86~[0.71, 1.05]}  % experiments/FIGS/data/paired_v1err/probes/ratios.json :: ratios[395].ratio, ratios[395].ci95 :: sha256 14dd51402305d147 :: random draw 3 over random draw 1, visible_content linear 1-auc: paired geometric mean [95% interval]
\newcommand{\PairedProbeVisibleContentLinearRandomDrawThreeOverRandomDrawTwo}{1.09~[0.88, 1.35]}  % experiments/FIGS/data/paired_v1err/probes/ratios.json :: ratios[396].ratio, ratios[396].ci95 :: sha256 14dd51402305d147 :: random draw 3 over random draw 2, visible_content linear 1-auc: paired geometric mean [95% interval]
\newcommand{\PairedProbeVisibleContentMlpOnesixtwoOverOneeighteight}{0.92~[0.85, 1.01]}  % experiments/FIGS/data/paired_v1err/probes/ratios.json :: ratios[442].ratio, ratios[442].ci95 :: sha256 14dd51402305d147 :: 162 over 188, visible_content mlp 1-auc: paired geometric mean [95% interval]
\newcommand{\PairedProbeVisibleContentMlpFourthreeOverOneeighteight}{0.93~[0.80, 1.07]}  % experiments/FIGS/data/paired_v1err/probes/ratios.json :: ratios[443].ratio, ratios[443].ci95 :: sha256 14dd51402305d147 :: 43 over 188, visible_content mlp 1-auc: paired geometric mean [95% interval]
\newcommand{\PairedProbeVisibleContentMlpOnesevenOverOneeighteight}{2.3~[1.7, 3.1]}  % experiments/FIGS/data/paired_v1err/probes/ratios.json :: ratios[444].ratio, ratios[444].ci95 :: sha256 14dd51402305d147 :: 17 over 188, visible_content mlp 1-auc: paired geometric mean [95% interval]
\newcommand{\PairedProbeVisibleContentMlpFourthreeOverOnesixtwo}{1.01~[0.89, 1.14]}  % experiments/FIGS/data/paired_v1err/probes/ratios.json :: ratios[445].ratio, ratios[445].ci95 :: sha256 14dd51402305d147 :: 43 over 162, visible_content mlp 1-auc: paired geometric mean [95% interval]
\newcommand{\PairedProbeVisibleContentMlpOnesevenOverOnesixtwo}{2.5~[1.9, 3.2]}  % experiments/FIGS/data/paired_v1err/probes/ratios.json :: ratios[446].ratio, ratios[446].ci95 :: sha256 14dd51402305d147 :: 17 over 162, visible_content mlp 1-auc: paired geometric mean [95% interval]
\newcommand{\PairedProbeVisibleContentMlpOnesevenOverFourthree}{2.5~[2.0, 3.1]}  % experiments/FIGS/data/paired_v1err/probes/ratios.json :: ratios[447].ratio, ratios[447].ci95 :: sha256 14dd51402305d147 :: 17 over 43, visible_content mlp 1-auc: paired geometric mean [95% interval]
\newcommand{\PairedProbeVisibleContentMlpOnesixtwoMassOverOnesixtwo}{1.08~[0.94, 1.24]}  % experiments/FIGS/data/paired_v1err/probes/ratios.json :: ratios[448].ratio, ratios[448].ci95 :: sha256 14dd51402305d147 :: 162+mass over 162, visible_content mlp 1-auc: paired geometric mean [95% interval]
\newcommand{\PairedProbeVisibleContentMlpOnesevenMassOverOneseven}{1.41~[1.19, 1.67]}  % experiments/FIGS/data/paired_v1err/probes/ratios.json :: ratios[449].ratio, ratios[449].ci95 :: sha256 14dd51402305d147 :: 17+mass over 17, visible_content mlp 1-auc: paired geometric mean [95% interval]
\newcommand{\PairedProbeVisibleContentMlpRandomOverOneeighteight}{1.27~[1.01, 1.59]}  % experiments/FIGS/data/paired_v1err/probes/ratios.json :: ratios[450].ratio, ratios[450].ci95 :: sha256 14dd51402305d147 :: random over 188, visible_content mlp 1-auc: paired geometric mean [95% interval]
\newcommand{\PairedProbeVisibleContentMlpRandomOverOnesixtwo}{1.37~[1.12, 1.68]}  % experiments/FIGS/data/paired_v1err/probes/ratios.json :: ratios[451].ratio, ratios[451].ci95 :: sha256 14dd51402305d147 :: random over 162, visible_content mlp 1-auc: paired geometric mean [95% interval]
\newcommand{\PairedProbeVisibleContentMlpRandomOverFourthree}{1.35~[1.13, 1.61]}  % experiments/FIGS/data/paired_v1err/probes/ratios.json :: ratios[452].ratio, ratios[452].ci95 :: sha256 14dd51402305d147 :: random over 43, visible_content mlp 1-auc: paired geometric mean [95% interval]
\newcommand{\PairedProbeVisibleContentMlpRandomOverOneseven}{0.56~[0.45, 0.70]}  % experiments/FIGS/data/paired_v1err/probes/ratios.json :: ratios[453].ratio, ratios[453].ci95 :: sha256 14dd51402305d147 :: random over 17, visible_content mlp 1-auc: paired geometric mean [95% interval]
\newcommand{\PairedProbeVisibleContentMlpRandomDrawTwoOverRandomDrawOne}{0.88~[0.71, 1.09]}  % experiments/FIGS/data/paired_v1err/probes/ratios.json :: ratios[454].ratio, ratios[454].ci95 :: sha256 14dd51402305d147 :: random draw 2 over random draw 1, visible_content mlp 1-auc: paired geometric mean [95% interval]
\newcommand{\PairedProbeVisibleContentMlpRandomDrawThreeOverRandomDrawOne}{0.94~[0.75, 1.18]}  % experiments/FIGS/data/paired_v1err/probes/ratios.json :: ratios[455].ratio, ratios[455].ci95 :: sha256 14dd51402305d147 :: random draw 3 over random draw 1, visible_content mlp 1-auc: paired geometric mean [95% interval]
\newcommand{\PairedProbeVisibleContentMlpRandomDrawThreeOverRandomDrawTwo}{1.06~[0.84, 1.35]}  % experiments/FIGS/data/paired_v1err/probes/ratios.json :: ratios[456].ratio, ratios[456].ci95 :: sha256 14dd51402305d147 :: random draw 3 over random draw 2, visible_content mlp 1-auc: paired geometric mean [95% interval]
\newcommand{\PairedFtJciiEThreeOnesixtwoOverOneeighteight}{0.95~[0.87, 1.05]}  % experiments/FIGS/data/paired_v1err/finetune/ratios.json :: ratios[0].ratio, ratios[0].ci95 :: sha256 3e5cf25e859b68c2 :: 162 over 188, leg1 N1000 1-macro_auc: paired geometric mean [95% interval]
\newcommand{\PairedFtJciiEThreeFourthreeOverOneeighteight}{0.84~[0.77, 0.91]}  % experiments/FIGS/data/paired_v1err/finetune/ratios.json :: ratios[1].ratio, ratios[1].ci95 :: sha256 3e5cf25e859b68c2 :: 43 over 188, leg1 N1000 1-macro_auc: paired geometric mean [95% interval]
\newcommand{\PairedFtJciiEThreeOnesevenOverOneeighteight}{0.97~[0.86, 1.09]}  % experiments/FIGS/data/paired_v1err/finetune/ratios.json :: ratios[2].ratio, ratios[2].ci95 :: sha256 3e5cf25e859b68c2 :: 17 over 188, leg1 N1000 1-macro_auc: paired geometric mean [95% interval]
\newcommand{\PairedFtJciiEThreeFourthreeOverOnesixtwo}{0.88~[0.80, 0.97]}  % experiments/FIGS/data/paired_v1err/finetune/ratios.json :: ratios[3].ratio, ratios[3].ci95 :: sha256 3e5cf25e859b68c2 :: 43 over 162, leg1 N1000 1-macro_auc: paired geometric mean [95% interval]
\newcommand{\PairedFtJciiEThreeOnesevenOverOnesixtwo}{1.02~[0.88, 1.18]}  % experiments/FIGS/data/paired_v1err/finetune/ratios.json :: ratios[4].ratio, ratios[4].ci95 :: sha256 3e5cf25e859b68c2 :: 17 over 162, leg1 N1000 1-macro_auc: paired geometric mean [95% interval]
\newcommand{\PairedFtJciiEThreeOnesevenOverFourthree}{1.16~[1.08, 1.24]}  % experiments/FIGS/data/paired_v1err/finetune/ratios.json :: ratios[5].ratio, ratios[5].ci95 :: sha256 3e5cf25e859b68c2 :: 17 over 43, leg1 N1000 1-macro_auc: paired geometric mean [95% interval]
\newcommand{\PairedFtJciiEThreeOnesixtwoMassOverOnesixtwo}{0.85~[0.77, 0.93]}  % experiments/FIGS/data/paired_v1err/finetune/ratios.json :: ratios[6].ratio, ratios[6].ci95 :: sha256 3e5cf25e859b68c2 :: 162+mass over 162, leg1 N1000 1-macro_auc: paired geometric mean [95% interval]
\newcommand{\PairedFtJciiEThreeOnesevenMassOverOneseven}{1.00~[0.86, 1.17]}  % experiments/FIGS/data/paired_v1err/finetune/ratios.json :: ratios[7].ratio, ratios[7].ci95 :: sha256 3e5cf25e859b68c2 :: 17+mass over 17, leg1 N1000 1-macro_auc: paired geometric mean [95% interval]
\newcommand{\PairedFtJciiEThreeRandomOverOneeighteight}{0.82~[0.73, 0.93]}  % experiments/FIGS/data/paired_v1err/finetune/ratios.json :: ratios[8].ratio, ratios[8].ci95 :: sha256 3e5cf25e859b68c2 :: random over 188, leg1 N1000 1-macro_auc: paired geometric mean [95% interval]
\newcommand{\PairedFtJciiEThreeRandomOverOnesixtwo}{0.90~[0.85, 0.96]}  % experiments/FIGS/data/paired_v1err/finetune/ratios.json :: ratios[9].ratio, ratios[9].ci95 :: sha256 3e5cf25e859b68c2 :: random over 162, leg1 N1000 1-macro_auc: paired geometric mean [95% interval]
\newcommand{\PairedFtJciiEThreeRandomOverFourthree}{0.99~[0.75, 1.30]}  % experiments/FIGS/data/paired_v1err/finetune/ratios.json :: ratios[10].ratio, ratios[10].ci95 :: sha256 3e5cf25e859b68c2 :: random over 43, leg1 N1000 1-macro_auc: paired geometric mean [95% interval]
\newcommand{\PairedFtJciiEThreeRandomOverOneseven}{0.86~[0.58, 1.27]}  % experiments/FIGS/data/paired_v1err/finetune/ratios.json :: ratios[11].ratio, ratios[11].ci95 :: sha256 3e5cf25e859b68c2 :: random over 17, leg1 N1000 1-macro_auc: paired geometric mean [95% interval]
\newcommand{\PairedFtJciiEThreeRandomDrawTwoOverRandomDrawOne}{1.09~[1.07, 1.10]}  % experiments/FIGS/data/paired_v1err/finetune/ratios.json :: ratios[12].ratio, ratios[12].ci95 :: sha256 3e5cf25e859b68c2 :: random draw 2 over random draw 1, leg1 N1000 1-macro_auc: paired geometric mean [95% interval]
\newcommand{\PairedFtJciiEThreeRandomDrawThreeOverRandomDrawOne}{1.04~[1.02, 1.05]}  % experiments/FIGS/data/paired_v1err/finetune/ratios.json :: ratios[13].ratio, ratios[13].ci95 :: sha256 3e5cf25e859b68c2 :: random draw 3 over random draw 1, leg1 N1000 1-macro_auc: paired geometric mean [95% interval]
\newcommand{\PairedFtJciiEThreeRandomDrawThreeOverRandomDrawTwo}{0.95~[0.94, 0.97]}  % experiments/FIGS/data/paired_v1err/finetune/ratios.json :: ratios[14].ratio, ratios[14].ci95 :: sha256 3e5cf25e859b68c2 :: random draw 3 over random draw 2, leg1 N1000 1-macro_auc: paired geometric mean [95% interval]
\newcommand{\PairedFtJciiEFourOnesixtwoOverOneeighteight}{0.99~[0.97, 1.02]}  % experiments/FIGS/data/paired_v1err/finetune/ratios.json :: ratios[15].ratio, ratios[15].ci95 :: sha256 3e5cf25e859b68c2 :: 162 over 188, leg1 N10000 1-macro_auc: paired geometric mean [95% interval]
\newcommand{\PairedFtJciiEFourFourthreeOverOneeighteight}{1.20~[1.17, 1.23]}  % experiments/FIGS/data/paired_v1err/finetune/ratios.json :: ratios[16].ratio, ratios[16].ci95 :: sha256 3e5cf25e859b68c2 :: 43 over 188, leg1 N10000 1-macro_auc: paired geometric mean [95% interval]
\newcommand{\PairedFtJciiEFourOnesevenOverOneeighteight}{1.32~[1.29, 1.34]}  % experiments/FIGS/data/paired_v1err/finetune/ratios.json :: ratios[17].ratio, ratios[17].ci95 :: sha256 3e5cf25e859b68c2 :: 17 over 188, leg1 N10000 1-macro_auc: paired geometric mean [95% interval]
\newcommand{\PairedFtJciiEFourFourthreeOverOnesixtwo}{1.21~[1.19, 1.23]}  % experiments/FIGS/data/paired_v1err/finetune/ratios.json :: ratios[18].ratio, ratios[18].ci95 :: sha256 3e5cf25e859b68c2 :: 43 over 162, leg1 N10000 1-macro_auc: paired geometric mean [95% interval]
\newcommand{\PairedFtJciiEFourOnesevenOverOnesixtwo}{1.32~[1.29, 1.36]}  % experiments/FIGS/data/paired_v1err/finetune/ratios.json :: ratios[19].ratio, ratios[19].ci95 :: sha256 3e5cf25e859b68c2 :: 17 over 162, leg1 N10000 1-macro_auc: paired geometric mean [95% interval]
\newcommand{\PairedFtJciiEFourOnesevenOverFourthree}{1.10~[1.07, 1.13]}  % experiments/FIGS/data/paired_v1err/finetune/ratios.json :: ratios[20].ratio, ratios[20].ci95 :: sha256 3e5cf25e859b68c2 :: 17 over 43, leg1 N10000 1-macro_auc: paired geometric mean [95% interval]
\newcommand{\PairedFtJciiEFourOnesixtwoMassOverOnesixtwo}{1.00~[0.94, 1.06]}  % experiments/FIGS/data/paired_v1err/finetune/ratios.json :: ratios[21].ratio, ratios[21].ci95 :: sha256 3e5cf25e859b68c2 :: 162+mass over 162, leg1 N10000 1-macro_auc: paired geometric mean [95% interval]
\newcommand{\PairedFtJciiEFourOnesevenMassOverOneseven}{0.990~[0.976, 1.004]}  % experiments/FIGS/data/paired_v1err/finetune/ratios.json :: ratios[22].ratio, ratios[22].ci95 :: sha256 3e5cf25e859b68c2 :: 17+mass over 17, leg1 N10000 1-macro_auc: paired geometric mean [95% interval]
\newcommand{\PairedFtJciiEFourRandomOverOneeighteight}{1.15~[1.09, 1.22]}  % experiments/FIGS/data/paired_v1err/finetune/ratios.json :: ratios[23].ratio, ratios[23].ci95 :: sha256 3e5cf25e859b68c2 :: random over 188, leg1 N10000 1-macro_auc: paired geometric mean [95% interval]
\newcommand{\PairedFtJciiEFourRandomOverOnesixtwo}{1.18~[1.13, 1.22]}  % experiments/FIGS/data/paired_v1err/finetune/ratios.json :: ratios[24].ratio, ratios[24].ci95 :: sha256 3e5cf25e859b68c2 :: random over 162, leg1 N10000 1-macro_auc: paired geometric mean [95% interval]
\newcommand{\PairedFtJciiEFourRandomOverFourthree}{0.97~[0.96, 0.99]}  % experiments/FIGS/data/paired_v1err/finetune/ratios.json :: ratios[25].ratio, ratios[25].ci95 :: sha256 3e5cf25e859b68c2 :: random over 43, leg1 N10000 1-macro_auc: paired geometric mean [95% interval]
\newcommand{\PairedFtJciiEFourRandomOverOneseven}{0.88~[0.82, 0.95]}  % experiments/FIGS/data/paired_v1err/finetune/ratios.json :: ratios[26].ratio, ratios[26].ci95 :: sha256 3e5cf25e859b68c2 :: random over 17, leg1 N10000 1-macro_auc: paired geometric mean [95% interval]
\newcommand{\PairedFtJciiEFourRandomDrawTwoOverRandomDrawOne}{1.00~[0.98, 1.02]}  % experiments/FIGS/data/paired_v1err/finetune/ratios.json :: ratios[27].ratio, ratios[27].ci95 :: sha256 3e5cf25e859b68c2 :: random draw 2 over random draw 1, leg1 N10000 1-macro_auc: paired geometric mean [95% interval]
\newcommand{\PairedFtJciiEFourRandomDrawThreeOverRandomDrawOne}{0.97~[0.95, 0.98]}  % experiments/FIGS/data/paired_v1err/finetune/ratios.json :: ratios[28].ratio, ratios[28].ci95 :: sha256 3e5cf25e859b68c2 :: random draw 3 over random draw 1, leg1 N10000 1-macro_auc: paired geometric mean [95% interval]
\newcommand{\PairedFtJciiEFourRandomDrawThreeOverRandomDrawTwo}{0.97~[0.95, 0.98]}  % experiments/FIGS/data/paired_v1err/finetune/ratios.json :: ratios[29].ratio, ratios[29].ci95 :: sha256 3e5cf25e859b68c2 :: random draw 3 over random draw 2, leg1 N10000 1-macro_auc: paired geometric mean [95% interval]
\newcommand{\PairedFtJciiEFiveOnesixtwoOverOneeighteight}{1.00~[0.98, 1.01]}  % experiments/FIGS/data/paired_v1err/finetune/ratios.json :: ratios[30].ratio, ratios[30].ci95 :: sha256 3e5cf25e859b68c2 :: 162 over 188, leg1 N100000 1-macro_auc: paired geometric mean [95% interval]
\newcommand{\PairedFtJciiEFiveFourthreeOverOneeighteight}{1.09~[1.06, 1.11]}  % experiments/FIGS/data/paired_v1err/finetune/ratios.json :: ratios[31].ratio, ratios[31].ci95 :: sha256 3e5cf25e859b68c2 :: 43 over 188, leg1 N100000 1-macro_auc: paired geometric mean [95% interval]
\newcommand{\PairedFtJciiEFiveOnesevenOverOneeighteight}{1.22~[1.20, 1.23]}  % experiments/FIGS/data/paired_v1err/finetune/ratios.json :: ratios[32].ratio, ratios[32].ci95 :: sha256 3e5cf25e859b68c2 :: 17 over 188, leg1 N100000 1-macro_auc: paired geometric mean [95% interval]
\newcommand{\PairedFtJciiEFiveFourthreeOverOnesixtwo}{1.09~[1.07, 1.11]}  % experiments/FIGS/data/paired_v1err/finetune/ratios.json :: ratios[33].ratio, ratios[33].ci95 :: sha256 3e5cf25e859b68c2 :: 43 over 162, leg1 N100000 1-macro_auc: paired geometric mean [95% interval]
\newcommand{\PairedFtJciiEFiveOnesevenOverOnesixtwo}{1.22~[1.20, 1.24]}  % experiments/FIGS/data/paired_v1err/finetune/ratios.json :: ratios[34].ratio, ratios[34].ci95 :: sha256 3e5cf25e859b68c2 :: 17 over 162, leg1 N100000 1-macro_auc: paired geometric mean [95% interval]
\newcommand{\PairedFtJciiEFiveOnesevenOverFourthree}{1.12~[1.10, 1.14]}  % experiments/FIGS/data/paired_v1err/finetune/ratios.json :: ratios[35].ratio, ratios[35].ci95 :: sha256 3e5cf25e859b68c2 :: 17 over 43, leg1 N100000 1-macro_auc: paired geometric mean [95% interval]
\newcommand{\PairedFtJciiEFiveOnesixtwoMassOverOnesixtwo}{1.018~[1.002, 1.035]}  % experiments/FIGS/data/paired_v1err/finetune/ratios.json :: ratios[36].ratio, ratios[36].ci95 :: sha256 3e5cf25e859b68c2 :: 162+mass over 162, leg1 N100000 1-macro_auc: paired geometric mean [95% interval]
\newcommand{\PairedFtJciiEFiveOnesevenMassOverOneseven}{1.03~[1.01, 1.06]}  % experiments/FIGS/data/paired_v1err/finetune/ratios.json :: ratios[37].ratio, ratios[37].ci95 :: sha256 3e5cf25e859b68c2 :: 17+mass over 17, leg1 N100000 1-macro_auc: paired geometric mean [95% interval]
\newcommand{\PairedFtJciiEFiveRandomOverOneeighteight}{1.15~[1.09, 1.20]}  % experiments/FIGS/data/paired_v1err/finetune/ratios.json :: ratios[38].ratio, ratios[38].ci95 :: sha256 3e5cf25e859b68c2 :: random over 188, leg1 N100000 1-macro_auc: paired geometric mean [95% interval]
\newcommand{\PairedFtJciiEFiveRandomOverOnesixtwo}{1.15~[1.09, 1.20]}  % experiments/FIGS/data/paired_v1err/finetune/ratios.json :: ratios[39].ratio, ratios[39].ci95 :: sha256 3e5cf25e859b68c2 :: random over 162, leg1 N100000 1-macro_auc: paired geometric mean [95% interval]
\newcommand{\PairedFtJciiEFiveRandomOverFourthree}{1.05~[1.04, 1.06]}  % experiments/FIGS/data/paired_v1err/finetune/ratios.json :: ratios[40].ratio, ratios[40].ci95 :: sha256 3e5cf25e859b68c2 :: random over 43, leg1 N100000 1-macro_auc: paired geometric mean [95% interval]
\newcommand{\PairedFtJciiEFiveRandomOverOneseven}{0.95~[0.93, 0.96]}  % experiments/FIGS/data/paired_v1err/finetune/ratios.json :: ratios[41].ratio, ratios[41].ci95 :: sha256 3e5cf25e859b68c2 :: random over 17, leg1 N100000 1-macro_auc: paired geometric mean [95% interval]
\newcommand{\PairedFtJciiEFiveRandomDrawTwoOverRandomDrawOne}{0.98~[0.97, 0.99]}  % experiments/FIGS/data/paired_v1err/finetune/ratios.json :: ratios[42].ratio, ratios[42].ci95 :: sha256 3e5cf25e859b68c2 :: random draw 2 over random draw 1, leg1 N100000 1-macro_auc: paired geometric mean [95% interval]
\newcommand{\PairedFtJciiEFiveRandomDrawThreeOverRandomDrawOne}{1.00~[0.99, 1.02]}  % experiments/FIGS/data/paired_v1err/finetune/ratios.json :: ratios[43].ratio, ratios[43].ci95 :: sha256 3e5cf25e859b68c2 :: random draw 3 over random draw 1, leg1 N100000 1-macro_auc: paired geometric mean [95% interval]
\newcommand{\PairedFtJciiEFiveRandomDrawThreeOverRandomDrawTwo}{1.03~[1.01, 1.04]}  % experiments/FIGS/data/paired_v1err/finetune/ratios.json :: ratios[44].ratio, ratios[44].ci95 :: sha256 3e5cf25e859b68c2 :: random draw 3 over random draw 2, leg1 N100000 1-macro_auc: paired geometric mean [95% interval]
\newcommand{\PairedFtJciiESixOnesixtwoOverOneeighteight}{1.002~[0.998, 1.005]}  % experiments/FIGS/data/paired_v1err/finetune/ratios.json :: ratios[45].ratio, ratios[45].ci95 :: sha256 3e5cf25e859b68c2 :: 162 over 188, leg1 N1000000 1-macro_auc: paired geometric mean [95% interval]
\newcommand{\PairedFtJciiESixFourthreeOverOneeighteight}{1.045~[1.037, 1.053]}  % experiments/FIGS/data/paired_v1err/finetune/ratios.json :: ratios[46].ratio, ratios[46].ci95 :: sha256 3e5cf25e859b68c2 :: 43 over 188, leg1 N1000000 1-macro_auc: paired geometric mean [95% interval]
\newcommand{\PairedFtJciiESixOnesevenOverOneeighteight}{1.14~[1.13, 1.16]}  % experiments/FIGS/data/paired_v1err/finetune/ratios.json :: ratios[47].ratio, ratios[47].ci95 :: sha256 3e5cf25e859b68c2 :: 17 over 188, leg1 N1000000 1-macro_auc: paired geometric mean [95% interval]
\newcommand{\PairedFtJciiESixFourthreeOverOnesixtwo}{1.044~[1.036, 1.051]}  % experiments/FIGS/data/paired_v1err/finetune/ratios.json :: ratios[48].ratio, ratios[48].ci95 :: sha256 3e5cf25e859b68c2 :: 43 over 162, leg1 N1000000 1-macro_auc: paired geometric mean [95% interval]
\newcommand{\PairedFtJciiESixOnesevenOverOnesixtwo}{1.14~[1.13, 1.15]}  % experiments/FIGS/data/paired_v1err/finetune/ratios.json :: ratios[49].ratio, ratios[49].ci95 :: sha256 3e5cf25e859b68c2 :: 17 over 162, leg1 N1000000 1-macro_auc: paired geometric mean [95% interval]
\newcommand{\PairedFtJciiESixOnesevenOverFourthree}{1.09~[1.08, 1.11]}  % experiments/FIGS/data/paired_v1err/finetune/ratios.json :: ratios[50].ratio, ratios[50].ci95 :: sha256 3e5cf25e859b68c2 :: 17 over 43, leg1 N1000000 1-macro_auc: paired geometric mean [95% interval]
\newcommand{\PairedFtJciiESixOnesixtwoMassOverOnesixtwo}{1.019~[1.014, 1.023]}  % experiments/FIGS/data/paired_v1err/finetune/ratios.json :: ratios[51].ratio, ratios[51].ci95 :: sha256 3e5cf25e859b68c2 :: 162+mass over 162, leg1 N1000000 1-macro_auc: paired geometric mean [95% interval]
\newcommand{\PairedFtJciiESixOnesevenMassOverOneseven}{1.03~[1.02, 1.04]}  % experiments/FIGS/data/paired_v1err/finetune/ratios.json :: ratios[52].ratio, ratios[52].ci95 :: sha256 3e5cf25e859b68c2 :: 17+mass over 17, leg1 N1000000 1-macro_auc: paired geometric mean [95% interval]
\newcommand{\PairedFtJciiESixRandomOverOneeighteight}{1.073~[1.067, 1.080]}  % experiments/FIGS/data/paired_v1err/finetune/ratios.json :: ratios[53].ratio, ratios[53].ci95 :: sha256 3e5cf25e859b68c2 :: random over 188, leg1 N1000000 1-macro_auc: paired geometric mean [95% interval]
\newcommand{\PairedFtJciiESixRandomOverOnesixtwo}{1.072~[1.065, 1.078]}  % experiments/FIGS/data/paired_v1err/finetune/ratios.json :: ratios[54].ratio, ratios[54].ci95 :: sha256 3e5cf25e859b68c2 :: random over 162, leg1 N1000000 1-macro_auc: paired geometric mean [95% interval]
\newcommand{\PairedFtJciiESixRandomOverFourthree}{1.03~[1.02, 1.04]}  % experiments/FIGS/data/paired_v1err/finetune/ratios.json :: ratios[55].ratio, ratios[55].ci95 :: sha256 3e5cf25e859b68c2 :: random over 43, leg1 N1000000 1-macro_auc: paired geometric mean [95% interval]
\newcommand{\PairedFtJciiESixRandomOverOneseven}{0.94~[0.93, 0.95]}  % experiments/FIGS/data/paired_v1err/finetune/ratios.json :: ratios[56].ratio, ratios[56].ci95 :: sha256 3e5cf25e859b68c2 :: random over 17, leg1 N1000000 1-macro_auc: paired geometric mean [95% interval]
\newcommand{\PairedFtJciiESixRandomDrawTwoOverRandomDrawOne}{1.00~[0.99, 1.01]}  % experiments/FIGS/data/paired_v1err/finetune/ratios.json :: ratios[57].ratio, ratios[57].ci95 :: sha256 3e5cf25e859b68c2 :: random draw 2 over random draw 1, leg1 N1000000 1-macro_auc: paired geometric mean [95% interval]
\newcommand{\PairedFtJciiESixRandomDrawThreeOverRandomDrawOne}{0.996~[0.989, 1.003]}  % experiments/FIGS/data/paired_v1err/finetune/ratios.json :: ratios[58].ratio, ratios[58].ci95 :: sha256 3e5cf25e859b68c2 :: random draw 3 over random draw 1, leg1 N1000000 1-macro_auc: paired geometric mean [95% interval]
\newcommand{\PairedFtJciiESixRandomDrawThreeOverRandomDrawTwo}{0.997~[0.989, 1.004]}  % experiments/FIGS/data/paired_v1err/finetune/ratios.json :: ratios[59].ratio, ratios[59].ci95 :: sha256 3e5cf25e859b68c2 :: random draw 3 over random draw 2, leg1 N1000000 1-macro_auc: paired geometric mean [95% interval]
\newcommand{\PairedFtJcEThreeOnesixtwoOverOneeighteight}{0.96~[0.83, 1.11]}  % experiments/FIGS/data/paired_v1err/finetune/ratios.json :: ratios[60].ratio, ratios[60].ci95 :: sha256 3e5cf25e859b68c2 :: 162 over 188, leg2 N1000 1-macro_auc: paired geometric mean [95% interval]
\newcommand{\PairedFtJcEThreeFourthreeOverOneeighteight}{1.01~[0.84, 1.21]}  % experiments/FIGS/data/paired_v1err/finetune/ratios.json :: ratios[61].ratio, ratios[61].ci95 :: sha256 3e5cf25e859b68c2 :: 43 over 188, leg2 N1000 1-macro_auc: paired geometric mean [95% interval]
\newcommand{\PairedFtJcEThreeOnesevenOverOneeighteight}{1.07~[0.93, 1.23]}  % experiments/FIGS/data/paired_v1err/finetune/ratios.json :: ratios[62].ratio, ratios[62].ci95 :: sha256 3e5cf25e859b68c2 :: 17 over 188, leg2 N1000 1-macro_auc: paired geometric mean [95% interval]
\newcommand{\PairedFtJcEThreeFourthreeOverOnesixtwo}{1.05~[0.96, 1.15]}  % experiments/FIGS/data/paired_v1err/finetune/ratios.json :: ratios[63].ratio, ratios[63].ci95 :: sha256 3e5cf25e859b68c2 :: 43 over 162, leg2 N1000 1-macro_auc: paired geometric mean [95% interval]
\newcommand{\PairedFtJcEThreeOnesevenOverOnesixtwo}{1.12~[1.07, 1.17]}  % experiments/FIGS/data/paired_v1err/finetune/ratios.json :: ratios[64].ratio, ratios[64].ci95 :: sha256 3e5cf25e859b68c2 :: 17 over 162, leg2 N1000 1-macro_auc: paired geometric mean [95% interval]
\newcommand{\PairedFtJcEThreeOnesevenOverFourthree}{1.06~[0.98, 1.15]}  % experiments/FIGS/data/paired_v1err/finetune/ratios.json :: ratios[65].ratio, ratios[65].ci95 :: sha256 3e5cf25e859b68c2 :: 17 over 43, leg2 N1000 1-macro_auc: paired geometric mean [95% interval]
\newcommand{\PairedFtJcEThreeOnesixtwoMassOverOnesixtwo}{0.91~[0.88, 0.94]}  % experiments/FIGS/data/paired_v1err/finetune/ratios.json :: ratios[66].ratio, ratios[66].ci95 :: sha256 3e5cf25e859b68c2 :: 162+mass over 162, leg2 N1000 1-macro_auc: paired geometric mean [95% interval]
\newcommand{\PairedFtJcEThreeOnesevenMassOverOneseven}{1.02~[0.92, 1.13]}  % experiments/FIGS/data/paired_v1err/finetune/ratios.json :: ratios[67].ratio, ratios[67].ci95 :: sha256 3e5cf25e859b68c2 :: 17+mass over 17, leg2 N1000 1-macro_auc: paired geometric mean [95% interval]
\newcommand{\PairedFtJcEThreeRandomOverOneeighteight}{1.01~[0.78, 1.32]}  % experiments/FIGS/data/paired_v1err/finetune/ratios.json :: ratios[68].ratio, ratios[68].ci95 :: sha256 3e5cf25e859b68c2 :: random over 188, leg2 N1000 1-macro_auc: paired geometric mean [95% interval]
\newcommand{\PairedFtJcEThreeRandomOverOnesixtwo}{1.12~[1.04, 1.20]}  % experiments/FIGS/data/paired_v1err/finetune/ratios.json :: ratios[69].ratio, ratios[69].ci95 :: sha256 3e5cf25e859b68c2 :: random over 162, leg2 N1000 1-macro_auc: paired geometric mean [95% interval]
\newcommand{\PairedFtJcEThreeRandomOverFourthree}{1.06~[0.79, 1.41]}  % experiments/FIGS/data/paired_v1err/finetune/ratios.json :: ratios[70].ratio, ratios[70].ci95 :: sha256 3e5cf25e859b68c2 :: random over 43, leg2 N1000 1-macro_auc: paired geometric mean [95% interval]
\newcommand{\PairedFtJcEThreeRandomOverOneseven}{0.98~[0.84, 1.13]}  % experiments/FIGS/data/paired_v1err/finetune/ratios.json :: ratios[71].ratio, ratios[71].ci95 :: sha256 3e5cf25e859b68c2 :: random over 17, leg2 N1000 1-macro_auc: paired geometric mean [95% interval]
\newcommand{\PairedFtJcEThreeRandomDrawTwoOverRandomDrawOne}{0.954~[0.949, 0.958]}  % experiments/FIGS/data/paired_v1err/finetune/ratios.json :: ratios[72].ratio, ratios[72].ci95 :: sha256 3e5cf25e859b68c2 :: random draw 2 over random draw 1, leg2 N1000 1-macro_auc: paired geometric mean [95% interval]
\newcommand{\PairedFtJcEThreeRandomDrawThreeOverRandomDrawOne}{1.013~[1.009, 1.017]}  % experiments/FIGS/data/paired_v1err/finetune/ratios.json :: ratios[73].ratio, ratios[73].ci95 :: sha256 3e5cf25e859b68c2 :: random draw 3 over random draw 1, leg2 N1000 1-macro_auc: paired geometric mean [95% interval]
\newcommand{\PairedFtJcEThreeRandomDrawThreeOverRandomDrawTwo}{1.062~[1.057, 1.067]}  % experiments/FIGS/data/paired_v1err/finetune/ratios.json :: ratios[74].ratio, ratios[74].ci95 :: sha256 3e5cf25e859b68c2 :: random draw 3 over random draw 2, leg2 N1000 1-macro_auc: paired geometric mean [95% interval]
\newcommand{\PairedFtJcEFourOnesixtwoOverOneeighteight}{0.98~[0.93, 1.02]}  % experiments/FIGS/data/paired_v1err/finetune/ratios.json :: ratios[75].ratio, ratios[75].ci95 :: sha256 3e5cf25e859b68c2 :: 162 over 188, leg2 N10000 1-macro_auc: paired geometric mean [95% interval]
\newcommand{\PairedFtJcEFourFourthreeOverOneeighteight}{0.98~[0.93, 1.04]}  % experiments/FIGS/data/paired_v1err/finetune/ratios.json :: ratios[76].ratio, ratios[76].ci95 :: sha256 3e5cf25e859b68c2 :: 43 over 188, leg2 N10000 1-macro_auc: paired geometric mean [95% interval]
\newcommand{\PairedFtJcEFourOnesevenOverOneeighteight}{1.07~[1.02, 1.11]}  % experiments/FIGS/data/paired_v1err/finetune/ratios.json :: ratios[77].ratio, ratios[77].ci95 :: sha256 3e5cf25e859b68c2 :: 17 over 188, leg2 N10000 1-macro_auc: paired geometric mean [95% interval]
\newcommand{\PairedFtJcEFourFourthreeOverOnesixtwo}{1.01~[0.97, 1.04]}  % experiments/FIGS/data/paired_v1err/finetune/ratios.json :: ratios[78].ratio, ratios[78].ci95 :: sha256 3e5cf25e859b68c2 :: 43 over 162, leg2 N10000 1-macro_auc: paired geometric mean [95% interval]
\newcommand{\PairedFtJcEFourOnesevenOverOnesixtwo}{1.09~[1.06, 1.12]}  % experiments/FIGS/data/paired_v1err/finetune/ratios.json :: ratios[79].ratio, ratios[79].ci95 :: sha256 3e5cf25e859b68c2 :: 17 over 162, leg2 N10000 1-macro_auc: paired geometric mean [95% interval]
\newcommand{\PairedFtJcEFourOnesevenOverFourthree}{1.08~[1.07, 1.10]}  % experiments/FIGS/data/paired_v1err/finetune/ratios.json :: ratios[80].ratio, ratios[80].ci95 :: sha256 3e5cf25e859b68c2 :: 17 over 43, leg2 N10000 1-macro_auc: paired geometric mean [95% interval]
\newcommand{\PairedFtJcEFourOnesixtwoMassOverOnesixtwo}{1.01~[0.98, 1.04]}  % experiments/FIGS/data/paired_v1err/finetune/ratios.json :: ratios[81].ratio, ratios[81].ci95 :: sha256 3e5cf25e859b68c2 :: 162+mass over 162, leg2 N10000 1-macro_auc: paired geometric mean [95% interval]
\newcommand{\PairedFtJcEFourOnesevenMassOverOneseven}{1.05~[1.02, 1.08]}  % experiments/FIGS/data/paired_v1err/finetune/ratios.json :: ratios[82].ratio, ratios[82].ci95 :: sha256 3e5cf25e859b68c2 :: 17+mass over 17, leg2 N10000 1-macro_auc: paired geometric mean [95% interval]
\newcommand{\PairedFtJcEFourRandomOverOneeighteight}{1.08~[1.03, 1.14]}  % experiments/FIGS/data/paired_v1err/finetune/ratios.json :: ratios[83].ratio, ratios[83].ci95 :: sha256 3e5cf25e859b68c2 :: random over 188, leg2 N10000 1-macro_auc: paired geometric mean [95% interval]
\newcommand{\PairedFtJcEFourRandomOverOnesixtwo}{1.11~[0.99, 1.23]}  % experiments/FIGS/data/paired_v1err/finetune/ratios.json :: ratios[84].ratio, ratios[84].ci95 :: sha256 3e5cf25e859b68c2 :: random over 162, leg2 N10000 1-macro_auc: paired geometric mean [95% interval]
\newcommand{\PairedFtJcEFourRandomOverFourthree}{1.10~[0.99, 1.21]}  % experiments/FIGS/data/paired_v1err/finetune/ratios.json :: ratios[85].ratio, ratios[85].ci95 :: sha256 3e5cf25e859b68c2 :: random over 43, leg2 N10000 1-macro_auc: paired geometric mean [95% interval]
\newcommand{\PairedFtJcEFourRandomOverOneseven}{1.02~[0.94, 1.10]}  % experiments/FIGS/data/paired_v1err/finetune/ratios.json :: ratios[86].ratio, ratios[86].ci95 :: sha256 3e5cf25e859b68c2 :: random over 17, leg2 N10000 1-macro_auc: paired geometric mean [95% interval]
\newcommand{\PairedFtJcEFourRandomDrawTwoOverRandomDrawOne}{0.960~[0.955, 0.965]}  % experiments/FIGS/data/paired_v1err/finetune/ratios.json :: ratios[87].ratio, ratios[87].ci95 :: sha256 3e5cf25e859b68c2 :: random draw 2 over random draw 1, leg2 N10000 1-macro_auc: paired geometric mean [95% interval]
\newcommand{\PairedFtJcEFourRandomDrawThreeOverRandomDrawOne}{1.008~[1.004, 1.012]}  % experiments/FIGS/data/paired_v1err/finetune/ratios.json :: ratios[88].ratio, ratios[88].ci95 :: sha256 3e5cf25e859b68c2 :: random draw 3 over random draw 1, leg2 N10000 1-macro_auc: paired geometric mean [95% interval]
\newcommand{\PairedFtJcEFourRandomDrawThreeOverRandomDrawTwo}{1.050~[1.045, 1.055]}  % experiments/FIGS/data/paired_v1err/finetune/ratios.json :: ratios[89].ratio, ratios[89].ci95 :: sha256 3e5cf25e859b68c2 :: random draw 3 over random draw 2, leg2 N10000 1-macro_auc: paired geometric mean [95% interval]
\newcommand{\PairedFtJcEFiveOnesixtwoOverOneeighteight}{0.983~[0.968, 0.998]}  % experiments/FIGS/data/paired_v1err/finetune/ratios.json :: ratios[90].ratio, ratios[90].ci95 :: sha256 3e5cf25e859b68c2 :: 162 over 188, leg2 N100000 1-macro_auc: paired geometric mean [95% interval]
\newcommand{\PairedFtJcEFiveFourthreeOverOneeighteight}{1.015~[1.002, 1.029]}  % experiments/FIGS/data/paired_v1err/finetune/ratios.json :: ratios[91].ratio, ratios[91].ci95 :: sha256 3e5cf25e859b68c2 :: 43 over 188, leg2 N100000 1-macro_auc: paired geometric mean [95% interval]
\newcommand{\PairedFtJcEFiveOnesevenOverOneeighteight}{1.07~[1.06, 1.09]}  % experiments/FIGS/data/paired_v1err/finetune/ratios.json :: ratios[92].ratio, ratios[92].ci95 :: sha256 3e5cf25e859b68c2 :: 17 over 188, leg2 N100000 1-macro_auc: paired geometric mean [95% interval]
\newcommand{\PairedFtJcEFiveFourthreeOverOnesixtwo}{1.03~[1.02, 1.05]}  % experiments/FIGS/data/paired_v1err/finetune/ratios.json :: ratios[93].ratio, ratios[93].ci95 :: sha256 3e5cf25e859b68c2 :: 43 over 162, leg2 N100000 1-macro_auc: paired geometric mean [95% interval]
\newcommand{\PairedFtJcEFiveOnesevenOverOnesixtwo}{1.091~[1.088, 1.094]}  % experiments/FIGS/data/paired_v1err/finetune/ratios.json :: ratios[94].ratio, ratios[94].ci95 :: sha256 3e5cf25e859b68c2 :: 17 over 162, leg2 N100000 1-macro_auc: paired geometric mean [95% interval]
\newcommand{\PairedFtJcEFiveOnesevenOverFourthree}{1.06~[1.04, 1.07]}  % experiments/FIGS/data/paired_v1err/finetune/ratios.json :: ratios[95].ratio, ratios[95].ci95 :: sha256 3e5cf25e859b68c2 :: 17 over 43, leg2 N100000 1-macro_auc: paired geometric mean [95% interval]
\newcommand{\PairedFtJcEFiveOnesixtwoMassOverOnesixtwo}{1.016~[1.008, 1.024]}  % experiments/FIGS/data/paired_v1err/finetune/ratios.json :: ratios[96].ratio, ratios[96].ci95 :: sha256 3e5cf25e859b68c2 :: 162+mass over 162, leg2 N100000 1-macro_auc: paired geometric mean [95% interval]
\newcommand{\PairedFtJcEFiveOnesevenMassOverOneseven}{1.02~[1.01, 1.03]}  % experiments/FIGS/data/paired_v1err/finetune/ratios.json :: ratios[97].ratio, ratios[97].ci95 :: sha256 3e5cf25e859b68c2 :: 17+mass over 17, leg2 N100000 1-macro_auc: paired geometric mean [95% interval]
\newcommand{\PairedFtJcEFiveRandomOverOneeighteight}{1.044~[1.039, 1.049]}  % experiments/FIGS/data/paired_v1err/finetune/ratios.json :: ratios[98].ratio, ratios[98].ci95 :: sha256 3e5cf25e859b68c2 :: random over 188, leg2 N100000 1-macro_auc: paired geometric mean [95% interval]
\newcommand{\PairedFtJcEFiveRandomOverOnesixtwo}{1.06~[1.04, 1.09]}  % experiments/FIGS/data/paired_v1err/finetune/ratios.json :: ratios[99].ratio, ratios[99].ci95 :: sha256 3e5cf25e859b68c2 :: random over 162, leg2 N100000 1-macro_auc: paired geometric mean [95% interval]
\newcommand{\PairedFtJcEFiveRandomOverFourthree}{1.03~[1.01, 1.04]}  % experiments/FIGS/data/paired_v1err/finetune/ratios.json :: ratios[100].ratio, ratios[100].ci95 :: sha256 3e5cf25e859b68c2 :: random over 43, leg2 N100000 1-macro_auc: paired geometric mean [95% interval]
\newcommand{\PairedFtJcEFiveRandomOverOneseven}{0.97~[0.95, 0.99]}  % experiments/FIGS/data/paired_v1err/finetune/ratios.json :: ratios[101].ratio, ratios[101].ci95 :: sha256 3e5cf25e859b68c2 :: random over 17, leg2 N100000 1-macro_auc: paired geometric mean [95% interval]
\newcommand{\PairedFtJcEFiveRandomDrawTwoOverRandomDrawOne}{0.992~[0.988, 0.997]}  % experiments/FIGS/data/paired_v1err/finetune/ratios.json :: ratios[102].ratio, ratios[102].ci95 :: sha256 3e5cf25e859b68c2 :: random draw 2 over random draw 1, leg2 N100000 1-macro_auc: paired geometric mean [95% interval]
\newcommand{\PairedFtJcEFiveRandomDrawThreeOverRandomDrawOne}{1.007~[1.003, 1.011]}  % experiments/FIGS/data/paired_v1err/finetune/ratios.json :: ratios[103].ratio, ratios[103].ci95 :: sha256 3e5cf25e859b68c2 :: random draw 3 over random draw 1, leg2 N100000 1-macro_auc: paired geometric mean [95% interval]
\newcommand{\PairedFtJcEFiveRandomDrawThreeOverRandomDrawTwo}{1.015~[1.011, 1.020]}  % experiments/FIGS/data/paired_v1err/finetune/ratios.json :: ratios[104].ratio, ratios[104].ci95 :: sha256 3e5cf25e859b68c2 :: random draw 3 over random draw 2, leg2 N100000 1-macro_auc: paired geometric mean [95% interval]
\newcommand{\PairedFtJcESixOnesixtwoOverOneeighteight}{1.00~[0.99, 1.01]}  % experiments/FIGS/data/paired_v1err/finetune/ratios.json :: ratios[105].ratio, ratios[105].ci95 :: sha256 3e5cf25e859b68c2 :: 162 over 188, leg2 N1000000 1-macro_auc: paired geometric mean [95% interval]
\newcommand{\PairedFtJcESixFourthreeOverOneeighteight}{1.02~[1.01, 1.03]}  % experiments/FIGS/data/paired_v1err/finetune/ratios.json :: ratios[106].ratio, ratios[106].ci95 :: sha256 3e5cf25e859b68c2 :: 43 over 188, leg2 N1000000 1-macro_auc: paired geometric mean [95% interval]
\newcommand{\PairedFtJcESixOnesevenOverOneeighteight}{1.08~[1.07, 1.09]}  % experiments/FIGS/data/paired_v1err/finetune/ratios.json :: ratios[107].ratio, ratios[107].ci95 :: sha256 3e5cf25e859b68c2 :: 17 over 188, leg2 N1000000 1-macro_auc: paired geometric mean [95% interval]
\newcommand{\PairedFtJcESixFourthreeOverOnesixtwo}{1.02~[1.01, 1.03]}  % experiments/FIGS/data/paired_v1err/finetune/ratios.json :: ratios[108].ratio, ratios[108].ci95 :: sha256 3e5cf25e859b68c2 :: 43 over 162, leg2 N1000000 1-macro_auc: paired geometric mean [95% interval]
\newcommand{\PairedFtJcESixOnesevenOverOnesixtwo}{1.08~[1.07, 1.09]}  % experiments/FIGS/data/paired_v1err/finetune/ratios.json :: ratios[109].ratio, ratios[109].ci95 :: sha256 3e5cf25e859b68c2 :: 17 over 162, leg2 N1000000 1-macro_auc: paired geometric mean [95% interval]
\newcommand{\PairedFtJcESixOnesevenOverFourthree}{1.06~[1.05, 1.07]}  % experiments/FIGS/data/paired_v1err/finetune/ratios.json :: ratios[110].ratio, ratios[110].ci95 :: sha256 3e5cf25e859b68c2 :: 17 over 43, leg2 N1000000 1-macro_auc: paired geometric mean [95% interval]
\newcommand{\PairedFtJcESixOnesixtwoMassOverOnesixtwo}{1.012~[1.004, 1.020]}  % experiments/FIGS/data/paired_v1err/finetune/ratios.json :: ratios[111].ratio, ratios[111].ci95 :: sha256 3e5cf25e859b68c2 :: 162+mass over 162, leg2 N1000000 1-macro_auc: paired geometric mean [95% interval]
\newcommand{\PairedFtJcESixOnesevenMassOverOneseven}{1.02~[1.01, 1.04]}  % experiments/FIGS/data/paired_v1err/finetune/ratios.json :: ratios[112].ratio, ratios[112].ci95 :: sha256 3e5cf25e859b68c2 :: 17+mass over 17, leg2 N1000000 1-macro_auc: paired geometric mean [95% interval]
\newcommand{\PairedFtJcESixRandomOverOneeighteight}{1.046~[1.038, 1.053]}  % experiments/FIGS/data/paired_v1err/finetune/ratios.json :: ratios[113].ratio, ratios[113].ci95 :: sha256 3e5cf25e859b68c2 :: random over 188, leg2 N1000000 1-macro_auc: paired geometric mean [95% interval]
\newcommand{\PairedFtJcESixRandomOverOnesixtwo}{1.047~[1.038, 1.055]}  % experiments/FIGS/data/paired_v1err/finetune/ratios.json :: ratios[114].ratio, ratios[114].ci95 :: sha256 3e5cf25e859b68c2 :: random over 162, leg2 N1000000 1-macro_auc: paired geometric mean [95% interval]
\newcommand{\PairedFtJcESixRandomOverFourthree}{1.03~[1.02, 1.04]}  % experiments/FIGS/data/paired_v1err/finetune/ratios.json :: ratios[115].ratio, ratios[115].ci95 :: sha256 3e5cf25e859b68c2 :: random over 43, leg2 N1000000 1-macro_auc: paired geometric mean [95% interval]
\newcommand{\PairedFtJcESixRandomOverOneseven}{0.97~[0.96, 0.98]}  % experiments/FIGS/data/paired_v1err/finetune/ratios.json :: ratios[116].ratio, ratios[116].ci95 :: sha256 3e5cf25e859b68c2 :: random over 17, leg2 N1000000 1-macro_auc: paired geometric mean [95% interval]
\newcommand{\PairedFtJcESixRandomDrawTwoOverRandomDrawOne}{1.001~[0.998, 1.004]}  % experiments/FIGS/data/paired_v1err/finetune/ratios.json :: ratios[117].ratio, ratios[117].ci95 :: sha256 3e5cf25e859b68c2 :: random draw 2 over random draw 1, leg2 N1000000 1-macro_auc: paired geometric mean [95% interval]
\newcommand{\PairedFtJcESixRandomDrawThreeOverRandomDrawOne}{1.005~[1.002, 1.008]}  % experiments/FIGS/data/paired_v1err/finetune/ratios.json :: ratios[118].ratio, ratios[118].ci95 :: sha256 3e5cf25e859b68c2 :: random draw 3 over random draw 1, leg2 N1000000 1-macro_auc: paired geometric mean [95% interval]
\newcommand{\PairedFtJcESixRandomDrawThreeOverRandomDrawTwo}{1.004~[1.001, 1.007]}  % experiments/FIGS/data/paired_v1err/finetune/ratios.json :: ratios[119].ratio, ratios[119].ci95 :: sha256 3e5cf25e859b68c2 :: random draw 3 over random draw 2, leg2 N1000000 1-macro_auc: paired geometric mean [95% interval]
\newcommand{\PairedMassResMlpOnesixtwoOverOneeighteight}{0.990~[0.977, 1.003]}  % experiments/FIGS/data/paired_v1err/mass/ratios.json :: ratios[0].ratio, ratios[0].ci95 :: sha256 d877a7f087260f64 :: 162 over 188, resolution mlp sigma_eff: paired geometric mean [95% interval]
\newcommand{\PairedMassResMlpFourthreeOverOneeighteight}{0.98~[0.96, 0.99]}  % experiments/FIGS/data/paired_v1err/mass/ratios.json :: ratios[1].ratio, ratios[1].ci95 :: sha256 d877a7f087260f64 :: 43 over 188, resolution mlp sigma_eff: paired geometric mean [95% interval]
\newcommand{\PairedMassResMlpOnesevenOverOneeighteight}{0.96~[0.94, 0.98]}  % experiments/FIGS/data/paired_v1err/mass/ratios.json :: ratios[2].ratio, ratios[2].ci95 :: sha256 d877a7f087260f64 :: 17 over 188, resolution mlp sigma_eff: paired geometric mean [95% interval]
\newcommand{\PairedMassResMlpFourthreeOverOnesixtwo}{0.989~[0.978, 0.999]}  % experiments/FIGS/data/paired_v1err/mass/ratios.json :: ratios[3].ratio, ratios[3].ci95 :: sha256 d877a7f087260f64 :: 43 over 162, resolution mlp sigma_eff: paired geometric mean [95% interval]
\newcommand{\PairedMassResMlpOnesevenOverOnesixtwo}{0.97~[0.94, 0.99]}  % experiments/FIGS/data/paired_v1err/mass/ratios.json :: ratios[4].ratio, ratios[4].ci95 :: sha256 d877a7f087260f64 :: 17 over 162, resolution mlp sigma_eff: paired geometric mean [95% interval]
\newcommand{\PairedMassResMlpOnesevenOverFourthree}{0.979~[0.961, 0.997]}  % experiments/FIGS/data/paired_v1err/mass/ratios.json :: ratios[5].ratio, ratios[5].ci95 :: sha256 d877a7f087260f64 :: 17 over 43, resolution mlp sigma_eff: paired geometric mean [95% interval]
\newcommand{\PairedMassResMlpOnesixtwoMassOverOnesixtwo}{0.80~[0.79, 0.81]}  % experiments/FIGS/data/paired_v1err/mass/ratios.json :: ratios[6].ratio, ratios[6].ci95 :: sha256 d877a7f087260f64 :: 162+mass over 162, resolution mlp sigma_eff: paired geometric mean [95% interval]
\newcommand{\PairedMassResMlpOnesevenMassOverOneseven}{0.83~[0.82, 0.84]}  % experiments/FIGS/data/paired_v1err/mass/ratios.json :: ratios[7].ratio, ratios[7].ci95 :: sha256 d877a7f087260f64 :: 17+mass over 17, resolution mlp sigma_eff: paired geometric mean [95% interval]
\newcommand{\PairedMassResRidgeOnesixtwoOverOneeighteight}{0.97~[0.96, 0.98]}  % experiments/FIGS/data/paired_v1err/mass/ratios.json :: ratios[8].ratio, ratios[8].ci95 :: sha256 d877a7f087260f64 :: 162 over 188, resolution ridge sigma_eff: paired geometric mean [95% interval]
\newcommand{\PairedMassResRidgeFourthreeOverOneeighteight}{0.95~[0.94, 0.96]}  % experiments/FIGS/data/paired_v1err/mass/ratios.json :: ratios[9].ratio, ratios[9].ci95 :: sha256 d877a7f087260f64 :: 43 over 188, resolution ridge sigma_eff: paired geometric mean [95% interval]
\newcommand{\PairedMassResRidgeOnesevenOverOneeighteight}{0.88~[0.85, 0.92]}  % experiments/FIGS/data/paired_v1err/mass/ratios.json :: ratios[10].ratio, ratios[10].ci95 :: sha256 d877a7f087260f64 :: 17 over 188, resolution ridge sigma_eff: paired geometric mean [95% interval]
\newcommand{\PairedMassResRidgeFourthreeOverOnesixtwo}{0.98~[0.97, 0.99]}  % experiments/FIGS/data/paired_v1err/mass/ratios.json :: ratios[11].ratio, ratios[11].ci95 :: sha256 d877a7f087260f64 :: 43 over 162, resolution ridge sigma_eff: paired geometric mean [95% interval]
\newcommand{\PairedMassResRidgeOnesevenOverOnesixtwo}{0.91~[0.87, 0.95]}  % experiments/FIGS/data/paired_v1err/mass/ratios.json :: ratios[12].ratio, ratios[12].ci95 :: sha256 d877a7f087260f64 :: 17 over 162, resolution ridge sigma_eff: paired geometric mean [95% interval]
\newcommand{\PairedMassResRidgeOnesevenOverFourthree}{0.93~[0.89, 0.97]}  % experiments/FIGS/data/paired_v1err/mass/ratios.json :: ratios[13].ratio, ratios[13].ci95 :: sha256 d877a7f087260f64 :: 17 over 43, resolution ridge sigma_eff: paired geometric mean [95% interval]
\newcommand{\PairedMassResRidgeOnesixtwoMassOverOnesixtwo}{0.678~[0.672, 0.684]}  % experiments/FIGS/data/paired_v1err/mass/ratios.json :: ratios[14].ratio, ratios[14].ci95 :: sha256 d877a7f087260f64 :: 162+mass over 162, resolution ridge sigma_eff: paired geometric mean [95% interval]
\newcommand{\PairedMassResRidgeOnesevenMassOverOneseven}{0.72~[0.69, 0.76]}  % experiments/FIGS/data/paired_v1err/mass/ratios.json :: ratios[15].ratio, ratios[15].ci95 :: sha256 d877a7f087260f64 :: 17+mass over 17, resolution ridge sigma_eff: paired geometric mean [95% interval]
\newcommand{\VcbNSignal}{1{,}659}  % experiments/FIGS/data/probe_ladder_vcbwindow_mlp2/sall.json :: tasks.bc_vs_rest.n_signal_test :: sha256 5b8b5c7c28602f93 :: X->bc test jets in the window
\newcommand{\VcbNBackground}{19{,}244}  % experiments/FIGS/data/probe_ladder_vcbwindow_mlp2/sall.json :: tasks.bc_vs_rest.n_background_test :: sha256 5b8b5c7c28602f93 :: bq, cs, bqq and QCD test jets in the window
\newcommand{\VcbEpsSixty}{60}  % experiments/FIGS/data/probe_ladder_vcbwindow_mlp2/sall.json :: tasks.bc_vs_rest.eps_s :: sha256 5b8b5c7c28602f93 :: signal efficiency, percent
\newcommand{\VcbOmaLinearOnesixtwo}{\ensuremath{(3.17 \pm 0.05)\times10^{-2}}}  % experiments/FIGS/data/probe_ladder_vcbwindow_mlp2/sall.json :: tasks.bc_vs_rest.arms.l162-s*.linear.log1m_auc :: sha256 5b8b5c7c28602f93 :: 1-AUC = exp of it, mean +- SD over 5 seeds
\newcommand{\VcbLogOmaLinearOnesixtwo}{\ensuremath{-}3.451}  % experiments/FIGS/data/probe_ladder_vcbwindow_mlp2/sall.json :: tasks.bc_vs_rest.arms.l162-s*.linear.log1m_auc :: sha256 5b8b5c7c28602f93 :: mean over seeds
\newcommand{\VcbRejLinearSixtyOnesixtwo}{207\,\ensuremath{\pm}\,19}  % experiments/FIGS/data/probe_ladder_vcbwindow_mlp2/sall.json :: tasks.bc_vs_rest.arms.l162-s*.linear.rejection_at['0.60'].rejection :: sha256 5b8b5c7c28602f93 :: background rejection at this signal efficiency, mean +- SD over seeds
\newcommand{\VcbBkgLeftLinearSixtyOnesixtwo}{93.4}  % experiments/FIGS/data/probe_ladder_vcbwindow_mlp2/sall.json :: tasks.bc_vs_rest.arms.l162-s*.linear.rejection_at['0.60'].n_bkg_pass :: sha256 5b8b5c7c28602f93 :: background jets passing the cut, mean over seeds
\newcommand{\VcbOmaLinearOneseven}{\ensuremath{(5.21 \pm 0.23)\times10^{-2}}}  % experiments/FIGS/data/probe_ladder_vcbwindow_mlp2/sall.json :: tasks.bc_vs_rest.arms.r16q1-s*.linear.log1m_auc :: sha256 5b8b5c7c28602f93 :: 1-AUC = exp of it, mean +- SD over 5 seeds
\newcommand{\VcbLogOmaLinearOneseven}{\ensuremath{-}2.955}  % experiments/FIGS/data/probe_ladder_vcbwindow_mlp2/sall.json :: tasks.bc_vs_rest.arms.r16q1-s*.linear.log1m_auc :: sha256 5b8b5c7c28602f93 :: mean over seeds
\newcommand{\VcbRejLinearSixtyOneseven}{67\,\ensuremath{\pm}\,14}  % experiments/FIGS/data/probe_ladder_vcbwindow_mlp2/sall.json :: tasks.bc_vs_rest.arms.r16q1-s*.linear.rejection_at['0.60'].rejection :: sha256 5b8b5c7c28602f93 :: background rejection at this signal efficiency, mean +- SD over seeds
\newcommand{\VcbBkgLeftLinearSixtyOneseven}{295.0}  % experiments/FIGS/data/probe_ladder_vcbwindow_mlp2/sall.json :: tasks.bc_vs_rest.arms.r16q1-s*.linear.rejection_at['0.60'].n_bkg_pass :: sha256 5b8b5c7c28602f93 :: background jets passing the cut, mean over seeds
\newcommand{\VcbOmaRatioLinear}{1.64}  % experiments/FIGS/data/probe_ladder_vcbwindow_mlp2/sall.json :: tasks.bc_vs_rest.arms.*.linear.log1m_auc :: sha256 5b8b5c7c28602f93 :: seed mean of 1-AUC, 17 classes over 162
\newcommand{\VcbFactorLinear}{1.64}  % experiments/FIGS/data/probe_ladder_vcbwindow_mlp2/sall.json :: tasks.bc_vs_rest.arms.*.linear.log1m_auc :: sha256 5b8b5c7c28602f93 :: exp of the difference of the seed means of log(1-AUC), 17 classes over 162
\newcommand{\VcbOmaMlpOnesixtwo}{\ensuremath{(3.102 \pm 0.023)\times10^{-2}}}  % experiments/FIGS/data/probe_ladder_vcbwindow_mlp2/sall.json :: tasks.bc_vs_rest.arms.l162-s*.mlp.log1m_auc :: sha256 5b8b5c7c28602f93 :: 1-AUC = exp of it, mean +- SD over 5 seeds
\newcommand{\VcbLogOmaMlpOnesixtwo}{\ensuremath{-}3.473}  % experiments/FIGS/data/probe_ladder_vcbwindow_mlp2/sall.json :: tasks.bc_vs_rest.arms.l162-s*.mlp.log1m_auc :: sha256 5b8b5c7c28602f93 :: mean over seeds
\newcommand{\VcbRejMlpSixtyOnesixtwo}{230\,\ensuremath{\pm}\,31}  % experiments/FIGS/data/probe_ladder_vcbwindow_mlp2/sall.json :: tasks.bc_vs_rest.arms.l162-s*.mlp.rejection_at['0.60'].rejection :: sha256 5b8b5c7c28602f93 :: background rejection at this signal efficiency, mean +- SD over seeds
\newcommand{\VcbBkgLeftMlpSixtyOnesixtwo}{84.8}  % experiments/FIGS/data/probe_ladder_vcbwindow_mlp2/sall.json :: tasks.bc_vs_rest.arms.l162-s*.mlp.rejection_at['0.60'].n_bkg_pass :: sha256 5b8b5c7c28602f93 :: background jets passing the cut, mean over seeds
\newcommand{\VcbOmaMlpOneseven}{\ensuremath{(4.66 \pm 0.19)\times10^{-2}}}  % experiments/FIGS/data/probe_ladder_vcbwindow_mlp2/sall.json :: tasks.bc_vs_rest.arms.r16q1-s*.mlp.log1m_auc :: sha256 5b8b5c7c28602f93 :: 1-AUC = exp of it, mean +- SD over 5 seeds
\newcommand{\VcbLogOmaMlpOneseven}{\ensuremath{-}3.066}  % experiments/FIGS/data/probe_ladder_vcbwindow_mlp2/sall.json :: tasks.bc_vs_rest.arms.r16q1-s*.mlp.log1m_auc :: sha256 5b8b5c7c28602f93 :: mean over seeds
\newcommand{\VcbRejMlpSixtyOneseven}{78\,\ensuremath{\pm}\,13}  % experiments/FIGS/data/probe_ladder_vcbwindow_mlp2/sall.json :: tasks.bc_vs_rest.arms.r16q1-s*.mlp.rejection_at['0.60'].rejection :: sha256 5b8b5c7c28602f93 :: background rejection at this signal efficiency, mean +- SD over seeds
\newcommand{\VcbBkgLeftMlpSixtyOneseven}{250.8}  % experiments/FIGS/data/probe_ladder_vcbwindow_mlp2/sall.json :: tasks.bc_vs_rest.arms.r16q1-s*.mlp.rejection_at['0.60'].n_bkg_pass :: sha256 5b8b5c7c28602f93 :: background jets passing the cut, mean over seeds
\newcommand{\VcbOmaRatioMlp}{1.50}  % experiments/FIGS/data/probe_ladder_vcbwindow_mlp2/sall.json :: tasks.bc_vs_rest.arms.*.mlp.log1m_auc :: sha256 5b8b5c7c28602f93 :: seed mean of 1-AUC, 17 classes over 162
\newcommand{\VcbFactorMlp}{1.50}  % experiments/FIGS/data/probe_ladder_vcbwindow_mlp2/sall.json :: tasks.bc_vs_rest.arms.*.mlp.log1m_auc :: sha256 5b8b5c7c28602f93 :: exp of the difference of the seed means of log(1-AUC), 17 classes over 162
\newcommand{\AccLegOneOnesixtwoSonebEFour}{0.32273}  % experiments/FIGS/data/leg1_metrics.json :: summary.l162-s1b.N10000.accuracy_mean :: sha256 3f564e960b2c86a7 :: wave 1, superseded; 3 fine-tuning seeds
\newcommand{\AccSdLegOneOnesixtwoSonebEFour}{0.01894}  % experiments/FIGS/data/leg1_metrics.json :: summary.l162-s1b.N10000.accuracy_sd :: sha256 3f564e960b2c86a7 :: wave 1, superseded
\newcommand{\AccLegOneOnesixtwoSonebEFive}{0.53506}  % experiments/FIGS/data/leg1_metrics.json :: summary.l162-s1b.N100000.accuracy_mean :: sha256 3f564e960b2c86a7 :: wave 1, superseded; 3 fine-tuning seeds
\newcommand{\AccSdLegOneOnesixtwoSonebEFive}{0.00044}  % experiments/FIGS/data/leg1_metrics.json :: summary.l162-s1b.N100000.accuracy_sd :: sha256 3f564e960b2c86a7 :: wave 1, superseded
\newcommand{\AccLegOneOnesixtwoSonebESix}{0.54337}  % experiments/FIGS/data/leg1_metrics.json :: summary.l162-s1b.N1000000.accuracy_mean :: sha256 3f564e960b2c86a7 :: wave 1, superseded; 3 fine-tuning seeds
\newcommand{\AccSdLegOneOnesixtwoSonebESix}{0.00074}  % experiments/FIGS/data/leg1_metrics.json :: summary.l162-s1b.N1000000.accuracy_sd :: sha256 3f564e960b2c86a7 :: wave 1, superseded
\newcommand{\AccLegOneOnesevenStwoEFour}{0.22684}  % experiments/FIGS/data/leg1_metrics.json :: summary.r16q1-s2.N10000.accuracy_mean :: sha256 3f564e960b2c86a7 :: wave 1, superseded; 3 fine-tuning seeds
\newcommand{\AccSdLegOneOnesevenStwoEFour}{0.00805}  % experiments/FIGS/data/leg1_metrics.json :: summary.r16q1-s2.N10000.accuracy_sd :: sha256 3f564e960b2c86a7 :: wave 1, superseded
\newcommand{\AccLegOneOnesevenStwoEFive}{0.49650}  % experiments/FIGS/data/leg1_metrics.json :: summary.r16q1-s2.N100000.accuracy_mean :: sha256 3f564e960b2c86a7 :: wave 1, superseded; 3 fine-tuning seeds
\newcommand{\AccSdLegOneOnesevenStwoEFive}{0.00110}  % experiments/FIGS/data/leg1_metrics.json :: summary.r16q1-s2.N100000.accuracy_sd :: sha256 3f564e960b2c86a7 :: wave 1, superseded
\newcommand{\AccLegOneOnesevenStwoESix}{0.51969}  % experiments/FIGS/data/leg1_metrics.json :: summary.r16q1-s2.N1000000.accuracy_mean :: sha256 3f564e960b2c86a7 :: wave 1, superseded; 3 fine-tuning seeds
\newcommand{\AccSdLegOneOnesevenStwoESix}{0.00072}  % experiments/FIGS/data/leg1_metrics.json :: summary.r16q1-s2.N1000000.accuracy_sd :: sha256 3f564e960b2c86a7 :: wave 1, superseded
\newcommand{\AccLegOneOnesevenSthreeEFour}{0.24203}  % experiments/FIGS/data/leg1_metrics.json :: summary.r16q1-s3.N10000.accuracy_mean :: sha256 3f564e960b2c86a7 :: wave 1, superseded; 3 fine-tuning seeds
\newcommand{\AccSdLegOneOnesevenSthreeEFour}{0.00213}  % experiments/FIGS/data/leg1_metrics.json :: summary.r16q1-s3.N10000.accuracy_sd :: sha256 3f564e960b2c86a7 :: wave 1, superseded
\newcommand{\AccLegOneOnesevenSthreeEFive}{0.49313}  % experiments/FIGS/data/leg1_metrics.json :: summary.r16q1-s3.N100000.accuracy_mean :: sha256 3f564e960b2c86a7 :: wave 1, superseded; 3 fine-tuning seeds
\newcommand{\AccSdLegOneOnesevenSthreeEFive}{0.00294}  % experiments/FIGS/data/leg1_metrics.json :: summary.r16q1-s3.N100000.accuracy_sd :: sha256 3f564e960b2c86a7 :: wave 1, superseded
\newcommand{\AccLegOneOnesevenSthreeESix}{0.51804}  % experiments/FIGS/data/leg1_metrics.json :: summary.r16q1-s3.N1000000.accuracy_mean :: sha256 3f564e960b2c86a7 :: wave 1, superseded; 3 fine-tuning seeds
\newcommand{\AccSdLegOneOnesevenSthreeESix}{0.00049}  % experiments/FIGS/data/leg1_metrics.json :: summary.r16q1-s3.N1000000.accuracy_sd :: sha256 3f564e960b2c86a7 :: wave 1, superseded
\newcommand{\AccLegOneOnesevenSfourEFour}{0.24487}  % experiments/FIGS/data/leg1_metrics.json :: summary.r16q1-s4.N10000.accuracy_mean :: sha256 3f564e960b2c86a7 :: wave 1, superseded; 3 fine-tuning seeds
\newcommand{\AccSdLegOneOnesevenSfourEFour}{0.00645}  % experiments/FIGS/data/leg1_metrics.json :: summary.r16q1-s4.N10000.accuracy_sd :: sha256 3f564e960b2c86a7 :: wave 1, superseded
\newcommand{\AccLegOneOnesevenSfourEFive}{0.49510}  % experiments/FIGS/data/leg1_metrics.json :: summary.r16q1-s4.N100000.accuracy_mean :: sha256 3f564e960b2c86a7 :: wave 1, superseded; 3 fine-tuning seeds
\newcommand{\AccSdLegOneOnesevenSfourEFive}{0.00133}  % experiments/FIGS/data/leg1_metrics.json :: summary.r16q1-s4.N100000.accuracy_sd :: sha256 3f564e960b2c86a7 :: wave 1, superseded
\newcommand{\AccLegOneOnesevenSfourESix}{0.51872}  % experiments/FIGS/data/leg1_metrics.json :: summary.r16q1-s4.N1000000.accuracy_mean :: sha256 3f564e960b2c86a7 :: wave 1, superseded; 3 fine-tuning seeds
\newcommand{\AccSdLegOneOnesevenSfourESix}{0.00049}  % experiments/FIGS/data/leg1_metrics.json :: summary.r16q1-s4.N1000000.accuracy_sd :: sha256 3f564e960b2c86a7 :: wave 1, superseded
\newcommand{\AccLegOneScratchEFour}{0.31369}  % experiments/FIGS/data/leg1_metrics.json :: summary.scratch.N10000.accuracy_mean :: sha256 3f564e960b2c86a7 :: wave 1, superseded; 3 fine-tuning seeds
\newcommand{\AccSdLegOneScratchEFour}{0.00712}  % experiments/FIGS/data/leg1_metrics.json :: summary.scratch.N10000.accuracy_sd :: sha256 3f564e960b2c86a7 :: wave 1, superseded
\newcommand{\AccLegOneScratchEFive}{0.42015}  % experiments/FIGS/data/leg1_metrics.json :: summary.scratch.N100000.accuracy_mean :: sha256 3f564e960b2c86a7 :: wave 1, superseded; 3 fine-tuning seeds
\newcommand{\AccSdLegOneScratchEFive}{0.00267}  % experiments/FIGS/data/leg1_metrics.json :: summary.scratch.N100000.accuracy_sd :: sha256 3f564e960b2c86a7 :: wave 1, superseded
\newcommand{\AccLegOneScratchESix}{0.48621}  % experiments/FIGS/data/leg1_metrics.json :: summary.scratch.N1000000.accuracy_mean :: sha256 3f564e960b2c86a7 :: wave 1, superseded; 3 fine-tuning seeds
\newcommand{\AccSdLegOneScratchESix}{0.00100}  % experiments/FIGS/data/leg1_metrics.json :: summary.scratch.N1000000.accuracy_sd :: sha256 3f564e960b2c86a7 :: wave 1, superseded
\newcommand{\AccLegOneSophonPublicEFour}{0.21613}  % experiments/FIGS/data/leg1_metrics.json :: summary.sophon-public.N10000.accuracy_mean :: sha256 3f564e960b2c86a7 :: wave 1, superseded; 3 fine-tuning seeds
\newcommand{\AccSdLegOneSophonPublicEFour}{0.00734}  % experiments/FIGS/data/leg1_metrics.json :: summary.sophon-public.N10000.accuracy_sd :: sha256 3f564e960b2c86a7 :: wave 1, superseded
\newcommand{\AccLegOneSophonPublicEFive}{0.53187}  % experiments/FIGS/data/leg1_metrics.json :: summary.sophon-public.N100000.accuracy_mean :: sha256 3f564e960b2c86a7 :: wave 1, superseded; 3 fine-tuning seeds
\newcommand{\AccSdLegOneSophonPublicEFive}{0.00231}  % experiments/FIGS/data/leg1_metrics.json :: summary.sophon-public.N100000.accuracy_sd :: sha256 3f564e960b2c86a7 :: wave 1, superseded
\newcommand{\AccLegOneSophonPublicESix}{0.54273}  % experiments/FIGS/data/leg1_metrics.json :: summary.sophon-public.N1000000.accuracy_mean :: sha256 3f564e960b2c86a7 :: wave 1, superseded; 3 fine-tuning seeds
\newcommand{\AccSdLegOneSophonPublicESix}{0.00064}  % experiments/FIGS/data/leg1_metrics.json :: summary.sophon-public.N1000000.accuracy_sd :: sha256 3f564e960b2c86a7 :: wave 1, superseded
\newcommand{\AccLegTwoOnesixtwoSonebEFour}{0.78606}  % experiments/FIGS/data/leg2_metrics.json :: summary.l162-s1b.N10000.accuracy_mean :: sha256 701c58c3240274bb :: wave 1, superseded; 3 fine-tuning seeds
\newcommand{\AccSdLegTwoOnesixtwoSonebEFour}{0.00111}  % experiments/FIGS/data/leg2_metrics.json :: summary.l162-s1b.N10000.accuracy_sd :: sha256 701c58c3240274bb :: wave 1, superseded
\newcommand{\AccLegTwoOnesixtwoSonebEFive}{0.83598}  % experiments/FIGS/data/leg2_metrics.json :: summary.l162-s1b.N100000.accuracy_mean :: sha256 701c58c3240274bb :: wave 1, superseded; 3 fine-tuning seeds
\newcommand{\AccSdLegTwoOnesixtwoSonebEFive}{0.00037}  % experiments/FIGS/data/leg2_metrics.json :: summary.l162-s1b.N100000.accuracy_sd :: sha256 701c58c3240274bb :: wave 1, superseded
\newcommand{\AccLegTwoOnesixtwoSonebESix}{0.84704}  % experiments/FIGS/data/leg2_metrics.json :: summary.l162-s1b.N1000000.accuracy_mean :: sha256 701c58c3240274bb :: wave 1, superseded; 3 fine-tuning seeds
\newcommand{\AccSdLegTwoOnesixtwoSonebESix}{0.00027}  % experiments/FIGS/data/leg2_metrics.json :: summary.l162-s1b.N1000000.accuracy_sd :: sha256 701c58c3240274bb :: wave 1, superseded
\newcommand{\AccLegTwoOnesevenStwoEFour}{0.72156}  % experiments/FIGS/data/leg2_metrics.json :: summary.r16q1-s2.N10000.accuracy_mean :: sha256 701c58c3240274bb :: wave 1, superseded; 3 fine-tuning seeds
\newcommand{\AccSdLegTwoOnesevenStwoEFour}{0.01049}  % experiments/FIGS/data/leg2_metrics.json :: summary.r16q1-s2.N10000.accuracy_sd :: sha256 701c58c3240274bb :: wave 1, superseded
\newcommand{\AccLegTwoOnesevenStwoEFive}{0.82614}  % experiments/FIGS/data/leg2_metrics.json :: summary.r16q1-s2.N100000.accuracy_mean :: sha256 701c58c3240274bb :: wave 1, superseded; 3 fine-tuning seeds
\newcommand{\AccSdLegTwoOnesevenStwoEFive}{0.00058}  % experiments/FIGS/data/leg2_metrics.json :: summary.r16q1-s2.N100000.accuracy_sd :: sha256 701c58c3240274bb :: wave 1, superseded
\newcommand{\AccLegTwoOnesevenStwoESix}{0.83943}  % experiments/FIGS/data/leg2_metrics.json :: summary.r16q1-s2.N1000000.accuracy_mean :: sha256 701c58c3240274bb :: wave 1, superseded; 3 fine-tuning seeds
\newcommand{\AccSdLegTwoOnesevenStwoESix}{0.00052}  % experiments/FIGS/data/leg2_metrics.json :: summary.r16q1-s2.N1000000.accuracy_sd :: sha256 701c58c3240274bb :: wave 1, superseded
\newcommand{\AccLegTwoOnesevenSthreeEFour}{0.72103}  % experiments/FIGS/data/leg2_metrics.json :: summary.r16q1-s3.N10000.accuracy_mean :: sha256 701c58c3240274bb :: wave 1, superseded; 3 fine-tuning seeds
\newcommand{\AccSdLegTwoOnesevenSthreeEFour}{0.01651}  % experiments/FIGS/data/leg2_metrics.json :: summary.r16q1-s3.N10000.accuracy_sd :: sha256 701c58c3240274bb :: wave 1, superseded
\newcommand{\AccLegTwoOnesevenSthreeEFive}{0.82629}  % experiments/FIGS/data/leg2_metrics.json :: summary.r16q1-s3.N100000.accuracy_mean :: sha256 701c58c3240274bb :: wave 1, superseded; 3 fine-tuning seeds
\newcommand{\AccSdLegTwoOnesevenSthreeEFive}{0.00035}  % experiments/FIGS/data/leg2_metrics.json :: summary.r16q1-s3.N100000.accuracy_sd :: sha256 701c58c3240274bb :: wave 1, superseded
\newcommand{\AccLegTwoOnesevenSthreeESix}{0.83944}  % experiments/FIGS/data/leg2_metrics.json :: summary.r16q1-s3.N1000000.accuracy_mean :: sha256 701c58c3240274bb :: wave 1, superseded; 3 fine-tuning seeds
\newcommand{\AccSdLegTwoOnesevenSthreeESix}{0.00019}  % experiments/FIGS/data/leg2_metrics.json :: summary.r16q1-s3.N1000000.accuracy_sd :: sha256 701c58c3240274bb :: wave 1, superseded
\newcommand{\AccLegTwoOnesevenSfourEFour}{0.72953}  % experiments/FIGS/data/leg2_metrics.json :: summary.r16q1-s4.N10000.accuracy_mean :: sha256 701c58c3240274bb :: wave 1, superseded; 3 fine-tuning seeds
\newcommand{\AccSdLegTwoOnesevenSfourEFour}{0.01069}  % experiments/FIGS/data/leg2_metrics.json :: summary.r16q1-s4.N10000.accuracy_sd :: sha256 701c58c3240274bb :: wave 1, superseded
\newcommand{\AccLegTwoOnesevenSfourEFive}{0.82564}  % experiments/FIGS/data/leg2_metrics.json :: summary.r16q1-s4.N100000.accuracy_mean :: sha256 701c58c3240274bb :: wave 1, superseded; 3 fine-tuning seeds
\newcommand{\AccSdLegTwoOnesevenSfourEFive}{0.00041}  % experiments/FIGS/data/leg2_metrics.json :: summary.r16q1-s4.N100000.accuracy_sd :: sha256 701c58c3240274bb :: wave 1, superseded
\newcommand{\AccLegTwoOnesevenSfourESix}{0.83909}  % experiments/FIGS/data/leg2_metrics.json :: summary.r16q1-s4.N1000000.accuracy_mean :: sha256 701c58c3240274bb :: wave 1, superseded; 3 fine-tuning seeds
\newcommand{\AccSdLegTwoOnesevenSfourESix}{0.00037}  % experiments/FIGS/data/leg2_metrics.json :: summary.r16q1-s4.N1000000.accuracy_sd :: sha256 701c58c3240274bb :: wave 1, superseded
\newcommand{\AccLegTwoRefEonearmsSoneRef}{0.81770}  % experiments/FIGS/data/leg2_metrics.json :: summary.ref_e1arms-s1.ref.accuracy_mean :: sha256 701c58c3240274bb :: wave 1, superseded; 1 fine-tuning seeds
\newcommand{\AccSdLegTwoRefEonearmsSoneRef}{---}  % experiments/FIGS/data/leg2_metrics.json :: summary.ref_e1arms-s1.ref.accuracy_sd :: sha256 701c58c3240274bb :: wave 1, superseded
\newcommand{\AccLegTwoRefEonearmsStwoRef}{0.81602}  % experiments/FIGS/data/leg2_metrics.json :: summary.ref_e1arms-s2.ref.accuracy_mean :: sha256 701c58c3240274bb :: wave 1, superseded; 1 fine-tuning seeds
\newcommand{\AccSdLegTwoRefEonearmsStwoRef}{---}  % experiments/FIGS/data/leg2_metrics.json :: summary.ref_e1arms-s2.ref.accuracy_sd :: sha256 701c58c3240274bb :: wave 1, superseded
\newcommand{\AccLegTwoRefEonearmsSthreeRef}{0.81625}  % experiments/FIGS/data/leg2_metrics.json :: summary.ref_e1arms-s3.ref.accuracy_mean :: sha256 701c58c3240274bb :: wave 1, superseded; 1 fine-tuning seeds
\newcommand{\AccSdLegTwoRefEonearmsSthreeRef}{---}  % experiments/FIGS/data/leg2_metrics.json :: summary.ref_e1arms-s3.ref.accuracy_sd :: sha256 701c58c3240274bb :: wave 1, superseded
\newcommand{\AccLegTwoScratchEFour}{0.68230}  % experiments/FIGS/data/leg2_metrics.json :: summary.scratch.N10000.accuracy_mean :: sha256 701c58c3240274bb :: wave 1, superseded; 3 fine-tuning seeds
\newcommand{\AccSdLegTwoScratchEFour}{0.00282}  % experiments/FIGS/data/leg2_metrics.json :: summary.scratch.N10000.accuracy_sd :: sha256 701c58c3240274bb :: wave 1, superseded
\newcommand{\AccLegTwoScratchEFive}{0.77769}  % experiments/FIGS/data/leg2_metrics.json :: summary.scratch.N100000.accuracy_mean :: sha256 701c58c3240274bb :: wave 1, superseded; 3 fine-tuning seeds
\newcommand{\AccSdLegTwoScratchEFive}{0.00208}  % experiments/FIGS/data/leg2_metrics.json :: summary.scratch.N100000.accuracy_sd :: sha256 701c58c3240274bb :: wave 1, superseded
\newcommand{\AccLegTwoScratchESix}{0.82303}  % experiments/FIGS/data/leg2_metrics.json :: summary.scratch.N1000000.accuracy_mean :: sha256 701c58c3240274bb :: wave 1, superseded; 3 fine-tuning seeds
\newcommand{\AccSdLegTwoScratchESix}{0.00134}  % experiments/FIGS/data/leg2_metrics.json :: summary.scratch.N1000000.accuracy_sd :: sha256 701c58c3240274bb :: wave 1, superseded
\newcommand{\AccLegTwoSophonPublicEFour}{0.72739}  % experiments/FIGS/data/leg2_metrics.json :: summary.sophon-public.N10000.accuracy_mean :: sha256 701c58c3240274bb :: wave 1, superseded; 3 fine-tuning seeds
\newcommand{\AccSdLegTwoSophonPublicEFour}{0.00470}  % experiments/FIGS/data/leg2_metrics.json :: summary.sophon-public.N10000.accuracy_sd :: sha256 701c58c3240274bb :: wave 1, superseded
\newcommand{\AccLegTwoSophonPublicEFive}{0.83492}  % experiments/FIGS/data/leg2_metrics.json :: summary.sophon-public.N100000.accuracy_mean :: sha256 701c58c3240274bb :: wave 1, superseded; 3 fine-tuning seeds
\newcommand{\AccSdLegTwoSophonPublicEFive}{0.00049}  % experiments/FIGS/data/leg2_metrics.json :: summary.sophon-public.N100000.accuracy_sd :: sha256 701c58c3240274bb :: wave 1, superseded
\newcommand{\AccLegTwoSophonPublicESix}{0.84632}  % experiments/FIGS/data/leg2_metrics.json :: summary.sophon-public.N1000000.accuracy_mean :: sha256 701c58c3240274bb :: wave 1, superseded; 3 fine-tuning seeds
\newcommand{\AccSdLegTwoSophonPublicESix}{0.00031}  % experiments/FIGS/data/leg2_metrics.json :: summary.sophon-public.N1000000.accuracy_sd :: sha256 701c58c3240274bb :: wave 1, superseded
\newcommand{\FtBudgetSteps}{20{,}000}  % scripts/build_ft_budget_jobs.py :: EPOCHS * STEPS :: sha256 f2ed9faa8e3bd622 :: the largest number of optimizer steps below the full training set
\newcommand{\FtLongSteps}{80{,}000}  % scripts/build_ft_budget_jobs.py :: LONG_EPOCHS * STEPS :: sha256 f2ed9faa8e3bd622 :: the largest number of optimizer steps of the long budget at LONG_N jets (JetClass, JetClass-II)
\newcommand{\DesignNTrainSelected}{109{,}938{,}036}  % experiments/FIGS/data/v2/jc2_train/counts.json :: n_selected :: sha256 b0b5a8b63aa808c0 :: selected jets of the pretrain_train files (experiments/FT/count_jc2_train.py)
\newcommand{\FtBudgetEval}{500}  % scripts/build_ft_budget_jobs.py :: STEPS :: sha256 f2ed9faa8e3bd622 :: steps between validation passes at and above SMALL_N training jets
\newcommand{\FtBudgetEvalSmall}{50}  % scripts/build_ft_budget_jobs.py :: SMALL_STEPS :: sha256 f2ed9faa8e3bd622 :: steps between validation passes below SMALL_N training jets
\newcommand{\FtBudgetEvalSmallBelow}{\ensuremath{10^{5}}}  % scripts/build_ft_budget_jobs.py :: SMALL_N :: sha256 f2ed9faa8e3bd622 :: training jets below which the validation runs every SMALL_STEPS steps
\newcommand{\FtBudgetPatience}{10}  % scripts/build_ft_budget_jobs.py :: PATIENCE :: sha256 f2ed9faa8e3bd622 :: validation passes without improvement before training stops
\newcommand{\FtHeadMult}{50}  % scripts/build_ft_budget_jobs.py :: HEAD_MULT :: sha256 f2ed9faa8e3bd622 :: the head's learning-rate multiplier
\newcommand{\FtHeadCap}{\ensuremath{5\times10^{-3}}}  % scripts/build_ft_budget_jobs.py :: HEAD_CAP :: sha256 f2ed9faa8e3bd622 :: the head's largest learning rate
\newcommand{\FtBudgetBatch}{512}  % scripts/build_ft_budget_jobs.py :: BATCH :: sha256 f2ed9faa8e3bd622
\newcommand{\FtBudgetWeightDecay}{0.01}  % scripts/build_ft_budget_jobs.py :: --optimizer-option weight_decay (every cell) :: sha256 f2ed9faa8e3bd622
\newcommand{\FtSweepPre}{\ensuremath{10^{-4}}, \ensuremath{3\times10^{-4}} and \ensuremath{10^{-3}}}  % scripts/build_ft_budget_jobs.py :: SWEEP_PRE :: sha256 f2ed9faa8e3bd622 :: backbone rates of the sweep
\newcommand{\FtSweepScratch}{\ensuremath{3\times10^{-4}}, \ensuremath{10^{-3}} and \ensuremath{3\times10^{-3}}}  % scripts/build_ft_budget_jobs.py :: SWEEP_SCRATCH :: sha256 f2ed9faa8e3bd622 :: from-scratch rates of the sweep
\newcommand{\FtSweepN}{10{,}000}  % scripts/build_ft_budget_jobs.py :: SWEEP_N :: sha256 f2ed9faa8e3bd622 :: training jets of the sweep
\newcommand{\FtValJets}{200{,}000}  % scripts/build_ft_budget_jobs.py :: VAL_N :: sha256 f2ed9faa8e3bd622 :: the fixed validation set
\newcommand{\FtValJetsBal}{20{,}480}  % scripts/build_ft_budget_jobs.py :: VAL_SIZE[jc2b] :: sha256 f2ed9faa8e3bd622 :: the balanced JetClass-II validation set
\newcommand{\LpVfourCMax}{\ensuremath{10^{4}}}  % experiments/EVAL/probe_bal.py :: C_GRID_V4 :: sha256 a7804216e3179d9d :: largest inverse regularisation of the v4 probes
\newcommand{\JcBalTestPerClass}{5{,}000}  % experiments/FT/stage_jc2_balanced.py :: TEST_PER_CLASS :: sha256 0260341d5ed75fec
\newcommand{\JcBalValPerClass}{109}  % experiments/FT/stage_jc2_balanced.py :: VAL_PER_CLASS :: sha256 0260341d5ed75fec
\newcommand{\JcBalPairTrainPerSide}{\ensuremath{10^{5}}}  % experiments/FT/stage_jc2_balanced.py :: PAIR_TRAIN_PER_SIDE :: sha256 0260341d5ed75fec
\newcommand{\JcBalPairTestPerSide}{\ensuremath{2\times10^{5}}}  % experiments/FT/stage_jc2_balanced.py :: PAIR_TEST_PER_SIDE :: sha256 0260341d5ed75fec
\newcommand{\JcBalNShort}{11}  % experiments/FIGS/data/v2/jc2bal/census.json :: counts.ft_test :: sha256 a4379ee0817254cf :: classes with fewer than 5000 selected test jets
\newcommand{\JcBalRarestTest}{224}  % experiments/FIGS/data/v2/jc2bal/census.json :: counts.ft_test (minimum) :: sha256 a4379ee0817254cf :: label_QCD_bcs
\newcommand{\JcBalShortTrainESix}{11}  % experiments/FIGS/data/v2/jc2bal/census.json :: counts.ft_heldout :: sha256 a4379ee0817254cf :: classes with fewer held-out jets than the 1e6 quota
\newcommand{\JcBalNTest}{903{,}922}  % experiments/FIGS/data/v2/jc2bal/manifest.json :: outputs[test.parquet].rows :: sha256 9c7775be6ff56ac3 :: jets of the balanced test set
\newcommand{\FtLrPre}{\ensuremath{10^{-3}}}  % experiments/FIGS/data/v2/budget/lrsweep_metrics.json :: argmax of the summed validation metric, pretrained :: sha256 7cdc148054ca902c
\newcommand{\FtLrScratchTop}{\ensuremath{10^{-3}}}  % experiments/FIGS/data/v2/budget/lrsweep_metrics.json :: argmax of the validation metric on top, from scratch :: sha256 7cdc148054ca902c :: agrees with LR of the launched job-ft-budget-bench, -jc1, -jc2b, -scratch-full and -long-scratch-full from-scratch cells
\newcommand{\FtLrScratchQg}{\ensuremath{3\times10^{-3}}}  % experiments/FIGS/data/v2/budget/lrsweep_metrics.json :: argmax of the validation metric on quark/gluon, from scratch :: sha256 7cdc148054ca902c :: agrees with LR of the launched job-ft-budget-bench and -scratch-full from-scratch cells
\newcommand{\MlstNParamsFt}{2{,}209{,}270}  % experiments/FIGS/mlst_v4.py :: our_params(2, 17) :: sha256 fa536b4866ffd651 :: trainable parameters of our fine-tuned network, two outputs (top, quark/gluon)
\newcommand{\MlstNParamsFtKin}{2{,}207{,}950}  % experiments/FIGS/mlst_v4.py :: our_params(2, 7) :: sha256 fa536b4866ffd651 :: trainable parameters of our fine-tuned network, two outputs, the 7 kinematic inputs
\newcommand{\MlstNParamsJc}{2{,}213{,}374}  % experiments/FIGS/mlst_v4.py :: our_params(10, 17) :: sha256 fa536b4866ffd651 :: trainable parameters of our fine-tuned network, the ten JetClass classes
\newcommand{\MlstProbeRunsDone}{27}  % experiments/FIGS/data/v2/index.json :: runs with a linear-probe fit (count) :: sha256 c6c4138ee5c86423 :: of the paper's runs
\newcommand{\MlstPaperRuns}{27}  % configs/arms/v2_paper_runs.txt :: runs (count) :: sha256 b04562b14ea8cbc7
\newcommand{\MlstSdRelUnc}{50\%}  % configs/arms/v2_paper_runs.txt :: relative standard error of a standard deviation from the runs of a hierarchy level, 1/sqrt(2(n-1)) :: sha256 b04562b14ea8cbc7
\newcommand{\MlstCmpJcTwoResOverFull}{1.026\,\ensuremath{\pm}\,0.005}  % experiments/FIGS/data/v2/index.json :: linear probe jc2_188 N=1000000: (1 - macro_auc_ovr) of L162 over L188, run by run :: sha256 c6c4138ee5c86423
\newcommand{\MlstCmpJcTwoResOverFullPr}{2.6\% more errors}  % experiments/FIGS/data/v2/index.json :: linear probe jc2_188 N=1000000: (1 - macro_auc_ovr) of L162 over L188, run by run :: sha256 c6c4138ee5c86423 :: prose twin: two significant figures, direction in words
\newcommand{\MlstCmpJcTwoMidOverFull}{1.227\,\ensuremath{\pm}\,0.009}  % experiments/FIGS/data/v2/index.json :: linear probe jc2_188 N=1000000: (1 - macro_auc_ovr) of R42_Q1 over L188, run by run :: sha256 c6c4138ee5c86423
\newcommand{\MlstCmpJcTwoMidOverFullPr}{23\% more errors}  % experiments/FIGS/data/v2/index.json :: linear probe jc2_188 N=1000000: (1 - macro_auc_ovr) of R42_Q1 over L188, run by run :: sha256 c6c4138ee5c86423 :: prose twin: two significant figures, direction in words
\newcommand{\MlstCmpJcTwoCoarseOverFull}{1.557\,\ensuremath{\pm}\,0.019}  % experiments/FIGS/data/v2/index.json :: linear probe jc2_188 N=1000000: (1 - macro_auc_ovr) of R16_Q1 over L188, run by run :: sha256 c6c4138ee5c86423
\newcommand{\MlstCmpJcTwoCoarseOverFullPr}{56\% more errors}  % experiments/FIGS/data/v2/index.json :: linear probe jc2_188 N=1000000: (1 - macro_auc_ovr) of R16_Q1 over L188, run by run :: sha256 c6c4138ee5c86423 :: prose twin: two significant figures, direction in words
\newcommand{\MlstCmpJcTwoResMassOverRes}{1.022 and 1.016}  % experiments/FIGS/data/v2/index.json :: linear probe jc2_188 N=1000000: (1 - macro_auc_ovr) of L162_MASS over L162, run by run :: sha256 c6c4138ee5c86423
\newcommand{\MlstCmpJcTwoResMassOverResPr}{equal within 3\%}  % experiments/FIGS/data/v2/index.json :: linear probe jc2_188 N=1000000: (1 - macro_auc_ovr) of L162_MASS over L162, run by run :: sha256 c6c4138ee5c86423 :: prose twin: two significant figures, direction in words
\newcommand{\MlstCmpJcTwoResMassOverResPrMax}{equal within 3\%}  % experiments/FIGS/data/v2/index.json :: linear probe jc2_188 N=1000000: (1 - macro_auc_ovr) of L162_MASS over L162, run by run :: sha256 c6c4138ee5c86423 :: prose twin of the run farther from 1
\newcommand{\MlstCmpJcTwoCoarseMassOverCoarse}{1.054 and 0.977}  % experiments/FIGS/data/v2/index.json :: linear probe jc2_188 N=1000000: (1 - macro_auc_ovr) of R16_Q1_MASS over R16_Q1, run by run :: sha256 c6c4138ee5c86423
\newcommand{\MlstCmpJcTwoCoarseMassOverCoarsePr}{equal within 6\%}  % experiments/FIGS/data/v2/index.json :: linear probe jc2_188 N=1000000: (1 - macro_auc_ovr) of R16_Q1_MASS over R16_Q1, run by run :: sha256 c6c4138ee5c86423 :: prose twin: two significant figures, direction in words
\newcommand{\MlstCmpJcTwoCoarseMassOverCoarsePrMax}{equal within 6\%}  % experiments/FIGS/data/v2/index.json :: linear probe jc2_188 N=1000000: (1 - macro_auc_ovr) of R16_Q1_MASS over R16_Q1, run by run :: sha256 c6c4138ee5c86423 :: prose twin of the run farther from 1
\newcommand{\MlstCmpJcTwoResMassOverFull}{1.053 and 1.037}  % experiments/FIGS/data/v2/index.json :: linear probe jc2_188 N=1000000: (1 - macro_auc_ovr) of L162_MASS over L188, run by run :: sha256 c6c4138ee5c86423
\newcommand{\MlstCmpJcTwoResMassOverFullPr}{at least 3.7\% more errors}  % experiments/FIGS/data/v2/index.json :: linear probe jc2_188 N=1000000: (1 - macro_auc_ovr) of L162_MASS over L188, run by run :: sha256 c6c4138ee5c86423 :: prose twin: two significant figures, direction in words
\newcommand{\MlstCmpJcTwoResMassOverFullPrMax}{up to 5.3\% more errors}  % experiments/FIGS/data/v2/index.json :: linear probe jc2_188 N=1000000: (1 - macro_auc_ovr) of L162_MASS over L188, run by run :: sha256 c6c4138ee5c86423 :: prose twin of the run farther from 1
\newcommand{\MlstCmpJcTwoCoarseMassOverFull}{1.617 and 1.528}  % experiments/FIGS/data/v2/index.json :: linear probe jc2_188 N=1000000: (1 - macro_auc_ovr) of R16_Q1_MASS over L188, run by run :: sha256 c6c4138ee5c86423
\newcommand{\MlstCmpJcTwoCoarseMassOverFullPr}{at least 53\% more errors}  % experiments/FIGS/data/v2/index.json :: linear probe jc2_188 N=1000000: (1 - macro_auc_ovr) of R16_Q1_MASS over L188, run by run :: sha256 c6c4138ee5c86423 :: prose twin: two significant figures, direction in words
\newcommand{\MlstCmpJcTwoCoarseMassOverFullPrMax}{up to 62\% more errors}  % experiments/FIGS/data/v2/index.json :: linear probe jc2_188 N=1000000: (1 - macro_auc_ovr) of R16_Q1_MASS over L188, run by run :: sha256 c6c4138ee5c86423 :: prose twin of the run farther from 1
\newcommand{\MlstCmpJcTwoFzeroOverFull}{1.424 and 1.462}  % experiments/FIGS/data/v2/index.json :: linear probe jc2_188 N=1000000: (1 - macro_auc_ovr) of FLAV_F0 over L188, run by run :: sha256 c6c4138ee5c86423
\newcommand{\MlstCmpJcTwoFzeroOverFullPr}{at least 42\% more errors}  % experiments/FIGS/data/v2/index.json :: linear probe jc2_188 N=1000000: (1 - macro_auc_ovr) of FLAV_F0 over L188, run by run :: sha256 c6c4138ee5c86423 :: prose twin: two significant figures, direction in words
\newcommand{\MlstCmpJcTwoFzeroOverFullPrMax}{up to 46\% more errors}  % experiments/FIGS/data/v2/index.json :: linear probe jc2_188 N=1000000: (1 - macro_auc_ovr) of FLAV_F0 over L188, run by run :: sha256 c6c4138ee5c86423 :: prose twin of the run farther from 1
\newcommand{\MlstCmpJcTwoFoneOverFull}{1.218 and 1.212}  % experiments/FIGS/data/v2/index.json :: linear probe jc2_188 N=1000000: (1 - macro_auc_ovr) of FLAV_F1 over L188, run by run :: sha256 c6c4138ee5c86423
\newcommand{\MlstCmpJcTwoFoneOverFullPr}{at least 21\% more errors}  % experiments/FIGS/data/v2/index.json :: linear probe jc2_188 N=1000000: (1 - macro_auc_ovr) of FLAV_F1 over L188, run by run :: sha256 c6c4138ee5c86423 :: prose twin: two significant figures, direction in words
\newcommand{\MlstCmpJcTwoFoneOverFullPrMax}{up to 22\% more errors}  % experiments/FIGS/data/v2/index.json :: linear probe jc2_188 N=1000000: (1 - macro_auc_ovr) of FLAV_F1 over L188, run by run :: sha256 c6c4138ee5c86423 :: prose twin of the run farther from 1
\newcommand{\MlstCmpJcTwoFonerOverFull}{1.206 and 1.208}  % experiments/FIGS/data/v2/index.json :: linear probe jc2_188 N=1000000: (1 - macro_auc_ovr) of FLAV_F1R over L188, run by run :: sha256 c6c4138ee5c86423
\newcommand{\MlstCmpJcTwoFonerOverFullPr}{at least 21\% more errors}  % experiments/FIGS/data/v2/index.json :: linear probe jc2_188 N=1000000: (1 - macro_auc_ovr) of FLAV_F1R over L188, run by run :: sha256 c6c4138ee5c86423 :: prose twin: two significant figures, direction in words
\newcommand{\MlstCmpJcTwoFonerOverFullPrMax}{up to 21\% more errors}  % experiments/FIGS/data/v2/index.json :: linear probe jc2_188 N=1000000: (1 - macro_auc_ovr) of FLAV_F1R over L188, run by run :: sha256 c6c4138ee5c86423 :: prose twin of the run farther from 1
\newcommand{\MlstCmpJcTwoFzeroOverCoarse}{0.9275 and 0.9345}  % experiments/FIGS/data/v2/index.json :: linear probe jc2_188 N=1000000: (1 - macro_auc_ovr) of FLAV_F0 over R16_Q1, run by run :: sha256 c6c4138ee5c86423
\newcommand{\MlstCmpJcTwoFzeroOverCoarsePr}{at least 6.5\% fewer errors}  % experiments/FIGS/data/v2/index.json :: linear probe jc2_188 N=1000000: (1 - macro_auc_ovr) of FLAV_F0 over R16_Q1, run by run :: sha256 c6c4138ee5c86423 :: prose twin: two significant figures, direction in words
\newcommand{\MlstCmpJcTwoFzeroOverCoarsePrMax}{up to 7.2\% fewer errors}  % experiments/FIGS/data/v2/index.json :: linear probe jc2_188 N=1000000: (1 - macro_auc_ovr) of FLAV_F0 over R16_Q1, run by run :: sha256 c6c4138ee5c86423 :: prose twin of the run farther from 1
\newcommand{\MlstCmpJcTwoFoneOverCoarse}{0.794 and 0.775}  % experiments/FIGS/data/v2/index.json :: linear probe jc2_188 N=1000000: (1 - macro_auc_ovr) of FLAV_F1 over R16_Q1, run by run :: sha256 c6c4138ee5c86423
\newcommand{\MlstCmpJcTwoFoneOverCoarsePr}{at least 21\% fewer errors}  % experiments/FIGS/data/v2/index.json :: linear probe jc2_188 N=1000000: (1 - macro_auc_ovr) of FLAV_F1 over R16_Q1, run by run :: sha256 c6c4138ee5c86423 :: prose twin: two significant figures, direction in words
\newcommand{\MlstCmpJcTwoFoneOverCoarsePrMax}{up to 23\% fewer errors}  % experiments/FIGS/data/v2/index.json :: linear probe jc2_188 N=1000000: (1 - macro_auc_ovr) of FLAV_F1 over R16_Q1, run by run :: sha256 c6c4138ee5c86423 :: prose twin of the run farther from 1
\newcommand{\MlstCmpJcTwoFonerOverCoarse}{0.786 and 0.772}  % experiments/FIGS/data/v2/index.json :: linear probe jc2_188 N=1000000: (1 - macro_auc_ovr) of FLAV_F1R over R16_Q1, run by run :: sha256 c6c4138ee5c86423
\newcommand{\MlstCmpJcTwoFonerOverCoarsePr}{at least 21\% fewer errors}  % experiments/FIGS/data/v2/index.json :: linear probe jc2_188 N=1000000: (1 - macro_auc_ovr) of FLAV_F1R over R16_Q1, run by run :: sha256 c6c4138ee5c86423 :: prose twin: two significant figures, direction in words
\newcommand{\MlstCmpJcTwoFonerOverCoarsePrMax}{up to 23\% fewer errors}  % experiments/FIGS/data/v2/index.json :: linear probe jc2_188 N=1000000: (1 - macro_auc_ovr) of FLAV_F1R over R16_Q1, run by run :: sha256 c6c4138ee5c86423 :: prose twin of the run farther from 1
\newcommand{\MlstCmpJcTwoFoneOverFzero}{0.856 and 0.829}  % experiments/FIGS/data/v2/index.json :: linear probe jc2_188 N=1000000: (1 - macro_auc_ovr) of FLAV_F1 over FLAV_F0, run by run :: sha256 c6c4138ee5c86423
\newcommand{\MlstCmpJcTwoFoneOverFzeroPr}{at least 14\% fewer errors}  % experiments/FIGS/data/v2/index.json :: linear probe jc2_188 N=1000000: (1 - macro_auc_ovr) of FLAV_F1 over FLAV_F0, run by run :: sha256 c6c4138ee5c86423 :: prose twin: two significant figures, direction in words
\newcommand{\MlstCmpJcTwoFoneOverFzeroPrMax}{up to 17\% fewer errors}  % experiments/FIGS/data/v2/index.json :: linear probe jc2_188 N=1000000: (1 - macro_auc_ovr) of FLAV_F1 over FLAV_F0, run by run :: sha256 c6c4138ee5c86423 :: prose twin of the run farther from 1
\newcommand{\MlstCmpJcTwoFonerOverFzero}{0.847 and 0.826}  % experiments/FIGS/data/v2/index.json :: linear probe jc2_188 N=1000000: (1 - macro_auc_ovr) of FLAV_F1R over FLAV_F0, run by run :: sha256 c6c4138ee5c86423
\newcommand{\MlstCmpJcTwoFonerOverFzeroPr}{at least 15\% fewer errors}  % experiments/FIGS/data/v2/index.json :: linear probe jc2_188 N=1000000: (1 - macro_auc_ovr) of FLAV_F1R over FLAV_F0, run by run :: sha256 c6c4138ee5c86423 :: prose twin: two significant figures, direction in words
\newcommand{\MlstCmpJcTwoFonerOverFzeroPrMax}{up to 17\% fewer errors}  % experiments/FIGS/data/v2/index.json :: linear probe jc2_188 N=1000000: (1 - macro_auc_ovr) of FLAV_F1R over FLAV_F0, run by run :: sha256 c6c4138ee5c86423 :: prose twin of the run farther from 1
\newcommand{\MlstCmpJcTwoFonerOverFone}{0.9899 and 0.9966}  % experiments/FIGS/data/v2/index.json :: linear probe jc2_188 N=1000000: (1 - macro_auc_ovr) of FLAV_F1R over FLAV_F1, run by run :: sha256 c6c4138ee5c86423
\newcommand{\MlstCmpJcTwoFonerOverFonePr}{equal within 2\%}  % experiments/FIGS/data/v2/index.json :: linear probe jc2_188 N=1000000: (1 - macro_auc_ovr) of FLAV_F1R over FLAV_F1, run by run :: sha256 c6c4138ee5c86423 :: prose twin: two significant figures, direction in words
\newcommand{\MlstCmpJcTwoFonerOverFonePrMax}{equal within 2\%}  % experiments/FIGS/data/v2/index.json :: linear probe jc2_188 N=1000000: (1 - macro_auc_ovr) of FLAV_F1R over FLAV_F1, run by run :: sha256 c6c4138ee5c86423 :: prose twin of the run farther from 1
\newcommand{\MlstCmpJcTwoMidOverRes}{1.197\,\ensuremath{\pm}\,0.013}  % experiments/FIGS/data/v2/index.json :: linear probe jc2_188 N=1000000: (1 - macro_auc_ovr) of R42_Q1 over L162, run by run :: sha256 c6c4138ee5c86423
\newcommand{\MlstCmpJcTwoMidOverResPr}{20\% more errors}  % experiments/FIGS/data/v2/index.json :: linear probe jc2_188 N=1000000: (1 - macro_auc_ovr) of R42_Q1 over L162, run by run :: sha256 c6c4138ee5c86423 :: prose twin: two significant figures, direction in words
\newcommand{\MlstCmpJcTwoCoarseOverRes}{1.518\,\ensuremath{\pm}\,0.025}  % experiments/FIGS/data/v2/index.json :: linear probe jc2_188 N=1000000: (1 - macro_auc_ovr) of R16_Q1 over L162, run by run :: sha256 c6c4138ee5c86423
\newcommand{\MlstCmpJcTwoCoarseOverResPr}{52\% more errors}  % experiments/FIGS/data/v2/index.json :: linear probe jc2_188 N=1000000: (1 - macro_auc_ovr) of R16_Q1 over L162, run by run :: sha256 c6c4138ee5c86423 :: prose twin: two significant figures, direction in words
\newcommand{\MlstCmpJcTwoCoarseOverMid}{1.269\,\ensuremath{\pm}\,0.007}  % experiments/FIGS/data/v2/index.json :: linear probe jc2_188 N=1000000: (1 - macro_auc_ovr) of R16_Q1 over R42_Q1, run by run :: sha256 c6c4138ee5c86423
\newcommand{\MlstCmpJcTwoCoarseOverMidPr}{27\% more errors}  % experiments/FIGS/data/v2/index.json :: linear probe jc2_188 N=1000000: (1 - macro_auc_ovr) of R16_Q1 over R42_Q1, run by run :: sha256 c6c4138ee5c86423 :: prose twin: two significant figures, direction in words
\newcommand{\MlstCmpJcTwoRandOverFull}{1.125\,\ensuremath{\pm}\,0.008}  % experiments/FIGS/data/v2/index.json :: linear probe jc2_188 N=1000000: (1 - macro_auc_ovr) of each random partition over the L188 run of its seed, mean +- SD over the partitions :: sha256 c6c4138ee5c86423
\newcommand{\MlstCmpJcTwoRandOverFullPr}{12\% more errors}  % experiments/FIGS/data/v2/index.json :: linear probe jc2_188 N=1000000: (1 - macro_auc_ovr) of each random partition over the L188 run of its seed, mean +- SD over the partitions :: sha256 c6c4138ee5c86423 :: prose twin: two significant figures, direction in words
\newcommand{\MlstCmpJcTwoRandOverMid}{0.918\,\ensuremath{\pm}\,0.007}  % experiments/FIGS/data/v2/index.json :: linear probe jc2_188 N=1000000: (1 - macro_auc_ovr) of each random partition over the R42_Q1 run of its seed, mean +- SD over the partitions :: sha256 c6c4138ee5c86423
\newcommand{\MlstCmpJcTwoRandOverMidPr}{8.2\% fewer errors}  % experiments/FIGS/data/v2/index.json :: linear probe jc2_188 N=1000000: (1 - macro_auc_ovr) of each random partition over the R42_Q1 run of its seed, mean +- SD over the partitions :: sha256 c6c4138ee5c86423 :: prose twin: two significant figures, direction in words
\newcommand{\MlstCmpJcTwoRandOverCoarse}{0.724\,\ensuremath{\pm}\,0.008}  % experiments/FIGS/data/v2/index.json :: linear probe jc2_188 N=1000000: (1 - macro_auc_ovr) of each random partition over the R16_Q1 run of its seed, mean +- SD over the partitions :: sha256 c6c4138ee5c86423
\newcommand{\MlstCmpJcTwoRandOverCoarsePr}{28\% fewer errors}  % experiments/FIGS/data/v2/index.json :: linear probe jc2_188 N=1000000: (1 - macro_auc_ovr) of each random partition over the R16_Q1 run of its seed, mean +- SD over the partitions :: sha256 c6c4138ee5c86423 :: prose twin: two significant figures, direction in words
\newcommand{\MlstCmpJcTwoTrainedResOverFull}{1.026\,\ensuremath{\pm}\,0.005}  % experiments/FIGS/data/v2/index.json :: linear probe jc2_188 N=1000000: (1 - macro_auc_trained) of L162 over L188, run by run :: sha256 c6c4138ee5c86423
\newcommand{\MlstCmpJcTwoTrainedResOverFullPr}{2.6\% more errors}  % experiments/FIGS/data/v2/index.json :: linear probe jc2_188 N=1000000: (1 - macro_auc_trained) of L162 over L188, run by run :: sha256 c6c4138ee5c86423 :: prose twin: two significant figures, direction in words
\newcommand{\MlstCmpJcTwoTrainedMidOverFull}{1.227\,\ensuremath{\pm}\,0.009}  % experiments/FIGS/data/v2/index.json :: linear probe jc2_188 N=1000000: (1 - macro_auc_trained) of R42_Q1 over L188, run by run :: sha256 c6c4138ee5c86423
\newcommand{\MlstCmpJcTwoTrainedMidOverFullPr}{23\% more errors}  % experiments/FIGS/data/v2/index.json :: linear probe jc2_188 N=1000000: (1 - macro_auc_trained) of R42_Q1 over L188, run by run :: sha256 c6c4138ee5c86423 :: prose twin: two significant figures, direction in words
\newcommand{\MlstCmpJcTwoTrainedCoarseOverFull}{1.557\,\ensuremath{\pm}\,0.019}  % experiments/FIGS/data/v2/index.json :: linear probe jc2_188 N=1000000: (1 - macro_auc_trained) of R16_Q1 over L188, run by run :: sha256 c6c4138ee5c86423
\newcommand{\MlstCmpJcTwoTrainedCoarseOverFullPr}{56\% more errors}  % experiments/FIGS/data/v2/index.json :: linear probe jc2_188 N=1000000: (1 - macro_auc_trained) of R16_Q1 over L188, run by run :: sha256 c6c4138ee5c86423 :: prose twin: two significant figures, direction in words
\newcommand{\MlstCmpJcTwoTrainedResMassOverRes}{1.022 and 1.016}  % experiments/FIGS/data/v2/index.json :: linear probe jc2_188 N=1000000: (1 - macro_auc_trained) of L162_MASS over L162, run by run :: sha256 c6c4138ee5c86423
\newcommand{\MlstCmpJcTwoTrainedResMassOverResPr}{equal within 3\%}  % experiments/FIGS/data/v2/index.json :: linear probe jc2_188 N=1000000: (1 - macro_auc_trained) of L162_MASS over L162, run by run :: sha256 c6c4138ee5c86423 :: prose twin: two significant figures, direction in words
\newcommand{\MlstCmpJcTwoTrainedResMassOverResPrMax}{equal within 3\%}  % experiments/FIGS/data/v2/index.json :: linear probe jc2_188 N=1000000: (1 - macro_auc_trained) of L162_MASS over L162, run by run :: sha256 c6c4138ee5c86423 :: prose twin of the run farther from 1
\newcommand{\MlstCmpJcTwoTrainedCoarseMassOverCoarse}{1.054 and 0.977}  % experiments/FIGS/data/v2/index.json :: linear probe jc2_188 N=1000000: (1 - macro_auc_trained) of R16_Q1_MASS over R16_Q1, run by run :: sha256 c6c4138ee5c86423
\newcommand{\MlstCmpJcTwoTrainedCoarseMassOverCoarsePr}{equal within 6\%}  % experiments/FIGS/data/v2/index.json :: linear probe jc2_188 N=1000000: (1 - macro_auc_trained) of R16_Q1_MASS over R16_Q1, run by run :: sha256 c6c4138ee5c86423 :: prose twin: two significant figures, direction in words
\newcommand{\MlstCmpJcTwoTrainedCoarseMassOverCoarsePrMax}{equal within 6\%}  % experiments/FIGS/data/v2/index.json :: linear probe jc2_188 N=1000000: (1 - macro_auc_trained) of R16_Q1_MASS over R16_Q1, run by run :: sha256 c6c4138ee5c86423 :: prose twin of the run farther from 1
\newcommand{\MlstCmpJcTwoTrainedResMassOverFull}{1.053 and 1.037}  % experiments/FIGS/data/v2/index.json :: linear probe jc2_188 N=1000000: (1 - macro_auc_trained) of L162_MASS over L188, run by run :: sha256 c6c4138ee5c86423
\newcommand{\MlstCmpJcTwoTrainedResMassOverFullPr}{at least 3.7\% more errors}  % experiments/FIGS/data/v2/index.json :: linear probe jc2_188 N=1000000: (1 - macro_auc_trained) of L162_MASS over L188, run by run :: sha256 c6c4138ee5c86423 :: prose twin: two significant figures, direction in words
\newcommand{\MlstCmpJcTwoTrainedResMassOverFullPrMax}{up to 5.3\% more errors}  % experiments/FIGS/data/v2/index.json :: linear probe jc2_188 N=1000000: (1 - macro_auc_trained) of L162_MASS over L188, run by run :: sha256 c6c4138ee5c86423 :: prose twin of the run farther from 1
\newcommand{\MlstCmpJcTwoTrainedCoarseMassOverFull}{1.617 and 1.528}  % experiments/FIGS/data/v2/index.json :: linear probe jc2_188 N=1000000: (1 - macro_auc_trained) of R16_Q1_MASS over L188, run by run :: sha256 c6c4138ee5c86423
\newcommand{\MlstCmpJcTwoTrainedCoarseMassOverFullPr}{at least 53\% more errors}  % experiments/FIGS/data/v2/index.json :: linear probe jc2_188 N=1000000: (1 - macro_auc_trained) of R16_Q1_MASS over L188, run by run :: sha256 c6c4138ee5c86423 :: prose twin: two significant figures, direction in words
\newcommand{\MlstCmpJcTwoTrainedCoarseMassOverFullPrMax}{up to 62\% more errors}  % experiments/FIGS/data/v2/index.json :: linear probe jc2_188 N=1000000: (1 - macro_auc_trained) of R16_Q1_MASS over L188, run by run :: sha256 c6c4138ee5c86423 :: prose twin of the run farther from 1
\newcommand{\MlstCmpJcTwoTrainedFzeroOverFull}{1.424 and 1.462}  % experiments/FIGS/data/v2/index.json :: linear probe jc2_188 N=1000000: (1 - macro_auc_trained) of FLAV_F0 over L188, run by run :: sha256 c6c4138ee5c86423
\newcommand{\MlstCmpJcTwoTrainedFzeroOverFullPr}{at least 42\% more errors}  % experiments/FIGS/data/v2/index.json :: linear probe jc2_188 N=1000000: (1 - macro_auc_trained) of FLAV_F0 over L188, run by run :: sha256 c6c4138ee5c86423 :: prose twin: two significant figures, direction in words
\newcommand{\MlstCmpJcTwoTrainedFzeroOverFullPrMax}{up to 46\% more errors}  % experiments/FIGS/data/v2/index.json :: linear probe jc2_188 N=1000000: (1 - macro_auc_trained) of FLAV_F0 over L188, run by run :: sha256 c6c4138ee5c86423 :: prose twin of the run farther from 1
\newcommand{\MlstCmpJcTwoTrainedFoneOverFull}{1.218 and 1.212}  % experiments/FIGS/data/v2/index.json :: linear probe jc2_188 N=1000000: (1 - macro_auc_trained) of FLAV_F1 over L188, run by run :: sha256 c6c4138ee5c86423
\newcommand{\MlstCmpJcTwoTrainedFoneOverFullPr}{at least 21\% more errors}  % experiments/FIGS/data/v2/index.json :: linear probe jc2_188 N=1000000: (1 - macro_auc_trained) of FLAV_F1 over L188, run by run :: sha256 c6c4138ee5c86423 :: prose twin: two significant figures, direction in words
\newcommand{\MlstCmpJcTwoTrainedFoneOverFullPrMax}{up to 22\% more errors}  % experiments/FIGS/data/v2/index.json :: linear probe jc2_188 N=1000000: (1 - macro_auc_trained) of FLAV_F1 over L188, run by run :: sha256 c6c4138ee5c86423 :: prose twin of the run farther from 1
\newcommand{\MlstCmpJcTwoTrainedFonerOverFull}{1.206 and 1.208}  % experiments/FIGS/data/v2/index.json :: linear probe jc2_188 N=1000000: (1 - macro_auc_trained) of FLAV_F1R over L188, run by run :: sha256 c6c4138ee5c86423
\newcommand{\MlstCmpJcTwoTrainedFonerOverFullPr}{at least 21\% more errors}  % experiments/FIGS/data/v2/index.json :: linear probe jc2_188 N=1000000: (1 - macro_auc_trained) of FLAV_F1R over L188, run by run :: sha256 c6c4138ee5c86423 :: prose twin: two significant figures, direction in words
\newcommand{\MlstCmpJcTwoTrainedFonerOverFullPrMax}{up to 21\% more errors}  % experiments/FIGS/data/v2/index.json :: linear probe jc2_188 N=1000000: (1 - macro_auc_trained) of FLAV_F1R over L188, run by run :: sha256 c6c4138ee5c86423 :: prose twin of the run farther from 1
\newcommand{\MlstCmpJcTwoTrainedFzeroOverCoarse}{0.9275 and 0.9345}  % experiments/FIGS/data/v2/index.json :: linear probe jc2_188 N=1000000: (1 - macro_auc_trained) of FLAV_F0 over R16_Q1, run by run :: sha256 c6c4138ee5c86423
\newcommand{\MlstCmpJcTwoTrainedFzeroOverCoarsePr}{at least 6.5\% fewer errors}  % experiments/FIGS/data/v2/index.json :: linear probe jc2_188 N=1000000: (1 - macro_auc_trained) of FLAV_F0 over R16_Q1, run by run :: sha256 c6c4138ee5c86423 :: prose twin: two significant figures, direction in words
\newcommand{\MlstCmpJcTwoTrainedFzeroOverCoarsePrMax}{up to 7.2\% fewer errors}  % experiments/FIGS/data/v2/index.json :: linear probe jc2_188 N=1000000: (1 - macro_auc_trained) of FLAV_F0 over R16_Q1, run by run :: sha256 c6c4138ee5c86423 :: prose twin of the run farther from 1
\newcommand{\MlstCmpJcTwoTrainedFoneOverCoarse}{0.794 and 0.775}  % experiments/FIGS/data/v2/index.json :: linear probe jc2_188 N=1000000: (1 - macro_auc_trained) of FLAV_F1 over R16_Q1, run by run :: sha256 c6c4138ee5c86423
\newcommand{\MlstCmpJcTwoTrainedFoneOverCoarsePr}{at least 21\% fewer errors}  % experiments/FIGS/data/v2/index.json :: linear probe jc2_188 N=1000000: (1 - macro_auc_trained) of FLAV_F1 over R16_Q1, run by run :: sha256 c6c4138ee5c86423 :: prose twin: two significant figures, direction in words
\newcommand{\MlstCmpJcTwoTrainedFoneOverCoarsePrMax}{up to 23\% fewer errors}  % experiments/FIGS/data/v2/index.json :: linear probe jc2_188 N=1000000: (1 - macro_auc_trained) of FLAV_F1 over R16_Q1, run by run :: sha256 c6c4138ee5c86423 :: prose twin of the run farther from 1
\newcommand{\MlstCmpJcTwoTrainedFonerOverCoarse}{0.786 and 0.772}  % experiments/FIGS/data/v2/index.json :: linear probe jc2_188 N=1000000: (1 - macro_auc_trained) of FLAV_F1R over R16_Q1, run by run :: sha256 c6c4138ee5c86423
\newcommand{\MlstCmpJcTwoTrainedFonerOverCoarsePr}{at least 21\% fewer errors}  % experiments/FIGS/data/v2/index.json :: linear probe jc2_188 N=1000000: (1 - macro_auc_trained) of FLAV_F1R over R16_Q1, run by run :: sha256 c6c4138ee5c86423 :: prose twin: two significant figures, direction in words
\newcommand{\MlstCmpJcTwoTrainedFonerOverCoarsePrMax}{up to 23\% fewer errors}  % experiments/FIGS/data/v2/index.json :: linear probe jc2_188 N=1000000: (1 - macro_auc_trained) of FLAV_F1R over R16_Q1, run by run :: sha256 c6c4138ee5c86423 :: prose twin of the run farther from 1
\newcommand{\MlstCmpJcTwoTrainedFoneOverFzero}{0.856 and 0.829}  % experiments/FIGS/data/v2/index.json :: linear probe jc2_188 N=1000000: (1 - macro_auc_trained) of FLAV_F1 over FLAV_F0, run by run :: sha256 c6c4138ee5c86423
\newcommand{\MlstCmpJcTwoTrainedFoneOverFzeroPr}{at least 14\% fewer errors}  % experiments/FIGS/data/v2/index.json :: linear probe jc2_188 N=1000000: (1 - macro_auc_trained) of FLAV_F1 over FLAV_F0, run by run :: sha256 c6c4138ee5c86423 :: prose twin: two significant figures, direction in words
\newcommand{\MlstCmpJcTwoTrainedFoneOverFzeroPrMax}{up to 17\% fewer errors}  % experiments/FIGS/data/v2/index.json :: linear probe jc2_188 N=1000000: (1 - macro_auc_trained) of FLAV_F1 over FLAV_F0, run by run :: sha256 c6c4138ee5c86423 :: prose twin of the run farther from 1
\newcommand{\MlstCmpJcTwoTrainedFonerOverFzero}{0.847 and 0.826}  % experiments/FIGS/data/v2/index.json :: linear probe jc2_188 N=1000000: (1 - macro_auc_trained) of FLAV_F1R over FLAV_F0, run by run :: sha256 c6c4138ee5c86423
\newcommand{\MlstCmpJcTwoTrainedFonerOverFzeroPr}{at least 15\% fewer errors}  % experiments/FIGS/data/v2/index.json :: linear probe jc2_188 N=1000000: (1 - macro_auc_trained) of FLAV_F1R over FLAV_F0, run by run :: sha256 c6c4138ee5c86423 :: prose twin: two significant figures, direction in words
\newcommand{\MlstCmpJcTwoTrainedFonerOverFzeroPrMax}{up to 17\% fewer errors}  % experiments/FIGS/data/v2/index.json :: linear probe jc2_188 N=1000000: (1 - macro_auc_trained) of FLAV_F1R over FLAV_F0, run by run :: sha256 c6c4138ee5c86423 :: prose twin of the run farther from 1
\newcommand{\MlstCmpJcTwoTrainedFonerOverFone}{0.9899 and 0.9966}  % experiments/FIGS/data/v2/index.json :: linear probe jc2_188 N=1000000: (1 - macro_auc_trained) of FLAV_F1R over FLAV_F1, run by run :: sha256 c6c4138ee5c86423
\newcommand{\MlstCmpJcTwoTrainedFonerOverFonePr}{equal within 2\%}  % experiments/FIGS/data/v2/index.json :: linear probe jc2_188 N=1000000: (1 - macro_auc_trained) of FLAV_F1R over FLAV_F1, run by run :: sha256 c6c4138ee5c86423 :: prose twin: two significant figures, direction in words
\newcommand{\MlstCmpJcTwoTrainedFonerOverFonePrMax}{equal within 2\%}  % experiments/FIGS/data/v2/index.json :: linear probe jc2_188 N=1000000: (1 - macro_auc_trained) of FLAV_F1R over FLAV_F1, run by run :: sha256 c6c4138ee5c86423 :: prose twin of the run farther from 1
\newcommand{\MlstCmpJcTwoTrainedMidOverRes}{1.197\,\ensuremath{\pm}\,0.013}  % experiments/FIGS/data/v2/index.json :: linear probe jc2_188 N=1000000: (1 - macro_auc_trained) of R42_Q1 over L162, run by run :: sha256 c6c4138ee5c86423
\newcommand{\MlstCmpJcTwoTrainedMidOverResPr}{20\% more errors}  % experiments/FIGS/data/v2/index.json :: linear probe jc2_188 N=1000000: (1 - macro_auc_trained) of R42_Q1 over L162, run by run :: sha256 c6c4138ee5c86423 :: prose twin: two significant figures, direction in words
\newcommand{\MlstCmpJcTwoTrainedCoarseOverRes}{1.518\,\ensuremath{\pm}\,0.025}  % experiments/FIGS/data/v2/index.json :: linear probe jc2_188 N=1000000: (1 - macro_auc_trained) of R16_Q1 over L162, run by run :: sha256 c6c4138ee5c86423
\newcommand{\MlstCmpJcTwoTrainedCoarseOverResPr}{52\% more errors}  % experiments/FIGS/data/v2/index.json :: linear probe jc2_188 N=1000000: (1 - macro_auc_trained) of R16_Q1 over L162, run by run :: sha256 c6c4138ee5c86423 :: prose twin: two significant figures, direction in words
\newcommand{\MlstCmpJcTwoTrainedCoarseOverMid}{1.269\,\ensuremath{\pm}\,0.007}  % experiments/FIGS/data/v2/index.json :: linear probe jc2_188 N=1000000: (1 - macro_auc_trained) of R16_Q1 over R42_Q1, run by run :: sha256 c6c4138ee5c86423
\newcommand{\MlstCmpJcTwoTrainedCoarseOverMidPr}{27\% more errors}  % experiments/FIGS/data/v2/index.json :: linear probe jc2_188 N=1000000: (1 - macro_auc_trained) of R16_Q1 over R42_Q1, run by run :: sha256 c6c4138ee5c86423 :: prose twin: two significant figures, direction in words
\newcommand{\MlstCmpJcTwoTrainedRandOverFull}{1.125\,\ensuremath{\pm}\,0.008}  % experiments/FIGS/data/v2/index.json :: linear probe jc2_188 N=1000000: (1 - macro_auc_trained) of each random partition over the L188 run of its seed, mean +- SD over the partitions :: sha256 c6c4138ee5c86423
\newcommand{\MlstCmpJcTwoTrainedRandOverFullPr}{12\% more errors}  % experiments/FIGS/data/v2/index.json :: linear probe jc2_188 N=1000000: (1 - macro_auc_trained) of each random partition over the L188 run of its seed, mean +- SD over the partitions :: sha256 c6c4138ee5c86423 :: prose twin: two significant figures, direction in words
\newcommand{\MlstCmpJcTwoTrainedRandOverMid}{0.918\,\ensuremath{\pm}\,0.007}  % experiments/FIGS/data/v2/index.json :: linear probe jc2_188 N=1000000: (1 - macro_auc_trained) of each random partition over the R42_Q1 run of its seed, mean +- SD over the partitions :: sha256 c6c4138ee5c86423
\newcommand{\MlstCmpJcTwoTrainedRandOverMidPr}{8.2\% fewer errors}  % experiments/FIGS/data/v2/index.json :: linear probe jc2_188 N=1000000: (1 - macro_auc_trained) of each random partition over the R42_Q1 run of its seed, mean +- SD over the partitions :: sha256 c6c4138ee5c86423 :: prose twin: two significant figures, direction in words
\newcommand{\MlstCmpJcTwoTrainedRandOverCoarse}{0.724\,\ensuremath{\pm}\,0.008}  % experiments/FIGS/data/v2/index.json :: linear probe jc2_188 N=1000000: (1 - macro_auc_trained) of each random partition over the R16_Q1 run of its seed, mean +- SD over the partitions :: sha256 c6c4138ee5c86423
\newcommand{\MlstCmpJcTwoTrainedRandOverCoarsePr}{28\% fewer errors}  % experiments/FIGS/data/v2/index.json :: linear probe jc2_188 N=1000000: (1 - macro_auc_trained) of each random partition over the R16_Q1 run of its seed, mean +- SD over the partitions :: sha256 c6c4138ee5c86423 :: prose twin: two significant figures, direction in words
\newcommand{\MlstCmpJcTwoResResOverFull}{1.003\,\ensuremath{\pm}\,0.005}  % experiments/FIGS/data/v2/index.json :: linear probe jc2_188 N=1000000: (1 - macro_auc_resonant) of L162 over L188, run by run :: sha256 c6c4138ee5c86423
\newcommand{\MlstCmpJcTwoResResOverFullPr}{equal within 1\%}  % experiments/FIGS/data/v2/index.json :: linear probe jc2_188 N=1000000: (1 - macro_auc_resonant) of L162 over L188, run by run :: sha256 c6c4138ee5c86423 :: prose twin: two significant figures, direction in words
\newcommand{\MlstCmpJcTwoResMidOverFull}{1.262\,\ensuremath{\pm}\,0.012}  % experiments/FIGS/data/v2/index.json :: linear probe jc2_188 N=1000000: (1 - macro_auc_resonant) of R42_Q1 over L188, run by run :: sha256 c6c4138ee5c86423
\newcommand{\MlstCmpJcTwoResMidOverFullPr}{26\% more errors}  % experiments/FIGS/data/v2/index.json :: linear probe jc2_188 N=1000000: (1 - macro_auc_resonant) of R42_Q1 over L188, run by run :: sha256 c6c4138ee5c86423 :: prose twin: two significant figures, direction in words
\newcommand{\MlstCmpJcTwoResCoarseOverFull}{1.590\,\ensuremath{\pm}\,0.032}  % experiments/FIGS/data/v2/index.json :: linear probe jc2_188 N=1000000: (1 - macro_auc_resonant) of R16_Q1 over L188, run by run :: sha256 c6c4138ee5c86423
\newcommand{\MlstCmpJcTwoResCoarseOverFullPr}{59\% more errors}  % experiments/FIGS/data/v2/index.json :: linear probe jc2_188 N=1000000: (1 - macro_auc_resonant) of R16_Q1 over L188, run by run :: sha256 c6c4138ee5c86423 :: prose twin: two significant figures, direction in words
\newcommand{\MlstCmpJcTwoResResMassOverRes}{1.025 and 1.020}  % experiments/FIGS/data/v2/index.json :: linear probe jc2_188 N=1000000: (1 - macro_auc_resonant) of L162_MASS over L162, run by run :: sha256 c6c4138ee5c86423
\newcommand{\MlstCmpJcTwoResResMassOverResPr}{equal within 3\%}  % experiments/FIGS/data/v2/index.json :: linear probe jc2_188 N=1000000: (1 - macro_auc_resonant) of L162_MASS over L162, run by run :: sha256 c6c4138ee5c86423 :: prose twin: two significant figures, direction in words
\newcommand{\MlstCmpJcTwoResResMassOverResPrMax}{equal within 3\%}  % experiments/FIGS/data/v2/index.json :: linear probe jc2_188 N=1000000: (1 - macro_auc_resonant) of L162_MASS over L162, run by run :: sha256 c6c4138ee5c86423 :: prose twin of the run farther from 1
\newcommand{\MlstCmpJcTwoResCoarseMassOverCoarse}{1.08 and 0.96}  % experiments/FIGS/data/v2/index.json :: linear probe jc2_188 N=1000000: (1 - macro_auc_resonant) of R16_Q1_MASS over R16_Q1, run by run :: sha256 c6c4138ee5c86423
\newcommand{\MlstCmpJcTwoResCoarseMassOverCoarsePr}{equal within 8\%}  % experiments/FIGS/data/v2/index.json :: linear probe jc2_188 N=1000000: (1 - macro_auc_resonant) of R16_Q1_MASS over R16_Q1, run by run :: sha256 c6c4138ee5c86423 :: prose twin: two significant figures, direction in words
\newcommand{\MlstCmpJcTwoResCoarseMassOverCoarsePrMax}{equal within 8\%}  % experiments/FIGS/data/v2/index.json :: linear probe jc2_188 N=1000000: (1 - macro_auc_resonant) of R16_Q1_MASS over R16_Q1, run by run :: sha256 c6c4138ee5c86423 :: prose twin of the run farther from 1
\newcommand{\MlstCmpJcTwoResResMassOverFull}{1.033 and 1.018}  % experiments/FIGS/data/v2/index.json :: linear probe jc2_188 N=1000000: (1 - macro_auc_resonant) of L162_MASS over L188, run by run :: sha256 c6c4138ee5c86423
\newcommand{\MlstCmpJcTwoResResMassOverFullPr}{equal within 4\%}  % experiments/FIGS/data/v2/index.json :: linear probe jc2_188 N=1000000: (1 - macro_auc_resonant) of L162_MASS over L188, run by run :: sha256 c6c4138ee5c86423 :: prose twin: two significant figures, direction in words
\newcommand{\MlstCmpJcTwoResResMassOverFullPrMax}{equal within 4\%}  % experiments/FIGS/data/v2/index.json :: linear probe jc2_188 N=1000000: (1 - macro_auc_resonant) of L162_MASS over L188, run by run :: sha256 c6c4138ee5c86423 :: prose twin of the run farther from 1
\newcommand{\MlstCmpJcTwoResCoarseMassOverFull}{1.67 and 1.55}  % experiments/FIGS/data/v2/index.json :: linear probe jc2_188 N=1000000: (1 - macro_auc_resonant) of R16_Q1_MASS over L188, run by run :: sha256 c6c4138ee5c86423
\newcommand{\MlstCmpJcTwoResCoarseMassOverFullPr}{at least 55\% more errors}  % experiments/FIGS/data/v2/index.json :: linear probe jc2_188 N=1000000: (1 - macro_auc_resonant) of R16_Q1_MASS over L188, run by run :: sha256 c6c4138ee5c86423 :: prose twin: two significant figures, direction in words
\newcommand{\MlstCmpJcTwoResCoarseMassOverFullPrMax}{up to 67\% more errors}  % experiments/FIGS/data/v2/index.json :: linear probe jc2_188 N=1000000: (1 - macro_auc_resonant) of R16_Q1_MASS over L188, run by run :: sha256 c6c4138ee5c86423 :: prose twin of the run farther from 1
\newcommand{\MlstCmpJcTwoResFzeroOverFull}{1.408 and 1.449}  % experiments/FIGS/data/v2/index.json :: linear probe jc2_188 N=1000000: (1 - macro_auc_resonant) of FLAV_F0 over L188, run by run :: sha256 c6c4138ee5c86423
\newcommand{\MlstCmpJcTwoResFzeroOverFullPr}{at least 41\% more errors}  % experiments/FIGS/data/v2/index.json :: linear probe jc2_188 N=1000000: (1 - macro_auc_resonant) of FLAV_F0 over L188, run by run :: sha256 c6c4138ee5c86423 :: prose twin: two significant figures, direction in words
\newcommand{\MlstCmpJcTwoResFzeroOverFullPrMax}{up to 45\% more errors}  % experiments/FIGS/data/v2/index.json :: linear probe jc2_188 N=1000000: (1 - macro_auc_resonant) of FLAV_F0 over L188, run by run :: sha256 c6c4138ee5c86423 :: prose twin of the run farther from 1
\newcommand{\MlstCmpJcTwoResFoneOverFull}{1.210 and 1.204}  % experiments/FIGS/data/v2/index.json :: linear probe jc2_188 N=1000000: (1 - macro_auc_resonant) of FLAV_F1 over L188, run by run :: sha256 c6c4138ee5c86423
\newcommand{\MlstCmpJcTwoResFoneOverFullPr}{at least 20\% more errors}  % experiments/FIGS/data/v2/index.json :: linear probe jc2_188 N=1000000: (1 - macro_auc_resonant) of FLAV_F1 over L188, run by run :: sha256 c6c4138ee5c86423 :: prose twin: two significant figures, direction in words
\newcommand{\MlstCmpJcTwoResFoneOverFullPrMax}{up to 21\% more errors}  % experiments/FIGS/data/v2/index.json :: linear probe jc2_188 N=1000000: (1 - macro_auc_resonant) of FLAV_F1 over L188, run by run :: sha256 c6c4138ee5c86423 :: prose twin of the run farther from 1
\newcommand{\MlstCmpJcTwoResFonerOverFull}{1.187 and 1.188}  % experiments/FIGS/data/v2/index.json :: linear probe jc2_188 N=1000000: (1 - macro_auc_resonant) of FLAV_F1R over L188, run by run :: sha256 c6c4138ee5c86423
\newcommand{\MlstCmpJcTwoResFonerOverFullPr}{at least 19\% more errors}  % experiments/FIGS/data/v2/index.json :: linear probe jc2_188 N=1000000: (1 - macro_auc_resonant) of FLAV_F1R over L188, run by run :: sha256 c6c4138ee5c86423 :: prose twin: two significant figures, direction in words
\newcommand{\MlstCmpJcTwoResFonerOverFullPrMax}{up to 19\% more errors}  % experiments/FIGS/data/v2/index.json :: linear probe jc2_188 N=1000000: (1 - macro_auc_resonant) of FLAV_F1R over L188, run by run :: sha256 c6c4138ee5c86423 :: prose twin of the run farther from 1
\newcommand{\MlstCmpJcTwoResFzeroOverCoarse}{0.9069 and 0.9029}  % experiments/FIGS/data/v2/index.json :: linear probe jc2_188 N=1000000: (1 - macro_auc_resonant) of FLAV_F0 over R16_Q1, run by run :: sha256 c6c4138ee5c86423
\newcommand{\MlstCmpJcTwoResFzeroOverCoarsePr}{at least 9.3\% fewer errors}  % experiments/FIGS/data/v2/index.json :: linear probe jc2_188 N=1000000: (1 - macro_auc_resonant) of FLAV_F0 over R16_Q1, run by run :: sha256 c6c4138ee5c86423 :: prose twin: two significant figures, direction in words
\newcommand{\MlstCmpJcTwoResFzeroOverCoarsePrMax}{up to 9.7\% fewer errors}  % experiments/FIGS/data/v2/index.json :: linear probe jc2_188 N=1000000: (1 - macro_auc_resonant) of FLAV_F0 over R16_Q1, run by run :: sha256 c6c4138ee5c86423 :: prose twin of the run farther from 1
\newcommand{\MlstCmpJcTwoResFoneOverCoarse}{0.779 and 0.750}  % experiments/FIGS/data/v2/index.json :: linear probe jc2_188 N=1000000: (1 - macro_auc_resonant) of FLAV_F1 over R16_Q1, run by run :: sha256 c6c4138ee5c86423
\newcommand{\MlstCmpJcTwoResFoneOverCoarsePr}{at least 22\% fewer errors}  % experiments/FIGS/data/v2/index.json :: linear probe jc2_188 N=1000000: (1 - macro_auc_resonant) of FLAV_F1 over R16_Q1, run by run :: sha256 c6c4138ee5c86423 :: prose twin: two significant figures, direction in words
\newcommand{\MlstCmpJcTwoResFoneOverCoarsePrMax}{up to 25\% fewer errors}  % experiments/FIGS/data/v2/index.json :: linear probe jc2_188 N=1000000: (1 - macro_auc_resonant) of FLAV_F1 over R16_Q1, run by run :: sha256 c6c4138ee5c86423 :: prose twin of the run farther from 1
\newcommand{\MlstCmpJcTwoResFonerOverCoarse}{0.764 and 0.740}  % experiments/FIGS/data/v2/index.json :: linear probe jc2_188 N=1000000: (1 - macro_auc_resonant) of FLAV_F1R over R16_Q1, run by run :: sha256 c6c4138ee5c86423
\newcommand{\MlstCmpJcTwoResFonerOverCoarsePr}{at least 24\% fewer errors}  % experiments/FIGS/data/v2/index.json :: linear probe jc2_188 N=1000000: (1 - macro_auc_resonant) of FLAV_F1R over R16_Q1, run by run :: sha256 c6c4138ee5c86423 :: prose twin: two significant figures, direction in words
\newcommand{\MlstCmpJcTwoResFonerOverCoarsePrMax}{up to 26\% fewer errors}  % experiments/FIGS/data/v2/index.json :: linear probe jc2_188 N=1000000: (1 - macro_auc_resonant) of FLAV_F1R over R16_Q1, run by run :: sha256 c6c4138ee5c86423 :: prose twin of the run farther from 1
\newcommand{\MlstCmpJcTwoResFoneOverFzero}{0.859 and 0.831}  % experiments/FIGS/data/v2/index.json :: linear probe jc2_188 N=1000000: (1 - macro_auc_resonant) of FLAV_F1 over FLAV_F0, run by run :: sha256 c6c4138ee5c86423
\newcommand{\MlstCmpJcTwoResFoneOverFzeroPr}{at least 14\% fewer errors}  % experiments/FIGS/data/v2/index.json :: linear probe jc2_188 N=1000000: (1 - macro_auc_resonant) of FLAV_F1 over FLAV_F0, run by run :: sha256 c6c4138ee5c86423 :: prose twin: two significant figures, direction in words
\newcommand{\MlstCmpJcTwoResFoneOverFzeroPrMax}{up to 17\% fewer errors}  % experiments/FIGS/data/v2/index.json :: linear probe jc2_188 N=1000000: (1 - macro_auc_resonant) of FLAV_F1 over FLAV_F0, run by run :: sha256 c6c4138ee5c86423 :: prose twin of the run farther from 1
\newcommand{\MlstCmpJcTwoResFonerOverFzero}{0.843 and 0.820}  % experiments/FIGS/data/v2/index.json :: linear probe jc2_188 N=1000000: (1 - macro_auc_resonant) of FLAV_F1R over FLAV_F0, run by run :: sha256 c6c4138ee5c86423
\newcommand{\MlstCmpJcTwoResFonerOverFzeroPr}{at least 16\% fewer errors}  % experiments/FIGS/data/v2/index.json :: linear probe jc2_188 N=1000000: (1 - macro_auc_resonant) of FLAV_F1R over FLAV_F0, run by run :: sha256 c6c4138ee5c86423 :: prose twin: two significant figures, direction in words
\newcommand{\MlstCmpJcTwoResFonerOverFzeroPrMax}{up to 18\% fewer errors}  % experiments/FIGS/data/v2/index.json :: linear probe jc2_188 N=1000000: (1 - macro_auc_resonant) of FLAV_F1R over FLAV_F0, run by run :: sha256 c6c4138ee5c86423 :: prose twin of the run farther from 1
\newcommand{\MlstCmpJcTwoResFonerOverFone}{0.9806 and 0.9867}  % experiments/FIGS/data/v2/index.json :: linear probe jc2_188 N=1000000: (1 - macro_auc_resonant) of FLAV_F1R over FLAV_F1, run by run :: sha256 c6c4138ee5c86423
\newcommand{\MlstCmpJcTwoResFonerOverFonePr}{equal within 2\%}  % experiments/FIGS/data/v2/index.json :: linear probe jc2_188 N=1000000: (1 - macro_auc_resonant) of FLAV_F1R over FLAV_F1, run by run :: sha256 c6c4138ee5c86423 :: prose twin: two significant figures, direction in words
\newcommand{\MlstCmpJcTwoResFonerOverFonePrMax}{equal within 2\%}  % experiments/FIGS/data/v2/index.json :: linear probe jc2_188 N=1000000: (1 - macro_auc_resonant) of FLAV_F1R over FLAV_F1, run by run :: sha256 c6c4138ee5c86423 :: prose twin of the run farther from 1
\newcommand{\MlstCmpJcTwoResMidOverRes}{1.258\,\ensuremath{\pm}\,0.017}  % experiments/FIGS/data/v2/index.json :: linear probe jc2_188 N=1000000: (1 - macro_auc_resonant) of R42_Q1 over L162, run by run :: sha256 c6c4138ee5c86423
\newcommand{\MlstCmpJcTwoResMidOverResPr}{26\% more errors}  % experiments/FIGS/data/v2/index.json :: linear probe jc2_188 N=1000000: (1 - macro_auc_resonant) of R42_Q1 over L162, run by run :: sha256 c6c4138ee5c86423 :: prose twin: two significant figures, direction in words
\newcommand{\MlstCmpJcTwoResCoarseOverRes}{1.58\,\ensuremath{\pm}\,0.04}  % experiments/FIGS/data/v2/index.json :: linear probe jc2_188 N=1000000: (1 - macro_auc_resonant) of R16_Q1 over L162, run by run :: sha256 c6c4138ee5c86423
\newcommand{\MlstCmpJcTwoResCoarseOverResPr}{58\% more errors}  % experiments/FIGS/data/v2/index.json :: linear probe jc2_188 N=1000000: (1 - macro_auc_resonant) of R16_Q1 over L162, run by run :: sha256 c6c4138ee5c86423 :: prose twin: two significant figures, direction in words
\newcommand{\MlstCmpJcTwoResCoarseOverMid}{1.260\,\ensuremath{\pm}\,0.015}  % experiments/FIGS/data/v2/index.json :: linear probe jc2_188 N=1000000: (1 - macro_auc_resonant) of R16_Q1 over R42_Q1, run by run :: sha256 c6c4138ee5c86423
\newcommand{\MlstCmpJcTwoResCoarseOverMidPr}{26\% more errors}  % experiments/FIGS/data/v2/index.json :: linear probe jc2_188 N=1000000: (1 - macro_auc_resonant) of R16_Q1 over R42_Q1, run by run :: sha256 c6c4138ee5c86423 :: prose twin: two significant figures, direction in words
\newcommand{\MlstCmpJcTwoResRandOverFull}{1.098\,\ensuremath{\pm}\,0.008}  % experiments/FIGS/data/v2/index.json :: linear probe jc2_188 N=1000000: (1 - macro_auc_resonant) of each random partition over the L188 run of its seed, mean +- SD over the partitions :: sha256 c6c4138ee5c86423
\newcommand{\MlstCmpJcTwoResRandOverFullPr}{9.8\% more errors}  % experiments/FIGS/data/v2/index.json :: linear probe jc2_188 N=1000000: (1 - macro_auc_resonant) of each random partition over the L188 run of its seed, mean +- SD over the partitions :: sha256 c6c4138ee5c86423 :: prose twin: two significant figures, direction in words
\newcommand{\MlstCmpJcTwoResRandOverMid}{0.870\,\ensuremath{\pm}\,0.008}  % experiments/FIGS/data/v2/index.json :: linear probe jc2_188 N=1000000: (1 - macro_auc_resonant) of each random partition over the R42_Q1 run of its seed, mean +- SD over the partitions :: sha256 c6c4138ee5c86423
\newcommand{\MlstCmpJcTwoResRandOverMidPr}{13\% fewer errors}  % experiments/FIGS/data/v2/index.json :: linear probe jc2_188 N=1000000: (1 - macro_auc_resonant) of each random partition over the R42_Q1 run of its seed, mean +- SD over the partitions :: sha256 c6c4138ee5c86423 :: prose twin: two significant figures, direction in words
\newcommand{\MlstCmpJcTwoResRandOverCoarse}{0.693\,\ensuremath{\pm}\,0.011}  % experiments/FIGS/data/v2/index.json :: linear probe jc2_188 N=1000000: (1 - macro_auc_resonant) of each random partition over the R16_Q1 run of its seed, mean +- SD over the partitions :: sha256 c6c4138ee5c86423
\newcommand{\MlstCmpJcTwoResRandOverCoarsePr}{31\% fewer errors}  % experiments/FIGS/data/v2/index.json :: linear probe jc2_188 N=1000000: (1 - macro_auc_resonant) of each random partition over the R16_Q1 run of its seed, mean +- SD over the partitions :: sha256 c6c4138ee5c86423 :: prose twin: two significant figures, direction in words
\newcommand{\MlstCmpJcTwoQcdResOverFull}{1.087\,\ensuremath{\pm}\,0.006}  % experiments/FIGS/data/v2/index.json :: linear probe jc2_188 N=1000000: (1 - macro_auc_qcd) of L162 over L188, run by run :: sha256 c6c4138ee5c86423
\newcommand{\MlstCmpJcTwoQcdResOverFullPr}{8.7\% more errors}  % experiments/FIGS/data/v2/index.json :: linear probe jc2_188 N=1000000: (1 - macro_auc_qcd) of L162 over L188, run by run :: sha256 c6c4138ee5c86423 :: prose twin: two significant figures, direction in words
\newcommand{\MlstCmpJcTwoQcdMidOverFull}{1.131\,\ensuremath{\pm}\,0.007}  % experiments/FIGS/data/v2/index.json :: linear probe jc2_188 N=1000000: (1 - macro_auc_qcd) of R42_Q1 over L188, run by run :: sha256 c6c4138ee5c86423
\newcommand{\MlstCmpJcTwoQcdMidOverFullPr}{13\% more errors}  % experiments/FIGS/data/v2/index.json :: linear probe jc2_188 N=1000000: (1 - macro_auc_qcd) of R42_Q1 over L188, run by run :: sha256 c6c4138ee5c86423 :: prose twin: two significant figures, direction in words
\newcommand{\MlstCmpJcTwoQcdCoarseOverFull}{1.466\,\ensuremath{\pm}\,0.017}  % experiments/FIGS/data/v2/index.json :: linear probe jc2_188 N=1000000: (1 - macro_auc_qcd) of R16_Q1 over L188, run by run :: sha256 c6c4138ee5c86423
\newcommand{\MlstCmpJcTwoQcdCoarseOverFullPr}{47\% more errors}  % experiments/FIGS/data/v2/index.json :: linear probe jc2_188 N=1000000: (1 - macro_auc_qcd) of R16_Q1 over L188, run by run :: sha256 c6c4138ee5c86423 :: prose twin: two significant figures, direction in words
\newcommand{\MlstCmpJcTwoQcdResMassOverRes}{1.012 and 1.006}  % experiments/FIGS/data/v2/index.json :: linear probe jc2_188 N=1000000: (1 - macro_auc_qcd) of L162_MASS over L162, run by run :: sha256 c6c4138ee5c86423
\newcommand{\MlstCmpJcTwoQcdResMassOverResPr}{equal within 2\%}  % experiments/FIGS/data/v2/index.json :: linear probe jc2_188 N=1000000: (1 - macro_auc_qcd) of L162_MASS over L162, run by run :: sha256 c6c4138ee5c86423 :: prose twin: two significant figures, direction in words
\newcommand{\MlstCmpJcTwoQcdResMassOverResPrMax}{equal within 2\%}  % experiments/FIGS/data/v2/index.json :: linear probe jc2_188 N=1000000: (1 - macro_auc_qcd) of L162_MASS over L162, run by run :: sha256 c6c4138ee5c86423 :: prose twin of the run farther from 1
\newcommand{\MlstCmpJcTwoQcdCoarseMassOverCoarse}{0.991 and 1.014}  % experiments/FIGS/data/v2/index.json :: linear probe jc2_188 N=1000000: (1 - macro_auc_qcd) of R16_Q1_MASS over R16_Q1, run by run :: sha256 c6c4138ee5c86423
\newcommand{\MlstCmpJcTwoQcdCoarseMassOverCoarsePr}{equal within 2\%}  % experiments/FIGS/data/v2/index.json :: linear probe jc2_188 N=1000000: (1 - macro_auc_qcd) of R16_Q1_MASS over R16_Q1, run by run :: sha256 c6c4138ee5c86423 :: prose twin: two significant figures, direction in words
\newcommand{\MlstCmpJcTwoQcdCoarseMassOverCoarsePrMax}{equal within 2\%}  % experiments/FIGS/data/v2/index.json :: linear probe jc2_188 N=1000000: (1 - macro_auc_qcd) of R16_Q1_MASS over R16_Q1, run by run :: sha256 c6c4138ee5c86423 :: prose twin of the run farther from 1
\newcommand{\MlstCmpJcTwoQcdResMassOverFull}{1.107 and 1.092}  % experiments/FIGS/data/v2/index.json :: linear probe jc2_188 N=1000000: (1 - macro_auc_qcd) of L162_MASS over L188, run by run :: sha256 c6c4138ee5c86423
\newcommand{\MlstCmpJcTwoQcdResMassOverFullPr}{at least 9.2\% more errors}  % experiments/FIGS/data/v2/index.json :: linear probe jc2_188 N=1000000: (1 - macro_auc_qcd) of L162_MASS over L188, run by run :: sha256 c6c4138ee5c86423 :: prose twin: two significant figures, direction in words
\newcommand{\MlstCmpJcTwoQcdResMassOverFullPrMax}{up to 11\% more errors}  % experiments/FIGS/data/v2/index.json :: linear probe jc2_188 N=1000000: (1 - macro_auc_qcd) of L162_MASS over L188, run by run :: sha256 c6c4138ee5c86423 :: prose twin of the run farther from 1
\newcommand{\MlstCmpJcTwoQcdCoarseMassOverFull}{1.472 and 1.474}  % experiments/FIGS/data/v2/index.json :: linear probe jc2_188 N=1000000: (1 - macro_auc_qcd) of R16_Q1_MASS over L188, run by run :: sha256 c6c4138ee5c86423
\newcommand{\MlstCmpJcTwoQcdCoarseMassOverFullPr}{at least 47\% more errors}  % experiments/FIGS/data/v2/index.json :: linear probe jc2_188 N=1000000: (1 - macro_auc_qcd) of R16_Q1_MASS over L188, run by run :: sha256 c6c4138ee5c86423 :: prose twin: two significant figures, direction in words
\newcommand{\MlstCmpJcTwoQcdCoarseMassOverFullPrMax}{up to 47\% more errors}  % experiments/FIGS/data/v2/index.json :: linear probe jc2_188 N=1000000: (1 - macro_auc_qcd) of R16_Q1_MASS over L188, run by run :: sha256 c6c4138ee5c86423 :: prose twin of the run farther from 1
\newcommand{\MlstCmpJcTwoQcdFzeroOverFull}{1.466 and 1.497}  % experiments/FIGS/data/v2/index.json :: linear probe jc2_188 N=1000000: (1 - macro_auc_qcd) of FLAV_F0 over L188, run by run :: sha256 c6c4138ee5c86423
\newcommand{\MlstCmpJcTwoQcdFzeroOverFullPr}{at least 47\% more errors}  % experiments/FIGS/data/v2/index.json :: linear probe jc2_188 N=1000000: (1 - macro_auc_qcd) of FLAV_F0 over L188, run by run :: sha256 c6c4138ee5c86423 :: prose twin: two significant figures, direction in words
\newcommand{\MlstCmpJcTwoQcdFzeroOverFullPrMax}{up to 50\% more errors}  % experiments/FIGS/data/v2/index.json :: linear probe jc2_188 N=1000000: (1 - macro_auc_qcd) of FLAV_F0 over L188, run by run :: sha256 c6c4138ee5c86423 :: prose twin of the run farther from 1
\newcommand{\MlstCmpJcTwoQcdFoneOverFull}{1.241 and 1.235}  % experiments/FIGS/data/v2/index.json :: linear probe jc2_188 N=1000000: (1 - macro_auc_qcd) of FLAV_F1 over L188, run by run :: sha256 c6c4138ee5c86423
\newcommand{\MlstCmpJcTwoQcdFoneOverFullPr}{at least 23\% more errors}  % experiments/FIGS/data/v2/index.json :: linear probe jc2_188 N=1000000: (1 - macro_auc_qcd) of FLAV_F1 over L188, run by run :: sha256 c6c4138ee5c86423 :: prose twin: two significant figures, direction in words
\newcommand{\MlstCmpJcTwoQcdFoneOverFullPrMax}{up to 24\% more errors}  % experiments/FIGS/data/v2/index.json :: linear probe jc2_188 N=1000000: (1 - macro_auc_qcd) of FLAV_F1 over L188, run by run :: sha256 c6c4138ee5c86423 :: prose twin of the run farther from 1
\newcommand{\MlstCmpJcTwoQcdFonerOverFull}{1.260 and 1.263}  % experiments/FIGS/data/v2/index.json :: linear probe jc2_188 N=1000000: (1 - macro_auc_qcd) of FLAV_F1R over L188, run by run :: sha256 c6c4138ee5c86423
\newcommand{\MlstCmpJcTwoQcdFonerOverFullPr}{at least 26\% more errors}  % experiments/FIGS/data/v2/index.json :: linear probe jc2_188 N=1000000: (1 - macro_auc_qcd) of FLAV_F1R over L188, run by run :: sha256 c6c4138ee5c86423 :: prose twin: two significant figures, direction in words
\newcommand{\MlstCmpJcTwoQcdFonerOverFullPrMax}{up to 26\% more errors}  % experiments/FIGS/data/v2/index.json :: linear probe jc2_188 N=1000000: (1 - macro_auc_qcd) of FLAV_F1R over L188, run by run :: sha256 c6c4138ee5c86423 :: prose twin of the run farther from 1
\newcommand{\MlstCmpJcTwoQcdFzeroOverCoarse}{0.987 and 1.031}  % experiments/FIGS/data/v2/index.json :: linear probe jc2_188 N=1000000: (1 - macro_auc_qcd) of FLAV_F0 over R16_Q1, run by run :: sha256 c6c4138ee5c86423
\newcommand{\MlstCmpJcTwoQcdFzeroOverCoarsePr}{equal within 4\%}  % experiments/FIGS/data/v2/index.json :: linear probe jc2_188 N=1000000: (1 - macro_auc_qcd) of FLAV_F0 over R16_Q1, run by run :: sha256 c6c4138ee5c86423 :: prose twin: two significant figures, direction in words
\newcommand{\MlstCmpJcTwoQcdFzeroOverCoarsePrMax}{equal within 4\%}  % experiments/FIGS/data/v2/index.json :: linear probe jc2_188 N=1000000: (1 - macro_auc_qcd) of FLAV_F0 over R16_Q1, run by run :: sha256 c6c4138ee5c86423 :: prose twin of the run farther from 1
\newcommand{\MlstCmpJcTwoQcdFoneOverCoarse}{0.836 and 0.850}  % experiments/FIGS/data/v2/index.json :: linear probe jc2_188 N=1000000: (1 - macro_auc_qcd) of FLAV_F1 over R16_Q1, run by run :: sha256 c6c4138ee5c86423
\newcommand{\MlstCmpJcTwoQcdFoneOverCoarsePr}{at least 15\% fewer errors}  % experiments/FIGS/data/v2/index.json :: linear probe jc2_188 N=1000000: (1 - macro_auc_qcd) of FLAV_F1 over R16_Q1, run by run :: sha256 c6c4138ee5c86423 :: prose twin: two significant figures, direction in words
\newcommand{\MlstCmpJcTwoQcdFoneOverCoarsePrMax}{up to 16\% fewer errors}  % experiments/FIGS/data/v2/index.json :: linear probe jc2_188 N=1000000: (1 - macro_auc_qcd) of FLAV_F1 over R16_Q1, run by run :: sha256 c6c4138ee5c86423 :: prose twin of the run farther from 1
\newcommand{\MlstCmpJcTwoQcdFonerOverCoarse}{0.848 and 0.870}  % experiments/FIGS/data/v2/index.json :: linear probe jc2_188 N=1000000: (1 - macro_auc_qcd) of FLAV_F1R over R16_Q1, run by run :: sha256 c6c4138ee5c86423
\newcommand{\MlstCmpJcTwoQcdFonerOverCoarsePr}{at least 13\% fewer errors}  % experiments/FIGS/data/v2/index.json :: linear probe jc2_188 N=1000000: (1 - macro_auc_qcd) of FLAV_F1R over R16_Q1, run by run :: sha256 c6c4138ee5c86423 :: prose twin: two significant figures, direction in words
\newcommand{\MlstCmpJcTwoQcdFonerOverCoarsePrMax}{up to 15\% fewer errors}  % experiments/FIGS/data/v2/index.json :: linear probe jc2_188 N=1000000: (1 - macro_auc_qcd) of FLAV_F1R over R16_Q1, run by run :: sha256 c6c4138ee5c86423 :: prose twin of the run farther from 1
\newcommand{\MlstCmpJcTwoQcdFoneOverFzero}{0.847 and 0.825}  % experiments/FIGS/data/v2/index.json :: linear probe jc2_188 N=1000000: (1 - macro_auc_qcd) of FLAV_F1 over FLAV_F0, run by run :: sha256 c6c4138ee5c86423
\newcommand{\MlstCmpJcTwoQcdFoneOverFzeroPr}{at least 15\% fewer errors}  % experiments/FIGS/data/v2/index.json :: linear probe jc2_188 N=1000000: (1 - macro_auc_qcd) of FLAV_F1 over FLAV_F0, run by run :: sha256 c6c4138ee5c86423 :: prose twin: two significant figures, direction in words
\newcommand{\MlstCmpJcTwoQcdFoneOverFzeroPrMax}{up to 18\% fewer errors}  % experiments/FIGS/data/v2/index.json :: linear probe jc2_188 N=1000000: (1 - macro_auc_qcd) of FLAV_F1 over FLAV_F0, run by run :: sha256 c6c4138ee5c86423 :: prose twin of the run farther from 1
\newcommand{\MlstCmpJcTwoQcdFonerOverFzero}{0.859 and 0.844}  % experiments/FIGS/data/v2/index.json :: linear probe jc2_188 N=1000000: (1 - macro_auc_qcd) of FLAV_F1R over FLAV_F0, run by run :: sha256 c6c4138ee5c86423
\newcommand{\MlstCmpJcTwoQcdFonerOverFzeroPr}{at least 14\% fewer errors}  % experiments/FIGS/data/v2/index.json :: linear probe jc2_188 N=1000000: (1 - macro_auc_qcd) of FLAV_F1R over FLAV_F0, run by run :: sha256 c6c4138ee5c86423 :: prose twin: two significant figures, direction in words
\newcommand{\MlstCmpJcTwoQcdFonerOverFzeroPrMax}{up to 16\% fewer errors}  % experiments/FIGS/data/v2/index.json :: linear probe jc2_188 N=1000000: (1 - macro_auc_qcd) of FLAV_F1R over FLAV_F0, run by run :: sha256 c6c4138ee5c86423 :: prose twin of the run farther from 1
\newcommand{\MlstCmpJcTwoQcdFonerOverFone}{1.015 and 1.023}  % experiments/FIGS/data/v2/index.json :: linear probe jc2_188 N=1000000: (1 - macro_auc_qcd) of FLAV_F1R over FLAV_F1, run by run :: sha256 c6c4138ee5c86423
\newcommand{\MlstCmpJcTwoQcdFonerOverFonePr}{equal within 3\%}  % experiments/FIGS/data/v2/index.json :: linear probe jc2_188 N=1000000: (1 - macro_auc_qcd) of FLAV_F1R over FLAV_F1, run by run :: sha256 c6c4138ee5c86423 :: prose twin: two significant figures, direction in words
\newcommand{\MlstCmpJcTwoQcdFonerOverFonePrMax}{equal within 3\%}  % experiments/FIGS/data/v2/index.json :: linear probe jc2_188 N=1000000: (1 - macro_auc_qcd) of FLAV_F1R over FLAV_F1, run by run :: sha256 c6c4138ee5c86423 :: prose twin of the run farther from 1
\newcommand{\MlstCmpJcTwoQcdMidOverRes}{1.040\,\ensuremath{\pm}\,0.011}  % experiments/FIGS/data/v2/index.json :: linear probe jc2_188 N=1000000: (1 - macro_auc_qcd) of R42_Q1 over L162, run by run :: sha256 c6c4138ee5c86423
\newcommand{\MlstCmpJcTwoQcdMidOverResPr}{4.0\% more errors}  % experiments/FIGS/data/v2/index.json :: linear probe jc2_188 N=1000000: (1 - macro_auc_qcd) of R42_Q1 over L162, run by run :: sha256 c6c4138ee5c86423 :: prose twin: two significant figures, direction in words
\newcommand{\MlstCmpJcTwoQcdCoarseOverRes}{1.348\,\ensuremath{\pm}\,0.010}  % experiments/FIGS/data/v2/index.json :: linear probe jc2_188 N=1000000: (1 - macro_auc_qcd) of R16_Q1 over L162, run by run :: sha256 c6c4138ee5c86423
\newcommand{\MlstCmpJcTwoQcdCoarseOverResPr}{35\% more errors}  % experiments/FIGS/data/v2/index.json :: linear probe jc2_188 N=1000000: (1 - macro_auc_qcd) of R16_Q1 over L162, run by run :: sha256 c6c4138ee5c86423 :: prose twin: two significant figures, direction in words
\newcommand{\MlstCmpJcTwoQcdCoarseOverMid}{1.296\,\ensuremath{\pm}\,0.016}  % experiments/FIGS/data/v2/index.json :: linear probe jc2_188 N=1000000: (1 - macro_auc_qcd) of R16_Q1 over R42_Q1, run by run :: sha256 c6c4138ee5c86423
\newcommand{\MlstCmpJcTwoQcdCoarseOverMidPr}{30\% more errors}  % experiments/FIGS/data/v2/index.json :: linear probe jc2_188 N=1000000: (1 - macro_auc_qcd) of R16_Q1 over R42_Q1, run by run :: sha256 c6c4138ee5c86423 :: prose twin: two significant figures, direction in words
\newcommand{\MlstCmpJcTwoQcdRandOverFull}{1.198\,\ensuremath{\pm}\,0.011}  % experiments/FIGS/data/v2/index.json :: linear probe jc2_188 N=1000000: (1 - macro_auc_qcd) of each random partition over the L188 run of its seed, mean +- SD over the partitions :: sha256 c6c4138ee5c86423
\newcommand{\MlstCmpJcTwoQcdRandOverFullPr}{20\% more errors}  % experiments/FIGS/data/v2/index.json :: linear probe jc2_188 N=1000000: (1 - macro_auc_qcd) of each random partition over the L188 run of its seed, mean +- SD over the partitions :: sha256 c6c4138ee5c86423 :: prose twin: two significant figures, direction in words
\newcommand{\MlstCmpJcTwoQcdRandOverMid}{1.064\,\ensuremath{\pm}\,0.011}  % experiments/FIGS/data/v2/index.json :: linear probe jc2_188 N=1000000: (1 - macro_auc_qcd) of each random partition over the R42_Q1 run of its seed, mean +- SD over the partitions :: sha256 c6c4138ee5c86423
\newcommand{\MlstCmpJcTwoQcdRandOverMidPr}{6.4\% more errors}  % experiments/FIGS/data/v2/index.json :: linear probe jc2_188 N=1000000: (1 - macro_auc_qcd) of each random partition over the R42_Q1 run of its seed, mean +- SD over the partitions :: sha256 c6c4138ee5c86423 :: prose twin: two significant figures, direction in words
\newcommand{\MlstCmpJcTwoQcdRandOverCoarse}{0.818\,\ensuremath{\pm}\,0.013}  % experiments/FIGS/data/v2/index.json :: linear probe jc2_188 N=1000000: (1 - macro_auc_qcd) of each random partition over the R16_Q1 run of its seed, mean +- SD over the partitions :: sha256 c6c4138ee5c86423
\newcommand{\MlstCmpJcTwoQcdRandOverCoarsePr}{18\% fewer errors}  % experiments/FIGS/data/v2/index.json :: linear probe jc2_188 N=1000000: (1 - macro_auc_qcd) of each random partition over the R16_Q1 run of its seed, mean +- SD over the partitions :: sha256 c6c4138ee5c86423 :: prose twin: two significant figures, direction in words
\newcommand{\MlstCmpBcResOverFull}{0.998\,\ensuremath{\pm}\,0.007}  % experiments/FIGS/data/v2/index.json :: binary probe bc_vs_bq_cs: (1 - AUC) of L162 over L188, run by run :: sha256 c6c4138ee5c86423
\newcommand{\MlstCmpBcResOverFullPr}{equal within 1\%}  % experiments/FIGS/data/v2/index.json :: binary probe bc_vs_bq_cs: (1 - AUC) of L162 over L188, run by run :: sha256 c6c4138ee5c86423 :: prose twin: two significant figures, direction in words
\newcommand{\MlstCmpBcMidOverFull}{1.151\,\ensuremath{\pm}\,0.029}  % experiments/FIGS/data/v2/index.json :: binary probe bc_vs_bq_cs: (1 - AUC) of R42_Q1 over L188, run by run :: sha256 c6c4138ee5c86423
\newcommand{\MlstCmpBcMidOverFullPr}{15\% more errors}  % experiments/FIGS/data/v2/index.json :: binary probe bc_vs_bq_cs: (1 - AUC) of R42_Q1 over L188, run by run :: sha256 c6c4138ee5c86423 :: prose twin: two significant figures, direction in words
\newcommand{\MlstCmpBcCoarseOverFull}{1.95\,\ensuremath{\pm}\,0.06}  % experiments/FIGS/data/v2/index.json :: binary probe bc_vs_bq_cs: (1 - AUC) of R16_Q1 over L188, run by run :: sha256 c6c4138ee5c86423
\newcommand{\MlstCmpBcCoarseOverFullPr}{95\% more errors}  % experiments/FIGS/data/v2/index.json :: binary probe bc_vs_bq_cs: (1 - AUC) of R16_Q1 over L188, run by run :: sha256 c6c4138ee5c86423 :: prose twin: two significant figures, direction in words
\newcommand{\MlstCmpBcFoneOverFzero}{0.759 and 0.698}  % experiments/FIGS/data/v2/index.json :: binary probe bc_vs_bq_cs: (1 - AUC) of FLAV_F1 over FLAV_F0, run by run :: sha256 c6c4138ee5c86423
\newcommand{\MlstCmpBcFoneOverFzeroPr}{at least 24\% fewer errors}  % experiments/FIGS/data/v2/index.json :: binary probe bc_vs_bq_cs: (1 - AUC) of FLAV_F1 over FLAV_F0, run by run :: sha256 c6c4138ee5c86423 :: prose twin: two significant figures, direction in words
\newcommand{\MlstCmpBcFoneOverFzeroPrMax}{up to 30\% fewer errors}  % experiments/FIGS/data/v2/index.json :: binary probe bc_vs_bq_cs: (1 - AUC) of FLAV_F1 over FLAV_F0, run by run :: sha256 c6c4138ee5c86423 :: prose twin of the run farther from 1
\newcommand{\MlstCmpBcFonerOverFone}{1.008 and 0.990}  % experiments/FIGS/data/v2/index.json :: binary probe bc_vs_bq_cs: (1 - AUC) of FLAV_F1R over FLAV_F1, run by run :: sha256 c6c4138ee5c86423
\newcommand{\MlstCmpBcFonerOverFonePr}{equal within 1\%}  % experiments/FIGS/data/v2/index.json :: binary probe bc_vs_bq_cs: (1 - AUC) of FLAV_F1R over FLAV_F1, run by run :: sha256 c6c4138ee5c86423 :: prose twin: two significant figures, direction in words
\newcommand{\MlstCmpBcFonerOverFonePrMax}{equal within 1\%}  % experiments/FIGS/data/v2/index.json :: binary probe bc_vs_bq_cs: (1 - AUC) of FLAV_F1R over FLAV_F1, run by run :: sha256 c6c4138ee5c86423 :: prose twin of the run farther from 1
\newcommand{\MlstCmpBcFonerOverFzero}{0.765 and 0.691}  % experiments/FIGS/data/v2/index.json :: binary probe bc_vs_bq_cs: (1 - AUC) of FLAV_F1R over FLAV_F0, run by run :: sha256 c6c4138ee5c86423
\newcommand{\MlstCmpBcFonerOverFzeroPr}{at least 23\% fewer errors}  % experiments/FIGS/data/v2/index.json :: binary probe bc_vs_bq_cs: (1 - AUC) of FLAV_F1R over FLAV_F0, run by run :: sha256 c6c4138ee5c86423 :: prose twin: two significant figures, direction in words
\newcommand{\MlstCmpBcFonerOverFzeroPrMax}{up to 31\% fewer errors}  % experiments/FIGS/data/v2/index.json :: binary probe bc_vs_bq_cs: (1 - AUC) of FLAV_F1R over FLAV_F0, run by run :: sha256 c6c4138ee5c86423 :: prose twin of the run farther from 1
\newcommand{\MlstCmpBcFzeroOverCoarse}{0.99 and 1.13}  % experiments/FIGS/data/v2/index.json :: binary probe bc_vs_bq_cs: (1 - AUC) of FLAV_F0 over R16_Q1, run by run :: sha256 c6c4138ee5c86423
\newcommand{\MlstCmpBcFzeroOverCoarsePr}{equal within 20\%}  % experiments/FIGS/data/v2/index.json :: binary probe bc_vs_bq_cs: (1 - AUC) of FLAV_F0 over R16_Q1, run by run :: sha256 c6c4138ee5c86423 :: prose twin: two significant figures, direction in words
\newcommand{\MlstCmpBcFzeroOverCoarsePrMax}{equal within 20\%}  % experiments/FIGS/data/v2/index.json :: binary probe bc_vs_bq_cs: (1 - AUC) of FLAV_F0 over R16_Q1, run by run :: sha256 c6c4138ee5c86423 :: prose twin of the run farther from 1
\newcommand{\MlstCmpBcFoneOverCoarse}{0.752 and 0.786}  % experiments/FIGS/data/v2/index.json :: binary probe bc_vs_bq_cs: (1 - AUC) of FLAV_F1 over R16_Q1, run by run :: sha256 c6c4138ee5c86423
\newcommand{\MlstCmpBcFoneOverCoarsePr}{at least 21\% fewer errors}  % experiments/FIGS/data/v2/index.json :: binary probe bc_vs_bq_cs: (1 - AUC) of FLAV_F1 over R16_Q1, run by run :: sha256 c6c4138ee5c86423 :: prose twin: two significant figures, direction in words
\newcommand{\MlstCmpBcFoneOverCoarsePrMax}{up to 25\% fewer errors}  % experiments/FIGS/data/v2/index.json :: binary probe bc_vs_bq_cs: (1 - AUC) of FLAV_F1 over R16_Q1, run by run :: sha256 c6c4138ee5c86423 :: prose twin of the run farther from 1
\newcommand{\MlstCmpBcFonerOverCoarse}{0.758 and 0.778}  % experiments/FIGS/data/v2/index.json :: binary probe bc_vs_bq_cs: (1 - AUC) of FLAV_F1R over R16_Q1, run by run :: sha256 c6c4138ee5c86423
\newcommand{\MlstCmpBcFonerOverCoarsePr}{at least 22\% fewer errors}  % experiments/FIGS/data/v2/index.json :: binary probe bc_vs_bq_cs: (1 - AUC) of FLAV_F1R over R16_Q1, run by run :: sha256 c6c4138ee5c86423 :: prose twin: two significant figures, direction in words
\newcommand{\MlstCmpBcFonerOverCoarsePrMax}{up to 24\% fewer errors}  % experiments/FIGS/data/v2/index.json :: binary probe bc_vs_bq_cs: (1 - AUC) of FLAV_F1R over R16_Q1, run by run :: sha256 c6c4138ee5c86423 :: prose twin of the run farther from 1
\newcommand{\MlstCmpBcFzeroOverFull}{2.00 and 2.19}  % experiments/FIGS/data/v2/index.json :: binary probe bc_vs_bq_cs: (1 - AUC) of FLAV_F0 over L188, run by run :: sha256 c6c4138ee5c86423
\newcommand{\MlstCmpBcFzeroOverFullPr}{at least 2.0 times the errors}  % experiments/FIGS/data/v2/index.json :: binary probe bc_vs_bq_cs: (1 - AUC) of FLAV_F0 over L188, run by run :: sha256 c6c4138ee5c86423 :: prose twin: two significant figures, direction in words
\newcommand{\MlstCmpBcFzeroOverFullPrMax}{up to 2.2 times the errors}  % experiments/FIGS/data/v2/index.json :: binary probe bc_vs_bq_cs: (1 - AUC) of FLAV_F0 over L188, run by run :: sha256 c6c4138ee5c86423 :: prose twin of the run farther from 1
\newcommand{\MlstCmpBcFoneOverFull}{1.518 and 1.530}  % experiments/FIGS/data/v2/index.json :: binary probe bc_vs_bq_cs: (1 - AUC) of FLAV_F1 over L188, run by run :: sha256 c6c4138ee5c86423
\newcommand{\MlstCmpBcFoneOverFullPr}{at least 52\% more errors}  % experiments/FIGS/data/v2/index.json :: binary probe bc_vs_bq_cs: (1 - AUC) of FLAV_F1 over L188, run by run :: sha256 c6c4138ee5c86423 :: prose twin: two significant figures, direction in words
\newcommand{\MlstCmpBcFoneOverFullPrMax}{up to 53\% more errors}  % experiments/FIGS/data/v2/index.json :: binary probe bc_vs_bq_cs: (1 - AUC) of FLAV_F1 over L188, run by run :: sha256 c6c4138ee5c86423 :: prose twin of the run farther from 1
\newcommand{\MlstCmpBcFonerOverFull}{1.530 and 1.515}  % experiments/FIGS/data/v2/index.json :: binary probe bc_vs_bq_cs: (1 - AUC) of FLAV_F1R over L188, run by run :: sha256 c6c4138ee5c86423
\newcommand{\MlstCmpBcFonerOverFullPr}{at least 51\% more errors}  % experiments/FIGS/data/v2/index.json :: binary probe bc_vs_bq_cs: (1 - AUC) of FLAV_F1R over L188, run by run :: sha256 c6c4138ee5c86423 :: prose twin: two significant figures, direction in words
\newcommand{\MlstCmpBcFonerOverFullPrMax}{up to 53\% more errors}  % experiments/FIGS/data/v2/index.json :: binary probe bc_vs_bq_cs: (1 - AUC) of FLAV_F1R over L188, run by run :: sha256 c6c4138ee5c86423 :: prose twin of the run farther from 1
\newcommand{\MlstCmpBcRandOverFull}{1.15\,\ensuremath{\pm}\,0.07}  % experiments/FIGS/data/v2/index.json :: binary probe bc_vs_bq_cs: (1 - AUC) of each random partition over the 188-class run of its seed :: sha256 c6c4138ee5c86423
\newcommand{\MlstCmpBcRandOverFullPr}{15\% more errors}  % experiments/FIGS/data/v2/index.json :: binary probe bc_vs_bq_cs: (1 - AUC) of each random partition over the 188-class run of its seed :: sha256 c6c4138ee5c86423 :: prose twin: two significant figures, direction in words
\newcommand{\MlstPairRandMergedBcN}{2}  % experiments/FIGS/data/v2/index.json :: random partitions that merge bc_vs_bq_cs (count) :: sha256 c6c4138ee5c86423
\newcommand{\MlstPairRandApartBcN}{2}  % experiments/FIGS/data/v2/index.json :: random partitions that keep apart bc_vs_bq_cs (count) :: sha256 c6c4138ee5c86423
\newcommand{\MlstPairRandMergedOverApartBc}{1.12}  % experiments/FIGS/data/v2/index.json :: binary probe bc_vs_bq_cs: mean (1 - AUC) of the partitions that merge it over that of those that keep it apart (one run each) :: sha256 c6c4138ee5c86423
\newcommand{\MlstPairRandMergedOverApartBcPr}{12\% more errors}  % experiments/FIGS/data/v2/index.json :: binary probe bc_vs_bq_cs: mean (1 - AUC) of the partitions that merge it over that of those that keep it apart (one run each) :: sha256 c6c4138ee5c86423 :: prose twin: two significant figures, direction in words
\newcommand{\MlstCmpBbResOverFull}{1.015\,\ensuremath{\pm}\,0.027}  % experiments/FIGS/data/v2/index.json :: binary probe bb_vs_cc: (1 - AUC) of L162 over L188, run by run :: sha256 c6c4138ee5c86423
\newcommand{\MlstCmpBbResOverFullPr}{equal within 3\%}  % experiments/FIGS/data/v2/index.json :: binary probe bb_vs_cc: (1 - AUC) of L162 over L188, run by run :: sha256 c6c4138ee5c86423 :: prose twin: two significant figures, direction in words
\newcommand{\MlstCmpBbMidOverFull}{1.102\,\ensuremath{\pm}\,0.013}  % experiments/FIGS/data/v2/index.json :: binary probe bb_vs_cc: (1 - AUC) of R42_Q1 over L188, run by run :: sha256 c6c4138ee5c86423
\newcommand{\MlstCmpBbMidOverFullPr}{10\% more errors}  % experiments/FIGS/data/v2/index.json :: binary probe bb_vs_cc: (1 - AUC) of R42_Q1 over L188, run by run :: sha256 c6c4138ee5c86423 :: prose twin: two significant figures, direction in words
\newcommand{\MlstCmpBbCoarseOverFull}{4.22\,\ensuremath{\pm}\,0.29}  % experiments/FIGS/data/v2/index.json :: binary probe bb_vs_cc: (1 - AUC) of R16_Q1 over L188, run by run :: sha256 c6c4138ee5c86423
\newcommand{\MlstCmpBbCoarseOverFullPr}{4.2 times the errors}  % experiments/FIGS/data/v2/index.json :: binary probe bb_vs_cc: (1 - AUC) of R16_Q1 over L188, run by run :: sha256 c6c4138ee5c86423 :: prose twin: two significant figures, direction in words
\newcommand{\MlstCmpBbFoneOverFzero}{0.31 and 0.13}  % experiments/FIGS/data/v2/index.json :: binary probe bb_vs_cc: (1 - AUC) of FLAV_F1 over FLAV_F0, run by run :: sha256 c6c4138ee5c86423
\newcommand{\MlstCmpBbFoneOverFzeroPr}{at most 0.31 times the errors}  % experiments/FIGS/data/v2/index.json :: binary probe bb_vs_cc: (1 - AUC) of FLAV_F1 over FLAV_F0, run by run :: sha256 c6c4138ee5c86423 :: prose twin: two significant figures, direction in words
\newcommand{\MlstCmpBbFoneOverFzeroPrMax}{down to 0.13 times the errors}  % experiments/FIGS/data/v2/index.json :: binary probe bb_vs_cc: (1 - AUC) of FLAV_F1 over FLAV_F0, run by run :: sha256 c6c4138ee5c86423 :: prose twin of the run farther from 1
\newcommand{\MlstCmpBbFonerOverFone}{1.40 and 1.86}  % experiments/FIGS/data/v2/index.json :: binary probe bb_vs_cc: (1 - AUC) of FLAV_F1R over FLAV_F1, run by run :: sha256 c6c4138ee5c86423
\newcommand{\MlstCmpBbFonerOverFonePr}{at least 40\% more errors}  % experiments/FIGS/data/v2/index.json :: binary probe bb_vs_cc: (1 - AUC) of FLAV_F1R over FLAV_F1, run by run :: sha256 c6c4138ee5c86423 :: prose twin: two significant figures, direction in words
\newcommand{\MlstCmpBbFonerOverFonePrMax}{up to 86\% more errors}  % experiments/FIGS/data/v2/index.json :: binary probe bb_vs_cc: (1 - AUC) of FLAV_F1R over FLAV_F1, run by run :: sha256 c6c4138ee5c86423 :: prose twin of the run farther from 1
\newcommand{\MlstCmpBbFonerOverFzero}{0.44 and 0.24}  % experiments/FIGS/data/v2/index.json :: binary probe bb_vs_cc: (1 - AUC) of FLAV_F1R over FLAV_F0, run by run :: sha256 c6c4138ee5c86423
\newcommand{\MlstCmpBbFonerOverFzeroPr}{at most 0.44 times the errors}  % experiments/FIGS/data/v2/index.json :: binary probe bb_vs_cc: (1 - AUC) of FLAV_F1R over FLAV_F0, run by run :: sha256 c6c4138ee5c86423 :: prose twin: two significant figures, direction in words
\newcommand{\MlstCmpBbFonerOverFzeroPrMax}{down to 0.24 times the errors}  % experiments/FIGS/data/v2/index.json :: binary probe bb_vs_cc: (1 - AUC) of FLAV_F1R over FLAV_F0, run by run :: sha256 c6c4138ee5c86423 :: prose twin of the run farther from 1
\newcommand{\MlstCmpBbFzeroOverCoarse}{1.3 and 3.2}  % experiments/FIGS/data/v2/index.json :: binary probe bb_vs_cc: (1 - AUC) of FLAV_F0 over R16_Q1, run by run :: sha256 c6c4138ee5c86423
\newcommand{\MlstCmpBbFzeroOverCoarsePr}{at least 30\% more errors}  % experiments/FIGS/data/v2/index.json :: binary probe bb_vs_cc: (1 - AUC) of FLAV_F0 over R16_Q1, run by run :: sha256 c6c4138ee5c86423 :: prose twin: two significant figures, direction in words
\newcommand{\MlstCmpBbFzeroOverCoarsePrMax}{up to 3.2 times the errors}  % experiments/FIGS/data/v2/index.json :: binary probe bb_vs_cc: (1 - AUC) of FLAV_F0 over R16_Q1, run by run :: sha256 c6c4138ee5c86423 :: prose twin of the run farther from 1
\newcommand{\MlstCmpBbFoneOverCoarse}{0.4077 and 0.4126}  % experiments/FIGS/data/v2/index.json :: binary probe bb_vs_cc: (1 - AUC) of FLAV_F1 over R16_Q1, run by run :: sha256 c6c4138ee5c86423
\newcommand{\MlstCmpBbFoneOverCoarsePr}{at most 0.41 times the errors}  % experiments/FIGS/data/v2/index.json :: binary probe bb_vs_cc: (1 - AUC) of FLAV_F1 over R16_Q1, run by run :: sha256 c6c4138ee5c86423 :: prose twin: two significant figures, direction in words
\newcommand{\MlstCmpBbFoneOverCoarsePrMax}{down to 0.41 times the errors}  % experiments/FIGS/data/v2/index.json :: binary probe bb_vs_cc: (1 - AUC) of FLAV_F1 over R16_Q1, run by run :: sha256 c6c4138ee5c86423 :: prose twin of the run farther from 1
\newcommand{\MlstCmpBbFonerOverCoarse}{0.57 and 0.77}  % experiments/FIGS/data/v2/index.json :: binary probe bb_vs_cc: (1 - AUC) of FLAV_F1R over R16_Q1, run by run :: sha256 c6c4138ee5c86423
\newcommand{\MlstCmpBbFonerOverCoarsePr}{at least 23\% fewer errors}  % experiments/FIGS/data/v2/index.json :: binary probe bb_vs_cc: (1 - AUC) of FLAV_F1R over R16_Q1, run by run :: sha256 c6c4138ee5c86423 :: prose twin: two significant figures, direction in words
\newcommand{\MlstCmpBbFonerOverCoarsePrMax}{up to 43\% fewer errors}  % experiments/FIGS/data/v2/index.json :: binary probe bb_vs_cc: (1 - AUC) of FLAV_F1R over R16_Q1, run by run :: sha256 c6c4138ee5c86423 :: prose twin of the run farther from 1
\newcommand{\MlstCmpBbFzeroOverFull}{5.9 and 13.4}  % experiments/FIGS/data/v2/index.json :: binary probe bb_vs_cc: (1 - AUC) of FLAV_F0 over L188, run by run :: sha256 c6c4138ee5c86423
\newcommand{\MlstCmpBbFzeroOverFullPr}{at least 5.9 times the errors}  % experiments/FIGS/data/v2/index.json :: binary probe bb_vs_cc: (1 - AUC) of FLAV_F0 over L188, run by run :: sha256 c6c4138ee5c86423 :: prose twin: two significant figures, direction in words
\newcommand{\MlstCmpBbFzeroOverFullPrMax}{up to 13 times the errors}  % experiments/FIGS/data/v2/index.json :: binary probe bb_vs_cc: (1 - AUC) of FLAV_F0 over L188, run by run :: sha256 c6c4138ee5c86423 :: prose twin of the run farther from 1
\newcommand{\MlstCmpBbFoneOverFull}{1.837 and 1.748}  % experiments/FIGS/data/v2/index.json :: binary probe bb_vs_cc: (1 - AUC) of FLAV_F1 over L188, run by run :: sha256 c6c4138ee5c86423
\newcommand{\MlstCmpBbFoneOverFullPr}{at least 75\% more errors}  % experiments/FIGS/data/v2/index.json :: binary probe bb_vs_cc: (1 - AUC) of FLAV_F1 over L188, run by run :: sha256 c6c4138ee5c86423 :: prose twin: two significant figures, direction in words
\newcommand{\MlstCmpBbFoneOverFullPrMax}{up to 84\% more errors}  % experiments/FIGS/data/v2/index.json :: binary probe bb_vs_cc: (1 - AUC) of FLAV_F1 over L188, run by run :: sha256 c6c4138ee5c86423 :: prose twin of the run farther from 1
\newcommand{\MlstCmpBbFonerOverFull}{2.58 and 3.25}  % experiments/FIGS/data/v2/index.json :: binary probe bb_vs_cc: (1 - AUC) of FLAV_F1R over L188, run by run :: sha256 c6c4138ee5c86423
\newcommand{\MlstCmpBbFonerOverFullPr}{at least 2.6 times the errors}  % experiments/FIGS/data/v2/index.json :: binary probe bb_vs_cc: (1 - AUC) of FLAV_F1R over L188, run by run :: sha256 c6c4138ee5c86423 :: prose twin: two significant figures, direction in words
\newcommand{\MlstCmpBbFonerOverFullPrMax}{up to 3.2 times the errors}  % experiments/FIGS/data/v2/index.json :: binary probe bb_vs_cc: (1 - AUC) of FLAV_F1R over L188, run by run :: sha256 c6c4138ee5c86423 :: prose twin of the run farther from 1
\newcommand{\MlstCmpBbRandOverFull}{1.53\,\ensuremath{\pm}\,0.11}  % experiments/FIGS/data/v2/index.json :: binary probe bb_vs_cc: (1 - AUC) of each random partition over the 188-class run of its seed :: sha256 c6c4138ee5c86423
\newcommand{\MlstCmpBbRandOverFullPr}{53\% more errors}  % experiments/FIGS/data/v2/index.json :: binary probe bb_vs_cc: (1 - AUC) of each random partition over the 188-class run of its seed :: sha256 c6c4138ee5c86423 :: prose twin: two significant figures, direction in words
\newcommand{\MlstPairRandMergedBbN}{3}  % experiments/FIGS/data/v2/index.json :: random partitions that merge bb_vs_cc (count) :: sha256 c6c4138ee5c86423
\newcommand{\MlstPairRandApartBbN}{2}  % experiments/FIGS/data/v2/index.json :: random partitions that keep apart bb_vs_cc (count) :: sha256 c6c4138ee5c86423
\newcommand{\MlstPairRandMergedOverApartBb}{1.07}  % experiments/FIGS/data/v2/index.json :: binary probe bb_vs_cc: mean (1 - AUC) of the partitions that merge it over that of those that keep it apart (one run each) :: sha256 c6c4138ee5c86423
\newcommand{\MlstPairRandMergedOverApartBbPr}{7.3\% more errors}  % experiments/FIGS/data/v2/index.json :: binary probe bb_vs_cc: mean (1 - AUC) of the partitions that merge it over that of those that keep it apart (one run each) :: sha256 c6c4138ee5c86423 :: prose twin: two significant figures, direction in words
\newcommand{\MlstCmpBbqqResOverFull}{1.005\,\ensuremath{\pm}\,0.032}  % experiments/FIGS/data/v2/index.json :: binary probe bbqq_vs_ccqq: (1 - AUC) of L162 over L188, run by run :: sha256 c6c4138ee5c86423
\newcommand{\MlstCmpBbqqResOverFullPr}{equal within 4\%}  % experiments/FIGS/data/v2/index.json :: binary probe bbqq_vs_ccqq: (1 - AUC) of L162 over L188, run by run :: sha256 c6c4138ee5c86423 :: prose twin: two significant figures, direction in words
\newcommand{\MlstCmpBbqqMidOverFull}{1.117\,\ensuremath{\pm}\,0.008}  % experiments/FIGS/data/v2/index.json :: binary probe bbqq_vs_ccqq: (1 - AUC) of R42_Q1 over L188, run by run :: sha256 c6c4138ee5c86423
\newcommand{\MlstCmpBbqqMidOverFullPr}{12\% more errors}  % experiments/FIGS/data/v2/index.json :: binary probe bbqq_vs_ccqq: (1 - AUC) of R42_Q1 over L188, run by run :: sha256 c6c4138ee5c86423 :: prose twin: two significant figures, direction in words
\newcommand{\MlstCmpBbqqCoarseOverFull}{4.7\,\ensuremath{\pm}\,0.5}  % experiments/FIGS/data/v2/index.json :: binary probe bbqq_vs_ccqq: (1 - AUC) of R16_Q1 over L188, run by run :: sha256 c6c4138ee5c86423
\newcommand{\MlstCmpBbqqCoarseOverFullPr}{4.7 times the errors}  % experiments/FIGS/data/v2/index.json :: binary probe bbqq_vs_ccqq: (1 - AUC) of R16_Q1 over L188, run by run :: sha256 c6c4138ee5c86423 :: prose twin: two significant figures, direction in words
\newcommand{\MlstCmpBbqqFoneOverFzero}{0.236 and 0.154}  % experiments/FIGS/data/v2/index.json :: binary probe bbqq_vs_ccqq: (1 - AUC) of FLAV_F1 over FLAV_F0, run by run :: sha256 c6c4138ee5c86423
\newcommand{\MlstCmpBbqqFoneOverFzeroPr}{at most 0.24 times the errors}  % experiments/FIGS/data/v2/index.json :: binary probe bbqq_vs_ccqq: (1 - AUC) of FLAV_F1 over FLAV_F0, run by run :: sha256 c6c4138ee5c86423 :: prose twin: two significant figures, direction in words
\newcommand{\MlstCmpBbqqFoneOverFzeroPrMax}{down to 0.15 times the errors}  % experiments/FIGS/data/v2/index.json :: binary probe bbqq_vs_ccqq: (1 - AUC) of FLAV_F1 over FLAV_F0, run by run :: sha256 c6c4138ee5c86423 :: prose twin of the run farther from 1
\newcommand{\MlstCmpBbqqFonerOverFone}{1.505 and 1.577}  % experiments/FIGS/data/v2/index.json :: binary probe bbqq_vs_ccqq: (1 - AUC) of FLAV_F1R over FLAV_F1, run by run :: sha256 c6c4138ee5c86423
\newcommand{\MlstCmpBbqqFonerOverFonePr}{at least 50\% more errors}  % experiments/FIGS/data/v2/index.json :: binary probe bbqq_vs_ccqq: (1 - AUC) of FLAV_F1R over FLAV_F1, run by run :: sha256 c6c4138ee5c86423 :: prose twin: two significant figures, direction in words
\newcommand{\MlstCmpBbqqFonerOverFonePrMax}{up to 58\% more errors}  % experiments/FIGS/data/v2/index.json :: binary probe bbqq_vs_ccqq: (1 - AUC) of FLAV_F1R over FLAV_F1, run by run :: sha256 c6c4138ee5c86423 :: prose twin of the run farther from 1
\newcommand{\MlstCmpBbqqFonerOverFzero}{0.36 and 0.24}  % experiments/FIGS/data/v2/index.json :: binary probe bbqq_vs_ccqq: (1 - AUC) of FLAV_F1R over FLAV_F0, run by run :: sha256 c6c4138ee5c86423
\newcommand{\MlstCmpBbqqFonerOverFzeroPr}{at most 0.36 times the errors}  % experiments/FIGS/data/v2/index.json :: binary probe bbqq_vs_ccqq: (1 - AUC) of FLAV_F1R over FLAV_F0, run by run :: sha256 c6c4138ee5c86423 :: prose twin: two significant figures, direction in words
\newcommand{\MlstCmpBbqqFonerOverFzeroPrMax}{down to 0.24 times the errors}  % experiments/FIGS/data/v2/index.json :: binary probe bbqq_vs_ccqq: (1 - AUC) of FLAV_F1R over FLAV_F0, run by run :: sha256 c6c4138ee5c86423 :: prose twin of the run farther from 1
\newcommand{\MlstCmpBbqqFzeroOverCoarse}{1.51 and 1.79}  % experiments/FIGS/data/v2/index.json :: binary probe bbqq_vs_ccqq: (1 - AUC) of FLAV_F0 over R16_Q1, run by run :: sha256 c6c4138ee5c86423
\newcommand{\MlstCmpBbqqFzeroOverCoarsePr}{at least 51\% more errors}  % experiments/FIGS/data/v2/index.json :: binary probe bbqq_vs_ccqq: (1 - AUC) of FLAV_F0 over R16_Q1, run by run :: sha256 c6c4138ee5c86423 :: prose twin: two significant figures, direction in words
\newcommand{\MlstCmpBbqqFzeroOverCoarsePrMax}{up to 79\% more errors}  % experiments/FIGS/data/v2/index.json :: binary probe bbqq_vs_ccqq: (1 - AUC) of FLAV_F0 over R16_Q1, run by run :: sha256 c6c4138ee5c86423 :: prose twin of the run farther from 1
\newcommand{\MlstCmpBbqqFoneOverCoarse}{0.356 and 0.276}  % experiments/FIGS/data/v2/index.json :: binary probe bbqq_vs_ccqq: (1 - AUC) of FLAV_F1 over R16_Q1, run by run :: sha256 c6c4138ee5c86423
\newcommand{\MlstCmpBbqqFoneOverCoarsePr}{at most 0.36 times the errors}  % experiments/FIGS/data/v2/index.json :: binary probe bbqq_vs_ccqq: (1 - AUC) of FLAV_F1 over R16_Q1, run by run :: sha256 c6c4138ee5c86423 :: prose twin: two significant figures, direction in words
\newcommand{\MlstCmpBbqqFoneOverCoarsePrMax}{down to 0.28 times the errors}  % experiments/FIGS/data/v2/index.json :: binary probe bbqq_vs_ccqq: (1 - AUC) of FLAV_F1 over R16_Q1, run by run :: sha256 c6c4138ee5c86423 :: prose twin of the run farther from 1
\newcommand{\MlstCmpBbqqFonerOverCoarse}{0.54 and 0.44}  % experiments/FIGS/data/v2/index.json :: binary probe bbqq_vs_ccqq: (1 - AUC) of FLAV_F1R over R16_Q1, run by run :: sha256 c6c4138ee5c86423
\newcommand{\MlstCmpBbqqFonerOverCoarsePr}{at least 46\% fewer errors}  % experiments/FIGS/data/v2/index.json :: binary probe bbqq_vs_ccqq: (1 - AUC) of FLAV_F1R over R16_Q1, run by run :: sha256 c6c4138ee5c86423 :: prose twin: two significant figures, direction in words
\newcommand{\MlstCmpBbqqFonerOverCoarsePrMax}{down to 0.44 times the errors}  % experiments/FIGS/data/v2/index.json :: binary probe bbqq_vs_ccqq: (1 - AUC) of FLAV_F1R over R16_Q1, run by run :: sha256 c6c4138ee5c86423 :: prose twin of the run farther from 1
\newcommand{\MlstCmpBbqqFzeroOverFull}{6.4 and 9.4}  % experiments/FIGS/data/v2/index.json :: binary probe bbqq_vs_ccqq: (1 - AUC) of FLAV_F0 over L188, run by run :: sha256 c6c4138ee5c86423
\newcommand{\MlstCmpBbqqFzeroOverFullPr}{at least 6.4 times the errors}  % experiments/FIGS/data/v2/index.json :: binary probe bbqq_vs_ccqq: (1 - AUC) of FLAV_F0 over L188, run by run :: sha256 c6c4138ee5c86423 :: prose twin: two significant figures, direction in words
\newcommand{\MlstCmpBbqqFzeroOverFullPrMax}{up to 9.4 times the errors}  % experiments/FIGS/data/v2/index.json :: binary probe bbqq_vs_ccqq: (1 - AUC) of FLAV_F0 over L188, run by run :: sha256 c6c4138ee5c86423 :: prose twin of the run farther from 1
\newcommand{\MlstCmpBbqqFoneOverFull}{1.510 and 1.447}  % experiments/FIGS/data/v2/index.json :: binary probe bbqq_vs_ccqq: (1 - AUC) of FLAV_F1 over L188, run by run :: sha256 c6c4138ee5c86423
\newcommand{\MlstCmpBbqqFoneOverFullPr}{at least 45\% more errors}  % experiments/FIGS/data/v2/index.json :: binary probe bbqq_vs_ccqq: (1 - AUC) of FLAV_F1 over L188, run by run :: sha256 c6c4138ee5c86423 :: prose twin: two significant figures, direction in words
\newcommand{\MlstCmpBbqqFoneOverFullPrMax}{up to 51\% more errors}  % experiments/FIGS/data/v2/index.json :: binary probe bbqq_vs_ccqq: (1 - AUC) of FLAV_F1 over L188, run by run :: sha256 c6c4138ee5c86423 :: prose twin of the run farther from 1
\newcommand{\MlstCmpBbqqFonerOverFull}{2.272 and 2.281}  % experiments/FIGS/data/v2/index.json :: binary probe bbqq_vs_ccqq: (1 - AUC) of FLAV_F1R over L188, run by run :: sha256 c6c4138ee5c86423
\newcommand{\MlstCmpBbqqFonerOverFullPr}{at least 2.3 times the errors}  % experiments/FIGS/data/v2/index.json :: binary probe bbqq_vs_ccqq: (1 - AUC) of FLAV_F1R over L188, run by run :: sha256 c6c4138ee5c86423 :: prose twin: two significant figures, direction in words
\newcommand{\MlstCmpBbqqFonerOverFullPrMax}{up to 2.3 times the errors}  % experiments/FIGS/data/v2/index.json :: binary probe bbqq_vs_ccqq: (1 - AUC) of FLAV_F1R over L188, run by run :: sha256 c6c4138ee5c86423 :: prose twin of the run farther from 1
\newcommand{\MlstCmpBbqqRandOverFull}{1.42\,\ensuremath{\pm}\,0.06}  % experiments/FIGS/data/v2/index.json :: binary probe bbqq_vs_ccqq: (1 - AUC) of each random partition over the 188-class run of its seed :: sha256 c6c4138ee5c86423
\newcommand{\MlstCmpBbqqRandOverFullPr}{42\% more errors}  % experiments/FIGS/data/v2/index.json :: binary probe bbqq_vs_ccqq: (1 - AUC) of each random partition over the 188-class run of its seed :: sha256 c6c4138ee5c86423 :: prose twin: two significant figures, direction in words
\newcommand{\MlstPairRandMergedBbqqN}{2}  % experiments/FIGS/data/v2/index.json :: random partitions that merge bbqq_vs_ccqq (count) :: sha256 c6c4138ee5c86423
\newcommand{\MlstPairRandApartBbqqN}{3}  % experiments/FIGS/data/v2/index.json :: random partitions that keep apart bbqq_vs_ccqq (count) :: sha256 c6c4138ee5c86423
\newcommand{\MlstPairRandMergedOverApartBbqq}{0.996}  % experiments/FIGS/data/v2/index.json :: binary probe bbqq_vs_ccqq: mean (1 - AUC) of the partitions that merge it over that of those that keep it apart (one run each) :: sha256 c6c4138ee5c86423
\newcommand{\MlstPairRandMergedOverApartBbqqPr}{0.44\% fewer errors}  % experiments/FIGS/data/v2/index.json :: binary probe bbqq_vs_ccqq: mean (1 - AUC) of the partitions that merge it over that of those that keep it apart (one run each) :: sha256 c6c4138ee5c86423 :: prose twin: two significant figures, direction in words
\newcommand{\MlstCmpEemmResOverFull}{--}  % experiments/FIGS/data/v2/index.json :: binary probe ee_vs_mm: (1 - AUC) of L162 over L188, run by run :: sha256 c6c4138ee5c86423 :: the reference's 1 - AUC is 0
\newcommand{\MlstCmpEemmResOverFullPr}{undefined, since the reference separates the test jets completely}  % experiments/FIGS/data/v2/index.json :: binary probe ee_vs_mm: (1 - AUC) of L162 over L188, run by run :: sha256 c6c4138ee5c86423 :: the reference's 1 - AUC is 0
\newcommand{\MlstCmpEemmMidOverFull}{--}  % experiments/FIGS/data/v2/index.json :: binary probe ee_vs_mm: (1 - AUC) of R42_Q1 over L188, run by run :: sha256 c6c4138ee5c86423 :: the reference's 1 - AUC is 0
\newcommand{\MlstCmpEemmMidOverFullPr}{undefined, since the reference separates the test jets completely}  % experiments/FIGS/data/v2/index.json :: binary probe ee_vs_mm: (1 - AUC) of R42_Q1 over L188, run by run :: sha256 c6c4138ee5c86423 :: the reference's 1 - AUC is 0
\newcommand{\MlstCmpEemmCoarseOverFull}{--}  % experiments/FIGS/data/v2/index.json :: binary probe ee_vs_mm: (1 - AUC) of R16_Q1 over L188, run by run :: sha256 c6c4138ee5c86423 :: the reference's 1 - AUC is 0
\newcommand{\MlstCmpEemmCoarseOverFullPr}{undefined, since the reference separates the test jets completely}  % experiments/FIGS/data/v2/index.json :: binary probe ee_vs_mm: (1 - AUC) of R16_Q1 over L188, run by run :: sha256 c6c4138ee5c86423 :: the reference's 1 - AUC is 0
\newcommand{\MlstCmpEemmFoneOverFzero}{--}  % experiments/FIGS/data/v2/index.json :: binary probe ee_vs_mm: (1 - AUC) of FLAV_F1 over FLAV_F0, run by run :: sha256 c6c4138ee5c86423 :: the reference's 1 - AUC is 0
\newcommand{\MlstCmpEemmFoneOverFzeroPr}{undefined, since the reference separates the test jets completely}  % experiments/FIGS/data/v2/index.json :: binary probe ee_vs_mm: (1 - AUC) of FLAV_F1 over FLAV_F0, run by run :: sha256 c6c4138ee5c86423 :: the reference's 1 - AUC is 0
\newcommand{\MlstCmpEemmFoneOverFzeroPrMax}{undefined, since the reference separates the test jets completely}  % experiments/FIGS/data/v2/index.json :: binary probe ee_vs_mm: (1 - AUC) of FLAV_F1 over FLAV_F0, run by run :: sha256 c6c4138ee5c86423 :: the reference's 1 - AUC is 0
\newcommand{\MlstCmpEemmFonerOverFone}{--}  % experiments/FIGS/data/v2/index.json :: binary probe ee_vs_mm: (1 - AUC) of FLAV_F1R over FLAV_F1, run by run :: sha256 c6c4138ee5c86423 :: the reference's 1 - AUC is 0
\newcommand{\MlstCmpEemmFonerOverFonePr}{undefined, since the reference separates the test jets completely}  % experiments/FIGS/data/v2/index.json :: binary probe ee_vs_mm: (1 - AUC) of FLAV_F1R over FLAV_F1, run by run :: sha256 c6c4138ee5c86423 :: the reference's 1 - AUC is 0
\newcommand{\MlstCmpEemmFonerOverFonePrMax}{undefined, since the reference separates the test jets completely}  % experiments/FIGS/data/v2/index.json :: binary probe ee_vs_mm: (1 - AUC) of FLAV_F1R over FLAV_F1, run by run :: sha256 c6c4138ee5c86423 :: the reference's 1 - AUC is 0
\newcommand{\MlstCmpEemmFonerOverFzero}{--}  % experiments/FIGS/data/v2/index.json :: binary probe ee_vs_mm: (1 - AUC) of FLAV_F1R over FLAV_F0, run by run :: sha256 c6c4138ee5c86423 :: the reference's 1 - AUC is 0
\newcommand{\MlstCmpEemmFonerOverFzeroPr}{undefined, since the reference separates the test jets completely}  % experiments/FIGS/data/v2/index.json :: binary probe ee_vs_mm: (1 - AUC) of FLAV_F1R over FLAV_F0, run by run :: sha256 c6c4138ee5c86423 :: the reference's 1 - AUC is 0
\newcommand{\MlstCmpEemmFonerOverFzeroPrMax}{undefined, since the reference separates the test jets completely}  % experiments/FIGS/data/v2/index.json :: binary probe ee_vs_mm: (1 - AUC) of FLAV_F1R over FLAV_F0, run by run :: sha256 c6c4138ee5c86423 :: the reference's 1 - AUC is 0
\newcommand{\MlstCmpEemmFzeroOverCoarse}{0.0000022 and 0.0000003}  % experiments/FIGS/data/v2/index.json :: binary probe ee_vs_mm: (1 - AUC) of FLAV_F0 over R16_Q1, run by run :: sha256 c6c4138ee5c86423
\newcommand{\MlstCmpEemmFzeroOverCoarsePr}{at most 0.0000022 times the errors}  % experiments/FIGS/data/v2/index.json :: binary probe ee_vs_mm: (1 - AUC) of FLAV_F0 over R16_Q1, run by run :: sha256 c6c4138ee5c86423 :: prose twin: two significant figures, direction in words
\newcommand{\MlstCmpEemmFzeroOverCoarsePrMax}{down to 0.00000034 times the errors}  % experiments/FIGS/data/v2/index.json :: binary probe ee_vs_mm: (1 - AUC) of FLAV_F0 over R16_Q1, run by run :: sha256 c6c4138ee5c86423 :: prose twin of the run farther from 1
\newcommand{\MlstCmpEemmFoneOverCoarse}{0.000000 and 0.000038}  % experiments/FIGS/data/v2/index.json :: binary probe ee_vs_mm: (1 - AUC) of FLAV_F1 over R16_Q1, run by run :: sha256 c6c4138ee5c86423
\newcommand{\MlstCmpEemmFoneOverCoarsePr}{at most 0.000038 times the errors}  % experiments/FIGS/data/v2/index.json :: binary probe ee_vs_mm: (1 - AUC) of FLAV_F1 over R16_Q1, run by run :: sha256 c6c4138ee5c86423 :: prose twin: two significant figures, direction in words
\newcommand{\MlstCmpEemmFoneOverCoarsePrMax}{down to 0.000000042 times the errors}  % experiments/FIGS/data/v2/index.json :: binary probe ee_vs_mm: (1 - AUC) of FLAV_F1 over R16_Q1, run by run :: sha256 c6c4138ee5c86423 :: prose twin of the run farther from 1
\newcommand{\MlstCmpEemmFonerOverCoarse}{0.00000005 and 0.00000017}  % experiments/FIGS/data/v2/index.json :: binary probe ee_vs_mm: (1 - AUC) of FLAV_F1R over R16_Q1, run by run :: sha256 c6c4138ee5c86423
\newcommand{\MlstCmpEemmFonerOverCoarsePr}{at most 0.00000017 times the errors}  % experiments/FIGS/data/v2/index.json :: binary probe ee_vs_mm: (1 - AUC) of FLAV_F1R over R16_Q1, run by run :: sha256 c6c4138ee5c86423 :: prose twin: two significant figures, direction in words
\newcommand{\MlstCmpEemmFonerOverCoarsePrMax}{down to 0.000000055 times the errors}  % experiments/FIGS/data/v2/index.json :: binary probe ee_vs_mm: (1 - AUC) of FLAV_F1R over R16_Q1, run by run :: sha256 c6c4138ee5c86423 :: prose twin of the run farther from 1
\newcommand{\MlstCmpEemmFzeroOverFull}{--}  % experiments/FIGS/data/v2/index.json :: binary probe ee_vs_mm: (1 - AUC) of FLAV_F0 over L188, run by run :: sha256 c6c4138ee5c86423 :: the reference's 1 - AUC is 0
\newcommand{\MlstCmpEemmFzeroOverFullPr}{undefined, since the reference separates the test jets completely}  % experiments/FIGS/data/v2/index.json :: binary probe ee_vs_mm: (1 - AUC) of FLAV_F0 over L188, run by run :: sha256 c6c4138ee5c86423 :: the reference's 1 - AUC is 0
\newcommand{\MlstCmpEemmFzeroOverFullPrMax}{undefined, since the reference separates the test jets completely}  % experiments/FIGS/data/v2/index.json :: binary probe ee_vs_mm: (1 - AUC) of FLAV_F0 over L188, run by run :: sha256 c6c4138ee5c86423 :: the reference's 1 - AUC is 0
\newcommand{\MlstCmpEemmFoneOverFull}{--}  % experiments/FIGS/data/v2/index.json :: binary probe ee_vs_mm: (1 - AUC) of FLAV_F1 over L188, run by run :: sha256 c6c4138ee5c86423 :: the reference's 1 - AUC is 0
\newcommand{\MlstCmpEemmFoneOverFullPr}{undefined, since the reference separates the test jets completely}  % experiments/FIGS/data/v2/index.json :: binary probe ee_vs_mm: (1 - AUC) of FLAV_F1 over L188, run by run :: sha256 c6c4138ee5c86423 :: the reference's 1 - AUC is 0
\newcommand{\MlstCmpEemmFoneOverFullPrMax}{undefined, since the reference separates the test jets completely}  % experiments/FIGS/data/v2/index.json :: binary probe ee_vs_mm: (1 - AUC) of FLAV_F1 over L188, run by run :: sha256 c6c4138ee5c86423 :: the reference's 1 - AUC is 0
\newcommand{\MlstCmpEemmFonerOverFull}{--}  % experiments/FIGS/data/v2/index.json :: binary probe ee_vs_mm: (1 - AUC) of FLAV_F1R over L188, run by run :: sha256 c6c4138ee5c86423 :: the reference's 1 - AUC is 0
\newcommand{\MlstCmpEemmFonerOverFullPr}{undefined, since the reference separates the test jets completely}  % experiments/FIGS/data/v2/index.json :: binary probe ee_vs_mm: (1 - AUC) of FLAV_F1R over L188, run by run :: sha256 c6c4138ee5c86423 :: the reference's 1 - AUC is 0
\newcommand{\MlstCmpEemmFonerOverFullPrMax}{undefined, since the reference separates the test jets completely}  % experiments/FIGS/data/v2/index.json :: binary probe ee_vs_mm: (1 - AUC) of FLAV_F1R over L188, run by run :: sha256 c6c4138ee5c86423 :: the reference's 1 - AUC is 0
\newcommand{\MlstCmpEemmRandOverFull}{--}  % experiments/FIGS/data/v2/index.json :: binary probe ee_vs_mm: (1 - AUC) of each random partition over the 188-class run of its seed :: sha256 c6c4138ee5c86423 :: the reference's 1 - AUC is 0
\newcommand{\MlstCmpEemmRandOverFullPr}{undefined, since the reference separates the test jets completely}  % experiments/FIGS/data/v2/index.json :: binary probe ee_vs_mm: (1 - AUC) of each random partition over the 188-class run of its seed :: sha256 c6c4138ee5c86423 :: the reference's 1 - AUC is 0
\newcommand{\MlstPairRandMergedEemmN}{3}  % experiments/FIGS/data/v2/index.json :: random partitions that merge ee_vs_mm (count) :: sha256 c6c4138ee5c86423
\newcommand{\MlstPairRandApartEemmN}{2}  % experiments/FIGS/data/v2/index.json :: random partitions that keep apart ee_vs_mm (count) :: sha256 c6c4138ee5c86423
\newcommand{\MlstPairRandMergedOverApartEemm}{75.3}  % experiments/FIGS/data/v2/index.json :: binary probe ee_vs_mm: mean (1 - AUC) of the partitions that merge it over that of those that keep it apart (one run each) :: sha256 c6c4138ee5c86423
\newcommand{\MlstPairRandMergedOverApartEemmPr}{75 times the errors}  % experiments/FIGS/data/v2/index.json :: binary probe ee_vs_mm: mean (1 - AUC) of the partitions that merge it over that of those that keep it apart (one run each) :: sha256 c6c4138ee5c86423 :: prose twin: two significant figures, direction in words
\newcommand{\MlstCmpQcdBbResOverFull}{1.15\,\ensuremath{\pm}\,0.04}  % experiments/FIGS/data/v2/index.json :: binary probe qcd_bb_vs_cc: (1 - AUC) of L162 over L188, run by run :: sha256 c6c4138ee5c86423
\newcommand{\MlstCmpQcdBbResOverFullPr}{15\% more errors}  % experiments/FIGS/data/v2/index.json :: binary probe qcd_bb_vs_cc: (1 - AUC) of L162 over L188, run by run :: sha256 c6c4138ee5c86423 :: prose twin: two significant figures, direction in words
\newcommand{\MlstCmpQcdBbMidOverFull}{1.163\,\ensuremath{\pm}\,0.016}  % experiments/FIGS/data/v2/index.json :: binary probe qcd_bb_vs_cc: (1 - AUC) of R42_Q1 over L188, run by run :: sha256 c6c4138ee5c86423
\newcommand{\MlstCmpQcdBbMidOverFullPr}{16\% more errors}  % experiments/FIGS/data/v2/index.json :: binary probe qcd_bb_vs_cc: (1 - AUC) of R42_Q1 over L188, run by run :: sha256 c6c4138ee5c86423 :: prose twin: two significant figures, direction in words
\newcommand{\MlstCmpQcdBbCoarseOverFull}{3.294\,\ensuremath{\pm}\,0.033}  % experiments/FIGS/data/v2/index.json :: binary probe qcd_bb_vs_cc: (1 - AUC) of R16_Q1 over L188, run by run :: sha256 c6c4138ee5c86423
\newcommand{\MlstCmpQcdBbCoarseOverFullPr}{3.3 times the errors}  % experiments/FIGS/data/v2/index.json :: binary probe qcd_bb_vs_cc: (1 - AUC) of R16_Q1 over L188, run by run :: sha256 c6c4138ee5c86423 :: prose twin: two significant figures, direction in words
\newcommand{\MlstCmpQcdBbFoneOverFzero}{0.368 and 0.294}  % experiments/FIGS/data/v2/index.json :: binary probe qcd_bb_vs_cc: (1 - AUC) of FLAV_F1 over FLAV_F0, run by run :: sha256 c6c4138ee5c86423
\newcommand{\MlstCmpQcdBbFoneOverFzeroPr}{at most 0.37 times the errors}  % experiments/FIGS/data/v2/index.json :: binary probe qcd_bb_vs_cc: (1 - AUC) of FLAV_F1 over FLAV_F0, run by run :: sha256 c6c4138ee5c86423 :: prose twin: two significant figures, direction in words
\newcommand{\MlstCmpQcdBbFoneOverFzeroPrMax}{down to 0.29 times the errors}  % experiments/FIGS/data/v2/index.json :: binary probe qcd_bb_vs_cc: (1 - AUC) of FLAV_F1 over FLAV_F0, run by run :: sha256 c6c4138ee5c86423 :: prose twin of the run farther from 1
\newcommand{\MlstCmpQcdBbFonerOverFone}{1.361 and 1.429}  % experiments/FIGS/data/v2/index.json :: binary probe qcd_bb_vs_cc: (1 - AUC) of FLAV_F1R over FLAV_F1, run by run :: sha256 c6c4138ee5c86423
\newcommand{\MlstCmpQcdBbFonerOverFonePr}{at least 36\% more errors}  % experiments/FIGS/data/v2/index.json :: binary probe qcd_bb_vs_cc: (1 - AUC) of FLAV_F1R over FLAV_F1, run by run :: sha256 c6c4138ee5c86423 :: prose twin: two significant figures, direction in words
\newcommand{\MlstCmpQcdBbFonerOverFonePrMax}{up to 43\% more errors}  % experiments/FIGS/data/v2/index.json :: binary probe qcd_bb_vs_cc: (1 - AUC) of FLAV_F1R over FLAV_F1, run by run :: sha256 c6c4138ee5c86423 :: prose twin of the run farther from 1
\newcommand{\MlstCmpQcdBbFonerOverFzero}{0.501 and 0.420}  % experiments/FIGS/data/v2/index.json :: binary probe qcd_bb_vs_cc: (1 - AUC) of FLAV_F1R over FLAV_F0, run by run :: sha256 c6c4138ee5c86423
\newcommand{\MlstCmpQcdBbFonerOverFzeroPr}{at least 50\% fewer errors}  % experiments/FIGS/data/v2/index.json :: binary probe qcd_bb_vs_cc: (1 - AUC) of FLAV_F1R over FLAV_F0, run by run :: sha256 c6c4138ee5c86423 :: prose twin: two significant figures, direction in words
\newcommand{\MlstCmpQcdBbFonerOverFzeroPrMax}{down to 0.42 times the errors}  % experiments/FIGS/data/v2/index.json :: binary probe qcd_bb_vs_cc: (1 - AUC) of FLAV_F1R over FLAV_F0, run by run :: sha256 c6c4138ee5c86423 :: prose twin of the run farther from 1
\newcommand{\MlstCmpQcdBbFzeroOverCoarse}{1.28 and 1.58}  % experiments/FIGS/data/v2/index.json :: binary probe qcd_bb_vs_cc: (1 - AUC) of FLAV_F0 over R16_Q1, run by run :: sha256 c6c4138ee5c86423
\newcommand{\MlstCmpQcdBbFzeroOverCoarsePr}{at least 28\% more errors}  % experiments/FIGS/data/v2/index.json :: binary probe qcd_bb_vs_cc: (1 - AUC) of FLAV_F0 over R16_Q1, run by run :: sha256 c6c4138ee5c86423 :: prose twin: two significant figures, direction in words
\newcommand{\MlstCmpQcdBbFzeroOverCoarsePrMax}{up to 58\% more errors}  % experiments/FIGS/data/v2/index.json :: binary probe qcd_bb_vs_cc: (1 - AUC) of FLAV_F0 over R16_Q1, run by run :: sha256 c6c4138ee5c86423 :: prose twin of the run farther from 1
\newcommand{\MlstCmpQcdBbFoneOverCoarse}{0.4692 and 0.4648}  % experiments/FIGS/data/v2/index.json :: binary probe qcd_bb_vs_cc: (1 - AUC) of FLAV_F1 over R16_Q1, run by run :: sha256 c6c4138ee5c86423
\newcommand{\MlstCmpQcdBbFoneOverCoarsePr}{at most 0.47 times the errors}  % experiments/FIGS/data/v2/index.json :: binary probe qcd_bb_vs_cc: (1 - AUC) of FLAV_F1 over R16_Q1, run by run :: sha256 c6c4138ee5c86423 :: prose twin: two significant figures, direction in words
\newcommand{\MlstCmpQcdBbFoneOverCoarsePrMax}{down to 0.46 times the errors}  % experiments/FIGS/data/v2/index.json :: binary probe qcd_bb_vs_cc: (1 - AUC) of FLAV_F1 over R16_Q1, run by run :: sha256 c6c4138ee5c86423 :: prose twin of the run farther from 1
\newcommand{\MlstCmpQcdBbFonerOverCoarse}{0.639 and 0.664}  % experiments/FIGS/data/v2/index.json :: binary probe qcd_bb_vs_cc: (1 - AUC) of FLAV_F1R over R16_Q1, run by run :: sha256 c6c4138ee5c86423
\newcommand{\MlstCmpQcdBbFonerOverCoarsePr}{at least 34\% fewer errors}  % experiments/FIGS/data/v2/index.json :: binary probe qcd_bb_vs_cc: (1 - AUC) of FLAV_F1R over R16_Q1, run by run :: sha256 c6c4138ee5c86423 :: prose twin: two significant figures, direction in words
\newcommand{\MlstCmpQcdBbFonerOverCoarsePrMax}{up to 36\% fewer errors}  % experiments/FIGS/data/v2/index.json :: binary probe qcd_bb_vs_cc: (1 - AUC) of FLAV_F1R over R16_Q1, run by run :: sha256 c6c4138ee5c86423 :: prose twin of the run farther from 1
\newcommand{\MlstCmpQcdBbFzeroOverFull}{4.24 and 5.21}  % experiments/FIGS/data/v2/index.json :: binary probe qcd_bb_vs_cc: (1 - AUC) of FLAV_F0 over L188, run by run :: sha256 c6c4138ee5c86423
\newcommand{\MlstCmpQcdBbFzeroOverFullPr}{at least 4.2 times the errors}  % experiments/FIGS/data/v2/index.json :: binary probe qcd_bb_vs_cc: (1 - AUC) of FLAV_F0 over L188, run by run :: sha256 c6c4138ee5c86423 :: prose twin: two significant figures, direction in words
\newcommand{\MlstCmpQcdBbFzeroOverFullPrMax}{up to 5.2 times the errors}  % experiments/FIGS/data/v2/index.json :: binary probe qcd_bb_vs_cc: (1 - AUC) of FLAV_F0 over L188, run by run :: sha256 c6c4138ee5c86423 :: prose twin of the run farther from 1
\newcommand{\MlstCmpQcdBbFoneOverFull}{1.561 and 1.532}  % experiments/FIGS/data/v2/index.json :: binary probe qcd_bb_vs_cc: (1 - AUC) of FLAV_F1 over L188, run by run :: sha256 c6c4138ee5c86423
\newcommand{\MlstCmpQcdBbFoneOverFullPr}{at least 53\% more errors}  % experiments/FIGS/data/v2/index.json :: binary probe qcd_bb_vs_cc: (1 - AUC) of FLAV_F1 over L188, run by run :: sha256 c6c4138ee5c86423 :: prose twin: two significant figures, direction in words
\newcommand{\MlstCmpQcdBbFoneOverFullPrMax}{up to 56\% more errors}  % experiments/FIGS/data/v2/index.json :: binary probe qcd_bb_vs_cc: (1 - AUC) of FLAV_F1 over L188, run by run :: sha256 c6c4138ee5c86423 :: prose twin of the run farther from 1
\newcommand{\MlstCmpQcdBbFonerOverFull}{2.124 and 2.189}  % experiments/FIGS/data/v2/index.json :: binary probe qcd_bb_vs_cc: (1 - AUC) of FLAV_F1R over L188, run by run :: sha256 c6c4138ee5c86423
\newcommand{\MlstCmpQcdBbFonerOverFullPr}{at least 2.1 times the errors}  % experiments/FIGS/data/v2/index.json :: binary probe qcd_bb_vs_cc: (1 - AUC) of FLAV_F1R over L188, run by run :: sha256 c6c4138ee5c86423 :: prose twin: two significant figures, direction in words
\newcommand{\MlstCmpQcdBbFonerOverFullPrMax}{up to 2.2 times the errors}  % experiments/FIGS/data/v2/index.json :: binary probe qcd_bb_vs_cc: (1 - AUC) of FLAV_F1R over L188, run by run :: sha256 c6c4138ee5c86423 :: prose twin of the run farther from 1
\newcommand{\MlstCmpQcdBResOverFull}{1.085\,\ensuremath{\pm}\,0.004}  % experiments/FIGS/data/v2/index.json :: binary probe qcd_b_vs_light: (1 - AUC) of L162 over L188, run by run :: sha256 c6c4138ee5c86423
\newcommand{\MlstCmpQcdBResOverFullPr}{8.5\% more errors}  % experiments/FIGS/data/v2/index.json :: binary probe qcd_b_vs_light: (1 - AUC) of L162 over L188, run by run :: sha256 c6c4138ee5c86423 :: prose twin: two significant figures, direction in words
\newcommand{\MlstCmpQcdBMidOverFull}{1.081\,\ensuremath{\pm}\,0.008}  % experiments/FIGS/data/v2/index.json :: binary probe qcd_b_vs_light: (1 - AUC) of R42_Q1 over L188, run by run :: sha256 c6c4138ee5c86423
\newcommand{\MlstCmpQcdBMidOverFullPr}{8.1\% more errors}  % experiments/FIGS/data/v2/index.json :: binary probe qcd_b_vs_light: (1 - AUC) of R42_Q1 over L188, run by run :: sha256 c6c4138ee5c86423 :: prose twin: two significant figures, direction in words
\newcommand{\MlstCmpQcdBCoarseOverFull}{1.54\,\ensuremath{\pm}\,0.07}  % experiments/FIGS/data/v2/index.json :: binary probe qcd_b_vs_light: (1 - AUC) of R16_Q1 over L188, run by run :: sha256 c6c4138ee5c86423
\newcommand{\MlstCmpQcdBCoarseOverFullPr}{54\% more errors}  % experiments/FIGS/data/v2/index.json :: binary probe qcd_b_vs_light: (1 - AUC) of R16_Q1 over L188, run by run :: sha256 c6c4138ee5c86423 :: prose twin: two significant figures, direction in words
\newcommand{\MlstCmpQcdBFoneOverFzero}{0.778 and 0.835}  % experiments/FIGS/data/v2/index.json :: binary probe qcd_b_vs_light: (1 - AUC) of FLAV_F1 over FLAV_F0, run by run :: sha256 c6c4138ee5c86423
\newcommand{\MlstCmpQcdBFoneOverFzeroPr}{at least 17\% fewer errors}  % experiments/FIGS/data/v2/index.json :: binary probe qcd_b_vs_light: (1 - AUC) of FLAV_F1 over FLAV_F0, run by run :: sha256 c6c4138ee5c86423 :: prose twin: two significant figures, direction in words
\newcommand{\MlstCmpQcdBFoneOverFzeroPrMax}{up to 22\% fewer errors}  % experiments/FIGS/data/v2/index.json :: binary probe qcd_b_vs_light: (1 - AUC) of FLAV_F1 over FLAV_F0, run by run :: sha256 c6c4138ee5c86423 :: prose twin of the run farther from 1
\newcommand{\MlstCmpQcdBFonerOverFone}{1.054 and 1.061}  % experiments/FIGS/data/v2/index.json :: binary probe qcd_b_vs_light: (1 - AUC) of FLAV_F1R over FLAV_F1, run by run :: sha256 c6c4138ee5c86423
\newcommand{\MlstCmpQcdBFonerOverFonePr}{equal within 7\%}  % experiments/FIGS/data/v2/index.json :: binary probe qcd_b_vs_light: (1 - AUC) of FLAV_F1R over FLAV_F1, run by run :: sha256 c6c4138ee5c86423 :: prose twin: two significant figures, direction in words
\newcommand{\MlstCmpQcdBFonerOverFonePrMax}{equal within 7\%}  % experiments/FIGS/data/v2/index.json :: binary probe qcd_b_vs_light: (1 - AUC) of FLAV_F1R over FLAV_F1, run by run :: sha256 c6c4138ee5c86423 :: prose twin of the run farther from 1
\newcommand{\MlstCmpQcdBFonerOverFzero}{0.820 and 0.885}  % experiments/FIGS/data/v2/index.json :: binary probe qcd_b_vs_light: (1 - AUC) of FLAV_F1R over FLAV_F0, run by run :: sha256 c6c4138ee5c86423
\newcommand{\MlstCmpQcdBFonerOverFzeroPr}{at least 11\% fewer errors}  % experiments/FIGS/data/v2/index.json :: binary probe qcd_b_vs_light: (1 - AUC) of FLAV_F1R over FLAV_F0, run by run :: sha256 c6c4138ee5c86423 :: prose twin: two significant figures, direction in words
\newcommand{\MlstCmpQcdBFonerOverFzeroPrMax}{up to 18\% fewer errors}  % experiments/FIGS/data/v2/index.json :: binary probe qcd_b_vs_light: (1 - AUC) of FLAV_F1R over FLAV_F0, run by run :: sha256 c6c4138ee5c86423 :: prose twin of the run farther from 1
\newcommand{\MlstCmpQcdBFzeroOverCoarse}{0.962 and 0.981}  % experiments/FIGS/data/v2/index.json :: binary probe qcd_b_vs_light: (1 - AUC) of FLAV_F0 over R16_Q1, run by run :: sha256 c6c4138ee5c86423
\newcommand{\MlstCmpQcdBFzeroOverCoarsePr}{equal within 4\%}  % experiments/FIGS/data/v2/index.json :: binary probe qcd_b_vs_light: (1 - AUC) of FLAV_F0 over R16_Q1, run by run :: sha256 c6c4138ee5c86423 :: prose twin: two significant figures, direction in words
\newcommand{\MlstCmpQcdBFzeroOverCoarsePrMax}{equal within 4\%}  % experiments/FIGS/data/v2/index.json :: binary probe qcd_b_vs_light: (1 - AUC) of FLAV_F0 over R16_Q1, run by run :: sha256 c6c4138ee5c86423 :: prose twin of the run farther from 1
\newcommand{\MlstCmpQcdBFoneOverCoarse}{0.749 and 0.819}  % experiments/FIGS/data/v2/index.json :: binary probe qcd_b_vs_light: (1 - AUC) of FLAV_F1 over R16_Q1, run by run :: sha256 c6c4138ee5c86423
\newcommand{\MlstCmpQcdBFoneOverCoarsePr}{at least 18\% fewer errors}  % experiments/FIGS/data/v2/index.json :: binary probe qcd_b_vs_light: (1 - AUC) of FLAV_F1 over R16_Q1, run by run :: sha256 c6c4138ee5c86423 :: prose twin: two significant figures, direction in words
\newcommand{\MlstCmpQcdBFoneOverCoarsePrMax}{up to 25\% fewer errors}  % experiments/FIGS/data/v2/index.json :: binary probe qcd_b_vs_light: (1 - AUC) of FLAV_F1 over R16_Q1, run by run :: sha256 c6c4138ee5c86423 :: prose twin of the run farther from 1
\newcommand{\MlstCmpQcdBFonerOverCoarse}{0.789 and 0.869}  % experiments/FIGS/data/v2/index.json :: binary probe qcd_b_vs_light: (1 - AUC) of FLAV_F1R over R16_Q1, run by run :: sha256 c6c4138ee5c86423
\newcommand{\MlstCmpQcdBFonerOverCoarsePr}{at least 13\% fewer errors}  % experiments/FIGS/data/v2/index.json :: binary probe qcd_b_vs_light: (1 - AUC) of FLAV_F1R over R16_Q1, run by run :: sha256 c6c4138ee5c86423 :: prose twin: two significant figures, direction in words
\newcommand{\MlstCmpQcdBFonerOverCoarsePrMax}{up to 21\% fewer errors}  % experiments/FIGS/data/v2/index.json :: binary probe qcd_b_vs_light: (1 - AUC) of FLAV_F1R over R16_Q1, run by run :: sha256 c6c4138ee5c86423 :: prose twin of the run farther from 1
\newcommand{\MlstCmpQcdBFzeroOverFull}{1.56 and 1.45}  % experiments/FIGS/data/v2/index.json :: binary probe qcd_b_vs_light: (1 - AUC) of FLAV_F0 over L188, run by run :: sha256 c6c4138ee5c86423
\newcommand{\MlstCmpQcdBFzeroOverFullPr}{at least 45\% more errors}  % experiments/FIGS/data/v2/index.json :: binary probe qcd_b_vs_light: (1 - AUC) of FLAV_F0 over L188, run by run :: sha256 c6c4138ee5c86423 :: prose twin: two significant figures, direction in words
\newcommand{\MlstCmpQcdBFzeroOverFullPrMax}{up to 56\% more errors}  % experiments/FIGS/data/v2/index.json :: binary probe qcd_b_vs_light: (1 - AUC) of FLAV_F0 over L188, run by run :: sha256 c6c4138ee5c86423 :: prose twin of the run farther from 1
\newcommand{\MlstCmpQcdBFoneOverFull}{1.212 and 1.213}  % experiments/FIGS/data/v2/index.json :: binary probe qcd_b_vs_light: (1 - AUC) of FLAV_F1 over L188, run by run :: sha256 c6c4138ee5c86423
\newcommand{\MlstCmpQcdBFoneOverFullPr}{at least 21\% more errors}  % experiments/FIGS/data/v2/index.json :: binary probe qcd_b_vs_light: (1 - AUC) of FLAV_F1 over L188, run by run :: sha256 c6c4138ee5c86423 :: prose twin: two significant figures, direction in words
\newcommand{\MlstCmpQcdBFoneOverFullPrMax}{up to 21\% more errors}  % experiments/FIGS/data/v2/index.json :: binary probe qcd_b_vs_light: (1 - AUC) of FLAV_F1 over L188, run by run :: sha256 c6c4138ee5c86423 :: prose twin of the run farther from 1
\newcommand{\MlstCmpQcdBFonerOverFull}{1.277 and 1.286}  % experiments/FIGS/data/v2/index.json :: binary probe qcd_b_vs_light: (1 - AUC) of FLAV_F1R over L188, run by run :: sha256 c6c4138ee5c86423
\newcommand{\MlstCmpQcdBFonerOverFullPr}{at least 28\% more errors}  % experiments/FIGS/data/v2/index.json :: binary probe qcd_b_vs_light: (1 - AUC) of FLAV_F1R over L188, run by run :: sha256 c6c4138ee5c86423 :: prose twin: two significant figures, direction in words
\newcommand{\MlstCmpQcdBFonerOverFullPrMax}{up to 29\% more errors}  % experiments/FIGS/data/v2/index.json :: binary probe qcd_b_vs_light: (1 - AUC) of FLAV_F1R over L188, run by run :: sha256 c6c4138ee5c86423 :: prose twin of the run farther from 1
\newcommand{\MlstCmpQcdCResOverFull}{1.061\,\ensuremath{\pm}\,0.005}  % experiments/FIGS/data/v2/index.json :: binary probe qcd_c_vs_light: (1 - AUC) of L162 over L188, run by run :: sha256 c6c4138ee5c86423
\newcommand{\MlstCmpQcdCResOverFullPr}{6.1\% more errors}  % experiments/FIGS/data/v2/index.json :: binary probe qcd_c_vs_light: (1 - AUC) of L162 over L188, run by run :: sha256 c6c4138ee5c86423 :: prose twin: two significant figures, direction in words
\newcommand{\MlstCmpQcdCMidOverFull}{1.072\,\ensuremath{\pm}\,0.008}  % experiments/FIGS/data/v2/index.json :: binary probe qcd_c_vs_light: (1 - AUC) of R42_Q1 over L188, run by run :: sha256 c6c4138ee5c86423
\newcommand{\MlstCmpQcdCMidOverFullPr}{7.2\% more errors}  % experiments/FIGS/data/v2/index.json :: binary probe qcd_c_vs_light: (1 - AUC) of R42_Q1 over L188, run by run :: sha256 c6c4138ee5c86423 :: prose twin: two significant figures, direction in words
\newcommand{\MlstCmpQcdCCoarseOverFull}{1.407\,\ensuremath{\pm}\,0.021}  % experiments/FIGS/data/v2/index.json :: binary probe qcd_c_vs_light: (1 - AUC) of R16_Q1 over L188, run by run :: sha256 c6c4138ee5c86423
\newcommand{\MlstCmpQcdCCoarseOverFullPr}{41\% more errors}  % experiments/FIGS/data/v2/index.json :: binary probe qcd_c_vs_light: (1 - AUC) of R16_Q1 over L188, run by run :: sha256 c6c4138ee5c86423 :: prose twin: two significant figures, direction in words
\newcommand{\MlstCmpQcdCFoneOverFzero}{0.915 and 0.944}  % experiments/FIGS/data/v2/index.json :: binary probe qcd_c_vs_light: (1 - AUC) of FLAV_F1 over FLAV_F0, run by run :: sha256 c6c4138ee5c86423
\newcommand{\MlstCmpQcdCFoneOverFzeroPr}{at least 5.6\% fewer errors}  % experiments/FIGS/data/v2/index.json :: binary probe qcd_c_vs_light: (1 - AUC) of FLAV_F1 over FLAV_F0, run by run :: sha256 c6c4138ee5c86423 :: prose twin: two significant figures, direction in words
\newcommand{\MlstCmpQcdCFoneOverFzeroPrMax}{up to 8.5\% fewer errors}  % experiments/FIGS/data/v2/index.json :: binary probe qcd_c_vs_light: (1 - AUC) of FLAV_F1 over FLAV_F0, run by run :: sha256 c6c4138ee5c86423 :: prose twin of the run farther from 1
\newcommand{\MlstCmpQcdCFonerOverFone}{0.9535 and 0.9520}  % experiments/FIGS/data/v2/index.json :: binary probe qcd_c_vs_light: (1 - AUC) of FLAV_F1R over FLAV_F1, run by run :: sha256 c6c4138ee5c86423
\newcommand{\MlstCmpQcdCFonerOverFonePr}{at least 4.6\% fewer errors}  % experiments/FIGS/data/v2/index.json :: binary probe qcd_c_vs_light: (1 - AUC) of FLAV_F1R over FLAV_F1, run by run :: sha256 c6c4138ee5c86423 :: prose twin: two significant figures, direction in words
\newcommand{\MlstCmpQcdCFonerOverFonePrMax}{up to 4.8\% fewer errors}  % experiments/FIGS/data/v2/index.json :: binary probe qcd_c_vs_light: (1 - AUC) of FLAV_F1R over FLAV_F1, run by run :: sha256 c6c4138ee5c86423 :: prose twin of the run farther from 1
\newcommand{\MlstCmpQcdCFonerOverFzero}{0.872 and 0.899}  % experiments/FIGS/data/v2/index.json :: binary probe qcd_c_vs_light: (1 - AUC) of FLAV_F1R over FLAV_F0, run by run :: sha256 c6c4138ee5c86423
\newcommand{\MlstCmpQcdCFonerOverFzeroPr}{at least 10\% fewer errors}  % experiments/FIGS/data/v2/index.json :: binary probe qcd_c_vs_light: (1 - AUC) of FLAV_F1R over FLAV_F0, run by run :: sha256 c6c4138ee5c86423 :: prose twin: two significant figures, direction in words
\newcommand{\MlstCmpQcdCFonerOverFzeroPrMax}{up to 13\% fewer errors}  % experiments/FIGS/data/v2/index.json :: binary probe qcd_c_vs_light: (1 - AUC) of FLAV_F1R over FLAV_F0, run by run :: sha256 c6c4138ee5c86423 :: prose twin of the run farther from 1
\newcommand{\MlstCmpQcdCFzeroOverCoarse}{0.952 and 0.934}  % experiments/FIGS/data/v2/index.json :: binary probe qcd_c_vs_light: (1 - AUC) of FLAV_F0 over R16_Q1, run by run :: sha256 c6c4138ee5c86423
\newcommand{\MlstCmpQcdCFzeroOverCoarsePr}{at least 4.8\% fewer errors}  % experiments/FIGS/data/v2/index.json :: binary probe qcd_c_vs_light: (1 - AUC) of FLAV_F0 over R16_Q1, run by run :: sha256 c6c4138ee5c86423 :: prose twin: two significant figures, direction in words
\newcommand{\MlstCmpQcdCFzeroOverCoarsePrMax}{up to 6.6\% fewer errors}  % experiments/FIGS/data/v2/index.json :: binary probe qcd_c_vs_light: (1 - AUC) of FLAV_F0 over R16_Q1, run by run :: sha256 c6c4138ee5c86423 :: prose twin of the run farther from 1
\newcommand{\MlstCmpQcdCFoneOverCoarse}{0.871 and 0.882}  % experiments/FIGS/data/v2/index.json :: binary probe qcd_c_vs_light: (1 - AUC) of FLAV_F1 over R16_Q1, run by run :: sha256 c6c4138ee5c86423
\newcommand{\MlstCmpQcdCFoneOverCoarsePr}{at least 12\% fewer errors}  % experiments/FIGS/data/v2/index.json :: binary probe qcd_c_vs_light: (1 - AUC) of FLAV_F1 over R16_Q1, run by run :: sha256 c6c4138ee5c86423 :: prose twin: two significant figures, direction in words
\newcommand{\MlstCmpQcdCFoneOverCoarsePrMax}{up to 13\% fewer errors}  % experiments/FIGS/data/v2/index.json :: binary probe qcd_c_vs_light: (1 - AUC) of FLAV_F1 over R16_Q1, run by run :: sha256 c6c4138ee5c86423 :: prose twin of the run farther from 1
\newcommand{\MlstCmpQcdCFonerOverCoarse}{0.8301 and 0.8396}  % experiments/FIGS/data/v2/index.json :: binary probe qcd_c_vs_light: (1 - AUC) of FLAV_F1R over R16_Q1, run by run :: sha256 c6c4138ee5c86423
\newcommand{\MlstCmpQcdCFonerOverCoarsePr}{at least 16\% fewer errors}  % experiments/FIGS/data/v2/index.json :: binary probe qcd_c_vs_light: (1 - AUC) of FLAV_F1R over R16_Q1, run by run :: sha256 c6c4138ee5c86423 :: prose twin: two significant figures, direction in words
\newcommand{\MlstCmpQcdCFonerOverCoarsePrMax}{up to 17\% fewer errors}  % experiments/FIGS/data/v2/index.json :: binary probe qcd_c_vs_light: (1 - AUC) of FLAV_F1R over R16_Q1, run by run :: sha256 c6c4138ee5c86423 :: prose twin of the run farther from 1
\newcommand{\MlstCmpQcdCFzeroOverFull}{1.352 and 1.293}  % experiments/FIGS/data/v2/index.json :: binary probe qcd_c_vs_light: (1 - AUC) of FLAV_F0 over L188, run by run :: sha256 c6c4138ee5c86423
\newcommand{\MlstCmpQcdCFzeroOverFullPr}{at least 29\% more errors}  % experiments/FIGS/data/v2/index.json :: binary probe qcd_c_vs_light: (1 - AUC) of FLAV_F0 over L188, run by run :: sha256 c6c4138ee5c86423 :: prose twin: two significant figures, direction in words
\newcommand{\MlstCmpQcdCFzeroOverFullPrMax}{up to 35\% more errors}  % experiments/FIGS/data/v2/index.json :: binary probe qcd_c_vs_light: (1 - AUC) of FLAV_F0 over L188, run by run :: sha256 c6c4138ee5c86423 :: prose twin of the run farther from 1
\newcommand{\MlstCmpQcdCFoneOverFull}{1.237 and 1.220}  % experiments/FIGS/data/v2/index.json :: binary probe qcd_c_vs_light: (1 - AUC) of FLAV_F1 over L188, run by run :: sha256 c6c4138ee5c86423
\newcommand{\MlstCmpQcdCFoneOverFullPr}{at least 22\% more errors}  % experiments/FIGS/data/v2/index.json :: binary probe qcd_c_vs_light: (1 - AUC) of FLAV_F1 over L188, run by run :: sha256 c6c4138ee5c86423 :: prose twin: two significant figures, direction in words
\newcommand{\MlstCmpQcdCFoneOverFullPrMax}{up to 24\% more errors}  % experiments/FIGS/data/v2/index.json :: binary probe qcd_c_vs_light: (1 - AUC) of FLAV_F1 over L188, run by run :: sha256 c6c4138ee5c86423 :: prose twin of the run farther from 1
\newcommand{\MlstCmpQcdCFonerOverFull}{1.179 and 1.162}  % experiments/FIGS/data/v2/index.json :: binary probe qcd_c_vs_light: (1 - AUC) of FLAV_F1R over L188, run by run :: sha256 c6c4138ee5c86423
\newcommand{\MlstCmpQcdCFonerOverFullPr}{at least 16\% more errors}  % experiments/FIGS/data/v2/index.json :: binary probe qcd_c_vs_light: (1 - AUC) of FLAV_F1R over L188, run by run :: sha256 c6c4138ee5c86423 :: prose twin: two significant figures, direction in words
\newcommand{\MlstCmpQcdCFonerOverFullPrMax}{up to 18\% more errors}  % experiments/FIGS/data/v2/index.json :: binary probe qcd_c_vs_light: (1 - AUC) of FLAV_F1R over L188, run by run :: sha256 c6c4138ee5c86423 :: prose twin of the run farther from 1
\newcommand{\MlstLpEemmAucFull}{1.0000\,\ensuremath{\pm}\,\ensuremath{<}0.0001}  % experiments/FIGS/data/v2/index.json :: linear probe ee_vs_mm (fits_v4_sib): AUC, L188 :: sha256 c6c4138ee5c86423
\newcommand{\MlstLpEemmAucRes}{1.0000\,\ensuremath{\pm}\,\ensuremath{<}0.0001}  % experiments/FIGS/data/v2/index.json :: linear probe ee_vs_mm (fits_v4_sib): AUC, L162 :: sha256 c6c4138ee5c86423
\newcommand{\MlstLpEemmAucMid}{0.9610\,\ensuremath{\pm}\,0.0030}  % experiments/FIGS/data/v2/index.json :: linear probe ee_vs_mm (fits_v4_sib): AUC, R42_Q1 :: sha256 c6c4138ee5c86423
\newcommand{\MlstLpEemmAucCoarse}{0.960\,\ensuremath{\pm}\,0.020}  % experiments/FIGS/data/v2/index.json :: linear probe ee_vs_mm (fits_v4_sib): AUC, R16_Q1 :: sha256 c6c4138ee5c86423
\newcommand{\MlstLpEemmAucFzero}{1.0000 and 1.0000}  % experiments/FIGS/data/v2/index.json :: linear probe ee_vs_mm (fits_v4_sib): AUC, FLAV_F0 :: sha256 c6c4138ee5c86423
\newcommand{\MlstLpEemmAucFone}{1.0000 and 1.0000}  % experiments/FIGS/data/v2/index.json :: linear probe ee_vs_mm (fits_v4_sib): AUC, FLAV_F1 :: sha256 c6c4138ee5c86423
\newcommand{\MlstLpEemmAucFoner}{1.0000 and 1.0000}  % experiments/FIGS/data/v2/index.json :: linear probe ee_vs_mm (fits_v4_sib): AUC, FLAV_F1R :: sha256 c6c4138ee5c86423
\newcommand{\MlstLpEemmAucInit}{0.9995\,\ensuremath{\pm}\,0.0003}  % experiments/FIGS/data/v2/index.json :: linear probe ee_vs_mm (fits_v4_sib): AUC, INIT :: sha256 c6c4138ee5c86423
\newcommand{\MlstLpEemmAucInitBn}{1.0000\,\ensuremath{\pm}\,\ensuremath{<}0.0001}  % experiments/FIGS/data/v2/index.json :: linear probe ee_vs_mm (fits_v4_sib): AUC, INIT_BN :: sha256 c6c4138ee5c86423
\newcommand{\MlstLpEemmAucRandMerged}{1.0000\,\ensuremath{\pm}\,\ensuremath{<}0.0001}  % experiments/FIGS/data/v2/index.json :: linear probe ee_vs_mm (fits_v4_sib): AUC, the random partitions that merge the pair :: sha256 c6c4138ee5c86423
\newcommand{\MlstPairEemmNTrain}{9{,}571}  % experiments/FIGS/data/v2/index.json :: fits_v4_sib ee_vs_mm: n_train :: sha256 c6c4138ee5c86423
\newcommand{\MlstPairEemmNVal}{1{,}068}  % experiments/FIGS/data/v2/index.json :: fits_v4_sib ee_vs_mm: n_val :: sha256 c6c4138ee5c86423
\newcommand{\MlstPairEemmNTestPerSide}{128{,}833--135{,}245}  % experiments/FIGS/data/v2/index.json :: fits_v4_sib ee_vs_mm: n_test_signal and n_test_background :: sha256 c6c4138ee5c86423
\newcommand{\MlstLpJcTwoAucFull}{0.9801\,\ensuremath{\pm}\,\ensuremath{<}0.0001}  % experiments/FIGS/data/v2/index.json :: linear probe JetClass-II N=1000000: AUC averaged over the classes, L188 :: sha256 c6c4138ee5c86423
\newcommand{\MlstLpJcTwoAccFull}{0.4348\,\ensuremath{\pm}\,0.0003}  % experiments/FIGS/data/v2/index.json :: linear probe JetClass-II N=1000000: accuracy, L188 :: sha256 c6c4138ee5c86423
\newcommand{\MlstLpJcTwoAucRes}{0.9796\,\ensuremath{\pm}\,0.0001}  % experiments/FIGS/data/v2/index.json :: linear probe JetClass-II N=1000000: AUC averaged over the classes, L162 :: sha256 c6c4138ee5c86423
\newcommand{\MlstLpJcTwoAccRes}{0.4328\,\ensuremath{\pm}\,0.0008}  % experiments/FIGS/data/v2/index.json :: linear probe JetClass-II N=1000000: accuracy, L162 :: sha256 c6c4138ee5c86423
\newcommand{\MlstLpJcTwoAucMid}{0.9756\,\ensuremath{\pm}\,0.0002}  % experiments/FIGS/data/v2/index.json :: linear probe JetClass-II N=1000000: AUC averaged over the classes, R42_Q1 :: sha256 c6c4138ee5c86423
\newcommand{\MlstLpJcTwoAccMid}{0.361\,\ensuremath{\pm}\,0.005}  % experiments/FIGS/data/v2/index.json :: linear probe JetClass-II N=1000000: accuracy, R42_Q1 :: sha256 c6c4138ee5c86423
\newcommand{\MlstLpJcTwoAucCoarse}{0.9691\,\ensuremath{\pm}\,0.0004}  % experiments/FIGS/data/v2/index.json :: linear probe JetClass-II N=1000000: AUC averaged over the classes, R16_Q1 :: sha256 c6c4138ee5c86423
\newcommand{\MlstLpJcTwoAccCoarse}{0.334\,\ensuremath{\pm}\,0.007}  % experiments/FIGS/data/v2/index.json :: linear probe JetClass-II N=1000000: accuracy, R16_Q1 :: sha256 c6c4138ee5c86423
\newcommand{\MlstLpJcTwoAucResMass}{0.9791 and 0.9794}  % experiments/FIGS/data/v2/index.json :: linear probe JetClass-II N=1000000: AUC averaged over the classes, L162_MASS :: sha256 c6c4138ee5c86423
\newcommand{\MlstLpJcTwoAccResMass}{0.4271 and 0.4305}  % experiments/FIGS/data/v2/index.json :: linear probe JetClass-II N=1000000: accuracy, L162_MASS :: sha256 c6c4138ee5c86423
\newcommand{\MlstLpJcTwoAucCoarseMass}{0.9679 and 0.9696}  % experiments/FIGS/data/v2/index.json :: linear probe JetClass-II N=1000000: AUC averaged over the classes, R16_Q1_MASS :: sha256 c6c4138ee5c86423
\newcommand{\MlstLpJcTwoAccCoarseMass}{0.312 and 0.350}  % experiments/FIGS/data/v2/index.json :: linear probe JetClass-II N=1000000: accuracy, R16_Q1_MASS :: sha256 c6c4138ee5c86423
\newcommand{\MlstLpJcTwoAucFzero}{0.9717 and 0.9710}  % experiments/FIGS/data/v2/index.json :: linear probe JetClass-II N=1000000: AUC averaged over the classes, FLAV_F0 :: sha256 c6c4138ee5c86423
\newcommand{\MlstLpJcTwoAccFzero}{0.3769 and 0.3723}  % experiments/FIGS/data/v2/index.json :: linear probe JetClass-II N=1000000: accuracy, FLAV_F0 :: sha256 c6c4138ee5c86423
\newcommand{\MlstLpJcTwoAucFone}{0.9758 and 0.9759}  % experiments/FIGS/data/v2/index.json :: linear probe JetClass-II N=1000000: AUC averaged over the classes, FLAV_F1 :: sha256 c6c4138ee5c86423
\newcommand{\MlstLpJcTwoAccFone}{0.3945 and 0.3963}  % experiments/FIGS/data/v2/index.json :: linear probe JetClass-II N=1000000: accuracy, FLAV_F1 :: sha256 c6c4138ee5c86423
\newcommand{\MlstLpJcTwoAucFoner}{0.9760 and 0.9760}  % experiments/FIGS/data/v2/index.json :: linear probe JetClass-II N=1000000: AUC averaged over the classes, FLAV_F1R :: sha256 c6c4138ee5c86423
\newcommand{\MlstLpJcTwoAccFoner}{0.4071 and 0.4089}  % experiments/FIGS/data/v2/index.json :: linear probe JetClass-II N=1000000: accuracy, FLAV_F1R :: sha256 c6c4138ee5c86423
\newcommand{\MlstLpJcTwoAucInit}{0.8121\,\ensuremath{\pm}\,0.0019}  % experiments/FIGS/data/v2/index.json :: linear probe JetClass-II N=1000000: AUC averaged over the classes, INIT :: sha256 c6c4138ee5c86423
\newcommand{\MlstLpJcTwoAccInit}{0.0724\,\ensuremath{\pm}\,0.0015}  % experiments/FIGS/data/v2/index.json :: linear probe JetClass-II N=1000000: accuracy, INIT :: sha256 c6c4138ee5c86423
\newcommand{\MlstLpJcTwoAucInitBn}{0.866\,\ensuremath{\pm}\,0.004}  % experiments/FIGS/data/v2/index.json :: linear probe JetClass-II N=1000000: AUC averaged over the classes, INIT_BN :: sha256 c6c4138ee5c86423
\newcommand{\MlstLpJcTwoAccInitBn}{0.1017\,\ensuremath{\pm}\,0.0017}  % experiments/FIGS/data/v2/index.json :: linear probe JetClass-II N=1000000: accuracy, INIT_BN :: sha256 c6c4138ee5c86423
\newcommand{\MlstLpJcTwoAucRand}{0.9777\,\ensuremath{\pm}\,0.0002}  % experiments/FIGS/data/v2/index.json :: linear probe JetClass-II N=1000000: AUC averaged over the classes, the random partitions :: sha256 c6c4138ee5c86423
\newcommand{\MlstLpJcOneAccInitBn}{0.475\,\ensuremath{\pm}\,0.005}  % experiments/FIGS/data/v2/index.json :: linear probe JetClass N=1000000: accuracy, INIT_BN :: sha256 c6c4138ee5c86423
\newcommand{\MlstBnMaxDev}{0.2\%}  % experiments/FIGS/data/v2/index.json :: largest |r(best70_bn) / r(best70) - 1| of a JetClass-II ratio of 1 - AUC to the 188-class run, per cent :: sha256 c6c4138ee5c86423
\newcommand{\MlstBnMaxDevPair}{0.7\%}  % experiments/FIGS/data/v2/index.json :: largest |r(best70_bn) / r(best70) - 1| of a ratio of 1 - AUC on a class pair to the 188-class run, levels of the hierarchy, mass models and controlled partitions, per cent :: sha256 c6c4138ee5c86423 :: pairs with a check: bc_vs_bq_cs, bb_vs_cc, bbqq_vs_ccqq, qcd_bb_vs_cc, qcd_b_vs_light, qcd_c_vs_light; without: ee_vs_mm
\newcommand{\MlstGpuTypeBound}{1.4\%}  % experiments/FIGS/data/v2/index.json :: largest |label-set ratio of 1 - AUC of index 3 / mean of the other indices - 1|, JetClass-II largest set, per cent :: sha256 c6c4138ee5c86423 :: mtx-l188-s3 trained on NVIDIA GeForce RTX 3090; largest for R16_Q1; label sets L162 0.09%, R42_Q1 0.70%, R16_Q1 1.37%
\newcommand{\MlstGpuOtherSeedsDiff}{1.9\%}  % experiments/FIGS/data/v2/index.json :: largest (max - min) / mean of a label-set ratio of 1 - AUC over the indices other than 3, JetClass-II largest set, per cent :: sha256 c6c4138ee5c86423 :: label sets L162 0.90%, R42_Q1 1.14%, R16_Q1 1.93%
\newcommand{\MlstPerClassResNDoubled}{0}  % experiments/FIGS/data/v2/index.json :: JetClass-II classes whose mean ratio of 1 - AUC (L162 over 188) is at least 2 (count) :: sha256 c6c4138ee5c86423
\newcommand{\MlstPerClassResSeedSdMedian}{0.011}  % experiments/FIGS/data/v2/index.json :: median over classes of the SD over runs of the per-class ratio, L162 over 188 :: sha256 c6c4138ee5c86423
\newcommand{\MlstPerClassResTopClass}{QCD, \ensuremath{ccss}}  % experiments/FIGS/data/v2/index.json :: class with the largest mean ratio, L162 :: sha256 c6c4138ee5c86423
\newcommand{\MlstPerClassResTopRatio}{1.129\,\ensuremath{\pm}\,0.010}  % experiments/FIGS/data/v2/index.json :: per-class ratio of 1 - AUC of that class, L162 over 188, run by run :: sha256 c6c4138ee5c86423
\newcommand{\MlstPerClassResTopRatioPr}{13\% more errors}  % experiments/FIGS/data/v2/index.json :: per-class ratio of 1 - AUC of that class, L162 over 188, run by run :: sha256 c6c4138ee5c86423 :: prose twin: two significant figures, direction in words
\newcommand{\MlstPerClassResLeptonN}{0}  % experiments/FIGS/data/v2/index.json :: resonant classes sharing their L162 group with a decay that differs only in lepton flavor (count) :: sha256 c6c4138ee5c86423
\newcommand{\MlstPerClassResLeptonMedian}{--}  % experiments/FIGS/data/v2/index.json :: median over those classes of the mean per-class ratio, L162 over 188 :: sha256 c6c4138ee5c86423
\newcommand{\MlstPerClassResLeptonMedianPr}{\pending{no run yet}}  % experiments/FIGS/data/v2/index.json :: median over those classes of the mean per-class ratio, L162 over 188 :: sha256 c6c4138ee5c86423 :: prose twin: two significant figures, direction in words
\newcommand{\MlstPerClassResLeptonShare}{--}  % experiments/FIGS/data/v2/index.json :: their share of the summed ln ratio over the resonant classes, L162 :: sha256 c6c4138ee5c86423
\newcommand{\MlstPerClassResQuarkN}{0}  % experiments/FIGS/data/v2/index.json :: resonant classes sharing their L162 group with a decay that differs only in quark flavor (count) :: sha256 c6c4138ee5c86423
\newcommand{\MlstPerClassResQuarkMedian}{--}  % experiments/FIGS/data/v2/index.json :: median over those classes of the mean per-class ratio, L162 over 188 :: sha256 c6c4138ee5c86423
\newcommand{\MlstPerClassResQuarkMedianPr}{\pending{no run yet}}  % experiments/FIGS/data/v2/index.json :: median over those classes of the mean per-class ratio, L162 over 188 :: sha256 c6c4138ee5c86423 :: prose twin: two significant figures, direction in words
\newcommand{\MlstPerClassResQuarkShare}{--}  % experiments/FIGS/data/v2/index.json :: their share of the summed ln ratio over the resonant classes, L162 :: sha256 c6c4138ee5c86423
\newcommand{\MlstPerClassResQuarkLeptonN}{0}  % experiments/FIGS/data/v2/index.json :: resonant classes counted in both the quark- and the lepton-flavor set, L162 (count) :: sha256 c6c4138ee5c86423
\newcommand{\MlstPerClassMidNDoubled}{40}  % experiments/FIGS/data/v2/index.json :: JetClass-II classes whose mean ratio of 1 - AUC (R42_Q1 over 188) is at least 2 (count) :: sha256 c6c4138ee5c86423
\newcommand{\MlstPerClassMidSeedSdMedian}{0.023}  % experiments/FIGS/data/v2/index.json :: median over classes of the SD over runs of the per-class ratio, R42_Q1 over 188 :: sha256 c6c4138ee5c86423
\newcommand{\MlstPerClassMidTopClass}{\ensuremath{X\to YY\to bb\mu\mu}}  % experiments/FIGS/data/v2/index.json :: class with the largest mean ratio, R42_Q1 :: sha256 c6c4138ee5c86423
\newcommand{\MlstPerClassMidTopRatio}{4.06\,\ensuremath{\pm}\,0.14}  % experiments/FIGS/data/v2/index.json :: per-class ratio of 1 - AUC of that class, R42_Q1 over 188, run by run :: sha256 c6c4138ee5c86423
\newcommand{\MlstPerClassMidTopRatioPr}{4.1 times the errors}  % experiments/FIGS/data/v2/index.json :: per-class ratio of 1 - AUC of that class, R42_Q1 over 188, run by run :: sha256 c6c4138ee5c86423 :: prose twin: two significant figures, direction in words
\newcommand{\MlstPerClassMidLeptonN}{78}  % experiments/FIGS/data/v2/index.json :: resonant classes sharing their R42_Q1 group with a decay that differs only in lepton flavor (count) :: sha256 c6c4138ee5c86423
\newcommand{\MlstPerClassMidLeptonMedian}{2.02}  % experiments/FIGS/data/v2/index.json :: median over those classes of the mean per-class ratio, R42_Q1 over 188 :: sha256 c6c4138ee5c86423
\newcommand{\MlstPerClassMidLeptonMedianPr}{2.0 times the errors}  % experiments/FIGS/data/v2/index.json :: median over those classes of the mean per-class ratio, R42_Q1 over 188 :: sha256 c6c4138ee5c86423 :: prose twin: two significant figures, direction in words
\newcommand{\MlstPerClassMidLeptonShare}{90\%}  % experiments/FIGS/data/v2/index.json :: their share of the summed ln ratio over the resonant classes, R42_Q1 :: sha256 c6c4138ee5c86423
\newcommand{\MlstPerClassMidQuarkN}{111}  % experiments/FIGS/data/v2/index.json :: resonant classes sharing their R42_Q1 group with a decay that differs only in quark flavor (count) :: sha256 c6c4138ee5c86423
\newcommand{\MlstPerClassMidQuarkMedian}{1.14}  % experiments/FIGS/data/v2/index.json :: median over those classes of the mean per-class ratio, R42_Q1 over 188 :: sha256 c6c4138ee5c86423
\newcommand{\MlstPerClassMidQuarkMedianPr}{14\% more errors}  % experiments/FIGS/data/v2/index.json :: median over those classes of the mean per-class ratio, R42_Q1 over 188 :: sha256 c6c4138ee5c86423 :: prose twin: two significant figures, direction in words
\newcommand{\MlstPerClassMidQuarkShare}{53\%}  % experiments/FIGS/data/v2/index.json :: their share of the summed ln ratio over the resonant classes, R42_Q1 :: sha256 c6c4138ee5c86423
\newcommand{\MlstPerClassMidQuarkLeptonN}{44}  % experiments/FIGS/data/v2/index.json :: resonant classes counted in both the quark- and the lepton-flavor set, R42_Q1 (count) :: sha256 c6c4138ee5c86423
\newcommand{\MlstPerClassCoarseNDoubled}{60}  % experiments/FIGS/data/v2/index.json :: JetClass-II classes whose mean ratio of 1 - AUC (R16_Q1 over 188) is at least 2 (count) :: sha256 c6c4138ee5c86423
\newcommand{\MlstPerClassCoarseSeedSdMedian}{0.055}  % experiments/FIGS/data/v2/index.json :: median over classes of the SD over runs of the per-class ratio, R16_Q1 over 188 :: sha256 c6c4138ee5c86423
\newcommand{\MlstPerClassCoarseTopClass}{\ensuremath{X\to YY\to bbee}}  % experiments/FIGS/data/v2/index.json :: class with the largest mean ratio, R16_Q1 :: sha256 c6c4138ee5c86423
\newcommand{\MlstPerClassCoarseTopRatio}{5.2\,\ensuremath{\pm}\,0.4}  % experiments/FIGS/data/v2/index.json :: per-class ratio of 1 - AUC of that class, R16_Q1 over 188, run by run :: sha256 c6c4138ee5c86423
\newcommand{\MlstPerClassCoarseTopRatioPr}{5.2 times the errors}  % experiments/FIGS/data/v2/index.json :: per-class ratio of 1 - AUC of that class, R16_Q1 over 188, run by run :: sha256 c6c4138ee5c86423 :: prose twin: two significant figures, direction in words
\newcommand{\MlstPerClassCoarseLeptonN}{78}  % experiments/FIGS/data/v2/index.json :: resonant classes sharing their R16_Q1 group with a decay that differs only in lepton flavor (count) :: sha256 c6c4138ee5c86423
\newcommand{\MlstPerClassCoarseLeptonMedian}{2.27}  % experiments/FIGS/data/v2/index.json :: median over those classes of the mean per-class ratio, R16_Q1 over 188 :: sha256 c6c4138ee5c86423
\newcommand{\MlstPerClassCoarseLeptonMedianPr}{2.3 times the errors}  % experiments/FIGS/data/v2/index.json :: median over those classes of the mean per-class ratio, R16_Q1 over 188 :: sha256 c6c4138ee5c86423 :: prose twin: two significant figures, direction in words
\newcommand{\MlstPerClassCoarseLeptonShare}{68\%}  % experiments/FIGS/data/v2/index.json :: their share of the summed ln ratio over the resonant classes, R16_Q1 :: sha256 c6c4138ee5c86423
\newcommand{\MlstPerClassCoarseQuarkN}{141}  % experiments/FIGS/data/v2/index.json :: resonant classes sharing their R16_Q1 group with a decay that differs only in quark flavor (count) :: sha256 c6c4138ee5c86423
\newcommand{\MlstPerClassCoarseQuarkMedian}{1.65}  % experiments/FIGS/data/v2/index.json :: median over those classes of the mean per-class ratio, R16_Q1 over 188 :: sha256 c6c4138ee5c86423
\newcommand{\MlstPerClassCoarseQuarkMedianPr}{65\% more errors}  % experiments/FIGS/data/v2/index.json :: median over those classes of the mean per-class ratio, R16_Q1 over 188 :: sha256 c6c4138ee5c86423 :: prose twin: two significant figures, direction in words
\newcommand{\MlstPerClassCoarseQuarkShare}{87\%}  % experiments/FIGS/data/v2/index.json :: their share of the summed ln ratio over the resonant classes, R16_Q1 :: sha256 c6c4138ee5c86423
\newcommand{\MlstPerClassCoarseQuarkLeptonN}{64}  % experiments/FIGS/data/v2/index.json :: resonant classes counted in both the quark- and the lepton-flavor set, R16_Q1 (count) :: sha256 c6c4138ee5c86423
\newcommand{\MlstPerClassRandFlavMergedOverApart}{1.010\,\ensuremath{\pm}\,0.008}  % experiments/FIGS/data/v2/index.json :: random partitions, linear probe: median per-class ratio of the classes merged with a decay that differs only in quark flavor over that of the classes kept apart from all of them, mean +- SD over the partitions :: sha256 c6c4138ee5c86423
\newcommand{\MlstPerClassRandFlavMergedOverApartPr}{equal within 1\%}  % experiments/FIGS/data/v2/index.json :: random partitions, linear probe: median per-class ratio of the classes merged with a decay that differs only in quark flavor over that of the classes kept apart from all of them, mean +- SD over the partitions :: sha256 c6c4138ee5c86423 :: prose twin: two significant figures, direction in words
\newcommand{\MlstPerClassRandLeptonMergedOverApart}{0.997\,\ensuremath{\pm}\,0.020}  % experiments/FIGS/data/v2/index.json :: random partitions, linear probe: median per-class ratio of the classes merged with a decay that differs only in lepton flavor over that of the classes kept apart from all of them, mean +- SD over the partitions :: sha256 c6c4138ee5c86423
\newcommand{\MlstPerClassRandLeptonMergedOverApartPr}{equal within 3\%}  % experiments/FIGS/data/v2/index.json :: random partitions, linear probe: median per-class ratio of the classes merged with a decay that differs only in lepton flavor over that of the classes kept apart from all of them, mean +- SD over the partitions :: sha256 c6c4138ee5c86423 :: prose twin: two significant figures, direction in words
\newcommand{\MlstMlpNatPerClassRandFlavMergedOverApart}{1.003\,\ensuremath{\pm}\,0.007}  % experiments/FIGS/data/v2/index.json :: random partitions, nonlinear probe, natural composition: median per-class ratio of the classes merged with a decay that differs only in quark flavor over that of the classes kept apart from all of them, mean +- SD over the partitions :: sha256 c6c4138ee5c86423
\newcommand{\MlstMlpNatPerClassRandFlavMergedOverApartPr}{equal within 1\%}  % experiments/FIGS/data/v2/index.json :: random partitions, nonlinear probe, natural composition: median per-class ratio of the classes merged with a decay that differs only in quark flavor over that of the classes kept apart from all of them, mean +- SD over the partitions :: sha256 c6c4138ee5c86423 :: prose twin: two significant figures, direction in words
\newcommand{\MlstMlpNatPerClassRandLeptonMergedOverApart}{1.01\,\ensuremath{\pm}\,0.07}  % experiments/FIGS/data/v2/index.json :: random partitions, nonlinear probe, natural composition: median per-class ratio of the classes merged with a decay that differs only in lepton flavor over that of the classes kept apart from all of them, mean +- SD over the partitions :: sha256 c6c4138ee5c86423
\newcommand{\MlstMlpNatPerClassRandLeptonMergedOverApartPr}{equal within 8\%}  % experiments/FIGS/data/v2/index.json :: random partitions, nonlinear probe, natural composition: median per-class ratio of the classes merged with a decay that differs only in lepton flavor over that of the classes kept apart from all of them, mean +- SD over the partitions :: sha256 c6c4138ee5c86423 :: prose twin: two significant figures, direction in words
\newcommand{\MlstPerClassAucDiffResSeedSdMedian}{0.015}  % experiments/FIGS/data/v2/index.json :: median over classes of the SD over indices of 100 (AUC_c of L162 - AUC_c of 188), pp :: sha256 c6c4138ee5c86423
\newcommand{\MlstPerClassAucDiffResMin}{\ensuremath{-}0.73}  % experiments/FIGS/data/v2/index.json :: most negative mean over indices of 100 (AUC_c of L162 - AUC_c of 188), pp :: sha256 c6c4138ee5c86423 :: label_QCD_bcss
\newcommand{\MlstPerClassAucDiffResMinClass}{QCD, \ensuremath{bcss}}  % experiments/FIGS/data/v2/index.json :: its class, L162 :: sha256 c6c4138ee5c86423
\newcommand{\MlstPerClassAucDiffMidSeedSdMedian}{0.032}  % experiments/FIGS/data/v2/index.json :: median over classes of the SD over indices of 100 (AUC_c of R42_Q1 - AUC_c of 188), pp :: sha256 c6c4138ee5c86423
\newcommand{\MlstPerClassAucDiffMidMin}{\ensuremath{-}1.4}  % experiments/FIGS/data/v2/index.json :: most negative mean over indices of 100 (AUC_c of R42_Q1 - AUC_c of 188), pp :: sha256 c6c4138ee5c86423 :: label_X_YY_cqtauev
\newcommand{\MlstPerClassAucDiffMidMinClass}{\ensuremath{X\to YY\to cq\tau_e\nu}}  % experiments/FIGS/data/v2/index.json :: its class, R42_Q1 :: sha256 c6c4138ee5c86423
\newcommand{\MlstPerClassAucDiffCoarseSeedSdMedian}{0.077}  % experiments/FIGS/data/v2/index.json :: median over classes of the SD over indices of 100 (AUC_c of R16_Q1 - AUC_c of 188), pp :: sha256 c6c4138ee5c86423
\newcommand{\MlstPerClassAucDiffCoarseMin}{\ensuremath{-}3.1}  % experiments/FIGS/data/v2/index.json :: most negative mean over indices of 100 (AUC_c of R16_Q1 - AUC_c of 188), pp :: sha256 c6c4138ee5c86423 :: label_QCD_bcss
\newcommand{\MlstPerClassAucDiffCoarseMinClass}{QCD, \ensuremath{bcss}}  % experiments/FIGS/data/v2/index.json :: its class, R16_Q1 :: sha256 c6c4138ee5c86423
\newcommand{\MlstRocJcTwoHurtList}{\ensuremath{X\to YY\to ccq}; \ensuremath{X\to YY\to ssb}; \ensuremath{X\to YY\to ccss}; QCD, \ensuremath{bcc}; QCD, \ensuremath{bccss}; QCD, \ensuremath{bss}; QCD, \ensuremath{ccs} and QCD, \ensuremath{ccss}}  % experiments/FIGS/data/v2/index.json :: jc2_roc_classes on the class-balanced probes; candidates with at least 1000 test jets :: sha256 c6c4138ee5c86423 :: label_X_YY_ccq, label_X_YY_ssb, label_X_YY_ccss, label_QCD_bcc, label_QCD_bccss, label_QCD_bss, label_QCD_ccs, label_QCD_ccss
\newcommand{\MlstRocJcTwoStableList}{\ensuremath{X\to gg}; \ensuremath{X\to YY\to ggg}; \ensuremath{X\to YY\to g\tau_h\tau_h} and QCD, light}  % experiments/FIGS/data/v2/index.json :: jc2_roc_classes on the class-balanced probes; candidates with at least 1000 test jets :: sha256 c6c4138ee5c86423 :: label_X_gg, label_X_YY_ggg, label_X_YY_gtauhtauh, label_QCD_light
\newcommand{\MlstRocJcTwoHurtMinDrop}{2.1}  % experiments/FIGS/data/v2/index.json :: jc2_roc_classes on the class-balanced probes; candidates with at least 1000 test jets :: sha256 c6c4138ee5c86423 :: smallest AUC loss at 17 classes of the hurt, pp
\newcommand{\MlstRocJcTwoHurtMaxDrop}{2.8}  % experiments/FIGS/data/v2/index.json :: jc2_roc_classes on the class-balanced probes; candidates with at least 1000 test jets :: sha256 c6c4138ee5c86423 :: largest AUC loss at 17 classes of the hurt, pp
\newcommand{\MlstRocJcTwoStableMaxChange}{15}  % experiments/FIGS/data/v2/index.json :: jc2_roc_classes on the class-balanced probes; candidates with at least 1000 test jets :: sha256 c6c4138ee5c86423 :: largest change of 1 - AUC at any label set of the stable classes, per cent of the 188-class value
\newcommand{\MlstJcBalNThin}{4}  % experiments/FIGS/data/v2/jc2bal/census.json :: counts.ft_test :: sha256 a4379ee0817254cf :: classes with fewer than 1000 class-balanced test jets (F3 hatching)
\newcommand{\MlstMinBkgPass}{20}  % experiments/FIGS/mlst_v4.py :: JC1_MIN_BKG :: sha256 fa536b4866ffd651 :: fewest background jets behind a drawn ROC point or a quoted rejection
\newcommand{\MlstMassMinJets}{100}  % experiments/EVAL/atlas_resolution.py :: MIN_JETS :: sha256 b82cc219f79f0f14 :: jets below which a mass bin is null
\newcommand{\MlstRocCutBbFull}{0.79}  % experiments/FIGS/data/v2/index.json :: curves_v2/pairs curve bb_vs_cc: largest cut.eps_s over the L188 runs ('not reached' when a run's cut is null or leaves fewer than 2 grid points) :: sha256 c6c4138ee5c86423
\newcommand{\MlstRocCutBbRes}{0.77}  % experiments/FIGS/data/v2/index.json :: curves_v2/pairs curve bb_vs_cc: largest cut.eps_s over the L162 runs ('not reached' when a run's cut is null or leaves fewer than 2 grid points) :: sha256 c6c4138ee5c86423
\newcommand{\MlstRocCutBbMid}{0.79}  % experiments/FIGS/data/v2/index.json :: curves_v2/pairs curve bb_vs_cc: largest cut.eps_s over the R42_Q1 runs ('not reached' when a run's cut is null or leaves fewer than 2 grid points) :: sha256 c6c4138ee5c86423
\newcommand{\MlstRocCutBbCoarse}{0.13}  % experiments/FIGS/data/v2/index.json :: curves_v2/pairs curve bb_vs_cc: largest cut.eps_s over the R16_Q1 runs ('not reached' when a run's cut is null or leaves fewer than 2 grid points) :: sha256 c6c4138ee5c86423
\newcommand{\MlstRocCutBbFzero}{0.03}  % experiments/FIGS/data/v2/index.json :: curves_v2/pairs curve bb_vs_cc: largest cut.eps_s over the FLAV_F0 runs ('not reached' when a run's cut is null or leaves fewer than 2 grid points) :: sha256 c6c4138ee5c86423
\newcommand{\MlstRocCutBbFone}{0.62}  % experiments/FIGS/data/v2/index.json :: curves_v2/pairs curve bb_vs_cc: largest cut.eps_s over the FLAV_F1 runs ('not reached' when a run's cut is null or leaves fewer than 2 grid points) :: sha256 c6c4138ee5c86423
\newcommand{\MlstRocCutBbFoner}{0.51}  % experiments/FIGS/data/v2/index.json :: curves_v2/pairs curve bb_vs_cc: largest cut.eps_s over the FLAV_F1R runs ('not reached' when a run's cut is null or leaves fewer than 2 grid points) :: sha256 c6c4138ee5c86423
\newcommand{\MlstRocNBkgBb}{200{,}000}  % experiments/FIGS/data/v2/index.json :: curves_v2/pairs curve bb_vs_cc: n_background :: sha256 c6c4138ee5c86423
\newcommand{\MlstRocRejMaxBb}{10{,}000}  % experiments/FIGS/data/v2/index.json :: curves_v2/pairs curve bb_vs_cc: n_background / 20 :: sha256 c6c4138ee5c86423
\newcommand{\MlstRocCutBbqqFull}{0.67}  % experiments/FIGS/data/v2/index.json :: curves_v2/pairs curve bbqq_vs_ccqq: largest cut.eps_s over the L188 runs ('not reached' when a run's cut is null or leaves fewer than 2 grid points) :: sha256 c6c4138ee5c86423
\newcommand{\MlstRocCutBbqqRes}{0.66}  % experiments/FIGS/data/v2/index.json :: curves_v2/pairs curve bbqq_vs_ccqq: largest cut.eps_s over the L162 runs ('not reached' when a run's cut is null or leaves fewer than 2 grid points) :: sha256 c6c4138ee5c86423
\newcommand{\MlstRocCutBbqqMid}{0.62}  % experiments/FIGS/data/v2/index.json :: curves_v2/pairs curve bbqq_vs_ccqq: largest cut.eps_s over the R42_Q1 runs ('not reached' when a run's cut is null or leaves fewer than 2 grid points) :: sha256 c6c4138ee5c86423
\newcommand{\MlstRocCutBbqqCoarse}{0.06}  % experiments/FIGS/data/v2/index.json :: curves_v2/pairs curve bbqq_vs_ccqq: largest cut.eps_s over the R16_Q1 runs ('not reached' when a run's cut is null or leaves fewer than 2 grid points) :: sha256 c6c4138ee5c86423
\newcommand{\MlstRocCutBbqqFzero}{0.06}  % experiments/FIGS/data/v2/index.json :: curves_v2/pairs curve bbqq_vs_ccqq: largest cut.eps_s over the FLAV_F0 runs ('not reached' when a run's cut is null or leaves fewer than 2 grid points) :: sha256 c6c4138ee5c86423
\newcommand{\MlstRocCutBbqqFone}{0.56}  % experiments/FIGS/data/v2/index.json :: curves_v2/pairs curve bbqq_vs_ccqq: largest cut.eps_s over the FLAV_F1 runs ('not reached' when a run's cut is null or leaves fewer than 2 grid points) :: sha256 c6c4138ee5c86423
\newcommand{\MlstRocCutBbqqFoner}{0.32}  % experiments/FIGS/data/v2/index.json :: curves_v2/pairs curve bbqq_vs_ccqq: largest cut.eps_s over the FLAV_F1R runs ('not reached' when a run's cut is null or leaves fewer than 2 grid points) :: sha256 c6c4138ee5c86423
\newcommand{\MlstRocNBkgBbqq}{200{,}000}  % experiments/FIGS/data/v2/index.json :: curves_v2/pairs curve bbqq_vs_ccqq: n_background :: sha256 c6c4138ee5c86423
\newcommand{\MlstRocRejMaxBbqq}{10{,}000}  % experiments/FIGS/data/v2/index.json :: curves_v2/pairs curve bbqq_vs_ccqq: n_background / 20 :: sha256 c6c4138ee5c86423
\newcommand{\MlstRocCutBcFull}{0.24}  % experiments/FIGS/data/v2/index.json :: curves_v2/pairs curve bc_vs_bq_cs: largest cut.eps_s over the L188 runs ('not reached' when a run's cut is null or leaves fewer than 2 grid points) :: sha256 c6c4138ee5c86423
\newcommand{\MlstRocCutBcRes}{0.25}  % experiments/FIGS/data/v2/index.json :: curves_v2/pairs curve bc_vs_bq_cs: largest cut.eps_s over the L162 runs ('not reached' when a run's cut is null or leaves fewer than 2 grid points) :: sha256 c6c4138ee5c86423
\newcommand{\MlstRocCutBcMid}{0.14}  % experiments/FIGS/data/v2/index.json :: curves_v2/pairs curve bc_vs_bq_cs: largest cut.eps_s over the R42_Q1 runs ('not reached' when a run's cut is null or leaves fewer than 2 grid points) :: sha256 c6c4138ee5c86423
\newcommand{\MlstRocCutBcCoarse}{0.00}  % experiments/FIGS/data/v2/index.json :: curves_v2/pairs curve bc_vs_bq_cs: largest cut.eps_s over the R16_Q1 runs ('not reached' when a run's cut is null or leaves fewer than 2 grid points) :: sha256 c6c4138ee5c86423
\newcommand{\MlstRocCutBcFzero}{0.00}  % experiments/FIGS/data/v2/index.json :: curves_v2/pairs curve bc_vs_bq_cs: largest cut.eps_s over the FLAV_F0 runs ('not reached' when a run's cut is null or leaves fewer than 2 grid points) :: sha256 c6c4138ee5c86423
\newcommand{\MlstRocCutBcFone}{0.01}  % experiments/FIGS/data/v2/index.json :: curves_v2/pairs curve bc_vs_bq_cs: largest cut.eps_s over the FLAV_F1 runs ('not reached' when a run's cut is null or leaves fewer than 2 grid points) :: sha256 c6c4138ee5c86423
\newcommand{\MlstRocCutBcFoner}{0.01}  % experiments/FIGS/data/v2/index.json :: curves_v2/pairs curve bc_vs_bq_cs: largest cut.eps_s over the FLAV_F1R runs ('not reached' when a run's cut is null or leaves fewer than 2 grid points) :: sha256 c6c4138ee5c86423
\newcommand{\MlstRocNBkgBc}{200{,}000}  % experiments/FIGS/data/v2/index.json :: curves_v2/pairs curve bc_vs_bq_cs: n_background :: sha256 c6c4138ee5c86423
\newcommand{\MlstRocRejMaxBc}{10{,}000}  % experiments/FIGS/data/v2/index.json :: curves_v2/pairs curve bc_vs_bq_cs: n_background / 20 :: sha256 c6c4138ee5c86423
\newcommand{\MlstRocCutEemmFull}{not reached}  % experiments/FIGS/data/v2/index.json :: curves_v2/pairs curve ee_vs_mm: largest cut.eps_s over the L188 runs ('not reached' when a run's cut is null or leaves fewer than 2 grid points) :: sha256 c6c4138ee5c86423
\newcommand{\MlstRocCutEemmRes}{not reached}  % experiments/FIGS/data/v2/index.json :: curves_v2/pairs curve ee_vs_mm: largest cut.eps_s over the L162 runs ('not reached' when a run's cut is null or leaves fewer than 2 grid points) :: sha256 c6c4138ee5c86423
\newcommand{\MlstRocCutEemmMid}{0.00}  % experiments/FIGS/data/v2/index.json :: curves_v2/pairs curve ee_vs_mm: largest cut.eps_s over the R42_Q1 runs ('not reached' when a run's cut is null or leaves fewer than 2 grid points) :: sha256 c6c4138ee5c86423
\newcommand{\MlstRocCutEemmCoarse}{0.00}  % experiments/FIGS/data/v2/index.json :: curves_v2/pairs curve ee_vs_mm: largest cut.eps_s over the R16_Q1 runs ('not reached' when a run's cut is null or leaves fewer than 2 grid points) :: sha256 c6c4138ee5c86423
\newcommand{\MlstRocCutEemmFzero}{not reached}  % experiments/FIGS/data/v2/index.json :: curves_v2/pairs curve ee_vs_mm: largest cut.eps_s over the FLAV_F0 runs ('not reached' when a run's cut is null or leaves fewer than 2 grid points) :: sha256 c6c4138ee5c86423
\newcommand{\MlstRocCutEemmFone}{not reached}  % experiments/FIGS/data/v2/index.json :: curves_v2/pairs curve ee_vs_mm: largest cut.eps_s over the FLAV_F1 runs ('not reached' when a run's cut is null or leaves fewer than 2 grid points) :: sha256 c6c4138ee5c86423
\newcommand{\MlstRocCutEemmFoner}{not reached}  % experiments/FIGS/data/v2/index.json :: curves_v2/pairs curve ee_vs_mm: largest cut.eps_s over the FLAV_F1R runs ('not reached' when a run's cut is null or leaves fewer than 2 grid points) :: sha256 c6c4138ee5c86423
\newcommand{\MlstRocNBkgEemm}{128{,}833}  % experiments/FIGS/data/v2/index.json :: curves_v2/pairs curve ee_vs_mm: n_background :: sha256 c6c4138ee5c86423
\newcommand{\MlstRocRejMaxEemm}{6{,}442}  % experiments/FIGS/data/v2/index.json :: curves_v2/pairs curve ee_vs_mm: n_background / 20 :: sha256 c6c4138ee5c86423
\newcommand{\MlstRocNoCurveEemm}{the \Kfull- and \Kres-class models and the three random partitions with controlled flavor boundaries}  % experiments/FIGS/data/v2/index.json :: curves_v2/pairs curve ee_vs_mm: configurations of F4 with a run whose cut is null or leaves fewer than 2 grid points (F4 draws no curve for them) :: sha256 c6c4138ee5c86423 :: L188, L162, FLAV_F0, FLAV_F1, FLAV_F1R
\newcommand{\MlstRocCutJcOneHbbFull}{0.40}  % experiments/FIGS/data/v2/index.json :: curves_v2/jc1_ft curve label_Hbb: largest cut.eps_s over the 188-class runs ('not reached' when a run's cut is null or leaves fewer than 2 grid points) :: sha256 c6c4138ee5c86423
\newcommand{\MlstRocCutJcOneHccFull}{0.37}  % experiments/FIGS/data/v2/index.json :: curves_v2/jc1_ft curve label_Hcc: largest cut.eps_s over the 188-class runs ('not reached' when a run's cut is null or leaves fewer than 2 grid points) :: sha256 c6c4138ee5c86423
\newcommand{\MlstRocCutJcOneHggFull}{0.04}  % experiments/FIGS/data/v2/index.json :: curves_v2/jc1_ft curve label_Hgg: largest cut.eps_s over the 188-class runs ('not reached' when a run's cut is null or leaves fewer than 2 grid points) :: sha256 c6c4138ee5c86423
\newcommand{\MlstRocCutJcOneHfourqFull}{0.26}  % experiments/FIGS/data/v2/index.json :: curves_v2/jc1_ft curve label_H4q: largest cut.eps_s over the 188-class runs ('not reached' when a run's cut is null or leaves fewer than 2 grid points) :: sha256 c6c4138ee5c86423
\newcommand{\MlstRocCutJcOneHqqlFull}{0.98}  % experiments/FIGS/data/v2/index.json :: curves_v2/jc1_ft curve label_Hqql: largest cut.eps_s over the 188-class runs ('not reached' when a run's cut is null or leaves fewer than 2 grid points) :: sha256 c6c4138ee5c86423
\newcommand{\MlstRocCutJcOneTbqqFull}{0.53}  % experiments/FIGS/data/v2/index.json :: curves_v2/jc1_ft curve label_Tbqq: largest cut.eps_s over the 188-class runs ('not reached' when a run's cut is null or leaves fewer than 2 grid points) :: sha256 c6c4138ee5c86423
\newcommand{\MlstRocCutJcOneTblFull}{not reached}  % experiments/FIGS/data/v2/index.json :: curves_v2/jc1_ft curve label_Tbl: largest cut.eps_s over the 188-class runs ('not reached' when a run's cut is null or leaves fewer than 2 grid points) :: sha256 c6c4138ee5c86423
\newcommand{\MlstRocCutJcOneWqqFull}{0.15}  % experiments/FIGS/data/v2/index.json :: curves_v2/jc1_ft curve label_Wqq: largest cut.eps_s over the 188-class runs ('not reached' when a run's cut is null or leaves fewer than 2 grid points) :: sha256 c6c4138ee5c86423
\newcommand{\MlstRocCutJcOneZqqFull}{0.12}  % experiments/FIGS/data/v2/index.json :: curves_v2/jc1_ft curve label_Zqq: largest cut.eps_s over the 188-class runs ('not reached' when a run's cut is null or leaves fewer than 2 grid points) :: sha256 c6c4138ee5c86423
\newcommand{\MlstFtTopRfiftyCoarseOverFull}{0.932\,\ensuremath{\pm}\,0.018}  % experiments/FIGS/data/v2/index.json :: curves_v2/top_ft curve top: rejection at eps_S = 0.5 of R16_Q1 over L188, run by run :: sha256 c6c4138ee5c86423
\newcommand{\MlstFtTopRfiftyCoarseOverFullPr}{6.8\% less rejection}  % experiments/FIGS/data/v2/index.json :: curves_v2/top_ft curve top: rejection at eps_S = 0.5 of R16_Q1 over L188, run by run :: sha256 c6c4138ee5c86423 :: prose twin: two significant figures, direction in words
\newcommand{\MlstProbeUnconverged}{0}  % experiments/FIGS/data/v2/index.json :: linear probes of pretrained models that did not converge (count) :: sha256 c6c4138ee5c86423
\newcommand{\MlstProbeUnconvergedInit}{0}  % experiments/FIGS/data/v2/index.json :: linear probes of the randomly initialized networks that did not converge (count) :: sha256 c6c4138ee5c86423
\newcommand{\MlstProbeFits}{895}  % experiments/FIGS/data/v2/index.json :: selected linear-probe fits of the reported checkpoints (count) :: sha256 c6c4138ee5c86423
\newcommand{\MlstProbeEdgeC}{16}  % experiments/FIGS/data/v2/index.json :: of these, fits whose selected C is the largest of the grid (count) :: sha256 c6c4138ee5c86423
\newcommand{\MlstMlpUnconverged}{0}  % experiments/FIGS/data/v2/index.json :: nonlinear probes that did not reach their plateau within the epoch cap (count) :: sha256 c6c4138ee5c86423
\newcommand{\MlstMlpCmpJcTwoResOverFull}{1.019\,\ensuremath{\pm}\,0.007}  % experiments/FIGS/data/v2/index.json :: nonlinear probe jc2_188: (1 - AUC) of L162 over L188, run by run :: sha256 c6c4138ee5c86423
\newcommand{\MlstMlpCmpJcTwoResOverFullPr}{1.9\% more errors}  % experiments/FIGS/data/v2/index.json :: nonlinear probe jc2_188: (1 - AUC) of L162 over L188, run by run :: sha256 c6c4138ee5c86423 :: prose twin: two significant figures, direction in words
\newcommand{\MlstMlpCmpJcTwoMidOverFull}{1.170\,\ensuremath{\pm}\,0.006}  % experiments/FIGS/data/v2/index.json :: nonlinear probe jc2_188: (1 - AUC) of R42_Q1 over L188, run by run :: sha256 c6c4138ee5c86423
\newcommand{\MlstMlpCmpJcTwoMidOverFullPr}{17\% more errors}  % experiments/FIGS/data/v2/index.json :: nonlinear probe jc2_188: (1 - AUC) of R42_Q1 over L188, run by run :: sha256 c6c4138ee5c86423 :: prose twin: two significant figures, direction in words
\newcommand{\MlstMlpCmpJcTwoCoarseOverFull}{1.388\,\ensuremath{\pm}\,0.021}  % experiments/FIGS/data/v2/index.json :: nonlinear probe jc2_188: (1 - AUC) of R16_Q1 over L188, run by run :: sha256 c6c4138ee5c86423
\newcommand{\MlstMlpCmpJcTwoCoarseOverFullPr}{39\% more errors}  % experiments/FIGS/data/v2/index.json :: nonlinear probe jc2_188: (1 - AUC) of R16_Q1 over L188, run by run :: sha256 c6c4138ee5c86423 :: prose twin: two significant figures, direction in words
\newcommand{\MlstMlpCmpJcTwoResMassOverRes}{\pending{no run yet}}  % experiments/FIGS/data/v2/index.json :: nonlinear probe jc2_188: (1 - AUC) of L162_MASS over L162, run by run :: sha256 c6c4138ee5c86423
\newcommand{\MlstMlpCmpJcTwoResMassOverResPr}{\pending{no run yet}}  % experiments/FIGS/data/v2/index.json :: nonlinear probe jc2_188: (1 - AUC) of L162_MASS over L162, run by run :: sha256 c6c4138ee5c86423 :: prose twin: two significant figures, direction in words
\newcommand{\MlstMlpCmpJcTwoResMassOverResPrMax}{\pending{no run yet}}  % experiments/FIGS/data/v2/index.json :: nonlinear probe jc2_188: (1 - AUC) of L162_MASS over L162, run by run :: sha256 c6c4138ee5c86423 :: prose twin of the run farther from 1
\newcommand{\MlstMlpCmpJcTwoCoarseMassOverCoarse}{\pending{no run yet}}  % experiments/FIGS/data/v2/index.json :: nonlinear probe jc2_188: (1 - AUC) of R16_Q1_MASS over R16_Q1, run by run :: sha256 c6c4138ee5c86423
\newcommand{\MlstMlpCmpJcTwoCoarseMassOverCoarsePr}{\pending{no run yet}}  % experiments/FIGS/data/v2/index.json :: nonlinear probe jc2_188: (1 - AUC) of R16_Q1_MASS over R16_Q1, run by run :: sha256 c6c4138ee5c86423 :: prose twin: two significant figures, direction in words
\newcommand{\MlstMlpCmpJcTwoCoarseMassOverCoarsePrMax}{\pending{no run yet}}  % experiments/FIGS/data/v2/index.json :: nonlinear probe jc2_188: (1 - AUC) of R16_Q1_MASS over R16_Q1, run by run :: sha256 c6c4138ee5c86423 :: prose twin of the run farther from 1
\newcommand{\MlstMlpCmpJcTwoFzeroOverFull}{1.294 and 1.310}  % experiments/FIGS/data/v2/index.json :: nonlinear probe jc2_188: (1 - AUC) of FLAV_F0 over L188, run by run :: sha256 c6c4138ee5c86423
\newcommand{\MlstMlpCmpJcTwoFzeroOverFullPr}{at least 29\% more errors}  % experiments/FIGS/data/v2/index.json :: nonlinear probe jc2_188: (1 - AUC) of FLAV_F0 over L188, run by run :: sha256 c6c4138ee5c86423 :: prose twin: two significant figures, direction in words
\newcommand{\MlstMlpCmpJcTwoFzeroOverFullPrMax}{up to 31\% more errors}  % experiments/FIGS/data/v2/index.json :: nonlinear probe jc2_188: (1 - AUC) of FLAV_F0 over L188, run by run :: sha256 c6c4138ee5c86423 :: prose twin of the run farther from 1
\newcommand{\MlstMlpCmpJcTwoFoneOverFzero}{0.903 and 0.885}  % experiments/FIGS/data/v2/index.json :: nonlinear probe jc2_188: (1 - AUC) of FLAV_F1 over FLAV_F0, run by run :: sha256 c6c4138ee5c86423
\newcommand{\MlstMlpCmpJcTwoFoneOverFzeroPr}{at least 9.7\% fewer errors}  % experiments/FIGS/data/v2/index.json :: nonlinear probe jc2_188: (1 - AUC) of FLAV_F1 over FLAV_F0, run by run :: sha256 c6c4138ee5c86423 :: prose twin: two significant figures, direction in words
\newcommand{\MlstMlpCmpJcTwoFoneOverFzeroPrMax}{up to 12\% fewer errors}  % experiments/FIGS/data/v2/index.json :: nonlinear probe jc2_188: (1 - AUC) of FLAV_F1 over FLAV_F0, run by run :: sha256 c6c4138ee5c86423 :: prose twin of the run farther from 1
\newcommand{\MlstMlpCmpJcTwoFonerOverFone}{0.9916 and 0.9918}  % experiments/FIGS/data/v2/index.json :: nonlinear probe jc2_188: (1 - AUC) of FLAV_F1R over FLAV_F1, run by run :: sha256 c6c4138ee5c86423
\newcommand{\MlstMlpCmpJcTwoFonerOverFonePr}{equal within 1\%}  % experiments/FIGS/data/v2/index.json :: nonlinear probe jc2_188: (1 - AUC) of FLAV_F1R over FLAV_F1, run by run :: sha256 c6c4138ee5c86423 :: prose twin: two significant figures, direction in words
\newcommand{\MlstMlpCmpJcTwoFonerOverFonePrMax}{equal within 1\%}  % experiments/FIGS/data/v2/index.json :: nonlinear probe jc2_188: (1 - AUC) of FLAV_F1R over FLAV_F1, run by run :: sha256 c6c4138ee5c86423 :: prose twin of the run farther from 1
\newcommand{\MlstMlpCmpJcTwoFonerOverFzero}{0.895 and 0.877}  % experiments/FIGS/data/v2/index.json :: nonlinear probe jc2_188: (1 - AUC) of FLAV_F1R over FLAV_F0, run by run :: sha256 c6c4138ee5c86423
\newcommand{\MlstMlpCmpJcTwoFonerOverFzeroPr}{at least 10\% fewer errors}  % experiments/FIGS/data/v2/index.json :: nonlinear probe jc2_188: (1 - AUC) of FLAV_F1R over FLAV_F0, run by run :: sha256 c6c4138ee5c86423 :: prose twin: two significant figures, direction in words
\newcommand{\MlstMlpCmpJcTwoFonerOverFzeroPrMax}{up to 12\% fewer errors}  % experiments/FIGS/data/v2/index.json :: nonlinear probe jc2_188: (1 - AUC) of FLAV_F1R over FLAV_F0, run by run :: sha256 c6c4138ee5c86423 :: prose twin of the run farther from 1
\newcommand{\MlstMlpCmpJcTwoFzeroOverCoarse}{0.9461 and 0.9456}  % experiments/FIGS/data/v2/index.json :: nonlinear probe jc2_188: (1 - AUC) of FLAV_F0 over R16_Q1, run by run :: sha256 c6c4138ee5c86423
\newcommand{\MlstMlpCmpJcTwoFzeroOverCoarsePr}{at least 5.4\% fewer errors}  % experiments/FIGS/data/v2/index.json :: nonlinear probe jc2_188: (1 - AUC) of FLAV_F0 over R16_Q1, run by run :: sha256 c6c4138ee5c86423 :: prose twin: two significant figures, direction in words
\newcommand{\MlstMlpCmpJcTwoFzeroOverCoarsePrMax}{up to 5.4\% fewer errors}  % experiments/FIGS/data/v2/index.json :: nonlinear probe jc2_188: (1 - AUC) of FLAV_F0 over R16_Q1, run by run :: sha256 c6c4138ee5c86423 :: prose twin of the run farther from 1
\newcommand{\MlstMlpCmpJcTwoRandOverFull}{\pending{no run yet}}  % experiments/FIGS/data/v2/index.json :: nonlinear probe jc2_188: (1 - AUC) of each random partition over the 188-class run of its seed :: sha256 c6c4138ee5c86423
\newcommand{\MlstMlpCmpJcTwoRandOverFullPr}{\pending{no run yet}}  % experiments/FIGS/data/v2/index.json :: nonlinear probe jc2_188: (1 - AUC) of each random partition over the 188-class run of its seed :: sha256 c6c4138ee5c86423 :: prose twin: two significant figures, direction in words
\newcommand{\MlstMlpCmpJcTwoResResOverFull}{1.001\,\ensuremath{\pm}\,0.007}  % experiments/FIGS/data/v2/index.json :: nonlinear probe jc2_188_res: (1 - AUC) of L162 over L188, run by run :: sha256 c6c4138ee5c86423
\newcommand{\MlstMlpCmpJcTwoResResOverFullPr}{equal within 1\%}  % experiments/FIGS/data/v2/index.json :: nonlinear probe jc2_188_res: (1 - AUC) of L162 over L188, run by run :: sha256 c6c4138ee5c86423 :: prose twin: two significant figures, direction in words
\newcommand{\MlstMlpCmpJcTwoResMidOverFull}{1.192\,\ensuremath{\pm}\,0.008}  % experiments/FIGS/data/v2/index.json :: nonlinear probe jc2_188_res: (1 - AUC) of R42_Q1 over L188, run by run :: sha256 c6c4138ee5c86423
\newcommand{\MlstMlpCmpJcTwoResMidOverFullPr}{19\% more errors}  % experiments/FIGS/data/v2/index.json :: nonlinear probe jc2_188_res: (1 - AUC) of R42_Q1 over L188, run by run :: sha256 c6c4138ee5c86423 :: prose twin: two significant figures, direction in words
\newcommand{\MlstMlpCmpJcTwoResCoarseOverFull}{1.408\,\ensuremath{\pm}\,0.026}  % experiments/FIGS/data/v2/index.json :: nonlinear probe jc2_188_res: (1 - AUC) of R16_Q1 over L188, run by run :: sha256 c6c4138ee5c86423
\newcommand{\MlstMlpCmpJcTwoResCoarseOverFullPr}{41\% more errors}  % experiments/FIGS/data/v2/index.json :: nonlinear probe jc2_188_res: (1 - AUC) of R16_Q1 over L188, run by run :: sha256 c6c4138ee5c86423 :: prose twin: two significant figures, direction in words
\newcommand{\MlstMlpCmpJcTwoResResMassOverRes}{\pending{no run yet}}  % experiments/FIGS/data/v2/index.json :: nonlinear probe jc2_188_res: (1 - AUC) of L162_MASS over L162, run by run :: sha256 c6c4138ee5c86423
\newcommand{\MlstMlpCmpJcTwoResResMassOverResPr}{\pending{no run yet}}  % experiments/FIGS/data/v2/index.json :: nonlinear probe jc2_188_res: (1 - AUC) of L162_MASS over L162, run by run :: sha256 c6c4138ee5c86423 :: prose twin: two significant figures, direction in words
\newcommand{\MlstMlpCmpJcTwoResResMassOverResPrMax}{\pending{no run yet}}  % experiments/FIGS/data/v2/index.json :: nonlinear probe jc2_188_res: (1 - AUC) of L162_MASS over L162, run by run :: sha256 c6c4138ee5c86423 :: prose twin of the run farther from 1
\newcommand{\MlstMlpCmpJcTwoResCoarseMassOverCoarse}{\pending{no run yet}}  % experiments/FIGS/data/v2/index.json :: nonlinear probe jc2_188_res: (1 - AUC) of R16_Q1_MASS over R16_Q1, run by run :: sha256 c6c4138ee5c86423
\newcommand{\MlstMlpCmpJcTwoResCoarseMassOverCoarsePr}{\pending{no run yet}}  % experiments/FIGS/data/v2/index.json :: nonlinear probe jc2_188_res: (1 - AUC) of R16_Q1_MASS over R16_Q1, run by run :: sha256 c6c4138ee5c86423 :: prose twin: two significant figures, direction in words
\newcommand{\MlstMlpCmpJcTwoResCoarseMassOverCoarsePrMax}{\pending{no run yet}}  % experiments/FIGS/data/v2/index.json :: nonlinear probe jc2_188_res: (1 - AUC) of R16_Q1_MASS over R16_Q1, run by run :: sha256 c6c4138ee5c86423 :: prose twin of the run farther from 1
\newcommand{\MlstMlpCmpJcTwoResFzeroOverFull}{1.285 and 1.300}  % experiments/FIGS/data/v2/index.json :: nonlinear probe jc2_188_res: (1 - AUC) of FLAV_F0 over L188, run by run :: sha256 c6c4138ee5c86423
\newcommand{\MlstMlpCmpJcTwoResFzeroOverFullPr}{at least 29\% more errors}  % experiments/FIGS/data/v2/index.json :: nonlinear probe jc2_188_res: (1 - AUC) of FLAV_F0 over L188, run by run :: sha256 c6c4138ee5c86423 :: prose twin: two significant figures, direction in words
\newcommand{\MlstMlpCmpJcTwoResFzeroOverFullPrMax}{up to 30\% more errors}  % experiments/FIGS/data/v2/index.json :: nonlinear probe jc2_188_res: (1 - AUC) of FLAV_F0 over L188, run by run :: sha256 c6c4138ee5c86423 :: prose twin of the run farther from 1
\newcommand{\MlstMlpCmpJcTwoResFoneOverFzero}{0.909 and 0.889}  % experiments/FIGS/data/v2/index.json :: nonlinear probe jc2_188_res: (1 - AUC) of FLAV_F1 over FLAV_F0, run by run :: sha256 c6c4138ee5c86423
\newcommand{\MlstMlpCmpJcTwoResFoneOverFzeroPr}{at least 9.1\% fewer errors}  % experiments/FIGS/data/v2/index.json :: nonlinear probe jc2_188_res: (1 - AUC) of FLAV_F1 over FLAV_F0, run by run :: sha256 c6c4138ee5c86423 :: prose twin: two significant figures, direction in words
\newcommand{\MlstMlpCmpJcTwoResFoneOverFzeroPrMax}{up to 11\% fewer errors}  % experiments/FIGS/data/v2/index.json :: nonlinear probe jc2_188_res: (1 - AUC) of FLAV_F1 over FLAV_F0, run by run :: sha256 c6c4138ee5c86423 :: prose twin of the run farther from 1
\newcommand{\MlstMlpCmpJcTwoResFonerOverFone}{0.9793 and 0.9825}  % experiments/FIGS/data/v2/index.json :: nonlinear probe jc2_188_res: (1 - AUC) of FLAV_F1R over FLAV_F1, run by run :: sha256 c6c4138ee5c86423
\newcommand{\MlstMlpCmpJcTwoResFonerOverFonePr}{equal within 3\%}  % experiments/FIGS/data/v2/index.json :: nonlinear probe jc2_188_res: (1 - AUC) of FLAV_F1R over FLAV_F1, run by run :: sha256 c6c4138ee5c86423 :: prose twin: two significant figures, direction in words
\newcommand{\MlstMlpCmpJcTwoResFonerOverFonePrMax}{equal within 3\%}  % experiments/FIGS/data/v2/index.json :: nonlinear probe jc2_188_res: (1 - AUC) of FLAV_F1R over FLAV_F1, run by run :: sha256 c6c4138ee5c86423 :: prose twin of the run farther from 1
\newcommand{\MlstMlpCmpJcTwoResFonerOverFzero}{0.890 and 0.874}  % experiments/FIGS/data/v2/index.json :: nonlinear probe jc2_188_res: (1 - AUC) of FLAV_F1R over FLAV_F0, run by run :: sha256 c6c4138ee5c86423
\newcommand{\MlstMlpCmpJcTwoResFonerOverFzeroPr}{at least 11\% fewer errors}  % experiments/FIGS/data/v2/index.json :: nonlinear probe jc2_188_res: (1 - AUC) of FLAV_F1R over FLAV_F0, run by run :: sha256 c6c4138ee5c86423 :: prose twin: two significant figures, direction in words
\newcommand{\MlstMlpCmpJcTwoResFonerOverFzeroPrMax}{up to 13\% fewer errors}  % experiments/FIGS/data/v2/index.json :: nonlinear probe jc2_188_res: (1 - AUC) of FLAV_F1R over FLAV_F0, run by run :: sha256 c6c4138ee5c86423 :: prose twin of the run farther from 1
\newcommand{\MlstMlpCmpJcTwoResFzeroOverCoarse}{0.9299 and 0.9233}  % experiments/FIGS/data/v2/index.json :: nonlinear probe jc2_188_res: (1 - AUC) of FLAV_F0 over R16_Q1, run by run :: sha256 c6c4138ee5c86423
\newcommand{\MlstMlpCmpJcTwoResFzeroOverCoarsePr}{at least 7.0\% fewer errors}  % experiments/FIGS/data/v2/index.json :: nonlinear probe jc2_188_res: (1 - AUC) of FLAV_F0 over R16_Q1, run by run :: sha256 c6c4138ee5c86423 :: prose twin: two significant figures, direction in words
\newcommand{\MlstMlpCmpJcTwoResFzeroOverCoarsePrMax}{up to 7.7\% fewer errors}  % experiments/FIGS/data/v2/index.json :: nonlinear probe jc2_188_res: (1 - AUC) of FLAV_F0 over R16_Q1, run by run :: sha256 c6c4138ee5c86423 :: prose twin of the run farther from 1
\newcommand{\MlstMlpCmpJcTwoResRandOverFull}{\pending{no run yet}}  % experiments/FIGS/data/v2/index.json :: nonlinear probe jc2_188_res: (1 - AUC) of each random partition over the 188-class run of its seed :: sha256 c6c4138ee5c86423
\newcommand{\MlstMlpCmpJcTwoResRandOverFullPr}{\pending{no run yet}}  % experiments/FIGS/data/v2/index.json :: nonlinear probe jc2_188_res: (1 - AUC) of each random partition over the 188-class run of its seed :: sha256 c6c4138ee5c86423 :: prose twin: two significant figures, direction in words
\newcommand{\MlstMlpCmpJcTwoQcdResOverFull}{1.070\,\ensuremath{\pm}\,0.006}  % experiments/FIGS/data/v2/index.json :: nonlinear probe jc2_188_qcd: (1 - AUC) of L162 over L188, run by run :: sha256 c6c4138ee5c86423
\newcommand{\MlstMlpCmpJcTwoQcdResOverFullPr}{7.0\% more errors}  % experiments/FIGS/data/v2/index.json :: nonlinear probe jc2_188_qcd: (1 - AUC) of L162 over L188, run by run :: sha256 c6c4138ee5c86423 :: prose twin: two significant figures, direction in words
\newcommand{\MlstMlpCmpJcTwoQcdMidOverFull}{1.107\,\ensuremath{\pm}\,0.002}  % experiments/FIGS/data/v2/index.json :: nonlinear probe jc2_188_qcd: (1 - AUC) of R42_Q1 over L188, run by run :: sha256 c6c4138ee5c86423
\newcommand{\MlstMlpCmpJcTwoQcdMidOverFullPr}{11\% more errors}  % experiments/FIGS/data/v2/index.json :: nonlinear probe jc2_188_qcd: (1 - AUC) of R42_Q1 over L188, run by run :: sha256 c6c4138ee5c86423 :: prose twin: two significant figures, direction in words
\newcommand{\MlstMlpCmpJcTwoQcdCoarseOverFull}{1.333\,\ensuremath{\pm}\,0.012}  % experiments/FIGS/data/v2/index.json :: nonlinear probe jc2_188_qcd: (1 - AUC) of R16_Q1 over L188, run by run :: sha256 c6c4138ee5c86423
\newcommand{\MlstMlpCmpJcTwoQcdCoarseOverFullPr}{33\% more errors}  % experiments/FIGS/data/v2/index.json :: nonlinear probe jc2_188_qcd: (1 - AUC) of R16_Q1 over L188, run by run :: sha256 c6c4138ee5c86423 :: prose twin: two significant figures, direction in words
\newcommand{\MlstMlpCmpJcTwoQcdResMassOverRes}{\pending{no run yet}}  % experiments/FIGS/data/v2/index.json :: nonlinear probe jc2_188_qcd: (1 - AUC) of L162_MASS over L162, run by run :: sha256 c6c4138ee5c86423
\newcommand{\MlstMlpCmpJcTwoQcdResMassOverResPr}{\pending{no run yet}}  % experiments/FIGS/data/v2/index.json :: nonlinear probe jc2_188_qcd: (1 - AUC) of L162_MASS over L162, run by run :: sha256 c6c4138ee5c86423 :: prose twin: two significant figures, direction in words
\newcommand{\MlstMlpCmpJcTwoQcdResMassOverResPrMax}{\pending{no run yet}}  % experiments/FIGS/data/v2/index.json :: nonlinear probe jc2_188_qcd: (1 - AUC) of L162_MASS over L162, run by run :: sha256 c6c4138ee5c86423 :: prose twin of the run farther from 1
\newcommand{\MlstMlpCmpJcTwoQcdCoarseMassOverCoarse}{\pending{no run yet}}  % experiments/FIGS/data/v2/index.json :: nonlinear probe jc2_188_qcd: (1 - AUC) of R16_Q1_MASS over R16_Q1, run by run :: sha256 c6c4138ee5c86423
\newcommand{\MlstMlpCmpJcTwoQcdCoarseMassOverCoarsePr}{\pending{no run yet}}  % experiments/FIGS/data/v2/index.json :: nonlinear probe jc2_188_qcd: (1 - AUC) of R16_Q1_MASS over R16_Q1, run by run :: sha256 c6c4138ee5c86423 :: prose twin: two significant figures, direction in words
\newcommand{\MlstMlpCmpJcTwoQcdCoarseMassOverCoarsePrMax}{\pending{no run yet}}  % experiments/FIGS/data/v2/index.json :: nonlinear probe jc2_188_qcd: (1 - AUC) of R16_Q1_MASS over R16_Q1, run by run :: sha256 c6c4138ee5c86423 :: prose twin of the run farther from 1
\newcommand{\MlstMlpCmpJcTwoQcdFzeroOverFull}{1.319 and 1.339}  % experiments/FIGS/data/v2/index.json :: nonlinear probe jc2_188_qcd: (1 - AUC) of FLAV_F0 over L188, run by run :: sha256 c6c4138ee5c86423
\newcommand{\MlstMlpCmpJcTwoQcdFzeroOverFullPr}{at least 32\% more errors}  % experiments/FIGS/data/v2/index.json :: nonlinear probe jc2_188_qcd: (1 - AUC) of FLAV_F0 over L188, run by run :: sha256 c6c4138ee5c86423 :: prose twin: two significant figures, direction in words
\newcommand{\MlstMlpCmpJcTwoQcdFzeroOverFullPrMax}{up to 34\% more errors}  % experiments/FIGS/data/v2/index.json :: nonlinear probe jc2_188_qcd: (1 - AUC) of FLAV_F0 over L188, run by run :: sha256 c6c4138ee5c86423 :: prose twin of the run farther from 1
\newcommand{\MlstMlpCmpJcTwoQcdFoneOverFzero}{0.886 and 0.872}  % experiments/FIGS/data/v2/index.json :: nonlinear probe jc2_188_qcd: (1 - AUC) of FLAV_F1 over FLAV_F0, run by run :: sha256 c6c4138ee5c86423
\newcommand{\MlstMlpCmpJcTwoQcdFoneOverFzeroPr}{at least 11\% fewer errors}  % experiments/FIGS/data/v2/index.json :: nonlinear probe jc2_188_qcd: (1 - AUC) of FLAV_F1 over FLAV_F0, run by run :: sha256 c6c4138ee5c86423 :: prose twin: two significant figures, direction in words
\newcommand{\MlstMlpCmpJcTwoQcdFoneOverFzeroPrMax}{up to 13\% fewer errors}  % experiments/FIGS/data/v2/index.json :: nonlinear probe jc2_188_qcd: (1 - AUC) of FLAV_F1 over FLAV_F0, run by run :: sha256 c6c4138ee5c86423 :: prose twin of the run farther from 1
\newcommand{\MlstMlpCmpJcTwoQcdFonerOverFone}{1.026 and 1.018}  % experiments/FIGS/data/v2/index.json :: nonlinear probe jc2_188_qcd: (1 - AUC) of FLAV_F1R over FLAV_F1, run by run :: sha256 c6c4138ee5c86423
\newcommand{\MlstMlpCmpJcTwoQcdFonerOverFonePr}{at least 1.8\% more errors}  % experiments/FIGS/data/v2/index.json :: nonlinear probe jc2_188_qcd: (1 - AUC) of FLAV_F1R over FLAV_F1, run by run :: sha256 c6c4138ee5c86423 :: prose twin: two significant figures, direction in words
\newcommand{\MlstMlpCmpJcTwoQcdFonerOverFonePrMax}{up to 2.6\% more errors}  % experiments/FIGS/data/v2/index.json :: nonlinear probe jc2_188_qcd: (1 - AUC) of FLAV_F1R over FLAV_F1, run by run :: sha256 c6c4138ee5c86423 :: prose twin of the run farther from 1
\newcommand{\MlstMlpCmpJcTwoQcdFonerOverFzero}{0.909 and 0.887}  % experiments/FIGS/data/v2/index.json :: nonlinear probe jc2_188_qcd: (1 - AUC) of FLAV_F1R over FLAV_F0, run by run :: sha256 c6c4138ee5c86423
\newcommand{\MlstMlpCmpJcTwoQcdFonerOverFzeroPr}{at least 9.1\% fewer errors}  % experiments/FIGS/data/v2/index.json :: nonlinear probe jc2_188_qcd: (1 - AUC) of FLAV_F1R over FLAV_F0, run by run :: sha256 c6c4138ee5c86423 :: prose twin: two significant figures, direction in words
\newcommand{\MlstMlpCmpJcTwoQcdFonerOverFzeroPrMax}{up to 11\% fewer errors}  % experiments/FIGS/data/v2/index.json :: nonlinear probe jc2_188_qcd: (1 - AUC) of FLAV_F1R over FLAV_F0, run by run :: sha256 c6c4138ee5c86423 :: prose twin of the run farther from 1
\newcommand{\MlstMlpCmpJcTwoQcdFzeroOverCoarse}{0.993 and 1.012}  % experiments/FIGS/data/v2/index.json :: nonlinear probe jc2_188_qcd: (1 - AUC) of FLAV_F0 over R16_Q1, run by run :: sha256 c6c4138ee5c86423
\newcommand{\MlstMlpCmpJcTwoQcdFzeroOverCoarsePr}{equal within 2\%}  % experiments/FIGS/data/v2/index.json :: nonlinear probe jc2_188_qcd: (1 - AUC) of FLAV_F0 over R16_Q1, run by run :: sha256 c6c4138ee5c86423 :: prose twin: two significant figures, direction in words
\newcommand{\MlstMlpCmpJcTwoQcdFzeroOverCoarsePrMax}{equal within 2\%}  % experiments/FIGS/data/v2/index.json :: nonlinear probe jc2_188_qcd: (1 - AUC) of FLAV_F0 over R16_Q1, run by run :: sha256 c6c4138ee5c86423 :: prose twin of the run farther from 1
\newcommand{\MlstMlpCmpJcTwoQcdRandOverFull}{\pending{no run yet}}  % experiments/FIGS/data/v2/index.json :: nonlinear probe jc2_188_qcd: (1 - AUC) of each random partition over the 188-class run of its seed :: sha256 c6c4138ee5c86423
\newcommand{\MlstMlpCmpJcTwoQcdRandOverFullPr}{\pending{no run yet}}  % experiments/FIGS/data/v2/index.json :: nonlinear probe jc2_188_qcd: (1 - AUC) of each random partition over the 188-class run of its seed :: sha256 c6c4138ee5c86423 :: prose twin: two significant figures, direction in words
\newcommand{\MlstMlpCmpBcResOverFull}{1.002\,\ensuremath{\pm}\,0.006}  % experiments/FIGS/data/v2/index.json :: nonlinear probe bc_vs_bq_cs: (1 - AUC) of L162 over L188, run by run :: sha256 c6c4138ee5c86423
\newcommand{\MlstMlpCmpBcResOverFullPr}{equal within 1\%}  % experiments/FIGS/data/v2/index.json :: nonlinear probe bc_vs_bq_cs: (1 - AUC) of L162 over L188, run by run :: sha256 c6c4138ee5c86423 :: prose twin: two significant figures, direction in words
\newcommand{\MlstMlpCmpBcMidOverFull}{1.104\,\ensuremath{\pm}\,0.017}  % experiments/FIGS/data/v2/index.json :: nonlinear probe bc_vs_bq_cs: (1 - AUC) of R42_Q1 over L188, run by run :: sha256 c6c4138ee5c86423
\newcommand{\MlstMlpCmpBcMidOverFullPr}{10\% more errors}  % experiments/FIGS/data/v2/index.json :: nonlinear probe bc_vs_bq_cs: (1 - AUC) of R42_Q1 over L188, run by run :: sha256 c6c4138ee5c86423 :: prose twin: two significant figures, direction in words
\newcommand{\MlstMlpCmpBcCoarseOverFull}{1.624\,\ensuremath{\pm}\,0.022}  % experiments/FIGS/data/v2/index.json :: nonlinear probe bc_vs_bq_cs: (1 - AUC) of R16_Q1 over L188, run by run :: sha256 c6c4138ee5c86423
\newcommand{\MlstMlpCmpBcCoarseOverFullPr}{62\% more errors}  % experiments/FIGS/data/v2/index.json :: nonlinear probe bc_vs_bq_cs: (1 - AUC) of R16_Q1 over L188, run by run :: sha256 c6c4138ee5c86423 :: prose twin: two significant figures, direction in words
\newcommand{\MlstMlpCmpBcResMassOverRes}{\pending{no run yet}}  % experiments/FIGS/data/v2/index.json :: nonlinear probe bc_vs_bq_cs: (1 - AUC) of L162_MASS over L162, run by run :: sha256 c6c4138ee5c86423
\newcommand{\MlstMlpCmpBcResMassOverResPr}{\pending{no run yet}}  % experiments/FIGS/data/v2/index.json :: nonlinear probe bc_vs_bq_cs: (1 - AUC) of L162_MASS over L162, run by run :: sha256 c6c4138ee5c86423 :: prose twin: two significant figures, direction in words
\newcommand{\MlstMlpCmpBcResMassOverResPrMax}{\pending{no run yet}}  % experiments/FIGS/data/v2/index.json :: nonlinear probe bc_vs_bq_cs: (1 - AUC) of L162_MASS over L162, run by run :: sha256 c6c4138ee5c86423 :: prose twin of the run farther from 1
\newcommand{\MlstMlpCmpBcCoarseMassOverCoarse}{\pending{no run yet}}  % experiments/FIGS/data/v2/index.json :: nonlinear probe bc_vs_bq_cs: (1 - AUC) of R16_Q1_MASS over R16_Q1, run by run :: sha256 c6c4138ee5c86423
\newcommand{\MlstMlpCmpBcCoarseMassOverCoarsePr}{\pending{no run yet}}  % experiments/FIGS/data/v2/index.json :: nonlinear probe bc_vs_bq_cs: (1 - AUC) of R16_Q1_MASS over R16_Q1, run by run :: sha256 c6c4138ee5c86423 :: prose twin: two significant figures, direction in words
\newcommand{\MlstMlpCmpBcCoarseMassOverCoarsePrMax}{\pending{no run yet}}  % experiments/FIGS/data/v2/index.json :: nonlinear probe bc_vs_bq_cs: (1 - AUC) of R16_Q1_MASS over R16_Q1, run by run :: sha256 c6c4138ee5c86423 :: prose twin of the run farther from 1
\newcommand{\MlstMlpCmpBcFzeroOverFull}{1.633 and 1.679}  % experiments/FIGS/data/v2/index.json :: nonlinear probe bc_vs_bq_cs: (1 - AUC) of FLAV_F0 over L188, run by run :: sha256 c6c4138ee5c86423
\newcommand{\MlstMlpCmpBcFzeroOverFullPr}{at least 63\% more errors}  % experiments/FIGS/data/v2/index.json :: nonlinear probe bc_vs_bq_cs: (1 - AUC) of FLAV_F0 over L188, run by run :: sha256 c6c4138ee5c86423 :: prose twin: two significant figures, direction in words
\newcommand{\MlstMlpCmpBcFzeroOverFullPrMax}{up to 68\% more errors}  % experiments/FIGS/data/v2/index.json :: nonlinear probe bc_vs_bq_cs: (1 - AUC) of FLAV_F0 over L188, run by run :: sha256 c6c4138ee5c86423 :: prose twin of the run farther from 1
\newcommand{\MlstMlpCmpBcFoneOverFzero}{0.848 and 0.810}  % experiments/FIGS/data/v2/index.json :: nonlinear probe bc_vs_bq_cs: (1 - AUC) of FLAV_F1 over FLAV_F0, run by run :: sha256 c6c4138ee5c86423
\newcommand{\MlstMlpCmpBcFoneOverFzeroPr}{at least 15\% fewer errors}  % experiments/FIGS/data/v2/index.json :: nonlinear probe bc_vs_bq_cs: (1 - AUC) of FLAV_F1 over FLAV_F0, run by run :: sha256 c6c4138ee5c86423 :: prose twin: two significant figures, direction in words
\newcommand{\MlstMlpCmpBcFoneOverFzeroPrMax}{up to 19\% fewer errors}  % experiments/FIGS/data/v2/index.json :: nonlinear probe bc_vs_bq_cs: (1 - AUC) of FLAV_F1 over FLAV_F0, run by run :: sha256 c6c4138ee5c86423 :: prose twin of the run farther from 1
\newcommand{\MlstMlpCmpBcFonerOverFone}{0.9812 and 0.9815}  % experiments/FIGS/data/v2/index.json :: nonlinear probe bc_vs_bq_cs: (1 - AUC) of FLAV_F1R over FLAV_F1, run by run :: sha256 c6c4138ee5c86423
\newcommand{\MlstMlpCmpBcFonerOverFonePr}{equal within 2\%}  % experiments/FIGS/data/v2/index.json :: nonlinear probe bc_vs_bq_cs: (1 - AUC) of FLAV_F1R over FLAV_F1, run by run :: sha256 c6c4138ee5c86423 :: prose twin: two significant figures, direction in words
\newcommand{\MlstMlpCmpBcFonerOverFonePrMax}{equal within 2\%}  % experiments/FIGS/data/v2/index.json :: nonlinear probe bc_vs_bq_cs: (1 - AUC) of FLAV_F1R over FLAV_F1, run by run :: sha256 c6c4138ee5c86423 :: prose twin of the run farther from 1
\newcommand{\MlstMlpCmpBcFonerOverFzero}{0.833 and 0.795}  % experiments/FIGS/data/v2/index.json :: nonlinear probe bc_vs_bq_cs: (1 - AUC) of FLAV_F1R over FLAV_F0, run by run :: sha256 c6c4138ee5c86423
\newcommand{\MlstMlpCmpBcFonerOverFzeroPr}{at least 17\% fewer errors}  % experiments/FIGS/data/v2/index.json :: nonlinear probe bc_vs_bq_cs: (1 - AUC) of FLAV_F1R over FLAV_F0, run by run :: sha256 c6c4138ee5c86423 :: prose twin: two significant figures, direction in words
\newcommand{\MlstMlpCmpBcFonerOverFzeroPrMax}{up to 21\% fewer errors}  % experiments/FIGS/data/v2/index.json :: nonlinear probe bc_vs_bq_cs: (1 - AUC) of FLAV_F1R over FLAV_F0, run by run :: sha256 c6c4138ee5c86423 :: prose twin of the run farther from 1
\newcommand{\MlstMlpCmpBcFzeroOverCoarse}{0.994 and 1.030}  % experiments/FIGS/data/v2/index.json :: nonlinear probe bc_vs_bq_cs: (1 - AUC) of FLAV_F0 over R16_Q1, run by run :: sha256 c6c4138ee5c86423
\newcommand{\MlstMlpCmpBcFzeroOverCoarsePr}{equal within 3\%}  % experiments/FIGS/data/v2/index.json :: nonlinear probe bc_vs_bq_cs: (1 - AUC) of FLAV_F0 over R16_Q1, run by run :: sha256 c6c4138ee5c86423 :: prose twin: two significant figures, direction in words
\newcommand{\MlstMlpCmpBcFzeroOverCoarsePrMax}{equal within 3\%}  % experiments/FIGS/data/v2/index.json :: nonlinear probe bc_vs_bq_cs: (1 - AUC) of FLAV_F0 over R16_Q1, run by run :: sha256 c6c4138ee5c86423 :: prose twin of the run farther from 1
\newcommand{\MlstMlpCmpBcRandOverFull}{\pending{no run yet}}  % experiments/FIGS/data/v2/index.json :: nonlinear probe bc_vs_bq_cs: (1 - AUC) of each random partition over the 188-class run of its seed :: sha256 c6c4138ee5c86423
\newcommand{\MlstMlpCmpBcRandOverFullPr}{\pending{no run yet}}  % experiments/FIGS/data/v2/index.json :: nonlinear probe bc_vs_bq_cs: (1 - AUC) of each random partition over the 188-class run of its seed :: sha256 c6c4138ee5c86423 :: prose twin: two significant figures, direction in words
\newcommand{\MlstMlpCmpBbResOverFull}{1.023\,\ensuremath{\pm}\,0.021}  % experiments/FIGS/data/v2/index.json :: nonlinear probe bb_vs_cc: (1 - AUC) of L162 over L188, run by run :: sha256 c6c4138ee5c86423
\newcommand{\MlstMlpCmpBbResOverFullPr}{equal within 3\%}  % experiments/FIGS/data/v2/index.json :: nonlinear probe bb_vs_cc: (1 - AUC) of L162 over L188, run by run :: sha256 c6c4138ee5c86423 :: prose twin: two significant figures, direction in words
\newcommand{\MlstMlpCmpBbMidOverFull}{1.125\,\ensuremath{\pm}\,0.011}  % experiments/FIGS/data/v2/index.json :: nonlinear probe bb_vs_cc: (1 - AUC) of R42_Q1 over L188, run by run :: sha256 c6c4138ee5c86423
\newcommand{\MlstMlpCmpBbMidOverFullPr}{13\% more errors}  % experiments/FIGS/data/v2/index.json :: nonlinear probe bb_vs_cc: (1 - AUC) of R42_Q1 over L188, run by run :: sha256 c6c4138ee5c86423 :: prose twin: two significant figures, direction in words
\newcommand{\MlstMlpCmpBbCoarseOverFull}{3.04\,\ensuremath{\pm}\,0.09}  % experiments/FIGS/data/v2/index.json :: nonlinear probe bb_vs_cc: (1 - AUC) of R16_Q1 over L188, run by run :: sha256 c6c4138ee5c86423
\newcommand{\MlstMlpCmpBbCoarseOverFullPr}{3.0 times the errors}  % experiments/FIGS/data/v2/index.json :: nonlinear probe bb_vs_cc: (1 - AUC) of R16_Q1 over L188, run by run :: sha256 c6c4138ee5c86423 :: prose twin: two significant figures, direction in words
\newcommand{\MlstMlpCmpBbResMassOverRes}{\pending{no run yet}}  % experiments/FIGS/data/v2/index.json :: nonlinear probe bb_vs_cc: (1 - AUC) of L162_MASS over L162, run by run :: sha256 c6c4138ee5c86423
\newcommand{\MlstMlpCmpBbResMassOverResPr}{\pending{no run yet}}  % experiments/FIGS/data/v2/index.json :: nonlinear probe bb_vs_cc: (1 - AUC) of L162_MASS over L162, run by run :: sha256 c6c4138ee5c86423 :: prose twin: two significant figures, direction in words
\newcommand{\MlstMlpCmpBbResMassOverResPrMax}{\pending{no run yet}}  % experiments/FIGS/data/v2/index.json :: nonlinear probe bb_vs_cc: (1 - AUC) of L162_MASS over L162, run by run :: sha256 c6c4138ee5c86423 :: prose twin of the run farther from 1
\newcommand{\MlstMlpCmpBbCoarseMassOverCoarse}{\pending{no run yet}}  % experiments/FIGS/data/v2/index.json :: nonlinear probe bb_vs_cc: (1 - AUC) of R16_Q1_MASS over R16_Q1, run by run :: sha256 c6c4138ee5c86423
\newcommand{\MlstMlpCmpBbCoarseMassOverCoarsePr}{\pending{no run yet}}  % experiments/FIGS/data/v2/index.json :: nonlinear probe bb_vs_cc: (1 - AUC) of R16_Q1_MASS over R16_Q1, run by run :: sha256 c6c4138ee5c86423 :: prose twin: two significant figures, direction in words
\newcommand{\MlstMlpCmpBbCoarseMassOverCoarsePrMax}{\pending{no run yet}}  % experiments/FIGS/data/v2/index.json :: nonlinear probe bb_vs_cc: (1 - AUC) of R16_Q1_MASS over R16_Q1, run by run :: sha256 c6c4138ee5c86423 :: prose twin of the run farther from 1
\newcommand{\MlstMlpCmpBbFzeroOverFull}{3.3 and 4.3}  % experiments/FIGS/data/v2/index.json :: nonlinear probe bb_vs_cc: (1 - AUC) of FLAV_F0 over L188, run by run :: sha256 c6c4138ee5c86423
\newcommand{\MlstMlpCmpBbFzeroOverFullPr}{at least 3.3 times the errors}  % experiments/FIGS/data/v2/index.json :: nonlinear probe bb_vs_cc: (1 - AUC) of FLAV_F0 over L188, run by run :: sha256 c6c4138ee5c86423 :: prose twin: two significant figures, direction in words
\newcommand{\MlstMlpCmpBbFzeroOverFullPrMax}{up to 4.3 times the errors}  % experiments/FIGS/data/v2/index.json :: nonlinear probe bb_vs_cc: (1 - AUC) of FLAV_F0 over L188, run by run :: sha256 c6c4138ee5c86423 :: prose twin of the run farther from 1
\newcommand{\MlstMlpCmpBbFoneOverFzero}{0.52 and 0.38}  % experiments/FIGS/data/v2/index.json :: nonlinear probe bb_vs_cc: (1 - AUC) of FLAV_F1 over FLAV_F0, run by run :: sha256 c6c4138ee5c86423
\newcommand{\MlstMlpCmpBbFoneOverFzeroPr}{at least 48\% fewer errors}  % experiments/FIGS/data/v2/index.json :: nonlinear probe bb_vs_cc: (1 - AUC) of FLAV_F1 over FLAV_F0, run by run :: sha256 c6c4138ee5c86423 :: prose twin: two significant figures, direction in words
\newcommand{\MlstMlpCmpBbFoneOverFzeroPrMax}{down to 0.38 times the errors}  % experiments/FIGS/data/v2/index.json :: nonlinear probe bb_vs_cc: (1 - AUC) of FLAV_F1 over FLAV_F0, run by run :: sha256 c6c4138ee5c86423 :: prose twin of the run farther from 1
\newcommand{\MlstMlpCmpBbFonerOverFone}{1.27 and 1.45}  % experiments/FIGS/data/v2/index.json :: nonlinear probe bb_vs_cc: (1 - AUC) of FLAV_F1R over FLAV_F1, run by run :: sha256 c6c4138ee5c86423
\newcommand{\MlstMlpCmpBbFonerOverFonePr}{at least 27\% more errors}  % experiments/FIGS/data/v2/index.json :: nonlinear probe bb_vs_cc: (1 - AUC) of FLAV_F1R over FLAV_F1, run by run :: sha256 c6c4138ee5c86423 :: prose twin: two significant figures, direction in words
\newcommand{\MlstMlpCmpBbFonerOverFonePrMax}{up to 45\% more errors}  % experiments/FIGS/data/v2/index.json :: nonlinear probe bb_vs_cc: (1 - AUC) of FLAV_F1R over FLAV_F1, run by run :: sha256 c6c4138ee5c86423 :: prose twin of the run farther from 1
\newcommand{\MlstMlpCmpBbFonerOverFzero}{0.66 and 0.55}  % experiments/FIGS/data/v2/index.json :: nonlinear probe bb_vs_cc: (1 - AUC) of FLAV_F1R over FLAV_F0, run by run :: sha256 c6c4138ee5c86423
\newcommand{\MlstMlpCmpBbFonerOverFzeroPr}{at least 34\% fewer errors}  % experiments/FIGS/data/v2/index.json :: nonlinear probe bb_vs_cc: (1 - AUC) of FLAV_F1R over FLAV_F0, run by run :: sha256 c6c4138ee5c86423 :: prose twin: two significant figures, direction in words
\newcommand{\MlstMlpCmpBbFonerOverFzeroPrMax}{up to 45\% fewer errors}  % experiments/FIGS/data/v2/index.json :: nonlinear probe bb_vs_cc: (1 - AUC) of FLAV_F1R over FLAV_F0, run by run :: sha256 c6c4138ee5c86423 :: prose twin of the run farther from 1
\newcommand{\MlstMlpCmpBbFzeroOverCoarse}{1.07 and 1.41}  % experiments/FIGS/data/v2/index.json :: nonlinear probe bb_vs_cc: (1 - AUC) of FLAV_F0 over R16_Q1, run by run :: sha256 c6c4138ee5c86423
\newcommand{\MlstMlpCmpBbFzeroOverCoarsePr}{at least 7.4\% more errors}  % experiments/FIGS/data/v2/index.json :: nonlinear probe bb_vs_cc: (1 - AUC) of FLAV_F0 over R16_Q1, run by run :: sha256 c6c4138ee5c86423 :: prose twin: two significant figures, direction in words
\newcommand{\MlstMlpCmpBbFzeroOverCoarsePrMax}{up to 41\% more errors}  % experiments/FIGS/data/v2/index.json :: nonlinear probe bb_vs_cc: (1 - AUC) of FLAV_F0 over R16_Q1, run by run :: sha256 c6c4138ee5c86423 :: prose twin of the run farther from 1
\newcommand{\MlstMlpCmpBbRandOverFull}{\pending{no run yet}}  % experiments/FIGS/data/v2/index.json :: nonlinear probe bb_vs_cc: (1 - AUC) of each random partition over the 188-class run of its seed :: sha256 c6c4138ee5c86423
\newcommand{\MlstMlpCmpBbRandOverFullPr}{\pending{no run yet}}  % experiments/FIGS/data/v2/index.json :: nonlinear probe bb_vs_cc: (1 - AUC) of each random partition over the 188-class run of its seed :: sha256 c6c4138ee5c86423 :: prose twin: two significant figures, direction in words
\newcommand{\MlstMlpCmpBbqqResOverFull}{1.011\,\ensuremath{\pm}\,0.030}  % experiments/FIGS/data/v2/index.json :: nonlinear probe bbqq_vs_ccqq: (1 - AUC) of L162 over L188, run by run :: sha256 c6c4138ee5c86423
\newcommand{\MlstMlpCmpBbqqResOverFullPr}{equal within 3\%}  % experiments/FIGS/data/v2/index.json :: nonlinear probe bbqq_vs_ccqq: (1 - AUC) of L162 over L188, run by run :: sha256 c6c4138ee5c86423 :: prose twin: two significant figures, direction in words
\newcommand{\MlstMlpCmpBbqqMidOverFull}{1.112\,\ensuremath{\pm}\,0.006}  % experiments/FIGS/data/v2/index.json :: nonlinear probe bbqq_vs_ccqq: (1 - AUC) of R42_Q1 over L188, run by run :: sha256 c6c4138ee5c86423
\newcommand{\MlstMlpCmpBbqqMidOverFullPr}{11\% more errors}  % experiments/FIGS/data/v2/index.json :: nonlinear probe bbqq_vs_ccqq: (1 - AUC) of R42_Q1 over L188, run by run :: sha256 c6c4138ee5c86423 :: prose twin: two significant figures, direction in words
\newcommand{\MlstMlpCmpBbqqCoarseOverFull}{2.93\,\ensuremath{\pm}\,0.17}  % experiments/FIGS/data/v2/index.json :: nonlinear probe bbqq_vs_ccqq: (1 - AUC) of R16_Q1 over L188, run by run :: sha256 c6c4138ee5c86423
\newcommand{\MlstMlpCmpBbqqCoarseOverFullPr}{2.9 times the errors}  % experiments/FIGS/data/v2/index.json :: nonlinear probe bbqq_vs_ccqq: (1 - AUC) of R16_Q1 over L188, run by run :: sha256 c6c4138ee5c86423 :: prose twin: two significant figures, direction in words
\newcommand{\MlstMlpCmpBbqqResMassOverRes}{\pending{no run yet}}  % experiments/FIGS/data/v2/index.json :: nonlinear probe bbqq_vs_ccqq: (1 - AUC) of L162_MASS over L162, run by run :: sha256 c6c4138ee5c86423
\newcommand{\MlstMlpCmpBbqqResMassOverResPr}{\pending{no run yet}}  % experiments/FIGS/data/v2/index.json :: nonlinear probe bbqq_vs_ccqq: (1 - AUC) of L162_MASS over L162, run by run :: sha256 c6c4138ee5c86423 :: prose twin: two significant figures, direction in words
\newcommand{\MlstMlpCmpBbqqResMassOverResPrMax}{\pending{no run yet}}  % experiments/FIGS/data/v2/index.json :: nonlinear probe bbqq_vs_ccqq: (1 - AUC) of L162_MASS over L162, run by run :: sha256 c6c4138ee5c86423 :: prose twin of the run farther from 1
\newcommand{\MlstMlpCmpBbqqCoarseMassOverCoarse}{\pending{no run yet}}  % experiments/FIGS/data/v2/index.json :: nonlinear probe bbqq_vs_ccqq: (1 - AUC) of R16_Q1_MASS over R16_Q1, run by run :: sha256 c6c4138ee5c86423
\newcommand{\MlstMlpCmpBbqqCoarseMassOverCoarsePr}{\pending{no run yet}}  % experiments/FIGS/data/v2/index.json :: nonlinear probe bbqq_vs_ccqq: (1 - AUC) of R16_Q1_MASS over R16_Q1, run by run :: sha256 c6c4138ee5c86423 :: prose twin: two significant figures, direction in words
\newcommand{\MlstMlpCmpBbqqCoarseMassOverCoarsePrMax}{\pending{no run yet}}  % experiments/FIGS/data/v2/index.json :: nonlinear probe bbqq_vs_ccqq: (1 - AUC) of R16_Q1_MASS over R16_Q1, run by run :: sha256 c6c4138ee5c86423 :: prose twin of the run farther from 1
\newcommand{\MlstMlpCmpBbqqFzeroOverFull}{3.20 and 3.78}  % experiments/FIGS/data/v2/index.json :: nonlinear probe bbqq_vs_ccqq: (1 - AUC) of FLAV_F0 over L188, run by run :: sha256 c6c4138ee5c86423
\newcommand{\MlstMlpCmpBbqqFzeroOverFullPr}{at least 3.2 times the errors}  % experiments/FIGS/data/v2/index.json :: nonlinear probe bbqq_vs_ccqq: (1 - AUC) of FLAV_F0 over L188, run by run :: sha256 c6c4138ee5c86423 :: prose twin: two significant figures, direction in words
\newcommand{\MlstMlpCmpBbqqFzeroOverFullPrMax}{up to 3.8 times the errors}  % experiments/FIGS/data/v2/index.json :: nonlinear probe bbqq_vs_ccqq: (1 - AUC) of FLAV_F0 over L188, run by run :: sha256 c6c4138ee5c86423 :: prose twin of the run farther from 1
\newcommand{\MlstMlpCmpBbqqFoneOverFzero}{0.450 and 0.363}  % experiments/FIGS/data/v2/index.json :: nonlinear probe bbqq_vs_ccqq: (1 - AUC) of FLAV_F1 over FLAV_F0, run by run :: sha256 c6c4138ee5c86423
\newcommand{\MlstMlpCmpBbqqFoneOverFzeroPr}{at most 0.45 times the errors}  % experiments/FIGS/data/v2/index.json :: nonlinear probe bbqq_vs_ccqq: (1 - AUC) of FLAV_F1 over FLAV_F0, run by run :: sha256 c6c4138ee5c86423 :: prose twin: two significant figures, direction in words
\newcommand{\MlstMlpCmpBbqqFoneOverFzeroPrMax}{down to 0.36 times the errors}  % experiments/FIGS/data/v2/index.json :: nonlinear probe bbqq_vs_ccqq: (1 - AUC) of FLAV_F1 over FLAV_F0, run by run :: sha256 c6c4138ee5c86423 :: prose twin of the run farther from 1
\newcommand{\MlstMlpCmpBbqqFonerOverFone}{1.311 and 1.343}  % experiments/FIGS/data/v2/index.json :: nonlinear probe bbqq_vs_ccqq: (1 - AUC) of FLAV_F1R over FLAV_F1, run by run :: sha256 c6c4138ee5c86423
\newcommand{\MlstMlpCmpBbqqFonerOverFonePr}{at least 31\% more errors}  % experiments/FIGS/data/v2/index.json :: nonlinear probe bbqq_vs_ccqq: (1 - AUC) of FLAV_F1R over FLAV_F1, run by run :: sha256 c6c4138ee5c86423 :: prose twin: two significant figures, direction in words
\newcommand{\MlstMlpCmpBbqqFonerOverFonePrMax}{up to 34\% more errors}  % experiments/FIGS/data/v2/index.json :: nonlinear probe bbqq_vs_ccqq: (1 - AUC) of FLAV_F1R over FLAV_F1, run by run :: sha256 c6c4138ee5c86423 :: prose twin of the run farther from 1
\newcommand{\MlstMlpCmpBbqqFonerOverFzero}{0.59 and 0.49}  % experiments/FIGS/data/v2/index.json :: nonlinear probe bbqq_vs_ccqq: (1 - AUC) of FLAV_F1R over FLAV_F0, run by run :: sha256 c6c4138ee5c86423
\newcommand{\MlstMlpCmpBbqqFonerOverFzeroPr}{at least 41\% fewer errors}  % experiments/FIGS/data/v2/index.json :: nonlinear probe bbqq_vs_ccqq: (1 - AUC) of FLAV_F1R over FLAV_F0, run by run :: sha256 c6c4138ee5c86423 :: prose twin: two significant figures, direction in words
\newcommand{\MlstMlpCmpBbqqFonerOverFzeroPrMax}{down to 0.49 times the errors}  % experiments/FIGS/data/v2/index.json :: nonlinear probe bbqq_vs_ccqq: (1 - AUC) of FLAV_F1R over FLAV_F0, run by run :: sha256 c6c4138ee5c86423 :: prose twin of the run farther from 1
\newcommand{\MlstMlpCmpBbqqFzeroOverCoarse}{1.170 and 1.236}  % experiments/FIGS/data/v2/index.json :: nonlinear probe bbqq_vs_ccqq: (1 - AUC) of FLAV_F0 over R16_Q1, run by run :: sha256 c6c4138ee5c86423
\newcommand{\MlstMlpCmpBbqqFzeroOverCoarsePr}{at least 17\% more errors}  % experiments/FIGS/data/v2/index.json :: nonlinear probe bbqq_vs_ccqq: (1 - AUC) of FLAV_F0 over R16_Q1, run by run :: sha256 c6c4138ee5c86423 :: prose twin: two significant figures, direction in words
\newcommand{\MlstMlpCmpBbqqFzeroOverCoarsePrMax}{up to 24\% more errors}  % experiments/FIGS/data/v2/index.json :: nonlinear probe bbqq_vs_ccqq: (1 - AUC) of FLAV_F0 over R16_Q1, run by run :: sha256 c6c4138ee5c86423 :: prose twin of the run farther from 1
\newcommand{\MlstMlpCmpBbqqRandOverFull}{\pending{no run yet}}  % experiments/FIGS/data/v2/index.json :: nonlinear probe bbqq_vs_ccqq: (1 - AUC) of each random partition over the 188-class run of its seed :: sha256 c6c4138ee5c86423
\newcommand{\MlstMlpCmpBbqqRandOverFullPr}{\pending{no run yet}}  % experiments/FIGS/data/v2/index.json :: nonlinear probe bbqq_vs_ccqq: (1 - AUC) of each random partition over the 188-class run of its seed :: sha256 c6c4138ee5c86423 :: prose twin: two significant figures, direction in words
\newcommand{\MlstMlpCmpQcdBbResOverFull}{1.085\,\ensuremath{\pm}\,0.030}  % experiments/FIGS/data/v2/index.json :: nonlinear probe qcd_bb_vs_cc: (1 - AUC) of L162 over L188, run by run :: sha256 c6c4138ee5c86423
\newcommand{\MlstMlpCmpQcdBbResOverFullPr}{8.5\% more errors}  % experiments/FIGS/data/v2/index.json :: nonlinear probe qcd_bb_vs_cc: (1 - AUC) of L162 over L188, run by run :: sha256 c6c4138ee5c86423 :: prose twin: two significant figures, direction in words
\newcommand{\MlstMlpCmpQcdBbMidOverFull}{1.121\,\ensuremath{\pm}\,0.017}  % experiments/FIGS/data/v2/index.json :: nonlinear probe qcd_bb_vs_cc: (1 - AUC) of R42_Q1 over L188, run by run :: sha256 c6c4138ee5c86423
\newcommand{\MlstMlpCmpQcdBbMidOverFullPr}{12\% more errors}  % experiments/FIGS/data/v2/index.json :: nonlinear probe qcd_bb_vs_cc: (1 - AUC) of R42_Q1 over L188, run by run :: sha256 c6c4138ee5c86423 :: prose twin: two significant figures, direction in words
\newcommand{\MlstMlpCmpQcdBbCoarseOverFull}{2.31\,\ensuremath{\pm}\,0.04}  % experiments/FIGS/data/v2/index.json :: nonlinear probe qcd_bb_vs_cc: (1 - AUC) of R16_Q1 over L188, run by run :: sha256 c6c4138ee5c86423
\newcommand{\MlstMlpCmpQcdBbCoarseOverFullPr}{2.3 times the errors}  % experiments/FIGS/data/v2/index.json :: nonlinear probe qcd_bb_vs_cc: (1 - AUC) of R16_Q1 over L188, run by run :: sha256 c6c4138ee5c86423 :: prose twin: two significant figures, direction in words
\newcommand{\MlstMlpCmpQcdBbResMassOverRes}{\pending{no run yet}}  % experiments/FIGS/data/v2/index.json :: nonlinear probe qcd_bb_vs_cc: (1 - AUC) of L162_MASS over L162, run by run :: sha256 c6c4138ee5c86423
\newcommand{\MlstMlpCmpQcdBbResMassOverResPr}{\pending{no run yet}}  % experiments/FIGS/data/v2/index.json :: nonlinear probe qcd_bb_vs_cc: (1 - AUC) of L162_MASS over L162, run by run :: sha256 c6c4138ee5c86423 :: prose twin: two significant figures, direction in words
\newcommand{\MlstMlpCmpQcdBbResMassOverResPrMax}{\pending{no run yet}}  % experiments/FIGS/data/v2/index.json :: nonlinear probe qcd_bb_vs_cc: (1 - AUC) of L162_MASS over L162, run by run :: sha256 c6c4138ee5c86423 :: prose twin of the run farther from 1
\newcommand{\MlstMlpCmpQcdBbCoarseMassOverCoarse}{\pending{no run yet}}  % experiments/FIGS/data/v2/index.json :: nonlinear probe qcd_bb_vs_cc: (1 - AUC) of R16_Q1_MASS over R16_Q1, run by run :: sha256 c6c4138ee5c86423
\newcommand{\MlstMlpCmpQcdBbCoarseMassOverCoarsePr}{\pending{no run yet}}  % experiments/FIGS/data/v2/index.json :: nonlinear probe qcd_bb_vs_cc: (1 - AUC) of R16_Q1_MASS over R16_Q1, run by run :: sha256 c6c4138ee5c86423 :: prose twin: two significant figures, direction in words
\newcommand{\MlstMlpCmpQcdBbCoarseMassOverCoarsePrMax}{\pending{no run yet}}  % experiments/FIGS/data/v2/index.json :: nonlinear probe qcd_bb_vs_cc: (1 - AUC) of R16_Q1_MASS over R16_Q1, run by run :: sha256 c6c4138ee5c86423 :: prose twin of the run farther from 1
\newcommand{\MlstMlpCmpQcdBbFzeroOverFull}{2.50 and 2.72}  % experiments/FIGS/data/v2/index.json :: nonlinear probe qcd_bb_vs_cc: (1 - AUC) of FLAV_F0 over L188, run by run :: sha256 c6c4138ee5c86423
\newcommand{\MlstMlpCmpQcdBbFzeroOverFullPr}{at least 2.5 times the errors}  % experiments/FIGS/data/v2/index.json :: nonlinear probe qcd_bb_vs_cc: (1 - AUC) of FLAV_F0 over L188, run by run :: sha256 c6c4138ee5c86423 :: prose twin: two significant figures, direction in words
\newcommand{\MlstMlpCmpQcdBbFzeroOverFullPrMax}{up to 2.7 times the errors}  % experiments/FIGS/data/v2/index.json :: nonlinear probe qcd_bb_vs_cc: (1 - AUC) of FLAV_F0 over L188, run by run :: sha256 c6c4138ee5c86423 :: prose twin of the run farther from 1
\newcommand{\MlstMlpCmpQcdBbFoneOverFzero}{0.573 and 0.505}  % experiments/FIGS/data/v2/index.json :: nonlinear probe qcd_bb_vs_cc: (1 - AUC) of FLAV_F1 over FLAV_F0, run by run :: sha256 c6c4138ee5c86423
\newcommand{\MlstMlpCmpQcdBbFoneOverFzeroPr}{at least 43\% fewer errors}  % experiments/FIGS/data/v2/index.json :: nonlinear probe qcd_bb_vs_cc: (1 - AUC) of FLAV_F1 over FLAV_F0, run by run :: sha256 c6c4138ee5c86423 :: prose twin: two significant figures, direction in words
\newcommand{\MlstMlpCmpQcdBbFoneOverFzeroPrMax}{up to 49\% fewer errors}  % experiments/FIGS/data/v2/index.json :: nonlinear probe qcd_bb_vs_cc: (1 - AUC) of FLAV_F1 over FLAV_F0, run by run :: sha256 c6c4138ee5c86423 :: prose twin of the run farther from 1
\newcommand{\MlstMlpCmpQcdBbFonerOverFone}{1.219 and 1.258}  % experiments/FIGS/data/v2/index.json :: nonlinear probe qcd_bb_vs_cc: (1 - AUC) of FLAV_F1R over FLAV_F1, run by run :: sha256 c6c4138ee5c86423
\newcommand{\MlstMlpCmpQcdBbFonerOverFonePr}{at least 22\% more errors}  % experiments/FIGS/data/v2/index.json :: nonlinear probe qcd_bb_vs_cc: (1 - AUC) of FLAV_F1R over FLAV_F1, run by run :: sha256 c6c4138ee5c86423 :: prose twin: two significant figures, direction in words
\newcommand{\MlstMlpCmpQcdBbFonerOverFonePrMax}{up to 26\% more errors}  % experiments/FIGS/data/v2/index.json :: nonlinear probe qcd_bb_vs_cc: (1 - AUC) of FLAV_F1R over FLAV_F1, run by run :: sha256 c6c4138ee5c86423 :: prose twin of the run farther from 1
\newcommand{\MlstMlpCmpQcdBbFonerOverFzero}{0.699 and 0.635}  % experiments/FIGS/data/v2/index.json :: nonlinear probe qcd_bb_vs_cc: (1 - AUC) of FLAV_F1R over FLAV_F0, run by run :: sha256 c6c4138ee5c86423
\newcommand{\MlstMlpCmpQcdBbFonerOverFzeroPr}{at least 30\% fewer errors}  % experiments/FIGS/data/v2/index.json :: nonlinear probe qcd_bb_vs_cc: (1 - AUC) of FLAV_F1R over FLAV_F0, run by run :: sha256 c6c4138ee5c86423 :: prose twin: two significant figures, direction in words
\newcommand{\MlstMlpCmpQcdBbFonerOverFzeroPrMax}{up to 36\% fewer errors}  % experiments/FIGS/data/v2/index.json :: nonlinear probe qcd_bb_vs_cc: (1 - AUC) of FLAV_F1R over FLAV_F0, run by run :: sha256 c6c4138ee5c86423 :: prose twin of the run farther from 1
\newcommand{\MlstMlpCmpQcdBbFzeroOverCoarse}{1.103 and 1.167}  % experiments/FIGS/data/v2/index.json :: nonlinear probe qcd_bb_vs_cc: (1 - AUC) of FLAV_F0 over R16_Q1, run by run :: sha256 c6c4138ee5c86423
\newcommand{\MlstMlpCmpQcdBbFzeroOverCoarsePr}{at least 10\% more errors}  % experiments/FIGS/data/v2/index.json :: nonlinear probe qcd_bb_vs_cc: (1 - AUC) of FLAV_F0 over R16_Q1, run by run :: sha256 c6c4138ee5c86423 :: prose twin: two significant figures, direction in words
\newcommand{\MlstMlpCmpQcdBbFzeroOverCoarsePrMax}{up to 17\% more errors}  % experiments/FIGS/data/v2/index.json :: nonlinear probe qcd_bb_vs_cc: (1 - AUC) of FLAV_F0 over R16_Q1, run by run :: sha256 c6c4138ee5c86423 :: prose twin of the run farther from 1
\newcommand{\MlstMlpCmpQcdBbRandOverFull}{\pending{no run yet}}  % experiments/FIGS/data/v2/index.json :: nonlinear probe qcd_bb_vs_cc: (1 - AUC) of each random partition over the 188-class run of its seed :: sha256 c6c4138ee5c86423
\newcommand{\MlstMlpCmpQcdBbRandOverFullPr}{\pending{no run yet}}  % experiments/FIGS/data/v2/index.json :: nonlinear probe qcd_bb_vs_cc: (1 - AUC) of each random partition over the 188-class run of its seed :: sha256 c6c4138ee5c86423 :: prose twin: two significant figures, direction in words
\newcommand{\MlstMlpCmpQcdBResOverFull}{1.050\,\ensuremath{\pm}\,0.004}  % experiments/FIGS/data/v2/index.json :: nonlinear probe qcd_b_vs_light: (1 - AUC) of L162 over L188, run by run :: sha256 c6c4138ee5c86423
\newcommand{\MlstMlpCmpQcdBResOverFullPr}{5.0\% more errors}  % experiments/FIGS/data/v2/index.json :: nonlinear probe qcd_b_vs_light: (1 - AUC) of L162 over L188, run by run :: sha256 c6c4138ee5c86423 :: prose twin: two significant figures, direction in words
\newcommand{\MlstMlpCmpQcdBMidOverFull}{1.055\,\ensuremath{\pm}\,0.005}  % experiments/FIGS/data/v2/index.json :: nonlinear probe qcd_b_vs_light: (1 - AUC) of R42_Q1 over L188, run by run :: sha256 c6c4138ee5c86423
\newcommand{\MlstMlpCmpQcdBMidOverFullPr}{5.5\% more errors}  % experiments/FIGS/data/v2/index.json :: nonlinear probe qcd_b_vs_light: (1 - AUC) of R42_Q1 over L188, run by run :: sha256 c6c4138ee5c86423 :: prose twin: two significant figures, direction in words
\newcommand{\MlstMlpCmpQcdBCoarseOverFull}{1.337\,\ensuremath{\pm}\,0.021}  % experiments/FIGS/data/v2/index.json :: nonlinear probe qcd_b_vs_light: (1 - AUC) of R16_Q1 over L188, run by run :: sha256 c6c4138ee5c86423
\newcommand{\MlstMlpCmpQcdBCoarseOverFullPr}{34\% more errors}  % experiments/FIGS/data/v2/index.json :: nonlinear probe qcd_b_vs_light: (1 - AUC) of R16_Q1 over L188, run by run :: sha256 c6c4138ee5c86423 :: prose twin: two significant figures, direction in words
\newcommand{\MlstMlpCmpQcdBResMassOverRes}{\pending{no run yet}}  % experiments/FIGS/data/v2/index.json :: nonlinear probe qcd_b_vs_light: (1 - AUC) of L162_MASS over L162, run by run :: sha256 c6c4138ee5c86423
\newcommand{\MlstMlpCmpQcdBResMassOverResPr}{\pending{no run yet}}  % experiments/FIGS/data/v2/index.json :: nonlinear probe qcd_b_vs_light: (1 - AUC) of L162_MASS over L162, run by run :: sha256 c6c4138ee5c86423 :: prose twin: two significant figures, direction in words
\newcommand{\MlstMlpCmpQcdBResMassOverResPrMax}{\pending{no run yet}}  % experiments/FIGS/data/v2/index.json :: nonlinear probe qcd_b_vs_light: (1 - AUC) of L162_MASS over L162, run by run :: sha256 c6c4138ee5c86423 :: prose twin of the run farther from 1
\newcommand{\MlstMlpCmpQcdBCoarseMassOverCoarse}{\pending{no run yet}}  % experiments/FIGS/data/v2/index.json :: nonlinear probe qcd_b_vs_light: (1 - AUC) of R16_Q1_MASS over R16_Q1, run by run :: sha256 c6c4138ee5c86423
\newcommand{\MlstMlpCmpQcdBCoarseMassOverCoarsePr}{\pending{no run yet}}  % experiments/FIGS/data/v2/index.json :: nonlinear probe qcd_b_vs_light: (1 - AUC) of R16_Q1_MASS over R16_Q1, run by run :: sha256 c6c4138ee5c86423 :: prose twin: two significant figures, direction in words
\newcommand{\MlstMlpCmpQcdBCoarseMassOverCoarsePrMax}{\pending{no run yet}}  % experiments/FIGS/data/v2/index.json :: nonlinear probe qcd_b_vs_light: (1 - AUC) of R16_Q1_MASS over R16_Q1, run by run :: sha256 c6c4138ee5c86423 :: prose twin of the run farther from 1
\newcommand{\MlstMlpCmpQcdBFzeroOverFull}{1.315 and 1.288}  % experiments/FIGS/data/v2/index.json :: nonlinear probe qcd_b_vs_light: (1 - AUC) of FLAV_F0 over L188, run by run :: sha256 c6c4138ee5c86423
\newcommand{\MlstMlpCmpQcdBFzeroOverFullPr}{at least 29\% more errors}  % experiments/FIGS/data/v2/index.json :: nonlinear probe qcd_b_vs_light: (1 - AUC) of FLAV_F0 over L188, run by run :: sha256 c6c4138ee5c86423 :: prose twin: two significant figures, direction in words
\newcommand{\MlstMlpCmpQcdBFzeroOverFullPrMax}{up to 32\% more errors}  % experiments/FIGS/data/v2/index.json :: nonlinear probe qcd_b_vs_light: (1 - AUC) of FLAV_F0 over L188, run by run :: sha256 c6c4138ee5c86423 :: prose twin of the run farther from 1
\newcommand{\MlstMlpCmpQcdBFoneOverFzero}{0.8804 and 0.8849}  % experiments/FIGS/data/v2/index.json :: nonlinear probe qcd_b_vs_light: (1 - AUC) of FLAV_F1 over FLAV_F0, run by run :: sha256 c6c4138ee5c86423
\newcommand{\MlstMlpCmpQcdBFoneOverFzeroPr}{at least 12\% fewer errors}  % experiments/FIGS/data/v2/index.json :: nonlinear probe qcd_b_vs_light: (1 - AUC) of FLAV_F1 over FLAV_F0, run by run :: sha256 c6c4138ee5c86423 :: prose twin: two significant figures, direction in words
\newcommand{\MlstMlpCmpQcdBFoneOverFzeroPrMax}{up to 12\% fewer errors}  % experiments/FIGS/data/v2/index.json :: nonlinear probe qcd_b_vs_light: (1 - AUC) of FLAV_F1 over FLAV_F0, run by run :: sha256 c6c4138ee5c86423 :: prose twin of the run farther from 1
\newcommand{\MlstMlpCmpQcdBFonerOverFone}{1.037 and 1.042}  % experiments/FIGS/data/v2/index.json :: nonlinear probe qcd_b_vs_light: (1 - AUC) of FLAV_F1R over FLAV_F1, run by run :: sha256 c6c4138ee5c86423
\newcommand{\MlstMlpCmpQcdBFonerOverFonePr}{at least 3.7\% more errors}  % experiments/FIGS/data/v2/index.json :: nonlinear probe qcd_b_vs_light: (1 - AUC) of FLAV_F1R over FLAV_F1, run by run :: sha256 c6c4138ee5c86423 :: prose twin: two significant figures, direction in words
\newcommand{\MlstMlpCmpQcdBFonerOverFonePrMax}{up to 4.2\% more errors}  % experiments/FIGS/data/v2/index.json :: nonlinear probe qcd_b_vs_light: (1 - AUC) of FLAV_F1R over FLAV_F1, run by run :: sha256 c6c4138ee5c86423 :: prose twin of the run farther from 1
\newcommand{\MlstMlpCmpQcdBFonerOverFzero}{0.9130 and 0.9217}  % experiments/FIGS/data/v2/index.json :: nonlinear probe qcd_b_vs_light: (1 - AUC) of FLAV_F1R over FLAV_F0, run by run :: sha256 c6c4138ee5c86423
\newcommand{\MlstMlpCmpQcdBFonerOverFzeroPr}{at least 7.8\% fewer errors}  % experiments/FIGS/data/v2/index.json :: nonlinear probe qcd_b_vs_light: (1 - AUC) of FLAV_F1R over FLAV_F0, run by run :: sha256 c6c4138ee5c86423 :: prose twin: two significant figures, direction in words
\newcommand{\MlstMlpCmpQcdBFonerOverFzeroPrMax}{up to 8.7\% fewer errors}  % experiments/FIGS/data/v2/index.json :: nonlinear probe qcd_b_vs_light: (1 - AUC) of FLAV_F1R over FLAV_F0, run by run :: sha256 c6c4138ee5c86423 :: prose twin of the run farther from 1
\newcommand{\MlstMlpCmpQcdBFzeroOverCoarse}{0.968 and 0.978}  % experiments/FIGS/data/v2/index.json :: nonlinear probe qcd_b_vs_light: (1 - AUC) of FLAV_F0 over R16_Q1, run by run :: sha256 c6c4138ee5c86423
\newcommand{\MlstMlpCmpQcdBFzeroOverCoarsePr}{at least 2.2\% fewer errors}  % experiments/FIGS/data/v2/index.json :: nonlinear probe qcd_b_vs_light: (1 - AUC) of FLAV_F0 over R16_Q1, run by run :: sha256 c6c4138ee5c86423 :: prose twin: two significant figures, direction in words
\newcommand{\MlstMlpCmpQcdBFzeroOverCoarsePrMax}{up to 3.2\% fewer errors}  % experiments/FIGS/data/v2/index.json :: nonlinear probe qcd_b_vs_light: (1 - AUC) of FLAV_F0 over R16_Q1, run by run :: sha256 c6c4138ee5c86423 :: prose twin of the run farther from 1
\newcommand{\MlstMlpCmpQcdBRandOverFull}{\pending{no run yet}}  % experiments/FIGS/data/v2/index.json :: nonlinear probe qcd_b_vs_light: (1 - AUC) of each random partition over the 188-class run of its seed :: sha256 c6c4138ee5c86423
\newcommand{\MlstMlpCmpQcdBRandOverFullPr}{\pending{no run yet}}  % experiments/FIGS/data/v2/index.json :: nonlinear probe qcd_b_vs_light: (1 - AUC) of each random partition over the 188-class run of its seed :: sha256 c6c4138ee5c86423 :: prose twin: two significant figures, direction in words
\newcommand{\MlstMlpCmpQcdCResOverFull}{1.039\,\ensuremath{\pm}\,0.005}  % experiments/FIGS/data/v2/index.json :: nonlinear probe qcd_c_vs_light: (1 - AUC) of L162 over L188, run by run :: sha256 c6c4138ee5c86423
\newcommand{\MlstMlpCmpQcdCResOverFullPr}{3.9\% more errors}  % experiments/FIGS/data/v2/index.json :: nonlinear probe qcd_c_vs_light: (1 - AUC) of L162 over L188, run by run :: sha256 c6c4138ee5c86423 :: prose twin: two significant figures, direction in words
\newcommand{\MlstMlpCmpQcdCMidOverFull}{1.055\,\ensuremath{\pm}\,0.006}  % experiments/FIGS/data/v2/index.json :: nonlinear probe qcd_c_vs_light: (1 - AUC) of R42_Q1 over L188, run by run :: sha256 c6c4138ee5c86423
\newcommand{\MlstMlpCmpQcdCMidOverFullPr}{5.5\% more errors}  % experiments/FIGS/data/v2/index.json :: nonlinear probe qcd_c_vs_light: (1 - AUC) of R42_Q1 over L188, run by run :: sha256 c6c4138ee5c86423 :: prose twin: two significant figures, direction in words
\newcommand{\MlstMlpCmpQcdCCoarseOverFull}{1.274\,\ensuremath{\pm}\,0.008}  % experiments/FIGS/data/v2/index.json :: nonlinear probe qcd_c_vs_light: (1 - AUC) of R16_Q1 over L188, run by run :: sha256 c6c4138ee5c86423
\newcommand{\MlstMlpCmpQcdCCoarseOverFullPr}{27\% more errors}  % experiments/FIGS/data/v2/index.json :: nonlinear probe qcd_c_vs_light: (1 - AUC) of R16_Q1 over L188, run by run :: sha256 c6c4138ee5c86423 :: prose twin: two significant figures, direction in words
\newcommand{\MlstMlpCmpQcdCResMassOverRes}{\pending{no run yet}}  % experiments/FIGS/data/v2/index.json :: nonlinear probe qcd_c_vs_light: (1 - AUC) of L162_MASS over L162, run by run :: sha256 c6c4138ee5c86423
\newcommand{\MlstMlpCmpQcdCResMassOverResPr}{\pending{no run yet}}  % experiments/FIGS/data/v2/index.json :: nonlinear probe qcd_c_vs_light: (1 - AUC) of L162_MASS over L162, run by run :: sha256 c6c4138ee5c86423 :: prose twin: two significant figures, direction in words
\newcommand{\MlstMlpCmpQcdCResMassOverResPrMax}{\pending{no run yet}}  % experiments/FIGS/data/v2/index.json :: nonlinear probe qcd_c_vs_light: (1 - AUC) of L162_MASS over L162, run by run :: sha256 c6c4138ee5c86423 :: prose twin of the run farther from 1
\newcommand{\MlstMlpCmpQcdCCoarseMassOverCoarse}{\pending{no run yet}}  % experiments/FIGS/data/v2/index.json :: nonlinear probe qcd_c_vs_light: (1 - AUC) of R16_Q1_MASS over R16_Q1, run by run :: sha256 c6c4138ee5c86423
\newcommand{\MlstMlpCmpQcdCCoarseMassOverCoarsePr}{\pending{no run yet}}  % experiments/FIGS/data/v2/index.json :: nonlinear probe qcd_c_vs_light: (1 - AUC) of R16_Q1_MASS over R16_Q1, run by run :: sha256 c6c4138ee5c86423 :: prose twin: two significant figures, direction in words
\newcommand{\MlstMlpCmpQcdCCoarseMassOverCoarsePrMax}{\pending{no run yet}}  % experiments/FIGS/data/v2/index.json :: nonlinear probe qcd_c_vs_light: (1 - AUC) of R16_Q1_MASS over R16_Q1, run by run :: sha256 c6c4138ee5c86423 :: prose twin of the run farther from 1
\newcommand{\MlstMlpCmpQcdCFzeroOverFull}{1.234 and 1.211}  % experiments/FIGS/data/v2/index.json :: nonlinear probe qcd_c_vs_light: (1 - AUC) of FLAV_F0 over L188, run by run :: sha256 c6c4138ee5c86423
\newcommand{\MlstMlpCmpQcdCFzeroOverFullPr}{at least 21\% more errors}  % experiments/FIGS/data/v2/index.json :: nonlinear probe qcd_c_vs_light: (1 - AUC) of FLAV_F0 over L188, run by run :: sha256 c6c4138ee5c86423 :: prose twin: two significant figures, direction in words
\newcommand{\MlstMlpCmpQcdCFzeroOverFullPrMax}{up to 23\% more errors}  % experiments/FIGS/data/v2/index.json :: nonlinear probe qcd_c_vs_light: (1 - AUC) of FLAV_F0 over L188, run by run :: sha256 c6c4138ee5c86423 :: prose twin of the run farther from 1
\newcommand{\MlstMlpCmpQcdCFoneOverFzero}{0.9395 and 0.9472}  % experiments/FIGS/data/v2/index.json :: nonlinear probe qcd_c_vs_light: (1 - AUC) of FLAV_F1 over FLAV_F0, run by run :: sha256 c6c4138ee5c86423
\newcommand{\MlstMlpCmpQcdCFoneOverFzeroPr}{at least 5.3\% fewer errors}  % experiments/FIGS/data/v2/index.json :: nonlinear probe qcd_c_vs_light: (1 - AUC) of FLAV_F1 over FLAV_F0, run by run :: sha256 c6c4138ee5c86423 :: prose twin: two significant figures, direction in words
\newcommand{\MlstMlpCmpQcdCFoneOverFzeroPrMax}{up to 6.0\% fewer errors}  % experiments/FIGS/data/v2/index.json :: nonlinear probe qcd_c_vs_light: (1 - AUC) of FLAV_F1 over FLAV_F0, run by run :: sha256 c6c4138ee5c86423 :: prose twin of the run farther from 1
\newcommand{\MlstMlpCmpQcdCFonerOverFone}{0.9846 and 0.9767}  % experiments/FIGS/data/v2/index.json :: nonlinear probe qcd_c_vs_light: (1 - AUC) of FLAV_F1R over FLAV_F1, run by run :: sha256 c6c4138ee5c86423
\newcommand{\MlstMlpCmpQcdCFonerOverFonePr}{at least 1.5\% fewer errors}  % experiments/FIGS/data/v2/index.json :: nonlinear probe qcd_c_vs_light: (1 - AUC) of FLAV_F1R over FLAV_F1, run by run :: sha256 c6c4138ee5c86423 :: prose twin: two significant figures, direction in words
\newcommand{\MlstMlpCmpQcdCFonerOverFonePrMax}{up to 2.3\% fewer errors}  % experiments/FIGS/data/v2/index.json :: nonlinear probe qcd_c_vs_light: (1 - AUC) of FLAV_F1R over FLAV_F1, run by run :: sha256 c6c4138ee5c86423 :: prose twin of the run farther from 1
\newcommand{\MlstMlpCmpQcdCFonerOverFzero}{0.9251 and 0.9251}  % experiments/FIGS/data/v2/index.json :: nonlinear probe qcd_c_vs_light: (1 - AUC) of FLAV_F1R over FLAV_F0, run by run :: sha256 c6c4138ee5c86423
\newcommand{\MlstMlpCmpQcdCFonerOverFzeroPr}{at least 7.5\% fewer errors}  % experiments/FIGS/data/v2/index.json :: nonlinear probe qcd_c_vs_light: (1 - AUC) of FLAV_F1R over FLAV_F0, run by run :: sha256 c6c4138ee5c86423 :: prose twin: two significant figures, direction in words
\newcommand{\MlstMlpCmpQcdCFonerOverFzeroPrMax}{up to 7.5\% fewer errors}  % experiments/FIGS/data/v2/index.json :: nonlinear probe qcd_c_vs_light: (1 - AUC) of FLAV_F1R over FLAV_F0, run by run :: sha256 c6c4138ee5c86423 :: prose twin of the run farther from 1
\newcommand{\MlstMlpCmpQcdCFzeroOverCoarse}{0.968 and 0.956}  % experiments/FIGS/data/v2/index.json :: nonlinear probe qcd_c_vs_light: (1 - AUC) of FLAV_F0 over R16_Q1, run by run :: sha256 c6c4138ee5c86423
\newcommand{\MlstMlpCmpQcdCFzeroOverCoarsePr}{at least 3.2\% fewer errors}  % experiments/FIGS/data/v2/index.json :: nonlinear probe qcd_c_vs_light: (1 - AUC) of FLAV_F0 over R16_Q1, run by run :: sha256 c6c4138ee5c86423 :: prose twin: two significant figures, direction in words
\newcommand{\MlstMlpCmpQcdCFzeroOverCoarsePrMax}{up to 4.4\% fewer errors}  % experiments/FIGS/data/v2/index.json :: nonlinear probe qcd_c_vs_light: (1 - AUC) of FLAV_F0 over R16_Q1, run by run :: sha256 c6c4138ee5c86423 :: prose twin of the run farther from 1
\newcommand{\MlstMlpCmpQcdCRandOverFull}{\pending{no run yet}}  % experiments/FIGS/data/v2/index.json :: nonlinear probe qcd_c_vs_light: (1 - AUC) of each random partition over the 188-class run of its seed :: sha256 c6c4138ee5c86423
\newcommand{\MlstMlpCmpQcdCRandOverFullPr}{\pending{no run yet}}  % experiments/FIGS/data/v2/index.json :: nonlinear probe qcd_c_vs_light: (1 - AUC) of each random partition over the 188-class run of its seed :: sha256 c6c4138ee5c86423 :: prose twin: two significant figures, direction in words
\newcommand{\MlstMlpTopRthirtyFullOverInitBn}{1.00\,\ensuremath{\pm}\,0.31}  % experiments/FIGS/data/v2/index.json :: nonlinear probe top, missing inputs zero: R30 of L188 over the BatchNorm-recomputed random network of the same index :: sha256 c6c4138ee5c86423
\newcommand{\MlstMlpTopRthirtyFullOverInitBnPr}{equal within 40\%}  % experiments/FIGS/data/v2/index.json :: nonlinear probe top, missing inputs zero: R30 of L188 over the BatchNorm-recomputed random network of the same index :: sha256 c6c4138ee5c86423 :: prose twin: two significant figures, direction in words
\newcommand{\MlstMlpTopRthirtyResOverInitBn}{0.94\,\ensuremath{\pm}\,0.26}  % experiments/FIGS/data/v2/index.json :: nonlinear probe top, missing inputs zero: R30 of L162 over the BatchNorm-recomputed random network of the same index :: sha256 c6c4138ee5c86423
\newcommand{\MlstMlpTopRthirtyResOverInitBnPr}{equal within 30\%}  % experiments/FIGS/data/v2/index.json :: nonlinear probe top, missing inputs zero: R30 of L162 over the BatchNorm-recomputed random network of the same index :: sha256 c6c4138ee5c86423 :: prose twin: two significant figures, direction in words
\newcommand{\MlstMlpTopRthirtyMidOverInitBn}{0.84\,\ensuremath{\pm}\,0.17}  % experiments/FIGS/data/v2/index.json :: nonlinear probe top, missing inputs zero: R30 of R42_Q1 over the BatchNorm-recomputed random network of the same index :: sha256 c6c4138ee5c86423
\newcommand{\MlstMlpTopRthirtyMidOverInitBnPr}{equal within 20\%}  % experiments/FIGS/data/v2/index.json :: nonlinear probe top, missing inputs zero: R30 of R42_Q1 over the BatchNorm-recomputed random network of the same index :: sha256 c6c4138ee5c86423 :: prose twin: two significant figures, direction in words
\newcommand{\MlstMlpTopRthirtyCoarseOverInitBn}{0.75\,\ensuremath{\pm}\,0.05}  % experiments/FIGS/data/v2/index.json :: nonlinear probe top, missing inputs zero: R30 of R16_Q1 over the BatchNorm-recomputed random network of the same index :: sha256 c6c4138ee5c86423
\newcommand{\MlstMlpTopRthirtyCoarseOverInitBnPr}{25\% less rejection}  % experiments/FIGS/data/v2/index.json :: nonlinear probe top, missing inputs zero: R30 of R16_Q1 over the BatchNorm-recomputed random network of the same index :: sha256 c6c4138ee5c86423 :: prose twin: two significant figures, direction in words
\newcommand{\MlstMlpQgRthirtyFullOverInitBn}{1.083\,\ensuremath{\pm}\,0.026}  % experiments/FIGS/data/v2/index.json :: nonlinear probe qg, missing inputs zero: R30 of L188 over the BatchNorm-recomputed random network of the same index :: sha256 c6c4138ee5c86423
\newcommand{\MlstMlpQgRthirtyFullOverInitBnPr}{8.3\% more rejection}  % experiments/FIGS/data/v2/index.json :: nonlinear probe qg, missing inputs zero: R30 of L188 over the BatchNorm-recomputed random network of the same index :: sha256 c6c4138ee5c86423 :: prose twin: two significant figures, direction in words
\newcommand{\MlstMlpQgRthirtyResOverInitBn}{1.010\,\ensuremath{\pm}\,0.016}  % experiments/FIGS/data/v2/index.json :: nonlinear probe qg, missing inputs zero: R30 of L162 over the BatchNorm-recomputed random network of the same index :: sha256 c6c4138ee5c86423
\newcommand{\MlstMlpQgRthirtyResOverInitBnPr}{equal within 2\%}  % experiments/FIGS/data/v2/index.json :: nonlinear probe qg, missing inputs zero: R30 of L162 over the BatchNorm-recomputed random network of the same index :: sha256 c6c4138ee5c86423 :: prose twin: two significant figures, direction in words
\newcommand{\MlstMlpQgRthirtyMidOverInitBn}{1.06\,\ensuremath{\pm}\,0.05}  % experiments/FIGS/data/v2/index.json :: nonlinear probe qg, missing inputs zero: R30 of R42_Q1 over the BatchNorm-recomputed random network of the same index :: sha256 c6c4138ee5c86423
\newcommand{\MlstMlpQgRthirtyMidOverInitBnPr}{equal within 7\%}  % experiments/FIGS/data/v2/index.json :: nonlinear probe qg, missing inputs zero: R30 of R42_Q1 over the BatchNorm-recomputed random network of the same index :: sha256 c6c4138ee5c86423 :: prose twin: two significant figures, direction in words
\newcommand{\MlstMlpQgRthirtyCoarseOverInitBn}{1.04\,\ensuremath{\pm}\,0.04}  % experiments/FIGS/data/v2/index.json :: nonlinear probe qg, missing inputs zero: R30 of R16_Q1 over the BatchNorm-recomputed random network of the same index :: sha256 c6c4138ee5c86423
\newcommand{\MlstMlpQgRthirtyCoarseOverInitBnPr}{equal within 5\%}  % experiments/FIGS/data/v2/index.json :: nonlinear probe qg, missing inputs zero: R30 of R16_Q1 over the BatchNorm-recomputed random network of the same index :: sha256 c6c4138ee5c86423 :: prose twin: two significant figures, direction in words
\newcommand{\MlstMlpQgHerwigRthirtyFullOverInitBn}{0.971\,\ensuremath{\pm}\,0.034}  % experiments/FIGS/data/v2/index.json :: nonlinear probe qg_herwig, missing inputs zero: R30 of L188 over the BatchNorm-recomputed random network of the same index :: sha256 c6c4138ee5c86423
\newcommand{\MlstMlpQgHerwigRthirtyFullOverInitBnPr}{equal within 4\%}  % experiments/FIGS/data/v2/index.json :: nonlinear probe qg_herwig, missing inputs zero: R30 of L188 over the BatchNorm-recomputed random network of the same index :: sha256 c6c4138ee5c86423 :: prose twin: two significant figures, direction in words
\newcommand{\MlstMlpQgHerwigRthirtyResOverInitBn}{0.94\,\ensuremath{\pm}\,0.05}  % experiments/FIGS/data/v2/index.json :: nonlinear probe qg_herwig, missing inputs zero: R30 of L162 over the BatchNorm-recomputed random network of the same index :: sha256 c6c4138ee5c86423
\newcommand{\MlstMlpQgHerwigRthirtyResOverInitBnPr}{equal within 6\%}  % experiments/FIGS/data/v2/index.json :: nonlinear probe qg_herwig, missing inputs zero: R30 of L162 over the BatchNorm-recomputed random network of the same index :: sha256 c6c4138ee5c86423 :: prose twin: two significant figures, direction in words
\newcommand{\MlstMlpQgHerwigRthirtyMidOverInitBn}{0.945\,\ensuremath{\pm}\,0.030}  % experiments/FIGS/data/v2/index.json :: nonlinear probe qg_herwig, missing inputs zero: R30 of R42_Q1 over the BatchNorm-recomputed random network of the same index :: sha256 c6c4138ee5c86423
\newcommand{\MlstMlpQgHerwigRthirtyMidOverInitBnPr}{equal within 6\%}  % experiments/FIGS/data/v2/index.json :: nonlinear probe qg_herwig, missing inputs zero: R30 of R42_Q1 over the BatchNorm-recomputed random network of the same index :: sha256 c6c4138ee5c86423 :: prose twin: two significant figures, direction in words
\newcommand{\MlstMlpQgHerwigRthirtyCoarseOverInitBn}{0.96\,\ensuremath{\pm}\,0.04}  % experiments/FIGS/data/v2/index.json :: nonlinear probe qg_herwig, missing inputs zero: R30 of R16_Q1 over the BatchNorm-recomputed random network of the same index :: sha256 c6c4138ee5c86423
\newcommand{\MlstMlpQgHerwigRthirtyCoarseOverInitBnPr}{equal within 5\%}  % experiments/FIGS/data/v2/index.json :: nonlinear probe qg_herwig, missing inputs zero: R30 of R16_Q1 over the BatchNorm-recomputed random network of the same index :: sha256 c6c4138ee5c86423 :: prose twin: two significant figures, direction in words
\newcommand{\MlstMlpNatJcTwoResMassOverRes}{1.016 and 1.009}  % experiments/FIGS/data/v2/index.json :: nonlinear probe, natural-composition JetClass-II (fits_v3_mlp2): (1 - AUC) of L162_MASS over L162, run by run :: sha256 c6c4138ee5c86423
\newcommand{\MlstMlpNatJcTwoResMassOverResPr}{equal within 2\%}  % experiments/FIGS/data/v2/index.json :: nonlinear probe, natural-composition JetClass-II (fits_v3_mlp2): (1 - AUC) of L162_MASS over L162, run by run :: sha256 c6c4138ee5c86423 :: prose twin: two significant figures, direction in words
\newcommand{\MlstMlpNatJcTwoResMassOverResPrMax}{equal within 2\%}  % experiments/FIGS/data/v2/index.json :: nonlinear probe, natural-composition JetClass-II (fits_v3_mlp2): (1 - AUC) of L162_MASS over L162, run by run :: sha256 c6c4138ee5c86423 :: prose twin of the run farther from 1
\newcommand{\MlstMlpNatJcTwoCoarseMassOverCoarse}{1.026 and 1.019}  % experiments/FIGS/data/v2/index.json :: nonlinear probe, natural-composition JetClass-II (fits_v3_mlp2): (1 - AUC) of R16_Q1_MASS over R16_Q1, run by run :: sha256 c6c4138ee5c86423
\newcommand{\MlstMlpNatJcTwoCoarseMassOverCoarsePr}{equal within 3\%}  % experiments/FIGS/data/v2/index.json :: nonlinear probe, natural-composition JetClass-II (fits_v3_mlp2): (1 - AUC) of R16_Q1_MASS over R16_Q1, run by run :: sha256 c6c4138ee5c86423 :: prose twin: two significant figures, direction in words
\newcommand{\MlstMlpNatJcTwoCoarseMassOverCoarsePrMax}{equal within 3\%}  % experiments/FIGS/data/v2/index.json :: nonlinear probe, natural-composition JetClass-II (fits_v3_mlp2): (1 - AUC) of R16_Q1_MASS over R16_Q1, run by run :: sha256 c6c4138ee5c86423 :: prose twin of the run farther from 1
\newcommand{\LpSubstitutionTop}{random draws}  % experiments/FIGS/data/v2/index.json :: linprobe/fits_v4_fill: lowest summed validation cross-entropy over every model, largest training set :: sha256 c6c4138ee5c86423 :: zero 8.8428; mean 8.99394; draws 7.00838
\newcommand{\LpSubstitutionQg}{the training mean}  % experiments/FIGS/data/v2/index.json :: linprobe/fits_v4_fill: lowest summed validation cross-entropy over every model, largest training set :: sha256 c6c4138ee5c86423 :: zero 13.6224; mean 13.1964; draws 13.3584
\newcommand{\MlstLpTopRthirtyFullOverInitBn}{4.67\,\ensuremath{\pm}\,0.23}  % experiments/FIGS/data/v2/index.json :: linear probe top, largest set, chosen substitution: R30 of L188 over the BatchNorm-recomputed random network of the same index :: sha256 c6c4138ee5c86423
\newcommand{\MlstLpTopRthirtyFullOverInitBnPr}{4.7 times the rejection}  % experiments/FIGS/data/v2/index.json :: linear probe top, largest set, chosen substitution: R30 of L188 over the BatchNorm-recomputed random network of the same index :: sha256 c6c4138ee5c86423 :: prose twin: two significant figures, direction in words
\newcommand{\MlstLpTopRthirtyFullOverInitBnZero}{1.04\,\ensuremath{\pm}\,0.05}  % experiments/FIGS/data/v2/index.json :: linear probe top, largest set, zero: R30 of L188 over the BatchNorm-recomputed random network of the same index :: sha256 c6c4138ee5c86423
\newcommand{\MlstLpTopRthirtyFullOverInitBnZeroPr}{equal within 5\%}  % experiments/FIGS/data/v2/index.json :: linear probe top, largest set, zero: R30 of L188 over the BatchNorm-recomputed random network of the same index :: sha256 c6c4138ee5c86423 :: prose twin: two significant figures, direction in words
\newcommand{\MlstLpTopRthirtyResOverInitBn}{4.08\,\ensuremath{\pm}\,0.28}  % experiments/FIGS/data/v2/index.json :: linear probe top, largest set, chosen substitution: R30 of L162 over the BatchNorm-recomputed random network of the same index :: sha256 c6c4138ee5c86423
\newcommand{\MlstLpTopRthirtyResOverInitBnPr}{4.1 times the rejection}  % experiments/FIGS/data/v2/index.json :: linear probe top, largest set, chosen substitution: R30 of L162 over the BatchNorm-recomputed random network of the same index :: sha256 c6c4138ee5c86423 :: prose twin: two significant figures, direction in words
\newcommand{\MlstLpTopRthirtyResOverInitBnZero}{1.33\,\ensuremath{\pm}\,0.30}  % experiments/FIGS/data/v2/index.json :: linear probe top, largest set, zero: R30 of L162 over the BatchNorm-recomputed random network of the same index :: sha256 c6c4138ee5c86423
\newcommand{\MlstLpTopRthirtyResOverInitBnZeroPr}{equal within 40\%}  % experiments/FIGS/data/v2/index.json :: linear probe top, largest set, zero: R30 of L162 over the BatchNorm-recomputed random network of the same index :: sha256 c6c4138ee5c86423 :: prose twin: two significant figures, direction in words
\newcommand{\MlstLpTopRthirtyMidOverInitBn}{4.76\,\ensuremath{\pm}\,0.15}  % experiments/FIGS/data/v2/index.json :: linear probe top, largest set, chosen substitution: R30 of R42_Q1 over the BatchNorm-recomputed random network of the same index :: sha256 c6c4138ee5c86423
\newcommand{\MlstLpTopRthirtyMidOverInitBnPr}{4.8 times the rejection}  % experiments/FIGS/data/v2/index.json :: linear probe top, largest set, chosen substitution: R30 of R42_Q1 over the BatchNorm-recomputed random network of the same index :: sha256 c6c4138ee5c86423 :: prose twin: two significant figures, direction in words
\newcommand{\MlstLpTopRthirtyMidOverInitBnZero}{1.01\,\ensuremath{\pm}\,0.28}  % experiments/FIGS/data/v2/index.json :: linear probe top, largest set, zero: R30 of R42_Q1 over the BatchNorm-recomputed random network of the same index :: sha256 c6c4138ee5c86423
\newcommand{\MlstLpTopRthirtyMidOverInitBnZeroPr}{equal within 30\%}  % experiments/FIGS/data/v2/index.json :: linear probe top, largest set, zero: R30 of R42_Q1 over the BatchNorm-recomputed random network of the same index :: sha256 c6c4138ee5c86423 :: prose twin: two significant figures, direction in words
\newcommand{\MlstLpTopRthirtyCoarseOverInitBn}{5.24\,\ensuremath{\pm}\,0.06}  % experiments/FIGS/data/v2/index.json :: linear probe top, largest set, chosen substitution: R30 of R16_Q1 over the BatchNorm-recomputed random network of the same index :: sha256 c6c4138ee5c86423
\newcommand{\MlstLpTopRthirtyCoarseOverInitBnPr}{5.2 times the rejection}  % experiments/FIGS/data/v2/index.json :: linear probe top, largest set, chosen substitution: R30 of R16_Q1 over the BatchNorm-recomputed random network of the same index :: sha256 c6c4138ee5c86423 :: prose twin: two significant figures, direction in words
\newcommand{\MlstLpTopRthirtyCoarseOverInitBnZero}{0.97\,\ensuremath{\pm}\,0.16}  % experiments/FIGS/data/v2/index.json :: linear probe top, largest set, zero: R30 of R16_Q1 over the BatchNorm-recomputed random network of the same index :: sha256 c6c4138ee5c86423
\newcommand{\MlstLpTopRthirtyCoarseOverInitBnZeroPr}{equal within 20\%}  % experiments/FIGS/data/v2/index.json :: linear probe top, largest set, zero: R30 of R16_Q1 over the BatchNorm-recomputed random network of the same index :: sha256 c6c4138ee5c86423 :: prose twin: two significant figures, direction in words
\newcommand{\MlstLpTopRthirtyResMassOverInitBn}{5.71 and 6.25}  % experiments/FIGS/data/v2/index.json :: linear probe top, largest set, chosen substitution: R30 of L162_MASS over the BatchNorm-recomputed random network of the same index :: sha256 c6c4138ee5c86423
\newcommand{\MlstLpTopRthirtyResMassOverInitBnPr}{at least 5.7 times the rejection}  % experiments/FIGS/data/v2/index.json :: linear probe top, largest set, chosen substitution: R30 of L162_MASS over the BatchNorm-recomputed random network of the same index :: sha256 c6c4138ee5c86423 :: prose twin: two significant figures, direction in words
\newcommand{\MlstLpTopRthirtyResMassOverInitBnPrMax}{up to 6.3 times the rejection}  % experiments/FIGS/data/v2/index.json :: linear probe top, largest set, chosen substitution: R30 of L162_MASS over the BatchNorm-recomputed random network of the same index :: sha256 c6c4138ee5c86423 :: prose twin of the run farther from 1
\newcommand{\MlstLpTopRthirtyResMassOverInitBnZero}{1.683 and 1.678}  % experiments/FIGS/data/v2/index.json :: linear probe top, largest set, zero: R30 of L162_MASS over the BatchNorm-recomputed random network of the same index :: sha256 c6c4138ee5c86423
\newcommand{\MlstLpTopRthirtyResMassOverInitBnZeroPr}{at least 68\% more rejection}  % experiments/FIGS/data/v2/index.json :: linear probe top, largest set, zero: R30 of L162_MASS over the BatchNorm-recomputed random network of the same index :: sha256 c6c4138ee5c86423 :: prose twin: two significant figures, direction in words
\newcommand{\MlstLpTopRthirtyResMassOverInitBnZeroPrMax}{up to 68\% more rejection}  % experiments/FIGS/data/v2/index.json :: linear probe top, largest set, zero: R30 of L162_MASS over the BatchNorm-recomputed random network of the same index :: sha256 c6c4138ee5c86423 :: prose twin of the run farther from 1
\newcommand{\MlstLpTopRthirtyCoarseMassOverInitBn}{6.73 and 7.72}  % experiments/FIGS/data/v2/index.json :: linear probe top, largest set, chosen substitution: R30 of R16_Q1_MASS over the BatchNorm-recomputed random network of the same index :: sha256 c6c4138ee5c86423
\newcommand{\MlstLpTopRthirtyCoarseMassOverInitBnPr}{at least 6.7 times the rejection}  % experiments/FIGS/data/v2/index.json :: linear probe top, largest set, chosen substitution: R30 of R16_Q1_MASS over the BatchNorm-recomputed random network of the same index :: sha256 c6c4138ee5c86423 :: prose twin: two significant figures, direction in words
\newcommand{\MlstLpTopRthirtyCoarseMassOverInitBnPrMax}{up to 7.7 times the rejection}  % experiments/FIGS/data/v2/index.json :: linear probe top, largest set, chosen substitution: R30 of R16_Q1_MASS over the BatchNorm-recomputed random network of the same index :: sha256 c6c4138ee5c86423 :: prose twin of the run farther from 1
\newcommand{\MlstLpTopRthirtyCoarseMassOverInitBnZero}{2.3 and 1.2}  % experiments/FIGS/data/v2/index.json :: linear probe top, largest set, zero: R30 of R16_Q1_MASS over the BatchNorm-recomputed random network of the same index :: sha256 c6c4138ee5c86423
\newcommand{\MlstLpTopRthirtyCoarseMassOverInitBnZeroPr}{equal within 200\%}  % experiments/FIGS/data/v2/index.json :: linear probe top, largest set, zero: R30 of R16_Q1_MASS over the BatchNorm-recomputed random network of the same index :: sha256 c6c4138ee5c86423 :: prose twin: two significant figures, direction in words
\newcommand{\MlstLpTopRthirtyCoarseMassOverInitBnZeroPrMax}{equal within 200\%}  % experiments/FIGS/data/v2/index.json :: linear probe top, largest set, zero: R30 of R16_Q1_MASS over the BatchNorm-recomputed random network of the same index :: sha256 c6c4138ee5c86423 :: prose twin of the run farther from 1
\newcommand{\MlstLpTopRthirtyResOverFull}{0.874\,\ensuremath{\pm}\,0.018}  % experiments/FIGS/data/v2/index.json :: linear probe top: R30 of L162 over the 188-class run :: sha256 c6c4138ee5c86423
\newcommand{\MlstLpTopRthirtyResOverFullPr}{13\% less rejection}  % experiments/FIGS/data/v2/index.json :: linear probe top: R30 of L162 over the 188-class run :: sha256 c6c4138ee5c86423 :: prose twin: two significant figures, direction in words
\newcommand{\MlstLpTopRthirtyResOverFullZero}{1.29\,\ensuremath{\pm}\,0.31}  % experiments/FIGS/data/v2/index.json :: linear probe top: R30 of L162 over the 188-class run :: sha256 c6c4138ee5c86423
\newcommand{\MlstLpTopRthirtyResOverFullZeroPr}{equal within 40\%}  % experiments/FIGS/data/v2/index.json :: linear probe top: R30 of L162 over the 188-class run :: sha256 c6c4138ee5c86423 :: prose twin: two significant figures, direction in words
\newcommand{\MlstLpTopRthirtyMidOverFull}{1.02\,\ensuremath{\pm}\,0.08}  % experiments/FIGS/data/v2/index.json :: linear probe top: R30 of R42_Q1 over the 188-class run :: sha256 c6c4138ee5c86423
\newcommand{\MlstLpTopRthirtyMidOverFullPr}{equal within 9\%}  % experiments/FIGS/data/v2/index.json :: linear probe top: R30 of R42_Q1 over the 188-class run :: sha256 c6c4138ee5c86423 :: prose twin: two significant figures, direction in words
\newcommand{\MlstLpTopRthirtyMidOverFullZero}{0.98\,\ensuremath{\pm}\,0.28}  % experiments/FIGS/data/v2/index.json :: linear probe top: R30 of R42_Q1 over the 188-class run :: sha256 c6c4138ee5c86423
\newcommand{\MlstLpTopRthirtyMidOverFullZeroPr}{equal within 30\%}  % experiments/FIGS/data/v2/index.json :: linear probe top: R30 of R42_Q1 over the 188-class run :: sha256 c6c4138ee5c86423 :: prose twin: two significant figures, direction in words
\newcommand{\MlstLpTopRthirtyCoarseOverFull}{1.12\,\ensuremath{\pm}\,0.06}  % experiments/FIGS/data/v2/index.json :: linear probe top: R30 of R16_Q1 over the 188-class run :: sha256 c6c4138ee5c86423
\newcommand{\MlstLpTopRthirtyCoarseOverFullPr}{12\% more rejection}  % experiments/FIGS/data/v2/index.json :: linear probe top: R30 of R16_Q1 over the 188-class run :: sha256 c6c4138ee5c86423 :: prose twin: two significant figures, direction in words
\newcommand{\MlstLpTopRthirtyCoarseOverFullZero}{0.93\,\ensuremath{\pm}\,0.15}  % experiments/FIGS/data/v2/index.json :: linear probe top: R30 of R16_Q1 over the 188-class run :: sha256 c6c4138ee5c86423
\newcommand{\MlstLpTopRthirtyCoarseOverFullZeroPr}{equal within 20\%}  % experiments/FIGS/data/v2/index.json :: linear probe top: R30 of R16_Q1 over the 188-class run :: sha256 c6c4138ee5c86423 :: prose twin: two significant figures, direction in words
\newcommand{\MlstLpTopRthirtyMeanOverZeroFull}{1.25\,\ensuremath{\pm}\,0.27}  % experiments/FIGS/data/v2/index.json :: linear probe top: R30 with mean over R30 with zero, L188 :: sha256 c6c4138ee5c86423
\newcommand{\MlstLpTopRthirtyMeanOverZeroFullPr}{equal within 30\%}  % experiments/FIGS/data/v2/index.json :: linear probe top: R30 with mean over R30 with zero, L188 :: sha256 c6c4138ee5c86423 :: prose twin: two significant figures, direction in words
\newcommand{\MlstLpTopRthirtyDrawOverZeroFull}{3.9\,\ensuremath{\pm}\,0.4}  % experiments/FIGS/data/v2/index.json :: linear probe top: R30 with draws over R30 with zero, L188 :: sha256 c6c4138ee5c86423
\newcommand{\MlstLpTopRthirtyDrawOverZeroFullPr}{3.9 times the rejection}  % experiments/FIGS/data/v2/index.json :: linear probe top: R30 with draws over R30 with zero, L188 :: sha256 c6c4138ee5c86423 :: prose twin: two significant figures, direction in words
\newcommand{\MlstLpTopRthirtyMeanOverZeroInitBn}{1.08\,\ensuremath{\pm}\,0.20}  % experiments/FIGS/data/v2/index.json :: linear probe top: R30 with mean over R30 with zero, INIT_BN :: sha256 c6c4138ee5c86423
\newcommand{\MlstLpTopRthirtyMeanOverZeroInitBnPr}{equal within 30\%}  % experiments/FIGS/data/v2/index.json :: linear probe top: R30 with mean over R30 with zero, INIT_BN :: sha256 c6c4138ee5c86423 :: prose twin: two significant figures, direction in words
\newcommand{\MlstLpTopRthirtyDrawOverZeroInitBn}{0.86\,\ensuremath{\pm}\,0.14}  % experiments/FIGS/data/v2/index.json :: linear probe top: R30 with draws over R30 with zero, INIT_BN :: sha256 c6c4138ee5c86423
\newcommand{\MlstLpTopRthirtyDrawOverZeroInitBnPr}{equal within 20\%}  % experiments/FIGS/data/v2/index.json :: linear probe top: R30 with draws over R30 with zero, INIT_BN :: sha256 c6c4138ee5c86423 :: prose twin: two significant figures, direction in words
\newcommand{\MlstLpTopRthirtyKinOverFull}{3.88}  % experiments/FIGS/data/v2/index.json :: linear probe top: R30 of the kinematic-only model over the mean of the 188-class runs (zero) :: sha256 c6c4138ee5c86423
\newcommand{\MlstLpTopRthirtyKinOverFullPr}{3.9 times the rejection}  % experiments/FIGS/data/v2/index.json :: linear probe top: R30 of the kinematic-only model over the mean of the 188-class runs (zero) :: sha256 c6c4138ee5c86423 :: prose twin: two significant figures, direction in words
\newcommand{\MlstLpTopRthirtyKin}{231}  % experiments/FIGS/data/v2/index.json :: linear probe top, largest set: R30 of the kinematic-only model (one run) :: sha256 c6c4138ee5c86423
\newcommand{\MlstLpTopRthirtyFullZero}{60\,\ensuremath{\pm}\,12}  % experiments/FIGS/data/v2/index.json :: linear probe top, largest set: R30 with zero, L188 :: sha256 c6c4138ee5c86423
\newcommand{\MlstLpTopRthirtyFullMean}{73\,\ensuremath{\pm}\,15}  % experiments/FIGS/data/v2/index.json :: linear probe top, largest set: R30 with mean, L188 :: sha256 c6c4138ee5c86423
\newcommand{\MlstLpTopRthirtyFullDraws}{228\,\ensuremath{\pm}\,25}  % experiments/FIGS/data/v2/index.json :: linear probe top, largest set: R30 with draws, L188 :: sha256 c6c4138ee5c86423
\newcommand{\MlstLpTopRthirtyInitBnZero}{58\,\ensuremath{\pm}\,13}  % experiments/FIGS/data/v2/index.json :: linear probe top, largest set: R30 with zero, INIT_BN :: sha256 c6c4138ee5c86423
\newcommand{\MlstLpTopRthirtyInitBnMean}{60.4\,\ensuremath{\pm}\,2.4}  % experiments/FIGS/data/v2/index.json :: linear probe top, largest set: R30 with mean, INIT_BN :: sha256 c6c4138ee5c86423
\newcommand{\MlstLpTopRthirtyInitBnDraws}{48.6\,\ensuremath{\pm}\,3.0}  % experiments/FIGS/data/v2/index.json :: linear probe top, largest set: R30 with draws, INIT_BN :: sha256 c6c4138ee5c86423
\newcommand{\MlstLpQgRthirtyFullOverInitBn}{1.39\,\ensuremath{\pm}\,0.06}  % experiments/FIGS/data/v2/index.json :: linear probe qg, largest set, chosen substitution: R30 of L188 over the BatchNorm-recomputed random network of the same index :: sha256 c6c4138ee5c86423
\newcommand{\MlstLpQgRthirtyFullOverInitBnPr}{39\% more rejection}  % experiments/FIGS/data/v2/index.json :: linear probe qg, largest set, chosen substitution: R30 of L188 over the BatchNorm-recomputed random network of the same index :: sha256 c6c4138ee5c86423 :: prose twin: two significant figures, direction in words
\newcommand{\MlstLpQgRthirtyFullOverInitBnZero}{1.21\,\ensuremath{\pm}\,0.15}  % experiments/FIGS/data/v2/index.json :: linear probe qg, largest set, zero: R30 of L188 over the BatchNorm-recomputed random network of the same index :: sha256 c6c4138ee5c86423
\newcommand{\MlstLpQgRthirtyFullOverInitBnZeroPr}{equal within 30\%}  % experiments/FIGS/data/v2/index.json :: linear probe qg, largest set, zero: R30 of L188 over the BatchNorm-recomputed random network of the same index :: sha256 c6c4138ee5c86423 :: prose twin: two significant figures, direction in words
\newcommand{\MlstLpQgRthirtyResOverInitBn}{1.360\,\ensuremath{\pm}\,0.025}  % experiments/FIGS/data/v2/index.json :: linear probe qg, largest set, chosen substitution: R30 of L162 over the BatchNorm-recomputed random network of the same index :: sha256 c6c4138ee5c86423
\newcommand{\MlstLpQgRthirtyResOverInitBnPr}{36\% more rejection}  % experiments/FIGS/data/v2/index.json :: linear probe qg, largest set, chosen substitution: R30 of L162 over the BatchNorm-recomputed random network of the same index :: sha256 c6c4138ee5c86423 :: prose twin: two significant figures, direction in words
\newcommand{\MlstLpQgRthirtyResOverInitBnZero}{1.07\,\ensuremath{\pm}\,0.13}  % experiments/FIGS/data/v2/index.json :: linear probe qg, largest set, zero: R30 of L162 over the BatchNorm-recomputed random network of the same index :: sha256 c6c4138ee5c86423
\newcommand{\MlstLpQgRthirtyResOverInitBnZeroPr}{equal within 20\%}  % experiments/FIGS/data/v2/index.json :: linear probe qg, largest set, zero: R30 of L162 over the BatchNorm-recomputed random network of the same index :: sha256 c6c4138ee5c86423 :: prose twin: two significant figures, direction in words
\newcommand{\MlstLpQgRthirtyMidOverInitBn}{1.357\,\ensuremath{\pm}\,0.033}  % experiments/FIGS/data/v2/index.json :: linear probe qg, largest set, chosen substitution: R30 of R42_Q1 over the BatchNorm-recomputed random network of the same index :: sha256 c6c4138ee5c86423
\newcommand{\MlstLpQgRthirtyMidOverInitBnPr}{36\% more rejection}  % experiments/FIGS/data/v2/index.json :: linear probe qg, largest set, chosen substitution: R30 of R42_Q1 over the BatchNorm-recomputed random network of the same index :: sha256 c6c4138ee5c86423 :: prose twin: two significant figures, direction in words
\newcommand{\MlstLpQgRthirtyMidOverInitBnZero}{1.12\,\ensuremath{\pm}\,0.11}  % experiments/FIGS/data/v2/index.json :: linear probe qg, largest set, zero: R30 of R42_Q1 over the BatchNorm-recomputed random network of the same index :: sha256 c6c4138ee5c86423
\newcommand{\MlstLpQgRthirtyMidOverInitBnZeroPr}{equal within 20\%}  % experiments/FIGS/data/v2/index.json :: linear probe qg, largest set, zero: R30 of R42_Q1 over the BatchNorm-recomputed random network of the same index :: sha256 c6c4138ee5c86423 :: prose twin: two significant figures, direction in words
\newcommand{\MlstLpQgRthirtyCoarseOverInitBn}{1.42\,\ensuremath{\pm}\,0.04}  % experiments/FIGS/data/v2/index.json :: linear probe qg, largest set, chosen substitution: R30 of R16_Q1 over the BatchNorm-recomputed random network of the same index :: sha256 c6c4138ee5c86423
\newcommand{\MlstLpQgRthirtyCoarseOverInitBnPr}{42\% more rejection}  % experiments/FIGS/data/v2/index.json :: linear probe qg, largest set, chosen substitution: R30 of R16_Q1 over the BatchNorm-recomputed random network of the same index :: sha256 c6c4138ee5c86423 :: prose twin: two significant figures, direction in words
\newcommand{\MlstLpQgRthirtyCoarseOverInitBnZero}{1.15\,\ensuremath{\pm}\,0.17}  % experiments/FIGS/data/v2/index.json :: linear probe qg, largest set, zero: R30 of R16_Q1 over the BatchNorm-recomputed random network of the same index :: sha256 c6c4138ee5c86423
\newcommand{\MlstLpQgRthirtyCoarseOverInitBnZeroPr}{equal within 20\%}  % experiments/FIGS/data/v2/index.json :: linear probe qg, largest set, zero: R30 of R16_Q1 over the BatchNorm-recomputed random network of the same index :: sha256 c6c4138ee5c86423 :: prose twin: two significant figures, direction in words
\newcommand{\MlstLpQgRthirtyResMassOverInitBn}{1.430 and 1.364}  % experiments/FIGS/data/v2/index.json :: linear probe qg, largest set, chosen substitution: R30 of L162_MASS over the BatchNorm-recomputed random network of the same index :: sha256 c6c4138ee5c86423
\newcommand{\MlstLpQgRthirtyResMassOverInitBnPr}{at least 36\% more rejection}  % experiments/FIGS/data/v2/index.json :: linear probe qg, largest set, chosen substitution: R30 of L162_MASS over the BatchNorm-recomputed random network of the same index :: sha256 c6c4138ee5c86423 :: prose twin: two significant figures, direction in words
\newcommand{\MlstLpQgRthirtyResMassOverInitBnPrMax}{up to 43\% more rejection}  % experiments/FIGS/data/v2/index.json :: linear probe qg, largest set, chosen substitution: R30 of L162_MASS over the BatchNorm-recomputed random network of the same index :: sha256 c6c4138ee5c86423 :: prose twin of the run farther from 1
\newcommand{\MlstLpQgRthirtyResMassOverInitBnZero}{1.213 and 1.175}  % experiments/FIGS/data/v2/index.json :: linear probe qg, largest set, zero: R30 of L162_MASS over the BatchNorm-recomputed random network of the same index :: sha256 c6c4138ee5c86423
\newcommand{\MlstLpQgRthirtyResMassOverInitBnZeroPr}{equal within 30\%}  % experiments/FIGS/data/v2/index.json :: linear probe qg, largest set, zero: R30 of L162_MASS over the BatchNorm-recomputed random network of the same index :: sha256 c6c4138ee5c86423 :: prose twin: two significant figures, direction in words
\newcommand{\MlstLpQgRthirtyResMassOverInitBnZeroPrMax}{equal within 30\%}  % experiments/FIGS/data/v2/index.json :: linear probe qg, largest set, zero: R30 of L162_MASS over the BatchNorm-recomputed random network of the same index :: sha256 c6c4138ee5c86423 :: prose twin of the run farther from 1
\newcommand{\MlstLpQgRthirtyCoarseMassOverInitBn}{1.324 and 1.394}  % experiments/FIGS/data/v2/index.json :: linear probe qg, largest set, chosen substitution: R30 of R16_Q1_MASS over the BatchNorm-recomputed random network of the same index :: sha256 c6c4138ee5c86423
\newcommand{\MlstLpQgRthirtyCoarseMassOverInitBnPr}{at least 32\% more rejection}  % experiments/FIGS/data/v2/index.json :: linear probe qg, largest set, chosen substitution: R30 of R16_Q1_MASS over the BatchNorm-recomputed random network of the same index :: sha256 c6c4138ee5c86423 :: prose twin: two significant figures, direction in words
\newcommand{\MlstLpQgRthirtyCoarseMassOverInitBnPrMax}{up to 39\% more rejection}  % experiments/FIGS/data/v2/index.json :: linear probe qg, largest set, chosen substitution: R30 of R16_Q1_MASS over the BatchNorm-recomputed random network of the same index :: sha256 c6c4138ee5c86423 :: prose twin of the run farther from 1
\newcommand{\MlstLpQgRthirtyCoarseMassOverInitBnZero}{1.224 and 1.261}  % experiments/FIGS/data/v2/index.json :: linear probe qg, largest set, zero: R30 of R16_Q1_MASS over the BatchNorm-recomputed random network of the same index :: sha256 c6c4138ee5c86423
\newcommand{\MlstLpQgRthirtyCoarseMassOverInitBnZeroPr}{equal within 30\%}  % experiments/FIGS/data/v2/index.json :: linear probe qg, largest set, zero: R30 of R16_Q1_MASS over the BatchNorm-recomputed random network of the same index :: sha256 c6c4138ee5c86423 :: prose twin: two significant figures, direction in words
\newcommand{\MlstLpQgRthirtyCoarseMassOverInitBnZeroPrMax}{equal within 30\%}  % experiments/FIGS/data/v2/index.json :: linear probe qg, largest set, zero: R30 of R16_Q1_MASS over the BatchNorm-recomputed random network of the same index :: sha256 c6c4138ee5c86423 :: prose twin of the run farther from 1
\newcommand{\MlstLpQgRthirtyResOverFull}{0.977\,\ensuremath{\pm}\,0.028}  % experiments/FIGS/data/v2/index.json :: linear probe qg: R30 of L162 over the 188-class run :: sha256 c6c4138ee5c86423
\newcommand{\MlstLpQgRthirtyResOverFullPr}{equal within 3\%}  % experiments/FIGS/data/v2/index.json :: linear probe qg: R30 of L162 over the 188-class run :: sha256 c6c4138ee5c86423 :: prose twin: two significant figures, direction in words
\newcommand{\MlstLpQgRthirtyResOverFullZero}{0.89\,\ensuremath{\pm}\,0.20}  % experiments/FIGS/data/v2/index.json :: linear probe qg: R30 of L162 over the 188-class run :: sha256 c6c4138ee5c86423
\newcommand{\MlstLpQgRthirtyResOverFullZeroPr}{equal within 20\%}  % experiments/FIGS/data/v2/index.json :: linear probe qg: R30 of L162 over the 188-class run :: sha256 c6c4138ee5c86423 :: prose twin: two significant figures, direction in words
\newcommand{\MlstLpQgRthirtyMidOverFull}{0.974\,\ensuremath{\pm}\,0.031}  % experiments/FIGS/data/v2/index.json :: linear probe qg: R30 of R42_Q1 over the 188-class run :: sha256 c6c4138ee5c86423
\newcommand{\MlstLpQgRthirtyMidOverFullPr}{equal within 4\%}  % experiments/FIGS/data/v2/index.json :: linear probe qg: R30 of R42_Q1 over the 188-class run :: sha256 c6c4138ee5c86423 :: prose twin: two significant figures, direction in words
\newcommand{\MlstLpQgRthirtyMidOverFullZero}{0.94\,\ensuremath{\pm}\,0.18}  % experiments/FIGS/data/v2/index.json :: linear probe qg: R30 of R42_Q1 over the 188-class run :: sha256 c6c4138ee5c86423
\newcommand{\MlstLpQgRthirtyMidOverFullZeroPr}{equal within 20\%}  % experiments/FIGS/data/v2/index.json :: linear probe qg: R30 of R42_Q1 over the 188-class run :: sha256 c6c4138ee5c86423 :: prose twin: two significant figures, direction in words
\newcommand{\MlstLpQgRthirtyCoarseOverFull}{1.02\,\ensuremath{\pm}\,0.06}  % experiments/FIGS/data/v2/index.json :: linear probe qg: R30 of R16_Q1 over the 188-class run :: sha256 c6c4138ee5c86423
\newcommand{\MlstLpQgRthirtyCoarseOverFullPr}{equal within 6\%}  % experiments/FIGS/data/v2/index.json :: linear probe qg: R30 of R16_Q1 over the 188-class run :: sha256 c6c4138ee5c86423 :: prose twin: two significant figures, direction in words
\newcommand{\MlstLpQgRthirtyCoarseOverFullZero}{0.95\,\ensuremath{\pm}\,0.09}  % experiments/FIGS/data/v2/index.json :: linear probe qg: R30 of R16_Q1 over the 188-class run :: sha256 c6c4138ee5c86423
\newcommand{\MlstLpQgRthirtyCoarseOverFullZeroPr}{equal within 9\%}  % experiments/FIGS/data/v2/index.json :: linear probe qg: R30 of R16_Q1 over the 188-class run :: sha256 c6c4138ee5c86423 :: prose twin: two significant figures, direction in words
\newcommand{\MlstLpQgRthirtyMeanOverZeroFull}{1.12\,\ensuremath{\pm}\,0.17}  % experiments/FIGS/data/v2/index.json :: linear probe qg: R30 with mean over R30 with zero, L188 :: sha256 c6c4138ee5c86423
\newcommand{\MlstLpQgRthirtyMeanOverZeroFullPr}{equal within 20\%}  % experiments/FIGS/data/v2/index.json :: linear probe qg: R30 with mean over R30 with zero, L188 :: sha256 c6c4138ee5c86423 :: prose twin: two significant figures, direction in words
\newcommand{\MlstLpQgRthirtyDrawOverZeroFull}{1.05\,\ensuremath{\pm}\,0.12}  % experiments/FIGS/data/v2/index.json :: linear probe qg: R30 with draws over R30 with zero, L188 :: sha256 c6c4138ee5c86423
\newcommand{\MlstLpQgRthirtyDrawOverZeroFullPr}{equal within 20\%}  % experiments/FIGS/data/v2/index.json :: linear probe qg: R30 with draws over R30 with zero, L188 :: sha256 c6c4138ee5c86423 :: prose twin: two significant figures, direction in words
\newcommand{\MlstLpQgRthirtyMeanOverZeroInitBn}{0.964\,\ensuremath{\pm}\,0.030}  % experiments/FIGS/data/v2/index.json :: linear probe qg: R30 with mean over R30 with zero, INIT_BN :: sha256 c6c4138ee5c86423
\newcommand{\MlstLpQgRthirtyMeanOverZeroInitBnPr}{equal within 4\%}  % experiments/FIGS/data/v2/index.json :: linear probe qg: R30 with mean over R30 with zero, INIT_BN :: sha256 c6c4138ee5c86423 :: prose twin: two significant figures, direction in words
\newcommand{\MlstLpQgRthirtyDrawOverZeroInitBn}{0.97\,\ensuremath{\pm}\,0.04}  % experiments/FIGS/data/v2/index.json :: linear probe qg: R30 with draws over R30 with zero, INIT_BN :: sha256 c6c4138ee5c86423
\newcommand{\MlstLpQgRthirtyDrawOverZeroInitBnPr}{equal within 4\%}  % experiments/FIGS/data/v2/index.json :: linear probe qg: R30 with draws over R30 with zero, INIT_BN :: sha256 c6c4138ee5c86423 :: prose twin: two significant figures, direction in words
\newcommand{\MlstLpQgRthirtyFullZero}{66\,\ensuremath{\pm}\,7}  % experiments/FIGS/data/v2/index.json :: linear probe qg, largest set: R30 with zero, L188 :: sha256 c6c4138ee5c86423
\newcommand{\MlstLpQgRthirtyFullMean}{72.7\,\ensuremath{\pm}\,2.4}  % experiments/FIGS/data/v2/index.json :: linear probe qg, largest set: R30 with mean, L188 :: sha256 c6c4138ee5c86423
\newcommand{\MlstLpQgRthirtyFullDraws}{68.1\,\ensuremath{\pm}\,1.8}  % experiments/FIGS/data/v2/index.json :: linear probe qg, largest set: R30 with draws, L188 :: sha256 c6c4138ee5c86423
\newcommand{\MlstLpQgRthirtyInitBnZero}{54.2\,\ensuremath{\pm}\,2.4}  % experiments/FIGS/data/v2/index.json :: linear probe qg, largest set: R30 with zero, INIT_BN :: sha256 c6c4138ee5c86423
\newcommand{\MlstLpQgRthirtyInitBnMean}{52.2\,\ensuremath{\pm}\,0.8}  % experiments/FIGS/data/v2/index.json :: linear probe qg, largest set: R30 with mean, INIT_BN :: sha256 c6c4138ee5c86423
\newcommand{\MlstLpQgRthirtyInitBnDraws}{52.4\,\ensuremath{\pm}\,0.4}  % experiments/FIGS/data/v2/index.json :: linear probe qg, largest set: R30 with draws, INIT_BN :: sha256 c6c4138ee5c86423
\newcommand{\MlstLpQgHerwigRthirtyFullOverInitBn}{1.14\,\ensuremath{\pm}\,0.10}  % experiments/FIGS/data/v2/index.json :: linear probe qg_herwig, largest set, chosen substitution: R30 of L188 over the BatchNorm-recomputed random network of the same index :: sha256 c6c4138ee5c86423
\newcommand{\MlstLpQgHerwigRthirtyFullOverInitBnPr}{equal within 20\%}  % experiments/FIGS/data/v2/index.json :: linear probe qg_herwig, largest set, chosen substitution: R30 of L188 over the BatchNorm-recomputed random network of the same index :: sha256 c6c4138ee5c86423 :: prose twin: two significant figures, direction in words
\newcommand{\MlstLpQgHerwigRthirtyFullOverInitBnZero}{1.054\,\ensuremath{\pm}\,0.028}  % experiments/FIGS/data/v2/index.json :: linear probe qg_herwig, largest set, zero: R30 of L188 over the BatchNorm-recomputed random network of the same index :: sha256 c6c4138ee5c86423
\newcommand{\MlstLpQgHerwigRthirtyFullOverInitBnZeroPr}{equal within 6\%}  % experiments/FIGS/data/v2/index.json :: linear probe qg_herwig, largest set, zero: R30 of L188 over the BatchNorm-recomputed random network of the same index :: sha256 c6c4138ee5c86423 :: prose twin: two significant figures, direction in words
\newcommand{\MlstLpQgHerwigRthirtyResOverInitBn}{1.17\,\ensuremath{\pm}\,0.04}  % experiments/FIGS/data/v2/index.json :: linear probe qg_herwig, largest set, chosen substitution: R30 of L162 over the BatchNorm-recomputed random network of the same index :: sha256 c6c4138ee5c86423
\newcommand{\MlstLpQgHerwigRthirtyResOverInitBnPr}{17\% more rejection}  % experiments/FIGS/data/v2/index.json :: linear probe qg_herwig, largest set, chosen substitution: R30 of L162 over the BatchNorm-recomputed random network of the same index :: sha256 c6c4138ee5c86423 :: prose twin: two significant figures, direction in words
\newcommand{\MlstLpQgHerwigRthirtyResOverInitBnZero}{0.98\,\ensuremath{\pm}\,0.10}  % experiments/FIGS/data/v2/index.json :: linear probe qg_herwig, largest set, zero: R30 of L162 over the BatchNorm-recomputed random network of the same index :: sha256 c6c4138ee5c86423
\newcommand{\MlstLpQgHerwigRthirtyResOverInitBnZeroPr}{equal within 10\%}  % experiments/FIGS/data/v2/index.json :: linear probe qg_herwig, largest set, zero: R30 of L162 over the BatchNorm-recomputed random network of the same index :: sha256 c6c4138ee5c86423 :: prose twin: two significant figures, direction in words
\newcommand{\MlstLpQgHerwigRthirtyMidOverInitBn}{1.12\,\ensuremath{\pm}\,0.07}  % experiments/FIGS/data/v2/index.json :: linear probe qg_herwig, largest set, chosen substitution: R30 of R42_Q1 over the BatchNorm-recomputed random network of the same index :: sha256 c6c4138ee5c86423
\newcommand{\MlstLpQgHerwigRthirtyMidOverInitBnPr}{equal within 20\%}  % experiments/FIGS/data/v2/index.json :: linear probe qg_herwig, largest set, chosen substitution: R30 of R42_Q1 over the BatchNorm-recomputed random network of the same index :: sha256 c6c4138ee5c86423 :: prose twin: two significant figures, direction in words
\newcommand{\MlstLpQgHerwigRthirtyMidOverInitBnZero}{0.99\,\ensuremath{\pm}\,0.07}  % experiments/FIGS/data/v2/index.json :: linear probe qg_herwig, largest set, zero: R30 of R42_Q1 over the BatchNorm-recomputed random network of the same index :: sha256 c6c4138ee5c86423
\newcommand{\MlstLpQgHerwigRthirtyMidOverInitBnZeroPr}{equal within 7\%}  % experiments/FIGS/data/v2/index.json :: linear probe qg_herwig, largest set, zero: R30 of R42_Q1 over the BatchNorm-recomputed random network of the same index :: sha256 c6c4138ee5c86423 :: prose twin: two significant figures, direction in words
\newcommand{\MlstLpQgHerwigRthirtyCoarseOverInitBn}{1.154\,\ensuremath{\pm}\,0.015}  % experiments/FIGS/data/v2/index.json :: linear probe qg_herwig, largest set, chosen substitution: R30 of R16_Q1 over the BatchNorm-recomputed random network of the same index :: sha256 c6c4138ee5c86423
\newcommand{\MlstLpQgHerwigRthirtyCoarseOverInitBnPr}{15\% more rejection}  % experiments/FIGS/data/v2/index.json :: linear probe qg_herwig, largest set, chosen substitution: R30 of R16_Q1 over the BatchNorm-recomputed random network of the same index :: sha256 c6c4138ee5c86423 :: prose twin: two significant figures, direction in words
\newcommand{\MlstLpQgHerwigRthirtyCoarseOverInitBnZero}{1.07\,\ensuremath{\pm}\,0.05}  % experiments/FIGS/data/v2/index.json :: linear probe qg_herwig, largest set, zero: R30 of R16_Q1 over the BatchNorm-recomputed random network of the same index :: sha256 c6c4138ee5c86423
\newcommand{\MlstLpQgHerwigRthirtyCoarseOverInitBnZeroPr}{equal within 8\%}  % experiments/FIGS/data/v2/index.json :: linear probe qg_herwig, largest set, zero: R30 of R16_Q1 over the BatchNorm-recomputed random network of the same index :: sha256 c6c4138ee5c86423 :: prose twin: two significant figures, direction in words
\newcommand{\MlstLpQgHerwigRthirtyResMassOverInitBn}{1.152 and 1.138}  % experiments/FIGS/data/v2/index.json :: linear probe qg_herwig, largest set, chosen substitution: R30 of L162_MASS over the BatchNorm-recomputed random network of the same index :: sha256 c6c4138ee5c86423
\newcommand{\MlstLpQgHerwigRthirtyResMassOverInitBnPr}{equal within 20\%}  % experiments/FIGS/data/v2/index.json :: linear probe qg_herwig, largest set, chosen substitution: R30 of L162_MASS over the BatchNorm-recomputed random network of the same index :: sha256 c6c4138ee5c86423 :: prose twin: two significant figures, direction in words
\newcommand{\MlstLpQgHerwigRthirtyResMassOverInitBnPrMax}{equal within 20\%}  % experiments/FIGS/data/v2/index.json :: linear probe qg_herwig, largest set, chosen substitution: R30 of L162_MASS over the BatchNorm-recomputed random network of the same index :: sha256 c6c4138ee5c86423 :: prose twin of the run farther from 1
\newcommand{\MlstLpQgHerwigRthirtyResMassOverInitBnZero}{1.13 and 1.03}  % experiments/FIGS/data/v2/index.json :: linear probe qg_herwig, largest set, zero: R30 of L162_MASS over the BatchNorm-recomputed random network of the same index :: sha256 c6c4138ee5c86423
\newcommand{\MlstLpQgHerwigRthirtyResMassOverInitBnZeroPr}{equal within 20\%}  % experiments/FIGS/data/v2/index.json :: linear probe qg_herwig, largest set, zero: R30 of L162_MASS over the BatchNorm-recomputed random network of the same index :: sha256 c6c4138ee5c86423 :: prose twin: two significant figures, direction in words
\newcommand{\MlstLpQgHerwigRthirtyResMassOverInitBnZeroPrMax}{equal within 20\%}  % experiments/FIGS/data/v2/index.json :: linear probe qg_herwig, largest set, zero: R30 of L162_MASS over the BatchNorm-recomputed random network of the same index :: sha256 c6c4138ee5c86423 :: prose twin of the run farther from 1
\newcommand{\MlstLpQgHerwigRthirtyCoarseMassOverInitBn}{1.146 and 1.159}  % experiments/FIGS/data/v2/index.json :: linear probe qg_herwig, largest set, chosen substitution: R30 of R16_Q1_MASS over the BatchNorm-recomputed random network of the same index :: sha256 c6c4138ee5c86423
\newcommand{\MlstLpQgHerwigRthirtyCoarseMassOverInitBnPr}{equal within 20\%}  % experiments/FIGS/data/v2/index.json :: linear probe qg_herwig, largest set, chosen substitution: R30 of R16_Q1_MASS over the BatchNorm-recomputed random network of the same index :: sha256 c6c4138ee5c86423 :: prose twin: two significant figures, direction in words
\newcommand{\MlstLpQgHerwigRthirtyCoarseMassOverInitBnPrMax}{equal within 20\%}  % experiments/FIGS/data/v2/index.json :: linear probe qg_herwig, largest set, chosen substitution: R30 of R16_Q1_MASS over the BatchNorm-recomputed random network of the same index :: sha256 c6c4138ee5c86423 :: prose twin of the run farther from 1
\newcommand{\MlstLpQgHerwigRthirtyCoarseMassOverInitBnZero}{1.070 and 1.100}  % experiments/FIGS/data/v2/index.json :: linear probe qg_herwig, largest set, zero: R30 of R16_Q1_MASS over the BatchNorm-recomputed random network of the same index :: sha256 c6c4138ee5c86423
\newcommand{\MlstLpQgHerwigRthirtyCoarseMassOverInitBnZeroPr}{equal within 20\%}  % experiments/FIGS/data/v2/index.json :: linear probe qg_herwig, largest set, zero: R30 of R16_Q1_MASS over the BatchNorm-recomputed random network of the same index :: sha256 c6c4138ee5c86423 :: prose twin: two significant figures, direction in words
\newcommand{\MlstLpQgHerwigRthirtyCoarseMassOverInitBnZeroPrMax}{equal within 20\%}  % experiments/FIGS/data/v2/index.json :: linear probe qg_herwig, largest set, zero: R30 of R16_Q1_MASS over the BatchNorm-recomputed random network of the same index :: sha256 c6c4138ee5c86423 :: prose twin of the run farther from 1
\newcommand{\MlstLpQgHerwigRthirtyResOverFull}{1.03\,\ensuremath{\pm}\,0.11}  % experiments/FIGS/data/v2/index.json :: linear probe qg_herwig: R30 of L162 over the 188-class run :: sha256 c6c4138ee5c86423
\newcommand{\MlstLpQgHerwigRthirtyResOverFullPr}{equal within 20\%}  % experiments/FIGS/data/v2/index.json :: linear probe qg_herwig: R30 of L162 over the 188-class run :: sha256 c6c4138ee5c86423 :: prose twin: two significant figures, direction in words
\newcommand{\MlstLpQgHerwigRthirtyResOverFullZero}{0.93\,\ensuremath{\pm}\,0.11}  % experiments/FIGS/data/v2/index.json :: linear probe qg_herwig: R30 of L162 over the 188-class run :: sha256 c6c4138ee5c86423
\newcommand{\MlstLpQgHerwigRthirtyResOverFullZeroPr}{equal within 20\%}  % experiments/FIGS/data/v2/index.json :: linear probe qg_herwig: R30 of L162 over the 188-class run :: sha256 c6c4138ee5c86423 :: prose twin: two significant figures, direction in words
\newcommand{\MlstLpQgHerwigRthirtyMidOverFull}{0.98\,\ensuremath{\pm}\,0.06}  % experiments/FIGS/data/v2/index.json :: linear probe qg_herwig: R30 of R42_Q1 over the 188-class run :: sha256 c6c4138ee5c86423
\newcommand{\MlstLpQgHerwigRthirtyMidOverFullPr}{equal within 6\%}  % experiments/FIGS/data/v2/index.json :: linear probe qg_herwig: R30 of R42_Q1 over the 188-class run :: sha256 c6c4138ee5c86423 :: prose twin: two significant figures, direction in words
\newcommand{\MlstLpQgHerwigRthirtyMidOverFullZero}{0.94\,\ensuremath{\pm}\,0.09}  % experiments/FIGS/data/v2/index.json :: linear probe qg_herwig: R30 of R42_Q1 over the 188-class run :: sha256 c6c4138ee5c86423
\newcommand{\MlstLpQgHerwigRthirtyMidOverFullZeroPr}{equal within 9\%}  % experiments/FIGS/data/v2/index.json :: linear probe qg_herwig: R30 of R42_Q1 over the 188-class run :: sha256 c6c4138ee5c86423 :: prose twin: two significant figures, direction in words
\newcommand{\MlstLpQgHerwigRthirtyCoarseOverFull}{1.02\,\ensuremath{\pm}\,0.09}  % experiments/FIGS/data/v2/index.json :: linear probe qg_herwig: R30 of R16_Q1 over the 188-class run :: sha256 c6c4138ee5c86423
\newcommand{\MlstLpQgHerwigRthirtyCoarseOverFullPr}{equal within 10\%}  % experiments/FIGS/data/v2/index.json :: linear probe qg_herwig: R30 of R16_Q1 over the 188-class run :: sha256 c6c4138ee5c86423 :: prose twin: two significant figures, direction in words
\newcommand{\MlstLpQgHerwigRthirtyCoarseOverFullZero}{1.018\,\ensuremath{\pm}\,0.024}  % experiments/FIGS/data/v2/index.json :: linear probe qg_herwig: R30 of R16_Q1 over the 188-class run :: sha256 c6c4138ee5c86423
\newcommand{\MlstLpQgHerwigRthirtyCoarseOverFullZeroPr}{equal within 3\%}  % experiments/FIGS/data/v2/index.json :: linear probe qg_herwig: R30 of R16_Q1 over the 188-class run :: sha256 c6c4138ee5c86423 :: prose twin: two significant figures, direction in words
\newcommand{\MlstLpQgHerwigRthirtyMeanOverZeroFull}{1.07\,\ensuremath{\pm}\,0.11}  % experiments/FIGS/data/v2/index.json :: linear probe qg_herwig: R30 with mean over R30 with zero, L188 :: sha256 c6c4138ee5c86423
\newcommand{\MlstLpQgHerwigRthirtyMeanOverZeroFullPr}{equal within 20\%}  % experiments/FIGS/data/v2/index.json :: linear probe qg_herwig: R30 with mean over R30 with zero, L188 :: sha256 c6c4138ee5c86423 :: prose twin: two significant figures, direction in words
\newcommand{\MlstLpQgHerwigRthirtyDrawOverZeroFull}{1.07\,\ensuremath{\pm}\,0.06}  % experiments/FIGS/data/v2/index.json :: linear probe qg_herwig: R30 with draws over R30 with zero, L188 :: sha256 c6c4138ee5c86423
\newcommand{\MlstLpQgHerwigRthirtyDrawOverZeroFullPr}{equal within 7\%}  % experiments/FIGS/data/v2/index.json :: linear probe qg_herwig: R30 with draws over R30 with zero, L188 :: sha256 c6c4138ee5c86423 :: prose twin: two significant figures, direction in words
\newcommand{\MlstLpQgHerwigRthirtyMeanOverZeroInitBn}{0.988\,\ensuremath{\pm}\,0.019}  % experiments/FIGS/data/v2/index.json :: linear probe qg_herwig: R30 with mean over R30 with zero, INIT_BN :: sha256 c6c4138ee5c86423
\newcommand{\MlstLpQgHerwigRthirtyMeanOverZeroInitBnPr}{equal within 2\%}  % experiments/FIGS/data/v2/index.json :: linear probe qg_herwig: R30 with mean over R30 with zero, INIT_BN :: sha256 c6c4138ee5c86423 :: prose twin: two significant figures, direction in words
\newcommand{\MlstLpQgHerwigRthirtyDrawOverZeroInitBn}{0.983\,\ensuremath{\pm}\,0.032}  % experiments/FIGS/data/v2/index.json :: linear probe qg_herwig: R30 with draws over R30 with zero, INIT_BN :: sha256 c6c4138ee5c86423
\newcommand{\MlstLpQgHerwigRthirtyDrawOverZeroInitBnPr}{equal within 4\%}  % experiments/FIGS/data/v2/index.json :: linear probe qg_herwig: R30 with draws over R30 with zero, INIT_BN :: sha256 c6c4138ee5c86423 :: prose twin: two significant figures, direction in words
\newcommand{\MlstLpQgHerwigRthirtyFullZero}{32.8\,\ensuremath{\pm}\,1.0}  % experiments/FIGS/data/v2/index.json :: linear probe qg_herwig, largest set: R30 with zero, L188 :: sha256 c6c4138ee5c86423
\newcommand{\MlstLpQgHerwigRthirtyFullMean}{35.0\,\ensuremath{\pm}\,2.6}  % experiments/FIGS/data/v2/index.json :: linear probe qg_herwig, largest set: R30 with mean, L188 :: sha256 c6c4138ee5c86423
\newcommand{\MlstLpQgHerwigRthirtyFullDraws}{35.0\,\ensuremath{\pm}\,0.8}  % experiments/FIGS/data/v2/index.json :: linear probe qg_herwig, largest set: R30 with draws, L188 :: sha256 c6c4138ee5c86423
\newcommand{\MlstLpQgHerwigRthirtyInitBnZero}{31.1\,\ensuremath{\pm}\,0.5}  % experiments/FIGS/data/v2/index.json :: linear probe qg_herwig, largest set: R30 with zero, INIT_BN :: sha256 c6c4138ee5c86423
\newcommand{\MlstLpQgHerwigRthirtyInitBnMean}{30.7\,\ensuremath{\pm}\,0.5}  % experiments/FIGS/data/v2/index.json :: linear probe qg_herwig, largest set: R30 with mean, INIT_BN :: sha256 c6c4138ee5c86423
\newcommand{\MlstLpQgHerwigRthirtyInitBnDraws}{30.6\,\ensuremath{\pm}\,0.6}  % experiments/FIGS/data/v2/index.json :: linear probe qg_herwig, largest set: R30 with draws, INIT_BN :: sha256 c6c4138ee5c86423
\newcommand{\MlstLpJcTwoTopZeroAccDropFull}{88.9\,\ensuremath{\pm}\,1.0\%}  % experiments/FIGS/data/v2/index.json :: linear probe JetClass-II, largest set: relative drop in accuracy (%) with the top data's missing slots set to zero, L188 :: sha256 c6c4138ee5c86423
\newcommand{\MlstLpJcTwoTopZeroAccDropFullPr}{89\%}  % experiments/FIGS/data/v2/index.json :: linear probe JetClass-II, largest set: relative drop in accuracy (%) with the top data's missing slots set to zero, L188 :: sha256 c6c4138ee5c86423 :: prose twin, two significant figures
\newcommand{\MlstLpJcTwoTopMeanAccDropFull}{87.9\,\ensuremath{\pm}\,1.1\%}  % experiments/FIGS/data/v2/index.json :: linear probe JetClass-II, largest set: relative drop in accuracy (%) with the top data's missing slots set to mean, L188 :: sha256 c6c4138ee5c86423
\newcommand{\MlstLpJcTwoTopMeanAccDropFullPr}{88\%}  % experiments/FIGS/data/v2/index.json :: linear probe JetClass-II, largest set: relative drop in accuracy (%) with the top data's missing slots set to mean, L188 :: sha256 c6c4138ee5c86423 :: prose twin, two significant figures
\newcommand{\MlstLpJcTwoTopDrawAccDropFull}{85.98\,\ensuremath{\pm}\,0.07\%}  % experiments/FIGS/data/v2/index.json :: linear probe JetClass-II, largest set: relative drop in accuracy (%) with the top data's missing slots set to draws, L188 :: sha256 c6c4138ee5c86423
\newcommand{\MlstLpJcTwoTopDrawAccDropFullPr}{86\%}  % experiments/FIGS/data/v2/index.json :: linear probe JetClass-II, largest set: relative drop in accuracy (%) with the top data's missing slots set to draws, L188 :: sha256 c6c4138ee5c86423 :: prose twin, two significant figures
\newcommand{\MlstLpJcTwoTopZeroAccDropRes}{89.0\,\ensuremath{\pm}\,0.9\%}  % experiments/FIGS/data/v2/index.json :: linear probe JetClass-II, largest set: relative drop in accuracy (%) with the top data's missing slots set to zero, L162 :: sha256 c6c4138ee5c86423
\newcommand{\MlstLpJcTwoTopZeroAccDropResPr}{89\%}  % experiments/FIGS/data/v2/index.json :: linear probe JetClass-II, largest set: relative drop in accuracy (%) with the top data's missing slots set to zero, L162 :: sha256 c6c4138ee5c86423 :: prose twin, two significant figures
\newcommand{\MlstLpJcTwoTopMeanAccDropRes}{87.8\,\ensuremath{\pm}\,0.7\%}  % experiments/FIGS/data/v2/index.json :: linear probe JetClass-II, largest set: relative drop in accuracy (%) with the top data's missing slots set to mean, L162 :: sha256 c6c4138ee5c86423
\newcommand{\MlstLpJcTwoTopMeanAccDropResPr}{88\%}  % experiments/FIGS/data/v2/index.json :: linear probe JetClass-II, largest set: relative drop in accuracy (%) with the top data's missing slots set to mean, L162 :: sha256 c6c4138ee5c86423 :: prose twin, two significant figures
\newcommand{\MlstLpJcTwoTopDrawAccDropRes}{85.89\,\ensuremath{\pm}\,0.11\%}  % experiments/FIGS/data/v2/index.json :: linear probe JetClass-II, largest set: relative drop in accuracy (%) with the top data's missing slots set to draws, L162 :: sha256 c6c4138ee5c86423
\newcommand{\MlstLpJcTwoTopDrawAccDropResPr}{86\%}  % experiments/FIGS/data/v2/index.json :: linear probe JetClass-II, largest set: relative drop in accuracy (%) with the top data's missing slots set to draws, L162 :: sha256 c6c4138ee5c86423 :: prose twin, two significant figures
\newcommand{\MlstLpJcTwoTopZeroAccDropMid}{87.81\,\ensuremath{\pm}\,0.28\%}  % experiments/FIGS/data/v2/index.json :: linear probe JetClass-II, largest set: relative drop in accuracy (%) with the top data's missing slots set to zero, R42_Q1 :: sha256 c6c4138ee5c86423
\newcommand{\MlstLpJcTwoTopZeroAccDropMidPr}{88\%}  % experiments/FIGS/data/v2/index.json :: linear probe JetClass-II, largest set: relative drop in accuracy (%) with the top data's missing slots set to zero, R42_Q1 :: sha256 c6c4138ee5c86423 :: prose twin, two significant figures
\newcommand{\MlstLpJcTwoTopMeanAccDropMid}{86.7\,\ensuremath{\pm}\,0.9\%}  % experiments/FIGS/data/v2/index.json :: linear probe JetClass-II, largest set: relative drop in accuracy (%) with the top data's missing slots set to mean, R42_Q1 :: sha256 c6c4138ee5c86423
\newcommand{\MlstLpJcTwoTopMeanAccDropMidPr}{87\%}  % experiments/FIGS/data/v2/index.json :: linear probe JetClass-II, largest set: relative drop in accuracy (%) with the top data's missing slots set to mean, R42_Q1 :: sha256 c6c4138ee5c86423 :: prose twin, two significant figures
\newcommand{\MlstLpJcTwoTopDrawAccDropMid}{83.17\,\ensuremath{\pm}\,0.29\%}  % experiments/FIGS/data/v2/index.json :: linear probe JetClass-II, largest set: relative drop in accuracy (%) with the top data's missing slots set to draws, R42_Q1 :: sha256 c6c4138ee5c86423
\newcommand{\MlstLpJcTwoTopDrawAccDropMidPr}{83\%}  % experiments/FIGS/data/v2/index.json :: linear probe JetClass-II, largest set: relative drop in accuracy (%) with the top data's missing slots set to draws, R42_Q1 :: sha256 c6c4138ee5c86423 :: prose twin, two significant figures
\newcommand{\MlstLpJcTwoTopZeroAccDropCoarse}{86.7\,\ensuremath{\pm}\,0.7\%}  % experiments/FIGS/data/v2/index.json :: linear probe JetClass-II, largest set: relative drop in accuracy (%) with the top data's missing slots set to zero, R16_Q1 :: sha256 c6c4138ee5c86423
\newcommand{\MlstLpJcTwoTopZeroAccDropCoarsePr}{87\%}  % experiments/FIGS/data/v2/index.json :: linear probe JetClass-II, largest set: relative drop in accuracy (%) with the top data's missing slots set to zero, R16_Q1 :: sha256 c6c4138ee5c86423 :: prose twin, two significant figures
\newcommand{\MlstLpJcTwoTopMeanAccDropCoarse}{86.1\,\ensuremath{\pm}\,0.8\%}  % experiments/FIGS/data/v2/index.json :: linear probe JetClass-II, largest set: relative drop in accuracy (%) with the top data's missing slots set to mean, R16_Q1 :: sha256 c6c4138ee5c86423
\newcommand{\MlstLpJcTwoTopMeanAccDropCoarsePr}{86\%}  % experiments/FIGS/data/v2/index.json :: linear probe JetClass-II, largest set: relative drop in accuracy (%) with the top data's missing slots set to mean, R16_Q1 :: sha256 c6c4138ee5c86423 :: prose twin, two significant figures
\newcommand{\MlstLpJcTwoTopDrawAccDropCoarse}{81.4\,\ensuremath{\pm}\,0.5\%}  % experiments/FIGS/data/v2/index.json :: linear probe JetClass-II, largest set: relative drop in accuracy (%) with the top data's missing slots set to draws, R16_Q1 :: sha256 c6c4138ee5c86423
\newcommand{\MlstLpJcTwoTopDrawAccDropCoarsePr}{81\%}  % experiments/FIGS/data/v2/index.json :: linear probe JetClass-II, largest set: relative drop in accuracy (%) with the top data's missing slots set to draws, R16_Q1 :: sha256 c6c4138ee5c86423 :: prose twin, two significant figures
\newcommand{\MlstLpJcTwoTopZeroAccDropResMass}{89.2, 88.3\%}  % experiments/FIGS/data/v2/index.json :: linear probe JetClass-II, largest set: relative drop in accuracy (%) with the top data's missing slots set to zero, L162_MASS :: sha256 c6c4138ee5c86423
\newcommand{\MlstLpJcTwoTopZeroAccDropResMassPr}{88--89\%}  % experiments/FIGS/data/v2/index.json :: linear probe JetClass-II, largest set: relative drop in accuracy (%) with the top data's missing slots set to zero, L162_MASS :: sha256 c6c4138ee5c86423 :: prose twin, two significant figures
\newcommand{\MlstLpJcTwoTopMeanAccDropResMass}{88.5, 87.8\%}  % experiments/FIGS/data/v2/index.json :: linear probe JetClass-II, largest set: relative drop in accuracy (%) with the top data's missing slots set to mean, L162_MASS :: sha256 c6c4138ee5c86423
\newcommand{\MlstLpJcTwoTopMeanAccDropResMassPr}{88--88\%}  % experiments/FIGS/data/v2/index.json :: linear probe JetClass-II, largest set: relative drop in accuracy (%) with the top data's missing slots set to mean, L162_MASS :: sha256 c6c4138ee5c86423 :: prose twin, two significant figures
\newcommand{\MlstLpJcTwoTopDrawAccDropResMass}{85.4, 85.4\%}  % experiments/FIGS/data/v2/index.json :: linear probe JetClass-II, largest set: relative drop in accuracy (%) with the top data's missing slots set to draws, L162_MASS :: sha256 c6c4138ee5c86423
\newcommand{\MlstLpJcTwoTopDrawAccDropResMassPr}{85--85\%}  % experiments/FIGS/data/v2/index.json :: linear probe JetClass-II, largest set: relative drop in accuracy (%) with the top data's missing slots set to draws, L162_MASS :: sha256 c6c4138ee5c86423 :: prose twin, two significant figures
\newcommand{\MlstLpJcTwoTopZeroAccDropCoarseMass}{84.8, 87.4\%}  % experiments/FIGS/data/v2/index.json :: linear probe JetClass-II, largest set: relative drop in accuracy (%) with the top data's missing slots set to zero, R16_Q1_MASS :: sha256 c6c4138ee5c86423
\newcommand{\MlstLpJcTwoTopZeroAccDropCoarseMassPr}{85--87\%}  % experiments/FIGS/data/v2/index.json :: linear probe JetClass-II, largest set: relative drop in accuracy (%) with the top data's missing slots set to zero, R16_Q1_MASS :: sha256 c6c4138ee5c86423 :: prose twin, two significant figures
\newcommand{\MlstLpJcTwoTopMeanAccDropCoarseMass}{85.9, 86.9\%}  % experiments/FIGS/data/v2/index.json :: linear probe JetClass-II, largest set: relative drop in accuracy (%) with the top data's missing slots set to mean, R16_Q1_MASS :: sha256 c6c4138ee5c86423
\newcommand{\MlstLpJcTwoTopMeanAccDropCoarseMassPr}{86--87\%}  % experiments/FIGS/data/v2/index.json :: linear probe JetClass-II, largest set: relative drop in accuracy (%) with the top data's missing slots set to mean, R16_Q1_MASS :: sha256 c6c4138ee5c86423 :: prose twin, two significant figures
\newcommand{\MlstLpJcTwoTopDrawAccDropCoarseMass}{79.1, 81.7\%}  % experiments/FIGS/data/v2/index.json :: linear probe JetClass-II, largest set: relative drop in accuracy (%) with the top data's missing slots set to draws, R16_Q1_MASS :: sha256 c6c4138ee5c86423
\newcommand{\MlstLpJcTwoTopDrawAccDropCoarseMassPr}{79--82\%}  % experiments/FIGS/data/v2/index.json :: linear probe JetClass-II, largest set: relative drop in accuracy (%) with the top data's missing slots set to draws, R16_Q1_MASS :: sha256 c6c4138ee5c86423 :: prose twin, two significant figures
\newcommand{\MlstLpJcTwoTopZeroAccDropFzero}{88.9, 88.5\%}  % experiments/FIGS/data/v2/index.json :: linear probe JetClass-II, largest set: relative drop in accuracy (%) with the top data's missing slots set to zero, FLAV_F0 :: sha256 c6c4138ee5c86423
\newcommand{\MlstLpJcTwoTopZeroAccDropFzeroPr}{88--89\%}  % experiments/FIGS/data/v2/index.json :: linear probe JetClass-II, largest set: relative drop in accuracy (%) with the top data's missing slots set to zero, FLAV_F0 :: sha256 c6c4138ee5c86423 :: prose twin, two significant figures
\newcommand{\MlstLpJcTwoTopMeanAccDropFzero}{88.3, 87.0\%}  % experiments/FIGS/data/v2/index.json :: linear probe JetClass-II, largest set: relative drop in accuracy (%) with the top data's missing slots set to mean, FLAV_F0 :: sha256 c6c4138ee5c86423
\newcommand{\MlstLpJcTwoTopMeanAccDropFzeroPr}{87--88\%}  % experiments/FIGS/data/v2/index.json :: linear probe JetClass-II, largest set: relative drop in accuracy (%) with the top data's missing slots set to mean, FLAV_F0 :: sha256 c6c4138ee5c86423 :: prose twin, two significant figures
\newcommand{\MlstLpJcTwoTopDrawAccDropFzero}{83.3, 83.9\%}  % experiments/FIGS/data/v2/index.json :: linear probe JetClass-II, largest set: relative drop in accuracy (%) with the top data's missing slots set to draws, FLAV_F0 :: sha256 c6c4138ee5c86423
\newcommand{\MlstLpJcTwoTopDrawAccDropFzeroPr}{83--84\%}  % experiments/FIGS/data/v2/index.json :: linear probe JetClass-II, largest set: relative drop in accuracy (%) with the top data's missing slots set to draws, FLAV_F0 :: sha256 c6c4138ee5c86423 :: prose twin, two significant figures
\newcommand{\MlstLpJcTwoTopZeroAccDropFone}{87.8, 87.6\%}  % experiments/FIGS/data/v2/index.json :: linear probe JetClass-II, largest set: relative drop in accuracy (%) with the top data's missing slots set to zero, FLAV_F1 :: sha256 c6c4138ee5c86423
\newcommand{\MlstLpJcTwoTopZeroAccDropFonePr}{88--88\%}  % experiments/FIGS/data/v2/index.json :: linear probe JetClass-II, largest set: relative drop in accuracy (%) with the top data's missing slots set to zero, FLAV_F1 :: sha256 c6c4138ee5c86423 :: prose twin, two significant figures
\newcommand{\MlstLpJcTwoTopMeanAccDropFone}{88.9, 86.9\%}  % experiments/FIGS/data/v2/index.json :: linear probe JetClass-II, largest set: relative drop in accuracy (%) with the top data's missing slots set to mean, FLAV_F1 :: sha256 c6c4138ee5c86423
\newcommand{\MlstLpJcTwoTopMeanAccDropFonePr}{87--89\%}  % experiments/FIGS/data/v2/index.json :: linear probe JetClass-II, largest set: relative drop in accuracy (%) with the top data's missing slots set to mean, FLAV_F1 :: sha256 c6c4138ee5c86423 :: prose twin, two significant figures
\newcommand{\MlstLpJcTwoTopDrawAccDropFone}{84.4, 85.2\%}  % experiments/FIGS/data/v2/index.json :: linear probe JetClass-II, largest set: relative drop in accuracy (%) with the top data's missing slots set to draws, FLAV_F1 :: sha256 c6c4138ee5c86423
\newcommand{\MlstLpJcTwoTopDrawAccDropFonePr}{84--85\%}  % experiments/FIGS/data/v2/index.json :: linear probe JetClass-II, largest set: relative drop in accuracy (%) with the top data's missing slots set to draws, FLAV_F1 :: sha256 c6c4138ee5c86423 :: prose twin, two significant figures
\newcommand{\MlstLpJcTwoTopZeroAccDropFoner}{89.6, 90.3\%}  % experiments/FIGS/data/v2/index.json :: linear probe JetClass-II, largest set: relative drop in accuracy (%) with the top data's missing slots set to zero, FLAV_F1R :: sha256 c6c4138ee5c86423
\newcommand{\MlstLpJcTwoTopZeroAccDropFonerPr}{90--90\%}  % experiments/FIGS/data/v2/index.json :: linear probe JetClass-II, largest set: relative drop in accuracy (%) with the top data's missing slots set to zero, FLAV_F1R :: sha256 c6c4138ee5c86423 :: prose twin, two significant figures
\newcommand{\MlstLpJcTwoTopMeanAccDropFoner}{89.4, 89.7\%}  % experiments/FIGS/data/v2/index.json :: linear probe JetClass-II, largest set: relative drop in accuracy (%) with the top data's missing slots set to mean, FLAV_F1R :: sha256 c6c4138ee5c86423
\newcommand{\MlstLpJcTwoTopMeanAccDropFonerPr}{89--90\%}  % experiments/FIGS/data/v2/index.json :: linear probe JetClass-II, largest set: relative drop in accuracy (%) with the top data's missing slots set to mean, FLAV_F1R :: sha256 c6c4138ee5c86423 :: prose twin, two significant figures
\newcommand{\MlstLpJcTwoTopDrawAccDropFoner}{84.9, 85.6\%}  % experiments/FIGS/data/v2/index.json :: linear probe JetClass-II, largest set: relative drop in accuracy (%) with the top data's missing slots set to draws, FLAV_F1R :: sha256 c6c4138ee5c86423
\newcommand{\MlstLpJcTwoTopDrawAccDropFonerPr}{85--86\%}  % experiments/FIGS/data/v2/index.json :: linear probe JetClass-II, largest set: relative drop in accuracy (%) with the top data's missing slots set to draws, FLAV_F1R :: sha256 c6c4138ee5c86423 :: prose twin, two significant figures
\newcommand{\MlstLpJcTwoTopZeroAccDropInitBn}{56.4\,\ensuremath{\pm}\,1.4\%}  % experiments/FIGS/data/v2/index.json :: linear probe JetClass-II, largest set: relative drop in accuracy (%) with the top data's missing slots set to zero, INIT_BN :: sha256 c6c4138ee5c86423
\newcommand{\MlstLpJcTwoTopZeroAccDropInitBnPr}{56\%}  % experiments/FIGS/data/v2/index.json :: linear probe JetClass-II, largest set: relative drop in accuracy (%) with the top data's missing slots set to zero, INIT_BN :: sha256 c6c4138ee5c86423 :: prose twin, two significant figures
\newcommand{\MlstLpJcTwoTopMeanAccDropInitBn}{59.8\,\ensuremath{\pm}\,0.8\%}  % experiments/FIGS/data/v2/index.json :: linear probe JetClass-II, largest set: relative drop in accuracy (%) with the top data's missing slots set to mean, INIT_BN :: sha256 c6c4138ee5c86423
\newcommand{\MlstLpJcTwoTopMeanAccDropInitBnPr}{60\%}  % experiments/FIGS/data/v2/index.json :: linear probe JetClass-II, largest set: relative drop in accuracy (%) with the top data's missing slots set to mean, INIT_BN :: sha256 c6c4138ee5c86423 :: prose twin, two significant figures
\newcommand{\MlstLpJcTwoTopDrawAccDropInitBn}{59.8\,\ensuremath{\pm}\,1.2\%}  % experiments/FIGS/data/v2/index.json :: linear probe JetClass-II, largest set: relative drop in accuracy (%) with the top data's missing slots set to draws, INIT_BN :: sha256 c6c4138ee5c86423
\newcommand{\MlstLpJcTwoTopDrawAccDropInitBnPr}{60\%}  % experiments/FIGS/data/v2/index.json :: linear probe JetClass-II, largest set: relative drop in accuracy (%) with the top data's missing slots set to draws, INIT_BN :: sha256 c6c4138ee5c86423 :: prose twin, two significant figures
\newcommand{\MlstLpJcTwoQgZeroAccDropFull}{61.2\,\ensuremath{\pm}\,3.0\%}  % experiments/FIGS/data/v2/index.json :: linear probe JetClass-II, largest set: relative drop in accuracy (%) with the qg data's missing slots set to zero, L188 :: sha256 c6c4138ee5c86423
\newcommand{\MlstLpJcTwoQgZeroAccDropFullPr}{61\%}  % experiments/FIGS/data/v2/index.json :: linear probe JetClass-II, largest set: relative drop in accuracy (%) with the qg data's missing slots set to zero, L188 :: sha256 c6c4138ee5c86423 :: prose twin, two significant figures
\newcommand{\MlstLpJcTwoQgMeanAccDropFull}{48.8\,\ensuremath{\pm}\,1.1\%}  % experiments/FIGS/data/v2/index.json :: linear probe JetClass-II, largest set: relative drop in accuracy (%) with the qg data's missing slots set to mean, L188 :: sha256 c6c4138ee5c86423
\newcommand{\MlstLpJcTwoQgMeanAccDropFullPr}{49\%}  % experiments/FIGS/data/v2/index.json :: linear probe JetClass-II, largest set: relative drop in accuracy (%) with the qg data's missing slots set to mean, L188 :: sha256 c6c4138ee5c86423 :: prose twin, two significant figures
\newcommand{\MlstLpJcTwoQgDrawAccDropFull}{54.30\,\ensuremath{\pm}\,0.21\%}  % experiments/FIGS/data/v2/index.json :: linear probe JetClass-II, largest set: relative drop in accuracy (%) with the qg data's missing slots set to draws, L188 :: sha256 c6c4138ee5c86423
\newcommand{\MlstLpJcTwoQgDrawAccDropFullPr}{54\%}  % experiments/FIGS/data/v2/index.json :: linear probe JetClass-II, largest set: relative drop in accuracy (%) with the qg data's missing slots set to draws, L188 :: sha256 c6c4138ee5c86423 :: prose twin, two significant figures
\newcommand{\MlstLpJcTwoQgZeroAccDropRes}{64\,\ensuremath{\pm}\,5\%}  % experiments/FIGS/data/v2/index.json :: linear probe JetClass-II, largest set: relative drop in accuracy (%) with the qg data's missing slots set to zero, L162 :: sha256 c6c4138ee5c86423
\newcommand{\MlstLpJcTwoQgZeroAccDropResPr}{64\%}  % experiments/FIGS/data/v2/index.json :: linear probe JetClass-II, largest set: relative drop in accuracy (%) with the qg data's missing slots set to zero, L162 :: sha256 c6c4138ee5c86423 :: prose twin, two significant figures
\newcommand{\MlstLpJcTwoQgMeanAccDropRes}{49.5\,\ensuremath{\pm}\,0.5\%}  % experiments/FIGS/data/v2/index.json :: linear probe JetClass-II, largest set: relative drop in accuracy (%) with the qg data's missing slots set to mean, L162 :: sha256 c6c4138ee5c86423
\newcommand{\MlstLpJcTwoQgMeanAccDropResPr}{49\%}  % experiments/FIGS/data/v2/index.json :: linear probe JetClass-II, largest set: relative drop in accuracy (%) with the qg data's missing slots set to mean, L162 :: sha256 c6c4138ee5c86423 :: prose twin, two significant figures
\newcommand{\MlstLpJcTwoQgDrawAccDropRes}{54.6\,\ensuremath{\pm}\,0.6\%}  % experiments/FIGS/data/v2/index.json :: linear probe JetClass-II, largest set: relative drop in accuracy (%) with the qg data's missing slots set to draws, L162 :: sha256 c6c4138ee5c86423
\newcommand{\MlstLpJcTwoQgDrawAccDropResPr}{55\%}  % experiments/FIGS/data/v2/index.json :: linear probe JetClass-II, largest set: relative drop in accuracy (%) with the qg data's missing slots set to draws, L162 :: sha256 c6c4138ee5c86423 :: prose twin, two significant figures
\newcommand{\MlstLpJcTwoQgZeroAccDropMid}{68\,\ensuremath{\pm}\,5\%}  % experiments/FIGS/data/v2/index.json :: linear probe JetClass-II, largest set: relative drop in accuracy (%) with the qg data's missing slots set to zero, R42_Q1 :: sha256 c6c4138ee5c86423
\newcommand{\MlstLpJcTwoQgZeroAccDropMidPr}{68\%}  % experiments/FIGS/data/v2/index.json :: linear probe JetClass-II, largest set: relative drop in accuracy (%) with the qg data's missing slots set to zero, R42_Q1 :: sha256 c6c4138ee5c86423 :: prose twin, two significant figures
\newcommand{\MlstLpJcTwoQgMeanAccDropMid}{47.6\,\ensuremath{\pm}\,1.6\%}  % experiments/FIGS/data/v2/index.json :: linear probe JetClass-II, largest set: relative drop in accuracy (%) with the qg data's missing slots set to mean, R42_Q1 :: sha256 c6c4138ee5c86423
\newcommand{\MlstLpJcTwoQgMeanAccDropMidPr}{48\%}  % experiments/FIGS/data/v2/index.json :: linear probe JetClass-II, largest set: relative drop in accuracy (%) with the qg data's missing slots set to mean, R42_Q1 :: sha256 c6c4138ee5c86423 :: prose twin, two significant figures
\newcommand{\MlstLpJcTwoQgDrawAccDropMid}{53.5\,\ensuremath{\pm}\,1.4\%}  % experiments/FIGS/data/v2/index.json :: linear probe JetClass-II, largest set: relative drop in accuracy (%) with the qg data's missing slots set to draws, R42_Q1 :: sha256 c6c4138ee5c86423
\newcommand{\MlstLpJcTwoQgDrawAccDropMidPr}{54\%}  % experiments/FIGS/data/v2/index.json :: linear probe JetClass-II, largest set: relative drop in accuracy (%) with the qg data's missing slots set to draws, R42_Q1 :: sha256 c6c4138ee5c86423 :: prose twin, two significant figures
\newcommand{\MlstLpJcTwoQgZeroAccDropCoarse}{70.9\,\ensuremath{\pm}\,2.3\%}  % experiments/FIGS/data/v2/index.json :: linear probe JetClass-II, largest set: relative drop in accuracy (%) with the qg data's missing slots set to zero, R16_Q1 :: sha256 c6c4138ee5c86423
\newcommand{\MlstLpJcTwoQgZeroAccDropCoarsePr}{71\%}  % experiments/FIGS/data/v2/index.json :: linear probe JetClass-II, largest set: relative drop in accuracy (%) with the qg data's missing slots set to zero, R16_Q1 :: sha256 c6c4138ee5c86423 :: prose twin, two significant figures
\newcommand{\MlstLpJcTwoQgMeanAccDropCoarse}{50.0\,\ensuremath{\pm}\,1.0\%}  % experiments/FIGS/data/v2/index.json :: linear probe JetClass-II, largest set: relative drop in accuracy (%) with the qg data's missing slots set to mean, R16_Q1 :: sha256 c6c4138ee5c86423
\newcommand{\MlstLpJcTwoQgMeanAccDropCoarsePr}{50\%}  % experiments/FIGS/data/v2/index.json :: linear probe JetClass-II, largest set: relative drop in accuracy (%) with the qg data's missing slots set to mean, R16_Q1 :: sha256 c6c4138ee5c86423 :: prose twin, two significant figures
\newcommand{\MlstLpJcTwoQgDrawAccDropCoarse}{51.7\,\ensuremath{\pm}\,0.4\%}  % experiments/FIGS/data/v2/index.json :: linear probe JetClass-II, largest set: relative drop in accuracy (%) with the qg data's missing slots set to draws, R16_Q1 :: sha256 c6c4138ee5c86423
\newcommand{\MlstLpJcTwoQgDrawAccDropCoarsePr}{52\%}  % experiments/FIGS/data/v2/index.json :: linear probe JetClass-II, largest set: relative drop in accuracy (%) with the qg data's missing slots set to draws, R16_Q1 :: sha256 c6c4138ee5c86423 :: prose twin, two significant figures
\newcommand{\MlstLpJcTwoQgZeroAccDropResMass}{61.8, 67.9\%}  % experiments/FIGS/data/v2/index.json :: linear probe JetClass-II, largest set: relative drop in accuracy (%) with the qg data's missing slots set to zero, L162_MASS :: sha256 c6c4138ee5c86423
\newcommand{\MlstLpJcTwoQgZeroAccDropResMassPr}{62--68\%}  % experiments/FIGS/data/v2/index.json :: linear probe JetClass-II, largest set: relative drop in accuracy (%) with the qg data's missing slots set to zero, L162_MASS :: sha256 c6c4138ee5c86423 :: prose twin, two significant figures
\newcommand{\MlstLpJcTwoQgMeanAccDropResMass}{49.7, 52.2\%}  % experiments/FIGS/data/v2/index.json :: linear probe JetClass-II, largest set: relative drop in accuracy (%) with the qg data's missing slots set to mean, L162_MASS :: sha256 c6c4138ee5c86423
\newcommand{\MlstLpJcTwoQgMeanAccDropResMassPr}{50--52\%}  % experiments/FIGS/data/v2/index.json :: linear probe JetClass-II, largest set: relative drop in accuracy (%) with the qg data's missing slots set to mean, L162_MASS :: sha256 c6c4138ee5c86423 :: prose twin, two significant figures
\newcommand{\MlstLpJcTwoQgDrawAccDropResMass}{54.3, 55.1\%}  % experiments/FIGS/data/v2/index.json :: linear probe JetClass-II, largest set: relative drop in accuracy (%) with the qg data's missing slots set to draws, L162_MASS :: sha256 c6c4138ee5c86423
\newcommand{\MlstLpJcTwoQgDrawAccDropResMassPr}{54--55\%}  % experiments/FIGS/data/v2/index.json :: linear probe JetClass-II, largest set: relative drop in accuracy (%) with the qg data's missing slots set to draws, L162_MASS :: sha256 c6c4138ee5c86423 :: prose twin, two significant figures
\newcommand{\MlstLpJcTwoQgZeroAccDropCoarseMass}{66.7, 65.4\%}  % experiments/FIGS/data/v2/index.json :: linear probe JetClass-II, largest set: relative drop in accuracy (%) with the qg data's missing slots set to zero, R16_Q1_MASS :: sha256 c6c4138ee5c86423
\newcommand{\MlstLpJcTwoQgZeroAccDropCoarseMassPr}{65--67\%}  % experiments/FIGS/data/v2/index.json :: linear probe JetClass-II, largest set: relative drop in accuracy (%) with the qg data's missing slots set to zero, R16_Q1_MASS :: sha256 c6c4138ee5c86423 :: prose twin, two significant figures
\newcommand{\MlstLpJcTwoQgMeanAccDropCoarseMass}{49.0, 49.6\%}  % experiments/FIGS/data/v2/index.json :: linear probe JetClass-II, largest set: relative drop in accuracy (%) with the qg data's missing slots set to mean, R16_Q1_MASS :: sha256 c6c4138ee5c86423
\newcommand{\MlstLpJcTwoQgMeanAccDropCoarseMassPr}{49--50\%}  % experiments/FIGS/data/v2/index.json :: linear probe JetClass-II, largest set: relative drop in accuracy (%) with the qg data's missing slots set to mean, R16_Q1_MASS :: sha256 c6c4138ee5c86423 :: prose twin, two significant figures
\newcommand{\MlstLpJcTwoQgDrawAccDropCoarseMass}{50.8, 51.1\%}  % experiments/FIGS/data/v2/index.json :: linear probe JetClass-II, largest set: relative drop in accuracy (%) with the qg data's missing slots set to draws, R16_Q1_MASS :: sha256 c6c4138ee5c86423
\newcommand{\MlstLpJcTwoQgDrawAccDropCoarseMassPr}{51--51\%}  % experiments/FIGS/data/v2/index.json :: linear probe JetClass-II, largest set: relative drop in accuracy (%) with the qg data's missing slots set to draws, R16_Q1_MASS :: sha256 c6c4138ee5c86423 :: prose twin, two significant figures
\newcommand{\MlstLpJcTwoQgZeroAccDropFzero}{62.7, 62.2\%}  % experiments/FIGS/data/v2/index.json :: linear probe JetClass-II, largest set: relative drop in accuracy (%) with the qg data's missing slots set to zero, FLAV_F0 :: sha256 c6c4138ee5c86423
\newcommand{\MlstLpJcTwoQgZeroAccDropFzeroPr}{62--63\%}  % experiments/FIGS/data/v2/index.json :: linear probe JetClass-II, largest set: relative drop in accuracy (%) with the qg data's missing slots set to zero, FLAV_F0 :: sha256 c6c4138ee5c86423 :: prose twin, two significant figures
\newcommand{\MlstLpJcTwoQgMeanAccDropFzero}{51.9, 51.3\%}  % experiments/FIGS/data/v2/index.json :: linear probe JetClass-II, largest set: relative drop in accuracy (%) with the qg data's missing slots set to mean, FLAV_F0 :: sha256 c6c4138ee5c86423
\newcommand{\MlstLpJcTwoQgMeanAccDropFzeroPr}{51--52\%}  % experiments/FIGS/data/v2/index.json :: linear probe JetClass-II, largest set: relative drop in accuracy (%) with the qg data's missing slots set to mean, FLAV_F0 :: sha256 c6c4138ee5c86423 :: prose twin, two significant figures
\newcommand{\MlstLpJcTwoQgDrawAccDropFzero}{53.9, 53.4\%}  % experiments/FIGS/data/v2/index.json :: linear probe JetClass-II, largest set: relative drop in accuracy (%) with the qg data's missing slots set to draws, FLAV_F0 :: sha256 c6c4138ee5c86423
\newcommand{\MlstLpJcTwoQgDrawAccDropFzeroPr}{53--54\%}  % experiments/FIGS/data/v2/index.json :: linear probe JetClass-II, largest set: relative drop in accuracy (%) with the qg data's missing slots set to draws, FLAV_F0 :: sha256 c6c4138ee5c86423 :: prose twin, two significant figures
\newcommand{\MlstLpJcTwoQgZeroAccDropFone}{66.2, 62.8\%}  % experiments/FIGS/data/v2/index.json :: linear probe JetClass-II, largest set: relative drop in accuracy (%) with the qg data's missing slots set to zero, FLAV_F1 :: sha256 c6c4138ee5c86423
\newcommand{\MlstLpJcTwoQgZeroAccDropFonePr}{63--66\%}  % experiments/FIGS/data/v2/index.json :: linear probe JetClass-II, largest set: relative drop in accuracy (%) with the qg data's missing slots set to zero, FLAV_F1 :: sha256 c6c4138ee5c86423 :: prose twin, two significant figures
\newcommand{\MlstLpJcTwoQgMeanAccDropFone}{52.7, 52.8\%}  % experiments/FIGS/data/v2/index.json :: linear probe JetClass-II, largest set: relative drop in accuracy (%) with the qg data's missing slots set to mean, FLAV_F1 :: sha256 c6c4138ee5c86423
\newcommand{\MlstLpJcTwoQgMeanAccDropFonePr}{53--53\%}  % experiments/FIGS/data/v2/index.json :: linear probe JetClass-II, largest set: relative drop in accuracy (%) with the qg data's missing slots set to mean, FLAV_F1 :: sha256 c6c4138ee5c86423 :: prose twin, two significant figures
\newcommand{\MlstLpJcTwoQgDrawAccDropFone}{55.6, 56.5\%}  % experiments/FIGS/data/v2/index.json :: linear probe JetClass-II, largest set: relative drop in accuracy (%) with the qg data's missing slots set to draws, FLAV_F1 :: sha256 c6c4138ee5c86423
\newcommand{\MlstLpJcTwoQgDrawAccDropFonePr}{56--56\%}  % experiments/FIGS/data/v2/index.json :: linear probe JetClass-II, largest set: relative drop in accuracy (%) with the qg data's missing slots set to draws, FLAV_F1 :: sha256 c6c4138ee5c86423 :: prose twin, two significant figures
\newcommand{\MlstLpJcTwoQgZeroAccDropFoner}{67.2, 69.0\%}  % experiments/FIGS/data/v2/index.json :: linear probe JetClass-II, largest set: relative drop in accuracy (%) with the qg data's missing slots set to zero, FLAV_F1R :: sha256 c6c4138ee5c86423
\newcommand{\MlstLpJcTwoQgZeroAccDropFonerPr}{67--69\%}  % experiments/FIGS/data/v2/index.json :: linear probe JetClass-II, largest set: relative drop in accuracy (%) with the qg data's missing slots set to zero, FLAV_F1R :: sha256 c6c4138ee5c86423 :: prose twin, two significant figures
\newcommand{\MlstLpJcTwoQgMeanAccDropFoner}{52.9, 51.8\%}  % experiments/FIGS/data/v2/index.json :: linear probe JetClass-II, largest set: relative drop in accuracy (%) with the qg data's missing slots set to mean, FLAV_F1R :: sha256 c6c4138ee5c86423
\newcommand{\MlstLpJcTwoQgMeanAccDropFonerPr}{52--53\%}  % experiments/FIGS/data/v2/index.json :: linear probe JetClass-II, largest set: relative drop in accuracy (%) with the qg data's missing slots set to mean, FLAV_F1R :: sha256 c6c4138ee5c86423 :: prose twin, two significant figures
\newcommand{\MlstLpJcTwoQgDrawAccDropFoner}{57.2, 55.4\%}  % experiments/FIGS/data/v2/index.json :: linear probe JetClass-II, largest set: relative drop in accuracy (%) with the qg data's missing slots set to draws, FLAV_F1R :: sha256 c6c4138ee5c86423
\newcommand{\MlstLpJcTwoQgDrawAccDropFonerPr}{55--57\%}  % experiments/FIGS/data/v2/index.json :: linear probe JetClass-II, largest set: relative drop in accuracy (%) with the qg data's missing slots set to draws, FLAV_F1R :: sha256 c6c4138ee5c86423 :: prose twin, two significant figures
\newcommand{\MlstLpJcTwoQgZeroAccDropInitBn}{8.7\,\ensuremath{\pm}\,2.6\%}  % experiments/FIGS/data/v2/index.json :: linear probe JetClass-II, largest set: relative drop in accuracy (%) with the qg data's missing slots set to zero, INIT_BN :: sha256 c6c4138ee5c86423
\newcommand{\MlstLpJcTwoQgZeroAccDropInitBnPr}{8.7\%}  % experiments/FIGS/data/v2/index.json :: linear probe JetClass-II, largest set: relative drop in accuracy (%) with the qg data's missing slots set to zero, INIT_BN :: sha256 c6c4138ee5c86423 :: prose twin, two significant figures
\newcommand{\MlstLpJcTwoQgMeanAccDropInitBn}{10.4\,\ensuremath{\pm}\,1.4\%}  % experiments/FIGS/data/v2/index.json :: linear probe JetClass-II, largest set: relative drop in accuracy (%) with the qg data's missing slots set to mean, INIT_BN :: sha256 c6c4138ee5c86423
\newcommand{\MlstLpJcTwoQgMeanAccDropInitBnPr}{10\%}  % experiments/FIGS/data/v2/index.json :: linear probe JetClass-II, largest set: relative drop in accuracy (%) with the qg data's missing slots set to mean, INIT_BN :: sha256 c6c4138ee5c86423 :: prose twin, two significant figures
\newcommand{\MlstLpJcTwoQgDrawAccDropInitBn}{11.4\,\ensuremath{\pm}\,2.1\%}  % experiments/FIGS/data/v2/index.json :: linear probe JetClass-II, largest set: relative drop in accuracy (%) with the qg data's missing slots set to draws, INIT_BN :: sha256 c6c4138ee5c86423
\newcommand{\MlstLpJcTwoQgDrawAccDropInitBnPr}{11\%}  % experiments/FIGS/data/v2/index.json :: linear probe JetClass-II, largest set: relative drop in accuracy (%) with the qg data's missing slots set to draws, INIT_BN :: sha256 c6c4138ee5c86423 :: prose twin, two significant figures
\newcommand{\MlstMassLambda}{5}  % configs/arms/v2_grid.json :: arms[L162_MASS, R16_Q1_MASS].mass_lambda :: sha256 d8a31b1486992302 :: weight of the mass-regression term in the pretraining loss, both mass models
\newcommand{\MlstMassLossShareResMass}{13.78 and 13.87}  % experiments/FIGS/data/v2/index.json :: share of the mass-regression term in the training loss of L162_MASS, mean over epochs, percent :: sha256 c6c4138ee5c86423
\newcommand{\MlstMassLossShareCoarseMass}{28.65 and 28.67}  % experiments/FIGS/data/v2/index.json :: share of the mass-regression term in the training loss of R16_Q1_MASS, mean over epochs, percent :: sha256 c6c4138ee5c86423
\newcommand{\MlstMassUncorrected}{0.245}  % experiments/FIGS/data/v2/index.json :: head/atlas_res: groomed mass, all jets, ATLAS resolution :: sha256 c6c4138ee5c86423
\newcommand{\MlstMassUncorrectedResponse}{0.929}  % experiments/FIGS/data/v2/index.json :: head/atlas_res: groomed mass, all jets, median response :: sha256 c6c4138ee5c86423
\newcommand{\MlstMassHeadResMass}{0.1471 and 0.1463}  % experiments/FIGS/data/v2/index.json :: regression output of L162_MASS, all jets, ATLAS resolution :: sha256 c6c4138ee5c86423
\newcommand{\MlstMassHeadResMassResponse}{0.995 and 1.003}  % experiments/FIGS/data/v2/index.json :: regression output of L162_MASS, all jets, median response :: sha256 c6c4138ee5c86423
\newcommand{\MlstMassHeadResMassOverUncorr}{0.6009 and 0.5976}  % experiments/FIGS/data/v2/index.json :: resolution of the regression output of L162_MASS over the groomed mass's, all jets :: sha256 c6c4138ee5c86423
\newcommand{\MlstMassHeadResMassOverUncorrPr}{at most 0.60 times the resolution}  % experiments/FIGS/data/v2/index.json :: resolution of the regression output of L162_MASS over the groomed mass's, all jets :: sha256 c6c4138ee5c86423 :: prose twin: two significant figures, direction in words
\newcommand{\MlstMassHeadResMassOverUncorrPrMax}{down to 0.60 times the resolution}  % experiments/FIGS/data/v2/index.json :: resolution of the regression output of L162_MASS over the groomed mass's, all jets :: sha256 c6c4138ee5c86423 :: prose twin of the run farther from 1
\newcommand{\MlstMassHeadResMassOverMlpRes}{0.6645 and 0.6691}  % experiments/FIGS/data/v2/index.json :: resolution of the regression output of L162_MASS over the frozen-feature MLP of L162, run by run, all jets :: sha256 c6c4138ee5c86423
\newcommand{\MlstMassHeadResMassOverMlpResPr}{at most 0.67 times the resolution}  % experiments/FIGS/data/v2/index.json :: resolution of the regression output of L162_MASS over the frozen-feature MLP of L162, run by run, all jets :: sha256 c6c4138ee5c86423 :: prose twin: two significant figures, direction in words
\newcommand{\MlstMassHeadResMassOverMlpResPrMax}{down to 0.66 times the resolution}  % experiments/FIGS/data/v2/index.json :: resolution of the regression output of L162_MASS over the frozen-feature MLP of L162, run by run, all jets :: sha256 c6c4138ee5c86423 :: prose twin of the run farther from 1
\newcommand{\MlstMassHeadCoarseMass}{0.1450 and 0.1461}  % experiments/FIGS/data/v2/index.json :: regression output of R16_Q1_MASS, all jets, ATLAS resolution :: sha256 c6c4138ee5c86423
\newcommand{\MlstMassHeadCoarseMassResponse}{1.005 and 1.005}  % experiments/FIGS/data/v2/index.json :: regression output of R16_Q1_MASS, all jets, median response :: sha256 c6c4138ee5c86423
\newcommand{\MlstMassHeadCoarseMassOverUncorr}{0.5921 and 0.5966}  % experiments/FIGS/data/v2/index.json :: resolution of the regression output of R16_Q1_MASS over the groomed mass's, all jets :: sha256 c6c4138ee5c86423
\newcommand{\MlstMassHeadCoarseMassOverUncorrPr}{at most 0.60 times the resolution}  % experiments/FIGS/data/v2/index.json :: resolution of the regression output of R16_Q1_MASS over the groomed mass's, all jets :: sha256 c6c4138ee5c86423 :: prose twin: two significant figures, direction in words
\newcommand{\MlstMassHeadCoarseMassOverUncorrPrMax}{down to 0.59 times the resolution}  % experiments/FIGS/data/v2/index.json :: resolution of the regression output of R16_Q1_MASS over the groomed mass's, all jets :: sha256 c6c4138ee5c86423 :: prose twin of the run farther from 1
\newcommand{\MlstMassHeadCoarseMassOverMlpCoarse}{0.703 and 0.690}  % experiments/FIGS/data/v2/index.json :: resolution of the regression output of R16_Q1_MASS over the frozen-feature MLP of R16_Q1, run by run, all jets :: sha256 c6c4138ee5c86423
\newcommand{\MlstMassHeadCoarseMassOverMlpCoarsePr}{at most 0.70 times the resolution}  % experiments/FIGS/data/v2/index.json :: resolution of the regression output of R16_Q1_MASS over the frozen-feature MLP of R16_Q1, run by run, all jets :: sha256 c6c4138ee5c86423 :: prose twin: two significant figures, direction in words
\newcommand{\MlstMassHeadCoarseMassOverMlpCoarsePrMax}{down to 0.69 times the resolution}  % experiments/FIGS/data/v2/index.json :: resolution of the regression output of R16_Q1_MASS over the frozen-feature MLP of R16_Q1, run by run, all jets :: sha256 c6c4138ee5c86423 :: prose twin of the run farther from 1
\newcommand{\MlstMassMlpInit}{0.2421\,\ensuremath{\pm}\,0.0018}  % experiments/FIGS/data/v2/index.json :: frozen-feature MLP of INIT, all jets, ATLAS resolution :: sha256 c6c4138ee5c86423
\newcommand{\MlstMassMlpFull}{0.2200\,\ensuremath{\pm}\,0.0014}  % experiments/FIGS/data/v2/index.json :: frozen-feature MLP of L188, all jets, ATLAS resolution :: sha256 c6c4138ee5c86423
\newcommand{\MlstMassMlpRes}{0.2195\,\ensuremath{\pm}\,0.0017}  % experiments/FIGS/data/v2/index.json :: frozen-feature MLP of L162, all jets, ATLAS resolution :: sha256 c6c4138ee5c86423
\newcommand{\MlstMassMlpMid}{0.2158\,\ensuremath{\pm}\,0.0029}  % experiments/FIGS/data/v2/index.json :: frozen-feature MLP of R42_Q1, all jets, ATLAS resolution :: sha256 c6c4138ee5c86423
\newcommand{\MlstMassMlpCoarse}{0.2098\,\ensuremath{\pm}\,0.0031}  % experiments/FIGS/data/v2/index.json :: frozen-feature MLP of R16_Q1, all jets, ATLAS resolution :: sha256 c6c4138ee5c86423
\newcommand{\MlstMassMlpResMass}{0.1591 and 0.1586}  % experiments/FIGS/data/v2/index.json :: frozen-feature MLP of L162_MASS, all jets, ATLAS resolution :: sha256 c6c4138ee5c86423
\newcommand{\MlstMassMlpCoarseMass}{0.1583 and 0.1588}  % experiments/FIGS/data/v2/index.json :: frozen-feature MLP of R16_Q1_MASS, all jets, ATLAS resolution :: sha256 c6c4138ee5c86423
\newcommand{\MlstCmpMlpResOverFull}{0.998\,\ensuremath{\pm}\,0.006}  % experiments/FIGS/data/v2/index.json :: frozen-feature MLP resolution of L162 over L188, run by run, all jets :: sha256 c6c4138ee5c86423
\newcommand{\MlstCmpMlpResOverFullPr}{equal within 1\%}  % experiments/FIGS/data/v2/index.json :: frozen-feature MLP resolution of L162 over L188, run by run, all jets :: sha256 c6c4138ee5c86423 :: prose twin: two significant figures, direction in words
\newcommand{\MlstCmpMlpMidOverFull}{0.981\,\ensuremath{\pm}\,0.017}  % experiments/FIGS/data/v2/index.json :: frozen-feature MLP resolution of R42_Q1 over L188, run by run, all jets :: sha256 c6c4138ee5c86423
\newcommand{\MlstCmpMlpMidOverFullPr}{equal within 2\%}  % experiments/FIGS/data/v2/index.json :: frozen-feature MLP resolution of R42_Q1 over L188, run by run, all jets :: sha256 c6c4138ee5c86423 :: prose twin: two significant figures, direction in words
\newcommand{\MlstCmpMlpCoarseOverFull}{0.954\,\ensuremath{\pm}\,0.018}  % experiments/FIGS/data/v2/index.json :: frozen-feature MLP resolution of R16_Q1 over L188, run by run, all jets :: sha256 c6c4138ee5c86423
\newcommand{\MlstCmpMlpCoarseOverFullPr}{0.95 times the resolution}  % experiments/FIGS/data/v2/index.json :: frozen-feature MLP resolution of R16_Q1 over L188, run by run, all jets :: sha256 c6c4138ee5c86423 :: prose twin: two significant figures, direction in words
\newcommand{\MlstCmpMlpResMassOverRes}{0.7186 and 0.7253}  % experiments/FIGS/data/v2/index.json :: frozen-feature MLP resolution of L162_MASS over L162, run by run, all jets :: sha256 c6c4138ee5c86423
\newcommand{\MlstCmpMlpResMassOverResPr}{at most 0.73 times the resolution}  % experiments/FIGS/data/v2/index.json :: frozen-feature MLP resolution of L162_MASS over L162, run by run, all jets :: sha256 c6c4138ee5c86423 :: prose twin: two significant figures, direction in words
\newcommand{\MlstCmpMlpResMassOverResPrMax}{down to 0.72 times the resolution}  % experiments/FIGS/data/v2/index.json :: frozen-feature MLP resolution of L162_MASS over L162, run by run, all jets :: sha256 c6c4138ee5c86423 :: prose twin of the run farther from 1
\newcommand{\MlstCmpMlpCoarseMassOverCoarse}{0.768 and 0.750}  % experiments/FIGS/data/v2/index.json :: frozen-feature MLP resolution of R16_Q1_MASS over R16_Q1, run by run, all jets :: sha256 c6c4138ee5c86423
\newcommand{\MlstCmpMlpCoarseMassOverCoarsePr}{at most 0.77 times the resolution}  % experiments/FIGS/data/v2/index.json :: frozen-feature MLP resolution of R16_Q1_MASS over R16_Q1, run by run, all jets :: sha256 c6c4138ee5c86423 :: prose twin: two significant figures, direction in words
\newcommand{\MlstCmpMlpCoarseMassOverCoarsePrMax}{down to 0.75 times the resolution}  % experiments/FIGS/data/v2/index.json :: frozen-feature MLP resolution of R16_Q1_MASS over R16_Q1, run by run, all jets :: sha256 c6c4138ee5c86423 :: prose twin of the run farther from 1
\newcommand{\MlstMassUncorrectedResonant}{0.213}  % experiments/FIGS/data/v2/index.json :: head/atlas_res: groomed mass, resonant jets, ATLAS resolution :: sha256 c6c4138ee5c86423
\newcommand{\MlstMassUncorrectedResponseResonant}{0.929}  % experiments/FIGS/data/v2/index.json :: head/atlas_res: groomed mass, resonant jets, median response :: sha256 c6c4138ee5c86423
\newcommand{\MlstMassHeadResMassResonant}{0.1332 and 0.1320}  % experiments/FIGS/data/v2/index.json :: regression output of L162_MASS, resonant jets, ATLAS resolution :: sha256 c6c4138ee5c86423
\newcommand{\MlstMassHeadResMassResponseResonant}{0.996 and 1.004}  % experiments/FIGS/data/v2/index.json :: regression output of L162_MASS, resonant jets, median response :: sha256 c6c4138ee5c86423
\newcommand{\MlstMassHeadResMassOverUncorrResonant}{0.6259 and 0.6205}  % experiments/FIGS/data/v2/index.json :: resolution of the regression output of L162_MASS over the groomed mass's, resonant jets :: sha256 c6c4138ee5c86423
\newcommand{\MlstMassHeadResMassOverUncorrResonantPr}{at most 0.63 times the resolution}  % experiments/FIGS/data/v2/index.json :: resolution of the regression output of L162_MASS over the groomed mass's, resonant jets :: sha256 c6c4138ee5c86423 :: prose twin: two significant figures, direction in words
\newcommand{\MlstMassHeadResMassOverUncorrResonantPrMax}{down to 0.62 times the resolution}  % experiments/FIGS/data/v2/index.json :: resolution of the regression output of L162_MASS over the groomed mass's, resonant jets :: sha256 c6c4138ee5c86423 :: prose twin of the run farther from 1
\newcommand{\MlstMassHeadResMassOverMlpResResonant}{0.6620 and 0.6662}  % experiments/FIGS/data/v2/index.json :: resolution of the regression output of L162_MASS over the frozen-feature MLP of L162, run by run, resonant jets :: sha256 c6c4138ee5c86423
\newcommand{\MlstMassHeadResMassOverMlpResResonantPr}{at most 0.67 times the resolution}  % experiments/FIGS/data/v2/index.json :: resolution of the regression output of L162_MASS over the frozen-feature MLP of L162, run by run, resonant jets :: sha256 c6c4138ee5c86423 :: prose twin: two significant figures, direction in words
\newcommand{\MlstMassHeadResMassOverMlpResResonantPrMax}{down to 0.66 times the resolution}  % experiments/FIGS/data/v2/index.json :: resolution of the regression output of L162_MASS over the frozen-feature MLP of L162, run by run, resonant jets :: sha256 c6c4138ee5c86423 :: prose twin of the run farther from 1
\newcommand{\MlstMassHeadCoarseMassResonant}{0.1312 and 0.1320}  % experiments/FIGS/data/v2/index.json :: regression output of R16_Q1_MASS, resonant jets, ATLAS resolution :: sha256 c6c4138ee5c86423
\newcommand{\MlstMassHeadCoarseMassResponseResonant}{1.006 and 1.006}  % experiments/FIGS/data/v2/index.json :: regression output of R16_Q1_MASS, resonant jets, median response :: sha256 c6c4138ee5c86423
\newcommand{\MlstMassHeadCoarseMassOverUncorrResonant}{0.6165 and 0.6206}  % experiments/FIGS/data/v2/index.json :: resolution of the regression output of R16_Q1_MASS over the groomed mass's, resonant jets :: sha256 c6c4138ee5c86423
\newcommand{\MlstMassHeadCoarseMassOverUncorrResonantPr}{at most 0.62 times the resolution}  % experiments/FIGS/data/v2/index.json :: resolution of the regression output of R16_Q1_MASS over the groomed mass's, resonant jets :: sha256 c6c4138ee5c86423 :: prose twin: two significant figures, direction in words
\newcommand{\MlstMassHeadCoarseMassOverUncorrResonantPrMax}{down to 0.62 times the resolution}  % experiments/FIGS/data/v2/index.json :: resolution of the regression output of R16_Q1_MASS over the groomed mass's, resonant jets :: sha256 c6c4138ee5c86423 :: prose twin of the run farther from 1
\newcommand{\MlstMassHeadCoarseMassOverMlpCoarseResonant}{0.707 and 0.688}  % experiments/FIGS/data/v2/index.json :: resolution of the regression output of R16_Q1_MASS over the frozen-feature MLP of R16_Q1, run by run, resonant jets :: sha256 c6c4138ee5c86423
\newcommand{\MlstMassHeadCoarseMassOverMlpCoarseResonantPr}{at most 0.71 times the resolution}  % experiments/FIGS/data/v2/index.json :: resolution of the regression output of R16_Q1_MASS over the frozen-feature MLP of R16_Q1, run by run, resonant jets :: sha256 c6c4138ee5c86423 :: prose twin: two significant figures, direction in words
\newcommand{\MlstMassHeadCoarseMassOverMlpCoarseResonantPrMax}{down to 0.69 times the resolution}  % experiments/FIGS/data/v2/index.json :: resolution of the regression output of R16_Q1_MASS over the frozen-feature MLP of R16_Q1, run by run, resonant jets :: sha256 c6c4138ee5c86423 :: prose twin of the run farther from 1
\newcommand{\MlstMassMlpInitResonant}{0.2214\,\ensuremath{\pm}\,0.0017}  % experiments/FIGS/data/v2/index.json :: frozen-feature MLP of INIT, resonant jets, ATLAS resolution :: sha256 c6c4138ee5c86423
\newcommand{\MlstMassMlpFullResonant}{0.1990\,\ensuremath{\pm}\,0.0016}  % experiments/FIGS/data/v2/index.json :: frozen-feature MLP of L188, resonant jets, ATLAS resolution :: sha256 c6c4138ee5c86423
\newcommand{\MlstMassMlpResResonant}{0.1989\,\ensuremath{\pm}\,0.0020}  % experiments/FIGS/data/v2/index.json :: frozen-feature MLP of L162, resonant jets, ATLAS resolution :: sha256 c6c4138ee5c86423
\newcommand{\MlstMassMlpMidResonant}{0.1949\,\ensuremath{\pm}\,0.0023}  % experiments/FIGS/data/v2/index.json :: frozen-feature MLP of R42_Q1, resonant jets, ATLAS resolution :: sha256 c6c4138ee5c86423
\newcommand{\MlstMassMlpCoarseResonant}{0.1894\,\ensuremath{\pm}\,0.0035}  % experiments/FIGS/data/v2/index.json :: frozen-feature MLP of R16_Q1, resonant jets, ATLAS resolution :: sha256 c6c4138ee5c86423
\newcommand{\MlstMassMlpResMassResonant}{0.1436 and 0.1430}  % experiments/FIGS/data/v2/index.json :: frozen-feature MLP of L162_MASS, resonant jets, ATLAS resolution :: sha256 c6c4138ee5c86423
\newcommand{\MlstMassMlpCoarseMassResonant}{0.1420 and 0.1431}  % experiments/FIGS/data/v2/index.json :: frozen-feature MLP of R16_Q1_MASS, resonant jets, ATLAS resolution :: sha256 c6c4138ee5c86423
\newcommand{\MlstCmpMlpResOverFullResonant}{1.000\,\ensuremath{\pm}\,0.005}  % experiments/FIGS/data/v2/index.json :: frozen-feature MLP resolution of L162 over L188, run by run, resonant jets :: sha256 c6c4138ee5c86423
\newcommand{\MlstCmpMlpResOverFullResonantPr}{equal within 1\%}  % experiments/FIGS/data/v2/index.json :: frozen-feature MLP resolution of L162 over L188, run by run, resonant jets :: sha256 c6c4138ee5c86423 :: prose twin: two significant figures, direction in words
\newcommand{\MlstCmpMlpMidOverFullResonant}{0.980\,\ensuremath{\pm}\,0.018}  % experiments/FIGS/data/v2/index.json :: frozen-feature MLP resolution of R42_Q1 over L188, run by run, resonant jets :: sha256 c6c4138ee5c86423
\newcommand{\MlstCmpMlpMidOverFullResonantPr}{equal within 3\%}  % experiments/FIGS/data/v2/index.json :: frozen-feature MLP resolution of R42_Q1 over L188, run by run, resonant jets :: sha256 c6c4138ee5c86423 :: prose twin: two significant figures, direction in words
\newcommand{\MlstCmpMlpCoarseOverFullResonant}{0.952\,\ensuremath{\pm}\,0.023}  % experiments/FIGS/data/v2/index.json :: frozen-feature MLP resolution of R16_Q1 over L188, run by run, resonant jets :: sha256 c6c4138ee5c86423
\newcommand{\MlstCmpMlpCoarseOverFullResonantPr}{0.95 times the resolution}  % experiments/FIGS/data/v2/index.json :: frozen-feature MLP resolution of R16_Q1 over L188, run by run, resonant jets :: sha256 c6c4138ee5c86423 :: prose twin: two significant figures, direction in words
\newcommand{\MlstCmpMlpResMassOverResResonant}{0.7137 and 0.7219}  % experiments/FIGS/data/v2/index.json :: frozen-feature MLP resolution of L162_MASS over L162, run by run, resonant jets :: sha256 c6c4138ee5c86423
\newcommand{\MlstCmpMlpResMassOverResResonantPr}{at most 0.72 times the resolution}  % experiments/FIGS/data/v2/index.json :: frozen-feature MLP resolution of L162_MASS over L162, run by run, resonant jets :: sha256 c6c4138ee5c86423 :: prose twin: two significant figures, direction in words
\newcommand{\MlstCmpMlpResMassOverResResonantPrMax}{down to 0.71 times the resolution}  % experiments/FIGS/data/v2/index.json :: frozen-feature MLP resolution of L162_MASS over L162, run by run, resonant jets :: sha256 c6c4138ee5c86423 :: prose twin of the run farther from 1
\newcommand{\MlstCmpMlpCoarseMassOverCoarseResonant}{0.766 and 0.745}  % experiments/FIGS/data/v2/index.json :: frozen-feature MLP resolution of R16_Q1_MASS over R16_Q1, run by run, resonant jets :: sha256 c6c4138ee5c86423
\newcommand{\MlstCmpMlpCoarseMassOverCoarseResonantPr}{at most 0.77 times the resolution}  % experiments/FIGS/data/v2/index.json :: frozen-feature MLP resolution of R16_Q1_MASS over R16_Q1, run by run, resonant jets :: sha256 c6c4138ee5c86423 :: prose twin: two significant figures, direction in words
\newcommand{\MlstCmpMlpCoarseMassOverCoarseResonantPrMax}{down to 0.75 times the resolution}  % experiments/FIGS/data/v2/index.json :: frozen-feature MLP resolution of R16_Q1_MASS over R16_Q1, run by run, resonant jets :: sha256 c6c4138ee5c86423 :: prose twin of the run farther from 1
\newcommand{\MlstMassUncorrectedQcd}{1.52}  % experiments/FIGS/data/v2/index.json :: head/atlas_res: groomed mass, qcd jets, ATLAS resolution :: sha256 c6c4138ee5c86423
\newcommand{\MlstMassUncorrectedResponseQcd}{0.931}  % experiments/FIGS/data/v2/index.json :: head/atlas_res: groomed mass, qcd jets, median response :: sha256 c6c4138ee5c86423
\newcommand{\MlstMassHeadResMassQcd}{0.2263 and 0.2256}  % experiments/FIGS/data/v2/index.json :: regression output of L162_MASS, qcd jets, ATLAS resolution :: sha256 c6c4138ee5c86423
\newcommand{\MlstMassHeadResMassResponseQcd}{0.9897 and 0.9996}  % experiments/FIGS/data/v2/index.json :: regression output of L162_MASS, qcd jets, median response :: sha256 c6c4138ee5c86423
\newcommand{\MlstMassHeadResMassOverUncorrQcd}{0.1485 and 0.1480}  % experiments/FIGS/data/v2/index.json :: resolution of the regression output of L162_MASS over the groomed mass's, qcd jets :: sha256 c6c4138ee5c86423
\newcommand{\MlstMassHeadResMassOverUncorrQcdPr}{at most 0.15 times the resolution}  % experiments/FIGS/data/v2/index.json :: resolution of the regression output of L162_MASS over the groomed mass's, qcd jets :: sha256 c6c4138ee5c86423 :: prose twin: two significant figures, direction in words
\newcommand{\MlstMassHeadResMassOverUncorrQcdPrMax}{down to 0.15 times the resolution}  % experiments/FIGS/data/v2/index.json :: resolution of the regression output of L162_MASS over the groomed mass's, qcd jets :: sha256 c6c4138ee5c86423 :: prose twin of the run farther from 1
\newcommand{\MlstMassHeadResMassOverMlpResQcd}{0.6542 and 0.6632}  % experiments/FIGS/data/v2/index.json :: resolution of the regression output of L162_MASS over the frozen-feature MLP of L162, run by run, qcd jets :: sha256 c6c4138ee5c86423
\newcommand{\MlstMassHeadResMassOverMlpResQcdPr}{at most 0.66 times the resolution}  % experiments/FIGS/data/v2/index.json :: resolution of the regression output of L162_MASS over the frozen-feature MLP of L162, run by run, qcd jets :: sha256 c6c4138ee5c86423 :: prose twin: two significant figures, direction in words
\newcommand{\MlstMassHeadResMassOverMlpResQcdPrMax}{down to 0.65 times the resolution}  % experiments/FIGS/data/v2/index.json :: resolution of the regression output of L162_MASS over the frozen-feature MLP of L162, run by run, qcd jets :: sha256 c6c4138ee5c86423 :: prose twin of the run farther from 1
\newcommand{\MlstMassHeadCoarseMassQcd}{0.2239 and 0.2240}  % experiments/FIGS/data/v2/index.json :: regression output of R16_Q1_MASS, qcd jets, ATLAS resolution :: sha256 c6c4138ee5c86423
\newcommand{\MlstMassHeadCoarseMassResponseQcd}{1.001 and 1.000}  % experiments/FIGS/data/v2/index.json :: regression output of R16_Q1_MASS, qcd jets, median response :: sha256 c6c4138ee5c86423
\newcommand{\MlstMassHeadCoarseMassOverUncorrQcd}{0.1469 and 0.1469}  % experiments/FIGS/data/v2/index.json :: resolution of the regression output of R16_Q1_MASS over the groomed mass's, qcd jets :: sha256 c6c4138ee5c86423
\newcommand{\MlstMassHeadCoarseMassOverUncorrQcdPr}{at most 0.15 times the resolution}  % experiments/FIGS/data/v2/index.json :: resolution of the regression output of R16_Q1_MASS over the groomed mass's, qcd jets :: sha256 c6c4138ee5c86423 :: prose twin: two significant figures, direction in words
\newcommand{\MlstMassHeadCoarseMassOverUncorrQcdPrMax}{down to 0.15 times the resolution}  % experiments/FIGS/data/v2/index.json :: resolution of the regression output of R16_Q1_MASS over the groomed mass's, qcd jets :: sha256 c6c4138ee5c86423 :: prose twin of the run farther from 1
\newcommand{\MlstMassHeadCoarseMassOverMlpCoarseQcd}{0.679 and 0.698}  % experiments/FIGS/data/v2/index.json :: resolution of the regression output of R16_Q1_MASS over the frozen-feature MLP of R16_Q1, run by run, qcd jets :: sha256 c6c4138ee5c86423
\newcommand{\MlstMassHeadCoarseMassOverMlpCoarseQcdPr}{at most 0.70 times the resolution}  % experiments/FIGS/data/v2/index.json :: resolution of the regression output of R16_Q1_MASS over the frozen-feature MLP of R16_Q1, run by run, qcd jets :: sha256 c6c4138ee5c86423 :: prose twin: two significant figures, direction in words
\newcommand{\MlstMassHeadCoarseMassOverMlpCoarseQcdPrMax}{down to 0.68 times the resolution}  % experiments/FIGS/data/v2/index.json :: resolution of the regression output of R16_Q1_MASS over the frozen-feature MLP of R16_Q1, run by run, qcd jets :: sha256 c6c4138ee5c86423 :: prose twin of the run farther from 1
\newcommand{\MlstMassMlpInitQcd}{0.502\,\ensuremath{\pm}\,0.014}  % experiments/FIGS/data/v2/index.json :: frozen-feature MLP of INIT, qcd jets, ATLAS resolution :: sha256 c6c4138ee5c86423
\newcommand{\MlstMassMlpFullQcd}{0.3487\,\ensuremath{\pm}\,0.0029}  % experiments/FIGS/data/v2/index.json :: frozen-feature MLP of L188, qcd jets, ATLAS resolution :: sha256 c6c4138ee5c86423
\newcommand{\MlstMassMlpResQcd}{0.3425\,\ensuremath{\pm}\,0.0031}  % experiments/FIGS/data/v2/index.json :: frozen-feature MLP of L162, qcd jets, ATLAS resolution :: sha256 c6c4138ee5c86423
\newcommand{\MlstMassMlpMidQcd}{0.338\,\ensuremath{\pm}\,0.007}  % experiments/FIGS/data/v2/index.json :: frozen-feature MLP of R42_Q1, qcd jets, ATLAS resolution :: sha256 c6c4138ee5c86423
\newcommand{\MlstMassMlpCoarseQcd}{0.327\,\ensuremath{\pm}\,0.006}  % experiments/FIGS/data/v2/index.json :: frozen-feature MLP of R16_Q1, qcd jets, ATLAS resolution :: sha256 c6c4138ee5c86423
\newcommand{\MlstMassMlpResMassQcd}{0.2473 and 0.2452}  % experiments/FIGS/data/v2/index.json :: frozen-feature MLP of L162_MASS, qcd jets, ATLAS resolution :: sha256 c6c4138ee5c86423
\newcommand{\MlstMassMlpCoarseMassQcd}{0.2465 and 0.2463}  % experiments/FIGS/data/v2/index.json :: frozen-feature MLP of R16_Q1_MASS, qcd jets, ATLAS resolution :: sha256 c6c4138ee5c86423
\newcommand{\MlstCmpMlpResOverFullQcd}{0.982\,\ensuremath{\pm}\,0.017}  % experiments/FIGS/data/v2/index.json :: frozen-feature MLP resolution of L162 over L188, run by run, qcd jets :: sha256 c6c4138ee5c86423
\newcommand{\MlstCmpMlpResOverFullQcdPr}{equal within 2\%}  % experiments/FIGS/data/v2/index.json :: frozen-feature MLP resolution of L162 over L188, run by run, qcd jets :: sha256 c6c4138ee5c86423 :: prose twin: two significant figures, direction in words
\newcommand{\MlstCmpMlpMidOverFullQcd}{0.969\,\ensuremath{\pm}\,0.013}  % experiments/FIGS/data/v2/index.json :: frozen-feature MLP resolution of R42_Q1 over L188, run by run, qcd jets :: sha256 c6c4138ee5c86423
\newcommand{\MlstCmpMlpMidOverFullQcdPr}{0.97 times the resolution}  % experiments/FIGS/data/v2/index.json :: frozen-feature MLP resolution of R42_Q1 over L188, run by run, qcd jets :: sha256 c6c4138ee5c86423 :: prose twin: two significant figures, direction in words
\newcommand{\MlstCmpMlpCoarseOverFullQcd}{0.939\,\ensuremath{\pm}\,0.023}  % experiments/FIGS/data/v2/index.json :: frozen-feature MLP resolution of R16_Q1 over L188, run by run, qcd jets :: sha256 c6c4138ee5c86423
\newcommand{\MlstCmpMlpCoarseOverFullQcdPr}{0.94 times the resolution}  % experiments/FIGS/data/v2/index.json :: frozen-feature MLP resolution of R16_Q1 over L188, run by run, qcd jets :: sha256 c6c4138ee5c86423 :: prose twin: two significant figures, direction in words
\newcommand{\MlstCmpMlpResMassOverResQcd}{0.7147 and 0.7209}  % experiments/FIGS/data/v2/index.json :: frozen-feature MLP resolution of L162_MASS over L162, run by run, qcd jets :: sha256 c6c4138ee5c86423
\newcommand{\MlstCmpMlpResMassOverResQcdPr}{at most 0.72 times the resolution}  % experiments/FIGS/data/v2/index.json :: frozen-feature MLP resolution of L162_MASS over L162, run by run, qcd jets :: sha256 c6c4138ee5c86423 :: prose twin: two significant figures, direction in words
\newcommand{\MlstCmpMlpResMassOverResQcdPrMax}{down to 0.71 times the resolution}  % experiments/FIGS/data/v2/index.json :: frozen-feature MLP resolution of L162_MASS over L162, run by run, qcd jets :: sha256 c6c4138ee5c86423 :: prose twin of the run farther from 1
\newcommand{\MlstCmpMlpCoarseMassOverCoarseQcd}{0.747 and 0.768}  % experiments/FIGS/data/v2/index.json :: frozen-feature MLP resolution of R16_Q1_MASS over R16_Q1, run by run, qcd jets :: sha256 c6c4138ee5c86423
\newcommand{\MlstCmpMlpCoarseMassOverCoarseQcdPr}{at most 0.77 times the resolution}  % experiments/FIGS/data/v2/index.json :: frozen-feature MLP resolution of R16_Q1_MASS over R16_Q1, run by run, qcd jets :: sha256 c6c4138ee5c86423 :: prose twin: two significant figures, direction in words
\newcommand{\MlstCmpMlpCoarseMassOverCoarseQcdPrMax}{down to 0.75 times the resolution}  % experiments/FIGS/data/v2/index.json :: frozen-feature MLP resolution of R16_Q1_MASS over R16_Q1, run by run, qcd jets :: sha256 c6c4138ee5c86423 :: prose twin of the run farther from 1
\newcommand{\MlstMassRespAxesShareQcd}{76\%}  % experiments/FIGS/data/v2/index.json :: head/atlas_res: share of QCD jets whose groomed-mass response lies inside F9a(b)'s x axis :: sha256 c6c4138ee5c86423
\newcommand{\MlstSculptNPassTen}{41{,}028}  % experiments/FIGS/data/v2/index.json :: QCD test jets passing the cut at QCD efficiency 0.1, every run :: sha256 c6c4138ee5c86423
\newcommand{\MlstSculptFloorTen}{0.00944}  % experiments/FIGS/data/v2/index.json :: Jensen-Shannon distance expected for a random selection of as many QCD jets as pass at 0.1 :: sha256 c6c4138ee5c86423
\newcommand{\MlstSculptFullTen}{0.0524\,\ensuremath{\pm}\,0.0034}  % experiments/FIGS/data/v2/index.json :: Jensen-Shannon distance of the QCD m_SD distribution at QCD efficiency 0.1, L188 :: sha256 c6c4138ee5c86423
\newcommand{\MlstSculptFullOverFloorTen}{5.6\,\ensuremath{\pm}\,0.4}  % experiments/FIGS/data/v2/index.json :: distance at 0.1 over the random-selection floor, L188 :: sha256 c6c4138ee5c86423
\newcommand{\MlstSculptFullOverFloorTenPr}{5.6 times the distance}  % experiments/FIGS/data/v2/index.json :: distance at 0.1 over the random-selection floor, L188 :: sha256 c6c4138ee5c86423 :: prose twin: two significant figures, direction in words
\newcommand{\MlstSculptResTen}{0.0532\,\ensuremath{\pm}\,0.0009}  % experiments/FIGS/data/v2/index.json :: Jensen-Shannon distance of the QCD m_SD distribution at QCD efficiency 0.1, L162 :: sha256 c6c4138ee5c86423
\newcommand{\MlstSculptResOverFloorTen}{5.64\,\ensuremath{\pm}\,0.10}  % experiments/FIGS/data/v2/index.json :: distance at 0.1 over the random-selection floor, L162 :: sha256 c6c4138ee5c86423
\newcommand{\MlstSculptResOverFloorTenPr}{5.6 times the distance}  % experiments/FIGS/data/v2/index.json :: distance at 0.1 over the random-selection floor, L162 :: sha256 c6c4138ee5c86423 :: prose twin: two significant figures, direction in words
\newcommand{\MlstSculptMidTen}{0.0560\,\ensuremath{\pm}\,0.0027}  % experiments/FIGS/data/v2/index.json :: Jensen-Shannon distance of the QCD m_SD distribution at QCD efficiency 0.1, R42_Q1 :: sha256 c6c4138ee5c86423
\newcommand{\MlstSculptMidOverFloorTen}{5.94\,\ensuremath{\pm}\,0.28}  % experiments/FIGS/data/v2/index.json :: distance at 0.1 over the random-selection floor, R42_Q1 :: sha256 c6c4138ee5c86423
\newcommand{\MlstSculptMidOverFloorTenPr}{5.9 times the distance}  % experiments/FIGS/data/v2/index.json :: distance at 0.1 over the random-selection floor, R42_Q1 :: sha256 c6c4138ee5c86423 :: prose twin: two significant figures, direction in words
\newcommand{\MlstSculptCoarseTen}{0.0543\,\ensuremath{\pm}\,0.0010}  % experiments/FIGS/data/v2/index.json :: Jensen-Shannon distance of the QCD m_SD distribution at QCD efficiency 0.1, R16_Q1 :: sha256 c6c4138ee5c86423
\newcommand{\MlstSculptCoarseOverFloorTen}{5.76\,\ensuremath{\pm}\,0.10}  % experiments/FIGS/data/v2/index.json :: distance at 0.1 over the random-selection floor, R16_Q1 :: sha256 c6c4138ee5c86423
\newcommand{\MlstSculptCoarseOverFloorTenPr}{5.8 times the distance}  % experiments/FIGS/data/v2/index.json :: distance at 0.1 over the random-selection floor, R16_Q1 :: sha256 c6c4138ee5c86423 :: prose twin: two significant figures, direction in words
\newcommand{\MlstSculptResMassTen}{0.0422 and 0.0464}  % experiments/FIGS/data/v2/index.json :: Jensen-Shannon distance of the QCD m_SD distribution at QCD efficiency 0.1, L162_MASS :: sha256 c6c4138ee5c86423
\newcommand{\MlstSculptResMassOverFloorTen}{4.47 and 4.91}  % experiments/FIGS/data/v2/index.json :: distance at 0.1 over the random-selection floor, L162_MASS :: sha256 c6c4138ee5c86423
\newcommand{\MlstSculptResMassOverFloorTenPr}{at least 4.5 times the distance}  % experiments/FIGS/data/v2/index.json :: distance at 0.1 over the random-selection floor, L162_MASS :: sha256 c6c4138ee5c86423 :: prose twin: two significant figures, direction in words
\newcommand{\MlstSculptResMassOverFloorTenPrMax}{up to 4.9 times the distance}  % experiments/FIGS/data/v2/index.json :: distance at 0.1 over the random-selection floor, L162_MASS :: sha256 c6c4138ee5c86423 :: prose twin of the run farther from 1
\newcommand{\MlstSculptCoarseMassTen}{0.0522 and 0.0423}  % experiments/FIGS/data/v2/index.json :: Jensen-Shannon distance of the QCD m_SD distribution at QCD efficiency 0.1, R16_Q1_MASS :: sha256 c6c4138ee5c86423
\newcommand{\MlstSculptCoarseMassOverFloorTen}{5.5 and 4.5}  % experiments/FIGS/data/v2/index.json :: distance at 0.1 over the random-selection floor, R16_Q1_MASS :: sha256 c6c4138ee5c86423
\newcommand{\MlstSculptCoarseMassOverFloorTenPr}{at least 4.5 times the distance}  % experiments/FIGS/data/v2/index.json :: distance at 0.1 over the random-selection floor, R16_Q1_MASS :: sha256 c6c4138ee5c86423 :: prose twin: two significant figures, direction in words
\newcommand{\MlstSculptCoarseMassOverFloorTenPrMax}{up to 5.5 times the distance}  % experiments/FIGS/data/v2/index.json :: distance at 0.1 over the random-selection floor, R16_Q1_MASS :: sha256 c6c4138ee5c86423 :: prose twin of the run farther from 1
\newcommand{\MlstSculptResMassOverResTen}{0.789 and 0.860}  % experiments/FIGS/data/v2/index.json :: distance at QCD efficiency 0.1: L162_MASS over L162, run by run :: sha256 c6c4138ee5c86423
\newcommand{\MlstSculptResMassOverResTenPr}{at most 0.86 times the distance}  % experiments/FIGS/data/v2/index.json :: distance at QCD efficiency 0.1: L162_MASS over L162, run by run :: sha256 c6c4138ee5c86423 :: prose twin: two significant figures, direction in words
\newcommand{\MlstSculptResMassOverResTenPrMax}{down to 0.79 times the distance}  % experiments/FIGS/data/v2/index.json :: distance at QCD efficiency 0.1: L162_MASS over L162, run by run :: sha256 c6c4138ee5c86423 :: prose twin of the run farther from 1
\newcommand{\MlstSculptCoarseMassOverCoarseTen}{0.94 and 0.79}  % experiments/FIGS/data/v2/index.json :: distance at QCD efficiency 0.1: R16_Q1_MASS over R16_Q1, run by run :: sha256 c6c4138ee5c86423
\newcommand{\MlstSculptCoarseMassOverCoarseTenPr}{equal within 30\%}  % experiments/FIGS/data/v2/index.json :: distance at QCD efficiency 0.1: R16_Q1_MASS over R16_Q1, run by run :: sha256 c6c4138ee5c86423 :: prose twin: two significant figures, direction in words
\newcommand{\MlstSculptCoarseMassOverCoarseTenPrMax}{equal within 30\%}  % experiments/FIGS/data/v2/index.json :: distance at QCD efficiency 0.1: R16_Q1_MASS over R16_Q1, run by run :: sha256 c6c4138ee5c86423 :: prose twin of the run farther from 1
\newcommand{\MlstSculptCoarseOverFullTen}{1.04\,\ensuremath{\pm}\,0.05}  % experiments/FIGS/data/v2/index.json :: distance at QCD efficiency 0.1: R16_Q1 over L188, run by run :: sha256 c6c4138ee5c86423
\newcommand{\MlstSculptCoarseOverFullTenPr}{equal within 6\%}  % experiments/FIGS/data/v2/index.json :: distance at QCD efficiency 0.1: R16_Q1 over L188, run by run :: sha256 c6c4138ee5c86423 :: prose twin: two significant figures, direction in words
\newcommand{\MlstSculptResOverFullTen}{1.02\,\ensuremath{\pm}\,0.06}  % experiments/FIGS/data/v2/index.json :: distance at QCD efficiency 0.1: L162 over L188, run by run :: sha256 c6c4138ee5c86423
\newcommand{\MlstSculptResOverFullTenPr}{equal within 6\%}  % experiments/FIGS/data/v2/index.json :: distance at QCD efficiency 0.1: L162 over L188, run by run :: sha256 c6c4138ee5c86423 :: prose twin: two significant figures, direction in words
\newcommand{\MlstSculptMidOverFullTen}{1.069\,\ensuremath{\pm}\,0.020}  % experiments/FIGS/data/v2/index.json :: distance at QCD efficiency 0.1: R42_Q1 over L188, run by run :: sha256 c6c4138ee5c86423
\newcommand{\MlstSculptMidOverFullTenPr}{1.1 times the distance}  % experiments/FIGS/data/v2/index.json :: distance at QCD efficiency 0.1: R42_Q1 over L188, run by run :: sha256 c6c4138ee5c86423 :: prose twin: two significant figures, direction in words
\newcommand{\MlstSculptCoarseOverMidTen}{0.971\,\ensuremath{\pm}\,0.035}  % experiments/FIGS/data/v2/index.json :: distance at QCD efficiency 0.1: R16_Q1 over R42_Q1, run by run :: sha256 c6c4138ee5c86423
\newcommand{\MlstSculptCoarseOverMidTenPr}{equal within 4\%}  % experiments/FIGS/data/v2/index.json :: distance at QCD efficiency 0.1: R16_Q1 over R42_Q1, run by run :: sha256 c6c4138ee5c86423 :: prose twin: two significant figures, direction in words
\newcommand{\MlstSculptNPassOne}{4{,}103}  % experiments/FIGS/data/v2/index.json :: QCD test jets passing the cut at QCD efficiency 0.01, every run :: sha256 c6c4138ee5c86423
\newcommand{\MlstSculptFloorOne}{0.0313}  % experiments/FIGS/data/v2/index.json :: Jensen-Shannon distance expected for a random selection of as many QCD jets as pass at 0.01 :: sha256 c6c4138ee5c86423
\newcommand{\MlstSculptFullOne}{0.0439\,\ensuremath{\pm}\,0.0024}  % experiments/FIGS/data/v2/index.json :: Jensen-Shannon distance of the QCD m_SD distribution at QCD efficiency 0.01, L188 :: sha256 c6c4138ee5c86423
\newcommand{\MlstSculptFullOverFloorOne}{1.40\,\ensuremath{\pm}\,0.08}  % experiments/FIGS/data/v2/index.json :: distance at 0.01 over the random-selection floor, L188 :: sha256 c6c4138ee5c86423
\newcommand{\MlstSculptFullOverFloorOnePr}{1.4 times the distance}  % experiments/FIGS/data/v2/index.json :: distance at 0.01 over the random-selection floor, L188 :: sha256 c6c4138ee5c86423 :: prose twin: two significant figures, direction in words
\newcommand{\MlstSculptResOne}{0.0429\,\ensuremath{\pm}\,0.0029}  % experiments/FIGS/data/v2/index.json :: Jensen-Shannon distance of the QCD m_SD distribution at QCD efficiency 0.01, L162 :: sha256 c6c4138ee5c86423
\newcommand{\MlstSculptResOverFloorOne}{1.37\,\ensuremath{\pm}\,0.09}  % experiments/FIGS/data/v2/index.json :: distance at 0.01 over the random-selection floor, L162 :: sha256 c6c4138ee5c86423
\newcommand{\MlstSculptResOverFloorOnePr}{1.4 times the distance}  % experiments/FIGS/data/v2/index.json :: distance at 0.01 over the random-selection floor, L162 :: sha256 c6c4138ee5c86423 :: prose twin: two significant figures, direction in words
\newcommand{\MlstSculptMidOne}{0.051\,\ensuremath{\pm}\,0.007}  % experiments/FIGS/data/v2/index.json :: Jensen-Shannon distance of the QCD m_SD distribution at QCD efficiency 0.01, R42_Q1 :: sha256 c6c4138ee5c86423
\newcommand{\MlstSculptMidOverFloorOne}{1.62\,\ensuremath{\pm}\,0.23}  % experiments/FIGS/data/v2/index.json :: distance at 0.01 over the random-selection floor, R42_Q1 :: sha256 c6c4138ee5c86423
\newcommand{\MlstSculptMidOverFloorOnePr}{1.6 times the distance}  % experiments/FIGS/data/v2/index.json :: distance at 0.01 over the random-selection floor, R42_Q1 :: sha256 c6c4138ee5c86423 :: prose twin: two significant figures, direction in words
\newcommand{\MlstSculptCoarseOne}{0.0561\,\ensuremath{\pm}\,0.0017}  % experiments/FIGS/data/v2/index.json :: Jensen-Shannon distance of the QCD m_SD distribution at QCD efficiency 0.01, R16_Q1 :: sha256 c6c4138ee5c86423
\newcommand{\MlstSculptCoarseOverFloorOne}{1.79\,\ensuremath{\pm}\,0.06}  % experiments/FIGS/data/v2/index.json :: distance at 0.01 over the random-selection floor, R16_Q1 :: sha256 c6c4138ee5c86423
\newcommand{\MlstSculptCoarseOverFloorOnePr}{1.8 times the distance}  % experiments/FIGS/data/v2/index.json :: distance at 0.01 over the random-selection floor, R16_Q1 :: sha256 c6c4138ee5c86423 :: prose twin: two significant figures, direction in words
\newcommand{\MlstSculptResMassOne}{0.0424 and 0.0442}  % experiments/FIGS/data/v2/index.json :: Jensen-Shannon distance of the QCD m_SD distribution at QCD efficiency 0.01, L162_MASS :: sha256 c6c4138ee5c86423
\newcommand{\MlstSculptResMassOverFloorOne}{1.356 and 1.411}  % experiments/FIGS/data/v2/index.json :: distance at 0.01 over the random-selection floor, L162_MASS :: sha256 c6c4138ee5c86423
\newcommand{\MlstSculptResMassOverFloorOnePr}{at least 1.4 times the distance}  % experiments/FIGS/data/v2/index.json :: distance at 0.01 over the random-selection floor, L162_MASS :: sha256 c6c4138ee5c86423 :: prose twin: two significant figures, direction in words
\newcommand{\MlstSculptResMassOverFloorOnePrMax}{up to 1.4 times the distance}  % experiments/FIGS/data/v2/index.json :: distance at 0.01 over the random-selection floor, L162_MASS :: sha256 c6c4138ee5c86423 :: prose twin of the run farther from 1
\newcommand{\MlstSculptCoarseMassOne}{0.0405 and 0.0452}  % experiments/FIGS/data/v2/index.json :: Jensen-Shannon distance of the QCD m_SD distribution at QCD efficiency 0.01, R16_Q1_MASS :: sha256 c6c4138ee5c86423
\newcommand{\MlstSculptCoarseMassOverFloorOne}{1.29 and 1.44}  % experiments/FIGS/data/v2/index.json :: distance at 0.01 over the random-selection floor, R16_Q1_MASS :: sha256 c6c4138ee5c86423
\newcommand{\MlstSculptCoarseMassOverFloorOnePr}{at least 1.3 times the distance}  % experiments/FIGS/data/v2/index.json :: distance at 0.01 over the random-selection floor, R16_Q1_MASS :: sha256 c6c4138ee5c86423 :: prose twin: two significant figures, direction in words
\newcommand{\MlstSculptCoarseMassOverFloorOnePrMax}{up to 1.4 times the distance}  % experiments/FIGS/data/v2/index.json :: distance at 0.01 over the random-selection floor, R16_Q1_MASS :: sha256 c6c4138ee5c86423 :: prose twin of the run farther from 1
\newcommand{\MlstSculptResMassOverResOne}{0.962 and 0.983}  % experiments/FIGS/data/v2/index.json :: distance at QCD efficiency 0.01: L162_MASS over L162, run by run :: sha256 c6c4138ee5c86423
\newcommand{\MlstSculptResMassOverResOnePr}{equal within 4\%}  % experiments/FIGS/data/v2/index.json :: distance at QCD efficiency 0.01: L162_MASS over L162, run by run :: sha256 c6c4138ee5c86423 :: prose twin: two significant figures, direction in words
\newcommand{\MlstSculptResMassOverResOnePrMax}{equal within 4\%}  % experiments/FIGS/data/v2/index.json :: distance at QCD efficiency 0.01: L162_MASS over L162, run by run :: sha256 c6c4138ee5c86423 :: prose twin of the run farther from 1
\newcommand{\MlstSculptCoarseMassOverCoarseOne}{0.70 and 0.82}  % experiments/FIGS/data/v2/index.json :: distance at QCD efficiency 0.01: R16_Q1_MASS over R16_Q1, run by run :: sha256 c6c4138ee5c86423
\newcommand{\MlstSculptCoarseMassOverCoarseOnePr}{at most 0.82 times the distance}  % experiments/FIGS/data/v2/index.json :: distance at QCD efficiency 0.01: R16_Q1_MASS over R16_Q1, run by run :: sha256 c6c4138ee5c86423 :: prose twin: two significant figures, direction in words
\newcommand{\MlstSculptCoarseMassOverCoarseOnePrMax}{down to 0.70 times the distance}  % experiments/FIGS/data/v2/index.json :: distance at QCD efficiency 0.01: R16_Q1_MASS over R16_Q1, run by run :: sha256 c6c4138ee5c86423 :: prose twin of the run farther from 1
\newcommand{\MlstSculptCoarseOverFullOne}{1.28\,\ensuremath{\pm}\,0.04}  % experiments/FIGS/data/v2/index.json :: distance at QCD efficiency 0.01: R16_Q1 over L188, run by run :: sha256 c6c4138ee5c86423
\newcommand{\MlstSculptCoarseOverFullOnePr}{1.3 times the distance}  % experiments/FIGS/data/v2/index.json :: distance at QCD efficiency 0.01: R16_Q1 over L188, run by run :: sha256 c6c4138ee5c86423 :: prose twin: two significant figures, direction in words
\newcommand{\MlstSculptResOverFullOne}{0.98\,\ensuremath{\pm}\,0.05}  % experiments/FIGS/data/v2/index.json :: distance at QCD efficiency 0.01: L162 over L188, run by run :: sha256 c6c4138ee5c86423
\newcommand{\MlstSculptResOverFullOnePr}{equal within 5\%}  % experiments/FIGS/data/v2/index.json :: distance at QCD efficiency 0.01: L162 over L188, run by run :: sha256 c6c4138ee5c86423 :: prose twin: two significant figures, direction in words
\newcommand{\MlstSculptMidOverFullOne}{1.15\,\ensuremath{\pm}\,0.12}  % experiments/FIGS/data/v2/index.json :: distance at QCD efficiency 0.01: R42_Q1 over L188, run by run :: sha256 c6c4138ee5c86423
\newcommand{\MlstSculptMidOverFullOnePr}{equal within 20\%}  % experiments/FIGS/data/v2/index.json :: distance at QCD efficiency 0.01: R42_Q1 over L188, run by run :: sha256 c6c4138ee5c86423 :: prose twin: two significant figures, direction in words
\newcommand{\MlstSculptCoarseOverMidOne}{1.13\,\ensuremath{\pm}\,0.16}  % experiments/FIGS/data/v2/index.json :: distance at QCD efficiency 0.01: R16_Q1 over R42_Q1, run by run :: sha256 c6c4138ee5c86423
\newcommand{\MlstSculptCoarseOverMidOnePr}{equal within 20\%}  % experiments/FIGS/data/v2/index.json :: distance at QCD efficiency 0.01: R16_Q1 over R42_Q1, run by run :: sha256 c6c4138ee5c86423 :: prose twin: two significant figures, direction in words
\newcommand{\MlstSculptDiffEffLo}{0.52\%}  % experiments/FIGS/data/v2/index.json :: lowest stored QCD efficiency of the unbroken stretch around the tighter working point at which the 17/188 distance ratio exceeds 1 by twice its SD :: sha256 c6c4138ee5c86423 :: passing grid points: 0.001, 0.005245, 0.007937, 0.012011
\newcommand{\MlstSculptDiffEffHi}{1.2\%}  % experiments/FIGS/data/v2/index.json :: highest efficiency of that stretch :: sha256 c6c4138ee5c86423
\newcommand{\MlstSculptLowestEff}{0.10\%}  % experiments/FIGS/data/v2/index.json :: lowest stored QCD efficiency of the sculpting curves :: sha256 c6c4138ee5c86423
\newcommand{\MlstSculptLowestNPass}{410}  % experiments/FIGS/data/v2/index.json :: QCD test jets passing at the lowest stored efficiency :: sha256 c6c4138ee5c86423
\newcommand{\MlstSculptLowestMaxFloorSd}{1.7}  % experiments/FIGS/data/v2/index.json :: largest distance of a configuration's mean from the random-selection floor at the lowest stored efficiency, in SDs of the floor over random selections :: sha256 c6c4138ee5c86423
\newcommand{\MlstSculptNConfigs}{6}  % experiments/FIGS/data/v2/index.json :: configurations with a sculpting curve (count) :: sha256 c6c4138ee5c86423
\newcommand{\LpJciiTestQcdPercent}{20.5\%}  % experiments/FIGS/data/v2/jc2_test/counts.json :: QCD share of the jets, per cent :: sha256 3574ba1f6f2a59f0
\newcommand{\LpJciiTestTwoProngPercent}{38.7\%}  % experiments/FIGS/data/v2/jc2_test/counts.json :: two-prong share of the jets, per cent :: sha256 3574ba1f6f2a59f0
\newcommand{\LpJciiTrainQcdPercent}{13.5\%}  % experiments/FT/data/jc2_v2_manifest.json :: QCD share of the jets, per cent :: sha256 4bd68655919b14e9
\newcommand{\LpJciiTrainTwoProngPercent}{15.4\%}  % experiments/FT/data/jc2_v2_manifest.json :: two-prong share of the jets, per cent :: sha256 4bd68655919b14e9
\newcommand{\LpJciiTestMinAucJets}{10}  % experiments/FIGS/data/v2/jc2_test/counts.json :: fewest jets of a class among the macro-AUC rows :: sha256 3574ba1f6f2a59f0
\newcommand{\LpJciiTrainMinJets}{11}  % experiments/FT/data/jc2_v2_manifest.json :: fewest jets of a class in the 10^6 training subset :: sha256 4bd68655919b14e9
\newcommand{\LpNTestJci}{2{,}000{,}000}  % experiments/FIGS/data/v2/index.json :: jc1 test jets the probes score :: sha256 c6c4138ee5c86423
\newcommand{\LpNTrainMaxJci}{1{,}000{,}000}  % experiments/FIGS/data/v2/index.json :: jc1 largest training subset :: sha256 c6c4138ee5c86423
\newcommand{\LpNTestTop}{404{,}000}  % experiments/FIGS/data/v2/index.json :: top test jets the probes score :: sha256 c6c4138ee5c86423
\newcommand{\LpNTrainMaxTop}{1{,}200{,}000}  % experiments/FIGS/data/v2/index.json :: top largest training subset :: sha256 c6c4138ee5c86423
\newcommand{\LpNTestQg}{200{,}000}  % experiments/FIGS/data/v2/index.json :: qg test jets the probes score :: sha256 c6c4138ee5c86423
\newcommand{\LpNTrainMaxQg}{1{,}600{,}000}  % experiments/FIGS/data/v2/index.json :: qg largest training subset :: sha256 c6c4138ee5c86423
\newcommand{\LpNTestAucJci}{2{,}000{,}000}  % experiments/FIGS/data/v2/index.json :: JetClass test jets of the macro AUC :: sha256 c6c4138ee5c86423
\newcommand{\LpPairNValMin}{1{,}068}  % experiments/FIGS/data/v2/index.json :: validation jets of the pair probe with the fewest :: sha256 c6c4138ee5c86423
\newcommand{\LpPairNValMax}{20{,}195}  % experiments/FIGS/data/v2/index.json :: validation jets of the pair probe with the most :: sha256 c6c4138ee5c86423
\newcommand{\LpNValJci}{200{,}000}  % experiments/FIGS/data/v2/index.json :: JetClass probe validation jets :: sha256 c6c4138ee5c86423
\newcommand{\SculptEffHigh}{10\%}  % experiments/EVAL/mass_sculpting.py :: WORKING_POINTS :: sha256 e85d125c85dac054 :: QCD efficiency of a sculpting cut
\newcommand{\SculptEffLow}{1\%}  % experiments/EVAL/mass_sculpting.py :: WORKING_POINTS :: sha256 e85d125c85dac054 :: QCD efficiency of a sculpting cut
\newcommand{\LpMlpHidden}{256}  % experiments/EVAL/probe.py :: fit_mlp hidden layer :: sha256 485db0966e660588 :: units of the MLP probe's hidden layer
\newcommand{\LpMlpSeedsWord}{three}  % experiments/EVAL/probe.py :: MLP_SEEDS :: sha256 485db0966e660588 :: seeds of the MLP probe whose scores are averaged
\newcommand{\MassMinTrainPerClass}{50}  % experiments/EVAL/mass_resolution.py :: MIN_TRAIN_PER_CLASS :: sha256 26acb0e315094bd1 :: fewest training jets of a class the mass regression keeps
\newcommand{\MlstFtLr}{\ensuremath{10^{-4}}}  % experiments/FIGS/data/v2/ft_mlst/manifests.json :: cells[best70].lr :: sha256 a689f21e894f48ae
\newcommand{\MlstFtHeadMult}{50}  % experiments/FIGS/data/v2/ft_mlst/manifests.json :: cells[best70].head_lr_mult :: sha256 a689f21e894f48ae
\newcommand{\MlstFtWeightDecay}{0.01}  % experiments/FIGS/data/v2/ft_mlst/manifests.json :: cells[best70].weight_decay :: sha256 a689f21e894f48ae
\newcommand{\MlstFtBatch}{512}  % experiments/FIGS/data/v2/ft_mlst/manifests.json :: cells[best70].batch_size :: sha256 a689f21e894f48ae
\newcommand{\MlstFtBenchEpochs}{20}  % experiments/FIGS/data/v2/ft_mlst/manifests.json :: cells[best70, top/qg].epochs :: sha256 a689f21e894f48ae
\newcommand{\MlstFtEpochsEThree}{50}  % experiments/FIGS/data/v2/ft_mlst/manifests.json :: cells[best70, leg 1/2, n_train 1000].epochs :: sha256 a689f21e894f48ae
\newcommand{\MlstFtEpochsEFour}{50}  % experiments/FIGS/data/v2/ft_mlst/manifests.json :: cells[best70, leg 1/2, n_train 10000].epochs :: sha256 a689f21e894f48ae
\newcommand{\MlstFtEpochsEFive}{30}  % experiments/FIGS/data/v2/ft_mlst/manifests.json :: cells[best70, leg 1/2, n_train 100000].epochs :: sha256 a689f21e894f48ae
\newcommand{\MlstFtEpochsESix}{10}  % experiments/FIGS/data/v2/ft_mlst/manifests.json :: cells[best70, leg 1/2, n_train 1000000].epochs :: sha256 a689f21e894f48ae
\newcommand{\MlstFtHeadDecayShare}{30\%}  % experiments/FT/k8s/job-ft-v2-legs-best70-l188-s3-raunav.yaml :: no --lr-scheduler: weaver 0.4.17 flat+decay, the head's last epochs :: sha256 19f59aa54a385c11
\newcommand{\MlstFtHeadDecayTo}{1\%}  % experiments/FT/k8s/job-ft-v2-legs-best70-l188-s3-raunav.yaml :: no --lr-scheduler: weaver 0.4.17 flat+decay, final factor :: sha256 19f59aa54a385c11
\newcommand{\MlstFtHeadDropout}{0.1}  % experiments/FT/k8s/job-ft-v2-legs-best70-l188-s3-raunav.yaml :: fc_params dropout of the fine-tuning head :: sha256 19f59aa54a385c11
\newcommand{\MlstFtMinEpochJets}{10{,}000}  % experiments/FIGS/data/v2/ft_mlst/manifests.json :: cells[best70].samples_per_epoch (smallest) :: sha256 a689f21e894f48ae
\newcommand{\MlstFtValJcii}{20{,}480}  % experiments/FIGS/data/v2/ft_mlst/manifests.json :: cells[best70, leg 1].val_jets :: sha256 a689f21e894f48ae
\newcommand{\MlstFtValSmall}{20{,}000}  % experiments/FT/k8s/job-ft-v2-legs-best70-l188-s3-raunav.yaml :: JetClass --samples-per-epoch-val; top and q/g at N <= 1e4 :: sha256 19f59aa54a385c11
\newcommand{\MlstFtValLarge}{200{,}000}  % experiments/FT/k8s/job-ft-v2-bench-best70-l188-s1-top-raunav.yaml :: top and q/g --samples-per-epoch-val at N > 1e4 :: sha256 1e38d4f542119ea3
\newcommand{\MlstFtGpuTypes}{V100}  % experiments/FIGS/data/v2/ft_mlst/manifests.json :: cells[best70].gpu_device_name :: sha256 a689f21e894f48ae
\newcommand{\MlstFtCellsDone}{445}  % experiments/FIGS/data/v2/ft_mlst/manifests.json :: DONE paper-run fine-tuning cells (count) :: sha256 a689f21e894f48ae
\newcommand{\MlstFtRefLr}{\ensuremath{10^{-3}}}  % experiments/FIGS/data/v2/ft_mlst/manifests.json :: cells[ref_v100].lr :: sha256 a689f21e894f48ae
\newcommand{\MlstFtRefSeedsWord}{three}  % experiments/FIGS/data/v2/ft_mlst/manifests.json :: cells[ref_v100].ft_seed (count) :: sha256 a689f21e894f48ae
\newcommand{\MlstFtRefGpuTypes}{V100}  % experiments/FIGS/data/v2/ft_mlst/manifests.json :: cells[ref_v100].gpu_device_name :: sha256 a689f21e894f48ae
\newcommand{\MlstFtRefWeightDecayBench}{0.01}  % experiments/FIGS/data/v2/ft_mlst/manifests.json :: cells[ref_v100, top/qg].weight_decay :: sha256 a689f21e894f48ae
\newcommand{\MlstBootResamples}{200}  % experiments/FIGS/data/v2/index.json :: bootstrap resamples of the test jets :: sha256 c6c4138ee5c86423
\newcommand{\MlstBootTestMaxProbe}{0.21}  % experiments/FIGS/data/v2/index.json :: largest test-sample SD of the ln ratio of 1 - AUC on JetClass-II, linear probe, label sets over the 188-class model, per cent :: sha256 c6c4138ee5c86423
\newcommand{\MlstBootTestMaxFt}{0.15}  % experiments/FIGS/data/v2/index.json :: largest test-sample SD of the ln ratio of 1 - AUC on JetClass-II, fine-tuning, label sets over the 188-class model, per cent :: sha256 c6c4138ee5c86423
\newcommand{\MlstBootTestMaxRthirtyProbe}{3.6}  % experiments/FIGS/data/v2/index.json :: largest test-sample SD of the ln R30 ratio on top and quark/gluon, linear probe (counting), label sets over the 188-class model, per cent :: sha256 c6c4138ee5c86423
\newcommand{\MlstBootTestMaxRthirtyFt}{13}  % experiments/FIGS/data/v2/index.json :: largest test-sample SD of the ln R30 ratio on top and quark/gluon, fine-tuning (counting), label sets over the 188-class model, per cent :: sha256 c6c4138ee5c86423
\newcommand{\MlstFtCmpJcTwoResOverFullEFour}{0.989\,\ensuremath{\pm}\,0.011}  % experiments/FIGS/data/v2/index.json :: fine-tuning jc2 N=10000: (1 - macro_auc_ovr) of L162 over L188, run by run :: sha256 c6c4138ee5c86423
\newcommand{\MlstFtCmpJcTwoResOverFullEFourPr}{equal within 2\%}  % experiments/FIGS/data/v2/index.json :: fine-tuning jc2 N=10000: (1 - macro_auc_ovr) of L162 over L188, run by run :: sha256 c6c4138ee5c86423 :: prose twin: two significant figures, direction in words
\newcommand{\MlstFtCmpJcTwoMidOverFullEFour}{1.191\,\ensuremath{\pm}\,0.030}  % experiments/FIGS/data/v2/index.json :: fine-tuning jc2 N=10000: (1 - macro_auc_ovr) of R42_Q1 over L188, run by run :: sha256 c6c4138ee5c86423
\newcommand{\MlstFtCmpJcTwoMidOverFullEFourPr}{19\% more errors}  % experiments/FIGS/data/v2/index.json :: fine-tuning jc2 N=10000: (1 - macro_auc_ovr) of R42_Q1 over L188, run by run :: sha256 c6c4138ee5c86423 :: prose twin: two significant figures, direction in words
\newcommand{\MlstFtCmpJcTwoCoarseOverFullEFour}{1.353\,\ensuremath{\pm}\,0.003}  % experiments/FIGS/data/v2/index.json :: fine-tuning jc2 N=10000: (1 - macro_auc_ovr) of R16_Q1 over L188, run by run :: sha256 c6c4138ee5c86423
\newcommand{\MlstFtCmpJcTwoCoarseOverFullEFourPr}{35\% more errors}  % experiments/FIGS/data/v2/index.json :: fine-tuning jc2 N=10000: (1 - macro_auc_ovr) of R16_Q1 over L188, run by run :: sha256 c6c4138ee5c86423 :: prose twin: two significant figures, direction in words
\newcommand{\MlstFtCmpJcTwoResMassOverResEFour}{1.016 and 1.005}  % experiments/FIGS/data/v2/index.json :: fine-tuning jc2 N=10000: (1 - macro_auc_ovr) of L162_MASS over L162, run by run :: sha256 c6c4138ee5c86423
\newcommand{\MlstFtCmpJcTwoResMassOverResEFourPr}{equal within 2\%}  % experiments/FIGS/data/v2/index.json :: fine-tuning jc2 N=10000: (1 - macro_auc_ovr) of L162_MASS over L162, run by run :: sha256 c6c4138ee5c86423 :: prose twin: two significant figures, direction in words
\newcommand{\MlstFtCmpJcTwoResMassOverResEFourPrMax}{equal within 2\%}  % experiments/FIGS/data/v2/index.json :: fine-tuning jc2 N=10000: (1 - macro_auc_ovr) of L162_MASS over L162, run by run :: sha256 c6c4138ee5c86423 :: prose twin of the run farther from 1
\newcommand{\MlstFtCmpJcTwoCoarseMassOverCoarseEFour}{0.989 and 1.000}  % experiments/FIGS/data/v2/index.json :: fine-tuning jc2 N=10000: (1 - macro_auc_ovr) of R16_Q1_MASS over R16_Q1, run by run :: sha256 c6c4138ee5c86423
\newcommand{\MlstFtCmpJcTwoCoarseMassOverCoarseEFourPr}{equal within 2\%}  % experiments/FIGS/data/v2/index.json :: fine-tuning jc2 N=10000: (1 - macro_auc_ovr) of R16_Q1_MASS over R16_Q1, run by run :: sha256 c6c4138ee5c86423 :: prose twin: two significant figures, direction in words
\newcommand{\MlstFtCmpJcTwoCoarseMassOverCoarseEFourPrMax}{equal within 2\%}  % experiments/FIGS/data/v2/index.json :: fine-tuning jc2 N=10000: (1 - macro_auc_ovr) of R16_Q1_MASS over R16_Q1, run by run :: sha256 c6c4138ee5c86423 :: prose twin of the run farther from 1
\newcommand{\MlstFtCmpJcTwoResOverFullEFive}{1.011\,\ensuremath{\pm}\,0.005}  % experiments/FIGS/data/v2/index.json :: fine-tuning jc2 N=100000: (1 - macro_auc_ovr) of L162 over L188, run by run :: sha256 c6c4138ee5c86423
\newcommand{\MlstFtCmpJcTwoResOverFullEFivePr}{1.1\% more errors}  % experiments/FIGS/data/v2/index.json :: fine-tuning jc2 N=100000: (1 - macro_auc_ovr) of L162 over L188, run by run :: sha256 c6c4138ee5c86423 :: prose twin: two significant figures, direction in words
\newcommand{\MlstFtCmpJcTwoMidOverFullEFive}{1.064\,\ensuremath{\pm}\,0.009}  % experiments/FIGS/data/v2/index.json :: fine-tuning jc2 N=100000: (1 - macro_auc_ovr) of R42_Q1 over L188, run by run :: sha256 c6c4138ee5c86423
\newcommand{\MlstFtCmpJcTwoMidOverFullEFivePr}{6.4\% more errors}  % experiments/FIGS/data/v2/index.json :: fine-tuning jc2 N=100000: (1 - macro_auc_ovr) of R42_Q1 over L188, run by run :: sha256 c6c4138ee5c86423 :: prose twin: two significant figures, direction in words
\newcommand{\MlstFtCmpJcTwoCoarseOverFullEFive}{1.157\,\ensuremath{\pm}\,0.001}  % experiments/FIGS/data/v2/index.json :: fine-tuning jc2 N=100000: (1 - macro_auc_ovr) of R16_Q1 over L188, run by run :: sha256 c6c4138ee5c86423
\newcommand{\MlstFtCmpJcTwoCoarseOverFullEFivePr}{16\% more errors}  % experiments/FIGS/data/v2/index.json :: fine-tuning jc2 N=100000: (1 - macro_auc_ovr) of R16_Q1 over L188, run by run :: sha256 c6c4138ee5c86423 :: prose twin: two significant figures, direction in words
\newcommand{\MlstFtCmpJcTwoResMassOverResEFive}{1.027 and 1.021}  % experiments/FIGS/data/v2/index.json :: fine-tuning jc2 N=100000: (1 - macro_auc_ovr) of L162_MASS over L162, run by run :: sha256 c6c4138ee5c86423
\newcommand{\MlstFtCmpJcTwoResMassOverResEFivePr}{at least 2.1\% more errors}  % experiments/FIGS/data/v2/index.json :: fine-tuning jc2 N=100000: (1 - macro_auc_ovr) of L162_MASS over L162, run by run :: sha256 c6c4138ee5c86423 :: prose twin: two significant figures, direction in words
\newcommand{\MlstFtCmpJcTwoResMassOverResEFivePrMax}{up to 2.7\% more errors}  % experiments/FIGS/data/v2/index.json :: fine-tuning jc2 N=100000: (1 - macro_auc_ovr) of L162_MASS over L162, run by run :: sha256 c6c4138ee5c86423 :: prose twin of the run farther from 1
\newcommand{\MlstFtCmpJcTwoCoarseMassOverCoarseEFive}{1.029 and 1.022}  % experiments/FIGS/data/v2/index.json :: fine-tuning jc2 N=100000: (1 - macro_auc_ovr) of R16_Q1_MASS over R16_Q1, run by run :: sha256 c6c4138ee5c86423
\newcommand{\MlstFtCmpJcTwoCoarseMassOverCoarseEFivePr}{at least 2.2\% more errors}  % experiments/FIGS/data/v2/index.json :: fine-tuning jc2 N=100000: (1 - macro_auc_ovr) of R16_Q1_MASS over R16_Q1, run by run :: sha256 c6c4138ee5c86423 :: prose twin: two significant figures, direction in words
\newcommand{\MlstFtCmpJcTwoCoarseMassOverCoarseEFivePrMax}{up to 2.9\% more errors}  % experiments/FIGS/data/v2/index.json :: fine-tuning jc2 N=100000: (1 - macro_auc_ovr) of R16_Q1_MASS over R16_Q1, run by run :: sha256 c6c4138ee5c86423 :: prose twin of the run farther from 1
\newcommand{\MlstFtCmpJcTwoResOverFullESix}{1.014\,\ensuremath{\pm}\,0.011}  % experiments/FIGS/data/v2/index.json :: fine-tuning jc2 N=1000000: (1 - macro_auc_ovr) of L162 over L188, run by run :: sha256 c6c4138ee5c86423
\newcommand{\MlstFtCmpJcTwoResOverFullESixPr}{equal within 2\%}  % experiments/FIGS/data/v2/index.json :: fine-tuning jc2 N=1000000: (1 - macro_auc_ovr) of L162 over L188, run by run :: sha256 c6c4138ee5c86423 :: prose twin: two significant figures, direction in words
\newcommand{\MlstFtCmpJcTwoMidOverFullESix}{1.043\,\ensuremath{\pm}\,0.002}  % experiments/FIGS/data/v2/index.json :: fine-tuning jc2 N=1000000: (1 - macro_auc_ovr) of R42_Q1 over L188, run by run :: sha256 c6c4138ee5c86423
\newcommand{\MlstFtCmpJcTwoMidOverFullESixPr}{4.3\% more errors}  % experiments/FIGS/data/v2/index.json :: fine-tuning jc2 N=1000000: (1 - macro_auc_ovr) of R42_Q1 over L188, run by run :: sha256 c6c4138ee5c86423 :: prose twin: two significant figures, direction in words
\newcommand{\MlstFtCmpJcTwoCoarseOverFullESix}{1.123\,\ensuremath{\pm}\,0.003}  % experiments/FIGS/data/v2/index.json :: fine-tuning jc2 N=1000000: (1 - macro_auc_ovr) of R16_Q1 over L188, run by run :: sha256 c6c4138ee5c86423
\newcommand{\MlstFtCmpJcTwoCoarseOverFullESixPr}{12\% more errors}  % experiments/FIGS/data/v2/index.json :: fine-tuning jc2 N=1000000: (1 - macro_auc_ovr) of R16_Q1 over L188, run by run :: sha256 c6c4138ee5c86423 :: prose twin: two significant figures, direction in words
\newcommand{\MlstFtCmpJcTwoResMassOverResESix}{1.009 and 1.016}  % experiments/FIGS/data/v2/index.json :: fine-tuning jc2 N=1000000: (1 - macro_auc_ovr) of L162_MASS over L162, run by run :: sha256 c6c4138ee5c86423
\newcommand{\MlstFtCmpJcTwoResMassOverResESixPr}{equal within 2\%}  % experiments/FIGS/data/v2/index.json :: fine-tuning jc2 N=1000000: (1 - macro_auc_ovr) of L162_MASS over L162, run by run :: sha256 c6c4138ee5c86423 :: prose twin: two significant figures, direction in words
\newcommand{\MlstFtCmpJcTwoResMassOverResESixPrMax}{equal within 2\%}  % experiments/FIGS/data/v2/index.json :: fine-tuning jc2 N=1000000: (1 - macro_auc_ovr) of L162_MASS over L162, run by run :: sha256 c6c4138ee5c86423 :: prose twin of the run farther from 1
\newcommand{\MlstFtCmpJcTwoCoarseMassOverCoarseESix}{1.023 and 1.017}  % experiments/FIGS/data/v2/index.json :: fine-tuning jc2 N=1000000: (1 - macro_auc_ovr) of R16_Q1_MASS over R16_Q1, run by run :: sha256 c6c4138ee5c86423
\newcommand{\MlstFtCmpJcTwoCoarseMassOverCoarseESixPr}{at least 1.7\% more errors}  % experiments/FIGS/data/v2/index.json :: fine-tuning jc2 N=1000000: (1 - macro_auc_ovr) of R16_Q1_MASS over R16_Q1, run by run :: sha256 c6c4138ee5c86423 :: prose twin: two significant figures, direction in words
\newcommand{\MlstFtCmpJcTwoCoarseMassOverCoarseESixPrMax}{up to 2.3\% more errors}  % experiments/FIGS/data/v2/index.json :: fine-tuning jc2 N=1000000: (1 - macro_auc_ovr) of R16_Q1_MASS over R16_Q1, run by run :: sha256 c6c4138ee5c86423 :: prose twin of the run farther from 1
\newcommand{\MlstFtCmpJcTwoQcdResOverFullEFour}{1.038\,\ensuremath{\pm}\,0.005}  % experiments/FIGS/data/v2/index.json :: fine-tuning jc2_qcd N=10000: (1 - macro_auc_qcd) of L162 over L188, run by run :: sha256 c6c4138ee5c86423
\newcommand{\MlstFtCmpJcTwoQcdResOverFullEFourPr}{3.8\% more errors}  % experiments/FIGS/data/v2/index.json :: fine-tuning jc2_qcd N=10000: (1 - macro_auc_qcd) of L162 over L188, run by run :: sha256 c6c4138ee5c86423 :: prose twin: two significant figures, direction in words
\newcommand{\MlstFtCmpJcTwoQcdMidOverFullEFour}{1.164\,\ensuremath{\pm}\,0.018}  % experiments/FIGS/data/v2/index.json :: fine-tuning jc2_qcd N=10000: (1 - macro_auc_qcd) of R42_Q1 over L188, run by run :: sha256 c6c4138ee5c86423
\newcommand{\MlstFtCmpJcTwoQcdMidOverFullEFourPr}{16\% more errors}  % experiments/FIGS/data/v2/index.json :: fine-tuning jc2_qcd N=10000: (1 - macro_auc_qcd) of R42_Q1 over L188, run by run :: sha256 c6c4138ee5c86423 :: prose twin: two significant figures, direction in words
\newcommand{\MlstFtCmpJcTwoQcdCoarseOverFullEFour}{1.258\,\ensuremath{\pm}\,0.025}  % experiments/FIGS/data/v2/index.json :: fine-tuning jc2_qcd N=10000: (1 - macro_auc_qcd) of R16_Q1 over L188, run by run :: sha256 c6c4138ee5c86423
\newcommand{\MlstFtCmpJcTwoQcdCoarseOverFullEFourPr}{26\% more errors}  % experiments/FIGS/data/v2/index.json :: fine-tuning jc2_qcd N=10000: (1 - macro_auc_qcd) of R16_Q1 over L188, run by run :: sha256 c6c4138ee5c86423 :: prose twin: two significant figures, direction in words
\newcommand{\MlstFtCmpJcTwoQcdResMassOverResEFour}{1.027 and 0.997}  % experiments/FIGS/data/v2/index.json :: fine-tuning jc2_qcd N=10000: (1 - macro_auc_qcd) of L162_MASS over L162, run by run :: sha256 c6c4138ee5c86423
\newcommand{\MlstFtCmpJcTwoQcdResMassOverResEFourPr}{equal within 3\%}  % experiments/FIGS/data/v2/index.json :: fine-tuning jc2_qcd N=10000: (1 - macro_auc_qcd) of L162_MASS over L162, run by run :: sha256 c6c4138ee5c86423 :: prose twin: two significant figures, direction in words
\newcommand{\MlstFtCmpJcTwoQcdResMassOverResEFourPrMax}{equal within 3\%}  % experiments/FIGS/data/v2/index.json :: fine-tuning jc2_qcd N=10000: (1 - macro_auc_qcd) of L162_MASS over L162, run by run :: sha256 c6c4138ee5c86423 :: prose twin of the run farther from 1
\newcommand{\MlstFtCmpJcTwoQcdCoarseMassOverCoarseEFour}{1.028 and 1.048}  % experiments/FIGS/data/v2/index.json :: fine-tuning jc2_qcd N=10000: (1 - macro_auc_qcd) of R16_Q1_MASS over R16_Q1, run by run :: sha256 c6c4138ee5c86423
\newcommand{\MlstFtCmpJcTwoQcdCoarseMassOverCoarseEFourPr}{equal within 5\%}  % experiments/FIGS/data/v2/index.json :: fine-tuning jc2_qcd N=10000: (1 - macro_auc_qcd) of R16_Q1_MASS over R16_Q1, run by run :: sha256 c6c4138ee5c86423 :: prose twin: two significant figures, direction in words
\newcommand{\MlstFtCmpJcTwoQcdCoarseMassOverCoarseEFourPrMax}{equal within 5\%}  % experiments/FIGS/data/v2/index.json :: fine-tuning jc2_qcd N=10000: (1 - macro_auc_qcd) of R16_Q1_MASS over R16_Q1, run by run :: sha256 c6c4138ee5c86423 :: prose twin of the run farther from 1
\newcommand{\MlstFtCmpJcTwoQcdResOverFullEFive}{1.060\,\ensuremath{\pm}\,0.002}  % experiments/FIGS/data/v2/index.json :: fine-tuning jc2_qcd N=100000: (1 - macro_auc_qcd) of L162 over L188, run by run :: sha256 c6c4138ee5c86423
\newcommand{\MlstFtCmpJcTwoQcdResOverFullEFivePr}{6.0\% more errors}  % experiments/FIGS/data/v2/index.json :: fine-tuning jc2_qcd N=100000: (1 - macro_auc_qcd) of L162 over L188, run by run :: sha256 c6c4138ee5c86423 :: prose twin: two significant figures, direction in words
\newcommand{\MlstFtCmpJcTwoQcdMidOverFullEFive}{1.062\,\ensuremath{\pm}\,0.010}  % experiments/FIGS/data/v2/index.json :: fine-tuning jc2_qcd N=100000: (1 - macro_auc_qcd) of R42_Q1 over L188, run by run :: sha256 c6c4138ee5c86423
\newcommand{\MlstFtCmpJcTwoQcdMidOverFullEFivePr}{6.2\% more errors}  % experiments/FIGS/data/v2/index.json :: fine-tuning jc2_qcd N=100000: (1 - macro_auc_qcd) of R42_Q1 over L188, run by run :: sha256 c6c4138ee5c86423 :: prose twin: two significant figures, direction in words
\newcommand{\MlstFtCmpJcTwoQcdCoarseOverFullEFive}{1.137\,\ensuremath{\pm}\,0.001}  % experiments/FIGS/data/v2/index.json :: fine-tuning jc2_qcd N=100000: (1 - macro_auc_qcd) of R16_Q1 over L188, run by run :: sha256 c6c4138ee5c86423
\newcommand{\MlstFtCmpJcTwoQcdCoarseOverFullEFivePr}{14\% more errors}  % experiments/FIGS/data/v2/index.json :: fine-tuning jc2_qcd N=100000: (1 - macro_auc_qcd) of R16_Q1 over L188, run by run :: sha256 c6c4138ee5c86423 :: prose twin: two significant figures, direction in words
\newcommand{\MlstFtCmpJcTwoQcdResMassOverResEFive}{1.022 and 1.009}  % experiments/FIGS/data/v2/index.json :: fine-tuning jc2_qcd N=100000: (1 - macro_auc_qcd) of L162_MASS over L162, run by run :: sha256 c6c4138ee5c86423
\newcommand{\MlstFtCmpJcTwoQcdResMassOverResEFivePr}{equal within 3\%}  % experiments/FIGS/data/v2/index.json :: fine-tuning jc2_qcd N=100000: (1 - macro_auc_qcd) of L162_MASS over L162, run by run :: sha256 c6c4138ee5c86423 :: prose twin: two significant figures, direction in words
\newcommand{\MlstFtCmpJcTwoQcdResMassOverResEFivePrMax}{equal within 3\%}  % experiments/FIGS/data/v2/index.json :: fine-tuning jc2_qcd N=100000: (1 - macro_auc_qcd) of L162_MASS over L162, run by run :: sha256 c6c4138ee5c86423 :: prose twin of the run farther from 1
\newcommand{\MlstFtCmpJcTwoQcdCoarseMassOverCoarseEFive}{1.028 and 1.027}  % experiments/FIGS/data/v2/index.json :: fine-tuning jc2_qcd N=100000: (1 - macro_auc_qcd) of R16_Q1_MASS over R16_Q1, run by run :: sha256 c6c4138ee5c86423
\newcommand{\MlstFtCmpJcTwoQcdCoarseMassOverCoarseEFivePr}{at least 2.7\% more errors}  % experiments/FIGS/data/v2/index.json :: fine-tuning jc2_qcd N=100000: (1 - macro_auc_qcd) of R16_Q1_MASS over R16_Q1, run by run :: sha256 c6c4138ee5c86423 :: prose twin: two significant figures, direction in words
\newcommand{\MlstFtCmpJcTwoQcdCoarseMassOverCoarseEFivePrMax}{up to 2.8\% more errors}  % experiments/FIGS/data/v2/index.json :: fine-tuning jc2_qcd N=100000: (1 - macro_auc_qcd) of R16_Q1_MASS over R16_Q1, run by run :: sha256 c6c4138ee5c86423 :: prose twin of the run farther from 1
\newcommand{\MlstFtCmpJcTwoQcdResOverFullESix}{1.035\,\ensuremath{\pm}\,0.010}  % experiments/FIGS/data/v2/index.json :: fine-tuning jc2_qcd N=1000000: (1 - macro_auc_qcd) of L162 over L188, run by run :: sha256 c6c4138ee5c86423
\newcommand{\MlstFtCmpJcTwoQcdResOverFullESixPr}{3.5\% more errors}  % experiments/FIGS/data/v2/index.json :: fine-tuning jc2_qcd N=1000000: (1 - macro_auc_qcd) of L162 over L188, run by run :: sha256 c6c4138ee5c86423 :: prose twin: two significant figures, direction in words
\newcommand{\MlstFtCmpJcTwoQcdMidOverFullESix}{1.050\,\ensuremath{\pm}\,0.004}  % experiments/FIGS/data/v2/index.json :: fine-tuning jc2_qcd N=1000000: (1 - macro_auc_qcd) of R42_Q1 over L188, run by run :: sha256 c6c4138ee5c86423
\newcommand{\MlstFtCmpJcTwoQcdMidOverFullESixPr}{5.0\% more errors}  % experiments/FIGS/data/v2/index.json :: fine-tuning jc2_qcd N=1000000: (1 - macro_auc_qcd) of R42_Q1 over L188, run by run :: sha256 c6c4138ee5c86423 :: prose twin: two significant figures, direction in words
\newcommand{\MlstFtCmpJcTwoQcdCoarseOverFullESix}{1.122\,\ensuremath{\pm}\,0.002}  % experiments/FIGS/data/v2/index.json :: fine-tuning jc2_qcd N=1000000: (1 - macro_auc_qcd) of R16_Q1 over L188, run by run :: sha256 c6c4138ee5c86423
\newcommand{\MlstFtCmpJcTwoQcdCoarseOverFullESixPr}{12\% more errors}  % experiments/FIGS/data/v2/index.json :: fine-tuning jc2_qcd N=1000000: (1 - macro_auc_qcd) of R16_Q1 over L188, run by run :: sha256 c6c4138ee5c86423 :: prose twin: two significant figures, direction in words
\newcommand{\MlstFtCmpJcTwoQcdResMassOverResESix}{1.011 and 1.025}  % experiments/FIGS/data/v2/index.json :: fine-tuning jc2_qcd N=1000000: (1 - macro_auc_qcd) of L162_MASS over L162, run by run :: sha256 c6c4138ee5c86423
\newcommand{\MlstFtCmpJcTwoQcdResMassOverResESixPr}{equal within 3\%}  % experiments/FIGS/data/v2/index.json :: fine-tuning jc2_qcd N=1000000: (1 - macro_auc_qcd) of L162_MASS over L162, run by run :: sha256 c6c4138ee5c86423 :: prose twin: two significant figures, direction in words
\newcommand{\MlstFtCmpJcTwoQcdResMassOverResESixPrMax}{equal within 3\%}  % experiments/FIGS/data/v2/index.json :: fine-tuning jc2_qcd N=1000000: (1 - macro_auc_qcd) of L162_MASS over L162, run by run :: sha256 c6c4138ee5c86423 :: prose twin of the run farther from 1
\newcommand{\MlstFtCmpJcTwoQcdCoarseMassOverCoarseESix}{1.015 and 1.015}  % experiments/FIGS/data/v2/index.json :: fine-tuning jc2_qcd N=1000000: (1 - macro_auc_qcd) of R16_Q1_MASS over R16_Q1, run by run :: sha256 c6c4138ee5c86423
\newcommand{\MlstFtCmpJcTwoQcdCoarseMassOverCoarseESixPr}{at least 1.5\% more errors}  % experiments/FIGS/data/v2/index.json :: fine-tuning jc2_qcd N=1000000: (1 - macro_auc_qcd) of R16_Q1_MASS over R16_Q1, run by run :: sha256 c6c4138ee5c86423 :: prose twin: two significant figures, direction in words
\newcommand{\MlstFtCmpJcTwoQcdCoarseMassOverCoarseESixPrMax}{up to 1.5\% more errors}  % experiments/FIGS/data/v2/index.json :: fine-tuning jc2_qcd N=1000000: (1 - macro_auc_qcd) of R16_Q1_MASS over R16_Q1, run by run :: sha256 c6c4138ee5c86423 :: prose twin of the run farther from 1
\newcommand{\MlstFtJcTwoAccFullEFour}{0.3914\,\ensuremath{\pm}\,0.0011}  % experiments/FIGS/data/v2/index.json :: fine-tuning jc2 N=10000: accuracy, L188 :: sha256 c6c4138ee5c86423
\newcommand{\MlstFtJcTwoAccResEFour}{0.3934\,\ensuremath{\pm}\,0.0018}  % experiments/FIGS/data/v2/index.json :: fine-tuning jc2 N=10000: accuracy, L162 :: sha256 c6c4138ee5c86423
\newcommand{\MlstFtJcTwoAccMidEFour}{0.3563\,\ensuremath{\pm}\,0.0025}  % experiments/FIGS/data/v2/index.json :: fine-tuning jc2 N=10000: accuracy, R42_Q1 :: sha256 c6c4138ee5c86423
\newcommand{\MlstFtJcTwoAccCoarseEFour}{0.3219\,\ensuremath{\pm}\,0.0011}  % experiments/FIGS/data/v2/index.json :: fine-tuning jc2 N=10000: accuracy, R16_Q1 :: sha256 c6c4138ee5c86423
\newcommand{\MlstFtJcTwoAccResMassEFour}{0.381 and 0.392}  % experiments/FIGS/data/v2/index.json :: fine-tuning jc2 N=10000: accuracy, L162_MASS :: sha256 c6c4138ee5c86423
\newcommand{\MlstFtJcTwoAccCoarseMassEFour}{0.3156 and 0.3187}  % experiments/FIGS/data/v2/index.json :: fine-tuning jc2 N=10000: accuracy, R16_Q1_MASS :: sha256 c6c4138ee5c86423
\newcommand{\MlstFtJcTwoAccScratchEFour}{0.1492\,\ensuremath{\pm}\,0.0011}  % experiments/FIGS/data/v2/index.json :: fine-tuning jc2 N=10000: accuracy from scratch :: sha256 c6c4138ee5c86423
\newcommand{\MlstFtJcTwoAccFullMinusScratchEFour}{0.2423\,\ensuremath{\pm}\,0.0011}  % experiments/FIGS/data/v2/index.json :: fine-tuning jc2 N=10000: accuracy of each 188-class run minus the from-scratch mean (absolute) :: sha256 c6c4138ee5c86423
\newcommand{\MlstFtJcTwoAccFullEFive}{0.4164\,\ensuremath{\pm}\,0.0015}  % experiments/FIGS/data/v2/index.json :: fine-tuning jc2 N=100000: accuracy, L188 :: sha256 c6c4138ee5c86423
\newcommand{\MlstFtJcTwoAccResEFive}{0.4161\,\ensuremath{\pm}\,0.0015}  % experiments/FIGS/data/v2/index.json :: fine-tuning jc2 N=100000: accuracy, L162 :: sha256 c6c4138ee5c86423
\newcommand{\MlstFtJcTwoAccMidEFive}{0.4007\,\ensuremath{\pm}\,0.0013}  % experiments/FIGS/data/v2/index.json :: fine-tuning jc2 N=100000: accuracy, R42_Q1 :: sha256 c6c4138ee5c86423
\newcommand{\MlstFtJcTwoAccCoarseEFive}{0.3833\,\ensuremath{\pm}\,0.0020}  % experiments/FIGS/data/v2/index.json :: fine-tuning jc2 N=100000: accuracy, R16_Q1 :: sha256 c6c4138ee5c86423
\newcommand{\MlstFtJcTwoAccResMassEFive}{0.4093 and 0.4115}  % experiments/FIGS/data/v2/index.json :: fine-tuning jc2 N=100000: accuracy, L162_MASS :: sha256 c6c4138ee5c86423
\newcommand{\MlstFtJcTwoAccCoarseMassEFive}{0.3758 and 0.3781}  % experiments/FIGS/data/v2/index.json :: fine-tuning jc2 N=100000: accuracy, R16_Q1_MASS :: sha256 c6c4138ee5c86423
\newcommand{\MlstFtJcTwoAccScratchEFive}{0.292\,\ensuremath{\pm}\,0.004}  % experiments/FIGS/data/v2/index.json :: fine-tuning jc2 N=100000: accuracy from scratch :: sha256 c6c4138ee5c86423
\newcommand{\MlstFtJcTwoAccFullMinusScratchEFive}{0.1243\,\ensuremath{\pm}\,0.0015}  % experiments/FIGS/data/v2/index.json :: fine-tuning jc2 N=100000: accuracy of each 188-class run minus the from-scratch mean (absolute) :: sha256 c6c4138ee5c86423
\newcommand{\MlstFtJcTwoAccFullESix}{0.4263\,\ensuremath{\pm}\,0.0005}  % experiments/FIGS/data/v2/index.json :: fine-tuning jc2 N=1000000: accuracy, L188 :: sha256 c6c4138ee5c86423
\newcommand{\MlstFtJcTwoAccResESix}{0.4249\,\ensuremath{\pm}\,0.0021}  % experiments/FIGS/data/v2/index.json :: fine-tuning jc2 N=1000000: accuracy, L162 :: sha256 c6c4138ee5c86423
\newcommand{\MlstFtJcTwoAccMidESix}{0.4166\,\ensuremath{\pm}\,0.0011}  % experiments/FIGS/data/v2/index.json :: fine-tuning jc2 N=1000000: accuracy, R42_Q1 :: sha256 c6c4138ee5c86423
\newcommand{\MlstFtJcTwoAccCoarseESix}{0.4054\,\ensuremath{\pm}\,0.0020}  % experiments/FIGS/data/v2/index.json :: fine-tuning jc2 N=1000000: accuracy, R16_Q1 :: sha256 c6c4138ee5c86423
\newcommand{\MlstFtJcTwoAccResMassESix}{0.4198 and 0.4213}  % experiments/FIGS/data/v2/index.json :: fine-tuning jc2 N=1000000: accuracy, L162_MASS :: sha256 c6c4138ee5c86423
\newcommand{\MlstFtJcTwoAccCoarseMassESix}{0.4012 and 0.4000}  % experiments/FIGS/data/v2/index.json :: fine-tuning jc2 N=1000000: accuracy, R16_Q1_MASS :: sha256 c6c4138ee5c86423
\newcommand{\MlstFtJcTwoAccScratchESix}{0.3677\,\ensuremath{\pm}\,0.0012}  % experiments/FIGS/data/v2/index.json :: fine-tuning jc2 N=1000000: accuracy from scratch :: sha256 c6c4138ee5c86423
\newcommand{\MlstFtJcTwoAccFullMinusScratchESix}{0.0585\,\ensuremath{\pm}\,0.0005}  % experiments/FIGS/data/v2/index.json :: fine-tuning jc2 N=1000000: accuracy of each 188-class run minus the from-scratch mean (absolute) :: sha256 c6c4138ee5c86423
\newcommand{\MlstFtJcTwoAucFull}{0.9797\,\ensuremath{\pm}\,0.0001}  % experiments/FIGS/data/v2/index.json :: fine-tuning class-balanced JetClass-II N=1000000: AUC averaged over the classes, Full :: sha256 c6c4138ee5c86423
\newcommand{\MlstFtJcTwoAucScratch}{0.9734\,\ensuremath{\pm}\,0.0003}  % experiments/FIGS/data/v2/index.json :: fine-tuning class-balanced JetClass-II N=1000000: AUC averaged over the classes, Scratch :: sha256 c6c4138ee5c86423
\newcommand{\MlstFtJcOneAccFullEFour}{0.8117\,\ensuremath{\pm}\,0.0031}  % experiments/FIGS/data/v2/index.json :: fine-tuning jc1 N=10000: accuracy, L188 :: sha256 c6c4138ee5c86423
\newcommand{\MlstFtJcOneAccResEFour}{0.8117\,\ensuremath{\pm}\,0.0013}  % experiments/FIGS/data/v2/index.json :: fine-tuning jc1 N=10000: accuracy, L162 :: sha256 c6c4138ee5c86423
\newcommand{\MlstFtJcOneAccMidEFour}{0.8086\,\ensuremath{\pm}\,0.0031}  % experiments/FIGS/data/v2/index.json :: fine-tuning jc1 N=10000: accuracy, R42_Q1 :: sha256 c6c4138ee5c86423
\newcommand{\MlstFtJcOneAccCoarseEFour}{0.7982\,\ensuremath{\pm}\,0.0023}  % experiments/FIGS/data/v2/index.json :: fine-tuning jc1 N=10000: accuracy, R16_Q1 :: sha256 c6c4138ee5c86423
\newcommand{\MlstFtJcOneAccResMassEFour}{0.8107 and 0.8098}  % experiments/FIGS/data/v2/index.json :: fine-tuning jc1 N=10000: accuracy, L162_MASS :: sha256 c6c4138ee5c86423
\newcommand{\MlstFtJcOneAccCoarseMassEFour}{0.7944 and 0.7990}  % experiments/FIGS/data/v2/index.json :: fine-tuning jc1 N=10000: accuracy, R16_Q1_MASS :: sha256 c6c4138ee5c86423
\newcommand{\MlstFtJcOneAccScratchEFour}{0.677\,\ensuremath{\pm}\,0.004}  % experiments/FIGS/data/v2/index.json :: fine-tuning jc1 N=10000: accuracy from scratch :: sha256 c6c4138ee5c86423
\newcommand{\MlstFtJcOneAccFullMinusScratchEFour}{0.1343\,\ensuremath{\pm}\,0.0031}  % experiments/FIGS/data/v2/index.json :: fine-tuning jc1 N=10000: accuracy of each 188-class run minus the from-scratch mean (absolute) :: sha256 c6c4138ee5c86423
\newcommand{\MlstFtJcOneAccFullEFive}{0.8343\,\ensuremath{\pm}\,0.0035}  % experiments/FIGS/data/v2/index.json :: fine-tuning jc1 N=100000: accuracy, L188 :: sha256 c6c4138ee5c86423
\newcommand{\MlstFtJcOneAccResEFive}{0.8328\,\ensuremath{\pm}\,0.0005}  % experiments/FIGS/data/v2/index.json :: fine-tuning jc1 N=100000: accuracy, L162 :: sha256 c6c4138ee5c86423
\newcommand{\MlstFtJcOneAccMidEFive}{0.8324\,\ensuremath{\pm}\,0.0004}  % experiments/FIGS/data/v2/index.json :: fine-tuning jc1 N=100000: accuracy, R42_Q1 :: sha256 c6c4138ee5c86423
\newcommand{\MlstFtJcOneAccCoarseEFive}{0.8288\,\ensuremath{\pm}\,0.0005}  % experiments/FIGS/data/v2/index.json :: fine-tuning jc1 N=100000: accuracy, R16_Q1 :: sha256 c6c4138ee5c86423
\newcommand{\MlstFtJcOneAccResMassEFive}{0.8326 and 0.8348}  % experiments/FIGS/data/v2/index.json :: fine-tuning jc1 N=100000: accuracy, L162_MASS :: sha256 c6c4138ee5c86423
\newcommand{\MlstFtJcOneAccCoarseMassEFive}{0.8253 and 0.8262}  % experiments/FIGS/data/v2/index.json :: fine-tuning jc1 N=100000: accuracy, R16_Q1_MASS :: sha256 c6c4138ee5c86423
\newcommand{\MlstFtJcOneAccScratchEFive}{0.777\,\ensuremath{\pm}\,0.005}  % experiments/FIGS/data/v2/index.json :: fine-tuning jc1 N=100000: accuracy from scratch :: sha256 c6c4138ee5c86423
\newcommand{\MlstFtJcOneAccFullMinusScratchEFive}{0.0571\,\ensuremath{\pm}\,0.0035}  % experiments/FIGS/data/v2/index.json :: fine-tuning jc1 N=100000: accuracy of each 188-class run minus the from-scratch mean (absolute) :: sha256 c6c4138ee5c86423
\newcommand{\MlstFtJcOneAccFullESix}{0.8468\,\ensuremath{\pm}\,0.0006}  % experiments/FIGS/data/v2/index.json :: fine-tuning jc1 N=1000000: accuracy, L188 :: sha256 c6c4138ee5c86423
\newcommand{\MlstFtJcOneAccResESix}{0.8475\,\ensuremath{\pm}\,0.0009}  % experiments/FIGS/data/v2/index.json :: fine-tuning jc1 N=1000000: accuracy, L162 :: sha256 c6c4138ee5c86423
\newcommand{\MlstFtJcOneAccMidESix}{0.8459\,\ensuremath{\pm}\,0.0002}  % experiments/FIGS/data/v2/index.json :: fine-tuning jc1 N=1000000: accuracy, R42_Q1 :: sha256 c6c4138ee5c86423
\newcommand{\MlstFtJcOneAccCoarseESix}{0.8422\,\ensuremath{\pm}\,0.0005}  % experiments/FIGS/data/v2/index.json :: fine-tuning jc1 N=1000000: accuracy, R16_Q1 :: sha256 c6c4138ee5c86423
\newcommand{\MlstFtJcOneAccResMassESix}{0.8459 and 0.8471}  % experiments/FIGS/data/v2/index.json :: fine-tuning jc1 N=1000000: accuracy, L162_MASS :: sha256 c6c4138ee5c86423
\newcommand{\MlstFtJcOneAccCoarseMassESix}{0.8400 and 0.8398}  % experiments/FIGS/data/v2/index.json :: fine-tuning jc1 N=1000000: accuracy, R16_Q1_MASS :: sha256 c6c4138ee5c86423
\newcommand{\MlstFtJcOneAccScratchESix}{0.8267\,\ensuremath{\pm}\,0.0011}  % experiments/FIGS/data/v2/index.json :: fine-tuning jc1 N=1000000: accuracy from scratch :: sha256 c6c4138ee5c86423
\newcommand{\MlstFtJcOneAccFullMinusScratchESix}{0.0201\,\ensuremath{\pm}\,0.0006}  % experiments/FIGS/data/v2/index.json :: fine-tuning jc1 N=1000000: accuracy of each 188-class run minus the from-scratch mean (absolute) :: sha256 c6c4138ee5c86423
\newcommand{\MlstFtTopRthirtyScratchEThree}{49\,\ensuremath{\pm}\,12}  % experiments/FIGS/data/v2/index.json :: fine-tuning top N=1000: R30 from scratch, every seed :: sha256 c6c4138ee5c86423
\newcommand{\MlstFtTopRthirtyFullEThree}{72\,\ensuremath{\pm}\,19}  % experiments/FIGS/data/v2/index.json :: fine-tuning top N=1000: R30, L188 :: sha256 c6c4138ee5c86423
\newcommand{\MlstFtTopRthirtyFullOverScratchEThree}{1.5\,\ensuremath{\pm}\,0.4}  % experiments/FIGS/data/v2/index.json :: fine-tuning top N=1000: R30 of each L188 run over the from-scratch mean, SD over runs and standard error of that mean in quadrature :: sha256 c6c4138ee5c86423
\newcommand{\MlstFtTopRthirtyFullOverScratchEThreePr}{equal within 50\%}  % experiments/FIGS/data/v2/index.json :: fine-tuning top N=1000: R30 of each L188 run over the from-scratch mean, SD over runs and standard error of that mean in quadrature :: sha256 c6c4138ee5c86423 :: prose twin: two significant figures, direction in words
\newcommand{\MlstFtTopRthirtyResEThree}{135\,\ensuremath{\pm}\,26}  % experiments/FIGS/data/v2/index.json :: fine-tuning top N=1000: R30, L162 :: sha256 c6c4138ee5c86423
\newcommand{\MlstFtTopRthirtyResOverScratchEThree}{2.8\,\ensuremath{\pm}\,0.7}  % experiments/FIGS/data/v2/index.json :: fine-tuning top N=1000: R30 of each L162 run over the from-scratch mean, SD over runs and standard error of that mean in quadrature :: sha256 c6c4138ee5c86423
\newcommand{\MlstFtTopRthirtyResOverScratchEThreePr}{2.8 times the rejection}  % experiments/FIGS/data/v2/index.json :: fine-tuning top N=1000: R30 of each L162 run over the from-scratch mean, SD over runs and standard error of that mean in quadrature :: sha256 c6c4138ee5c86423 :: prose twin: two significant figures, direction in words
\newcommand{\MlstFtTopRthirtyMidEThree}{104\,\ensuremath{\pm}\,20}  % experiments/FIGS/data/v2/index.json :: fine-tuning top N=1000: R30, R42_Q1 :: sha256 c6c4138ee5c86423
\newcommand{\MlstFtTopRthirtyMidOverScratchEThree}{2.1\,\ensuremath{\pm}\,0.5}  % experiments/FIGS/data/v2/index.json :: fine-tuning top N=1000: R30 of each R42_Q1 run over the from-scratch mean, SD over runs and standard error of that mean in quadrature :: sha256 c6c4138ee5c86423
\newcommand{\MlstFtTopRthirtyMidOverScratchEThreePr}{2.1 times the rejection}  % experiments/FIGS/data/v2/index.json :: fine-tuning top N=1000: R30 of each R42_Q1 run over the from-scratch mean, SD over runs and standard error of that mean in quadrature :: sha256 c6c4138ee5c86423 :: prose twin: two significant figures, direction in words
\newcommand{\MlstFtTopRthirtyCoarseEThree}{83\,\ensuremath{\pm}\,25}  % experiments/FIGS/data/v2/index.json :: fine-tuning top N=1000: R30, R16_Q1 :: sha256 c6c4138ee5c86423
\newcommand{\MlstFtTopRthirtyCoarseOverScratchEThree}{1.7\,\ensuremath{\pm}\,0.6}  % experiments/FIGS/data/v2/index.json :: fine-tuning top N=1000: R30 of each R16_Q1 run over the from-scratch mean, SD over runs and standard error of that mean in quadrature :: sha256 c6c4138ee5c86423
\newcommand{\MlstFtTopRthirtyCoarseOverScratchEThreePr}{equal within 80\%}  % experiments/FIGS/data/v2/index.json :: fine-tuning top N=1000: R30 of each R16_Q1 run over the from-scratch mean, SD over runs and standard error of that mean in quadrature :: sha256 c6c4138ee5c86423 :: prose twin: two significant figures, direction in words
\newcommand{\MlstFtTopRthirtyResMassEThree}{105.9 and 112.7}  % experiments/FIGS/data/v2/index.json :: fine-tuning top N=1000: R30, L162_MASS :: sha256 c6c4138ee5c86423
\newcommand{\MlstFtTopRthirtyResMassOverScratchEThree}{2.18 and 2.32}  % experiments/FIGS/data/v2/index.json :: fine-tuning top N=1000: R30 of each L162_MASS run over the from-scratch mean, SD over runs and standard error of that mean in quadrature :: sha256 c6c4138ee5c86423
\newcommand{\MlstFtTopRthirtyResMassOverScratchEThreePr}{at least 2.2 times the rejection}  % experiments/FIGS/data/v2/index.json :: fine-tuning top N=1000: R30 of each L162_MASS run over the from-scratch mean, SD over runs and standard error of that mean in quadrature :: sha256 c6c4138ee5c86423 :: prose twin: two significant figures, direction in words
\newcommand{\MlstFtTopRthirtyResMassOverScratchEThreePrMax}{up to 2.3 times the rejection}  % experiments/FIGS/data/v2/index.json :: fine-tuning top N=1000: R30 of each L162_MASS run over the from-scratch mean, SD over runs and standard error of that mean in quadrature :: sha256 c6c4138ee5c86423 :: prose twin of the run farther from 1
\newcommand{\MlstFtTopRthirtyCoarseMassEThree}{59.1 and 63.9}  % experiments/FIGS/data/v2/index.json :: fine-tuning top N=1000: R30, R16_Q1_MASS :: sha256 c6c4138ee5c86423
\newcommand{\MlstFtTopRthirtyCoarseMassOverScratchEThree}{1.216 and 1.314}  % experiments/FIGS/data/v2/index.json :: fine-tuning top N=1000: R30 of each R16_Q1_MASS run over the from-scratch mean, SD over runs and standard error of that mean in quadrature :: sha256 c6c4138ee5c86423
\newcommand{\MlstFtTopRthirtyCoarseMassOverScratchEThreePr}{at least 22\% more rejection}  % experiments/FIGS/data/v2/index.json :: fine-tuning top N=1000: R30 of each R16_Q1_MASS run over the from-scratch mean, SD over runs and standard error of that mean in quadrature :: sha256 c6c4138ee5c86423 :: prose twin: two significant figures, direction in words
\newcommand{\MlstFtTopRthirtyCoarseMassOverScratchEThreePrMax}{up to 31\% more rejection}  % experiments/FIGS/data/v2/index.json :: fine-tuning top N=1000: R30 of each R16_Q1_MASS run over the from-scratch mean, SD over runs and standard error of that mean in quadrature :: sha256 c6c4138ee5c86423 :: prose twin of the run farther from 1
\newcommand{\MlstFtTopRthirtyResOverFullEThree}{2.0\,\ensuremath{\pm}\,0.8}  % experiments/FIGS/data/v2/index.json :: fine-tuning top N=1000: R30 of L162 over L188, run by run :: sha256 c6c4138ee5c86423
\newcommand{\MlstFtTopRthirtyResOverFullEThreePr}{equal within 100\%}  % experiments/FIGS/data/v2/index.json :: fine-tuning top N=1000: R30 of L162 over L188, run by run :: sha256 c6c4138ee5c86423 :: prose twin: two significant figures, direction in words
\newcommand{\MlstFtTopRthirtyMidOverFullEThree}{1.46\,\ensuremath{\pm}\,0.12}  % experiments/FIGS/data/v2/index.json :: fine-tuning top N=1000: R30 of R42_Q1 over L188, run by run :: sha256 c6c4138ee5c86423
\newcommand{\MlstFtTopRthirtyMidOverFullEThreePr}{46\% more rejection}  % experiments/FIGS/data/v2/index.json :: fine-tuning top N=1000: R30 of R42_Q1 over L188, run by run :: sha256 c6c4138ee5c86423 :: prose twin: two significant figures, direction in words
\newcommand{\MlstFtTopRthirtyCoarseOverFullEThree}{1.3\,\ensuremath{\pm}\,0.7}  % experiments/FIGS/data/v2/index.json :: fine-tuning top N=1000: R30 of R16_Q1 over L188, run by run :: sha256 c6c4138ee5c86423
\newcommand{\MlstFtTopRthirtyCoarseOverFullEThreePr}{equal within 80\%}  % experiments/FIGS/data/v2/index.json :: fine-tuning top N=1000: R30 of R16_Q1 over L188, run by run :: sha256 c6c4138ee5c86423 :: prose twin: two significant figures, direction in words
\newcommand{\MlstFtTopRthirtyResMassOverResEThree}{0.73 and 1.07}  % experiments/FIGS/data/v2/index.json :: fine-tuning top N=1000: R30 of L162_MASS over L162, run by run :: sha256 c6c4138ee5c86423
\newcommand{\MlstFtTopRthirtyResMassOverResEThreePr}{equal within 30\%}  % experiments/FIGS/data/v2/index.json :: fine-tuning top N=1000: R30 of L162_MASS over L162, run by run :: sha256 c6c4138ee5c86423 :: prose twin: two significant figures, direction in words
\newcommand{\MlstFtTopRthirtyResMassOverResEThreePrMax}{equal within 30\%}  % experiments/FIGS/data/v2/index.json :: fine-tuning top N=1000: R30 of L162_MASS over L162, run by run :: sha256 c6c4138ee5c86423 :: prose twin of the run farther from 1
\newcommand{\MlstFtTopRthirtyCoarseMassOverCoarseEThree}{0.56 and 0.75}  % experiments/FIGS/data/v2/index.json :: fine-tuning top N=1000: R30 of R16_Q1_MASS over R16_Q1, run by run :: sha256 c6c4138ee5c86423
\newcommand{\MlstFtTopRthirtyCoarseMassOverCoarseEThreePr}{equal within 50\%}  % experiments/FIGS/data/v2/index.json :: fine-tuning top N=1000: R30 of R16_Q1_MASS over R16_Q1, run by run :: sha256 c6c4138ee5c86423 :: prose twin: two significant figures, direction in words
\newcommand{\MlstFtTopRthirtyCoarseMassOverCoarseEThreePrMax}{equal within 50\%}  % experiments/FIGS/data/v2/index.json :: fine-tuning top N=1000: R30 of R16_Q1_MASS over R16_Q1, run by run :: sha256 c6c4138ee5c86423 :: prose twin of the run farther from 1
\newcommand{\MlstFtTopRthirtyRandOverFullEThree}{1.3\,\ensuremath{\pm}\,0.5}  % experiments/FIGS/data/v2/index.json :: fine-tuning top N=1000: R30 of each random partition over the L188 run of its seed :: sha256 c6c4138ee5c86423
\newcommand{\MlstFtTopRthirtyRandOverFullEThreePr}{equal within 50\%}  % experiments/FIGS/data/v2/index.json :: fine-tuning top N=1000: R30 of each random partition over the L188 run of its seed :: sha256 c6c4138ee5c86423 :: prose twin: two significant figures, direction in words
\newcommand{\MlstFtTopRthirtyRandOverCoarseEThree}{0.87\,\ensuremath{\pm}\,0.12}  % experiments/FIGS/data/v2/index.json :: fine-tuning top N=1000: R30 of each random partition over the R16_Q1 run of its seed :: sha256 c6c4138ee5c86423
\newcommand{\MlstFtTopRthirtyRandOverCoarseEThreePr}{equal within 20\%}  % experiments/FIGS/data/v2/index.json :: fine-tuning top N=1000: R30 of each random partition over the R16_Q1 run of its seed :: sha256 c6c4138ee5c86423 :: prose twin: two significant figures, direction in words
\newcommand{\MlstFtTopRthirtyScratchEFour}{340\,\ensuremath{\pm}\,50}  % experiments/FIGS/data/v2/index.json :: fine-tuning top N=10000: R30 from scratch, every seed :: sha256 c6c4138ee5c86423
\newcommand{\MlstFtTopRthirtyFullEFour}{420\,\ensuremath{\pm}\,9}  % experiments/FIGS/data/v2/index.json :: fine-tuning top N=10000: R30, L188 :: sha256 c6c4138ee5c86423
\newcommand{\MlstFtTopRthirtyFullOverScratchEFour}{1.24\,\ensuremath{\pm}\,0.11}  % experiments/FIGS/data/v2/index.json :: fine-tuning top N=10000: R30 of each L188 run over the from-scratch mean, SD over runs and standard error of that mean in quadrature :: sha256 c6c4138ee5c86423
\newcommand{\MlstFtTopRthirtyFullOverScratchEFourPr}{24\% more rejection}  % experiments/FIGS/data/v2/index.json :: fine-tuning top N=10000: R30 of each L188 run over the from-scratch mean, SD over runs and standard error of that mean in quadrature :: sha256 c6c4138ee5c86423 :: prose twin: two significant figures, direction in words
\newcommand{\MlstFtTopRthirtyResEFour}{430\,\ensuremath{\pm}\,60}  % experiments/FIGS/data/v2/index.json :: fine-tuning top N=10000: R30, L162 :: sha256 c6c4138ee5c86423
\newcommand{\MlstFtTopRthirtyResOverScratchEFour}{1.27\,\ensuremath{\pm}\,0.21}  % experiments/FIGS/data/v2/index.json :: fine-tuning top N=10000: R30 of each L162 run over the from-scratch mean, SD over runs and standard error of that mean in quadrature :: sha256 c6c4138ee5c86423
\newcommand{\MlstFtTopRthirtyResOverScratchEFourPr}{equal within 30\%}  % experiments/FIGS/data/v2/index.json :: fine-tuning top N=10000: R30 of each L162 run over the from-scratch mean, SD over runs and standard error of that mean in quadrature :: sha256 c6c4138ee5c86423 :: prose twin: two significant figures, direction in words
\newcommand{\MlstFtTopRthirtyMidEFour}{455\,\ensuremath{\pm}\,30}  % experiments/FIGS/data/v2/index.json :: fine-tuning top N=10000: R30, R42_Q1 :: sha256 c6c4138ee5c86423
\newcommand{\MlstFtTopRthirtyMidOverScratchEFour}{1.34\,\ensuremath{\pm}\,0.14}  % experiments/FIGS/data/v2/index.json :: fine-tuning top N=10000: R30 of each R42_Q1 run over the from-scratch mean, SD over runs and standard error of that mean in quadrature :: sha256 c6c4138ee5c86423
\newcommand{\MlstFtTopRthirtyMidOverScratchEFourPr}{34\% more rejection}  % experiments/FIGS/data/v2/index.json :: fine-tuning top N=10000: R30 of each R42_Q1 run over the from-scratch mean, SD over runs and standard error of that mean in quadrature :: sha256 c6c4138ee5c86423 :: prose twin: two significant figures, direction in words
\newcommand{\MlstFtTopRthirtyCoarseEFour}{350\,\ensuremath{\pm}\,40}  % experiments/FIGS/data/v2/index.json :: fine-tuning top N=10000: R30, R16_Q1 :: sha256 c6c4138ee5c86423
\newcommand{\MlstFtTopRthirtyCoarseOverScratchEFour}{1.04\,\ensuremath{\pm}\,0.14}  % experiments/FIGS/data/v2/index.json :: fine-tuning top N=10000: R30 of each R16_Q1 run over the from-scratch mean, SD over runs and standard error of that mean in quadrature :: sha256 c6c4138ee5c86423
\newcommand{\MlstFtTopRthirtyCoarseOverScratchEFourPr}{equal within 20\%}  % experiments/FIGS/data/v2/index.json :: fine-tuning top N=10000: R30 of each R16_Q1 run over the from-scratch mean, SD over runs and standard error of that mean in quadrature :: sha256 c6c4138ee5c86423 :: prose twin: two significant figures, direction in words
\newcommand{\MlstFtTopRthirtyResMassEFour}{373.2 and 376.7}  % experiments/FIGS/data/v2/index.json :: fine-tuning top N=10000: R30, L162_MASS :: sha256 c6c4138ee5c86423
\newcommand{\MlstFtTopRthirtyResMassOverScratchEFour}{1.101 and 1.111}  % experiments/FIGS/data/v2/index.json :: fine-tuning top N=10000: R30 of each L162_MASS run over the from-scratch mean, SD over runs and standard error of that mean in quadrature :: sha256 c6c4138ee5c86423
\newcommand{\MlstFtTopRthirtyResMassOverScratchEFourPr}{at least 10\% more rejection}  % experiments/FIGS/data/v2/index.json :: fine-tuning top N=10000: R30 of each L162_MASS run over the from-scratch mean, SD over runs and standard error of that mean in quadrature :: sha256 c6c4138ee5c86423 :: prose twin: two significant figures, direction in words
\newcommand{\MlstFtTopRthirtyResMassOverScratchEFourPrMax}{up to 11\% more rejection}  % experiments/FIGS/data/v2/index.json :: fine-tuning top N=10000: R30 of each L162_MASS run over the from-scratch mean, SD over runs and standard error of that mean in quadrature :: sha256 c6c4138ee5c86423 :: prose twin of the run farther from 1
\newcommand{\MlstFtTopRthirtyCoarseMassEFour}{365.8 and 359.9}  % experiments/FIGS/data/v2/index.json :: fine-tuning top N=10000: R30, R16_Q1_MASS :: sha256 c6c4138ee5c86423
\newcommand{\MlstFtTopRthirtyCoarseMassOverScratchEFour}{1.079 and 1.062}  % experiments/FIGS/data/v2/index.json :: fine-tuning top N=10000: R30 of each R16_Q1_MASS run over the from-scratch mean, SD over runs and standard error of that mean in quadrature :: sha256 c6c4138ee5c86423
\newcommand{\MlstFtTopRthirtyCoarseMassOverScratchEFourPr}{at least 6.2\% more rejection}  % experiments/FIGS/data/v2/index.json :: fine-tuning top N=10000: R30 of each R16_Q1_MASS run over the from-scratch mean, SD over runs and standard error of that mean in quadrature :: sha256 c6c4138ee5c86423 :: prose twin: two significant figures, direction in words
\newcommand{\MlstFtTopRthirtyCoarseMassOverScratchEFourPrMax}{up to 7.9\% more rejection}  % experiments/FIGS/data/v2/index.json :: fine-tuning top N=10000: R30 of each R16_Q1_MASS run over the from-scratch mean, SD over runs and standard error of that mean in quadrature :: sha256 c6c4138ee5c86423 :: prose twin of the run farther from 1
\newcommand{\MlstFtTopRthirtyResOverFullEFour}{1.02\,\ensuremath{\pm}\,0.12}  % experiments/FIGS/data/v2/index.json :: fine-tuning top N=10000: R30 of L162 over L188, run by run :: sha256 c6c4138ee5c86423
\newcommand{\MlstFtTopRthirtyResOverFullEFourPr}{equal within 20\%}  % experiments/FIGS/data/v2/index.json :: fine-tuning top N=10000: R30 of L162 over L188, run by run :: sha256 c6c4138ee5c86423 :: prose twin: two significant figures, direction in words
\newcommand{\MlstFtTopRthirtyMidOverFullEFour}{1.08\,\ensuremath{\pm}\,0.08}  % experiments/FIGS/data/v2/index.json :: fine-tuning top N=10000: R30 of R42_Q1 over L188, run by run :: sha256 c6c4138ee5c86423
\newcommand{\MlstFtTopRthirtyMidOverFullEFourPr}{equal within 9\%}  % experiments/FIGS/data/v2/index.json :: fine-tuning top N=10000: R30 of R42_Q1 over L188, run by run :: sha256 c6c4138ee5c86423 :: prose twin: two significant figures, direction in words
\newcommand{\MlstFtTopRthirtyCoarseOverFullEFour}{0.84\,\ensuremath{\pm}\,0.07}  % experiments/FIGS/data/v2/index.json :: fine-tuning top N=10000: R30 of R16_Q1 over L188, run by run :: sha256 c6c4138ee5c86423
\newcommand{\MlstFtTopRthirtyCoarseOverFullEFourPr}{16\% less rejection}  % experiments/FIGS/data/v2/index.json :: fine-tuning top N=10000: R30 of R16_Q1 over L188, run by run :: sha256 c6c4138ee5c86423 :: prose twin: two significant figures, direction in words
\newcommand{\MlstFtTopRthirtyResMassOverResEFour}{1.00 and 0.76}  % experiments/FIGS/data/v2/index.json :: fine-tuning top N=10000: R30 of L162_MASS over L162, run by run :: sha256 c6c4138ee5c86423
\newcommand{\MlstFtTopRthirtyResMassOverResEFourPr}{equal within 30\%}  % experiments/FIGS/data/v2/index.json :: fine-tuning top N=10000: R30 of L162_MASS over L162, run by run :: sha256 c6c4138ee5c86423 :: prose twin: two significant figures, direction in words
\newcommand{\MlstFtTopRthirtyResMassOverResEFourPrMax}{equal within 30\%}  % experiments/FIGS/data/v2/index.json :: fine-tuning top N=10000: R30 of L162_MASS over L162, run by run :: sha256 c6c4138ee5c86423 :: prose twin of the run farther from 1
\newcommand{\MlstFtTopRthirtyCoarseMassOverCoarseEFour}{1.13 and 0.92}  % experiments/FIGS/data/v2/index.json :: fine-tuning top N=10000: R30 of R16_Q1_MASS over R16_Q1, run by run :: sha256 c6c4138ee5c86423
\newcommand{\MlstFtTopRthirtyCoarseMassOverCoarseEFourPr}{equal within 20\%}  % experiments/FIGS/data/v2/index.json :: fine-tuning top N=10000: R30 of R16_Q1_MASS over R16_Q1, run by run :: sha256 c6c4138ee5c86423 :: prose twin: two significant figures, direction in words
\newcommand{\MlstFtTopRthirtyCoarseMassOverCoarseEFourPrMax}{equal within 20\%}  % experiments/FIGS/data/v2/index.json :: fine-tuning top N=10000: R30 of R16_Q1_MASS over R16_Q1, run by run :: sha256 c6c4138ee5c86423 :: prose twin of the run farther from 1
\newcommand{\MlstFtTopRthirtyRandOverFullEFour}{0.83\,\ensuremath{\pm}\,0.10}  % experiments/FIGS/data/v2/index.json :: fine-tuning top N=10000: R30 of each random partition over the L188 run of its seed :: sha256 c6c4138ee5c86423
\newcommand{\MlstFtTopRthirtyRandOverFullEFourPr}{equal within 20\%}  % experiments/FIGS/data/v2/index.json :: fine-tuning top N=10000: R30 of each random partition over the L188 run of its seed :: sha256 c6c4138ee5c86423 :: prose twin: two significant figures, direction in words
\newcommand{\MlstFtTopRthirtyRandOverCoarseEFour}{0.97\,\ensuremath{\pm}\,0.15}  % experiments/FIGS/data/v2/index.json :: fine-tuning top N=10000: R30 of each random partition over the R16_Q1 run of its seed :: sha256 c6c4138ee5c86423
\newcommand{\MlstFtTopRthirtyRandOverCoarseEFourPr}{equal within 20\%}  % experiments/FIGS/data/v2/index.json :: fine-tuning top N=10000: R30 of each random partition over the R16_Q1 run of its seed :: sha256 c6c4138ee5c86423 :: prose twin: two significant figures, direction in words
\newcommand{\MlstFtTopRthirtyScratchEFive}{560\,\ensuremath{\pm}\,90}  % experiments/FIGS/data/v2/index.json :: fine-tuning top N=100000: R30 from scratch, every seed :: sha256 c6c4138ee5c86423
\newcommand{\MlstFtTopRthirtyFullEFive}{850\,\ensuremath{\pm}\,90}  % experiments/FIGS/data/v2/index.json :: fine-tuning top N=100000: R30, L188 :: sha256 c6c4138ee5c86423
\newcommand{\MlstFtTopRthirtyFullOverScratchEFive}{1.50\,\ensuremath{\pm}\,0.21}  % experiments/FIGS/data/v2/index.json :: fine-tuning top N=100000: R30 of each L188 run over the from-scratch mean, SD over runs and standard error of that mean in quadrature :: sha256 c6c4138ee5c86423
\newcommand{\MlstFtTopRthirtyFullOverScratchEFivePr}{50\% more rejection}  % experiments/FIGS/data/v2/index.json :: fine-tuning top N=100000: R30 of each L188 run over the from-scratch mean, SD over runs and standard error of that mean in quadrature :: sha256 c6c4138ee5c86423 :: prose twin: two significant figures, direction in words
\newcommand{\MlstFtTopRthirtyResEFive}{810\,\ensuremath{\pm}\,60}  % experiments/FIGS/data/v2/index.json :: fine-tuning top N=100000: R30, L162 :: sha256 c6c4138ee5c86423
\newcommand{\MlstFtTopRthirtyResOverScratchEFive}{1.44\,\ensuremath{\pm}\,0.18}  % experiments/FIGS/data/v2/index.json :: fine-tuning top N=100000: R30 of each L162 run over the from-scratch mean, SD over runs and standard error of that mean in quadrature :: sha256 c6c4138ee5c86423
\newcommand{\MlstFtTopRthirtyResOverScratchEFivePr}{44\% more rejection}  % experiments/FIGS/data/v2/index.json :: fine-tuning top N=100000: R30 of each L162 run over the from-scratch mean, SD over runs and standard error of that mean in quadrature :: sha256 c6c4138ee5c86423 :: prose twin: two significant figures, direction in words
\newcommand{\MlstFtTopRthirtyMidEFive}{850\,\ensuremath{\pm}\,80}  % experiments/FIGS/data/v2/index.json :: fine-tuning top N=100000: R30, R42_Q1 :: sha256 c6c4138ee5c86423
\newcommand{\MlstFtTopRthirtyMidOverScratchEFive}{1.51\,\ensuremath{\pm}\,0.20}  % experiments/FIGS/data/v2/index.json :: fine-tuning top N=100000: R30 of each R42_Q1 run over the from-scratch mean, SD over runs and standard error of that mean in quadrature :: sha256 c6c4138ee5c86423
\newcommand{\MlstFtTopRthirtyMidOverScratchEFivePr}{51\% more rejection}  % experiments/FIGS/data/v2/index.json :: fine-tuning top N=100000: R30 of each R42_Q1 run over the from-scratch mean, SD over runs and standard error of that mean in quadrature :: sha256 c6c4138ee5c86423 :: prose twin: two significant figures, direction in words
\newcommand{\MlstFtTopRthirtyCoarseEFive}{830\,\ensuremath{\pm}\,60}  % experiments/FIGS/data/v2/index.json :: fine-tuning top N=100000: R30, R16_Q1 :: sha256 c6c4138ee5c86423
\newcommand{\MlstFtTopRthirtyCoarseOverScratchEFive}{1.47\,\ensuremath{\pm}\,0.17}  % experiments/FIGS/data/v2/index.json :: fine-tuning top N=100000: R30 of each R16_Q1 run over the from-scratch mean, SD over runs and standard error of that mean in quadrature :: sha256 c6c4138ee5c86423
\newcommand{\MlstFtTopRthirtyCoarseOverScratchEFivePr}{47\% more rejection}  % experiments/FIGS/data/v2/index.json :: fine-tuning top N=100000: R30 of each R16_Q1 run over the from-scratch mean, SD over runs and standard error of that mean in quadrature :: sha256 c6c4138ee5c86423 :: prose twin: two significant figures, direction in words
\newcommand{\MlstFtTopRthirtyResMassEFive}{848 and 848}  % experiments/FIGS/data/v2/index.json :: fine-tuning top N=100000: R30, L162_MASS :: sha256 c6c4138ee5c86423
\newcommand{\MlstFtTopRthirtyResMassOverScratchEFive}{1.50 and 1.50}  % experiments/FIGS/data/v2/index.json :: fine-tuning top N=100000: R30 of each L162_MASS run over the from-scratch mean, SD over runs and standard error of that mean in quadrature :: sha256 c6c4138ee5c86423
\newcommand{\MlstFtTopRthirtyResMassOverScratchEFivePr}{at least 50\% more rejection}  % experiments/FIGS/data/v2/index.json :: fine-tuning top N=100000: R30 of each L162_MASS run over the from-scratch mean, SD over runs and standard error of that mean in quadrature :: sha256 c6c4138ee5c86423 :: prose twin: two significant figures, direction in words
\newcommand{\MlstFtTopRthirtyResMassOverScratchEFivePrMax}{up to 50\% more rejection}  % experiments/FIGS/data/v2/index.json :: fine-tuning top N=100000: R30 of each L162_MASS run over the from-scratch mean, SD over runs and standard error of that mean in quadrature :: sha256 c6c4138ee5c86423 :: prose twin of the run farther from 1
\newcommand{\MlstFtTopRthirtyCoarseMassEFive}{856 and 774}  % experiments/FIGS/data/v2/index.json :: fine-tuning top N=100000: R30, R16_Q1_MASS :: sha256 c6c4138ee5c86423
\newcommand{\MlstFtTopRthirtyCoarseMassOverScratchEFive}{1.51 and 1.37}  % experiments/FIGS/data/v2/index.json :: fine-tuning top N=100000: R30 of each R16_Q1_MASS run over the from-scratch mean, SD over runs and standard error of that mean in quadrature :: sha256 c6c4138ee5c86423
\newcommand{\MlstFtTopRthirtyCoarseMassOverScratchEFivePr}{at least 37\% more rejection}  % experiments/FIGS/data/v2/index.json :: fine-tuning top N=100000: R30 of each R16_Q1_MASS run over the from-scratch mean, SD over runs and standard error of that mean in quadrature :: sha256 c6c4138ee5c86423 :: prose twin: two significant figures, direction in words
\newcommand{\MlstFtTopRthirtyCoarseMassOverScratchEFivePrMax}{up to 51\% more rejection}  % experiments/FIGS/data/v2/index.json :: fine-tuning top N=100000: R30 of each R16_Q1_MASS run over the from-scratch mean, SD over runs and standard error of that mean in quadrature :: sha256 c6c4138ee5c86423 :: prose twin of the run farther from 1
\newcommand{\MlstFtTopRthirtyResOverFullEFive}{0.96\,\ensuremath{\pm}\,0.05}  % experiments/FIGS/data/v2/index.json :: fine-tuning top N=100000: R30 of L162 over L188, run by run :: sha256 c6c4138ee5c86423
\newcommand{\MlstFtTopRthirtyResOverFullEFivePr}{equal within 6\%}  % experiments/FIGS/data/v2/index.json :: fine-tuning top N=100000: R30 of L162 over L188, run by run :: sha256 c6c4138ee5c86423 :: prose twin: two significant figures, direction in words
\newcommand{\MlstFtTopRthirtyMidOverFullEFive}{1.02\,\ensuremath{\pm}\,0.18}  % experiments/FIGS/data/v2/index.json :: fine-tuning top N=100000: R30 of R42_Q1 over L188, run by run :: sha256 c6c4138ee5c86423
\newcommand{\MlstFtTopRthirtyMidOverFullEFivePr}{equal within 20\%}  % experiments/FIGS/data/v2/index.json :: fine-tuning top N=100000: R30 of R42_Q1 over L188, run by run :: sha256 c6c4138ee5c86423 :: prose twin: two significant figures, direction in words
\newcommand{\MlstFtTopRthirtyCoarseOverFullEFive}{0.99\,\ensuremath{\pm}\,0.17}  % experiments/FIGS/data/v2/index.json :: fine-tuning top N=100000: R30 of R16_Q1 over L188, run by run :: sha256 c6c4138ee5c86423
\newcommand{\MlstFtTopRthirtyCoarseOverFullEFivePr}{equal within 20\%}  % experiments/FIGS/data/v2/index.json :: fine-tuning top N=100000: R30 of R16_Q1 over L188, run by run :: sha256 c6c4138ee5c86423 :: prose twin: two significant figures, direction in words
\newcommand{\MlstFtTopRthirtyResMassOverResEFive}{1.00 and 1.15}  % experiments/FIGS/data/v2/index.json :: fine-tuning top N=100000: R30 of L162_MASS over L162, run by run :: sha256 c6c4138ee5c86423
\newcommand{\MlstFtTopRthirtyResMassOverResEFivePr}{equal within 20\%}  % experiments/FIGS/data/v2/index.json :: fine-tuning top N=100000: R30 of L162_MASS over L162, run by run :: sha256 c6c4138ee5c86423 :: prose twin: two significant figures, direction in words
\newcommand{\MlstFtTopRthirtyResMassOverResEFivePrMax}{equal within 20\%}  % experiments/FIGS/data/v2/index.json :: fine-tuning top N=100000: R30 of L162_MASS over L162, run by run :: sha256 c6c4138ee5c86423 :: prose twin of the run farther from 1
\newcommand{\MlstFtTopRthirtyCoarseMassOverCoarseEFive}{1.08 and 0.87}  % experiments/FIGS/data/v2/index.json :: fine-tuning top N=100000: R30 of R16_Q1_MASS over R16_Q1, run by run :: sha256 c6c4138ee5c86423
\newcommand{\MlstFtTopRthirtyCoarseMassOverCoarseEFivePr}{equal within 20\%}  % experiments/FIGS/data/v2/index.json :: fine-tuning top N=100000: R30 of R16_Q1_MASS over R16_Q1, run by run :: sha256 c6c4138ee5c86423 :: prose twin: two significant figures, direction in words
\newcommand{\MlstFtTopRthirtyCoarseMassOverCoarseEFivePrMax}{equal within 20\%}  % experiments/FIGS/data/v2/index.json :: fine-tuning top N=100000: R30 of R16_Q1_MASS over R16_Q1, run by run :: sha256 c6c4138ee5c86423 :: prose twin of the run farther from 1
\newcommand{\MlstFtTopRthirtyRandOverFullEFive}{1.01\,\ensuremath{\pm}\,0.22}  % experiments/FIGS/data/v2/index.json :: fine-tuning top N=100000: R30 of each random partition over the L188 run of its seed :: sha256 c6c4138ee5c86423
\newcommand{\MlstFtTopRthirtyRandOverFullEFivePr}{equal within 30\%}  % experiments/FIGS/data/v2/index.json :: fine-tuning top N=100000: R30 of each random partition over the L188 run of its seed :: sha256 c6c4138ee5c86423 :: prose twin: two significant figures, direction in words
\newcommand{\MlstFtTopRthirtyRandOverCoarseEFive}{0.97\,\ensuremath{\pm}\,0.10}  % experiments/FIGS/data/v2/index.json :: fine-tuning top N=100000: R30 of each random partition over the R16_Q1 run of its seed :: sha256 c6c4138ee5c86423
\newcommand{\MlstFtTopRthirtyRandOverCoarseEFivePr}{equal within 10\%}  % experiments/FIGS/data/v2/index.json :: fine-tuning top N=100000: R30 of each random partition over the R16_Q1 run of its seed :: sha256 c6c4138ee5c86423 :: prose twin: two significant figures, direction in words
\newcommand{\MlstFtTopRthirtyScratch}{1{,}412\,\ensuremath{\pm}\,20}  % experiments/FIGS/data/v2/index.json :: fine-tuning top N=1200000: R30 from scratch, every seed :: sha256 c6c4138ee5c86423
\newcommand{\MlstFtTopRthirtyFull}{1{,}800\,\ensuremath{\pm}\,60}  % experiments/FIGS/data/v2/index.json :: fine-tuning top N=1200000: R30, L188 :: sha256 c6c4138ee5c86423
\newcommand{\MlstFtTopRthirtyFullOverScratch}{1.27\,\ensuremath{\pm}\,0.04}  % experiments/FIGS/data/v2/index.json :: fine-tuning top N=1200000: R30 of each L188 run over the from-scratch mean, SD over runs and standard error of that mean in quadrature :: sha256 c6c4138ee5c86423
\newcommand{\MlstFtTopRthirtyFullOverScratchPr}{27\% more rejection}  % experiments/FIGS/data/v2/index.json :: fine-tuning top N=1200000: R30 of each L188 run over the from-scratch mean, SD over runs and standard error of that mean in quadrature :: sha256 c6c4138ee5c86423 :: prose twin: two significant figures, direction in words
\newcommand{\MlstFtTopRthirtyRes}{1{,}790\,\ensuremath{\pm}\,60}  % experiments/FIGS/data/v2/index.json :: fine-tuning top N=1200000: R30, L162 :: sha256 c6c4138ee5c86423
\newcommand{\MlstFtTopRthirtyResOverScratch}{1.27\,\ensuremath{\pm}\,0.04}  % experiments/FIGS/data/v2/index.json :: fine-tuning top N=1200000: R30 of each L162 run over the from-scratch mean, SD over runs and standard error of that mean in quadrature :: sha256 c6c4138ee5c86423
\newcommand{\MlstFtTopRthirtyResOverScratchPr}{27\% more rejection}  % experiments/FIGS/data/v2/index.json :: fine-tuning top N=1200000: R30 of each L162 run over the from-scratch mean, SD over runs and standard error of that mean in quadrature :: sha256 c6c4138ee5c86423 :: prose twin: two significant figures, direction in words
\newcommand{\MlstFtTopRthirtyMid}{1{,}780\,\ensuremath{\pm}\,150}  % experiments/FIGS/data/v2/index.json :: fine-tuning top N=1200000: R30, R42_Q1 :: sha256 c6c4138ee5c86423
\newcommand{\MlstFtTopRthirtyMidOverScratch}{1.26\,\ensuremath{\pm}\,0.11}  % experiments/FIGS/data/v2/index.json :: fine-tuning top N=1200000: R30 of each R42_Q1 run over the from-scratch mean, SD over runs and standard error of that mean in quadrature :: sha256 c6c4138ee5c86423
\newcommand{\MlstFtTopRthirtyMidOverScratchPr}{26\% more rejection}  % experiments/FIGS/data/v2/index.json :: fine-tuning top N=1200000: R30 of each R42_Q1 run over the from-scratch mean, SD over runs and standard error of that mean in quadrature :: sha256 c6c4138ee5c86423 :: prose twin: two significant figures, direction in words
\newcommand{\MlstFtTopRthirtyCoarse}{1{,}740\,\ensuremath{\pm}\,140}  % experiments/FIGS/data/v2/index.json :: fine-tuning top N=1200000: R30, R16_Q1 :: sha256 c6c4138ee5c86423
\newcommand{\MlstFtTopRthirtyCoarseOverScratch}{1.23\,\ensuremath{\pm}\,0.10}  % experiments/FIGS/data/v2/index.json :: fine-tuning top N=1200000: R30 of each R16_Q1 run over the from-scratch mean, SD over runs and standard error of that mean in quadrature :: sha256 c6c4138ee5c86423
\newcommand{\MlstFtTopRthirtyCoarseOverScratchPr}{23\% more rejection}  % experiments/FIGS/data/v2/index.json :: fine-tuning top N=1200000: R30 of each R16_Q1 run over the from-scratch mean, SD over runs and standard error of that mean in quadrature :: sha256 c6c4138ee5c86423 :: prose twin: two significant figures, direction in words
\newcommand{\MlstFtTopRthirtyResMass}{1{,}941 and 1{,}590}  % experiments/FIGS/data/v2/index.json :: fine-tuning top N=1200000: R30, L162_MASS :: sha256 c6c4138ee5c86423
\newcommand{\MlstFtTopRthirtyResMassOverScratch}{1.37 and 1.13}  % experiments/FIGS/data/v2/index.json :: fine-tuning top N=1200000: R30 of each L162_MASS run over the from-scratch mean, SD over runs and standard error of that mean in quadrature :: sha256 c6c4138ee5c86423
\newcommand{\MlstFtTopRthirtyResMassOverScratchPr}{at least 13\% more rejection}  % experiments/FIGS/data/v2/index.json :: fine-tuning top N=1200000: R30 of each L162_MASS run over the from-scratch mean, SD over runs and standard error of that mean in quadrature :: sha256 c6c4138ee5c86423 :: prose twin: two significant figures, direction in words
\newcommand{\MlstFtTopRthirtyResMassOverScratchPrMax}{up to 37\% more rejection}  % experiments/FIGS/data/v2/index.json :: fine-tuning top N=1200000: R30 of each L162_MASS run over the from-scratch mean, SD over runs and standard error of that mean in quadrature :: sha256 c6c4138ee5c86423 :: prose twin of the run farther from 1
\newcommand{\MlstFtTopRthirtyCoarseMass}{1{,}683 and 1{,}628}  % experiments/FIGS/data/v2/index.json :: fine-tuning top N=1200000: R30, R16_Q1_MASS :: sha256 c6c4138ee5c86423
\newcommand{\MlstFtTopRthirtyCoarseMassOverScratch}{1.192 and 1.153}  % experiments/FIGS/data/v2/index.json :: fine-tuning top N=1200000: R30 of each R16_Q1_MASS run over the from-scratch mean, SD over runs and standard error of that mean in quadrature :: sha256 c6c4138ee5c86423
\newcommand{\MlstFtTopRthirtyCoarseMassOverScratchPr}{at least 15\% more rejection}  % experiments/FIGS/data/v2/index.json :: fine-tuning top N=1200000: R30 of each R16_Q1_MASS run over the from-scratch mean, SD over runs and standard error of that mean in quadrature :: sha256 c6c4138ee5c86423 :: prose twin: two significant figures, direction in words
\newcommand{\MlstFtTopRthirtyCoarseMassOverScratchPrMax}{up to 19\% more rejection}  % experiments/FIGS/data/v2/index.json :: fine-tuning top N=1200000: R30 of each R16_Q1_MASS run over the from-scratch mean, SD over runs and standard error of that mean in quadrature :: sha256 c6c4138ee5c86423 :: prose twin of the run farther from 1
\newcommand{\MlstFtTopRthirtyResOverFull}{0.997\,\ensuremath{\pm}\,0.032}  % experiments/FIGS/data/v2/index.json :: fine-tuning top N=1200000: R30 of L162 over L188, run by run :: sha256 c6c4138ee5c86423
\newcommand{\MlstFtTopRthirtyResOverFullPr}{equal within 4\%}  % experiments/FIGS/data/v2/index.json :: fine-tuning top N=1200000: R30 of L162 over L188, run by run :: sha256 c6c4138ee5c86423 :: prose twin: two significant figures, direction in words
\newcommand{\MlstFtTopRthirtyMidOverFull}{0.99\,\ensuremath{\pm}\,0.07}  % experiments/FIGS/data/v2/index.json :: fine-tuning top N=1200000: R30 of R42_Q1 over L188, run by run :: sha256 c6c4138ee5c86423
\newcommand{\MlstFtTopRthirtyMidOverFullPr}{equal within 8\%}  % experiments/FIGS/data/v2/index.json :: fine-tuning top N=1200000: R30 of R42_Q1 over L188, run by run :: sha256 c6c4138ee5c86423 :: prose twin: two significant figures, direction in words
\newcommand{\MlstFtTopRthirtyCoarseOverFull}{0.97\,\ensuremath{\pm}\,0.10}  % experiments/FIGS/data/v2/index.json :: fine-tuning top N=1200000: R30 of R16_Q1 over L188, run by run :: sha256 c6c4138ee5c86423
\newcommand{\MlstFtTopRthirtyCoarseOverFullPr}{equal within 10\%}  % experiments/FIGS/data/v2/index.json :: fine-tuning top N=1200000: R30 of R16_Q1 over L188, run by run :: sha256 c6c4138ee5c86423 :: prose twin: two significant figures, direction in words
\newcommand{\MlstFtTopRthirtyResMassOverRes}{1.12 and 0.89}  % experiments/FIGS/data/v2/index.json :: fine-tuning top N=1200000: R30 of L162_MASS over L162, run by run :: sha256 c6c4138ee5c86423
\newcommand{\MlstFtTopRthirtyResMassOverResPr}{equal within 20\%}  % experiments/FIGS/data/v2/index.json :: fine-tuning top N=1200000: R30 of L162_MASS over L162, run by run :: sha256 c6c4138ee5c86423 :: prose twin: two significant figures, direction in words
\newcommand{\MlstFtTopRthirtyResMassOverResPrMax}{equal within 20\%}  % experiments/FIGS/data/v2/index.json :: fine-tuning top N=1200000: R30 of L162_MASS over L162, run by run :: sha256 c6c4138ee5c86423 :: prose twin of the run farther from 1
\newcommand{\MlstFtTopRthirtyCoarseMassOverCoarse}{0.900 and 0.919}  % experiments/FIGS/data/v2/index.json :: fine-tuning top N=1200000: R30 of R16_Q1_MASS over R16_Q1, run by run :: sha256 c6c4138ee5c86423
\newcommand{\MlstFtTopRthirtyCoarseMassOverCoarsePr}{equal within 10\%}  % experiments/FIGS/data/v2/index.json :: fine-tuning top N=1200000: R30 of R16_Q1_MASS over R16_Q1, run by run :: sha256 c6c4138ee5c86423 :: prose twin: two significant figures, direction in words
\newcommand{\MlstFtTopRthirtyCoarseMassOverCoarsePrMax}{equal within 10\%}  % experiments/FIGS/data/v2/index.json :: fine-tuning top N=1200000: R30 of R16_Q1_MASS over R16_Q1, run by run :: sha256 c6c4138ee5c86423 :: prose twin of the run farther from 1
\newcommand{\MlstFtTopRthirtyRandOverFull}{0.97\,\ensuremath{\pm}\,0.04}  % experiments/FIGS/data/v2/index.json :: fine-tuning top N=1200000: R30 of each random partition over the L188 run of its seed :: sha256 c6c4138ee5c86423
\newcommand{\MlstFtTopRthirtyRandOverFullPr}{equal within 4\%}  % experiments/FIGS/data/v2/index.json :: fine-tuning top N=1200000: R30 of each random partition over the L188 run of its seed :: sha256 c6c4138ee5c86423 :: prose twin: two significant figures, direction in words
\newcommand{\MlstFtTopRthirtyRandOverCoarse}{0.97\,\ensuremath{\pm}\,0.09}  % experiments/FIGS/data/v2/index.json :: fine-tuning top N=1200000: R30 of each random partition over the R16_Q1 run of its seed :: sha256 c6c4138ee5c86423
\newcommand{\MlstFtTopRthirtyRandOverCoarsePr}{equal within 9\%}  % experiments/FIGS/data/v2/index.json :: fine-tuning top N=1200000: R30 of each random partition over the R16_Q1 run of its seed :: sha256 c6c4138ee5c86423 :: prose twin: two significant figures, direction in words
\newcommand{\MlstFtTopRthirtyFullOverPart}{0.650\,\ensuremath{\pm}\,0.020}  % experiments/FIGS/data/v2/index.json :: fine-tuning top N=1200000: R30 of each 188-class run over ParT-f.t.'s published value :: sha256 c6c4138ee5c86423
\newcommand{\MlstFtTopRthirtyFullOverPartPr}{35\% less rejection}  % experiments/FIGS/data/v2/index.json :: fine-tuning top N=1200000: R30 of each 188-class run over ParT-f.t.'s published value :: sha256 c6c4138ee5c86423 :: prose twin: two significant figures, direction in words
\newcommand{\MlstFtTopAucFull}{0.9861\,\ensuremath{\pm}\,0.0001}  % experiments/FIGS/data/v2/index.json :: fine-tuning top N=1200000: AUC of the 188-class runs :: sha256 c6c4138ee5c86423
\newcommand{\MlstFtTopAucScratch}{0.9852\,\ensuremath{\pm}\,0.0001}  % experiments/FIGS/data/v2/index.json :: fine-tuning top N=1200000: AUC from scratch :: sha256 c6c4138ee5c86423
\newcommand{\MlstFtTopRthirtyFullRelSd}{3\%}  % experiments/FIGS/data/v2/index.json :: fine-tuning top, full set: relative SD of R30 over the 188-class runs :: sha256 c6c4138ee5c86423
\newcommand{\MlstFtTopRthirtyKin}{1{,}800}  % experiments/FIGS/data/v2/index.json :: fine-tuning top, full set: R30 of the kinematic-only model (one run) :: sha256 c6c4138ee5c86423
\newcommand{\MlstFtTopRthirtyKinOverFull}{1.00}  % experiments/FIGS/data/v2/index.json :: fine-tuning top, full set: R30 of the kinematic-only model over the mean of the 188-class runs :: sha256 c6c4138ee5c86423
\newcommand{\MlstFtTopRthirtyKinOverFullPr}{0.23\% more rejection}  % experiments/FIGS/data/v2/index.json :: fine-tuning top, full set: R30 of the kinematic-only model over the mean of the 188-class runs :: sha256 c6c4138ee5c86423 :: prose twin: two significant figures, direction in words
\newcommand{\MlstFtTopDataEquivJets}{1{,}200}  % experiments/FIGS/data/v2/index.json :: fine-tuning top: from-scratch training jets whose R30 equals the 188-class mean at 1000 jets (interpolated linearly in R30 against log10 N) :: sha256 c6c4138ee5c86423
\newcommand{\MlstFtTopDataEquivJetsEFour}{23{,}000}  % experiments/FIGS/data/v2/index.json :: fine-tuning top: from-scratch training jets whose R30 equals the 188-class mean at 10000 jets (interpolated linearly in R30 against log10 N) :: sha256 c6c4138ee5c86423
\newcommand{\MlstFtQgRthirtyScratchEThree}{49.7\,\ensuremath{\pm}\,2.1}  % experiments/FIGS/data/v2/index.json :: fine-tuning qg N=1000: R30 from scratch, every seed :: sha256 c6c4138ee5c86423
\newcommand{\MlstFtQgRthirtyFullEThree}{47\,\ensuremath{\pm}\,10}  % experiments/FIGS/data/v2/index.json :: fine-tuning qg N=1000: R30, L188 :: sha256 c6c4138ee5c86423
\newcommand{\MlstFtQgRthirtyFullOverScratchEThree}{0.95\,\ensuremath{\pm}\,0.19}  % experiments/FIGS/data/v2/index.json :: fine-tuning qg N=1000: R30 of each L188 run over the from-scratch mean, SD over runs and standard error of that mean in quadrature :: sha256 c6c4138ee5c86423
\newcommand{\MlstFtQgRthirtyFullOverScratchEThreePr}{equal within 20\%}  % experiments/FIGS/data/v2/index.json :: fine-tuning qg N=1000: R30 of each L188 run over the from-scratch mean, SD over runs and standard error of that mean in quadrature :: sha256 c6c4138ee5c86423 :: prose twin: two significant figures, direction in words
\newcommand{\MlstFtQgRthirtyResEThree}{52.4\,\ensuremath{\pm}\,3.1}  % experiments/FIGS/data/v2/index.json :: fine-tuning qg N=1000: R30, L162 :: sha256 c6c4138ee5c86423
\newcommand{\MlstFtQgRthirtyResOverScratchEThree}{1.06\,\ensuremath{\pm}\,0.07}  % experiments/FIGS/data/v2/index.json :: fine-tuning qg N=1000: R30 of each L162 run over the from-scratch mean, SD over runs and standard error of that mean in quadrature :: sha256 c6c4138ee5c86423
\newcommand{\MlstFtQgRthirtyResOverScratchEThreePr}{equal within 7\%}  % experiments/FIGS/data/v2/index.json :: fine-tuning qg N=1000: R30 of each L162 run over the from-scratch mean, SD over runs and standard error of that mean in quadrature :: sha256 c6c4138ee5c86423 :: prose twin: two significant figures, direction in words
\newcommand{\MlstFtQgRthirtyMidEThree}{54\,\ensuremath{\pm}\,6}  % experiments/FIGS/data/v2/index.json :: fine-tuning qg N=1000: R30, R42_Q1 :: sha256 c6c4138ee5c86423
\newcommand{\MlstFtQgRthirtyMidOverScratchEThree}{1.09\,\ensuremath{\pm}\,0.12}  % experiments/FIGS/data/v2/index.json :: fine-tuning qg N=1000: R30 of each R42_Q1 run over the from-scratch mean, SD over runs and standard error of that mean in quadrature :: sha256 c6c4138ee5c86423
\newcommand{\MlstFtQgRthirtyMidOverScratchEThreePr}{equal within 20\%}  % experiments/FIGS/data/v2/index.json :: fine-tuning qg N=1000: R30 of each R42_Q1 run over the from-scratch mean, SD over runs and standard error of that mean in quadrature :: sha256 c6c4138ee5c86423 :: prose twin: two significant figures, direction in words
\newcommand{\MlstFtQgRthirtyCoarseEThree}{55\,\ensuremath{\pm}\,5}  % experiments/FIGS/data/v2/index.json :: fine-tuning qg N=1000: R30, R16_Q1 :: sha256 c6c4138ee5c86423
\newcommand{\MlstFtQgRthirtyCoarseOverScratchEThree}{1.12\,\ensuremath{\pm}\,0.10}  % experiments/FIGS/data/v2/index.json :: fine-tuning qg N=1000: R30 of each R16_Q1 run over the from-scratch mean, SD over runs and standard error of that mean in quadrature :: sha256 c6c4138ee5c86423
\newcommand{\MlstFtQgRthirtyCoarseOverScratchEThreePr}{equal within 20\%}  % experiments/FIGS/data/v2/index.json :: fine-tuning qg N=1000: R30 of each R16_Q1 run over the from-scratch mean, SD over runs and standard error of that mean in quadrature :: sha256 c6c4138ee5c86423 :: prose twin: two significant figures, direction in words
\newcommand{\MlstFtQgRthirtyResMassEThree}{59.3 and 57.1}  % experiments/FIGS/data/v2/index.json :: fine-tuning qg N=1000: R30, L162_MASS :: sha256 c6c4138ee5c86423
\newcommand{\MlstFtQgRthirtyResMassOverScratchEThree}{1.193 and 1.149}  % experiments/FIGS/data/v2/index.json :: fine-tuning qg N=1000: R30 of each L162_MASS run over the from-scratch mean, SD over runs and standard error of that mean in quadrature :: sha256 c6c4138ee5c86423
\newcommand{\MlstFtQgRthirtyResMassOverScratchEThreePr}{at least 15\% more rejection}  % experiments/FIGS/data/v2/index.json :: fine-tuning qg N=1000: R30 of each L162_MASS run over the from-scratch mean, SD over runs and standard error of that mean in quadrature :: sha256 c6c4138ee5c86423 :: prose twin: two significant figures, direction in words
\newcommand{\MlstFtQgRthirtyResMassOverScratchEThreePrMax}{up to 19\% more rejection}  % experiments/FIGS/data/v2/index.json :: fine-tuning qg N=1000: R30 of each L162_MASS run over the from-scratch mean, SD over runs and standard error of that mean in quadrature :: sha256 c6c4138ee5c86423 :: prose twin of the run farther from 1
\newcommand{\MlstFtQgRthirtyCoarseMassEThree}{54.2 and 57.7}  % experiments/FIGS/data/v2/index.json :: fine-tuning qg N=1000: R30, R16_Q1_MASS :: sha256 c6c4138ee5c86423
\newcommand{\MlstFtQgRthirtyCoarseMassOverScratchEThree}{1.091 and 1.161}  % experiments/FIGS/data/v2/index.json :: fine-tuning qg N=1000: R30 of each R16_Q1_MASS run over the from-scratch mean, SD over runs and standard error of that mean in quadrature :: sha256 c6c4138ee5c86423
\newcommand{\MlstFtQgRthirtyCoarseMassOverScratchEThreePr}{at least 9.1\% more rejection}  % experiments/FIGS/data/v2/index.json :: fine-tuning qg N=1000: R30 of each R16_Q1_MASS run over the from-scratch mean, SD over runs and standard error of that mean in quadrature :: sha256 c6c4138ee5c86423 :: prose twin: two significant figures, direction in words
\newcommand{\MlstFtQgRthirtyCoarseMassOverScratchEThreePrMax}{up to 16\% more rejection}  % experiments/FIGS/data/v2/index.json :: fine-tuning qg N=1000: R30 of each R16_Q1_MASS run over the from-scratch mean, SD over runs and standard error of that mean in quadrature :: sha256 c6c4138ee5c86423 :: prose twin of the run farther from 1
\newcommand{\MlstFtQgRthirtyResOverFullEThree}{1.14\,\ensuremath{\pm}\,0.17}  % experiments/FIGS/data/v2/index.json :: fine-tuning qg N=1000: R30 of L162 over L188, run by run :: sha256 c6c4138ee5c86423
\newcommand{\MlstFtQgRthirtyResOverFullEThreePr}{equal within 20\%}  % experiments/FIGS/data/v2/index.json :: fine-tuning qg N=1000: R30 of L162 over L188, run by run :: sha256 c6c4138ee5c86423 :: prose twin: two significant figures, direction in words
\newcommand{\MlstFtQgRthirtyMidOverFullEThree}{1.17\,\ensuremath{\pm}\,0.14}  % experiments/FIGS/data/v2/index.json :: fine-tuning qg N=1000: R30 of R42_Q1 over L188, run by run :: sha256 c6c4138ee5c86423
\newcommand{\MlstFtQgRthirtyMidOverFullEThreePr}{equal within 20\%}  % experiments/FIGS/data/v2/index.json :: fine-tuning qg N=1000: R30 of R42_Q1 over L188, run by run :: sha256 c6c4138ee5c86423 :: prose twin: two significant figures, direction in words
\newcommand{\MlstFtQgRthirtyCoarseOverFullEThree}{1.21\,\ensuremath{\pm}\,0.27}  % experiments/FIGS/data/v2/index.json :: fine-tuning qg N=1000: R30 of R16_Q1 over L188, run by run :: sha256 c6c4138ee5c86423
\newcommand{\MlstFtQgRthirtyCoarseOverFullEThreePr}{equal within 30\%}  % experiments/FIGS/data/v2/index.json :: fine-tuning qg N=1000: R30 of R16_Q1 over L188, run by run :: sha256 c6c4138ee5c86423 :: prose twin: two significant figures, direction in words
\newcommand{\MlstFtQgRthirtyResMassOverResEThree}{1.108 and 1.042}  % experiments/FIGS/data/v2/index.json :: fine-tuning qg N=1000: R30 of L162_MASS over L162, run by run :: sha256 c6c4138ee5c86423
\newcommand{\MlstFtQgRthirtyResMassOverResEThreePr}{equal within 20\%}  % experiments/FIGS/data/v2/index.json :: fine-tuning qg N=1000: R30 of L162_MASS over L162, run by run :: sha256 c6c4138ee5c86423 :: prose twin: two significant figures, direction in words
\newcommand{\MlstFtQgRthirtyResMassOverResEThreePrMax}{equal within 20\%}  % experiments/FIGS/data/v2/index.json :: fine-tuning qg N=1000: R30 of L162_MASS over L162, run by run :: sha256 c6c4138ee5c86423 :: prose twin of the run farther from 1
\newcommand{\MlstFtQgRthirtyCoarseMassOverCoarseEThree}{0.89 and 1.11}  % experiments/FIGS/data/v2/index.json :: fine-tuning qg N=1000: R30 of R16_Q1_MASS over R16_Q1, run by run :: sha256 c6c4138ee5c86423
\newcommand{\MlstFtQgRthirtyCoarseMassOverCoarseEThreePr}{equal within 20\%}  % experiments/FIGS/data/v2/index.json :: fine-tuning qg N=1000: R30 of R16_Q1_MASS over R16_Q1, run by run :: sha256 c6c4138ee5c86423 :: prose twin: two significant figures, direction in words
\newcommand{\MlstFtQgRthirtyCoarseMassOverCoarseEThreePrMax}{equal within 20\%}  % experiments/FIGS/data/v2/index.json :: fine-tuning qg N=1000: R30 of R16_Q1_MASS over R16_Q1, run by run :: sha256 c6c4138ee5c86423 :: prose twin of the run farther from 1
\newcommand{\MlstFtQgRthirtyRandOverFullEThree}{0.91\,\ensuremath{\pm}\,0.07}  % experiments/FIGS/data/v2/index.json :: fine-tuning qg N=1000: R30 of each random partition over the L188 run of its seed :: sha256 c6c4138ee5c86423
\newcommand{\MlstFtQgRthirtyRandOverFullEThreePr}{equal within 10\%}  % experiments/FIGS/data/v2/index.json :: fine-tuning qg N=1000: R30 of each random partition over the L188 run of its seed :: sha256 c6c4138ee5c86423 :: prose twin: two significant figures, direction in words
\newcommand{\MlstFtQgRthirtyRandOverCoarseEThree}{0.87\,\ensuremath{\pm}\,0.14}  % experiments/FIGS/data/v2/index.json :: fine-tuning qg N=1000: R30 of each random partition over the R16_Q1 run of its seed :: sha256 c6c4138ee5c86423
\newcommand{\MlstFtQgRthirtyRandOverCoarseEThreePr}{equal within 20\%}  % experiments/FIGS/data/v2/index.json :: fine-tuning qg N=1000: R30 of each random partition over the R16_Q1 run of its seed :: sha256 c6c4138ee5c86423 :: prose twin: two significant figures, direction in words
\newcommand{\MlstFtQgRthirtyScratchEFour}{66\,\ensuremath{\pm}\,9}  % experiments/FIGS/data/v2/index.json :: fine-tuning qg N=10000: R30 from scratch, every seed :: sha256 c6c4138ee5c86423
\newcommand{\MlstFtQgRthirtyFullEFour}{75\,\ensuremath{\pm}\,7}  % experiments/FIGS/data/v2/index.json :: fine-tuning qg N=10000: R30, L188 :: sha256 c6c4138ee5c86423
\newcommand{\MlstFtQgRthirtyFullOverScratchEFour}{1.14\,\ensuremath{\pm}\,0.14}  % experiments/FIGS/data/v2/index.json :: fine-tuning qg N=10000: R30 of each L188 run over the from-scratch mean, SD over runs and standard error of that mean in quadrature :: sha256 c6c4138ee5c86423
\newcommand{\MlstFtQgRthirtyFullOverScratchEFourPr}{equal within 20\%}  % experiments/FIGS/data/v2/index.json :: fine-tuning qg N=10000: R30 of each L188 run over the from-scratch mean, SD over runs and standard error of that mean in quadrature :: sha256 c6c4138ee5c86423 :: prose twin: two significant figures, direction in words
\newcommand{\MlstFtQgRthirtyResEFour}{76.4\,\ensuremath{\pm}\,0.6}  % experiments/FIGS/data/v2/index.json :: fine-tuning qg N=10000: R30, L162 :: sha256 c6c4138ee5c86423
\newcommand{\MlstFtQgRthirtyResOverScratchEFour}{1.15\,\ensuremath{\pm}\,0.09}  % experiments/FIGS/data/v2/index.json :: fine-tuning qg N=10000: R30 of each L162 run over the from-scratch mean, SD over runs and standard error of that mean in quadrature :: sha256 c6c4138ee5c86423
\newcommand{\MlstFtQgRthirtyResOverScratchEFourPr}{15\% more rejection}  % experiments/FIGS/data/v2/index.json :: fine-tuning qg N=10000: R30 of each L162 run over the from-scratch mean, SD over runs and standard error of that mean in quadrature :: sha256 c6c4138ee5c86423 :: prose twin: two significant figures, direction in words
\newcommand{\MlstFtQgRthirtyMidEFour}{74\,\ensuremath{\pm}\,5}  % experiments/FIGS/data/v2/index.json :: fine-tuning qg N=10000: R30, R42_Q1 :: sha256 c6c4138ee5c86423
\newcommand{\MlstFtQgRthirtyMidOverScratchEFour}{1.12\,\ensuremath{\pm}\,0.11}  % experiments/FIGS/data/v2/index.json :: fine-tuning qg N=10000: R30 of each R42_Q1 run over the from-scratch mean, SD over runs and standard error of that mean in quadrature :: sha256 c6c4138ee5c86423
\newcommand{\MlstFtQgRthirtyMidOverScratchEFourPr}{equal within 20\%}  % experiments/FIGS/data/v2/index.json :: fine-tuning qg N=10000: R30 of each R42_Q1 run over the from-scratch mean, SD over runs and standard error of that mean in quadrature :: sha256 c6c4138ee5c86423 :: prose twin: two significant figures, direction in words
\newcommand{\MlstFtQgRthirtyCoarseEFour}{80.3\,\ensuremath{\pm}\,1.1}  % experiments/FIGS/data/v2/index.json :: fine-tuning qg N=10000: R30, R16_Q1 :: sha256 c6c4138ee5c86423
\newcommand{\MlstFtQgRthirtyCoarseOverScratchEFour}{1.21\,\ensuremath{\pm}\,0.10}  % experiments/FIGS/data/v2/index.json :: fine-tuning qg N=10000: R30 of each R16_Q1 run over the from-scratch mean, SD over runs and standard error of that mean in quadrature :: sha256 c6c4138ee5c86423
\newcommand{\MlstFtQgRthirtyCoarseOverScratchEFourPr}{21\% more rejection}  % experiments/FIGS/data/v2/index.json :: fine-tuning qg N=10000: R30 of each R16_Q1 run over the from-scratch mean, SD over runs and standard error of that mean in quadrature :: sha256 c6c4138ee5c86423 :: prose twin: two significant figures, direction in words
\newcommand{\MlstFtQgRthirtyResMassEFour}{80.1 and 72.6}  % experiments/FIGS/data/v2/index.json :: fine-tuning qg N=10000: R30, L162_MASS :: sha256 c6c4138ee5c86423
\newcommand{\MlstFtQgRthirtyResMassOverScratchEFour}{1.21 and 1.09}  % experiments/FIGS/data/v2/index.json :: fine-tuning qg N=10000: R30 of each L162_MASS run over the from-scratch mean, SD over runs and standard error of that mean in quadrature :: sha256 c6c4138ee5c86423
\newcommand{\MlstFtQgRthirtyResMassOverScratchEFourPr}{at least 9.3\% more rejection}  % experiments/FIGS/data/v2/index.json :: fine-tuning qg N=10000: R30 of each L162_MASS run over the from-scratch mean, SD over runs and standard error of that mean in quadrature :: sha256 c6c4138ee5c86423 :: prose twin: two significant figures, direction in words
\newcommand{\MlstFtQgRthirtyResMassOverScratchEFourPrMax}{up to 21\% more rejection}  % experiments/FIGS/data/v2/index.json :: fine-tuning qg N=10000: R30 of each L162_MASS run over the from-scratch mean, SD over runs and standard error of that mean in quadrature :: sha256 c6c4138ee5c86423 :: prose twin of the run farther from 1
\newcommand{\MlstFtQgRthirtyCoarseMassEFour}{79.5 and 77.9}  % experiments/FIGS/data/v2/index.json :: fine-tuning qg N=10000: R30, R16_Q1_MASS :: sha256 c6c4138ee5c86423
\newcommand{\MlstFtQgRthirtyCoarseMassOverScratchEFour}{1.198 and 1.173}  % experiments/FIGS/data/v2/index.json :: fine-tuning qg N=10000: R30 of each R16_Q1_MASS run over the from-scratch mean, SD over runs and standard error of that mean in quadrature :: sha256 c6c4138ee5c86423
\newcommand{\MlstFtQgRthirtyCoarseMassOverScratchEFourPr}{at least 17\% more rejection}  % experiments/FIGS/data/v2/index.json :: fine-tuning qg N=10000: R30 of each R16_Q1_MASS run over the from-scratch mean, SD over runs and standard error of that mean in quadrature :: sha256 c6c4138ee5c86423 :: prose twin: two significant figures, direction in words
\newcommand{\MlstFtQgRthirtyCoarseMassOverScratchEFourPrMax}{up to 20\% more rejection}  % experiments/FIGS/data/v2/index.json :: fine-tuning qg N=10000: R30 of each R16_Q1_MASS run over the from-scratch mean, SD over runs and standard error of that mean in quadrature :: sha256 c6c4138ee5c86423 :: prose twin of the run farther from 1
\newcommand{\MlstFtQgRthirtyResOverFullEFour}{1.02\,\ensuremath{\pm}\,0.10}  % experiments/FIGS/data/v2/index.json :: fine-tuning qg N=10000: R30 of L162 over L188, run by run :: sha256 c6c4138ee5c86423
\newcommand{\MlstFtQgRthirtyResOverFullEFourPr}{equal within 10\%}  % experiments/FIGS/data/v2/index.json :: fine-tuning qg N=10000: R30 of L162 over L188, run by run :: sha256 c6c4138ee5c86423 :: prose twin: two significant figures, direction in words
\newcommand{\MlstFtQgRthirtyMidOverFullEFour}{0.99\,\ensuremath{\pm}\,0.14}  % experiments/FIGS/data/v2/index.json :: fine-tuning qg N=10000: R30 of R42_Q1 over L188, run by run :: sha256 c6c4138ee5c86423
\newcommand{\MlstFtQgRthirtyMidOverFullEFourPr}{equal within 20\%}  % experiments/FIGS/data/v2/index.json :: fine-tuning qg N=10000: R30 of R42_Q1 over L188, run by run :: sha256 c6c4138ee5c86423 :: prose twin: two significant figures, direction in words
\newcommand{\MlstFtQgRthirtyCoarseOverFullEFour}{1.07\,\ensuremath{\pm}\,0.13}  % experiments/FIGS/data/v2/index.json :: fine-tuning qg N=10000: R30 of R16_Q1 over L188, run by run :: sha256 c6c4138ee5c86423
\newcommand{\MlstFtQgRthirtyCoarseOverFullEFourPr}{equal within 20\%}  % experiments/FIGS/data/v2/index.json :: fine-tuning qg N=10000: R30 of R16_Q1 over L188, run by run :: sha256 c6c4138ee5c86423 :: prose twin: two significant figures, direction in words
\newcommand{\MlstFtQgRthirtyResMassOverResEFour}{1.06 and 0.95}  % experiments/FIGS/data/v2/index.json :: fine-tuning qg N=10000: R30 of L162_MASS over L162, run by run :: sha256 c6c4138ee5c86423
\newcommand{\MlstFtQgRthirtyResMassOverResEFourPr}{equal within 6\%}  % experiments/FIGS/data/v2/index.json :: fine-tuning qg N=10000: R30 of L162_MASS over L162, run by run :: sha256 c6c4138ee5c86423 :: prose twin: two significant figures, direction in words
\newcommand{\MlstFtQgRthirtyResMassOverResEFourPrMax}{equal within 6\%}  % experiments/FIGS/data/v2/index.json :: fine-tuning qg N=10000: R30 of L162_MASS over L162, run by run :: sha256 c6c4138ee5c86423 :: prose twin of the run farther from 1
\newcommand{\MlstFtQgRthirtyCoarseMassOverCoarseEFour}{0.9769 and 0.9696}  % experiments/FIGS/data/v2/index.json :: fine-tuning qg N=10000: R30 of R16_Q1_MASS over R16_Q1, run by run :: sha256 c6c4138ee5c86423
\newcommand{\MlstFtQgRthirtyCoarseMassOverCoarseEFourPr}{equal within 4\%}  % experiments/FIGS/data/v2/index.json :: fine-tuning qg N=10000: R30 of R16_Q1_MASS over R16_Q1, run by run :: sha256 c6c4138ee5c86423 :: prose twin: two significant figures, direction in words
\newcommand{\MlstFtQgRthirtyCoarseMassOverCoarseEFourPrMax}{equal within 4\%}  % experiments/FIGS/data/v2/index.json :: fine-tuning qg N=10000: R30 of R16_Q1_MASS over R16_Q1, run by run :: sha256 c6c4138ee5c86423 :: prose twin of the run farther from 1
\newcommand{\MlstFtQgRthirtyRandOverFullEFour}{1.03\,\ensuremath{\pm}\,0.18}  % experiments/FIGS/data/v2/index.json :: fine-tuning qg N=10000: R30 of each random partition over the L188 run of its seed :: sha256 c6c4138ee5c86423
\newcommand{\MlstFtQgRthirtyRandOverFullEFourPr}{equal within 20\%}  % experiments/FIGS/data/v2/index.json :: fine-tuning qg N=10000: R30 of each random partition over the L188 run of its seed :: sha256 c6c4138ee5c86423 :: prose twin: two significant figures, direction in words
\newcommand{\MlstFtQgRthirtyRandOverCoarseEFour}{0.93\,\ensuremath{\pm}\,0.12}  % experiments/FIGS/data/v2/index.json :: fine-tuning qg N=10000: R30 of each random partition over the R16_Q1 run of its seed :: sha256 c6c4138ee5c86423
\newcommand{\MlstFtQgRthirtyRandOverCoarseEFourPr}{equal within 20\%}  % experiments/FIGS/data/v2/index.json :: fine-tuning qg N=10000: R30 of each random partition over the R16_Q1 run of its seed :: sha256 c6c4138ee5c86423 :: prose twin: two significant figures, direction in words
\newcommand{\MlstFtQgRthirtyScratchEFive}{88.4\,\ensuremath{\pm}\,0.8}  % experiments/FIGS/data/v2/index.json :: fine-tuning qg N=100000: R30 from scratch, every seed :: sha256 c6c4138ee5c86423
\newcommand{\MlstFtQgRthirtyFullEFive}{99.3\,\ensuremath{\pm}\,1.4}  % experiments/FIGS/data/v2/index.json :: fine-tuning qg N=100000: R30, L188 :: sha256 c6c4138ee5c86423
\newcommand{\MlstFtQgRthirtyFullOverScratchEFive}{1.124\,\ensuremath{\pm}\,0.017}  % experiments/FIGS/data/v2/index.json :: fine-tuning qg N=100000: R30 of each L188 run over the from-scratch mean, SD over runs and standard error of that mean in quadrature :: sha256 c6c4138ee5c86423
\newcommand{\MlstFtQgRthirtyFullOverScratchEFivePr}{12\% more rejection}  % experiments/FIGS/data/v2/index.json :: fine-tuning qg N=100000: R30 of each L188 run over the from-scratch mean, SD over runs and standard error of that mean in quadrature :: sha256 c6c4138ee5c86423 :: prose twin: two significant figures, direction in words
\newcommand{\MlstFtQgRthirtyResEFive}{97.9\,\ensuremath{\pm}\,3.3}  % experiments/FIGS/data/v2/index.json :: fine-tuning qg N=100000: R30, L162 :: sha256 c6c4138ee5c86423
\newcommand{\MlstFtQgRthirtyResOverScratchEFive}{1.11\,\ensuremath{\pm}\,0.04}  % experiments/FIGS/data/v2/index.json :: fine-tuning qg N=100000: R30 of each L162 run over the from-scratch mean, SD over runs and standard error of that mean in quadrature :: sha256 c6c4138ee5c86423
\newcommand{\MlstFtQgRthirtyResOverScratchEFivePr}{11\% more rejection}  % experiments/FIGS/data/v2/index.json :: fine-tuning qg N=100000: R30 of each L162 run over the from-scratch mean, SD over runs and standard error of that mean in quadrature :: sha256 c6c4138ee5c86423 :: prose twin: two significant figures, direction in words
\newcommand{\MlstFtQgRthirtyMidEFive}{96.1\,\ensuremath{\pm}\,1.8}  % experiments/FIGS/data/v2/index.json :: fine-tuning qg N=100000: R30, R42_Q1 :: sha256 c6c4138ee5c86423
\newcommand{\MlstFtQgRthirtyMidOverScratchEFive}{1.088\,\ensuremath{\pm}\,0.021}  % experiments/FIGS/data/v2/index.json :: fine-tuning qg N=100000: R30 of each R42_Q1 run over the from-scratch mean, SD over runs and standard error of that mean in quadrature :: sha256 c6c4138ee5c86423
\newcommand{\MlstFtQgRthirtyMidOverScratchEFivePr}{8.8\% more rejection}  % experiments/FIGS/data/v2/index.json :: fine-tuning qg N=100000: R30 of each R42_Q1 run over the from-scratch mean, SD over runs and standard error of that mean in quadrature :: sha256 c6c4138ee5c86423 :: prose twin: two significant figures, direction in words
\newcommand{\MlstFtQgRthirtyCoarseEFive}{96.4\,\ensuremath{\pm}\,1.0}  % experiments/FIGS/data/v2/index.json :: fine-tuning qg N=100000: R30, R16_Q1 :: sha256 c6c4138ee5c86423
\newcommand{\MlstFtQgRthirtyCoarseOverScratchEFive}{1.091\,\ensuremath{\pm}\,0.013}  % experiments/FIGS/data/v2/index.json :: fine-tuning qg N=100000: R30 of each R16_Q1 run over the from-scratch mean, SD over runs and standard error of that mean in quadrature :: sha256 c6c4138ee5c86423
\newcommand{\MlstFtQgRthirtyCoarseOverScratchEFivePr}{9.1\% more rejection}  % experiments/FIGS/data/v2/index.json :: fine-tuning qg N=100000: R30 of each R16_Q1 run over the from-scratch mean, SD over runs and standard error of that mean in quadrature :: sha256 c6c4138ee5c86423 :: prose twin: two significant figures, direction in words
\newcommand{\MlstFtQgRthirtyResMassEFive}{95.6 and 97.6}  % experiments/FIGS/data/v2/index.json :: fine-tuning qg N=100000: R30, L162_MASS :: sha256 c6c4138ee5c86423
\newcommand{\MlstFtQgRthirtyResMassOverScratchEFive}{1.082 and 1.104}  % experiments/FIGS/data/v2/index.json :: fine-tuning qg N=100000: R30 of each L162_MASS run over the from-scratch mean, SD over runs and standard error of that mean in quadrature :: sha256 c6c4138ee5c86423
\newcommand{\MlstFtQgRthirtyResMassOverScratchEFivePr}{at least 8.2\% more rejection}  % experiments/FIGS/data/v2/index.json :: fine-tuning qg N=100000: R30 of each L162_MASS run over the from-scratch mean, SD over runs and standard error of that mean in quadrature :: sha256 c6c4138ee5c86423 :: prose twin: two significant figures, direction in words
\newcommand{\MlstFtQgRthirtyResMassOverScratchEFivePrMax}{up to 10\% more rejection}  % experiments/FIGS/data/v2/index.json :: fine-tuning qg N=100000: R30 of each L162_MASS run over the from-scratch mean, SD over runs and standard error of that mean in quadrature :: sha256 c6c4138ee5c86423 :: prose twin of the run farther from 1
\newcommand{\MlstFtQgRthirtyCoarseMassEFive}{97.85 and 97.94}  % experiments/FIGS/data/v2/index.json :: fine-tuning qg N=100000: R30, R16_Q1_MASS :: sha256 c6c4138ee5c86423
\newcommand{\MlstFtQgRthirtyCoarseMassOverScratchEFive}{1.107 and 1.108}  % experiments/FIGS/data/v2/index.json :: fine-tuning qg N=100000: R30 of each R16_Q1_MASS run over the from-scratch mean, SD over runs and standard error of that mean in quadrature :: sha256 c6c4138ee5c86423
\newcommand{\MlstFtQgRthirtyCoarseMassOverScratchEFivePr}{at least 11\% more rejection}  % experiments/FIGS/data/v2/index.json :: fine-tuning qg N=100000: R30 of each R16_Q1_MASS run over the from-scratch mean, SD over runs and standard error of that mean in quadrature :: sha256 c6c4138ee5c86423 :: prose twin: two significant figures, direction in words
\newcommand{\MlstFtQgRthirtyCoarseMassOverScratchEFivePrMax}{up to 11\% more rejection}  % experiments/FIGS/data/v2/index.json :: fine-tuning qg N=100000: R30 of each R16_Q1_MASS run over the from-scratch mean, SD over runs and standard error of that mean in quadrature :: sha256 c6c4138ee5c86423 :: prose twin of the run farther from 1
\newcommand{\MlstFtQgRthirtyResOverFullEFive}{0.985\,\ensuremath{\pm}\,0.020}  % experiments/FIGS/data/v2/index.json :: fine-tuning qg N=100000: R30 of L162 over L188, run by run :: sha256 c6c4138ee5c86423
\newcommand{\MlstFtQgRthirtyResOverFullEFivePr}{equal within 2\%}  % experiments/FIGS/data/v2/index.json :: fine-tuning qg N=100000: R30 of L162 over L188, run by run :: sha256 c6c4138ee5c86423 :: prose twin: two significant figures, direction in words
\newcommand{\MlstFtQgRthirtyMidOverFullEFive}{0.969\,\ensuremath{\pm}\,0.030}  % experiments/FIGS/data/v2/index.json :: fine-tuning qg N=100000: R30 of R42_Q1 over L188, run by run :: sha256 c6c4138ee5c86423
\newcommand{\MlstFtQgRthirtyMidOverFullEFivePr}{equal within 4\%}  % experiments/FIGS/data/v2/index.json :: fine-tuning qg N=100000: R30 of R42_Q1 over L188, run by run :: sha256 c6c4138ee5c86423 :: prose twin: two significant figures, direction in words
\newcommand{\MlstFtQgRthirtyCoarseOverFullEFive}{0.972\,\ensuremath{\pm}\,0.020}  % experiments/FIGS/data/v2/index.json :: fine-tuning qg N=100000: R30 of R16_Q1 over L188, run by run :: sha256 c6c4138ee5c86423
\newcommand{\MlstFtQgRthirtyCoarseOverFullEFivePr}{equal within 3\%}  % experiments/FIGS/data/v2/index.json :: fine-tuning qg N=100000: R30 of R16_Q1 over L188, run by run :: sha256 c6c4138ee5c86423 :: prose twin: two significant figures, direction in words
\newcommand{\MlstFtQgRthirtyResMassOverResEFive}{0.942 and 1.027}  % experiments/FIGS/data/v2/index.json :: fine-tuning qg N=100000: R30 of L162_MASS over L162, run by run :: sha256 c6c4138ee5c86423
\newcommand{\MlstFtQgRthirtyResMassOverResEFivePr}{equal within 6\%}  % experiments/FIGS/data/v2/index.json :: fine-tuning qg N=100000: R30 of L162_MASS over L162, run by run :: sha256 c6c4138ee5c86423 :: prose twin: two significant figures, direction in words
\newcommand{\MlstFtQgRthirtyResMassOverResEFivePrMax}{equal within 6\%}  % experiments/FIGS/data/v2/index.json :: fine-tuning qg N=100000: R30 of L162_MASS over L162, run by run :: sha256 c6c4138ee5c86423 :: prose twin of the run farther from 1
\newcommand{\MlstFtQgRthirtyCoarseMassOverCoarseEFive}{1.018 and 1.004}  % experiments/FIGS/data/v2/index.json :: fine-tuning qg N=100000: R30 of R16_Q1_MASS over R16_Q1, run by run :: sha256 c6c4138ee5c86423
\newcommand{\MlstFtQgRthirtyCoarseMassOverCoarseEFivePr}{equal within 2\%}  % experiments/FIGS/data/v2/index.json :: fine-tuning qg N=100000: R30 of R16_Q1_MASS over R16_Q1, run by run :: sha256 c6c4138ee5c86423 :: prose twin: two significant figures, direction in words
\newcommand{\MlstFtQgRthirtyCoarseMassOverCoarseEFivePrMax}{equal within 2\%}  % experiments/FIGS/data/v2/index.json :: fine-tuning qg N=100000: R30 of R16_Q1_MASS over R16_Q1, run by run :: sha256 c6c4138ee5c86423 :: prose twin of the run farther from 1
\newcommand{\MlstFtQgRthirtyRandOverFullEFive}{0.986\,\ensuremath{\pm}\,0.025}  % experiments/FIGS/data/v2/index.json :: fine-tuning qg N=100000: R30 of each random partition over the L188 run of its seed :: sha256 c6c4138ee5c86423
\newcommand{\MlstFtQgRthirtyRandOverFullEFivePr}{equal within 3\%}  % experiments/FIGS/data/v2/index.json :: fine-tuning qg N=100000: R30 of each random partition over the L188 run of its seed :: sha256 c6c4138ee5c86423 :: prose twin: two significant figures, direction in words
\newcommand{\MlstFtQgRthirtyRandOverCoarseEFive}{1.009\,\ensuremath{\pm}\,0.020}  % experiments/FIGS/data/v2/index.json :: fine-tuning qg N=100000: R30 of each random partition over the R16_Q1 run of its seed :: sha256 c6c4138ee5c86423
\newcommand{\MlstFtQgRthirtyRandOverCoarseEFivePr}{equal within 3\%}  % experiments/FIGS/data/v2/index.json :: fine-tuning qg N=100000: R30 of each random partition over the R16_Q1 run of its seed :: sha256 c6c4138ee5c86423 :: prose twin: two significant figures, direction in words
\newcommand{\MlstFtQgRthirtyScratch}{104.0\,\ensuremath{\pm}\,1.0}  % experiments/FIGS/data/v2/index.json :: fine-tuning qg N=1600000: R30 from scratch, every seed :: sha256 c6c4138ee5c86423
\newcommand{\MlstFtQgRthirtyFull}{107.3\,\ensuremath{\pm}\,2.2}  % experiments/FIGS/data/v2/index.json :: fine-tuning qg N=1600000: R30, L188 :: sha256 c6c4138ee5c86423
\newcommand{\MlstFtQgRthirtyFullOverScratch}{1.031\,\ensuremath{\pm}\,0.022}  % experiments/FIGS/data/v2/index.json :: fine-tuning qg N=1600000: R30 of each L188 run over the from-scratch mean, SD over runs and standard error of that mean in quadrature :: sha256 c6c4138ee5c86423
\newcommand{\MlstFtQgRthirtyFullOverScratchPr}{equal within 4\%}  % experiments/FIGS/data/v2/index.json :: fine-tuning qg N=1600000: R30 of each L188 run over the from-scratch mean, SD over runs and standard error of that mean in quadrature :: sha256 c6c4138ee5c86423 :: prose twin: two significant figures, direction in words
\newcommand{\MlstFtQgRthirtyRes}{108.0\,\ensuremath{\pm}\,1.4}  % experiments/FIGS/data/v2/index.json :: fine-tuning qg N=1600000: R30, L162 :: sha256 c6c4138ee5c86423
\newcommand{\MlstFtQgRthirtyResOverScratch}{1.038\,\ensuremath{\pm}\,0.015}  % experiments/FIGS/data/v2/index.json :: fine-tuning qg N=1600000: R30 of each L162 run over the from-scratch mean, SD over runs and standard error of that mean in quadrature :: sha256 c6c4138ee5c86423
\newcommand{\MlstFtQgRthirtyResOverScratchPr}{3.8\% more rejection}  % experiments/FIGS/data/v2/index.json :: fine-tuning qg N=1600000: R30 of each L162 run over the from-scratch mean, SD over runs and standard error of that mean in quadrature :: sha256 c6c4138ee5c86423 :: prose twin: two significant figures, direction in words
\newcommand{\MlstFtQgRthirtyMid}{108.3\,\ensuremath{\pm}\,1.0}  % experiments/FIGS/data/v2/index.json :: fine-tuning qg N=1600000: R30, R42_Q1 :: sha256 c6c4138ee5c86423
\newcommand{\MlstFtQgRthirtyMidOverScratch}{1.041\,\ensuremath{\pm}\,0.012}  % experiments/FIGS/data/v2/index.json :: fine-tuning qg N=1600000: R30 of each R42_Q1 run over the from-scratch mean, SD over runs and standard error of that mean in quadrature :: sha256 c6c4138ee5c86423
\newcommand{\MlstFtQgRthirtyMidOverScratchPr}{4.1\% more rejection}  % experiments/FIGS/data/v2/index.json :: fine-tuning qg N=1600000: R30 of each R42_Q1 run over the from-scratch mean, SD over runs and standard error of that mean in quadrature :: sha256 c6c4138ee5c86423 :: prose twin: two significant figures, direction in words
\newcommand{\MlstFtQgRthirtyCoarse}{106.5\,\ensuremath{\pm}\,0.8}  % experiments/FIGS/data/v2/index.json :: fine-tuning qg N=1600000: R30, R16_Q1 :: sha256 c6c4138ee5c86423
\newcommand{\MlstFtQgRthirtyCoarseOverScratch}{1.024\,\ensuremath{\pm}\,0.010}  % experiments/FIGS/data/v2/index.json :: fine-tuning qg N=1600000: R30 of each R16_Q1 run over the from-scratch mean, SD over runs and standard error of that mean in quadrature :: sha256 c6c4138ee5c86423
\newcommand{\MlstFtQgRthirtyCoarseOverScratchPr}{2.4\% more rejection}  % experiments/FIGS/data/v2/index.json :: fine-tuning qg N=1600000: R30 of each R16_Q1 run over the from-scratch mean, SD over runs and standard error of that mean in quadrature :: sha256 c6c4138ee5c86423 :: prose twin: two significant figures, direction in words
\newcommand{\MlstFtQgRthirtyResMass}{107.0 and 106.6}  % experiments/FIGS/data/v2/index.json :: fine-tuning qg N=1600000: R30, L162_MASS :: sha256 c6c4138ee5c86423
\newcommand{\MlstFtQgRthirtyResMassOverScratch}{1.028 and 1.025}  % experiments/FIGS/data/v2/index.json :: fine-tuning qg N=1600000: R30 of each L162_MASS run over the from-scratch mean, SD over runs and standard error of that mean in quadrature :: sha256 c6c4138ee5c86423
\newcommand{\MlstFtQgRthirtyResMassOverScratchPr}{at least 2.5\% more rejection}  % experiments/FIGS/data/v2/index.json :: fine-tuning qg N=1600000: R30 of each L162_MASS run over the from-scratch mean, SD over runs and standard error of that mean in quadrature :: sha256 c6c4138ee5c86423 :: prose twin: two significant figures, direction in words
\newcommand{\MlstFtQgRthirtyResMassOverScratchPrMax}{up to 2.8\% more rejection}  % experiments/FIGS/data/v2/index.json :: fine-tuning qg N=1600000: R30 of each L162_MASS run over the from-scratch mean, SD over runs and standard error of that mean in quadrature :: sha256 c6c4138ee5c86423 :: prose twin of the run farther from 1
\newcommand{\MlstFtQgRthirtyCoarseMass}{104.9 and 105.6}  % experiments/FIGS/data/v2/index.json :: fine-tuning qg N=1600000: R30, R16_Q1_MASS :: sha256 c6c4138ee5c86423
\newcommand{\MlstFtQgRthirtyCoarseMassOverScratch}{1.009 and 1.015}  % experiments/FIGS/data/v2/index.json :: fine-tuning qg N=1600000: R30 of each R16_Q1_MASS run over the from-scratch mean, SD over runs and standard error of that mean in quadrature :: sha256 c6c4138ee5c86423
\newcommand{\MlstFtQgRthirtyCoarseMassOverScratchPr}{at least 0.87\% more rejection}  % experiments/FIGS/data/v2/index.json :: fine-tuning qg N=1600000: R30 of each R16_Q1_MASS run over the from-scratch mean, SD over runs and standard error of that mean in quadrature :: sha256 c6c4138ee5c86423 :: prose twin: two significant figures, direction in words
\newcommand{\MlstFtQgRthirtyCoarseMassOverScratchPrMax}{up to 1.5\% more rejection}  % experiments/FIGS/data/v2/index.json :: fine-tuning qg N=1600000: R30 of each R16_Q1_MASS run over the from-scratch mean, SD over runs and standard error of that mean in quadrature :: sha256 c6c4138ee5c86423 :: prose twin of the run farther from 1
\newcommand{\MlstFtQgRthirtyResOverFull}{1.007\,\ensuremath{\pm}\,0.026}  % experiments/FIGS/data/v2/index.json :: fine-tuning qg N=1600000: R30 of L162 over L188, run by run :: sha256 c6c4138ee5c86423
\newcommand{\MlstFtQgRthirtyResOverFullPr}{equal within 3\%}  % experiments/FIGS/data/v2/index.json :: fine-tuning qg N=1600000: R30 of L162 over L188, run by run :: sha256 c6c4138ee5c86423 :: prose twin: two significant figures, direction in words
\newcommand{\MlstFtQgRthirtyMidOverFull}{1.010\,\ensuremath{\pm}\,0.029}  % experiments/FIGS/data/v2/index.json :: fine-tuning qg N=1600000: R30 of R42_Q1 over L188, run by run :: sha256 c6c4138ee5c86423
\newcommand{\MlstFtQgRthirtyMidOverFullPr}{equal within 3\%}  % experiments/FIGS/data/v2/index.json :: fine-tuning qg N=1600000: R30 of R42_Q1 over L188, run by run :: sha256 c6c4138ee5c86423 :: prose twin: two significant figures, direction in words
\newcommand{\MlstFtQgRthirtyCoarseOverFull}{0.993\,\ensuremath{\pm}\,0.026}  % experiments/FIGS/data/v2/index.json :: fine-tuning qg N=1600000: R30 of R16_Q1 over L188, run by run :: sha256 c6c4138ee5c86423
\newcommand{\MlstFtQgRthirtyCoarseOverFullPr}{equal within 3\%}  % experiments/FIGS/data/v2/index.json :: fine-tuning qg N=1600000: R30 of R16_Q1 over L188, run by run :: sha256 c6c4138ee5c86423 :: prose twin: two significant figures, direction in words
\newcommand{\MlstFtQgRthirtyResMassOverRes}{0.976 and 0.999}  % experiments/FIGS/data/v2/index.json :: fine-tuning qg N=1600000: R30 of L162_MASS over L162, run by run :: sha256 c6c4138ee5c86423
\newcommand{\MlstFtQgRthirtyResMassOverResPr}{equal within 3\%}  % experiments/FIGS/data/v2/index.json :: fine-tuning qg N=1600000: R30 of L162_MASS over L162, run by run :: sha256 c6c4138ee5c86423 :: prose twin: two significant figures, direction in words
\newcommand{\MlstFtQgRthirtyResMassOverResPrMax}{equal within 3\%}  % experiments/FIGS/data/v2/index.json :: fine-tuning qg N=1600000: R30 of L162_MASS over L162, run by run :: sha256 c6c4138ee5c86423 :: prose twin of the run farther from 1
\newcommand{\MlstFtQgRthirtyCoarseMassOverCoarse}{0.9916 and 0.9926}  % experiments/FIGS/data/v2/index.json :: fine-tuning qg N=1600000: R30 of R16_Q1_MASS over R16_Q1, run by run :: sha256 c6c4138ee5c86423
\newcommand{\MlstFtQgRthirtyCoarseMassOverCoarsePr}{equal within 1\%}  % experiments/FIGS/data/v2/index.json :: fine-tuning qg N=1600000: R30 of R16_Q1_MASS over R16_Q1, run by run :: sha256 c6c4138ee5c86423 :: prose twin: two significant figures, direction in words
\newcommand{\MlstFtQgRthirtyCoarseMassOverCoarsePrMax}{equal within 1\%}  % experiments/FIGS/data/v2/index.json :: fine-tuning qg N=1600000: R30 of R16_Q1_MASS over R16_Q1, run by run :: sha256 c6c4138ee5c86423 :: prose twin of the run farther from 1
\newcommand{\MlstFtQgRthirtyRandOverFull}{0.985\,\ensuremath{\pm}\,0.011}  % experiments/FIGS/data/v2/index.json :: fine-tuning qg N=1600000: R30 of each random partition over the L188 run of its seed :: sha256 c6c4138ee5c86423
\newcommand{\MlstFtQgRthirtyRandOverFullPr}{equal within 2\%}  % experiments/FIGS/data/v2/index.json :: fine-tuning qg N=1600000: R30 of each random partition over the L188 run of its seed :: sha256 c6c4138ee5c86423 :: prose twin: two significant figures, direction in words
\newcommand{\MlstFtQgRthirtyRandOverCoarse}{1.008\,\ensuremath{\pm}\,0.009}  % experiments/FIGS/data/v2/index.json :: fine-tuning qg N=1600000: R30 of each random partition over the R16_Q1 run of its seed :: sha256 c6c4138ee5c86423
\newcommand{\MlstFtQgRthirtyRandOverCoarsePr}{equal within 1\%}  % experiments/FIGS/data/v2/index.json :: fine-tuning qg N=1600000: R30 of each random partition over the R16_Q1 run of its seed :: sha256 c6c4138ee5c86423 :: prose twin: two significant figures, direction in words
\newcommand{\MlstFtQgRthirtyFullOverPart}{0.994\,\ensuremath{\pm}\,0.020}  % experiments/FIGS/data/v2/index.json :: fine-tuning qg N=1600000: R30 of each 188-class run over ParT-f.t.'s published value :: sha256 c6c4138ee5c86423
\newcommand{\MlstFtQgRthirtyFullOverPartPr}{equal within 3\%}  % experiments/FIGS/data/v2/index.json :: fine-tuning qg N=1600000: R30 of each 188-class run over ParT-f.t.'s published value :: sha256 c6c4138ee5c86423 :: prose twin: two significant figures, direction in words
\newcommand{\MlstFtQgAucFull}{0.9148\,\ensuremath{\pm}\,0.0001}  % experiments/FIGS/data/v2/index.json :: fine-tuning qg N=1600000: AUC of the 188-class runs :: sha256 c6c4138ee5c86423
\newcommand{\MlstFtQgAucScratch}{0.9121\,\ensuremath{\pm}\,0.0002}  % experiments/FIGS/data/v2/index.json :: fine-tuning qg N=1600000: AUC from scratch :: sha256 c6c4138ee5c86423
\newcommand{\MlstFtQgDataEquivJets}{\ensuremath{<}1{,}000}  % experiments/FIGS/data/v2/index.json :: fine-tuning qg: from-scratch training jets whose R30 equals the 188-class mean at 1000 jets (interpolated linearly in R30 against log10 N) :: sha256 c6c4138ee5c86423
\newcommand{\MlstFtQgDataEquivJetsEFour}{26{,}000}  % experiments/FIGS/data/v2/index.json :: fine-tuning qg: from-scratch training jets whose R30 equals the 188-class mean at 10000 jets (interpolated linearly in R30 against log10 N) :: sha256 c6c4138ee5c86423
\newcommand{\MlstFtQgHerwigRthirtyScratchEThree}{27.0\,\ensuremath{\pm}\,1.3}  % experiments/FIGS/data/v2/index.json :: fine-tuning qg_herwig N=1000: R30 from scratch, every seed :: sha256 c6c4138ee5c86423
\newcommand{\MlstFtQgHerwigRthirtyFullEThree}{21\,\ensuremath{\pm}\,6}  % experiments/FIGS/data/v2/index.json :: fine-tuning qg_herwig N=1000: R30, L188 :: sha256 c6c4138ee5c86423
\newcommand{\MlstFtQgHerwigRthirtyFullOverScratchEThree}{0.77\,\ensuremath{\pm}\,0.22}  % experiments/FIGS/data/v2/index.json :: fine-tuning qg_herwig N=1000: R30 of each L188 run over the from-scratch mean, SD over runs and standard error of that mean in quadrature :: sha256 c6c4138ee5c86423
\newcommand{\MlstFtQgHerwigRthirtyFullOverScratchEThreePr}{equal within 30\%}  % experiments/FIGS/data/v2/index.json :: fine-tuning qg_herwig N=1000: R30 of each L188 run over the from-scratch mean, SD over runs and standard error of that mean in quadrature :: sha256 c6c4138ee5c86423 :: prose twin: two significant figures, direction in words
\newcommand{\MlstFtQgHerwigRthirtyResEThree}{26.2\,\ensuremath{\pm}\,0.9}  % experiments/FIGS/data/v2/index.json :: fine-tuning qg_herwig N=1000: R30, L162 :: sha256 c6c4138ee5c86423
\newcommand{\MlstFtQgHerwigRthirtyResOverScratchEThree}{0.97\,\ensuremath{\pm}\,0.04}  % experiments/FIGS/data/v2/index.json :: fine-tuning qg_herwig N=1000: R30 of each L162 run over the from-scratch mean, SD over runs and standard error of that mean in quadrature :: sha256 c6c4138ee5c86423
\newcommand{\MlstFtQgHerwigRthirtyResOverScratchEThreePr}{equal within 4\%}  % experiments/FIGS/data/v2/index.json :: fine-tuning qg_herwig N=1000: R30 of each L162 run over the from-scratch mean, SD over runs and standard error of that mean in quadrature :: sha256 c6c4138ee5c86423 :: prose twin: two significant figures, direction in words
\newcommand{\MlstFtQgHerwigRthirtyMidEThree}{26.9\,\ensuremath{\pm}\,2.7}  % experiments/FIGS/data/v2/index.json :: fine-tuning qg_herwig N=1000: R30, R42_Q1 :: sha256 c6c4138ee5c86423
\newcommand{\MlstFtQgHerwigRthirtyMidOverScratchEThree}{1.00\,\ensuremath{\pm}\,0.10}  % experiments/FIGS/data/v2/index.json :: fine-tuning qg_herwig N=1000: R30 of each R42_Q1 run over the from-scratch mean, SD over runs and standard error of that mean in quadrature :: sha256 c6c4138ee5c86423
\newcommand{\MlstFtQgHerwigRthirtyMidOverScratchEThreePr}{equal within 10\%}  % experiments/FIGS/data/v2/index.json :: fine-tuning qg_herwig N=1000: R30 of each R42_Q1 run over the from-scratch mean, SD over runs and standard error of that mean in quadrature :: sha256 c6c4138ee5c86423 :: prose twin: two significant figures, direction in words
\newcommand{\MlstFtQgHerwigRthirtyCoarseEThree}{25.2\,\ensuremath{\pm}\,2.5}  % experiments/FIGS/data/v2/index.json :: fine-tuning qg_herwig N=1000: R30, R16_Q1 :: sha256 c6c4138ee5c86423
\newcommand{\MlstFtQgHerwigRthirtyCoarseOverScratchEThree}{0.94\,\ensuremath{\pm}\,0.10}  % experiments/FIGS/data/v2/index.json :: fine-tuning qg_herwig N=1000: R30 of each R16_Q1 run over the from-scratch mean, SD over runs and standard error of that mean in quadrature :: sha256 c6c4138ee5c86423
\newcommand{\MlstFtQgHerwigRthirtyCoarseOverScratchEThreePr}{equal within 10\%}  % experiments/FIGS/data/v2/index.json :: fine-tuning qg_herwig N=1000: R30 of each R16_Q1 run over the from-scratch mean, SD over runs and standard error of that mean in quadrature :: sha256 c6c4138ee5c86423 :: prose twin: two significant figures, direction in words
\newcommand{\MlstFtQgHerwigRthirtyResMassEThree}{28.0 and 29.0}  % experiments/FIGS/data/v2/index.json :: fine-tuning qg_herwig N=1000: R30, L162_MASS :: sha256 c6c4138ee5c86423
\newcommand{\MlstFtQgHerwigRthirtyResMassOverScratchEThree}{1.038 and 1.076}  % experiments/FIGS/data/v2/index.json :: fine-tuning qg_herwig N=1000: R30 of each L162_MASS run over the from-scratch mean, SD over runs and standard error of that mean in quadrature :: sha256 c6c4138ee5c86423
\newcommand{\MlstFtQgHerwigRthirtyResMassOverScratchEThreePr}{at least 3.8\% more rejection}  % experiments/FIGS/data/v2/index.json :: fine-tuning qg_herwig N=1000: R30 of each L162_MASS run over the from-scratch mean, SD over runs and standard error of that mean in quadrature :: sha256 c6c4138ee5c86423 :: prose twin: two significant figures, direction in words
\newcommand{\MlstFtQgHerwigRthirtyResMassOverScratchEThreePrMax}{up to 7.6\% more rejection}  % experiments/FIGS/data/v2/index.json :: fine-tuning qg_herwig N=1000: R30 of each L162_MASS run over the from-scratch mean, SD over runs and standard error of that mean in quadrature :: sha256 c6c4138ee5c86423 :: prose twin of the run farther from 1
\newcommand{\MlstFtQgHerwigRthirtyCoarseMassEThree}{25.8 and 27.4}  % experiments/FIGS/data/v2/index.json :: fine-tuning qg_herwig N=1000: R30, R16_Q1_MASS :: sha256 c6c4138ee5c86423
\newcommand{\MlstFtQgHerwigRthirtyCoarseMassOverScratchEThree}{0.955 and 1.015}  % experiments/FIGS/data/v2/index.json :: fine-tuning qg_herwig N=1000: R30 of each R16_Q1_MASS run over the from-scratch mean, SD over runs and standard error of that mean in quadrature :: sha256 c6c4138ee5c86423
\newcommand{\MlstFtQgHerwigRthirtyCoarseMassOverScratchEThreePr}{equal within 5\%}  % experiments/FIGS/data/v2/index.json :: fine-tuning qg_herwig N=1000: R30 of each R16_Q1_MASS run over the from-scratch mean, SD over runs and standard error of that mean in quadrature :: sha256 c6c4138ee5c86423 :: prose twin: two significant figures, direction in words
\newcommand{\MlstFtQgHerwigRthirtyCoarseMassOverScratchEThreePrMax}{equal within 5\%}  % experiments/FIGS/data/v2/index.json :: fine-tuning qg_herwig N=1000: R30 of each R16_Q1_MASS run over the from-scratch mean, SD over runs and standard error of that mean in quadrature :: sha256 c6c4138ee5c86423 :: prose twin of the run farther from 1
\newcommand{\MlstFtQgHerwigRthirtyResOverFullEThree}{1.33\,\ensuremath{\pm}\,0.33}  % experiments/FIGS/data/v2/index.json :: fine-tuning qg_herwig N=1000: R30 of L162 over L188, run by run :: sha256 c6c4138ee5c86423
\newcommand{\MlstFtQgHerwigRthirtyResOverFullEThreePr}{equal within 40\%}  % experiments/FIGS/data/v2/index.json :: fine-tuning qg_herwig N=1000: R30 of L162 over L188, run by run :: sha256 c6c4138ee5c86423 :: prose twin: two significant figures, direction in words
\newcommand{\MlstFtQgHerwigRthirtyMidOverFullEThree}{1.36\,\ensuremath{\pm}\,0.31}  % experiments/FIGS/data/v2/index.json :: fine-tuning qg_herwig N=1000: R30 of R42_Q1 over L188, run by run :: sha256 c6c4138ee5c86423
\newcommand{\MlstFtQgHerwigRthirtyMidOverFullEThreePr}{equal within 40\%}  % experiments/FIGS/data/v2/index.json :: fine-tuning qg_herwig N=1000: R30 of R42_Q1 over L188, run by run :: sha256 c6c4138ee5c86423 :: prose twin: two significant figures, direction in words
\newcommand{\MlstFtQgHerwigRthirtyCoarseOverFullEThree}{1.3\,\ensuremath{\pm}\,0.4}  % experiments/FIGS/data/v2/index.json :: fine-tuning qg_herwig N=1000: R30 of R16_Q1 over L188, run by run :: sha256 c6c4138ee5c86423
\newcommand{\MlstFtQgHerwigRthirtyCoarseOverFullEThreePr}{equal within 50\%}  % experiments/FIGS/data/v2/index.json :: fine-tuning qg_herwig N=1000: R30 of R16_Q1 over L188, run by run :: sha256 c6c4138ee5c86423 :: prose twin: two significant figures, direction in words
\newcommand{\MlstFtQgHerwigRthirtyResMassOverResEThree}{1.062 and 1.073}  % experiments/FIGS/data/v2/index.json :: fine-tuning qg_herwig N=1000: R30 of L162_MASS over L162, run by run :: sha256 c6c4138ee5c86423
\newcommand{\MlstFtQgHerwigRthirtyResMassOverResEThreePr}{equal within 8\%}  % experiments/FIGS/data/v2/index.json :: fine-tuning qg_herwig N=1000: R30 of L162_MASS over L162, run by run :: sha256 c6c4138ee5c86423 :: prose twin: two significant figures, direction in words
\newcommand{\MlstFtQgHerwigRthirtyResMassOverResEThreePrMax}{equal within 8\%}  % experiments/FIGS/data/v2/index.json :: fine-tuning qg_herwig N=1000: R30 of L162_MASS over L162, run by run :: sha256 c6c4138ee5c86423 :: prose twin of the run farther from 1
\newcommand{\MlstFtQgHerwigRthirtyCoarseMassOverCoarseEThree}{0.93 and 1.20}  % experiments/FIGS/data/v2/index.json :: fine-tuning qg_herwig N=1000: R30 of R16_Q1_MASS over R16_Q1, run by run :: sha256 c6c4138ee5c86423
\newcommand{\MlstFtQgHerwigRthirtyCoarseMassOverCoarseEThreePr}{equal within 30\%}  % experiments/FIGS/data/v2/index.json :: fine-tuning qg_herwig N=1000: R30 of R16_Q1_MASS over R16_Q1, run by run :: sha256 c6c4138ee5c86423 :: prose twin: two significant figures, direction in words
\newcommand{\MlstFtQgHerwigRthirtyCoarseMassOverCoarseEThreePrMax}{equal within 30\%}  % experiments/FIGS/data/v2/index.json :: fine-tuning qg_herwig N=1000: R30 of R16_Q1_MASS over R16_Q1, run by run :: sha256 c6c4138ee5c86423 :: prose twin of the run farther from 1
\newcommand{\MlstFtQgHerwigRthirtyRandOverFullEThree}{0.99\,\ensuremath{\pm}\,0.20}  % experiments/FIGS/data/v2/index.json :: fine-tuning qg_herwig N=1000: R30 of each random partition over the L188 run of its seed :: sha256 c6c4138ee5c86423
\newcommand{\MlstFtQgHerwigRthirtyRandOverFullEThreePr}{equal within 30\%}  % experiments/FIGS/data/v2/index.json :: fine-tuning qg_herwig N=1000: R30 of each random partition over the L188 run of its seed :: sha256 c6c4138ee5c86423 :: prose twin: two significant figures, direction in words
\newcommand{\MlstFtQgHerwigRthirtyRandOverCoarseEThree}{0.95\,\ensuremath{\pm}\,0.13}  % experiments/FIGS/data/v2/index.json :: fine-tuning qg_herwig N=1000: R30 of each random partition over the R16_Q1 run of its seed :: sha256 c6c4138ee5c86423
\newcommand{\MlstFtQgHerwigRthirtyRandOverCoarseEThreePr}{equal within 20\%}  % experiments/FIGS/data/v2/index.json :: fine-tuning qg_herwig N=1000: R30 of each random partition over the R16_Q1 run of its seed :: sha256 c6c4138ee5c86423 :: prose twin: two significant figures, direction in words
\newcommand{\MlstFtQgHerwigRthirtyScratchEFour}{31\,\ensuremath{\pm}\,7}  % experiments/FIGS/data/v2/index.json :: fine-tuning qg_herwig N=10000: R30 from scratch, every seed :: sha256 c6c4138ee5c86423
\newcommand{\MlstFtQgHerwigRthirtyFullEFour}{27.9\,\ensuremath{\pm}\,2.8}  % experiments/FIGS/data/v2/index.json :: fine-tuning qg_herwig N=10000: R30, L188 :: sha256 c6c4138ee5c86423
\newcommand{\MlstFtQgHerwigRthirtyFullOverScratchEFour}{0.89\,\ensuremath{\pm}\,0.14}  % experiments/FIGS/data/v2/index.json :: fine-tuning qg_herwig N=10000: R30 of each L188 run over the from-scratch mean, SD over runs and standard error of that mean in quadrature :: sha256 c6c4138ee5c86423
\newcommand{\MlstFtQgHerwigRthirtyFullOverScratchEFourPr}{equal within 20\%}  % experiments/FIGS/data/v2/index.json :: fine-tuning qg_herwig N=10000: R30 of each L188 run over the from-scratch mean, SD over runs and standard error of that mean in quadrature :: sha256 c6c4138ee5c86423 :: prose twin: two significant figures, direction in words
\newcommand{\MlstFtQgHerwigRthirtyResEFour}{27.4\,\ensuremath{\pm}\,0.5}  % experiments/FIGS/data/v2/index.json :: fine-tuning qg_herwig N=10000: R30, L162 :: sha256 c6c4138ee5c86423
\newcommand{\MlstFtQgHerwigRthirtyResOverScratchEFour}{0.87\,\ensuremath{\pm}\,0.11}  % experiments/FIGS/data/v2/index.json :: fine-tuning qg_herwig N=10000: R30 of each L162 run over the from-scratch mean, SD over runs and standard error of that mean in quadrature :: sha256 c6c4138ee5c86423
\newcommand{\MlstFtQgHerwigRthirtyResOverScratchEFourPr}{13\% less rejection}  % experiments/FIGS/data/v2/index.json :: fine-tuning qg_herwig N=10000: R30 of each L162 run over the from-scratch mean, SD over runs and standard error of that mean in quadrature :: sha256 c6c4138ee5c86423 :: prose twin: two significant figures, direction in words
\newcommand{\MlstFtQgHerwigRthirtyMidEFour}{28.4\,\ensuremath{\pm}\,2.1}  % experiments/FIGS/data/v2/index.json :: fine-tuning qg_herwig N=10000: R30, R42_Q1 :: sha256 c6c4138ee5c86423
\newcommand{\MlstFtQgHerwigRthirtyMidOverScratchEFour}{0.91\,\ensuremath{\pm}\,0.13}  % experiments/FIGS/data/v2/index.json :: fine-tuning qg_herwig N=10000: R30 of each R42_Q1 run over the from-scratch mean, SD over runs and standard error of that mean in quadrature :: sha256 c6c4138ee5c86423
\newcommand{\MlstFtQgHerwigRthirtyMidOverScratchEFourPr}{equal within 10\%}  % experiments/FIGS/data/v2/index.json :: fine-tuning qg_herwig N=10000: R30 of each R42_Q1 run over the from-scratch mean, SD over runs and standard error of that mean in quadrature :: sha256 c6c4138ee5c86423 :: prose twin: two significant figures, direction in words
\newcommand{\MlstFtQgHerwigRthirtyCoarseEFour}{32.7\,\ensuremath{\pm}\,1.8}  % experiments/FIGS/data/v2/index.json :: fine-tuning qg_herwig N=10000: R30, R16_Q1 :: sha256 c6c4138ee5c86423
\newcommand{\MlstFtQgHerwigRthirtyCoarseOverScratchEFour}{1.05\,\ensuremath{\pm}\,0.14}  % experiments/FIGS/data/v2/index.json :: fine-tuning qg_herwig N=10000: R30 of each R16_Q1 run over the from-scratch mean, SD over runs and standard error of that mean in quadrature :: sha256 c6c4138ee5c86423
\newcommand{\MlstFtQgHerwigRthirtyCoarseOverScratchEFourPr}{equal within 6\%}  % experiments/FIGS/data/v2/index.json :: fine-tuning qg_herwig N=10000: R30 of each R16_Q1 run over the from-scratch mean, SD over runs and standard error of that mean in quadrature :: sha256 c6c4138ee5c86423 :: prose twin: two significant figures, direction in words
\newcommand{\MlstFtQgHerwigRthirtyResMassEFour}{34.3 and 28.5}  % experiments/FIGS/data/v2/index.json :: fine-tuning qg_herwig N=10000: R30, L162_MASS :: sha256 c6c4138ee5c86423
\newcommand{\MlstFtQgHerwigRthirtyResMassOverScratchEFour}{1.09 and 0.91}  % experiments/FIGS/data/v2/index.json :: fine-tuning qg_herwig N=10000: R30 of each L162_MASS run over the from-scratch mean, SD over runs and standard error of that mean in quadrature :: sha256 c6c4138ee5c86423
\newcommand{\MlstFtQgHerwigRthirtyResMassOverScratchEFourPr}{equal within 10\%}  % experiments/FIGS/data/v2/index.json :: fine-tuning qg_herwig N=10000: R30 of each L162_MASS run over the from-scratch mean, SD over runs and standard error of that mean in quadrature :: sha256 c6c4138ee5c86423 :: prose twin: two significant figures, direction in words
\newcommand{\MlstFtQgHerwigRthirtyResMassOverScratchEFourPrMax}{equal within 10\%}  % experiments/FIGS/data/v2/index.json :: fine-tuning qg_herwig N=10000: R30 of each L162_MASS run over the from-scratch mean, SD over runs and standard error of that mean in quadrature :: sha256 c6c4138ee5c86423 :: prose twin of the run farther from 1
\newcommand{\MlstFtQgHerwigRthirtyCoarseMassEFour}{34.4 and 31.2}  % experiments/FIGS/data/v2/index.json :: fine-tuning qg_herwig N=10000: R30, R16_Q1_MASS :: sha256 c6c4138ee5c86423
\newcommand{\MlstFtQgHerwigRthirtyCoarseMassOverScratchEFour}{1.10 and 1.00}  % experiments/FIGS/data/v2/index.json :: fine-tuning qg_herwig N=10000: R30 of each R16_Q1_MASS run over the from-scratch mean, SD over runs and standard error of that mean in quadrature :: sha256 c6c4138ee5c86423
\newcommand{\MlstFtQgHerwigRthirtyCoarseMassOverScratchEFourPr}{equal within 20\%}  % experiments/FIGS/data/v2/index.json :: fine-tuning qg_herwig N=10000: R30 of each R16_Q1_MASS run over the from-scratch mean, SD over runs and standard error of that mean in quadrature :: sha256 c6c4138ee5c86423 :: prose twin: two significant figures, direction in words
\newcommand{\MlstFtQgHerwigRthirtyCoarseMassOverScratchEFourPrMax}{equal within 20\%}  % experiments/FIGS/data/v2/index.json :: fine-tuning qg_herwig N=10000: R30 of each R16_Q1_MASS run over the from-scratch mean, SD over runs and standard error of that mean in quadrature :: sha256 c6c4138ee5c86423 :: prose twin of the run farther from 1
\newcommand{\MlstFtQgHerwigRthirtyResOverFullEFour}{0.99\,\ensuremath{\pm}\,0.12}  % experiments/FIGS/data/v2/index.json :: fine-tuning qg_herwig N=10000: R30 of L162 over L188, run by run :: sha256 c6c4138ee5c86423
\newcommand{\MlstFtQgHerwigRthirtyResOverFullEFourPr}{equal within 20\%}  % experiments/FIGS/data/v2/index.json :: fine-tuning qg_herwig N=10000: R30 of L162 over L188, run by run :: sha256 c6c4138ee5c86423 :: prose twin: two significant figures, direction in words
\newcommand{\MlstFtQgHerwigRthirtyMidOverFullEFour}{1.03\,\ensuremath{\pm}\,0.18}  % experiments/FIGS/data/v2/index.json :: fine-tuning qg_herwig N=10000: R30 of R42_Q1 over L188, run by run :: sha256 c6c4138ee5c86423
\newcommand{\MlstFtQgHerwigRthirtyMidOverFullEFourPr}{equal within 20\%}  % experiments/FIGS/data/v2/index.json :: fine-tuning qg_herwig N=10000: R30 of R42_Q1 over L188, run by run :: sha256 c6c4138ee5c86423 :: prose twin: two significant figures, direction in words
\newcommand{\MlstFtQgHerwigRthirtyCoarseOverFullEFour}{1.18\,\ensuremath{\pm}\,0.08}  % experiments/FIGS/data/v2/index.json :: fine-tuning qg_herwig N=10000: R30 of R16_Q1 over L188, run by run :: sha256 c6c4138ee5c86423
\newcommand{\MlstFtQgHerwigRthirtyCoarseOverFullEFourPr}{18\% more rejection}  % experiments/FIGS/data/v2/index.json :: fine-tuning qg_herwig N=10000: R30 of R16_Q1 over L188, run by run :: sha256 c6c4138ee5c86423 :: prose twin: two significant figures, direction in words
\newcommand{\MlstFtQgHerwigRthirtyResMassOverResEFour}{1.23 and 1.03}  % experiments/FIGS/data/v2/index.json :: fine-tuning qg_herwig N=10000: R30 of L162_MASS over L162, run by run :: sha256 c6c4138ee5c86423
\newcommand{\MlstFtQgHerwigRthirtyResMassOverResEFourPr}{equal within 30\%}  % experiments/FIGS/data/v2/index.json :: fine-tuning qg_herwig N=10000: R30 of L162_MASS over L162, run by run :: sha256 c6c4138ee5c86423 :: prose twin: two significant figures, direction in words
\newcommand{\MlstFtQgHerwigRthirtyResMassOverResEFourPrMax}{equal within 30\%}  % experiments/FIGS/data/v2/index.json :: fine-tuning qg_herwig N=10000: R30 of L162_MASS over L162, run by run :: sha256 c6c4138ee5c86423 :: prose twin of the run farther from 1
\newcommand{\MlstFtQgHerwigRthirtyCoarseMassOverCoarseEFour}{1.12 and 0.91}  % experiments/FIGS/data/v2/index.json :: fine-tuning qg_herwig N=10000: R30 of R16_Q1_MASS over R16_Q1, run by run :: sha256 c6c4138ee5c86423
\newcommand{\MlstFtQgHerwigRthirtyCoarseMassOverCoarseEFourPr}{equal within 20\%}  % experiments/FIGS/data/v2/index.json :: fine-tuning qg_herwig N=10000: R30 of R16_Q1_MASS over R16_Q1, run by run :: sha256 c6c4138ee5c86423 :: prose twin: two significant figures, direction in words
\newcommand{\MlstFtQgHerwigRthirtyCoarseMassOverCoarseEFourPrMax}{equal within 20\%}  % experiments/FIGS/data/v2/index.json :: fine-tuning qg_herwig N=10000: R30 of R16_Q1_MASS over R16_Q1, run by run :: sha256 c6c4138ee5c86423 :: prose twin of the run farther from 1
\newcommand{\MlstFtQgHerwigRthirtyRandOverFullEFour}{1.13\,\ensuremath{\pm}\,0.18}  % experiments/FIGS/data/v2/index.json :: fine-tuning qg_herwig N=10000: R30 of each random partition over the L188 run of its seed :: sha256 c6c4138ee5c86423
\newcommand{\MlstFtQgHerwigRthirtyRandOverFullEFourPr}{equal within 20\%}  % experiments/FIGS/data/v2/index.json :: fine-tuning qg_herwig N=10000: R30 of each random partition over the L188 run of its seed :: sha256 c6c4138ee5c86423 :: prose twin: two significant figures, direction in words
\newcommand{\MlstFtQgHerwigRthirtyRandOverCoarseEFour}{0.93\,\ensuremath{\pm}\,0.14}  % experiments/FIGS/data/v2/index.json :: fine-tuning qg_herwig N=10000: R30 of each random partition over the R16_Q1 run of its seed :: sha256 c6c4138ee5c86423
\newcommand{\MlstFtQgHerwigRthirtyRandOverCoarseEFourPr}{equal within 20\%}  % experiments/FIGS/data/v2/index.json :: fine-tuning qg_herwig N=10000: R30 of each random partition over the R16_Q1 run of its seed :: sha256 c6c4138ee5c86423 :: prose twin: two significant figures, direction in words
\newcommand{\MlstFtQgHerwigRthirtyScratchEFive}{36.7\,\ensuremath{\pm}\,0.9}  % experiments/FIGS/data/v2/index.json :: fine-tuning qg_herwig N=100000: R30 from scratch, every seed :: sha256 c6c4138ee5c86423
\newcommand{\MlstFtQgHerwigRthirtyFullEFive}{34.3\,\ensuremath{\pm}\,1.0}  % experiments/FIGS/data/v2/index.json :: fine-tuning qg_herwig N=100000: R30, L188 :: sha256 c6c4138ee5c86423
\newcommand{\MlstFtQgHerwigRthirtyFullOverScratchEFive}{0.935\,\ensuremath{\pm}\,0.031}  % experiments/FIGS/data/v2/index.json :: fine-tuning qg_herwig N=100000: R30 of each L188 run over the from-scratch mean, SD over runs and standard error of that mean in quadrature :: sha256 c6c4138ee5c86423
\newcommand{\MlstFtQgHerwigRthirtyFullOverScratchEFivePr}{6.5\% less rejection}  % experiments/FIGS/data/v2/index.json :: fine-tuning qg_herwig N=100000: R30 of each L188 run over the from-scratch mean, SD over runs and standard error of that mean in quadrature :: sha256 c6c4138ee5c86423 :: prose twin: two significant figures, direction in words
\newcommand{\MlstFtQgHerwigRthirtyResEFive}{35.5\,\ensuremath{\pm}\,0.5}  % experiments/FIGS/data/v2/index.json :: fine-tuning qg_herwig N=100000: R30, L162 :: sha256 c6c4138ee5c86423
\newcommand{\MlstFtQgHerwigRthirtyResOverScratchEFive}{0.966\,\ensuremath{\pm}\,0.019}  % experiments/FIGS/data/v2/index.json :: fine-tuning qg_herwig N=100000: R30 of each L162 run over the from-scratch mean, SD over runs and standard error of that mean in quadrature :: sha256 c6c4138ee5c86423
\newcommand{\MlstFtQgHerwigRthirtyResOverScratchEFivePr}{3.4\% less rejection}  % experiments/FIGS/data/v2/index.json :: fine-tuning qg_herwig N=100000: R30 of each L162 run over the from-scratch mean, SD over runs and standard error of that mean in quadrature :: sha256 c6c4138ee5c86423 :: prose twin: two significant figures, direction in words
\newcommand{\MlstFtQgHerwigRthirtyMidEFive}{35.5\,\ensuremath{\pm}\,2.1}  % experiments/FIGS/data/v2/index.json :: fine-tuning qg_herwig N=100000: R30, R42_Q1 :: sha256 c6c4138ee5c86423
\newcommand{\MlstFtQgHerwigRthirtyMidOverScratchEFive}{0.97\,\ensuremath{\pm}\,0.06}  % experiments/FIGS/data/v2/index.json :: fine-tuning qg_herwig N=100000: R30 of each R42_Q1 run over the from-scratch mean, SD over runs and standard error of that mean in quadrature :: sha256 c6c4138ee5c86423
\newcommand{\MlstFtQgHerwigRthirtyMidOverScratchEFivePr}{equal within 6\%}  % experiments/FIGS/data/v2/index.json :: fine-tuning qg_herwig N=100000: R30 of each R42_Q1 run over the from-scratch mean, SD over runs and standard error of that mean in quadrature :: sha256 c6c4138ee5c86423 :: prose twin: two significant figures, direction in words
\newcommand{\MlstFtQgHerwigRthirtyCoarseEFive}{34.9\,\ensuremath{\pm}\,1.1}  % experiments/FIGS/data/v2/index.json :: fine-tuning qg_herwig N=100000: R30, R16_Q1 :: sha256 c6c4138ee5c86423
\newcommand{\MlstFtQgHerwigRthirtyCoarseOverScratchEFive}{0.950\,\ensuremath{\pm}\,0.033}  % experiments/FIGS/data/v2/index.json :: fine-tuning qg_herwig N=100000: R30 of each R16_Q1 run over the from-scratch mean, SD over runs and standard error of that mean in quadrature :: sha256 c6c4138ee5c86423
\newcommand{\MlstFtQgHerwigRthirtyCoarseOverScratchEFivePr}{equal within 6\%}  % experiments/FIGS/data/v2/index.json :: fine-tuning qg_herwig N=100000: R30 of each R16_Q1 run over the from-scratch mean, SD over runs and standard error of that mean in quadrature :: sha256 c6c4138ee5c86423 :: prose twin: two significant figures, direction in words
\newcommand{\MlstFtQgHerwigRthirtyResMassEFive}{34.0 and 36.5}  % experiments/FIGS/data/v2/index.json :: fine-tuning qg_herwig N=100000: R30, L162_MASS :: sha256 c6c4138ee5c86423
\newcommand{\MlstFtQgHerwigRthirtyResMassOverScratchEFive}{0.925 and 0.994}  % experiments/FIGS/data/v2/index.json :: fine-tuning qg_herwig N=100000: R30 of each L162_MASS run over the from-scratch mean, SD over runs and standard error of that mean in quadrature :: sha256 c6c4138ee5c86423
\newcommand{\MlstFtQgHerwigRthirtyResMassOverScratchEFivePr}{at least 0.63\% less rejection}  % experiments/FIGS/data/v2/index.json :: fine-tuning qg_herwig N=100000: R30 of each L162_MASS run over the from-scratch mean, SD over runs and standard error of that mean in quadrature :: sha256 c6c4138ee5c86423 :: prose twin: two significant figures, direction in words
\newcommand{\MlstFtQgHerwigRthirtyResMassOverScratchEFivePrMax}{up to 7.5\% less rejection}  % experiments/FIGS/data/v2/index.json :: fine-tuning qg_herwig N=100000: R30 of each L162_MASS run over the from-scratch mean, SD over runs and standard error of that mean in quadrature :: sha256 c6c4138ee5c86423 :: prose twin of the run farther from 1
\newcommand{\MlstFtQgHerwigRthirtyCoarseMassEFive}{35.9 and 39.0}  % experiments/FIGS/data/v2/index.json :: fine-tuning qg_herwig N=100000: R30, R16_Q1_MASS :: sha256 c6c4138ee5c86423
\newcommand{\MlstFtQgHerwigRthirtyCoarseMassOverScratchEFive}{0.979 and 1.064}  % experiments/FIGS/data/v2/index.json :: fine-tuning qg_herwig N=100000: R30 of each R16_Q1_MASS run over the from-scratch mean, SD over runs and standard error of that mean in quadrature :: sha256 c6c4138ee5c86423
\newcommand{\MlstFtQgHerwigRthirtyCoarseMassOverScratchEFivePr}{equal within 7\%}  % experiments/FIGS/data/v2/index.json :: fine-tuning qg_herwig N=100000: R30 of each R16_Q1_MASS run over the from-scratch mean, SD over runs and standard error of that mean in quadrature :: sha256 c6c4138ee5c86423 :: prose twin: two significant figures, direction in words
\newcommand{\MlstFtQgHerwigRthirtyCoarseMassOverScratchEFivePrMax}{equal within 7\%}  % experiments/FIGS/data/v2/index.json :: fine-tuning qg_herwig N=100000: R30 of each R16_Q1_MASS run over the from-scratch mean, SD over runs and standard error of that mean in quadrature :: sha256 c6c4138ee5c86423 :: prose twin of the run farther from 1
\newcommand{\MlstFtQgHerwigRthirtyResOverFullEFive}{1.03\,\ensuremath{\pm}\,0.04}  % experiments/FIGS/data/v2/index.json :: fine-tuning qg_herwig N=100000: R30 of L162 over L188, run by run :: sha256 c6c4138ee5c86423
\newcommand{\MlstFtQgHerwigRthirtyResOverFullEFivePr}{equal within 5\%}  % experiments/FIGS/data/v2/index.json :: fine-tuning qg_herwig N=100000: R30 of L162 over L188, run by run :: sha256 c6c4138ee5c86423 :: prose twin: two significant figures, direction in words
\newcommand{\MlstFtQgHerwigRthirtyMidOverFullEFive}{1.03\,\ensuremath{\pm}\,0.05}  % experiments/FIGS/data/v2/index.json :: fine-tuning qg_herwig N=100000: R30 of R42_Q1 over L188, run by run :: sha256 c6c4138ee5c86423
\newcommand{\MlstFtQgHerwigRthirtyMidOverFullEFivePr}{equal within 6\%}  % experiments/FIGS/data/v2/index.json :: fine-tuning qg_herwig N=100000: R30 of R42_Q1 over L188, run by run :: sha256 c6c4138ee5c86423 :: prose twin: two significant figures, direction in words
\newcommand{\MlstFtQgHerwigRthirtyCoarseOverFullEFive}{1.02\,\ensuremath{\pm}\,0.06}  % experiments/FIGS/data/v2/index.json :: fine-tuning qg_herwig N=100000: R30 of R16_Q1 over L188, run by run :: sha256 c6c4138ee5c86423
\newcommand{\MlstFtQgHerwigRthirtyCoarseOverFullEFivePr}{equal within 7\%}  % experiments/FIGS/data/v2/index.json :: fine-tuning qg_herwig N=100000: R30 of R16_Q1 over L188, run by run :: sha256 c6c4138ee5c86423 :: prose twin: two significant figures, direction in words
\newcommand{\MlstFtQgHerwigRthirtyResMassOverResEFive}{0.971 and 1.028}  % experiments/FIGS/data/v2/index.json :: fine-tuning qg_herwig N=100000: R30 of L162_MASS over L162, run by run :: sha256 c6c4138ee5c86423
\newcommand{\MlstFtQgHerwigRthirtyResMassOverResEFivePr}{equal within 3\%}  % experiments/FIGS/data/v2/index.json :: fine-tuning qg_herwig N=100000: R30 of L162_MASS over L162, run by run :: sha256 c6c4138ee5c86423 :: prose twin: two significant figures, direction in words
\newcommand{\MlstFtQgHerwigRthirtyResMassOverResEFivePrMax}{equal within 3\%}  % experiments/FIGS/data/v2/index.json :: fine-tuning qg_herwig N=100000: R30 of L162_MASS over L162, run by run :: sha256 c6c4138ee5c86423 :: prose twin of the run farther from 1
\newcommand{\MlstFtQgHerwigRthirtyCoarseMassOverCoarseEFive}{1.064 and 1.121}  % experiments/FIGS/data/v2/index.json :: fine-tuning qg_herwig N=100000: R30 of R16_Q1_MASS over R16_Q1, run by run :: sha256 c6c4138ee5c86423
\newcommand{\MlstFtQgHerwigRthirtyCoarseMassOverCoarseEFivePr}{equal within 20\%}  % experiments/FIGS/data/v2/index.json :: fine-tuning qg_herwig N=100000: R30 of R16_Q1_MASS over R16_Q1, run by run :: sha256 c6c4138ee5c86423 :: prose twin: two significant figures, direction in words
\newcommand{\MlstFtQgHerwigRthirtyCoarseMassOverCoarseEFivePrMax}{equal within 20\%}  % experiments/FIGS/data/v2/index.json :: fine-tuning qg_herwig N=100000: R30 of R16_Q1_MASS over R16_Q1, run by run :: sha256 c6c4138ee5c86423 :: prose twin of the run farther from 1
\newcommand{\MlstFtQgHerwigRthirtyRandOverFullEFive}{1.01\,\ensuremath{\pm}\,0.05}  % experiments/FIGS/data/v2/index.json :: fine-tuning qg_herwig N=100000: R30 of each random partition over the L188 run of its seed :: sha256 c6c4138ee5c86423
\newcommand{\MlstFtQgHerwigRthirtyRandOverFullEFivePr}{equal within 5\%}  % experiments/FIGS/data/v2/index.json :: fine-tuning qg_herwig N=100000: R30 of each random partition over the L188 run of its seed :: sha256 c6c4138ee5c86423 :: prose twin: two significant figures, direction in words
\newcommand{\MlstFtQgHerwigRthirtyRandOverCoarseEFive}{1.009\,\ensuremath{\pm}\,0.024}  % experiments/FIGS/data/v2/index.json :: fine-tuning qg_herwig N=100000: R30 of each random partition over the R16_Q1 run of its seed :: sha256 c6c4138ee5c86423
\newcommand{\MlstFtQgHerwigRthirtyRandOverCoarseEFivePr}{equal within 3\%}  % experiments/FIGS/data/v2/index.json :: fine-tuning qg_herwig N=100000: R30 of each random partition over the R16_Q1 run of its seed :: sha256 c6c4138ee5c86423 :: prose twin: two significant figures, direction in words
\newcommand{\MlstFtQgHerwigRthirtyScratch}{35.5\,\ensuremath{\pm}\,1.2}  % experiments/FIGS/data/v2/index.json :: fine-tuning qg_herwig N=1600000: R30 from scratch, every seed :: sha256 c6c4138ee5c86423
\newcommand{\MlstFtQgHerwigRthirtyFull}{35.2\,\ensuremath{\pm}\,2.1}  % experiments/FIGS/data/v2/index.json :: fine-tuning qg_herwig N=1600000: R30, L188 :: sha256 c6c4138ee5c86423
\newcommand{\MlstFtQgHerwigRthirtyFullOverScratch}{0.99\,\ensuremath{\pm}\,0.06}  % experiments/FIGS/data/v2/index.json :: fine-tuning qg_herwig N=1600000: R30 of each L188 run over the from-scratch mean, SD over runs and standard error of that mean in quadrature :: sha256 c6c4138ee5c86423
\newcommand{\MlstFtQgHerwigRthirtyFullOverScratchPr}{equal within 6\%}  % experiments/FIGS/data/v2/index.json :: fine-tuning qg_herwig N=1600000: R30 of each L188 run over the from-scratch mean, SD over runs and standard error of that mean in quadrature :: sha256 c6c4138ee5c86423 :: prose twin: two significant figures, direction in words
\newcommand{\MlstFtQgHerwigRthirtyRes}{34.6\,\ensuremath{\pm}\,2.0}  % experiments/FIGS/data/v2/index.json :: fine-tuning qg_herwig N=1600000: R30, L162 :: sha256 c6c4138ee5c86423
\newcommand{\MlstFtQgHerwigRthirtyResOverScratch}{0.97\,\ensuremath{\pm}\,0.06}  % experiments/FIGS/data/v2/index.json :: fine-tuning qg_herwig N=1600000: R30 of each L162 run over the from-scratch mean, SD over runs and standard error of that mean in quadrature :: sha256 c6c4138ee5c86423
\newcommand{\MlstFtQgHerwigRthirtyResOverScratchPr}{equal within 6\%}  % experiments/FIGS/data/v2/index.json :: fine-tuning qg_herwig N=1600000: R30 of each L162 run over the from-scratch mean, SD over runs and standard error of that mean in quadrature :: sha256 c6c4138ee5c86423 :: prose twin: two significant figures, direction in words
\newcommand{\MlstFtQgHerwigRthirtyMid}{36.1\,\ensuremath{\pm}\,2.7}  % experiments/FIGS/data/v2/index.json :: fine-tuning qg_herwig N=1600000: R30, R42_Q1 :: sha256 c6c4138ee5c86423
\newcommand{\MlstFtQgHerwigRthirtyMidOverScratch}{1.02\,\ensuremath{\pm}\,0.08}  % experiments/FIGS/data/v2/index.json :: fine-tuning qg_herwig N=1600000: R30 of each R42_Q1 run over the from-scratch mean, SD over runs and standard error of that mean in quadrature :: sha256 c6c4138ee5c86423
\newcommand{\MlstFtQgHerwigRthirtyMidOverScratchPr}{equal within 8\%}  % experiments/FIGS/data/v2/index.json :: fine-tuning qg_herwig N=1600000: R30 of each R42_Q1 run over the from-scratch mean, SD over runs and standard error of that mean in quadrature :: sha256 c6c4138ee5c86423 :: prose twin: two significant figures, direction in words
\newcommand{\MlstFtQgHerwigRthirtyCoarse}{36.9\,\ensuremath{\pm}\,0.9}  % experiments/FIGS/data/v2/index.json :: fine-tuning qg_herwig N=1600000: R30, R16_Q1 :: sha256 c6c4138ee5c86423
\newcommand{\MlstFtQgHerwigRthirtyCoarseOverScratch}{1.040\,\ensuremath{\pm}\,0.033}  % experiments/FIGS/data/v2/index.json :: fine-tuning qg_herwig N=1600000: R30 of each R16_Q1 run over the from-scratch mean, SD over runs and standard error of that mean in quadrature :: sha256 c6c4138ee5c86423
\newcommand{\MlstFtQgHerwigRthirtyCoarseOverScratchPr}{equal within 4\%}  % experiments/FIGS/data/v2/index.json :: fine-tuning qg_herwig N=1600000: R30 of each R16_Q1 run over the from-scratch mean, SD over runs and standard error of that mean in quadrature :: sha256 c6c4138ee5c86423 :: prose twin: two significant figures, direction in words
\newcommand{\MlstFtQgHerwigRthirtyResMass}{37.8 and 34.4}  % experiments/FIGS/data/v2/index.json :: fine-tuning qg_herwig N=1600000: R30, L162_MASS :: sha256 c6c4138ee5c86423
\newcommand{\MlstFtQgHerwigRthirtyResMassOverScratch}{1.065 and 0.970}  % experiments/FIGS/data/v2/index.json :: fine-tuning qg_herwig N=1600000: R30 of each L162_MASS run over the from-scratch mean, SD over runs and standard error of that mean in quadrature :: sha256 c6c4138ee5c86423
\newcommand{\MlstFtQgHerwigRthirtyResMassOverScratchPr}{equal within 7\%}  % experiments/FIGS/data/v2/index.json :: fine-tuning qg_herwig N=1600000: R30 of each L162_MASS run over the from-scratch mean, SD over runs and standard error of that mean in quadrature :: sha256 c6c4138ee5c86423 :: prose twin: two significant figures, direction in words
\newcommand{\MlstFtQgHerwigRthirtyResMassOverScratchPrMax}{equal within 7\%}  % experiments/FIGS/data/v2/index.json :: fine-tuning qg_herwig N=1600000: R30 of each L162_MASS run over the from-scratch mean, SD over runs and standard error of that mean in quadrature :: sha256 c6c4138ee5c86423 :: prose twin of the run farther from 1
\newcommand{\MlstFtQgHerwigRthirtyCoarseMass}{37.7 and 32.9}  % experiments/FIGS/data/v2/index.json :: fine-tuning qg_herwig N=1600000: R30, R16_Q1_MASS :: sha256 c6c4138ee5c86423
\newcommand{\MlstFtQgHerwigRthirtyCoarseMassOverScratch}{1.06 and 0.93}  % experiments/FIGS/data/v2/index.json :: fine-tuning qg_herwig N=1600000: R30 of each R16_Q1_MASS run over the from-scratch mean, SD over runs and standard error of that mean in quadrature :: sha256 c6c4138ee5c86423
\newcommand{\MlstFtQgHerwigRthirtyCoarseMassOverScratchPr}{equal within 8\%}  % experiments/FIGS/data/v2/index.json :: fine-tuning qg_herwig N=1600000: R30 of each R16_Q1_MASS run over the from-scratch mean, SD over runs and standard error of that mean in quadrature :: sha256 c6c4138ee5c86423 :: prose twin: two significant figures, direction in words
\newcommand{\MlstFtQgHerwigRthirtyCoarseMassOverScratchPrMax}{equal within 8\%}  % experiments/FIGS/data/v2/index.json :: fine-tuning qg_herwig N=1600000: R30 of each R16_Q1_MASS run over the from-scratch mean, SD over runs and standard error of that mean in quadrature :: sha256 c6c4138ee5c86423 :: prose twin of the run farther from 1
\newcommand{\MlstFtQgHerwigRthirtyResOverFull}{0.9832\,\ensuremath{\pm}\,0.0022}  % experiments/FIGS/data/v2/index.json :: fine-tuning qg_herwig N=1600000: R30 of L162 over L188, run by run :: sha256 c6c4138ee5c86423
\newcommand{\MlstFtQgHerwigRthirtyResOverFullPr}{1.7\% less rejection}  % experiments/FIGS/data/v2/index.json :: fine-tuning qg_herwig N=1600000: R30 of L162 over L188, run by run :: sha256 c6c4138ee5c86423 :: prose twin: two significant figures, direction in words
\newcommand{\MlstFtQgHerwigRthirtyMidOverFull}{1.03\,\ensuremath{\pm}\,0.05}  % experiments/FIGS/data/v2/index.json :: fine-tuning qg_herwig N=1600000: R30 of R42_Q1 over L188, run by run :: sha256 c6c4138ee5c86423
\newcommand{\MlstFtQgHerwigRthirtyMidOverFullPr}{equal within 6\%}  % experiments/FIGS/data/v2/index.json :: fine-tuning qg_herwig N=1600000: R30 of R42_Q1 over L188, run by run :: sha256 c6c4138ee5c86423 :: prose twin: two significant figures, direction in words
\newcommand{\MlstFtQgHerwigRthirtyCoarseOverFull}{1.05\,\ensuremath{\pm}\,0.06}  % experiments/FIGS/data/v2/index.json :: fine-tuning qg_herwig N=1600000: R30 of R16_Q1 over L188, run by run :: sha256 c6c4138ee5c86423
\newcommand{\MlstFtQgHerwigRthirtyCoarseOverFullPr}{equal within 6\%}  % experiments/FIGS/data/v2/index.json :: fine-tuning qg_herwig N=1600000: R30 of R16_Q1 over L188, run by run :: sha256 c6c4138ee5c86423 :: prose twin: two significant figures, direction in words
\newcommand{\MlstFtQgHerwigRthirtyResMassOverRes}{1.14 and 0.93}  % experiments/FIGS/data/v2/index.json :: fine-tuning qg_herwig N=1600000: R30 of L162_MASS over L162, run by run :: sha256 c6c4138ee5c86423
\newcommand{\MlstFtQgHerwigRthirtyResMassOverResPr}{equal within 20\%}  % experiments/FIGS/data/v2/index.json :: fine-tuning qg_herwig N=1600000: R30 of L162_MASS over L162, run by run :: sha256 c6c4138ee5c86423 :: prose twin: two significant figures, direction in words
\newcommand{\MlstFtQgHerwigRthirtyResMassOverResPrMax}{equal within 20\%}  % experiments/FIGS/data/v2/index.json :: fine-tuning qg_herwig N=1600000: R30 of L162_MASS over L162, run by run :: sha256 c6c4138ee5c86423 :: prose twin of the run farther from 1
\newcommand{\MlstFtQgHerwigRthirtyCoarseMassOverCoarse}{1.01 and 0.88}  % experiments/FIGS/data/v2/index.json :: fine-tuning qg_herwig N=1600000: R30 of R16_Q1_MASS over R16_Q1, run by run :: sha256 c6c4138ee5c86423
\newcommand{\MlstFtQgHerwigRthirtyCoarseMassOverCoarsePr}{equal within 20\%}  % experiments/FIGS/data/v2/index.json :: fine-tuning qg_herwig N=1600000: R30 of R16_Q1_MASS over R16_Q1, run by run :: sha256 c6c4138ee5c86423 :: prose twin: two significant figures, direction in words
\newcommand{\MlstFtQgHerwigRthirtyCoarseMassOverCoarsePrMax}{equal within 20\%}  % experiments/FIGS/data/v2/index.json :: fine-tuning qg_herwig N=1600000: R30 of R16_Q1_MASS over R16_Q1, run by run :: sha256 c6c4138ee5c86423 :: prose twin of the run farther from 1
\newcommand{\MlstFtQgHerwigRthirtyRandOverFull}{0.98\,\ensuremath{\pm}\,0.04}  % experiments/FIGS/data/v2/index.json :: fine-tuning qg_herwig N=1600000: R30 of each random partition over the L188 run of its seed :: sha256 c6c4138ee5c86423
\newcommand{\MlstFtQgHerwigRthirtyRandOverFullPr}{equal within 5\%}  % experiments/FIGS/data/v2/index.json :: fine-tuning qg_herwig N=1600000: R30 of each random partition over the L188 run of its seed :: sha256 c6c4138ee5c86423 :: prose twin: two significant figures, direction in words
\newcommand{\MlstFtQgHerwigRthirtyRandOverCoarse}{0.94\,\ensuremath{\pm}\,0.04}  % experiments/FIGS/data/v2/index.json :: fine-tuning qg_herwig N=1600000: R30 of each random partition over the R16_Q1 run of its seed :: sha256 c6c4138ee5c86423
\newcommand{\MlstFtQgHerwigRthirtyRandOverCoarsePr}{equal within 6\%}  % experiments/FIGS/data/v2/index.json :: fine-tuning qg_herwig N=1600000: R30 of each random partition over the R16_Q1 run of its seed :: sha256 c6c4138ee5c86423 :: prose twin: two significant figures, direction in words
\newcommand{\MlstFtQgRthirtyHerwigOverPythiaFull}{0.328\,\ensuremath{\pm}\,0.015}  % experiments/FIGS/data/v2/index.json :: fine-tuning q/g N=1600000: R30 on Herwig over R30 on Pythia, each run of L188 :: sha256 c6c4138ee5c86423
\newcommand{\MlstFtQgRthirtyHerwigOverPythiaFullPr}{0.33 times the rejection}  % experiments/FIGS/data/v2/index.json :: fine-tuning q/g N=1600000: R30 on Herwig over R30 on Pythia, each run of L188 :: sha256 c6c4138ee5c86423 :: prose twin: two significant figures, direction in words
\newcommand{\MlstFtQgRthirtyHerwigOverPythiaRes}{0.320\,\ensuremath{\pm}\,0.023}  % experiments/FIGS/data/v2/index.json :: fine-tuning q/g N=1600000: R30 on Herwig over R30 on Pythia, each run of L162 :: sha256 c6c4138ee5c86423
\newcommand{\MlstFtQgRthirtyHerwigOverPythiaResPr}{0.32 times the rejection}  % experiments/FIGS/data/v2/index.json :: fine-tuning q/g N=1600000: R30 on Herwig over R30 on Pythia, each run of L162 :: sha256 c6c4138ee5c86423 :: prose twin: two significant figures, direction in words
\newcommand{\MlstFtQgRthirtyHerwigOverPythiaMid}{0.334\,\ensuremath{\pm}\,0.023}  % experiments/FIGS/data/v2/index.json :: fine-tuning q/g N=1600000: R30 on Herwig over R30 on Pythia, each run of R42_Q1 :: sha256 c6c4138ee5c86423
\newcommand{\MlstFtQgRthirtyHerwigOverPythiaMidPr}{0.33 times the rejection}  % experiments/FIGS/data/v2/index.json :: fine-tuning q/g N=1600000: R30 on Herwig over R30 on Pythia, each run of R42_Q1 :: sha256 c6c4138ee5c86423 :: prose twin: two significant figures, direction in words
\newcommand{\MlstFtQgRthirtyHerwigOverPythiaCoarse}{0.346\,\ensuremath{\pm}\,0.011}  % experiments/FIGS/data/v2/index.json :: fine-tuning q/g N=1600000: R30 on Herwig over R30 on Pythia, each run of R16_Q1 :: sha256 c6c4138ee5c86423
\newcommand{\MlstFtQgRthirtyHerwigOverPythiaCoarsePr}{0.35 times the rejection}  % experiments/FIGS/data/v2/index.json :: fine-tuning q/g N=1600000: R30 on Herwig over R30 on Pythia, each run of R16_Q1 :: sha256 c6c4138ee5c86423 :: prose twin: two significant figures, direction in words
\newcommand{\MlstFtQgRthirtyHerwigOverPythiaResMass}{0.353 and 0.323}  % experiments/FIGS/data/v2/index.json :: fine-tuning q/g N=1600000: R30 on Herwig over R30 on Pythia, each run of L162_MASS :: sha256 c6c4138ee5c86423
\newcommand{\MlstFtQgRthirtyHerwigOverPythiaResMassPr}{at most 0.35 times the rejection}  % experiments/FIGS/data/v2/index.json :: fine-tuning q/g N=1600000: R30 on Herwig over R30 on Pythia, each run of L162_MASS :: sha256 c6c4138ee5c86423 :: prose twin: two significant figures, direction in words
\newcommand{\MlstFtQgRthirtyHerwigOverPythiaResMassPrMax}{down to 0.32 times the rejection}  % experiments/FIGS/data/v2/index.json :: fine-tuning q/g N=1600000: R30 on Herwig over R30 on Pythia, each run of L162_MASS :: sha256 c6c4138ee5c86423 :: prose twin of the run farther from 1
\newcommand{\MlstFtQgRthirtyHerwigOverPythiaCoarseMass}{0.359 and 0.311}  % experiments/FIGS/data/v2/index.json :: fine-tuning q/g N=1600000: R30 on Herwig over R30 on Pythia, each run of R16_Q1_MASS :: sha256 c6c4138ee5c86423
\newcommand{\MlstFtQgRthirtyHerwigOverPythiaCoarseMassPr}{at most 0.36 times the rejection}  % experiments/FIGS/data/v2/index.json :: fine-tuning q/g N=1600000: R30 on Herwig over R30 on Pythia, each run of R16_Q1_MASS :: sha256 c6c4138ee5c86423 :: prose twin: two significant figures, direction in words
\newcommand{\MlstFtQgRthirtyHerwigOverPythiaCoarseMassPrMax}{down to 0.31 times the rejection}  % experiments/FIGS/data/v2/index.json :: fine-tuning q/g N=1600000: R30 on Herwig over R30 on Pythia, each run of R16_Q1_MASS :: sha256 c6c4138ee5c86423 :: prose twin of the run farther from 1
\newcommand{\MlstFtQgRthirtyHerwigOverPythiaScratch}{0.341\,\ensuremath{\pm}\,0.009}  % experiments/FIGS/data/v2/index.json :: fine-tuning q/g N=1600000: R30 on Herwig over R30 on Pythia, each from-scratch seed :: sha256 c6c4138ee5c86423
\newcommand{\MlstFtQgRthirtyHerwigOverPythiaScratchPr}{0.34 times the rejection}  % experiments/FIGS/data/v2/index.json :: fine-tuning q/g N=1600000: R30 on Herwig over R30 on Pythia, each from-scratch seed :: sha256 c6c4138ee5c86423 :: prose twin: two significant figures, direction in words
\newcommand{\MlstTopCheckFullLong}{\pending{recipe check}}  % experiments/FIGS/data/v2/index.json :: the recipe check is not read out yet :: sha256 c6c4138ee5c86423
\newcommand{\MlstTopCheckScratchLong}{\pending{recipe check}}  % experiments/FIGS/data/v2/index.json :: the recipe check is not read out yet :: sha256 c6c4138ee5c86423
\newcommand{\MlstTopCheckFullShort}{\pending{recipe check}}  % experiments/FIGS/data/v2/index.json :: the recipe check is not read out yet :: sha256 c6c4138ee5c86423
\newcommand{\MlstTopCheckScratchShort}{\pending{recipe check}}  % experiments/FIGS/data/v2/index.json :: the recipe check is not read out yet :: sha256 c6c4138ee5c86423
\newcommand{\MlstLpJcOneGapFracEFour}{56.3\,\ensuremath{\pm}\,1.1\%}  % experiments/FIGS/data/v2/index.json :: jc1 N=10000: (probe - scratch) / (fine-tuned - scratch) accuracy, 188-class runs, per cent :: sha256 c6c4138ee5c86423
\newcommand{\MlstLpJcOneGapFracEFourPr}{56\%}  % experiments/FIGS/data/v2/index.json :: jc1 N=10000: (probe - scratch) / (fine-tuned - scratch) accuracy, 188-class runs, per cent :: sha256 c6c4138ee5c86423 :: prose twin, two significant figures
\newcommand{\MlstLpJcOneGapFracEFive}{\ensuremath{-}15\,\ensuremath{\pm}\,5\%}  % experiments/FIGS/data/v2/index.json :: jc1 N=100000: (probe - scratch) / (fine-tuned - scratch) accuracy, 188-class runs, per cent :: sha256 c6c4138ee5c86423
\newcommand{\MlstLpJcOneGapFracEFivePr}{\ensuremath{-}15\%}  % experiments/FIGS/data/v2/index.json :: jc1 N=100000: (probe - scratch) / (fine-tuned - scratch) accuracy, 188-class runs, per cent :: sha256 c6c4138ee5c86423 :: prose twin, two significant figures
\newcommand{\MlstLpJcTwoGapFracEFour}{96.2\,\ensuremath{\pm}\,0.6\%}  % experiments/FIGS/data/v2/index.json :: jc2_188 N=10000: (probe - scratch) / (fine-tuned - scratch) accuracy, 188-class runs, per cent :: sha256 c6c4138ee5c86423
\newcommand{\MlstLpJcTwoGapFracEFourPr}{96\%}  % experiments/FIGS/data/v2/index.json :: jc2_188 N=10000: (probe - scratch) / (fine-tuned - scratch) accuracy, 188-class runs, per cent :: sha256 c6c4138ee5c86423 :: prose twin, two significant figures
\newcommand{\MlstLpJcTwoGapFracEFive}{102.2\,\ensuremath{\pm}\,1.0\%}  % experiments/FIGS/data/v2/index.json :: jc2_188 N=100000: (probe - scratch) / (fine-tuned - scratch) accuracy, 188-class runs, per cent :: sha256 c6c4138ee5c86423
\newcommand{\MlstLpJcTwoGapFracEFivePr}{100\%}  % experiments/FIGS/data/v2/index.json :: jc2_188 N=100000: (probe - scratch) / (fine-tuned - scratch) accuracy, 188-class runs, per cent :: sha256 c6c4138ee5c86423 :: prose twin, two significant figures
\newcommand{\XbxMx}{\pending{samples not in}}  % experiments/XBS/xbx_report.py :: pending :: sha256 cb6dffe774ff5b1e
\newcommand{\XbxWindowLo}{\pending{samples not in}}  % experiments/XBS/xbx_report.py :: pending :: sha256 cb6dffe774ff5b1e
\newcommand{\XbxWindowHi}{\pending{samples not in}}  % experiments/XBS/xbx_report.py :: pending :: sha256 cb6dffe774ff5b1e
\newcommand{\XbbNTrainPerClass}{\pending{samples not in}}  % experiments/XBS/xbx_report.py :: pending :: sha256 cb6dffe774ff5b1e
\newcommand{\XbbRepetition}{\pending{samples not in}}  % experiments/XBS/xbx_report.py :: pending :: sha256 cb6dffe774ff5b1e
\newcommand{\XbbNEvalSignal}{\pending{samples not in}}  % experiments/XBS/xbx_report.py :: pending :: sha256 cb6dffe774ff5b1e
\newcommand{\XbbNEvalQcd}{\pending{samples not in}}  % experiments/XBS/xbx_report.py :: pending :: sha256 cb6dffe774ff5b1e
\newcommand{\XbsNTrainPerClass}{\pending{samples not in}}  % experiments/XBS/xbx_report.py :: pending :: sha256 cb6dffe774ff5b1e
\newcommand{\XbsRepetition}{\pending{samples not in}}  % experiments/XBS/xbx_report.py :: pending :: sha256 cb6dffe774ff5b1e
\newcommand{\XbsNEvalSignal}{\pending{samples not in}}  % experiments/XBS/xbx_report.py :: pending :: sha256 cb6dffe774ff5b1e
\newcommand{\XbsNEvalQcd}{\pending{samples not in}}  % experiments/XBS/xbx_report.py :: pending :: sha256 cb6dffe774ff5b1e
\newcommand{\MlstPartreproRthirtyPart}{1{,}711}  % experiments/FIGS/data/v2/partrepro/part.json :: r30 :: sha256 0d44b997af9e6448 :: ParT reproduction on top, part: R30 (one run)
\newcommand{\MlstPartreproRthirtyHead}{1{,}577}  % experiments/FIGS/data/v2/partrepro/head.json :: r30 :: sha256 37b82ba0a0fe9299 :: ParT reproduction on top, head: R30 (one run)
\newcommand{\MlstPartreproRthirtyRescale}{1{,}224}  % experiments/FIGS/data/v2/partrepro/rescale.json :: r30 :: sha256 0a9f1401047d07ff :: ParT reproduction on top, rescale: R30 (one run)
\newcommand{\MlstPartreproRthirtyHeadOverPart}{0.922}  % experiments/FIGS/data/v2/partrepro/head.json :: r30 / part r30 :: sha256 37b82ba0a0fe9299
\newcommand{\MlstPartreproRthirtyHeadOverPartPr}{7.8\% less rejection}  % experiments/FIGS/data/v2/partrepro/head.json :: r30 / part r30 :: sha256 37b82ba0a0fe9299 :: prose twin: two significant figures, direction in words
\newcommand{\MlstPartreproRthirtyRescaleOverPart}{0.715}  % experiments/FIGS/data/v2/partrepro/rescale.json :: r30 / part r30 :: sha256 0a9f1401047d07ff
\newcommand{\MlstPartreproRthirtyRescaleOverPartPr}{28\% less rejection}  % experiments/FIGS/data/v2/partrepro/rescale.json :: r30 / part r30 :: sha256 0a9f1401047d07ff :: prose twin: two significant figures, direction in words
\newcommand{\MlstPartreproSliceDropRatio}{\pending{jet-pT slice test not run}}  % experiments/FIGS/data/v2/index.json :: open item O10: R30 in narrow pT slices of partrepro part and rescale :: sha256 c6c4138ee5c86423
\newcommand{\MlstInputBnZeroSigmaMax}{0.96}  % experiments/FIGS/data/v2/head/input_bn_offsets.json :: offsets :: sha256 d6b821509fccecff :: largest |offset| over the five type flags
\newcommand{\MlstFtJcOneNSignal}{nine}  % experiments/FIGS/data/v2/curves/full_budget_long_leg2.json :: JC1_COLS (count) :: sha256 588e935996d1006a :: JetClass signal classes
\newcommand{\MlstFtJcOneRejHbbResOverFull}{0.92\,\ensuremath{\pm}\,0.05}  % experiments/FIGS/data/v2/curves/full_budget_long_leg2.json :: cells.<run>.N1000000.s1.jc1.per_class.label_Hbb: rejection of L162 over L188, run by run :: sha256 588e935996d1006a
\newcommand{\MlstFtJcOneRejHbbResOverFullPr}{equal within 9\%}  % experiments/FIGS/data/v2/curves/full_budget_long_leg2.json :: cells.<run>.N1000000.s1.jc1.per_class.label_Hbb: rejection of L162 over L188, run by run :: sha256 588e935996d1006a :: prose twin: two significant figures, direction in words
\newcommand{\MlstFtJcOneRejHbbMidOverFull}{0.98\,\ensuremath{\pm}\,0.12}  % experiments/FIGS/data/v2/curves/full_budget_long_leg2.json :: cells.<run>.N1000000.s1.jc1.per_class.label_Hbb: rejection of R42_Q1 over L188, run by run :: sha256 588e935996d1006a
\newcommand{\MlstFtJcOneRejHbbMidOverFullPr}{equal within 20\%}  % experiments/FIGS/data/v2/curves/full_budget_long_leg2.json :: cells.<run>.N1000000.s1.jc1.per_class.label_Hbb: rejection of R42_Q1 over L188, run by run :: sha256 588e935996d1006a :: prose twin: two significant figures, direction in words
\newcommand{\MlstFtJcOneRejHbbCoarseOverFull}{1.01\,\ensuremath{\pm}\,0.07}  % experiments/FIGS/data/v2/curves/full_budget_long_leg2.json :: cells.<run>.N1000000.s1.jc1.per_class.label_Hbb: rejection of R16_Q1 over L188, run by run :: sha256 588e935996d1006a
\newcommand{\MlstFtJcOneRejHbbCoarseOverFullPr}{equal within 7\%}  % experiments/FIGS/data/v2/curves/full_budget_long_leg2.json :: cells.<run>.N1000000.s1.jc1.per_class.label_Hbb: rejection of R16_Q1 over L188, run by run :: sha256 588e935996d1006a :: prose twin: two significant figures, direction in words
\newcommand{\MlstFtJcOneRejHbbResMassOverRes}{1.065 and 1.100}  % experiments/FIGS/data/v2/curves/full_budget_long_leg2.json :: cells.<run>.N1000000.s1.jc1.per_class.label_Hbb: rejection of L162_MASS over L162, run by run :: sha256 588e935996d1006a
\newcommand{\MlstFtJcOneRejHbbResMassOverResPr}{equal within 10\%}  % experiments/FIGS/data/v2/curves/full_budget_long_leg2.json :: cells.<run>.N1000000.s1.jc1.per_class.label_Hbb: rejection of L162_MASS over L162, run by run :: sha256 588e935996d1006a :: prose twin: two significant figures, direction in words
\newcommand{\MlstFtJcOneRejHbbResMassOverResPrMax}{equal within 10\%}  % experiments/FIGS/data/v2/curves/full_budget_long_leg2.json :: cells.<run>.N1000000.s1.jc1.per_class.label_Hbb: rejection of L162_MASS over L162, run by run :: sha256 588e935996d1006a :: prose twin of the run farther from 1
\newcommand{\MlstFtJcOneRejHbbCoarseMassOverCoarse}{0.97 and 0.86}  % experiments/FIGS/data/v2/curves/full_budget_long_leg2.json :: cells.<run>.N1000000.s1.jc1.per_class.label_Hbb: rejection of R16_Q1_MASS over R16_Q1, run by run :: sha256 588e935996d1006a
\newcommand{\MlstFtJcOneRejHbbCoarseMassOverCoarsePr}{equal within 20\%}  % experiments/FIGS/data/v2/curves/full_budget_long_leg2.json :: cells.<run>.N1000000.s1.jc1.per_class.label_Hbb: rejection of R16_Q1_MASS over R16_Q1, run by run :: sha256 588e935996d1006a :: prose twin: two significant figures, direction in words
\newcommand{\MlstFtJcOneRejHbbCoarseMassOverCoarsePrMax}{equal within 20\%}  % experiments/FIGS/data/v2/curves/full_budget_long_leg2.json :: cells.<run>.N1000000.s1.jc1.per_class.label_Hbb: rejection of R16_Q1_MASS over R16_Q1, run by run :: sha256 588e935996d1006a :: prose twin of the run farther from 1
\newcommand{\MlstFtJcOneRejHbbRandOverFull}{0.92\,\ensuremath{\pm}\,0.11}  % experiments/FIGS/data/v2/curves/full_budget_long_leg2.json :: per_class.label_Hbb: rejection of each random partition over the 188-class run of its seed :: sha256 588e935996d1006a
\newcommand{\MlstFtJcOneRejHbbRandOverFullPr}{equal within 20\%}  % experiments/FIGS/data/v2/curves/full_budget_long_leg2.json :: per_class.label_Hbb: rejection of each random partition over the 188-class run of its seed :: sha256 588e935996d1006a :: prose twin: two significant figures, direction in words
\newcommand{\MlstFtJcOneRejHbbRandOverFullRange}{0.84--1.1}  % experiments/FIGS/data/v2/curves/full_budget_long_leg2.json :: per_class.label_Hbb: smallest and largest of the partitions' ratios :: sha256 588e935996d1006a
\newcommand{\MlstFtJcOneRejHccResOverFull}{1.07\,\ensuremath{\pm}\,0.12}  % experiments/FIGS/data/v2/curves/full_budget_long_leg2.json :: cells.<run>.N1000000.s1.jc1.per_class.label_Hcc: rejection of L162 over L188, run by run :: sha256 588e935996d1006a
\newcommand{\MlstFtJcOneRejHccResOverFullPr}{equal within 20\%}  % experiments/FIGS/data/v2/curves/full_budget_long_leg2.json :: cells.<run>.N1000000.s1.jc1.per_class.label_Hcc: rejection of L162 over L188, run by run :: sha256 588e935996d1006a :: prose twin: two significant figures, direction in words
\newcommand{\MlstFtJcOneRejHccMidOverFull}{0.99\,\ensuremath{\pm}\,0.10}  % experiments/FIGS/data/v2/curves/full_budget_long_leg2.json :: cells.<run>.N1000000.s1.jc1.per_class.label_Hcc: rejection of R42_Q1 over L188, run by run :: sha256 588e935996d1006a
\newcommand{\MlstFtJcOneRejHccMidOverFullPr}{equal within 20\%}  % experiments/FIGS/data/v2/curves/full_budget_long_leg2.json :: cells.<run>.N1000000.s1.jc1.per_class.label_Hcc: rejection of R42_Q1 over L188, run by run :: sha256 588e935996d1006a :: prose twin: two significant figures, direction in words
\newcommand{\MlstFtJcOneRejHccCoarseOverFull}{0.91\,\ensuremath{\pm}\,0.07}  % experiments/FIGS/data/v2/curves/full_budget_long_leg2.json :: cells.<run>.N1000000.s1.jc1.per_class.label_Hcc: rejection of R16_Q1 over L188, run by run :: sha256 588e935996d1006a
\newcommand{\MlstFtJcOneRejHccCoarseOverFullPr}{equal within 10\%}  % experiments/FIGS/data/v2/curves/full_budget_long_leg2.json :: cells.<run>.N1000000.s1.jc1.per_class.label_Hcc: rejection of R16_Q1 over L188, run by run :: sha256 588e935996d1006a :: prose twin: two significant figures, direction in words
\newcommand{\MlstFtJcOneRejHccResMassOverRes}{0.886 and 0.816}  % experiments/FIGS/data/v2/curves/full_budget_long_leg2.json :: cells.<run>.N1000000.s1.jc1.per_class.label_Hcc: rejection of L162_MASS over L162, run by run :: sha256 588e935996d1006a
\newcommand{\MlstFtJcOneRejHccResMassOverResPr}{equal within 20\%}  % experiments/FIGS/data/v2/curves/full_budget_long_leg2.json :: cells.<run>.N1000000.s1.jc1.per_class.label_Hcc: rejection of L162_MASS over L162, run by run :: sha256 588e935996d1006a :: prose twin: two significant figures, direction in words
\newcommand{\MlstFtJcOneRejHccResMassOverResPrMax}{equal within 20\%}  % experiments/FIGS/data/v2/curves/full_budget_long_leg2.json :: cells.<run>.N1000000.s1.jc1.per_class.label_Hcc: rejection of L162_MASS over L162, run by run :: sha256 588e935996d1006a :: prose twin of the run farther from 1
\newcommand{\MlstFtJcOneRejHccCoarseMassOverCoarse}{0.98 and 0.85}  % experiments/FIGS/data/v2/curves/full_budget_long_leg2.json :: cells.<run>.N1000000.s1.jc1.per_class.label_Hcc: rejection of R16_Q1_MASS over R16_Q1, run by run :: sha256 588e935996d1006a
\newcommand{\MlstFtJcOneRejHccCoarseMassOverCoarsePr}{equal within 20\%}  % experiments/FIGS/data/v2/curves/full_budget_long_leg2.json :: cells.<run>.N1000000.s1.jc1.per_class.label_Hcc: rejection of R16_Q1_MASS over R16_Q1, run by run :: sha256 588e935996d1006a :: prose twin: two significant figures, direction in words
\newcommand{\MlstFtJcOneRejHccCoarseMassOverCoarsePrMax}{equal within 20\%}  % experiments/FIGS/data/v2/curves/full_budget_long_leg2.json :: cells.<run>.N1000000.s1.jc1.per_class.label_Hcc: rejection of R16_Q1_MASS over R16_Q1, run by run :: sha256 588e935996d1006a :: prose twin of the run farther from 1
\newcommand{\MlstFtJcOneRejHccRandOverFull}{1.03\,\ensuremath{\pm}\,0.07}  % experiments/FIGS/data/v2/curves/full_budget_long_leg2.json :: per_class.label_Hcc: rejection of each random partition over the 188-class run of its seed :: sha256 588e935996d1006a
\newcommand{\MlstFtJcOneRejHccRandOverFullPr}{equal within 8\%}  % experiments/FIGS/data/v2/curves/full_budget_long_leg2.json :: per_class.label_Hcc: rejection of each random partition over the 188-class run of its seed :: sha256 588e935996d1006a :: prose twin: two significant figures, direction in words
\newcommand{\MlstFtJcOneRejHccRandOverFullRange}{0.94--1.1}  % experiments/FIGS/data/v2/curves/full_budget_long_leg2.json :: per_class.label_Hcc: smallest and largest of the partitions' ratios :: sha256 588e935996d1006a
\newcommand{\MlstFtJcOneRejHggResOverFull}{0.98\,\ensuremath{\pm}\,0.05}  % experiments/FIGS/data/v2/curves/full_budget_long_leg2.json :: cells.<run>.N1000000.s1.jc1.per_class.label_Hgg: rejection of L162 over L188, run by run :: sha256 588e935996d1006a
\newcommand{\MlstFtJcOneRejHggResOverFullPr}{equal within 6\%}  % experiments/FIGS/data/v2/curves/full_budget_long_leg2.json :: cells.<run>.N1000000.s1.jc1.per_class.label_Hgg: rejection of L162 over L188, run by run :: sha256 588e935996d1006a :: prose twin: two significant figures, direction in words
\newcommand{\MlstFtJcOneRejHggMidOverFull}{1.000\,\ensuremath{\pm}\,0.025}  % experiments/FIGS/data/v2/curves/full_budget_long_leg2.json :: cells.<run>.N1000000.s1.jc1.per_class.label_Hgg: rejection of R42_Q1 over L188, run by run :: sha256 588e935996d1006a
\newcommand{\MlstFtJcOneRejHggMidOverFullPr}{equal within 3\%}  % experiments/FIGS/data/v2/curves/full_budget_long_leg2.json :: cells.<run>.N1000000.s1.jc1.per_class.label_Hgg: rejection of R42_Q1 over L188, run by run :: sha256 588e935996d1006a :: prose twin: two significant figures, direction in words
\newcommand{\MlstFtJcOneRejHggCoarseOverFull}{0.97\,\ensuremath{\pm}\,0.09}  % experiments/FIGS/data/v2/curves/full_budget_long_leg2.json :: cells.<run>.N1000000.s1.jc1.per_class.label_Hgg: rejection of R16_Q1 over L188, run by run :: sha256 588e935996d1006a
\newcommand{\MlstFtJcOneRejHggCoarseOverFullPr}{equal within 10\%}  % experiments/FIGS/data/v2/curves/full_budget_long_leg2.json :: cells.<run>.N1000000.s1.jc1.per_class.label_Hgg: rejection of R16_Q1 over L188, run by run :: sha256 588e935996d1006a :: prose twin: two significant figures, direction in words
\newcommand{\MlstFtJcOneRejHggResMassOverRes}{1.046 and 1.008}  % experiments/FIGS/data/v2/curves/full_budget_long_leg2.json :: cells.<run>.N1000000.s1.jc1.per_class.label_Hgg: rejection of L162_MASS over L162, run by run :: sha256 588e935996d1006a
\newcommand{\MlstFtJcOneRejHggResMassOverResPr}{equal within 5\%}  % experiments/FIGS/data/v2/curves/full_budget_long_leg2.json :: cells.<run>.N1000000.s1.jc1.per_class.label_Hgg: rejection of L162_MASS over L162, run by run :: sha256 588e935996d1006a :: prose twin: two significant figures, direction in words
\newcommand{\MlstFtJcOneRejHggResMassOverResPrMax}{equal within 5\%}  % experiments/FIGS/data/v2/curves/full_budget_long_leg2.json :: cells.<run>.N1000000.s1.jc1.per_class.label_Hgg: rejection of L162_MASS over L162, run by run :: sha256 588e935996d1006a :: prose twin of the run farther from 1
\newcommand{\MlstFtJcOneRejHggCoarseMassOverCoarse}{1.12 and 1.01}  % experiments/FIGS/data/v2/curves/full_budget_long_leg2.json :: cells.<run>.N1000000.s1.jc1.per_class.label_Hgg: rejection of R16_Q1_MASS over R16_Q1, run by run :: sha256 588e935996d1006a
\newcommand{\MlstFtJcOneRejHggCoarseMassOverCoarsePr}{equal within 20\%}  % experiments/FIGS/data/v2/curves/full_budget_long_leg2.json :: cells.<run>.N1000000.s1.jc1.per_class.label_Hgg: rejection of R16_Q1_MASS over R16_Q1, run by run :: sha256 588e935996d1006a :: prose twin: two significant figures, direction in words
\newcommand{\MlstFtJcOneRejHggCoarseMassOverCoarsePrMax}{equal within 20\%}  % experiments/FIGS/data/v2/curves/full_budget_long_leg2.json :: cells.<run>.N1000000.s1.jc1.per_class.label_Hgg: rejection of R16_Q1_MASS over R16_Q1, run by run :: sha256 588e935996d1006a :: prose twin of the run farther from 1
\newcommand{\MlstFtJcOneRejHggRandOverFull}{0.968\,\ensuremath{\pm}\,0.016}  % experiments/FIGS/data/v2/curves/full_budget_long_leg2.json :: per_class.label_Hgg: rejection of each random partition over the 188-class run of its seed :: sha256 588e935996d1006a
\newcommand{\MlstFtJcOneRejHggRandOverFullPr}{equal within 4\%}  % experiments/FIGS/data/v2/curves/full_budget_long_leg2.json :: per_class.label_Hgg: rejection of each random partition over the 188-class run of its seed :: sha256 588e935996d1006a :: prose twin: two significant figures, direction in words
\newcommand{\MlstFtJcOneRejHggRandOverFullRange}{0.94--0.98}  % experiments/FIGS/data/v2/curves/full_budget_long_leg2.json :: per_class.label_Hgg: smallest and largest of the partitions' ratios :: sha256 588e935996d1006a
\newcommand{\MlstFtJcOneRejHfourqResOverFull}{1.05\,\ensuremath{\pm}\,0.06}  % experiments/FIGS/data/v2/curves/full_budget_long_leg2.json :: cells.<run>.N1000000.s1.jc1.per_class.label_H4q: rejection of L162 over L188, run by run :: sha256 588e935996d1006a
\newcommand{\MlstFtJcOneRejHfourqResOverFullPr}{equal within 7\%}  % experiments/FIGS/data/v2/curves/full_budget_long_leg2.json :: cells.<run>.N1000000.s1.jc1.per_class.label_H4q: rejection of L162 over L188, run by run :: sha256 588e935996d1006a :: prose twin: two significant figures, direction in words
\newcommand{\MlstFtJcOneRejHfourqMidOverFull}{0.95\,\ensuremath{\pm}\,0.04}  % experiments/FIGS/data/v2/curves/full_budget_long_leg2.json :: cells.<run>.N1000000.s1.jc1.per_class.label_H4q: rejection of R42_Q1 over L188, run by run :: sha256 588e935996d1006a
\newcommand{\MlstFtJcOneRejHfourqMidOverFullPr}{equal within 6\%}  % experiments/FIGS/data/v2/curves/full_budget_long_leg2.json :: cells.<run>.N1000000.s1.jc1.per_class.label_H4q: rejection of R42_Q1 over L188, run by run :: sha256 588e935996d1006a :: prose twin: two significant figures, direction in words
\newcommand{\MlstFtJcOneRejHfourqCoarseOverFull}{0.962\,\ensuremath{\pm}\,0.025}  % experiments/FIGS/data/v2/curves/full_budget_long_leg2.json :: cells.<run>.N1000000.s1.jc1.per_class.label_H4q: rejection of R16_Q1 over L188, run by run :: sha256 588e935996d1006a
\newcommand{\MlstFtJcOneRejHfourqCoarseOverFullPr}{equal within 4\%}  % experiments/FIGS/data/v2/curves/full_budget_long_leg2.json :: cells.<run>.N1000000.s1.jc1.per_class.label_H4q: rejection of R16_Q1 over L188, run by run :: sha256 588e935996d1006a :: prose twin: two significant figures, direction in words
\newcommand{\MlstFtJcOneRejHfourqResMassOverRes}{0.912 and 0.966}  % experiments/FIGS/data/v2/curves/full_budget_long_leg2.json :: cells.<run>.N1000000.s1.jc1.per_class.label_H4q: rejection of L162_MASS over L162, run by run :: sha256 588e935996d1006a
\newcommand{\MlstFtJcOneRejHfourqResMassOverResPr}{equal within 9\%}  % experiments/FIGS/data/v2/curves/full_budget_long_leg2.json :: cells.<run>.N1000000.s1.jc1.per_class.label_H4q: rejection of L162_MASS over L162, run by run :: sha256 588e935996d1006a :: prose twin: two significant figures, direction in words
\newcommand{\MlstFtJcOneRejHfourqResMassOverResPrMax}{equal within 9\%}  % experiments/FIGS/data/v2/curves/full_budget_long_leg2.json :: cells.<run>.N1000000.s1.jc1.per_class.label_H4q: rejection of L162_MASS over L162, run by run :: sha256 588e935996d1006a :: prose twin of the run farther from 1
\newcommand{\MlstFtJcOneRejHfourqCoarseMassOverCoarse}{0.937 and 0.888}  % experiments/FIGS/data/v2/curves/full_budget_long_leg2.json :: cells.<run>.N1000000.s1.jc1.per_class.label_H4q: rejection of R16_Q1_MASS over R16_Q1, run by run :: sha256 588e935996d1006a
\newcommand{\MlstFtJcOneRejHfourqCoarseMassOverCoarsePr}{equal within 20\%}  % experiments/FIGS/data/v2/curves/full_budget_long_leg2.json :: cells.<run>.N1000000.s1.jc1.per_class.label_H4q: rejection of R16_Q1_MASS over R16_Q1, run by run :: sha256 588e935996d1006a :: prose twin: two significant figures, direction in words
\newcommand{\MlstFtJcOneRejHfourqCoarseMassOverCoarsePrMax}{equal within 20\%}  % experiments/FIGS/data/v2/curves/full_budget_long_leg2.json :: cells.<run>.N1000000.s1.jc1.per_class.label_H4q: rejection of R16_Q1_MASS over R16_Q1, run by run :: sha256 588e935996d1006a :: prose twin of the run farther from 1
\newcommand{\MlstFtJcOneRejHfourqRandOverFull}{0.93\,\ensuremath{\pm}\,0.04}  % experiments/FIGS/data/v2/curves/full_budget_long_leg2.json :: per_class.label_H4q: rejection of each random partition over the 188-class run of its seed :: sha256 588e935996d1006a
\newcommand{\MlstFtJcOneRejHfourqRandOverFullPr}{equal within 8\%}  % experiments/FIGS/data/v2/curves/full_budget_long_leg2.json :: per_class.label_H4q: rejection of each random partition over the 188-class run of its seed :: sha256 588e935996d1006a :: prose twin: two significant figures, direction in words
\newcommand{\MlstFtJcOneRejHfourqRandOverFullRange}{0.86--0.96}  % experiments/FIGS/data/v2/curves/full_budget_long_leg2.json :: per_class.label_H4q: smallest and largest of the partitions' ratios :: sha256 588e935996d1006a
\newcommand{\MlstFtJcOneRejHqqlResOverFull}{1.004\,\ensuremath{\pm}\,0.018}  % experiments/FIGS/data/v2/curves/full_budget_long_leg2.json :: cells.<run>.N1000000.s1.jc1.per_class.label_Hqql: rejection of L162 over L188, run by run :: sha256 588e935996d1006a
\newcommand{\MlstFtJcOneRejHqqlResOverFullPr}{equal within 2\%}  % experiments/FIGS/data/v2/curves/full_budget_long_leg2.json :: cells.<run>.N1000000.s1.jc1.per_class.label_Hqql: rejection of L162 over L188, run by run :: sha256 588e935996d1006a :: prose twin: two significant figures, direction in words
\newcommand{\MlstFtJcOneRejHqqlMidOverFull}{1.04\,\ensuremath{\pm}\,0.18}  % experiments/FIGS/data/v2/curves/full_budget_long_leg2.json :: cells.<run>.N1000000.s1.jc1.per_class.label_Hqql: rejection of R42_Q1 over L188, run by run :: sha256 588e935996d1006a
\newcommand{\MlstFtJcOneRejHqqlMidOverFullPr}{equal within 20\%}  % experiments/FIGS/data/v2/curves/full_budget_long_leg2.json :: cells.<run>.N1000000.s1.jc1.per_class.label_Hqql: rejection of R42_Q1 over L188, run by run :: sha256 588e935996d1006a :: prose twin: two significant figures, direction in words
\newcommand{\MlstFtJcOneRejHqqlCoarseOverFull}{0.93\,\ensuremath{\pm}\,0.16}  % experiments/FIGS/data/v2/curves/full_budget_long_leg2.json :: cells.<run>.N1000000.s1.jc1.per_class.label_Hqql: rejection of R16_Q1 over L188, run by run :: sha256 588e935996d1006a
\newcommand{\MlstFtJcOneRejHqqlCoarseOverFullPr}{equal within 20\%}  % experiments/FIGS/data/v2/curves/full_budget_long_leg2.json :: cells.<run>.N1000000.s1.jc1.per_class.label_Hqql: rejection of R16_Q1 over L188, run by run :: sha256 588e935996d1006a :: prose twin: two significant figures, direction in words
\newcommand{\MlstFtJcOneRejHqqlResMassOverRes}{0.986 and 1.000}  % experiments/FIGS/data/v2/curves/full_budget_long_leg2.json :: cells.<run>.N1000000.s1.jc1.per_class.label_Hqql: rejection of L162_MASS over L162, run by run :: sha256 588e935996d1006a
\newcommand{\MlstFtJcOneRejHqqlResMassOverResPr}{equal within 2\%}  % experiments/FIGS/data/v2/curves/full_budget_long_leg2.json :: cells.<run>.N1000000.s1.jc1.per_class.label_Hqql: rejection of L162_MASS over L162, run by run :: sha256 588e935996d1006a :: prose twin: two significant figures, direction in words
\newcommand{\MlstFtJcOneRejHqqlResMassOverResPrMax}{equal within 2\%}  % experiments/FIGS/data/v2/curves/full_budget_long_leg2.json :: cells.<run>.N1000000.s1.jc1.per_class.label_Hqql: rejection of L162_MASS over L162, run by run :: sha256 588e935996d1006a :: prose twin of the run farther from 1
\newcommand{\MlstFtJcOneRejHqqlCoarseMassOverCoarse}{0.97 and 1.09}  % experiments/FIGS/data/v2/curves/full_budget_long_leg2.json :: cells.<run>.N1000000.s1.jc1.per_class.label_Hqql: rejection of R16_Q1_MASS over R16_Q1, run by run :: sha256 588e935996d1006a
\newcommand{\MlstFtJcOneRejHqqlCoarseMassOverCoarsePr}{equal within 9\%}  % experiments/FIGS/data/v2/curves/full_budget_long_leg2.json :: cells.<run>.N1000000.s1.jc1.per_class.label_Hqql: rejection of R16_Q1_MASS over R16_Q1, run by run :: sha256 588e935996d1006a :: prose twin: two significant figures, direction in words
\newcommand{\MlstFtJcOneRejHqqlCoarseMassOverCoarsePrMax}{equal within 9\%}  % experiments/FIGS/data/v2/curves/full_budget_long_leg2.json :: cells.<run>.N1000000.s1.jc1.per_class.label_Hqql: rejection of R16_Q1_MASS over R16_Q1, run by run :: sha256 588e935996d1006a :: prose twin of the run farther from 1
\newcommand{\MlstFtJcOneRejHqqlRandOverFull}{0.94\,\ensuremath{\pm}\,0.12}  % experiments/FIGS/data/v2/curves/full_budget_long_leg2.json :: per_class.label_Hqql: rejection of each random partition over the 188-class run of its seed :: sha256 588e935996d1006a
\newcommand{\MlstFtJcOneRejHqqlRandOverFullPr}{equal within 20\%}  % experiments/FIGS/data/v2/curves/full_budget_long_leg2.json :: per_class.label_Hqql: rejection of each random partition over the 188-class run of its seed :: sha256 588e935996d1006a :: prose twin: two significant figures, direction in words
\newcommand{\MlstFtJcOneRejHqqlRandOverFullRange}{0.79--1.1}  % experiments/FIGS/data/v2/curves/full_budget_long_leg2.json :: per_class.label_Hqql: smallest and largest of the partitions' ratios :: sha256 588e935996d1006a
\newcommand{\MlstFtJcOneRejTbqqResOverFull}{1.1\,\ensuremath{\pm}\,0.4}  % experiments/FIGS/data/v2/curves/full_budget_long_leg2.json :: cells.<run>.N1000000.s1.jc1.per_class.label_Tbqq: rejection of L162 over L188, run by run :: sha256 588e935996d1006a
\newcommand{\MlstFtJcOneRejTbqqResOverFullPr}{equal within 50\%}  % experiments/FIGS/data/v2/curves/full_budget_long_leg2.json :: cells.<run>.N1000000.s1.jc1.per_class.label_Tbqq: rejection of L162 over L188, run by run :: sha256 588e935996d1006a :: prose twin: two significant figures, direction in words
\newcommand{\MlstFtJcOneRejTbqqMidOverFull}{1.10\,\ensuremath{\pm}\,0.05}  % experiments/FIGS/data/v2/curves/full_budget_long_leg2.json :: cells.<run>.N1000000.s1.jc1.per_class.label_Tbqq: rejection of R42_Q1 over L188, run by run :: sha256 588e935996d1006a
\newcommand{\MlstFtJcOneRejTbqqMidOverFullPr}{equal within 10\%}  % experiments/FIGS/data/v2/curves/full_budget_long_leg2.json :: cells.<run>.N1000000.s1.jc1.per_class.label_Tbqq: rejection of R42_Q1 over L188, run by run :: sha256 588e935996d1006a :: prose twin: two significant figures, direction in words
\newcommand{\MlstFtJcOneRejTbqqCoarseOverFull}{0.85\,\ensuremath{\pm}\,0.32}  % experiments/FIGS/data/v2/curves/full_budget_long_leg2.json :: cells.<run>.N1000000.s1.jc1.per_class.label_Tbqq: rejection of R16_Q1 over L188, run by run :: sha256 588e935996d1006a
\newcommand{\MlstFtJcOneRejTbqqCoarseOverFullPr}{equal within 40\%}  % experiments/FIGS/data/v2/curves/full_budget_long_leg2.json :: cells.<run>.N1000000.s1.jc1.per_class.label_Tbqq: rejection of R16_Q1 over L188, run by run :: sha256 588e935996d1006a :: prose twin: two significant figures, direction in words
\newcommand{\MlstFtJcOneRejTbqqResMassOverRes}{0.89 and 1.21}  % experiments/FIGS/data/v2/curves/full_budget_long_leg2.json :: cells.<run>.N1000000.s1.jc1.per_class.label_Tbqq: rejection of L162_MASS over L162, run by run :: sha256 588e935996d1006a
\newcommand{\MlstFtJcOneRejTbqqResMassOverResPr}{equal within 30\%}  % experiments/FIGS/data/v2/curves/full_budget_long_leg2.json :: cells.<run>.N1000000.s1.jc1.per_class.label_Tbqq: rejection of L162_MASS over L162, run by run :: sha256 588e935996d1006a :: prose twin: two significant figures, direction in words
\newcommand{\MlstFtJcOneRejTbqqResMassOverResPrMax}{equal within 30\%}  % experiments/FIGS/data/v2/curves/full_budget_long_leg2.json :: cells.<run>.N1000000.s1.jc1.per_class.label_Tbqq: rejection of L162_MASS over L162, run by run :: sha256 588e935996d1006a :: prose twin of the run farther from 1
\newcommand{\MlstFtJcOneRejTbqqCoarseMassOverCoarse}{1.00 and 1.79}  % experiments/FIGS/data/v2/curves/full_budget_long_leg2.json :: cells.<run>.N1000000.s1.jc1.per_class.label_Tbqq: rejection of R16_Q1_MASS over R16_Q1, run by run :: sha256 588e935996d1006a
\newcommand{\MlstFtJcOneRejTbqqCoarseMassOverCoarsePr}{equal within 80\%}  % experiments/FIGS/data/v2/curves/full_budget_long_leg2.json :: cells.<run>.N1000000.s1.jc1.per_class.label_Tbqq: rejection of R16_Q1_MASS over R16_Q1, run by run :: sha256 588e935996d1006a :: prose twin: two significant figures, direction in words
\newcommand{\MlstFtJcOneRejTbqqCoarseMassOverCoarsePrMax}{equal within 80\%}  % experiments/FIGS/data/v2/curves/full_budget_long_leg2.json :: cells.<run>.N1000000.s1.jc1.per_class.label_Tbqq: rejection of R16_Q1_MASS over R16_Q1, run by run :: sha256 588e935996d1006a :: prose twin of the run farther from 1
\newcommand{\MlstFtJcOneRejTbqqRandOverFull}{0.79\,\ensuremath{\pm}\,0.15}  % experiments/FIGS/data/v2/curves/full_budget_long_leg2.json :: per_class.label_Tbqq: rejection of each random partition over the 188-class run of its seed :: sha256 588e935996d1006a
\newcommand{\MlstFtJcOneRejTbqqRandOverFullPr}{equal within 30\%}  % experiments/FIGS/data/v2/curves/full_budget_long_leg2.json :: per_class.label_Tbqq: rejection of each random partition over the 188-class run of its seed :: sha256 588e935996d1006a :: prose twin: two significant figures, direction in words
\newcommand{\MlstFtJcOneRejTbqqRandOverFullRange}{0.64--1.0}  % experiments/FIGS/data/v2/curves/full_budget_long_leg2.json :: per_class.label_Tbqq: smallest and largest of the partitions' ratios :: sha256 588e935996d1006a
\newcommand{\MlstFtJcOneRejTblResOverFull}{0.85\,\ensuremath{\pm}\,0.17}  % experiments/FIGS/data/v2/curves/full_budget_long_leg2.json :: cells.<run>.N1000000.s1.jc1.per_class.label_Tbl: rejection of L162 over L188, run by run :: sha256 588e935996d1006a
\newcommand{\MlstFtJcOneRejTblResOverFullPr}{equal within 20\%}  % experiments/FIGS/data/v2/curves/full_budget_long_leg2.json :: cells.<run>.N1000000.s1.jc1.per_class.label_Tbl: rejection of L162 over L188, run by run :: sha256 588e935996d1006a :: prose twin: two significant figures, direction in words
\newcommand{\MlstFtJcOneRejTblMidOverFull}{0.91\,\ensuremath{\pm}\,0.15}  % experiments/FIGS/data/v2/curves/full_budget_long_leg2.json :: cells.<run>.N1000000.s1.jc1.per_class.label_Tbl: rejection of R42_Q1 over L188, run by run :: sha256 588e935996d1006a
\newcommand{\MlstFtJcOneRejTblMidOverFullPr}{equal within 20\%}  % experiments/FIGS/data/v2/curves/full_budget_long_leg2.json :: cells.<run>.N1000000.s1.jc1.per_class.label_Tbl: rejection of R42_Q1 over L188, run by run :: sha256 588e935996d1006a :: prose twin: two significant figures, direction in words
\newcommand{\MlstFtJcOneRejTblCoarseOverFull}{0.94\,\ensuremath{\pm}\,0.28}  % experiments/FIGS/data/v2/curves/full_budget_long_leg2.json :: cells.<run>.N1000000.s1.jc1.per_class.label_Tbl: rejection of R16_Q1 over L188, run by run :: sha256 588e935996d1006a
\newcommand{\MlstFtJcOneRejTblCoarseOverFullPr}{equal within 30\%}  % experiments/FIGS/data/v2/curves/full_budget_long_leg2.json :: cells.<run>.N1000000.s1.jc1.per_class.label_Tbl: rejection of R16_Q1 over L188, run by run :: sha256 588e935996d1006a :: prose twin: two significant figures, direction in words
\newcommand{\MlstFtJcOneRejTblResMassOverRes}{1.04 and 0.93}  % experiments/FIGS/data/v2/curves/full_budget_long_leg2.json :: cells.<run>.N1000000.s1.jc1.per_class.label_Tbl: rejection of L162_MASS over L162, run by run :: sha256 588e935996d1006a
\newcommand{\MlstFtJcOneRejTblResMassOverResPr}{equal within 7\%}  % experiments/FIGS/data/v2/curves/full_budget_long_leg2.json :: cells.<run>.N1000000.s1.jc1.per_class.label_Tbl: rejection of L162_MASS over L162, run by run :: sha256 588e935996d1006a :: prose twin: two significant figures, direction in words
\newcommand{\MlstFtJcOneRejTblResMassOverResPrMax}{equal within 7\%}  % experiments/FIGS/data/v2/curves/full_budget_long_leg2.json :: cells.<run>.N1000000.s1.jc1.per_class.label_Tbl: rejection of L162_MASS over L162, run by run :: sha256 588e935996d1006a :: prose twin of the run farther from 1
\newcommand{\MlstFtJcOneRejTblCoarseMassOverCoarse}{0.84 and 1.22}  % experiments/FIGS/data/v2/curves/full_budget_long_leg2.json :: cells.<run>.N1000000.s1.jc1.per_class.label_Tbl: rejection of R16_Q1_MASS over R16_Q1, run by run :: sha256 588e935996d1006a
\newcommand{\MlstFtJcOneRejTblCoarseMassOverCoarsePr}{equal within 30\%}  % experiments/FIGS/data/v2/curves/full_budget_long_leg2.json :: cells.<run>.N1000000.s1.jc1.per_class.label_Tbl: rejection of R16_Q1_MASS over R16_Q1, run by run :: sha256 588e935996d1006a :: prose twin: two significant figures, direction in words
\newcommand{\MlstFtJcOneRejTblCoarseMassOverCoarsePrMax}{equal within 30\%}  % experiments/FIGS/data/v2/curves/full_budget_long_leg2.json :: cells.<run>.N1000000.s1.jc1.per_class.label_Tbl: rejection of R16_Q1_MASS over R16_Q1, run by run :: sha256 588e935996d1006a :: prose twin of the run farther from 1
\newcommand{\MlstFtJcOneRejTblRandOverFull}{0.88\,\ensuremath{\pm}\,0.26}  % experiments/FIGS/data/v2/curves/full_budget_long_leg2.json :: per_class.label_Tbl: rejection of each random partition over the 188-class run of its seed :: sha256 588e935996d1006a
\newcommand{\MlstFtJcOneRejTblRandOverFullPr}{equal within 30\%}  % experiments/FIGS/data/v2/curves/full_budget_long_leg2.json :: per_class.label_Tbl: rejection of each random partition over the 188-class run of its seed :: sha256 588e935996d1006a :: prose twin: two significant figures, direction in words
\newcommand{\MlstFtJcOneRejTblRandOverFullRange}{0.61--1.2}  % experiments/FIGS/data/v2/curves/full_budget_long_leg2.json :: per_class.label_Tbl: smallest and largest of the partitions' ratios :: sha256 588e935996d1006a
\newcommand{\MlstFtJcOneRejWqqResOverFull}{1.039\,\ensuremath{\pm}\,0.035}  % experiments/FIGS/data/v2/curves/full_budget_long_leg2.json :: cells.<run>.N1000000.s1.jc1.per_class.label_Wqq: rejection of L162 over L188, run by run :: sha256 588e935996d1006a
\newcommand{\MlstFtJcOneRejWqqResOverFullPr}{equal within 4\%}  % experiments/FIGS/data/v2/curves/full_budget_long_leg2.json :: cells.<run>.N1000000.s1.jc1.per_class.label_Wqq: rejection of L162 over L188, run by run :: sha256 588e935996d1006a :: prose twin: two significant figures, direction in words
\newcommand{\MlstFtJcOneRejWqqMidOverFull}{1.005\,\ensuremath{\pm}\,0.032}  % experiments/FIGS/data/v2/curves/full_budget_long_leg2.json :: cells.<run>.N1000000.s1.jc1.per_class.label_Wqq: rejection of R42_Q1 over L188, run by run :: sha256 588e935996d1006a
\newcommand{\MlstFtJcOneRejWqqMidOverFullPr}{equal within 4\%}  % experiments/FIGS/data/v2/curves/full_budget_long_leg2.json :: cells.<run>.N1000000.s1.jc1.per_class.label_Wqq: rejection of R42_Q1 over L188, run by run :: sha256 588e935996d1006a :: prose twin: two significant figures, direction in words
\newcommand{\MlstFtJcOneRejWqqCoarseOverFull}{1.01\,\ensuremath{\pm}\,0.07}  % experiments/FIGS/data/v2/curves/full_budget_long_leg2.json :: cells.<run>.N1000000.s1.jc1.per_class.label_Wqq: rejection of R16_Q1 over L188, run by run :: sha256 588e935996d1006a
\newcommand{\MlstFtJcOneRejWqqCoarseOverFullPr}{equal within 8\%}  % experiments/FIGS/data/v2/curves/full_budget_long_leg2.json :: cells.<run>.N1000000.s1.jc1.per_class.label_Wqq: rejection of R16_Q1 over L188, run by run :: sha256 588e935996d1006a :: prose twin: two significant figures, direction in words
\newcommand{\MlstFtJcOneRejWqqResMassOverRes}{0.996 and 0.964}  % experiments/FIGS/data/v2/curves/full_budget_long_leg2.json :: cells.<run>.N1000000.s1.jc1.per_class.label_Wqq: rejection of L162_MASS over L162, run by run :: sha256 588e935996d1006a
\newcommand{\MlstFtJcOneRejWqqResMassOverResPr}{equal within 4\%}  % experiments/FIGS/data/v2/curves/full_budget_long_leg2.json :: cells.<run>.N1000000.s1.jc1.per_class.label_Wqq: rejection of L162_MASS over L162, run by run :: sha256 588e935996d1006a :: prose twin: two significant figures, direction in words
\newcommand{\MlstFtJcOneRejWqqResMassOverResPrMax}{equal within 4\%}  % experiments/FIGS/data/v2/curves/full_budget_long_leg2.json :: cells.<run>.N1000000.s1.jc1.per_class.label_Wqq: rejection of L162_MASS over L162, run by run :: sha256 588e935996d1006a :: prose twin of the run farther from 1
\newcommand{\MlstFtJcOneRejWqqCoarseMassOverCoarse}{1.025 and 0.926}  % experiments/FIGS/data/v2/curves/full_budget_long_leg2.json :: cells.<run>.N1000000.s1.jc1.per_class.label_Wqq: rejection of R16_Q1_MASS over R16_Q1, run by run :: sha256 588e935996d1006a
\newcommand{\MlstFtJcOneRejWqqCoarseMassOverCoarsePr}{equal within 8\%}  % experiments/FIGS/data/v2/curves/full_budget_long_leg2.json :: cells.<run>.N1000000.s1.jc1.per_class.label_Wqq: rejection of R16_Q1_MASS over R16_Q1, run by run :: sha256 588e935996d1006a :: prose twin: two significant figures, direction in words
\newcommand{\MlstFtJcOneRejWqqCoarseMassOverCoarsePrMax}{equal within 8\%}  % experiments/FIGS/data/v2/curves/full_budget_long_leg2.json :: cells.<run>.N1000000.s1.jc1.per_class.label_Wqq: rejection of R16_Q1_MASS over R16_Q1, run by run :: sha256 588e935996d1006a :: prose twin of the run farther from 1
\newcommand{\MlstFtJcOneRejWqqRandOverFull}{1.02\,\ensuremath{\pm}\,0.07}  % experiments/FIGS/data/v2/curves/full_budget_long_leg2.json :: per_class.label_Wqq: rejection of each random partition over the 188-class run of its seed :: sha256 588e935996d1006a
\newcommand{\MlstFtJcOneRejWqqRandOverFullPr}{equal within 7\%}  % experiments/FIGS/data/v2/curves/full_budget_long_leg2.json :: per_class.label_Wqq: rejection of each random partition over the 188-class run of its seed :: sha256 588e935996d1006a :: prose twin: two significant figures, direction in words
\newcommand{\MlstFtJcOneRejWqqRandOverFullRange}{0.93--1.1}  % experiments/FIGS/data/v2/curves/full_budget_long_leg2.json :: per_class.label_Wqq: smallest and largest of the partitions' ratios :: sha256 588e935996d1006a
\newcommand{\MlstFtJcOneRejZqqResOverFull}{0.987\,\ensuremath{\pm}\,0.019}  % experiments/FIGS/data/v2/curves/full_budget_long_leg2.json :: cells.<run>.N1000000.s1.jc1.per_class.label_Zqq: rejection of L162 over L188, run by run :: sha256 588e935996d1006a
\newcommand{\MlstFtJcOneRejZqqResOverFullPr}{equal within 2\%}  % experiments/FIGS/data/v2/curves/full_budget_long_leg2.json :: cells.<run>.N1000000.s1.jc1.per_class.label_Zqq: rejection of L162 over L188, run by run :: sha256 588e935996d1006a :: prose twin: two significant figures, direction in words
\newcommand{\MlstFtJcOneRejZqqMidOverFull}{0.984\,\ensuremath{\pm}\,0.010}  % experiments/FIGS/data/v2/curves/full_budget_long_leg2.json :: cells.<run>.N1000000.s1.jc1.per_class.label_Zqq: rejection of R42_Q1 over L188, run by run :: sha256 588e935996d1006a
\newcommand{\MlstFtJcOneRejZqqMidOverFullPr}{equal within 2\%}  % experiments/FIGS/data/v2/curves/full_budget_long_leg2.json :: cells.<run>.N1000000.s1.jc1.per_class.label_Zqq: rejection of R42_Q1 over L188, run by run :: sha256 588e935996d1006a :: prose twin: two significant figures, direction in words
\newcommand{\MlstFtJcOneRejZqqCoarseOverFull}{0.91\,\ensuremath{\pm}\,0.05}  % experiments/FIGS/data/v2/curves/full_budget_long_leg2.json :: cells.<run>.N1000000.s1.jc1.per_class.label_Zqq: rejection of R16_Q1 over L188, run by run :: sha256 588e935996d1006a
\newcommand{\MlstFtJcOneRejZqqCoarseOverFullPr}{equal within 9\%}  % experiments/FIGS/data/v2/curves/full_budget_long_leg2.json :: cells.<run>.N1000000.s1.jc1.per_class.label_Zqq: rejection of R16_Q1 over L188, run by run :: sha256 588e935996d1006a :: prose twin: two significant figures, direction in words
\newcommand{\MlstFtJcOneRejZqqResMassOverRes}{0.9803 and 0.9820}  % experiments/FIGS/data/v2/curves/full_budget_long_leg2.json :: cells.<run>.N1000000.s1.jc1.per_class.label_Zqq: rejection of L162_MASS over L162, run by run :: sha256 588e935996d1006a
\newcommand{\MlstFtJcOneRejZqqResMassOverResPr}{equal within 2\%}  % experiments/FIGS/data/v2/curves/full_budget_long_leg2.json :: cells.<run>.N1000000.s1.jc1.per_class.label_Zqq: rejection of L162_MASS over L162, run by run :: sha256 588e935996d1006a :: prose twin: two significant figures, direction in words
\newcommand{\MlstFtJcOneRejZqqResMassOverResPrMax}{equal within 2\%}  % experiments/FIGS/data/v2/curves/full_budget_long_leg2.json :: cells.<run>.N1000000.s1.jc1.per_class.label_Zqq: rejection of L162_MASS over L162, run by run :: sha256 588e935996d1006a :: prose twin of the run farther from 1
\newcommand{\MlstFtJcOneRejZqqCoarseMassOverCoarse}{1.06 and 0.94}  % experiments/FIGS/data/v2/curves/full_budget_long_leg2.json :: cells.<run>.N1000000.s1.jc1.per_class.label_Zqq: rejection of R16_Q1_MASS over R16_Q1, run by run :: sha256 588e935996d1006a
\newcommand{\MlstFtJcOneRejZqqCoarseMassOverCoarsePr}{equal within 7\%}  % experiments/FIGS/data/v2/curves/full_budget_long_leg2.json :: cells.<run>.N1000000.s1.jc1.per_class.label_Zqq: rejection of R16_Q1_MASS over R16_Q1, run by run :: sha256 588e935996d1006a :: prose twin: two significant figures, direction in words
\newcommand{\MlstFtJcOneRejZqqCoarseMassOverCoarsePrMax}{equal within 7\%}  % experiments/FIGS/data/v2/curves/full_budget_long_leg2.json :: cells.<run>.N1000000.s1.jc1.per_class.label_Zqq: rejection of R16_Q1_MASS over R16_Q1, run by run :: sha256 588e935996d1006a :: prose twin of the run farther from 1
\newcommand{\MlstFtJcOneRejZqqRandOverFull}{0.952\,\ensuremath{\pm}\,0.028}  % experiments/FIGS/data/v2/curves/full_budget_long_leg2.json :: per_class.label_Zqq: rejection of each random partition over the 188-class run of its seed :: sha256 588e935996d1006a
\newcommand{\MlstFtJcOneRejZqqRandOverFullPr}{equal within 5\%}  % experiments/FIGS/data/v2/curves/full_budget_long_leg2.json :: per_class.label_Zqq: rejection of each random partition over the 188-class run of its seed :: sha256 588e935996d1006a :: prose twin: two significant figures, direction in words
\newcommand{\MlstFtJcOneRejZqqRandOverFullRange}{0.93--0.99}  % experiments/FIGS/data/v2/curves/full_budget_long_leg2.json :: per_class.label_Zqq: smallest and largest of the partitions' ratios :: sha256 588e935996d1006a
\newcommand{\MlstFtJcOneResNRejDrop}{zero}  % experiments/FIGS/data/v2/curves/full_budget_long_leg2.json :: per_class rejections :: sha256 588e935996d1006a :: classes whose paired rejection ratio to the 188-class model departs from 1 by more than twice its SD and its counting uncertainty, columns with fewer than JC1_MIN_BKG passing jets excluded, ratio below 1 (count)
\newcommand{\MlstFtJcOneResNRejDiff}{zero}  % experiments/FIGS/data/v2/curves/full_budget_long_leg2.json :: per_class rejections :: sha256 588e935996d1006a :: classes whose paired rejection ratio to the 188-class model departs from 1 by more than twice its SD and its counting uncertainty, columns with fewer than JC1_MIN_BKG passing jets excluded (count)
\newcommand{\MlstFtJcOneResRejDropClasses}{none}  % experiments/FIGS/data/v2/curves/full_budget_long_leg2.json :: per_class rejections :: sha256 588e935996d1006a :: the classes of ...NRejDrop, in ParT's column order
\newcommand{\MlstFtJcOneMidNRejDrop}{zero}  % experiments/FIGS/data/v2/curves/full_budget_long_leg2.json :: per_class rejections :: sha256 588e935996d1006a :: classes whose paired rejection ratio to the 188-class model departs from 1 by more than twice its SD and its counting uncertainty, columns with fewer than JC1_MIN_BKG passing jets excluded, ratio below 1 (count)
\newcommand{\MlstFtJcOneMidNRejDiff}{zero}  % experiments/FIGS/data/v2/curves/full_budget_long_leg2.json :: per_class rejections :: sha256 588e935996d1006a :: classes whose paired rejection ratio to the 188-class model departs from 1 by more than twice its SD and its counting uncertainty, columns with fewer than JC1_MIN_BKG passing jets excluded (count)
\newcommand{\MlstFtJcOneMidRejDropClasses}{none}  % experiments/FIGS/data/v2/curves/full_budget_long_leg2.json :: per_class rejections :: sha256 588e935996d1006a :: the classes of ...NRejDrop, in ParT's column order
\newcommand{\MlstFtJcOneCoarseNRejDrop}{zero}  % experiments/FIGS/data/v2/curves/full_budget_long_leg2.json :: per_class rejections :: sha256 588e935996d1006a :: classes whose paired rejection ratio to the 188-class model departs from 1 by more than twice its SD and its counting uncertainty, columns with fewer than JC1_MIN_BKG passing jets excluded, ratio below 1 (count)
\newcommand{\MlstFtJcOneCoarseNRejDiff}{zero}  % experiments/FIGS/data/v2/curves/full_budget_long_leg2.json :: per_class rejections :: sha256 588e935996d1006a :: classes whose paired rejection ratio to the 188-class model departs from 1 by more than twice its SD and its counting uncertainty, columns with fewer than JC1_MIN_BKG passing jets excluded (count)
\newcommand{\MlstFtJcOneCoarseRejDropClasses}{none}  % experiments/FIGS/data/v2/curves/full_budget_long_leg2.json :: per_class rejections :: sha256 588e935996d1006a :: the classes of ...NRejDrop, in ParT's column order
\newcommand{\MlstFtJcOneAccRandMinusCoarse}{+0.002}  % experiments/FIGS/data/v2/curves/full_budget_long_leg2.json :: cells.*.N1000000.s1.jc1.accuracy :: sha256 588e935996d1006a :: mean over the random partitions minus the mean over the 17-class runs
\newcommand{\MlstFtJcOneAccFullBoot}{0.847}  % experiments/FIGS/data/v2/budget/jc1_acc_boot_N1000000.json :: diffs.r16q1-l188.accuracy_ref (mean) :: sha256 7cccb7b0ff900fb8
\newcommand{\MlstFtJcOneAccCoarseBoot}{0.842}  % experiments/FIGS/data/v2/budget/jc1_acc_boot_N1000000.json :: diffs.r16q1-l188.accuracy_config (mean) :: sha256 7cccb7b0ff900fb8
\newcommand{\MlstFtJcOneErrResOverFullBoot}{consistent}  % experiments/FIGS/data/v2/budget/jc1_acc_boot_N1000000.json :: diffs.l162-l188.error_ratio :: sha256 7cccb7b0ff900fb8 :: delta and sigma = hypot(run SD, bootstrap SD) in per cent; 'consistent' when |delta| <= 3 sigma
\newcommand{\MlstFtJcOneErrResOverFullZ}{0.6}  % experiments/FIGS/data/v2/budget/jc1_acc_boot_N1000000.json :: diffs.l162-l188.error_ratio.z :: sha256 7cccb7b0ff900fb8 :: |delta| / sigma
\newcommand{\MlstFtJcOneErrMidOverFullBoot}{consistent}  % experiments/FIGS/data/v2/budget/jc1_acc_boot_N1000000.json :: diffs.r42q1-l188.error_ratio :: sha256 7cccb7b0ff900fb8 :: delta and sigma = hypot(run SD, bootstrap SD) in per cent; 'consistent' when |delta| <= 3 sigma
\newcommand{\MlstFtJcOneErrMidOverFullZ}{1.4}  % experiments/FIGS/data/v2/budget/jc1_acc_boot_N1000000.json :: diffs.r42q1-l188.error_ratio.z :: sha256 7cccb7b0ff900fb8 :: |delta| / sigma
\newcommand{\MlstFtJcOneErrCoarseOverFullBoot}{(3.0\,\ensuremath{\pm}\,0.6)\% more errors}  % experiments/FIGS/data/v2/budget/jc1_acc_boot_N1000000.json :: diffs.r16q1-l188.error_ratio :: sha256 7cccb7b0ff900fb8 :: delta and sigma = hypot(run SD, bootstrap SD) in per cent; 'consistent' when |delta| <= 3 sigma
\newcommand{\MlstFtJcOneErrCoarseOverFullZ}{4.8}  % experiments/FIGS/data/v2/budget/jc1_acc_boot_N1000000.json :: diffs.r16q1-l188.error_ratio.z :: sha256 7cccb7b0ff900fb8 :: |delta| / sigma
\newcommand{\MlstFtJcOneBootReplicas}{200}  % experiments/FIGS/data/v2/budget/jc1_acc_boot_N1000000.json :: replicas :: sha256 7cccb7b0ff900fb8
\newcommand{\MlstFtJcOneBootJets}{2{,}000{,}000}  % experiments/FIGS/data/v2/budget/jc1_acc_boot_N1000000.json :: n_test_jets :: sha256 7cccb7b0ff900fb8
\newcommand{\ComputeTotalGpuHours}{3{,}435}  % experiments/FIGS/data/v2/index.json :: GPU-hours of every reported pretraining run, training and validation passes :: sha256 c6c4138ee5c86423
\newcommand{\MlstComputeTotalGpuHours}{3{,}510}  % experiments/FIGS/data/v2/index.json :: GPU-hours of every pretraining run the paper reports, the kinematic-only model included :: sha256 c6c4138ee5c86423
\newcommand{\ComputeMainGpu}{V100}  % experiments/FIGS/data/v2/index.json :: GPU type with the most hours over the four label sets :: sha256 c6c4138ee5c86423
\newcommand{\ComputeMainGpuShareMin}{56\%}  % experiments/FIGS/data/v2/index.json :: lowest share of a label set's GPU-hours on that type :: sha256 c6c4138ee5c86423
\newcommand{\ComputeMainGpuShareMax}{71\%}  % experiments/FIGS/data/v2/index.json :: highest share of a label set's GPU-hours on that type :: sha256 c6c4138ee5c86423
\newcommand{\ComputeRunsDone}{27}  % experiments/FIGS/data/v2/index.json :: paper runs with a compute record (count) :: sha256 c6c4138ee5c86423
\newcommand{\ComputeGpuTypes}{L4, RTX 3090, RTX 4090 and V100}  % experiments/FIGS/data/v2/index.json :: GPU types any finished paper run trained on :: sha256 c6c4138ee5c86423
\newcommand{\MlstComputeGpuTypes}{L4, RTX 3090, RTX 4090 and V100}  % experiments/FIGS/data/v2/index.json :: GPU types of the pretraining runs the paper reports :: sha256 c6c4138ee5c86423
\newcommand{\MlstFtJcOneErrCoarseOverFullBootEFour}{\pending{the JetClass bootstrap at 1e4 (after the equal-stopping read-out)}}  % experiments/FIGS/data/v2/index.json :: not computed yet :: sha256 c6c4138ee5c86423 :: waits on the JetClass bootstrap at 1e4 (after the equal-stopping read-out)
\newcommand{\MlstLpBbQgZeroAccDropFull}{\pending{pair probes with the quark/gluon inputs zeroed (open item O10, not run)}}  % experiments/FIGS/data/v2/index.json :: not computed yet :: sha256 c6c4138ee5c86423 :: waits on pair probes with the quark/gluon inputs zeroed (open item O10, not run)
\newcommand{\MlstLpBbQgZeroAccDropCoarse}{\pending{pair probes with the quark/gluon inputs zeroed (open item O10, not run)}}  % experiments/FIGS/data/v2/index.json :: not computed yet :: sha256 c6c4138ee5c86423 :: waits on pair probes with the quark/gluon inputs zeroed (open item O10, not run)
\newcommand{\XbsRthirtyFullOverScratch}{\pending{X->bs fits and from-scratch runs}}  % experiments/FIGS/data/v2/index.json :: not computed yet :: sha256 c6c4138ee5c86423 :: waits on X->bs fits and from-scratch runs
\newcommand{\XbsRthirtyFullOverScratchPr}{\pending{X->bs fits and from-scratch runs}}  % experiments/FIGS/data/v2/index.json :: not computed yet :: sha256 c6c4138ee5c86423 :: waits on X->bs fits and from-scratch runs
\newcommand{\XbsRthirtyCoarseOverFull}{\pending{X->bs fits}}  % experiments/FIGS/data/v2/index.json :: not computed yet :: sha256 c6c4138ee5c86423 :: waits on X->bs fits
\newcommand{\XbsRthirtyCoarseOverFullPr}{\pending{X->bs fits}}  % experiments/FIGS/data/v2/index.json :: not computed yet :: sha256 c6c4138ee5c86423 :: waits on X->bs fits
\newcommand{\XbsRthirtyNodeOverTransferFull}{\pending{X->bs fits}}  % experiments/FIGS/data/v2/index.json :: not computed yet :: sha256 c6c4138ee5c86423 :: waits on X->bs fits
\newcommand{\XbsRthirtyNodeOverTransferFullPr}{\pending{X->bs fits}}  % experiments/FIGS/data/v2/index.json :: not computed yet :: sha256 c6c4138ee5c86423 :: waits on X->bs fits
\newcommand{\XbbRthirtyDirectOverTransferFull}{\pending{X->bb fits}}  % experiments/FIGS/data/v2/index.json :: not computed yet :: sha256 c6c4138ee5c86423 :: waits on X->bb fits
\newcommand{\XbbRthirtyDirectOverTransferFullPr}{\pending{X->bb fits}}  % experiments/FIGS/data/v2/index.json :: not computed yet :: sha256 c6c4138ee5c86423 :: waits on X->bb fits
\newcommand{\XbbRthirtyTransferCoarseOverFull}{\pending{X->bb fits}}  % experiments/FIGS/data/v2/index.json :: not computed yet :: sha256 c6c4138ee5c86423 :: waits on X->bb fits
\newcommand{\XbbRthirtyTransferCoarseOverFullPr}{\pending{X->bb fits}}  % experiments/FIGS/data/v2/index.json :: not computed yet :: sha256 c6c4138ee5c86423 :: waits on X->bb fits
             % the counts below are its macros
% Visual language: Sophon's Fig. 1 (arXiv:2405.12972): the document's serif font, light fills,
% rounded boxes, light grey panels behind the stages (PI 2026-11-02). The text is set in the
% document's font at the printed sizes 8.5-10 pt (MLST 8-12 pt), never scaled.
% Text width: 15 cm (MLST full width). The picture is 15.0 cm wide and is drawn at its final
% size, never scaled, so the font sizes hold.
% Optional: \def\FoneDropXbx{} before % Figure 1 (F1 of docs/FIGURE_PLAN.md): the design of the study, in TikZ.
% % Figure 1 (F1 of docs/FIGURE_PLAN.md): the design of the study, in TikZ.
% % Figure 1 (F1 of docs/FIGURE_PLAN.md): the design of the study, in TikZ.
% \input{fig_f1} from main.tex; this file holds one figure environment and nothing else.
%
% WHAT main.tex MUST LOAD (this file loads nothing itself):
%   \usepackage{tikz}
%   \usetikzlibrary{arrows.meta}          % the only TikZ library used (arrow tips)
%   \definecolor{c188}{HTML}{1F77B4}      % the configuration colours of experiments/FIGS/
%   \definecolor{c162}{HTML}{FF7F0E}      % mlst_style.py (matplotlib's tab10 in order, as in
%   \definecolor{c43}{HTML}{2CA02C}       % Sophon's figures); the variants are drawn in their
%   \definecolor{c17}{HTML}{D62728}       % label set's colour with a dotted border, as the
%   \definecolor{crand}{HTML}{17BECF}     % plotted figures draw variants dotted
%   \definecolor{cflav}{HTML}{9467BD}
%   \input{results_generated}             % the counts below are its macros
% Visual language: Sophon's Fig. 1 (arXiv:2405.12972): the document's serif font, light fills,
% rounded boxes, light grey panels behind the stages (PI 2026-11-02). The text is set in the
% document's font at the printed sizes 8.5-10 pt (MLST 8-12 pt), never scaled.
% Text width: 15 cm (MLST full width). The picture is 15.0 cm wide and is drawn at its final
% size, never scaled, so the font sizes hold.
% Optional: \def\FoneDropXbx{} before \input{fig_f1} removes X->bb/X->bs from panel (c)
% (PI 2026-11-02: if the samples are not staged by 12 Nov, the X->bb/X->bs section goes).
%
% Every class box is drawn from the label maps (checked by hand, 2026-11-02):
%   (a) configs/labelmaps/rung_label_maps.v1.csv, columns L188, L162, R42_Q1, R16_Q1. Rows: the ten
%       decays X->QQ/gg in the order bb bc bq | cc cs cq | ss sq qq gg (the 43-class b / c / light
%       groups are then contiguous), then ee mm | tauh-taue tauh-taum, then the 27 QCD classes.
%   (b) column: configs/labelmaps/rand_label_map.v2.csv, RAND2_p1 on the same rows: groups
%       0 0 0 0 0 9 9 9 0 3 | 11 11 | 15 4 | QCD, lettered A..F in order of first appearance.
%       strips: configs/labelmaps/flavour_pair_map.v2.csv, the 22 classes of orbit X_YY_QQQQ.
%       Order with a b quark: bbbb bbcc bbss bbqq qqqb bbcq ccbs ccbq qqbc qqbs bcsq;
%       without: cccc ccss ccqq ssss ssqq qqqq qqqc qqqs ccsq sscq qqcs.
%       FLAV_F0: all 22 in group 9. FLAV_F1: the 11 with a b in group 11, the rest 9.
%       FLAV_F1R: group 11 = bbss bbqq qqqb bbcq ccbq qqbc (with b, positions 2-5, 7, 8 from 0)
%       and ccss ccqq ssss qqqc sscq (without, positions 1-3, 6, 9); filled square = group 11.
% Geometry (cm): row height 0.35 from y = -1.18, gaps of 0.14 between the quark/gluon, lepton
% and QCD blocks, the QCD row 0.70; boxes inset 0.025 top and bottom; columns 1.45 wide with
% left edges 1.75, 3.28, 4.81, 6.34 (hierarchy) and 8.25 (random partition).
\providecommand{\Kfull}{\VocabSizeLoneeighteight}
\providecommand{\Kres}{\VocabSizeLonesixtwo}
\providecommand{\Kmid}{\VocabSizeRfourtwoQone}
\providecommand{\Kcoarse}{\VocabSizeRonesixQone}
\begin{figure}[tbp]
\centering
\ifdim\linewidth<14.99cm\typeout{fig_f1.tex: WARNING: the drawing is 15 cm wide but \string\linewidth\space
  is \the\linewidth; set a 15 cm text width (it is not scaled, to keep its 8.5 pt text)}\fi
\def\FoneText{\rmfamily\fontsize{8.5}{9.5}\selectfont}%
\def\FoneMain{\rmfamily\fontsize{9}{10.5}\selectfont}%
\def\FoneBold{\rmfamily\bfseries\fontsize{9}{10.5}\selectfont}%
\def\FoneTitle{\rmfamily\bfseries\fontsize{10}{12}\selectfont}%
\ifdefined\FoneDropXbx\def\FoneXbx{}\else\def\FoneXbx{\\$X\to b\bar b$, $X\to bs$}\fi
\begin{tikzpicture}[
  x=1cm, y=1cm,
  font=\FoneMain,
  ttl/.style={anchor=base west, font=\FoneTitle, inner sep=0pt},
  cls/.style={draw=black!55, line width=0.45pt, rounded corners=1.2pt},
  hdr/.style={draw=#1!70!black, fill=#1!22, line width=0.6pt, rounded corners=2.5pt, align=center,
              font=\FoneText, inner sep=0pt},
  inbox/.style={align=center, font=\FoneText, inner sep=0pt},
  letter/.style={font=\FoneBold, inner sep=0pt},
  rowlab/.style={anchor=east, inner sep=0pt},
  txt/.style={align=left, font=\FoneText, inner sep=0pt},
  sq/.style={draw=black!65, line width=0.4pt},
  chip/.style={draw=#1!70!black, fill=#1!22, line width=0.6pt, rounded corners=3pt,
               align=center, font=\FoneText, inner sep=0pt, anchor=north west},
  stage/.style={draw=black!45, fill=black!4, line width=0.6pt, rounded corners=3pt, align=left,
                font=\FoneText, inner sep=0.1cm},
  flow/.style={-{Latex[length=1.8mm, width=1.4mm]}, line width=0.7pt, black!75},
]
% ================================================================ (a) the class hierarchy
\node[ttl] at (0,0) {(a) Class hierarchy};
% faint bands on every other row, drawn first so the boxes sit on top
\foreach \rowtop in {-1.53,-2.23,-2.93,-3.63,-4.33,-5.17,-5.87}
  \fill[black!4] (0,\rowtop) rectangle (9.70,\rowtop-0.35);
% column headers
\node[rowlab, align=right, font=\FoneText] at (1.65,-0.68) {JetClass-II\\class};
\node[hdr=c188, minimum width=1.45cm, minimum height=0.80cm, anchor=north west] at (1.75,-0.28) {\Kfull\\classes};
\node[hdr=c162, minimum width=1.45cm, minimum height=0.80cm, anchor=north west] at (3.28,-0.28) {\Kres\\classes};
\node[hdr=c43,  minimum width=1.45cm, minimum height=0.80cm, anchor=north west] at (4.81,-0.28) {\Kmid\\classes};
\node[hdr=c17,  minimum width=1.45cm, minimum height=0.80cm, anchor=north west] at (6.34,-0.28) {\Kcoarse\\classes};
% row labels
\node[rowlab] at (1.65,-1.355) {$X\to b\bar b$};
\node[rowlab] at (1.65,-1.705) {$X\to bc$};
\node[rowlab] at (1.65,-2.055) {$X\to bq$};
\node[rowlab] at (1.65,-2.405) {$X\to c\bar c$};
\node[rowlab] at (1.65,-2.755) {$X\to cs$};
\node[rowlab] at (1.65,-3.105) {$X\to cq$};
\node[rowlab] at (1.65,-3.455) {$X\to s\bar s$};
\node[rowlab] at (1.65,-3.805) {$X\to sq$};
\node[rowlab] at (1.65,-4.155) {$X\to q\bar q$};
\node[rowlab] at (1.65,-4.505) {$X\to gg$};
\node[rowlab] at (1.65,-4.995) {$X\to ee$};
\node[rowlab] at (1.65,-5.345) {$X\to\mu\mu$};
\node[rowlab] at (1.65,-5.695) {$X\to\tau_h\tau_e$};
\node[rowlab] at (1.65,-6.045) {$X\to\tau_h\tau_\mu$};
\node[rowlab] at (1.65,-6.710) {QCD jets};
% 188 and 162 classes: every resonant decay is its own class
\foreach \rowtop in {-1.18,-1.53,-1.88,-2.23,-2.58,-2.93,-3.28,-3.63,-3.98,-4.33,-4.82,-5.17,-5.52,-5.87} {
  \draw[cls, fill=c188!15] (1.75,\rowtop-0.025) rectangle (3.20,\rowtop-0.325);
  \draw[cls, fill=c162!15] (3.28,\rowtop-0.025) rectangle (4.73,\rowtop-0.325);
}
\draw[cls, fill=c188!15] (1.75,-6.385) rectangle (3.20,-7.035);
\node[inbox] at (2.475,-6.71) {\LitJetClassTwoNQcdClasses\\classes};
\draw[cls, fill=c162!15] (3.28,-6.385) rectangle (4.73,-7.035);
\node[inbox] at (4.005,-6.71) {1 class};
% 43 classes: decay topology x {b, else c, else light quark or gluon}; e and mu merged
\draw[cls, fill=c43!15] (4.81,-1.205) rectangle (6.26,-2.205);
\node[inbox] at (5.535,-1.705) {with $b$};
\draw[cls, fill=c43!15] (4.81,-2.255) rectangle (6.26,-3.255);
\node[inbox] at (5.535,-2.755) {with $c$,\\no $b$};
\draw[cls, fill=c43!15] (4.81,-3.305) rectangle (6.26,-4.655);
\node[inbox] at (5.535,-3.980) {$u$, $d$, $s$\\or $g$};
\draw[cls, fill=c43!15] (4.81,-4.845) rectangle (6.26,-5.495);
\node[inbox] at (5.535,-5.170) {$\ell\ell$};
\draw[cls, fill=c43!15] (4.81,-5.545) rectangle (6.26,-6.195);
\node[inbox] at (5.535,-5.870) {$\tau_h\tau_\ell$};
\draw[cls, fill=c43!15] (4.81,-6.385) rectangle (6.26,-7.035);
\node[inbox] at (5.535,-6.710) {1 class};
% 17 classes: number and type of prongs, no flavor
\draw[cls, fill=c17!15] (6.34,-1.205) rectangle (7.79,-4.655);
\node[inbox] at (7.065,-2.930) {two\\quarks\\or gluons};
\draw[cls, fill=c17!15] (6.34,-4.845) rectangle (7.79,-5.495);
\node[inbox] at (7.065,-5.170) {$\ell\ell$};
\draw[cls, fill=c17!15] (6.34,-5.545) rectangle (7.79,-6.195);
\node[inbox] at (7.065,-5.870) {$\tau_h\tau_\ell$};
\draw[cls, fill=c17!15] (6.34,-6.385) rectangle (7.79,-7.035);
\node[inbox] at (7.065,-6.710) {1 class};
% what each merge removes
\node[anchor=north west, font=\FoneBold, inner sep=0pt] at (0,-7.27) {What each merge removes};
\node[txt, anchor=north west, text width=6.04cm] (m1) at (1.75,-7.65)
  {the $b$, $c$ and $s$ content of the QCD jets};
\node[txt, anchor=north west, text width=6.04cm] (m2) at ([yshift=-0.10cm]m1.south west)
  {the number of $b$ and $c$ quarks, light-quark flavor, quark against gluon, and lepton
   flavor ($e$ against $\mu$)};
\node[txt, anchor=north west, text width=6.04cm] (m3) at ([yshift=-0.10cm]m2.south west)
  {quark flavor ($b$, $c$ or light); the number and type of prongs remain};
\node[txt, anchor=north east] at ([xshift=-0.10cm]m1.north west) {\Kfull\,$\to$\,\Kres};
\node[txt, anchor=north east] at ([xshift=-0.10cm]m2.north west) {\Kres\,$\to$\,\Kmid};
\node[txt, anchor=north east] at ([xshift=-0.10cm]m3.north west) {\Kmid\,$\to$\,\Kcoarse};
% ================================================================ (b) the controls
\node[ttl] at (8.25,0) {(b) Controls, \Kcoarse{} classes each};
% the first random partition on the rows of (a); one letter = one class
\node[hdr=crand, minimum width=1.45cm, minimum height=0.80cm, anchor=north west] at (8.25,-0.28) {random\\partition};
\draw[cls, fill=crand!15] (8.25,-1.205) rectangle (9.70,-2.905);
\node[letter] at (8.975,-2.055) {A};
\draw[cls, fill=crand!15] (8.25,-2.955) rectangle (9.70,-3.955);
\node[letter] at (8.975,-3.455) {B};
\draw[cls, fill=crand!15] (8.25,-4.005) rectangle (9.70,-4.305);
\node[letter] at (8.975,-4.155) {A};
\draw[cls, fill=crand!15] (8.25,-4.355) rectangle (9.70,-4.655);
\node[letter] at (8.975,-4.505) {C};
\draw[cls, fill=crand!15] (8.25,-4.845) rectangle (9.70,-5.495);
\node[letter] at (8.975,-5.170) {D};
\draw[cls, fill=crand!15] (8.25,-5.545) rectangle (9.70,-5.845);
\node[letter] at (8.975,-5.695) {E};
\draw[cls, fill=crand!15] (8.25,-5.895) rectangle (9.70,-6.195);
\node[letter] at (8.975,-6.045) {F};
\draw[cls, fill=crand!15] (8.25,-6.385) rectangle (9.70,-7.035);
\node[inbox] at (8.975,-6.710) {1 class};
% how the random partitions are built
\node[hdr=crand, anchor=north west, text width=4.75cm, inner sep=0.1cm, align=left, font=\FoneBold]
  (rh) at (10.05,-0.28) {Random partitions (\RandVtwoNPartitions)};
\node[txt, anchor=north west, text width=4.95cm] (rb) at ([yshift=-0.12cm]rh.south west)
  {\textbullet\ groups of resonant classes drawn at random, plus one class of QCD jets\\[2pt]
   \textbullet\ group shares of training jets within \RandVtwoShareTol{} of the \Kcoarse-class groups\\[2pt]
   \textbullet\ best of \RandVtwoTrials{} draws: the fewest merges shared with the \Kcoarse-class level\\[2pt]
   \textbullet\ left: the first partition; boxes with the same letter are one class};
% how the random partitions with controlled flavor boundaries are built
\node[hdr=cflav, anchor=north west, text width=4.75cm, inner sep=0.1cm, align=left, font=\FoneBold]
  (fh) at ([yshift=-0.22cm]rb.south west) {Random partitions with controlled flavor boundaries};
\node[txt, anchor=north west, text width=4.95cm] (fn) at ([yshift=-0.12cm]fh.south west)
  {whole families of decays that differ only in quark flavor; two of them split one family,
   the \RandFlavSplitOrbitClasses{} decays \RandFlavSplitOrbit:};
\begin{scope}[shift={(fn.south west)}]
  % brackets over the decays with and without a b quark (squares at a 0.2 cm pitch)
  \node[txt, anchor=south] at (1.075,-0.45) {with $b$};
  \node[txt, anchor=south] at (3.425,-0.45) {without $b$};
  \draw[line width=0.5pt] (0,-0.54) -- (0,-0.47) -- (2.15,-0.47) -- (2.15,-0.54);
  \draw[line width=0.5pt] (2.35,-0.54) -- (2.35,-0.47) -- (4.50,-0.47) -- (4.50,-0.54);
  % no flavor boundary: all 22 in one class
  \node[txt, anchor=north west] at (0,-0.62) {no flavor boundary};
  \foreach \i in {0,...,10} {
    \draw[sq, fill=white] (0.2*\i,-1.10) rectangle ++(0.15,0.15);
    \draw[sq, fill=white] (2.35+0.2*\i,-1.10) rectangle ++(0.15,0.15);
  }
  % b-quark boundary: the decays with a b quark moved to another class
  \node[txt, anchor=north west] at (0,-1.22) {$b$-quark boundary};
  \foreach \i in {0,...,10} {
    \draw[sq, fill=cflav!70!black] (0.2*\i,-1.70) rectangle ++(0.15,0.15);
    \draw[sq, fill=white] (2.35+0.2*\i,-1.70) rectangle ++(0.15,0.15);
  }
  % random flavor boundary: the same number of decays, drawn at random
  \node[txt, anchor=north west] at (0,-1.82) {random flavor boundary};
  \foreach \i in {0,1,6,9,10}
    \draw[sq, fill=white] (0.2*\i,-2.30) rectangle ++(0.15,0.15);
  \foreach \i in {2,3,4,5,7,8}
    \draw[sq, fill=cflav!70!black] (0.2*\i,-2.30) rectangle ++(0.15,0.15);
  \foreach \i in {0,4,5,7,8,10}
    \draw[sq, fill=white] (2.35+0.2*\i,-2.30) rectangle ++(0.15,0.15);
  \foreach \i in {1,2,3,6,9}
    \draw[sq, fill=cflav!70!black] (2.35+0.2*\i,-2.30) rectangle ++(0.15,0.15);
  % key
  \draw[sq, fill=cflav!70!black] (0,-2.62) rectangle ++(0.15,0.15);
  \node[txt, anchor=west] (fkey) at (0.25,-2.545) {moved to another class};
\end{scope}
% ================================================================ (c) the pipeline
% (c) starts below the lower of the two columns above, whatever their text heights
\pgfpointanchor{m3}{south}\pgfgetlastxy{\FoneXa}{\FoneYa}%
\pgfpointanchor{fkey}{south}\pgfgetlastxy{\FoneXb}{\FoneYb}%
\pgfmathsetlengthmacro{\FoneYc}{min(\FoneYa,\FoneYb)-0.45cm}%
\coordinate (ctop) at (0pt,\FoneYc);
\begin{scope}[shift={(ctop)}]
  \node[ttl, anchor=north west] at (0,0) {(c) Pretraining and evaluation};
  % what every model of one seed shares
  \node[stage, text width=2.15cm] (seed) at (1.175,-1.96)
    {\textbf{Same for every model of one seed:} the training jets, their order and the
     initial weights, except the output layer};
  % one pretraining per label set: the configuration colours of every later figure
  \node[draw=black!45, fill=black!4, line width=0.6pt, rounded corners=3pt, minimum width=5.15cm,
        minimum height=3.08cm, anchor=north west, inner sep=0pt] (pre) at (2.81,-0.42) {};
  \node[anchor=north west, font=\FoneBold, inner sep=0pt] at (2.96,-0.52) {Pretraining on JetClass-II};
  \node[chip=c188, minimum width=1.13cm, minimum height=0.42cm] at (2.96,-0.95) {\Kfull};
  \node[chip=c162, minimum width=1.13cm, minimum height=0.42cm] at (4.20,-0.95) {\Kres};
  \node[chip=c43,  minimum width=1.13cm, minimum height=0.42cm] at (5.44,-0.95) {\Kmid};
  \node[chip=c17,  minimum width=1.13cm, minimum height=0.42cm] at (6.68,-0.95) {\Kcoarse};
  \node[chip=crand, minimum width=1.70cm, minimum height=0.70cm] at (2.96,-1.47) {random\\partitions};
  \node[chip=cflav, minimum width=3.04cm, minimum height=0.70cm] at (4.77,-1.47) {controlled flavor\\boundaries};
  \node[chip=c162, densely dotted, minimum width=1.13cm, minimum height=0.42cm] at (2.96,-2.32) {\Kres};
  \node[chip=c17,  densely dotted, minimum width=1.13cm, minimum height=0.42cm] at (4.20,-2.32) {\Kcoarse};
  \node[txt, anchor=west] at (5.44,-2.53) {with auxiliary\\mass regression};
  \node[chip=c188, densely dotted, minimum width=1.13cm, minimum height=0.42cm] at (2.96,-2.94) {\Kfull};
  \node[txt, anchor=west] at (4.20,-3.15) {kinematic inputs only};
  % evaluation
  \node[stage, text width=2.4cm] (fz) at (9.68,-0.95) {frozen features:\\linear probe,\\MLP probe};
  \node[stage, text width=2.4cm] (ft) at (9.68,-2.05) {fine-tuning\\of all weights};
  \node[stage, text width=2.4cm, draw=black, fill=white, line width=1.1pt] (sc) at (9.68,-3.15) {trained from\\scratch};
  \node[stage, text width=3.4cm, minimum height=3.2cm] (task) at (13.20,-2.00)
    {\textbf{Downstream tasks}\\JetClass-II: all classes\\JetClass-II: class pairs\\\quad(both class-balanced)\\JetClass\\
     top tagging\\quark/gluon, Pythia\\quark/gluon, Herwig\\jet mass\FoneXbx};
  \draw[flow] (seed.east) -- (seed.east -| pre.west);
  \draw[flow] (fz.west -| pre.east) -- (fz.west);
  \draw[flow] (ft.west -| pre.east) -- (ft.west);
  \draw[flow] (fz.east) -- (fz.east -| task.west);
  \draw[flow] (ft.east) -- (ft.east -| task.west);
  \draw[flow] (sc.east) -- (sc.east -| task.west);
\end{scope}
\end{tikzpicture}
\caption{Design of the study.
(a)~The class hierarchy on the decays of a heavy resonance~$X$ and on the jets from quantum
chromodynamics (QCD); $q$ is a $u$ or $d$ quark, $\ell$ an electron or a muon, and $\tau_h$,
$\tau_e$, $\tau_\mu$ a $\tau$ lepton decaying to hadrons, an electron or a muon. Each box is one
class of the label set above its column, except the QCD box of the \Kfull-class column (Sophon's
label set~\cite{sophon}), which stands for \LitJetClassTwoNQcdClasses{} classes. Each level merges
classes of the level to its left; the lines below name what each merge removes.
(b)~The controls: the first random partition, where boxes with the same letter are one class, and
the random partitions with controlled flavor boundaries on the decays \RandFlavSplitOrbit{} (one
square per decay; filled squares are moved to a second class).
(c)~Models of one seed share the training jets, their order and the initial weights, except the
output layer, the mass regression target, and the first layer and initial weights of the model with
kinematic inputs only (the particle four-momenta). The hierarchy levels have \MlstNRunsWord{} seeds
each, the models with the auxiliary mass regression \MlstMassNRunsWord{}, the random partitions with
controlled flavor boundaries \RandFlavRuns{} and each random partition \RandVtwoRuns{}. Frozen
features, the output of the backbone (the network below the classification head) with its weights held
fixed, feed linear probes (logistic regressions) and multilayer perceptron (MLP) probes; fine-tuning
trains all weights. Colors identify the label sets and controls in every figure; the plots give the three partitions with
controlled flavor boundaries three shades; dotted borders mark the two
variants, which the other figures draw as dotted lines in the color of their label set.}
\label{fig:design}
\end{figure}
 from main.tex; this file holds one figure environment and nothing else.
%
% WHAT main.tex MUST LOAD (this file loads nothing itself):
%   \usepackage{tikz}
%   \usetikzlibrary{arrows.meta}          % the only TikZ library used (arrow tips)
%   \definecolor{c188}{HTML}{1F77B4}      % the configuration colours of experiments/FIGS/
%   \definecolor{c162}{HTML}{FF7F0E}      % mlst_style.py (matplotlib's tab10 in order, as in
%   \definecolor{c43}{HTML}{2CA02C}       % Sophon's figures); the variants are drawn in their
%   \definecolor{c17}{HTML}{D62728}       % label set's colour with a dotted border, as the
%   \definecolor{crand}{HTML}{17BECF}     % plotted figures draw variants dotted
%   \definecolor{cflav}{HTML}{9467BD}
%   % GENERATED by experiments/FIGS/make_tables.py -- do not edit.
% Every line carries: source file :: path inside that file :: first 16 hex of its
% sha256. paper/journal/provenance.json carries the full hash beside the value.
% Regenerate with:  python3 experiments/FIGS/make_tables.py
% Check in CI with: python3 experiments/FIGS/make_tables.py --check
%
% Macro names spell digits as words (162 -> Onesixtwo) because \newcommand takes
% letters only. A number whose source file does not exist yet is NOT emitted --
% see the missing list below, and leave a \pending{} in the text for it.
%
% MISSING INPUTS -- no macro was emitted for these:
%   community benchmarks, bench v3 readout -- experiments/FIGS/data/bench_v3_metrics*/bench_metrics_last.json
%   v2 pretraining rerun: data order and dropout fixed by seed and epoch, identical initial trunk weights, each epoch reading one of five disjoint random parts of the rows of nearly every file, fixed validation sample, checkpoint at the first maximum within the last ten epochs -- experiments/FIGS/data/v2/pretraining/
%   v2 frozen probes on the rerun models, at the first maximum within the last ten epochs and on the model whose weights are the average of the last ten epochs, BatchNorm statistics recomputed, with untrained-network and pooled-embedding reference rows -- experiments/FIGS/data/v2/probe_ladder/
%   v2 paired ratios with run and test-sample errors on the rerun models -- experiments/FIGS/data/v2/paired_errors/
%   v2 the 64- and 30-class vocabularies -- experiments/FIGS/data/v2/levels_64_30/
%   v2 the leave-one-family-out vocabularies -- experiments/FIGS/data/v2/leave_one_family_out/
%   v2 label-tree recovery on the rerun models -- experiments/FIGS/data/v2/label_recovery/
%   v2 the rerun's random partitions and flavour pair -- experiments/FIGS/data/v2/random_partitions/
%   v2 JetClass-II fine-tuning on files held out from pretraining -- experiments/FIGS/data/v2/finetune/
%   v2 fine-tuning from random initialisation and from the self-supervised model -- experiments/FIGS/data/v2/finetune_references/
%   v2 top tagging and quark/gluon tagging (Pythia and Herwig) fine-tuned from every rerun model -- experiments/FIGS/data/v2/benchmarks/
%   v2 the second and third self-supervised pretraining runs -- experiments/FIGS/data/v2/self_supervised/
%   v2 the mass-output models with the loss weight matched in loss share -- experiments/FIGS/data/v2/mass_lambda_matched/
%   v2 anomaly detection on the rerun models, per-run values -- experiments/FIGS/data/v2/anomaly/
%   v2 the open-data fits on the rerun models -- experiments/FIGS/data/v2/real_data/
%   MassLossShareVtwoOnesixtwoMass -- experiments/FIGS/data/v2/mass_lambda_matched/loss_share.json :: shares.162+mass
%   MassLossShareVtwoOnesevenMass -- experiments/FIGS/data/v2/mass_lambda_matched/loss_share.json :: shares.17+mass
%   MassLossShareVtwoOnesevenMassMatched -- experiments/FIGS/data/v2/mass_lambda_matched/loss_share.json :: shares.17+mass_matched
%   DesignVtwoSecondGpu, DesignVtwoSecondGpuRuns -- the grid's job specs put every run index on one product
%   per-arm parameters and MACs -- flops_by_arm.json (see docs/RECORD.md 2.1)
%   fine-tuning read-outs -- experiments/FIGS/data/v2/ft_mlst/: jc2, jc2_qcd
%   top-tagging recipe check -- ft_mlst/topcheck_bench_metrics.json
%   X->bb and X->bs fits -- experiments/FIGS/data/v2/xbx/
%   X->bs from scratch -- experiments/FIGS/data/v2/xbx_scratch/
%   JetClass accuracy bootstrap -- experiments/FIGS/data/v2/budget/jc1_acc_boot_N10000.json
%   MlstFtJcOneErrCoarseOverFullBootEFour -- waits on the JetClass bootstrap at 1e4 (after the equal-stopping read-out)
%   MlstLpBbQgZeroAccDropFull -- waits on pair probes with the quark/gluon inputs zeroed (open item O10, not run)
%   MlstLpBbQgZeroAccDropCoarse -- waits on pair probes with the quark/gluon inputs zeroed (open item O10, not run)
%   XbsRthirtyFullOverScratch -- waits on X->bs fits and from-scratch runs
%   XbsRthirtyFullOverScratchPr -- waits on X->bs fits and from-scratch runs
%   XbsRthirtyCoarseOverFull -- waits on X->bs fits
%   XbsRthirtyCoarseOverFullPr -- waits on X->bs fits
%   XbsRthirtyNodeOverTransferFull -- waits on X->bs fits
%   XbsRthirtyNodeOverTransferFullPr -- waits on X->bs fits
%   XbbRthirtyDirectOverTransferFull -- waits on X->bb fits
%   XbbRthirtyDirectOverTransferFullPr -- waits on X->bb fits
%   XbbRthirtyTransferCoarseOverFull -- waits on X->bb fits
%   XbbRthirtyTransferCoarseOverFullPr -- waits on X->bb fits
%
% SKIPPED TABLES:
%   T4 compute/parameters: no machine-readable source in the repository; docs/RECORD.md 2.1 names /data/results/mtx/flops/flops_by_arm.json, which is not committed

\newcommand{\PendingPretrain}{\pending{pending v2: pretraining rerun: data order and dropout fixed by seed and epoch, identical initial trunk weights, each epoch reading one of five disjoint random parts of the rows of nearly every file, fixed validation sample, checkpoint at the first maximum within the last ten epochs}}  % experiments/FIGS/make_tables.py :: V2_PENDING['Pretrain'] :: sha256 10e6ec08356a4a2b :: a marker, not a number: red in the draft until the input exists
\newcommand{\PendingProbes}{\pending{pending v2: frozen probes on the rerun models, at the first maximum within the last ten epochs and on the model whose weights are the average of the last ten epochs, BatchNorm statistics recomputed, with untrained-network and pooled-embedding reference rows}}  % experiments/FIGS/make_tables.py :: V2_PENDING['Probes'] :: sha256 10e6ec08356a4a2b :: a marker, not a number: red in the draft until the input exists
\newcommand{\PendingPaired}{\pending{pending v2: paired ratios with run and test-sample errors on the rerun models}}  % experiments/FIGS/make_tables.py :: V2_PENDING['Paired'] :: sha256 10e6ec08356a4a2b :: a marker, not a number: red in the draft until the input exists
\newcommand{\PendingLevels}{\pending{pending v2: the 64- and 30-class vocabularies}}  % experiments/FIGS/make_tables.py :: V2_PENDING['Levels'] :: sha256 10e6ec08356a4a2b :: a marker, not a number: red in the draft until the input exists
\newcommand{\PendingLofo}{\pending{pending v2: the leave-one-family-out vocabularies}}  % experiments/FIGS/make_tables.py :: V2_PENDING['Lofo'] :: sha256 10e6ec08356a4a2b :: a marker, not a number: red in the draft until the input exists
\newcommand{\PendingRecovery}{\pending{pending v2: label-tree recovery on the rerun models}}  % experiments/FIGS/make_tables.py :: V2_PENDING['Recovery'] :: sha256 10e6ec08356a4a2b :: a marker, not a number: red in the draft until the input exists
\newcommand{\PendingRandom}{\pending{pending v2: the rerun's random partitions and flavour pair}}  % experiments/FIGS/make_tables.py :: V2_PENDING['Random'] :: sha256 10e6ec08356a4a2b :: a marker, not a number: red in the draft until the input exists
\newcommand{\PendingFtHeldout}{\pending{pending v2: JetClass-II fine-tuning on files held out from pretraining}}  % experiments/FIGS/make_tables.py :: V2_PENDING['FtHeldout'] :: sha256 10e6ec08356a4a2b :: a marker, not a number: red in the draft until the input exists
\newcommand{\PendingFtRefs}{\pending{pending v2: fine-tuning from random initialisation and from the self-supervised model}}  % experiments/FIGS/make_tables.py :: V2_PENDING['FtRefs'] :: sha256 10e6ec08356a4a2b :: a marker, not a number: red in the draft until the input exists
\newcommand{\PendingBench}{\pending{pending v2: top tagging and quark/gluon tagging (Pythia and Herwig) fine-tuned from every rerun model}}  % experiments/FIGS/make_tables.py :: V2_PENDING['Bench'] :: sha256 10e6ec08356a4a2b :: a marker, not a number: red in the draft until the input exists
\newcommand{\PendingSsl}{\pending{pending v2: the second and third self-supervised pretraining runs}}  % experiments/FIGS/make_tables.py :: V2_PENDING['Ssl'] :: sha256 10e6ec08356a4a2b :: a marker, not a number: red in the draft until the input exists
\newcommand{\PendingMassLambda}{\pending{pending v2: the mass-output models with the loss weight matched in loss share}}  % experiments/FIGS/make_tables.py :: V2_PENDING['MassLambda'] :: sha256 10e6ec08356a4a2b :: a marker, not a number: red in the draft until the input exists
\newcommand{\PendingAnomaly}{\pending{pending v2: anomaly detection on the rerun models, per-run values}}  % experiments/FIGS/make_tables.py :: V2_PENDING['Anomaly'] :: sha256 10e6ec08356a4a2b :: a marker, not a number: red in the draft until the input exists
\newcommand{\PendingRealData}{\pending{pending v2: the open-data fits on the rerun models}}  % experiments/FIGS/make_tables.py :: V2_PENDING['RealData'] :: sha256 10e6ec08356a4a2b :: a marker, not a number: red in the draft until the input exists
\newcommand{\PendingConclusion}{\pending{pending v2: written once every input above is in}}  % experiments/FIGS/make_tables.py :: V2_PENDING['Conclusion'] :: sha256 10e6ec08356a4a2b :: a marker, not a number: red in the draft until the input exists
\newcommand{\AojChainDataset}{178}  % experiments/FIGS/data/aoj_checks_v2/selection_chain_v2.json :: n_dataset :: sha256 729d83085aba25f8
\newcommand{\AojChainStaged}{12{,}589{,}467}  % experiments/FIGS/data/aoj_checks_v2/selection_chain_v2.json :: n_staged :: sha256 729d83085aba25f8
\newcommand{\AojChainRhoWindow}{9{,}441{,}209}  % experiments/FIGS/data/aoj_checks_v2/selection_chain_v2.json :: n_rho_window :: sha256 729d83085aba25f8
\newcommand{\AojNJetsFit}{4{,}194{,}846}  % experiments/FIGS/data/aoj_checks_v2/selection_chain_v2.json :: n_fit :: sha256 729d83085aba25f8
\newcommand{\AojChainPtBins}{8}  % experiments/FIGS/data/aoj_checks_v2/selection_chain_v2.json :: n_pt_bins :: sha256 729d83085aba25f8
\newcommand{\AojChainEtaMax}{2.40}  % experiments/FIGS/data/aoj_checks_v2/selection_chain_v2.json :: abs_eta_max :: sha256 729d83085aba25f8
\newcommand{\MassLossShareVtwoOnesixtwoMass}{\pending{pending: realised mass-loss share in the rerun, 162 classes}}  % experiments/FIGS/make_tables.py :: SLOTS[MassLossShareVtwoOnesixtwoMass] :: sha256 10e6ec08356a4a2b :: a marker, not a number: experiments/FIGS/data/v2/mass_lambda_matched/loss_share.json :: shares.162+mass does not exist yet
\newcommand{\MassLossShareVtwoOnesevenMass}{\pending{pending: realised mass-loss share in the rerun, 17 classes}}  % experiments/FIGS/make_tables.py :: SLOTS[MassLossShareVtwoOnesevenMass] :: sha256 10e6ec08356a4a2b :: a marker, not a number: experiments/FIGS/data/v2/mass_lambda_matched/loss_share.json :: shares.17+mass does not exist yet
\newcommand{\MassLossShareVtwoOnesevenMassMatched}{\pending{pending: realised mass-loss share in the rerun, 17 classes, matched weight}}  % experiments/FIGS/make_tables.py :: SLOTS[MassLossShareVtwoOnesevenMassMatched] :: sha256 10e6ec08356a4a2b :: a marker, not a number: experiments/FIGS/data/v2/mass_lambda_matched/loss_share.json :: shares.17+mass_matched does not exist yet
\newcommand{\FtCoverageEThreeNClasses}{160}  % experiments/FT/data/jc2_v2_manifest.json :: class_coverage['train_N1000_s1.parquet'].n_classes_present :: sha256 4bd68655919b14e9
\newcommand{\FtCoverageEFourNClasses}{183}  % experiments/FT/data/jc2_v2_manifest.json :: class_coverage['train_N10000_s1.parquet'].n_classes_present :: sha256 4bd68655919b14e9
\newcommand{\FtCoverageEThreeMedian}{4}  % experiments/FT/data/jc2_v2_manifest.json :: class_coverage['train_N1000_s1.parquet'].median_per_present_class :: sha256 4bd68655919b14e9
\newcommand{\FtCoverageEThreeOneJet}{23}  % experiments/FT/data/jc2_v2_manifest.json :: class_coverage['train_N1000_s1.parquet'].n_classes_with_one_jet :: sha256 4bd68655919b14e9
\newcommand{\FtCoverageVoneEThreeNClasses}{161}  % experiments/FT/data/jc2_v1_coverage.json :: class_coverage['train_N1000_s1.parquet'].n_classes_present :: sha256 82652f57f708dc23
\newcommand{\FtCoverageVoneEThreeMedian}{5}  % experiments/FT/data/jc2_v1_coverage.json :: class_coverage['train_N1000_s1.parquet'].median_per_present_class :: sha256 82652f57f708dc23
\newcommand{\FtCoverageVoneEThreeOneJet}{25}  % experiments/FT/data/jc2_v1_coverage.json :: class_coverage['train_N1000_s1.parquet'].n_classes_with_one_jet :: sha256 82652f57f708dc23
\newcommand{\ProbeNJets}{2{,}000{,}000}  % experiments/FIGS/data/probe_ladder_v2_mlp2/analysis/seed_level_results.json :: provenance.n_jets_total :: sha256 53e1ec81343ac2c1
\newcommand{\ProbeNSeeds}{5}  % experiments/FIGS/data/probe_ladder_v2_mlp2/analysis/seed_level_results.json :: seeds_used (length) :: sha256 53e1ec81343ac2c1
\newcommand{\ProbeNSeedsWord}{five}  % experiments/FIGS/data/probe_ladder_v2_mlp2/analysis/seed_level_results.json :: seeds_used (length) :: sha256 53e1ec81343ac2c1 :: pretraining runs per vocabulary, as a word for the prose
\newcommand{\ProbeRowAlign}{\texttt{a9449fd420c473f6}}  % experiments/FIGS/data/probe_ladder_v2_mlp2/analysis/seed_level_results.json :: provenance.row_alignment_sha256 (first 16) :: sha256 53e1ec81343ac2c1
\newcommand{\ProbeEpsS}{90\%}  % experiments/FIGS/data/probe_ladder_v2_mlp2/analysis/seed_level_results.json :: levels.bvc_4prong.linear[0].headline_eps_s :: sha256 53e1ec81343ac2c1
\newcommand{\ProbeAucBvcFourprongLinearOneeighteight}{0.99320}  % experiments/FIGS/data/probe_ladder_v2_mlp2/analysis/seed_level_results.json :: levels.bvc_4prong.linear[0].mean_auc :: sha256 53e1ec81343ac2c1 :: 5 seeds, 0 saturated at AUC=1
\newcommand{\ProbeAucSdBvcFourprongLinearOneeighteight}{0.00027}  % experiments/FIGS/data/probe_ladder_v2_mlp2/analysis/seed_level_results.json :: table[bvc_4prong,linear,188].auc (SD over seeds, ddof=1) :: sha256 53e1ec81343ac2c1
\newcommand{\ProbeLogOneMinusAucBvcFourprongLinearOneeighteight}{\ensuremath{-}4.9915}  % experiments/FIGS/data/probe_ladder_v2_mlp2/analysis/seed_level_results.json :: levels.bvc_4prong.linear[0].mean :: sha256 53e1ec81343ac2c1 :: natural log of 1-AUC, lower is better
\newcommand{\ProbeLogOneMinusAucSdBvcFourprongLinearOneeighteight}{0.0389}  % experiments/FIGS/data/probe_ladder_v2_mlp2/analysis/seed_level_results.json :: levels.bvc_4prong.linear[0].seed_sd :: sha256 53e1ec81343ac2c1
\newcommand{\ProbeOmaBvcFourprongLinearOneeighteight}{\ensuremath{(6.80 \pm 0.27)\times10^{-3}}}  % experiments/FIGS/data/probe_ladder_v2_mlp2/analysis/seed_level_results.json :: table[bvc_4prong,linear,188].auc :: sha256 53e1ec81343ac2c1 :: 1-AUC, mean +- SD over 5 seeds
\newcommand{\ProbeBkgLeftBvcFourprongLinearOneeighteight}{14.0}  % experiments/FIGS/data/probe_ladder_v2_mlp2/analysis/seed_level_results.json :: levels.bvc_4prong.linear[0].rejection_points['0.90'].mean_n_bkg_pass :: sha256 53e1ec81343ac2c1 :: background jets passing the cut, mean over seeds
\newcommand{\ProbeRejBvcFourprongLinearOneeighteight}{227\,\ensuremath{\pm}\,23}  % experiments/FIGS/data/probe_ladder_v2_mlp2/analysis/seed_level_results.json :: table[bvc_4prong,linear,188].rejection_points['0.90'].rejection :: sha256 53e1ec81343ac2c1 :: at 90% signal efficiency, mean +- SD over seeds; 0 of 5 seeds at the cap
\newcommand{\ProbeAucBvcFourprongLinearOnesixtwo}{0.99311}  % experiments/FIGS/data/probe_ladder_v2_mlp2/analysis/seed_level_results.json :: levels.bvc_4prong.linear[1].mean_auc :: sha256 53e1ec81343ac2c1 :: 5 seeds, 0 saturated at AUC=1
\newcommand{\ProbeAucSdBvcFourprongLinearOnesixtwo}{0.00018}  % experiments/FIGS/data/probe_ladder_v2_mlp2/analysis/seed_level_results.json :: table[bvc_4prong,linear,162].auc (SD over seeds, ddof=1) :: sha256 53e1ec81343ac2c1
\newcommand{\ProbeLogOneMinusAucBvcFourprongLinearOnesixtwo}{\ensuremath{-}4.9777}  % experiments/FIGS/data/probe_ladder_v2_mlp2/analysis/seed_level_results.json :: levels.bvc_4prong.linear[1].mean :: sha256 53e1ec81343ac2c1 :: natural log of 1-AUC, lower is better
\newcommand{\ProbeLogOneMinusAucSdBvcFourprongLinearOnesixtwo}{0.0264}  % experiments/FIGS/data/probe_ladder_v2_mlp2/analysis/seed_level_results.json :: levels.bvc_4prong.linear[1].seed_sd :: sha256 53e1ec81343ac2c1
\newcommand{\ProbeOmaBvcFourprongLinearOnesixtwo}{\ensuremath{(6.89 \pm 0.18)\times10^{-3}}}  % experiments/FIGS/data/probe_ladder_v2_mlp2/analysis/seed_level_results.json :: table[bvc_4prong,linear,162].auc :: sha256 53e1ec81343ac2c1 :: 1-AUC, mean +- SD over 5 seeds
\newcommand{\ProbeBkgLeftBvcFourprongLinearOnesixtwo}{14.4}  % experiments/FIGS/data/probe_ladder_v2_mlp2/analysis/seed_level_results.json :: levels.bvc_4prong.linear[1].rejection_points['0.90'].mean_n_bkg_pass :: sha256 53e1ec81343ac2c1 :: background jets passing the cut, mean over seeds
\newcommand{\ProbeRejBvcFourprongLinearOnesixtwo}{220\,\ensuremath{\pm}\,40}  % experiments/FIGS/data/probe_ladder_v2_mlp2/analysis/seed_level_results.json :: table[bvc_4prong,linear,162].rejection_points['0.90'].rejection :: sha256 53e1ec81343ac2c1 :: at 90% signal efficiency, mean +- SD over seeds; 0 of 5 seeds at the cap
\newcommand{\ProbeAucBvcFourprongLinearFourthree}{0.99250}  % experiments/FIGS/data/probe_ladder_v2_mlp2/analysis/seed_level_results.json :: levels.bvc_4prong.linear[2].mean_auc :: sha256 53e1ec81343ac2c1 :: 5 seeds, 0 saturated at AUC=1
\newcommand{\ProbeAucSdBvcFourprongLinearFourthree}{0.00049}  % experiments/FIGS/data/probe_ladder_v2_mlp2/analysis/seed_level_results.json :: table[bvc_4prong,linear,43].auc (SD over seeds, ddof=1) :: sha256 53e1ec81343ac2c1
\newcommand{\ProbeLogOneMinusAucBvcFourprongLinearFourthree}{\ensuremath{-}4.8946}  % experiments/FIGS/data/probe_ladder_v2_mlp2/analysis/seed_level_results.json :: levels.bvc_4prong.linear[2].mean :: sha256 53e1ec81343ac2c1 :: natural log of 1-AUC, lower is better
\newcommand{\ProbeLogOneMinusAucSdBvcFourprongLinearFourthree}{0.0646}  % experiments/FIGS/data/probe_ladder_v2_mlp2/analysis/seed_level_results.json :: levels.bvc_4prong.linear[2].seed_sd :: sha256 53e1ec81343ac2c1
\newcommand{\ProbeOmaBvcFourprongLinearFourthree}{\ensuremath{(7.5 \pm 0.5)\times10^{-3}}}  % experiments/FIGS/data/probe_ladder_v2_mlp2/analysis/seed_level_results.json :: table[bvc_4prong,linear,43].auc :: sha256 53e1ec81343ac2c1 :: 1-AUC, mean +- SD over 5 seeds
\newcommand{\ProbeBkgLeftBvcFourprongLinearFourthree}{16.6}  % experiments/FIGS/data/probe_ladder_v2_mlp2/analysis/seed_level_results.json :: levels.bvc_4prong.linear[2].rejection_points['0.90'].mean_n_bkg_pass :: sha256 53e1ec81343ac2c1 :: background jets passing the cut, mean over seeds
\newcommand{\ProbeRejBvcFourprongLinearFourthree}{192\,\ensuremath{\pm}\,20}  % experiments/FIGS/data/probe_ladder_v2_mlp2/analysis/seed_level_results.json :: table[bvc_4prong,linear,43].rejection_points['0.90'].rejection :: sha256 53e1ec81343ac2c1 :: at 90% signal efficiency, mean +- SD over seeds; 0 of 5 seeds at the cap
\newcommand{\ProbeAucBvcFourprongLinearOneseven}{0.96936}  % experiments/FIGS/data/probe_ladder_v2_mlp2/analysis/seed_level_results.json :: levels.bvc_4prong.linear[3].mean_auc :: sha256 53e1ec81343ac2c1 :: 5 seeds, 0 saturated at AUC=1
\newcommand{\ProbeAucSdBvcFourprongLinearOneseven}{0.00514}  % experiments/FIGS/data/probe_ladder_v2_mlp2/analysis/seed_level_results.json :: table[bvc_4prong,linear,17].auc (SD over seeds, ddof=1) :: sha256 53e1ec81343ac2c1
\newcommand{\ProbeLogOneMinusAucBvcFourprongLinearOneseven}{\ensuremath{-}3.4967}  % experiments/FIGS/data/probe_ladder_v2_mlp2/analysis/seed_level_results.json :: levels.bvc_4prong.linear[3].mean :: sha256 53e1ec81343ac2c1 :: natural log of 1-AUC, lower is better
\newcommand{\ProbeLogOneMinusAucSdBvcFourprongLinearOneseven}{0.1667}  % experiments/FIGS/data/probe_ladder_v2_mlp2/analysis/seed_level_results.json :: levels.bvc_4prong.linear[3].seed_sd :: sha256 53e1ec81343ac2c1
\newcommand{\ProbeOmaBvcFourprongLinearOneseven}{\ensuremath{(3.1 \pm 0.5)\times10^{-2}}}  % experiments/FIGS/data/probe_ladder_v2_mlp2/analysis/seed_level_results.json :: table[bvc_4prong,linear,17].auc :: sha256 53e1ec81343ac2c1 :: 1-AUC, mean +- SD over 5 seeds
\newcommand{\ProbeBkgLeftBvcFourprongLinearOneseven}{137.4}  % experiments/FIGS/data/probe_ladder_v2_mlp2/analysis/seed_level_results.json :: levels.bvc_4prong.linear[3].rejection_points['0.90'].mean_n_bkg_pass :: sha256 53e1ec81343ac2c1 :: background jets passing the cut, mean over seeds
\newcommand{\ProbeRejBvcFourprongLinearOneseven}{24\,\ensuremath{\pm}\,6}  % experiments/FIGS/data/probe_ladder_v2_mlp2/analysis/seed_level_results.json :: table[bvc_4prong,linear,17].rejection_points['0.90'].rejection :: sha256 53e1ec81343ac2c1 :: at 90% signal efficiency, mean +- SD over seeds; 0 of 5 seeds at the cap
\newcommand{\ProbeOmaRatioBvcFourprongLinearOnesixtwoOverOneeighteight}{1.01}  % experiments/FIGS/data/probe_ladder_v2_mlp2/analysis/seed_level_results.json :: table[bvc_4prong,linear,162/188].auc :: sha256 53e1ec81343ac2c1 :: ratio of the seed means of 1-AUC, coarser over finer
\newcommand{\ProbeOmaRatioBvcFourprongLinearFourthreeOverOneeighteight}{1.10}  % experiments/FIGS/data/probe_ladder_v2_mlp2/analysis/seed_level_results.json :: table[bvc_4prong,linear,43/188].auc :: sha256 53e1ec81343ac2c1 :: ratio of the seed means of 1-AUC, coarser over finer
\newcommand{\ProbeOmaRatioBvcFourprongLinearOnesevenOverOneeighteight}{4.5}  % experiments/FIGS/data/probe_ladder_v2_mlp2/analysis/seed_level_results.json :: table[bvc_4prong,linear,17/188].auc :: sha256 53e1ec81343ac2c1 :: ratio of the seed means of 1-AUC, coarser over finer
\newcommand{\ProbeOmaRatioBvcFourprongLinearFourthreeOverOnesixtwo}{1.09}  % experiments/FIGS/data/probe_ladder_v2_mlp2/analysis/seed_level_results.json :: table[bvc_4prong,linear,43/162].auc :: sha256 53e1ec81343ac2c1 :: ratio of the seed means of 1-AUC, coarser over finer
\newcommand{\ProbeOmaRatioBvcFourprongLinearOnesevenOverOnesixtwo}{4.4}  % experiments/FIGS/data/probe_ladder_v2_mlp2/analysis/seed_level_results.json :: table[bvc_4prong,linear,17/162].auc :: sha256 53e1ec81343ac2c1 :: ratio of the seed means of 1-AUC, coarser over finer
\newcommand{\ProbeOmaRatioBvcFourprongLinearOnesevenOverFourthree}{4.1}  % experiments/FIGS/data/probe_ladder_v2_mlp2/analysis/seed_level_results.json :: table[bvc_4prong,linear,17/43].auc :: sha256 53e1ec81343ac2c1 :: ratio of the seed means of 1-AUC, coarser over finer
\newcommand{\ProbeAucBvcFourprongMlpOneeighteight}{0.99316}  % experiments/FIGS/data/probe_ladder_v2_mlp2/analysis/seed_level_results.json :: levels.bvc_4prong.mlp[0].mean_auc :: sha256 53e1ec81343ac2c1 :: 5 seeds, 0 saturated at AUC=1
\newcommand{\ProbeAucSdBvcFourprongMlpOneeighteight}{0.00020}  % experiments/FIGS/data/probe_ladder_v2_mlp2/analysis/seed_level_results.json :: table[bvc_4prong,mlp,188].auc (SD over seeds, ddof=1) :: sha256 53e1ec81343ac2c1
\newcommand{\ProbeLogOneMinusAucBvcFourprongMlpOneeighteight}{\ensuremath{-}4.9857}  % experiments/FIGS/data/probe_ladder_v2_mlp2/analysis/seed_level_results.json :: levels.bvc_4prong.mlp[0].mean :: sha256 53e1ec81343ac2c1 :: natural log of 1-AUC, lower is better
\newcommand{\ProbeLogOneMinusAucSdBvcFourprongMlpOneeighteight}{0.0290}  % experiments/FIGS/data/probe_ladder_v2_mlp2/analysis/seed_level_results.json :: levels.bvc_4prong.mlp[0].seed_sd :: sha256 53e1ec81343ac2c1
\newcommand{\ProbeOmaBvcFourprongMlpOneeighteight}{\ensuremath{(6.84 \pm 0.20)\times10^{-3}}}  % experiments/FIGS/data/probe_ladder_v2_mlp2/analysis/seed_level_results.json :: table[bvc_4prong,mlp,188].auc :: sha256 53e1ec81343ac2c1 :: 1-AUC, mean +- SD over 5 seeds
\newcommand{\ProbeBkgLeftBvcFourprongMlpOneeighteight}{15.2}  % experiments/FIGS/data/probe_ladder_v2_mlp2/analysis/seed_level_results.json :: levels.bvc_4prong.mlp[0].rejection_points['0.90'].mean_n_bkg_pass :: sha256 53e1ec81343ac2c1 :: background jets passing the cut, mean over seeds
\newcommand{\ProbeRejBvcFourprongMlpOneeighteight}{212\,\ensuremath{\pm}\,34}  % experiments/FIGS/data/probe_ladder_v2_mlp2/analysis/seed_level_results.json :: table[bvc_4prong,mlp,188].rejection_points['0.90'].rejection :: sha256 53e1ec81343ac2c1 :: at 90% signal efficiency, mean +- SD over seeds; 0 of 5 seeds at the cap
\newcommand{\ProbeAucBvcFourprongMlpOnesixtwo}{0.99318}  % experiments/FIGS/data/probe_ladder_v2_mlp2/analysis/seed_level_results.json :: levels.bvc_4prong.mlp[1].mean_auc :: sha256 53e1ec81343ac2c1 :: 5 seeds, 0 saturated at AUC=1
\newcommand{\ProbeAucSdBvcFourprongMlpOnesixtwo}{0.00020}  % experiments/FIGS/data/probe_ladder_v2_mlp2/analysis/seed_level_results.json :: table[bvc_4prong,mlp,162].auc (SD over seeds, ddof=1) :: sha256 53e1ec81343ac2c1
\newcommand{\ProbeLogOneMinusAucBvcFourprongMlpOnesixtwo}{\ensuremath{-}4.9886}  % experiments/FIGS/data/probe_ladder_v2_mlp2/analysis/seed_level_results.json :: levels.bvc_4prong.mlp[1].mean :: sha256 53e1ec81343ac2c1 :: natural log of 1-AUC, lower is better
\newcommand{\ProbeLogOneMinusAucSdBvcFourprongMlpOnesixtwo}{0.0298}  % experiments/FIGS/data/probe_ladder_v2_mlp2/analysis/seed_level_results.json :: levels.bvc_4prong.mlp[1].seed_sd :: sha256 53e1ec81343ac2c1
\newcommand{\ProbeOmaBvcFourprongMlpOnesixtwo}{\ensuremath{(6.82 \pm 0.20)\times10^{-3}}}  % experiments/FIGS/data/probe_ladder_v2_mlp2/analysis/seed_level_results.json :: table[bvc_4prong,mlp,162].auc :: sha256 53e1ec81343ac2c1 :: 1-AUC, mean +- SD over 5 seeds
\newcommand{\ProbeBkgLeftBvcFourprongMlpOnesixtwo}{16.6}  % experiments/FIGS/data/probe_ladder_v2_mlp2/analysis/seed_level_results.json :: levels.bvc_4prong.mlp[1].rejection_points['0.90'].mean_n_bkg_pass :: sha256 53e1ec81343ac2c1 :: background jets passing the cut, mean over seeds
\newcommand{\ProbeRejBvcFourprongMlpOnesixtwo}{200\,\ensuremath{\pm}\,40}  % experiments/FIGS/data/probe_ladder_v2_mlp2/analysis/seed_level_results.json :: table[bvc_4prong,mlp,162].rejection_points['0.90'].rejection :: sha256 53e1ec81343ac2c1 :: at 90% signal efficiency, mean +- SD over seeds; 0 of 5 seeds at the cap
\newcommand{\ProbeAucBvcFourprongMlpFourthree}{0.99232}  % experiments/FIGS/data/probe_ladder_v2_mlp2/analysis/seed_level_results.json :: levels.bvc_4prong.mlp[2].mean_auc :: sha256 53e1ec81343ac2c1 :: 5 seeds, 0 saturated at AUC=1
\newcommand{\ProbeAucSdBvcFourprongMlpFourthree}{0.00058}  % experiments/FIGS/data/probe_ladder_v2_mlp2/analysis/seed_level_results.json :: table[bvc_4prong,mlp,43].auc (SD over seeds, ddof=1) :: sha256 53e1ec81343ac2c1
\newcommand{\ProbeLogOneMinusAucBvcFourprongMlpFourthree}{\ensuremath{-}4.8719}  % experiments/FIGS/data/probe_ladder_v2_mlp2/analysis/seed_level_results.json :: levels.bvc_4prong.mlp[2].mean :: sha256 53e1ec81343ac2c1 :: natural log of 1-AUC, lower is better
\newcommand{\ProbeLogOneMinusAucSdBvcFourprongMlpFourthree}{0.0734}  % experiments/FIGS/data/probe_ladder_v2_mlp2/analysis/seed_level_results.json :: levels.bvc_4prong.mlp[2].seed_sd :: sha256 53e1ec81343ac2c1
\newcommand{\ProbeOmaBvcFourprongMlpFourthree}{\ensuremath{(7.7 \pm 0.6)\times10^{-3}}}  % experiments/FIGS/data/probe_ladder_v2_mlp2/analysis/seed_level_results.json :: table[bvc_4prong,mlp,43].auc :: sha256 53e1ec81343ac2c1 :: 1-AUC, mean +- SD over 5 seeds
\newcommand{\ProbeBkgLeftBvcFourprongMlpFourthree}{17.4}  % experiments/FIGS/data/probe_ladder_v2_mlp2/analysis/seed_level_results.json :: levels.bvc_4prong.mlp[2].rejection_points['0.90'].mean_n_bkg_pass :: sha256 53e1ec81343ac2c1 :: background jets passing the cut, mean over seeds
\newcommand{\ProbeRejBvcFourprongMlpFourthree}{184\,\ensuremath{\pm}\,23}  % experiments/FIGS/data/probe_ladder_v2_mlp2/analysis/seed_level_results.json :: table[bvc_4prong,mlp,43].rejection_points['0.90'].rejection :: sha256 53e1ec81343ac2c1 :: at 90% signal efficiency, mean +- SD over seeds; 0 of 5 seeds at the cap
\newcommand{\ProbeAucBvcFourprongMlpOneseven}{0.97766}  % experiments/FIGS/data/probe_ladder_v2_mlp2/analysis/seed_level_results.json :: levels.bvc_4prong.mlp[3].mean_auc :: sha256 53e1ec81343ac2c1 :: 5 seeds, 0 saturated at AUC=1
\newcommand{\ProbeAucSdBvcFourprongMlpOneseven}{0.00310}  % experiments/FIGS/data/probe_ladder_v2_mlp2/analysis/seed_level_results.json :: table[bvc_4prong,mlp,17].auc (SD over seeds, ddof=1) :: sha256 53e1ec81343ac2c1
\newcommand{\ProbeLogOneMinusAucBvcFourprongMlpOneseven}{\ensuremath{-}3.8087}  % experiments/FIGS/data/probe_ladder_v2_mlp2/analysis/seed_level_results.json :: levels.bvc_4prong.mlp[3].mean :: sha256 53e1ec81343ac2c1 :: natural log of 1-AUC, lower is better
\newcommand{\ProbeLogOneMinusAucSdBvcFourprongMlpOneseven}{0.1366}  % experiments/FIGS/data/probe_ladder_v2_mlp2/analysis/seed_level_results.json :: levels.bvc_4prong.mlp[3].seed_sd :: sha256 53e1ec81343ac2c1
\newcommand{\ProbeOmaBvcFourprongMlpOneseven}{\ensuremath{(2.23 \pm 0.31)\times10^{-2}}}  % experiments/FIGS/data/probe_ladder_v2_mlp2/analysis/seed_level_results.json :: table[bvc_4prong,mlp,17].auc :: sha256 53e1ec81343ac2c1 :: 1-AUC, mean +- SD over 5 seeds
\newcommand{\ProbeBkgLeftBvcFourprongMlpOneseven}{99.0}  % experiments/FIGS/data/probe_ladder_v2_mlp2/analysis/seed_level_results.json :: levels.bvc_4prong.mlp[3].rejection_points['0.90'].mean_n_bkg_pass :: sha256 53e1ec81343ac2c1 :: background jets passing the cut, mean over seeds
\newcommand{\ProbeRejBvcFourprongMlpOneseven}{33\,\ensuremath{\pm}\,5}  % experiments/FIGS/data/probe_ladder_v2_mlp2/analysis/seed_level_results.json :: table[bvc_4prong,mlp,17].rejection_points['0.90'].rejection :: sha256 53e1ec81343ac2c1 :: at 90% signal efficiency, mean +- SD over seeds; 0 of 5 seeds at the cap
\newcommand{\ProbeOmaRatioBvcFourprongMlpOnesixtwoOverOneeighteight}{1.00}  % experiments/FIGS/data/probe_ladder_v2_mlp2/analysis/seed_level_results.json :: table[bvc_4prong,mlp,162/188].auc :: sha256 53e1ec81343ac2c1 :: ratio of the seed means of 1-AUC, coarser over finer
\newcommand{\ProbeOmaRatioBvcFourprongMlpFourthreeOverOneeighteight}{1.12}  % experiments/FIGS/data/probe_ladder_v2_mlp2/analysis/seed_level_results.json :: table[bvc_4prong,mlp,43/188].auc :: sha256 53e1ec81343ac2c1 :: ratio of the seed means of 1-AUC, coarser over finer
\newcommand{\ProbeOmaRatioBvcFourprongMlpOnesevenOverOneeighteight}{3.3}  % experiments/FIGS/data/probe_ladder_v2_mlp2/analysis/seed_level_results.json :: table[bvc_4prong,mlp,17/188].auc :: sha256 53e1ec81343ac2c1 :: ratio of the seed means of 1-AUC, coarser over finer
\newcommand{\ProbeOmaRatioBvcFourprongMlpFourthreeOverOnesixtwo}{1.13}  % experiments/FIGS/data/probe_ladder_v2_mlp2/analysis/seed_level_results.json :: table[bvc_4prong,mlp,43/162].auc :: sha256 53e1ec81343ac2c1 :: ratio of the seed means of 1-AUC, coarser over finer
\newcommand{\ProbeOmaRatioBvcFourprongMlpOnesevenOverOnesixtwo}{3.3}  % experiments/FIGS/data/probe_ladder_v2_mlp2/analysis/seed_level_results.json :: table[bvc_4prong,mlp,17/162].auc :: sha256 53e1ec81343ac2c1 :: ratio of the seed means of 1-AUC, coarser over finer
\newcommand{\ProbeOmaRatioBvcFourprongMlpOnesevenOverFourthree}{2.9}  % experiments/FIGS/data/probe_ladder_v2_mlp2/analysis/seed_level_results.json :: table[bvc_4prong,mlp,17/43].auc :: sha256 53e1ec81343ac2c1 :: ratio of the seed means of 1-AUC, coarser over finer
\newcommand{\ProbeAucBvcQcdLinearOneeighteight}{0.98519}  % experiments/FIGS/data/probe_ladder_v2_mlp2/analysis/seed_level_results.json :: levels.bvc_qcd.linear[0].mean_auc :: sha256 53e1ec81343ac2c1 :: 5 seeds, 0 saturated at AUC=1
\newcommand{\ProbeAucSdBvcQcdLinearOneeighteight}{0.00061}  % experiments/FIGS/data/probe_ladder_v2_mlp2/analysis/seed_level_results.json :: table[bvc_qcd,linear,188].auc (SD over seeds, ddof=1) :: sha256 53e1ec81343ac2c1
\newcommand{\ProbeLogOneMinusAucBvcQcdLinearOneeighteight}{\ensuremath{-}4.2133}  % experiments/FIGS/data/probe_ladder_v2_mlp2/analysis/seed_level_results.json :: levels.bvc_qcd.linear[0].mean :: sha256 53e1ec81343ac2c1 :: natural log of 1-AUC, lower is better
\newcommand{\ProbeLogOneMinusAucSdBvcQcdLinearOneeighteight}{0.0422}  % experiments/FIGS/data/probe_ladder_v2_mlp2/analysis/seed_level_results.json :: levels.bvc_qcd.linear[0].seed_sd :: sha256 53e1ec81343ac2c1
\newcommand{\ProbeOmaBvcQcdLinearOneeighteight}{\ensuremath{(1.48 \pm 0.06)\times10^{-2}}}  % experiments/FIGS/data/probe_ladder_v2_mlp2/analysis/seed_level_results.json :: table[bvc_qcd,linear,188].auc :: sha256 53e1ec81343ac2c1 :: 1-AUC, mean +- SD over 5 seeds
\newcommand{\ProbeBkgLeftBvcQcdLinearOneeighteight}{96.4}  % experiments/FIGS/data/probe_ladder_v2_mlp2/analysis/seed_level_results.json :: levels.bvc_qcd.linear[0].rejection_points['0.90'].mean_n_bkg_pass :: sha256 53e1ec81343ac2c1 :: background jets passing the cut, mean over seeds
\newcommand{\ProbeRejBvcQcdLinearOneeighteight}{56\,\ensuremath{\pm}\,5}  % experiments/FIGS/data/probe_ladder_v2_mlp2/analysis/seed_level_results.json :: table[bvc_qcd,linear,188].rejection_points['0.90'].rejection :: sha256 53e1ec81343ac2c1 :: at 90% signal efficiency, mean +- SD over seeds; 0 of 5 seeds at the cap
\newcommand{\ProbeAucBvcQcdLinearOnesixtwo}{0.98305}  % experiments/FIGS/data/probe_ladder_v2_mlp2/analysis/seed_level_results.json :: levels.bvc_qcd.linear[1].mean_auc :: sha256 53e1ec81343ac2c1 :: 5 seeds, 0 saturated at AUC=1
\newcommand{\ProbeAucSdBvcQcdLinearOnesixtwo}{0.00024}  % experiments/FIGS/data/probe_ladder_v2_mlp2/analysis/seed_level_results.json :: table[bvc_qcd,linear,162].auc (SD over seeds, ddof=1) :: sha256 53e1ec81343ac2c1
\newcommand{\ProbeLogOneMinusAucBvcQcdLinearOnesixtwo}{\ensuremath{-}4.0774}  % experiments/FIGS/data/probe_ladder_v2_mlp2/analysis/seed_level_results.json :: levels.bvc_qcd.linear[1].mean :: sha256 53e1ec81343ac2c1 :: natural log of 1-AUC, lower is better
\newcommand{\ProbeLogOneMinusAucSdBvcQcdLinearOnesixtwo}{0.0142}  % experiments/FIGS/data/probe_ladder_v2_mlp2/analysis/seed_level_results.json :: levels.bvc_qcd.linear[1].seed_sd :: sha256 53e1ec81343ac2c1
\newcommand{\ProbeOmaBvcQcdLinearOnesixtwo}{\ensuremath{(1.695 \pm 0.024)\times10^{-2}}}  % experiments/FIGS/data/probe_ladder_v2_mlp2/analysis/seed_level_results.json :: table[bvc_qcd,linear,162].auc :: sha256 53e1ec81343ac2c1 :: 1-AUC, mean +- SD over 5 seeds
\newcommand{\ProbeBkgLeftBvcQcdLinearOnesixtwo}{121.4}  % experiments/FIGS/data/probe_ladder_v2_mlp2/analysis/seed_level_results.json :: levels.bvc_qcd.linear[1].rejection_points['0.90'].mean_n_bkg_pass :: sha256 53e1ec81343ac2c1 :: background jets passing the cut, mean over seeds
\newcommand{\ProbeRejBvcQcdLinearOnesixtwo}{44.3\,\ensuremath{\pm}\,2.1}  % experiments/FIGS/data/probe_ladder_v2_mlp2/analysis/seed_level_results.json :: table[bvc_qcd,linear,162].rejection_points['0.90'].rejection :: sha256 53e1ec81343ac2c1 :: at 90% signal efficiency, mean +- SD over seeds; 0 of 5 seeds at the cap
\newcommand{\ProbeAucBvcQcdLinearFourthree}{0.98281}  % experiments/FIGS/data/probe_ladder_v2_mlp2/analysis/seed_level_results.json :: levels.bvc_qcd.linear[2].mean_auc :: sha256 53e1ec81343ac2c1 :: 5 seeds, 0 saturated at AUC=1
\newcommand{\ProbeAucSdBvcQcdLinearFourthree}{0.00064}  % experiments/FIGS/data/probe_ladder_v2_mlp2/analysis/seed_level_results.json :: table[bvc_qcd,linear,43].auc (SD over seeds, ddof=1) :: sha256 53e1ec81343ac2c1
\newcommand{\ProbeLogOneMinusAucBvcQcdLinearFourthree}{\ensuremath{-}4.0638}  % experiments/FIGS/data/probe_ladder_v2_mlp2/analysis/seed_level_results.json :: levels.bvc_qcd.linear[2].mean :: sha256 53e1ec81343ac2c1 :: natural log of 1-AUC, lower is better
\newcommand{\ProbeLogOneMinusAucSdBvcQcdLinearFourthree}{0.0370}  % experiments/FIGS/data/probe_ladder_v2_mlp2/analysis/seed_level_results.json :: levels.bvc_qcd.linear[2].seed_sd :: sha256 53e1ec81343ac2c1
\newcommand{\ProbeOmaBvcQcdLinearFourthree}{\ensuremath{(1.72 \pm 0.06)\times10^{-2}}}  % experiments/FIGS/data/probe_ladder_v2_mlp2/analysis/seed_level_results.json :: table[bvc_qcd,linear,43].auc :: sha256 53e1ec81343ac2c1 :: 1-AUC, mean +- SD over 5 seeds
\newcommand{\ProbeBkgLeftBvcQcdLinearFourthree}{126.8}  % experiments/FIGS/data/probe_ladder_v2_mlp2/analysis/seed_level_results.json :: levels.bvc_qcd.linear[2].rejection_points['0.90'].mean_n_bkg_pass :: sha256 53e1ec81343ac2c1 :: background jets passing the cut, mean over seeds
\newcommand{\ProbeRejBvcQcdLinearFourthree}{43\,\ensuremath{\pm}\,4}  % experiments/FIGS/data/probe_ladder_v2_mlp2/analysis/seed_level_results.json :: table[bvc_qcd,linear,43].rejection_points['0.90'].rejection :: sha256 53e1ec81343ac2c1 :: at 90% signal efficiency, mean +- SD over seeds; 0 of 5 seeds at the cap
\newcommand{\ProbeAucBvcQcdLinearOneseven}{0.95504}  % experiments/FIGS/data/probe_ladder_v2_mlp2/analysis/seed_level_results.json :: levels.bvc_qcd.linear[3].mean_auc :: sha256 53e1ec81343ac2c1 :: 5 seeds, 0 saturated at AUC=1
\newcommand{\ProbeAucSdBvcQcdLinearOneseven}{0.00527}  % experiments/FIGS/data/probe_ladder_v2_mlp2/analysis/seed_level_results.json :: table[bvc_qcd,linear,17].auc (SD over seeds, ddof=1) :: sha256 53e1ec81343ac2c1
\newcommand{\ProbeLogOneMinusAucBvcQcdLinearOneseven}{\ensuremath{-}3.1076}  % experiments/FIGS/data/probe_ladder_v2_mlp2/analysis/seed_level_results.json :: levels.bvc_qcd.linear[3].mean :: sha256 53e1ec81343ac2c1 :: natural log of 1-AUC, lower is better
\newcommand{\ProbeLogOneMinusAucSdBvcQcdLinearOneseven}{0.1198}  % experiments/FIGS/data/probe_ladder_v2_mlp2/analysis/seed_level_results.json :: levels.bvc_qcd.linear[3].seed_sd :: sha256 53e1ec81343ac2c1
\newcommand{\ProbeOmaBvcQcdLinearOneseven}{\ensuremath{(4.5 \pm 0.5)\times10^{-2}}}  % experiments/FIGS/data/probe_ladder_v2_mlp2/analysis/seed_level_results.json :: table[bvc_qcd,linear,17].auc :: sha256 53e1ec81343ac2c1 :: 1-AUC, mean +- SD over 5 seeds
\newcommand{\ProbeBkgLeftBvcQcdLinearOneseven}{544.6}  % experiments/FIGS/data/probe_ladder_v2_mlp2/analysis/seed_level_results.json :: levels.bvc_qcd.linear[3].rejection_points['0.90'].mean_n_bkg_pass :: sha256 53e1ec81343ac2c1 :: background jets passing the cut, mean over seeds
\newcommand{\ProbeRejBvcQcdLinearOneseven}{10.1\,\ensuremath{\pm}\,1.9}  % experiments/FIGS/data/probe_ladder_v2_mlp2/analysis/seed_level_results.json :: table[bvc_qcd,linear,17].rejection_points['0.90'].rejection :: sha256 53e1ec81343ac2c1 :: at 90% signal efficiency, mean +- SD over seeds; 0 of 5 seeds at the cap
\newcommand{\ProbeOmaRatioBvcQcdLinearOnesixtwoOverOneeighteight}{1.14}  % experiments/FIGS/data/probe_ladder_v2_mlp2/analysis/seed_level_results.json :: table[bvc_qcd,linear,162/188].auc :: sha256 53e1ec81343ac2c1 :: ratio of the seed means of 1-AUC, coarser over finer
\newcommand{\ProbeOmaRatioBvcQcdLinearFourthreeOverOneeighteight}{1.16}  % experiments/FIGS/data/probe_ladder_v2_mlp2/analysis/seed_level_results.json :: table[bvc_qcd,linear,43/188].auc :: sha256 53e1ec81343ac2c1 :: ratio of the seed means of 1-AUC, coarser over finer
\newcommand{\ProbeOmaRatioBvcQcdLinearOnesevenOverOneeighteight}{3.0}  % experiments/FIGS/data/probe_ladder_v2_mlp2/analysis/seed_level_results.json :: table[bvc_qcd,linear,17/188].auc :: sha256 53e1ec81343ac2c1 :: ratio of the seed means of 1-AUC, coarser over finer
\newcommand{\ProbeOmaRatioBvcQcdLinearFourthreeOverOnesixtwo}{1.01}  % experiments/FIGS/data/probe_ladder_v2_mlp2/analysis/seed_level_results.json :: table[bvc_qcd,linear,43/162].auc :: sha256 53e1ec81343ac2c1 :: ratio of the seed means of 1-AUC, coarser over finer
\newcommand{\ProbeOmaRatioBvcQcdLinearOnesevenOverOnesixtwo}{2.7}  % experiments/FIGS/data/probe_ladder_v2_mlp2/analysis/seed_level_results.json :: table[bvc_qcd,linear,17/162].auc :: sha256 53e1ec81343ac2c1 :: ratio of the seed means of 1-AUC, coarser over finer
\newcommand{\ProbeOmaRatioBvcQcdLinearOnesevenOverFourthree}{2.6}  % experiments/FIGS/data/probe_ladder_v2_mlp2/analysis/seed_level_results.json :: table[bvc_qcd,linear,17/43].auc :: sha256 53e1ec81343ac2c1 :: ratio of the seed means of 1-AUC, coarser over finer
\newcommand{\ProbeAucBvcQcdMlpOneeighteight}{0.98527}  % experiments/FIGS/data/probe_ladder_v2_mlp2/analysis/seed_level_results.json :: levels.bvc_qcd.mlp[0].mean_auc :: sha256 53e1ec81343ac2c1 :: 5 seeds, 0 saturated at AUC=1
\newcommand{\ProbeAucSdBvcQcdMlpOneeighteight}{0.00058}  % experiments/FIGS/data/probe_ladder_v2_mlp2/analysis/seed_level_results.json :: table[bvc_qcd,mlp,188].auc (SD over seeds, ddof=1) :: sha256 53e1ec81343ac2c1
\newcommand{\ProbeLogOneMinusAucBvcQcdMlpOneeighteight}{\ensuremath{-}4.2188}  % experiments/FIGS/data/probe_ladder_v2_mlp2/analysis/seed_level_results.json :: levels.bvc_qcd.mlp[0].mean :: sha256 53e1ec81343ac2c1 :: natural log of 1-AUC, lower is better
\newcommand{\ProbeLogOneMinusAucSdBvcQcdMlpOneeighteight}{0.0400}  % experiments/FIGS/data/probe_ladder_v2_mlp2/analysis/seed_level_results.json :: levels.bvc_qcd.mlp[0].seed_sd :: sha256 53e1ec81343ac2c1
\newcommand{\ProbeOmaBvcQcdMlpOneeighteight}{\ensuremath{(1.47 \pm 0.06)\times10^{-2}}}  % experiments/FIGS/data/probe_ladder_v2_mlp2/analysis/seed_level_results.json :: table[bvc_qcd,mlp,188].auc :: sha256 53e1ec81343ac2c1 :: 1-AUC, mean +- SD over 5 seeds
\newcommand{\ProbeBkgLeftBvcQcdMlpOneeighteight}{94.6}  % experiments/FIGS/data/probe_ladder_v2_mlp2/analysis/seed_level_results.json :: levels.bvc_qcd.mlp[0].rejection_points['0.90'].mean_n_bkg_pass :: sha256 53e1ec81343ac2c1 :: background jets passing the cut, mean over seeds
\newcommand{\ProbeRejBvcQcdMlpOneeighteight}{57\,\ensuremath{\pm}\,6}  % experiments/FIGS/data/probe_ladder_v2_mlp2/analysis/seed_level_results.json :: table[bvc_qcd,mlp,188].rejection_points['0.90'].rejection :: sha256 53e1ec81343ac2c1 :: at 90% signal efficiency, mean +- SD over seeds; 0 of 5 seeds at the cap
\newcommand{\ProbeAucBvcQcdMlpOnesixtwo}{0.98361}  % experiments/FIGS/data/probe_ladder_v2_mlp2/analysis/seed_level_results.json :: levels.bvc_qcd.mlp[1].mean_auc :: sha256 53e1ec81343ac2c1 :: 5 seeds, 0 saturated at AUC=1
\newcommand{\ProbeAucSdBvcQcdMlpOnesixtwo}{0.00039}  % experiments/FIGS/data/probe_ladder_v2_mlp2/analysis/seed_level_results.json :: table[bvc_qcd,mlp,162].auc (SD over seeds, ddof=1) :: sha256 53e1ec81343ac2c1
\newcommand{\ProbeLogOneMinusAucBvcQcdMlpOnesixtwo}{\ensuremath{-}4.1113}  % experiments/FIGS/data/probe_ladder_v2_mlp2/analysis/seed_level_results.json :: levels.bvc_qcd.mlp[1].mean :: sha256 53e1ec81343ac2c1 :: natural log of 1-AUC, lower is better
\newcommand{\ProbeLogOneMinusAucSdBvcQcdMlpOnesixtwo}{0.0241}  % experiments/FIGS/data/probe_ladder_v2_mlp2/analysis/seed_level_results.json :: levels.bvc_qcd.mlp[1].seed_sd :: sha256 53e1ec81343ac2c1
\newcommand{\ProbeOmaBvcQcdMlpOnesixtwo}{\ensuremath{(1.64 \pm 0.04)\times10^{-2}}}  % experiments/FIGS/data/probe_ladder_v2_mlp2/analysis/seed_level_results.json :: table[bvc_qcd,mlp,162].auc :: sha256 53e1ec81343ac2c1 :: 1-AUC, mean +- SD over 5 seeds
\newcommand{\ProbeBkgLeftBvcQcdMlpOnesixtwo}{119.6}  % experiments/FIGS/data/probe_ladder_v2_mlp2/analysis/seed_level_results.json :: levels.bvc_qcd.mlp[1].rejection_points['0.90'].mean_n_bkg_pass :: sha256 53e1ec81343ac2c1 :: background jets passing the cut, mean over seeds
\newcommand{\ProbeRejBvcQcdMlpOnesixtwo}{45.1\,\ensuremath{\pm}\,3.2}  % experiments/FIGS/data/probe_ladder_v2_mlp2/analysis/seed_level_results.json :: table[bvc_qcd,mlp,162].rejection_points['0.90'].rejection :: sha256 53e1ec81343ac2c1 :: at 90% signal efficiency, mean +- SD over seeds; 0 of 5 seeds at the cap
\newcommand{\ProbeAucBvcQcdMlpFourthree}{0.98341}  % experiments/FIGS/data/probe_ladder_v2_mlp2/analysis/seed_level_results.json :: levels.bvc_qcd.mlp[2].mean_auc :: sha256 53e1ec81343ac2c1 :: 5 seeds, 0 saturated at AUC=1
\newcommand{\ProbeAucSdBvcQcdMlpFourthree}{0.00052}  % experiments/FIGS/data/probe_ladder_v2_mlp2/analysis/seed_level_results.json :: table[bvc_qcd,mlp,43].auc (SD over seeds, ddof=1) :: sha256 53e1ec81343ac2c1
\newcommand{\ProbeLogOneMinusAucBvcQcdMlpFourthree}{\ensuremath{-}4.0995}  % experiments/FIGS/data/probe_ladder_v2_mlp2/analysis/seed_level_results.json :: levels.bvc_qcd.mlp[2].mean :: sha256 53e1ec81343ac2c1 :: natural log of 1-AUC, lower is better
\newcommand{\ProbeLogOneMinusAucSdBvcQcdMlpFourthree}{0.0318}  % experiments/FIGS/data/probe_ladder_v2_mlp2/analysis/seed_level_results.json :: levels.bvc_qcd.mlp[2].seed_sd :: sha256 53e1ec81343ac2c1
\newcommand{\ProbeOmaBvcQcdMlpFourthree}{\ensuremath{(1.66 \pm 0.05)\times10^{-2}}}  % experiments/FIGS/data/probe_ladder_v2_mlp2/analysis/seed_level_results.json :: table[bvc_qcd,mlp,43].auc :: sha256 53e1ec81343ac2c1 :: 1-AUC, mean +- SD over 5 seeds
\newcommand{\ProbeBkgLeftBvcQcdMlpFourthree}{119.2}  % experiments/FIGS/data/probe_ladder_v2_mlp2/analysis/seed_level_results.json :: levels.bvc_qcd.mlp[2].rejection_points['0.90'].mean_n_bkg_pass :: sha256 53e1ec81343ac2c1 :: background jets passing the cut, mean over seeds
\newcommand{\ProbeRejBvcQcdMlpFourthree}{45.1\,\ensuremath{\pm}\,1.3}  % experiments/FIGS/data/probe_ladder_v2_mlp2/analysis/seed_level_results.json :: table[bvc_qcd,mlp,43].rejection_points['0.90'].rejection :: sha256 53e1ec81343ac2c1 :: at 90% signal efficiency, mean +- SD over seeds; 0 of 5 seeds at the cap
\newcommand{\ProbeAucBvcQcdMlpOneseven}{0.96226}  % experiments/FIGS/data/probe_ladder_v2_mlp2/analysis/seed_level_results.json :: levels.bvc_qcd.mlp[3].mean_auc :: sha256 53e1ec81343ac2c1 :: 5 seeds, 0 saturated at AUC=1
\newcommand{\ProbeAucSdBvcQcdMlpOneseven}{0.00373}  % experiments/FIGS/data/probe_ladder_v2_mlp2/analysis/seed_level_results.json :: table[bvc_qcd,mlp,17].auc (SD over seeds, ddof=1) :: sha256 53e1ec81343ac2c1
\newcommand{\ProbeLogOneMinusAucBvcQcdMlpOneseven}{\ensuremath{-}3.2811}  % experiments/FIGS/data/probe_ladder_v2_mlp2/analysis/seed_level_results.json :: levels.bvc_qcd.mlp[3].mean :: sha256 53e1ec81343ac2c1 :: natural log of 1-AUC, lower is better
\newcommand{\ProbeLogOneMinusAucSdBvcQcdMlpOneseven}{0.1010}  % experiments/FIGS/data/probe_ladder_v2_mlp2/analysis/seed_level_results.json :: levels.bvc_qcd.mlp[3].seed_sd :: sha256 53e1ec81343ac2c1
\newcommand{\ProbeOmaBvcQcdMlpOneseven}{\ensuremath{(3.8 \pm 0.4)\times10^{-2}}}  % experiments/FIGS/data/probe_ladder_v2_mlp2/analysis/seed_level_results.json :: table[bvc_qcd,mlp,17].auc :: sha256 53e1ec81343ac2c1 :: 1-AUC, mean +- SD over 5 seeds
\newcommand{\ProbeBkgLeftBvcQcdMlpOneseven}{458.8}  % experiments/FIGS/data/probe_ladder_v2_mlp2/analysis/seed_level_results.json :: levels.bvc_qcd.mlp[3].rejection_points['0.90'].mean_n_bkg_pass :: sha256 53e1ec81343ac2c1 :: background jets passing the cut, mean over seeds
\newcommand{\ProbeRejBvcQcdMlpOneseven}{11.8\,\ensuremath{\pm}\,1.5}  % experiments/FIGS/data/probe_ladder_v2_mlp2/analysis/seed_level_results.json :: table[bvc_qcd,mlp,17].rejection_points['0.90'].rejection :: sha256 53e1ec81343ac2c1 :: at 90% signal efficiency, mean +- SD over seeds; 0 of 5 seeds at the cap
\newcommand{\ProbeOmaRatioBvcQcdMlpOnesixtwoOverOneeighteight}{1.11}  % experiments/FIGS/data/probe_ladder_v2_mlp2/analysis/seed_level_results.json :: table[bvc_qcd,mlp,162/188].auc :: sha256 53e1ec81343ac2c1 :: ratio of the seed means of 1-AUC, coarser over finer
\newcommand{\ProbeOmaRatioBvcQcdMlpFourthreeOverOneeighteight}{1.13}  % experiments/FIGS/data/probe_ladder_v2_mlp2/analysis/seed_level_results.json :: table[bvc_qcd,mlp,43/188].auc :: sha256 53e1ec81343ac2c1 :: ratio of the seed means of 1-AUC, coarser over finer
\newcommand{\ProbeOmaRatioBvcQcdMlpOnesevenOverOneeighteight}{2.6}  % experiments/FIGS/data/probe_ladder_v2_mlp2/analysis/seed_level_results.json :: table[bvc_qcd,mlp,17/188].auc :: sha256 53e1ec81343ac2c1 :: ratio of the seed means of 1-AUC, coarser over finer
\newcommand{\ProbeOmaRatioBvcQcdMlpFourthreeOverOnesixtwo}{1.01}  % experiments/FIGS/data/probe_ladder_v2_mlp2/analysis/seed_level_results.json :: table[bvc_qcd,mlp,43/162].auc :: sha256 53e1ec81343ac2c1 :: ratio of the seed means of 1-AUC, coarser over finer
\newcommand{\ProbeOmaRatioBvcQcdMlpOnesevenOverOnesixtwo}{2.3}  % experiments/FIGS/data/probe_ladder_v2_mlp2/analysis/seed_level_results.json :: table[bvc_qcd,mlp,17/162].auc :: sha256 53e1ec81343ac2c1 :: ratio of the seed means of 1-AUC, coarser over finer
\newcommand{\ProbeOmaRatioBvcQcdMlpOnesevenOverFourthree}{2.3}  % experiments/FIGS/data/probe_ladder_v2_mlp2/analysis/seed_level_results.json :: table[bvc_qcd,mlp,17/43].auc :: sha256 53e1ec81343ac2c1 :: ratio of the seed means of 1-AUC, coarser over finer
\newcommand{\ProbeAucBvcResonantLinearOneeighteight}{0.99768}  % experiments/FIGS/data/probe_ladder_v2_mlp2/analysis/seed_level_results.json :: levels.bvc_resonant.linear[0].mean_auc :: sha256 53e1ec81343ac2c1 :: 5 seeds, 0 saturated at AUC=1
\newcommand{\ProbeAucSdBvcResonantLinearOneeighteight}{0.00007}  % experiments/FIGS/data/probe_ladder_v2_mlp2/analysis/seed_level_results.json :: table[bvc_resonant,linear,188].auc (SD over seeds, ddof=1) :: sha256 53e1ec81343ac2c1
\newcommand{\ProbeLogOneMinusAucBvcResonantLinearOneeighteight}{\ensuremath{-}6.0683}  % experiments/FIGS/data/probe_ladder_v2_mlp2/analysis/seed_level_results.json :: levels.bvc_resonant.linear[0].mean :: sha256 53e1ec81343ac2c1 :: natural log of 1-AUC, lower is better
\newcommand{\ProbeLogOneMinusAucSdBvcResonantLinearOneeighteight}{0.0317}  % experiments/FIGS/data/probe_ladder_v2_mlp2/analysis/seed_level_results.json :: levels.bvc_resonant.linear[0].seed_sd :: sha256 53e1ec81343ac2c1
\newcommand{\ProbeOmaBvcResonantLinearOneeighteight}{\ensuremath{(2.32 \pm 0.07)\times10^{-3}}}  % experiments/FIGS/data/probe_ladder_v2_mlp2/analysis/seed_level_results.json :: table[bvc_resonant,linear,188].auc :: sha256 53e1ec81343ac2c1 :: 1-AUC, mean +- SD over 5 seeds
\newcommand{\ProbeBkgLeftBvcResonantLinearOneeighteight}{9.0}  % experiments/FIGS/data/probe_ladder_v2_mlp2/analysis/seed_level_results.json :: levels.bvc_resonant.linear[0].rejection_points['0.90'].mean_n_bkg_pass :: sha256 53e1ec81343ac2c1 :: background jets passing the cut, mean over seeds
\newcommand{\ProbeRejBvcResonantLinearOneeighteight}{1{,}380\,\ensuremath{\pm}\,350}  % experiments/FIGS/data/probe_ladder_v2_mlp2/analysis/seed_level_results.json :: table[bvc_resonant,linear,188].rejection_points['0.90'].rejection :: sha256 53e1ec81343ac2c1 :: at 90% signal efficiency, mean +- SD over seeds; 0 of 5 seeds at the cap
\newcommand{\ProbeAucBvcResonantLinearOnesixtwo}{0.99764}  % experiments/FIGS/data/probe_ladder_v2_mlp2/analysis/seed_level_results.json :: levels.bvc_resonant.linear[1].mean_auc :: sha256 53e1ec81343ac2c1 :: 5 seeds, 0 saturated at AUC=1
\newcommand{\ProbeAucSdBvcResonantLinearOnesixtwo}{0.00006}  % experiments/FIGS/data/probe_ladder_v2_mlp2/analysis/seed_level_results.json :: table[bvc_resonant,linear,162].auc (SD over seeds, ddof=1) :: sha256 53e1ec81343ac2c1
\newcommand{\ProbeLogOneMinusAucBvcResonantLinearOnesixtwo}{\ensuremath{-}6.0493}  % experiments/FIGS/data/probe_ladder_v2_mlp2/analysis/seed_level_results.json :: levels.bvc_resonant.linear[1].mean :: sha256 53e1ec81343ac2c1 :: natural log of 1-AUC, lower is better
\newcommand{\ProbeLogOneMinusAucSdBvcResonantLinearOnesixtwo}{0.0261}  % experiments/FIGS/data/probe_ladder_v2_mlp2/analysis/seed_level_results.json :: levels.bvc_resonant.linear[1].seed_sd :: sha256 53e1ec81343ac2c1
\newcommand{\ProbeOmaBvcResonantLinearOnesixtwo}{\ensuremath{(2.36 \pm 0.06)\times10^{-3}}}  % experiments/FIGS/data/probe_ladder_v2_mlp2/analysis/seed_level_results.json :: table[bvc_resonant,linear,162].auc :: sha256 53e1ec81343ac2c1 :: 1-AUC, mean +- SD over 5 seeds
\newcommand{\ProbeBkgLeftBvcResonantLinearOnesixtwo}{9.4}  % experiments/FIGS/data/probe_ladder_v2_mlp2/analysis/seed_level_results.json :: levels.bvc_resonant.linear[1].rejection_points['0.90'].mean_n_bkg_pass :: sha256 53e1ec81343ac2c1 :: background jets passing the cut, mean over seeds
\newcommand{\ProbeRejBvcResonantLinearOnesixtwo}{1{,}290\,\ensuremath{\pm}\,240}  % experiments/FIGS/data/probe_ladder_v2_mlp2/analysis/seed_level_results.json :: table[bvc_resonant,linear,162].rejection_points['0.90'].rejection :: sha256 53e1ec81343ac2c1 :: at 90% signal efficiency, mean +- SD over seeds; 0 of 5 seeds at the cap
\newcommand{\ProbeAucBvcResonantLinearFourthree}{0.99748}  % experiments/FIGS/data/probe_ladder_v2_mlp2/analysis/seed_level_results.json :: levels.bvc_resonant.linear[2].mean_auc :: sha256 53e1ec81343ac2c1 :: 5 seeds, 0 saturated at AUC=1
\newcommand{\ProbeAucSdBvcResonantLinearFourthree}{0.00007}  % experiments/FIGS/data/probe_ladder_v2_mlp2/analysis/seed_level_results.json :: table[bvc_resonant,linear,43].auc (SD over seeds, ddof=1) :: sha256 53e1ec81343ac2c1
\newcommand{\ProbeLogOneMinusAucBvcResonantLinearFourthree}{\ensuremath{-}5.9834}  % experiments/FIGS/data/probe_ladder_v2_mlp2/analysis/seed_level_results.json :: levels.bvc_resonant.linear[2].mean :: sha256 53e1ec81343ac2c1 :: natural log of 1-AUC, lower is better
\newcommand{\ProbeLogOneMinusAucSdBvcResonantLinearFourthree}{0.0264}  % experiments/FIGS/data/probe_ladder_v2_mlp2/analysis/seed_level_results.json :: levels.bvc_resonant.linear[2].seed_sd :: sha256 53e1ec81343ac2c1
\newcommand{\ProbeOmaBvcResonantLinearFourthree}{\ensuremath{(2.52 \pm 0.07)\times10^{-3}}}  % experiments/FIGS/data/probe_ladder_v2_mlp2/analysis/seed_level_results.json :: table[bvc_resonant,linear,43].auc :: sha256 53e1ec81343ac2c1 :: 1-AUC, mean +- SD over 5 seeds
\newcommand{\ProbeBkgLeftBvcResonantLinearFourthree}{11.0}  % experiments/FIGS/data/probe_ladder_v2_mlp2/analysis/seed_level_results.json :: levels.bvc_resonant.linear[2].rejection_points['0.90'].mean_n_bkg_pass :: sha256 53e1ec81343ac2c1 :: background jets passing the cut, mean over seeds
\newcommand{\ProbeRejBvcResonantLinearFourthree}{1{,}130\,\ensuremath{\pm}\,260}  % experiments/FIGS/data/probe_ladder_v2_mlp2/analysis/seed_level_results.json :: table[bvc_resonant,linear,43].rejection_points['0.90'].rejection :: sha256 53e1ec81343ac2c1 :: at 90% signal efficiency, mean +- SD over seeds; 0 of 5 seeds at the cap
\newcommand{\ProbeAucBvcResonantLinearOneseven}{0.99067}  % experiments/FIGS/data/probe_ladder_v2_mlp2/analysis/seed_level_results.json :: levels.bvc_resonant.linear[3].mean_auc :: sha256 53e1ec81343ac2c1 :: 5 seeds, 0 saturated at AUC=1
\newcommand{\ProbeAucSdBvcResonantLinearOneseven}{0.00147}  % experiments/FIGS/data/probe_ladder_v2_mlp2/analysis/seed_level_results.json :: table[bvc_resonant,linear,17].auc (SD over seeds, ddof=1) :: sha256 53e1ec81343ac2c1
\newcommand{\ProbeLogOneMinusAucBvcResonantLinearOneseven}{\ensuremath{-}4.6848}  % experiments/FIGS/data/probe_ladder_v2_mlp2/analysis/seed_level_results.json :: levels.bvc_resonant.linear[3].mean :: sha256 53e1ec81343ac2c1 :: natural log of 1-AUC, lower is better
\newcommand{\ProbeLogOneMinusAucSdBvcResonantLinearOneseven}{0.1619}  % experiments/FIGS/data/probe_ladder_v2_mlp2/analysis/seed_level_results.json :: levels.bvc_resonant.linear[3].seed_sd :: sha256 53e1ec81343ac2c1
\newcommand{\ProbeOmaBvcResonantLinearOneseven}{\ensuremath{(9.3 \pm 1.5)\times10^{-3}}}  % experiments/FIGS/data/probe_ladder_v2_mlp2/analysis/seed_level_results.json :: table[bvc_resonant,linear,17].auc :: sha256 53e1ec81343ac2c1 :: 1-AUC, mean +- SD over 5 seeds
\newcommand{\ProbeBkgLeftBvcResonantLinearOneseven}{90.8}  % experiments/FIGS/data/probe_ladder_v2_mlp2/analysis/seed_level_results.json :: levels.bvc_resonant.linear[3].rejection_points['0.90'].mean_n_bkg_pass :: sha256 53e1ec81343ac2c1 :: background jets passing the cut, mean over seeds
\newcommand{\ProbeRejBvcResonantLinearOneseven}{134\,\ensuremath{\pm}\,21}  % experiments/FIGS/data/probe_ladder_v2_mlp2/analysis/seed_level_results.json :: table[bvc_resonant,linear,17].rejection_points['0.90'].rejection :: sha256 53e1ec81343ac2c1 :: at 90% signal efficiency, mean +- SD over seeds; 0 of 5 seeds at the cap
\newcommand{\ProbeOmaRatioBvcResonantLinearOnesixtwoOverOneeighteight}{1.02}  % experiments/FIGS/data/probe_ladder_v2_mlp2/analysis/seed_level_results.json :: table[bvc_resonant,linear,162/188].auc :: sha256 53e1ec81343ac2c1 :: ratio of the seed means of 1-AUC, coarser over finer
\newcommand{\ProbeOmaRatioBvcResonantLinearFourthreeOverOneeighteight}{1.09}  % experiments/FIGS/data/probe_ladder_v2_mlp2/analysis/seed_level_results.json :: table[bvc_resonant,linear,43/188].auc :: sha256 53e1ec81343ac2c1 :: ratio of the seed means of 1-AUC, coarser over finer
\newcommand{\ProbeOmaRatioBvcResonantLinearOnesevenOverOneeighteight}{4.0}  % experiments/FIGS/data/probe_ladder_v2_mlp2/analysis/seed_level_results.json :: table[bvc_resonant,linear,17/188].auc :: sha256 53e1ec81343ac2c1 :: ratio of the seed means of 1-AUC, coarser over finer
\newcommand{\ProbeOmaRatioBvcResonantLinearFourthreeOverOnesixtwo}{1.07}  % experiments/FIGS/data/probe_ladder_v2_mlp2/analysis/seed_level_results.json :: table[bvc_resonant,linear,43/162].auc :: sha256 53e1ec81343ac2c1 :: ratio of the seed means of 1-AUC, coarser over finer
\newcommand{\ProbeOmaRatioBvcResonantLinearOnesevenOverOnesixtwo}{4.0}  % experiments/FIGS/data/probe_ladder_v2_mlp2/analysis/seed_level_results.json :: table[bvc_resonant,linear,17/162].auc :: sha256 53e1ec81343ac2c1 :: ratio of the seed means of 1-AUC, coarser over finer
\newcommand{\ProbeOmaRatioBvcResonantLinearOnesevenOverFourthree}{3.7}  % experiments/FIGS/data/probe_ladder_v2_mlp2/analysis/seed_level_results.json :: table[bvc_resonant,linear,17/43].auc :: sha256 53e1ec81343ac2c1 :: ratio of the seed means of 1-AUC, coarser over finer
\newcommand{\ProbeAucBvcResonantMlpOneeighteight}{0.99778}  % experiments/FIGS/data/probe_ladder_v2_mlp2/analysis/seed_level_results.json :: levels.bvc_resonant.mlp[0].mean_auc :: sha256 53e1ec81343ac2c1 :: 5 seeds, 0 saturated at AUC=1
\newcommand{\ProbeAucSdBvcResonantMlpOneeighteight}{0.00009}  % experiments/FIGS/data/probe_ladder_v2_mlp2/analysis/seed_level_results.json :: table[bvc_resonant,mlp,188].auc (SD over seeds, ddof=1) :: sha256 53e1ec81343ac2c1
\newcommand{\ProbeLogOneMinusAucBvcResonantMlpOneeighteight}{\ensuremath{-}6.1102}  % experiments/FIGS/data/probe_ladder_v2_mlp2/analysis/seed_level_results.json :: levels.bvc_resonant.mlp[0].mean :: sha256 53e1ec81343ac2c1 :: natural log of 1-AUC, lower is better
\newcommand{\ProbeLogOneMinusAucSdBvcResonantMlpOneeighteight}{0.0383}  % experiments/FIGS/data/probe_ladder_v2_mlp2/analysis/seed_level_results.json :: levels.bvc_resonant.mlp[0].seed_sd :: sha256 53e1ec81343ac2c1
\newcommand{\ProbeOmaBvcResonantMlpOneeighteight}{\ensuremath{(2.22 \pm 0.09)\times10^{-3}}}  % experiments/FIGS/data/probe_ladder_v2_mlp2/analysis/seed_level_results.json :: table[bvc_resonant,mlp,188].auc :: sha256 53e1ec81343ac2c1 :: 1-AUC, mean +- SD over 5 seeds
\newcommand{\ProbeBkgLeftBvcResonantMlpOneeighteight}{9.8}  % experiments/FIGS/data/probe_ladder_v2_mlp2/analysis/seed_level_results.json :: levels.bvc_resonant.mlp[0].rejection_points['0.90'].mean_n_bkg_pass :: sha256 53e1ec81343ac2c1 :: background jets passing the cut, mean over seeds
\newcommand{\ProbeRejBvcResonantMlpOneeighteight}{1{,}230\,\ensuremath{\pm}\,180}  % experiments/FIGS/data/probe_ladder_v2_mlp2/analysis/seed_level_results.json :: table[bvc_resonant,mlp,188].rejection_points['0.90'].rejection :: sha256 53e1ec81343ac2c1 :: at 90% signal efficiency, mean +- SD over seeds; 0 of 5 seeds at the cap
\newcommand{\ProbeAucBvcResonantMlpOnesixtwo}{0.99768}  % experiments/FIGS/data/probe_ladder_v2_mlp2/analysis/seed_level_results.json :: levels.bvc_resonant.mlp[1].mean_auc :: sha256 53e1ec81343ac2c1 :: 5 seeds, 0 saturated at AUC=1
\newcommand{\ProbeAucSdBvcResonantMlpOnesixtwo}{0.00005}  % experiments/FIGS/data/probe_ladder_v2_mlp2/analysis/seed_level_results.json :: table[bvc_resonant,mlp,162].auc (SD over seeds, ddof=1) :: sha256 53e1ec81343ac2c1
\newcommand{\ProbeLogOneMinusAucBvcResonantMlpOnesixtwo}{\ensuremath{-}6.0667}  % experiments/FIGS/data/probe_ladder_v2_mlp2/analysis/seed_level_results.json :: levels.bvc_resonant.mlp[1].mean :: sha256 53e1ec81343ac2c1 :: natural log of 1-AUC, lower is better
\newcommand{\ProbeLogOneMinusAucSdBvcResonantMlpOnesixtwo}{0.0220}  % experiments/FIGS/data/probe_ladder_v2_mlp2/analysis/seed_level_results.json :: levels.bvc_resonant.mlp[1].seed_sd :: sha256 53e1ec81343ac2c1
\newcommand{\ProbeOmaBvcResonantMlpOnesixtwo}{\ensuremath{(2.32 \pm 0.05)\times10^{-3}}}  % experiments/FIGS/data/probe_ladder_v2_mlp2/analysis/seed_level_results.json :: table[bvc_resonant,mlp,162].auc :: sha256 53e1ec81343ac2c1 :: 1-AUC, mean +- SD over 5 seeds
\newcommand{\ProbeBkgLeftBvcResonantMlpOnesixtwo}{10.4}  % experiments/FIGS/data/probe_ladder_v2_mlp2/analysis/seed_level_results.json :: levels.bvc_resonant.mlp[1].rejection_points['0.90'].mean_n_bkg_pass :: sha256 53e1ec81343ac2c1 :: background jets passing the cut, mean over seeds
\newcommand{\ProbeRejBvcResonantMlpOnesixtwo}{1{,}170\,\ensuremath{\pm}\,220}  % experiments/FIGS/data/probe_ladder_v2_mlp2/analysis/seed_level_results.json :: table[bvc_resonant,mlp,162].rejection_points['0.90'].rejection :: sha256 53e1ec81343ac2c1 :: at 90% signal efficiency, mean +- SD over seeds; 0 of 5 seeds at the cap
\newcommand{\ProbeAucBvcResonantMlpFourthree}{0.99747}  % experiments/FIGS/data/probe_ladder_v2_mlp2/analysis/seed_level_results.json :: levels.bvc_resonant.mlp[2].mean_auc :: sha256 53e1ec81343ac2c1 :: 5 seeds, 0 saturated at AUC=1
\newcommand{\ProbeAucSdBvcResonantMlpFourthree}{0.00004}  % experiments/FIGS/data/probe_ladder_v2_mlp2/analysis/seed_level_results.json :: table[bvc_resonant,mlp,43].auc (SD over seeds, ddof=1) :: sha256 53e1ec81343ac2c1
\newcommand{\ProbeLogOneMinusAucBvcResonantMlpFourthree}{\ensuremath{-}5.9792}  % experiments/FIGS/data/probe_ladder_v2_mlp2/analysis/seed_level_results.json :: levels.bvc_resonant.mlp[2].mean :: sha256 53e1ec81343ac2c1 :: natural log of 1-AUC, lower is better
\newcommand{\ProbeLogOneMinusAucSdBvcResonantMlpFourthree}{0.0164}  % experiments/FIGS/data/probe_ladder_v2_mlp2/analysis/seed_level_results.json :: levels.bvc_resonant.mlp[2].seed_sd :: sha256 53e1ec81343ac2c1
\newcommand{\ProbeOmaBvcResonantMlpFourthree}{\ensuremath{(2.53 \pm 0.04)\times10^{-3}}}  % experiments/FIGS/data/probe_ladder_v2_mlp2/analysis/seed_level_results.json :: table[bvc_resonant,mlp,43].auc :: sha256 53e1ec81343ac2c1 :: 1-AUC, mean +- SD over 5 seeds
\newcommand{\ProbeBkgLeftBvcResonantMlpFourthree}{11.8}  % experiments/FIGS/data/probe_ladder_v2_mlp2/analysis/seed_level_results.json :: levels.bvc_resonant.mlp[2].rejection_points['0.90'].mean_n_bkg_pass :: sha256 53e1ec81343ac2c1 :: background jets passing the cut, mean over seeds
\newcommand{\ProbeRejBvcResonantMlpFourthree}{1{,}030\,\ensuremath{\pm}\,170}  % experiments/FIGS/data/probe_ladder_v2_mlp2/analysis/seed_level_results.json :: table[bvc_resonant,mlp,43].rejection_points['0.90'].rejection :: sha256 53e1ec81343ac2c1 :: at 90% signal efficiency, mean +- SD over seeds; 0 of 5 seeds at the cap
\newcommand{\ProbeAucBvcResonantMlpOneseven}{0.99274}  % experiments/FIGS/data/probe_ladder_v2_mlp2/analysis/seed_level_results.json :: levels.bvc_resonant.mlp[3].mean_auc :: sha256 53e1ec81343ac2c1 :: 5 seeds, 0 saturated at AUC=1
\newcommand{\ProbeAucSdBvcResonantMlpOneseven}{0.00089}  % experiments/FIGS/data/probe_ladder_v2_mlp2/analysis/seed_level_results.json :: table[bvc_resonant,mlp,17].auc (SD over seeds, ddof=1) :: sha256 53e1ec81343ac2c1
\newcommand{\ProbeLogOneMinusAucBvcResonantMlpOneseven}{\ensuremath{-}4.9309}  % experiments/FIGS/data/probe_ladder_v2_mlp2/analysis/seed_level_results.json :: levels.bvc_resonant.mlp[3].mean :: sha256 53e1ec81343ac2c1 :: natural log of 1-AUC, lower is better
\newcommand{\ProbeLogOneMinusAucSdBvcResonantMlpOneseven}{0.1224}  % experiments/FIGS/data/probe_ladder_v2_mlp2/analysis/seed_level_results.json :: levels.bvc_resonant.mlp[3].seed_sd :: sha256 53e1ec81343ac2c1
\newcommand{\ProbeOmaBvcResonantMlpOneseven}{\ensuremath{(7.3 \pm 0.9)\times10^{-3}}}  % experiments/FIGS/data/probe_ladder_v2_mlp2/analysis/seed_level_results.json :: table[bvc_resonant,mlp,17].auc :: sha256 53e1ec81343ac2c1 :: 1-AUC, mean +- SD over 5 seeds
\newcommand{\ProbeBkgLeftBvcResonantMlpOneseven}{62.6}  % experiments/FIGS/data/probe_ladder_v2_mlp2/analysis/seed_level_results.json :: levels.bvc_resonant.mlp[3].rejection_points['0.90'].mean_n_bkg_pass :: sha256 53e1ec81343ac2c1 :: background jets passing the cut, mean over seeds
\newcommand{\ProbeRejBvcResonantMlpOneseven}{192\,\ensuremath{\pm}\,21}  % experiments/FIGS/data/probe_ladder_v2_mlp2/analysis/seed_level_results.json :: table[bvc_resonant,mlp,17].rejection_points['0.90'].rejection :: sha256 53e1ec81343ac2c1 :: at 90% signal efficiency, mean +- SD over seeds; 0 of 5 seeds at the cap
\newcommand{\ProbeOmaRatioBvcResonantMlpOnesixtwoOverOneeighteight}{1.04}  % experiments/FIGS/data/probe_ladder_v2_mlp2/analysis/seed_level_results.json :: table[bvc_resonant,mlp,162/188].auc :: sha256 53e1ec81343ac2c1 :: ratio of the seed means of 1-AUC, coarser over finer
\newcommand{\ProbeOmaRatioBvcResonantMlpFourthreeOverOneeighteight}{1.14}  % experiments/FIGS/data/probe_ladder_v2_mlp2/analysis/seed_level_results.json :: table[bvc_resonant,mlp,43/188].auc :: sha256 53e1ec81343ac2c1 :: ratio of the seed means of 1-AUC, coarser over finer
\newcommand{\ProbeOmaRatioBvcResonantMlpOnesevenOverOneeighteight}{3.3}  % experiments/FIGS/data/probe_ladder_v2_mlp2/analysis/seed_level_results.json :: table[bvc_resonant,mlp,17/188].auc :: sha256 53e1ec81343ac2c1 :: ratio of the seed means of 1-AUC, coarser over finer
\newcommand{\ProbeOmaRatioBvcResonantMlpFourthreeOverOnesixtwo}{1.09}  % experiments/FIGS/data/probe_ladder_v2_mlp2/analysis/seed_level_results.json :: table[bvc_resonant,mlp,43/162].auc :: sha256 53e1ec81343ac2c1 :: ratio of the seed means of 1-AUC, coarser over finer
\newcommand{\ProbeOmaRatioBvcResonantMlpOnesevenOverOnesixtwo}{3.1}  % experiments/FIGS/data/probe_ladder_v2_mlp2/analysis/seed_level_results.json :: table[bvc_resonant,mlp,17/162].auc :: sha256 53e1ec81343ac2c1 :: ratio of the seed means of 1-AUC, coarser over finer
\newcommand{\ProbeOmaRatioBvcResonantMlpOnesevenOverFourthree}{2.9}  % experiments/FIGS/data/probe_ladder_v2_mlp2/analysis/seed_level_results.json :: table[bvc_resonant,mlp,17/43].auc :: sha256 53e1ec81343ac2c1 :: ratio of the seed means of 1-AUC, coarser over finer
\newcommand{\ProbeAucEeVsMmLinearOneeighteight}{\ensuremath{1^{\dagger}}}  % experiments/FIGS/data/probe_ladder_v2_mlp2/analysis/seed_level_results.json :: levels.ee_vs_mm.linear[0].mean_auc :: sha256 53e1ec81343ac2c1 :: 5 seeds, 5 saturated at AUC=1
\newcommand{\ProbeAucSdEeVsMmLinearOneeighteight}{0.00000}  % experiments/FIGS/data/probe_ladder_v2_mlp2/analysis/seed_level_results.json :: table[ee_vs_mm,linear,188].auc (SD over seeds, ddof=1) :: sha256 53e1ec81343ac2c1
\newcommand{\ProbeLogOneMinusAucEeVsMmLinearOneeighteight}{\ensuremath{-}16.9579}  % experiments/FIGS/data/probe_ladder_v2_mlp2/analysis/seed_level_results.json :: levels.ee_vs_mm.linear[0].mean :: sha256 53e1ec81343ac2c1 :: natural log of 1-AUC, lower is better
\newcommand{\ProbeLogOneMinusAucSdEeVsMmLinearOneeighteight}{0.0000}  % experiments/FIGS/data/probe_ladder_v2_mlp2/analysis/seed_level_results.json :: levels.ee_vs_mm.linear[0].seed_sd :: sha256 53e1ec81343ac2c1
\newcommand{\ProbeOmaEeVsMmLinearOneeighteight}{---}  % experiments/FIGS/data/probe_ladder_v2_mlp2/analysis/seed_level_results.json :: table[ee_vs_mm,linear,188].auc :: sha256 53e1ec81343ac2c1 :: not quoted: the AUC reached 1 in a seed, so 1-AUC is only a bound
\newcommand{\ProbeBkgLeftEeVsMmLinearOneeighteight}{0.0}  % experiments/FIGS/data/probe_ladder_v2_mlp2/analysis/seed_level_results.json :: levels.ee_vs_mm.linear[0].rejection_points['0.90'].mean_n_bkg_pass :: sha256 53e1ec81343ac2c1 :: background jets passing the cut, mean over seeds
\newcommand{\ProbeRejEeVsMmLinearOneeighteight}{\ensuremath{>}1{,}571}  % experiments/FIGS/data/probe_ladder_v2_mlp2/analysis/seed_level_results.json :: table[ee_vs_mm,linear,188].rejection_points['0.90'].rejection :: sha256 53e1ec81343ac2c1 :: at 90% signal efficiency, mean +- SD over seeds; 5 of 5 seeds at the cap
\newcommand{\ProbeAucEeVsMmLinearOnesixtwo}{\ensuremath{1^{\dagger}}}  % experiments/FIGS/data/probe_ladder_v2_mlp2/analysis/seed_level_results.json :: levels.ee_vs_mm.linear[1].mean_auc :: sha256 53e1ec81343ac2c1 :: 5 seeds, 5 saturated at AUC=1
\newcommand{\ProbeAucSdEeVsMmLinearOnesixtwo}{0.00000}  % experiments/FIGS/data/probe_ladder_v2_mlp2/analysis/seed_level_results.json :: table[ee_vs_mm,linear,162].auc (SD over seeds, ddof=1) :: sha256 53e1ec81343ac2c1
\newcommand{\ProbeLogOneMinusAucEeVsMmLinearOnesixtwo}{\ensuremath{-}16.9579}  % experiments/FIGS/data/probe_ladder_v2_mlp2/analysis/seed_level_results.json :: levels.ee_vs_mm.linear[1].mean :: sha256 53e1ec81343ac2c1 :: natural log of 1-AUC, lower is better
\newcommand{\ProbeLogOneMinusAucSdEeVsMmLinearOnesixtwo}{0.0000}  % experiments/FIGS/data/probe_ladder_v2_mlp2/analysis/seed_level_results.json :: levels.ee_vs_mm.linear[1].seed_sd :: sha256 53e1ec81343ac2c1
\newcommand{\ProbeOmaEeVsMmLinearOnesixtwo}{---}  % experiments/FIGS/data/probe_ladder_v2_mlp2/analysis/seed_level_results.json :: table[ee_vs_mm,linear,162].auc :: sha256 53e1ec81343ac2c1 :: not quoted: the AUC reached 1 in a seed, so 1-AUC is only a bound
\newcommand{\ProbeBkgLeftEeVsMmLinearOnesixtwo}{0.0}  % experiments/FIGS/data/probe_ladder_v2_mlp2/analysis/seed_level_results.json :: levels.ee_vs_mm.linear[1].rejection_points['0.90'].mean_n_bkg_pass :: sha256 53e1ec81343ac2c1 :: background jets passing the cut, mean over seeds
\newcommand{\ProbeRejEeVsMmLinearOnesixtwo}{\ensuremath{>}1{,}571}  % experiments/FIGS/data/probe_ladder_v2_mlp2/analysis/seed_level_results.json :: table[ee_vs_mm,linear,162].rejection_points['0.90'].rejection :: sha256 53e1ec81343ac2c1 :: at 90% signal efficiency, mean +- SD over seeds; 5 of 5 seeds at the cap
\newcommand{\ProbeAucEeVsMmLinearFourthree}{0.97519}  % experiments/FIGS/data/probe_ladder_v2_mlp2/analysis/seed_level_results.json :: levels.ee_vs_mm.linear[2].mean_auc :: sha256 53e1ec81343ac2c1 :: 5 seeds, 0 saturated at AUC=1
\newcommand{\ProbeAucSdEeVsMmLinearFourthree}{0.02073}  % experiments/FIGS/data/probe_ladder_v2_mlp2/analysis/seed_level_results.json :: table[ee_vs_mm,linear,43].auc (SD over seeds, ddof=1) :: sha256 53e1ec81343ac2c1
\newcommand{\ProbeLogOneMinusAucEeVsMmLinearFourthree}{\ensuremath{-}4.0326}  % experiments/FIGS/data/probe_ladder_v2_mlp2/analysis/seed_level_results.json :: levels.ee_vs_mm.linear[2].mean :: sha256 53e1ec81343ac2c1 :: natural log of 1-AUC, lower is better
\newcommand{\ProbeLogOneMinusAucSdEeVsMmLinearFourthree}{0.9467}  % experiments/FIGS/data/probe_ladder_v2_mlp2/analysis/seed_level_results.json :: levels.ee_vs_mm.linear[2].seed_sd :: sha256 53e1ec81343ac2c1
\newcommand{\ProbeOmaEeVsMmLinearFourthree}{\ensuremath{(2.5 \pm 2.1)\times10^{-2}}}  % experiments/FIGS/data/probe_ladder_v2_mlp2/analysis/seed_level_results.json :: table[ee_vs_mm,linear,43].auc :: sha256 53e1ec81343ac2c1 :: 1-AUC, mean +- SD over 5 seeds
\newcommand{\ProbeBkgLeftEeVsMmLinearFourthree}{155.8}  % experiments/FIGS/data/probe_ladder_v2_mlp2/analysis/seed_level_results.json :: levels.ee_vs_mm.linear[2].rejection_points['0.90'].mean_n_bkg_pass :: sha256 53e1ec81343ac2c1 :: background jets passing the cut, mean over seeds
\newcommand{\ProbeRejEeVsMmLinearFourthree}{80\,\ensuremath{\pm}\,70}  % experiments/FIGS/data/probe_ladder_v2_mlp2/analysis/seed_level_results.json :: table[ee_vs_mm,linear,43].rejection_points['0.90'].rejection :: sha256 53e1ec81343ac2c1 :: at 90% signal efficiency, mean +- SD over seeds; 0 of 5 seeds at the cap
\newcommand{\ProbeAucEeVsMmLinearOneseven}{0.95832}  % experiments/FIGS/data/probe_ladder_v2_mlp2/analysis/seed_level_results.json :: levels.ee_vs_mm.linear[3].mean_auc :: sha256 53e1ec81343ac2c1 :: 5 seeds, 0 saturated at AUC=1
\newcommand{\ProbeAucSdEeVsMmLinearOneseven}{0.00970}  % experiments/FIGS/data/probe_ladder_v2_mlp2/analysis/seed_level_results.json :: table[ee_vs_mm,linear,17].auc (SD over seeds, ddof=1) :: sha256 53e1ec81343ac2c1
\newcommand{\ProbeLogOneMinusAucEeVsMmLinearOneseven}{\ensuremath{-}3.2013}  % experiments/FIGS/data/probe_ladder_v2_mlp2/analysis/seed_level_results.json :: levels.ee_vs_mm.linear[3].mean :: sha256 53e1ec81343ac2c1 :: natural log of 1-AUC, lower is better
\newcommand{\ProbeLogOneMinusAucSdEeVsMmLinearOneseven}{0.2488}  % experiments/FIGS/data/probe_ladder_v2_mlp2/analysis/seed_level_results.json :: levels.ee_vs_mm.linear[3].seed_sd :: sha256 53e1ec81343ac2c1
\newcommand{\ProbeOmaEeVsMmLinearOneseven}{\ensuremath{(4.2 \pm 1.0)\times10^{-2}}}  % experiments/FIGS/data/probe_ladder_v2_mlp2/analysis/seed_level_results.json :: table[ee_vs_mm,linear,17].auc :: sha256 53e1ec81343ac2c1 :: 1-AUC, mean +- SD over 5 seeds
\newcommand{\ProbeBkgLeftEeVsMmLinearOneseven}{289.6}  % experiments/FIGS/data/probe_ladder_v2_mlp2/analysis/seed_level_results.json :: levels.ee_vs_mm.linear[3].rejection_points['0.90'].mean_n_bkg_pass :: sha256 53e1ec81343ac2c1 :: background jets passing the cut, mean over seeds
\newcommand{\ProbeRejEeVsMmLinearOneseven}{17\,\ensuremath{\pm}\,6}  % experiments/FIGS/data/probe_ladder_v2_mlp2/analysis/seed_level_results.json :: table[ee_vs_mm,linear,17].rejection_points['0.90'].rejection :: sha256 53e1ec81343ac2c1 :: at 90% signal efficiency, mean +- SD over seeds; 0 of 5 seeds at the cap
\newcommand{\ProbeOmaRatioEeVsMmLinearOnesevenOverFourthree}{1.68}  % experiments/FIGS/data/probe_ladder_v2_mlp2/analysis/seed_level_results.json :: table[ee_vs_mm,linear,17/43].auc :: sha256 53e1ec81343ac2c1 :: ratio of the seed means of 1-AUC, coarser over finer
\newcommand{\ProbeAucEeVsMmMlpOneeighteight}{\ensuremath{1^{\dagger}}}  % experiments/FIGS/data/probe_ladder_v2_mlp2/analysis/seed_level_results.json :: levels.ee_vs_mm.mlp[0].mean_auc :: sha256 53e1ec81343ac2c1 :: 5 seeds, 5 saturated at AUC=1
\newcommand{\ProbeAucSdEeVsMmMlpOneeighteight}{0.00000}  % experiments/FIGS/data/probe_ladder_v2_mlp2/analysis/seed_level_results.json :: table[ee_vs_mm,mlp,188].auc (SD over seeds, ddof=1) :: sha256 53e1ec81343ac2c1
\newcommand{\ProbeLogOneMinusAucEeVsMmMlpOneeighteight}{\ensuremath{-}16.9579}  % experiments/FIGS/data/probe_ladder_v2_mlp2/analysis/seed_level_results.json :: levels.ee_vs_mm.mlp[0].mean :: sha256 53e1ec81343ac2c1 :: natural log of 1-AUC, lower is better
\newcommand{\ProbeLogOneMinusAucSdEeVsMmMlpOneeighteight}{0.0000}  % experiments/FIGS/data/probe_ladder_v2_mlp2/analysis/seed_level_results.json :: levels.ee_vs_mm.mlp[0].seed_sd :: sha256 53e1ec81343ac2c1
\newcommand{\ProbeOmaEeVsMmMlpOneeighteight}{---}  % experiments/FIGS/data/probe_ladder_v2_mlp2/analysis/seed_level_results.json :: table[ee_vs_mm,mlp,188].auc :: sha256 53e1ec81343ac2c1 :: not quoted: the AUC reached 1 in a seed, so 1-AUC is only a bound
\newcommand{\ProbeBkgLeftEeVsMmMlpOneeighteight}{0.0}  % experiments/FIGS/data/probe_ladder_v2_mlp2/analysis/seed_level_results.json :: levels.ee_vs_mm.mlp[0].rejection_points['0.90'].mean_n_bkg_pass :: sha256 53e1ec81343ac2c1 :: background jets passing the cut, mean over seeds
\newcommand{\ProbeRejEeVsMmMlpOneeighteight}{\ensuremath{>}1{,}571}  % experiments/FIGS/data/probe_ladder_v2_mlp2/analysis/seed_level_results.json :: table[ee_vs_mm,mlp,188].rejection_points['0.90'].rejection :: sha256 53e1ec81343ac2c1 :: at 90% signal efficiency, mean +- SD over seeds; 5 of 5 seeds at the cap
\newcommand{\ProbeAucEeVsMmMlpOnesixtwo}{\ensuremath{1^{\dagger}}}  % experiments/FIGS/data/probe_ladder_v2_mlp2/analysis/seed_level_results.json :: levels.ee_vs_mm.mlp[1].mean_auc :: sha256 53e1ec81343ac2c1 :: 5 seeds, 5 saturated at AUC=1
\newcommand{\ProbeAucSdEeVsMmMlpOnesixtwo}{0.00000}  % experiments/FIGS/data/probe_ladder_v2_mlp2/analysis/seed_level_results.json :: table[ee_vs_mm,mlp,162].auc (SD over seeds, ddof=1) :: sha256 53e1ec81343ac2c1
\newcommand{\ProbeLogOneMinusAucEeVsMmMlpOnesixtwo}{\ensuremath{-}16.9579}  % experiments/FIGS/data/probe_ladder_v2_mlp2/analysis/seed_level_results.json :: levels.ee_vs_mm.mlp[1].mean :: sha256 53e1ec81343ac2c1 :: natural log of 1-AUC, lower is better
\newcommand{\ProbeLogOneMinusAucSdEeVsMmMlpOnesixtwo}{0.0000}  % experiments/FIGS/data/probe_ladder_v2_mlp2/analysis/seed_level_results.json :: levels.ee_vs_mm.mlp[1].seed_sd :: sha256 53e1ec81343ac2c1
\newcommand{\ProbeOmaEeVsMmMlpOnesixtwo}{---}  % experiments/FIGS/data/probe_ladder_v2_mlp2/analysis/seed_level_results.json :: table[ee_vs_mm,mlp,162].auc :: sha256 53e1ec81343ac2c1 :: not quoted: the AUC reached 1 in a seed, so 1-AUC is only a bound
\newcommand{\ProbeBkgLeftEeVsMmMlpOnesixtwo}{0.0}  % experiments/FIGS/data/probe_ladder_v2_mlp2/analysis/seed_level_results.json :: levels.ee_vs_mm.mlp[1].rejection_points['0.90'].mean_n_bkg_pass :: sha256 53e1ec81343ac2c1 :: background jets passing the cut, mean over seeds
\newcommand{\ProbeRejEeVsMmMlpOnesixtwo}{\ensuremath{>}1{,}571}  % experiments/FIGS/data/probe_ladder_v2_mlp2/analysis/seed_level_results.json :: table[ee_vs_mm,mlp,162].rejection_points['0.90'].rejection :: sha256 53e1ec81343ac2c1 :: at 90% signal efficiency, mean +- SD over seeds; 5 of 5 seeds at the cap
\newcommand{\ProbeAucEeVsMmMlpFourthree}{0.99657}  % experiments/FIGS/data/probe_ladder_v2_mlp2/analysis/seed_level_results.json :: levels.ee_vs_mm.mlp[2].mean_auc :: sha256 53e1ec81343ac2c1 :: 5 seeds, 0 saturated at AUC=1
\newcommand{\ProbeAucSdEeVsMmMlpFourthree}{0.00300}  % experiments/FIGS/data/probe_ladder_v2_mlp2/analysis/seed_level_results.json :: table[ee_vs_mm,mlp,43].auc (SD over seeds, ddof=1) :: sha256 53e1ec81343ac2c1
\newcommand{\ProbeLogOneMinusAucEeVsMmMlpFourthree}{\ensuremath{-}6.0371}  % experiments/FIGS/data/probe_ladder_v2_mlp2/analysis/seed_level_results.json :: levels.ee_vs_mm.mlp[2].mean :: sha256 53e1ec81343ac2c1 :: natural log of 1-AUC, lower is better
\newcommand{\ProbeLogOneMinusAucSdEeVsMmMlpFourthree}{0.9962}  % experiments/FIGS/data/probe_ladder_v2_mlp2/analysis/seed_level_results.json :: levels.ee_vs_mm.mlp[2].seed_sd :: sha256 53e1ec81343ac2c1
\newcommand{\ProbeOmaEeVsMmMlpFourthree}{\ensuremath{(3.4 \pm 3.0)\times10^{-3}}}  % experiments/FIGS/data/probe_ladder_v2_mlp2/analysis/seed_level_results.json :: table[ee_vs_mm,mlp,43].auc :: sha256 53e1ec81343ac2c1 :: 1-AUC, mean +- SD over 5 seeds
\newcommand{\ProbeBkgLeftEeVsMmMlpFourthree}{33.4}  % experiments/FIGS/data/probe_ladder_v2_mlp2/analysis/seed_level_results.json :: levels.ee_vs_mm.mlp[2].rejection_points['0.90'].mean_n_bkg_pass :: sha256 53e1ec81343ac2c1 :: background jets passing the cut, mean over seeds
\newcommand{\ProbeRejEeVsMmMlpFourthree}{400\,\ensuremath{\pm}\,500}  % experiments/FIGS/data/probe_ladder_v2_mlp2/analysis/seed_level_results.json :: table[ee_vs_mm,mlp,43].rejection_points['0.90'].rejection :: sha256 53e1ec81343ac2c1 :: at 90% signal efficiency, mean +- SD over seeds; 0 of 5 seeds at the cap
\newcommand{\ProbeAucEeVsMmMlpOneseven}{0.99338}  % experiments/FIGS/data/probe_ladder_v2_mlp2/analysis/seed_level_results.json :: levels.ee_vs_mm.mlp[3].mean_auc :: sha256 53e1ec81343ac2c1 :: 5 seeds, 0 saturated at AUC=1
\newcommand{\ProbeAucSdEeVsMmMlpOneseven}{0.00065}  % experiments/FIGS/data/probe_ladder_v2_mlp2/analysis/seed_level_results.json :: table[ee_vs_mm,mlp,17].auc (SD over seeds, ddof=1) :: sha256 53e1ec81343ac2c1
\newcommand{\ProbeLogOneMinusAucEeVsMmMlpOneseven}{\ensuremath{-}5.0211}  % experiments/FIGS/data/probe_ladder_v2_mlp2/analysis/seed_level_results.json :: levels.ee_vs_mm.mlp[3].mean :: sha256 53e1ec81343ac2c1 :: natural log of 1-AUC, lower is better
\newcommand{\ProbeLogOneMinusAucSdEeVsMmMlpOneseven}{0.0997}  % experiments/FIGS/data/probe_ladder_v2_mlp2/analysis/seed_level_results.json :: levels.ee_vs_mm.mlp[3].seed_sd :: sha256 53e1ec81343ac2c1
\newcommand{\ProbeOmaEeVsMmMlpOneseven}{\ensuremath{(6.6 \pm 0.7)\times10^{-3}}}  % experiments/FIGS/data/probe_ladder_v2_mlp2/analysis/seed_level_results.json :: table[ee_vs_mm,mlp,17].auc :: sha256 53e1ec81343ac2c1 :: 1-AUC, mean +- SD over 5 seeds
\newcommand{\ProbeBkgLeftEeVsMmMlpOneseven}{59.2}  % experiments/FIGS/data/probe_ladder_v2_mlp2/analysis/seed_level_results.json :: levels.ee_vs_mm.mlp[3].rejection_points['0.90'].mean_n_bkg_pass :: sha256 53e1ec81343ac2c1 :: background jets passing the cut, mean over seeds
\newcommand{\ProbeRejEeVsMmMlpOneseven}{81\,\ensuremath{\pm}\,12}  % experiments/FIGS/data/probe_ladder_v2_mlp2/analysis/seed_level_results.json :: table[ee_vs_mm,mlp,17].rejection_points['0.90'].rejection :: sha256 53e1ec81343ac2c1 :: at 90% signal efficiency, mean +- SD over seeds; 0 of 5 seeds at the cap
\newcommand{\ProbeOmaRatioEeVsMmMlpOnesevenOverFourthree}{1.93}  % experiments/FIGS/data/probe_ladder_v2_mlp2/analysis/seed_level_results.json :: table[ee_vs_mm,mlp,17/43].auc :: sha256 53e1ec81343ac2c1 :: ratio of the seed means of 1-AUC, coarser over finer
\newcommand{\ProbeAucRetainedTopologyLinearOneeighteight}{0.99413}  % experiments/FIGS/data/probe_ladder_v2_mlp2/analysis/seed_level_results.json :: levels.retained_topology.linear[0].mean_auc :: sha256 53e1ec81343ac2c1 :: 5 seeds, 0 saturated at AUC=1
\newcommand{\ProbeAucSdRetainedTopologyLinearOneeighteight}{0.00011}  % experiments/FIGS/data/probe_ladder_v2_mlp2/analysis/seed_level_results.json :: table[retained_topology,linear,188].auc (SD over seeds, ddof=1) :: sha256 53e1ec81343ac2c1
\newcommand{\ProbeLogOneMinusAucRetainedTopologyLinearOneeighteight}{\ensuremath{-}5.1388}  % experiments/FIGS/data/probe_ladder_v2_mlp2/analysis/seed_level_results.json :: levels.retained_topology.linear[0].mean :: sha256 53e1ec81343ac2c1 :: natural log of 1-AUC, lower is better
\newcommand{\ProbeLogOneMinusAucSdRetainedTopologyLinearOneeighteight}{0.0187}  % experiments/FIGS/data/probe_ladder_v2_mlp2/analysis/seed_level_results.json :: levels.retained_topology.linear[0].seed_sd :: sha256 53e1ec81343ac2c1
\newcommand{\ProbeOmaRetainedTopologyLinearOneeighteight}{\ensuremath{(5.87 \pm 0.11)\times10^{-3}}}  % experiments/FIGS/data/probe_ladder_v2_mlp2/analysis/seed_level_results.json :: table[retained_topology,linear,188].auc :: sha256 53e1ec81343ac2c1 :: 1-AUC, mean +- SD over 5 seeds
\newcommand{\ProbeBkgLeftRetainedTopologyLinearOneeighteight}{15.4}  % experiments/FIGS/data/probe_ladder_v2_mlp2/analysis/seed_level_results.json :: levels.retained_topology.linear[0].rejection_points['0.90'].mean_n_bkg_pass :: sha256 53e1ec81343ac2c1 :: background jets passing the cut, mean over seeds
\newcommand{\ProbeRejRetainedTopologyLinearOneeighteight}{71\,\ensuremath{\pm}\,7}  % experiments/FIGS/data/probe_ladder_v2_mlp2/analysis/seed_level_results.json :: table[retained_topology,linear,188].rejection_points['0.90'].rejection :: sha256 53e1ec81343ac2c1 :: at 90% signal efficiency, mean +- SD over seeds; 0 of 5 seeds at the cap
\newcommand{\ProbeAucRetainedTopologyLinearOnesixtwo}{0.99413}  % experiments/FIGS/data/probe_ladder_v2_mlp2/analysis/seed_level_results.json :: levels.retained_topology.linear[1].mean_auc :: sha256 53e1ec81343ac2c1 :: 5 seeds, 0 saturated at AUC=1
\newcommand{\ProbeAucSdRetainedTopologyLinearOnesixtwo}{0.00030}  % experiments/FIGS/data/probe_ladder_v2_mlp2/analysis/seed_level_results.json :: table[retained_topology,linear,162].auc (SD over seeds, ddof=1) :: sha256 53e1ec81343ac2c1
\newcommand{\ProbeLogOneMinusAucRetainedTopologyLinearOnesixtwo}{\ensuremath{-}5.1395}  % experiments/FIGS/data/probe_ladder_v2_mlp2/analysis/seed_level_results.json :: levels.retained_topology.linear[1].mean :: sha256 53e1ec81343ac2c1 :: natural log of 1-AUC, lower is better
\newcommand{\ProbeLogOneMinusAucSdRetainedTopologyLinearOnesixtwo}{0.0510}  % experiments/FIGS/data/probe_ladder_v2_mlp2/analysis/seed_level_results.json :: levels.retained_topology.linear[1].seed_sd :: sha256 53e1ec81343ac2c1
\newcommand{\ProbeOmaRetainedTopologyLinearOnesixtwo}{\ensuremath{(5.87 \pm 0.30)\times10^{-3}}}  % experiments/FIGS/data/probe_ladder_v2_mlp2/analysis/seed_level_results.json :: table[retained_topology,linear,162].auc :: sha256 53e1ec81343ac2c1 :: 1-AUC, mean +- SD over 5 seeds
\newcommand{\ProbeBkgLeftRetainedTopologyLinearOnesixtwo}{14.0}  % experiments/FIGS/data/probe_ladder_v2_mlp2/analysis/seed_level_results.json :: levels.retained_topology.linear[1].rejection_points['0.90'].mean_n_bkg_pass :: sha256 53e1ec81343ac2c1 :: background jets passing the cut, mean over seeds
\newcommand{\ProbeRejRetainedTopologyLinearOnesixtwo}{79\,\ensuremath{\pm}\,13}  % experiments/FIGS/data/probe_ladder_v2_mlp2/analysis/seed_level_results.json :: table[retained_topology,linear,162].rejection_points['0.90'].rejection :: sha256 53e1ec81343ac2c1 :: at 90% signal efficiency, mean +- SD over seeds; 0 of 5 seeds at the cap
\newcommand{\ProbeAucRetainedTopologyLinearFourthree}{0.99274}  % experiments/FIGS/data/probe_ladder_v2_mlp2/analysis/seed_level_results.json :: levels.retained_topology.linear[2].mean_auc :: sha256 53e1ec81343ac2c1 :: 5 seeds, 0 saturated at AUC=1
\newcommand{\ProbeAucSdRetainedTopologyLinearFourthree}{0.00012}  % experiments/FIGS/data/probe_ladder_v2_mlp2/analysis/seed_level_results.json :: table[retained_topology,linear,43].auc (SD over seeds, ddof=1) :: sha256 53e1ec81343ac2c1
\newcommand{\ProbeLogOneMinusAucRetainedTopologyLinearFourthree}{\ensuremath{-}4.9253}  % experiments/FIGS/data/probe_ladder_v2_mlp2/analysis/seed_level_results.json :: levels.retained_topology.linear[2].mean :: sha256 53e1ec81343ac2c1 :: natural log of 1-AUC, lower is better
\newcommand{\ProbeLogOneMinusAucSdRetainedTopologyLinearFourthree}{0.0169}  % experiments/FIGS/data/probe_ladder_v2_mlp2/analysis/seed_level_results.json :: levels.retained_topology.linear[2].seed_sd :: sha256 53e1ec81343ac2c1
\newcommand{\ProbeOmaRetainedTopologyLinearFourthree}{\ensuremath{(7.26 \pm 0.12)\times10^{-3}}}  % experiments/FIGS/data/probe_ladder_v2_mlp2/analysis/seed_level_results.json :: table[retained_topology,linear,43].auc :: sha256 53e1ec81343ac2c1 :: 1-AUC, mean +- SD over 5 seeds
\newcommand{\ProbeBkgLeftRetainedTopologyLinearFourthree}{16.4}  % experiments/FIGS/data/probe_ladder_v2_mlp2/analysis/seed_level_results.json :: levels.retained_topology.linear[2].rejection_points['0.90'].mean_n_bkg_pass :: sha256 53e1ec81343ac2c1 :: background jets passing the cut, mean over seeds
\newcommand{\ProbeRejRetainedTopologyLinearFourthree}{66\,\ensuremath{\pm}\,5}  % experiments/FIGS/data/probe_ladder_v2_mlp2/analysis/seed_level_results.json :: table[retained_topology,linear,43].rejection_points['0.90'].rejection :: sha256 53e1ec81343ac2c1 :: at 90% signal efficiency, mean +- SD over seeds; 0 of 5 seeds at the cap
\newcommand{\ProbeAucRetainedTopologyLinearOneseven}{0.99096}  % experiments/FIGS/data/probe_ladder_v2_mlp2/analysis/seed_level_results.json :: levels.retained_topology.linear[3].mean_auc :: sha256 53e1ec81343ac2c1 :: 5 seeds, 0 saturated at AUC=1
\newcommand{\ProbeAucSdRetainedTopologyLinearOneseven}{0.00092}  % experiments/FIGS/data/probe_ladder_v2_mlp2/analysis/seed_level_results.json :: table[retained_topology,linear,17].auc (SD over seeds, ddof=1) :: sha256 53e1ec81343ac2c1
\newcommand{\ProbeLogOneMinusAucRetainedTopologyLinearOneseven}{\ensuremath{-}4.7095}  % experiments/FIGS/data/probe_ladder_v2_mlp2/analysis/seed_level_results.json :: levels.retained_topology.linear[3].mean :: sha256 53e1ec81343ac2c1 :: natural log of 1-AUC, lower is better
\newcommand{\ProbeLogOneMinusAucSdRetainedTopologyLinearOneseven}{0.0979}  % experiments/FIGS/data/probe_ladder_v2_mlp2/analysis/seed_level_results.json :: levels.retained_topology.linear[3].seed_sd :: sha256 53e1ec81343ac2c1
\newcommand{\ProbeOmaRetainedTopologyLinearOneseven}{\ensuremath{(9.0 \pm 0.9)\times10^{-3}}}  % experiments/FIGS/data/probe_ladder_v2_mlp2/analysis/seed_level_results.json :: table[retained_topology,linear,17].auc :: sha256 53e1ec81343ac2c1 :: 1-AUC, mean +- SD over 5 seeds
\newcommand{\ProbeBkgLeftRetainedTopologyLinearOneseven}{21.8}  % experiments/FIGS/data/probe_ladder_v2_mlp2/analysis/seed_level_results.json :: levels.retained_topology.linear[3].rejection_points['0.90'].mean_n_bkg_pass :: sha256 53e1ec81343ac2c1 :: background jets passing the cut, mean over seeds
\newcommand{\ProbeRejRetainedTopologyLinearOneseven}{51\,\ensuremath{\pm}\,11}  % experiments/FIGS/data/probe_ladder_v2_mlp2/analysis/seed_level_results.json :: table[retained_topology,linear,17].rejection_points['0.90'].rejection :: sha256 53e1ec81343ac2c1 :: at 90% signal efficiency, mean +- SD over seeds; 0 of 5 seeds at the cap
\newcommand{\ProbeOmaRatioRetainedTopologyLinearOnesixtwoOverOneeighteight}{1.00}  % experiments/FIGS/data/probe_ladder_v2_mlp2/analysis/seed_level_results.json :: table[retained_topology,linear,162/188].auc :: sha256 53e1ec81343ac2c1 :: ratio of the seed means of 1-AUC, coarser over finer
\newcommand{\ProbeOmaRatioRetainedTopologyLinearFourthreeOverOneeighteight}{1.24}  % experiments/FIGS/data/probe_ladder_v2_mlp2/analysis/seed_level_results.json :: table[retained_topology,linear,43/188].auc :: sha256 53e1ec81343ac2c1 :: ratio of the seed means of 1-AUC, coarser over finer
\newcommand{\ProbeOmaRatioRetainedTopologyLinearOnesevenOverOneeighteight}{1.54}  % experiments/FIGS/data/probe_ladder_v2_mlp2/analysis/seed_level_results.json :: table[retained_topology,linear,17/188].auc :: sha256 53e1ec81343ac2c1 :: ratio of the seed means of 1-AUC, coarser over finer
\newcommand{\ProbeOmaRatioRetainedTopologyLinearFourthreeOverOnesixtwo}{1.24}  % experiments/FIGS/data/probe_ladder_v2_mlp2/analysis/seed_level_results.json :: table[retained_topology,linear,43/162].auc :: sha256 53e1ec81343ac2c1 :: ratio of the seed means of 1-AUC, coarser over finer
\newcommand{\ProbeOmaRatioRetainedTopologyLinearOnesevenOverOnesixtwo}{1.54}  % experiments/FIGS/data/probe_ladder_v2_mlp2/analysis/seed_level_results.json :: table[retained_topology,linear,17/162].auc :: sha256 53e1ec81343ac2c1 :: ratio of the seed means of 1-AUC, coarser over finer
\newcommand{\ProbeOmaRatioRetainedTopologyLinearOnesevenOverFourthree}{1.25}  % experiments/FIGS/data/probe_ladder_v2_mlp2/analysis/seed_level_results.json :: table[retained_topology,linear,17/43].auc :: sha256 53e1ec81343ac2c1 :: ratio of the seed means of 1-AUC, coarser over finer
\newcommand{\ProbeAucRetainedTopologyMlpOneeighteight}{0.99394}  % experiments/FIGS/data/probe_ladder_v2_mlp2/analysis/seed_level_results.json :: levels.retained_topology.mlp[0].mean_auc :: sha256 53e1ec81343ac2c1 :: 5 seeds, 0 saturated at AUC=1
\newcommand{\ProbeAucSdRetainedTopologyMlpOneeighteight}{0.00032}  % experiments/FIGS/data/probe_ladder_v2_mlp2/analysis/seed_level_results.json :: table[retained_topology,mlp,188].auc (SD over seeds, ddof=1) :: sha256 53e1ec81343ac2c1
\newcommand{\ProbeLogOneMinusAucRetainedTopologyMlpOneeighteight}{\ensuremath{-}5.1078}  % experiments/FIGS/data/probe_ladder_v2_mlp2/analysis/seed_level_results.json :: levels.retained_topology.mlp[0].mean :: sha256 53e1ec81343ac2c1 :: natural log of 1-AUC, lower is better
\newcommand{\ProbeLogOneMinusAucSdRetainedTopologyMlpOneeighteight}{0.0522}  % experiments/FIGS/data/probe_ladder_v2_mlp2/analysis/seed_level_results.json :: levels.retained_topology.mlp[0].seed_sd :: sha256 53e1ec81343ac2c1
\newcommand{\ProbeOmaRetainedTopologyMlpOneeighteight}{\ensuremath{(6.06 \pm 0.32)\times10^{-3}}}  % experiments/FIGS/data/probe_ladder_v2_mlp2/analysis/seed_level_results.json :: table[retained_topology,mlp,188].auc :: sha256 53e1ec81343ac2c1 :: 1-AUC, mean +- SD over 5 seeds
\newcommand{\ProbeBkgLeftRetainedTopologyMlpOneeighteight}{15.8}  % experiments/FIGS/data/probe_ladder_v2_mlp2/analysis/seed_level_results.json :: levels.retained_topology.mlp[0].rejection_points['0.90'].mean_n_bkg_pass :: sha256 53e1ec81343ac2c1 :: background jets passing the cut, mean over seeds
\newcommand{\ProbeRejRetainedTopologyMlpOneeighteight}{69\,\ensuremath{\pm}\,9}  % experiments/FIGS/data/probe_ladder_v2_mlp2/analysis/seed_level_results.json :: table[retained_topology,mlp,188].rejection_points['0.90'].rejection :: sha256 53e1ec81343ac2c1 :: at 90% signal efficiency, mean +- SD over seeds; 0 of 5 seeds at the cap
\newcommand{\ProbeAucRetainedTopologyMlpOnesixtwo}{0.99387}  % experiments/FIGS/data/probe_ladder_v2_mlp2/analysis/seed_level_results.json :: levels.retained_topology.mlp[1].mean_auc :: sha256 53e1ec81343ac2c1 :: 5 seeds, 0 saturated at AUC=1
\newcommand{\ProbeAucSdRetainedTopologyMlpOnesixtwo}{0.00048}  % experiments/FIGS/data/probe_ladder_v2_mlp2/analysis/seed_level_results.json :: table[retained_topology,mlp,162].auc (SD over seeds, ddof=1) :: sha256 53e1ec81343ac2c1
\newcommand{\ProbeLogOneMinusAucRetainedTopologyMlpOnesixtwo}{\ensuremath{-}5.0964}  % experiments/FIGS/data/probe_ladder_v2_mlp2/analysis/seed_level_results.json :: levels.retained_topology.mlp[1].mean :: sha256 53e1ec81343ac2c1 :: natural log of 1-AUC, lower is better
\newcommand{\ProbeLogOneMinusAucSdRetainedTopologyMlpOnesixtwo}{0.0773}  % experiments/FIGS/data/probe_ladder_v2_mlp2/analysis/seed_level_results.json :: levels.retained_topology.mlp[1].seed_sd :: sha256 53e1ec81343ac2c1
\newcommand{\ProbeOmaRetainedTopologyMlpOnesixtwo}{\ensuremath{(6.1 \pm 0.5)\times10^{-3}}}  % experiments/FIGS/data/probe_ladder_v2_mlp2/analysis/seed_level_results.json :: table[retained_topology,mlp,162].auc :: sha256 53e1ec81343ac2c1 :: 1-AUC, mean +- SD over 5 seeds
\newcommand{\ProbeBkgLeftRetainedTopologyMlpOnesixtwo}{16.6}  % experiments/FIGS/data/probe_ladder_v2_mlp2/analysis/seed_level_results.json :: levels.retained_topology.mlp[1].rejection_points['0.90'].mean_n_bkg_pass :: sha256 53e1ec81343ac2c1 :: background jets passing the cut, mean over seeds
\newcommand{\ProbeRejRetainedTopologyMlpOnesixtwo}{67\,\ensuremath{\pm}\,14}  % experiments/FIGS/data/probe_ladder_v2_mlp2/analysis/seed_level_results.json :: table[retained_topology,mlp,162].rejection_points['0.90'].rejection :: sha256 53e1ec81343ac2c1 :: at 90% signal efficiency, mean +- SD over seeds; 0 of 5 seeds at the cap
\newcommand{\ProbeAucRetainedTopologyMlpFourthree}{0.99284}  % experiments/FIGS/data/probe_ladder_v2_mlp2/analysis/seed_level_results.json :: levels.retained_topology.mlp[2].mean_auc :: sha256 53e1ec81343ac2c1 :: 5 seeds, 0 saturated at AUC=1
\newcommand{\ProbeAucSdRetainedTopologyMlpFourthree}{0.00040}  % experiments/FIGS/data/probe_ladder_v2_mlp2/analysis/seed_level_results.json :: table[retained_topology,mlp,43].auc (SD over seeds, ddof=1) :: sha256 53e1ec81343ac2c1
\newcommand{\ProbeLogOneMinusAucRetainedTopologyMlpFourthree}{\ensuremath{-}4.9400}  % experiments/FIGS/data/probe_ladder_v2_mlp2/analysis/seed_level_results.json :: levels.retained_topology.mlp[2].mean :: sha256 53e1ec81343ac2c1 :: natural log of 1-AUC, lower is better
\newcommand{\ProbeLogOneMinusAucSdRetainedTopologyMlpFourthree}{0.0570}  % experiments/FIGS/data/probe_ladder_v2_mlp2/analysis/seed_level_results.json :: levels.retained_topology.mlp[2].seed_sd :: sha256 53e1ec81343ac2c1
\newcommand{\ProbeOmaRetainedTopologyMlpFourthree}{\ensuremath{(7.2 \pm 0.4)\times10^{-3}}}  % experiments/FIGS/data/probe_ladder_v2_mlp2/analysis/seed_level_results.json :: table[retained_topology,mlp,43].auc :: sha256 53e1ec81343ac2c1 :: 1-AUC, mean +- SD over 5 seeds
\newcommand{\ProbeBkgLeftRetainedTopologyMlpFourthree}{17.8}  % experiments/FIGS/data/probe_ladder_v2_mlp2/analysis/seed_level_results.json :: levels.retained_topology.mlp[2].rejection_points['0.90'].mean_n_bkg_pass :: sha256 53e1ec81343ac2c1 :: background jets passing the cut, mean over seeds
\newcommand{\ProbeRejRetainedTopologyMlpFourthree}{61\,\ensuremath{\pm}\,7}  % experiments/FIGS/data/probe_ladder_v2_mlp2/analysis/seed_level_results.json :: table[retained_topology,mlp,43].rejection_points['0.90'].rejection :: sha256 53e1ec81343ac2c1 :: at 90% signal efficiency, mean +- SD over seeds; 0 of 5 seeds at the cap
\newcommand{\ProbeAucRetainedTopologyMlpOneseven}{0.99081}  % experiments/FIGS/data/probe_ladder_v2_mlp2/analysis/seed_level_results.json :: levels.retained_topology.mlp[3].mean_auc :: sha256 53e1ec81343ac2c1 :: 5 seeds, 0 saturated at AUC=1
\newcommand{\ProbeAucSdRetainedTopologyMlpOneseven}{0.00062}  % experiments/FIGS/data/probe_ladder_v2_mlp2/analysis/seed_level_results.json :: table[retained_topology,mlp,17].auc (SD over seeds, ddof=1) :: sha256 53e1ec81343ac2c1
\newcommand{\ProbeLogOneMinusAucRetainedTopologyMlpOneseven}{\ensuremath{-}4.6913}  % experiments/FIGS/data/probe_ladder_v2_mlp2/analysis/seed_level_results.json :: levels.retained_topology.mlp[3].mean :: sha256 53e1ec81343ac2c1 :: natural log of 1-AUC, lower is better
\newcommand{\ProbeLogOneMinusAucSdRetainedTopologyMlpOneseven}{0.0658}  % experiments/FIGS/data/probe_ladder_v2_mlp2/analysis/seed_level_results.json :: levels.retained_topology.mlp[3].seed_sd :: sha256 53e1ec81343ac2c1
\newcommand{\ProbeOmaRetainedTopologyMlpOneseven}{\ensuremath{(9.2 \pm 0.6)\times10^{-3}}}  % experiments/FIGS/data/probe_ladder_v2_mlp2/analysis/seed_level_results.json :: table[retained_topology,mlp,17].auc :: sha256 53e1ec81343ac2c1 :: 1-AUC, mean +- SD over 5 seeds
\newcommand{\ProbeBkgLeftRetainedTopologyMlpOneseven}{21.2}  % experiments/FIGS/data/probe_ladder_v2_mlp2/analysis/seed_level_results.json :: levels.retained_topology.mlp[3].rejection_points['0.90'].mean_n_bkg_pass :: sha256 53e1ec81343ac2c1 :: background jets passing the cut, mean over seeds
\newcommand{\ProbeRejRetainedTopologyMlpOneseven}{51\,\ensuremath{\pm}\,4}  % experiments/FIGS/data/probe_ladder_v2_mlp2/analysis/seed_level_results.json :: table[retained_topology,mlp,17].rejection_points['0.90'].rejection :: sha256 53e1ec81343ac2c1 :: at 90% signal efficiency, mean +- SD over seeds; 0 of 5 seeds at the cap
\newcommand{\ProbeOmaRatioRetainedTopologyMlpOnesixtwoOverOneeighteight}{1.01}  % experiments/FIGS/data/probe_ladder_v2_mlp2/analysis/seed_level_results.json :: table[retained_topology,mlp,162/188].auc :: sha256 53e1ec81343ac2c1 :: ratio of the seed means of 1-AUC, coarser over finer
\newcommand{\ProbeOmaRatioRetainedTopologyMlpFourthreeOverOneeighteight}{1.18}  % experiments/FIGS/data/probe_ladder_v2_mlp2/analysis/seed_level_results.json :: table[retained_topology,mlp,43/188].auc :: sha256 53e1ec81343ac2c1 :: ratio of the seed means of 1-AUC, coarser over finer
\newcommand{\ProbeOmaRatioRetainedTopologyMlpOnesevenOverOneeighteight}{1.52}  % experiments/FIGS/data/probe_ladder_v2_mlp2/analysis/seed_level_results.json :: table[retained_topology,mlp,17/188].auc :: sha256 53e1ec81343ac2c1 :: ratio of the seed means of 1-AUC, coarser over finer
\newcommand{\ProbeOmaRatioRetainedTopologyMlpFourthreeOverOnesixtwo}{1.17}  % experiments/FIGS/data/probe_ladder_v2_mlp2/analysis/seed_level_results.json :: table[retained_topology,mlp,43/162].auc :: sha256 53e1ec81343ac2c1 :: ratio of the seed means of 1-AUC, coarser over finer
\newcommand{\ProbeOmaRatioRetainedTopologyMlpOnesevenOverOnesixtwo}{1.50}  % experiments/FIGS/data/probe_ladder_v2_mlp2/analysis/seed_level_results.json :: table[retained_topology,mlp,17/162].auc :: sha256 53e1ec81343ac2c1 :: ratio of the seed means of 1-AUC, coarser over finer
\newcommand{\ProbeOmaRatioRetainedTopologyMlpOnesevenOverFourthree}{1.28}  % experiments/FIGS/data/probe_ladder_v2_mlp2/analysis/seed_level_results.json :: table[retained_topology,mlp,17/43].auc :: sha256 53e1ec81343ac2c1 :: ratio of the seed means of 1-AUC, coarser over finer
\newcommand{\ProbeAucVisibleContentLinearOneeighteight}{0.99917}  % experiments/FIGS/data/probe_ladder_v2_mlp2/analysis/seed_level_results.json :: levels.visible_content.linear[0].mean_auc :: sha256 53e1ec81343ac2c1 :: 5 seeds, 0 saturated at AUC=1
\newcommand{\ProbeAucSdVisibleContentLinearOneeighteight}{0.00003}  % experiments/FIGS/data/probe_ladder_v2_mlp2/analysis/seed_level_results.json :: table[visible_content,linear,188].auc (SD over seeds, ddof=1) :: sha256 53e1ec81343ac2c1
\newcommand{\ProbeLogOneMinusAucVisibleContentLinearOneeighteight}{\ensuremath{-}7.0999}  % experiments/FIGS/data/probe_ladder_v2_mlp2/analysis/seed_level_results.json :: levels.visible_content.linear[0].mean :: sha256 53e1ec81343ac2c1 :: natural log of 1-AUC, lower is better
\newcommand{\ProbeLogOneMinusAucSdVisibleContentLinearOneeighteight}{0.0358}  % experiments/FIGS/data/probe_ladder_v2_mlp2/analysis/seed_level_results.json :: levels.visible_content.linear[0].seed_sd :: sha256 53e1ec81343ac2c1
\newcommand{\ProbeOmaVisibleContentLinearOneeighteight}{\ensuremath{(8.26 \pm 0.30)\times10^{-4}}}  % experiments/FIGS/data/probe_ladder_v2_mlp2/analysis/seed_level_results.json :: table[visible_content,linear,188].auc :: sha256 53e1ec81343ac2c1 :: 1-AUC, mean +- SD over 5 seeds
\newcommand{\ProbeBkgLeftVisibleContentLinearOneeighteight}{0.4}  % experiments/FIGS/data/probe_ladder_v2_mlp2/analysis/seed_level_results.json :: levels.visible_content.linear[0].rejection_points['0.90'].mean_n_bkg_pass :: sha256 53e1ec81343ac2c1 :: background jets passing the cut, mean over seeds
\newcommand{\ProbeRejVisibleContentLinearOneeighteight}{\ensuremath{>}538}  % experiments/FIGS/data/probe_ladder_v2_mlp2/analysis/seed_level_results.json :: table[visible_content,linear,188].rejection_points['0.90'].rejection :: sha256 53e1ec81343ac2c1 :: at 90% signal efficiency, mean +- SD over seeds; 5 of 5 seeds at the cap
\newcommand{\ProbeAucVisibleContentLinearOnesixtwo}{0.99921}  % experiments/FIGS/data/probe_ladder_v2_mlp2/analysis/seed_level_results.json :: levels.visible_content.linear[1].mean_auc :: sha256 53e1ec81343ac2c1 :: 5 seeds, 0 saturated at AUC=1
\newcommand{\ProbeAucSdVisibleContentLinearOnesixtwo}{0.00004}  % experiments/FIGS/data/probe_ladder_v2_mlp2/analysis/seed_level_results.json :: table[visible_content,linear,162].auc (SD over seeds, ddof=1) :: sha256 53e1ec81343ac2c1
\newcommand{\ProbeLogOneMinusAucVisibleContentLinearOnesixtwo}{\ensuremath{-}7.1480}  % experiments/FIGS/data/probe_ladder_v2_mlp2/analysis/seed_level_results.json :: levels.visible_content.linear[1].mean :: sha256 53e1ec81343ac2c1 :: natural log of 1-AUC, lower is better
\newcommand{\ProbeLogOneMinusAucSdVisibleContentLinearOnesixtwo}{0.0449}  % experiments/FIGS/data/probe_ladder_v2_mlp2/analysis/seed_level_results.json :: levels.visible_content.linear[1].seed_sd :: sha256 53e1ec81343ac2c1
\newcommand{\ProbeOmaVisibleContentLinearOnesixtwo}{\ensuremath{(7.87 \pm 0.35)\times10^{-4}}}  % experiments/FIGS/data/probe_ladder_v2_mlp2/analysis/seed_level_results.json :: table[visible_content,linear,162].auc :: sha256 53e1ec81343ac2c1 :: 1-AUC, mean +- SD over 5 seeds
\newcommand{\ProbeBkgLeftVisibleContentLinearOnesixtwo}{0.6}  % experiments/FIGS/data/probe_ladder_v2_mlp2/analysis/seed_level_results.json :: levels.visible_content.linear[1].rejection_points['0.90'].mean_n_bkg_pass :: sha256 53e1ec81343ac2c1 :: background jets passing the cut, mean over seeds
\newcommand{\ProbeRejVisibleContentLinearOnesixtwo}{\ensuremath{>}538}  % experiments/FIGS/data/probe_ladder_v2_mlp2/analysis/seed_level_results.json :: table[visible_content,linear,162].rejection_points['0.90'].rejection :: sha256 53e1ec81343ac2c1 :: at 90% signal efficiency, mean +- SD over seeds; 5 of 5 seeds at the cap
\newcommand{\ProbeAucVisibleContentLinearFourthree}{0.99918}  % experiments/FIGS/data/probe_ladder_v2_mlp2/analysis/seed_level_results.json :: levels.visible_content.linear[2].mean_auc :: sha256 53e1ec81343ac2c1 :: 5 seeds, 0 saturated at AUC=1
\newcommand{\ProbeAucSdVisibleContentLinearFourthree}{0.00009}  % experiments/FIGS/data/probe_ladder_v2_mlp2/analysis/seed_level_results.json :: table[visible_content,linear,43].auc (SD over seeds, ddof=1) :: sha256 53e1ec81343ac2c1
\newcommand{\ProbeLogOneMinusAucVisibleContentLinearFourthree}{\ensuremath{-}7.1124}  % experiments/FIGS/data/probe_ladder_v2_mlp2/analysis/seed_level_results.json :: levels.visible_content.linear[2].mean :: sha256 53e1ec81343ac2c1 :: natural log of 1-AUC, lower is better
\newcommand{\ProbeLogOneMinusAucSdVisibleContentLinearFourthree}{0.1137}  % experiments/FIGS/data/probe_ladder_v2_mlp2/analysis/seed_level_results.json :: levels.visible_content.linear[2].seed_sd :: sha256 53e1ec81343ac2c1
\newcommand{\ProbeOmaVisibleContentLinearFourthree}{\ensuremath{(8.2 \pm 1.0)\times10^{-4}}}  % experiments/FIGS/data/probe_ladder_v2_mlp2/analysis/seed_level_results.json :: table[visible_content,linear,43].auc :: sha256 53e1ec81343ac2c1 :: 1-AUC, mean +- SD over 5 seeds
\newcommand{\ProbeBkgLeftVisibleContentLinearFourthree}{0.4}  % experiments/FIGS/data/probe_ladder_v2_mlp2/analysis/seed_level_results.json :: levels.visible_content.linear[2].rejection_points['0.90'].mean_n_bkg_pass :: sha256 53e1ec81343ac2c1 :: background jets passing the cut, mean over seeds
\newcommand{\ProbeRejVisibleContentLinearFourthree}{\ensuremath{>}538}  % experiments/FIGS/data/probe_ladder_v2_mlp2/analysis/seed_level_results.json :: table[visible_content,linear,43].rejection_points['0.90'].rejection :: sha256 53e1ec81343ac2c1 :: at 90% signal efficiency, mean +- SD over seeds; 5 of 5 seeds at the cap
\newcommand{\ProbeAucVisibleContentLinearOneseven}{0.99768}  % experiments/FIGS/data/probe_ladder_v2_mlp2/analysis/seed_level_results.json :: levels.visible_content.linear[3].mean_auc :: sha256 53e1ec81343ac2c1 :: 5 seeds, 0 saturated at AUC=1
\newcommand{\ProbeAucSdVisibleContentLinearOneseven}{0.00020}  % experiments/FIGS/data/probe_ladder_v2_mlp2/analysis/seed_level_results.json :: table[visible_content,linear,17].auc (SD over seeds, ddof=1) :: sha256 53e1ec81343ac2c1
\newcommand{\ProbeLogOneMinusAucVisibleContentLinearOneseven}{\ensuremath{-}6.0684}  % experiments/FIGS/data/probe_ladder_v2_mlp2/analysis/seed_level_results.json :: levels.visible_content.linear[3].mean :: sha256 53e1ec81343ac2c1 :: natural log of 1-AUC, lower is better
\newcommand{\ProbeLogOneMinusAucSdVisibleContentLinearOneseven}{0.0889}  % experiments/FIGS/data/probe_ladder_v2_mlp2/analysis/seed_level_results.json :: levels.visible_content.linear[3].seed_sd :: sha256 53e1ec81343ac2c1
\newcommand{\ProbeOmaVisibleContentLinearOneseven}{\ensuremath{(2.32 \pm 0.20)\times10^{-3}}}  % experiments/FIGS/data/probe_ladder_v2_mlp2/analysis/seed_level_results.json :: table[visible_content,linear,17].auc :: sha256 53e1ec81343ac2c1 :: 1-AUC, mean +- SD over 5 seeds
\newcommand{\ProbeBkgLeftVisibleContentLinearOneseven}{3.4}  % experiments/FIGS/data/probe_ladder_v2_mlp2/analysis/seed_level_results.json :: levels.visible_content.linear[3].rejection_points['0.90'].mean_n_bkg_pass :: sha256 53e1ec81343ac2c1 :: background jets passing the cut, mean over seeds
\newcommand{\ProbeRejVisibleContentLinearOneseven}{900\,\ensuremath{\pm}\,400}  % experiments/FIGS/data/probe_ladder_v2_mlp2/analysis/seed_level_results.json :: table[visible_content,linear,17].rejection_points['0.90'].rejection :: sha256 53e1ec81343ac2c1 :: at 90% signal efficiency, mean +- SD over seeds; 0 of 5 seeds at the cap
\newcommand{\ProbeOmaRatioVisibleContentLinearOnesixtwoOverOneeighteight}{0.95}  % experiments/FIGS/data/probe_ladder_v2_mlp2/analysis/seed_level_results.json :: table[visible_content,linear,162/188].auc :: sha256 53e1ec81343ac2c1 :: ratio of the seed means of 1-AUC, coarser over finer
\newcommand{\ProbeOmaRatioVisibleContentLinearFourthreeOverOneeighteight}{0.99}  % experiments/FIGS/data/probe_ladder_v2_mlp2/analysis/seed_level_results.json :: table[visible_content,linear,43/188].auc :: sha256 53e1ec81343ac2c1 :: ratio of the seed means of 1-AUC, coarser over finer
\newcommand{\ProbeOmaRatioVisibleContentLinearOnesevenOverOneeighteight}{2.8}  % experiments/FIGS/data/probe_ladder_v2_mlp2/analysis/seed_level_results.json :: table[visible_content,linear,17/188].auc :: sha256 53e1ec81343ac2c1 :: ratio of the seed means of 1-AUC, coarser over finer
\newcommand{\ProbeOmaRatioVisibleContentLinearFourthreeOverOnesixtwo}{1.04}  % experiments/FIGS/data/probe_ladder_v2_mlp2/analysis/seed_level_results.json :: table[visible_content,linear,43/162].auc :: sha256 53e1ec81343ac2c1 :: ratio of the seed means of 1-AUC, coarser over finer
\newcommand{\ProbeOmaRatioVisibleContentLinearOnesevenOverOnesixtwo}{3.0}  % experiments/FIGS/data/probe_ladder_v2_mlp2/analysis/seed_level_results.json :: table[visible_content,linear,17/162].auc :: sha256 53e1ec81343ac2c1 :: ratio of the seed means of 1-AUC, coarser over finer
\newcommand{\ProbeOmaRatioVisibleContentLinearOnesevenOverFourthree}{2.8}  % experiments/FIGS/data/probe_ladder_v2_mlp2/analysis/seed_level_results.json :: table[visible_content,linear,17/43].auc :: sha256 53e1ec81343ac2c1 :: ratio of the seed means of 1-AUC, coarser over finer
\newcommand{\ProbeAucVisibleContentMlpOneeighteight}{0.99911}  % experiments/FIGS/data/probe_ladder_v2_mlp2/analysis/seed_level_results.json :: levels.visible_content.mlp[0].mean_auc :: sha256 53e1ec81343ac2c1 :: 5 seeds, 0 saturated at AUC=1
\newcommand{\ProbeAucSdVisibleContentMlpOneeighteight}{0.00005}  % experiments/FIGS/data/probe_ladder_v2_mlp2/analysis/seed_level_results.json :: table[visible_content,mlp,188].auc (SD over seeds, ddof=1) :: sha256 53e1ec81343ac2c1
\newcommand{\ProbeLogOneMinusAucVisibleContentMlpOneeighteight}{\ensuremath{-}7.0202}  % experiments/FIGS/data/probe_ladder_v2_mlp2/analysis/seed_level_results.json :: levels.visible_content.mlp[0].mean :: sha256 53e1ec81343ac2c1 :: natural log of 1-AUC, lower is better
\newcommand{\ProbeLogOneMinusAucSdVisibleContentMlpOneeighteight}{0.0570}  % experiments/FIGS/data/probe_ladder_v2_mlp2/analysis/seed_level_results.json :: levels.visible_content.mlp[0].seed_sd :: sha256 53e1ec81343ac2c1
\newcommand{\ProbeOmaVisibleContentMlpOneeighteight}{\ensuremath{(8.9 \pm 0.5)\times10^{-4}}}  % experiments/FIGS/data/probe_ladder_v2_mlp2/analysis/seed_level_results.json :: table[visible_content,mlp,188].auc :: sha256 53e1ec81343ac2c1 :: 1-AUC, mean +- SD over 5 seeds
\newcommand{\ProbeBkgLeftVisibleContentMlpOneeighteight}{0.4}  % experiments/FIGS/data/probe_ladder_v2_mlp2/analysis/seed_level_results.json :: levels.visible_content.mlp[0].rejection_points['0.90'].mean_n_bkg_pass :: sha256 53e1ec81343ac2c1 :: background jets passing the cut, mean over seeds
\newcommand{\ProbeRejVisibleContentMlpOneeighteight}{\ensuremath{>}538}  % experiments/FIGS/data/probe_ladder_v2_mlp2/analysis/seed_level_results.json :: table[visible_content,mlp,188].rejection_points['0.90'].rejection :: sha256 53e1ec81343ac2c1 :: at 90% signal efficiency, mean +- SD over seeds; 5 of 5 seeds at the cap
\newcommand{\ProbeAucVisibleContentMlpOnesixtwo}{0.99918}  % experiments/FIGS/data/probe_ladder_v2_mlp2/analysis/seed_level_results.json :: levels.visible_content.mlp[1].mean_auc :: sha256 53e1ec81343ac2c1 :: 5 seeds, 0 saturated at AUC=1
\newcommand{\ProbeAucSdVisibleContentMlpOnesixtwo}{0.00003}  % experiments/FIGS/data/probe_ladder_v2_mlp2/analysis/seed_level_results.json :: table[visible_content,mlp,162].auc (SD over seeds, ddof=1) :: sha256 53e1ec81343ac2c1
\newcommand{\ProbeLogOneMinusAucVisibleContentMlpOnesixtwo}{\ensuremath{-}7.1010}  % experiments/FIGS/data/probe_ladder_v2_mlp2/analysis/seed_level_results.json :: levels.visible_content.mlp[1].mean :: sha256 53e1ec81343ac2c1 :: natural log of 1-AUC, lower is better
\newcommand{\ProbeLogOneMinusAucSdVisibleContentMlpOnesixtwo}{0.0371}  % experiments/FIGS/data/probe_ladder_v2_mlp2/analysis/seed_level_results.json :: levels.visible_content.mlp[1].seed_sd :: sha256 53e1ec81343ac2c1
\newcommand{\ProbeOmaVisibleContentMlpOnesixtwo}{\ensuremath{(8.25 \pm 0.30)\times10^{-4}}}  % experiments/FIGS/data/probe_ladder_v2_mlp2/analysis/seed_level_results.json :: table[visible_content,mlp,162].auc :: sha256 53e1ec81343ac2c1 :: 1-AUC, mean +- SD over 5 seeds
\newcommand{\ProbeBkgLeftVisibleContentMlpOnesixtwo}{0.2}  % experiments/FIGS/data/probe_ladder_v2_mlp2/analysis/seed_level_results.json :: levels.visible_content.mlp[1].rejection_points['0.90'].mean_n_bkg_pass :: sha256 53e1ec81343ac2c1 :: background jets passing the cut, mean over seeds
\newcommand{\ProbeRejVisibleContentMlpOnesixtwo}{\ensuremath{>}538}  % experiments/FIGS/data/probe_ladder_v2_mlp2/analysis/seed_level_results.json :: table[visible_content,mlp,162].rejection_points['0.90'].rejection :: sha256 53e1ec81343ac2c1 :: at 90% signal efficiency, mean +- SD over seeds; 5 of 5 seeds at the cap
\newcommand{\ProbeAucVisibleContentMlpFourthree}{0.99917}  % experiments/FIGS/data/probe_ladder_v2_mlp2/analysis/seed_level_results.json :: levels.visible_content.mlp[2].mean_auc :: sha256 53e1ec81343ac2c1 :: 5 seeds, 0 saturated at AUC=1
\newcommand{\ProbeAucSdVisibleContentMlpFourthree}{0.00006}  % experiments/FIGS/data/probe_ladder_v2_mlp2/analysis/seed_level_results.json :: table[visible_content,mlp,43].auc (SD over seeds, ddof=1) :: sha256 53e1ec81343ac2c1
\newcommand{\ProbeLogOneMinusAucVisibleContentMlpFourthree}{\ensuremath{-}7.0957}  % experiments/FIGS/data/probe_ladder_v2_mlp2/analysis/seed_level_results.json :: levels.visible_content.mlp[2].mean :: sha256 53e1ec81343ac2c1 :: natural log of 1-AUC, lower is better
\newcommand{\ProbeLogOneMinusAucSdVisibleContentMlpFourthree}{0.0773}  % experiments/FIGS/data/probe_ladder_v2_mlp2/analysis/seed_level_results.json :: levels.visible_content.mlp[2].seed_sd :: sha256 53e1ec81343ac2c1
\newcommand{\ProbeOmaVisibleContentMlpFourthree}{\ensuremath{(8.3 \pm 0.6)\times10^{-4}}}  % experiments/FIGS/data/probe_ladder_v2_mlp2/analysis/seed_level_results.json :: table[visible_content,mlp,43].auc :: sha256 53e1ec81343ac2c1 :: 1-AUC, mean +- SD over 5 seeds
\newcommand{\ProbeBkgLeftVisibleContentMlpFourthree}{0.4}  % experiments/FIGS/data/probe_ladder_v2_mlp2/analysis/seed_level_results.json :: levels.visible_content.mlp[2].rejection_points['0.90'].mean_n_bkg_pass :: sha256 53e1ec81343ac2c1 :: background jets passing the cut, mean over seeds
\newcommand{\ProbeRejVisibleContentMlpFourthree}{\ensuremath{>}538}  % experiments/FIGS/data/probe_ladder_v2_mlp2/analysis/seed_level_results.json :: table[visible_content,mlp,43].rejection_points['0.90'].rejection :: sha256 53e1ec81343ac2c1 :: at 90% signal efficiency, mean +- SD over seeds; 5 of 5 seeds at the cap
\newcommand{\ProbeAucVisibleContentMlpOneseven}{0.99793}  % experiments/FIGS/data/probe_ladder_v2_mlp2/analysis/seed_level_results.json :: levels.visible_content.mlp[3].mean_auc :: sha256 53e1ec81343ac2c1 :: 5 seeds, 0 saturated at AUC=1
\newcommand{\ProbeAucSdVisibleContentMlpOneseven}{0.00020}  % experiments/FIGS/data/probe_ladder_v2_mlp2/analysis/seed_level_results.json :: table[visible_content,mlp,17].auc (SD over seeds, ddof=1) :: sha256 53e1ec81343ac2c1
\newcommand{\ProbeLogOneMinusAucVisibleContentMlpOneseven}{\ensuremath{-}6.1839}  % experiments/FIGS/data/probe_ladder_v2_mlp2/analysis/seed_level_results.json :: levels.visible_content.mlp[3].mean :: sha256 53e1ec81343ac2c1 :: natural log of 1-AUC, lower is better
\newcommand{\ProbeLogOneMinusAucSdVisibleContentMlpOneseven}{0.1014}  % experiments/FIGS/data/probe_ladder_v2_mlp2/analysis/seed_level_results.json :: levels.visible_content.mlp[3].seed_sd :: sha256 53e1ec81343ac2c1
\newcommand{\ProbeOmaVisibleContentMlpOneseven}{\ensuremath{(2.07 \pm 0.20)\times10^{-3}}}  % experiments/FIGS/data/probe_ladder_v2_mlp2/analysis/seed_level_results.json :: table[visible_content,mlp,17].auc :: sha256 53e1ec81343ac2c1 :: 1-AUC, mean +- SD over 5 seeds
\newcommand{\ProbeBkgLeftVisibleContentMlpOneseven}{3.4}  % experiments/FIGS/data/probe_ladder_v2_mlp2/analysis/seed_level_results.json :: levels.visible_content.mlp[3].rejection_points['0.90'].mean_n_bkg_pass :: sha256 53e1ec81343ac2c1 :: background jets passing the cut, mean over seeds
\newcommand{\ProbeRejVisibleContentMlpOneseven}{\ensuremath{\geq}638\ensuremath{^{\ast}}}  % experiments/FIGS/data/probe_ladder_v2_mlp2/analysis/seed_level_results.json :: table[visible_content,mlp,17].rejection_points['0.90'].rejection :: sha256 53e1ec81343ac2c1 :: at 90% signal efficiency, mean +- SD over seeds; 1 of 5 seeds at the cap
\newcommand{\ProbeOmaRatioVisibleContentMlpOnesixtwoOverOneeighteight}{0.92}  % experiments/FIGS/data/probe_ladder_v2_mlp2/analysis/seed_level_results.json :: table[visible_content,mlp,162/188].auc :: sha256 53e1ec81343ac2c1 :: ratio of the seed means of 1-AUC, coarser over finer
\newcommand{\ProbeOmaRatioVisibleContentMlpFourthreeOverOneeighteight}{0.93}  % experiments/FIGS/data/probe_ladder_v2_mlp2/analysis/seed_level_results.json :: table[visible_content,mlp,43/188].auc :: sha256 53e1ec81343ac2c1 :: ratio of the seed means of 1-AUC, coarser over finer
\newcommand{\ProbeOmaRatioVisibleContentMlpOnesevenOverOneeighteight}{2.3}  % experiments/FIGS/data/probe_ladder_v2_mlp2/analysis/seed_level_results.json :: table[visible_content,mlp,17/188].auc :: sha256 53e1ec81343ac2c1 :: ratio of the seed means of 1-AUC, coarser over finer
\newcommand{\ProbeOmaRatioVisibleContentMlpFourthreeOverOnesixtwo}{1.01}  % experiments/FIGS/data/probe_ladder_v2_mlp2/analysis/seed_level_results.json :: table[visible_content,mlp,43/162].auc :: sha256 53e1ec81343ac2c1 :: ratio of the seed means of 1-AUC, coarser over finer
\newcommand{\ProbeOmaRatioVisibleContentMlpOnesevenOverOnesixtwo}{2.5}  % experiments/FIGS/data/probe_ladder_v2_mlp2/analysis/seed_level_results.json :: table[visible_content,mlp,17/162].auc :: sha256 53e1ec81343ac2c1 :: ratio of the seed means of 1-AUC, coarser over finer
\newcommand{\ProbeOmaRatioVisibleContentMlpOnesevenOverFourthree}{2.5}  % experiments/FIGS/data/probe_ladder_v2_mlp2/analysis/seed_level_results.json :: table[visible_content,mlp,17/43].auc :: sha256 53e1ec81343ac2c1 :: ratio of the seed means of 1-AUC, coarser over finer
\newcommand{\VocabSizeLoneeighteight}{188}  % configs/labelmaps/rung_label_maps.v1.csv :: column L188 (distinct group count) :: sha256 412f9aa3b58db412
\newcommand{\VocabSizeLonesixtwo}{162}  % configs/labelmaps/rung_label_maps.v1.csv :: column L162 (distinct group count) :: sha256 412f9aa3b58db412
\newcommand{\VocabSizeRsixthreeQone}{64}  % configs/labelmaps/rung_label_maps.v1.csv :: column R63_Q1 (distinct group count) :: sha256 412f9aa3b58db412
\newcommand{\VocabSizeRfourtwoQone}{43}  % configs/labelmaps/rung_label_maps.v1.csv :: column R42_Q1 (distinct group count) :: sha256 412f9aa3b58db412
\newcommand{\VocabSizeRtwonineQone}{30}  % configs/labelmaps/rung_label_maps.v1.csv :: column R29_Q1 (distinct group count) :: sha256 412f9aa3b58db412
\newcommand{\VocabSizeRonesixQone}{17}  % configs/labelmaps/rung_label_maps.v1.csv :: column R16_Q1 (distinct group count) :: sha256 412f9aa3b58db412
\newcommand{\VocabSizeRthreeVIS}{4}  % configs/labelmaps/rung_label_maps.v1.csv :: column R3_VIS (distinct group count) :: sha256 412f9aa3b58db412
\newcommand{\VocabSizeRoneQone}{2}  % configs/labelmaps/rung_label_maps.v1.csv :: column R1_Q1 (distinct group count) :: sha256 412f9aa3b58db412
\newcommand{\ProbeSplitTrain}{60\%}  % experiments/EVAL/probe.py :: SPLIT_FRACTIONS :: sha256 485db0966e660588 :: share of the probe sample; the same jets for every model
\newcommand{\ProbeSplitVal}{20\%}  % experiments/EVAL/probe.py :: SPLIT_FRACTIONS :: sha256 485db0966e660588 :: share of the probe sample; the same jets for every model
\newcommand{\ProbeSplitTest}{20\%}  % experiments/EVAL/probe.py :: SPLIT_FRACTIONS :: sha256 485db0966e660588 :: share of the probe sample; the same jets for every model
\newcommand{\ProbeMlpHidden}{256}  % experiments/EVAL/probe.py :: _fit_mlp :: sha256 485db0966e660588 :: hidden width of the MLP probe
\newcommand{\ProbeMlpInits}{3}  % experiments/EVAL/probe.py :: MLP_SEEDS :: sha256 485db0966e660588 :: MLP initialisations whose scores are averaged per cell
\newcommand{\ProbeMlpLrDrop}{10}  % experiments/EVAL/probe.py :: MLP_SCHEDULE['lr_factor'] :: sha256 485db0966e660588 :: factor the MLP's learning rate drops by on a plateau
\newcommand{\ProbeMlpPlateaus}{three}  % experiments/EVAL/probe.py :: MLP_SCHEDULE: stop_patience / (lr_patience + 1) :: sha256 485db0966e660588 :: patience windows without gain after which the MLP stops
\newcommand{\ProbeCMin}{\ensuremath{10^{-2}}}  % experiments/EVAL/probe.py :: C_GRID :: sha256 485db0966e660588 :: end of the logistic-regression inverse-regularisation grid
\newcommand{\ProbeCMax}{\ensuremath{10^{2}}}  % experiments/EVAL/probe.py :: C_GRID :: sha256 485db0966e660588 :: end of the logistic-regression inverse-regularisation grid
\newcommand{\ProbeNBkgTestBvcFourprong}{3{,}158}  % experiments/FIGS/data/probe_ladder_v2_mlp2/s1.json :: tasks.bvc_4prong.n_background_test :: sha256 f7e1de0e43be89b8 :: test-sample background jets
\newcommand{\ProbeNBkgTestBvcQcd}{5{,}371}  % experiments/FIGS/data/probe_ladder_v2_mlp2/s1.json :: tasks.bvc_qcd.n_background_test :: sha256 f7e1de0e43be89b8 :: test-sample background jets
\newcommand{\ProbeNBkgTestBvcResonant}{11{,}876}  % experiments/FIGS/data/probe_ladder_v2_mlp2/s1.json :: tasks.bvc_resonant.n_background_test :: sha256 f7e1de0e43be89b8 :: test-sample background jets
\newcommand{\ProbeNBkgTestEeVsMm}{4{,}708}  % experiments/FIGS/data/probe_ladder_v2_mlp2/s1.json :: tasks.ee_vs_mm.n_background_test :: sha256 f7e1de0e43be89b8 :: test-sample background jets
\newcommand{\ProbeNBkgTestRetainedTopology}{1{,}078}  % experiments/FIGS/data/probe_ladder_v2_mlp2/s1.json :: tasks.retained_topology.n_background_test :: sha256 f7e1de0e43be89b8 :: test-sample background jets
\newcommand{\ProbeNBkgTestVisibleContent}{2{,}554}  % experiments/FIGS/data/probe_ladder_v2_mlp2/s1.json :: tasks.visible_content.n_background_test :: sha256 f7e1de0e43be89b8 :: test-sample background jets
\newcommand{\UseDiesAtSophonEqfour}{43}  % configs/labelmaps/usecase_survival.v1.json :: sophon_eq4.dies_at :: sha256 122275213c349593 :: first vocabulary size at which it is not constructible
\newcommand{\UseLastAliveSophonEqfour}{64}  % configs/labelmaps/usecase_survival.v1.json :: sophon_eq4.last_rung_alive :: sha256 122275213c349593
\newcommand{\UseDiesAtSophonEqsixAone}{64}  % configs/labelmaps/usecase_survival.v1.json :: sophon_eq6_a1.dies_at :: sha256 122275213c349593 :: first vocabulary size at which it is not constructible
\newcommand{\UseLastAliveSophonEqsixAone}{162}  % configs/labelmaps/usecase_survival.v1.json :: sophon_eq6_a1.last_rung_alive :: sha256 122275213c349593
\newcommand{\UseDiesAtSophonEqsixAthree}{64}  % configs/labelmaps/usecase_survival.v1.json :: sophon_eq6_a3.dies_at :: sha256 122275213c349593 :: first vocabulary size at which it is not constructible
\newcommand{\UseLastAliveSophonEqsixAthree}{162}  % configs/labelmaps/usecase_survival.v1.json :: sophon_eq6_a3.last_rung_alive :: sha256 122275213c349593
\newcommand{\UseDiesAtSophonEqsevenNine}{64}  % configs/labelmaps/usecase_survival.v1.json :: sophon_eq7_9.dies_at :: sha256 122275213c349593 :: first vocabulary size at which it is not constructible
\newcommand{\UseLastAliveSophonEqsevenNine}{162}  % configs/labelmaps/usecase_survival.v1.json :: sophon_eq7_9.last_rung_alive :: sha256 122275213c349593
\newcommand{\UseDiesAtVcbEqone}{64}  % configs/labelmaps/usecase_survival.v1.json :: vcb_eq1.dies_at :: sha256 122275213c349593 :: first vocabulary size at which it is not constructible
\newcommand{\UseLastAliveVcbEqone}{162}  % configs/labelmaps/usecase_survival.v1.json :: vcb_eq1.last_rung_alive :: sha256 122275213c349593
\newcommand{\UseDiesAtGnthreexQcdSubclass}{162}  % configs/labelmaps/usecase_survival.v1.json :: gn3x_qcd_subclass.dies_at :: sha256 122275213c349593 :: first vocabulary size at which it is not constructible
\newcommand{\UseLastAliveGnthreexQcdSubclass}{188}  % configs/labelmaps/usecase_survival.v1.json :: gn3x_qcd_subclass.last_rung_alive :: sha256 122275213c349593
\newcommand{\UseWlikeLastAlive}{162}  % configs/labelmaps/rung_label_maps.v1.csv :: groups at each level vs {label_X_cs, label_X_qq} and that plus the QCD classes :: sha256 412f9aa3b58db412 :: last level at which Sophon's W-like score g_W(2) is constructible
\newcommand{\UseWlikeDiesAt}{64}  % configs/labelmaps/rung_label_maps.v1.csv :: groups at each level vs {label_X_cs, label_X_qq} and that plus the QCD classes :: sha256 412f9aa3b58db412 :: first level at which it is not
\newcommand{\DesignEpochs}{80}  % experiments/MTX/k8s/job-mtx-l188-s1-raunav.yaml :: --num-epochs :: sha256 96fef4b6337aaa27
\newcommand{\DesignSamplesPerEpoch}{10{,}240{,}000}  % experiments/MTX/k8s/job-mtx-l188-s1-raunav.yaml :: --samples-per-epoch :: sha256 96fef4b6337aaa27
\newcommand{\DesignExamplesSeen}{\ensuremath{8.192\times10^{8}}}  % experiments/MTX/k8s/job-mtx-l188-s1-raunav.yaml :: --num-epochs x --samples-per-epoch :: sha256 96fef4b6337aaa27 :: training examples drawn over the run
\newcommand{\DesignPassesSelected}{7.45}  % experiments/FIGS/data/v2/jc2_train/counts.json :: n_selected :: sha256 b0b5a8b63aa808c0 :: --num-epochs x --samples-per-epoch of job-mtx-l188-s1-raunav.yaml / n_selected: passes over the selected training jets
\newcommand{\DesignBatchSize}{512}  % experiments/MTX/k8s/job-mtx-l188-s1-raunav.yaml :: --batch-size :: sha256 96fef4b6337aaa27
\newcommand{\DesignStartLr}{\ensuremath{5\times10^{-4}}}  % experiments/MTX/k8s/job-mtx-l188-s1-raunav.yaml :: --start-lr :: sha256 96fef4b6337aaa27
\newcommand{\DesignHeadHidden}{512}  % experiments/MTX/k8s/job-mtx-l188-s1-raunav.yaml :: -o fc_params (hidden width) :: sha256 96fef4b6337aaa27
\newcommand{\DesignParticleBlocks}{8}  % experiments/E1/ParT_sophon_arch_10c.py :: num_layers=(\d+) :: sha256 5793faa7d55a8b0d
\newcommand{\DesignClassBlocks}{2}  % experiments/E1/ParT_sophon_arch_10c.py :: num_cls_layers=(\d+) :: sha256 5793faa7d55a8b0d
\newcommand{\DesignHeads}{8}  % experiments/E1/ParT_sophon_arch_10c.py :: num_heads=(\d+) :: sha256 5793faa7d55a8b0d
\newcommand{\DesignEmbedDim}{128}  % experiments/E1/ParT_sophon_arch_10c.py :: \bembed_dims=\[\d+, \d+, (\d+)\] :: sha256 5793faa7d55a8b0d
\newcommand{\DesignMpmMaskRate}{40\%}  % experiments/MTX/mpm.py :: DEFAULT_MASK_RATE :: sha256 49f0c8c6666d0211 :: share of each jet's particles hidden from the encoder
\newcommand{\DesignPtMin}{200}  % configs/arms/L188.yaml :: selection :: sha256 d5a188dc80885515 :: GeV
\newcommand{\DesignPtMax}{2{,}500}  % configs/arms/L188.yaml :: selection :: sha256 d5a188dc80885515 :: GeV
\newcommand{\DesignMsdMin}{20}  % configs/arms/L188.yaml :: selection :: sha256 d5a188dc80885515 :: GeV
\newcommand{\DesignMsdMax}{500}  % configs/arms/L188.yaml :: selection :: sha256 d5a188dc80885515 :: GeV
\newcommand{\LitJetClassTwoNJetsMillion}{139}  % experiments/FIGS/data/literature_facts.json :: facts[0].value :: sha256 713dba2a5e6e1dd1 :: arXiv:2405.12972v1 App. A, p. 11
\newcommand{\LitJetClassTwoNResonantClasses}{161}  % experiments/FIGS/data/literature_facts.json :: facts[1].value :: sha256 713dba2a5e6e1dd1 :: arXiv:2405.12972v1 App. A, p. 11
\newcommand{\LitJetClassTwoNQcdClasses}{27}  % experiments/FIGS/data/literature_facts.json :: facts[2].value :: sha256 713dba2a5e6e1dd1 :: arXiv:2405.12972v1 App. A, p. 11
\newcommand{\LitJetClassTwoMaxResonanceMass}{500}  % experiments/FIGS/data/literature_facts.json :: facts[3].value :: sha256 713dba2a5e6e1dd1 :: arXiv:2405.12972v1 Sec. II, p. 2
\newcommand{\LitMadgraphVersion}{2.9.18}  % experiments/FIGS/data/literature_facts.json :: facts[4].value :: sha256 713dba2a5e6e1dd1 :: arXiv:2405.12972v1 App. A, p. 11
\newcommand{\LitPythiaVersion}{8.3}  % experiments/FIGS/data/literature_facts.json :: facts[5].value :: sha256 713dba2a5e6e1dd1 :: arXiv:2405.12972v1 App. A, p. 11
\newcommand{\LitDelphesVersion}{3}  % experiments/FIGS/data/literature_facts.json :: facts[6].value :: sha256 713dba2a5e6e1dd1 :: arXiv:2405.12972v1 App. A, p. 11
\newcommand{\LitJetClassTwoPileupMean}{50}  % experiments/FIGS/data/literature_facts.json :: facts[7].value :: sha256 713dba2a5e6e1dd1 :: arXiv:2405.12972v1 Sec. II, p. 2
\newcommand{\LitJetClassTwoJetRadius}{0.8}  % experiments/FIGS/data/literature_facts.json :: facts[8].value :: sha256 713dba2a5e6e1dd1 :: arXiv:2405.12972v1 Sec. II, p. 2
\newcommand{\LitAtlasUnifiedEffDropPercent}{10}  % experiments/FIGS/data/literature_facts.json :: facts[9].value :: sha256 713dba2a5e6e1dd1 :: ATL-PHYS-PUB-2026-013 Sec. 6.5, p. 22
\newcommand{\LitSophonPaperParticleBlocks}{6}  % experiments/FIGS/data/literature_facts.json :: facts[10].value :: sha256 713dba2a5e6e1dd1 :: arXiv:2405.12972v1 App. B.1, p. 12
\newcommand{\LitSophonSamplesPerEpochMillion}{10}  % experiments/FIGS/data/literature_facts.json :: facts[11].value :: sha256 713dba2a5e6e1dd1 :: arXiv:2405.12972v1 App. B.1, p. 12
\newcommand{\LitSophonSubsetSamplesPerEpochMillion}{2.5}  % experiments/FIGS/data/literature_facts.json :: facts[12].value :: sha256 713dba2a5e6e1dd1 :: arXiv:2405.12972v1 App. B.2, p. 12
\newcommand{\LitGlopartNClasses}{314}  % experiments/FIGS/data/literature_facts.json :: facts[13].value :: sha256 713dba2a5e6e1dd1 :: C. Li, PhD thesis, CERN-THESIS-2024-281 Sec. 11.2.1, p. 188
\newcommand{\LitPartFtEpochs}{20}  % experiments/FIGS/data/literature_facts.json :: facts[14].value :: sha256 713dba2a5e6e1dd1 :: arXiv:2202.03772 Sec. 5.2, p. 8
\newcommand{\LitMehtaNCategories}{eight}  % experiments/FIGS/data/literature_facts.json :: facts[15].value :: sha256 713dba2a5e6e1dd1 :: arXiv:2605.05556v2 Abstract
\newcommand{\LitHongPretrainRuns}{2}  % experiments/FIGS/data/literature_facts.json :: facts[16].value :: sha256 713dba2a5e6e1dd1 :: arXiv:2410.23129 App. A.1.1, p. 18
\newcommand{\LitSophonLookaheadK}{6}  % experiments/FIGS/data/literature_facts.json :: facts[17].value :: sha256 713dba2a5e6e1dd1 :: arXiv:2405.12972v1 App. B.1, p. 12
\newcommand{\LitSophonLookaheadAlpha}{0.5}  % experiments/FIGS/data/literature_facts.json :: facts[18].value :: sha256 713dba2a5e6e1dd1 :: arXiv:2405.12972v1 App. B.1, p. 12
\newcommand{\LitSophonRadamBetaOne}{0.95}  % experiments/FIGS/data/literature_facts.json :: facts[19].value :: sha256 713dba2a5e6e1dd1 :: arXiv:2405.12972v1 App. B.1, p. 12
\newcommand{\LitSophonRadamBetaTwo}{0.999}  % experiments/FIGS/data/literature_facts.json :: facts[20].value :: sha256 713dba2a5e6e1dd1 :: arXiv:2405.12972v1 App. B.1, p. 12
\newcommand{\LitSophonRadamEps}{\ensuremath{10^{-5}}}  % experiments/FIGS/data/literature_facts.json :: facts[21].value :: sha256 713dba2a5e6e1dd1 :: arXiv:2405.12972v1 App. B.1, p. 12
\newcommand{\LitGnTwoAuxRejectionDropPercent}{30}  % experiments/FIGS/data/literature_facts.json :: facts[22].value :: sha256 713dba2a5e6e1dd1 :: arXiv:2505.19689v2 (Nature Commun. 17 (2026) 541) p. 9
\newcommand{\LitCmsPartNClasses}{37}  % experiments/FIGS/data/literature_facts.json :: facts[23].value :: sha256 713dba2a5e6e1dd1 :: arXiv:2604.09809v1 (CMS-JME-25-001) Sec. 5.2, p. 7
\newcommand{\LitCmsPartMassLambda}{0.05}  % experiments/FIGS/data/literature_facts.json :: facts[24].value :: sha256 713dba2a5e6e1dd1 :: arXiv:2604.09809v1 (CMS-JME-25-001) Sec. 5.2, p. 7
\newcommand{\LitGlopartVthreeNClasses}{374}  % experiments/FIGS/data/literature_facts.json :: facts[25].value :: sha256 713dba2a5e6e1dd1 :: CMS-DP-2026-104 (CMS public page) Architecture and outputs
\newcommand{\LitHuhCoarseClasses}{127}  % experiments/FIGS/data/literature_facts.json :: facts[26].value :: sha256 713dba2a5e6e1dd1 :: arXiv:1608.08614v2 Sec. 5.3, p. 5
\newcommand{\LitHuhKnnGapPercent}{15}  % experiments/FIGS/data/literature_facts.json :: facts[27].value :: sha256 713dba2a5e6e1dd1 :: arXiv:1608.08614v2 Sec. 5.3, p. 5
\newcommand{\LitRsthreelRuns}{3}  % experiments/FIGS/data/literature_facts.json :: facts[28].value :: sha256 713dba2a5e6e1dd1 :: arXiv:2403.07066v2 (Phys. Rev. D 111 (2025) 032010) Sec. IV, p. 5
\newcommand{\LitSophonParamsMillion}{2.3}  % experiments/FIGS/data/literature_facts.json :: facts[29].value :: sha256 713dba2a5e6e1dd1 :: arXiv:2405.12972v1 App. B.1, p. 12
\newcommand{\LitGlopartVoneNClasses}{37}  % experiments/FIGS/data/literature_facts.json :: facts[30].value :: sha256 713dba2a5e6e1dd1 :: CMS-DP-2026-104 (CMS public page) Model generations
\newcommand{\LitSophonXMass}{200}  % experiments/FIGS/data/literature_facts.json :: facts[31].value :: sha256 713dba2a5e6e1dd1 :: arXiv:2405.12972v1 Sec. III, p. 3
\newcommand{\LitSophonXLumi}{40}  % experiments/FIGS/data/literature_facts.json :: facts[32].value :: sha256 713dba2a5e6e1dd1 :: arXiv:2405.12972v1 Sec. III, Fig. 2
\newcommand{\LitSophonXWindowLo}{150}  % experiments/FIGS/data/literature_facts.json :: facts[33].value :: sha256 713dba2a5e6e1dd1 :: arXiv:2405.12972v1 Sec. III, Fig. 2
\newcommand{\LitSophonXWindowHi}{230}  % experiments/FIGS/data/literature_facts.json :: facts[34].value :: sha256 713dba2a5e6e1dd1 :: arXiv:2405.12972v1 Sec. III, Fig. 2
\newcommand{\LitSophonAblationNClasses}{42}  % experiments/FIGS/data/literature_facts.json :: facts[35].value :: sha256 713dba2a5e6e1dd1 :: arXiv:2405.12972v1 Sec. III
\newcommand{\LitSophonTransferSamplesMillion}{2.5}  % experiments/FIGS/data/literature_facts.json :: facts[36].value :: sha256 713dba2a5e6e1dd1 :: arXiv:2405.12972v1 App. B.5, p. 13
\newcommand{\LitSophonTransferBatch}{1024}  % experiments/FIGS/data/literature_facts.json :: facts[37].value :: sha256 713dba2a5e6e1dd1 :: arXiv:2405.12972v1 App. B.5, p. 13
\newcommand{\LitSophonScratchEpochs}{50}  % experiments/FIGS/data/literature_facts.json :: facts[38].value :: sha256 713dba2a5e6e1dd1 :: arXiv:2405.12972v1 App. B.3
\newcommand{\LitHuhFineClasses}{1000}  % experiments/FIGS/data/literature_facts.json :: facts[39].value :: sha256 713dba2a5e6e1dd1 :: arXiv:1608.08614v2 Sec. 1, p. 1
\newcommand{\LitHuhPascalDropMap}{2.8}  % experiments/FIGS/data/literature_facts.json :: facts[40].value :: sha256 713dba2a5e6e1dd1 :: arXiv:1608.08614v2 Sec. 1, p. 1
\newcommand{\LitHongSuperclasses}{11}  % experiments/FIGS/data/literature_facts.json :: facts[41].value :: sha256 713dba2a5e6e1dd1 :: arXiv:2410.23129 Sec. 6.2
\newcommand{\LitHongDeclineClasses}{\ensuremath{10^{3}}}  % experiments/FIGS/data/literature_facts.json :: facts[42].value :: sha256 713dba2a5e6e1dd1 :: arXiv:2410.23129 Sec. 6.2
\newcommand{\LitObjectivesFullJets}{\ensuremath{1.2\times10^{6}}}  % experiments/FIGS/data/literature_facts.json :: facts[43].value :: sha256 713dba2a5e6e1dd1 :: arXiv:2606.14870 Sec. IV A
\newcommand{\LitObjectivesFewJets}{\ensuremath{10^{4}}}  % experiments/FIGS/data/literature_facts.json :: facts[44].value :: sha256 713dba2a5e6e1dd1 :: arXiv:2606.14870 Sec. IV A
\newcommand{\LitMpmBenefitJets}{10k}  % experiments/FIGS/data/literature_facts.json :: facts[45].value :: sha256 713dba2a5e6e1dd1 :: arXiv:2401.13537 Sec. V
\newcommand{\LitPartTopRejScratch}{1602}  % experiments/FIGS/data/literature_facts.json :: facts[46].value :: sha256 713dba2a5e6e1dd1 :: arXiv:2202.03772 Table 5
\newcommand{\LitPartTopRejFt}{2766}  % experiments/FIGS/data/literature_facts.json :: facts[47].value :: sha256 713dba2a5e6e1dd1 :: arXiv:2202.03772 Table 5
\newcommand{\LitMahajanLinearGapPercent}{7}  % experiments/FIGS/data/literature_facts.json :: facts[48].value :: sha256 713dba2a5e6e1dd1 :: arXiv:1805.00932 Sec. 3.1.2
\newcommand{\LitBukchinCoarse}{33.83}  % experiments/FIGS/data/literature_facts.json :: facts[49].value :: sha256 713dba2a5e6e1dd1 :: arXiv:2012.03515 Table 2, LIVING-17 all-way
\newcommand{\LitBukchinFine}{58.95}  % experiments/FIGS/data/literature_facts.json :: facts[50].value :: sha256 713dba2a5e6e1dd1 :: arXiv:2012.03515 Table 2, LIVING-17 all-way
\newcommand{\LitKornblithRho}{\ensuremath{-0.93}}  % experiments/FIGS/data/literature_facts.json :: facts[51].value :: sha256 713dba2a5e6e1dd1 :: arXiv:2010.16402 Sec. 3.4
\newcommand{\LitPartTopSigEffPercent}{30}  % experiments/FIGS/data/literature_facts.json :: facts[52].value :: sha256 713dba2a5e6e1dd1 :: arXiv:2202.03772 Table 5 (column Rej30%)
\newcommand{\LitPartTopKinFeatures}{7}  % experiments/FIGS/data/literature_facts.json :: facts[53].value :: sha256 713dba2a5e6e1dd1 :: arXiv:2202.03772 Sec. 5.2
\newcommand{\LitBukchinAncor}{45.14}  % experiments/FIGS/data/literature_facts.json :: facts[54].value :: sha256 713dba2a5e6e1dd1 :: arXiv:2012.03515 Table 2, LIVING-17 all-way
\newcommand{\LitMahajanFtGapPp}{0.1}  % experiments/FIGS/data/literature_facts.json :: facts[55].value :: sha256 713dba2a5e6e1dd1 :: arXiv:1805.00932 Fig. 1, ImageNet-9k target
\newcommand{\LitSophonTransferLr}{5 × 10−4}  % experiments/FIGS/data/literature_facts.json :: facts[56].value :: sha256 713dba2a5e6e1dd1 :: arXiv:2405.12972v1 App. B.5
\newcommand{\LitQgJetRadius}{0.4}  % experiments/FIGS/data/literature_facts.json :: facts[57].value :: sha256 713dba2a5e6e1dd1 :: arXiv:1810.05165 Sec. 4.1
\newcommand{\LitQgPtMin}{500}  % experiments/FIGS/data/literature_facts.json :: facts[58].value :: sha256 713dba2a5e6e1dd1 :: arXiv:1810.05165 Sec. 4.1
\newcommand{\LitQgPtMax}{550}  % experiments/FIGS/data/literature_facts.json :: facts[59].value :: sha256 713dba2a5e6e1dd1 :: arXiv:1810.05165 Sec. 4.1
\newcommand{\LitPartTopAucScratch}{0.9858}  % experiments/FIGS/data/literature_facts.json :: facts[60].value :: sha256 713dba2a5e6e1dd1 :: arXiv:2202.03772 Table 5
\newcommand{\LitPartTopAucFt}{0.9877}  % experiments/FIGS/data/literature_facts.json :: facts[61].value :: sha256 713dba2a5e6e1dd1 :: arXiv:2202.03772 Table 5
\newcommand{\LitPartQgAucScratch}{0.9121}  % experiments/FIGS/data/literature_facts.json :: facts[62].value :: sha256 713dba2a5e6e1dd1 :: arXiv:2202.03772 Table 6 (exp)
\newcommand{\LitPartQgRejScratch}{101.2}  % experiments/FIGS/data/literature_facts.json :: facts[63].value :: sha256 713dba2a5e6e1dd1 :: arXiv:2202.03772 Table 6 (exp)
\newcommand{\LitPartQgAucFt}{0.9151}  % experiments/FIGS/data/literature_facts.json :: facts[64].value :: sha256 713dba2a5e6e1dd1 :: arXiv:2202.03772 Table 6 (exp)
\newcommand{\LitPartQgRejFt}{107.9}  % experiments/FIGS/data/literature_facts.json :: facts[65].value :: sha256 713dba2a5e6e1dd1 :: arXiv:2202.03772 Table 6 (exp)
\newcommand{\LitPartJcTestMillion}{20}  % experiments/FIGS/data/literature_facts.json :: facts[66].value :: sha256 713dba2a5e6e1dd1 :: arXiv:2202.03772 Sec. 3
\newcommand{\LitPartJcTwoMAcc}{0.836}  % experiments/FIGS/data/literature_facts.json :: facts[67].value :: sha256 713dba2a5e6e1dd1 :: arXiv:2202.03772 Table 3
\newcommand{\LitPartSmallSetIters}{100 k}  % experiments/FIGS/data/literature_facts.json :: facts[68].value :: sha256 713dba2a5e6e1dd1 :: arXiv:2202.03772 Sec. 5.1 (training dataset size)
\newcommand{\LitPartSmallSetPercent}{2}  % experiments/FIGS/data/literature_facts.json :: facts[69].value :: sha256 713dba2a5e6e1dd1 :: arXiv:2202.03772 Sec. 5.1 (training dataset size)
\newcommand{\LitSophonEtaMax}{2.5}  % experiments/FIGS/data/literature_facts.json :: facts[70].value :: sha256 713dba2a5e6e1dd1 :: arXiv:2405.12972 App. A
\newcommand{\LitBHadronLifetimePs}{1.5}  % experiments/FIGS/data/literature_facts.json :: facts[71].value :: sha256 713dba2a5e6e1dd1 :: arXiv:1712.07158 Sec. 4.1
\newcommand{\LitCHadronLifetimePs}{1}  % experiments/FIGS/data/literature_facts.json :: facts[72].value :: sha256 713dba2a5e6e1dd1 :: arXiv:1712.07158 Sec. 4.1
\newcommand{\LitOmnilearnPidDropProb}{0.2}  % experiments/FIGS/data/literature_facts.json :: facts[73].value :: sha256 713dba2a5e6e1dd1 :: arXiv:2502.14652 Sec. II.A
\newcommand{\LitMahajanFtNineKLarge}{48.5}  % experiments/FIGS/data/literature_facts.json :: facts[74].value :: sha256 713dba2a5e6e1dd1 :: arXiv:1805.00932 Fig. 1 (IG-1B-17k, IN-9k, full fine-tuning)
\newcommand{\LitMahajanFtNineKSmall}{48.4}  % experiments/FIGS/data/literature_facts.json :: facts[75].value :: sha256 713dba2a5e6e1dd1 :: arXiv:1805.00932 Fig. 1 (IG-940M-1.5k, IN-9k, full fine-tuning)
\newcommand{\LitMahajanFtOneKSmall}{84.2}  % experiments/FIGS/data/literature_facts.json :: facts[76].value :: sha256 713dba2a5e6e1dd1 :: arXiv:1805.00932 Fig. 1 (IG-940M-1.5k, IN-1k, full fine-tuning)
\newcommand{\LitMahajanFtOneKLarge}{83.6}  % experiments/FIGS/data/literature_facts.json :: facts[77].value :: sha256 713dba2a5e6e1dd1 :: arXiv:1805.00932 Fig. 1 (IG-1B-17k, IN-1k, full fine-tuning)
\newcommand{\LitChenGainFifteen}{0.16}  % experiments/FIGS/data/literature_facts.json :: facts[78].value :: sha256 713dba2a5e6e1dd1 :: arXiv:1901.07012 Table VIII (15 coarse classes, CIFAR-100)
\newcommand{\LitChenGainTwenty}{0.01}  % experiments/FIGS/data/literature_facts.json :: facts[79].value :: sha256 713dba2a5e6e1dd1 :: arXiv:1901.07012 Table VIII (20 coarse classes, CIFAR-100)
\newcommand{\LitChenCoarseFifteen}{15}  % experiments/FIGS/data/literature_facts.json :: facts[80].value :: sha256 713dba2a5e6e1dd1 :: arXiv:1901.07012 Table VIII
\newcommand{\LitChenCoarseTwenty}{20}  % experiments/FIGS/data/literature_facts.json :: facts[81].value :: sha256 713dba2a5e6e1dd1 :: arXiv:1901.07012 Table VIII
\newcommand{\LitChenRandomMaxGain}{2.78}  % experiments/FIGS/data/literature_facts.json :: facts[82].value :: sha256 713dba2a5e6e1dd1 :: arXiv:1901.07012 Sec. VI-A, Table VI
\newcommand{\LitGlopartMassLambda}{5}  % experiments/FIGS/data/literature_facts.json :: facts[83].value :: sha256 713dba2a5e6e1dd1 :: CERN-THESIS-2024-281 Eq. (11.7), Sec. 11.2.1
\newcommand{\LitGlopartAddedAucBase}{0.9555}  % experiments/FIGS/data/literature_facts.json :: facts[84].value :: sha256 713dba2a5e6e1dd1 :: CERN-THESIS-2024-281 Sec. 11.4.1, Fig. 11.19 (H->WW 4q vs QCD)
\newcommand{\LitGlopartAddedAucTop}{0.9593}  % experiments/FIGS/data/literature_facts.json :: facts[85].value :: sha256 713dba2a5e6e1dd1 :: CERN-THESIS-2024-281 Sec. 11.4.1, Fig. 11.19 (H->WW 4q vs QCD)
\newcommand{\LitCmsSdZcut}{0.1}  % experiments/FIGS/data/literature_facts.json :: facts[86].value :: sha256 713dba2a5e6e1dd1 :: arXiv:2004.08262 Sec. 6.1
\newcommand{\LitCmsSdBeta}{0}  % experiments/FIGS/data/literature_facts.json :: facts[87].value :: sha256 713dba2a5e6e1dd1 :: arXiv:2004.08262 Sec. 6.1
\newcommand{\LitQgNJetsMillion}{2}  % experiments/FIGS/data/literature_facts.json :: facts[88].value :: sha256 713dba2a5e6e1dd1 :: arXiv:2202.03772 Sec. 5.2
\newcommand{\LitQgTrainMillion}{1.6}  % experiments/FIGS/data/literature_facts.json :: facts[89].value :: sha256 713dba2a5e6e1dd1 :: arXiv:2202.03772 Sec. 5.2
\newcommand{\LitQgValMillion}{0.2}  % experiments/FIGS/data/literature_facts.json :: facts[90].value :: sha256 713dba2a5e6e1dd1 :: arXiv:2202.03772 Sec. 5.2
\newcommand{\LitQgRapidityMax}{1.7}  % experiments/FIGS/data/literature_facts.json :: facts[91].value :: sha256 713dba2a5e6e1dd1 :: Zenodo 10.5281/zenodo.3164691 Zenodo record 3164691, description
\newcommand{\LitHerwigVersion}{7.1.4}  % experiments/FIGS/data/literature_facts.json :: facts[92].value :: sha256 713dba2a5e6e1dd1 :: Zenodo 10.5281/zenodo.3066475 record description
\newcommand{\LitObjectivesFewestJets}{\ensuremath{10^{3}}}  % experiments/FIGS/data/literature_facts.json :: facts[93].value :: sha256 713dba2a5e6e1dd1 :: arXiv:2606.14870 Table XVII (Micro model, Npt = 10^8: Classifier 136 vs MPM 218 at 10^3, 331 vs 539 at 10^4)
\newcommand{\LitPartTopNTrainMillion}{1.2}  % experiments/FIGS/data/literature_facts.json :: facts[94].value :: sha256 713dba2a5e6e1dd1 :: arXiv:2202.03772 Sec. 5.2
\newcommand{\LitLandscapeParticlenetEnsMean}{1412}  % experiments/FIGS/data/literature_facts.json :: facts[95].value :: sha256 713dba2a5e6e1dd1 :: arXiv:1902.09914 Table 1 (ParticleNet, ensemble mean of 1/eps_B at eps_S=0.3)
\newcommand{\LitPartFtHeadReinits}{9}  % experiments/FIGS/data/literature_facts.json :: facts[96].value :: sha256 713dba2a5e6e1dd1 :: arXiv:2202.03772 Sec. 5.2
\newcommand{\LitSophonBlocksCode}{8}  % experiments/FIGS/data/literature_facts.json :: facts[97].value :: sha256 713dba2a5e6e1dd1 :: github.com/jet-universe/sophon@9dd6dd6 networks/example_ParticleTransformer_sophon.py line 67
\newcommand{\LitGrasHerwigVersion}{2.7.1}  % experiments/FIGS/data/literature_facts.json :: facts[98].value :: sha256 713dba2a5e6e1dd1 :: arXiv:1704.03878 Sec. 2 (list of generators)
\newcommand{\DesignMassLambda}{5.0}  % experiments/MTX/k8s/job-mtx-l162_mass-s1-raunav.yaml :: --mass-lambda :: sha256 8882459fe5520795 :: identical in all ten mass-output pretraining specs
\newcommand{\FtRecipeNRuns}{264}  % experiments/FIGS/data/ft_recipes/recipes_w2b_bench_v2_v3.json :: runs (count), reported cells :: sha256 7dc8c791343e1b6f
\newcommand{\FtRecipeLr}{\ensuremath{1\times10^{-4}}}  % experiments/FIGS/data/ft_recipes/recipes_w2b_bench_v2_v3.json :: runs[*].lr, pretrained starts :: sha256 7dc8c791343e1b6f
\newcommand{\FtRecipeHeadMult}{50}  % experiments/FIGS/data/ft_recipes/recipes_w2b_bench_v2_v3.json :: runs[*].head_lr_mult, pretrained starts :: sha256 7dc8c791343e1b6f
\newcommand{\FtRecipeLrHead}{\ensuremath{5\times10^{-3}}}  % experiments/FIGS/data/ft_recipes/recipes_w2b_bench_v2_v3.json :: lr x head_lr_mult :: sha256 7dc8c791343e1b6f
\newcommand{\FtRecipeWeightDecay}{0.01}  % experiments/FIGS/data/ft_recipes/recipes_w2b_bench_v2_v3.json :: runs[*].weight_decay :: sha256 7dc8c791343e1b6f
\newcommand{\FtRecipeBatch}{512}  % experiments/FIGS/data/ft_recipes/recipes_w2b_bench_v2_v3.json :: runs[*].batch_size :: sha256 7dc8c791343e1b6f
\newcommand{\FtRecipeEpochsEThree}{50}  % experiments/FIGS/data/ft_recipes/recipes_w2b_bench_v2_v3.json :: runs[leg 1,2; n_train 1000].epochs :: sha256 7dc8c791343e1b6f
\newcommand{\FtRecipeStepsEThree}{19}  % experiments/FT/k8s/job-ft-legs-w2-raunav.yaml :: job commands: samples per epoch at n_train 1000 // batch size :: sha256 f0962eddb451b7b0 :: optimizer steps per epoch, as weaver floors them
\newcommand{\FtRecipePassesEThree}{10}  % experiments/FT/k8s/job-ft-legs-w2-raunav.yaml :: job commands: samples per epoch at n_train 1000 // n_train :: sha256 f0962eddb451b7b0 :: passes over the training subset per epoch
\newcommand{\FtRecipeEpochsEFour}{50}  % experiments/FIGS/data/ft_recipes/recipes_w2b_bench_v2_v3.json :: runs[leg 1,2; n_train 10000].epochs :: sha256 7dc8c791343e1b6f
\newcommand{\FtRecipeStepsEFour}{19}  % experiments/FT/k8s/job-ft-legs-w2-raunav.yaml :: job commands: samples per epoch at n_train 10000 // batch size :: sha256 f0962eddb451b7b0 :: optimizer steps per epoch, as weaver floors them
\newcommand{\FtRecipePassesEFour}{one}  % experiments/FT/k8s/job-ft-legs-w2-raunav.yaml :: job commands: samples per epoch at n_train 10000 // n_train :: sha256 f0962eddb451b7b0 :: passes over the training subset per epoch
\newcommand{\FtRecipeEpochsEFive}{30}  % experiments/FIGS/data/ft_recipes/recipes_w2b_bench_v2_v3.json :: runs[leg 1,2; n_train 100000].epochs :: sha256 7dc8c791343e1b6f
\newcommand{\FtRecipeStepsEFive}{195}  % experiments/FT/k8s/job-ft-legs-w2-raunav.yaml :: job commands: samples per epoch at n_train 100000 // batch size :: sha256 f0962eddb451b7b0 :: optimizer steps per epoch, as weaver floors them
\newcommand{\FtRecipePassesEFive}{one}  % experiments/FT/k8s/job-ft-legs-w2-raunav.yaml :: job commands: samples per epoch at n_train 100000 // n_train :: sha256 f0962eddb451b7b0 :: passes over the training subset per epoch
\newcommand{\FtRecipeEpochsESix}{10}  % experiments/FIGS/data/ft_recipes/recipes_w2b_bench_v2_v3.json :: runs[leg 1,2; n_train 1000000].epochs :: sha256 7dc8c791343e1b6f
\newcommand{\FtRecipeStepsESix}{1{,}953}  % experiments/FT/k8s/job-ft-legs-w2-raunav.yaml :: job commands: samples per epoch at n_train 1000000 // batch size :: sha256 f0962eddb451b7b0 :: optimizer steps per epoch, as weaver floors them
\newcommand{\FtRecipePassesESix}{one}  % experiments/FT/k8s/job-ft-legs-w2-raunav.yaml :: job commands: samples per epoch at n_train 1000000 // n_train :: sha256 f0962eddb451b7b0 :: passes over the training subset per epoch
\newcommand{\FtRecipeValJets}{19{,}968}  % experiments/FT/k8s/job-ft-legs-w2-raunav.yaml :: --samples-per-epoch-val 20000 // batch x batch :: sha256 f0962eddb451b7b0 :: validation jets per epoch, floored to whole batches
\newcommand{\RecoveryAccOneeighteightByOneeighteight}{0.4123\,\ensuremath{\pm}\,0.0011}  % experiments/FIGS/data/label_recovery_curve_v1err/analysis/label_recovery_curve_summary.json :: table[rung=L188,level=188,probe=linear,n_train=max].accuracy :: sha256 1250d1f00537e0bd :: balanced accuracy, mean +- SD over 5 runs
\newcommand{\RecoveryAccOneeighteightByOnesixtwo}{0.4090\,\ensuremath{\pm}\,0.0012}  % experiments/FIGS/data/label_recovery_curve_v1err/analysis/label_recovery_curve_summary.json :: table[rung=L188,level=162,probe=linear,n_train=max].accuracy :: sha256 1250d1f00537e0bd :: balanced accuracy, mean +- SD over 5 runs
\newcommand{\RecoveryAccOneeighteightByFourthree}{0.356\,\ensuremath{\pm}\,0.008}  % experiments/FIGS/data/label_recovery_curve_v1err/analysis/label_recovery_curve_summary.json :: table[rung=L188,level=43,probe=linear,n_train=max].accuracy :: sha256 1250d1f00537e0bd :: balanced accuracy, mean +- SD over 5 runs
\newcommand{\RecoveryAccOneeighteightByOneseven}{0.306\,\ensuremath{\pm}\,0.008}  % experiments/FIGS/data/label_recovery_curve_v1err/analysis/label_recovery_curve_summary.json :: table[rung=L188,level=17,probe=linear,n_train=max].accuracy :: sha256 1250d1f00537e0bd :: balanced accuracy, mean +- SD over 5 runs
\newcommand{\RecoveryAccOnesixtwoByOneeighteight}{0.4623\,\ensuremath{\pm}\,0.0006}  % experiments/FIGS/data/label_recovery_curve_v1err/analysis/label_recovery_curve_summary.json :: table[rung=L162,level=188,probe=linear,n_train=max].accuracy :: sha256 1250d1f00537e0bd :: balanced accuracy, mean +- SD over 5 runs
\newcommand{\RecoveryAccOnesixtwoByOnesixtwo}{0.4622\,\ensuremath{\pm}\,0.0007}  % experiments/FIGS/data/label_recovery_curve_v1err/analysis/label_recovery_curve_summary.json :: table[rung=L162,level=162,probe=linear,n_train=max].accuracy :: sha256 1250d1f00537e0bd :: balanced accuracy, mean +- SD over 5 runs
\newcommand{\RecoveryAccOnesixtwoByFourthree}{0.402\,\ensuremath{\pm}\,0.009}  % experiments/FIGS/data/label_recovery_curve_v1err/analysis/label_recovery_curve_summary.json :: table[rung=L162,level=43,probe=linear,n_train=max].accuracy :: sha256 1250d1f00537e0bd :: balanced accuracy, mean +- SD over 5 runs
\newcommand{\RecoveryAccOnesixtwoByOneseven}{0.348\,\ensuremath{\pm}\,0.010}  % experiments/FIGS/data/label_recovery_curve_v1err/analysis/label_recovery_curve_summary.json :: table[rung=L162,level=17,probe=linear,n_train=max].accuracy :: sha256 1250d1f00537e0bd :: balanced accuracy, mean +- SD over 5 runs
\newcommand{\RecoveryAccSixfourByOneeighteight}{0.5979\,\ensuremath{\pm}\,0.0008}  % experiments/FIGS/data/label_recovery_curve_v1err/analysis/label_recovery_curve_summary.json :: table[rung=R63_Q1,level=188,probe=linear,n_train=max].accuracy :: sha256 1250d1f00537e0bd :: balanced accuracy, mean +- SD over 5 runs
\newcommand{\RecoveryAccSixfourByOnesixtwo}{0.5973\,\ensuremath{\pm}\,0.0013}  % experiments/FIGS/data/label_recovery_curve_v1err/analysis/label_recovery_curve_summary.json :: table[rung=R63_Q1,level=162,probe=linear,n_train=max].accuracy :: sha256 1250d1f00537e0bd :: balanced accuracy, mean +- SD over 5 runs
\newcommand{\RecoveryAccSixfourByFourthree}{0.5881\,\ensuremath{\pm}\,0.0007}  % experiments/FIGS/data/label_recovery_curve_v1err/analysis/label_recovery_curve_summary.json :: table[rung=R63_Q1,level=43,probe=linear,n_train=max].accuracy :: sha256 1250d1f00537e0bd :: balanced accuracy, mean +- SD over 5 runs
\newcommand{\RecoveryAccSixfourByOneseven}{0.531\,\ensuremath{\pm}\,0.006}  % experiments/FIGS/data/label_recovery_curve_v1err/analysis/label_recovery_curve_summary.json :: table[rung=R63_Q1,level=17,probe=linear,n_train=max].accuracy :: sha256 1250d1f00537e0bd :: balanced accuracy, mean +- SD over 5 runs
\newcommand{\RecoveryAccFourthreeByOneeighteight}{0.6770\,\ensuremath{\pm}\,0.0006}  % experiments/FIGS/data/label_recovery_curve_v1err/analysis/label_recovery_curve_summary.json :: table[rung=R42_Q1,level=188,probe=linear,n_train=max].accuracy :: sha256 1250d1f00537e0bd :: balanced accuracy, mean +- SD over 5 runs
\newcommand{\RecoveryAccFourthreeByOnesixtwo}{0.6766\,\ensuremath{\pm}\,0.0008}  % experiments/FIGS/data/label_recovery_curve_v1err/analysis/label_recovery_curve_summary.json :: table[rung=R42_Q1,level=162,probe=linear,n_train=max].accuracy :: sha256 1250d1f00537e0bd :: balanced accuracy, mean +- SD over 5 runs
\newcommand{\RecoveryAccFourthreeByFourthree}{0.6751\,\ensuremath{\pm}\,0.0011}  % experiments/FIGS/data/label_recovery_curve_v1err/analysis/label_recovery_curve_summary.json :: table[rung=R42_Q1,level=43,probe=linear,n_train=max].accuracy :: sha256 1250d1f00537e0bd :: balanced accuracy, mean +- SD over 5 runs
\newcommand{\RecoveryAccFourthreeByOneseven}{0.633\,\ensuremath{\pm}\,0.006}  % experiments/FIGS/data/label_recovery_curve_v1err/analysis/label_recovery_curve_summary.json :: table[rung=R42_Q1,level=17,probe=linear,n_train=max].accuracy :: sha256 1250d1f00537e0bd :: balanced accuracy, mean +- SD over 5 runs
\newcommand{\RecoveryAccThreezeroByOneeighteight}{0.7001\,\ensuremath{\pm}\,0.0005}  % experiments/FIGS/data/label_recovery_curve_v1err/analysis/label_recovery_curve_summary.json :: table[rung=R29_Q1,level=188,probe=linear,n_train=max].accuracy :: sha256 1250d1f00537e0bd :: balanced accuracy, mean +- SD over 5 runs
\newcommand{\RecoveryAccThreezeroByOnesixtwo}{0.7004\,\ensuremath{\pm}\,0.0007}  % experiments/FIGS/data/label_recovery_curve_v1err/analysis/label_recovery_curve_summary.json :: table[rung=R29_Q1,level=162,probe=linear,n_train=max].accuracy :: sha256 1250d1f00537e0bd :: balanced accuracy, mean +- SD over 5 runs
\newcommand{\RecoveryAccThreezeroByFourthree}{0.6998\,\ensuremath{\pm}\,0.0008}  % experiments/FIGS/data/label_recovery_curve_v1err/analysis/label_recovery_curve_summary.json :: table[rung=R29_Q1,level=43,probe=linear,n_train=max].accuracy :: sha256 1250d1f00537e0bd :: balanced accuracy, mean +- SD over 5 runs
\newcommand{\RecoveryAccThreezeroByOneseven}{0.6778\,\ensuremath{\pm}\,0.0030}  % experiments/FIGS/data/label_recovery_curve_v1err/analysis/label_recovery_curve_summary.json :: table[rung=R29_Q1,level=17,probe=linear,n_train=max].accuracy :: sha256 1250d1f00537e0bd :: balanced accuracy, mean +- SD over 5 runs
\newcommand{\RecoveryAccOnesevenByOneeighteight}{0.7675\,\ensuremath{\pm}\,0.0010}  % experiments/FIGS/data/label_recovery_curve_v1err/analysis/label_recovery_curve_summary.json :: table[rung=R16_Q1,level=188,probe=linear,n_train=max].accuracy :: sha256 1250d1f00537e0bd :: balanced accuracy, mean +- SD over 5 runs
\newcommand{\RecoveryAccOnesevenByOnesixtwo}{0.7681\,\ensuremath{\pm}\,0.0009}  % experiments/FIGS/data/label_recovery_curve_v1err/analysis/label_recovery_curve_summary.json :: table[rung=R16_Q1,level=162,probe=linear,n_train=max].accuracy :: sha256 1250d1f00537e0bd :: balanced accuracy, mean +- SD over 5 runs
\newcommand{\RecoveryAccOnesevenByFourthree}{0.7680\,\ensuremath{\pm}\,0.0010}  % experiments/FIGS/data/label_recovery_curve_v1err/analysis/label_recovery_curve_summary.json :: table[rung=R16_Q1,level=43,probe=linear,n_train=max].accuracy :: sha256 1250d1f00537e0bd :: balanced accuracy, mean +- SD over 5 runs
\newcommand{\RecoveryAccOnesevenByOneseven}{0.7681\,\ensuremath{\pm}\,0.0003}  % experiments/FIGS/data/label_recovery_curve_v1err/analysis/label_recovery_curve_summary.json :: table[rung=R16_Q1,level=17,probe=linear,n_train=max].accuracy :: sha256 1250d1f00537e0bd :: balanced accuracy, mean +- SD over 5 runs
\newcommand{\RecoveryAccFourByOneeighteight}{0.7526\,\ensuremath{\pm}\,0.0002}  % experiments/FIGS/data/label_recovery_curve_v1err/analysis/label_recovery_curve_summary.json :: table[rung=R3_VIS,level=188,probe=linear,n_train=max].accuracy :: sha256 1250d1f00537e0bd :: balanced accuracy, mean +- SD over 5 runs
\newcommand{\RecoveryAccFourByOnesixtwo}{0.7538\,\ensuremath{\pm}\,0.0017}  % experiments/FIGS/data/label_recovery_curve_v1err/analysis/label_recovery_curve_summary.json :: table[rung=R3_VIS,level=162,probe=linear,n_train=max].accuracy :: sha256 1250d1f00537e0bd :: balanced accuracy, mean +- SD over 5 runs
\newcommand{\RecoveryAccFourByFourthree}{0.7557\,\ensuremath{\pm}\,0.0015}  % experiments/FIGS/data/label_recovery_curve_v1err/analysis/label_recovery_curve_summary.json :: table[rung=R3_VIS,level=43,probe=linear,n_train=max].accuracy :: sha256 1250d1f00537e0bd :: balanced accuracy, mean +- SD over 5 runs
\newcommand{\RecoveryAccFourByOneseven}{0.7586\,\ensuremath{\pm}\,0.0008}  % experiments/FIGS/data/label_recovery_curve_v1err/analysis/label_recovery_curve_summary.json :: table[rung=R3_VIS,level=17,probe=linear,n_train=max].accuracy :: sha256 1250d1f00537e0bd :: balanced accuracy, mean +- SD over 5 runs
\newcommand{\RecoveryAccTwoByOneeighteight}{0.8688\,\ensuremath{\pm}\,0.0005}  % experiments/FIGS/data/label_recovery_curve_v1err/analysis/label_recovery_curve_summary.json :: table[rung=R1_Q1,level=188,probe=linear,n_train=max].accuracy :: sha256 1250d1f00537e0bd :: balanced accuracy, mean +- SD over 5 runs
\newcommand{\RecoveryAccTwoByOnesixtwo}{0.8703\,\ensuremath{\pm}\,0.0008}  % experiments/FIGS/data/label_recovery_curve_v1err/analysis/label_recovery_curve_summary.json :: table[rung=R1_Q1,level=162,probe=linear,n_train=max].accuracy :: sha256 1250d1f00537e0bd :: balanced accuracy, mean +- SD over 5 runs
\newcommand{\RecoveryAccTwoByFourthree}{0.8709\,\ensuremath{\pm}\,0.0011}  % experiments/FIGS/data/label_recovery_curve_v1err/analysis/label_recovery_curve_summary.json :: table[rung=R1_Q1,level=43,probe=linear,n_train=max].accuracy :: sha256 1250d1f00537e0bd :: balanced accuracy, mean +- SD over 5 runs
\newcommand{\RecoveryAccTwoByOneseven}{0.8713\,\ensuremath{\pm}\,0.0002}  % experiments/FIGS/data/label_recovery_curve_v1err/analysis/label_recovery_curve_summary.json :: table[rung=R1_Q1,level=17,probe=linear,n_train=max].accuracy :: sha256 1250d1f00537e0bd :: balanced accuracy, mean +- SD over 5 runs
\newcommand{\RecoveryOneeighteightOnesixtwoMaxAbs}{0.0033}  % experiments/FIGS/data/label_recovery_curve_v1err/analysis/label_recovery_curve_summary.json :: table[rung=*,level=188/162,probe=linear,n_train=max].accuracy :: sha256 1250d1f00537e0bd :: max |difference of the run means| over the tree
\newcommand{\RecoveryPairedOnesevenLevelOnesevenMinusOneeighteight}{+0.0006~[\ensuremath{-}0.0004, +0.0015]}  % experiments/FIGS/data/label_recovery_curve_v1err/analysis/label_recovery_curve_summary.json :: table[rung=R16_Q1,level=17 minus 188,probe=linear,n_train=max].accuracy :: sha256 1250d1f00537e0bd :: paired by run, mean difference [95% Student-t interval]
\newcommand{\RecoveryPairedOnesevenLevelOnesevenMinusOnesixtwo}{+0.00004~[\ensuremath{-}0.00123, +0.00131]}  % experiments/FIGS/data/label_recovery_curve_v1err/analysis/label_recovery_curve_summary.json :: table[rung=R16_Q1,level=17 minus 162,probe=linear,n_train=max].accuracy :: sha256 1250d1f00537e0bd :: paired by run, mean difference [95% Student-t interval]
\newcommand{\RecoveryPairedFourLevelOnesevenMinusOneeighteight}{+0.0060~[+0.0051, +0.0069]}  % experiments/FIGS/data/label_recovery_curve_v1err/analysis/label_recovery_curve_summary.json :: table[rung=R3_VIS,level=17 minus 188,probe=linear,n_train=max].accuracy :: sha256 1250d1f00537e0bd :: paired by run, mean difference [95% Student-t interval]
\newcommand{\RecoveryPairedFourLevelOnesevenMinusOnesixtwo}{+0.0048~[+0.0028, +0.0068]}  % experiments/FIGS/data/label_recovery_curve_v1err/analysis/label_recovery_curve_summary.json :: table[rung=R3_VIS,level=17 minus 162,probe=linear,n_train=max].accuracy :: sha256 1250d1f00537e0bd :: paired by run, mean difference [95% Student-t interval]
\newcommand{\RecoveryPairedTwoLevelOnesevenMinusOneeighteight}{+0.0025~[+0.0019, +0.0032]}  % experiments/FIGS/data/label_recovery_curve_v1err/analysis/label_recovery_curve_summary.json :: table[rung=R1_Q1,level=17 minus 188,probe=linear,n_train=max].accuracy :: sha256 1250d1f00537e0bd :: paired by run, mean difference [95% Student-t interval]
\newcommand{\RecoveryPairedTwoLevelOnesevenMinusOnesixtwo}{+0.00101~[+0.00003, +0.00198]}  % experiments/FIGS/data/label_recovery_curve_v1err/analysis/label_recovery_curve_summary.json :: table[rung=R1_Q1,level=17 minus 162,probe=linear,n_train=max].accuracy :: sha256 1250d1f00537e0bd :: paired by run, mean difference [95% Student-t interval]
\newcommand{\RecoveryMlpGainOneseven}{+0.0099~[+0.0078, +0.0120]}  % experiments/FIGS/data/label_recovery_curve_v1err/analysis/label_recovery_curve_summary.json :: mlp_minus_linear_at_largest[rung=L188,level=17].per_run :: sha256 1250d1f00537e0bd :: nonlinear minus linear probe, 188 native labels from the 17-class models
\newcommand{\RecoveryLastGainOneseven}{0.0125\,\ensuremath{\pm}\,0.0021}  % experiments/FIGS/data/label_recovery_curve_v1err/analysis/label_recovery_curve_summary.json :: data_gain_last_step[rung=L188,level=17] :: sha256 1250d1f00537e0bd :: balanced accuracy gained from 443000 to 1400000 training jets
\newcommand{\RecoveryLastGainOneeighteight}{0.0147\,\ensuremath{\pm}\,0.0014}  % experiments/FIGS/data/label_recovery_curve_v1err/analysis/label_recovery_curve_summary.json :: data_gain_last_step[rung=L188,level=188] :: sha256 1250d1f00537e0bd :: balanced accuracy gained from 443000 to 1400000 training jets
\newcommand{\RecoveryNTrain}{1{,}400{,}000}  % experiments/FIGS/data/label_recovery_curve_v1err/analysis/label_recovery_curve_summary.json :: sizes (max) :: sha256 1250d1f00537e0bd :: jets the linear probe is fit on at the end of the curve
\newcommand{\RecoveryNTrainMin}{14{,}000}  % experiments/FIGS/data/label_recovery_curve_v1err/analysis/label_recovery_curve_summary.json :: sizes (min) :: sha256 1250d1f00537e0bd
\newcommand{\RecoveryNTrainPrev}{443{,}000}  % experiments/FIGS/data/label_recovery_curve_v1err/analysis/label_recovery_curve_summary.json :: sizes (second largest) :: sha256 1250d1f00537e0bd
\newcommand{\RecoveryNTest}{400{,}000}  % experiments/FIGS/data/label_recovery_curve_v1err/l188-s1/label_recovery_curve.json :: n_test :: sha256 65651e4c06c7766f :: jets the probes are scored on, recorded in all 20 per-run files
\newcommand{\RecoveryMlpHidden}{512}  % experiments/FIGS/data/label_recovery_curve_v1err/l188-s1/label_recovery_curve.json :: mlp.hidden :: sha256 65651e4c06c7766f :: hidden units of the nonlinear probe
\newcommand{\RandNDraws}{three}  % experiments/FIGS/data/probe_ladder_randcontrol_mlp2/analysis/c4_random_control.json :: table[*].draw :: sha256 c675e609ff4fe3c9 :: random partitions drawn
\newcommand{\RandNDrawsCap}{Three}  % experiments/FIGS/data/probe_ladder_randcontrol_mlp2/analysis/c4_random_control.json :: table[*].draw :: sha256 c675e609ff4fe3c9 :: random partitions drawn, capitalised to open a sentence
\newcommand{\RandOmaBvcFourprongDrawOneLinear}{\ensuremath{1.05\times10^{-2}}}  % experiments/FIGS/data/probe_ladder_randcontrol_mlp2/analysis/c4_random_control.json :: table[0].control_log1m_auc :: sha256 c675e609ff4fe3c9 :: 1-AUC = exp of it, one run
\newcommand{\RandOmaBvcFourprongDrawTwoLinear}{\ensuremath{1.11\times10^{-2}}}  % experiments/FIGS/data/probe_ladder_randcontrol_mlp2/analysis/c4_random_control.json :: table[4].control_log1m_auc :: sha256 c675e609ff4fe3c9 :: 1-AUC = exp of it, one run
\newcommand{\RandOmaBvcFourprongDrawThreeLinear}{\ensuremath{9.98\times10^{-3}}}  % experiments/FIGS/data/probe_ladder_randcontrol_mlp2/analysis/c4_random_control.json :: table[8].control_log1m_auc :: sha256 c675e609ff4fe3c9 :: 1-AUC = exp of it, one run
\newcommand{\RandOmaMeanBvcFourprongLinear}{\ensuremath{(1.05 \pm 0.06)\times10^{-2}}}  % experiments/FIGS/data/probe_ladder_randcontrol_mlp2/analysis/c4_random_control.json :: table[task=bvc_4prong,probe=linear].control_log1m_auc :: sha256 c675e609ff4fe3c9 :: 1-AUC, mean +- SD over 3 draws
\newcommand{\RandOmaBvcFourprongDrawOneMlp}{\ensuremath{1.05\times10^{-2}}}  % experiments/FIGS/data/probe_ladder_randcontrol_mlp2/analysis/c4_random_control.json :: table[1].control_log1m_auc :: sha256 c675e609ff4fe3c9 :: 1-AUC = exp of it, one run
\newcommand{\RandOmaBvcFourprongDrawTwoMlp}{\ensuremath{1.04\times10^{-2}}}  % experiments/FIGS/data/probe_ladder_randcontrol_mlp2/analysis/c4_random_control.json :: table[5].control_log1m_auc :: sha256 c675e609ff4fe3c9 :: 1-AUC = exp of it, one run
\newcommand{\RandOmaBvcFourprongDrawThreeMlp}{\ensuremath{9.28\times10^{-3}}}  % experiments/FIGS/data/probe_ladder_randcontrol_mlp2/analysis/c4_random_control.json :: table[9].control_log1m_auc :: sha256 c675e609ff4fe3c9 :: 1-AUC = exp of it, one run
\newcommand{\RandOmaMeanBvcFourprongMlp}{\ensuremath{(1.01 \pm 0.07)\times10^{-2}}}  % experiments/FIGS/data/probe_ladder_randcontrol_mlp2/analysis/c4_random_control.json :: table[task=bvc_4prong,probe=mlp].control_log1m_auc :: sha256 c675e609ff4fe3c9 :: 1-AUC, mean +- SD over 3 draws
\newcommand{\RandOmaVisibleContentDrawOneLinear}{\ensuremath{1.25\times10^{-3}}}  % experiments/FIGS/data/probe_ladder_randcontrol_mlp2/analysis/c4_random_control.json :: table[2].control_log1m_auc :: sha256 c675e609ff4fe3c9 :: 1-AUC = exp of it, one run
\newcommand{\RandOmaVisibleContentDrawTwoLinear}{\ensuremath{9.89\times10^{-4}}}  % experiments/FIGS/data/probe_ladder_randcontrol_mlp2/analysis/c4_random_control.json :: table[6].control_log1m_auc :: sha256 c675e609ff4fe3c9 :: 1-AUC = exp of it, one run
\newcommand{\RandOmaVisibleContentDrawThreeLinear}{\ensuremath{1.08\times10^{-3}}}  % experiments/FIGS/data/probe_ladder_randcontrol_mlp2/analysis/c4_random_control.json :: table[10].control_log1m_auc :: sha256 c675e609ff4fe3c9 :: 1-AUC = exp of it, one run
\newcommand{\RandOmaMeanVisibleContentLinear}{\ensuremath{(1.11 \pm 0.13)\times10^{-3}}}  % experiments/FIGS/data/probe_ladder_randcontrol_mlp2/analysis/c4_random_control.json :: table[task=visible_content,probe=linear].control_log1m_auc :: sha256 c675e609ff4fe3c9 :: 1-AUC, mean +- SD over 3 draws
\newcommand{\RandOmaVisibleContentDrawOneMlp}{\ensuremath{1.21\times10^{-3}}}  % experiments/FIGS/data/probe_ladder_randcontrol_mlp2/analysis/c4_random_control.json :: table[3].control_log1m_auc :: sha256 c675e609ff4fe3c9 :: 1-AUC = exp of it, one run
\newcommand{\RandOmaVisibleContentDrawTwoMlp}{\ensuremath{1.07\times10^{-3}}}  % experiments/FIGS/data/probe_ladder_randcontrol_mlp2/analysis/c4_random_control.json :: table[7].control_log1m_auc :: sha256 c675e609ff4fe3c9 :: 1-AUC = exp of it, one run
\newcommand{\RandOmaVisibleContentDrawThreeMlp}{\ensuremath{1.13\times10^{-3}}}  % experiments/FIGS/data/probe_ladder_randcontrol_mlp2/analysis/c4_random_control.json :: table[11].control_log1m_auc :: sha256 c675e609ff4fe3c9 :: 1-AUC = exp of it, one run
\newcommand{\RandOmaMeanVisibleContentMlp}{\ensuremath{(1.14 \pm 0.07)\times10^{-3}}}  % experiments/FIGS/data/probe_ladder_randcontrol_mlp2/analysis/c4_random_control.json :: table[task=visible_content,probe=mlp].control_log1m_auc :: sha256 c675e609ff4fe3c9 :: 1-AUC, mean +- SD over 3 draws
\newcommand{\RandCostFactorBvcFourprongLinear}{1.49}  % experiments/FIGS/data/probe_ladder_randcontrol_mlp2/analysis/c4_random_control.json :: grouping_cost_post_hoc.tasks.bvc_4prong.linear.factor_control :: sha256 c675e609ff4fe3c9 :: POST HOC: random grouping vs finer models, in 1-AUC
\newcommand{\RandSemCostFactorBvcFourprongLinear}{4.42}  % experiments/FIGS/data/probe_ladder_randcontrol_mlp2/analysis/c4_random_control.json :: grouping_cost_post_hoc.tasks.bvc_4prong.linear.factor_semantic :: sha256 c675e609ff4fe3c9 :: POST HOC: semantic 17-class vs finer models
\newcommand{\RandCostFactorBvcFourprongMlp}{1.41}  % experiments/FIGS/data/probe_ladder_randcontrol_mlp2/analysis/c4_random_control.json :: grouping_cost_post_hoc.tasks.bvc_4prong.mlp.factor_control :: sha256 c675e609ff4fe3c9 :: POST HOC: random grouping vs finer models, in 1-AUC
\newcommand{\RandSemCostFactorBvcFourprongMlp}{3.26}  % experiments/FIGS/data/probe_ladder_randcontrol_mlp2/analysis/c4_random_control.json :: grouping_cost_post_hoc.tasks.bvc_4prong.mlp.factor_semantic :: sha256 c675e609ff4fe3c9 :: POST HOC: semantic 17-class vs finer models
\newcommand{\RandCostFactorVisibleContentLinear}{1.35}  % experiments/FIGS/data/probe_ladder_randcontrol_mlp2/analysis/c4_random_control.json :: grouping_cost_post_hoc.tasks.visible_content.linear.factor_control :: sha256 c675e609ff4fe3c9 :: POST HOC: random grouping vs finer models, in 1-AUC
\newcommand{\RandSemCostFactorVisibleContentLinear}{2.82}  % experiments/FIGS/data/probe_ladder_randcontrol_mlp2/analysis/c4_random_control.json :: grouping_cost_post_hoc.tasks.visible_content.linear.factor_semantic :: sha256 c675e609ff4fe3c9 :: POST HOC: semantic 17-class vs finer models
\newcommand{\RandCostFactorVisibleContentMlp}{1.33}  % experiments/FIGS/data/probe_ladder_randcontrol_mlp2/analysis/c4_random_control.json :: grouping_cost_post_hoc.tasks.visible_content.mlp.factor_control :: sha256 c675e609ff4fe3c9 :: POST HOC: random grouping vs finer models, in 1-AUC
\newcommand{\RandSemCostFactorVisibleContentMlp}{2.37}  % experiments/FIGS/data/probe_ladder_randcontrol_mlp2/analysis/c4_random_control.json :: grouping_cost_post_hoc.tasks.visible_content.mlp.factor_semantic :: sha256 c675e609ff4fe3c9 :: POST HOC: semantic 17-class vs finer models
\newcommand{\MassResNTest}{398{,}625}  % experiments/FIGS/data/mass_resolution/s1.json :: centering_detail.split[2] :: sha256 88a57bbad270d51c :: test jets sigma_eff is computed on, the same in every per-seed file and arm
\newcommand{\MassResNClasses}{184}  % experiments/FIGS/data/mass_resolution/analysis_holm/s7_mass_resolution.json :: provenance.n_classes_used :: sha256 e03cde45ece9f385 :: native classes with enough training jets to centre
\newcommand{\MassResSigmaEffRidgeOneeighteight}{0.3393\,\ensuremath{\pm}\,0.0019}  % experiments/FIGS/data/mass_resolution/analysis_holm/s7_mass_resolution.json :: table[cell=188,probe=ridge].sigma_eff :: sha256 e03cde45ece9f385 :: mean +- SD over 5 seeds
\newcommand{\MassResSigmaEffRidgeOnesixtwo}{0.3302\,\ensuremath{\pm}\,0.0029}  % experiments/FIGS/data/mass_resolution/analysis_holm/s7_mass_resolution.json :: table[cell=162,probe=ridge].sigma_eff :: sha256 e03cde45ece9f385 :: mean +- SD over 5 seeds
\newcommand{\MassResSigmaEffRidgeFourthree}{0.3228\,\ensuremath{\pm}\,0.0019}  % experiments/FIGS/data/mass_resolution/analysis_holm/s7_mass_resolution.json :: table[cell=43,probe=ridge].sigma_eff :: sha256 e03cde45ece9f385 :: mean +- SD over 5 seeds
\newcommand{\MassResSigmaEffRidgeOneseven}{0.300\,\ensuremath{\pm}\,0.010}  % experiments/FIGS/data/mass_resolution/analysis_holm/s7_mass_resolution.json :: table[cell=17,probe=ridge].sigma_eff :: sha256 e03cde45ece9f385 :: mean +- SD over 5 seeds
\newcommand{\MassResSigmaEffRidgeOnesixtwoMass}{0.2238\,\ensuremath{\pm}\,0.0019}  % experiments/FIGS/data/mass_resolution/analysis_holm/s7_mass_resolution.json :: table[cell=162+mass,probe=ridge].sigma_eff :: sha256 e03cde45ece9f385 :: mean +- SD over 5 seeds
\newcommand{\MassResSigmaEffRidgeOnesevenMass}{0.2170\,\ensuremath{\pm}\,0.0034}  % experiments/FIGS/data/mass_resolution/analysis_holm/s7_mass_resolution.json :: table[cell=17+mass,probe=ridge].sigma_eff :: sha256 e03cde45ece9f385 :: mean +- SD over 5 seeds
\newcommand{\MassResGainRidgeOnesixtwo}{\ensuremath{-}0.1063}  % experiments/FIGS/data/mass_resolution/analysis_holm/s7_mass_resolution.json :: table[cell=162+mass/162,probe=ridge].sigma_eff :: sha256 e03cde45ece9f385 :: sigma_eff with the mass output minus without, difference of the seed means; negative is better
\newcommand{\MassResGainRidgeOneseven}{\ensuremath{-}0.0831}  % experiments/FIGS/data/mass_resolution/analysis_holm/s7_mass_resolution.json :: table[cell=17+mass/17,probe=ridge].sigma_eff :: sha256 e03cde45ece9f385 :: sigma_eff with the mass output minus without, difference of the seed means; negative is better
\newcommand{\MassResDidRidge}{\ensuremath{-}0.0233}  % experiments/FIGS/data/mass_resolution/analysis_holm/s7_mass_resolution.json :: table[cell=162+mass/162/17+mass/17,probe=ridge].sigma_eff :: sha256 e03cde45ece9f385 :: gain at 162 classes minus gain at 17
\newcommand{\MassResSigmaEffMlpOneeighteight}{0.2463\,\ensuremath{\pm}\,0.0022}  % experiments/FIGS/data/mass_resolution/analysis_holm/s7_mass_resolution.json :: table[cell=188,probe=mlp].sigma_eff :: sha256 e03cde45ece9f385 :: mean +- SD over 5 seeds
\newcommand{\MassResSigmaEffMlpOnesixtwo}{0.2438\,\ensuremath{\pm}\,0.0018}  % experiments/FIGS/data/mass_resolution/analysis_holm/s7_mass_resolution.json :: table[cell=162,probe=mlp].sigma_eff :: sha256 e03cde45ece9f385 :: mean +- SD over 5 seeds
\newcommand{\MassResSigmaEffMlpFourthree}{0.2407\,\ensuremath{\pm}\,0.0009}  % experiments/FIGS/data/mass_resolution/analysis_holm/s7_mass_resolution.json :: table[cell=43,probe=mlp].sigma_eff :: sha256 e03cde45ece9f385 :: mean +- SD over 5 seeds
\newcommand{\MassResSigmaEffMlpOneseven}{0.2358\,\ensuremath{\pm}\,0.0025}  % experiments/FIGS/data/mass_resolution/analysis_holm/s7_mass_resolution.json :: table[cell=17,probe=mlp].sigma_eff :: sha256 e03cde45ece9f385 :: mean +- SD over 5 seeds
\newcommand{\MassResSigmaEffMlpOnesixtwoMass}{0.1943\,\ensuremath{\pm}\,0.0001}  % experiments/FIGS/data/mass_resolution/analysis_holm/s7_mass_resolution.json :: table[cell=162+mass,probe=mlp].sigma_eff :: sha256 e03cde45ece9f385 :: mean +- SD over 5 seeds
\newcommand{\MassResSigmaEffMlpOnesevenMass}{0.1949\,\ensuremath{\pm}\,0.0006}  % experiments/FIGS/data/mass_resolution/analysis_holm/s7_mass_resolution.json :: table[cell=17+mass,probe=mlp].sigma_eff :: sha256 e03cde45ece9f385 :: mean +- SD over 5 seeds
\newcommand{\MassResGainMlpOnesixtwo}{\ensuremath{-}0.0496}  % experiments/FIGS/data/mass_resolution/analysis_holm/s7_mass_resolution.json :: table[cell=162+mass/162,probe=mlp].sigma_eff :: sha256 e03cde45ece9f385 :: sigma_eff with the mass output minus without, difference of the seed means; negative is better
\newcommand{\MassResGainMlpOneseven}{\ensuremath{-}0.0409}  % experiments/FIGS/data/mass_resolution/analysis_holm/s7_mass_resolution.json :: table[cell=17+mass/17,probe=mlp].sigma_eff :: sha256 e03cde45ece9f385 :: sigma_eff with the mass output minus without, difference of the seed means; negative is better
\newcommand{\MassResDidMlp}{\ensuremath{-}0.0087}  % experiments/FIGS/data/mass_resolution/analysis_holm/s7_mass_resolution.json :: table[cell=162+mass/162/17+mass/17,probe=mlp].sigma_eff :: sha256 e03cde45ece9f385 :: gain at 162 classes minus gain at 17
\newcommand{\MassResTargetSigmaEff}{0.2986}  % experiments/FIGS/data/mass_resolution/analysis_holm/s7_mass_resolution.json :: table[*].target_sigma_eff :: sha256 e03cde45ece9f385 :: sigma_eff of the class-centred target itself (a model that predicts the class mean)
\newcommand{\MassResRidgeRtwoMin}{20\%}  % experiments/FIGS/data/mass_resolution/s1.json :: arms.{l188,l162,r42q1,r16q1}-s*.ridge.val_r2 :: sha256 88a57bbad270d51c :: validation R^2 of the ridge probe, smallest of 20 models without the mass output, all 5 per-seed files
\newcommand{\MassResRidgeRtwoMax}{30\%}  % experiments/FIGS/data/mass_resolution/s1.json :: arms.{l188,l162,r42q1,r16q1}-s*.ridge.val_r2 :: sha256 88a57bbad270d51c :: validation R^2 of the ridge probe, largest of 20 models without the mass output, all 5 per-seed files
\newcommand{\MassResSigmaEffOnesixtwo}{0.3302}  % experiments/FIGS/data/mass_resolution/analysis_holm/s7_mass_resolution.json :: table[cell=162,probe=ridge].sigma_eff :: sha256 e03cde45ece9f385 :: mean over 5 seeds
\newcommand{\MassResSdOnesixtwo}{0.626}  % experiments/FIGS/data/mass_resolution/analysis_holm/s7_mass_resolution.json :: table[cell=162,probe=ridge].sd :: sha256 e03cde45ece9f385 :: mean over 5 seeds
\newcommand{\MassResTailFractionOnesixtwo}{0.062}  % experiments/FIGS/data/mass_resolution/analysis_holm/s7_mass_resolution.json :: table[cell=162,probe=ridge].tail_fraction :: sha256 e03cde45ece9f385 :: mean over 5 seeds
\newcommand{\MassResSigmaEffOnesixtwoMass}{0.2238}  % experiments/FIGS/data/mass_resolution/analysis_holm/s7_mass_resolution.json :: table[cell=162+mass,probe=ridge].sigma_eff :: sha256 e03cde45ece9f385 :: mean over 5 seeds
\newcommand{\MassResSdOnesixtwoMass}{0.494}  % experiments/FIGS/data/mass_resolution/analysis_holm/s7_mass_resolution.json :: table[cell=162+mass,probe=ridge].sd :: sha256 e03cde45ece9f385 :: mean over 5 seeds
\newcommand{\MassResTailFractionOnesixtwoMass}{0.034}  % experiments/FIGS/data/mass_resolution/analysis_holm/s7_mass_resolution.json :: table[cell=162+mass,probe=ridge].tail_fraction :: sha256 e03cde45ece9f385 :: mean over 5 seeds
\newcommand{\MassResSigmaEffOneseven}{0.3001}  % experiments/FIGS/data/mass_resolution/analysis_holm/s7_mass_resolution.json :: table[cell=17,probe=ridge].sigma_eff :: sha256 e03cde45ece9f385 :: mean over 5 seeds
\newcommand{\MassResSdOneseven}{0.601}  % experiments/FIGS/data/mass_resolution/analysis_holm/s7_mass_resolution.json :: table[cell=17,probe=ridge].sd :: sha256 e03cde45ece9f385 :: mean over 5 seeds
\newcommand{\MassResTailFractionOneseven}{0.056}  % experiments/FIGS/data/mass_resolution/analysis_holm/s7_mass_resolution.json :: table[cell=17,probe=ridge].tail_fraction :: sha256 e03cde45ece9f385 :: mean over 5 seeds
\newcommand{\MassResSigmaEffOnesevenMass}{0.2170}  % experiments/FIGS/data/mass_resolution/analysis_holm/s7_mass_resolution.json :: table[cell=17+mass,probe=ridge].sigma_eff :: sha256 e03cde45ece9f385 :: mean over 5 seeds
\newcommand{\MassResSdOnesevenMass}{0.490}  % experiments/FIGS/data/mass_resolution/analysis_holm/s7_mass_resolution.json :: table[cell=17+mass,probe=ridge].sd :: sha256 e03cde45ece9f385 :: mean over 5 seeds
\newcommand{\MassResTailFractionOnesevenMass}{0.033}  % experiments/FIGS/data/mass_resolution/analysis_holm/s7_mass_resolution.json :: table[cell=17+mass,probe=ridge].tail_fraction :: sha256 e03cde45ece9f385 :: mean over 5 seeds
\newcommand{\MassResSigmaEffOneeighteight}{0.3393}  % experiments/FIGS/data/mass_resolution/analysis_holm/s7_mass_resolution.json :: table[cell=188,probe=ridge].sigma_eff :: sha256 e03cde45ece9f385 :: mean over 5 seeds
\newcommand{\MassResSdOneeighteight}{0.631}  % experiments/FIGS/data/mass_resolution/analysis_holm/s7_mass_resolution.json :: table[cell=188,probe=ridge].sd :: sha256 e03cde45ece9f385 :: mean over 5 seeds
\newcommand{\MassResTailFractionOneeighteight}{0.064}  % experiments/FIGS/data/mass_resolution/analysis_holm/s7_mass_resolution.json :: table[cell=188,probe=ridge].tail_fraction :: sha256 e03cde45ece9f385 :: mean over 5 seeds
\newcommand{\MassResSigmaEffFourthree}{0.3228}  % experiments/FIGS/data/mass_resolution/analysis_holm/s7_mass_resolution.json :: table[cell=43,probe=ridge].sigma_eff :: sha256 e03cde45ece9f385 :: mean over 5 seeds
\newcommand{\MassResSdFourthree}{0.618}  % experiments/FIGS/data/mass_resolution/analysis_holm/s7_mass_resolution.json :: table[cell=43,probe=ridge].sd :: sha256 e03cde45ece9f385 :: mean over 5 seeds
\newcommand{\MassResTailFractionFourthree}{0.060}  % experiments/FIGS/data/mass_resolution/analysis_holm/s7_mass_resolution.json :: table[cell=43,probe=ridge].tail_fraction :: sha256 e03cde45ece9f385 :: mean over 5 seeds
\newcommand{\AojEffPercent}{1\%}  % experiments/FIGS/data/aoj_full_v1/fit_v6/results.json :: eff :: sha256 2416fc21f93fcfa6 :: the score cut's pass fraction outside the mass windows
\newcommand{\AojPtMin}{500}  % experiments/AOJ/peak_fit.py :: PT_RANGE[0] :: sha256 cad06229a74532be :: GeV
\newcommand{\AojPtMax}{2{,}500}  % experiments/AOJ/peak_fit.py :: PT_RANGE[1] :: sha256 cad06229a74532be :: GeV
\newcommand{\AojRhoMin}{\ensuremath{-}5.5}  % experiments/AOJ/peak_fit.py :: RHO_RANGE[0] :: sha256 cad06229a74532be
\newcommand{\AojRhoMax}{\ensuremath{-}2.0}  % experiments/AOJ/peak_fit.py :: RHO_RANGE[1] :: sha256 cad06229a74532be
\newcommand{\AojWindowWLo}{65}  % experiments/AOJ/peak_fit.py :: PEAKS[W].window[0] :: sha256 cad06229a74532be :: GeV
\newcommand{\AojWindowWHi}{105}  % experiments/AOJ/peak_fit.py :: PEAKS[W].window[1] :: sha256 cad06229a74532be :: GeV
\newcommand{\AojWindowTopLo}{140}  % experiments/AOJ/peak_fit.py :: PEAKS[top].window[0] :: sha256 cad06229a74532be :: GeV
\newcommand{\AojWindowTopHi}{220}  % experiments/AOJ/peak_fit.py :: PEAKS[top].window[1] :: sha256 cad06229a74532be :: GeV
\newcommand{\AojFitLo}{105}  % experiments/AOJ/peak_fit.py :: PEAKS[top].fit_range[0] :: sha256 cad06229a74532be :: GeV
\newcommand{\AojFitHi}{300}  % experiments/AOJ/peak_fit.py :: PEAKS[top].fit_range[1] :: sha256 cad06229a74532be :: GeV
\newcommand{\AojMassBin}{5}  % experiments/AOJ/peak_fit.py :: MASS_BIN :: sha256 cad06229a74532be :: GeV
\newcommand{\AojMapOrderRho}{3}  % experiments/AOJ/peak_fit.py :: MAP_ORDER[0] :: sha256 cad06229a74532be :: polynomial order in rho
\newcommand{\AojMapOrderPt}{2}  % experiments/AOJ/peak_fit.py :: MAP_ORDER[1] :: sha256 cad06229a74532be :: polynomial order in ln pT
\newcommand{\AojNRunsWord}{five}  % experiments/FIGS/data/aoj_full_v1/analysis_v6/aoj_top.json :: per_label_set.label_sets.*.signal_yield.n :: sha256 5dfef902a5fe75ba :: pretraining runs per vocabulary
\newcommand{\AojNJetsRhoWindow}{9{,}441{,}209}  % experiments/FIGS/data/aoj_full_v1/fit_v6/results.json :: n_jets :: sha256 2416fc21f93fcfa6 :: staged jets inside the rho window, all files; not the jets in the mass fit
\newcommand{\AojYieldOneeighteight}{2{,}400\,\ensuremath{\pm}\,500}  % experiments/FIGS/data/aoj_full_v1/analysis_v6/aoj_top.json :: per_label_set.label_sets.188.signal_yield :: sha256 5dfef902a5fe75ba :: mean +- SD over 5 runs, shared shape
\newcommand{\AojStatErrOneeighteight}{211}  % experiments/FIGS/data/aoj_full_v1/analysis_v6/aoj_top.json :: per_label_set.label_sets.188.median_stat_err :: sha256 5dfef902a5fe75ba :: median per-fit statistical error
\newcommand{\AojYieldOwnShapeOneeighteight}{2{,}400\,\ensuremath{\pm}\,600}  % experiments/FIGS/data/aoj_full_v1/analysis_v6/aoj_top.json :: per_label_set.label_sets.188.by_shape.own_shape :: sha256 5dfef902a5fe75ba :: every fit with its own floated shape
\newcommand{\AojSpreadOverStatOneeighteight}{2.4}  % experiments/FIGS/data/aoj_full_v1/analysis_v6/aoj_top.json :: per_label_set.label_sets.188.signal_yield.sd / median_stat_err :: sha256 5dfef902a5fe75ba :: spread between runs over the median per-fit statistical error
\newcommand{\AojYieldOnesixtwo}{2{,}400\,\ensuremath{\pm}\,900}  % experiments/FIGS/data/aoj_full_v1/analysis_v6/aoj_top.json :: per_label_set.label_sets.162.signal_yield :: sha256 5dfef902a5fe75ba :: mean +- SD over 5 runs, shared shape
\newcommand{\AojStatErrOnesixtwo}{199}  % experiments/FIGS/data/aoj_full_v1/analysis_v6/aoj_top.json :: per_label_set.label_sets.162.median_stat_err :: sha256 5dfef902a5fe75ba :: median per-fit statistical error
\newcommand{\AojYieldOwnShapeOnesixtwo}{2{,}200\,\ensuremath{\pm}\,800}  % experiments/FIGS/data/aoj_full_v1/analysis_v6/aoj_top.json :: per_label_set.label_sets.162.by_shape.own_shape :: sha256 5dfef902a5fe75ba :: every fit with its own floated shape
\newcommand{\AojSpreadOverStatOnesixtwo}{4.5}  % experiments/FIGS/data/aoj_full_v1/analysis_v6/aoj_top.json :: per_label_set.label_sets.162.signal_yield.sd / median_stat_err :: sha256 5dfef902a5fe75ba :: spread between runs over the median per-fit statistical error
\newcommand{\AojYieldFourthree}{2{,}400\,\ensuremath{\pm}\,800}  % experiments/FIGS/data/aoj_full_v1/analysis_v6/aoj_top.json :: per_label_set.label_sets.43.signal_yield :: sha256 5dfef902a5fe75ba :: mean +- SD over 5 runs, shared shape
\newcommand{\AojStatErrFourthree}{184}  % experiments/FIGS/data/aoj_full_v1/analysis_v6/aoj_top.json :: per_label_set.label_sets.43.median_stat_err :: sha256 5dfef902a5fe75ba :: median per-fit statistical error
\newcommand{\AojYieldOwnShapeFourthree}{2{,}300\,\ensuremath{\pm}\,800}  % experiments/FIGS/data/aoj_full_v1/analysis_v6/aoj_top.json :: per_label_set.label_sets.43.by_shape.own_shape :: sha256 5dfef902a5fe75ba :: every fit with its own floated shape
\newcommand{\AojSpreadOverStatFourthree}{4.2}  % experiments/FIGS/data/aoj_full_v1/analysis_v6/aoj_top.json :: per_label_set.label_sets.43.signal_yield.sd / median_stat_err :: sha256 5dfef902a5fe75ba :: spread between runs over the median per-fit statistical error
\newcommand{\AojYieldOneseven}{2{,}350\,\ensuremath{\pm}\,240}  % experiments/FIGS/data/aoj_full_v1/analysis_v6/aoj_top.json :: per_label_set.label_sets.17.signal_yield :: sha256 5dfef902a5fe75ba :: mean +- SD over 5 runs, shared shape
\newcommand{\AojStatErrOneseven}{201}  % experiments/FIGS/data/aoj_full_v1/analysis_v6/aoj_top.json :: per_label_set.label_sets.17.median_stat_err :: sha256 5dfef902a5fe75ba :: median per-fit statistical error
\newcommand{\AojYieldOwnShapeOneseven}{2{,}480\,\ensuremath{\pm}\,290}  % experiments/FIGS/data/aoj_full_v1/analysis_v6/aoj_top.json :: per_label_set.label_sets.17.by_shape.own_shape :: sha256 5dfef902a5fe75ba :: every fit with its own floated shape
\newcommand{\AojSpreadOverStatOneseven}{1.2}  % experiments/FIGS/data/aoj_full_v1/analysis_v6/aoj_top.json :: per_label_set.label_sets.17.signal_yield.sd / median_stat_err :: sha256 5dfef902a5fe75ba :: spread between runs over the median per-fit statistical error
\newcommand{\AojYieldOnesixtwoMass}{2{,}420\,\ensuremath{\pm}\,240}  % experiments/FIGS/data/aoj_full_v1/analysis_v6/aoj_top.json :: per_label_set.label_sets.162+mass.signal_yield :: sha256 5dfef902a5fe75ba :: mean +- SD over 5 runs, shared shape
\newcommand{\AojStatErrOnesixtwoMass}{179}  % experiments/FIGS/data/aoj_full_v1/analysis_v6/aoj_top.json :: per_label_set.label_sets.162+mass.median_stat_err :: sha256 5dfef902a5fe75ba :: median per-fit statistical error
\newcommand{\AojYieldOwnShapeOnesixtwoMass}{2{,}350\,\ensuremath{\pm}\,320}  % experiments/FIGS/data/aoj_full_v1/analysis_v6/aoj_top.json :: per_label_set.label_sets.162+mass.by_shape.own_shape :: sha256 5dfef902a5fe75ba :: every fit with its own floated shape
\newcommand{\AojSpreadOverStatOnesixtwoMass}{1.3}  % experiments/FIGS/data/aoj_full_v1/analysis_v6/aoj_top.json :: per_label_set.label_sets.162+mass.signal_yield.sd / median_stat_err :: sha256 5dfef902a5fe75ba :: spread between runs over the median per-fit statistical error
\newcommand{\AojYieldOnesevenMass}{2{,}500\,\ensuremath{\pm}\,600}  % experiments/FIGS/data/aoj_full_v1/analysis_v6/aoj_top.json :: per_label_set.label_sets.17+mass.signal_yield :: sha256 5dfef902a5fe75ba :: mean +- SD over 5 runs, shared shape
\newcommand{\AojStatErrOnesevenMass}{200}  % experiments/FIGS/data/aoj_full_v1/analysis_v6/aoj_top.json :: per_label_set.label_sets.17+mass.median_stat_err :: sha256 5dfef902a5fe75ba :: median per-fit statistical error
\newcommand{\AojYieldOwnShapeOnesevenMass}{2{,}900\,\ensuremath{\pm}\,500}  % experiments/FIGS/data/aoj_full_v1/analysis_v6/aoj_top.json :: per_label_set.label_sets.17+mass.by_shape.own_shape :: sha256 5dfef902a5fe75ba :: every fit with its own floated shape
\newcommand{\AojSpreadOverStatOnesevenMass}{2.9}  % experiments/FIGS/data/aoj_full_v1/analysis_v6/aoj_top.json :: per_label_set.label_sets.17+mass.signal_yield.sd / median_stat_err :: sha256 5dfef902a5fe75ba :: spread between runs over the median per-fit statistical error
\newcommand{\AojOwnShapeShiftMin}{\ensuremath{-}7\%}  % experiments/FIGS/data/aoj_full_v1/analysis_v6/aoj_top.json :: per_label_set.label_sets.*.by_shape.own_shape.mean / signal_yield.mean - 1 :: sha256 5dfef902a5fe75ba :: the vocabulary means with each model's own shape, relative to the shared shape
\newcommand{\AojOwnShapeShiftMax}{+13\%}  % experiments/FIGS/data/aoj_full_v1/analysis_v6/aoj_top.json :: per_label_set.label_sets.*.by_shape.own_shape.mean / signal_yield.mean - 1 :: sha256 5dfef902a5fe75ba :: the vocabulary means with each model's own shape, relative to the shared shape
\newcommand{\AojPooledMean}{182.8}  % experiments/FIGS/data/aoj_full_v1/analysis_v6/aoj_top.json :: per_label_set.shapes.pooled :: sha256 5dfef902a5fe75ba :: GeV
\newcommand{\AojPooledWidth}{11.7}  % experiments/FIGS/data/aoj_full_v1/analysis_v6/aoj_top.json :: per_label_set.shapes.pooled :: sha256 5dfef902a5fe75ba :: GeV
\newcommand{\AojPoolSize}{30}  % experiments/FIGS/data/aoj_full_v1/analysis_v6/aoj_top.json :: per_label_set.shapes.pool (count) :: sha256 5dfef902a5fe75ba :: pretrained models the shared shape is fitted over
\newcommand{\AojNModels}{30}  % experiments/FIGS/data/aoj_full_v1/fit_v6/results.json :: models (pretrained) :: sha256 2416fc21f93fcfa6 :: fits
\newcommand{\AojOwnMeanMin}{181.6}  % experiments/FIGS/data/aoj_full_v1/fit_v6/results.json :: models.*.top.own_shape_fit.mean (min) :: sha256 2416fc21f93fcfa6 :: GeV
\newcommand{\AojOwnMeanMax}{184.0}  % experiments/FIGS/data/aoj_full_v1/fit_v6/results.json :: models.*.top.own_shape_fit.mean (max) :: sha256 2416fc21f93fcfa6 :: GeV
\newcommand{\AojOwnWidthMin}{10.0}  % experiments/FIGS/data/aoj_full_v1/fit_v6/results.json :: models.*.top.own_shape_fit.width (min) :: sha256 2416fc21f93fcfa6 :: GeV
\newcommand{\AojOwnWidthMax}{13.8}  % experiments/FIGS/data/aoj_full_v1/fit_v6/results.json :: models.*.top.own_shape_fit.width (max) :: sha256 2416fc21f93fcfa6 :: GeV
\newcommand{\AojShapeStepMean}{0.7}  % experiments/FIGS/data/aoj_full_v1/fit_v6/results.json :: pooled_shape.systematic_step.mean :: sha256 2416fc21f93fcfa6 :: GeV: the SD of the models' own floated values, one shift for every model together
\newcommand{\AojShapeSystMeanMax}{1\%}  % experiments/FIGS/data/aoj_full_v1/fit_v6/results.json :: models.*.top.shape_systematic.pooled_mean_(down|up) / signal_yield - 1 (largest |.|) :: sha256 2416fc21f93fcfa6 :: the largest yield change of any model under that shift
\newcommand{\AojShapeStepWidth}{1.0}  % experiments/FIGS/data/aoj_full_v1/fit_v6/results.json :: pooled_shape.systematic_step.width :: sha256 2416fc21f93fcfa6 :: GeV: the SD of the models' own floated values, one shift for every model together
\newcommand{\AojShapeSystWidthMax}{12\%}  % experiments/FIGS/data/aoj_full_v1/fit_v6/results.json :: models.*.top.shape_systematic.pooled_width_(down|up) / signal_yield - 1 (largest |.|) :: sha256 2416fc21f93fcfa6 :: the largest yield change of any model under that shift
\newcommand{\AojLeakEpsHi}{0.6}  % experiments/FIGS/data/aoj_full_v1/fit_v6/results.json :: models.*.top.leak_systematic.eps_ref (max) :: sha256 2416fc21f93fcfa6 :: the reference's assumed efficiency for data tops, upper variation
\newcommand{\AojLeakEpsLo}{0.4}  % experiments/FIGS/data/aoj_full_v1/fit_v6/results.json :: models.*.top.leak_systematic.eps_ref (min) :: sha256 2416fc21f93fcfa6 :: the same, lower variation
\newcommand{\AojLeakShiftMin}{+97}  % experiments/FIGS/data/aoj_full_v1/fit_v6/results.json :: models.*.top.leak_systematic.shift['0.6'] (min) :: sha256 2416fc21f93fcfa6 :: tops added to a yield
\newcommand{\AojLeakShiftMax}{+271}  % experiments/FIGS/data/aoj_full_v1/fit_v6/results.json :: models.*.top.leak_systematic.shift['0.4'] (max) :: sha256 2416fc21f93fcfa6 :: tops added to a yield
\newcommand{\AojFitPMin}{0.010}  % experiments/FIGS/data/aoj_full_v1/fit_v6/results.json :: models.*.top.asymptotic_p (min) :: sha256 2416fc21f93fcfa6 :: goodness of fit of the working-point fit, pretrained models
\newcommand{\AojFitPMax}{0.974}  % experiments/FIGS/data/aoj_full_v1/fit_v6/results.json :: models.*.top.asymptotic_p (max) :: sha256 2416fc21f93fcfa6
\newcommand{\AojWeakestYield}{885}  % experiments/FIGS/data/aoj_full_v1/fit_v6/results.json :: models.l162-s2.top.signal_yield :: sha256 2416fc21f93fcfa6 :: the smallest yield of any pretrained model
\newcommand{\AojWeakestYieldErr}{211}  % experiments/FIGS/data/aoj_full_v1/fit_v6/results.json :: models.l162-s2.top.signal_yield_err :: sha256 2416fc21f93fcfa6
\newcommand{\AojWeakestFitP}{0.94}  % experiments/FIGS/data/aoj_full_v1/fit_v6/results.json :: models.l162-s2.top.asymptotic_p :: sha256 2416fc21f93fcfa6
\newcommand{\AojSidebandZMin}{\ensuremath{-}0.43}  % experiments/FIGS/data/aoj_full_v1/fit_v6/histograms.npz :: experiments/AOJ/sideband_residuals.py: z (min over the pretrained models) :: sha256 2bd1e3d5b32d4a12 :: passing jets minus the fit, summed from the fit's lower edge to the top window
\newcommand{\AojSidebandZMax}{+0.58}  % experiments/FIGS/data/aoj_full_v1/fit_v6/histograms.npz :: experiments/AOJ/sideband_residuals.py: z (max over the pretrained models) :: sha256 2bd1e3d5b32d4a12
\newcommand{\AojWeakestSidebandZ}{0.00}  % experiments/FIGS/data/aoj_full_v1/fit_v6/histograms.npz :: experiments/AOJ/sideband_residuals.py: z of l162-s2 :: sha256 2bd1e3d5b32d4a12 :: the smallest-yield model
\newcommand{\AojSidebandMaxPull}{2.7}  % experiments/FIGS/data/aoj_full_v1/fit_v6/histograms.npz :: experiments/AOJ/sideband_residuals.py: max_abs_pull (max over the pretrained models) :: sha256 2bd1e3d5b32d4a12 :: largest single-bin pull below the top window, any pretrained model
\newcommand{\AojSharedBinLo}{125}  % experiments/FIGS/data/aoj_full_v1/fit_v6/histograms.npz :: the bin holding the most models' largest pull :: sha256 2bd1e3d5b32d4a12 :: GeV
\newcommand{\AojSharedBinHi}{130}  % experiments/FIGS/data/aoj_full_v1/fit_v6/histograms.npz :: its upper edge :: sha256 2bd1e3d5b32d4a12 :: GeV
\newcommand{\AojSharedBinNModels}{18}  % experiments/FIGS/data/aoj_full_v1/fit_v6/histograms.npz :: models whose largest pull is in that bin :: sha256 2bd1e3d5b32d4a12
\newcommand{\AojSharedBinMeanPull}{+1.4}  % experiments/FIGS/data/aoj_full_v1/fit_v6/histograms.npz :: experiments/AOJ/sideband_residuals.py mean_pulls: that bin :: sha256 2bd1e3d5b32d4a12 :: mean pull over the pretrained models
\newcommand{\AojSharedBinRank}{third}  % experiments/FIGS/data/aoj_full_v1/fit_v6/histograms.npz :: experiments/AOJ/sideband_residuals.py mean_pulls: rank of that bin's |mean pull| :: sha256 2bd1e3d5b32d4a12
\newcommand{\AojFitNBins}{39}  % experiments/FIGS/data/aoj_full_v1/fit_v6/histograms.npz :: mass bins of the fit range :: sha256 2bd1e3d5b32d4a12
\newcommand{\AojPairedYieldOnesixtwoOverOneeighteight}{0.94~[0.45, 1.95]}  % experiments/FIGS/data/aoj_full_v1/analysis_v6/aoj_top.json :: per_label_set.label_sets.{162,188}.signal_yields :: sha256 5dfef902a5fe75ba :: yield ratio paired by run, geometric mean [95% Student-t interval]
\newcommand{\AojPairedYieldFourthreeOverOneeighteight}{0.97~[0.81, 1.16]}  % experiments/FIGS/data/aoj_full_v1/analysis_v6/aoj_top.json :: per_label_set.label_sets.{43,188}.signal_yields :: sha256 5dfef902a5fe75ba :: yield ratio paired by run, geometric mean [95% Student-t interval]
\newcommand{\AojPairedYieldOnesevenOverOneeighteight}{1.01~[0.75, 1.35]}  % experiments/FIGS/data/aoj_full_v1/analysis_v6/aoj_top.json :: per_label_set.label_sets.{17,188}.signal_yields :: sha256 5dfef902a5fe75ba :: yield ratio paired by run, geometric mean [95% Student-t interval]
\newcommand{\AojPairedYieldOnesevenOverOnesixtwo}{1.07~[0.51, 2.24]}  % experiments/FIGS/data/aoj_full_v1/analysis_v6/aoj_top.json :: per_label_set.label_sets.{17,162}.signal_yields :: sha256 5dfef902a5fe75ba :: yield ratio paired by run, geometric mean [95% Student-t interval]
\newcommand{\AojPairedYieldOnesixtwoMassOverOnesixtwo}{1.10~[0.57, 2.14]}  % experiments/FIGS/data/aoj_full_v1/analysis_v6/aoj_top.json :: per_label_set.label_sets.{162+mass,162}.signal_yields :: sha256 5dfef902a5fe75ba :: yield ratio paired by run, geometric mean [95% Student-t interval]
\newcommand{\AojPairedYieldOnesevenMassOverOneseven}{1.05~[0.74, 1.49]}  % experiments/FIGS/data/aoj_full_v1/analysis_v6/aoj_top.json :: per_label_set.label_sets.{17+mass,17}.signal_yields :: sha256 5dfef902a5fe75ba :: yield ratio paired by run, geometric mean [95% Student-t interval]
\newcommand{\AojLeakPairedShiftMax}{2\%}  % experiments/FIGS/data/aoj_full_v1/fit_v6/results.json :: models.*.top.leak_systematic.signal_yield against signal_yield, paired ratios :: sha256 2416fc21f93fcfa6 :: largest change of a paired yield ratio when the fail-region tops are scaled up
\newcommand{\AojRefEffPercent}{1.38\%}  % experiments/FIGS/data/aoj_full_v1/fit_v6/results.json :: reference.top.data_efficiency :: sha256 2416fc21f93fcfa6 :: fraction of all data jets passing its cut
\newcommand{\AojRefFitP}{\ensuremath{2.56\times10^{-11}}}  % experiments/FIGS/data/aoj_full_v1/fit_v6/results.json :: reference.top.asymptotic_p :: sha256 2416fc21f93fcfa6 :: goodness of fit of the reference's own top fit
\newcommand{\AojRefValidationP}{\ensuremath{p \leq 0.005} (0 of 200 toys reach the observed deviance)}  % experiments/FIGS/data/aoj_full_v1/fit_v6/results.json :: reference.top.validation.toy_p, n_toys :: sha256 2416fc21f93fcfa6 :: toy p-value of the background validation, (k + 1) / (n + 1)
\newcommand{\AojRefBandSignalZ}{2.5}  % experiments/FIGS/data/aoj_full_v1/fit_v6/results.json :: reference.top.validation.band_signal_z :: sha256 2416fc21f93fcfa6 :: significance of a peak of the score's own shape in the validation band
\newcommand{\AojModelEffMin}{1.01\%}  % experiments/FIGS/data/aoj_full_v1/fit_v6/results.json :: models.*.top.data_efficiency (min) :: sha256 2416fc21f93fcfa6 :: pretrained models
\newcommand{\AojModelEffMax}{1.05\%}  % experiments/FIGS/data/aoj_full_v1/fit_v6/results.json :: models.*.top.data_efficiency (max) :: sha256 2416fc21f93fcfa6 :: pretrained models
\newcommand{\AojModelsValidationPMin}{0.11}  % experiments/FIGS/data/aoj_full_v1/fit_v6/results.json :: models.*.top.validation.toy_p (min) :: sha256 2416fc21f93fcfa6 :: pretrained models
\newcommand{\AojRefYield}{29{,}143}  % experiments/FIGS/data/aoj_full_v1/fit_v6/results.json :: reference.top.signal_yield :: sha256 2416fc21f93fcfa6
\newcommand{\AojRefYieldErr}{865}  % experiments/FIGS/data/aoj_full_v1/fit_v6/results.json :: reference.top.signal_yield_err :: sha256 2416fc21f93fcfa6
\newcommand{\AojRefMass}{181.0}  % experiments/FIGS/data/aoj_full_v1/fit_v6/results.json :: reference.top.floated_mean :: sha256 2416fc21f93fcfa6 :: GeV
\newcommand{\AojRefWidth}{17.8}  % experiments/FIGS/data/aoj_full_v1/fit_v6/results.json :: reference.top.floated_width :: sha256 2416fc21f93fcfa6 :: GeV
\newcommand{\AojPublicYield}{1{,}813}  % experiments/FIGS/data/aoj_full_v1/fit_v6/results.json :: models.sophon-public.top.signal_yield :: sha256 2416fc21f93fcfa6
\newcommand{\AojPublicYieldErr}{193}  % experiments/FIGS/data/aoj_full_v1/fit_v6/results.json :: models.sophon-public.top.signal_yield_err :: sha256 2416fc21f93fcfa6
\newcommand{\AojPublicMass}{182.8}  % experiments/FIGS/data/aoj_full_v1/fit_v6/results.json :: models.sophon-public.top.floated_mean :: sha256 2416fc21f93fcfa6 :: GeV
\newcommand{\AojPublicWidth}{11.7}  % experiments/FIGS/data/aoj_full_v1/fit_v6/results.json :: models.sophon-public.top.floated_width :: sha256 2416fc21f93fcfa6 :: GeV
\newcommand{\AojInjSizeMin}{1{,}000}  % experiments/FIGS/data/aoj_injection_v1/summary.json :: groups[region=top,mode=ensemble].size (min) :: sha256 d9b66466ae24b584 :: injected tops
\newcommand{\AojInjSizeMax}{4{,}000}  % experiments/FIGS/data/aoj_injection_v1/summary.json :: groups[region=top,mode=ensemble].size (max) :: sha256 d9b66466ae24b584
\newcommand{\AojInjFloatPullSdMax}{2.0}  % experiments/FIGS/data/aoj_injection_v1/summary.json :: groups[region=top,mode=bootstrap,variant=float].pull.sd (max) :: sha256 d9b66466ae24b584 :: pull width with each model's own floated shape
\newcommand{\AojInjEnsembleRecoveryMin}{0.983}  % experiments/FIGS/data/aoj_injection_v1/summary.json :: groups[region=top,mode=ensemble].ratio (min) :: sha256 d9b66466ae24b584 :: the whole procedure, shape re-derived per toy set
\newcommand{\AojInjEnsembleRecoveryMax}{0.998}  % experiments/FIGS/data/aoj_injection_v1/summary.json :: groups[region=top,mode=ensemble].ratio (max) :: sha256 d9b66466ae24b584
\newcommand{\AojInjEnsemblePullSdMin}{1.02}  % experiments/FIGS/data/aoj_injection_v1/summary.json :: groups[region=top,mode=ensemble].pull (min) :: sha256 d9b66466ae24b584 :: the whole procedure, shape re-derived per toy set
\newcommand{\AojInjEnsemblePullSdMax}{1.09}  % experiments/FIGS/data/aoj_injection_v1/summary.json :: groups[region=top,mode=ensemble].pull (max) :: sha256 d9b66466ae24b584
\newcommand{\AojClosureNShards}{10}  % experiments/FIGS/data/aoj_checks_v2/closure_all_shards.json :: n_shards :: sha256 6a6018351d7048d2 :: staged shards of the open data
\newcommand{\AojClosurePtLo}{500}  % experiments/FIGS/data/aoj_checks_v2/closure_all_shards.json :: pt_window[0] :: sha256 6a6018351d7048d2 :: GeV
\newcommand{\AojClosurePtHi}{700}  % experiments/FIGS/data/aoj_checks_v2/closure_all_shards.json :: pt_window[1] :: sha256 6a6018351d7048d2 :: GeV
\newcommand{\AojClosureDzeroIqr}{8.7}  % experiments/FIGS/data/aoj_checks_v2/closure_all_shards.json :: domain_shift.part_d0.iqr_ratio :: sha256 6a6018351d7048d2 :: open data over JetClass-II QCD, pooled over the shards
\newcommand{\AojClosureDzIqr}{10.6}  % experiments/FIGS/data/aoj_checks_v2/closure_all_shards.json :: domain_shift.part_dz.iqr_ratio :: sha256 6a6018351d7048d2 :: open data over JetClass-II QCD, pooled over the shards
\newcommand{\AojClosureNeutral}{0.65}  % experiments/FIGS/data/aoj_checks_v2/closure_all_shards.json :: features.part_isNeutralHadron.ratio.per_shard (mean) :: sha256 6a6018351d7048d2 :: open data over JetClass-II QCD
\newcommand{\AojClosureCharged}{1.22}  % experiments/FIGS/data/aoj_checks_v2/closure_all_shards.json :: features.part_isChargedHadron.ratio.per_shard (mean) :: sha256 6a6018351d7048d2 :: open data over JetClass-II QCD
\newcommand{\AojClosureNPartAoj}{61}  % experiments/FIGS/data/aoj_checks_v2/closure_all_shards.json :: features.n_particles.aoj.per_shard (mean of the shard medians) :: sha256 6a6018351d7048d2 :: median
\newcommand{\AojClosureNPartRef}{48}  % experiments/FIGS/data/aoj_checks_v2/closure_all_shards.json :: features.n_particles.reference.per_shard (mean) :: sha256 6a6018351d7048d2 :: median
\newcommand{\RandMergingDrawsBvcFourprong}{three}  % configs/labelmaps/rand_label_map.v1.csv :: RAND_d* of label_X_YY_bbqq and label_X_YY_ccqq :: sha256 f395a25f52e3cfad :: first-grid partitions that merge the task's pair
\newcommand{\RandSplittingDrawsBvcFourprong}{one and two}  % configs/labelmaps/rand_label_map.v1.csv :: RAND_d* of label_X_YY_bbqq and label_X_YY_ccqq :: sha256 f395a25f52e3cfad :: first-grid partitions that keep the pair apart
\newcommand{\RandMergingDrawsVisibleContent}{one and three}  % configs/labelmaps/rand_label_map.v1.csv :: RAND_d* of label_X_YY_bbqq and label_X_YY_cqtauhv :: sha256 f395a25f52e3cfad :: first-grid partitions that merge the task's pair
\newcommand{\RandSplittingDrawsVisibleContent}{two}  % configs/labelmaps/rand_label_map.v1.csv :: RAND_d* of label_X_YY_bbqq and label_X_YY_cqtauhv :: sha256 f395a25f52e3cfad :: first-grid partitions that keep the pair apart
\newcommand{\RandVtwoNPartitions}{five}  % configs/arms/v2_grid.json :: arms[RAND2_*] (count) :: sha256 d8a31b1486992302 :: random partitions in the rerun
\newcommand{\RandVtwoRuns}{one}  % configs/arms/v2_paper_runs.txt :: runs of RAND2_* :: sha256 b04562b14ea8cbc7 :: runs per partition
\newcommand{\RandFlavRuns}{two}  % configs/arms/v2_paper_runs.txt :: runs of FLAV_* :: sha256 b04562b14ea8cbc7 :: runs per flavour-pair model
\newcommand{\RandVtwoNPairs}{seven}  % configs/labelmaps/rand_v2_selection.json :: pairs (count) :: sha256 98cf52d49c7d2f02 :: probe pairs the balance rule covers
\newcommand{\RandVtwoMergedMin}{two}  % configs/labelmaps/rand_v2_selection.json :: merged_in :: sha256 98cf52d49c7d2f02 :: partitions merging each probe pair, checked on the map
\newcommand{\RandVtwoMergedMax}{three}  % configs/labelmaps/rand_v2_selection.json :: merged_in :: sha256 98cf52d49c7d2f02 :: partitions merging each probe pair, checked on the map
\newcommand{\RandVtwoShareTol}{5\%}  % configs/labelmaps/rand_v2_selection.json :: realised_tolerance :: sha256 98cf52d49c7d2f02 :: largest relative deviation of a realised group share from its target (rule 3)
\newcommand{\RandVtwoTrials}{32}  % scripts/build_rand_control.py :: TRIALS :: sha256 5bc91d77a3f2e81f :: share-matched assignments drawn per random partition; the one sharing the fewest merged pairs with the 17-class label set is kept
\newcommand{\RandVtwoPoolDraws}{6{,}600}  % configs/labelmaps/rand_v2_selection.json :: pool.draws :: sha256 98cf52d49c7d2f02 :: candidate partitions drawn (seeds) before the realised-share rule
\newcommand{\RandVtwoPoolPass}{27}  % configs/labelmaps/rand_v2_selection.json :: rule3_seeds (count) :: sha256 98cf52d49c7d2f02 :: candidate partitions meeting the realised-share rule (rule 3)
\newcommand{\RandFlavSplitOrbit}{\ensuremath{X\to YY\to QQQQ}}  % configs/labelmaps/flavour_pair.v2.json :: split_orbit :: sha256 3b79ce5e44c1e2f4 :: the decays F1 cuts, Q any quark
\newcommand{\RandFlavSplitOrbitClasses}{22}  % configs/labelmaps/flavour_pair_map.v2.csv :: classes whose orbit is split_orbit (count) :: sha256 9aab235099ceac43 :: classes of the split orbit
\newcommand{\RandFlavCutClasses}{11}  % configs/labelmaps/flavour_pair.v2.json :: split_classes_moved (count) :: sha256 3b79ce5e44c1e2f4 :: classes of the split orbit that F1 moves: those with a b quark
\newcommand{\RandFlavOrbitsMoved}{three}  % configs/labelmaps/flavour_pair.v2.json :: orbits_moved_B_to_A (count) :: sha256 3b79ce5e44c1e2f4 :: whole orbits F1 moves the other way
\newcommand{\RandFlavOrbitsMovedList}{\ensuremath{X\to YY\to QQ\mu\mu}, \ensuremath{X\to YY\to QQ\tau_h\tau_h} and \ensuremath{X\to \mu\mu}}  % configs/labelmaps/flavour_pair.v2.json :: orbits_moved_B_to_A :: sha256 3b79ce5e44c1e2f4 :: the orbits F1 moves the other way, Q any quark
\newcommand{\RandFlavOrbitClassesMoved}{9}  % configs/labelmaps/flavour_pair_map.v2.csv :: classes whose orbit is in orbits_moved_B_to_A :: sha256 9aab235099ceac43 :: classes in those orbits
\newcommand{\RandFlavFonerCutWithB}{6}  % configs/labelmaps/flavour_pair_map.v2.csv :: FLAV_F1R != FLAV_F0 within split_orbit, with a b quark (count) :: sha256 9aab235099ceac43 :: classes with a b quark among those the F1r cut moves
\newcommand{\RandFlavBcPairs}{11}  % configs/labelmaps/flavour_pair.v2.json :: b_c_pairs_in_split_orbit (count) :: sha256 3b79ce5e44c1e2f4 :: pairs of the split orbit's classes that differ only by exchanging b and c
\newcommand{\RandFlavFoneBcSplit}{8}  % configs/labelmaps/flavour_pair_map.v2.csv :: b_c_pairs_in_split_orbit split by FLAV_F1 (count) :: sha256 9aab235099ceac43
\newcommand{\RandFlavFonerBcSplit}{none}  % configs/labelmaps/flavour_pair_map.v2.csv :: b_c_pairs_in_split_orbit split by FLAV_F1R (count) :: sha256 9aab235099ceac43 :: PRESPEC A14: none
\newcommand{\RandFlavFonerShareTol}{1\%}  % configs/labelmaps/flavour_pair.v2.json :: F1R.realised_tolerance_vs_F1 :: sha256 3b79ce5e44c1e2f4 :: largest allowed relative difference of a group's realised share from F1's
\newcommand{\RandFlavFonerShareDev}{0.77\%}  % configs/labelmaps/realised_native_shares.v2.json :: share summed per FLAV_F1R and FLAV_F1 group of flavour_pair_map.v2.csv (largest relative difference) :: sha256 5728c12e88fa3683 :: realised shares from the final loader's dry run
\newcommand{\MlstNRunsWord}{three}  % configs/arms/v2_paper_runs.txt :: runs of L188/L162/R42_Q1/R16_Q1 :: sha256 b04562b14ea8cbc7 :: pretraining runs per configuration
\newcommand{\MlstMassNRunsWord}{two}  % configs/arms/v2_paper_runs.txt :: runs of L162_MASS/R16_Q1_MASS :: sha256 b04562b14ea8cbc7 :: pretraining runs per configuration
\newcommand{\LpCMin}{\ensuremath{1\times10^{-3}}}  % experiments/EVAL/linear_probe_v2.py :: C_GRID :: sha256 4c576b44812c18a2 :: smallest inverse regularisation
\newcommand{\LpCMax}{100}  % experiments/EVAL/linear_probe_v2.py :: C_GRID :: sha256 4c576b44812c18a2 :: largest inverse regularisation
\newcommand{\LpNTestJcii}{2{,}000{,}000}  % experiments/EVAL/linear_probe_v2.py :: TEST_FIRST.jc2 :: sha256 4c576b44812c18a2 :: JetClass-II test jets
\newcommand{\LpNTestAucJcii}{500{,}000}  % experiments/EVAL/linear_probe_v2.py :: TEST_FIRST.jc2 / AUC_STRIDE :: sha256 4c576b44812c18a2 :: JetClass-II test jets of the macro AUC
\newcommand{\LpConvergedGrad}{\ensuremath{1\times10^{-5}}}  % experiments/EVAL/linear_probe_v2.py :: CONVERGED_GRAD :: sha256 4c576b44812c18a2 :: largest gradient component a converged probe ends below
\newcommand{\FillNFeatures}{17}  % configs/finetune/TopReference.yaml :: inputs.pf_features.vars (count) :: sha256 501a6ffe5d0a8051 :: per-particle input features of the network
\newcommand{\FillTopNZero}{10}  % configs/finetune/TopReference.yaml :: inputs.pf_features.vars, part_zero_* (count) :: sha256 501a6ffe5d0a8051 :: features the data set lacks, set to zero
\newcommand{\FillTopNPresent}{7}  % configs/finetune/TopReference.yaml :: inputs.pf_features.vars, not part_zero_* (count) :: sha256 501a6ffe5d0a8051 :: features the data set provides
\newcommand{\FillQgNZero}{4}  % configs/finetune/EnergyFlowQG.yaml :: inputs.pf_features.vars, part_zero_* (count) :: sha256 db42d6471e8b865c :: features the data set lacks, set to zero
\newcommand{\FillQgNPresent}{13}  % configs/finetune/EnergyFlowQG.yaml :: inputs.pf_features.vars, not part_zero_* (count) :: sha256 db42d6471e8b865c :: features the data set provides
\newcommand{\DesignVtwoDataFraction}{20\%}  % experiments/MTX/k8s/v2/grid/job-mtx2-l188-s1-raunav.yaml :: --data-fraction :: sha256 d28ce02267a6b87a :: share of every file read per epoch
\newcommand{\DesignVtwoNValFiles}{25}  % experiments/MTX/k8s/v2/grid/job-mtx2-l188-s1-raunav.yaml :: --data-val (files, brace ranges expanded) :: sha256 d28ce02267a6b87a :: validation files of the fixed sample
\newcommand{\DesignWavgFirstEpoch}{70}  % experiments/MTX/pretrain_v2.py :: --num-epochs - WAVG_EPOCHS :: sha256 95da5a6e029855ee :: first epoch of the weight average (epochs count from 0)
\newcommand{\DesignWavgLastEpoch}{79}  % experiments/MTX/k8s/v2/grid/job-mtx2-l188-s1-raunav.yaml :: --num-epochs - 1 :: sha256 d28ce02267a6b87a :: last epoch
\newcommand{\DesignBnJets}{200{,}000}  % experiments/MTX/pretrain_v2.py :: BN_JETS :: sha256 95da5a6e029855ee :: training jets the BatchNorm statistics are recomputed on
\newcommand{\DesignVtwoNWindows}{five}  % experiments/MTX/k8s/v2/grid/job-mtx2-l188-s1-raunav.yaml :: 1 / --data-fraction :: sha256 d28ce02267a6b87a :: disjoint windows of each file's rows, one read per epoch
\newcommand{\DesignLofoNWindows}{three}  % experiments/MTX/k8s/v2/cut/grid/job-mtx2-l188lofo4p-s1-raunav.yaml :: --data-windows :: sha256 44a9573929efb7ac :: windows of the models trained without one family
\newcommand{\DesignVtwoPlanUsed}{97\%}  % experiments/FIGS/data/v2_loader/loader_dryrun/dryrun_s176_seed1.json :: epochs[*].(max_fetch_id + 1) / args.data_split_num, mean :: sha256 03b6586b26ad7a2f :: share of an epoch's planned reading used before its jets are reached, 20 epochs
\newcommand{\DesignLofoPlanUsed}{76\%}  % experiments/FIGS/data/v2_loader/loader_dryrun/dryrun_lofo4p_seed1.json :: epochs[*].(max_fetch_id + 1) / fetches_per_pass, mean :: sha256 2a1315e611d13ff5 :: the same, leaving the family out
\newcommand{\DesignLofoRowFactor}{1.29}  % experiments/FIGS/data/v2_loader/loader_dryrun/dryrun_lofo4p_seed1.json :: (mean plan used / --data-windows) over (mean plan used / k) of the full models' dry run :: sha256 2a1315e611d13ff5 :: how much more often a row of a file is read per epoch when the family is left out
\newcommand{\DesignDistinctJetsCycle}{\ensuremath{3.3\times10^{7}}}  % experiments/FIGS/data/v2_loader/loader_dryrun/dryrun_s176_seed1.json :: sum of epochs[c*k:(c+1)*k].distinct_jets, mean over 4 cycles :: sha256 03b6586b26ad7a2f :: distinct jets the full models train on in one cycle of windows
\newcommand{\DesignLofoFamilyShare}{22.4\%}  % experiments/FIGS/data/v2_loader/loader_dryrun/dryrun_s176_seed1.json :: epochs[*].native_counts over summary.labels_absent_every_epoch of the family-out run :: sha256 03b6586b26ad7a2f :: the left-out family's share of the full models' stream
\newcommand{\DesignLofoNClasses}{27}  % experiments/FIGS/data/v2_loader/loader_dryrun/dryrun_lofo4p_seed1.json :: summary.labels_absent_every_epoch (count) :: sha256 2a1315e611d13ff5 :: native classes of the left-out family
\newcommand{\DesignLofoExposureFactor}{1.29}  % experiments/FIGS/data/v2_loader/loader_dryrun/dryrun_s176_seed1.json :: 1 / (1 - family share) :: sha256 03b6586b26ad7a2f :: how much more often every other class is seen
\newcommand{\DesignLofoQcdShareFull}{14.4\%}  % experiments/FIGS/data/v2_loader/loader_dryrun/dryrun_s176_seed1.json :: summary.mean_qcd_share :: sha256 03b6586b26ad7a2f :: QCD share of the realised stream
\newcommand{\DesignLofoQcdShareLofo}{18.6\%}  % experiments/FIGS/data/v2_loader/loader_dryrun/dryrun_lofo4p_seed1.json :: summary.mean_qcd_share :: sha256 2a1315e611d13ff5 :: QCD share of the realised stream
\newcommand{\DesignLrFlatPercent}{70\%}  % experiments/MTX/pretrain_v2.py :: make_optimizer: 1 - the decay share :: sha256 95da5a6e029855ee :: share of the epochs at the initial rate
\newcommand{\DesignLrDecayFirstEpoch}{56}  % experiments/MTX/pretrain_v2.py :: num_epochs - int(num_epochs * decay share) :: sha256 95da5a6e029855ee :: first epoch at a reduced rate (epochs from 0)
\newcommand{\DesignLrEndPercent}{1\%}  % experiments/MTX/pretrain_v2.py :: make_optimizer: gamma ** n_decay :: sha256 95da5a6e029855ee :: the last epoch's rate over the initial rate
\newcommand{\DesignEndLr}{\ensuremath{5\times10^{-6}}}  % experiments/MTX/k8s/v2/grid/job-mtx2-l188-s1-raunav.yaml :: --start-lr x the final factor :: sha256 d28ce02267a6b87a :: the rate of the last epoch
\newcommand{\DesignWeaverVersion}{0.4.17}  % experiments/MTX/pretrain_v2.py :: the weaver release whose loop it reproduces :: sha256 95da5a6e029855ee
\newcommand{\DesignMeiFactor}{1.1}  % src/stats/mde.py :: MEI_LOG = ln(factor) :: sha256 626468dcc37940bb :: the smallest effect of interest, as a factor
\newcommand{\DesignVtwoGpu}{RTX 3090}  % experiments/MTX/k8s/v2/grid/job-mtx2-flavf0-s1-raunav.yaml :: nodeAffinity nvidia.com/gpu.product :: sha256 960cad1691314980 :: the product of run index 1, every configuration
\newcommand{\DesignVtwoSecondGpu}{\pending{pending: the GPU type of each run index}}  % experiments/FIGS/make_tables.py :: emit_v2_recipe[DesignVtwoSecondGpu] :: sha256 10e6ec08356a4a2b :: a marker, not a number: the grid's job specs put every run index on one product
\newcommand{\DesignVtwoSecondGpuRuns}{\pending{pending: the GPU type of each run index}}  % experiments/FIGS/make_tables.py :: emit_v2_recipe[DesignVtwoSecondGpuRuns] :: sha256 10e6ec08356a4a2b :: a marker, not a number: the grid's job specs put every run index on one product
\newcommand{\MassNRunsWord}{five}  % experiments/FIGS/data/probe_ladder_mass2x2_mlp2/s1.json :: per-seed files (count) :: sha256 906e778d7043fe67 :: pretraining runs of each mass-output configuration
\newcommand{\MassOmaLinearOnesixtwo}{\ensuremath{(2.36 \pm 0.06)\times10^{-3}}}  % experiments/FIGS/data/probe_ladder_v2_mlp2/analysis/seed_level_results.json :: table[bvc_resonant,linear,162].auc :: sha256 53e1ec81343ac2c1 :: 1-AUC, mean +- SD over 5 seeds
\newcommand{\MassOmaLinearOnesixtwoMass}{\ensuremath{(2.62 \pm 0.08)\times10^{-3}}}  % experiments/FIGS/data/probe_ladder_mass2x2_mlp2/s1.json :: tasks.bvc_resonant.arms.l162mass-s*.linear.auc :: sha256 906e778d7043fe67 :: 1-AUC, mean +- SD over the 5 per-seed files
\newcommand{\MassOmaRatioLinearOnesixtwo}{1.11}  % experiments/FIGS/data/probe_ladder_mass2x2_mlp2/s1.json :: tasks.bvc_resonant.arms.l162mass-s*.linear.auc over table[bvc_resonant,linear,162].auc :: sha256 906e778d7043fe67 :: seed mean of 1-AUC with the mass output over without
\newcommand{\MassOmaLinearOneseven}{\ensuremath{(9.3 \pm 1.5)\times10^{-3}}}  % experiments/FIGS/data/probe_ladder_v2_mlp2/analysis/seed_level_results.json :: table[bvc_resonant,linear,17].auc :: sha256 53e1ec81343ac2c1 :: 1-AUC, mean +- SD over 5 seeds
\newcommand{\MassOmaLinearOnesevenMass}{\ensuremath{(3.0 \pm 1.2)\times10^{-2}}}  % experiments/FIGS/data/probe_ladder_mass2x2_mlp2/s1.json :: tasks.bvc_resonant.arms.r16q1mass-s*.linear.auc :: sha256 906e778d7043fe67 :: 1-AUC, mean +- SD over the 5 per-seed files
\newcommand{\MassOmaRatioLinearOneseven}{3.2}  % experiments/FIGS/data/probe_ladder_mass2x2_mlp2/s1.json :: tasks.bvc_resonant.arms.r16q1mass-s*.linear.auc over table[bvc_resonant,linear,17].auc :: sha256 906e778d7043fe67 :: seed mean of 1-AUC with the mass output over without
\newcommand{\MassOmaMlpOnesixtwo}{\ensuremath{(2.32 \pm 0.05)\times10^{-3}}}  % experiments/FIGS/data/probe_ladder_v2_mlp2/analysis/seed_level_results.json :: table[bvc_resonant,mlp,162].auc :: sha256 53e1ec81343ac2c1 :: 1-AUC, mean +- SD over 5 seeds
\newcommand{\MassOmaMlpOnesixtwoMass}{\ensuremath{(2.52 \pm 0.06)\times10^{-3}}}  % experiments/FIGS/data/probe_ladder_mass2x2_mlp2/s1.json :: tasks.bvc_resonant.arms.l162mass-s*.mlp.auc :: sha256 906e778d7043fe67 :: 1-AUC, mean +- SD over the 5 per-seed files
\newcommand{\MassOmaRatioMlpOnesixtwo}{1.09}  % experiments/FIGS/data/probe_ladder_mass2x2_mlp2/s1.json :: tasks.bvc_resonant.arms.l162mass-s*.mlp.auc over table[bvc_resonant,mlp,162].auc :: sha256 906e778d7043fe67 :: seed mean of 1-AUC with the mass output over without
\newcommand{\MassOmaMlpOneseven}{\ensuremath{(7.3 \pm 0.9)\times10^{-3}}}  % experiments/FIGS/data/probe_ladder_v2_mlp2/analysis/seed_level_results.json :: table[bvc_resonant,mlp,17].auc :: sha256 53e1ec81343ac2c1 :: 1-AUC, mean +- SD over 5 seeds
\newcommand{\MassOmaMlpOnesevenMass}{\ensuremath{(1.5 \pm 0.4)\times10^{-2}}}  % experiments/FIGS/data/probe_ladder_mass2x2_mlp2/s1.json :: tasks.bvc_resonant.arms.r16q1mass-s*.mlp.auc :: sha256 906e778d7043fe67 :: 1-AUC, mean +- SD over the 5 per-seed files
\newcommand{\MassOmaRatioMlpOneseven}{2.0}  % experiments/FIGS/data/probe_ladder_mass2x2_mlp2/s1.json :: tasks.bvc_resonant.arms.r16q1mass-s*.mlp.auc over table[bvc_resonant,mlp,17].auc :: sha256 906e778d7043fe67 :: seed mean of 1-AUC with the mass output over without
\newcommand{\FtAucJciiEThreeOneeighteight}{0.902\,\ensuremath{\pm}\,0.004}  % experiments/FIGS/data/w2b_leg1_metrics_v2.json :: cells.{l188-s1,l188-s2,l188-s3,l188-s4,l188-s5}.N1000.s1.macro_auc_ovr :: sha256 b56bf70be9982044 :: mean +- SD over the row's 5 pretrained models
\newcommand{\FtAucJciiEFourOneeighteight}{0.9723\,\ensuremath{\pm}\,0.0004}  % experiments/FIGS/data/w2b_leg1_metrics_v2.json :: cells.{l188-s1,l188-s2,l188-s3,l188-s4,l188-s5}.N10000.s1.macro_auc_ovr :: sha256 b56bf70be9982044 :: mean +- SD over the row's 5 pretrained models
\newcommand{\FtAucJciiEFiveOneeighteight}{0.9793\,\ensuremath{\pm}\,0.0002}  % experiments/FIGS/data/w2b_leg1_metrics_v2.json :: cells.{l188-s1,l188-s2,l188-s3,l188-s4,l188-s5}.N100000.s1.macro_auc_ovr :: sha256 b56bf70be9982044 :: mean +- SD over the row's 5 pretrained models
\newcommand{\FtAucJciiESixOneeighteight}{0.9809\,\ensuremath{\pm}\,\ensuremath{<}0.0001}  % experiments/FIGS/data/w2b_leg1_metrics_v2.json :: cells.{l188-s1,l188-s2,l188-s3,l188-s4,l188-s5}.N1000000.s1.macro_auc_ovr :: sha256 b56bf70be9982044 :: mean +- SD over the row's 5 pretrained models
\newcommand{\FtAucJciiEThreeOnesixtwo}{0.906\,\ensuremath{\pm}\,0.007}  % experiments/FIGS/data/w2b_leg1_metrics_v2.json :: cells.{l162-s1b,l162-s2,l162-s3,l162-s4,l162-s5}.N1000.s1.macro_auc_ovr :: sha256 b56bf70be9982044 :: mean +- SD over the row's 5 pretrained models
\newcommand{\FtAucJciiEFourOnesixtwo}{0.9724\,\ensuremath{\pm}\,0.0005}  % experiments/FIGS/data/w2b_leg1_metrics_v2.json :: cells.{l162-s1b,l162-s2,l162-s3,l162-s4,l162-s5}.N10000.s1.macro_auc_ovr :: sha256 b56bf70be9982044 :: mean +- SD over the row's 5 pretrained models
\newcommand{\FtAucJciiEFiveOnesixtwo}{0.9794\,\ensuremath{\pm}\,0.0002}  % experiments/FIGS/data/w2b_leg1_metrics_v2.json :: cells.{l162-s1b,l162-s2,l162-s3,l162-s4,l162-s5}.N100000.s1.macro_auc_ovr :: sha256 b56bf70be9982044 :: mean +- SD over the row's 5 pretrained models
\newcommand{\FtAucJciiESixOnesixtwo}{0.9808\,\ensuremath{\pm}\,0.0001}  % experiments/FIGS/data/w2b_leg1_metrics_v2.json :: cells.{l162-s1b,l162-s2,l162-s3,l162-s4,l162-s5}.N1000000.s1.macro_auc_ovr :: sha256 b56bf70be9982044 :: mean +- SD over the row's 5 pretrained models
\newcommand{\FtAucJciiEThreeFourthree}{0.918\,\ensuremath{\pm}\,0.005}  % experiments/FIGS/data/w2b_leg1_metrics_v2.json :: cells.{r42q1-s1,r42q1-s2,r42q1-s3,r42q1-s4,r42q1-s5}.N1000.s1.macro_auc_ovr :: sha256 b56bf70be9982044 :: mean +- SD over the row's 5 pretrained models
\newcommand{\FtAucJciiEFourFourthree}{0.9667\,\ensuremath{\pm}\,0.0005}  % experiments/FIGS/data/w2b_leg1_metrics_v2.json :: cells.{r42q1-s1,r42q1-s2,r42q1-s3,r42q1-s4,r42q1-s5}.N10000.s1.macro_auc_ovr :: sha256 b56bf70be9982044 :: mean +- SD over the row's 5 pretrained models
\newcommand{\FtAucJciiEFiveFourthree}{0.9775\,\ensuremath{\pm}\,0.0003}  % experiments/FIGS/data/w2b_leg1_metrics_v2.json :: cells.{r42q1-s1,r42q1-s2,r42q1-s3,r42q1-s4,r42q1-s5}.N100000.s1.macro_auc_ovr :: sha256 b56bf70be9982044 :: mean +- SD over the row's 5 pretrained models
\newcommand{\FtAucJciiESixFourthree}{0.9800\,\ensuremath{\pm}\,0.0001}  % experiments/FIGS/data/w2b_leg1_metrics_v2.json :: cells.{r42q1-s1,r42q1-s2,r42q1-s3,r42q1-s4,r42q1-s5}.N1000000.s1.macro_auc_ovr :: sha256 b56bf70be9982044 :: mean +- SD over the row's 5 pretrained models
\newcommand{\FtAucJciiEThreeOneseven}{0.905\,\ensuremath{\pm}\,0.008}  % experiments/FIGS/data/w2b_leg1_metrics_v2.json :: cells.{r16q1-s1,r16q1-s2,r16q1-s3,r16q1-s4,r16q1-s5}.N1000.s1.macro_auc_ovr :: sha256 b56bf70be9982044 :: mean +- SD over the row's 5 pretrained models
\newcommand{\FtAucJciiEFourOneseven}{0.9635\,\ensuremath{\pm}\,0.0005}  % experiments/FIGS/data/w2b_leg1_metrics_v2.json :: cells.{r16q1-s1,r16q1-s2,r16q1-s3,r16q1-s4,r16q1-s5}.N10000.s1.macro_auc_ovr :: sha256 b56bf70be9982044 :: mean +- SD over the row's 5 pretrained models
\newcommand{\FtAucJciiEFiveOneseven}{0.9748\,\ensuremath{\pm}\,0.0003}  % experiments/FIGS/data/w2b_leg1_metrics_v2.json :: cells.{r16q1-s1,r16q1-s2,r16q1-s3,r16q1-s4,r16q1-s5}.N100000.s1.macro_auc_ovr :: sha256 b56bf70be9982044 :: mean +- SD over the row's 5 pretrained models
\newcommand{\FtAucJciiESixOneseven}{0.9781\,\ensuremath{\pm}\,0.0001}  % experiments/FIGS/data/w2b_leg1_metrics_v2.json :: cells.{r16q1-s1,r16q1-s2,r16q1-s3,r16q1-s4,r16q1-s5}.N1000000.s1.macro_auc_ovr :: sha256 b56bf70be9982044 :: mean +- SD over the row's 5 pretrained models
\newcommand{\FtAucJciiEThreeOnesixtwoMass}{0.921\,\ensuremath{\pm}\,0.005}  % experiments/FIGS/data/w2b_leg1_metrics_v2.json :: cells.{l162mass-s1,l162mass-s2,l162mass-s3,l162mass-s4,l162mass-s5}.N1000.s1.macro_auc_ovr :: sha256 b56bf70be9982044 :: mean +- SD over the row's 5 pretrained models
\newcommand{\FtAucJciiEFourOnesixtwoMass}{0.9724\,\ensuremath{\pm}\,0.0013}  % experiments/FIGS/data/w2b_leg1_metrics_v2.json :: cells.{l162mass-s1,l162mass-s2,l162mass-s3,l162mass-s4,l162mass-s5}.N10000.s1.macro_auc_ovr :: sha256 b56bf70be9982044 :: mean +- SD over the row's 5 pretrained models
\newcommand{\FtAucJciiEFiveOnesixtwoMass}{0.9790\,\ensuremath{\pm}\,0.0001}  % experiments/FIGS/data/w2b_leg1_metrics_v2.json :: cells.{l162mass-s1,l162mass-s2,l162mass-s3,l162mass-s4,l162mass-s5}.N100000.s1.macro_auc_ovr :: sha256 b56bf70be9982044 :: mean +- SD over the row's 5 pretrained models
\newcommand{\FtAucJciiESixOnesixtwoMass}{0.9805\,\ensuremath{\pm}\,0.0001}  % experiments/FIGS/data/w2b_leg1_metrics_v2.json :: cells.{l162mass-s1,l162mass-s2,l162mass-s3,l162mass-s4,l162mass-s5}.N1000000.s1.macro_auc_ovr :: sha256 b56bf70be9982044 :: mean +- SD over the row's 5 pretrained models
\newcommand{\FtAucJciiEThreeOnesevenMass}{0.904\,\ensuremath{\pm}\,0.007}  % experiments/FIGS/data/w2b_leg1_metrics_v2.json :: cells.{r16q1mass-s1,r16q1mass-s2,r16q1mass-s3,r16q1mass-s4,r16q1mass-s5}.N1000.s1.macro_auc_ovr :: sha256 b56bf70be9982044 :: mean +- SD over the row's 5 pretrained models
\newcommand{\FtAucJciiEFourOnesevenMass}{0.9639\,\ensuremath{\pm}\,0.0007}  % experiments/FIGS/data/w2b_leg1_metrics_v2.json :: cells.{r16q1mass-s1,r16q1mass-s2,r16q1mass-s3,r16q1mass-s4,r16q1mass-s5}.N10000.s1.macro_auc_ovr :: sha256 b56bf70be9982044 :: mean +- SD over the row's 5 pretrained models
\newcommand{\FtAucJciiEFiveOnesevenMass}{0.9740\,\ensuremath{\pm}\,0.0004}  % experiments/FIGS/data/w2b_leg1_metrics_v2.json :: cells.{r16q1mass-s1,r16q1mass-s2,r16q1mass-s3,r16q1mass-s4,r16q1mass-s5}.N100000.s1.macro_auc_ovr :: sha256 b56bf70be9982044 :: mean +- SD over the row's 5 pretrained models
\newcommand{\FtAucJciiESixOnesevenMass}{0.9774\,\ensuremath{\pm}\,0.0001}  % experiments/FIGS/data/w2b_leg1_metrics_v2.json :: cells.{r16q1mass-s1,r16q1mass-s2,r16q1mass-s3,r16q1mass-s4,r16q1mass-s5}.N1000000.s1.macro_auc_ovr :: sha256 b56bf70be9982044 :: mean +- SD over the row's 5 pretrained models
\newcommand{\FtAucJciiEThreeRandDrawOne}{0.921}  % experiments/FIGS/data/w2b_leg1_metrics_v2.json :: cells.{rand-d1-s1b}.N1000.s1.macro_auc_ovr :: sha256 b56bf70be9982044 :: one run
\newcommand{\FtAucJciiEFourRandDrawOne}{0.9675}  % experiments/FIGS/data/w2b_leg1_metrics_v2.json :: cells.{rand-d1-s1b}.N10000.s1.macro_auc_ovr :: sha256 b56bf70be9982044 :: one run
\newcommand{\FtAucJciiEFiveRandDrawOne}{0.97617}  % experiments/FIGS/data/w2b_leg1_metrics_v2.json :: cells.{rand-d1-s1b}.N100000.s1.macro_auc_ovr :: sha256 b56bf70be9982044 :: one run
\newcommand{\FtAucJciiESixRandDrawOne}{0.97941}  % experiments/FIGS/data/w2b_leg1_metrics_v2.json :: cells.{rand-d1-s1b}.N1000000.s1.macro_auc_ovr :: sha256 b56bf70be9982044 :: one run
\newcommand{\FtAucJciiEThreeRandDrawTwo}{0.914}  % experiments/FIGS/data/w2b_leg1_metrics_v2.json :: cells.{rand-d2-s2}.N1000.s1.macro_auc_ovr :: sha256 b56bf70be9982044 :: one run
\newcommand{\FtAucJciiEFourRandDrawTwo}{0.9675}  % experiments/FIGS/data/w2b_leg1_metrics_v2.json :: cells.{rand-d2-s2}.N10000.s1.macro_auc_ovr :: sha256 b56bf70be9982044 :: one run
\newcommand{\FtAucJciiEFiveRandDrawTwo}{0.97668}  % experiments/FIGS/data/w2b_leg1_metrics_v2.json :: cells.{rand-d2-s2}.N100000.s1.macro_auc_ovr :: sha256 b56bf70be9982044 :: one run
\newcommand{\FtAucJciiESixRandDrawTwo}{0.97943}  % experiments/FIGS/data/w2b_leg1_metrics_v2.json :: cells.{rand-d2-s2}.N1000000.s1.macro_auc_ovr :: sha256 b56bf70be9982044 :: one run
\newcommand{\FtAucJciiEThreeRandDrawThree}{0.918}  % experiments/FIGS/data/w2b_leg1_metrics_v2.json :: cells.{rand-d3-s3}.N1000.s1.macro_auc_ovr :: sha256 b56bf70be9982044 :: one run
\newcommand{\FtAucJciiEFourRandDrawThree}{0.9686}  % experiments/FIGS/data/w2b_leg1_metrics_v2.json :: cells.{rand-d3-s3}.N10000.s1.macro_auc_ovr :: sha256 b56bf70be9982044 :: one run
\newcommand{\FtAucJciiEFiveRandDrawThree}{0.97605}  % experiments/FIGS/data/w2b_leg1_metrics_v2.json :: cells.{rand-d3-s3}.N100000.s1.macro_auc_ovr :: sha256 b56bf70be9982044 :: one run
\newcommand{\FtAucJciiESixRandDrawThree}{0.97950}  % experiments/FIGS/data/w2b_leg1_metrics_v2.json :: cells.{rand-d3-s3}.N1000000.s1.macro_auc_ovr :: sha256 b56bf70be9982044 :: one run
\newcommand{\FtAccJciiEThreeOneeighteight}{0.445\,\ensuremath{\pm}\,0.004}  % experiments/FIGS/data/w2b_leg1_metrics_v2.json :: cells.{l188-s1,l188-s2,l188-s3,l188-s4,l188-s5}.N1000.s1.accuracy :: sha256 b56bf70be9982044 :: mean +- SD over the row's 5 pretrained models
\newcommand{\FtAccJciiEFourOneeighteight}{0.5199\,\ensuremath{\pm}\,0.0018}  % experiments/FIGS/data/w2b_leg1_metrics_v2.json :: cells.{l188-s1,l188-s2,l188-s3,l188-s4,l188-s5}.N10000.s1.accuracy :: sha256 b56bf70be9982044 :: mean +- SD over the row's 5 pretrained models
\newcommand{\FtAccJciiEFiveOneeighteight}{0.5379\,\ensuremath{\pm}\,0.0010}  % experiments/FIGS/data/w2b_leg1_metrics_v2.json :: cells.{l188-s1,l188-s2,l188-s3,l188-s4,l188-s5}.N100000.s1.accuracy :: sha256 b56bf70be9982044 :: mean +- SD over the row's 5 pretrained models
\newcommand{\FtAccJciiESixOneeighteight}{0.5443\,\ensuremath{\pm}\,0.0003}  % experiments/FIGS/data/w2b_leg1_metrics_v2.json :: cells.{l188-s1,l188-s2,l188-s3,l188-s4,l188-s5}.N1000000.s1.accuracy :: sha256 b56bf70be9982044 :: mean +- SD over the row's 5 pretrained models
\newcommand{\FtAccJciiEThreeOnesixtwo}{0.448\,\ensuremath{\pm}\,0.006}  % experiments/FIGS/data/w2b_leg1_metrics_v2.json :: cells.{l162-s1b,l162-s2,l162-s3,l162-s4,l162-s5}.N1000.s1.accuracy :: sha256 b56bf70be9982044 :: mean +- SD over the row's 5 pretrained models
\newcommand{\FtAccJciiEFourOnesixtwo}{0.5212\,\ensuremath{\pm}\,0.0013}  % experiments/FIGS/data/w2b_leg1_metrics_v2.json :: cells.{l162-s1b,l162-s2,l162-s3,l162-s4,l162-s5}.N10000.s1.accuracy :: sha256 b56bf70be9982044 :: mean +- SD over the row's 5 pretrained models
\newcommand{\FtAccJciiEFiveOnesixtwo}{0.5387\,\ensuremath{\pm}\,0.0008}  % experiments/FIGS/data/w2b_leg1_metrics_v2.json :: cells.{l162-s1b,l162-s2,l162-s3,l162-s4,l162-s5}.N100000.s1.accuracy :: sha256 b56bf70be9982044 :: mean +- SD over the row's 5 pretrained models
\newcommand{\FtAccJciiESixOnesixtwo}{0.5440\,\ensuremath{\pm}\,0.0006}  % experiments/FIGS/data/w2b_leg1_metrics_v2.json :: cells.{l162-s1b,l162-s2,l162-s3,l162-s4,l162-s5}.N1000000.s1.accuracy :: sha256 b56bf70be9982044 :: mean +- SD over the row's 5 pretrained models
\newcommand{\FtAccJciiEThreeFourthree}{0.380\,\ensuremath{\pm}\,0.005}  % experiments/FIGS/data/w2b_leg1_metrics_v2.json :: cells.{r42q1-s1,r42q1-s2,r42q1-s3,r42q1-s4,r42q1-s5}.N1000.s1.accuracy :: sha256 b56bf70be9982044 :: mean +- SD over the row's 5 pretrained models
\newcommand{\FtAccJciiEFourFourthree}{0.481\,\ensuremath{\pm}\,0.004}  % experiments/FIGS/data/w2b_leg1_metrics_v2.json :: cells.{r42q1-s1,r42q1-s2,r42q1-s3,r42q1-s4,r42q1-s5}.N10000.s1.accuracy :: sha256 b56bf70be9982044 :: mean +- SD over the row's 5 pretrained models
\newcommand{\FtAccJciiEFiveFourthree}{0.5219\,\ensuremath{\pm}\,0.0017}  % experiments/FIGS/data/w2b_leg1_metrics_v2.json :: cells.{r42q1-s1,r42q1-s2,r42q1-s3,r42q1-s4,r42q1-s5}.N100000.s1.accuracy :: sha256 b56bf70be9982044 :: mean +- SD over the row's 5 pretrained models
\newcommand{\FtAccJciiESixFourthree}{0.5362\,\ensuremath{\pm}\,0.0007}  % experiments/FIGS/data/w2b_leg1_metrics_v2.json :: cells.{r42q1-s1,r42q1-s2,r42q1-s3,r42q1-s4,r42q1-s5}.N1000000.s1.accuracy :: sha256 b56bf70be9982044 :: mean +- SD over the row's 5 pretrained models
\newcommand{\FtAccJciiEThreeOneseven}{0.332\,\ensuremath{\pm}\,0.004}  % experiments/FIGS/data/w2b_leg1_metrics_v2.json :: cells.{r16q1-s1,r16q1-s2,r16q1-s3,r16q1-s4,r16q1-s5}.N1000.s1.accuracy :: sha256 b56bf70be9982044 :: mean +- SD over the row's 5 pretrained models
\newcommand{\FtAccJciiEFourOneseven}{0.447\,\ensuremath{\pm}\,0.013}  % experiments/FIGS/data/w2b_leg1_metrics_v2.json :: cells.{r16q1-s1,r16q1-s2,r16q1-s3,r16q1-s4,r16q1-s5}.N10000.s1.accuracy :: sha256 b56bf70be9982044 :: mean +- SD over the row's 5 pretrained models
\newcommand{\FtAccJciiEFiveOneseven}{0.5078\,\ensuremath{\pm}\,0.0011}  % experiments/FIGS/data/w2b_leg1_metrics_v2.json :: cells.{r16q1-s1,r16q1-s2,r16q1-s3,r16q1-s4,r16q1-s5}.N100000.s1.accuracy :: sha256 b56bf70be9982044 :: mean +- SD over the row's 5 pretrained models
\newcommand{\FtAccJciiESixOneseven}{0.5227\,\ensuremath{\pm}\,0.0010}  % experiments/FIGS/data/w2b_leg1_metrics_v2.json :: cells.{r16q1-s1,r16q1-s2,r16q1-s3,r16q1-s4,r16q1-s5}.N1000000.s1.accuracy :: sha256 b56bf70be9982044 :: mean +- SD over the row's 5 pretrained models
\newcommand{\FtAccJciiEThreeOnesixtwoMass}{0.451\,\ensuremath{\pm}\,0.007}  % experiments/FIGS/data/w2b_leg1_metrics_v2.json :: cells.{l162mass-s1,l162mass-s2,l162mass-s3,l162mass-s4,l162mass-s5}.N1000.s1.accuracy :: sha256 b56bf70be9982044 :: mean +- SD over the row's 5 pretrained models
\newcommand{\FtAccJciiEFourOnesixtwoMass}{0.515\,\ensuremath{\pm}\,0.004}  % experiments/FIGS/data/w2b_leg1_metrics_v2.json :: cells.{l162mass-s1,l162mass-s2,l162mass-s3,l162mass-s4,l162mass-s5}.N10000.s1.accuracy :: sha256 b56bf70be9982044 :: mean +- SD over the row's 5 pretrained models
\newcommand{\FtAccJciiEFiveOnesixtwoMass}{0.5357\,\ensuremath{\pm}\,0.0007}  % experiments/FIGS/data/w2b_leg1_metrics_v2.json :: cells.{l162mass-s1,l162mass-s2,l162mass-s3,l162mass-s4,l162mass-s5}.N100000.s1.accuracy :: sha256 b56bf70be9982044 :: mean +- SD over the row's 5 pretrained models
\newcommand{\FtAccJciiESixOnesixtwoMass}{0.5409\,\ensuremath{\pm}\,0.0010}  % experiments/FIGS/data/w2b_leg1_metrics_v2.json :: cells.{l162mass-s1,l162mass-s2,l162mass-s3,l162mass-s4,l162mass-s5}.N1000000.s1.accuracy :: sha256 b56bf70be9982044 :: mean +- SD over the row's 5 pretrained models
\newcommand{\FtAccJciiEThreeOnesevenMass}{0.331\,\ensuremath{\pm}\,0.005}  % experiments/FIGS/data/w2b_leg1_metrics_v2.json :: cells.{r16q1mass-s1,r16q1mass-s2,r16q1mass-s3,r16q1mass-s4,r16q1mass-s5}.N1000.s1.accuracy :: sha256 b56bf70be9982044 :: mean +- SD over the row's 5 pretrained models
\newcommand{\FtAccJciiEFourOnesevenMass}{0.4507\,\ensuremath{\pm}\,0.0034}  % experiments/FIGS/data/w2b_leg1_metrics_v2.json :: cells.{r16q1mass-s1,r16q1mass-s2,r16q1mass-s3,r16q1mass-s4,r16q1mass-s5}.N10000.s1.accuracy :: sha256 b56bf70be9982044 :: mean +- SD over the row's 5 pretrained models
\newcommand{\FtAccJciiEFiveOnesevenMass}{0.4993\,\ensuremath{\pm}\,0.0022}  % experiments/FIGS/data/w2b_leg1_metrics_v2.json :: cells.{r16q1mass-s1,r16q1mass-s2,r16q1mass-s3,r16q1mass-s4,r16q1mass-s5}.N100000.s1.accuracy :: sha256 b56bf70be9982044 :: mean +- SD over the row's 5 pretrained models
\newcommand{\FtAccJciiESixOnesevenMass}{0.5166\,\ensuremath{\pm}\,0.0013}  % experiments/FIGS/data/w2b_leg1_metrics_v2.json :: cells.{r16q1mass-s1,r16q1mass-s2,r16q1mass-s3,r16q1mass-s4,r16q1mass-s5}.N1000000.s1.accuracy :: sha256 b56bf70be9982044 :: mean +- SD over the row's 5 pretrained models
\newcommand{\FtAccJciiEThreeRandDrawOne}{0.424}  % experiments/FIGS/data/w2b_leg1_metrics_v2.json :: cells.{rand-d1-s1b}.N1000.s1.accuracy :: sha256 b56bf70be9982044 :: one run
\newcommand{\FtAccJciiEFourRandDrawOne}{0.5070}  % experiments/FIGS/data/w2b_leg1_metrics_v2.json :: cells.{rand-d1-s1b}.N10000.s1.accuracy :: sha256 b56bf70be9982044 :: one run
\newcommand{\FtAccJciiEFiveRandDrawOne}{0.5226}  % experiments/FIGS/data/w2b_leg1_metrics_v2.json :: cells.{rand-d1-s1b}.N100000.s1.accuracy :: sha256 b56bf70be9982044 :: one run
\newcommand{\FtAccJciiESixRandDrawOne}{0.53354}  % experiments/FIGS/data/w2b_leg1_metrics_v2.json :: cells.{rand-d1-s1b}.N1000000.s1.accuracy :: sha256 b56bf70be9982044 :: one run
\newcommand{\FtAccJciiEThreeRandDrawTwo}{0.420}  % experiments/FIGS/data/w2b_leg1_metrics_v2.json :: cells.{rand-d2-s2}.N1000.s1.accuracy :: sha256 b56bf70be9982044 :: one run
\newcommand{\FtAccJciiEFourRandDrawTwo}{0.5074}  % experiments/FIGS/data/w2b_leg1_metrics_v2.json :: cells.{rand-d2-s2}.N10000.s1.accuracy :: sha256 b56bf70be9982044 :: one run
\newcommand{\FtAccJciiEFiveRandDrawTwo}{0.5243}  % experiments/FIGS/data/w2b_leg1_metrics_v2.json :: cells.{rand-d2-s2}.N100000.s1.accuracy :: sha256 b56bf70be9982044 :: one run
\newcommand{\FtAccJciiESixRandDrawTwo}{0.53441}  % experiments/FIGS/data/w2b_leg1_metrics_v2.json :: cells.{rand-d2-s2}.N1000000.s1.accuracy :: sha256 b56bf70be9982044 :: one run
\newcommand{\FtAccJciiEThreeRandDrawThree}{0.423}  % experiments/FIGS/data/w2b_leg1_metrics_v2.json :: cells.{rand-d3-s3}.N1000.s1.accuracy :: sha256 b56bf70be9982044 :: one run
\newcommand{\FtAccJciiEFourRandDrawThree}{0.5114}  % experiments/FIGS/data/w2b_leg1_metrics_v2.json :: cells.{rand-d3-s3}.N10000.s1.accuracy :: sha256 b56bf70be9982044 :: one run
\newcommand{\FtAccJciiEFiveRandDrawThree}{0.5245}  % experiments/FIGS/data/w2b_leg1_metrics_v2.json :: cells.{rand-d3-s3}.N100000.s1.accuracy :: sha256 b56bf70be9982044 :: one run
\newcommand{\FtAccJciiESixRandDrawThree}{0.53473}  % experiments/FIGS/data/w2b_leg1_metrics_v2.json :: cells.{rand-d3-s3}.N1000000.s1.accuracy :: sha256 b56bf70be9982044 :: one run
\newcommand{\FtOmaRatioJciiEThreeOnesixtwoOverOneeighteight}{0.96}  % experiments/FIGS/data/w2b_leg1_metrics_v2.json :: cells.*.N1000.s1.macro_auc_ovr :: sha256 b56bf70be9982044 :: ratio of the seed means of 1 - macro AUC, this row over the reference
\newcommand{\FtOmaRatioJciiEThreeFourthreeOverOneeighteight}{0.84}  % experiments/FIGS/data/w2b_leg1_metrics_v2.json :: cells.*.N1000.s1.macro_auc_ovr :: sha256 b56bf70be9982044 :: ratio of the seed means of 1 - macro AUC, this row over the reference
\newcommand{\FtOmaRatioJciiEThreeOnesevenOverOneeighteight}{0.97}  % experiments/FIGS/data/w2b_leg1_metrics_v2.json :: cells.*.N1000.s1.macro_auc_ovr :: sha256 b56bf70be9982044 :: ratio of the seed means of 1 - macro AUC, this row over the reference
\newcommand{\FtOmaRatioJciiEThreeOnesixtwoMassOverOneeighteight}{0.81}  % experiments/FIGS/data/w2b_leg1_metrics_v2.json :: cells.*.N1000.s1.macro_auc_ovr :: sha256 b56bf70be9982044 :: ratio of the seed means of 1 - macro AUC, this row over the reference
\newcommand{\FtOmaRatioJciiEThreeOnesevenMassOverOneeighteight}{0.98}  % experiments/FIGS/data/w2b_leg1_metrics_v2.json :: cells.*.N1000.s1.macro_auc_ovr :: sha256 b56bf70be9982044 :: ratio of the seed means of 1 - macro AUC, this row over the reference
\newcommand{\FtOmaRatioJciiEThreeRandDrawOneOverOneeighteight}{0.80}  % experiments/FIGS/data/w2b_leg1_metrics_v2.json :: cells.*.N1000.s1.macro_auc_ovr :: sha256 b56bf70be9982044 :: ratio of the seed means of 1 - macro AUC, this row over the reference
\newcommand{\FtOmaRatioJciiEThreeRandDrawTwoOverOneeighteight}{0.87}  % experiments/FIGS/data/w2b_leg1_metrics_v2.json :: cells.*.N1000.s1.macro_auc_ovr :: sha256 b56bf70be9982044 :: ratio of the seed means of 1 - macro AUC, this row over the reference
\newcommand{\FtOmaRatioJciiEThreeRandDrawThreeOverOneeighteight}{0.83}  % experiments/FIGS/data/w2b_leg1_metrics_v2.json :: cells.*.N1000.s1.macro_auc_ovr :: sha256 b56bf70be9982044 :: ratio of the seed means of 1 - macro AUC, this row over the reference
\newcommand{\FtOmaRatioJciiEThreeOneeighteightOverOnesixtwo}{1.05}  % experiments/FIGS/data/w2b_leg1_metrics_v2.json :: cells.*.N1000.s1.macro_auc_ovr :: sha256 b56bf70be9982044 :: ratio of the seed means of 1 - macro AUC, this row over the reference
\newcommand{\FtOmaRatioJciiEThreeFourthreeOverOnesixtwo}{0.88}  % experiments/FIGS/data/w2b_leg1_metrics_v2.json :: cells.*.N1000.s1.macro_auc_ovr :: sha256 b56bf70be9982044 :: ratio of the seed means of 1 - macro AUC, this row over the reference
\newcommand{\FtOmaRatioJciiEThreeOnesevenOverOnesixtwo}{1.02}  % experiments/FIGS/data/w2b_leg1_metrics_v2.json :: cells.*.N1000.s1.macro_auc_ovr :: sha256 b56bf70be9982044 :: ratio of the seed means of 1 - macro AUC, this row over the reference
\newcommand{\FtOmaRatioJciiEThreeOnesixtwoMassOverOnesixtwo}{0.85}  % experiments/FIGS/data/w2b_leg1_metrics_v2.json :: cells.*.N1000.s1.macro_auc_ovr :: sha256 b56bf70be9982044 :: ratio of the seed means of 1 - macro AUC, this row over the reference
\newcommand{\FtOmaRatioJciiEThreeOnesevenMassOverOnesixtwo}{1.02}  % experiments/FIGS/data/w2b_leg1_metrics_v2.json :: cells.*.N1000.s1.macro_auc_ovr :: sha256 b56bf70be9982044 :: ratio of the seed means of 1 - macro AUC, this row over the reference
\newcommand{\FtOmaRatioJciiEThreeRandDrawOneOverOnesixtwo}{0.84}  % experiments/FIGS/data/w2b_leg1_metrics_v2.json :: cells.*.N1000.s1.macro_auc_ovr :: sha256 b56bf70be9982044 :: ratio of the seed means of 1 - macro AUC, this row over the reference
\newcommand{\FtOmaRatioJciiEThreeRandDrawTwoOverOnesixtwo}{0.91}  % experiments/FIGS/data/w2b_leg1_metrics_v2.json :: cells.*.N1000.s1.macro_auc_ovr :: sha256 b56bf70be9982044 :: ratio of the seed means of 1 - macro AUC, this row over the reference
\newcommand{\FtOmaRatioJciiEThreeRandDrawThreeOverOnesixtwo}{0.87}  % experiments/FIGS/data/w2b_leg1_metrics_v2.json :: cells.*.N1000.s1.macro_auc_ovr :: sha256 b56bf70be9982044 :: ratio of the seed means of 1 - macro AUC, this row over the reference
\newcommand{\FtOmaRatioJciiEFourOnesixtwoOverOneeighteight}{0.99}  % experiments/FIGS/data/w2b_leg1_metrics_v2.json :: cells.*.N10000.s1.macro_auc_ovr :: sha256 b56bf70be9982044 :: ratio of the seed means of 1 - macro AUC, this row over the reference
\newcommand{\FtOmaRatioJciiEFourFourthreeOverOneeighteight}{1.20}  % experiments/FIGS/data/w2b_leg1_metrics_v2.json :: cells.*.N10000.s1.macro_auc_ovr :: sha256 b56bf70be9982044 :: ratio of the seed means of 1 - macro AUC, this row over the reference
\newcommand{\FtOmaRatioJciiEFourOnesevenOverOneeighteight}{1.32}  % experiments/FIGS/data/w2b_leg1_metrics_v2.json :: cells.*.N10000.s1.macro_auc_ovr :: sha256 b56bf70be9982044 :: ratio of the seed means of 1 - macro AUC, this row over the reference
\newcommand{\FtOmaRatioJciiEFourOnesixtwoMassOverOneeighteight}{0.99}  % experiments/FIGS/data/w2b_leg1_metrics_v2.json :: cells.*.N10000.s1.macro_auc_ovr :: sha256 b56bf70be9982044 :: ratio of the seed means of 1 - macro AUC, this row over the reference
\newcommand{\FtOmaRatioJciiEFourOnesevenMassOverOneeighteight}{1.30}  % experiments/FIGS/data/w2b_leg1_metrics_v2.json :: cells.*.N10000.s1.macro_auc_ovr :: sha256 b56bf70be9982044 :: ratio of the seed means of 1 - macro AUC, this row over the reference
\newcommand{\FtOmaRatioJciiEFourRandDrawOneOverOneeighteight}{1.17}  % experiments/FIGS/data/w2b_leg1_metrics_v2.json :: cells.*.N10000.s1.macro_auc_ovr :: sha256 b56bf70be9982044 :: ratio of the seed means of 1 - macro AUC, this row over the reference
\newcommand{\FtOmaRatioJciiEFourRandDrawTwoOverOneeighteight}{1.17}  % experiments/FIGS/data/w2b_leg1_metrics_v2.json :: cells.*.N10000.s1.macro_auc_ovr :: sha256 b56bf70be9982044 :: ratio of the seed means of 1 - macro AUC, this row over the reference
\newcommand{\FtOmaRatioJciiEFourRandDrawThreeOverOneeighteight}{1.13}  % experiments/FIGS/data/w2b_leg1_metrics_v2.json :: cells.*.N10000.s1.macro_auc_ovr :: sha256 b56bf70be9982044 :: ratio of the seed means of 1 - macro AUC, this row over the reference
\newcommand{\FtOmaRatioJciiEFourOneeighteightOverOnesixtwo}{1.01}  % experiments/FIGS/data/w2b_leg1_metrics_v2.json :: cells.*.N10000.s1.macro_auc_ovr :: sha256 b56bf70be9982044 :: ratio of the seed means of 1 - macro AUC, this row over the reference
\newcommand{\FtOmaRatioJciiEFourFourthreeOverOnesixtwo}{1.21}  % experiments/FIGS/data/w2b_leg1_metrics_v2.json :: cells.*.N10000.s1.macro_auc_ovr :: sha256 b56bf70be9982044 :: ratio of the seed means of 1 - macro AUC, this row over the reference
\newcommand{\FtOmaRatioJciiEFourOnesevenOverOnesixtwo}{1.32}  % experiments/FIGS/data/w2b_leg1_metrics_v2.json :: cells.*.N10000.s1.macro_auc_ovr :: sha256 b56bf70be9982044 :: ratio of the seed means of 1 - macro AUC, this row over the reference
\newcommand{\FtOmaRatioJciiEFourOnesixtwoMassOverOnesixtwo}{1.00}  % experiments/FIGS/data/w2b_leg1_metrics_v2.json :: cells.*.N10000.s1.macro_auc_ovr :: sha256 b56bf70be9982044 :: ratio of the seed means of 1 - macro AUC, this row over the reference
\newcommand{\FtOmaRatioJciiEFourOnesevenMassOverOnesixtwo}{1.31}  % experiments/FIGS/data/w2b_leg1_metrics_v2.json :: cells.*.N10000.s1.macro_auc_ovr :: sha256 b56bf70be9982044 :: ratio of the seed means of 1 - macro AUC, this row over the reference
\newcommand{\FtOmaRatioJciiEFourRandDrawOneOverOnesixtwo}{1.18}  % experiments/FIGS/data/w2b_leg1_metrics_v2.json :: cells.*.N10000.s1.macro_auc_ovr :: sha256 b56bf70be9982044 :: ratio of the seed means of 1 - macro AUC, this row over the reference
\newcommand{\FtOmaRatioJciiEFourRandDrawTwoOverOnesixtwo}{1.18}  % experiments/FIGS/data/w2b_leg1_metrics_v2.json :: cells.*.N10000.s1.macro_auc_ovr :: sha256 b56bf70be9982044 :: ratio of the seed means of 1 - macro AUC, this row over the reference
\newcommand{\FtOmaRatioJciiEFourRandDrawThreeOverOnesixtwo}{1.14}  % experiments/FIGS/data/w2b_leg1_metrics_v2.json :: cells.*.N10000.s1.macro_auc_ovr :: sha256 b56bf70be9982044 :: ratio of the seed means of 1 - macro AUC, this row over the reference
\newcommand{\FtOmaRatioJciiEFiveOnesixtwoOverOneeighteight}{1.00}  % experiments/FIGS/data/w2b_leg1_metrics_v2.json :: cells.*.N100000.s1.macro_auc_ovr :: sha256 b56bf70be9982044 :: ratio of the seed means of 1 - macro AUC, this row over the reference
\newcommand{\FtOmaRatioJciiEFiveFourthreeOverOneeighteight}{1.09}  % experiments/FIGS/data/w2b_leg1_metrics_v2.json :: cells.*.N100000.s1.macro_auc_ovr :: sha256 b56bf70be9982044 :: ratio of the seed means of 1 - macro AUC, this row over the reference
\newcommand{\FtOmaRatioJciiEFiveOnesevenOverOneeighteight}{1.22}  % experiments/FIGS/data/w2b_leg1_metrics_v2.json :: cells.*.N100000.s1.macro_auc_ovr :: sha256 b56bf70be9982044 :: ratio of the seed means of 1 - macro AUC, this row over the reference
\newcommand{\FtOmaRatioJciiEFiveOnesixtwoMassOverOneeighteight}{1.01}  % experiments/FIGS/data/w2b_leg1_metrics_v2.json :: cells.*.N100000.s1.macro_auc_ovr :: sha256 b56bf70be9982044 :: ratio of the seed means of 1 - macro AUC, this row over the reference
\newcommand{\FtOmaRatioJciiEFiveOnesevenMassOverOneeighteight}{1.26}  % experiments/FIGS/data/w2b_leg1_metrics_v2.json :: cells.*.N100000.s1.macro_auc_ovr :: sha256 b56bf70be9982044 :: ratio of the seed means of 1 - macro AUC, this row over the reference
\newcommand{\FtOmaRatioJciiEFiveRandDrawOneOverOneeighteight}{1.15}  % experiments/FIGS/data/w2b_leg1_metrics_v2.json :: cells.*.N100000.s1.macro_auc_ovr :: sha256 b56bf70be9982044 :: ratio of the seed means of 1 - macro AUC, this row over the reference
\newcommand{\FtOmaRatioJciiEFiveRandDrawTwoOverOneeighteight}{1.12}  % experiments/FIGS/data/w2b_leg1_metrics_v2.json :: cells.*.N100000.s1.macro_auc_ovr :: sha256 b56bf70be9982044 :: ratio of the seed means of 1 - macro AUC, this row over the reference
\newcommand{\FtOmaRatioJciiEFiveRandDrawThreeOverOneeighteight}{1.15}  % experiments/FIGS/data/w2b_leg1_metrics_v2.json :: cells.*.N100000.s1.macro_auc_ovr :: sha256 b56bf70be9982044 :: ratio of the seed means of 1 - macro AUC, this row over the reference
\newcommand{\FtOmaRatioJciiEFiveOneeighteightOverOnesixtwo}{1.00}  % experiments/FIGS/data/w2b_leg1_metrics_v2.json :: cells.*.N100000.s1.macro_auc_ovr :: sha256 b56bf70be9982044 :: ratio of the seed means of 1 - macro AUC, this row over the reference
\newcommand{\FtOmaRatioJciiEFiveFourthreeOverOnesixtwo}{1.09}  % experiments/FIGS/data/w2b_leg1_metrics_v2.json :: cells.*.N100000.s1.macro_auc_ovr :: sha256 b56bf70be9982044 :: ratio of the seed means of 1 - macro AUC, this row over the reference
\newcommand{\FtOmaRatioJciiEFiveOnesevenOverOnesixtwo}{1.22}  % experiments/FIGS/data/w2b_leg1_metrics_v2.json :: cells.*.N100000.s1.macro_auc_ovr :: sha256 b56bf70be9982044 :: ratio of the seed means of 1 - macro AUC, this row over the reference
\newcommand{\FtOmaRatioJciiEFiveOnesixtwoMassOverOnesixtwo}{1.02}  % experiments/FIGS/data/w2b_leg1_metrics_v2.json :: cells.*.N100000.s1.macro_auc_ovr :: sha256 b56bf70be9982044 :: ratio of the seed means of 1 - macro AUC, this row over the reference
\newcommand{\FtOmaRatioJciiEFiveOnesevenMassOverOnesixtwo}{1.26}  % experiments/FIGS/data/w2b_leg1_metrics_v2.json :: cells.*.N100000.s1.macro_auc_ovr :: sha256 b56bf70be9982044 :: ratio of the seed means of 1 - macro AUC, this row over the reference
\newcommand{\FtOmaRatioJciiEFiveRandDrawOneOverOnesixtwo}{1.15}  % experiments/FIGS/data/w2b_leg1_metrics_v2.json :: cells.*.N100000.s1.macro_auc_ovr :: sha256 b56bf70be9982044 :: ratio of the seed means of 1 - macro AUC, this row over the reference
\newcommand{\FtOmaRatioJciiEFiveRandDrawTwoOverOnesixtwo}{1.13}  % experiments/FIGS/data/w2b_leg1_metrics_v2.json :: cells.*.N100000.s1.macro_auc_ovr :: sha256 b56bf70be9982044 :: ratio of the seed means of 1 - macro AUC, this row over the reference
\newcommand{\FtOmaRatioJciiEFiveRandDrawThreeOverOnesixtwo}{1.16}  % experiments/FIGS/data/w2b_leg1_metrics_v2.json :: cells.*.N100000.s1.macro_auc_ovr :: sha256 b56bf70be9982044 :: ratio of the seed means of 1 - macro AUC, this row over the reference
\newcommand{\FtOmaRatioJciiESixOnesixtwoOverOneeighteight}{1.00}  % experiments/FIGS/data/w2b_leg1_metrics_v2.json :: cells.*.N1000000.s1.macro_auc_ovr :: sha256 b56bf70be9982044 :: ratio of the seed means of 1 - macro AUC, this row over the reference
\newcommand{\FtOmaRatioJciiESixFourthreeOverOneeighteight}{1.05}  % experiments/FIGS/data/w2b_leg1_metrics_v2.json :: cells.*.N1000000.s1.macro_auc_ovr :: sha256 b56bf70be9982044 :: ratio of the seed means of 1 - macro AUC, this row over the reference
\newcommand{\FtOmaRatioJciiESixOnesevenOverOneeighteight}{1.14}  % experiments/FIGS/data/w2b_leg1_metrics_v2.json :: cells.*.N1000000.s1.macro_auc_ovr :: sha256 b56bf70be9982044 :: ratio of the seed means of 1 - macro AUC, this row over the reference
\newcommand{\FtOmaRatioJciiESixOnesixtwoMassOverOneeighteight}{1.02}  % experiments/FIGS/data/w2b_leg1_metrics_v2.json :: cells.*.N1000000.s1.macro_auc_ovr :: sha256 b56bf70be9982044 :: ratio of the seed means of 1 - macro AUC, this row over the reference
\newcommand{\FtOmaRatioJciiESixOnesevenMassOverOneeighteight}{1.18}  % experiments/FIGS/data/w2b_leg1_metrics_v2.json :: cells.*.N1000000.s1.macro_auc_ovr :: sha256 b56bf70be9982044 :: ratio of the seed means of 1 - macro AUC, this row over the reference
\newcommand{\FtOmaRatioJciiESixRandDrawOneOverOneeighteight}{1.08}  % experiments/FIGS/data/w2b_leg1_metrics_v2.json :: cells.*.N1000000.s1.macro_auc_ovr :: sha256 b56bf70be9982044 :: ratio of the seed means of 1 - macro AUC, this row over the reference
\newcommand{\FtOmaRatioJciiESixRandDrawTwoOverOneeighteight}{1.07}  % experiments/FIGS/data/w2b_leg1_metrics_v2.json :: cells.*.N1000000.s1.macro_auc_ovr :: sha256 b56bf70be9982044 :: ratio of the seed means of 1 - macro AUC, this row over the reference
\newcommand{\FtOmaRatioJciiESixRandDrawThreeOverOneeighteight}{1.07}  % experiments/FIGS/data/w2b_leg1_metrics_v2.json :: cells.*.N1000000.s1.macro_auc_ovr :: sha256 b56bf70be9982044 :: ratio of the seed means of 1 - macro AUC, this row over the reference
\newcommand{\FtOmaRatioJciiESixOneeighteightOverOnesixtwo}{1.00}  % experiments/FIGS/data/w2b_leg1_metrics_v2.json :: cells.*.N1000000.s1.macro_auc_ovr :: sha256 b56bf70be9982044 :: ratio of the seed means of 1 - macro AUC, this row over the reference
\newcommand{\FtOmaRatioJciiESixFourthreeOverOnesixtwo}{1.04}  % experiments/FIGS/data/w2b_leg1_metrics_v2.json :: cells.*.N1000000.s1.macro_auc_ovr :: sha256 b56bf70be9982044 :: ratio of the seed means of 1 - macro AUC, this row over the reference
\newcommand{\FtOmaRatioJciiESixOnesevenOverOnesixtwo}{1.14}  % experiments/FIGS/data/w2b_leg1_metrics_v2.json :: cells.*.N1000000.s1.macro_auc_ovr :: sha256 b56bf70be9982044 :: ratio of the seed means of 1 - macro AUC, this row over the reference
\newcommand{\FtOmaRatioJciiESixOnesixtwoMassOverOnesixtwo}{1.02}  % experiments/FIGS/data/w2b_leg1_metrics_v2.json :: cells.*.N1000000.s1.macro_auc_ovr :: sha256 b56bf70be9982044 :: ratio of the seed means of 1 - macro AUC, this row over the reference
\newcommand{\FtOmaRatioJciiESixOnesevenMassOverOnesixtwo}{1.18}  % experiments/FIGS/data/w2b_leg1_metrics_v2.json :: cells.*.N1000000.s1.macro_auc_ovr :: sha256 b56bf70be9982044 :: ratio of the seed means of 1 - macro AUC, this row over the reference
\newcommand{\FtOmaRatioJciiESixRandDrawOneOverOnesixtwo}{1.07}  % experiments/FIGS/data/w2b_leg1_metrics_v2.json :: cells.*.N1000000.s1.macro_auc_ovr :: sha256 b56bf70be9982044 :: ratio of the seed means of 1 - macro AUC, this row over the reference
\newcommand{\FtOmaRatioJciiESixRandDrawTwoOverOnesixtwo}{1.07}  % experiments/FIGS/data/w2b_leg1_metrics_v2.json :: cells.*.N1000000.s1.macro_auc_ovr :: sha256 b56bf70be9982044 :: ratio of the seed means of 1 - macro AUC, this row over the reference
\newcommand{\FtOmaRatioJciiESixRandDrawThreeOverOnesixtwo}{1.07}  % experiments/FIGS/data/w2b_leg1_metrics_v2.json :: cells.*.N1000000.s1.macro_auc_ovr :: sha256 b56bf70be9982044 :: ratio of the seed means of 1 - macro AUC, this row over the reference
\newcommand{\PairedFtAccJciiEThreeOnesevenMinusOneeighteight}{\ensuremath{-}0.1132~[\ensuremath{-}0.1214, \ensuremath{-}0.1051]}  % experiments/FIGS/data/w2b_leg1_metrics_v2.json :: cells.{r16q1-s1,r16q1-s2,r16q1-s3,r16q1-s4,r16q1-s5,l188-s1,l188-s2,l188-s3,l188-s4,l188-s5}.N1000.s1.accuracy :: sha256 b56bf70be9982044 :: accuracy, coarse minus fine, paired by run [95% Student-t interval]
\newcommand{\PairedFtAccJciiEFourOnesevenMinusOneeighteight}{\ensuremath{-}0.0730~[\ensuremath{-}0.0902, \ensuremath{-}0.0558]}  % experiments/FIGS/data/w2b_leg1_metrics_v2.json :: cells.{r16q1-s1,r16q1-s2,r16q1-s3,r16q1-s4,r16q1-s5,l188-s1,l188-s2,l188-s3,l188-s4,l188-s5}.N10000.s1.accuracy :: sha256 b56bf70be9982044 :: accuracy, coarse minus fine, paired by run [95% Student-t interval]
\newcommand{\PairedFtAccJciiEFiveOnesevenMinusOneeighteight}{\ensuremath{-}0.0301~[\ensuremath{-}0.0319, \ensuremath{-}0.0284]}  % experiments/FIGS/data/w2b_leg1_metrics_v2.json :: cells.{r16q1-s1,r16q1-s2,r16q1-s3,r16q1-s4,r16q1-s5,l188-s1,l188-s2,l188-s3,l188-s4,l188-s5}.N100000.s1.accuracy :: sha256 b56bf70be9982044 :: accuracy, coarse minus fine, paired by run [95% Student-t interval]
\newcommand{\PairedFtAccJciiESixOnesevenMinusOneeighteight}{\ensuremath{-}0.0216~[\ensuremath{-}0.0232, \ensuremath{-}0.0199]}  % experiments/FIGS/data/w2b_leg1_metrics_v2.json :: cells.{r16q1-s1,r16q1-s2,r16q1-s3,r16q1-s4,r16q1-s5,l188-s1,l188-s2,l188-s3,l188-s4,l188-s5}.N1000000.s1.accuracy :: sha256 b56bf70be9982044 :: accuracy, coarse minus fine, paired by run [95% Student-t interval]
\newcommand{\FtNClassesJcii}{162}  % experiments/FIGS/data/w2b_leg1_metrics_v2.json :: cells.*.*.s1.n_classes_present :: sha256 b56bf70be9982044 :: classes of the fine-tuning task
\newcommand{\FtNTestAucJcii}{125{,}000}  % experiments/FIGS/data/w2b_leg1_metrics_v2.json :: cells.*.*.s1.n_jets_auc :: sha256 b56bf70be9982044 :: test jets the macro AUC is computed on
\newcommand{\FtNTestAccJcii}{500{,}000}  % experiments/FIGS/data/w2b_leg1_metrics_v2.json :: cells.*.*.s1.n_jets :: sha256 b56bf70be9982044 :: test jets the accuracy is computed on
\newcommand{\FtAucJcEThreeOneeighteight}{0.950\,\ensuremath{\pm}\,0.004}  % experiments/FIGS/data/w2b_leg2_metrics_v2.json :: cells.{l188-s1,l188-s2,l188-s3,l188-s4,l188-s5}.N1000.s1.macro_auc_ovr :: sha256 9cd059d3e2f51798 :: mean +- SD over the row's 5 pretrained models
\newcommand{\FtAucJcEFourOneeighteight}{0.9765\,\ensuremath{\pm}\,0.0008}  % experiments/FIGS/data/w2b_leg2_metrics_v2.json :: cells.{l188-s1,l188-s2,l188-s3,l188-s4,l188-s5}.N10000.s1.macro_auc_ovr :: sha256 9cd059d3e2f51798 :: mean +- SD over the row's 5 pretrained models
\newcommand{\FtAucJcEFiveOneeighteight}{0.9833\,\ensuremath{\pm}\,0.0002}  % experiments/FIGS/data/w2b_leg2_metrics_v2.json :: cells.{l188-s1,l188-s2,l188-s3,l188-s4,l188-s5}.N100000.s1.macro_auc_ovr :: sha256 9cd059d3e2f51798 :: mean +- SD over the row's 5 pretrained models
\newcommand{\FtAucJcESixOneeighteight}{0.9858\,\ensuremath{\pm}\,0.0001}  % experiments/FIGS/data/w2b_leg2_metrics_v2.json :: cells.{l188-s1,l188-s2,l188-s3,l188-s4,l188-s5}.N1000000.s1.macro_auc_ovr :: sha256 9cd059d3e2f51798 :: mean +- SD over the row's 5 pretrained models
\newcommand{\FtAucJcEThreeOnesixtwo}{0.9526\,\ensuremath{\pm}\,0.0020}  % experiments/FIGS/data/w2b_leg2_metrics_v2.json :: cells.{l162-s1b,l162-s2,l162-s3,l162-s4,l162-s5}.N1000.s1.macro_auc_ovr :: sha256 9cd059d3e2f51798 :: mean +- SD over the row's 5 pretrained models
\newcommand{\FtAucJcEFourOnesixtwo}{0.9770\,\ensuremath{\pm}\,0.0007}  % experiments/FIGS/data/w2b_leg2_metrics_v2.json :: cells.{l162-s1b,l162-s2,l162-s3,l162-s4,l162-s5}.N10000.s1.macro_auc_ovr :: sha256 9cd059d3e2f51798 :: mean +- SD over the row's 5 pretrained models
\newcommand{\FtAucJcEFiveOnesixtwo}{0.9836\,\ensuremath{\pm}\,0.0001}  % experiments/FIGS/data/w2b_leg2_metrics_v2.json :: cells.{l162-s1b,l162-s2,l162-s3,l162-s4,l162-s5}.N100000.s1.macro_auc_ovr :: sha256 9cd059d3e2f51798 :: mean +- SD over the row's 5 pretrained models
\newcommand{\FtAucJcESixOnesixtwo}{0.9858\,\ensuremath{\pm}\,0.0001}  % experiments/FIGS/data/w2b_leg2_metrics_v2.json :: cells.{l162-s1b,l162-s2,l162-s3,l162-s4,l162-s5}.N1000000.s1.macro_auc_ovr :: sha256 9cd059d3e2f51798 :: mean +- SD over the row's 5 pretrained models
\newcommand{\FtAucJcEThreeFourthree}{0.950\,\ensuremath{\pm}\,0.004}  % experiments/FIGS/data/w2b_leg2_metrics_v2.json :: cells.{r42q1-s1,r42q1-s2,r42q1-s3,r42q1-s4,r42q1-s5}.N1000.s1.macro_auc_ovr :: sha256 9cd059d3e2f51798 :: mean +- SD over the row's 5 pretrained models
\newcommand{\FtAucJcEFourFourthree}{0.9769\,\ensuremath{\pm}\,0.0005}  % experiments/FIGS/data/w2b_leg2_metrics_v2.json :: cells.{r42q1-s1,r42q1-s2,r42q1-s3,r42q1-s4,r42q1-s5}.N10000.s1.macro_auc_ovr :: sha256 9cd059d3e2f51798 :: mean +- SD over the row's 5 pretrained models
\newcommand{\FtAucJcEFiveFourthree}{0.9831\,\ensuremath{\pm}\,0.0002}  % experiments/FIGS/data/w2b_leg2_metrics_v2.json :: cells.{r42q1-s1,r42q1-s2,r42q1-s3,r42q1-s4,r42q1-s5}.N100000.s1.macro_auc_ovr :: sha256 9cd059d3e2f51798 :: mean +- SD over the row's 5 pretrained models
\newcommand{\FtAucJcESixFourthree}{0.9855\,\ensuremath{\pm}\,0.0001}  % experiments/FIGS/data/w2b_leg2_metrics_v2.json :: cells.{r42q1-s1,r42q1-s2,r42q1-s3,r42q1-s4,r42q1-s5}.N1000000.s1.macro_auc_ovr :: sha256 9cd059d3e2f51798 :: mean +- SD over the row's 5 pretrained models
\newcommand{\FtAucJcEThreeOneseven}{0.9470\,\ensuremath{\pm}\,0.0024}  % experiments/FIGS/data/w2b_leg2_metrics_v2.json :: cells.{r16q1-s1,r16q1-s2,r16q1-s3,r16q1-s4,r16q1-s5}.N1000.s1.macro_auc_ovr :: sha256 9cd059d3e2f51798 :: mean +- SD over the row's 5 pretrained models
\newcommand{\FtAucJcEFourOneseven}{0.9749\,\ensuremath{\pm}\,0.0006}  % experiments/FIGS/data/w2b_leg2_metrics_v2.json :: cells.{r16q1-s1,r16q1-s2,r16q1-s3,r16q1-s4,r16q1-s5}.N10000.s1.macro_auc_ovr :: sha256 9cd059d3e2f51798 :: mean +- SD over the row's 5 pretrained models
\newcommand{\FtAucJcEFiveOneseven}{0.9821\,\ensuremath{\pm}\,\ensuremath{<}0.0001}  % experiments/FIGS/data/w2b_leg2_metrics_v2.json :: cells.{r16q1-s1,r16q1-s2,r16q1-s3,r16q1-s4,r16q1-s5}.N100000.s1.macro_auc_ovr :: sha256 9cd059d3e2f51798 :: mean +- SD over the row's 5 pretrained models
\newcommand{\FtAucJcESixOneseven}{0.9846\,\ensuremath{\pm}\,0.0001}  % experiments/FIGS/data/w2b_leg2_metrics_v2.json :: cells.{r16q1-s1,r16q1-s2,r16q1-s3,r16q1-s4,r16q1-s5}.N1000000.s1.macro_auc_ovr :: sha256 9cd059d3e2f51798 :: mean +- SD over the row's 5 pretrained models
\newcommand{\FtAucJcEThreeOnesixtwoMass}{0.9569\,\ensuremath{\pm}\,0.0016}  % experiments/FIGS/data/w2b_leg2_metrics_v2.json :: cells.{l162mass-s1,l162mass-s2,l162mass-s3,l162mass-s4,l162mass-s5}.N1000.s1.macro_auc_ovr :: sha256 9cd059d3e2f51798 :: mean +- SD over the row's 5 pretrained models
\newcommand{\FtAucJcEFourOnesixtwoMass}{0.9768\,\ensuremath{\pm}\,0.0005}  % experiments/FIGS/data/w2b_leg2_metrics_v2.json :: cells.{l162mass-s1,l162mass-s2,l162mass-s3,l162mass-s4,l162mass-s5}.N10000.s1.macro_auc_ovr :: sha256 9cd059d3e2f51798 :: mean +- SD over the row's 5 pretrained models
\newcommand{\FtAucJcEFiveOnesixtwoMass}{0.9834\,\ensuremath{\pm}\,0.0001}  % experiments/FIGS/data/w2b_leg2_metrics_v2.json :: cells.{l162mass-s1,l162mass-s2,l162mass-s3,l162mass-s4,l162mass-s5}.N100000.s1.macro_auc_ovr :: sha256 9cd059d3e2f51798 :: mean +- SD over the row's 5 pretrained models
\newcommand{\FtAucJcESixOnesixtwoMass}{0.9856\,\ensuremath{\pm}\,0.0001}  % experiments/FIGS/data/w2b_leg2_metrics_v2.json :: cells.{l162mass-s1,l162mass-s2,l162mass-s3,l162mass-s4,l162mass-s5}.N1000000.s1.macro_auc_ovr :: sha256 9cd059d3e2f51798 :: mean +- SD over the row's 5 pretrained models
\newcommand{\FtAucJcEThreeOnesevenMass}{0.9460\,\ensuremath{\pm}\,0.0029}  % experiments/FIGS/data/w2b_leg2_metrics_v2.json :: cells.{r16q1mass-s1,r16q1mass-s2,r16q1mass-s3,r16q1mass-s4,r16q1mass-s5}.N1000.s1.macro_auc_ovr :: sha256 9cd059d3e2f51798 :: mean +- SD over the row's 5 pretrained models
\newcommand{\FtAucJcEFourOnesevenMass}{0.9737\,\ensuremath{\pm}\,0.0008}  % experiments/FIGS/data/w2b_leg2_metrics_v2.json :: cells.{r16q1mass-s1,r16q1mass-s2,r16q1mass-s3,r16q1mass-s4,r16q1mass-s5}.N10000.s1.macro_auc_ovr :: sha256 9cd059d3e2f51798 :: mean +- SD over the row's 5 pretrained models
\newcommand{\FtAucJcEFiveOnesevenMass}{0.9817\,\ensuremath{\pm}\,0.0002}  % experiments/FIGS/data/w2b_leg2_metrics_v2.json :: cells.{r16q1mass-s1,r16q1mass-s2,r16q1mass-s3,r16q1mass-s4,r16q1mass-s5}.N100000.s1.macro_auc_ovr :: sha256 9cd059d3e2f51798 :: mean +- SD over the row's 5 pretrained models
\newcommand{\FtAucJcESixOnesevenMass}{0.9842\,\ensuremath{\pm}\,0.0002}  % experiments/FIGS/data/w2b_leg2_metrics_v2.json :: cells.{r16q1mass-s1,r16q1mass-s2,r16q1mass-s3,r16q1mass-s4,r16q1mass-s5}.N1000000.s1.macro_auc_ovr :: sha256 9cd059d3e2f51798 :: mean +- SD over the row's 5 pretrained models
\newcommand{\FtAucJcEThreeRandDrawOne}{0.9477}  % experiments/FIGS/data/w2b_leg2_metrics_v2.json :: cells.{rand-d1-s1b}.N1000.s1.macro_auc_ovr :: sha256 9cd059d3e2f51798 :: one run
\newcommand{\FtAucJcEFourRandDrawOne}{0.9745}  % experiments/FIGS/data/w2b_leg2_metrics_v2.json :: cells.{rand-d1-s1b}.N10000.s1.macro_auc_ovr :: sha256 9cd059d3e2f51798 :: one run
\newcommand{\FtAucJcEFiveRandDrawOne}{0.98265}  % experiments/FIGS/data/w2b_leg2_metrics_v2.json :: cells.{rand-d1-s1b}.N100000.s1.macro_auc_ovr :: sha256 9cd059d3e2f51798 :: one run
\newcommand{\FtAucJcESixRandDrawOne}{0.98513}  % experiments/FIGS/data/w2b_leg2_metrics_v2.json :: cells.{rand-d1-s1b}.N1000000.s1.macro_auc_ovr :: sha256 9cd059d3e2f51798 :: one run
\newcommand{\FtAucJcEThreeRandDrawTwo}{0.9501}  % experiments/FIGS/data/w2b_leg2_metrics_v2.json :: cells.{rand-d2-s2}.N1000.s1.macro_auc_ovr :: sha256 9cd059d3e2f51798 :: one run
\newcommand{\FtAucJcEFourRandDrawTwo}{0.9756}  % experiments/FIGS/data/w2b_leg2_metrics_v2.json :: cells.{rand-d2-s2}.N10000.s1.macro_auc_ovr :: sha256 9cd059d3e2f51798 :: one run
\newcommand{\FtAucJcEFiveRandDrawTwo}{0.98278}  % experiments/FIGS/data/w2b_leg2_metrics_v2.json :: cells.{rand-d2-s2}.N100000.s1.macro_auc_ovr :: sha256 9cd059d3e2f51798 :: one run
\newcommand{\FtAucJcESixRandDrawTwo}{0.98512}  % experiments/FIGS/data/w2b_leg2_metrics_v2.json :: cells.{rand-d2-s2}.N1000000.s1.macro_auc_ovr :: sha256 9cd059d3e2f51798 :: one run
\newcommand{\FtAucJcEThreeRandDrawThree}{0.9470}  % experiments/FIGS/data/w2b_leg2_metrics_v2.json :: cells.{rand-d3-s3}.N1000.s1.macro_auc_ovr :: sha256 9cd059d3e2f51798 :: one run
\newcommand{\FtAucJcEFourRandDrawThree}{0.9743}  % experiments/FIGS/data/w2b_leg2_metrics_v2.json :: cells.{rand-d3-s3}.N10000.s1.macro_auc_ovr :: sha256 9cd059d3e2f51798 :: one run
\newcommand{\FtAucJcEFiveRandDrawThree}{0.98253}  % experiments/FIGS/data/w2b_leg2_metrics_v2.json :: cells.{rand-d3-s3}.N100000.s1.macro_auc_ovr :: sha256 9cd059d3e2f51798 :: one run
\newcommand{\FtAucJcESixRandDrawThree}{0.98506}  % experiments/FIGS/data/w2b_leg2_metrics_v2.json :: cells.{rand-d3-s3}.N1000000.s1.macro_auc_ovr :: sha256 9cd059d3e2f51798 :: one run
\newcommand{\FtAccJcEThreeOneeighteight}{0.7183\,\ensuremath{\pm}\,0.0032}  % experiments/FIGS/data/w2b_leg2_metrics_v2.json :: cells.{l188-s1,l188-s2,l188-s3,l188-s4,l188-s5}.N1000.s1.accuracy :: sha256 9cd059d3e2f51798 :: mean +- SD over the row's 5 pretrained models
\newcommand{\FtAccJcEFourOneeighteight}{0.8039\,\ensuremath{\pm}\,0.0030}  % experiments/FIGS/data/w2b_leg2_metrics_v2.json :: cells.{l188-s1,l188-s2,l188-s3,l188-s4,l188-s5}.N10000.s1.accuracy :: sha256 9cd059d3e2f51798 :: mean +- SD over the row's 5 pretrained models
\newcommand{\FtAccJcEFiveOneeighteight}{0.8351\,\ensuremath{\pm}\,0.0008}  % experiments/FIGS/data/w2b_leg2_metrics_v2.json :: cells.{l188-s1,l188-s2,l188-s3,l188-s4,l188-s5}.N100000.s1.accuracy :: sha256 9cd059d3e2f51798 :: mean +- SD over the row's 5 pretrained models
\newcommand{\FtAccJcESixOneeighteight}{0.8486\,\ensuremath{\pm}\,0.0005}  % experiments/FIGS/data/w2b_leg2_metrics_v2.json :: cells.{l188-s1,l188-s2,l188-s3,l188-s4,l188-s5}.N1000000.s1.accuracy :: sha256 9cd059d3e2f51798 :: mean +- SD over the row's 5 pretrained models
\newcommand{\FtAccJcEThreeOnesixtwo}{0.7268\,\ensuremath{\pm}\,0.0017}  % experiments/FIGS/data/w2b_leg2_metrics_v2.json :: cells.{l162-s1b,l162-s2,l162-s3,l162-s4,l162-s5}.N1000.s1.accuracy :: sha256 9cd059d3e2f51798 :: mean +- SD over the row's 5 pretrained models
\newcommand{\FtAccJcEFourOnesixtwo}{0.8069\,\ensuremath{\pm}\,0.0025}  % experiments/FIGS/data/w2b_leg2_metrics_v2.json :: cells.{l162-s1b,l162-s2,l162-s3,l162-s4,l162-s5}.N10000.s1.accuracy :: sha256 9cd059d3e2f51798 :: mean +- SD over the row's 5 pretrained models
\newcommand{\FtAccJcEFiveOnesixtwo}{0.8363\,\ensuremath{\pm}\,0.0008}  % experiments/FIGS/data/w2b_leg2_metrics_v2.json :: cells.{l162-s1b,l162-s2,l162-s3,l162-s4,l162-s5}.N100000.s1.accuracy :: sha256 9cd059d3e2f51798 :: mean +- SD over the row's 5 pretrained models
\newcommand{\FtAccJcESixOnesixtwo}{0.8485\,\ensuremath{\pm}\,0.0005}  % experiments/FIGS/data/w2b_leg2_metrics_v2.json :: cells.{l162-s1b,l162-s2,l162-s3,l162-s4,l162-s5}.N1000000.s1.accuracy :: sha256 9cd059d3e2f51798 :: mean +- SD over the row's 5 pretrained models
\newcommand{\FtAccJcEThreeFourthree}{0.722\,\ensuremath{\pm}\,0.007}  % experiments/FIGS/data/w2b_leg2_metrics_v2.json :: cells.{r42q1-s1,r42q1-s2,r42q1-s3,r42q1-s4,r42q1-s5}.N1000.s1.accuracy :: sha256 9cd059d3e2f51798 :: mean +- SD over the row's 5 pretrained models
\newcommand{\FtAccJcEFourFourthree}{0.8050\,\ensuremath{\pm}\,0.0026}  % experiments/FIGS/data/w2b_leg2_metrics_v2.json :: cells.{r42q1-s1,r42q1-s2,r42q1-s3,r42q1-s4,r42q1-s5}.N10000.s1.accuracy :: sha256 9cd059d3e2f51798 :: mean +- SD over the row's 5 pretrained models
\newcommand{\FtAccJcEFiveFourthree}{0.8336\,\ensuremath{\pm}\,0.0008}  % experiments/FIGS/data/w2b_leg2_metrics_v2.json :: cells.{r42q1-s1,r42q1-s2,r42q1-s3,r42q1-s4,r42q1-s5}.N100000.s1.accuracy :: sha256 9cd059d3e2f51798 :: mean +- SD over the row's 5 pretrained models
\newcommand{\FtAccJcESixFourthree}{0.8468\,\ensuremath{\pm}\,0.0003}  % experiments/FIGS/data/w2b_leg2_metrics_v2.json :: cells.{r42q1-s1,r42q1-s2,r42q1-s3,r42q1-s4,r42q1-s5}.N1000000.s1.accuracy :: sha256 9cd059d3e2f51798 :: mean +- SD over the row's 5 pretrained models
\newcommand{\FtAccJcEThreeOneseven}{0.714\,\ensuremath{\pm}\,0.005}  % experiments/FIGS/data/w2b_leg2_metrics_v2.json :: cells.{r16q1-s1,r16q1-s2,r16q1-s3,r16q1-s4,r16q1-s5}.N1000.s1.accuracy :: sha256 9cd059d3e2f51798 :: mean +- SD over the row's 5 pretrained models
\newcommand{\FtAccJcEFourOneseven}{0.7964\,\ensuremath{\pm}\,0.0019}  % experiments/FIGS/data/w2b_leg2_metrics_v2.json :: cells.{r16q1-s1,r16q1-s2,r16q1-s3,r16q1-s4,r16q1-s5}.N10000.s1.accuracy :: sha256 9cd059d3e2f51798 :: mean +- SD over the row's 5 pretrained models
\newcommand{\FtAccJcEFiveOneseven}{0.8286\,\ensuremath{\pm}\,0.0006}  % experiments/FIGS/data/w2b_leg2_metrics_v2.json :: cells.{r16q1-s1,r16q1-s2,r16q1-s3,r16q1-s4,r16q1-s5}.N100000.s1.accuracy :: sha256 9cd059d3e2f51798 :: mean +- SD over the row's 5 pretrained models
\newcommand{\FtAccJcESixOneseven}{0.8415\,\ensuremath{\pm}\,0.0007}  % experiments/FIGS/data/w2b_leg2_metrics_v2.json :: cells.{r16q1-s1,r16q1-s2,r16q1-s3,r16q1-s4,r16q1-s5}.N1000000.s1.accuracy :: sha256 9cd059d3e2f51798 :: mean +- SD over the row's 5 pretrained models
\newcommand{\FtAccJcEThreeOnesixtwoMass}{0.728\,\ensuremath{\pm}\,0.004}  % experiments/FIGS/data/w2b_leg2_metrics_v2.json :: cells.{l162mass-s1,l162mass-s2,l162mass-s3,l162mass-s4,l162mass-s5}.N1000.s1.accuracy :: sha256 9cd059d3e2f51798 :: mean +- SD over the row's 5 pretrained models
\newcommand{\FtAccJcEFourOnesixtwoMass}{0.8043\,\ensuremath{\pm}\,0.0020}  % experiments/FIGS/data/w2b_leg2_metrics_v2.json :: cells.{l162mass-s1,l162mass-s2,l162mass-s3,l162mass-s4,l162mass-s5}.N10000.s1.accuracy :: sha256 9cd059d3e2f51798 :: mean +- SD over the row's 5 pretrained models
\newcommand{\FtAccJcEFiveOnesixtwoMass}{0.8349\,\ensuremath{\pm}\,0.0005}  % experiments/FIGS/data/w2b_leg2_metrics_v2.json :: cells.{l162mass-s1,l162mass-s2,l162mass-s3,l162mass-s4,l162mass-s5}.N100000.s1.accuracy :: sha256 9cd059d3e2f51798 :: mean +- SD over the row's 5 pretrained models
\newcommand{\FtAccJcESixOnesixtwoMass}{0.8475\,\ensuremath{\pm}\,0.0009}  % experiments/FIGS/data/w2b_leg2_metrics_v2.json :: cells.{l162mass-s1,l162mass-s2,l162mass-s3,l162mass-s4,l162mass-s5}.N1000000.s1.accuracy :: sha256 9cd059d3e2f51798 :: mean +- SD over the row's 5 pretrained models
\newcommand{\FtAccJcEThreeOnesevenMass}{0.702\,\ensuremath{\pm}\,0.005}  % experiments/FIGS/data/w2b_leg2_metrics_v2.json :: cells.{r16q1mass-s1,r16q1mass-s2,r16q1mass-s3,r16q1mass-s4,r16q1mass-s5}.N1000.s1.accuracy :: sha256 9cd059d3e2f51798 :: mean +- SD over the row's 5 pretrained models
\newcommand{\FtAccJcEFourOnesevenMass}{0.7909\,\ensuremath{\pm}\,0.0031}  % experiments/FIGS/data/w2b_leg2_metrics_v2.json :: cells.{r16q1mass-s1,r16q1mass-s2,r16q1mass-s3,r16q1mass-s4,r16q1mass-s5}.N10000.s1.accuracy :: sha256 9cd059d3e2f51798 :: mean +- SD over the row's 5 pretrained models
\newcommand{\FtAccJcEFiveOnesevenMass}{0.8265\,\ensuremath{\pm}\,0.0007}  % experiments/FIGS/data/w2b_leg2_metrics_v2.json :: cells.{r16q1mass-s1,r16q1mass-s2,r16q1mass-s3,r16q1mass-s4,r16q1mass-s5}.N100000.s1.accuracy :: sha256 9cd059d3e2f51798 :: mean +- SD over the row's 5 pretrained models
\newcommand{\FtAccJcESixOnesevenMass}{0.8390\,\ensuremath{\pm}\,0.0011}  % experiments/FIGS/data/w2b_leg2_metrics_v2.json :: cells.{r16q1mass-s1,r16q1mass-s2,r16q1mass-s3,r16q1mass-s4,r16q1mass-s5}.N1000000.s1.accuracy :: sha256 9cd059d3e2f51798 :: mean +- SD over the row's 5 pretrained models
\newcommand{\FtAccJcEThreeRandDrawOne}{0.7049}  % experiments/FIGS/data/w2b_leg2_metrics_v2.json :: cells.{rand-d1-s1b}.N1000.s1.accuracy :: sha256 9cd059d3e2f51798 :: one run
\newcommand{\FtAccJcEFourRandDrawOne}{0.7974}  % experiments/FIGS/data/w2b_leg2_metrics_v2.json :: cells.{rand-d1-s1b}.N10000.s1.accuracy :: sha256 9cd059d3e2f51798 :: one run
\newcommand{\FtAccJcEFiveRandDrawOne}{0.8326}  % experiments/FIGS/data/w2b_leg2_metrics_v2.json :: cells.{rand-d1-s1b}.N100000.s1.accuracy :: sha256 9cd059d3e2f51798 :: one run
\newcommand{\FtAccJcESixRandDrawOne}{0.84467}  % experiments/FIGS/data/w2b_leg2_metrics_v2.json :: cells.{rand-d1-s1b}.N1000000.s1.accuracy :: sha256 9cd059d3e2f51798 :: one run
\newcommand{\FtAccJcEThreeRandDrawTwo}{0.7135}  % experiments/FIGS/data/w2b_leg2_metrics_v2.json :: cells.{rand-d2-s2}.N1000.s1.accuracy :: sha256 9cd059d3e2f51798 :: one run
\newcommand{\FtAccJcEFourRandDrawTwo}{0.7998}  % experiments/FIGS/data/w2b_leg2_metrics_v2.json :: cells.{rand-d2-s2}.N10000.s1.accuracy :: sha256 9cd059d3e2f51798 :: one run
\newcommand{\FtAccJcEFiveRandDrawTwo}{0.8332}  % experiments/FIGS/data/w2b_leg2_metrics_v2.json :: cells.{rand-d2-s2}.N100000.s1.accuracy :: sha256 9cd059d3e2f51798 :: one run
\newcommand{\FtAccJcESixRandDrawTwo}{0.84449}  % experiments/FIGS/data/w2b_leg2_metrics_v2.json :: cells.{rand-d2-s2}.N1000000.s1.accuracy :: sha256 9cd059d3e2f51798 :: one run
\newcommand{\FtAccJcEThreeRandDrawThree}{0.7090}  % experiments/FIGS/data/w2b_leg2_metrics_v2.json :: cells.{rand-d3-s3}.N1000.s1.accuracy :: sha256 9cd059d3e2f51798 :: one run
\newcommand{\FtAccJcEFourRandDrawThree}{0.7948}  % experiments/FIGS/data/w2b_leg2_metrics_v2.json :: cells.{rand-d3-s3}.N10000.s1.accuracy :: sha256 9cd059d3e2f51798 :: one run
\newcommand{\FtAccJcEFiveRandDrawThree}{0.8310}  % experiments/FIGS/data/w2b_leg2_metrics_v2.json :: cells.{rand-d3-s3}.N100000.s1.accuracy :: sha256 9cd059d3e2f51798 :: one run
\newcommand{\FtAccJcESixRandDrawThree}{0.84409}  % experiments/FIGS/data/w2b_leg2_metrics_v2.json :: cells.{rand-d3-s3}.N1000000.s1.accuracy :: sha256 9cd059d3e2f51798 :: one run
\newcommand{\FtOmaRatioJcEThreeOnesixtwoOverOneeighteight}{0.96}  % experiments/FIGS/data/w2b_leg2_metrics_v2.json :: cells.*.N1000.s1.macro_auc_ovr :: sha256 9cd059d3e2f51798 :: ratio of the seed means of 1 - macro AUC, this row over the reference
\newcommand{\FtOmaRatioJcEThreeFourthreeOverOneeighteight}{1.01}  % experiments/FIGS/data/w2b_leg2_metrics_v2.json :: cells.*.N1000.s1.macro_auc_ovr :: sha256 9cd059d3e2f51798 :: ratio of the seed means of 1 - macro AUC, this row over the reference
\newcommand{\FtOmaRatioJcEThreeOnesevenOverOneeighteight}{1.07}  % experiments/FIGS/data/w2b_leg2_metrics_v2.json :: cells.*.N1000.s1.macro_auc_ovr :: sha256 9cd059d3e2f51798 :: ratio of the seed means of 1 - macro AUC, this row over the reference
\newcommand{\FtOmaRatioJcEThreeOnesixtwoMassOverOneeighteight}{0.87}  % experiments/FIGS/data/w2b_leg2_metrics_v2.json :: cells.*.N1000.s1.macro_auc_ovr :: sha256 9cd059d3e2f51798 :: ratio of the seed means of 1 - macro AUC, this row over the reference
\newcommand{\FtOmaRatioJcEThreeOnesevenMassOverOneeighteight}{1.09}  % experiments/FIGS/data/w2b_leg2_metrics_v2.json :: cells.*.N1000.s1.macro_auc_ovr :: sha256 9cd059d3e2f51798 :: ratio of the seed means of 1 - macro AUC, this row over the reference
\newcommand{\FtOmaRatioJcEThreeRandDrawOneOverOneeighteight}{1.06}  % experiments/FIGS/data/w2b_leg2_metrics_v2.json :: cells.*.N1000.s1.macro_auc_ovr :: sha256 9cd059d3e2f51798 :: ratio of the seed means of 1 - macro AUC, this row over the reference
\newcommand{\FtOmaRatioJcEThreeRandDrawTwoOverOneeighteight}{1.01}  % experiments/FIGS/data/w2b_leg2_metrics_v2.json :: cells.*.N1000.s1.macro_auc_ovr :: sha256 9cd059d3e2f51798 :: ratio of the seed means of 1 - macro AUC, this row over the reference
\newcommand{\FtOmaRatioJcEThreeRandDrawThreeOverOneeighteight}{1.07}  % experiments/FIGS/data/w2b_leg2_metrics_v2.json :: cells.*.N1000.s1.macro_auc_ovr :: sha256 9cd059d3e2f51798 :: ratio of the seed means of 1 - macro AUC, this row over the reference
\newcommand{\FtOmaRatioJcEThreeOneeighteightOverOnesixtwo}{1.05}  % experiments/FIGS/data/w2b_leg2_metrics_v2.json :: cells.*.N1000.s1.macro_auc_ovr :: sha256 9cd059d3e2f51798 :: ratio of the seed means of 1 - macro AUC, this row over the reference
\newcommand{\FtOmaRatioJcEThreeFourthreeOverOnesixtwo}{1.05}  % experiments/FIGS/data/w2b_leg2_metrics_v2.json :: cells.*.N1000.s1.macro_auc_ovr :: sha256 9cd059d3e2f51798 :: ratio of the seed means of 1 - macro AUC, this row over the reference
\newcommand{\FtOmaRatioJcEThreeOnesevenOverOnesixtwo}{1.12}  % experiments/FIGS/data/w2b_leg2_metrics_v2.json :: cells.*.N1000.s1.macro_auc_ovr :: sha256 9cd059d3e2f51798 :: ratio of the seed means of 1 - macro AUC, this row over the reference
\newcommand{\FtOmaRatioJcEThreeOnesixtwoMassOverOnesixtwo}{0.91}  % experiments/FIGS/data/w2b_leg2_metrics_v2.json :: cells.*.N1000.s1.macro_auc_ovr :: sha256 9cd059d3e2f51798 :: ratio of the seed means of 1 - macro AUC, this row over the reference
\newcommand{\FtOmaRatioJcEThreeOnesevenMassOverOnesixtwo}{1.14}  % experiments/FIGS/data/w2b_leg2_metrics_v2.json :: cells.*.N1000.s1.macro_auc_ovr :: sha256 9cd059d3e2f51798 :: ratio of the seed means of 1 - macro AUC, this row over the reference
\newcommand{\FtOmaRatioJcEThreeRandDrawOneOverOnesixtwo}{1.10}  % experiments/FIGS/data/w2b_leg2_metrics_v2.json :: cells.*.N1000.s1.macro_auc_ovr :: sha256 9cd059d3e2f51798 :: ratio of the seed means of 1 - macro AUC, this row over the reference
\newcommand{\FtOmaRatioJcEThreeRandDrawTwoOverOnesixtwo}{1.05}  % experiments/FIGS/data/w2b_leg2_metrics_v2.json :: cells.*.N1000.s1.macro_auc_ovr :: sha256 9cd059d3e2f51798 :: ratio of the seed means of 1 - macro AUC, this row over the reference
\newcommand{\FtOmaRatioJcEThreeRandDrawThreeOverOnesixtwo}{1.12}  % experiments/FIGS/data/w2b_leg2_metrics_v2.json :: cells.*.N1000.s1.macro_auc_ovr :: sha256 9cd059d3e2f51798 :: ratio of the seed means of 1 - macro AUC, this row over the reference
\newcommand{\FtOmaRatioJcEFourOnesixtwoOverOneeighteight}{0.98}  % experiments/FIGS/data/w2b_leg2_metrics_v2.json :: cells.*.N10000.s1.macro_auc_ovr :: sha256 9cd059d3e2f51798 :: ratio of the seed means of 1 - macro AUC, this row over the reference
\newcommand{\FtOmaRatioJcEFourFourthreeOverOneeighteight}{0.98}  % experiments/FIGS/data/w2b_leg2_metrics_v2.json :: cells.*.N10000.s1.macro_auc_ovr :: sha256 9cd059d3e2f51798 :: ratio of the seed means of 1 - macro AUC, this row over the reference
\newcommand{\FtOmaRatioJcEFourOnesevenOverOneeighteight}{1.07}  % experiments/FIGS/data/w2b_leg2_metrics_v2.json :: cells.*.N10000.s1.macro_auc_ovr :: sha256 9cd059d3e2f51798 :: ratio of the seed means of 1 - macro AUC, this row over the reference
\newcommand{\FtOmaRatioJcEFourOnesixtwoMassOverOneeighteight}{0.99}  % experiments/FIGS/data/w2b_leg2_metrics_v2.json :: cells.*.N10000.s1.macro_auc_ovr :: sha256 9cd059d3e2f51798 :: ratio of the seed means of 1 - macro AUC, this row over the reference
\newcommand{\FtOmaRatioJcEFourOnesevenMassOverOneeighteight}{1.12}  % experiments/FIGS/data/w2b_leg2_metrics_v2.json :: cells.*.N10000.s1.macro_auc_ovr :: sha256 9cd059d3e2f51798 :: ratio of the seed means of 1 - macro AUC, this row over the reference
\newcommand{\FtOmaRatioJcEFourRandDrawOneOverOneeighteight}{1.08}  % experiments/FIGS/data/w2b_leg2_metrics_v2.json :: cells.*.N10000.s1.macro_auc_ovr :: sha256 9cd059d3e2f51798 :: ratio of the seed means of 1 - macro AUC, this row over the reference
\newcommand{\FtOmaRatioJcEFourRandDrawTwoOverOneeighteight}{1.04}  % experiments/FIGS/data/w2b_leg2_metrics_v2.json :: cells.*.N10000.s1.macro_auc_ovr :: sha256 9cd059d3e2f51798 :: ratio of the seed means of 1 - macro AUC, this row over the reference
\newcommand{\FtOmaRatioJcEFourRandDrawThreeOverOneeighteight}{1.09}  % experiments/FIGS/data/w2b_leg2_metrics_v2.json :: cells.*.N10000.s1.macro_auc_ovr :: sha256 9cd059d3e2f51798 :: ratio of the seed means of 1 - macro AUC, this row over the reference
\newcommand{\FtOmaRatioJcEFourOneeighteightOverOnesixtwo}{1.02}  % experiments/FIGS/data/w2b_leg2_metrics_v2.json :: cells.*.N10000.s1.macro_auc_ovr :: sha256 9cd059d3e2f51798 :: ratio of the seed means of 1 - macro AUC, this row over the reference
\newcommand{\FtOmaRatioJcEFourFourthreeOverOnesixtwo}{1.01}  % experiments/FIGS/data/w2b_leg2_metrics_v2.json :: cells.*.N10000.s1.macro_auc_ovr :: sha256 9cd059d3e2f51798 :: ratio of the seed means of 1 - macro AUC, this row over the reference
\newcommand{\FtOmaRatioJcEFourOnesevenOverOnesixtwo}{1.09}  % experiments/FIGS/data/w2b_leg2_metrics_v2.json :: cells.*.N10000.s1.macro_auc_ovr :: sha256 9cd059d3e2f51798 :: ratio of the seed means of 1 - macro AUC, this row over the reference
\newcommand{\FtOmaRatioJcEFourOnesixtwoMassOverOnesixtwo}{1.01}  % experiments/FIGS/data/w2b_leg2_metrics_v2.json :: cells.*.N10000.s1.macro_auc_ovr :: sha256 9cd059d3e2f51798 :: ratio of the seed means of 1 - macro AUC, this row over the reference
\newcommand{\FtOmaRatioJcEFourOnesevenMassOverOnesixtwo}{1.14}  % experiments/FIGS/data/w2b_leg2_metrics_v2.json :: cells.*.N10000.s1.macro_auc_ovr :: sha256 9cd059d3e2f51798 :: ratio of the seed means of 1 - macro AUC, this row over the reference
\newcommand{\FtOmaRatioJcEFourRandDrawOneOverOnesixtwo}{1.11}  % experiments/FIGS/data/w2b_leg2_metrics_v2.json :: cells.*.N10000.s1.macro_auc_ovr :: sha256 9cd059d3e2f51798 :: ratio of the seed means of 1 - macro AUC, this row over the reference
\newcommand{\FtOmaRatioJcEFourRandDrawTwoOverOnesixtwo}{1.06}  % experiments/FIGS/data/w2b_leg2_metrics_v2.json :: cells.*.N10000.s1.macro_auc_ovr :: sha256 9cd059d3e2f51798 :: ratio of the seed means of 1 - macro AUC, this row over the reference
\newcommand{\FtOmaRatioJcEFourRandDrawThreeOverOnesixtwo}{1.12}  % experiments/FIGS/data/w2b_leg2_metrics_v2.json :: cells.*.N10000.s1.macro_auc_ovr :: sha256 9cd059d3e2f51798 :: ratio of the seed means of 1 - macro AUC, this row over the reference
\newcommand{\FtOmaRatioJcEFiveOnesixtwoOverOneeighteight}{0.98}  % experiments/FIGS/data/w2b_leg2_metrics_v2.json :: cells.*.N100000.s1.macro_auc_ovr :: sha256 9cd059d3e2f51798 :: ratio of the seed means of 1 - macro AUC, this row over the reference
\newcommand{\FtOmaRatioJcEFiveFourthreeOverOneeighteight}{1.02}  % experiments/FIGS/data/w2b_leg2_metrics_v2.json :: cells.*.N100000.s1.macro_auc_ovr :: sha256 9cd059d3e2f51798 :: ratio of the seed means of 1 - macro AUC, this row over the reference
\newcommand{\FtOmaRatioJcEFiveOnesevenOverOneeighteight}{1.07}  % experiments/FIGS/data/w2b_leg2_metrics_v2.json :: cells.*.N100000.s1.macro_auc_ovr :: sha256 9cd059d3e2f51798 :: ratio of the seed means of 1 - macro AUC, this row over the reference
\newcommand{\FtOmaRatioJcEFiveOnesixtwoMassOverOneeighteight}{1.00}  % experiments/FIGS/data/w2b_leg2_metrics_v2.json :: cells.*.N100000.s1.macro_auc_ovr :: sha256 9cd059d3e2f51798 :: ratio of the seed means of 1 - macro AUC, this row over the reference
\newcommand{\FtOmaRatioJcEFiveOnesevenMassOverOneeighteight}{1.10}  % experiments/FIGS/data/w2b_leg2_metrics_v2.json :: cells.*.N100000.s1.macro_auc_ovr :: sha256 9cd059d3e2f51798 :: ratio of the seed means of 1 - macro AUC, this row over the reference
\newcommand{\FtOmaRatioJcEFiveRandDrawOneOverOneeighteight}{1.04}  % experiments/FIGS/data/w2b_leg2_metrics_v2.json :: cells.*.N100000.s1.macro_auc_ovr :: sha256 9cd059d3e2f51798 :: ratio of the seed means of 1 - macro AUC, this row over the reference
\newcommand{\FtOmaRatioJcEFiveRandDrawTwoOverOneeighteight}{1.03}  % experiments/FIGS/data/w2b_leg2_metrics_v2.json :: cells.*.N100000.s1.macro_auc_ovr :: sha256 9cd059d3e2f51798 :: ratio of the seed means of 1 - macro AUC, this row over the reference
\newcommand{\FtOmaRatioJcEFiveRandDrawThreeOverOneeighteight}{1.05}  % experiments/FIGS/data/w2b_leg2_metrics_v2.json :: cells.*.N100000.s1.macro_auc_ovr :: sha256 9cd059d3e2f51798 :: ratio of the seed means of 1 - macro AUC, this row over the reference
\newcommand{\FtOmaRatioJcEFiveOneeighteightOverOnesixtwo}{1.02}  % experiments/FIGS/data/w2b_leg2_metrics_v2.json :: cells.*.N100000.s1.macro_auc_ovr :: sha256 9cd059d3e2f51798 :: ratio of the seed means of 1 - macro AUC, this row over the reference
\newcommand{\FtOmaRatioJcEFiveFourthreeOverOnesixtwo}{1.03}  % experiments/FIGS/data/w2b_leg2_metrics_v2.json :: cells.*.N100000.s1.macro_auc_ovr :: sha256 9cd059d3e2f51798 :: ratio of the seed means of 1 - macro AUC, this row over the reference
\newcommand{\FtOmaRatioJcEFiveOnesevenOverOnesixtwo}{1.09}  % experiments/FIGS/data/w2b_leg2_metrics_v2.json :: cells.*.N100000.s1.macro_auc_ovr :: sha256 9cd059d3e2f51798 :: ratio of the seed means of 1 - macro AUC, this row over the reference
\newcommand{\FtOmaRatioJcEFiveOnesixtwoMassOverOnesixtwo}{1.02}  % experiments/FIGS/data/w2b_leg2_metrics_v2.json :: cells.*.N100000.s1.macro_auc_ovr :: sha256 9cd059d3e2f51798 :: ratio of the seed means of 1 - macro AUC, this row over the reference
\newcommand{\FtOmaRatioJcEFiveOnesevenMassOverOnesixtwo}{1.11}  % experiments/FIGS/data/w2b_leg2_metrics_v2.json :: cells.*.N100000.s1.macro_auc_ovr :: sha256 9cd059d3e2f51798 :: ratio of the seed means of 1 - macro AUC, this row over the reference
\newcommand{\FtOmaRatioJcEFiveRandDrawOneOverOnesixtwo}{1.06}  % experiments/FIGS/data/w2b_leg2_metrics_v2.json :: cells.*.N100000.s1.macro_auc_ovr :: sha256 9cd059d3e2f51798 :: ratio of the seed means of 1 - macro AUC, this row over the reference
\newcommand{\FtOmaRatioJcEFiveRandDrawTwoOverOnesixtwo}{1.05}  % experiments/FIGS/data/w2b_leg2_metrics_v2.json :: cells.*.N100000.s1.macro_auc_ovr :: sha256 9cd059d3e2f51798 :: ratio of the seed means of 1 - macro AUC, this row over the reference
\newcommand{\FtOmaRatioJcEFiveRandDrawThreeOverOnesixtwo}{1.07}  % experiments/FIGS/data/w2b_leg2_metrics_v2.json :: cells.*.N100000.s1.macro_auc_ovr :: sha256 9cd059d3e2f51798 :: ratio of the seed means of 1 - macro AUC, this row over the reference
\newcommand{\FtOmaRatioJcESixOnesixtwoOverOneeighteight}{1.00}  % experiments/FIGS/data/w2b_leg2_metrics_v2.json :: cells.*.N1000000.s1.macro_auc_ovr :: sha256 9cd059d3e2f51798 :: ratio of the seed means of 1 - macro AUC, this row over the reference
\newcommand{\FtOmaRatioJcESixFourthreeOverOneeighteight}{1.02}  % experiments/FIGS/data/w2b_leg2_metrics_v2.json :: cells.*.N1000000.s1.macro_auc_ovr :: sha256 9cd059d3e2f51798 :: ratio of the seed means of 1 - macro AUC, this row over the reference
\newcommand{\FtOmaRatioJcESixOnesevenOverOneeighteight}{1.08}  % experiments/FIGS/data/w2b_leg2_metrics_v2.json :: cells.*.N1000000.s1.macro_auc_ovr :: sha256 9cd059d3e2f51798 :: ratio of the seed means of 1 - macro AUC, this row over the reference
\newcommand{\FtOmaRatioJcESixOnesixtwoMassOverOneeighteight}{1.01}  % experiments/FIGS/data/w2b_leg2_metrics_v2.json :: cells.*.N1000000.s1.macro_auc_ovr :: sha256 9cd059d3e2f51798 :: ratio of the seed means of 1 - macro AUC, this row over the reference
\newcommand{\FtOmaRatioJcESixOnesevenMassOverOneeighteight}{1.11}  % experiments/FIGS/data/w2b_leg2_metrics_v2.json :: cells.*.N1000000.s1.macro_auc_ovr :: sha256 9cd059d3e2f51798 :: ratio of the seed means of 1 - macro AUC, this row over the reference
\newcommand{\FtOmaRatioJcESixRandDrawOneOverOneeighteight}{1.05}  % experiments/FIGS/data/w2b_leg2_metrics_v2.json :: cells.*.N1000000.s1.macro_auc_ovr :: sha256 9cd059d3e2f51798 :: ratio of the seed means of 1 - macro AUC, this row over the reference
\newcommand{\FtOmaRatioJcESixRandDrawTwoOverOneeighteight}{1.05}  % experiments/FIGS/data/w2b_leg2_metrics_v2.json :: cells.*.N1000000.s1.macro_auc_ovr :: sha256 9cd059d3e2f51798 :: ratio of the seed means of 1 - macro AUC, this row over the reference
\newcommand{\FtOmaRatioJcESixRandDrawThreeOverOneeighteight}{1.05}  % experiments/FIGS/data/w2b_leg2_metrics_v2.json :: cells.*.N1000000.s1.macro_auc_ovr :: sha256 9cd059d3e2f51798 :: ratio of the seed means of 1 - macro AUC, this row over the reference
\newcommand{\FtOmaRatioJcESixOneeighteightOverOnesixtwo}{1.00}  % experiments/FIGS/data/w2b_leg2_metrics_v2.json :: cells.*.N1000000.s1.macro_auc_ovr :: sha256 9cd059d3e2f51798 :: ratio of the seed means of 1 - macro AUC, this row over the reference
\newcommand{\FtOmaRatioJcESixFourthreeOverOnesixtwo}{1.02}  % experiments/FIGS/data/w2b_leg2_metrics_v2.json :: cells.*.N1000000.s1.macro_auc_ovr :: sha256 9cd059d3e2f51798 :: ratio of the seed means of 1 - macro AUC, this row over the reference
\newcommand{\FtOmaRatioJcESixOnesevenOverOnesixtwo}{1.08}  % experiments/FIGS/data/w2b_leg2_metrics_v2.json :: cells.*.N1000000.s1.macro_auc_ovr :: sha256 9cd059d3e2f51798 :: ratio of the seed means of 1 - macro AUC, this row over the reference
\newcommand{\FtOmaRatioJcESixOnesixtwoMassOverOnesixtwo}{1.01}  % experiments/FIGS/data/w2b_leg2_metrics_v2.json :: cells.*.N1000000.s1.macro_auc_ovr :: sha256 9cd059d3e2f51798 :: ratio of the seed means of 1 - macro AUC, this row over the reference
\newcommand{\FtOmaRatioJcESixOnesevenMassOverOnesixtwo}{1.11}  % experiments/FIGS/data/w2b_leg2_metrics_v2.json :: cells.*.N1000000.s1.macro_auc_ovr :: sha256 9cd059d3e2f51798 :: ratio of the seed means of 1 - macro AUC, this row over the reference
\newcommand{\FtOmaRatioJcESixRandDrawOneOverOnesixtwo}{1.05}  % experiments/FIGS/data/w2b_leg2_metrics_v2.json :: cells.*.N1000000.s1.macro_auc_ovr :: sha256 9cd059d3e2f51798 :: ratio of the seed means of 1 - macro AUC, this row over the reference
\newcommand{\FtOmaRatioJcESixRandDrawTwoOverOnesixtwo}{1.05}  % experiments/FIGS/data/w2b_leg2_metrics_v2.json :: cells.*.N1000000.s1.macro_auc_ovr :: sha256 9cd059d3e2f51798 :: ratio of the seed means of 1 - macro AUC, this row over the reference
\newcommand{\FtOmaRatioJcESixRandDrawThreeOverOnesixtwo}{1.05}  % experiments/FIGS/data/w2b_leg2_metrics_v2.json :: cells.*.N1000000.s1.macro_auc_ovr :: sha256 9cd059d3e2f51798 :: ratio of the seed means of 1 - macro AUC, this row over the reference
\newcommand{\PairedFtAccJcEThreeOnesevenMinusOneeighteight}{\ensuremath{-}0.0040~[\ensuremath{-}0.0144, +0.0064]}  % experiments/FIGS/data/w2b_leg2_metrics_v2.json :: cells.{r16q1-s1,r16q1-s2,r16q1-s3,r16q1-s4,r16q1-s5,l188-s1,l188-s2,l188-s3,l188-s4,l188-s5}.N1000.s1.accuracy :: sha256 9cd059d3e2f51798 :: accuracy, coarse minus fine, paired by run [95% Student-t interval]
\newcommand{\PairedFtAccJcEFourOnesevenMinusOneeighteight}{\ensuremath{-}0.0075~[\ensuremath{-}0.0118, \ensuremath{-}0.0033]}  % experiments/FIGS/data/w2b_leg2_metrics_v2.json :: cells.{r16q1-s1,r16q1-s2,r16q1-s3,r16q1-s4,r16q1-s5,l188-s1,l188-s2,l188-s3,l188-s4,l188-s5}.N10000.s1.accuracy :: sha256 9cd059d3e2f51798 :: accuracy, coarse minus fine, paired by run [95% Student-t interval]
\newcommand{\PairedFtAccJcEFiveOnesevenMinusOneeighteight}{\ensuremath{-}0.0066~[\ensuremath{-}0.0083, \ensuremath{-}0.0049]}  % experiments/FIGS/data/w2b_leg2_metrics_v2.json :: cells.{r16q1-s1,r16q1-s2,r16q1-s3,r16q1-s4,r16q1-s5,l188-s1,l188-s2,l188-s3,l188-s4,l188-s5}.N100000.s1.accuracy :: sha256 9cd059d3e2f51798 :: accuracy, coarse minus fine, paired by run [95% Student-t interval]
\newcommand{\PairedFtAccJcESixOnesevenMinusOneeighteight}{\ensuremath{-}0.0071~[\ensuremath{-}0.0080, \ensuremath{-}0.0062]}  % experiments/FIGS/data/w2b_leg2_metrics_v2.json :: cells.{r16q1-s1,r16q1-s2,r16q1-s3,r16q1-s4,r16q1-s5,l188-s1,l188-s2,l188-s3,l188-s4,l188-s5}.N1000000.s1.accuracy :: sha256 9cd059d3e2f51798 :: accuracy, coarse minus fine, paired by run [95% Student-t interval]
\newcommand{\FtNClassesJc}{10}  % experiments/FIGS/data/w2b_leg2_metrics_v2.json :: cells.*.*.s1.n_classes_present :: sha256 9cd059d3e2f51798 :: classes of the fine-tuning task
\newcommand{\FtNTestAucJc}{500{,}000}  % experiments/FIGS/data/w2b_leg2_metrics_v2.json :: cells.*.*.s1.n_jets_auc :: sha256 9cd059d3e2f51798 :: test jets the macro AUC is computed on
\newcommand{\FtNTestAccJc}{2{,}000{,}000}  % experiments/FIGS/data/w2b_leg2_metrics_v2.json :: cells.*.*.s1.n_jets :: sha256 9cd059d3e2f51798 :: test jets the accuracy is computed on
\newcommand{\FtSizeEThree}{\ensuremath{10^{3}}}  % experiments/FIGS/data/w2b_leg1_metrics_v2.json :: cells.*.N1000 :: sha256 b56bf70be9982044 :: fine-tuning training jets
\newcommand{\FtSizeEFour}{\ensuremath{10^{4}}}  % experiments/FIGS/data/w2b_leg1_metrics_v2.json :: cells.*.N10000 :: sha256 b56bf70be9982044 :: fine-tuning training jets
\newcommand{\FtSizeEFive}{\ensuremath{10^{5}}}  % experiments/FIGS/data/w2b_leg1_metrics_v2.json :: cells.*.N100000 :: sha256 b56bf70be9982044 :: fine-tuning training jets
\newcommand{\FtSizeESix}{\ensuremath{10^{6}}}  % experiments/FIGS/data/w2b_leg1_metrics_v2.json :: cells.*.N1000000 :: sha256 b56bf70be9982044 :: fine-tuning training jets
\newcommand{\FtNRunsWord}{five}  % experiments/FIGS/data/w2b_leg1_metrics_v2.json :: cells (pretrained models per vocabulary row) :: sha256 b56bf70be9982044 :: pretraining runs per vocabulary
\newcommand{\PairedProbeBcVsRestLinearOnesevenOverOnesixtwo}{1.64~[1.52, 1.78]}  % experiments/FIGS/data/paired_v1err/probes/ratios.json :: ratios[0].ratio, ratios[0].ci95 :: sha256 14dd51402305d147 :: 17 over 162, bc_vs_rest linear 1-auc: paired geometric mean [95% interval]
\newcommand{\PairedProbeBcVsRestMlpOnesevenOverOnesixtwo}{1.50~[1.39, 1.62]}  % experiments/FIGS/data/paired_v1err/probes/ratios.json :: ratios[3].ratio, ratios[3].ci95 :: sha256 14dd51402305d147 :: 17 over 162, bc_vs_rest mlp 1-auc: paired geometric mean [95% interval]
\newcommand{\PairedProbeBvcFourprongLinearOnesixtwoOverOneeighteight}{1.01~[0.95, 1.08]}  % experiments/FIGS/data/paired_v1err/probes/ratios.json :: ratios[6].ratio, ratios[6].ci95 :: sha256 14dd51402305d147 :: 162 over 188, bvc_4prong linear 1-auc: paired geometric mean [95% interval]
\newcommand{\PairedProbeBvcFourprongLinearFourthreeOverOneeighteight}{1.10~[1.01, 1.20]}  % experiments/FIGS/data/paired_v1err/probes/ratios.json :: ratios[7].ratio, ratios[7].ci95 :: sha256 14dd51402305d147 :: 43 over 188, bvc_4prong linear 1-auc: paired geometric mean [95% interval]
\newcommand{\PairedProbeBvcFourprongLinearOnesevenOverOneeighteight}{4.5~[3.5, 5.7]}  % experiments/FIGS/data/paired_v1err/probes/ratios.json :: ratios[8].ratio, ratios[8].ci95 :: sha256 14dd51402305d147 :: 17 over 188, bvc_4prong linear 1-auc: paired geometric mean [95% interval]
\newcommand{\PairedProbeBvcFourprongLinearFourthreeOverOnesixtwo}{1.09~[0.98, 1.20]}  % experiments/FIGS/data/paired_v1err/probes/ratios.json :: ratios[9].ratio, ratios[9].ci95 :: sha256 14dd51402305d147 :: 43 over 162, bvc_4prong linear 1-auc: paired geometric mean [95% interval]
\newcommand{\PairedProbeBvcFourprongLinearOnesevenOverOnesixtwo}{4.4~[3.4, 5.6]}  % experiments/FIGS/data/paired_v1err/probes/ratios.json :: ratios[10].ratio, ratios[10].ci95 :: sha256 14dd51402305d147 :: 17 over 162, bvc_4prong linear 1-auc: paired geometric mean [95% interval]
\newcommand{\PairedProbeBvcFourprongLinearOnesevenOverFourthree}{4.0~[3.1, 5.4]}  % experiments/FIGS/data/paired_v1err/probes/ratios.json :: ratios[11].ratio, ratios[11].ci95 :: sha256 14dd51402305d147 :: 17 over 43, bvc_4prong linear 1-auc: paired geometric mean [95% interval]
\newcommand{\PairedProbeBvcFourprongLinearOnesixtwoMassOverOnesixtwo}{1.05~[0.97, 1.13]}  % experiments/FIGS/data/paired_v1err/probes/ratios.json :: ratios[12].ratio, ratios[12].ci95 :: sha256 14dd51402305d147 :: 162+mass over 162, bvc_4prong linear 1-auc: paired geometric mean [95% interval]
\newcommand{\PairedProbeBvcFourprongLinearOnesevenMassOverOneseven}{2.1~[1.6, 2.7]}  % experiments/FIGS/data/paired_v1err/probes/ratios.json :: ratios[13].ratio, ratios[13].ci95 :: sha256 14dd51402305d147 :: 17+mass over 17, bvc_4prong linear 1-auc: paired geometric mean [95% interval]
\newcommand{\PairedProbeBvcFourprongLinearRandomOverOneeighteight}{1.54~[1.37, 1.74]}  % experiments/FIGS/data/paired_v1err/probes/ratios.json :: ratios[14].ratio, ratios[14].ci95 :: sha256 14dd51402305d147 :: random over 188, bvc_4prong linear 1-auc: paired geometric mean [95% interval]
\newcommand{\PairedProbeBvcFourprongLinearRandomOverOnesixtwo}{1.55~[1.36, 1.77]}  % experiments/FIGS/data/paired_v1err/probes/ratios.json :: ratios[15].ratio, ratios[15].ci95 :: sha256 14dd51402305d147 :: random over 162, bvc_4prong linear 1-auc: paired geometric mean [95% interval]
\newcommand{\PairedProbeBvcFourprongLinearRandomOverFourthree}{1.39~[1.27, 1.52]}  % experiments/FIGS/data/paired_v1err/probes/ratios.json :: ratios[16].ratio, ratios[16].ci95 :: sha256 14dd51402305d147 :: random over 43, bvc_4prong linear 1-auc: paired geometric mean [95% interval]
\newcommand{\PairedProbeBvcFourprongLinearRandomOverOneseven}{0.34~[0.20, 0.58]}  % experiments/FIGS/data/paired_v1err/probes/ratios.json :: ratios[17].ratio, ratios[17].ci95 :: sha256 14dd51402305d147 :: random over 17, bvc_4prong linear 1-auc: paired geometric mean [95% interval]
\newcommand{\PairedProbeBvcFourprongLinearRandomDrawTwoOverRandomDrawOne}{1.05~[0.94, 1.18]}  % experiments/FIGS/data/paired_v1err/probes/ratios.json :: ratios[18].ratio, ratios[18].ci95 :: sha256 14dd51402305d147 :: random draw 2 over random draw 1, bvc_4prong linear 1-auc: paired geometric mean [95% interval]
\newcommand{\PairedProbeBvcFourprongLinearRandomDrawThreeOverRandomDrawOne}{0.95~[0.85, 1.06]}  % experiments/FIGS/data/paired_v1err/probes/ratios.json :: ratios[19].ratio, ratios[19].ci95 :: sha256 14dd51402305d147 :: random draw 3 over random draw 1, bvc_4prong linear 1-auc: paired geometric mean [95% interval]
\newcommand{\PairedProbeBvcFourprongLinearRandomDrawThreeOverRandomDrawTwo}{0.90~[0.79, 1.02]}  % experiments/FIGS/data/paired_v1err/probes/ratios.json :: ratios[20].ratio, ratios[20].ci95 :: sha256 14dd51402305d147 :: random draw 3 over random draw 2, bvc_4prong linear 1-auc: paired geometric mean [95% interval]
\newcommand{\PairedProbeBvcFourprongMlpOnesixtwoOverOneeighteight}{1.00~[0.93, 1.07]}  % experiments/FIGS/data/paired_v1err/probes/ratios.json :: ratios[66].ratio, ratios[66].ci95 :: sha256 14dd51402305d147 :: 162 over 188, bvc_4prong mlp 1-auc: paired geometric mean [95% interval]
\newcommand{\PairedProbeBvcFourprongMlpFourthreeOverOneeighteight}{1.12~[1.03, 1.22]}  % experiments/FIGS/data/paired_v1err/probes/ratios.json :: ratios[67].ratio, ratios[67].ci95 :: sha256 14dd51402305d147 :: 43 over 188, bvc_4prong mlp 1-auc: paired geometric mean [95% interval]
\newcommand{\PairedProbeBvcFourprongMlpOnesevenOverOneeighteight}{3.2~[2.6, 4.1]}  % experiments/FIGS/data/paired_v1err/probes/ratios.json :: ratios[68].ratio, ratios[68].ci95 :: sha256 14dd51402305d147 :: 17 over 188, bvc_4prong mlp 1-auc: paired geometric mean [95% interval]
\newcommand{\PairedProbeBvcFourprongMlpFourthreeOverOnesixtwo}{1.12~[1.01, 1.26]}  % experiments/FIGS/data/paired_v1err/probes/ratios.json :: ratios[69].ratio, ratios[69].ci95 :: sha256 14dd51402305d147 :: 43 over 162, bvc_4prong mlp 1-auc: paired geometric mean [95% interval]
\newcommand{\PairedProbeBvcFourprongMlpOnesevenOverOnesixtwo}{3.3~[2.6, 4.0]}  % experiments/FIGS/data/paired_v1err/probes/ratios.json :: ratios[70].ratio, ratios[70].ci95 :: sha256 14dd51402305d147 :: 17 over 162, bvc_4prong mlp 1-auc: paired geometric mean [95% interval]
\newcommand{\PairedProbeBvcFourprongMlpOnesevenOverFourthree}{2.9~[2.3, 3.7]}  % experiments/FIGS/data/paired_v1err/probes/ratios.json :: ratios[71].ratio, ratios[71].ci95 :: sha256 14dd51402305d147 :: 17 over 43, bvc_4prong mlp 1-auc: paired geometric mean [95% interval]
\newcommand{\PairedProbeBvcFourprongMlpOnesixtwoMassOverOnesixtwo}{1.07~[0.99, 1.16]}  % experiments/FIGS/data/paired_v1err/probes/ratios.json :: ratios[72].ratio, ratios[72].ci95 :: sha256 14dd51402305d147 :: 162+mass over 162, bvc_4prong mlp 1-auc: paired geometric mean [95% interval]
\newcommand{\PairedProbeBvcFourprongMlpOnesevenMassOverOneseven}{1.97~[1.46, 2.67]}  % experiments/FIGS/data/paired_v1err/probes/ratios.json :: ratios[73].ratio, ratios[73].ci95 :: sha256 14dd51402305d147 :: 17+mass over 17, bvc_4prong mlp 1-auc: paired geometric mean [95% interval]
\newcommand{\PairedProbeBvcFourprongMlpRandomOverOneeighteight}{1.47~[1.30, 1.65]}  % experiments/FIGS/data/paired_v1err/probes/ratios.json :: ratios[74].ratio, ratios[74].ci95 :: sha256 14dd51402305d147 :: random over 188, bvc_4prong mlp 1-auc: paired geometric mean [95% interval]
\newcommand{\PairedProbeBvcFourprongMlpRandomOverOnesixtwo}{1.49~[1.24, 1.79]}  % experiments/FIGS/data/paired_v1err/probes/ratios.json :: ratios[75].ratio, ratios[75].ci95 :: sha256 14dd51402305d147 :: random over 162, bvc_4prong mlp 1-auc: paired geometric mean [95% interval]
\newcommand{\PairedProbeBvcFourprongMlpRandomOverFourthree}{1.28~[1.16, 1.43]}  % experiments/FIGS/data/paired_v1err/probes/ratios.json :: ratios[76].ratio, ratios[76].ci95 :: sha256 14dd51402305d147 :: random over 43, bvc_4prong mlp 1-auc: paired geometric mean [95% interval]
\newcommand{\PairedProbeBvcFourprongMlpRandomOverOneseven}{0.43~[0.28, 0.67]}  % experiments/FIGS/data/paired_v1err/probes/ratios.json :: ratios[77].ratio, ratios[77].ci95 :: sha256 14dd51402305d147 :: random over 17, bvc_4prong mlp 1-auc: paired geometric mean [95% interval]
\newcommand{\PairedProbeBvcFourprongMlpRandomDrawTwoOverRandomDrawOne}{1.00~[0.90, 1.10]}  % experiments/FIGS/data/paired_v1err/probes/ratios.json :: ratios[78].ratio, ratios[78].ci95 :: sha256 14dd51402305d147 :: random draw 2 over random draw 1, bvc_4prong mlp 1-auc: paired geometric mean [95% interval]
\newcommand{\PairedProbeBvcFourprongMlpRandomDrawThreeOverRandomDrawOne}{0.89~[0.79, 0.99]}  % experiments/FIGS/data/paired_v1err/probes/ratios.json :: ratios[79].ratio, ratios[79].ci95 :: sha256 14dd51402305d147 :: random draw 3 over random draw 1, bvc_4prong mlp 1-auc: paired geometric mean [95% interval]
\newcommand{\PairedProbeBvcFourprongMlpRandomDrawThreeOverRandomDrawTwo}{0.89~[0.79, 0.99]}  % experiments/FIGS/data/paired_v1err/probes/ratios.json :: ratios[80].ratio, ratios[80].ci95 :: sha256 14dd51402305d147 :: random draw 3 over random draw 2, bvc_4prong mlp 1-auc: paired geometric mean [95% interval]
\newcommand{\PairedProbeBvcQcdLinearOnesixtwoOverOneeighteight}{1.15~[1.07, 1.23]}  % experiments/FIGS/data/paired_v1err/probes/ratios.json :: ratios[126].ratio, ratios[126].ci95 :: sha256 14dd51402305d147 :: 162 over 188, bvc_qcd linear 1-auc: paired geometric mean [95% interval]
\newcommand{\PairedProbeBvcQcdLinearFourthreeOverOneeighteight}{1.16~[1.08, 1.25]}  % experiments/FIGS/data/paired_v1err/probes/ratios.json :: ratios[127].ratio, ratios[127].ci95 :: sha256 14dd51402305d147 :: 43 over 188, bvc_qcd linear 1-auc: paired geometric mean [95% interval]
\newcommand{\PairedProbeBvcQcdLinearOnesevenOverOneeighteight}{3.0~[2.5, 3.7]}  % experiments/FIGS/data/paired_v1err/probes/ratios.json :: ratios[128].ratio, ratios[128].ci95 :: sha256 14dd51402305d147 :: 17 over 188, bvc_qcd linear 1-auc: paired geometric mean [95% interval]
\newcommand{\PairedProbeBvcQcdLinearFourthreeOverOnesixtwo}{1.01~[0.96, 1.07]}  % experiments/FIGS/data/paired_v1err/probes/ratios.json :: ratios[129].ratio, ratios[129].ci95 :: sha256 14dd51402305d147 :: 43 over 162, bvc_qcd linear 1-auc: paired geometric mean [95% interval]
\newcommand{\PairedProbeBvcQcdLinearOnesevenOverOnesixtwo}{2.6~[2.2, 3.1]}  % experiments/FIGS/data/paired_v1err/probes/ratios.json :: ratios[130].ratio, ratios[130].ci95 :: sha256 14dd51402305d147 :: 17 over 162, bvc_qcd linear 1-auc: paired geometric mean [95% interval]
\newcommand{\PairedProbeBvcQcdLinearOnesevenOverFourthree}{2.6~[2.2, 3.0]}  % experiments/FIGS/data/paired_v1err/probes/ratios.json :: ratios[131].ratio, ratios[131].ci95 :: sha256 14dd51402305d147 :: 17 over 43, bvc_qcd linear 1-auc: paired geometric mean [95% interval]
\newcommand{\PairedProbeBvcQcdLinearOnesixtwoMassOverOnesixtwo}{1.060~[1.002, 1.123]}  % experiments/FIGS/data/paired_v1err/probes/ratios.json :: ratios[132].ratio, ratios[132].ci95 :: sha256 14dd51402305d147 :: 162+mass over 162, bvc_qcd linear 1-auc: paired geometric mean [95% interval]
\newcommand{\PairedProbeBvcQcdLinearOnesevenMassOverOneseven}{1.89~[1.52, 2.34]}  % experiments/FIGS/data/paired_v1err/probes/ratios.json :: ratios[133].ratio, ratios[133].ci95 :: sha256 14dd51402305d147 :: 17+mass over 17, bvc_qcd linear 1-auc: paired geometric mean [95% interval]
\newcommand{\PairedProbeBvcQcdMlpOnesixtwoOverOneeighteight}{1.11~[1.04, 1.19]}  % experiments/FIGS/data/paired_v1err/probes/ratios.json :: ratios[158].ratio, ratios[158].ci95 :: sha256 14dd51402305d147 :: 162 over 188, bvc_qcd mlp 1-auc: paired geometric mean [95% interval]
\newcommand{\PairedProbeBvcQcdMlpFourthreeOverOneeighteight}{1.13~[1.05, 1.21]}  % experiments/FIGS/data/paired_v1err/probes/ratios.json :: ratios[159].ratio, ratios[159].ci95 :: sha256 14dd51402305d147 :: 43 over 188, bvc_qcd mlp 1-auc: paired geometric mean [95% interval]
\newcommand{\PairedProbeBvcQcdMlpOnesevenOverOneeighteight}{2.6~[2.2, 3.0]}  % experiments/FIGS/data/paired_v1err/probes/ratios.json :: ratios[160].ratio, ratios[160].ci95 :: sha256 14dd51402305d147 :: 17 over 188, bvc_qcd mlp 1-auc: paired geometric mean [95% interval]
\newcommand{\PairedProbeBvcQcdMlpFourthreeOverOnesixtwo}{1.01~[0.96, 1.06]}  % experiments/FIGS/data/paired_v1err/probes/ratios.json :: ratios[161].ratio, ratios[161].ci95 :: sha256 14dd51402305d147 :: 43 over 162, bvc_qcd mlp 1-auc: paired geometric mean [95% interval]
\newcommand{\PairedProbeBvcQcdMlpOnesevenOverOnesixtwo}{2.3~[2.0, 2.6]}  % experiments/FIGS/data/paired_v1err/probes/ratios.json :: ratios[162].ratio, ratios[162].ci95 :: sha256 14dd51402305d147 :: 17 over 162, bvc_qcd mlp 1-auc: paired geometric mean [95% interval]
\newcommand{\PairedProbeBvcQcdMlpOnesevenOverFourthree}{2.3~[2.0, 2.6]}  % experiments/FIGS/data/paired_v1err/probes/ratios.json :: ratios[163].ratio, ratios[163].ci95 :: sha256 14dd51402305d147 :: 17 over 43, bvc_qcd mlp 1-auc: paired geometric mean [95% interval]
\newcommand{\PairedProbeBvcQcdMlpOnesixtwoMassOverOnesixtwo}{1.06~[1.01, 1.12]}  % experiments/FIGS/data/paired_v1err/probes/ratios.json :: ratios[164].ratio, ratios[164].ci95 :: sha256 14dd51402305d147 :: 162+mass over 162, bvc_qcd mlp 1-auc: paired geometric mean [95% interval]
\newcommand{\PairedProbeBvcQcdMlpOnesevenMassOverOneseven}{1.64~[1.33, 2.02]}  % experiments/FIGS/data/paired_v1err/probes/ratios.json :: ratios[165].ratio, ratios[165].ci95 :: sha256 14dd51402305d147 :: 17+mass over 17, bvc_qcd mlp 1-auc: paired geometric mean [95% interval]
\newcommand{\PairedProbeBvcResonantLinearOnesixtwoOverOneeighteight}{1.02~[0.96, 1.08]}  % experiments/FIGS/data/paired_v1err/probes/ratios.json :: ratios[190].ratio, ratios[190].ci95 :: sha256 14dd51402305d147 :: 162 over 188, bvc_resonant linear 1-auc: paired geometric mean [95% interval]
\newcommand{\PairedProbeBvcResonantLinearFourthreeOverOneeighteight}{1.09~[1.03, 1.15]}  % experiments/FIGS/data/paired_v1err/probes/ratios.json :: ratios[191].ratio, ratios[191].ci95 :: sha256 14dd51402305d147 :: 43 over 188, bvc_resonant linear 1-auc: paired geometric mean [95% interval]
\newcommand{\PairedProbeBvcResonantLinearOnesevenOverOneeighteight}{4.0~[3.3, 4.9]}  % experiments/FIGS/data/paired_v1err/probes/ratios.json :: ratios[192].ratio, ratios[192].ci95 :: sha256 14dd51402305d147 :: 17 over 188, bvc_resonant linear 1-auc: paired geometric mean [95% interval]
\newcommand{\PairedProbeBvcResonantLinearFourthreeOverOnesixtwo}{1.07~[1.02, 1.12]}  % experiments/FIGS/data/paired_v1err/probes/ratios.json :: ratios[193].ratio, ratios[193].ci95 :: sha256 14dd51402305d147 :: 43 over 162, bvc_resonant linear 1-auc: paired geometric mean [95% interval]
\newcommand{\PairedProbeBvcResonantLinearOnesevenOverOnesixtwo}{3.9~[3.1, 4.9]}  % experiments/FIGS/data/paired_v1err/probes/ratios.json :: ratios[194].ratio, ratios[194].ci95 :: sha256 14dd51402305d147 :: 17 over 162, bvc_resonant linear 1-auc: paired geometric mean [95% interval]
\newcommand{\PairedProbeBvcResonantLinearOnesevenOverFourthree}{3.7~[3.0, 4.5]}  % experiments/FIGS/data/paired_v1err/probes/ratios.json :: ratios[195].ratio, ratios[195].ci95 :: sha256 14dd51402305d147 :: 17 over 43, bvc_resonant linear 1-auc: paired geometric mean [95% interval]
\newcommand{\PairedProbeBvcResonantLinearOnesixtwoMassOverOnesixtwo}{1.11~[1.06, 1.16]}  % experiments/FIGS/data/paired_v1err/probes/ratios.json :: ratios[196].ratio, ratios[196].ci95 :: sha256 14dd51402305d147 :: 162+mass over 162, bvc_resonant linear 1-auc: paired geometric mean [95% interval]
\newcommand{\PairedProbeBvcResonantLinearOnesevenMassOverOneseven}{3.0~[1.9, 4.7]}  % experiments/FIGS/data/paired_v1err/probes/ratios.json :: ratios[197].ratio, ratios[197].ci95 :: sha256 14dd51402305d147 :: 17+mass over 17, bvc_resonant linear 1-auc: paired geometric mean [95% interval]
\newcommand{\PairedProbeBvcResonantMlpOnesixtwoOverOneeighteight}{1.04~[0.98, 1.12]}  % experiments/FIGS/data/paired_v1err/probes/ratios.json :: ratios[222].ratio, ratios[222].ci95 :: sha256 14dd51402305d147 :: 162 over 188, bvc_resonant mlp 1-auc: paired geometric mean [95% interval]
\newcommand{\PairedProbeBvcResonantMlpFourthreeOverOneeighteight}{1.14~[1.08, 1.21]}  % experiments/FIGS/data/paired_v1err/probes/ratios.json :: ratios[223].ratio, ratios[223].ci95 :: sha256 14dd51402305d147 :: 43 over 188, bvc_resonant mlp 1-auc: paired geometric mean [95% interval]
\newcommand{\PairedProbeBvcResonantMlpOnesevenOverOneeighteight}{3.3~[2.8, 3.8]}  % experiments/FIGS/data/paired_v1err/probes/ratios.json :: ratios[224].ratio, ratios[224].ci95 :: sha256 14dd51402305d147 :: 17 over 188, bvc_resonant mlp 1-auc: paired geometric mean [95% interval]
\newcommand{\PairedProbeBvcResonantMlpFourthreeOverOnesixtwo}{1.09~[1.04, 1.14]}  % experiments/FIGS/data/paired_v1err/probes/ratios.json :: ratios[225].ratio, ratios[225].ci95 :: sha256 14dd51402305d147 :: 43 over 162, bvc_resonant mlp 1-auc: paired geometric mean [95% interval]
\newcommand{\PairedProbeBvcResonantMlpOnesevenOverOnesixtwo}{3.1~[2.6, 3.8]}  % experiments/FIGS/data/paired_v1err/probes/ratios.json :: ratios[226].ratio, ratios[226].ci95 :: sha256 14dd51402305d147 :: 17 over 162, bvc_resonant mlp 1-auc: paired geometric mean [95% interval]
\newcommand{\PairedProbeBvcResonantMlpOnesevenOverFourthree}{2.9~[2.4, 3.4]}  % experiments/FIGS/data/paired_v1err/probes/ratios.json :: ratios[227].ratio, ratios[227].ci95 :: sha256 14dd51402305d147 :: 17 over 43, bvc_resonant mlp 1-auc: paired geometric mean [95% interval]
\newcommand{\PairedProbeBvcResonantMlpOnesixtwoMassOverOnesixtwo}{1.09~[1.04, 1.14]}  % experiments/FIGS/data/paired_v1err/probes/ratios.json :: ratios[228].ratio, ratios[228].ci95 :: sha256 14dd51402305d147 :: 162+mass over 162, bvc_resonant mlp 1-auc: paired geometric mean [95% interval]
\newcommand{\PairedProbeBvcResonantMlpOnesevenMassOverOneseven}{1.95~[1.45, 2.61]}  % experiments/FIGS/data/paired_v1err/probes/ratios.json :: ratios[229].ratio, ratios[229].ci95 :: sha256 14dd51402305d147 :: 17+mass over 17, bvc_resonant mlp 1-auc: paired geometric mean [95% interval]
\newcommand{\PairedProbeEeVsMmLinearOnesevenOverFourthree}{2.3~[0.8, 6.5]}  % experiments/FIGS/data/paired_v1err/probes/ratios.json :: ratios[259].ratio, ratios[259].ci95 :: sha256 14dd51402305d147 :: 17 over 43, ee_vs_mm linear 1-auc: paired geometric mean [95% interval]
\newcommand{\PairedProbeEeVsMmLinearOnesevenMassOverOneseven}{0.23~[0.08, 0.64]}  % experiments/FIGS/data/paired_v1err/probes/ratios.json :: ratios[261].ratio, ratios[261].ci95 :: sha256 14dd51402305d147 :: 17+mass over 17, ee_vs_mm linear 1-auc: paired geometric mean [95% interval]
\newcommand{\PairedProbeEeVsMmMlpOnesevenOverFourthree}{2.8~[0.8, 9.0]}  % experiments/FIGS/data/paired_v1err/probes/ratios.json :: ratios[291].ratio, ratios[291].ci95 :: sha256 14dd51402305d147 :: 17 over 43, ee_vs_mm mlp 1-auc: paired geometric mean [95% interval]
\newcommand{\PairedProbeEeVsMmMlpOnesevenMassOverOneseven}{0.172~[0.026, 1.125]}  % experiments/FIGS/data/paired_v1err/probes/ratios.json :: ratios[293].ratio, ratios[293].ci95 :: sha256 14dd51402305d147 :: 17+mass over 17, ee_vs_mm mlp 1-auc: paired geometric mean [95% interval]
\newcommand{\PairedProbeRetainedTopologyLinearOnesixtwoOverOneeighteight}{1.00~[0.92, 1.08]}  % experiments/FIGS/data/paired_v1err/probes/ratios.json :: ratios[318].ratio, ratios[318].ci95 :: sha256 14dd51402305d147 :: 162 over 188, retained_topology linear 1-auc: paired geometric mean [95% interval]
\newcommand{\PairedProbeRetainedTopologyLinearFourthreeOverOneeighteight}{1.24~[1.15, 1.34]}  % experiments/FIGS/data/paired_v1err/probes/ratios.json :: ratios[319].ratio, ratios[319].ci95 :: sha256 14dd51402305d147 :: 43 over 188, retained_topology linear 1-auc: paired geometric mean [95% interval]
\newcommand{\PairedProbeRetainedTopologyLinearOnesevenOverOneeighteight}{1.54~[1.35, 1.75]}  % experiments/FIGS/data/paired_v1err/probes/ratios.json :: ratios[320].ratio, ratios[320].ci95 :: sha256 14dd51402305d147 :: 17 over 188, retained_topology linear 1-auc: paired geometric mean [95% interval]
\newcommand{\PairedProbeRetainedTopologyLinearFourthreeOverOnesixtwo}{1.24~[1.14, 1.34]}  % experiments/FIGS/data/paired_v1err/probes/ratios.json :: ratios[321].ratio, ratios[321].ci95 :: sha256 14dd51402305d147 :: 43 over 162, retained_topology linear 1-auc: paired geometric mean [95% interval]
\newcommand{\PairedProbeRetainedTopologyLinearOnesevenOverOnesixtwo}{1.54~[1.33, 1.78]}  % experiments/FIGS/data/paired_v1err/probes/ratios.json :: ratios[322].ratio, ratios[322].ci95 :: sha256 14dd51402305d147 :: 17 over 162, retained_topology linear 1-auc: paired geometric mean [95% interval]
\newcommand{\PairedProbeRetainedTopologyLinearOnesevenOverFourthree}{1.24~[1.08, 1.42]}  % experiments/FIGS/data/paired_v1err/probes/ratios.json :: ratios[323].ratio, ratios[323].ci95 :: sha256 14dd51402305d147 :: 17 over 43, retained_topology linear 1-auc: paired geometric mean [95% interval]
\newcommand{\PairedProbeRetainedTopologyLinearOnesixtwoMassOverOnesixtwo}{1.06~[0.99, 1.14]}  % experiments/FIGS/data/paired_v1err/probes/ratios.json :: ratios[324].ratio, ratios[324].ci95 :: sha256 14dd51402305d147 :: 162+mass over 162, retained_topology linear 1-auc: paired geometric mean [95% interval]
\newcommand{\PairedProbeRetainedTopologyLinearOnesevenMassOverOneseven}{1.12~[0.97, 1.28]}  % experiments/FIGS/data/paired_v1err/probes/ratios.json :: ratios[325].ratio, ratios[325].ci95 :: sha256 14dd51402305d147 :: 17+mass over 17, retained_topology linear 1-auc: paired geometric mean [95% interval]
\newcommand{\PairedProbeRetainedTopologyMlpOnesixtwoOverOneeighteight}{1.01~[0.87, 1.17]}  % experiments/FIGS/data/paired_v1err/probes/ratios.json :: ratios[350].ratio, ratios[350].ci95 :: sha256 14dd51402305d147 :: 162 over 188, retained_topology mlp 1-auc: paired geometric mean [95% interval]
\newcommand{\PairedProbeRetainedTopologyMlpFourthreeOverOneeighteight}{1.18~[1.03, 1.36]}  % experiments/FIGS/data/paired_v1err/probes/ratios.json :: ratios[351].ratio, ratios[351].ci95 :: sha256 14dd51402305d147 :: 43 over 188, retained_topology mlp 1-auc: paired geometric mean [95% interval]
\newcommand{\PairedProbeRetainedTopologyMlpOnesevenOverOneeighteight}{1.52~[1.34, 1.72]}  % experiments/FIGS/data/paired_v1err/probes/ratios.json :: ratios[352].ratio, ratios[352].ci95 :: sha256 14dd51402305d147 :: 17 over 188, retained_topology mlp 1-auc: paired geometric mean [95% interval]
\newcommand{\PairedProbeRetainedTopologyMlpFourthreeOverOnesixtwo}{1.17~[1.07, 1.28]}  % experiments/FIGS/data/paired_v1err/probes/ratios.json :: ratios[353].ratio, ratios[353].ci95 :: sha256 14dd51402305d147 :: 43 over 162, retained_topology mlp 1-auc: paired geometric mean [95% interval]
\newcommand{\PairedProbeRetainedTopologyMlpOnesevenOverOnesixtwo}{1.50~[1.29, 1.75]}  % experiments/FIGS/data/paired_v1err/probes/ratios.json :: ratios[354].ratio, ratios[354].ci95 :: sha256 14dd51402305d147 :: 17 over 162, retained_topology mlp 1-auc: paired geometric mean [95% interval]
\newcommand{\PairedProbeRetainedTopologyMlpOnesevenOverFourthree}{1.28~[1.12, 1.47]}  % experiments/FIGS/data/paired_v1err/probes/ratios.json :: ratios[355].ratio, ratios[355].ci95 :: sha256 14dd51402305d147 :: 17 over 43, retained_topology mlp 1-auc: paired geometric mean [95% interval]
\newcommand{\PairedProbeRetainedTopologyMlpOnesixtwoMassOverOnesixtwo}{1.04~[0.92, 1.17]}  % experiments/FIGS/data/paired_v1err/probes/ratios.json :: ratios[356].ratio, ratios[356].ci95 :: sha256 14dd51402305d147 :: 162+mass over 162, retained_topology mlp 1-auc: paired geometric mean [95% interval]
\newcommand{\PairedProbeRetainedTopologyMlpOnesevenMassOverOneseven}{1.11~[0.99, 1.23]}  % experiments/FIGS/data/paired_v1err/probes/ratios.json :: ratios[357].ratio, ratios[357].ci95 :: sha256 14dd51402305d147 :: 17+mass over 17, retained_topology mlp 1-auc: paired geometric mean [95% interval]
\newcommand{\PairedProbeVisibleContentLinearOnesixtwoOverOneeighteight}{0.95~[0.88, 1.03]}  % experiments/FIGS/data/paired_v1err/probes/ratios.json :: ratios[382].ratio, ratios[382].ci95 :: sha256 14dd51402305d147 :: 162 over 188, visible_content linear 1-auc: paired geometric mean [95% interval]
\newcommand{\PairedProbeVisibleContentLinearFourthreeOverOneeighteight}{0.99~[0.85, 1.15]}  % experiments/FIGS/data/paired_v1err/probes/ratios.json :: ratios[383].ratio, ratios[383].ci95 :: sha256 14dd51402305d147 :: 43 over 188, visible_content linear 1-auc: paired geometric mean [95% interval]
\newcommand{\PairedProbeVisibleContentLinearOnesevenOverOneeighteight}{2.8~[2.1, 3.7]}  % experiments/FIGS/data/paired_v1err/probes/ratios.json :: ratios[384].ratio, ratios[384].ci95 :: sha256 14dd51402305d147 :: 17 over 188, visible_content linear 1-auc: paired geometric mean [95% interval]
\newcommand{\PairedProbeVisibleContentLinearFourthreeOverOnesixtwo}{1.04~[0.89, 1.20]}  % experiments/FIGS/data/paired_v1err/probes/ratios.json :: ratios[385].ratio, ratios[385].ci95 :: sha256 14dd51402305d147 :: 43 over 162, visible_content linear 1-auc: paired geometric mean [95% interval]
\newcommand{\PairedProbeVisibleContentLinearOnesevenOverOnesixtwo}{2.9~[2.3, 3.8]}  % experiments/FIGS/data/paired_v1err/probes/ratios.json :: ratios[386].ratio, ratios[386].ci95 :: sha256 14dd51402305d147 :: 17 over 162, visible_content linear 1-auc: paired geometric mean [95% interval]
\newcommand{\PairedProbeVisibleContentLinearOnesevenOverFourthree}{2.8~[2.2, 3.7]}  % experiments/FIGS/data/paired_v1err/probes/ratios.json :: ratios[387].ratio, ratios[387].ci95 :: sha256 14dd51402305d147 :: 17 over 43, visible_content linear 1-auc: paired geometric mean [95% interval]
\newcommand{\PairedProbeVisibleContentLinearOnesixtwoMassOverOnesixtwo}{1.10~[0.98, 1.23]}  % experiments/FIGS/data/paired_v1err/probes/ratios.json :: ratios[388].ratio, ratios[388].ci95 :: sha256 14dd51402305d147 :: 162+mass over 162, visible_content linear 1-auc: paired geometric mean [95% interval]
\newcommand{\PairedProbeVisibleContentLinearOnesevenMassOverOneseven}{1.44~[1.26, 1.65]}  % experiments/FIGS/data/paired_v1err/probes/ratios.json :: ratios[389].ratio, ratios[389].ci95 :: sha256 14dd51402305d147 :: 17+mass over 17, visible_content linear 1-auc: paired geometric mean [95% interval]
\newcommand{\PairedProbeVisibleContentLinearRandomOverOneeighteight}{1.31~[1.01, 1.70]}  % experiments/FIGS/data/paired_v1err/probes/ratios.json :: ratios[390].ratio, ratios[390].ci95 :: sha256 14dd51402305d147 :: random over 188, visible_content linear 1-auc: paired geometric mean [95% interval]
\newcommand{\PairedProbeVisibleContentLinearRandomOverOnesixtwo}{1.41~[1.17, 1.70]}  % experiments/FIGS/data/paired_v1err/probes/ratios.json :: ratios[391].ratio, ratios[391].ci95 :: sha256 14dd51402305d147 :: random over 162, visible_content linear 1-auc: paired geometric mean [95% interval]
\newcommand{\PairedProbeVisibleContentLinearRandomOverFourthree}{1.33~[1.12, 1.57]}  % experiments/FIGS/data/paired_v1err/probes/ratios.json :: ratios[392].ratio, ratios[392].ci95 :: sha256 14dd51402305d147 :: random over 43, visible_content linear 1-auc: paired geometric mean [95% interval]
\newcommand{\PairedProbeVisibleContentLinearRandomOverOneseven}{0.48~[0.38, 0.60]}  % experiments/FIGS/data/paired_v1err/probes/ratios.json :: ratios[393].ratio, ratios[393].ci95 :: sha256 14dd51402305d147 :: random over 17, visible_content linear 1-auc: paired geometric mean [95% interval]
\newcommand{\PairedProbeVisibleContentLinearRandomDrawTwoOverRandomDrawOne}{0.79~[0.64, 0.98]}  % experiments/FIGS/data/paired_v1err/probes/ratios.json :: ratios[394].ratio, ratios[394].ci95 :: sha256 14dd51402305d147 :: random draw 2 over random draw 1, visible_content linear 1-auc: paired geometric mean [95% interval]
\newcommand{\PairedProbeVisibleContentLinearRandomDrawThreeOverRandomDrawOne}{0.86~[0.71, 1.05]}  % experiments/FIGS/data/paired_v1err/probes/ratios.json :: ratios[395].ratio, ratios[395].ci95 :: sha256 14dd51402305d147 :: random draw 3 over random draw 1, visible_content linear 1-auc: paired geometric mean [95% interval]
\newcommand{\PairedProbeVisibleContentLinearRandomDrawThreeOverRandomDrawTwo}{1.09~[0.88, 1.35]}  % experiments/FIGS/data/paired_v1err/probes/ratios.json :: ratios[396].ratio, ratios[396].ci95 :: sha256 14dd51402305d147 :: random draw 3 over random draw 2, visible_content linear 1-auc: paired geometric mean [95% interval]
\newcommand{\PairedProbeVisibleContentMlpOnesixtwoOverOneeighteight}{0.92~[0.85, 1.01]}  % experiments/FIGS/data/paired_v1err/probes/ratios.json :: ratios[442].ratio, ratios[442].ci95 :: sha256 14dd51402305d147 :: 162 over 188, visible_content mlp 1-auc: paired geometric mean [95% interval]
\newcommand{\PairedProbeVisibleContentMlpFourthreeOverOneeighteight}{0.93~[0.80, 1.07]}  % experiments/FIGS/data/paired_v1err/probes/ratios.json :: ratios[443].ratio, ratios[443].ci95 :: sha256 14dd51402305d147 :: 43 over 188, visible_content mlp 1-auc: paired geometric mean [95% interval]
\newcommand{\PairedProbeVisibleContentMlpOnesevenOverOneeighteight}{2.3~[1.7, 3.1]}  % experiments/FIGS/data/paired_v1err/probes/ratios.json :: ratios[444].ratio, ratios[444].ci95 :: sha256 14dd51402305d147 :: 17 over 188, visible_content mlp 1-auc: paired geometric mean [95% interval]
\newcommand{\PairedProbeVisibleContentMlpFourthreeOverOnesixtwo}{1.01~[0.89, 1.14]}  % experiments/FIGS/data/paired_v1err/probes/ratios.json :: ratios[445].ratio, ratios[445].ci95 :: sha256 14dd51402305d147 :: 43 over 162, visible_content mlp 1-auc: paired geometric mean [95% interval]
\newcommand{\PairedProbeVisibleContentMlpOnesevenOverOnesixtwo}{2.5~[1.9, 3.2]}  % experiments/FIGS/data/paired_v1err/probes/ratios.json :: ratios[446].ratio, ratios[446].ci95 :: sha256 14dd51402305d147 :: 17 over 162, visible_content mlp 1-auc: paired geometric mean [95% interval]
\newcommand{\PairedProbeVisibleContentMlpOnesevenOverFourthree}{2.5~[2.0, 3.1]}  % experiments/FIGS/data/paired_v1err/probes/ratios.json :: ratios[447].ratio, ratios[447].ci95 :: sha256 14dd51402305d147 :: 17 over 43, visible_content mlp 1-auc: paired geometric mean [95% interval]
\newcommand{\PairedProbeVisibleContentMlpOnesixtwoMassOverOnesixtwo}{1.08~[0.94, 1.24]}  % experiments/FIGS/data/paired_v1err/probes/ratios.json :: ratios[448].ratio, ratios[448].ci95 :: sha256 14dd51402305d147 :: 162+mass over 162, visible_content mlp 1-auc: paired geometric mean [95% interval]
\newcommand{\PairedProbeVisibleContentMlpOnesevenMassOverOneseven}{1.41~[1.19, 1.67]}  % experiments/FIGS/data/paired_v1err/probes/ratios.json :: ratios[449].ratio, ratios[449].ci95 :: sha256 14dd51402305d147 :: 17+mass over 17, visible_content mlp 1-auc: paired geometric mean [95% interval]
\newcommand{\PairedProbeVisibleContentMlpRandomOverOneeighteight}{1.27~[1.01, 1.59]}  % experiments/FIGS/data/paired_v1err/probes/ratios.json :: ratios[450].ratio, ratios[450].ci95 :: sha256 14dd51402305d147 :: random over 188, visible_content mlp 1-auc: paired geometric mean [95% interval]
\newcommand{\PairedProbeVisibleContentMlpRandomOverOnesixtwo}{1.37~[1.12, 1.68]}  % experiments/FIGS/data/paired_v1err/probes/ratios.json :: ratios[451].ratio, ratios[451].ci95 :: sha256 14dd51402305d147 :: random over 162, visible_content mlp 1-auc: paired geometric mean [95% interval]
\newcommand{\PairedProbeVisibleContentMlpRandomOverFourthree}{1.35~[1.13, 1.61]}  % experiments/FIGS/data/paired_v1err/probes/ratios.json :: ratios[452].ratio, ratios[452].ci95 :: sha256 14dd51402305d147 :: random over 43, visible_content mlp 1-auc: paired geometric mean [95% interval]
\newcommand{\PairedProbeVisibleContentMlpRandomOverOneseven}{0.56~[0.45, 0.70]}  % experiments/FIGS/data/paired_v1err/probes/ratios.json :: ratios[453].ratio, ratios[453].ci95 :: sha256 14dd51402305d147 :: random over 17, visible_content mlp 1-auc: paired geometric mean [95% interval]
\newcommand{\PairedProbeVisibleContentMlpRandomDrawTwoOverRandomDrawOne}{0.88~[0.71, 1.09]}  % experiments/FIGS/data/paired_v1err/probes/ratios.json :: ratios[454].ratio, ratios[454].ci95 :: sha256 14dd51402305d147 :: random draw 2 over random draw 1, visible_content mlp 1-auc: paired geometric mean [95% interval]
\newcommand{\PairedProbeVisibleContentMlpRandomDrawThreeOverRandomDrawOne}{0.94~[0.75, 1.18]}  % experiments/FIGS/data/paired_v1err/probes/ratios.json :: ratios[455].ratio, ratios[455].ci95 :: sha256 14dd51402305d147 :: random draw 3 over random draw 1, visible_content mlp 1-auc: paired geometric mean [95% interval]
\newcommand{\PairedProbeVisibleContentMlpRandomDrawThreeOverRandomDrawTwo}{1.06~[0.84, 1.35]}  % experiments/FIGS/data/paired_v1err/probes/ratios.json :: ratios[456].ratio, ratios[456].ci95 :: sha256 14dd51402305d147 :: random draw 3 over random draw 2, visible_content mlp 1-auc: paired geometric mean [95% interval]
\newcommand{\PairedFtJciiEThreeOnesixtwoOverOneeighteight}{0.95~[0.87, 1.05]}  % experiments/FIGS/data/paired_v1err/finetune/ratios.json :: ratios[0].ratio, ratios[0].ci95 :: sha256 3e5cf25e859b68c2 :: 162 over 188, leg1 N1000 1-macro_auc: paired geometric mean [95% interval]
\newcommand{\PairedFtJciiEThreeFourthreeOverOneeighteight}{0.84~[0.77, 0.91]}  % experiments/FIGS/data/paired_v1err/finetune/ratios.json :: ratios[1].ratio, ratios[1].ci95 :: sha256 3e5cf25e859b68c2 :: 43 over 188, leg1 N1000 1-macro_auc: paired geometric mean [95% interval]
\newcommand{\PairedFtJciiEThreeOnesevenOverOneeighteight}{0.97~[0.86, 1.09]}  % experiments/FIGS/data/paired_v1err/finetune/ratios.json :: ratios[2].ratio, ratios[2].ci95 :: sha256 3e5cf25e859b68c2 :: 17 over 188, leg1 N1000 1-macro_auc: paired geometric mean [95% interval]
\newcommand{\PairedFtJciiEThreeFourthreeOverOnesixtwo}{0.88~[0.80, 0.97]}  % experiments/FIGS/data/paired_v1err/finetune/ratios.json :: ratios[3].ratio, ratios[3].ci95 :: sha256 3e5cf25e859b68c2 :: 43 over 162, leg1 N1000 1-macro_auc: paired geometric mean [95% interval]
\newcommand{\PairedFtJciiEThreeOnesevenOverOnesixtwo}{1.02~[0.88, 1.18]}  % experiments/FIGS/data/paired_v1err/finetune/ratios.json :: ratios[4].ratio, ratios[4].ci95 :: sha256 3e5cf25e859b68c2 :: 17 over 162, leg1 N1000 1-macro_auc: paired geometric mean [95% interval]
\newcommand{\PairedFtJciiEThreeOnesevenOverFourthree}{1.16~[1.08, 1.24]}  % experiments/FIGS/data/paired_v1err/finetune/ratios.json :: ratios[5].ratio, ratios[5].ci95 :: sha256 3e5cf25e859b68c2 :: 17 over 43, leg1 N1000 1-macro_auc: paired geometric mean [95% interval]
\newcommand{\PairedFtJciiEThreeOnesixtwoMassOverOnesixtwo}{0.85~[0.77, 0.93]}  % experiments/FIGS/data/paired_v1err/finetune/ratios.json :: ratios[6].ratio, ratios[6].ci95 :: sha256 3e5cf25e859b68c2 :: 162+mass over 162, leg1 N1000 1-macro_auc: paired geometric mean [95% interval]
\newcommand{\PairedFtJciiEThreeOnesevenMassOverOneseven}{1.00~[0.86, 1.17]}  % experiments/FIGS/data/paired_v1err/finetune/ratios.json :: ratios[7].ratio, ratios[7].ci95 :: sha256 3e5cf25e859b68c2 :: 17+mass over 17, leg1 N1000 1-macro_auc: paired geometric mean [95% interval]
\newcommand{\PairedFtJciiEThreeRandomOverOneeighteight}{0.82~[0.73, 0.93]}  % experiments/FIGS/data/paired_v1err/finetune/ratios.json :: ratios[8].ratio, ratios[8].ci95 :: sha256 3e5cf25e859b68c2 :: random over 188, leg1 N1000 1-macro_auc: paired geometric mean [95% interval]
\newcommand{\PairedFtJciiEThreeRandomOverOnesixtwo}{0.90~[0.85, 0.96]}  % experiments/FIGS/data/paired_v1err/finetune/ratios.json :: ratios[9].ratio, ratios[9].ci95 :: sha256 3e5cf25e859b68c2 :: random over 162, leg1 N1000 1-macro_auc: paired geometric mean [95% interval]
\newcommand{\PairedFtJciiEThreeRandomOverFourthree}{0.99~[0.75, 1.30]}  % experiments/FIGS/data/paired_v1err/finetune/ratios.json :: ratios[10].ratio, ratios[10].ci95 :: sha256 3e5cf25e859b68c2 :: random over 43, leg1 N1000 1-macro_auc: paired geometric mean [95% interval]
\newcommand{\PairedFtJciiEThreeRandomOverOneseven}{0.86~[0.58, 1.27]}  % experiments/FIGS/data/paired_v1err/finetune/ratios.json :: ratios[11].ratio, ratios[11].ci95 :: sha256 3e5cf25e859b68c2 :: random over 17, leg1 N1000 1-macro_auc: paired geometric mean [95% interval]
\newcommand{\PairedFtJciiEThreeRandomDrawTwoOverRandomDrawOne}{1.09~[1.07, 1.10]}  % experiments/FIGS/data/paired_v1err/finetune/ratios.json :: ratios[12].ratio, ratios[12].ci95 :: sha256 3e5cf25e859b68c2 :: random draw 2 over random draw 1, leg1 N1000 1-macro_auc: paired geometric mean [95% interval]
\newcommand{\PairedFtJciiEThreeRandomDrawThreeOverRandomDrawOne}{1.04~[1.02, 1.05]}  % experiments/FIGS/data/paired_v1err/finetune/ratios.json :: ratios[13].ratio, ratios[13].ci95 :: sha256 3e5cf25e859b68c2 :: random draw 3 over random draw 1, leg1 N1000 1-macro_auc: paired geometric mean [95% interval]
\newcommand{\PairedFtJciiEThreeRandomDrawThreeOverRandomDrawTwo}{0.95~[0.94, 0.97]}  % experiments/FIGS/data/paired_v1err/finetune/ratios.json :: ratios[14].ratio, ratios[14].ci95 :: sha256 3e5cf25e859b68c2 :: random draw 3 over random draw 2, leg1 N1000 1-macro_auc: paired geometric mean [95% interval]
\newcommand{\PairedFtJciiEFourOnesixtwoOverOneeighteight}{0.99~[0.97, 1.02]}  % experiments/FIGS/data/paired_v1err/finetune/ratios.json :: ratios[15].ratio, ratios[15].ci95 :: sha256 3e5cf25e859b68c2 :: 162 over 188, leg1 N10000 1-macro_auc: paired geometric mean [95% interval]
\newcommand{\PairedFtJciiEFourFourthreeOverOneeighteight}{1.20~[1.17, 1.23]}  % experiments/FIGS/data/paired_v1err/finetune/ratios.json :: ratios[16].ratio, ratios[16].ci95 :: sha256 3e5cf25e859b68c2 :: 43 over 188, leg1 N10000 1-macro_auc: paired geometric mean [95% interval]
\newcommand{\PairedFtJciiEFourOnesevenOverOneeighteight}{1.32~[1.29, 1.34]}  % experiments/FIGS/data/paired_v1err/finetune/ratios.json :: ratios[17].ratio, ratios[17].ci95 :: sha256 3e5cf25e859b68c2 :: 17 over 188, leg1 N10000 1-macro_auc: paired geometric mean [95% interval]
\newcommand{\PairedFtJciiEFourFourthreeOverOnesixtwo}{1.21~[1.19, 1.23]}  % experiments/FIGS/data/paired_v1err/finetune/ratios.json :: ratios[18].ratio, ratios[18].ci95 :: sha256 3e5cf25e859b68c2 :: 43 over 162, leg1 N10000 1-macro_auc: paired geometric mean [95% interval]
\newcommand{\PairedFtJciiEFourOnesevenOverOnesixtwo}{1.32~[1.29, 1.36]}  % experiments/FIGS/data/paired_v1err/finetune/ratios.json :: ratios[19].ratio, ratios[19].ci95 :: sha256 3e5cf25e859b68c2 :: 17 over 162, leg1 N10000 1-macro_auc: paired geometric mean [95% interval]
\newcommand{\PairedFtJciiEFourOnesevenOverFourthree}{1.10~[1.07, 1.13]}  % experiments/FIGS/data/paired_v1err/finetune/ratios.json :: ratios[20].ratio, ratios[20].ci95 :: sha256 3e5cf25e859b68c2 :: 17 over 43, leg1 N10000 1-macro_auc: paired geometric mean [95% interval]
\newcommand{\PairedFtJciiEFourOnesixtwoMassOverOnesixtwo}{1.00~[0.94, 1.06]}  % experiments/FIGS/data/paired_v1err/finetune/ratios.json :: ratios[21].ratio, ratios[21].ci95 :: sha256 3e5cf25e859b68c2 :: 162+mass over 162, leg1 N10000 1-macro_auc: paired geometric mean [95% interval]
\newcommand{\PairedFtJciiEFourOnesevenMassOverOneseven}{0.990~[0.976, 1.004]}  % experiments/FIGS/data/paired_v1err/finetune/ratios.json :: ratios[22].ratio, ratios[22].ci95 :: sha256 3e5cf25e859b68c2 :: 17+mass over 17, leg1 N10000 1-macro_auc: paired geometric mean [95% interval]
\newcommand{\PairedFtJciiEFourRandomOverOneeighteight}{1.15~[1.09, 1.22]}  % experiments/FIGS/data/paired_v1err/finetune/ratios.json :: ratios[23].ratio, ratios[23].ci95 :: sha256 3e5cf25e859b68c2 :: random over 188, leg1 N10000 1-macro_auc: paired geometric mean [95% interval]
\newcommand{\PairedFtJciiEFourRandomOverOnesixtwo}{1.18~[1.13, 1.22]}  % experiments/FIGS/data/paired_v1err/finetune/ratios.json :: ratios[24].ratio, ratios[24].ci95 :: sha256 3e5cf25e859b68c2 :: random over 162, leg1 N10000 1-macro_auc: paired geometric mean [95% interval]
\newcommand{\PairedFtJciiEFourRandomOverFourthree}{0.97~[0.96, 0.99]}  % experiments/FIGS/data/paired_v1err/finetune/ratios.json :: ratios[25].ratio, ratios[25].ci95 :: sha256 3e5cf25e859b68c2 :: random over 43, leg1 N10000 1-macro_auc: paired geometric mean [95% interval]
\newcommand{\PairedFtJciiEFourRandomOverOneseven}{0.88~[0.82, 0.95]}  % experiments/FIGS/data/paired_v1err/finetune/ratios.json :: ratios[26].ratio, ratios[26].ci95 :: sha256 3e5cf25e859b68c2 :: random over 17, leg1 N10000 1-macro_auc: paired geometric mean [95% interval]
\newcommand{\PairedFtJciiEFourRandomDrawTwoOverRandomDrawOne}{1.00~[0.98, 1.02]}  % experiments/FIGS/data/paired_v1err/finetune/ratios.json :: ratios[27].ratio, ratios[27].ci95 :: sha256 3e5cf25e859b68c2 :: random draw 2 over random draw 1, leg1 N10000 1-macro_auc: paired geometric mean [95% interval]
\newcommand{\PairedFtJciiEFourRandomDrawThreeOverRandomDrawOne}{0.97~[0.95, 0.98]}  % experiments/FIGS/data/paired_v1err/finetune/ratios.json :: ratios[28].ratio, ratios[28].ci95 :: sha256 3e5cf25e859b68c2 :: random draw 3 over random draw 1, leg1 N10000 1-macro_auc: paired geometric mean [95% interval]
\newcommand{\PairedFtJciiEFourRandomDrawThreeOverRandomDrawTwo}{0.97~[0.95, 0.98]}  % experiments/FIGS/data/paired_v1err/finetune/ratios.json :: ratios[29].ratio, ratios[29].ci95 :: sha256 3e5cf25e859b68c2 :: random draw 3 over random draw 2, leg1 N10000 1-macro_auc: paired geometric mean [95% interval]
\newcommand{\PairedFtJciiEFiveOnesixtwoOverOneeighteight}{1.00~[0.98, 1.01]}  % experiments/FIGS/data/paired_v1err/finetune/ratios.json :: ratios[30].ratio, ratios[30].ci95 :: sha256 3e5cf25e859b68c2 :: 162 over 188, leg1 N100000 1-macro_auc: paired geometric mean [95% interval]
\newcommand{\PairedFtJciiEFiveFourthreeOverOneeighteight}{1.09~[1.06, 1.11]}  % experiments/FIGS/data/paired_v1err/finetune/ratios.json :: ratios[31].ratio, ratios[31].ci95 :: sha256 3e5cf25e859b68c2 :: 43 over 188, leg1 N100000 1-macro_auc: paired geometric mean [95% interval]
\newcommand{\PairedFtJciiEFiveOnesevenOverOneeighteight}{1.22~[1.20, 1.23]}  % experiments/FIGS/data/paired_v1err/finetune/ratios.json :: ratios[32].ratio, ratios[32].ci95 :: sha256 3e5cf25e859b68c2 :: 17 over 188, leg1 N100000 1-macro_auc: paired geometric mean [95% interval]
\newcommand{\PairedFtJciiEFiveFourthreeOverOnesixtwo}{1.09~[1.07, 1.11]}  % experiments/FIGS/data/paired_v1err/finetune/ratios.json :: ratios[33].ratio, ratios[33].ci95 :: sha256 3e5cf25e859b68c2 :: 43 over 162, leg1 N100000 1-macro_auc: paired geometric mean [95% interval]
\newcommand{\PairedFtJciiEFiveOnesevenOverOnesixtwo}{1.22~[1.20, 1.24]}  % experiments/FIGS/data/paired_v1err/finetune/ratios.json :: ratios[34].ratio, ratios[34].ci95 :: sha256 3e5cf25e859b68c2 :: 17 over 162, leg1 N100000 1-macro_auc: paired geometric mean [95% interval]
\newcommand{\PairedFtJciiEFiveOnesevenOverFourthree}{1.12~[1.10, 1.14]}  % experiments/FIGS/data/paired_v1err/finetune/ratios.json :: ratios[35].ratio, ratios[35].ci95 :: sha256 3e5cf25e859b68c2 :: 17 over 43, leg1 N100000 1-macro_auc: paired geometric mean [95% interval]
\newcommand{\PairedFtJciiEFiveOnesixtwoMassOverOnesixtwo}{1.018~[1.002, 1.035]}  % experiments/FIGS/data/paired_v1err/finetune/ratios.json :: ratios[36].ratio, ratios[36].ci95 :: sha256 3e5cf25e859b68c2 :: 162+mass over 162, leg1 N100000 1-macro_auc: paired geometric mean [95% interval]
\newcommand{\PairedFtJciiEFiveOnesevenMassOverOneseven}{1.03~[1.01, 1.06]}  % experiments/FIGS/data/paired_v1err/finetune/ratios.json :: ratios[37].ratio, ratios[37].ci95 :: sha256 3e5cf25e859b68c2 :: 17+mass over 17, leg1 N100000 1-macro_auc: paired geometric mean [95% interval]
\newcommand{\PairedFtJciiEFiveRandomOverOneeighteight}{1.15~[1.09, 1.20]}  % experiments/FIGS/data/paired_v1err/finetune/ratios.json :: ratios[38].ratio, ratios[38].ci95 :: sha256 3e5cf25e859b68c2 :: random over 188, leg1 N100000 1-macro_auc: paired geometric mean [95% interval]
\newcommand{\PairedFtJciiEFiveRandomOverOnesixtwo}{1.15~[1.09, 1.20]}  % experiments/FIGS/data/paired_v1err/finetune/ratios.json :: ratios[39].ratio, ratios[39].ci95 :: sha256 3e5cf25e859b68c2 :: random over 162, leg1 N100000 1-macro_auc: paired geometric mean [95% interval]
\newcommand{\PairedFtJciiEFiveRandomOverFourthree}{1.05~[1.04, 1.06]}  % experiments/FIGS/data/paired_v1err/finetune/ratios.json :: ratios[40].ratio, ratios[40].ci95 :: sha256 3e5cf25e859b68c2 :: random over 43, leg1 N100000 1-macro_auc: paired geometric mean [95% interval]
\newcommand{\PairedFtJciiEFiveRandomOverOneseven}{0.95~[0.93, 0.96]}  % experiments/FIGS/data/paired_v1err/finetune/ratios.json :: ratios[41].ratio, ratios[41].ci95 :: sha256 3e5cf25e859b68c2 :: random over 17, leg1 N100000 1-macro_auc: paired geometric mean [95% interval]
\newcommand{\PairedFtJciiEFiveRandomDrawTwoOverRandomDrawOne}{0.98~[0.97, 0.99]}  % experiments/FIGS/data/paired_v1err/finetune/ratios.json :: ratios[42].ratio, ratios[42].ci95 :: sha256 3e5cf25e859b68c2 :: random draw 2 over random draw 1, leg1 N100000 1-macro_auc: paired geometric mean [95% interval]
\newcommand{\PairedFtJciiEFiveRandomDrawThreeOverRandomDrawOne}{1.00~[0.99, 1.02]}  % experiments/FIGS/data/paired_v1err/finetune/ratios.json :: ratios[43].ratio, ratios[43].ci95 :: sha256 3e5cf25e859b68c2 :: random draw 3 over random draw 1, leg1 N100000 1-macro_auc: paired geometric mean [95% interval]
\newcommand{\PairedFtJciiEFiveRandomDrawThreeOverRandomDrawTwo}{1.03~[1.01, 1.04]}  % experiments/FIGS/data/paired_v1err/finetune/ratios.json :: ratios[44].ratio, ratios[44].ci95 :: sha256 3e5cf25e859b68c2 :: random draw 3 over random draw 2, leg1 N100000 1-macro_auc: paired geometric mean [95% interval]
\newcommand{\PairedFtJciiESixOnesixtwoOverOneeighteight}{1.002~[0.998, 1.005]}  % experiments/FIGS/data/paired_v1err/finetune/ratios.json :: ratios[45].ratio, ratios[45].ci95 :: sha256 3e5cf25e859b68c2 :: 162 over 188, leg1 N1000000 1-macro_auc: paired geometric mean [95% interval]
\newcommand{\PairedFtJciiESixFourthreeOverOneeighteight}{1.045~[1.037, 1.053]}  % experiments/FIGS/data/paired_v1err/finetune/ratios.json :: ratios[46].ratio, ratios[46].ci95 :: sha256 3e5cf25e859b68c2 :: 43 over 188, leg1 N1000000 1-macro_auc: paired geometric mean [95% interval]
\newcommand{\PairedFtJciiESixOnesevenOverOneeighteight}{1.14~[1.13, 1.16]}  % experiments/FIGS/data/paired_v1err/finetune/ratios.json :: ratios[47].ratio, ratios[47].ci95 :: sha256 3e5cf25e859b68c2 :: 17 over 188, leg1 N1000000 1-macro_auc: paired geometric mean [95% interval]
\newcommand{\PairedFtJciiESixFourthreeOverOnesixtwo}{1.044~[1.036, 1.051]}  % experiments/FIGS/data/paired_v1err/finetune/ratios.json :: ratios[48].ratio, ratios[48].ci95 :: sha256 3e5cf25e859b68c2 :: 43 over 162, leg1 N1000000 1-macro_auc: paired geometric mean [95% interval]
\newcommand{\PairedFtJciiESixOnesevenOverOnesixtwo}{1.14~[1.13, 1.15]}  % experiments/FIGS/data/paired_v1err/finetune/ratios.json :: ratios[49].ratio, ratios[49].ci95 :: sha256 3e5cf25e859b68c2 :: 17 over 162, leg1 N1000000 1-macro_auc: paired geometric mean [95% interval]
\newcommand{\PairedFtJciiESixOnesevenOverFourthree}{1.09~[1.08, 1.11]}  % experiments/FIGS/data/paired_v1err/finetune/ratios.json :: ratios[50].ratio, ratios[50].ci95 :: sha256 3e5cf25e859b68c2 :: 17 over 43, leg1 N1000000 1-macro_auc: paired geometric mean [95% interval]
\newcommand{\PairedFtJciiESixOnesixtwoMassOverOnesixtwo}{1.019~[1.014, 1.023]}  % experiments/FIGS/data/paired_v1err/finetune/ratios.json :: ratios[51].ratio, ratios[51].ci95 :: sha256 3e5cf25e859b68c2 :: 162+mass over 162, leg1 N1000000 1-macro_auc: paired geometric mean [95% interval]
\newcommand{\PairedFtJciiESixOnesevenMassOverOneseven}{1.03~[1.02, 1.04]}  % experiments/FIGS/data/paired_v1err/finetune/ratios.json :: ratios[52].ratio, ratios[52].ci95 :: sha256 3e5cf25e859b68c2 :: 17+mass over 17, leg1 N1000000 1-macro_auc: paired geometric mean [95% interval]
\newcommand{\PairedFtJciiESixRandomOverOneeighteight}{1.073~[1.067, 1.080]}  % experiments/FIGS/data/paired_v1err/finetune/ratios.json :: ratios[53].ratio, ratios[53].ci95 :: sha256 3e5cf25e859b68c2 :: random over 188, leg1 N1000000 1-macro_auc: paired geometric mean [95% interval]
\newcommand{\PairedFtJciiESixRandomOverOnesixtwo}{1.072~[1.065, 1.078]}  % experiments/FIGS/data/paired_v1err/finetune/ratios.json :: ratios[54].ratio, ratios[54].ci95 :: sha256 3e5cf25e859b68c2 :: random over 162, leg1 N1000000 1-macro_auc: paired geometric mean [95% interval]
\newcommand{\PairedFtJciiESixRandomOverFourthree}{1.03~[1.02, 1.04]}  % experiments/FIGS/data/paired_v1err/finetune/ratios.json :: ratios[55].ratio, ratios[55].ci95 :: sha256 3e5cf25e859b68c2 :: random over 43, leg1 N1000000 1-macro_auc: paired geometric mean [95% interval]
\newcommand{\PairedFtJciiESixRandomOverOneseven}{0.94~[0.93, 0.95]}  % experiments/FIGS/data/paired_v1err/finetune/ratios.json :: ratios[56].ratio, ratios[56].ci95 :: sha256 3e5cf25e859b68c2 :: random over 17, leg1 N1000000 1-macro_auc: paired geometric mean [95% interval]
\newcommand{\PairedFtJciiESixRandomDrawTwoOverRandomDrawOne}{1.00~[0.99, 1.01]}  % experiments/FIGS/data/paired_v1err/finetune/ratios.json :: ratios[57].ratio, ratios[57].ci95 :: sha256 3e5cf25e859b68c2 :: random draw 2 over random draw 1, leg1 N1000000 1-macro_auc: paired geometric mean [95% interval]
\newcommand{\PairedFtJciiESixRandomDrawThreeOverRandomDrawOne}{0.996~[0.989, 1.003]}  % experiments/FIGS/data/paired_v1err/finetune/ratios.json :: ratios[58].ratio, ratios[58].ci95 :: sha256 3e5cf25e859b68c2 :: random draw 3 over random draw 1, leg1 N1000000 1-macro_auc: paired geometric mean [95% interval]
\newcommand{\PairedFtJciiESixRandomDrawThreeOverRandomDrawTwo}{0.997~[0.989, 1.004]}  % experiments/FIGS/data/paired_v1err/finetune/ratios.json :: ratios[59].ratio, ratios[59].ci95 :: sha256 3e5cf25e859b68c2 :: random draw 3 over random draw 2, leg1 N1000000 1-macro_auc: paired geometric mean [95% interval]
\newcommand{\PairedFtJcEThreeOnesixtwoOverOneeighteight}{0.96~[0.83, 1.11]}  % experiments/FIGS/data/paired_v1err/finetune/ratios.json :: ratios[60].ratio, ratios[60].ci95 :: sha256 3e5cf25e859b68c2 :: 162 over 188, leg2 N1000 1-macro_auc: paired geometric mean [95% interval]
\newcommand{\PairedFtJcEThreeFourthreeOverOneeighteight}{1.01~[0.84, 1.21]}  % experiments/FIGS/data/paired_v1err/finetune/ratios.json :: ratios[61].ratio, ratios[61].ci95 :: sha256 3e5cf25e859b68c2 :: 43 over 188, leg2 N1000 1-macro_auc: paired geometric mean [95% interval]
\newcommand{\PairedFtJcEThreeOnesevenOverOneeighteight}{1.07~[0.93, 1.23]}  % experiments/FIGS/data/paired_v1err/finetune/ratios.json :: ratios[62].ratio, ratios[62].ci95 :: sha256 3e5cf25e859b68c2 :: 17 over 188, leg2 N1000 1-macro_auc: paired geometric mean [95% interval]
\newcommand{\PairedFtJcEThreeFourthreeOverOnesixtwo}{1.05~[0.96, 1.15]}  % experiments/FIGS/data/paired_v1err/finetune/ratios.json :: ratios[63].ratio, ratios[63].ci95 :: sha256 3e5cf25e859b68c2 :: 43 over 162, leg2 N1000 1-macro_auc: paired geometric mean [95% interval]
\newcommand{\PairedFtJcEThreeOnesevenOverOnesixtwo}{1.12~[1.07, 1.17]}  % experiments/FIGS/data/paired_v1err/finetune/ratios.json :: ratios[64].ratio, ratios[64].ci95 :: sha256 3e5cf25e859b68c2 :: 17 over 162, leg2 N1000 1-macro_auc: paired geometric mean [95% interval]
\newcommand{\PairedFtJcEThreeOnesevenOverFourthree}{1.06~[0.98, 1.15]}  % experiments/FIGS/data/paired_v1err/finetune/ratios.json :: ratios[65].ratio, ratios[65].ci95 :: sha256 3e5cf25e859b68c2 :: 17 over 43, leg2 N1000 1-macro_auc: paired geometric mean [95% interval]
\newcommand{\PairedFtJcEThreeOnesixtwoMassOverOnesixtwo}{0.91~[0.88, 0.94]}  % experiments/FIGS/data/paired_v1err/finetune/ratios.json :: ratios[66].ratio, ratios[66].ci95 :: sha256 3e5cf25e859b68c2 :: 162+mass over 162, leg2 N1000 1-macro_auc: paired geometric mean [95% interval]
\newcommand{\PairedFtJcEThreeOnesevenMassOverOneseven}{1.02~[0.92, 1.13]}  % experiments/FIGS/data/paired_v1err/finetune/ratios.json :: ratios[67].ratio, ratios[67].ci95 :: sha256 3e5cf25e859b68c2 :: 17+mass over 17, leg2 N1000 1-macro_auc: paired geometric mean [95% interval]
\newcommand{\PairedFtJcEThreeRandomOverOneeighteight}{1.01~[0.78, 1.32]}  % experiments/FIGS/data/paired_v1err/finetune/ratios.json :: ratios[68].ratio, ratios[68].ci95 :: sha256 3e5cf25e859b68c2 :: random over 188, leg2 N1000 1-macro_auc: paired geometric mean [95% interval]
\newcommand{\PairedFtJcEThreeRandomOverOnesixtwo}{1.12~[1.04, 1.20]}  % experiments/FIGS/data/paired_v1err/finetune/ratios.json :: ratios[69].ratio, ratios[69].ci95 :: sha256 3e5cf25e859b68c2 :: random over 162, leg2 N1000 1-macro_auc: paired geometric mean [95% interval]
\newcommand{\PairedFtJcEThreeRandomOverFourthree}{1.06~[0.79, 1.41]}  % experiments/FIGS/data/paired_v1err/finetune/ratios.json :: ratios[70].ratio, ratios[70].ci95 :: sha256 3e5cf25e859b68c2 :: random over 43, leg2 N1000 1-macro_auc: paired geometric mean [95% interval]
\newcommand{\PairedFtJcEThreeRandomOverOneseven}{0.98~[0.84, 1.13]}  % experiments/FIGS/data/paired_v1err/finetune/ratios.json :: ratios[71].ratio, ratios[71].ci95 :: sha256 3e5cf25e859b68c2 :: random over 17, leg2 N1000 1-macro_auc: paired geometric mean [95% interval]
\newcommand{\PairedFtJcEThreeRandomDrawTwoOverRandomDrawOne}{0.954~[0.949, 0.958]}  % experiments/FIGS/data/paired_v1err/finetune/ratios.json :: ratios[72].ratio, ratios[72].ci95 :: sha256 3e5cf25e859b68c2 :: random draw 2 over random draw 1, leg2 N1000 1-macro_auc: paired geometric mean [95% interval]
\newcommand{\PairedFtJcEThreeRandomDrawThreeOverRandomDrawOne}{1.013~[1.009, 1.017]}  % experiments/FIGS/data/paired_v1err/finetune/ratios.json :: ratios[73].ratio, ratios[73].ci95 :: sha256 3e5cf25e859b68c2 :: random draw 3 over random draw 1, leg2 N1000 1-macro_auc: paired geometric mean [95% interval]
\newcommand{\PairedFtJcEThreeRandomDrawThreeOverRandomDrawTwo}{1.062~[1.057, 1.067]}  % experiments/FIGS/data/paired_v1err/finetune/ratios.json :: ratios[74].ratio, ratios[74].ci95 :: sha256 3e5cf25e859b68c2 :: random draw 3 over random draw 2, leg2 N1000 1-macro_auc: paired geometric mean [95% interval]
\newcommand{\PairedFtJcEFourOnesixtwoOverOneeighteight}{0.98~[0.93, 1.02]}  % experiments/FIGS/data/paired_v1err/finetune/ratios.json :: ratios[75].ratio, ratios[75].ci95 :: sha256 3e5cf25e859b68c2 :: 162 over 188, leg2 N10000 1-macro_auc: paired geometric mean [95% interval]
\newcommand{\PairedFtJcEFourFourthreeOverOneeighteight}{0.98~[0.93, 1.04]}  % experiments/FIGS/data/paired_v1err/finetune/ratios.json :: ratios[76].ratio, ratios[76].ci95 :: sha256 3e5cf25e859b68c2 :: 43 over 188, leg2 N10000 1-macro_auc: paired geometric mean [95% interval]
\newcommand{\PairedFtJcEFourOnesevenOverOneeighteight}{1.07~[1.02, 1.11]}  % experiments/FIGS/data/paired_v1err/finetune/ratios.json :: ratios[77].ratio, ratios[77].ci95 :: sha256 3e5cf25e859b68c2 :: 17 over 188, leg2 N10000 1-macro_auc: paired geometric mean [95% interval]
\newcommand{\PairedFtJcEFourFourthreeOverOnesixtwo}{1.01~[0.97, 1.04]}  % experiments/FIGS/data/paired_v1err/finetune/ratios.json :: ratios[78].ratio, ratios[78].ci95 :: sha256 3e5cf25e859b68c2 :: 43 over 162, leg2 N10000 1-macro_auc: paired geometric mean [95% interval]
\newcommand{\PairedFtJcEFourOnesevenOverOnesixtwo}{1.09~[1.06, 1.12]}  % experiments/FIGS/data/paired_v1err/finetune/ratios.json :: ratios[79].ratio, ratios[79].ci95 :: sha256 3e5cf25e859b68c2 :: 17 over 162, leg2 N10000 1-macro_auc: paired geometric mean [95% interval]
\newcommand{\PairedFtJcEFourOnesevenOverFourthree}{1.08~[1.07, 1.10]}  % experiments/FIGS/data/paired_v1err/finetune/ratios.json :: ratios[80].ratio, ratios[80].ci95 :: sha256 3e5cf25e859b68c2 :: 17 over 43, leg2 N10000 1-macro_auc: paired geometric mean [95% interval]
\newcommand{\PairedFtJcEFourOnesixtwoMassOverOnesixtwo}{1.01~[0.98, 1.04]}  % experiments/FIGS/data/paired_v1err/finetune/ratios.json :: ratios[81].ratio, ratios[81].ci95 :: sha256 3e5cf25e859b68c2 :: 162+mass over 162, leg2 N10000 1-macro_auc: paired geometric mean [95% interval]
\newcommand{\PairedFtJcEFourOnesevenMassOverOneseven}{1.05~[1.02, 1.08]}  % experiments/FIGS/data/paired_v1err/finetune/ratios.json :: ratios[82].ratio, ratios[82].ci95 :: sha256 3e5cf25e859b68c2 :: 17+mass over 17, leg2 N10000 1-macro_auc: paired geometric mean [95% interval]
\newcommand{\PairedFtJcEFourRandomOverOneeighteight}{1.08~[1.03, 1.14]}  % experiments/FIGS/data/paired_v1err/finetune/ratios.json :: ratios[83].ratio, ratios[83].ci95 :: sha256 3e5cf25e859b68c2 :: random over 188, leg2 N10000 1-macro_auc: paired geometric mean [95% interval]
\newcommand{\PairedFtJcEFourRandomOverOnesixtwo}{1.11~[0.99, 1.23]}  % experiments/FIGS/data/paired_v1err/finetune/ratios.json :: ratios[84].ratio, ratios[84].ci95 :: sha256 3e5cf25e859b68c2 :: random over 162, leg2 N10000 1-macro_auc: paired geometric mean [95% interval]
\newcommand{\PairedFtJcEFourRandomOverFourthree}{1.10~[0.99, 1.21]}  % experiments/FIGS/data/paired_v1err/finetune/ratios.json :: ratios[85].ratio, ratios[85].ci95 :: sha256 3e5cf25e859b68c2 :: random over 43, leg2 N10000 1-macro_auc: paired geometric mean [95% interval]
\newcommand{\PairedFtJcEFourRandomOverOneseven}{1.02~[0.94, 1.10]}  % experiments/FIGS/data/paired_v1err/finetune/ratios.json :: ratios[86].ratio, ratios[86].ci95 :: sha256 3e5cf25e859b68c2 :: random over 17, leg2 N10000 1-macro_auc: paired geometric mean [95% interval]
\newcommand{\PairedFtJcEFourRandomDrawTwoOverRandomDrawOne}{0.960~[0.955, 0.965]}  % experiments/FIGS/data/paired_v1err/finetune/ratios.json :: ratios[87].ratio, ratios[87].ci95 :: sha256 3e5cf25e859b68c2 :: random draw 2 over random draw 1, leg2 N10000 1-macro_auc: paired geometric mean [95% interval]
\newcommand{\PairedFtJcEFourRandomDrawThreeOverRandomDrawOne}{1.008~[1.004, 1.012]}  % experiments/FIGS/data/paired_v1err/finetune/ratios.json :: ratios[88].ratio, ratios[88].ci95 :: sha256 3e5cf25e859b68c2 :: random draw 3 over random draw 1, leg2 N10000 1-macro_auc: paired geometric mean [95% interval]
\newcommand{\PairedFtJcEFourRandomDrawThreeOverRandomDrawTwo}{1.050~[1.045, 1.055]}  % experiments/FIGS/data/paired_v1err/finetune/ratios.json :: ratios[89].ratio, ratios[89].ci95 :: sha256 3e5cf25e859b68c2 :: random draw 3 over random draw 2, leg2 N10000 1-macro_auc: paired geometric mean [95% interval]
\newcommand{\PairedFtJcEFiveOnesixtwoOverOneeighteight}{0.983~[0.968, 0.998]}  % experiments/FIGS/data/paired_v1err/finetune/ratios.json :: ratios[90].ratio, ratios[90].ci95 :: sha256 3e5cf25e859b68c2 :: 162 over 188, leg2 N100000 1-macro_auc: paired geometric mean [95% interval]
\newcommand{\PairedFtJcEFiveFourthreeOverOneeighteight}{1.015~[1.002, 1.029]}  % experiments/FIGS/data/paired_v1err/finetune/ratios.json :: ratios[91].ratio, ratios[91].ci95 :: sha256 3e5cf25e859b68c2 :: 43 over 188, leg2 N100000 1-macro_auc: paired geometric mean [95% interval]
\newcommand{\PairedFtJcEFiveOnesevenOverOneeighteight}{1.07~[1.06, 1.09]}  % experiments/FIGS/data/paired_v1err/finetune/ratios.json :: ratios[92].ratio, ratios[92].ci95 :: sha256 3e5cf25e859b68c2 :: 17 over 188, leg2 N100000 1-macro_auc: paired geometric mean [95% interval]
\newcommand{\PairedFtJcEFiveFourthreeOverOnesixtwo}{1.03~[1.02, 1.05]}  % experiments/FIGS/data/paired_v1err/finetune/ratios.json :: ratios[93].ratio, ratios[93].ci95 :: sha256 3e5cf25e859b68c2 :: 43 over 162, leg2 N100000 1-macro_auc: paired geometric mean [95% interval]
\newcommand{\PairedFtJcEFiveOnesevenOverOnesixtwo}{1.091~[1.088, 1.094]}  % experiments/FIGS/data/paired_v1err/finetune/ratios.json :: ratios[94].ratio, ratios[94].ci95 :: sha256 3e5cf25e859b68c2 :: 17 over 162, leg2 N100000 1-macro_auc: paired geometric mean [95% interval]
\newcommand{\PairedFtJcEFiveOnesevenOverFourthree}{1.06~[1.04, 1.07]}  % experiments/FIGS/data/paired_v1err/finetune/ratios.json :: ratios[95].ratio, ratios[95].ci95 :: sha256 3e5cf25e859b68c2 :: 17 over 43, leg2 N100000 1-macro_auc: paired geometric mean [95% interval]
\newcommand{\PairedFtJcEFiveOnesixtwoMassOverOnesixtwo}{1.016~[1.008, 1.024]}  % experiments/FIGS/data/paired_v1err/finetune/ratios.json :: ratios[96].ratio, ratios[96].ci95 :: sha256 3e5cf25e859b68c2 :: 162+mass over 162, leg2 N100000 1-macro_auc: paired geometric mean [95% interval]
\newcommand{\PairedFtJcEFiveOnesevenMassOverOneseven}{1.02~[1.01, 1.03]}  % experiments/FIGS/data/paired_v1err/finetune/ratios.json :: ratios[97].ratio, ratios[97].ci95 :: sha256 3e5cf25e859b68c2 :: 17+mass over 17, leg2 N100000 1-macro_auc: paired geometric mean [95% interval]
\newcommand{\PairedFtJcEFiveRandomOverOneeighteight}{1.044~[1.039, 1.049]}  % experiments/FIGS/data/paired_v1err/finetune/ratios.json :: ratios[98].ratio, ratios[98].ci95 :: sha256 3e5cf25e859b68c2 :: random over 188, leg2 N100000 1-macro_auc: paired geometric mean [95% interval]
\newcommand{\PairedFtJcEFiveRandomOverOnesixtwo}{1.06~[1.04, 1.09]}  % experiments/FIGS/data/paired_v1err/finetune/ratios.json :: ratios[99].ratio, ratios[99].ci95 :: sha256 3e5cf25e859b68c2 :: random over 162, leg2 N100000 1-macro_auc: paired geometric mean [95% interval]
\newcommand{\PairedFtJcEFiveRandomOverFourthree}{1.03~[1.01, 1.04]}  % experiments/FIGS/data/paired_v1err/finetune/ratios.json :: ratios[100].ratio, ratios[100].ci95 :: sha256 3e5cf25e859b68c2 :: random over 43, leg2 N100000 1-macro_auc: paired geometric mean [95% interval]
\newcommand{\PairedFtJcEFiveRandomOverOneseven}{0.97~[0.95, 0.99]}  % experiments/FIGS/data/paired_v1err/finetune/ratios.json :: ratios[101].ratio, ratios[101].ci95 :: sha256 3e5cf25e859b68c2 :: random over 17, leg2 N100000 1-macro_auc: paired geometric mean [95% interval]
\newcommand{\PairedFtJcEFiveRandomDrawTwoOverRandomDrawOne}{0.992~[0.988, 0.997]}  % experiments/FIGS/data/paired_v1err/finetune/ratios.json :: ratios[102].ratio, ratios[102].ci95 :: sha256 3e5cf25e859b68c2 :: random draw 2 over random draw 1, leg2 N100000 1-macro_auc: paired geometric mean [95% interval]
\newcommand{\PairedFtJcEFiveRandomDrawThreeOverRandomDrawOne}{1.007~[1.003, 1.011]}  % experiments/FIGS/data/paired_v1err/finetune/ratios.json :: ratios[103].ratio, ratios[103].ci95 :: sha256 3e5cf25e859b68c2 :: random draw 3 over random draw 1, leg2 N100000 1-macro_auc: paired geometric mean [95% interval]
\newcommand{\PairedFtJcEFiveRandomDrawThreeOverRandomDrawTwo}{1.015~[1.011, 1.020]}  % experiments/FIGS/data/paired_v1err/finetune/ratios.json :: ratios[104].ratio, ratios[104].ci95 :: sha256 3e5cf25e859b68c2 :: random draw 3 over random draw 2, leg2 N100000 1-macro_auc: paired geometric mean [95% interval]
\newcommand{\PairedFtJcESixOnesixtwoOverOneeighteight}{1.00~[0.99, 1.01]}  % experiments/FIGS/data/paired_v1err/finetune/ratios.json :: ratios[105].ratio, ratios[105].ci95 :: sha256 3e5cf25e859b68c2 :: 162 over 188, leg2 N1000000 1-macro_auc: paired geometric mean [95% interval]
\newcommand{\PairedFtJcESixFourthreeOverOneeighteight}{1.02~[1.01, 1.03]}  % experiments/FIGS/data/paired_v1err/finetune/ratios.json :: ratios[106].ratio, ratios[106].ci95 :: sha256 3e5cf25e859b68c2 :: 43 over 188, leg2 N1000000 1-macro_auc: paired geometric mean [95% interval]
\newcommand{\PairedFtJcESixOnesevenOverOneeighteight}{1.08~[1.07, 1.09]}  % experiments/FIGS/data/paired_v1err/finetune/ratios.json :: ratios[107].ratio, ratios[107].ci95 :: sha256 3e5cf25e859b68c2 :: 17 over 188, leg2 N1000000 1-macro_auc: paired geometric mean [95% interval]
\newcommand{\PairedFtJcESixFourthreeOverOnesixtwo}{1.02~[1.01, 1.03]}  % experiments/FIGS/data/paired_v1err/finetune/ratios.json :: ratios[108].ratio, ratios[108].ci95 :: sha256 3e5cf25e859b68c2 :: 43 over 162, leg2 N1000000 1-macro_auc: paired geometric mean [95% interval]
\newcommand{\PairedFtJcESixOnesevenOverOnesixtwo}{1.08~[1.07, 1.09]}  % experiments/FIGS/data/paired_v1err/finetune/ratios.json :: ratios[109].ratio, ratios[109].ci95 :: sha256 3e5cf25e859b68c2 :: 17 over 162, leg2 N1000000 1-macro_auc: paired geometric mean [95% interval]
\newcommand{\PairedFtJcESixOnesevenOverFourthree}{1.06~[1.05, 1.07]}  % experiments/FIGS/data/paired_v1err/finetune/ratios.json :: ratios[110].ratio, ratios[110].ci95 :: sha256 3e5cf25e859b68c2 :: 17 over 43, leg2 N1000000 1-macro_auc: paired geometric mean [95% interval]
\newcommand{\PairedFtJcESixOnesixtwoMassOverOnesixtwo}{1.012~[1.004, 1.020]}  % experiments/FIGS/data/paired_v1err/finetune/ratios.json :: ratios[111].ratio, ratios[111].ci95 :: sha256 3e5cf25e859b68c2 :: 162+mass over 162, leg2 N1000000 1-macro_auc: paired geometric mean [95% interval]
\newcommand{\PairedFtJcESixOnesevenMassOverOneseven}{1.02~[1.01, 1.04]}  % experiments/FIGS/data/paired_v1err/finetune/ratios.json :: ratios[112].ratio, ratios[112].ci95 :: sha256 3e5cf25e859b68c2 :: 17+mass over 17, leg2 N1000000 1-macro_auc: paired geometric mean [95% interval]
\newcommand{\PairedFtJcESixRandomOverOneeighteight}{1.046~[1.038, 1.053]}  % experiments/FIGS/data/paired_v1err/finetune/ratios.json :: ratios[113].ratio, ratios[113].ci95 :: sha256 3e5cf25e859b68c2 :: random over 188, leg2 N1000000 1-macro_auc: paired geometric mean [95% interval]
\newcommand{\PairedFtJcESixRandomOverOnesixtwo}{1.047~[1.038, 1.055]}  % experiments/FIGS/data/paired_v1err/finetune/ratios.json :: ratios[114].ratio, ratios[114].ci95 :: sha256 3e5cf25e859b68c2 :: random over 162, leg2 N1000000 1-macro_auc: paired geometric mean [95% interval]
\newcommand{\PairedFtJcESixRandomOverFourthree}{1.03~[1.02, 1.04]}  % experiments/FIGS/data/paired_v1err/finetune/ratios.json :: ratios[115].ratio, ratios[115].ci95 :: sha256 3e5cf25e859b68c2 :: random over 43, leg2 N1000000 1-macro_auc: paired geometric mean [95% interval]
\newcommand{\PairedFtJcESixRandomOverOneseven}{0.97~[0.96, 0.98]}  % experiments/FIGS/data/paired_v1err/finetune/ratios.json :: ratios[116].ratio, ratios[116].ci95 :: sha256 3e5cf25e859b68c2 :: random over 17, leg2 N1000000 1-macro_auc: paired geometric mean [95% interval]
\newcommand{\PairedFtJcESixRandomDrawTwoOverRandomDrawOne}{1.001~[0.998, 1.004]}  % experiments/FIGS/data/paired_v1err/finetune/ratios.json :: ratios[117].ratio, ratios[117].ci95 :: sha256 3e5cf25e859b68c2 :: random draw 2 over random draw 1, leg2 N1000000 1-macro_auc: paired geometric mean [95% interval]
\newcommand{\PairedFtJcESixRandomDrawThreeOverRandomDrawOne}{1.005~[1.002, 1.008]}  % experiments/FIGS/data/paired_v1err/finetune/ratios.json :: ratios[118].ratio, ratios[118].ci95 :: sha256 3e5cf25e859b68c2 :: random draw 3 over random draw 1, leg2 N1000000 1-macro_auc: paired geometric mean [95% interval]
\newcommand{\PairedFtJcESixRandomDrawThreeOverRandomDrawTwo}{1.004~[1.001, 1.007]}  % experiments/FIGS/data/paired_v1err/finetune/ratios.json :: ratios[119].ratio, ratios[119].ci95 :: sha256 3e5cf25e859b68c2 :: random draw 3 over random draw 2, leg2 N1000000 1-macro_auc: paired geometric mean [95% interval]
\newcommand{\PairedMassResMlpOnesixtwoOverOneeighteight}{0.990~[0.977, 1.003]}  % experiments/FIGS/data/paired_v1err/mass/ratios.json :: ratios[0].ratio, ratios[0].ci95 :: sha256 d877a7f087260f64 :: 162 over 188, resolution mlp sigma_eff: paired geometric mean [95% interval]
\newcommand{\PairedMassResMlpFourthreeOverOneeighteight}{0.98~[0.96, 0.99]}  % experiments/FIGS/data/paired_v1err/mass/ratios.json :: ratios[1].ratio, ratios[1].ci95 :: sha256 d877a7f087260f64 :: 43 over 188, resolution mlp sigma_eff: paired geometric mean [95% interval]
\newcommand{\PairedMassResMlpOnesevenOverOneeighteight}{0.96~[0.94, 0.98]}  % experiments/FIGS/data/paired_v1err/mass/ratios.json :: ratios[2].ratio, ratios[2].ci95 :: sha256 d877a7f087260f64 :: 17 over 188, resolution mlp sigma_eff: paired geometric mean [95% interval]
\newcommand{\PairedMassResMlpFourthreeOverOnesixtwo}{0.989~[0.978, 0.999]}  % experiments/FIGS/data/paired_v1err/mass/ratios.json :: ratios[3].ratio, ratios[3].ci95 :: sha256 d877a7f087260f64 :: 43 over 162, resolution mlp sigma_eff: paired geometric mean [95% interval]
\newcommand{\PairedMassResMlpOnesevenOverOnesixtwo}{0.97~[0.94, 0.99]}  % experiments/FIGS/data/paired_v1err/mass/ratios.json :: ratios[4].ratio, ratios[4].ci95 :: sha256 d877a7f087260f64 :: 17 over 162, resolution mlp sigma_eff: paired geometric mean [95% interval]
\newcommand{\PairedMassResMlpOnesevenOverFourthree}{0.979~[0.961, 0.997]}  % experiments/FIGS/data/paired_v1err/mass/ratios.json :: ratios[5].ratio, ratios[5].ci95 :: sha256 d877a7f087260f64 :: 17 over 43, resolution mlp sigma_eff: paired geometric mean [95% interval]
\newcommand{\PairedMassResMlpOnesixtwoMassOverOnesixtwo}{0.80~[0.79, 0.81]}  % experiments/FIGS/data/paired_v1err/mass/ratios.json :: ratios[6].ratio, ratios[6].ci95 :: sha256 d877a7f087260f64 :: 162+mass over 162, resolution mlp sigma_eff: paired geometric mean [95% interval]
\newcommand{\PairedMassResMlpOnesevenMassOverOneseven}{0.83~[0.82, 0.84]}  % experiments/FIGS/data/paired_v1err/mass/ratios.json :: ratios[7].ratio, ratios[7].ci95 :: sha256 d877a7f087260f64 :: 17+mass over 17, resolution mlp sigma_eff: paired geometric mean [95% interval]
\newcommand{\PairedMassResRidgeOnesixtwoOverOneeighteight}{0.97~[0.96, 0.98]}  % experiments/FIGS/data/paired_v1err/mass/ratios.json :: ratios[8].ratio, ratios[8].ci95 :: sha256 d877a7f087260f64 :: 162 over 188, resolution ridge sigma_eff: paired geometric mean [95% interval]
\newcommand{\PairedMassResRidgeFourthreeOverOneeighteight}{0.95~[0.94, 0.96]}  % experiments/FIGS/data/paired_v1err/mass/ratios.json :: ratios[9].ratio, ratios[9].ci95 :: sha256 d877a7f087260f64 :: 43 over 188, resolution ridge sigma_eff: paired geometric mean [95% interval]
\newcommand{\PairedMassResRidgeOnesevenOverOneeighteight}{0.88~[0.85, 0.92]}  % experiments/FIGS/data/paired_v1err/mass/ratios.json :: ratios[10].ratio, ratios[10].ci95 :: sha256 d877a7f087260f64 :: 17 over 188, resolution ridge sigma_eff: paired geometric mean [95% interval]
\newcommand{\PairedMassResRidgeFourthreeOverOnesixtwo}{0.98~[0.97, 0.99]}  % experiments/FIGS/data/paired_v1err/mass/ratios.json :: ratios[11].ratio, ratios[11].ci95 :: sha256 d877a7f087260f64 :: 43 over 162, resolution ridge sigma_eff: paired geometric mean [95% interval]
\newcommand{\PairedMassResRidgeOnesevenOverOnesixtwo}{0.91~[0.87, 0.95]}  % experiments/FIGS/data/paired_v1err/mass/ratios.json :: ratios[12].ratio, ratios[12].ci95 :: sha256 d877a7f087260f64 :: 17 over 162, resolution ridge sigma_eff: paired geometric mean [95% interval]
\newcommand{\PairedMassResRidgeOnesevenOverFourthree}{0.93~[0.89, 0.97]}  % experiments/FIGS/data/paired_v1err/mass/ratios.json :: ratios[13].ratio, ratios[13].ci95 :: sha256 d877a7f087260f64 :: 17 over 43, resolution ridge sigma_eff: paired geometric mean [95% interval]
\newcommand{\PairedMassResRidgeOnesixtwoMassOverOnesixtwo}{0.678~[0.672, 0.684]}  % experiments/FIGS/data/paired_v1err/mass/ratios.json :: ratios[14].ratio, ratios[14].ci95 :: sha256 d877a7f087260f64 :: 162+mass over 162, resolution ridge sigma_eff: paired geometric mean [95% interval]
\newcommand{\PairedMassResRidgeOnesevenMassOverOneseven}{0.72~[0.69, 0.76]}  % experiments/FIGS/data/paired_v1err/mass/ratios.json :: ratios[15].ratio, ratios[15].ci95 :: sha256 d877a7f087260f64 :: 17+mass over 17, resolution ridge sigma_eff: paired geometric mean [95% interval]
\newcommand{\VcbNSignal}{1{,}659}  % experiments/FIGS/data/probe_ladder_vcbwindow_mlp2/sall.json :: tasks.bc_vs_rest.n_signal_test :: sha256 5b8b5c7c28602f93 :: X->bc test jets in the window
\newcommand{\VcbNBackground}{19{,}244}  % experiments/FIGS/data/probe_ladder_vcbwindow_mlp2/sall.json :: tasks.bc_vs_rest.n_background_test :: sha256 5b8b5c7c28602f93 :: bq, cs, bqq and QCD test jets in the window
\newcommand{\VcbEpsSixty}{60}  % experiments/FIGS/data/probe_ladder_vcbwindow_mlp2/sall.json :: tasks.bc_vs_rest.eps_s :: sha256 5b8b5c7c28602f93 :: signal efficiency, percent
\newcommand{\VcbOmaLinearOnesixtwo}{\ensuremath{(3.17 \pm 0.05)\times10^{-2}}}  % experiments/FIGS/data/probe_ladder_vcbwindow_mlp2/sall.json :: tasks.bc_vs_rest.arms.l162-s*.linear.log1m_auc :: sha256 5b8b5c7c28602f93 :: 1-AUC = exp of it, mean +- SD over 5 seeds
\newcommand{\VcbLogOmaLinearOnesixtwo}{\ensuremath{-}3.451}  % experiments/FIGS/data/probe_ladder_vcbwindow_mlp2/sall.json :: tasks.bc_vs_rest.arms.l162-s*.linear.log1m_auc :: sha256 5b8b5c7c28602f93 :: mean over seeds
\newcommand{\VcbRejLinearSixtyOnesixtwo}{207\,\ensuremath{\pm}\,19}  % experiments/FIGS/data/probe_ladder_vcbwindow_mlp2/sall.json :: tasks.bc_vs_rest.arms.l162-s*.linear.rejection_at['0.60'].rejection :: sha256 5b8b5c7c28602f93 :: background rejection at this signal efficiency, mean +- SD over seeds
\newcommand{\VcbBkgLeftLinearSixtyOnesixtwo}{93.4}  % experiments/FIGS/data/probe_ladder_vcbwindow_mlp2/sall.json :: tasks.bc_vs_rest.arms.l162-s*.linear.rejection_at['0.60'].n_bkg_pass :: sha256 5b8b5c7c28602f93 :: background jets passing the cut, mean over seeds
\newcommand{\VcbOmaLinearOneseven}{\ensuremath{(5.21 \pm 0.23)\times10^{-2}}}  % experiments/FIGS/data/probe_ladder_vcbwindow_mlp2/sall.json :: tasks.bc_vs_rest.arms.r16q1-s*.linear.log1m_auc :: sha256 5b8b5c7c28602f93 :: 1-AUC = exp of it, mean +- SD over 5 seeds
\newcommand{\VcbLogOmaLinearOneseven}{\ensuremath{-}2.955}  % experiments/FIGS/data/probe_ladder_vcbwindow_mlp2/sall.json :: tasks.bc_vs_rest.arms.r16q1-s*.linear.log1m_auc :: sha256 5b8b5c7c28602f93 :: mean over seeds
\newcommand{\VcbRejLinearSixtyOneseven}{67\,\ensuremath{\pm}\,14}  % experiments/FIGS/data/probe_ladder_vcbwindow_mlp2/sall.json :: tasks.bc_vs_rest.arms.r16q1-s*.linear.rejection_at['0.60'].rejection :: sha256 5b8b5c7c28602f93 :: background rejection at this signal efficiency, mean +- SD over seeds
\newcommand{\VcbBkgLeftLinearSixtyOneseven}{295.0}  % experiments/FIGS/data/probe_ladder_vcbwindow_mlp2/sall.json :: tasks.bc_vs_rest.arms.r16q1-s*.linear.rejection_at['0.60'].n_bkg_pass :: sha256 5b8b5c7c28602f93 :: background jets passing the cut, mean over seeds
\newcommand{\VcbOmaRatioLinear}{1.64}  % experiments/FIGS/data/probe_ladder_vcbwindow_mlp2/sall.json :: tasks.bc_vs_rest.arms.*.linear.log1m_auc :: sha256 5b8b5c7c28602f93 :: seed mean of 1-AUC, 17 classes over 162
\newcommand{\VcbFactorLinear}{1.64}  % experiments/FIGS/data/probe_ladder_vcbwindow_mlp2/sall.json :: tasks.bc_vs_rest.arms.*.linear.log1m_auc :: sha256 5b8b5c7c28602f93 :: exp of the difference of the seed means of log(1-AUC), 17 classes over 162
\newcommand{\VcbOmaMlpOnesixtwo}{\ensuremath{(3.102 \pm 0.023)\times10^{-2}}}  % experiments/FIGS/data/probe_ladder_vcbwindow_mlp2/sall.json :: tasks.bc_vs_rest.arms.l162-s*.mlp.log1m_auc :: sha256 5b8b5c7c28602f93 :: 1-AUC = exp of it, mean +- SD over 5 seeds
\newcommand{\VcbLogOmaMlpOnesixtwo}{\ensuremath{-}3.473}  % experiments/FIGS/data/probe_ladder_vcbwindow_mlp2/sall.json :: tasks.bc_vs_rest.arms.l162-s*.mlp.log1m_auc :: sha256 5b8b5c7c28602f93 :: mean over seeds
\newcommand{\VcbRejMlpSixtyOnesixtwo}{230\,\ensuremath{\pm}\,31}  % experiments/FIGS/data/probe_ladder_vcbwindow_mlp2/sall.json :: tasks.bc_vs_rest.arms.l162-s*.mlp.rejection_at['0.60'].rejection :: sha256 5b8b5c7c28602f93 :: background rejection at this signal efficiency, mean +- SD over seeds
\newcommand{\VcbBkgLeftMlpSixtyOnesixtwo}{84.8}  % experiments/FIGS/data/probe_ladder_vcbwindow_mlp2/sall.json :: tasks.bc_vs_rest.arms.l162-s*.mlp.rejection_at['0.60'].n_bkg_pass :: sha256 5b8b5c7c28602f93 :: background jets passing the cut, mean over seeds
\newcommand{\VcbOmaMlpOneseven}{\ensuremath{(4.66 \pm 0.19)\times10^{-2}}}  % experiments/FIGS/data/probe_ladder_vcbwindow_mlp2/sall.json :: tasks.bc_vs_rest.arms.r16q1-s*.mlp.log1m_auc :: sha256 5b8b5c7c28602f93 :: 1-AUC = exp of it, mean +- SD over 5 seeds
\newcommand{\VcbLogOmaMlpOneseven}{\ensuremath{-}3.066}  % experiments/FIGS/data/probe_ladder_vcbwindow_mlp2/sall.json :: tasks.bc_vs_rest.arms.r16q1-s*.mlp.log1m_auc :: sha256 5b8b5c7c28602f93 :: mean over seeds
\newcommand{\VcbRejMlpSixtyOneseven}{78\,\ensuremath{\pm}\,13}  % experiments/FIGS/data/probe_ladder_vcbwindow_mlp2/sall.json :: tasks.bc_vs_rest.arms.r16q1-s*.mlp.rejection_at['0.60'].rejection :: sha256 5b8b5c7c28602f93 :: background rejection at this signal efficiency, mean +- SD over seeds
\newcommand{\VcbBkgLeftMlpSixtyOneseven}{250.8}  % experiments/FIGS/data/probe_ladder_vcbwindow_mlp2/sall.json :: tasks.bc_vs_rest.arms.r16q1-s*.mlp.rejection_at['0.60'].n_bkg_pass :: sha256 5b8b5c7c28602f93 :: background jets passing the cut, mean over seeds
\newcommand{\VcbOmaRatioMlp}{1.50}  % experiments/FIGS/data/probe_ladder_vcbwindow_mlp2/sall.json :: tasks.bc_vs_rest.arms.*.mlp.log1m_auc :: sha256 5b8b5c7c28602f93 :: seed mean of 1-AUC, 17 classes over 162
\newcommand{\VcbFactorMlp}{1.50}  % experiments/FIGS/data/probe_ladder_vcbwindow_mlp2/sall.json :: tasks.bc_vs_rest.arms.*.mlp.log1m_auc :: sha256 5b8b5c7c28602f93 :: exp of the difference of the seed means of log(1-AUC), 17 classes over 162
\newcommand{\AccLegOneOnesixtwoSonebEFour}{0.32273}  % experiments/FIGS/data/leg1_metrics.json :: summary.l162-s1b.N10000.accuracy_mean :: sha256 3f564e960b2c86a7 :: wave 1, superseded; 3 fine-tuning seeds
\newcommand{\AccSdLegOneOnesixtwoSonebEFour}{0.01894}  % experiments/FIGS/data/leg1_metrics.json :: summary.l162-s1b.N10000.accuracy_sd :: sha256 3f564e960b2c86a7 :: wave 1, superseded
\newcommand{\AccLegOneOnesixtwoSonebEFive}{0.53506}  % experiments/FIGS/data/leg1_metrics.json :: summary.l162-s1b.N100000.accuracy_mean :: sha256 3f564e960b2c86a7 :: wave 1, superseded; 3 fine-tuning seeds
\newcommand{\AccSdLegOneOnesixtwoSonebEFive}{0.00044}  % experiments/FIGS/data/leg1_metrics.json :: summary.l162-s1b.N100000.accuracy_sd :: sha256 3f564e960b2c86a7 :: wave 1, superseded
\newcommand{\AccLegOneOnesixtwoSonebESix}{0.54337}  % experiments/FIGS/data/leg1_metrics.json :: summary.l162-s1b.N1000000.accuracy_mean :: sha256 3f564e960b2c86a7 :: wave 1, superseded; 3 fine-tuning seeds
\newcommand{\AccSdLegOneOnesixtwoSonebESix}{0.00074}  % experiments/FIGS/data/leg1_metrics.json :: summary.l162-s1b.N1000000.accuracy_sd :: sha256 3f564e960b2c86a7 :: wave 1, superseded
\newcommand{\AccLegOneOnesevenStwoEFour}{0.22684}  % experiments/FIGS/data/leg1_metrics.json :: summary.r16q1-s2.N10000.accuracy_mean :: sha256 3f564e960b2c86a7 :: wave 1, superseded; 3 fine-tuning seeds
\newcommand{\AccSdLegOneOnesevenStwoEFour}{0.00805}  % experiments/FIGS/data/leg1_metrics.json :: summary.r16q1-s2.N10000.accuracy_sd :: sha256 3f564e960b2c86a7 :: wave 1, superseded
\newcommand{\AccLegOneOnesevenStwoEFive}{0.49650}  % experiments/FIGS/data/leg1_metrics.json :: summary.r16q1-s2.N100000.accuracy_mean :: sha256 3f564e960b2c86a7 :: wave 1, superseded; 3 fine-tuning seeds
\newcommand{\AccSdLegOneOnesevenStwoEFive}{0.00110}  % experiments/FIGS/data/leg1_metrics.json :: summary.r16q1-s2.N100000.accuracy_sd :: sha256 3f564e960b2c86a7 :: wave 1, superseded
\newcommand{\AccLegOneOnesevenStwoESix}{0.51969}  % experiments/FIGS/data/leg1_metrics.json :: summary.r16q1-s2.N1000000.accuracy_mean :: sha256 3f564e960b2c86a7 :: wave 1, superseded; 3 fine-tuning seeds
\newcommand{\AccSdLegOneOnesevenStwoESix}{0.00072}  % experiments/FIGS/data/leg1_metrics.json :: summary.r16q1-s2.N1000000.accuracy_sd :: sha256 3f564e960b2c86a7 :: wave 1, superseded
\newcommand{\AccLegOneOnesevenSthreeEFour}{0.24203}  % experiments/FIGS/data/leg1_metrics.json :: summary.r16q1-s3.N10000.accuracy_mean :: sha256 3f564e960b2c86a7 :: wave 1, superseded; 3 fine-tuning seeds
\newcommand{\AccSdLegOneOnesevenSthreeEFour}{0.00213}  % experiments/FIGS/data/leg1_metrics.json :: summary.r16q1-s3.N10000.accuracy_sd :: sha256 3f564e960b2c86a7 :: wave 1, superseded
\newcommand{\AccLegOneOnesevenSthreeEFive}{0.49313}  % experiments/FIGS/data/leg1_metrics.json :: summary.r16q1-s3.N100000.accuracy_mean :: sha256 3f564e960b2c86a7 :: wave 1, superseded; 3 fine-tuning seeds
\newcommand{\AccSdLegOneOnesevenSthreeEFive}{0.00294}  % experiments/FIGS/data/leg1_metrics.json :: summary.r16q1-s3.N100000.accuracy_sd :: sha256 3f564e960b2c86a7 :: wave 1, superseded
\newcommand{\AccLegOneOnesevenSthreeESix}{0.51804}  % experiments/FIGS/data/leg1_metrics.json :: summary.r16q1-s3.N1000000.accuracy_mean :: sha256 3f564e960b2c86a7 :: wave 1, superseded; 3 fine-tuning seeds
\newcommand{\AccSdLegOneOnesevenSthreeESix}{0.00049}  % experiments/FIGS/data/leg1_metrics.json :: summary.r16q1-s3.N1000000.accuracy_sd :: sha256 3f564e960b2c86a7 :: wave 1, superseded
\newcommand{\AccLegOneOnesevenSfourEFour}{0.24487}  % experiments/FIGS/data/leg1_metrics.json :: summary.r16q1-s4.N10000.accuracy_mean :: sha256 3f564e960b2c86a7 :: wave 1, superseded; 3 fine-tuning seeds
\newcommand{\AccSdLegOneOnesevenSfourEFour}{0.00645}  % experiments/FIGS/data/leg1_metrics.json :: summary.r16q1-s4.N10000.accuracy_sd :: sha256 3f564e960b2c86a7 :: wave 1, superseded
\newcommand{\AccLegOneOnesevenSfourEFive}{0.49510}  % experiments/FIGS/data/leg1_metrics.json :: summary.r16q1-s4.N100000.accuracy_mean :: sha256 3f564e960b2c86a7 :: wave 1, superseded; 3 fine-tuning seeds
\newcommand{\AccSdLegOneOnesevenSfourEFive}{0.00133}  % experiments/FIGS/data/leg1_metrics.json :: summary.r16q1-s4.N100000.accuracy_sd :: sha256 3f564e960b2c86a7 :: wave 1, superseded
\newcommand{\AccLegOneOnesevenSfourESix}{0.51872}  % experiments/FIGS/data/leg1_metrics.json :: summary.r16q1-s4.N1000000.accuracy_mean :: sha256 3f564e960b2c86a7 :: wave 1, superseded; 3 fine-tuning seeds
\newcommand{\AccSdLegOneOnesevenSfourESix}{0.00049}  % experiments/FIGS/data/leg1_metrics.json :: summary.r16q1-s4.N1000000.accuracy_sd :: sha256 3f564e960b2c86a7 :: wave 1, superseded
\newcommand{\AccLegOneScratchEFour}{0.31369}  % experiments/FIGS/data/leg1_metrics.json :: summary.scratch.N10000.accuracy_mean :: sha256 3f564e960b2c86a7 :: wave 1, superseded; 3 fine-tuning seeds
\newcommand{\AccSdLegOneScratchEFour}{0.00712}  % experiments/FIGS/data/leg1_metrics.json :: summary.scratch.N10000.accuracy_sd :: sha256 3f564e960b2c86a7 :: wave 1, superseded
\newcommand{\AccLegOneScratchEFive}{0.42015}  % experiments/FIGS/data/leg1_metrics.json :: summary.scratch.N100000.accuracy_mean :: sha256 3f564e960b2c86a7 :: wave 1, superseded; 3 fine-tuning seeds
\newcommand{\AccSdLegOneScratchEFive}{0.00267}  % experiments/FIGS/data/leg1_metrics.json :: summary.scratch.N100000.accuracy_sd :: sha256 3f564e960b2c86a7 :: wave 1, superseded
\newcommand{\AccLegOneScratchESix}{0.48621}  % experiments/FIGS/data/leg1_metrics.json :: summary.scratch.N1000000.accuracy_mean :: sha256 3f564e960b2c86a7 :: wave 1, superseded; 3 fine-tuning seeds
\newcommand{\AccSdLegOneScratchESix}{0.00100}  % experiments/FIGS/data/leg1_metrics.json :: summary.scratch.N1000000.accuracy_sd :: sha256 3f564e960b2c86a7 :: wave 1, superseded
\newcommand{\AccLegOneSophonPublicEFour}{0.21613}  % experiments/FIGS/data/leg1_metrics.json :: summary.sophon-public.N10000.accuracy_mean :: sha256 3f564e960b2c86a7 :: wave 1, superseded; 3 fine-tuning seeds
\newcommand{\AccSdLegOneSophonPublicEFour}{0.00734}  % experiments/FIGS/data/leg1_metrics.json :: summary.sophon-public.N10000.accuracy_sd :: sha256 3f564e960b2c86a7 :: wave 1, superseded
\newcommand{\AccLegOneSophonPublicEFive}{0.53187}  % experiments/FIGS/data/leg1_metrics.json :: summary.sophon-public.N100000.accuracy_mean :: sha256 3f564e960b2c86a7 :: wave 1, superseded; 3 fine-tuning seeds
\newcommand{\AccSdLegOneSophonPublicEFive}{0.00231}  % experiments/FIGS/data/leg1_metrics.json :: summary.sophon-public.N100000.accuracy_sd :: sha256 3f564e960b2c86a7 :: wave 1, superseded
\newcommand{\AccLegOneSophonPublicESix}{0.54273}  % experiments/FIGS/data/leg1_metrics.json :: summary.sophon-public.N1000000.accuracy_mean :: sha256 3f564e960b2c86a7 :: wave 1, superseded; 3 fine-tuning seeds
\newcommand{\AccSdLegOneSophonPublicESix}{0.00064}  % experiments/FIGS/data/leg1_metrics.json :: summary.sophon-public.N1000000.accuracy_sd :: sha256 3f564e960b2c86a7 :: wave 1, superseded
\newcommand{\AccLegTwoOnesixtwoSonebEFour}{0.78606}  % experiments/FIGS/data/leg2_metrics.json :: summary.l162-s1b.N10000.accuracy_mean :: sha256 701c58c3240274bb :: wave 1, superseded; 3 fine-tuning seeds
\newcommand{\AccSdLegTwoOnesixtwoSonebEFour}{0.00111}  % experiments/FIGS/data/leg2_metrics.json :: summary.l162-s1b.N10000.accuracy_sd :: sha256 701c58c3240274bb :: wave 1, superseded
\newcommand{\AccLegTwoOnesixtwoSonebEFive}{0.83598}  % experiments/FIGS/data/leg2_metrics.json :: summary.l162-s1b.N100000.accuracy_mean :: sha256 701c58c3240274bb :: wave 1, superseded; 3 fine-tuning seeds
\newcommand{\AccSdLegTwoOnesixtwoSonebEFive}{0.00037}  % experiments/FIGS/data/leg2_metrics.json :: summary.l162-s1b.N100000.accuracy_sd :: sha256 701c58c3240274bb :: wave 1, superseded
\newcommand{\AccLegTwoOnesixtwoSonebESix}{0.84704}  % experiments/FIGS/data/leg2_metrics.json :: summary.l162-s1b.N1000000.accuracy_mean :: sha256 701c58c3240274bb :: wave 1, superseded; 3 fine-tuning seeds
\newcommand{\AccSdLegTwoOnesixtwoSonebESix}{0.00027}  % experiments/FIGS/data/leg2_metrics.json :: summary.l162-s1b.N1000000.accuracy_sd :: sha256 701c58c3240274bb :: wave 1, superseded
\newcommand{\AccLegTwoOnesevenStwoEFour}{0.72156}  % experiments/FIGS/data/leg2_metrics.json :: summary.r16q1-s2.N10000.accuracy_mean :: sha256 701c58c3240274bb :: wave 1, superseded; 3 fine-tuning seeds
\newcommand{\AccSdLegTwoOnesevenStwoEFour}{0.01049}  % experiments/FIGS/data/leg2_metrics.json :: summary.r16q1-s2.N10000.accuracy_sd :: sha256 701c58c3240274bb :: wave 1, superseded
\newcommand{\AccLegTwoOnesevenStwoEFive}{0.82614}  % experiments/FIGS/data/leg2_metrics.json :: summary.r16q1-s2.N100000.accuracy_mean :: sha256 701c58c3240274bb :: wave 1, superseded; 3 fine-tuning seeds
\newcommand{\AccSdLegTwoOnesevenStwoEFive}{0.00058}  % experiments/FIGS/data/leg2_metrics.json :: summary.r16q1-s2.N100000.accuracy_sd :: sha256 701c58c3240274bb :: wave 1, superseded
\newcommand{\AccLegTwoOnesevenStwoESix}{0.83943}  % experiments/FIGS/data/leg2_metrics.json :: summary.r16q1-s2.N1000000.accuracy_mean :: sha256 701c58c3240274bb :: wave 1, superseded; 3 fine-tuning seeds
\newcommand{\AccSdLegTwoOnesevenStwoESix}{0.00052}  % experiments/FIGS/data/leg2_metrics.json :: summary.r16q1-s2.N1000000.accuracy_sd :: sha256 701c58c3240274bb :: wave 1, superseded
\newcommand{\AccLegTwoOnesevenSthreeEFour}{0.72103}  % experiments/FIGS/data/leg2_metrics.json :: summary.r16q1-s3.N10000.accuracy_mean :: sha256 701c58c3240274bb :: wave 1, superseded; 3 fine-tuning seeds
\newcommand{\AccSdLegTwoOnesevenSthreeEFour}{0.01651}  % experiments/FIGS/data/leg2_metrics.json :: summary.r16q1-s3.N10000.accuracy_sd :: sha256 701c58c3240274bb :: wave 1, superseded
\newcommand{\AccLegTwoOnesevenSthreeEFive}{0.82629}  % experiments/FIGS/data/leg2_metrics.json :: summary.r16q1-s3.N100000.accuracy_mean :: sha256 701c58c3240274bb :: wave 1, superseded; 3 fine-tuning seeds
\newcommand{\AccSdLegTwoOnesevenSthreeEFive}{0.00035}  % experiments/FIGS/data/leg2_metrics.json :: summary.r16q1-s3.N100000.accuracy_sd :: sha256 701c58c3240274bb :: wave 1, superseded
\newcommand{\AccLegTwoOnesevenSthreeESix}{0.83944}  % experiments/FIGS/data/leg2_metrics.json :: summary.r16q1-s3.N1000000.accuracy_mean :: sha256 701c58c3240274bb :: wave 1, superseded; 3 fine-tuning seeds
\newcommand{\AccSdLegTwoOnesevenSthreeESix}{0.00019}  % experiments/FIGS/data/leg2_metrics.json :: summary.r16q1-s3.N1000000.accuracy_sd :: sha256 701c58c3240274bb :: wave 1, superseded
\newcommand{\AccLegTwoOnesevenSfourEFour}{0.72953}  % experiments/FIGS/data/leg2_metrics.json :: summary.r16q1-s4.N10000.accuracy_mean :: sha256 701c58c3240274bb :: wave 1, superseded; 3 fine-tuning seeds
\newcommand{\AccSdLegTwoOnesevenSfourEFour}{0.01069}  % experiments/FIGS/data/leg2_metrics.json :: summary.r16q1-s4.N10000.accuracy_sd :: sha256 701c58c3240274bb :: wave 1, superseded
\newcommand{\AccLegTwoOnesevenSfourEFive}{0.82564}  % experiments/FIGS/data/leg2_metrics.json :: summary.r16q1-s4.N100000.accuracy_mean :: sha256 701c58c3240274bb :: wave 1, superseded; 3 fine-tuning seeds
\newcommand{\AccSdLegTwoOnesevenSfourEFive}{0.00041}  % experiments/FIGS/data/leg2_metrics.json :: summary.r16q1-s4.N100000.accuracy_sd :: sha256 701c58c3240274bb :: wave 1, superseded
\newcommand{\AccLegTwoOnesevenSfourESix}{0.83909}  % experiments/FIGS/data/leg2_metrics.json :: summary.r16q1-s4.N1000000.accuracy_mean :: sha256 701c58c3240274bb :: wave 1, superseded; 3 fine-tuning seeds
\newcommand{\AccSdLegTwoOnesevenSfourESix}{0.00037}  % experiments/FIGS/data/leg2_metrics.json :: summary.r16q1-s4.N1000000.accuracy_sd :: sha256 701c58c3240274bb :: wave 1, superseded
\newcommand{\AccLegTwoRefEonearmsSoneRef}{0.81770}  % experiments/FIGS/data/leg2_metrics.json :: summary.ref_e1arms-s1.ref.accuracy_mean :: sha256 701c58c3240274bb :: wave 1, superseded; 1 fine-tuning seeds
\newcommand{\AccSdLegTwoRefEonearmsSoneRef}{---}  % experiments/FIGS/data/leg2_metrics.json :: summary.ref_e1arms-s1.ref.accuracy_sd :: sha256 701c58c3240274bb :: wave 1, superseded
\newcommand{\AccLegTwoRefEonearmsStwoRef}{0.81602}  % experiments/FIGS/data/leg2_metrics.json :: summary.ref_e1arms-s2.ref.accuracy_mean :: sha256 701c58c3240274bb :: wave 1, superseded; 1 fine-tuning seeds
\newcommand{\AccSdLegTwoRefEonearmsStwoRef}{---}  % experiments/FIGS/data/leg2_metrics.json :: summary.ref_e1arms-s2.ref.accuracy_sd :: sha256 701c58c3240274bb :: wave 1, superseded
\newcommand{\AccLegTwoRefEonearmsSthreeRef}{0.81625}  % experiments/FIGS/data/leg2_metrics.json :: summary.ref_e1arms-s3.ref.accuracy_mean :: sha256 701c58c3240274bb :: wave 1, superseded; 1 fine-tuning seeds
\newcommand{\AccSdLegTwoRefEonearmsSthreeRef}{---}  % experiments/FIGS/data/leg2_metrics.json :: summary.ref_e1arms-s3.ref.accuracy_sd :: sha256 701c58c3240274bb :: wave 1, superseded
\newcommand{\AccLegTwoScratchEFour}{0.68230}  % experiments/FIGS/data/leg2_metrics.json :: summary.scratch.N10000.accuracy_mean :: sha256 701c58c3240274bb :: wave 1, superseded; 3 fine-tuning seeds
\newcommand{\AccSdLegTwoScratchEFour}{0.00282}  % experiments/FIGS/data/leg2_metrics.json :: summary.scratch.N10000.accuracy_sd :: sha256 701c58c3240274bb :: wave 1, superseded
\newcommand{\AccLegTwoScratchEFive}{0.77769}  % experiments/FIGS/data/leg2_metrics.json :: summary.scratch.N100000.accuracy_mean :: sha256 701c58c3240274bb :: wave 1, superseded; 3 fine-tuning seeds
\newcommand{\AccSdLegTwoScratchEFive}{0.00208}  % experiments/FIGS/data/leg2_metrics.json :: summary.scratch.N100000.accuracy_sd :: sha256 701c58c3240274bb :: wave 1, superseded
\newcommand{\AccLegTwoScratchESix}{0.82303}  % experiments/FIGS/data/leg2_metrics.json :: summary.scratch.N1000000.accuracy_mean :: sha256 701c58c3240274bb :: wave 1, superseded; 3 fine-tuning seeds
\newcommand{\AccSdLegTwoScratchESix}{0.00134}  % experiments/FIGS/data/leg2_metrics.json :: summary.scratch.N1000000.accuracy_sd :: sha256 701c58c3240274bb :: wave 1, superseded
\newcommand{\AccLegTwoSophonPublicEFour}{0.72739}  % experiments/FIGS/data/leg2_metrics.json :: summary.sophon-public.N10000.accuracy_mean :: sha256 701c58c3240274bb :: wave 1, superseded; 3 fine-tuning seeds
\newcommand{\AccSdLegTwoSophonPublicEFour}{0.00470}  % experiments/FIGS/data/leg2_metrics.json :: summary.sophon-public.N10000.accuracy_sd :: sha256 701c58c3240274bb :: wave 1, superseded
\newcommand{\AccLegTwoSophonPublicEFive}{0.83492}  % experiments/FIGS/data/leg2_metrics.json :: summary.sophon-public.N100000.accuracy_mean :: sha256 701c58c3240274bb :: wave 1, superseded; 3 fine-tuning seeds
\newcommand{\AccSdLegTwoSophonPublicEFive}{0.00049}  % experiments/FIGS/data/leg2_metrics.json :: summary.sophon-public.N100000.accuracy_sd :: sha256 701c58c3240274bb :: wave 1, superseded
\newcommand{\AccLegTwoSophonPublicESix}{0.84632}  % experiments/FIGS/data/leg2_metrics.json :: summary.sophon-public.N1000000.accuracy_mean :: sha256 701c58c3240274bb :: wave 1, superseded; 3 fine-tuning seeds
\newcommand{\AccSdLegTwoSophonPublicESix}{0.00031}  % experiments/FIGS/data/leg2_metrics.json :: summary.sophon-public.N1000000.accuracy_sd :: sha256 701c58c3240274bb :: wave 1, superseded
\newcommand{\FtBudgetSteps}{20{,}000}  % scripts/build_ft_budget_jobs.py :: EPOCHS * STEPS :: sha256 f2ed9faa8e3bd622 :: the largest number of optimizer steps below the full training set
\newcommand{\FtLongSteps}{80{,}000}  % scripts/build_ft_budget_jobs.py :: LONG_EPOCHS * STEPS :: sha256 f2ed9faa8e3bd622 :: the largest number of optimizer steps of the long budget at LONG_N jets (JetClass, JetClass-II)
\newcommand{\DesignNTrainSelected}{109{,}938{,}036}  % experiments/FIGS/data/v2/jc2_train/counts.json :: n_selected :: sha256 b0b5a8b63aa808c0 :: selected jets of the pretrain_train files (experiments/FT/count_jc2_train.py)
\newcommand{\FtBudgetEval}{500}  % scripts/build_ft_budget_jobs.py :: STEPS :: sha256 f2ed9faa8e3bd622 :: steps between validation passes at and above SMALL_N training jets
\newcommand{\FtBudgetEvalSmall}{50}  % scripts/build_ft_budget_jobs.py :: SMALL_STEPS :: sha256 f2ed9faa8e3bd622 :: steps between validation passes below SMALL_N training jets
\newcommand{\FtBudgetEvalSmallBelow}{\ensuremath{10^{5}}}  % scripts/build_ft_budget_jobs.py :: SMALL_N :: sha256 f2ed9faa8e3bd622 :: training jets below which the validation runs every SMALL_STEPS steps
\newcommand{\FtBudgetPatience}{10}  % scripts/build_ft_budget_jobs.py :: PATIENCE :: sha256 f2ed9faa8e3bd622 :: validation passes without improvement before training stops
\newcommand{\FtHeadMult}{50}  % scripts/build_ft_budget_jobs.py :: HEAD_MULT :: sha256 f2ed9faa8e3bd622 :: the head's learning-rate multiplier
\newcommand{\FtHeadCap}{\ensuremath{5\times10^{-3}}}  % scripts/build_ft_budget_jobs.py :: HEAD_CAP :: sha256 f2ed9faa8e3bd622 :: the head's largest learning rate
\newcommand{\FtBudgetBatch}{512}  % scripts/build_ft_budget_jobs.py :: BATCH :: sha256 f2ed9faa8e3bd622
\newcommand{\FtBudgetWeightDecay}{0.01}  % scripts/build_ft_budget_jobs.py :: --optimizer-option weight_decay (every cell) :: sha256 f2ed9faa8e3bd622
\newcommand{\FtSweepPre}{\ensuremath{10^{-4}}, \ensuremath{3\times10^{-4}} and \ensuremath{10^{-3}}}  % scripts/build_ft_budget_jobs.py :: SWEEP_PRE :: sha256 f2ed9faa8e3bd622 :: backbone rates of the sweep
\newcommand{\FtSweepScratch}{\ensuremath{3\times10^{-4}}, \ensuremath{10^{-3}} and \ensuremath{3\times10^{-3}}}  % scripts/build_ft_budget_jobs.py :: SWEEP_SCRATCH :: sha256 f2ed9faa8e3bd622 :: from-scratch rates of the sweep
\newcommand{\FtSweepN}{10{,}000}  % scripts/build_ft_budget_jobs.py :: SWEEP_N :: sha256 f2ed9faa8e3bd622 :: training jets of the sweep
\newcommand{\FtValJets}{200{,}000}  % scripts/build_ft_budget_jobs.py :: VAL_N :: sha256 f2ed9faa8e3bd622 :: the fixed validation set
\newcommand{\FtValJetsBal}{20{,}480}  % scripts/build_ft_budget_jobs.py :: VAL_SIZE[jc2b] :: sha256 f2ed9faa8e3bd622 :: the balanced JetClass-II validation set
\newcommand{\LpVfourCMax}{\ensuremath{10^{4}}}  % experiments/EVAL/probe_bal.py :: C_GRID_V4 :: sha256 a7804216e3179d9d :: largest inverse regularisation of the v4 probes
\newcommand{\JcBalTestPerClass}{5{,}000}  % experiments/FT/stage_jc2_balanced.py :: TEST_PER_CLASS :: sha256 0260341d5ed75fec
\newcommand{\JcBalValPerClass}{109}  % experiments/FT/stage_jc2_balanced.py :: VAL_PER_CLASS :: sha256 0260341d5ed75fec
\newcommand{\JcBalPairTrainPerSide}{\ensuremath{10^{5}}}  % experiments/FT/stage_jc2_balanced.py :: PAIR_TRAIN_PER_SIDE :: sha256 0260341d5ed75fec
\newcommand{\JcBalPairTestPerSide}{\ensuremath{2\times10^{5}}}  % experiments/FT/stage_jc2_balanced.py :: PAIR_TEST_PER_SIDE :: sha256 0260341d5ed75fec
\newcommand{\JcBalNShort}{11}  % experiments/FIGS/data/v2/jc2bal/census.json :: counts.ft_test :: sha256 a4379ee0817254cf :: classes with fewer than 5000 selected test jets
\newcommand{\JcBalRarestTest}{224}  % experiments/FIGS/data/v2/jc2bal/census.json :: counts.ft_test (minimum) :: sha256 a4379ee0817254cf :: label_QCD_bcs
\newcommand{\JcBalShortTrainESix}{11}  % experiments/FIGS/data/v2/jc2bal/census.json :: counts.ft_heldout :: sha256 a4379ee0817254cf :: classes with fewer held-out jets than the 1e6 quota
\newcommand{\JcBalNTest}{903{,}922}  % experiments/FIGS/data/v2/jc2bal/manifest.json :: outputs[test.parquet].rows :: sha256 9c7775be6ff56ac3 :: jets of the balanced test set
\newcommand{\FtLrPre}{\ensuremath{10^{-3}}}  % experiments/FIGS/data/v2/budget/lrsweep_metrics.json :: argmax of the summed validation metric, pretrained :: sha256 7cdc148054ca902c
\newcommand{\FtLrScratchTop}{\ensuremath{10^{-3}}}  % experiments/FIGS/data/v2/budget/lrsweep_metrics.json :: argmax of the validation metric on top, from scratch :: sha256 7cdc148054ca902c :: agrees with LR of the launched job-ft-budget-bench, -jc1, -jc2b, -scratch-full and -long-scratch-full from-scratch cells
\newcommand{\FtLrScratchQg}{\ensuremath{3\times10^{-3}}}  % experiments/FIGS/data/v2/budget/lrsweep_metrics.json :: argmax of the validation metric on quark/gluon, from scratch :: sha256 7cdc148054ca902c :: agrees with LR of the launched job-ft-budget-bench and -scratch-full from-scratch cells
\newcommand{\MlstNParamsFt}{2{,}209{,}270}  % experiments/FIGS/mlst_v4.py :: our_params(2, 17) :: sha256 fa536b4866ffd651 :: trainable parameters of our fine-tuned network, two outputs (top, quark/gluon)
\newcommand{\MlstNParamsFtKin}{2{,}207{,}950}  % experiments/FIGS/mlst_v4.py :: our_params(2, 7) :: sha256 fa536b4866ffd651 :: trainable parameters of our fine-tuned network, two outputs, the 7 kinematic inputs
\newcommand{\MlstNParamsJc}{2{,}213{,}374}  % experiments/FIGS/mlst_v4.py :: our_params(10, 17) :: sha256 fa536b4866ffd651 :: trainable parameters of our fine-tuned network, the ten JetClass classes
\newcommand{\MlstProbeRunsDone}{27}  % experiments/FIGS/data/v2/index.json :: runs with a linear-probe fit (count) :: sha256 c6c4138ee5c86423 :: of the paper's runs
\newcommand{\MlstPaperRuns}{27}  % configs/arms/v2_paper_runs.txt :: runs (count) :: sha256 b04562b14ea8cbc7
\newcommand{\MlstSdRelUnc}{50\%}  % configs/arms/v2_paper_runs.txt :: relative standard error of a standard deviation from the runs of a hierarchy level, 1/sqrt(2(n-1)) :: sha256 b04562b14ea8cbc7
\newcommand{\MlstCmpJcTwoResOverFull}{1.026\,\ensuremath{\pm}\,0.005}  % experiments/FIGS/data/v2/index.json :: linear probe jc2_188 N=1000000: (1 - macro_auc_ovr) of L162 over L188, run by run :: sha256 c6c4138ee5c86423
\newcommand{\MlstCmpJcTwoResOverFullPr}{2.6\% more errors}  % experiments/FIGS/data/v2/index.json :: linear probe jc2_188 N=1000000: (1 - macro_auc_ovr) of L162 over L188, run by run :: sha256 c6c4138ee5c86423 :: prose twin: two significant figures, direction in words
\newcommand{\MlstCmpJcTwoMidOverFull}{1.227\,\ensuremath{\pm}\,0.009}  % experiments/FIGS/data/v2/index.json :: linear probe jc2_188 N=1000000: (1 - macro_auc_ovr) of R42_Q1 over L188, run by run :: sha256 c6c4138ee5c86423
\newcommand{\MlstCmpJcTwoMidOverFullPr}{23\% more errors}  % experiments/FIGS/data/v2/index.json :: linear probe jc2_188 N=1000000: (1 - macro_auc_ovr) of R42_Q1 over L188, run by run :: sha256 c6c4138ee5c86423 :: prose twin: two significant figures, direction in words
\newcommand{\MlstCmpJcTwoCoarseOverFull}{1.557\,\ensuremath{\pm}\,0.019}  % experiments/FIGS/data/v2/index.json :: linear probe jc2_188 N=1000000: (1 - macro_auc_ovr) of R16_Q1 over L188, run by run :: sha256 c6c4138ee5c86423
\newcommand{\MlstCmpJcTwoCoarseOverFullPr}{56\% more errors}  % experiments/FIGS/data/v2/index.json :: linear probe jc2_188 N=1000000: (1 - macro_auc_ovr) of R16_Q1 over L188, run by run :: sha256 c6c4138ee5c86423 :: prose twin: two significant figures, direction in words
\newcommand{\MlstCmpJcTwoResMassOverRes}{1.022 and 1.016}  % experiments/FIGS/data/v2/index.json :: linear probe jc2_188 N=1000000: (1 - macro_auc_ovr) of L162_MASS over L162, run by run :: sha256 c6c4138ee5c86423
\newcommand{\MlstCmpJcTwoResMassOverResPr}{equal within 3\%}  % experiments/FIGS/data/v2/index.json :: linear probe jc2_188 N=1000000: (1 - macro_auc_ovr) of L162_MASS over L162, run by run :: sha256 c6c4138ee5c86423 :: prose twin: two significant figures, direction in words
\newcommand{\MlstCmpJcTwoResMassOverResPrMax}{equal within 3\%}  % experiments/FIGS/data/v2/index.json :: linear probe jc2_188 N=1000000: (1 - macro_auc_ovr) of L162_MASS over L162, run by run :: sha256 c6c4138ee5c86423 :: prose twin of the run farther from 1
\newcommand{\MlstCmpJcTwoCoarseMassOverCoarse}{1.054 and 0.977}  % experiments/FIGS/data/v2/index.json :: linear probe jc2_188 N=1000000: (1 - macro_auc_ovr) of R16_Q1_MASS over R16_Q1, run by run :: sha256 c6c4138ee5c86423
\newcommand{\MlstCmpJcTwoCoarseMassOverCoarsePr}{equal within 6\%}  % experiments/FIGS/data/v2/index.json :: linear probe jc2_188 N=1000000: (1 - macro_auc_ovr) of R16_Q1_MASS over R16_Q1, run by run :: sha256 c6c4138ee5c86423 :: prose twin: two significant figures, direction in words
\newcommand{\MlstCmpJcTwoCoarseMassOverCoarsePrMax}{equal within 6\%}  % experiments/FIGS/data/v2/index.json :: linear probe jc2_188 N=1000000: (1 - macro_auc_ovr) of R16_Q1_MASS over R16_Q1, run by run :: sha256 c6c4138ee5c86423 :: prose twin of the run farther from 1
\newcommand{\MlstCmpJcTwoResMassOverFull}{1.053 and 1.037}  % experiments/FIGS/data/v2/index.json :: linear probe jc2_188 N=1000000: (1 - macro_auc_ovr) of L162_MASS over L188, run by run :: sha256 c6c4138ee5c86423
\newcommand{\MlstCmpJcTwoResMassOverFullPr}{at least 3.7\% more errors}  % experiments/FIGS/data/v2/index.json :: linear probe jc2_188 N=1000000: (1 - macro_auc_ovr) of L162_MASS over L188, run by run :: sha256 c6c4138ee5c86423 :: prose twin: two significant figures, direction in words
\newcommand{\MlstCmpJcTwoResMassOverFullPrMax}{up to 5.3\% more errors}  % experiments/FIGS/data/v2/index.json :: linear probe jc2_188 N=1000000: (1 - macro_auc_ovr) of L162_MASS over L188, run by run :: sha256 c6c4138ee5c86423 :: prose twin of the run farther from 1
\newcommand{\MlstCmpJcTwoCoarseMassOverFull}{1.617 and 1.528}  % experiments/FIGS/data/v2/index.json :: linear probe jc2_188 N=1000000: (1 - macro_auc_ovr) of R16_Q1_MASS over L188, run by run :: sha256 c6c4138ee5c86423
\newcommand{\MlstCmpJcTwoCoarseMassOverFullPr}{at least 53\% more errors}  % experiments/FIGS/data/v2/index.json :: linear probe jc2_188 N=1000000: (1 - macro_auc_ovr) of R16_Q1_MASS over L188, run by run :: sha256 c6c4138ee5c86423 :: prose twin: two significant figures, direction in words
\newcommand{\MlstCmpJcTwoCoarseMassOverFullPrMax}{up to 62\% more errors}  % experiments/FIGS/data/v2/index.json :: linear probe jc2_188 N=1000000: (1 - macro_auc_ovr) of R16_Q1_MASS over L188, run by run :: sha256 c6c4138ee5c86423 :: prose twin of the run farther from 1
\newcommand{\MlstCmpJcTwoFzeroOverFull}{1.424 and 1.462}  % experiments/FIGS/data/v2/index.json :: linear probe jc2_188 N=1000000: (1 - macro_auc_ovr) of FLAV_F0 over L188, run by run :: sha256 c6c4138ee5c86423
\newcommand{\MlstCmpJcTwoFzeroOverFullPr}{at least 42\% more errors}  % experiments/FIGS/data/v2/index.json :: linear probe jc2_188 N=1000000: (1 - macro_auc_ovr) of FLAV_F0 over L188, run by run :: sha256 c6c4138ee5c86423 :: prose twin: two significant figures, direction in words
\newcommand{\MlstCmpJcTwoFzeroOverFullPrMax}{up to 46\% more errors}  % experiments/FIGS/data/v2/index.json :: linear probe jc2_188 N=1000000: (1 - macro_auc_ovr) of FLAV_F0 over L188, run by run :: sha256 c6c4138ee5c86423 :: prose twin of the run farther from 1
\newcommand{\MlstCmpJcTwoFoneOverFull}{1.218 and 1.212}  % experiments/FIGS/data/v2/index.json :: linear probe jc2_188 N=1000000: (1 - macro_auc_ovr) of FLAV_F1 over L188, run by run :: sha256 c6c4138ee5c86423
\newcommand{\MlstCmpJcTwoFoneOverFullPr}{at least 21\% more errors}  % experiments/FIGS/data/v2/index.json :: linear probe jc2_188 N=1000000: (1 - macro_auc_ovr) of FLAV_F1 over L188, run by run :: sha256 c6c4138ee5c86423 :: prose twin: two significant figures, direction in words
\newcommand{\MlstCmpJcTwoFoneOverFullPrMax}{up to 22\% more errors}  % experiments/FIGS/data/v2/index.json :: linear probe jc2_188 N=1000000: (1 - macro_auc_ovr) of FLAV_F1 over L188, run by run :: sha256 c6c4138ee5c86423 :: prose twin of the run farther from 1
\newcommand{\MlstCmpJcTwoFonerOverFull}{1.206 and 1.208}  % experiments/FIGS/data/v2/index.json :: linear probe jc2_188 N=1000000: (1 - macro_auc_ovr) of FLAV_F1R over L188, run by run :: sha256 c6c4138ee5c86423
\newcommand{\MlstCmpJcTwoFonerOverFullPr}{at least 21\% more errors}  % experiments/FIGS/data/v2/index.json :: linear probe jc2_188 N=1000000: (1 - macro_auc_ovr) of FLAV_F1R over L188, run by run :: sha256 c6c4138ee5c86423 :: prose twin: two significant figures, direction in words
\newcommand{\MlstCmpJcTwoFonerOverFullPrMax}{up to 21\% more errors}  % experiments/FIGS/data/v2/index.json :: linear probe jc2_188 N=1000000: (1 - macro_auc_ovr) of FLAV_F1R over L188, run by run :: sha256 c6c4138ee5c86423 :: prose twin of the run farther from 1
\newcommand{\MlstCmpJcTwoFzeroOverCoarse}{0.9275 and 0.9345}  % experiments/FIGS/data/v2/index.json :: linear probe jc2_188 N=1000000: (1 - macro_auc_ovr) of FLAV_F0 over R16_Q1, run by run :: sha256 c6c4138ee5c86423
\newcommand{\MlstCmpJcTwoFzeroOverCoarsePr}{at least 6.5\% fewer errors}  % experiments/FIGS/data/v2/index.json :: linear probe jc2_188 N=1000000: (1 - macro_auc_ovr) of FLAV_F0 over R16_Q1, run by run :: sha256 c6c4138ee5c86423 :: prose twin: two significant figures, direction in words
\newcommand{\MlstCmpJcTwoFzeroOverCoarsePrMax}{up to 7.2\% fewer errors}  % experiments/FIGS/data/v2/index.json :: linear probe jc2_188 N=1000000: (1 - macro_auc_ovr) of FLAV_F0 over R16_Q1, run by run :: sha256 c6c4138ee5c86423 :: prose twin of the run farther from 1
\newcommand{\MlstCmpJcTwoFoneOverCoarse}{0.794 and 0.775}  % experiments/FIGS/data/v2/index.json :: linear probe jc2_188 N=1000000: (1 - macro_auc_ovr) of FLAV_F1 over R16_Q1, run by run :: sha256 c6c4138ee5c86423
\newcommand{\MlstCmpJcTwoFoneOverCoarsePr}{at least 21\% fewer errors}  % experiments/FIGS/data/v2/index.json :: linear probe jc2_188 N=1000000: (1 - macro_auc_ovr) of FLAV_F1 over R16_Q1, run by run :: sha256 c6c4138ee5c86423 :: prose twin: two significant figures, direction in words
\newcommand{\MlstCmpJcTwoFoneOverCoarsePrMax}{up to 23\% fewer errors}  % experiments/FIGS/data/v2/index.json :: linear probe jc2_188 N=1000000: (1 - macro_auc_ovr) of FLAV_F1 over R16_Q1, run by run :: sha256 c6c4138ee5c86423 :: prose twin of the run farther from 1
\newcommand{\MlstCmpJcTwoFonerOverCoarse}{0.786 and 0.772}  % experiments/FIGS/data/v2/index.json :: linear probe jc2_188 N=1000000: (1 - macro_auc_ovr) of FLAV_F1R over R16_Q1, run by run :: sha256 c6c4138ee5c86423
\newcommand{\MlstCmpJcTwoFonerOverCoarsePr}{at least 21\% fewer errors}  % experiments/FIGS/data/v2/index.json :: linear probe jc2_188 N=1000000: (1 - macro_auc_ovr) of FLAV_F1R over R16_Q1, run by run :: sha256 c6c4138ee5c86423 :: prose twin: two significant figures, direction in words
\newcommand{\MlstCmpJcTwoFonerOverCoarsePrMax}{up to 23\% fewer errors}  % experiments/FIGS/data/v2/index.json :: linear probe jc2_188 N=1000000: (1 - macro_auc_ovr) of FLAV_F1R over R16_Q1, run by run :: sha256 c6c4138ee5c86423 :: prose twin of the run farther from 1
\newcommand{\MlstCmpJcTwoFoneOverFzero}{0.856 and 0.829}  % experiments/FIGS/data/v2/index.json :: linear probe jc2_188 N=1000000: (1 - macro_auc_ovr) of FLAV_F1 over FLAV_F0, run by run :: sha256 c6c4138ee5c86423
\newcommand{\MlstCmpJcTwoFoneOverFzeroPr}{at least 14\% fewer errors}  % experiments/FIGS/data/v2/index.json :: linear probe jc2_188 N=1000000: (1 - macro_auc_ovr) of FLAV_F1 over FLAV_F0, run by run :: sha256 c6c4138ee5c86423 :: prose twin: two significant figures, direction in words
\newcommand{\MlstCmpJcTwoFoneOverFzeroPrMax}{up to 17\% fewer errors}  % experiments/FIGS/data/v2/index.json :: linear probe jc2_188 N=1000000: (1 - macro_auc_ovr) of FLAV_F1 over FLAV_F0, run by run :: sha256 c6c4138ee5c86423 :: prose twin of the run farther from 1
\newcommand{\MlstCmpJcTwoFonerOverFzero}{0.847 and 0.826}  % experiments/FIGS/data/v2/index.json :: linear probe jc2_188 N=1000000: (1 - macro_auc_ovr) of FLAV_F1R over FLAV_F0, run by run :: sha256 c6c4138ee5c86423
\newcommand{\MlstCmpJcTwoFonerOverFzeroPr}{at least 15\% fewer errors}  % experiments/FIGS/data/v2/index.json :: linear probe jc2_188 N=1000000: (1 - macro_auc_ovr) of FLAV_F1R over FLAV_F0, run by run :: sha256 c6c4138ee5c86423 :: prose twin: two significant figures, direction in words
\newcommand{\MlstCmpJcTwoFonerOverFzeroPrMax}{up to 17\% fewer errors}  % experiments/FIGS/data/v2/index.json :: linear probe jc2_188 N=1000000: (1 - macro_auc_ovr) of FLAV_F1R over FLAV_F0, run by run :: sha256 c6c4138ee5c86423 :: prose twin of the run farther from 1
\newcommand{\MlstCmpJcTwoFonerOverFone}{0.9899 and 0.9966}  % experiments/FIGS/data/v2/index.json :: linear probe jc2_188 N=1000000: (1 - macro_auc_ovr) of FLAV_F1R over FLAV_F1, run by run :: sha256 c6c4138ee5c86423
\newcommand{\MlstCmpJcTwoFonerOverFonePr}{equal within 2\%}  % experiments/FIGS/data/v2/index.json :: linear probe jc2_188 N=1000000: (1 - macro_auc_ovr) of FLAV_F1R over FLAV_F1, run by run :: sha256 c6c4138ee5c86423 :: prose twin: two significant figures, direction in words
\newcommand{\MlstCmpJcTwoFonerOverFonePrMax}{equal within 2\%}  % experiments/FIGS/data/v2/index.json :: linear probe jc2_188 N=1000000: (1 - macro_auc_ovr) of FLAV_F1R over FLAV_F1, run by run :: sha256 c6c4138ee5c86423 :: prose twin of the run farther from 1
\newcommand{\MlstCmpJcTwoMidOverRes}{1.197\,\ensuremath{\pm}\,0.013}  % experiments/FIGS/data/v2/index.json :: linear probe jc2_188 N=1000000: (1 - macro_auc_ovr) of R42_Q1 over L162, run by run :: sha256 c6c4138ee5c86423
\newcommand{\MlstCmpJcTwoMidOverResPr}{20\% more errors}  % experiments/FIGS/data/v2/index.json :: linear probe jc2_188 N=1000000: (1 - macro_auc_ovr) of R42_Q1 over L162, run by run :: sha256 c6c4138ee5c86423 :: prose twin: two significant figures, direction in words
\newcommand{\MlstCmpJcTwoCoarseOverRes}{1.518\,\ensuremath{\pm}\,0.025}  % experiments/FIGS/data/v2/index.json :: linear probe jc2_188 N=1000000: (1 - macro_auc_ovr) of R16_Q1 over L162, run by run :: sha256 c6c4138ee5c86423
\newcommand{\MlstCmpJcTwoCoarseOverResPr}{52\% more errors}  % experiments/FIGS/data/v2/index.json :: linear probe jc2_188 N=1000000: (1 - macro_auc_ovr) of R16_Q1 over L162, run by run :: sha256 c6c4138ee5c86423 :: prose twin: two significant figures, direction in words
\newcommand{\MlstCmpJcTwoCoarseOverMid}{1.269\,\ensuremath{\pm}\,0.007}  % experiments/FIGS/data/v2/index.json :: linear probe jc2_188 N=1000000: (1 - macro_auc_ovr) of R16_Q1 over R42_Q1, run by run :: sha256 c6c4138ee5c86423
\newcommand{\MlstCmpJcTwoCoarseOverMidPr}{27\% more errors}  % experiments/FIGS/data/v2/index.json :: linear probe jc2_188 N=1000000: (1 - macro_auc_ovr) of R16_Q1 over R42_Q1, run by run :: sha256 c6c4138ee5c86423 :: prose twin: two significant figures, direction in words
\newcommand{\MlstCmpJcTwoRandOverFull}{1.125\,\ensuremath{\pm}\,0.008}  % experiments/FIGS/data/v2/index.json :: linear probe jc2_188 N=1000000: (1 - macro_auc_ovr) of each random partition over the L188 run of its seed, mean +- SD over the partitions :: sha256 c6c4138ee5c86423
\newcommand{\MlstCmpJcTwoRandOverFullPr}{12\% more errors}  % experiments/FIGS/data/v2/index.json :: linear probe jc2_188 N=1000000: (1 - macro_auc_ovr) of each random partition over the L188 run of its seed, mean +- SD over the partitions :: sha256 c6c4138ee5c86423 :: prose twin: two significant figures, direction in words
\newcommand{\MlstCmpJcTwoRandOverMid}{0.918\,\ensuremath{\pm}\,0.007}  % experiments/FIGS/data/v2/index.json :: linear probe jc2_188 N=1000000: (1 - macro_auc_ovr) of each random partition over the R42_Q1 run of its seed, mean +- SD over the partitions :: sha256 c6c4138ee5c86423
\newcommand{\MlstCmpJcTwoRandOverMidPr}{8.2\% fewer errors}  % experiments/FIGS/data/v2/index.json :: linear probe jc2_188 N=1000000: (1 - macro_auc_ovr) of each random partition over the R42_Q1 run of its seed, mean +- SD over the partitions :: sha256 c6c4138ee5c86423 :: prose twin: two significant figures, direction in words
\newcommand{\MlstCmpJcTwoRandOverCoarse}{0.724\,\ensuremath{\pm}\,0.008}  % experiments/FIGS/data/v2/index.json :: linear probe jc2_188 N=1000000: (1 - macro_auc_ovr) of each random partition over the R16_Q1 run of its seed, mean +- SD over the partitions :: sha256 c6c4138ee5c86423
\newcommand{\MlstCmpJcTwoRandOverCoarsePr}{28\% fewer errors}  % experiments/FIGS/data/v2/index.json :: linear probe jc2_188 N=1000000: (1 - macro_auc_ovr) of each random partition over the R16_Q1 run of its seed, mean +- SD over the partitions :: sha256 c6c4138ee5c86423 :: prose twin: two significant figures, direction in words
\newcommand{\MlstCmpJcTwoTrainedResOverFull}{1.026\,\ensuremath{\pm}\,0.005}  % experiments/FIGS/data/v2/index.json :: linear probe jc2_188 N=1000000: (1 - macro_auc_trained) of L162 over L188, run by run :: sha256 c6c4138ee5c86423
\newcommand{\MlstCmpJcTwoTrainedResOverFullPr}{2.6\% more errors}  % experiments/FIGS/data/v2/index.json :: linear probe jc2_188 N=1000000: (1 - macro_auc_trained) of L162 over L188, run by run :: sha256 c6c4138ee5c86423 :: prose twin: two significant figures, direction in words
\newcommand{\MlstCmpJcTwoTrainedMidOverFull}{1.227\,\ensuremath{\pm}\,0.009}  % experiments/FIGS/data/v2/index.json :: linear probe jc2_188 N=1000000: (1 - macro_auc_trained) of R42_Q1 over L188, run by run :: sha256 c6c4138ee5c86423
\newcommand{\MlstCmpJcTwoTrainedMidOverFullPr}{23\% more errors}  % experiments/FIGS/data/v2/index.json :: linear probe jc2_188 N=1000000: (1 - macro_auc_trained) of R42_Q1 over L188, run by run :: sha256 c6c4138ee5c86423 :: prose twin: two significant figures, direction in words
\newcommand{\MlstCmpJcTwoTrainedCoarseOverFull}{1.557\,\ensuremath{\pm}\,0.019}  % experiments/FIGS/data/v2/index.json :: linear probe jc2_188 N=1000000: (1 - macro_auc_trained) of R16_Q1 over L188, run by run :: sha256 c6c4138ee5c86423
\newcommand{\MlstCmpJcTwoTrainedCoarseOverFullPr}{56\% more errors}  % experiments/FIGS/data/v2/index.json :: linear probe jc2_188 N=1000000: (1 - macro_auc_trained) of R16_Q1 over L188, run by run :: sha256 c6c4138ee5c86423 :: prose twin: two significant figures, direction in words
\newcommand{\MlstCmpJcTwoTrainedResMassOverRes}{1.022 and 1.016}  % experiments/FIGS/data/v2/index.json :: linear probe jc2_188 N=1000000: (1 - macro_auc_trained) of L162_MASS over L162, run by run :: sha256 c6c4138ee5c86423
\newcommand{\MlstCmpJcTwoTrainedResMassOverResPr}{equal within 3\%}  % experiments/FIGS/data/v2/index.json :: linear probe jc2_188 N=1000000: (1 - macro_auc_trained) of L162_MASS over L162, run by run :: sha256 c6c4138ee5c86423 :: prose twin: two significant figures, direction in words
\newcommand{\MlstCmpJcTwoTrainedResMassOverResPrMax}{equal within 3\%}  % experiments/FIGS/data/v2/index.json :: linear probe jc2_188 N=1000000: (1 - macro_auc_trained) of L162_MASS over L162, run by run :: sha256 c6c4138ee5c86423 :: prose twin of the run farther from 1
\newcommand{\MlstCmpJcTwoTrainedCoarseMassOverCoarse}{1.054 and 0.977}  % experiments/FIGS/data/v2/index.json :: linear probe jc2_188 N=1000000: (1 - macro_auc_trained) of R16_Q1_MASS over R16_Q1, run by run :: sha256 c6c4138ee5c86423
\newcommand{\MlstCmpJcTwoTrainedCoarseMassOverCoarsePr}{equal within 6\%}  % experiments/FIGS/data/v2/index.json :: linear probe jc2_188 N=1000000: (1 - macro_auc_trained) of R16_Q1_MASS over R16_Q1, run by run :: sha256 c6c4138ee5c86423 :: prose twin: two significant figures, direction in words
\newcommand{\MlstCmpJcTwoTrainedCoarseMassOverCoarsePrMax}{equal within 6\%}  % experiments/FIGS/data/v2/index.json :: linear probe jc2_188 N=1000000: (1 - macro_auc_trained) of R16_Q1_MASS over R16_Q1, run by run :: sha256 c6c4138ee5c86423 :: prose twin of the run farther from 1
\newcommand{\MlstCmpJcTwoTrainedResMassOverFull}{1.053 and 1.037}  % experiments/FIGS/data/v2/index.json :: linear probe jc2_188 N=1000000: (1 - macro_auc_trained) of L162_MASS over L188, run by run :: sha256 c6c4138ee5c86423
\newcommand{\MlstCmpJcTwoTrainedResMassOverFullPr}{at least 3.7\% more errors}  % experiments/FIGS/data/v2/index.json :: linear probe jc2_188 N=1000000: (1 - macro_auc_trained) of L162_MASS over L188, run by run :: sha256 c6c4138ee5c86423 :: prose twin: two significant figures, direction in words
\newcommand{\MlstCmpJcTwoTrainedResMassOverFullPrMax}{up to 5.3\% more errors}  % experiments/FIGS/data/v2/index.json :: linear probe jc2_188 N=1000000: (1 - macro_auc_trained) of L162_MASS over L188, run by run :: sha256 c6c4138ee5c86423 :: prose twin of the run farther from 1
\newcommand{\MlstCmpJcTwoTrainedCoarseMassOverFull}{1.617 and 1.528}  % experiments/FIGS/data/v2/index.json :: linear probe jc2_188 N=1000000: (1 - macro_auc_trained) of R16_Q1_MASS over L188, run by run :: sha256 c6c4138ee5c86423
\newcommand{\MlstCmpJcTwoTrainedCoarseMassOverFullPr}{at least 53\% more errors}  % experiments/FIGS/data/v2/index.json :: linear probe jc2_188 N=1000000: (1 - macro_auc_trained) of R16_Q1_MASS over L188, run by run :: sha256 c6c4138ee5c86423 :: prose twin: two significant figures, direction in words
\newcommand{\MlstCmpJcTwoTrainedCoarseMassOverFullPrMax}{up to 62\% more errors}  % experiments/FIGS/data/v2/index.json :: linear probe jc2_188 N=1000000: (1 - macro_auc_trained) of R16_Q1_MASS over L188, run by run :: sha256 c6c4138ee5c86423 :: prose twin of the run farther from 1
\newcommand{\MlstCmpJcTwoTrainedFzeroOverFull}{1.424 and 1.462}  % experiments/FIGS/data/v2/index.json :: linear probe jc2_188 N=1000000: (1 - macro_auc_trained) of FLAV_F0 over L188, run by run :: sha256 c6c4138ee5c86423
\newcommand{\MlstCmpJcTwoTrainedFzeroOverFullPr}{at least 42\% more errors}  % experiments/FIGS/data/v2/index.json :: linear probe jc2_188 N=1000000: (1 - macro_auc_trained) of FLAV_F0 over L188, run by run :: sha256 c6c4138ee5c86423 :: prose twin: two significant figures, direction in words
\newcommand{\MlstCmpJcTwoTrainedFzeroOverFullPrMax}{up to 46\% more errors}  % experiments/FIGS/data/v2/index.json :: linear probe jc2_188 N=1000000: (1 - macro_auc_trained) of FLAV_F0 over L188, run by run :: sha256 c6c4138ee5c86423 :: prose twin of the run farther from 1
\newcommand{\MlstCmpJcTwoTrainedFoneOverFull}{1.218 and 1.212}  % experiments/FIGS/data/v2/index.json :: linear probe jc2_188 N=1000000: (1 - macro_auc_trained) of FLAV_F1 over L188, run by run :: sha256 c6c4138ee5c86423
\newcommand{\MlstCmpJcTwoTrainedFoneOverFullPr}{at least 21\% more errors}  % experiments/FIGS/data/v2/index.json :: linear probe jc2_188 N=1000000: (1 - macro_auc_trained) of FLAV_F1 over L188, run by run :: sha256 c6c4138ee5c86423 :: prose twin: two significant figures, direction in words
\newcommand{\MlstCmpJcTwoTrainedFoneOverFullPrMax}{up to 22\% more errors}  % experiments/FIGS/data/v2/index.json :: linear probe jc2_188 N=1000000: (1 - macro_auc_trained) of FLAV_F1 over L188, run by run :: sha256 c6c4138ee5c86423 :: prose twin of the run farther from 1
\newcommand{\MlstCmpJcTwoTrainedFonerOverFull}{1.206 and 1.208}  % experiments/FIGS/data/v2/index.json :: linear probe jc2_188 N=1000000: (1 - macro_auc_trained) of FLAV_F1R over L188, run by run :: sha256 c6c4138ee5c86423
\newcommand{\MlstCmpJcTwoTrainedFonerOverFullPr}{at least 21\% more errors}  % experiments/FIGS/data/v2/index.json :: linear probe jc2_188 N=1000000: (1 - macro_auc_trained) of FLAV_F1R over L188, run by run :: sha256 c6c4138ee5c86423 :: prose twin: two significant figures, direction in words
\newcommand{\MlstCmpJcTwoTrainedFonerOverFullPrMax}{up to 21\% more errors}  % experiments/FIGS/data/v2/index.json :: linear probe jc2_188 N=1000000: (1 - macro_auc_trained) of FLAV_F1R over L188, run by run :: sha256 c6c4138ee5c86423 :: prose twin of the run farther from 1
\newcommand{\MlstCmpJcTwoTrainedFzeroOverCoarse}{0.9275 and 0.9345}  % experiments/FIGS/data/v2/index.json :: linear probe jc2_188 N=1000000: (1 - macro_auc_trained) of FLAV_F0 over R16_Q1, run by run :: sha256 c6c4138ee5c86423
\newcommand{\MlstCmpJcTwoTrainedFzeroOverCoarsePr}{at least 6.5\% fewer errors}  % experiments/FIGS/data/v2/index.json :: linear probe jc2_188 N=1000000: (1 - macro_auc_trained) of FLAV_F0 over R16_Q1, run by run :: sha256 c6c4138ee5c86423 :: prose twin: two significant figures, direction in words
\newcommand{\MlstCmpJcTwoTrainedFzeroOverCoarsePrMax}{up to 7.2\% fewer errors}  % experiments/FIGS/data/v2/index.json :: linear probe jc2_188 N=1000000: (1 - macro_auc_trained) of FLAV_F0 over R16_Q1, run by run :: sha256 c6c4138ee5c86423 :: prose twin of the run farther from 1
\newcommand{\MlstCmpJcTwoTrainedFoneOverCoarse}{0.794 and 0.775}  % experiments/FIGS/data/v2/index.json :: linear probe jc2_188 N=1000000: (1 - macro_auc_trained) of FLAV_F1 over R16_Q1, run by run :: sha256 c6c4138ee5c86423
\newcommand{\MlstCmpJcTwoTrainedFoneOverCoarsePr}{at least 21\% fewer errors}  % experiments/FIGS/data/v2/index.json :: linear probe jc2_188 N=1000000: (1 - macro_auc_trained) of FLAV_F1 over R16_Q1, run by run :: sha256 c6c4138ee5c86423 :: prose twin: two significant figures, direction in words
\newcommand{\MlstCmpJcTwoTrainedFoneOverCoarsePrMax}{up to 23\% fewer errors}  % experiments/FIGS/data/v2/index.json :: linear probe jc2_188 N=1000000: (1 - macro_auc_trained) of FLAV_F1 over R16_Q1, run by run :: sha256 c6c4138ee5c86423 :: prose twin of the run farther from 1
\newcommand{\MlstCmpJcTwoTrainedFonerOverCoarse}{0.786 and 0.772}  % experiments/FIGS/data/v2/index.json :: linear probe jc2_188 N=1000000: (1 - macro_auc_trained) of FLAV_F1R over R16_Q1, run by run :: sha256 c6c4138ee5c86423
\newcommand{\MlstCmpJcTwoTrainedFonerOverCoarsePr}{at least 21\% fewer errors}  % experiments/FIGS/data/v2/index.json :: linear probe jc2_188 N=1000000: (1 - macro_auc_trained) of FLAV_F1R over R16_Q1, run by run :: sha256 c6c4138ee5c86423 :: prose twin: two significant figures, direction in words
\newcommand{\MlstCmpJcTwoTrainedFonerOverCoarsePrMax}{up to 23\% fewer errors}  % experiments/FIGS/data/v2/index.json :: linear probe jc2_188 N=1000000: (1 - macro_auc_trained) of FLAV_F1R over R16_Q1, run by run :: sha256 c6c4138ee5c86423 :: prose twin of the run farther from 1
\newcommand{\MlstCmpJcTwoTrainedFoneOverFzero}{0.856 and 0.829}  % experiments/FIGS/data/v2/index.json :: linear probe jc2_188 N=1000000: (1 - macro_auc_trained) of FLAV_F1 over FLAV_F0, run by run :: sha256 c6c4138ee5c86423
\newcommand{\MlstCmpJcTwoTrainedFoneOverFzeroPr}{at least 14\% fewer errors}  % experiments/FIGS/data/v2/index.json :: linear probe jc2_188 N=1000000: (1 - macro_auc_trained) of FLAV_F1 over FLAV_F0, run by run :: sha256 c6c4138ee5c86423 :: prose twin: two significant figures, direction in words
\newcommand{\MlstCmpJcTwoTrainedFoneOverFzeroPrMax}{up to 17\% fewer errors}  % experiments/FIGS/data/v2/index.json :: linear probe jc2_188 N=1000000: (1 - macro_auc_trained) of FLAV_F1 over FLAV_F0, run by run :: sha256 c6c4138ee5c86423 :: prose twin of the run farther from 1
\newcommand{\MlstCmpJcTwoTrainedFonerOverFzero}{0.847 and 0.826}  % experiments/FIGS/data/v2/index.json :: linear probe jc2_188 N=1000000: (1 - macro_auc_trained) of FLAV_F1R over FLAV_F0, run by run :: sha256 c6c4138ee5c86423
\newcommand{\MlstCmpJcTwoTrainedFonerOverFzeroPr}{at least 15\% fewer errors}  % experiments/FIGS/data/v2/index.json :: linear probe jc2_188 N=1000000: (1 - macro_auc_trained) of FLAV_F1R over FLAV_F0, run by run :: sha256 c6c4138ee5c86423 :: prose twin: two significant figures, direction in words
\newcommand{\MlstCmpJcTwoTrainedFonerOverFzeroPrMax}{up to 17\% fewer errors}  % experiments/FIGS/data/v2/index.json :: linear probe jc2_188 N=1000000: (1 - macro_auc_trained) of FLAV_F1R over FLAV_F0, run by run :: sha256 c6c4138ee5c86423 :: prose twin of the run farther from 1
\newcommand{\MlstCmpJcTwoTrainedFonerOverFone}{0.9899 and 0.9966}  % experiments/FIGS/data/v2/index.json :: linear probe jc2_188 N=1000000: (1 - macro_auc_trained) of FLAV_F1R over FLAV_F1, run by run :: sha256 c6c4138ee5c86423
\newcommand{\MlstCmpJcTwoTrainedFonerOverFonePr}{equal within 2\%}  % experiments/FIGS/data/v2/index.json :: linear probe jc2_188 N=1000000: (1 - macro_auc_trained) of FLAV_F1R over FLAV_F1, run by run :: sha256 c6c4138ee5c86423 :: prose twin: two significant figures, direction in words
\newcommand{\MlstCmpJcTwoTrainedFonerOverFonePrMax}{equal within 2\%}  % experiments/FIGS/data/v2/index.json :: linear probe jc2_188 N=1000000: (1 - macro_auc_trained) of FLAV_F1R over FLAV_F1, run by run :: sha256 c6c4138ee5c86423 :: prose twin of the run farther from 1
\newcommand{\MlstCmpJcTwoTrainedMidOverRes}{1.197\,\ensuremath{\pm}\,0.013}  % experiments/FIGS/data/v2/index.json :: linear probe jc2_188 N=1000000: (1 - macro_auc_trained) of R42_Q1 over L162, run by run :: sha256 c6c4138ee5c86423
\newcommand{\MlstCmpJcTwoTrainedMidOverResPr}{20\% more errors}  % experiments/FIGS/data/v2/index.json :: linear probe jc2_188 N=1000000: (1 - macro_auc_trained) of R42_Q1 over L162, run by run :: sha256 c6c4138ee5c86423 :: prose twin: two significant figures, direction in words
\newcommand{\MlstCmpJcTwoTrainedCoarseOverRes}{1.518\,\ensuremath{\pm}\,0.025}  % experiments/FIGS/data/v2/index.json :: linear probe jc2_188 N=1000000: (1 - macro_auc_trained) of R16_Q1 over L162, run by run :: sha256 c6c4138ee5c86423
\newcommand{\MlstCmpJcTwoTrainedCoarseOverResPr}{52\% more errors}  % experiments/FIGS/data/v2/index.json :: linear probe jc2_188 N=1000000: (1 - macro_auc_trained) of R16_Q1 over L162, run by run :: sha256 c6c4138ee5c86423 :: prose twin: two significant figures, direction in words
\newcommand{\MlstCmpJcTwoTrainedCoarseOverMid}{1.269\,\ensuremath{\pm}\,0.007}  % experiments/FIGS/data/v2/index.json :: linear probe jc2_188 N=1000000: (1 - macro_auc_trained) of R16_Q1 over R42_Q1, run by run :: sha256 c6c4138ee5c86423
\newcommand{\MlstCmpJcTwoTrainedCoarseOverMidPr}{27\% more errors}  % experiments/FIGS/data/v2/index.json :: linear probe jc2_188 N=1000000: (1 - macro_auc_trained) of R16_Q1 over R42_Q1, run by run :: sha256 c6c4138ee5c86423 :: prose twin: two significant figures, direction in words
\newcommand{\MlstCmpJcTwoTrainedRandOverFull}{1.125\,\ensuremath{\pm}\,0.008}  % experiments/FIGS/data/v2/index.json :: linear probe jc2_188 N=1000000: (1 - macro_auc_trained) of each random partition over the L188 run of its seed, mean +- SD over the partitions :: sha256 c6c4138ee5c86423
\newcommand{\MlstCmpJcTwoTrainedRandOverFullPr}{12\% more errors}  % experiments/FIGS/data/v2/index.json :: linear probe jc2_188 N=1000000: (1 - macro_auc_trained) of each random partition over the L188 run of its seed, mean +- SD over the partitions :: sha256 c6c4138ee5c86423 :: prose twin: two significant figures, direction in words
\newcommand{\MlstCmpJcTwoTrainedRandOverMid}{0.918\,\ensuremath{\pm}\,0.007}  % experiments/FIGS/data/v2/index.json :: linear probe jc2_188 N=1000000: (1 - macro_auc_trained) of each random partition over the R42_Q1 run of its seed, mean +- SD over the partitions :: sha256 c6c4138ee5c86423
\newcommand{\MlstCmpJcTwoTrainedRandOverMidPr}{8.2\% fewer errors}  % experiments/FIGS/data/v2/index.json :: linear probe jc2_188 N=1000000: (1 - macro_auc_trained) of each random partition over the R42_Q1 run of its seed, mean +- SD over the partitions :: sha256 c6c4138ee5c86423 :: prose twin: two significant figures, direction in words
\newcommand{\MlstCmpJcTwoTrainedRandOverCoarse}{0.724\,\ensuremath{\pm}\,0.008}  % experiments/FIGS/data/v2/index.json :: linear probe jc2_188 N=1000000: (1 - macro_auc_trained) of each random partition over the R16_Q1 run of its seed, mean +- SD over the partitions :: sha256 c6c4138ee5c86423
\newcommand{\MlstCmpJcTwoTrainedRandOverCoarsePr}{28\% fewer errors}  % experiments/FIGS/data/v2/index.json :: linear probe jc2_188 N=1000000: (1 - macro_auc_trained) of each random partition over the R16_Q1 run of its seed, mean +- SD over the partitions :: sha256 c6c4138ee5c86423 :: prose twin: two significant figures, direction in words
\newcommand{\MlstCmpJcTwoResResOverFull}{1.003\,\ensuremath{\pm}\,0.005}  % experiments/FIGS/data/v2/index.json :: linear probe jc2_188 N=1000000: (1 - macro_auc_resonant) of L162 over L188, run by run :: sha256 c6c4138ee5c86423
\newcommand{\MlstCmpJcTwoResResOverFullPr}{equal within 1\%}  % experiments/FIGS/data/v2/index.json :: linear probe jc2_188 N=1000000: (1 - macro_auc_resonant) of L162 over L188, run by run :: sha256 c6c4138ee5c86423 :: prose twin: two significant figures, direction in words
\newcommand{\MlstCmpJcTwoResMidOverFull}{1.262\,\ensuremath{\pm}\,0.012}  % experiments/FIGS/data/v2/index.json :: linear probe jc2_188 N=1000000: (1 - macro_auc_resonant) of R42_Q1 over L188, run by run :: sha256 c6c4138ee5c86423
\newcommand{\MlstCmpJcTwoResMidOverFullPr}{26\% more errors}  % experiments/FIGS/data/v2/index.json :: linear probe jc2_188 N=1000000: (1 - macro_auc_resonant) of R42_Q1 over L188, run by run :: sha256 c6c4138ee5c86423 :: prose twin: two significant figures, direction in words
\newcommand{\MlstCmpJcTwoResCoarseOverFull}{1.590\,\ensuremath{\pm}\,0.032}  % experiments/FIGS/data/v2/index.json :: linear probe jc2_188 N=1000000: (1 - macro_auc_resonant) of R16_Q1 over L188, run by run :: sha256 c6c4138ee5c86423
\newcommand{\MlstCmpJcTwoResCoarseOverFullPr}{59\% more errors}  % experiments/FIGS/data/v2/index.json :: linear probe jc2_188 N=1000000: (1 - macro_auc_resonant) of R16_Q1 over L188, run by run :: sha256 c6c4138ee5c86423 :: prose twin: two significant figures, direction in words
\newcommand{\MlstCmpJcTwoResResMassOverRes}{1.025 and 1.020}  % experiments/FIGS/data/v2/index.json :: linear probe jc2_188 N=1000000: (1 - macro_auc_resonant) of L162_MASS over L162, run by run :: sha256 c6c4138ee5c86423
\newcommand{\MlstCmpJcTwoResResMassOverResPr}{equal within 3\%}  % experiments/FIGS/data/v2/index.json :: linear probe jc2_188 N=1000000: (1 - macro_auc_resonant) of L162_MASS over L162, run by run :: sha256 c6c4138ee5c86423 :: prose twin: two significant figures, direction in words
\newcommand{\MlstCmpJcTwoResResMassOverResPrMax}{equal within 3\%}  % experiments/FIGS/data/v2/index.json :: linear probe jc2_188 N=1000000: (1 - macro_auc_resonant) of L162_MASS over L162, run by run :: sha256 c6c4138ee5c86423 :: prose twin of the run farther from 1
\newcommand{\MlstCmpJcTwoResCoarseMassOverCoarse}{1.08 and 0.96}  % experiments/FIGS/data/v2/index.json :: linear probe jc2_188 N=1000000: (1 - macro_auc_resonant) of R16_Q1_MASS over R16_Q1, run by run :: sha256 c6c4138ee5c86423
\newcommand{\MlstCmpJcTwoResCoarseMassOverCoarsePr}{equal within 8\%}  % experiments/FIGS/data/v2/index.json :: linear probe jc2_188 N=1000000: (1 - macro_auc_resonant) of R16_Q1_MASS over R16_Q1, run by run :: sha256 c6c4138ee5c86423 :: prose twin: two significant figures, direction in words
\newcommand{\MlstCmpJcTwoResCoarseMassOverCoarsePrMax}{equal within 8\%}  % experiments/FIGS/data/v2/index.json :: linear probe jc2_188 N=1000000: (1 - macro_auc_resonant) of R16_Q1_MASS over R16_Q1, run by run :: sha256 c6c4138ee5c86423 :: prose twin of the run farther from 1
\newcommand{\MlstCmpJcTwoResResMassOverFull}{1.033 and 1.018}  % experiments/FIGS/data/v2/index.json :: linear probe jc2_188 N=1000000: (1 - macro_auc_resonant) of L162_MASS over L188, run by run :: sha256 c6c4138ee5c86423
\newcommand{\MlstCmpJcTwoResResMassOverFullPr}{equal within 4\%}  % experiments/FIGS/data/v2/index.json :: linear probe jc2_188 N=1000000: (1 - macro_auc_resonant) of L162_MASS over L188, run by run :: sha256 c6c4138ee5c86423 :: prose twin: two significant figures, direction in words
\newcommand{\MlstCmpJcTwoResResMassOverFullPrMax}{equal within 4\%}  % experiments/FIGS/data/v2/index.json :: linear probe jc2_188 N=1000000: (1 - macro_auc_resonant) of L162_MASS over L188, run by run :: sha256 c6c4138ee5c86423 :: prose twin of the run farther from 1
\newcommand{\MlstCmpJcTwoResCoarseMassOverFull}{1.67 and 1.55}  % experiments/FIGS/data/v2/index.json :: linear probe jc2_188 N=1000000: (1 - macro_auc_resonant) of R16_Q1_MASS over L188, run by run :: sha256 c6c4138ee5c86423
\newcommand{\MlstCmpJcTwoResCoarseMassOverFullPr}{at least 55\% more errors}  % experiments/FIGS/data/v2/index.json :: linear probe jc2_188 N=1000000: (1 - macro_auc_resonant) of R16_Q1_MASS over L188, run by run :: sha256 c6c4138ee5c86423 :: prose twin: two significant figures, direction in words
\newcommand{\MlstCmpJcTwoResCoarseMassOverFullPrMax}{up to 67\% more errors}  % experiments/FIGS/data/v2/index.json :: linear probe jc2_188 N=1000000: (1 - macro_auc_resonant) of R16_Q1_MASS over L188, run by run :: sha256 c6c4138ee5c86423 :: prose twin of the run farther from 1
\newcommand{\MlstCmpJcTwoResFzeroOverFull}{1.408 and 1.449}  % experiments/FIGS/data/v2/index.json :: linear probe jc2_188 N=1000000: (1 - macro_auc_resonant) of FLAV_F0 over L188, run by run :: sha256 c6c4138ee5c86423
\newcommand{\MlstCmpJcTwoResFzeroOverFullPr}{at least 41\% more errors}  % experiments/FIGS/data/v2/index.json :: linear probe jc2_188 N=1000000: (1 - macro_auc_resonant) of FLAV_F0 over L188, run by run :: sha256 c6c4138ee5c86423 :: prose twin: two significant figures, direction in words
\newcommand{\MlstCmpJcTwoResFzeroOverFullPrMax}{up to 45\% more errors}  % experiments/FIGS/data/v2/index.json :: linear probe jc2_188 N=1000000: (1 - macro_auc_resonant) of FLAV_F0 over L188, run by run :: sha256 c6c4138ee5c86423 :: prose twin of the run farther from 1
\newcommand{\MlstCmpJcTwoResFoneOverFull}{1.210 and 1.204}  % experiments/FIGS/data/v2/index.json :: linear probe jc2_188 N=1000000: (1 - macro_auc_resonant) of FLAV_F1 over L188, run by run :: sha256 c6c4138ee5c86423
\newcommand{\MlstCmpJcTwoResFoneOverFullPr}{at least 20\% more errors}  % experiments/FIGS/data/v2/index.json :: linear probe jc2_188 N=1000000: (1 - macro_auc_resonant) of FLAV_F1 over L188, run by run :: sha256 c6c4138ee5c86423 :: prose twin: two significant figures, direction in words
\newcommand{\MlstCmpJcTwoResFoneOverFullPrMax}{up to 21\% more errors}  % experiments/FIGS/data/v2/index.json :: linear probe jc2_188 N=1000000: (1 - macro_auc_resonant) of FLAV_F1 over L188, run by run :: sha256 c6c4138ee5c86423 :: prose twin of the run farther from 1
\newcommand{\MlstCmpJcTwoResFonerOverFull}{1.187 and 1.188}  % experiments/FIGS/data/v2/index.json :: linear probe jc2_188 N=1000000: (1 - macro_auc_resonant) of FLAV_F1R over L188, run by run :: sha256 c6c4138ee5c86423
\newcommand{\MlstCmpJcTwoResFonerOverFullPr}{at least 19\% more errors}  % experiments/FIGS/data/v2/index.json :: linear probe jc2_188 N=1000000: (1 - macro_auc_resonant) of FLAV_F1R over L188, run by run :: sha256 c6c4138ee5c86423 :: prose twin: two significant figures, direction in words
\newcommand{\MlstCmpJcTwoResFonerOverFullPrMax}{up to 19\% more errors}  % experiments/FIGS/data/v2/index.json :: linear probe jc2_188 N=1000000: (1 - macro_auc_resonant) of FLAV_F1R over L188, run by run :: sha256 c6c4138ee5c86423 :: prose twin of the run farther from 1
\newcommand{\MlstCmpJcTwoResFzeroOverCoarse}{0.9069 and 0.9029}  % experiments/FIGS/data/v2/index.json :: linear probe jc2_188 N=1000000: (1 - macro_auc_resonant) of FLAV_F0 over R16_Q1, run by run :: sha256 c6c4138ee5c86423
\newcommand{\MlstCmpJcTwoResFzeroOverCoarsePr}{at least 9.3\% fewer errors}  % experiments/FIGS/data/v2/index.json :: linear probe jc2_188 N=1000000: (1 - macro_auc_resonant) of FLAV_F0 over R16_Q1, run by run :: sha256 c6c4138ee5c86423 :: prose twin: two significant figures, direction in words
\newcommand{\MlstCmpJcTwoResFzeroOverCoarsePrMax}{up to 9.7\% fewer errors}  % experiments/FIGS/data/v2/index.json :: linear probe jc2_188 N=1000000: (1 - macro_auc_resonant) of FLAV_F0 over R16_Q1, run by run :: sha256 c6c4138ee5c86423 :: prose twin of the run farther from 1
\newcommand{\MlstCmpJcTwoResFoneOverCoarse}{0.779 and 0.750}  % experiments/FIGS/data/v2/index.json :: linear probe jc2_188 N=1000000: (1 - macro_auc_resonant) of FLAV_F1 over R16_Q1, run by run :: sha256 c6c4138ee5c86423
\newcommand{\MlstCmpJcTwoResFoneOverCoarsePr}{at least 22\% fewer errors}  % experiments/FIGS/data/v2/index.json :: linear probe jc2_188 N=1000000: (1 - macro_auc_resonant) of FLAV_F1 over R16_Q1, run by run :: sha256 c6c4138ee5c86423 :: prose twin: two significant figures, direction in words
\newcommand{\MlstCmpJcTwoResFoneOverCoarsePrMax}{up to 25\% fewer errors}  % experiments/FIGS/data/v2/index.json :: linear probe jc2_188 N=1000000: (1 - macro_auc_resonant) of FLAV_F1 over R16_Q1, run by run :: sha256 c6c4138ee5c86423 :: prose twin of the run farther from 1
\newcommand{\MlstCmpJcTwoResFonerOverCoarse}{0.764 and 0.740}  % experiments/FIGS/data/v2/index.json :: linear probe jc2_188 N=1000000: (1 - macro_auc_resonant) of FLAV_F1R over R16_Q1, run by run :: sha256 c6c4138ee5c86423
\newcommand{\MlstCmpJcTwoResFonerOverCoarsePr}{at least 24\% fewer errors}  % experiments/FIGS/data/v2/index.json :: linear probe jc2_188 N=1000000: (1 - macro_auc_resonant) of FLAV_F1R over R16_Q1, run by run :: sha256 c6c4138ee5c86423 :: prose twin: two significant figures, direction in words
\newcommand{\MlstCmpJcTwoResFonerOverCoarsePrMax}{up to 26\% fewer errors}  % experiments/FIGS/data/v2/index.json :: linear probe jc2_188 N=1000000: (1 - macro_auc_resonant) of FLAV_F1R over R16_Q1, run by run :: sha256 c6c4138ee5c86423 :: prose twin of the run farther from 1
\newcommand{\MlstCmpJcTwoResFoneOverFzero}{0.859 and 0.831}  % experiments/FIGS/data/v2/index.json :: linear probe jc2_188 N=1000000: (1 - macro_auc_resonant) of FLAV_F1 over FLAV_F0, run by run :: sha256 c6c4138ee5c86423
\newcommand{\MlstCmpJcTwoResFoneOverFzeroPr}{at least 14\% fewer errors}  % experiments/FIGS/data/v2/index.json :: linear probe jc2_188 N=1000000: (1 - macro_auc_resonant) of FLAV_F1 over FLAV_F0, run by run :: sha256 c6c4138ee5c86423 :: prose twin: two significant figures, direction in words
\newcommand{\MlstCmpJcTwoResFoneOverFzeroPrMax}{up to 17\% fewer errors}  % experiments/FIGS/data/v2/index.json :: linear probe jc2_188 N=1000000: (1 - macro_auc_resonant) of FLAV_F1 over FLAV_F0, run by run :: sha256 c6c4138ee5c86423 :: prose twin of the run farther from 1
\newcommand{\MlstCmpJcTwoResFonerOverFzero}{0.843 and 0.820}  % experiments/FIGS/data/v2/index.json :: linear probe jc2_188 N=1000000: (1 - macro_auc_resonant) of FLAV_F1R over FLAV_F0, run by run :: sha256 c6c4138ee5c86423
\newcommand{\MlstCmpJcTwoResFonerOverFzeroPr}{at least 16\% fewer errors}  % experiments/FIGS/data/v2/index.json :: linear probe jc2_188 N=1000000: (1 - macro_auc_resonant) of FLAV_F1R over FLAV_F0, run by run :: sha256 c6c4138ee5c86423 :: prose twin: two significant figures, direction in words
\newcommand{\MlstCmpJcTwoResFonerOverFzeroPrMax}{up to 18\% fewer errors}  % experiments/FIGS/data/v2/index.json :: linear probe jc2_188 N=1000000: (1 - macro_auc_resonant) of FLAV_F1R over FLAV_F0, run by run :: sha256 c6c4138ee5c86423 :: prose twin of the run farther from 1
\newcommand{\MlstCmpJcTwoResFonerOverFone}{0.9806 and 0.9867}  % experiments/FIGS/data/v2/index.json :: linear probe jc2_188 N=1000000: (1 - macro_auc_resonant) of FLAV_F1R over FLAV_F1, run by run :: sha256 c6c4138ee5c86423
\newcommand{\MlstCmpJcTwoResFonerOverFonePr}{equal within 2\%}  % experiments/FIGS/data/v2/index.json :: linear probe jc2_188 N=1000000: (1 - macro_auc_resonant) of FLAV_F1R over FLAV_F1, run by run :: sha256 c6c4138ee5c86423 :: prose twin: two significant figures, direction in words
\newcommand{\MlstCmpJcTwoResFonerOverFonePrMax}{equal within 2\%}  % experiments/FIGS/data/v2/index.json :: linear probe jc2_188 N=1000000: (1 - macro_auc_resonant) of FLAV_F1R over FLAV_F1, run by run :: sha256 c6c4138ee5c86423 :: prose twin of the run farther from 1
\newcommand{\MlstCmpJcTwoResMidOverRes}{1.258\,\ensuremath{\pm}\,0.017}  % experiments/FIGS/data/v2/index.json :: linear probe jc2_188 N=1000000: (1 - macro_auc_resonant) of R42_Q1 over L162, run by run :: sha256 c6c4138ee5c86423
\newcommand{\MlstCmpJcTwoResMidOverResPr}{26\% more errors}  % experiments/FIGS/data/v2/index.json :: linear probe jc2_188 N=1000000: (1 - macro_auc_resonant) of R42_Q1 over L162, run by run :: sha256 c6c4138ee5c86423 :: prose twin: two significant figures, direction in words
\newcommand{\MlstCmpJcTwoResCoarseOverRes}{1.58\,\ensuremath{\pm}\,0.04}  % experiments/FIGS/data/v2/index.json :: linear probe jc2_188 N=1000000: (1 - macro_auc_resonant) of R16_Q1 over L162, run by run :: sha256 c6c4138ee5c86423
\newcommand{\MlstCmpJcTwoResCoarseOverResPr}{58\% more errors}  % experiments/FIGS/data/v2/index.json :: linear probe jc2_188 N=1000000: (1 - macro_auc_resonant) of R16_Q1 over L162, run by run :: sha256 c6c4138ee5c86423 :: prose twin: two significant figures, direction in words
\newcommand{\MlstCmpJcTwoResCoarseOverMid}{1.260\,\ensuremath{\pm}\,0.015}  % experiments/FIGS/data/v2/index.json :: linear probe jc2_188 N=1000000: (1 - macro_auc_resonant) of R16_Q1 over R42_Q1, run by run :: sha256 c6c4138ee5c86423
\newcommand{\MlstCmpJcTwoResCoarseOverMidPr}{26\% more errors}  % experiments/FIGS/data/v2/index.json :: linear probe jc2_188 N=1000000: (1 - macro_auc_resonant) of R16_Q1 over R42_Q1, run by run :: sha256 c6c4138ee5c86423 :: prose twin: two significant figures, direction in words
\newcommand{\MlstCmpJcTwoResRandOverFull}{1.098\,\ensuremath{\pm}\,0.008}  % experiments/FIGS/data/v2/index.json :: linear probe jc2_188 N=1000000: (1 - macro_auc_resonant) of each random partition over the L188 run of its seed, mean +- SD over the partitions :: sha256 c6c4138ee5c86423
\newcommand{\MlstCmpJcTwoResRandOverFullPr}{9.8\% more errors}  % experiments/FIGS/data/v2/index.json :: linear probe jc2_188 N=1000000: (1 - macro_auc_resonant) of each random partition over the L188 run of its seed, mean +- SD over the partitions :: sha256 c6c4138ee5c86423 :: prose twin: two significant figures, direction in words
\newcommand{\MlstCmpJcTwoResRandOverMid}{0.870\,\ensuremath{\pm}\,0.008}  % experiments/FIGS/data/v2/index.json :: linear probe jc2_188 N=1000000: (1 - macro_auc_resonant) of each random partition over the R42_Q1 run of its seed, mean +- SD over the partitions :: sha256 c6c4138ee5c86423
\newcommand{\MlstCmpJcTwoResRandOverMidPr}{13\% fewer errors}  % experiments/FIGS/data/v2/index.json :: linear probe jc2_188 N=1000000: (1 - macro_auc_resonant) of each random partition over the R42_Q1 run of its seed, mean +- SD over the partitions :: sha256 c6c4138ee5c86423 :: prose twin: two significant figures, direction in words
\newcommand{\MlstCmpJcTwoResRandOverCoarse}{0.693\,\ensuremath{\pm}\,0.011}  % experiments/FIGS/data/v2/index.json :: linear probe jc2_188 N=1000000: (1 - macro_auc_resonant) of each random partition over the R16_Q1 run of its seed, mean +- SD over the partitions :: sha256 c6c4138ee5c86423
\newcommand{\MlstCmpJcTwoResRandOverCoarsePr}{31\% fewer errors}  % experiments/FIGS/data/v2/index.json :: linear probe jc2_188 N=1000000: (1 - macro_auc_resonant) of each random partition over the R16_Q1 run of its seed, mean +- SD over the partitions :: sha256 c6c4138ee5c86423 :: prose twin: two significant figures, direction in words
\newcommand{\MlstCmpJcTwoQcdResOverFull}{1.087\,\ensuremath{\pm}\,0.006}  % experiments/FIGS/data/v2/index.json :: linear probe jc2_188 N=1000000: (1 - macro_auc_qcd) of L162 over L188, run by run :: sha256 c6c4138ee5c86423
\newcommand{\MlstCmpJcTwoQcdResOverFullPr}{8.7\% more errors}  % experiments/FIGS/data/v2/index.json :: linear probe jc2_188 N=1000000: (1 - macro_auc_qcd) of L162 over L188, run by run :: sha256 c6c4138ee5c86423 :: prose twin: two significant figures, direction in words
\newcommand{\MlstCmpJcTwoQcdMidOverFull}{1.131\,\ensuremath{\pm}\,0.007}  % experiments/FIGS/data/v2/index.json :: linear probe jc2_188 N=1000000: (1 - macro_auc_qcd) of R42_Q1 over L188, run by run :: sha256 c6c4138ee5c86423
\newcommand{\MlstCmpJcTwoQcdMidOverFullPr}{13\% more errors}  % experiments/FIGS/data/v2/index.json :: linear probe jc2_188 N=1000000: (1 - macro_auc_qcd) of R42_Q1 over L188, run by run :: sha256 c6c4138ee5c86423 :: prose twin: two significant figures, direction in words
\newcommand{\MlstCmpJcTwoQcdCoarseOverFull}{1.466\,\ensuremath{\pm}\,0.017}  % experiments/FIGS/data/v2/index.json :: linear probe jc2_188 N=1000000: (1 - macro_auc_qcd) of R16_Q1 over L188, run by run :: sha256 c6c4138ee5c86423
\newcommand{\MlstCmpJcTwoQcdCoarseOverFullPr}{47\% more errors}  % experiments/FIGS/data/v2/index.json :: linear probe jc2_188 N=1000000: (1 - macro_auc_qcd) of R16_Q1 over L188, run by run :: sha256 c6c4138ee5c86423 :: prose twin: two significant figures, direction in words
\newcommand{\MlstCmpJcTwoQcdResMassOverRes}{1.012 and 1.006}  % experiments/FIGS/data/v2/index.json :: linear probe jc2_188 N=1000000: (1 - macro_auc_qcd) of L162_MASS over L162, run by run :: sha256 c6c4138ee5c86423
\newcommand{\MlstCmpJcTwoQcdResMassOverResPr}{equal within 2\%}  % experiments/FIGS/data/v2/index.json :: linear probe jc2_188 N=1000000: (1 - macro_auc_qcd) of L162_MASS over L162, run by run :: sha256 c6c4138ee5c86423 :: prose twin: two significant figures, direction in words
\newcommand{\MlstCmpJcTwoQcdResMassOverResPrMax}{equal within 2\%}  % experiments/FIGS/data/v2/index.json :: linear probe jc2_188 N=1000000: (1 - macro_auc_qcd) of L162_MASS over L162, run by run :: sha256 c6c4138ee5c86423 :: prose twin of the run farther from 1
\newcommand{\MlstCmpJcTwoQcdCoarseMassOverCoarse}{0.991 and 1.014}  % experiments/FIGS/data/v2/index.json :: linear probe jc2_188 N=1000000: (1 - macro_auc_qcd) of R16_Q1_MASS over R16_Q1, run by run :: sha256 c6c4138ee5c86423
\newcommand{\MlstCmpJcTwoQcdCoarseMassOverCoarsePr}{equal within 2\%}  % experiments/FIGS/data/v2/index.json :: linear probe jc2_188 N=1000000: (1 - macro_auc_qcd) of R16_Q1_MASS over R16_Q1, run by run :: sha256 c6c4138ee5c86423 :: prose twin: two significant figures, direction in words
\newcommand{\MlstCmpJcTwoQcdCoarseMassOverCoarsePrMax}{equal within 2\%}  % experiments/FIGS/data/v2/index.json :: linear probe jc2_188 N=1000000: (1 - macro_auc_qcd) of R16_Q1_MASS over R16_Q1, run by run :: sha256 c6c4138ee5c86423 :: prose twin of the run farther from 1
\newcommand{\MlstCmpJcTwoQcdResMassOverFull}{1.107 and 1.092}  % experiments/FIGS/data/v2/index.json :: linear probe jc2_188 N=1000000: (1 - macro_auc_qcd) of L162_MASS over L188, run by run :: sha256 c6c4138ee5c86423
\newcommand{\MlstCmpJcTwoQcdResMassOverFullPr}{at least 9.2\% more errors}  % experiments/FIGS/data/v2/index.json :: linear probe jc2_188 N=1000000: (1 - macro_auc_qcd) of L162_MASS over L188, run by run :: sha256 c6c4138ee5c86423 :: prose twin: two significant figures, direction in words
\newcommand{\MlstCmpJcTwoQcdResMassOverFullPrMax}{up to 11\% more errors}  % experiments/FIGS/data/v2/index.json :: linear probe jc2_188 N=1000000: (1 - macro_auc_qcd) of L162_MASS over L188, run by run :: sha256 c6c4138ee5c86423 :: prose twin of the run farther from 1
\newcommand{\MlstCmpJcTwoQcdCoarseMassOverFull}{1.472 and 1.474}  % experiments/FIGS/data/v2/index.json :: linear probe jc2_188 N=1000000: (1 - macro_auc_qcd) of R16_Q1_MASS over L188, run by run :: sha256 c6c4138ee5c86423
\newcommand{\MlstCmpJcTwoQcdCoarseMassOverFullPr}{at least 47\% more errors}  % experiments/FIGS/data/v2/index.json :: linear probe jc2_188 N=1000000: (1 - macro_auc_qcd) of R16_Q1_MASS over L188, run by run :: sha256 c6c4138ee5c86423 :: prose twin: two significant figures, direction in words
\newcommand{\MlstCmpJcTwoQcdCoarseMassOverFullPrMax}{up to 47\% more errors}  % experiments/FIGS/data/v2/index.json :: linear probe jc2_188 N=1000000: (1 - macro_auc_qcd) of R16_Q1_MASS over L188, run by run :: sha256 c6c4138ee5c86423 :: prose twin of the run farther from 1
\newcommand{\MlstCmpJcTwoQcdFzeroOverFull}{1.466 and 1.497}  % experiments/FIGS/data/v2/index.json :: linear probe jc2_188 N=1000000: (1 - macro_auc_qcd) of FLAV_F0 over L188, run by run :: sha256 c6c4138ee5c86423
\newcommand{\MlstCmpJcTwoQcdFzeroOverFullPr}{at least 47\% more errors}  % experiments/FIGS/data/v2/index.json :: linear probe jc2_188 N=1000000: (1 - macro_auc_qcd) of FLAV_F0 over L188, run by run :: sha256 c6c4138ee5c86423 :: prose twin: two significant figures, direction in words
\newcommand{\MlstCmpJcTwoQcdFzeroOverFullPrMax}{up to 50\% more errors}  % experiments/FIGS/data/v2/index.json :: linear probe jc2_188 N=1000000: (1 - macro_auc_qcd) of FLAV_F0 over L188, run by run :: sha256 c6c4138ee5c86423 :: prose twin of the run farther from 1
\newcommand{\MlstCmpJcTwoQcdFoneOverFull}{1.241 and 1.235}  % experiments/FIGS/data/v2/index.json :: linear probe jc2_188 N=1000000: (1 - macro_auc_qcd) of FLAV_F1 over L188, run by run :: sha256 c6c4138ee5c86423
\newcommand{\MlstCmpJcTwoQcdFoneOverFullPr}{at least 23\% more errors}  % experiments/FIGS/data/v2/index.json :: linear probe jc2_188 N=1000000: (1 - macro_auc_qcd) of FLAV_F1 over L188, run by run :: sha256 c6c4138ee5c86423 :: prose twin: two significant figures, direction in words
\newcommand{\MlstCmpJcTwoQcdFoneOverFullPrMax}{up to 24\% more errors}  % experiments/FIGS/data/v2/index.json :: linear probe jc2_188 N=1000000: (1 - macro_auc_qcd) of FLAV_F1 over L188, run by run :: sha256 c6c4138ee5c86423 :: prose twin of the run farther from 1
\newcommand{\MlstCmpJcTwoQcdFonerOverFull}{1.260 and 1.263}  % experiments/FIGS/data/v2/index.json :: linear probe jc2_188 N=1000000: (1 - macro_auc_qcd) of FLAV_F1R over L188, run by run :: sha256 c6c4138ee5c86423
\newcommand{\MlstCmpJcTwoQcdFonerOverFullPr}{at least 26\% more errors}  % experiments/FIGS/data/v2/index.json :: linear probe jc2_188 N=1000000: (1 - macro_auc_qcd) of FLAV_F1R over L188, run by run :: sha256 c6c4138ee5c86423 :: prose twin: two significant figures, direction in words
\newcommand{\MlstCmpJcTwoQcdFonerOverFullPrMax}{up to 26\% more errors}  % experiments/FIGS/data/v2/index.json :: linear probe jc2_188 N=1000000: (1 - macro_auc_qcd) of FLAV_F1R over L188, run by run :: sha256 c6c4138ee5c86423 :: prose twin of the run farther from 1
\newcommand{\MlstCmpJcTwoQcdFzeroOverCoarse}{0.987 and 1.031}  % experiments/FIGS/data/v2/index.json :: linear probe jc2_188 N=1000000: (1 - macro_auc_qcd) of FLAV_F0 over R16_Q1, run by run :: sha256 c6c4138ee5c86423
\newcommand{\MlstCmpJcTwoQcdFzeroOverCoarsePr}{equal within 4\%}  % experiments/FIGS/data/v2/index.json :: linear probe jc2_188 N=1000000: (1 - macro_auc_qcd) of FLAV_F0 over R16_Q1, run by run :: sha256 c6c4138ee5c86423 :: prose twin: two significant figures, direction in words
\newcommand{\MlstCmpJcTwoQcdFzeroOverCoarsePrMax}{equal within 4\%}  % experiments/FIGS/data/v2/index.json :: linear probe jc2_188 N=1000000: (1 - macro_auc_qcd) of FLAV_F0 over R16_Q1, run by run :: sha256 c6c4138ee5c86423 :: prose twin of the run farther from 1
\newcommand{\MlstCmpJcTwoQcdFoneOverCoarse}{0.836 and 0.850}  % experiments/FIGS/data/v2/index.json :: linear probe jc2_188 N=1000000: (1 - macro_auc_qcd) of FLAV_F1 over R16_Q1, run by run :: sha256 c6c4138ee5c86423
\newcommand{\MlstCmpJcTwoQcdFoneOverCoarsePr}{at least 15\% fewer errors}  % experiments/FIGS/data/v2/index.json :: linear probe jc2_188 N=1000000: (1 - macro_auc_qcd) of FLAV_F1 over R16_Q1, run by run :: sha256 c6c4138ee5c86423 :: prose twin: two significant figures, direction in words
\newcommand{\MlstCmpJcTwoQcdFoneOverCoarsePrMax}{up to 16\% fewer errors}  % experiments/FIGS/data/v2/index.json :: linear probe jc2_188 N=1000000: (1 - macro_auc_qcd) of FLAV_F1 over R16_Q1, run by run :: sha256 c6c4138ee5c86423 :: prose twin of the run farther from 1
\newcommand{\MlstCmpJcTwoQcdFonerOverCoarse}{0.848 and 0.870}  % experiments/FIGS/data/v2/index.json :: linear probe jc2_188 N=1000000: (1 - macro_auc_qcd) of FLAV_F1R over R16_Q1, run by run :: sha256 c6c4138ee5c86423
\newcommand{\MlstCmpJcTwoQcdFonerOverCoarsePr}{at least 13\% fewer errors}  % experiments/FIGS/data/v2/index.json :: linear probe jc2_188 N=1000000: (1 - macro_auc_qcd) of FLAV_F1R over R16_Q1, run by run :: sha256 c6c4138ee5c86423 :: prose twin: two significant figures, direction in words
\newcommand{\MlstCmpJcTwoQcdFonerOverCoarsePrMax}{up to 15\% fewer errors}  % experiments/FIGS/data/v2/index.json :: linear probe jc2_188 N=1000000: (1 - macro_auc_qcd) of FLAV_F1R over R16_Q1, run by run :: sha256 c6c4138ee5c86423 :: prose twin of the run farther from 1
\newcommand{\MlstCmpJcTwoQcdFoneOverFzero}{0.847 and 0.825}  % experiments/FIGS/data/v2/index.json :: linear probe jc2_188 N=1000000: (1 - macro_auc_qcd) of FLAV_F1 over FLAV_F0, run by run :: sha256 c6c4138ee5c86423
\newcommand{\MlstCmpJcTwoQcdFoneOverFzeroPr}{at least 15\% fewer errors}  % experiments/FIGS/data/v2/index.json :: linear probe jc2_188 N=1000000: (1 - macro_auc_qcd) of FLAV_F1 over FLAV_F0, run by run :: sha256 c6c4138ee5c86423 :: prose twin: two significant figures, direction in words
\newcommand{\MlstCmpJcTwoQcdFoneOverFzeroPrMax}{up to 18\% fewer errors}  % experiments/FIGS/data/v2/index.json :: linear probe jc2_188 N=1000000: (1 - macro_auc_qcd) of FLAV_F1 over FLAV_F0, run by run :: sha256 c6c4138ee5c86423 :: prose twin of the run farther from 1
\newcommand{\MlstCmpJcTwoQcdFonerOverFzero}{0.859 and 0.844}  % experiments/FIGS/data/v2/index.json :: linear probe jc2_188 N=1000000: (1 - macro_auc_qcd) of FLAV_F1R over FLAV_F0, run by run :: sha256 c6c4138ee5c86423
\newcommand{\MlstCmpJcTwoQcdFonerOverFzeroPr}{at least 14\% fewer errors}  % experiments/FIGS/data/v2/index.json :: linear probe jc2_188 N=1000000: (1 - macro_auc_qcd) of FLAV_F1R over FLAV_F0, run by run :: sha256 c6c4138ee5c86423 :: prose twin: two significant figures, direction in words
\newcommand{\MlstCmpJcTwoQcdFonerOverFzeroPrMax}{up to 16\% fewer errors}  % experiments/FIGS/data/v2/index.json :: linear probe jc2_188 N=1000000: (1 - macro_auc_qcd) of FLAV_F1R over FLAV_F0, run by run :: sha256 c6c4138ee5c86423 :: prose twin of the run farther from 1
\newcommand{\MlstCmpJcTwoQcdFonerOverFone}{1.015 and 1.023}  % experiments/FIGS/data/v2/index.json :: linear probe jc2_188 N=1000000: (1 - macro_auc_qcd) of FLAV_F1R over FLAV_F1, run by run :: sha256 c6c4138ee5c86423
\newcommand{\MlstCmpJcTwoQcdFonerOverFonePr}{equal within 3\%}  % experiments/FIGS/data/v2/index.json :: linear probe jc2_188 N=1000000: (1 - macro_auc_qcd) of FLAV_F1R over FLAV_F1, run by run :: sha256 c6c4138ee5c86423 :: prose twin: two significant figures, direction in words
\newcommand{\MlstCmpJcTwoQcdFonerOverFonePrMax}{equal within 3\%}  % experiments/FIGS/data/v2/index.json :: linear probe jc2_188 N=1000000: (1 - macro_auc_qcd) of FLAV_F1R over FLAV_F1, run by run :: sha256 c6c4138ee5c86423 :: prose twin of the run farther from 1
\newcommand{\MlstCmpJcTwoQcdMidOverRes}{1.040\,\ensuremath{\pm}\,0.011}  % experiments/FIGS/data/v2/index.json :: linear probe jc2_188 N=1000000: (1 - macro_auc_qcd) of R42_Q1 over L162, run by run :: sha256 c6c4138ee5c86423
\newcommand{\MlstCmpJcTwoQcdMidOverResPr}{4.0\% more errors}  % experiments/FIGS/data/v2/index.json :: linear probe jc2_188 N=1000000: (1 - macro_auc_qcd) of R42_Q1 over L162, run by run :: sha256 c6c4138ee5c86423 :: prose twin: two significant figures, direction in words
\newcommand{\MlstCmpJcTwoQcdCoarseOverRes}{1.348\,\ensuremath{\pm}\,0.010}  % experiments/FIGS/data/v2/index.json :: linear probe jc2_188 N=1000000: (1 - macro_auc_qcd) of R16_Q1 over L162, run by run :: sha256 c6c4138ee5c86423
\newcommand{\MlstCmpJcTwoQcdCoarseOverResPr}{35\% more errors}  % experiments/FIGS/data/v2/index.json :: linear probe jc2_188 N=1000000: (1 - macro_auc_qcd) of R16_Q1 over L162, run by run :: sha256 c6c4138ee5c86423 :: prose twin: two significant figures, direction in words
\newcommand{\MlstCmpJcTwoQcdCoarseOverMid}{1.296\,\ensuremath{\pm}\,0.016}  % experiments/FIGS/data/v2/index.json :: linear probe jc2_188 N=1000000: (1 - macro_auc_qcd) of R16_Q1 over R42_Q1, run by run :: sha256 c6c4138ee5c86423
\newcommand{\MlstCmpJcTwoQcdCoarseOverMidPr}{30\% more errors}  % experiments/FIGS/data/v2/index.json :: linear probe jc2_188 N=1000000: (1 - macro_auc_qcd) of R16_Q1 over R42_Q1, run by run :: sha256 c6c4138ee5c86423 :: prose twin: two significant figures, direction in words
\newcommand{\MlstCmpJcTwoQcdRandOverFull}{1.198\,\ensuremath{\pm}\,0.011}  % experiments/FIGS/data/v2/index.json :: linear probe jc2_188 N=1000000: (1 - macro_auc_qcd) of each random partition over the L188 run of its seed, mean +- SD over the partitions :: sha256 c6c4138ee5c86423
\newcommand{\MlstCmpJcTwoQcdRandOverFullPr}{20\% more errors}  % experiments/FIGS/data/v2/index.json :: linear probe jc2_188 N=1000000: (1 - macro_auc_qcd) of each random partition over the L188 run of its seed, mean +- SD over the partitions :: sha256 c6c4138ee5c86423 :: prose twin: two significant figures, direction in words
\newcommand{\MlstCmpJcTwoQcdRandOverMid}{1.064\,\ensuremath{\pm}\,0.011}  % experiments/FIGS/data/v2/index.json :: linear probe jc2_188 N=1000000: (1 - macro_auc_qcd) of each random partition over the R42_Q1 run of its seed, mean +- SD over the partitions :: sha256 c6c4138ee5c86423
\newcommand{\MlstCmpJcTwoQcdRandOverMidPr}{6.4\% more errors}  % experiments/FIGS/data/v2/index.json :: linear probe jc2_188 N=1000000: (1 - macro_auc_qcd) of each random partition over the R42_Q1 run of its seed, mean +- SD over the partitions :: sha256 c6c4138ee5c86423 :: prose twin: two significant figures, direction in words
\newcommand{\MlstCmpJcTwoQcdRandOverCoarse}{0.818\,\ensuremath{\pm}\,0.013}  % experiments/FIGS/data/v2/index.json :: linear probe jc2_188 N=1000000: (1 - macro_auc_qcd) of each random partition over the R16_Q1 run of its seed, mean +- SD over the partitions :: sha256 c6c4138ee5c86423
\newcommand{\MlstCmpJcTwoQcdRandOverCoarsePr}{18\% fewer errors}  % experiments/FIGS/data/v2/index.json :: linear probe jc2_188 N=1000000: (1 - macro_auc_qcd) of each random partition over the R16_Q1 run of its seed, mean +- SD over the partitions :: sha256 c6c4138ee5c86423 :: prose twin: two significant figures, direction in words
\newcommand{\MlstCmpBcResOverFull}{0.998\,\ensuremath{\pm}\,0.007}  % experiments/FIGS/data/v2/index.json :: binary probe bc_vs_bq_cs: (1 - AUC) of L162 over L188, run by run :: sha256 c6c4138ee5c86423
\newcommand{\MlstCmpBcResOverFullPr}{equal within 1\%}  % experiments/FIGS/data/v2/index.json :: binary probe bc_vs_bq_cs: (1 - AUC) of L162 over L188, run by run :: sha256 c6c4138ee5c86423 :: prose twin: two significant figures, direction in words
\newcommand{\MlstCmpBcMidOverFull}{1.151\,\ensuremath{\pm}\,0.029}  % experiments/FIGS/data/v2/index.json :: binary probe bc_vs_bq_cs: (1 - AUC) of R42_Q1 over L188, run by run :: sha256 c6c4138ee5c86423
\newcommand{\MlstCmpBcMidOverFullPr}{15\% more errors}  % experiments/FIGS/data/v2/index.json :: binary probe bc_vs_bq_cs: (1 - AUC) of R42_Q1 over L188, run by run :: sha256 c6c4138ee5c86423 :: prose twin: two significant figures, direction in words
\newcommand{\MlstCmpBcCoarseOverFull}{1.95\,\ensuremath{\pm}\,0.06}  % experiments/FIGS/data/v2/index.json :: binary probe bc_vs_bq_cs: (1 - AUC) of R16_Q1 over L188, run by run :: sha256 c6c4138ee5c86423
\newcommand{\MlstCmpBcCoarseOverFullPr}{95\% more errors}  % experiments/FIGS/data/v2/index.json :: binary probe bc_vs_bq_cs: (1 - AUC) of R16_Q1 over L188, run by run :: sha256 c6c4138ee5c86423 :: prose twin: two significant figures, direction in words
\newcommand{\MlstCmpBcFoneOverFzero}{0.759 and 0.698}  % experiments/FIGS/data/v2/index.json :: binary probe bc_vs_bq_cs: (1 - AUC) of FLAV_F1 over FLAV_F0, run by run :: sha256 c6c4138ee5c86423
\newcommand{\MlstCmpBcFoneOverFzeroPr}{at least 24\% fewer errors}  % experiments/FIGS/data/v2/index.json :: binary probe bc_vs_bq_cs: (1 - AUC) of FLAV_F1 over FLAV_F0, run by run :: sha256 c6c4138ee5c86423 :: prose twin: two significant figures, direction in words
\newcommand{\MlstCmpBcFoneOverFzeroPrMax}{up to 30\% fewer errors}  % experiments/FIGS/data/v2/index.json :: binary probe bc_vs_bq_cs: (1 - AUC) of FLAV_F1 over FLAV_F0, run by run :: sha256 c6c4138ee5c86423 :: prose twin of the run farther from 1
\newcommand{\MlstCmpBcFonerOverFone}{1.008 and 0.990}  % experiments/FIGS/data/v2/index.json :: binary probe bc_vs_bq_cs: (1 - AUC) of FLAV_F1R over FLAV_F1, run by run :: sha256 c6c4138ee5c86423
\newcommand{\MlstCmpBcFonerOverFonePr}{equal within 1\%}  % experiments/FIGS/data/v2/index.json :: binary probe bc_vs_bq_cs: (1 - AUC) of FLAV_F1R over FLAV_F1, run by run :: sha256 c6c4138ee5c86423 :: prose twin: two significant figures, direction in words
\newcommand{\MlstCmpBcFonerOverFonePrMax}{equal within 1\%}  % experiments/FIGS/data/v2/index.json :: binary probe bc_vs_bq_cs: (1 - AUC) of FLAV_F1R over FLAV_F1, run by run :: sha256 c6c4138ee5c86423 :: prose twin of the run farther from 1
\newcommand{\MlstCmpBcFonerOverFzero}{0.765 and 0.691}  % experiments/FIGS/data/v2/index.json :: binary probe bc_vs_bq_cs: (1 - AUC) of FLAV_F1R over FLAV_F0, run by run :: sha256 c6c4138ee5c86423
\newcommand{\MlstCmpBcFonerOverFzeroPr}{at least 23\% fewer errors}  % experiments/FIGS/data/v2/index.json :: binary probe bc_vs_bq_cs: (1 - AUC) of FLAV_F1R over FLAV_F0, run by run :: sha256 c6c4138ee5c86423 :: prose twin: two significant figures, direction in words
\newcommand{\MlstCmpBcFonerOverFzeroPrMax}{up to 31\% fewer errors}  % experiments/FIGS/data/v2/index.json :: binary probe bc_vs_bq_cs: (1 - AUC) of FLAV_F1R over FLAV_F0, run by run :: sha256 c6c4138ee5c86423 :: prose twin of the run farther from 1
\newcommand{\MlstCmpBcFzeroOverCoarse}{0.99 and 1.13}  % experiments/FIGS/data/v2/index.json :: binary probe bc_vs_bq_cs: (1 - AUC) of FLAV_F0 over R16_Q1, run by run :: sha256 c6c4138ee5c86423
\newcommand{\MlstCmpBcFzeroOverCoarsePr}{equal within 20\%}  % experiments/FIGS/data/v2/index.json :: binary probe bc_vs_bq_cs: (1 - AUC) of FLAV_F0 over R16_Q1, run by run :: sha256 c6c4138ee5c86423 :: prose twin: two significant figures, direction in words
\newcommand{\MlstCmpBcFzeroOverCoarsePrMax}{equal within 20\%}  % experiments/FIGS/data/v2/index.json :: binary probe bc_vs_bq_cs: (1 - AUC) of FLAV_F0 over R16_Q1, run by run :: sha256 c6c4138ee5c86423 :: prose twin of the run farther from 1
\newcommand{\MlstCmpBcFoneOverCoarse}{0.752 and 0.786}  % experiments/FIGS/data/v2/index.json :: binary probe bc_vs_bq_cs: (1 - AUC) of FLAV_F1 over R16_Q1, run by run :: sha256 c6c4138ee5c86423
\newcommand{\MlstCmpBcFoneOverCoarsePr}{at least 21\% fewer errors}  % experiments/FIGS/data/v2/index.json :: binary probe bc_vs_bq_cs: (1 - AUC) of FLAV_F1 over R16_Q1, run by run :: sha256 c6c4138ee5c86423 :: prose twin: two significant figures, direction in words
\newcommand{\MlstCmpBcFoneOverCoarsePrMax}{up to 25\% fewer errors}  % experiments/FIGS/data/v2/index.json :: binary probe bc_vs_bq_cs: (1 - AUC) of FLAV_F1 over R16_Q1, run by run :: sha256 c6c4138ee5c86423 :: prose twin of the run farther from 1
\newcommand{\MlstCmpBcFonerOverCoarse}{0.758 and 0.778}  % experiments/FIGS/data/v2/index.json :: binary probe bc_vs_bq_cs: (1 - AUC) of FLAV_F1R over R16_Q1, run by run :: sha256 c6c4138ee5c86423
\newcommand{\MlstCmpBcFonerOverCoarsePr}{at least 22\% fewer errors}  % experiments/FIGS/data/v2/index.json :: binary probe bc_vs_bq_cs: (1 - AUC) of FLAV_F1R over R16_Q1, run by run :: sha256 c6c4138ee5c86423 :: prose twin: two significant figures, direction in words
\newcommand{\MlstCmpBcFonerOverCoarsePrMax}{up to 24\% fewer errors}  % experiments/FIGS/data/v2/index.json :: binary probe bc_vs_bq_cs: (1 - AUC) of FLAV_F1R over R16_Q1, run by run :: sha256 c6c4138ee5c86423 :: prose twin of the run farther from 1
\newcommand{\MlstCmpBcFzeroOverFull}{2.00 and 2.19}  % experiments/FIGS/data/v2/index.json :: binary probe bc_vs_bq_cs: (1 - AUC) of FLAV_F0 over L188, run by run :: sha256 c6c4138ee5c86423
\newcommand{\MlstCmpBcFzeroOverFullPr}{at least 2.0 times the errors}  % experiments/FIGS/data/v2/index.json :: binary probe bc_vs_bq_cs: (1 - AUC) of FLAV_F0 over L188, run by run :: sha256 c6c4138ee5c86423 :: prose twin: two significant figures, direction in words
\newcommand{\MlstCmpBcFzeroOverFullPrMax}{up to 2.2 times the errors}  % experiments/FIGS/data/v2/index.json :: binary probe bc_vs_bq_cs: (1 - AUC) of FLAV_F0 over L188, run by run :: sha256 c6c4138ee5c86423 :: prose twin of the run farther from 1
\newcommand{\MlstCmpBcFoneOverFull}{1.518 and 1.530}  % experiments/FIGS/data/v2/index.json :: binary probe bc_vs_bq_cs: (1 - AUC) of FLAV_F1 over L188, run by run :: sha256 c6c4138ee5c86423
\newcommand{\MlstCmpBcFoneOverFullPr}{at least 52\% more errors}  % experiments/FIGS/data/v2/index.json :: binary probe bc_vs_bq_cs: (1 - AUC) of FLAV_F1 over L188, run by run :: sha256 c6c4138ee5c86423 :: prose twin: two significant figures, direction in words
\newcommand{\MlstCmpBcFoneOverFullPrMax}{up to 53\% more errors}  % experiments/FIGS/data/v2/index.json :: binary probe bc_vs_bq_cs: (1 - AUC) of FLAV_F1 over L188, run by run :: sha256 c6c4138ee5c86423 :: prose twin of the run farther from 1
\newcommand{\MlstCmpBcFonerOverFull}{1.530 and 1.515}  % experiments/FIGS/data/v2/index.json :: binary probe bc_vs_bq_cs: (1 - AUC) of FLAV_F1R over L188, run by run :: sha256 c6c4138ee5c86423
\newcommand{\MlstCmpBcFonerOverFullPr}{at least 51\% more errors}  % experiments/FIGS/data/v2/index.json :: binary probe bc_vs_bq_cs: (1 - AUC) of FLAV_F1R over L188, run by run :: sha256 c6c4138ee5c86423 :: prose twin: two significant figures, direction in words
\newcommand{\MlstCmpBcFonerOverFullPrMax}{up to 53\% more errors}  % experiments/FIGS/data/v2/index.json :: binary probe bc_vs_bq_cs: (1 - AUC) of FLAV_F1R over L188, run by run :: sha256 c6c4138ee5c86423 :: prose twin of the run farther from 1
\newcommand{\MlstCmpBcRandOverFull}{1.15\,\ensuremath{\pm}\,0.07}  % experiments/FIGS/data/v2/index.json :: binary probe bc_vs_bq_cs: (1 - AUC) of each random partition over the 188-class run of its seed :: sha256 c6c4138ee5c86423
\newcommand{\MlstCmpBcRandOverFullPr}{15\% more errors}  % experiments/FIGS/data/v2/index.json :: binary probe bc_vs_bq_cs: (1 - AUC) of each random partition over the 188-class run of its seed :: sha256 c6c4138ee5c86423 :: prose twin: two significant figures, direction in words
\newcommand{\MlstPairRandMergedBcN}{2}  % experiments/FIGS/data/v2/index.json :: random partitions that merge bc_vs_bq_cs (count) :: sha256 c6c4138ee5c86423
\newcommand{\MlstPairRandApartBcN}{2}  % experiments/FIGS/data/v2/index.json :: random partitions that keep apart bc_vs_bq_cs (count) :: sha256 c6c4138ee5c86423
\newcommand{\MlstPairRandMergedOverApartBc}{1.12}  % experiments/FIGS/data/v2/index.json :: binary probe bc_vs_bq_cs: mean (1 - AUC) of the partitions that merge it over that of those that keep it apart (one run each) :: sha256 c6c4138ee5c86423
\newcommand{\MlstPairRandMergedOverApartBcPr}{12\% more errors}  % experiments/FIGS/data/v2/index.json :: binary probe bc_vs_bq_cs: mean (1 - AUC) of the partitions that merge it over that of those that keep it apart (one run each) :: sha256 c6c4138ee5c86423 :: prose twin: two significant figures, direction in words
\newcommand{\MlstCmpBbResOverFull}{1.015\,\ensuremath{\pm}\,0.027}  % experiments/FIGS/data/v2/index.json :: binary probe bb_vs_cc: (1 - AUC) of L162 over L188, run by run :: sha256 c6c4138ee5c86423
\newcommand{\MlstCmpBbResOverFullPr}{equal within 3\%}  % experiments/FIGS/data/v2/index.json :: binary probe bb_vs_cc: (1 - AUC) of L162 over L188, run by run :: sha256 c6c4138ee5c86423 :: prose twin: two significant figures, direction in words
\newcommand{\MlstCmpBbMidOverFull}{1.102\,\ensuremath{\pm}\,0.013}  % experiments/FIGS/data/v2/index.json :: binary probe bb_vs_cc: (1 - AUC) of R42_Q1 over L188, run by run :: sha256 c6c4138ee5c86423
\newcommand{\MlstCmpBbMidOverFullPr}{10\% more errors}  % experiments/FIGS/data/v2/index.json :: binary probe bb_vs_cc: (1 - AUC) of R42_Q1 over L188, run by run :: sha256 c6c4138ee5c86423 :: prose twin: two significant figures, direction in words
\newcommand{\MlstCmpBbCoarseOverFull}{4.22\,\ensuremath{\pm}\,0.29}  % experiments/FIGS/data/v2/index.json :: binary probe bb_vs_cc: (1 - AUC) of R16_Q1 over L188, run by run :: sha256 c6c4138ee5c86423
\newcommand{\MlstCmpBbCoarseOverFullPr}{4.2 times the errors}  % experiments/FIGS/data/v2/index.json :: binary probe bb_vs_cc: (1 - AUC) of R16_Q1 over L188, run by run :: sha256 c6c4138ee5c86423 :: prose twin: two significant figures, direction in words
\newcommand{\MlstCmpBbFoneOverFzero}{0.31 and 0.13}  % experiments/FIGS/data/v2/index.json :: binary probe bb_vs_cc: (1 - AUC) of FLAV_F1 over FLAV_F0, run by run :: sha256 c6c4138ee5c86423
\newcommand{\MlstCmpBbFoneOverFzeroPr}{at most 0.31 times the errors}  % experiments/FIGS/data/v2/index.json :: binary probe bb_vs_cc: (1 - AUC) of FLAV_F1 over FLAV_F0, run by run :: sha256 c6c4138ee5c86423 :: prose twin: two significant figures, direction in words
\newcommand{\MlstCmpBbFoneOverFzeroPrMax}{down to 0.13 times the errors}  % experiments/FIGS/data/v2/index.json :: binary probe bb_vs_cc: (1 - AUC) of FLAV_F1 over FLAV_F0, run by run :: sha256 c6c4138ee5c86423 :: prose twin of the run farther from 1
\newcommand{\MlstCmpBbFonerOverFone}{1.40 and 1.86}  % experiments/FIGS/data/v2/index.json :: binary probe bb_vs_cc: (1 - AUC) of FLAV_F1R over FLAV_F1, run by run :: sha256 c6c4138ee5c86423
\newcommand{\MlstCmpBbFonerOverFonePr}{at least 40\% more errors}  % experiments/FIGS/data/v2/index.json :: binary probe bb_vs_cc: (1 - AUC) of FLAV_F1R over FLAV_F1, run by run :: sha256 c6c4138ee5c86423 :: prose twin: two significant figures, direction in words
\newcommand{\MlstCmpBbFonerOverFonePrMax}{up to 86\% more errors}  % experiments/FIGS/data/v2/index.json :: binary probe bb_vs_cc: (1 - AUC) of FLAV_F1R over FLAV_F1, run by run :: sha256 c6c4138ee5c86423 :: prose twin of the run farther from 1
\newcommand{\MlstCmpBbFonerOverFzero}{0.44 and 0.24}  % experiments/FIGS/data/v2/index.json :: binary probe bb_vs_cc: (1 - AUC) of FLAV_F1R over FLAV_F0, run by run :: sha256 c6c4138ee5c86423
\newcommand{\MlstCmpBbFonerOverFzeroPr}{at most 0.44 times the errors}  % experiments/FIGS/data/v2/index.json :: binary probe bb_vs_cc: (1 - AUC) of FLAV_F1R over FLAV_F0, run by run :: sha256 c6c4138ee5c86423 :: prose twin: two significant figures, direction in words
\newcommand{\MlstCmpBbFonerOverFzeroPrMax}{down to 0.24 times the errors}  % experiments/FIGS/data/v2/index.json :: binary probe bb_vs_cc: (1 - AUC) of FLAV_F1R over FLAV_F0, run by run :: sha256 c6c4138ee5c86423 :: prose twin of the run farther from 1
\newcommand{\MlstCmpBbFzeroOverCoarse}{1.3 and 3.2}  % experiments/FIGS/data/v2/index.json :: binary probe bb_vs_cc: (1 - AUC) of FLAV_F0 over R16_Q1, run by run :: sha256 c6c4138ee5c86423
\newcommand{\MlstCmpBbFzeroOverCoarsePr}{at least 30\% more errors}  % experiments/FIGS/data/v2/index.json :: binary probe bb_vs_cc: (1 - AUC) of FLAV_F0 over R16_Q1, run by run :: sha256 c6c4138ee5c86423 :: prose twin: two significant figures, direction in words
\newcommand{\MlstCmpBbFzeroOverCoarsePrMax}{up to 3.2 times the errors}  % experiments/FIGS/data/v2/index.json :: binary probe bb_vs_cc: (1 - AUC) of FLAV_F0 over R16_Q1, run by run :: sha256 c6c4138ee5c86423 :: prose twin of the run farther from 1
\newcommand{\MlstCmpBbFoneOverCoarse}{0.4077 and 0.4126}  % experiments/FIGS/data/v2/index.json :: binary probe bb_vs_cc: (1 - AUC) of FLAV_F1 over R16_Q1, run by run :: sha256 c6c4138ee5c86423
\newcommand{\MlstCmpBbFoneOverCoarsePr}{at most 0.41 times the errors}  % experiments/FIGS/data/v2/index.json :: binary probe bb_vs_cc: (1 - AUC) of FLAV_F1 over R16_Q1, run by run :: sha256 c6c4138ee5c86423 :: prose twin: two significant figures, direction in words
\newcommand{\MlstCmpBbFoneOverCoarsePrMax}{down to 0.41 times the errors}  % experiments/FIGS/data/v2/index.json :: binary probe bb_vs_cc: (1 - AUC) of FLAV_F1 over R16_Q1, run by run :: sha256 c6c4138ee5c86423 :: prose twin of the run farther from 1
\newcommand{\MlstCmpBbFonerOverCoarse}{0.57 and 0.77}  % experiments/FIGS/data/v2/index.json :: binary probe bb_vs_cc: (1 - AUC) of FLAV_F1R over R16_Q1, run by run :: sha256 c6c4138ee5c86423
\newcommand{\MlstCmpBbFonerOverCoarsePr}{at least 23\% fewer errors}  % experiments/FIGS/data/v2/index.json :: binary probe bb_vs_cc: (1 - AUC) of FLAV_F1R over R16_Q1, run by run :: sha256 c6c4138ee5c86423 :: prose twin: two significant figures, direction in words
\newcommand{\MlstCmpBbFonerOverCoarsePrMax}{up to 43\% fewer errors}  % experiments/FIGS/data/v2/index.json :: binary probe bb_vs_cc: (1 - AUC) of FLAV_F1R over R16_Q1, run by run :: sha256 c6c4138ee5c86423 :: prose twin of the run farther from 1
\newcommand{\MlstCmpBbFzeroOverFull}{5.9 and 13.4}  % experiments/FIGS/data/v2/index.json :: binary probe bb_vs_cc: (1 - AUC) of FLAV_F0 over L188, run by run :: sha256 c6c4138ee5c86423
\newcommand{\MlstCmpBbFzeroOverFullPr}{at least 5.9 times the errors}  % experiments/FIGS/data/v2/index.json :: binary probe bb_vs_cc: (1 - AUC) of FLAV_F0 over L188, run by run :: sha256 c6c4138ee5c86423 :: prose twin: two significant figures, direction in words
\newcommand{\MlstCmpBbFzeroOverFullPrMax}{up to 13 times the errors}  % experiments/FIGS/data/v2/index.json :: binary probe bb_vs_cc: (1 - AUC) of FLAV_F0 over L188, run by run :: sha256 c6c4138ee5c86423 :: prose twin of the run farther from 1
\newcommand{\MlstCmpBbFoneOverFull}{1.837 and 1.748}  % experiments/FIGS/data/v2/index.json :: binary probe bb_vs_cc: (1 - AUC) of FLAV_F1 over L188, run by run :: sha256 c6c4138ee5c86423
\newcommand{\MlstCmpBbFoneOverFullPr}{at least 75\% more errors}  % experiments/FIGS/data/v2/index.json :: binary probe bb_vs_cc: (1 - AUC) of FLAV_F1 over L188, run by run :: sha256 c6c4138ee5c86423 :: prose twin: two significant figures, direction in words
\newcommand{\MlstCmpBbFoneOverFullPrMax}{up to 84\% more errors}  % experiments/FIGS/data/v2/index.json :: binary probe bb_vs_cc: (1 - AUC) of FLAV_F1 over L188, run by run :: sha256 c6c4138ee5c86423 :: prose twin of the run farther from 1
\newcommand{\MlstCmpBbFonerOverFull}{2.58 and 3.25}  % experiments/FIGS/data/v2/index.json :: binary probe bb_vs_cc: (1 - AUC) of FLAV_F1R over L188, run by run :: sha256 c6c4138ee5c86423
\newcommand{\MlstCmpBbFonerOverFullPr}{at least 2.6 times the errors}  % experiments/FIGS/data/v2/index.json :: binary probe bb_vs_cc: (1 - AUC) of FLAV_F1R over L188, run by run :: sha256 c6c4138ee5c86423 :: prose twin: two significant figures, direction in words
\newcommand{\MlstCmpBbFonerOverFullPrMax}{up to 3.2 times the errors}  % experiments/FIGS/data/v2/index.json :: binary probe bb_vs_cc: (1 - AUC) of FLAV_F1R over L188, run by run :: sha256 c6c4138ee5c86423 :: prose twin of the run farther from 1
\newcommand{\MlstCmpBbRandOverFull}{1.53\,\ensuremath{\pm}\,0.11}  % experiments/FIGS/data/v2/index.json :: binary probe bb_vs_cc: (1 - AUC) of each random partition over the 188-class run of its seed :: sha256 c6c4138ee5c86423
\newcommand{\MlstCmpBbRandOverFullPr}{53\% more errors}  % experiments/FIGS/data/v2/index.json :: binary probe bb_vs_cc: (1 - AUC) of each random partition over the 188-class run of its seed :: sha256 c6c4138ee5c86423 :: prose twin: two significant figures, direction in words
\newcommand{\MlstPairRandMergedBbN}{3}  % experiments/FIGS/data/v2/index.json :: random partitions that merge bb_vs_cc (count) :: sha256 c6c4138ee5c86423
\newcommand{\MlstPairRandApartBbN}{2}  % experiments/FIGS/data/v2/index.json :: random partitions that keep apart bb_vs_cc (count) :: sha256 c6c4138ee5c86423
\newcommand{\MlstPairRandMergedOverApartBb}{1.07}  % experiments/FIGS/data/v2/index.json :: binary probe bb_vs_cc: mean (1 - AUC) of the partitions that merge it over that of those that keep it apart (one run each) :: sha256 c6c4138ee5c86423
\newcommand{\MlstPairRandMergedOverApartBbPr}{7.3\% more errors}  % experiments/FIGS/data/v2/index.json :: binary probe bb_vs_cc: mean (1 - AUC) of the partitions that merge it over that of those that keep it apart (one run each) :: sha256 c6c4138ee5c86423 :: prose twin: two significant figures, direction in words
\newcommand{\MlstCmpBbqqResOverFull}{1.005\,\ensuremath{\pm}\,0.032}  % experiments/FIGS/data/v2/index.json :: binary probe bbqq_vs_ccqq: (1 - AUC) of L162 over L188, run by run :: sha256 c6c4138ee5c86423
\newcommand{\MlstCmpBbqqResOverFullPr}{equal within 4\%}  % experiments/FIGS/data/v2/index.json :: binary probe bbqq_vs_ccqq: (1 - AUC) of L162 over L188, run by run :: sha256 c6c4138ee5c86423 :: prose twin: two significant figures, direction in words
\newcommand{\MlstCmpBbqqMidOverFull}{1.117\,\ensuremath{\pm}\,0.008}  % experiments/FIGS/data/v2/index.json :: binary probe bbqq_vs_ccqq: (1 - AUC) of R42_Q1 over L188, run by run :: sha256 c6c4138ee5c86423
\newcommand{\MlstCmpBbqqMidOverFullPr}{12\% more errors}  % experiments/FIGS/data/v2/index.json :: binary probe bbqq_vs_ccqq: (1 - AUC) of R42_Q1 over L188, run by run :: sha256 c6c4138ee5c86423 :: prose twin: two significant figures, direction in words
\newcommand{\MlstCmpBbqqCoarseOverFull}{4.7\,\ensuremath{\pm}\,0.5}  % experiments/FIGS/data/v2/index.json :: binary probe bbqq_vs_ccqq: (1 - AUC) of R16_Q1 over L188, run by run :: sha256 c6c4138ee5c86423
\newcommand{\MlstCmpBbqqCoarseOverFullPr}{4.7 times the errors}  % experiments/FIGS/data/v2/index.json :: binary probe bbqq_vs_ccqq: (1 - AUC) of R16_Q1 over L188, run by run :: sha256 c6c4138ee5c86423 :: prose twin: two significant figures, direction in words
\newcommand{\MlstCmpBbqqFoneOverFzero}{0.236 and 0.154}  % experiments/FIGS/data/v2/index.json :: binary probe bbqq_vs_ccqq: (1 - AUC) of FLAV_F1 over FLAV_F0, run by run :: sha256 c6c4138ee5c86423
\newcommand{\MlstCmpBbqqFoneOverFzeroPr}{at most 0.24 times the errors}  % experiments/FIGS/data/v2/index.json :: binary probe bbqq_vs_ccqq: (1 - AUC) of FLAV_F1 over FLAV_F0, run by run :: sha256 c6c4138ee5c86423 :: prose twin: two significant figures, direction in words
\newcommand{\MlstCmpBbqqFoneOverFzeroPrMax}{down to 0.15 times the errors}  % experiments/FIGS/data/v2/index.json :: binary probe bbqq_vs_ccqq: (1 - AUC) of FLAV_F1 over FLAV_F0, run by run :: sha256 c6c4138ee5c86423 :: prose twin of the run farther from 1
\newcommand{\MlstCmpBbqqFonerOverFone}{1.505 and 1.577}  % experiments/FIGS/data/v2/index.json :: binary probe bbqq_vs_ccqq: (1 - AUC) of FLAV_F1R over FLAV_F1, run by run :: sha256 c6c4138ee5c86423
\newcommand{\MlstCmpBbqqFonerOverFonePr}{at least 50\% more errors}  % experiments/FIGS/data/v2/index.json :: binary probe bbqq_vs_ccqq: (1 - AUC) of FLAV_F1R over FLAV_F1, run by run :: sha256 c6c4138ee5c86423 :: prose twin: two significant figures, direction in words
\newcommand{\MlstCmpBbqqFonerOverFonePrMax}{up to 58\% more errors}  % experiments/FIGS/data/v2/index.json :: binary probe bbqq_vs_ccqq: (1 - AUC) of FLAV_F1R over FLAV_F1, run by run :: sha256 c6c4138ee5c86423 :: prose twin of the run farther from 1
\newcommand{\MlstCmpBbqqFonerOverFzero}{0.36 and 0.24}  % experiments/FIGS/data/v2/index.json :: binary probe bbqq_vs_ccqq: (1 - AUC) of FLAV_F1R over FLAV_F0, run by run :: sha256 c6c4138ee5c86423
\newcommand{\MlstCmpBbqqFonerOverFzeroPr}{at most 0.36 times the errors}  % experiments/FIGS/data/v2/index.json :: binary probe bbqq_vs_ccqq: (1 - AUC) of FLAV_F1R over FLAV_F0, run by run :: sha256 c6c4138ee5c86423 :: prose twin: two significant figures, direction in words
\newcommand{\MlstCmpBbqqFonerOverFzeroPrMax}{down to 0.24 times the errors}  % experiments/FIGS/data/v2/index.json :: binary probe bbqq_vs_ccqq: (1 - AUC) of FLAV_F1R over FLAV_F0, run by run :: sha256 c6c4138ee5c86423 :: prose twin of the run farther from 1
\newcommand{\MlstCmpBbqqFzeroOverCoarse}{1.51 and 1.79}  % experiments/FIGS/data/v2/index.json :: binary probe bbqq_vs_ccqq: (1 - AUC) of FLAV_F0 over R16_Q1, run by run :: sha256 c6c4138ee5c86423
\newcommand{\MlstCmpBbqqFzeroOverCoarsePr}{at least 51\% more errors}  % experiments/FIGS/data/v2/index.json :: binary probe bbqq_vs_ccqq: (1 - AUC) of FLAV_F0 over R16_Q1, run by run :: sha256 c6c4138ee5c86423 :: prose twin: two significant figures, direction in words
\newcommand{\MlstCmpBbqqFzeroOverCoarsePrMax}{up to 79\% more errors}  % experiments/FIGS/data/v2/index.json :: binary probe bbqq_vs_ccqq: (1 - AUC) of FLAV_F0 over R16_Q1, run by run :: sha256 c6c4138ee5c86423 :: prose twin of the run farther from 1
\newcommand{\MlstCmpBbqqFoneOverCoarse}{0.356 and 0.276}  % experiments/FIGS/data/v2/index.json :: binary probe bbqq_vs_ccqq: (1 - AUC) of FLAV_F1 over R16_Q1, run by run :: sha256 c6c4138ee5c86423
\newcommand{\MlstCmpBbqqFoneOverCoarsePr}{at most 0.36 times the errors}  % experiments/FIGS/data/v2/index.json :: binary probe bbqq_vs_ccqq: (1 - AUC) of FLAV_F1 over R16_Q1, run by run :: sha256 c6c4138ee5c86423 :: prose twin: two significant figures, direction in words
\newcommand{\MlstCmpBbqqFoneOverCoarsePrMax}{down to 0.28 times the errors}  % experiments/FIGS/data/v2/index.json :: binary probe bbqq_vs_ccqq: (1 - AUC) of FLAV_F1 over R16_Q1, run by run :: sha256 c6c4138ee5c86423 :: prose twin of the run farther from 1
\newcommand{\MlstCmpBbqqFonerOverCoarse}{0.54 and 0.44}  % experiments/FIGS/data/v2/index.json :: binary probe bbqq_vs_ccqq: (1 - AUC) of FLAV_F1R over R16_Q1, run by run :: sha256 c6c4138ee5c86423
\newcommand{\MlstCmpBbqqFonerOverCoarsePr}{at least 46\% fewer errors}  % experiments/FIGS/data/v2/index.json :: binary probe bbqq_vs_ccqq: (1 - AUC) of FLAV_F1R over R16_Q1, run by run :: sha256 c6c4138ee5c86423 :: prose twin: two significant figures, direction in words
\newcommand{\MlstCmpBbqqFonerOverCoarsePrMax}{down to 0.44 times the errors}  % experiments/FIGS/data/v2/index.json :: binary probe bbqq_vs_ccqq: (1 - AUC) of FLAV_F1R over R16_Q1, run by run :: sha256 c6c4138ee5c86423 :: prose twin of the run farther from 1
\newcommand{\MlstCmpBbqqFzeroOverFull}{6.4 and 9.4}  % experiments/FIGS/data/v2/index.json :: binary probe bbqq_vs_ccqq: (1 - AUC) of FLAV_F0 over L188, run by run :: sha256 c6c4138ee5c86423
\newcommand{\MlstCmpBbqqFzeroOverFullPr}{at least 6.4 times the errors}  % experiments/FIGS/data/v2/index.json :: binary probe bbqq_vs_ccqq: (1 - AUC) of FLAV_F0 over L188, run by run :: sha256 c6c4138ee5c86423 :: prose twin: two significant figures, direction in words
\newcommand{\MlstCmpBbqqFzeroOverFullPrMax}{up to 9.4 times the errors}  % experiments/FIGS/data/v2/index.json :: binary probe bbqq_vs_ccqq: (1 - AUC) of FLAV_F0 over L188, run by run :: sha256 c6c4138ee5c86423 :: prose twin of the run farther from 1
\newcommand{\MlstCmpBbqqFoneOverFull}{1.510 and 1.447}  % experiments/FIGS/data/v2/index.json :: binary probe bbqq_vs_ccqq: (1 - AUC) of FLAV_F1 over L188, run by run :: sha256 c6c4138ee5c86423
\newcommand{\MlstCmpBbqqFoneOverFullPr}{at least 45\% more errors}  % experiments/FIGS/data/v2/index.json :: binary probe bbqq_vs_ccqq: (1 - AUC) of FLAV_F1 over L188, run by run :: sha256 c6c4138ee5c86423 :: prose twin: two significant figures, direction in words
\newcommand{\MlstCmpBbqqFoneOverFullPrMax}{up to 51\% more errors}  % experiments/FIGS/data/v2/index.json :: binary probe bbqq_vs_ccqq: (1 - AUC) of FLAV_F1 over L188, run by run :: sha256 c6c4138ee5c86423 :: prose twin of the run farther from 1
\newcommand{\MlstCmpBbqqFonerOverFull}{2.272 and 2.281}  % experiments/FIGS/data/v2/index.json :: binary probe bbqq_vs_ccqq: (1 - AUC) of FLAV_F1R over L188, run by run :: sha256 c6c4138ee5c86423
\newcommand{\MlstCmpBbqqFonerOverFullPr}{at least 2.3 times the errors}  % experiments/FIGS/data/v2/index.json :: binary probe bbqq_vs_ccqq: (1 - AUC) of FLAV_F1R over L188, run by run :: sha256 c6c4138ee5c86423 :: prose twin: two significant figures, direction in words
\newcommand{\MlstCmpBbqqFonerOverFullPrMax}{up to 2.3 times the errors}  % experiments/FIGS/data/v2/index.json :: binary probe bbqq_vs_ccqq: (1 - AUC) of FLAV_F1R over L188, run by run :: sha256 c6c4138ee5c86423 :: prose twin of the run farther from 1
\newcommand{\MlstCmpBbqqRandOverFull}{1.42\,\ensuremath{\pm}\,0.06}  % experiments/FIGS/data/v2/index.json :: binary probe bbqq_vs_ccqq: (1 - AUC) of each random partition over the 188-class run of its seed :: sha256 c6c4138ee5c86423
\newcommand{\MlstCmpBbqqRandOverFullPr}{42\% more errors}  % experiments/FIGS/data/v2/index.json :: binary probe bbqq_vs_ccqq: (1 - AUC) of each random partition over the 188-class run of its seed :: sha256 c6c4138ee5c86423 :: prose twin: two significant figures, direction in words
\newcommand{\MlstPairRandMergedBbqqN}{2}  % experiments/FIGS/data/v2/index.json :: random partitions that merge bbqq_vs_ccqq (count) :: sha256 c6c4138ee5c86423
\newcommand{\MlstPairRandApartBbqqN}{3}  % experiments/FIGS/data/v2/index.json :: random partitions that keep apart bbqq_vs_ccqq (count) :: sha256 c6c4138ee5c86423
\newcommand{\MlstPairRandMergedOverApartBbqq}{0.996}  % experiments/FIGS/data/v2/index.json :: binary probe bbqq_vs_ccqq: mean (1 - AUC) of the partitions that merge it over that of those that keep it apart (one run each) :: sha256 c6c4138ee5c86423
\newcommand{\MlstPairRandMergedOverApartBbqqPr}{0.44\% fewer errors}  % experiments/FIGS/data/v2/index.json :: binary probe bbqq_vs_ccqq: mean (1 - AUC) of the partitions that merge it over that of those that keep it apart (one run each) :: sha256 c6c4138ee5c86423 :: prose twin: two significant figures, direction in words
\newcommand{\MlstCmpEemmResOverFull}{--}  % experiments/FIGS/data/v2/index.json :: binary probe ee_vs_mm: (1 - AUC) of L162 over L188, run by run :: sha256 c6c4138ee5c86423 :: the reference's 1 - AUC is 0
\newcommand{\MlstCmpEemmResOverFullPr}{undefined, since the reference separates the test jets completely}  % experiments/FIGS/data/v2/index.json :: binary probe ee_vs_mm: (1 - AUC) of L162 over L188, run by run :: sha256 c6c4138ee5c86423 :: the reference's 1 - AUC is 0
\newcommand{\MlstCmpEemmMidOverFull}{--}  % experiments/FIGS/data/v2/index.json :: binary probe ee_vs_mm: (1 - AUC) of R42_Q1 over L188, run by run :: sha256 c6c4138ee5c86423 :: the reference's 1 - AUC is 0
\newcommand{\MlstCmpEemmMidOverFullPr}{undefined, since the reference separates the test jets completely}  % experiments/FIGS/data/v2/index.json :: binary probe ee_vs_mm: (1 - AUC) of R42_Q1 over L188, run by run :: sha256 c6c4138ee5c86423 :: the reference's 1 - AUC is 0
\newcommand{\MlstCmpEemmCoarseOverFull}{--}  % experiments/FIGS/data/v2/index.json :: binary probe ee_vs_mm: (1 - AUC) of R16_Q1 over L188, run by run :: sha256 c6c4138ee5c86423 :: the reference's 1 - AUC is 0
\newcommand{\MlstCmpEemmCoarseOverFullPr}{undefined, since the reference separates the test jets completely}  % experiments/FIGS/data/v2/index.json :: binary probe ee_vs_mm: (1 - AUC) of R16_Q1 over L188, run by run :: sha256 c6c4138ee5c86423 :: the reference's 1 - AUC is 0
\newcommand{\MlstCmpEemmFoneOverFzero}{--}  % experiments/FIGS/data/v2/index.json :: binary probe ee_vs_mm: (1 - AUC) of FLAV_F1 over FLAV_F0, run by run :: sha256 c6c4138ee5c86423 :: the reference's 1 - AUC is 0
\newcommand{\MlstCmpEemmFoneOverFzeroPr}{undefined, since the reference separates the test jets completely}  % experiments/FIGS/data/v2/index.json :: binary probe ee_vs_mm: (1 - AUC) of FLAV_F1 over FLAV_F0, run by run :: sha256 c6c4138ee5c86423 :: the reference's 1 - AUC is 0
\newcommand{\MlstCmpEemmFoneOverFzeroPrMax}{undefined, since the reference separates the test jets completely}  % experiments/FIGS/data/v2/index.json :: binary probe ee_vs_mm: (1 - AUC) of FLAV_F1 over FLAV_F0, run by run :: sha256 c6c4138ee5c86423 :: the reference's 1 - AUC is 0
\newcommand{\MlstCmpEemmFonerOverFone}{--}  % experiments/FIGS/data/v2/index.json :: binary probe ee_vs_mm: (1 - AUC) of FLAV_F1R over FLAV_F1, run by run :: sha256 c6c4138ee5c86423 :: the reference's 1 - AUC is 0
\newcommand{\MlstCmpEemmFonerOverFonePr}{undefined, since the reference separates the test jets completely}  % experiments/FIGS/data/v2/index.json :: binary probe ee_vs_mm: (1 - AUC) of FLAV_F1R over FLAV_F1, run by run :: sha256 c6c4138ee5c86423 :: the reference's 1 - AUC is 0
\newcommand{\MlstCmpEemmFonerOverFonePrMax}{undefined, since the reference separates the test jets completely}  % experiments/FIGS/data/v2/index.json :: binary probe ee_vs_mm: (1 - AUC) of FLAV_F1R over FLAV_F1, run by run :: sha256 c6c4138ee5c86423 :: the reference's 1 - AUC is 0
\newcommand{\MlstCmpEemmFonerOverFzero}{--}  % experiments/FIGS/data/v2/index.json :: binary probe ee_vs_mm: (1 - AUC) of FLAV_F1R over FLAV_F0, run by run :: sha256 c6c4138ee5c86423 :: the reference's 1 - AUC is 0
\newcommand{\MlstCmpEemmFonerOverFzeroPr}{undefined, since the reference separates the test jets completely}  % experiments/FIGS/data/v2/index.json :: binary probe ee_vs_mm: (1 - AUC) of FLAV_F1R over FLAV_F0, run by run :: sha256 c6c4138ee5c86423 :: the reference's 1 - AUC is 0
\newcommand{\MlstCmpEemmFonerOverFzeroPrMax}{undefined, since the reference separates the test jets completely}  % experiments/FIGS/data/v2/index.json :: binary probe ee_vs_mm: (1 - AUC) of FLAV_F1R over FLAV_F0, run by run :: sha256 c6c4138ee5c86423 :: the reference's 1 - AUC is 0
\newcommand{\MlstCmpEemmFzeroOverCoarse}{0.0000022 and 0.0000003}  % experiments/FIGS/data/v2/index.json :: binary probe ee_vs_mm: (1 - AUC) of FLAV_F0 over R16_Q1, run by run :: sha256 c6c4138ee5c86423
\newcommand{\MlstCmpEemmFzeroOverCoarsePr}{at most 0.0000022 times the errors}  % experiments/FIGS/data/v2/index.json :: binary probe ee_vs_mm: (1 - AUC) of FLAV_F0 over R16_Q1, run by run :: sha256 c6c4138ee5c86423 :: prose twin: two significant figures, direction in words
\newcommand{\MlstCmpEemmFzeroOverCoarsePrMax}{down to 0.00000034 times the errors}  % experiments/FIGS/data/v2/index.json :: binary probe ee_vs_mm: (1 - AUC) of FLAV_F0 over R16_Q1, run by run :: sha256 c6c4138ee5c86423 :: prose twin of the run farther from 1
\newcommand{\MlstCmpEemmFoneOverCoarse}{0.000000 and 0.000038}  % experiments/FIGS/data/v2/index.json :: binary probe ee_vs_mm: (1 - AUC) of FLAV_F1 over R16_Q1, run by run :: sha256 c6c4138ee5c86423
\newcommand{\MlstCmpEemmFoneOverCoarsePr}{at most 0.000038 times the errors}  % experiments/FIGS/data/v2/index.json :: binary probe ee_vs_mm: (1 - AUC) of FLAV_F1 over R16_Q1, run by run :: sha256 c6c4138ee5c86423 :: prose twin: two significant figures, direction in words
\newcommand{\MlstCmpEemmFoneOverCoarsePrMax}{down to 0.000000042 times the errors}  % experiments/FIGS/data/v2/index.json :: binary probe ee_vs_mm: (1 - AUC) of FLAV_F1 over R16_Q1, run by run :: sha256 c6c4138ee5c86423 :: prose twin of the run farther from 1
\newcommand{\MlstCmpEemmFonerOverCoarse}{0.00000005 and 0.00000017}  % experiments/FIGS/data/v2/index.json :: binary probe ee_vs_mm: (1 - AUC) of FLAV_F1R over R16_Q1, run by run :: sha256 c6c4138ee5c86423
\newcommand{\MlstCmpEemmFonerOverCoarsePr}{at most 0.00000017 times the errors}  % experiments/FIGS/data/v2/index.json :: binary probe ee_vs_mm: (1 - AUC) of FLAV_F1R over R16_Q1, run by run :: sha256 c6c4138ee5c86423 :: prose twin: two significant figures, direction in words
\newcommand{\MlstCmpEemmFonerOverCoarsePrMax}{down to 0.000000055 times the errors}  % experiments/FIGS/data/v2/index.json :: binary probe ee_vs_mm: (1 - AUC) of FLAV_F1R over R16_Q1, run by run :: sha256 c6c4138ee5c86423 :: prose twin of the run farther from 1
\newcommand{\MlstCmpEemmFzeroOverFull}{--}  % experiments/FIGS/data/v2/index.json :: binary probe ee_vs_mm: (1 - AUC) of FLAV_F0 over L188, run by run :: sha256 c6c4138ee5c86423 :: the reference's 1 - AUC is 0
\newcommand{\MlstCmpEemmFzeroOverFullPr}{undefined, since the reference separates the test jets completely}  % experiments/FIGS/data/v2/index.json :: binary probe ee_vs_mm: (1 - AUC) of FLAV_F0 over L188, run by run :: sha256 c6c4138ee5c86423 :: the reference's 1 - AUC is 0
\newcommand{\MlstCmpEemmFzeroOverFullPrMax}{undefined, since the reference separates the test jets completely}  % experiments/FIGS/data/v2/index.json :: binary probe ee_vs_mm: (1 - AUC) of FLAV_F0 over L188, run by run :: sha256 c6c4138ee5c86423 :: the reference's 1 - AUC is 0
\newcommand{\MlstCmpEemmFoneOverFull}{--}  % experiments/FIGS/data/v2/index.json :: binary probe ee_vs_mm: (1 - AUC) of FLAV_F1 over L188, run by run :: sha256 c6c4138ee5c86423 :: the reference's 1 - AUC is 0
\newcommand{\MlstCmpEemmFoneOverFullPr}{undefined, since the reference separates the test jets completely}  % experiments/FIGS/data/v2/index.json :: binary probe ee_vs_mm: (1 - AUC) of FLAV_F1 over L188, run by run :: sha256 c6c4138ee5c86423 :: the reference's 1 - AUC is 0
\newcommand{\MlstCmpEemmFoneOverFullPrMax}{undefined, since the reference separates the test jets completely}  % experiments/FIGS/data/v2/index.json :: binary probe ee_vs_mm: (1 - AUC) of FLAV_F1 over L188, run by run :: sha256 c6c4138ee5c86423 :: the reference's 1 - AUC is 0
\newcommand{\MlstCmpEemmFonerOverFull}{--}  % experiments/FIGS/data/v2/index.json :: binary probe ee_vs_mm: (1 - AUC) of FLAV_F1R over L188, run by run :: sha256 c6c4138ee5c86423 :: the reference's 1 - AUC is 0
\newcommand{\MlstCmpEemmFonerOverFullPr}{undefined, since the reference separates the test jets completely}  % experiments/FIGS/data/v2/index.json :: binary probe ee_vs_mm: (1 - AUC) of FLAV_F1R over L188, run by run :: sha256 c6c4138ee5c86423 :: the reference's 1 - AUC is 0
\newcommand{\MlstCmpEemmFonerOverFullPrMax}{undefined, since the reference separates the test jets completely}  % experiments/FIGS/data/v2/index.json :: binary probe ee_vs_mm: (1 - AUC) of FLAV_F1R over L188, run by run :: sha256 c6c4138ee5c86423 :: the reference's 1 - AUC is 0
\newcommand{\MlstCmpEemmRandOverFull}{--}  % experiments/FIGS/data/v2/index.json :: binary probe ee_vs_mm: (1 - AUC) of each random partition over the 188-class run of its seed :: sha256 c6c4138ee5c86423 :: the reference's 1 - AUC is 0
\newcommand{\MlstCmpEemmRandOverFullPr}{undefined, since the reference separates the test jets completely}  % experiments/FIGS/data/v2/index.json :: binary probe ee_vs_mm: (1 - AUC) of each random partition over the 188-class run of its seed :: sha256 c6c4138ee5c86423 :: the reference's 1 - AUC is 0
\newcommand{\MlstPairRandMergedEemmN}{3}  % experiments/FIGS/data/v2/index.json :: random partitions that merge ee_vs_mm (count) :: sha256 c6c4138ee5c86423
\newcommand{\MlstPairRandApartEemmN}{2}  % experiments/FIGS/data/v2/index.json :: random partitions that keep apart ee_vs_mm (count) :: sha256 c6c4138ee5c86423
\newcommand{\MlstPairRandMergedOverApartEemm}{75.3}  % experiments/FIGS/data/v2/index.json :: binary probe ee_vs_mm: mean (1 - AUC) of the partitions that merge it over that of those that keep it apart (one run each) :: sha256 c6c4138ee5c86423
\newcommand{\MlstPairRandMergedOverApartEemmPr}{75 times the errors}  % experiments/FIGS/data/v2/index.json :: binary probe ee_vs_mm: mean (1 - AUC) of the partitions that merge it over that of those that keep it apart (one run each) :: sha256 c6c4138ee5c86423 :: prose twin: two significant figures, direction in words
\newcommand{\MlstCmpQcdBbResOverFull}{1.15\,\ensuremath{\pm}\,0.04}  % experiments/FIGS/data/v2/index.json :: binary probe qcd_bb_vs_cc: (1 - AUC) of L162 over L188, run by run :: sha256 c6c4138ee5c86423
\newcommand{\MlstCmpQcdBbResOverFullPr}{15\% more errors}  % experiments/FIGS/data/v2/index.json :: binary probe qcd_bb_vs_cc: (1 - AUC) of L162 over L188, run by run :: sha256 c6c4138ee5c86423 :: prose twin: two significant figures, direction in words
\newcommand{\MlstCmpQcdBbMidOverFull}{1.163\,\ensuremath{\pm}\,0.016}  % experiments/FIGS/data/v2/index.json :: binary probe qcd_bb_vs_cc: (1 - AUC) of R42_Q1 over L188, run by run :: sha256 c6c4138ee5c86423
\newcommand{\MlstCmpQcdBbMidOverFullPr}{16\% more errors}  % experiments/FIGS/data/v2/index.json :: binary probe qcd_bb_vs_cc: (1 - AUC) of R42_Q1 over L188, run by run :: sha256 c6c4138ee5c86423 :: prose twin: two significant figures, direction in words
\newcommand{\MlstCmpQcdBbCoarseOverFull}{3.294\,\ensuremath{\pm}\,0.033}  % experiments/FIGS/data/v2/index.json :: binary probe qcd_bb_vs_cc: (1 - AUC) of R16_Q1 over L188, run by run :: sha256 c6c4138ee5c86423
\newcommand{\MlstCmpQcdBbCoarseOverFullPr}{3.3 times the errors}  % experiments/FIGS/data/v2/index.json :: binary probe qcd_bb_vs_cc: (1 - AUC) of R16_Q1 over L188, run by run :: sha256 c6c4138ee5c86423 :: prose twin: two significant figures, direction in words
\newcommand{\MlstCmpQcdBbFoneOverFzero}{0.368 and 0.294}  % experiments/FIGS/data/v2/index.json :: binary probe qcd_bb_vs_cc: (1 - AUC) of FLAV_F1 over FLAV_F0, run by run :: sha256 c6c4138ee5c86423
\newcommand{\MlstCmpQcdBbFoneOverFzeroPr}{at most 0.37 times the errors}  % experiments/FIGS/data/v2/index.json :: binary probe qcd_bb_vs_cc: (1 - AUC) of FLAV_F1 over FLAV_F0, run by run :: sha256 c6c4138ee5c86423 :: prose twin: two significant figures, direction in words
\newcommand{\MlstCmpQcdBbFoneOverFzeroPrMax}{down to 0.29 times the errors}  % experiments/FIGS/data/v2/index.json :: binary probe qcd_bb_vs_cc: (1 - AUC) of FLAV_F1 over FLAV_F0, run by run :: sha256 c6c4138ee5c86423 :: prose twin of the run farther from 1
\newcommand{\MlstCmpQcdBbFonerOverFone}{1.361 and 1.429}  % experiments/FIGS/data/v2/index.json :: binary probe qcd_bb_vs_cc: (1 - AUC) of FLAV_F1R over FLAV_F1, run by run :: sha256 c6c4138ee5c86423
\newcommand{\MlstCmpQcdBbFonerOverFonePr}{at least 36\% more errors}  % experiments/FIGS/data/v2/index.json :: binary probe qcd_bb_vs_cc: (1 - AUC) of FLAV_F1R over FLAV_F1, run by run :: sha256 c6c4138ee5c86423 :: prose twin: two significant figures, direction in words
\newcommand{\MlstCmpQcdBbFonerOverFonePrMax}{up to 43\% more errors}  % experiments/FIGS/data/v2/index.json :: binary probe qcd_bb_vs_cc: (1 - AUC) of FLAV_F1R over FLAV_F1, run by run :: sha256 c6c4138ee5c86423 :: prose twin of the run farther from 1
\newcommand{\MlstCmpQcdBbFonerOverFzero}{0.501 and 0.420}  % experiments/FIGS/data/v2/index.json :: binary probe qcd_bb_vs_cc: (1 - AUC) of FLAV_F1R over FLAV_F0, run by run :: sha256 c6c4138ee5c86423
\newcommand{\MlstCmpQcdBbFonerOverFzeroPr}{at least 50\% fewer errors}  % experiments/FIGS/data/v2/index.json :: binary probe qcd_bb_vs_cc: (1 - AUC) of FLAV_F1R over FLAV_F0, run by run :: sha256 c6c4138ee5c86423 :: prose twin: two significant figures, direction in words
\newcommand{\MlstCmpQcdBbFonerOverFzeroPrMax}{down to 0.42 times the errors}  % experiments/FIGS/data/v2/index.json :: binary probe qcd_bb_vs_cc: (1 - AUC) of FLAV_F1R over FLAV_F0, run by run :: sha256 c6c4138ee5c86423 :: prose twin of the run farther from 1
\newcommand{\MlstCmpQcdBbFzeroOverCoarse}{1.28 and 1.58}  % experiments/FIGS/data/v2/index.json :: binary probe qcd_bb_vs_cc: (1 - AUC) of FLAV_F0 over R16_Q1, run by run :: sha256 c6c4138ee5c86423
\newcommand{\MlstCmpQcdBbFzeroOverCoarsePr}{at least 28\% more errors}  % experiments/FIGS/data/v2/index.json :: binary probe qcd_bb_vs_cc: (1 - AUC) of FLAV_F0 over R16_Q1, run by run :: sha256 c6c4138ee5c86423 :: prose twin: two significant figures, direction in words
\newcommand{\MlstCmpQcdBbFzeroOverCoarsePrMax}{up to 58\% more errors}  % experiments/FIGS/data/v2/index.json :: binary probe qcd_bb_vs_cc: (1 - AUC) of FLAV_F0 over R16_Q1, run by run :: sha256 c6c4138ee5c86423 :: prose twin of the run farther from 1
\newcommand{\MlstCmpQcdBbFoneOverCoarse}{0.4692 and 0.4648}  % experiments/FIGS/data/v2/index.json :: binary probe qcd_bb_vs_cc: (1 - AUC) of FLAV_F1 over R16_Q1, run by run :: sha256 c6c4138ee5c86423
\newcommand{\MlstCmpQcdBbFoneOverCoarsePr}{at most 0.47 times the errors}  % experiments/FIGS/data/v2/index.json :: binary probe qcd_bb_vs_cc: (1 - AUC) of FLAV_F1 over R16_Q1, run by run :: sha256 c6c4138ee5c86423 :: prose twin: two significant figures, direction in words
\newcommand{\MlstCmpQcdBbFoneOverCoarsePrMax}{down to 0.46 times the errors}  % experiments/FIGS/data/v2/index.json :: binary probe qcd_bb_vs_cc: (1 - AUC) of FLAV_F1 over R16_Q1, run by run :: sha256 c6c4138ee5c86423 :: prose twin of the run farther from 1
\newcommand{\MlstCmpQcdBbFonerOverCoarse}{0.639 and 0.664}  % experiments/FIGS/data/v2/index.json :: binary probe qcd_bb_vs_cc: (1 - AUC) of FLAV_F1R over R16_Q1, run by run :: sha256 c6c4138ee5c86423
\newcommand{\MlstCmpQcdBbFonerOverCoarsePr}{at least 34\% fewer errors}  % experiments/FIGS/data/v2/index.json :: binary probe qcd_bb_vs_cc: (1 - AUC) of FLAV_F1R over R16_Q1, run by run :: sha256 c6c4138ee5c86423 :: prose twin: two significant figures, direction in words
\newcommand{\MlstCmpQcdBbFonerOverCoarsePrMax}{up to 36\% fewer errors}  % experiments/FIGS/data/v2/index.json :: binary probe qcd_bb_vs_cc: (1 - AUC) of FLAV_F1R over R16_Q1, run by run :: sha256 c6c4138ee5c86423 :: prose twin of the run farther from 1
\newcommand{\MlstCmpQcdBbFzeroOverFull}{4.24 and 5.21}  % experiments/FIGS/data/v2/index.json :: binary probe qcd_bb_vs_cc: (1 - AUC) of FLAV_F0 over L188, run by run :: sha256 c6c4138ee5c86423
\newcommand{\MlstCmpQcdBbFzeroOverFullPr}{at least 4.2 times the errors}  % experiments/FIGS/data/v2/index.json :: binary probe qcd_bb_vs_cc: (1 - AUC) of FLAV_F0 over L188, run by run :: sha256 c6c4138ee5c86423 :: prose twin: two significant figures, direction in words
\newcommand{\MlstCmpQcdBbFzeroOverFullPrMax}{up to 5.2 times the errors}  % experiments/FIGS/data/v2/index.json :: binary probe qcd_bb_vs_cc: (1 - AUC) of FLAV_F0 over L188, run by run :: sha256 c6c4138ee5c86423 :: prose twin of the run farther from 1
\newcommand{\MlstCmpQcdBbFoneOverFull}{1.561 and 1.532}  % experiments/FIGS/data/v2/index.json :: binary probe qcd_bb_vs_cc: (1 - AUC) of FLAV_F1 over L188, run by run :: sha256 c6c4138ee5c86423
\newcommand{\MlstCmpQcdBbFoneOverFullPr}{at least 53\% more errors}  % experiments/FIGS/data/v2/index.json :: binary probe qcd_bb_vs_cc: (1 - AUC) of FLAV_F1 over L188, run by run :: sha256 c6c4138ee5c86423 :: prose twin: two significant figures, direction in words
\newcommand{\MlstCmpQcdBbFoneOverFullPrMax}{up to 56\% more errors}  % experiments/FIGS/data/v2/index.json :: binary probe qcd_bb_vs_cc: (1 - AUC) of FLAV_F1 over L188, run by run :: sha256 c6c4138ee5c86423 :: prose twin of the run farther from 1
\newcommand{\MlstCmpQcdBbFonerOverFull}{2.124 and 2.189}  % experiments/FIGS/data/v2/index.json :: binary probe qcd_bb_vs_cc: (1 - AUC) of FLAV_F1R over L188, run by run :: sha256 c6c4138ee5c86423
\newcommand{\MlstCmpQcdBbFonerOverFullPr}{at least 2.1 times the errors}  % experiments/FIGS/data/v2/index.json :: binary probe qcd_bb_vs_cc: (1 - AUC) of FLAV_F1R over L188, run by run :: sha256 c6c4138ee5c86423 :: prose twin: two significant figures, direction in words
\newcommand{\MlstCmpQcdBbFonerOverFullPrMax}{up to 2.2 times the errors}  % experiments/FIGS/data/v2/index.json :: binary probe qcd_bb_vs_cc: (1 - AUC) of FLAV_F1R over L188, run by run :: sha256 c6c4138ee5c86423 :: prose twin of the run farther from 1
\newcommand{\MlstCmpQcdBResOverFull}{1.085\,\ensuremath{\pm}\,0.004}  % experiments/FIGS/data/v2/index.json :: binary probe qcd_b_vs_light: (1 - AUC) of L162 over L188, run by run :: sha256 c6c4138ee5c86423
\newcommand{\MlstCmpQcdBResOverFullPr}{8.5\% more errors}  % experiments/FIGS/data/v2/index.json :: binary probe qcd_b_vs_light: (1 - AUC) of L162 over L188, run by run :: sha256 c6c4138ee5c86423 :: prose twin: two significant figures, direction in words
\newcommand{\MlstCmpQcdBMidOverFull}{1.081\,\ensuremath{\pm}\,0.008}  % experiments/FIGS/data/v2/index.json :: binary probe qcd_b_vs_light: (1 - AUC) of R42_Q1 over L188, run by run :: sha256 c6c4138ee5c86423
\newcommand{\MlstCmpQcdBMidOverFullPr}{8.1\% more errors}  % experiments/FIGS/data/v2/index.json :: binary probe qcd_b_vs_light: (1 - AUC) of R42_Q1 over L188, run by run :: sha256 c6c4138ee5c86423 :: prose twin: two significant figures, direction in words
\newcommand{\MlstCmpQcdBCoarseOverFull}{1.54\,\ensuremath{\pm}\,0.07}  % experiments/FIGS/data/v2/index.json :: binary probe qcd_b_vs_light: (1 - AUC) of R16_Q1 over L188, run by run :: sha256 c6c4138ee5c86423
\newcommand{\MlstCmpQcdBCoarseOverFullPr}{54\% more errors}  % experiments/FIGS/data/v2/index.json :: binary probe qcd_b_vs_light: (1 - AUC) of R16_Q1 over L188, run by run :: sha256 c6c4138ee5c86423 :: prose twin: two significant figures, direction in words
\newcommand{\MlstCmpQcdBFoneOverFzero}{0.778 and 0.835}  % experiments/FIGS/data/v2/index.json :: binary probe qcd_b_vs_light: (1 - AUC) of FLAV_F1 over FLAV_F0, run by run :: sha256 c6c4138ee5c86423
\newcommand{\MlstCmpQcdBFoneOverFzeroPr}{at least 17\% fewer errors}  % experiments/FIGS/data/v2/index.json :: binary probe qcd_b_vs_light: (1 - AUC) of FLAV_F1 over FLAV_F0, run by run :: sha256 c6c4138ee5c86423 :: prose twin: two significant figures, direction in words
\newcommand{\MlstCmpQcdBFoneOverFzeroPrMax}{up to 22\% fewer errors}  % experiments/FIGS/data/v2/index.json :: binary probe qcd_b_vs_light: (1 - AUC) of FLAV_F1 over FLAV_F0, run by run :: sha256 c6c4138ee5c86423 :: prose twin of the run farther from 1
\newcommand{\MlstCmpQcdBFonerOverFone}{1.054 and 1.061}  % experiments/FIGS/data/v2/index.json :: binary probe qcd_b_vs_light: (1 - AUC) of FLAV_F1R over FLAV_F1, run by run :: sha256 c6c4138ee5c86423
\newcommand{\MlstCmpQcdBFonerOverFonePr}{equal within 7\%}  % experiments/FIGS/data/v2/index.json :: binary probe qcd_b_vs_light: (1 - AUC) of FLAV_F1R over FLAV_F1, run by run :: sha256 c6c4138ee5c86423 :: prose twin: two significant figures, direction in words
\newcommand{\MlstCmpQcdBFonerOverFonePrMax}{equal within 7\%}  % experiments/FIGS/data/v2/index.json :: binary probe qcd_b_vs_light: (1 - AUC) of FLAV_F1R over FLAV_F1, run by run :: sha256 c6c4138ee5c86423 :: prose twin of the run farther from 1
\newcommand{\MlstCmpQcdBFonerOverFzero}{0.820 and 0.885}  % experiments/FIGS/data/v2/index.json :: binary probe qcd_b_vs_light: (1 - AUC) of FLAV_F1R over FLAV_F0, run by run :: sha256 c6c4138ee5c86423
\newcommand{\MlstCmpQcdBFonerOverFzeroPr}{at least 11\% fewer errors}  % experiments/FIGS/data/v2/index.json :: binary probe qcd_b_vs_light: (1 - AUC) of FLAV_F1R over FLAV_F0, run by run :: sha256 c6c4138ee5c86423 :: prose twin: two significant figures, direction in words
\newcommand{\MlstCmpQcdBFonerOverFzeroPrMax}{up to 18\% fewer errors}  % experiments/FIGS/data/v2/index.json :: binary probe qcd_b_vs_light: (1 - AUC) of FLAV_F1R over FLAV_F0, run by run :: sha256 c6c4138ee5c86423 :: prose twin of the run farther from 1
\newcommand{\MlstCmpQcdBFzeroOverCoarse}{0.962 and 0.981}  % experiments/FIGS/data/v2/index.json :: binary probe qcd_b_vs_light: (1 - AUC) of FLAV_F0 over R16_Q1, run by run :: sha256 c6c4138ee5c86423
\newcommand{\MlstCmpQcdBFzeroOverCoarsePr}{equal within 4\%}  % experiments/FIGS/data/v2/index.json :: binary probe qcd_b_vs_light: (1 - AUC) of FLAV_F0 over R16_Q1, run by run :: sha256 c6c4138ee5c86423 :: prose twin: two significant figures, direction in words
\newcommand{\MlstCmpQcdBFzeroOverCoarsePrMax}{equal within 4\%}  % experiments/FIGS/data/v2/index.json :: binary probe qcd_b_vs_light: (1 - AUC) of FLAV_F0 over R16_Q1, run by run :: sha256 c6c4138ee5c86423 :: prose twin of the run farther from 1
\newcommand{\MlstCmpQcdBFoneOverCoarse}{0.749 and 0.819}  % experiments/FIGS/data/v2/index.json :: binary probe qcd_b_vs_light: (1 - AUC) of FLAV_F1 over R16_Q1, run by run :: sha256 c6c4138ee5c86423
\newcommand{\MlstCmpQcdBFoneOverCoarsePr}{at least 18\% fewer errors}  % experiments/FIGS/data/v2/index.json :: binary probe qcd_b_vs_light: (1 - AUC) of FLAV_F1 over R16_Q1, run by run :: sha256 c6c4138ee5c86423 :: prose twin: two significant figures, direction in words
\newcommand{\MlstCmpQcdBFoneOverCoarsePrMax}{up to 25\% fewer errors}  % experiments/FIGS/data/v2/index.json :: binary probe qcd_b_vs_light: (1 - AUC) of FLAV_F1 over R16_Q1, run by run :: sha256 c6c4138ee5c86423 :: prose twin of the run farther from 1
\newcommand{\MlstCmpQcdBFonerOverCoarse}{0.789 and 0.869}  % experiments/FIGS/data/v2/index.json :: binary probe qcd_b_vs_light: (1 - AUC) of FLAV_F1R over R16_Q1, run by run :: sha256 c6c4138ee5c86423
\newcommand{\MlstCmpQcdBFonerOverCoarsePr}{at least 13\% fewer errors}  % experiments/FIGS/data/v2/index.json :: binary probe qcd_b_vs_light: (1 - AUC) of FLAV_F1R over R16_Q1, run by run :: sha256 c6c4138ee5c86423 :: prose twin: two significant figures, direction in words
\newcommand{\MlstCmpQcdBFonerOverCoarsePrMax}{up to 21\% fewer errors}  % experiments/FIGS/data/v2/index.json :: binary probe qcd_b_vs_light: (1 - AUC) of FLAV_F1R over R16_Q1, run by run :: sha256 c6c4138ee5c86423 :: prose twin of the run farther from 1
\newcommand{\MlstCmpQcdBFzeroOverFull}{1.56 and 1.45}  % experiments/FIGS/data/v2/index.json :: binary probe qcd_b_vs_light: (1 - AUC) of FLAV_F0 over L188, run by run :: sha256 c6c4138ee5c86423
\newcommand{\MlstCmpQcdBFzeroOverFullPr}{at least 45\% more errors}  % experiments/FIGS/data/v2/index.json :: binary probe qcd_b_vs_light: (1 - AUC) of FLAV_F0 over L188, run by run :: sha256 c6c4138ee5c86423 :: prose twin: two significant figures, direction in words
\newcommand{\MlstCmpQcdBFzeroOverFullPrMax}{up to 56\% more errors}  % experiments/FIGS/data/v2/index.json :: binary probe qcd_b_vs_light: (1 - AUC) of FLAV_F0 over L188, run by run :: sha256 c6c4138ee5c86423 :: prose twin of the run farther from 1
\newcommand{\MlstCmpQcdBFoneOverFull}{1.212 and 1.213}  % experiments/FIGS/data/v2/index.json :: binary probe qcd_b_vs_light: (1 - AUC) of FLAV_F1 over L188, run by run :: sha256 c6c4138ee5c86423
\newcommand{\MlstCmpQcdBFoneOverFullPr}{at least 21\% more errors}  % experiments/FIGS/data/v2/index.json :: binary probe qcd_b_vs_light: (1 - AUC) of FLAV_F1 over L188, run by run :: sha256 c6c4138ee5c86423 :: prose twin: two significant figures, direction in words
\newcommand{\MlstCmpQcdBFoneOverFullPrMax}{up to 21\% more errors}  % experiments/FIGS/data/v2/index.json :: binary probe qcd_b_vs_light: (1 - AUC) of FLAV_F1 over L188, run by run :: sha256 c6c4138ee5c86423 :: prose twin of the run farther from 1
\newcommand{\MlstCmpQcdBFonerOverFull}{1.277 and 1.286}  % experiments/FIGS/data/v2/index.json :: binary probe qcd_b_vs_light: (1 - AUC) of FLAV_F1R over L188, run by run :: sha256 c6c4138ee5c86423
\newcommand{\MlstCmpQcdBFonerOverFullPr}{at least 28\% more errors}  % experiments/FIGS/data/v2/index.json :: binary probe qcd_b_vs_light: (1 - AUC) of FLAV_F1R over L188, run by run :: sha256 c6c4138ee5c86423 :: prose twin: two significant figures, direction in words
\newcommand{\MlstCmpQcdBFonerOverFullPrMax}{up to 29\% more errors}  % experiments/FIGS/data/v2/index.json :: binary probe qcd_b_vs_light: (1 - AUC) of FLAV_F1R over L188, run by run :: sha256 c6c4138ee5c86423 :: prose twin of the run farther from 1
\newcommand{\MlstCmpQcdCResOverFull}{1.061\,\ensuremath{\pm}\,0.005}  % experiments/FIGS/data/v2/index.json :: binary probe qcd_c_vs_light: (1 - AUC) of L162 over L188, run by run :: sha256 c6c4138ee5c86423
\newcommand{\MlstCmpQcdCResOverFullPr}{6.1\% more errors}  % experiments/FIGS/data/v2/index.json :: binary probe qcd_c_vs_light: (1 - AUC) of L162 over L188, run by run :: sha256 c6c4138ee5c86423 :: prose twin: two significant figures, direction in words
\newcommand{\MlstCmpQcdCMidOverFull}{1.072\,\ensuremath{\pm}\,0.008}  % experiments/FIGS/data/v2/index.json :: binary probe qcd_c_vs_light: (1 - AUC) of R42_Q1 over L188, run by run :: sha256 c6c4138ee5c86423
\newcommand{\MlstCmpQcdCMidOverFullPr}{7.2\% more errors}  % experiments/FIGS/data/v2/index.json :: binary probe qcd_c_vs_light: (1 - AUC) of R42_Q1 over L188, run by run :: sha256 c6c4138ee5c86423 :: prose twin: two significant figures, direction in words
\newcommand{\MlstCmpQcdCCoarseOverFull}{1.407\,\ensuremath{\pm}\,0.021}  % experiments/FIGS/data/v2/index.json :: binary probe qcd_c_vs_light: (1 - AUC) of R16_Q1 over L188, run by run :: sha256 c6c4138ee5c86423
\newcommand{\MlstCmpQcdCCoarseOverFullPr}{41\% more errors}  % experiments/FIGS/data/v2/index.json :: binary probe qcd_c_vs_light: (1 - AUC) of R16_Q1 over L188, run by run :: sha256 c6c4138ee5c86423 :: prose twin: two significant figures, direction in words
\newcommand{\MlstCmpQcdCFoneOverFzero}{0.915 and 0.944}  % experiments/FIGS/data/v2/index.json :: binary probe qcd_c_vs_light: (1 - AUC) of FLAV_F1 over FLAV_F0, run by run :: sha256 c6c4138ee5c86423
\newcommand{\MlstCmpQcdCFoneOverFzeroPr}{at least 5.6\% fewer errors}  % experiments/FIGS/data/v2/index.json :: binary probe qcd_c_vs_light: (1 - AUC) of FLAV_F1 over FLAV_F0, run by run :: sha256 c6c4138ee5c86423 :: prose twin: two significant figures, direction in words
\newcommand{\MlstCmpQcdCFoneOverFzeroPrMax}{up to 8.5\% fewer errors}  % experiments/FIGS/data/v2/index.json :: binary probe qcd_c_vs_light: (1 - AUC) of FLAV_F1 over FLAV_F0, run by run :: sha256 c6c4138ee5c86423 :: prose twin of the run farther from 1
\newcommand{\MlstCmpQcdCFonerOverFone}{0.9535 and 0.9520}  % experiments/FIGS/data/v2/index.json :: binary probe qcd_c_vs_light: (1 - AUC) of FLAV_F1R over FLAV_F1, run by run :: sha256 c6c4138ee5c86423
\newcommand{\MlstCmpQcdCFonerOverFonePr}{at least 4.6\% fewer errors}  % experiments/FIGS/data/v2/index.json :: binary probe qcd_c_vs_light: (1 - AUC) of FLAV_F1R over FLAV_F1, run by run :: sha256 c6c4138ee5c86423 :: prose twin: two significant figures, direction in words
\newcommand{\MlstCmpQcdCFonerOverFonePrMax}{up to 4.8\% fewer errors}  % experiments/FIGS/data/v2/index.json :: binary probe qcd_c_vs_light: (1 - AUC) of FLAV_F1R over FLAV_F1, run by run :: sha256 c6c4138ee5c86423 :: prose twin of the run farther from 1
\newcommand{\MlstCmpQcdCFonerOverFzero}{0.872 and 0.899}  % experiments/FIGS/data/v2/index.json :: binary probe qcd_c_vs_light: (1 - AUC) of FLAV_F1R over FLAV_F0, run by run :: sha256 c6c4138ee5c86423
\newcommand{\MlstCmpQcdCFonerOverFzeroPr}{at least 10\% fewer errors}  % experiments/FIGS/data/v2/index.json :: binary probe qcd_c_vs_light: (1 - AUC) of FLAV_F1R over FLAV_F0, run by run :: sha256 c6c4138ee5c86423 :: prose twin: two significant figures, direction in words
\newcommand{\MlstCmpQcdCFonerOverFzeroPrMax}{up to 13\% fewer errors}  % experiments/FIGS/data/v2/index.json :: binary probe qcd_c_vs_light: (1 - AUC) of FLAV_F1R over FLAV_F0, run by run :: sha256 c6c4138ee5c86423 :: prose twin of the run farther from 1
\newcommand{\MlstCmpQcdCFzeroOverCoarse}{0.952 and 0.934}  % experiments/FIGS/data/v2/index.json :: binary probe qcd_c_vs_light: (1 - AUC) of FLAV_F0 over R16_Q1, run by run :: sha256 c6c4138ee5c86423
\newcommand{\MlstCmpQcdCFzeroOverCoarsePr}{at least 4.8\% fewer errors}  % experiments/FIGS/data/v2/index.json :: binary probe qcd_c_vs_light: (1 - AUC) of FLAV_F0 over R16_Q1, run by run :: sha256 c6c4138ee5c86423 :: prose twin: two significant figures, direction in words
\newcommand{\MlstCmpQcdCFzeroOverCoarsePrMax}{up to 6.6\% fewer errors}  % experiments/FIGS/data/v2/index.json :: binary probe qcd_c_vs_light: (1 - AUC) of FLAV_F0 over R16_Q1, run by run :: sha256 c6c4138ee5c86423 :: prose twin of the run farther from 1
\newcommand{\MlstCmpQcdCFoneOverCoarse}{0.871 and 0.882}  % experiments/FIGS/data/v2/index.json :: binary probe qcd_c_vs_light: (1 - AUC) of FLAV_F1 over R16_Q1, run by run :: sha256 c6c4138ee5c86423
\newcommand{\MlstCmpQcdCFoneOverCoarsePr}{at least 12\% fewer errors}  % experiments/FIGS/data/v2/index.json :: binary probe qcd_c_vs_light: (1 - AUC) of FLAV_F1 over R16_Q1, run by run :: sha256 c6c4138ee5c86423 :: prose twin: two significant figures, direction in words
\newcommand{\MlstCmpQcdCFoneOverCoarsePrMax}{up to 13\% fewer errors}  % experiments/FIGS/data/v2/index.json :: binary probe qcd_c_vs_light: (1 - AUC) of FLAV_F1 over R16_Q1, run by run :: sha256 c6c4138ee5c86423 :: prose twin of the run farther from 1
\newcommand{\MlstCmpQcdCFonerOverCoarse}{0.8301 and 0.8396}  % experiments/FIGS/data/v2/index.json :: binary probe qcd_c_vs_light: (1 - AUC) of FLAV_F1R over R16_Q1, run by run :: sha256 c6c4138ee5c86423
\newcommand{\MlstCmpQcdCFonerOverCoarsePr}{at least 16\% fewer errors}  % experiments/FIGS/data/v2/index.json :: binary probe qcd_c_vs_light: (1 - AUC) of FLAV_F1R over R16_Q1, run by run :: sha256 c6c4138ee5c86423 :: prose twin: two significant figures, direction in words
\newcommand{\MlstCmpQcdCFonerOverCoarsePrMax}{up to 17\% fewer errors}  % experiments/FIGS/data/v2/index.json :: binary probe qcd_c_vs_light: (1 - AUC) of FLAV_F1R over R16_Q1, run by run :: sha256 c6c4138ee5c86423 :: prose twin of the run farther from 1
\newcommand{\MlstCmpQcdCFzeroOverFull}{1.352 and 1.293}  % experiments/FIGS/data/v2/index.json :: binary probe qcd_c_vs_light: (1 - AUC) of FLAV_F0 over L188, run by run :: sha256 c6c4138ee5c86423
\newcommand{\MlstCmpQcdCFzeroOverFullPr}{at least 29\% more errors}  % experiments/FIGS/data/v2/index.json :: binary probe qcd_c_vs_light: (1 - AUC) of FLAV_F0 over L188, run by run :: sha256 c6c4138ee5c86423 :: prose twin: two significant figures, direction in words
\newcommand{\MlstCmpQcdCFzeroOverFullPrMax}{up to 35\% more errors}  % experiments/FIGS/data/v2/index.json :: binary probe qcd_c_vs_light: (1 - AUC) of FLAV_F0 over L188, run by run :: sha256 c6c4138ee5c86423 :: prose twin of the run farther from 1
\newcommand{\MlstCmpQcdCFoneOverFull}{1.237 and 1.220}  % experiments/FIGS/data/v2/index.json :: binary probe qcd_c_vs_light: (1 - AUC) of FLAV_F1 over L188, run by run :: sha256 c6c4138ee5c86423
\newcommand{\MlstCmpQcdCFoneOverFullPr}{at least 22\% more errors}  % experiments/FIGS/data/v2/index.json :: binary probe qcd_c_vs_light: (1 - AUC) of FLAV_F1 over L188, run by run :: sha256 c6c4138ee5c86423 :: prose twin: two significant figures, direction in words
\newcommand{\MlstCmpQcdCFoneOverFullPrMax}{up to 24\% more errors}  % experiments/FIGS/data/v2/index.json :: binary probe qcd_c_vs_light: (1 - AUC) of FLAV_F1 over L188, run by run :: sha256 c6c4138ee5c86423 :: prose twin of the run farther from 1
\newcommand{\MlstCmpQcdCFonerOverFull}{1.179 and 1.162}  % experiments/FIGS/data/v2/index.json :: binary probe qcd_c_vs_light: (1 - AUC) of FLAV_F1R over L188, run by run :: sha256 c6c4138ee5c86423
\newcommand{\MlstCmpQcdCFonerOverFullPr}{at least 16\% more errors}  % experiments/FIGS/data/v2/index.json :: binary probe qcd_c_vs_light: (1 - AUC) of FLAV_F1R over L188, run by run :: sha256 c6c4138ee5c86423 :: prose twin: two significant figures, direction in words
\newcommand{\MlstCmpQcdCFonerOverFullPrMax}{up to 18\% more errors}  % experiments/FIGS/data/v2/index.json :: binary probe qcd_c_vs_light: (1 - AUC) of FLAV_F1R over L188, run by run :: sha256 c6c4138ee5c86423 :: prose twin of the run farther from 1
\newcommand{\MlstLpEemmAucFull}{1.0000\,\ensuremath{\pm}\,\ensuremath{<}0.0001}  % experiments/FIGS/data/v2/index.json :: linear probe ee_vs_mm (fits_v4_sib): AUC, L188 :: sha256 c6c4138ee5c86423
\newcommand{\MlstLpEemmAucRes}{1.0000\,\ensuremath{\pm}\,\ensuremath{<}0.0001}  % experiments/FIGS/data/v2/index.json :: linear probe ee_vs_mm (fits_v4_sib): AUC, L162 :: sha256 c6c4138ee5c86423
\newcommand{\MlstLpEemmAucMid}{0.9610\,\ensuremath{\pm}\,0.0030}  % experiments/FIGS/data/v2/index.json :: linear probe ee_vs_mm (fits_v4_sib): AUC, R42_Q1 :: sha256 c6c4138ee5c86423
\newcommand{\MlstLpEemmAucCoarse}{0.960\,\ensuremath{\pm}\,0.020}  % experiments/FIGS/data/v2/index.json :: linear probe ee_vs_mm (fits_v4_sib): AUC, R16_Q1 :: sha256 c6c4138ee5c86423
\newcommand{\MlstLpEemmAucFzero}{1.0000 and 1.0000}  % experiments/FIGS/data/v2/index.json :: linear probe ee_vs_mm (fits_v4_sib): AUC, FLAV_F0 :: sha256 c6c4138ee5c86423
\newcommand{\MlstLpEemmAucFone}{1.0000 and 1.0000}  % experiments/FIGS/data/v2/index.json :: linear probe ee_vs_mm (fits_v4_sib): AUC, FLAV_F1 :: sha256 c6c4138ee5c86423
\newcommand{\MlstLpEemmAucFoner}{1.0000 and 1.0000}  % experiments/FIGS/data/v2/index.json :: linear probe ee_vs_mm (fits_v4_sib): AUC, FLAV_F1R :: sha256 c6c4138ee5c86423
\newcommand{\MlstLpEemmAucInit}{0.9995\,\ensuremath{\pm}\,0.0003}  % experiments/FIGS/data/v2/index.json :: linear probe ee_vs_mm (fits_v4_sib): AUC, INIT :: sha256 c6c4138ee5c86423
\newcommand{\MlstLpEemmAucInitBn}{1.0000\,\ensuremath{\pm}\,\ensuremath{<}0.0001}  % experiments/FIGS/data/v2/index.json :: linear probe ee_vs_mm (fits_v4_sib): AUC, INIT_BN :: sha256 c6c4138ee5c86423
\newcommand{\MlstLpEemmAucRandMerged}{1.0000\,\ensuremath{\pm}\,\ensuremath{<}0.0001}  % experiments/FIGS/data/v2/index.json :: linear probe ee_vs_mm (fits_v4_sib): AUC, the random partitions that merge the pair :: sha256 c6c4138ee5c86423
\newcommand{\MlstPairEemmNTrain}{9{,}571}  % experiments/FIGS/data/v2/index.json :: fits_v4_sib ee_vs_mm: n_train :: sha256 c6c4138ee5c86423
\newcommand{\MlstPairEemmNVal}{1{,}068}  % experiments/FIGS/data/v2/index.json :: fits_v4_sib ee_vs_mm: n_val :: sha256 c6c4138ee5c86423
\newcommand{\MlstPairEemmNTestPerSide}{128{,}833--135{,}245}  % experiments/FIGS/data/v2/index.json :: fits_v4_sib ee_vs_mm: n_test_signal and n_test_background :: sha256 c6c4138ee5c86423
\newcommand{\MlstLpJcTwoAucFull}{0.9801\,\ensuremath{\pm}\,\ensuremath{<}0.0001}  % experiments/FIGS/data/v2/index.json :: linear probe JetClass-II N=1000000: AUC averaged over the classes, L188 :: sha256 c6c4138ee5c86423
\newcommand{\MlstLpJcTwoAccFull}{0.4348\,\ensuremath{\pm}\,0.0003}  % experiments/FIGS/data/v2/index.json :: linear probe JetClass-II N=1000000: accuracy, L188 :: sha256 c6c4138ee5c86423
\newcommand{\MlstLpJcTwoAucRes}{0.9796\,\ensuremath{\pm}\,0.0001}  % experiments/FIGS/data/v2/index.json :: linear probe JetClass-II N=1000000: AUC averaged over the classes, L162 :: sha256 c6c4138ee5c86423
\newcommand{\MlstLpJcTwoAccRes}{0.4328\,\ensuremath{\pm}\,0.0008}  % experiments/FIGS/data/v2/index.json :: linear probe JetClass-II N=1000000: accuracy, L162 :: sha256 c6c4138ee5c86423
\newcommand{\MlstLpJcTwoAucMid}{0.9756\,\ensuremath{\pm}\,0.0002}  % experiments/FIGS/data/v2/index.json :: linear probe JetClass-II N=1000000: AUC averaged over the classes, R42_Q1 :: sha256 c6c4138ee5c86423
\newcommand{\MlstLpJcTwoAccMid}{0.361\,\ensuremath{\pm}\,0.005}  % experiments/FIGS/data/v2/index.json :: linear probe JetClass-II N=1000000: accuracy, R42_Q1 :: sha256 c6c4138ee5c86423
\newcommand{\MlstLpJcTwoAucCoarse}{0.9691\,\ensuremath{\pm}\,0.0004}  % experiments/FIGS/data/v2/index.json :: linear probe JetClass-II N=1000000: AUC averaged over the classes, R16_Q1 :: sha256 c6c4138ee5c86423
\newcommand{\MlstLpJcTwoAccCoarse}{0.334\,\ensuremath{\pm}\,0.007}  % experiments/FIGS/data/v2/index.json :: linear probe JetClass-II N=1000000: accuracy, R16_Q1 :: sha256 c6c4138ee5c86423
\newcommand{\MlstLpJcTwoAucResMass}{0.9791 and 0.9794}  % experiments/FIGS/data/v2/index.json :: linear probe JetClass-II N=1000000: AUC averaged over the classes, L162_MASS :: sha256 c6c4138ee5c86423
\newcommand{\MlstLpJcTwoAccResMass}{0.4271 and 0.4305}  % experiments/FIGS/data/v2/index.json :: linear probe JetClass-II N=1000000: accuracy, L162_MASS :: sha256 c6c4138ee5c86423
\newcommand{\MlstLpJcTwoAucCoarseMass}{0.9679 and 0.9696}  % experiments/FIGS/data/v2/index.json :: linear probe JetClass-II N=1000000: AUC averaged over the classes, R16_Q1_MASS :: sha256 c6c4138ee5c86423
\newcommand{\MlstLpJcTwoAccCoarseMass}{0.312 and 0.350}  % experiments/FIGS/data/v2/index.json :: linear probe JetClass-II N=1000000: accuracy, R16_Q1_MASS :: sha256 c6c4138ee5c86423
\newcommand{\MlstLpJcTwoAucFzero}{0.9717 and 0.9710}  % experiments/FIGS/data/v2/index.json :: linear probe JetClass-II N=1000000: AUC averaged over the classes, FLAV_F0 :: sha256 c6c4138ee5c86423
\newcommand{\MlstLpJcTwoAccFzero}{0.3769 and 0.3723}  % experiments/FIGS/data/v2/index.json :: linear probe JetClass-II N=1000000: accuracy, FLAV_F0 :: sha256 c6c4138ee5c86423
\newcommand{\MlstLpJcTwoAucFone}{0.9758 and 0.9759}  % experiments/FIGS/data/v2/index.json :: linear probe JetClass-II N=1000000: AUC averaged over the classes, FLAV_F1 :: sha256 c6c4138ee5c86423
\newcommand{\MlstLpJcTwoAccFone}{0.3945 and 0.3963}  % experiments/FIGS/data/v2/index.json :: linear probe JetClass-II N=1000000: accuracy, FLAV_F1 :: sha256 c6c4138ee5c86423
\newcommand{\MlstLpJcTwoAucFoner}{0.9760 and 0.9760}  % experiments/FIGS/data/v2/index.json :: linear probe JetClass-II N=1000000: AUC averaged over the classes, FLAV_F1R :: sha256 c6c4138ee5c86423
\newcommand{\MlstLpJcTwoAccFoner}{0.4071 and 0.4089}  % experiments/FIGS/data/v2/index.json :: linear probe JetClass-II N=1000000: accuracy, FLAV_F1R :: sha256 c6c4138ee5c86423
\newcommand{\MlstLpJcTwoAucInit}{0.8121\,\ensuremath{\pm}\,0.0019}  % experiments/FIGS/data/v2/index.json :: linear probe JetClass-II N=1000000: AUC averaged over the classes, INIT :: sha256 c6c4138ee5c86423
\newcommand{\MlstLpJcTwoAccInit}{0.0724\,\ensuremath{\pm}\,0.0015}  % experiments/FIGS/data/v2/index.json :: linear probe JetClass-II N=1000000: accuracy, INIT :: sha256 c6c4138ee5c86423
\newcommand{\MlstLpJcTwoAucInitBn}{0.866\,\ensuremath{\pm}\,0.004}  % experiments/FIGS/data/v2/index.json :: linear probe JetClass-II N=1000000: AUC averaged over the classes, INIT_BN :: sha256 c6c4138ee5c86423
\newcommand{\MlstLpJcTwoAccInitBn}{0.1017\,\ensuremath{\pm}\,0.0017}  % experiments/FIGS/data/v2/index.json :: linear probe JetClass-II N=1000000: accuracy, INIT_BN :: sha256 c6c4138ee5c86423
\newcommand{\MlstLpJcTwoAucRand}{0.9777\,\ensuremath{\pm}\,0.0002}  % experiments/FIGS/data/v2/index.json :: linear probe JetClass-II N=1000000: AUC averaged over the classes, the random partitions :: sha256 c6c4138ee5c86423
\newcommand{\MlstLpJcOneAccInitBn}{0.475\,\ensuremath{\pm}\,0.005}  % experiments/FIGS/data/v2/index.json :: linear probe JetClass N=1000000: accuracy, INIT_BN :: sha256 c6c4138ee5c86423
\newcommand{\MlstBnMaxDev}{0.2\%}  % experiments/FIGS/data/v2/index.json :: largest |r(best70_bn) / r(best70) - 1| of a JetClass-II ratio of 1 - AUC to the 188-class run, per cent :: sha256 c6c4138ee5c86423
\newcommand{\MlstBnMaxDevPair}{0.7\%}  % experiments/FIGS/data/v2/index.json :: largest |r(best70_bn) / r(best70) - 1| of a ratio of 1 - AUC on a class pair to the 188-class run, levels of the hierarchy, mass models and controlled partitions, per cent :: sha256 c6c4138ee5c86423 :: pairs with a check: bc_vs_bq_cs, bb_vs_cc, bbqq_vs_ccqq, qcd_bb_vs_cc, qcd_b_vs_light, qcd_c_vs_light; without: ee_vs_mm
\newcommand{\MlstGpuTypeBound}{1.4\%}  % experiments/FIGS/data/v2/index.json :: largest |label-set ratio of 1 - AUC of index 3 / mean of the other indices - 1|, JetClass-II largest set, per cent :: sha256 c6c4138ee5c86423 :: mtx-l188-s3 trained on NVIDIA GeForce RTX 3090; largest for R16_Q1; label sets L162 0.09%, R42_Q1 0.70%, R16_Q1 1.37%
\newcommand{\MlstGpuOtherSeedsDiff}{1.9\%}  % experiments/FIGS/data/v2/index.json :: largest (max - min) / mean of a label-set ratio of 1 - AUC over the indices other than 3, JetClass-II largest set, per cent :: sha256 c6c4138ee5c86423 :: label sets L162 0.90%, R42_Q1 1.14%, R16_Q1 1.93%
\newcommand{\MlstPerClassResNDoubled}{0}  % experiments/FIGS/data/v2/index.json :: JetClass-II classes whose mean ratio of 1 - AUC (L162 over 188) is at least 2 (count) :: sha256 c6c4138ee5c86423
\newcommand{\MlstPerClassResSeedSdMedian}{0.011}  % experiments/FIGS/data/v2/index.json :: median over classes of the SD over runs of the per-class ratio, L162 over 188 :: sha256 c6c4138ee5c86423
\newcommand{\MlstPerClassResTopClass}{QCD, \ensuremath{ccss}}  % experiments/FIGS/data/v2/index.json :: class with the largest mean ratio, L162 :: sha256 c6c4138ee5c86423
\newcommand{\MlstPerClassResTopRatio}{1.129\,\ensuremath{\pm}\,0.010}  % experiments/FIGS/data/v2/index.json :: per-class ratio of 1 - AUC of that class, L162 over 188, run by run :: sha256 c6c4138ee5c86423
\newcommand{\MlstPerClassResTopRatioPr}{13\% more errors}  % experiments/FIGS/data/v2/index.json :: per-class ratio of 1 - AUC of that class, L162 over 188, run by run :: sha256 c6c4138ee5c86423 :: prose twin: two significant figures, direction in words
\newcommand{\MlstPerClassResLeptonN}{0}  % experiments/FIGS/data/v2/index.json :: resonant classes sharing their L162 group with a decay that differs only in lepton flavor (count) :: sha256 c6c4138ee5c86423
\newcommand{\MlstPerClassResLeptonMedian}{--}  % experiments/FIGS/data/v2/index.json :: median over those classes of the mean per-class ratio, L162 over 188 :: sha256 c6c4138ee5c86423
\newcommand{\MlstPerClassResLeptonMedianPr}{\pending{no run yet}}  % experiments/FIGS/data/v2/index.json :: median over those classes of the mean per-class ratio, L162 over 188 :: sha256 c6c4138ee5c86423 :: prose twin: two significant figures, direction in words
\newcommand{\MlstPerClassResLeptonShare}{--}  % experiments/FIGS/data/v2/index.json :: their share of the summed ln ratio over the resonant classes, L162 :: sha256 c6c4138ee5c86423
\newcommand{\MlstPerClassResQuarkN}{0}  % experiments/FIGS/data/v2/index.json :: resonant classes sharing their L162 group with a decay that differs only in quark flavor (count) :: sha256 c6c4138ee5c86423
\newcommand{\MlstPerClassResQuarkMedian}{--}  % experiments/FIGS/data/v2/index.json :: median over those classes of the mean per-class ratio, L162 over 188 :: sha256 c6c4138ee5c86423
\newcommand{\MlstPerClassResQuarkMedianPr}{\pending{no run yet}}  % experiments/FIGS/data/v2/index.json :: median over those classes of the mean per-class ratio, L162 over 188 :: sha256 c6c4138ee5c86423 :: prose twin: two significant figures, direction in words
\newcommand{\MlstPerClassResQuarkShare}{--}  % experiments/FIGS/data/v2/index.json :: their share of the summed ln ratio over the resonant classes, L162 :: sha256 c6c4138ee5c86423
\newcommand{\MlstPerClassResQuarkLeptonN}{0}  % experiments/FIGS/data/v2/index.json :: resonant classes counted in both the quark- and the lepton-flavor set, L162 (count) :: sha256 c6c4138ee5c86423
\newcommand{\MlstPerClassMidNDoubled}{40}  % experiments/FIGS/data/v2/index.json :: JetClass-II classes whose mean ratio of 1 - AUC (R42_Q1 over 188) is at least 2 (count) :: sha256 c6c4138ee5c86423
\newcommand{\MlstPerClassMidSeedSdMedian}{0.023}  % experiments/FIGS/data/v2/index.json :: median over classes of the SD over runs of the per-class ratio, R42_Q1 over 188 :: sha256 c6c4138ee5c86423
\newcommand{\MlstPerClassMidTopClass}{\ensuremath{X\to YY\to bb\mu\mu}}  % experiments/FIGS/data/v2/index.json :: class with the largest mean ratio, R42_Q1 :: sha256 c6c4138ee5c86423
\newcommand{\MlstPerClassMidTopRatio}{4.06\,\ensuremath{\pm}\,0.14}  % experiments/FIGS/data/v2/index.json :: per-class ratio of 1 - AUC of that class, R42_Q1 over 188, run by run :: sha256 c6c4138ee5c86423
\newcommand{\MlstPerClassMidTopRatioPr}{4.1 times the errors}  % experiments/FIGS/data/v2/index.json :: per-class ratio of 1 - AUC of that class, R42_Q1 over 188, run by run :: sha256 c6c4138ee5c86423 :: prose twin: two significant figures, direction in words
\newcommand{\MlstPerClassMidLeptonN}{78}  % experiments/FIGS/data/v2/index.json :: resonant classes sharing their R42_Q1 group with a decay that differs only in lepton flavor (count) :: sha256 c6c4138ee5c86423
\newcommand{\MlstPerClassMidLeptonMedian}{2.02}  % experiments/FIGS/data/v2/index.json :: median over those classes of the mean per-class ratio, R42_Q1 over 188 :: sha256 c6c4138ee5c86423
\newcommand{\MlstPerClassMidLeptonMedianPr}{2.0 times the errors}  % experiments/FIGS/data/v2/index.json :: median over those classes of the mean per-class ratio, R42_Q1 over 188 :: sha256 c6c4138ee5c86423 :: prose twin: two significant figures, direction in words
\newcommand{\MlstPerClassMidLeptonShare}{90\%}  % experiments/FIGS/data/v2/index.json :: their share of the summed ln ratio over the resonant classes, R42_Q1 :: sha256 c6c4138ee5c86423
\newcommand{\MlstPerClassMidQuarkN}{111}  % experiments/FIGS/data/v2/index.json :: resonant classes sharing their R42_Q1 group with a decay that differs only in quark flavor (count) :: sha256 c6c4138ee5c86423
\newcommand{\MlstPerClassMidQuarkMedian}{1.14}  % experiments/FIGS/data/v2/index.json :: median over those classes of the mean per-class ratio, R42_Q1 over 188 :: sha256 c6c4138ee5c86423
\newcommand{\MlstPerClassMidQuarkMedianPr}{14\% more errors}  % experiments/FIGS/data/v2/index.json :: median over those classes of the mean per-class ratio, R42_Q1 over 188 :: sha256 c6c4138ee5c86423 :: prose twin: two significant figures, direction in words
\newcommand{\MlstPerClassMidQuarkShare}{53\%}  % experiments/FIGS/data/v2/index.json :: their share of the summed ln ratio over the resonant classes, R42_Q1 :: sha256 c6c4138ee5c86423
\newcommand{\MlstPerClassMidQuarkLeptonN}{44}  % experiments/FIGS/data/v2/index.json :: resonant classes counted in both the quark- and the lepton-flavor set, R42_Q1 (count) :: sha256 c6c4138ee5c86423
\newcommand{\MlstPerClassCoarseNDoubled}{60}  % experiments/FIGS/data/v2/index.json :: JetClass-II classes whose mean ratio of 1 - AUC (R16_Q1 over 188) is at least 2 (count) :: sha256 c6c4138ee5c86423
\newcommand{\MlstPerClassCoarseSeedSdMedian}{0.055}  % experiments/FIGS/data/v2/index.json :: median over classes of the SD over runs of the per-class ratio, R16_Q1 over 188 :: sha256 c6c4138ee5c86423
\newcommand{\MlstPerClassCoarseTopClass}{\ensuremath{X\to YY\to bbee}}  % experiments/FIGS/data/v2/index.json :: class with the largest mean ratio, R16_Q1 :: sha256 c6c4138ee5c86423
\newcommand{\MlstPerClassCoarseTopRatio}{5.2\,\ensuremath{\pm}\,0.4}  % experiments/FIGS/data/v2/index.json :: per-class ratio of 1 - AUC of that class, R16_Q1 over 188, run by run :: sha256 c6c4138ee5c86423
\newcommand{\MlstPerClassCoarseTopRatioPr}{5.2 times the errors}  % experiments/FIGS/data/v2/index.json :: per-class ratio of 1 - AUC of that class, R16_Q1 over 188, run by run :: sha256 c6c4138ee5c86423 :: prose twin: two significant figures, direction in words
\newcommand{\MlstPerClassCoarseLeptonN}{78}  % experiments/FIGS/data/v2/index.json :: resonant classes sharing their R16_Q1 group with a decay that differs only in lepton flavor (count) :: sha256 c6c4138ee5c86423
\newcommand{\MlstPerClassCoarseLeptonMedian}{2.27}  % experiments/FIGS/data/v2/index.json :: median over those classes of the mean per-class ratio, R16_Q1 over 188 :: sha256 c6c4138ee5c86423
\newcommand{\MlstPerClassCoarseLeptonMedianPr}{2.3 times the errors}  % experiments/FIGS/data/v2/index.json :: median over those classes of the mean per-class ratio, R16_Q1 over 188 :: sha256 c6c4138ee5c86423 :: prose twin: two significant figures, direction in words
\newcommand{\MlstPerClassCoarseLeptonShare}{68\%}  % experiments/FIGS/data/v2/index.json :: their share of the summed ln ratio over the resonant classes, R16_Q1 :: sha256 c6c4138ee5c86423
\newcommand{\MlstPerClassCoarseQuarkN}{141}  % experiments/FIGS/data/v2/index.json :: resonant classes sharing their R16_Q1 group with a decay that differs only in quark flavor (count) :: sha256 c6c4138ee5c86423
\newcommand{\MlstPerClassCoarseQuarkMedian}{1.65}  % experiments/FIGS/data/v2/index.json :: median over those classes of the mean per-class ratio, R16_Q1 over 188 :: sha256 c6c4138ee5c86423
\newcommand{\MlstPerClassCoarseQuarkMedianPr}{65\% more errors}  % experiments/FIGS/data/v2/index.json :: median over those classes of the mean per-class ratio, R16_Q1 over 188 :: sha256 c6c4138ee5c86423 :: prose twin: two significant figures, direction in words
\newcommand{\MlstPerClassCoarseQuarkShare}{87\%}  % experiments/FIGS/data/v2/index.json :: their share of the summed ln ratio over the resonant classes, R16_Q1 :: sha256 c6c4138ee5c86423
\newcommand{\MlstPerClassCoarseQuarkLeptonN}{64}  % experiments/FIGS/data/v2/index.json :: resonant classes counted in both the quark- and the lepton-flavor set, R16_Q1 (count) :: sha256 c6c4138ee5c86423
\newcommand{\MlstPerClassRandFlavMergedOverApart}{1.010\,\ensuremath{\pm}\,0.008}  % experiments/FIGS/data/v2/index.json :: random partitions, linear probe: median per-class ratio of the classes merged with a decay that differs only in quark flavor over that of the classes kept apart from all of them, mean +- SD over the partitions :: sha256 c6c4138ee5c86423
\newcommand{\MlstPerClassRandFlavMergedOverApartPr}{equal within 1\%}  % experiments/FIGS/data/v2/index.json :: random partitions, linear probe: median per-class ratio of the classes merged with a decay that differs only in quark flavor over that of the classes kept apart from all of them, mean +- SD over the partitions :: sha256 c6c4138ee5c86423 :: prose twin: two significant figures, direction in words
\newcommand{\MlstPerClassRandLeptonMergedOverApart}{0.997\,\ensuremath{\pm}\,0.020}  % experiments/FIGS/data/v2/index.json :: random partitions, linear probe: median per-class ratio of the classes merged with a decay that differs only in lepton flavor over that of the classes kept apart from all of them, mean +- SD over the partitions :: sha256 c6c4138ee5c86423
\newcommand{\MlstPerClassRandLeptonMergedOverApartPr}{equal within 3\%}  % experiments/FIGS/data/v2/index.json :: random partitions, linear probe: median per-class ratio of the classes merged with a decay that differs only in lepton flavor over that of the classes kept apart from all of them, mean +- SD over the partitions :: sha256 c6c4138ee5c86423 :: prose twin: two significant figures, direction in words
\newcommand{\MlstMlpNatPerClassRandFlavMergedOverApart}{1.003\,\ensuremath{\pm}\,0.007}  % experiments/FIGS/data/v2/index.json :: random partitions, nonlinear probe, natural composition: median per-class ratio of the classes merged with a decay that differs only in quark flavor over that of the classes kept apart from all of them, mean +- SD over the partitions :: sha256 c6c4138ee5c86423
\newcommand{\MlstMlpNatPerClassRandFlavMergedOverApartPr}{equal within 1\%}  % experiments/FIGS/data/v2/index.json :: random partitions, nonlinear probe, natural composition: median per-class ratio of the classes merged with a decay that differs only in quark flavor over that of the classes kept apart from all of them, mean +- SD over the partitions :: sha256 c6c4138ee5c86423 :: prose twin: two significant figures, direction in words
\newcommand{\MlstMlpNatPerClassRandLeptonMergedOverApart}{1.01\,\ensuremath{\pm}\,0.07}  % experiments/FIGS/data/v2/index.json :: random partitions, nonlinear probe, natural composition: median per-class ratio of the classes merged with a decay that differs only in lepton flavor over that of the classes kept apart from all of them, mean +- SD over the partitions :: sha256 c6c4138ee5c86423
\newcommand{\MlstMlpNatPerClassRandLeptonMergedOverApartPr}{equal within 8\%}  % experiments/FIGS/data/v2/index.json :: random partitions, nonlinear probe, natural composition: median per-class ratio of the classes merged with a decay that differs only in lepton flavor over that of the classes kept apart from all of them, mean +- SD over the partitions :: sha256 c6c4138ee5c86423 :: prose twin: two significant figures, direction in words
\newcommand{\MlstPerClassAucDiffResSeedSdMedian}{0.015}  % experiments/FIGS/data/v2/index.json :: median over classes of the SD over indices of 100 (AUC_c of L162 - AUC_c of 188), pp :: sha256 c6c4138ee5c86423
\newcommand{\MlstPerClassAucDiffResMin}{\ensuremath{-}0.73}  % experiments/FIGS/data/v2/index.json :: most negative mean over indices of 100 (AUC_c of L162 - AUC_c of 188), pp :: sha256 c6c4138ee5c86423 :: label_QCD_bcss
\newcommand{\MlstPerClassAucDiffResMinClass}{QCD, \ensuremath{bcss}}  % experiments/FIGS/data/v2/index.json :: its class, L162 :: sha256 c6c4138ee5c86423
\newcommand{\MlstPerClassAucDiffMidSeedSdMedian}{0.032}  % experiments/FIGS/data/v2/index.json :: median over classes of the SD over indices of 100 (AUC_c of R42_Q1 - AUC_c of 188), pp :: sha256 c6c4138ee5c86423
\newcommand{\MlstPerClassAucDiffMidMin}{\ensuremath{-}1.4}  % experiments/FIGS/data/v2/index.json :: most negative mean over indices of 100 (AUC_c of R42_Q1 - AUC_c of 188), pp :: sha256 c6c4138ee5c86423 :: label_X_YY_cqtauev
\newcommand{\MlstPerClassAucDiffMidMinClass}{\ensuremath{X\to YY\to cq\tau_e\nu}}  % experiments/FIGS/data/v2/index.json :: its class, R42_Q1 :: sha256 c6c4138ee5c86423
\newcommand{\MlstPerClassAucDiffCoarseSeedSdMedian}{0.077}  % experiments/FIGS/data/v2/index.json :: median over classes of the SD over indices of 100 (AUC_c of R16_Q1 - AUC_c of 188), pp :: sha256 c6c4138ee5c86423
\newcommand{\MlstPerClassAucDiffCoarseMin}{\ensuremath{-}3.1}  % experiments/FIGS/data/v2/index.json :: most negative mean over indices of 100 (AUC_c of R16_Q1 - AUC_c of 188), pp :: sha256 c6c4138ee5c86423 :: label_QCD_bcss
\newcommand{\MlstPerClassAucDiffCoarseMinClass}{QCD, \ensuremath{bcss}}  % experiments/FIGS/data/v2/index.json :: its class, R16_Q1 :: sha256 c6c4138ee5c86423
\newcommand{\MlstRocJcTwoHurtList}{\ensuremath{X\to YY\to ccq}; \ensuremath{X\to YY\to ssb}; \ensuremath{X\to YY\to ccss}; QCD, \ensuremath{bcc}; QCD, \ensuremath{bccss}; QCD, \ensuremath{bss}; QCD, \ensuremath{ccs} and QCD, \ensuremath{ccss}}  % experiments/FIGS/data/v2/index.json :: jc2_roc_classes on the class-balanced probes; candidates with at least 1000 test jets :: sha256 c6c4138ee5c86423 :: label_X_YY_ccq, label_X_YY_ssb, label_X_YY_ccss, label_QCD_bcc, label_QCD_bccss, label_QCD_bss, label_QCD_ccs, label_QCD_ccss
\newcommand{\MlstRocJcTwoStableList}{\ensuremath{X\to gg}; \ensuremath{X\to YY\to ggg}; \ensuremath{X\to YY\to g\tau_h\tau_h} and QCD, light}  % experiments/FIGS/data/v2/index.json :: jc2_roc_classes on the class-balanced probes; candidates with at least 1000 test jets :: sha256 c6c4138ee5c86423 :: label_X_gg, label_X_YY_ggg, label_X_YY_gtauhtauh, label_QCD_light
\newcommand{\MlstRocJcTwoHurtMinDrop}{2.1}  % experiments/FIGS/data/v2/index.json :: jc2_roc_classes on the class-balanced probes; candidates with at least 1000 test jets :: sha256 c6c4138ee5c86423 :: smallest AUC loss at 17 classes of the hurt, pp
\newcommand{\MlstRocJcTwoHurtMaxDrop}{2.8}  % experiments/FIGS/data/v2/index.json :: jc2_roc_classes on the class-balanced probes; candidates with at least 1000 test jets :: sha256 c6c4138ee5c86423 :: largest AUC loss at 17 classes of the hurt, pp
\newcommand{\MlstRocJcTwoStableMaxChange}{15}  % experiments/FIGS/data/v2/index.json :: jc2_roc_classes on the class-balanced probes; candidates with at least 1000 test jets :: sha256 c6c4138ee5c86423 :: largest change of 1 - AUC at any label set of the stable classes, per cent of the 188-class value
\newcommand{\MlstJcBalNThin}{4}  % experiments/FIGS/data/v2/jc2bal/census.json :: counts.ft_test :: sha256 a4379ee0817254cf :: classes with fewer than 1000 class-balanced test jets (F3 hatching)
\newcommand{\MlstMinBkgPass}{20}  % experiments/FIGS/mlst_v4.py :: JC1_MIN_BKG :: sha256 fa536b4866ffd651 :: fewest background jets behind a drawn ROC point or a quoted rejection
\newcommand{\MlstMassMinJets}{100}  % experiments/EVAL/atlas_resolution.py :: MIN_JETS :: sha256 b82cc219f79f0f14 :: jets below which a mass bin is null
\newcommand{\MlstRocCutBbFull}{0.79}  % experiments/FIGS/data/v2/index.json :: curves_v2/pairs curve bb_vs_cc: largest cut.eps_s over the L188 runs ('not reached' when a run's cut is null or leaves fewer than 2 grid points) :: sha256 c6c4138ee5c86423
\newcommand{\MlstRocCutBbRes}{0.77}  % experiments/FIGS/data/v2/index.json :: curves_v2/pairs curve bb_vs_cc: largest cut.eps_s over the L162 runs ('not reached' when a run's cut is null or leaves fewer than 2 grid points) :: sha256 c6c4138ee5c86423
\newcommand{\MlstRocCutBbMid}{0.79}  % experiments/FIGS/data/v2/index.json :: curves_v2/pairs curve bb_vs_cc: largest cut.eps_s over the R42_Q1 runs ('not reached' when a run's cut is null or leaves fewer than 2 grid points) :: sha256 c6c4138ee5c86423
\newcommand{\MlstRocCutBbCoarse}{0.13}  % experiments/FIGS/data/v2/index.json :: curves_v2/pairs curve bb_vs_cc: largest cut.eps_s over the R16_Q1 runs ('not reached' when a run's cut is null or leaves fewer than 2 grid points) :: sha256 c6c4138ee5c86423
\newcommand{\MlstRocCutBbFzero}{0.03}  % experiments/FIGS/data/v2/index.json :: curves_v2/pairs curve bb_vs_cc: largest cut.eps_s over the FLAV_F0 runs ('not reached' when a run's cut is null or leaves fewer than 2 grid points) :: sha256 c6c4138ee5c86423
\newcommand{\MlstRocCutBbFone}{0.62}  % experiments/FIGS/data/v2/index.json :: curves_v2/pairs curve bb_vs_cc: largest cut.eps_s over the FLAV_F1 runs ('not reached' when a run's cut is null or leaves fewer than 2 grid points) :: sha256 c6c4138ee5c86423
\newcommand{\MlstRocCutBbFoner}{0.51}  % experiments/FIGS/data/v2/index.json :: curves_v2/pairs curve bb_vs_cc: largest cut.eps_s over the FLAV_F1R runs ('not reached' when a run's cut is null or leaves fewer than 2 grid points) :: sha256 c6c4138ee5c86423
\newcommand{\MlstRocNBkgBb}{200{,}000}  % experiments/FIGS/data/v2/index.json :: curves_v2/pairs curve bb_vs_cc: n_background :: sha256 c6c4138ee5c86423
\newcommand{\MlstRocRejMaxBb}{10{,}000}  % experiments/FIGS/data/v2/index.json :: curves_v2/pairs curve bb_vs_cc: n_background / 20 :: sha256 c6c4138ee5c86423
\newcommand{\MlstRocCutBbqqFull}{0.67}  % experiments/FIGS/data/v2/index.json :: curves_v2/pairs curve bbqq_vs_ccqq: largest cut.eps_s over the L188 runs ('not reached' when a run's cut is null or leaves fewer than 2 grid points) :: sha256 c6c4138ee5c86423
\newcommand{\MlstRocCutBbqqRes}{0.66}  % experiments/FIGS/data/v2/index.json :: curves_v2/pairs curve bbqq_vs_ccqq: largest cut.eps_s over the L162 runs ('not reached' when a run's cut is null or leaves fewer than 2 grid points) :: sha256 c6c4138ee5c86423
\newcommand{\MlstRocCutBbqqMid}{0.62}  % experiments/FIGS/data/v2/index.json :: curves_v2/pairs curve bbqq_vs_ccqq: largest cut.eps_s over the R42_Q1 runs ('not reached' when a run's cut is null or leaves fewer than 2 grid points) :: sha256 c6c4138ee5c86423
\newcommand{\MlstRocCutBbqqCoarse}{0.06}  % experiments/FIGS/data/v2/index.json :: curves_v2/pairs curve bbqq_vs_ccqq: largest cut.eps_s over the R16_Q1 runs ('not reached' when a run's cut is null or leaves fewer than 2 grid points) :: sha256 c6c4138ee5c86423
\newcommand{\MlstRocCutBbqqFzero}{0.06}  % experiments/FIGS/data/v2/index.json :: curves_v2/pairs curve bbqq_vs_ccqq: largest cut.eps_s over the FLAV_F0 runs ('not reached' when a run's cut is null or leaves fewer than 2 grid points) :: sha256 c6c4138ee5c86423
\newcommand{\MlstRocCutBbqqFone}{0.56}  % experiments/FIGS/data/v2/index.json :: curves_v2/pairs curve bbqq_vs_ccqq: largest cut.eps_s over the FLAV_F1 runs ('not reached' when a run's cut is null or leaves fewer than 2 grid points) :: sha256 c6c4138ee5c86423
\newcommand{\MlstRocCutBbqqFoner}{0.32}  % experiments/FIGS/data/v2/index.json :: curves_v2/pairs curve bbqq_vs_ccqq: largest cut.eps_s over the FLAV_F1R runs ('not reached' when a run's cut is null or leaves fewer than 2 grid points) :: sha256 c6c4138ee5c86423
\newcommand{\MlstRocNBkgBbqq}{200{,}000}  % experiments/FIGS/data/v2/index.json :: curves_v2/pairs curve bbqq_vs_ccqq: n_background :: sha256 c6c4138ee5c86423
\newcommand{\MlstRocRejMaxBbqq}{10{,}000}  % experiments/FIGS/data/v2/index.json :: curves_v2/pairs curve bbqq_vs_ccqq: n_background / 20 :: sha256 c6c4138ee5c86423
\newcommand{\MlstRocCutBcFull}{0.24}  % experiments/FIGS/data/v2/index.json :: curves_v2/pairs curve bc_vs_bq_cs: largest cut.eps_s over the L188 runs ('not reached' when a run's cut is null or leaves fewer than 2 grid points) :: sha256 c6c4138ee5c86423
\newcommand{\MlstRocCutBcRes}{0.25}  % experiments/FIGS/data/v2/index.json :: curves_v2/pairs curve bc_vs_bq_cs: largest cut.eps_s over the L162 runs ('not reached' when a run's cut is null or leaves fewer than 2 grid points) :: sha256 c6c4138ee5c86423
\newcommand{\MlstRocCutBcMid}{0.14}  % experiments/FIGS/data/v2/index.json :: curves_v2/pairs curve bc_vs_bq_cs: largest cut.eps_s over the R42_Q1 runs ('not reached' when a run's cut is null or leaves fewer than 2 grid points) :: sha256 c6c4138ee5c86423
\newcommand{\MlstRocCutBcCoarse}{0.00}  % experiments/FIGS/data/v2/index.json :: curves_v2/pairs curve bc_vs_bq_cs: largest cut.eps_s over the R16_Q1 runs ('not reached' when a run's cut is null or leaves fewer than 2 grid points) :: sha256 c6c4138ee5c86423
\newcommand{\MlstRocCutBcFzero}{0.00}  % experiments/FIGS/data/v2/index.json :: curves_v2/pairs curve bc_vs_bq_cs: largest cut.eps_s over the FLAV_F0 runs ('not reached' when a run's cut is null or leaves fewer than 2 grid points) :: sha256 c6c4138ee5c86423
\newcommand{\MlstRocCutBcFone}{0.01}  % experiments/FIGS/data/v2/index.json :: curves_v2/pairs curve bc_vs_bq_cs: largest cut.eps_s over the FLAV_F1 runs ('not reached' when a run's cut is null or leaves fewer than 2 grid points) :: sha256 c6c4138ee5c86423
\newcommand{\MlstRocCutBcFoner}{0.01}  % experiments/FIGS/data/v2/index.json :: curves_v2/pairs curve bc_vs_bq_cs: largest cut.eps_s over the FLAV_F1R runs ('not reached' when a run's cut is null or leaves fewer than 2 grid points) :: sha256 c6c4138ee5c86423
\newcommand{\MlstRocNBkgBc}{200{,}000}  % experiments/FIGS/data/v2/index.json :: curves_v2/pairs curve bc_vs_bq_cs: n_background :: sha256 c6c4138ee5c86423
\newcommand{\MlstRocRejMaxBc}{10{,}000}  % experiments/FIGS/data/v2/index.json :: curves_v2/pairs curve bc_vs_bq_cs: n_background / 20 :: sha256 c6c4138ee5c86423
\newcommand{\MlstRocCutEemmFull}{not reached}  % experiments/FIGS/data/v2/index.json :: curves_v2/pairs curve ee_vs_mm: largest cut.eps_s over the L188 runs ('not reached' when a run's cut is null or leaves fewer than 2 grid points) :: sha256 c6c4138ee5c86423
\newcommand{\MlstRocCutEemmRes}{not reached}  % experiments/FIGS/data/v2/index.json :: curves_v2/pairs curve ee_vs_mm: largest cut.eps_s over the L162 runs ('not reached' when a run's cut is null or leaves fewer than 2 grid points) :: sha256 c6c4138ee5c86423
\newcommand{\MlstRocCutEemmMid}{0.00}  % experiments/FIGS/data/v2/index.json :: curves_v2/pairs curve ee_vs_mm: largest cut.eps_s over the R42_Q1 runs ('not reached' when a run's cut is null or leaves fewer than 2 grid points) :: sha256 c6c4138ee5c86423
\newcommand{\MlstRocCutEemmCoarse}{0.00}  % experiments/FIGS/data/v2/index.json :: curves_v2/pairs curve ee_vs_mm: largest cut.eps_s over the R16_Q1 runs ('not reached' when a run's cut is null or leaves fewer than 2 grid points) :: sha256 c6c4138ee5c86423
\newcommand{\MlstRocCutEemmFzero}{not reached}  % experiments/FIGS/data/v2/index.json :: curves_v2/pairs curve ee_vs_mm: largest cut.eps_s over the FLAV_F0 runs ('not reached' when a run's cut is null or leaves fewer than 2 grid points) :: sha256 c6c4138ee5c86423
\newcommand{\MlstRocCutEemmFone}{not reached}  % experiments/FIGS/data/v2/index.json :: curves_v2/pairs curve ee_vs_mm: largest cut.eps_s over the FLAV_F1 runs ('not reached' when a run's cut is null or leaves fewer than 2 grid points) :: sha256 c6c4138ee5c86423
\newcommand{\MlstRocCutEemmFoner}{not reached}  % experiments/FIGS/data/v2/index.json :: curves_v2/pairs curve ee_vs_mm: largest cut.eps_s over the FLAV_F1R runs ('not reached' when a run's cut is null or leaves fewer than 2 grid points) :: sha256 c6c4138ee5c86423
\newcommand{\MlstRocNBkgEemm}{128{,}833}  % experiments/FIGS/data/v2/index.json :: curves_v2/pairs curve ee_vs_mm: n_background :: sha256 c6c4138ee5c86423
\newcommand{\MlstRocRejMaxEemm}{6{,}442}  % experiments/FIGS/data/v2/index.json :: curves_v2/pairs curve ee_vs_mm: n_background / 20 :: sha256 c6c4138ee5c86423
\newcommand{\MlstRocNoCurveEemm}{the \Kfull- and \Kres-class models and the three random partitions with controlled flavor boundaries}  % experiments/FIGS/data/v2/index.json :: curves_v2/pairs curve ee_vs_mm: configurations of F4 with a run whose cut is null or leaves fewer than 2 grid points (F4 draws no curve for them) :: sha256 c6c4138ee5c86423 :: L188, L162, FLAV_F0, FLAV_F1, FLAV_F1R
\newcommand{\MlstRocCutJcOneHbbFull}{0.40}  % experiments/FIGS/data/v2/index.json :: curves_v2/jc1_ft curve label_Hbb: largest cut.eps_s over the 188-class runs ('not reached' when a run's cut is null or leaves fewer than 2 grid points) :: sha256 c6c4138ee5c86423
\newcommand{\MlstRocCutJcOneHccFull}{0.37}  % experiments/FIGS/data/v2/index.json :: curves_v2/jc1_ft curve label_Hcc: largest cut.eps_s over the 188-class runs ('not reached' when a run's cut is null or leaves fewer than 2 grid points) :: sha256 c6c4138ee5c86423
\newcommand{\MlstRocCutJcOneHggFull}{0.04}  % experiments/FIGS/data/v2/index.json :: curves_v2/jc1_ft curve label_Hgg: largest cut.eps_s over the 188-class runs ('not reached' when a run's cut is null or leaves fewer than 2 grid points) :: sha256 c6c4138ee5c86423
\newcommand{\MlstRocCutJcOneHfourqFull}{0.26}  % experiments/FIGS/data/v2/index.json :: curves_v2/jc1_ft curve label_H4q: largest cut.eps_s over the 188-class runs ('not reached' when a run's cut is null or leaves fewer than 2 grid points) :: sha256 c6c4138ee5c86423
\newcommand{\MlstRocCutJcOneHqqlFull}{0.98}  % experiments/FIGS/data/v2/index.json :: curves_v2/jc1_ft curve label_Hqql: largest cut.eps_s over the 188-class runs ('not reached' when a run's cut is null or leaves fewer than 2 grid points) :: sha256 c6c4138ee5c86423
\newcommand{\MlstRocCutJcOneTbqqFull}{0.53}  % experiments/FIGS/data/v2/index.json :: curves_v2/jc1_ft curve label_Tbqq: largest cut.eps_s over the 188-class runs ('not reached' when a run's cut is null or leaves fewer than 2 grid points) :: sha256 c6c4138ee5c86423
\newcommand{\MlstRocCutJcOneTblFull}{not reached}  % experiments/FIGS/data/v2/index.json :: curves_v2/jc1_ft curve label_Tbl: largest cut.eps_s over the 188-class runs ('not reached' when a run's cut is null or leaves fewer than 2 grid points) :: sha256 c6c4138ee5c86423
\newcommand{\MlstRocCutJcOneWqqFull}{0.15}  % experiments/FIGS/data/v2/index.json :: curves_v2/jc1_ft curve label_Wqq: largest cut.eps_s over the 188-class runs ('not reached' when a run's cut is null or leaves fewer than 2 grid points) :: sha256 c6c4138ee5c86423
\newcommand{\MlstRocCutJcOneZqqFull}{0.12}  % experiments/FIGS/data/v2/index.json :: curves_v2/jc1_ft curve label_Zqq: largest cut.eps_s over the 188-class runs ('not reached' when a run's cut is null or leaves fewer than 2 grid points) :: sha256 c6c4138ee5c86423
\newcommand{\MlstFtTopRfiftyCoarseOverFull}{0.932\,\ensuremath{\pm}\,0.018}  % experiments/FIGS/data/v2/index.json :: curves_v2/top_ft curve top: rejection at eps_S = 0.5 of R16_Q1 over L188, run by run :: sha256 c6c4138ee5c86423
\newcommand{\MlstFtTopRfiftyCoarseOverFullPr}{6.8\% less rejection}  % experiments/FIGS/data/v2/index.json :: curves_v2/top_ft curve top: rejection at eps_S = 0.5 of R16_Q1 over L188, run by run :: sha256 c6c4138ee5c86423 :: prose twin: two significant figures, direction in words
\newcommand{\MlstProbeUnconverged}{0}  % experiments/FIGS/data/v2/index.json :: linear probes of pretrained models that did not converge (count) :: sha256 c6c4138ee5c86423
\newcommand{\MlstProbeUnconvergedInit}{0}  % experiments/FIGS/data/v2/index.json :: linear probes of the randomly initialized networks that did not converge (count) :: sha256 c6c4138ee5c86423
\newcommand{\MlstProbeFits}{895}  % experiments/FIGS/data/v2/index.json :: selected linear-probe fits of the reported checkpoints (count) :: sha256 c6c4138ee5c86423
\newcommand{\MlstProbeEdgeC}{16}  % experiments/FIGS/data/v2/index.json :: of these, fits whose selected C is the largest of the grid (count) :: sha256 c6c4138ee5c86423
\newcommand{\MlstMlpUnconverged}{0}  % experiments/FIGS/data/v2/index.json :: nonlinear probes that did not reach their plateau within the epoch cap (count) :: sha256 c6c4138ee5c86423
\newcommand{\MlstMlpCmpJcTwoResOverFull}{1.019\,\ensuremath{\pm}\,0.007}  % experiments/FIGS/data/v2/index.json :: nonlinear probe jc2_188: (1 - AUC) of L162 over L188, run by run :: sha256 c6c4138ee5c86423
\newcommand{\MlstMlpCmpJcTwoResOverFullPr}{1.9\% more errors}  % experiments/FIGS/data/v2/index.json :: nonlinear probe jc2_188: (1 - AUC) of L162 over L188, run by run :: sha256 c6c4138ee5c86423 :: prose twin: two significant figures, direction in words
\newcommand{\MlstMlpCmpJcTwoMidOverFull}{1.170\,\ensuremath{\pm}\,0.006}  % experiments/FIGS/data/v2/index.json :: nonlinear probe jc2_188: (1 - AUC) of R42_Q1 over L188, run by run :: sha256 c6c4138ee5c86423
\newcommand{\MlstMlpCmpJcTwoMidOverFullPr}{17\% more errors}  % experiments/FIGS/data/v2/index.json :: nonlinear probe jc2_188: (1 - AUC) of R42_Q1 over L188, run by run :: sha256 c6c4138ee5c86423 :: prose twin: two significant figures, direction in words
\newcommand{\MlstMlpCmpJcTwoCoarseOverFull}{1.388\,\ensuremath{\pm}\,0.021}  % experiments/FIGS/data/v2/index.json :: nonlinear probe jc2_188: (1 - AUC) of R16_Q1 over L188, run by run :: sha256 c6c4138ee5c86423
\newcommand{\MlstMlpCmpJcTwoCoarseOverFullPr}{39\% more errors}  % experiments/FIGS/data/v2/index.json :: nonlinear probe jc2_188: (1 - AUC) of R16_Q1 over L188, run by run :: sha256 c6c4138ee5c86423 :: prose twin: two significant figures, direction in words
\newcommand{\MlstMlpCmpJcTwoResMassOverRes}{\pending{no run yet}}  % experiments/FIGS/data/v2/index.json :: nonlinear probe jc2_188: (1 - AUC) of L162_MASS over L162, run by run :: sha256 c6c4138ee5c86423
\newcommand{\MlstMlpCmpJcTwoResMassOverResPr}{\pending{no run yet}}  % experiments/FIGS/data/v2/index.json :: nonlinear probe jc2_188: (1 - AUC) of L162_MASS over L162, run by run :: sha256 c6c4138ee5c86423 :: prose twin: two significant figures, direction in words
\newcommand{\MlstMlpCmpJcTwoResMassOverResPrMax}{\pending{no run yet}}  % experiments/FIGS/data/v2/index.json :: nonlinear probe jc2_188: (1 - AUC) of L162_MASS over L162, run by run :: sha256 c6c4138ee5c86423 :: prose twin of the run farther from 1
\newcommand{\MlstMlpCmpJcTwoCoarseMassOverCoarse}{\pending{no run yet}}  % experiments/FIGS/data/v2/index.json :: nonlinear probe jc2_188: (1 - AUC) of R16_Q1_MASS over R16_Q1, run by run :: sha256 c6c4138ee5c86423
\newcommand{\MlstMlpCmpJcTwoCoarseMassOverCoarsePr}{\pending{no run yet}}  % experiments/FIGS/data/v2/index.json :: nonlinear probe jc2_188: (1 - AUC) of R16_Q1_MASS over R16_Q1, run by run :: sha256 c6c4138ee5c86423 :: prose twin: two significant figures, direction in words
\newcommand{\MlstMlpCmpJcTwoCoarseMassOverCoarsePrMax}{\pending{no run yet}}  % experiments/FIGS/data/v2/index.json :: nonlinear probe jc2_188: (1 - AUC) of R16_Q1_MASS over R16_Q1, run by run :: sha256 c6c4138ee5c86423 :: prose twin of the run farther from 1
\newcommand{\MlstMlpCmpJcTwoFzeroOverFull}{1.294 and 1.310}  % experiments/FIGS/data/v2/index.json :: nonlinear probe jc2_188: (1 - AUC) of FLAV_F0 over L188, run by run :: sha256 c6c4138ee5c86423
\newcommand{\MlstMlpCmpJcTwoFzeroOverFullPr}{at least 29\% more errors}  % experiments/FIGS/data/v2/index.json :: nonlinear probe jc2_188: (1 - AUC) of FLAV_F0 over L188, run by run :: sha256 c6c4138ee5c86423 :: prose twin: two significant figures, direction in words
\newcommand{\MlstMlpCmpJcTwoFzeroOverFullPrMax}{up to 31\% more errors}  % experiments/FIGS/data/v2/index.json :: nonlinear probe jc2_188: (1 - AUC) of FLAV_F0 over L188, run by run :: sha256 c6c4138ee5c86423 :: prose twin of the run farther from 1
\newcommand{\MlstMlpCmpJcTwoFoneOverFzero}{0.903 and 0.885}  % experiments/FIGS/data/v2/index.json :: nonlinear probe jc2_188: (1 - AUC) of FLAV_F1 over FLAV_F0, run by run :: sha256 c6c4138ee5c86423
\newcommand{\MlstMlpCmpJcTwoFoneOverFzeroPr}{at least 9.7\% fewer errors}  % experiments/FIGS/data/v2/index.json :: nonlinear probe jc2_188: (1 - AUC) of FLAV_F1 over FLAV_F0, run by run :: sha256 c6c4138ee5c86423 :: prose twin: two significant figures, direction in words
\newcommand{\MlstMlpCmpJcTwoFoneOverFzeroPrMax}{up to 12\% fewer errors}  % experiments/FIGS/data/v2/index.json :: nonlinear probe jc2_188: (1 - AUC) of FLAV_F1 over FLAV_F0, run by run :: sha256 c6c4138ee5c86423 :: prose twin of the run farther from 1
\newcommand{\MlstMlpCmpJcTwoFonerOverFone}{0.9916 and 0.9918}  % experiments/FIGS/data/v2/index.json :: nonlinear probe jc2_188: (1 - AUC) of FLAV_F1R over FLAV_F1, run by run :: sha256 c6c4138ee5c86423
\newcommand{\MlstMlpCmpJcTwoFonerOverFonePr}{equal within 1\%}  % experiments/FIGS/data/v2/index.json :: nonlinear probe jc2_188: (1 - AUC) of FLAV_F1R over FLAV_F1, run by run :: sha256 c6c4138ee5c86423 :: prose twin: two significant figures, direction in words
\newcommand{\MlstMlpCmpJcTwoFonerOverFonePrMax}{equal within 1\%}  % experiments/FIGS/data/v2/index.json :: nonlinear probe jc2_188: (1 - AUC) of FLAV_F1R over FLAV_F1, run by run :: sha256 c6c4138ee5c86423 :: prose twin of the run farther from 1
\newcommand{\MlstMlpCmpJcTwoFonerOverFzero}{0.895 and 0.877}  % experiments/FIGS/data/v2/index.json :: nonlinear probe jc2_188: (1 - AUC) of FLAV_F1R over FLAV_F0, run by run :: sha256 c6c4138ee5c86423
\newcommand{\MlstMlpCmpJcTwoFonerOverFzeroPr}{at least 10\% fewer errors}  % experiments/FIGS/data/v2/index.json :: nonlinear probe jc2_188: (1 - AUC) of FLAV_F1R over FLAV_F0, run by run :: sha256 c6c4138ee5c86423 :: prose twin: two significant figures, direction in words
\newcommand{\MlstMlpCmpJcTwoFonerOverFzeroPrMax}{up to 12\% fewer errors}  % experiments/FIGS/data/v2/index.json :: nonlinear probe jc2_188: (1 - AUC) of FLAV_F1R over FLAV_F0, run by run :: sha256 c6c4138ee5c86423 :: prose twin of the run farther from 1
\newcommand{\MlstMlpCmpJcTwoFzeroOverCoarse}{0.9461 and 0.9456}  % experiments/FIGS/data/v2/index.json :: nonlinear probe jc2_188: (1 - AUC) of FLAV_F0 over R16_Q1, run by run :: sha256 c6c4138ee5c86423
\newcommand{\MlstMlpCmpJcTwoFzeroOverCoarsePr}{at least 5.4\% fewer errors}  % experiments/FIGS/data/v2/index.json :: nonlinear probe jc2_188: (1 - AUC) of FLAV_F0 over R16_Q1, run by run :: sha256 c6c4138ee5c86423 :: prose twin: two significant figures, direction in words
\newcommand{\MlstMlpCmpJcTwoFzeroOverCoarsePrMax}{up to 5.4\% fewer errors}  % experiments/FIGS/data/v2/index.json :: nonlinear probe jc2_188: (1 - AUC) of FLAV_F0 over R16_Q1, run by run :: sha256 c6c4138ee5c86423 :: prose twin of the run farther from 1
\newcommand{\MlstMlpCmpJcTwoRandOverFull}{\pending{no run yet}}  % experiments/FIGS/data/v2/index.json :: nonlinear probe jc2_188: (1 - AUC) of each random partition over the 188-class run of its seed :: sha256 c6c4138ee5c86423
\newcommand{\MlstMlpCmpJcTwoRandOverFullPr}{\pending{no run yet}}  % experiments/FIGS/data/v2/index.json :: nonlinear probe jc2_188: (1 - AUC) of each random partition over the 188-class run of its seed :: sha256 c6c4138ee5c86423 :: prose twin: two significant figures, direction in words
\newcommand{\MlstMlpCmpJcTwoResResOverFull}{1.001\,\ensuremath{\pm}\,0.007}  % experiments/FIGS/data/v2/index.json :: nonlinear probe jc2_188_res: (1 - AUC) of L162 over L188, run by run :: sha256 c6c4138ee5c86423
\newcommand{\MlstMlpCmpJcTwoResResOverFullPr}{equal within 1\%}  % experiments/FIGS/data/v2/index.json :: nonlinear probe jc2_188_res: (1 - AUC) of L162 over L188, run by run :: sha256 c6c4138ee5c86423 :: prose twin: two significant figures, direction in words
\newcommand{\MlstMlpCmpJcTwoResMidOverFull}{1.192\,\ensuremath{\pm}\,0.008}  % experiments/FIGS/data/v2/index.json :: nonlinear probe jc2_188_res: (1 - AUC) of R42_Q1 over L188, run by run :: sha256 c6c4138ee5c86423
\newcommand{\MlstMlpCmpJcTwoResMidOverFullPr}{19\% more errors}  % experiments/FIGS/data/v2/index.json :: nonlinear probe jc2_188_res: (1 - AUC) of R42_Q1 over L188, run by run :: sha256 c6c4138ee5c86423 :: prose twin: two significant figures, direction in words
\newcommand{\MlstMlpCmpJcTwoResCoarseOverFull}{1.408\,\ensuremath{\pm}\,0.026}  % experiments/FIGS/data/v2/index.json :: nonlinear probe jc2_188_res: (1 - AUC) of R16_Q1 over L188, run by run :: sha256 c6c4138ee5c86423
\newcommand{\MlstMlpCmpJcTwoResCoarseOverFullPr}{41\% more errors}  % experiments/FIGS/data/v2/index.json :: nonlinear probe jc2_188_res: (1 - AUC) of R16_Q1 over L188, run by run :: sha256 c6c4138ee5c86423 :: prose twin: two significant figures, direction in words
\newcommand{\MlstMlpCmpJcTwoResResMassOverRes}{\pending{no run yet}}  % experiments/FIGS/data/v2/index.json :: nonlinear probe jc2_188_res: (1 - AUC) of L162_MASS over L162, run by run :: sha256 c6c4138ee5c86423
\newcommand{\MlstMlpCmpJcTwoResResMassOverResPr}{\pending{no run yet}}  % experiments/FIGS/data/v2/index.json :: nonlinear probe jc2_188_res: (1 - AUC) of L162_MASS over L162, run by run :: sha256 c6c4138ee5c86423 :: prose twin: two significant figures, direction in words
\newcommand{\MlstMlpCmpJcTwoResResMassOverResPrMax}{\pending{no run yet}}  % experiments/FIGS/data/v2/index.json :: nonlinear probe jc2_188_res: (1 - AUC) of L162_MASS over L162, run by run :: sha256 c6c4138ee5c86423 :: prose twin of the run farther from 1
\newcommand{\MlstMlpCmpJcTwoResCoarseMassOverCoarse}{\pending{no run yet}}  % experiments/FIGS/data/v2/index.json :: nonlinear probe jc2_188_res: (1 - AUC) of R16_Q1_MASS over R16_Q1, run by run :: sha256 c6c4138ee5c86423
\newcommand{\MlstMlpCmpJcTwoResCoarseMassOverCoarsePr}{\pending{no run yet}}  % experiments/FIGS/data/v2/index.json :: nonlinear probe jc2_188_res: (1 - AUC) of R16_Q1_MASS over R16_Q1, run by run :: sha256 c6c4138ee5c86423 :: prose twin: two significant figures, direction in words
\newcommand{\MlstMlpCmpJcTwoResCoarseMassOverCoarsePrMax}{\pending{no run yet}}  % experiments/FIGS/data/v2/index.json :: nonlinear probe jc2_188_res: (1 - AUC) of R16_Q1_MASS over R16_Q1, run by run :: sha256 c6c4138ee5c86423 :: prose twin of the run farther from 1
\newcommand{\MlstMlpCmpJcTwoResFzeroOverFull}{1.285 and 1.300}  % experiments/FIGS/data/v2/index.json :: nonlinear probe jc2_188_res: (1 - AUC) of FLAV_F0 over L188, run by run :: sha256 c6c4138ee5c86423
\newcommand{\MlstMlpCmpJcTwoResFzeroOverFullPr}{at least 29\% more errors}  % experiments/FIGS/data/v2/index.json :: nonlinear probe jc2_188_res: (1 - AUC) of FLAV_F0 over L188, run by run :: sha256 c6c4138ee5c86423 :: prose twin: two significant figures, direction in words
\newcommand{\MlstMlpCmpJcTwoResFzeroOverFullPrMax}{up to 30\% more errors}  % experiments/FIGS/data/v2/index.json :: nonlinear probe jc2_188_res: (1 - AUC) of FLAV_F0 over L188, run by run :: sha256 c6c4138ee5c86423 :: prose twin of the run farther from 1
\newcommand{\MlstMlpCmpJcTwoResFoneOverFzero}{0.909 and 0.889}  % experiments/FIGS/data/v2/index.json :: nonlinear probe jc2_188_res: (1 - AUC) of FLAV_F1 over FLAV_F0, run by run :: sha256 c6c4138ee5c86423
\newcommand{\MlstMlpCmpJcTwoResFoneOverFzeroPr}{at least 9.1\% fewer errors}  % experiments/FIGS/data/v2/index.json :: nonlinear probe jc2_188_res: (1 - AUC) of FLAV_F1 over FLAV_F0, run by run :: sha256 c6c4138ee5c86423 :: prose twin: two significant figures, direction in words
\newcommand{\MlstMlpCmpJcTwoResFoneOverFzeroPrMax}{up to 11\% fewer errors}  % experiments/FIGS/data/v2/index.json :: nonlinear probe jc2_188_res: (1 - AUC) of FLAV_F1 over FLAV_F0, run by run :: sha256 c6c4138ee5c86423 :: prose twin of the run farther from 1
\newcommand{\MlstMlpCmpJcTwoResFonerOverFone}{0.9793 and 0.9825}  % experiments/FIGS/data/v2/index.json :: nonlinear probe jc2_188_res: (1 - AUC) of FLAV_F1R over FLAV_F1, run by run :: sha256 c6c4138ee5c86423
\newcommand{\MlstMlpCmpJcTwoResFonerOverFonePr}{equal within 3\%}  % experiments/FIGS/data/v2/index.json :: nonlinear probe jc2_188_res: (1 - AUC) of FLAV_F1R over FLAV_F1, run by run :: sha256 c6c4138ee5c86423 :: prose twin: two significant figures, direction in words
\newcommand{\MlstMlpCmpJcTwoResFonerOverFonePrMax}{equal within 3\%}  % experiments/FIGS/data/v2/index.json :: nonlinear probe jc2_188_res: (1 - AUC) of FLAV_F1R over FLAV_F1, run by run :: sha256 c6c4138ee5c86423 :: prose twin of the run farther from 1
\newcommand{\MlstMlpCmpJcTwoResFonerOverFzero}{0.890 and 0.874}  % experiments/FIGS/data/v2/index.json :: nonlinear probe jc2_188_res: (1 - AUC) of FLAV_F1R over FLAV_F0, run by run :: sha256 c6c4138ee5c86423
\newcommand{\MlstMlpCmpJcTwoResFonerOverFzeroPr}{at least 11\% fewer errors}  % experiments/FIGS/data/v2/index.json :: nonlinear probe jc2_188_res: (1 - AUC) of FLAV_F1R over FLAV_F0, run by run :: sha256 c6c4138ee5c86423 :: prose twin: two significant figures, direction in words
\newcommand{\MlstMlpCmpJcTwoResFonerOverFzeroPrMax}{up to 13\% fewer errors}  % experiments/FIGS/data/v2/index.json :: nonlinear probe jc2_188_res: (1 - AUC) of FLAV_F1R over FLAV_F0, run by run :: sha256 c6c4138ee5c86423 :: prose twin of the run farther from 1
\newcommand{\MlstMlpCmpJcTwoResFzeroOverCoarse}{0.9299 and 0.9233}  % experiments/FIGS/data/v2/index.json :: nonlinear probe jc2_188_res: (1 - AUC) of FLAV_F0 over R16_Q1, run by run :: sha256 c6c4138ee5c86423
\newcommand{\MlstMlpCmpJcTwoResFzeroOverCoarsePr}{at least 7.0\% fewer errors}  % experiments/FIGS/data/v2/index.json :: nonlinear probe jc2_188_res: (1 - AUC) of FLAV_F0 over R16_Q1, run by run :: sha256 c6c4138ee5c86423 :: prose twin: two significant figures, direction in words
\newcommand{\MlstMlpCmpJcTwoResFzeroOverCoarsePrMax}{up to 7.7\% fewer errors}  % experiments/FIGS/data/v2/index.json :: nonlinear probe jc2_188_res: (1 - AUC) of FLAV_F0 over R16_Q1, run by run :: sha256 c6c4138ee5c86423 :: prose twin of the run farther from 1
\newcommand{\MlstMlpCmpJcTwoResRandOverFull}{\pending{no run yet}}  % experiments/FIGS/data/v2/index.json :: nonlinear probe jc2_188_res: (1 - AUC) of each random partition over the 188-class run of its seed :: sha256 c6c4138ee5c86423
\newcommand{\MlstMlpCmpJcTwoResRandOverFullPr}{\pending{no run yet}}  % experiments/FIGS/data/v2/index.json :: nonlinear probe jc2_188_res: (1 - AUC) of each random partition over the 188-class run of its seed :: sha256 c6c4138ee5c86423 :: prose twin: two significant figures, direction in words
\newcommand{\MlstMlpCmpJcTwoQcdResOverFull}{1.070\,\ensuremath{\pm}\,0.006}  % experiments/FIGS/data/v2/index.json :: nonlinear probe jc2_188_qcd: (1 - AUC) of L162 over L188, run by run :: sha256 c6c4138ee5c86423
\newcommand{\MlstMlpCmpJcTwoQcdResOverFullPr}{7.0\% more errors}  % experiments/FIGS/data/v2/index.json :: nonlinear probe jc2_188_qcd: (1 - AUC) of L162 over L188, run by run :: sha256 c6c4138ee5c86423 :: prose twin: two significant figures, direction in words
\newcommand{\MlstMlpCmpJcTwoQcdMidOverFull}{1.107\,\ensuremath{\pm}\,0.002}  % experiments/FIGS/data/v2/index.json :: nonlinear probe jc2_188_qcd: (1 - AUC) of R42_Q1 over L188, run by run :: sha256 c6c4138ee5c86423
\newcommand{\MlstMlpCmpJcTwoQcdMidOverFullPr}{11\% more errors}  % experiments/FIGS/data/v2/index.json :: nonlinear probe jc2_188_qcd: (1 - AUC) of R42_Q1 over L188, run by run :: sha256 c6c4138ee5c86423 :: prose twin: two significant figures, direction in words
\newcommand{\MlstMlpCmpJcTwoQcdCoarseOverFull}{1.333\,\ensuremath{\pm}\,0.012}  % experiments/FIGS/data/v2/index.json :: nonlinear probe jc2_188_qcd: (1 - AUC) of R16_Q1 over L188, run by run :: sha256 c6c4138ee5c86423
\newcommand{\MlstMlpCmpJcTwoQcdCoarseOverFullPr}{33\% more errors}  % experiments/FIGS/data/v2/index.json :: nonlinear probe jc2_188_qcd: (1 - AUC) of R16_Q1 over L188, run by run :: sha256 c6c4138ee5c86423 :: prose twin: two significant figures, direction in words
\newcommand{\MlstMlpCmpJcTwoQcdResMassOverRes}{\pending{no run yet}}  % experiments/FIGS/data/v2/index.json :: nonlinear probe jc2_188_qcd: (1 - AUC) of L162_MASS over L162, run by run :: sha256 c6c4138ee5c86423
\newcommand{\MlstMlpCmpJcTwoQcdResMassOverResPr}{\pending{no run yet}}  % experiments/FIGS/data/v2/index.json :: nonlinear probe jc2_188_qcd: (1 - AUC) of L162_MASS over L162, run by run :: sha256 c6c4138ee5c86423 :: prose twin: two significant figures, direction in words
\newcommand{\MlstMlpCmpJcTwoQcdResMassOverResPrMax}{\pending{no run yet}}  % experiments/FIGS/data/v2/index.json :: nonlinear probe jc2_188_qcd: (1 - AUC) of L162_MASS over L162, run by run :: sha256 c6c4138ee5c86423 :: prose twin of the run farther from 1
\newcommand{\MlstMlpCmpJcTwoQcdCoarseMassOverCoarse}{\pending{no run yet}}  % experiments/FIGS/data/v2/index.json :: nonlinear probe jc2_188_qcd: (1 - AUC) of R16_Q1_MASS over R16_Q1, run by run :: sha256 c6c4138ee5c86423
\newcommand{\MlstMlpCmpJcTwoQcdCoarseMassOverCoarsePr}{\pending{no run yet}}  % experiments/FIGS/data/v2/index.json :: nonlinear probe jc2_188_qcd: (1 - AUC) of R16_Q1_MASS over R16_Q1, run by run :: sha256 c6c4138ee5c86423 :: prose twin: two significant figures, direction in words
\newcommand{\MlstMlpCmpJcTwoQcdCoarseMassOverCoarsePrMax}{\pending{no run yet}}  % experiments/FIGS/data/v2/index.json :: nonlinear probe jc2_188_qcd: (1 - AUC) of R16_Q1_MASS over R16_Q1, run by run :: sha256 c6c4138ee5c86423 :: prose twin of the run farther from 1
\newcommand{\MlstMlpCmpJcTwoQcdFzeroOverFull}{1.319 and 1.339}  % experiments/FIGS/data/v2/index.json :: nonlinear probe jc2_188_qcd: (1 - AUC) of FLAV_F0 over L188, run by run :: sha256 c6c4138ee5c86423
\newcommand{\MlstMlpCmpJcTwoQcdFzeroOverFullPr}{at least 32\% more errors}  % experiments/FIGS/data/v2/index.json :: nonlinear probe jc2_188_qcd: (1 - AUC) of FLAV_F0 over L188, run by run :: sha256 c6c4138ee5c86423 :: prose twin: two significant figures, direction in words
\newcommand{\MlstMlpCmpJcTwoQcdFzeroOverFullPrMax}{up to 34\% more errors}  % experiments/FIGS/data/v2/index.json :: nonlinear probe jc2_188_qcd: (1 - AUC) of FLAV_F0 over L188, run by run :: sha256 c6c4138ee5c86423 :: prose twin of the run farther from 1
\newcommand{\MlstMlpCmpJcTwoQcdFoneOverFzero}{0.886 and 0.872}  % experiments/FIGS/data/v2/index.json :: nonlinear probe jc2_188_qcd: (1 - AUC) of FLAV_F1 over FLAV_F0, run by run :: sha256 c6c4138ee5c86423
\newcommand{\MlstMlpCmpJcTwoQcdFoneOverFzeroPr}{at least 11\% fewer errors}  % experiments/FIGS/data/v2/index.json :: nonlinear probe jc2_188_qcd: (1 - AUC) of FLAV_F1 over FLAV_F0, run by run :: sha256 c6c4138ee5c86423 :: prose twin: two significant figures, direction in words
\newcommand{\MlstMlpCmpJcTwoQcdFoneOverFzeroPrMax}{up to 13\% fewer errors}  % experiments/FIGS/data/v2/index.json :: nonlinear probe jc2_188_qcd: (1 - AUC) of FLAV_F1 over FLAV_F0, run by run :: sha256 c6c4138ee5c86423 :: prose twin of the run farther from 1
\newcommand{\MlstMlpCmpJcTwoQcdFonerOverFone}{1.026 and 1.018}  % experiments/FIGS/data/v2/index.json :: nonlinear probe jc2_188_qcd: (1 - AUC) of FLAV_F1R over FLAV_F1, run by run :: sha256 c6c4138ee5c86423
\newcommand{\MlstMlpCmpJcTwoQcdFonerOverFonePr}{at least 1.8\% more errors}  % experiments/FIGS/data/v2/index.json :: nonlinear probe jc2_188_qcd: (1 - AUC) of FLAV_F1R over FLAV_F1, run by run :: sha256 c6c4138ee5c86423 :: prose twin: two significant figures, direction in words
\newcommand{\MlstMlpCmpJcTwoQcdFonerOverFonePrMax}{up to 2.6\% more errors}  % experiments/FIGS/data/v2/index.json :: nonlinear probe jc2_188_qcd: (1 - AUC) of FLAV_F1R over FLAV_F1, run by run :: sha256 c6c4138ee5c86423 :: prose twin of the run farther from 1
\newcommand{\MlstMlpCmpJcTwoQcdFonerOverFzero}{0.909 and 0.887}  % experiments/FIGS/data/v2/index.json :: nonlinear probe jc2_188_qcd: (1 - AUC) of FLAV_F1R over FLAV_F0, run by run :: sha256 c6c4138ee5c86423
\newcommand{\MlstMlpCmpJcTwoQcdFonerOverFzeroPr}{at least 9.1\% fewer errors}  % experiments/FIGS/data/v2/index.json :: nonlinear probe jc2_188_qcd: (1 - AUC) of FLAV_F1R over FLAV_F0, run by run :: sha256 c6c4138ee5c86423 :: prose twin: two significant figures, direction in words
\newcommand{\MlstMlpCmpJcTwoQcdFonerOverFzeroPrMax}{up to 11\% fewer errors}  % experiments/FIGS/data/v2/index.json :: nonlinear probe jc2_188_qcd: (1 - AUC) of FLAV_F1R over FLAV_F0, run by run :: sha256 c6c4138ee5c86423 :: prose twin of the run farther from 1
\newcommand{\MlstMlpCmpJcTwoQcdFzeroOverCoarse}{0.993 and 1.012}  % experiments/FIGS/data/v2/index.json :: nonlinear probe jc2_188_qcd: (1 - AUC) of FLAV_F0 over R16_Q1, run by run :: sha256 c6c4138ee5c86423
\newcommand{\MlstMlpCmpJcTwoQcdFzeroOverCoarsePr}{equal within 2\%}  % experiments/FIGS/data/v2/index.json :: nonlinear probe jc2_188_qcd: (1 - AUC) of FLAV_F0 over R16_Q1, run by run :: sha256 c6c4138ee5c86423 :: prose twin: two significant figures, direction in words
\newcommand{\MlstMlpCmpJcTwoQcdFzeroOverCoarsePrMax}{equal within 2\%}  % experiments/FIGS/data/v2/index.json :: nonlinear probe jc2_188_qcd: (1 - AUC) of FLAV_F0 over R16_Q1, run by run :: sha256 c6c4138ee5c86423 :: prose twin of the run farther from 1
\newcommand{\MlstMlpCmpJcTwoQcdRandOverFull}{\pending{no run yet}}  % experiments/FIGS/data/v2/index.json :: nonlinear probe jc2_188_qcd: (1 - AUC) of each random partition over the 188-class run of its seed :: sha256 c6c4138ee5c86423
\newcommand{\MlstMlpCmpJcTwoQcdRandOverFullPr}{\pending{no run yet}}  % experiments/FIGS/data/v2/index.json :: nonlinear probe jc2_188_qcd: (1 - AUC) of each random partition over the 188-class run of its seed :: sha256 c6c4138ee5c86423 :: prose twin: two significant figures, direction in words
\newcommand{\MlstMlpCmpBcResOverFull}{1.002\,\ensuremath{\pm}\,0.006}  % experiments/FIGS/data/v2/index.json :: nonlinear probe bc_vs_bq_cs: (1 - AUC) of L162 over L188, run by run :: sha256 c6c4138ee5c86423
\newcommand{\MlstMlpCmpBcResOverFullPr}{equal within 1\%}  % experiments/FIGS/data/v2/index.json :: nonlinear probe bc_vs_bq_cs: (1 - AUC) of L162 over L188, run by run :: sha256 c6c4138ee5c86423 :: prose twin: two significant figures, direction in words
\newcommand{\MlstMlpCmpBcMidOverFull}{1.104\,\ensuremath{\pm}\,0.017}  % experiments/FIGS/data/v2/index.json :: nonlinear probe bc_vs_bq_cs: (1 - AUC) of R42_Q1 over L188, run by run :: sha256 c6c4138ee5c86423
\newcommand{\MlstMlpCmpBcMidOverFullPr}{10\% more errors}  % experiments/FIGS/data/v2/index.json :: nonlinear probe bc_vs_bq_cs: (1 - AUC) of R42_Q1 over L188, run by run :: sha256 c6c4138ee5c86423 :: prose twin: two significant figures, direction in words
\newcommand{\MlstMlpCmpBcCoarseOverFull}{1.624\,\ensuremath{\pm}\,0.022}  % experiments/FIGS/data/v2/index.json :: nonlinear probe bc_vs_bq_cs: (1 - AUC) of R16_Q1 over L188, run by run :: sha256 c6c4138ee5c86423
\newcommand{\MlstMlpCmpBcCoarseOverFullPr}{62\% more errors}  % experiments/FIGS/data/v2/index.json :: nonlinear probe bc_vs_bq_cs: (1 - AUC) of R16_Q1 over L188, run by run :: sha256 c6c4138ee5c86423 :: prose twin: two significant figures, direction in words
\newcommand{\MlstMlpCmpBcResMassOverRes}{\pending{no run yet}}  % experiments/FIGS/data/v2/index.json :: nonlinear probe bc_vs_bq_cs: (1 - AUC) of L162_MASS over L162, run by run :: sha256 c6c4138ee5c86423
\newcommand{\MlstMlpCmpBcResMassOverResPr}{\pending{no run yet}}  % experiments/FIGS/data/v2/index.json :: nonlinear probe bc_vs_bq_cs: (1 - AUC) of L162_MASS over L162, run by run :: sha256 c6c4138ee5c86423 :: prose twin: two significant figures, direction in words
\newcommand{\MlstMlpCmpBcResMassOverResPrMax}{\pending{no run yet}}  % experiments/FIGS/data/v2/index.json :: nonlinear probe bc_vs_bq_cs: (1 - AUC) of L162_MASS over L162, run by run :: sha256 c6c4138ee5c86423 :: prose twin of the run farther from 1
\newcommand{\MlstMlpCmpBcCoarseMassOverCoarse}{\pending{no run yet}}  % experiments/FIGS/data/v2/index.json :: nonlinear probe bc_vs_bq_cs: (1 - AUC) of R16_Q1_MASS over R16_Q1, run by run :: sha256 c6c4138ee5c86423
\newcommand{\MlstMlpCmpBcCoarseMassOverCoarsePr}{\pending{no run yet}}  % experiments/FIGS/data/v2/index.json :: nonlinear probe bc_vs_bq_cs: (1 - AUC) of R16_Q1_MASS over R16_Q1, run by run :: sha256 c6c4138ee5c86423 :: prose twin: two significant figures, direction in words
\newcommand{\MlstMlpCmpBcCoarseMassOverCoarsePrMax}{\pending{no run yet}}  % experiments/FIGS/data/v2/index.json :: nonlinear probe bc_vs_bq_cs: (1 - AUC) of R16_Q1_MASS over R16_Q1, run by run :: sha256 c6c4138ee5c86423 :: prose twin of the run farther from 1
\newcommand{\MlstMlpCmpBcFzeroOverFull}{1.633 and 1.679}  % experiments/FIGS/data/v2/index.json :: nonlinear probe bc_vs_bq_cs: (1 - AUC) of FLAV_F0 over L188, run by run :: sha256 c6c4138ee5c86423
\newcommand{\MlstMlpCmpBcFzeroOverFullPr}{at least 63\% more errors}  % experiments/FIGS/data/v2/index.json :: nonlinear probe bc_vs_bq_cs: (1 - AUC) of FLAV_F0 over L188, run by run :: sha256 c6c4138ee5c86423 :: prose twin: two significant figures, direction in words
\newcommand{\MlstMlpCmpBcFzeroOverFullPrMax}{up to 68\% more errors}  % experiments/FIGS/data/v2/index.json :: nonlinear probe bc_vs_bq_cs: (1 - AUC) of FLAV_F0 over L188, run by run :: sha256 c6c4138ee5c86423 :: prose twin of the run farther from 1
\newcommand{\MlstMlpCmpBcFoneOverFzero}{0.848 and 0.810}  % experiments/FIGS/data/v2/index.json :: nonlinear probe bc_vs_bq_cs: (1 - AUC) of FLAV_F1 over FLAV_F0, run by run :: sha256 c6c4138ee5c86423
\newcommand{\MlstMlpCmpBcFoneOverFzeroPr}{at least 15\% fewer errors}  % experiments/FIGS/data/v2/index.json :: nonlinear probe bc_vs_bq_cs: (1 - AUC) of FLAV_F1 over FLAV_F0, run by run :: sha256 c6c4138ee5c86423 :: prose twin: two significant figures, direction in words
\newcommand{\MlstMlpCmpBcFoneOverFzeroPrMax}{up to 19\% fewer errors}  % experiments/FIGS/data/v2/index.json :: nonlinear probe bc_vs_bq_cs: (1 - AUC) of FLAV_F1 over FLAV_F0, run by run :: sha256 c6c4138ee5c86423 :: prose twin of the run farther from 1
\newcommand{\MlstMlpCmpBcFonerOverFone}{0.9812 and 0.9815}  % experiments/FIGS/data/v2/index.json :: nonlinear probe bc_vs_bq_cs: (1 - AUC) of FLAV_F1R over FLAV_F1, run by run :: sha256 c6c4138ee5c86423
\newcommand{\MlstMlpCmpBcFonerOverFonePr}{equal within 2\%}  % experiments/FIGS/data/v2/index.json :: nonlinear probe bc_vs_bq_cs: (1 - AUC) of FLAV_F1R over FLAV_F1, run by run :: sha256 c6c4138ee5c86423 :: prose twin: two significant figures, direction in words
\newcommand{\MlstMlpCmpBcFonerOverFonePrMax}{equal within 2\%}  % experiments/FIGS/data/v2/index.json :: nonlinear probe bc_vs_bq_cs: (1 - AUC) of FLAV_F1R over FLAV_F1, run by run :: sha256 c6c4138ee5c86423 :: prose twin of the run farther from 1
\newcommand{\MlstMlpCmpBcFonerOverFzero}{0.833 and 0.795}  % experiments/FIGS/data/v2/index.json :: nonlinear probe bc_vs_bq_cs: (1 - AUC) of FLAV_F1R over FLAV_F0, run by run :: sha256 c6c4138ee5c86423
\newcommand{\MlstMlpCmpBcFonerOverFzeroPr}{at least 17\% fewer errors}  % experiments/FIGS/data/v2/index.json :: nonlinear probe bc_vs_bq_cs: (1 - AUC) of FLAV_F1R over FLAV_F0, run by run :: sha256 c6c4138ee5c86423 :: prose twin: two significant figures, direction in words
\newcommand{\MlstMlpCmpBcFonerOverFzeroPrMax}{up to 21\% fewer errors}  % experiments/FIGS/data/v2/index.json :: nonlinear probe bc_vs_bq_cs: (1 - AUC) of FLAV_F1R over FLAV_F0, run by run :: sha256 c6c4138ee5c86423 :: prose twin of the run farther from 1
\newcommand{\MlstMlpCmpBcFzeroOverCoarse}{0.994 and 1.030}  % experiments/FIGS/data/v2/index.json :: nonlinear probe bc_vs_bq_cs: (1 - AUC) of FLAV_F0 over R16_Q1, run by run :: sha256 c6c4138ee5c86423
\newcommand{\MlstMlpCmpBcFzeroOverCoarsePr}{equal within 3\%}  % experiments/FIGS/data/v2/index.json :: nonlinear probe bc_vs_bq_cs: (1 - AUC) of FLAV_F0 over R16_Q1, run by run :: sha256 c6c4138ee5c86423 :: prose twin: two significant figures, direction in words
\newcommand{\MlstMlpCmpBcFzeroOverCoarsePrMax}{equal within 3\%}  % experiments/FIGS/data/v2/index.json :: nonlinear probe bc_vs_bq_cs: (1 - AUC) of FLAV_F0 over R16_Q1, run by run :: sha256 c6c4138ee5c86423 :: prose twin of the run farther from 1
\newcommand{\MlstMlpCmpBcRandOverFull}{\pending{no run yet}}  % experiments/FIGS/data/v2/index.json :: nonlinear probe bc_vs_bq_cs: (1 - AUC) of each random partition over the 188-class run of its seed :: sha256 c6c4138ee5c86423
\newcommand{\MlstMlpCmpBcRandOverFullPr}{\pending{no run yet}}  % experiments/FIGS/data/v2/index.json :: nonlinear probe bc_vs_bq_cs: (1 - AUC) of each random partition over the 188-class run of its seed :: sha256 c6c4138ee5c86423 :: prose twin: two significant figures, direction in words
\newcommand{\MlstMlpCmpBbResOverFull}{1.023\,\ensuremath{\pm}\,0.021}  % experiments/FIGS/data/v2/index.json :: nonlinear probe bb_vs_cc: (1 - AUC) of L162 over L188, run by run :: sha256 c6c4138ee5c86423
\newcommand{\MlstMlpCmpBbResOverFullPr}{equal within 3\%}  % experiments/FIGS/data/v2/index.json :: nonlinear probe bb_vs_cc: (1 - AUC) of L162 over L188, run by run :: sha256 c6c4138ee5c86423 :: prose twin: two significant figures, direction in words
\newcommand{\MlstMlpCmpBbMidOverFull}{1.125\,\ensuremath{\pm}\,0.011}  % experiments/FIGS/data/v2/index.json :: nonlinear probe bb_vs_cc: (1 - AUC) of R42_Q1 over L188, run by run :: sha256 c6c4138ee5c86423
\newcommand{\MlstMlpCmpBbMidOverFullPr}{13\% more errors}  % experiments/FIGS/data/v2/index.json :: nonlinear probe bb_vs_cc: (1 - AUC) of R42_Q1 over L188, run by run :: sha256 c6c4138ee5c86423 :: prose twin: two significant figures, direction in words
\newcommand{\MlstMlpCmpBbCoarseOverFull}{3.04\,\ensuremath{\pm}\,0.09}  % experiments/FIGS/data/v2/index.json :: nonlinear probe bb_vs_cc: (1 - AUC) of R16_Q1 over L188, run by run :: sha256 c6c4138ee5c86423
\newcommand{\MlstMlpCmpBbCoarseOverFullPr}{3.0 times the errors}  % experiments/FIGS/data/v2/index.json :: nonlinear probe bb_vs_cc: (1 - AUC) of R16_Q1 over L188, run by run :: sha256 c6c4138ee5c86423 :: prose twin: two significant figures, direction in words
\newcommand{\MlstMlpCmpBbResMassOverRes}{\pending{no run yet}}  % experiments/FIGS/data/v2/index.json :: nonlinear probe bb_vs_cc: (1 - AUC) of L162_MASS over L162, run by run :: sha256 c6c4138ee5c86423
\newcommand{\MlstMlpCmpBbResMassOverResPr}{\pending{no run yet}}  % experiments/FIGS/data/v2/index.json :: nonlinear probe bb_vs_cc: (1 - AUC) of L162_MASS over L162, run by run :: sha256 c6c4138ee5c86423 :: prose twin: two significant figures, direction in words
\newcommand{\MlstMlpCmpBbResMassOverResPrMax}{\pending{no run yet}}  % experiments/FIGS/data/v2/index.json :: nonlinear probe bb_vs_cc: (1 - AUC) of L162_MASS over L162, run by run :: sha256 c6c4138ee5c86423 :: prose twin of the run farther from 1
\newcommand{\MlstMlpCmpBbCoarseMassOverCoarse}{\pending{no run yet}}  % experiments/FIGS/data/v2/index.json :: nonlinear probe bb_vs_cc: (1 - AUC) of R16_Q1_MASS over R16_Q1, run by run :: sha256 c6c4138ee5c86423
\newcommand{\MlstMlpCmpBbCoarseMassOverCoarsePr}{\pending{no run yet}}  % experiments/FIGS/data/v2/index.json :: nonlinear probe bb_vs_cc: (1 - AUC) of R16_Q1_MASS over R16_Q1, run by run :: sha256 c6c4138ee5c86423 :: prose twin: two significant figures, direction in words
\newcommand{\MlstMlpCmpBbCoarseMassOverCoarsePrMax}{\pending{no run yet}}  % experiments/FIGS/data/v2/index.json :: nonlinear probe bb_vs_cc: (1 - AUC) of R16_Q1_MASS over R16_Q1, run by run :: sha256 c6c4138ee5c86423 :: prose twin of the run farther from 1
\newcommand{\MlstMlpCmpBbFzeroOverFull}{3.3 and 4.3}  % experiments/FIGS/data/v2/index.json :: nonlinear probe bb_vs_cc: (1 - AUC) of FLAV_F0 over L188, run by run :: sha256 c6c4138ee5c86423
\newcommand{\MlstMlpCmpBbFzeroOverFullPr}{at least 3.3 times the errors}  % experiments/FIGS/data/v2/index.json :: nonlinear probe bb_vs_cc: (1 - AUC) of FLAV_F0 over L188, run by run :: sha256 c6c4138ee5c86423 :: prose twin: two significant figures, direction in words
\newcommand{\MlstMlpCmpBbFzeroOverFullPrMax}{up to 4.3 times the errors}  % experiments/FIGS/data/v2/index.json :: nonlinear probe bb_vs_cc: (1 - AUC) of FLAV_F0 over L188, run by run :: sha256 c6c4138ee5c86423 :: prose twin of the run farther from 1
\newcommand{\MlstMlpCmpBbFoneOverFzero}{0.52 and 0.38}  % experiments/FIGS/data/v2/index.json :: nonlinear probe bb_vs_cc: (1 - AUC) of FLAV_F1 over FLAV_F0, run by run :: sha256 c6c4138ee5c86423
\newcommand{\MlstMlpCmpBbFoneOverFzeroPr}{at least 48\% fewer errors}  % experiments/FIGS/data/v2/index.json :: nonlinear probe bb_vs_cc: (1 - AUC) of FLAV_F1 over FLAV_F0, run by run :: sha256 c6c4138ee5c86423 :: prose twin: two significant figures, direction in words
\newcommand{\MlstMlpCmpBbFoneOverFzeroPrMax}{down to 0.38 times the errors}  % experiments/FIGS/data/v2/index.json :: nonlinear probe bb_vs_cc: (1 - AUC) of FLAV_F1 over FLAV_F0, run by run :: sha256 c6c4138ee5c86423 :: prose twin of the run farther from 1
\newcommand{\MlstMlpCmpBbFonerOverFone}{1.27 and 1.45}  % experiments/FIGS/data/v2/index.json :: nonlinear probe bb_vs_cc: (1 - AUC) of FLAV_F1R over FLAV_F1, run by run :: sha256 c6c4138ee5c86423
\newcommand{\MlstMlpCmpBbFonerOverFonePr}{at least 27\% more errors}  % experiments/FIGS/data/v2/index.json :: nonlinear probe bb_vs_cc: (1 - AUC) of FLAV_F1R over FLAV_F1, run by run :: sha256 c6c4138ee5c86423 :: prose twin: two significant figures, direction in words
\newcommand{\MlstMlpCmpBbFonerOverFonePrMax}{up to 45\% more errors}  % experiments/FIGS/data/v2/index.json :: nonlinear probe bb_vs_cc: (1 - AUC) of FLAV_F1R over FLAV_F1, run by run :: sha256 c6c4138ee5c86423 :: prose twin of the run farther from 1
\newcommand{\MlstMlpCmpBbFonerOverFzero}{0.66 and 0.55}  % experiments/FIGS/data/v2/index.json :: nonlinear probe bb_vs_cc: (1 - AUC) of FLAV_F1R over FLAV_F0, run by run :: sha256 c6c4138ee5c86423
\newcommand{\MlstMlpCmpBbFonerOverFzeroPr}{at least 34\% fewer errors}  % experiments/FIGS/data/v2/index.json :: nonlinear probe bb_vs_cc: (1 - AUC) of FLAV_F1R over FLAV_F0, run by run :: sha256 c6c4138ee5c86423 :: prose twin: two significant figures, direction in words
\newcommand{\MlstMlpCmpBbFonerOverFzeroPrMax}{up to 45\% fewer errors}  % experiments/FIGS/data/v2/index.json :: nonlinear probe bb_vs_cc: (1 - AUC) of FLAV_F1R over FLAV_F0, run by run :: sha256 c6c4138ee5c86423 :: prose twin of the run farther from 1
\newcommand{\MlstMlpCmpBbFzeroOverCoarse}{1.07 and 1.41}  % experiments/FIGS/data/v2/index.json :: nonlinear probe bb_vs_cc: (1 - AUC) of FLAV_F0 over R16_Q1, run by run :: sha256 c6c4138ee5c86423
\newcommand{\MlstMlpCmpBbFzeroOverCoarsePr}{at least 7.4\% more errors}  % experiments/FIGS/data/v2/index.json :: nonlinear probe bb_vs_cc: (1 - AUC) of FLAV_F0 over R16_Q1, run by run :: sha256 c6c4138ee5c86423 :: prose twin: two significant figures, direction in words
\newcommand{\MlstMlpCmpBbFzeroOverCoarsePrMax}{up to 41\% more errors}  % experiments/FIGS/data/v2/index.json :: nonlinear probe bb_vs_cc: (1 - AUC) of FLAV_F0 over R16_Q1, run by run :: sha256 c6c4138ee5c86423 :: prose twin of the run farther from 1
\newcommand{\MlstMlpCmpBbRandOverFull}{\pending{no run yet}}  % experiments/FIGS/data/v2/index.json :: nonlinear probe bb_vs_cc: (1 - AUC) of each random partition over the 188-class run of its seed :: sha256 c6c4138ee5c86423
\newcommand{\MlstMlpCmpBbRandOverFullPr}{\pending{no run yet}}  % experiments/FIGS/data/v2/index.json :: nonlinear probe bb_vs_cc: (1 - AUC) of each random partition over the 188-class run of its seed :: sha256 c6c4138ee5c86423 :: prose twin: two significant figures, direction in words
\newcommand{\MlstMlpCmpBbqqResOverFull}{1.011\,\ensuremath{\pm}\,0.030}  % experiments/FIGS/data/v2/index.json :: nonlinear probe bbqq_vs_ccqq: (1 - AUC) of L162 over L188, run by run :: sha256 c6c4138ee5c86423
\newcommand{\MlstMlpCmpBbqqResOverFullPr}{equal within 3\%}  % experiments/FIGS/data/v2/index.json :: nonlinear probe bbqq_vs_ccqq: (1 - AUC) of L162 over L188, run by run :: sha256 c6c4138ee5c86423 :: prose twin: two significant figures, direction in words
\newcommand{\MlstMlpCmpBbqqMidOverFull}{1.112\,\ensuremath{\pm}\,0.006}  % experiments/FIGS/data/v2/index.json :: nonlinear probe bbqq_vs_ccqq: (1 - AUC) of R42_Q1 over L188, run by run :: sha256 c6c4138ee5c86423
\newcommand{\MlstMlpCmpBbqqMidOverFullPr}{11\% more errors}  % experiments/FIGS/data/v2/index.json :: nonlinear probe bbqq_vs_ccqq: (1 - AUC) of R42_Q1 over L188, run by run :: sha256 c6c4138ee5c86423 :: prose twin: two significant figures, direction in words
\newcommand{\MlstMlpCmpBbqqCoarseOverFull}{2.93\,\ensuremath{\pm}\,0.17}  % experiments/FIGS/data/v2/index.json :: nonlinear probe bbqq_vs_ccqq: (1 - AUC) of R16_Q1 over L188, run by run :: sha256 c6c4138ee5c86423
\newcommand{\MlstMlpCmpBbqqCoarseOverFullPr}{2.9 times the errors}  % experiments/FIGS/data/v2/index.json :: nonlinear probe bbqq_vs_ccqq: (1 - AUC) of R16_Q1 over L188, run by run :: sha256 c6c4138ee5c86423 :: prose twin: two significant figures, direction in words
\newcommand{\MlstMlpCmpBbqqResMassOverRes}{\pending{no run yet}}  % experiments/FIGS/data/v2/index.json :: nonlinear probe bbqq_vs_ccqq: (1 - AUC) of L162_MASS over L162, run by run :: sha256 c6c4138ee5c86423
\newcommand{\MlstMlpCmpBbqqResMassOverResPr}{\pending{no run yet}}  % experiments/FIGS/data/v2/index.json :: nonlinear probe bbqq_vs_ccqq: (1 - AUC) of L162_MASS over L162, run by run :: sha256 c6c4138ee5c86423 :: prose twin: two significant figures, direction in words
\newcommand{\MlstMlpCmpBbqqResMassOverResPrMax}{\pending{no run yet}}  % experiments/FIGS/data/v2/index.json :: nonlinear probe bbqq_vs_ccqq: (1 - AUC) of L162_MASS over L162, run by run :: sha256 c6c4138ee5c86423 :: prose twin of the run farther from 1
\newcommand{\MlstMlpCmpBbqqCoarseMassOverCoarse}{\pending{no run yet}}  % experiments/FIGS/data/v2/index.json :: nonlinear probe bbqq_vs_ccqq: (1 - AUC) of R16_Q1_MASS over R16_Q1, run by run :: sha256 c6c4138ee5c86423
\newcommand{\MlstMlpCmpBbqqCoarseMassOverCoarsePr}{\pending{no run yet}}  % experiments/FIGS/data/v2/index.json :: nonlinear probe bbqq_vs_ccqq: (1 - AUC) of R16_Q1_MASS over R16_Q1, run by run :: sha256 c6c4138ee5c86423 :: prose twin: two significant figures, direction in words
\newcommand{\MlstMlpCmpBbqqCoarseMassOverCoarsePrMax}{\pending{no run yet}}  % experiments/FIGS/data/v2/index.json :: nonlinear probe bbqq_vs_ccqq: (1 - AUC) of R16_Q1_MASS over R16_Q1, run by run :: sha256 c6c4138ee5c86423 :: prose twin of the run farther from 1
\newcommand{\MlstMlpCmpBbqqFzeroOverFull}{3.20 and 3.78}  % experiments/FIGS/data/v2/index.json :: nonlinear probe bbqq_vs_ccqq: (1 - AUC) of FLAV_F0 over L188, run by run :: sha256 c6c4138ee5c86423
\newcommand{\MlstMlpCmpBbqqFzeroOverFullPr}{at least 3.2 times the errors}  % experiments/FIGS/data/v2/index.json :: nonlinear probe bbqq_vs_ccqq: (1 - AUC) of FLAV_F0 over L188, run by run :: sha256 c6c4138ee5c86423 :: prose twin: two significant figures, direction in words
\newcommand{\MlstMlpCmpBbqqFzeroOverFullPrMax}{up to 3.8 times the errors}  % experiments/FIGS/data/v2/index.json :: nonlinear probe bbqq_vs_ccqq: (1 - AUC) of FLAV_F0 over L188, run by run :: sha256 c6c4138ee5c86423 :: prose twin of the run farther from 1
\newcommand{\MlstMlpCmpBbqqFoneOverFzero}{0.450 and 0.363}  % experiments/FIGS/data/v2/index.json :: nonlinear probe bbqq_vs_ccqq: (1 - AUC) of FLAV_F1 over FLAV_F0, run by run :: sha256 c6c4138ee5c86423
\newcommand{\MlstMlpCmpBbqqFoneOverFzeroPr}{at most 0.45 times the errors}  % experiments/FIGS/data/v2/index.json :: nonlinear probe bbqq_vs_ccqq: (1 - AUC) of FLAV_F1 over FLAV_F0, run by run :: sha256 c6c4138ee5c86423 :: prose twin: two significant figures, direction in words
\newcommand{\MlstMlpCmpBbqqFoneOverFzeroPrMax}{down to 0.36 times the errors}  % experiments/FIGS/data/v2/index.json :: nonlinear probe bbqq_vs_ccqq: (1 - AUC) of FLAV_F1 over FLAV_F0, run by run :: sha256 c6c4138ee5c86423 :: prose twin of the run farther from 1
\newcommand{\MlstMlpCmpBbqqFonerOverFone}{1.311 and 1.343}  % experiments/FIGS/data/v2/index.json :: nonlinear probe bbqq_vs_ccqq: (1 - AUC) of FLAV_F1R over FLAV_F1, run by run :: sha256 c6c4138ee5c86423
\newcommand{\MlstMlpCmpBbqqFonerOverFonePr}{at least 31\% more errors}  % experiments/FIGS/data/v2/index.json :: nonlinear probe bbqq_vs_ccqq: (1 - AUC) of FLAV_F1R over FLAV_F1, run by run :: sha256 c6c4138ee5c86423 :: prose twin: two significant figures, direction in words
\newcommand{\MlstMlpCmpBbqqFonerOverFonePrMax}{up to 34\% more errors}  % experiments/FIGS/data/v2/index.json :: nonlinear probe bbqq_vs_ccqq: (1 - AUC) of FLAV_F1R over FLAV_F1, run by run :: sha256 c6c4138ee5c86423 :: prose twin of the run farther from 1
\newcommand{\MlstMlpCmpBbqqFonerOverFzero}{0.59 and 0.49}  % experiments/FIGS/data/v2/index.json :: nonlinear probe bbqq_vs_ccqq: (1 - AUC) of FLAV_F1R over FLAV_F0, run by run :: sha256 c6c4138ee5c86423
\newcommand{\MlstMlpCmpBbqqFonerOverFzeroPr}{at least 41\% fewer errors}  % experiments/FIGS/data/v2/index.json :: nonlinear probe bbqq_vs_ccqq: (1 - AUC) of FLAV_F1R over FLAV_F0, run by run :: sha256 c6c4138ee5c86423 :: prose twin: two significant figures, direction in words
\newcommand{\MlstMlpCmpBbqqFonerOverFzeroPrMax}{down to 0.49 times the errors}  % experiments/FIGS/data/v2/index.json :: nonlinear probe bbqq_vs_ccqq: (1 - AUC) of FLAV_F1R over FLAV_F0, run by run :: sha256 c6c4138ee5c86423 :: prose twin of the run farther from 1
\newcommand{\MlstMlpCmpBbqqFzeroOverCoarse}{1.170 and 1.236}  % experiments/FIGS/data/v2/index.json :: nonlinear probe bbqq_vs_ccqq: (1 - AUC) of FLAV_F0 over R16_Q1, run by run :: sha256 c6c4138ee5c86423
\newcommand{\MlstMlpCmpBbqqFzeroOverCoarsePr}{at least 17\% more errors}  % experiments/FIGS/data/v2/index.json :: nonlinear probe bbqq_vs_ccqq: (1 - AUC) of FLAV_F0 over R16_Q1, run by run :: sha256 c6c4138ee5c86423 :: prose twin: two significant figures, direction in words
\newcommand{\MlstMlpCmpBbqqFzeroOverCoarsePrMax}{up to 24\% more errors}  % experiments/FIGS/data/v2/index.json :: nonlinear probe bbqq_vs_ccqq: (1 - AUC) of FLAV_F0 over R16_Q1, run by run :: sha256 c6c4138ee5c86423 :: prose twin of the run farther from 1
\newcommand{\MlstMlpCmpBbqqRandOverFull}{\pending{no run yet}}  % experiments/FIGS/data/v2/index.json :: nonlinear probe bbqq_vs_ccqq: (1 - AUC) of each random partition over the 188-class run of its seed :: sha256 c6c4138ee5c86423
\newcommand{\MlstMlpCmpBbqqRandOverFullPr}{\pending{no run yet}}  % experiments/FIGS/data/v2/index.json :: nonlinear probe bbqq_vs_ccqq: (1 - AUC) of each random partition over the 188-class run of its seed :: sha256 c6c4138ee5c86423 :: prose twin: two significant figures, direction in words
\newcommand{\MlstMlpCmpQcdBbResOverFull}{1.085\,\ensuremath{\pm}\,0.030}  % experiments/FIGS/data/v2/index.json :: nonlinear probe qcd_bb_vs_cc: (1 - AUC) of L162 over L188, run by run :: sha256 c6c4138ee5c86423
\newcommand{\MlstMlpCmpQcdBbResOverFullPr}{8.5\% more errors}  % experiments/FIGS/data/v2/index.json :: nonlinear probe qcd_bb_vs_cc: (1 - AUC) of L162 over L188, run by run :: sha256 c6c4138ee5c86423 :: prose twin: two significant figures, direction in words
\newcommand{\MlstMlpCmpQcdBbMidOverFull}{1.121\,\ensuremath{\pm}\,0.017}  % experiments/FIGS/data/v2/index.json :: nonlinear probe qcd_bb_vs_cc: (1 - AUC) of R42_Q1 over L188, run by run :: sha256 c6c4138ee5c86423
\newcommand{\MlstMlpCmpQcdBbMidOverFullPr}{12\% more errors}  % experiments/FIGS/data/v2/index.json :: nonlinear probe qcd_bb_vs_cc: (1 - AUC) of R42_Q1 over L188, run by run :: sha256 c6c4138ee5c86423 :: prose twin: two significant figures, direction in words
\newcommand{\MlstMlpCmpQcdBbCoarseOverFull}{2.31\,\ensuremath{\pm}\,0.04}  % experiments/FIGS/data/v2/index.json :: nonlinear probe qcd_bb_vs_cc: (1 - AUC) of R16_Q1 over L188, run by run :: sha256 c6c4138ee5c86423
\newcommand{\MlstMlpCmpQcdBbCoarseOverFullPr}{2.3 times the errors}  % experiments/FIGS/data/v2/index.json :: nonlinear probe qcd_bb_vs_cc: (1 - AUC) of R16_Q1 over L188, run by run :: sha256 c6c4138ee5c86423 :: prose twin: two significant figures, direction in words
\newcommand{\MlstMlpCmpQcdBbResMassOverRes}{\pending{no run yet}}  % experiments/FIGS/data/v2/index.json :: nonlinear probe qcd_bb_vs_cc: (1 - AUC) of L162_MASS over L162, run by run :: sha256 c6c4138ee5c86423
\newcommand{\MlstMlpCmpQcdBbResMassOverResPr}{\pending{no run yet}}  % experiments/FIGS/data/v2/index.json :: nonlinear probe qcd_bb_vs_cc: (1 - AUC) of L162_MASS over L162, run by run :: sha256 c6c4138ee5c86423 :: prose twin: two significant figures, direction in words
\newcommand{\MlstMlpCmpQcdBbResMassOverResPrMax}{\pending{no run yet}}  % experiments/FIGS/data/v2/index.json :: nonlinear probe qcd_bb_vs_cc: (1 - AUC) of L162_MASS over L162, run by run :: sha256 c6c4138ee5c86423 :: prose twin of the run farther from 1
\newcommand{\MlstMlpCmpQcdBbCoarseMassOverCoarse}{\pending{no run yet}}  % experiments/FIGS/data/v2/index.json :: nonlinear probe qcd_bb_vs_cc: (1 - AUC) of R16_Q1_MASS over R16_Q1, run by run :: sha256 c6c4138ee5c86423
\newcommand{\MlstMlpCmpQcdBbCoarseMassOverCoarsePr}{\pending{no run yet}}  % experiments/FIGS/data/v2/index.json :: nonlinear probe qcd_bb_vs_cc: (1 - AUC) of R16_Q1_MASS over R16_Q1, run by run :: sha256 c6c4138ee5c86423 :: prose twin: two significant figures, direction in words
\newcommand{\MlstMlpCmpQcdBbCoarseMassOverCoarsePrMax}{\pending{no run yet}}  % experiments/FIGS/data/v2/index.json :: nonlinear probe qcd_bb_vs_cc: (1 - AUC) of R16_Q1_MASS over R16_Q1, run by run :: sha256 c6c4138ee5c86423 :: prose twin of the run farther from 1
\newcommand{\MlstMlpCmpQcdBbFzeroOverFull}{2.50 and 2.72}  % experiments/FIGS/data/v2/index.json :: nonlinear probe qcd_bb_vs_cc: (1 - AUC) of FLAV_F0 over L188, run by run :: sha256 c6c4138ee5c86423
\newcommand{\MlstMlpCmpQcdBbFzeroOverFullPr}{at least 2.5 times the errors}  % experiments/FIGS/data/v2/index.json :: nonlinear probe qcd_bb_vs_cc: (1 - AUC) of FLAV_F0 over L188, run by run :: sha256 c6c4138ee5c86423 :: prose twin: two significant figures, direction in words
\newcommand{\MlstMlpCmpQcdBbFzeroOverFullPrMax}{up to 2.7 times the errors}  % experiments/FIGS/data/v2/index.json :: nonlinear probe qcd_bb_vs_cc: (1 - AUC) of FLAV_F0 over L188, run by run :: sha256 c6c4138ee5c86423 :: prose twin of the run farther from 1
\newcommand{\MlstMlpCmpQcdBbFoneOverFzero}{0.573 and 0.505}  % experiments/FIGS/data/v2/index.json :: nonlinear probe qcd_bb_vs_cc: (1 - AUC) of FLAV_F1 over FLAV_F0, run by run :: sha256 c6c4138ee5c86423
\newcommand{\MlstMlpCmpQcdBbFoneOverFzeroPr}{at least 43\% fewer errors}  % experiments/FIGS/data/v2/index.json :: nonlinear probe qcd_bb_vs_cc: (1 - AUC) of FLAV_F1 over FLAV_F0, run by run :: sha256 c6c4138ee5c86423 :: prose twin: two significant figures, direction in words
\newcommand{\MlstMlpCmpQcdBbFoneOverFzeroPrMax}{up to 49\% fewer errors}  % experiments/FIGS/data/v2/index.json :: nonlinear probe qcd_bb_vs_cc: (1 - AUC) of FLAV_F1 over FLAV_F0, run by run :: sha256 c6c4138ee5c86423 :: prose twin of the run farther from 1
\newcommand{\MlstMlpCmpQcdBbFonerOverFone}{1.219 and 1.258}  % experiments/FIGS/data/v2/index.json :: nonlinear probe qcd_bb_vs_cc: (1 - AUC) of FLAV_F1R over FLAV_F1, run by run :: sha256 c6c4138ee5c86423
\newcommand{\MlstMlpCmpQcdBbFonerOverFonePr}{at least 22\% more errors}  % experiments/FIGS/data/v2/index.json :: nonlinear probe qcd_bb_vs_cc: (1 - AUC) of FLAV_F1R over FLAV_F1, run by run :: sha256 c6c4138ee5c86423 :: prose twin: two significant figures, direction in words
\newcommand{\MlstMlpCmpQcdBbFonerOverFonePrMax}{up to 26\% more errors}  % experiments/FIGS/data/v2/index.json :: nonlinear probe qcd_bb_vs_cc: (1 - AUC) of FLAV_F1R over FLAV_F1, run by run :: sha256 c6c4138ee5c86423 :: prose twin of the run farther from 1
\newcommand{\MlstMlpCmpQcdBbFonerOverFzero}{0.699 and 0.635}  % experiments/FIGS/data/v2/index.json :: nonlinear probe qcd_bb_vs_cc: (1 - AUC) of FLAV_F1R over FLAV_F0, run by run :: sha256 c6c4138ee5c86423
\newcommand{\MlstMlpCmpQcdBbFonerOverFzeroPr}{at least 30\% fewer errors}  % experiments/FIGS/data/v2/index.json :: nonlinear probe qcd_bb_vs_cc: (1 - AUC) of FLAV_F1R over FLAV_F0, run by run :: sha256 c6c4138ee5c86423 :: prose twin: two significant figures, direction in words
\newcommand{\MlstMlpCmpQcdBbFonerOverFzeroPrMax}{up to 36\% fewer errors}  % experiments/FIGS/data/v2/index.json :: nonlinear probe qcd_bb_vs_cc: (1 - AUC) of FLAV_F1R over FLAV_F0, run by run :: sha256 c6c4138ee5c86423 :: prose twin of the run farther from 1
\newcommand{\MlstMlpCmpQcdBbFzeroOverCoarse}{1.103 and 1.167}  % experiments/FIGS/data/v2/index.json :: nonlinear probe qcd_bb_vs_cc: (1 - AUC) of FLAV_F0 over R16_Q1, run by run :: sha256 c6c4138ee5c86423
\newcommand{\MlstMlpCmpQcdBbFzeroOverCoarsePr}{at least 10\% more errors}  % experiments/FIGS/data/v2/index.json :: nonlinear probe qcd_bb_vs_cc: (1 - AUC) of FLAV_F0 over R16_Q1, run by run :: sha256 c6c4138ee5c86423 :: prose twin: two significant figures, direction in words
\newcommand{\MlstMlpCmpQcdBbFzeroOverCoarsePrMax}{up to 17\% more errors}  % experiments/FIGS/data/v2/index.json :: nonlinear probe qcd_bb_vs_cc: (1 - AUC) of FLAV_F0 over R16_Q1, run by run :: sha256 c6c4138ee5c86423 :: prose twin of the run farther from 1
\newcommand{\MlstMlpCmpQcdBbRandOverFull}{\pending{no run yet}}  % experiments/FIGS/data/v2/index.json :: nonlinear probe qcd_bb_vs_cc: (1 - AUC) of each random partition over the 188-class run of its seed :: sha256 c6c4138ee5c86423
\newcommand{\MlstMlpCmpQcdBbRandOverFullPr}{\pending{no run yet}}  % experiments/FIGS/data/v2/index.json :: nonlinear probe qcd_bb_vs_cc: (1 - AUC) of each random partition over the 188-class run of its seed :: sha256 c6c4138ee5c86423 :: prose twin: two significant figures, direction in words
\newcommand{\MlstMlpCmpQcdBResOverFull}{1.050\,\ensuremath{\pm}\,0.004}  % experiments/FIGS/data/v2/index.json :: nonlinear probe qcd_b_vs_light: (1 - AUC) of L162 over L188, run by run :: sha256 c6c4138ee5c86423
\newcommand{\MlstMlpCmpQcdBResOverFullPr}{5.0\% more errors}  % experiments/FIGS/data/v2/index.json :: nonlinear probe qcd_b_vs_light: (1 - AUC) of L162 over L188, run by run :: sha256 c6c4138ee5c86423 :: prose twin: two significant figures, direction in words
\newcommand{\MlstMlpCmpQcdBMidOverFull}{1.055\,\ensuremath{\pm}\,0.005}  % experiments/FIGS/data/v2/index.json :: nonlinear probe qcd_b_vs_light: (1 - AUC) of R42_Q1 over L188, run by run :: sha256 c6c4138ee5c86423
\newcommand{\MlstMlpCmpQcdBMidOverFullPr}{5.5\% more errors}  % experiments/FIGS/data/v2/index.json :: nonlinear probe qcd_b_vs_light: (1 - AUC) of R42_Q1 over L188, run by run :: sha256 c6c4138ee5c86423 :: prose twin: two significant figures, direction in words
\newcommand{\MlstMlpCmpQcdBCoarseOverFull}{1.337\,\ensuremath{\pm}\,0.021}  % experiments/FIGS/data/v2/index.json :: nonlinear probe qcd_b_vs_light: (1 - AUC) of R16_Q1 over L188, run by run :: sha256 c6c4138ee5c86423
\newcommand{\MlstMlpCmpQcdBCoarseOverFullPr}{34\% more errors}  % experiments/FIGS/data/v2/index.json :: nonlinear probe qcd_b_vs_light: (1 - AUC) of R16_Q1 over L188, run by run :: sha256 c6c4138ee5c86423 :: prose twin: two significant figures, direction in words
\newcommand{\MlstMlpCmpQcdBResMassOverRes}{\pending{no run yet}}  % experiments/FIGS/data/v2/index.json :: nonlinear probe qcd_b_vs_light: (1 - AUC) of L162_MASS over L162, run by run :: sha256 c6c4138ee5c86423
\newcommand{\MlstMlpCmpQcdBResMassOverResPr}{\pending{no run yet}}  % experiments/FIGS/data/v2/index.json :: nonlinear probe qcd_b_vs_light: (1 - AUC) of L162_MASS over L162, run by run :: sha256 c6c4138ee5c86423 :: prose twin: two significant figures, direction in words
\newcommand{\MlstMlpCmpQcdBResMassOverResPrMax}{\pending{no run yet}}  % experiments/FIGS/data/v2/index.json :: nonlinear probe qcd_b_vs_light: (1 - AUC) of L162_MASS over L162, run by run :: sha256 c6c4138ee5c86423 :: prose twin of the run farther from 1
\newcommand{\MlstMlpCmpQcdBCoarseMassOverCoarse}{\pending{no run yet}}  % experiments/FIGS/data/v2/index.json :: nonlinear probe qcd_b_vs_light: (1 - AUC) of R16_Q1_MASS over R16_Q1, run by run :: sha256 c6c4138ee5c86423
\newcommand{\MlstMlpCmpQcdBCoarseMassOverCoarsePr}{\pending{no run yet}}  % experiments/FIGS/data/v2/index.json :: nonlinear probe qcd_b_vs_light: (1 - AUC) of R16_Q1_MASS over R16_Q1, run by run :: sha256 c6c4138ee5c86423 :: prose twin: two significant figures, direction in words
\newcommand{\MlstMlpCmpQcdBCoarseMassOverCoarsePrMax}{\pending{no run yet}}  % experiments/FIGS/data/v2/index.json :: nonlinear probe qcd_b_vs_light: (1 - AUC) of R16_Q1_MASS over R16_Q1, run by run :: sha256 c6c4138ee5c86423 :: prose twin of the run farther from 1
\newcommand{\MlstMlpCmpQcdBFzeroOverFull}{1.315 and 1.288}  % experiments/FIGS/data/v2/index.json :: nonlinear probe qcd_b_vs_light: (1 - AUC) of FLAV_F0 over L188, run by run :: sha256 c6c4138ee5c86423
\newcommand{\MlstMlpCmpQcdBFzeroOverFullPr}{at least 29\% more errors}  % experiments/FIGS/data/v2/index.json :: nonlinear probe qcd_b_vs_light: (1 - AUC) of FLAV_F0 over L188, run by run :: sha256 c6c4138ee5c86423 :: prose twin: two significant figures, direction in words
\newcommand{\MlstMlpCmpQcdBFzeroOverFullPrMax}{up to 32\% more errors}  % experiments/FIGS/data/v2/index.json :: nonlinear probe qcd_b_vs_light: (1 - AUC) of FLAV_F0 over L188, run by run :: sha256 c6c4138ee5c86423 :: prose twin of the run farther from 1
\newcommand{\MlstMlpCmpQcdBFoneOverFzero}{0.8804 and 0.8849}  % experiments/FIGS/data/v2/index.json :: nonlinear probe qcd_b_vs_light: (1 - AUC) of FLAV_F1 over FLAV_F0, run by run :: sha256 c6c4138ee5c86423
\newcommand{\MlstMlpCmpQcdBFoneOverFzeroPr}{at least 12\% fewer errors}  % experiments/FIGS/data/v2/index.json :: nonlinear probe qcd_b_vs_light: (1 - AUC) of FLAV_F1 over FLAV_F0, run by run :: sha256 c6c4138ee5c86423 :: prose twin: two significant figures, direction in words
\newcommand{\MlstMlpCmpQcdBFoneOverFzeroPrMax}{up to 12\% fewer errors}  % experiments/FIGS/data/v2/index.json :: nonlinear probe qcd_b_vs_light: (1 - AUC) of FLAV_F1 over FLAV_F0, run by run :: sha256 c6c4138ee5c86423 :: prose twin of the run farther from 1
\newcommand{\MlstMlpCmpQcdBFonerOverFone}{1.037 and 1.042}  % experiments/FIGS/data/v2/index.json :: nonlinear probe qcd_b_vs_light: (1 - AUC) of FLAV_F1R over FLAV_F1, run by run :: sha256 c6c4138ee5c86423
\newcommand{\MlstMlpCmpQcdBFonerOverFonePr}{at least 3.7\% more errors}  % experiments/FIGS/data/v2/index.json :: nonlinear probe qcd_b_vs_light: (1 - AUC) of FLAV_F1R over FLAV_F1, run by run :: sha256 c6c4138ee5c86423 :: prose twin: two significant figures, direction in words
\newcommand{\MlstMlpCmpQcdBFonerOverFonePrMax}{up to 4.2\% more errors}  % experiments/FIGS/data/v2/index.json :: nonlinear probe qcd_b_vs_light: (1 - AUC) of FLAV_F1R over FLAV_F1, run by run :: sha256 c6c4138ee5c86423 :: prose twin of the run farther from 1
\newcommand{\MlstMlpCmpQcdBFonerOverFzero}{0.9130 and 0.9217}  % experiments/FIGS/data/v2/index.json :: nonlinear probe qcd_b_vs_light: (1 - AUC) of FLAV_F1R over FLAV_F0, run by run :: sha256 c6c4138ee5c86423
\newcommand{\MlstMlpCmpQcdBFonerOverFzeroPr}{at least 7.8\% fewer errors}  % experiments/FIGS/data/v2/index.json :: nonlinear probe qcd_b_vs_light: (1 - AUC) of FLAV_F1R over FLAV_F0, run by run :: sha256 c6c4138ee5c86423 :: prose twin: two significant figures, direction in words
\newcommand{\MlstMlpCmpQcdBFonerOverFzeroPrMax}{up to 8.7\% fewer errors}  % experiments/FIGS/data/v2/index.json :: nonlinear probe qcd_b_vs_light: (1 - AUC) of FLAV_F1R over FLAV_F0, run by run :: sha256 c6c4138ee5c86423 :: prose twin of the run farther from 1
\newcommand{\MlstMlpCmpQcdBFzeroOverCoarse}{0.968 and 0.978}  % experiments/FIGS/data/v2/index.json :: nonlinear probe qcd_b_vs_light: (1 - AUC) of FLAV_F0 over R16_Q1, run by run :: sha256 c6c4138ee5c86423
\newcommand{\MlstMlpCmpQcdBFzeroOverCoarsePr}{at least 2.2\% fewer errors}  % experiments/FIGS/data/v2/index.json :: nonlinear probe qcd_b_vs_light: (1 - AUC) of FLAV_F0 over R16_Q1, run by run :: sha256 c6c4138ee5c86423 :: prose twin: two significant figures, direction in words
\newcommand{\MlstMlpCmpQcdBFzeroOverCoarsePrMax}{up to 3.2\% fewer errors}  % experiments/FIGS/data/v2/index.json :: nonlinear probe qcd_b_vs_light: (1 - AUC) of FLAV_F0 over R16_Q1, run by run :: sha256 c6c4138ee5c86423 :: prose twin of the run farther from 1
\newcommand{\MlstMlpCmpQcdBRandOverFull}{\pending{no run yet}}  % experiments/FIGS/data/v2/index.json :: nonlinear probe qcd_b_vs_light: (1 - AUC) of each random partition over the 188-class run of its seed :: sha256 c6c4138ee5c86423
\newcommand{\MlstMlpCmpQcdBRandOverFullPr}{\pending{no run yet}}  % experiments/FIGS/data/v2/index.json :: nonlinear probe qcd_b_vs_light: (1 - AUC) of each random partition over the 188-class run of its seed :: sha256 c6c4138ee5c86423 :: prose twin: two significant figures, direction in words
\newcommand{\MlstMlpCmpQcdCResOverFull}{1.039\,\ensuremath{\pm}\,0.005}  % experiments/FIGS/data/v2/index.json :: nonlinear probe qcd_c_vs_light: (1 - AUC) of L162 over L188, run by run :: sha256 c6c4138ee5c86423
\newcommand{\MlstMlpCmpQcdCResOverFullPr}{3.9\% more errors}  % experiments/FIGS/data/v2/index.json :: nonlinear probe qcd_c_vs_light: (1 - AUC) of L162 over L188, run by run :: sha256 c6c4138ee5c86423 :: prose twin: two significant figures, direction in words
\newcommand{\MlstMlpCmpQcdCMidOverFull}{1.055\,\ensuremath{\pm}\,0.006}  % experiments/FIGS/data/v2/index.json :: nonlinear probe qcd_c_vs_light: (1 - AUC) of R42_Q1 over L188, run by run :: sha256 c6c4138ee5c86423
\newcommand{\MlstMlpCmpQcdCMidOverFullPr}{5.5\% more errors}  % experiments/FIGS/data/v2/index.json :: nonlinear probe qcd_c_vs_light: (1 - AUC) of R42_Q1 over L188, run by run :: sha256 c6c4138ee5c86423 :: prose twin: two significant figures, direction in words
\newcommand{\MlstMlpCmpQcdCCoarseOverFull}{1.274\,\ensuremath{\pm}\,0.008}  % experiments/FIGS/data/v2/index.json :: nonlinear probe qcd_c_vs_light: (1 - AUC) of R16_Q1 over L188, run by run :: sha256 c6c4138ee5c86423
\newcommand{\MlstMlpCmpQcdCCoarseOverFullPr}{27\% more errors}  % experiments/FIGS/data/v2/index.json :: nonlinear probe qcd_c_vs_light: (1 - AUC) of R16_Q1 over L188, run by run :: sha256 c6c4138ee5c86423 :: prose twin: two significant figures, direction in words
\newcommand{\MlstMlpCmpQcdCResMassOverRes}{\pending{no run yet}}  % experiments/FIGS/data/v2/index.json :: nonlinear probe qcd_c_vs_light: (1 - AUC) of L162_MASS over L162, run by run :: sha256 c6c4138ee5c86423
\newcommand{\MlstMlpCmpQcdCResMassOverResPr}{\pending{no run yet}}  % experiments/FIGS/data/v2/index.json :: nonlinear probe qcd_c_vs_light: (1 - AUC) of L162_MASS over L162, run by run :: sha256 c6c4138ee5c86423 :: prose twin: two significant figures, direction in words
\newcommand{\MlstMlpCmpQcdCResMassOverResPrMax}{\pending{no run yet}}  % experiments/FIGS/data/v2/index.json :: nonlinear probe qcd_c_vs_light: (1 - AUC) of L162_MASS over L162, run by run :: sha256 c6c4138ee5c86423 :: prose twin of the run farther from 1
\newcommand{\MlstMlpCmpQcdCCoarseMassOverCoarse}{\pending{no run yet}}  % experiments/FIGS/data/v2/index.json :: nonlinear probe qcd_c_vs_light: (1 - AUC) of R16_Q1_MASS over R16_Q1, run by run :: sha256 c6c4138ee5c86423
\newcommand{\MlstMlpCmpQcdCCoarseMassOverCoarsePr}{\pending{no run yet}}  % experiments/FIGS/data/v2/index.json :: nonlinear probe qcd_c_vs_light: (1 - AUC) of R16_Q1_MASS over R16_Q1, run by run :: sha256 c6c4138ee5c86423 :: prose twin: two significant figures, direction in words
\newcommand{\MlstMlpCmpQcdCCoarseMassOverCoarsePrMax}{\pending{no run yet}}  % experiments/FIGS/data/v2/index.json :: nonlinear probe qcd_c_vs_light: (1 - AUC) of R16_Q1_MASS over R16_Q1, run by run :: sha256 c6c4138ee5c86423 :: prose twin of the run farther from 1
\newcommand{\MlstMlpCmpQcdCFzeroOverFull}{1.234 and 1.211}  % experiments/FIGS/data/v2/index.json :: nonlinear probe qcd_c_vs_light: (1 - AUC) of FLAV_F0 over L188, run by run :: sha256 c6c4138ee5c86423
\newcommand{\MlstMlpCmpQcdCFzeroOverFullPr}{at least 21\% more errors}  % experiments/FIGS/data/v2/index.json :: nonlinear probe qcd_c_vs_light: (1 - AUC) of FLAV_F0 over L188, run by run :: sha256 c6c4138ee5c86423 :: prose twin: two significant figures, direction in words
\newcommand{\MlstMlpCmpQcdCFzeroOverFullPrMax}{up to 23\% more errors}  % experiments/FIGS/data/v2/index.json :: nonlinear probe qcd_c_vs_light: (1 - AUC) of FLAV_F0 over L188, run by run :: sha256 c6c4138ee5c86423 :: prose twin of the run farther from 1
\newcommand{\MlstMlpCmpQcdCFoneOverFzero}{0.9395 and 0.9472}  % experiments/FIGS/data/v2/index.json :: nonlinear probe qcd_c_vs_light: (1 - AUC) of FLAV_F1 over FLAV_F0, run by run :: sha256 c6c4138ee5c86423
\newcommand{\MlstMlpCmpQcdCFoneOverFzeroPr}{at least 5.3\% fewer errors}  % experiments/FIGS/data/v2/index.json :: nonlinear probe qcd_c_vs_light: (1 - AUC) of FLAV_F1 over FLAV_F0, run by run :: sha256 c6c4138ee5c86423 :: prose twin: two significant figures, direction in words
\newcommand{\MlstMlpCmpQcdCFoneOverFzeroPrMax}{up to 6.0\% fewer errors}  % experiments/FIGS/data/v2/index.json :: nonlinear probe qcd_c_vs_light: (1 - AUC) of FLAV_F1 over FLAV_F0, run by run :: sha256 c6c4138ee5c86423 :: prose twin of the run farther from 1
\newcommand{\MlstMlpCmpQcdCFonerOverFone}{0.9846 and 0.9767}  % experiments/FIGS/data/v2/index.json :: nonlinear probe qcd_c_vs_light: (1 - AUC) of FLAV_F1R over FLAV_F1, run by run :: sha256 c6c4138ee5c86423
\newcommand{\MlstMlpCmpQcdCFonerOverFonePr}{at least 1.5\% fewer errors}  % experiments/FIGS/data/v2/index.json :: nonlinear probe qcd_c_vs_light: (1 - AUC) of FLAV_F1R over FLAV_F1, run by run :: sha256 c6c4138ee5c86423 :: prose twin: two significant figures, direction in words
\newcommand{\MlstMlpCmpQcdCFonerOverFonePrMax}{up to 2.3\% fewer errors}  % experiments/FIGS/data/v2/index.json :: nonlinear probe qcd_c_vs_light: (1 - AUC) of FLAV_F1R over FLAV_F1, run by run :: sha256 c6c4138ee5c86423 :: prose twin of the run farther from 1
\newcommand{\MlstMlpCmpQcdCFonerOverFzero}{0.9251 and 0.9251}  % experiments/FIGS/data/v2/index.json :: nonlinear probe qcd_c_vs_light: (1 - AUC) of FLAV_F1R over FLAV_F0, run by run :: sha256 c6c4138ee5c86423
\newcommand{\MlstMlpCmpQcdCFonerOverFzeroPr}{at least 7.5\% fewer errors}  % experiments/FIGS/data/v2/index.json :: nonlinear probe qcd_c_vs_light: (1 - AUC) of FLAV_F1R over FLAV_F0, run by run :: sha256 c6c4138ee5c86423 :: prose twin: two significant figures, direction in words
\newcommand{\MlstMlpCmpQcdCFonerOverFzeroPrMax}{up to 7.5\% fewer errors}  % experiments/FIGS/data/v2/index.json :: nonlinear probe qcd_c_vs_light: (1 - AUC) of FLAV_F1R over FLAV_F0, run by run :: sha256 c6c4138ee5c86423 :: prose twin of the run farther from 1
\newcommand{\MlstMlpCmpQcdCFzeroOverCoarse}{0.968 and 0.956}  % experiments/FIGS/data/v2/index.json :: nonlinear probe qcd_c_vs_light: (1 - AUC) of FLAV_F0 over R16_Q1, run by run :: sha256 c6c4138ee5c86423
\newcommand{\MlstMlpCmpQcdCFzeroOverCoarsePr}{at least 3.2\% fewer errors}  % experiments/FIGS/data/v2/index.json :: nonlinear probe qcd_c_vs_light: (1 - AUC) of FLAV_F0 over R16_Q1, run by run :: sha256 c6c4138ee5c86423 :: prose twin: two significant figures, direction in words
\newcommand{\MlstMlpCmpQcdCFzeroOverCoarsePrMax}{up to 4.4\% fewer errors}  % experiments/FIGS/data/v2/index.json :: nonlinear probe qcd_c_vs_light: (1 - AUC) of FLAV_F0 over R16_Q1, run by run :: sha256 c6c4138ee5c86423 :: prose twin of the run farther from 1
\newcommand{\MlstMlpCmpQcdCRandOverFull}{\pending{no run yet}}  % experiments/FIGS/data/v2/index.json :: nonlinear probe qcd_c_vs_light: (1 - AUC) of each random partition over the 188-class run of its seed :: sha256 c6c4138ee5c86423
\newcommand{\MlstMlpCmpQcdCRandOverFullPr}{\pending{no run yet}}  % experiments/FIGS/data/v2/index.json :: nonlinear probe qcd_c_vs_light: (1 - AUC) of each random partition over the 188-class run of its seed :: sha256 c6c4138ee5c86423 :: prose twin: two significant figures, direction in words
\newcommand{\MlstMlpTopRthirtyFullOverInitBn}{1.00\,\ensuremath{\pm}\,0.31}  % experiments/FIGS/data/v2/index.json :: nonlinear probe top, missing inputs zero: R30 of L188 over the BatchNorm-recomputed random network of the same index :: sha256 c6c4138ee5c86423
\newcommand{\MlstMlpTopRthirtyFullOverInitBnPr}{equal within 40\%}  % experiments/FIGS/data/v2/index.json :: nonlinear probe top, missing inputs zero: R30 of L188 over the BatchNorm-recomputed random network of the same index :: sha256 c6c4138ee5c86423 :: prose twin: two significant figures, direction in words
\newcommand{\MlstMlpTopRthirtyResOverInitBn}{0.94\,\ensuremath{\pm}\,0.26}  % experiments/FIGS/data/v2/index.json :: nonlinear probe top, missing inputs zero: R30 of L162 over the BatchNorm-recomputed random network of the same index :: sha256 c6c4138ee5c86423
\newcommand{\MlstMlpTopRthirtyResOverInitBnPr}{equal within 30\%}  % experiments/FIGS/data/v2/index.json :: nonlinear probe top, missing inputs zero: R30 of L162 over the BatchNorm-recomputed random network of the same index :: sha256 c6c4138ee5c86423 :: prose twin: two significant figures, direction in words
\newcommand{\MlstMlpTopRthirtyMidOverInitBn}{0.84\,\ensuremath{\pm}\,0.17}  % experiments/FIGS/data/v2/index.json :: nonlinear probe top, missing inputs zero: R30 of R42_Q1 over the BatchNorm-recomputed random network of the same index :: sha256 c6c4138ee5c86423
\newcommand{\MlstMlpTopRthirtyMidOverInitBnPr}{equal within 20\%}  % experiments/FIGS/data/v2/index.json :: nonlinear probe top, missing inputs zero: R30 of R42_Q1 over the BatchNorm-recomputed random network of the same index :: sha256 c6c4138ee5c86423 :: prose twin: two significant figures, direction in words
\newcommand{\MlstMlpTopRthirtyCoarseOverInitBn}{0.75\,\ensuremath{\pm}\,0.05}  % experiments/FIGS/data/v2/index.json :: nonlinear probe top, missing inputs zero: R30 of R16_Q1 over the BatchNorm-recomputed random network of the same index :: sha256 c6c4138ee5c86423
\newcommand{\MlstMlpTopRthirtyCoarseOverInitBnPr}{25\% less rejection}  % experiments/FIGS/data/v2/index.json :: nonlinear probe top, missing inputs zero: R30 of R16_Q1 over the BatchNorm-recomputed random network of the same index :: sha256 c6c4138ee5c86423 :: prose twin: two significant figures, direction in words
\newcommand{\MlstMlpQgRthirtyFullOverInitBn}{1.083\,\ensuremath{\pm}\,0.026}  % experiments/FIGS/data/v2/index.json :: nonlinear probe qg, missing inputs zero: R30 of L188 over the BatchNorm-recomputed random network of the same index :: sha256 c6c4138ee5c86423
\newcommand{\MlstMlpQgRthirtyFullOverInitBnPr}{8.3\% more rejection}  % experiments/FIGS/data/v2/index.json :: nonlinear probe qg, missing inputs zero: R30 of L188 over the BatchNorm-recomputed random network of the same index :: sha256 c6c4138ee5c86423 :: prose twin: two significant figures, direction in words
\newcommand{\MlstMlpQgRthirtyResOverInitBn}{1.010\,\ensuremath{\pm}\,0.016}  % experiments/FIGS/data/v2/index.json :: nonlinear probe qg, missing inputs zero: R30 of L162 over the BatchNorm-recomputed random network of the same index :: sha256 c6c4138ee5c86423
\newcommand{\MlstMlpQgRthirtyResOverInitBnPr}{equal within 2\%}  % experiments/FIGS/data/v2/index.json :: nonlinear probe qg, missing inputs zero: R30 of L162 over the BatchNorm-recomputed random network of the same index :: sha256 c6c4138ee5c86423 :: prose twin: two significant figures, direction in words
\newcommand{\MlstMlpQgRthirtyMidOverInitBn}{1.06\,\ensuremath{\pm}\,0.05}  % experiments/FIGS/data/v2/index.json :: nonlinear probe qg, missing inputs zero: R30 of R42_Q1 over the BatchNorm-recomputed random network of the same index :: sha256 c6c4138ee5c86423
\newcommand{\MlstMlpQgRthirtyMidOverInitBnPr}{equal within 7\%}  % experiments/FIGS/data/v2/index.json :: nonlinear probe qg, missing inputs zero: R30 of R42_Q1 over the BatchNorm-recomputed random network of the same index :: sha256 c6c4138ee5c86423 :: prose twin: two significant figures, direction in words
\newcommand{\MlstMlpQgRthirtyCoarseOverInitBn}{1.04\,\ensuremath{\pm}\,0.04}  % experiments/FIGS/data/v2/index.json :: nonlinear probe qg, missing inputs zero: R30 of R16_Q1 over the BatchNorm-recomputed random network of the same index :: sha256 c6c4138ee5c86423
\newcommand{\MlstMlpQgRthirtyCoarseOverInitBnPr}{equal within 5\%}  % experiments/FIGS/data/v2/index.json :: nonlinear probe qg, missing inputs zero: R30 of R16_Q1 over the BatchNorm-recomputed random network of the same index :: sha256 c6c4138ee5c86423 :: prose twin: two significant figures, direction in words
\newcommand{\MlstMlpQgHerwigRthirtyFullOverInitBn}{0.971\,\ensuremath{\pm}\,0.034}  % experiments/FIGS/data/v2/index.json :: nonlinear probe qg_herwig, missing inputs zero: R30 of L188 over the BatchNorm-recomputed random network of the same index :: sha256 c6c4138ee5c86423
\newcommand{\MlstMlpQgHerwigRthirtyFullOverInitBnPr}{equal within 4\%}  % experiments/FIGS/data/v2/index.json :: nonlinear probe qg_herwig, missing inputs zero: R30 of L188 over the BatchNorm-recomputed random network of the same index :: sha256 c6c4138ee5c86423 :: prose twin: two significant figures, direction in words
\newcommand{\MlstMlpQgHerwigRthirtyResOverInitBn}{0.94\,\ensuremath{\pm}\,0.05}  % experiments/FIGS/data/v2/index.json :: nonlinear probe qg_herwig, missing inputs zero: R30 of L162 over the BatchNorm-recomputed random network of the same index :: sha256 c6c4138ee5c86423
\newcommand{\MlstMlpQgHerwigRthirtyResOverInitBnPr}{equal within 6\%}  % experiments/FIGS/data/v2/index.json :: nonlinear probe qg_herwig, missing inputs zero: R30 of L162 over the BatchNorm-recomputed random network of the same index :: sha256 c6c4138ee5c86423 :: prose twin: two significant figures, direction in words
\newcommand{\MlstMlpQgHerwigRthirtyMidOverInitBn}{0.945\,\ensuremath{\pm}\,0.030}  % experiments/FIGS/data/v2/index.json :: nonlinear probe qg_herwig, missing inputs zero: R30 of R42_Q1 over the BatchNorm-recomputed random network of the same index :: sha256 c6c4138ee5c86423
\newcommand{\MlstMlpQgHerwigRthirtyMidOverInitBnPr}{equal within 6\%}  % experiments/FIGS/data/v2/index.json :: nonlinear probe qg_herwig, missing inputs zero: R30 of R42_Q1 over the BatchNorm-recomputed random network of the same index :: sha256 c6c4138ee5c86423 :: prose twin: two significant figures, direction in words
\newcommand{\MlstMlpQgHerwigRthirtyCoarseOverInitBn}{0.96\,\ensuremath{\pm}\,0.04}  % experiments/FIGS/data/v2/index.json :: nonlinear probe qg_herwig, missing inputs zero: R30 of R16_Q1 over the BatchNorm-recomputed random network of the same index :: sha256 c6c4138ee5c86423
\newcommand{\MlstMlpQgHerwigRthirtyCoarseOverInitBnPr}{equal within 5\%}  % experiments/FIGS/data/v2/index.json :: nonlinear probe qg_herwig, missing inputs zero: R30 of R16_Q1 over the BatchNorm-recomputed random network of the same index :: sha256 c6c4138ee5c86423 :: prose twin: two significant figures, direction in words
\newcommand{\MlstMlpNatJcTwoResMassOverRes}{1.016 and 1.009}  % experiments/FIGS/data/v2/index.json :: nonlinear probe, natural-composition JetClass-II (fits_v3_mlp2): (1 - AUC) of L162_MASS over L162, run by run :: sha256 c6c4138ee5c86423
\newcommand{\MlstMlpNatJcTwoResMassOverResPr}{equal within 2\%}  % experiments/FIGS/data/v2/index.json :: nonlinear probe, natural-composition JetClass-II (fits_v3_mlp2): (1 - AUC) of L162_MASS over L162, run by run :: sha256 c6c4138ee5c86423 :: prose twin: two significant figures, direction in words
\newcommand{\MlstMlpNatJcTwoResMassOverResPrMax}{equal within 2\%}  % experiments/FIGS/data/v2/index.json :: nonlinear probe, natural-composition JetClass-II (fits_v3_mlp2): (1 - AUC) of L162_MASS over L162, run by run :: sha256 c6c4138ee5c86423 :: prose twin of the run farther from 1
\newcommand{\MlstMlpNatJcTwoCoarseMassOverCoarse}{1.026 and 1.019}  % experiments/FIGS/data/v2/index.json :: nonlinear probe, natural-composition JetClass-II (fits_v3_mlp2): (1 - AUC) of R16_Q1_MASS over R16_Q1, run by run :: sha256 c6c4138ee5c86423
\newcommand{\MlstMlpNatJcTwoCoarseMassOverCoarsePr}{equal within 3\%}  % experiments/FIGS/data/v2/index.json :: nonlinear probe, natural-composition JetClass-II (fits_v3_mlp2): (1 - AUC) of R16_Q1_MASS over R16_Q1, run by run :: sha256 c6c4138ee5c86423 :: prose twin: two significant figures, direction in words
\newcommand{\MlstMlpNatJcTwoCoarseMassOverCoarsePrMax}{equal within 3\%}  % experiments/FIGS/data/v2/index.json :: nonlinear probe, natural-composition JetClass-II (fits_v3_mlp2): (1 - AUC) of R16_Q1_MASS over R16_Q1, run by run :: sha256 c6c4138ee5c86423 :: prose twin of the run farther from 1
\newcommand{\LpSubstitutionTop}{random draws}  % experiments/FIGS/data/v2/index.json :: linprobe/fits_v4_fill: lowest summed validation cross-entropy over every model, largest training set :: sha256 c6c4138ee5c86423 :: zero 8.8428; mean 8.99394; draws 7.00838
\newcommand{\LpSubstitutionQg}{the training mean}  % experiments/FIGS/data/v2/index.json :: linprobe/fits_v4_fill: lowest summed validation cross-entropy over every model, largest training set :: sha256 c6c4138ee5c86423 :: zero 13.6224; mean 13.1964; draws 13.3584
\newcommand{\MlstLpTopRthirtyFullOverInitBn}{4.67\,\ensuremath{\pm}\,0.23}  % experiments/FIGS/data/v2/index.json :: linear probe top, largest set, chosen substitution: R30 of L188 over the BatchNorm-recomputed random network of the same index :: sha256 c6c4138ee5c86423
\newcommand{\MlstLpTopRthirtyFullOverInitBnPr}{4.7 times the rejection}  % experiments/FIGS/data/v2/index.json :: linear probe top, largest set, chosen substitution: R30 of L188 over the BatchNorm-recomputed random network of the same index :: sha256 c6c4138ee5c86423 :: prose twin: two significant figures, direction in words
\newcommand{\MlstLpTopRthirtyFullOverInitBnZero}{1.04\,\ensuremath{\pm}\,0.05}  % experiments/FIGS/data/v2/index.json :: linear probe top, largest set, zero: R30 of L188 over the BatchNorm-recomputed random network of the same index :: sha256 c6c4138ee5c86423
\newcommand{\MlstLpTopRthirtyFullOverInitBnZeroPr}{equal within 5\%}  % experiments/FIGS/data/v2/index.json :: linear probe top, largest set, zero: R30 of L188 over the BatchNorm-recomputed random network of the same index :: sha256 c6c4138ee5c86423 :: prose twin: two significant figures, direction in words
\newcommand{\MlstLpTopRthirtyResOverInitBn}{4.08\,\ensuremath{\pm}\,0.28}  % experiments/FIGS/data/v2/index.json :: linear probe top, largest set, chosen substitution: R30 of L162 over the BatchNorm-recomputed random network of the same index :: sha256 c6c4138ee5c86423
\newcommand{\MlstLpTopRthirtyResOverInitBnPr}{4.1 times the rejection}  % experiments/FIGS/data/v2/index.json :: linear probe top, largest set, chosen substitution: R30 of L162 over the BatchNorm-recomputed random network of the same index :: sha256 c6c4138ee5c86423 :: prose twin: two significant figures, direction in words
\newcommand{\MlstLpTopRthirtyResOverInitBnZero}{1.33\,\ensuremath{\pm}\,0.30}  % experiments/FIGS/data/v2/index.json :: linear probe top, largest set, zero: R30 of L162 over the BatchNorm-recomputed random network of the same index :: sha256 c6c4138ee5c86423
\newcommand{\MlstLpTopRthirtyResOverInitBnZeroPr}{equal within 40\%}  % experiments/FIGS/data/v2/index.json :: linear probe top, largest set, zero: R30 of L162 over the BatchNorm-recomputed random network of the same index :: sha256 c6c4138ee5c86423 :: prose twin: two significant figures, direction in words
\newcommand{\MlstLpTopRthirtyMidOverInitBn}{4.76\,\ensuremath{\pm}\,0.15}  % experiments/FIGS/data/v2/index.json :: linear probe top, largest set, chosen substitution: R30 of R42_Q1 over the BatchNorm-recomputed random network of the same index :: sha256 c6c4138ee5c86423
\newcommand{\MlstLpTopRthirtyMidOverInitBnPr}{4.8 times the rejection}  % experiments/FIGS/data/v2/index.json :: linear probe top, largest set, chosen substitution: R30 of R42_Q1 over the BatchNorm-recomputed random network of the same index :: sha256 c6c4138ee5c86423 :: prose twin: two significant figures, direction in words
\newcommand{\MlstLpTopRthirtyMidOverInitBnZero}{1.01\,\ensuremath{\pm}\,0.28}  % experiments/FIGS/data/v2/index.json :: linear probe top, largest set, zero: R30 of R42_Q1 over the BatchNorm-recomputed random network of the same index :: sha256 c6c4138ee5c86423
\newcommand{\MlstLpTopRthirtyMidOverInitBnZeroPr}{equal within 30\%}  % experiments/FIGS/data/v2/index.json :: linear probe top, largest set, zero: R30 of R42_Q1 over the BatchNorm-recomputed random network of the same index :: sha256 c6c4138ee5c86423 :: prose twin: two significant figures, direction in words
\newcommand{\MlstLpTopRthirtyCoarseOverInitBn}{5.24\,\ensuremath{\pm}\,0.06}  % experiments/FIGS/data/v2/index.json :: linear probe top, largest set, chosen substitution: R30 of R16_Q1 over the BatchNorm-recomputed random network of the same index :: sha256 c6c4138ee5c86423
\newcommand{\MlstLpTopRthirtyCoarseOverInitBnPr}{5.2 times the rejection}  % experiments/FIGS/data/v2/index.json :: linear probe top, largest set, chosen substitution: R30 of R16_Q1 over the BatchNorm-recomputed random network of the same index :: sha256 c6c4138ee5c86423 :: prose twin: two significant figures, direction in words
\newcommand{\MlstLpTopRthirtyCoarseOverInitBnZero}{0.97\,\ensuremath{\pm}\,0.16}  % experiments/FIGS/data/v2/index.json :: linear probe top, largest set, zero: R30 of R16_Q1 over the BatchNorm-recomputed random network of the same index :: sha256 c6c4138ee5c86423
\newcommand{\MlstLpTopRthirtyCoarseOverInitBnZeroPr}{equal within 20\%}  % experiments/FIGS/data/v2/index.json :: linear probe top, largest set, zero: R30 of R16_Q1 over the BatchNorm-recomputed random network of the same index :: sha256 c6c4138ee5c86423 :: prose twin: two significant figures, direction in words
\newcommand{\MlstLpTopRthirtyResMassOverInitBn}{5.71 and 6.25}  % experiments/FIGS/data/v2/index.json :: linear probe top, largest set, chosen substitution: R30 of L162_MASS over the BatchNorm-recomputed random network of the same index :: sha256 c6c4138ee5c86423
\newcommand{\MlstLpTopRthirtyResMassOverInitBnPr}{at least 5.7 times the rejection}  % experiments/FIGS/data/v2/index.json :: linear probe top, largest set, chosen substitution: R30 of L162_MASS over the BatchNorm-recomputed random network of the same index :: sha256 c6c4138ee5c86423 :: prose twin: two significant figures, direction in words
\newcommand{\MlstLpTopRthirtyResMassOverInitBnPrMax}{up to 6.3 times the rejection}  % experiments/FIGS/data/v2/index.json :: linear probe top, largest set, chosen substitution: R30 of L162_MASS over the BatchNorm-recomputed random network of the same index :: sha256 c6c4138ee5c86423 :: prose twin of the run farther from 1
\newcommand{\MlstLpTopRthirtyResMassOverInitBnZero}{1.683 and 1.678}  % experiments/FIGS/data/v2/index.json :: linear probe top, largest set, zero: R30 of L162_MASS over the BatchNorm-recomputed random network of the same index :: sha256 c6c4138ee5c86423
\newcommand{\MlstLpTopRthirtyResMassOverInitBnZeroPr}{at least 68\% more rejection}  % experiments/FIGS/data/v2/index.json :: linear probe top, largest set, zero: R30 of L162_MASS over the BatchNorm-recomputed random network of the same index :: sha256 c6c4138ee5c86423 :: prose twin: two significant figures, direction in words
\newcommand{\MlstLpTopRthirtyResMassOverInitBnZeroPrMax}{up to 68\% more rejection}  % experiments/FIGS/data/v2/index.json :: linear probe top, largest set, zero: R30 of L162_MASS over the BatchNorm-recomputed random network of the same index :: sha256 c6c4138ee5c86423 :: prose twin of the run farther from 1
\newcommand{\MlstLpTopRthirtyCoarseMassOverInitBn}{6.73 and 7.72}  % experiments/FIGS/data/v2/index.json :: linear probe top, largest set, chosen substitution: R30 of R16_Q1_MASS over the BatchNorm-recomputed random network of the same index :: sha256 c6c4138ee5c86423
\newcommand{\MlstLpTopRthirtyCoarseMassOverInitBnPr}{at least 6.7 times the rejection}  % experiments/FIGS/data/v2/index.json :: linear probe top, largest set, chosen substitution: R30 of R16_Q1_MASS over the BatchNorm-recomputed random network of the same index :: sha256 c6c4138ee5c86423 :: prose twin: two significant figures, direction in words
\newcommand{\MlstLpTopRthirtyCoarseMassOverInitBnPrMax}{up to 7.7 times the rejection}  % experiments/FIGS/data/v2/index.json :: linear probe top, largest set, chosen substitution: R30 of R16_Q1_MASS over the BatchNorm-recomputed random network of the same index :: sha256 c6c4138ee5c86423 :: prose twin of the run farther from 1
\newcommand{\MlstLpTopRthirtyCoarseMassOverInitBnZero}{2.3 and 1.2}  % experiments/FIGS/data/v2/index.json :: linear probe top, largest set, zero: R30 of R16_Q1_MASS over the BatchNorm-recomputed random network of the same index :: sha256 c6c4138ee5c86423
\newcommand{\MlstLpTopRthirtyCoarseMassOverInitBnZeroPr}{equal within 200\%}  % experiments/FIGS/data/v2/index.json :: linear probe top, largest set, zero: R30 of R16_Q1_MASS over the BatchNorm-recomputed random network of the same index :: sha256 c6c4138ee5c86423 :: prose twin: two significant figures, direction in words
\newcommand{\MlstLpTopRthirtyCoarseMassOverInitBnZeroPrMax}{equal within 200\%}  % experiments/FIGS/data/v2/index.json :: linear probe top, largest set, zero: R30 of R16_Q1_MASS over the BatchNorm-recomputed random network of the same index :: sha256 c6c4138ee5c86423 :: prose twin of the run farther from 1
\newcommand{\MlstLpTopRthirtyResOverFull}{0.874\,\ensuremath{\pm}\,0.018}  % experiments/FIGS/data/v2/index.json :: linear probe top: R30 of L162 over the 188-class run :: sha256 c6c4138ee5c86423
\newcommand{\MlstLpTopRthirtyResOverFullPr}{13\% less rejection}  % experiments/FIGS/data/v2/index.json :: linear probe top: R30 of L162 over the 188-class run :: sha256 c6c4138ee5c86423 :: prose twin: two significant figures, direction in words
\newcommand{\MlstLpTopRthirtyResOverFullZero}{1.29\,\ensuremath{\pm}\,0.31}  % experiments/FIGS/data/v2/index.json :: linear probe top: R30 of L162 over the 188-class run :: sha256 c6c4138ee5c86423
\newcommand{\MlstLpTopRthirtyResOverFullZeroPr}{equal within 40\%}  % experiments/FIGS/data/v2/index.json :: linear probe top: R30 of L162 over the 188-class run :: sha256 c6c4138ee5c86423 :: prose twin: two significant figures, direction in words
\newcommand{\MlstLpTopRthirtyMidOverFull}{1.02\,\ensuremath{\pm}\,0.08}  % experiments/FIGS/data/v2/index.json :: linear probe top: R30 of R42_Q1 over the 188-class run :: sha256 c6c4138ee5c86423
\newcommand{\MlstLpTopRthirtyMidOverFullPr}{equal within 9\%}  % experiments/FIGS/data/v2/index.json :: linear probe top: R30 of R42_Q1 over the 188-class run :: sha256 c6c4138ee5c86423 :: prose twin: two significant figures, direction in words
\newcommand{\MlstLpTopRthirtyMidOverFullZero}{0.98\,\ensuremath{\pm}\,0.28}  % experiments/FIGS/data/v2/index.json :: linear probe top: R30 of R42_Q1 over the 188-class run :: sha256 c6c4138ee5c86423
\newcommand{\MlstLpTopRthirtyMidOverFullZeroPr}{equal within 30\%}  % experiments/FIGS/data/v2/index.json :: linear probe top: R30 of R42_Q1 over the 188-class run :: sha256 c6c4138ee5c86423 :: prose twin: two significant figures, direction in words
\newcommand{\MlstLpTopRthirtyCoarseOverFull}{1.12\,\ensuremath{\pm}\,0.06}  % experiments/FIGS/data/v2/index.json :: linear probe top: R30 of R16_Q1 over the 188-class run :: sha256 c6c4138ee5c86423
\newcommand{\MlstLpTopRthirtyCoarseOverFullPr}{12\% more rejection}  % experiments/FIGS/data/v2/index.json :: linear probe top: R30 of R16_Q1 over the 188-class run :: sha256 c6c4138ee5c86423 :: prose twin: two significant figures, direction in words
\newcommand{\MlstLpTopRthirtyCoarseOverFullZero}{0.93\,\ensuremath{\pm}\,0.15}  % experiments/FIGS/data/v2/index.json :: linear probe top: R30 of R16_Q1 over the 188-class run :: sha256 c6c4138ee5c86423
\newcommand{\MlstLpTopRthirtyCoarseOverFullZeroPr}{equal within 20\%}  % experiments/FIGS/data/v2/index.json :: linear probe top: R30 of R16_Q1 over the 188-class run :: sha256 c6c4138ee5c86423 :: prose twin: two significant figures, direction in words
\newcommand{\MlstLpTopRthirtyMeanOverZeroFull}{1.25\,\ensuremath{\pm}\,0.27}  % experiments/FIGS/data/v2/index.json :: linear probe top: R30 with mean over R30 with zero, L188 :: sha256 c6c4138ee5c86423
\newcommand{\MlstLpTopRthirtyMeanOverZeroFullPr}{equal within 30\%}  % experiments/FIGS/data/v2/index.json :: linear probe top: R30 with mean over R30 with zero, L188 :: sha256 c6c4138ee5c86423 :: prose twin: two significant figures, direction in words
\newcommand{\MlstLpTopRthirtyDrawOverZeroFull}{3.9\,\ensuremath{\pm}\,0.4}  % experiments/FIGS/data/v2/index.json :: linear probe top: R30 with draws over R30 with zero, L188 :: sha256 c6c4138ee5c86423
\newcommand{\MlstLpTopRthirtyDrawOverZeroFullPr}{3.9 times the rejection}  % experiments/FIGS/data/v2/index.json :: linear probe top: R30 with draws over R30 with zero, L188 :: sha256 c6c4138ee5c86423 :: prose twin: two significant figures, direction in words
\newcommand{\MlstLpTopRthirtyMeanOverZeroInitBn}{1.08\,\ensuremath{\pm}\,0.20}  % experiments/FIGS/data/v2/index.json :: linear probe top: R30 with mean over R30 with zero, INIT_BN :: sha256 c6c4138ee5c86423
\newcommand{\MlstLpTopRthirtyMeanOverZeroInitBnPr}{equal within 30\%}  % experiments/FIGS/data/v2/index.json :: linear probe top: R30 with mean over R30 with zero, INIT_BN :: sha256 c6c4138ee5c86423 :: prose twin: two significant figures, direction in words
\newcommand{\MlstLpTopRthirtyDrawOverZeroInitBn}{0.86\,\ensuremath{\pm}\,0.14}  % experiments/FIGS/data/v2/index.json :: linear probe top: R30 with draws over R30 with zero, INIT_BN :: sha256 c6c4138ee5c86423
\newcommand{\MlstLpTopRthirtyDrawOverZeroInitBnPr}{equal within 20\%}  % experiments/FIGS/data/v2/index.json :: linear probe top: R30 with draws over R30 with zero, INIT_BN :: sha256 c6c4138ee5c86423 :: prose twin: two significant figures, direction in words
\newcommand{\MlstLpTopRthirtyKinOverFull}{3.88}  % experiments/FIGS/data/v2/index.json :: linear probe top: R30 of the kinematic-only model over the mean of the 188-class runs (zero) :: sha256 c6c4138ee5c86423
\newcommand{\MlstLpTopRthirtyKinOverFullPr}{3.9 times the rejection}  % experiments/FIGS/data/v2/index.json :: linear probe top: R30 of the kinematic-only model over the mean of the 188-class runs (zero) :: sha256 c6c4138ee5c86423 :: prose twin: two significant figures, direction in words
\newcommand{\MlstLpTopRthirtyKin}{231}  % experiments/FIGS/data/v2/index.json :: linear probe top, largest set: R30 of the kinematic-only model (one run) :: sha256 c6c4138ee5c86423
\newcommand{\MlstLpTopRthirtyFullZero}{60\,\ensuremath{\pm}\,12}  % experiments/FIGS/data/v2/index.json :: linear probe top, largest set: R30 with zero, L188 :: sha256 c6c4138ee5c86423
\newcommand{\MlstLpTopRthirtyFullMean}{73\,\ensuremath{\pm}\,15}  % experiments/FIGS/data/v2/index.json :: linear probe top, largest set: R30 with mean, L188 :: sha256 c6c4138ee5c86423
\newcommand{\MlstLpTopRthirtyFullDraws}{228\,\ensuremath{\pm}\,25}  % experiments/FIGS/data/v2/index.json :: linear probe top, largest set: R30 with draws, L188 :: sha256 c6c4138ee5c86423
\newcommand{\MlstLpTopRthirtyInitBnZero}{58\,\ensuremath{\pm}\,13}  % experiments/FIGS/data/v2/index.json :: linear probe top, largest set: R30 with zero, INIT_BN :: sha256 c6c4138ee5c86423
\newcommand{\MlstLpTopRthirtyInitBnMean}{60.4\,\ensuremath{\pm}\,2.4}  % experiments/FIGS/data/v2/index.json :: linear probe top, largest set: R30 with mean, INIT_BN :: sha256 c6c4138ee5c86423
\newcommand{\MlstLpTopRthirtyInitBnDraws}{48.6\,\ensuremath{\pm}\,3.0}  % experiments/FIGS/data/v2/index.json :: linear probe top, largest set: R30 with draws, INIT_BN :: sha256 c6c4138ee5c86423
\newcommand{\MlstLpQgRthirtyFullOverInitBn}{1.39\,\ensuremath{\pm}\,0.06}  % experiments/FIGS/data/v2/index.json :: linear probe qg, largest set, chosen substitution: R30 of L188 over the BatchNorm-recomputed random network of the same index :: sha256 c6c4138ee5c86423
\newcommand{\MlstLpQgRthirtyFullOverInitBnPr}{39\% more rejection}  % experiments/FIGS/data/v2/index.json :: linear probe qg, largest set, chosen substitution: R30 of L188 over the BatchNorm-recomputed random network of the same index :: sha256 c6c4138ee5c86423 :: prose twin: two significant figures, direction in words
\newcommand{\MlstLpQgRthirtyFullOverInitBnZero}{1.21\,\ensuremath{\pm}\,0.15}  % experiments/FIGS/data/v2/index.json :: linear probe qg, largest set, zero: R30 of L188 over the BatchNorm-recomputed random network of the same index :: sha256 c6c4138ee5c86423
\newcommand{\MlstLpQgRthirtyFullOverInitBnZeroPr}{equal within 30\%}  % experiments/FIGS/data/v2/index.json :: linear probe qg, largest set, zero: R30 of L188 over the BatchNorm-recomputed random network of the same index :: sha256 c6c4138ee5c86423 :: prose twin: two significant figures, direction in words
\newcommand{\MlstLpQgRthirtyResOverInitBn}{1.360\,\ensuremath{\pm}\,0.025}  % experiments/FIGS/data/v2/index.json :: linear probe qg, largest set, chosen substitution: R30 of L162 over the BatchNorm-recomputed random network of the same index :: sha256 c6c4138ee5c86423
\newcommand{\MlstLpQgRthirtyResOverInitBnPr}{36\% more rejection}  % experiments/FIGS/data/v2/index.json :: linear probe qg, largest set, chosen substitution: R30 of L162 over the BatchNorm-recomputed random network of the same index :: sha256 c6c4138ee5c86423 :: prose twin: two significant figures, direction in words
\newcommand{\MlstLpQgRthirtyResOverInitBnZero}{1.07\,\ensuremath{\pm}\,0.13}  % experiments/FIGS/data/v2/index.json :: linear probe qg, largest set, zero: R30 of L162 over the BatchNorm-recomputed random network of the same index :: sha256 c6c4138ee5c86423
\newcommand{\MlstLpQgRthirtyResOverInitBnZeroPr}{equal within 20\%}  % experiments/FIGS/data/v2/index.json :: linear probe qg, largest set, zero: R30 of L162 over the BatchNorm-recomputed random network of the same index :: sha256 c6c4138ee5c86423 :: prose twin: two significant figures, direction in words
\newcommand{\MlstLpQgRthirtyMidOverInitBn}{1.357\,\ensuremath{\pm}\,0.033}  % experiments/FIGS/data/v2/index.json :: linear probe qg, largest set, chosen substitution: R30 of R42_Q1 over the BatchNorm-recomputed random network of the same index :: sha256 c6c4138ee5c86423
\newcommand{\MlstLpQgRthirtyMidOverInitBnPr}{36\% more rejection}  % experiments/FIGS/data/v2/index.json :: linear probe qg, largest set, chosen substitution: R30 of R42_Q1 over the BatchNorm-recomputed random network of the same index :: sha256 c6c4138ee5c86423 :: prose twin: two significant figures, direction in words
\newcommand{\MlstLpQgRthirtyMidOverInitBnZero}{1.12\,\ensuremath{\pm}\,0.11}  % experiments/FIGS/data/v2/index.json :: linear probe qg, largest set, zero: R30 of R42_Q1 over the BatchNorm-recomputed random network of the same index :: sha256 c6c4138ee5c86423
\newcommand{\MlstLpQgRthirtyMidOverInitBnZeroPr}{equal within 20\%}  % experiments/FIGS/data/v2/index.json :: linear probe qg, largest set, zero: R30 of R42_Q1 over the BatchNorm-recomputed random network of the same index :: sha256 c6c4138ee5c86423 :: prose twin: two significant figures, direction in words
\newcommand{\MlstLpQgRthirtyCoarseOverInitBn}{1.42\,\ensuremath{\pm}\,0.04}  % experiments/FIGS/data/v2/index.json :: linear probe qg, largest set, chosen substitution: R30 of R16_Q1 over the BatchNorm-recomputed random network of the same index :: sha256 c6c4138ee5c86423
\newcommand{\MlstLpQgRthirtyCoarseOverInitBnPr}{42\% more rejection}  % experiments/FIGS/data/v2/index.json :: linear probe qg, largest set, chosen substitution: R30 of R16_Q1 over the BatchNorm-recomputed random network of the same index :: sha256 c6c4138ee5c86423 :: prose twin: two significant figures, direction in words
\newcommand{\MlstLpQgRthirtyCoarseOverInitBnZero}{1.15\,\ensuremath{\pm}\,0.17}  % experiments/FIGS/data/v2/index.json :: linear probe qg, largest set, zero: R30 of R16_Q1 over the BatchNorm-recomputed random network of the same index :: sha256 c6c4138ee5c86423
\newcommand{\MlstLpQgRthirtyCoarseOverInitBnZeroPr}{equal within 20\%}  % experiments/FIGS/data/v2/index.json :: linear probe qg, largest set, zero: R30 of R16_Q1 over the BatchNorm-recomputed random network of the same index :: sha256 c6c4138ee5c86423 :: prose twin: two significant figures, direction in words
\newcommand{\MlstLpQgRthirtyResMassOverInitBn}{1.430 and 1.364}  % experiments/FIGS/data/v2/index.json :: linear probe qg, largest set, chosen substitution: R30 of L162_MASS over the BatchNorm-recomputed random network of the same index :: sha256 c6c4138ee5c86423
\newcommand{\MlstLpQgRthirtyResMassOverInitBnPr}{at least 36\% more rejection}  % experiments/FIGS/data/v2/index.json :: linear probe qg, largest set, chosen substitution: R30 of L162_MASS over the BatchNorm-recomputed random network of the same index :: sha256 c6c4138ee5c86423 :: prose twin: two significant figures, direction in words
\newcommand{\MlstLpQgRthirtyResMassOverInitBnPrMax}{up to 43\% more rejection}  % experiments/FIGS/data/v2/index.json :: linear probe qg, largest set, chosen substitution: R30 of L162_MASS over the BatchNorm-recomputed random network of the same index :: sha256 c6c4138ee5c86423 :: prose twin of the run farther from 1
\newcommand{\MlstLpQgRthirtyResMassOverInitBnZero}{1.213 and 1.175}  % experiments/FIGS/data/v2/index.json :: linear probe qg, largest set, zero: R30 of L162_MASS over the BatchNorm-recomputed random network of the same index :: sha256 c6c4138ee5c86423
\newcommand{\MlstLpQgRthirtyResMassOverInitBnZeroPr}{equal within 30\%}  % experiments/FIGS/data/v2/index.json :: linear probe qg, largest set, zero: R30 of L162_MASS over the BatchNorm-recomputed random network of the same index :: sha256 c6c4138ee5c86423 :: prose twin: two significant figures, direction in words
\newcommand{\MlstLpQgRthirtyResMassOverInitBnZeroPrMax}{equal within 30\%}  % experiments/FIGS/data/v2/index.json :: linear probe qg, largest set, zero: R30 of L162_MASS over the BatchNorm-recomputed random network of the same index :: sha256 c6c4138ee5c86423 :: prose twin of the run farther from 1
\newcommand{\MlstLpQgRthirtyCoarseMassOverInitBn}{1.324 and 1.394}  % experiments/FIGS/data/v2/index.json :: linear probe qg, largest set, chosen substitution: R30 of R16_Q1_MASS over the BatchNorm-recomputed random network of the same index :: sha256 c6c4138ee5c86423
\newcommand{\MlstLpQgRthirtyCoarseMassOverInitBnPr}{at least 32\% more rejection}  % experiments/FIGS/data/v2/index.json :: linear probe qg, largest set, chosen substitution: R30 of R16_Q1_MASS over the BatchNorm-recomputed random network of the same index :: sha256 c6c4138ee5c86423 :: prose twin: two significant figures, direction in words
\newcommand{\MlstLpQgRthirtyCoarseMassOverInitBnPrMax}{up to 39\% more rejection}  % experiments/FIGS/data/v2/index.json :: linear probe qg, largest set, chosen substitution: R30 of R16_Q1_MASS over the BatchNorm-recomputed random network of the same index :: sha256 c6c4138ee5c86423 :: prose twin of the run farther from 1
\newcommand{\MlstLpQgRthirtyCoarseMassOverInitBnZero}{1.224 and 1.261}  % experiments/FIGS/data/v2/index.json :: linear probe qg, largest set, zero: R30 of R16_Q1_MASS over the BatchNorm-recomputed random network of the same index :: sha256 c6c4138ee5c86423
\newcommand{\MlstLpQgRthirtyCoarseMassOverInitBnZeroPr}{equal within 30\%}  % experiments/FIGS/data/v2/index.json :: linear probe qg, largest set, zero: R30 of R16_Q1_MASS over the BatchNorm-recomputed random network of the same index :: sha256 c6c4138ee5c86423 :: prose twin: two significant figures, direction in words
\newcommand{\MlstLpQgRthirtyCoarseMassOverInitBnZeroPrMax}{equal within 30\%}  % experiments/FIGS/data/v2/index.json :: linear probe qg, largest set, zero: R30 of R16_Q1_MASS over the BatchNorm-recomputed random network of the same index :: sha256 c6c4138ee5c86423 :: prose twin of the run farther from 1
\newcommand{\MlstLpQgRthirtyResOverFull}{0.977\,\ensuremath{\pm}\,0.028}  % experiments/FIGS/data/v2/index.json :: linear probe qg: R30 of L162 over the 188-class run :: sha256 c6c4138ee5c86423
\newcommand{\MlstLpQgRthirtyResOverFullPr}{equal within 3\%}  % experiments/FIGS/data/v2/index.json :: linear probe qg: R30 of L162 over the 188-class run :: sha256 c6c4138ee5c86423 :: prose twin: two significant figures, direction in words
\newcommand{\MlstLpQgRthirtyResOverFullZero}{0.89\,\ensuremath{\pm}\,0.20}  % experiments/FIGS/data/v2/index.json :: linear probe qg: R30 of L162 over the 188-class run :: sha256 c6c4138ee5c86423
\newcommand{\MlstLpQgRthirtyResOverFullZeroPr}{equal within 20\%}  % experiments/FIGS/data/v2/index.json :: linear probe qg: R30 of L162 over the 188-class run :: sha256 c6c4138ee5c86423 :: prose twin: two significant figures, direction in words
\newcommand{\MlstLpQgRthirtyMidOverFull}{0.974\,\ensuremath{\pm}\,0.031}  % experiments/FIGS/data/v2/index.json :: linear probe qg: R30 of R42_Q1 over the 188-class run :: sha256 c6c4138ee5c86423
\newcommand{\MlstLpQgRthirtyMidOverFullPr}{equal within 4\%}  % experiments/FIGS/data/v2/index.json :: linear probe qg: R30 of R42_Q1 over the 188-class run :: sha256 c6c4138ee5c86423 :: prose twin: two significant figures, direction in words
\newcommand{\MlstLpQgRthirtyMidOverFullZero}{0.94\,\ensuremath{\pm}\,0.18}  % experiments/FIGS/data/v2/index.json :: linear probe qg: R30 of R42_Q1 over the 188-class run :: sha256 c6c4138ee5c86423
\newcommand{\MlstLpQgRthirtyMidOverFullZeroPr}{equal within 20\%}  % experiments/FIGS/data/v2/index.json :: linear probe qg: R30 of R42_Q1 over the 188-class run :: sha256 c6c4138ee5c86423 :: prose twin: two significant figures, direction in words
\newcommand{\MlstLpQgRthirtyCoarseOverFull}{1.02\,\ensuremath{\pm}\,0.06}  % experiments/FIGS/data/v2/index.json :: linear probe qg: R30 of R16_Q1 over the 188-class run :: sha256 c6c4138ee5c86423
\newcommand{\MlstLpQgRthirtyCoarseOverFullPr}{equal within 6\%}  % experiments/FIGS/data/v2/index.json :: linear probe qg: R30 of R16_Q1 over the 188-class run :: sha256 c6c4138ee5c86423 :: prose twin: two significant figures, direction in words
\newcommand{\MlstLpQgRthirtyCoarseOverFullZero}{0.95\,\ensuremath{\pm}\,0.09}  % experiments/FIGS/data/v2/index.json :: linear probe qg: R30 of R16_Q1 over the 188-class run :: sha256 c6c4138ee5c86423
\newcommand{\MlstLpQgRthirtyCoarseOverFullZeroPr}{equal within 9\%}  % experiments/FIGS/data/v2/index.json :: linear probe qg: R30 of R16_Q1 over the 188-class run :: sha256 c6c4138ee5c86423 :: prose twin: two significant figures, direction in words
\newcommand{\MlstLpQgRthirtyMeanOverZeroFull}{1.12\,\ensuremath{\pm}\,0.17}  % experiments/FIGS/data/v2/index.json :: linear probe qg: R30 with mean over R30 with zero, L188 :: sha256 c6c4138ee5c86423
\newcommand{\MlstLpQgRthirtyMeanOverZeroFullPr}{equal within 20\%}  % experiments/FIGS/data/v2/index.json :: linear probe qg: R30 with mean over R30 with zero, L188 :: sha256 c6c4138ee5c86423 :: prose twin: two significant figures, direction in words
\newcommand{\MlstLpQgRthirtyDrawOverZeroFull}{1.05\,\ensuremath{\pm}\,0.12}  % experiments/FIGS/data/v2/index.json :: linear probe qg: R30 with draws over R30 with zero, L188 :: sha256 c6c4138ee5c86423
\newcommand{\MlstLpQgRthirtyDrawOverZeroFullPr}{equal within 20\%}  % experiments/FIGS/data/v2/index.json :: linear probe qg: R30 with draws over R30 with zero, L188 :: sha256 c6c4138ee5c86423 :: prose twin: two significant figures, direction in words
\newcommand{\MlstLpQgRthirtyMeanOverZeroInitBn}{0.964\,\ensuremath{\pm}\,0.030}  % experiments/FIGS/data/v2/index.json :: linear probe qg: R30 with mean over R30 with zero, INIT_BN :: sha256 c6c4138ee5c86423
\newcommand{\MlstLpQgRthirtyMeanOverZeroInitBnPr}{equal within 4\%}  % experiments/FIGS/data/v2/index.json :: linear probe qg: R30 with mean over R30 with zero, INIT_BN :: sha256 c6c4138ee5c86423 :: prose twin: two significant figures, direction in words
\newcommand{\MlstLpQgRthirtyDrawOverZeroInitBn}{0.97\,\ensuremath{\pm}\,0.04}  % experiments/FIGS/data/v2/index.json :: linear probe qg: R30 with draws over R30 with zero, INIT_BN :: sha256 c6c4138ee5c86423
\newcommand{\MlstLpQgRthirtyDrawOverZeroInitBnPr}{equal within 4\%}  % experiments/FIGS/data/v2/index.json :: linear probe qg: R30 with draws over R30 with zero, INIT_BN :: sha256 c6c4138ee5c86423 :: prose twin: two significant figures, direction in words
\newcommand{\MlstLpQgRthirtyFullZero}{66\,\ensuremath{\pm}\,7}  % experiments/FIGS/data/v2/index.json :: linear probe qg, largest set: R30 with zero, L188 :: sha256 c6c4138ee5c86423
\newcommand{\MlstLpQgRthirtyFullMean}{72.7\,\ensuremath{\pm}\,2.4}  % experiments/FIGS/data/v2/index.json :: linear probe qg, largest set: R30 with mean, L188 :: sha256 c6c4138ee5c86423
\newcommand{\MlstLpQgRthirtyFullDraws}{68.1\,\ensuremath{\pm}\,1.8}  % experiments/FIGS/data/v2/index.json :: linear probe qg, largest set: R30 with draws, L188 :: sha256 c6c4138ee5c86423
\newcommand{\MlstLpQgRthirtyInitBnZero}{54.2\,\ensuremath{\pm}\,2.4}  % experiments/FIGS/data/v2/index.json :: linear probe qg, largest set: R30 with zero, INIT_BN :: sha256 c6c4138ee5c86423
\newcommand{\MlstLpQgRthirtyInitBnMean}{52.2\,\ensuremath{\pm}\,0.8}  % experiments/FIGS/data/v2/index.json :: linear probe qg, largest set: R30 with mean, INIT_BN :: sha256 c6c4138ee5c86423
\newcommand{\MlstLpQgRthirtyInitBnDraws}{52.4\,\ensuremath{\pm}\,0.4}  % experiments/FIGS/data/v2/index.json :: linear probe qg, largest set: R30 with draws, INIT_BN :: sha256 c6c4138ee5c86423
\newcommand{\MlstLpQgHerwigRthirtyFullOverInitBn}{1.14\,\ensuremath{\pm}\,0.10}  % experiments/FIGS/data/v2/index.json :: linear probe qg_herwig, largest set, chosen substitution: R30 of L188 over the BatchNorm-recomputed random network of the same index :: sha256 c6c4138ee5c86423
\newcommand{\MlstLpQgHerwigRthirtyFullOverInitBnPr}{equal within 20\%}  % experiments/FIGS/data/v2/index.json :: linear probe qg_herwig, largest set, chosen substitution: R30 of L188 over the BatchNorm-recomputed random network of the same index :: sha256 c6c4138ee5c86423 :: prose twin: two significant figures, direction in words
\newcommand{\MlstLpQgHerwigRthirtyFullOverInitBnZero}{1.054\,\ensuremath{\pm}\,0.028}  % experiments/FIGS/data/v2/index.json :: linear probe qg_herwig, largest set, zero: R30 of L188 over the BatchNorm-recomputed random network of the same index :: sha256 c6c4138ee5c86423
\newcommand{\MlstLpQgHerwigRthirtyFullOverInitBnZeroPr}{equal within 6\%}  % experiments/FIGS/data/v2/index.json :: linear probe qg_herwig, largest set, zero: R30 of L188 over the BatchNorm-recomputed random network of the same index :: sha256 c6c4138ee5c86423 :: prose twin: two significant figures, direction in words
\newcommand{\MlstLpQgHerwigRthirtyResOverInitBn}{1.17\,\ensuremath{\pm}\,0.04}  % experiments/FIGS/data/v2/index.json :: linear probe qg_herwig, largest set, chosen substitution: R30 of L162 over the BatchNorm-recomputed random network of the same index :: sha256 c6c4138ee5c86423
\newcommand{\MlstLpQgHerwigRthirtyResOverInitBnPr}{17\% more rejection}  % experiments/FIGS/data/v2/index.json :: linear probe qg_herwig, largest set, chosen substitution: R30 of L162 over the BatchNorm-recomputed random network of the same index :: sha256 c6c4138ee5c86423 :: prose twin: two significant figures, direction in words
\newcommand{\MlstLpQgHerwigRthirtyResOverInitBnZero}{0.98\,\ensuremath{\pm}\,0.10}  % experiments/FIGS/data/v2/index.json :: linear probe qg_herwig, largest set, zero: R30 of L162 over the BatchNorm-recomputed random network of the same index :: sha256 c6c4138ee5c86423
\newcommand{\MlstLpQgHerwigRthirtyResOverInitBnZeroPr}{equal within 10\%}  % experiments/FIGS/data/v2/index.json :: linear probe qg_herwig, largest set, zero: R30 of L162 over the BatchNorm-recomputed random network of the same index :: sha256 c6c4138ee5c86423 :: prose twin: two significant figures, direction in words
\newcommand{\MlstLpQgHerwigRthirtyMidOverInitBn}{1.12\,\ensuremath{\pm}\,0.07}  % experiments/FIGS/data/v2/index.json :: linear probe qg_herwig, largest set, chosen substitution: R30 of R42_Q1 over the BatchNorm-recomputed random network of the same index :: sha256 c6c4138ee5c86423
\newcommand{\MlstLpQgHerwigRthirtyMidOverInitBnPr}{equal within 20\%}  % experiments/FIGS/data/v2/index.json :: linear probe qg_herwig, largest set, chosen substitution: R30 of R42_Q1 over the BatchNorm-recomputed random network of the same index :: sha256 c6c4138ee5c86423 :: prose twin: two significant figures, direction in words
\newcommand{\MlstLpQgHerwigRthirtyMidOverInitBnZero}{0.99\,\ensuremath{\pm}\,0.07}  % experiments/FIGS/data/v2/index.json :: linear probe qg_herwig, largest set, zero: R30 of R42_Q1 over the BatchNorm-recomputed random network of the same index :: sha256 c6c4138ee5c86423
\newcommand{\MlstLpQgHerwigRthirtyMidOverInitBnZeroPr}{equal within 7\%}  % experiments/FIGS/data/v2/index.json :: linear probe qg_herwig, largest set, zero: R30 of R42_Q1 over the BatchNorm-recomputed random network of the same index :: sha256 c6c4138ee5c86423 :: prose twin: two significant figures, direction in words
\newcommand{\MlstLpQgHerwigRthirtyCoarseOverInitBn}{1.154\,\ensuremath{\pm}\,0.015}  % experiments/FIGS/data/v2/index.json :: linear probe qg_herwig, largest set, chosen substitution: R30 of R16_Q1 over the BatchNorm-recomputed random network of the same index :: sha256 c6c4138ee5c86423
\newcommand{\MlstLpQgHerwigRthirtyCoarseOverInitBnPr}{15\% more rejection}  % experiments/FIGS/data/v2/index.json :: linear probe qg_herwig, largest set, chosen substitution: R30 of R16_Q1 over the BatchNorm-recomputed random network of the same index :: sha256 c6c4138ee5c86423 :: prose twin: two significant figures, direction in words
\newcommand{\MlstLpQgHerwigRthirtyCoarseOverInitBnZero}{1.07\,\ensuremath{\pm}\,0.05}  % experiments/FIGS/data/v2/index.json :: linear probe qg_herwig, largest set, zero: R30 of R16_Q1 over the BatchNorm-recomputed random network of the same index :: sha256 c6c4138ee5c86423
\newcommand{\MlstLpQgHerwigRthirtyCoarseOverInitBnZeroPr}{equal within 8\%}  % experiments/FIGS/data/v2/index.json :: linear probe qg_herwig, largest set, zero: R30 of R16_Q1 over the BatchNorm-recomputed random network of the same index :: sha256 c6c4138ee5c86423 :: prose twin: two significant figures, direction in words
\newcommand{\MlstLpQgHerwigRthirtyResMassOverInitBn}{1.152 and 1.138}  % experiments/FIGS/data/v2/index.json :: linear probe qg_herwig, largest set, chosen substitution: R30 of L162_MASS over the BatchNorm-recomputed random network of the same index :: sha256 c6c4138ee5c86423
\newcommand{\MlstLpQgHerwigRthirtyResMassOverInitBnPr}{equal within 20\%}  % experiments/FIGS/data/v2/index.json :: linear probe qg_herwig, largest set, chosen substitution: R30 of L162_MASS over the BatchNorm-recomputed random network of the same index :: sha256 c6c4138ee5c86423 :: prose twin: two significant figures, direction in words
\newcommand{\MlstLpQgHerwigRthirtyResMassOverInitBnPrMax}{equal within 20\%}  % experiments/FIGS/data/v2/index.json :: linear probe qg_herwig, largest set, chosen substitution: R30 of L162_MASS over the BatchNorm-recomputed random network of the same index :: sha256 c6c4138ee5c86423 :: prose twin of the run farther from 1
\newcommand{\MlstLpQgHerwigRthirtyResMassOverInitBnZero}{1.13 and 1.03}  % experiments/FIGS/data/v2/index.json :: linear probe qg_herwig, largest set, zero: R30 of L162_MASS over the BatchNorm-recomputed random network of the same index :: sha256 c6c4138ee5c86423
\newcommand{\MlstLpQgHerwigRthirtyResMassOverInitBnZeroPr}{equal within 20\%}  % experiments/FIGS/data/v2/index.json :: linear probe qg_herwig, largest set, zero: R30 of L162_MASS over the BatchNorm-recomputed random network of the same index :: sha256 c6c4138ee5c86423 :: prose twin: two significant figures, direction in words
\newcommand{\MlstLpQgHerwigRthirtyResMassOverInitBnZeroPrMax}{equal within 20\%}  % experiments/FIGS/data/v2/index.json :: linear probe qg_herwig, largest set, zero: R30 of L162_MASS over the BatchNorm-recomputed random network of the same index :: sha256 c6c4138ee5c86423 :: prose twin of the run farther from 1
\newcommand{\MlstLpQgHerwigRthirtyCoarseMassOverInitBn}{1.146 and 1.159}  % experiments/FIGS/data/v2/index.json :: linear probe qg_herwig, largest set, chosen substitution: R30 of R16_Q1_MASS over the BatchNorm-recomputed random network of the same index :: sha256 c6c4138ee5c86423
\newcommand{\MlstLpQgHerwigRthirtyCoarseMassOverInitBnPr}{equal within 20\%}  % experiments/FIGS/data/v2/index.json :: linear probe qg_herwig, largest set, chosen substitution: R30 of R16_Q1_MASS over the BatchNorm-recomputed random network of the same index :: sha256 c6c4138ee5c86423 :: prose twin: two significant figures, direction in words
\newcommand{\MlstLpQgHerwigRthirtyCoarseMassOverInitBnPrMax}{equal within 20\%}  % experiments/FIGS/data/v2/index.json :: linear probe qg_herwig, largest set, chosen substitution: R30 of R16_Q1_MASS over the BatchNorm-recomputed random network of the same index :: sha256 c6c4138ee5c86423 :: prose twin of the run farther from 1
\newcommand{\MlstLpQgHerwigRthirtyCoarseMassOverInitBnZero}{1.070 and 1.100}  % experiments/FIGS/data/v2/index.json :: linear probe qg_herwig, largest set, zero: R30 of R16_Q1_MASS over the BatchNorm-recomputed random network of the same index :: sha256 c6c4138ee5c86423
\newcommand{\MlstLpQgHerwigRthirtyCoarseMassOverInitBnZeroPr}{equal within 20\%}  % experiments/FIGS/data/v2/index.json :: linear probe qg_herwig, largest set, zero: R30 of R16_Q1_MASS over the BatchNorm-recomputed random network of the same index :: sha256 c6c4138ee5c86423 :: prose twin: two significant figures, direction in words
\newcommand{\MlstLpQgHerwigRthirtyCoarseMassOverInitBnZeroPrMax}{equal within 20\%}  % experiments/FIGS/data/v2/index.json :: linear probe qg_herwig, largest set, zero: R30 of R16_Q1_MASS over the BatchNorm-recomputed random network of the same index :: sha256 c6c4138ee5c86423 :: prose twin of the run farther from 1
\newcommand{\MlstLpQgHerwigRthirtyResOverFull}{1.03\,\ensuremath{\pm}\,0.11}  % experiments/FIGS/data/v2/index.json :: linear probe qg_herwig: R30 of L162 over the 188-class run :: sha256 c6c4138ee5c86423
\newcommand{\MlstLpQgHerwigRthirtyResOverFullPr}{equal within 20\%}  % experiments/FIGS/data/v2/index.json :: linear probe qg_herwig: R30 of L162 over the 188-class run :: sha256 c6c4138ee5c86423 :: prose twin: two significant figures, direction in words
\newcommand{\MlstLpQgHerwigRthirtyResOverFullZero}{0.93\,\ensuremath{\pm}\,0.11}  % experiments/FIGS/data/v2/index.json :: linear probe qg_herwig: R30 of L162 over the 188-class run :: sha256 c6c4138ee5c86423
\newcommand{\MlstLpQgHerwigRthirtyResOverFullZeroPr}{equal within 20\%}  % experiments/FIGS/data/v2/index.json :: linear probe qg_herwig: R30 of L162 over the 188-class run :: sha256 c6c4138ee5c86423 :: prose twin: two significant figures, direction in words
\newcommand{\MlstLpQgHerwigRthirtyMidOverFull}{0.98\,\ensuremath{\pm}\,0.06}  % experiments/FIGS/data/v2/index.json :: linear probe qg_herwig: R30 of R42_Q1 over the 188-class run :: sha256 c6c4138ee5c86423
\newcommand{\MlstLpQgHerwigRthirtyMidOverFullPr}{equal within 6\%}  % experiments/FIGS/data/v2/index.json :: linear probe qg_herwig: R30 of R42_Q1 over the 188-class run :: sha256 c6c4138ee5c86423 :: prose twin: two significant figures, direction in words
\newcommand{\MlstLpQgHerwigRthirtyMidOverFullZero}{0.94\,\ensuremath{\pm}\,0.09}  % experiments/FIGS/data/v2/index.json :: linear probe qg_herwig: R30 of R42_Q1 over the 188-class run :: sha256 c6c4138ee5c86423
\newcommand{\MlstLpQgHerwigRthirtyMidOverFullZeroPr}{equal within 9\%}  % experiments/FIGS/data/v2/index.json :: linear probe qg_herwig: R30 of R42_Q1 over the 188-class run :: sha256 c6c4138ee5c86423 :: prose twin: two significant figures, direction in words
\newcommand{\MlstLpQgHerwigRthirtyCoarseOverFull}{1.02\,\ensuremath{\pm}\,0.09}  % experiments/FIGS/data/v2/index.json :: linear probe qg_herwig: R30 of R16_Q1 over the 188-class run :: sha256 c6c4138ee5c86423
\newcommand{\MlstLpQgHerwigRthirtyCoarseOverFullPr}{equal within 10\%}  % experiments/FIGS/data/v2/index.json :: linear probe qg_herwig: R30 of R16_Q1 over the 188-class run :: sha256 c6c4138ee5c86423 :: prose twin: two significant figures, direction in words
\newcommand{\MlstLpQgHerwigRthirtyCoarseOverFullZero}{1.018\,\ensuremath{\pm}\,0.024}  % experiments/FIGS/data/v2/index.json :: linear probe qg_herwig: R30 of R16_Q1 over the 188-class run :: sha256 c6c4138ee5c86423
\newcommand{\MlstLpQgHerwigRthirtyCoarseOverFullZeroPr}{equal within 3\%}  % experiments/FIGS/data/v2/index.json :: linear probe qg_herwig: R30 of R16_Q1 over the 188-class run :: sha256 c6c4138ee5c86423 :: prose twin: two significant figures, direction in words
\newcommand{\MlstLpQgHerwigRthirtyMeanOverZeroFull}{1.07\,\ensuremath{\pm}\,0.11}  % experiments/FIGS/data/v2/index.json :: linear probe qg_herwig: R30 with mean over R30 with zero, L188 :: sha256 c6c4138ee5c86423
\newcommand{\MlstLpQgHerwigRthirtyMeanOverZeroFullPr}{equal within 20\%}  % experiments/FIGS/data/v2/index.json :: linear probe qg_herwig: R30 with mean over R30 with zero, L188 :: sha256 c6c4138ee5c86423 :: prose twin: two significant figures, direction in words
\newcommand{\MlstLpQgHerwigRthirtyDrawOverZeroFull}{1.07\,\ensuremath{\pm}\,0.06}  % experiments/FIGS/data/v2/index.json :: linear probe qg_herwig: R30 with draws over R30 with zero, L188 :: sha256 c6c4138ee5c86423
\newcommand{\MlstLpQgHerwigRthirtyDrawOverZeroFullPr}{equal within 7\%}  % experiments/FIGS/data/v2/index.json :: linear probe qg_herwig: R30 with draws over R30 with zero, L188 :: sha256 c6c4138ee5c86423 :: prose twin: two significant figures, direction in words
\newcommand{\MlstLpQgHerwigRthirtyMeanOverZeroInitBn}{0.988\,\ensuremath{\pm}\,0.019}  % experiments/FIGS/data/v2/index.json :: linear probe qg_herwig: R30 with mean over R30 with zero, INIT_BN :: sha256 c6c4138ee5c86423
\newcommand{\MlstLpQgHerwigRthirtyMeanOverZeroInitBnPr}{equal within 2\%}  % experiments/FIGS/data/v2/index.json :: linear probe qg_herwig: R30 with mean over R30 with zero, INIT_BN :: sha256 c6c4138ee5c86423 :: prose twin: two significant figures, direction in words
\newcommand{\MlstLpQgHerwigRthirtyDrawOverZeroInitBn}{0.983\,\ensuremath{\pm}\,0.032}  % experiments/FIGS/data/v2/index.json :: linear probe qg_herwig: R30 with draws over R30 with zero, INIT_BN :: sha256 c6c4138ee5c86423
\newcommand{\MlstLpQgHerwigRthirtyDrawOverZeroInitBnPr}{equal within 4\%}  % experiments/FIGS/data/v2/index.json :: linear probe qg_herwig: R30 with draws over R30 with zero, INIT_BN :: sha256 c6c4138ee5c86423 :: prose twin: two significant figures, direction in words
\newcommand{\MlstLpQgHerwigRthirtyFullZero}{32.8\,\ensuremath{\pm}\,1.0}  % experiments/FIGS/data/v2/index.json :: linear probe qg_herwig, largest set: R30 with zero, L188 :: sha256 c6c4138ee5c86423
\newcommand{\MlstLpQgHerwigRthirtyFullMean}{35.0\,\ensuremath{\pm}\,2.6}  % experiments/FIGS/data/v2/index.json :: linear probe qg_herwig, largest set: R30 with mean, L188 :: sha256 c6c4138ee5c86423
\newcommand{\MlstLpQgHerwigRthirtyFullDraws}{35.0\,\ensuremath{\pm}\,0.8}  % experiments/FIGS/data/v2/index.json :: linear probe qg_herwig, largest set: R30 with draws, L188 :: sha256 c6c4138ee5c86423
\newcommand{\MlstLpQgHerwigRthirtyInitBnZero}{31.1\,\ensuremath{\pm}\,0.5}  % experiments/FIGS/data/v2/index.json :: linear probe qg_herwig, largest set: R30 with zero, INIT_BN :: sha256 c6c4138ee5c86423
\newcommand{\MlstLpQgHerwigRthirtyInitBnMean}{30.7\,\ensuremath{\pm}\,0.5}  % experiments/FIGS/data/v2/index.json :: linear probe qg_herwig, largest set: R30 with mean, INIT_BN :: sha256 c6c4138ee5c86423
\newcommand{\MlstLpQgHerwigRthirtyInitBnDraws}{30.6\,\ensuremath{\pm}\,0.6}  % experiments/FIGS/data/v2/index.json :: linear probe qg_herwig, largest set: R30 with draws, INIT_BN :: sha256 c6c4138ee5c86423
\newcommand{\MlstLpJcTwoTopZeroAccDropFull}{88.9\,\ensuremath{\pm}\,1.0\%}  % experiments/FIGS/data/v2/index.json :: linear probe JetClass-II, largest set: relative drop in accuracy (%) with the top data's missing slots set to zero, L188 :: sha256 c6c4138ee5c86423
\newcommand{\MlstLpJcTwoTopZeroAccDropFullPr}{89\%}  % experiments/FIGS/data/v2/index.json :: linear probe JetClass-II, largest set: relative drop in accuracy (%) with the top data's missing slots set to zero, L188 :: sha256 c6c4138ee5c86423 :: prose twin, two significant figures
\newcommand{\MlstLpJcTwoTopMeanAccDropFull}{87.9\,\ensuremath{\pm}\,1.1\%}  % experiments/FIGS/data/v2/index.json :: linear probe JetClass-II, largest set: relative drop in accuracy (%) with the top data's missing slots set to mean, L188 :: sha256 c6c4138ee5c86423
\newcommand{\MlstLpJcTwoTopMeanAccDropFullPr}{88\%}  % experiments/FIGS/data/v2/index.json :: linear probe JetClass-II, largest set: relative drop in accuracy (%) with the top data's missing slots set to mean, L188 :: sha256 c6c4138ee5c86423 :: prose twin, two significant figures
\newcommand{\MlstLpJcTwoTopDrawAccDropFull}{85.98\,\ensuremath{\pm}\,0.07\%}  % experiments/FIGS/data/v2/index.json :: linear probe JetClass-II, largest set: relative drop in accuracy (%) with the top data's missing slots set to draws, L188 :: sha256 c6c4138ee5c86423
\newcommand{\MlstLpJcTwoTopDrawAccDropFullPr}{86\%}  % experiments/FIGS/data/v2/index.json :: linear probe JetClass-II, largest set: relative drop in accuracy (%) with the top data's missing slots set to draws, L188 :: sha256 c6c4138ee5c86423 :: prose twin, two significant figures
\newcommand{\MlstLpJcTwoTopZeroAccDropRes}{89.0\,\ensuremath{\pm}\,0.9\%}  % experiments/FIGS/data/v2/index.json :: linear probe JetClass-II, largest set: relative drop in accuracy (%) with the top data's missing slots set to zero, L162 :: sha256 c6c4138ee5c86423
\newcommand{\MlstLpJcTwoTopZeroAccDropResPr}{89\%}  % experiments/FIGS/data/v2/index.json :: linear probe JetClass-II, largest set: relative drop in accuracy (%) with the top data's missing slots set to zero, L162 :: sha256 c6c4138ee5c86423 :: prose twin, two significant figures
\newcommand{\MlstLpJcTwoTopMeanAccDropRes}{87.8\,\ensuremath{\pm}\,0.7\%}  % experiments/FIGS/data/v2/index.json :: linear probe JetClass-II, largest set: relative drop in accuracy (%) with the top data's missing slots set to mean, L162 :: sha256 c6c4138ee5c86423
\newcommand{\MlstLpJcTwoTopMeanAccDropResPr}{88\%}  % experiments/FIGS/data/v2/index.json :: linear probe JetClass-II, largest set: relative drop in accuracy (%) with the top data's missing slots set to mean, L162 :: sha256 c6c4138ee5c86423 :: prose twin, two significant figures
\newcommand{\MlstLpJcTwoTopDrawAccDropRes}{85.89\,\ensuremath{\pm}\,0.11\%}  % experiments/FIGS/data/v2/index.json :: linear probe JetClass-II, largest set: relative drop in accuracy (%) with the top data's missing slots set to draws, L162 :: sha256 c6c4138ee5c86423
\newcommand{\MlstLpJcTwoTopDrawAccDropResPr}{86\%}  % experiments/FIGS/data/v2/index.json :: linear probe JetClass-II, largest set: relative drop in accuracy (%) with the top data's missing slots set to draws, L162 :: sha256 c6c4138ee5c86423 :: prose twin, two significant figures
\newcommand{\MlstLpJcTwoTopZeroAccDropMid}{87.81\,\ensuremath{\pm}\,0.28\%}  % experiments/FIGS/data/v2/index.json :: linear probe JetClass-II, largest set: relative drop in accuracy (%) with the top data's missing slots set to zero, R42_Q1 :: sha256 c6c4138ee5c86423
\newcommand{\MlstLpJcTwoTopZeroAccDropMidPr}{88\%}  % experiments/FIGS/data/v2/index.json :: linear probe JetClass-II, largest set: relative drop in accuracy (%) with the top data's missing slots set to zero, R42_Q1 :: sha256 c6c4138ee5c86423 :: prose twin, two significant figures
\newcommand{\MlstLpJcTwoTopMeanAccDropMid}{86.7\,\ensuremath{\pm}\,0.9\%}  % experiments/FIGS/data/v2/index.json :: linear probe JetClass-II, largest set: relative drop in accuracy (%) with the top data's missing slots set to mean, R42_Q1 :: sha256 c6c4138ee5c86423
\newcommand{\MlstLpJcTwoTopMeanAccDropMidPr}{87\%}  % experiments/FIGS/data/v2/index.json :: linear probe JetClass-II, largest set: relative drop in accuracy (%) with the top data's missing slots set to mean, R42_Q1 :: sha256 c6c4138ee5c86423 :: prose twin, two significant figures
\newcommand{\MlstLpJcTwoTopDrawAccDropMid}{83.17\,\ensuremath{\pm}\,0.29\%}  % experiments/FIGS/data/v2/index.json :: linear probe JetClass-II, largest set: relative drop in accuracy (%) with the top data's missing slots set to draws, R42_Q1 :: sha256 c6c4138ee5c86423
\newcommand{\MlstLpJcTwoTopDrawAccDropMidPr}{83\%}  % experiments/FIGS/data/v2/index.json :: linear probe JetClass-II, largest set: relative drop in accuracy (%) with the top data's missing slots set to draws, R42_Q1 :: sha256 c6c4138ee5c86423 :: prose twin, two significant figures
\newcommand{\MlstLpJcTwoTopZeroAccDropCoarse}{86.7\,\ensuremath{\pm}\,0.7\%}  % experiments/FIGS/data/v2/index.json :: linear probe JetClass-II, largest set: relative drop in accuracy (%) with the top data's missing slots set to zero, R16_Q1 :: sha256 c6c4138ee5c86423
\newcommand{\MlstLpJcTwoTopZeroAccDropCoarsePr}{87\%}  % experiments/FIGS/data/v2/index.json :: linear probe JetClass-II, largest set: relative drop in accuracy (%) with the top data's missing slots set to zero, R16_Q1 :: sha256 c6c4138ee5c86423 :: prose twin, two significant figures
\newcommand{\MlstLpJcTwoTopMeanAccDropCoarse}{86.1\,\ensuremath{\pm}\,0.8\%}  % experiments/FIGS/data/v2/index.json :: linear probe JetClass-II, largest set: relative drop in accuracy (%) with the top data's missing slots set to mean, R16_Q1 :: sha256 c6c4138ee5c86423
\newcommand{\MlstLpJcTwoTopMeanAccDropCoarsePr}{86\%}  % experiments/FIGS/data/v2/index.json :: linear probe JetClass-II, largest set: relative drop in accuracy (%) with the top data's missing slots set to mean, R16_Q1 :: sha256 c6c4138ee5c86423 :: prose twin, two significant figures
\newcommand{\MlstLpJcTwoTopDrawAccDropCoarse}{81.4\,\ensuremath{\pm}\,0.5\%}  % experiments/FIGS/data/v2/index.json :: linear probe JetClass-II, largest set: relative drop in accuracy (%) with the top data's missing slots set to draws, R16_Q1 :: sha256 c6c4138ee5c86423
\newcommand{\MlstLpJcTwoTopDrawAccDropCoarsePr}{81\%}  % experiments/FIGS/data/v2/index.json :: linear probe JetClass-II, largest set: relative drop in accuracy (%) with the top data's missing slots set to draws, R16_Q1 :: sha256 c6c4138ee5c86423 :: prose twin, two significant figures
\newcommand{\MlstLpJcTwoTopZeroAccDropResMass}{89.2, 88.3\%}  % experiments/FIGS/data/v2/index.json :: linear probe JetClass-II, largest set: relative drop in accuracy (%) with the top data's missing slots set to zero, L162_MASS :: sha256 c6c4138ee5c86423
\newcommand{\MlstLpJcTwoTopZeroAccDropResMassPr}{88--89\%}  % experiments/FIGS/data/v2/index.json :: linear probe JetClass-II, largest set: relative drop in accuracy (%) with the top data's missing slots set to zero, L162_MASS :: sha256 c6c4138ee5c86423 :: prose twin, two significant figures
\newcommand{\MlstLpJcTwoTopMeanAccDropResMass}{88.5, 87.8\%}  % experiments/FIGS/data/v2/index.json :: linear probe JetClass-II, largest set: relative drop in accuracy (%) with the top data's missing slots set to mean, L162_MASS :: sha256 c6c4138ee5c86423
\newcommand{\MlstLpJcTwoTopMeanAccDropResMassPr}{88--88\%}  % experiments/FIGS/data/v2/index.json :: linear probe JetClass-II, largest set: relative drop in accuracy (%) with the top data's missing slots set to mean, L162_MASS :: sha256 c6c4138ee5c86423 :: prose twin, two significant figures
\newcommand{\MlstLpJcTwoTopDrawAccDropResMass}{85.4, 85.4\%}  % experiments/FIGS/data/v2/index.json :: linear probe JetClass-II, largest set: relative drop in accuracy (%) with the top data's missing slots set to draws, L162_MASS :: sha256 c6c4138ee5c86423
\newcommand{\MlstLpJcTwoTopDrawAccDropResMassPr}{85--85\%}  % experiments/FIGS/data/v2/index.json :: linear probe JetClass-II, largest set: relative drop in accuracy (%) with the top data's missing slots set to draws, L162_MASS :: sha256 c6c4138ee5c86423 :: prose twin, two significant figures
\newcommand{\MlstLpJcTwoTopZeroAccDropCoarseMass}{84.8, 87.4\%}  % experiments/FIGS/data/v2/index.json :: linear probe JetClass-II, largest set: relative drop in accuracy (%) with the top data's missing slots set to zero, R16_Q1_MASS :: sha256 c6c4138ee5c86423
\newcommand{\MlstLpJcTwoTopZeroAccDropCoarseMassPr}{85--87\%}  % experiments/FIGS/data/v2/index.json :: linear probe JetClass-II, largest set: relative drop in accuracy (%) with the top data's missing slots set to zero, R16_Q1_MASS :: sha256 c6c4138ee5c86423 :: prose twin, two significant figures
\newcommand{\MlstLpJcTwoTopMeanAccDropCoarseMass}{85.9, 86.9\%}  % experiments/FIGS/data/v2/index.json :: linear probe JetClass-II, largest set: relative drop in accuracy (%) with the top data's missing slots set to mean, R16_Q1_MASS :: sha256 c6c4138ee5c86423
\newcommand{\MlstLpJcTwoTopMeanAccDropCoarseMassPr}{86--87\%}  % experiments/FIGS/data/v2/index.json :: linear probe JetClass-II, largest set: relative drop in accuracy (%) with the top data's missing slots set to mean, R16_Q1_MASS :: sha256 c6c4138ee5c86423 :: prose twin, two significant figures
\newcommand{\MlstLpJcTwoTopDrawAccDropCoarseMass}{79.1, 81.7\%}  % experiments/FIGS/data/v2/index.json :: linear probe JetClass-II, largest set: relative drop in accuracy (%) with the top data's missing slots set to draws, R16_Q1_MASS :: sha256 c6c4138ee5c86423
\newcommand{\MlstLpJcTwoTopDrawAccDropCoarseMassPr}{79--82\%}  % experiments/FIGS/data/v2/index.json :: linear probe JetClass-II, largest set: relative drop in accuracy (%) with the top data's missing slots set to draws, R16_Q1_MASS :: sha256 c6c4138ee5c86423 :: prose twin, two significant figures
\newcommand{\MlstLpJcTwoTopZeroAccDropFzero}{88.9, 88.5\%}  % experiments/FIGS/data/v2/index.json :: linear probe JetClass-II, largest set: relative drop in accuracy (%) with the top data's missing slots set to zero, FLAV_F0 :: sha256 c6c4138ee5c86423
\newcommand{\MlstLpJcTwoTopZeroAccDropFzeroPr}{88--89\%}  % experiments/FIGS/data/v2/index.json :: linear probe JetClass-II, largest set: relative drop in accuracy (%) with the top data's missing slots set to zero, FLAV_F0 :: sha256 c6c4138ee5c86423 :: prose twin, two significant figures
\newcommand{\MlstLpJcTwoTopMeanAccDropFzero}{88.3, 87.0\%}  % experiments/FIGS/data/v2/index.json :: linear probe JetClass-II, largest set: relative drop in accuracy (%) with the top data's missing slots set to mean, FLAV_F0 :: sha256 c6c4138ee5c86423
\newcommand{\MlstLpJcTwoTopMeanAccDropFzeroPr}{87--88\%}  % experiments/FIGS/data/v2/index.json :: linear probe JetClass-II, largest set: relative drop in accuracy (%) with the top data's missing slots set to mean, FLAV_F0 :: sha256 c6c4138ee5c86423 :: prose twin, two significant figures
\newcommand{\MlstLpJcTwoTopDrawAccDropFzero}{83.3, 83.9\%}  % experiments/FIGS/data/v2/index.json :: linear probe JetClass-II, largest set: relative drop in accuracy (%) with the top data's missing slots set to draws, FLAV_F0 :: sha256 c6c4138ee5c86423
\newcommand{\MlstLpJcTwoTopDrawAccDropFzeroPr}{83--84\%}  % experiments/FIGS/data/v2/index.json :: linear probe JetClass-II, largest set: relative drop in accuracy (%) with the top data's missing slots set to draws, FLAV_F0 :: sha256 c6c4138ee5c86423 :: prose twin, two significant figures
\newcommand{\MlstLpJcTwoTopZeroAccDropFone}{87.8, 87.6\%}  % experiments/FIGS/data/v2/index.json :: linear probe JetClass-II, largest set: relative drop in accuracy (%) with the top data's missing slots set to zero, FLAV_F1 :: sha256 c6c4138ee5c86423
\newcommand{\MlstLpJcTwoTopZeroAccDropFonePr}{88--88\%}  % experiments/FIGS/data/v2/index.json :: linear probe JetClass-II, largest set: relative drop in accuracy (%) with the top data's missing slots set to zero, FLAV_F1 :: sha256 c6c4138ee5c86423 :: prose twin, two significant figures
\newcommand{\MlstLpJcTwoTopMeanAccDropFone}{88.9, 86.9\%}  % experiments/FIGS/data/v2/index.json :: linear probe JetClass-II, largest set: relative drop in accuracy (%) with the top data's missing slots set to mean, FLAV_F1 :: sha256 c6c4138ee5c86423
\newcommand{\MlstLpJcTwoTopMeanAccDropFonePr}{87--89\%}  % experiments/FIGS/data/v2/index.json :: linear probe JetClass-II, largest set: relative drop in accuracy (%) with the top data's missing slots set to mean, FLAV_F1 :: sha256 c6c4138ee5c86423 :: prose twin, two significant figures
\newcommand{\MlstLpJcTwoTopDrawAccDropFone}{84.4, 85.2\%}  % experiments/FIGS/data/v2/index.json :: linear probe JetClass-II, largest set: relative drop in accuracy (%) with the top data's missing slots set to draws, FLAV_F1 :: sha256 c6c4138ee5c86423
\newcommand{\MlstLpJcTwoTopDrawAccDropFonePr}{84--85\%}  % experiments/FIGS/data/v2/index.json :: linear probe JetClass-II, largest set: relative drop in accuracy (%) with the top data's missing slots set to draws, FLAV_F1 :: sha256 c6c4138ee5c86423 :: prose twin, two significant figures
\newcommand{\MlstLpJcTwoTopZeroAccDropFoner}{89.6, 90.3\%}  % experiments/FIGS/data/v2/index.json :: linear probe JetClass-II, largest set: relative drop in accuracy (%) with the top data's missing slots set to zero, FLAV_F1R :: sha256 c6c4138ee5c86423
\newcommand{\MlstLpJcTwoTopZeroAccDropFonerPr}{90--90\%}  % experiments/FIGS/data/v2/index.json :: linear probe JetClass-II, largest set: relative drop in accuracy (%) with the top data's missing slots set to zero, FLAV_F1R :: sha256 c6c4138ee5c86423 :: prose twin, two significant figures
\newcommand{\MlstLpJcTwoTopMeanAccDropFoner}{89.4, 89.7\%}  % experiments/FIGS/data/v2/index.json :: linear probe JetClass-II, largest set: relative drop in accuracy (%) with the top data's missing slots set to mean, FLAV_F1R :: sha256 c6c4138ee5c86423
\newcommand{\MlstLpJcTwoTopMeanAccDropFonerPr}{89--90\%}  % experiments/FIGS/data/v2/index.json :: linear probe JetClass-II, largest set: relative drop in accuracy (%) with the top data's missing slots set to mean, FLAV_F1R :: sha256 c6c4138ee5c86423 :: prose twin, two significant figures
\newcommand{\MlstLpJcTwoTopDrawAccDropFoner}{84.9, 85.6\%}  % experiments/FIGS/data/v2/index.json :: linear probe JetClass-II, largest set: relative drop in accuracy (%) with the top data's missing slots set to draws, FLAV_F1R :: sha256 c6c4138ee5c86423
\newcommand{\MlstLpJcTwoTopDrawAccDropFonerPr}{85--86\%}  % experiments/FIGS/data/v2/index.json :: linear probe JetClass-II, largest set: relative drop in accuracy (%) with the top data's missing slots set to draws, FLAV_F1R :: sha256 c6c4138ee5c86423 :: prose twin, two significant figures
\newcommand{\MlstLpJcTwoTopZeroAccDropInitBn}{56.4\,\ensuremath{\pm}\,1.4\%}  % experiments/FIGS/data/v2/index.json :: linear probe JetClass-II, largest set: relative drop in accuracy (%) with the top data's missing slots set to zero, INIT_BN :: sha256 c6c4138ee5c86423
\newcommand{\MlstLpJcTwoTopZeroAccDropInitBnPr}{56\%}  % experiments/FIGS/data/v2/index.json :: linear probe JetClass-II, largest set: relative drop in accuracy (%) with the top data's missing slots set to zero, INIT_BN :: sha256 c6c4138ee5c86423 :: prose twin, two significant figures
\newcommand{\MlstLpJcTwoTopMeanAccDropInitBn}{59.8\,\ensuremath{\pm}\,0.8\%}  % experiments/FIGS/data/v2/index.json :: linear probe JetClass-II, largest set: relative drop in accuracy (%) with the top data's missing slots set to mean, INIT_BN :: sha256 c6c4138ee5c86423
\newcommand{\MlstLpJcTwoTopMeanAccDropInitBnPr}{60\%}  % experiments/FIGS/data/v2/index.json :: linear probe JetClass-II, largest set: relative drop in accuracy (%) with the top data's missing slots set to mean, INIT_BN :: sha256 c6c4138ee5c86423 :: prose twin, two significant figures
\newcommand{\MlstLpJcTwoTopDrawAccDropInitBn}{59.8\,\ensuremath{\pm}\,1.2\%}  % experiments/FIGS/data/v2/index.json :: linear probe JetClass-II, largest set: relative drop in accuracy (%) with the top data's missing slots set to draws, INIT_BN :: sha256 c6c4138ee5c86423
\newcommand{\MlstLpJcTwoTopDrawAccDropInitBnPr}{60\%}  % experiments/FIGS/data/v2/index.json :: linear probe JetClass-II, largest set: relative drop in accuracy (%) with the top data's missing slots set to draws, INIT_BN :: sha256 c6c4138ee5c86423 :: prose twin, two significant figures
\newcommand{\MlstLpJcTwoQgZeroAccDropFull}{61.2\,\ensuremath{\pm}\,3.0\%}  % experiments/FIGS/data/v2/index.json :: linear probe JetClass-II, largest set: relative drop in accuracy (%) with the qg data's missing slots set to zero, L188 :: sha256 c6c4138ee5c86423
\newcommand{\MlstLpJcTwoQgZeroAccDropFullPr}{61\%}  % experiments/FIGS/data/v2/index.json :: linear probe JetClass-II, largest set: relative drop in accuracy (%) with the qg data's missing slots set to zero, L188 :: sha256 c6c4138ee5c86423 :: prose twin, two significant figures
\newcommand{\MlstLpJcTwoQgMeanAccDropFull}{48.8\,\ensuremath{\pm}\,1.1\%}  % experiments/FIGS/data/v2/index.json :: linear probe JetClass-II, largest set: relative drop in accuracy (%) with the qg data's missing slots set to mean, L188 :: sha256 c6c4138ee5c86423
\newcommand{\MlstLpJcTwoQgMeanAccDropFullPr}{49\%}  % experiments/FIGS/data/v2/index.json :: linear probe JetClass-II, largest set: relative drop in accuracy (%) with the qg data's missing slots set to mean, L188 :: sha256 c6c4138ee5c86423 :: prose twin, two significant figures
\newcommand{\MlstLpJcTwoQgDrawAccDropFull}{54.30\,\ensuremath{\pm}\,0.21\%}  % experiments/FIGS/data/v2/index.json :: linear probe JetClass-II, largest set: relative drop in accuracy (%) with the qg data's missing slots set to draws, L188 :: sha256 c6c4138ee5c86423
\newcommand{\MlstLpJcTwoQgDrawAccDropFullPr}{54\%}  % experiments/FIGS/data/v2/index.json :: linear probe JetClass-II, largest set: relative drop in accuracy (%) with the qg data's missing slots set to draws, L188 :: sha256 c6c4138ee5c86423 :: prose twin, two significant figures
\newcommand{\MlstLpJcTwoQgZeroAccDropRes}{64\,\ensuremath{\pm}\,5\%}  % experiments/FIGS/data/v2/index.json :: linear probe JetClass-II, largest set: relative drop in accuracy (%) with the qg data's missing slots set to zero, L162 :: sha256 c6c4138ee5c86423
\newcommand{\MlstLpJcTwoQgZeroAccDropResPr}{64\%}  % experiments/FIGS/data/v2/index.json :: linear probe JetClass-II, largest set: relative drop in accuracy (%) with the qg data's missing slots set to zero, L162 :: sha256 c6c4138ee5c86423 :: prose twin, two significant figures
\newcommand{\MlstLpJcTwoQgMeanAccDropRes}{49.5\,\ensuremath{\pm}\,0.5\%}  % experiments/FIGS/data/v2/index.json :: linear probe JetClass-II, largest set: relative drop in accuracy (%) with the qg data's missing slots set to mean, L162 :: sha256 c6c4138ee5c86423
\newcommand{\MlstLpJcTwoQgMeanAccDropResPr}{49\%}  % experiments/FIGS/data/v2/index.json :: linear probe JetClass-II, largest set: relative drop in accuracy (%) with the qg data's missing slots set to mean, L162 :: sha256 c6c4138ee5c86423 :: prose twin, two significant figures
\newcommand{\MlstLpJcTwoQgDrawAccDropRes}{54.6\,\ensuremath{\pm}\,0.6\%}  % experiments/FIGS/data/v2/index.json :: linear probe JetClass-II, largest set: relative drop in accuracy (%) with the qg data's missing slots set to draws, L162 :: sha256 c6c4138ee5c86423
\newcommand{\MlstLpJcTwoQgDrawAccDropResPr}{55\%}  % experiments/FIGS/data/v2/index.json :: linear probe JetClass-II, largest set: relative drop in accuracy (%) with the qg data's missing slots set to draws, L162 :: sha256 c6c4138ee5c86423 :: prose twin, two significant figures
\newcommand{\MlstLpJcTwoQgZeroAccDropMid}{68\,\ensuremath{\pm}\,5\%}  % experiments/FIGS/data/v2/index.json :: linear probe JetClass-II, largest set: relative drop in accuracy (%) with the qg data's missing slots set to zero, R42_Q1 :: sha256 c6c4138ee5c86423
\newcommand{\MlstLpJcTwoQgZeroAccDropMidPr}{68\%}  % experiments/FIGS/data/v2/index.json :: linear probe JetClass-II, largest set: relative drop in accuracy (%) with the qg data's missing slots set to zero, R42_Q1 :: sha256 c6c4138ee5c86423 :: prose twin, two significant figures
\newcommand{\MlstLpJcTwoQgMeanAccDropMid}{47.6\,\ensuremath{\pm}\,1.6\%}  % experiments/FIGS/data/v2/index.json :: linear probe JetClass-II, largest set: relative drop in accuracy (%) with the qg data's missing slots set to mean, R42_Q1 :: sha256 c6c4138ee5c86423
\newcommand{\MlstLpJcTwoQgMeanAccDropMidPr}{48\%}  % experiments/FIGS/data/v2/index.json :: linear probe JetClass-II, largest set: relative drop in accuracy (%) with the qg data's missing slots set to mean, R42_Q1 :: sha256 c6c4138ee5c86423 :: prose twin, two significant figures
\newcommand{\MlstLpJcTwoQgDrawAccDropMid}{53.5\,\ensuremath{\pm}\,1.4\%}  % experiments/FIGS/data/v2/index.json :: linear probe JetClass-II, largest set: relative drop in accuracy (%) with the qg data's missing slots set to draws, R42_Q1 :: sha256 c6c4138ee5c86423
\newcommand{\MlstLpJcTwoQgDrawAccDropMidPr}{54\%}  % experiments/FIGS/data/v2/index.json :: linear probe JetClass-II, largest set: relative drop in accuracy (%) with the qg data's missing slots set to draws, R42_Q1 :: sha256 c6c4138ee5c86423 :: prose twin, two significant figures
\newcommand{\MlstLpJcTwoQgZeroAccDropCoarse}{70.9\,\ensuremath{\pm}\,2.3\%}  % experiments/FIGS/data/v2/index.json :: linear probe JetClass-II, largest set: relative drop in accuracy (%) with the qg data's missing slots set to zero, R16_Q1 :: sha256 c6c4138ee5c86423
\newcommand{\MlstLpJcTwoQgZeroAccDropCoarsePr}{71\%}  % experiments/FIGS/data/v2/index.json :: linear probe JetClass-II, largest set: relative drop in accuracy (%) with the qg data's missing slots set to zero, R16_Q1 :: sha256 c6c4138ee5c86423 :: prose twin, two significant figures
\newcommand{\MlstLpJcTwoQgMeanAccDropCoarse}{50.0\,\ensuremath{\pm}\,1.0\%}  % experiments/FIGS/data/v2/index.json :: linear probe JetClass-II, largest set: relative drop in accuracy (%) with the qg data's missing slots set to mean, R16_Q1 :: sha256 c6c4138ee5c86423
\newcommand{\MlstLpJcTwoQgMeanAccDropCoarsePr}{50\%}  % experiments/FIGS/data/v2/index.json :: linear probe JetClass-II, largest set: relative drop in accuracy (%) with the qg data's missing slots set to mean, R16_Q1 :: sha256 c6c4138ee5c86423 :: prose twin, two significant figures
\newcommand{\MlstLpJcTwoQgDrawAccDropCoarse}{51.7\,\ensuremath{\pm}\,0.4\%}  % experiments/FIGS/data/v2/index.json :: linear probe JetClass-II, largest set: relative drop in accuracy (%) with the qg data's missing slots set to draws, R16_Q1 :: sha256 c6c4138ee5c86423
\newcommand{\MlstLpJcTwoQgDrawAccDropCoarsePr}{52\%}  % experiments/FIGS/data/v2/index.json :: linear probe JetClass-II, largest set: relative drop in accuracy (%) with the qg data's missing slots set to draws, R16_Q1 :: sha256 c6c4138ee5c86423 :: prose twin, two significant figures
\newcommand{\MlstLpJcTwoQgZeroAccDropResMass}{61.8, 67.9\%}  % experiments/FIGS/data/v2/index.json :: linear probe JetClass-II, largest set: relative drop in accuracy (%) with the qg data's missing slots set to zero, L162_MASS :: sha256 c6c4138ee5c86423
\newcommand{\MlstLpJcTwoQgZeroAccDropResMassPr}{62--68\%}  % experiments/FIGS/data/v2/index.json :: linear probe JetClass-II, largest set: relative drop in accuracy (%) with the qg data's missing slots set to zero, L162_MASS :: sha256 c6c4138ee5c86423 :: prose twin, two significant figures
\newcommand{\MlstLpJcTwoQgMeanAccDropResMass}{49.7, 52.2\%}  % experiments/FIGS/data/v2/index.json :: linear probe JetClass-II, largest set: relative drop in accuracy (%) with the qg data's missing slots set to mean, L162_MASS :: sha256 c6c4138ee5c86423
\newcommand{\MlstLpJcTwoQgMeanAccDropResMassPr}{50--52\%}  % experiments/FIGS/data/v2/index.json :: linear probe JetClass-II, largest set: relative drop in accuracy (%) with the qg data's missing slots set to mean, L162_MASS :: sha256 c6c4138ee5c86423 :: prose twin, two significant figures
\newcommand{\MlstLpJcTwoQgDrawAccDropResMass}{54.3, 55.1\%}  % experiments/FIGS/data/v2/index.json :: linear probe JetClass-II, largest set: relative drop in accuracy (%) with the qg data's missing slots set to draws, L162_MASS :: sha256 c6c4138ee5c86423
\newcommand{\MlstLpJcTwoQgDrawAccDropResMassPr}{54--55\%}  % experiments/FIGS/data/v2/index.json :: linear probe JetClass-II, largest set: relative drop in accuracy (%) with the qg data's missing slots set to draws, L162_MASS :: sha256 c6c4138ee5c86423 :: prose twin, two significant figures
\newcommand{\MlstLpJcTwoQgZeroAccDropCoarseMass}{66.7, 65.4\%}  % experiments/FIGS/data/v2/index.json :: linear probe JetClass-II, largest set: relative drop in accuracy (%) with the qg data's missing slots set to zero, R16_Q1_MASS :: sha256 c6c4138ee5c86423
\newcommand{\MlstLpJcTwoQgZeroAccDropCoarseMassPr}{65--67\%}  % experiments/FIGS/data/v2/index.json :: linear probe JetClass-II, largest set: relative drop in accuracy (%) with the qg data's missing slots set to zero, R16_Q1_MASS :: sha256 c6c4138ee5c86423 :: prose twin, two significant figures
\newcommand{\MlstLpJcTwoQgMeanAccDropCoarseMass}{49.0, 49.6\%}  % experiments/FIGS/data/v2/index.json :: linear probe JetClass-II, largest set: relative drop in accuracy (%) with the qg data's missing slots set to mean, R16_Q1_MASS :: sha256 c6c4138ee5c86423
\newcommand{\MlstLpJcTwoQgMeanAccDropCoarseMassPr}{49--50\%}  % experiments/FIGS/data/v2/index.json :: linear probe JetClass-II, largest set: relative drop in accuracy (%) with the qg data's missing slots set to mean, R16_Q1_MASS :: sha256 c6c4138ee5c86423 :: prose twin, two significant figures
\newcommand{\MlstLpJcTwoQgDrawAccDropCoarseMass}{50.8, 51.1\%}  % experiments/FIGS/data/v2/index.json :: linear probe JetClass-II, largest set: relative drop in accuracy (%) with the qg data's missing slots set to draws, R16_Q1_MASS :: sha256 c6c4138ee5c86423
\newcommand{\MlstLpJcTwoQgDrawAccDropCoarseMassPr}{51--51\%}  % experiments/FIGS/data/v2/index.json :: linear probe JetClass-II, largest set: relative drop in accuracy (%) with the qg data's missing slots set to draws, R16_Q1_MASS :: sha256 c6c4138ee5c86423 :: prose twin, two significant figures
\newcommand{\MlstLpJcTwoQgZeroAccDropFzero}{62.7, 62.2\%}  % experiments/FIGS/data/v2/index.json :: linear probe JetClass-II, largest set: relative drop in accuracy (%) with the qg data's missing slots set to zero, FLAV_F0 :: sha256 c6c4138ee5c86423
\newcommand{\MlstLpJcTwoQgZeroAccDropFzeroPr}{62--63\%}  % experiments/FIGS/data/v2/index.json :: linear probe JetClass-II, largest set: relative drop in accuracy (%) with the qg data's missing slots set to zero, FLAV_F0 :: sha256 c6c4138ee5c86423 :: prose twin, two significant figures
\newcommand{\MlstLpJcTwoQgMeanAccDropFzero}{51.9, 51.3\%}  % experiments/FIGS/data/v2/index.json :: linear probe JetClass-II, largest set: relative drop in accuracy (%) with the qg data's missing slots set to mean, FLAV_F0 :: sha256 c6c4138ee5c86423
\newcommand{\MlstLpJcTwoQgMeanAccDropFzeroPr}{51--52\%}  % experiments/FIGS/data/v2/index.json :: linear probe JetClass-II, largest set: relative drop in accuracy (%) with the qg data's missing slots set to mean, FLAV_F0 :: sha256 c6c4138ee5c86423 :: prose twin, two significant figures
\newcommand{\MlstLpJcTwoQgDrawAccDropFzero}{53.9, 53.4\%}  % experiments/FIGS/data/v2/index.json :: linear probe JetClass-II, largest set: relative drop in accuracy (%) with the qg data's missing slots set to draws, FLAV_F0 :: sha256 c6c4138ee5c86423
\newcommand{\MlstLpJcTwoQgDrawAccDropFzeroPr}{53--54\%}  % experiments/FIGS/data/v2/index.json :: linear probe JetClass-II, largest set: relative drop in accuracy (%) with the qg data's missing slots set to draws, FLAV_F0 :: sha256 c6c4138ee5c86423 :: prose twin, two significant figures
\newcommand{\MlstLpJcTwoQgZeroAccDropFone}{66.2, 62.8\%}  % experiments/FIGS/data/v2/index.json :: linear probe JetClass-II, largest set: relative drop in accuracy (%) with the qg data's missing slots set to zero, FLAV_F1 :: sha256 c6c4138ee5c86423
\newcommand{\MlstLpJcTwoQgZeroAccDropFonePr}{63--66\%}  % experiments/FIGS/data/v2/index.json :: linear probe JetClass-II, largest set: relative drop in accuracy (%) with the qg data's missing slots set to zero, FLAV_F1 :: sha256 c6c4138ee5c86423 :: prose twin, two significant figures
\newcommand{\MlstLpJcTwoQgMeanAccDropFone}{52.7, 52.8\%}  % experiments/FIGS/data/v2/index.json :: linear probe JetClass-II, largest set: relative drop in accuracy (%) with the qg data's missing slots set to mean, FLAV_F1 :: sha256 c6c4138ee5c86423
\newcommand{\MlstLpJcTwoQgMeanAccDropFonePr}{53--53\%}  % experiments/FIGS/data/v2/index.json :: linear probe JetClass-II, largest set: relative drop in accuracy (%) with the qg data's missing slots set to mean, FLAV_F1 :: sha256 c6c4138ee5c86423 :: prose twin, two significant figures
\newcommand{\MlstLpJcTwoQgDrawAccDropFone}{55.6, 56.5\%}  % experiments/FIGS/data/v2/index.json :: linear probe JetClass-II, largest set: relative drop in accuracy (%) with the qg data's missing slots set to draws, FLAV_F1 :: sha256 c6c4138ee5c86423
\newcommand{\MlstLpJcTwoQgDrawAccDropFonePr}{56--56\%}  % experiments/FIGS/data/v2/index.json :: linear probe JetClass-II, largest set: relative drop in accuracy (%) with the qg data's missing slots set to draws, FLAV_F1 :: sha256 c6c4138ee5c86423 :: prose twin, two significant figures
\newcommand{\MlstLpJcTwoQgZeroAccDropFoner}{67.2, 69.0\%}  % experiments/FIGS/data/v2/index.json :: linear probe JetClass-II, largest set: relative drop in accuracy (%) with the qg data's missing slots set to zero, FLAV_F1R :: sha256 c6c4138ee5c86423
\newcommand{\MlstLpJcTwoQgZeroAccDropFonerPr}{67--69\%}  % experiments/FIGS/data/v2/index.json :: linear probe JetClass-II, largest set: relative drop in accuracy (%) with the qg data's missing slots set to zero, FLAV_F1R :: sha256 c6c4138ee5c86423 :: prose twin, two significant figures
\newcommand{\MlstLpJcTwoQgMeanAccDropFoner}{52.9, 51.8\%}  % experiments/FIGS/data/v2/index.json :: linear probe JetClass-II, largest set: relative drop in accuracy (%) with the qg data's missing slots set to mean, FLAV_F1R :: sha256 c6c4138ee5c86423
\newcommand{\MlstLpJcTwoQgMeanAccDropFonerPr}{52--53\%}  % experiments/FIGS/data/v2/index.json :: linear probe JetClass-II, largest set: relative drop in accuracy (%) with the qg data's missing slots set to mean, FLAV_F1R :: sha256 c6c4138ee5c86423 :: prose twin, two significant figures
\newcommand{\MlstLpJcTwoQgDrawAccDropFoner}{57.2, 55.4\%}  % experiments/FIGS/data/v2/index.json :: linear probe JetClass-II, largest set: relative drop in accuracy (%) with the qg data's missing slots set to draws, FLAV_F1R :: sha256 c6c4138ee5c86423
\newcommand{\MlstLpJcTwoQgDrawAccDropFonerPr}{55--57\%}  % experiments/FIGS/data/v2/index.json :: linear probe JetClass-II, largest set: relative drop in accuracy (%) with the qg data's missing slots set to draws, FLAV_F1R :: sha256 c6c4138ee5c86423 :: prose twin, two significant figures
\newcommand{\MlstLpJcTwoQgZeroAccDropInitBn}{8.7\,\ensuremath{\pm}\,2.6\%}  % experiments/FIGS/data/v2/index.json :: linear probe JetClass-II, largest set: relative drop in accuracy (%) with the qg data's missing slots set to zero, INIT_BN :: sha256 c6c4138ee5c86423
\newcommand{\MlstLpJcTwoQgZeroAccDropInitBnPr}{8.7\%}  % experiments/FIGS/data/v2/index.json :: linear probe JetClass-II, largest set: relative drop in accuracy (%) with the qg data's missing slots set to zero, INIT_BN :: sha256 c6c4138ee5c86423 :: prose twin, two significant figures
\newcommand{\MlstLpJcTwoQgMeanAccDropInitBn}{10.4\,\ensuremath{\pm}\,1.4\%}  % experiments/FIGS/data/v2/index.json :: linear probe JetClass-II, largest set: relative drop in accuracy (%) with the qg data's missing slots set to mean, INIT_BN :: sha256 c6c4138ee5c86423
\newcommand{\MlstLpJcTwoQgMeanAccDropInitBnPr}{10\%}  % experiments/FIGS/data/v2/index.json :: linear probe JetClass-II, largest set: relative drop in accuracy (%) with the qg data's missing slots set to mean, INIT_BN :: sha256 c6c4138ee5c86423 :: prose twin, two significant figures
\newcommand{\MlstLpJcTwoQgDrawAccDropInitBn}{11.4\,\ensuremath{\pm}\,2.1\%}  % experiments/FIGS/data/v2/index.json :: linear probe JetClass-II, largest set: relative drop in accuracy (%) with the qg data's missing slots set to draws, INIT_BN :: sha256 c6c4138ee5c86423
\newcommand{\MlstLpJcTwoQgDrawAccDropInitBnPr}{11\%}  % experiments/FIGS/data/v2/index.json :: linear probe JetClass-II, largest set: relative drop in accuracy (%) with the qg data's missing slots set to draws, INIT_BN :: sha256 c6c4138ee5c86423 :: prose twin, two significant figures
\newcommand{\MlstMassLambda}{5}  % configs/arms/v2_grid.json :: arms[L162_MASS, R16_Q1_MASS].mass_lambda :: sha256 d8a31b1486992302 :: weight of the mass-regression term in the pretraining loss, both mass models
\newcommand{\MlstMassLossShareResMass}{13.78 and 13.87}  % experiments/FIGS/data/v2/index.json :: share of the mass-regression term in the training loss of L162_MASS, mean over epochs, percent :: sha256 c6c4138ee5c86423
\newcommand{\MlstMassLossShareCoarseMass}{28.65 and 28.67}  % experiments/FIGS/data/v2/index.json :: share of the mass-regression term in the training loss of R16_Q1_MASS, mean over epochs, percent :: sha256 c6c4138ee5c86423
\newcommand{\MlstMassUncorrected}{0.245}  % experiments/FIGS/data/v2/index.json :: head/atlas_res: groomed mass, all jets, ATLAS resolution :: sha256 c6c4138ee5c86423
\newcommand{\MlstMassUncorrectedResponse}{0.929}  % experiments/FIGS/data/v2/index.json :: head/atlas_res: groomed mass, all jets, median response :: sha256 c6c4138ee5c86423
\newcommand{\MlstMassHeadResMass}{0.1471 and 0.1463}  % experiments/FIGS/data/v2/index.json :: regression output of L162_MASS, all jets, ATLAS resolution :: sha256 c6c4138ee5c86423
\newcommand{\MlstMassHeadResMassResponse}{0.995 and 1.003}  % experiments/FIGS/data/v2/index.json :: regression output of L162_MASS, all jets, median response :: sha256 c6c4138ee5c86423
\newcommand{\MlstMassHeadResMassOverUncorr}{0.6009 and 0.5976}  % experiments/FIGS/data/v2/index.json :: resolution of the regression output of L162_MASS over the groomed mass's, all jets :: sha256 c6c4138ee5c86423
\newcommand{\MlstMassHeadResMassOverUncorrPr}{at most 0.60 times the resolution}  % experiments/FIGS/data/v2/index.json :: resolution of the regression output of L162_MASS over the groomed mass's, all jets :: sha256 c6c4138ee5c86423 :: prose twin: two significant figures, direction in words
\newcommand{\MlstMassHeadResMassOverUncorrPrMax}{down to 0.60 times the resolution}  % experiments/FIGS/data/v2/index.json :: resolution of the regression output of L162_MASS over the groomed mass's, all jets :: sha256 c6c4138ee5c86423 :: prose twin of the run farther from 1
\newcommand{\MlstMassHeadResMassOverMlpRes}{0.6645 and 0.6691}  % experiments/FIGS/data/v2/index.json :: resolution of the regression output of L162_MASS over the frozen-feature MLP of L162, run by run, all jets :: sha256 c6c4138ee5c86423
\newcommand{\MlstMassHeadResMassOverMlpResPr}{at most 0.67 times the resolution}  % experiments/FIGS/data/v2/index.json :: resolution of the regression output of L162_MASS over the frozen-feature MLP of L162, run by run, all jets :: sha256 c6c4138ee5c86423 :: prose twin: two significant figures, direction in words
\newcommand{\MlstMassHeadResMassOverMlpResPrMax}{down to 0.66 times the resolution}  % experiments/FIGS/data/v2/index.json :: resolution of the regression output of L162_MASS over the frozen-feature MLP of L162, run by run, all jets :: sha256 c6c4138ee5c86423 :: prose twin of the run farther from 1
\newcommand{\MlstMassHeadCoarseMass}{0.1450 and 0.1461}  % experiments/FIGS/data/v2/index.json :: regression output of R16_Q1_MASS, all jets, ATLAS resolution :: sha256 c6c4138ee5c86423
\newcommand{\MlstMassHeadCoarseMassResponse}{1.005 and 1.005}  % experiments/FIGS/data/v2/index.json :: regression output of R16_Q1_MASS, all jets, median response :: sha256 c6c4138ee5c86423
\newcommand{\MlstMassHeadCoarseMassOverUncorr}{0.5921 and 0.5966}  % experiments/FIGS/data/v2/index.json :: resolution of the regression output of R16_Q1_MASS over the groomed mass's, all jets :: sha256 c6c4138ee5c86423
\newcommand{\MlstMassHeadCoarseMassOverUncorrPr}{at most 0.60 times the resolution}  % experiments/FIGS/data/v2/index.json :: resolution of the regression output of R16_Q1_MASS over the groomed mass's, all jets :: sha256 c6c4138ee5c86423 :: prose twin: two significant figures, direction in words
\newcommand{\MlstMassHeadCoarseMassOverUncorrPrMax}{down to 0.59 times the resolution}  % experiments/FIGS/data/v2/index.json :: resolution of the regression output of R16_Q1_MASS over the groomed mass's, all jets :: sha256 c6c4138ee5c86423 :: prose twin of the run farther from 1
\newcommand{\MlstMassHeadCoarseMassOverMlpCoarse}{0.703 and 0.690}  % experiments/FIGS/data/v2/index.json :: resolution of the regression output of R16_Q1_MASS over the frozen-feature MLP of R16_Q1, run by run, all jets :: sha256 c6c4138ee5c86423
\newcommand{\MlstMassHeadCoarseMassOverMlpCoarsePr}{at most 0.70 times the resolution}  % experiments/FIGS/data/v2/index.json :: resolution of the regression output of R16_Q1_MASS over the frozen-feature MLP of R16_Q1, run by run, all jets :: sha256 c6c4138ee5c86423 :: prose twin: two significant figures, direction in words
\newcommand{\MlstMassHeadCoarseMassOverMlpCoarsePrMax}{down to 0.69 times the resolution}  % experiments/FIGS/data/v2/index.json :: resolution of the regression output of R16_Q1_MASS over the frozen-feature MLP of R16_Q1, run by run, all jets :: sha256 c6c4138ee5c86423 :: prose twin of the run farther from 1
\newcommand{\MlstMassMlpInit}{0.2421\,\ensuremath{\pm}\,0.0018}  % experiments/FIGS/data/v2/index.json :: frozen-feature MLP of INIT, all jets, ATLAS resolution :: sha256 c6c4138ee5c86423
\newcommand{\MlstMassMlpFull}{0.2200\,\ensuremath{\pm}\,0.0014}  % experiments/FIGS/data/v2/index.json :: frozen-feature MLP of L188, all jets, ATLAS resolution :: sha256 c6c4138ee5c86423
\newcommand{\MlstMassMlpRes}{0.2195\,\ensuremath{\pm}\,0.0017}  % experiments/FIGS/data/v2/index.json :: frozen-feature MLP of L162, all jets, ATLAS resolution :: sha256 c6c4138ee5c86423
\newcommand{\MlstMassMlpMid}{0.2158\,\ensuremath{\pm}\,0.0029}  % experiments/FIGS/data/v2/index.json :: frozen-feature MLP of R42_Q1, all jets, ATLAS resolution :: sha256 c6c4138ee5c86423
\newcommand{\MlstMassMlpCoarse}{0.2098\,\ensuremath{\pm}\,0.0031}  % experiments/FIGS/data/v2/index.json :: frozen-feature MLP of R16_Q1, all jets, ATLAS resolution :: sha256 c6c4138ee5c86423
\newcommand{\MlstMassMlpResMass}{0.1591 and 0.1586}  % experiments/FIGS/data/v2/index.json :: frozen-feature MLP of L162_MASS, all jets, ATLAS resolution :: sha256 c6c4138ee5c86423
\newcommand{\MlstMassMlpCoarseMass}{0.1583 and 0.1588}  % experiments/FIGS/data/v2/index.json :: frozen-feature MLP of R16_Q1_MASS, all jets, ATLAS resolution :: sha256 c6c4138ee5c86423
\newcommand{\MlstCmpMlpResOverFull}{0.998\,\ensuremath{\pm}\,0.006}  % experiments/FIGS/data/v2/index.json :: frozen-feature MLP resolution of L162 over L188, run by run, all jets :: sha256 c6c4138ee5c86423
\newcommand{\MlstCmpMlpResOverFullPr}{equal within 1\%}  % experiments/FIGS/data/v2/index.json :: frozen-feature MLP resolution of L162 over L188, run by run, all jets :: sha256 c6c4138ee5c86423 :: prose twin: two significant figures, direction in words
\newcommand{\MlstCmpMlpMidOverFull}{0.981\,\ensuremath{\pm}\,0.017}  % experiments/FIGS/data/v2/index.json :: frozen-feature MLP resolution of R42_Q1 over L188, run by run, all jets :: sha256 c6c4138ee5c86423
\newcommand{\MlstCmpMlpMidOverFullPr}{equal within 2\%}  % experiments/FIGS/data/v2/index.json :: frozen-feature MLP resolution of R42_Q1 over L188, run by run, all jets :: sha256 c6c4138ee5c86423 :: prose twin: two significant figures, direction in words
\newcommand{\MlstCmpMlpCoarseOverFull}{0.954\,\ensuremath{\pm}\,0.018}  % experiments/FIGS/data/v2/index.json :: frozen-feature MLP resolution of R16_Q1 over L188, run by run, all jets :: sha256 c6c4138ee5c86423
\newcommand{\MlstCmpMlpCoarseOverFullPr}{0.95 times the resolution}  % experiments/FIGS/data/v2/index.json :: frozen-feature MLP resolution of R16_Q1 over L188, run by run, all jets :: sha256 c6c4138ee5c86423 :: prose twin: two significant figures, direction in words
\newcommand{\MlstCmpMlpResMassOverRes}{0.7186 and 0.7253}  % experiments/FIGS/data/v2/index.json :: frozen-feature MLP resolution of L162_MASS over L162, run by run, all jets :: sha256 c6c4138ee5c86423
\newcommand{\MlstCmpMlpResMassOverResPr}{at most 0.73 times the resolution}  % experiments/FIGS/data/v2/index.json :: frozen-feature MLP resolution of L162_MASS over L162, run by run, all jets :: sha256 c6c4138ee5c86423 :: prose twin: two significant figures, direction in words
\newcommand{\MlstCmpMlpResMassOverResPrMax}{down to 0.72 times the resolution}  % experiments/FIGS/data/v2/index.json :: frozen-feature MLP resolution of L162_MASS over L162, run by run, all jets :: sha256 c6c4138ee5c86423 :: prose twin of the run farther from 1
\newcommand{\MlstCmpMlpCoarseMassOverCoarse}{0.768 and 0.750}  % experiments/FIGS/data/v2/index.json :: frozen-feature MLP resolution of R16_Q1_MASS over R16_Q1, run by run, all jets :: sha256 c6c4138ee5c86423
\newcommand{\MlstCmpMlpCoarseMassOverCoarsePr}{at most 0.77 times the resolution}  % experiments/FIGS/data/v2/index.json :: frozen-feature MLP resolution of R16_Q1_MASS over R16_Q1, run by run, all jets :: sha256 c6c4138ee5c86423 :: prose twin: two significant figures, direction in words
\newcommand{\MlstCmpMlpCoarseMassOverCoarsePrMax}{down to 0.75 times the resolution}  % experiments/FIGS/data/v2/index.json :: frozen-feature MLP resolution of R16_Q1_MASS over R16_Q1, run by run, all jets :: sha256 c6c4138ee5c86423 :: prose twin of the run farther from 1
\newcommand{\MlstMassUncorrectedResonant}{0.213}  % experiments/FIGS/data/v2/index.json :: head/atlas_res: groomed mass, resonant jets, ATLAS resolution :: sha256 c6c4138ee5c86423
\newcommand{\MlstMassUncorrectedResponseResonant}{0.929}  % experiments/FIGS/data/v2/index.json :: head/atlas_res: groomed mass, resonant jets, median response :: sha256 c6c4138ee5c86423
\newcommand{\MlstMassHeadResMassResonant}{0.1332 and 0.1320}  % experiments/FIGS/data/v2/index.json :: regression output of L162_MASS, resonant jets, ATLAS resolution :: sha256 c6c4138ee5c86423
\newcommand{\MlstMassHeadResMassResponseResonant}{0.996 and 1.004}  % experiments/FIGS/data/v2/index.json :: regression output of L162_MASS, resonant jets, median response :: sha256 c6c4138ee5c86423
\newcommand{\MlstMassHeadResMassOverUncorrResonant}{0.6259 and 0.6205}  % experiments/FIGS/data/v2/index.json :: resolution of the regression output of L162_MASS over the groomed mass's, resonant jets :: sha256 c6c4138ee5c86423
\newcommand{\MlstMassHeadResMassOverUncorrResonantPr}{at most 0.63 times the resolution}  % experiments/FIGS/data/v2/index.json :: resolution of the regression output of L162_MASS over the groomed mass's, resonant jets :: sha256 c6c4138ee5c86423 :: prose twin: two significant figures, direction in words
\newcommand{\MlstMassHeadResMassOverUncorrResonantPrMax}{down to 0.62 times the resolution}  % experiments/FIGS/data/v2/index.json :: resolution of the regression output of L162_MASS over the groomed mass's, resonant jets :: sha256 c6c4138ee5c86423 :: prose twin of the run farther from 1
\newcommand{\MlstMassHeadResMassOverMlpResResonant}{0.6620 and 0.6662}  % experiments/FIGS/data/v2/index.json :: resolution of the regression output of L162_MASS over the frozen-feature MLP of L162, run by run, resonant jets :: sha256 c6c4138ee5c86423
\newcommand{\MlstMassHeadResMassOverMlpResResonantPr}{at most 0.67 times the resolution}  % experiments/FIGS/data/v2/index.json :: resolution of the regression output of L162_MASS over the frozen-feature MLP of L162, run by run, resonant jets :: sha256 c6c4138ee5c86423 :: prose twin: two significant figures, direction in words
\newcommand{\MlstMassHeadResMassOverMlpResResonantPrMax}{down to 0.66 times the resolution}  % experiments/FIGS/data/v2/index.json :: resolution of the regression output of L162_MASS over the frozen-feature MLP of L162, run by run, resonant jets :: sha256 c6c4138ee5c86423 :: prose twin of the run farther from 1
\newcommand{\MlstMassHeadCoarseMassResonant}{0.1312 and 0.1320}  % experiments/FIGS/data/v2/index.json :: regression output of R16_Q1_MASS, resonant jets, ATLAS resolution :: sha256 c6c4138ee5c86423
\newcommand{\MlstMassHeadCoarseMassResponseResonant}{1.006 and 1.006}  % experiments/FIGS/data/v2/index.json :: regression output of R16_Q1_MASS, resonant jets, median response :: sha256 c6c4138ee5c86423
\newcommand{\MlstMassHeadCoarseMassOverUncorrResonant}{0.6165 and 0.6206}  % experiments/FIGS/data/v2/index.json :: resolution of the regression output of R16_Q1_MASS over the groomed mass's, resonant jets :: sha256 c6c4138ee5c86423
\newcommand{\MlstMassHeadCoarseMassOverUncorrResonantPr}{at most 0.62 times the resolution}  % experiments/FIGS/data/v2/index.json :: resolution of the regression output of R16_Q1_MASS over the groomed mass's, resonant jets :: sha256 c6c4138ee5c86423 :: prose twin: two significant figures, direction in words
\newcommand{\MlstMassHeadCoarseMassOverUncorrResonantPrMax}{down to 0.62 times the resolution}  % experiments/FIGS/data/v2/index.json :: resolution of the regression output of R16_Q1_MASS over the groomed mass's, resonant jets :: sha256 c6c4138ee5c86423 :: prose twin of the run farther from 1
\newcommand{\MlstMassHeadCoarseMassOverMlpCoarseResonant}{0.707 and 0.688}  % experiments/FIGS/data/v2/index.json :: resolution of the regression output of R16_Q1_MASS over the frozen-feature MLP of R16_Q1, run by run, resonant jets :: sha256 c6c4138ee5c86423
\newcommand{\MlstMassHeadCoarseMassOverMlpCoarseResonantPr}{at most 0.71 times the resolution}  % experiments/FIGS/data/v2/index.json :: resolution of the regression output of R16_Q1_MASS over the frozen-feature MLP of R16_Q1, run by run, resonant jets :: sha256 c6c4138ee5c86423 :: prose twin: two significant figures, direction in words
\newcommand{\MlstMassHeadCoarseMassOverMlpCoarseResonantPrMax}{down to 0.69 times the resolution}  % experiments/FIGS/data/v2/index.json :: resolution of the regression output of R16_Q1_MASS over the frozen-feature MLP of R16_Q1, run by run, resonant jets :: sha256 c6c4138ee5c86423 :: prose twin of the run farther from 1
\newcommand{\MlstMassMlpInitResonant}{0.2214\,\ensuremath{\pm}\,0.0017}  % experiments/FIGS/data/v2/index.json :: frozen-feature MLP of INIT, resonant jets, ATLAS resolution :: sha256 c6c4138ee5c86423
\newcommand{\MlstMassMlpFullResonant}{0.1990\,\ensuremath{\pm}\,0.0016}  % experiments/FIGS/data/v2/index.json :: frozen-feature MLP of L188, resonant jets, ATLAS resolution :: sha256 c6c4138ee5c86423
\newcommand{\MlstMassMlpResResonant}{0.1989\,\ensuremath{\pm}\,0.0020}  % experiments/FIGS/data/v2/index.json :: frozen-feature MLP of L162, resonant jets, ATLAS resolution :: sha256 c6c4138ee5c86423
\newcommand{\MlstMassMlpMidResonant}{0.1949\,\ensuremath{\pm}\,0.0023}  % experiments/FIGS/data/v2/index.json :: frozen-feature MLP of R42_Q1, resonant jets, ATLAS resolution :: sha256 c6c4138ee5c86423
\newcommand{\MlstMassMlpCoarseResonant}{0.1894\,\ensuremath{\pm}\,0.0035}  % experiments/FIGS/data/v2/index.json :: frozen-feature MLP of R16_Q1, resonant jets, ATLAS resolution :: sha256 c6c4138ee5c86423
\newcommand{\MlstMassMlpResMassResonant}{0.1436 and 0.1430}  % experiments/FIGS/data/v2/index.json :: frozen-feature MLP of L162_MASS, resonant jets, ATLAS resolution :: sha256 c6c4138ee5c86423
\newcommand{\MlstMassMlpCoarseMassResonant}{0.1420 and 0.1431}  % experiments/FIGS/data/v2/index.json :: frozen-feature MLP of R16_Q1_MASS, resonant jets, ATLAS resolution :: sha256 c6c4138ee5c86423
\newcommand{\MlstCmpMlpResOverFullResonant}{1.000\,\ensuremath{\pm}\,0.005}  % experiments/FIGS/data/v2/index.json :: frozen-feature MLP resolution of L162 over L188, run by run, resonant jets :: sha256 c6c4138ee5c86423
\newcommand{\MlstCmpMlpResOverFullResonantPr}{equal within 1\%}  % experiments/FIGS/data/v2/index.json :: frozen-feature MLP resolution of L162 over L188, run by run, resonant jets :: sha256 c6c4138ee5c86423 :: prose twin: two significant figures, direction in words
\newcommand{\MlstCmpMlpMidOverFullResonant}{0.980\,\ensuremath{\pm}\,0.018}  % experiments/FIGS/data/v2/index.json :: frozen-feature MLP resolution of R42_Q1 over L188, run by run, resonant jets :: sha256 c6c4138ee5c86423
\newcommand{\MlstCmpMlpMidOverFullResonantPr}{equal within 3\%}  % experiments/FIGS/data/v2/index.json :: frozen-feature MLP resolution of R42_Q1 over L188, run by run, resonant jets :: sha256 c6c4138ee5c86423 :: prose twin: two significant figures, direction in words
\newcommand{\MlstCmpMlpCoarseOverFullResonant}{0.952\,\ensuremath{\pm}\,0.023}  % experiments/FIGS/data/v2/index.json :: frozen-feature MLP resolution of R16_Q1 over L188, run by run, resonant jets :: sha256 c6c4138ee5c86423
\newcommand{\MlstCmpMlpCoarseOverFullResonantPr}{0.95 times the resolution}  % experiments/FIGS/data/v2/index.json :: frozen-feature MLP resolution of R16_Q1 over L188, run by run, resonant jets :: sha256 c6c4138ee5c86423 :: prose twin: two significant figures, direction in words
\newcommand{\MlstCmpMlpResMassOverResResonant}{0.7137 and 0.7219}  % experiments/FIGS/data/v2/index.json :: frozen-feature MLP resolution of L162_MASS over L162, run by run, resonant jets :: sha256 c6c4138ee5c86423
\newcommand{\MlstCmpMlpResMassOverResResonantPr}{at most 0.72 times the resolution}  % experiments/FIGS/data/v2/index.json :: frozen-feature MLP resolution of L162_MASS over L162, run by run, resonant jets :: sha256 c6c4138ee5c86423 :: prose twin: two significant figures, direction in words
\newcommand{\MlstCmpMlpResMassOverResResonantPrMax}{down to 0.71 times the resolution}  % experiments/FIGS/data/v2/index.json :: frozen-feature MLP resolution of L162_MASS over L162, run by run, resonant jets :: sha256 c6c4138ee5c86423 :: prose twin of the run farther from 1
\newcommand{\MlstCmpMlpCoarseMassOverCoarseResonant}{0.766 and 0.745}  % experiments/FIGS/data/v2/index.json :: frozen-feature MLP resolution of R16_Q1_MASS over R16_Q1, run by run, resonant jets :: sha256 c6c4138ee5c86423
\newcommand{\MlstCmpMlpCoarseMassOverCoarseResonantPr}{at most 0.77 times the resolution}  % experiments/FIGS/data/v2/index.json :: frozen-feature MLP resolution of R16_Q1_MASS over R16_Q1, run by run, resonant jets :: sha256 c6c4138ee5c86423 :: prose twin: two significant figures, direction in words
\newcommand{\MlstCmpMlpCoarseMassOverCoarseResonantPrMax}{down to 0.75 times the resolution}  % experiments/FIGS/data/v2/index.json :: frozen-feature MLP resolution of R16_Q1_MASS over R16_Q1, run by run, resonant jets :: sha256 c6c4138ee5c86423 :: prose twin of the run farther from 1
\newcommand{\MlstMassUncorrectedQcd}{1.52}  % experiments/FIGS/data/v2/index.json :: head/atlas_res: groomed mass, qcd jets, ATLAS resolution :: sha256 c6c4138ee5c86423
\newcommand{\MlstMassUncorrectedResponseQcd}{0.931}  % experiments/FIGS/data/v2/index.json :: head/atlas_res: groomed mass, qcd jets, median response :: sha256 c6c4138ee5c86423
\newcommand{\MlstMassHeadResMassQcd}{0.2263 and 0.2256}  % experiments/FIGS/data/v2/index.json :: regression output of L162_MASS, qcd jets, ATLAS resolution :: sha256 c6c4138ee5c86423
\newcommand{\MlstMassHeadResMassResponseQcd}{0.9897 and 0.9996}  % experiments/FIGS/data/v2/index.json :: regression output of L162_MASS, qcd jets, median response :: sha256 c6c4138ee5c86423
\newcommand{\MlstMassHeadResMassOverUncorrQcd}{0.1485 and 0.1480}  % experiments/FIGS/data/v2/index.json :: resolution of the regression output of L162_MASS over the groomed mass's, qcd jets :: sha256 c6c4138ee5c86423
\newcommand{\MlstMassHeadResMassOverUncorrQcdPr}{at most 0.15 times the resolution}  % experiments/FIGS/data/v2/index.json :: resolution of the regression output of L162_MASS over the groomed mass's, qcd jets :: sha256 c6c4138ee5c86423 :: prose twin: two significant figures, direction in words
\newcommand{\MlstMassHeadResMassOverUncorrQcdPrMax}{down to 0.15 times the resolution}  % experiments/FIGS/data/v2/index.json :: resolution of the regression output of L162_MASS over the groomed mass's, qcd jets :: sha256 c6c4138ee5c86423 :: prose twin of the run farther from 1
\newcommand{\MlstMassHeadResMassOverMlpResQcd}{0.6542 and 0.6632}  % experiments/FIGS/data/v2/index.json :: resolution of the regression output of L162_MASS over the frozen-feature MLP of L162, run by run, qcd jets :: sha256 c6c4138ee5c86423
\newcommand{\MlstMassHeadResMassOverMlpResQcdPr}{at most 0.66 times the resolution}  % experiments/FIGS/data/v2/index.json :: resolution of the regression output of L162_MASS over the frozen-feature MLP of L162, run by run, qcd jets :: sha256 c6c4138ee5c86423 :: prose twin: two significant figures, direction in words
\newcommand{\MlstMassHeadResMassOverMlpResQcdPrMax}{down to 0.65 times the resolution}  % experiments/FIGS/data/v2/index.json :: resolution of the regression output of L162_MASS over the frozen-feature MLP of L162, run by run, qcd jets :: sha256 c6c4138ee5c86423 :: prose twin of the run farther from 1
\newcommand{\MlstMassHeadCoarseMassQcd}{0.2239 and 0.2240}  % experiments/FIGS/data/v2/index.json :: regression output of R16_Q1_MASS, qcd jets, ATLAS resolution :: sha256 c6c4138ee5c86423
\newcommand{\MlstMassHeadCoarseMassResponseQcd}{1.001 and 1.000}  % experiments/FIGS/data/v2/index.json :: regression output of R16_Q1_MASS, qcd jets, median response :: sha256 c6c4138ee5c86423
\newcommand{\MlstMassHeadCoarseMassOverUncorrQcd}{0.1469 and 0.1469}  % experiments/FIGS/data/v2/index.json :: resolution of the regression output of R16_Q1_MASS over the groomed mass's, qcd jets :: sha256 c6c4138ee5c86423
\newcommand{\MlstMassHeadCoarseMassOverUncorrQcdPr}{at most 0.15 times the resolution}  % experiments/FIGS/data/v2/index.json :: resolution of the regression output of R16_Q1_MASS over the groomed mass's, qcd jets :: sha256 c6c4138ee5c86423 :: prose twin: two significant figures, direction in words
\newcommand{\MlstMassHeadCoarseMassOverUncorrQcdPrMax}{down to 0.15 times the resolution}  % experiments/FIGS/data/v2/index.json :: resolution of the regression output of R16_Q1_MASS over the groomed mass's, qcd jets :: sha256 c6c4138ee5c86423 :: prose twin of the run farther from 1
\newcommand{\MlstMassHeadCoarseMassOverMlpCoarseQcd}{0.679 and 0.698}  % experiments/FIGS/data/v2/index.json :: resolution of the regression output of R16_Q1_MASS over the frozen-feature MLP of R16_Q1, run by run, qcd jets :: sha256 c6c4138ee5c86423
\newcommand{\MlstMassHeadCoarseMassOverMlpCoarseQcdPr}{at most 0.70 times the resolution}  % experiments/FIGS/data/v2/index.json :: resolution of the regression output of R16_Q1_MASS over the frozen-feature MLP of R16_Q1, run by run, qcd jets :: sha256 c6c4138ee5c86423 :: prose twin: two significant figures, direction in words
\newcommand{\MlstMassHeadCoarseMassOverMlpCoarseQcdPrMax}{down to 0.68 times the resolution}  % experiments/FIGS/data/v2/index.json :: resolution of the regression output of R16_Q1_MASS over the frozen-feature MLP of R16_Q1, run by run, qcd jets :: sha256 c6c4138ee5c86423 :: prose twin of the run farther from 1
\newcommand{\MlstMassMlpInitQcd}{0.502\,\ensuremath{\pm}\,0.014}  % experiments/FIGS/data/v2/index.json :: frozen-feature MLP of INIT, qcd jets, ATLAS resolution :: sha256 c6c4138ee5c86423
\newcommand{\MlstMassMlpFullQcd}{0.3487\,\ensuremath{\pm}\,0.0029}  % experiments/FIGS/data/v2/index.json :: frozen-feature MLP of L188, qcd jets, ATLAS resolution :: sha256 c6c4138ee5c86423
\newcommand{\MlstMassMlpResQcd}{0.3425\,\ensuremath{\pm}\,0.0031}  % experiments/FIGS/data/v2/index.json :: frozen-feature MLP of L162, qcd jets, ATLAS resolution :: sha256 c6c4138ee5c86423
\newcommand{\MlstMassMlpMidQcd}{0.338\,\ensuremath{\pm}\,0.007}  % experiments/FIGS/data/v2/index.json :: frozen-feature MLP of R42_Q1, qcd jets, ATLAS resolution :: sha256 c6c4138ee5c86423
\newcommand{\MlstMassMlpCoarseQcd}{0.327\,\ensuremath{\pm}\,0.006}  % experiments/FIGS/data/v2/index.json :: frozen-feature MLP of R16_Q1, qcd jets, ATLAS resolution :: sha256 c6c4138ee5c86423
\newcommand{\MlstMassMlpResMassQcd}{0.2473 and 0.2452}  % experiments/FIGS/data/v2/index.json :: frozen-feature MLP of L162_MASS, qcd jets, ATLAS resolution :: sha256 c6c4138ee5c86423
\newcommand{\MlstMassMlpCoarseMassQcd}{0.2465 and 0.2463}  % experiments/FIGS/data/v2/index.json :: frozen-feature MLP of R16_Q1_MASS, qcd jets, ATLAS resolution :: sha256 c6c4138ee5c86423
\newcommand{\MlstCmpMlpResOverFullQcd}{0.982\,\ensuremath{\pm}\,0.017}  % experiments/FIGS/data/v2/index.json :: frozen-feature MLP resolution of L162 over L188, run by run, qcd jets :: sha256 c6c4138ee5c86423
\newcommand{\MlstCmpMlpResOverFullQcdPr}{equal within 2\%}  % experiments/FIGS/data/v2/index.json :: frozen-feature MLP resolution of L162 over L188, run by run, qcd jets :: sha256 c6c4138ee5c86423 :: prose twin: two significant figures, direction in words
\newcommand{\MlstCmpMlpMidOverFullQcd}{0.969\,\ensuremath{\pm}\,0.013}  % experiments/FIGS/data/v2/index.json :: frozen-feature MLP resolution of R42_Q1 over L188, run by run, qcd jets :: sha256 c6c4138ee5c86423
\newcommand{\MlstCmpMlpMidOverFullQcdPr}{0.97 times the resolution}  % experiments/FIGS/data/v2/index.json :: frozen-feature MLP resolution of R42_Q1 over L188, run by run, qcd jets :: sha256 c6c4138ee5c86423 :: prose twin: two significant figures, direction in words
\newcommand{\MlstCmpMlpCoarseOverFullQcd}{0.939\,\ensuremath{\pm}\,0.023}  % experiments/FIGS/data/v2/index.json :: frozen-feature MLP resolution of R16_Q1 over L188, run by run, qcd jets :: sha256 c6c4138ee5c86423
\newcommand{\MlstCmpMlpCoarseOverFullQcdPr}{0.94 times the resolution}  % experiments/FIGS/data/v2/index.json :: frozen-feature MLP resolution of R16_Q1 over L188, run by run, qcd jets :: sha256 c6c4138ee5c86423 :: prose twin: two significant figures, direction in words
\newcommand{\MlstCmpMlpResMassOverResQcd}{0.7147 and 0.7209}  % experiments/FIGS/data/v2/index.json :: frozen-feature MLP resolution of L162_MASS over L162, run by run, qcd jets :: sha256 c6c4138ee5c86423
\newcommand{\MlstCmpMlpResMassOverResQcdPr}{at most 0.72 times the resolution}  % experiments/FIGS/data/v2/index.json :: frozen-feature MLP resolution of L162_MASS over L162, run by run, qcd jets :: sha256 c6c4138ee5c86423 :: prose twin: two significant figures, direction in words
\newcommand{\MlstCmpMlpResMassOverResQcdPrMax}{down to 0.71 times the resolution}  % experiments/FIGS/data/v2/index.json :: frozen-feature MLP resolution of L162_MASS over L162, run by run, qcd jets :: sha256 c6c4138ee5c86423 :: prose twin of the run farther from 1
\newcommand{\MlstCmpMlpCoarseMassOverCoarseQcd}{0.747 and 0.768}  % experiments/FIGS/data/v2/index.json :: frozen-feature MLP resolution of R16_Q1_MASS over R16_Q1, run by run, qcd jets :: sha256 c6c4138ee5c86423
\newcommand{\MlstCmpMlpCoarseMassOverCoarseQcdPr}{at most 0.77 times the resolution}  % experiments/FIGS/data/v2/index.json :: frozen-feature MLP resolution of R16_Q1_MASS over R16_Q1, run by run, qcd jets :: sha256 c6c4138ee5c86423 :: prose twin: two significant figures, direction in words
\newcommand{\MlstCmpMlpCoarseMassOverCoarseQcdPrMax}{down to 0.75 times the resolution}  % experiments/FIGS/data/v2/index.json :: frozen-feature MLP resolution of R16_Q1_MASS over R16_Q1, run by run, qcd jets :: sha256 c6c4138ee5c86423 :: prose twin of the run farther from 1
\newcommand{\MlstMassRespAxesShareQcd}{76\%}  % experiments/FIGS/data/v2/index.json :: head/atlas_res: share of QCD jets whose groomed-mass response lies inside F9a(b)'s x axis :: sha256 c6c4138ee5c86423
\newcommand{\MlstSculptNPassTen}{41{,}028}  % experiments/FIGS/data/v2/index.json :: QCD test jets passing the cut at QCD efficiency 0.1, every run :: sha256 c6c4138ee5c86423
\newcommand{\MlstSculptFloorTen}{0.00944}  % experiments/FIGS/data/v2/index.json :: Jensen-Shannon distance expected for a random selection of as many QCD jets as pass at 0.1 :: sha256 c6c4138ee5c86423
\newcommand{\MlstSculptFullTen}{0.0524\,\ensuremath{\pm}\,0.0034}  % experiments/FIGS/data/v2/index.json :: Jensen-Shannon distance of the QCD m_SD distribution at QCD efficiency 0.1, L188 :: sha256 c6c4138ee5c86423
\newcommand{\MlstSculptFullOverFloorTen}{5.6\,\ensuremath{\pm}\,0.4}  % experiments/FIGS/data/v2/index.json :: distance at 0.1 over the random-selection floor, L188 :: sha256 c6c4138ee5c86423
\newcommand{\MlstSculptFullOverFloorTenPr}{5.6 times the distance}  % experiments/FIGS/data/v2/index.json :: distance at 0.1 over the random-selection floor, L188 :: sha256 c6c4138ee5c86423 :: prose twin: two significant figures, direction in words
\newcommand{\MlstSculptResTen}{0.0532\,\ensuremath{\pm}\,0.0009}  % experiments/FIGS/data/v2/index.json :: Jensen-Shannon distance of the QCD m_SD distribution at QCD efficiency 0.1, L162 :: sha256 c6c4138ee5c86423
\newcommand{\MlstSculptResOverFloorTen}{5.64\,\ensuremath{\pm}\,0.10}  % experiments/FIGS/data/v2/index.json :: distance at 0.1 over the random-selection floor, L162 :: sha256 c6c4138ee5c86423
\newcommand{\MlstSculptResOverFloorTenPr}{5.6 times the distance}  % experiments/FIGS/data/v2/index.json :: distance at 0.1 over the random-selection floor, L162 :: sha256 c6c4138ee5c86423 :: prose twin: two significant figures, direction in words
\newcommand{\MlstSculptMidTen}{0.0560\,\ensuremath{\pm}\,0.0027}  % experiments/FIGS/data/v2/index.json :: Jensen-Shannon distance of the QCD m_SD distribution at QCD efficiency 0.1, R42_Q1 :: sha256 c6c4138ee5c86423
\newcommand{\MlstSculptMidOverFloorTen}{5.94\,\ensuremath{\pm}\,0.28}  % experiments/FIGS/data/v2/index.json :: distance at 0.1 over the random-selection floor, R42_Q1 :: sha256 c6c4138ee5c86423
\newcommand{\MlstSculptMidOverFloorTenPr}{5.9 times the distance}  % experiments/FIGS/data/v2/index.json :: distance at 0.1 over the random-selection floor, R42_Q1 :: sha256 c6c4138ee5c86423 :: prose twin: two significant figures, direction in words
\newcommand{\MlstSculptCoarseTen}{0.0543\,\ensuremath{\pm}\,0.0010}  % experiments/FIGS/data/v2/index.json :: Jensen-Shannon distance of the QCD m_SD distribution at QCD efficiency 0.1, R16_Q1 :: sha256 c6c4138ee5c86423
\newcommand{\MlstSculptCoarseOverFloorTen}{5.76\,\ensuremath{\pm}\,0.10}  % experiments/FIGS/data/v2/index.json :: distance at 0.1 over the random-selection floor, R16_Q1 :: sha256 c6c4138ee5c86423
\newcommand{\MlstSculptCoarseOverFloorTenPr}{5.8 times the distance}  % experiments/FIGS/data/v2/index.json :: distance at 0.1 over the random-selection floor, R16_Q1 :: sha256 c6c4138ee5c86423 :: prose twin: two significant figures, direction in words
\newcommand{\MlstSculptResMassTen}{0.0422 and 0.0464}  % experiments/FIGS/data/v2/index.json :: Jensen-Shannon distance of the QCD m_SD distribution at QCD efficiency 0.1, L162_MASS :: sha256 c6c4138ee5c86423
\newcommand{\MlstSculptResMassOverFloorTen}{4.47 and 4.91}  % experiments/FIGS/data/v2/index.json :: distance at 0.1 over the random-selection floor, L162_MASS :: sha256 c6c4138ee5c86423
\newcommand{\MlstSculptResMassOverFloorTenPr}{at least 4.5 times the distance}  % experiments/FIGS/data/v2/index.json :: distance at 0.1 over the random-selection floor, L162_MASS :: sha256 c6c4138ee5c86423 :: prose twin: two significant figures, direction in words
\newcommand{\MlstSculptResMassOverFloorTenPrMax}{up to 4.9 times the distance}  % experiments/FIGS/data/v2/index.json :: distance at 0.1 over the random-selection floor, L162_MASS :: sha256 c6c4138ee5c86423 :: prose twin of the run farther from 1
\newcommand{\MlstSculptCoarseMassTen}{0.0522 and 0.0423}  % experiments/FIGS/data/v2/index.json :: Jensen-Shannon distance of the QCD m_SD distribution at QCD efficiency 0.1, R16_Q1_MASS :: sha256 c6c4138ee5c86423
\newcommand{\MlstSculptCoarseMassOverFloorTen}{5.5 and 4.5}  % experiments/FIGS/data/v2/index.json :: distance at 0.1 over the random-selection floor, R16_Q1_MASS :: sha256 c6c4138ee5c86423
\newcommand{\MlstSculptCoarseMassOverFloorTenPr}{at least 4.5 times the distance}  % experiments/FIGS/data/v2/index.json :: distance at 0.1 over the random-selection floor, R16_Q1_MASS :: sha256 c6c4138ee5c86423 :: prose twin: two significant figures, direction in words
\newcommand{\MlstSculptCoarseMassOverFloorTenPrMax}{up to 5.5 times the distance}  % experiments/FIGS/data/v2/index.json :: distance at 0.1 over the random-selection floor, R16_Q1_MASS :: sha256 c6c4138ee5c86423 :: prose twin of the run farther from 1
\newcommand{\MlstSculptResMassOverResTen}{0.789 and 0.860}  % experiments/FIGS/data/v2/index.json :: distance at QCD efficiency 0.1: L162_MASS over L162, run by run :: sha256 c6c4138ee5c86423
\newcommand{\MlstSculptResMassOverResTenPr}{at most 0.86 times the distance}  % experiments/FIGS/data/v2/index.json :: distance at QCD efficiency 0.1: L162_MASS over L162, run by run :: sha256 c6c4138ee5c86423 :: prose twin: two significant figures, direction in words
\newcommand{\MlstSculptResMassOverResTenPrMax}{down to 0.79 times the distance}  % experiments/FIGS/data/v2/index.json :: distance at QCD efficiency 0.1: L162_MASS over L162, run by run :: sha256 c6c4138ee5c86423 :: prose twin of the run farther from 1
\newcommand{\MlstSculptCoarseMassOverCoarseTen}{0.94 and 0.79}  % experiments/FIGS/data/v2/index.json :: distance at QCD efficiency 0.1: R16_Q1_MASS over R16_Q1, run by run :: sha256 c6c4138ee5c86423
\newcommand{\MlstSculptCoarseMassOverCoarseTenPr}{equal within 30\%}  % experiments/FIGS/data/v2/index.json :: distance at QCD efficiency 0.1: R16_Q1_MASS over R16_Q1, run by run :: sha256 c6c4138ee5c86423 :: prose twin: two significant figures, direction in words
\newcommand{\MlstSculptCoarseMassOverCoarseTenPrMax}{equal within 30\%}  % experiments/FIGS/data/v2/index.json :: distance at QCD efficiency 0.1: R16_Q1_MASS over R16_Q1, run by run :: sha256 c6c4138ee5c86423 :: prose twin of the run farther from 1
\newcommand{\MlstSculptCoarseOverFullTen}{1.04\,\ensuremath{\pm}\,0.05}  % experiments/FIGS/data/v2/index.json :: distance at QCD efficiency 0.1: R16_Q1 over L188, run by run :: sha256 c6c4138ee5c86423
\newcommand{\MlstSculptCoarseOverFullTenPr}{equal within 6\%}  % experiments/FIGS/data/v2/index.json :: distance at QCD efficiency 0.1: R16_Q1 over L188, run by run :: sha256 c6c4138ee5c86423 :: prose twin: two significant figures, direction in words
\newcommand{\MlstSculptResOverFullTen}{1.02\,\ensuremath{\pm}\,0.06}  % experiments/FIGS/data/v2/index.json :: distance at QCD efficiency 0.1: L162 over L188, run by run :: sha256 c6c4138ee5c86423
\newcommand{\MlstSculptResOverFullTenPr}{equal within 6\%}  % experiments/FIGS/data/v2/index.json :: distance at QCD efficiency 0.1: L162 over L188, run by run :: sha256 c6c4138ee5c86423 :: prose twin: two significant figures, direction in words
\newcommand{\MlstSculptMidOverFullTen}{1.069\,\ensuremath{\pm}\,0.020}  % experiments/FIGS/data/v2/index.json :: distance at QCD efficiency 0.1: R42_Q1 over L188, run by run :: sha256 c6c4138ee5c86423
\newcommand{\MlstSculptMidOverFullTenPr}{1.1 times the distance}  % experiments/FIGS/data/v2/index.json :: distance at QCD efficiency 0.1: R42_Q1 over L188, run by run :: sha256 c6c4138ee5c86423 :: prose twin: two significant figures, direction in words
\newcommand{\MlstSculptCoarseOverMidTen}{0.971\,\ensuremath{\pm}\,0.035}  % experiments/FIGS/data/v2/index.json :: distance at QCD efficiency 0.1: R16_Q1 over R42_Q1, run by run :: sha256 c6c4138ee5c86423
\newcommand{\MlstSculptCoarseOverMidTenPr}{equal within 4\%}  % experiments/FIGS/data/v2/index.json :: distance at QCD efficiency 0.1: R16_Q1 over R42_Q1, run by run :: sha256 c6c4138ee5c86423 :: prose twin: two significant figures, direction in words
\newcommand{\MlstSculptNPassOne}{4{,}103}  % experiments/FIGS/data/v2/index.json :: QCD test jets passing the cut at QCD efficiency 0.01, every run :: sha256 c6c4138ee5c86423
\newcommand{\MlstSculptFloorOne}{0.0313}  % experiments/FIGS/data/v2/index.json :: Jensen-Shannon distance expected for a random selection of as many QCD jets as pass at 0.01 :: sha256 c6c4138ee5c86423
\newcommand{\MlstSculptFullOne}{0.0439\,\ensuremath{\pm}\,0.0024}  % experiments/FIGS/data/v2/index.json :: Jensen-Shannon distance of the QCD m_SD distribution at QCD efficiency 0.01, L188 :: sha256 c6c4138ee5c86423
\newcommand{\MlstSculptFullOverFloorOne}{1.40\,\ensuremath{\pm}\,0.08}  % experiments/FIGS/data/v2/index.json :: distance at 0.01 over the random-selection floor, L188 :: sha256 c6c4138ee5c86423
\newcommand{\MlstSculptFullOverFloorOnePr}{1.4 times the distance}  % experiments/FIGS/data/v2/index.json :: distance at 0.01 over the random-selection floor, L188 :: sha256 c6c4138ee5c86423 :: prose twin: two significant figures, direction in words
\newcommand{\MlstSculptResOne}{0.0429\,\ensuremath{\pm}\,0.0029}  % experiments/FIGS/data/v2/index.json :: Jensen-Shannon distance of the QCD m_SD distribution at QCD efficiency 0.01, L162 :: sha256 c6c4138ee5c86423
\newcommand{\MlstSculptResOverFloorOne}{1.37\,\ensuremath{\pm}\,0.09}  % experiments/FIGS/data/v2/index.json :: distance at 0.01 over the random-selection floor, L162 :: sha256 c6c4138ee5c86423
\newcommand{\MlstSculptResOverFloorOnePr}{1.4 times the distance}  % experiments/FIGS/data/v2/index.json :: distance at 0.01 over the random-selection floor, L162 :: sha256 c6c4138ee5c86423 :: prose twin: two significant figures, direction in words
\newcommand{\MlstSculptMidOne}{0.051\,\ensuremath{\pm}\,0.007}  % experiments/FIGS/data/v2/index.json :: Jensen-Shannon distance of the QCD m_SD distribution at QCD efficiency 0.01, R42_Q1 :: sha256 c6c4138ee5c86423
\newcommand{\MlstSculptMidOverFloorOne}{1.62\,\ensuremath{\pm}\,0.23}  % experiments/FIGS/data/v2/index.json :: distance at 0.01 over the random-selection floor, R42_Q1 :: sha256 c6c4138ee5c86423
\newcommand{\MlstSculptMidOverFloorOnePr}{1.6 times the distance}  % experiments/FIGS/data/v2/index.json :: distance at 0.01 over the random-selection floor, R42_Q1 :: sha256 c6c4138ee5c86423 :: prose twin: two significant figures, direction in words
\newcommand{\MlstSculptCoarseOne}{0.0561\,\ensuremath{\pm}\,0.0017}  % experiments/FIGS/data/v2/index.json :: Jensen-Shannon distance of the QCD m_SD distribution at QCD efficiency 0.01, R16_Q1 :: sha256 c6c4138ee5c86423
\newcommand{\MlstSculptCoarseOverFloorOne}{1.79\,\ensuremath{\pm}\,0.06}  % experiments/FIGS/data/v2/index.json :: distance at 0.01 over the random-selection floor, R16_Q1 :: sha256 c6c4138ee5c86423
\newcommand{\MlstSculptCoarseOverFloorOnePr}{1.8 times the distance}  % experiments/FIGS/data/v2/index.json :: distance at 0.01 over the random-selection floor, R16_Q1 :: sha256 c6c4138ee5c86423 :: prose twin: two significant figures, direction in words
\newcommand{\MlstSculptResMassOne}{0.0424 and 0.0442}  % experiments/FIGS/data/v2/index.json :: Jensen-Shannon distance of the QCD m_SD distribution at QCD efficiency 0.01, L162_MASS :: sha256 c6c4138ee5c86423
\newcommand{\MlstSculptResMassOverFloorOne}{1.356 and 1.411}  % experiments/FIGS/data/v2/index.json :: distance at 0.01 over the random-selection floor, L162_MASS :: sha256 c6c4138ee5c86423
\newcommand{\MlstSculptResMassOverFloorOnePr}{at least 1.4 times the distance}  % experiments/FIGS/data/v2/index.json :: distance at 0.01 over the random-selection floor, L162_MASS :: sha256 c6c4138ee5c86423 :: prose twin: two significant figures, direction in words
\newcommand{\MlstSculptResMassOverFloorOnePrMax}{up to 1.4 times the distance}  % experiments/FIGS/data/v2/index.json :: distance at 0.01 over the random-selection floor, L162_MASS :: sha256 c6c4138ee5c86423 :: prose twin of the run farther from 1
\newcommand{\MlstSculptCoarseMassOne}{0.0405 and 0.0452}  % experiments/FIGS/data/v2/index.json :: Jensen-Shannon distance of the QCD m_SD distribution at QCD efficiency 0.01, R16_Q1_MASS :: sha256 c6c4138ee5c86423
\newcommand{\MlstSculptCoarseMassOverFloorOne}{1.29 and 1.44}  % experiments/FIGS/data/v2/index.json :: distance at 0.01 over the random-selection floor, R16_Q1_MASS :: sha256 c6c4138ee5c86423
\newcommand{\MlstSculptCoarseMassOverFloorOnePr}{at least 1.3 times the distance}  % experiments/FIGS/data/v2/index.json :: distance at 0.01 over the random-selection floor, R16_Q1_MASS :: sha256 c6c4138ee5c86423 :: prose twin: two significant figures, direction in words
\newcommand{\MlstSculptCoarseMassOverFloorOnePrMax}{up to 1.4 times the distance}  % experiments/FIGS/data/v2/index.json :: distance at 0.01 over the random-selection floor, R16_Q1_MASS :: sha256 c6c4138ee5c86423 :: prose twin of the run farther from 1
\newcommand{\MlstSculptResMassOverResOne}{0.962 and 0.983}  % experiments/FIGS/data/v2/index.json :: distance at QCD efficiency 0.01: L162_MASS over L162, run by run :: sha256 c6c4138ee5c86423
\newcommand{\MlstSculptResMassOverResOnePr}{equal within 4\%}  % experiments/FIGS/data/v2/index.json :: distance at QCD efficiency 0.01: L162_MASS over L162, run by run :: sha256 c6c4138ee5c86423 :: prose twin: two significant figures, direction in words
\newcommand{\MlstSculptResMassOverResOnePrMax}{equal within 4\%}  % experiments/FIGS/data/v2/index.json :: distance at QCD efficiency 0.01: L162_MASS over L162, run by run :: sha256 c6c4138ee5c86423 :: prose twin of the run farther from 1
\newcommand{\MlstSculptCoarseMassOverCoarseOne}{0.70 and 0.82}  % experiments/FIGS/data/v2/index.json :: distance at QCD efficiency 0.01: R16_Q1_MASS over R16_Q1, run by run :: sha256 c6c4138ee5c86423
\newcommand{\MlstSculptCoarseMassOverCoarseOnePr}{at most 0.82 times the distance}  % experiments/FIGS/data/v2/index.json :: distance at QCD efficiency 0.01: R16_Q1_MASS over R16_Q1, run by run :: sha256 c6c4138ee5c86423 :: prose twin: two significant figures, direction in words
\newcommand{\MlstSculptCoarseMassOverCoarseOnePrMax}{down to 0.70 times the distance}  % experiments/FIGS/data/v2/index.json :: distance at QCD efficiency 0.01: R16_Q1_MASS over R16_Q1, run by run :: sha256 c6c4138ee5c86423 :: prose twin of the run farther from 1
\newcommand{\MlstSculptCoarseOverFullOne}{1.28\,\ensuremath{\pm}\,0.04}  % experiments/FIGS/data/v2/index.json :: distance at QCD efficiency 0.01: R16_Q1 over L188, run by run :: sha256 c6c4138ee5c86423
\newcommand{\MlstSculptCoarseOverFullOnePr}{1.3 times the distance}  % experiments/FIGS/data/v2/index.json :: distance at QCD efficiency 0.01: R16_Q1 over L188, run by run :: sha256 c6c4138ee5c86423 :: prose twin: two significant figures, direction in words
\newcommand{\MlstSculptResOverFullOne}{0.98\,\ensuremath{\pm}\,0.05}  % experiments/FIGS/data/v2/index.json :: distance at QCD efficiency 0.01: L162 over L188, run by run :: sha256 c6c4138ee5c86423
\newcommand{\MlstSculptResOverFullOnePr}{equal within 5\%}  % experiments/FIGS/data/v2/index.json :: distance at QCD efficiency 0.01: L162 over L188, run by run :: sha256 c6c4138ee5c86423 :: prose twin: two significant figures, direction in words
\newcommand{\MlstSculptMidOverFullOne}{1.15\,\ensuremath{\pm}\,0.12}  % experiments/FIGS/data/v2/index.json :: distance at QCD efficiency 0.01: R42_Q1 over L188, run by run :: sha256 c6c4138ee5c86423
\newcommand{\MlstSculptMidOverFullOnePr}{equal within 20\%}  % experiments/FIGS/data/v2/index.json :: distance at QCD efficiency 0.01: R42_Q1 over L188, run by run :: sha256 c6c4138ee5c86423 :: prose twin: two significant figures, direction in words
\newcommand{\MlstSculptCoarseOverMidOne}{1.13\,\ensuremath{\pm}\,0.16}  % experiments/FIGS/data/v2/index.json :: distance at QCD efficiency 0.01: R16_Q1 over R42_Q1, run by run :: sha256 c6c4138ee5c86423
\newcommand{\MlstSculptCoarseOverMidOnePr}{equal within 20\%}  % experiments/FIGS/data/v2/index.json :: distance at QCD efficiency 0.01: R16_Q1 over R42_Q1, run by run :: sha256 c6c4138ee5c86423 :: prose twin: two significant figures, direction in words
\newcommand{\MlstSculptDiffEffLo}{0.52\%}  % experiments/FIGS/data/v2/index.json :: lowest stored QCD efficiency of the unbroken stretch around the tighter working point at which the 17/188 distance ratio exceeds 1 by twice its SD :: sha256 c6c4138ee5c86423 :: passing grid points: 0.001, 0.005245, 0.007937, 0.012011
\newcommand{\MlstSculptDiffEffHi}{1.2\%}  % experiments/FIGS/data/v2/index.json :: highest efficiency of that stretch :: sha256 c6c4138ee5c86423
\newcommand{\MlstSculptLowestEff}{0.10\%}  % experiments/FIGS/data/v2/index.json :: lowest stored QCD efficiency of the sculpting curves :: sha256 c6c4138ee5c86423
\newcommand{\MlstSculptLowestNPass}{410}  % experiments/FIGS/data/v2/index.json :: QCD test jets passing at the lowest stored efficiency :: sha256 c6c4138ee5c86423
\newcommand{\MlstSculptLowestMaxFloorSd}{1.7}  % experiments/FIGS/data/v2/index.json :: largest distance of a configuration's mean from the random-selection floor at the lowest stored efficiency, in SDs of the floor over random selections :: sha256 c6c4138ee5c86423
\newcommand{\MlstSculptNConfigs}{6}  % experiments/FIGS/data/v2/index.json :: configurations with a sculpting curve (count) :: sha256 c6c4138ee5c86423
\newcommand{\LpJciiTestQcdPercent}{20.5\%}  % experiments/FIGS/data/v2/jc2_test/counts.json :: QCD share of the jets, per cent :: sha256 3574ba1f6f2a59f0
\newcommand{\LpJciiTestTwoProngPercent}{38.7\%}  % experiments/FIGS/data/v2/jc2_test/counts.json :: two-prong share of the jets, per cent :: sha256 3574ba1f6f2a59f0
\newcommand{\LpJciiTrainQcdPercent}{13.5\%}  % experiments/FT/data/jc2_v2_manifest.json :: QCD share of the jets, per cent :: sha256 4bd68655919b14e9
\newcommand{\LpJciiTrainTwoProngPercent}{15.4\%}  % experiments/FT/data/jc2_v2_manifest.json :: two-prong share of the jets, per cent :: sha256 4bd68655919b14e9
\newcommand{\LpJciiTestMinAucJets}{10}  % experiments/FIGS/data/v2/jc2_test/counts.json :: fewest jets of a class among the macro-AUC rows :: sha256 3574ba1f6f2a59f0
\newcommand{\LpJciiTrainMinJets}{11}  % experiments/FT/data/jc2_v2_manifest.json :: fewest jets of a class in the 10^6 training subset :: sha256 4bd68655919b14e9
\newcommand{\LpNTestJci}{2{,}000{,}000}  % experiments/FIGS/data/v2/index.json :: jc1 test jets the probes score :: sha256 c6c4138ee5c86423
\newcommand{\LpNTrainMaxJci}{1{,}000{,}000}  % experiments/FIGS/data/v2/index.json :: jc1 largest training subset :: sha256 c6c4138ee5c86423
\newcommand{\LpNTestTop}{404{,}000}  % experiments/FIGS/data/v2/index.json :: top test jets the probes score :: sha256 c6c4138ee5c86423
\newcommand{\LpNTrainMaxTop}{1{,}200{,}000}  % experiments/FIGS/data/v2/index.json :: top largest training subset :: sha256 c6c4138ee5c86423
\newcommand{\LpNTestQg}{200{,}000}  % experiments/FIGS/data/v2/index.json :: qg test jets the probes score :: sha256 c6c4138ee5c86423
\newcommand{\LpNTrainMaxQg}{1{,}600{,}000}  % experiments/FIGS/data/v2/index.json :: qg largest training subset :: sha256 c6c4138ee5c86423
\newcommand{\LpNTestAucJci}{2{,}000{,}000}  % experiments/FIGS/data/v2/index.json :: JetClass test jets of the macro AUC :: sha256 c6c4138ee5c86423
\newcommand{\LpPairNValMin}{1{,}068}  % experiments/FIGS/data/v2/index.json :: validation jets of the pair probe with the fewest :: sha256 c6c4138ee5c86423
\newcommand{\LpPairNValMax}{20{,}195}  % experiments/FIGS/data/v2/index.json :: validation jets of the pair probe with the most :: sha256 c6c4138ee5c86423
\newcommand{\LpNValJci}{200{,}000}  % experiments/FIGS/data/v2/index.json :: JetClass probe validation jets :: sha256 c6c4138ee5c86423
\newcommand{\SculptEffHigh}{10\%}  % experiments/EVAL/mass_sculpting.py :: WORKING_POINTS :: sha256 e85d125c85dac054 :: QCD efficiency of a sculpting cut
\newcommand{\SculptEffLow}{1\%}  % experiments/EVAL/mass_sculpting.py :: WORKING_POINTS :: sha256 e85d125c85dac054 :: QCD efficiency of a sculpting cut
\newcommand{\LpMlpHidden}{256}  % experiments/EVAL/probe.py :: fit_mlp hidden layer :: sha256 485db0966e660588 :: units of the MLP probe's hidden layer
\newcommand{\LpMlpSeedsWord}{three}  % experiments/EVAL/probe.py :: MLP_SEEDS :: sha256 485db0966e660588 :: seeds of the MLP probe whose scores are averaged
\newcommand{\MassMinTrainPerClass}{50}  % experiments/EVAL/mass_resolution.py :: MIN_TRAIN_PER_CLASS :: sha256 26acb0e315094bd1 :: fewest training jets of a class the mass regression keeps
\newcommand{\MlstFtLr}{\ensuremath{10^{-4}}}  % experiments/FIGS/data/v2/ft_mlst/manifests.json :: cells[best70].lr :: sha256 a689f21e894f48ae
\newcommand{\MlstFtHeadMult}{50}  % experiments/FIGS/data/v2/ft_mlst/manifests.json :: cells[best70].head_lr_mult :: sha256 a689f21e894f48ae
\newcommand{\MlstFtWeightDecay}{0.01}  % experiments/FIGS/data/v2/ft_mlst/manifests.json :: cells[best70].weight_decay :: sha256 a689f21e894f48ae
\newcommand{\MlstFtBatch}{512}  % experiments/FIGS/data/v2/ft_mlst/manifests.json :: cells[best70].batch_size :: sha256 a689f21e894f48ae
\newcommand{\MlstFtBenchEpochs}{20}  % experiments/FIGS/data/v2/ft_mlst/manifests.json :: cells[best70, top/qg].epochs :: sha256 a689f21e894f48ae
\newcommand{\MlstFtEpochsEThree}{50}  % experiments/FIGS/data/v2/ft_mlst/manifests.json :: cells[best70, leg 1/2, n_train 1000].epochs :: sha256 a689f21e894f48ae
\newcommand{\MlstFtEpochsEFour}{50}  % experiments/FIGS/data/v2/ft_mlst/manifests.json :: cells[best70, leg 1/2, n_train 10000].epochs :: sha256 a689f21e894f48ae
\newcommand{\MlstFtEpochsEFive}{30}  % experiments/FIGS/data/v2/ft_mlst/manifests.json :: cells[best70, leg 1/2, n_train 100000].epochs :: sha256 a689f21e894f48ae
\newcommand{\MlstFtEpochsESix}{10}  % experiments/FIGS/data/v2/ft_mlst/manifests.json :: cells[best70, leg 1/2, n_train 1000000].epochs :: sha256 a689f21e894f48ae
\newcommand{\MlstFtHeadDecayShare}{30\%}  % experiments/FT/k8s/job-ft-v2-legs-best70-l188-s3-raunav.yaml :: no --lr-scheduler: weaver 0.4.17 flat+decay, the head's last epochs :: sha256 19f59aa54a385c11
\newcommand{\MlstFtHeadDecayTo}{1\%}  % experiments/FT/k8s/job-ft-v2-legs-best70-l188-s3-raunav.yaml :: no --lr-scheduler: weaver 0.4.17 flat+decay, final factor :: sha256 19f59aa54a385c11
\newcommand{\MlstFtHeadDropout}{0.1}  % experiments/FT/k8s/job-ft-v2-legs-best70-l188-s3-raunav.yaml :: fc_params dropout of the fine-tuning head :: sha256 19f59aa54a385c11
\newcommand{\MlstFtMinEpochJets}{10{,}000}  % experiments/FIGS/data/v2/ft_mlst/manifests.json :: cells[best70].samples_per_epoch (smallest) :: sha256 a689f21e894f48ae
\newcommand{\MlstFtValJcii}{20{,}480}  % experiments/FIGS/data/v2/ft_mlst/manifests.json :: cells[best70, leg 1].val_jets :: sha256 a689f21e894f48ae
\newcommand{\MlstFtValSmall}{20{,}000}  % experiments/FT/k8s/job-ft-v2-legs-best70-l188-s3-raunav.yaml :: JetClass --samples-per-epoch-val; top and q/g at N <= 1e4 :: sha256 19f59aa54a385c11
\newcommand{\MlstFtValLarge}{200{,}000}  % experiments/FT/k8s/job-ft-v2-bench-best70-l188-s1-top-raunav.yaml :: top and q/g --samples-per-epoch-val at N > 1e4 :: sha256 1e38d4f542119ea3
\newcommand{\MlstFtGpuTypes}{V100}  % experiments/FIGS/data/v2/ft_mlst/manifests.json :: cells[best70].gpu_device_name :: sha256 a689f21e894f48ae
\newcommand{\MlstFtCellsDone}{445}  % experiments/FIGS/data/v2/ft_mlst/manifests.json :: DONE paper-run fine-tuning cells (count) :: sha256 a689f21e894f48ae
\newcommand{\MlstFtRefLr}{\ensuremath{10^{-3}}}  % experiments/FIGS/data/v2/ft_mlst/manifests.json :: cells[ref_v100].lr :: sha256 a689f21e894f48ae
\newcommand{\MlstFtRefSeedsWord}{three}  % experiments/FIGS/data/v2/ft_mlst/manifests.json :: cells[ref_v100].ft_seed (count) :: sha256 a689f21e894f48ae
\newcommand{\MlstFtRefGpuTypes}{V100}  % experiments/FIGS/data/v2/ft_mlst/manifests.json :: cells[ref_v100].gpu_device_name :: sha256 a689f21e894f48ae
\newcommand{\MlstFtRefWeightDecayBench}{0.01}  % experiments/FIGS/data/v2/ft_mlst/manifests.json :: cells[ref_v100, top/qg].weight_decay :: sha256 a689f21e894f48ae
\newcommand{\MlstBootResamples}{200}  % experiments/FIGS/data/v2/index.json :: bootstrap resamples of the test jets :: sha256 c6c4138ee5c86423
\newcommand{\MlstBootTestMaxProbe}{0.21}  % experiments/FIGS/data/v2/index.json :: largest test-sample SD of the ln ratio of 1 - AUC on JetClass-II, linear probe, label sets over the 188-class model, per cent :: sha256 c6c4138ee5c86423
\newcommand{\MlstBootTestMaxFt}{0.15}  % experiments/FIGS/data/v2/index.json :: largest test-sample SD of the ln ratio of 1 - AUC on JetClass-II, fine-tuning, label sets over the 188-class model, per cent :: sha256 c6c4138ee5c86423
\newcommand{\MlstBootTestMaxRthirtyProbe}{3.6}  % experiments/FIGS/data/v2/index.json :: largest test-sample SD of the ln R30 ratio on top and quark/gluon, linear probe (counting), label sets over the 188-class model, per cent :: sha256 c6c4138ee5c86423
\newcommand{\MlstBootTestMaxRthirtyFt}{13}  % experiments/FIGS/data/v2/index.json :: largest test-sample SD of the ln R30 ratio on top and quark/gluon, fine-tuning (counting), label sets over the 188-class model, per cent :: sha256 c6c4138ee5c86423
\newcommand{\MlstFtCmpJcTwoResOverFullEFour}{0.989\,\ensuremath{\pm}\,0.011}  % experiments/FIGS/data/v2/index.json :: fine-tuning jc2 N=10000: (1 - macro_auc_ovr) of L162 over L188, run by run :: sha256 c6c4138ee5c86423
\newcommand{\MlstFtCmpJcTwoResOverFullEFourPr}{equal within 2\%}  % experiments/FIGS/data/v2/index.json :: fine-tuning jc2 N=10000: (1 - macro_auc_ovr) of L162 over L188, run by run :: sha256 c6c4138ee5c86423 :: prose twin: two significant figures, direction in words
\newcommand{\MlstFtCmpJcTwoMidOverFullEFour}{1.191\,\ensuremath{\pm}\,0.030}  % experiments/FIGS/data/v2/index.json :: fine-tuning jc2 N=10000: (1 - macro_auc_ovr) of R42_Q1 over L188, run by run :: sha256 c6c4138ee5c86423
\newcommand{\MlstFtCmpJcTwoMidOverFullEFourPr}{19\% more errors}  % experiments/FIGS/data/v2/index.json :: fine-tuning jc2 N=10000: (1 - macro_auc_ovr) of R42_Q1 over L188, run by run :: sha256 c6c4138ee5c86423 :: prose twin: two significant figures, direction in words
\newcommand{\MlstFtCmpJcTwoCoarseOverFullEFour}{1.353\,\ensuremath{\pm}\,0.003}  % experiments/FIGS/data/v2/index.json :: fine-tuning jc2 N=10000: (1 - macro_auc_ovr) of R16_Q1 over L188, run by run :: sha256 c6c4138ee5c86423
\newcommand{\MlstFtCmpJcTwoCoarseOverFullEFourPr}{35\% more errors}  % experiments/FIGS/data/v2/index.json :: fine-tuning jc2 N=10000: (1 - macro_auc_ovr) of R16_Q1 over L188, run by run :: sha256 c6c4138ee5c86423 :: prose twin: two significant figures, direction in words
\newcommand{\MlstFtCmpJcTwoResMassOverResEFour}{1.016 and 1.005}  % experiments/FIGS/data/v2/index.json :: fine-tuning jc2 N=10000: (1 - macro_auc_ovr) of L162_MASS over L162, run by run :: sha256 c6c4138ee5c86423
\newcommand{\MlstFtCmpJcTwoResMassOverResEFourPr}{equal within 2\%}  % experiments/FIGS/data/v2/index.json :: fine-tuning jc2 N=10000: (1 - macro_auc_ovr) of L162_MASS over L162, run by run :: sha256 c6c4138ee5c86423 :: prose twin: two significant figures, direction in words
\newcommand{\MlstFtCmpJcTwoResMassOverResEFourPrMax}{equal within 2\%}  % experiments/FIGS/data/v2/index.json :: fine-tuning jc2 N=10000: (1 - macro_auc_ovr) of L162_MASS over L162, run by run :: sha256 c6c4138ee5c86423 :: prose twin of the run farther from 1
\newcommand{\MlstFtCmpJcTwoCoarseMassOverCoarseEFour}{0.989 and 1.000}  % experiments/FIGS/data/v2/index.json :: fine-tuning jc2 N=10000: (1 - macro_auc_ovr) of R16_Q1_MASS over R16_Q1, run by run :: sha256 c6c4138ee5c86423
\newcommand{\MlstFtCmpJcTwoCoarseMassOverCoarseEFourPr}{equal within 2\%}  % experiments/FIGS/data/v2/index.json :: fine-tuning jc2 N=10000: (1 - macro_auc_ovr) of R16_Q1_MASS over R16_Q1, run by run :: sha256 c6c4138ee5c86423 :: prose twin: two significant figures, direction in words
\newcommand{\MlstFtCmpJcTwoCoarseMassOverCoarseEFourPrMax}{equal within 2\%}  % experiments/FIGS/data/v2/index.json :: fine-tuning jc2 N=10000: (1 - macro_auc_ovr) of R16_Q1_MASS over R16_Q1, run by run :: sha256 c6c4138ee5c86423 :: prose twin of the run farther from 1
\newcommand{\MlstFtCmpJcTwoResOverFullEFive}{1.011\,\ensuremath{\pm}\,0.005}  % experiments/FIGS/data/v2/index.json :: fine-tuning jc2 N=100000: (1 - macro_auc_ovr) of L162 over L188, run by run :: sha256 c6c4138ee5c86423
\newcommand{\MlstFtCmpJcTwoResOverFullEFivePr}{1.1\% more errors}  % experiments/FIGS/data/v2/index.json :: fine-tuning jc2 N=100000: (1 - macro_auc_ovr) of L162 over L188, run by run :: sha256 c6c4138ee5c86423 :: prose twin: two significant figures, direction in words
\newcommand{\MlstFtCmpJcTwoMidOverFullEFive}{1.064\,\ensuremath{\pm}\,0.009}  % experiments/FIGS/data/v2/index.json :: fine-tuning jc2 N=100000: (1 - macro_auc_ovr) of R42_Q1 over L188, run by run :: sha256 c6c4138ee5c86423
\newcommand{\MlstFtCmpJcTwoMidOverFullEFivePr}{6.4\% more errors}  % experiments/FIGS/data/v2/index.json :: fine-tuning jc2 N=100000: (1 - macro_auc_ovr) of R42_Q1 over L188, run by run :: sha256 c6c4138ee5c86423 :: prose twin: two significant figures, direction in words
\newcommand{\MlstFtCmpJcTwoCoarseOverFullEFive}{1.157\,\ensuremath{\pm}\,0.001}  % experiments/FIGS/data/v2/index.json :: fine-tuning jc2 N=100000: (1 - macro_auc_ovr) of R16_Q1 over L188, run by run :: sha256 c6c4138ee5c86423
\newcommand{\MlstFtCmpJcTwoCoarseOverFullEFivePr}{16\% more errors}  % experiments/FIGS/data/v2/index.json :: fine-tuning jc2 N=100000: (1 - macro_auc_ovr) of R16_Q1 over L188, run by run :: sha256 c6c4138ee5c86423 :: prose twin: two significant figures, direction in words
\newcommand{\MlstFtCmpJcTwoResMassOverResEFive}{1.027 and 1.021}  % experiments/FIGS/data/v2/index.json :: fine-tuning jc2 N=100000: (1 - macro_auc_ovr) of L162_MASS over L162, run by run :: sha256 c6c4138ee5c86423
\newcommand{\MlstFtCmpJcTwoResMassOverResEFivePr}{at least 2.1\% more errors}  % experiments/FIGS/data/v2/index.json :: fine-tuning jc2 N=100000: (1 - macro_auc_ovr) of L162_MASS over L162, run by run :: sha256 c6c4138ee5c86423 :: prose twin: two significant figures, direction in words
\newcommand{\MlstFtCmpJcTwoResMassOverResEFivePrMax}{up to 2.7\% more errors}  % experiments/FIGS/data/v2/index.json :: fine-tuning jc2 N=100000: (1 - macro_auc_ovr) of L162_MASS over L162, run by run :: sha256 c6c4138ee5c86423 :: prose twin of the run farther from 1
\newcommand{\MlstFtCmpJcTwoCoarseMassOverCoarseEFive}{1.029 and 1.022}  % experiments/FIGS/data/v2/index.json :: fine-tuning jc2 N=100000: (1 - macro_auc_ovr) of R16_Q1_MASS over R16_Q1, run by run :: sha256 c6c4138ee5c86423
\newcommand{\MlstFtCmpJcTwoCoarseMassOverCoarseEFivePr}{at least 2.2\% more errors}  % experiments/FIGS/data/v2/index.json :: fine-tuning jc2 N=100000: (1 - macro_auc_ovr) of R16_Q1_MASS over R16_Q1, run by run :: sha256 c6c4138ee5c86423 :: prose twin: two significant figures, direction in words
\newcommand{\MlstFtCmpJcTwoCoarseMassOverCoarseEFivePrMax}{up to 2.9\% more errors}  % experiments/FIGS/data/v2/index.json :: fine-tuning jc2 N=100000: (1 - macro_auc_ovr) of R16_Q1_MASS over R16_Q1, run by run :: sha256 c6c4138ee5c86423 :: prose twin of the run farther from 1
\newcommand{\MlstFtCmpJcTwoResOverFullESix}{1.014\,\ensuremath{\pm}\,0.011}  % experiments/FIGS/data/v2/index.json :: fine-tuning jc2 N=1000000: (1 - macro_auc_ovr) of L162 over L188, run by run :: sha256 c6c4138ee5c86423
\newcommand{\MlstFtCmpJcTwoResOverFullESixPr}{equal within 2\%}  % experiments/FIGS/data/v2/index.json :: fine-tuning jc2 N=1000000: (1 - macro_auc_ovr) of L162 over L188, run by run :: sha256 c6c4138ee5c86423 :: prose twin: two significant figures, direction in words
\newcommand{\MlstFtCmpJcTwoMidOverFullESix}{1.043\,\ensuremath{\pm}\,0.002}  % experiments/FIGS/data/v2/index.json :: fine-tuning jc2 N=1000000: (1 - macro_auc_ovr) of R42_Q1 over L188, run by run :: sha256 c6c4138ee5c86423
\newcommand{\MlstFtCmpJcTwoMidOverFullESixPr}{4.3\% more errors}  % experiments/FIGS/data/v2/index.json :: fine-tuning jc2 N=1000000: (1 - macro_auc_ovr) of R42_Q1 over L188, run by run :: sha256 c6c4138ee5c86423 :: prose twin: two significant figures, direction in words
\newcommand{\MlstFtCmpJcTwoCoarseOverFullESix}{1.123\,\ensuremath{\pm}\,0.003}  % experiments/FIGS/data/v2/index.json :: fine-tuning jc2 N=1000000: (1 - macro_auc_ovr) of R16_Q1 over L188, run by run :: sha256 c6c4138ee5c86423
\newcommand{\MlstFtCmpJcTwoCoarseOverFullESixPr}{12\% more errors}  % experiments/FIGS/data/v2/index.json :: fine-tuning jc2 N=1000000: (1 - macro_auc_ovr) of R16_Q1 over L188, run by run :: sha256 c6c4138ee5c86423 :: prose twin: two significant figures, direction in words
\newcommand{\MlstFtCmpJcTwoResMassOverResESix}{1.009 and 1.016}  % experiments/FIGS/data/v2/index.json :: fine-tuning jc2 N=1000000: (1 - macro_auc_ovr) of L162_MASS over L162, run by run :: sha256 c6c4138ee5c86423
\newcommand{\MlstFtCmpJcTwoResMassOverResESixPr}{equal within 2\%}  % experiments/FIGS/data/v2/index.json :: fine-tuning jc2 N=1000000: (1 - macro_auc_ovr) of L162_MASS over L162, run by run :: sha256 c6c4138ee5c86423 :: prose twin: two significant figures, direction in words
\newcommand{\MlstFtCmpJcTwoResMassOverResESixPrMax}{equal within 2\%}  % experiments/FIGS/data/v2/index.json :: fine-tuning jc2 N=1000000: (1 - macro_auc_ovr) of L162_MASS over L162, run by run :: sha256 c6c4138ee5c86423 :: prose twin of the run farther from 1
\newcommand{\MlstFtCmpJcTwoCoarseMassOverCoarseESix}{1.023 and 1.017}  % experiments/FIGS/data/v2/index.json :: fine-tuning jc2 N=1000000: (1 - macro_auc_ovr) of R16_Q1_MASS over R16_Q1, run by run :: sha256 c6c4138ee5c86423
\newcommand{\MlstFtCmpJcTwoCoarseMassOverCoarseESixPr}{at least 1.7\% more errors}  % experiments/FIGS/data/v2/index.json :: fine-tuning jc2 N=1000000: (1 - macro_auc_ovr) of R16_Q1_MASS over R16_Q1, run by run :: sha256 c6c4138ee5c86423 :: prose twin: two significant figures, direction in words
\newcommand{\MlstFtCmpJcTwoCoarseMassOverCoarseESixPrMax}{up to 2.3\% more errors}  % experiments/FIGS/data/v2/index.json :: fine-tuning jc2 N=1000000: (1 - macro_auc_ovr) of R16_Q1_MASS over R16_Q1, run by run :: sha256 c6c4138ee5c86423 :: prose twin of the run farther from 1
\newcommand{\MlstFtCmpJcTwoQcdResOverFullEFour}{1.038\,\ensuremath{\pm}\,0.005}  % experiments/FIGS/data/v2/index.json :: fine-tuning jc2_qcd N=10000: (1 - macro_auc_qcd) of L162 over L188, run by run :: sha256 c6c4138ee5c86423
\newcommand{\MlstFtCmpJcTwoQcdResOverFullEFourPr}{3.8\% more errors}  % experiments/FIGS/data/v2/index.json :: fine-tuning jc2_qcd N=10000: (1 - macro_auc_qcd) of L162 over L188, run by run :: sha256 c6c4138ee5c86423 :: prose twin: two significant figures, direction in words
\newcommand{\MlstFtCmpJcTwoQcdMidOverFullEFour}{1.164\,\ensuremath{\pm}\,0.018}  % experiments/FIGS/data/v2/index.json :: fine-tuning jc2_qcd N=10000: (1 - macro_auc_qcd) of R42_Q1 over L188, run by run :: sha256 c6c4138ee5c86423
\newcommand{\MlstFtCmpJcTwoQcdMidOverFullEFourPr}{16\% more errors}  % experiments/FIGS/data/v2/index.json :: fine-tuning jc2_qcd N=10000: (1 - macro_auc_qcd) of R42_Q1 over L188, run by run :: sha256 c6c4138ee5c86423 :: prose twin: two significant figures, direction in words
\newcommand{\MlstFtCmpJcTwoQcdCoarseOverFullEFour}{1.258\,\ensuremath{\pm}\,0.025}  % experiments/FIGS/data/v2/index.json :: fine-tuning jc2_qcd N=10000: (1 - macro_auc_qcd) of R16_Q1 over L188, run by run :: sha256 c6c4138ee5c86423
\newcommand{\MlstFtCmpJcTwoQcdCoarseOverFullEFourPr}{26\% more errors}  % experiments/FIGS/data/v2/index.json :: fine-tuning jc2_qcd N=10000: (1 - macro_auc_qcd) of R16_Q1 over L188, run by run :: sha256 c6c4138ee5c86423 :: prose twin: two significant figures, direction in words
\newcommand{\MlstFtCmpJcTwoQcdResMassOverResEFour}{1.027 and 0.997}  % experiments/FIGS/data/v2/index.json :: fine-tuning jc2_qcd N=10000: (1 - macro_auc_qcd) of L162_MASS over L162, run by run :: sha256 c6c4138ee5c86423
\newcommand{\MlstFtCmpJcTwoQcdResMassOverResEFourPr}{equal within 3\%}  % experiments/FIGS/data/v2/index.json :: fine-tuning jc2_qcd N=10000: (1 - macro_auc_qcd) of L162_MASS over L162, run by run :: sha256 c6c4138ee5c86423 :: prose twin: two significant figures, direction in words
\newcommand{\MlstFtCmpJcTwoQcdResMassOverResEFourPrMax}{equal within 3\%}  % experiments/FIGS/data/v2/index.json :: fine-tuning jc2_qcd N=10000: (1 - macro_auc_qcd) of L162_MASS over L162, run by run :: sha256 c6c4138ee5c86423 :: prose twin of the run farther from 1
\newcommand{\MlstFtCmpJcTwoQcdCoarseMassOverCoarseEFour}{1.028 and 1.048}  % experiments/FIGS/data/v2/index.json :: fine-tuning jc2_qcd N=10000: (1 - macro_auc_qcd) of R16_Q1_MASS over R16_Q1, run by run :: sha256 c6c4138ee5c86423
\newcommand{\MlstFtCmpJcTwoQcdCoarseMassOverCoarseEFourPr}{equal within 5\%}  % experiments/FIGS/data/v2/index.json :: fine-tuning jc2_qcd N=10000: (1 - macro_auc_qcd) of R16_Q1_MASS over R16_Q1, run by run :: sha256 c6c4138ee5c86423 :: prose twin: two significant figures, direction in words
\newcommand{\MlstFtCmpJcTwoQcdCoarseMassOverCoarseEFourPrMax}{equal within 5\%}  % experiments/FIGS/data/v2/index.json :: fine-tuning jc2_qcd N=10000: (1 - macro_auc_qcd) of R16_Q1_MASS over R16_Q1, run by run :: sha256 c6c4138ee5c86423 :: prose twin of the run farther from 1
\newcommand{\MlstFtCmpJcTwoQcdResOverFullEFive}{1.060\,\ensuremath{\pm}\,0.002}  % experiments/FIGS/data/v2/index.json :: fine-tuning jc2_qcd N=100000: (1 - macro_auc_qcd) of L162 over L188, run by run :: sha256 c6c4138ee5c86423
\newcommand{\MlstFtCmpJcTwoQcdResOverFullEFivePr}{6.0\% more errors}  % experiments/FIGS/data/v2/index.json :: fine-tuning jc2_qcd N=100000: (1 - macro_auc_qcd) of L162 over L188, run by run :: sha256 c6c4138ee5c86423 :: prose twin: two significant figures, direction in words
\newcommand{\MlstFtCmpJcTwoQcdMidOverFullEFive}{1.062\,\ensuremath{\pm}\,0.010}  % experiments/FIGS/data/v2/index.json :: fine-tuning jc2_qcd N=100000: (1 - macro_auc_qcd) of R42_Q1 over L188, run by run :: sha256 c6c4138ee5c86423
\newcommand{\MlstFtCmpJcTwoQcdMidOverFullEFivePr}{6.2\% more errors}  % experiments/FIGS/data/v2/index.json :: fine-tuning jc2_qcd N=100000: (1 - macro_auc_qcd) of R42_Q1 over L188, run by run :: sha256 c6c4138ee5c86423 :: prose twin: two significant figures, direction in words
\newcommand{\MlstFtCmpJcTwoQcdCoarseOverFullEFive}{1.137\,\ensuremath{\pm}\,0.001}  % experiments/FIGS/data/v2/index.json :: fine-tuning jc2_qcd N=100000: (1 - macro_auc_qcd) of R16_Q1 over L188, run by run :: sha256 c6c4138ee5c86423
\newcommand{\MlstFtCmpJcTwoQcdCoarseOverFullEFivePr}{14\% more errors}  % experiments/FIGS/data/v2/index.json :: fine-tuning jc2_qcd N=100000: (1 - macro_auc_qcd) of R16_Q1 over L188, run by run :: sha256 c6c4138ee5c86423 :: prose twin: two significant figures, direction in words
\newcommand{\MlstFtCmpJcTwoQcdResMassOverResEFive}{1.022 and 1.009}  % experiments/FIGS/data/v2/index.json :: fine-tuning jc2_qcd N=100000: (1 - macro_auc_qcd) of L162_MASS over L162, run by run :: sha256 c6c4138ee5c86423
\newcommand{\MlstFtCmpJcTwoQcdResMassOverResEFivePr}{equal within 3\%}  % experiments/FIGS/data/v2/index.json :: fine-tuning jc2_qcd N=100000: (1 - macro_auc_qcd) of L162_MASS over L162, run by run :: sha256 c6c4138ee5c86423 :: prose twin: two significant figures, direction in words
\newcommand{\MlstFtCmpJcTwoQcdResMassOverResEFivePrMax}{equal within 3\%}  % experiments/FIGS/data/v2/index.json :: fine-tuning jc2_qcd N=100000: (1 - macro_auc_qcd) of L162_MASS over L162, run by run :: sha256 c6c4138ee5c86423 :: prose twin of the run farther from 1
\newcommand{\MlstFtCmpJcTwoQcdCoarseMassOverCoarseEFive}{1.028 and 1.027}  % experiments/FIGS/data/v2/index.json :: fine-tuning jc2_qcd N=100000: (1 - macro_auc_qcd) of R16_Q1_MASS over R16_Q1, run by run :: sha256 c6c4138ee5c86423
\newcommand{\MlstFtCmpJcTwoQcdCoarseMassOverCoarseEFivePr}{at least 2.7\% more errors}  % experiments/FIGS/data/v2/index.json :: fine-tuning jc2_qcd N=100000: (1 - macro_auc_qcd) of R16_Q1_MASS over R16_Q1, run by run :: sha256 c6c4138ee5c86423 :: prose twin: two significant figures, direction in words
\newcommand{\MlstFtCmpJcTwoQcdCoarseMassOverCoarseEFivePrMax}{up to 2.8\% more errors}  % experiments/FIGS/data/v2/index.json :: fine-tuning jc2_qcd N=100000: (1 - macro_auc_qcd) of R16_Q1_MASS over R16_Q1, run by run :: sha256 c6c4138ee5c86423 :: prose twin of the run farther from 1
\newcommand{\MlstFtCmpJcTwoQcdResOverFullESix}{1.035\,\ensuremath{\pm}\,0.010}  % experiments/FIGS/data/v2/index.json :: fine-tuning jc2_qcd N=1000000: (1 - macro_auc_qcd) of L162 over L188, run by run :: sha256 c6c4138ee5c86423
\newcommand{\MlstFtCmpJcTwoQcdResOverFullESixPr}{3.5\% more errors}  % experiments/FIGS/data/v2/index.json :: fine-tuning jc2_qcd N=1000000: (1 - macro_auc_qcd) of L162 over L188, run by run :: sha256 c6c4138ee5c86423 :: prose twin: two significant figures, direction in words
\newcommand{\MlstFtCmpJcTwoQcdMidOverFullESix}{1.050\,\ensuremath{\pm}\,0.004}  % experiments/FIGS/data/v2/index.json :: fine-tuning jc2_qcd N=1000000: (1 - macro_auc_qcd) of R42_Q1 over L188, run by run :: sha256 c6c4138ee5c86423
\newcommand{\MlstFtCmpJcTwoQcdMidOverFullESixPr}{5.0\% more errors}  % experiments/FIGS/data/v2/index.json :: fine-tuning jc2_qcd N=1000000: (1 - macro_auc_qcd) of R42_Q1 over L188, run by run :: sha256 c6c4138ee5c86423 :: prose twin: two significant figures, direction in words
\newcommand{\MlstFtCmpJcTwoQcdCoarseOverFullESix}{1.122\,\ensuremath{\pm}\,0.002}  % experiments/FIGS/data/v2/index.json :: fine-tuning jc2_qcd N=1000000: (1 - macro_auc_qcd) of R16_Q1 over L188, run by run :: sha256 c6c4138ee5c86423
\newcommand{\MlstFtCmpJcTwoQcdCoarseOverFullESixPr}{12\% more errors}  % experiments/FIGS/data/v2/index.json :: fine-tuning jc2_qcd N=1000000: (1 - macro_auc_qcd) of R16_Q1 over L188, run by run :: sha256 c6c4138ee5c86423 :: prose twin: two significant figures, direction in words
\newcommand{\MlstFtCmpJcTwoQcdResMassOverResESix}{1.011 and 1.025}  % experiments/FIGS/data/v2/index.json :: fine-tuning jc2_qcd N=1000000: (1 - macro_auc_qcd) of L162_MASS over L162, run by run :: sha256 c6c4138ee5c86423
\newcommand{\MlstFtCmpJcTwoQcdResMassOverResESixPr}{equal within 3\%}  % experiments/FIGS/data/v2/index.json :: fine-tuning jc2_qcd N=1000000: (1 - macro_auc_qcd) of L162_MASS over L162, run by run :: sha256 c6c4138ee5c86423 :: prose twin: two significant figures, direction in words
\newcommand{\MlstFtCmpJcTwoQcdResMassOverResESixPrMax}{equal within 3\%}  % experiments/FIGS/data/v2/index.json :: fine-tuning jc2_qcd N=1000000: (1 - macro_auc_qcd) of L162_MASS over L162, run by run :: sha256 c6c4138ee5c86423 :: prose twin of the run farther from 1
\newcommand{\MlstFtCmpJcTwoQcdCoarseMassOverCoarseESix}{1.015 and 1.015}  % experiments/FIGS/data/v2/index.json :: fine-tuning jc2_qcd N=1000000: (1 - macro_auc_qcd) of R16_Q1_MASS over R16_Q1, run by run :: sha256 c6c4138ee5c86423
\newcommand{\MlstFtCmpJcTwoQcdCoarseMassOverCoarseESixPr}{at least 1.5\% more errors}  % experiments/FIGS/data/v2/index.json :: fine-tuning jc2_qcd N=1000000: (1 - macro_auc_qcd) of R16_Q1_MASS over R16_Q1, run by run :: sha256 c6c4138ee5c86423 :: prose twin: two significant figures, direction in words
\newcommand{\MlstFtCmpJcTwoQcdCoarseMassOverCoarseESixPrMax}{up to 1.5\% more errors}  % experiments/FIGS/data/v2/index.json :: fine-tuning jc2_qcd N=1000000: (1 - macro_auc_qcd) of R16_Q1_MASS over R16_Q1, run by run :: sha256 c6c4138ee5c86423 :: prose twin of the run farther from 1
\newcommand{\MlstFtJcTwoAccFullEFour}{0.3914\,\ensuremath{\pm}\,0.0011}  % experiments/FIGS/data/v2/index.json :: fine-tuning jc2 N=10000: accuracy, L188 :: sha256 c6c4138ee5c86423
\newcommand{\MlstFtJcTwoAccResEFour}{0.3934\,\ensuremath{\pm}\,0.0018}  % experiments/FIGS/data/v2/index.json :: fine-tuning jc2 N=10000: accuracy, L162 :: sha256 c6c4138ee5c86423
\newcommand{\MlstFtJcTwoAccMidEFour}{0.3563\,\ensuremath{\pm}\,0.0025}  % experiments/FIGS/data/v2/index.json :: fine-tuning jc2 N=10000: accuracy, R42_Q1 :: sha256 c6c4138ee5c86423
\newcommand{\MlstFtJcTwoAccCoarseEFour}{0.3219\,\ensuremath{\pm}\,0.0011}  % experiments/FIGS/data/v2/index.json :: fine-tuning jc2 N=10000: accuracy, R16_Q1 :: sha256 c6c4138ee5c86423
\newcommand{\MlstFtJcTwoAccResMassEFour}{0.381 and 0.392}  % experiments/FIGS/data/v2/index.json :: fine-tuning jc2 N=10000: accuracy, L162_MASS :: sha256 c6c4138ee5c86423
\newcommand{\MlstFtJcTwoAccCoarseMassEFour}{0.3156 and 0.3187}  % experiments/FIGS/data/v2/index.json :: fine-tuning jc2 N=10000: accuracy, R16_Q1_MASS :: sha256 c6c4138ee5c86423
\newcommand{\MlstFtJcTwoAccScratchEFour}{0.1492\,\ensuremath{\pm}\,0.0011}  % experiments/FIGS/data/v2/index.json :: fine-tuning jc2 N=10000: accuracy from scratch :: sha256 c6c4138ee5c86423
\newcommand{\MlstFtJcTwoAccFullMinusScratchEFour}{0.2423\,\ensuremath{\pm}\,0.0011}  % experiments/FIGS/data/v2/index.json :: fine-tuning jc2 N=10000: accuracy of each 188-class run minus the from-scratch mean (absolute) :: sha256 c6c4138ee5c86423
\newcommand{\MlstFtJcTwoAccFullEFive}{0.4164\,\ensuremath{\pm}\,0.0015}  % experiments/FIGS/data/v2/index.json :: fine-tuning jc2 N=100000: accuracy, L188 :: sha256 c6c4138ee5c86423
\newcommand{\MlstFtJcTwoAccResEFive}{0.4161\,\ensuremath{\pm}\,0.0015}  % experiments/FIGS/data/v2/index.json :: fine-tuning jc2 N=100000: accuracy, L162 :: sha256 c6c4138ee5c86423
\newcommand{\MlstFtJcTwoAccMidEFive}{0.4007\,\ensuremath{\pm}\,0.0013}  % experiments/FIGS/data/v2/index.json :: fine-tuning jc2 N=100000: accuracy, R42_Q1 :: sha256 c6c4138ee5c86423
\newcommand{\MlstFtJcTwoAccCoarseEFive}{0.3833\,\ensuremath{\pm}\,0.0020}  % experiments/FIGS/data/v2/index.json :: fine-tuning jc2 N=100000: accuracy, R16_Q1 :: sha256 c6c4138ee5c86423
\newcommand{\MlstFtJcTwoAccResMassEFive}{0.4093 and 0.4115}  % experiments/FIGS/data/v2/index.json :: fine-tuning jc2 N=100000: accuracy, L162_MASS :: sha256 c6c4138ee5c86423
\newcommand{\MlstFtJcTwoAccCoarseMassEFive}{0.3758 and 0.3781}  % experiments/FIGS/data/v2/index.json :: fine-tuning jc2 N=100000: accuracy, R16_Q1_MASS :: sha256 c6c4138ee5c86423
\newcommand{\MlstFtJcTwoAccScratchEFive}{0.292\,\ensuremath{\pm}\,0.004}  % experiments/FIGS/data/v2/index.json :: fine-tuning jc2 N=100000: accuracy from scratch :: sha256 c6c4138ee5c86423
\newcommand{\MlstFtJcTwoAccFullMinusScratchEFive}{0.1243\,\ensuremath{\pm}\,0.0015}  % experiments/FIGS/data/v2/index.json :: fine-tuning jc2 N=100000: accuracy of each 188-class run minus the from-scratch mean (absolute) :: sha256 c6c4138ee5c86423
\newcommand{\MlstFtJcTwoAccFullESix}{0.4263\,\ensuremath{\pm}\,0.0005}  % experiments/FIGS/data/v2/index.json :: fine-tuning jc2 N=1000000: accuracy, L188 :: sha256 c6c4138ee5c86423
\newcommand{\MlstFtJcTwoAccResESix}{0.4249\,\ensuremath{\pm}\,0.0021}  % experiments/FIGS/data/v2/index.json :: fine-tuning jc2 N=1000000: accuracy, L162 :: sha256 c6c4138ee5c86423
\newcommand{\MlstFtJcTwoAccMidESix}{0.4166\,\ensuremath{\pm}\,0.0011}  % experiments/FIGS/data/v2/index.json :: fine-tuning jc2 N=1000000: accuracy, R42_Q1 :: sha256 c6c4138ee5c86423
\newcommand{\MlstFtJcTwoAccCoarseESix}{0.4054\,\ensuremath{\pm}\,0.0020}  % experiments/FIGS/data/v2/index.json :: fine-tuning jc2 N=1000000: accuracy, R16_Q1 :: sha256 c6c4138ee5c86423
\newcommand{\MlstFtJcTwoAccResMassESix}{0.4198 and 0.4213}  % experiments/FIGS/data/v2/index.json :: fine-tuning jc2 N=1000000: accuracy, L162_MASS :: sha256 c6c4138ee5c86423
\newcommand{\MlstFtJcTwoAccCoarseMassESix}{0.4012 and 0.4000}  % experiments/FIGS/data/v2/index.json :: fine-tuning jc2 N=1000000: accuracy, R16_Q1_MASS :: sha256 c6c4138ee5c86423
\newcommand{\MlstFtJcTwoAccScratchESix}{0.3677\,\ensuremath{\pm}\,0.0012}  % experiments/FIGS/data/v2/index.json :: fine-tuning jc2 N=1000000: accuracy from scratch :: sha256 c6c4138ee5c86423
\newcommand{\MlstFtJcTwoAccFullMinusScratchESix}{0.0585\,\ensuremath{\pm}\,0.0005}  % experiments/FIGS/data/v2/index.json :: fine-tuning jc2 N=1000000: accuracy of each 188-class run minus the from-scratch mean (absolute) :: sha256 c6c4138ee5c86423
\newcommand{\MlstFtJcTwoAucFull}{0.9797\,\ensuremath{\pm}\,0.0001}  % experiments/FIGS/data/v2/index.json :: fine-tuning class-balanced JetClass-II N=1000000: AUC averaged over the classes, Full :: sha256 c6c4138ee5c86423
\newcommand{\MlstFtJcTwoAucScratch}{0.9734\,\ensuremath{\pm}\,0.0003}  % experiments/FIGS/data/v2/index.json :: fine-tuning class-balanced JetClass-II N=1000000: AUC averaged over the classes, Scratch :: sha256 c6c4138ee5c86423
\newcommand{\MlstFtJcOneAccFullEFour}{0.8117\,\ensuremath{\pm}\,0.0031}  % experiments/FIGS/data/v2/index.json :: fine-tuning jc1 N=10000: accuracy, L188 :: sha256 c6c4138ee5c86423
\newcommand{\MlstFtJcOneAccResEFour}{0.8117\,\ensuremath{\pm}\,0.0013}  % experiments/FIGS/data/v2/index.json :: fine-tuning jc1 N=10000: accuracy, L162 :: sha256 c6c4138ee5c86423
\newcommand{\MlstFtJcOneAccMidEFour}{0.8086\,\ensuremath{\pm}\,0.0031}  % experiments/FIGS/data/v2/index.json :: fine-tuning jc1 N=10000: accuracy, R42_Q1 :: sha256 c6c4138ee5c86423
\newcommand{\MlstFtJcOneAccCoarseEFour}{0.7982\,\ensuremath{\pm}\,0.0023}  % experiments/FIGS/data/v2/index.json :: fine-tuning jc1 N=10000: accuracy, R16_Q1 :: sha256 c6c4138ee5c86423
\newcommand{\MlstFtJcOneAccResMassEFour}{0.8107 and 0.8098}  % experiments/FIGS/data/v2/index.json :: fine-tuning jc1 N=10000: accuracy, L162_MASS :: sha256 c6c4138ee5c86423
\newcommand{\MlstFtJcOneAccCoarseMassEFour}{0.7944 and 0.7990}  % experiments/FIGS/data/v2/index.json :: fine-tuning jc1 N=10000: accuracy, R16_Q1_MASS :: sha256 c6c4138ee5c86423
\newcommand{\MlstFtJcOneAccScratchEFour}{0.677\,\ensuremath{\pm}\,0.004}  % experiments/FIGS/data/v2/index.json :: fine-tuning jc1 N=10000: accuracy from scratch :: sha256 c6c4138ee5c86423
\newcommand{\MlstFtJcOneAccFullMinusScratchEFour}{0.1343\,\ensuremath{\pm}\,0.0031}  % experiments/FIGS/data/v2/index.json :: fine-tuning jc1 N=10000: accuracy of each 188-class run minus the from-scratch mean (absolute) :: sha256 c6c4138ee5c86423
\newcommand{\MlstFtJcOneAccFullEFive}{0.8343\,\ensuremath{\pm}\,0.0035}  % experiments/FIGS/data/v2/index.json :: fine-tuning jc1 N=100000: accuracy, L188 :: sha256 c6c4138ee5c86423
\newcommand{\MlstFtJcOneAccResEFive}{0.8328\,\ensuremath{\pm}\,0.0005}  % experiments/FIGS/data/v2/index.json :: fine-tuning jc1 N=100000: accuracy, L162 :: sha256 c6c4138ee5c86423
\newcommand{\MlstFtJcOneAccMidEFive}{0.8324\,\ensuremath{\pm}\,0.0004}  % experiments/FIGS/data/v2/index.json :: fine-tuning jc1 N=100000: accuracy, R42_Q1 :: sha256 c6c4138ee5c86423
\newcommand{\MlstFtJcOneAccCoarseEFive}{0.8288\,\ensuremath{\pm}\,0.0005}  % experiments/FIGS/data/v2/index.json :: fine-tuning jc1 N=100000: accuracy, R16_Q1 :: sha256 c6c4138ee5c86423
\newcommand{\MlstFtJcOneAccResMassEFive}{0.8326 and 0.8348}  % experiments/FIGS/data/v2/index.json :: fine-tuning jc1 N=100000: accuracy, L162_MASS :: sha256 c6c4138ee5c86423
\newcommand{\MlstFtJcOneAccCoarseMassEFive}{0.8253 and 0.8262}  % experiments/FIGS/data/v2/index.json :: fine-tuning jc1 N=100000: accuracy, R16_Q1_MASS :: sha256 c6c4138ee5c86423
\newcommand{\MlstFtJcOneAccScratchEFive}{0.777\,\ensuremath{\pm}\,0.005}  % experiments/FIGS/data/v2/index.json :: fine-tuning jc1 N=100000: accuracy from scratch :: sha256 c6c4138ee5c86423
\newcommand{\MlstFtJcOneAccFullMinusScratchEFive}{0.0571\,\ensuremath{\pm}\,0.0035}  % experiments/FIGS/data/v2/index.json :: fine-tuning jc1 N=100000: accuracy of each 188-class run minus the from-scratch mean (absolute) :: sha256 c6c4138ee5c86423
\newcommand{\MlstFtJcOneAccFullESix}{0.8468\,\ensuremath{\pm}\,0.0006}  % experiments/FIGS/data/v2/index.json :: fine-tuning jc1 N=1000000: accuracy, L188 :: sha256 c6c4138ee5c86423
\newcommand{\MlstFtJcOneAccResESix}{0.8475\,\ensuremath{\pm}\,0.0009}  % experiments/FIGS/data/v2/index.json :: fine-tuning jc1 N=1000000: accuracy, L162 :: sha256 c6c4138ee5c86423
\newcommand{\MlstFtJcOneAccMidESix}{0.8459\,\ensuremath{\pm}\,0.0002}  % experiments/FIGS/data/v2/index.json :: fine-tuning jc1 N=1000000: accuracy, R42_Q1 :: sha256 c6c4138ee5c86423
\newcommand{\MlstFtJcOneAccCoarseESix}{0.8422\,\ensuremath{\pm}\,0.0005}  % experiments/FIGS/data/v2/index.json :: fine-tuning jc1 N=1000000: accuracy, R16_Q1 :: sha256 c6c4138ee5c86423
\newcommand{\MlstFtJcOneAccResMassESix}{0.8459 and 0.8471}  % experiments/FIGS/data/v2/index.json :: fine-tuning jc1 N=1000000: accuracy, L162_MASS :: sha256 c6c4138ee5c86423
\newcommand{\MlstFtJcOneAccCoarseMassESix}{0.8400 and 0.8398}  % experiments/FIGS/data/v2/index.json :: fine-tuning jc1 N=1000000: accuracy, R16_Q1_MASS :: sha256 c6c4138ee5c86423
\newcommand{\MlstFtJcOneAccScratchESix}{0.8267\,\ensuremath{\pm}\,0.0011}  % experiments/FIGS/data/v2/index.json :: fine-tuning jc1 N=1000000: accuracy from scratch :: sha256 c6c4138ee5c86423
\newcommand{\MlstFtJcOneAccFullMinusScratchESix}{0.0201\,\ensuremath{\pm}\,0.0006}  % experiments/FIGS/data/v2/index.json :: fine-tuning jc1 N=1000000: accuracy of each 188-class run minus the from-scratch mean (absolute) :: sha256 c6c4138ee5c86423
\newcommand{\MlstFtTopRthirtyScratchEThree}{49\,\ensuremath{\pm}\,12}  % experiments/FIGS/data/v2/index.json :: fine-tuning top N=1000: R30 from scratch, every seed :: sha256 c6c4138ee5c86423
\newcommand{\MlstFtTopRthirtyFullEThree}{72\,\ensuremath{\pm}\,19}  % experiments/FIGS/data/v2/index.json :: fine-tuning top N=1000: R30, L188 :: sha256 c6c4138ee5c86423
\newcommand{\MlstFtTopRthirtyFullOverScratchEThree}{1.5\,\ensuremath{\pm}\,0.4}  % experiments/FIGS/data/v2/index.json :: fine-tuning top N=1000: R30 of each L188 run over the from-scratch mean, SD over runs and standard error of that mean in quadrature :: sha256 c6c4138ee5c86423
\newcommand{\MlstFtTopRthirtyFullOverScratchEThreePr}{equal within 50\%}  % experiments/FIGS/data/v2/index.json :: fine-tuning top N=1000: R30 of each L188 run over the from-scratch mean, SD over runs and standard error of that mean in quadrature :: sha256 c6c4138ee5c86423 :: prose twin: two significant figures, direction in words
\newcommand{\MlstFtTopRthirtyResEThree}{135\,\ensuremath{\pm}\,26}  % experiments/FIGS/data/v2/index.json :: fine-tuning top N=1000: R30, L162 :: sha256 c6c4138ee5c86423
\newcommand{\MlstFtTopRthirtyResOverScratchEThree}{2.8\,\ensuremath{\pm}\,0.7}  % experiments/FIGS/data/v2/index.json :: fine-tuning top N=1000: R30 of each L162 run over the from-scratch mean, SD over runs and standard error of that mean in quadrature :: sha256 c6c4138ee5c86423
\newcommand{\MlstFtTopRthirtyResOverScratchEThreePr}{2.8 times the rejection}  % experiments/FIGS/data/v2/index.json :: fine-tuning top N=1000: R30 of each L162 run over the from-scratch mean, SD over runs and standard error of that mean in quadrature :: sha256 c6c4138ee5c86423 :: prose twin: two significant figures, direction in words
\newcommand{\MlstFtTopRthirtyMidEThree}{104\,\ensuremath{\pm}\,20}  % experiments/FIGS/data/v2/index.json :: fine-tuning top N=1000: R30, R42_Q1 :: sha256 c6c4138ee5c86423
\newcommand{\MlstFtTopRthirtyMidOverScratchEThree}{2.1\,\ensuremath{\pm}\,0.5}  % experiments/FIGS/data/v2/index.json :: fine-tuning top N=1000: R30 of each R42_Q1 run over the from-scratch mean, SD over runs and standard error of that mean in quadrature :: sha256 c6c4138ee5c86423
\newcommand{\MlstFtTopRthirtyMidOverScratchEThreePr}{2.1 times the rejection}  % experiments/FIGS/data/v2/index.json :: fine-tuning top N=1000: R30 of each R42_Q1 run over the from-scratch mean, SD over runs and standard error of that mean in quadrature :: sha256 c6c4138ee5c86423 :: prose twin: two significant figures, direction in words
\newcommand{\MlstFtTopRthirtyCoarseEThree}{83\,\ensuremath{\pm}\,25}  % experiments/FIGS/data/v2/index.json :: fine-tuning top N=1000: R30, R16_Q1 :: sha256 c6c4138ee5c86423
\newcommand{\MlstFtTopRthirtyCoarseOverScratchEThree}{1.7\,\ensuremath{\pm}\,0.6}  % experiments/FIGS/data/v2/index.json :: fine-tuning top N=1000: R30 of each R16_Q1 run over the from-scratch mean, SD over runs and standard error of that mean in quadrature :: sha256 c6c4138ee5c86423
\newcommand{\MlstFtTopRthirtyCoarseOverScratchEThreePr}{equal within 80\%}  % experiments/FIGS/data/v2/index.json :: fine-tuning top N=1000: R30 of each R16_Q1 run over the from-scratch mean, SD over runs and standard error of that mean in quadrature :: sha256 c6c4138ee5c86423 :: prose twin: two significant figures, direction in words
\newcommand{\MlstFtTopRthirtyResMassEThree}{105.9 and 112.7}  % experiments/FIGS/data/v2/index.json :: fine-tuning top N=1000: R30, L162_MASS :: sha256 c6c4138ee5c86423
\newcommand{\MlstFtTopRthirtyResMassOverScratchEThree}{2.18 and 2.32}  % experiments/FIGS/data/v2/index.json :: fine-tuning top N=1000: R30 of each L162_MASS run over the from-scratch mean, SD over runs and standard error of that mean in quadrature :: sha256 c6c4138ee5c86423
\newcommand{\MlstFtTopRthirtyResMassOverScratchEThreePr}{at least 2.2 times the rejection}  % experiments/FIGS/data/v2/index.json :: fine-tuning top N=1000: R30 of each L162_MASS run over the from-scratch mean, SD over runs and standard error of that mean in quadrature :: sha256 c6c4138ee5c86423 :: prose twin: two significant figures, direction in words
\newcommand{\MlstFtTopRthirtyResMassOverScratchEThreePrMax}{up to 2.3 times the rejection}  % experiments/FIGS/data/v2/index.json :: fine-tuning top N=1000: R30 of each L162_MASS run over the from-scratch mean, SD over runs and standard error of that mean in quadrature :: sha256 c6c4138ee5c86423 :: prose twin of the run farther from 1
\newcommand{\MlstFtTopRthirtyCoarseMassEThree}{59.1 and 63.9}  % experiments/FIGS/data/v2/index.json :: fine-tuning top N=1000: R30, R16_Q1_MASS :: sha256 c6c4138ee5c86423
\newcommand{\MlstFtTopRthirtyCoarseMassOverScratchEThree}{1.216 and 1.314}  % experiments/FIGS/data/v2/index.json :: fine-tuning top N=1000: R30 of each R16_Q1_MASS run over the from-scratch mean, SD over runs and standard error of that mean in quadrature :: sha256 c6c4138ee5c86423
\newcommand{\MlstFtTopRthirtyCoarseMassOverScratchEThreePr}{at least 22\% more rejection}  % experiments/FIGS/data/v2/index.json :: fine-tuning top N=1000: R30 of each R16_Q1_MASS run over the from-scratch mean, SD over runs and standard error of that mean in quadrature :: sha256 c6c4138ee5c86423 :: prose twin: two significant figures, direction in words
\newcommand{\MlstFtTopRthirtyCoarseMassOverScratchEThreePrMax}{up to 31\% more rejection}  % experiments/FIGS/data/v2/index.json :: fine-tuning top N=1000: R30 of each R16_Q1_MASS run over the from-scratch mean, SD over runs and standard error of that mean in quadrature :: sha256 c6c4138ee5c86423 :: prose twin of the run farther from 1
\newcommand{\MlstFtTopRthirtyResOverFullEThree}{2.0\,\ensuremath{\pm}\,0.8}  % experiments/FIGS/data/v2/index.json :: fine-tuning top N=1000: R30 of L162 over L188, run by run :: sha256 c6c4138ee5c86423
\newcommand{\MlstFtTopRthirtyResOverFullEThreePr}{equal within 100\%}  % experiments/FIGS/data/v2/index.json :: fine-tuning top N=1000: R30 of L162 over L188, run by run :: sha256 c6c4138ee5c86423 :: prose twin: two significant figures, direction in words
\newcommand{\MlstFtTopRthirtyMidOverFullEThree}{1.46\,\ensuremath{\pm}\,0.12}  % experiments/FIGS/data/v2/index.json :: fine-tuning top N=1000: R30 of R42_Q1 over L188, run by run :: sha256 c6c4138ee5c86423
\newcommand{\MlstFtTopRthirtyMidOverFullEThreePr}{46\% more rejection}  % experiments/FIGS/data/v2/index.json :: fine-tuning top N=1000: R30 of R42_Q1 over L188, run by run :: sha256 c6c4138ee5c86423 :: prose twin: two significant figures, direction in words
\newcommand{\MlstFtTopRthirtyCoarseOverFullEThree}{1.3\,\ensuremath{\pm}\,0.7}  % experiments/FIGS/data/v2/index.json :: fine-tuning top N=1000: R30 of R16_Q1 over L188, run by run :: sha256 c6c4138ee5c86423
\newcommand{\MlstFtTopRthirtyCoarseOverFullEThreePr}{equal within 80\%}  % experiments/FIGS/data/v2/index.json :: fine-tuning top N=1000: R30 of R16_Q1 over L188, run by run :: sha256 c6c4138ee5c86423 :: prose twin: two significant figures, direction in words
\newcommand{\MlstFtTopRthirtyResMassOverResEThree}{0.73 and 1.07}  % experiments/FIGS/data/v2/index.json :: fine-tuning top N=1000: R30 of L162_MASS over L162, run by run :: sha256 c6c4138ee5c86423
\newcommand{\MlstFtTopRthirtyResMassOverResEThreePr}{equal within 30\%}  % experiments/FIGS/data/v2/index.json :: fine-tuning top N=1000: R30 of L162_MASS over L162, run by run :: sha256 c6c4138ee5c86423 :: prose twin: two significant figures, direction in words
\newcommand{\MlstFtTopRthirtyResMassOverResEThreePrMax}{equal within 30\%}  % experiments/FIGS/data/v2/index.json :: fine-tuning top N=1000: R30 of L162_MASS over L162, run by run :: sha256 c6c4138ee5c86423 :: prose twin of the run farther from 1
\newcommand{\MlstFtTopRthirtyCoarseMassOverCoarseEThree}{0.56 and 0.75}  % experiments/FIGS/data/v2/index.json :: fine-tuning top N=1000: R30 of R16_Q1_MASS over R16_Q1, run by run :: sha256 c6c4138ee5c86423
\newcommand{\MlstFtTopRthirtyCoarseMassOverCoarseEThreePr}{equal within 50\%}  % experiments/FIGS/data/v2/index.json :: fine-tuning top N=1000: R30 of R16_Q1_MASS over R16_Q1, run by run :: sha256 c6c4138ee5c86423 :: prose twin: two significant figures, direction in words
\newcommand{\MlstFtTopRthirtyCoarseMassOverCoarseEThreePrMax}{equal within 50\%}  % experiments/FIGS/data/v2/index.json :: fine-tuning top N=1000: R30 of R16_Q1_MASS over R16_Q1, run by run :: sha256 c6c4138ee5c86423 :: prose twin of the run farther from 1
\newcommand{\MlstFtTopRthirtyRandOverFullEThree}{1.3\,\ensuremath{\pm}\,0.5}  % experiments/FIGS/data/v2/index.json :: fine-tuning top N=1000: R30 of each random partition over the L188 run of its seed :: sha256 c6c4138ee5c86423
\newcommand{\MlstFtTopRthirtyRandOverFullEThreePr}{equal within 50\%}  % experiments/FIGS/data/v2/index.json :: fine-tuning top N=1000: R30 of each random partition over the L188 run of its seed :: sha256 c6c4138ee5c86423 :: prose twin: two significant figures, direction in words
\newcommand{\MlstFtTopRthirtyRandOverCoarseEThree}{0.87\,\ensuremath{\pm}\,0.12}  % experiments/FIGS/data/v2/index.json :: fine-tuning top N=1000: R30 of each random partition over the R16_Q1 run of its seed :: sha256 c6c4138ee5c86423
\newcommand{\MlstFtTopRthirtyRandOverCoarseEThreePr}{equal within 20\%}  % experiments/FIGS/data/v2/index.json :: fine-tuning top N=1000: R30 of each random partition over the R16_Q1 run of its seed :: sha256 c6c4138ee5c86423 :: prose twin: two significant figures, direction in words
\newcommand{\MlstFtTopRthirtyScratchEFour}{340\,\ensuremath{\pm}\,50}  % experiments/FIGS/data/v2/index.json :: fine-tuning top N=10000: R30 from scratch, every seed :: sha256 c6c4138ee5c86423
\newcommand{\MlstFtTopRthirtyFullEFour}{420\,\ensuremath{\pm}\,9}  % experiments/FIGS/data/v2/index.json :: fine-tuning top N=10000: R30, L188 :: sha256 c6c4138ee5c86423
\newcommand{\MlstFtTopRthirtyFullOverScratchEFour}{1.24\,\ensuremath{\pm}\,0.11}  % experiments/FIGS/data/v2/index.json :: fine-tuning top N=10000: R30 of each L188 run over the from-scratch mean, SD over runs and standard error of that mean in quadrature :: sha256 c6c4138ee5c86423
\newcommand{\MlstFtTopRthirtyFullOverScratchEFourPr}{24\% more rejection}  % experiments/FIGS/data/v2/index.json :: fine-tuning top N=10000: R30 of each L188 run over the from-scratch mean, SD over runs and standard error of that mean in quadrature :: sha256 c6c4138ee5c86423 :: prose twin: two significant figures, direction in words
\newcommand{\MlstFtTopRthirtyResEFour}{430\,\ensuremath{\pm}\,60}  % experiments/FIGS/data/v2/index.json :: fine-tuning top N=10000: R30, L162 :: sha256 c6c4138ee5c86423
\newcommand{\MlstFtTopRthirtyResOverScratchEFour}{1.27\,\ensuremath{\pm}\,0.21}  % experiments/FIGS/data/v2/index.json :: fine-tuning top N=10000: R30 of each L162 run over the from-scratch mean, SD over runs and standard error of that mean in quadrature :: sha256 c6c4138ee5c86423
\newcommand{\MlstFtTopRthirtyResOverScratchEFourPr}{equal within 30\%}  % experiments/FIGS/data/v2/index.json :: fine-tuning top N=10000: R30 of each L162 run over the from-scratch mean, SD over runs and standard error of that mean in quadrature :: sha256 c6c4138ee5c86423 :: prose twin: two significant figures, direction in words
\newcommand{\MlstFtTopRthirtyMidEFour}{455\,\ensuremath{\pm}\,30}  % experiments/FIGS/data/v2/index.json :: fine-tuning top N=10000: R30, R42_Q1 :: sha256 c6c4138ee5c86423
\newcommand{\MlstFtTopRthirtyMidOverScratchEFour}{1.34\,\ensuremath{\pm}\,0.14}  % experiments/FIGS/data/v2/index.json :: fine-tuning top N=10000: R30 of each R42_Q1 run over the from-scratch mean, SD over runs and standard error of that mean in quadrature :: sha256 c6c4138ee5c86423
\newcommand{\MlstFtTopRthirtyMidOverScratchEFourPr}{34\% more rejection}  % experiments/FIGS/data/v2/index.json :: fine-tuning top N=10000: R30 of each R42_Q1 run over the from-scratch mean, SD over runs and standard error of that mean in quadrature :: sha256 c6c4138ee5c86423 :: prose twin: two significant figures, direction in words
\newcommand{\MlstFtTopRthirtyCoarseEFour}{350\,\ensuremath{\pm}\,40}  % experiments/FIGS/data/v2/index.json :: fine-tuning top N=10000: R30, R16_Q1 :: sha256 c6c4138ee5c86423
\newcommand{\MlstFtTopRthirtyCoarseOverScratchEFour}{1.04\,\ensuremath{\pm}\,0.14}  % experiments/FIGS/data/v2/index.json :: fine-tuning top N=10000: R30 of each R16_Q1 run over the from-scratch mean, SD over runs and standard error of that mean in quadrature :: sha256 c6c4138ee5c86423
\newcommand{\MlstFtTopRthirtyCoarseOverScratchEFourPr}{equal within 20\%}  % experiments/FIGS/data/v2/index.json :: fine-tuning top N=10000: R30 of each R16_Q1 run over the from-scratch mean, SD over runs and standard error of that mean in quadrature :: sha256 c6c4138ee5c86423 :: prose twin: two significant figures, direction in words
\newcommand{\MlstFtTopRthirtyResMassEFour}{373.2 and 376.7}  % experiments/FIGS/data/v2/index.json :: fine-tuning top N=10000: R30, L162_MASS :: sha256 c6c4138ee5c86423
\newcommand{\MlstFtTopRthirtyResMassOverScratchEFour}{1.101 and 1.111}  % experiments/FIGS/data/v2/index.json :: fine-tuning top N=10000: R30 of each L162_MASS run over the from-scratch mean, SD over runs and standard error of that mean in quadrature :: sha256 c6c4138ee5c86423
\newcommand{\MlstFtTopRthirtyResMassOverScratchEFourPr}{at least 10\% more rejection}  % experiments/FIGS/data/v2/index.json :: fine-tuning top N=10000: R30 of each L162_MASS run over the from-scratch mean, SD over runs and standard error of that mean in quadrature :: sha256 c6c4138ee5c86423 :: prose twin: two significant figures, direction in words
\newcommand{\MlstFtTopRthirtyResMassOverScratchEFourPrMax}{up to 11\% more rejection}  % experiments/FIGS/data/v2/index.json :: fine-tuning top N=10000: R30 of each L162_MASS run over the from-scratch mean, SD over runs and standard error of that mean in quadrature :: sha256 c6c4138ee5c86423 :: prose twin of the run farther from 1
\newcommand{\MlstFtTopRthirtyCoarseMassEFour}{365.8 and 359.9}  % experiments/FIGS/data/v2/index.json :: fine-tuning top N=10000: R30, R16_Q1_MASS :: sha256 c6c4138ee5c86423
\newcommand{\MlstFtTopRthirtyCoarseMassOverScratchEFour}{1.079 and 1.062}  % experiments/FIGS/data/v2/index.json :: fine-tuning top N=10000: R30 of each R16_Q1_MASS run over the from-scratch mean, SD over runs and standard error of that mean in quadrature :: sha256 c6c4138ee5c86423
\newcommand{\MlstFtTopRthirtyCoarseMassOverScratchEFourPr}{at least 6.2\% more rejection}  % experiments/FIGS/data/v2/index.json :: fine-tuning top N=10000: R30 of each R16_Q1_MASS run over the from-scratch mean, SD over runs and standard error of that mean in quadrature :: sha256 c6c4138ee5c86423 :: prose twin: two significant figures, direction in words
\newcommand{\MlstFtTopRthirtyCoarseMassOverScratchEFourPrMax}{up to 7.9\% more rejection}  % experiments/FIGS/data/v2/index.json :: fine-tuning top N=10000: R30 of each R16_Q1_MASS run over the from-scratch mean, SD over runs and standard error of that mean in quadrature :: sha256 c6c4138ee5c86423 :: prose twin of the run farther from 1
\newcommand{\MlstFtTopRthirtyResOverFullEFour}{1.02\,\ensuremath{\pm}\,0.12}  % experiments/FIGS/data/v2/index.json :: fine-tuning top N=10000: R30 of L162 over L188, run by run :: sha256 c6c4138ee5c86423
\newcommand{\MlstFtTopRthirtyResOverFullEFourPr}{equal within 20\%}  % experiments/FIGS/data/v2/index.json :: fine-tuning top N=10000: R30 of L162 over L188, run by run :: sha256 c6c4138ee5c86423 :: prose twin: two significant figures, direction in words
\newcommand{\MlstFtTopRthirtyMidOverFullEFour}{1.08\,\ensuremath{\pm}\,0.08}  % experiments/FIGS/data/v2/index.json :: fine-tuning top N=10000: R30 of R42_Q1 over L188, run by run :: sha256 c6c4138ee5c86423
\newcommand{\MlstFtTopRthirtyMidOverFullEFourPr}{equal within 9\%}  % experiments/FIGS/data/v2/index.json :: fine-tuning top N=10000: R30 of R42_Q1 over L188, run by run :: sha256 c6c4138ee5c86423 :: prose twin: two significant figures, direction in words
\newcommand{\MlstFtTopRthirtyCoarseOverFullEFour}{0.84\,\ensuremath{\pm}\,0.07}  % experiments/FIGS/data/v2/index.json :: fine-tuning top N=10000: R30 of R16_Q1 over L188, run by run :: sha256 c6c4138ee5c86423
\newcommand{\MlstFtTopRthirtyCoarseOverFullEFourPr}{16\% less rejection}  % experiments/FIGS/data/v2/index.json :: fine-tuning top N=10000: R30 of R16_Q1 over L188, run by run :: sha256 c6c4138ee5c86423 :: prose twin: two significant figures, direction in words
\newcommand{\MlstFtTopRthirtyResMassOverResEFour}{1.00 and 0.76}  % experiments/FIGS/data/v2/index.json :: fine-tuning top N=10000: R30 of L162_MASS over L162, run by run :: sha256 c6c4138ee5c86423
\newcommand{\MlstFtTopRthirtyResMassOverResEFourPr}{equal within 30\%}  % experiments/FIGS/data/v2/index.json :: fine-tuning top N=10000: R30 of L162_MASS over L162, run by run :: sha256 c6c4138ee5c86423 :: prose twin: two significant figures, direction in words
\newcommand{\MlstFtTopRthirtyResMassOverResEFourPrMax}{equal within 30\%}  % experiments/FIGS/data/v2/index.json :: fine-tuning top N=10000: R30 of L162_MASS over L162, run by run :: sha256 c6c4138ee5c86423 :: prose twin of the run farther from 1
\newcommand{\MlstFtTopRthirtyCoarseMassOverCoarseEFour}{1.13 and 0.92}  % experiments/FIGS/data/v2/index.json :: fine-tuning top N=10000: R30 of R16_Q1_MASS over R16_Q1, run by run :: sha256 c6c4138ee5c86423
\newcommand{\MlstFtTopRthirtyCoarseMassOverCoarseEFourPr}{equal within 20\%}  % experiments/FIGS/data/v2/index.json :: fine-tuning top N=10000: R30 of R16_Q1_MASS over R16_Q1, run by run :: sha256 c6c4138ee5c86423 :: prose twin: two significant figures, direction in words
\newcommand{\MlstFtTopRthirtyCoarseMassOverCoarseEFourPrMax}{equal within 20\%}  % experiments/FIGS/data/v2/index.json :: fine-tuning top N=10000: R30 of R16_Q1_MASS over R16_Q1, run by run :: sha256 c6c4138ee5c86423 :: prose twin of the run farther from 1
\newcommand{\MlstFtTopRthirtyRandOverFullEFour}{0.83\,\ensuremath{\pm}\,0.10}  % experiments/FIGS/data/v2/index.json :: fine-tuning top N=10000: R30 of each random partition over the L188 run of its seed :: sha256 c6c4138ee5c86423
\newcommand{\MlstFtTopRthirtyRandOverFullEFourPr}{equal within 20\%}  % experiments/FIGS/data/v2/index.json :: fine-tuning top N=10000: R30 of each random partition over the L188 run of its seed :: sha256 c6c4138ee5c86423 :: prose twin: two significant figures, direction in words
\newcommand{\MlstFtTopRthirtyRandOverCoarseEFour}{0.97\,\ensuremath{\pm}\,0.15}  % experiments/FIGS/data/v2/index.json :: fine-tuning top N=10000: R30 of each random partition over the R16_Q1 run of its seed :: sha256 c6c4138ee5c86423
\newcommand{\MlstFtTopRthirtyRandOverCoarseEFourPr}{equal within 20\%}  % experiments/FIGS/data/v2/index.json :: fine-tuning top N=10000: R30 of each random partition over the R16_Q1 run of its seed :: sha256 c6c4138ee5c86423 :: prose twin: two significant figures, direction in words
\newcommand{\MlstFtTopRthirtyScratchEFive}{560\,\ensuremath{\pm}\,90}  % experiments/FIGS/data/v2/index.json :: fine-tuning top N=100000: R30 from scratch, every seed :: sha256 c6c4138ee5c86423
\newcommand{\MlstFtTopRthirtyFullEFive}{850\,\ensuremath{\pm}\,90}  % experiments/FIGS/data/v2/index.json :: fine-tuning top N=100000: R30, L188 :: sha256 c6c4138ee5c86423
\newcommand{\MlstFtTopRthirtyFullOverScratchEFive}{1.50\,\ensuremath{\pm}\,0.21}  % experiments/FIGS/data/v2/index.json :: fine-tuning top N=100000: R30 of each L188 run over the from-scratch mean, SD over runs and standard error of that mean in quadrature :: sha256 c6c4138ee5c86423
\newcommand{\MlstFtTopRthirtyFullOverScratchEFivePr}{50\% more rejection}  % experiments/FIGS/data/v2/index.json :: fine-tuning top N=100000: R30 of each L188 run over the from-scratch mean, SD over runs and standard error of that mean in quadrature :: sha256 c6c4138ee5c86423 :: prose twin: two significant figures, direction in words
\newcommand{\MlstFtTopRthirtyResEFive}{810\,\ensuremath{\pm}\,60}  % experiments/FIGS/data/v2/index.json :: fine-tuning top N=100000: R30, L162 :: sha256 c6c4138ee5c86423
\newcommand{\MlstFtTopRthirtyResOverScratchEFive}{1.44\,\ensuremath{\pm}\,0.18}  % experiments/FIGS/data/v2/index.json :: fine-tuning top N=100000: R30 of each L162 run over the from-scratch mean, SD over runs and standard error of that mean in quadrature :: sha256 c6c4138ee5c86423
\newcommand{\MlstFtTopRthirtyResOverScratchEFivePr}{44\% more rejection}  % experiments/FIGS/data/v2/index.json :: fine-tuning top N=100000: R30 of each L162 run over the from-scratch mean, SD over runs and standard error of that mean in quadrature :: sha256 c6c4138ee5c86423 :: prose twin: two significant figures, direction in words
\newcommand{\MlstFtTopRthirtyMidEFive}{850\,\ensuremath{\pm}\,80}  % experiments/FIGS/data/v2/index.json :: fine-tuning top N=100000: R30, R42_Q1 :: sha256 c6c4138ee5c86423
\newcommand{\MlstFtTopRthirtyMidOverScratchEFive}{1.51\,\ensuremath{\pm}\,0.20}  % experiments/FIGS/data/v2/index.json :: fine-tuning top N=100000: R30 of each R42_Q1 run over the from-scratch mean, SD over runs and standard error of that mean in quadrature :: sha256 c6c4138ee5c86423
\newcommand{\MlstFtTopRthirtyMidOverScratchEFivePr}{51\% more rejection}  % experiments/FIGS/data/v2/index.json :: fine-tuning top N=100000: R30 of each R42_Q1 run over the from-scratch mean, SD over runs and standard error of that mean in quadrature :: sha256 c6c4138ee5c86423 :: prose twin: two significant figures, direction in words
\newcommand{\MlstFtTopRthirtyCoarseEFive}{830\,\ensuremath{\pm}\,60}  % experiments/FIGS/data/v2/index.json :: fine-tuning top N=100000: R30, R16_Q1 :: sha256 c6c4138ee5c86423
\newcommand{\MlstFtTopRthirtyCoarseOverScratchEFive}{1.47\,\ensuremath{\pm}\,0.17}  % experiments/FIGS/data/v2/index.json :: fine-tuning top N=100000: R30 of each R16_Q1 run over the from-scratch mean, SD over runs and standard error of that mean in quadrature :: sha256 c6c4138ee5c86423
\newcommand{\MlstFtTopRthirtyCoarseOverScratchEFivePr}{47\% more rejection}  % experiments/FIGS/data/v2/index.json :: fine-tuning top N=100000: R30 of each R16_Q1 run over the from-scratch mean, SD over runs and standard error of that mean in quadrature :: sha256 c6c4138ee5c86423 :: prose twin: two significant figures, direction in words
\newcommand{\MlstFtTopRthirtyResMassEFive}{848 and 848}  % experiments/FIGS/data/v2/index.json :: fine-tuning top N=100000: R30, L162_MASS :: sha256 c6c4138ee5c86423
\newcommand{\MlstFtTopRthirtyResMassOverScratchEFive}{1.50 and 1.50}  % experiments/FIGS/data/v2/index.json :: fine-tuning top N=100000: R30 of each L162_MASS run over the from-scratch mean, SD over runs and standard error of that mean in quadrature :: sha256 c6c4138ee5c86423
\newcommand{\MlstFtTopRthirtyResMassOverScratchEFivePr}{at least 50\% more rejection}  % experiments/FIGS/data/v2/index.json :: fine-tuning top N=100000: R30 of each L162_MASS run over the from-scratch mean, SD over runs and standard error of that mean in quadrature :: sha256 c6c4138ee5c86423 :: prose twin: two significant figures, direction in words
\newcommand{\MlstFtTopRthirtyResMassOverScratchEFivePrMax}{up to 50\% more rejection}  % experiments/FIGS/data/v2/index.json :: fine-tuning top N=100000: R30 of each L162_MASS run over the from-scratch mean, SD over runs and standard error of that mean in quadrature :: sha256 c6c4138ee5c86423 :: prose twin of the run farther from 1
\newcommand{\MlstFtTopRthirtyCoarseMassEFive}{856 and 774}  % experiments/FIGS/data/v2/index.json :: fine-tuning top N=100000: R30, R16_Q1_MASS :: sha256 c6c4138ee5c86423
\newcommand{\MlstFtTopRthirtyCoarseMassOverScratchEFive}{1.51 and 1.37}  % experiments/FIGS/data/v2/index.json :: fine-tuning top N=100000: R30 of each R16_Q1_MASS run over the from-scratch mean, SD over runs and standard error of that mean in quadrature :: sha256 c6c4138ee5c86423
\newcommand{\MlstFtTopRthirtyCoarseMassOverScratchEFivePr}{at least 37\% more rejection}  % experiments/FIGS/data/v2/index.json :: fine-tuning top N=100000: R30 of each R16_Q1_MASS run over the from-scratch mean, SD over runs and standard error of that mean in quadrature :: sha256 c6c4138ee5c86423 :: prose twin: two significant figures, direction in words
\newcommand{\MlstFtTopRthirtyCoarseMassOverScratchEFivePrMax}{up to 51\% more rejection}  % experiments/FIGS/data/v2/index.json :: fine-tuning top N=100000: R30 of each R16_Q1_MASS run over the from-scratch mean, SD over runs and standard error of that mean in quadrature :: sha256 c6c4138ee5c86423 :: prose twin of the run farther from 1
\newcommand{\MlstFtTopRthirtyResOverFullEFive}{0.96\,\ensuremath{\pm}\,0.05}  % experiments/FIGS/data/v2/index.json :: fine-tuning top N=100000: R30 of L162 over L188, run by run :: sha256 c6c4138ee5c86423
\newcommand{\MlstFtTopRthirtyResOverFullEFivePr}{equal within 6\%}  % experiments/FIGS/data/v2/index.json :: fine-tuning top N=100000: R30 of L162 over L188, run by run :: sha256 c6c4138ee5c86423 :: prose twin: two significant figures, direction in words
\newcommand{\MlstFtTopRthirtyMidOverFullEFive}{1.02\,\ensuremath{\pm}\,0.18}  % experiments/FIGS/data/v2/index.json :: fine-tuning top N=100000: R30 of R42_Q1 over L188, run by run :: sha256 c6c4138ee5c86423
\newcommand{\MlstFtTopRthirtyMidOverFullEFivePr}{equal within 20\%}  % experiments/FIGS/data/v2/index.json :: fine-tuning top N=100000: R30 of R42_Q1 over L188, run by run :: sha256 c6c4138ee5c86423 :: prose twin: two significant figures, direction in words
\newcommand{\MlstFtTopRthirtyCoarseOverFullEFive}{0.99\,\ensuremath{\pm}\,0.17}  % experiments/FIGS/data/v2/index.json :: fine-tuning top N=100000: R30 of R16_Q1 over L188, run by run :: sha256 c6c4138ee5c86423
\newcommand{\MlstFtTopRthirtyCoarseOverFullEFivePr}{equal within 20\%}  % experiments/FIGS/data/v2/index.json :: fine-tuning top N=100000: R30 of R16_Q1 over L188, run by run :: sha256 c6c4138ee5c86423 :: prose twin: two significant figures, direction in words
\newcommand{\MlstFtTopRthirtyResMassOverResEFive}{1.00 and 1.15}  % experiments/FIGS/data/v2/index.json :: fine-tuning top N=100000: R30 of L162_MASS over L162, run by run :: sha256 c6c4138ee5c86423
\newcommand{\MlstFtTopRthirtyResMassOverResEFivePr}{equal within 20\%}  % experiments/FIGS/data/v2/index.json :: fine-tuning top N=100000: R30 of L162_MASS over L162, run by run :: sha256 c6c4138ee5c86423 :: prose twin: two significant figures, direction in words
\newcommand{\MlstFtTopRthirtyResMassOverResEFivePrMax}{equal within 20\%}  % experiments/FIGS/data/v2/index.json :: fine-tuning top N=100000: R30 of L162_MASS over L162, run by run :: sha256 c6c4138ee5c86423 :: prose twin of the run farther from 1
\newcommand{\MlstFtTopRthirtyCoarseMassOverCoarseEFive}{1.08 and 0.87}  % experiments/FIGS/data/v2/index.json :: fine-tuning top N=100000: R30 of R16_Q1_MASS over R16_Q1, run by run :: sha256 c6c4138ee5c86423
\newcommand{\MlstFtTopRthirtyCoarseMassOverCoarseEFivePr}{equal within 20\%}  % experiments/FIGS/data/v2/index.json :: fine-tuning top N=100000: R30 of R16_Q1_MASS over R16_Q1, run by run :: sha256 c6c4138ee5c86423 :: prose twin: two significant figures, direction in words
\newcommand{\MlstFtTopRthirtyCoarseMassOverCoarseEFivePrMax}{equal within 20\%}  % experiments/FIGS/data/v2/index.json :: fine-tuning top N=100000: R30 of R16_Q1_MASS over R16_Q1, run by run :: sha256 c6c4138ee5c86423 :: prose twin of the run farther from 1
\newcommand{\MlstFtTopRthirtyRandOverFullEFive}{1.01\,\ensuremath{\pm}\,0.22}  % experiments/FIGS/data/v2/index.json :: fine-tuning top N=100000: R30 of each random partition over the L188 run of its seed :: sha256 c6c4138ee5c86423
\newcommand{\MlstFtTopRthirtyRandOverFullEFivePr}{equal within 30\%}  % experiments/FIGS/data/v2/index.json :: fine-tuning top N=100000: R30 of each random partition over the L188 run of its seed :: sha256 c6c4138ee5c86423 :: prose twin: two significant figures, direction in words
\newcommand{\MlstFtTopRthirtyRandOverCoarseEFive}{0.97\,\ensuremath{\pm}\,0.10}  % experiments/FIGS/data/v2/index.json :: fine-tuning top N=100000: R30 of each random partition over the R16_Q1 run of its seed :: sha256 c6c4138ee5c86423
\newcommand{\MlstFtTopRthirtyRandOverCoarseEFivePr}{equal within 10\%}  % experiments/FIGS/data/v2/index.json :: fine-tuning top N=100000: R30 of each random partition over the R16_Q1 run of its seed :: sha256 c6c4138ee5c86423 :: prose twin: two significant figures, direction in words
\newcommand{\MlstFtTopRthirtyScratch}{1{,}412\,\ensuremath{\pm}\,20}  % experiments/FIGS/data/v2/index.json :: fine-tuning top N=1200000: R30 from scratch, every seed :: sha256 c6c4138ee5c86423
\newcommand{\MlstFtTopRthirtyFull}{1{,}800\,\ensuremath{\pm}\,60}  % experiments/FIGS/data/v2/index.json :: fine-tuning top N=1200000: R30, L188 :: sha256 c6c4138ee5c86423
\newcommand{\MlstFtTopRthirtyFullOverScratch}{1.27\,\ensuremath{\pm}\,0.04}  % experiments/FIGS/data/v2/index.json :: fine-tuning top N=1200000: R30 of each L188 run over the from-scratch mean, SD over runs and standard error of that mean in quadrature :: sha256 c6c4138ee5c86423
\newcommand{\MlstFtTopRthirtyFullOverScratchPr}{27\% more rejection}  % experiments/FIGS/data/v2/index.json :: fine-tuning top N=1200000: R30 of each L188 run over the from-scratch mean, SD over runs and standard error of that mean in quadrature :: sha256 c6c4138ee5c86423 :: prose twin: two significant figures, direction in words
\newcommand{\MlstFtTopRthirtyRes}{1{,}790\,\ensuremath{\pm}\,60}  % experiments/FIGS/data/v2/index.json :: fine-tuning top N=1200000: R30, L162 :: sha256 c6c4138ee5c86423
\newcommand{\MlstFtTopRthirtyResOverScratch}{1.27\,\ensuremath{\pm}\,0.04}  % experiments/FIGS/data/v2/index.json :: fine-tuning top N=1200000: R30 of each L162 run over the from-scratch mean, SD over runs and standard error of that mean in quadrature :: sha256 c6c4138ee5c86423
\newcommand{\MlstFtTopRthirtyResOverScratchPr}{27\% more rejection}  % experiments/FIGS/data/v2/index.json :: fine-tuning top N=1200000: R30 of each L162 run over the from-scratch mean, SD over runs and standard error of that mean in quadrature :: sha256 c6c4138ee5c86423 :: prose twin: two significant figures, direction in words
\newcommand{\MlstFtTopRthirtyMid}{1{,}780\,\ensuremath{\pm}\,150}  % experiments/FIGS/data/v2/index.json :: fine-tuning top N=1200000: R30, R42_Q1 :: sha256 c6c4138ee5c86423
\newcommand{\MlstFtTopRthirtyMidOverScratch}{1.26\,\ensuremath{\pm}\,0.11}  % experiments/FIGS/data/v2/index.json :: fine-tuning top N=1200000: R30 of each R42_Q1 run over the from-scratch mean, SD over runs and standard error of that mean in quadrature :: sha256 c6c4138ee5c86423
\newcommand{\MlstFtTopRthirtyMidOverScratchPr}{26\% more rejection}  % experiments/FIGS/data/v2/index.json :: fine-tuning top N=1200000: R30 of each R42_Q1 run over the from-scratch mean, SD over runs and standard error of that mean in quadrature :: sha256 c6c4138ee5c86423 :: prose twin: two significant figures, direction in words
\newcommand{\MlstFtTopRthirtyCoarse}{1{,}740\,\ensuremath{\pm}\,140}  % experiments/FIGS/data/v2/index.json :: fine-tuning top N=1200000: R30, R16_Q1 :: sha256 c6c4138ee5c86423
\newcommand{\MlstFtTopRthirtyCoarseOverScratch}{1.23\,\ensuremath{\pm}\,0.10}  % experiments/FIGS/data/v2/index.json :: fine-tuning top N=1200000: R30 of each R16_Q1 run over the from-scratch mean, SD over runs and standard error of that mean in quadrature :: sha256 c6c4138ee5c86423
\newcommand{\MlstFtTopRthirtyCoarseOverScratchPr}{23\% more rejection}  % experiments/FIGS/data/v2/index.json :: fine-tuning top N=1200000: R30 of each R16_Q1 run over the from-scratch mean, SD over runs and standard error of that mean in quadrature :: sha256 c6c4138ee5c86423 :: prose twin: two significant figures, direction in words
\newcommand{\MlstFtTopRthirtyResMass}{1{,}941 and 1{,}590}  % experiments/FIGS/data/v2/index.json :: fine-tuning top N=1200000: R30, L162_MASS :: sha256 c6c4138ee5c86423
\newcommand{\MlstFtTopRthirtyResMassOverScratch}{1.37 and 1.13}  % experiments/FIGS/data/v2/index.json :: fine-tuning top N=1200000: R30 of each L162_MASS run over the from-scratch mean, SD over runs and standard error of that mean in quadrature :: sha256 c6c4138ee5c86423
\newcommand{\MlstFtTopRthirtyResMassOverScratchPr}{at least 13\% more rejection}  % experiments/FIGS/data/v2/index.json :: fine-tuning top N=1200000: R30 of each L162_MASS run over the from-scratch mean, SD over runs and standard error of that mean in quadrature :: sha256 c6c4138ee5c86423 :: prose twin: two significant figures, direction in words
\newcommand{\MlstFtTopRthirtyResMassOverScratchPrMax}{up to 37\% more rejection}  % experiments/FIGS/data/v2/index.json :: fine-tuning top N=1200000: R30 of each L162_MASS run over the from-scratch mean, SD over runs and standard error of that mean in quadrature :: sha256 c6c4138ee5c86423 :: prose twin of the run farther from 1
\newcommand{\MlstFtTopRthirtyCoarseMass}{1{,}683 and 1{,}628}  % experiments/FIGS/data/v2/index.json :: fine-tuning top N=1200000: R30, R16_Q1_MASS :: sha256 c6c4138ee5c86423
\newcommand{\MlstFtTopRthirtyCoarseMassOverScratch}{1.192 and 1.153}  % experiments/FIGS/data/v2/index.json :: fine-tuning top N=1200000: R30 of each R16_Q1_MASS run over the from-scratch mean, SD over runs and standard error of that mean in quadrature :: sha256 c6c4138ee5c86423
\newcommand{\MlstFtTopRthirtyCoarseMassOverScratchPr}{at least 15\% more rejection}  % experiments/FIGS/data/v2/index.json :: fine-tuning top N=1200000: R30 of each R16_Q1_MASS run over the from-scratch mean, SD over runs and standard error of that mean in quadrature :: sha256 c6c4138ee5c86423 :: prose twin: two significant figures, direction in words
\newcommand{\MlstFtTopRthirtyCoarseMassOverScratchPrMax}{up to 19\% more rejection}  % experiments/FIGS/data/v2/index.json :: fine-tuning top N=1200000: R30 of each R16_Q1_MASS run over the from-scratch mean, SD over runs and standard error of that mean in quadrature :: sha256 c6c4138ee5c86423 :: prose twin of the run farther from 1
\newcommand{\MlstFtTopRthirtyResOverFull}{0.997\,\ensuremath{\pm}\,0.032}  % experiments/FIGS/data/v2/index.json :: fine-tuning top N=1200000: R30 of L162 over L188, run by run :: sha256 c6c4138ee5c86423
\newcommand{\MlstFtTopRthirtyResOverFullPr}{equal within 4\%}  % experiments/FIGS/data/v2/index.json :: fine-tuning top N=1200000: R30 of L162 over L188, run by run :: sha256 c6c4138ee5c86423 :: prose twin: two significant figures, direction in words
\newcommand{\MlstFtTopRthirtyMidOverFull}{0.99\,\ensuremath{\pm}\,0.07}  % experiments/FIGS/data/v2/index.json :: fine-tuning top N=1200000: R30 of R42_Q1 over L188, run by run :: sha256 c6c4138ee5c86423
\newcommand{\MlstFtTopRthirtyMidOverFullPr}{equal within 8\%}  % experiments/FIGS/data/v2/index.json :: fine-tuning top N=1200000: R30 of R42_Q1 over L188, run by run :: sha256 c6c4138ee5c86423 :: prose twin: two significant figures, direction in words
\newcommand{\MlstFtTopRthirtyCoarseOverFull}{0.97\,\ensuremath{\pm}\,0.10}  % experiments/FIGS/data/v2/index.json :: fine-tuning top N=1200000: R30 of R16_Q1 over L188, run by run :: sha256 c6c4138ee5c86423
\newcommand{\MlstFtTopRthirtyCoarseOverFullPr}{equal within 10\%}  % experiments/FIGS/data/v2/index.json :: fine-tuning top N=1200000: R30 of R16_Q1 over L188, run by run :: sha256 c6c4138ee5c86423 :: prose twin: two significant figures, direction in words
\newcommand{\MlstFtTopRthirtyResMassOverRes}{1.12 and 0.89}  % experiments/FIGS/data/v2/index.json :: fine-tuning top N=1200000: R30 of L162_MASS over L162, run by run :: sha256 c6c4138ee5c86423
\newcommand{\MlstFtTopRthirtyResMassOverResPr}{equal within 20\%}  % experiments/FIGS/data/v2/index.json :: fine-tuning top N=1200000: R30 of L162_MASS over L162, run by run :: sha256 c6c4138ee5c86423 :: prose twin: two significant figures, direction in words
\newcommand{\MlstFtTopRthirtyResMassOverResPrMax}{equal within 20\%}  % experiments/FIGS/data/v2/index.json :: fine-tuning top N=1200000: R30 of L162_MASS over L162, run by run :: sha256 c6c4138ee5c86423 :: prose twin of the run farther from 1
\newcommand{\MlstFtTopRthirtyCoarseMassOverCoarse}{0.900 and 0.919}  % experiments/FIGS/data/v2/index.json :: fine-tuning top N=1200000: R30 of R16_Q1_MASS over R16_Q1, run by run :: sha256 c6c4138ee5c86423
\newcommand{\MlstFtTopRthirtyCoarseMassOverCoarsePr}{equal within 10\%}  % experiments/FIGS/data/v2/index.json :: fine-tuning top N=1200000: R30 of R16_Q1_MASS over R16_Q1, run by run :: sha256 c6c4138ee5c86423 :: prose twin: two significant figures, direction in words
\newcommand{\MlstFtTopRthirtyCoarseMassOverCoarsePrMax}{equal within 10\%}  % experiments/FIGS/data/v2/index.json :: fine-tuning top N=1200000: R30 of R16_Q1_MASS over R16_Q1, run by run :: sha256 c6c4138ee5c86423 :: prose twin of the run farther from 1
\newcommand{\MlstFtTopRthirtyRandOverFull}{0.97\,\ensuremath{\pm}\,0.04}  % experiments/FIGS/data/v2/index.json :: fine-tuning top N=1200000: R30 of each random partition over the L188 run of its seed :: sha256 c6c4138ee5c86423
\newcommand{\MlstFtTopRthirtyRandOverFullPr}{equal within 4\%}  % experiments/FIGS/data/v2/index.json :: fine-tuning top N=1200000: R30 of each random partition over the L188 run of its seed :: sha256 c6c4138ee5c86423 :: prose twin: two significant figures, direction in words
\newcommand{\MlstFtTopRthirtyRandOverCoarse}{0.97\,\ensuremath{\pm}\,0.09}  % experiments/FIGS/data/v2/index.json :: fine-tuning top N=1200000: R30 of each random partition over the R16_Q1 run of its seed :: sha256 c6c4138ee5c86423
\newcommand{\MlstFtTopRthirtyRandOverCoarsePr}{equal within 9\%}  % experiments/FIGS/data/v2/index.json :: fine-tuning top N=1200000: R30 of each random partition over the R16_Q1 run of its seed :: sha256 c6c4138ee5c86423 :: prose twin: two significant figures, direction in words
\newcommand{\MlstFtTopRthirtyFullOverPart}{0.650\,\ensuremath{\pm}\,0.020}  % experiments/FIGS/data/v2/index.json :: fine-tuning top N=1200000: R30 of each 188-class run over ParT-f.t.'s published value :: sha256 c6c4138ee5c86423
\newcommand{\MlstFtTopRthirtyFullOverPartPr}{35\% less rejection}  % experiments/FIGS/data/v2/index.json :: fine-tuning top N=1200000: R30 of each 188-class run over ParT-f.t.'s published value :: sha256 c6c4138ee5c86423 :: prose twin: two significant figures, direction in words
\newcommand{\MlstFtTopAucFull}{0.9861\,\ensuremath{\pm}\,0.0001}  % experiments/FIGS/data/v2/index.json :: fine-tuning top N=1200000: AUC of the 188-class runs :: sha256 c6c4138ee5c86423
\newcommand{\MlstFtTopAucScratch}{0.9852\,\ensuremath{\pm}\,0.0001}  % experiments/FIGS/data/v2/index.json :: fine-tuning top N=1200000: AUC from scratch :: sha256 c6c4138ee5c86423
\newcommand{\MlstFtTopRthirtyFullRelSd}{3\%}  % experiments/FIGS/data/v2/index.json :: fine-tuning top, full set: relative SD of R30 over the 188-class runs :: sha256 c6c4138ee5c86423
\newcommand{\MlstFtTopRthirtyKin}{1{,}800}  % experiments/FIGS/data/v2/index.json :: fine-tuning top, full set: R30 of the kinematic-only model (one run) :: sha256 c6c4138ee5c86423
\newcommand{\MlstFtTopRthirtyKinOverFull}{1.00}  % experiments/FIGS/data/v2/index.json :: fine-tuning top, full set: R30 of the kinematic-only model over the mean of the 188-class runs :: sha256 c6c4138ee5c86423
\newcommand{\MlstFtTopRthirtyKinOverFullPr}{0.23\% more rejection}  % experiments/FIGS/data/v2/index.json :: fine-tuning top, full set: R30 of the kinematic-only model over the mean of the 188-class runs :: sha256 c6c4138ee5c86423 :: prose twin: two significant figures, direction in words
\newcommand{\MlstFtTopDataEquivJets}{1{,}200}  % experiments/FIGS/data/v2/index.json :: fine-tuning top: from-scratch training jets whose R30 equals the 188-class mean at 1000 jets (interpolated linearly in R30 against log10 N) :: sha256 c6c4138ee5c86423
\newcommand{\MlstFtTopDataEquivJetsEFour}{23{,}000}  % experiments/FIGS/data/v2/index.json :: fine-tuning top: from-scratch training jets whose R30 equals the 188-class mean at 10000 jets (interpolated linearly in R30 against log10 N) :: sha256 c6c4138ee5c86423
\newcommand{\MlstFtQgRthirtyScratchEThree}{49.7\,\ensuremath{\pm}\,2.1}  % experiments/FIGS/data/v2/index.json :: fine-tuning qg N=1000: R30 from scratch, every seed :: sha256 c6c4138ee5c86423
\newcommand{\MlstFtQgRthirtyFullEThree}{47\,\ensuremath{\pm}\,10}  % experiments/FIGS/data/v2/index.json :: fine-tuning qg N=1000: R30, L188 :: sha256 c6c4138ee5c86423
\newcommand{\MlstFtQgRthirtyFullOverScratchEThree}{0.95\,\ensuremath{\pm}\,0.19}  % experiments/FIGS/data/v2/index.json :: fine-tuning qg N=1000: R30 of each L188 run over the from-scratch mean, SD over runs and standard error of that mean in quadrature :: sha256 c6c4138ee5c86423
\newcommand{\MlstFtQgRthirtyFullOverScratchEThreePr}{equal within 20\%}  % experiments/FIGS/data/v2/index.json :: fine-tuning qg N=1000: R30 of each L188 run over the from-scratch mean, SD over runs and standard error of that mean in quadrature :: sha256 c6c4138ee5c86423 :: prose twin: two significant figures, direction in words
\newcommand{\MlstFtQgRthirtyResEThree}{52.4\,\ensuremath{\pm}\,3.1}  % experiments/FIGS/data/v2/index.json :: fine-tuning qg N=1000: R30, L162 :: sha256 c6c4138ee5c86423
\newcommand{\MlstFtQgRthirtyResOverScratchEThree}{1.06\,\ensuremath{\pm}\,0.07}  % experiments/FIGS/data/v2/index.json :: fine-tuning qg N=1000: R30 of each L162 run over the from-scratch mean, SD over runs and standard error of that mean in quadrature :: sha256 c6c4138ee5c86423
\newcommand{\MlstFtQgRthirtyResOverScratchEThreePr}{equal within 7\%}  % experiments/FIGS/data/v2/index.json :: fine-tuning qg N=1000: R30 of each L162 run over the from-scratch mean, SD over runs and standard error of that mean in quadrature :: sha256 c6c4138ee5c86423 :: prose twin: two significant figures, direction in words
\newcommand{\MlstFtQgRthirtyMidEThree}{54\,\ensuremath{\pm}\,6}  % experiments/FIGS/data/v2/index.json :: fine-tuning qg N=1000: R30, R42_Q1 :: sha256 c6c4138ee5c86423
\newcommand{\MlstFtQgRthirtyMidOverScratchEThree}{1.09\,\ensuremath{\pm}\,0.12}  % experiments/FIGS/data/v2/index.json :: fine-tuning qg N=1000: R30 of each R42_Q1 run over the from-scratch mean, SD over runs and standard error of that mean in quadrature :: sha256 c6c4138ee5c86423
\newcommand{\MlstFtQgRthirtyMidOverScratchEThreePr}{equal within 20\%}  % experiments/FIGS/data/v2/index.json :: fine-tuning qg N=1000: R30 of each R42_Q1 run over the from-scratch mean, SD over runs and standard error of that mean in quadrature :: sha256 c6c4138ee5c86423 :: prose twin: two significant figures, direction in words
\newcommand{\MlstFtQgRthirtyCoarseEThree}{55\,\ensuremath{\pm}\,5}  % experiments/FIGS/data/v2/index.json :: fine-tuning qg N=1000: R30, R16_Q1 :: sha256 c6c4138ee5c86423
\newcommand{\MlstFtQgRthirtyCoarseOverScratchEThree}{1.12\,\ensuremath{\pm}\,0.10}  % experiments/FIGS/data/v2/index.json :: fine-tuning qg N=1000: R30 of each R16_Q1 run over the from-scratch mean, SD over runs and standard error of that mean in quadrature :: sha256 c6c4138ee5c86423
\newcommand{\MlstFtQgRthirtyCoarseOverScratchEThreePr}{equal within 20\%}  % experiments/FIGS/data/v2/index.json :: fine-tuning qg N=1000: R30 of each R16_Q1 run over the from-scratch mean, SD over runs and standard error of that mean in quadrature :: sha256 c6c4138ee5c86423 :: prose twin: two significant figures, direction in words
\newcommand{\MlstFtQgRthirtyResMassEThree}{59.3 and 57.1}  % experiments/FIGS/data/v2/index.json :: fine-tuning qg N=1000: R30, L162_MASS :: sha256 c6c4138ee5c86423
\newcommand{\MlstFtQgRthirtyResMassOverScratchEThree}{1.193 and 1.149}  % experiments/FIGS/data/v2/index.json :: fine-tuning qg N=1000: R30 of each L162_MASS run over the from-scratch mean, SD over runs and standard error of that mean in quadrature :: sha256 c6c4138ee5c86423
\newcommand{\MlstFtQgRthirtyResMassOverScratchEThreePr}{at least 15\% more rejection}  % experiments/FIGS/data/v2/index.json :: fine-tuning qg N=1000: R30 of each L162_MASS run over the from-scratch mean, SD over runs and standard error of that mean in quadrature :: sha256 c6c4138ee5c86423 :: prose twin: two significant figures, direction in words
\newcommand{\MlstFtQgRthirtyResMassOverScratchEThreePrMax}{up to 19\% more rejection}  % experiments/FIGS/data/v2/index.json :: fine-tuning qg N=1000: R30 of each L162_MASS run over the from-scratch mean, SD over runs and standard error of that mean in quadrature :: sha256 c6c4138ee5c86423 :: prose twin of the run farther from 1
\newcommand{\MlstFtQgRthirtyCoarseMassEThree}{54.2 and 57.7}  % experiments/FIGS/data/v2/index.json :: fine-tuning qg N=1000: R30, R16_Q1_MASS :: sha256 c6c4138ee5c86423
\newcommand{\MlstFtQgRthirtyCoarseMassOverScratchEThree}{1.091 and 1.161}  % experiments/FIGS/data/v2/index.json :: fine-tuning qg N=1000: R30 of each R16_Q1_MASS run over the from-scratch mean, SD over runs and standard error of that mean in quadrature :: sha256 c6c4138ee5c86423
\newcommand{\MlstFtQgRthirtyCoarseMassOverScratchEThreePr}{at least 9.1\% more rejection}  % experiments/FIGS/data/v2/index.json :: fine-tuning qg N=1000: R30 of each R16_Q1_MASS run over the from-scratch mean, SD over runs and standard error of that mean in quadrature :: sha256 c6c4138ee5c86423 :: prose twin: two significant figures, direction in words
\newcommand{\MlstFtQgRthirtyCoarseMassOverScratchEThreePrMax}{up to 16\% more rejection}  % experiments/FIGS/data/v2/index.json :: fine-tuning qg N=1000: R30 of each R16_Q1_MASS run over the from-scratch mean, SD over runs and standard error of that mean in quadrature :: sha256 c6c4138ee5c86423 :: prose twin of the run farther from 1
\newcommand{\MlstFtQgRthirtyResOverFullEThree}{1.14\,\ensuremath{\pm}\,0.17}  % experiments/FIGS/data/v2/index.json :: fine-tuning qg N=1000: R30 of L162 over L188, run by run :: sha256 c6c4138ee5c86423
\newcommand{\MlstFtQgRthirtyResOverFullEThreePr}{equal within 20\%}  % experiments/FIGS/data/v2/index.json :: fine-tuning qg N=1000: R30 of L162 over L188, run by run :: sha256 c6c4138ee5c86423 :: prose twin: two significant figures, direction in words
\newcommand{\MlstFtQgRthirtyMidOverFullEThree}{1.17\,\ensuremath{\pm}\,0.14}  % experiments/FIGS/data/v2/index.json :: fine-tuning qg N=1000: R30 of R42_Q1 over L188, run by run :: sha256 c6c4138ee5c86423
\newcommand{\MlstFtQgRthirtyMidOverFullEThreePr}{equal within 20\%}  % experiments/FIGS/data/v2/index.json :: fine-tuning qg N=1000: R30 of R42_Q1 over L188, run by run :: sha256 c6c4138ee5c86423 :: prose twin: two significant figures, direction in words
\newcommand{\MlstFtQgRthirtyCoarseOverFullEThree}{1.21\,\ensuremath{\pm}\,0.27}  % experiments/FIGS/data/v2/index.json :: fine-tuning qg N=1000: R30 of R16_Q1 over L188, run by run :: sha256 c6c4138ee5c86423
\newcommand{\MlstFtQgRthirtyCoarseOverFullEThreePr}{equal within 30\%}  % experiments/FIGS/data/v2/index.json :: fine-tuning qg N=1000: R30 of R16_Q1 over L188, run by run :: sha256 c6c4138ee5c86423 :: prose twin: two significant figures, direction in words
\newcommand{\MlstFtQgRthirtyResMassOverResEThree}{1.108 and 1.042}  % experiments/FIGS/data/v2/index.json :: fine-tuning qg N=1000: R30 of L162_MASS over L162, run by run :: sha256 c6c4138ee5c86423
\newcommand{\MlstFtQgRthirtyResMassOverResEThreePr}{equal within 20\%}  % experiments/FIGS/data/v2/index.json :: fine-tuning qg N=1000: R30 of L162_MASS over L162, run by run :: sha256 c6c4138ee5c86423 :: prose twin: two significant figures, direction in words
\newcommand{\MlstFtQgRthirtyResMassOverResEThreePrMax}{equal within 20\%}  % experiments/FIGS/data/v2/index.json :: fine-tuning qg N=1000: R30 of L162_MASS over L162, run by run :: sha256 c6c4138ee5c86423 :: prose twin of the run farther from 1
\newcommand{\MlstFtQgRthirtyCoarseMassOverCoarseEThree}{0.89 and 1.11}  % experiments/FIGS/data/v2/index.json :: fine-tuning qg N=1000: R30 of R16_Q1_MASS over R16_Q1, run by run :: sha256 c6c4138ee5c86423
\newcommand{\MlstFtQgRthirtyCoarseMassOverCoarseEThreePr}{equal within 20\%}  % experiments/FIGS/data/v2/index.json :: fine-tuning qg N=1000: R30 of R16_Q1_MASS over R16_Q1, run by run :: sha256 c6c4138ee5c86423 :: prose twin: two significant figures, direction in words
\newcommand{\MlstFtQgRthirtyCoarseMassOverCoarseEThreePrMax}{equal within 20\%}  % experiments/FIGS/data/v2/index.json :: fine-tuning qg N=1000: R30 of R16_Q1_MASS over R16_Q1, run by run :: sha256 c6c4138ee5c86423 :: prose twin of the run farther from 1
\newcommand{\MlstFtQgRthirtyRandOverFullEThree}{0.91\,\ensuremath{\pm}\,0.07}  % experiments/FIGS/data/v2/index.json :: fine-tuning qg N=1000: R30 of each random partition over the L188 run of its seed :: sha256 c6c4138ee5c86423
\newcommand{\MlstFtQgRthirtyRandOverFullEThreePr}{equal within 10\%}  % experiments/FIGS/data/v2/index.json :: fine-tuning qg N=1000: R30 of each random partition over the L188 run of its seed :: sha256 c6c4138ee5c86423 :: prose twin: two significant figures, direction in words
\newcommand{\MlstFtQgRthirtyRandOverCoarseEThree}{0.87\,\ensuremath{\pm}\,0.14}  % experiments/FIGS/data/v2/index.json :: fine-tuning qg N=1000: R30 of each random partition over the R16_Q1 run of its seed :: sha256 c6c4138ee5c86423
\newcommand{\MlstFtQgRthirtyRandOverCoarseEThreePr}{equal within 20\%}  % experiments/FIGS/data/v2/index.json :: fine-tuning qg N=1000: R30 of each random partition over the R16_Q1 run of its seed :: sha256 c6c4138ee5c86423 :: prose twin: two significant figures, direction in words
\newcommand{\MlstFtQgRthirtyScratchEFour}{66\,\ensuremath{\pm}\,9}  % experiments/FIGS/data/v2/index.json :: fine-tuning qg N=10000: R30 from scratch, every seed :: sha256 c6c4138ee5c86423
\newcommand{\MlstFtQgRthirtyFullEFour}{75\,\ensuremath{\pm}\,7}  % experiments/FIGS/data/v2/index.json :: fine-tuning qg N=10000: R30, L188 :: sha256 c6c4138ee5c86423
\newcommand{\MlstFtQgRthirtyFullOverScratchEFour}{1.14\,\ensuremath{\pm}\,0.14}  % experiments/FIGS/data/v2/index.json :: fine-tuning qg N=10000: R30 of each L188 run over the from-scratch mean, SD over runs and standard error of that mean in quadrature :: sha256 c6c4138ee5c86423
\newcommand{\MlstFtQgRthirtyFullOverScratchEFourPr}{equal within 20\%}  % experiments/FIGS/data/v2/index.json :: fine-tuning qg N=10000: R30 of each L188 run over the from-scratch mean, SD over runs and standard error of that mean in quadrature :: sha256 c6c4138ee5c86423 :: prose twin: two significant figures, direction in words
\newcommand{\MlstFtQgRthirtyResEFour}{76.4\,\ensuremath{\pm}\,0.6}  % experiments/FIGS/data/v2/index.json :: fine-tuning qg N=10000: R30, L162 :: sha256 c6c4138ee5c86423
\newcommand{\MlstFtQgRthirtyResOverScratchEFour}{1.15\,\ensuremath{\pm}\,0.09}  % experiments/FIGS/data/v2/index.json :: fine-tuning qg N=10000: R30 of each L162 run over the from-scratch mean, SD over runs and standard error of that mean in quadrature :: sha256 c6c4138ee5c86423
\newcommand{\MlstFtQgRthirtyResOverScratchEFourPr}{15\% more rejection}  % experiments/FIGS/data/v2/index.json :: fine-tuning qg N=10000: R30 of each L162 run over the from-scratch mean, SD over runs and standard error of that mean in quadrature :: sha256 c6c4138ee5c86423 :: prose twin: two significant figures, direction in words
\newcommand{\MlstFtQgRthirtyMidEFour}{74\,\ensuremath{\pm}\,5}  % experiments/FIGS/data/v2/index.json :: fine-tuning qg N=10000: R30, R42_Q1 :: sha256 c6c4138ee5c86423
\newcommand{\MlstFtQgRthirtyMidOverScratchEFour}{1.12\,\ensuremath{\pm}\,0.11}  % experiments/FIGS/data/v2/index.json :: fine-tuning qg N=10000: R30 of each R42_Q1 run over the from-scratch mean, SD over runs and standard error of that mean in quadrature :: sha256 c6c4138ee5c86423
\newcommand{\MlstFtQgRthirtyMidOverScratchEFourPr}{equal within 20\%}  % experiments/FIGS/data/v2/index.json :: fine-tuning qg N=10000: R30 of each R42_Q1 run over the from-scratch mean, SD over runs and standard error of that mean in quadrature :: sha256 c6c4138ee5c86423 :: prose twin: two significant figures, direction in words
\newcommand{\MlstFtQgRthirtyCoarseEFour}{80.3\,\ensuremath{\pm}\,1.1}  % experiments/FIGS/data/v2/index.json :: fine-tuning qg N=10000: R30, R16_Q1 :: sha256 c6c4138ee5c86423
\newcommand{\MlstFtQgRthirtyCoarseOverScratchEFour}{1.21\,\ensuremath{\pm}\,0.10}  % experiments/FIGS/data/v2/index.json :: fine-tuning qg N=10000: R30 of each R16_Q1 run over the from-scratch mean, SD over runs and standard error of that mean in quadrature :: sha256 c6c4138ee5c86423
\newcommand{\MlstFtQgRthirtyCoarseOverScratchEFourPr}{21\% more rejection}  % experiments/FIGS/data/v2/index.json :: fine-tuning qg N=10000: R30 of each R16_Q1 run over the from-scratch mean, SD over runs and standard error of that mean in quadrature :: sha256 c6c4138ee5c86423 :: prose twin: two significant figures, direction in words
\newcommand{\MlstFtQgRthirtyResMassEFour}{80.1 and 72.6}  % experiments/FIGS/data/v2/index.json :: fine-tuning qg N=10000: R30, L162_MASS :: sha256 c6c4138ee5c86423
\newcommand{\MlstFtQgRthirtyResMassOverScratchEFour}{1.21 and 1.09}  % experiments/FIGS/data/v2/index.json :: fine-tuning qg N=10000: R30 of each L162_MASS run over the from-scratch mean, SD over runs and standard error of that mean in quadrature :: sha256 c6c4138ee5c86423
\newcommand{\MlstFtQgRthirtyResMassOverScratchEFourPr}{at least 9.3\% more rejection}  % experiments/FIGS/data/v2/index.json :: fine-tuning qg N=10000: R30 of each L162_MASS run over the from-scratch mean, SD over runs and standard error of that mean in quadrature :: sha256 c6c4138ee5c86423 :: prose twin: two significant figures, direction in words
\newcommand{\MlstFtQgRthirtyResMassOverScratchEFourPrMax}{up to 21\% more rejection}  % experiments/FIGS/data/v2/index.json :: fine-tuning qg N=10000: R30 of each L162_MASS run over the from-scratch mean, SD over runs and standard error of that mean in quadrature :: sha256 c6c4138ee5c86423 :: prose twin of the run farther from 1
\newcommand{\MlstFtQgRthirtyCoarseMassEFour}{79.5 and 77.9}  % experiments/FIGS/data/v2/index.json :: fine-tuning qg N=10000: R30, R16_Q1_MASS :: sha256 c6c4138ee5c86423
\newcommand{\MlstFtQgRthirtyCoarseMassOverScratchEFour}{1.198 and 1.173}  % experiments/FIGS/data/v2/index.json :: fine-tuning qg N=10000: R30 of each R16_Q1_MASS run over the from-scratch mean, SD over runs and standard error of that mean in quadrature :: sha256 c6c4138ee5c86423
\newcommand{\MlstFtQgRthirtyCoarseMassOverScratchEFourPr}{at least 17\% more rejection}  % experiments/FIGS/data/v2/index.json :: fine-tuning qg N=10000: R30 of each R16_Q1_MASS run over the from-scratch mean, SD over runs and standard error of that mean in quadrature :: sha256 c6c4138ee5c86423 :: prose twin: two significant figures, direction in words
\newcommand{\MlstFtQgRthirtyCoarseMassOverScratchEFourPrMax}{up to 20\% more rejection}  % experiments/FIGS/data/v2/index.json :: fine-tuning qg N=10000: R30 of each R16_Q1_MASS run over the from-scratch mean, SD over runs and standard error of that mean in quadrature :: sha256 c6c4138ee5c86423 :: prose twin of the run farther from 1
\newcommand{\MlstFtQgRthirtyResOverFullEFour}{1.02\,\ensuremath{\pm}\,0.10}  % experiments/FIGS/data/v2/index.json :: fine-tuning qg N=10000: R30 of L162 over L188, run by run :: sha256 c6c4138ee5c86423
\newcommand{\MlstFtQgRthirtyResOverFullEFourPr}{equal within 10\%}  % experiments/FIGS/data/v2/index.json :: fine-tuning qg N=10000: R30 of L162 over L188, run by run :: sha256 c6c4138ee5c86423 :: prose twin: two significant figures, direction in words
\newcommand{\MlstFtQgRthirtyMidOverFullEFour}{0.99\,\ensuremath{\pm}\,0.14}  % experiments/FIGS/data/v2/index.json :: fine-tuning qg N=10000: R30 of R42_Q1 over L188, run by run :: sha256 c6c4138ee5c86423
\newcommand{\MlstFtQgRthirtyMidOverFullEFourPr}{equal within 20\%}  % experiments/FIGS/data/v2/index.json :: fine-tuning qg N=10000: R30 of R42_Q1 over L188, run by run :: sha256 c6c4138ee5c86423 :: prose twin: two significant figures, direction in words
\newcommand{\MlstFtQgRthirtyCoarseOverFullEFour}{1.07\,\ensuremath{\pm}\,0.13}  % experiments/FIGS/data/v2/index.json :: fine-tuning qg N=10000: R30 of R16_Q1 over L188, run by run :: sha256 c6c4138ee5c86423
\newcommand{\MlstFtQgRthirtyCoarseOverFullEFourPr}{equal within 20\%}  % experiments/FIGS/data/v2/index.json :: fine-tuning qg N=10000: R30 of R16_Q1 over L188, run by run :: sha256 c6c4138ee5c86423 :: prose twin: two significant figures, direction in words
\newcommand{\MlstFtQgRthirtyResMassOverResEFour}{1.06 and 0.95}  % experiments/FIGS/data/v2/index.json :: fine-tuning qg N=10000: R30 of L162_MASS over L162, run by run :: sha256 c6c4138ee5c86423
\newcommand{\MlstFtQgRthirtyResMassOverResEFourPr}{equal within 6\%}  % experiments/FIGS/data/v2/index.json :: fine-tuning qg N=10000: R30 of L162_MASS over L162, run by run :: sha256 c6c4138ee5c86423 :: prose twin: two significant figures, direction in words
\newcommand{\MlstFtQgRthirtyResMassOverResEFourPrMax}{equal within 6\%}  % experiments/FIGS/data/v2/index.json :: fine-tuning qg N=10000: R30 of L162_MASS over L162, run by run :: sha256 c6c4138ee5c86423 :: prose twin of the run farther from 1
\newcommand{\MlstFtQgRthirtyCoarseMassOverCoarseEFour}{0.9769 and 0.9696}  % experiments/FIGS/data/v2/index.json :: fine-tuning qg N=10000: R30 of R16_Q1_MASS over R16_Q1, run by run :: sha256 c6c4138ee5c86423
\newcommand{\MlstFtQgRthirtyCoarseMassOverCoarseEFourPr}{equal within 4\%}  % experiments/FIGS/data/v2/index.json :: fine-tuning qg N=10000: R30 of R16_Q1_MASS over R16_Q1, run by run :: sha256 c6c4138ee5c86423 :: prose twin: two significant figures, direction in words
\newcommand{\MlstFtQgRthirtyCoarseMassOverCoarseEFourPrMax}{equal within 4\%}  % experiments/FIGS/data/v2/index.json :: fine-tuning qg N=10000: R30 of R16_Q1_MASS over R16_Q1, run by run :: sha256 c6c4138ee5c86423 :: prose twin of the run farther from 1
\newcommand{\MlstFtQgRthirtyRandOverFullEFour}{1.03\,\ensuremath{\pm}\,0.18}  % experiments/FIGS/data/v2/index.json :: fine-tuning qg N=10000: R30 of each random partition over the L188 run of its seed :: sha256 c6c4138ee5c86423
\newcommand{\MlstFtQgRthirtyRandOverFullEFourPr}{equal within 20\%}  % experiments/FIGS/data/v2/index.json :: fine-tuning qg N=10000: R30 of each random partition over the L188 run of its seed :: sha256 c6c4138ee5c86423 :: prose twin: two significant figures, direction in words
\newcommand{\MlstFtQgRthirtyRandOverCoarseEFour}{0.93\,\ensuremath{\pm}\,0.12}  % experiments/FIGS/data/v2/index.json :: fine-tuning qg N=10000: R30 of each random partition over the R16_Q1 run of its seed :: sha256 c6c4138ee5c86423
\newcommand{\MlstFtQgRthirtyRandOverCoarseEFourPr}{equal within 20\%}  % experiments/FIGS/data/v2/index.json :: fine-tuning qg N=10000: R30 of each random partition over the R16_Q1 run of its seed :: sha256 c6c4138ee5c86423 :: prose twin: two significant figures, direction in words
\newcommand{\MlstFtQgRthirtyScratchEFive}{88.4\,\ensuremath{\pm}\,0.8}  % experiments/FIGS/data/v2/index.json :: fine-tuning qg N=100000: R30 from scratch, every seed :: sha256 c6c4138ee5c86423
\newcommand{\MlstFtQgRthirtyFullEFive}{99.3\,\ensuremath{\pm}\,1.4}  % experiments/FIGS/data/v2/index.json :: fine-tuning qg N=100000: R30, L188 :: sha256 c6c4138ee5c86423
\newcommand{\MlstFtQgRthirtyFullOverScratchEFive}{1.124\,\ensuremath{\pm}\,0.017}  % experiments/FIGS/data/v2/index.json :: fine-tuning qg N=100000: R30 of each L188 run over the from-scratch mean, SD over runs and standard error of that mean in quadrature :: sha256 c6c4138ee5c86423
\newcommand{\MlstFtQgRthirtyFullOverScratchEFivePr}{12\% more rejection}  % experiments/FIGS/data/v2/index.json :: fine-tuning qg N=100000: R30 of each L188 run over the from-scratch mean, SD over runs and standard error of that mean in quadrature :: sha256 c6c4138ee5c86423 :: prose twin: two significant figures, direction in words
\newcommand{\MlstFtQgRthirtyResEFive}{97.9\,\ensuremath{\pm}\,3.3}  % experiments/FIGS/data/v2/index.json :: fine-tuning qg N=100000: R30, L162 :: sha256 c6c4138ee5c86423
\newcommand{\MlstFtQgRthirtyResOverScratchEFive}{1.11\,\ensuremath{\pm}\,0.04}  % experiments/FIGS/data/v2/index.json :: fine-tuning qg N=100000: R30 of each L162 run over the from-scratch mean, SD over runs and standard error of that mean in quadrature :: sha256 c6c4138ee5c86423
\newcommand{\MlstFtQgRthirtyResOverScratchEFivePr}{11\% more rejection}  % experiments/FIGS/data/v2/index.json :: fine-tuning qg N=100000: R30 of each L162 run over the from-scratch mean, SD over runs and standard error of that mean in quadrature :: sha256 c6c4138ee5c86423 :: prose twin: two significant figures, direction in words
\newcommand{\MlstFtQgRthirtyMidEFive}{96.1\,\ensuremath{\pm}\,1.8}  % experiments/FIGS/data/v2/index.json :: fine-tuning qg N=100000: R30, R42_Q1 :: sha256 c6c4138ee5c86423
\newcommand{\MlstFtQgRthirtyMidOverScratchEFive}{1.088\,\ensuremath{\pm}\,0.021}  % experiments/FIGS/data/v2/index.json :: fine-tuning qg N=100000: R30 of each R42_Q1 run over the from-scratch mean, SD over runs and standard error of that mean in quadrature :: sha256 c6c4138ee5c86423
\newcommand{\MlstFtQgRthirtyMidOverScratchEFivePr}{8.8\% more rejection}  % experiments/FIGS/data/v2/index.json :: fine-tuning qg N=100000: R30 of each R42_Q1 run over the from-scratch mean, SD over runs and standard error of that mean in quadrature :: sha256 c6c4138ee5c86423 :: prose twin: two significant figures, direction in words
\newcommand{\MlstFtQgRthirtyCoarseEFive}{96.4\,\ensuremath{\pm}\,1.0}  % experiments/FIGS/data/v2/index.json :: fine-tuning qg N=100000: R30, R16_Q1 :: sha256 c6c4138ee5c86423
\newcommand{\MlstFtQgRthirtyCoarseOverScratchEFive}{1.091\,\ensuremath{\pm}\,0.013}  % experiments/FIGS/data/v2/index.json :: fine-tuning qg N=100000: R30 of each R16_Q1 run over the from-scratch mean, SD over runs and standard error of that mean in quadrature :: sha256 c6c4138ee5c86423
\newcommand{\MlstFtQgRthirtyCoarseOverScratchEFivePr}{9.1\% more rejection}  % experiments/FIGS/data/v2/index.json :: fine-tuning qg N=100000: R30 of each R16_Q1 run over the from-scratch mean, SD over runs and standard error of that mean in quadrature :: sha256 c6c4138ee5c86423 :: prose twin: two significant figures, direction in words
\newcommand{\MlstFtQgRthirtyResMassEFive}{95.6 and 97.6}  % experiments/FIGS/data/v2/index.json :: fine-tuning qg N=100000: R30, L162_MASS :: sha256 c6c4138ee5c86423
\newcommand{\MlstFtQgRthirtyResMassOverScratchEFive}{1.082 and 1.104}  % experiments/FIGS/data/v2/index.json :: fine-tuning qg N=100000: R30 of each L162_MASS run over the from-scratch mean, SD over runs and standard error of that mean in quadrature :: sha256 c6c4138ee5c86423
\newcommand{\MlstFtQgRthirtyResMassOverScratchEFivePr}{at least 8.2\% more rejection}  % experiments/FIGS/data/v2/index.json :: fine-tuning qg N=100000: R30 of each L162_MASS run over the from-scratch mean, SD over runs and standard error of that mean in quadrature :: sha256 c6c4138ee5c86423 :: prose twin: two significant figures, direction in words
\newcommand{\MlstFtQgRthirtyResMassOverScratchEFivePrMax}{up to 10\% more rejection}  % experiments/FIGS/data/v2/index.json :: fine-tuning qg N=100000: R30 of each L162_MASS run over the from-scratch mean, SD over runs and standard error of that mean in quadrature :: sha256 c6c4138ee5c86423 :: prose twin of the run farther from 1
\newcommand{\MlstFtQgRthirtyCoarseMassEFive}{97.85 and 97.94}  % experiments/FIGS/data/v2/index.json :: fine-tuning qg N=100000: R30, R16_Q1_MASS :: sha256 c6c4138ee5c86423
\newcommand{\MlstFtQgRthirtyCoarseMassOverScratchEFive}{1.107 and 1.108}  % experiments/FIGS/data/v2/index.json :: fine-tuning qg N=100000: R30 of each R16_Q1_MASS run over the from-scratch mean, SD over runs and standard error of that mean in quadrature :: sha256 c6c4138ee5c86423
\newcommand{\MlstFtQgRthirtyCoarseMassOverScratchEFivePr}{at least 11\% more rejection}  % experiments/FIGS/data/v2/index.json :: fine-tuning qg N=100000: R30 of each R16_Q1_MASS run over the from-scratch mean, SD over runs and standard error of that mean in quadrature :: sha256 c6c4138ee5c86423 :: prose twin: two significant figures, direction in words
\newcommand{\MlstFtQgRthirtyCoarseMassOverScratchEFivePrMax}{up to 11\% more rejection}  % experiments/FIGS/data/v2/index.json :: fine-tuning qg N=100000: R30 of each R16_Q1_MASS run over the from-scratch mean, SD over runs and standard error of that mean in quadrature :: sha256 c6c4138ee5c86423 :: prose twin of the run farther from 1
\newcommand{\MlstFtQgRthirtyResOverFullEFive}{0.985\,\ensuremath{\pm}\,0.020}  % experiments/FIGS/data/v2/index.json :: fine-tuning qg N=100000: R30 of L162 over L188, run by run :: sha256 c6c4138ee5c86423
\newcommand{\MlstFtQgRthirtyResOverFullEFivePr}{equal within 2\%}  % experiments/FIGS/data/v2/index.json :: fine-tuning qg N=100000: R30 of L162 over L188, run by run :: sha256 c6c4138ee5c86423 :: prose twin: two significant figures, direction in words
\newcommand{\MlstFtQgRthirtyMidOverFullEFive}{0.969\,\ensuremath{\pm}\,0.030}  % experiments/FIGS/data/v2/index.json :: fine-tuning qg N=100000: R30 of R42_Q1 over L188, run by run :: sha256 c6c4138ee5c86423
\newcommand{\MlstFtQgRthirtyMidOverFullEFivePr}{equal within 4\%}  % experiments/FIGS/data/v2/index.json :: fine-tuning qg N=100000: R30 of R42_Q1 over L188, run by run :: sha256 c6c4138ee5c86423 :: prose twin: two significant figures, direction in words
\newcommand{\MlstFtQgRthirtyCoarseOverFullEFive}{0.972\,\ensuremath{\pm}\,0.020}  % experiments/FIGS/data/v2/index.json :: fine-tuning qg N=100000: R30 of R16_Q1 over L188, run by run :: sha256 c6c4138ee5c86423
\newcommand{\MlstFtQgRthirtyCoarseOverFullEFivePr}{equal within 3\%}  % experiments/FIGS/data/v2/index.json :: fine-tuning qg N=100000: R30 of R16_Q1 over L188, run by run :: sha256 c6c4138ee5c86423 :: prose twin: two significant figures, direction in words
\newcommand{\MlstFtQgRthirtyResMassOverResEFive}{0.942 and 1.027}  % experiments/FIGS/data/v2/index.json :: fine-tuning qg N=100000: R30 of L162_MASS over L162, run by run :: sha256 c6c4138ee5c86423
\newcommand{\MlstFtQgRthirtyResMassOverResEFivePr}{equal within 6\%}  % experiments/FIGS/data/v2/index.json :: fine-tuning qg N=100000: R30 of L162_MASS over L162, run by run :: sha256 c6c4138ee5c86423 :: prose twin: two significant figures, direction in words
\newcommand{\MlstFtQgRthirtyResMassOverResEFivePrMax}{equal within 6\%}  % experiments/FIGS/data/v2/index.json :: fine-tuning qg N=100000: R30 of L162_MASS over L162, run by run :: sha256 c6c4138ee5c86423 :: prose twin of the run farther from 1
\newcommand{\MlstFtQgRthirtyCoarseMassOverCoarseEFive}{1.018 and 1.004}  % experiments/FIGS/data/v2/index.json :: fine-tuning qg N=100000: R30 of R16_Q1_MASS over R16_Q1, run by run :: sha256 c6c4138ee5c86423
\newcommand{\MlstFtQgRthirtyCoarseMassOverCoarseEFivePr}{equal within 2\%}  % experiments/FIGS/data/v2/index.json :: fine-tuning qg N=100000: R30 of R16_Q1_MASS over R16_Q1, run by run :: sha256 c6c4138ee5c86423 :: prose twin: two significant figures, direction in words
\newcommand{\MlstFtQgRthirtyCoarseMassOverCoarseEFivePrMax}{equal within 2\%}  % experiments/FIGS/data/v2/index.json :: fine-tuning qg N=100000: R30 of R16_Q1_MASS over R16_Q1, run by run :: sha256 c6c4138ee5c86423 :: prose twin of the run farther from 1
\newcommand{\MlstFtQgRthirtyRandOverFullEFive}{0.986\,\ensuremath{\pm}\,0.025}  % experiments/FIGS/data/v2/index.json :: fine-tuning qg N=100000: R30 of each random partition over the L188 run of its seed :: sha256 c6c4138ee5c86423
\newcommand{\MlstFtQgRthirtyRandOverFullEFivePr}{equal within 3\%}  % experiments/FIGS/data/v2/index.json :: fine-tuning qg N=100000: R30 of each random partition over the L188 run of its seed :: sha256 c6c4138ee5c86423 :: prose twin: two significant figures, direction in words
\newcommand{\MlstFtQgRthirtyRandOverCoarseEFive}{1.009\,\ensuremath{\pm}\,0.020}  % experiments/FIGS/data/v2/index.json :: fine-tuning qg N=100000: R30 of each random partition over the R16_Q1 run of its seed :: sha256 c6c4138ee5c86423
\newcommand{\MlstFtQgRthirtyRandOverCoarseEFivePr}{equal within 3\%}  % experiments/FIGS/data/v2/index.json :: fine-tuning qg N=100000: R30 of each random partition over the R16_Q1 run of its seed :: sha256 c6c4138ee5c86423 :: prose twin: two significant figures, direction in words
\newcommand{\MlstFtQgRthirtyScratch}{104.0\,\ensuremath{\pm}\,1.0}  % experiments/FIGS/data/v2/index.json :: fine-tuning qg N=1600000: R30 from scratch, every seed :: sha256 c6c4138ee5c86423
\newcommand{\MlstFtQgRthirtyFull}{107.3\,\ensuremath{\pm}\,2.2}  % experiments/FIGS/data/v2/index.json :: fine-tuning qg N=1600000: R30, L188 :: sha256 c6c4138ee5c86423
\newcommand{\MlstFtQgRthirtyFullOverScratch}{1.031\,\ensuremath{\pm}\,0.022}  % experiments/FIGS/data/v2/index.json :: fine-tuning qg N=1600000: R30 of each L188 run over the from-scratch mean, SD over runs and standard error of that mean in quadrature :: sha256 c6c4138ee5c86423
\newcommand{\MlstFtQgRthirtyFullOverScratchPr}{equal within 4\%}  % experiments/FIGS/data/v2/index.json :: fine-tuning qg N=1600000: R30 of each L188 run over the from-scratch mean, SD over runs and standard error of that mean in quadrature :: sha256 c6c4138ee5c86423 :: prose twin: two significant figures, direction in words
\newcommand{\MlstFtQgRthirtyRes}{108.0\,\ensuremath{\pm}\,1.4}  % experiments/FIGS/data/v2/index.json :: fine-tuning qg N=1600000: R30, L162 :: sha256 c6c4138ee5c86423
\newcommand{\MlstFtQgRthirtyResOverScratch}{1.038\,\ensuremath{\pm}\,0.015}  % experiments/FIGS/data/v2/index.json :: fine-tuning qg N=1600000: R30 of each L162 run over the from-scratch mean, SD over runs and standard error of that mean in quadrature :: sha256 c6c4138ee5c86423
\newcommand{\MlstFtQgRthirtyResOverScratchPr}{3.8\% more rejection}  % experiments/FIGS/data/v2/index.json :: fine-tuning qg N=1600000: R30 of each L162 run over the from-scratch mean, SD over runs and standard error of that mean in quadrature :: sha256 c6c4138ee5c86423 :: prose twin: two significant figures, direction in words
\newcommand{\MlstFtQgRthirtyMid}{108.3\,\ensuremath{\pm}\,1.0}  % experiments/FIGS/data/v2/index.json :: fine-tuning qg N=1600000: R30, R42_Q1 :: sha256 c6c4138ee5c86423
\newcommand{\MlstFtQgRthirtyMidOverScratch}{1.041\,\ensuremath{\pm}\,0.012}  % experiments/FIGS/data/v2/index.json :: fine-tuning qg N=1600000: R30 of each R42_Q1 run over the from-scratch mean, SD over runs and standard error of that mean in quadrature :: sha256 c6c4138ee5c86423
\newcommand{\MlstFtQgRthirtyMidOverScratchPr}{4.1\% more rejection}  % experiments/FIGS/data/v2/index.json :: fine-tuning qg N=1600000: R30 of each R42_Q1 run over the from-scratch mean, SD over runs and standard error of that mean in quadrature :: sha256 c6c4138ee5c86423 :: prose twin: two significant figures, direction in words
\newcommand{\MlstFtQgRthirtyCoarse}{106.5\,\ensuremath{\pm}\,0.8}  % experiments/FIGS/data/v2/index.json :: fine-tuning qg N=1600000: R30, R16_Q1 :: sha256 c6c4138ee5c86423
\newcommand{\MlstFtQgRthirtyCoarseOverScratch}{1.024\,\ensuremath{\pm}\,0.010}  % experiments/FIGS/data/v2/index.json :: fine-tuning qg N=1600000: R30 of each R16_Q1 run over the from-scratch mean, SD over runs and standard error of that mean in quadrature :: sha256 c6c4138ee5c86423
\newcommand{\MlstFtQgRthirtyCoarseOverScratchPr}{2.4\% more rejection}  % experiments/FIGS/data/v2/index.json :: fine-tuning qg N=1600000: R30 of each R16_Q1 run over the from-scratch mean, SD over runs and standard error of that mean in quadrature :: sha256 c6c4138ee5c86423 :: prose twin: two significant figures, direction in words
\newcommand{\MlstFtQgRthirtyResMass}{107.0 and 106.6}  % experiments/FIGS/data/v2/index.json :: fine-tuning qg N=1600000: R30, L162_MASS :: sha256 c6c4138ee5c86423
\newcommand{\MlstFtQgRthirtyResMassOverScratch}{1.028 and 1.025}  % experiments/FIGS/data/v2/index.json :: fine-tuning qg N=1600000: R30 of each L162_MASS run over the from-scratch mean, SD over runs and standard error of that mean in quadrature :: sha256 c6c4138ee5c86423
\newcommand{\MlstFtQgRthirtyResMassOverScratchPr}{at least 2.5\% more rejection}  % experiments/FIGS/data/v2/index.json :: fine-tuning qg N=1600000: R30 of each L162_MASS run over the from-scratch mean, SD over runs and standard error of that mean in quadrature :: sha256 c6c4138ee5c86423 :: prose twin: two significant figures, direction in words
\newcommand{\MlstFtQgRthirtyResMassOverScratchPrMax}{up to 2.8\% more rejection}  % experiments/FIGS/data/v2/index.json :: fine-tuning qg N=1600000: R30 of each L162_MASS run over the from-scratch mean, SD over runs and standard error of that mean in quadrature :: sha256 c6c4138ee5c86423 :: prose twin of the run farther from 1
\newcommand{\MlstFtQgRthirtyCoarseMass}{104.9 and 105.6}  % experiments/FIGS/data/v2/index.json :: fine-tuning qg N=1600000: R30, R16_Q1_MASS :: sha256 c6c4138ee5c86423
\newcommand{\MlstFtQgRthirtyCoarseMassOverScratch}{1.009 and 1.015}  % experiments/FIGS/data/v2/index.json :: fine-tuning qg N=1600000: R30 of each R16_Q1_MASS run over the from-scratch mean, SD over runs and standard error of that mean in quadrature :: sha256 c6c4138ee5c86423
\newcommand{\MlstFtQgRthirtyCoarseMassOverScratchPr}{at least 0.87\% more rejection}  % experiments/FIGS/data/v2/index.json :: fine-tuning qg N=1600000: R30 of each R16_Q1_MASS run over the from-scratch mean, SD over runs and standard error of that mean in quadrature :: sha256 c6c4138ee5c86423 :: prose twin: two significant figures, direction in words
\newcommand{\MlstFtQgRthirtyCoarseMassOverScratchPrMax}{up to 1.5\% more rejection}  % experiments/FIGS/data/v2/index.json :: fine-tuning qg N=1600000: R30 of each R16_Q1_MASS run over the from-scratch mean, SD over runs and standard error of that mean in quadrature :: sha256 c6c4138ee5c86423 :: prose twin of the run farther from 1
\newcommand{\MlstFtQgRthirtyResOverFull}{1.007\,\ensuremath{\pm}\,0.026}  % experiments/FIGS/data/v2/index.json :: fine-tuning qg N=1600000: R30 of L162 over L188, run by run :: sha256 c6c4138ee5c86423
\newcommand{\MlstFtQgRthirtyResOverFullPr}{equal within 3\%}  % experiments/FIGS/data/v2/index.json :: fine-tuning qg N=1600000: R30 of L162 over L188, run by run :: sha256 c6c4138ee5c86423 :: prose twin: two significant figures, direction in words
\newcommand{\MlstFtQgRthirtyMidOverFull}{1.010\,\ensuremath{\pm}\,0.029}  % experiments/FIGS/data/v2/index.json :: fine-tuning qg N=1600000: R30 of R42_Q1 over L188, run by run :: sha256 c6c4138ee5c86423
\newcommand{\MlstFtQgRthirtyMidOverFullPr}{equal within 3\%}  % experiments/FIGS/data/v2/index.json :: fine-tuning qg N=1600000: R30 of R42_Q1 over L188, run by run :: sha256 c6c4138ee5c86423 :: prose twin: two significant figures, direction in words
\newcommand{\MlstFtQgRthirtyCoarseOverFull}{0.993\,\ensuremath{\pm}\,0.026}  % experiments/FIGS/data/v2/index.json :: fine-tuning qg N=1600000: R30 of R16_Q1 over L188, run by run :: sha256 c6c4138ee5c86423
\newcommand{\MlstFtQgRthirtyCoarseOverFullPr}{equal within 3\%}  % experiments/FIGS/data/v2/index.json :: fine-tuning qg N=1600000: R30 of R16_Q1 over L188, run by run :: sha256 c6c4138ee5c86423 :: prose twin: two significant figures, direction in words
\newcommand{\MlstFtQgRthirtyResMassOverRes}{0.976 and 0.999}  % experiments/FIGS/data/v2/index.json :: fine-tuning qg N=1600000: R30 of L162_MASS over L162, run by run :: sha256 c6c4138ee5c86423
\newcommand{\MlstFtQgRthirtyResMassOverResPr}{equal within 3\%}  % experiments/FIGS/data/v2/index.json :: fine-tuning qg N=1600000: R30 of L162_MASS over L162, run by run :: sha256 c6c4138ee5c86423 :: prose twin: two significant figures, direction in words
\newcommand{\MlstFtQgRthirtyResMassOverResPrMax}{equal within 3\%}  % experiments/FIGS/data/v2/index.json :: fine-tuning qg N=1600000: R30 of L162_MASS over L162, run by run :: sha256 c6c4138ee5c86423 :: prose twin of the run farther from 1
\newcommand{\MlstFtQgRthirtyCoarseMassOverCoarse}{0.9916 and 0.9926}  % experiments/FIGS/data/v2/index.json :: fine-tuning qg N=1600000: R30 of R16_Q1_MASS over R16_Q1, run by run :: sha256 c6c4138ee5c86423
\newcommand{\MlstFtQgRthirtyCoarseMassOverCoarsePr}{equal within 1\%}  % experiments/FIGS/data/v2/index.json :: fine-tuning qg N=1600000: R30 of R16_Q1_MASS over R16_Q1, run by run :: sha256 c6c4138ee5c86423 :: prose twin: two significant figures, direction in words
\newcommand{\MlstFtQgRthirtyCoarseMassOverCoarsePrMax}{equal within 1\%}  % experiments/FIGS/data/v2/index.json :: fine-tuning qg N=1600000: R30 of R16_Q1_MASS over R16_Q1, run by run :: sha256 c6c4138ee5c86423 :: prose twin of the run farther from 1
\newcommand{\MlstFtQgRthirtyRandOverFull}{0.985\,\ensuremath{\pm}\,0.011}  % experiments/FIGS/data/v2/index.json :: fine-tuning qg N=1600000: R30 of each random partition over the L188 run of its seed :: sha256 c6c4138ee5c86423
\newcommand{\MlstFtQgRthirtyRandOverFullPr}{equal within 2\%}  % experiments/FIGS/data/v2/index.json :: fine-tuning qg N=1600000: R30 of each random partition over the L188 run of its seed :: sha256 c6c4138ee5c86423 :: prose twin: two significant figures, direction in words
\newcommand{\MlstFtQgRthirtyRandOverCoarse}{1.008\,\ensuremath{\pm}\,0.009}  % experiments/FIGS/data/v2/index.json :: fine-tuning qg N=1600000: R30 of each random partition over the R16_Q1 run of its seed :: sha256 c6c4138ee5c86423
\newcommand{\MlstFtQgRthirtyRandOverCoarsePr}{equal within 1\%}  % experiments/FIGS/data/v2/index.json :: fine-tuning qg N=1600000: R30 of each random partition over the R16_Q1 run of its seed :: sha256 c6c4138ee5c86423 :: prose twin: two significant figures, direction in words
\newcommand{\MlstFtQgRthirtyFullOverPart}{0.994\,\ensuremath{\pm}\,0.020}  % experiments/FIGS/data/v2/index.json :: fine-tuning qg N=1600000: R30 of each 188-class run over ParT-f.t.'s published value :: sha256 c6c4138ee5c86423
\newcommand{\MlstFtQgRthirtyFullOverPartPr}{equal within 3\%}  % experiments/FIGS/data/v2/index.json :: fine-tuning qg N=1600000: R30 of each 188-class run over ParT-f.t.'s published value :: sha256 c6c4138ee5c86423 :: prose twin: two significant figures, direction in words
\newcommand{\MlstFtQgAucFull}{0.9148\,\ensuremath{\pm}\,0.0001}  % experiments/FIGS/data/v2/index.json :: fine-tuning qg N=1600000: AUC of the 188-class runs :: sha256 c6c4138ee5c86423
\newcommand{\MlstFtQgAucScratch}{0.9121\,\ensuremath{\pm}\,0.0002}  % experiments/FIGS/data/v2/index.json :: fine-tuning qg N=1600000: AUC from scratch :: sha256 c6c4138ee5c86423
\newcommand{\MlstFtQgDataEquivJets}{\ensuremath{<}1{,}000}  % experiments/FIGS/data/v2/index.json :: fine-tuning qg: from-scratch training jets whose R30 equals the 188-class mean at 1000 jets (interpolated linearly in R30 against log10 N) :: sha256 c6c4138ee5c86423
\newcommand{\MlstFtQgDataEquivJetsEFour}{26{,}000}  % experiments/FIGS/data/v2/index.json :: fine-tuning qg: from-scratch training jets whose R30 equals the 188-class mean at 10000 jets (interpolated linearly in R30 against log10 N) :: sha256 c6c4138ee5c86423
\newcommand{\MlstFtQgHerwigRthirtyScratchEThree}{27.0\,\ensuremath{\pm}\,1.3}  % experiments/FIGS/data/v2/index.json :: fine-tuning qg_herwig N=1000: R30 from scratch, every seed :: sha256 c6c4138ee5c86423
\newcommand{\MlstFtQgHerwigRthirtyFullEThree}{21\,\ensuremath{\pm}\,6}  % experiments/FIGS/data/v2/index.json :: fine-tuning qg_herwig N=1000: R30, L188 :: sha256 c6c4138ee5c86423
\newcommand{\MlstFtQgHerwigRthirtyFullOverScratchEThree}{0.77\,\ensuremath{\pm}\,0.22}  % experiments/FIGS/data/v2/index.json :: fine-tuning qg_herwig N=1000: R30 of each L188 run over the from-scratch mean, SD over runs and standard error of that mean in quadrature :: sha256 c6c4138ee5c86423
\newcommand{\MlstFtQgHerwigRthirtyFullOverScratchEThreePr}{equal within 30\%}  % experiments/FIGS/data/v2/index.json :: fine-tuning qg_herwig N=1000: R30 of each L188 run over the from-scratch mean, SD over runs and standard error of that mean in quadrature :: sha256 c6c4138ee5c86423 :: prose twin: two significant figures, direction in words
\newcommand{\MlstFtQgHerwigRthirtyResEThree}{26.2\,\ensuremath{\pm}\,0.9}  % experiments/FIGS/data/v2/index.json :: fine-tuning qg_herwig N=1000: R30, L162 :: sha256 c6c4138ee5c86423
\newcommand{\MlstFtQgHerwigRthirtyResOverScratchEThree}{0.97\,\ensuremath{\pm}\,0.04}  % experiments/FIGS/data/v2/index.json :: fine-tuning qg_herwig N=1000: R30 of each L162 run over the from-scratch mean, SD over runs and standard error of that mean in quadrature :: sha256 c6c4138ee5c86423
\newcommand{\MlstFtQgHerwigRthirtyResOverScratchEThreePr}{equal within 4\%}  % experiments/FIGS/data/v2/index.json :: fine-tuning qg_herwig N=1000: R30 of each L162 run over the from-scratch mean, SD over runs and standard error of that mean in quadrature :: sha256 c6c4138ee5c86423 :: prose twin: two significant figures, direction in words
\newcommand{\MlstFtQgHerwigRthirtyMidEThree}{26.9\,\ensuremath{\pm}\,2.7}  % experiments/FIGS/data/v2/index.json :: fine-tuning qg_herwig N=1000: R30, R42_Q1 :: sha256 c6c4138ee5c86423
\newcommand{\MlstFtQgHerwigRthirtyMidOverScratchEThree}{1.00\,\ensuremath{\pm}\,0.10}  % experiments/FIGS/data/v2/index.json :: fine-tuning qg_herwig N=1000: R30 of each R42_Q1 run over the from-scratch mean, SD over runs and standard error of that mean in quadrature :: sha256 c6c4138ee5c86423
\newcommand{\MlstFtQgHerwigRthirtyMidOverScratchEThreePr}{equal within 10\%}  % experiments/FIGS/data/v2/index.json :: fine-tuning qg_herwig N=1000: R30 of each R42_Q1 run over the from-scratch mean, SD over runs and standard error of that mean in quadrature :: sha256 c6c4138ee5c86423 :: prose twin: two significant figures, direction in words
\newcommand{\MlstFtQgHerwigRthirtyCoarseEThree}{25.2\,\ensuremath{\pm}\,2.5}  % experiments/FIGS/data/v2/index.json :: fine-tuning qg_herwig N=1000: R30, R16_Q1 :: sha256 c6c4138ee5c86423
\newcommand{\MlstFtQgHerwigRthirtyCoarseOverScratchEThree}{0.94\,\ensuremath{\pm}\,0.10}  % experiments/FIGS/data/v2/index.json :: fine-tuning qg_herwig N=1000: R30 of each R16_Q1 run over the from-scratch mean, SD over runs and standard error of that mean in quadrature :: sha256 c6c4138ee5c86423
\newcommand{\MlstFtQgHerwigRthirtyCoarseOverScratchEThreePr}{equal within 10\%}  % experiments/FIGS/data/v2/index.json :: fine-tuning qg_herwig N=1000: R30 of each R16_Q1 run over the from-scratch mean, SD over runs and standard error of that mean in quadrature :: sha256 c6c4138ee5c86423 :: prose twin: two significant figures, direction in words
\newcommand{\MlstFtQgHerwigRthirtyResMassEThree}{28.0 and 29.0}  % experiments/FIGS/data/v2/index.json :: fine-tuning qg_herwig N=1000: R30, L162_MASS :: sha256 c6c4138ee5c86423
\newcommand{\MlstFtQgHerwigRthirtyResMassOverScratchEThree}{1.038 and 1.076}  % experiments/FIGS/data/v2/index.json :: fine-tuning qg_herwig N=1000: R30 of each L162_MASS run over the from-scratch mean, SD over runs and standard error of that mean in quadrature :: sha256 c6c4138ee5c86423
\newcommand{\MlstFtQgHerwigRthirtyResMassOverScratchEThreePr}{at least 3.8\% more rejection}  % experiments/FIGS/data/v2/index.json :: fine-tuning qg_herwig N=1000: R30 of each L162_MASS run over the from-scratch mean, SD over runs and standard error of that mean in quadrature :: sha256 c6c4138ee5c86423 :: prose twin: two significant figures, direction in words
\newcommand{\MlstFtQgHerwigRthirtyResMassOverScratchEThreePrMax}{up to 7.6\% more rejection}  % experiments/FIGS/data/v2/index.json :: fine-tuning qg_herwig N=1000: R30 of each L162_MASS run over the from-scratch mean, SD over runs and standard error of that mean in quadrature :: sha256 c6c4138ee5c86423 :: prose twin of the run farther from 1
\newcommand{\MlstFtQgHerwigRthirtyCoarseMassEThree}{25.8 and 27.4}  % experiments/FIGS/data/v2/index.json :: fine-tuning qg_herwig N=1000: R30, R16_Q1_MASS :: sha256 c6c4138ee5c86423
\newcommand{\MlstFtQgHerwigRthirtyCoarseMassOverScratchEThree}{0.955 and 1.015}  % experiments/FIGS/data/v2/index.json :: fine-tuning qg_herwig N=1000: R30 of each R16_Q1_MASS run over the from-scratch mean, SD over runs and standard error of that mean in quadrature :: sha256 c6c4138ee5c86423
\newcommand{\MlstFtQgHerwigRthirtyCoarseMassOverScratchEThreePr}{equal within 5\%}  % experiments/FIGS/data/v2/index.json :: fine-tuning qg_herwig N=1000: R30 of each R16_Q1_MASS run over the from-scratch mean, SD over runs and standard error of that mean in quadrature :: sha256 c6c4138ee5c86423 :: prose twin: two significant figures, direction in words
\newcommand{\MlstFtQgHerwigRthirtyCoarseMassOverScratchEThreePrMax}{equal within 5\%}  % experiments/FIGS/data/v2/index.json :: fine-tuning qg_herwig N=1000: R30 of each R16_Q1_MASS run over the from-scratch mean, SD over runs and standard error of that mean in quadrature :: sha256 c6c4138ee5c86423 :: prose twin of the run farther from 1
\newcommand{\MlstFtQgHerwigRthirtyResOverFullEThree}{1.33\,\ensuremath{\pm}\,0.33}  % experiments/FIGS/data/v2/index.json :: fine-tuning qg_herwig N=1000: R30 of L162 over L188, run by run :: sha256 c6c4138ee5c86423
\newcommand{\MlstFtQgHerwigRthirtyResOverFullEThreePr}{equal within 40\%}  % experiments/FIGS/data/v2/index.json :: fine-tuning qg_herwig N=1000: R30 of L162 over L188, run by run :: sha256 c6c4138ee5c86423 :: prose twin: two significant figures, direction in words
\newcommand{\MlstFtQgHerwigRthirtyMidOverFullEThree}{1.36\,\ensuremath{\pm}\,0.31}  % experiments/FIGS/data/v2/index.json :: fine-tuning qg_herwig N=1000: R30 of R42_Q1 over L188, run by run :: sha256 c6c4138ee5c86423
\newcommand{\MlstFtQgHerwigRthirtyMidOverFullEThreePr}{equal within 40\%}  % experiments/FIGS/data/v2/index.json :: fine-tuning qg_herwig N=1000: R30 of R42_Q1 over L188, run by run :: sha256 c6c4138ee5c86423 :: prose twin: two significant figures, direction in words
\newcommand{\MlstFtQgHerwigRthirtyCoarseOverFullEThree}{1.3\,\ensuremath{\pm}\,0.4}  % experiments/FIGS/data/v2/index.json :: fine-tuning qg_herwig N=1000: R30 of R16_Q1 over L188, run by run :: sha256 c6c4138ee5c86423
\newcommand{\MlstFtQgHerwigRthirtyCoarseOverFullEThreePr}{equal within 50\%}  % experiments/FIGS/data/v2/index.json :: fine-tuning qg_herwig N=1000: R30 of R16_Q1 over L188, run by run :: sha256 c6c4138ee5c86423 :: prose twin: two significant figures, direction in words
\newcommand{\MlstFtQgHerwigRthirtyResMassOverResEThree}{1.062 and 1.073}  % experiments/FIGS/data/v2/index.json :: fine-tuning qg_herwig N=1000: R30 of L162_MASS over L162, run by run :: sha256 c6c4138ee5c86423
\newcommand{\MlstFtQgHerwigRthirtyResMassOverResEThreePr}{equal within 8\%}  % experiments/FIGS/data/v2/index.json :: fine-tuning qg_herwig N=1000: R30 of L162_MASS over L162, run by run :: sha256 c6c4138ee5c86423 :: prose twin: two significant figures, direction in words
\newcommand{\MlstFtQgHerwigRthirtyResMassOverResEThreePrMax}{equal within 8\%}  % experiments/FIGS/data/v2/index.json :: fine-tuning qg_herwig N=1000: R30 of L162_MASS over L162, run by run :: sha256 c6c4138ee5c86423 :: prose twin of the run farther from 1
\newcommand{\MlstFtQgHerwigRthirtyCoarseMassOverCoarseEThree}{0.93 and 1.20}  % experiments/FIGS/data/v2/index.json :: fine-tuning qg_herwig N=1000: R30 of R16_Q1_MASS over R16_Q1, run by run :: sha256 c6c4138ee5c86423
\newcommand{\MlstFtQgHerwigRthirtyCoarseMassOverCoarseEThreePr}{equal within 30\%}  % experiments/FIGS/data/v2/index.json :: fine-tuning qg_herwig N=1000: R30 of R16_Q1_MASS over R16_Q1, run by run :: sha256 c6c4138ee5c86423 :: prose twin: two significant figures, direction in words
\newcommand{\MlstFtQgHerwigRthirtyCoarseMassOverCoarseEThreePrMax}{equal within 30\%}  % experiments/FIGS/data/v2/index.json :: fine-tuning qg_herwig N=1000: R30 of R16_Q1_MASS over R16_Q1, run by run :: sha256 c6c4138ee5c86423 :: prose twin of the run farther from 1
\newcommand{\MlstFtQgHerwigRthirtyRandOverFullEThree}{0.99\,\ensuremath{\pm}\,0.20}  % experiments/FIGS/data/v2/index.json :: fine-tuning qg_herwig N=1000: R30 of each random partition over the L188 run of its seed :: sha256 c6c4138ee5c86423
\newcommand{\MlstFtQgHerwigRthirtyRandOverFullEThreePr}{equal within 30\%}  % experiments/FIGS/data/v2/index.json :: fine-tuning qg_herwig N=1000: R30 of each random partition over the L188 run of its seed :: sha256 c6c4138ee5c86423 :: prose twin: two significant figures, direction in words
\newcommand{\MlstFtQgHerwigRthirtyRandOverCoarseEThree}{0.95\,\ensuremath{\pm}\,0.13}  % experiments/FIGS/data/v2/index.json :: fine-tuning qg_herwig N=1000: R30 of each random partition over the R16_Q1 run of its seed :: sha256 c6c4138ee5c86423
\newcommand{\MlstFtQgHerwigRthirtyRandOverCoarseEThreePr}{equal within 20\%}  % experiments/FIGS/data/v2/index.json :: fine-tuning qg_herwig N=1000: R30 of each random partition over the R16_Q1 run of its seed :: sha256 c6c4138ee5c86423 :: prose twin: two significant figures, direction in words
\newcommand{\MlstFtQgHerwigRthirtyScratchEFour}{31\,\ensuremath{\pm}\,7}  % experiments/FIGS/data/v2/index.json :: fine-tuning qg_herwig N=10000: R30 from scratch, every seed :: sha256 c6c4138ee5c86423
\newcommand{\MlstFtQgHerwigRthirtyFullEFour}{27.9\,\ensuremath{\pm}\,2.8}  % experiments/FIGS/data/v2/index.json :: fine-tuning qg_herwig N=10000: R30, L188 :: sha256 c6c4138ee5c86423
\newcommand{\MlstFtQgHerwigRthirtyFullOverScratchEFour}{0.89\,\ensuremath{\pm}\,0.14}  % experiments/FIGS/data/v2/index.json :: fine-tuning qg_herwig N=10000: R30 of each L188 run over the from-scratch mean, SD over runs and standard error of that mean in quadrature :: sha256 c6c4138ee5c86423
\newcommand{\MlstFtQgHerwigRthirtyFullOverScratchEFourPr}{equal within 20\%}  % experiments/FIGS/data/v2/index.json :: fine-tuning qg_herwig N=10000: R30 of each L188 run over the from-scratch mean, SD over runs and standard error of that mean in quadrature :: sha256 c6c4138ee5c86423 :: prose twin: two significant figures, direction in words
\newcommand{\MlstFtQgHerwigRthirtyResEFour}{27.4\,\ensuremath{\pm}\,0.5}  % experiments/FIGS/data/v2/index.json :: fine-tuning qg_herwig N=10000: R30, L162 :: sha256 c6c4138ee5c86423
\newcommand{\MlstFtQgHerwigRthirtyResOverScratchEFour}{0.87\,\ensuremath{\pm}\,0.11}  % experiments/FIGS/data/v2/index.json :: fine-tuning qg_herwig N=10000: R30 of each L162 run over the from-scratch mean, SD over runs and standard error of that mean in quadrature :: sha256 c6c4138ee5c86423
\newcommand{\MlstFtQgHerwigRthirtyResOverScratchEFourPr}{13\% less rejection}  % experiments/FIGS/data/v2/index.json :: fine-tuning qg_herwig N=10000: R30 of each L162 run over the from-scratch mean, SD over runs and standard error of that mean in quadrature :: sha256 c6c4138ee5c86423 :: prose twin: two significant figures, direction in words
\newcommand{\MlstFtQgHerwigRthirtyMidEFour}{28.4\,\ensuremath{\pm}\,2.1}  % experiments/FIGS/data/v2/index.json :: fine-tuning qg_herwig N=10000: R30, R42_Q1 :: sha256 c6c4138ee5c86423
\newcommand{\MlstFtQgHerwigRthirtyMidOverScratchEFour}{0.91\,\ensuremath{\pm}\,0.13}  % experiments/FIGS/data/v2/index.json :: fine-tuning qg_herwig N=10000: R30 of each R42_Q1 run over the from-scratch mean, SD over runs and standard error of that mean in quadrature :: sha256 c6c4138ee5c86423
\newcommand{\MlstFtQgHerwigRthirtyMidOverScratchEFourPr}{equal within 10\%}  % experiments/FIGS/data/v2/index.json :: fine-tuning qg_herwig N=10000: R30 of each R42_Q1 run over the from-scratch mean, SD over runs and standard error of that mean in quadrature :: sha256 c6c4138ee5c86423 :: prose twin: two significant figures, direction in words
\newcommand{\MlstFtQgHerwigRthirtyCoarseEFour}{32.7\,\ensuremath{\pm}\,1.8}  % experiments/FIGS/data/v2/index.json :: fine-tuning qg_herwig N=10000: R30, R16_Q1 :: sha256 c6c4138ee5c86423
\newcommand{\MlstFtQgHerwigRthirtyCoarseOverScratchEFour}{1.05\,\ensuremath{\pm}\,0.14}  % experiments/FIGS/data/v2/index.json :: fine-tuning qg_herwig N=10000: R30 of each R16_Q1 run over the from-scratch mean, SD over runs and standard error of that mean in quadrature :: sha256 c6c4138ee5c86423
\newcommand{\MlstFtQgHerwigRthirtyCoarseOverScratchEFourPr}{equal within 6\%}  % experiments/FIGS/data/v2/index.json :: fine-tuning qg_herwig N=10000: R30 of each R16_Q1 run over the from-scratch mean, SD over runs and standard error of that mean in quadrature :: sha256 c6c4138ee5c86423 :: prose twin: two significant figures, direction in words
\newcommand{\MlstFtQgHerwigRthirtyResMassEFour}{34.3 and 28.5}  % experiments/FIGS/data/v2/index.json :: fine-tuning qg_herwig N=10000: R30, L162_MASS :: sha256 c6c4138ee5c86423
\newcommand{\MlstFtQgHerwigRthirtyResMassOverScratchEFour}{1.09 and 0.91}  % experiments/FIGS/data/v2/index.json :: fine-tuning qg_herwig N=10000: R30 of each L162_MASS run over the from-scratch mean, SD over runs and standard error of that mean in quadrature :: sha256 c6c4138ee5c86423
\newcommand{\MlstFtQgHerwigRthirtyResMassOverScratchEFourPr}{equal within 10\%}  % experiments/FIGS/data/v2/index.json :: fine-tuning qg_herwig N=10000: R30 of each L162_MASS run over the from-scratch mean, SD over runs and standard error of that mean in quadrature :: sha256 c6c4138ee5c86423 :: prose twin: two significant figures, direction in words
\newcommand{\MlstFtQgHerwigRthirtyResMassOverScratchEFourPrMax}{equal within 10\%}  % experiments/FIGS/data/v2/index.json :: fine-tuning qg_herwig N=10000: R30 of each L162_MASS run over the from-scratch mean, SD over runs and standard error of that mean in quadrature :: sha256 c6c4138ee5c86423 :: prose twin of the run farther from 1
\newcommand{\MlstFtQgHerwigRthirtyCoarseMassEFour}{34.4 and 31.2}  % experiments/FIGS/data/v2/index.json :: fine-tuning qg_herwig N=10000: R30, R16_Q1_MASS :: sha256 c6c4138ee5c86423
\newcommand{\MlstFtQgHerwigRthirtyCoarseMassOverScratchEFour}{1.10 and 1.00}  % experiments/FIGS/data/v2/index.json :: fine-tuning qg_herwig N=10000: R30 of each R16_Q1_MASS run over the from-scratch mean, SD over runs and standard error of that mean in quadrature :: sha256 c6c4138ee5c86423
\newcommand{\MlstFtQgHerwigRthirtyCoarseMassOverScratchEFourPr}{equal within 20\%}  % experiments/FIGS/data/v2/index.json :: fine-tuning qg_herwig N=10000: R30 of each R16_Q1_MASS run over the from-scratch mean, SD over runs and standard error of that mean in quadrature :: sha256 c6c4138ee5c86423 :: prose twin: two significant figures, direction in words
\newcommand{\MlstFtQgHerwigRthirtyCoarseMassOverScratchEFourPrMax}{equal within 20\%}  % experiments/FIGS/data/v2/index.json :: fine-tuning qg_herwig N=10000: R30 of each R16_Q1_MASS run over the from-scratch mean, SD over runs and standard error of that mean in quadrature :: sha256 c6c4138ee5c86423 :: prose twin of the run farther from 1
\newcommand{\MlstFtQgHerwigRthirtyResOverFullEFour}{0.99\,\ensuremath{\pm}\,0.12}  % experiments/FIGS/data/v2/index.json :: fine-tuning qg_herwig N=10000: R30 of L162 over L188, run by run :: sha256 c6c4138ee5c86423
\newcommand{\MlstFtQgHerwigRthirtyResOverFullEFourPr}{equal within 20\%}  % experiments/FIGS/data/v2/index.json :: fine-tuning qg_herwig N=10000: R30 of L162 over L188, run by run :: sha256 c6c4138ee5c86423 :: prose twin: two significant figures, direction in words
\newcommand{\MlstFtQgHerwigRthirtyMidOverFullEFour}{1.03\,\ensuremath{\pm}\,0.18}  % experiments/FIGS/data/v2/index.json :: fine-tuning qg_herwig N=10000: R30 of R42_Q1 over L188, run by run :: sha256 c6c4138ee5c86423
\newcommand{\MlstFtQgHerwigRthirtyMidOverFullEFourPr}{equal within 20\%}  % experiments/FIGS/data/v2/index.json :: fine-tuning qg_herwig N=10000: R30 of R42_Q1 over L188, run by run :: sha256 c6c4138ee5c86423 :: prose twin: two significant figures, direction in words
\newcommand{\MlstFtQgHerwigRthirtyCoarseOverFullEFour}{1.18\,\ensuremath{\pm}\,0.08}  % experiments/FIGS/data/v2/index.json :: fine-tuning qg_herwig N=10000: R30 of R16_Q1 over L188, run by run :: sha256 c6c4138ee5c86423
\newcommand{\MlstFtQgHerwigRthirtyCoarseOverFullEFourPr}{18\% more rejection}  % experiments/FIGS/data/v2/index.json :: fine-tuning qg_herwig N=10000: R30 of R16_Q1 over L188, run by run :: sha256 c6c4138ee5c86423 :: prose twin: two significant figures, direction in words
\newcommand{\MlstFtQgHerwigRthirtyResMassOverResEFour}{1.23 and 1.03}  % experiments/FIGS/data/v2/index.json :: fine-tuning qg_herwig N=10000: R30 of L162_MASS over L162, run by run :: sha256 c6c4138ee5c86423
\newcommand{\MlstFtQgHerwigRthirtyResMassOverResEFourPr}{equal within 30\%}  % experiments/FIGS/data/v2/index.json :: fine-tuning qg_herwig N=10000: R30 of L162_MASS over L162, run by run :: sha256 c6c4138ee5c86423 :: prose twin: two significant figures, direction in words
\newcommand{\MlstFtQgHerwigRthirtyResMassOverResEFourPrMax}{equal within 30\%}  % experiments/FIGS/data/v2/index.json :: fine-tuning qg_herwig N=10000: R30 of L162_MASS over L162, run by run :: sha256 c6c4138ee5c86423 :: prose twin of the run farther from 1
\newcommand{\MlstFtQgHerwigRthirtyCoarseMassOverCoarseEFour}{1.12 and 0.91}  % experiments/FIGS/data/v2/index.json :: fine-tuning qg_herwig N=10000: R30 of R16_Q1_MASS over R16_Q1, run by run :: sha256 c6c4138ee5c86423
\newcommand{\MlstFtQgHerwigRthirtyCoarseMassOverCoarseEFourPr}{equal within 20\%}  % experiments/FIGS/data/v2/index.json :: fine-tuning qg_herwig N=10000: R30 of R16_Q1_MASS over R16_Q1, run by run :: sha256 c6c4138ee5c86423 :: prose twin: two significant figures, direction in words
\newcommand{\MlstFtQgHerwigRthirtyCoarseMassOverCoarseEFourPrMax}{equal within 20\%}  % experiments/FIGS/data/v2/index.json :: fine-tuning qg_herwig N=10000: R30 of R16_Q1_MASS over R16_Q1, run by run :: sha256 c6c4138ee5c86423 :: prose twin of the run farther from 1
\newcommand{\MlstFtQgHerwigRthirtyRandOverFullEFour}{1.13\,\ensuremath{\pm}\,0.18}  % experiments/FIGS/data/v2/index.json :: fine-tuning qg_herwig N=10000: R30 of each random partition over the L188 run of its seed :: sha256 c6c4138ee5c86423
\newcommand{\MlstFtQgHerwigRthirtyRandOverFullEFourPr}{equal within 20\%}  % experiments/FIGS/data/v2/index.json :: fine-tuning qg_herwig N=10000: R30 of each random partition over the L188 run of its seed :: sha256 c6c4138ee5c86423 :: prose twin: two significant figures, direction in words
\newcommand{\MlstFtQgHerwigRthirtyRandOverCoarseEFour}{0.93\,\ensuremath{\pm}\,0.14}  % experiments/FIGS/data/v2/index.json :: fine-tuning qg_herwig N=10000: R30 of each random partition over the R16_Q1 run of its seed :: sha256 c6c4138ee5c86423
\newcommand{\MlstFtQgHerwigRthirtyRandOverCoarseEFourPr}{equal within 20\%}  % experiments/FIGS/data/v2/index.json :: fine-tuning qg_herwig N=10000: R30 of each random partition over the R16_Q1 run of its seed :: sha256 c6c4138ee5c86423 :: prose twin: two significant figures, direction in words
\newcommand{\MlstFtQgHerwigRthirtyScratchEFive}{36.7\,\ensuremath{\pm}\,0.9}  % experiments/FIGS/data/v2/index.json :: fine-tuning qg_herwig N=100000: R30 from scratch, every seed :: sha256 c6c4138ee5c86423
\newcommand{\MlstFtQgHerwigRthirtyFullEFive}{34.3\,\ensuremath{\pm}\,1.0}  % experiments/FIGS/data/v2/index.json :: fine-tuning qg_herwig N=100000: R30, L188 :: sha256 c6c4138ee5c86423
\newcommand{\MlstFtQgHerwigRthirtyFullOverScratchEFive}{0.935\,\ensuremath{\pm}\,0.031}  % experiments/FIGS/data/v2/index.json :: fine-tuning qg_herwig N=100000: R30 of each L188 run over the from-scratch mean, SD over runs and standard error of that mean in quadrature :: sha256 c6c4138ee5c86423
\newcommand{\MlstFtQgHerwigRthirtyFullOverScratchEFivePr}{6.5\% less rejection}  % experiments/FIGS/data/v2/index.json :: fine-tuning qg_herwig N=100000: R30 of each L188 run over the from-scratch mean, SD over runs and standard error of that mean in quadrature :: sha256 c6c4138ee5c86423 :: prose twin: two significant figures, direction in words
\newcommand{\MlstFtQgHerwigRthirtyResEFive}{35.5\,\ensuremath{\pm}\,0.5}  % experiments/FIGS/data/v2/index.json :: fine-tuning qg_herwig N=100000: R30, L162 :: sha256 c6c4138ee5c86423
\newcommand{\MlstFtQgHerwigRthirtyResOverScratchEFive}{0.966\,\ensuremath{\pm}\,0.019}  % experiments/FIGS/data/v2/index.json :: fine-tuning qg_herwig N=100000: R30 of each L162 run over the from-scratch mean, SD over runs and standard error of that mean in quadrature :: sha256 c6c4138ee5c86423
\newcommand{\MlstFtQgHerwigRthirtyResOverScratchEFivePr}{3.4\% less rejection}  % experiments/FIGS/data/v2/index.json :: fine-tuning qg_herwig N=100000: R30 of each L162 run over the from-scratch mean, SD over runs and standard error of that mean in quadrature :: sha256 c6c4138ee5c86423 :: prose twin: two significant figures, direction in words
\newcommand{\MlstFtQgHerwigRthirtyMidEFive}{35.5\,\ensuremath{\pm}\,2.1}  % experiments/FIGS/data/v2/index.json :: fine-tuning qg_herwig N=100000: R30, R42_Q1 :: sha256 c6c4138ee5c86423
\newcommand{\MlstFtQgHerwigRthirtyMidOverScratchEFive}{0.97\,\ensuremath{\pm}\,0.06}  % experiments/FIGS/data/v2/index.json :: fine-tuning qg_herwig N=100000: R30 of each R42_Q1 run over the from-scratch mean, SD over runs and standard error of that mean in quadrature :: sha256 c6c4138ee5c86423
\newcommand{\MlstFtQgHerwigRthirtyMidOverScratchEFivePr}{equal within 6\%}  % experiments/FIGS/data/v2/index.json :: fine-tuning qg_herwig N=100000: R30 of each R42_Q1 run over the from-scratch mean, SD over runs and standard error of that mean in quadrature :: sha256 c6c4138ee5c86423 :: prose twin: two significant figures, direction in words
\newcommand{\MlstFtQgHerwigRthirtyCoarseEFive}{34.9\,\ensuremath{\pm}\,1.1}  % experiments/FIGS/data/v2/index.json :: fine-tuning qg_herwig N=100000: R30, R16_Q1 :: sha256 c6c4138ee5c86423
\newcommand{\MlstFtQgHerwigRthirtyCoarseOverScratchEFive}{0.950\,\ensuremath{\pm}\,0.033}  % experiments/FIGS/data/v2/index.json :: fine-tuning qg_herwig N=100000: R30 of each R16_Q1 run over the from-scratch mean, SD over runs and standard error of that mean in quadrature :: sha256 c6c4138ee5c86423
\newcommand{\MlstFtQgHerwigRthirtyCoarseOverScratchEFivePr}{equal within 6\%}  % experiments/FIGS/data/v2/index.json :: fine-tuning qg_herwig N=100000: R30 of each R16_Q1 run over the from-scratch mean, SD over runs and standard error of that mean in quadrature :: sha256 c6c4138ee5c86423 :: prose twin: two significant figures, direction in words
\newcommand{\MlstFtQgHerwigRthirtyResMassEFive}{34.0 and 36.5}  % experiments/FIGS/data/v2/index.json :: fine-tuning qg_herwig N=100000: R30, L162_MASS :: sha256 c6c4138ee5c86423
\newcommand{\MlstFtQgHerwigRthirtyResMassOverScratchEFive}{0.925 and 0.994}  % experiments/FIGS/data/v2/index.json :: fine-tuning qg_herwig N=100000: R30 of each L162_MASS run over the from-scratch mean, SD over runs and standard error of that mean in quadrature :: sha256 c6c4138ee5c86423
\newcommand{\MlstFtQgHerwigRthirtyResMassOverScratchEFivePr}{at least 0.63\% less rejection}  % experiments/FIGS/data/v2/index.json :: fine-tuning qg_herwig N=100000: R30 of each L162_MASS run over the from-scratch mean, SD over runs and standard error of that mean in quadrature :: sha256 c6c4138ee5c86423 :: prose twin: two significant figures, direction in words
\newcommand{\MlstFtQgHerwigRthirtyResMassOverScratchEFivePrMax}{up to 7.5\% less rejection}  % experiments/FIGS/data/v2/index.json :: fine-tuning qg_herwig N=100000: R30 of each L162_MASS run over the from-scratch mean, SD over runs and standard error of that mean in quadrature :: sha256 c6c4138ee5c86423 :: prose twin of the run farther from 1
\newcommand{\MlstFtQgHerwigRthirtyCoarseMassEFive}{35.9 and 39.0}  % experiments/FIGS/data/v2/index.json :: fine-tuning qg_herwig N=100000: R30, R16_Q1_MASS :: sha256 c6c4138ee5c86423
\newcommand{\MlstFtQgHerwigRthirtyCoarseMassOverScratchEFive}{0.979 and 1.064}  % experiments/FIGS/data/v2/index.json :: fine-tuning qg_herwig N=100000: R30 of each R16_Q1_MASS run over the from-scratch mean, SD over runs and standard error of that mean in quadrature :: sha256 c6c4138ee5c86423
\newcommand{\MlstFtQgHerwigRthirtyCoarseMassOverScratchEFivePr}{equal within 7\%}  % experiments/FIGS/data/v2/index.json :: fine-tuning qg_herwig N=100000: R30 of each R16_Q1_MASS run over the from-scratch mean, SD over runs and standard error of that mean in quadrature :: sha256 c6c4138ee5c86423 :: prose twin: two significant figures, direction in words
\newcommand{\MlstFtQgHerwigRthirtyCoarseMassOverScratchEFivePrMax}{equal within 7\%}  % experiments/FIGS/data/v2/index.json :: fine-tuning qg_herwig N=100000: R30 of each R16_Q1_MASS run over the from-scratch mean, SD over runs and standard error of that mean in quadrature :: sha256 c6c4138ee5c86423 :: prose twin of the run farther from 1
\newcommand{\MlstFtQgHerwigRthirtyResOverFullEFive}{1.03\,\ensuremath{\pm}\,0.04}  % experiments/FIGS/data/v2/index.json :: fine-tuning qg_herwig N=100000: R30 of L162 over L188, run by run :: sha256 c6c4138ee5c86423
\newcommand{\MlstFtQgHerwigRthirtyResOverFullEFivePr}{equal within 5\%}  % experiments/FIGS/data/v2/index.json :: fine-tuning qg_herwig N=100000: R30 of L162 over L188, run by run :: sha256 c6c4138ee5c86423 :: prose twin: two significant figures, direction in words
\newcommand{\MlstFtQgHerwigRthirtyMidOverFullEFive}{1.03\,\ensuremath{\pm}\,0.05}  % experiments/FIGS/data/v2/index.json :: fine-tuning qg_herwig N=100000: R30 of R42_Q1 over L188, run by run :: sha256 c6c4138ee5c86423
\newcommand{\MlstFtQgHerwigRthirtyMidOverFullEFivePr}{equal within 6\%}  % experiments/FIGS/data/v2/index.json :: fine-tuning qg_herwig N=100000: R30 of R42_Q1 over L188, run by run :: sha256 c6c4138ee5c86423 :: prose twin: two significant figures, direction in words
\newcommand{\MlstFtQgHerwigRthirtyCoarseOverFullEFive}{1.02\,\ensuremath{\pm}\,0.06}  % experiments/FIGS/data/v2/index.json :: fine-tuning qg_herwig N=100000: R30 of R16_Q1 over L188, run by run :: sha256 c6c4138ee5c86423
\newcommand{\MlstFtQgHerwigRthirtyCoarseOverFullEFivePr}{equal within 7\%}  % experiments/FIGS/data/v2/index.json :: fine-tuning qg_herwig N=100000: R30 of R16_Q1 over L188, run by run :: sha256 c6c4138ee5c86423 :: prose twin: two significant figures, direction in words
\newcommand{\MlstFtQgHerwigRthirtyResMassOverResEFive}{0.971 and 1.028}  % experiments/FIGS/data/v2/index.json :: fine-tuning qg_herwig N=100000: R30 of L162_MASS over L162, run by run :: sha256 c6c4138ee5c86423
\newcommand{\MlstFtQgHerwigRthirtyResMassOverResEFivePr}{equal within 3\%}  % experiments/FIGS/data/v2/index.json :: fine-tuning qg_herwig N=100000: R30 of L162_MASS over L162, run by run :: sha256 c6c4138ee5c86423 :: prose twin: two significant figures, direction in words
\newcommand{\MlstFtQgHerwigRthirtyResMassOverResEFivePrMax}{equal within 3\%}  % experiments/FIGS/data/v2/index.json :: fine-tuning qg_herwig N=100000: R30 of L162_MASS over L162, run by run :: sha256 c6c4138ee5c86423 :: prose twin of the run farther from 1
\newcommand{\MlstFtQgHerwigRthirtyCoarseMassOverCoarseEFive}{1.064 and 1.121}  % experiments/FIGS/data/v2/index.json :: fine-tuning qg_herwig N=100000: R30 of R16_Q1_MASS over R16_Q1, run by run :: sha256 c6c4138ee5c86423
\newcommand{\MlstFtQgHerwigRthirtyCoarseMassOverCoarseEFivePr}{equal within 20\%}  % experiments/FIGS/data/v2/index.json :: fine-tuning qg_herwig N=100000: R30 of R16_Q1_MASS over R16_Q1, run by run :: sha256 c6c4138ee5c86423 :: prose twin: two significant figures, direction in words
\newcommand{\MlstFtQgHerwigRthirtyCoarseMassOverCoarseEFivePrMax}{equal within 20\%}  % experiments/FIGS/data/v2/index.json :: fine-tuning qg_herwig N=100000: R30 of R16_Q1_MASS over R16_Q1, run by run :: sha256 c6c4138ee5c86423 :: prose twin of the run farther from 1
\newcommand{\MlstFtQgHerwigRthirtyRandOverFullEFive}{1.01\,\ensuremath{\pm}\,0.05}  % experiments/FIGS/data/v2/index.json :: fine-tuning qg_herwig N=100000: R30 of each random partition over the L188 run of its seed :: sha256 c6c4138ee5c86423
\newcommand{\MlstFtQgHerwigRthirtyRandOverFullEFivePr}{equal within 5\%}  % experiments/FIGS/data/v2/index.json :: fine-tuning qg_herwig N=100000: R30 of each random partition over the L188 run of its seed :: sha256 c6c4138ee5c86423 :: prose twin: two significant figures, direction in words
\newcommand{\MlstFtQgHerwigRthirtyRandOverCoarseEFive}{1.009\,\ensuremath{\pm}\,0.024}  % experiments/FIGS/data/v2/index.json :: fine-tuning qg_herwig N=100000: R30 of each random partition over the R16_Q1 run of its seed :: sha256 c6c4138ee5c86423
\newcommand{\MlstFtQgHerwigRthirtyRandOverCoarseEFivePr}{equal within 3\%}  % experiments/FIGS/data/v2/index.json :: fine-tuning qg_herwig N=100000: R30 of each random partition over the R16_Q1 run of its seed :: sha256 c6c4138ee5c86423 :: prose twin: two significant figures, direction in words
\newcommand{\MlstFtQgHerwigRthirtyScratch}{35.5\,\ensuremath{\pm}\,1.2}  % experiments/FIGS/data/v2/index.json :: fine-tuning qg_herwig N=1600000: R30 from scratch, every seed :: sha256 c6c4138ee5c86423
\newcommand{\MlstFtQgHerwigRthirtyFull}{35.2\,\ensuremath{\pm}\,2.1}  % experiments/FIGS/data/v2/index.json :: fine-tuning qg_herwig N=1600000: R30, L188 :: sha256 c6c4138ee5c86423
\newcommand{\MlstFtQgHerwigRthirtyFullOverScratch}{0.99\,\ensuremath{\pm}\,0.06}  % experiments/FIGS/data/v2/index.json :: fine-tuning qg_herwig N=1600000: R30 of each L188 run over the from-scratch mean, SD over runs and standard error of that mean in quadrature :: sha256 c6c4138ee5c86423
\newcommand{\MlstFtQgHerwigRthirtyFullOverScratchPr}{equal within 6\%}  % experiments/FIGS/data/v2/index.json :: fine-tuning qg_herwig N=1600000: R30 of each L188 run over the from-scratch mean, SD over runs and standard error of that mean in quadrature :: sha256 c6c4138ee5c86423 :: prose twin: two significant figures, direction in words
\newcommand{\MlstFtQgHerwigRthirtyRes}{34.6\,\ensuremath{\pm}\,2.0}  % experiments/FIGS/data/v2/index.json :: fine-tuning qg_herwig N=1600000: R30, L162 :: sha256 c6c4138ee5c86423
\newcommand{\MlstFtQgHerwigRthirtyResOverScratch}{0.97\,\ensuremath{\pm}\,0.06}  % experiments/FIGS/data/v2/index.json :: fine-tuning qg_herwig N=1600000: R30 of each L162 run over the from-scratch mean, SD over runs and standard error of that mean in quadrature :: sha256 c6c4138ee5c86423
\newcommand{\MlstFtQgHerwigRthirtyResOverScratchPr}{equal within 6\%}  % experiments/FIGS/data/v2/index.json :: fine-tuning qg_herwig N=1600000: R30 of each L162 run over the from-scratch mean, SD over runs and standard error of that mean in quadrature :: sha256 c6c4138ee5c86423 :: prose twin: two significant figures, direction in words
\newcommand{\MlstFtQgHerwigRthirtyMid}{36.1\,\ensuremath{\pm}\,2.7}  % experiments/FIGS/data/v2/index.json :: fine-tuning qg_herwig N=1600000: R30, R42_Q1 :: sha256 c6c4138ee5c86423
\newcommand{\MlstFtQgHerwigRthirtyMidOverScratch}{1.02\,\ensuremath{\pm}\,0.08}  % experiments/FIGS/data/v2/index.json :: fine-tuning qg_herwig N=1600000: R30 of each R42_Q1 run over the from-scratch mean, SD over runs and standard error of that mean in quadrature :: sha256 c6c4138ee5c86423
\newcommand{\MlstFtQgHerwigRthirtyMidOverScratchPr}{equal within 8\%}  % experiments/FIGS/data/v2/index.json :: fine-tuning qg_herwig N=1600000: R30 of each R42_Q1 run over the from-scratch mean, SD over runs and standard error of that mean in quadrature :: sha256 c6c4138ee5c86423 :: prose twin: two significant figures, direction in words
\newcommand{\MlstFtQgHerwigRthirtyCoarse}{36.9\,\ensuremath{\pm}\,0.9}  % experiments/FIGS/data/v2/index.json :: fine-tuning qg_herwig N=1600000: R30, R16_Q1 :: sha256 c6c4138ee5c86423
\newcommand{\MlstFtQgHerwigRthirtyCoarseOverScratch}{1.040\,\ensuremath{\pm}\,0.033}  % experiments/FIGS/data/v2/index.json :: fine-tuning qg_herwig N=1600000: R30 of each R16_Q1 run over the from-scratch mean, SD over runs and standard error of that mean in quadrature :: sha256 c6c4138ee5c86423
\newcommand{\MlstFtQgHerwigRthirtyCoarseOverScratchPr}{equal within 4\%}  % experiments/FIGS/data/v2/index.json :: fine-tuning qg_herwig N=1600000: R30 of each R16_Q1 run over the from-scratch mean, SD over runs and standard error of that mean in quadrature :: sha256 c6c4138ee5c86423 :: prose twin: two significant figures, direction in words
\newcommand{\MlstFtQgHerwigRthirtyResMass}{37.8 and 34.4}  % experiments/FIGS/data/v2/index.json :: fine-tuning qg_herwig N=1600000: R30, L162_MASS :: sha256 c6c4138ee5c86423
\newcommand{\MlstFtQgHerwigRthirtyResMassOverScratch}{1.065 and 0.970}  % experiments/FIGS/data/v2/index.json :: fine-tuning qg_herwig N=1600000: R30 of each L162_MASS run over the from-scratch mean, SD over runs and standard error of that mean in quadrature :: sha256 c6c4138ee5c86423
\newcommand{\MlstFtQgHerwigRthirtyResMassOverScratchPr}{equal within 7\%}  % experiments/FIGS/data/v2/index.json :: fine-tuning qg_herwig N=1600000: R30 of each L162_MASS run over the from-scratch mean, SD over runs and standard error of that mean in quadrature :: sha256 c6c4138ee5c86423 :: prose twin: two significant figures, direction in words
\newcommand{\MlstFtQgHerwigRthirtyResMassOverScratchPrMax}{equal within 7\%}  % experiments/FIGS/data/v2/index.json :: fine-tuning qg_herwig N=1600000: R30 of each L162_MASS run over the from-scratch mean, SD over runs and standard error of that mean in quadrature :: sha256 c6c4138ee5c86423 :: prose twin of the run farther from 1
\newcommand{\MlstFtQgHerwigRthirtyCoarseMass}{37.7 and 32.9}  % experiments/FIGS/data/v2/index.json :: fine-tuning qg_herwig N=1600000: R30, R16_Q1_MASS :: sha256 c6c4138ee5c86423
\newcommand{\MlstFtQgHerwigRthirtyCoarseMassOverScratch}{1.06 and 0.93}  % experiments/FIGS/data/v2/index.json :: fine-tuning qg_herwig N=1600000: R30 of each R16_Q1_MASS run over the from-scratch mean, SD over runs and standard error of that mean in quadrature :: sha256 c6c4138ee5c86423
\newcommand{\MlstFtQgHerwigRthirtyCoarseMassOverScratchPr}{equal within 8\%}  % experiments/FIGS/data/v2/index.json :: fine-tuning qg_herwig N=1600000: R30 of each R16_Q1_MASS run over the from-scratch mean, SD over runs and standard error of that mean in quadrature :: sha256 c6c4138ee5c86423 :: prose twin: two significant figures, direction in words
\newcommand{\MlstFtQgHerwigRthirtyCoarseMassOverScratchPrMax}{equal within 8\%}  % experiments/FIGS/data/v2/index.json :: fine-tuning qg_herwig N=1600000: R30 of each R16_Q1_MASS run over the from-scratch mean, SD over runs and standard error of that mean in quadrature :: sha256 c6c4138ee5c86423 :: prose twin of the run farther from 1
\newcommand{\MlstFtQgHerwigRthirtyResOverFull}{0.9832\,\ensuremath{\pm}\,0.0022}  % experiments/FIGS/data/v2/index.json :: fine-tuning qg_herwig N=1600000: R30 of L162 over L188, run by run :: sha256 c6c4138ee5c86423
\newcommand{\MlstFtQgHerwigRthirtyResOverFullPr}{1.7\% less rejection}  % experiments/FIGS/data/v2/index.json :: fine-tuning qg_herwig N=1600000: R30 of L162 over L188, run by run :: sha256 c6c4138ee5c86423 :: prose twin: two significant figures, direction in words
\newcommand{\MlstFtQgHerwigRthirtyMidOverFull}{1.03\,\ensuremath{\pm}\,0.05}  % experiments/FIGS/data/v2/index.json :: fine-tuning qg_herwig N=1600000: R30 of R42_Q1 over L188, run by run :: sha256 c6c4138ee5c86423
\newcommand{\MlstFtQgHerwigRthirtyMidOverFullPr}{equal within 6\%}  % experiments/FIGS/data/v2/index.json :: fine-tuning qg_herwig N=1600000: R30 of R42_Q1 over L188, run by run :: sha256 c6c4138ee5c86423 :: prose twin: two significant figures, direction in words
\newcommand{\MlstFtQgHerwigRthirtyCoarseOverFull}{1.05\,\ensuremath{\pm}\,0.06}  % experiments/FIGS/data/v2/index.json :: fine-tuning qg_herwig N=1600000: R30 of R16_Q1 over L188, run by run :: sha256 c6c4138ee5c86423
\newcommand{\MlstFtQgHerwigRthirtyCoarseOverFullPr}{equal within 6\%}  % experiments/FIGS/data/v2/index.json :: fine-tuning qg_herwig N=1600000: R30 of R16_Q1 over L188, run by run :: sha256 c6c4138ee5c86423 :: prose twin: two significant figures, direction in words
\newcommand{\MlstFtQgHerwigRthirtyResMassOverRes}{1.14 and 0.93}  % experiments/FIGS/data/v2/index.json :: fine-tuning qg_herwig N=1600000: R30 of L162_MASS over L162, run by run :: sha256 c6c4138ee5c86423
\newcommand{\MlstFtQgHerwigRthirtyResMassOverResPr}{equal within 20\%}  % experiments/FIGS/data/v2/index.json :: fine-tuning qg_herwig N=1600000: R30 of L162_MASS over L162, run by run :: sha256 c6c4138ee5c86423 :: prose twin: two significant figures, direction in words
\newcommand{\MlstFtQgHerwigRthirtyResMassOverResPrMax}{equal within 20\%}  % experiments/FIGS/data/v2/index.json :: fine-tuning qg_herwig N=1600000: R30 of L162_MASS over L162, run by run :: sha256 c6c4138ee5c86423 :: prose twin of the run farther from 1
\newcommand{\MlstFtQgHerwigRthirtyCoarseMassOverCoarse}{1.01 and 0.88}  % experiments/FIGS/data/v2/index.json :: fine-tuning qg_herwig N=1600000: R30 of R16_Q1_MASS over R16_Q1, run by run :: sha256 c6c4138ee5c86423
\newcommand{\MlstFtQgHerwigRthirtyCoarseMassOverCoarsePr}{equal within 20\%}  % experiments/FIGS/data/v2/index.json :: fine-tuning qg_herwig N=1600000: R30 of R16_Q1_MASS over R16_Q1, run by run :: sha256 c6c4138ee5c86423 :: prose twin: two significant figures, direction in words
\newcommand{\MlstFtQgHerwigRthirtyCoarseMassOverCoarsePrMax}{equal within 20\%}  % experiments/FIGS/data/v2/index.json :: fine-tuning qg_herwig N=1600000: R30 of R16_Q1_MASS over R16_Q1, run by run :: sha256 c6c4138ee5c86423 :: prose twin of the run farther from 1
\newcommand{\MlstFtQgHerwigRthirtyRandOverFull}{0.98\,\ensuremath{\pm}\,0.04}  % experiments/FIGS/data/v2/index.json :: fine-tuning qg_herwig N=1600000: R30 of each random partition over the L188 run of its seed :: sha256 c6c4138ee5c86423
\newcommand{\MlstFtQgHerwigRthirtyRandOverFullPr}{equal within 5\%}  % experiments/FIGS/data/v2/index.json :: fine-tuning qg_herwig N=1600000: R30 of each random partition over the L188 run of its seed :: sha256 c6c4138ee5c86423 :: prose twin: two significant figures, direction in words
\newcommand{\MlstFtQgHerwigRthirtyRandOverCoarse}{0.94\,\ensuremath{\pm}\,0.04}  % experiments/FIGS/data/v2/index.json :: fine-tuning qg_herwig N=1600000: R30 of each random partition over the R16_Q1 run of its seed :: sha256 c6c4138ee5c86423
\newcommand{\MlstFtQgHerwigRthirtyRandOverCoarsePr}{equal within 6\%}  % experiments/FIGS/data/v2/index.json :: fine-tuning qg_herwig N=1600000: R30 of each random partition over the R16_Q1 run of its seed :: sha256 c6c4138ee5c86423 :: prose twin: two significant figures, direction in words
\newcommand{\MlstFtQgRthirtyHerwigOverPythiaFull}{0.328\,\ensuremath{\pm}\,0.015}  % experiments/FIGS/data/v2/index.json :: fine-tuning q/g N=1600000: R30 on Herwig over R30 on Pythia, each run of L188 :: sha256 c6c4138ee5c86423
\newcommand{\MlstFtQgRthirtyHerwigOverPythiaFullPr}{0.33 times the rejection}  % experiments/FIGS/data/v2/index.json :: fine-tuning q/g N=1600000: R30 on Herwig over R30 on Pythia, each run of L188 :: sha256 c6c4138ee5c86423 :: prose twin: two significant figures, direction in words
\newcommand{\MlstFtQgRthirtyHerwigOverPythiaRes}{0.320\,\ensuremath{\pm}\,0.023}  % experiments/FIGS/data/v2/index.json :: fine-tuning q/g N=1600000: R30 on Herwig over R30 on Pythia, each run of L162 :: sha256 c6c4138ee5c86423
\newcommand{\MlstFtQgRthirtyHerwigOverPythiaResPr}{0.32 times the rejection}  % experiments/FIGS/data/v2/index.json :: fine-tuning q/g N=1600000: R30 on Herwig over R30 on Pythia, each run of L162 :: sha256 c6c4138ee5c86423 :: prose twin: two significant figures, direction in words
\newcommand{\MlstFtQgRthirtyHerwigOverPythiaMid}{0.334\,\ensuremath{\pm}\,0.023}  % experiments/FIGS/data/v2/index.json :: fine-tuning q/g N=1600000: R30 on Herwig over R30 on Pythia, each run of R42_Q1 :: sha256 c6c4138ee5c86423
\newcommand{\MlstFtQgRthirtyHerwigOverPythiaMidPr}{0.33 times the rejection}  % experiments/FIGS/data/v2/index.json :: fine-tuning q/g N=1600000: R30 on Herwig over R30 on Pythia, each run of R42_Q1 :: sha256 c6c4138ee5c86423 :: prose twin: two significant figures, direction in words
\newcommand{\MlstFtQgRthirtyHerwigOverPythiaCoarse}{0.346\,\ensuremath{\pm}\,0.011}  % experiments/FIGS/data/v2/index.json :: fine-tuning q/g N=1600000: R30 on Herwig over R30 on Pythia, each run of R16_Q1 :: sha256 c6c4138ee5c86423
\newcommand{\MlstFtQgRthirtyHerwigOverPythiaCoarsePr}{0.35 times the rejection}  % experiments/FIGS/data/v2/index.json :: fine-tuning q/g N=1600000: R30 on Herwig over R30 on Pythia, each run of R16_Q1 :: sha256 c6c4138ee5c86423 :: prose twin: two significant figures, direction in words
\newcommand{\MlstFtQgRthirtyHerwigOverPythiaResMass}{0.353 and 0.323}  % experiments/FIGS/data/v2/index.json :: fine-tuning q/g N=1600000: R30 on Herwig over R30 on Pythia, each run of L162_MASS :: sha256 c6c4138ee5c86423
\newcommand{\MlstFtQgRthirtyHerwigOverPythiaResMassPr}{at most 0.35 times the rejection}  % experiments/FIGS/data/v2/index.json :: fine-tuning q/g N=1600000: R30 on Herwig over R30 on Pythia, each run of L162_MASS :: sha256 c6c4138ee5c86423 :: prose twin: two significant figures, direction in words
\newcommand{\MlstFtQgRthirtyHerwigOverPythiaResMassPrMax}{down to 0.32 times the rejection}  % experiments/FIGS/data/v2/index.json :: fine-tuning q/g N=1600000: R30 on Herwig over R30 on Pythia, each run of L162_MASS :: sha256 c6c4138ee5c86423 :: prose twin of the run farther from 1
\newcommand{\MlstFtQgRthirtyHerwigOverPythiaCoarseMass}{0.359 and 0.311}  % experiments/FIGS/data/v2/index.json :: fine-tuning q/g N=1600000: R30 on Herwig over R30 on Pythia, each run of R16_Q1_MASS :: sha256 c6c4138ee5c86423
\newcommand{\MlstFtQgRthirtyHerwigOverPythiaCoarseMassPr}{at most 0.36 times the rejection}  % experiments/FIGS/data/v2/index.json :: fine-tuning q/g N=1600000: R30 on Herwig over R30 on Pythia, each run of R16_Q1_MASS :: sha256 c6c4138ee5c86423 :: prose twin: two significant figures, direction in words
\newcommand{\MlstFtQgRthirtyHerwigOverPythiaCoarseMassPrMax}{down to 0.31 times the rejection}  % experiments/FIGS/data/v2/index.json :: fine-tuning q/g N=1600000: R30 on Herwig over R30 on Pythia, each run of R16_Q1_MASS :: sha256 c6c4138ee5c86423 :: prose twin of the run farther from 1
\newcommand{\MlstFtQgRthirtyHerwigOverPythiaScratch}{0.341\,\ensuremath{\pm}\,0.009}  % experiments/FIGS/data/v2/index.json :: fine-tuning q/g N=1600000: R30 on Herwig over R30 on Pythia, each from-scratch seed :: sha256 c6c4138ee5c86423
\newcommand{\MlstFtQgRthirtyHerwigOverPythiaScratchPr}{0.34 times the rejection}  % experiments/FIGS/data/v2/index.json :: fine-tuning q/g N=1600000: R30 on Herwig over R30 on Pythia, each from-scratch seed :: sha256 c6c4138ee5c86423 :: prose twin: two significant figures, direction in words
\newcommand{\MlstTopCheckFullLong}{\pending{recipe check}}  % experiments/FIGS/data/v2/index.json :: the recipe check is not read out yet :: sha256 c6c4138ee5c86423
\newcommand{\MlstTopCheckScratchLong}{\pending{recipe check}}  % experiments/FIGS/data/v2/index.json :: the recipe check is not read out yet :: sha256 c6c4138ee5c86423
\newcommand{\MlstTopCheckFullShort}{\pending{recipe check}}  % experiments/FIGS/data/v2/index.json :: the recipe check is not read out yet :: sha256 c6c4138ee5c86423
\newcommand{\MlstTopCheckScratchShort}{\pending{recipe check}}  % experiments/FIGS/data/v2/index.json :: the recipe check is not read out yet :: sha256 c6c4138ee5c86423
\newcommand{\MlstLpJcOneGapFracEFour}{56.3\,\ensuremath{\pm}\,1.1\%}  % experiments/FIGS/data/v2/index.json :: jc1 N=10000: (probe - scratch) / (fine-tuned - scratch) accuracy, 188-class runs, per cent :: sha256 c6c4138ee5c86423
\newcommand{\MlstLpJcOneGapFracEFourPr}{56\%}  % experiments/FIGS/data/v2/index.json :: jc1 N=10000: (probe - scratch) / (fine-tuned - scratch) accuracy, 188-class runs, per cent :: sha256 c6c4138ee5c86423 :: prose twin, two significant figures
\newcommand{\MlstLpJcOneGapFracEFive}{\ensuremath{-}15\,\ensuremath{\pm}\,5\%}  % experiments/FIGS/data/v2/index.json :: jc1 N=100000: (probe - scratch) / (fine-tuned - scratch) accuracy, 188-class runs, per cent :: sha256 c6c4138ee5c86423
\newcommand{\MlstLpJcOneGapFracEFivePr}{\ensuremath{-}15\%}  % experiments/FIGS/data/v2/index.json :: jc1 N=100000: (probe - scratch) / (fine-tuned - scratch) accuracy, 188-class runs, per cent :: sha256 c6c4138ee5c86423 :: prose twin, two significant figures
\newcommand{\MlstLpJcTwoGapFracEFour}{96.2\,\ensuremath{\pm}\,0.6\%}  % experiments/FIGS/data/v2/index.json :: jc2_188 N=10000: (probe - scratch) / (fine-tuned - scratch) accuracy, 188-class runs, per cent :: sha256 c6c4138ee5c86423
\newcommand{\MlstLpJcTwoGapFracEFourPr}{96\%}  % experiments/FIGS/data/v2/index.json :: jc2_188 N=10000: (probe - scratch) / (fine-tuned - scratch) accuracy, 188-class runs, per cent :: sha256 c6c4138ee5c86423 :: prose twin, two significant figures
\newcommand{\MlstLpJcTwoGapFracEFive}{102.2\,\ensuremath{\pm}\,1.0\%}  % experiments/FIGS/data/v2/index.json :: jc2_188 N=100000: (probe - scratch) / (fine-tuned - scratch) accuracy, 188-class runs, per cent :: sha256 c6c4138ee5c86423
\newcommand{\MlstLpJcTwoGapFracEFivePr}{100\%}  % experiments/FIGS/data/v2/index.json :: jc2_188 N=100000: (probe - scratch) / (fine-tuned - scratch) accuracy, 188-class runs, per cent :: sha256 c6c4138ee5c86423 :: prose twin, two significant figures
\newcommand{\XbxMx}{\pending{samples not in}}  % experiments/XBS/xbx_report.py :: pending :: sha256 cb6dffe774ff5b1e
\newcommand{\XbxWindowLo}{\pending{samples not in}}  % experiments/XBS/xbx_report.py :: pending :: sha256 cb6dffe774ff5b1e
\newcommand{\XbxWindowHi}{\pending{samples not in}}  % experiments/XBS/xbx_report.py :: pending :: sha256 cb6dffe774ff5b1e
\newcommand{\XbbNTrainPerClass}{\pending{samples not in}}  % experiments/XBS/xbx_report.py :: pending :: sha256 cb6dffe774ff5b1e
\newcommand{\XbbRepetition}{\pending{samples not in}}  % experiments/XBS/xbx_report.py :: pending :: sha256 cb6dffe774ff5b1e
\newcommand{\XbbNEvalSignal}{\pending{samples not in}}  % experiments/XBS/xbx_report.py :: pending :: sha256 cb6dffe774ff5b1e
\newcommand{\XbbNEvalQcd}{\pending{samples not in}}  % experiments/XBS/xbx_report.py :: pending :: sha256 cb6dffe774ff5b1e
\newcommand{\XbsNTrainPerClass}{\pending{samples not in}}  % experiments/XBS/xbx_report.py :: pending :: sha256 cb6dffe774ff5b1e
\newcommand{\XbsRepetition}{\pending{samples not in}}  % experiments/XBS/xbx_report.py :: pending :: sha256 cb6dffe774ff5b1e
\newcommand{\XbsNEvalSignal}{\pending{samples not in}}  % experiments/XBS/xbx_report.py :: pending :: sha256 cb6dffe774ff5b1e
\newcommand{\XbsNEvalQcd}{\pending{samples not in}}  % experiments/XBS/xbx_report.py :: pending :: sha256 cb6dffe774ff5b1e
\newcommand{\MlstPartreproRthirtyPart}{1{,}711}  % experiments/FIGS/data/v2/partrepro/part.json :: r30 :: sha256 0d44b997af9e6448 :: ParT reproduction on top, part: R30 (one run)
\newcommand{\MlstPartreproRthirtyHead}{1{,}577}  % experiments/FIGS/data/v2/partrepro/head.json :: r30 :: sha256 37b82ba0a0fe9299 :: ParT reproduction on top, head: R30 (one run)
\newcommand{\MlstPartreproRthirtyRescale}{1{,}224}  % experiments/FIGS/data/v2/partrepro/rescale.json :: r30 :: sha256 0a9f1401047d07ff :: ParT reproduction on top, rescale: R30 (one run)
\newcommand{\MlstPartreproRthirtyHeadOverPart}{0.922}  % experiments/FIGS/data/v2/partrepro/head.json :: r30 / part r30 :: sha256 37b82ba0a0fe9299
\newcommand{\MlstPartreproRthirtyHeadOverPartPr}{7.8\% less rejection}  % experiments/FIGS/data/v2/partrepro/head.json :: r30 / part r30 :: sha256 37b82ba0a0fe9299 :: prose twin: two significant figures, direction in words
\newcommand{\MlstPartreproRthirtyRescaleOverPart}{0.715}  % experiments/FIGS/data/v2/partrepro/rescale.json :: r30 / part r30 :: sha256 0a9f1401047d07ff
\newcommand{\MlstPartreproRthirtyRescaleOverPartPr}{28\% less rejection}  % experiments/FIGS/data/v2/partrepro/rescale.json :: r30 / part r30 :: sha256 0a9f1401047d07ff :: prose twin: two significant figures, direction in words
\newcommand{\MlstPartreproSliceDropRatio}{\pending{jet-pT slice test not run}}  % experiments/FIGS/data/v2/index.json :: open item O10: R30 in narrow pT slices of partrepro part and rescale :: sha256 c6c4138ee5c86423
\newcommand{\MlstInputBnZeroSigmaMax}{0.96}  % experiments/FIGS/data/v2/head/input_bn_offsets.json :: offsets :: sha256 d6b821509fccecff :: largest |offset| over the five type flags
\newcommand{\MlstFtJcOneNSignal}{nine}  % experiments/FIGS/data/v2/curves/full_budget_long_leg2.json :: JC1_COLS (count) :: sha256 588e935996d1006a :: JetClass signal classes
\newcommand{\MlstFtJcOneRejHbbResOverFull}{0.92\,\ensuremath{\pm}\,0.05}  % experiments/FIGS/data/v2/curves/full_budget_long_leg2.json :: cells.<run>.N1000000.s1.jc1.per_class.label_Hbb: rejection of L162 over L188, run by run :: sha256 588e935996d1006a
\newcommand{\MlstFtJcOneRejHbbResOverFullPr}{equal within 9\%}  % experiments/FIGS/data/v2/curves/full_budget_long_leg2.json :: cells.<run>.N1000000.s1.jc1.per_class.label_Hbb: rejection of L162 over L188, run by run :: sha256 588e935996d1006a :: prose twin: two significant figures, direction in words
\newcommand{\MlstFtJcOneRejHbbMidOverFull}{0.98\,\ensuremath{\pm}\,0.12}  % experiments/FIGS/data/v2/curves/full_budget_long_leg2.json :: cells.<run>.N1000000.s1.jc1.per_class.label_Hbb: rejection of R42_Q1 over L188, run by run :: sha256 588e935996d1006a
\newcommand{\MlstFtJcOneRejHbbMidOverFullPr}{equal within 20\%}  % experiments/FIGS/data/v2/curves/full_budget_long_leg2.json :: cells.<run>.N1000000.s1.jc1.per_class.label_Hbb: rejection of R42_Q1 over L188, run by run :: sha256 588e935996d1006a :: prose twin: two significant figures, direction in words
\newcommand{\MlstFtJcOneRejHbbCoarseOverFull}{1.01\,\ensuremath{\pm}\,0.07}  % experiments/FIGS/data/v2/curves/full_budget_long_leg2.json :: cells.<run>.N1000000.s1.jc1.per_class.label_Hbb: rejection of R16_Q1 over L188, run by run :: sha256 588e935996d1006a
\newcommand{\MlstFtJcOneRejHbbCoarseOverFullPr}{equal within 7\%}  % experiments/FIGS/data/v2/curves/full_budget_long_leg2.json :: cells.<run>.N1000000.s1.jc1.per_class.label_Hbb: rejection of R16_Q1 over L188, run by run :: sha256 588e935996d1006a :: prose twin: two significant figures, direction in words
\newcommand{\MlstFtJcOneRejHbbResMassOverRes}{1.065 and 1.100}  % experiments/FIGS/data/v2/curves/full_budget_long_leg2.json :: cells.<run>.N1000000.s1.jc1.per_class.label_Hbb: rejection of L162_MASS over L162, run by run :: sha256 588e935996d1006a
\newcommand{\MlstFtJcOneRejHbbResMassOverResPr}{equal within 10\%}  % experiments/FIGS/data/v2/curves/full_budget_long_leg2.json :: cells.<run>.N1000000.s1.jc1.per_class.label_Hbb: rejection of L162_MASS over L162, run by run :: sha256 588e935996d1006a :: prose twin: two significant figures, direction in words
\newcommand{\MlstFtJcOneRejHbbResMassOverResPrMax}{equal within 10\%}  % experiments/FIGS/data/v2/curves/full_budget_long_leg2.json :: cells.<run>.N1000000.s1.jc1.per_class.label_Hbb: rejection of L162_MASS over L162, run by run :: sha256 588e935996d1006a :: prose twin of the run farther from 1
\newcommand{\MlstFtJcOneRejHbbCoarseMassOverCoarse}{0.97 and 0.86}  % experiments/FIGS/data/v2/curves/full_budget_long_leg2.json :: cells.<run>.N1000000.s1.jc1.per_class.label_Hbb: rejection of R16_Q1_MASS over R16_Q1, run by run :: sha256 588e935996d1006a
\newcommand{\MlstFtJcOneRejHbbCoarseMassOverCoarsePr}{equal within 20\%}  % experiments/FIGS/data/v2/curves/full_budget_long_leg2.json :: cells.<run>.N1000000.s1.jc1.per_class.label_Hbb: rejection of R16_Q1_MASS over R16_Q1, run by run :: sha256 588e935996d1006a :: prose twin: two significant figures, direction in words
\newcommand{\MlstFtJcOneRejHbbCoarseMassOverCoarsePrMax}{equal within 20\%}  % experiments/FIGS/data/v2/curves/full_budget_long_leg2.json :: cells.<run>.N1000000.s1.jc1.per_class.label_Hbb: rejection of R16_Q1_MASS over R16_Q1, run by run :: sha256 588e935996d1006a :: prose twin of the run farther from 1
\newcommand{\MlstFtJcOneRejHbbRandOverFull}{0.92\,\ensuremath{\pm}\,0.11}  % experiments/FIGS/data/v2/curves/full_budget_long_leg2.json :: per_class.label_Hbb: rejection of each random partition over the 188-class run of its seed :: sha256 588e935996d1006a
\newcommand{\MlstFtJcOneRejHbbRandOverFullPr}{equal within 20\%}  % experiments/FIGS/data/v2/curves/full_budget_long_leg2.json :: per_class.label_Hbb: rejection of each random partition over the 188-class run of its seed :: sha256 588e935996d1006a :: prose twin: two significant figures, direction in words
\newcommand{\MlstFtJcOneRejHbbRandOverFullRange}{0.84--1.1}  % experiments/FIGS/data/v2/curves/full_budget_long_leg2.json :: per_class.label_Hbb: smallest and largest of the partitions' ratios :: sha256 588e935996d1006a
\newcommand{\MlstFtJcOneRejHccResOverFull}{1.07\,\ensuremath{\pm}\,0.12}  % experiments/FIGS/data/v2/curves/full_budget_long_leg2.json :: cells.<run>.N1000000.s1.jc1.per_class.label_Hcc: rejection of L162 over L188, run by run :: sha256 588e935996d1006a
\newcommand{\MlstFtJcOneRejHccResOverFullPr}{equal within 20\%}  % experiments/FIGS/data/v2/curves/full_budget_long_leg2.json :: cells.<run>.N1000000.s1.jc1.per_class.label_Hcc: rejection of L162 over L188, run by run :: sha256 588e935996d1006a :: prose twin: two significant figures, direction in words
\newcommand{\MlstFtJcOneRejHccMidOverFull}{0.99\,\ensuremath{\pm}\,0.10}  % experiments/FIGS/data/v2/curves/full_budget_long_leg2.json :: cells.<run>.N1000000.s1.jc1.per_class.label_Hcc: rejection of R42_Q1 over L188, run by run :: sha256 588e935996d1006a
\newcommand{\MlstFtJcOneRejHccMidOverFullPr}{equal within 20\%}  % experiments/FIGS/data/v2/curves/full_budget_long_leg2.json :: cells.<run>.N1000000.s1.jc1.per_class.label_Hcc: rejection of R42_Q1 over L188, run by run :: sha256 588e935996d1006a :: prose twin: two significant figures, direction in words
\newcommand{\MlstFtJcOneRejHccCoarseOverFull}{0.91\,\ensuremath{\pm}\,0.07}  % experiments/FIGS/data/v2/curves/full_budget_long_leg2.json :: cells.<run>.N1000000.s1.jc1.per_class.label_Hcc: rejection of R16_Q1 over L188, run by run :: sha256 588e935996d1006a
\newcommand{\MlstFtJcOneRejHccCoarseOverFullPr}{equal within 10\%}  % experiments/FIGS/data/v2/curves/full_budget_long_leg2.json :: cells.<run>.N1000000.s1.jc1.per_class.label_Hcc: rejection of R16_Q1 over L188, run by run :: sha256 588e935996d1006a :: prose twin: two significant figures, direction in words
\newcommand{\MlstFtJcOneRejHccResMassOverRes}{0.886 and 0.816}  % experiments/FIGS/data/v2/curves/full_budget_long_leg2.json :: cells.<run>.N1000000.s1.jc1.per_class.label_Hcc: rejection of L162_MASS over L162, run by run :: sha256 588e935996d1006a
\newcommand{\MlstFtJcOneRejHccResMassOverResPr}{equal within 20\%}  % experiments/FIGS/data/v2/curves/full_budget_long_leg2.json :: cells.<run>.N1000000.s1.jc1.per_class.label_Hcc: rejection of L162_MASS over L162, run by run :: sha256 588e935996d1006a :: prose twin: two significant figures, direction in words
\newcommand{\MlstFtJcOneRejHccResMassOverResPrMax}{equal within 20\%}  % experiments/FIGS/data/v2/curves/full_budget_long_leg2.json :: cells.<run>.N1000000.s1.jc1.per_class.label_Hcc: rejection of L162_MASS over L162, run by run :: sha256 588e935996d1006a :: prose twin of the run farther from 1
\newcommand{\MlstFtJcOneRejHccCoarseMassOverCoarse}{0.98 and 0.85}  % experiments/FIGS/data/v2/curves/full_budget_long_leg2.json :: cells.<run>.N1000000.s1.jc1.per_class.label_Hcc: rejection of R16_Q1_MASS over R16_Q1, run by run :: sha256 588e935996d1006a
\newcommand{\MlstFtJcOneRejHccCoarseMassOverCoarsePr}{equal within 20\%}  % experiments/FIGS/data/v2/curves/full_budget_long_leg2.json :: cells.<run>.N1000000.s1.jc1.per_class.label_Hcc: rejection of R16_Q1_MASS over R16_Q1, run by run :: sha256 588e935996d1006a :: prose twin: two significant figures, direction in words
\newcommand{\MlstFtJcOneRejHccCoarseMassOverCoarsePrMax}{equal within 20\%}  % experiments/FIGS/data/v2/curves/full_budget_long_leg2.json :: cells.<run>.N1000000.s1.jc1.per_class.label_Hcc: rejection of R16_Q1_MASS over R16_Q1, run by run :: sha256 588e935996d1006a :: prose twin of the run farther from 1
\newcommand{\MlstFtJcOneRejHccRandOverFull}{1.03\,\ensuremath{\pm}\,0.07}  % experiments/FIGS/data/v2/curves/full_budget_long_leg2.json :: per_class.label_Hcc: rejection of each random partition over the 188-class run of its seed :: sha256 588e935996d1006a
\newcommand{\MlstFtJcOneRejHccRandOverFullPr}{equal within 8\%}  % experiments/FIGS/data/v2/curves/full_budget_long_leg2.json :: per_class.label_Hcc: rejection of each random partition over the 188-class run of its seed :: sha256 588e935996d1006a :: prose twin: two significant figures, direction in words
\newcommand{\MlstFtJcOneRejHccRandOverFullRange}{0.94--1.1}  % experiments/FIGS/data/v2/curves/full_budget_long_leg2.json :: per_class.label_Hcc: smallest and largest of the partitions' ratios :: sha256 588e935996d1006a
\newcommand{\MlstFtJcOneRejHggResOverFull}{0.98\,\ensuremath{\pm}\,0.05}  % experiments/FIGS/data/v2/curves/full_budget_long_leg2.json :: cells.<run>.N1000000.s1.jc1.per_class.label_Hgg: rejection of L162 over L188, run by run :: sha256 588e935996d1006a
\newcommand{\MlstFtJcOneRejHggResOverFullPr}{equal within 6\%}  % experiments/FIGS/data/v2/curves/full_budget_long_leg2.json :: cells.<run>.N1000000.s1.jc1.per_class.label_Hgg: rejection of L162 over L188, run by run :: sha256 588e935996d1006a :: prose twin: two significant figures, direction in words
\newcommand{\MlstFtJcOneRejHggMidOverFull}{1.000\,\ensuremath{\pm}\,0.025}  % experiments/FIGS/data/v2/curves/full_budget_long_leg2.json :: cells.<run>.N1000000.s1.jc1.per_class.label_Hgg: rejection of R42_Q1 over L188, run by run :: sha256 588e935996d1006a
\newcommand{\MlstFtJcOneRejHggMidOverFullPr}{equal within 3\%}  % experiments/FIGS/data/v2/curves/full_budget_long_leg2.json :: cells.<run>.N1000000.s1.jc1.per_class.label_Hgg: rejection of R42_Q1 over L188, run by run :: sha256 588e935996d1006a :: prose twin: two significant figures, direction in words
\newcommand{\MlstFtJcOneRejHggCoarseOverFull}{0.97\,\ensuremath{\pm}\,0.09}  % experiments/FIGS/data/v2/curves/full_budget_long_leg2.json :: cells.<run>.N1000000.s1.jc1.per_class.label_Hgg: rejection of R16_Q1 over L188, run by run :: sha256 588e935996d1006a
\newcommand{\MlstFtJcOneRejHggCoarseOverFullPr}{equal within 10\%}  % experiments/FIGS/data/v2/curves/full_budget_long_leg2.json :: cells.<run>.N1000000.s1.jc1.per_class.label_Hgg: rejection of R16_Q1 over L188, run by run :: sha256 588e935996d1006a :: prose twin: two significant figures, direction in words
\newcommand{\MlstFtJcOneRejHggResMassOverRes}{1.046 and 1.008}  % experiments/FIGS/data/v2/curves/full_budget_long_leg2.json :: cells.<run>.N1000000.s1.jc1.per_class.label_Hgg: rejection of L162_MASS over L162, run by run :: sha256 588e935996d1006a
\newcommand{\MlstFtJcOneRejHggResMassOverResPr}{equal within 5\%}  % experiments/FIGS/data/v2/curves/full_budget_long_leg2.json :: cells.<run>.N1000000.s1.jc1.per_class.label_Hgg: rejection of L162_MASS over L162, run by run :: sha256 588e935996d1006a :: prose twin: two significant figures, direction in words
\newcommand{\MlstFtJcOneRejHggResMassOverResPrMax}{equal within 5\%}  % experiments/FIGS/data/v2/curves/full_budget_long_leg2.json :: cells.<run>.N1000000.s1.jc1.per_class.label_Hgg: rejection of L162_MASS over L162, run by run :: sha256 588e935996d1006a :: prose twin of the run farther from 1
\newcommand{\MlstFtJcOneRejHggCoarseMassOverCoarse}{1.12 and 1.01}  % experiments/FIGS/data/v2/curves/full_budget_long_leg2.json :: cells.<run>.N1000000.s1.jc1.per_class.label_Hgg: rejection of R16_Q1_MASS over R16_Q1, run by run :: sha256 588e935996d1006a
\newcommand{\MlstFtJcOneRejHggCoarseMassOverCoarsePr}{equal within 20\%}  % experiments/FIGS/data/v2/curves/full_budget_long_leg2.json :: cells.<run>.N1000000.s1.jc1.per_class.label_Hgg: rejection of R16_Q1_MASS over R16_Q1, run by run :: sha256 588e935996d1006a :: prose twin: two significant figures, direction in words
\newcommand{\MlstFtJcOneRejHggCoarseMassOverCoarsePrMax}{equal within 20\%}  % experiments/FIGS/data/v2/curves/full_budget_long_leg2.json :: cells.<run>.N1000000.s1.jc1.per_class.label_Hgg: rejection of R16_Q1_MASS over R16_Q1, run by run :: sha256 588e935996d1006a :: prose twin of the run farther from 1
\newcommand{\MlstFtJcOneRejHggRandOverFull}{0.968\,\ensuremath{\pm}\,0.016}  % experiments/FIGS/data/v2/curves/full_budget_long_leg2.json :: per_class.label_Hgg: rejection of each random partition over the 188-class run of its seed :: sha256 588e935996d1006a
\newcommand{\MlstFtJcOneRejHggRandOverFullPr}{equal within 4\%}  % experiments/FIGS/data/v2/curves/full_budget_long_leg2.json :: per_class.label_Hgg: rejection of each random partition over the 188-class run of its seed :: sha256 588e935996d1006a :: prose twin: two significant figures, direction in words
\newcommand{\MlstFtJcOneRejHggRandOverFullRange}{0.94--0.98}  % experiments/FIGS/data/v2/curves/full_budget_long_leg2.json :: per_class.label_Hgg: smallest and largest of the partitions' ratios :: sha256 588e935996d1006a
\newcommand{\MlstFtJcOneRejHfourqResOverFull}{1.05\,\ensuremath{\pm}\,0.06}  % experiments/FIGS/data/v2/curves/full_budget_long_leg2.json :: cells.<run>.N1000000.s1.jc1.per_class.label_H4q: rejection of L162 over L188, run by run :: sha256 588e935996d1006a
\newcommand{\MlstFtJcOneRejHfourqResOverFullPr}{equal within 7\%}  % experiments/FIGS/data/v2/curves/full_budget_long_leg2.json :: cells.<run>.N1000000.s1.jc1.per_class.label_H4q: rejection of L162 over L188, run by run :: sha256 588e935996d1006a :: prose twin: two significant figures, direction in words
\newcommand{\MlstFtJcOneRejHfourqMidOverFull}{0.95\,\ensuremath{\pm}\,0.04}  % experiments/FIGS/data/v2/curves/full_budget_long_leg2.json :: cells.<run>.N1000000.s1.jc1.per_class.label_H4q: rejection of R42_Q1 over L188, run by run :: sha256 588e935996d1006a
\newcommand{\MlstFtJcOneRejHfourqMidOverFullPr}{equal within 6\%}  % experiments/FIGS/data/v2/curves/full_budget_long_leg2.json :: cells.<run>.N1000000.s1.jc1.per_class.label_H4q: rejection of R42_Q1 over L188, run by run :: sha256 588e935996d1006a :: prose twin: two significant figures, direction in words
\newcommand{\MlstFtJcOneRejHfourqCoarseOverFull}{0.962\,\ensuremath{\pm}\,0.025}  % experiments/FIGS/data/v2/curves/full_budget_long_leg2.json :: cells.<run>.N1000000.s1.jc1.per_class.label_H4q: rejection of R16_Q1 over L188, run by run :: sha256 588e935996d1006a
\newcommand{\MlstFtJcOneRejHfourqCoarseOverFullPr}{equal within 4\%}  % experiments/FIGS/data/v2/curves/full_budget_long_leg2.json :: cells.<run>.N1000000.s1.jc1.per_class.label_H4q: rejection of R16_Q1 over L188, run by run :: sha256 588e935996d1006a :: prose twin: two significant figures, direction in words
\newcommand{\MlstFtJcOneRejHfourqResMassOverRes}{0.912 and 0.966}  % experiments/FIGS/data/v2/curves/full_budget_long_leg2.json :: cells.<run>.N1000000.s1.jc1.per_class.label_H4q: rejection of L162_MASS over L162, run by run :: sha256 588e935996d1006a
\newcommand{\MlstFtJcOneRejHfourqResMassOverResPr}{equal within 9\%}  % experiments/FIGS/data/v2/curves/full_budget_long_leg2.json :: cells.<run>.N1000000.s1.jc1.per_class.label_H4q: rejection of L162_MASS over L162, run by run :: sha256 588e935996d1006a :: prose twin: two significant figures, direction in words
\newcommand{\MlstFtJcOneRejHfourqResMassOverResPrMax}{equal within 9\%}  % experiments/FIGS/data/v2/curves/full_budget_long_leg2.json :: cells.<run>.N1000000.s1.jc1.per_class.label_H4q: rejection of L162_MASS over L162, run by run :: sha256 588e935996d1006a :: prose twin of the run farther from 1
\newcommand{\MlstFtJcOneRejHfourqCoarseMassOverCoarse}{0.937 and 0.888}  % experiments/FIGS/data/v2/curves/full_budget_long_leg2.json :: cells.<run>.N1000000.s1.jc1.per_class.label_H4q: rejection of R16_Q1_MASS over R16_Q1, run by run :: sha256 588e935996d1006a
\newcommand{\MlstFtJcOneRejHfourqCoarseMassOverCoarsePr}{equal within 20\%}  % experiments/FIGS/data/v2/curves/full_budget_long_leg2.json :: cells.<run>.N1000000.s1.jc1.per_class.label_H4q: rejection of R16_Q1_MASS over R16_Q1, run by run :: sha256 588e935996d1006a :: prose twin: two significant figures, direction in words
\newcommand{\MlstFtJcOneRejHfourqCoarseMassOverCoarsePrMax}{equal within 20\%}  % experiments/FIGS/data/v2/curves/full_budget_long_leg2.json :: cells.<run>.N1000000.s1.jc1.per_class.label_H4q: rejection of R16_Q1_MASS over R16_Q1, run by run :: sha256 588e935996d1006a :: prose twin of the run farther from 1
\newcommand{\MlstFtJcOneRejHfourqRandOverFull}{0.93\,\ensuremath{\pm}\,0.04}  % experiments/FIGS/data/v2/curves/full_budget_long_leg2.json :: per_class.label_H4q: rejection of each random partition over the 188-class run of its seed :: sha256 588e935996d1006a
\newcommand{\MlstFtJcOneRejHfourqRandOverFullPr}{equal within 8\%}  % experiments/FIGS/data/v2/curves/full_budget_long_leg2.json :: per_class.label_H4q: rejection of each random partition over the 188-class run of its seed :: sha256 588e935996d1006a :: prose twin: two significant figures, direction in words
\newcommand{\MlstFtJcOneRejHfourqRandOverFullRange}{0.86--0.96}  % experiments/FIGS/data/v2/curves/full_budget_long_leg2.json :: per_class.label_H4q: smallest and largest of the partitions' ratios :: sha256 588e935996d1006a
\newcommand{\MlstFtJcOneRejHqqlResOverFull}{1.004\,\ensuremath{\pm}\,0.018}  % experiments/FIGS/data/v2/curves/full_budget_long_leg2.json :: cells.<run>.N1000000.s1.jc1.per_class.label_Hqql: rejection of L162 over L188, run by run :: sha256 588e935996d1006a
\newcommand{\MlstFtJcOneRejHqqlResOverFullPr}{equal within 2\%}  % experiments/FIGS/data/v2/curves/full_budget_long_leg2.json :: cells.<run>.N1000000.s1.jc1.per_class.label_Hqql: rejection of L162 over L188, run by run :: sha256 588e935996d1006a :: prose twin: two significant figures, direction in words
\newcommand{\MlstFtJcOneRejHqqlMidOverFull}{1.04\,\ensuremath{\pm}\,0.18}  % experiments/FIGS/data/v2/curves/full_budget_long_leg2.json :: cells.<run>.N1000000.s1.jc1.per_class.label_Hqql: rejection of R42_Q1 over L188, run by run :: sha256 588e935996d1006a
\newcommand{\MlstFtJcOneRejHqqlMidOverFullPr}{equal within 20\%}  % experiments/FIGS/data/v2/curves/full_budget_long_leg2.json :: cells.<run>.N1000000.s1.jc1.per_class.label_Hqql: rejection of R42_Q1 over L188, run by run :: sha256 588e935996d1006a :: prose twin: two significant figures, direction in words
\newcommand{\MlstFtJcOneRejHqqlCoarseOverFull}{0.93\,\ensuremath{\pm}\,0.16}  % experiments/FIGS/data/v2/curves/full_budget_long_leg2.json :: cells.<run>.N1000000.s1.jc1.per_class.label_Hqql: rejection of R16_Q1 over L188, run by run :: sha256 588e935996d1006a
\newcommand{\MlstFtJcOneRejHqqlCoarseOverFullPr}{equal within 20\%}  % experiments/FIGS/data/v2/curves/full_budget_long_leg2.json :: cells.<run>.N1000000.s1.jc1.per_class.label_Hqql: rejection of R16_Q1 over L188, run by run :: sha256 588e935996d1006a :: prose twin: two significant figures, direction in words
\newcommand{\MlstFtJcOneRejHqqlResMassOverRes}{0.986 and 1.000}  % experiments/FIGS/data/v2/curves/full_budget_long_leg2.json :: cells.<run>.N1000000.s1.jc1.per_class.label_Hqql: rejection of L162_MASS over L162, run by run :: sha256 588e935996d1006a
\newcommand{\MlstFtJcOneRejHqqlResMassOverResPr}{equal within 2\%}  % experiments/FIGS/data/v2/curves/full_budget_long_leg2.json :: cells.<run>.N1000000.s1.jc1.per_class.label_Hqql: rejection of L162_MASS over L162, run by run :: sha256 588e935996d1006a :: prose twin: two significant figures, direction in words
\newcommand{\MlstFtJcOneRejHqqlResMassOverResPrMax}{equal within 2\%}  % experiments/FIGS/data/v2/curves/full_budget_long_leg2.json :: cells.<run>.N1000000.s1.jc1.per_class.label_Hqql: rejection of L162_MASS over L162, run by run :: sha256 588e935996d1006a :: prose twin of the run farther from 1
\newcommand{\MlstFtJcOneRejHqqlCoarseMassOverCoarse}{0.97 and 1.09}  % experiments/FIGS/data/v2/curves/full_budget_long_leg2.json :: cells.<run>.N1000000.s1.jc1.per_class.label_Hqql: rejection of R16_Q1_MASS over R16_Q1, run by run :: sha256 588e935996d1006a
\newcommand{\MlstFtJcOneRejHqqlCoarseMassOverCoarsePr}{equal within 9\%}  % experiments/FIGS/data/v2/curves/full_budget_long_leg2.json :: cells.<run>.N1000000.s1.jc1.per_class.label_Hqql: rejection of R16_Q1_MASS over R16_Q1, run by run :: sha256 588e935996d1006a :: prose twin: two significant figures, direction in words
\newcommand{\MlstFtJcOneRejHqqlCoarseMassOverCoarsePrMax}{equal within 9\%}  % experiments/FIGS/data/v2/curves/full_budget_long_leg2.json :: cells.<run>.N1000000.s1.jc1.per_class.label_Hqql: rejection of R16_Q1_MASS over R16_Q1, run by run :: sha256 588e935996d1006a :: prose twin of the run farther from 1
\newcommand{\MlstFtJcOneRejHqqlRandOverFull}{0.94\,\ensuremath{\pm}\,0.12}  % experiments/FIGS/data/v2/curves/full_budget_long_leg2.json :: per_class.label_Hqql: rejection of each random partition over the 188-class run of its seed :: sha256 588e935996d1006a
\newcommand{\MlstFtJcOneRejHqqlRandOverFullPr}{equal within 20\%}  % experiments/FIGS/data/v2/curves/full_budget_long_leg2.json :: per_class.label_Hqql: rejection of each random partition over the 188-class run of its seed :: sha256 588e935996d1006a :: prose twin: two significant figures, direction in words
\newcommand{\MlstFtJcOneRejHqqlRandOverFullRange}{0.79--1.1}  % experiments/FIGS/data/v2/curves/full_budget_long_leg2.json :: per_class.label_Hqql: smallest and largest of the partitions' ratios :: sha256 588e935996d1006a
\newcommand{\MlstFtJcOneRejTbqqResOverFull}{1.1\,\ensuremath{\pm}\,0.4}  % experiments/FIGS/data/v2/curves/full_budget_long_leg2.json :: cells.<run>.N1000000.s1.jc1.per_class.label_Tbqq: rejection of L162 over L188, run by run :: sha256 588e935996d1006a
\newcommand{\MlstFtJcOneRejTbqqResOverFullPr}{equal within 50\%}  % experiments/FIGS/data/v2/curves/full_budget_long_leg2.json :: cells.<run>.N1000000.s1.jc1.per_class.label_Tbqq: rejection of L162 over L188, run by run :: sha256 588e935996d1006a :: prose twin: two significant figures, direction in words
\newcommand{\MlstFtJcOneRejTbqqMidOverFull}{1.10\,\ensuremath{\pm}\,0.05}  % experiments/FIGS/data/v2/curves/full_budget_long_leg2.json :: cells.<run>.N1000000.s1.jc1.per_class.label_Tbqq: rejection of R42_Q1 over L188, run by run :: sha256 588e935996d1006a
\newcommand{\MlstFtJcOneRejTbqqMidOverFullPr}{equal within 10\%}  % experiments/FIGS/data/v2/curves/full_budget_long_leg2.json :: cells.<run>.N1000000.s1.jc1.per_class.label_Tbqq: rejection of R42_Q1 over L188, run by run :: sha256 588e935996d1006a :: prose twin: two significant figures, direction in words
\newcommand{\MlstFtJcOneRejTbqqCoarseOverFull}{0.85\,\ensuremath{\pm}\,0.32}  % experiments/FIGS/data/v2/curves/full_budget_long_leg2.json :: cells.<run>.N1000000.s1.jc1.per_class.label_Tbqq: rejection of R16_Q1 over L188, run by run :: sha256 588e935996d1006a
\newcommand{\MlstFtJcOneRejTbqqCoarseOverFullPr}{equal within 40\%}  % experiments/FIGS/data/v2/curves/full_budget_long_leg2.json :: cells.<run>.N1000000.s1.jc1.per_class.label_Tbqq: rejection of R16_Q1 over L188, run by run :: sha256 588e935996d1006a :: prose twin: two significant figures, direction in words
\newcommand{\MlstFtJcOneRejTbqqResMassOverRes}{0.89 and 1.21}  % experiments/FIGS/data/v2/curves/full_budget_long_leg2.json :: cells.<run>.N1000000.s1.jc1.per_class.label_Tbqq: rejection of L162_MASS over L162, run by run :: sha256 588e935996d1006a
\newcommand{\MlstFtJcOneRejTbqqResMassOverResPr}{equal within 30\%}  % experiments/FIGS/data/v2/curves/full_budget_long_leg2.json :: cells.<run>.N1000000.s1.jc1.per_class.label_Tbqq: rejection of L162_MASS over L162, run by run :: sha256 588e935996d1006a :: prose twin: two significant figures, direction in words
\newcommand{\MlstFtJcOneRejTbqqResMassOverResPrMax}{equal within 30\%}  % experiments/FIGS/data/v2/curves/full_budget_long_leg2.json :: cells.<run>.N1000000.s1.jc1.per_class.label_Tbqq: rejection of L162_MASS over L162, run by run :: sha256 588e935996d1006a :: prose twin of the run farther from 1
\newcommand{\MlstFtJcOneRejTbqqCoarseMassOverCoarse}{1.00 and 1.79}  % experiments/FIGS/data/v2/curves/full_budget_long_leg2.json :: cells.<run>.N1000000.s1.jc1.per_class.label_Tbqq: rejection of R16_Q1_MASS over R16_Q1, run by run :: sha256 588e935996d1006a
\newcommand{\MlstFtJcOneRejTbqqCoarseMassOverCoarsePr}{equal within 80\%}  % experiments/FIGS/data/v2/curves/full_budget_long_leg2.json :: cells.<run>.N1000000.s1.jc1.per_class.label_Tbqq: rejection of R16_Q1_MASS over R16_Q1, run by run :: sha256 588e935996d1006a :: prose twin: two significant figures, direction in words
\newcommand{\MlstFtJcOneRejTbqqCoarseMassOverCoarsePrMax}{equal within 80\%}  % experiments/FIGS/data/v2/curves/full_budget_long_leg2.json :: cells.<run>.N1000000.s1.jc1.per_class.label_Tbqq: rejection of R16_Q1_MASS over R16_Q1, run by run :: sha256 588e935996d1006a :: prose twin of the run farther from 1
\newcommand{\MlstFtJcOneRejTbqqRandOverFull}{0.79\,\ensuremath{\pm}\,0.15}  % experiments/FIGS/data/v2/curves/full_budget_long_leg2.json :: per_class.label_Tbqq: rejection of each random partition over the 188-class run of its seed :: sha256 588e935996d1006a
\newcommand{\MlstFtJcOneRejTbqqRandOverFullPr}{equal within 30\%}  % experiments/FIGS/data/v2/curves/full_budget_long_leg2.json :: per_class.label_Tbqq: rejection of each random partition over the 188-class run of its seed :: sha256 588e935996d1006a :: prose twin: two significant figures, direction in words
\newcommand{\MlstFtJcOneRejTbqqRandOverFullRange}{0.64--1.0}  % experiments/FIGS/data/v2/curves/full_budget_long_leg2.json :: per_class.label_Tbqq: smallest and largest of the partitions' ratios :: sha256 588e935996d1006a
\newcommand{\MlstFtJcOneRejTblResOverFull}{0.85\,\ensuremath{\pm}\,0.17}  % experiments/FIGS/data/v2/curves/full_budget_long_leg2.json :: cells.<run>.N1000000.s1.jc1.per_class.label_Tbl: rejection of L162 over L188, run by run :: sha256 588e935996d1006a
\newcommand{\MlstFtJcOneRejTblResOverFullPr}{equal within 20\%}  % experiments/FIGS/data/v2/curves/full_budget_long_leg2.json :: cells.<run>.N1000000.s1.jc1.per_class.label_Tbl: rejection of L162 over L188, run by run :: sha256 588e935996d1006a :: prose twin: two significant figures, direction in words
\newcommand{\MlstFtJcOneRejTblMidOverFull}{0.91\,\ensuremath{\pm}\,0.15}  % experiments/FIGS/data/v2/curves/full_budget_long_leg2.json :: cells.<run>.N1000000.s1.jc1.per_class.label_Tbl: rejection of R42_Q1 over L188, run by run :: sha256 588e935996d1006a
\newcommand{\MlstFtJcOneRejTblMidOverFullPr}{equal within 20\%}  % experiments/FIGS/data/v2/curves/full_budget_long_leg2.json :: cells.<run>.N1000000.s1.jc1.per_class.label_Tbl: rejection of R42_Q1 over L188, run by run :: sha256 588e935996d1006a :: prose twin: two significant figures, direction in words
\newcommand{\MlstFtJcOneRejTblCoarseOverFull}{0.94\,\ensuremath{\pm}\,0.28}  % experiments/FIGS/data/v2/curves/full_budget_long_leg2.json :: cells.<run>.N1000000.s1.jc1.per_class.label_Tbl: rejection of R16_Q1 over L188, run by run :: sha256 588e935996d1006a
\newcommand{\MlstFtJcOneRejTblCoarseOverFullPr}{equal within 30\%}  % experiments/FIGS/data/v2/curves/full_budget_long_leg2.json :: cells.<run>.N1000000.s1.jc1.per_class.label_Tbl: rejection of R16_Q1 over L188, run by run :: sha256 588e935996d1006a :: prose twin: two significant figures, direction in words
\newcommand{\MlstFtJcOneRejTblResMassOverRes}{1.04 and 0.93}  % experiments/FIGS/data/v2/curves/full_budget_long_leg2.json :: cells.<run>.N1000000.s1.jc1.per_class.label_Tbl: rejection of L162_MASS over L162, run by run :: sha256 588e935996d1006a
\newcommand{\MlstFtJcOneRejTblResMassOverResPr}{equal within 7\%}  % experiments/FIGS/data/v2/curves/full_budget_long_leg2.json :: cells.<run>.N1000000.s1.jc1.per_class.label_Tbl: rejection of L162_MASS over L162, run by run :: sha256 588e935996d1006a :: prose twin: two significant figures, direction in words
\newcommand{\MlstFtJcOneRejTblResMassOverResPrMax}{equal within 7\%}  % experiments/FIGS/data/v2/curves/full_budget_long_leg2.json :: cells.<run>.N1000000.s1.jc1.per_class.label_Tbl: rejection of L162_MASS over L162, run by run :: sha256 588e935996d1006a :: prose twin of the run farther from 1
\newcommand{\MlstFtJcOneRejTblCoarseMassOverCoarse}{0.84 and 1.22}  % experiments/FIGS/data/v2/curves/full_budget_long_leg2.json :: cells.<run>.N1000000.s1.jc1.per_class.label_Tbl: rejection of R16_Q1_MASS over R16_Q1, run by run :: sha256 588e935996d1006a
\newcommand{\MlstFtJcOneRejTblCoarseMassOverCoarsePr}{equal within 30\%}  % experiments/FIGS/data/v2/curves/full_budget_long_leg2.json :: cells.<run>.N1000000.s1.jc1.per_class.label_Tbl: rejection of R16_Q1_MASS over R16_Q1, run by run :: sha256 588e935996d1006a :: prose twin: two significant figures, direction in words
\newcommand{\MlstFtJcOneRejTblCoarseMassOverCoarsePrMax}{equal within 30\%}  % experiments/FIGS/data/v2/curves/full_budget_long_leg2.json :: cells.<run>.N1000000.s1.jc1.per_class.label_Tbl: rejection of R16_Q1_MASS over R16_Q1, run by run :: sha256 588e935996d1006a :: prose twin of the run farther from 1
\newcommand{\MlstFtJcOneRejTblRandOverFull}{0.88\,\ensuremath{\pm}\,0.26}  % experiments/FIGS/data/v2/curves/full_budget_long_leg2.json :: per_class.label_Tbl: rejection of each random partition over the 188-class run of its seed :: sha256 588e935996d1006a
\newcommand{\MlstFtJcOneRejTblRandOverFullPr}{equal within 30\%}  % experiments/FIGS/data/v2/curves/full_budget_long_leg2.json :: per_class.label_Tbl: rejection of each random partition over the 188-class run of its seed :: sha256 588e935996d1006a :: prose twin: two significant figures, direction in words
\newcommand{\MlstFtJcOneRejTblRandOverFullRange}{0.61--1.2}  % experiments/FIGS/data/v2/curves/full_budget_long_leg2.json :: per_class.label_Tbl: smallest and largest of the partitions' ratios :: sha256 588e935996d1006a
\newcommand{\MlstFtJcOneRejWqqResOverFull}{1.039\,\ensuremath{\pm}\,0.035}  % experiments/FIGS/data/v2/curves/full_budget_long_leg2.json :: cells.<run>.N1000000.s1.jc1.per_class.label_Wqq: rejection of L162 over L188, run by run :: sha256 588e935996d1006a
\newcommand{\MlstFtJcOneRejWqqResOverFullPr}{equal within 4\%}  % experiments/FIGS/data/v2/curves/full_budget_long_leg2.json :: cells.<run>.N1000000.s1.jc1.per_class.label_Wqq: rejection of L162 over L188, run by run :: sha256 588e935996d1006a :: prose twin: two significant figures, direction in words
\newcommand{\MlstFtJcOneRejWqqMidOverFull}{1.005\,\ensuremath{\pm}\,0.032}  % experiments/FIGS/data/v2/curves/full_budget_long_leg2.json :: cells.<run>.N1000000.s1.jc1.per_class.label_Wqq: rejection of R42_Q1 over L188, run by run :: sha256 588e935996d1006a
\newcommand{\MlstFtJcOneRejWqqMidOverFullPr}{equal within 4\%}  % experiments/FIGS/data/v2/curves/full_budget_long_leg2.json :: cells.<run>.N1000000.s1.jc1.per_class.label_Wqq: rejection of R42_Q1 over L188, run by run :: sha256 588e935996d1006a :: prose twin: two significant figures, direction in words
\newcommand{\MlstFtJcOneRejWqqCoarseOverFull}{1.01\,\ensuremath{\pm}\,0.07}  % experiments/FIGS/data/v2/curves/full_budget_long_leg2.json :: cells.<run>.N1000000.s1.jc1.per_class.label_Wqq: rejection of R16_Q1 over L188, run by run :: sha256 588e935996d1006a
\newcommand{\MlstFtJcOneRejWqqCoarseOverFullPr}{equal within 8\%}  % experiments/FIGS/data/v2/curves/full_budget_long_leg2.json :: cells.<run>.N1000000.s1.jc1.per_class.label_Wqq: rejection of R16_Q1 over L188, run by run :: sha256 588e935996d1006a :: prose twin: two significant figures, direction in words
\newcommand{\MlstFtJcOneRejWqqResMassOverRes}{0.996 and 0.964}  % experiments/FIGS/data/v2/curves/full_budget_long_leg2.json :: cells.<run>.N1000000.s1.jc1.per_class.label_Wqq: rejection of L162_MASS over L162, run by run :: sha256 588e935996d1006a
\newcommand{\MlstFtJcOneRejWqqResMassOverResPr}{equal within 4\%}  % experiments/FIGS/data/v2/curves/full_budget_long_leg2.json :: cells.<run>.N1000000.s1.jc1.per_class.label_Wqq: rejection of L162_MASS over L162, run by run :: sha256 588e935996d1006a :: prose twin: two significant figures, direction in words
\newcommand{\MlstFtJcOneRejWqqResMassOverResPrMax}{equal within 4\%}  % experiments/FIGS/data/v2/curves/full_budget_long_leg2.json :: cells.<run>.N1000000.s1.jc1.per_class.label_Wqq: rejection of L162_MASS over L162, run by run :: sha256 588e935996d1006a :: prose twin of the run farther from 1
\newcommand{\MlstFtJcOneRejWqqCoarseMassOverCoarse}{1.025 and 0.926}  % experiments/FIGS/data/v2/curves/full_budget_long_leg2.json :: cells.<run>.N1000000.s1.jc1.per_class.label_Wqq: rejection of R16_Q1_MASS over R16_Q1, run by run :: sha256 588e935996d1006a
\newcommand{\MlstFtJcOneRejWqqCoarseMassOverCoarsePr}{equal within 8\%}  % experiments/FIGS/data/v2/curves/full_budget_long_leg2.json :: cells.<run>.N1000000.s1.jc1.per_class.label_Wqq: rejection of R16_Q1_MASS over R16_Q1, run by run :: sha256 588e935996d1006a :: prose twin: two significant figures, direction in words
\newcommand{\MlstFtJcOneRejWqqCoarseMassOverCoarsePrMax}{equal within 8\%}  % experiments/FIGS/data/v2/curves/full_budget_long_leg2.json :: cells.<run>.N1000000.s1.jc1.per_class.label_Wqq: rejection of R16_Q1_MASS over R16_Q1, run by run :: sha256 588e935996d1006a :: prose twin of the run farther from 1
\newcommand{\MlstFtJcOneRejWqqRandOverFull}{1.02\,\ensuremath{\pm}\,0.07}  % experiments/FIGS/data/v2/curves/full_budget_long_leg2.json :: per_class.label_Wqq: rejection of each random partition over the 188-class run of its seed :: sha256 588e935996d1006a
\newcommand{\MlstFtJcOneRejWqqRandOverFullPr}{equal within 7\%}  % experiments/FIGS/data/v2/curves/full_budget_long_leg2.json :: per_class.label_Wqq: rejection of each random partition over the 188-class run of its seed :: sha256 588e935996d1006a :: prose twin: two significant figures, direction in words
\newcommand{\MlstFtJcOneRejWqqRandOverFullRange}{0.93--1.1}  % experiments/FIGS/data/v2/curves/full_budget_long_leg2.json :: per_class.label_Wqq: smallest and largest of the partitions' ratios :: sha256 588e935996d1006a
\newcommand{\MlstFtJcOneRejZqqResOverFull}{0.987\,\ensuremath{\pm}\,0.019}  % experiments/FIGS/data/v2/curves/full_budget_long_leg2.json :: cells.<run>.N1000000.s1.jc1.per_class.label_Zqq: rejection of L162 over L188, run by run :: sha256 588e935996d1006a
\newcommand{\MlstFtJcOneRejZqqResOverFullPr}{equal within 2\%}  % experiments/FIGS/data/v2/curves/full_budget_long_leg2.json :: cells.<run>.N1000000.s1.jc1.per_class.label_Zqq: rejection of L162 over L188, run by run :: sha256 588e935996d1006a :: prose twin: two significant figures, direction in words
\newcommand{\MlstFtJcOneRejZqqMidOverFull}{0.984\,\ensuremath{\pm}\,0.010}  % experiments/FIGS/data/v2/curves/full_budget_long_leg2.json :: cells.<run>.N1000000.s1.jc1.per_class.label_Zqq: rejection of R42_Q1 over L188, run by run :: sha256 588e935996d1006a
\newcommand{\MlstFtJcOneRejZqqMidOverFullPr}{equal within 2\%}  % experiments/FIGS/data/v2/curves/full_budget_long_leg2.json :: cells.<run>.N1000000.s1.jc1.per_class.label_Zqq: rejection of R42_Q1 over L188, run by run :: sha256 588e935996d1006a :: prose twin: two significant figures, direction in words
\newcommand{\MlstFtJcOneRejZqqCoarseOverFull}{0.91\,\ensuremath{\pm}\,0.05}  % experiments/FIGS/data/v2/curves/full_budget_long_leg2.json :: cells.<run>.N1000000.s1.jc1.per_class.label_Zqq: rejection of R16_Q1 over L188, run by run :: sha256 588e935996d1006a
\newcommand{\MlstFtJcOneRejZqqCoarseOverFullPr}{equal within 9\%}  % experiments/FIGS/data/v2/curves/full_budget_long_leg2.json :: cells.<run>.N1000000.s1.jc1.per_class.label_Zqq: rejection of R16_Q1 over L188, run by run :: sha256 588e935996d1006a :: prose twin: two significant figures, direction in words
\newcommand{\MlstFtJcOneRejZqqResMassOverRes}{0.9803 and 0.9820}  % experiments/FIGS/data/v2/curves/full_budget_long_leg2.json :: cells.<run>.N1000000.s1.jc1.per_class.label_Zqq: rejection of L162_MASS over L162, run by run :: sha256 588e935996d1006a
\newcommand{\MlstFtJcOneRejZqqResMassOverResPr}{equal within 2\%}  % experiments/FIGS/data/v2/curves/full_budget_long_leg2.json :: cells.<run>.N1000000.s1.jc1.per_class.label_Zqq: rejection of L162_MASS over L162, run by run :: sha256 588e935996d1006a :: prose twin: two significant figures, direction in words
\newcommand{\MlstFtJcOneRejZqqResMassOverResPrMax}{equal within 2\%}  % experiments/FIGS/data/v2/curves/full_budget_long_leg2.json :: cells.<run>.N1000000.s1.jc1.per_class.label_Zqq: rejection of L162_MASS over L162, run by run :: sha256 588e935996d1006a :: prose twin of the run farther from 1
\newcommand{\MlstFtJcOneRejZqqCoarseMassOverCoarse}{1.06 and 0.94}  % experiments/FIGS/data/v2/curves/full_budget_long_leg2.json :: cells.<run>.N1000000.s1.jc1.per_class.label_Zqq: rejection of R16_Q1_MASS over R16_Q1, run by run :: sha256 588e935996d1006a
\newcommand{\MlstFtJcOneRejZqqCoarseMassOverCoarsePr}{equal within 7\%}  % experiments/FIGS/data/v2/curves/full_budget_long_leg2.json :: cells.<run>.N1000000.s1.jc1.per_class.label_Zqq: rejection of R16_Q1_MASS over R16_Q1, run by run :: sha256 588e935996d1006a :: prose twin: two significant figures, direction in words
\newcommand{\MlstFtJcOneRejZqqCoarseMassOverCoarsePrMax}{equal within 7\%}  % experiments/FIGS/data/v2/curves/full_budget_long_leg2.json :: cells.<run>.N1000000.s1.jc1.per_class.label_Zqq: rejection of R16_Q1_MASS over R16_Q1, run by run :: sha256 588e935996d1006a :: prose twin of the run farther from 1
\newcommand{\MlstFtJcOneRejZqqRandOverFull}{0.952\,\ensuremath{\pm}\,0.028}  % experiments/FIGS/data/v2/curves/full_budget_long_leg2.json :: per_class.label_Zqq: rejection of each random partition over the 188-class run of its seed :: sha256 588e935996d1006a
\newcommand{\MlstFtJcOneRejZqqRandOverFullPr}{equal within 5\%}  % experiments/FIGS/data/v2/curves/full_budget_long_leg2.json :: per_class.label_Zqq: rejection of each random partition over the 188-class run of its seed :: sha256 588e935996d1006a :: prose twin: two significant figures, direction in words
\newcommand{\MlstFtJcOneRejZqqRandOverFullRange}{0.93--0.99}  % experiments/FIGS/data/v2/curves/full_budget_long_leg2.json :: per_class.label_Zqq: smallest and largest of the partitions' ratios :: sha256 588e935996d1006a
\newcommand{\MlstFtJcOneResNRejDrop}{zero}  % experiments/FIGS/data/v2/curves/full_budget_long_leg2.json :: per_class rejections :: sha256 588e935996d1006a :: classes whose paired rejection ratio to the 188-class model departs from 1 by more than twice its SD and its counting uncertainty, columns with fewer than JC1_MIN_BKG passing jets excluded, ratio below 1 (count)
\newcommand{\MlstFtJcOneResNRejDiff}{zero}  % experiments/FIGS/data/v2/curves/full_budget_long_leg2.json :: per_class rejections :: sha256 588e935996d1006a :: classes whose paired rejection ratio to the 188-class model departs from 1 by more than twice its SD and its counting uncertainty, columns with fewer than JC1_MIN_BKG passing jets excluded (count)
\newcommand{\MlstFtJcOneResRejDropClasses}{none}  % experiments/FIGS/data/v2/curves/full_budget_long_leg2.json :: per_class rejections :: sha256 588e935996d1006a :: the classes of ...NRejDrop, in ParT's column order
\newcommand{\MlstFtJcOneMidNRejDrop}{zero}  % experiments/FIGS/data/v2/curves/full_budget_long_leg2.json :: per_class rejections :: sha256 588e935996d1006a :: classes whose paired rejection ratio to the 188-class model departs from 1 by more than twice its SD and its counting uncertainty, columns with fewer than JC1_MIN_BKG passing jets excluded, ratio below 1 (count)
\newcommand{\MlstFtJcOneMidNRejDiff}{zero}  % experiments/FIGS/data/v2/curves/full_budget_long_leg2.json :: per_class rejections :: sha256 588e935996d1006a :: classes whose paired rejection ratio to the 188-class model departs from 1 by more than twice its SD and its counting uncertainty, columns with fewer than JC1_MIN_BKG passing jets excluded (count)
\newcommand{\MlstFtJcOneMidRejDropClasses}{none}  % experiments/FIGS/data/v2/curves/full_budget_long_leg2.json :: per_class rejections :: sha256 588e935996d1006a :: the classes of ...NRejDrop, in ParT's column order
\newcommand{\MlstFtJcOneCoarseNRejDrop}{zero}  % experiments/FIGS/data/v2/curves/full_budget_long_leg2.json :: per_class rejections :: sha256 588e935996d1006a :: classes whose paired rejection ratio to the 188-class model departs from 1 by more than twice its SD and its counting uncertainty, columns with fewer than JC1_MIN_BKG passing jets excluded, ratio below 1 (count)
\newcommand{\MlstFtJcOneCoarseNRejDiff}{zero}  % experiments/FIGS/data/v2/curves/full_budget_long_leg2.json :: per_class rejections :: sha256 588e935996d1006a :: classes whose paired rejection ratio to the 188-class model departs from 1 by more than twice its SD and its counting uncertainty, columns with fewer than JC1_MIN_BKG passing jets excluded (count)
\newcommand{\MlstFtJcOneCoarseRejDropClasses}{none}  % experiments/FIGS/data/v2/curves/full_budget_long_leg2.json :: per_class rejections :: sha256 588e935996d1006a :: the classes of ...NRejDrop, in ParT's column order
\newcommand{\MlstFtJcOneAccRandMinusCoarse}{+0.002}  % experiments/FIGS/data/v2/curves/full_budget_long_leg2.json :: cells.*.N1000000.s1.jc1.accuracy :: sha256 588e935996d1006a :: mean over the random partitions minus the mean over the 17-class runs
\newcommand{\MlstFtJcOneAccFullBoot}{0.847}  % experiments/FIGS/data/v2/budget/jc1_acc_boot_N1000000.json :: diffs.r16q1-l188.accuracy_ref (mean) :: sha256 7cccb7b0ff900fb8
\newcommand{\MlstFtJcOneAccCoarseBoot}{0.842}  % experiments/FIGS/data/v2/budget/jc1_acc_boot_N1000000.json :: diffs.r16q1-l188.accuracy_config (mean) :: sha256 7cccb7b0ff900fb8
\newcommand{\MlstFtJcOneErrResOverFullBoot}{consistent}  % experiments/FIGS/data/v2/budget/jc1_acc_boot_N1000000.json :: diffs.l162-l188.error_ratio :: sha256 7cccb7b0ff900fb8 :: delta and sigma = hypot(run SD, bootstrap SD) in per cent; 'consistent' when |delta| <= 3 sigma
\newcommand{\MlstFtJcOneErrResOverFullZ}{0.6}  % experiments/FIGS/data/v2/budget/jc1_acc_boot_N1000000.json :: diffs.l162-l188.error_ratio.z :: sha256 7cccb7b0ff900fb8 :: |delta| / sigma
\newcommand{\MlstFtJcOneErrMidOverFullBoot}{consistent}  % experiments/FIGS/data/v2/budget/jc1_acc_boot_N1000000.json :: diffs.r42q1-l188.error_ratio :: sha256 7cccb7b0ff900fb8 :: delta and sigma = hypot(run SD, bootstrap SD) in per cent; 'consistent' when |delta| <= 3 sigma
\newcommand{\MlstFtJcOneErrMidOverFullZ}{1.4}  % experiments/FIGS/data/v2/budget/jc1_acc_boot_N1000000.json :: diffs.r42q1-l188.error_ratio.z :: sha256 7cccb7b0ff900fb8 :: |delta| / sigma
\newcommand{\MlstFtJcOneErrCoarseOverFullBoot}{(3.0\,\ensuremath{\pm}\,0.6)\% more errors}  % experiments/FIGS/data/v2/budget/jc1_acc_boot_N1000000.json :: diffs.r16q1-l188.error_ratio :: sha256 7cccb7b0ff900fb8 :: delta and sigma = hypot(run SD, bootstrap SD) in per cent; 'consistent' when |delta| <= 3 sigma
\newcommand{\MlstFtJcOneErrCoarseOverFullZ}{4.8}  % experiments/FIGS/data/v2/budget/jc1_acc_boot_N1000000.json :: diffs.r16q1-l188.error_ratio.z :: sha256 7cccb7b0ff900fb8 :: |delta| / sigma
\newcommand{\MlstFtJcOneBootReplicas}{200}  % experiments/FIGS/data/v2/budget/jc1_acc_boot_N1000000.json :: replicas :: sha256 7cccb7b0ff900fb8
\newcommand{\MlstFtJcOneBootJets}{2{,}000{,}000}  % experiments/FIGS/data/v2/budget/jc1_acc_boot_N1000000.json :: n_test_jets :: sha256 7cccb7b0ff900fb8
\newcommand{\ComputeTotalGpuHours}{3{,}435}  % experiments/FIGS/data/v2/index.json :: GPU-hours of every reported pretraining run, training and validation passes :: sha256 c6c4138ee5c86423
\newcommand{\MlstComputeTotalGpuHours}{3{,}510}  % experiments/FIGS/data/v2/index.json :: GPU-hours of every pretraining run the paper reports, the kinematic-only model included :: sha256 c6c4138ee5c86423
\newcommand{\ComputeMainGpu}{V100}  % experiments/FIGS/data/v2/index.json :: GPU type with the most hours over the four label sets :: sha256 c6c4138ee5c86423
\newcommand{\ComputeMainGpuShareMin}{56\%}  % experiments/FIGS/data/v2/index.json :: lowest share of a label set's GPU-hours on that type :: sha256 c6c4138ee5c86423
\newcommand{\ComputeMainGpuShareMax}{71\%}  % experiments/FIGS/data/v2/index.json :: highest share of a label set's GPU-hours on that type :: sha256 c6c4138ee5c86423
\newcommand{\ComputeRunsDone}{27}  % experiments/FIGS/data/v2/index.json :: paper runs with a compute record (count) :: sha256 c6c4138ee5c86423
\newcommand{\ComputeGpuTypes}{L4, RTX 3090, RTX 4090 and V100}  % experiments/FIGS/data/v2/index.json :: GPU types any finished paper run trained on :: sha256 c6c4138ee5c86423
\newcommand{\MlstComputeGpuTypes}{L4, RTX 3090, RTX 4090 and V100}  % experiments/FIGS/data/v2/index.json :: GPU types of the pretraining runs the paper reports :: sha256 c6c4138ee5c86423
\newcommand{\MlstFtJcOneErrCoarseOverFullBootEFour}{\pending{the JetClass bootstrap at 1e4 (after the equal-stopping read-out)}}  % experiments/FIGS/data/v2/index.json :: not computed yet :: sha256 c6c4138ee5c86423 :: waits on the JetClass bootstrap at 1e4 (after the equal-stopping read-out)
\newcommand{\MlstLpBbQgZeroAccDropFull}{\pending{pair probes with the quark/gluon inputs zeroed (open item O10, not run)}}  % experiments/FIGS/data/v2/index.json :: not computed yet :: sha256 c6c4138ee5c86423 :: waits on pair probes with the quark/gluon inputs zeroed (open item O10, not run)
\newcommand{\MlstLpBbQgZeroAccDropCoarse}{\pending{pair probes with the quark/gluon inputs zeroed (open item O10, not run)}}  % experiments/FIGS/data/v2/index.json :: not computed yet :: sha256 c6c4138ee5c86423 :: waits on pair probes with the quark/gluon inputs zeroed (open item O10, not run)
\newcommand{\XbsRthirtyFullOverScratch}{\pending{X->bs fits and from-scratch runs}}  % experiments/FIGS/data/v2/index.json :: not computed yet :: sha256 c6c4138ee5c86423 :: waits on X->bs fits and from-scratch runs
\newcommand{\XbsRthirtyFullOverScratchPr}{\pending{X->bs fits and from-scratch runs}}  % experiments/FIGS/data/v2/index.json :: not computed yet :: sha256 c6c4138ee5c86423 :: waits on X->bs fits and from-scratch runs
\newcommand{\XbsRthirtyCoarseOverFull}{\pending{X->bs fits}}  % experiments/FIGS/data/v2/index.json :: not computed yet :: sha256 c6c4138ee5c86423 :: waits on X->bs fits
\newcommand{\XbsRthirtyCoarseOverFullPr}{\pending{X->bs fits}}  % experiments/FIGS/data/v2/index.json :: not computed yet :: sha256 c6c4138ee5c86423 :: waits on X->bs fits
\newcommand{\XbsRthirtyNodeOverTransferFull}{\pending{X->bs fits}}  % experiments/FIGS/data/v2/index.json :: not computed yet :: sha256 c6c4138ee5c86423 :: waits on X->bs fits
\newcommand{\XbsRthirtyNodeOverTransferFullPr}{\pending{X->bs fits}}  % experiments/FIGS/data/v2/index.json :: not computed yet :: sha256 c6c4138ee5c86423 :: waits on X->bs fits
\newcommand{\XbbRthirtyDirectOverTransferFull}{\pending{X->bb fits}}  % experiments/FIGS/data/v2/index.json :: not computed yet :: sha256 c6c4138ee5c86423 :: waits on X->bb fits
\newcommand{\XbbRthirtyDirectOverTransferFullPr}{\pending{X->bb fits}}  % experiments/FIGS/data/v2/index.json :: not computed yet :: sha256 c6c4138ee5c86423 :: waits on X->bb fits
\newcommand{\XbbRthirtyTransferCoarseOverFull}{\pending{X->bb fits}}  % experiments/FIGS/data/v2/index.json :: not computed yet :: sha256 c6c4138ee5c86423 :: waits on X->bb fits
\newcommand{\XbbRthirtyTransferCoarseOverFullPr}{\pending{X->bb fits}}  % experiments/FIGS/data/v2/index.json :: not computed yet :: sha256 c6c4138ee5c86423 :: waits on X->bb fits
             % the counts below are its macros
% Visual language: Sophon's Fig. 1 (arXiv:2405.12972): the document's serif font, light fills,
% rounded boxes, light grey panels behind the stages (PI 2026-11-02). The text is set in the
% document's font at the printed sizes 8.5-10 pt (MLST 8-12 pt), never scaled.
% Text width: 15 cm (MLST full width). The picture is 15.0 cm wide and is drawn at its final
% size, never scaled, so the font sizes hold.
% Optional: \def\FoneDropXbx{} before % Figure 1 (F1 of docs/FIGURE_PLAN.md): the design of the study, in TikZ.
% \input{fig_f1} from main.tex; this file holds one figure environment and nothing else.
%
% WHAT main.tex MUST LOAD (this file loads nothing itself):
%   \usepackage{tikz}
%   \usetikzlibrary{arrows.meta}          % the only TikZ library used (arrow tips)
%   \definecolor{c188}{HTML}{1F77B4}      % the configuration colours of experiments/FIGS/
%   \definecolor{c162}{HTML}{FF7F0E}      % mlst_style.py (matplotlib's tab10 in order, as in
%   \definecolor{c43}{HTML}{2CA02C}       % Sophon's figures); the variants are drawn in their
%   \definecolor{c17}{HTML}{D62728}       % label set's colour with a dotted border, as the
%   \definecolor{crand}{HTML}{17BECF}     % plotted figures draw variants dotted
%   \definecolor{cflav}{HTML}{9467BD}
%   \input{results_generated}             % the counts below are its macros
% Visual language: Sophon's Fig. 1 (arXiv:2405.12972): the document's serif font, light fills,
% rounded boxes, light grey panels behind the stages (PI 2026-11-02). The text is set in the
% document's font at the printed sizes 8.5-10 pt (MLST 8-12 pt), never scaled.
% Text width: 15 cm (MLST full width). The picture is 15.0 cm wide and is drawn at its final
% size, never scaled, so the font sizes hold.
% Optional: \def\FoneDropXbx{} before \input{fig_f1} removes X->bb/X->bs from panel (c)
% (PI 2026-11-02: if the samples are not staged by 12 Nov, the X->bb/X->bs section goes).
%
% Every class box is drawn from the label maps (checked by hand, 2026-11-02):
%   (a) configs/labelmaps/rung_label_maps.v1.csv, columns L188, L162, R42_Q1, R16_Q1. Rows: the ten
%       decays X->QQ/gg in the order bb bc bq | cc cs cq | ss sq qq gg (the 43-class b / c / light
%       groups are then contiguous), then ee mm | tauh-taue tauh-taum, then the 27 QCD classes.
%   (b) column: configs/labelmaps/rand_label_map.v2.csv, RAND2_p1 on the same rows: groups
%       0 0 0 0 0 9 9 9 0 3 | 11 11 | 15 4 | QCD, lettered A..F in order of first appearance.
%       strips: configs/labelmaps/flavour_pair_map.v2.csv, the 22 classes of orbit X_YY_QQQQ.
%       Order with a b quark: bbbb bbcc bbss bbqq qqqb bbcq ccbs ccbq qqbc qqbs bcsq;
%       without: cccc ccss ccqq ssss ssqq qqqq qqqc qqqs ccsq sscq qqcs.
%       FLAV_F0: all 22 in group 9. FLAV_F1: the 11 with a b in group 11, the rest 9.
%       FLAV_F1R: group 11 = bbss bbqq qqqb bbcq ccbq qqbc (with b, positions 2-5, 7, 8 from 0)
%       and ccss ccqq ssss qqqc sscq (without, positions 1-3, 6, 9); filled square = group 11.
% Geometry (cm): row height 0.35 from y = -1.18, gaps of 0.14 between the quark/gluon, lepton
% and QCD blocks, the QCD row 0.70; boxes inset 0.025 top and bottom; columns 1.45 wide with
% left edges 1.75, 3.28, 4.81, 6.34 (hierarchy) and 8.25 (random partition).
\providecommand{\Kfull}{\VocabSizeLoneeighteight}
\providecommand{\Kres}{\VocabSizeLonesixtwo}
\providecommand{\Kmid}{\VocabSizeRfourtwoQone}
\providecommand{\Kcoarse}{\VocabSizeRonesixQone}
\begin{figure}[tbp]
\centering
\ifdim\linewidth<14.99cm\typeout{fig_f1.tex: WARNING: the drawing is 15 cm wide but \string\linewidth\space
  is \the\linewidth; set a 15 cm text width (it is not scaled, to keep its 8.5 pt text)}\fi
\def\FoneText{\rmfamily\fontsize{8.5}{9.5}\selectfont}%
\def\FoneMain{\rmfamily\fontsize{9}{10.5}\selectfont}%
\def\FoneBold{\rmfamily\bfseries\fontsize{9}{10.5}\selectfont}%
\def\FoneTitle{\rmfamily\bfseries\fontsize{10}{12}\selectfont}%
\ifdefined\FoneDropXbx\def\FoneXbx{}\else\def\FoneXbx{\\$X\to b\bar b$, $X\to bs$}\fi
\begin{tikzpicture}[
  x=1cm, y=1cm,
  font=\FoneMain,
  ttl/.style={anchor=base west, font=\FoneTitle, inner sep=0pt},
  cls/.style={draw=black!55, line width=0.45pt, rounded corners=1.2pt},
  hdr/.style={draw=#1!70!black, fill=#1!22, line width=0.6pt, rounded corners=2.5pt, align=center,
              font=\FoneText, inner sep=0pt},
  inbox/.style={align=center, font=\FoneText, inner sep=0pt},
  letter/.style={font=\FoneBold, inner sep=0pt},
  rowlab/.style={anchor=east, inner sep=0pt},
  txt/.style={align=left, font=\FoneText, inner sep=0pt},
  sq/.style={draw=black!65, line width=0.4pt},
  chip/.style={draw=#1!70!black, fill=#1!22, line width=0.6pt, rounded corners=3pt,
               align=center, font=\FoneText, inner sep=0pt, anchor=north west},
  stage/.style={draw=black!45, fill=black!4, line width=0.6pt, rounded corners=3pt, align=left,
                font=\FoneText, inner sep=0.1cm},
  flow/.style={-{Latex[length=1.8mm, width=1.4mm]}, line width=0.7pt, black!75},
]
% ================================================================ (a) the class hierarchy
\node[ttl] at (0,0) {(a) Class hierarchy};
% faint bands on every other row, drawn first so the boxes sit on top
\foreach \rowtop in {-1.53,-2.23,-2.93,-3.63,-4.33,-5.17,-5.87}
  \fill[black!4] (0,\rowtop) rectangle (9.70,\rowtop-0.35);
% column headers
\node[rowlab, align=right, font=\FoneText] at (1.65,-0.68) {JetClass-II\\class};
\node[hdr=c188, minimum width=1.45cm, minimum height=0.80cm, anchor=north west] at (1.75,-0.28) {\Kfull\\classes};
\node[hdr=c162, minimum width=1.45cm, minimum height=0.80cm, anchor=north west] at (3.28,-0.28) {\Kres\\classes};
\node[hdr=c43,  minimum width=1.45cm, minimum height=0.80cm, anchor=north west] at (4.81,-0.28) {\Kmid\\classes};
\node[hdr=c17,  minimum width=1.45cm, minimum height=0.80cm, anchor=north west] at (6.34,-0.28) {\Kcoarse\\classes};
% row labels
\node[rowlab] at (1.65,-1.355) {$X\to b\bar b$};
\node[rowlab] at (1.65,-1.705) {$X\to bc$};
\node[rowlab] at (1.65,-2.055) {$X\to bq$};
\node[rowlab] at (1.65,-2.405) {$X\to c\bar c$};
\node[rowlab] at (1.65,-2.755) {$X\to cs$};
\node[rowlab] at (1.65,-3.105) {$X\to cq$};
\node[rowlab] at (1.65,-3.455) {$X\to s\bar s$};
\node[rowlab] at (1.65,-3.805) {$X\to sq$};
\node[rowlab] at (1.65,-4.155) {$X\to q\bar q$};
\node[rowlab] at (1.65,-4.505) {$X\to gg$};
\node[rowlab] at (1.65,-4.995) {$X\to ee$};
\node[rowlab] at (1.65,-5.345) {$X\to\mu\mu$};
\node[rowlab] at (1.65,-5.695) {$X\to\tau_h\tau_e$};
\node[rowlab] at (1.65,-6.045) {$X\to\tau_h\tau_\mu$};
\node[rowlab] at (1.65,-6.710) {QCD jets};
% 188 and 162 classes: every resonant decay is its own class
\foreach \rowtop in {-1.18,-1.53,-1.88,-2.23,-2.58,-2.93,-3.28,-3.63,-3.98,-4.33,-4.82,-5.17,-5.52,-5.87} {
  \draw[cls, fill=c188!15] (1.75,\rowtop-0.025) rectangle (3.20,\rowtop-0.325);
  \draw[cls, fill=c162!15] (3.28,\rowtop-0.025) rectangle (4.73,\rowtop-0.325);
}
\draw[cls, fill=c188!15] (1.75,-6.385) rectangle (3.20,-7.035);
\node[inbox] at (2.475,-6.71) {\LitJetClassTwoNQcdClasses\\classes};
\draw[cls, fill=c162!15] (3.28,-6.385) rectangle (4.73,-7.035);
\node[inbox] at (4.005,-6.71) {1 class};
% 43 classes: decay topology x {b, else c, else light quark or gluon}; e and mu merged
\draw[cls, fill=c43!15] (4.81,-1.205) rectangle (6.26,-2.205);
\node[inbox] at (5.535,-1.705) {with $b$};
\draw[cls, fill=c43!15] (4.81,-2.255) rectangle (6.26,-3.255);
\node[inbox] at (5.535,-2.755) {with $c$,\\no $b$};
\draw[cls, fill=c43!15] (4.81,-3.305) rectangle (6.26,-4.655);
\node[inbox] at (5.535,-3.980) {$u$, $d$, $s$\\or $g$};
\draw[cls, fill=c43!15] (4.81,-4.845) rectangle (6.26,-5.495);
\node[inbox] at (5.535,-5.170) {$\ell\ell$};
\draw[cls, fill=c43!15] (4.81,-5.545) rectangle (6.26,-6.195);
\node[inbox] at (5.535,-5.870) {$\tau_h\tau_\ell$};
\draw[cls, fill=c43!15] (4.81,-6.385) rectangle (6.26,-7.035);
\node[inbox] at (5.535,-6.710) {1 class};
% 17 classes: number and type of prongs, no flavor
\draw[cls, fill=c17!15] (6.34,-1.205) rectangle (7.79,-4.655);
\node[inbox] at (7.065,-2.930) {two\\quarks\\or gluons};
\draw[cls, fill=c17!15] (6.34,-4.845) rectangle (7.79,-5.495);
\node[inbox] at (7.065,-5.170) {$\ell\ell$};
\draw[cls, fill=c17!15] (6.34,-5.545) rectangle (7.79,-6.195);
\node[inbox] at (7.065,-5.870) {$\tau_h\tau_\ell$};
\draw[cls, fill=c17!15] (6.34,-6.385) rectangle (7.79,-7.035);
\node[inbox] at (7.065,-6.710) {1 class};
% what each merge removes
\node[anchor=north west, font=\FoneBold, inner sep=0pt] at (0,-7.27) {What each merge removes};
\node[txt, anchor=north west, text width=6.04cm] (m1) at (1.75,-7.65)
  {the $b$, $c$ and $s$ content of the QCD jets};
\node[txt, anchor=north west, text width=6.04cm] (m2) at ([yshift=-0.10cm]m1.south west)
  {the number of $b$ and $c$ quarks, light-quark flavor, quark against gluon, and lepton
   flavor ($e$ against $\mu$)};
\node[txt, anchor=north west, text width=6.04cm] (m3) at ([yshift=-0.10cm]m2.south west)
  {quark flavor ($b$, $c$ or light); the number and type of prongs remain};
\node[txt, anchor=north east] at ([xshift=-0.10cm]m1.north west) {\Kfull\,$\to$\,\Kres};
\node[txt, anchor=north east] at ([xshift=-0.10cm]m2.north west) {\Kres\,$\to$\,\Kmid};
\node[txt, anchor=north east] at ([xshift=-0.10cm]m3.north west) {\Kmid\,$\to$\,\Kcoarse};
% ================================================================ (b) the controls
\node[ttl] at (8.25,0) {(b) Controls, \Kcoarse{} classes each};
% the first random partition on the rows of (a); one letter = one class
\node[hdr=crand, minimum width=1.45cm, minimum height=0.80cm, anchor=north west] at (8.25,-0.28) {random\\partition};
\draw[cls, fill=crand!15] (8.25,-1.205) rectangle (9.70,-2.905);
\node[letter] at (8.975,-2.055) {A};
\draw[cls, fill=crand!15] (8.25,-2.955) rectangle (9.70,-3.955);
\node[letter] at (8.975,-3.455) {B};
\draw[cls, fill=crand!15] (8.25,-4.005) rectangle (9.70,-4.305);
\node[letter] at (8.975,-4.155) {A};
\draw[cls, fill=crand!15] (8.25,-4.355) rectangle (9.70,-4.655);
\node[letter] at (8.975,-4.505) {C};
\draw[cls, fill=crand!15] (8.25,-4.845) rectangle (9.70,-5.495);
\node[letter] at (8.975,-5.170) {D};
\draw[cls, fill=crand!15] (8.25,-5.545) rectangle (9.70,-5.845);
\node[letter] at (8.975,-5.695) {E};
\draw[cls, fill=crand!15] (8.25,-5.895) rectangle (9.70,-6.195);
\node[letter] at (8.975,-6.045) {F};
\draw[cls, fill=crand!15] (8.25,-6.385) rectangle (9.70,-7.035);
\node[inbox] at (8.975,-6.710) {1 class};
% how the random partitions are built
\node[hdr=crand, anchor=north west, text width=4.75cm, inner sep=0.1cm, align=left, font=\FoneBold]
  (rh) at (10.05,-0.28) {Random partitions (\RandVtwoNPartitions)};
\node[txt, anchor=north west, text width=4.95cm] (rb) at ([yshift=-0.12cm]rh.south west)
  {\textbullet\ groups of resonant classes drawn at random, plus one class of QCD jets\\[2pt]
   \textbullet\ group shares of training jets within \RandVtwoShareTol{} of the \Kcoarse-class groups\\[2pt]
   \textbullet\ best of \RandVtwoTrials{} draws: the fewest merges shared with the \Kcoarse-class level\\[2pt]
   \textbullet\ left: the first partition; boxes with the same letter are one class};
% how the random partitions with controlled flavor boundaries are built
\node[hdr=cflav, anchor=north west, text width=4.75cm, inner sep=0.1cm, align=left, font=\FoneBold]
  (fh) at ([yshift=-0.22cm]rb.south west) {Random partitions with controlled flavor boundaries};
\node[txt, anchor=north west, text width=4.95cm] (fn) at ([yshift=-0.12cm]fh.south west)
  {whole families of decays that differ only in quark flavor; two of them split one family,
   the \RandFlavSplitOrbitClasses{} decays \RandFlavSplitOrbit:};
\begin{scope}[shift={(fn.south west)}]
  % brackets over the decays with and without a b quark (squares at a 0.2 cm pitch)
  \node[txt, anchor=south] at (1.075,-0.45) {with $b$};
  \node[txt, anchor=south] at (3.425,-0.45) {without $b$};
  \draw[line width=0.5pt] (0,-0.54) -- (0,-0.47) -- (2.15,-0.47) -- (2.15,-0.54);
  \draw[line width=0.5pt] (2.35,-0.54) -- (2.35,-0.47) -- (4.50,-0.47) -- (4.50,-0.54);
  % no flavor boundary: all 22 in one class
  \node[txt, anchor=north west] at (0,-0.62) {no flavor boundary};
  \foreach \i in {0,...,10} {
    \draw[sq, fill=white] (0.2*\i,-1.10) rectangle ++(0.15,0.15);
    \draw[sq, fill=white] (2.35+0.2*\i,-1.10) rectangle ++(0.15,0.15);
  }
  % b-quark boundary: the decays with a b quark moved to another class
  \node[txt, anchor=north west] at (0,-1.22) {$b$-quark boundary};
  \foreach \i in {0,...,10} {
    \draw[sq, fill=cflav!70!black] (0.2*\i,-1.70) rectangle ++(0.15,0.15);
    \draw[sq, fill=white] (2.35+0.2*\i,-1.70) rectangle ++(0.15,0.15);
  }
  % random flavor boundary: the same number of decays, drawn at random
  \node[txt, anchor=north west] at (0,-1.82) {random flavor boundary};
  \foreach \i in {0,1,6,9,10}
    \draw[sq, fill=white] (0.2*\i,-2.30) rectangle ++(0.15,0.15);
  \foreach \i in {2,3,4,5,7,8}
    \draw[sq, fill=cflav!70!black] (0.2*\i,-2.30) rectangle ++(0.15,0.15);
  \foreach \i in {0,4,5,7,8,10}
    \draw[sq, fill=white] (2.35+0.2*\i,-2.30) rectangle ++(0.15,0.15);
  \foreach \i in {1,2,3,6,9}
    \draw[sq, fill=cflav!70!black] (2.35+0.2*\i,-2.30) rectangle ++(0.15,0.15);
  % key
  \draw[sq, fill=cflav!70!black] (0,-2.62) rectangle ++(0.15,0.15);
  \node[txt, anchor=west] (fkey) at (0.25,-2.545) {moved to another class};
\end{scope}
% ================================================================ (c) the pipeline
% (c) starts below the lower of the two columns above, whatever their text heights
\pgfpointanchor{m3}{south}\pgfgetlastxy{\FoneXa}{\FoneYa}%
\pgfpointanchor{fkey}{south}\pgfgetlastxy{\FoneXb}{\FoneYb}%
\pgfmathsetlengthmacro{\FoneYc}{min(\FoneYa,\FoneYb)-0.45cm}%
\coordinate (ctop) at (0pt,\FoneYc);
\begin{scope}[shift={(ctop)}]
  \node[ttl, anchor=north west] at (0,0) {(c) Pretraining and evaluation};
  % what every model of one seed shares
  \node[stage, text width=2.15cm] (seed) at (1.175,-1.96)
    {\textbf{Same for every model of one seed:} the training jets, their order and the
     initial weights, except the output layer};
  % one pretraining per label set: the configuration colours of every later figure
  \node[draw=black!45, fill=black!4, line width=0.6pt, rounded corners=3pt, minimum width=5.15cm,
        minimum height=3.08cm, anchor=north west, inner sep=0pt] (pre) at (2.81,-0.42) {};
  \node[anchor=north west, font=\FoneBold, inner sep=0pt] at (2.96,-0.52) {Pretraining on JetClass-II};
  \node[chip=c188, minimum width=1.13cm, minimum height=0.42cm] at (2.96,-0.95) {\Kfull};
  \node[chip=c162, minimum width=1.13cm, minimum height=0.42cm] at (4.20,-0.95) {\Kres};
  \node[chip=c43,  minimum width=1.13cm, minimum height=0.42cm] at (5.44,-0.95) {\Kmid};
  \node[chip=c17,  minimum width=1.13cm, minimum height=0.42cm] at (6.68,-0.95) {\Kcoarse};
  \node[chip=crand, minimum width=1.70cm, minimum height=0.70cm] at (2.96,-1.47) {random\\partitions};
  \node[chip=cflav, minimum width=3.04cm, minimum height=0.70cm] at (4.77,-1.47) {controlled flavor\\boundaries};
  \node[chip=c162, densely dotted, minimum width=1.13cm, minimum height=0.42cm] at (2.96,-2.32) {\Kres};
  \node[chip=c17,  densely dotted, minimum width=1.13cm, minimum height=0.42cm] at (4.20,-2.32) {\Kcoarse};
  \node[txt, anchor=west] at (5.44,-2.53) {with auxiliary\\mass regression};
  \node[chip=c188, densely dotted, minimum width=1.13cm, minimum height=0.42cm] at (2.96,-2.94) {\Kfull};
  \node[txt, anchor=west] at (4.20,-3.15) {kinematic inputs only};
  % evaluation
  \node[stage, text width=2.4cm] (fz) at (9.68,-0.95) {frozen features:\\linear probe,\\MLP probe};
  \node[stage, text width=2.4cm] (ft) at (9.68,-2.05) {fine-tuning\\of all weights};
  \node[stage, text width=2.4cm, draw=black, fill=white, line width=1.1pt] (sc) at (9.68,-3.15) {trained from\\scratch};
  \node[stage, text width=3.4cm, minimum height=3.2cm] (task) at (13.20,-2.00)
    {\textbf{Downstream tasks}\\JetClass-II: all classes\\JetClass-II: class pairs\\\quad(both class-balanced)\\JetClass\\
     top tagging\\quark/gluon, Pythia\\quark/gluon, Herwig\\jet mass\FoneXbx};
  \draw[flow] (seed.east) -- (seed.east -| pre.west);
  \draw[flow] (fz.west -| pre.east) -- (fz.west);
  \draw[flow] (ft.west -| pre.east) -- (ft.west);
  \draw[flow] (fz.east) -- (fz.east -| task.west);
  \draw[flow] (ft.east) -- (ft.east -| task.west);
  \draw[flow] (sc.east) -- (sc.east -| task.west);
\end{scope}
\end{tikzpicture}
\caption{Design of the study.
(a)~The class hierarchy on the decays of a heavy resonance~$X$ and on the jets from quantum
chromodynamics (QCD); $q$ is a $u$ or $d$ quark, $\ell$ an electron or a muon, and $\tau_h$,
$\tau_e$, $\tau_\mu$ a $\tau$ lepton decaying to hadrons, an electron or a muon. Each box is one
class of the label set above its column, except the QCD box of the \Kfull-class column (Sophon's
label set~\cite{sophon}), which stands for \LitJetClassTwoNQcdClasses{} classes. Each level merges
classes of the level to its left; the lines below name what each merge removes.
(b)~The controls: the first random partition, where boxes with the same letter are one class, and
the random partitions with controlled flavor boundaries on the decays \RandFlavSplitOrbit{} (one
square per decay; filled squares are moved to a second class).
(c)~Models of one seed share the training jets, their order and the initial weights, except the
output layer, the mass regression target, and the first layer and initial weights of the model with
kinematic inputs only (the particle four-momenta). The hierarchy levels have \MlstNRunsWord{} seeds
each, the models with the auxiliary mass regression \MlstMassNRunsWord{}, the random partitions with
controlled flavor boundaries \RandFlavRuns{} and each random partition \RandVtwoRuns{}. Frozen
features, the output of the backbone (the network below the classification head) with its weights held
fixed, feed linear probes (logistic regressions) and multilayer perceptron (MLP) probes; fine-tuning
trains all weights. Colors identify the label sets and controls in every figure; the plots give the three partitions with
controlled flavor boundaries three shades; dotted borders mark the two
variants, which the other figures draw as dotted lines in the color of their label set.}
\label{fig:design}
\end{figure}
 removes X->bb/X->bs from panel (c)
% (PI 2026-11-02: if the samples are not staged by 12 Nov, the X->bb/X->bs section goes).
%
% Every class box is drawn from the label maps (checked by hand, 2026-11-02):
%   (a) configs/labelmaps/rung_label_maps.v1.csv, columns L188, L162, R42_Q1, R16_Q1. Rows: the ten
%       decays X->QQ/gg in the order bb bc bq | cc cs cq | ss sq qq gg (the 43-class b / c / light
%       groups are then contiguous), then ee mm | tauh-taue tauh-taum, then the 27 QCD classes.
%   (b) column: configs/labelmaps/rand_label_map.v2.csv, RAND2_p1 on the same rows: groups
%       0 0 0 0 0 9 9 9 0 3 | 11 11 | 15 4 | QCD, lettered A..F in order of first appearance.
%       strips: configs/labelmaps/flavour_pair_map.v2.csv, the 22 classes of orbit X_YY_QQQQ.
%       Order with a b quark: bbbb bbcc bbss bbqq qqqb bbcq ccbs ccbq qqbc qqbs bcsq;
%       without: cccc ccss ccqq ssss ssqq qqqq qqqc qqqs ccsq sscq qqcs.
%       FLAV_F0: all 22 in group 9. FLAV_F1: the 11 with a b in group 11, the rest 9.
%       FLAV_F1R: group 11 = bbss bbqq qqqb bbcq ccbq qqbc (with b, positions 2-5, 7, 8 from 0)
%       and ccss ccqq ssss qqqc sscq (without, positions 1-3, 6, 9); filled square = group 11.
% Geometry (cm): row height 0.35 from y = -1.18, gaps of 0.14 between the quark/gluon, lepton
% and QCD blocks, the QCD row 0.70; boxes inset 0.025 top and bottom; columns 1.45 wide with
% left edges 1.75, 3.28, 4.81, 6.34 (hierarchy) and 8.25 (random partition).
\providecommand{\Kfull}{\VocabSizeLoneeighteight}
\providecommand{\Kres}{\VocabSizeLonesixtwo}
\providecommand{\Kmid}{\VocabSizeRfourtwoQone}
\providecommand{\Kcoarse}{\VocabSizeRonesixQone}
\begin{figure}[tbp]
\centering
\ifdim\linewidth<14.99cm\typeout{fig_f1.tex: WARNING: the drawing is 15 cm wide but \string\linewidth\space
  is \the\linewidth; set a 15 cm text width (it is not scaled, to keep its 8.5 pt text)}\fi
\def\FoneText{\rmfamily\fontsize{8.5}{9.5}\selectfont}%
\def\FoneMain{\rmfamily\fontsize{9}{10.5}\selectfont}%
\def\FoneBold{\rmfamily\bfseries\fontsize{9}{10.5}\selectfont}%
\def\FoneTitle{\rmfamily\bfseries\fontsize{10}{12}\selectfont}%
\ifdefined\FoneDropXbx\def\FoneXbx{}\else\def\FoneXbx{\\$X\to b\bar b$, $X\to bs$}\fi
\begin{tikzpicture}[
  x=1cm, y=1cm,
  font=\FoneMain,
  ttl/.style={anchor=base west, font=\FoneTitle, inner sep=0pt},
  cls/.style={draw=black!55, line width=0.45pt, rounded corners=1.2pt},
  hdr/.style={draw=#1!70!black, fill=#1!22, line width=0.6pt, rounded corners=2.5pt, align=center,
              font=\FoneText, inner sep=0pt},
  inbox/.style={align=center, font=\FoneText, inner sep=0pt},
  letter/.style={font=\FoneBold, inner sep=0pt},
  rowlab/.style={anchor=east, inner sep=0pt},
  txt/.style={align=left, font=\FoneText, inner sep=0pt},
  sq/.style={draw=black!65, line width=0.4pt},
  chip/.style={draw=#1!70!black, fill=#1!22, line width=0.6pt, rounded corners=3pt,
               align=center, font=\FoneText, inner sep=0pt, anchor=north west},
  stage/.style={draw=black!45, fill=black!4, line width=0.6pt, rounded corners=3pt, align=left,
                font=\FoneText, inner sep=0.1cm},
  flow/.style={-{Latex[length=1.8mm, width=1.4mm]}, line width=0.7pt, black!75},
]
% ================================================================ (a) the class hierarchy
\node[ttl] at (0,0) {(a) Class hierarchy};
% faint bands on every other row, drawn first so the boxes sit on top
\foreach \rowtop in {-1.53,-2.23,-2.93,-3.63,-4.33,-5.17,-5.87}
  \fill[black!4] (0,\rowtop) rectangle (9.70,\rowtop-0.35);
% column headers
\node[rowlab, align=right, font=\FoneText] at (1.65,-0.68) {JetClass-II\\class};
\node[hdr=c188, minimum width=1.45cm, minimum height=0.80cm, anchor=north west] at (1.75,-0.28) {\Kfull\\classes};
\node[hdr=c162, minimum width=1.45cm, minimum height=0.80cm, anchor=north west] at (3.28,-0.28) {\Kres\\classes};
\node[hdr=c43,  minimum width=1.45cm, minimum height=0.80cm, anchor=north west] at (4.81,-0.28) {\Kmid\\classes};
\node[hdr=c17,  minimum width=1.45cm, minimum height=0.80cm, anchor=north west] at (6.34,-0.28) {\Kcoarse\\classes};
% row labels
\node[rowlab] at (1.65,-1.355) {$X\to b\bar b$};
\node[rowlab] at (1.65,-1.705) {$X\to bc$};
\node[rowlab] at (1.65,-2.055) {$X\to bq$};
\node[rowlab] at (1.65,-2.405) {$X\to c\bar c$};
\node[rowlab] at (1.65,-2.755) {$X\to cs$};
\node[rowlab] at (1.65,-3.105) {$X\to cq$};
\node[rowlab] at (1.65,-3.455) {$X\to s\bar s$};
\node[rowlab] at (1.65,-3.805) {$X\to sq$};
\node[rowlab] at (1.65,-4.155) {$X\to q\bar q$};
\node[rowlab] at (1.65,-4.505) {$X\to gg$};
\node[rowlab] at (1.65,-4.995) {$X\to ee$};
\node[rowlab] at (1.65,-5.345) {$X\to\mu\mu$};
\node[rowlab] at (1.65,-5.695) {$X\to\tau_h\tau_e$};
\node[rowlab] at (1.65,-6.045) {$X\to\tau_h\tau_\mu$};
\node[rowlab] at (1.65,-6.710) {QCD jets};
% 188 and 162 classes: every resonant decay is its own class
\foreach \rowtop in {-1.18,-1.53,-1.88,-2.23,-2.58,-2.93,-3.28,-3.63,-3.98,-4.33,-4.82,-5.17,-5.52,-5.87} {
  \draw[cls, fill=c188!15] (1.75,\rowtop-0.025) rectangle (3.20,\rowtop-0.325);
  \draw[cls, fill=c162!15] (3.28,\rowtop-0.025) rectangle (4.73,\rowtop-0.325);
}
\draw[cls, fill=c188!15] (1.75,-6.385) rectangle (3.20,-7.035);
\node[inbox] at (2.475,-6.71) {\LitJetClassTwoNQcdClasses\\classes};
\draw[cls, fill=c162!15] (3.28,-6.385) rectangle (4.73,-7.035);
\node[inbox] at (4.005,-6.71) {1 class};
% 43 classes: decay topology x {b, else c, else light quark or gluon}; e and mu merged
\draw[cls, fill=c43!15] (4.81,-1.205) rectangle (6.26,-2.205);
\node[inbox] at (5.535,-1.705) {with $b$};
\draw[cls, fill=c43!15] (4.81,-2.255) rectangle (6.26,-3.255);
\node[inbox] at (5.535,-2.755) {with $c$,\\no $b$};
\draw[cls, fill=c43!15] (4.81,-3.305) rectangle (6.26,-4.655);
\node[inbox] at (5.535,-3.980) {$u$, $d$, $s$\\or $g$};
\draw[cls, fill=c43!15] (4.81,-4.845) rectangle (6.26,-5.495);
\node[inbox] at (5.535,-5.170) {$\ell\ell$};
\draw[cls, fill=c43!15] (4.81,-5.545) rectangle (6.26,-6.195);
\node[inbox] at (5.535,-5.870) {$\tau_h\tau_\ell$};
\draw[cls, fill=c43!15] (4.81,-6.385) rectangle (6.26,-7.035);
\node[inbox] at (5.535,-6.710) {1 class};
% 17 classes: number and type of prongs, no flavor
\draw[cls, fill=c17!15] (6.34,-1.205) rectangle (7.79,-4.655);
\node[inbox] at (7.065,-2.930) {two\\quarks\\or gluons};
\draw[cls, fill=c17!15] (6.34,-4.845) rectangle (7.79,-5.495);
\node[inbox] at (7.065,-5.170) {$\ell\ell$};
\draw[cls, fill=c17!15] (6.34,-5.545) rectangle (7.79,-6.195);
\node[inbox] at (7.065,-5.870) {$\tau_h\tau_\ell$};
\draw[cls, fill=c17!15] (6.34,-6.385) rectangle (7.79,-7.035);
\node[inbox] at (7.065,-6.710) {1 class};
% what each merge removes
\node[anchor=north west, font=\FoneBold, inner sep=0pt] at (0,-7.27) {What each merge removes};
\node[txt, anchor=north west, text width=6.04cm] (m1) at (1.75,-7.65)
  {the $b$, $c$ and $s$ content of the QCD jets};
\node[txt, anchor=north west, text width=6.04cm] (m2) at ([yshift=-0.10cm]m1.south west)
  {the number of $b$ and $c$ quarks, light-quark flavor, quark against gluon, and lepton
   flavor ($e$ against $\mu$)};
\node[txt, anchor=north west, text width=6.04cm] (m3) at ([yshift=-0.10cm]m2.south west)
  {quark flavor ($b$, $c$ or light); the number and type of prongs remain};
\node[txt, anchor=north east] at ([xshift=-0.10cm]m1.north west) {\Kfull\,$\to$\,\Kres};
\node[txt, anchor=north east] at ([xshift=-0.10cm]m2.north west) {\Kres\,$\to$\,\Kmid};
\node[txt, anchor=north east] at ([xshift=-0.10cm]m3.north west) {\Kmid\,$\to$\,\Kcoarse};
% ================================================================ (b) the controls
\node[ttl] at (8.25,0) {(b) Controls, \Kcoarse{} classes each};
% the first random partition on the rows of (a); one letter = one class
\node[hdr=crand, minimum width=1.45cm, minimum height=0.80cm, anchor=north west] at (8.25,-0.28) {random\\partition};
\draw[cls, fill=crand!15] (8.25,-1.205) rectangle (9.70,-2.905);
\node[letter] at (8.975,-2.055) {A};
\draw[cls, fill=crand!15] (8.25,-2.955) rectangle (9.70,-3.955);
\node[letter] at (8.975,-3.455) {B};
\draw[cls, fill=crand!15] (8.25,-4.005) rectangle (9.70,-4.305);
\node[letter] at (8.975,-4.155) {A};
\draw[cls, fill=crand!15] (8.25,-4.355) rectangle (9.70,-4.655);
\node[letter] at (8.975,-4.505) {C};
\draw[cls, fill=crand!15] (8.25,-4.845) rectangle (9.70,-5.495);
\node[letter] at (8.975,-5.170) {D};
\draw[cls, fill=crand!15] (8.25,-5.545) rectangle (9.70,-5.845);
\node[letter] at (8.975,-5.695) {E};
\draw[cls, fill=crand!15] (8.25,-5.895) rectangle (9.70,-6.195);
\node[letter] at (8.975,-6.045) {F};
\draw[cls, fill=crand!15] (8.25,-6.385) rectangle (9.70,-7.035);
\node[inbox] at (8.975,-6.710) {1 class};
% how the random partitions are built
\node[hdr=crand, anchor=north west, text width=4.75cm, inner sep=0.1cm, align=left, font=\FoneBold]
  (rh) at (10.05,-0.28) {Random partitions (\RandVtwoNPartitions)};
\node[txt, anchor=north west, text width=4.95cm] (rb) at ([yshift=-0.12cm]rh.south west)
  {\textbullet\ groups of resonant classes drawn at random, plus one class of QCD jets\\[2pt]
   \textbullet\ group shares of training jets within \RandVtwoShareTol{} of the \Kcoarse-class groups\\[2pt]
   \textbullet\ best of \RandVtwoTrials{} draws: the fewest merges shared with the \Kcoarse-class level\\[2pt]
   \textbullet\ left: the first partition; boxes with the same letter are one class};
% how the random partitions with controlled flavor boundaries are built
\node[hdr=cflav, anchor=north west, text width=4.75cm, inner sep=0.1cm, align=left, font=\FoneBold]
  (fh) at ([yshift=-0.22cm]rb.south west) {Random partitions with controlled flavor boundaries};
\node[txt, anchor=north west, text width=4.95cm] (fn) at ([yshift=-0.12cm]fh.south west)
  {whole families of decays that differ only in quark flavor; two of them split one family,
   the \RandFlavSplitOrbitClasses{} decays \RandFlavSplitOrbit:};
\begin{scope}[shift={(fn.south west)}]
  % brackets over the decays with and without a b quark (squares at a 0.2 cm pitch)
  \node[txt, anchor=south] at (1.075,-0.45) {with $b$};
  \node[txt, anchor=south] at (3.425,-0.45) {without $b$};
  \draw[line width=0.5pt] (0,-0.54) -- (0,-0.47) -- (2.15,-0.47) -- (2.15,-0.54);
  \draw[line width=0.5pt] (2.35,-0.54) -- (2.35,-0.47) -- (4.50,-0.47) -- (4.50,-0.54);
  % no flavor boundary: all 22 in one class
  \node[txt, anchor=north west] at (0,-0.62) {no flavor boundary};
  \foreach \i in {0,...,10} {
    \draw[sq, fill=white] (0.2*\i,-1.10) rectangle ++(0.15,0.15);
    \draw[sq, fill=white] (2.35+0.2*\i,-1.10) rectangle ++(0.15,0.15);
  }
  % b-quark boundary: the decays with a b quark moved to another class
  \node[txt, anchor=north west] at (0,-1.22) {$b$-quark boundary};
  \foreach \i in {0,...,10} {
    \draw[sq, fill=cflav!70!black] (0.2*\i,-1.70) rectangle ++(0.15,0.15);
    \draw[sq, fill=white] (2.35+0.2*\i,-1.70) rectangle ++(0.15,0.15);
  }
  % random flavor boundary: the same number of decays, drawn at random
  \node[txt, anchor=north west] at (0,-1.82) {random flavor boundary};
  \foreach \i in {0,1,6,9,10}
    \draw[sq, fill=white] (0.2*\i,-2.30) rectangle ++(0.15,0.15);
  \foreach \i in {2,3,4,5,7,8}
    \draw[sq, fill=cflav!70!black] (0.2*\i,-2.30) rectangle ++(0.15,0.15);
  \foreach \i in {0,4,5,7,8,10}
    \draw[sq, fill=white] (2.35+0.2*\i,-2.30) rectangle ++(0.15,0.15);
  \foreach \i in {1,2,3,6,9}
    \draw[sq, fill=cflav!70!black] (2.35+0.2*\i,-2.30) rectangle ++(0.15,0.15);
  % key
  \draw[sq, fill=cflav!70!black] (0,-2.62) rectangle ++(0.15,0.15);
  \node[txt, anchor=west] (fkey) at (0.25,-2.545) {moved to another class};
\end{scope}
% ================================================================ (c) the pipeline
% (c) starts below the lower of the two columns above, whatever their text heights
\pgfpointanchor{m3}{south}\pgfgetlastxy{\FoneXa}{\FoneYa}%
\pgfpointanchor{fkey}{south}\pgfgetlastxy{\FoneXb}{\FoneYb}%
\pgfmathsetlengthmacro{\FoneYc}{min(\FoneYa,\FoneYb)-0.45cm}%
\coordinate (ctop) at (0pt,\FoneYc);
\begin{scope}[shift={(ctop)}]
  \node[ttl, anchor=north west] at (0,0) {(c) Pretraining and evaluation};
  % what every model of one seed shares
  \node[stage, text width=2.15cm] (seed) at (1.175,-1.96)
    {\textbf{Same for every model of one seed:} the training jets, their order and the
     initial weights, except the output layer};
  % one pretraining per label set: the configuration colours of every later figure
  \node[draw=black!45, fill=black!4, line width=0.6pt, rounded corners=3pt, minimum width=5.15cm,
        minimum height=3.08cm, anchor=north west, inner sep=0pt] (pre) at (2.81,-0.42) {};
  \node[anchor=north west, font=\FoneBold, inner sep=0pt] at (2.96,-0.52) {Pretraining on JetClass-II};
  \node[chip=c188, minimum width=1.13cm, minimum height=0.42cm] at (2.96,-0.95) {\Kfull};
  \node[chip=c162, minimum width=1.13cm, minimum height=0.42cm] at (4.20,-0.95) {\Kres};
  \node[chip=c43,  minimum width=1.13cm, minimum height=0.42cm] at (5.44,-0.95) {\Kmid};
  \node[chip=c17,  minimum width=1.13cm, minimum height=0.42cm] at (6.68,-0.95) {\Kcoarse};
  \node[chip=crand, minimum width=1.70cm, minimum height=0.70cm] at (2.96,-1.47) {random\\partitions};
  \node[chip=cflav, minimum width=3.04cm, minimum height=0.70cm] at (4.77,-1.47) {controlled flavor\\boundaries};
  \node[chip=c162, densely dotted, minimum width=1.13cm, minimum height=0.42cm] at (2.96,-2.32) {\Kres};
  \node[chip=c17,  densely dotted, minimum width=1.13cm, minimum height=0.42cm] at (4.20,-2.32) {\Kcoarse};
  \node[txt, anchor=west] at (5.44,-2.53) {with auxiliary\\mass regression};
  \node[chip=c188, densely dotted, minimum width=1.13cm, minimum height=0.42cm] at (2.96,-2.94) {\Kfull};
  \node[txt, anchor=west] at (4.20,-3.15) {kinematic inputs only};
  % evaluation
  \node[stage, text width=2.4cm] (fz) at (9.68,-0.95) {frozen features:\\linear probe,\\MLP probe};
  \node[stage, text width=2.4cm] (ft) at (9.68,-2.05) {fine-tuning\\of all weights};
  \node[stage, text width=2.4cm, draw=black, fill=white, line width=1.1pt] (sc) at (9.68,-3.15) {trained from\\scratch};
  \node[stage, text width=3.4cm, minimum height=3.2cm] (task) at (13.20,-2.00)
    {\textbf{Downstream tasks}\\JetClass-II: all classes\\JetClass-II: class pairs\\\quad(both class-balanced)\\JetClass\\
     top tagging\\quark/gluon, Pythia\\quark/gluon, Herwig\\jet mass\FoneXbx};
  \draw[flow] (seed.east) -- (seed.east -| pre.west);
  \draw[flow] (fz.west -| pre.east) -- (fz.west);
  \draw[flow] (ft.west -| pre.east) -- (ft.west);
  \draw[flow] (fz.east) -- (fz.east -| task.west);
  \draw[flow] (ft.east) -- (ft.east -| task.west);
  \draw[flow] (sc.east) -- (sc.east -| task.west);
\end{scope}
\end{tikzpicture}
\caption{Design of the study.
(a)~The class hierarchy on the decays of a heavy resonance~$X$ and on the jets from quantum
chromodynamics (QCD); $q$ is a $u$ or $d$ quark, $\ell$ an electron or a muon, and $\tau_h$,
$\tau_e$, $\tau_\mu$ a $\tau$ lepton decaying to hadrons, an electron or a muon. Each box is one
class of the label set above its column, except the QCD box of the \Kfull-class column (Sophon's
label set~\cite{sophon}), which stands for \LitJetClassTwoNQcdClasses{} classes. Each level merges
classes of the level to its left; the lines below name what each merge removes.
(b)~The controls: the first random partition, where boxes with the same letter are one class, and
the random partitions with controlled flavor boundaries on the decays \RandFlavSplitOrbit{} (one
square per decay; filled squares are moved to a second class).
(c)~Models of one seed share the training jets, their order and the initial weights, except the
output layer, the mass regression target, and the first layer and initial weights of the model with
kinematic inputs only (the particle four-momenta). The hierarchy levels have \MlstNRunsWord{} seeds
each, the models with the auxiliary mass regression \MlstMassNRunsWord{}, the random partitions with
controlled flavor boundaries \RandFlavRuns{} and each random partition \RandVtwoRuns{}. Frozen
features, the output of the backbone (the network below the classification head) with its weights held
fixed, feed linear probes (logistic regressions) and multilayer perceptron (MLP) probes; fine-tuning
trains all weights. Colors identify the label sets and controls in every figure; the plots give the three partitions with
controlled flavor boundaries three shades; dotted borders mark the two
variants, which the other figures draw as dotted lines in the color of their label set.}
\label{fig:design}
\end{figure}
 from main.tex; this file holds one figure environment and nothing else.
%
% WHAT main.tex MUST LOAD (this file loads nothing itself):
%   \usepackage{tikz}
%   \usetikzlibrary{arrows.meta}          % the only TikZ library used (arrow tips)
%   \definecolor{c188}{HTML}{1F77B4}      % the configuration colours of experiments/FIGS/
%   \definecolor{c162}{HTML}{FF7F0E}      % mlst_style.py (matplotlib's tab10 in order, as in
%   \definecolor{c43}{HTML}{2CA02C}       % Sophon's figures); the variants are drawn in their
%   \definecolor{c17}{HTML}{D62728}       % label set's colour with a dotted border, as the
%   \definecolor{crand}{HTML}{17BECF}     % plotted figures draw variants dotted
%   \definecolor{cflav}{HTML}{9467BD}
%   % GENERATED by experiments/FIGS/make_tables.py -- do not edit.
% Every line carries: source file :: path inside that file :: first 16 hex of its
% sha256. paper/journal/provenance.json carries the full hash beside the value.
% Regenerate with:  python3 experiments/FIGS/make_tables.py
% Check in CI with: python3 experiments/FIGS/make_tables.py --check
%
% Macro names spell digits as words (162 -> Onesixtwo) because \newcommand takes
% letters only. A number whose source file does not exist yet is NOT emitted --
% see the missing list below, and leave a \pending{} in the text for it.
%
% MISSING INPUTS -- no macro was emitted for these:
%   community benchmarks, bench v3 readout -- experiments/FIGS/data/bench_v3_metrics*/bench_metrics_last.json
%   v2 pretraining rerun: data order and dropout fixed by seed and epoch, identical initial trunk weights, each epoch reading one of five disjoint random parts of the rows of nearly every file, fixed validation sample, checkpoint at the first maximum within the last ten epochs -- experiments/FIGS/data/v2/pretraining/
%   v2 frozen probes on the rerun models, at the first maximum within the last ten epochs and on the model whose weights are the average of the last ten epochs, BatchNorm statistics recomputed, with untrained-network and pooled-embedding reference rows -- experiments/FIGS/data/v2/probe_ladder/
%   v2 paired ratios with run and test-sample errors on the rerun models -- experiments/FIGS/data/v2/paired_errors/
%   v2 the 64- and 30-class vocabularies -- experiments/FIGS/data/v2/levels_64_30/
%   v2 the leave-one-family-out vocabularies -- experiments/FIGS/data/v2/leave_one_family_out/
%   v2 label-tree recovery on the rerun models -- experiments/FIGS/data/v2/label_recovery/
%   v2 the rerun's random partitions and flavour pair -- experiments/FIGS/data/v2/random_partitions/
%   v2 JetClass-II fine-tuning on files held out from pretraining -- experiments/FIGS/data/v2/finetune/
%   v2 fine-tuning from random initialisation and from the self-supervised model -- experiments/FIGS/data/v2/finetune_references/
%   v2 top tagging and quark/gluon tagging (Pythia and Herwig) fine-tuned from every rerun model -- experiments/FIGS/data/v2/benchmarks/
%   v2 the second and third self-supervised pretraining runs -- experiments/FIGS/data/v2/self_supervised/
%   v2 the mass-output models with the loss weight matched in loss share -- experiments/FIGS/data/v2/mass_lambda_matched/
%   v2 anomaly detection on the rerun models, per-run values -- experiments/FIGS/data/v2/anomaly/
%   v2 the open-data fits on the rerun models -- experiments/FIGS/data/v2/real_data/
%   MassLossShareVtwoOnesixtwoMass -- experiments/FIGS/data/v2/mass_lambda_matched/loss_share.json :: shares.162+mass
%   MassLossShareVtwoOnesevenMass -- experiments/FIGS/data/v2/mass_lambda_matched/loss_share.json :: shares.17+mass
%   MassLossShareVtwoOnesevenMassMatched -- experiments/FIGS/data/v2/mass_lambda_matched/loss_share.json :: shares.17+mass_matched
%   DesignVtwoSecondGpu, DesignVtwoSecondGpuRuns -- the grid's job specs put every run index on one product
%   per-arm parameters and MACs -- flops_by_arm.json (see docs/RECORD.md 2.1)
%   fine-tuning read-outs -- experiments/FIGS/data/v2/ft_mlst/: jc2, jc2_qcd
%   top-tagging recipe check -- ft_mlst/topcheck_bench_metrics.json
%   X->bb and X->bs fits -- experiments/FIGS/data/v2/xbx/
%   X->bs from scratch -- experiments/FIGS/data/v2/xbx_scratch/
%   JetClass accuracy bootstrap -- experiments/FIGS/data/v2/budget/jc1_acc_boot_N10000.json
%   MlstFtJcOneErrCoarseOverFullBootEFour -- waits on the JetClass bootstrap at 1e4 (after the equal-stopping read-out)
%   MlstLpBbQgZeroAccDropFull -- waits on pair probes with the quark/gluon inputs zeroed (open item O10, not run)
%   MlstLpBbQgZeroAccDropCoarse -- waits on pair probes with the quark/gluon inputs zeroed (open item O10, not run)
%   XbsRthirtyFullOverScratch -- waits on X->bs fits and from-scratch runs
%   XbsRthirtyFullOverScratchPr -- waits on X->bs fits and from-scratch runs
%   XbsRthirtyCoarseOverFull -- waits on X->bs fits
%   XbsRthirtyCoarseOverFullPr -- waits on X->bs fits
%   XbsRthirtyNodeOverTransferFull -- waits on X->bs fits
%   XbsRthirtyNodeOverTransferFullPr -- waits on X->bs fits
%   XbbRthirtyDirectOverTransferFull -- waits on X->bb fits
%   XbbRthirtyDirectOverTransferFullPr -- waits on X->bb fits
%   XbbRthirtyTransferCoarseOverFull -- waits on X->bb fits
%   XbbRthirtyTransferCoarseOverFullPr -- waits on X->bb fits
%
% SKIPPED TABLES:
%   T4 compute/parameters: no machine-readable source in the repository; docs/RECORD.md 2.1 names /data/results/mtx/flops/flops_by_arm.json, which is not committed

\newcommand{\PendingPretrain}{\pending{pending v2: pretraining rerun: data order and dropout fixed by seed and epoch, identical initial trunk weights, each epoch reading one of five disjoint random parts of the rows of nearly every file, fixed validation sample, checkpoint at the first maximum within the last ten epochs}}  % experiments/FIGS/make_tables.py :: V2_PENDING['Pretrain'] :: sha256 10e6ec08356a4a2b :: a marker, not a number: red in the draft until the input exists
\newcommand{\PendingProbes}{\pending{pending v2: frozen probes on the rerun models, at the first maximum within the last ten epochs and on the model whose weights are the average of the last ten epochs, BatchNorm statistics recomputed, with untrained-network and pooled-embedding reference rows}}  % experiments/FIGS/make_tables.py :: V2_PENDING['Probes'] :: sha256 10e6ec08356a4a2b :: a marker, not a number: red in the draft until the input exists
\newcommand{\PendingPaired}{\pending{pending v2: paired ratios with run and test-sample errors on the rerun models}}  % experiments/FIGS/make_tables.py :: V2_PENDING['Paired'] :: sha256 10e6ec08356a4a2b :: a marker, not a number: red in the draft until the input exists
\newcommand{\PendingLevels}{\pending{pending v2: the 64- and 30-class vocabularies}}  % experiments/FIGS/make_tables.py :: V2_PENDING['Levels'] :: sha256 10e6ec08356a4a2b :: a marker, not a number: red in the draft until the input exists
\newcommand{\PendingLofo}{\pending{pending v2: the leave-one-family-out vocabularies}}  % experiments/FIGS/make_tables.py :: V2_PENDING['Lofo'] :: sha256 10e6ec08356a4a2b :: a marker, not a number: red in the draft until the input exists
\newcommand{\PendingRecovery}{\pending{pending v2: label-tree recovery on the rerun models}}  % experiments/FIGS/make_tables.py :: V2_PENDING['Recovery'] :: sha256 10e6ec08356a4a2b :: a marker, not a number: red in the draft until the input exists
\newcommand{\PendingRandom}{\pending{pending v2: the rerun's random partitions and flavour pair}}  % experiments/FIGS/make_tables.py :: V2_PENDING['Random'] :: sha256 10e6ec08356a4a2b :: a marker, not a number: red in the draft until the input exists
\newcommand{\PendingFtHeldout}{\pending{pending v2: JetClass-II fine-tuning on files held out from pretraining}}  % experiments/FIGS/make_tables.py :: V2_PENDING['FtHeldout'] :: sha256 10e6ec08356a4a2b :: a marker, not a number: red in the draft until the input exists
\newcommand{\PendingFtRefs}{\pending{pending v2: fine-tuning from random initialisation and from the self-supervised model}}  % experiments/FIGS/make_tables.py :: V2_PENDING['FtRefs'] :: sha256 10e6ec08356a4a2b :: a marker, not a number: red in the draft until the input exists
\newcommand{\PendingBench}{\pending{pending v2: top tagging and quark/gluon tagging (Pythia and Herwig) fine-tuned from every rerun model}}  % experiments/FIGS/make_tables.py :: V2_PENDING['Bench'] :: sha256 10e6ec08356a4a2b :: a marker, not a number: red in the draft until the input exists
\newcommand{\PendingSsl}{\pending{pending v2: the second and third self-supervised pretraining runs}}  % experiments/FIGS/make_tables.py :: V2_PENDING['Ssl'] :: sha256 10e6ec08356a4a2b :: a marker, not a number: red in the draft until the input exists
\newcommand{\PendingMassLambda}{\pending{pending v2: the mass-output models with the loss weight matched in loss share}}  % experiments/FIGS/make_tables.py :: V2_PENDING['MassLambda'] :: sha256 10e6ec08356a4a2b :: a marker, not a number: red in the draft until the input exists
\newcommand{\PendingAnomaly}{\pending{pending v2: anomaly detection on the rerun models, per-run values}}  % experiments/FIGS/make_tables.py :: V2_PENDING['Anomaly'] :: sha256 10e6ec08356a4a2b :: a marker, not a number: red in the draft until the input exists
\newcommand{\PendingRealData}{\pending{pending v2: the open-data fits on the rerun models}}  % experiments/FIGS/make_tables.py :: V2_PENDING['RealData'] :: sha256 10e6ec08356a4a2b :: a marker, not a number: red in the draft until the input exists
\newcommand{\PendingConclusion}{\pending{pending v2: written once every input above is in}}  % experiments/FIGS/make_tables.py :: V2_PENDING['Conclusion'] :: sha256 10e6ec08356a4a2b :: a marker, not a number: red in the draft until the input exists
\newcommand{\AojChainDataset}{178}  % experiments/FIGS/data/aoj_checks_v2/selection_chain_v2.json :: n_dataset :: sha256 729d83085aba25f8
\newcommand{\AojChainStaged}{12{,}589{,}467}  % experiments/FIGS/data/aoj_checks_v2/selection_chain_v2.json :: n_staged :: sha256 729d83085aba25f8
\newcommand{\AojChainRhoWindow}{9{,}441{,}209}  % experiments/FIGS/data/aoj_checks_v2/selection_chain_v2.json :: n_rho_window :: sha256 729d83085aba25f8
\newcommand{\AojNJetsFit}{4{,}194{,}846}  % experiments/FIGS/data/aoj_checks_v2/selection_chain_v2.json :: n_fit :: sha256 729d83085aba25f8
\newcommand{\AojChainPtBins}{8}  % experiments/FIGS/data/aoj_checks_v2/selection_chain_v2.json :: n_pt_bins :: sha256 729d83085aba25f8
\newcommand{\AojChainEtaMax}{2.40}  % experiments/FIGS/data/aoj_checks_v2/selection_chain_v2.json :: abs_eta_max :: sha256 729d83085aba25f8
\newcommand{\MassLossShareVtwoOnesixtwoMass}{\pending{pending: realised mass-loss share in the rerun, 162 classes}}  % experiments/FIGS/make_tables.py :: SLOTS[MassLossShareVtwoOnesixtwoMass] :: sha256 10e6ec08356a4a2b :: a marker, not a number: experiments/FIGS/data/v2/mass_lambda_matched/loss_share.json :: shares.162+mass does not exist yet
\newcommand{\MassLossShareVtwoOnesevenMass}{\pending{pending: realised mass-loss share in the rerun, 17 classes}}  % experiments/FIGS/make_tables.py :: SLOTS[MassLossShareVtwoOnesevenMass] :: sha256 10e6ec08356a4a2b :: a marker, not a number: experiments/FIGS/data/v2/mass_lambda_matched/loss_share.json :: shares.17+mass does not exist yet
\newcommand{\MassLossShareVtwoOnesevenMassMatched}{\pending{pending: realised mass-loss share in the rerun, 17 classes, matched weight}}  % experiments/FIGS/make_tables.py :: SLOTS[MassLossShareVtwoOnesevenMassMatched] :: sha256 10e6ec08356a4a2b :: a marker, not a number: experiments/FIGS/data/v2/mass_lambda_matched/loss_share.json :: shares.17+mass_matched does not exist yet
\newcommand{\FtCoverageEThreeNClasses}{160}  % experiments/FT/data/jc2_v2_manifest.json :: class_coverage['train_N1000_s1.parquet'].n_classes_present :: sha256 4bd68655919b14e9
\newcommand{\FtCoverageEFourNClasses}{183}  % experiments/FT/data/jc2_v2_manifest.json :: class_coverage['train_N10000_s1.parquet'].n_classes_present :: sha256 4bd68655919b14e9
\newcommand{\FtCoverageEThreeMedian}{4}  % experiments/FT/data/jc2_v2_manifest.json :: class_coverage['train_N1000_s1.parquet'].median_per_present_class :: sha256 4bd68655919b14e9
\newcommand{\FtCoverageEThreeOneJet}{23}  % experiments/FT/data/jc2_v2_manifest.json :: class_coverage['train_N1000_s1.parquet'].n_classes_with_one_jet :: sha256 4bd68655919b14e9
\newcommand{\FtCoverageVoneEThreeNClasses}{161}  % experiments/FT/data/jc2_v1_coverage.json :: class_coverage['train_N1000_s1.parquet'].n_classes_present :: sha256 82652f57f708dc23
\newcommand{\FtCoverageVoneEThreeMedian}{5}  % experiments/FT/data/jc2_v1_coverage.json :: class_coverage['train_N1000_s1.parquet'].median_per_present_class :: sha256 82652f57f708dc23
\newcommand{\FtCoverageVoneEThreeOneJet}{25}  % experiments/FT/data/jc2_v1_coverage.json :: class_coverage['train_N1000_s1.parquet'].n_classes_with_one_jet :: sha256 82652f57f708dc23
\newcommand{\ProbeNJets}{2{,}000{,}000}  % experiments/FIGS/data/probe_ladder_v2_mlp2/analysis/seed_level_results.json :: provenance.n_jets_total :: sha256 53e1ec81343ac2c1
\newcommand{\ProbeNSeeds}{5}  % experiments/FIGS/data/probe_ladder_v2_mlp2/analysis/seed_level_results.json :: seeds_used (length) :: sha256 53e1ec81343ac2c1
\newcommand{\ProbeNSeedsWord}{five}  % experiments/FIGS/data/probe_ladder_v2_mlp2/analysis/seed_level_results.json :: seeds_used (length) :: sha256 53e1ec81343ac2c1 :: pretraining runs per vocabulary, as a word for the prose
\newcommand{\ProbeRowAlign}{\texttt{a9449fd420c473f6}}  % experiments/FIGS/data/probe_ladder_v2_mlp2/analysis/seed_level_results.json :: provenance.row_alignment_sha256 (first 16) :: sha256 53e1ec81343ac2c1
\newcommand{\ProbeEpsS}{90\%}  % experiments/FIGS/data/probe_ladder_v2_mlp2/analysis/seed_level_results.json :: levels.bvc_4prong.linear[0].headline_eps_s :: sha256 53e1ec81343ac2c1
\newcommand{\ProbeAucBvcFourprongLinearOneeighteight}{0.99320}  % experiments/FIGS/data/probe_ladder_v2_mlp2/analysis/seed_level_results.json :: levels.bvc_4prong.linear[0].mean_auc :: sha256 53e1ec81343ac2c1 :: 5 seeds, 0 saturated at AUC=1
\newcommand{\ProbeAucSdBvcFourprongLinearOneeighteight}{0.00027}  % experiments/FIGS/data/probe_ladder_v2_mlp2/analysis/seed_level_results.json :: table[bvc_4prong,linear,188].auc (SD over seeds, ddof=1) :: sha256 53e1ec81343ac2c1
\newcommand{\ProbeLogOneMinusAucBvcFourprongLinearOneeighteight}{\ensuremath{-}4.9915}  % experiments/FIGS/data/probe_ladder_v2_mlp2/analysis/seed_level_results.json :: levels.bvc_4prong.linear[0].mean :: sha256 53e1ec81343ac2c1 :: natural log of 1-AUC, lower is better
\newcommand{\ProbeLogOneMinusAucSdBvcFourprongLinearOneeighteight}{0.0389}  % experiments/FIGS/data/probe_ladder_v2_mlp2/analysis/seed_level_results.json :: levels.bvc_4prong.linear[0].seed_sd :: sha256 53e1ec81343ac2c1
\newcommand{\ProbeOmaBvcFourprongLinearOneeighteight}{\ensuremath{(6.80 \pm 0.27)\times10^{-3}}}  % experiments/FIGS/data/probe_ladder_v2_mlp2/analysis/seed_level_results.json :: table[bvc_4prong,linear,188].auc :: sha256 53e1ec81343ac2c1 :: 1-AUC, mean +- SD over 5 seeds
\newcommand{\ProbeBkgLeftBvcFourprongLinearOneeighteight}{14.0}  % experiments/FIGS/data/probe_ladder_v2_mlp2/analysis/seed_level_results.json :: levels.bvc_4prong.linear[0].rejection_points['0.90'].mean_n_bkg_pass :: sha256 53e1ec81343ac2c1 :: background jets passing the cut, mean over seeds
\newcommand{\ProbeRejBvcFourprongLinearOneeighteight}{227\,\ensuremath{\pm}\,23}  % experiments/FIGS/data/probe_ladder_v2_mlp2/analysis/seed_level_results.json :: table[bvc_4prong,linear,188].rejection_points['0.90'].rejection :: sha256 53e1ec81343ac2c1 :: at 90% signal efficiency, mean +- SD over seeds; 0 of 5 seeds at the cap
\newcommand{\ProbeAucBvcFourprongLinearOnesixtwo}{0.99311}  % experiments/FIGS/data/probe_ladder_v2_mlp2/analysis/seed_level_results.json :: levels.bvc_4prong.linear[1].mean_auc :: sha256 53e1ec81343ac2c1 :: 5 seeds, 0 saturated at AUC=1
\newcommand{\ProbeAucSdBvcFourprongLinearOnesixtwo}{0.00018}  % experiments/FIGS/data/probe_ladder_v2_mlp2/analysis/seed_level_results.json :: table[bvc_4prong,linear,162].auc (SD over seeds, ddof=1) :: sha256 53e1ec81343ac2c1
\newcommand{\ProbeLogOneMinusAucBvcFourprongLinearOnesixtwo}{\ensuremath{-}4.9777}  % experiments/FIGS/data/probe_ladder_v2_mlp2/analysis/seed_level_results.json :: levels.bvc_4prong.linear[1].mean :: sha256 53e1ec81343ac2c1 :: natural log of 1-AUC, lower is better
\newcommand{\ProbeLogOneMinusAucSdBvcFourprongLinearOnesixtwo}{0.0264}  % experiments/FIGS/data/probe_ladder_v2_mlp2/analysis/seed_level_results.json :: levels.bvc_4prong.linear[1].seed_sd :: sha256 53e1ec81343ac2c1
\newcommand{\ProbeOmaBvcFourprongLinearOnesixtwo}{\ensuremath{(6.89 \pm 0.18)\times10^{-3}}}  % experiments/FIGS/data/probe_ladder_v2_mlp2/analysis/seed_level_results.json :: table[bvc_4prong,linear,162].auc :: sha256 53e1ec81343ac2c1 :: 1-AUC, mean +- SD over 5 seeds
\newcommand{\ProbeBkgLeftBvcFourprongLinearOnesixtwo}{14.4}  % experiments/FIGS/data/probe_ladder_v2_mlp2/analysis/seed_level_results.json :: levels.bvc_4prong.linear[1].rejection_points['0.90'].mean_n_bkg_pass :: sha256 53e1ec81343ac2c1 :: background jets passing the cut, mean over seeds
\newcommand{\ProbeRejBvcFourprongLinearOnesixtwo}{220\,\ensuremath{\pm}\,40}  % experiments/FIGS/data/probe_ladder_v2_mlp2/analysis/seed_level_results.json :: table[bvc_4prong,linear,162].rejection_points['0.90'].rejection :: sha256 53e1ec81343ac2c1 :: at 90% signal efficiency, mean +- SD over seeds; 0 of 5 seeds at the cap
\newcommand{\ProbeAucBvcFourprongLinearFourthree}{0.99250}  % experiments/FIGS/data/probe_ladder_v2_mlp2/analysis/seed_level_results.json :: levels.bvc_4prong.linear[2].mean_auc :: sha256 53e1ec81343ac2c1 :: 5 seeds, 0 saturated at AUC=1
\newcommand{\ProbeAucSdBvcFourprongLinearFourthree}{0.00049}  % experiments/FIGS/data/probe_ladder_v2_mlp2/analysis/seed_level_results.json :: table[bvc_4prong,linear,43].auc (SD over seeds, ddof=1) :: sha256 53e1ec81343ac2c1
\newcommand{\ProbeLogOneMinusAucBvcFourprongLinearFourthree}{\ensuremath{-}4.8946}  % experiments/FIGS/data/probe_ladder_v2_mlp2/analysis/seed_level_results.json :: levels.bvc_4prong.linear[2].mean :: sha256 53e1ec81343ac2c1 :: natural log of 1-AUC, lower is better
\newcommand{\ProbeLogOneMinusAucSdBvcFourprongLinearFourthree}{0.0646}  % experiments/FIGS/data/probe_ladder_v2_mlp2/analysis/seed_level_results.json :: levels.bvc_4prong.linear[2].seed_sd :: sha256 53e1ec81343ac2c1
\newcommand{\ProbeOmaBvcFourprongLinearFourthree}{\ensuremath{(7.5 \pm 0.5)\times10^{-3}}}  % experiments/FIGS/data/probe_ladder_v2_mlp2/analysis/seed_level_results.json :: table[bvc_4prong,linear,43].auc :: sha256 53e1ec81343ac2c1 :: 1-AUC, mean +- SD over 5 seeds
\newcommand{\ProbeBkgLeftBvcFourprongLinearFourthree}{16.6}  % experiments/FIGS/data/probe_ladder_v2_mlp2/analysis/seed_level_results.json :: levels.bvc_4prong.linear[2].rejection_points['0.90'].mean_n_bkg_pass :: sha256 53e1ec81343ac2c1 :: background jets passing the cut, mean over seeds
\newcommand{\ProbeRejBvcFourprongLinearFourthree}{192\,\ensuremath{\pm}\,20}  % experiments/FIGS/data/probe_ladder_v2_mlp2/analysis/seed_level_results.json :: table[bvc_4prong,linear,43].rejection_points['0.90'].rejection :: sha256 53e1ec81343ac2c1 :: at 90% signal efficiency, mean +- SD over seeds; 0 of 5 seeds at the cap
\newcommand{\ProbeAucBvcFourprongLinearOneseven}{0.96936}  % experiments/FIGS/data/probe_ladder_v2_mlp2/analysis/seed_level_results.json :: levels.bvc_4prong.linear[3].mean_auc :: sha256 53e1ec81343ac2c1 :: 5 seeds, 0 saturated at AUC=1
\newcommand{\ProbeAucSdBvcFourprongLinearOneseven}{0.00514}  % experiments/FIGS/data/probe_ladder_v2_mlp2/analysis/seed_level_results.json :: table[bvc_4prong,linear,17].auc (SD over seeds, ddof=1) :: sha256 53e1ec81343ac2c1
\newcommand{\ProbeLogOneMinusAucBvcFourprongLinearOneseven}{\ensuremath{-}3.4967}  % experiments/FIGS/data/probe_ladder_v2_mlp2/analysis/seed_level_results.json :: levels.bvc_4prong.linear[3].mean :: sha256 53e1ec81343ac2c1 :: natural log of 1-AUC, lower is better
\newcommand{\ProbeLogOneMinusAucSdBvcFourprongLinearOneseven}{0.1667}  % experiments/FIGS/data/probe_ladder_v2_mlp2/analysis/seed_level_results.json :: levels.bvc_4prong.linear[3].seed_sd :: sha256 53e1ec81343ac2c1
\newcommand{\ProbeOmaBvcFourprongLinearOneseven}{\ensuremath{(3.1 \pm 0.5)\times10^{-2}}}  % experiments/FIGS/data/probe_ladder_v2_mlp2/analysis/seed_level_results.json :: table[bvc_4prong,linear,17].auc :: sha256 53e1ec81343ac2c1 :: 1-AUC, mean +- SD over 5 seeds
\newcommand{\ProbeBkgLeftBvcFourprongLinearOneseven}{137.4}  % experiments/FIGS/data/probe_ladder_v2_mlp2/analysis/seed_level_results.json :: levels.bvc_4prong.linear[3].rejection_points['0.90'].mean_n_bkg_pass :: sha256 53e1ec81343ac2c1 :: background jets passing the cut, mean over seeds
\newcommand{\ProbeRejBvcFourprongLinearOneseven}{24\,\ensuremath{\pm}\,6}  % experiments/FIGS/data/probe_ladder_v2_mlp2/analysis/seed_level_results.json :: table[bvc_4prong,linear,17].rejection_points['0.90'].rejection :: sha256 53e1ec81343ac2c1 :: at 90% signal efficiency, mean +- SD over seeds; 0 of 5 seeds at the cap
\newcommand{\ProbeOmaRatioBvcFourprongLinearOnesixtwoOverOneeighteight}{1.01}  % experiments/FIGS/data/probe_ladder_v2_mlp2/analysis/seed_level_results.json :: table[bvc_4prong,linear,162/188].auc :: sha256 53e1ec81343ac2c1 :: ratio of the seed means of 1-AUC, coarser over finer
\newcommand{\ProbeOmaRatioBvcFourprongLinearFourthreeOverOneeighteight}{1.10}  % experiments/FIGS/data/probe_ladder_v2_mlp2/analysis/seed_level_results.json :: table[bvc_4prong,linear,43/188].auc :: sha256 53e1ec81343ac2c1 :: ratio of the seed means of 1-AUC, coarser over finer
\newcommand{\ProbeOmaRatioBvcFourprongLinearOnesevenOverOneeighteight}{4.5}  % experiments/FIGS/data/probe_ladder_v2_mlp2/analysis/seed_level_results.json :: table[bvc_4prong,linear,17/188].auc :: sha256 53e1ec81343ac2c1 :: ratio of the seed means of 1-AUC, coarser over finer
\newcommand{\ProbeOmaRatioBvcFourprongLinearFourthreeOverOnesixtwo}{1.09}  % experiments/FIGS/data/probe_ladder_v2_mlp2/analysis/seed_level_results.json :: table[bvc_4prong,linear,43/162].auc :: sha256 53e1ec81343ac2c1 :: ratio of the seed means of 1-AUC, coarser over finer
\newcommand{\ProbeOmaRatioBvcFourprongLinearOnesevenOverOnesixtwo}{4.4}  % experiments/FIGS/data/probe_ladder_v2_mlp2/analysis/seed_level_results.json :: table[bvc_4prong,linear,17/162].auc :: sha256 53e1ec81343ac2c1 :: ratio of the seed means of 1-AUC, coarser over finer
\newcommand{\ProbeOmaRatioBvcFourprongLinearOnesevenOverFourthree}{4.1}  % experiments/FIGS/data/probe_ladder_v2_mlp2/analysis/seed_level_results.json :: table[bvc_4prong,linear,17/43].auc :: sha256 53e1ec81343ac2c1 :: ratio of the seed means of 1-AUC, coarser over finer
\newcommand{\ProbeAucBvcFourprongMlpOneeighteight}{0.99316}  % experiments/FIGS/data/probe_ladder_v2_mlp2/analysis/seed_level_results.json :: levels.bvc_4prong.mlp[0].mean_auc :: sha256 53e1ec81343ac2c1 :: 5 seeds, 0 saturated at AUC=1
\newcommand{\ProbeAucSdBvcFourprongMlpOneeighteight}{0.00020}  % experiments/FIGS/data/probe_ladder_v2_mlp2/analysis/seed_level_results.json :: table[bvc_4prong,mlp,188].auc (SD over seeds, ddof=1) :: sha256 53e1ec81343ac2c1
\newcommand{\ProbeLogOneMinusAucBvcFourprongMlpOneeighteight}{\ensuremath{-}4.9857}  % experiments/FIGS/data/probe_ladder_v2_mlp2/analysis/seed_level_results.json :: levels.bvc_4prong.mlp[0].mean :: sha256 53e1ec81343ac2c1 :: natural log of 1-AUC, lower is better
\newcommand{\ProbeLogOneMinusAucSdBvcFourprongMlpOneeighteight}{0.0290}  % experiments/FIGS/data/probe_ladder_v2_mlp2/analysis/seed_level_results.json :: levels.bvc_4prong.mlp[0].seed_sd :: sha256 53e1ec81343ac2c1
\newcommand{\ProbeOmaBvcFourprongMlpOneeighteight}{\ensuremath{(6.84 \pm 0.20)\times10^{-3}}}  % experiments/FIGS/data/probe_ladder_v2_mlp2/analysis/seed_level_results.json :: table[bvc_4prong,mlp,188].auc :: sha256 53e1ec81343ac2c1 :: 1-AUC, mean +- SD over 5 seeds
\newcommand{\ProbeBkgLeftBvcFourprongMlpOneeighteight}{15.2}  % experiments/FIGS/data/probe_ladder_v2_mlp2/analysis/seed_level_results.json :: levels.bvc_4prong.mlp[0].rejection_points['0.90'].mean_n_bkg_pass :: sha256 53e1ec81343ac2c1 :: background jets passing the cut, mean over seeds
\newcommand{\ProbeRejBvcFourprongMlpOneeighteight}{212\,\ensuremath{\pm}\,34}  % experiments/FIGS/data/probe_ladder_v2_mlp2/analysis/seed_level_results.json :: table[bvc_4prong,mlp,188].rejection_points['0.90'].rejection :: sha256 53e1ec81343ac2c1 :: at 90% signal efficiency, mean +- SD over seeds; 0 of 5 seeds at the cap
\newcommand{\ProbeAucBvcFourprongMlpOnesixtwo}{0.99318}  % experiments/FIGS/data/probe_ladder_v2_mlp2/analysis/seed_level_results.json :: levels.bvc_4prong.mlp[1].mean_auc :: sha256 53e1ec81343ac2c1 :: 5 seeds, 0 saturated at AUC=1
\newcommand{\ProbeAucSdBvcFourprongMlpOnesixtwo}{0.00020}  % experiments/FIGS/data/probe_ladder_v2_mlp2/analysis/seed_level_results.json :: table[bvc_4prong,mlp,162].auc (SD over seeds, ddof=1) :: sha256 53e1ec81343ac2c1
\newcommand{\ProbeLogOneMinusAucBvcFourprongMlpOnesixtwo}{\ensuremath{-}4.9886}  % experiments/FIGS/data/probe_ladder_v2_mlp2/analysis/seed_level_results.json :: levels.bvc_4prong.mlp[1].mean :: sha256 53e1ec81343ac2c1 :: natural log of 1-AUC, lower is better
\newcommand{\ProbeLogOneMinusAucSdBvcFourprongMlpOnesixtwo}{0.0298}  % experiments/FIGS/data/probe_ladder_v2_mlp2/analysis/seed_level_results.json :: levels.bvc_4prong.mlp[1].seed_sd :: sha256 53e1ec81343ac2c1
\newcommand{\ProbeOmaBvcFourprongMlpOnesixtwo}{\ensuremath{(6.82 \pm 0.20)\times10^{-3}}}  % experiments/FIGS/data/probe_ladder_v2_mlp2/analysis/seed_level_results.json :: table[bvc_4prong,mlp,162].auc :: sha256 53e1ec81343ac2c1 :: 1-AUC, mean +- SD over 5 seeds
\newcommand{\ProbeBkgLeftBvcFourprongMlpOnesixtwo}{16.6}  % experiments/FIGS/data/probe_ladder_v2_mlp2/analysis/seed_level_results.json :: levels.bvc_4prong.mlp[1].rejection_points['0.90'].mean_n_bkg_pass :: sha256 53e1ec81343ac2c1 :: background jets passing the cut, mean over seeds
\newcommand{\ProbeRejBvcFourprongMlpOnesixtwo}{200\,\ensuremath{\pm}\,40}  % experiments/FIGS/data/probe_ladder_v2_mlp2/analysis/seed_level_results.json :: table[bvc_4prong,mlp,162].rejection_points['0.90'].rejection :: sha256 53e1ec81343ac2c1 :: at 90% signal efficiency, mean +- SD over seeds; 0 of 5 seeds at the cap
\newcommand{\ProbeAucBvcFourprongMlpFourthree}{0.99232}  % experiments/FIGS/data/probe_ladder_v2_mlp2/analysis/seed_level_results.json :: levels.bvc_4prong.mlp[2].mean_auc :: sha256 53e1ec81343ac2c1 :: 5 seeds, 0 saturated at AUC=1
\newcommand{\ProbeAucSdBvcFourprongMlpFourthree}{0.00058}  % experiments/FIGS/data/probe_ladder_v2_mlp2/analysis/seed_level_results.json :: table[bvc_4prong,mlp,43].auc (SD over seeds, ddof=1) :: sha256 53e1ec81343ac2c1
\newcommand{\ProbeLogOneMinusAucBvcFourprongMlpFourthree}{\ensuremath{-}4.8719}  % experiments/FIGS/data/probe_ladder_v2_mlp2/analysis/seed_level_results.json :: levels.bvc_4prong.mlp[2].mean :: sha256 53e1ec81343ac2c1 :: natural log of 1-AUC, lower is better
\newcommand{\ProbeLogOneMinusAucSdBvcFourprongMlpFourthree}{0.0734}  % experiments/FIGS/data/probe_ladder_v2_mlp2/analysis/seed_level_results.json :: levels.bvc_4prong.mlp[2].seed_sd :: sha256 53e1ec81343ac2c1
\newcommand{\ProbeOmaBvcFourprongMlpFourthree}{\ensuremath{(7.7 \pm 0.6)\times10^{-3}}}  % experiments/FIGS/data/probe_ladder_v2_mlp2/analysis/seed_level_results.json :: table[bvc_4prong,mlp,43].auc :: sha256 53e1ec81343ac2c1 :: 1-AUC, mean +- SD over 5 seeds
\newcommand{\ProbeBkgLeftBvcFourprongMlpFourthree}{17.4}  % experiments/FIGS/data/probe_ladder_v2_mlp2/analysis/seed_level_results.json :: levels.bvc_4prong.mlp[2].rejection_points['0.90'].mean_n_bkg_pass :: sha256 53e1ec81343ac2c1 :: background jets passing the cut, mean over seeds
\newcommand{\ProbeRejBvcFourprongMlpFourthree}{184\,\ensuremath{\pm}\,23}  % experiments/FIGS/data/probe_ladder_v2_mlp2/analysis/seed_level_results.json :: table[bvc_4prong,mlp,43].rejection_points['0.90'].rejection :: sha256 53e1ec81343ac2c1 :: at 90% signal efficiency, mean +- SD over seeds; 0 of 5 seeds at the cap
\newcommand{\ProbeAucBvcFourprongMlpOneseven}{0.97766}  % experiments/FIGS/data/probe_ladder_v2_mlp2/analysis/seed_level_results.json :: levels.bvc_4prong.mlp[3].mean_auc :: sha256 53e1ec81343ac2c1 :: 5 seeds, 0 saturated at AUC=1
\newcommand{\ProbeAucSdBvcFourprongMlpOneseven}{0.00310}  % experiments/FIGS/data/probe_ladder_v2_mlp2/analysis/seed_level_results.json :: table[bvc_4prong,mlp,17].auc (SD over seeds, ddof=1) :: sha256 53e1ec81343ac2c1
\newcommand{\ProbeLogOneMinusAucBvcFourprongMlpOneseven}{\ensuremath{-}3.8087}  % experiments/FIGS/data/probe_ladder_v2_mlp2/analysis/seed_level_results.json :: levels.bvc_4prong.mlp[3].mean :: sha256 53e1ec81343ac2c1 :: natural log of 1-AUC, lower is better
\newcommand{\ProbeLogOneMinusAucSdBvcFourprongMlpOneseven}{0.1366}  % experiments/FIGS/data/probe_ladder_v2_mlp2/analysis/seed_level_results.json :: levels.bvc_4prong.mlp[3].seed_sd :: sha256 53e1ec81343ac2c1
\newcommand{\ProbeOmaBvcFourprongMlpOneseven}{\ensuremath{(2.23 \pm 0.31)\times10^{-2}}}  % experiments/FIGS/data/probe_ladder_v2_mlp2/analysis/seed_level_results.json :: table[bvc_4prong,mlp,17].auc :: sha256 53e1ec81343ac2c1 :: 1-AUC, mean +- SD over 5 seeds
\newcommand{\ProbeBkgLeftBvcFourprongMlpOneseven}{99.0}  % experiments/FIGS/data/probe_ladder_v2_mlp2/analysis/seed_level_results.json :: levels.bvc_4prong.mlp[3].rejection_points['0.90'].mean_n_bkg_pass :: sha256 53e1ec81343ac2c1 :: background jets passing the cut, mean over seeds
\newcommand{\ProbeRejBvcFourprongMlpOneseven}{33\,\ensuremath{\pm}\,5}  % experiments/FIGS/data/probe_ladder_v2_mlp2/analysis/seed_level_results.json :: table[bvc_4prong,mlp,17].rejection_points['0.90'].rejection :: sha256 53e1ec81343ac2c1 :: at 90% signal efficiency, mean +- SD over seeds; 0 of 5 seeds at the cap
\newcommand{\ProbeOmaRatioBvcFourprongMlpOnesixtwoOverOneeighteight}{1.00}  % experiments/FIGS/data/probe_ladder_v2_mlp2/analysis/seed_level_results.json :: table[bvc_4prong,mlp,162/188].auc :: sha256 53e1ec81343ac2c1 :: ratio of the seed means of 1-AUC, coarser over finer
\newcommand{\ProbeOmaRatioBvcFourprongMlpFourthreeOverOneeighteight}{1.12}  % experiments/FIGS/data/probe_ladder_v2_mlp2/analysis/seed_level_results.json :: table[bvc_4prong,mlp,43/188].auc :: sha256 53e1ec81343ac2c1 :: ratio of the seed means of 1-AUC, coarser over finer
\newcommand{\ProbeOmaRatioBvcFourprongMlpOnesevenOverOneeighteight}{3.3}  % experiments/FIGS/data/probe_ladder_v2_mlp2/analysis/seed_level_results.json :: table[bvc_4prong,mlp,17/188].auc :: sha256 53e1ec81343ac2c1 :: ratio of the seed means of 1-AUC, coarser over finer
\newcommand{\ProbeOmaRatioBvcFourprongMlpFourthreeOverOnesixtwo}{1.13}  % experiments/FIGS/data/probe_ladder_v2_mlp2/analysis/seed_level_results.json :: table[bvc_4prong,mlp,43/162].auc :: sha256 53e1ec81343ac2c1 :: ratio of the seed means of 1-AUC, coarser over finer
\newcommand{\ProbeOmaRatioBvcFourprongMlpOnesevenOverOnesixtwo}{3.3}  % experiments/FIGS/data/probe_ladder_v2_mlp2/analysis/seed_level_results.json :: table[bvc_4prong,mlp,17/162].auc :: sha256 53e1ec81343ac2c1 :: ratio of the seed means of 1-AUC, coarser over finer
\newcommand{\ProbeOmaRatioBvcFourprongMlpOnesevenOverFourthree}{2.9}  % experiments/FIGS/data/probe_ladder_v2_mlp2/analysis/seed_level_results.json :: table[bvc_4prong,mlp,17/43].auc :: sha256 53e1ec81343ac2c1 :: ratio of the seed means of 1-AUC, coarser over finer
\newcommand{\ProbeAucBvcQcdLinearOneeighteight}{0.98519}  % experiments/FIGS/data/probe_ladder_v2_mlp2/analysis/seed_level_results.json :: levels.bvc_qcd.linear[0].mean_auc :: sha256 53e1ec81343ac2c1 :: 5 seeds, 0 saturated at AUC=1
\newcommand{\ProbeAucSdBvcQcdLinearOneeighteight}{0.00061}  % experiments/FIGS/data/probe_ladder_v2_mlp2/analysis/seed_level_results.json :: table[bvc_qcd,linear,188].auc (SD over seeds, ddof=1) :: sha256 53e1ec81343ac2c1
\newcommand{\ProbeLogOneMinusAucBvcQcdLinearOneeighteight}{\ensuremath{-}4.2133}  % experiments/FIGS/data/probe_ladder_v2_mlp2/analysis/seed_level_results.json :: levels.bvc_qcd.linear[0].mean :: sha256 53e1ec81343ac2c1 :: natural log of 1-AUC, lower is better
\newcommand{\ProbeLogOneMinusAucSdBvcQcdLinearOneeighteight}{0.0422}  % experiments/FIGS/data/probe_ladder_v2_mlp2/analysis/seed_level_results.json :: levels.bvc_qcd.linear[0].seed_sd :: sha256 53e1ec81343ac2c1
\newcommand{\ProbeOmaBvcQcdLinearOneeighteight}{\ensuremath{(1.48 \pm 0.06)\times10^{-2}}}  % experiments/FIGS/data/probe_ladder_v2_mlp2/analysis/seed_level_results.json :: table[bvc_qcd,linear,188].auc :: sha256 53e1ec81343ac2c1 :: 1-AUC, mean +- SD over 5 seeds
\newcommand{\ProbeBkgLeftBvcQcdLinearOneeighteight}{96.4}  % experiments/FIGS/data/probe_ladder_v2_mlp2/analysis/seed_level_results.json :: levels.bvc_qcd.linear[0].rejection_points['0.90'].mean_n_bkg_pass :: sha256 53e1ec81343ac2c1 :: background jets passing the cut, mean over seeds
\newcommand{\ProbeRejBvcQcdLinearOneeighteight}{56\,\ensuremath{\pm}\,5}  % experiments/FIGS/data/probe_ladder_v2_mlp2/analysis/seed_level_results.json :: table[bvc_qcd,linear,188].rejection_points['0.90'].rejection :: sha256 53e1ec81343ac2c1 :: at 90% signal efficiency, mean +- SD over seeds; 0 of 5 seeds at the cap
\newcommand{\ProbeAucBvcQcdLinearOnesixtwo}{0.98305}  % experiments/FIGS/data/probe_ladder_v2_mlp2/analysis/seed_level_results.json :: levels.bvc_qcd.linear[1].mean_auc :: sha256 53e1ec81343ac2c1 :: 5 seeds, 0 saturated at AUC=1
\newcommand{\ProbeAucSdBvcQcdLinearOnesixtwo}{0.00024}  % experiments/FIGS/data/probe_ladder_v2_mlp2/analysis/seed_level_results.json :: table[bvc_qcd,linear,162].auc (SD over seeds, ddof=1) :: sha256 53e1ec81343ac2c1
\newcommand{\ProbeLogOneMinusAucBvcQcdLinearOnesixtwo}{\ensuremath{-}4.0774}  % experiments/FIGS/data/probe_ladder_v2_mlp2/analysis/seed_level_results.json :: levels.bvc_qcd.linear[1].mean :: sha256 53e1ec81343ac2c1 :: natural log of 1-AUC, lower is better
\newcommand{\ProbeLogOneMinusAucSdBvcQcdLinearOnesixtwo}{0.0142}  % experiments/FIGS/data/probe_ladder_v2_mlp2/analysis/seed_level_results.json :: levels.bvc_qcd.linear[1].seed_sd :: sha256 53e1ec81343ac2c1
\newcommand{\ProbeOmaBvcQcdLinearOnesixtwo}{\ensuremath{(1.695 \pm 0.024)\times10^{-2}}}  % experiments/FIGS/data/probe_ladder_v2_mlp2/analysis/seed_level_results.json :: table[bvc_qcd,linear,162].auc :: sha256 53e1ec81343ac2c1 :: 1-AUC, mean +- SD over 5 seeds
\newcommand{\ProbeBkgLeftBvcQcdLinearOnesixtwo}{121.4}  % experiments/FIGS/data/probe_ladder_v2_mlp2/analysis/seed_level_results.json :: levels.bvc_qcd.linear[1].rejection_points['0.90'].mean_n_bkg_pass :: sha256 53e1ec81343ac2c1 :: background jets passing the cut, mean over seeds
\newcommand{\ProbeRejBvcQcdLinearOnesixtwo}{44.3\,\ensuremath{\pm}\,2.1}  % experiments/FIGS/data/probe_ladder_v2_mlp2/analysis/seed_level_results.json :: table[bvc_qcd,linear,162].rejection_points['0.90'].rejection :: sha256 53e1ec81343ac2c1 :: at 90% signal efficiency, mean +- SD over seeds; 0 of 5 seeds at the cap
\newcommand{\ProbeAucBvcQcdLinearFourthree}{0.98281}  % experiments/FIGS/data/probe_ladder_v2_mlp2/analysis/seed_level_results.json :: levels.bvc_qcd.linear[2].mean_auc :: sha256 53e1ec81343ac2c1 :: 5 seeds, 0 saturated at AUC=1
\newcommand{\ProbeAucSdBvcQcdLinearFourthree}{0.00064}  % experiments/FIGS/data/probe_ladder_v2_mlp2/analysis/seed_level_results.json :: table[bvc_qcd,linear,43].auc (SD over seeds, ddof=1) :: sha256 53e1ec81343ac2c1
\newcommand{\ProbeLogOneMinusAucBvcQcdLinearFourthree}{\ensuremath{-}4.0638}  % experiments/FIGS/data/probe_ladder_v2_mlp2/analysis/seed_level_results.json :: levels.bvc_qcd.linear[2].mean :: sha256 53e1ec81343ac2c1 :: natural log of 1-AUC, lower is better
\newcommand{\ProbeLogOneMinusAucSdBvcQcdLinearFourthree}{0.0370}  % experiments/FIGS/data/probe_ladder_v2_mlp2/analysis/seed_level_results.json :: levels.bvc_qcd.linear[2].seed_sd :: sha256 53e1ec81343ac2c1
\newcommand{\ProbeOmaBvcQcdLinearFourthree}{\ensuremath{(1.72 \pm 0.06)\times10^{-2}}}  % experiments/FIGS/data/probe_ladder_v2_mlp2/analysis/seed_level_results.json :: table[bvc_qcd,linear,43].auc :: sha256 53e1ec81343ac2c1 :: 1-AUC, mean +- SD over 5 seeds
\newcommand{\ProbeBkgLeftBvcQcdLinearFourthree}{126.8}  % experiments/FIGS/data/probe_ladder_v2_mlp2/analysis/seed_level_results.json :: levels.bvc_qcd.linear[2].rejection_points['0.90'].mean_n_bkg_pass :: sha256 53e1ec81343ac2c1 :: background jets passing the cut, mean over seeds
\newcommand{\ProbeRejBvcQcdLinearFourthree}{43\,\ensuremath{\pm}\,4}  % experiments/FIGS/data/probe_ladder_v2_mlp2/analysis/seed_level_results.json :: table[bvc_qcd,linear,43].rejection_points['0.90'].rejection :: sha256 53e1ec81343ac2c1 :: at 90% signal efficiency, mean +- SD over seeds; 0 of 5 seeds at the cap
\newcommand{\ProbeAucBvcQcdLinearOneseven}{0.95504}  % experiments/FIGS/data/probe_ladder_v2_mlp2/analysis/seed_level_results.json :: levels.bvc_qcd.linear[3].mean_auc :: sha256 53e1ec81343ac2c1 :: 5 seeds, 0 saturated at AUC=1
\newcommand{\ProbeAucSdBvcQcdLinearOneseven}{0.00527}  % experiments/FIGS/data/probe_ladder_v2_mlp2/analysis/seed_level_results.json :: table[bvc_qcd,linear,17].auc (SD over seeds, ddof=1) :: sha256 53e1ec81343ac2c1
\newcommand{\ProbeLogOneMinusAucBvcQcdLinearOneseven}{\ensuremath{-}3.1076}  % experiments/FIGS/data/probe_ladder_v2_mlp2/analysis/seed_level_results.json :: levels.bvc_qcd.linear[3].mean :: sha256 53e1ec81343ac2c1 :: natural log of 1-AUC, lower is better
\newcommand{\ProbeLogOneMinusAucSdBvcQcdLinearOneseven}{0.1198}  % experiments/FIGS/data/probe_ladder_v2_mlp2/analysis/seed_level_results.json :: levels.bvc_qcd.linear[3].seed_sd :: sha256 53e1ec81343ac2c1
\newcommand{\ProbeOmaBvcQcdLinearOneseven}{\ensuremath{(4.5 \pm 0.5)\times10^{-2}}}  % experiments/FIGS/data/probe_ladder_v2_mlp2/analysis/seed_level_results.json :: table[bvc_qcd,linear,17].auc :: sha256 53e1ec81343ac2c1 :: 1-AUC, mean +- SD over 5 seeds
\newcommand{\ProbeBkgLeftBvcQcdLinearOneseven}{544.6}  % experiments/FIGS/data/probe_ladder_v2_mlp2/analysis/seed_level_results.json :: levels.bvc_qcd.linear[3].rejection_points['0.90'].mean_n_bkg_pass :: sha256 53e1ec81343ac2c1 :: background jets passing the cut, mean over seeds
\newcommand{\ProbeRejBvcQcdLinearOneseven}{10.1\,\ensuremath{\pm}\,1.9}  % experiments/FIGS/data/probe_ladder_v2_mlp2/analysis/seed_level_results.json :: table[bvc_qcd,linear,17].rejection_points['0.90'].rejection :: sha256 53e1ec81343ac2c1 :: at 90% signal efficiency, mean +- SD over seeds; 0 of 5 seeds at the cap
\newcommand{\ProbeOmaRatioBvcQcdLinearOnesixtwoOverOneeighteight}{1.14}  % experiments/FIGS/data/probe_ladder_v2_mlp2/analysis/seed_level_results.json :: table[bvc_qcd,linear,162/188].auc :: sha256 53e1ec81343ac2c1 :: ratio of the seed means of 1-AUC, coarser over finer
\newcommand{\ProbeOmaRatioBvcQcdLinearFourthreeOverOneeighteight}{1.16}  % experiments/FIGS/data/probe_ladder_v2_mlp2/analysis/seed_level_results.json :: table[bvc_qcd,linear,43/188].auc :: sha256 53e1ec81343ac2c1 :: ratio of the seed means of 1-AUC, coarser over finer
\newcommand{\ProbeOmaRatioBvcQcdLinearOnesevenOverOneeighteight}{3.0}  % experiments/FIGS/data/probe_ladder_v2_mlp2/analysis/seed_level_results.json :: table[bvc_qcd,linear,17/188].auc :: sha256 53e1ec81343ac2c1 :: ratio of the seed means of 1-AUC, coarser over finer
\newcommand{\ProbeOmaRatioBvcQcdLinearFourthreeOverOnesixtwo}{1.01}  % experiments/FIGS/data/probe_ladder_v2_mlp2/analysis/seed_level_results.json :: table[bvc_qcd,linear,43/162].auc :: sha256 53e1ec81343ac2c1 :: ratio of the seed means of 1-AUC, coarser over finer
\newcommand{\ProbeOmaRatioBvcQcdLinearOnesevenOverOnesixtwo}{2.7}  % experiments/FIGS/data/probe_ladder_v2_mlp2/analysis/seed_level_results.json :: table[bvc_qcd,linear,17/162].auc :: sha256 53e1ec81343ac2c1 :: ratio of the seed means of 1-AUC, coarser over finer
\newcommand{\ProbeOmaRatioBvcQcdLinearOnesevenOverFourthree}{2.6}  % experiments/FIGS/data/probe_ladder_v2_mlp2/analysis/seed_level_results.json :: table[bvc_qcd,linear,17/43].auc :: sha256 53e1ec81343ac2c1 :: ratio of the seed means of 1-AUC, coarser over finer
\newcommand{\ProbeAucBvcQcdMlpOneeighteight}{0.98527}  % experiments/FIGS/data/probe_ladder_v2_mlp2/analysis/seed_level_results.json :: levels.bvc_qcd.mlp[0].mean_auc :: sha256 53e1ec81343ac2c1 :: 5 seeds, 0 saturated at AUC=1
\newcommand{\ProbeAucSdBvcQcdMlpOneeighteight}{0.00058}  % experiments/FIGS/data/probe_ladder_v2_mlp2/analysis/seed_level_results.json :: table[bvc_qcd,mlp,188].auc (SD over seeds, ddof=1) :: sha256 53e1ec81343ac2c1
\newcommand{\ProbeLogOneMinusAucBvcQcdMlpOneeighteight}{\ensuremath{-}4.2188}  % experiments/FIGS/data/probe_ladder_v2_mlp2/analysis/seed_level_results.json :: levels.bvc_qcd.mlp[0].mean :: sha256 53e1ec81343ac2c1 :: natural log of 1-AUC, lower is better
\newcommand{\ProbeLogOneMinusAucSdBvcQcdMlpOneeighteight}{0.0400}  % experiments/FIGS/data/probe_ladder_v2_mlp2/analysis/seed_level_results.json :: levels.bvc_qcd.mlp[0].seed_sd :: sha256 53e1ec81343ac2c1
\newcommand{\ProbeOmaBvcQcdMlpOneeighteight}{\ensuremath{(1.47 \pm 0.06)\times10^{-2}}}  % experiments/FIGS/data/probe_ladder_v2_mlp2/analysis/seed_level_results.json :: table[bvc_qcd,mlp,188].auc :: sha256 53e1ec81343ac2c1 :: 1-AUC, mean +- SD over 5 seeds
\newcommand{\ProbeBkgLeftBvcQcdMlpOneeighteight}{94.6}  % experiments/FIGS/data/probe_ladder_v2_mlp2/analysis/seed_level_results.json :: levels.bvc_qcd.mlp[0].rejection_points['0.90'].mean_n_bkg_pass :: sha256 53e1ec81343ac2c1 :: background jets passing the cut, mean over seeds
\newcommand{\ProbeRejBvcQcdMlpOneeighteight}{57\,\ensuremath{\pm}\,6}  % experiments/FIGS/data/probe_ladder_v2_mlp2/analysis/seed_level_results.json :: table[bvc_qcd,mlp,188].rejection_points['0.90'].rejection :: sha256 53e1ec81343ac2c1 :: at 90% signal efficiency, mean +- SD over seeds; 0 of 5 seeds at the cap
\newcommand{\ProbeAucBvcQcdMlpOnesixtwo}{0.98361}  % experiments/FIGS/data/probe_ladder_v2_mlp2/analysis/seed_level_results.json :: levels.bvc_qcd.mlp[1].mean_auc :: sha256 53e1ec81343ac2c1 :: 5 seeds, 0 saturated at AUC=1
\newcommand{\ProbeAucSdBvcQcdMlpOnesixtwo}{0.00039}  % experiments/FIGS/data/probe_ladder_v2_mlp2/analysis/seed_level_results.json :: table[bvc_qcd,mlp,162].auc (SD over seeds, ddof=1) :: sha256 53e1ec81343ac2c1
\newcommand{\ProbeLogOneMinusAucBvcQcdMlpOnesixtwo}{\ensuremath{-}4.1113}  % experiments/FIGS/data/probe_ladder_v2_mlp2/analysis/seed_level_results.json :: levels.bvc_qcd.mlp[1].mean :: sha256 53e1ec81343ac2c1 :: natural log of 1-AUC, lower is better
\newcommand{\ProbeLogOneMinusAucSdBvcQcdMlpOnesixtwo}{0.0241}  % experiments/FIGS/data/probe_ladder_v2_mlp2/analysis/seed_level_results.json :: levels.bvc_qcd.mlp[1].seed_sd :: sha256 53e1ec81343ac2c1
\newcommand{\ProbeOmaBvcQcdMlpOnesixtwo}{\ensuremath{(1.64 \pm 0.04)\times10^{-2}}}  % experiments/FIGS/data/probe_ladder_v2_mlp2/analysis/seed_level_results.json :: table[bvc_qcd,mlp,162].auc :: sha256 53e1ec81343ac2c1 :: 1-AUC, mean +- SD over 5 seeds
\newcommand{\ProbeBkgLeftBvcQcdMlpOnesixtwo}{119.6}  % experiments/FIGS/data/probe_ladder_v2_mlp2/analysis/seed_level_results.json :: levels.bvc_qcd.mlp[1].rejection_points['0.90'].mean_n_bkg_pass :: sha256 53e1ec81343ac2c1 :: background jets passing the cut, mean over seeds
\newcommand{\ProbeRejBvcQcdMlpOnesixtwo}{45.1\,\ensuremath{\pm}\,3.2}  % experiments/FIGS/data/probe_ladder_v2_mlp2/analysis/seed_level_results.json :: table[bvc_qcd,mlp,162].rejection_points['0.90'].rejection :: sha256 53e1ec81343ac2c1 :: at 90% signal efficiency, mean +- SD over seeds; 0 of 5 seeds at the cap
\newcommand{\ProbeAucBvcQcdMlpFourthree}{0.98341}  % experiments/FIGS/data/probe_ladder_v2_mlp2/analysis/seed_level_results.json :: levels.bvc_qcd.mlp[2].mean_auc :: sha256 53e1ec81343ac2c1 :: 5 seeds, 0 saturated at AUC=1
\newcommand{\ProbeAucSdBvcQcdMlpFourthree}{0.00052}  % experiments/FIGS/data/probe_ladder_v2_mlp2/analysis/seed_level_results.json :: table[bvc_qcd,mlp,43].auc (SD over seeds, ddof=1) :: sha256 53e1ec81343ac2c1
\newcommand{\ProbeLogOneMinusAucBvcQcdMlpFourthree}{\ensuremath{-}4.0995}  % experiments/FIGS/data/probe_ladder_v2_mlp2/analysis/seed_level_results.json :: levels.bvc_qcd.mlp[2].mean :: sha256 53e1ec81343ac2c1 :: natural log of 1-AUC, lower is better
\newcommand{\ProbeLogOneMinusAucSdBvcQcdMlpFourthree}{0.0318}  % experiments/FIGS/data/probe_ladder_v2_mlp2/analysis/seed_level_results.json :: levels.bvc_qcd.mlp[2].seed_sd :: sha256 53e1ec81343ac2c1
\newcommand{\ProbeOmaBvcQcdMlpFourthree}{\ensuremath{(1.66 \pm 0.05)\times10^{-2}}}  % experiments/FIGS/data/probe_ladder_v2_mlp2/analysis/seed_level_results.json :: table[bvc_qcd,mlp,43].auc :: sha256 53e1ec81343ac2c1 :: 1-AUC, mean +- SD over 5 seeds
\newcommand{\ProbeBkgLeftBvcQcdMlpFourthree}{119.2}  % experiments/FIGS/data/probe_ladder_v2_mlp2/analysis/seed_level_results.json :: levels.bvc_qcd.mlp[2].rejection_points['0.90'].mean_n_bkg_pass :: sha256 53e1ec81343ac2c1 :: background jets passing the cut, mean over seeds
\newcommand{\ProbeRejBvcQcdMlpFourthree}{45.1\,\ensuremath{\pm}\,1.3}  % experiments/FIGS/data/probe_ladder_v2_mlp2/analysis/seed_level_results.json :: table[bvc_qcd,mlp,43].rejection_points['0.90'].rejection :: sha256 53e1ec81343ac2c1 :: at 90% signal efficiency, mean +- SD over seeds; 0 of 5 seeds at the cap
\newcommand{\ProbeAucBvcQcdMlpOneseven}{0.96226}  % experiments/FIGS/data/probe_ladder_v2_mlp2/analysis/seed_level_results.json :: levels.bvc_qcd.mlp[3].mean_auc :: sha256 53e1ec81343ac2c1 :: 5 seeds, 0 saturated at AUC=1
\newcommand{\ProbeAucSdBvcQcdMlpOneseven}{0.00373}  % experiments/FIGS/data/probe_ladder_v2_mlp2/analysis/seed_level_results.json :: table[bvc_qcd,mlp,17].auc (SD over seeds, ddof=1) :: sha256 53e1ec81343ac2c1
\newcommand{\ProbeLogOneMinusAucBvcQcdMlpOneseven}{\ensuremath{-}3.2811}  % experiments/FIGS/data/probe_ladder_v2_mlp2/analysis/seed_level_results.json :: levels.bvc_qcd.mlp[3].mean :: sha256 53e1ec81343ac2c1 :: natural log of 1-AUC, lower is better
\newcommand{\ProbeLogOneMinusAucSdBvcQcdMlpOneseven}{0.1010}  % experiments/FIGS/data/probe_ladder_v2_mlp2/analysis/seed_level_results.json :: levels.bvc_qcd.mlp[3].seed_sd :: sha256 53e1ec81343ac2c1
\newcommand{\ProbeOmaBvcQcdMlpOneseven}{\ensuremath{(3.8 \pm 0.4)\times10^{-2}}}  % experiments/FIGS/data/probe_ladder_v2_mlp2/analysis/seed_level_results.json :: table[bvc_qcd,mlp,17].auc :: sha256 53e1ec81343ac2c1 :: 1-AUC, mean +- SD over 5 seeds
\newcommand{\ProbeBkgLeftBvcQcdMlpOneseven}{458.8}  % experiments/FIGS/data/probe_ladder_v2_mlp2/analysis/seed_level_results.json :: levels.bvc_qcd.mlp[3].rejection_points['0.90'].mean_n_bkg_pass :: sha256 53e1ec81343ac2c1 :: background jets passing the cut, mean over seeds
\newcommand{\ProbeRejBvcQcdMlpOneseven}{11.8\,\ensuremath{\pm}\,1.5}  % experiments/FIGS/data/probe_ladder_v2_mlp2/analysis/seed_level_results.json :: table[bvc_qcd,mlp,17].rejection_points['0.90'].rejection :: sha256 53e1ec81343ac2c1 :: at 90% signal efficiency, mean +- SD over seeds; 0 of 5 seeds at the cap
\newcommand{\ProbeOmaRatioBvcQcdMlpOnesixtwoOverOneeighteight}{1.11}  % experiments/FIGS/data/probe_ladder_v2_mlp2/analysis/seed_level_results.json :: table[bvc_qcd,mlp,162/188].auc :: sha256 53e1ec81343ac2c1 :: ratio of the seed means of 1-AUC, coarser over finer
\newcommand{\ProbeOmaRatioBvcQcdMlpFourthreeOverOneeighteight}{1.13}  % experiments/FIGS/data/probe_ladder_v2_mlp2/analysis/seed_level_results.json :: table[bvc_qcd,mlp,43/188].auc :: sha256 53e1ec81343ac2c1 :: ratio of the seed means of 1-AUC, coarser over finer
\newcommand{\ProbeOmaRatioBvcQcdMlpOnesevenOverOneeighteight}{2.6}  % experiments/FIGS/data/probe_ladder_v2_mlp2/analysis/seed_level_results.json :: table[bvc_qcd,mlp,17/188].auc :: sha256 53e1ec81343ac2c1 :: ratio of the seed means of 1-AUC, coarser over finer
\newcommand{\ProbeOmaRatioBvcQcdMlpFourthreeOverOnesixtwo}{1.01}  % experiments/FIGS/data/probe_ladder_v2_mlp2/analysis/seed_level_results.json :: table[bvc_qcd,mlp,43/162].auc :: sha256 53e1ec81343ac2c1 :: ratio of the seed means of 1-AUC, coarser over finer
\newcommand{\ProbeOmaRatioBvcQcdMlpOnesevenOverOnesixtwo}{2.3}  % experiments/FIGS/data/probe_ladder_v2_mlp2/analysis/seed_level_results.json :: table[bvc_qcd,mlp,17/162].auc :: sha256 53e1ec81343ac2c1 :: ratio of the seed means of 1-AUC, coarser over finer
\newcommand{\ProbeOmaRatioBvcQcdMlpOnesevenOverFourthree}{2.3}  % experiments/FIGS/data/probe_ladder_v2_mlp2/analysis/seed_level_results.json :: table[bvc_qcd,mlp,17/43].auc :: sha256 53e1ec81343ac2c1 :: ratio of the seed means of 1-AUC, coarser over finer
\newcommand{\ProbeAucBvcResonantLinearOneeighteight}{0.99768}  % experiments/FIGS/data/probe_ladder_v2_mlp2/analysis/seed_level_results.json :: levels.bvc_resonant.linear[0].mean_auc :: sha256 53e1ec81343ac2c1 :: 5 seeds, 0 saturated at AUC=1
\newcommand{\ProbeAucSdBvcResonantLinearOneeighteight}{0.00007}  % experiments/FIGS/data/probe_ladder_v2_mlp2/analysis/seed_level_results.json :: table[bvc_resonant,linear,188].auc (SD over seeds, ddof=1) :: sha256 53e1ec81343ac2c1
\newcommand{\ProbeLogOneMinusAucBvcResonantLinearOneeighteight}{\ensuremath{-}6.0683}  % experiments/FIGS/data/probe_ladder_v2_mlp2/analysis/seed_level_results.json :: levels.bvc_resonant.linear[0].mean :: sha256 53e1ec81343ac2c1 :: natural log of 1-AUC, lower is better
\newcommand{\ProbeLogOneMinusAucSdBvcResonantLinearOneeighteight}{0.0317}  % experiments/FIGS/data/probe_ladder_v2_mlp2/analysis/seed_level_results.json :: levels.bvc_resonant.linear[0].seed_sd :: sha256 53e1ec81343ac2c1
\newcommand{\ProbeOmaBvcResonantLinearOneeighteight}{\ensuremath{(2.32 \pm 0.07)\times10^{-3}}}  % experiments/FIGS/data/probe_ladder_v2_mlp2/analysis/seed_level_results.json :: table[bvc_resonant,linear,188].auc :: sha256 53e1ec81343ac2c1 :: 1-AUC, mean +- SD over 5 seeds
\newcommand{\ProbeBkgLeftBvcResonantLinearOneeighteight}{9.0}  % experiments/FIGS/data/probe_ladder_v2_mlp2/analysis/seed_level_results.json :: levels.bvc_resonant.linear[0].rejection_points['0.90'].mean_n_bkg_pass :: sha256 53e1ec81343ac2c1 :: background jets passing the cut, mean over seeds
\newcommand{\ProbeRejBvcResonantLinearOneeighteight}{1{,}380\,\ensuremath{\pm}\,350}  % experiments/FIGS/data/probe_ladder_v2_mlp2/analysis/seed_level_results.json :: table[bvc_resonant,linear,188].rejection_points['0.90'].rejection :: sha256 53e1ec81343ac2c1 :: at 90% signal efficiency, mean +- SD over seeds; 0 of 5 seeds at the cap
\newcommand{\ProbeAucBvcResonantLinearOnesixtwo}{0.99764}  % experiments/FIGS/data/probe_ladder_v2_mlp2/analysis/seed_level_results.json :: levels.bvc_resonant.linear[1].mean_auc :: sha256 53e1ec81343ac2c1 :: 5 seeds, 0 saturated at AUC=1
\newcommand{\ProbeAucSdBvcResonantLinearOnesixtwo}{0.00006}  % experiments/FIGS/data/probe_ladder_v2_mlp2/analysis/seed_level_results.json :: table[bvc_resonant,linear,162].auc (SD over seeds, ddof=1) :: sha256 53e1ec81343ac2c1
\newcommand{\ProbeLogOneMinusAucBvcResonantLinearOnesixtwo}{\ensuremath{-}6.0493}  % experiments/FIGS/data/probe_ladder_v2_mlp2/analysis/seed_level_results.json :: levels.bvc_resonant.linear[1].mean :: sha256 53e1ec81343ac2c1 :: natural log of 1-AUC, lower is better
\newcommand{\ProbeLogOneMinusAucSdBvcResonantLinearOnesixtwo}{0.0261}  % experiments/FIGS/data/probe_ladder_v2_mlp2/analysis/seed_level_results.json :: levels.bvc_resonant.linear[1].seed_sd :: sha256 53e1ec81343ac2c1
\newcommand{\ProbeOmaBvcResonantLinearOnesixtwo}{\ensuremath{(2.36 \pm 0.06)\times10^{-3}}}  % experiments/FIGS/data/probe_ladder_v2_mlp2/analysis/seed_level_results.json :: table[bvc_resonant,linear,162].auc :: sha256 53e1ec81343ac2c1 :: 1-AUC, mean +- SD over 5 seeds
\newcommand{\ProbeBkgLeftBvcResonantLinearOnesixtwo}{9.4}  % experiments/FIGS/data/probe_ladder_v2_mlp2/analysis/seed_level_results.json :: levels.bvc_resonant.linear[1].rejection_points['0.90'].mean_n_bkg_pass :: sha256 53e1ec81343ac2c1 :: background jets passing the cut, mean over seeds
\newcommand{\ProbeRejBvcResonantLinearOnesixtwo}{1{,}290\,\ensuremath{\pm}\,240}  % experiments/FIGS/data/probe_ladder_v2_mlp2/analysis/seed_level_results.json :: table[bvc_resonant,linear,162].rejection_points['0.90'].rejection :: sha256 53e1ec81343ac2c1 :: at 90% signal efficiency, mean +- SD over seeds; 0 of 5 seeds at the cap
\newcommand{\ProbeAucBvcResonantLinearFourthree}{0.99748}  % experiments/FIGS/data/probe_ladder_v2_mlp2/analysis/seed_level_results.json :: levels.bvc_resonant.linear[2].mean_auc :: sha256 53e1ec81343ac2c1 :: 5 seeds, 0 saturated at AUC=1
\newcommand{\ProbeAucSdBvcResonantLinearFourthree}{0.00007}  % experiments/FIGS/data/probe_ladder_v2_mlp2/analysis/seed_level_results.json :: table[bvc_resonant,linear,43].auc (SD over seeds, ddof=1) :: sha256 53e1ec81343ac2c1
\newcommand{\ProbeLogOneMinusAucBvcResonantLinearFourthree}{\ensuremath{-}5.9834}  % experiments/FIGS/data/probe_ladder_v2_mlp2/analysis/seed_level_results.json :: levels.bvc_resonant.linear[2].mean :: sha256 53e1ec81343ac2c1 :: natural log of 1-AUC, lower is better
\newcommand{\ProbeLogOneMinusAucSdBvcResonantLinearFourthree}{0.0264}  % experiments/FIGS/data/probe_ladder_v2_mlp2/analysis/seed_level_results.json :: levels.bvc_resonant.linear[2].seed_sd :: sha256 53e1ec81343ac2c1
\newcommand{\ProbeOmaBvcResonantLinearFourthree}{\ensuremath{(2.52 \pm 0.07)\times10^{-3}}}  % experiments/FIGS/data/probe_ladder_v2_mlp2/analysis/seed_level_results.json :: table[bvc_resonant,linear,43].auc :: sha256 53e1ec81343ac2c1 :: 1-AUC, mean +- SD over 5 seeds
\newcommand{\ProbeBkgLeftBvcResonantLinearFourthree}{11.0}  % experiments/FIGS/data/probe_ladder_v2_mlp2/analysis/seed_level_results.json :: levels.bvc_resonant.linear[2].rejection_points['0.90'].mean_n_bkg_pass :: sha256 53e1ec81343ac2c1 :: background jets passing the cut, mean over seeds
\newcommand{\ProbeRejBvcResonantLinearFourthree}{1{,}130\,\ensuremath{\pm}\,260}  % experiments/FIGS/data/probe_ladder_v2_mlp2/analysis/seed_level_results.json :: table[bvc_resonant,linear,43].rejection_points['0.90'].rejection :: sha256 53e1ec81343ac2c1 :: at 90% signal efficiency, mean +- SD over seeds; 0 of 5 seeds at the cap
\newcommand{\ProbeAucBvcResonantLinearOneseven}{0.99067}  % experiments/FIGS/data/probe_ladder_v2_mlp2/analysis/seed_level_results.json :: levels.bvc_resonant.linear[3].mean_auc :: sha256 53e1ec81343ac2c1 :: 5 seeds, 0 saturated at AUC=1
\newcommand{\ProbeAucSdBvcResonantLinearOneseven}{0.00147}  % experiments/FIGS/data/probe_ladder_v2_mlp2/analysis/seed_level_results.json :: table[bvc_resonant,linear,17].auc (SD over seeds, ddof=1) :: sha256 53e1ec81343ac2c1
\newcommand{\ProbeLogOneMinusAucBvcResonantLinearOneseven}{\ensuremath{-}4.6848}  % experiments/FIGS/data/probe_ladder_v2_mlp2/analysis/seed_level_results.json :: levels.bvc_resonant.linear[3].mean :: sha256 53e1ec81343ac2c1 :: natural log of 1-AUC, lower is better
\newcommand{\ProbeLogOneMinusAucSdBvcResonantLinearOneseven}{0.1619}  % experiments/FIGS/data/probe_ladder_v2_mlp2/analysis/seed_level_results.json :: levels.bvc_resonant.linear[3].seed_sd :: sha256 53e1ec81343ac2c1
\newcommand{\ProbeOmaBvcResonantLinearOneseven}{\ensuremath{(9.3 \pm 1.5)\times10^{-3}}}  % experiments/FIGS/data/probe_ladder_v2_mlp2/analysis/seed_level_results.json :: table[bvc_resonant,linear,17].auc :: sha256 53e1ec81343ac2c1 :: 1-AUC, mean +- SD over 5 seeds
\newcommand{\ProbeBkgLeftBvcResonantLinearOneseven}{90.8}  % experiments/FIGS/data/probe_ladder_v2_mlp2/analysis/seed_level_results.json :: levels.bvc_resonant.linear[3].rejection_points['0.90'].mean_n_bkg_pass :: sha256 53e1ec81343ac2c1 :: background jets passing the cut, mean over seeds
\newcommand{\ProbeRejBvcResonantLinearOneseven}{134\,\ensuremath{\pm}\,21}  % experiments/FIGS/data/probe_ladder_v2_mlp2/analysis/seed_level_results.json :: table[bvc_resonant,linear,17].rejection_points['0.90'].rejection :: sha256 53e1ec81343ac2c1 :: at 90% signal efficiency, mean +- SD over seeds; 0 of 5 seeds at the cap
\newcommand{\ProbeOmaRatioBvcResonantLinearOnesixtwoOverOneeighteight}{1.02}  % experiments/FIGS/data/probe_ladder_v2_mlp2/analysis/seed_level_results.json :: table[bvc_resonant,linear,162/188].auc :: sha256 53e1ec81343ac2c1 :: ratio of the seed means of 1-AUC, coarser over finer
\newcommand{\ProbeOmaRatioBvcResonantLinearFourthreeOverOneeighteight}{1.09}  % experiments/FIGS/data/probe_ladder_v2_mlp2/analysis/seed_level_results.json :: table[bvc_resonant,linear,43/188].auc :: sha256 53e1ec81343ac2c1 :: ratio of the seed means of 1-AUC, coarser over finer
\newcommand{\ProbeOmaRatioBvcResonantLinearOnesevenOverOneeighteight}{4.0}  % experiments/FIGS/data/probe_ladder_v2_mlp2/analysis/seed_level_results.json :: table[bvc_resonant,linear,17/188].auc :: sha256 53e1ec81343ac2c1 :: ratio of the seed means of 1-AUC, coarser over finer
\newcommand{\ProbeOmaRatioBvcResonantLinearFourthreeOverOnesixtwo}{1.07}  % experiments/FIGS/data/probe_ladder_v2_mlp2/analysis/seed_level_results.json :: table[bvc_resonant,linear,43/162].auc :: sha256 53e1ec81343ac2c1 :: ratio of the seed means of 1-AUC, coarser over finer
\newcommand{\ProbeOmaRatioBvcResonantLinearOnesevenOverOnesixtwo}{4.0}  % experiments/FIGS/data/probe_ladder_v2_mlp2/analysis/seed_level_results.json :: table[bvc_resonant,linear,17/162].auc :: sha256 53e1ec81343ac2c1 :: ratio of the seed means of 1-AUC, coarser over finer
\newcommand{\ProbeOmaRatioBvcResonantLinearOnesevenOverFourthree}{3.7}  % experiments/FIGS/data/probe_ladder_v2_mlp2/analysis/seed_level_results.json :: table[bvc_resonant,linear,17/43].auc :: sha256 53e1ec81343ac2c1 :: ratio of the seed means of 1-AUC, coarser over finer
\newcommand{\ProbeAucBvcResonantMlpOneeighteight}{0.99778}  % experiments/FIGS/data/probe_ladder_v2_mlp2/analysis/seed_level_results.json :: levels.bvc_resonant.mlp[0].mean_auc :: sha256 53e1ec81343ac2c1 :: 5 seeds, 0 saturated at AUC=1
\newcommand{\ProbeAucSdBvcResonantMlpOneeighteight}{0.00009}  % experiments/FIGS/data/probe_ladder_v2_mlp2/analysis/seed_level_results.json :: table[bvc_resonant,mlp,188].auc (SD over seeds, ddof=1) :: sha256 53e1ec81343ac2c1
\newcommand{\ProbeLogOneMinusAucBvcResonantMlpOneeighteight}{\ensuremath{-}6.1102}  % experiments/FIGS/data/probe_ladder_v2_mlp2/analysis/seed_level_results.json :: levels.bvc_resonant.mlp[0].mean :: sha256 53e1ec81343ac2c1 :: natural log of 1-AUC, lower is better
\newcommand{\ProbeLogOneMinusAucSdBvcResonantMlpOneeighteight}{0.0383}  % experiments/FIGS/data/probe_ladder_v2_mlp2/analysis/seed_level_results.json :: levels.bvc_resonant.mlp[0].seed_sd :: sha256 53e1ec81343ac2c1
\newcommand{\ProbeOmaBvcResonantMlpOneeighteight}{\ensuremath{(2.22 \pm 0.09)\times10^{-3}}}  % experiments/FIGS/data/probe_ladder_v2_mlp2/analysis/seed_level_results.json :: table[bvc_resonant,mlp,188].auc :: sha256 53e1ec81343ac2c1 :: 1-AUC, mean +- SD over 5 seeds
\newcommand{\ProbeBkgLeftBvcResonantMlpOneeighteight}{9.8}  % experiments/FIGS/data/probe_ladder_v2_mlp2/analysis/seed_level_results.json :: levels.bvc_resonant.mlp[0].rejection_points['0.90'].mean_n_bkg_pass :: sha256 53e1ec81343ac2c1 :: background jets passing the cut, mean over seeds
\newcommand{\ProbeRejBvcResonantMlpOneeighteight}{1{,}230\,\ensuremath{\pm}\,180}  % experiments/FIGS/data/probe_ladder_v2_mlp2/analysis/seed_level_results.json :: table[bvc_resonant,mlp,188].rejection_points['0.90'].rejection :: sha256 53e1ec81343ac2c1 :: at 90% signal efficiency, mean +- SD over seeds; 0 of 5 seeds at the cap
\newcommand{\ProbeAucBvcResonantMlpOnesixtwo}{0.99768}  % experiments/FIGS/data/probe_ladder_v2_mlp2/analysis/seed_level_results.json :: levels.bvc_resonant.mlp[1].mean_auc :: sha256 53e1ec81343ac2c1 :: 5 seeds, 0 saturated at AUC=1
\newcommand{\ProbeAucSdBvcResonantMlpOnesixtwo}{0.00005}  % experiments/FIGS/data/probe_ladder_v2_mlp2/analysis/seed_level_results.json :: table[bvc_resonant,mlp,162].auc (SD over seeds, ddof=1) :: sha256 53e1ec81343ac2c1
\newcommand{\ProbeLogOneMinusAucBvcResonantMlpOnesixtwo}{\ensuremath{-}6.0667}  % experiments/FIGS/data/probe_ladder_v2_mlp2/analysis/seed_level_results.json :: levels.bvc_resonant.mlp[1].mean :: sha256 53e1ec81343ac2c1 :: natural log of 1-AUC, lower is better
\newcommand{\ProbeLogOneMinusAucSdBvcResonantMlpOnesixtwo}{0.0220}  % experiments/FIGS/data/probe_ladder_v2_mlp2/analysis/seed_level_results.json :: levels.bvc_resonant.mlp[1].seed_sd :: sha256 53e1ec81343ac2c1
\newcommand{\ProbeOmaBvcResonantMlpOnesixtwo}{\ensuremath{(2.32 \pm 0.05)\times10^{-3}}}  % experiments/FIGS/data/probe_ladder_v2_mlp2/analysis/seed_level_results.json :: table[bvc_resonant,mlp,162].auc :: sha256 53e1ec81343ac2c1 :: 1-AUC, mean +- SD over 5 seeds
\newcommand{\ProbeBkgLeftBvcResonantMlpOnesixtwo}{10.4}  % experiments/FIGS/data/probe_ladder_v2_mlp2/analysis/seed_level_results.json :: levels.bvc_resonant.mlp[1].rejection_points['0.90'].mean_n_bkg_pass :: sha256 53e1ec81343ac2c1 :: background jets passing the cut, mean over seeds
\newcommand{\ProbeRejBvcResonantMlpOnesixtwo}{1{,}170\,\ensuremath{\pm}\,220}  % experiments/FIGS/data/probe_ladder_v2_mlp2/analysis/seed_level_results.json :: table[bvc_resonant,mlp,162].rejection_points['0.90'].rejection :: sha256 53e1ec81343ac2c1 :: at 90% signal efficiency, mean +- SD over seeds; 0 of 5 seeds at the cap
\newcommand{\ProbeAucBvcResonantMlpFourthree}{0.99747}  % experiments/FIGS/data/probe_ladder_v2_mlp2/analysis/seed_level_results.json :: levels.bvc_resonant.mlp[2].mean_auc :: sha256 53e1ec81343ac2c1 :: 5 seeds, 0 saturated at AUC=1
\newcommand{\ProbeAucSdBvcResonantMlpFourthree}{0.00004}  % experiments/FIGS/data/probe_ladder_v2_mlp2/analysis/seed_level_results.json :: table[bvc_resonant,mlp,43].auc (SD over seeds, ddof=1) :: sha256 53e1ec81343ac2c1
\newcommand{\ProbeLogOneMinusAucBvcResonantMlpFourthree}{\ensuremath{-}5.9792}  % experiments/FIGS/data/probe_ladder_v2_mlp2/analysis/seed_level_results.json :: levels.bvc_resonant.mlp[2].mean :: sha256 53e1ec81343ac2c1 :: natural log of 1-AUC, lower is better
\newcommand{\ProbeLogOneMinusAucSdBvcResonantMlpFourthree}{0.0164}  % experiments/FIGS/data/probe_ladder_v2_mlp2/analysis/seed_level_results.json :: levels.bvc_resonant.mlp[2].seed_sd :: sha256 53e1ec81343ac2c1
\newcommand{\ProbeOmaBvcResonantMlpFourthree}{\ensuremath{(2.53 \pm 0.04)\times10^{-3}}}  % experiments/FIGS/data/probe_ladder_v2_mlp2/analysis/seed_level_results.json :: table[bvc_resonant,mlp,43].auc :: sha256 53e1ec81343ac2c1 :: 1-AUC, mean +- SD over 5 seeds
\newcommand{\ProbeBkgLeftBvcResonantMlpFourthree}{11.8}  % experiments/FIGS/data/probe_ladder_v2_mlp2/analysis/seed_level_results.json :: levels.bvc_resonant.mlp[2].rejection_points['0.90'].mean_n_bkg_pass :: sha256 53e1ec81343ac2c1 :: background jets passing the cut, mean over seeds
\newcommand{\ProbeRejBvcResonantMlpFourthree}{1{,}030\,\ensuremath{\pm}\,170}  % experiments/FIGS/data/probe_ladder_v2_mlp2/analysis/seed_level_results.json :: table[bvc_resonant,mlp,43].rejection_points['0.90'].rejection :: sha256 53e1ec81343ac2c1 :: at 90% signal efficiency, mean +- SD over seeds; 0 of 5 seeds at the cap
\newcommand{\ProbeAucBvcResonantMlpOneseven}{0.99274}  % experiments/FIGS/data/probe_ladder_v2_mlp2/analysis/seed_level_results.json :: levels.bvc_resonant.mlp[3].mean_auc :: sha256 53e1ec81343ac2c1 :: 5 seeds, 0 saturated at AUC=1
\newcommand{\ProbeAucSdBvcResonantMlpOneseven}{0.00089}  % experiments/FIGS/data/probe_ladder_v2_mlp2/analysis/seed_level_results.json :: table[bvc_resonant,mlp,17].auc (SD over seeds, ddof=1) :: sha256 53e1ec81343ac2c1
\newcommand{\ProbeLogOneMinusAucBvcResonantMlpOneseven}{\ensuremath{-}4.9309}  % experiments/FIGS/data/probe_ladder_v2_mlp2/analysis/seed_level_results.json :: levels.bvc_resonant.mlp[3].mean :: sha256 53e1ec81343ac2c1 :: natural log of 1-AUC, lower is better
\newcommand{\ProbeLogOneMinusAucSdBvcResonantMlpOneseven}{0.1224}  % experiments/FIGS/data/probe_ladder_v2_mlp2/analysis/seed_level_results.json :: levels.bvc_resonant.mlp[3].seed_sd :: sha256 53e1ec81343ac2c1
\newcommand{\ProbeOmaBvcResonantMlpOneseven}{\ensuremath{(7.3 \pm 0.9)\times10^{-3}}}  % experiments/FIGS/data/probe_ladder_v2_mlp2/analysis/seed_level_results.json :: table[bvc_resonant,mlp,17].auc :: sha256 53e1ec81343ac2c1 :: 1-AUC, mean +- SD over 5 seeds
\newcommand{\ProbeBkgLeftBvcResonantMlpOneseven}{62.6}  % experiments/FIGS/data/probe_ladder_v2_mlp2/analysis/seed_level_results.json :: levels.bvc_resonant.mlp[3].rejection_points['0.90'].mean_n_bkg_pass :: sha256 53e1ec81343ac2c1 :: background jets passing the cut, mean over seeds
\newcommand{\ProbeRejBvcResonantMlpOneseven}{192\,\ensuremath{\pm}\,21}  % experiments/FIGS/data/probe_ladder_v2_mlp2/analysis/seed_level_results.json :: table[bvc_resonant,mlp,17].rejection_points['0.90'].rejection :: sha256 53e1ec81343ac2c1 :: at 90% signal efficiency, mean +- SD over seeds; 0 of 5 seeds at the cap
\newcommand{\ProbeOmaRatioBvcResonantMlpOnesixtwoOverOneeighteight}{1.04}  % experiments/FIGS/data/probe_ladder_v2_mlp2/analysis/seed_level_results.json :: table[bvc_resonant,mlp,162/188].auc :: sha256 53e1ec81343ac2c1 :: ratio of the seed means of 1-AUC, coarser over finer
\newcommand{\ProbeOmaRatioBvcResonantMlpFourthreeOverOneeighteight}{1.14}  % experiments/FIGS/data/probe_ladder_v2_mlp2/analysis/seed_level_results.json :: table[bvc_resonant,mlp,43/188].auc :: sha256 53e1ec81343ac2c1 :: ratio of the seed means of 1-AUC, coarser over finer
\newcommand{\ProbeOmaRatioBvcResonantMlpOnesevenOverOneeighteight}{3.3}  % experiments/FIGS/data/probe_ladder_v2_mlp2/analysis/seed_level_results.json :: table[bvc_resonant,mlp,17/188].auc :: sha256 53e1ec81343ac2c1 :: ratio of the seed means of 1-AUC, coarser over finer
\newcommand{\ProbeOmaRatioBvcResonantMlpFourthreeOverOnesixtwo}{1.09}  % experiments/FIGS/data/probe_ladder_v2_mlp2/analysis/seed_level_results.json :: table[bvc_resonant,mlp,43/162].auc :: sha256 53e1ec81343ac2c1 :: ratio of the seed means of 1-AUC, coarser over finer
\newcommand{\ProbeOmaRatioBvcResonantMlpOnesevenOverOnesixtwo}{3.1}  % experiments/FIGS/data/probe_ladder_v2_mlp2/analysis/seed_level_results.json :: table[bvc_resonant,mlp,17/162].auc :: sha256 53e1ec81343ac2c1 :: ratio of the seed means of 1-AUC, coarser over finer
\newcommand{\ProbeOmaRatioBvcResonantMlpOnesevenOverFourthree}{2.9}  % experiments/FIGS/data/probe_ladder_v2_mlp2/analysis/seed_level_results.json :: table[bvc_resonant,mlp,17/43].auc :: sha256 53e1ec81343ac2c1 :: ratio of the seed means of 1-AUC, coarser over finer
\newcommand{\ProbeAucEeVsMmLinearOneeighteight}{\ensuremath{1^{\dagger}}}  % experiments/FIGS/data/probe_ladder_v2_mlp2/analysis/seed_level_results.json :: levels.ee_vs_mm.linear[0].mean_auc :: sha256 53e1ec81343ac2c1 :: 5 seeds, 5 saturated at AUC=1
\newcommand{\ProbeAucSdEeVsMmLinearOneeighteight}{0.00000}  % experiments/FIGS/data/probe_ladder_v2_mlp2/analysis/seed_level_results.json :: table[ee_vs_mm,linear,188].auc (SD over seeds, ddof=1) :: sha256 53e1ec81343ac2c1
\newcommand{\ProbeLogOneMinusAucEeVsMmLinearOneeighteight}{\ensuremath{-}16.9579}  % experiments/FIGS/data/probe_ladder_v2_mlp2/analysis/seed_level_results.json :: levels.ee_vs_mm.linear[0].mean :: sha256 53e1ec81343ac2c1 :: natural log of 1-AUC, lower is better
\newcommand{\ProbeLogOneMinusAucSdEeVsMmLinearOneeighteight}{0.0000}  % experiments/FIGS/data/probe_ladder_v2_mlp2/analysis/seed_level_results.json :: levels.ee_vs_mm.linear[0].seed_sd :: sha256 53e1ec81343ac2c1
\newcommand{\ProbeOmaEeVsMmLinearOneeighteight}{---}  % experiments/FIGS/data/probe_ladder_v2_mlp2/analysis/seed_level_results.json :: table[ee_vs_mm,linear,188].auc :: sha256 53e1ec81343ac2c1 :: not quoted: the AUC reached 1 in a seed, so 1-AUC is only a bound
\newcommand{\ProbeBkgLeftEeVsMmLinearOneeighteight}{0.0}  % experiments/FIGS/data/probe_ladder_v2_mlp2/analysis/seed_level_results.json :: levels.ee_vs_mm.linear[0].rejection_points['0.90'].mean_n_bkg_pass :: sha256 53e1ec81343ac2c1 :: background jets passing the cut, mean over seeds
\newcommand{\ProbeRejEeVsMmLinearOneeighteight}{\ensuremath{>}1{,}571}  % experiments/FIGS/data/probe_ladder_v2_mlp2/analysis/seed_level_results.json :: table[ee_vs_mm,linear,188].rejection_points['0.90'].rejection :: sha256 53e1ec81343ac2c1 :: at 90% signal efficiency, mean +- SD over seeds; 5 of 5 seeds at the cap
\newcommand{\ProbeAucEeVsMmLinearOnesixtwo}{\ensuremath{1^{\dagger}}}  % experiments/FIGS/data/probe_ladder_v2_mlp2/analysis/seed_level_results.json :: levels.ee_vs_mm.linear[1].mean_auc :: sha256 53e1ec81343ac2c1 :: 5 seeds, 5 saturated at AUC=1
\newcommand{\ProbeAucSdEeVsMmLinearOnesixtwo}{0.00000}  % experiments/FIGS/data/probe_ladder_v2_mlp2/analysis/seed_level_results.json :: table[ee_vs_mm,linear,162].auc (SD over seeds, ddof=1) :: sha256 53e1ec81343ac2c1
\newcommand{\ProbeLogOneMinusAucEeVsMmLinearOnesixtwo}{\ensuremath{-}16.9579}  % experiments/FIGS/data/probe_ladder_v2_mlp2/analysis/seed_level_results.json :: levels.ee_vs_mm.linear[1].mean :: sha256 53e1ec81343ac2c1 :: natural log of 1-AUC, lower is better
\newcommand{\ProbeLogOneMinusAucSdEeVsMmLinearOnesixtwo}{0.0000}  % experiments/FIGS/data/probe_ladder_v2_mlp2/analysis/seed_level_results.json :: levels.ee_vs_mm.linear[1].seed_sd :: sha256 53e1ec81343ac2c1
\newcommand{\ProbeOmaEeVsMmLinearOnesixtwo}{---}  % experiments/FIGS/data/probe_ladder_v2_mlp2/analysis/seed_level_results.json :: table[ee_vs_mm,linear,162].auc :: sha256 53e1ec81343ac2c1 :: not quoted: the AUC reached 1 in a seed, so 1-AUC is only a bound
\newcommand{\ProbeBkgLeftEeVsMmLinearOnesixtwo}{0.0}  % experiments/FIGS/data/probe_ladder_v2_mlp2/analysis/seed_level_results.json :: levels.ee_vs_mm.linear[1].rejection_points['0.90'].mean_n_bkg_pass :: sha256 53e1ec81343ac2c1 :: background jets passing the cut, mean over seeds
\newcommand{\ProbeRejEeVsMmLinearOnesixtwo}{\ensuremath{>}1{,}571}  % experiments/FIGS/data/probe_ladder_v2_mlp2/analysis/seed_level_results.json :: table[ee_vs_mm,linear,162].rejection_points['0.90'].rejection :: sha256 53e1ec81343ac2c1 :: at 90% signal efficiency, mean +- SD over seeds; 5 of 5 seeds at the cap
\newcommand{\ProbeAucEeVsMmLinearFourthree}{0.97519}  % experiments/FIGS/data/probe_ladder_v2_mlp2/analysis/seed_level_results.json :: levels.ee_vs_mm.linear[2].mean_auc :: sha256 53e1ec81343ac2c1 :: 5 seeds, 0 saturated at AUC=1
\newcommand{\ProbeAucSdEeVsMmLinearFourthree}{0.02073}  % experiments/FIGS/data/probe_ladder_v2_mlp2/analysis/seed_level_results.json :: table[ee_vs_mm,linear,43].auc (SD over seeds, ddof=1) :: sha256 53e1ec81343ac2c1
\newcommand{\ProbeLogOneMinusAucEeVsMmLinearFourthree}{\ensuremath{-}4.0326}  % experiments/FIGS/data/probe_ladder_v2_mlp2/analysis/seed_level_results.json :: levels.ee_vs_mm.linear[2].mean :: sha256 53e1ec81343ac2c1 :: natural log of 1-AUC, lower is better
\newcommand{\ProbeLogOneMinusAucSdEeVsMmLinearFourthree}{0.9467}  % experiments/FIGS/data/probe_ladder_v2_mlp2/analysis/seed_level_results.json :: levels.ee_vs_mm.linear[2].seed_sd :: sha256 53e1ec81343ac2c1
\newcommand{\ProbeOmaEeVsMmLinearFourthree}{\ensuremath{(2.5 \pm 2.1)\times10^{-2}}}  % experiments/FIGS/data/probe_ladder_v2_mlp2/analysis/seed_level_results.json :: table[ee_vs_mm,linear,43].auc :: sha256 53e1ec81343ac2c1 :: 1-AUC, mean +- SD over 5 seeds
\newcommand{\ProbeBkgLeftEeVsMmLinearFourthree}{155.8}  % experiments/FIGS/data/probe_ladder_v2_mlp2/analysis/seed_level_results.json :: levels.ee_vs_mm.linear[2].rejection_points['0.90'].mean_n_bkg_pass :: sha256 53e1ec81343ac2c1 :: background jets passing the cut, mean over seeds
\newcommand{\ProbeRejEeVsMmLinearFourthree}{80\,\ensuremath{\pm}\,70}  % experiments/FIGS/data/probe_ladder_v2_mlp2/analysis/seed_level_results.json :: table[ee_vs_mm,linear,43].rejection_points['0.90'].rejection :: sha256 53e1ec81343ac2c1 :: at 90% signal efficiency, mean +- SD over seeds; 0 of 5 seeds at the cap
\newcommand{\ProbeAucEeVsMmLinearOneseven}{0.95832}  % experiments/FIGS/data/probe_ladder_v2_mlp2/analysis/seed_level_results.json :: levels.ee_vs_mm.linear[3].mean_auc :: sha256 53e1ec81343ac2c1 :: 5 seeds, 0 saturated at AUC=1
\newcommand{\ProbeAucSdEeVsMmLinearOneseven}{0.00970}  % experiments/FIGS/data/probe_ladder_v2_mlp2/analysis/seed_level_results.json :: table[ee_vs_mm,linear,17].auc (SD over seeds, ddof=1) :: sha256 53e1ec81343ac2c1
\newcommand{\ProbeLogOneMinusAucEeVsMmLinearOneseven}{\ensuremath{-}3.2013}  % experiments/FIGS/data/probe_ladder_v2_mlp2/analysis/seed_level_results.json :: levels.ee_vs_mm.linear[3].mean :: sha256 53e1ec81343ac2c1 :: natural log of 1-AUC, lower is better
\newcommand{\ProbeLogOneMinusAucSdEeVsMmLinearOneseven}{0.2488}  % experiments/FIGS/data/probe_ladder_v2_mlp2/analysis/seed_level_results.json :: levels.ee_vs_mm.linear[3].seed_sd :: sha256 53e1ec81343ac2c1
\newcommand{\ProbeOmaEeVsMmLinearOneseven}{\ensuremath{(4.2 \pm 1.0)\times10^{-2}}}  % experiments/FIGS/data/probe_ladder_v2_mlp2/analysis/seed_level_results.json :: table[ee_vs_mm,linear,17].auc :: sha256 53e1ec81343ac2c1 :: 1-AUC, mean +- SD over 5 seeds
\newcommand{\ProbeBkgLeftEeVsMmLinearOneseven}{289.6}  % experiments/FIGS/data/probe_ladder_v2_mlp2/analysis/seed_level_results.json :: levels.ee_vs_mm.linear[3].rejection_points['0.90'].mean_n_bkg_pass :: sha256 53e1ec81343ac2c1 :: background jets passing the cut, mean over seeds
\newcommand{\ProbeRejEeVsMmLinearOneseven}{17\,\ensuremath{\pm}\,6}  % experiments/FIGS/data/probe_ladder_v2_mlp2/analysis/seed_level_results.json :: table[ee_vs_mm,linear,17].rejection_points['0.90'].rejection :: sha256 53e1ec81343ac2c1 :: at 90% signal efficiency, mean +- SD over seeds; 0 of 5 seeds at the cap
\newcommand{\ProbeOmaRatioEeVsMmLinearOnesevenOverFourthree}{1.68}  % experiments/FIGS/data/probe_ladder_v2_mlp2/analysis/seed_level_results.json :: table[ee_vs_mm,linear,17/43].auc :: sha256 53e1ec81343ac2c1 :: ratio of the seed means of 1-AUC, coarser over finer
\newcommand{\ProbeAucEeVsMmMlpOneeighteight}{\ensuremath{1^{\dagger}}}  % experiments/FIGS/data/probe_ladder_v2_mlp2/analysis/seed_level_results.json :: levels.ee_vs_mm.mlp[0].mean_auc :: sha256 53e1ec81343ac2c1 :: 5 seeds, 5 saturated at AUC=1
\newcommand{\ProbeAucSdEeVsMmMlpOneeighteight}{0.00000}  % experiments/FIGS/data/probe_ladder_v2_mlp2/analysis/seed_level_results.json :: table[ee_vs_mm,mlp,188].auc (SD over seeds, ddof=1) :: sha256 53e1ec81343ac2c1
\newcommand{\ProbeLogOneMinusAucEeVsMmMlpOneeighteight}{\ensuremath{-}16.9579}  % experiments/FIGS/data/probe_ladder_v2_mlp2/analysis/seed_level_results.json :: levels.ee_vs_mm.mlp[0].mean :: sha256 53e1ec81343ac2c1 :: natural log of 1-AUC, lower is better
\newcommand{\ProbeLogOneMinusAucSdEeVsMmMlpOneeighteight}{0.0000}  % experiments/FIGS/data/probe_ladder_v2_mlp2/analysis/seed_level_results.json :: levels.ee_vs_mm.mlp[0].seed_sd :: sha256 53e1ec81343ac2c1
\newcommand{\ProbeOmaEeVsMmMlpOneeighteight}{---}  % experiments/FIGS/data/probe_ladder_v2_mlp2/analysis/seed_level_results.json :: table[ee_vs_mm,mlp,188].auc :: sha256 53e1ec81343ac2c1 :: not quoted: the AUC reached 1 in a seed, so 1-AUC is only a bound
\newcommand{\ProbeBkgLeftEeVsMmMlpOneeighteight}{0.0}  % experiments/FIGS/data/probe_ladder_v2_mlp2/analysis/seed_level_results.json :: levels.ee_vs_mm.mlp[0].rejection_points['0.90'].mean_n_bkg_pass :: sha256 53e1ec81343ac2c1 :: background jets passing the cut, mean over seeds
\newcommand{\ProbeRejEeVsMmMlpOneeighteight}{\ensuremath{>}1{,}571}  % experiments/FIGS/data/probe_ladder_v2_mlp2/analysis/seed_level_results.json :: table[ee_vs_mm,mlp,188].rejection_points['0.90'].rejection :: sha256 53e1ec81343ac2c1 :: at 90% signal efficiency, mean +- SD over seeds; 5 of 5 seeds at the cap
\newcommand{\ProbeAucEeVsMmMlpOnesixtwo}{\ensuremath{1^{\dagger}}}  % experiments/FIGS/data/probe_ladder_v2_mlp2/analysis/seed_level_results.json :: levels.ee_vs_mm.mlp[1].mean_auc :: sha256 53e1ec81343ac2c1 :: 5 seeds, 5 saturated at AUC=1
\newcommand{\ProbeAucSdEeVsMmMlpOnesixtwo}{0.00000}  % experiments/FIGS/data/probe_ladder_v2_mlp2/analysis/seed_level_results.json :: table[ee_vs_mm,mlp,162].auc (SD over seeds, ddof=1) :: sha256 53e1ec81343ac2c1
\newcommand{\ProbeLogOneMinusAucEeVsMmMlpOnesixtwo}{\ensuremath{-}16.9579}  % experiments/FIGS/data/probe_ladder_v2_mlp2/analysis/seed_level_results.json :: levels.ee_vs_mm.mlp[1].mean :: sha256 53e1ec81343ac2c1 :: natural log of 1-AUC, lower is better
\newcommand{\ProbeLogOneMinusAucSdEeVsMmMlpOnesixtwo}{0.0000}  % experiments/FIGS/data/probe_ladder_v2_mlp2/analysis/seed_level_results.json :: levels.ee_vs_mm.mlp[1].seed_sd :: sha256 53e1ec81343ac2c1
\newcommand{\ProbeOmaEeVsMmMlpOnesixtwo}{---}  % experiments/FIGS/data/probe_ladder_v2_mlp2/analysis/seed_level_results.json :: table[ee_vs_mm,mlp,162].auc :: sha256 53e1ec81343ac2c1 :: not quoted: the AUC reached 1 in a seed, so 1-AUC is only a bound
\newcommand{\ProbeBkgLeftEeVsMmMlpOnesixtwo}{0.0}  % experiments/FIGS/data/probe_ladder_v2_mlp2/analysis/seed_level_results.json :: levels.ee_vs_mm.mlp[1].rejection_points['0.90'].mean_n_bkg_pass :: sha256 53e1ec81343ac2c1 :: background jets passing the cut, mean over seeds
\newcommand{\ProbeRejEeVsMmMlpOnesixtwo}{\ensuremath{>}1{,}571}  % experiments/FIGS/data/probe_ladder_v2_mlp2/analysis/seed_level_results.json :: table[ee_vs_mm,mlp,162].rejection_points['0.90'].rejection :: sha256 53e1ec81343ac2c1 :: at 90% signal efficiency, mean +- SD over seeds; 5 of 5 seeds at the cap
\newcommand{\ProbeAucEeVsMmMlpFourthree}{0.99657}  % experiments/FIGS/data/probe_ladder_v2_mlp2/analysis/seed_level_results.json :: levels.ee_vs_mm.mlp[2].mean_auc :: sha256 53e1ec81343ac2c1 :: 5 seeds, 0 saturated at AUC=1
\newcommand{\ProbeAucSdEeVsMmMlpFourthree}{0.00300}  % experiments/FIGS/data/probe_ladder_v2_mlp2/analysis/seed_level_results.json :: table[ee_vs_mm,mlp,43].auc (SD over seeds, ddof=1) :: sha256 53e1ec81343ac2c1
\newcommand{\ProbeLogOneMinusAucEeVsMmMlpFourthree}{\ensuremath{-}6.0371}  % experiments/FIGS/data/probe_ladder_v2_mlp2/analysis/seed_level_results.json :: levels.ee_vs_mm.mlp[2].mean :: sha256 53e1ec81343ac2c1 :: natural log of 1-AUC, lower is better
\newcommand{\ProbeLogOneMinusAucSdEeVsMmMlpFourthree}{0.9962}  % experiments/FIGS/data/probe_ladder_v2_mlp2/analysis/seed_level_results.json :: levels.ee_vs_mm.mlp[2].seed_sd :: sha256 53e1ec81343ac2c1
\newcommand{\ProbeOmaEeVsMmMlpFourthree}{\ensuremath{(3.4 \pm 3.0)\times10^{-3}}}  % experiments/FIGS/data/probe_ladder_v2_mlp2/analysis/seed_level_results.json :: table[ee_vs_mm,mlp,43].auc :: sha256 53e1ec81343ac2c1 :: 1-AUC, mean +- SD over 5 seeds
\newcommand{\ProbeBkgLeftEeVsMmMlpFourthree}{33.4}  % experiments/FIGS/data/probe_ladder_v2_mlp2/analysis/seed_level_results.json :: levels.ee_vs_mm.mlp[2].rejection_points['0.90'].mean_n_bkg_pass :: sha256 53e1ec81343ac2c1 :: background jets passing the cut, mean over seeds
\newcommand{\ProbeRejEeVsMmMlpFourthree}{400\,\ensuremath{\pm}\,500}  % experiments/FIGS/data/probe_ladder_v2_mlp2/analysis/seed_level_results.json :: table[ee_vs_mm,mlp,43].rejection_points['0.90'].rejection :: sha256 53e1ec81343ac2c1 :: at 90% signal efficiency, mean +- SD over seeds; 0 of 5 seeds at the cap
\newcommand{\ProbeAucEeVsMmMlpOneseven}{0.99338}  % experiments/FIGS/data/probe_ladder_v2_mlp2/analysis/seed_level_results.json :: levels.ee_vs_mm.mlp[3].mean_auc :: sha256 53e1ec81343ac2c1 :: 5 seeds, 0 saturated at AUC=1
\newcommand{\ProbeAucSdEeVsMmMlpOneseven}{0.00065}  % experiments/FIGS/data/probe_ladder_v2_mlp2/analysis/seed_level_results.json :: table[ee_vs_mm,mlp,17].auc (SD over seeds, ddof=1) :: sha256 53e1ec81343ac2c1
\newcommand{\ProbeLogOneMinusAucEeVsMmMlpOneseven}{\ensuremath{-}5.0211}  % experiments/FIGS/data/probe_ladder_v2_mlp2/analysis/seed_level_results.json :: levels.ee_vs_mm.mlp[3].mean :: sha256 53e1ec81343ac2c1 :: natural log of 1-AUC, lower is better
\newcommand{\ProbeLogOneMinusAucSdEeVsMmMlpOneseven}{0.0997}  % experiments/FIGS/data/probe_ladder_v2_mlp2/analysis/seed_level_results.json :: levels.ee_vs_mm.mlp[3].seed_sd :: sha256 53e1ec81343ac2c1
\newcommand{\ProbeOmaEeVsMmMlpOneseven}{\ensuremath{(6.6 \pm 0.7)\times10^{-3}}}  % experiments/FIGS/data/probe_ladder_v2_mlp2/analysis/seed_level_results.json :: table[ee_vs_mm,mlp,17].auc :: sha256 53e1ec81343ac2c1 :: 1-AUC, mean +- SD over 5 seeds
\newcommand{\ProbeBkgLeftEeVsMmMlpOneseven}{59.2}  % experiments/FIGS/data/probe_ladder_v2_mlp2/analysis/seed_level_results.json :: levels.ee_vs_mm.mlp[3].rejection_points['0.90'].mean_n_bkg_pass :: sha256 53e1ec81343ac2c1 :: background jets passing the cut, mean over seeds
\newcommand{\ProbeRejEeVsMmMlpOneseven}{81\,\ensuremath{\pm}\,12}  % experiments/FIGS/data/probe_ladder_v2_mlp2/analysis/seed_level_results.json :: table[ee_vs_mm,mlp,17].rejection_points['0.90'].rejection :: sha256 53e1ec81343ac2c1 :: at 90% signal efficiency, mean +- SD over seeds; 0 of 5 seeds at the cap
\newcommand{\ProbeOmaRatioEeVsMmMlpOnesevenOverFourthree}{1.93}  % experiments/FIGS/data/probe_ladder_v2_mlp2/analysis/seed_level_results.json :: table[ee_vs_mm,mlp,17/43].auc :: sha256 53e1ec81343ac2c1 :: ratio of the seed means of 1-AUC, coarser over finer
\newcommand{\ProbeAucRetainedTopologyLinearOneeighteight}{0.99413}  % experiments/FIGS/data/probe_ladder_v2_mlp2/analysis/seed_level_results.json :: levels.retained_topology.linear[0].mean_auc :: sha256 53e1ec81343ac2c1 :: 5 seeds, 0 saturated at AUC=1
\newcommand{\ProbeAucSdRetainedTopologyLinearOneeighteight}{0.00011}  % experiments/FIGS/data/probe_ladder_v2_mlp2/analysis/seed_level_results.json :: table[retained_topology,linear,188].auc (SD over seeds, ddof=1) :: sha256 53e1ec81343ac2c1
\newcommand{\ProbeLogOneMinusAucRetainedTopologyLinearOneeighteight}{\ensuremath{-}5.1388}  % experiments/FIGS/data/probe_ladder_v2_mlp2/analysis/seed_level_results.json :: levels.retained_topology.linear[0].mean :: sha256 53e1ec81343ac2c1 :: natural log of 1-AUC, lower is better
\newcommand{\ProbeLogOneMinusAucSdRetainedTopologyLinearOneeighteight}{0.0187}  % experiments/FIGS/data/probe_ladder_v2_mlp2/analysis/seed_level_results.json :: levels.retained_topology.linear[0].seed_sd :: sha256 53e1ec81343ac2c1
\newcommand{\ProbeOmaRetainedTopologyLinearOneeighteight}{\ensuremath{(5.87 \pm 0.11)\times10^{-3}}}  % experiments/FIGS/data/probe_ladder_v2_mlp2/analysis/seed_level_results.json :: table[retained_topology,linear,188].auc :: sha256 53e1ec81343ac2c1 :: 1-AUC, mean +- SD over 5 seeds
\newcommand{\ProbeBkgLeftRetainedTopologyLinearOneeighteight}{15.4}  % experiments/FIGS/data/probe_ladder_v2_mlp2/analysis/seed_level_results.json :: levels.retained_topology.linear[0].rejection_points['0.90'].mean_n_bkg_pass :: sha256 53e1ec81343ac2c1 :: background jets passing the cut, mean over seeds
\newcommand{\ProbeRejRetainedTopologyLinearOneeighteight}{71\,\ensuremath{\pm}\,7}  % experiments/FIGS/data/probe_ladder_v2_mlp2/analysis/seed_level_results.json :: table[retained_topology,linear,188].rejection_points['0.90'].rejection :: sha256 53e1ec81343ac2c1 :: at 90% signal efficiency, mean +- SD over seeds; 0 of 5 seeds at the cap
\newcommand{\ProbeAucRetainedTopologyLinearOnesixtwo}{0.99413}  % experiments/FIGS/data/probe_ladder_v2_mlp2/analysis/seed_level_results.json :: levels.retained_topology.linear[1].mean_auc :: sha256 53e1ec81343ac2c1 :: 5 seeds, 0 saturated at AUC=1
\newcommand{\ProbeAucSdRetainedTopologyLinearOnesixtwo}{0.00030}  % experiments/FIGS/data/probe_ladder_v2_mlp2/analysis/seed_level_results.json :: table[retained_topology,linear,162].auc (SD over seeds, ddof=1) :: sha256 53e1ec81343ac2c1
\newcommand{\ProbeLogOneMinusAucRetainedTopologyLinearOnesixtwo}{\ensuremath{-}5.1395}  % experiments/FIGS/data/probe_ladder_v2_mlp2/analysis/seed_level_results.json :: levels.retained_topology.linear[1].mean :: sha256 53e1ec81343ac2c1 :: natural log of 1-AUC, lower is better
\newcommand{\ProbeLogOneMinusAucSdRetainedTopologyLinearOnesixtwo}{0.0510}  % experiments/FIGS/data/probe_ladder_v2_mlp2/analysis/seed_level_results.json :: levels.retained_topology.linear[1].seed_sd :: sha256 53e1ec81343ac2c1
\newcommand{\ProbeOmaRetainedTopologyLinearOnesixtwo}{\ensuremath{(5.87 \pm 0.30)\times10^{-3}}}  % experiments/FIGS/data/probe_ladder_v2_mlp2/analysis/seed_level_results.json :: table[retained_topology,linear,162].auc :: sha256 53e1ec81343ac2c1 :: 1-AUC, mean +- SD over 5 seeds
\newcommand{\ProbeBkgLeftRetainedTopologyLinearOnesixtwo}{14.0}  % experiments/FIGS/data/probe_ladder_v2_mlp2/analysis/seed_level_results.json :: levels.retained_topology.linear[1].rejection_points['0.90'].mean_n_bkg_pass :: sha256 53e1ec81343ac2c1 :: background jets passing the cut, mean over seeds
\newcommand{\ProbeRejRetainedTopologyLinearOnesixtwo}{79\,\ensuremath{\pm}\,13}  % experiments/FIGS/data/probe_ladder_v2_mlp2/analysis/seed_level_results.json :: table[retained_topology,linear,162].rejection_points['0.90'].rejection :: sha256 53e1ec81343ac2c1 :: at 90% signal efficiency, mean +- SD over seeds; 0 of 5 seeds at the cap
\newcommand{\ProbeAucRetainedTopologyLinearFourthree}{0.99274}  % experiments/FIGS/data/probe_ladder_v2_mlp2/analysis/seed_level_results.json :: levels.retained_topology.linear[2].mean_auc :: sha256 53e1ec81343ac2c1 :: 5 seeds, 0 saturated at AUC=1
\newcommand{\ProbeAucSdRetainedTopologyLinearFourthree}{0.00012}  % experiments/FIGS/data/probe_ladder_v2_mlp2/analysis/seed_level_results.json :: table[retained_topology,linear,43].auc (SD over seeds, ddof=1) :: sha256 53e1ec81343ac2c1
\newcommand{\ProbeLogOneMinusAucRetainedTopologyLinearFourthree}{\ensuremath{-}4.9253}  % experiments/FIGS/data/probe_ladder_v2_mlp2/analysis/seed_level_results.json :: levels.retained_topology.linear[2].mean :: sha256 53e1ec81343ac2c1 :: natural log of 1-AUC, lower is better
\newcommand{\ProbeLogOneMinusAucSdRetainedTopologyLinearFourthree}{0.0169}  % experiments/FIGS/data/probe_ladder_v2_mlp2/analysis/seed_level_results.json :: levels.retained_topology.linear[2].seed_sd :: sha256 53e1ec81343ac2c1
\newcommand{\ProbeOmaRetainedTopologyLinearFourthree}{\ensuremath{(7.26 \pm 0.12)\times10^{-3}}}  % experiments/FIGS/data/probe_ladder_v2_mlp2/analysis/seed_level_results.json :: table[retained_topology,linear,43].auc :: sha256 53e1ec81343ac2c1 :: 1-AUC, mean +- SD over 5 seeds
\newcommand{\ProbeBkgLeftRetainedTopologyLinearFourthree}{16.4}  % experiments/FIGS/data/probe_ladder_v2_mlp2/analysis/seed_level_results.json :: levels.retained_topology.linear[2].rejection_points['0.90'].mean_n_bkg_pass :: sha256 53e1ec81343ac2c1 :: background jets passing the cut, mean over seeds
\newcommand{\ProbeRejRetainedTopologyLinearFourthree}{66\,\ensuremath{\pm}\,5}  % experiments/FIGS/data/probe_ladder_v2_mlp2/analysis/seed_level_results.json :: table[retained_topology,linear,43].rejection_points['0.90'].rejection :: sha256 53e1ec81343ac2c1 :: at 90% signal efficiency, mean +- SD over seeds; 0 of 5 seeds at the cap
\newcommand{\ProbeAucRetainedTopologyLinearOneseven}{0.99096}  % experiments/FIGS/data/probe_ladder_v2_mlp2/analysis/seed_level_results.json :: levels.retained_topology.linear[3].mean_auc :: sha256 53e1ec81343ac2c1 :: 5 seeds, 0 saturated at AUC=1
\newcommand{\ProbeAucSdRetainedTopologyLinearOneseven}{0.00092}  % experiments/FIGS/data/probe_ladder_v2_mlp2/analysis/seed_level_results.json :: table[retained_topology,linear,17].auc (SD over seeds, ddof=1) :: sha256 53e1ec81343ac2c1
\newcommand{\ProbeLogOneMinusAucRetainedTopologyLinearOneseven}{\ensuremath{-}4.7095}  % experiments/FIGS/data/probe_ladder_v2_mlp2/analysis/seed_level_results.json :: levels.retained_topology.linear[3].mean :: sha256 53e1ec81343ac2c1 :: natural log of 1-AUC, lower is better
\newcommand{\ProbeLogOneMinusAucSdRetainedTopologyLinearOneseven}{0.0979}  % experiments/FIGS/data/probe_ladder_v2_mlp2/analysis/seed_level_results.json :: levels.retained_topology.linear[3].seed_sd :: sha256 53e1ec81343ac2c1
\newcommand{\ProbeOmaRetainedTopologyLinearOneseven}{\ensuremath{(9.0 \pm 0.9)\times10^{-3}}}  % experiments/FIGS/data/probe_ladder_v2_mlp2/analysis/seed_level_results.json :: table[retained_topology,linear,17].auc :: sha256 53e1ec81343ac2c1 :: 1-AUC, mean +- SD over 5 seeds
\newcommand{\ProbeBkgLeftRetainedTopologyLinearOneseven}{21.8}  % experiments/FIGS/data/probe_ladder_v2_mlp2/analysis/seed_level_results.json :: levels.retained_topology.linear[3].rejection_points['0.90'].mean_n_bkg_pass :: sha256 53e1ec81343ac2c1 :: background jets passing the cut, mean over seeds
\newcommand{\ProbeRejRetainedTopologyLinearOneseven}{51\,\ensuremath{\pm}\,11}  % experiments/FIGS/data/probe_ladder_v2_mlp2/analysis/seed_level_results.json :: table[retained_topology,linear,17].rejection_points['0.90'].rejection :: sha256 53e1ec81343ac2c1 :: at 90% signal efficiency, mean +- SD over seeds; 0 of 5 seeds at the cap
\newcommand{\ProbeOmaRatioRetainedTopologyLinearOnesixtwoOverOneeighteight}{1.00}  % experiments/FIGS/data/probe_ladder_v2_mlp2/analysis/seed_level_results.json :: table[retained_topology,linear,162/188].auc :: sha256 53e1ec81343ac2c1 :: ratio of the seed means of 1-AUC, coarser over finer
\newcommand{\ProbeOmaRatioRetainedTopologyLinearFourthreeOverOneeighteight}{1.24}  % experiments/FIGS/data/probe_ladder_v2_mlp2/analysis/seed_level_results.json :: table[retained_topology,linear,43/188].auc :: sha256 53e1ec81343ac2c1 :: ratio of the seed means of 1-AUC, coarser over finer
\newcommand{\ProbeOmaRatioRetainedTopologyLinearOnesevenOverOneeighteight}{1.54}  % experiments/FIGS/data/probe_ladder_v2_mlp2/analysis/seed_level_results.json :: table[retained_topology,linear,17/188].auc :: sha256 53e1ec81343ac2c1 :: ratio of the seed means of 1-AUC, coarser over finer
\newcommand{\ProbeOmaRatioRetainedTopologyLinearFourthreeOverOnesixtwo}{1.24}  % experiments/FIGS/data/probe_ladder_v2_mlp2/analysis/seed_level_results.json :: table[retained_topology,linear,43/162].auc :: sha256 53e1ec81343ac2c1 :: ratio of the seed means of 1-AUC, coarser over finer
\newcommand{\ProbeOmaRatioRetainedTopologyLinearOnesevenOverOnesixtwo}{1.54}  % experiments/FIGS/data/probe_ladder_v2_mlp2/analysis/seed_level_results.json :: table[retained_topology,linear,17/162].auc :: sha256 53e1ec81343ac2c1 :: ratio of the seed means of 1-AUC, coarser over finer
\newcommand{\ProbeOmaRatioRetainedTopologyLinearOnesevenOverFourthree}{1.25}  % experiments/FIGS/data/probe_ladder_v2_mlp2/analysis/seed_level_results.json :: table[retained_topology,linear,17/43].auc :: sha256 53e1ec81343ac2c1 :: ratio of the seed means of 1-AUC, coarser over finer
\newcommand{\ProbeAucRetainedTopologyMlpOneeighteight}{0.99394}  % experiments/FIGS/data/probe_ladder_v2_mlp2/analysis/seed_level_results.json :: levels.retained_topology.mlp[0].mean_auc :: sha256 53e1ec81343ac2c1 :: 5 seeds, 0 saturated at AUC=1
\newcommand{\ProbeAucSdRetainedTopologyMlpOneeighteight}{0.00032}  % experiments/FIGS/data/probe_ladder_v2_mlp2/analysis/seed_level_results.json :: table[retained_topology,mlp,188].auc (SD over seeds, ddof=1) :: sha256 53e1ec81343ac2c1
\newcommand{\ProbeLogOneMinusAucRetainedTopologyMlpOneeighteight}{\ensuremath{-}5.1078}  % experiments/FIGS/data/probe_ladder_v2_mlp2/analysis/seed_level_results.json :: levels.retained_topology.mlp[0].mean :: sha256 53e1ec81343ac2c1 :: natural log of 1-AUC, lower is better
\newcommand{\ProbeLogOneMinusAucSdRetainedTopologyMlpOneeighteight}{0.0522}  % experiments/FIGS/data/probe_ladder_v2_mlp2/analysis/seed_level_results.json :: levels.retained_topology.mlp[0].seed_sd :: sha256 53e1ec81343ac2c1
\newcommand{\ProbeOmaRetainedTopologyMlpOneeighteight}{\ensuremath{(6.06 \pm 0.32)\times10^{-3}}}  % experiments/FIGS/data/probe_ladder_v2_mlp2/analysis/seed_level_results.json :: table[retained_topology,mlp,188].auc :: sha256 53e1ec81343ac2c1 :: 1-AUC, mean +- SD over 5 seeds
\newcommand{\ProbeBkgLeftRetainedTopologyMlpOneeighteight}{15.8}  % experiments/FIGS/data/probe_ladder_v2_mlp2/analysis/seed_level_results.json :: levels.retained_topology.mlp[0].rejection_points['0.90'].mean_n_bkg_pass :: sha256 53e1ec81343ac2c1 :: background jets passing the cut, mean over seeds
\newcommand{\ProbeRejRetainedTopologyMlpOneeighteight}{69\,\ensuremath{\pm}\,9}  % experiments/FIGS/data/probe_ladder_v2_mlp2/analysis/seed_level_results.json :: table[retained_topology,mlp,188].rejection_points['0.90'].rejection :: sha256 53e1ec81343ac2c1 :: at 90% signal efficiency, mean +- SD over seeds; 0 of 5 seeds at the cap
\newcommand{\ProbeAucRetainedTopologyMlpOnesixtwo}{0.99387}  % experiments/FIGS/data/probe_ladder_v2_mlp2/analysis/seed_level_results.json :: levels.retained_topology.mlp[1].mean_auc :: sha256 53e1ec81343ac2c1 :: 5 seeds, 0 saturated at AUC=1
\newcommand{\ProbeAucSdRetainedTopologyMlpOnesixtwo}{0.00048}  % experiments/FIGS/data/probe_ladder_v2_mlp2/analysis/seed_level_results.json :: table[retained_topology,mlp,162].auc (SD over seeds, ddof=1) :: sha256 53e1ec81343ac2c1
\newcommand{\ProbeLogOneMinusAucRetainedTopologyMlpOnesixtwo}{\ensuremath{-}5.0964}  % experiments/FIGS/data/probe_ladder_v2_mlp2/analysis/seed_level_results.json :: levels.retained_topology.mlp[1].mean :: sha256 53e1ec81343ac2c1 :: natural log of 1-AUC, lower is better
\newcommand{\ProbeLogOneMinusAucSdRetainedTopologyMlpOnesixtwo}{0.0773}  % experiments/FIGS/data/probe_ladder_v2_mlp2/analysis/seed_level_results.json :: levels.retained_topology.mlp[1].seed_sd :: sha256 53e1ec81343ac2c1
\newcommand{\ProbeOmaRetainedTopologyMlpOnesixtwo}{\ensuremath{(6.1 \pm 0.5)\times10^{-3}}}  % experiments/FIGS/data/probe_ladder_v2_mlp2/analysis/seed_level_results.json :: table[retained_topology,mlp,162].auc :: sha256 53e1ec81343ac2c1 :: 1-AUC, mean +- SD over 5 seeds
\newcommand{\ProbeBkgLeftRetainedTopologyMlpOnesixtwo}{16.6}  % experiments/FIGS/data/probe_ladder_v2_mlp2/analysis/seed_level_results.json :: levels.retained_topology.mlp[1].rejection_points['0.90'].mean_n_bkg_pass :: sha256 53e1ec81343ac2c1 :: background jets passing the cut, mean over seeds
\newcommand{\ProbeRejRetainedTopologyMlpOnesixtwo}{67\,\ensuremath{\pm}\,14}  % experiments/FIGS/data/probe_ladder_v2_mlp2/analysis/seed_level_results.json :: table[retained_topology,mlp,162].rejection_points['0.90'].rejection :: sha256 53e1ec81343ac2c1 :: at 90% signal efficiency, mean +- SD over seeds; 0 of 5 seeds at the cap
\newcommand{\ProbeAucRetainedTopologyMlpFourthree}{0.99284}  % experiments/FIGS/data/probe_ladder_v2_mlp2/analysis/seed_level_results.json :: levels.retained_topology.mlp[2].mean_auc :: sha256 53e1ec81343ac2c1 :: 5 seeds, 0 saturated at AUC=1
\newcommand{\ProbeAucSdRetainedTopologyMlpFourthree}{0.00040}  % experiments/FIGS/data/probe_ladder_v2_mlp2/analysis/seed_level_results.json :: table[retained_topology,mlp,43].auc (SD over seeds, ddof=1) :: sha256 53e1ec81343ac2c1
\newcommand{\ProbeLogOneMinusAucRetainedTopologyMlpFourthree}{\ensuremath{-}4.9400}  % experiments/FIGS/data/probe_ladder_v2_mlp2/analysis/seed_level_results.json :: levels.retained_topology.mlp[2].mean :: sha256 53e1ec81343ac2c1 :: natural log of 1-AUC, lower is better
\newcommand{\ProbeLogOneMinusAucSdRetainedTopologyMlpFourthree}{0.0570}  % experiments/FIGS/data/probe_ladder_v2_mlp2/analysis/seed_level_results.json :: levels.retained_topology.mlp[2].seed_sd :: sha256 53e1ec81343ac2c1
\newcommand{\ProbeOmaRetainedTopologyMlpFourthree}{\ensuremath{(7.2 \pm 0.4)\times10^{-3}}}  % experiments/FIGS/data/probe_ladder_v2_mlp2/analysis/seed_level_results.json :: table[retained_topology,mlp,43].auc :: sha256 53e1ec81343ac2c1 :: 1-AUC, mean +- SD over 5 seeds
\newcommand{\ProbeBkgLeftRetainedTopologyMlpFourthree}{17.8}  % experiments/FIGS/data/probe_ladder_v2_mlp2/analysis/seed_level_results.json :: levels.retained_topology.mlp[2].rejection_points['0.90'].mean_n_bkg_pass :: sha256 53e1ec81343ac2c1 :: background jets passing the cut, mean over seeds
\newcommand{\ProbeRejRetainedTopologyMlpFourthree}{61\,\ensuremath{\pm}\,7}  % experiments/FIGS/data/probe_ladder_v2_mlp2/analysis/seed_level_results.json :: table[retained_topology,mlp,43].rejection_points['0.90'].rejection :: sha256 53e1ec81343ac2c1 :: at 90% signal efficiency, mean +- SD over seeds; 0 of 5 seeds at the cap
\newcommand{\ProbeAucRetainedTopologyMlpOneseven}{0.99081}  % experiments/FIGS/data/probe_ladder_v2_mlp2/analysis/seed_level_results.json :: levels.retained_topology.mlp[3].mean_auc :: sha256 53e1ec81343ac2c1 :: 5 seeds, 0 saturated at AUC=1
\newcommand{\ProbeAucSdRetainedTopologyMlpOneseven}{0.00062}  % experiments/FIGS/data/probe_ladder_v2_mlp2/analysis/seed_level_results.json :: table[retained_topology,mlp,17].auc (SD over seeds, ddof=1) :: sha256 53e1ec81343ac2c1
\newcommand{\ProbeLogOneMinusAucRetainedTopologyMlpOneseven}{\ensuremath{-}4.6913}  % experiments/FIGS/data/probe_ladder_v2_mlp2/analysis/seed_level_results.json :: levels.retained_topology.mlp[3].mean :: sha256 53e1ec81343ac2c1 :: natural log of 1-AUC, lower is better
\newcommand{\ProbeLogOneMinusAucSdRetainedTopologyMlpOneseven}{0.0658}  % experiments/FIGS/data/probe_ladder_v2_mlp2/analysis/seed_level_results.json :: levels.retained_topology.mlp[3].seed_sd :: sha256 53e1ec81343ac2c1
\newcommand{\ProbeOmaRetainedTopologyMlpOneseven}{\ensuremath{(9.2 \pm 0.6)\times10^{-3}}}  % experiments/FIGS/data/probe_ladder_v2_mlp2/analysis/seed_level_results.json :: table[retained_topology,mlp,17].auc :: sha256 53e1ec81343ac2c1 :: 1-AUC, mean +- SD over 5 seeds
\newcommand{\ProbeBkgLeftRetainedTopologyMlpOneseven}{21.2}  % experiments/FIGS/data/probe_ladder_v2_mlp2/analysis/seed_level_results.json :: levels.retained_topology.mlp[3].rejection_points['0.90'].mean_n_bkg_pass :: sha256 53e1ec81343ac2c1 :: background jets passing the cut, mean over seeds
\newcommand{\ProbeRejRetainedTopologyMlpOneseven}{51\,\ensuremath{\pm}\,4}  % experiments/FIGS/data/probe_ladder_v2_mlp2/analysis/seed_level_results.json :: table[retained_topology,mlp,17].rejection_points['0.90'].rejection :: sha256 53e1ec81343ac2c1 :: at 90% signal efficiency, mean +- SD over seeds; 0 of 5 seeds at the cap
\newcommand{\ProbeOmaRatioRetainedTopologyMlpOnesixtwoOverOneeighteight}{1.01}  % experiments/FIGS/data/probe_ladder_v2_mlp2/analysis/seed_level_results.json :: table[retained_topology,mlp,162/188].auc :: sha256 53e1ec81343ac2c1 :: ratio of the seed means of 1-AUC, coarser over finer
\newcommand{\ProbeOmaRatioRetainedTopologyMlpFourthreeOverOneeighteight}{1.18}  % experiments/FIGS/data/probe_ladder_v2_mlp2/analysis/seed_level_results.json :: table[retained_topology,mlp,43/188].auc :: sha256 53e1ec81343ac2c1 :: ratio of the seed means of 1-AUC, coarser over finer
\newcommand{\ProbeOmaRatioRetainedTopologyMlpOnesevenOverOneeighteight}{1.52}  % experiments/FIGS/data/probe_ladder_v2_mlp2/analysis/seed_level_results.json :: table[retained_topology,mlp,17/188].auc :: sha256 53e1ec81343ac2c1 :: ratio of the seed means of 1-AUC, coarser over finer
\newcommand{\ProbeOmaRatioRetainedTopologyMlpFourthreeOverOnesixtwo}{1.17}  % experiments/FIGS/data/probe_ladder_v2_mlp2/analysis/seed_level_results.json :: table[retained_topology,mlp,43/162].auc :: sha256 53e1ec81343ac2c1 :: ratio of the seed means of 1-AUC, coarser over finer
\newcommand{\ProbeOmaRatioRetainedTopologyMlpOnesevenOverOnesixtwo}{1.50}  % experiments/FIGS/data/probe_ladder_v2_mlp2/analysis/seed_level_results.json :: table[retained_topology,mlp,17/162].auc :: sha256 53e1ec81343ac2c1 :: ratio of the seed means of 1-AUC, coarser over finer
\newcommand{\ProbeOmaRatioRetainedTopologyMlpOnesevenOverFourthree}{1.28}  % experiments/FIGS/data/probe_ladder_v2_mlp2/analysis/seed_level_results.json :: table[retained_topology,mlp,17/43].auc :: sha256 53e1ec81343ac2c1 :: ratio of the seed means of 1-AUC, coarser over finer
\newcommand{\ProbeAucVisibleContentLinearOneeighteight}{0.99917}  % experiments/FIGS/data/probe_ladder_v2_mlp2/analysis/seed_level_results.json :: levels.visible_content.linear[0].mean_auc :: sha256 53e1ec81343ac2c1 :: 5 seeds, 0 saturated at AUC=1
\newcommand{\ProbeAucSdVisibleContentLinearOneeighteight}{0.00003}  % experiments/FIGS/data/probe_ladder_v2_mlp2/analysis/seed_level_results.json :: table[visible_content,linear,188].auc (SD over seeds, ddof=1) :: sha256 53e1ec81343ac2c1
\newcommand{\ProbeLogOneMinusAucVisibleContentLinearOneeighteight}{\ensuremath{-}7.0999}  % experiments/FIGS/data/probe_ladder_v2_mlp2/analysis/seed_level_results.json :: levels.visible_content.linear[0].mean :: sha256 53e1ec81343ac2c1 :: natural log of 1-AUC, lower is better
\newcommand{\ProbeLogOneMinusAucSdVisibleContentLinearOneeighteight}{0.0358}  % experiments/FIGS/data/probe_ladder_v2_mlp2/analysis/seed_level_results.json :: levels.visible_content.linear[0].seed_sd :: sha256 53e1ec81343ac2c1
\newcommand{\ProbeOmaVisibleContentLinearOneeighteight}{\ensuremath{(8.26 \pm 0.30)\times10^{-4}}}  % experiments/FIGS/data/probe_ladder_v2_mlp2/analysis/seed_level_results.json :: table[visible_content,linear,188].auc :: sha256 53e1ec81343ac2c1 :: 1-AUC, mean +- SD over 5 seeds
\newcommand{\ProbeBkgLeftVisibleContentLinearOneeighteight}{0.4}  % experiments/FIGS/data/probe_ladder_v2_mlp2/analysis/seed_level_results.json :: levels.visible_content.linear[0].rejection_points['0.90'].mean_n_bkg_pass :: sha256 53e1ec81343ac2c1 :: background jets passing the cut, mean over seeds
\newcommand{\ProbeRejVisibleContentLinearOneeighteight}{\ensuremath{>}538}  % experiments/FIGS/data/probe_ladder_v2_mlp2/analysis/seed_level_results.json :: table[visible_content,linear,188].rejection_points['0.90'].rejection :: sha256 53e1ec81343ac2c1 :: at 90% signal efficiency, mean +- SD over seeds; 5 of 5 seeds at the cap
\newcommand{\ProbeAucVisibleContentLinearOnesixtwo}{0.99921}  % experiments/FIGS/data/probe_ladder_v2_mlp2/analysis/seed_level_results.json :: levels.visible_content.linear[1].mean_auc :: sha256 53e1ec81343ac2c1 :: 5 seeds, 0 saturated at AUC=1
\newcommand{\ProbeAucSdVisibleContentLinearOnesixtwo}{0.00004}  % experiments/FIGS/data/probe_ladder_v2_mlp2/analysis/seed_level_results.json :: table[visible_content,linear,162].auc (SD over seeds, ddof=1) :: sha256 53e1ec81343ac2c1
\newcommand{\ProbeLogOneMinusAucVisibleContentLinearOnesixtwo}{\ensuremath{-}7.1480}  % experiments/FIGS/data/probe_ladder_v2_mlp2/analysis/seed_level_results.json :: levels.visible_content.linear[1].mean :: sha256 53e1ec81343ac2c1 :: natural log of 1-AUC, lower is better
\newcommand{\ProbeLogOneMinusAucSdVisibleContentLinearOnesixtwo}{0.0449}  % experiments/FIGS/data/probe_ladder_v2_mlp2/analysis/seed_level_results.json :: levels.visible_content.linear[1].seed_sd :: sha256 53e1ec81343ac2c1
\newcommand{\ProbeOmaVisibleContentLinearOnesixtwo}{\ensuremath{(7.87 \pm 0.35)\times10^{-4}}}  % experiments/FIGS/data/probe_ladder_v2_mlp2/analysis/seed_level_results.json :: table[visible_content,linear,162].auc :: sha256 53e1ec81343ac2c1 :: 1-AUC, mean +- SD over 5 seeds
\newcommand{\ProbeBkgLeftVisibleContentLinearOnesixtwo}{0.6}  % experiments/FIGS/data/probe_ladder_v2_mlp2/analysis/seed_level_results.json :: levels.visible_content.linear[1].rejection_points['0.90'].mean_n_bkg_pass :: sha256 53e1ec81343ac2c1 :: background jets passing the cut, mean over seeds
\newcommand{\ProbeRejVisibleContentLinearOnesixtwo}{\ensuremath{>}538}  % experiments/FIGS/data/probe_ladder_v2_mlp2/analysis/seed_level_results.json :: table[visible_content,linear,162].rejection_points['0.90'].rejection :: sha256 53e1ec81343ac2c1 :: at 90% signal efficiency, mean +- SD over seeds; 5 of 5 seeds at the cap
\newcommand{\ProbeAucVisibleContentLinearFourthree}{0.99918}  % experiments/FIGS/data/probe_ladder_v2_mlp2/analysis/seed_level_results.json :: levels.visible_content.linear[2].mean_auc :: sha256 53e1ec81343ac2c1 :: 5 seeds, 0 saturated at AUC=1
\newcommand{\ProbeAucSdVisibleContentLinearFourthree}{0.00009}  % experiments/FIGS/data/probe_ladder_v2_mlp2/analysis/seed_level_results.json :: table[visible_content,linear,43].auc (SD over seeds, ddof=1) :: sha256 53e1ec81343ac2c1
\newcommand{\ProbeLogOneMinusAucVisibleContentLinearFourthree}{\ensuremath{-}7.1124}  % experiments/FIGS/data/probe_ladder_v2_mlp2/analysis/seed_level_results.json :: levels.visible_content.linear[2].mean :: sha256 53e1ec81343ac2c1 :: natural log of 1-AUC, lower is better
\newcommand{\ProbeLogOneMinusAucSdVisibleContentLinearFourthree}{0.1137}  % experiments/FIGS/data/probe_ladder_v2_mlp2/analysis/seed_level_results.json :: levels.visible_content.linear[2].seed_sd :: sha256 53e1ec81343ac2c1
\newcommand{\ProbeOmaVisibleContentLinearFourthree}{\ensuremath{(8.2 \pm 1.0)\times10^{-4}}}  % experiments/FIGS/data/probe_ladder_v2_mlp2/analysis/seed_level_results.json :: table[visible_content,linear,43].auc :: sha256 53e1ec81343ac2c1 :: 1-AUC, mean +- SD over 5 seeds
\newcommand{\ProbeBkgLeftVisibleContentLinearFourthree}{0.4}  % experiments/FIGS/data/probe_ladder_v2_mlp2/analysis/seed_level_results.json :: levels.visible_content.linear[2].rejection_points['0.90'].mean_n_bkg_pass :: sha256 53e1ec81343ac2c1 :: background jets passing the cut, mean over seeds
\newcommand{\ProbeRejVisibleContentLinearFourthree}{\ensuremath{>}538}  % experiments/FIGS/data/probe_ladder_v2_mlp2/analysis/seed_level_results.json :: table[visible_content,linear,43].rejection_points['0.90'].rejection :: sha256 53e1ec81343ac2c1 :: at 90% signal efficiency, mean +- SD over seeds; 5 of 5 seeds at the cap
\newcommand{\ProbeAucVisibleContentLinearOneseven}{0.99768}  % experiments/FIGS/data/probe_ladder_v2_mlp2/analysis/seed_level_results.json :: levels.visible_content.linear[3].mean_auc :: sha256 53e1ec81343ac2c1 :: 5 seeds, 0 saturated at AUC=1
\newcommand{\ProbeAucSdVisibleContentLinearOneseven}{0.00020}  % experiments/FIGS/data/probe_ladder_v2_mlp2/analysis/seed_level_results.json :: table[visible_content,linear,17].auc (SD over seeds, ddof=1) :: sha256 53e1ec81343ac2c1
\newcommand{\ProbeLogOneMinusAucVisibleContentLinearOneseven}{\ensuremath{-}6.0684}  % experiments/FIGS/data/probe_ladder_v2_mlp2/analysis/seed_level_results.json :: levels.visible_content.linear[3].mean :: sha256 53e1ec81343ac2c1 :: natural log of 1-AUC, lower is better
\newcommand{\ProbeLogOneMinusAucSdVisibleContentLinearOneseven}{0.0889}  % experiments/FIGS/data/probe_ladder_v2_mlp2/analysis/seed_level_results.json :: levels.visible_content.linear[3].seed_sd :: sha256 53e1ec81343ac2c1
\newcommand{\ProbeOmaVisibleContentLinearOneseven}{\ensuremath{(2.32 \pm 0.20)\times10^{-3}}}  % experiments/FIGS/data/probe_ladder_v2_mlp2/analysis/seed_level_results.json :: table[visible_content,linear,17].auc :: sha256 53e1ec81343ac2c1 :: 1-AUC, mean +- SD over 5 seeds
\newcommand{\ProbeBkgLeftVisibleContentLinearOneseven}{3.4}  % experiments/FIGS/data/probe_ladder_v2_mlp2/analysis/seed_level_results.json :: levels.visible_content.linear[3].rejection_points['0.90'].mean_n_bkg_pass :: sha256 53e1ec81343ac2c1 :: background jets passing the cut, mean over seeds
\newcommand{\ProbeRejVisibleContentLinearOneseven}{900\,\ensuremath{\pm}\,400}  % experiments/FIGS/data/probe_ladder_v2_mlp2/analysis/seed_level_results.json :: table[visible_content,linear,17].rejection_points['0.90'].rejection :: sha256 53e1ec81343ac2c1 :: at 90% signal efficiency, mean +- SD over seeds; 0 of 5 seeds at the cap
\newcommand{\ProbeOmaRatioVisibleContentLinearOnesixtwoOverOneeighteight}{0.95}  % experiments/FIGS/data/probe_ladder_v2_mlp2/analysis/seed_level_results.json :: table[visible_content,linear,162/188].auc :: sha256 53e1ec81343ac2c1 :: ratio of the seed means of 1-AUC, coarser over finer
\newcommand{\ProbeOmaRatioVisibleContentLinearFourthreeOverOneeighteight}{0.99}  % experiments/FIGS/data/probe_ladder_v2_mlp2/analysis/seed_level_results.json :: table[visible_content,linear,43/188].auc :: sha256 53e1ec81343ac2c1 :: ratio of the seed means of 1-AUC, coarser over finer
\newcommand{\ProbeOmaRatioVisibleContentLinearOnesevenOverOneeighteight}{2.8}  % experiments/FIGS/data/probe_ladder_v2_mlp2/analysis/seed_level_results.json :: table[visible_content,linear,17/188].auc :: sha256 53e1ec81343ac2c1 :: ratio of the seed means of 1-AUC, coarser over finer
\newcommand{\ProbeOmaRatioVisibleContentLinearFourthreeOverOnesixtwo}{1.04}  % experiments/FIGS/data/probe_ladder_v2_mlp2/analysis/seed_level_results.json :: table[visible_content,linear,43/162].auc :: sha256 53e1ec81343ac2c1 :: ratio of the seed means of 1-AUC, coarser over finer
\newcommand{\ProbeOmaRatioVisibleContentLinearOnesevenOverOnesixtwo}{3.0}  % experiments/FIGS/data/probe_ladder_v2_mlp2/analysis/seed_level_results.json :: table[visible_content,linear,17/162].auc :: sha256 53e1ec81343ac2c1 :: ratio of the seed means of 1-AUC, coarser over finer
\newcommand{\ProbeOmaRatioVisibleContentLinearOnesevenOverFourthree}{2.8}  % experiments/FIGS/data/probe_ladder_v2_mlp2/analysis/seed_level_results.json :: table[visible_content,linear,17/43].auc :: sha256 53e1ec81343ac2c1 :: ratio of the seed means of 1-AUC, coarser over finer
\newcommand{\ProbeAucVisibleContentMlpOneeighteight}{0.99911}  % experiments/FIGS/data/probe_ladder_v2_mlp2/analysis/seed_level_results.json :: levels.visible_content.mlp[0].mean_auc :: sha256 53e1ec81343ac2c1 :: 5 seeds, 0 saturated at AUC=1
\newcommand{\ProbeAucSdVisibleContentMlpOneeighteight}{0.00005}  % experiments/FIGS/data/probe_ladder_v2_mlp2/analysis/seed_level_results.json :: table[visible_content,mlp,188].auc (SD over seeds, ddof=1) :: sha256 53e1ec81343ac2c1
\newcommand{\ProbeLogOneMinusAucVisibleContentMlpOneeighteight}{\ensuremath{-}7.0202}  % experiments/FIGS/data/probe_ladder_v2_mlp2/analysis/seed_level_results.json :: levels.visible_content.mlp[0].mean :: sha256 53e1ec81343ac2c1 :: natural log of 1-AUC, lower is better
\newcommand{\ProbeLogOneMinusAucSdVisibleContentMlpOneeighteight}{0.0570}  % experiments/FIGS/data/probe_ladder_v2_mlp2/analysis/seed_level_results.json :: levels.visible_content.mlp[0].seed_sd :: sha256 53e1ec81343ac2c1
\newcommand{\ProbeOmaVisibleContentMlpOneeighteight}{\ensuremath{(8.9 \pm 0.5)\times10^{-4}}}  % experiments/FIGS/data/probe_ladder_v2_mlp2/analysis/seed_level_results.json :: table[visible_content,mlp,188].auc :: sha256 53e1ec81343ac2c1 :: 1-AUC, mean +- SD over 5 seeds
\newcommand{\ProbeBkgLeftVisibleContentMlpOneeighteight}{0.4}  % experiments/FIGS/data/probe_ladder_v2_mlp2/analysis/seed_level_results.json :: levels.visible_content.mlp[0].rejection_points['0.90'].mean_n_bkg_pass :: sha256 53e1ec81343ac2c1 :: background jets passing the cut, mean over seeds
\newcommand{\ProbeRejVisibleContentMlpOneeighteight}{\ensuremath{>}538}  % experiments/FIGS/data/probe_ladder_v2_mlp2/analysis/seed_level_results.json :: table[visible_content,mlp,188].rejection_points['0.90'].rejection :: sha256 53e1ec81343ac2c1 :: at 90% signal efficiency, mean +- SD over seeds; 5 of 5 seeds at the cap
\newcommand{\ProbeAucVisibleContentMlpOnesixtwo}{0.99918}  % experiments/FIGS/data/probe_ladder_v2_mlp2/analysis/seed_level_results.json :: levels.visible_content.mlp[1].mean_auc :: sha256 53e1ec81343ac2c1 :: 5 seeds, 0 saturated at AUC=1
\newcommand{\ProbeAucSdVisibleContentMlpOnesixtwo}{0.00003}  % experiments/FIGS/data/probe_ladder_v2_mlp2/analysis/seed_level_results.json :: table[visible_content,mlp,162].auc (SD over seeds, ddof=1) :: sha256 53e1ec81343ac2c1
\newcommand{\ProbeLogOneMinusAucVisibleContentMlpOnesixtwo}{\ensuremath{-}7.1010}  % experiments/FIGS/data/probe_ladder_v2_mlp2/analysis/seed_level_results.json :: levels.visible_content.mlp[1].mean :: sha256 53e1ec81343ac2c1 :: natural log of 1-AUC, lower is better
\newcommand{\ProbeLogOneMinusAucSdVisibleContentMlpOnesixtwo}{0.0371}  % experiments/FIGS/data/probe_ladder_v2_mlp2/analysis/seed_level_results.json :: levels.visible_content.mlp[1].seed_sd :: sha256 53e1ec81343ac2c1
\newcommand{\ProbeOmaVisibleContentMlpOnesixtwo}{\ensuremath{(8.25 \pm 0.30)\times10^{-4}}}  % experiments/FIGS/data/probe_ladder_v2_mlp2/analysis/seed_level_results.json :: table[visible_content,mlp,162].auc :: sha256 53e1ec81343ac2c1 :: 1-AUC, mean +- SD over 5 seeds
\newcommand{\ProbeBkgLeftVisibleContentMlpOnesixtwo}{0.2}  % experiments/FIGS/data/probe_ladder_v2_mlp2/analysis/seed_level_results.json :: levels.visible_content.mlp[1].rejection_points['0.90'].mean_n_bkg_pass :: sha256 53e1ec81343ac2c1 :: background jets passing the cut, mean over seeds
\newcommand{\ProbeRejVisibleContentMlpOnesixtwo}{\ensuremath{>}538}  % experiments/FIGS/data/probe_ladder_v2_mlp2/analysis/seed_level_results.json :: table[visible_content,mlp,162].rejection_points['0.90'].rejection :: sha256 53e1ec81343ac2c1 :: at 90% signal efficiency, mean +- SD over seeds; 5 of 5 seeds at the cap
\newcommand{\ProbeAucVisibleContentMlpFourthree}{0.99917}  % experiments/FIGS/data/probe_ladder_v2_mlp2/analysis/seed_level_results.json :: levels.visible_content.mlp[2].mean_auc :: sha256 53e1ec81343ac2c1 :: 5 seeds, 0 saturated at AUC=1
\newcommand{\ProbeAucSdVisibleContentMlpFourthree}{0.00006}  % experiments/FIGS/data/probe_ladder_v2_mlp2/analysis/seed_level_results.json :: table[visible_content,mlp,43].auc (SD over seeds, ddof=1) :: sha256 53e1ec81343ac2c1
\newcommand{\ProbeLogOneMinusAucVisibleContentMlpFourthree}{\ensuremath{-}7.0957}  % experiments/FIGS/data/probe_ladder_v2_mlp2/analysis/seed_level_results.json :: levels.visible_content.mlp[2].mean :: sha256 53e1ec81343ac2c1 :: natural log of 1-AUC, lower is better
\newcommand{\ProbeLogOneMinusAucSdVisibleContentMlpFourthree}{0.0773}  % experiments/FIGS/data/probe_ladder_v2_mlp2/analysis/seed_level_results.json :: levels.visible_content.mlp[2].seed_sd :: sha256 53e1ec81343ac2c1
\newcommand{\ProbeOmaVisibleContentMlpFourthree}{\ensuremath{(8.3 \pm 0.6)\times10^{-4}}}  % experiments/FIGS/data/probe_ladder_v2_mlp2/analysis/seed_level_results.json :: table[visible_content,mlp,43].auc :: sha256 53e1ec81343ac2c1 :: 1-AUC, mean +- SD over 5 seeds
\newcommand{\ProbeBkgLeftVisibleContentMlpFourthree}{0.4}  % experiments/FIGS/data/probe_ladder_v2_mlp2/analysis/seed_level_results.json :: levels.visible_content.mlp[2].rejection_points['0.90'].mean_n_bkg_pass :: sha256 53e1ec81343ac2c1 :: background jets passing the cut, mean over seeds
\newcommand{\ProbeRejVisibleContentMlpFourthree}{\ensuremath{>}538}  % experiments/FIGS/data/probe_ladder_v2_mlp2/analysis/seed_level_results.json :: table[visible_content,mlp,43].rejection_points['0.90'].rejection :: sha256 53e1ec81343ac2c1 :: at 90% signal efficiency, mean +- SD over seeds; 5 of 5 seeds at the cap
\newcommand{\ProbeAucVisibleContentMlpOneseven}{0.99793}  % experiments/FIGS/data/probe_ladder_v2_mlp2/analysis/seed_level_results.json :: levels.visible_content.mlp[3].mean_auc :: sha256 53e1ec81343ac2c1 :: 5 seeds, 0 saturated at AUC=1
\newcommand{\ProbeAucSdVisibleContentMlpOneseven}{0.00020}  % experiments/FIGS/data/probe_ladder_v2_mlp2/analysis/seed_level_results.json :: table[visible_content,mlp,17].auc (SD over seeds, ddof=1) :: sha256 53e1ec81343ac2c1
\newcommand{\ProbeLogOneMinusAucVisibleContentMlpOneseven}{\ensuremath{-}6.1839}  % experiments/FIGS/data/probe_ladder_v2_mlp2/analysis/seed_level_results.json :: levels.visible_content.mlp[3].mean :: sha256 53e1ec81343ac2c1 :: natural log of 1-AUC, lower is better
\newcommand{\ProbeLogOneMinusAucSdVisibleContentMlpOneseven}{0.1014}  % experiments/FIGS/data/probe_ladder_v2_mlp2/analysis/seed_level_results.json :: levels.visible_content.mlp[3].seed_sd :: sha256 53e1ec81343ac2c1
\newcommand{\ProbeOmaVisibleContentMlpOneseven}{\ensuremath{(2.07 \pm 0.20)\times10^{-3}}}  % experiments/FIGS/data/probe_ladder_v2_mlp2/analysis/seed_level_results.json :: table[visible_content,mlp,17].auc :: sha256 53e1ec81343ac2c1 :: 1-AUC, mean +- SD over 5 seeds
\newcommand{\ProbeBkgLeftVisibleContentMlpOneseven}{3.4}  % experiments/FIGS/data/probe_ladder_v2_mlp2/analysis/seed_level_results.json :: levels.visible_content.mlp[3].rejection_points['0.90'].mean_n_bkg_pass :: sha256 53e1ec81343ac2c1 :: background jets passing the cut, mean over seeds
\newcommand{\ProbeRejVisibleContentMlpOneseven}{\ensuremath{\geq}638\ensuremath{^{\ast}}}  % experiments/FIGS/data/probe_ladder_v2_mlp2/analysis/seed_level_results.json :: table[visible_content,mlp,17].rejection_points['0.90'].rejection :: sha256 53e1ec81343ac2c1 :: at 90% signal efficiency, mean +- SD over seeds; 1 of 5 seeds at the cap
\newcommand{\ProbeOmaRatioVisibleContentMlpOnesixtwoOverOneeighteight}{0.92}  % experiments/FIGS/data/probe_ladder_v2_mlp2/analysis/seed_level_results.json :: table[visible_content,mlp,162/188].auc :: sha256 53e1ec81343ac2c1 :: ratio of the seed means of 1-AUC, coarser over finer
\newcommand{\ProbeOmaRatioVisibleContentMlpFourthreeOverOneeighteight}{0.93}  % experiments/FIGS/data/probe_ladder_v2_mlp2/analysis/seed_level_results.json :: table[visible_content,mlp,43/188].auc :: sha256 53e1ec81343ac2c1 :: ratio of the seed means of 1-AUC, coarser over finer
\newcommand{\ProbeOmaRatioVisibleContentMlpOnesevenOverOneeighteight}{2.3}  % experiments/FIGS/data/probe_ladder_v2_mlp2/analysis/seed_level_results.json :: table[visible_content,mlp,17/188].auc :: sha256 53e1ec81343ac2c1 :: ratio of the seed means of 1-AUC, coarser over finer
\newcommand{\ProbeOmaRatioVisibleContentMlpFourthreeOverOnesixtwo}{1.01}  % experiments/FIGS/data/probe_ladder_v2_mlp2/analysis/seed_level_results.json :: table[visible_content,mlp,43/162].auc :: sha256 53e1ec81343ac2c1 :: ratio of the seed means of 1-AUC, coarser over finer
\newcommand{\ProbeOmaRatioVisibleContentMlpOnesevenOverOnesixtwo}{2.5}  % experiments/FIGS/data/probe_ladder_v2_mlp2/analysis/seed_level_results.json :: table[visible_content,mlp,17/162].auc :: sha256 53e1ec81343ac2c1 :: ratio of the seed means of 1-AUC, coarser over finer
\newcommand{\ProbeOmaRatioVisibleContentMlpOnesevenOverFourthree}{2.5}  % experiments/FIGS/data/probe_ladder_v2_mlp2/analysis/seed_level_results.json :: table[visible_content,mlp,17/43].auc :: sha256 53e1ec81343ac2c1 :: ratio of the seed means of 1-AUC, coarser over finer
\newcommand{\VocabSizeLoneeighteight}{188}  % configs/labelmaps/rung_label_maps.v1.csv :: column L188 (distinct group count) :: sha256 412f9aa3b58db412
\newcommand{\VocabSizeLonesixtwo}{162}  % configs/labelmaps/rung_label_maps.v1.csv :: column L162 (distinct group count) :: sha256 412f9aa3b58db412
\newcommand{\VocabSizeRsixthreeQone}{64}  % configs/labelmaps/rung_label_maps.v1.csv :: column R63_Q1 (distinct group count) :: sha256 412f9aa3b58db412
\newcommand{\VocabSizeRfourtwoQone}{43}  % configs/labelmaps/rung_label_maps.v1.csv :: column R42_Q1 (distinct group count) :: sha256 412f9aa3b58db412
\newcommand{\VocabSizeRtwonineQone}{30}  % configs/labelmaps/rung_label_maps.v1.csv :: column R29_Q1 (distinct group count) :: sha256 412f9aa3b58db412
\newcommand{\VocabSizeRonesixQone}{17}  % configs/labelmaps/rung_label_maps.v1.csv :: column R16_Q1 (distinct group count) :: sha256 412f9aa3b58db412
\newcommand{\VocabSizeRthreeVIS}{4}  % configs/labelmaps/rung_label_maps.v1.csv :: column R3_VIS (distinct group count) :: sha256 412f9aa3b58db412
\newcommand{\VocabSizeRoneQone}{2}  % configs/labelmaps/rung_label_maps.v1.csv :: column R1_Q1 (distinct group count) :: sha256 412f9aa3b58db412
\newcommand{\ProbeSplitTrain}{60\%}  % experiments/EVAL/probe.py :: SPLIT_FRACTIONS :: sha256 485db0966e660588 :: share of the probe sample; the same jets for every model
\newcommand{\ProbeSplitVal}{20\%}  % experiments/EVAL/probe.py :: SPLIT_FRACTIONS :: sha256 485db0966e660588 :: share of the probe sample; the same jets for every model
\newcommand{\ProbeSplitTest}{20\%}  % experiments/EVAL/probe.py :: SPLIT_FRACTIONS :: sha256 485db0966e660588 :: share of the probe sample; the same jets for every model
\newcommand{\ProbeMlpHidden}{256}  % experiments/EVAL/probe.py :: _fit_mlp :: sha256 485db0966e660588 :: hidden width of the MLP probe
\newcommand{\ProbeMlpInits}{3}  % experiments/EVAL/probe.py :: MLP_SEEDS :: sha256 485db0966e660588 :: MLP initialisations whose scores are averaged per cell
\newcommand{\ProbeMlpLrDrop}{10}  % experiments/EVAL/probe.py :: MLP_SCHEDULE['lr_factor'] :: sha256 485db0966e660588 :: factor the MLP's learning rate drops by on a plateau
\newcommand{\ProbeMlpPlateaus}{three}  % experiments/EVAL/probe.py :: MLP_SCHEDULE: stop_patience / (lr_patience + 1) :: sha256 485db0966e660588 :: patience windows without gain after which the MLP stops
\newcommand{\ProbeCMin}{\ensuremath{10^{-2}}}  % experiments/EVAL/probe.py :: C_GRID :: sha256 485db0966e660588 :: end of the logistic-regression inverse-regularisation grid
\newcommand{\ProbeCMax}{\ensuremath{10^{2}}}  % experiments/EVAL/probe.py :: C_GRID :: sha256 485db0966e660588 :: end of the logistic-regression inverse-regularisation grid
\newcommand{\ProbeNBkgTestBvcFourprong}{3{,}158}  % experiments/FIGS/data/probe_ladder_v2_mlp2/s1.json :: tasks.bvc_4prong.n_background_test :: sha256 f7e1de0e43be89b8 :: test-sample background jets
\newcommand{\ProbeNBkgTestBvcQcd}{5{,}371}  % experiments/FIGS/data/probe_ladder_v2_mlp2/s1.json :: tasks.bvc_qcd.n_background_test :: sha256 f7e1de0e43be89b8 :: test-sample background jets
\newcommand{\ProbeNBkgTestBvcResonant}{11{,}876}  % experiments/FIGS/data/probe_ladder_v2_mlp2/s1.json :: tasks.bvc_resonant.n_background_test :: sha256 f7e1de0e43be89b8 :: test-sample background jets
\newcommand{\ProbeNBkgTestEeVsMm}{4{,}708}  % experiments/FIGS/data/probe_ladder_v2_mlp2/s1.json :: tasks.ee_vs_mm.n_background_test :: sha256 f7e1de0e43be89b8 :: test-sample background jets
\newcommand{\ProbeNBkgTestRetainedTopology}{1{,}078}  % experiments/FIGS/data/probe_ladder_v2_mlp2/s1.json :: tasks.retained_topology.n_background_test :: sha256 f7e1de0e43be89b8 :: test-sample background jets
\newcommand{\ProbeNBkgTestVisibleContent}{2{,}554}  % experiments/FIGS/data/probe_ladder_v2_mlp2/s1.json :: tasks.visible_content.n_background_test :: sha256 f7e1de0e43be89b8 :: test-sample background jets
\newcommand{\UseDiesAtSophonEqfour}{43}  % configs/labelmaps/usecase_survival.v1.json :: sophon_eq4.dies_at :: sha256 122275213c349593 :: first vocabulary size at which it is not constructible
\newcommand{\UseLastAliveSophonEqfour}{64}  % configs/labelmaps/usecase_survival.v1.json :: sophon_eq4.last_rung_alive :: sha256 122275213c349593
\newcommand{\UseDiesAtSophonEqsixAone}{64}  % configs/labelmaps/usecase_survival.v1.json :: sophon_eq6_a1.dies_at :: sha256 122275213c349593 :: first vocabulary size at which it is not constructible
\newcommand{\UseLastAliveSophonEqsixAone}{162}  % configs/labelmaps/usecase_survival.v1.json :: sophon_eq6_a1.last_rung_alive :: sha256 122275213c349593
\newcommand{\UseDiesAtSophonEqsixAthree}{64}  % configs/labelmaps/usecase_survival.v1.json :: sophon_eq6_a3.dies_at :: sha256 122275213c349593 :: first vocabulary size at which it is not constructible
\newcommand{\UseLastAliveSophonEqsixAthree}{162}  % configs/labelmaps/usecase_survival.v1.json :: sophon_eq6_a3.last_rung_alive :: sha256 122275213c349593
\newcommand{\UseDiesAtSophonEqsevenNine}{64}  % configs/labelmaps/usecase_survival.v1.json :: sophon_eq7_9.dies_at :: sha256 122275213c349593 :: first vocabulary size at which it is not constructible
\newcommand{\UseLastAliveSophonEqsevenNine}{162}  % configs/labelmaps/usecase_survival.v1.json :: sophon_eq7_9.last_rung_alive :: sha256 122275213c349593
\newcommand{\UseDiesAtVcbEqone}{64}  % configs/labelmaps/usecase_survival.v1.json :: vcb_eq1.dies_at :: sha256 122275213c349593 :: first vocabulary size at which it is not constructible
\newcommand{\UseLastAliveVcbEqone}{162}  % configs/labelmaps/usecase_survival.v1.json :: vcb_eq1.last_rung_alive :: sha256 122275213c349593
\newcommand{\UseDiesAtGnthreexQcdSubclass}{162}  % configs/labelmaps/usecase_survival.v1.json :: gn3x_qcd_subclass.dies_at :: sha256 122275213c349593 :: first vocabulary size at which it is not constructible
\newcommand{\UseLastAliveGnthreexQcdSubclass}{188}  % configs/labelmaps/usecase_survival.v1.json :: gn3x_qcd_subclass.last_rung_alive :: sha256 122275213c349593
\newcommand{\UseWlikeLastAlive}{162}  % configs/labelmaps/rung_label_maps.v1.csv :: groups at each level vs {label_X_cs, label_X_qq} and that plus the QCD classes :: sha256 412f9aa3b58db412 :: last level at which Sophon's W-like score g_W(2) is constructible
\newcommand{\UseWlikeDiesAt}{64}  % configs/labelmaps/rung_label_maps.v1.csv :: groups at each level vs {label_X_cs, label_X_qq} and that plus the QCD classes :: sha256 412f9aa3b58db412 :: first level at which it is not
\newcommand{\DesignEpochs}{80}  % experiments/MTX/k8s/job-mtx-l188-s1-raunav.yaml :: --num-epochs :: sha256 96fef4b6337aaa27
\newcommand{\DesignSamplesPerEpoch}{10{,}240{,}000}  % experiments/MTX/k8s/job-mtx-l188-s1-raunav.yaml :: --samples-per-epoch :: sha256 96fef4b6337aaa27
\newcommand{\DesignExamplesSeen}{\ensuremath{8.192\times10^{8}}}  % experiments/MTX/k8s/job-mtx-l188-s1-raunav.yaml :: --num-epochs x --samples-per-epoch :: sha256 96fef4b6337aaa27 :: training examples drawn over the run
\newcommand{\DesignPassesSelected}{7.45}  % experiments/FIGS/data/v2/jc2_train/counts.json :: n_selected :: sha256 b0b5a8b63aa808c0 :: --num-epochs x --samples-per-epoch of job-mtx-l188-s1-raunav.yaml / n_selected: passes over the selected training jets
\newcommand{\DesignBatchSize}{512}  % experiments/MTX/k8s/job-mtx-l188-s1-raunav.yaml :: --batch-size :: sha256 96fef4b6337aaa27
\newcommand{\DesignStartLr}{\ensuremath{5\times10^{-4}}}  % experiments/MTX/k8s/job-mtx-l188-s1-raunav.yaml :: --start-lr :: sha256 96fef4b6337aaa27
\newcommand{\DesignHeadHidden}{512}  % experiments/MTX/k8s/job-mtx-l188-s1-raunav.yaml :: -o fc_params (hidden width) :: sha256 96fef4b6337aaa27
\newcommand{\DesignParticleBlocks}{8}  % experiments/E1/ParT_sophon_arch_10c.py :: num_layers=(\d+) :: sha256 5793faa7d55a8b0d
\newcommand{\DesignClassBlocks}{2}  % experiments/E1/ParT_sophon_arch_10c.py :: num_cls_layers=(\d+) :: sha256 5793faa7d55a8b0d
\newcommand{\DesignHeads}{8}  % experiments/E1/ParT_sophon_arch_10c.py :: num_heads=(\d+) :: sha256 5793faa7d55a8b0d
\newcommand{\DesignEmbedDim}{128}  % experiments/E1/ParT_sophon_arch_10c.py :: \bembed_dims=\[\d+, \d+, (\d+)\] :: sha256 5793faa7d55a8b0d
\newcommand{\DesignMpmMaskRate}{40\%}  % experiments/MTX/mpm.py :: DEFAULT_MASK_RATE :: sha256 49f0c8c6666d0211 :: share of each jet's particles hidden from the encoder
\newcommand{\DesignPtMin}{200}  % configs/arms/L188.yaml :: selection :: sha256 d5a188dc80885515 :: GeV
\newcommand{\DesignPtMax}{2{,}500}  % configs/arms/L188.yaml :: selection :: sha256 d5a188dc80885515 :: GeV
\newcommand{\DesignMsdMin}{20}  % configs/arms/L188.yaml :: selection :: sha256 d5a188dc80885515 :: GeV
\newcommand{\DesignMsdMax}{500}  % configs/arms/L188.yaml :: selection :: sha256 d5a188dc80885515 :: GeV
\newcommand{\LitJetClassTwoNJetsMillion}{139}  % experiments/FIGS/data/literature_facts.json :: facts[0].value :: sha256 713dba2a5e6e1dd1 :: arXiv:2405.12972v1 App. A, p. 11
\newcommand{\LitJetClassTwoNResonantClasses}{161}  % experiments/FIGS/data/literature_facts.json :: facts[1].value :: sha256 713dba2a5e6e1dd1 :: arXiv:2405.12972v1 App. A, p. 11
\newcommand{\LitJetClassTwoNQcdClasses}{27}  % experiments/FIGS/data/literature_facts.json :: facts[2].value :: sha256 713dba2a5e6e1dd1 :: arXiv:2405.12972v1 App. A, p. 11
\newcommand{\LitJetClassTwoMaxResonanceMass}{500}  % experiments/FIGS/data/literature_facts.json :: facts[3].value :: sha256 713dba2a5e6e1dd1 :: arXiv:2405.12972v1 Sec. II, p. 2
\newcommand{\LitMadgraphVersion}{2.9.18}  % experiments/FIGS/data/literature_facts.json :: facts[4].value :: sha256 713dba2a5e6e1dd1 :: arXiv:2405.12972v1 App. A, p. 11
\newcommand{\LitPythiaVersion}{8.3}  % experiments/FIGS/data/literature_facts.json :: facts[5].value :: sha256 713dba2a5e6e1dd1 :: arXiv:2405.12972v1 App. A, p. 11
\newcommand{\LitDelphesVersion}{3}  % experiments/FIGS/data/literature_facts.json :: facts[6].value :: sha256 713dba2a5e6e1dd1 :: arXiv:2405.12972v1 App. A, p. 11
\newcommand{\LitJetClassTwoPileupMean}{50}  % experiments/FIGS/data/literature_facts.json :: facts[7].value :: sha256 713dba2a5e6e1dd1 :: arXiv:2405.12972v1 Sec. II, p. 2
\newcommand{\LitJetClassTwoJetRadius}{0.8}  % experiments/FIGS/data/literature_facts.json :: facts[8].value :: sha256 713dba2a5e6e1dd1 :: arXiv:2405.12972v1 Sec. II, p. 2
\newcommand{\LitAtlasUnifiedEffDropPercent}{10}  % experiments/FIGS/data/literature_facts.json :: facts[9].value :: sha256 713dba2a5e6e1dd1 :: ATL-PHYS-PUB-2026-013 Sec. 6.5, p. 22
\newcommand{\LitSophonPaperParticleBlocks}{6}  % experiments/FIGS/data/literature_facts.json :: facts[10].value :: sha256 713dba2a5e6e1dd1 :: arXiv:2405.12972v1 App. B.1, p. 12
\newcommand{\LitSophonSamplesPerEpochMillion}{10}  % experiments/FIGS/data/literature_facts.json :: facts[11].value :: sha256 713dba2a5e6e1dd1 :: arXiv:2405.12972v1 App. B.1, p. 12
\newcommand{\LitSophonSubsetSamplesPerEpochMillion}{2.5}  % experiments/FIGS/data/literature_facts.json :: facts[12].value :: sha256 713dba2a5e6e1dd1 :: arXiv:2405.12972v1 App. B.2, p. 12
\newcommand{\LitGlopartNClasses}{314}  % experiments/FIGS/data/literature_facts.json :: facts[13].value :: sha256 713dba2a5e6e1dd1 :: C. Li, PhD thesis, CERN-THESIS-2024-281 Sec. 11.2.1, p. 188
\newcommand{\LitPartFtEpochs}{20}  % experiments/FIGS/data/literature_facts.json :: facts[14].value :: sha256 713dba2a5e6e1dd1 :: arXiv:2202.03772 Sec. 5.2, p. 8
\newcommand{\LitMehtaNCategories}{eight}  % experiments/FIGS/data/literature_facts.json :: facts[15].value :: sha256 713dba2a5e6e1dd1 :: arXiv:2605.05556v2 Abstract
\newcommand{\LitHongPretrainRuns}{2}  % experiments/FIGS/data/literature_facts.json :: facts[16].value :: sha256 713dba2a5e6e1dd1 :: arXiv:2410.23129 App. A.1.1, p. 18
\newcommand{\LitSophonLookaheadK}{6}  % experiments/FIGS/data/literature_facts.json :: facts[17].value :: sha256 713dba2a5e6e1dd1 :: arXiv:2405.12972v1 App. B.1, p. 12
\newcommand{\LitSophonLookaheadAlpha}{0.5}  % experiments/FIGS/data/literature_facts.json :: facts[18].value :: sha256 713dba2a5e6e1dd1 :: arXiv:2405.12972v1 App. B.1, p. 12
\newcommand{\LitSophonRadamBetaOne}{0.95}  % experiments/FIGS/data/literature_facts.json :: facts[19].value :: sha256 713dba2a5e6e1dd1 :: arXiv:2405.12972v1 App. B.1, p. 12
\newcommand{\LitSophonRadamBetaTwo}{0.999}  % experiments/FIGS/data/literature_facts.json :: facts[20].value :: sha256 713dba2a5e6e1dd1 :: arXiv:2405.12972v1 App. B.1, p. 12
\newcommand{\LitSophonRadamEps}{\ensuremath{10^{-5}}}  % experiments/FIGS/data/literature_facts.json :: facts[21].value :: sha256 713dba2a5e6e1dd1 :: arXiv:2405.12972v1 App. B.1, p. 12
\newcommand{\LitGnTwoAuxRejectionDropPercent}{30}  % experiments/FIGS/data/literature_facts.json :: facts[22].value :: sha256 713dba2a5e6e1dd1 :: arXiv:2505.19689v2 (Nature Commun. 17 (2026) 541) p. 9
\newcommand{\LitCmsPartNClasses}{37}  % experiments/FIGS/data/literature_facts.json :: facts[23].value :: sha256 713dba2a5e6e1dd1 :: arXiv:2604.09809v1 (CMS-JME-25-001) Sec. 5.2, p. 7
\newcommand{\LitCmsPartMassLambda}{0.05}  % experiments/FIGS/data/literature_facts.json :: facts[24].value :: sha256 713dba2a5e6e1dd1 :: arXiv:2604.09809v1 (CMS-JME-25-001) Sec. 5.2, p. 7
\newcommand{\LitGlopartVthreeNClasses}{374}  % experiments/FIGS/data/literature_facts.json :: facts[25].value :: sha256 713dba2a5e6e1dd1 :: CMS-DP-2026-104 (CMS public page) Architecture and outputs
\newcommand{\LitHuhCoarseClasses}{127}  % experiments/FIGS/data/literature_facts.json :: facts[26].value :: sha256 713dba2a5e6e1dd1 :: arXiv:1608.08614v2 Sec. 5.3, p. 5
\newcommand{\LitHuhKnnGapPercent}{15}  % experiments/FIGS/data/literature_facts.json :: facts[27].value :: sha256 713dba2a5e6e1dd1 :: arXiv:1608.08614v2 Sec. 5.3, p. 5
\newcommand{\LitRsthreelRuns}{3}  % experiments/FIGS/data/literature_facts.json :: facts[28].value :: sha256 713dba2a5e6e1dd1 :: arXiv:2403.07066v2 (Phys. Rev. D 111 (2025) 032010) Sec. IV, p. 5
\newcommand{\LitSophonParamsMillion}{2.3}  % experiments/FIGS/data/literature_facts.json :: facts[29].value :: sha256 713dba2a5e6e1dd1 :: arXiv:2405.12972v1 App. B.1, p. 12
\newcommand{\LitGlopartVoneNClasses}{37}  % experiments/FIGS/data/literature_facts.json :: facts[30].value :: sha256 713dba2a5e6e1dd1 :: CMS-DP-2026-104 (CMS public page) Model generations
\newcommand{\LitSophonXMass}{200}  % experiments/FIGS/data/literature_facts.json :: facts[31].value :: sha256 713dba2a5e6e1dd1 :: arXiv:2405.12972v1 Sec. III, p. 3
\newcommand{\LitSophonXLumi}{40}  % experiments/FIGS/data/literature_facts.json :: facts[32].value :: sha256 713dba2a5e6e1dd1 :: arXiv:2405.12972v1 Sec. III, Fig. 2
\newcommand{\LitSophonXWindowLo}{150}  % experiments/FIGS/data/literature_facts.json :: facts[33].value :: sha256 713dba2a5e6e1dd1 :: arXiv:2405.12972v1 Sec. III, Fig. 2
\newcommand{\LitSophonXWindowHi}{230}  % experiments/FIGS/data/literature_facts.json :: facts[34].value :: sha256 713dba2a5e6e1dd1 :: arXiv:2405.12972v1 Sec. III, Fig. 2
\newcommand{\LitSophonAblationNClasses}{42}  % experiments/FIGS/data/literature_facts.json :: facts[35].value :: sha256 713dba2a5e6e1dd1 :: arXiv:2405.12972v1 Sec. III
\newcommand{\LitSophonTransferSamplesMillion}{2.5}  % experiments/FIGS/data/literature_facts.json :: facts[36].value :: sha256 713dba2a5e6e1dd1 :: arXiv:2405.12972v1 App. B.5, p. 13
\newcommand{\LitSophonTransferBatch}{1024}  % experiments/FIGS/data/literature_facts.json :: facts[37].value :: sha256 713dba2a5e6e1dd1 :: arXiv:2405.12972v1 App. B.5, p. 13
\newcommand{\LitSophonScratchEpochs}{50}  % experiments/FIGS/data/literature_facts.json :: facts[38].value :: sha256 713dba2a5e6e1dd1 :: arXiv:2405.12972v1 App. B.3
\newcommand{\LitHuhFineClasses}{1000}  % experiments/FIGS/data/literature_facts.json :: facts[39].value :: sha256 713dba2a5e6e1dd1 :: arXiv:1608.08614v2 Sec. 1, p. 1
\newcommand{\LitHuhPascalDropMap}{2.8}  % experiments/FIGS/data/literature_facts.json :: facts[40].value :: sha256 713dba2a5e6e1dd1 :: arXiv:1608.08614v2 Sec. 1, p. 1
\newcommand{\LitHongSuperclasses}{11}  % experiments/FIGS/data/literature_facts.json :: facts[41].value :: sha256 713dba2a5e6e1dd1 :: arXiv:2410.23129 Sec. 6.2
\newcommand{\LitHongDeclineClasses}{\ensuremath{10^{3}}}  % experiments/FIGS/data/literature_facts.json :: facts[42].value :: sha256 713dba2a5e6e1dd1 :: arXiv:2410.23129 Sec. 6.2
\newcommand{\LitObjectivesFullJets}{\ensuremath{1.2\times10^{6}}}  % experiments/FIGS/data/literature_facts.json :: facts[43].value :: sha256 713dba2a5e6e1dd1 :: arXiv:2606.14870 Sec. IV A
\newcommand{\LitObjectivesFewJets}{\ensuremath{10^{4}}}  % experiments/FIGS/data/literature_facts.json :: facts[44].value :: sha256 713dba2a5e6e1dd1 :: arXiv:2606.14870 Sec. IV A
\newcommand{\LitMpmBenefitJets}{10k}  % experiments/FIGS/data/literature_facts.json :: facts[45].value :: sha256 713dba2a5e6e1dd1 :: arXiv:2401.13537 Sec. V
\newcommand{\LitPartTopRejScratch}{1602}  % experiments/FIGS/data/literature_facts.json :: facts[46].value :: sha256 713dba2a5e6e1dd1 :: arXiv:2202.03772 Table 5
\newcommand{\LitPartTopRejFt}{2766}  % experiments/FIGS/data/literature_facts.json :: facts[47].value :: sha256 713dba2a5e6e1dd1 :: arXiv:2202.03772 Table 5
\newcommand{\LitMahajanLinearGapPercent}{7}  % experiments/FIGS/data/literature_facts.json :: facts[48].value :: sha256 713dba2a5e6e1dd1 :: arXiv:1805.00932 Sec. 3.1.2
\newcommand{\LitBukchinCoarse}{33.83}  % experiments/FIGS/data/literature_facts.json :: facts[49].value :: sha256 713dba2a5e6e1dd1 :: arXiv:2012.03515 Table 2, LIVING-17 all-way
\newcommand{\LitBukchinFine}{58.95}  % experiments/FIGS/data/literature_facts.json :: facts[50].value :: sha256 713dba2a5e6e1dd1 :: arXiv:2012.03515 Table 2, LIVING-17 all-way
\newcommand{\LitKornblithRho}{\ensuremath{-0.93}}  % experiments/FIGS/data/literature_facts.json :: facts[51].value :: sha256 713dba2a5e6e1dd1 :: arXiv:2010.16402 Sec. 3.4
\newcommand{\LitPartTopSigEffPercent}{30}  % experiments/FIGS/data/literature_facts.json :: facts[52].value :: sha256 713dba2a5e6e1dd1 :: arXiv:2202.03772 Table 5 (column Rej30%)
\newcommand{\LitPartTopKinFeatures}{7}  % experiments/FIGS/data/literature_facts.json :: facts[53].value :: sha256 713dba2a5e6e1dd1 :: arXiv:2202.03772 Sec. 5.2
\newcommand{\LitBukchinAncor}{45.14}  % experiments/FIGS/data/literature_facts.json :: facts[54].value :: sha256 713dba2a5e6e1dd1 :: arXiv:2012.03515 Table 2, LIVING-17 all-way
\newcommand{\LitMahajanFtGapPp}{0.1}  % experiments/FIGS/data/literature_facts.json :: facts[55].value :: sha256 713dba2a5e6e1dd1 :: arXiv:1805.00932 Fig. 1, ImageNet-9k target
\newcommand{\LitSophonTransferLr}{5 × 10−4}  % experiments/FIGS/data/literature_facts.json :: facts[56].value :: sha256 713dba2a5e6e1dd1 :: arXiv:2405.12972v1 App. B.5
\newcommand{\LitQgJetRadius}{0.4}  % experiments/FIGS/data/literature_facts.json :: facts[57].value :: sha256 713dba2a5e6e1dd1 :: arXiv:1810.05165 Sec. 4.1
\newcommand{\LitQgPtMin}{500}  % experiments/FIGS/data/literature_facts.json :: facts[58].value :: sha256 713dba2a5e6e1dd1 :: arXiv:1810.05165 Sec. 4.1
\newcommand{\LitQgPtMax}{550}  % experiments/FIGS/data/literature_facts.json :: facts[59].value :: sha256 713dba2a5e6e1dd1 :: arXiv:1810.05165 Sec. 4.1
\newcommand{\LitPartTopAucScratch}{0.9858}  % experiments/FIGS/data/literature_facts.json :: facts[60].value :: sha256 713dba2a5e6e1dd1 :: arXiv:2202.03772 Table 5
\newcommand{\LitPartTopAucFt}{0.9877}  % experiments/FIGS/data/literature_facts.json :: facts[61].value :: sha256 713dba2a5e6e1dd1 :: arXiv:2202.03772 Table 5
\newcommand{\LitPartQgAucScratch}{0.9121}  % experiments/FIGS/data/literature_facts.json :: facts[62].value :: sha256 713dba2a5e6e1dd1 :: arXiv:2202.03772 Table 6 (exp)
\newcommand{\LitPartQgRejScratch}{101.2}  % experiments/FIGS/data/literature_facts.json :: facts[63].value :: sha256 713dba2a5e6e1dd1 :: arXiv:2202.03772 Table 6 (exp)
\newcommand{\LitPartQgAucFt}{0.9151}  % experiments/FIGS/data/literature_facts.json :: facts[64].value :: sha256 713dba2a5e6e1dd1 :: arXiv:2202.03772 Table 6 (exp)
\newcommand{\LitPartQgRejFt}{107.9}  % experiments/FIGS/data/literature_facts.json :: facts[65].value :: sha256 713dba2a5e6e1dd1 :: arXiv:2202.03772 Table 6 (exp)
\newcommand{\LitPartJcTestMillion}{20}  % experiments/FIGS/data/literature_facts.json :: facts[66].value :: sha256 713dba2a5e6e1dd1 :: arXiv:2202.03772 Sec. 3
\newcommand{\LitPartJcTwoMAcc}{0.836}  % experiments/FIGS/data/literature_facts.json :: facts[67].value :: sha256 713dba2a5e6e1dd1 :: arXiv:2202.03772 Table 3
\newcommand{\LitPartSmallSetIters}{100 k}  % experiments/FIGS/data/literature_facts.json :: facts[68].value :: sha256 713dba2a5e6e1dd1 :: arXiv:2202.03772 Sec. 5.1 (training dataset size)
\newcommand{\LitPartSmallSetPercent}{2}  % experiments/FIGS/data/literature_facts.json :: facts[69].value :: sha256 713dba2a5e6e1dd1 :: arXiv:2202.03772 Sec. 5.1 (training dataset size)
\newcommand{\LitSophonEtaMax}{2.5}  % experiments/FIGS/data/literature_facts.json :: facts[70].value :: sha256 713dba2a5e6e1dd1 :: arXiv:2405.12972 App. A
\newcommand{\LitBHadronLifetimePs}{1.5}  % experiments/FIGS/data/literature_facts.json :: facts[71].value :: sha256 713dba2a5e6e1dd1 :: arXiv:1712.07158 Sec. 4.1
\newcommand{\LitCHadronLifetimePs}{1}  % experiments/FIGS/data/literature_facts.json :: facts[72].value :: sha256 713dba2a5e6e1dd1 :: arXiv:1712.07158 Sec. 4.1
\newcommand{\LitOmnilearnPidDropProb}{0.2}  % experiments/FIGS/data/literature_facts.json :: facts[73].value :: sha256 713dba2a5e6e1dd1 :: arXiv:2502.14652 Sec. II.A
\newcommand{\LitMahajanFtNineKLarge}{48.5}  % experiments/FIGS/data/literature_facts.json :: facts[74].value :: sha256 713dba2a5e6e1dd1 :: arXiv:1805.00932 Fig. 1 (IG-1B-17k, IN-9k, full fine-tuning)
\newcommand{\LitMahajanFtNineKSmall}{48.4}  % experiments/FIGS/data/literature_facts.json :: facts[75].value :: sha256 713dba2a5e6e1dd1 :: arXiv:1805.00932 Fig. 1 (IG-940M-1.5k, IN-9k, full fine-tuning)
\newcommand{\LitMahajanFtOneKSmall}{84.2}  % experiments/FIGS/data/literature_facts.json :: facts[76].value :: sha256 713dba2a5e6e1dd1 :: arXiv:1805.00932 Fig. 1 (IG-940M-1.5k, IN-1k, full fine-tuning)
\newcommand{\LitMahajanFtOneKLarge}{83.6}  % experiments/FIGS/data/literature_facts.json :: facts[77].value :: sha256 713dba2a5e6e1dd1 :: arXiv:1805.00932 Fig. 1 (IG-1B-17k, IN-1k, full fine-tuning)
\newcommand{\LitChenGainFifteen}{0.16}  % experiments/FIGS/data/literature_facts.json :: facts[78].value :: sha256 713dba2a5e6e1dd1 :: arXiv:1901.07012 Table VIII (15 coarse classes, CIFAR-100)
\newcommand{\LitChenGainTwenty}{0.01}  % experiments/FIGS/data/literature_facts.json :: facts[79].value :: sha256 713dba2a5e6e1dd1 :: arXiv:1901.07012 Table VIII (20 coarse classes, CIFAR-100)
\newcommand{\LitChenCoarseFifteen}{15}  % experiments/FIGS/data/literature_facts.json :: facts[80].value :: sha256 713dba2a5e6e1dd1 :: arXiv:1901.07012 Table VIII
\newcommand{\LitChenCoarseTwenty}{20}  % experiments/FIGS/data/literature_facts.json :: facts[81].value :: sha256 713dba2a5e6e1dd1 :: arXiv:1901.07012 Table VIII
\newcommand{\LitChenRandomMaxGain}{2.78}  % experiments/FIGS/data/literature_facts.json :: facts[82].value :: sha256 713dba2a5e6e1dd1 :: arXiv:1901.07012 Sec. VI-A, Table VI
\newcommand{\LitGlopartMassLambda}{5}  % experiments/FIGS/data/literature_facts.json :: facts[83].value :: sha256 713dba2a5e6e1dd1 :: CERN-THESIS-2024-281 Eq. (11.7), Sec. 11.2.1
\newcommand{\LitGlopartAddedAucBase}{0.9555}  % experiments/FIGS/data/literature_facts.json :: facts[84].value :: sha256 713dba2a5e6e1dd1 :: CERN-THESIS-2024-281 Sec. 11.4.1, Fig. 11.19 (H->WW 4q vs QCD)
\newcommand{\LitGlopartAddedAucTop}{0.9593}  % experiments/FIGS/data/literature_facts.json :: facts[85].value :: sha256 713dba2a5e6e1dd1 :: CERN-THESIS-2024-281 Sec. 11.4.1, Fig. 11.19 (H->WW 4q vs QCD)
\newcommand{\LitCmsSdZcut}{0.1}  % experiments/FIGS/data/literature_facts.json :: facts[86].value :: sha256 713dba2a5e6e1dd1 :: arXiv:2004.08262 Sec. 6.1
\newcommand{\LitCmsSdBeta}{0}  % experiments/FIGS/data/literature_facts.json :: facts[87].value :: sha256 713dba2a5e6e1dd1 :: arXiv:2004.08262 Sec. 6.1
\newcommand{\LitQgNJetsMillion}{2}  % experiments/FIGS/data/literature_facts.json :: facts[88].value :: sha256 713dba2a5e6e1dd1 :: arXiv:2202.03772 Sec. 5.2
\newcommand{\LitQgTrainMillion}{1.6}  % experiments/FIGS/data/literature_facts.json :: facts[89].value :: sha256 713dba2a5e6e1dd1 :: arXiv:2202.03772 Sec. 5.2
\newcommand{\LitQgValMillion}{0.2}  % experiments/FIGS/data/literature_facts.json :: facts[90].value :: sha256 713dba2a5e6e1dd1 :: arXiv:2202.03772 Sec. 5.2
\newcommand{\LitQgRapidityMax}{1.7}  % experiments/FIGS/data/literature_facts.json :: facts[91].value :: sha256 713dba2a5e6e1dd1 :: Zenodo 10.5281/zenodo.3164691 Zenodo record 3164691, description
\newcommand{\LitHerwigVersion}{7.1.4}  % experiments/FIGS/data/literature_facts.json :: facts[92].value :: sha256 713dba2a5e6e1dd1 :: Zenodo 10.5281/zenodo.3066475 record description
\newcommand{\LitObjectivesFewestJets}{\ensuremath{10^{3}}}  % experiments/FIGS/data/literature_facts.json :: facts[93].value :: sha256 713dba2a5e6e1dd1 :: arXiv:2606.14870 Table XVII (Micro model, Npt = 10^8: Classifier 136 vs MPM 218 at 10^3, 331 vs 539 at 10^4)
\newcommand{\LitPartTopNTrainMillion}{1.2}  % experiments/FIGS/data/literature_facts.json :: facts[94].value :: sha256 713dba2a5e6e1dd1 :: arXiv:2202.03772 Sec. 5.2
\newcommand{\LitLandscapeParticlenetEnsMean}{1412}  % experiments/FIGS/data/literature_facts.json :: facts[95].value :: sha256 713dba2a5e6e1dd1 :: arXiv:1902.09914 Table 1 (ParticleNet, ensemble mean of 1/eps_B at eps_S=0.3)
\newcommand{\LitPartFtHeadReinits}{9}  % experiments/FIGS/data/literature_facts.json :: facts[96].value :: sha256 713dba2a5e6e1dd1 :: arXiv:2202.03772 Sec. 5.2
\newcommand{\LitSophonBlocksCode}{8}  % experiments/FIGS/data/literature_facts.json :: facts[97].value :: sha256 713dba2a5e6e1dd1 :: github.com/jet-universe/sophon@9dd6dd6 networks/example_ParticleTransformer_sophon.py line 67
\newcommand{\LitGrasHerwigVersion}{2.7.1}  % experiments/FIGS/data/literature_facts.json :: facts[98].value :: sha256 713dba2a5e6e1dd1 :: arXiv:1704.03878 Sec. 2 (list of generators)
\newcommand{\DesignMassLambda}{5.0}  % experiments/MTX/k8s/job-mtx-l162_mass-s1-raunav.yaml :: --mass-lambda :: sha256 8882459fe5520795 :: identical in all ten mass-output pretraining specs
\newcommand{\FtRecipeNRuns}{264}  % experiments/FIGS/data/ft_recipes/recipes_w2b_bench_v2_v3.json :: runs (count), reported cells :: sha256 7dc8c791343e1b6f
\newcommand{\FtRecipeLr}{\ensuremath{1\times10^{-4}}}  % experiments/FIGS/data/ft_recipes/recipes_w2b_bench_v2_v3.json :: runs[*].lr, pretrained starts :: sha256 7dc8c791343e1b6f
\newcommand{\FtRecipeHeadMult}{50}  % experiments/FIGS/data/ft_recipes/recipes_w2b_bench_v2_v3.json :: runs[*].head_lr_mult, pretrained starts :: sha256 7dc8c791343e1b6f
\newcommand{\FtRecipeLrHead}{\ensuremath{5\times10^{-3}}}  % experiments/FIGS/data/ft_recipes/recipes_w2b_bench_v2_v3.json :: lr x head_lr_mult :: sha256 7dc8c791343e1b6f
\newcommand{\FtRecipeWeightDecay}{0.01}  % experiments/FIGS/data/ft_recipes/recipes_w2b_bench_v2_v3.json :: runs[*].weight_decay :: sha256 7dc8c791343e1b6f
\newcommand{\FtRecipeBatch}{512}  % experiments/FIGS/data/ft_recipes/recipes_w2b_bench_v2_v3.json :: runs[*].batch_size :: sha256 7dc8c791343e1b6f
\newcommand{\FtRecipeEpochsEThree}{50}  % experiments/FIGS/data/ft_recipes/recipes_w2b_bench_v2_v3.json :: runs[leg 1,2; n_train 1000].epochs :: sha256 7dc8c791343e1b6f
\newcommand{\FtRecipeStepsEThree}{19}  % experiments/FT/k8s/job-ft-legs-w2-raunav.yaml :: job commands: samples per epoch at n_train 1000 // batch size :: sha256 f0962eddb451b7b0 :: optimizer steps per epoch, as weaver floors them
\newcommand{\FtRecipePassesEThree}{10}  % experiments/FT/k8s/job-ft-legs-w2-raunav.yaml :: job commands: samples per epoch at n_train 1000 // n_train :: sha256 f0962eddb451b7b0 :: passes over the training subset per epoch
\newcommand{\FtRecipeEpochsEFour}{50}  % experiments/FIGS/data/ft_recipes/recipes_w2b_bench_v2_v3.json :: runs[leg 1,2; n_train 10000].epochs :: sha256 7dc8c791343e1b6f
\newcommand{\FtRecipeStepsEFour}{19}  % experiments/FT/k8s/job-ft-legs-w2-raunav.yaml :: job commands: samples per epoch at n_train 10000 // batch size :: sha256 f0962eddb451b7b0 :: optimizer steps per epoch, as weaver floors them
\newcommand{\FtRecipePassesEFour}{one}  % experiments/FT/k8s/job-ft-legs-w2-raunav.yaml :: job commands: samples per epoch at n_train 10000 // n_train :: sha256 f0962eddb451b7b0 :: passes over the training subset per epoch
\newcommand{\FtRecipeEpochsEFive}{30}  % experiments/FIGS/data/ft_recipes/recipes_w2b_bench_v2_v3.json :: runs[leg 1,2; n_train 100000].epochs :: sha256 7dc8c791343e1b6f
\newcommand{\FtRecipeStepsEFive}{195}  % experiments/FT/k8s/job-ft-legs-w2-raunav.yaml :: job commands: samples per epoch at n_train 100000 // batch size :: sha256 f0962eddb451b7b0 :: optimizer steps per epoch, as weaver floors them
\newcommand{\FtRecipePassesEFive}{one}  % experiments/FT/k8s/job-ft-legs-w2-raunav.yaml :: job commands: samples per epoch at n_train 100000 // n_train :: sha256 f0962eddb451b7b0 :: passes over the training subset per epoch
\newcommand{\FtRecipeEpochsESix}{10}  % experiments/FIGS/data/ft_recipes/recipes_w2b_bench_v2_v3.json :: runs[leg 1,2; n_train 1000000].epochs :: sha256 7dc8c791343e1b6f
\newcommand{\FtRecipeStepsESix}{1{,}953}  % experiments/FT/k8s/job-ft-legs-w2-raunav.yaml :: job commands: samples per epoch at n_train 1000000 // batch size :: sha256 f0962eddb451b7b0 :: optimizer steps per epoch, as weaver floors them
\newcommand{\FtRecipePassesESix}{one}  % experiments/FT/k8s/job-ft-legs-w2-raunav.yaml :: job commands: samples per epoch at n_train 1000000 // n_train :: sha256 f0962eddb451b7b0 :: passes over the training subset per epoch
\newcommand{\FtRecipeValJets}{19{,}968}  % experiments/FT/k8s/job-ft-legs-w2-raunav.yaml :: --samples-per-epoch-val 20000 // batch x batch :: sha256 f0962eddb451b7b0 :: validation jets per epoch, floored to whole batches
\newcommand{\RecoveryAccOneeighteightByOneeighteight}{0.4123\,\ensuremath{\pm}\,0.0011}  % experiments/FIGS/data/label_recovery_curve_v1err/analysis/label_recovery_curve_summary.json :: table[rung=L188,level=188,probe=linear,n_train=max].accuracy :: sha256 1250d1f00537e0bd :: balanced accuracy, mean +- SD over 5 runs
\newcommand{\RecoveryAccOneeighteightByOnesixtwo}{0.4090\,\ensuremath{\pm}\,0.0012}  % experiments/FIGS/data/label_recovery_curve_v1err/analysis/label_recovery_curve_summary.json :: table[rung=L188,level=162,probe=linear,n_train=max].accuracy :: sha256 1250d1f00537e0bd :: balanced accuracy, mean +- SD over 5 runs
\newcommand{\RecoveryAccOneeighteightByFourthree}{0.356\,\ensuremath{\pm}\,0.008}  % experiments/FIGS/data/label_recovery_curve_v1err/analysis/label_recovery_curve_summary.json :: table[rung=L188,level=43,probe=linear,n_train=max].accuracy :: sha256 1250d1f00537e0bd :: balanced accuracy, mean +- SD over 5 runs
\newcommand{\RecoveryAccOneeighteightByOneseven}{0.306\,\ensuremath{\pm}\,0.008}  % experiments/FIGS/data/label_recovery_curve_v1err/analysis/label_recovery_curve_summary.json :: table[rung=L188,level=17,probe=linear,n_train=max].accuracy :: sha256 1250d1f00537e0bd :: balanced accuracy, mean +- SD over 5 runs
\newcommand{\RecoveryAccOnesixtwoByOneeighteight}{0.4623\,\ensuremath{\pm}\,0.0006}  % experiments/FIGS/data/label_recovery_curve_v1err/analysis/label_recovery_curve_summary.json :: table[rung=L162,level=188,probe=linear,n_train=max].accuracy :: sha256 1250d1f00537e0bd :: balanced accuracy, mean +- SD over 5 runs
\newcommand{\RecoveryAccOnesixtwoByOnesixtwo}{0.4622\,\ensuremath{\pm}\,0.0007}  % experiments/FIGS/data/label_recovery_curve_v1err/analysis/label_recovery_curve_summary.json :: table[rung=L162,level=162,probe=linear,n_train=max].accuracy :: sha256 1250d1f00537e0bd :: balanced accuracy, mean +- SD over 5 runs
\newcommand{\RecoveryAccOnesixtwoByFourthree}{0.402\,\ensuremath{\pm}\,0.009}  % experiments/FIGS/data/label_recovery_curve_v1err/analysis/label_recovery_curve_summary.json :: table[rung=L162,level=43,probe=linear,n_train=max].accuracy :: sha256 1250d1f00537e0bd :: balanced accuracy, mean +- SD over 5 runs
\newcommand{\RecoveryAccOnesixtwoByOneseven}{0.348\,\ensuremath{\pm}\,0.010}  % experiments/FIGS/data/label_recovery_curve_v1err/analysis/label_recovery_curve_summary.json :: table[rung=L162,level=17,probe=linear,n_train=max].accuracy :: sha256 1250d1f00537e0bd :: balanced accuracy, mean +- SD over 5 runs
\newcommand{\RecoveryAccSixfourByOneeighteight}{0.5979\,\ensuremath{\pm}\,0.0008}  % experiments/FIGS/data/label_recovery_curve_v1err/analysis/label_recovery_curve_summary.json :: table[rung=R63_Q1,level=188,probe=linear,n_train=max].accuracy :: sha256 1250d1f00537e0bd :: balanced accuracy, mean +- SD over 5 runs
\newcommand{\RecoveryAccSixfourByOnesixtwo}{0.5973\,\ensuremath{\pm}\,0.0013}  % experiments/FIGS/data/label_recovery_curve_v1err/analysis/label_recovery_curve_summary.json :: table[rung=R63_Q1,level=162,probe=linear,n_train=max].accuracy :: sha256 1250d1f00537e0bd :: balanced accuracy, mean +- SD over 5 runs
\newcommand{\RecoveryAccSixfourByFourthree}{0.5881\,\ensuremath{\pm}\,0.0007}  % experiments/FIGS/data/label_recovery_curve_v1err/analysis/label_recovery_curve_summary.json :: table[rung=R63_Q1,level=43,probe=linear,n_train=max].accuracy :: sha256 1250d1f00537e0bd :: balanced accuracy, mean +- SD over 5 runs
\newcommand{\RecoveryAccSixfourByOneseven}{0.531\,\ensuremath{\pm}\,0.006}  % experiments/FIGS/data/label_recovery_curve_v1err/analysis/label_recovery_curve_summary.json :: table[rung=R63_Q1,level=17,probe=linear,n_train=max].accuracy :: sha256 1250d1f00537e0bd :: balanced accuracy, mean +- SD over 5 runs
\newcommand{\RecoveryAccFourthreeByOneeighteight}{0.6770\,\ensuremath{\pm}\,0.0006}  % experiments/FIGS/data/label_recovery_curve_v1err/analysis/label_recovery_curve_summary.json :: table[rung=R42_Q1,level=188,probe=linear,n_train=max].accuracy :: sha256 1250d1f00537e0bd :: balanced accuracy, mean +- SD over 5 runs
\newcommand{\RecoveryAccFourthreeByOnesixtwo}{0.6766\,\ensuremath{\pm}\,0.0008}  % experiments/FIGS/data/label_recovery_curve_v1err/analysis/label_recovery_curve_summary.json :: table[rung=R42_Q1,level=162,probe=linear,n_train=max].accuracy :: sha256 1250d1f00537e0bd :: balanced accuracy, mean +- SD over 5 runs
\newcommand{\RecoveryAccFourthreeByFourthree}{0.6751\,\ensuremath{\pm}\,0.0011}  % experiments/FIGS/data/label_recovery_curve_v1err/analysis/label_recovery_curve_summary.json :: table[rung=R42_Q1,level=43,probe=linear,n_train=max].accuracy :: sha256 1250d1f00537e0bd :: balanced accuracy, mean +- SD over 5 runs
\newcommand{\RecoveryAccFourthreeByOneseven}{0.633\,\ensuremath{\pm}\,0.006}  % experiments/FIGS/data/label_recovery_curve_v1err/analysis/label_recovery_curve_summary.json :: table[rung=R42_Q1,level=17,probe=linear,n_train=max].accuracy :: sha256 1250d1f00537e0bd :: balanced accuracy, mean +- SD over 5 runs
\newcommand{\RecoveryAccThreezeroByOneeighteight}{0.7001\,\ensuremath{\pm}\,0.0005}  % experiments/FIGS/data/label_recovery_curve_v1err/analysis/label_recovery_curve_summary.json :: table[rung=R29_Q1,level=188,probe=linear,n_train=max].accuracy :: sha256 1250d1f00537e0bd :: balanced accuracy, mean +- SD over 5 runs
\newcommand{\RecoveryAccThreezeroByOnesixtwo}{0.7004\,\ensuremath{\pm}\,0.0007}  % experiments/FIGS/data/label_recovery_curve_v1err/analysis/label_recovery_curve_summary.json :: table[rung=R29_Q1,level=162,probe=linear,n_train=max].accuracy :: sha256 1250d1f00537e0bd :: balanced accuracy, mean +- SD over 5 runs
\newcommand{\RecoveryAccThreezeroByFourthree}{0.6998\,\ensuremath{\pm}\,0.0008}  % experiments/FIGS/data/label_recovery_curve_v1err/analysis/label_recovery_curve_summary.json :: table[rung=R29_Q1,level=43,probe=linear,n_train=max].accuracy :: sha256 1250d1f00537e0bd :: balanced accuracy, mean +- SD over 5 runs
\newcommand{\RecoveryAccThreezeroByOneseven}{0.6778\,\ensuremath{\pm}\,0.0030}  % experiments/FIGS/data/label_recovery_curve_v1err/analysis/label_recovery_curve_summary.json :: table[rung=R29_Q1,level=17,probe=linear,n_train=max].accuracy :: sha256 1250d1f00537e0bd :: balanced accuracy, mean +- SD over 5 runs
\newcommand{\RecoveryAccOnesevenByOneeighteight}{0.7675\,\ensuremath{\pm}\,0.0010}  % experiments/FIGS/data/label_recovery_curve_v1err/analysis/label_recovery_curve_summary.json :: table[rung=R16_Q1,level=188,probe=linear,n_train=max].accuracy :: sha256 1250d1f00537e0bd :: balanced accuracy, mean +- SD over 5 runs
\newcommand{\RecoveryAccOnesevenByOnesixtwo}{0.7681\,\ensuremath{\pm}\,0.0009}  % experiments/FIGS/data/label_recovery_curve_v1err/analysis/label_recovery_curve_summary.json :: table[rung=R16_Q1,level=162,probe=linear,n_train=max].accuracy :: sha256 1250d1f00537e0bd :: balanced accuracy, mean +- SD over 5 runs
\newcommand{\RecoveryAccOnesevenByFourthree}{0.7680\,\ensuremath{\pm}\,0.0010}  % experiments/FIGS/data/label_recovery_curve_v1err/analysis/label_recovery_curve_summary.json :: table[rung=R16_Q1,level=43,probe=linear,n_train=max].accuracy :: sha256 1250d1f00537e0bd :: balanced accuracy, mean +- SD over 5 runs
\newcommand{\RecoveryAccOnesevenByOneseven}{0.7681\,\ensuremath{\pm}\,0.0003}  % experiments/FIGS/data/label_recovery_curve_v1err/analysis/label_recovery_curve_summary.json :: table[rung=R16_Q1,level=17,probe=linear,n_train=max].accuracy :: sha256 1250d1f00537e0bd :: balanced accuracy, mean +- SD over 5 runs
\newcommand{\RecoveryAccFourByOneeighteight}{0.7526\,\ensuremath{\pm}\,0.0002}  % experiments/FIGS/data/label_recovery_curve_v1err/analysis/label_recovery_curve_summary.json :: table[rung=R3_VIS,level=188,probe=linear,n_train=max].accuracy :: sha256 1250d1f00537e0bd :: balanced accuracy, mean +- SD over 5 runs
\newcommand{\RecoveryAccFourByOnesixtwo}{0.7538\,\ensuremath{\pm}\,0.0017}  % experiments/FIGS/data/label_recovery_curve_v1err/analysis/label_recovery_curve_summary.json :: table[rung=R3_VIS,level=162,probe=linear,n_train=max].accuracy :: sha256 1250d1f00537e0bd :: balanced accuracy, mean +- SD over 5 runs
\newcommand{\RecoveryAccFourByFourthree}{0.7557\,\ensuremath{\pm}\,0.0015}  % experiments/FIGS/data/label_recovery_curve_v1err/analysis/label_recovery_curve_summary.json :: table[rung=R3_VIS,level=43,probe=linear,n_train=max].accuracy :: sha256 1250d1f00537e0bd :: balanced accuracy, mean +- SD over 5 runs
\newcommand{\RecoveryAccFourByOneseven}{0.7586\,\ensuremath{\pm}\,0.0008}  % experiments/FIGS/data/label_recovery_curve_v1err/analysis/label_recovery_curve_summary.json :: table[rung=R3_VIS,level=17,probe=linear,n_train=max].accuracy :: sha256 1250d1f00537e0bd :: balanced accuracy, mean +- SD over 5 runs
\newcommand{\RecoveryAccTwoByOneeighteight}{0.8688\,\ensuremath{\pm}\,0.0005}  % experiments/FIGS/data/label_recovery_curve_v1err/analysis/label_recovery_curve_summary.json :: table[rung=R1_Q1,level=188,probe=linear,n_train=max].accuracy :: sha256 1250d1f00537e0bd :: balanced accuracy, mean +- SD over 5 runs
\newcommand{\RecoveryAccTwoByOnesixtwo}{0.8703\,\ensuremath{\pm}\,0.0008}  % experiments/FIGS/data/label_recovery_curve_v1err/analysis/label_recovery_curve_summary.json :: table[rung=R1_Q1,level=162,probe=linear,n_train=max].accuracy :: sha256 1250d1f00537e0bd :: balanced accuracy, mean +- SD over 5 runs
\newcommand{\RecoveryAccTwoByFourthree}{0.8709\,\ensuremath{\pm}\,0.0011}  % experiments/FIGS/data/label_recovery_curve_v1err/analysis/label_recovery_curve_summary.json :: table[rung=R1_Q1,level=43,probe=linear,n_train=max].accuracy :: sha256 1250d1f00537e0bd :: balanced accuracy, mean +- SD over 5 runs
\newcommand{\RecoveryAccTwoByOneseven}{0.8713\,\ensuremath{\pm}\,0.0002}  % experiments/FIGS/data/label_recovery_curve_v1err/analysis/label_recovery_curve_summary.json :: table[rung=R1_Q1,level=17,probe=linear,n_train=max].accuracy :: sha256 1250d1f00537e0bd :: balanced accuracy, mean +- SD over 5 runs
\newcommand{\RecoveryOneeighteightOnesixtwoMaxAbs}{0.0033}  % experiments/FIGS/data/label_recovery_curve_v1err/analysis/label_recovery_curve_summary.json :: table[rung=*,level=188/162,probe=linear,n_train=max].accuracy :: sha256 1250d1f00537e0bd :: max |difference of the run means| over the tree
\newcommand{\RecoveryPairedOnesevenLevelOnesevenMinusOneeighteight}{+0.0006~[\ensuremath{-}0.0004, +0.0015]}  % experiments/FIGS/data/label_recovery_curve_v1err/analysis/label_recovery_curve_summary.json :: table[rung=R16_Q1,level=17 minus 188,probe=linear,n_train=max].accuracy :: sha256 1250d1f00537e0bd :: paired by run, mean difference [95% Student-t interval]
\newcommand{\RecoveryPairedOnesevenLevelOnesevenMinusOnesixtwo}{+0.00004~[\ensuremath{-}0.00123, +0.00131]}  % experiments/FIGS/data/label_recovery_curve_v1err/analysis/label_recovery_curve_summary.json :: table[rung=R16_Q1,level=17 minus 162,probe=linear,n_train=max].accuracy :: sha256 1250d1f00537e0bd :: paired by run, mean difference [95% Student-t interval]
\newcommand{\RecoveryPairedFourLevelOnesevenMinusOneeighteight}{+0.0060~[+0.0051, +0.0069]}  % experiments/FIGS/data/label_recovery_curve_v1err/analysis/label_recovery_curve_summary.json :: table[rung=R3_VIS,level=17 minus 188,probe=linear,n_train=max].accuracy :: sha256 1250d1f00537e0bd :: paired by run, mean difference [95% Student-t interval]
\newcommand{\RecoveryPairedFourLevelOnesevenMinusOnesixtwo}{+0.0048~[+0.0028, +0.0068]}  % experiments/FIGS/data/label_recovery_curve_v1err/analysis/label_recovery_curve_summary.json :: table[rung=R3_VIS,level=17 minus 162,probe=linear,n_train=max].accuracy :: sha256 1250d1f00537e0bd :: paired by run, mean difference [95% Student-t interval]
\newcommand{\RecoveryPairedTwoLevelOnesevenMinusOneeighteight}{+0.0025~[+0.0019, +0.0032]}  % experiments/FIGS/data/label_recovery_curve_v1err/analysis/label_recovery_curve_summary.json :: table[rung=R1_Q1,level=17 minus 188,probe=linear,n_train=max].accuracy :: sha256 1250d1f00537e0bd :: paired by run, mean difference [95% Student-t interval]
\newcommand{\RecoveryPairedTwoLevelOnesevenMinusOnesixtwo}{+0.00101~[+0.00003, +0.00198]}  % experiments/FIGS/data/label_recovery_curve_v1err/analysis/label_recovery_curve_summary.json :: table[rung=R1_Q1,level=17 minus 162,probe=linear,n_train=max].accuracy :: sha256 1250d1f00537e0bd :: paired by run, mean difference [95% Student-t interval]
\newcommand{\RecoveryMlpGainOneseven}{+0.0099~[+0.0078, +0.0120]}  % experiments/FIGS/data/label_recovery_curve_v1err/analysis/label_recovery_curve_summary.json :: mlp_minus_linear_at_largest[rung=L188,level=17].per_run :: sha256 1250d1f00537e0bd :: nonlinear minus linear probe, 188 native labels from the 17-class models
\newcommand{\RecoveryLastGainOneseven}{0.0125\,\ensuremath{\pm}\,0.0021}  % experiments/FIGS/data/label_recovery_curve_v1err/analysis/label_recovery_curve_summary.json :: data_gain_last_step[rung=L188,level=17] :: sha256 1250d1f00537e0bd :: balanced accuracy gained from 443000 to 1400000 training jets
\newcommand{\RecoveryLastGainOneeighteight}{0.0147\,\ensuremath{\pm}\,0.0014}  % experiments/FIGS/data/label_recovery_curve_v1err/analysis/label_recovery_curve_summary.json :: data_gain_last_step[rung=L188,level=188] :: sha256 1250d1f00537e0bd :: balanced accuracy gained from 443000 to 1400000 training jets
\newcommand{\RecoveryNTrain}{1{,}400{,}000}  % experiments/FIGS/data/label_recovery_curve_v1err/analysis/label_recovery_curve_summary.json :: sizes (max) :: sha256 1250d1f00537e0bd :: jets the linear probe is fit on at the end of the curve
\newcommand{\RecoveryNTrainMin}{14{,}000}  % experiments/FIGS/data/label_recovery_curve_v1err/analysis/label_recovery_curve_summary.json :: sizes (min) :: sha256 1250d1f00537e0bd
\newcommand{\RecoveryNTrainPrev}{443{,}000}  % experiments/FIGS/data/label_recovery_curve_v1err/analysis/label_recovery_curve_summary.json :: sizes (second largest) :: sha256 1250d1f00537e0bd
\newcommand{\RecoveryNTest}{400{,}000}  % experiments/FIGS/data/label_recovery_curve_v1err/l188-s1/label_recovery_curve.json :: n_test :: sha256 65651e4c06c7766f :: jets the probes are scored on, recorded in all 20 per-run files
\newcommand{\RecoveryMlpHidden}{512}  % experiments/FIGS/data/label_recovery_curve_v1err/l188-s1/label_recovery_curve.json :: mlp.hidden :: sha256 65651e4c06c7766f :: hidden units of the nonlinear probe
\newcommand{\RandNDraws}{three}  % experiments/FIGS/data/probe_ladder_randcontrol_mlp2/analysis/c4_random_control.json :: table[*].draw :: sha256 c675e609ff4fe3c9 :: random partitions drawn
\newcommand{\RandNDrawsCap}{Three}  % experiments/FIGS/data/probe_ladder_randcontrol_mlp2/analysis/c4_random_control.json :: table[*].draw :: sha256 c675e609ff4fe3c9 :: random partitions drawn, capitalised to open a sentence
\newcommand{\RandOmaBvcFourprongDrawOneLinear}{\ensuremath{1.05\times10^{-2}}}  % experiments/FIGS/data/probe_ladder_randcontrol_mlp2/analysis/c4_random_control.json :: table[0].control_log1m_auc :: sha256 c675e609ff4fe3c9 :: 1-AUC = exp of it, one run
\newcommand{\RandOmaBvcFourprongDrawTwoLinear}{\ensuremath{1.11\times10^{-2}}}  % experiments/FIGS/data/probe_ladder_randcontrol_mlp2/analysis/c4_random_control.json :: table[4].control_log1m_auc :: sha256 c675e609ff4fe3c9 :: 1-AUC = exp of it, one run
\newcommand{\RandOmaBvcFourprongDrawThreeLinear}{\ensuremath{9.98\times10^{-3}}}  % experiments/FIGS/data/probe_ladder_randcontrol_mlp2/analysis/c4_random_control.json :: table[8].control_log1m_auc :: sha256 c675e609ff4fe3c9 :: 1-AUC = exp of it, one run
\newcommand{\RandOmaMeanBvcFourprongLinear}{\ensuremath{(1.05 \pm 0.06)\times10^{-2}}}  % experiments/FIGS/data/probe_ladder_randcontrol_mlp2/analysis/c4_random_control.json :: table[task=bvc_4prong,probe=linear].control_log1m_auc :: sha256 c675e609ff4fe3c9 :: 1-AUC, mean +- SD over 3 draws
\newcommand{\RandOmaBvcFourprongDrawOneMlp}{\ensuremath{1.05\times10^{-2}}}  % experiments/FIGS/data/probe_ladder_randcontrol_mlp2/analysis/c4_random_control.json :: table[1].control_log1m_auc :: sha256 c675e609ff4fe3c9 :: 1-AUC = exp of it, one run
\newcommand{\RandOmaBvcFourprongDrawTwoMlp}{\ensuremath{1.04\times10^{-2}}}  % experiments/FIGS/data/probe_ladder_randcontrol_mlp2/analysis/c4_random_control.json :: table[5].control_log1m_auc :: sha256 c675e609ff4fe3c9 :: 1-AUC = exp of it, one run
\newcommand{\RandOmaBvcFourprongDrawThreeMlp}{\ensuremath{9.28\times10^{-3}}}  % experiments/FIGS/data/probe_ladder_randcontrol_mlp2/analysis/c4_random_control.json :: table[9].control_log1m_auc :: sha256 c675e609ff4fe3c9 :: 1-AUC = exp of it, one run
\newcommand{\RandOmaMeanBvcFourprongMlp}{\ensuremath{(1.01 \pm 0.07)\times10^{-2}}}  % experiments/FIGS/data/probe_ladder_randcontrol_mlp2/analysis/c4_random_control.json :: table[task=bvc_4prong,probe=mlp].control_log1m_auc :: sha256 c675e609ff4fe3c9 :: 1-AUC, mean +- SD over 3 draws
\newcommand{\RandOmaVisibleContentDrawOneLinear}{\ensuremath{1.25\times10^{-3}}}  % experiments/FIGS/data/probe_ladder_randcontrol_mlp2/analysis/c4_random_control.json :: table[2].control_log1m_auc :: sha256 c675e609ff4fe3c9 :: 1-AUC = exp of it, one run
\newcommand{\RandOmaVisibleContentDrawTwoLinear}{\ensuremath{9.89\times10^{-4}}}  % experiments/FIGS/data/probe_ladder_randcontrol_mlp2/analysis/c4_random_control.json :: table[6].control_log1m_auc :: sha256 c675e609ff4fe3c9 :: 1-AUC = exp of it, one run
\newcommand{\RandOmaVisibleContentDrawThreeLinear}{\ensuremath{1.08\times10^{-3}}}  % experiments/FIGS/data/probe_ladder_randcontrol_mlp2/analysis/c4_random_control.json :: table[10].control_log1m_auc :: sha256 c675e609ff4fe3c9 :: 1-AUC = exp of it, one run
\newcommand{\RandOmaMeanVisibleContentLinear}{\ensuremath{(1.11 \pm 0.13)\times10^{-3}}}  % experiments/FIGS/data/probe_ladder_randcontrol_mlp2/analysis/c4_random_control.json :: table[task=visible_content,probe=linear].control_log1m_auc :: sha256 c675e609ff4fe3c9 :: 1-AUC, mean +- SD over 3 draws
\newcommand{\RandOmaVisibleContentDrawOneMlp}{\ensuremath{1.21\times10^{-3}}}  % experiments/FIGS/data/probe_ladder_randcontrol_mlp2/analysis/c4_random_control.json :: table[3].control_log1m_auc :: sha256 c675e609ff4fe3c9 :: 1-AUC = exp of it, one run
\newcommand{\RandOmaVisibleContentDrawTwoMlp}{\ensuremath{1.07\times10^{-3}}}  % experiments/FIGS/data/probe_ladder_randcontrol_mlp2/analysis/c4_random_control.json :: table[7].control_log1m_auc :: sha256 c675e609ff4fe3c9 :: 1-AUC = exp of it, one run
\newcommand{\RandOmaVisibleContentDrawThreeMlp}{\ensuremath{1.13\times10^{-3}}}  % experiments/FIGS/data/probe_ladder_randcontrol_mlp2/analysis/c4_random_control.json :: table[11].control_log1m_auc :: sha256 c675e609ff4fe3c9 :: 1-AUC = exp of it, one run
\newcommand{\RandOmaMeanVisibleContentMlp}{\ensuremath{(1.14 \pm 0.07)\times10^{-3}}}  % experiments/FIGS/data/probe_ladder_randcontrol_mlp2/analysis/c4_random_control.json :: table[task=visible_content,probe=mlp].control_log1m_auc :: sha256 c675e609ff4fe3c9 :: 1-AUC, mean +- SD over 3 draws
\newcommand{\RandCostFactorBvcFourprongLinear}{1.49}  % experiments/FIGS/data/probe_ladder_randcontrol_mlp2/analysis/c4_random_control.json :: grouping_cost_post_hoc.tasks.bvc_4prong.linear.factor_control :: sha256 c675e609ff4fe3c9 :: POST HOC: random grouping vs finer models, in 1-AUC
\newcommand{\RandSemCostFactorBvcFourprongLinear}{4.42}  % experiments/FIGS/data/probe_ladder_randcontrol_mlp2/analysis/c4_random_control.json :: grouping_cost_post_hoc.tasks.bvc_4prong.linear.factor_semantic :: sha256 c675e609ff4fe3c9 :: POST HOC: semantic 17-class vs finer models
\newcommand{\RandCostFactorBvcFourprongMlp}{1.41}  % experiments/FIGS/data/probe_ladder_randcontrol_mlp2/analysis/c4_random_control.json :: grouping_cost_post_hoc.tasks.bvc_4prong.mlp.factor_control :: sha256 c675e609ff4fe3c9 :: POST HOC: random grouping vs finer models, in 1-AUC
\newcommand{\RandSemCostFactorBvcFourprongMlp}{3.26}  % experiments/FIGS/data/probe_ladder_randcontrol_mlp2/analysis/c4_random_control.json :: grouping_cost_post_hoc.tasks.bvc_4prong.mlp.factor_semantic :: sha256 c675e609ff4fe3c9 :: POST HOC: semantic 17-class vs finer models
\newcommand{\RandCostFactorVisibleContentLinear}{1.35}  % experiments/FIGS/data/probe_ladder_randcontrol_mlp2/analysis/c4_random_control.json :: grouping_cost_post_hoc.tasks.visible_content.linear.factor_control :: sha256 c675e609ff4fe3c9 :: POST HOC: random grouping vs finer models, in 1-AUC
\newcommand{\RandSemCostFactorVisibleContentLinear}{2.82}  % experiments/FIGS/data/probe_ladder_randcontrol_mlp2/analysis/c4_random_control.json :: grouping_cost_post_hoc.tasks.visible_content.linear.factor_semantic :: sha256 c675e609ff4fe3c9 :: POST HOC: semantic 17-class vs finer models
\newcommand{\RandCostFactorVisibleContentMlp}{1.33}  % experiments/FIGS/data/probe_ladder_randcontrol_mlp2/analysis/c4_random_control.json :: grouping_cost_post_hoc.tasks.visible_content.mlp.factor_control :: sha256 c675e609ff4fe3c9 :: POST HOC: random grouping vs finer models, in 1-AUC
\newcommand{\RandSemCostFactorVisibleContentMlp}{2.37}  % experiments/FIGS/data/probe_ladder_randcontrol_mlp2/analysis/c4_random_control.json :: grouping_cost_post_hoc.tasks.visible_content.mlp.factor_semantic :: sha256 c675e609ff4fe3c9 :: POST HOC: semantic 17-class vs finer models
\newcommand{\MassResNTest}{398{,}625}  % experiments/FIGS/data/mass_resolution/s1.json :: centering_detail.split[2] :: sha256 88a57bbad270d51c :: test jets sigma_eff is computed on, the same in every per-seed file and arm
\newcommand{\MassResNClasses}{184}  % experiments/FIGS/data/mass_resolution/analysis_holm/s7_mass_resolution.json :: provenance.n_classes_used :: sha256 e03cde45ece9f385 :: native classes with enough training jets to centre
\newcommand{\MassResSigmaEffRidgeOneeighteight}{0.3393\,\ensuremath{\pm}\,0.0019}  % experiments/FIGS/data/mass_resolution/analysis_holm/s7_mass_resolution.json :: table[cell=188,probe=ridge].sigma_eff :: sha256 e03cde45ece9f385 :: mean +- SD over 5 seeds
\newcommand{\MassResSigmaEffRidgeOnesixtwo}{0.3302\,\ensuremath{\pm}\,0.0029}  % experiments/FIGS/data/mass_resolution/analysis_holm/s7_mass_resolution.json :: table[cell=162,probe=ridge].sigma_eff :: sha256 e03cde45ece9f385 :: mean +- SD over 5 seeds
\newcommand{\MassResSigmaEffRidgeFourthree}{0.3228\,\ensuremath{\pm}\,0.0019}  % experiments/FIGS/data/mass_resolution/analysis_holm/s7_mass_resolution.json :: table[cell=43,probe=ridge].sigma_eff :: sha256 e03cde45ece9f385 :: mean +- SD over 5 seeds
\newcommand{\MassResSigmaEffRidgeOneseven}{0.300\,\ensuremath{\pm}\,0.010}  % experiments/FIGS/data/mass_resolution/analysis_holm/s7_mass_resolution.json :: table[cell=17,probe=ridge].sigma_eff :: sha256 e03cde45ece9f385 :: mean +- SD over 5 seeds
\newcommand{\MassResSigmaEffRidgeOnesixtwoMass}{0.2238\,\ensuremath{\pm}\,0.0019}  % experiments/FIGS/data/mass_resolution/analysis_holm/s7_mass_resolution.json :: table[cell=162+mass,probe=ridge].sigma_eff :: sha256 e03cde45ece9f385 :: mean +- SD over 5 seeds
\newcommand{\MassResSigmaEffRidgeOnesevenMass}{0.2170\,\ensuremath{\pm}\,0.0034}  % experiments/FIGS/data/mass_resolution/analysis_holm/s7_mass_resolution.json :: table[cell=17+mass,probe=ridge].sigma_eff :: sha256 e03cde45ece9f385 :: mean +- SD over 5 seeds
\newcommand{\MassResGainRidgeOnesixtwo}{\ensuremath{-}0.1063}  % experiments/FIGS/data/mass_resolution/analysis_holm/s7_mass_resolution.json :: table[cell=162+mass/162,probe=ridge].sigma_eff :: sha256 e03cde45ece9f385 :: sigma_eff with the mass output minus without, difference of the seed means; negative is better
\newcommand{\MassResGainRidgeOneseven}{\ensuremath{-}0.0831}  % experiments/FIGS/data/mass_resolution/analysis_holm/s7_mass_resolution.json :: table[cell=17+mass/17,probe=ridge].sigma_eff :: sha256 e03cde45ece9f385 :: sigma_eff with the mass output minus without, difference of the seed means; negative is better
\newcommand{\MassResDidRidge}{\ensuremath{-}0.0233}  % experiments/FIGS/data/mass_resolution/analysis_holm/s7_mass_resolution.json :: table[cell=162+mass/162/17+mass/17,probe=ridge].sigma_eff :: sha256 e03cde45ece9f385 :: gain at 162 classes minus gain at 17
\newcommand{\MassResSigmaEffMlpOneeighteight}{0.2463\,\ensuremath{\pm}\,0.0022}  % experiments/FIGS/data/mass_resolution/analysis_holm/s7_mass_resolution.json :: table[cell=188,probe=mlp].sigma_eff :: sha256 e03cde45ece9f385 :: mean +- SD over 5 seeds
\newcommand{\MassResSigmaEffMlpOnesixtwo}{0.2438\,\ensuremath{\pm}\,0.0018}  % experiments/FIGS/data/mass_resolution/analysis_holm/s7_mass_resolution.json :: table[cell=162,probe=mlp].sigma_eff :: sha256 e03cde45ece9f385 :: mean +- SD over 5 seeds
\newcommand{\MassResSigmaEffMlpFourthree}{0.2407\,\ensuremath{\pm}\,0.0009}  % experiments/FIGS/data/mass_resolution/analysis_holm/s7_mass_resolution.json :: table[cell=43,probe=mlp].sigma_eff :: sha256 e03cde45ece9f385 :: mean +- SD over 5 seeds
\newcommand{\MassResSigmaEffMlpOneseven}{0.2358\,\ensuremath{\pm}\,0.0025}  % experiments/FIGS/data/mass_resolution/analysis_holm/s7_mass_resolution.json :: table[cell=17,probe=mlp].sigma_eff :: sha256 e03cde45ece9f385 :: mean +- SD over 5 seeds
\newcommand{\MassResSigmaEffMlpOnesixtwoMass}{0.1943\,\ensuremath{\pm}\,0.0001}  % experiments/FIGS/data/mass_resolution/analysis_holm/s7_mass_resolution.json :: table[cell=162+mass,probe=mlp].sigma_eff :: sha256 e03cde45ece9f385 :: mean +- SD over 5 seeds
\newcommand{\MassResSigmaEffMlpOnesevenMass}{0.1949\,\ensuremath{\pm}\,0.0006}  % experiments/FIGS/data/mass_resolution/analysis_holm/s7_mass_resolution.json :: table[cell=17+mass,probe=mlp].sigma_eff :: sha256 e03cde45ece9f385 :: mean +- SD over 5 seeds
\newcommand{\MassResGainMlpOnesixtwo}{\ensuremath{-}0.0496}  % experiments/FIGS/data/mass_resolution/analysis_holm/s7_mass_resolution.json :: table[cell=162+mass/162,probe=mlp].sigma_eff :: sha256 e03cde45ece9f385 :: sigma_eff with the mass output minus without, difference of the seed means; negative is better
\newcommand{\MassResGainMlpOneseven}{\ensuremath{-}0.0409}  % experiments/FIGS/data/mass_resolution/analysis_holm/s7_mass_resolution.json :: table[cell=17+mass/17,probe=mlp].sigma_eff :: sha256 e03cde45ece9f385 :: sigma_eff with the mass output minus without, difference of the seed means; negative is better
\newcommand{\MassResDidMlp}{\ensuremath{-}0.0087}  % experiments/FIGS/data/mass_resolution/analysis_holm/s7_mass_resolution.json :: table[cell=162+mass/162/17+mass/17,probe=mlp].sigma_eff :: sha256 e03cde45ece9f385 :: gain at 162 classes minus gain at 17
\newcommand{\MassResTargetSigmaEff}{0.2986}  % experiments/FIGS/data/mass_resolution/analysis_holm/s7_mass_resolution.json :: table[*].target_sigma_eff :: sha256 e03cde45ece9f385 :: sigma_eff of the class-centred target itself (a model that predicts the class mean)
\newcommand{\MassResRidgeRtwoMin}{20\%}  % experiments/FIGS/data/mass_resolution/s1.json :: arms.{l188,l162,r42q1,r16q1}-s*.ridge.val_r2 :: sha256 88a57bbad270d51c :: validation R^2 of the ridge probe, smallest of 20 models without the mass output, all 5 per-seed files
\newcommand{\MassResRidgeRtwoMax}{30\%}  % experiments/FIGS/data/mass_resolution/s1.json :: arms.{l188,l162,r42q1,r16q1}-s*.ridge.val_r2 :: sha256 88a57bbad270d51c :: validation R^2 of the ridge probe, largest of 20 models without the mass output, all 5 per-seed files
\newcommand{\MassResSigmaEffOnesixtwo}{0.3302}  % experiments/FIGS/data/mass_resolution/analysis_holm/s7_mass_resolution.json :: table[cell=162,probe=ridge].sigma_eff :: sha256 e03cde45ece9f385 :: mean over 5 seeds
\newcommand{\MassResSdOnesixtwo}{0.626}  % experiments/FIGS/data/mass_resolution/analysis_holm/s7_mass_resolution.json :: table[cell=162,probe=ridge].sd :: sha256 e03cde45ece9f385 :: mean over 5 seeds
\newcommand{\MassResTailFractionOnesixtwo}{0.062}  % experiments/FIGS/data/mass_resolution/analysis_holm/s7_mass_resolution.json :: table[cell=162,probe=ridge].tail_fraction :: sha256 e03cde45ece9f385 :: mean over 5 seeds
\newcommand{\MassResSigmaEffOnesixtwoMass}{0.2238}  % experiments/FIGS/data/mass_resolution/analysis_holm/s7_mass_resolution.json :: table[cell=162+mass,probe=ridge].sigma_eff :: sha256 e03cde45ece9f385 :: mean over 5 seeds
\newcommand{\MassResSdOnesixtwoMass}{0.494}  % experiments/FIGS/data/mass_resolution/analysis_holm/s7_mass_resolution.json :: table[cell=162+mass,probe=ridge].sd :: sha256 e03cde45ece9f385 :: mean over 5 seeds
\newcommand{\MassResTailFractionOnesixtwoMass}{0.034}  % experiments/FIGS/data/mass_resolution/analysis_holm/s7_mass_resolution.json :: table[cell=162+mass,probe=ridge].tail_fraction :: sha256 e03cde45ece9f385 :: mean over 5 seeds
\newcommand{\MassResSigmaEffOneseven}{0.3001}  % experiments/FIGS/data/mass_resolution/analysis_holm/s7_mass_resolution.json :: table[cell=17,probe=ridge].sigma_eff :: sha256 e03cde45ece9f385 :: mean over 5 seeds
\newcommand{\MassResSdOneseven}{0.601}  % experiments/FIGS/data/mass_resolution/analysis_holm/s7_mass_resolution.json :: table[cell=17,probe=ridge].sd :: sha256 e03cde45ece9f385 :: mean over 5 seeds
\newcommand{\MassResTailFractionOneseven}{0.056}  % experiments/FIGS/data/mass_resolution/analysis_holm/s7_mass_resolution.json :: table[cell=17,probe=ridge].tail_fraction :: sha256 e03cde45ece9f385 :: mean over 5 seeds
\newcommand{\MassResSigmaEffOnesevenMass}{0.2170}  % experiments/FIGS/data/mass_resolution/analysis_holm/s7_mass_resolution.json :: table[cell=17+mass,probe=ridge].sigma_eff :: sha256 e03cde45ece9f385 :: mean over 5 seeds
\newcommand{\MassResSdOnesevenMass}{0.490}  % experiments/FIGS/data/mass_resolution/analysis_holm/s7_mass_resolution.json :: table[cell=17+mass,probe=ridge].sd :: sha256 e03cde45ece9f385 :: mean over 5 seeds
\newcommand{\MassResTailFractionOnesevenMass}{0.033}  % experiments/FIGS/data/mass_resolution/analysis_holm/s7_mass_resolution.json :: table[cell=17+mass,probe=ridge].tail_fraction :: sha256 e03cde45ece9f385 :: mean over 5 seeds
\newcommand{\MassResSigmaEffOneeighteight}{0.3393}  % experiments/FIGS/data/mass_resolution/analysis_holm/s7_mass_resolution.json :: table[cell=188,probe=ridge].sigma_eff :: sha256 e03cde45ece9f385 :: mean over 5 seeds
\newcommand{\MassResSdOneeighteight}{0.631}  % experiments/FIGS/data/mass_resolution/analysis_holm/s7_mass_resolution.json :: table[cell=188,probe=ridge].sd :: sha256 e03cde45ece9f385 :: mean over 5 seeds
\newcommand{\MassResTailFractionOneeighteight}{0.064}  % experiments/FIGS/data/mass_resolution/analysis_holm/s7_mass_resolution.json :: table[cell=188,probe=ridge].tail_fraction :: sha256 e03cde45ece9f385 :: mean over 5 seeds
\newcommand{\MassResSigmaEffFourthree}{0.3228}  % experiments/FIGS/data/mass_resolution/analysis_holm/s7_mass_resolution.json :: table[cell=43,probe=ridge].sigma_eff :: sha256 e03cde45ece9f385 :: mean over 5 seeds
\newcommand{\MassResSdFourthree}{0.618}  % experiments/FIGS/data/mass_resolution/analysis_holm/s7_mass_resolution.json :: table[cell=43,probe=ridge].sd :: sha256 e03cde45ece9f385 :: mean over 5 seeds
\newcommand{\MassResTailFractionFourthree}{0.060}  % experiments/FIGS/data/mass_resolution/analysis_holm/s7_mass_resolution.json :: table[cell=43,probe=ridge].tail_fraction :: sha256 e03cde45ece9f385 :: mean over 5 seeds
\newcommand{\AojEffPercent}{1\%}  % experiments/FIGS/data/aoj_full_v1/fit_v6/results.json :: eff :: sha256 2416fc21f93fcfa6 :: the score cut's pass fraction outside the mass windows
\newcommand{\AojPtMin}{500}  % experiments/AOJ/peak_fit.py :: PT_RANGE[0] :: sha256 cad06229a74532be :: GeV
\newcommand{\AojPtMax}{2{,}500}  % experiments/AOJ/peak_fit.py :: PT_RANGE[1] :: sha256 cad06229a74532be :: GeV
\newcommand{\AojRhoMin}{\ensuremath{-}5.5}  % experiments/AOJ/peak_fit.py :: RHO_RANGE[0] :: sha256 cad06229a74532be
\newcommand{\AojRhoMax}{\ensuremath{-}2.0}  % experiments/AOJ/peak_fit.py :: RHO_RANGE[1] :: sha256 cad06229a74532be
\newcommand{\AojWindowWLo}{65}  % experiments/AOJ/peak_fit.py :: PEAKS[W].window[0] :: sha256 cad06229a74532be :: GeV
\newcommand{\AojWindowWHi}{105}  % experiments/AOJ/peak_fit.py :: PEAKS[W].window[1] :: sha256 cad06229a74532be :: GeV
\newcommand{\AojWindowTopLo}{140}  % experiments/AOJ/peak_fit.py :: PEAKS[top].window[0] :: sha256 cad06229a74532be :: GeV
\newcommand{\AojWindowTopHi}{220}  % experiments/AOJ/peak_fit.py :: PEAKS[top].window[1] :: sha256 cad06229a74532be :: GeV
\newcommand{\AojFitLo}{105}  % experiments/AOJ/peak_fit.py :: PEAKS[top].fit_range[0] :: sha256 cad06229a74532be :: GeV
\newcommand{\AojFitHi}{300}  % experiments/AOJ/peak_fit.py :: PEAKS[top].fit_range[1] :: sha256 cad06229a74532be :: GeV
\newcommand{\AojMassBin}{5}  % experiments/AOJ/peak_fit.py :: MASS_BIN :: sha256 cad06229a74532be :: GeV
\newcommand{\AojMapOrderRho}{3}  % experiments/AOJ/peak_fit.py :: MAP_ORDER[0] :: sha256 cad06229a74532be :: polynomial order in rho
\newcommand{\AojMapOrderPt}{2}  % experiments/AOJ/peak_fit.py :: MAP_ORDER[1] :: sha256 cad06229a74532be :: polynomial order in ln pT
\newcommand{\AojNRunsWord}{five}  % experiments/FIGS/data/aoj_full_v1/analysis_v6/aoj_top.json :: per_label_set.label_sets.*.signal_yield.n :: sha256 5dfef902a5fe75ba :: pretraining runs per vocabulary
\newcommand{\AojNJetsRhoWindow}{9{,}441{,}209}  % experiments/FIGS/data/aoj_full_v1/fit_v6/results.json :: n_jets :: sha256 2416fc21f93fcfa6 :: staged jets inside the rho window, all files; not the jets in the mass fit
\newcommand{\AojYieldOneeighteight}{2{,}400\,\ensuremath{\pm}\,500}  % experiments/FIGS/data/aoj_full_v1/analysis_v6/aoj_top.json :: per_label_set.label_sets.188.signal_yield :: sha256 5dfef902a5fe75ba :: mean +- SD over 5 runs, shared shape
\newcommand{\AojStatErrOneeighteight}{211}  % experiments/FIGS/data/aoj_full_v1/analysis_v6/aoj_top.json :: per_label_set.label_sets.188.median_stat_err :: sha256 5dfef902a5fe75ba :: median per-fit statistical error
\newcommand{\AojYieldOwnShapeOneeighteight}{2{,}400\,\ensuremath{\pm}\,600}  % experiments/FIGS/data/aoj_full_v1/analysis_v6/aoj_top.json :: per_label_set.label_sets.188.by_shape.own_shape :: sha256 5dfef902a5fe75ba :: every fit with its own floated shape
\newcommand{\AojSpreadOverStatOneeighteight}{2.4}  % experiments/FIGS/data/aoj_full_v1/analysis_v6/aoj_top.json :: per_label_set.label_sets.188.signal_yield.sd / median_stat_err :: sha256 5dfef902a5fe75ba :: spread between runs over the median per-fit statistical error
\newcommand{\AojYieldOnesixtwo}{2{,}400\,\ensuremath{\pm}\,900}  % experiments/FIGS/data/aoj_full_v1/analysis_v6/aoj_top.json :: per_label_set.label_sets.162.signal_yield :: sha256 5dfef902a5fe75ba :: mean +- SD over 5 runs, shared shape
\newcommand{\AojStatErrOnesixtwo}{199}  % experiments/FIGS/data/aoj_full_v1/analysis_v6/aoj_top.json :: per_label_set.label_sets.162.median_stat_err :: sha256 5dfef902a5fe75ba :: median per-fit statistical error
\newcommand{\AojYieldOwnShapeOnesixtwo}{2{,}200\,\ensuremath{\pm}\,800}  % experiments/FIGS/data/aoj_full_v1/analysis_v6/aoj_top.json :: per_label_set.label_sets.162.by_shape.own_shape :: sha256 5dfef902a5fe75ba :: every fit with its own floated shape
\newcommand{\AojSpreadOverStatOnesixtwo}{4.5}  % experiments/FIGS/data/aoj_full_v1/analysis_v6/aoj_top.json :: per_label_set.label_sets.162.signal_yield.sd / median_stat_err :: sha256 5dfef902a5fe75ba :: spread between runs over the median per-fit statistical error
\newcommand{\AojYieldFourthree}{2{,}400\,\ensuremath{\pm}\,800}  % experiments/FIGS/data/aoj_full_v1/analysis_v6/aoj_top.json :: per_label_set.label_sets.43.signal_yield :: sha256 5dfef902a5fe75ba :: mean +- SD over 5 runs, shared shape
\newcommand{\AojStatErrFourthree}{184}  % experiments/FIGS/data/aoj_full_v1/analysis_v6/aoj_top.json :: per_label_set.label_sets.43.median_stat_err :: sha256 5dfef902a5fe75ba :: median per-fit statistical error
\newcommand{\AojYieldOwnShapeFourthree}{2{,}300\,\ensuremath{\pm}\,800}  % experiments/FIGS/data/aoj_full_v1/analysis_v6/aoj_top.json :: per_label_set.label_sets.43.by_shape.own_shape :: sha256 5dfef902a5fe75ba :: every fit with its own floated shape
\newcommand{\AojSpreadOverStatFourthree}{4.2}  % experiments/FIGS/data/aoj_full_v1/analysis_v6/aoj_top.json :: per_label_set.label_sets.43.signal_yield.sd / median_stat_err :: sha256 5dfef902a5fe75ba :: spread between runs over the median per-fit statistical error
\newcommand{\AojYieldOneseven}{2{,}350\,\ensuremath{\pm}\,240}  % experiments/FIGS/data/aoj_full_v1/analysis_v6/aoj_top.json :: per_label_set.label_sets.17.signal_yield :: sha256 5dfef902a5fe75ba :: mean +- SD over 5 runs, shared shape
\newcommand{\AojStatErrOneseven}{201}  % experiments/FIGS/data/aoj_full_v1/analysis_v6/aoj_top.json :: per_label_set.label_sets.17.median_stat_err :: sha256 5dfef902a5fe75ba :: median per-fit statistical error
\newcommand{\AojYieldOwnShapeOneseven}{2{,}480\,\ensuremath{\pm}\,290}  % experiments/FIGS/data/aoj_full_v1/analysis_v6/aoj_top.json :: per_label_set.label_sets.17.by_shape.own_shape :: sha256 5dfef902a5fe75ba :: every fit with its own floated shape
\newcommand{\AojSpreadOverStatOneseven}{1.2}  % experiments/FIGS/data/aoj_full_v1/analysis_v6/aoj_top.json :: per_label_set.label_sets.17.signal_yield.sd / median_stat_err :: sha256 5dfef902a5fe75ba :: spread between runs over the median per-fit statistical error
\newcommand{\AojYieldOnesixtwoMass}{2{,}420\,\ensuremath{\pm}\,240}  % experiments/FIGS/data/aoj_full_v1/analysis_v6/aoj_top.json :: per_label_set.label_sets.162+mass.signal_yield :: sha256 5dfef902a5fe75ba :: mean +- SD over 5 runs, shared shape
\newcommand{\AojStatErrOnesixtwoMass}{179}  % experiments/FIGS/data/aoj_full_v1/analysis_v6/aoj_top.json :: per_label_set.label_sets.162+mass.median_stat_err :: sha256 5dfef902a5fe75ba :: median per-fit statistical error
\newcommand{\AojYieldOwnShapeOnesixtwoMass}{2{,}350\,\ensuremath{\pm}\,320}  % experiments/FIGS/data/aoj_full_v1/analysis_v6/aoj_top.json :: per_label_set.label_sets.162+mass.by_shape.own_shape :: sha256 5dfef902a5fe75ba :: every fit with its own floated shape
\newcommand{\AojSpreadOverStatOnesixtwoMass}{1.3}  % experiments/FIGS/data/aoj_full_v1/analysis_v6/aoj_top.json :: per_label_set.label_sets.162+mass.signal_yield.sd / median_stat_err :: sha256 5dfef902a5fe75ba :: spread between runs over the median per-fit statistical error
\newcommand{\AojYieldOnesevenMass}{2{,}500\,\ensuremath{\pm}\,600}  % experiments/FIGS/data/aoj_full_v1/analysis_v6/aoj_top.json :: per_label_set.label_sets.17+mass.signal_yield :: sha256 5dfef902a5fe75ba :: mean +- SD over 5 runs, shared shape
\newcommand{\AojStatErrOnesevenMass}{200}  % experiments/FIGS/data/aoj_full_v1/analysis_v6/aoj_top.json :: per_label_set.label_sets.17+mass.median_stat_err :: sha256 5dfef902a5fe75ba :: median per-fit statistical error
\newcommand{\AojYieldOwnShapeOnesevenMass}{2{,}900\,\ensuremath{\pm}\,500}  % experiments/FIGS/data/aoj_full_v1/analysis_v6/aoj_top.json :: per_label_set.label_sets.17+mass.by_shape.own_shape :: sha256 5dfef902a5fe75ba :: every fit with its own floated shape
\newcommand{\AojSpreadOverStatOnesevenMass}{2.9}  % experiments/FIGS/data/aoj_full_v1/analysis_v6/aoj_top.json :: per_label_set.label_sets.17+mass.signal_yield.sd / median_stat_err :: sha256 5dfef902a5fe75ba :: spread between runs over the median per-fit statistical error
\newcommand{\AojOwnShapeShiftMin}{\ensuremath{-}7\%}  % experiments/FIGS/data/aoj_full_v1/analysis_v6/aoj_top.json :: per_label_set.label_sets.*.by_shape.own_shape.mean / signal_yield.mean - 1 :: sha256 5dfef902a5fe75ba :: the vocabulary means with each model's own shape, relative to the shared shape
\newcommand{\AojOwnShapeShiftMax}{+13\%}  % experiments/FIGS/data/aoj_full_v1/analysis_v6/aoj_top.json :: per_label_set.label_sets.*.by_shape.own_shape.mean / signal_yield.mean - 1 :: sha256 5dfef902a5fe75ba :: the vocabulary means with each model's own shape, relative to the shared shape
\newcommand{\AojPooledMean}{182.8}  % experiments/FIGS/data/aoj_full_v1/analysis_v6/aoj_top.json :: per_label_set.shapes.pooled :: sha256 5dfef902a5fe75ba :: GeV
\newcommand{\AojPooledWidth}{11.7}  % experiments/FIGS/data/aoj_full_v1/analysis_v6/aoj_top.json :: per_label_set.shapes.pooled :: sha256 5dfef902a5fe75ba :: GeV
\newcommand{\AojPoolSize}{30}  % experiments/FIGS/data/aoj_full_v1/analysis_v6/aoj_top.json :: per_label_set.shapes.pool (count) :: sha256 5dfef902a5fe75ba :: pretrained models the shared shape is fitted over
\newcommand{\AojNModels}{30}  % experiments/FIGS/data/aoj_full_v1/fit_v6/results.json :: models (pretrained) :: sha256 2416fc21f93fcfa6 :: fits
\newcommand{\AojOwnMeanMin}{181.6}  % experiments/FIGS/data/aoj_full_v1/fit_v6/results.json :: models.*.top.own_shape_fit.mean (min) :: sha256 2416fc21f93fcfa6 :: GeV
\newcommand{\AojOwnMeanMax}{184.0}  % experiments/FIGS/data/aoj_full_v1/fit_v6/results.json :: models.*.top.own_shape_fit.mean (max) :: sha256 2416fc21f93fcfa6 :: GeV
\newcommand{\AojOwnWidthMin}{10.0}  % experiments/FIGS/data/aoj_full_v1/fit_v6/results.json :: models.*.top.own_shape_fit.width (min) :: sha256 2416fc21f93fcfa6 :: GeV
\newcommand{\AojOwnWidthMax}{13.8}  % experiments/FIGS/data/aoj_full_v1/fit_v6/results.json :: models.*.top.own_shape_fit.width (max) :: sha256 2416fc21f93fcfa6 :: GeV
\newcommand{\AojShapeStepMean}{0.7}  % experiments/FIGS/data/aoj_full_v1/fit_v6/results.json :: pooled_shape.systematic_step.mean :: sha256 2416fc21f93fcfa6 :: GeV: the SD of the models' own floated values, one shift for every model together
\newcommand{\AojShapeSystMeanMax}{1\%}  % experiments/FIGS/data/aoj_full_v1/fit_v6/results.json :: models.*.top.shape_systematic.pooled_mean_(down|up) / signal_yield - 1 (largest |.|) :: sha256 2416fc21f93fcfa6 :: the largest yield change of any model under that shift
\newcommand{\AojShapeStepWidth}{1.0}  % experiments/FIGS/data/aoj_full_v1/fit_v6/results.json :: pooled_shape.systematic_step.width :: sha256 2416fc21f93fcfa6 :: GeV: the SD of the models' own floated values, one shift for every model together
\newcommand{\AojShapeSystWidthMax}{12\%}  % experiments/FIGS/data/aoj_full_v1/fit_v6/results.json :: models.*.top.shape_systematic.pooled_width_(down|up) / signal_yield - 1 (largest |.|) :: sha256 2416fc21f93fcfa6 :: the largest yield change of any model under that shift
\newcommand{\AojLeakEpsHi}{0.6}  % experiments/FIGS/data/aoj_full_v1/fit_v6/results.json :: models.*.top.leak_systematic.eps_ref (max) :: sha256 2416fc21f93fcfa6 :: the reference's assumed efficiency for data tops, upper variation
\newcommand{\AojLeakEpsLo}{0.4}  % experiments/FIGS/data/aoj_full_v1/fit_v6/results.json :: models.*.top.leak_systematic.eps_ref (min) :: sha256 2416fc21f93fcfa6 :: the same, lower variation
\newcommand{\AojLeakShiftMin}{+97}  % experiments/FIGS/data/aoj_full_v1/fit_v6/results.json :: models.*.top.leak_systematic.shift['0.6'] (min) :: sha256 2416fc21f93fcfa6 :: tops added to a yield
\newcommand{\AojLeakShiftMax}{+271}  % experiments/FIGS/data/aoj_full_v1/fit_v6/results.json :: models.*.top.leak_systematic.shift['0.4'] (max) :: sha256 2416fc21f93fcfa6 :: tops added to a yield
\newcommand{\AojFitPMin}{0.010}  % experiments/FIGS/data/aoj_full_v1/fit_v6/results.json :: models.*.top.asymptotic_p (min) :: sha256 2416fc21f93fcfa6 :: goodness of fit of the working-point fit, pretrained models
\newcommand{\AojFitPMax}{0.974}  % experiments/FIGS/data/aoj_full_v1/fit_v6/results.json :: models.*.top.asymptotic_p (max) :: sha256 2416fc21f93fcfa6
\newcommand{\AojWeakestYield}{885}  % experiments/FIGS/data/aoj_full_v1/fit_v6/results.json :: models.l162-s2.top.signal_yield :: sha256 2416fc21f93fcfa6 :: the smallest yield of any pretrained model
\newcommand{\AojWeakestYieldErr}{211}  % experiments/FIGS/data/aoj_full_v1/fit_v6/results.json :: models.l162-s2.top.signal_yield_err :: sha256 2416fc21f93fcfa6
\newcommand{\AojWeakestFitP}{0.94}  % experiments/FIGS/data/aoj_full_v1/fit_v6/results.json :: models.l162-s2.top.asymptotic_p :: sha256 2416fc21f93fcfa6
\newcommand{\AojSidebandZMin}{\ensuremath{-}0.43}  % experiments/FIGS/data/aoj_full_v1/fit_v6/histograms.npz :: experiments/AOJ/sideband_residuals.py: z (min over the pretrained models) :: sha256 2bd1e3d5b32d4a12 :: passing jets minus the fit, summed from the fit's lower edge to the top window
\newcommand{\AojSidebandZMax}{+0.58}  % experiments/FIGS/data/aoj_full_v1/fit_v6/histograms.npz :: experiments/AOJ/sideband_residuals.py: z (max over the pretrained models) :: sha256 2bd1e3d5b32d4a12
\newcommand{\AojWeakestSidebandZ}{0.00}  % experiments/FIGS/data/aoj_full_v1/fit_v6/histograms.npz :: experiments/AOJ/sideband_residuals.py: z of l162-s2 :: sha256 2bd1e3d5b32d4a12 :: the smallest-yield model
\newcommand{\AojSidebandMaxPull}{2.7}  % experiments/FIGS/data/aoj_full_v1/fit_v6/histograms.npz :: experiments/AOJ/sideband_residuals.py: max_abs_pull (max over the pretrained models) :: sha256 2bd1e3d5b32d4a12 :: largest single-bin pull below the top window, any pretrained model
\newcommand{\AojSharedBinLo}{125}  % experiments/FIGS/data/aoj_full_v1/fit_v6/histograms.npz :: the bin holding the most models' largest pull :: sha256 2bd1e3d5b32d4a12 :: GeV
\newcommand{\AojSharedBinHi}{130}  % experiments/FIGS/data/aoj_full_v1/fit_v6/histograms.npz :: its upper edge :: sha256 2bd1e3d5b32d4a12 :: GeV
\newcommand{\AojSharedBinNModels}{18}  % experiments/FIGS/data/aoj_full_v1/fit_v6/histograms.npz :: models whose largest pull is in that bin :: sha256 2bd1e3d5b32d4a12
\newcommand{\AojSharedBinMeanPull}{+1.4}  % experiments/FIGS/data/aoj_full_v1/fit_v6/histograms.npz :: experiments/AOJ/sideband_residuals.py mean_pulls: that bin :: sha256 2bd1e3d5b32d4a12 :: mean pull over the pretrained models
\newcommand{\AojSharedBinRank}{third}  % experiments/FIGS/data/aoj_full_v1/fit_v6/histograms.npz :: experiments/AOJ/sideband_residuals.py mean_pulls: rank of that bin's |mean pull| :: sha256 2bd1e3d5b32d4a12
\newcommand{\AojFitNBins}{39}  % experiments/FIGS/data/aoj_full_v1/fit_v6/histograms.npz :: mass bins of the fit range :: sha256 2bd1e3d5b32d4a12
\newcommand{\AojPairedYieldOnesixtwoOverOneeighteight}{0.94~[0.45, 1.95]}  % experiments/FIGS/data/aoj_full_v1/analysis_v6/aoj_top.json :: per_label_set.label_sets.{162,188}.signal_yields :: sha256 5dfef902a5fe75ba :: yield ratio paired by run, geometric mean [95% Student-t interval]
\newcommand{\AojPairedYieldFourthreeOverOneeighteight}{0.97~[0.81, 1.16]}  % experiments/FIGS/data/aoj_full_v1/analysis_v6/aoj_top.json :: per_label_set.label_sets.{43,188}.signal_yields :: sha256 5dfef902a5fe75ba :: yield ratio paired by run, geometric mean [95% Student-t interval]
\newcommand{\AojPairedYieldOnesevenOverOneeighteight}{1.01~[0.75, 1.35]}  % experiments/FIGS/data/aoj_full_v1/analysis_v6/aoj_top.json :: per_label_set.label_sets.{17,188}.signal_yields :: sha256 5dfef902a5fe75ba :: yield ratio paired by run, geometric mean [95% Student-t interval]
\newcommand{\AojPairedYieldOnesevenOverOnesixtwo}{1.07~[0.51, 2.24]}  % experiments/FIGS/data/aoj_full_v1/analysis_v6/aoj_top.json :: per_label_set.label_sets.{17,162}.signal_yields :: sha256 5dfef902a5fe75ba :: yield ratio paired by run, geometric mean [95% Student-t interval]
\newcommand{\AojPairedYieldOnesixtwoMassOverOnesixtwo}{1.10~[0.57, 2.14]}  % experiments/FIGS/data/aoj_full_v1/analysis_v6/aoj_top.json :: per_label_set.label_sets.{162+mass,162}.signal_yields :: sha256 5dfef902a5fe75ba :: yield ratio paired by run, geometric mean [95% Student-t interval]
\newcommand{\AojPairedYieldOnesevenMassOverOneseven}{1.05~[0.74, 1.49]}  % experiments/FIGS/data/aoj_full_v1/analysis_v6/aoj_top.json :: per_label_set.label_sets.{17+mass,17}.signal_yields :: sha256 5dfef902a5fe75ba :: yield ratio paired by run, geometric mean [95% Student-t interval]
\newcommand{\AojLeakPairedShiftMax}{2\%}  % experiments/FIGS/data/aoj_full_v1/fit_v6/results.json :: models.*.top.leak_systematic.signal_yield against signal_yield, paired ratios :: sha256 2416fc21f93fcfa6 :: largest change of a paired yield ratio when the fail-region tops are scaled up
\newcommand{\AojRefEffPercent}{1.38\%}  % experiments/FIGS/data/aoj_full_v1/fit_v6/results.json :: reference.top.data_efficiency :: sha256 2416fc21f93fcfa6 :: fraction of all data jets passing its cut
\newcommand{\AojRefFitP}{\ensuremath{2.56\times10^{-11}}}  % experiments/FIGS/data/aoj_full_v1/fit_v6/results.json :: reference.top.asymptotic_p :: sha256 2416fc21f93fcfa6 :: goodness of fit of the reference's own top fit
\newcommand{\AojRefValidationP}{\ensuremath{p \leq 0.005} (0 of 200 toys reach the observed deviance)}  % experiments/FIGS/data/aoj_full_v1/fit_v6/results.json :: reference.top.validation.toy_p, n_toys :: sha256 2416fc21f93fcfa6 :: toy p-value of the background validation, (k + 1) / (n + 1)
\newcommand{\AojRefBandSignalZ}{2.5}  % experiments/FIGS/data/aoj_full_v1/fit_v6/results.json :: reference.top.validation.band_signal_z :: sha256 2416fc21f93fcfa6 :: significance of a peak of the score's own shape in the validation band
\newcommand{\AojModelEffMin}{1.01\%}  % experiments/FIGS/data/aoj_full_v1/fit_v6/results.json :: models.*.top.data_efficiency (min) :: sha256 2416fc21f93fcfa6 :: pretrained models
\newcommand{\AojModelEffMax}{1.05\%}  % experiments/FIGS/data/aoj_full_v1/fit_v6/results.json :: models.*.top.data_efficiency (max) :: sha256 2416fc21f93fcfa6 :: pretrained models
\newcommand{\AojModelsValidationPMin}{0.11}  % experiments/FIGS/data/aoj_full_v1/fit_v6/results.json :: models.*.top.validation.toy_p (min) :: sha256 2416fc21f93fcfa6 :: pretrained models
\newcommand{\AojRefYield}{29{,}143}  % experiments/FIGS/data/aoj_full_v1/fit_v6/results.json :: reference.top.signal_yield :: sha256 2416fc21f93fcfa6
\newcommand{\AojRefYieldErr}{865}  % experiments/FIGS/data/aoj_full_v1/fit_v6/results.json :: reference.top.signal_yield_err :: sha256 2416fc21f93fcfa6
\newcommand{\AojRefMass}{181.0}  % experiments/FIGS/data/aoj_full_v1/fit_v6/results.json :: reference.top.floated_mean :: sha256 2416fc21f93fcfa6 :: GeV
\newcommand{\AojRefWidth}{17.8}  % experiments/FIGS/data/aoj_full_v1/fit_v6/results.json :: reference.top.floated_width :: sha256 2416fc21f93fcfa6 :: GeV
\newcommand{\AojPublicYield}{1{,}813}  % experiments/FIGS/data/aoj_full_v1/fit_v6/results.json :: models.sophon-public.top.signal_yield :: sha256 2416fc21f93fcfa6
\newcommand{\AojPublicYieldErr}{193}  % experiments/FIGS/data/aoj_full_v1/fit_v6/results.json :: models.sophon-public.top.signal_yield_err :: sha256 2416fc21f93fcfa6
\newcommand{\AojPublicMass}{182.8}  % experiments/FIGS/data/aoj_full_v1/fit_v6/results.json :: models.sophon-public.top.floated_mean :: sha256 2416fc21f93fcfa6 :: GeV
\newcommand{\AojPublicWidth}{11.7}  % experiments/FIGS/data/aoj_full_v1/fit_v6/results.json :: models.sophon-public.top.floated_width :: sha256 2416fc21f93fcfa6 :: GeV
\newcommand{\AojInjSizeMin}{1{,}000}  % experiments/FIGS/data/aoj_injection_v1/summary.json :: groups[region=top,mode=ensemble].size (min) :: sha256 d9b66466ae24b584 :: injected tops
\newcommand{\AojInjSizeMax}{4{,}000}  % experiments/FIGS/data/aoj_injection_v1/summary.json :: groups[region=top,mode=ensemble].size (max) :: sha256 d9b66466ae24b584
\newcommand{\AojInjFloatPullSdMax}{2.0}  % experiments/FIGS/data/aoj_injection_v1/summary.json :: groups[region=top,mode=bootstrap,variant=float].pull.sd (max) :: sha256 d9b66466ae24b584 :: pull width with each model's own floated shape
\newcommand{\AojInjEnsembleRecoveryMin}{0.983}  % experiments/FIGS/data/aoj_injection_v1/summary.json :: groups[region=top,mode=ensemble].ratio (min) :: sha256 d9b66466ae24b584 :: the whole procedure, shape re-derived per toy set
\newcommand{\AojInjEnsembleRecoveryMax}{0.998}  % experiments/FIGS/data/aoj_injection_v1/summary.json :: groups[region=top,mode=ensemble].ratio (max) :: sha256 d9b66466ae24b584
\newcommand{\AojInjEnsemblePullSdMin}{1.02}  % experiments/FIGS/data/aoj_injection_v1/summary.json :: groups[region=top,mode=ensemble].pull (min) :: sha256 d9b66466ae24b584 :: the whole procedure, shape re-derived per toy set
\newcommand{\AojInjEnsemblePullSdMax}{1.09}  % experiments/FIGS/data/aoj_injection_v1/summary.json :: groups[region=top,mode=ensemble].pull (max) :: sha256 d9b66466ae24b584
\newcommand{\AojClosureNShards}{10}  % experiments/FIGS/data/aoj_checks_v2/closure_all_shards.json :: n_shards :: sha256 6a6018351d7048d2 :: staged shards of the open data
\newcommand{\AojClosurePtLo}{500}  % experiments/FIGS/data/aoj_checks_v2/closure_all_shards.json :: pt_window[0] :: sha256 6a6018351d7048d2 :: GeV
\newcommand{\AojClosurePtHi}{700}  % experiments/FIGS/data/aoj_checks_v2/closure_all_shards.json :: pt_window[1] :: sha256 6a6018351d7048d2 :: GeV
\newcommand{\AojClosureDzeroIqr}{8.7}  % experiments/FIGS/data/aoj_checks_v2/closure_all_shards.json :: domain_shift.part_d0.iqr_ratio :: sha256 6a6018351d7048d2 :: open data over JetClass-II QCD, pooled over the shards
\newcommand{\AojClosureDzIqr}{10.6}  % experiments/FIGS/data/aoj_checks_v2/closure_all_shards.json :: domain_shift.part_dz.iqr_ratio :: sha256 6a6018351d7048d2 :: open data over JetClass-II QCD, pooled over the shards
\newcommand{\AojClosureNeutral}{0.65}  % experiments/FIGS/data/aoj_checks_v2/closure_all_shards.json :: features.part_isNeutralHadron.ratio.per_shard (mean) :: sha256 6a6018351d7048d2 :: open data over JetClass-II QCD
\newcommand{\AojClosureCharged}{1.22}  % experiments/FIGS/data/aoj_checks_v2/closure_all_shards.json :: features.part_isChargedHadron.ratio.per_shard (mean) :: sha256 6a6018351d7048d2 :: open data over JetClass-II QCD
\newcommand{\AojClosureNPartAoj}{61}  % experiments/FIGS/data/aoj_checks_v2/closure_all_shards.json :: features.n_particles.aoj.per_shard (mean of the shard medians) :: sha256 6a6018351d7048d2 :: median
\newcommand{\AojClosureNPartRef}{48}  % experiments/FIGS/data/aoj_checks_v2/closure_all_shards.json :: features.n_particles.reference.per_shard (mean) :: sha256 6a6018351d7048d2 :: median
\newcommand{\RandMergingDrawsBvcFourprong}{three}  % configs/labelmaps/rand_label_map.v1.csv :: RAND_d* of label_X_YY_bbqq and label_X_YY_ccqq :: sha256 f395a25f52e3cfad :: first-grid partitions that merge the task's pair
\newcommand{\RandSplittingDrawsBvcFourprong}{one and two}  % configs/labelmaps/rand_label_map.v1.csv :: RAND_d* of label_X_YY_bbqq and label_X_YY_ccqq :: sha256 f395a25f52e3cfad :: first-grid partitions that keep the pair apart
\newcommand{\RandMergingDrawsVisibleContent}{one and three}  % configs/labelmaps/rand_label_map.v1.csv :: RAND_d* of label_X_YY_bbqq and label_X_YY_cqtauhv :: sha256 f395a25f52e3cfad :: first-grid partitions that merge the task's pair
\newcommand{\RandSplittingDrawsVisibleContent}{two}  % configs/labelmaps/rand_label_map.v1.csv :: RAND_d* of label_X_YY_bbqq and label_X_YY_cqtauhv :: sha256 f395a25f52e3cfad :: first-grid partitions that keep the pair apart
\newcommand{\RandVtwoNPartitions}{five}  % configs/arms/v2_grid.json :: arms[RAND2_*] (count) :: sha256 d8a31b1486992302 :: random partitions in the rerun
\newcommand{\RandVtwoRuns}{one}  % configs/arms/v2_paper_runs.txt :: runs of RAND2_* :: sha256 b04562b14ea8cbc7 :: runs per partition
\newcommand{\RandFlavRuns}{two}  % configs/arms/v2_paper_runs.txt :: runs of FLAV_* :: sha256 b04562b14ea8cbc7 :: runs per flavour-pair model
\newcommand{\RandVtwoNPairs}{seven}  % configs/labelmaps/rand_v2_selection.json :: pairs (count) :: sha256 98cf52d49c7d2f02 :: probe pairs the balance rule covers
\newcommand{\RandVtwoMergedMin}{two}  % configs/labelmaps/rand_v2_selection.json :: merged_in :: sha256 98cf52d49c7d2f02 :: partitions merging each probe pair, checked on the map
\newcommand{\RandVtwoMergedMax}{three}  % configs/labelmaps/rand_v2_selection.json :: merged_in :: sha256 98cf52d49c7d2f02 :: partitions merging each probe pair, checked on the map
\newcommand{\RandVtwoShareTol}{5\%}  % configs/labelmaps/rand_v2_selection.json :: realised_tolerance :: sha256 98cf52d49c7d2f02 :: largest relative deviation of a realised group share from its target (rule 3)
\newcommand{\RandVtwoTrials}{32}  % scripts/build_rand_control.py :: TRIALS :: sha256 5bc91d77a3f2e81f :: share-matched assignments drawn per random partition; the one sharing the fewest merged pairs with the 17-class label set is kept
\newcommand{\RandVtwoPoolDraws}{6{,}600}  % configs/labelmaps/rand_v2_selection.json :: pool.draws :: sha256 98cf52d49c7d2f02 :: candidate partitions drawn (seeds) before the realised-share rule
\newcommand{\RandVtwoPoolPass}{27}  % configs/labelmaps/rand_v2_selection.json :: rule3_seeds (count) :: sha256 98cf52d49c7d2f02 :: candidate partitions meeting the realised-share rule (rule 3)
\newcommand{\RandFlavSplitOrbit}{\ensuremath{X\to YY\to QQQQ}}  % configs/labelmaps/flavour_pair.v2.json :: split_orbit :: sha256 3b79ce5e44c1e2f4 :: the decays F1 cuts, Q any quark
\newcommand{\RandFlavSplitOrbitClasses}{22}  % configs/labelmaps/flavour_pair_map.v2.csv :: classes whose orbit is split_orbit (count) :: sha256 9aab235099ceac43 :: classes of the split orbit
\newcommand{\RandFlavCutClasses}{11}  % configs/labelmaps/flavour_pair.v2.json :: split_classes_moved (count) :: sha256 3b79ce5e44c1e2f4 :: classes of the split orbit that F1 moves: those with a b quark
\newcommand{\RandFlavOrbitsMoved}{three}  % configs/labelmaps/flavour_pair.v2.json :: orbits_moved_B_to_A (count) :: sha256 3b79ce5e44c1e2f4 :: whole orbits F1 moves the other way
\newcommand{\RandFlavOrbitsMovedList}{\ensuremath{X\to YY\to QQ\mu\mu}, \ensuremath{X\to YY\to QQ\tau_h\tau_h} and \ensuremath{X\to \mu\mu}}  % configs/labelmaps/flavour_pair.v2.json :: orbits_moved_B_to_A :: sha256 3b79ce5e44c1e2f4 :: the orbits F1 moves the other way, Q any quark
\newcommand{\RandFlavOrbitClassesMoved}{9}  % configs/labelmaps/flavour_pair_map.v2.csv :: classes whose orbit is in orbits_moved_B_to_A :: sha256 9aab235099ceac43 :: classes in those orbits
\newcommand{\RandFlavFonerCutWithB}{6}  % configs/labelmaps/flavour_pair_map.v2.csv :: FLAV_F1R != FLAV_F0 within split_orbit, with a b quark (count) :: sha256 9aab235099ceac43 :: classes with a b quark among those the F1r cut moves
\newcommand{\RandFlavBcPairs}{11}  % configs/labelmaps/flavour_pair.v2.json :: b_c_pairs_in_split_orbit (count) :: sha256 3b79ce5e44c1e2f4 :: pairs of the split orbit's classes that differ only by exchanging b and c
\newcommand{\RandFlavFoneBcSplit}{8}  % configs/labelmaps/flavour_pair_map.v2.csv :: b_c_pairs_in_split_orbit split by FLAV_F1 (count) :: sha256 9aab235099ceac43
\newcommand{\RandFlavFonerBcSplit}{none}  % configs/labelmaps/flavour_pair_map.v2.csv :: b_c_pairs_in_split_orbit split by FLAV_F1R (count) :: sha256 9aab235099ceac43 :: PRESPEC A14: none
\newcommand{\RandFlavFonerShareTol}{1\%}  % configs/labelmaps/flavour_pair.v2.json :: F1R.realised_tolerance_vs_F1 :: sha256 3b79ce5e44c1e2f4 :: largest allowed relative difference of a group's realised share from F1's
\newcommand{\RandFlavFonerShareDev}{0.77\%}  % configs/labelmaps/realised_native_shares.v2.json :: share summed per FLAV_F1R and FLAV_F1 group of flavour_pair_map.v2.csv (largest relative difference) :: sha256 5728c12e88fa3683 :: realised shares from the final loader's dry run
\newcommand{\MlstNRunsWord}{three}  % configs/arms/v2_paper_runs.txt :: runs of L188/L162/R42_Q1/R16_Q1 :: sha256 b04562b14ea8cbc7 :: pretraining runs per configuration
\newcommand{\MlstMassNRunsWord}{two}  % configs/arms/v2_paper_runs.txt :: runs of L162_MASS/R16_Q1_MASS :: sha256 b04562b14ea8cbc7 :: pretraining runs per configuration
\newcommand{\LpCMin}{\ensuremath{1\times10^{-3}}}  % experiments/EVAL/linear_probe_v2.py :: C_GRID :: sha256 4c576b44812c18a2 :: smallest inverse regularisation
\newcommand{\LpCMax}{100}  % experiments/EVAL/linear_probe_v2.py :: C_GRID :: sha256 4c576b44812c18a2 :: largest inverse regularisation
\newcommand{\LpNTestJcii}{2{,}000{,}000}  % experiments/EVAL/linear_probe_v2.py :: TEST_FIRST.jc2 :: sha256 4c576b44812c18a2 :: JetClass-II test jets
\newcommand{\LpNTestAucJcii}{500{,}000}  % experiments/EVAL/linear_probe_v2.py :: TEST_FIRST.jc2 / AUC_STRIDE :: sha256 4c576b44812c18a2 :: JetClass-II test jets of the macro AUC
\newcommand{\LpConvergedGrad}{\ensuremath{1\times10^{-5}}}  % experiments/EVAL/linear_probe_v2.py :: CONVERGED_GRAD :: sha256 4c576b44812c18a2 :: largest gradient component a converged probe ends below
\newcommand{\FillNFeatures}{17}  % configs/finetune/TopReference.yaml :: inputs.pf_features.vars (count) :: sha256 501a6ffe5d0a8051 :: per-particle input features of the network
\newcommand{\FillTopNZero}{10}  % configs/finetune/TopReference.yaml :: inputs.pf_features.vars, part_zero_* (count) :: sha256 501a6ffe5d0a8051 :: features the data set lacks, set to zero
\newcommand{\FillTopNPresent}{7}  % configs/finetune/TopReference.yaml :: inputs.pf_features.vars, not part_zero_* (count) :: sha256 501a6ffe5d0a8051 :: features the data set provides
\newcommand{\FillQgNZero}{4}  % configs/finetune/EnergyFlowQG.yaml :: inputs.pf_features.vars, part_zero_* (count) :: sha256 db42d6471e8b865c :: features the data set lacks, set to zero
\newcommand{\FillQgNPresent}{13}  % configs/finetune/EnergyFlowQG.yaml :: inputs.pf_features.vars, not part_zero_* (count) :: sha256 db42d6471e8b865c :: features the data set provides
\newcommand{\DesignVtwoDataFraction}{20\%}  % experiments/MTX/k8s/v2/grid/job-mtx2-l188-s1-raunav.yaml :: --data-fraction :: sha256 d28ce02267a6b87a :: share of every file read per epoch
\newcommand{\DesignVtwoNValFiles}{25}  % experiments/MTX/k8s/v2/grid/job-mtx2-l188-s1-raunav.yaml :: --data-val (files, brace ranges expanded) :: sha256 d28ce02267a6b87a :: validation files of the fixed sample
\newcommand{\DesignWavgFirstEpoch}{70}  % experiments/MTX/pretrain_v2.py :: --num-epochs - WAVG_EPOCHS :: sha256 95da5a6e029855ee :: first epoch of the weight average (epochs count from 0)
\newcommand{\DesignWavgLastEpoch}{79}  % experiments/MTX/k8s/v2/grid/job-mtx2-l188-s1-raunav.yaml :: --num-epochs - 1 :: sha256 d28ce02267a6b87a :: last epoch
\newcommand{\DesignBnJets}{200{,}000}  % experiments/MTX/pretrain_v2.py :: BN_JETS :: sha256 95da5a6e029855ee :: training jets the BatchNorm statistics are recomputed on
\newcommand{\DesignVtwoNWindows}{five}  % experiments/MTX/k8s/v2/grid/job-mtx2-l188-s1-raunav.yaml :: 1 / --data-fraction :: sha256 d28ce02267a6b87a :: disjoint windows of each file's rows, one read per epoch
\newcommand{\DesignLofoNWindows}{three}  % experiments/MTX/k8s/v2/cut/grid/job-mtx2-l188lofo4p-s1-raunav.yaml :: --data-windows :: sha256 44a9573929efb7ac :: windows of the models trained without one family
\newcommand{\DesignVtwoPlanUsed}{97\%}  % experiments/FIGS/data/v2_loader/loader_dryrun/dryrun_s176_seed1.json :: epochs[*].(max_fetch_id + 1) / args.data_split_num, mean :: sha256 03b6586b26ad7a2f :: share of an epoch's planned reading used before its jets are reached, 20 epochs
\newcommand{\DesignLofoPlanUsed}{76\%}  % experiments/FIGS/data/v2_loader/loader_dryrun/dryrun_lofo4p_seed1.json :: epochs[*].(max_fetch_id + 1) / fetches_per_pass, mean :: sha256 2a1315e611d13ff5 :: the same, leaving the family out
\newcommand{\DesignLofoRowFactor}{1.29}  % experiments/FIGS/data/v2_loader/loader_dryrun/dryrun_lofo4p_seed1.json :: (mean plan used / --data-windows) over (mean plan used / k) of the full models' dry run :: sha256 2a1315e611d13ff5 :: how much more often a row of a file is read per epoch when the family is left out
\newcommand{\DesignDistinctJetsCycle}{\ensuremath{3.3\times10^{7}}}  % experiments/FIGS/data/v2_loader/loader_dryrun/dryrun_s176_seed1.json :: sum of epochs[c*k:(c+1)*k].distinct_jets, mean over 4 cycles :: sha256 03b6586b26ad7a2f :: distinct jets the full models train on in one cycle of windows
\newcommand{\DesignLofoFamilyShare}{22.4\%}  % experiments/FIGS/data/v2_loader/loader_dryrun/dryrun_s176_seed1.json :: epochs[*].native_counts over summary.labels_absent_every_epoch of the family-out run :: sha256 03b6586b26ad7a2f :: the left-out family's share of the full models' stream
\newcommand{\DesignLofoNClasses}{27}  % experiments/FIGS/data/v2_loader/loader_dryrun/dryrun_lofo4p_seed1.json :: summary.labels_absent_every_epoch (count) :: sha256 2a1315e611d13ff5 :: native classes of the left-out family
\newcommand{\DesignLofoExposureFactor}{1.29}  % experiments/FIGS/data/v2_loader/loader_dryrun/dryrun_s176_seed1.json :: 1 / (1 - family share) :: sha256 03b6586b26ad7a2f :: how much more often every other class is seen
\newcommand{\DesignLofoQcdShareFull}{14.4\%}  % experiments/FIGS/data/v2_loader/loader_dryrun/dryrun_s176_seed1.json :: summary.mean_qcd_share :: sha256 03b6586b26ad7a2f :: QCD share of the realised stream
\newcommand{\DesignLofoQcdShareLofo}{18.6\%}  % experiments/FIGS/data/v2_loader/loader_dryrun/dryrun_lofo4p_seed1.json :: summary.mean_qcd_share :: sha256 2a1315e611d13ff5 :: QCD share of the realised stream
\newcommand{\DesignLrFlatPercent}{70\%}  % experiments/MTX/pretrain_v2.py :: make_optimizer: 1 - the decay share :: sha256 95da5a6e029855ee :: share of the epochs at the initial rate
\newcommand{\DesignLrDecayFirstEpoch}{56}  % experiments/MTX/pretrain_v2.py :: num_epochs - int(num_epochs * decay share) :: sha256 95da5a6e029855ee :: first epoch at a reduced rate (epochs from 0)
\newcommand{\DesignLrEndPercent}{1\%}  % experiments/MTX/pretrain_v2.py :: make_optimizer: gamma ** n_decay :: sha256 95da5a6e029855ee :: the last epoch's rate over the initial rate
\newcommand{\DesignEndLr}{\ensuremath{5\times10^{-6}}}  % experiments/MTX/k8s/v2/grid/job-mtx2-l188-s1-raunav.yaml :: --start-lr x the final factor :: sha256 d28ce02267a6b87a :: the rate of the last epoch
\newcommand{\DesignWeaverVersion}{0.4.17}  % experiments/MTX/pretrain_v2.py :: the weaver release whose loop it reproduces :: sha256 95da5a6e029855ee
\newcommand{\DesignMeiFactor}{1.1}  % src/stats/mde.py :: MEI_LOG = ln(factor) :: sha256 626468dcc37940bb :: the smallest effect of interest, as a factor
\newcommand{\DesignVtwoGpu}{RTX 3090}  % experiments/MTX/k8s/v2/grid/job-mtx2-flavf0-s1-raunav.yaml :: nodeAffinity nvidia.com/gpu.product :: sha256 960cad1691314980 :: the product of run index 1, every configuration
\newcommand{\DesignVtwoSecondGpu}{\pending{pending: the GPU type of each run index}}  % experiments/FIGS/make_tables.py :: emit_v2_recipe[DesignVtwoSecondGpu] :: sha256 10e6ec08356a4a2b :: a marker, not a number: the grid's job specs put every run index on one product
\newcommand{\DesignVtwoSecondGpuRuns}{\pending{pending: the GPU type of each run index}}  % experiments/FIGS/make_tables.py :: emit_v2_recipe[DesignVtwoSecondGpuRuns] :: sha256 10e6ec08356a4a2b :: a marker, not a number: the grid's job specs put every run index on one product
\newcommand{\MassNRunsWord}{five}  % experiments/FIGS/data/probe_ladder_mass2x2_mlp2/s1.json :: per-seed files (count) :: sha256 906e778d7043fe67 :: pretraining runs of each mass-output configuration
\newcommand{\MassOmaLinearOnesixtwo}{\ensuremath{(2.36 \pm 0.06)\times10^{-3}}}  % experiments/FIGS/data/probe_ladder_v2_mlp2/analysis/seed_level_results.json :: table[bvc_resonant,linear,162].auc :: sha256 53e1ec81343ac2c1 :: 1-AUC, mean +- SD over 5 seeds
\newcommand{\MassOmaLinearOnesixtwoMass}{\ensuremath{(2.62 \pm 0.08)\times10^{-3}}}  % experiments/FIGS/data/probe_ladder_mass2x2_mlp2/s1.json :: tasks.bvc_resonant.arms.l162mass-s*.linear.auc :: sha256 906e778d7043fe67 :: 1-AUC, mean +- SD over the 5 per-seed files
\newcommand{\MassOmaRatioLinearOnesixtwo}{1.11}  % experiments/FIGS/data/probe_ladder_mass2x2_mlp2/s1.json :: tasks.bvc_resonant.arms.l162mass-s*.linear.auc over table[bvc_resonant,linear,162].auc :: sha256 906e778d7043fe67 :: seed mean of 1-AUC with the mass output over without
\newcommand{\MassOmaLinearOneseven}{\ensuremath{(9.3 \pm 1.5)\times10^{-3}}}  % experiments/FIGS/data/probe_ladder_v2_mlp2/analysis/seed_level_results.json :: table[bvc_resonant,linear,17].auc :: sha256 53e1ec81343ac2c1 :: 1-AUC, mean +- SD over 5 seeds
\newcommand{\MassOmaLinearOnesevenMass}{\ensuremath{(3.0 \pm 1.2)\times10^{-2}}}  % experiments/FIGS/data/probe_ladder_mass2x2_mlp2/s1.json :: tasks.bvc_resonant.arms.r16q1mass-s*.linear.auc :: sha256 906e778d7043fe67 :: 1-AUC, mean +- SD over the 5 per-seed files
\newcommand{\MassOmaRatioLinearOneseven}{3.2}  % experiments/FIGS/data/probe_ladder_mass2x2_mlp2/s1.json :: tasks.bvc_resonant.arms.r16q1mass-s*.linear.auc over table[bvc_resonant,linear,17].auc :: sha256 906e778d7043fe67 :: seed mean of 1-AUC with the mass output over without
\newcommand{\MassOmaMlpOnesixtwo}{\ensuremath{(2.32 \pm 0.05)\times10^{-3}}}  % experiments/FIGS/data/probe_ladder_v2_mlp2/analysis/seed_level_results.json :: table[bvc_resonant,mlp,162].auc :: sha256 53e1ec81343ac2c1 :: 1-AUC, mean +- SD over 5 seeds
\newcommand{\MassOmaMlpOnesixtwoMass}{\ensuremath{(2.52 \pm 0.06)\times10^{-3}}}  % experiments/FIGS/data/probe_ladder_mass2x2_mlp2/s1.json :: tasks.bvc_resonant.arms.l162mass-s*.mlp.auc :: sha256 906e778d7043fe67 :: 1-AUC, mean +- SD over the 5 per-seed files
\newcommand{\MassOmaRatioMlpOnesixtwo}{1.09}  % experiments/FIGS/data/probe_ladder_mass2x2_mlp2/s1.json :: tasks.bvc_resonant.arms.l162mass-s*.mlp.auc over table[bvc_resonant,mlp,162].auc :: sha256 906e778d7043fe67 :: seed mean of 1-AUC with the mass output over without
\newcommand{\MassOmaMlpOneseven}{\ensuremath{(7.3 \pm 0.9)\times10^{-3}}}  % experiments/FIGS/data/probe_ladder_v2_mlp2/analysis/seed_level_results.json :: table[bvc_resonant,mlp,17].auc :: sha256 53e1ec81343ac2c1 :: 1-AUC, mean +- SD over 5 seeds
\newcommand{\MassOmaMlpOnesevenMass}{\ensuremath{(1.5 \pm 0.4)\times10^{-2}}}  % experiments/FIGS/data/probe_ladder_mass2x2_mlp2/s1.json :: tasks.bvc_resonant.arms.r16q1mass-s*.mlp.auc :: sha256 906e778d7043fe67 :: 1-AUC, mean +- SD over the 5 per-seed files
\newcommand{\MassOmaRatioMlpOneseven}{2.0}  % experiments/FIGS/data/probe_ladder_mass2x2_mlp2/s1.json :: tasks.bvc_resonant.arms.r16q1mass-s*.mlp.auc over table[bvc_resonant,mlp,17].auc :: sha256 906e778d7043fe67 :: seed mean of 1-AUC with the mass output over without
\newcommand{\FtAucJciiEThreeOneeighteight}{0.902\,\ensuremath{\pm}\,0.004}  % experiments/FIGS/data/w2b_leg1_metrics_v2.json :: cells.{l188-s1,l188-s2,l188-s3,l188-s4,l188-s5}.N1000.s1.macro_auc_ovr :: sha256 b56bf70be9982044 :: mean +- SD over the row's 5 pretrained models
\newcommand{\FtAucJciiEFourOneeighteight}{0.9723\,\ensuremath{\pm}\,0.0004}  % experiments/FIGS/data/w2b_leg1_metrics_v2.json :: cells.{l188-s1,l188-s2,l188-s3,l188-s4,l188-s5}.N10000.s1.macro_auc_ovr :: sha256 b56bf70be9982044 :: mean +- SD over the row's 5 pretrained models
\newcommand{\FtAucJciiEFiveOneeighteight}{0.9793\,\ensuremath{\pm}\,0.0002}  % experiments/FIGS/data/w2b_leg1_metrics_v2.json :: cells.{l188-s1,l188-s2,l188-s3,l188-s4,l188-s5}.N100000.s1.macro_auc_ovr :: sha256 b56bf70be9982044 :: mean +- SD over the row's 5 pretrained models
\newcommand{\FtAucJciiESixOneeighteight}{0.9809\,\ensuremath{\pm}\,\ensuremath{<}0.0001}  % experiments/FIGS/data/w2b_leg1_metrics_v2.json :: cells.{l188-s1,l188-s2,l188-s3,l188-s4,l188-s5}.N1000000.s1.macro_auc_ovr :: sha256 b56bf70be9982044 :: mean +- SD over the row's 5 pretrained models
\newcommand{\FtAucJciiEThreeOnesixtwo}{0.906\,\ensuremath{\pm}\,0.007}  % experiments/FIGS/data/w2b_leg1_metrics_v2.json :: cells.{l162-s1b,l162-s2,l162-s3,l162-s4,l162-s5}.N1000.s1.macro_auc_ovr :: sha256 b56bf70be9982044 :: mean +- SD over the row's 5 pretrained models
\newcommand{\FtAucJciiEFourOnesixtwo}{0.9724\,\ensuremath{\pm}\,0.0005}  % experiments/FIGS/data/w2b_leg1_metrics_v2.json :: cells.{l162-s1b,l162-s2,l162-s3,l162-s4,l162-s5}.N10000.s1.macro_auc_ovr :: sha256 b56bf70be9982044 :: mean +- SD over the row's 5 pretrained models
\newcommand{\FtAucJciiEFiveOnesixtwo}{0.9794\,\ensuremath{\pm}\,0.0002}  % experiments/FIGS/data/w2b_leg1_metrics_v2.json :: cells.{l162-s1b,l162-s2,l162-s3,l162-s4,l162-s5}.N100000.s1.macro_auc_ovr :: sha256 b56bf70be9982044 :: mean +- SD over the row's 5 pretrained models
\newcommand{\FtAucJciiESixOnesixtwo}{0.9808\,\ensuremath{\pm}\,0.0001}  % experiments/FIGS/data/w2b_leg1_metrics_v2.json :: cells.{l162-s1b,l162-s2,l162-s3,l162-s4,l162-s5}.N1000000.s1.macro_auc_ovr :: sha256 b56bf70be9982044 :: mean +- SD over the row's 5 pretrained models
\newcommand{\FtAucJciiEThreeFourthree}{0.918\,\ensuremath{\pm}\,0.005}  % experiments/FIGS/data/w2b_leg1_metrics_v2.json :: cells.{r42q1-s1,r42q1-s2,r42q1-s3,r42q1-s4,r42q1-s5}.N1000.s1.macro_auc_ovr :: sha256 b56bf70be9982044 :: mean +- SD over the row's 5 pretrained models
\newcommand{\FtAucJciiEFourFourthree}{0.9667\,\ensuremath{\pm}\,0.0005}  % experiments/FIGS/data/w2b_leg1_metrics_v2.json :: cells.{r42q1-s1,r42q1-s2,r42q1-s3,r42q1-s4,r42q1-s5}.N10000.s1.macro_auc_ovr :: sha256 b56bf70be9982044 :: mean +- SD over the row's 5 pretrained models
\newcommand{\FtAucJciiEFiveFourthree}{0.9775\,\ensuremath{\pm}\,0.0003}  % experiments/FIGS/data/w2b_leg1_metrics_v2.json :: cells.{r42q1-s1,r42q1-s2,r42q1-s3,r42q1-s4,r42q1-s5}.N100000.s1.macro_auc_ovr :: sha256 b56bf70be9982044 :: mean +- SD over the row's 5 pretrained models
\newcommand{\FtAucJciiESixFourthree}{0.9800\,\ensuremath{\pm}\,0.0001}  % experiments/FIGS/data/w2b_leg1_metrics_v2.json :: cells.{r42q1-s1,r42q1-s2,r42q1-s3,r42q1-s4,r42q1-s5}.N1000000.s1.macro_auc_ovr :: sha256 b56bf70be9982044 :: mean +- SD over the row's 5 pretrained models
\newcommand{\FtAucJciiEThreeOneseven}{0.905\,\ensuremath{\pm}\,0.008}  % experiments/FIGS/data/w2b_leg1_metrics_v2.json :: cells.{r16q1-s1,r16q1-s2,r16q1-s3,r16q1-s4,r16q1-s5}.N1000.s1.macro_auc_ovr :: sha256 b56bf70be9982044 :: mean +- SD over the row's 5 pretrained models
\newcommand{\FtAucJciiEFourOneseven}{0.9635\,\ensuremath{\pm}\,0.0005}  % experiments/FIGS/data/w2b_leg1_metrics_v2.json :: cells.{r16q1-s1,r16q1-s2,r16q1-s3,r16q1-s4,r16q1-s5}.N10000.s1.macro_auc_ovr :: sha256 b56bf70be9982044 :: mean +- SD over the row's 5 pretrained models
\newcommand{\FtAucJciiEFiveOneseven}{0.9748\,\ensuremath{\pm}\,0.0003}  % experiments/FIGS/data/w2b_leg1_metrics_v2.json :: cells.{r16q1-s1,r16q1-s2,r16q1-s3,r16q1-s4,r16q1-s5}.N100000.s1.macro_auc_ovr :: sha256 b56bf70be9982044 :: mean +- SD over the row's 5 pretrained models
\newcommand{\FtAucJciiESixOneseven}{0.9781\,\ensuremath{\pm}\,0.0001}  % experiments/FIGS/data/w2b_leg1_metrics_v2.json :: cells.{r16q1-s1,r16q1-s2,r16q1-s3,r16q1-s4,r16q1-s5}.N1000000.s1.macro_auc_ovr :: sha256 b56bf70be9982044 :: mean +- SD over the row's 5 pretrained models
\newcommand{\FtAucJciiEThreeOnesixtwoMass}{0.921\,\ensuremath{\pm}\,0.005}  % experiments/FIGS/data/w2b_leg1_metrics_v2.json :: cells.{l162mass-s1,l162mass-s2,l162mass-s3,l162mass-s4,l162mass-s5}.N1000.s1.macro_auc_ovr :: sha256 b56bf70be9982044 :: mean +- SD over the row's 5 pretrained models
\newcommand{\FtAucJciiEFourOnesixtwoMass}{0.9724\,\ensuremath{\pm}\,0.0013}  % experiments/FIGS/data/w2b_leg1_metrics_v2.json :: cells.{l162mass-s1,l162mass-s2,l162mass-s3,l162mass-s4,l162mass-s5}.N10000.s1.macro_auc_ovr :: sha256 b56bf70be9982044 :: mean +- SD over the row's 5 pretrained models
\newcommand{\FtAucJciiEFiveOnesixtwoMass}{0.9790\,\ensuremath{\pm}\,0.0001}  % experiments/FIGS/data/w2b_leg1_metrics_v2.json :: cells.{l162mass-s1,l162mass-s2,l162mass-s3,l162mass-s4,l162mass-s5}.N100000.s1.macro_auc_ovr :: sha256 b56bf70be9982044 :: mean +- SD over the row's 5 pretrained models
\newcommand{\FtAucJciiESixOnesixtwoMass}{0.9805\,\ensuremath{\pm}\,0.0001}  % experiments/FIGS/data/w2b_leg1_metrics_v2.json :: cells.{l162mass-s1,l162mass-s2,l162mass-s3,l162mass-s4,l162mass-s5}.N1000000.s1.macro_auc_ovr :: sha256 b56bf70be9982044 :: mean +- SD over the row's 5 pretrained models
\newcommand{\FtAucJciiEThreeOnesevenMass}{0.904\,\ensuremath{\pm}\,0.007}  % experiments/FIGS/data/w2b_leg1_metrics_v2.json :: cells.{r16q1mass-s1,r16q1mass-s2,r16q1mass-s3,r16q1mass-s4,r16q1mass-s5}.N1000.s1.macro_auc_ovr :: sha256 b56bf70be9982044 :: mean +- SD over the row's 5 pretrained models
\newcommand{\FtAucJciiEFourOnesevenMass}{0.9639\,\ensuremath{\pm}\,0.0007}  % experiments/FIGS/data/w2b_leg1_metrics_v2.json :: cells.{r16q1mass-s1,r16q1mass-s2,r16q1mass-s3,r16q1mass-s4,r16q1mass-s5}.N10000.s1.macro_auc_ovr :: sha256 b56bf70be9982044 :: mean +- SD over the row's 5 pretrained models
\newcommand{\FtAucJciiEFiveOnesevenMass}{0.9740\,\ensuremath{\pm}\,0.0004}  % experiments/FIGS/data/w2b_leg1_metrics_v2.json :: cells.{r16q1mass-s1,r16q1mass-s2,r16q1mass-s3,r16q1mass-s4,r16q1mass-s5}.N100000.s1.macro_auc_ovr :: sha256 b56bf70be9982044 :: mean +- SD over the row's 5 pretrained models
\newcommand{\FtAucJciiESixOnesevenMass}{0.9774\,\ensuremath{\pm}\,0.0001}  % experiments/FIGS/data/w2b_leg1_metrics_v2.json :: cells.{r16q1mass-s1,r16q1mass-s2,r16q1mass-s3,r16q1mass-s4,r16q1mass-s5}.N1000000.s1.macro_auc_ovr :: sha256 b56bf70be9982044 :: mean +- SD over the row's 5 pretrained models
\newcommand{\FtAucJciiEThreeRandDrawOne}{0.921}  % experiments/FIGS/data/w2b_leg1_metrics_v2.json :: cells.{rand-d1-s1b}.N1000.s1.macro_auc_ovr :: sha256 b56bf70be9982044 :: one run
\newcommand{\FtAucJciiEFourRandDrawOne}{0.9675}  % experiments/FIGS/data/w2b_leg1_metrics_v2.json :: cells.{rand-d1-s1b}.N10000.s1.macro_auc_ovr :: sha256 b56bf70be9982044 :: one run
\newcommand{\FtAucJciiEFiveRandDrawOne}{0.97617}  % experiments/FIGS/data/w2b_leg1_metrics_v2.json :: cells.{rand-d1-s1b}.N100000.s1.macro_auc_ovr :: sha256 b56bf70be9982044 :: one run
\newcommand{\FtAucJciiESixRandDrawOne}{0.97941}  % experiments/FIGS/data/w2b_leg1_metrics_v2.json :: cells.{rand-d1-s1b}.N1000000.s1.macro_auc_ovr :: sha256 b56bf70be9982044 :: one run
\newcommand{\FtAucJciiEThreeRandDrawTwo}{0.914}  % experiments/FIGS/data/w2b_leg1_metrics_v2.json :: cells.{rand-d2-s2}.N1000.s1.macro_auc_ovr :: sha256 b56bf70be9982044 :: one run
\newcommand{\FtAucJciiEFourRandDrawTwo}{0.9675}  % experiments/FIGS/data/w2b_leg1_metrics_v2.json :: cells.{rand-d2-s2}.N10000.s1.macro_auc_ovr :: sha256 b56bf70be9982044 :: one run
\newcommand{\FtAucJciiEFiveRandDrawTwo}{0.97668}  % experiments/FIGS/data/w2b_leg1_metrics_v2.json :: cells.{rand-d2-s2}.N100000.s1.macro_auc_ovr :: sha256 b56bf70be9982044 :: one run
\newcommand{\FtAucJciiESixRandDrawTwo}{0.97943}  % experiments/FIGS/data/w2b_leg1_metrics_v2.json :: cells.{rand-d2-s2}.N1000000.s1.macro_auc_ovr :: sha256 b56bf70be9982044 :: one run
\newcommand{\FtAucJciiEThreeRandDrawThree}{0.918}  % experiments/FIGS/data/w2b_leg1_metrics_v2.json :: cells.{rand-d3-s3}.N1000.s1.macro_auc_ovr :: sha256 b56bf70be9982044 :: one run
\newcommand{\FtAucJciiEFourRandDrawThree}{0.9686}  % experiments/FIGS/data/w2b_leg1_metrics_v2.json :: cells.{rand-d3-s3}.N10000.s1.macro_auc_ovr :: sha256 b56bf70be9982044 :: one run
\newcommand{\FtAucJciiEFiveRandDrawThree}{0.97605}  % experiments/FIGS/data/w2b_leg1_metrics_v2.json :: cells.{rand-d3-s3}.N100000.s1.macro_auc_ovr :: sha256 b56bf70be9982044 :: one run
\newcommand{\FtAucJciiESixRandDrawThree}{0.97950}  % experiments/FIGS/data/w2b_leg1_metrics_v2.json :: cells.{rand-d3-s3}.N1000000.s1.macro_auc_ovr :: sha256 b56bf70be9982044 :: one run
\newcommand{\FtAccJciiEThreeOneeighteight}{0.445\,\ensuremath{\pm}\,0.004}  % experiments/FIGS/data/w2b_leg1_metrics_v2.json :: cells.{l188-s1,l188-s2,l188-s3,l188-s4,l188-s5}.N1000.s1.accuracy :: sha256 b56bf70be9982044 :: mean +- SD over the row's 5 pretrained models
\newcommand{\FtAccJciiEFourOneeighteight}{0.5199\,\ensuremath{\pm}\,0.0018}  % experiments/FIGS/data/w2b_leg1_metrics_v2.json :: cells.{l188-s1,l188-s2,l188-s3,l188-s4,l188-s5}.N10000.s1.accuracy :: sha256 b56bf70be9982044 :: mean +- SD over the row's 5 pretrained models
\newcommand{\FtAccJciiEFiveOneeighteight}{0.5379\,\ensuremath{\pm}\,0.0010}  % experiments/FIGS/data/w2b_leg1_metrics_v2.json :: cells.{l188-s1,l188-s2,l188-s3,l188-s4,l188-s5}.N100000.s1.accuracy :: sha256 b56bf70be9982044 :: mean +- SD over the row's 5 pretrained models
\newcommand{\FtAccJciiESixOneeighteight}{0.5443\,\ensuremath{\pm}\,0.0003}  % experiments/FIGS/data/w2b_leg1_metrics_v2.json :: cells.{l188-s1,l188-s2,l188-s3,l188-s4,l188-s5}.N1000000.s1.accuracy :: sha256 b56bf70be9982044 :: mean +- SD over the row's 5 pretrained models
\newcommand{\FtAccJciiEThreeOnesixtwo}{0.448\,\ensuremath{\pm}\,0.006}  % experiments/FIGS/data/w2b_leg1_metrics_v2.json :: cells.{l162-s1b,l162-s2,l162-s3,l162-s4,l162-s5}.N1000.s1.accuracy :: sha256 b56bf70be9982044 :: mean +- SD over the row's 5 pretrained models
\newcommand{\FtAccJciiEFourOnesixtwo}{0.5212\,\ensuremath{\pm}\,0.0013}  % experiments/FIGS/data/w2b_leg1_metrics_v2.json :: cells.{l162-s1b,l162-s2,l162-s3,l162-s4,l162-s5}.N10000.s1.accuracy :: sha256 b56bf70be9982044 :: mean +- SD over the row's 5 pretrained models
\newcommand{\FtAccJciiEFiveOnesixtwo}{0.5387\,\ensuremath{\pm}\,0.0008}  % experiments/FIGS/data/w2b_leg1_metrics_v2.json :: cells.{l162-s1b,l162-s2,l162-s3,l162-s4,l162-s5}.N100000.s1.accuracy :: sha256 b56bf70be9982044 :: mean +- SD over the row's 5 pretrained models
\newcommand{\FtAccJciiESixOnesixtwo}{0.5440\,\ensuremath{\pm}\,0.0006}  % experiments/FIGS/data/w2b_leg1_metrics_v2.json :: cells.{l162-s1b,l162-s2,l162-s3,l162-s4,l162-s5}.N1000000.s1.accuracy :: sha256 b56bf70be9982044 :: mean +- SD over the row's 5 pretrained models
\newcommand{\FtAccJciiEThreeFourthree}{0.380\,\ensuremath{\pm}\,0.005}  % experiments/FIGS/data/w2b_leg1_metrics_v2.json :: cells.{r42q1-s1,r42q1-s2,r42q1-s3,r42q1-s4,r42q1-s5}.N1000.s1.accuracy :: sha256 b56bf70be9982044 :: mean +- SD over the row's 5 pretrained models
\newcommand{\FtAccJciiEFourFourthree}{0.481\,\ensuremath{\pm}\,0.004}  % experiments/FIGS/data/w2b_leg1_metrics_v2.json :: cells.{r42q1-s1,r42q1-s2,r42q1-s3,r42q1-s4,r42q1-s5}.N10000.s1.accuracy :: sha256 b56bf70be9982044 :: mean +- SD over the row's 5 pretrained models
\newcommand{\FtAccJciiEFiveFourthree}{0.5219\,\ensuremath{\pm}\,0.0017}  % experiments/FIGS/data/w2b_leg1_metrics_v2.json :: cells.{r42q1-s1,r42q1-s2,r42q1-s3,r42q1-s4,r42q1-s5}.N100000.s1.accuracy :: sha256 b56bf70be9982044 :: mean +- SD over the row's 5 pretrained models
\newcommand{\FtAccJciiESixFourthree}{0.5362\,\ensuremath{\pm}\,0.0007}  % experiments/FIGS/data/w2b_leg1_metrics_v2.json :: cells.{r42q1-s1,r42q1-s2,r42q1-s3,r42q1-s4,r42q1-s5}.N1000000.s1.accuracy :: sha256 b56bf70be9982044 :: mean +- SD over the row's 5 pretrained models
\newcommand{\FtAccJciiEThreeOneseven}{0.332\,\ensuremath{\pm}\,0.004}  % experiments/FIGS/data/w2b_leg1_metrics_v2.json :: cells.{r16q1-s1,r16q1-s2,r16q1-s3,r16q1-s4,r16q1-s5}.N1000.s1.accuracy :: sha256 b56bf70be9982044 :: mean +- SD over the row's 5 pretrained models
\newcommand{\FtAccJciiEFourOneseven}{0.447\,\ensuremath{\pm}\,0.013}  % experiments/FIGS/data/w2b_leg1_metrics_v2.json :: cells.{r16q1-s1,r16q1-s2,r16q1-s3,r16q1-s4,r16q1-s5}.N10000.s1.accuracy :: sha256 b56bf70be9982044 :: mean +- SD over the row's 5 pretrained models
\newcommand{\FtAccJciiEFiveOneseven}{0.5078\,\ensuremath{\pm}\,0.0011}  % experiments/FIGS/data/w2b_leg1_metrics_v2.json :: cells.{r16q1-s1,r16q1-s2,r16q1-s3,r16q1-s4,r16q1-s5}.N100000.s1.accuracy :: sha256 b56bf70be9982044 :: mean +- SD over the row's 5 pretrained models
\newcommand{\FtAccJciiESixOneseven}{0.5227\,\ensuremath{\pm}\,0.0010}  % experiments/FIGS/data/w2b_leg1_metrics_v2.json :: cells.{r16q1-s1,r16q1-s2,r16q1-s3,r16q1-s4,r16q1-s5}.N1000000.s1.accuracy :: sha256 b56bf70be9982044 :: mean +- SD over the row's 5 pretrained models
\newcommand{\FtAccJciiEThreeOnesixtwoMass}{0.451\,\ensuremath{\pm}\,0.007}  % experiments/FIGS/data/w2b_leg1_metrics_v2.json :: cells.{l162mass-s1,l162mass-s2,l162mass-s3,l162mass-s4,l162mass-s5}.N1000.s1.accuracy :: sha256 b56bf70be9982044 :: mean +- SD over the row's 5 pretrained models
\newcommand{\FtAccJciiEFourOnesixtwoMass}{0.515\,\ensuremath{\pm}\,0.004}  % experiments/FIGS/data/w2b_leg1_metrics_v2.json :: cells.{l162mass-s1,l162mass-s2,l162mass-s3,l162mass-s4,l162mass-s5}.N10000.s1.accuracy :: sha256 b56bf70be9982044 :: mean +- SD over the row's 5 pretrained models
\newcommand{\FtAccJciiEFiveOnesixtwoMass}{0.5357\,\ensuremath{\pm}\,0.0007}  % experiments/FIGS/data/w2b_leg1_metrics_v2.json :: cells.{l162mass-s1,l162mass-s2,l162mass-s3,l162mass-s4,l162mass-s5}.N100000.s1.accuracy :: sha256 b56bf70be9982044 :: mean +- SD over the row's 5 pretrained models
\newcommand{\FtAccJciiESixOnesixtwoMass}{0.5409\,\ensuremath{\pm}\,0.0010}  % experiments/FIGS/data/w2b_leg1_metrics_v2.json :: cells.{l162mass-s1,l162mass-s2,l162mass-s3,l162mass-s4,l162mass-s5}.N1000000.s1.accuracy :: sha256 b56bf70be9982044 :: mean +- SD over the row's 5 pretrained models
\newcommand{\FtAccJciiEThreeOnesevenMass}{0.331\,\ensuremath{\pm}\,0.005}  % experiments/FIGS/data/w2b_leg1_metrics_v2.json :: cells.{r16q1mass-s1,r16q1mass-s2,r16q1mass-s3,r16q1mass-s4,r16q1mass-s5}.N1000.s1.accuracy :: sha256 b56bf70be9982044 :: mean +- SD over the row's 5 pretrained models
\newcommand{\FtAccJciiEFourOnesevenMass}{0.4507\,\ensuremath{\pm}\,0.0034}  % experiments/FIGS/data/w2b_leg1_metrics_v2.json :: cells.{r16q1mass-s1,r16q1mass-s2,r16q1mass-s3,r16q1mass-s4,r16q1mass-s5}.N10000.s1.accuracy :: sha256 b56bf70be9982044 :: mean +- SD over the row's 5 pretrained models
\newcommand{\FtAccJciiEFiveOnesevenMass}{0.4993\,\ensuremath{\pm}\,0.0022}  % experiments/FIGS/data/w2b_leg1_metrics_v2.json :: cells.{r16q1mass-s1,r16q1mass-s2,r16q1mass-s3,r16q1mass-s4,r16q1mass-s5}.N100000.s1.accuracy :: sha256 b56bf70be9982044 :: mean +- SD over the row's 5 pretrained models
\newcommand{\FtAccJciiESixOnesevenMass}{0.5166\,\ensuremath{\pm}\,0.0013}  % experiments/FIGS/data/w2b_leg1_metrics_v2.json :: cells.{r16q1mass-s1,r16q1mass-s2,r16q1mass-s3,r16q1mass-s4,r16q1mass-s5}.N1000000.s1.accuracy :: sha256 b56bf70be9982044 :: mean +- SD over the row's 5 pretrained models
\newcommand{\FtAccJciiEThreeRandDrawOne}{0.424}  % experiments/FIGS/data/w2b_leg1_metrics_v2.json :: cells.{rand-d1-s1b}.N1000.s1.accuracy :: sha256 b56bf70be9982044 :: one run
\newcommand{\FtAccJciiEFourRandDrawOne}{0.5070}  % experiments/FIGS/data/w2b_leg1_metrics_v2.json :: cells.{rand-d1-s1b}.N10000.s1.accuracy :: sha256 b56bf70be9982044 :: one run
\newcommand{\FtAccJciiEFiveRandDrawOne}{0.5226}  % experiments/FIGS/data/w2b_leg1_metrics_v2.json :: cells.{rand-d1-s1b}.N100000.s1.accuracy :: sha256 b56bf70be9982044 :: one run
\newcommand{\FtAccJciiESixRandDrawOne}{0.53354}  % experiments/FIGS/data/w2b_leg1_metrics_v2.json :: cells.{rand-d1-s1b}.N1000000.s1.accuracy :: sha256 b56bf70be9982044 :: one run
\newcommand{\FtAccJciiEThreeRandDrawTwo}{0.420}  % experiments/FIGS/data/w2b_leg1_metrics_v2.json :: cells.{rand-d2-s2}.N1000.s1.accuracy :: sha256 b56bf70be9982044 :: one run
\newcommand{\FtAccJciiEFourRandDrawTwo}{0.5074}  % experiments/FIGS/data/w2b_leg1_metrics_v2.json :: cells.{rand-d2-s2}.N10000.s1.accuracy :: sha256 b56bf70be9982044 :: one run
\newcommand{\FtAccJciiEFiveRandDrawTwo}{0.5243}  % experiments/FIGS/data/w2b_leg1_metrics_v2.json :: cells.{rand-d2-s2}.N100000.s1.accuracy :: sha256 b56bf70be9982044 :: one run
\newcommand{\FtAccJciiESixRandDrawTwo}{0.53441}  % experiments/FIGS/data/w2b_leg1_metrics_v2.json :: cells.{rand-d2-s2}.N1000000.s1.accuracy :: sha256 b56bf70be9982044 :: one run
\newcommand{\FtAccJciiEThreeRandDrawThree}{0.423}  % experiments/FIGS/data/w2b_leg1_metrics_v2.json :: cells.{rand-d3-s3}.N1000.s1.accuracy :: sha256 b56bf70be9982044 :: one run
\newcommand{\FtAccJciiEFourRandDrawThree}{0.5114}  % experiments/FIGS/data/w2b_leg1_metrics_v2.json :: cells.{rand-d3-s3}.N10000.s1.accuracy :: sha256 b56bf70be9982044 :: one run
\newcommand{\FtAccJciiEFiveRandDrawThree}{0.5245}  % experiments/FIGS/data/w2b_leg1_metrics_v2.json :: cells.{rand-d3-s3}.N100000.s1.accuracy :: sha256 b56bf70be9982044 :: one run
\newcommand{\FtAccJciiESixRandDrawThree}{0.53473}  % experiments/FIGS/data/w2b_leg1_metrics_v2.json :: cells.{rand-d3-s3}.N1000000.s1.accuracy :: sha256 b56bf70be9982044 :: one run
\newcommand{\FtOmaRatioJciiEThreeOnesixtwoOverOneeighteight}{0.96}  % experiments/FIGS/data/w2b_leg1_metrics_v2.json :: cells.*.N1000.s1.macro_auc_ovr :: sha256 b56bf70be9982044 :: ratio of the seed means of 1 - macro AUC, this row over the reference
\newcommand{\FtOmaRatioJciiEThreeFourthreeOverOneeighteight}{0.84}  % experiments/FIGS/data/w2b_leg1_metrics_v2.json :: cells.*.N1000.s1.macro_auc_ovr :: sha256 b56bf70be9982044 :: ratio of the seed means of 1 - macro AUC, this row over the reference
\newcommand{\FtOmaRatioJciiEThreeOnesevenOverOneeighteight}{0.97}  % experiments/FIGS/data/w2b_leg1_metrics_v2.json :: cells.*.N1000.s1.macro_auc_ovr :: sha256 b56bf70be9982044 :: ratio of the seed means of 1 - macro AUC, this row over the reference
\newcommand{\FtOmaRatioJciiEThreeOnesixtwoMassOverOneeighteight}{0.81}  % experiments/FIGS/data/w2b_leg1_metrics_v2.json :: cells.*.N1000.s1.macro_auc_ovr :: sha256 b56bf70be9982044 :: ratio of the seed means of 1 - macro AUC, this row over the reference
\newcommand{\FtOmaRatioJciiEThreeOnesevenMassOverOneeighteight}{0.98}  % experiments/FIGS/data/w2b_leg1_metrics_v2.json :: cells.*.N1000.s1.macro_auc_ovr :: sha256 b56bf70be9982044 :: ratio of the seed means of 1 - macro AUC, this row over the reference
\newcommand{\FtOmaRatioJciiEThreeRandDrawOneOverOneeighteight}{0.80}  % experiments/FIGS/data/w2b_leg1_metrics_v2.json :: cells.*.N1000.s1.macro_auc_ovr :: sha256 b56bf70be9982044 :: ratio of the seed means of 1 - macro AUC, this row over the reference
\newcommand{\FtOmaRatioJciiEThreeRandDrawTwoOverOneeighteight}{0.87}  % experiments/FIGS/data/w2b_leg1_metrics_v2.json :: cells.*.N1000.s1.macro_auc_ovr :: sha256 b56bf70be9982044 :: ratio of the seed means of 1 - macro AUC, this row over the reference
\newcommand{\FtOmaRatioJciiEThreeRandDrawThreeOverOneeighteight}{0.83}  % experiments/FIGS/data/w2b_leg1_metrics_v2.json :: cells.*.N1000.s1.macro_auc_ovr :: sha256 b56bf70be9982044 :: ratio of the seed means of 1 - macro AUC, this row over the reference
\newcommand{\FtOmaRatioJciiEThreeOneeighteightOverOnesixtwo}{1.05}  % experiments/FIGS/data/w2b_leg1_metrics_v2.json :: cells.*.N1000.s1.macro_auc_ovr :: sha256 b56bf70be9982044 :: ratio of the seed means of 1 - macro AUC, this row over the reference
\newcommand{\FtOmaRatioJciiEThreeFourthreeOverOnesixtwo}{0.88}  % experiments/FIGS/data/w2b_leg1_metrics_v2.json :: cells.*.N1000.s1.macro_auc_ovr :: sha256 b56bf70be9982044 :: ratio of the seed means of 1 - macro AUC, this row over the reference
\newcommand{\FtOmaRatioJciiEThreeOnesevenOverOnesixtwo}{1.02}  % experiments/FIGS/data/w2b_leg1_metrics_v2.json :: cells.*.N1000.s1.macro_auc_ovr :: sha256 b56bf70be9982044 :: ratio of the seed means of 1 - macro AUC, this row over the reference
\newcommand{\FtOmaRatioJciiEThreeOnesixtwoMassOverOnesixtwo}{0.85}  % experiments/FIGS/data/w2b_leg1_metrics_v2.json :: cells.*.N1000.s1.macro_auc_ovr :: sha256 b56bf70be9982044 :: ratio of the seed means of 1 - macro AUC, this row over the reference
\newcommand{\FtOmaRatioJciiEThreeOnesevenMassOverOnesixtwo}{1.02}  % experiments/FIGS/data/w2b_leg1_metrics_v2.json :: cells.*.N1000.s1.macro_auc_ovr :: sha256 b56bf70be9982044 :: ratio of the seed means of 1 - macro AUC, this row over the reference
\newcommand{\FtOmaRatioJciiEThreeRandDrawOneOverOnesixtwo}{0.84}  % experiments/FIGS/data/w2b_leg1_metrics_v2.json :: cells.*.N1000.s1.macro_auc_ovr :: sha256 b56bf70be9982044 :: ratio of the seed means of 1 - macro AUC, this row over the reference
\newcommand{\FtOmaRatioJciiEThreeRandDrawTwoOverOnesixtwo}{0.91}  % experiments/FIGS/data/w2b_leg1_metrics_v2.json :: cells.*.N1000.s1.macro_auc_ovr :: sha256 b56bf70be9982044 :: ratio of the seed means of 1 - macro AUC, this row over the reference
\newcommand{\FtOmaRatioJciiEThreeRandDrawThreeOverOnesixtwo}{0.87}  % experiments/FIGS/data/w2b_leg1_metrics_v2.json :: cells.*.N1000.s1.macro_auc_ovr :: sha256 b56bf70be9982044 :: ratio of the seed means of 1 - macro AUC, this row over the reference
\newcommand{\FtOmaRatioJciiEFourOnesixtwoOverOneeighteight}{0.99}  % experiments/FIGS/data/w2b_leg1_metrics_v2.json :: cells.*.N10000.s1.macro_auc_ovr :: sha256 b56bf70be9982044 :: ratio of the seed means of 1 - macro AUC, this row over the reference
\newcommand{\FtOmaRatioJciiEFourFourthreeOverOneeighteight}{1.20}  % experiments/FIGS/data/w2b_leg1_metrics_v2.json :: cells.*.N10000.s1.macro_auc_ovr :: sha256 b56bf70be9982044 :: ratio of the seed means of 1 - macro AUC, this row over the reference
\newcommand{\FtOmaRatioJciiEFourOnesevenOverOneeighteight}{1.32}  % experiments/FIGS/data/w2b_leg1_metrics_v2.json :: cells.*.N10000.s1.macro_auc_ovr :: sha256 b56bf70be9982044 :: ratio of the seed means of 1 - macro AUC, this row over the reference
\newcommand{\FtOmaRatioJciiEFourOnesixtwoMassOverOneeighteight}{0.99}  % experiments/FIGS/data/w2b_leg1_metrics_v2.json :: cells.*.N10000.s1.macro_auc_ovr :: sha256 b56bf70be9982044 :: ratio of the seed means of 1 - macro AUC, this row over the reference
\newcommand{\FtOmaRatioJciiEFourOnesevenMassOverOneeighteight}{1.30}  % experiments/FIGS/data/w2b_leg1_metrics_v2.json :: cells.*.N10000.s1.macro_auc_ovr :: sha256 b56bf70be9982044 :: ratio of the seed means of 1 - macro AUC, this row over the reference
\newcommand{\FtOmaRatioJciiEFourRandDrawOneOverOneeighteight}{1.17}  % experiments/FIGS/data/w2b_leg1_metrics_v2.json :: cells.*.N10000.s1.macro_auc_ovr :: sha256 b56bf70be9982044 :: ratio of the seed means of 1 - macro AUC, this row over the reference
\newcommand{\FtOmaRatioJciiEFourRandDrawTwoOverOneeighteight}{1.17}  % experiments/FIGS/data/w2b_leg1_metrics_v2.json :: cells.*.N10000.s1.macro_auc_ovr :: sha256 b56bf70be9982044 :: ratio of the seed means of 1 - macro AUC, this row over the reference
\newcommand{\FtOmaRatioJciiEFourRandDrawThreeOverOneeighteight}{1.13}  % experiments/FIGS/data/w2b_leg1_metrics_v2.json :: cells.*.N10000.s1.macro_auc_ovr :: sha256 b56bf70be9982044 :: ratio of the seed means of 1 - macro AUC, this row over the reference
\newcommand{\FtOmaRatioJciiEFourOneeighteightOverOnesixtwo}{1.01}  % experiments/FIGS/data/w2b_leg1_metrics_v2.json :: cells.*.N10000.s1.macro_auc_ovr :: sha256 b56bf70be9982044 :: ratio of the seed means of 1 - macro AUC, this row over the reference
\newcommand{\FtOmaRatioJciiEFourFourthreeOverOnesixtwo}{1.21}  % experiments/FIGS/data/w2b_leg1_metrics_v2.json :: cells.*.N10000.s1.macro_auc_ovr :: sha256 b56bf70be9982044 :: ratio of the seed means of 1 - macro AUC, this row over the reference
\newcommand{\FtOmaRatioJciiEFourOnesevenOverOnesixtwo}{1.32}  % experiments/FIGS/data/w2b_leg1_metrics_v2.json :: cells.*.N10000.s1.macro_auc_ovr :: sha256 b56bf70be9982044 :: ratio of the seed means of 1 - macro AUC, this row over the reference
\newcommand{\FtOmaRatioJciiEFourOnesixtwoMassOverOnesixtwo}{1.00}  % experiments/FIGS/data/w2b_leg1_metrics_v2.json :: cells.*.N10000.s1.macro_auc_ovr :: sha256 b56bf70be9982044 :: ratio of the seed means of 1 - macro AUC, this row over the reference
\newcommand{\FtOmaRatioJciiEFourOnesevenMassOverOnesixtwo}{1.31}  % experiments/FIGS/data/w2b_leg1_metrics_v2.json :: cells.*.N10000.s1.macro_auc_ovr :: sha256 b56bf70be9982044 :: ratio of the seed means of 1 - macro AUC, this row over the reference
\newcommand{\FtOmaRatioJciiEFourRandDrawOneOverOnesixtwo}{1.18}  % experiments/FIGS/data/w2b_leg1_metrics_v2.json :: cells.*.N10000.s1.macro_auc_ovr :: sha256 b56bf70be9982044 :: ratio of the seed means of 1 - macro AUC, this row over the reference
\newcommand{\FtOmaRatioJciiEFourRandDrawTwoOverOnesixtwo}{1.18}  % experiments/FIGS/data/w2b_leg1_metrics_v2.json :: cells.*.N10000.s1.macro_auc_ovr :: sha256 b56bf70be9982044 :: ratio of the seed means of 1 - macro AUC, this row over the reference
\newcommand{\FtOmaRatioJciiEFourRandDrawThreeOverOnesixtwo}{1.14}  % experiments/FIGS/data/w2b_leg1_metrics_v2.json :: cells.*.N10000.s1.macro_auc_ovr :: sha256 b56bf70be9982044 :: ratio of the seed means of 1 - macro AUC, this row over the reference
\newcommand{\FtOmaRatioJciiEFiveOnesixtwoOverOneeighteight}{1.00}  % experiments/FIGS/data/w2b_leg1_metrics_v2.json :: cells.*.N100000.s1.macro_auc_ovr :: sha256 b56bf70be9982044 :: ratio of the seed means of 1 - macro AUC, this row over the reference
\newcommand{\FtOmaRatioJciiEFiveFourthreeOverOneeighteight}{1.09}  % experiments/FIGS/data/w2b_leg1_metrics_v2.json :: cells.*.N100000.s1.macro_auc_ovr :: sha256 b56bf70be9982044 :: ratio of the seed means of 1 - macro AUC, this row over the reference
\newcommand{\FtOmaRatioJciiEFiveOnesevenOverOneeighteight}{1.22}  % experiments/FIGS/data/w2b_leg1_metrics_v2.json :: cells.*.N100000.s1.macro_auc_ovr :: sha256 b56bf70be9982044 :: ratio of the seed means of 1 - macro AUC, this row over the reference
\newcommand{\FtOmaRatioJciiEFiveOnesixtwoMassOverOneeighteight}{1.01}  % experiments/FIGS/data/w2b_leg1_metrics_v2.json :: cells.*.N100000.s1.macro_auc_ovr :: sha256 b56bf70be9982044 :: ratio of the seed means of 1 - macro AUC, this row over the reference
\newcommand{\FtOmaRatioJciiEFiveOnesevenMassOverOneeighteight}{1.26}  % experiments/FIGS/data/w2b_leg1_metrics_v2.json :: cells.*.N100000.s1.macro_auc_ovr :: sha256 b56bf70be9982044 :: ratio of the seed means of 1 - macro AUC, this row over the reference
\newcommand{\FtOmaRatioJciiEFiveRandDrawOneOverOneeighteight}{1.15}  % experiments/FIGS/data/w2b_leg1_metrics_v2.json :: cells.*.N100000.s1.macro_auc_ovr :: sha256 b56bf70be9982044 :: ratio of the seed means of 1 - macro AUC, this row over the reference
\newcommand{\FtOmaRatioJciiEFiveRandDrawTwoOverOneeighteight}{1.12}  % experiments/FIGS/data/w2b_leg1_metrics_v2.json :: cells.*.N100000.s1.macro_auc_ovr :: sha256 b56bf70be9982044 :: ratio of the seed means of 1 - macro AUC, this row over the reference
\newcommand{\FtOmaRatioJciiEFiveRandDrawThreeOverOneeighteight}{1.15}  % experiments/FIGS/data/w2b_leg1_metrics_v2.json :: cells.*.N100000.s1.macro_auc_ovr :: sha256 b56bf70be9982044 :: ratio of the seed means of 1 - macro AUC, this row over the reference
\newcommand{\FtOmaRatioJciiEFiveOneeighteightOverOnesixtwo}{1.00}  % experiments/FIGS/data/w2b_leg1_metrics_v2.json :: cells.*.N100000.s1.macro_auc_ovr :: sha256 b56bf70be9982044 :: ratio of the seed means of 1 - macro AUC, this row over the reference
\newcommand{\FtOmaRatioJciiEFiveFourthreeOverOnesixtwo}{1.09}  % experiments/FIGS/data/w2b_leg1_metrics_v2.json :: cells.*.N100000.s1.macro_auc_ovr :: sha256 b56bf70be9982044 :: ratio of the seed means of 1 - macro AUC, this row over the reference
\newcommand{\FtOmaRatioJciiEFiveOnesevenOverOnesixtwo}{1.22}  % experiments/FIGS/data/w2b_leg1_metrics_v2.json :: cells.*.N100000.s1.macro_auc_ovr :: sha256 b56bf70be9982044 :: ratio of the seed means of 1 - macro AUC, this row over the reference
\newcommand{\FtOmaRatioJciiEFiveOnesixtwoMassOverOnesixtwo}{1.02}  % experiments/FIGS/data/w2b_leg1_metrics_v2.json :: cells.*.N100000.s1.macro_auc_ovr :: sha256 b56bf70be9982044 :: ratio of the seed means of 1 - macro AUC, this row over the reference
\newcommand{\FtOmaRatioJciiEFiveOnesevenMassOverOnesixtwo}{1.26}  % experiments/FIGS/data/w2b_leg1_metrics_v2.json :: cells.*.N100000.s1.macro_auc_ovr :: sha256 b56bf70be9982044 :: ratio of the seed means of 1 - macro AUC, this row over the reference
\newcommand{\FtOmaRatioJciiEFiveRandDrawOneOverOnesixtwo}{1.15}  % experiments/FIGS/data/w2b_leg1_metrics_v2.json :: cells.*.N100000.s1.macro_auc_ovr :: sha256 b56bf70be9982044 :: ratio of the seed means of 1 - macro AUC, this row over the reference
\newcommand{\FtOmaRatioJciiEFiveRandDrawTwoOverOnesixtwo}{1.13}  % experiments/FIGS/data/w2b_leg1_metrics_v2.json :: cells.*.N100000.s1.macro_auc_ovr :: sha256 b56bf70be9982044 :: ratio of the seed means of 1 - macro AUC, this row over the reference
\newcommand{\FtOmaRatioJciiEFiveRandDrawThreeOverOnesixtwo}{1.16}  % experiments/FIGS/data/w2b_leg1_metrics_v2.json :: cells.*.N100000.s1.macro_auc_ovr :: sha256 b56bf70be9982044 :: ratio of the seed means of 1 - macro AUC, this row over the reference
\newcommand{\FtOmaRatioJciiESixOnesixtwoOverOneeighteight}{1.00}  % experiments/FIGS/data/w2b_leg1_metrics_v2.json :: cells.*.N1000000.s1.macro_auc_ovr :: sha256 b56bf70be9982044 :: ratio of the seed means of 1 - macro AUC, this row over the reference
\newcommand{\FtOmaRatioJciiESixFourthreeOverOneeighteight}{1.05}  % experiments/FIGS/data/w2b_leg1_metrics_v2.json :: cells.*.N1000000.s1.macro_auc_ovr :: sha256 b56bf70be9982044 :: ratio of the seed means of 1 - macro AUC, this row over the reference
\newcommand{\FtOmaRatioJciiESixOnesevenOverOneeighteight}{1.14}  % experiments/FIGS/data/w2b_leg1_metrics_v2.json :: cells.*.N1000000.s1.macro_auc_ovr :: sha256 b56bf70be9982044 :: ratio of the seed means of 1 - macro AUC, this row over the reference
\newcommand{\FtOmaRatioJciiESixOnesixtwoMassOverOneeighteight}{1.02}  % experiments/FIGS/data/w2b_leg1_metrics_v2.json :: cells.*.N1000000.s1.macro_auc_ovr :: sha256 b56bf70be9982044 :: ratio of the seed means of 1 - macro AUC, this row over the reference
\newcommand{\FtOmaRatioJciiESixOnesevenMassOverOneeighteight}{1.18}  % experiments/FIGS/data/w2b_leg1_metrics_v2.json :: cells.*.N1000000.s1.macro_auc_ovr :: sha256 b56bf70be9982044 :: ratio of the seed means of 1 - macro AUC, this row over the reference
\newcommand{\FtOmaRatioJciiESixRandDrawOneOverOneeighteight}{1.08}  % experiments/FIGS/data/w2b_leg1_metrics_v2.json :: cells.*.N1000000.s1.macro_auc_ovr :: sha256 b56bf70be9982044 :: ratio of the seed means of 1 - macro AUC, this row over the reference
\newcommand{\FtOmaRatioJciiESixRandDrawTwoOverOneeighteight}{1.07}  % experiments/FIGS/data/w2b_leg1_metrics_v2.json :: cells.*.N1000000.s1.macro_auc_ovr :: sha256 b56bf70be9982044 :: ratio of the seed means of 1 - macro AUC, this row over the reference
\newcommand{\FtOmaRatioJciiESixRandDrawThreeOverOneeighteight}{1.07}  % experiments/FIGS/data/w2b_leg1_metrics_v2.json :: cells.*.N1000000.s1.macro_auc_ovr :: sha256 b56bf70be9982044 :: ratio of the seed means of 1 - macro AUC, this row over the reference
\newcommand{\FtOmaRatioJciiESixOneeighteightOverOnesixtwo}{1.00}  % experiments/FIGS/data/w2b_leg1_metrics_v2.json :: cells.*.N1000000.s1.macro_auc_ovr :: sha256 b56bf70be9982044 :: ratio of the seed means of 1 - macro AUC, this row over the reference
\newcommand{\FtOmaRatioJciiESixFourthreeOverOnesixtwo}{1.04}  % experiments/FIGS/data/w2b_leg1_metrics_v2.json :: cells.*.N1000000.s1.macro_auc_ovr :: sha256 b56bf70be9982044 :: ratio of the seed means of 1 - macro AUC, this row over the reference
\newcommand{\FtOmaRatioJciiESixOnesevenOverOnesixtwo}{1.14}  % experiments/FIGS/data/w2b_leg1_metrics_v2.json :: cells.*.N1000000.s1.macro_auc_ovr :: sha256 b56bf70be9982044 :: ratio of the seed means of 1 - macro AUC, this row over the reference
\newcommand{\FtOmaRatioJciiESixOnesixtwoMassOverOnesixtwo}{1.02}  % experiments/FIGS/data/w2b_leg1_metrics_v2.json :: cells.*.N1000000.s1.macro_auc_ovr :: sha256 b56bf70be9982044 :: ratio of the seed means of 1 - macro AUC, this row over the reference
\newcommand{\FtOmaRatioJciiESixOnesevenMassOverOnesixtwo}{1.18}  % experiments/FIGS/data/w2b_leg1_metrics_v2.json :: cells.*.N1000000.s1.macro_auc_ovr :: sha256 b56bf70be9982044 :: ratio of the seed means of 1 - macro AUC, this row over the reference
\newcommand{\FtOmaRatioJciiESixRandDrawOneOverOnesixtwo}{1.07}  % experiments/FIGS/data/w2b_leg1_metrics_v2.json :: cells.*.N1000000.s1.macro_auc_ovr :: sha256 b56bf70be9982044 :: ratio of the seed means of 1 - macro AUC, this row over the reference
\newcommand{\FtOmaRatioJciiESixRandDrawTwoOverOnesixtwo}{1.07}  % experiments/FIGS/data/w2b_leg1_metrics_v2.json :: cells.*.N1000000.s1.macro_auc_ovr :: sha256 b56bf70be9982044 :: ratio of the seed means of 1 - macro AUC, this row over the reference
\newcommand{\FtOmaRatioJciiESixRandDrawThreeOverOnesixtwo}{1.07}  % experiments/FIGS/data/w2b_leg1_metrics_v2.json :: cells.*.N1000000.s1.macro_auc_ovr :: sha256 b56bf70be9982044 :: ratio of the seed means of 1 - macro AUC, this row over the reference
\newcommand{\PairedFtAccJciiEThreeOnesevenMinusOneeighteight}{\ensuremath{-}0.1132~[\ensuremath{-}0.1214, \ensuremath{-}0.1051]}  % experiments/FIGS/data/w2b_leg1_metrics_v2.json :: cells.{r16q1-s1,r16q1-s2,r16q1-s3,r16q1-s4,r16q1-s5,l188-s1,l188-s2,l188-s3,l188-s4,l188-s5}.N1000.s1.accuracy :: sha256 b56bf70be9982044 :: accuracy, coarse minus fine, paired by run [95% Student-t interval]
\newcommand{\PairedFtAccJciiEFourOnesevenMinusOneeighteight}{\ensuremath{-}0.0730~[\ensuremath{-}0.0902, \ensuremath{-}0.0558]}  % experiments/FIGS/data/w2b_leg1_metrics_v2.json :: cells.{r16q1-s1,r16q1-s2,r16q1-s3,r16q1-s4,r16q1-s5,l188-s1,l188-s2,l188-s3,l188-s4,l188-s5}.N10000.s1.accuracy :: sha256 b56bf70be9982044 :: accuracy, coarse minus fine, paired by run [95% Student-t interval]
\newcommand{\PairedFtAccJciiEFiveOnesevenMinusOneeighteight}{\ensuremath{-}0.0301~[\ensuremath{-}0.0319, \ensuremath{-}0.0284]}  % experiments/FIGS/data/w2b_leg1_metrics_v2.json :: cells.{r16q1-s1,r16q1-s2,r16q1-s3,r16q1-s4,r16q1-s5,l188-s1,l188-s2,l188-s3,l188-s4,l188-s5}.N100000.s1.accuracy :: sha256 b56bf70be9982044 :: accuracy, coarse minus fine, paired by run [95% Student-t interval]
\newcommand{\PairedFtAccJciiESixOnesevenMinusOneeighteight}{\ensuremath{-}0.0216~[\ensuremath{-}0.0232, \ensuremath{-}0.0199]}  % experiments/FIGS/data/w2b_leg1_metrics_v2.json :: cells.{r16q1-s1,r16q1-s2,r16q1-s3,r16q1-s4,r16q1-s5,l188-s1,l188-s2,l188-s3,l188-s4,l188-s5}.N1000000.s1.accuracy :: sha256 b56bf70be9982044 :: accuracy, coarse minus fine, paired by run [95% Student-t interval]
\newcommand{\FtNClassesJcii}{162}  % experiments/FIGS/data/w2b_leg1_metrics_v2.json :: cells.*.*.s1.n_classes_present :: sha256 b56bf70be9982044 :: classes of the fine-tuning task
\newcommand{\FtNTestAucJcii}{125{,}000}  % experiments/FIGS/data/w2b_leg1_metrics_v2.json :: cells.*.*.s1.n_jets_auc :: sha256 b56bf70be9982044 :: test jets the macro AUC is computed on
\newcommand{\FtNTestAccJcii}{500{,}000}  % experiments/FIGS/data/w2b_leg1_metrics_v2.json :: cells.*.*.s1.n_jets :: sha256 b56bf70be9982044 :: test jets the accuracy is computed on
\newcommand{\FtAucJcEThreeOneeighteight}{0.950\,\ensuremath{\pm}\,0.004}  % experiments/FIGS/data/w2b_leg2_metrics_v2.json :: cells.{l188-s1,l188-s2,l188-s3,l188-s4,l188-s5}.N1000.s1.macro_auc_ovr :: sha256 9cd059d3e2f51798 :: mean +- SD over the row's 5 pretrained models
\newcommand{\FtAucJcEFourOneeighteight}{0.9765\,\ensuremath{\pm}\,0.0008}  % experiments/FIGS/data/w2b_leg2_metrics_v2.json :: cells.{l188-s1,l188-s2,l188-s3,l188-s4,l188-s5}.N10000.s1.macro_auc_ovr :: sha256 9cd059d3e2f51798 :: mean +- SD over the row's 5 pretrained models
\newcommand{\FtAucJcEFiveOneeighteight}{0.9833\,\ensuremath{\pm}\,0.0002}  % experiments/FIGS/data/w2b_leg2_metrics_v2.json :: cells.{l188-s1,l188-s2,l188-s3,l188-s4,l188-s5}.N100000.s1.macro_auc_ovr :: sha256 9cd059d3e2f51798 :: mean +- SD over the row's 5 pretrained models
\newcommand{\FtAucJcESixOneeighteight}{0.9858\,\ensuremath{\pm}\,0.0001}  % experiments/FIGS/data/w2b_leg2_metrics_v2.json :: cells.{l188-s1,l188-s2,l188-s3,l188-s4,l188-s5}.N1000000.s1.macro_auc_ovr :: sha256 9cd059d3e2f51798 :: mean +- SD over the row's 5 pretrained models
\newcommand{\FtAucJcEThreeOnesixtwo}{0.9526\,\ensuremath{\pm}\,0.0020}  % experiments/FIGS/data/w2b_leg2_metrics_v2.json :: cells.{l162-s1b,l162-s2,l162-s3,l162-s4,l162-s5}.N1000.s1.macro_auc_ovr :: sha256 9cd059d3e2f51798 :: mean +- SD over the row's 5 pretrained models
\newcommand{\FtAucJcEFourOnesixtwo}{0.9770\,\ensuremath{\pm}\,0.0007}  % experiments/FIGS/data/w2b_leg2_metrics_v2.json :: cells.{l162-s1b,l162-s2,l162-s3,l162-s4,l162-s5}.N10000.s1.macro_auc_ovr :: sha256 9cd059d3e2f51798 :: mean +- SD over the row's 5 pretrained models
\newcommand{\FtAucJcEFiveOnesixtwo}{0.9836\,\ensuremath{\pm}\,0.0001}  % experiments/FIGS/data/w2b_leg2_metrics_v2.json :: cells.{l162-s1b,l162-s2,l162-s3,l162-s4,l162-s5}.N100000.s1.macro_auc_ovr :: sha256 9cd059d3e2f51798 :: mean +- SD over the row's 5 pretrained models
\newcommand{\FtAucJcESixOnesixtwo}{0.9858\,\ensuremath{\pm}\,0.0001}  % experiments/FIGS/data/w2b_leg2_metrics_v2.json :: cells.{l162-s1b,l162-s2,l162-s3,l162-s4,l162-s5}.N1000000.s1.macro_auc_ovr :: sha256 9cd059d3e2f51798 :: mean +- SD over the row's 5 pretrained models
\newcommand{\FtAucJcEThreeFourthree}{0.950\,\ensuremath{\pm}\,0.004}  % experiments/FIGS/data/w2b_leg2_metrics_v2.json :: cells.{r42q1-s1,r42q1-s2,r42q1-s3,r42q1-s4,r42q1-s5}.N1000.s1.macro_auc_ovr :: sha256 9cd059d3e2f51798 :: mean +- SD over the row's 5 pretrained models
\newcommand{\FtAucJcEFourFourthree}{0.9769\,\ensuremath{\pm}\,0.0005}  % experiments/FIGS/data/w2b_leg2_metrics_v2.json :: cells.{r42q1-s1,r42q1-s2,r42q1-s3,r42q1-s4,r42q1-s5}.N10000.s1.macro_auc_ovr :: sha256 9cd059d3e2f51798 :: mean +- SD over the row's 5 pretrained models
\newcommand{\FtAucJcEFiveFourthree}{0.9831\,\ensuremath{\pm}\,0.0002}  % experiments/FIGS/data/w2b_leg2_metrics_v2.json :: cells.{r42q1-s1,r42q1-s2,r42q1-s3,r42q1-s4,r42q1-s5}.N100000.s1.macro_auc_ovr :: sha256 9cd059d3e2f51798 :: mean +- SD over the row's 5 pretrained models
\newcommand{\FtAucJcESixFourthree}{0.9855\,\ensuremath{\pm}\,0.0001}  % experiments/FIGS/data/w2b_leg2_metrics_v2.json :: cells.{r42q1-s1,r42q1-s2,r42q1-s3,r42q1-s4,r42q1-s5}.N1000000.s1.macro_auc_ovr :: sha256 9cd059d3e2f51798 :: mean +- SD over the row's 5 pretrained models
\newcommand{\FtAucJcEThreeOneseven}{0.9470\,\ensuremath{\pm}\,0.0024}  % experiments/FIGS/data/w2b_leg2_metrics_v2.json :: cells.{r16q1-s1,r16q1-s2,r16q1-s3,r16q1-s4,r16q1-s5}.N1000.s1.macro_auc_ovr :: sha256 9cd059d3e2f51798 :: mean +- SD over the row's 5 pretrained models
\newcommand{\FtAucJcEFourOneseven}{0.9749\,\ensuremath{\pm}\,0.0006}  % experiments/FIGS/data/w2b_leg2_metrics_v2.json :: cells.{r16q1-s1,r16q1-s2,r16q1-s3,r16q1-s4,r16q1-s5}.N10000.s1.macro_auc_ovr :: sha256 9cd059d3e2f51798 :: mean +- SD over the row's 5 pretrained models
\newcommand{\FtAucJcEFiveOneseven}{0.9821\,\ensuremath{\pm}\,\ensuremath{<}0.0001}  % experiments/FIGS/data/w2b_leg2_metrics_v2.json :: cells.{r16q1-s1,r16q1-s2,r16q1-s3,r16q1-s4,r16q1-s5}.N100000.s1.macro_auc_ovr :: sha256 9cd059d3e2f51798 :: mean +- SD over the row's 5 pretrained models
\newcommand{\FtAucJcESixOneseven}{0.9846\,\ensuremath{\pm}\,0.0001}  % experiments/FIGS/data/w2b_leg2_metrics_v2.json :: cells.{r16q1-s1,r16q1-s2,r16q1-s3,r16q1-s4,r16q1-s5}.N1000000.s1.macro_auc_ovr :: sha256 9cd059d3e2f51798 :: mean +- SD over the row's 5 pretrained models
\newcommand{\FtAucJcEThreeOnesixtwoMass}{0.9569\,\ensuremath{\pm}\,0.0016}  % experiments/FIGS/data/w2b_leg2_metrics_v2.json :: cells.{l162mass-s1,l162mass-s2,l162mass-s3,l162mass-s4,l162mass-s5}.N1000.s1.macro_auc_ovr :: sha256 9cd059d3e2f51798 :: mean +- SD over the row's 5 pretrained models
\newcommand{\FtAucJcEFourOnesixtwoMass}{0.9768\,\ensuremath{\pm}\,0.0005}  % experiments/FIGS/data/w2b_leg2_metrics_v2.json :: cells.{l162mass-s1,l162mass-s2,l162mass-s3,l162mass-s4,l162mass-s5}.N10000.s1.macro_auc_ovr :: sha256 9cd059d3e2f51798 :: mean +- SD over the row's 5 pretrained models
\newcommand{\FtAucJcEFiveOnesixtwoMass}{0.9834\,\ensuremath{\pm}\,0.0001}  % experiments/FIGS/data/w2b_leg2_metrics_v2.json :: cells.{l162mass-s1,l162mass-s2,l162mass-s3,l162mass-s4,l162mass-s5}.N100000.s1.macro_auc_ovr :: sha256 9cd059d3e2f51798 :: mean +- SD over the row's 5 pretrained models
\newcommand{\FtAucJcESixOnesixtwoMass}{0.9856\,\ensuremath{\pm}\,0.0001}  % experiments/FIGS/data/w2b_leg2_metrics_v2.json :: cells.{l162mass-s1,l162mass-s2,l162mass-s3,l162mass-s4,l162mass-s5}.N1000000.s1.macro_auc_ovr :: sha256 9cd059d3e2f51798 :: mean +- SD over the row's 5 pretrained models
\newcommand{\FtAucJcEThreeOnesevenMass}{0.9460\,\ensuremath{\pm}\,0.0029}  % experiments/FIGS/data/w2b_leg2_metrics_v2.json :: cells.{r16q1mass-s1,r16q1mass-s2,r16q1mass-s3,r16q1mass-s4,r16q1mass-s5}.N1000.s1.macro_auc_ovr :: sha256 9cd059d3e2f51798 :: mean +- SD over the row's 5 pretrained models
\newcommand{\FtAucJcEFourOnesevenMass}{0.9737\,\ensuremath{\pm}\,0.0008}  % experiments/FIGS/data/w2b_leg2_metrics_v2.json :: cells.{r16q1mass-s1,r16q1mass-s2,r16q1mass-s3,r16q1mass-s4,r16q1mass-s5}.N10000.s1.macro_auc_ovr :: sha256 9cd059d3e2f51798 :: mean +- SD over the row's 5 pretrained models
\newcommand{\FtAucJcEFiveOnesevenMass}{0.9817\,\ensuremath{\pm}\,0.0002}  % experiments/FIGS/data/w2b_leg2_metrics_v2.json :: cells.{r16q1mass-s1,r16q1mass-s2,r16q1mass-s3,r16q1mass-s4,r16q1mass-s5}.N100000.s1.macro_auc_ovr :: sha256 9cd059d3e2f51798 :: mean +- SD over the row's 5 pretrained models
\newcommand{\FtAucJcESixOnesevenMass}{0.9842\,\ensuremath{\pm}\,0.0002}  % experiments/FIGS/data/w2b_leg2_metrics_v2.json :: cells.{r16q1mass-s1,r16q1mass-s2,r16q1mass-s3,r16q1mass-s4,r16q1mass-s5}.N1000000.s1.macro_auc_ovr :: sha256 9cd059d3e2f51798 :: mean +- SD over the row's 5 pretrained models
\newcommand{\FtAucJcEThreeRandDrawOne}{0.9477}  % experiments/FIGS/data/w2b_leg2_metrics_v2.json :: cells.{rand-d1-s1b}.N1000.s1.macro_auc_ovr :: sha256 9cd059d3e2f51798 :: one run
\newcommand{\FtAucJcEFourRandDrawOne}{0.9745}  % experiments/FIGS/data/w2b_leg2_metrics_v2.json :: cells.{rand-d1-s1b}.N10000.s1.macro_auc_ovr :: sha256 9cd059d3e2f51798 :: one run
\newcommand{\FtAucJcEFiveRandDrawOne}{0.98265}  % experiments/FIGS/data/w2b_leg2_metrics_v2.json :: cells.{rand-d1-s1b}.N100000.s1.macro_auc_ovr :: sha256 9cd059d3e2f51798 :: one run
\newcommand{\FtAucJcESixRandDrawOne}{0.98513}  % experiments/FIGS/data/w2b_leg2_metrics_v2.json :: cells.{rand-d1-s1b}.N1000000.s1.macro_auc_ovr :: sha256 9cd059d3e2f51798 :: one run
\newcommand{\FtAucJcEThreeRandDrawTwo}{0.9501}  % experiments/FIGS/data/w2b_leg2_metrics_v2.json :: cells.{rand-d2-s2}.N1000.s1.macro_auc_ovr :: sha256 9cd059d3e2f51798 :: one run
\newcommand{\FtAucJcEFourRandDrawTwo}{0.9756}  % experiments/FIGS/data/w2b_leg2_metrics_v2.json :: cells.{rand-d2-s2}.N10000.s1.macro_auc_ovr :: sha256 9cd059d3e2f51798 :: one run
\newcommand{\FtAucJcEFiveRandDrawTwo}{0.98278}  % experiments/FIGS/data/w2b_leg2_metrics_v2.json :: cells.{rand-d2-s2}.N100000.s1.macro_auc_ovr :: sha256 9cd059d3e2f51798 :: one run
\newcommand{\FtAucJcESixRandDrawTwo}{0.98512}  % experiments/FIGS/data/w2b_leg2_metrics_v2.json :: cells.{rand-d2-s2}.N1000000.s1.macro_auc_ovr :: sha256 9cd059d3e2f51798 :: one run
\newcommand{\FtAucJcEThreeRandDrawThree}{0.9470}  % experiments/FIGS/data/w2b_leg2_metrics_v2.json :: cells.{rand-d3-s3}.N1000.s1.macro_auc_ovr :: sha256 9cd059d3e2f51798 :: one run
\newcommand{\FtAucJcEFourRandDrawThree}{0.9743}  % experiments/FIGS/data/w2b_leg2_metrics_v2.json :: cells.{rand-d3-s3}.N10000.s1.macro_auc_ovr :: sha256 9cd059d3e2f51798 :: one run
\newcommand{\FtAucJcEFiveRandDrawThree}{0.98253}  % experiments/FIGS/data/w2b_leg2_metrics_v2.json :: cells.{rand-d3-s3}.N100000.s1.macro_auc_ovr :: sha256 9cd059d3e2f51798 :: one run
\newcommand{\FtAucJcESixRandDrawThree}{0.98506}  % experiments/FIGS/data/w2b_leg2_metrics_v2.json :: cells.{rand-d3-s3}.N1000000.s1.macro_auc_ovr :: sha256 9cd059d3e2f51798 :: one run
\newcommand{\FtAccJcEThreeOneeighteight}{0.7183\,\ensuremath{\pm}\,0.0032}  % experiments/FIGS/data/w2b_leg2_metrics_v2.json :: cells.{l188-s1,l188-s2,l188-s3,l188-s4,l188-s5}.N1000.s1.accuracy :: sha256 9cd059d3e2f51798 :: mean +- SD over the row's 5 pretrained models
\newcommand{\FtAccJcEFourOneeighteight}{0.8039\,\ensuremath{\pm}\,0.0030}  % experiments/FIGS/data/w2b_leg2_metrics_v2.json :: cells.{l188-s1,l188-s2,l188-s3,l188-s4,l188-s5}.N10000.s1.accuracy :: sha256 9cd059d3e2f51798 :: mean +- SD over the row's 5 pretrained models
\newcommand{\FtAccJcEFiveOneeighteight}{0.8351\,\ensuremath{\pm}\,0.0008}  % experiments/FIGS/data/w2b_leg2_metrics_v2.json :: cells.{l188-s1,l188-s2,l188-s3,l188-s4,l188-s5}.N100000.s1.accuracy :: sha256 9cd059d3e2f51798 :: mean +- SD over the row's 5 pretrained models
\newcommand{\FtAccJcESixOneeighteight}{0.8486\,\ensuremath{\pm}\,0.0005}  % experiments/FIGS/data/w2b_leg2_metrics_v2.json :: cells.{l188-s1,l188-s2,l188-s3,l188-s4,l188-s5}.N1000000.s1.accuracy :: sha256 9cd059d3e2f51798 :: mean +- SD over the row's 5 pretrained models
\newcommand{\FtAccJcEThreeOnesixtwo}{0.7268\,\ensuremath{\pm}\,0.0017}  % experiments/FIGS/data/w2b_leg2_metrics_v2.json :: cells.{l162-s1b,l162-s2,l162-s3,l162-s4,l162-s5}.N1000.s1.accuracy :: sha256 9cd059d3e2f51798 :: mean +- SD over the row's 5 pretrained models
\newcommand{\FtAccJcEFourOnesixtwo}{0.8069\,\ensuremath{\pm}\,0.0025}  % experiments/FIGS/data/w2b_leg2_metrics_v2.json :: cells.{l162-s1b,l162-s2,l162-s3,l162-s4,l162-s5}.N10000.s1.accuracy :: sha256 9cd059d3e2f51798 :: mean +- SD over the row's 5 pretrained models
\newcommand{\FtAccJcEFiveOnesixtwo}{0.8363\,\ensuremath{\pm}\,0.0008}  % experiments/FIGS/data/w2b_leg2_metrics_v2.json :: cells.{l162-s1b,l162-s2,l162-s3,l162-s4,l162-s5}.N100000.s1.accuracy :: sha256 9cd059d3e2f51798 :: mean +- SD over the row's 5 pretrained models
\newcommand{\FtAccJcESixOnesixtwo}{0.8485\,\ensuremath{\pm}\,0.0005}  % experiments/FIGS/data/w2b_leg2_metrics_v2.json :: cells.{l162-s1b,l162-s2,l162-s3,l162-s4,l162-s5}.N1000000.s1.accuracy :: sha256 9cd059d3e2f51798 :: mean +- SD over the row's 5 pretrained models
\newcommand{\FtAccJcEThreeFourthree}{0.722\,\ensuremath{\pm}\,0.007}  % experiments/FIGS/data/w2b_leg2_metrics_v2.json :: cells.{r42q1-s1,r42q1-s2,r42q1-s3,r42q1-s4,r42q1-s5}.N1000.s1.accuracy :: sha256 9cd059d3e2f51798 :: mean +- SD over the row's 5 pretrained models
\newcommand{\FtAccJcEFourFourthree}{0.8050\,\ensuremath{\pm}\,0.0026}  % experiments/FIGS/data/w2b_leg2_metrics_v2.json :: cells.{r42q1-s1,r42q1-s2,r42q1-s3,r42q1-s4,r42q1-s5}.N10000.s1.accuracy :: sha256 9cd059d3e2f51798 :: mean +- SD over the row's 5 pretrained models
\newcommand{\FtAccJcEFiveFourthree}{0.8336\,\ensuremath{\pm}\,0.0008}  % experiments/FIGS/data/w2b_leg2_metrics_v2.json :: cells.{r42q1-s1,r42q1-s2,r42q1-s3,r42q1-s4,r42q1-s5}.N100000.s1.accuracy :: sha256 9cd059d3e2f51798 :: mean +- SD over the row's 5 pretrained models
\newcommand{\FtAccJcESixFourthree}{0.8468\,\ensuremath{\pm}\,0.0003}  % experiments/FIGS/data/w2b_leg2_metrics_v2.json :: cells.{r42q1-s1,r42q1-s2,r42q1-s3,r42q1-s4,r42q1-s5}.N1000000.s1.accuracy :: sha256 9cd059d3e2f51798 :: mean +- SD over the row's 5 pretrained models
\newcommand{\FtAccJcEThreeOneseven}{0.714\,\ensuremath{\pm}\,0.005}  % experiments/FIGS/data/w2b_leg2_metrics_v2.json :: cells.{r16q1-s1,r16q1-s2,r16q1-s3,r16q1-s4,r16q1-s5}.N1000.s1.accuracy :: sha256 9cd059d3e2f51798 :: mean +- SD over the row's 5 pretrained models
\newcommand{\FtAccJcEFourOneseven}{0.7964\,\ensuremath{\pm}\,0.0019}  % experiments/FIGS/data/w2b_leg2_metrics_v2.json :: cells.{r16q1-s1,r16q1-s2,r16q1-s3,r16q1-s4,r16q1-s5}.N10000.s1.accuracy :: sha256 9cd059d3e2f51798 :: mean +- SD over the row's 5 pretrained models
\newcommand{\FtAccJcEFiveOneseven}{0.8286\,\ensuremath{\pm}\,0.0006}  % experiments/FIGS/data/w2b_leg2_metrics_v2.json :: cells.{r16q1-s1,r16q1-s2,r16q1-s3,r16q1-s4,r16q1-s5}.N100000.s1.accuracy :: sha256 9cd059d3e2f51798 :: mean +- SD over the row's 5 pretrained models
\newcommand{\FtAccJcESixOneseven}{0.8415\,\ensuremath{\pm}\,0.0007}  % experiments/FIGS/data/w2b_leg2_metrics_v2.json :: cells.{r16q1-s1,r16q1-s2,r16q1-s3,r16q1-s4,r16q1-s5}.N1000000.s1.accuracy :: sha256 9cd059d3e2f51798 :: mean +- SD over the row's 5 pretrained models
\newcommand{\FtAccJcEThreeOnesixtwoMass}{0.728\,\ensuremath{\pm}\,0.004}  % experiments/FIGS/data/w2b_leg2_metrics_v2.json :: cells.{l162mass-s1,l162mass-s2,l162mass-s3,l162mass-s4,l162mass-s5}.N1000.s1.accuracy :: sha256 9cd059d3e2f51798 :: mean +- SD over the row's 5 pretrained models
\newcommand{\FtAccJcEFourOnesixtwoMass}{0.8043\,\ensuremath{\pm}\,0.0020}  % experiments/FIGS/data/w2b_leg2_metrics_v2.json :: cells.{l162mass-s1,l162mass-s2,l162mass-s3,l162mass-s4,l162mass-s5}.N10000.s1.accuracy :: sha256 9cd059d3e2f51798 :: mean +- SD over the row's 5 pretrained models
\newcommand{\FtAccJcEFiveOnesixtwoMass}{0.8349\,\ensuremath{\pm}\,0.0005}  % experiments/FIGS/data/w2b_leg2_metrics_v2.json :: cells.{l162mass-s1,l162mass-s2,l162mass-s3,l162mass-s4,l162mass-s5}.N100000.s1.accuracy :: sha256 9cd059d3e2f51798 :: mean +- SD over the row's 5 pretrained models
\newcommand{\FtAccJcESixOnesixtwoMass}{0.8475\,\ensuremath{\pm}\,0.0009}  % experiments/FIGS/data/w2b_leg2_metrics_v2.json :: cells.{l162mass-s1,l162mass-s2,l162mass-s3,l162mass-s4,l162mass-s5}.N1000000.s1.accuracy :: sha256 9cd059d3e2f51798 :: mean +- SD over the row's 5 pretrained models
\newcommand{\FtAccJcEThreeOnesevenMass}{0.702\,\ensuremath{\pm}\,0.005}  % experiments/FIGS/data/w2b_leg2_metrics_v2.json :: cells.{r16q1mass-s1,r16q1mass-s2,r16q1mass-s3,r16q1mass-s4,r16q1mass-s5}.N1000.s1.accuracy :: sha256 9cd059d3e2f51798 :: mean +- SD over the row's 5 pretrained models
\newcommand{\FtAccJcEFourOnesevenMass}{0.7909\,\ensuremath{\pm}\,0.0031}  % experiments/FIGS/data/w2b_leg2_metrics_v2.json :: cells.{r16q1mass-s1,r16q1mass-s2,r16q1mass-s3,r16q1mass-s4,r16q1mass-s5}.N10000.s1.accuracy :: sha256 9cd059d3e2f51798 :: mean +- SD over the row's 5 pretrained models
\newcommand{\FtAccJcEFiveOnesevenMass}{0.8265\,\ensuremath{\pm}\,0.0007}  % experiments/FIGS/data/w2b_leg2_metrics_v2.json :: cells.{r16q1mass-s1,r16q1mass-s2,r16q1mass-s3,r16q1mass-s4,r16q1mass-s5}.N100000.s1.accuracy :: sha256 9cd059d3e2f51798 :: mean +- SD over the row's 5 pretrained models
\newcommand{\FtAccJcESixOnesevenMass}{0.8390\,\ensuremath{\pm}\,0.0011}  % experiments/FIGS/data/w2b_leg2_metrics_v2.json :: cells.{r16q1mass-s1,r16q1mass-s2,r16q1mass-s3,r16q1mass-s4,r16q1mass-s5}.N1000000.s1.accuracy :: sha256 9cd059d3e2f51798 :: mean +- SD over the row's 5 pretrained models
\newcommand{\FtAccJcEThreeRandDrawOne}{0.7049}  % experiments/FIGS/data/w2b_leg2_metrics_v2.json :: cells.{rand-d1-s1b}.N1000.s1.accuracy :: sha256 9cd059d3e2f51798 :: one run
\newcommand{\FtAccJcEFourRandDrawOne}{0.7974}  % experiments/FIGS/data/w2b_leg2_metrics_v2.json :: cells.{rand-d1-s1b}.N10000.s1.accuracy :: sha256 9cd059d3e2f51798 :: one run
\newcommand{\FtAccJcEFiveRandDrawOne}{0.8326}  % experiments/FIGS/data/w2b_leg2_metrics_v2.json :: cells.{rand-d1-s1b}.N100000.s1.accuracy :: sha256 9cd059d3e2f51798 :: one run
\newcommand{\FtAccJcESixRandDrawOne}{0.84467}  % experiments/FIGS/data/w2b_leg2_metrics_v2.json :: cells.{rand-d1-s1b}.N1000000.s1.accuracy :: sha256 9cd059d3e2f51798 :: one run
\newcommand{\FtAccJcEThreeRandDrawTwo}{0.7135}  % experiments/FIGS/data/w2b_leg2_metrics_v2.json :: cells.{rand-d2-s2}.N1000.s1.accuracy :: sha256 9cd059d3e2f51798 :: one run
\newcommand{\FtAccJcEFourRandDrawTwo}{0.7998}  % experiments/FIGS/data/w2b_leg2_metrics_v2.json :: cells.{rand-d2-s2}.N10000.s1.accuracy :: sha256 9cd059d3e2f51798 :: one run
\newcommand{\FtAccJcEFiveRandDrawTwo}{0.8332}  % experiments/FIGS/data/w2b_leg2_metrics_v2.json :: cells.{rand-d2-s2}.N100000.s1.accuracy :: sha256 9cd059d3e2f51798 :: one run
\newcommand{\FtAccJcESixRandDrawTwo}{0.84449}  % experiments/FIGS/data/w2b_leg2_metrics_v2.json :: cells.{rand-d2-s2}.N1000000.s1.accuracy :: sha256 9cd059d3e2f51798 :: one run
\newcommand{\FtAccJcEThreeRandDrawThree}{0.7090}  % experiments/FIGS/data/w2b_leg2_metrics_v2.json :: cells.{rand-d3-s3}.N1000.s1.accuracy :: sha256 9cd059d3e2f51798 :: one run
\newcommand{\FtAccJcEFourRandDrawThree}{0.7948}  % experiments/FIGS/data/w2b_leg2_metrics_v2.json :: cells.{rand-d3-s3}.N10000.s1.accuracy :: sha256 9cd059d3e2f51798 :: one run
\newcommand{\FtAccJcEFiveRandDrawThree}{0.8310}  % experiments/FIGS/data/w2b_leg2_metrics_v2.json :: cells.{rand-d3-s3}.N100000.s1.accuracy :: sha256 9cd059d3e2f51798 :: one run
\newcommand{\FtAccJcESixRandDrawThree}{0.84409}  % experiments/FIGS/data/w2b_leg2_metrics_v2.json :: cells.{rand-d3-s3}.N1000000.s1.accuracy :: sha256 9cd059d3e2f51798 :: one run
\newcommand{\FtOmaRatioJcEThreeOnesixtwoOverOneeighteight}{0.96}  % experiments/FIGS/data/w2b_leg2_metrics_v2.json :: cells.*.N1000.s1.macro_auc_ovr :: sha256 9cd059d3e2f51798 :: ratio of the seed means of 1 - macro AUC, this row over the reference
\newcommand{\FtOmaRatioJcEThreeFourthreeOverOneeighteight}{1.01}  % experiments/FIGS/data/w2b_leg2_metrics_v2.json :: cells.*.N1000.s1.macro_auc_ovr :: sha256 9cd059d3e2f51798 :: ratio of the seed means of 1 - macro AUC, this row over the reference
\newcommand{\FtOmaRatioJcEThreeOnesevenOverOneeighteight}{1.07}  % experiments/FIGS/data/w2b_leg2_metrics_v2.json :: cells.*.N1000.s1.macro_auc_ovr :: sha256 9cd059d3e2f51798 :: ratio of the seed means of 1 - macro AUC, this row over the reference
\newcommand{\FtOmaRatioJcEThreeOnesixtwoMassOverOneeighteight}{0.87}  % experiments/FIGS/data/w2b_leg2_metrics_v2.json :: cells.*.N1000.s1.macro_auc_ovr :: sha256 9cd059d3e2f51798 :: ratio of the seed means of 1 - macro AUC, this row over the reference
\newcommand{\FtOmaRatioJcEThreeOnesevenMassOverOneeighteight}{1.09}  % experiments/FIGS/data/w2b_leg2_metrics_v2.json :: cells.*.N1000.s1.macro_auc_ovr :: sha256 9cd059d3e2f51798 :: ratio of the seed means of 1 - macro AUC, this row over the reference
\newcommand{\FtOmaRatioJcEThreeRandDrawOneOverOneeighteight}{1.06}  % experiments/FIGS/data/w2b_leg2_metrics_v2.json :: cells.*.N1000.s1.macro_auc_ovr :: sha256 9cd059d3e2f51798 :: ratio of the seed means of 1 - macro AUC, this row over the reference
\newcommand{\FtOmaRatioJcEThreeRandDrawTwoOverOneeighteight}{1.01}  % experiments/FIGS/data/w2b_leg2_metrics_v2.json :: cells.*.N1000.s1.macro_auc_ovr :: sha256 9cd059d3e2f51798 :: ratio of the seed means of 1 - macro AUC, this row over the reference
\newcommand{\FtOmaRatioJcEThreeRandDrawThreeOverOneeighteight}{1.07}  % experiments/FIGS/data/w2b_leg2_metrics_v2.json :: cells.*.N1000.s1.macro_auc_ovr :: sha256 9cd059d3e2f51798 :: ratio of the seed means of 1 - macro AUC, this row over the reference
\newcommand{\FtOmaRatioJcEThreeOneeighteightOverOnesixtwo}{1.05}  % experiments/FIGS/data/w2b_leg2_metrics_v2.json :: cells.*.N1000.s1.macro_auc_ovr :: sha256 9cd059d3e2f51798 :: ratio of the seed means of 1 - macro AUC, this row over the reference
\newcommand{\FtOmaRatioJcEThreeFourthreeOverOnesixtwo}{1.05}  % experiments/FIGS/data/w2b_leg2_metrics_v2.json :: cells.*.N1000.s1.macro_auc_ovr :: sha256 9cd059d3e2f51798 :: ratio of the seed means of 1 - macro AUC, this row over the reference
\newcommand{\FtOmaRatioJcEThreeOnesevenOverOnesixtwo}{1.12}  % experiments/FIGS/data/w2b_leg2_metrics_v2.json :: cells.*.N1000.s1.macro_auc_ovr :: sha256 9cd059d3e2f51798 :: ratio of the seed means of 1 - macro AUC, this row over the reference
\newcommand{\FtOmaRatioJcEThreeOnesixtwoMassOverOnesixtwo}{0.91}  % experiments/FIGS/data/w2b_leg2_metrics_v2.json :: cells.*.N1000.s1.macro_auc_ovr :: sha256 9cd059d3e2f51798 :: ratio of the seed means of 1 - macro AUC, this row over the reference
\newcommand{\FtOmaRatioJcEThreeOnesevenMassOverOnesixtwo}{1.14}  % experiments/FIGS/data/w2b_leg2_metrics_v2.json :: cells.*.N1000.s1.macro_auc_ovr :: sha256 9cd059d3e2f51798 :: ratio of the seed means of 1 - macro AUC, this row over the reference
\newcommand{\FtOmaRatioJcEThreeRandDrawOneOverOnesixtwo}{1.10}  % experiments/FIGS/data/w2b_leg2_metrics_v2.json :: cells.*.N1000.s1.macro_auc_ovr :: sha256 9cd059d3e2f51798 :: ratio of the seed means of 1 - macro AUC, this row over the reference
\newcommand{\FtOmaRatioJcEThreeRandDrawTwoOverOnesixtwo}{1.05}  % experiments/FIGS/data/w2b_leg2_metrics_v2.json :: cells.*.N1000.s1.macro_auc_ovr :: sha256 9cd059d3e2f51798 :: ratio of the seed means of 1 - macro AUC, this row over the reference
\newcommand{\FtOmaRatioJcEThreeRandDrawThreeOverOnesixtwo}{1.12}  % experiments/FIGS/data/w2b_leg2_metrics_v2.json :: cells.*.N1000.s1.macro_auc_ovr :: sha256 9cd059d3e2f51798 :: ratio of the seed means of 1 - macro AUC, this row over the reference
\newcommand{\FtOmaRatioJcEFourOnesixtwoOverOneeighteight}{0.98}  % experiments/FIGS/data/w2b_leg2_metrics_v2.json :: cells.*.N10000.s1.macro_auc_ovr :: sha256 9cd059d3e2f51798 :: ratio of the seed means of 1 - macro AUC, this row over the reference
\newcommand{\FtOmaRatioJcEFourFourthreeOverOneeighteight}{0.98}  % experiments/FIGS/data/w2b_leg2_metrics_v2.json :: cells.*.N10000.s1.macro_auc_ovr :: sha256 9cd059d3e2f51798 :: ratio of the seed means of 1 - macro AUC, this row over the reference
\newcommand{\FtOmaRatioJcEFourOnesevenOverOneeighteight}{1.07}  % experiments/FIGS/data/w2b_leg2_metrics_v2.json :: cells.*.N10000.s1.macro_auc_ovr :: sha256 9cd059d3e2f51798 :: ratio of the seed means of 1 - macro AUC, this row over the reference
\newcommand{\FtOmaRatioJcEFourOnesixtwoMassOverOneeighteight}{0.99}  % experiments/FIGS/data/w2b_leg2_metrics_v2.json :: cells.*.N10000.s1.macro_auc_ovr :: sha256 9cd059d3e2f51798 :: ratio of the seed means of 1 - macro AUC, this row over the reference
\newcommand{\FtOmaRatioJcEFourOnesevenMassOverOneeighteight}{1.12}  % experiments/FIGS/data/w2b_leg2_metrics_v2.json :: cells.*.N10000.s1.macro_auc_ovr :: sha256 9cd059d3e2f51798 :: ratio of the seed means of 1 - macro AUC, this row over the reference
\newcommand{\FtOmaRatioJcEFourRandDrawOneOverOneeighteight}{1.08}  % experiments/FIGS/data/w2b_leg2_metrics_v2.json :: cells.*.N10000.s1.macro_auc_ovr :: sha256 9cd059d3e2f51798 :: ratio of the seed means of 1 - macro AUC, this row over the reference
\newcommand{\FtOmaRatioJcEFourRandDrawTwoOverOneeighteight}{1.04}  % experiments/FIGS/data/w2b_leg2_metrics_v2.json :: cells.*.N10000.s1.macro_auc_ovr :: sha256 9cd059d3e2f51798 :: ratio of the seed means of 1 - macro AUC, this row over the reference
\newcommand{\FtOmaRatioJcEFourRandDrawThreeOverOneeighteight}{1.09}  % experiments/FIGS/data/w2b_leg2_metrics_v2.json :: cells.*.N10000.s1.macro_auc_ovr :: sha256 9cd059d3e2f51798 :: ratio of the seed means of 1 - macro AUC, this row over the reference
\newcommand{\FtOmaRatioJcEFourOneeighteightOverOnesixtwo}{1.02}  % experiments/FIGS/data/w2b_leg2_metrics_v2.json :: cells.*.N10000.s1.macro_auc_ovr :: sha256 9cd059d3e2f51798 :: ratio of the seed means of 1 - macro AUC, this row over the reference
\newcommand{\FtOmaRatioJcEFourFourthreeOverOnesixtwo}{1.01}  % experiments/FIGS/data/w2b_leg2_metrics_v2.json :: cells.*.N10000.s1.macro_auc_ovr :: sha256 9cd059d3e2f51798 :: ratio of the seed means of 1 - macro AUC, this row over the reference
\newcommand{\FtOmaRatioJcEFourOnesevenOverOnesixtwo}{1.09}  % experiments/FIGS/data/w2b_leg2_metrics_v2.json :: cells.*.N10000.s1.macro_auc_ovr :: sha256 9cd059d3e2f51798 :: ratio of the seed means of 1 - macro AUC, this row over the reference
\newcommand{\FtOmaRatioJcEFourOnesixtwoMassOverOnesixtwo}{1.01}  % experiments/FIGS/data/w2b_leg2_metrics_v2.json :: cells.*.N10000.s1.macro_auc_ovr :: sha256 9cd059d3e2f51798 :: ratio of the seed means of 1 - macro AUC, this row over the reference
\newcommand{\FtOmaRatioJcEFourOnesevenMassOverOnesixtwo}{1.14}  % experiments/FIGS/data/w2b_leg2_metrics_v2.json :: cells.*.N10000.s1.macro_auc_ovr :: sha256 9cd059d3e2f51798 :: ratio of the seed means of 1 - macro AUC, this row over the reference
\newcommand{\FtOmaRatioJcEFourRandDrawOneOverOnesixtwo}{1.11}  % experiments/FIGS/data/w2b_leg2_metrics_v2.json :: cells.*.N10000.s1.macro_auc_ovr :: sha256 9cd059d3e2f51798 :: ratio of the seed means of 1 - macro AUC, this row over the reference
\newcommand{\FtOmaRatioJcEFourRandDrawTwoOverOnesixtwo}{1.06}  % experiments/FIGS/data/w2b_leg2_metrics_v2.json :: cells.*.N10000.s1.macro_auc_ovr :: sha256 9cd059d3e2f51798 :: ratio of the seed means of 1 - macro AUC, this row over the reference
\newcommand{\FtOmaRatioJcEFourRandDrawThreeOverOnesixtwo}{1.12}  % experiments/FIGS/data/w2b_leg2_metrics_v2.json :: cells.*.N10000.s1.macro_auc_ovr :: sha256 9cd059d3e2f51798 :: ratio of the seed means of 1 - macro AUC, this row over the reference
\newcommand{\FtOmaRatioJcEFiveOnesixtwoOverOneeighteight}{0.98}  % experiments/FIGS/data/w2b_leg2_metrics_v2.json :: cells.*.N100000.s1.macro_auc_ovr :: sha256 9cd059d3e2f51798 :: ratio of the seed means of 1 - macro AUC, this row over the reference
\newcommand{\FtOmaRatioJcEFiveFourthreeOverOneeighteight}{1.02}  % experiments/FIGS/data/w2b_leg2_metrics_v2.json :: cells.*.N100000.s1.macro_auc_ovr :: sha256 9cd059d3e2f51798 :: ratio of the seed means of 1 - macro AUC, this row over the reference
\newcommand{\FtOmaRatioJcEFiveOnesevenOverOneeighteight}{1.07}  % experiments/FIGS/data/w2b_leg2_metrics_v2.json :: cells.*.N100000.s1.macro_auc_ovr :: sha256 9cd059d3e2f51798 :: ratio of the seed means of 1 - macro AUC, this row over the reference
\newcommand{\FtOmaRatioJcEFiveOnesixtwoMassOverOneeighteight}{1.00}  % experiments/FIGS/data/w2b_leg2_metrics_v2.json :: cells.*.N100000.s1.macro_auc_ovr :: sha256 9cd059d3e2f51798 :: ratio of the seed means of 1 - macro AUC, this row over the reference
\newcommand{\FtOmaRatioJcEFiveOnesevenMassOverOneeighteight}{1.10}  % experiments/FIGS/data/w2b_leg2_metrics_v2.json :: cells.*.N100000.s1.macro_auc_ovr :: sha256 9cd059d3e2f51798 :: ratio of the seed means of 1 - macro AUC, this row over the reference
\newcommand{\FtOmaRatioJcEFiveRandDrawOneOverOneeighteight}{1.04}  % experiments/FIGS/data/w2b_leg2_metrics_v2.json :: cells.*.N100000.s1.macro_auc_ovr :: sha256 9cd059d3e2f51798 :: ratio of the seed means of 1 - macro AUC, this row over the reference
\newcommand{\FtOmaRatioJcEFiveRandDrawTwoOverOneeighteight}{1.03}  % experiments/FIGS/data/w2b_leg2_metrics_v2.json :: cells.*.N100000.s1.macro_auc_ovr :: sha256 9cd059d3e2f51798 :: ratio of the seed means of 1 - macro AUC, this row over the reference
\newcommand{\FtOmaRatioJcEFiveRandDrawThreeOverOneeighteight}{1.05}  % experiments/FIGS/data/w2b_leg2_metrics_v2.json :: cells.*.N100000.s1.macro_auc_ovr :: sha256 9cd059d3e2f51798 :: ratio of the seed means of 1 - macro AUC, this row over the reference
\newcommand{\FtOmaRatioJcEFiveOneeighteightOverOnesixtwo}{1.02}  % experiments/FIGS/data/w2b_leg2_metrics_v2.json :: cells.*.N100000.s1.macro_auc_ovr :: sha256 9cd059d3e2f51798 :: ratio of the seed means of 1 - macro AUC, this row over the reference
\newcommand{\FtOmaRatioJcEFiveFourthreeOverOnesixtwo}{1.03}  % experiments/FIGS/data/w2b_leg2_metrics_v2.json :: cells.*.N100000.s1.macro_auc_ovr :: sha256 9cd059d3e2f51798 :: ratio of the seed means of 1 - macro AUC, this row over the reference
\newcommand{\FtOmaRatioJcEFiveOnesevenOverOnesixtwo}{1.09}  % experiments/FIGS/data/w2b_leg2_metrics_v2.json :: cells.*.N100000.s1.macro_auc_ovr :: sha256 9cd059d3e2f51798 :: ratio of the seed means of 1 - macro AUC, this row over the reference
\newcommand{\FtOmaRatioJcEFiveOnesixtwoMassOverOnesixtwo}{1.02}  % experiments/FIGS/data/w2b_leg2_metrics_v2.json :: cells.*.N100000.s1.macro_auc_ovr :: sha256 9cd059d3e2f51798 :: ratio of the seed means of 1 - macro AUC, this row over the reference
\newcommand{\FtOmaRatioJcEFiveOnesevenMassOverOnesixtwo}{1.11}  % experiments/FIGS/data/w2b_leg2_metrics_v2.json :: cells.*.N100000.s1.macro_auc_ovr :: sha256 9cd059d3e2f51798 :: ratio of the seed means of 1 - macro AUC, this row over the reference
\newcommand{\FtOmaRatioJcEFiveRandDrawOneOverOnesixtwo}{1.06}  % experiments/FIGS/data/w2b_leg2_metrics_v2.json :: cells.*.N100000.s1.macro_auc_ovr :: sha256 9cd059d3e2f51798 :: ratio of the seed means of 1 - macro AUC, this row over the reference
\newcommand{\FtOmaRatioJcEFiveRandDrawTwoOverOnesixtwo}{1.05}  % experiments/FIGS/data/w2b_leg2_metrics_v2.json :: cells.*.N100000.s1.macro_auc_ovr :: sha256 9cd059d3e2f51798 :: ratio of the seed means of 1 - macro AUC, this row over the reference
\newcommand{\FtOmaRatioJcEFiveRandDrawThreeOverOnesixtwo}{1.07}  % experiments/FIGS/data/w2b_leg2_metrics_v2.json :: cells.*.N100000.s1.macro_auc_ovr :: sha256 9cd059d3e2f51798 :: ratio of the seed means of 1 - macro AUC, this row over the reference
\newcommand{\FtOmaRatioJcESixOnesixtwoOverOneeighteight}{1.00}  % experiments/FIGS/data/w2b_leg2_metrics_v2.json :: cells.*.N1000000.s1.macro_auc_ovr :: sha256 9cd059d3e2f51798 :: ratio of the seed means of 1 - macro AUC, this row over the reference
\newcommand{\FtOmaRatioJcESixFourthreeOverOneeighteight}{1.02}  % experiments/FIGS/data/w2b_leg2_metrics_v2.json :: cells.*.N1000000.s1.macro_auc_ovr :: sha256 9cd059d3e2f51798 :: ratio of the seed means of 1 - macro AUC, this row over the reference
\newcommand{\FtOmaRatioJcESixOnesevenOverOneeighteight}{1.08}  % experiments/FIGS/data/w2b_leg2_metrics_v2.json :: cells.*.N1000000.s1.macro_auc_ovr :: sha256 9cd059d3e2f51798 :: ratio of the seed means of 1 - macro AUC, this row over the reference
\newcommand{\FtOmaRatioJcESixOnesixtwoMassOverOneeighteight}{1.01}  % experiments/FIGS/data/w2b_leg2_metrics_v2.json :: cells.*.N1000000.s1.macro_auc_ovr :: sha256 9cd059d3e2f51798 :: ratio of the seed means of 1 - macro AUC, this row over the reference
\newcommand{\FtOmaRatioJcESixOnesevenMassOverOneeighteight}{1.11}  % experiments/FIGS/data/w2b_leg2_metrics_v2.json :: cells.*.N1000000.s1.macro_auc_ovr :: sha256 9cd059d3e2f51798 :: ratio of the seed means of 1 - macro AUC, this row over the reference
\newcommand{\FtOmaRatioJcESixRandDrawOneOverOneeighteight}{1.05}  % experiments/FIGS/data/w2b_leg2_metrics_v2.json :: cells.*.N1000000.s1.macro_auc_ovr :: sha256 9cd059d3e2f51798 :: ratio of the seed means of 1 - macro AUC, this row over the reference
\newcommand{\FtOmaRatioJcESixRandDrawTwoOverOneeighteight}{1.05}  % experiments/FIGS/data/w2b_leg2_metrics_v2.json :: cells.*.N1000000.s1.macro_auc_ovr :: sha256 9cd059d3e2f51798 :: ratio of the seed means of 1 - macro AUC, this row over the reference
\newcommand{\FtOmaRatioJcESixRandDrawThreeOverOneeighteight}{1.05}  % experiments/FIGS/data/w2b_leg2_metrics_v2.json :: cells.*.N1000000.s1.macro_auc_ovr :: sha256 9cd059d3e2f51798 :: ratio of the seed means of 1 - macro AUC, this row over the reference
\newcommand{\FtOmaRatioJcESixOneeighteightOverOnesixtwo}{1.00}  % experiments/FIGS/data/w2b_leg2_metrics_v2.json :: cells.*.N1000000.s1.macro_auc_ovr :: sha256 9cd059d3e2f51798 :: ratio of the seed means of 1 - macro AUC, this row over the reference
\newcommand{\FtOmaRatioJcESixFourthreeOverOnesixtwo}{1.02}  % experiments/FIGS/data/w2b_leg2_metrics_v2.json :: cells.*.N1000000.s1.macro_auc_ovr :: sha256 9cd059d3e2f51798 :: ratio of the seed means of 1 - macro AUC, this row over the reference
\newcommand{\FtOmaRatioJcESixOnesevenOverOnesixtwo}{1.08}  % experiments/FIGS/data/w2b_leg2_metrics_v2.json :: cells.*.N1000000.s1.macro_auc_ovr :: sha256 9cd059d3e2f51798 :: ratio of the seed means of 1 - macro AUC, this row over the reference
\newcommand{\FtOmaRatioJcESixOnesixtwoMassOverOnesixtwo}{1.01}  % experiments/FIGS/data/w2b_leg2_metrics_v2.json :: cells.*.N1000000.s1.macro_auc_ovr :: sha256 9cd059d3e2f51798 :: ratio of the seed means of 1 - macro AUC, this row over the reference
\newcommand{\FtOmaRatioJcESixOnesevenMassOverOnesixtwo}{1.11}  % experiments/FIGS/data/w2b_leg2_metrics_v2.json :: cells.*.N1000000.s1.macro_auc_ovr :: sha256 9cd059d3e2f51798 :: ratio of the seed means of 1 - macro AUC, this row over the reference
\newcommand{\FtOmaRatioJcESixRandDrawOneOverOnesixtwo}{1.05}  % experiments/FIGS/data/w2b_leg2_metrics_v2.json :: cells.*.N1000000.s1.macro_auc_ovr :: sha256 9cd059d3e2f51798 :: ratio of the seed means of 1 - macro AUC, this row over the reference
\newcommand{\FtOmaRatioJcESixRandDrawTwoOverOnesixtwo}{1.05}  % experiments/FIGS/data/w2b_leg2_metrics_v2.json :: cells.*.N1000000.s1.macro_auc_ovr :: sha256 9cd059d3e2f51798 :: ratio of the seed means of 1 - macro AUC, this row over the reference
\newcommand{\FtOmaRatioJcESixRandDrawThreeOverOnesixtwo}{1.05}  % experiments/FIGS/data/w2b_leg2_metrics_v2.json :: cells.*.N1000000.s1.macro_auc_ovr :: sha256 9cd059d3e2f51798 :: ratio of the seed means of 1 - macro AUC, this row over the reference
\newcommand{\PairedFtAccJcEThreeOnesevenMinusOneeighteight}{\ensuremath{-}0.0040~[\ensuremath{-}0.0144, +0.0064]}  % experiments/FIGS/data/w2b_leg2_metrics_v2.json :: cells.{r16q1-s1,r16q1-s2,r16q1-s3,r16q1-s4,r16q1-s5,l188-s1,l188-s2,l188-s3,l188-s4,l188-s5}.N1000.s1.accuracy :: sha256 9cd059d3e2f51798 :: accuracy, coarse minus fine, paired by run [95% Student-t interval]
\newcommand{\PairedFtAccJcEFourOnesevenMinusOneeighteight}{\ensuremath{-}0.0075~[\ensuremath{-}0.0118, \ensuremath{-}0.0033]}  % experiments/FIGS/data/w2b_leg2_metrics_v2.json :: cells.{r16q1-s1,r16q1-s2,r16q1-s3,r16q1-s4,r16q1-s5,l188-s1,l188-s2,l188-s3,l188-s4,l188-s5}.N10000.s1.accuracy :: sha256 9cd059d3e2f51798 :: accuracy, coarse minus fine, paired by run [95% Student-t interval]
\newcommand{\PairedFtAccJcEFiveOnesevenMinusOneeighteight}{\ensuremath{-}0.0066~[\ensuremath{-}0.0083, \ensuremath{-}0.0049]}  % experiments/FIGS/data/w2b_leg2_metrics_v2.json :: cells.{r16q1-s1,r16q1-s2,r16q1-s3,r16q1-s4,r16q1-s5,l188-s1,l188-s2,l188-s3,l188-s4,l188-s5}.N100000.s1.accuracy :: sha256 9cd059d3e2f51798 :: accuracy, coarse minus fine, paired by run [95% Student-t interval]
\newcommand{\PairedFtAccJcESixOnesevenMinusOneeighteight}{\ensuremath{-}0.0071~[\ensuremath{-}0.0080, \ensuremath{-}0.0062]}  % experiments/FIGS/data/w2b_leg2_metrics_v2.json :: cells.{r16q1-s1,r16q1-s2,r16q1-s3,r16q1-s4,r16q1-s5,l188-s1,l188-s2,l188-s3,l188-s4,l188-s5}.N1000000.s1.accuracy :: sha256 9cd059d3e2f51798 :: accuracy, coarse minus fine, paired by run [95% Student-t interval]
\newcommand{\FtNClassesJc}{10}  % experiments/FIGS/data/w2b_leg2_metrics_v2.json :: cells.*.*.s1.n_classes_present :: sha256 9cd059d3e2f51798 :: classes of the fine-tuning task
\newcommand{\FtNTestAucJc}{500{,}000}  % experiments/FIGS/data/w2b_leg2_metrics_v2.json :: cells.*.*.s1.n_jets_auc :: sha256 9cd059d3e2f51798 :: test jets the macro AUC is computed on
\newcommand{\FtNTestAccJc}{2{,}000{,}000}  % experiments/FIGS/data/w2b_leg2_metrics_v2.json :: cells.*.*.s1.n_jets :: sha256 9cd059d3e2f51798 :: test jets the accuracy is computed on
\newcommand{\FtSizeEThree}{\ensuremath{10^{3}}}  % experiments/FIGS/data/w2b_leg1_metrics_v2.json :: cells.*.N1000 :: sha256 b56bf70be9982044 :: fine-tuning training jets
\newcommand{\FtSizeEFour}{\ensuremath{10^{4}}}  % experiments/FIGS/data/w2b_leg1_metrics_v2.json :: cells.*.N10000 :: sha256 b56bf70be9982044 :: fine-tuning training jets
\newcommand{\FtSizeEFive}{\ensuremath{10^{5}}}  % experiments/FIGS/data/w2b_leg1_metrics_v2.json :: cells.*.N100000 :: sha256 b56bf70be9982044 :: fine-tuning training jets
\newcommand{\FtSizeESix}{\ensuremath{10^{6}}}  % experiments/FIGS/data/w2b_leg1_metrics_v2.json :: cells.*.N1000000 :: sha256 b56bf70be9982044 :: fine-tuning training jets
\newcommand{\FtNRunsWord}{five}  % experiments/FIGS/data/w2b_leg1_metrics_v2.json :: cells (pretrained models per vocabulary row) :: sha256 b56bf70be9982044 :: pretraining runs per vocabulary
\newcommand{\PairedProbeBcVsRestLinearOnesevenOverOnesixtwo}{1.64~[1.52, 1.78]}  % experiments/FIGS/data/paired_v1err/probes/ratios.json :: ratios[0].ratio, ratios[0].ci95 :: sha256 14dd51402305d147 :: 17 over 162, bc_vs_rest linear 1-auc: paired geometric mean [95% interval]
\newcommand{\PairedProbeBcVsRestMlpOnesevenOverOnesixtwo}{1.50~[1.39, 1.62]}  % experiments/FIGS/data/paired_v1err/probes/ratios.json :: ratios[3].ratio, ratios[3].ci95 :: sha256 14dd51402305d147 :: 17 over 162, bc_vs_rest mlp 1-auc: paired geometric mean [95% interval]
\newcommand{\PairedProbeBvcFourprongLinearOnesixtwoOverOneeighteight}{1.01~[0.95, 1.08]}  % experiments/FIGS/data/paired_v1err/probes/ratios.json :: ratios[6].ratio, ratios[6].ci95 :: sha256 14dd51402305d147 :: 162 over 188, bvc_4prong linear 1-auc: paired geometric mean [95% interval]
\newcommand{\PairedProbeBvcFourprongLinearFourthreeOverOneeighteight}{1.10~[1.01, 1.20]}  % experiments/FIGS/data/paired_v1err/probes/ratios.json :: ratios[7].ratio, ratios[7].ci95 :: sha256 14dd51402305d147 :: 43 over 188, bvc_4prong linear 1-auc: paired geometric mean [95% interval]
\newcommand{\PairedProbeBvcFourprongLinearOnesevenOverOneeighteight}{4.5~[3.5, 5.7]}  % experiments/FIGS/data/paired_v1err/probes/ratios.json :: ratios[8].ratio, ratios[8].ci95 :: sha256 14dd51402305d147 :: 17 over 188, bvc_4prong linear 1-auc: paired geometric mean [95% interval]
\newcommand{\PairedProbeBvcFourprongLinearFourthreeOverOnesixtwo}{1.09~[0.98, 1.20]}  % experiments/FIGS/data/paired_v1err/probes/ratios.json :: ratios[9].ratio, ratios[9].ci95 :: sha256 14dd51402305d147 :: 43 over 162, bvc_4prong linear 1-auc: paired geometric mean [95% interval]
\newcommand{\PairedProbeBvcFourprongLinearOnesevenOverOnesixtwo}{4.4~[3.4, 5.6]}  % experiments/FIGS/data/paired_v1err/probes/ratios.json :: ratios[10].ratio, ratios[10].ci95 :: sha256 14dd51402305d147 :: 17 over 162, bvc_4prong linear 1-auc: paired geometric mean [95% interval]
\newcommand{\PairedProbeBvcFourprongLinearOnesevenOverFourthree}{4.0~[3.1, 5.4]}  % experiments/FIGS/data/paired_v1err/probes/ratios.json :: ratios[11].ratio, ratios[11].ci95 :: sha256 14dd51402305d147 :: 17 over 43, bvc_4prong linear 1-auc: paired geometric mean [95% interval]
\newcommand{\PairedProbeBvcFourprongLinearOnesixtwoMassOverOnesixtwo}{1.05~[0.97, 1.13]}  % experiments/FIGS/data/paired_v1err/probes/ratios.json :: ratios[12].ratio, ratios[12].ci95 :: sha256 14dd51402305d147 :: 162+mass over 162, bvc_4prong linear 1-auc: paired geometric mean [95% interval]
\newcommand{\PairedProbeBvcFourprongLinearOnesevenMassOverOneseven}{2.1~[1.6, 2.7]}  % experiments/FIGS/data/paired_v1err/probes/ratios.json :: ratios[13].ratio, ratios[13].ci95 :: sha256 14dd51402305d147 :: 17+mass over 17, bvc_4prong linear 1-auc: paired geometric mean [95% interval]
\newcommand{\PairedProbeBvcFourprongLinearRandomOverOneeighteight}{1.54~[1.37, 1.74]}  % experiments/FIGS/data/paired_v1err/probes/ratios.json :: ratios[14].ratio, ratios[14].ci95 :: sha256 14dd51402305d147 :: random over 188, bvc_4prong linear 1-auc: paired geometric mean [95% interval]
\newcommand{\PairedProbeBvcFourprongLinearRandomOverOnesixtwo}{1.55~[1.36, 1.77]}  % experiments/FIGS/data/paired_v1err/probes/ratios.json :: ratios[15].ratio, ratios[15].ci95 :: sha256 14dd51402305d147 :: random over 162, bvc_4prong linear 1-auc: paired geometric mean [95% interval]
\newcommand{\PairedProbeBvcFourprongLinearRandomOverFourthree}{1.39~[1.27, 1.52]}  % experiments/FIGS/data/paired_v1err/probes/ratios.json :: ratios[16].ratio, ratios[16].ci95 :: sha256 14dd51402305d147 :: random over 43, bvc_4prong linear 1-auc: paired geometric mean [95% interval]
\newcommand{\PairedProbeBvcFourprongLinearRandomOverOneseven}{0.34~[0.20, 0.58]}  % experiments/FIGS/data/paired_v1err/probes/ratios.json :: ratios[17].ratio, ratios[17].ci95 :: sha256 14dd51402305d147 :: random over 17, bvc_4prong linear 1-auc: paired geometric mean [95% interval]
\newcommand{\PairedProbeBvcFourprongLinearRandomDrawTwoOverRandomDrawOne}{1.05~[0.94, 1.18]}  % experiments/FIGS/data/paired_v1err/probes/ratios.json :: ratios[18].ratio, ratios[18].ci95 :: sha256 14dd51402305d147 :: random draw 2 over random draw 1, bvc_4prong linear 1-auc: paired geometric mean [95% interval]
\newcommand{\PairedProbeBvcFourprongLinearRandomDrawThreeOverRandomDrawOne}{0.95~[0.85, 1.06]}  % experiments/FIGS/data/paired_v1err/probes/ratios.json :: ratios[19].ratio, ratios[19].ci95 :: sha256 14dd51402305d147 :: random draw 3 over random draw 1, bvc_4prong linear 1-auc: paired geometric mean [95% interval]
\newcommand{\PairedProbeBvcFourprongLinearRandomDrawThreeOverRandomDrawTwo}{0.90~[0.79, 1.02]}  % experiments/FIGS/data/paired_v1err/probes/ratios.json :: ratios[20].ratio, ratios[20].ci95 :: sha256 14dd51402305d147 :: random draw 3 over random draw 2, bvc_4prong linear 1-auc: paired geometric mean [95% interval]
\newcommand{\PairedProbeBvcFourprongMlpOnesixtwoOverOneeighteight}{1.00~[0.93, 1.07]}  % experiments/FIGS/data/paired_v1err/probes/ratios.json :: ratios[66].ratio, ratios[66].ci95 :: sha256 14dd51402305d147 :: 162 over 188, bvc_4prong mlp 1-auc: paired geometric mean [95% interval]
\newcommand{\PairedProbeBvcFourprongMlpFourthreeOverOneeighteight}{1.12~[1.03, 1.22]}  % experiments/FIGS/data/paired_v1err/probes/ratios.json :: ratios[67].ratio, ratios[67].ci95 :: sha256 14dd51402305d147 :: 43 over 188, bvc_4prong mlp 1-auc: paired geometric mean [95% interval]
\newcommand{\PairedProbeBvcFourprongMlpOnesevenOverOneeighteight}{3.2~[2.6, 4.1]}  % experiments/FIGS/data/paired_v1err/probes/ratios.json :: ratios[68].ratio, ratios[68].ci95 :: sha256 14dd51402305d147 :: 17 over 188, bvc_4prong mlp 1-auc: paired geometric mean [95% interval]
\newcommand{\PairedProbeBvcFourprongMlpFourthreeOverOnesixtwo}{1.12~[1.01, 1.26]}  % experiments/FIGS/data/paired_v1err/probes/ratios.json :: ratios[69].ratio, ratios[69].ci95 :: sha256 14dd51402305d147 :: 43 over 162, bvc_4prong mlp 1-auc: paired geometric mean [95% interval]
\newcommand{\PairedProbeBvcFourprongMlpOnesevenOverOnesixtwo}{3.3~[2.6, 4.0]}  % experiments/FIGS/data/paired_v1err/probes/ratios.json :: ratios[70].ratio, ratios[70].ci95 :: sha256 14dd51402305d147 :: 17 over 162, bvc_4prong mlp 1-auc: paired geometric mean [95% interval]
\newcommand{\PairedProbeBvcFourprongMlpOnesevenOverFourthree}{2.9~[2.3, 3.7]}  % experiments/FIGS/data/paired_v1err/probes/ratios.json :: ratios[71].ratio, ratios[71].ci95 :: sha256 14dd51402305d147 :: 17 over 43, bvc_4prong mlp 1-auc: paired geometric mean [95% interval]
\newcommand{\PairedProbeBvcFourprongMlpOnesixtwoMassOverOnesixtwo}{1.07~[0.99, 1.16]}  % experiments/FIGS/data/paired_v1err/probes/ratios.json :: ratios[72].ratio, ratios[72].ci95 :: sha256 14dd51402305d147 :: 162+mass over 162, bvc_4prong mlp 1-auc: paired geometric mean [95% interval]
\newcommand{\PairedProbeBvcFourprongMlpOnesevenMassOverOneseven}{1.97~[1.46, 2.67]}  % experiments/FIGS/data/paired_v1err/probes/ratios.json :: ratios[73].ratio, ratios[73].ci95 :: sha256 14dd51402305d147 :: 17+mass over 17, bvc_4prong mlp 1-auc: paired geometric mean [95% interval]
\newcommand{\PairedProbeBvcFourprongMlpRandomOverOneeighteight}{1.47~[1.30, 1.65]}  % experiments/FIGS/data/paired_v1err/probes/ratios.json :: ratios[74].ratio, ratios[74].ci95 :: sha256 14dd51402305d147 :: random over 188, bvc_4prong mlp 1-auc: paired geometric mean [95% interval]
\newcommand{\PairedProbeBvcFourprongMlpRandomOverOnesixtwo}{1.49~[1.24, 1.79]}  % experiments/FIGS/data/paired_v1err/probes/ratios.json :: ratios[75].ratio, ratios[75].ci95 :: sha256 14dd51402305d147 :: random over 162, bvc_4prong mlp 1-auc: paired geometric mean [95% interval]
\newcommand{\PairedProbeBvcFourprongMlpRandomOverFourthree}{1.28~[1.16, 1.43]}  % experiments/FIGS/data/paired_v1err/probes/ratios.json :: ratios[76].ratio, ratios[76].ci95 :: sha256 14dd51402305d147 :: random over 43, bvc_4prong mlp 1-auc: paired geometric mean [95% interval]
\newcommand{\PairedProbeBvcFourprongMlpRandomOverOneseven}{0.43~[0.28, 0.67]}  % experiments/FIGS/data/paired_v1err/probes/ratios.json :: ratios[77].ratio, ratios[77].ci95 :: sha256 14dd51402305d147 :: random over 17, bvc_4prong mlp 1-auc: paired geometric mean [95% interval]
\newcommand{\PairedProbeBvcFourprongMlpRandomDrawTwoOverRandomDrawOne}{1.00~[0.90, 1.10]}  % experiments/FIGS/data/paired_v1err/probes/ratios.json :: ratios[78].ratio, ratios[78].ci95 :: sha256 14dd51402305d147 :: random draw 2 over random draw 1, bvc_4prong mlp 1-auc: paired geometric mean [95% interval]
\newcommand{\PairedProbeBvcFourprongMlpRandomDrawThreeOverRandomDrawOne}{0.89~[0.79, 0.99]}  % experiments/FIGS/data/paired_v1err/probes/ratios.json :: ratios[79].ratio, ratios[79].ci95 :: sha256 14dd51402305d147 :: random draw 3 over random draw 1, bvc_4prong mlp 1-auc: paired geometric mean [95% interval]
\newcommand{\PairedProbeBvcFourprongMlpRandomDrawThreeOverRandomDrawTwo}{0.89~[0.79, 0.99]}  % experiments/FIGS/data/paired_v1err/probes/ratios.json :: ratios[80].ratio, ratios[80].ci95 :: sha256 14dd51402305d147 :: random draw 3 over random draw 2, bvc_4prong mlp 1-auc: paired geometric mean [95% interval]
\newcommand{\PairedProbeBvcQcdLinearOnesixtwoOverOneeighteight}{1.15~[1.07, 1.23]}  % experiments/FIGS/data/paired_v1err/probes/ratios.json :: ratios[126].ratio, ratios[126].ci95 :: sha256 14dd51402305d147 :: 162 over 188, bvc_qcd linear 1-auc: paired geometric mean [95% interval]
\newcommand{\PairedProbeBvcQcdLinearFourthreeOverOneeighteight}{1.16~[1.08, 1.25]}  % experiments/FIGS/data/paired_v1err/probes/ratios.json :: ratios[127].ratio, ratios[127].ci95 :: sha256 14dd51402305d147 :: 43 over 188, bvc_qcd linear 1-auc: paired geometric mean [95% interval]
\newcommand{\PairedProbeBvcQcdLinearOnesevenOverOneeighteight}{3.0~[2.5, 3.7]}  % experiments/FIGS/data/paired_v1err/probes/ratios.json :: ratios[128].ratio, ratios[128].ci95 :: sha256 14dd51402305d147 :: 17 over 188, bvc_qcd linear 1-auc: paired geometric mean [95% interval]
\newcommand{\PairedProbeBvcQcdLinearFourthreeOverOnesixtwo}{1.01~[0.96, 1.07]}  % experiments/FIGS/data/paired_v1err/probes/ratios.json :: ratios[129].ratio, ratios[129].ci95 :: sha256 14dd51402305d147 :: 43 over 162, bvc_qcd linear 1-auc: paired geometric mean [95% interval]
\newcommand{\PairedProbeBvcQcdLinearOnesevenOverOnesixtwo}{2.6~[2.2, 3.1]}  % experiments/FIGS/data/paired_v1err/probes/ratios.json :: ratios[130].ratio, ratios[130].ci95 :: sha256 14dd51402305d147 :: 17 over 162, bvc_qcd linear 1-auc: paired geometric mean [95% interval]
\newcommand{\PairedProbeBvcQcdLinearOnesevenOverFourthree}{2.6~[2.2, 3.0]}  % experiments/FIGS/data/paired_v1err/probes/ratios.json :: ratios[131].ratio, ratios[131].ci95 :: sha256 14dd51402305d147 :: 17 over 43, bvc_qcd linear 1-auc: paired geometric mean [95% interval]
\newcommand{\PairedProbeBvcQcdLinearOnesixtwoMassOverOnesixtwo}{1.060~[1.002, 1.123]}  % experiments/FIGS/data/paired_v1err/probes/ratios.json :: ratios[132].ratio, ratios[132].ci95 :: sha256 14dd51402305d147 :: 162+mass over 162, bvc_qcd linear 1-auc: paired geometric mean [95% interval]
\newcommand{\PairedProbeBvcQcdLinearOnesevenMassOverOneseven}{1.89~[1.52, 2.34]}  % experiments/FIGS/data/paired_v1err/probes/ratios.json :: ratios[133].ratio, ratios[133].ci95 :: sha256 14dd51402305d147 :: 17+mass over 17, bvc_qcd linear 1-auc: paired geometric mean [95% interval]
\newcommand{\PairedProbeBvcQcdMlpOnesixtwoOverOneeighteight}{1.11~[1.04, 1.19]}  % experiments/FIGS/data/paired_v1err/probes/ratios.json :: ratios[158].ratio, ratios[158].ci95 :: sha256 14dd51402305d147 :: 162 over 188, bvc_qcd mlp 1-auc: paired geometric mean [95% interval]
\newcommand{\PairedProbeBvcQcdMlpFourthreeOverOneeighteight}{1.13~[1.05, 1.21]}  % experiments/FIGS/data/paired_v1err/probes/ratios.json :: ratios[159].ratio, ratios[159].ci95 :: sha256 14dd51402305d147 :: 43 over 188, bvc_qcd mlp 1-auc: paired geometric mean [95% interval]
\newcommand{\PairedProbeBvcQcdMlpOnesevenOverOneeighteight}{2.6~[2.2, 3.0]}  % experiments/FIGS/data/paired_v1err/probes/ratios.json :: ratios[160].ratio, ratios[160].ci95 :: sha256 14dd51402305d147 :: 17 over 188, bvc_qcd mlp 1-auc: paired geometric mean [95% interval]
\newcommand{\PairedProbeBvcQcdMlpFourthreeOverOnesixtwo}{1.01~[0.96, 1.06]}  % experiments/FIGS/data/paired_v1err/probes/ratios.json :: ratios[161].ratio, ratios[161].ci95 :: sha256 14dd51402305d147 :: 43 over 162, bvc_qcd mlp 1-auc: paired geometric mean [95% interval]
\newcommand{\PairedProbeBvcQcdMlpOnesevenOverOnesixtwo}{2.3~[2.0, 2.6]}  % experiments/FIGS/data/paired_v1err/probes/ratios.json :: ratios[162].ratio, ratios[162].ci95 :: sha256 14dd51402305d147 :: 17 over 162, bvc_qcd mlp 1-auc: paired geometric mean [95% interval]
\newcommand{\PairedProbeBvcQcdMlpOnesevenOverFourthree}{2.3~[2.0, 2.6]}  % experiments/FIGS/data/paired_v1err/probes/ratios.json :: ratios[163].ratio, ratios[163].ci95 :: sha256 14dd51402305d147 :: 17 over 43, bvc_qcd mlp 1-auc: paired geometric mean [95% interval]
\newcommand{\PairedProbeBvcQcdMlpOnesixtwoMassOverOnesixtwo}{1.06~[1.01, 1.12]}  % experiments/FIGS/data/paired_v1err/probes/ratios.json :: ratios[164].ratio, ratios[164].ci95 :: sha256 14dd51402305d147 :: 162+mass over 162, bvc_qcd mlp 1-auc: paired geometric mean [95% interval]
\newcommand{\PairedProbeBvcQcdMlpOnesevenMassOverOneseven}{1.64~[1.33, 2.02]}  % experiments/FIGS/data/paired_v1err/probes/ratios.json :: ratios[165].ratio, ratios[165].ci95 :: sha256 14dd51402305d147 :: 17+mass over 17, bvc_qcd mlp 1-auc: paired geometric mean [95% interval]
\newcommand{\PairedProbeBvcResonantLinearOnesixtwoOverOneeighteight}{1.02~[0.96, 1.08]}  % experiments/FIGS/data/paired_v1err/probes/ratios.json :: ratios[190].ratio, ratios[190].ci95 :: sha256 14dd51402305d147 :: 162 over 188, bvc_resonant linear 1-auc: paired geometric mean [95% interval]
\newcommand{\PairedProbeBvcResonantLinearFourthreeOverOneeighteight}{1.09~[1.03, 1.15]}  % experiments/FIGS/data/paired_v1err/probes/ratios.json :: ratios[191].ratio, ratios[191].ci95 :: sha256 14dd51402305d147 :: 43 over 188, bvc_resonant linear 1-auc: paired geometric mean [95% interval]
\newcommand{\PairedProbeBvcResonantLinearOnesevenOverOneeighteight}{4.0~[3.3, 4.9]}  % experiments/FIGS/data/paired_v1err/probes/ratios.json :: ratios[192].ratio, ratios[192].ci95 :: sha256 14dd51402305d147 :: 17 over 188, bvc_resonant linear 1-auc: paired geometric mean [95% interval]
\newcommand{\PairedProbeBvcResonantLinearFourthreeOverOnesixtwo}{1.07~[1.02, 1.12]}  % experiments/FIGS/data/paired_v1err/probes/ratios.json :: ratios[193].ratio, ratios[193].ci95 :: sha256 14dd51402305d147 :: 43 over 162, bvc_resonant linear 1-auc: paired geometric mean [95% interval]
\newcommand{\PairedProbeBvcResonantLinearOnesevenOverOnesixtwo}{3.9~[3.1, 4.9]}  % experiments/FIGS/data/paired_v1err/probes/ratios.json :: ratios[194].ratio, ratios[194].ci95 :: sha256 14dd51402305d147 :: 17 over 162, bvc_resonant linear 1-auc: paired geometric mean [95% interval]
\newcommand{\PairedProbeBvcResonantLinearOnesevenOverFourthree}{3.7~[3.0, 4.5]}  % experiments/FIGS/data/paired_v1err/probes/ratios.json :: ratios[195].ratio, ratios[195].ci95 :: sha256 14dd51402305d147 :: 17 over 43, bvc_resonant linear 1-auc: paired geometric mean [95% interval]
\newcommand{\PairedProbeBvcResonantLinearOnesixtwoMassOverOnesixtwo}{1.11~[1.06, 1.16]}  % experiments/FIGS/data/paired_v1err/probes/ratios.json :: ratios[196].ratio, ratios[196].ci95 :: sha256 14dd51402305d147 :: 162+mass over 162, bvc_resonant linear 1-auc: paired geometric mean [95% interval]
\newcommand{\PairedProbeBvcResonantLinearOnesevenMassOverOneseven}{3.0~[1.9, 4.7]}  % experiments/FIGS/data/paired_v1err/probes/ratios.json :: ratios[197].ratio, ratios[197].ci95 :: sha256 14dd51402305d147 :: 17+mass over 17, bvc_resonant linear 1-auc: paired geometric mean [95% interval]
\newcommand{\PairedProbeBvcResonantMlpOnesixtwoOverOneeighteight}{1.04~[0.98, 1.12]}  % experiments/FIGS/data/paired_v1err/probes/ratios.json :: ratios[222].ratio, ratios[222].ci95 :: sha256 14dd51402305d147 :: 162 over 188, bvc_resonant mlp 1-auc: paired geometric mean [95% interval]
\newcommand{\PairedProbeBvcResonantMlpFourthreeOverOneeighteight}{1.14~[1.08, 1.21]}  % experiments/FIGS/data/paired_v1err/probes/ratios.json :: ratios[223].ratio, ratios[223].ci95 :: sha256 14dd51402305d147 :: 43 over 188, bvc_resonant mlp 1-auc: paired geometric mean [95% interval]
\newcommand{\PairedProbeBvcResonantMlpOnesevenOverOneeighteight}{3.3~[2.8, 3.8]}  % experiments/FIGS/data/paired_v1err/probes/ratios.json :: ratios[224].ratio, ratios[224].ci95 :: sha256 14dd51402305d147 :: 17 over 188, bvc_resonant mlp 1-auc: paired geometric mean [95% interval]
\newcommand{\PairedProbeBvcResonantMlpFourthreeOverOnesixtwo}{1.09~[1.04, 1.14]}  % experiments/FIGS/data/paired_v1err/probes/ratios.json :: ratios[225].ratio, ratios[225].ci95 :: sha256 14dd51402305d147 :: 43 over 162, bvc_resonant mlp 1-auc: paired geometric mean [95% interval]
\newcommand{\PairedProbeBvcResonantMlpOnesevenOverOnesixtwo}{3.1~[2.6, 3.8]}  % experiments/FIGS/data/paired_v1err/probes/ratios.json :: ratios[226].ratio, ratios[226].ci95 :: sha256 14dd51402305d147 :: 17 over 162, bvc_resonant mlp 1-auc: paired geometric mean [95% interval]
\newcommand{\PairedProbeBvcResonantMlpOnesevenOverFourthree}{2.9~[2.4, 3.4]}  % experiments/FIGS/data/paired_v1err/probes/ratios.json :: ratios[227].ratio, ratios[227].ci95 :: sha256 14dd51402305d147 :: 17 over 43, bvc_resonant mlp 1-auc: paired geometric mean [95% interval]
\newcommand{\PairedProbeBvcResonantMlpOnesixtwoMassOverOnesixtwo}{1.09~[1.04, 1.14]}  % experiments/FIGS/data/paired_v1err/probes/ratios.json :: ratios[228].ratio, ratios[228].ci95 :: sha256 14dd51402305d147 :: 162+mass over 162, bvc_resonant mlp 1-auc: paired geometric mean [95% interval]
\newcommand{\PairedProbeBvcResonantMlpOnesevenMassOverOneseven}{1.95~[1.45, 2.61]}  % experiments/FIGS/data/paired_v1err/probes/ratios.json :: ratios[229].ratio, ratios[229].ci95 :: sha256 14dd51402305d147 :: 17+mass over 17, bvc_resonant mlp 1-auc: paired geometric mean [95% interval]
\newcommand{\PairedProbeEeVsMmLinearOnesevenOverFourthree}{2.3~[0.8, 6.5]}  % experiments/FIGS/data/paired_v1err/probes/ratios.json :: ratios[259].ratio, ratios[259].ci95 :: sha256 14dd51402305d147 :: 17 over 43, ee_vs_mm linear 1-auc: paired geometric mean [95% interval]
\newcommand{\PairedProbeEeVsMmLinearOnesevenMassOverOneseven}{0.23~[0.08, 0.64]}  % experiments/FIGS/data/paired_v1err/probes/ratios.json :: ratios[261].ratio, ratios[261].ci95 :: sha256 14dd51402305d147 :: 17+mass over 17, ee_vs_mm linear 1-auc: paired geometric mean [95% interval]
\newcommand{\PairedProbeEeVsMmMlpOnesevenOverFourthree}{2.8~[0.8, 9.0]}  % experiments/FIGS/data/paired_v1err/probes/ratios.json :: ratios[291].ratio, ratios[291].ci95 :: sha256 14dd51402305d147 :: 17 over 43, ee_vs_mm mlp 1-auc: paired geometric mean [95% interval]
\newcommand{\PairedProbeEeVsMmMlpOnesevenMassOverOneseven}{0.172~[0.026, 1.125]}  % experiments/FIGS/data/paired_v1err/probes/ratios.json :: ratios[293].ratio, ratios[293].ci95 :: sha256 14dd51402305d147 :: 17+mass over 17, ee_vs_mm mlp 1-auc: paired geometric mean [95% interval]
\newcommand{\PairedProbeRetainedTopologyLinearOnesixtwoOverOneeighteight}{1.00~[0.92, 1.08]}  % experiments/FIGS/data/paired_v1err/probes/ratios.json :: ratios[318].ratio, ratios[318].ci95 :: sha256 14dd51402305d147 :: 162 over 188, retained_topology linear 1-auc: paired geometric mean [95% interval]
\newcommand{\PairedProbeRetainedTopologyLinearFourthreeOverOneeighteight}{1.24~[1.15, 1.34]}  % experiments/FIGS/data/paired_v1err/probes/ratios.json :: ratios[319].ratio, ratios[319].ci95 :: sha256 14dd51402305d147 :: 43 over 188, retained_topology linear 1-auc: paired geometric mean [95% interval]
\newcommand{\PairedProbeRetainedTopologyLinearOnesevenOverOneeighteight}{1.54~[1.35, 1.75]}  % experiments/FIGS/data/paired_v1err/probes/ratios.json :: ratios[320].ratio, ratios[320].ci95 :: sha256 14dd51402305d147 :: 17 over 188, retained_topology linear 1-auc: paired geometric mean [95% interval]
\newcommand{\PairedProbeRetainedTopologyLinearFourthreeOverOnesixtwo}{1.24~[1.14, 1.34]}  % experiments/FIGS/data/paired_v1err/probes/ratios.json :: ratios[321].ratio, ratios[321].ci95 :: sha256 14dd51402305d147 :: 43 over 162, retained_topology linear 1-auc: paired geometric mean [95% interval]
\newcommand{\PairedProbeRetainedTopologyLinearOnesevenOverOnesixtwo}{1.54~[1.33, 1.78]}  % experiments/FIGS/data/paired_v1err/probes/ratios.json :: ratios[322].ratio, ratios[322].ci95 :: sha256 14dd51402305d147 :: 17 over 162, retained_topology linear 1-auc: paired geometric mean [95% interval]
\newcommand{\PairedProbeRetainedTopologyLinearOnesevenOverFourthree}{1.24~[1.08, 1.42]}  % experiments/FIGS/data/paired_v1err/probes/ratios.json :: ratios[323].ratio, ratios[323].ci95 :: sha256 14dd51402305d147 :: 17 over 43, retained_topology linear 1-auc: paired geometric mean [95% interval]
\newcommand{\PairedProbeRetainedTopologyLinearOnesixtwoMassOverOnesixtwo}{1.06~[0.99, 1.14]}  % experiments/FIGS/data/paired_v1err/probes/ratios.json :: ratios[324].ratio, ratios[324].ci95 :: sha256 14dd51402305d147 :: 162+mass over 162, retained_topology linear 1-auc: paired geometric mean [95% interval]
\newcommand{\PairedProbeRetainedTopologyLinearOnesevenMassOverOneseven}{1.12~[0.97, 1.28]}  % experiments/FIGS/data/paired_v1err/probes/ratios.json :: ratios[325].ratio, ratios[325].ci95 :: sha256 14dd51402305d147 :: 17+mass over 17, retained_topology linear 1-auc: paired geometric mean [95% interval]
\newcommand{\PairedProbeRetainedTopologyMlpOnesixtwoOverOneeighteight}{1.01~[0.87, 1.17]}  % experiments/FIGS/data/paired_v1err/probes/ratios.json :: ratios[350].ratio, ratios[350].ci95 :: sha256 14dd51402305d147 :: 162 over 188, retained_topology mlp 1-auc: paired geometric mean [95% interval]
\newcommand{\PairedProbeRetainedTopologyMlpFourthreeOverOneeighteight}{1.18~[1.03, 1.36]}  % experiments/FIGS/data/paired_v1err/probes/ratios.json :: ratios[351].ratio, ratios[351].ci95 :: sha256 14dd51402305d147 :: 43 over 188, retained_topology mlp 1-auc: paired geometric mean [95% interval]
\newcommand{\PairedProbeRetainedTopologyMlpOnesevenOverOneeighteight}{1.52~[1.34, 1.72]}  % experiments/FIGS/data/paired_v1err/probes/ratios.json :: ratios[352].ratio, ratios[352].ci95 :: sha256 14dd51402305d147 :: 17 over 188, retained_topology mlp 1-auc: paired geometric mean [95% interval]
\newcommand{\PairedProbeRetainedTopologyMlpFourthreeOverOnesixtwo}{1.17~[1.07, 1.28]}  % experiments/FIGS/data/paired_v1err/probes/ratios.json :: ratios[353].ratio, ratios[353].ci95 :: sha256 14dd51402305d147 :: 43 over 162, retained_topology mlp 1-auc: paired geometric mean [95% interval]
\newcommand{\PairedProbeRetainedTopologyMlpOnesevenOverOnesixtwo}{1.50~[1.29, 1.75]}  % experiments/FIGS/data/paired_v1err/probes/ratios.json :: ratios[354].ratio, ratios[354].ci95 :: sha256 14dd51402305d147 :: 17 over 162, retained_topology mlp 1-auc: paired geometric mean [95% interval]
\newcommand{\PairedProbeRetainedTopologyMlpOnesevenOverFourthree}{1.28~[1.12, 1.47]}  % experiments/FIGS/data/paired_v1err/probes/ratios.json :: ratios[355].ratio, ratios[355].ci95 :: sha256 14dd51402305d147 :: 17 over 43, retained_topology mlp 1-auc: paired geometric mean [95% interval]
\newcommand{\PairedProbeRetainedTopologyMlpOnesixtwoMassOverOnesixtwo}{1.04~[0.92, 1.17]}  % experiments/FIGS/data/paired_v1err/probes/ratios.json :: ratios[356].ratio, ratios[356].ci95 :: sha256 14dd51402305d147 :: 162+mass over 162, retained_topology mlp 1-auc: paired geometric mean [95% interval]
\newcommand{\PairedProbeRetainedTopologyMlpOnesevenMassOverOneseven}{1.11~[0.99, 1.23]}  % experiments/FIGS/data/paired_v1err/probes/ratios.json :: ratios[357].ratio, ratios[357].ci95 :: sha256 14dd51402305d147 :: 17+mass over 17, retained_topology mlp 1-auc: paired geometric mean [95% interval]
\newcommand{\PairedProbeVisibleContentLinearOnesixtwoOverOneeighteight}{0.95~[0.88, 1.03]}  % experiments/FIGS/data/paired_v1err/probes/ratios.json :: ratios[382].ratio, ratios[382].ci95 :: sha256 14dd51402305d147 :: 162 over 188, visible_content linear 1-auc: paired geometric mean [95% interval]
\newcommand{\PairedProbeVisibleContentLinearFourthreeOverOneeighteight}{0.99~[0.85, 1.15]}  % experiments/FIGS/data/paired_v1err/probes/ratios.json :: ratios[383].ratio, ratios[383].ci95 :: sha256 14dd51402305d147 :: 43 over 188, visible_content linear 1-auc: paired geometric mean [95% interval]
\newcommand{\PairedProbeVisibleContentLinearOnesevenOverOneeighteight}{2.8~[2.1, 3.7]}  % experiments/FIGS/data/paired_v1err/probes/ratios.json :: ratios[384].ratio, ratios[384].ci95 :: sha256 14dd51402305d147 :: 17 over 188, visible_content linear 1-auc: paired geometric mean [95% interval]
\newcommand{\PairedProbeVisibleContentLinearFourthreeOverOnesixtwo}{1.04~[0.89, 1.20]}  % experiments/FIGS/data/paired_v1err/probes/ratios.json :: ratios[385].ratio, ratios[385].ci95 :: sha256 14dd51402305d147 :: 43 over 162, visible_content linear 1-auc: paired geometric mean [95% interval]
\newcommand{\PairedProbeVisibleContentLinearOnesevenOverOnesixtwo}{2.9~[2.3, 3.8]}  % experiments/FIGS/data/paired_v1err/probes/ratios.json :: ratios[386].ratio, ratios[386].ci95 :: sha256 14dd51402305d147 :: 17 over 162, visible_content linear 1-auc: paired geometric mean [95% interval]
\newcommand{\PairedProbeVisibleContentLinearOnesevenOverFourthree}{2.8~[2.2, 3.7]}  % experiments/FIGS/data/paired_v1err/probes/ratios.json :: ratios[387].ratio, ratios[387].ci95 :: sha256 14dd51402305d147 :: 17 over 43, visible_content linear 1-auc: paired geometric mean [95% interval]
\newcommand{\PairedProbeVisibleContentLinearOnesixtwoMassOverOnesixtwo}{1.10~[0.98, 1.23]}  % experiments/FIGS/data/paired_v1err/probes/ratios.json :: ratios[388].ratio, ratios[388].ci95 :: sha256 14dd51402305d147 :: 162+mass over 162, visible_content linear 1-auc: paired geometric mean [95% interval]
\newcommand{\PairedProbeVisibleContentLinearOnesevenMassOverOneseven}{1.44~[1.26, 1.65]}  % experiments/FIGS/data/paired_v1err/probes/ratios.json :: ratios[389].ratio, ratios[389].ci95 :: sha256 14dd51402305d147 :: 17+mass over 17, visible_content linear 1-auc: paired geometric mean [95% interval]
\newcommand{\PairedProbeVisibleContentLinearRandomOverOneeighteight}{1.31~[1.01, 1.70]}  % experiments/FIGS/data/paired_v1err/probes/ratios.json :: ratios[390].ratio, ratios[390].ci95 :: sha256 14dd51402305d147 :: random over 188, visible_content linear 1-auc: paired geometric mean [95% interval]
\newcommand{\PairedProbeVisibleContentLinearRandomOverOnesixtwo}{1.41~[1.17, 1.70]}  % experiments/FIGS/data/paired_v1err/probes/ratios.json :: ratios[391].ratio, ratios[391].ci95 :: sha256 14dd51402305d147 :: random over 162, visible_content linear 1-auc: paired geometric mean [95% interval]
\newcommand{\PairedProbeVisibleContentLinearRandomOverFourthree}{1.33~[1.12, 1.57]}  % experiments/FIGS/data/paired_v1err/probes/ratios.json :: ratios[392].ratio, ratios[392].ci95 :: sha256 14dd51402305d147 :: random over 43, visible_content linear 1-auc: paired geometric mean [95% interval]
\newcommand{\PairedProbeVisibleContentLinearRandomOverOneseven}{0.48~[0.38, 0.60]}  % experiments/FIGS/data/paired_v1err/probes/ratios.json :: ratios[393].ratio, ratios[393].ci95 :: sha256 14dd51402305d147 :: random over 17, visible_content linear 1-auc: paired geometric mean [95% interval]
\newcommand{\PairedProbeVisibleContentLinearRandomDrawTwoOverRandomDrawOne}{0.79~[0.64, 0.98]}  % experiments/FIGS/data/paired_v1err/probes/ratios.json :: ratios[394].ratio, ratios[394].ci95 :: sha256 14dd51402305d147 :: random draw 2 over random draw 1, visible_content linear 1-auc: paired geometric mean [95% interval]
\newcommand{\PairedProbeVisibleContentLinearRandomDrawThreeOverRandomDrawOne}{0.86~[0.71, 1.05]}  % experiments/FIGS/data/paired_v1err/probes/ratios.json :: ratios[395].ratio, ratios[395].ci95 :: sha256 14dd51402305d147 :: random draw 3 over random draw 1, visible_content linear 1-auc: paired geometric mean [95% interval]
\newcommand{\PairedProbeVisibleContentLinearRandomDrawThreeOverRandomDrawTwo}{1.09~[0.88, 1.35]}  % experiments/FIGS/data/paired_v1err/probes/ratios.json :: ratios[396].ratio, ratios[396].ci95 :: sha256 14dd51402305d147 :: random draw 3 over random draw 2, visible_content linear 1-auc: paired geometric mean [95% interval]
\newcommand{\PairedProbeVisibleContentMlpOnesixtwoOverOneeighteight}{0.92~[0.85, 1.01]}  % experiments/FIGS/data/paired_v1err/probes/ratios.json :: ratios[442].ratio, ratios[442].ci95 :: sha256 14dd51402305d147 :: 162 over 188, visible_content mlp 1-auc: paired geometric mean [95% interval]
\newcommand{\PairedProbeVisibleContentMlpFourthreeOverOneeighteight}{0.93~[0.80, 1.07]}  % experiments/FIGS/data/paired_v1err/probes/ratios.json :: ratios[443].ratio, ratios[443].ci95 :: sha256 14dd51402305d147 :: 43 over 188, visible_content mlp 1-auc: paired geometric mean [95% interval]
\newcommand{\PairedProbeVisibleContentMlpOnesevenOverOneeighteight}{2.3~[1.7, 3.1]}  % experiments/FIGS/data/paired_v1err/probes/ratios.json :: ratios[444].ratio, ratios[444].ci95 :: sha256 14dd51402305d147 :: 17 over 188, visible_content mlp 1-auc: paired geometric mean [95% interval]
\newcommand{\PairedProbeVisibleContentMlpFourthreeOverOnesixtwo}{1.01~[0.89, 1.14]}  % experiments/FIGS/data/paired_v1err/probes/ratios.json :: ratios[445].ratio, ratios[445].ci95 :: sha256 14dd51402305d147 :: 43 over 162, visible_content mlp 1-auc: paired geometric mean [95% interval]
\newcommand{\PairedProbeVisibleContentMlpOnesevenOverOnesixtwo}{2.5~[1.9, 3.2]}  % experiments/FIGS/data/paired_v1err/probes/ratios.json :: ratios[446].ratio, ratios[446].ci95 :: sha256 14dd51402305d147 :: 17 over 162, visible_content mlp 1-auc: paired geometric mean [95% interval]
\newcommand{\PairedProbeVisibleContentMlpOnesevenOverFourthree}{2.5~[2.0, 3.1]}  % experiments/FIGS/data/paired_v1err/probes/ratios.json :: ratios[447].ratio, ratios[447].ci95 :: sha256 14dd51402305d147 :: 17 over 43, visible_content mlp 1-auc: paired geometric mean [95% interval]
\newcommand{\PairedProbeVisibleContentMlpOnesixtwoMassOverOnesixtwo}{1.08~[0.94, 1.24]}  % experiments/FIGS/data/paired_v1err/probes/ratios.json :: ratios[448].ratio, ratios[448].ci95 :: sha256 14dd51402305d147 :: 162+mass over 162, visible_content mlp 1-auc: paired geometric mean [95% interval]
\newcommand{\PairedProbeVisibleContentMlpOnesevenMassOverOneseven}{1.41~[1.19, 1.67]}  % experiments/FIGS/data/paired_v1err/probes/ratios.json :: ratios[449].ratio, ratios[449].ci95 :: sha256 14dd51402305d147 :: 17+mass over 17, visible_content mlp 1-auc: paired geometric mean [95% interval]
\newcommand{\PairedProbeVisibleContentMlpRandomOverOneeighteight}{1.27~[1.01, 1.59]}  % experiments/FIGS/data/paired_v1err/probes/ratios.json :: ratios[450].ratio, ratios[450].ci95 :: sha256 14dd51402305d147 :: random over 188, visible_content mlp 1-auc: paired geometric mean [95% interval]
\newcommand{\PairedProbeVisibleContentMlpRandomOverOnesixtwo}{1.37~[1.12, 1.68]}  % experiments/FIGS/data/paired_v1err/probes/ratios.json :: ratios[451].ratio, ratios[451].ci95 :: sha256 14dd51402305d147 :: random over 162, visible_content mlp 1-auc: paired geometric mean [95% interval]
\newcommand{\PairedProbeVisibleContentMlpRandomOverFourthree}{1.35~[1.13, 1.61]}  % experiments/FIGS/data/paired_v1err/probes/ratios.json :: ratios[452].ratio, ratios[452].ci95 :: sha256 14dd51402305d147 :: random over 43, visible_content mlp 1-auc: paired geometric mean [95% interval]
\newcommand{\PairedProbeVisibleContentMlpRandomOverOneseven}{0.56~[0.45, 0.70]}  % experiments/FIGS/data/paired_v1err/probes/ratios.json :: ratios[453].ratio, ratios[453].ci95 :: sha256 14dd51402305d147 :: random over 17, visible_content mlp 1-auc: paired geometric mean [95% interval]
\newcommand{\PairedProbeVisibleContentMlpRandomDrawTwoOverRandomDrawOne}{0.88~[0.71, 1.09]}  % experiments/FIGS/data/paired_v1err/probes/ratios.json :: ratios[454].ratio, ratios[454].ci95 :: sha256 14dd51402305d147 :: random draw 2 over random draw 1, visible_content mlp 1-auc: paired geometric mean [95% interval]
\newcommand{\PairedProbeVisibleContentMlpRandomDrawThreeOverRandomDrawOne}{0.94~[0.75, 1.18]}  % experiments/FIGS/data/paired_v1err/probes/ratios.json :: ratios[455].ratio, ratios[455].ci95 :: sha256 14dd51402305d147 :: random draw 3 over random draw 1, visible_content mlp 1-auc: paired geometric mean [95% interval]
\newcommand{\PairedProbeVisibleContentMlpRandomDrawThreeOverRandomDrawTwo}{1.06~[0.84, 1.35]}  % experiments/FIGS/data/paired_v1err/probes/ratios.json :: ratios[456].ratio, ratios[456].ci95 :: sha256 14dd51402305d147 :: random draw 3 over random draw 2, visible_content mlp 1-auc: paired geometric mean [95% interval]
\newcommand{\PairedFtJciiEThreeOnesixtwoOverOneeighteight}{0.95~[0.87, 1.05]}  % experiments/FIGS/data/paired_v1err/finetune/ratios.json :: ratios[0].ratio, ratios[0].ci95 :: sha256 3e5cf25e859b68c2 :: 162 over 188, leg1 N1000 1-macro_auc: paired geometric mean [95% interval]
\newcommand{\PairedFtJciiEThreeFourthreeOverOneeighteight}{0.84~[0.77, 0.91]}  % experiments/FIGS/data/paired_v1err/finetune/ratios.json :: ratios[1].ratio, ratios[1].ci95 :: sha256 3e5cf25e859b68c2 :: 43 over 188, leg1 N1000 1-macro_auc: paired geometric mean [95% interval]
\newcommand{\PairedFtJciiEThreeOnesevenOverOneeighteight}{0.97~[0.86, 1.09]}  % experiments/FIGS/data/paired_v1err/finetune/ratios.json :: ratios[2].ratio, ratios[2].ci95 :: sha256 3e5cf25e859b68c2 :: 17 over 188, leg1 N1000 1-macro_auc: paired geometric mean [95% interval]
\newcommand{\PairedFtJciiEThreeFourthreeOverOnesixtwo}{0.88~[0.80, 0.97]}  % experiments/FIGS/data/paired_v1err/finetune/ratios.json :: ratios[3].ratio, ratios[3].ci95 :: sha256 3e5cf25e859b68c2 :: 43 over 162, leg1 N1000 1-macro_auc: paired geometric mean [95% interval]
\newcommand{\PairedFtJciiEThreeOnesevenOverOnesixtwo}{1.02~[0.88, 1.18]}  % experiments/FIGS/data/paired_v1err/finetune/ratios.json :: ratios[4].ratio, ratios[4].ci95 :: sha256 3e5cf25e859b68c2 :: 17 over 162, leg1 N1000 1-macro_auc: paired geometric mean [95% interval]
\newcommand{\PairedFtJciiEThreeOnesevenOverFourthree}{1.16~[1.08, 1.24]}  % experiments/FIGS/data/paired_v1err/finetune/ratios.json :: ratios[5].ratio, ratios[5].ci95 :: sha256 3e5cf25e859b68c2 :: 17 over 43, leg1 N1000 1-macro_auc: paired geometric mean [95% interval]
\newcommand{\PairedFtJciiEThreeOnesixtwoMassOverOnesixtwo}{0.85~[0.77, 0.93]}  % experiments/FIGS/data/paired_v1err/finetune/ratios.json :: ratios[6].ratio, ratios[6].ci95 :: sha256 3e5cf25e859b68c2 :: 162+mass over 162, leg1 N1000 1-macro_auc: paired geometric mean [95% interval]
\newcommand{\PairedFtJciiEThreeOnesevenMassOverOneseven}{1.00~[0.86, 1.17]}  % experiments/FIGS/data/paired_v1err/finetune/ratios.json :: ratios[7].ratio, ratios[7].ci95 :: sha256 3e5cf25e859b68c2 :: 17+mass over 17, leg1 N1000 1-macro_auc: paired geometric mean [95% interval]
\newcommand{\PairedFtJciiEThreeRandomOverOneeighteight}{0.82~[0.73, 0.93]}  % experiments/FIGS/data/paired_v1err/finetune/ratios.json :: ratios[8].ratio, ratios[8].ci95 :: sha256 3e5cf25e859b68c2 :: random over 188, leg1 N1000 1-macro_auc: paired geometric mean [95% interval]
\newcommand{\PairedFtJciiEThreeRandomOverOnesixtwo}{0.90~[0.85, 0.96]}  % experiments/FIGS/data/paired_v1err/finetune/ratios.json :: ratios[9].ratio, ratios[9].ci95 :: sha256 3e5cf25e859b68c2 :: random over 162, leg1 N1000 1-macro_auc: paired geometric mean [95% interval]
\newcommand{\PairedFtJciiEThreeRandomOverFourthree}{0.99~[0.75, 1.30]}  % experiments/FIGS/data/paired_v1err/finetune/ratios.json :: ratios[10].ratio, ratios[10].ci95 :: sha256 3e5cf25e859b68c2 :: random over 43, leg1 N1000 1-macro_auc: paired geometric mean [95% interval]
\newcommand{\PairedFtJciiEThreeRandomOverOneseven}{0.86~[0.58, 1.27]}  % experiments/FIGS/data/paired_v1err/finetune/ratios.json :: ratios[11].ratio, ratios[11].ci95 :: sha256 3e5cf25e859b68c2 :: random over 17, leg1 N1000 1-macro_auc: paired geometric mean [95% interval]
\newcommand{\PairedFtJciiEThreeRandomDrawTwoOverRandomDrawOne}{1.09~[1.07, 1.10]}  % experiments/FIGS/data/paired_v1err/finetune/ratios.json :: ratios[12].ratio, ratios[12].ci95 :: sha256 3e5cf25e859b68c2 :: random draw 2 over random draw 1, leg1 N1000 1-macro_auc: paired geometric mean [95% interval]
\newcommand{\PairedFtJciiEThreeRandomDrawThreeOverRandomDrawOne}{1.04~[1.02, 1.05]}  % experiments/FIGS/data/paired_v1err/finetune/ratios.json :: ratios[13].ratio, ratios[13].ci95 :: sha256 3e5cf25e859b68c2 :: random draw 3 over random draw 1, leg1 N1000 1-macro_auc: paired geometric mean [95% interval]
\newcommand{\PairedFtJciiEThreeRandomDrawThreeOverRandomDrawTwo}{0.95~[0.94, 0.97]}  % experiments/FIGS/data/paired_v1err/finetune/ratios.json :: ratios[14].ratio, ratios[14].ci95 :: sha256 3e5cf25e859b68c2 :: random draw 3 over random draw 2, leg1 N1000 1-macro_auc: paired geometric mean [95% interval]
\newcommand{\PairedFtJciiEFourOnesixtwoOverOneeighteight}{0.99~[0.97, 1.02]}  % experiments/FIGS/data/paired_v1err/finetune/ratios.json :: ratios[15].ratio, ratios[15].ci95 :: sha256 3e5cf25e859b68c2 :: 162 over 188, leg1 N10000 1-macro_auc: paired geometric mean [95% interval]
\newcommand{\PairedFtJciiEFourFourthreeOverOneeighteight}{1.20~[1.17, 1.23]}  % experiments/FIGS/data/paired_v1err/finetune/ratios.json :: ratios[16].ratio, ratios[16].ci95 :: sha256 3e5cf25e859b68c2 :: 43 over 188, leg1 N10000 1-macro_auc: paired geometric mean [95% interval]
\newcommand{\PairedFtJciiEFourOnesevenOverOneeighteight}{1.32~[1.29, 1.34]}  % experiments/FIGS/data/paired_v1err/finetune/ratios.json :: ratios[17].ratio, ratios[17].ci95 :: sha256 3e5cf25e859b68c2 :: 17 over 188, leg1 N10000 1-macro_auc: paired geometric mean [95% interval]
\newcommand{\PairedFtJciiEFourFourthreeOverOnesixtwo}{1.21~[1.19, 1.23]}  % experiments/FIGS/data/paired_v1err/finetune/ratios.json :: ratios[18].ratio, ratios[18].ci95 :: sha256 3e5cf25e859b68c2 :: 43 over 162, leg1 N10000 1-macro_auc: paired geometric mean [95% interval]
\newcommand{\PairedFtJciiEFourOnesevenOverOnesixtwo}{1.32~[1.29, 1.36]}  % experiments/FIGS/data/paired_v1err/finetune/ratios.json :: ratios[19].ratio, ratios[19].ci95 :: sha256 3e5cf25e859b68c2 :: 17 over 162, leg1 N10000 1-macro_auc: paired geometric mean [95% interval]
\newcommand{\PairedFtJciiEFourOnesevenOverFourthree}{1.10~[1.07, 1.13]}  % experiments/FIGS/data/paired_v1err/finetune/ratios.json :: ratios[20].ratio, ratios[20].ci95 :: sha256 3e5cf25e859b68c2 :: 17 over 43, leg1 N10000 1-macro_auc: paired geometric mean [95% interval]
\newcommand{\PairedFtJciiEFourOnesixtwoMassOverOnesixtwo}{1.00~[0.94, 1.06]}  % experiments/FIGS/data/paired_v1err/finetune/ratios.json :: ratios[21].ratio, ratios[21].ci95 :: sha256 3e5cf25e859b68c2 :: 162+mass over 162, leg1 N10000 1-macro_auc: paired geometric mean [95% interval]
\newcommand{\PairedFtJciiEFourOnesevenMassOverOneseven}{0.990~[0.976, 1.004]}  % experiments/FIGS/data/paired_v1err/finetune/ratios.json :: ratios[22].ratio, ratios[22].ci95 :: sha256 3e5cf25e859b68c2 :: 17+mass over 17, leg1 N10000 1-macro_auc: paired geometric mean [95% interval]
\newcommand{\PairedFtJciiEFourRandomOverOneeighteight}{1.15~[1.09, 1.22]}  % experiments/FIGS/data/paired_v1err/finetune/ratios.json :: ratios[23].ratio, ratios[23].ci95 :: sha256 3e5cf25e859b68c2 :: random over 188, leg1 N10000 1-macro_auc: paired geometric mean [95% interval]
\newcommand{\PairedFtJciiEFourRandomOverOnesixtwo}{1.18~[1.13, 1.22]}  % experiments/FIGS/data/paired_v1err/finetune/ratios.json :: ratios[24].ratio, ratios[24].ci95 :: sha256 3e5cf25e859b68c2 :: random over 162, leg1 N10000 1-macro_auc: paired geometric mean [95% interval]
\newcommand{\PairedFtJciiEFourRandomOverFourthree}{0.97~[0.96, 0.99]}  % experiments/FIGS/data/paired_v1err/finetune/ratios.json :: ratios[25].ratio, ratios[25].ci95 :: sha256 3e5cf25e859b68c2 :: random over 43, leg1 N10000 1-macro_auc: paired geometric mean [95% interval]
\newcommand{\PairedFtJciiEFourRandomOverOneseven}{0.88~[0.82, 0.95]}  % experiments/FIGS/data/paired_v1err/finetune/ratios.json :: ratios[26].ratio, ratios[26].ci95 :: sha256 3e5cf25e859b68c2 :: random over 17, leg1 N10000 1-macro_auc: paired geometric mean [95% interval]
\newcommand{\PairedFtJciiEFourRandomDrawTwoOverRandomDrawOne}{1.00~[0.98, 1.02]}  % experiments/FIGS/data/paired_v1err/finetune/ratios.json :: ratios[27].ratio, ratios[27].ci95 :: sha256 3e5cf25e859b68c2 :: random draw 2 over random draw 1, leg1 N10000 1-macro_auc: paired geometric mean [95% interval]
\newcommand{\PairedFtJciiEFourRandomDrawThreeOverRandomDrawOne}{0.97~[0.95, 0.98]}  % experiments/FIGS/data/paired_v1err/finetune/ratios.json :: ratios[28].ratio, ratios[28].ci95 :: sha256 3e5cf25e859b68c2 :: random draw 3 over random draw 1, leg1 N10000 1-macro_auc: paired geometric mean [95% interval]
\newcommand{\PairedFtJciiEFourRandomDrawThreeOverRandomDrawTwo}{0.97~[0.95, 0.98]}  % experiments/FIGS/data/paired_v1err/finetune/ratios.json :: ratios[29].ratio, ratios[29].ci95 :: sha256 3e5cf25e859b68c2 :: random draw 3 over random draw 2, leg1 N10000 1-macro_auc: paired geometric mean [95% interval]
\newcommand{\PairedFtJciiEFiveOnesixtwoOverOneeighteight}{1.00~[0.98, 1.01]}  % experiments/FIGS/data/paired_v1err/finetune/ratios.json :: ratios[30].ratio, ratios[30].ci95 :: sha256 3e5cf25e859b68c2 :: 162 over 188, leg1 N100000 1-macro_auc: paired geometric mean [95% interval]
\newcommand{\PairedFtJciiEFiveFourthreeOverOneeighteight}{1.09~[1.06, 1.11]}  % experiments/FIGS/data/paired_v1err/finetune/ratios.json :: ratios[31].ratio, ratios[31].ci95 :: sha256 3e5cf25e859b68c2 :: 43 over 188, leg1 N100000 1-macro_auc: paired geometric mean [95% interval]
\newcommand{\PairedFtJciiEFiveOnesevenOverOneeighteight}{1.22~[1.20, 1.23]}  % experiments/FIGS/data/paired_v1err/finetune/ratios.json :: ratios[32].ratio, ratios[32].ci95 :: sha256 3e5cf25e859b68c2 :: 17 over 188, leg1 N100000 1-macro_auc: paired geometric mean [95% interval]
\newcommand{\PairedFtJciiEFiveFourthreeOverOnesixtwo}{1.09~[1.07, 1.11]}  % experiments/FIGS/data/paired_v1err/finetune/ratios.json :: ratios[33].ratio, ratios[33].ci95 :: sha256 3e5cf25e859b68c2 :: 43 over 162, leg1 N100000 1-macro_auc: paired geometric mean [95% interval]
\newcommand{\PairedFtJciiEFiveOnesevenOverOnesixtwo}{1.22~[1.20, 1.24]}  % experiments/FIGS/data/paired_v1err/finetune/ratios.json :: ratios[34].ratio, ratios[34].ci95 :: sha256 3e5cf25e859b68c2 :: 17 over 162, leg1 N100000 1-macro_auc: paired geometric mean [95% interval]
\newcommand{\PairedFtJciiEFiveOnesevenOverFourthree}{1.12~[1.10, 1.14]}  % experiments/FIGS/data/paired_v1err/finetune/ratios.json :: ratios[35].ratio, ratios[35].ci95 :: sha256 3e5cf25e859b68c2 :: 17 over 43, leg1 N100000 1-macro_auc: paired geometric mean [95% interval]
\newcommand{\PairedFtJciiEFiveOnesixtwoMassOverOnesixtwo}{1.018~[1.002, 1.035]}  % experiments/FIGS/data/paired_v1err/finetune/ratios.json :: ratios[36].ratio, ratios[36].ci95 :: sha256 3e5cf25e859b68c2 :: 162+mass over 162, leg1 N100000 1-macro_auc: paired geometric mean [95% interval]
\newcommand{\PairedFtJciiEFiveOnesevenMassOverOneseven}{1.03~[1.01, 1.06]}  % experiments/FIGS/data/paired_v1err/finetune/ratios.json :: ratios[37].ratio, ratios[37].ci95 :: sha256 3e5cf25e859b68c2 :: 17+mass over 17, leg1 N100000 1-macro_auc: paired geometric mean [95% interval]
\newcommand{\PairedFtJciiEFiveRandomOverOneeighteight}{1.15~[1.09, 1.20]}  % experiments/FIGS/data/paired_v1err/finetune/ratios.json :: ratios[38].ratio, ratios[38].ci95 :: sha256 3e5cf25e859b68c2 :: random over 188, leg1 N100000 1-macro_auc: paired geometric mean [95% interval]
\newcommand{\PairedFtJciiEFiveRandomOverOnesixtwo}{1.15~[1.09, 1.20]}  % experiments/FIGS/data/paired_v1err/finetune/ratios.json :: ratios[39].ratio, ratios[39].ci95 :: sha256 3e5cf25e859b68c2 :: random over 162, leg1 N100000 1-macro_auc: paired geometric mean [95% interval]
\newcommand{\PairedFtJciiEFiveRandomOverFourthree}{1.05~[1.04, 1.06]}  % experiments/FIGS/data/paired_v1err/finetune/ratios.json :: ratios[40].ratio, ratios[40].ci95 :: sha256 3e5cf25e859b68c2 :: random over 43, leg1 N100000 1-macro_auc: paired geometric mean [95% interval]
\newcommand{\PairedFtJciiEFiveRandomOverOneseven}{0.95~[0.93, 0.96]}  % experiments/FIGS/data/paired_v1err/finetune/ratios.json :: ratios[41].ratio, ratios[41].ci95 :: sha256 3e5cf25e859b68c2 :: random over 17, leg1 N100000 1-macro_auc: paired geometric mean [95% interval]
\newcommand{\PairedFtJciiEFiveRandomDrawTwoOverRandomDrawOne}{0.98~[0.97, 0.99]}  % experiments/FIGS/data/paired_v1err/finetune/ratios.json :: ratios[42].ratio, ratios[42].ci95 :: sha256 3e5cf25e859b68c2 :: random draw 2 over random draw 1, leg1 N100000 1-macro_auc: paired geometric mean [95% interval]
\newcommand{\PairedFtJciiEFiveRandomDrawThreeOverRandomDrawOne}{1.00~[0.99, 1.02]}  % experiments/FIGS/data/paired_v1err/finetune/ratios.json :: ratios[43].ratio, ratios[43].ci95 :: sha256 3e5cf25e859b68c2 :: random draw 3 over random draw 1, leg1 N100000 1-macro_auc: paired geometric mean [95% interval]
\newcommand{\PairedFtJciiEFiveRandomDrawThreeOverRandomDrawTwo}{1.03~[1.01, 1.04]}  % experiments/FIGS/data/paired_v1err/finetune/ratios.json :: ratios[44].ratio, ratios[44].ci95 :: sha256 3e5cf25e859b68c2 :: random draw 3 over random draw 2, leg1 N100000 1-macro_auc: paired geometric mean [95% interval]
\newcommand{\PairedFtJciiESixOnesixtwoOverOneeighteight}{1.002~[0.998, 1.005]}  % experiments/FIGS/data/paired_v1err/finetune/ratios.json :: ratios[45].ratio, ratios[45].ci95 :: sha256 3e5cf25e859b68c2 :: 162 over 188, leg1 N1000000 1-macro_auc: paired geometric mean [95% interval]
\newcommand{\PairedFtJciiESixFourthreeOverOneeighteight}{1.045~[1.037, 1.053]}  % experiments/FIGS/data/paired_v1err/finetune/ratios.json :: ratios[46].ratio, ratios[46].ci95 :: sha256 3e5cf25e859b68c2 :: 43 over 188, leg1 N1000000 1-macro_auc: paired geometric mean [95% interval]
\newcommand{\PairedFtJciiESixOnesevenOverOneeighteight}{1.14~[1.13, 1.16]}  % experiments/FIGS/data/paired_v1err/finetune/ratios.json :: ratios[47].ratio, ratios[47].ci95 :: sha256 3e5cf25e859b68c2 :: 17 over 188, leg1 N1000000 1-macro_auc: paired geometric mean [95% interval]
\newcommand{\PairedFtJciiESixFourthreeOverOnesixtwo}{1.044~[1.036, 1.051]}  % experiments/FIGS/data/paired_v1err/finetune/ratios.json :: ratios[48].ratio, ratios[48].ci95 :: sha256 3e5cf25e859b68c2 :: 43 over 162, leg1 N1000000 1-macro_auc: paired geometric mean [95% interval]
\newcommand{\PairedFtJciiESixOnesevenOverOnesixtwo}{1.14~[1.13, 1.15]}  % experiments/FIGS/data/paired_v1err/finetune/ratios.json :: ratios[49].ratio, ratios[49].ci95 :: sha256 3e5cf25e859b68c2 :: 17 over 162, leg1 N1000000 1-macro_auc: paired geometric mean [95% interval]
\newcommand{\PairedFtJciiESixOnesevenOverFourthree}{1.09~[1.08, 1.11]}  % experiments/FIGS/data/paired_v1err/finetune/ratios.json :: ratios[50].ratio, ratios[50].ci95 :: sha256 3e5cf25e859b68c2 :: 17 over 43, leg1 N1000000 1-macro_auc: paired geometric mean [95% interval]
\newcommand{\PairedFtJciiESixOnesixtwoMassOverOnesixtwo}{1.019~[1.014, 1.023]}  % experiments/FIGS/data/paired_v1err/finetune/ratios.json :: ratios[51].ratio, ratios[51].ci95 :: sha256 3e5cf25e859b68c2 :: 162+mass over 162, leg1 N1000000 1-macro_auc: paired geometric mean [95% interval]
\newcommand{\PairedFtJciiESixOnesevenMassOverOneseven}{1.03~[1.02, 1.04]}  % experiments/FIGS/data/paired_v1err/finetune/ratios.json :: ratios[52].ratio, ratios[52].ci95 :: sha256 3e5cf25e859b68c2 :: 17+mass over 17, leg1 N1000000 1-macro_auc: paired geometric mean [95% interval]
\newcommand{\PairedFtJciiESixRandomOverOneeighteight}{1.073~[1.067, 1.080]}  % experiments/FIGS/data/paired_v1err/finetune/ratios.json :: ratios[53].ratio, ratios[53].ci95 :: sha256 3e5cf25e859b68c2 :: random over 188, leg1 N1000000 1-macro_auc: paired geometric mean [95% interval]
\newcommand{\PairedFtJciiESixRandomOverOnesixtwo}{1.072~[1.065, 1.078]}  % experiments/FIGS/data/paired_v1err/finetune/ratios.json :: ratios[54].ratio, ratios[54].ci95 :: sha256 3e5cf25e859b68c2 :: random over 162, leg1 N1000000 1-macro_auc: paired geometric mean [95% interval]
\newcommand{\PairedFtJciiESixRandomOverFourthree}{1.03~[1.02, 1.04]}  % experiments/FIGS/data/paired_v1err/finetune/ratios.json :: ratios[55].ratio, ratios[55].ci95 :: sha256 3e5cf25e859b68c2 :: random over 43, leg1 N1000000 1-macro_auc: paired geometric mean [95% interval]
\newcommand{\PairedFtJciiESixRandomOverOneseven}{0.94~[0.93, 0.95]}  % experiments/FIGS/data/paired_v1err/finetune/ratios.json :: ratios[56].ratio, ratios[56].ci95 :: sha256 3e5cf25e859b68c2 :: random over 17, leg1 N1000000 1-macro_auc: paired geometric mean [95% interval]
\newcommand{\PairedFtJciiESixRandomDrawTwoOverRandomDrawOne}{1.00~[0.99, 1.01]}  % experiments/FIGS/data/paired_v1err/finetune/ratios.json :: ratios[57].ratio, ratios[57].ci95 :: sha256 3e5cf25e859b68c2 :: random draw 2 over random draw 1, leg1 N1000000 1-macro_auc: paired geometric mean [95% interval]
\newcommand{\PairedFtJciiESixRandomDrawThreeOverRandomDrawOne}{0.996~[0.989, 1.003]}  % experiments/FIGS/data/paired_v1err/finetune/ratios.json :: ratios[58].ratio, ratios[58].ci95 :: sha256 3e5cf25e859b68c2 :: random draw 3 over random draw 1, leg1 N1000000 1-macro_auc: paired geometric mean [95% interval]
\newcommand{\PairedFtJciiESixRandomDrawThreeOverRandomDrawTwo}{0.997~[0.989, 1.004]}  % experiments/FIGS/data/paired_v1err/finetune/ratios.json :: ratios[59].ratio, ratios[59].ci95 :: sha256 3e5cf25e859b68c2 :: random draw 3 over random draw 2, leg1 N1000000 1-macro_auc: paired geometric mean [95% interval]
\newcommand{\PairedFtJcEThreeOnesixtwoOverOneeighteight}{0.96~[0.83, 1.11]}  % experiments/FIGS/data/paired_v1err/finetune/ratios.json :: ratios[60].ratio, ratios[60].ci95 :: sha256 3e5cf25e859b68c2 :: 162 over 188, leg2 N1000 1-macro_auc: paired geometric mean [95% interval]
\newcommand{\PairedFtJcEThreeFourthreeOverOneeighteight}{1.01~[0.84, 1.21]}  % experiments/FIGS/data/paired_v1err/finetune/ratios.json :: ratios[61].ratio, ratios[61].ci95 :: sha256 3e5cf25e859b68c2 :: 43 over 188, leg2 N1000 1-macro_auc: paired geometric mean [95% interval]
\newcommand{\PairedFtJcEThreeOnesevenOverOneeighteight}{1.07~[0.93, 1.23]}  % experiments/FIGS/data/paired_v1err/finetune/ratios.json :: ratios[62].ratio, ratios[62].ci95 :: sha256 3e5cf25e859b68c2 :: 17 over 188, leg2 N1000 1-macro_auc: paired geometric mean [95% interval]
\newcommand{\PairedFtJcEThreeFourthreeOverOnesixtwo}{1.05~[0.96, 1.15]}  % experiments/FIGS/data/paired_v1err/finetune/ratios.json :: ratios[63].ratio, ratios[63].ci95 :: sha256 3e5cf25e859b68c2 :: 43 over 162, leg2 N1000 1-macro_auc: paired geometric mean [95% interval]
\newcommand{\PairedFtJcEThreeOnesevenOverOnesixtwo}{1.12~[1.07, 1.17]}  % experiments/FIGS/data/paired_v1err/finetune/ratios.json :: ratios[64].ratio, ratios[64].ci95 :: sha256 3e5cf25e859b68c2 :: 17 over 162, leg2 N1000 1-macro_auc: paired geometric mean [95% interval]
\newcommand{\PairedFtJcEThreeOnesevenOverFourthree}{1.06~[0.98, 1.15]}  % experiments/FIGS/data/paired_v1err/finetune/ratios.json :: ratios[65].ratio, ratios[65].ci95 :: sha256 3e5cf25e859b68c2 :: 17 over 43, leg2 N1000 1-macro_auc: paired geometric mean [95% interval]
\newcommand{\PairedFtJcEThreeOnesixtwoMassOverOnesixtwo}{0.91~[0.88, 0.94]}  % experiments/FIGS/data/paired_v1err/finetune/ratios.json :: ratios[66].ratio, ratios[66].ci95 :: sha256 3e5cf25e859b68c2 :: 162+mass over 162, leg2 N1000 1-macro_auc: paired geometric mean [95% interval]
\newcommand{\PairedFtJcEThreeOnesevenMassOverOneseven}{1.02~[0.92, 1.13]}  % experiments/FIGS/data/paired_v1err/finetune/ratios.json :: ratios[67].ratio, ratios[67].ci95 :: sha256 3e5cf25e859b68c2 :: 17+mass over 17, leg2 N1000 1-macro_auc: paired geometric mean [95% interval]
\newcommand{\PairedFtJcEThreeRandomOverOneeighteight}{1.01~[0.78, 1.32]}  % experiments/FIGS/data/paired_v1err/finetune/ratios.json :: ratios[68].ratio, ratios[68].ci95 :: sha256 3e5cf25e859b68c2 :: random over 188, leg2 N1000 1-macro_auc: paired geometric mean [95% interval]
\newcommand{\PairedFtJcEThreeRandomOverOnesixtwo}{1.12~[1.04, 1.20]}  % experiments/FIGS/data/paired_v1err/finetune/ratios.json :: ratios[69].ratio, ratios[69].ci95 :: sha256 3e5cf25e859b68c2 :: random over 162, leg2 N1000 1-macro_auc: paired geometric mean [95% interval]
\newcommand{\PairedFtJcEThreeRandomOverFourthree}{1.06~[0.79, 1.41]}  % experiments/FIGS/data/paired_v1err/finetune/ratios.json :: ratios[70].ratio, ratios[70].ci95 :: sha256 3e5cf25e859b68c2 :: random over 43, leg2 N1000 1-macro_auc: paired geometric mean [95% interval]
\newcommand{\PairedFtJcEThreeRandomOverOneseven}{0.98~[0.84, 1.13]}  % experiments/FIGS/data/paired_v1err/finetune/ratios.json :: ratios[71].ratio, ratios[71].ci95 :: sha256 3e5cf25e859b68c2 :: random over 17, leg2 N1000 1-macro_auc: paired geometric mean [95% interval]
\newcommand{\PairedFtJcEThreeRandomDrawTwoOverRandomDrawOne}{0.954~[0.949, 0.958]}  % experiments/FIGS/data/paired_v1err/finetune/ratios.json :: ratios[72].ratio, ratios[72].ci95 :: sha256 3e5cf25e859b68c2 :: random draw 2 over random draw 1, leg2 N1000 1-macro_auc: paired geometric mean [95% interval]
\newcommand{\PairedFtJcEThreeRandomDrawThreeOverRandomDrawOne}{1.013~[1.009, 1.017]}  % experiments/FIGS/data/paired_v1err/finetune/ratios.json :: ratios[73].ratio, ratios[73].ci95 :: sha256 3e5cf25e859b68c2 :: random draw 3 over random draw 1, leg2 N1000 1-macro_auc: paired geometric mean [95% interval]
\newcommand{\PairedFtJcEThreeRandomDrawThreeOverRandomDrawTwo}{1.062~[1.057, 1.067]}  % experiments/FIGS/data/paired_v1err/finetune/ratios.json :: ratios[74].ratio, ratios[74].ci95 :: sha256 3e5cf25e859b68c2 :: random draw 3 over random draw 2, leg2 N1000 1-macro_auc: paired geometric mean [95% interval]
\newcommand{\PairedFtJcEFourOnesixtwoOverOneeighteight}{0.98~[0.93, 1.02]}  % experiments/FIGS/data/paired_v1err/finetune/ratios.json :: ratios[75].ratio, ratios[75].ci95 :: sha256 3e5cf25e859b68c2 :: 162 over 188, leg2 N10000 1-macro_auc: paired geometric mean [95% interval]
\newcommand{\PairedFtJcEFourFourthreeOverOneeighteight}{0.98~[0.93, 1.04]}  % experiments/FIGS/data/paired_v1err/finetune/ratios.json :: ratios[76].ratio, ratios[76].ci95 :: sha256 3e5cf25e859b68c2 :: 43 over 188, leg2 N10000 1-macro_auc: paired geometric mean [95% interval]
\newcommand{\PairedFtJcEFourOnesevenOverOneeighteight}{1.07~[1.02, 1.11]}  % experiments/FIGS/data/paired_v1err/finetune/ratios.json :: ratios[77].ratio, ratios[77].ci95 :: sha256 3e5cf25e859b68c2 :: 17 over 188, leg2 N10000 1-macro_auc: paired geometric mean [95% interval]
\newcommand{\PairedFtJcEFourFourthreeOverOnesixtwo}{1.01~[0.97, 1.04]}  % experiments/FIGS/data/paired_v1err/finetune/ratios.json :: ratios[78].ratio, ratios[78].ci95 :: sha256 3e5cf25e859b68c2 :: 43 over 162, leg2 N10000 1-macro_auc: paired geometric mean [95% interval]
\newcommand{\PairedFtJcEFourOnesevenOverOnesixtwo}{1.09~[1.06, 1.12]}  % experiments/FIGS/data/paired_v1err/finetune/ratios.json :: ratios[79].ratio, ratios[79].ci95 :: sha256 3e5cf25e859b68c2 :: 17 over 162, leg2 N10000 1-macro_auc: paired geometric mean [95% interval]
\newcommand{\PairedFtJcEFourOnesevenOverFourthree}{1.08~[1.07, 1.10]}  % experiments/FIGS/data/paired_v1err/finetune/ratios.json :: ratios[80].ratio, ratios[80].ci95 :: sha256 3e5cf25e859b68c2 :: 17 over 43, leg2 N10000 1-macro_auc: paired geometric mean [95% interval]
\newcommand{\PairedFtJcEFourOnesixtwoMassOverOnesixtwo}{1.01~[0.98, 1.04]}  % experiments/FIGS/data/paired_v1err/finetune/ratios.json :: ratios[81].ratio, ratios[81].ci95 :: sha256 3e5cf25e859b68c2 :: 162+mass over 162, leg2 N10000 1-macro_auc: paired geometric mean [95% interval]
\newcommand{\PairedFtJcEFourOnesevenMassOverOneseven}{1.05~[1.02, 1.08]}  % experiments/FIGS/data/paired_v1err/finetune/ratios.json :: ratios[82].ratio, ratios[82].ci95 :: sha256 3e5cf25e859b68c2 :: 17+mass over 17, leg2 N10000 1-macro_auc: paired geometric mean [95% interval]
\newcommand{\PairedFtJcEFourRandomOverOneeighteight}{1.08~[1.03, 1.14]}  % experiments/FIGS/data/paired_v1err/finetune/ratios.json :: ratios[83].ratio, ratios[83].ci95 :: sha256 3e5cf25e859b68c2 :: random over 188, leg2 N10000 1-macro_auc: paired geometric mean [95% interval]
\newcommand{\PairedFtJcEFourRandomOverOnesixtwo}{1.11~[0.99, 1.23]}  % experiments/FIGS/data/paired_v1err/finetune/ratios.json :: ratios[84].ratio, ratios[84].ci95 :: sha256 3e5cf25e859b68c2 :: random over 162, leg2 N10000 1-macro_auc: paired geometric mean [95% interval]
\newcommand{\PairedFtJcEFourRandomOverFourthree}{1.10~[0.99, 1.21]}  % experiments/FIGS/data/paired_v1err/finetune/ratios.json :: ratios[85].ratio, ratios[85].ci95 :: sha256 3e5cf25e859b68c2 :: random over 43, leg2 N10000 1-macro_auc: paired geometric mean [95% interval]
\newcommand{\PairedFtJcEFourRandomOverOneseven}{1.02~[0.94, 1.10]}  % experiments/FIGS/data/paired_v1err/finetune/ratios.json :: ratios[86].ratio, ratios[86].ci95 :: sha256 3e5cf25e859b68c2 :: random over 17, leg2 N10000 1-macro_auc: paired geometric mean [95% interval]
\newcommand{\PairedFtJcEFourRandomDrawTwoOverRandomDrawOne}{0.960~[0.955, 0.965]}  % experiments/FIGS/data/paired_v1err/finetune/ratios.json :: ratios[87].ratio, ratios[87].ci95 :: sha256 3e5cf25e859b68c2 :: random draw 2 over random draw 1, leg2 N10000 1-macro_auc: paired geometric mean [95% interval]
\newcommand{\PairedFtJcEFourRandomDrawThreeOverRandomDrawOne}{1.008~[1.004, 1.012]}  % experiments/FIGS/data/paired_v1err/finetune/ratios.json :: ratios[88].ratio, ratios[88].ci95 :: sha256 3e5cf25e859b68c2 :: random draw 3 over random draw 1, leg2 N10000 1-macro_auc: paired geometric mean [95% interval]
\newcommand{\PairedFtJcEFourRandomDrawThreeOverRandomDrawTwo}{1.050~[1.045, 1.055]}  % experiments/FIGS/data/paired_v1err/finetune/ratios.json :: ratios[89].ratio, ratios[89].ci95 :: sha256 3e5cf25e859b68c2 :: random draw 3 over random draw 2, leg2 N10000 1-macro_auc: paired geometric mean [95% interval]
\newcommand{\PairedFtJcEFiveOnesixtwoOverOneeighteight}{0.983~[0.968, 0.998]}  % experiments/FIGS/data/paired_v1err/finetune/ratios.json :: ratios[90].ratio, ratios[90].ci95 :: sha256 3e5cf25e859b68c2 :: 162 over 188, leg2 N100000 1-macro_auc: paired geometric mean [95% interval]
\newcommand{\PairedFtJcEFiveFourthreeOverOneeighteight}{1.015~[1.002, 1.029]}  % experiments/FIGS/data/paired_v1err/finetune/ratios.json :: ratios[91].ratio, ratios[91].ci95 :: sha256 3e5cf25e859b68c2 :: 43 over 188, leg2 N100000 1-macro_auc: paired geometric mean [95% interval]
\newcommand{\PairedFtJcEFiveOnesevenOverOneeighteight}{1.07~[1.06, 1.09]}  % experiments/FIGS/data/paired_v1err/finetune/ratios.json :: ratios[92].ratio, ratios[92].ci95 :: sha256 3e5cf25e859b68c2 :: 17 over 188, leg2 N100000 1-macro_auc: paired geometric mean [95% interval]
\newcommand{\PairedFtJcEFiveFourthreeOverOnesixtwo}{1.03~[1.02, 1.05]}  % experiments/FIGS/data/paired_v1err/finetune/ratios.json :: ratios[93].ratio, ratios[93].ci95 :: sha256 3e5cf25e859b68c2 :: 43 over 162, leg2 N100000 1-macro_auc: paired geometric mean [95% interval]
\newcommand{\PairedFtJcEFiveOnesevenOverOnesixtwo}{1.091~[1.088, 1.094]}  % experiments/FIGS/data/paired_v1err/finetune/ratios.json :: ratios[94].ratio, ratios[94].ci95 :: sha256 3e5cf25e859b68c2 :: 17 over 162, leg2 N100000 1-macro_auc: paired geometric mean [95% interval]
\newcommand{\PairedFtJcEFiveOnesevenOverFourthree}{1.06~[1.04, 1.07]}  % experiments/FIGS/data/paired_v1err/finetune/ratios.json :: ratios[95].ratio, ratios[95].ci95 :: sha256 3e5cf25e859b68c2 :: 17 over 43, leg2 N100000 1-macro_auc: paired geometric mean [95% interval]
\newcommand{\PairedFtJcEFiveOnesixtwoMassOverOnesixtwo}{1.016~[1.008, 1.024]}  % experiments/FIGS/data/paired_v1err/finetune/ratios.json :: ratios[96].ratio, ratios[96].ci95 :: sha256 3e5cf25e859b68c2 :: 162+mass over 162, leg2 N100000 1-macro_auc: paired geometric mean [95% interval]
\newcommand{\PairedFtJcEFiveOnesevenMassOverOneseven}{1.02~[1.01, 1.03]}  % experiments/FIGS/data/paired_v1err/finetune/ratios.json :: ratios[97].ratio, ratios[97].ci95 :: sha256 3e5cf25e859b68c2 :: 17+mass over 17, leg2 N100000 1-macro_auc: paired geometric mean [95% interval]
\newcommand{\PairedFtJcEFiveRandomOverOneeighteight}{1.044~[1.039, 1.049]}  % experiments/FIGS/data/paired_v1err/finetune/ratios.json :: ratios[98].ratio, ratios[98].ci95 :: sha256 3e5cf25e859b68c2 :: random over 188, leg2 N100000 1-macro_auc: paired geometric mean [95% interval]
\newcommand{\PairedFtJcEFiveRandomOverOnesixtwo}{1.06~[1.04, 1.09]}  % experiments/FIGS/data/paired_v1err/finetune/ratios.json :: ratios[99].ratio, ratios[99].ci95 :: sha256 3e5cf25e859b68c2 :: random over 162, leg2 N100000 1-macro_auc: paired geometric mean [95% interval]
\newcommand{\PairedFtJcEFiveRandomOverFourthree}{1.03~[1.01, 1.04]}  % experiments/FIGS/data/paired_v1err/finetune/ratios.json :: ratios[100].ratio, ratios[100].ci95 :: sha256 3e5cf25e859b68c2 :: random over 43, leg2 N100000 1-macro_auc: paired geometric mean [95% interval]
\newcommand{\PairedFtJcEFiveRandomOverOneseven}{0.97~[0.95, 0.99]}  % experiments/FIGS/data/paired_v1err/finetune/ratios.json :: ratios[101].ratio, ratios[101].ci95 :: sha256 3e5cf25e859b68c2 :: random over 17, leg2 N100000 1-macro_auc: paired geometric mean [95% interval]
\newcommand{\PairedFtJcEFiveRandomDrawTwoOverRandomDrawOne}{0.992~[0.988, 0.997]}  % experiments/FIGS/data/paired_v1err/finetune/ratios.json :: ratios[102].ratio, ratios[102].ci95 :: sha256 3e5cf25e859b68c2 :: random draw 2 over random draw 1, leg2 N100000 1-macro_auc: paired geometric mean [95% interval]
\newcommand{\PairedFtJcEFiveRandomDrawThreeOverRandomDrawOne}{1.007~[1.003, 1.011]}  % experiments/FIGS/data/paired_v1err/finetune/ratios.json :: ratios[103].ratio, ratios[103].ci95 :: sha256 3e5cf25e859b68c2 :: random draw 3 over random draw 1, leg2 N100000 1-macro_auc: paired geometric mean [95% interval]
\newcommand{\PairedFtJcEFiveRandomDrawThreeOverRandomDrawTwo}{1.015~[1.011, 1.020]}  % experiments/FIGS/data/paired_v1err/finetune/ratios.json :: ratios[104].ratio, ratios[104].ci95 :: sha256 3e5cf25e859b68c2 :: random draw 3 over random draw 2, leg2 N100000 1-macro_auc: paired geometric mean [95% interval]
\newcommand{\PairedFtJcESixOnesixtwoOverOneeighteight}{1.00~[0.99, 1.01]}  % experiments/FIGS/data/paired_v1err/finetune/ratios.json :: ratios[105].ratio, ratios[105].ci95 :: sha256 3e5cf25e859b68c2 :: 162 over 188, leg2 N1000000 1-macro_auc: paired geometric mean [95% interval]
\newcommand{\PairedFtJcESixFourthreeOverOneeighteight}{1.02~[1.01, 1.03]}  % experiments/FIGS/data/paired_v1err/finetune/ratios.json :: ratios[106].ratio, ratios[106].ci95 :: sha256 3e5cf25e859b68c2 :: 43 over 188, leg2 N1000000 1-macro_auc: paired geometric mean [95% interval]
\newcommand{\PairedFtJcESixOnesevenOverOneeighteight}{1.08~[1.07, 1.09]}  % experiments/FIGS/data/paired_v1err/finetune/ratios.json :: ratios[107].ratio, ratios[107].ci95 :: sha256 3e5cf25e859b68c2 :: 17 over 188, leg2 N1000000 1-macro_auc: paired geometric mean [95% interval]
\newcommand{\PairedFtJcESixFourthreeOverOnesixtwo}{1.02~[1.01, 1.03]}  % experiments/FIGS/data/paired_v1err/finetune/ratios.json :: ratios[108].ratio, ratios[108].ci95 :: sha256 3e5cf25e859b68c2 :: 43 over 162, leg2 N1000000 1-macro_auc: paired geometric mean [95% interval]
\newcommand{\PairedFtJcESixOnesevenOverOnesixtwo}{1.08~[1.07, 1.09]}  % experiments/FIGS/data/paired_v1err/finetune/ratios.json :: ratios[109].ratio, ratios[109].ci95 :: sha256 3e5cf25e859b68c2 :: 17 over 162, leg2 N1000000 1-macro_auc: paired geometric mean [95% interval]
\newcommand{\PairedFtJcESixOnesevenOverFourthree}{1.06~[1.05, 1.07]}  % experiments/FIGS/data/paired_v1err/finetune/ratios.json :: ratios[110].ratio, ratios[110].ci95 :: sha256 3e5cf25e859b68c2 :: 17 over 43, leg2 N1000000 1-macro_auc: paired geometric mean [95% interval]
\newcommand{\PairedFtJcESixOnesixtwoMassOverOnesixtwo}{1.012~[1.004, 1.020]}  % experiments/FIGS/data/paired_v1err/finetune/ratios.json :: ratios[111].ratio, ratios[111].ci95 :: sha256 3e5cf25e859b68c2 :: 162+mass over 162, leg2 N1000000 1-macro_auc: paired geometric mean [95% interval]
\newcommand{\PairedFtJcESixOnesevenMassOverOneseven}{1.02~[1.01, 1.04]}  % experiments/FIGS/data/paired_v1err/finetune/ratios.json :: ratios[112].ratio, ratios[112].ci95 :: sha256 3e5cf25e859b68c2 :: 17+mass over 17, leg2 N1000000 1-macro_auc: paired geometric mean [95% interval]
\newcommand{\PairedFtJcESixRandomOverOneeighteight}{1.046~[1.038, 1.053]}  % experiments/FIGS/data/paired_v1err/finetune/ratios.json :: ratios[113].ratio, ratios[113].ci95 :: sha256 3e5cf25e859b68c2 :: random over 188, leg2 N1000000 1-macro_auc: paired geometric mean [95% interval]
\newcommand{\PairedFtJcESixRandomOverOnesixtwo}{1.047~[1.038, 1.055]}  % experiments/FIGS/data/paired_v1err/finetune/ratios.json :: ratios[114].ratio, ratios[114].ci95 :: sha256 3e5cf25e859b68c2 :: random over 162, leg2 N1000000 1-macro_auc: paired geometric mean [95% interval]
\newcommand{\PairedFtJcESixRandomOverFourthree}{1.03~[1.02, 1.04]}  % experiments/FIGS/data/paired_v1err/finetune/ratios.json :: ratios[115].ratio, ratios[115].ci95 :: sha256 3e5cf25e859b68c2 :: random over 43, leg2 N1000000 1-macro_auc: paired geometric mean [95% interval]
\newcommand{\PairedFtJcESixRandomOverOneseven}{0.97~[0.96, 0.98]}  % experiments/FIGS/data/paired_v1err/finetune/ratios.json :: ratios[116].ratio, ratios[116].ci95 :: sha256 3e5cf25e859b68c2 :: random over 17, leg2 N1000000 1-macro_auc: paired geometric mean [95% interval]
\newcommand{\PairedFtJcESixRandomDrawTwoOverRandomDrawOne}{1.001~[0.998, 1.004]}  % experiments/FIGS/data/paired_v1err/finetune/ratios.json :: ratios[117].ratio, ratios[117].ci95 :: sha256 3e5cf25e859b68c2 :: random draw 2 over random draw 1, leg2 N1000000 1-macro_auc: paired geometric mean [95% interval]
\newcommand{\PairedFtJcESixRandomDrawThreeOverRandomDrawOne}{1.005~[1.002, 1.008]}  % experiments/FIGS/data/paired_v1err/finetune/ratios.json :: ratios[118].ratio, ratios[118].ci95 :: sha256 3e5cf25e859b68c2 :: random draw 3 over random draw 1, leg2 N1000000 1-macro_auc: paired geometric mean [95% interval]
\newcommand{\PairedFtJcESixRandomDrawThreeOverRandomDrawTwo}{1.004~[1.001, 1.007]}  % experiments/FIGS/data/paired_v1err/finetune/ratios.json :: ratios[119].ratio, ratios[119].ci95 :: sha256 3e5cf25e859b68c2 :: random draw 3 over random draw 2, leg2 N1000000 1-macro_auc: paired geometric mean [95% interval]
\newcommand{\PairedMassResMlpOnesixtwoOverOneeighteight}{0.990~[0.977, 1.003]}  % experiments/FIGS/data/paired_v1err/mass/ratios.json :: ratios[0].ratio, ratios[0].ci95 :: sha256 d877a7f087260f64 :: 162 over 188, resolution mlp sigma_eff: paired geometric mean [95% interval]
\newcommand{\PairedMassResMlpFourthreeOverOneeighteight}{0.98~[0.96, 0.99]}  % experiments/FIGS/data/paired_v1err/mass/ratios.json :: ratios[1].ratio, ratios[1].ci95 :: sha256 d877a7f087260f64 :: 43 over 188, resolution mlp sigma_eff: paired geometric mean [95% interval]
\newcommand{\PairedMassResMlpOnesevenOverOneeighteight}{0.96~[0.94, 0.98]}  % experiments/FIGS/data/paired_v1err/mass/ratios.json :: ratios[2].ratio, ratios[2].ci95 :: sha256 d877a7f087260f64 :: 17 over 188, resolution mlp sigma_eff: paired geometric mean [95% interval]
\newcommand{\PairedMassResMlpFourthreeOverOnesixtwo}{0.989~[0.978, 0.999]}  % experiments/FIGS/data/paired_v1err/mass/ratios.json :: ratios[3].ratio, ratios[3].ci95 :: sha256 d877a7f087260f64 :: 43 over 162, resolution mlp sigma_eff: paired geometric mean [95% interval]
\newcommand{\PairedMassResMlpOnesevenOverOnesixtwo}{0.97~[0.94, 0.99]}  % experiments/FIGS/data/paired_v1err/mass/ratios.json :: ratios[4].ratio, ratios[4].ci95 :: sha256 d877a7f087260f64 :: 17 over 162, resolution mlp sigma_eff: paired geometric mean [95% interval]
\newcommand{\PairedMassResMlpOnesevenOverFourthree}{0.979~[0.961, 0.997]}  % experiments/FIGS/data/paired_v1err/mass/ratios.json :: ratios[5].ratio, ratios[5].ci95 :: sha256 d877a7f087260f64 :: 17 over 43, resolution mlp sigma_eff: paired geometric mean [95% interval]
\newcommand{\PairedMassResMlpOnesixtwoMassOverOnesixtwo}{0.80~[0.79, 0.81]}  % experiments/FIGS/data/paired_v1err/mass/ratios.json :: ratios[6].ratio, ratios[6].ci95 :: sha256 d877a7f087260f64 :: 162+mass over 162, resolution mlp sigma_eff: paired geometric mean [95% interval]
\newcommand{\PairedMassResMlpOnesevenMassOverOneseven}{0.83~[0.82, 0.84]}  % experiments/FIGS/data/paired_v1err/mass/ratios.json :: ratios[7].ratio, ratios[7].ci95 :: sha256 d877a7f087260f64 :: 17+mass over 17, resolution mlp sigma_eff: paired geometric mean [95% interval]
\newcommand{\PairedMassResRidgeOnesixtwoOverOneeighteight}{0.97~[0.96, 0.98]}  % experiments/FIGS/data/paired_v1err/mass/ratios.json :: ratios[8].ratio, ratios[8].ci95 :: sha256 d877a7f087260f64 :: 162 over 188, resolution ridge sigma_eff: paired geometric mean [95% interval]
\newcommand{\PairedMassResRidgeFourthreeOverOneeighteight}{0.95~[0.94, 0.96]}  % experiments/FIGS/data/paired_v1err/mass/ratios.json :: ratios[9].ratio, ratios[9].ci95 :: sha256 d877a7f087260f64 :: 43 over 188, resolution ridge sigma_eff: paired geometric mean [95% interval]
\newcommand{\PairedMassResRidgeOnesevenOverOneeighteight}{0.88~[0.85, 0.92]}  % experiments/FIGS/data/paired_v1err/mass/ratios.json :: ratios[10].ratio, ratios[10].ci95 :: sha256 d877a7f087260f64 :: 17 over 188, resolution ridge sigma_eff: paired geometric mean [95% interval]
\newcommand{\PairedMassResRidgeFourthreeOverOnesixtwo}{0.98~[0.97, 0.99]}  % experiments/FIGS/data/paired_v1err/mass/ratios.json :: ratios[11].ratio, ratios[11].ci95 :: sha256 d877a7f087260f64 :: 43 over 162, resolution ridge sigma_eff: paired geometric mean [95% interval]
\newcommand{\PairedMassResRidgeOnesevenOverOnesixtwo}{0.91~[0.87, 0.95]}  % experiments/FIGS/data/paired_v1err/mass/ratios.json :: ratios[12].ratio, ratios[12].ci95 :: sha256 d877a7f087260f64 :: 17 over 162, resolution ridge sigma_eff: paired geometric mean [95% interval]
\newcommand{\PairedMassResRidgeOnesevenOverFourthree}{0.93~[0.89, 0.97]}  % experiments/FIGS/data/paired_v1err/mass/ratios.json :: ratios[13].ratio, ratios[13].ci95 :: sha256 d877a7f087260f64 :: 17 over 43, resolution ridge sigma_eff: paired geometric mean [95% interval]
\newcommand{\PairedMassResRidgeOnesixtwoMassOverOnesixtwo}{0.678~[0.672, 0.684]}  % experiments/FIGS/data/paired_v1err/mass/ratios.json :: ratios[14].ratio, ratios[14].ci95 :: sha256 d877a7f087260f64 :: 162+mass over 162, resolution ridge sigma_eff: paired geometric mean [95% interval]
\newcommand{\PairedMassResRidgeOnesevenMassOverOneseven}{0.72~[0.69, 0.76]}  % experiments/FIGS/data/paired_v1err/mass/ratios.json :: ratios[15].ratio, ratios[15].ci95 :: sha256 d877a7f087260f64 :: 17+mass over 17, resolution ridge sigma_eff: paired geometric mean [95% interval]
\newcommand{\VcbNSignal}{1{,}659}  % experiments/FIGS/data/probe_ladder_vcbwindow_mlp2/sall.json :: tasks.bc_vs_rest.n_signal_test :: sha256 5b8b5c7c28602f93 :: X->bc test jets in the window
\newcommand{\VcbNBackground}{19{,}244}  % experiments/FIGS/data/probe_ladder_vcbwindow_mlp2/sall.json :: tasks.bc_vs_rest.n_background_test :: sha256 5b8b5c7c28602f93 :: bq, cs, bqq and QCD test jets in the window
\newcommand{\VcbEpsSixty}{60}  % experiments/FIGS/data/probe_ladder_vcbwindow_mlp2/sall.json :: tasks.bc_vs_rest.eps_s :: sha256 5b8b5c7c28602f93 :: signal efficiency, percent
\newcommand{\VcbOmaLinearOnesixtwo}{\ensuremath{(3.17 \pm 0.05)\times10^{-2}}}  % experiments/FIGS/data/probe_ladder_vcbwindow_mlp2/sall.json :: tasks.bc_vs_rest.arms.l162-s*.linear.log1m_auc :: sha256 5b8b5c7c28602f93 :: 1-AUC = exp of it, mean +- SD over 5 seeds
\newcommand{\VcbLogOmaLinearOnesixtwo}{\ensuremath{-}3.451}  % experiments/FIGS/data/probe_ladder_vcbwindow_mlp2/sall.json :: tasks.bc_vs_rest.arms.l162-s*.linear.log1m_auc :: sha256 5b8b5c7c28602f93 :: mean over seeds
\newcommand{\VcbRejLinearSixtyOnesixtwo}{207\,\ensuremath{\pm}\,19}  % experiments/FIGS/data/probe_ladder_vcbwindow_mlp2/sall.json :: tasks.bc_vs_rest.arms.l162-s*.linear.rejection_at['0.60'].rejection :: sha256 5b8b5c7c28602f93 :: background rejection at this signal efficiency, mean +- SD over seeds
\newcommand{\VcbBkgLeftLinearSixtyOnesixtwo}{93.4}  % experiments/FIGS/data/probe_ladder_vcbwindow_mlp2/sall.json :: tasks.bc_vs_rest.arms.l162-s*.linear.rejection_at['0.60'].n_bkg_pass :: sha256 5b8b5c7c28602f93 :: background jets passing the cut, mean over seeds
\newcommand{\VcbOmaLinearOneseven}{\ensuremath{(5.21 \pm 0.23)\times10^{-2}}}  % experiments/FIGS/data/probe_ladder_vcbwindow_mlp2/sall.json :: tasks.bc_vs_rest.arms.r16q1-s*.linear.log1m_auc :: sha256 5b8b5c7c28602f93 :: 1-AUC = exp of it, mean +- SD over 5 seeds
\newcommand{\VcbLogOmaLinearOneseven}{\ensuremath{-}2.955}  % experiments/FIGS/data/probe_ladder_vcbwindow_mlp2/sall.json :: tasks.bc_vs_rest.arms.r16q1-s*.linear.log1m_auc :: sha256 5b8b5c7c28602f93 :: mean over seeds
\newcommand{\VcbRejLinearSixtyOneseven}{67\,\ensuremath{\pm}\,14}  % experiments/FIGS/data/probe_ladder_vcbwindow_mlp2/sall.json :: tasks.bc_vs_rest.arms.r16q1-s*.linear.rejection_at['0.60'].rejection :: sha256 5b8b5c7c28602f93 :: background rejection at this signal efficiency, mean +- SD over seeds
\newcommand{\VcbBkgLeftLinearSixtyOneseven}{295.0}  % experiments/FIGS/data/probe_ladder_vcbwindow_mlp2/sall.json :: tasks.bc_vs_rest.arms.r16q1-s*.linear.rejection_at['0.60'].n_bkg_pass :: sha256 5b8b5c7c28602f93 :: background jets passing the cut, mean over seeds
\newcommand{\VcbOmaRatioLinear}{1.64}  % experiments/FIGS/data/probe_ladder_vcbwindow_mlp2/sall.json :: tasks.bc_vs_rest.arms.*.linear.log1m_auc :: sha256 5b8b5c7c28602f93 :: seed mean of 1-AUC, 17 classes over 162
\newcommand{\VcbFactorLinear}{1.64}  % experiments/FIGS/data/probe_ladder_vcbwindow_mlp2/sall.json :: tasks.bc_vs_rest.arms.*.linear.log1m_auc :: sha256 5b8b5c7c28602f93 :: exp of the difference of the seed means of log(1-AUC), 17 classes over 162
\newcommand{\VcbOmaMlpOnesixtwo}{\ensuremath{(3.102 \pm 0.023)\times10^{-2}}}  % experiments/FIGS/data/probe_ladder_vcbwindow_mlp2/sall.json :: tasks.bc_vs_rest.arms.l162-s*.mlp.log1m_auc :: sha256 5b8b5c7c28602f93 :: 1-AUC = exp of it, mean +- SD over 5 seeds
\newcommand{\VcbLogOmaMlpOnesixtwo}{\ensuremath{-}3.473}  % experiments/FIGS/data/probe_ladder_vcbwindow_mlp2/sall.json :: tasks.bc_vs_rest.arms.l162-s*.mlp.log1m_auc :: sha256 5b8b5c7c28602f93 :: mean over seeds
\newcommand{\VcbRejMlpSixtyOnesixtwo}{230\,\ensuremath{\pm}\,31}  % experiments/FIGS/data/probe_ladder_vcbwindow_mlp2/sall.json :: tasks.bc_vs_rest.arms.l162-s*.mlp.rejection_at['0.60'].rejection :: sha256 5b8b5c7c28602f93 :: background rejection at this signal efficiency, mean +- SD over seeds
\newcommand{\VcbBkgLeftMlpSixtyOnesixtwo}{84.8}  % experiments/FIGS/data/probe_ladder_vcbwindow_mlp2/sall.json :: tasks.bc_vs_rest.arms.l162-s*.mlp.rejection_at['0.60'].n_bkg_pass :: sha256 5b8b5c7c28602f93 :: background jets passing the cut, mean over seeds
\newcommand{\VcbOmaMlpOneseven}{\ensuremath{(4.66 \pm 0.19)\times10^{-2}}}  % experiments/FIGS/data/probe_ladder_vcbwindow_mlp2/sall.json :: tasks.bc_vs_rest.arms.r16q1-s*.mlp.log1m_auc :: sha256 5b8b5c7c28602f93 :: 1-AUC = exp of it, mean +- SD over 5 seeds
\newcommand{\VcbLogOmaMlpOneseven}{\ensuremath{-}3.066}  % experiments/FIGS/data/probe_ladder_vcbwindow_mlp2/sall.json :: tasks.bc_vs_rest.arms.r16q1-s*.mlp.log1m_auc :: sha256 5b8b5c7c28602f93 :: mean over seeds
\newcommand{\VcbRejMlpSixtyOneseven}{78\,\ensuremath{\pm}\,13}  % experiments/FIGS/data/probe_ladder_vcbwindow_mlp2/sall.json :: tasks.bc_vs_rest.arms.r16q1-s*.mlp.rejection_at['0.60'].rejection :: sha256 5b8b5c7c28602f93 :: background rejection at this signal efficiency, mean +- SD over seeds
\newcommand{\VcbBkgLeftMlpSixtyOneseven}{250.8}  % experiments/FIGS/data/probe_ladder_vcbwindow_mlp2/sall.json :: tasks.bc_vs_rest.arms.r16q1-s*.mlp.rejection_at['0.60'].n_bkg_pass :: sha256 5b8b5c7c28602f93 :: background jets passing the cut, mean over seeds
\newcommand{\VcbOmaRatioMlp}{1.50}  % experiments/FIGS/data/probe_ladder_vcbwindow_mlp2/sall.json :: tasks.bc_vs_rest.arms.*.mlp.log1m_auc :: sha256 5b8b5c7c28602f93 :: seed mean of 1-AUC, 17 classes over 162
\newcommand{\VcbFactorMlp}{1.50}  % experiments/FIGS/data/probe_ladder_vcbwindow_mlp2/sall.json :: tasks.bc_vs_rest.arms.*.mlp.log1m_auc :: sha256 5b8b5c7c28602f93 :: exp of the difference of the seed means of log(1-AUC), 17 classes over 162
\newcommand{\AccLegOneOnesixtwoSonebEFour}{0.32273}  % experiments/FIGS/data/leg1_metrics.json :: summary.l162-s1b.N10000.accuracy_mean :: sha256 3f564e960b2c86a7 :: wave 1, superseded; 3 fine-tuning seeds
\newcommand{\AccSdLegOneOnesixtwoSonebEFour}{0.01894}  % experiments/FIGS/data/leg1_metrics.json :: summary.l162-s1b.N10000.accuracy_sd :: sha256 3f564e960b2c86a7 :: wave 1, superseded
\newcommand{\AccLegOneOnesixtwoSonebEFive}{0.53506}  % experiments/FIGS/data/leg1_metrics.json :: summary.l162-s1b.N100000.accuracy_mean :: sha256 3f564e960b2c86a7 :: wave 1, superseded; 3 fine-tuning seeds
\newcommand{\AccSdLegOneOnesixtwoSonebEFive}{0.00044}  % experiments/FIGS/data/leg1_metrics.json :: summary.l162-s1b.N100000.accuracy_sd :: sha256 3f564e960b2c86a7 :: wave 1, superseded
\newcommand{\AccLegOneOnesixtwoSonebESix}{0.54337}  % experiments/FIGS/data/leg1_metrics.json :: summary.l162-s1b.N1000000.accuracy_mean :: sha256 3f564e960b2c86a7 :: wave 1, superseded; 3 fine-tuning seeds
\newcommand{\AccSdLegOneOnesixtwoSonebESix}{0.00074}  % experiments/FIGS/data/leg1_metrics.json :: summary.l162-s1b.N1000000.accuracy_sd :: sha256 3f564e960b2c86a7 :: wave 1, superseded
\newcommand{\AccLegOneOnesevenStwoEFour}{0.22684}  % experiments/FIGS/data/leg1_metrics.json :: summary.r16q1-s2.N10000.accuracy_mean :: sha256 3f564e960b2c86a7 :: wave 1, superseded; 3 fine-tuning seeds
\newcommand{\AccSdLegOneOnesevenStwoEFour}{0.00805}  % experiments/FIGS/data/leg1_metrics.json :: summary.r16q1-s2.N10000.accuracy_sd :: sha256 3f564e960b2c86a7 :: wave 1, superseded
\newcommand{\AccLegOneOnesevenStwoEFive}{0.49650}  % experiments/FIGS/data/leg1_metrics.json :: summary.r16q1-s2.N100000.accuracy_mean :: sha256 3f564e960b2c86a7 :: wave 1, superseded; 3 fine-tuning seeds
\newcommand{\AccSdLegOneOnesevenStwoEFive}{0.00110}  % experiments/FIGS/data/leg1_metrics.json :: summary.r16q1-s2.N100000.accuracy_sd :: sha256 3f564e960b2c86a7 :: wave 1, superseded
\newcommand{\AccLegOneOnesevenStwoESix}{0.51969}  % experiments/FIGS/data/leg1_metrics.json :: summary.r16q1-s2.N1000000.accuracy_mean :: sha256 3f564e960b2c86a7 :: wave 1, superseded; 3 fine-tuning seeds
\newcommand{\AccSdLegOneOnesevenStwoESix}{0.00072}  % experiments/FIGS/data/leg1_metrics.json :: summary.r16q1-s2.N1000000.accuracy_sd :: sha256 3f564e960b2c86a7 :: wave 1, superseded
\newcommand{\AccLegOneOnesevenSthreeEFour}{0.24203}  % experiments/FIGS/data/leg1_metrics.json :: summary.r16q1-s3.N10000.accuracy_mean :: sha256 3f564e960b2c86a7 :: wave 1, superseded; 3 fine-tuning seeds
\newcommand{\AccSdLegOneOnesevenSthreeEFour}{0.00213}  % experiments/FIGS/data/leg1_metrics.json :: summary.r16q1-s3.N10000.accuracy_sd :: sha256 3f564e960b2c86a7 :: wave 1, superseded
\newcommand{\AccLegOneOnesevenSthreeEFive}{0.49313}  % experiments/FIGS/data/leg1_metrics.json :: summary.r16q1-s3.N100000.accuracy_mean :: sha256 3f564e960b2c86a7 :: wave 1, superseded; 3 fine-tuning seeds
\newcommand{\AccSdLegOneOnesevenSthreeEFive}{0.00294}  % experiments/FIGS/data/leg1_metrics.json :: summary.r16q1-s3.N100000.accuracy_sd :: sha256 3f564e960b2c86a7 :: wave 1, superseded
\newcommand{\AccLegOneOnesevenSthreeESix}{0.51804}  % experiments/FIGS/data/leg1_metrics.json :: summary.r16q1-s3.N1000000.accuracy_mean :: sha256 3f564e960b2c86a7 :: wave 1, superseded; 3 fine-tuning seeds
\newcommand{\AccSdLegOneOnesevenSthreeESix}{0.00049}  % experiments/FIGS/data/leg1_metrics.json :: summary.r16q1-s3.N1000000.accuracy_sd :: sha256 3f564e960b2c86a7 :: wave 1, superseded
\newcommand{\AccLegOneOnesevenSfourEFour}{0.24487}  % experiments/FIGS/data/leg1_metrics.json :: summary.r16q1-s4.N10000.accuracy_mean :: sha256 3f564e960b2c86a7 :: wave 1, superseded; 3 fine-tuning seeds
\newcommand{\AccSdLegOneOnesevenSfourEFour}{0.00645}  % experiments/FIGS/data/leg1_metrics.json :: summary.r16q1-s4.N10000.accuracy_sd :: sha256 3f564e960b2c86a7 :: wave 1, superseded
\newcommand{\AccLegOneOnesevenSfourEFive}{0.49510}  % experiments/FIGS/data/leg1_metrics.json :: summary.r16q1-s4.N100000.accuracy_mean :: sha256 3f564e960b2c86a7 :: wave 1, superseded; 3 fine-tuning seeds
\newcommand{\AccSdLegOneOnesevenSfourEFive}{0.00133}  % experiments/FIGS/data/leg1_metrics.json :: summary.r16q1-s4.N100000.accuracy_sd :: sha256 3f564e960b2c86a7 :: wave 1, superseded
\newcommand{\AccLegOneOnesevenSfourESix}{0.51872}  % experiments/FIGS/data/leg1_metrics.json :: summary.r16q1-s4.N1000000.accuracy_mean :: sha256 3f564e960b2c86a7 :: wave 1, superseded; 3 fine-tuning seeds
\newcommand{\AccSdLegOneOnesevenSfourESix}{0.00049}  % experiments/FIGS/data/leg1_metrics.json :: summary.r16q1-s4.N1000000.accuracy_sd :: sha256 3f564e960b2c86a7 :: wave 1, superseded
\newcommand{\AccLegOneScratchEFour}{0.31369}  % experiments/FIGS/data/leg1_metrics.json :: summary.scratch.N10000.accuracy_mean :: sha256 3f564e960b2c86a7 :: wave 1, superseded; 3 fine-tuning seeds
\newcommand{\AccSdLegOneScratchEFour}{0.00712}  % experiments/FIGS/data/leg1_metrics.json :: summary.scratch.N10000.accuracy_sd :: sha256 3f564e960b2c86a7 :: wave 1, superseded
\newcommand{\AccLegOneScratchEFive}{0.42015}  % experiments/FIGS/data/leg1_metrics.json :: summary.scratch.N100000.accuracy_mean :: sha256 3f564e960b2c86a7 :: wave 1, superseded; 3 fine-tuning seeds
\newcommand{\AccSdLegOneScratchEFive}{0.00267}  % experiments/FIGS/data/leg1_metrics.json :: summary.scratch.N100000.accuracy_sd :: sha256 3f564e960b2c86a7 :: wave 1, superseded
\newcommand{\AccLegOneScratchESix}{0.48621}  % experiments/FIGS/data/leg1_metrics.json :: summary.scratch.N1000000.accuracy_mean :: sha256 3f564e960b2c86a7 :: wave 1, superseded; 3 fine-tuning seeds
\newcommand{\AccSdLegOneScratchESix}{0.00100}  % experiments/FIGS/data/leg1_metrics.json :: summary.scratch.N1000000.accuracy_sd :: sha256 3f564e960b2c86a7 :: wave 1, superseded
\newcommand{\AccLegOneSophonPublicEFour}{0.21613}  % experiments/FIGS/data/leg1_metrics.json :: summary.sophon-public.N10000.accuracy_mean :: sha256 3f564e960b2c86a7 :: wave 1, superseded; 3 fine-tuning seeds
\newcommand{\AccSdLegOneSophonPublicEFour}{0.00734}  % experiments/FIGS/data/leg1_metrics.json :: summary.sophon-public.N10000.accuracy_sd :: sha256 3f564e960b2c86a7 :: wave 1, superseded
\newcommand{\AccLegOneSophonPublicEFive}{0.53187}  % experiments/FIGS/data/leg1_metrics.json :: summary.sophon-public.N100000.accuracy_mean :: sha256 3f564e960b2c86a7 :: wave 1, superseded; 3 fine-tuning seeds
\newcommand{\AccSdLegOneSophonPublicEFive}{0.00231}  % experiments/FIGS/data/leg1_metrics.json :: summary.sophon-public.N100000.accuracy_sd :: sha256 3f564e960b2c86a7 :: wave 1, superseded
\newcommand{\AccLegOneSophonPublicESix}{0.54273}  % experiments/FIGS/data/leg1_metrics.json :: summary.sophon-public.N1000000.accuracy_mean :: sha256 3f564e960b2c86a7 :: wave 1, superseded; 3 fine-tuning seeds
\newcommand{\AccSdLegOneSophonPublicESix}{0.00064}  % experiments/FIGS/data/leg1_metrics.json :: summary.sophon-public.N1000000.accuracy_sd :: sha256 3f564e960b2c86a7 :: wave 1, superseded
\newcommand{\AccLegTwoOnesixtwoSonebEFour}{0.78606}  % experiments/FIGS/data/leg2_metrics.json :: summary.l162-s1b.N10000.accuracy_mean :: sha256 701c58c3240274bb :: wave 1, superseded; 3 fine-tuning seeds
\newcommand{\AccSdLegTwoOnesixtwoSonebEFour}{0.00111}  % experiments/FIGS/data/leg2_metrics.json :: summary.l162-s1b.N10000.accuracy_sd :: sha256 701c58c3240274bb :: wave 1, superseded
\newcommand{\AccLegTwoOnesixtwoSonebEFive}{0.83598}  % experiments/FIGS/data/leg2_metrics.json :: summary.l162-s1b.N100000.accuracy_mean :: sha256 701c58c3240274bb :: wave 1, superseded; 3 fine-tuning seeds
\newcommand{\AccSdLegTwoOnesixtwoSonebEFive}{0.00037}  % experiments/FIGS/data/leg2_metrics.json :: summary.l162-s1b.N100000.accuracy_sd :: sha256 701c58c3240274bb :: wave 1, superseded
\newcommand{\AccLegTwoOnesixtwoSonebESix}{0.84704}  % experiments/FIGS/data/leg2_metrics.json :: summary.l162-s1b.N1000000.accuracy_mean :: sha256 701c58c3240274bb :: wave 1, superseded; 3 fine-tuning seeds
\newcommand{\AccSdLegTwoOnesixtwoSonebESix}{0.00027}  % experiments/FIGS/data/leg2_metrics.json :: summary.l162-s1b.N1000000.accuracy_sd :: sha256 701c58c3240274bb :: wave 1, superseded
\newcommand{\AccLegTwoOnesevenStwoEFour}{0.72156}  % experiments/FIGS/data/leg2_metrics.json :: summary.r16q1-s2.N10000.accuracy_mean :: sha256 701c58c3240274bb :: wave 1, superseded; 3 fine-tuning seeds
\newcommand{\AccSdLegTwoOnesevenStwoEFour}{0.01049}  % experiments/FIGS/data/leg2_metrics.json :: summary.r16q1-s2.N10000.accuracy_sd :: sha256 701c58c3240274bb :: wave 1, superseded
\newcommand{\AccLegTwoOnesevenStwoEFive}{0.82614}  % experiments/FIGS/data/leg2_metrics.json :: summary.r16q1-s2.N100000.accuracy_mean :: sha256 701c58c3240274bb :: wave 1, superseded; 3 fine-tuning seeds
\newcommand{\AccSdLegTwoOnesevenStwoEFive}{0.00058}  % experiments/FIGS/data/leg2_metrics.json :: summary.r16q1-s2.N100000.accuracy_sd :: sha256 701c58c3240274bb :: wave 1, superseded
\newcommand{\AccLegTwoOnesevenStwoESix}{0.83943}  % experiments/FIGS/data/leg2_metrics.json :: summary.r16q1-s2.N1000000.accuracy_mean :: sha256 701c58c3240274bb :: wave 1, superseded; 3 fine-tuning seeds
\newcommand{\AccSdLegTwoOnesevenStwoESix}{0.00052}  % experiments/FIGS/data/leg2_metrics.json :: summary.r16q1-s2.N1000000.accuracy_sd :: sha256 701c58c3240274bb :: wave 1, superseded
\newcommand{\AccLegTwoOnesevenSthreeEFour}{0.72103}  % experiments/FIGS/data/leg2_metrics.json :: summary.r16q1-s3.N10000.accuracy_mean :: sha256 701c58c3240274bb :: wave 1, superseded; 3 fine-tuning seeds
\newcommand{\AccSdLegTwoOnesevenSthreeEFour}{0.01651}  % experiments/FIGS/data/leg2_metrics.json :: summary.r16q1-s3.N10000.accuracy_sd :: sha256 701c58c3240274bb :: wave 1, superseded
\newcommand{\AccLegTwoOnesevenSthreeEFive}{0.82629}  % experiments/FIGS/data/leg2_metrics.json :: summary.r16q1-s3.N100000.accuracy_mean :: sha256 701c58c3240274bb :: wave 1, superseded; 3 fine-tuning seeds
\newcommand{\AccSdLegTwoOnesevenSthreeEFive}{0.00035}  % experiments/FIGS/data/leg2_metrics.json :: summary.r16q1-s3.N100000.accuracy_sd :: sha256 701c58c3240274bb :: wave 1, superseded
\newcommand{\AccLegTwoOnesevenSthreeESix}{0.83944}  % experiments/FIGS/data/leg2_metrics.json :: summary.r16q1-s3.N1000000.accuracy_mean :: sha256 701c58c3240274bb :: wave 1, superseded; 3 fine-tuning seeds
\newcommand{\AccSdLegTwoOnesevenSthreeESix}{0.00019}  % experiments/FIGS/data/leg2_metrics.json :: summary.r16q1-s3.N1000000.accuracy_sd :: sha256 701c58c3240274bb :: wave 1, superseded
\newcommand{\AccLegTwoOnesevenSfourEFour}{0.72953}  % experiments/FIGS/data/leg2_metrics.json :: summary.r16q1-s4.N10000.accuracy_mean :: sha256 701c58c3240274bb :: wave 1, superseded; 3 fine-tuning seeds
\newcommand{\AccSdLegTwoOnesevenSfourEFour}{0.01069}  % experiments/FIGS/data/leg2_metrics.json :: summary.r16q1-s4.N10000.accuracy_sd :: sha256 701c58c3240274bb :: wave 1, superseded
\newcommand{\AccLegTwoOnesevenSfourEFive}{0.82564}  % experiments/FIGS/data/leg2_metrics.json :: summary.r16q1-s4.N100000.accuracy_mean :: sha256 701c58c3240274bb :: wave 1, superseded; 3 fine-tuning seeds
\newcommand{\AccSdLegTwoOnesevenSfourEFive}{0.00041}  % experiments/FIGS/data/leg2_metrics.json :: summary.r16q1-s4.N100000.accuracy_sd :: sha256 701c58c3240274bb :: wave 1, superseded
\newcommand{\AccLegTwoOnesevenSfourESix}{0.83909}  % experiments/FIGS/data/leg2_metrics.json :: summary.r16q1-s4.N1000000.accuracy_mean :: sha256 701c58c3240274bb :: wave 1, superseded; 3 fine-tuning seeds
\newcommand{\AccSdLegTwoOnesevenSfourESix}{0.00037}  % experiments/FIGS/data/leg2_metrics.json :: summary.r16q1-s4.N1000000.accuracy_sd :: sha256 701c58c3240274bb :: wave 1, superseded
\newcommand{\AccLegTwoRefEonearmsSoneRef}{0.81770}  % experiments/FIGS/data/leg2_metrics.json :: summary.ref_e1arms-s1.ref.accuracy_mean :: sha256 701c58c3240274bb :: wave 1, superseded; 1 fine-tuning seeds
\newcommand{\AccSdLegTwoRefEonearmsSoneRef}{---}  % experiments/FIGS/data/leg2_metrics.json :: summary.ref_e1arms-s1.ref.accuracy_sd :: sha256 701c58c3240274bb :: wave 1, superseded
\newcommand{\AccLegTwoRefEonearmsStwoRef}{0.81602}  % experiments/FIGS/data/leg2_metrics.json :: summary.ref_e1arms-s2.ref.accuracy_mean :: sha256 701c58c3240274bb :: wave 1, superseded; 1 fine-tuning seeds
\newcommand{\AccSdLegTwoRefEonearmsStwoRef}{---}  % experiments/FIGS/data/leg2_metrics.json :: summary.ref_e1arms-s2.ref.accuracy_sd :: sha256 701c58c3240274bb :: wave 1, superseded
\newcommand{\AccLegTwoRefEonearmsSthreeRef}{0.81625}  % experiments/FIGS/data/leg2_metrics.json :: summary.ref_e1arms-s3.ref.accuracy_mean :: sha256 701c58c3240274bb :: wave 1, superseded; 1 fine-tuning seeds
\newcommand{\AccSdLegTwoRefEonearmsSthreeRef}{---}  % experiments/FIGS/data/leg2_metrics.json :: summary.ref_e1arms-s3.ref.accuracy_sd :: sha256 701c58c3240274bb :: wave 1, superseded
\newcommand{\AccLegTwoScratchEFour}{0.68230}  % experiments/FIGS/data/leg2_metrics.json :: summary.scratch.N10000.accuracy_mean :: sha256 701c58c3240274bb :: wave 1, superseded; 3 fine-tuning seeds
\newcommand{\AccSdLegTwoScratchEFour}{0.00282}  % experiments/FIGS/data/leg2_metrics.json :: summary.scratch.N10000.accuracy_sd :: sha256 701c58c3240274bb :: wave 1, superseded
\newcommand{\AccLegTwoScratchEFive}{0.77769}  % experiments/FIGS/data/leg2_metrics.json :: summary.scratch.N100000.accuracy_mean :: sha256 701c58c3240274bb :: wave 1, superseded; 3 fine-tuning seeds
\newcommand{\AccSdLegTwoScratchEFive}{0.00208}  % experiments/FIGS/data/leg2_metrics.json :: summary.scratch.N100000.accuracy_sd :: sha256 701c58c3240274bb :: wave 1, superseded
\newcommand{\AccLegTwoScratchESix}{0.82303}  % experiments/FIGS/data/leg2_metrics.json :: summary.scratch.N1000000.accuracy_mean :: sha256 701c58c3240274bb :: wave 1, superseded; 3 fine-tuning seeds
\newcommand{\AccSdLegTwoScratchESix}{0.00134}  % experiments/FIGS/data/leg2_metrics.json :: summary.scratch.N1000000.accuracy_sd :: sha256 701c58c3240274bb :: wave 1, superseded
\newcommand{\AccLegTwoSophonPublicEFour}{0.72739}  % experiments/FIGS/data/leg2_metrics.json :: summary.sophon-public.N10000.accuracy_mean :: sha256 701c58c3240274bb :: wave 1, superseded; 3 fine-tuning seeds
\newcommand{\AccSdLegTwoSophonPublicEFour}{0.00470}  % experiments/FIGS/data/leg2_metrics.json :: summary.sophon-public.N10000.accuracy_sd :: sha256 701c58c3240274bb :: wave 1, superseded
\newcommand{\AccLegTwoSophonPublicEFive}{0.83492}  % experiments/FIGS/data/leg2_metrics.json :: summary.sophon-public.N100000.accuracy_mean :: sha256 701c58c3240274bb :: wave 1, superseded; 3 fine-tuning seeds
\newcommand{\AccSdLegTwoSophonPublicEFive}{0.00049}  % experiments/FIGS/data/leg2_metrics.json :: summary.sophon-public.N100000.accuracy_sd :: sha256 701c58c3240274bb :: wave 1, superseded
\newcommand{\AccLegTwoSophonPublicESix}{0.84632}  % experiments/FIGS/data/leg2_metrics.json :: summary.sophon-public.N1000000.accuracy_mean :: sha256 701c58c3240274bb :: wave 1, superseded; 3 fine-tuning seeds
\newcommand{\AccSdLegTwoSophonPublicESix}{0.00031}  % experiments/FIGS/data/leg2_metrics.json :: summary.sophon-public.N1000000.accuracy_sd :: sha256 701c58c3240274bb :: wave 1, superseded
\newcommand{\FtBudgetSteps}{20{,}000}  % scripts/build_ft_budget_jobs.py :: EPOCHS * STEPS :: sha256 f2ed9faa8e3bd622 :: the largest number of optimizer steps below the full training set
\newcommand{\FtLongSteps}{80{,}000}  % scripts/build_ft_budget_jobs.py :: LONG_EPOCHS * STEPS :: sha256 f2ed9faa8e3bd622 :: the largest number of optimizer steps of the long budget at LONG_N jets (JetClass, JetClass-II)
\newcommand{\DesignNTrainSelected}{109{,}938{,}036}  % experiments/FIGS/data/v2/jc2_train/counts.json :: n_selected :: sha256 b0b5a8b63aa808c0 :: selected jets of the pretrain_train files (experiments/FT/count_jc2_train.py)
\newcommand{\FtBudgetEval}{500}  % scripts/build_ft_budget_jobs.py :: STEPS :: sha256 f2ed9faa8e3bd622 :: steps between validation passes at and above SMALL_N training jets
\newcommand{\FtBudgetEvalSmall}{50}  % scripts/build_ft_budget_jobs.py :: SMALL_STEPS :: sha256 f2ed9faa8e3bd622 :: steps between validation passes below SMALL_N training jets
\newcommand{\FtBudgetEvalSmallBelow}{\ensuremath{10^{5}}}  % scripts/build_ft_budget_jobs.py :: SMALL_N :: sha256 f2ed9faa8e3bd622 :: training jets below which the validation runs every SMALL_STEPS steps
\newcommand{\FtBudgetPatience}{10}  % scripts/build_ft_budget_jobs.py :: PATIENCE :: sha256 f2ed9faa8e3bd622 :: validation passes without improvement before training stops
\newcommand{\FtHeadMult}{50}  % scripts/build_ft_budget_jobs.py :: HEAD_MULT :: sha256 f2ed9faa8e3bd622 :: the head's learning-rate multiplier
\newcommand{\FtHeadCap}{\ensuremath{5\times10^{-3}}}  % scripts/build_ft_budget_jobs.py :: HEAD_CAP :: sha256 f2ed9faa8e3bd622 :: the head's largest learning rate
\newcommand{\FtBudgetBatch}{512}  % scripts/build_ft_budget_jobs.py :: BATCH :: sha256 f2ed9faa8e3bd622
\newcommand{\FtBudgetWeightDecay}{0.01}  % scripts/build_ft_budget_jobs.py :: --optimizer-option weight_decay (every cell) :: sha256 f2ed9faa8e3bd622
\newcommand{\FtSweepPre}{\ensuremath{10^{-4}}, \ensuremath{3\times10^{-4}} and \ensuremath{10^{-3}}}  % scripts/build_ft_budget_jobs.py :: SWEEP_PRE :: sha256 f2ed9faa8e3bd622 :: backbone rates of the sweep
\newcommand{\FtSweepScratch}{\ensuremath{3\times10^{-4}}, \ensuremath{10^{-3}} and \ensuremath{3\times10^{-3}}}  % scripts/build_ft_budget_jobs.py :: SWEEP_SCRATCH :: sha256 f2ed9faa8e3bd622 :: from-scratch rates of the sweep
\newcommand{\FtSweepN}{10{,}000}  % scripts/build_ft_budget_jobs.py :: SWEEP_N :: sha256 f2ed9faa8e3bd622 :: training jets of the sweep
\newcommand{\FtValJets}{200{,}000}  % scripts/build_ft_budget_jobs.py :: VAL_N :: sha256 f2ed9faa8e3bd622 :: the fixed validation set
\newcommand{\FtValJetsBal}{20{,}480}  % scripts/build_ft_budget_jobs.py :: VAL_SIZE[jc2b] :: sha256 f2ed9faa8e3bd622 :: the balanced JetClass-II validation set
\newcommand{\LpVfourCMax}{\ensuremath{10^{4}}}  % experiments/EVAL/probe_bal.py :: C_GRID_V4 :: sha256 a7804216e3179d9d :: largest inverse regularisation of the v4 probes
\newcommand{\JcBalTestPerClass}{5{,}000}  % experiments/FT/stage_jc2_balanced.py :: TEST_PER_CLASS :: sha256 0260341d5ed75fec
\newcommand{\JcBalValPerClass}{109}  % experiments/FT/stage_jc2_balanced.py :: VAL_PER_CLASS :: sha256 0260341d5ed75fec
\newcommand{\JcBalPairTrainPerSide}{\ensuremath{10^{5}}}  % experiments/FT/stage_jc2_balanced.py :: PAIR_TRAIN_PER_SIDE :: sha256 0260341d5ed75fec
\newcommand{\JcBalPairTestPerSide}{\ensuremath{2\times10^{5}}}  % experiments/FT/stage_jc2_balanced.py :: PAIR_TEST_PER_SIDE :: sha256 0260341d5ed75fec
\newcommand{\JcBalNShort}{11}  % experiments/FIGS/data/v2/jc2bal/census.json :: counts.ft_test :: sha256 a4379ee0817254cf :: classes with fewer than 5000 selected test jets
\newcommand{\JcBalRarestTest}{224}  % experiments/FIGS/data/v2/jc2bal/census.json :: counts.ft_test (minimum) :: sha256 a4379ee0817254cf :: label_QCD_bcs
\newcommand{\JcBalShortTrainESix}{11}  % experiments/FIGS/data/v2/jc2bal/census.json :: counts.ft_heldout :: sha256 a4379ee0817254cf :: classes with fewer held-out jets than the 1e6 quota
\newcommand{\JcBalNTest}{903{,}922}  % experiments/FIGS/data/v2/jc2bal/manifest.json :: outputs[test.parquet].rows :: sha256 9c7775be6ff56ac3 :: jets of the balanced test set
\newcommand{\FtLrPre}{\ensuremath{10^{-3}}}  % experiments/FIGS/data/v2/budget/lrsweep_metrics.json :: argmax of the summed validation metric, pretrained :: sha256 7cdc148054ca902c
\newcommand{\FtLrScratchTop}{\ensuremath{10^{-3}}}  % experiments/FIGS/data/v2/budget/lrsweep_metrics.json :: argmax of the validation metric on top, from scratch :: sha256 7cdc148054ca902c :: agrees with LR of the launched job-ft-budget-bench, -jc1, -jc2b, -scratch-full and -long-scratch-full from-scratch cells
\newcommand{\FtLrScratchQg}{\ensuremath{3\times10^{-3}}}  % experiments/FIGS/data/v2/budget/lrsweep_metrics.json :: argmax of the validation metric on quark/gluon, from scratch :: sha256 7cdc148054ca902c :: agrees with LR of the launched job-ft-budget-bench and -scratch-full from-scratch cells
\newcommand{\MlstNParamsFt}{2{,}209{,}270}  % experiments/FIGS/mlst_v4.py :: our_params(2, 17) :: sha256 fa536b4866ffd651 :: trainable parameters of our fine-tuned network, two outputs (top, quark/gluon)
\newcommand{\MlstNParamsFtKin}{2{,}207{,}950}  % experiments/FIGS/mlst_v4.py :: our_params(2, 7) :: sha256 fa536b4866ffd651 :: trainable parameters of our fine-tuned network, two outputs, the 7 kinematic inputs
\newcommand{\MlstNParamsJc}{2{,}213{,}374}  % experiments/FIGS/mlst_v4.py :: our_params(10, 17) :: sha256 fa536b4866ffd651 :: trainable parameters of our fine-tuned network, the ten JetClass classes
\newcommand{\MlstProbeRunsDone}{27}  % experiments/FIGS/data/v2/index.json :: runs with a linear-probe fit (count) :: sha256 c6c4138ee5c86423 :: of the paper's runs
\newcommand{\MlstPaperRuns}{27}  % configs/arms/v2_paper_runs.txt :: runs (count) :: sha256 b04562b14ea8cbc7
\newcommand{\MlstSdRelUnc}{50\%}  % configs/arms/v2_paper_runs.txt :: relative standard error of a standard deviation from the runs of a hierarchy level, 1/sqrt(2(n-1)) :: sha256 b04562b14ea8cbc7
\newcommand{\MlstCmpJcTwoResOverFull}{1.026\,\ensuremath{\pm}\,0.005}  % experiments/FIGS/data/v2/index.json :: linear probe jc2_188 N=1000000: (1 - macro_auc_ovr) of L162 over L188, run by run :: sha256 c6c4138ee5c86423
\newcommand{\MlstCmpJcTwoResOverFullPr}{2.6\% more errors}  % experiments/FIGS/data/v2/index.json :: linear probe jc2_188 N=1000000: (1 - macro_auc_ovr) of L162 over L188, run by run :: sha256 c6c4138ee5c86423 :: prose twin: two significant figures, direction in words
\newcommand{\MlstCmpJcTwoMidOverFull}{1.227\,\ensuremath{\pm}\,0.009}  % experiments/FIGS/data/v2/index.json :: linear probe jc2_188 N=1000000: (1 - macro_auc_ovr) of R42_Q1 over L188, run by run :: sha256 c6c4138ee5c86423
\newcommand{\MlstCmpJcTwoMidOverFullPr}{23\% more errors}  % experiments/FIGS/data/v2/index.json :: linear probe jc2_188 N=1000000: (1 - macro_auc_ovr) of R42_Q1 over L188, run by run :: sha256 c6c4138ee5c86423 :: prose twin: two significant figures, direction in words
\newcommand{\MlstCmpJcTwoCoarseOverFull}{1.557\,\ensuremath{\pm}\,0.019}  % experiments/FIGS/data/v2/index.json :: linear probe jc2_188 N=1000000: (1 - macro_auc_ovr) of R16_Q1 over L188, run by run :: sha256 c6c4138ee5c86423
\newcommand{\MlstCmpJcTwoCoarseOverFullPr}{56\% more errors}  % experiments/FIGS/data/v2/index.json :: linear probe jc2_188 N=1000000: (1 - macro_auc_ovr) of R16_Q1 over L188, run by run :: sha256 c6c4138ee5c86423 :: prose twin: two significant figures, direction in words
\newcommand{\MlstCmpJcTwoResMassOverRes}{1.022 and 1.016}  % experiments/FIGS/data/v2/index.json :: linear probe jc2_188 N=1000000: (1 - macro_auc_ovr) of L162_MASS over L162, run by run :: sha256 c6c4138ee5c86423
\newcommand{\MlstCmpJcTwoResMassOverResPr}{equal within 3\%}  % experiments/FIGS/data/v2/index.json :: linear probe jc2_188 N=1000000: (1 - macro_auc_ovr) of L162_MASS over L162, run by run :: sha256 c6c4138ee5c86423 :: prose twin: two significant figures, direction in words
\newcommand{\MlstCmpJcTwoResMassOverResPrMax}{equal within 3\%}  % experiments/FIGS/data/v2/index.json :: linear probe jc2_188 N=1000000: (1 - macro_auc_ovr) of L162_MASS over L162, run by run :: sha256 c6c4138ee5c86423 :: prose twin of the run farther from 1
\newcommand{\MlstCmpJcTwoCoarseMassOverCoarse}{1.054 and 0.977}  % experiments/FIGS/data/v2/index.json :: linear probe jc2_188 N=1000000: (1 - macro_auc_ovr) of R16_Q1_MASS over R16_Q1, run by run :: sha256 c6c4138ee5c86423
\newcommand{\MlstCmpJcTwoCoarseMassOverCoarsePr}{equal within 6\%}  % experiments/FIGS/data/v2/index.json :: linear probe jc2_188 N=1000000: (1 - macro_auc_ovr) of R16_Q1_MASS over R16_Q1, run by run :: sha256 c6c4138ee5c86423 :: prose twin: two significant figures, direction in words
\newcommand{\MlstCmpJcTwoCoarseMassOverCoarsePrMax}{equal within 6\%}  % experiments/FIGS/data/v2/index.json :: linear probe jc2_188 N=1000000: (1 - macro_auc_ovr) of R16_Q1_MASS over R16_Q1, run by run :: sha256 c6c4138ee5c86423 :: prose twin of the run farther from 1
\newcommand{\MlstCmpJcTwoResMassOverFull}{1.053 and 1.037}  % experiments/FIGS/data/v2/index.json :: linear probe jc2_188 N=1000000: (1 - macro_auc_ovr) of L162_MASS over L188, run by run :: sha256 c6c4138ee5c86423
\newcommand{\MlstCmpJcTwoResMassOverFullPr}{at least 3.7\% more errors}  % experiments/FIGS/data/v2/index.json :: linear probe jc2_188 N=1000000: (1 - macro_auc_ovr) of L162_MASS over L188, run by run :: sha256 c6c4138ee5c86423 :: prose twin: two significant figures, direction in words
\newcommand{\MlstCmpJcTwoResMassOverFullPrMax}{up to 5.3\% more errors}  % experiments/FIGS/data/v2/index.json :: linear probe jc2_188 N=1000000: (1 - macro_auc_ovr) of L162_MASS over L188, run by run :: sha256 c6c4138ee5c86423 :: prose twin of the run farther from 1
\newcommand{\MlstCmpJcTwoCoarseMassOverFull}{1.617 and 1.528}  % experiments/FIGS/data/v2/index.json :: linear probe jc2_188 N=1000000: (1 - macro_auc_ovr) of R16_Q1_MASS over L188, run by run :: sha256 c6c4138ee5c86423
\newcommand{\MlstCmpJcTwoCoarseMassOverFullPr}{at least 53\% more errors}  % experiments/FIGS/data/v2/index.json :: linear probe jc2_188 N=1000000: (1 - macro_auc_ovr) of R16_Q1_MASS over L188, run by run :: sha256 c6c4138ee5c86423 :: prose twin: two significant figures, direction in words
\newcommand{\MlstCmpJcTwoCoarseMassOverFullPrMax}{up to 62\% more errors}  % experiments/FIGS/data/v2/index.json :: linear probe jc2_188 N=1000000: (1 - macro_auc_ovr) of R16_Q1_MASS over L188, run by run :: sha256 c6c4138ee5c86423 :: prose twin of the run farther from 1
\newcommand{\MlstCmpJcTwoFzeroOverFull}{1.424 and 1.462}  % experiments/FIGS/data/v2/index.json :: linear probe jc2_188 N=1000000: (1 - macro_auc_ovr) of FLAV_F0 over L188, run by run :: sha256 c6c4138ee5c86423
\newcommand{\MlstCmpJcTwoFzeroOverFullPr}{at least 42\% more errors}  % experiments/FIGS/data/v2/index.json :: linear probe jc2_188 N=1000000: (1 - macro_auc_ovr) of FLAV_F0 over L188, run by run :: sha256 c6c4138ee5c86423 :: prose twin: two significant figures, direction in words
\newcommand{\MlstCmpJcTwoFzeroOverFullPrMax}{up to 46\% more errors}  % experiments/FIGS/data/v2/index.json :: linear probe jc2_188 N=1000000: (1 - macro_auc_ovr) of FLAV_F0 over L188, run by run :: sha256 c6c4138ee5c86423 :: prose twin of the run farther from 1
\newcommand{\MlstCmpJcTwoFoneOverFull}{1.218 and 1.212}  % experiments/FIGS/data/v2/index.json :: linear probe jc2_188 N=1000000: (1 - macro_auc_ovr) of FLAV_F1 over L188, run by run :: sha256 c6c4138ee5c86423
\newcommand{\MlstCmpJcTwoFoneOverFullPr}{at least 21\% more errors}  % experiments/FIGS/data/v2/index.json :: linear probe jc2_188 N=1000000: (1 - macro_auc_ovr) of FLAV_F1 over L188, run by run :: sha256 c6c4138ee5c86423 :: prose twin: two significant figures, direction in words
\newcommand{\MlstCmpJcTwoFoneOverFullPrMax}{up to 22\% more errors}  % experiments/FIGS/data/v2/index.json :: linear probe jc2_188 N=1000000: (1 - macro_auc_ovr) of FLAV_F1 over L188, run by run :: sha256 c6c4138ee5c86423 :: prose twin of the run farther from 1
\newcommand{\MlstCmpJcTwoFonerOverFull}{1.206 and 1.208}  % experiments/FIGS/data/v2/index.json :: linear probe jc2_188 N=1000000: (1 - macro_auc_ovr) of FLAV_F1R over L188, run by run :: sha256 c6c4138ee5c86423
\newcommand{\MlstCmpJcTwoFonerOverFullPr}{at least 21\% more errors}  % experiments/FIGS/data/v2/index.json :: linear probe jc2_188 N=1000000: (1 - macro_auc_ovr) of FLAV_F1R over L188, run by run :: sha256 c6c4138ee5c86423 :: prose twin: two significant figures, direction in words
\newcommand{\MlstCmpJcTwoFonerOverFullPrMax}{up to 21\% more errors}  % experiments/FIGS/data/v2/index.json :: linear probe jc2_188 N=1000000: (1 - macro_auc_ovr) of FLAV_F1R over L188, run by run :: sha256 c6c4138ee5c86423 :: prose twin of the run farther from 1
\newcommand{\MlstCmpJcTwoFzeroOverCoarse}{0.9275 and 0.9345}  % experiments/FIGS/data/v2/index.json :: linear probe jc2_188 N=1000000: (1 - macro_auc_ovr) of FLAV_F0 over R16_Q1, run by run :: sha256 c6c4138ee5c86423
\newcommand{\MlstCmpJcTwoFzeroOverCoarsePr}{at least 6.5\% fewer errors}  % experiments/FIGS/data/v2/index.json :: linear probe jc2_188 N=1000000: (1 - macro_auc_ovr) of FLAV_F0 over R16_Q1, run by run :: sha256 c6c4138ee5c86423 :: prose twin: two significant figures, direction in words
\newcommand{\MlstCmpJcTwoFzeroOverCoarsePrMax}{up to 7.2\% fewer errors}  % experiments/FIGS/data/v2/index.json :: linear probe jc2_188 N=1000000: (1 - macro_auc_ovr) of FLAV_F0 over R16_Q1, run by run :: sha256 c6c4138ee5c86423 :: prose twin of the run farther from 1
\newcommand{\MlstCmpJcTwoFoneOverCoarse}{0.794 and 0.775}  % experiments/FIGS/data/v2/index.json :: linear probe jc2_188 N=1000000: (1 - macro_auc_ovr) of FLAV_F1 over R16_Q1, run by run :: sha256 c6c4138ee5c86423
\newcommand{\MlstCmpJcTwoFoneOverCoarsePr}{at least 21\% fewer errors}  % experiments/FIGS/data/v2/index.json :: linear probe jc2_188 N=1000000: (1 - macro_auc_ovr) of FLAV_F1 over R16_Q1, run by run :: sha256 c6c4138ee5c86423 :: prose twin: two significant figures, direction in words
\newcommand{\MlstCmpJcTwoFoneOverCoarsePrMax}{up to 23\% fewer errors}  % experiments/FIGS/data/v2/index.json :: linear probe jc2_188 N=1000000: (1 - macro_auc_ovr) of FLAV_F1 over R16_Q1, run by run :: sha256 c6c4138ee5c86423 :: prose twin of the run farther from 1
\newcommand{\MlstCmpJcTwoFonerOverCoarse}{0.786 and 0.772}  % experiments/FIGS/data/v2/index.json :: linear probe jc2_188 N=1000000: (1 - macro_auc_ovr) of FLAV_F1R over R16_Q1, run by run :: sha256 c6c4138ee5c86423
\newcommand{\MlstCmpJcTwoFonerOverCoarsePr}{at least 21\% fewer errors}  % experiments/FIGS/data/v2/index.json :: linear probe jc2_188 N=1000000: (1 - macro_auc_ovr) of FLAV_F1R over R16_Q1, run by run :: sha256 c6c4138ee5c86423 :: prose twin: two significant figures, direction in words
\newcommand{\MlstCmpJcTwoFonerOverCoarsePrMax}{up to 23\% fewer errors}  % experiments/FIGS/data/v2/index.json :: linear probe jc2_188 N=1000000: (1 - macro_auc_ovr) of FLAV_F1R over R16_Q1, run by run :: sha256 c6c4138ee5c86423 :: prose twin of the run farther from 1
\newcommand{\MlstCmpJcTwoFoneOverFzero}{0.856 and 0.829}  % experiments/FIGS/data/v2/index.json :: linear probe jc2_188 N=1000000: (1 - macro_auc_ovr) of FLAV_F1 over FLAV_F0, run by run :: sha256 c6c4138ee5c86423
\newcommand{\MlstCmpJcTwoFoneOverFzeroPr}{at least 14\% fewer errors}  % experiments/FIGS/data/v2/index.json :: linear probe jc2_188 N=1000000: (1 - macro_auc_ovr) of FLAV_F1 over FLAV_F0, run by run :: sha256 c6c4138ee5c86423 :: prose twin: two significant figures, direction in words
\newcommand{\MlstCmpJcTwoFoneOverFzeroPrMax}{up to 17\% fewer errors}  % experiments/FIGS/data/v2/index.json :: linear probe jc2_188 N=1000000: (1 - macro_auc_ovr) of FLAV_F1 over FLAV_F0, run by run :: sha256 c6c4138ee5c86423 :: prose twin of the run farther from 1
\newcommand{\MlstCmpJcTwoFonerOverFzero}{0.847 and 0.826}  % experiments/FIGS/data/v2/index.json :: linear probe jc2_188 N=1000000: (1 - macro_auc_ovr) of FLAV_F1R over FLAV_F0, run by run :: sha256 c6c4138ee5c86423
\newcommand{\MlstCmpJcTwoFonerOverFzeroPr}{at least 15\% fewer errors}  % experiments/FIGS/data/v2/index.json :: linear probe jc2_188 N=1000000: (1 - macro_auc_ovr) of FLAV_F1R over FLAV_F0, run by run :: sha256 c6c4138ee5c86423 :: prose twin: two significant figures, direction in words
\newcommand{\MlstCmpJcTwoFonerOverFzeroPrMax}{up to 17\% fewer errors}  % experiments/FIGS/data/v2/index.json :: linear probe jc2_188 N=1000000: (1 - macro_auc_ovr) of FLAV_F1R over FLAV_F0, run by run :: sha256 c6c4138ee5c86423 :: prose twin of the run farther from 1
\newcommand{\MlstCmpJcTwoFonerOverFone}{0.9899 and 0.9966}  % experiments/FIGS/data/v2/index.json :: linear probe jc2_188 N=1000000: (1 - macro_auc_ovr) of FLAV_F1R over FLAV_F1, run by run :: sha256 c6c4138ee5c86423
\newcommand{\MlstCmpJcTwoFonerOverFonePr}{equal within 2\%}  % experiments/FIGS/data/v2/index.json :: linear probe jc2_188 N=1000000: (1 - macro_auc_ovr) of FLAV_F1R over FLAV_F1, run by run :: sha256 c6c4138ee5c86423 :: prose twin: two significant figures, direction in words
\newcommand{\MlstCmpJcTwoFonerOverFonePrMax}{equal within 2\%}  % experiments/FIGS/data/v2/index.json :: linear probe jc2_188 N=1000000: (1 - macro_auc_ovr) of FLAV_F1R over FLAV_F1, run by run :: sha256 c6c4138ee5c86423 :: prose twin of the run farther from 1
\newcommand{\MlstCmpJcTwoMidOverRes}{1.197\,\ensuremath{\pm}\,0.013}  % experiments/FIGS/data/v2/index.json :: linear probe jc2_188 N=1000000: (1 - macro_auc_ovr) of R42_Q1 over L162, run by run :: sha256 c6c4138ee5c86423
\newcommand{\MlstCmpJcTwoMidOverResPr}{20\% more errors}  % experiments/FIGS/data/v2/index.json :: linear probe jc2_188 N=1000000: (1 - macro_auc_ovr) of R42_Q1 over L162, run by run :: sha256 c6c4138ee5c86423 :: prose twin: two significant figures, direction in words
\newcommand{\MlstCmpJcTwoCoarseOverRes}{1.518\,\ensuremath{\pm}\,0.025}  % experiments/FIGS/data/v2/index.json :: linear probe jc2_188 N=1000000: (1 - macro_auc_ovr) of R16_Q1 over L162, run by run :: sha256 c6c4138ee5c86423
\newcommand{\MlstCmpJcTwoCoarseOverResPr}{52\% more errors}  % experiments/FIGS/data/v2/index.json :: linear probe jc2_188 N=1000000: (1 - macro_auc_ovr) of R16_Q1 over L162, run by run :: sha256 c6c4138ee5c86423 :: prose twin: two significant figures, direction in words
\newcommand{\MlstCmpJcTwoCoarseOverMid}{1.269\,\ensuremath{\pm}\,0.007}  % experiments/FIGS/data/v2/index.json :: linear probe jc2_188 N=1000000: (1 - macro_auc_ovr) of R16_Q1 over R42_Q1, run by run :: sha256 c6c4138ee5c86423
\newcommand{\MlstCmpJcTwoCoarseOverMidPr}{27\% more errors}  % experiments/FIGS/data/v2/index.json :: linear probe jc2_188 N=1000000: (1 - macro_auc_ovr) of R16_Q1 over R42_Q1, run by run :: sha256 c6c4138ee5c86423 :: prose twin: two significant figures, direction in words
\newcommand{\MlstCmpJcTwoRandOverFull}{1.125\,\ensuremath{\pm}\,0.008}  % experiments/FIGS/data/v2/index.json :: linear probe jc2_188 N=1000000: (1 - macro_auc_ovr) of each random partition over the L188 run of its seed, mean +- SD over the partitions :: sha256 c6c4138ee5c86423
\newcommand{\MlstCmpJcTwoRandOverFullPr}{12\% more errors}  % experiments/FIGS/data/v2/index.json :: linear probe jc2_188 N=1000000: (1 - macro_auc_ovr) of each random partition over the L188 run of its seed, mean +- SD over the partitions :: sha256 c6c4138ee5c86423 :: prose twin: two significant figures, direction in words
\newcommand{\MlstCmpJcTwoRandOverMid}{0.918\,\ensuremath{\pm}\,0.007}  % experiments/FIGS/data/v2/index.json :: linear probe jc2_188 N=1000000: (1 - macro_auc_ovr) of each random partition over the R42_Q1 run of its seed, mean +- SD over the partitions :: sha256 c6c4138ee5c86423
\newcommand{\MlstCmpJcTwoRandOverMidPr}{8.2\% fewer errors}  % experiments/FIGS/data/v2/index.json :: linear probe jc2_188 N=1000000: (1 - macro_auc_ovr) of each random partition over the R42_Q1 run of its seed, mean +- SD over the partitions :: sha256 c6c4138ee5c86423 :: prose twin: two significant figures, direction in words
\newcommand{\MlstCmpJcTwoRandOverCoarse}{0.724\,\ensuremath{\pm}\,0.008}  % experiments/FIGS/data/v2/index.json :: linear probe jc2_188 N=1000000: (1 - macro_auc_ovr) of each random partition over the R16_Q1 run of its seed, mean +- SD over the partitions :: sha256 c6c4138ee5c86423
\newcommand{\MlstCmpJcTwoRandOverCoarsePr}{28\% fewer errors}  % experiments/FIGS/data/v2/index.json :: linear probe jc2_188 N=1000000: (1 - macro_auc_ovr) of each random partition over the R16_Q1 run of its seed, mean +- SD over the partitions :: sha256 c6c4138ee5c86423 :: prose twin: two significant figures, direction in words
\newcommand{\MlstCmpJcTwoTrainedResOverFull}{1.026\,\ensuremath{\pm}\,0.005}  % experiments/FIGS/data/v2/index.json :: linear probe jc2_188 N=1000000: (1 - macro_auc_trained) of L162 over L188, run by run :: sha256 c6c4138ee5c86423
\newcommand{\MlstCmpJcTwoTrainedResOverFullPr}{2.6\% more errors}  % experiments/FIGS/data/v2/index.json :: linear probe jc2_188 N=1000000: (1 - macro_auc_trained) of L162 over L188, run by run :: sha256 c6c4138ee5c86423 :: prose twin: two significant figures, direction in words
\newcommand{\MlstCmpJcTwoTrainedMidOverFull}{1.227\,\ensuremath{\pm}\,0.009}  % experiments/FIGS/data/v2/index.json :: linear probe jc2_188 N=1000000: (1 - macro_auc_trained) of R42_Q1 over L188, run by run :: sha256 c6c4138ee5c86423
\newcommand{\MlstCmpJcTwoTrainedMidOverFullPr}{23\% more errors}  % experiments/FIGS/data/v2/index.json :: linear probe jc2_188 N=1000000: (1 - macro_auc_trained) of R42_Q1 over L188, run by run :: sha256 c6c4138ee5c86423 :: prose twin: two significant figures, direction in words
\newcommand{\MlstCmpJcTwoTrainedCoarseOverFull}{1.557\,\ensuremath{\pm}\,0.019}  % experiments/FIGS/data/v2/index.json :: linear probe jc2_188 N=1000000: (1 - macro_auc_trained) of R16_Q1 over L188, run by run :: sha256 c6c4138ee5c86423
\newcommand{\MlstCmpJcTwoTrainedCoarseOverFullPr}{56\% more errors}  % experiments/FIGS/data/v2/index.json :: linear probe jc2_188 N=1000000: (1 - macro_auc_trained) of R16_Q1 over L188, run by run :: sha256 c6c4138ee5c86423 :: prose twin: two significant figures, direction in words
\newcommand{\MlstCmpJcTwoTrainedResMassOverRes}{1.022 and 1.016}  % experiments/FIGS/data/v2/index.json :: linear probe jc2_188 N=1000000: (1 - macro_auc_trained) of L162_MASS over L162, run by run :: sha256 c6c4138ee5c86423
\newcommand{\MlstCmpJcTwoTrainedResMassOverResPr}{equal within 3\%}  % experiments/FIGS/data/v2/index.json :: linear probe jc2_188 N=1000000: (1 - macro_auc_trained) of L162_MASS over L162, run by run :: sha256 c6c4138ee5c86423 :: prose twin: two significant figures, direction in words
\newcommand{\MlstCmpJcTwoTrainedResMassOverResPrMax}{equal within 3\%}  % experiments/FIGS/data/v2/index.json :: linear probe jc2_188 N=1000000: (1 - macro_auc_trained) of L162_MASS over L162, run by run :: sha256 c6c4138ee5c86423 :: prose twin of the run farther from 1
\newcommand{\MlstCmpJcTwoTrainedCoarseMassOverCoarse}{1.054 and 0.977}  % experiments/FIGS/data/v2/index.json :: linear probe jc2_188 N=1000000: (1 - macro_auc_trained) of R16_Q1_MASS over R16_Q1, run by run :: sha256 c6c4138ee5c86423
\newcommand{\MlstCmpJcTwoTrainedCoarseMassOverCoarsePr}{equal within 6\%}  % experiments/FIGS/data/v2/index.json :: linear probe jc2_188 N=1000000: (1 - macro_auc_trained) of R16_Q1_MASS over R16_Q1, run by run :: sha256 c6c4138ee5c86423 :: prose twin: two significant figures, direction in words
\newcommand{\MlstCmpJcTwoTrainedCoarseMassOverCoarsePrMax}{equal within 6\%}  % experiments/FIGS/data/v2/index.json :: linear probe jc2_188 N=1000000: (1 - macro_auc_trained) of R16_Q1_MASS over R16_Q1, run by run :: sha256 c6c4138ee5c86423 :: prose twin of the run farther from 1
\newcommand{\MlstCmpJcTwoTrainedResMassOverFull}{1.053 and 1.037}  % experiments/FIGS/data/v2/index.json :: linear probe jc2_188 N=1000000: (1 - macro_auc_trained) of L162_MASS over L188, run by run :: sha256 c6c4138ee5c86423
\newcommand{\MlstCmpJcTwoTrainedResMassOverFullPr}{at least 3.7\% more errors}  % experiments/FIGS/data/v2/index.json :: linear probe jc2_188 N=1000000: (1 - macro_auc_trained) of L162_MASS over L188, run by run :: sha256 c6c4138ee5c86423 :: prose twin: two significant figures, direction in words
\newcommand{\MlstCmpJcTwoTrainedResMassOverFullPrMax}{up to 5.3\% more errors}  % experiments/FIGS/data/v2/index.json :: linear probe jc2_188 N=1000000: (1 - macro_auc_trained) of L162_MASS over L188, run by run :: sha256 c6c4138ee5c86423 :: prose twin of the run farther from 1
\newcommand{\MlstCmpJcTwoTrainedCoarseMassOverFull}{1.617 and 1.528}  % experiments/FIGS/data/v2/index.json :: linear probe jc2_188 N=1000000: (1 - macro_auc_trained) of R16_Q1_MASS over L188, run by run :: sha256 c6c4138ee5c86423
\newcommand{\MlstCmpJcTwoTrainedCoarseMassOverFullPr}{at least 53\% more errors}  % experiments/FIGS/data/v2/index.json :: linear probe jc2_188 N=1000000: (1 - macro_auc_trained) of R16_Q1_MASS over L188, run by run :: sha256 c6c4138ee5c86423 :: prose twin: two significant figures, direction in words
\newcommand{\MlstCmpJcTwoTrainedCoarseMassOverFullPrMax}{up to 62\% more errors}  % experiments/FIGS/data/v2/index.json :: linear probe jc2_188 N=1000000: (1 - macro_auc_trained) of R16_Q1_MASS over L188, run by run :: sha256 c6c4138ee5c86423 :: prose twin of the run farther from 1
\newcommand{\MlstCmpJcTwoTrainedFzeroOverFull}{1.424 and 1.462}  % experiments/FIGS/data/v2/index.json :: linear probe jc2_188 N=1000000: (1 - macro_auc_trained) of FLAV_F0 over L188, run by run :: sha256 c6c4138ee5c86423
\newcommand{\MlstCmpJcTwoTrainedFzeroOverFullPr}{at least 42\% more errors}  % experiments/FIGS/data/v2/index.json :: linear probe jc2_188 N=1000000: (1 - macro_auc_trained) of FLAV_F0 over L188, run by run :: sha256 c6c4138ee5c86423 :: prose twin: two significant figures, direction in words
\newcommand{\MlstCmpJcTwoTrainedFzeroOverFullPrMax}{up to 46\% more errors}  % experiments/FIGS/data/v2/index.json :: linear probe jc2_188 N=1000000: (1 - macro_auc_trained) of FLAV_F0 over L188, run by run :: sha256 c6c4138ee5c86423 :: prose twin of the run farther from 1
\newcommand{\MlstCmpJcTwoTrainedFoneOverFull}{1.218 and 1.212}  % experiments/FIGS/data/v2/index.json :: linear probe jc2_188 N=1000000: (1 - macro_auc_trained) of FLAV_F1 over L188, run by run :: sha256 c6c4138ee5c86423
\newcommand{\MlstCmpJcTwoTrainedFoneOverFullPr}{at least 21\% more errors}  % experiments/FIGS/data/v2/index.json :: linear probe jc2_188 N=1000000: (1 - macro_auc_trained) of FLAV_F1 over L188, run by run :: sha256 c6c4138ee5c86423 :: prose twin: two significant figures, direction in words
\newcommand{\MlstCmpJcTwoTrainedFoneOverFullPrMax}{up to 22\% more errors}  % experiments/FIGS/data/v2/index.json :: linear probe jc2_188 N=1000000: (1 - macro_auc_trained) of FLAV_F1 over L188, run by run :: sha256 c6c4138ee5c86423 :: prose twin of the run farther from 1
\newcommand{\MlstCmpJcTwoTrainedFonerOverFull}{1.206 and 1.208}  % experiments/FIGS/data/v2/index.json :: linear probe jc2_188 N=1000000: (1 - macro_auc_trained) of FLAV_F1R over L188, run by run :: sha256 c6c4138ee5c86423
\newcommand{\MlstCmpJcTwoTrainedFonerOverFullPr}{at least 21\% more errors}  % experiments/FIGS/data/v2/index.json :: linear probe jc2_188 N=1000000: (1 - macro_auc_trained) of FLAV_F1R over L188, run by run :: sha256 c6c4138ee5c86423 :: prose twin: two significant figures, direction in words
\newcommand{\MlstCmpJcTwoTrainedFonerOverFullPrMax}{up to 21\% more errors}  % experiments/FIGS/data/v2/index.json :: linear probe jc2_188 N=1000000: (1 - macro_auc_trained) of FLAV_F1R over L188, run by run :: sha256 c6c4138ee5c86423 :: prose twin of the run farther from 1
\newcommand{\MlstCmpJcTwoTrainedFzeroOverCoarse}{0.9275 and 0.9345}  % experiments/FIGS/data/v2/index.json :: linear probe jc2_188 N=1000000: (1 - macro_auc_trained) of FLAV_F0 over R16_Q1, run by run :: sha256 c6c4138ee5c86423
\newcommand{\MlstCmpJcTwoTrainedFzeroOverCoarsePr}{at least 6.5\% fewer errors}  % experiments/FIGS/data/v2/index.json :: linear probe jc2_188 N=1000000: (1 - macro_auc_trained) of FLAV_F0 over R16_Q1, run by run :: sha256 c6c4138ee5c86423 :: prose twin: two significant figures, direction in words
\newcommand{\MlstCmpJcTwoTrainedFzeroOverCoarsePrMax}{up to 7.2\% fewer errors}  % experiments/FIGS/data/v2/index.json :: linear probe jc2_188 N=1000000: (1 - macro_auc_trained) of FLAV_F0 over R16_Q1, run by run :: sha256 c6c4138ee5c86423 :: prose twin of the run farther from 1
\newcommand{\MlstCmpJcTwoTrainedFoneOverCoarse}{0.794 and 0.775}  % experiments/FIGS/data/v2/index.json :: linear probe jc2_188 N=1000000: (1 - macro_auc_trained) of FLAV_F1 over R16_Q1, run by run :: sha256 c6c4138ee5c86423
\newcommand{\MlstCmpJcTwoTrainedFoneOverCoarsePr}{at least 21\% fewer errors}  % experiments/FIGS/data/v2/index.json :: linear probe jc2_188 N=1000000: (1 - macro_auc_trained) of FLAV_F1 over R16_Q1, run by run :: sha256 c6c4138ee5c86423 :: prose twin: two significant figures, direction in words
\newcommand{\MlstCmpJcTwoTrainedFoneOverCoarsePrMax}{up to 23\% fewer errors}  % experiments/FIGS/data/v2/index.json :: linear probe jc2_188 N=1000000: (1 - macro_auc_trained) of FLAV_F1 over R16_Q1, run by run :: sha256 c6c4138ee5c86423 :: prose twin of the run farther from 1
\newcommand{\MlstCmpJcTwoTrainedFonerOverCoarse}{0.786 and 0.772}  % experiments/FIGS/data/v2/index.json :: linear probe jc2_188 N=1000000: (1 - macro_auc_trained) of FLAV_F1R over R16_Q1, run by run :: sha256 c6c4138ee5c86423
\newcommand{\MlstCmpJcTwoTrainedFonerOverCoarsePr}{at least 21\% fewer errors}  % experiments/FIGS/data/v2/index.json :: linear probe jc2_188 N=1000000: (1 - macro_auc_trained) of FLAV_F1R over R16_Q1, run by run :: sha256 c6c4138ee5c86423 :: prose twin: two significant figures, direction in words
\newcommand{\MlstCmpJcTwoTrainedFonerOverCoarsePrMax}{up to 23\% fewer errors}  % experiments/FIGS/data/v2/index.json :: linear probe jc2_188 N=1000000: (1 - macro_auc_trained) of FLAV_F1R over R16_Q1, run by run :: sha256 c6c4138ee5c86423 :: prose twin of the run farther from 1
\newcommand{\MlstCmpJcTwoTrainedFoneOverFzero}{0.856 and 0.829}  % experiments/FIGS/data/v2/index.json :: linear probe jc2_188 N=1000000: (1 - macro_auc_trained) of FLAV_F1 over FLAV_F0, run by run :: sha256 c6c4138ee5c86423
\newcommand{\MlstCmpJcTwoTrainedFoneOverFzeroPr}{at least 14\% fewer errors}  % experiments/FIGS/data/v2/index.json :: linear probe jc2_188 N=1000000: (1 - macro_auc_trained) of FLAV_F1 over FLAV_F0, run by run :: sha256 c6c4138ee5c86423 :: prose twin: two significant figures, direction in words
\newcommand{\MlstCmpJcTwoTrainedFoneOverFzeroPrMax}{up to 17\% fewer errors}  % experiments/FIGS/data/v2/index.json :: linear probe jc2_188 N=1000000: (1 - macro_auc_trained) of FLAV_F1 over FLAV_F0, run by run :: sha256 c6c4138ee5c86423 :: prose twin of the run farther from 1
\newcommand{\MlstCmpJcTwoTrainedFonerOverFzero}{0.847 and 0.826}  % experiments/FIGS/data/v2/index.json :: linear probe jc2_188 N=1000000: (1 - macro_auc_trained) of FLAV_F1R over FLAV_F0, run by run :: sha256 c6c4138ee5c86423
\newcommand{\MlstCmpJcTwoTrainedFonerOverFzeroPr}{at least 15\% fewer errors}  % experiments/FIGS/data/v2/index.json :: linear probe jc2_188 N=1000000: (1 - macro_auc_trained) of FLAV_F1R over FLAV_F0, run by run :: sha256 c6c4138ee5c86423 :: prose twin: two significant figures, direction in words
\newcommand{\MlstCmpJcTwoTrainedFonerOverFzeroPrMax}{up to 17\% fewer errors}  % experiments/FIGS/data/v2/index.json :: linear probe jc2_188 N=1000000: (1 - macro_auc_trained) of FLAV_F1R over FLAV_F0, run by run :: sha256 c6c4138ee5c86423 :: prose twin of the run farther from 1
\newcommand{\MlstCmpJcTwoTrainedFonerOverFone}{0.9899 and 0.9966}  % experiments/FIGS/data/v2/index.json :: linear probe jc2_188 N=1000000: (1 - macro_auc_trained) of FLAV_F1R over FLAV_F1, run by run :: sha256 c6c4138ee5c86423
\newcommand{\MlstCmpJcTwoTrainedFonerOverFonePr}{equal within 2\%}  % experiments/FIGS/data/v2/index.json :: linear probe jc2_188 N=1000000: (1 - macro_auc_trained) of FLAV_F1R over FLAV_F1, run by run :: sha256 c6c4138ee5c86423 :: prose twin: two significant figures, direction in words
\newcommand{\MlstCmpJcTwoTrainedFonerOverFonePrMax}{equal within 2\%}  % experiments/FIGS/data/v2/index.json :: linear probe jc2_188 N=1000000: (1 - macro_auc_trained) of FLAV_F1R over FLAV_F1, run by run :: sha256 c6c4138ee5c86423 :: prose twin of the run farther from 1
\newcommand{\MlstCmpJcTwoTrainedMidOverRes}{1.197\,\ensuremath{\pm}\,0.013}  % experiments/FIGS/data/v2/index.json :: linear probe jc2_188 N=1000000: (1 - macro_auc_trained) of R42_Q1 over L162, run by run :: sha256 c6c4138ee5c86423
\newcommand{\MlstCmpJcTwoTrainedMidOverResPr}{20\% more errors}  % experiments/FIGS/data/v2/index.json :: linear probe jc2_188 N=1000000: (1 - macro_auc_trained) of R42_Q1 over L162, run by run :: sha256 c6c4138ee5c86423 :: prose twin: two significant figures, direction in words
\newcommand{\MlstCmpJcTwoTrainedCoarseOverRes}{1.518\,\ensuremath{\pm}\,0.025}  % experiments/FIGS/data/v2/index.json :: linear probe jc2_188 N=1000000: (1 - macro_auc_trained) of R16_Q1 over L162, run by run :: sha256 c6c4138ee5c86423
\newcommand{\MlstCmpJcTwoTrainedCoarseOverResPr}{52\% more errors}  % experiments/FIGS/data/v2/index.json :: linear probe jc2_188 N=1000000: (1 - macro_auc_trained) of R16_Q1 over L162, run by run :: sha256 c6c4138ee5c86423 :: prose twin: two significant figures, direction in words
\newcommand{\MlstCmpJcTwoTrainedCoarseOverMid}{1.269\,\ensuremath{\pm}\,0.007}  % experiments/FIGS/data/v2/index.json :: linear probe jc2_188 N=1000000: (1 - macro_auc_trained) of R16_Q1 over R42_Q1, run by run :: sha256 c6c4138ee5c86423
\newcommand{\MlstCmpJcTwoTrainedCoarseOverMidPr}{27\% more errors}  % experiments/FIGS/data/v2/index.json :: linear probe jc2_188 N=1000000: (1 - macro_auc_trained) of R16_Q1 over R42_Q1, run by run :: sha256 c6c4138ee5c86423 :: prose twin: two significant figures, direction in words
\newcommand{\MlstCmpJcTwoTrainedRandOverFull}{1.125\,\ensuremath{\pm}\,0.008}  % experiments/FIGS/data/v2/index.json :: linear probe jc2_188 N=1000000: (1 - macro_auc_trained) of each random partition over the L188 run of its seed, mean +- SD over the partitions :: sha256 c6c4138ee5c86423
\newcommand{\MlstCmpJcTwoTrainedRandOverFullPr}{12\% more errors}  % experiments/FIGS/data/v2/index.json :: linear probe jc2_188 N=1000000: (1 - macro_auc_trained) of each random partition over the L188 run of its seed, mean +- SD over the partitions :: sha256 c6c4138ee5c86423 :: prose twin: two significant figures, direction in words
\newcommand{\MlstCmpJcTwoTrainedRandOverMid}{0.918\,\ensuremath{\pm}\,0.007}  % experiments/FIGS/data/v2/index.json :: linear probe jc2_188 N=1000000: (1 - macro_auc_trained) of each random partition over the R42_Q1 run of its seed, mean +- SD over the partitions :: sha256 c6c4138ee5c86423
\newcommand{\MlstCmpJcTwoTrainedRandOverMidPr}{8.2\% fewer errors}  % experiments/FIGS/data/v2/index.json :: linear probe jc2_188 N=1000000: (1 - macro_auc_trained) of each random partition over the R42_Q1 run of its seed, mean +- SD over the partitions :: sha256 c6c4138ee5c86423 :: prose twin: two significant figures, direction in words
\newcommand{\MlstCmpJcTwoTrainedRandOverCoarse}{0.724\,\ensuremath{\pm}\,0.008}  % experiments/FIGS/data/v2/index.json :: linear probe jc2_188 N=1000000: (1 - macro_auc_trained) of each random partition over the R16_Q1 run of its seed, mean +- SD over the partitions :: sha256 c6c4138ee5c86423
\newcommand{\MlstCmpJcTwoTrainedRandOverCoarsePr}{28\% fewer errors}  % experiments/FIGS/data/v2/index.json :: linear probe jc2_188 N=1000000: (1 - macro_auc_trained) of each random partition over the R16_Q1 run of its seed, mean +- SD over the partitions :: sha256 c6c4138ee5c86423 :: prose twin: two significant figures, direction in words
\newcommand{\MlstCmpJcTwoResResOverFull}{1.003\,\ensuremath{\pm}\,0.005}  % experiments/FIGS/data/v2/index.json :: linear probe jc2_188 N=1000000: (1 - macro_auc_resonant) of L162 over L188, run by run :: sha256 c6c4138ee5c86423
\newcommand{\MlstCmpJcTwoResResOverFullPr}{equal within 1\%}  % experiments/FIGS/data/v2/index.json :: linear probe jc2_188 N=1000000: (1 - macro_auc_resonant) of L162 over L188, run by run :: sha256 c6c4138ee5c86423 :: prose twin: two significant figures, direction in words
\newcommand{\MlstCmpJcTwoResMidOverFull}{1.262\,\ensuremath{\pm}\,0.012}  % experiments/FIGS/data/v2/index.json :: linear probe jc2_188 N=1000000: (1 - macro_auc_resonant) of R42_Q1 over L188, run by run :: sha256 c6c4138ee5c86423
\newcommand{\MlstCmpJcTwoResMidOverFullPr}{26\% more errors}  % experiments/FIGS/data/v2/index.json :: linear probe jc2_188 N=1000000: (1 - macro_auc_resonant) of R42_Q1 over L188, run by run :: sha256 c6c4138ee5c86423 :: prose twin: two significant figures, direction in words
\newcommand{\MlstCmpJcTwoResCoarseOverFull}{1.590\,\ensuremath{\pm}\,0.032}  % experiments/FIGS/data/v2/index.json :: linear probe jc2_188 N=1000000: (1 - macro_auc_resonant) of R16_Q1 over L188, run by run :: sha256 c6c4138ee5c86423
\newcommand{\MlstCmpJcTwoResCoarseOverFullPr}{59\% more errors}  % experiments/FIGS/data/v2/index.json :: linear probe jc2_188 N=1000000: (1 - macro_auc_resonant) of R16_Q1 over L188, run by run :: sha256 c6c4138ee5c86423 :: prose twin: two significant figures, direction in words
\newcommand{\MlstCmpJcTwoResResMassOverRes}{1.025 and 1.020}  % experiments/FIGS/data/v2/index.json :: linear probe jc2_188 N=1000000: (1 - macro_auc_resonant) of L162_MASS over L162, run by run :: sha256 c6c4138ee5c86423
\newcommand{\MlstCmpJcTwoResResMassOverResPr}{equal within 3\%}  % experiments/FIGS/data/v2/index.json :: linear probe jc2_188 N=1000000: (1 - macro_auc_resonant) of L162_MASS over L162, run by run :: sha256 c6c4138ee5c86423 :: prose twin: two significant figures, direction in words
\newcommand{\MlstCmpJcTwoResResMassOverResPrMax}{equal within 3\%}  % experiments/FIGS/data/v2/index.json :: linear probe jc2_188 N=1000000: (1 - macro_auc_resonant) of L162_MASS over L162, run by run :: sha256 c6c4138ee5c86423 :: prose twin of the run farther from 1
\newcommand{\MlstCmpJcTwoResCoarseMassOverCoarse}{1.08 and 0.96}  % experiments/FIGS/data/v2/index.json :: linear probe jc2_188 N=1000000: (1 - macro_auc_resonant) of R16_Q1_MASS over R16_Q1, run by run :: sha256 c6c4138ee5c86423
\newcommand{\MlstCmpJcTwoResCoarseMassOverCoarsePr}{equal within 8\%}  % experiments/FIGS/data/v2/index.json :: linear probe jc2_188 N=1000000: (1 - macro_auc_resonant) of R16_Q1_MASS over R16_Q1, run by run :: sha256 c6c4138ee5c86423 :: prose twin: two significant figures, direction in words
\newcommand{\MlstCmpJcTwoResCoarseMassOverCoarsePrMax}{equal within 8\%}  % experiments/FIGS/data/v2/index.json :: linear probe jc2_188 N=1000000: (1 - macro_auc_resonant) of R16_Q1_MASS over R16_Q1, run by run :: sha256 c6c4138ee5c86423 :: prose twin of the run farther from 1
\newcommand{\MlstCmpJcTwoResResMassOverFull}{1.033 and 1.018}  % experiments/FIGS/data/v2/index.json :: linear probe jc2_188 N=1000000: (1 - macro_auc_resonant) of L162_MASS over L188, run by run :: sha256 c6c4138ee5c86423
\newcommand{\MlstCmpJcTwoResResMassOverFullPr}{equal within 4\%}  % experiments/FIGS/data/v2/index.json :: linear probe jc2_188 N=1000000: (1 - macro_auc_resonant) of L162_MASS over L188, run by run :: sha256 c6c4138ee5c86423 :: prose twin: two significant figures, direction in words
\newcommand{\MlstCmpJcTwoResResMassOverFullPrMax}{equal within 4\%}  % experiments/FIGS/data/v2/index.json :: linear probe jc2_188 N=1000000: (1 - macro_auc_resonant) of L162_MASS over L188, run by run :: sha256 c6c4138ee5c86423 :: prose twin of the run farther from 1
\newcommand{\MlstCmpJcTwoResCoarseMassOverFull}{1.67 and 1.55}  % experiments/FIGS/data/v2/index.json :: linear probe jc2_188 N=1000000: (1 - macro_auc_resonant) of R16_Q1_MASS over L188, run by run :: sha256 c6c4138ee5c86423
\newcommand{\MlstCmpJcTwoResCoarseMassOverFullPr}{at least 55\% more errors}  % experiments/FIGS/data/v2/index.json :: linear probe jc2_188 N=1000000: (1 - macro_auc_resonant) of R16_Q1_MASS over L188, run by run :: sha256 c6c4138ee5c86423 :: prose twin: two significant figures, direction in words
\newcommand{\MlstCmpJcTwoResCoarseMassOverFullPrMax}{up to 67\% more errors}  % experiments/FIGS/data/v2/index.json :: linear probe jc2_188 N=1000000: (1 - macro_auc_resonant) of R16_Q1_MASS over L188, run by run :: sha256 c6c4138ee5c86423 :: prose twin of the run farther from 1
\newcommand{\MlstCmpJcTwoResFzeroOverFull}{1.408 and 1.449}  % experiments/FIGS/data/v2/index.json :: linear probe jc2_188 N=1000000: (1 - macro_auc_resonant) of FLAV_F0 over L188, run by run :: sha256 c6c4138ee5c86423
\newcommand{\MlstCmpJcTwoResFzeroOverFullPr}{at least 41\% more errors}  % experiments/FIGS/data/v2/index.json :: linear probe jc2_188 N=1000000: (1 - macro_auc_resonant) of FLAV_F0 over L188, run by run :: sha256 c6c4138ee5c86423 :: prose twin: two significant figures, direction in words
\newcommand{\MlstCmpJcTwoResFzeroOverFullPrMax}{up to 45\% more errors}  % experiments/FIGS/data/v2/index.json :: linear probe jc2_188 N=1000000: (1 - macro_auc_resonant) of FLAV_F0 over L188, run by run :: sha256 c6c4138ee5c86423 :: prose twin of the run farther from 1
\newcommand{\MlstCmpJcTwoResFoneOverFull}{1.210 and 1.204}  % experiments/FIGS/data/v2/index.json :: linear probe jc2_188 N=1000000: (1 - macro_auc_resonant) of FLAV_F1 over L188, run by run :: sha256 c6c4138ee5c86423
\newcommand{\MlstCmpJcTwoResFoneOverFullPr}{at least 20\% more errors}  % experiments/FIGS/data/v2/index.json :: linear probe jc2_188 N=1000000: (1 - macro_auc_resonant) of FLAV_F1 over L188, run by run :: sha256 c6c4138ee5c86423 :: prose twin: two significant figures, direction in words
\newcommand{\MlstCmpJcTwoResFoneOverFullPrMax}{up to 21\% more errors}  % experiments/FIGS/data/v2/index.json :: linear probe jc2_188 N=1000000: (1 - macro_auc_resonant) of FLAV_F1 over L188, run by run :: sha256 c6c4138ee5c86423 :: prose twin of the run farther from 1
\newcommand{\MlstCmpJcTwoResFonerOverFull}{1.187 and 1.188}  % experiments/FIGS/data/v2/index.json :: linear probe jc2_188 N=1000000: (1 - macro_auc_resonant) of FLAV_F1R over L188, run by run :: sha256 c6c4138ee5c86423
\newcommand{\MlstCmpJcTwoResFonerOverFullPr}{at least 19\% more errors}  % experiments/FIGS/data/v2/index.json :: linear probe jc2_188 N=1000000: (1 - macro_auc_resonant) of FLAV_F1R over L188, run by run :: sha256 c6c4138ee5c86423 :: prose twin: two significant figures, direction in words
\newcommand{\MlstCmpJcTwoResFonerOverFullPrMax}{up to 19\% more errors}  % experiments/FIGS/data/v2/index.json :: linear probe jc2_188 N=1000000: (1 - macro_auc_resonant) of FLAV_F1R over L188, run by run :: sha256 c6c4138ee5c86423 :: prose twin of the run farther from 1
\newcommand{\MlstCmpJcTwoResFzeroOverCoarse}{0.9069 and 0.9029}  % experiments/FIGS/data/v2/index.json :: linear probe jc2_188 N=1000000: (1 - macro_auc_resonant) of FLAV_F0 over R16_Q1, run by run :: sha256 c6c4138ee5c86423
\newcommand{\MlstCmpJcTwoResFzeroOverCoarsePr}{at least 9.3\% fewer errors}  % experiments/FIGS/data/v2/index.json :: linear probe jc2_188 N=1000000: (1 - macro_auc_resonant) of FLAV_F0 over R16_Q1, run by run :: sha256 c6c4138ee5c86423 :: prose twin: two significant figures, direction in words
\newcommand{\MlstCmpJcTwoResFzeroOverCoarsePrMax}{up to 9.7\% fewer errors}  % experiments/FIGS/data/v2/index.json :: linear probe jc2_188 N=1000000: (1 - macro_auc_resonant) of FLAV_F0 over R16_Q1, run by run :: sha256 c6c4138ee5c86423 :: prose twin of the run farther from 1
\newcommand{\MlstCmpJcTwoResFoneOverCoarse}{0.779 and 0.750}  % experiments/FIGS/data/v2/index.json :: linear probe jc2_188 N=1000000: (1 - macro_auc_resonant) of FLAV_F1 over R16_Q1, run by run :: sha256 c6c4138ee5c86423
\newcommand{\MlstCmpJcTwoResFoneOverCoarsePr}{at least 22\% fewer errors}  % experiments/FIGS/data/v2/index.json :: linear probe jc2_188 N=1000000: (1 - macro_auc_resonant) of FLAV_F1 over R16_Q1, run by run :: sha256 c6c4138ee5c86423 :: prose twin: two significant figures, direction in words
\newcommand{\MlstCmpJcTwoResFoneOverCoarsePrMax}{up to 25\% fewer errors}  % experiments/FIGS/data/v2/index.json :: linear probe jc2_188 N=1000000: (1 - macro_auc_resonant) of FLAV_F1 over R16_Q1, run by run :: sha256 c6c4138ee5c86423 :: prose twin of the run farther from 1
\newcommand{\MlstCmpJcTwoResFonerOverCoarse}{0.764 and 0.740}  % experiments/FIGS/data/v2/index.json :: linear probe jc2_188 N=1000000: (1 - macro_auc_resonant) of FLAV_F1R over R16_Q1, run by run :: sha256 c6c4138ee5c86423
\newcommand{\MlstCmpJcTwoResFonerOverCoarsePr}{at least 24\% fewer errors}  % experiments/FIGS/data/v2/index.json :: linear probe jc2_188 N=1000000: (1 - macro_auc_resonant) of FLAV_F1R over R16_Q1, run by run :: sha256 c6c4138ee5c86423 :: prose twin: two significant figures, direction in words
\newcommand{\MlstCmpJcTwoResFonerOverCoarsePrMax}{up to 26\% fewer errors}  % experiments/FIGS/data/v2/index.json :: linear probe jc2_188 N=1000000: (1 - macro_auc_resonant) of FLAV_F1R over R16_Q1, run by run :: sha256 c6c4138ee5c86423 :: prose twin of the run farther from 1
\newcommand{\MlstCmpJcTwoResFoneOverFzero}{0.859 and 0.831}  % experiments/FIGS/data/v2/index.json :: linear probe jc2_188 N=1000000: (1 - macro_auc_resonant) of FLAV_F1 over FLAV_F0, run by run :: sha256 c6c4138ee5c86423
\newcommand{\MlstCmpJcTwoResFoneOverFzeroPr}{at least 14\% fewer errors}  % experiments/FIGS/data/v2/index.json :: linear probe jc2_188 N=1000000: (1 - macro_auc_resonant) of FLAV_F1 over FLAV_F0, run by run :: sha256 c6c4138ee5c86423 :: prose twin: two significant figures, direction in words
\newcommand{\MlstCmpJcTwoResFoneOverFzeroPrMax}{up to 17\% fewer errors}  % experiments/FIGS/data/v2/index.json :: linear probe jc2_188 N=1000000: (1 - macro_auc_resonant) of FLAV_F1 over FLAV_F0, run by run :: sha256 c6c4138ee5c86423 :: prose twin of the run farther from 1
\newcommand{\MlstCmpJcTwoResFonerOverFzero}{0.843 and 0.820}  % experiments/FIGS/data/v2/index.json :: linear probe jc2_188 N=1000000: (1 - macro_auc_resonant) of FLAV_F1R over FLAV_F0, run by run :: sha256 c6c4138ee5c86423
\newcommand{\MlstCmpJcTwoResFonerOverFzeroPr}{at least 16\% fewer errors}  % experiments/FIGS/data/v2/index.json :: linear probe jc2_188 N=1000000: (1 - macro_auc_resonant) of FLAV_F1R over FLAV_F0, run by run :: sha256 c6c4138ee5c86423 :: prose twin: two significant figures, direction in words
\newcommand{\MlstCmpJcTwoResFonerOverFzeroPrMax}{up to 18\% fewer errors}  % experiments/FIGS/data/v2/index.json :: linear probe jc2_188 N=1000000: (1 - macro_auc_resonant) of FLAV_F1R over FLAV_F0, run by run :: sha256 c6c4138ee5c86423 :: prose twin of the run farther from 1
\newcommand{\MlstCmpJcTwoResFonerOverFone}{0.9806 and 0.9867}  % experiments/FIGS/data/v2/index.json :: linear probe jc2_188 N=1000000: (1 - macro_auc_resonant) of FLAV_F1R over FLAV_F1, run by run :: sha256 c6c4138ee5c86423
\newcommand{\MlstCmpJcTwoResFonerOverFonePr}{equal within 2\%}  % experiments/FIGS/data/v2/index.json :: linear probe jc2_188 N=1000000: (1 - macro_auc_resonant) of FLAV_F1R over FLAV_F1, run by run :: sha256 c6c4138ee5c86423 :: prose twin: two significant figures, direction in words
\newcommand{\MlstCmpJcTwoResFonerOverFonePrMax}{equal within 2\%}  % experiments/FIGS/data/v2/index.json :: linear probe jc2_188 N=1000000: (1 - macro_auc_resonant) of FLAV_F1R over FLAV_F1, run by run :: sha256 c6c4138ee5c86423 :: prose twin of the run farther from 1
\newcommand{\MlstCmpJcTwoResMidOverRes}{1.258\,\ensuremath{\pm}\,0.017}  % experiments/FIGS/data/v2/index.json :: linear probe jc2_188 N=1000000: (1 - macro_auc_resonant) of R42_Q1 over L162, run by run :: sha256 c6c4138ee5c86423
\newcommand{\MlstCmpJcTwoResMidOverResPr}{26\% more errors}  % experiments/FIGS/data/v2/index.json :: linear probe jc2_188 N=1000000: (1 - macro_auc_resonant) of R42_Q1 over L162, run by run :: sha256 c6c4138ee5c86423 :: prose twin: two significant figures, direction in words
\newcommand{\MlstCmpJcTwoResCoarseOverRes}{1.58\,\ensuremath{\pm}\,0.04}  % experiments/FIGS/data/v2/index.json :: linear probe jc2_188 N=1000000: (1 - macro_auc_resonant) of R16_Q1 over L162, run by run :: sha256 c6c4138ee5c86423
\newcommand{\MlstCmpJcTwoResCoarseOverResPr}{58\% more errors}  % experiments/FIGS/data/v2/index.json :: linear probe jc2_188 N=1000000: (1 - macro_auc_resonant) of R16_Q1 over L162, run by run :: sha256 c6c4138ee5c86423 :: prose twin: two significant figures, direction in words
\newcommand{\MlstCmpJcTwoResCoarseOverMid}{1.260\,\ensuremath{\pm}\,0.015}  % experiments/FIGS/data/v2/index.json :: linear probe jc2_188 N=1000000: (1 - macro_auc_resonant) of R16_Q1 over R42_Q1, run by run :: sha256 c6c4138ee5c86423
\newcommand{\MlstCmpJcTwoResCoarseOverMidPr}{26\% more errors}  % experiments/FIGS/data/v2/index.json :: linear probe jc2_188 N=1000000: (1 - macro_auc_resonant) of R16_Q1 over R42_Q1, run by run :: sha256 c6c4138ee5c86423 :: prose twin: two significant figures, direction in words
\newcommand{\MlstCmpJcTwoResRandOverFull}{1.098\,\ensuremath{\pm}\,0.008}  % experiments/FIGS/data/v2/index.json :: linear probe jc2_188 N=1000000: (1 - macro_auc_resonant) of each random partition over the L188 run of its seed, mean +- SD over the partitions :: sha256 c6c4138ee5c86423
\newcommand{\MlstCmpJcTwoResRandOverFullPr}{9.8\% more errors}  % experiments/FIGS/data/v2/index.json :: linear probe jc2_188 N=1000000: (1 - macro_auc_resonant) of each random partition over the L188 run of its seed, mean +- SD over the partitions :: sha256 c6c4138ee5c86423 :: prose twin: two significant figures, direction in words
\newcommand{\MlstCmpJcTwoResRandOverMid}{0.870\,\ensuremath{\pm}\,0.008}  % experiments/FIGS/data/v2/index.json :: linear probe jc2_188 N=1000000: (1 - macro_auc_resonant) of each random partition over the R42_Q1 run of its seed, mean +- SD over the partitions :: sha256 c6c4138ee5c86423
\newcommand{\MlstCmpJcTwoResRandOverMidPr}{13\% fewer errors}  % experiments/FIGS/data/v2/index.json :: linear probe jc2_188 N=1000000: (1 - macro_auc_resonant) of each random partition over the R42_Q1 run of its seed, mean +- SD over the partitions :: sha256 c6c4138ee5c86423 :: prose twin: two significant figures, direction in words
\newcommand{\MlstCmpJcTwoResRandOverCoarse}{0.693\,\ensuremath{\pm}\,0.011}  % experiments/FIGS/data/v2/index.json :: linear probe jc2_188 N=1000000: (1 - macro_auc_resonant) of each random partition over the R16_Q1 run of its seed, mean +- SD over the partitions :: sha256 c6c4138ee5c86423
\newcommand{\MlstCmpJcTwoResRandOverCoarsePr}{31\% fewer errors}  % experiments/FIGS/data/v2/index.json :: linear probe jc2_188 N=1000000: (1 - macro_auc_resonant) of each random partition over the R16_Q1 run of its seed, mean +- SD over the partitions :: sha256 c6c4138ee5c86423 :: prose twin: two significant figures, direction in words
\newcommand{\MlstCmpJcTwoQcdResOverFull}{1.087\,\ensuremath{\pm}\,0.006}  % experiments/FIGS/data/v2/index.json :: linear probe jc2_188 N=1000000: (1 - macro_auc_qcd) of L162 over L188, run by run :: sha256 c6c4138ee5c86423
\newcommand{\MlstCmpJcTwoQcdResOverFullPr}{8.7\% more errors}  % experiments/FIGS/data/v2/index.json :: linear probe jc2_188 N=1000000: (1 - macro_auc_qcd) of L162 over L188, run by run :: sha256 c6c4138ee5c86423 :: prose twin: two significant figures, direction in words
\newcommand{\MlstCmpJcTwoQcdMidOverFull}{1.131\,\ensuremath{\pm}\,0.007}  % experiments/FIGS/data/v2/index.json :: linear probe jc2_188 N=1000000: (1 - macro_auc_qcd) of R42_Q1 over L188, run by run :: sha256 c6c4138ee5c86423
\newcommand{\MlstCmpJcTwoQcdMidOverFullPr}{13\% more errors}  % experiments/FIGS/data/v2/index.json :: linear probe jc2_188 N=1000000: (1 - macro_auc_qcd) of R42_Q1 over L188, run by run :: sha256 c6c4138ee5c86423 :: prose twin: two significant figures, direction in words
\newcommand{\MlstCmpJcTwoQcdCoarseOverFull}{1.466\,\ensuremath{\pm}\,0.017}  % experiments/FIGS/data/v2/index.json :: linear probe jc2_188 N=1000000: (1 - macro_auc_qcd) of R16_Q1 over L188, run by run :: sha256 c6c4138ee5c86423
\newcommand{\MlstCmpJcTwoQcdCoarseOverFullPr}{47\% more errors}  % experiments/FIGS/data/v2/index.json :: linear probe jc2_188 N=1000000: (1 - macro_auc_qcd) of R16_Q1 over L188, run by run :: sha256 c6c4138ee5c86423 :: prose twin: two significant figures, direction in words
\newcommand{\MlstCmpJcTwoQcdResMassOverRes}{1.012 and 1.006}  % experiments/FIGS/data/v2/index.json :: linear probe jc2_188 N=1000000: (1 - macro_auc_qcd) of L162_MASS over L162, run by run :: sha256 c6c4138ee5c86423
\newcommand{\MlstCmpJcTwoQcdResMassOverResPr}{equal within 2\%}  % experiments/FIGS/data/v2/index.json :: linear probe jc2_188 N=1000000: (1 - macro_auc_qcd) of L162_MASS over L162, run by run :: sha256 c6c4138ee5c86423 :: prose twin: two significant figures, direction in words
\newcommand{\MlstCmpJcTwoQcdResMassOverResPrMax}{equal within 2\%}  % experiments/FIGS/data/v2/index.json :: linear probe jc2_188 N=1000000: (1 - macro_auc_qcd) of L162_MASS over L162, run by run :: sha256 c6c4138ee5c86423 :: prose twin of the run farther from 1
\newcommand{\MlstCmpJcTwoQcdCoarseMassOverCoarse}{0.991 and 1.014}  % experiments/FIGS/data/v2/index.json :: linear probe jc2_188 N=1000000: (1 - macro_auc_qcd) of R16_Q1_MASS over R16_Q1, run by run :: sha256 c6c4138ee5c86423
\newcommand{\MlstCmpJcTwoQcdCoarseMassOverCoarsePr}{equal within 2\%}  % experiments/FIGS/data/v2/index.json :: linear probe jc2_188 N=1000000: (1 - macro_auc_qcd) of R16_Q1_MASS over R16_Q1, run by run :: sha256 c6c4138ee5c86423 :: prose twin: two significant figures, direction in words
\newcommand{\MlstCmpJcTwoQcdCoarseMassOverCoarsePrMax}{equal within 2\%}  % experiments/FIGS/data/v2/index.json :: linear probe jc2_188 N=1000000: (1 - macro_auc_qcd) of R16_Q1_MASS over R16_Q1, run by run :: sha256 c6c4138ee5c86423 :: prose twin of the run farther from 1
\newcommand{\MlstCmpJcTwoQcdResMassOverFull}{1.107 and 1.092}  % experiments/FIGS/data/v2/index.json :: linear probe jc2_188 N=1000000: (1 - macro_auc_qcd) of L162_MASS over L188, run by run :: sha256 c6c4138ee5c86423
\newcommand{\MlstCmpJcTwoQcdResMassOverFullPr}{at least 9.2\% more errors}  % experiments/FIGS/data/v2/index.json :: linear probe jc2_188 N=1000000: (1 - macro_auc_qcd) of L162_MASS over L188, run by run :: sha256 c6c4138ee5c86423 :: prose twin: two significant figures, direction in words
\newcommand{\MlstCmpJcTwoQcdResMassOverFullPrMax}{up to 11\% more errors}  % experiments/FIGS/data/v2/index.json :: linear probe jc2_188 N=1000000: (1 - macro_auc_qcd) of L162_MASS over L188, run by run :: sha256 c6c4138ee5c86423 :: prose twin of the run farther from 1
\newcommand{\MlstCmpJcTwoQcdCoarseMassOverFull}{1.472 and 1.474}  % experiments/FIGS/data/v2/index.json :: linear probe jc2_188 N=1000000: (1 - macro_auc_qcd) of R16_Q1_MASS over L188, run by run :: sha256 c6c4138ee5c86423
\newcommand{\MlstCmpJcTwoQcdCoarseMassOverFullPr}{at least 47\% more errors}  % experiments/FIGS/data/v2/index.json :: linear probe jc2_188 N=1000000: (1 - macro_auc_qcd) of R16_Q1_MASS over L188, run by run :: sha256 c6c4138ee5c86423 :: prose twin: two significant figures, direction in words
\newcommand{\MlstCmpJcTwoQcdCoarseMassOverFullPrMax}{up to 47\% more errors}  % experiments/FIGS/data/v2/index.json :: linear probe jc2_188 N=1000000: (1 - macro_auc_qcd) of R16_Q1_MASS over L188, run by run :: sha256 c6c4138ee5c86423 :: prose twin of the run farther from 1
\newcommand{\MlstCmpJcTwoQcdFzeroOverFull}{1.466 and 1.497}  % experiments/FIGS/data/v2/index.json :: linear probe jc2_188 N=1000000: (1 - macro_auc_qcd) of FLAV_F0 over L188, run by run :: sha256 c6c4138ee5c86423
\newcommand{\MlstCmpJcTwoQcdFzeroOverFullPr}{at least 47\% more errors}  % experiments/FIGS/data/v2/index.json :: linear probe jc2_188 N=1000000: (1 - macro_auc_qcd) of FLAV_F0 over L188, run by run :: sha256 c6c4138ee5c86423 :: prose twin: two significant figures, direction in words
\newcommand{\MlstCmpJcTwoQcdFzeroOverFullPrMax}{up to 50\% more errors}  % experiments/FIGS/data/v2/index.json :: linear probe jc2_188 N=1000000: (1 - macro_auc_qcd) of FLAV_F0 over L188, run by run :: sha256 c6c4138ee5c86423 :: prose twin of the run farther from 1
\newcommand{\MlstCmpJcTwoQcdFoneOverFull}{1.241 and 1.235}  % experiments/FIGS/data/v2/index.json :: linear probe jc2_188 N=1000000: (1 - macro_auc_qcd) of FLAV_F1 over L188, run by run :: sha256 c6c4138ee5c86423
\newcommand{\MlstCmpJcTwoQcdFoneOverFullPr}{at least 23\% more errors}  % experiments/FIGS/data/v2/index.json :: linear probe jc2_188 N=1000000: (1 - macro_auc_qcd) of FLAV_F1 over L188, run by run :: sha256 c6c4138ee5c86423 :: prose twin: two significant figures, direction in words
\newcommand{\MlstCmpJcTwoQcdFoneOverFullPrMax}{up to 24\% more errors}  % experiments/FIGS/data/v2/index.json :: linear probe jc2_188 N=1000000: (1 - macro_auc_qcd) of FLAV_F1 over L188, run by run :: sha256 c6c4138ee5c86423 :: prose twin of the run farther from 1
\newcommand{\MlstCmpJcTwoQcdFonerOverFull}{1.260 and 1.263}  % experiments/FIGS/data/v2/index.json :: linear probe jc2_188 N=1000000: (1 - macro_auc_qcd) of FLAV_F1R over L188, run by run :: sha256 c6c4138ee5c86423
\newcommand{\MlstCmpJcTwoQcdFonerOverFullPr}{at least 26\% more errors}  % experiments/FIGS/data/v2/index.json :: linear probe jc2_188 N=1000000: (1 - macro_auc_qcd) of FLAV_F1R over L188, run by run :: sha256 c6c4138ee5c86423 :: prose twin: two significant figures, direction in words
\newcommand{\MlstCmpJcTwoQcdFonerOverFullPrMax}{up to 26\% more errors}  % experiments/FIGS/data/v2/index.json :: linear probe jc2_188 N=1000000: (1 - macro_auc_qcd) of FLAV_F1R over L188, run by run :: sha256 c6c4138ee5c86423 :: prose twin of the run farther from 1
\newcommand{\MlstCmpJcTwoQcdFzeroOverCoarse}{0.987 and 1.031}  % experiments/FIGS/data/v2/index.json :: linear probe jc2_188 N=1000000: (1 - macro_auc_qcd) of FLAV_F0 over R16_Q1, run by run :: sha256 c6c4138ee5c86423
\newcommand{\MlstCmpJcTwoQcdFzeroOverCoarsePr}{equal within 4\%}  % experiments/FIGS/data/v2/index.json :: linear probe jc2_188 N=1000000: (1 - macro_auc_qcd) of FLAV_F0 over R16_Q1, run by run :: sha256 c6c4138ee5c86423 :: prose twin: two significant figures, direction in words
\newcommand{\MlstCmpJcTwoQcdFzeroOverCoarsePrMax}{equal within 4\%}  % experiments/FIGS/data/v2/index.json :: linear probe jc2_188 N=1000000: (1 - macro_auc_qcd) of FLAV_F0 over R16_Q1, run by run :: sha256 c6c4138ee5c86423 :: prose twin of the run farther from 1
\newcommand{\MlstCmpJcTwoQcdFoneOverCoarse}{0.836 and 0.850}  % experiments/FIGS/data/v2/index.json :: linear probe jc2_188 N=1000000: (1 - macro_auc_qcd) of FLAV_F1 over R16_Q1, run by run :: sha256 c6c4138ee5c86423
\newcommand{\MlstCmpJcTwoQcdFoneOverCoarsePr}{at least 15\% fewer errors}  % experiments/FIGS/data/v2/index.json :: linear probe jc2_188 N=1000000: (1 - macro_auc_qcd) of FLAV_F1 over R16_Q1, run by run :: sha256 c6c4138ee5c86423 :: prose twin: two significant figures, direction in words
\newcommand{\MlstCmpJcTwoQcdFoneOverCoarsePrMax}{up to 16\% fewer errors}  % experiments/FIGS/data/v2/index.json :: linear probe jc2_188 N=1000000: (1 - macro_auc_qcd) of FLAV_F1 over R16_Q1, run by run :: sha256 c6c4138ee5c86423 :: prose twin of the run farther from 1
\newcommand{\MlstCmpJcTwoQcdFonerOverCoarse}{0.848 and 0.870}  % experiments/FIGS/data/v2/index.json :: linear probe jc2_188 N=1000000: (1 - macro_auc_qcd) of FLAV_F1R over R16_Q1, run by run :: sha256 c6c4138ee5c86423
\newcommand{\MlstCmpJcTwoQcdFonerOverCoarsePr}{at least 13\% fewer errors}  % experiments/FIGS/data/v2/index.json :: linear probe jc2_188 N=1000000: (1 - macro_auc_qcd) of FLAV_F1R over R16_Q1, run by run :: sha256 c6c4138ee5c86423 :: prose twin: two significant figures, direction in words
\newcommand{\MlstCmpJcTwoQcdFonerOverCoarsePrMax}{up to 15\% fewer errors}  % experiments/FIGS/data/v2/index.json :: linear probe jc2_188 N=1000000: (1 - macro_auc_qcd) of FLAV_F1R over R16_Q1, run by run :: sha256 c6c4138ee5c86423 :: prose twin of the run farther from 1
\newcommand{\MlstCmpJcTwoQcdFoneOverFzero}{0.847 and 0.825}  % experiments/FIGS/data/v2/index.json :: linear probe jc2_188 N=1000000: (1 - macro_auc_qcd) of FLAV_F1 over FLAV_F0, run by run :: sha256 c6c4138ee5c86423
\newcommand{\MlstCmpJcTwoQcdFoneOverFzeroPr}{at least 15\% fewer errors}  % experiments/FIGS/data/v2/index.json :: linear probe jc2_188 N=1000000: (1 - macro_auc_qcd) of FLAV_F1 over FLAV_F0, run by run :: sha256 c6c4138ee5c86423 :: prose twin: two significant figures, direction in words
\newcommand{\MlstCmpJcTwoQcdFoneOverFzeroPrMax}{up to 18\% fewer errors}  % experiments/FIGS/data/v2/index.json :: linear probe jc2_188 N=1000000: (1 - macro_auc_qcd) of FLAV_F1 over FLAV_F0, run by run :: sha256 c6c4138ee5c86423 :: prose twin of the run farther from 1
\newcommand{\MlstCmpJcTwoQcdFonerOverFzero}{0.859 and 0.844}  % experiments/FIGS/data/v2/index.json :: linear probe jc2_188 N=1000000: (1 - macro_auc_qcd) of FLAV_F1R over FLAV_F0, run by run :: sha256 c6c4138ee5c86423
\newcommand{\MlstCmpJcTwoQcdFonerOverFzeroPr}{at least 14\% fewer errors}  % experiments/FIGS/data/v2/index.json :: linear probe jc2_188 N=1000000: (1 - macro_auc_qcd) of FLAV_F1R over FLAV_F0, run by run :: sha256 c6c4138ee5c86423 :: prose twin: two significant figures, direction in words
\newcommand{\MlstCmpJcTwoQcdFonerOverFzeroPrMax}{up to 16\% fewer errors}  % experiments/FIGS/data/v2/index.json :: linear probe jc2_188 N=1000000: (1 - macro_auc_qcd) of FLAV_F1R over FLAV_F0, run by run :: sha256 c6c4138ee5c86423 :: prose twin of the run farther from 1
\newcommand{\MlstCmpJcTwoQcdFonerOverFone}{1.015 and 1.023}  % experiments/FIGS/data/v2/index.json :: linear probe jc2_188 N=1000000: (1 - macro_auc_qcd) of FLAV_F1R over FLAV_F1, run by run :: sha256 c6c4138ee5c86423
\newcommand{\MlstCmpJcTwoQcdFonerOverFonePr}{equal within 3\%}  % experiments/FIGS/data/v2/index.json :: linear probe jc2_188 N=1000000: (1 - macro_auc_qcd) of FLAV_F1R over FLAV_F1, run by run :: sha256 c6c4138ee5c86423 :: prose twin: two significant figures, direction in words
\newcommand{\MlstCmpJcTwoQcdFonerOverFonePrMax}{equal within 3\%}  % experiments/FIGS/data/v2/index.json :: linear probe jc2_188 N=1000000: (1 - macro_auc_qcd) of FLAV_F1R over FLAV_F1, run by run :: sha256 c6c4138ee5c86423 :: prose twin of the run farther from 1
\newcommand{\MlstCmpJcTwoQcdMidOverRes}{1.040\,\ensuremath{\pm}\,0.011}  % experiments/FIGS/data/v2/index.json :: linear probe jc2_188 N=1000000: (1 - macro_auc_qcd) of R42_Q1 over L162, run by run :: sha256 c6c4138ee5c86423
\newcommand{\MlstCmpJcTwoQcdMidOverResPr}{4.0\% more errors}  % experiments/FIGS/data/v2/index.json :: linear probe jc2_188 N=1000000: (1 - macro_auc_qcd) of R42_Q1 over L162, run by run :: sha256 c6c4138ee5c86423 :: prose twin: two significant figures, direction in words
\newcommand{\MlstCmpJcTwoQcdCoarseOverRes}{1.348\,\ensuremath{\pm}\,0.010}  % experiments/FIGS/data/v2/index.json :: linear probe jc2_188 N=1000000: (1 - macro_auc_qcd) of R16_Q1 over L162, run by run :: sha256 c6c4138ee5c86423
\newcommand{\MlstCmpJcTwoQcdCoarseOverResPr}{35\% more errors}  % experiments/FIGS/data/v2/index.json :: linear probe jc2_188 N=1000000: (1 - macro_auc_qcd) of R16_Q1 over L162, run by run :: sha256 c6c4138ee5c86423 :: prose twin: two significant figures, direction in words
\newcommand{\MlstCmpJcTwoQcdCoarseOverMid}{1.296\,\ensuremath{\pm}\,0.016}  % experiments/FIGS/data/v2/index.json :: linear probe jc2_188 N=1000000: (1 - macro_auc_qcd) of R16_Q1 over R42_Q1, run by run :: sha256 c6c4138ee5c86423
\newcommand{\MlstCmpJcTwoQcdCoarseOverMidPr}{30\% more errors}  % experiments/FIGS/data/v2/index.json :: linear probe jc2_188 N=1000000: (1 - macro_auc_qcd) of R16_Q1 over R42_Q1, run by run :: sha256 c6c4138ee5c86423 :: prose twin: two significant figures, direction in words
\newcommand{\MlstCmpJcTwoQcdRandOverFull}{1.198\,\ensuremath{\pm}\,0.011}  % experiments/FIGS/data/v2/index.json :: linear probe jc2_188 N=1000000: (1 - macro_auc_qcd) of each random partition over the L188 run of its seed, mean +- SD over the partitions :: sha256 c6c4138ee5c86423
\newcommand{\MlstCmpJcTwoQcdRandOverFullPr}{20\% more errors}  % experiments/FIGS/data/v2/index.json :: linear probe jc2_188 N=1000000: (1 - macro_auc_qcd) of each random partition over the L188 run of its seed, mean +- SD over the partitions :: sha256 c6c4138ee5c86423 :: prose twin: two significant figures, direction in words
\newcommand{\MlstCmpJcTwoQcdRandOverMid}{1.064\,\ensuremath{\pm}\,0.011}  % experiments/FIGS/data/v2/index.json :: linear probe jc2_188 N=1000000: (1 - macro_auc_qcd) of each random partition over the R42_Q1 run of its seed, mean +- SD over the partitions :: sha256 c6c4138ee5c86423
\newcommand{\MlstCmpJcTwoQcdRandOverMidPr}{6.4\% more errors}  % experiments/FIGS/data/v2/index.json :: linear probe jc2_188 N=1000000: (1 - macro_auc_qcd) of each random partition over the R42_Q1 run of its seed, mean +- SD over the partitions :: sha256 c6c4138ee5c86423 :: prose twin: two significant figures, direction in words
\newcommand{\MlstCmpJcTwoQcdRandOverCoarse}{0.818\,\ensuremath{\pm}\,0.013}  % experiments/FIGS/data/v2/index.json :: linear probe jc2_188 N=1000000: (1 - macro_auc_qcd) of each random partition over the R16_Q1 run of its seed, mean +- SD over the partitions :: sha256 c6c4138ee5c86423
\newcommand{\MlstCmpJcTwoQcdRandOverCoarsePr}{18\% fewer errors}  % experiments/FIGS/data/v2/index.json :: linear probe jc2_188 N=1000000: (1 - macro_auc_qcd) of each random partition over the R16_Q1 run of its seed, mean +- SD over the partitions :: sha256 c6c4138ee5c86423 :: prose twin: two significant figures, direction in words
\newcommand{\MlstCmpBcResOverFull}{0.998\,\ensuremath{\pm}\,0.007}  % experiments/FIGS/data/v2/index.json :: binary probe bc_vs_bq_cs: (1 - AUC) of L162 over L188, run by run :: sha256 c6c4138ee5c86423
\newcommand{\MlstCmpBcResOverFullPr}{equal within 1\%}  % experiments/FIGS/data/v2/index.json :: binary probe bc_vs_bq_cs: (1 - AUC) of L162 over L188, run by run :: sha256 c6c4138ee5c86423 :: prose twin: two significant figures, direction in words
\newcommand{\MlstCmpBcMidOverFull}{1.151\,\ensuremath{\pm}\,0.029}  % experiments/FIGS/data/v2/index.json :: binary probe bc_vs_bq_cs: (1 - AUC) of R42_Q1 over L188, run by run :: sha256 c6c4138ee5c86423
\newcommand{\MlstCmpBcMidOverFullPr}{15\% more errors}  % experiments/FIGS/data/v2/index.json :: binary probe bc_vs_bq_cs: (1 - AUC) of R42_Q1 over L188, run by run :: sha256 c6c4138ee5c86423 :: prose twin: two significant figures, direction in words
\newcommand{\MlstCmpBcCoarseOverFull}{1.95\,\ensuremath{\pm}\,0.06}  % experiments/FIGS/data/v2/index.json :: binary probe bc_vs_bq_cs: (1 - AUC) of R16_Q1 over L188, run by run :: sha256 c6c4138ee5c86423
\newcommand{\MlstCmpBcCoarseOverFullPr}{95\% more errors}  % experiments/FIGS/data/v2/index.json :: binary probe bc_vs_bq_cs: (1 - AUC) of R16_Q1 over L188, run by run :: sha256 c6c4138ee5c86423 :: prose twin: two significant figures, direction in words
\newcommand{\MlstCmpBcFoneOverFzero}{0.759 and 0.698}  % experiments/FIGS/data/v2/index.json :: binary probe bc_vs_bq_cs: (1 - AUC) of FLAV_F1 over FLAV_F0, run by run :: sha256 c6c4138ee5c86423
\newcommand{\MlstCmpBcFoneOverFzeroPr}{at least 24\% fewer errors}  % experiments/FIGS/data/v2/index.json :: binary probe bc_vs_bq_cs: (1 - AUC) of FLAV_F1 over FLAV_F0, run by run :: sha256 c6c4138ee5c86423 :: prose twin: two significant figures, direction in words
\newcommand{\MlstCmpBcFoneOverFzeroPrMax}{up to 30\% fewer errors}  % experiments/FIGS/data/v2/index.json :: binary probe bc_vs_bq_cs: (1 - AUC) of FLAV_F1 over FLAV_F0, run by run :: sha256 c6c4138ee5c86423 :: prose twin of the run farther from 1
\newcommand{\MlstCmpBcFonerOverFone}{1.008 and 0.990}  % experiments/FIGS/data/v2/index.json :: binary probe bc_vs_bq_cs: (1 - AUC) of FLAV_F1R over FLAV_F1, run by run :: sha256 c6c4138ee5c86423
\newcommand{\MlstCmpBcFonerOverFonePr}{equal within 1\%}  % experiments/FIGS/data/v2/index.json :: binary probe bc_vs_bq_cs: (1 - AUC) of FLAV_F1R over FLAV_F1, run by run :: sha256 c6c4138ee5c86423 :: prose twin: two significant figures, direction in words
\newcommand{\MlstCmpBcFonerOverFonePrMax}{equal within 1\%}  % experiments/FIGS/data/v2/index.json :: binary probe bc_vs_bq_cs: (1 - AUC) of FLAV_F1R over FLAV_F1, run by run :: sha256 c6c4138ee5c86423 :: prose twin of the run farther from 1
\newcommand{\MlstCmpBcFonerOverFzero}{0.765 and 0.691}  % experiments/FIGS/data/v2/index.json :: binary probe bc_vs_bq_cs: (1 - AUC) of FLAV_F1R over FLAV_F0, run by run :: sha256 c6c4138ee5c86423
\newcommand{\MlstCmpBcFonerOverFzeroPr}{at least 23\% fewer errors}  % experiments/FIGS/data/v2/index.json :: binary probe bc_vs_bq_cs: (1 - AUC) of FLAV_F1R over FLAV_F0, run by run :: sha256 c6c4138ee5c86423 :: prose twin: two significant figures, direction in words
\newcommand{\MlstCmpBcFonerOverFzeroPrMax}{up to 31\% fewer errors}  % experiments/FIGS/data/v2/index.json :: binary probe bc_vs_bq_cs: (1 - AUC) of FLAV_F1R over FLAV_F0, run by run :: sha256 c6c4138ee5c86423 :: prose twin of the run farther from 1
\newcommand{\MlstCmpBcFzeroOverCoarse}{0.99 and 1.13}  % experiments/FIGS/data/v2/index.json :: binary probe bc_vs_bq_cs: (1 - AUC) of FLAV_F0 over R16_Q1, run by run :: sha256 c6c4138ee5c86423
\newcommand{\MlstCmpBcFzeroOverCoarsePr}{equal within 20\%}  % experiments/FIGS/data/v2/index.json :: binary probe bc_vs_bq_cs: (1 - AUC) of FLAV_F0 over R16_Q1, run by run :: sha256 c6c4138ee5c86423 :: prose twin: two significant figures, direction in words
\newcommand{\MlstCmpBcFzeroOverCoarsePrMax}{equal within 20\%}  % experiments/FIGS/data/v2/index.json :: binary probe bc_vs_bq_cs: (1 - AUC) of FLAV_F0 over R16_Q1, run by run :: sha256 c6c4138ee5c86423 :: prose twin of the run farther from 1
\newcommand{\MlstCmpBcFoneOverCoarse}{0.752 and 0.786}  % experiments/FIGS/data/v2/index.json :: binary probe bc_vs_bq_cs: (1 - AUC) of FLAV_F1 over R16_Q1, run by run :: sha256 c6c4138ee5c86423
\newcommand{\MlstCmpBcFoneOverCoarsePr}{at least 21\% fewer errors}  % experiments/FIGS/data/v2/index.json :: binary probe bc_vs_bq_cs: (1 - AUC) of FLAV_F1 over R16_Q1, run by run :: sha256 c6c4138ee5c86423 :: prose twin: two significant figures, direction in words
\newcommand{\MlstCmpBcFoneOverCoarsePrMax}{up to 25\% fewer errors}  % experiments/FIGS/data/v2/index.json :: binary probe bc_vs_bq_cs: (1 - AUC) of FLAV_F1 over R16_Q1, run by run :: sha256 c6c4138ee5c86423 :: prose twin of the run farther from 1
\newcommand{\MlstCmpBcFonerOverCoarse}{0.758 and 0.778}  % experiments/FIGS/data/v2/index.json :: binary probe bc_vs_bq_cs: (1 - AUC) of FLAV_F1R over R16_Q1, run by run :: sha256 c6c4138ee5c86423
\newcommand{\MlstCmpBcFonerOverCoarsePr}{at least 22\% fewer errors}  % experiments/FIGS/data/v2/index.json :: binary probe bc_vs_bq_cs: (1 - AUC) of FLAV_F1R over R16_Q1, run by run :: sha256 c6c4138ee5c86423 :: prose twin: two significant figures, direction in words
\newcommand{\MlstCmpBcFonerOverCoarsePrMax}{up to 24\% fewer errors}  % experiments/FIGS/data/v2/index.json :: binary probe bc_vs_bq_cs: (1 - AUC) of FLAV_F1R over R16_Q1, run by run :: sha256 c6c4138ee5c86423 :: prose twin of the run farther from 1
\newcommand{\MlstCmpBcFzeroOverFull}{2.00 and 2.19}  % experiments/FIGS/data/v2/index.json :: binary probe bc_vs_bq_cs: (1 - AUC) of FLAV_F0 over L188, run by run :: sha256 c6c4138ee5c86423
\newcommand{\MlstCmpBcFzeroOverFullPr}{at least 2.0 times the errors}  % experiments/FIGS/data/v2/index.json :: binary probe bc_vs_bq_cs: (1 - AUC) of FLAV_F0 over L188, run by run :: sha256 c6c4138ee5c86423 :: prose twin: two significant figures, direction in words
\newcommand{\MlstCmpBcFzeroOverFullPrMax}{up to 2.2 times the errors}  % experiments/FIGS/data/v2/index.json :: binary probe bc_vs_bq_cs: (1 - AUC) of FLAV_F0 over L188, run by run :: sha256 c6c4138ee5c86423 :: prose twin of the run farther from 1
\newcommand{\MlstCmpBcFoneOverFull}{1.518 and 1.530}  % experiments/FIGS/data/v2/index.json :: binary probe bc_vs_bq_cs: (1 - AUC) of FLAV_F1 over L188, run by run :: sha256 c6c4138ee5c86423
\newcommand{\MlstCmpBcFoneOverFullPr}{at least 52\% more errors}  % experiments/FIGS/data/v2/index.json :: binary probe bc_vs_bq_cs: (1 - AUC) of FLAV_F1 over L188, run by run :: sha256 c6c4138ee5c86423 :: prose twin: two significant figures, direction in words
\newcommand{\MlstCmpBcFoneOverFullPrMax}{up to 53\% more errors}  % experiments/FIGS/data/v2/index.json :: binary probe bc_vs_bq_cs: (1 - AUC) of FLAV_F1 over L188, run by run :: sha256 c6c4138ee5c86423 :: prose twin of the run farther from 1
\newcommand{\MlstCmpBcFonerOverFull}{1.530 and 1.515}  % experiments/FIGS/data/v2/index.json :: binary probe bc_vs_bq_cs: (1 - AUC) of FLAV_F1R over L188, run by run :: sha256 c6c4138ee5c86423
\newcommand{\MlstCmpBcFonerOverFullPr}{at least 51\% more errors}  % experiments/FIGS/data/v2/index.json :: binary probe bc_vs_bq_cs: (1 - AUC) of FLAV_F1R over L188, run by run :: sha256 c6c4138ee5c86423 :: prose twin: two significant figures, direction in words
\newcommand{\MlstCmpBcFonerOverFullPrMax}{up to 53\% more errors}  % experiments/FIGS/data/v2/index.json :: binary probe bc_vs_bq_cs: (1 - AUC) of FLAV_F1R over L188, run by run :: sha256 c6c4138ee5c86423 :: prose twin of the run farther from 1
\newcommand{\MlstCmpBcRandOverFull}{1.15\,\ensuremath{\pm}\,0.07}  % experiments/FIGS/data/v2/index.json :: binary probe bc_vs_bq_cs: (1 - AUC) of each random partition over the 188-class run of its seed :: sha256 c6c4138ee5c86423
\newcommand{\MlstCmpBcRandOverFullPr}{15\% more errors}  % experiments/FIGS/data/v2/index.json :: binary probe bc_vs_bq_cs: (1 - AUC) of each random partition over the 188-class run of its seed :: sha256 c6c4138ee5c86423 :: prose twin: two significant figures, direction in words
\newcommand{\MlstPairRandMergedBcN}{2}  % experiments/FIGS/data/v2/index.json :: random partitions that merge bc_vs_bq_cs (count) :: sha256 c6c4138ee5c86423
\newcommand{\MlstPairRandApartBcN}{2}  % experiments/FIGS/data/v2/index.json :: random partitions that keep apart bc_vs_bq_cs (count) :: sha256 c6c4138ee5c86423
\newcommand{\MlstPairRandMergedOverApartBc}{1.12}  % experiments/FIGS/data/v2/index.json :: binary probe bc_vs_bq_cs: mean (1 - AUC) of the partitions that merge it over that of those that keep it apart (one run each) :: sha256 c6c4138ee5c86423
\newcommand{\MlstPairRandMergedOverApartBcPr}{12\% more errors}  % experiments/FIGS/data/v2/index.json :: binary probe bc_vs_bq_cs: mean (1 - AUC) of the partitions that merge it over that of those that keep it apart (one run each) :: sha256 c6c4138ee5c86423 :: prose twin: two significant figures, direction in words
\newcommand{\MlstCmpBbResOverFull}{1.015\,\ensuremath{\pm}\,0.027}  % experiments/FIGS/data/v2/index.json :: binary probe bb_vs_cc: (1 - AUC) of L162 over L188, run by run :: sha256 c6c4138ee5c86423
\newcommand{\MlstCmpBbResOverFullPr}{equal within 3\%}  % experiments/FIGS/data/v2/index.json :: binary probe bb_vs_cc: (1 - AUC) of L162 over L188, run by run :: sha256 c6c4138ee5c86423 :: prose twin: two significant figures, direction in words
\newcommand{\MlstCmpBbMidOverFull}{1.102\,\ensuremath{\pm}\,0.013}  % experiments/FIGS/data/v2/index.json :: binary probe bb_vs_cc: (1 - AUC) of R42_Q1 over L188, run by run :: sha256 c6c4138ee5c86423
\newcommand{\MlstCmpBbMidOverFullPr}{10\% more errors}  % experiments/FIGS/data/v2/index.json :: binary probe bb_vs_cc: (1 - AUC) of R42_Q1 over L188, run by run :: sha256 c6c4138ee5c86423 :: prose twin: two significant figures, direction in words
\newcommand{\MlstCmpBbCoarseOverFull}{4.22\,\ensuremath{\pm}\,0.29}  % experiments/FIGS/data/v2/index.json :: binary probe bb_vs_cc: (1 - AUC) of R16_Q1 over L188, run by run :: sha256 c6c4138ee5c86423
\newcommand{\MlstCmpBbCoarseOverFullPr}{4.2 times the errors}  % experiments/FIGS/data/v2/index.json :: binary probe bb_vs_cc: (1 - AUC) of R16_Q1 over L188, run by run :: sha256 c6c4138ee5c86423 :: prose twin: two significant figures, direction in words
\newcommand{\MlstCmpBbFoneOverFzero}{0.31 and 0.13}  % experiments/FIGS/data/v2/index.json :: binary probe bb_vs_cc: (1 - AUC) of FLAV_F1 over FLAV_F0, run by run :: sha256 c6c4138ee5c86423
\newcommand{\MlstCmpBbFoneOverFzeroPr}{at most 0.31 times the errors}  % experiments/FIGS/data/v2/index.json :: binary probe bb_vs_cc: (1 - AUC) of FLAV_F1 over FLAV_F0, run by run :: sha256 c6c4138ee5c86423 :: prose twin: two significant figures, direction in words
\newcommand{\MlstCmpBbFoneOverFzeroPrMax}{down to 0.13 times the errors}  % experiments/FIGS/data/v2/index.json :: binary probe bb_vs_cc: (1 - AUC) of FLAV_F1 over FLAV_F0, run by run :: sha256 c6c4138ee5c86423 :: prose twin of the run farther from 1
\newcommand{\MlstCmpBbFonerOverFone}{1.40 and 1.86}  % experiments/FIGS/data/v2/index.json :: binary probe bb_vs_cc: (1 - AUC) of FLAV_F1R over FLAV_F1, run by run :: sha256 c6c4138ee5c86423
\newcommand{\MlstCmpBbFonerOverFonePr}{at least 40\% more errors}  % experiments/FIGS/data/v2/index.json :: binary probe bb_vs_cc: (1 - AUC) of FLAV_F1R over FLAV_F1, run by run :: sha256 c6c4138ee5c86423 :: prose twin: two significant figures, direction in words
\newcommand{\MlstCmpBbFonerOverFonePrMax}{up to 86\% more errors}  % experiments/FIGS/data/v2/index.json :: binary probe bb_vs_cc: (1 - AUC) of FLAV_F1R over FLAV_F1, run by run :: sha256 c6c4138ee5c86423 :: prose twin of the run farther from 1
\newcommand{\MlstCmpBbFonerOverFzero}{0.44 and 0.24}  % experiments/FIGS/data/v2/index.json :: binary probe bb_vs_cc: (1 - AUC) of FLAV_F1R over FLAV_F0, run by run :: sha256 c6c4138ee5c86423
\newcommand{\MlstCmpBbFonerOverFzeroPr}{at most 0.44 times the errors}  % experiments/FIGS/data/v2/index.json :: binary probe bb_vs_cc: (1 - AUC) of FLAV_F1R over FLAV_F0, run by run :: sha256 c6c4138ee5c86423 :: prose twin: two significant figures, direction in words
\newcommand{\MlstCmpBbFonerOverFzeroPrMax}{down to 0.24 times the errors}  % experiments/FIGS/data/v2/index.json :: binary probe bb_vs_cc: (1 - AUC) of FLAV_F1R over FLAV_F0, run by run :: sha256 c6c4138ee5c86423 :: prose twin of the run farther from 1
\newcommand{\MlstCmpBbFzeroOverCoarse}{1.3 and 3.2}  % experiments/FIGS/data/v2/index.json :: binary probe bb_vs_cc: (1 - AUC) of FLAV_F0 over R16_Q1, run by run :: sha256 c6c4138ee5c86423
\newcommand{\MlstCmpBbFzeroOverCoarsePr}{at least 30\% more errors}  % experiments/FIGS/data/v2/index.json :: binary probe bb_vs_cc: (1 - AUC) of FLAV_F0 over R16_Q1, run by run :: sha256 c6c4138ee5c86423 :: prose twin: two significant figures, direction in words
\newcommand{\MlstCmpBbFzeroOverCoarsePrMax}{up to 3.2 times the errors}  % experiments/FIGS/data/v2/index.json :: binary probe bb_vs_cc: (1 - AUC) of FLAV_F0 over R16_Q1, run by run :: sha256 c6c4138ee5c86423 :: prose twin of the run farther from 1
\newcommand{\MlstCmpBbFoneOverCoarse}{0.4077 and 0.4126}  % experiments/FIGS/data/v2/index.json :: binary probe bb_vs_cc: (1 - AUC) of FLAV_F1 over R16_Q1, run by run :: sha256 c6c4138ee5c86423
\newcommand{\MlstCmpBbFoneOverCoarsePr}{at most 0.41 times the errors}  % experiments/FIGS/data/v2/index.json :: binary probe bb_vs_cc: (1 - AUC) of FLAV_F1 over R16_Q1, run by run :: sha256 c6c4138ee5c86423 :: prose twin: two significant figures, direction in words
\newcommand{\MlstCmpBbFoneOverCoarsePrMax}{down to 0.41 times the errors}  % experiments/FIGS/data/v2/index.json :: binary probe bb_vs_cc: (1 - AUC) of FLAV_F1 over R16_Q1, run by run :: sha256 c6c4138ee5c86423 :: prose twin of the run farther from 1
\newcommand{\MlstCmpBbFonerOverCoarse}{0.57 and 0.77}  % experiments/FIGS/data/v2/index.json :: binary probe bb_vs_cc: (1 - AUC) of FLAV_F1R over R16_Q1, run by run :: sha256 c6c4138ee5c86423
\newcommand{\MlstCmpBbFonerOverCoarsePr}{at least 23\% fewer errors}  % experiments/FIGS/data/v2/index.json :: binary probe bb_vs_cc: (1 - AUC) of FLAV_F1R over R16_Q1, run by run :: sha256 c6c4138ee5c86423 :: prose twin: two significant figures, direction in words
\newcommand{\MlstCmpBbFonerOverCoarsePrMax}{up to 43\% fewer errors}  % experiments/FIGS/data/v2/index.json :: binary probe bb_vs_cc: (1 - AUC) of FLAV_F1R over R16_Q1, run by run :: sha256 c6c4138ee5c86423 :: prose twin of the run farther from 1
\newcommand{\MlstCmpBbFzeroOverFull}{5.9 and 13.4}  % experiments/FIGS/data/v2/index.json :: binary probe bb_vs_cc: (1 - AUC) of FLAV_F0 over L188, run by run :: sha256 c6c4138ee5c86423
\newcommand{\MlstCmpBbFzeroOverFullPr}{at least 5.9 times the errors}  % experiments/FIGS/data/v2/index.json :: binary probe bb_vs_cc: (1 - AUC) of FLAV_F0 over L188, run by run :: sha256 c6c4138ee5c86423 :: prose twin: two significant figures, direction in words
\newcommand{\MlstCmpBbFzeroOverFullPrMax}{up to 13 times the errors}  % experiments/FIGS/data/v2/index.json :: binary probe bb_vs_cc: (1 - AUC) of FLAV_F0 over L188, run by run :: sha256 c6c4138ee5c86423 :: prose twin of the run farther from 1
\newcommand{\MlstCmpBbFoneOverFull}{1.837 and 1.748}  % experiments/FIGS/data/v2/index.json :: binary probe bb_vs_cc: (1 - AUC) of FLAV_F1 over L188, run by run :: sha256 c6c4138ee5c86423
\newcommand{\MlstCmpBbFoneOverFullPr}{at least 75\% more errors}  % experiments/FIGS/data/v2/index.json :: binary probe bb_vs_cc: (1 - AUC) of FLAV_F1 over L188, run by run :: sha256 c6c4138ee5c86423 :: prose twin: two significant figures, direction in words
\newcommand{\MlstCmpBbFoneOverFullPrMax}{up to 84\% more errors}  % experiments/FIGS/data/v2/index.json :: binary probe bb_vs_cc: (1 - AUC) of FLAV_F1 over L188, run by run :: sha256 c6c4138ee5c86423 :: prose twin of the run farther from 1
\newcommand{\MlstCmpBbFonerOverFull}{2.58 and 3.25}  % experiments/FIGS/data/v2/index.json :: binary probe bb_vs_cc: (1 - AUC) of FLAV_F1R over L188, run by run :: sha256 c6c4138ee5c86423
\newcommand{\MlstCmpBbFonerOverFullPr}{at least 2.6 times the errors}  % experiments/FIGS/data/v2/index.json :: binary probe bb_vs_cc: (1 - AUC) of FLAV_F1R over L188, run by run :: sha256 c6c4138ee5c86423 :: prose twin: two significant figures, direction in words
\newcommand{\MlstCmpBbFonerOverFullPrMax}{up to 3.2 times the errors}  % experiments/FIGS/data/v2/index.json :: binary probe bb_vs_cc: (1 - AUC) of FLAV_F1R over L188, run by run :: sha256 c6c4138ee5c86423 :: prose twin of the run farther from 1
\newcommand{\MlstCmpBbRandOverFull}{1.53\,\ensuremath{\pm}\,0.11}  % experiments/FIGS/data/v2/index.json :: binary probe bb_vs_cc: (1 - AUC) of each random partition over the 188-class run of its seed :: sha256 c6c4138ee5c86423
\newcommand{\MlstCmpBbRandOverFullPr}{53\% more errors}  % experiments/FIGS/data/v2/index.json :: binary probe bb_vs_cc: (1 - AUC) of each random partition over the 188-class run of its seed :: sha256 c6c4138ee5c86423 :: prose twin: two significant figures, direction in words
\newcommand{\MlstPairRandMergedBbN}{3}  % experiments/FIGS/data/v2/index.json :: random partitions that merge bb_vs_cc (count) :: sha256 c6c4138ee5c86423
\newcommand{\MlstPairRandApartBbN}{2}  % experiments/FIGS/data/v2/index.json :: random partitions that keep apart bb_vs_cc (count) :: sha256 c6c4138ee5c86423
\newcommand{\MlstPairRandMergedOverApartBb}{1.07}  % experiments/FIGS/data/v2/index.json :: binary probe bb_vs_cc: mean (1 - AUC) of the partitions that merge it over that of those that keep it apart (one run each) :: sha256 c6c4138ee5c86423
\newcommand{\MlstPairRandMergedOverApartBbPr}{7.3\% more errors}  % experiments/FIGS/data/v2/index.json :: binary probe bb_vs_cc: mean (1 - AUC) of the partitions that merge it over that of those that keep it apart (one run each) :: sha256 c6c4138ee5c86423 :: prose twin: two significant figures, direction in words
\newcommand{\MlstCmpBbqqResOverFull}{1.005\,\ensuremath{\pm}\,0.032}  % experiments/FIGS/data/v2/index.json :: binary probe bbqq_vs_ccqq: (1 - AUC) of L162 over L188, run by run :: sha256 c6c4138ee5c86423
\newcommand{\MlstCmpBbqqResOverFullPr}{equal within 4\%}  % experiments/FIGS/data/v2/index.json :: binary probe bbqq_vs_ccqq: (1 - AUC) of L162 over L188, run by run :: sha256 c6c4138ee5c86423 :: prose twin: two significant figures, direction in words
\newcommand{\MlstCmpBbqqMidOverFull}{1.117\,\ensuremath{\pm}\,0.008}  % experiments/FIGS/data/v2/index.json :: binary probe bbqq_vs_ccqq: (1 - AUC) of R42_Q1 over L188, run by run :: sha256 c6c4138ee5c86423
\newcommand{\MlstCmpBbqqMidOverFullPr}{12\% more errors}  % experiments/FIGS/data/v2/index.json :: binary probe bbqq_vs_ccqq: (1 - AUC) of R42_Q1 over L188, run by run :: sha256 c6c4138ee5c86423 :: prose twin: two significant figures, direction in words
\newcommand{\MlstCmpBbqqCoarseOverFull}{4.7\,\ensuremath{\pm}\,0.5}  % experiments/FIGS/data/v2/index.json :: binary probe bbqq_vs_ccqq: (1 - AUC) of R16_Q1 over L188, run by run :: sha256 c6c4138ee5c86423
\newcommand{\MlstCmpBbqqCoarseOverFullPr}{4.7 times the errors}  % experiments/FIGS/data/v2/index.json :: binary probe bbqq_vs_ccqq: (1 - AUC) of R16_Q1 over L188, run by run :: sha256 c6c4138ee5c86423 :: prose twin: two significant figures, direction in words
\newcommand{\MlstCmpBbqqFoneOverFzero}{0.236 and 0.154}  % experiments/FIGS/data/v2/index.json :: binary probe bbqq_vs_ccqq: (1 - AUC) of FLAV_F1 over FLAV_F0, run by run :: sha256 c6c4138ee5c86423
\newcommand{\MlstCmpBbqqFoneOverFzeroPr}{at most 0.24 times the errors}  % experiments/FIGS/data/v2/index.json :: binary probe bbqq_vs_ccqq: (1 - AUC) of FLAV_F1 over FLAV_F0, run by run :: sha256 c6c4138ee5c86423 :: prose twin: two significant figures, direction in words
\newcommand{\MlstCmpBbqqFoneOverFzeroPrMax}{down to 0.15 times the errors}  % experiments/FIGS/data/v2/index.json :: binary probe bbqq_vs_ccqq: (1 - AUC) of FLAV_F1 over FLAV_F0, run by run :: sha256 c6c4138ee5c86423 :: prose twin of the run farther from 1
\newcommand{\MlstCmpBbqqFonerOverFone}{1.505 and 1.577}  % experiments/FIGS/data/v2/index.json :: binary probe bbqq_vs_ccqq: (1 - AUC) of FLAV_F1R over FLAV_F1, run by run :: sha256 c6c4138ee5c86423
\newcommand{\MlstCmpBbqqFonerOverFonePr}{at least 50\% more errors}  % experiments/FIGS/data/v2/index.json :: binary probe bbqq_vs_ccqq: (1 - AUC) of FLAV_F1R over FLAV_F1, run by run :: sha256 c6c4138ee5c86423 :: prose twin: two significant figures, direction in words
\newcommand{\MlstCmpBbqqFonerOverFonePrMax}{up to 58\% more errors}  % experiments/FIGS/data/v2/index.json :: binary probe bbqq_vs_ccqq: (1 - AUC) of FLAV_F1R over FLAV_F1, run by run :: sha256 c6c4138ee5c86423 :: prose twin of the run farther from 1
\newcommand{\MlstCmpBbqqFonerOverFzero}{0.36 and 0.24}  % experiments/FIGS/data/v2/index.json :: binary probe bbqq_vs_ccqq: (1 - AUC) of FLAV_F1R over FLAV_F0, run by run :: sha256 c6c4138ee5c86423
\newcommand{\MlstCmpBbqqFonerOverFzeroPr}{at most 0.36 times the errors}  % experiments/FIGS/data/v2/index.json :: binary probe bbqq_vs_ccqq: (1 - AUC) of FLAV_F1R over FLAV_F0, run by run :: sha256 c6c4138ee5c86423 :: prose twin: two significant figures, direction in words
\newcommand{\MlstCmpBbqqFonerOverFzeroPrMax}{down to 0.24 times the errors}  % experiments/FIGS/data/v2/index.json :: binary probe bbqq_vs_ccqq: (1 - AUC) of FLAV_F1R over FLAV_F0, run by run :: sha256 c6c4138ee5c86423 :: prose twin of the run farther from 1
\newcommand{\MlstCmpBbqqFzeroOverCoarse}{1.51 and 1.79}  % experiments/FIGS/data/v2/index.json :: binary probe bbqq_vs_ccqq: (1 - AUC) of FLAV_F0 over R16_Q1, run by run :: sha256 c6c4138ee5c86423
\newcommand{\MlstCmpBbqqFzeroOverCoarsePr}{at least 51\% more errors}  % experiments/FIGS/data/v2/index.json :: binary probe bbqq_vs_ccqq: (1 - AUC) of FLAV_F0 over R16_Q1, run by run :: sha256 c6c4138ee5c86423 :: prose twin: two significant figures, direction in words
\newcommand{\MlstCmpBbqqFzeroOverCoarsePrMax}{up to 79\% more errors}  % experiments/FIGS/data/v2/index.json :: binary probe bbqq_vs_ccqq: (1 - AUC) of FLAV_F0 over R16_Q1, run by run :: sha256 c6c4138ee5c86423 :: prose twin of the run farther from 1
\newcommand{\MlstCmpBbqqFoneOverCoarse}{0.356 and 0.276}  % experiments/FIGS/data/v2/index.json :: binary probe bbqq_vs_ccqq: (1 - AUC) of FLAV_F1 over R16_Q1, run by run :: sha256 c6c4138ee5c86423
\newcommand{\MlstCmpBbqqFoneOverCoarsePr}{at most 0.36 times the errors}  % experiments/FIGS/data/v2/index.json :: binary probe bbqq_vs_ccqq: (1 - AUC) of FLAV_F1 over R16_Q1, run by run :: sha256 c6c4138ee5c86423 :: prose twin: two significant figures, direction in words
\newcommand{\MlstCmpBbqqFoneOverCoarsePrMax}{down to 0.28 times the errors}  % experiments/FIGS/data/v2/index.json :: binary probe bbqq_vs_ccqq: (1 - AUC) of FLAV_F1 over R16_Q1, run by run :: sha256 c6c4138ee5c86423 :: prose twin of the run farther from 1
\newcommand{\MlstCmpBbqqFonerOverCoarse}{0.54 and 0.44}  % experiments/FIGS/data/v2/index.json :: binary probe bbqq_vs_ccqq: (1 - AUC) of FLAV_F1R over R16_Q1, run by run :: sha256 c6c4138ee5c86423
\newcommand{\MlstCmpBbqqFonerOverCoarsePr}{at least 46\% fewer errors}  % experiments/FIGS/data/v2/index.json :: binary probe bbqq_vs_ccqq: (1 - AUC) of FLAV_F1R over R16_Q1, run by run :: sha256 c6c4138ee5c86423 :: prose twin: two significant figures, direction in words
\newcommand{\MlstCmpBbqqFonerOverCoarsePrMax}{down to 0.44 times the errors}  % experiments/FIGS/data/v2/index.json :: binary probe bbqq_vs_ccqq: (1 - AUC) of FLAV_F1R over R16_Q1, run by run :: sha256 c6c4138ee5c86423 :: prose twin of the run farther from 1
\newcommand{\MlstCmpBbqqFzeroOverFull}{6.4 and 9.4}  % experiments/FIGS/data/v2/index.json :: binary probe bbqq_vs_ccqq: (1 - AUC) of FLAV_F0 over L188, run by run :: sha256 c6c4138ee5c86423
\newcommand{\MlstCmpBbqqFzeroOverFullPr}{at least 6.4 times the errors}  % experiments/FIGS/data/v2/index.json :: binary probe bbqq_vs_ccqq: (1 - AUC) of FLAV_F0 over L188, run by run :: sha256 c6c4138ee5c86423 :: prose twin: two significant figures, direction in words
\newcommand{\MlstCmpBbqqFzeroOverFullPrMax}{up to 9.4 times the errors}  % experiments/FIGS/data/v2/index.json :: binary probe bbqq_vs_ccqq: (1 - AUC) of FLAV_F0 over L188, run by run :: sha256 c6c4138ee5c86423 :: prose twin of the run farther from 1
\newcommand{\MlstCmpBbqqFoneOverFull}{1.510 and 1.447}  % experiments/FIGS/data/v2/index.json :: binary probe bbqq_vs_ccqq: (1 - AUC) of FLAV_F1 over L188, run by run :: sha256 c6c4138ee5c86423
\newcommand{\MlstCmpBbqqFoneOverFullPr}{at least 45\% more errors}  % experiments/FIGS/data/v2/index.json :: binary probe bbqq_vs_ccqq: (1 - AUC) of FLAV_F1 over L188, run by run :: sha256 c6c4138ee5c86423 :: prose twin: two significant figures, direction in words
\newcommand{\MlstCmpBbqqFoneOverFullPrMax}{up to 51\% more errors}  % experiments/FIGS/data/v2/index.json :: binary probe bbqq_vs_ccqq: (1 - AUC) of FLAV_F1 over L188, run by run :: sha256 c6c4138ee5c86423 :: prose twin of the run farther from 1
\newcommand{\MlstCmpBbqqFonerOverFull}{2.272 and 2.281}  % experiments/FIGS/data/v2/index.json :: binary probe bbqq_vs_ccqq: (1 - AUC) of FLAV_F1R over L188, run by run :: sha256 c6c4138ee5c86423
\newcommand{\MlstCmpBbqqFonerOverFullPr}{at least 2.3 times the errors}  % experiments/FIGS/data/v2/index.json :: binary probe bbqq_vs_ccqq: (1 - AUC) of FLAV_F1R over L188, run by run :: sha256 c6c4138ee5c86423 :: prose twin: two significant figures, direction in words
\newcommand{\MlstCmpBbqqFonerOverFullPrMax}{up to 2.3 times the errors}  % experiments/FIGS/data/v2/index.json :: binary probe bbqq_vs_ccqq: (1 - AUC) of FLAV_F1R over L188, run by run :: sha256 c6c4138ee5c86423 :: prose twin of the run farther from 1
\newcommand{\MlstCmpBbqqRandOverFull}{1.42\,\ensuremath{\pm}\,0.06}  % experiments/FIGS/data/v2/index.json :: binary probe bbqq_vs_ccqq: (1 - AUC) of each random partition over the 188-class run of its seed :: sha256 c6c4138ee5c86423
\newcommand{\MlstCmpBbqqRandOverFullPr}{42\% more errors}  % experiments/FIGS/data/v2/index.json :: binary probe bbqq_vs_ccqq: (1 - AUC) of each random partition over the 188-class run of its seed :: sha256 c6c4138ee5c86423 :: prose twin: two significant figures, direction in words
\newcommand{\MlstPairRandMergedBbqqN}{2}  % experiments/FIGS/data/v2/index.json :: random partitions that merge bbqq_vs_ccqq (count) :: sha256 c6c4138ee5c86423
\newcommand{\MlstPairRandApartBbqqN}{3}  % experiments/FIGS/data/v2/index.json :: random partitions that keep apart bbqq_vs_ccqq (count) :: sha256 c6c4138ee5c86423
\newcommand{\MlstPairRandMergedOverApartBbqq}{0.996}  % experiments/FIGS/data/v2/index.json :: binary probe bbqq_vs_ccqq: mean (1 - AUC) of the partitions that merge it over that of those that keep it apart (one run each) :: sha256 c6c4138ee5c86423
\newcommand{\MlstPairRandMergedOverApartBbqqPr}{0.44\% fewer errors}  % experiments/FIGS/data/v2/index.json :: binary probe bbqq_vs_ccqq: mean (1 - AUC) of the partitions that merge it over that of those that keep it apart (one run each) :: sha256 c6c4138ee5c86423 :: prose twin: two significant figures, direction in words
\newcommand{\MlstCmpEemmResOverFull}{--}  % experiments/FIGS/data/v2/index.json :: binary probe ee_vs_mm: (1 - AUC) of L162 over L188, run by run :: sha256 c6c4138ee5c86423 :: the reference's 1 - AUC is 0
\newcommand{\MlstCmpEemmResOverFullPr}{undefined, since the reference separates the test jets completely}  % experiments/FIGS/data/v2/index.json :: binary probe ee_vs_mm: (1 - AUC) of L162 over L188, run by run :: sha256 c6c4138ee5c86423 :: the reference's 1 - AUC is 0
\newcommand{\MlstCmpEemmMidOverFull}{--}  % experiments/FIGS/data/v2/index.json :: binary probe ee_vs_mm: (1 - AUC) of R42_Q1 over L188, run by run :: sha256 c6c4138ee5c86423 :: the reference's 1 - AUC is 0
\newcommand{\MlstCmpEemmMidOverFullPr}{undefined, since the reference separates the test jets completely}  % experiments/FIGS/data/v2/index.json :: binary probe ee_vs_mm: (1 - AUC) of R42_Q1 over L188, run by run :: sha256 c6c4138ee5c86423 :: the reference's 1 - AUC is 0
\newcommand{\MlstCmpEemmCoarseOverFull}{--}  % experiments/FIGS/data/v2/index.json :: binary probe ee_vs_mm: (1 - AUC) of R16_Q1 over L188, run by run :: sha256 c6c4138ee5c86423 :: the reference's 1 - AUC is 0
\newcommand{\MlstCmpEemmCoarseOverFullPr}{undefined, since the reference separates the test jets completely}  % experiments/FIGS/data/v2/index.json :: binary probe ee_vs_mm: (1 - AUC) of R16_Q1 over L188, run by run :: sha256 c6c4138ee5c86423 :: the reference's 1 - AUC is 0
\newcommand{\MlstCmpEemmFoneOverFzero}{--}  % experiments/FIGS/data/v2/index.json :: binary probe ee_vs_mm: (1 - AUC) of FLAV_F1 over FLAV_F0, run by run :: sha256 c6c4138ee5c86423 :: the reference's 1 - AUC is 0
\newcommand{\MlstCmpEemmFoneOverFzeroPr}{undefined, since the reference separates the test jets completely}  % experiments/FIGS/data/v2/index.json :: binary probe ee_vs_mm: (1 - AUC) of FLAV_F1 over FLAV_F0, run by run :: sha256 c6c4138ee5c86423 :: the reference's 1 - AUC is 0
\newcommand{\MlstCmpEemmFoneOverFzeroPrMax}{undefined, since the reference separates the test jets completely}  % experiments/FIGS/data/v2/index.json :: binary probe ee_vs_mm: (1 - AUC) of FLAV_F1 over FLAV_F0, run by run :: sha256 c6c4138ee5c86423 :: the reference's 1 - AUC is 0
\newcommand{\MlstCmpEemmFonerOverFone}{--}  % experiments/FIGS/data/v2/index.json :: binary probe ee_vs_mm: (1 - AUC) of FLAV_F1R over FLAV_F1, run by run :: sha256 c6c4138ee5c86423 :: the reference's 1 - AUC is 0
\newcommand{\MlstCmpEemmFonerOverFonePr}{undefined, since the reference separates the test jets completely}  % experiments/FIGS/data/v2/index.json :: binary probe ee_vs_mm: (1 - AUC) of FLAV_F1R over FLAV_F1, run by run :: sha256 c6c4138ee5c86423 :: the reference's 1 - AUC is 0
\newcommand{\MlstCmpEemmFonerOverFonePrMax}{undefined, since the reference separates the test jets completely}  % experiments/FIGS/data/v2/index.json :: binary probe ee_vs_mm: (1 - AUC) of FLAV_F1R over FLAV_F1, run by run :: sha256 c6c4138ee5c86423 :: the reference's 1 - AUC is 0
\newcommand{\MlstCmpEemmFonerOverFzero}{--}  % experiments/FIGS/data/v2/index.json :: binary probe ee_vs_mm: (1 - AUC) of FLAV_F1R over FLAV_F0, run by run :: sha256 c6c4138ee5c86423 :: the reference's 1 - AUC is 0
\newcommand{\MlstCmpEemmFonerOverFzeroPr}{undefined, since the reference separates the test jets completely}  % experiments/FIGS/data/v2/index.json :: binary probe ee_vs_mm: (1 - AUC) of FLAV_F1R over FLAV_F0, run by run :: sha256 c6c4138ee5c86423 :: the reference's 1 - AUC is 0
\newcommand{\MlstCmpEemmFonerOverFzeroPrMax}{undefined, since the reference separates the test jets completely}  % experiments/FIGS/data/v2/index.json :: binary probe ee_vs_mm: (1 - AUC) of FLAV_F1R over FLAV_F0, run by run :: sha256 c6c4138ee5c86423 :: the reference's 1 - AUC is 0
\newcommand{\MlstCmpEemmFzeroOverCoarse}{0.0000022 and 0.0000003}  % experiments/FIGS/data/v2/index.json :: binary probe ee_vs_mm: (1 - AUC) of FLAV_F0 over R16_Q1, run by run :: sha256 c6c4138ee5c86423
\newcommand{\MlstCmpEemmFzeroOverCoarsePr}{at most 0.0000022 times the errors}  % experiments/FIGS/data/v2/index.json :: binary probe ee_vs_mm: (1 - AUC) of FLAV_F0 over R16_Q1, run by run :: sha256 c6c4138ee5c86423 :: prose twin: two significant figures, direction in words
\newcommand{\MlstCmpEemmFzeroOverCoarsePrMax}{down to 0.00000034 times the errors}  % experiments/FIGS/data/v2/index.json :: binary probe ee_vs_mm: (1 - AUC) of FLAV_F0 over R16_Q1, run by run :: sha256 c6c4138ee5c86423 :: prose twin of the run farther from 1
\newcommand{\MlstCmpEemmFoneOverCoarse}{0.000000 and 0.000038}  % experiments/FIGS/data/v2/index.json :: binary probe ee_vs_mm: (1 - AUC) of FLAV_F1 over R16_Q1, run by run :: sha256 c6c4138ee5c86423
\newcommand{\MlstCmpEemmFoneOverCoarsePr}{at most 0.000038 times the errors}  % experiments/FIGS/data/v2/index.json :: binary probe ee_vs_mm: (1 - AUC) of FLAV_F1 over R16_Q1, run by run :: sha256 c6c4138ee5c86423 :: prose twin: two significant figures, direction in words
\newcommand{\MlstCmpEemmFoneOverCoarsePrMax}{down to 0.000000042 times the errors}  % experiments/FIGS/data/v2/index.json :: binary probe ee_vs_mm: (1 - AUC) of FLAV_F1 over R16_Q1, run by run :: sha256 c6c4138ee5c86423 :: prose twin of the run farther from 1
\newcommand{\MlstCmpEemmFonerOverCoarse}{0.00000005 and 0.00000017}  % experiments/FIGS/data/v2/index.json :: binary probe ee_vs_mm: (1 - AUC) of FLAV_F1R over R16_Q1, run by run :: sha256 c6c4138ee5c86423
\newcommand{\MlstCmpEemmFonerOverCoarsePr}{at most 0.00000017 times the errors}  % experiments/FIGS/data/v2/index.json :: binary probe ee_vs_mm: (1 - AUC) of FLAV_F1R over R16_Q1, run by run :: sha256 c6c4138ee5c86423 :: prose twin: two significant figures, direction in words
\newcommand{\MlstCmpEemmFonerOverCoarsePrMax}{down to 0.000000055 times the errors}  % experiments/FIGS/data/v2/index.json :: binary probe ee_vs_mm: (1 - AUC) of FLAV_F1R over R16_Q1, run by run :: sha256 c6c4138ee5c86423 :: prose twin of the run farther from 1
\newcommand{\MlstCmpEemmFzeroOverFull}{--}  % experiments/FIGS/data/v2/index.json :: binary probe ee_vs_mm: (1 - AUC) of FLAV_F0 over L188, run by run :: sha256 c6c4138ee5c86423 :: the reference's 1 - AUC is 0
\newcommand{\MlstCmpEemmFzeroOverFullPr}{undefined, since the reference separates the test jets completely}  % experiments/FIGS/data/v2/index.json :: binary probe ee_vs_mm: (1 - AUC) of FLAV_F0 over L188, run by run :: sha256 c6c4138ee5c86423 :: the reference's 1 - AUC is 0
\newcommand{\MlstCmpEemmFzeroOverFullPrMax}{undefined, since the reference separates the test jets completely}  % experiments/FIGS/data/v2/index.json :: binary probe ee_vs_mm: (1 - AUC) of FLAV_F0 over L188, run by run :: sha256 c6c4138ee5c86423 :: the reference's 1 - AUC is 0
\newcommand{\MlstCmpEemmFoneOverFull}{--}  % experiments/FIGS/data/v2/index.json :: binary probe ee_vs_mm: (1 - AUC) of FLAV_F1 over L188, run by run :: sha256 c6c4138ee5c86423 :: the reference's 1 - AUC is 0
\newcommand{\MlstCmpEemmFoneOverFullPr}{undefined, since the reference separates the test jets completely}  % experiments/FIGS/data/v2/index.json :: binary probe ee_vs_mm: (1 - AUC) of FLAV_F1 over L188, run by run :: sha256 c6c4138ee5c86423 :: the reference's 1 - AUC is 0
\newcommand{\MlstCmpEemmFoneOverFullPrMax}{undefined, since the reference separates the test jets completely}  % experiments/FIGS/data/v2/index.json :: binary probe ee_vs_mm: (1 - AUC) of FLAV_F1 over L188, run by run :: sha256 c6c4138ee5c86423 :: the reference's 1 - AUC is 0
\newcommand{\MlstCmpEemmFonerOverFull}{--}  % experiments/FIGS/data/v2/index.json :: binary probe ee_vs_mm: (1 - AUC) of FLAV_F1R over L188, run by run :: sha256 c6c4138ee5c86423 :: the reference's 1 - AUC is 0
\newcommand{\MlstCmpEemmFonerOverFullPr}{undefined, since the reference separates the test jets completely}  % experiments/FIGS/data/v2/index.json :: binary probe ee_vs_mm: (1 - AUC) of FLAV_F1R over L188, run by run :: sha256 c6c4138ee5c86423 :: the reference's 1 - AUC is 0
\newcommand{\MlstCmpEemmFonerOverFullPrMax}{undefined, since the reference separates the test jets completely}  % experiments/FIGS/data/v2/index.json :: binary probe ee_vs_mm: (1 - AUC) of FLAV_F1R over L188, run by run :: sha256 c6c4138ee5c86423 :: the reference's 1 - AUC is 0
\newcommand{\MlstCmpEemmRandOverFull}{--}  % experiments/FIGS/data/v2/index.json :: binary probe ee_vs_mm: (1 - AUC) of each random partition over the 188-class run of its seed :: sha256 c6c4138ee5c86423 :: the reference's 1 - AUC is 0
\newcommand{\MlstCmpEemmRandOverFullPr}{undefined, since the reference separates the test jets completely}  % experiments/FIGS/data/v2/index.json :: binary probe ee_vs_mm: (1 - AUC) of each random partition over the 188-class run of its seed :: sha256 c6c4138ee5c86423 :: the reference's 1 - AUC is 0
\newcommand{\MlstPairRandMergedEemmN}{3}  % experiments/FIGS/data/v2/index.json :: random partitions that merge ee_vs_mm (count) :: sha256 c6c4138ee5c86423
\newcommand{\MlstPairRandApartEemmN}{2}  % experiments/FIGS/data/v2/index.json :: random partitions that keep apart ee_vs_mm (count) :: sha256 c6c4138ee5c86423
\newcommand{\MlstPairRandMergedOverApartEemm}{75.3}  % experiments/FIGS/data/v2/index.json :: binary probe ee_vs_mm: mean (1 - AUC) of the partitions that merge it over that of those that keep it apart (one run each) :: sha256 c6c4138ee5c86423
\newcommand{\MlstPairRandMergedOverApartEemmPr}{75 times the errors}  % experiments/FIGS/data/v2/index.json :: binary probe ee_vs_mm: mean (1 - AUC) of the partitions that merge it over that of those that keep it apart (one run each) :: sha256 c6c4138ee5c86423 :: prose twin: two significant figures, direction in words
\newcommand{\MlstCmpQcdBbResOverFull}{1.15\,\ensuremath{\pm}\,0.04}  % experiments/FIGS/data/v2/index.json :: binary probe qcd_bb_vs_cc: (1 - AUC) of L162 over L188, run by run :: sha256 c6c4138ee5c86423
\newcommand{\MlstCmpQcdBbResOverFullPr}{15\% more errors}  % experiments/FIGS/data/v2/index.json :: binary probe qcd_bb_vs_cc: (1 - AUC) of L162 over L188, run by run :: sha256 c6c4138ee5c86423 :: prose twin: two significant figures, direction in words
\newcommand{\MlstCmpQcdBbMidOverFull}{1.163\,\ensuremath{\pm}\,0.016}  % experiments/FIGS/data/v2/index.json :: binary probe qcd_bb_vs_cc: (1 - AUC) of R42_Q1 over L188, run by run :: sha256 c6c4138ee5c86423
\newcommand{\MlstCmpQcdBbMidOverFullPr}{16\% more errors}  % experiments/FIGS/data/v2/index.json :: binary probe qcd_bb_vs_cc: (1 - AUC) of R42_Q1 over L188, run by run :: sha256 c6c4138ee5c86423 :: prose twin: two significant figures, direction in words
\newcommand{\MlstCmpQcdBbCoarseOverFull}{3.294\,\ensuremath{\pm}\,0.033}  % experiments/FIGS/data/v2/index.json :: binary probe qcd_bb_vs_cc: (1 - AUC) of R16_Q1 over L188, run by run :: sha256 c6c4138ee5c86423
\newcommand{\MlstCmpQcdBbCoarseOverFullPr}{3.3 times the errors}  % experiments/FIGS/data/v2/index.json :: binary probe qcd_bb_vs_cc: (1 - AUC) of R16_Q1 over L188, run by run :: sha256 c6c4138ee5c86423 :: prose twin: two significant figures, direction in words
\newcommand{\MlstCmpQcdBbFoneOverFzero}{0.368 and 0.294}  % experiments/FIGS/data/v2/index.json :: binary probe qcd_bb_vs_cc: (1 - AUC) of FLAV_F1 over FLAV_F0, run by run :: sha256 c6c4138ee5c86423
\newcommand{\MlstCmpQcdBbFoneOverFzeroPr}{at most 0.37 times the errors}  % experiments/FIGS/data/v2/index.json :: binary probe qcd_bb_vs_cc: (1 - AUC) of FLAV_F1 over FLAV_F0, run by run :: sha256 c6c4138ee5c86423 :: prose twin: two significant figures, direction in words
\newcommand{\MlstCmpQcdBbFoneOverFzeroPrMax}{down to 0.29 times the errors}  % experiments/FIGS/data/v2/index.json :: binary probe qcd_bb_vs_cc: (1 - AUC) of FLAV_F1 over FLAV_F0, run by run :: sha256 c6c4138ee5c86423 :: prose twin of the run farther from 1
\newcommand{\MlstCmpQcdBbFonerOverFone}{1.361 and 1.429}  % experiments/FIGS/data/v2/index.json :: binary probe qcd_bb_vs_cc: (1 - AUC) of FLAV_F1R over FLAV_F1, run by run :: sha256 c6c4138ee5c86423
\newcommand{\MlstCmpQcdBbFonerOverFonePr}{at least 36\% more errors}  % experiments/FIGS/data/v2/index.json :: binary probe qcd_bb_vs_cc: (1 - AUC) of FLAV_F1R over FLAV_F1, run by run :: sha256 c6c4138ee5c86423 :: prose twin: two significant figures, direction in words
\newcommand{\MlstCmpQcdBbFonerOverFonePrMax}{up to 43\% more errors}  % experiments/FIGS/data/v2/index.json :: binary probe qcd_bb_vs_cc: (1 - AUC) of FLAV_F1R over FLAV_F1, run by run :: sha256 c6c4138ee5c86423 :: prose twin of the run farther from 1
\newcommand{\MlstCmpQcdBbFonerOverFzero}{0.501 and 0.420}  % experiments/FIGS/data/v2/index.json :: binary probe qcd_bb_vs_cc: (1 - AUC) of FLAV_F1R over FLAV_F0, run by run :: sha256 c6c4138ee5c86423
\newcommand{\MlstCmpQcdBbFonerOverFzeroPr}{at least 50\% fewer errors}  % experiments/FIGS/data/v2/index.json :: binary probe qcd_bb_vs_cc: (1 - AUC) of FLAV_F1R over FLAV_F0, run by run :: sha256 c6c4138ee5c86423 :: prose twin: two significant figures, direction in words
\newcommand{\MlstCmpQcdBbFonerOverFzeroPrMax}{down to 0.42 times the errors}  % experiments/FIGS/data/v2/index.json :: binary probe qcd_bb_vs_cc: (1 - AUC) of FLAV_F1R over FLAV_F0, run by run :: sha256 c6c4138ee5c86423 :: prose twin of the run farther from 1
\newcommand{\MlstCmpQcdBbFzeroOverCoarse}{1.28 and 1.58}  % experiments/FIGS/data/v2/index.json :: binary probe qcd_bb_vs_cc: (1 - AUC) of FLAV_F0 over R16_Q1, run by run :: sha256 c6c4138ee5c86423
\newcommand{\MlstCmpQcdBbFzeroOverCoarsePr}{at least 28\% more errors}  % experiments/FIGS/data/v2/index.json :: binary probe qcd_bb_vs_cc: (1 - AUC) of FLAV_F0 over R16_Q1, run by run :: sha256 c6c4138ee5c86423 :: prose twin: two significant figures, direction in words
\newcommand{\MlstCmpQcdBbFzeroOverCoarsePrMax}{up to 58\% more errors}  % experiments/FIGS/data/v2/index.json :: binary probe qcd_bb_vs_cc: (1 - AUC) of FLAV_F0 over R16_Q1, run by run :: sha256 c6c4138ee5c86423 :: prose twin of the run farther from 1
\newcommand{\MlstCmpQcdBbFoneOverCoarse}{0.4692 and 0.4648}  % experiments/FIGS/data/v2/index.json :: binary probe qcd_bb_vs_cc: (1 - AUC) of FLAV_F1 over R16_Q1, run by run :: sha256 c6c4138ee5c86423
\newcommand{\MlstCmpQcdBbFoneOverCoarsePr}{at most 0.47 times the errors}  % experiments/FIGS/data/v2/index.json :: binary probe qcd_bb_vs_cc: (1 - AUC) of FLAV_F1 over R16_Q1, run by run :: sha256 c6c4138ee5c86423 :: prose twin: two significant figures, direction in words
\newcommand{\MlstCmpQcdBbFoneOverCoarsePrMax}{down to 0.46 times the errors}  % experiments/FIGS/data/v2/index.json :: binary probe qcd_bb_vs_cc: (1 - AUC) of FLAV_F1 over R16_Q1, run by run :: sha256 c6c4138ee5c86423 :: prose twin of the run farther from 1
\newcommand{\MlstCmpQcdBbFonerOverCoarse}{0.639 and 0.664}  % experiments/FIGS/data/v2/index.json :: binary probe qcd_bb_vs_cc: (1 - AUC) of FLAV_F1R over R16_Q1, run by run :: sha256 c6c4138ee5c86423
\newcommand{\MlstCmpQcdBbFonerOverCoarsePr}{at least 34\% fewer errors}  % experiments/FIGS/data/v2/index.json :: binary probe qcd_bb_vs_cc: (1 - AUC) of FLAV_F1R over R16_Q1, run by run :: sha256 c6c4138ee5c86423 :: prose twin: two significant figures, direction in words
\newcommand{\MlstCmpQcdBbFonerOverCoarsePrMax}{up to 36\% fewer errors}  % experiments/FIGS/data/v2/index.json :: binary probe qcd_bb_vs_cc: (1 - AUC) of FLAV_F1R over R16_Q1, run by run :: sha256 c6c4138ee5c86423 :: prose twin of the run farther from 1
\newcommand{\MlstCmpQcdBbFzeroOverFull}{4.24 and 5.21}  % experiments/FIGS/data/v2/index.json :: binary probe qcd_bb_vs_cc: (1 - AUC) of FLAV_F0 over L188, run by run :: sha256 c6c4138ee5c86423
\newcommand{\MlstCmpQcdBbFzeroOverFullPr}{at least 4.2 times the errors}  % experiments/FIGS/data/v2/index.json :: binary probe qcd_bb_vs_cc: (1 - AUC) of FLAV_F0 over L188, run by run :: sha256 c6c4138ee5c86423 :: prose twin: two significant figures, direction in words
\newcommand{\MlstCmpQcdBbFzeroOverFullPrMax}{up to 5.2 times the errors}  % experiments/FIGS/data/v2/index.json :: binary probe qcd_bb_vs_cc: (1 - AUC) of FLAV_F0 over L188, run by run :: sha256 c6c4138ee5c86423 :: prose twin of the run farther from 1
\newcommand{\MlstCmpQcdBbFoneOverFull}{1.561 and 1.532}  % experiments/FIGS/data/v2/index.json :: binary probe qcd_bb_vs_cc: (1 - AUC) of FLAV_F1 over L188, run by run :: sha256 c6c4138ee5c86423
\newcommand{\MlstCmpQcdBbFoneOverFullPr}{at least 53\% more errors}  % experiments/FIGS/data/v2/index.json :: binary probe qcd_bb_vs_cc: (1 - AUC) of FLAV_F1 over L188, run by run :: sha256 c6c4138ee5c86423 :: prose twin: two significant figures, direction in words
\newcommand{\MlstCmpQcdBbFoneOverFullPrMax}{up to 56\% more errors}  % experiments/FIGS/data/v2/index.json :: binary probe qcd_bb_vs_cc: (1 - AUC) of FLAV_F1 over L188, run by run :: sha256 c6c4138ee5c86423 :: prose twin of the run farther from 1
\newcommand{\MlstCmpQcdBbFonerOverFull}{2.124 and 2.189}  % experiments/FIGS/data/v2/index.json :: binary probe qcd_bb_vs_cc: (1 - AUC) of FLAV_F1R over L188, run by run :: sha256 c6c4138ee5c86423
\newcommand{\MlstCmpQcdBbFonerOverFullPr}{at least 2.1 times the errors}  % experiments/FIGS/data/v2/index.json :: binary probe qcd_bb_vs_cc: (1 - AUC) of FLAV_F1R over L188, run by run :: sha256 c6c4138ee5c86423 :: prose twin: two significant figures, direction in words
\newcommand{\MlstCmpQcdBbFonerOverFullPrMax}{up to 2.2 times the errors}  % experiments/FIGS/data/v2/index.json :: binary probe qcd_bb_vs_cc: (1 - AUC) of FLAV_F1R over L188, run by run :: sha256 c6c4138ee5c86423 :: prose twin of the run farther from 1
\newcommand{\MlstCmpQcdBResOverFull}{1.085\,\ensuremath{\pm}\,0.004}  % experiments/FIGS/data/v2/index.json :: binary probe qcd_b_vs_light: (1 - AUC) of L162 over L188, run by run :: sha256 c6c4138ee5c86423
\newcommand{\MlstCmpQcdBResOverFullPr}{8.5\% more errors}  % experiments/FIGS/data/v2/index.json :: binary probe qcd_b_vs_light: (1 - AUC) of L162 over L188, run by run :: sha256 c6c4138ee5c86423 :: prose twin: two significant figures, direction in words
\newcommand{\MlstCmpQcdBMidOverFull}{1.081\,\ensuremath{\pm}\,0.008}  % experiments/FIGS/data/v2/index.json :: binary probe qcd_b_vs_light: (1 - AUC) of R42_Q1 over L188, run by run :: sha256 c6c4138ee5c86423
\newcommand{\MlstCmpQcdBMidOverFullPr}{8.1\% more errors}  % experiments/FIGS/data/v2/index.json :: binary probe qcd_b_vs_light: (1 - AUC) of R42_Q1 over L188, run by run :: sha256 c6c4138ee5c86423 :: prose twin: two significant figures, direction in words
\newcommand{\MlstCmpQcdBCoarseOverFull}{1.54\,\ensuremath{\pm}\,0.07}  % experiments/FIGS/data/v2/index.json :: binary probe qcd_b_vs_light: (1 - AUC) of R16_Q1 over L188, run by run :: sha256 c6c4138ee5c86423
\newcommand{\MlstCmpQcdBCoarseOverFullPr}{54\% more errors}  % experiments/FIGS/data/v2/index.json :: binary probe qcd_b_vs_light: (1 - AUC) of R16_Q1 over L188, run by run :: sha256 c6c4138ee5c86423 :: prose twin: two significant figures, direction in words
\newcommand{\MlstCmpQcdBFoneOverFzero}{0.778 and 0.835}  % experiments/FIGS/data/v2/index.json :: binary probe qcd_b_vs_light: (1 - AUC) of FLAV_F1 over FLAV_F0, run by run :: sha256 c6c4138ee5c86423
\newcommand{\MlstCmpQcdBFoneOverFzeroPr}{at least 17\% fewer errors}  % experiments/FIGS/data/v2/index.json :: binary probe qcd_b_vs_light: (1 - AUC) of FLAV_F1 over FLAV_F0, run by run :: sha256 c6c4138ee5c86423 :: prose twin: two significant figures, direction in words
\newcommand{\MlstCmpQcdBFoneOverFzeroPrMax}{up to 22\% fewer errors}  % experiments/FIGS/data/v2/index.json :: binary probe qcd_b_vs_light: (1 - AUC) of FLAV_F1 over FLAV_F0, run by run :: sha256 c6c4138ee5c86423 :: prose twin of the run farther from 1
\newcommand{\MlstCmpQcdBFonerOverFone}{1.054 and 1.061}  % experiments/FIGS/data/v2/index.json :: binary probe qcd_b_vs_light: (1 - AUC) of FLAV_F1R over FLAV_F1, run by run :: sha256 c6c4138ee5c86423
\newcommand{\MlstCmpQcdBFonerOverFonePr}{equal within 7\%}  % experiments/FIGS/data/v2/index.json :: binary probe qcd_b_vs_light: (1 - AUC) of FLAV_F1R over FLAV_F1, run by run :: sha256 c6c4138ee5c86423 :: prose twin: two significant figures, direction in words
\newcommand{\MlstCmpQcdBFonerOverFonePrMax}{equal within 7\%}  % experiments/FIGS/data/v2/index.json :: binary probe qcd_b_vs_light: (1 - AUC) of FLAV_F1R over FLAV_F1, run by run :: sha256 c6c4138ee5c86423 :: prose twin of the run farther from 1
\newcommand{\MlstCmpQcdBFonerOverFzero}{0.820 and 0.885}  % experiments/FIGS/data/v2/index.json :: binary probe qcd_b_vs_light: (1 - AUC) of FLAV_F1R over FLAV_F0, run by run :: sha256 c6c4138ee5c86423
\newcommand{\MlstCmpQcdBFonerOverFzeroPr}{at least 11\% fewer errors}  % experiments/FIGS/data/v2/index.json :: binary probe qcd_b_vs_light: (1 - AUC) of FLAV_F1R over FLAV_F0, run by run :: sha256 c6c4138ee5c86423 :: prose twin: two significant figures, direction in words
\newcommand{\MlstCmpQcdBFonerOverFzeroPrMax}{up to 18\% fewer errors}  % experiments/FIGS/data/v2/index.json :: binary probe qcd_b_vs_light: (1 - AUC) of FLAV_F1R over FLAV_F0, run by run :: sha256 c6c4138ee5c86423 :: prose twin of the run farther from 1
\newcommand{\MlstCmpQcdBFzeroOverCoarse}{0.962 and 0.981}  % experiments/FIGS/data/v2/index.json :: binary probe qcd_b_vs_light: (1 - AUC) of FLAV_F0 over R16_Q1, run by run :: sha256 c6c4138ee5c86423
\newcommand{\MlstCmpQcdBFzeroOverCoarsePr}{equal within 4\%}  % experiments/FIGS/data/v2/index.json :: binary probe qcd_b_vs_light: (1 - AUC) of FLAV_F0 over R16_Q1, run by run :: sha256 c6c4138ee5c86423 :: prose twin: two significant figures, direction in words
\newcommand{\MlstCmpQcdBFzeroOverCoarsePrMax}{equal within 4\%}  % experiments/FIGS/data/v2/index.json :: binary probe qcd_b_vs_light: (1 - AUC) of FLAV_F0 over R16_Q1, run by run :: sha256 c6c4138ee5c86423 :: prose twin of the run farther from 1
\newcommand{\MlstCmpQcdBFoneOverCoarse}{0.749 and 0.819}  % experiments/FIGS/data/v2/index.json :: binary probe qcd_b_vs_light: (1 - AUC) of FLAV_F1 over R16_Q1, run by run :: sha256 c6c4138ee5c86423
\newcommand{\MlstCmpQcdBFoneOverCoarsePr}{at least 18\% fewer errors}  % experiments/FIGS/data/v2/index.json :: binary probe qcd_b_vs_light: (1 - AUC) of FLAV_F1 over R16_Q1, run by run :: sha256 c6c4138ee5c86423 :: prose twin: two significant figures, direction in words
\newcommand{\MlstCmpQcdBFoneOverCoarsePrMax}{up to 25\% fewer errors}  % experiments/FIGS/data/v2/index.json :: binary probe qcd_b_vs_light: (1 - AUC) of FLAV_F1 over R16_Q1, run by run :: sha256 c6c4138ee5c86423 :: prose twin of the run farther from 1
\newcommand{\MlstCmpQcdBFonerOverCoarse}{0.789 and 0.869}  % experiments/FIGS/data/v2/index.json :: binary probe qcd_b_vs_light: (1 - AUC) of FLAV_F1R over R16_Q1, run by run :: sha256 c6c4138ee5c86423
\newcommand{\MlstCmpQcdBFonerOverCoarsePr}{at least 13\% fewer errors}  % experiments/FIGS/data/v2/index.json :: binary probe qcd_b_vs_light: (1 - AUC) of FLAV_F1R over R16_Q1, run by run :: sha256 c6c4138ee5c86423 :: prose twin: two significant figures, direction in words
\newcommand{\MlstCmpQcdBFonerOverCoarsePrMax}{up to 21\% fewer errors}  % experiments/FIGS/data/v2/index.json :: binary probe qcd_b_vs_light: (1 - AUC) of FLAV_F1R over R16_Q1, run by run :: sha256 c6c4138ee5c86423 :: prose twin of the run farther from 1
\newcommand{\MlstCmpQcdBFzeroOverFull}{1.56 and 1.45}  % experiments/FIGS/data/v2/index.json :: binary probe qcd_b_vs_light: (1 - AUC) of FLAV_F0 over L188, run by run :: sha256 c6c4138ee5c86423
\newcommand{\MlstCmpQcdBFzeroOverFullPr}{at least 45\% more errors}  % experiments/FIGS/data/v2/index.json :: binary probe qcd_b_vs_light: (1 - AUC) of FLAV_F0 over L188, run by run :: sha256 c6c4138ee5c86423 :: prose twin: two significant figures, direction in words
\newcommand{\MlstCmpQcdBFzeroOverFullPrMax}{up to 56\% more errors}  % experiments/FIGS/data/v2/index.json :: binary probe qcd_b_vs_light: (1 - AUC) of FLAV_F0 over L188, run by run :: sha256 c6c4138ee5c86423 :: prose twin of the run farther from 1
\newcommand{\MlstCmpQcdBFoneOverFull}{1.212 and 1.213}  % experiments/FIGS/data/v2/index.json :: binary probe qcd_b_vs_light: (1 - AUC) of FLAV_F1 over L188, run by run :: sha256 c6c4138ee5c86423
\newcommand{\MlstCmpQcdBFoneOverFullPr}{at least 21\% more errors}  % experiments/FIGS/data/v2/index.json :: binary probe qcd_b_vs_light: (1 - AUC) of FLAV_F1 over L188, run by run :: sha256 c6c4138ee5c86423 :: prose twin: two significant figures, direction in words
\newcommand{\MlstCmpQcdBFoneOverFullPrMax}{up to 21\% more errors}  % experiments/FIGS/data/v2/index.json :: binary probe qcd_b_vs_light: (1 - AUC) of FLAV_F1 over L188, run by run :: sha256 c6c4138ee5c86423 :: prose twin of the run farther from 1
\newcommand{\MlstCmpQcdBFonerOverFull}{1.277 and 1.286}  % experiments/FIGS/data/v2/index.json :: binary probe qcd_b_vs_light: (1 - AUC) of FLAV_F1R over L188, run by run :: sha256 c6c4138ee5c86423
\newcommand{\MlstCmpQcdBFonerOverFullPr}{at least 28\% more errors}  % experiments/FIGS/data/v2/index.json :: binary probe qcd_b_vs_light: (1 - AUC) of FLAV_F1R over L188, run by run :: sha256 c6c4138ee5c86423 :: prose twin: two significant figures, direction in words
\newcommand{\MlstCmpQcdBFonerOverFullPrMax}{up to 29\% more errors}  % experiments/FIGS/data/v2/index.json :: binary probe qcd_b_vs_light: (1 - AUC) of FLAV_F1R over L188, run by run :: sha256 c6c4138ee5c86423 :: prose twin of the run farther from 1
\newcommand{\MlstCmpQcdCResOverFull}{1.061\,\ensuremath{\pm}\,0.005}  % experiments/FIGS/data/v2/index.json :: binary probe qcd_c_vs_light: (1 - AUC) of L162 over L188, run by run :: sha256 c6c4138ee5c86423
\newcommand{\MlstCmpQcdCResOverFullPr}{6.1\% more errors}  % experiments/FIGS/data/v2/index.json :: binary probe qcd_c_vs_light: (1 - AUC) of L162 over L188, run by run :: sha256 c6c4138ee5c86423 :: prose twin: two significant figures, direction in words
\newcommand{\MlstCmpQcdCMidOverFull}{1.072\,\ensuremath{\pm}\,0.008}  % experiments/FIGS/data/v2/index.json :: binary probe qcd_c_vs_light: (1 - AUC) of R42_Q1 over L188, run by run :: sha256 c6c4138ee5c86423
\newcommand{\MlstCmpQcdCMidOverFullPr}{7.2\% more errors}  % experiments/FIGS/data/v2/index.json :: binary probe qcd_c_vs_light: (1 - AUC) of R42_Q1 over L188, run by run :: sha256 c6c4138ee5c86423 :: prose twin: two significant figures, direction in words
\newcommand{\MlstCmpQcdCCoarseOverFull}{1.407\,\ensuremath{\pm}\,0.021}  % experiments/FIGS/data/v2/index.json :: binary probe qcd_c_vs_light: (1 - AUC) of R16_Q1 over L188, run by run :: sha256 c6c4138ee5c86423
\newcommand{\MlstCmpQcdCCoarseOverFullPr}{41\% more errors}  % experiments/FIGS/data/v2/index.json :: binary probe qcd_c_vs_light: (1 - AUC) of R16_Q1 over L188, run by run :: sha256 c6c4138ee5c86423 :: prose twin: two significant figures, direction in words
\newcommand{\MlstCmpQcdCFoneOverFzero}{0.915 and 0.944}  % experiments/FIGS/data/v2/index.json :: binary probe qcd_c_vs_light: (1 - AUC) of FLAV_F1 over FLAV_F0, run by run :: sha256 c6c4138ee5c86423
\newcommand{\MlstCmpQcdCFoneOverFzeroPr}{at least 5.6\% fewer errors}  % experiments/FIGS/data/v2/index.json :: binary probe qcd_c_vs_light: (1 - AUC) of FLAV_F1 over FLAV_F0, run by run :: sha256 c6c4138ee5c86423 :: prose twin: two significant figures, direction in words
\newcommand{\MlstCmpQcdCFoneOverFzeroPrMax}{up to 8.5\% fewer errors}  % experiments/FIGS/data/v2/index.json :: binary probe qcd_c_vs_light: (1 - AUC) of FLAV_F1 over FLAV_F0, run by run :: sha256 c6c4138ee5c86423 :: prose twin of the run farther from 1
\newcommand{\MlstCmpQcdCFonerOverFone}{0.9535 and 0.9520}  % experiments/FIGS/data/v2/index.json :: binary probe qcd_c_vs_light: (1 - AUC) of FLAV_F1R over FLAV_F1, run by run :: sha256 c6c4138ee5c86423
\newcommand{\MlstCmpQcdCFonerOverFonePr}{at least 4.6\% fewer errors}  % experiments/FIGS/data/v2/index.json :: binary probe qcd_c_vs_light: (1 - AUC) of FLAV_F1R over FLAV_F1, run by run :: sha256 c6c4138ee5c86423 :: prose twin: two significant figures, direction in words
\newcommand{\MlstCmpQcdCFonerOverFonePrMax}{up to 4.8\% fewer errors}  % experiments/FIGS/data/v2/index.json :: binary probe qcd_c_vs_light: (1 - AUC) of FLAV_F1R over FLAV_F1, run by run :: sha256 c6c4138ee5c86423 :: prose twin of the run farther from 1
\newcommand{\MlstCmpQcdCFonerOverFzero}{0.872 and 0.899}  % experiments/FIGS/data/v2/index.json :: binary probe qcd_c_vs_light: (1 - AUC) of FLAV_F1R over FLAV_F0, run by run :: sha256 c6c4138ee5c86423
\newcommand{\MlstCmpQcdCFonerOverFzeroPr}{at least 10\% fewer errors}  % experiments/FIGS/data/v2/index.json :: binary probe qcd_c_vs_light: (1 - AUC) of FLAV_F1R over FLAV_F0, run by run :: sha256 c6c4138ee5c86423 :: prose twin: two significant figures, direction in words
\newcommand{\MlstCmpQcdCFonerOverFzeroPrMax}{up to 13\% fewer errors}  % experiments/FIGS/data/v2/index.json :: binary probe qcd_c_vs_light: (1 - AUC) of FLAV_F1R over FLAV_F0, run by run :: sha256 c6c4138ee5c86423 :: prose twin of the run farther from 1
\newcommand{\MlstCmpQcdCFzeroOverCoarse}{0.952 and 0.934}  % experiments/FIGS/data/v2/index.json :: binary probe qcd_c_vs_light: (1 - AUC) of FLAV_F0 over R16_Q1, run by run :: sha256 c6c4138ee5c86423
\newcommand{\MlstCmpQcdCFzeroOverCoarsePr}{at least 4.8\% fewer errors}  % experiments/FIGS/data/v2/index.json :: binary probe qcd_c_vs_light: (1 - AUC) of FLAV_F0 over R16_Q1, run by run :: sha256 c6c4138ee5c86423 :: prose twin: two significant figures, direction in words
\newcommand{\MlstCmpQcdCFzeroOverCoarsePrMax}{up to 6.6\% fewer errors}  % experiments/FIGS/data/v2/index.json :: binary probe qcd_c_vs_light: (1 - AUC) of FLAV_F0 over R16_Q1, run by run :: sha256 c6c4138ee5c86423 :: prose twin of the run farther from 1
\newcommand{\MlstCmpQcdCFoneOverCoarse}{0.871 and 0.882}  % experiments/FIGS/data/v2/index.json :: binary probe qcd_c_vs_light: (1 - AUC) of FLAV_F1 over R16_Q1, run by run :: sha256 c6c4138ee5c86423
\newcommand{\MlstCmpQcdCFoneOverCoarsePr}{at least 12\% fewer errors}  % experiments/FIGS/data/v2/index.json :: binary probe qcd_c_vs_light: (1 - AUC) of FLAV_F1 over R16_Q1, run by run :: sha256 c6c4138ee5c86423 :: prose twin: two significant figures, direction in words
\newcommand{\MlstCmpQcdCFoneOverCoarsePrMax}{up to 13\% fewer errors}  % experiments/FIGS/data/v2/index.json :: binary probe qcd_c_vs_light: (1 - AUC) of FLAV_F1 over R16_Q1, run by run :: sha256 c6c4138ee5c86423 :: prose twin of the run farther from 1
\newcommand{\MlstCmpQcdCFonerOverCoarse}{0.8301 and 0.8396}  % experiments/FIGS/data/v2/index.json :: binary probe qcd_c_vs_light: (1 - AUC) of FLAV_F1R over R16_Q1, run by run :: sha256 c6c4138ee5c86423
\newcommand{\MlstCmpQcdCFonerOverCoarsePr}{at least 16\% fewer errors}  % experiments/FIGS/data/v2/index.json :: binary probe qcd_c_vs_light: (1 - AUC) of FLAV_F1R over R16_Q1, run by run :: sha256 c6c4138ee5c86423 :: prose twin: two significant figures, direction in words
\newcommand{\MlstCmpQcdCFonerOverCoarsePrMax}{up to 17\% fewer errors}  % experiments/FIGS/data/v2/index.json :: binary probe qcd_c_vs_light: (1 - AUC) of FLAV_F1R over R16_Q1, run by run :: sha256 c6c4138ee5c86423 :: prose twin of the run farther from 1
\newcommand{\MlstCmpQcdCFzeroOverFull}{1.352 and 1.293}  % experiments/FIGS/data/v2/index.json :: binary probe qcd_c_vs_light: (1 - AUC) of FLAV_F0 over L188, run by run :: sha256 c6c4138ee5c86423
\newcommand{\MlstCmpQcdCFzeroOverFullPr}{at least 29\% more errors}  % experiments/FIGS/data/v2/index.json :: binary probe qcd_c_vs_light: (1 - AUC) of FLAV_F0 over L188, run by run :: sha256 c6c4138ee5c86423 :: prose twin: two significant figures, direction in words
\newcommand{\MlstCmpQcdCFzeroOverFullPrMax}{up to 35\% more errors}  % experiments/FIGS/data/v2/index.json :: binary probe qcd_c_vs_light: (1 - AUC) of FLAV_F0 over L188, run by run :: sha256 c6c4138ee5c86423 :: prose twin of the run farther from 1
\newcommand{\MlstCmpQcdCFoneOverFull}{1.237 and 1.220}  % experiments/FIGS/data/v2/index.json :: binary probe qcd_c_vs_light: (1 - AUC) of FLAV_F1 over L188, run by run :: sha256 c6c4138ee5c86423
\newcommand{\MlstCmpQcdCFoneOverFullPr}{at least 22\% more errors}  % experiments/FIGS/data/v2/index.json :: binary probe qcd_c_vs_light: (1 - AUC) of FLAV_F1 over L188, run by run :: sha256 c6c4138ee5c86423 :: prose twin: two significant figures, direction in words
\newcommand{\MlstCmpQcdCFoneOverFullPrMax}{up to 24\% more errors}  % experiments/FIGS/data/v2/index.json :: binary probe qcd_c_vs_light: (1 - AUC) of FLAV_F1 over L188, run by run :: sha256 c6c4138ee5c86423 :: prose twin of the run farther from 1
\newcommand{\MlstCmpQcdCFonerOverFull}{1.179 and 1.162}  % experiments/FIGS/data/v2/index.json :: binary probe qcd_c_vs_light: (1 - AUC) of FLAV_F1R over L188, run by run :: sha256 c6c4138ee5c86423
\newcommand{\MlstCmpQcdCFonerOverFullPr}{at least 16\% more errors}  % experiments/FIGS/data/v2/index.json :: binary probe qcd_c_vs_light: (1 - AUC) of FLAV_F1R over L188, run by run :: sha256 c6c4138ee5c86423 :: prose twin: two significant figures, direction in words
\newcommand{\MlstCmpQcdCFonerOverFullPrMax}{up to 18\% more errors}  % experiments/FIGS/data/v2/index.json :: binary probe qcd_c_vs_light: (1 - AUC) of FLAV_F1R over L188, run by run :: sha256 c6c4138ee5c86423 :: prose twin of the run farther from 1
\newcommand{\MlstLpEemmAucFull}{1.0000\,\ensuremath{\pm}\,\ensuremath{<}0.0001}  % experiments/FIGS/data/v2/index.json :: linear probe ee_vs_mm (fits_v4_sib): AUC, L188 :: sha256 c6c4138ee5c86423
\newcommand{\MlstLpEemmAucRes}{1.0000\,\ensuremath{\pm}\,\ensuremath{<}0.0001}  % experiments/FIGS/data/v2/index.json :: linear probe ee_vs_mm (fits_v4_sib): AUC, L162 :: sha256 c6c4138ee5c86423
\newcommand{\MlstLpEemmAucMid}{0.9610\,\ensuremath{\pm}\,0.0030}  % experiments/FIGS/data/v2/index.json :: linear probe ee_vs_mm (fits_v4_sib): AUC, R42_Q1 :: sha256 c6c4138ee5c86423
\newcommand{\MlstLpEemmAucCoarse}{0.960\,\ensuremath{\pm}\,0.020}  % experiments/FIGS/data/v2/index.json :: linear probe ee_vs_mm (fits_v4_sib): AUC, R16_Q1 :: sha256 c6c4138ee5c86423
\newcommand{\MlstLpEemmAucFzero}{1.0000 and 1.0000}  % experiments/FIGS/data/v2/index.json :: linear probe ee_vs_mm (fits_v4_sib): AUC, FLAV_F0 :: sha256 c6c4138ee5c86423
\newcommand{\MlstLpEemmAucFone}{1.0000 and 1.0000}  % experiments/FIGS/data/v2/index.json :: linear probe ee_vs_mm (fits_v4_sib): AUC, FLAV_F1 :: sha256 c6c4138ee5c86423
\newcommand{\MlstLpEemmAucFoner}{1.0000 and 1.0000}  % experiments/FIGS/data/v2/index.json :: linear probe ee_vs_mm (fits_v4_sib): AUC, FLAV_F1R :: sha256 c6c4138ee5c86423
\newcommand{\MlstLpEemmAucInit}{0.9995\,\ensuremath{\pm}\,0.0003}  % experiments/FIGS/data/v2/index.json :: linear probe ee_vs_mm (fits_v4_sib): AUC, INIT :: sha256 c6c4138ee5c86423
\newcommand{\MlstLpEemmAucInitBn}{1.0000\,\ensuremath{\pm}\,\ensuremath{<}0.0001}  % experiments/FIGS/data/v2/index.json :: linear probe ee_vs_mm (fits_v4_sib): AUC, INIT_BN :: sha256 c6c4138ee5c86423
\newcommand{\MlstLpEemmAucRandMerged}{1.0000\,\ensuremath{\pm}\,\ensuremath{<}0.0001}  % experiments/FIGS/data/v2/index.json :: linear probe ee_vs_mm (fits_v4_sib): AUC, the random partitions that merge the pair :: sha256 c6c4138ee5c86423
\newcommand{\MlstPairEemmNTrain}{9{,}571}  % experiments/FIGS/data/v2/index.json :: fits_v4_sib ee_vs_mm: n_train :: sha256 c6c4138ee5c86423
\newcommand{\MlstPairEemmNVal}{1{,}068}  % experiments/FIGS/data/v2/index.json :: fits_v4_sib ee_vs_mm: n_val :: sha256 c6c4138ee5c86423
\newcommand{\MlstPairEemmNTestPerSide}{128{,}833--135{,}245}  % experiments/FIGS/data/v2/index.json :: fits_v4_sib ee_vs_mm: n_test_signal and n_test_background :: sha256 c6c4138ee5c86423
\newcommand{\MlstLpJcTwoAucFull}{0.9801\,\ensuremath{\pm}\,\ensuremath{<}0.0001}  % experiments/FIGS/data/v2/index.json :: linear probe JetClass-II N=1000000: AUC averaged over the classes, L188 :: sha256 c6c4138ee5c86423
\newcommand{\MlstLpJcTwoAccFull}{0.4348\,\ensuremath{\pm}\,0.0003}  % experiments/FIGS/data/v2/index.json :: linear probe JetClass-II N=1000000: accuracy, L188 :: sha256 c6c4138ee5c86423
\newcommand{\MlstLpJcTwoAucRes}{0.9796\,\ensuremath{\pm}\,0.0001}  % experiments/FIGS/data/v2/index.json :: linear probe JetClass-II N=1000000: AUC averaged over the classes, L162 :: sha256 c6c4138ee5c86423
\newcommand{\MlstLpJcTwoAccRes}{0.4328\,\ensuremath{\pm}\,0.0008}  % experiments/FIGS/data/v2/index.json :: linear probe JetClass-II N=1000000: accuracy, L162 :: sha256 c6c4138ee5c86423
\newcommand{\MlstLpJcTwoAucMid}{0.9756\,\ensuremath{\pm}\,0.0002}  % experiments/FIGS/data/v2/index.json :: linear probe JetClass-II N=1000000: AUC averaged over the classes, R42_Q1 :: sha256 c6c4138ee5c86423
\newcommand{\MlstLpJcTwoAccMid}{0.361\,\ensuremath{\pm}\,0.005}  % experiments/FIGS/data/v2/index.json :: linear probe JetClass-II N=1000000: accuracy, R42_Q1 :: sha256 c6c4138ee5c86423
\newcommand{\MlstLpJcTwoAucCoarse}{0.9691\,\ensuremath{\pm}\,0.0004}  % experiments/FIGS/data/v2/index.json :: linear probe JetClass-II N=1000000: AUC averaged over the classes, R16_Q1 :: sha256 c6c4138ee5c86423
\newcommand{\MlstLpJcTwoAccCoarse}{0.334\,\ensuremath{\pm}\,0.007}  % experiments/FIGS/data/v2/index.json :: linear probe JetClass-II N=1000000: accuracy, R16_Q1 :: sha256 c6c4138ee5c86423
\newcommand{\MlstLpJcTwoAucResMass}{0.9791 and 0.9794}  % experiments/FIGS/data/v2/index.json :: linear probe JetClass-II N=1000000: AUC averaged over the classes, L162_MASS :: sha256 c6c4138ee5c86423
\newcommand{\MlstLpJcTwoAccResMass}{0.4271 and 0.4305}  % experiments/FIGS/data/v2/index.json :: linear probe JetClass-II N=1000000: accuracy, L162_MASS :: sha256 c6c4138ee5c86423
\newcommand{\MlstLpJcTwoAucCoarseMass}{0.9679 and 0.9696}  % experiments/FIGS/data/v2/index.json :: linear probe JetClass-II N=1000000: AUC averaged over the classes, R16_Q1_MASS :: sha256 c6c4138ee5c86423
\newcommand{\MlstLpJcTwoAccCoarseMass}{0.312 and 0.350}  % experiments/FIGS/data/v2/index.json :: linear probe JetClass-II N=1000000: accuracy, R16_Q1_MASS :: sha256 c6c4138ee5c86423
\newcommand{\MlstLpJcTwoAucFzero}{0.9717 and 0.9710}  % experiments/FIGS/data/v2/index.json :: linear probe JetClass-II N=1000000: AUC averaged over the classes, FLAV_F0 :: sha256 c6c4138ee5c86423
\newcommand{\MlstLpJcTwoAccFzero}{0.3769 and 0.3723}  % experiments/FIGS/data/v2/index.json :: linear probe JetClass-II N=1000000: accuracy, FLAV_F0 :: sha256 c6c4138ee5c86423
\newcommand{\MlstLpJcTwoAucFone}{0.9758 and 0.9759}  % experiments/FIGS/data/v2/index.json :: linear probe JetClass-II N=1000000: AUC averaged over the classes, FLAV_F1 :: sha256 c6c4138ee5c86423
\newcommand{\MlstLpJcTwoAccFone}{0.3945 and 0.3963}  % experiments/FIGS/data/v2/index.json :: linear probe JetClass-II N=1000000: accuracy, FLAV_F1 :: sha256 c6c4138ee5c86423
\newcommand{\MlstLpJcTwoAucFoner}{0.9760 and 0.9760}  % experiments/FIGS/data/v2/index.json :: linear probe JetClass-II N=1000000: AUC averaged over the classes, FLAV_F1R :: sha256 c6c4138ee5c86423
\newcommand{\MlstLpJcTwoAccFoner}{0.4071 and 0.4089}  % experiments/FIGS/data/v2/index.json :: linear probe JetClass-II N=1000000: accuracy, FLAV_F1R :: sha256 c6c4138ee5c86423
\newcommand{\MlstLpJcTwoAucInit}{0.8121\,\ensuremath{\pm}\,0.0019}  % experiments/FIGS/data/v2/index.json :: linear probe JetClass-II N=1000000: AUC averaged over the classes, INIT :: sha256 c6c4138ee5c86423
\newcommand{\MlstLpJcTwoAccInit}{0.0724\,\ensuremath{\pm}\,0.0015}  % experiments/FIGS/data/v2/index.json :: linear probe JetClass-II N=1000000: accuracy, INIT :: sha256 c6c4138ee5c86423
\newcommand{\MlstLpJcTwoAucInitBn}{0.866\,\ensuremath{\pm}\,0.004}  % experiments/FIGS/data/v2/index.json :: linear probe JetClass-II N=1000000: AUC averaged over the classes, INIT_BN :: sha256 c6c4138ee5c86423
\newcommand{\MlstLpJcTwoAccInitBn}{0.1017\,\ensuremath{\pm}\,0.0017}  % experiments/FIGS/data/v2/index.json :: linear probe JetClass-II N=1000000: accuracy, INIT_BN :: sha256 c6c4138ee5c86423
\newcommand{\MlstLpJcTwoAucRand}{0.9777\,\ensuremath{\pm}\,0.0002}  % experiments/FIGS/data/v2/index.json :: linear probe JetClass-II N=1000000: AUC averaged over the classes, the random partitions :: sha256 c6c4138ee5c86423
\newcommand{\MlstLpJcOneAccInitBn}{0.475\,\ensuremath{\pm}\,0.005}  % experiments/FIGS/data/v2/index.json :: linear probe JetClass N=1000000: accuracy, INIT_BN :: sha256 c6c4138ee5c86423
\newcommand{\MlstBnMaxDev}{0.2\%}  % experiments/FIGS/data/v2/index.json :: largest |r(best70_bn) / r(best70) - 1| of a JetClass-II ratio of 1 - AUC to the 188-class run, per cent :: sha256 c6c4138ee5c86423
\newcommand{\MlstBnMaxDevPair}{0.7\%}  % experiments/FIGS/data/v2/index.json :: largest |r(best70_bn) / r(best70) - 1| of a ratio of 1 - AUC on a class pair to the 188-class run, levels of the hierarchy, mass models and controlled partitions, per cent :: sha256 c6c4138ee5c86423 :: pairs with a check: bc_vs_bq_cs, bb_vs_cc, bbqq_vs_ccqq, qcd_bb_vs_cc, qcd_b_vs_light, qcd_c_vs_light; without: ee_vs_mm
\newcommand{\MlstGpuTypeBound}{1.4\%}  % experiments/FIGS/data/v2/index.json :: largest |label-set ratio of 1 - AUC of index 3 / mean of the other indices - 1|, JetClass-II largest set, per cent :: sha256 c6c4138ee5c86423 :: mtx-l188-s3 trained on NVIDIA GeForce RTX 3090; largest for R16_Q1; label sets L162 0.09%, R42_Q1 0.70%, R16_Q1 1.37%
\newcommand{\MlstGpuOtherSeedsDiff}{1.9\%}  % experiments/FIGS/data/v2/index.json :: largest (max - min) / mean of a label-set ratio of 1 - AUC over the indices other than 3, JetClass-II largest set, per cent :: sha256 c6c4138ee5c86423 :: label sets L162 0.90%, R42_Q1 1.14%, R16_Q1 1.93%
\newcommand{\MlstPerClassResNDoubled}{0}  % experiments/FIGS/data/v2/index.json :: JetClass-II classes whose mean ratio of 1 - AUC (L162 over 188) is at least 2 (count) :: sha256 c6c4138ee5c86423
\newcommand{\MlstPerClassResSeedSdMedian}{0.011}  % experiments/FIGS/data/v2/index.json :: median over classes of the SD over runs of the per-class ratio, L162 over 188 :: sha256 c6c4138ee5c86423
\newcommand{\MlstPerClassResTopClass}{QCD, \ensuremath{ccss}}  % experiments/FIGS/data/v2/index.json :: class with the largest mean ratio, L162 :: sha256 c6c4138ee5c86423
\newcommand{\MlstPerClassResTopRatio}{1.129\,\ensuremath{\pm}\,0.010}  % experiments/FIGS/data/v2/index.json :: per-class ratio of 1 - AUC of that class, L162 over 188, run by run :: sha256 c6c4138ee5c86423
\newcommand{\MlstPerClassResTopRatioPr}{13\% more errors}  % experiments/FIGS/data/v2/index.json :: per-class ratio of 1 - AUC of that class, L162 over 188, run by run :: sha256 c6c4138ee5c86423 :: prose twin: two significant figures, direction in words
\newcommand{\MlstPerClassResLeptonN}{0}  % experiments/FIGS/data/v2/index.json :: resonant classes sharing their L162 group with a decay that differs only in lepton flavor (count) :: sha256 c6c4138ee5c86423
\newcommand{\MlstPerClassResLeptonMedian}{--}  % experiments/FIGS/data/v2/index.json :: median over those classes of the mean per-class ratio, L162 over 188 :: sha256 c6c4138ee5c86423
\newcommand{\MlstPerClassResLeptonMedianPr}{\pending{no run yet}}  % experiments/FIGS/data/v2/index.json :: median over those classes of the mean per-class ratio, L162 over 188 :: sha256 c6c4138ee5c86423 :: prose twin: two significant figures, direction in words
\newcommand{\MlstPerClassResLeptonShare}{--}  % experiments/FIGS/data/v2/index.json :: their share of the summed ln ratio over the resonant classes, L162 :: sha256 c6c4138ee5c86423
\newcommand{\MlstPerClassResQuarkN}{0}  % experiments/FIGS/data/v2/index.json :: resonant classes sharing their L162 group with a decay that differs only in quark flavor (count) :: sha256 c6c4138ee5c86423
\newcommand{\MlstPerClassResQuarkMedian}{--}  % experiments/FIGS/data/v2/index.json :: median over those classes of the mean per-class ratio, L162 over 188 :: sha256 c6c4138ee5c86423
\newcommand{\MlstPerClassResQuarkMedianPr}{\pending{no run yet}}  % experiments/FIGS/data/v2/index.json :: median over those classes of the mean per-class ratio, L162 over 188 :: sha256 c6c4138ee5c86423 :: prose twin: two significant figures, direction in words
\newcommand{\MlstPerClassResQuarkShare}{--}  % experiments/FIGS/data/v2/index.json :: their share of the summed ln ratio over the resonant classes, L162 :: sha256 c6c4138ee5c86423
\newcommand{\MlstPerClassResQuarkLeptonN}{0}  % experiments/FIGS/data/v2/index.json :: resonant classes counted in both the quark- and the lepton-flavor set, L162 (count) :: sha256 c6c4138ee5c86423
\newcommand{\MlstPerClassMidNDoubled}{40}  % experiments/FIGS/data/v2/index.json :: JetClass-II classes whose mean ratio of 1 - AUC (R42_Q1 over 188) is at least 2 (count) :: sha256 c6c4138ee5c86423
\newcommand{\MlstPerClassMidSeedSdMedian}{0.023}  % experiments/FIGS/data/v2/index.json :: median over classes of the SD over runs of the per-class ratio, R42_Q1 over 188 :: sha256 c6c4138ee5c86423
\newcommand{\MlstPerClassMidTopClass}{\ensuremath{X\to YY\to bb\mu\mu}}  % experiments/FIGS/data/v2/index.json :: class with the largest mean ratio, R42_Q1 :: sha256 c6c4138ee5c86423
\newcommand{\MlstPerClassMidTopRatio}{4.06\,\ensuremath{\pm}\,0.14}  % experiments/FIGS/data/v2/index.json :: per-class ratio of 1 - AUC of that class, R42_Q1 over 188, run by run :: sha256 c6c4138ee5c86423
\newcommand{\MlstPerClassMidTopRatioPr}{4.1 times the errors}  % experiments/FIGS/data/v2/index.json :: per-class ratio of 1 - AUC of that class, R42_Q1 over 188, run by run :: sha256 c6c4138ee5c86423 :: prose twin: two significant figures, direction in words
\newcommand{\MlstPerClassMidLeptonN}{78}  % experiments/FIGS/data/v2/index.json :: resonant classes sharing their R42_Q1 group with a decay that differs only in lepton flavor (count) :: sha256 c6c4138ee5c86423
\newcommand{\MlstPerClassMidLeptonMedian}{2.02}  % experiments/FIGS/data/v2/index.json :: median over those classes of the mean per-class ratio, R42_Q1 over 188 :: sha256 c6c4138ee5c86423
\newcommand{\MlstPerClassMidLeptonMedianPr}{2.0 times the errors}  % experiments/FIGS/data/v2/index.json :: median over those classes of the mean per-class ratio, R42_Q1 over 188 :: sha256 c6c4138ee5c86423 :: prose twin: two significant figures, direction in words
\newcommand{\MlstPerClassMidLeptonShare}{90\%}  % experiments/FIGS/data/v2/index.json :: their share of the summed ln ratio over the resonant classes, R42_Q1 :: sha256 c6c4138ee5c86423
\newcommand{\MlstPerClassMidQuarkN}{111}  % experiments/FIGS/data/v2/index.json :: resonant classes sharing their R42_Q1 group with a decay that differs only in quark flavor (count) :: sha256 c6c4138ee5c86423
\newcommand{\MlstPerClassMidQuarkMedian}{1.14}  % experiments/FIGS/data/v2/index.json :: median over those classes of the mean per-class ratio, R42_Q1 over 188 :: sha256 c6c4138ee5c86423
\newcommand{\MlstPerClassMidQuarkMedianPr}{14\% more errors}  % experiments/FIGS/data/v2/index.json :: median over those classes of the mean per-class ratio, R42_Q1 over 188 :: sha256 c6c4138ee5c86423 :: prose twin: two significant figures, direction in words
\newcommand{\MlstPerClassMidQuarkShare}{53\%}  % experiments/FIGS/data/v2/index.json :: their share of the summed ln ratio over the resonant classes, R42_Q1 :: sha256 c6c4138ee5c86423
\newcommand{\MlstPerClassMidQuarkLeptonN}{44}  % experiments/FIGS/data/v2/index.json :: resonant classes counted in both the quark- and the lepton-flavor set, R42_Q1 (count) :: sha256 c6c4138ee5c86423
\newcommand{\MlstPerClassCoarseNDoubled}{60}  % experiments/FIGS/data/v2/index.json :: JetClass-II classes whose mean ratio of 1 - AUC (R16_Q1 over 188) is at least 2 (count) :: sha256 c6c4138ee5c86423
\newcommand{\MlstPerClassCoarseSeedSdMedian}{0.055}  % experiments/FIGS/data/v2/index.json :: median over classes of the SD over runs of the per-class ratio, R16_Q1 over 188 :: sha256 c6c4138ee5c86423
\newcommand{\MlstPerClassCoarseTopClass}{\ensuremath{X\to YY\to bbee}}  % experiments/FIGS/data/v2/index.json :: class with the largest mean ratio, R16_Q1 :: sha256 c6c4138ee5c86423
\newcommand{\MlstPerClassCoarseTopRatio}{5.2\,\ensuremath{\pm}\,0.4}  % experiments/FIGS/data/v2/index.json :: per-class ratio of 1 - AUC of that class, R16_Q1 over 188, run by run :: sha256 c6c4138ee5c86423
\newcommand{\MlstPerClassCoarseTopRatioPr}{5.2 times the errors}  % experiments/FIGS/data/v2/index.json :: per-class ratio of 1 - AUC of that class, R16_Q1 over 188, run by run :: sha256 c6c4138ee5c86423 :: prose twin: two significant figures, direction in words
\newcommand{\MlstPerClassCoarseLeptonN}{78}  % experiments/FIGS/data/v2/index.json :: resonant classes sharing their R16_Q1 group with a decay that differs only in lepton flavor (count) :: sha256 c6c4138ee5c86423
\newcommand{\MlstPerClassCoarseLeptonMedian}{2.27}  % experiments/FIGS/data/v2/index.json :: median over those classes of the mean per-class ratio, R16_Q1 over 188 :: sha256 c6c4138ee5c86423
\newcommand{\MlstPerClassCoarseLeptonMedianPr}{2.3 times the errors}  % experiments/FIGS/data/v2/index.json :: median over those classes of the mean per-class ratio, R16_Q1 over 188 :: sha256 c6c4138ee5c86423 :: prose twin: two significant figures, direction in words
\newcommand{\MlstPerClassCoarseLeptonShare}{68\%}  % experiments/FIGS/data/v2/index.json :: their share of the summed ln ratio over the resonant classes, R16_Q1 :: sha256 c6c4138ee5c86423
\newcommand{\MlstPerClassCoarseQuarkN}{141}  % experiments/FIGS/data/v2/index.json :: resonant classes sharing their R16_Q1 group with a decay that differs only in quark flavor (count) :: sha256 c6c4138ee5c86423
\newcommand{\MlstPerClassCoarseQuarkMedian}{1.65}  % experiments/FIGS/data/v2/index.json :: median over those classes of the mean per-class ratio, R16_Q1 over 188 :: sha256 c6c4138ee5c86423
\newcommand{\MlstPerClassCoarseQuarkMedianPr}{65\% more errors}  % experiments/FIGS/data/v2/index.json :: median over those classes of the mean per-class ratio, R16_Q1 over 188 :: sha256 c6c4138ee5c86423 :: prose twin: two significant figures, direction in words
\newcommand{\MlstPerClassCoarseQuarkShare}{87\%}  % experiments/FIGS/data/v2/index.json :: their share of the summed ln ratio over the resonant classes, R16_Q1 :: sha256 c6c4138ee5c86423
\newcommand{\MlstPerClassCoarseQuarkLeptonN}{64}  % experiments/FIGS/data/v2/index.json :: resonant classes counted in both the quark- and the lepton-flavor set, R16_Q1 (count) :: sha256 c6c4138ee5c86423
\newcommand{\MlstPerClassRandFlavMergedOverApart}{1.010\,\ensuremath{\pm}\,0.008}  % experiments/FIGS/data/v2/index.json :: random partitions, linear probe: median per-class ratio of the classes merged with a decay that differs only in quark flavor over that of the classes kept apart from all of them, mean +- SD over the partitions :: sha256 c6c4138ee5c86423
\newcommand{\MlstPerClassRandFlavMergedOverApartPr}{equal within 1\%}  % experiments/FIGS/data/v2/index.json :: random partitions, linear probe: median per-class ratio of the classes merged with a decay that differs only in quark flavor over that of the classes kept apart from all of them, mean +- SD over the partitions :: sha256 c6c4138ee5c86423 :: prose twin: two significant figures, direction in words
\newcommand{\MlstPerClassRandLeptonMergedOverApart}{0.997\,\ensuremath{\pm}\,0.020}  % experiments/FIGS/data/v2/index.json :: random partitions, linear probe: median per-class ratio of the classes merged with a decay that differs only in lepton flavor over that of the classes kept apart from all of them, mean +- SD over the partitions :: sha256 c6c4138ee5c86423
\newcommand{\MlstPerClassRandLeptonMergedOverApartPr}{equal within 3\%}  % experiments/FIGS/data/v2/index.json :: random partitions, linear probe: median per-class ratio of the classes merged with a decay that differs only in lepton flavor over that of the classes kept apart from all of them, mean +- SD over the partitions :: sha256 c6c4138ee5c86423 :: prose twin: two significant figures, direction in words
\newcommand{\MlstMlpNatPerClassRandFlavMergedOverApart}{1.003\,\ensuremath{\pm}\,0.007}  % experiments/FIGS/data/v2/index.json :: random partitions, nonlinear probe, natural composition: median per-class ratio of the classes merged with a decay that differs only in quark flavor over that of the classes kept apart from all of them, mean +- SD over the partitions :: sha256 c6c4138ee5c86423
\newcommand{\MlstMlpNatPerClassRandFlavMergedOverApartPr}{equal within 1\%}  % experiments/FIGS/data/v2/index.json :: random partitions, nonlinear probe, natural composition: median per-class ratio of the classes merged with a decay that differs only in quark flavor over that of the classes kept apart from all of them, mean +- SD over the partitions :: sha256 c6c4138ee5c86423 :: prose twin: two significant figures, direction in words
\newcommand{\MlstMlpNatPerClassRandLeptonMergedOverApart}{1.01\,\ensuremath{\pm}\,0.07}  % experiments/FIGS/data/v2/index.json :: random partitions, nonlinear probe, natural composition: median per-class ratio of the classes merged with a decay that differs only in lepton flavor over that of the classes kept apart from all of them, mean +- SD over the partitions :: sha256 c6c4138ee5c86423
\newcommand{\MlstMlpNatPerClassRandLeptonMergedOverApartPr}{equal within 8\%}  % experiments/FIGS/data/v2/index.json :: random partitions, nonlinear probe, natural composition: median per-class ratio of the classes merged with a decay that differs only in lepton flavor over that of the classes kept apart from all of them, mean +- SD over the partitions :: sha256 c6c4138ee5c86423 :: prose twin: two significant figures, direction in words
\newcommand{\MlstPerClassAucDiffResSeedSdMedian}{0.015}  % experiments/FIGS/data/v2/index.json :: median over classes of the SD over indices of 100 (AUC_c of L162 - AUC_c of 188), pp :: sha256 c6c4138ee5c86423
\newcommand{\MlstPerClassAucDiffResMin}{\ensuremath{-}0.73}  % experiments/FIGS/data/v2/index.json :: most negative mean over indices of 100 (AUC_c of L162 - AUC_c of 188), pp :: sha256 c6c4138ee5c86423 :: label_QCD_bcss
\newcommand{\MlstPerClassAucDiffResMinClass}{QCD, \ensuremath{bcss}}  % experiments/FIGS/data/v2/index.json :: its class, L162 :: sha256 c6c4138ee5c86423
\newcommand{\MlstPerClassAucDiffMidSeedSdMedian}{0.032}  % experiments/FIGS/data/v2/index.json :: median over classes of the SD over indices of 100 (AUC_c of R42_Q1 - AUC_c of 188), pp :: sha256 c6c4138ee5c86423
\newcommand{\MlstPerClassAucDiffMidMin}{\ensuremath{-}1.4}  % experiments/FIGS/data/v2/index.json :: most negative mean over indices of 100 (AUC_c of R42_Q1 - AUC_c of 188), pp :: sha256 c6c4138ee5c86423 :: label_X_YY_cqtauev
\newcommand{\MlstPerClassAucDiffMidMinClass}{\ensuremath{X\to YY\to cq\tau_e\nu}}  % experiments/FIGS/data/v2/index.json :: its class, R42_Q1 :: sha256 c6c4138ee5c86423
\newcommand{\MlstPerClassAucDiffCoarseSeedSdMedian}{0.077}  % experiments/FIGS/data/v2/index.json :: median over classes of the SD over indices of 100 (AUC_c of R16_Q1 - AUC_c of 188), pp :: sha256 c6c4138ee5c86423
\newcommand{\MlstPerClassAucDiffCoarseMin}{\ensuremath{-}3.1}  % experiments/FIGS/data/v2/index.json :: most negative mean over indices of 100 (AUC_c of R16_Q1 - AUC_c of 188), pp :: sha256 c6c4138ee5c86423 :: label_QCD_bcss
\newcommand{\MlstPerClassAucDiffCoarseMinClass}{QCD, \ensuremath{bcss}}  % experiments/FIGS/data/v2/index.json :: its class, R16_Q1 :: sha256 c6c4138ee5c86423
\newcommand{\MlstRocJcTwoHurtList}{\ensuremath{X\to YY\to ccq}; \ensuremath{X\to YY\to ssb}; \ensuremath{X\to YY\to ccss}; QCD, \ensuremath{bcc}; QCD, \ensuremath{bccss}; QCD, \ensuremath{bss}; QCD, \ensuremath{ccs} and QCD, \ensuremath{ccss}}  % experiments/FIGS/data/v2/index.json :: jc2_roc_classes on the class-balanced probes; candidates with at least 1000 test jets :: sha256 c6c4138ee5c86423 :: label_X_YY_ccq, label_X_YY_ssb, label_X_YY_ccss, label_QCD_bcc, label_QCD_bccss, label_QCD_bss, label_QCD_ccs, label_QCD_ccss
\newcommand{\MlstRocJcTwoStableList}{\ensuremath{X\to gg}; \ensuremath{X\to YY\to ggg}; \ensuremath{X\to YY\to g\tau_h\tau_h} and QCD, light}  % experiments/FIGS/data/v2/index.json :: jc2_roc_classes on the class-balanced probes; candidates with at least 1000 test jets :: sha256 c6c4138ee5c86423 :: label_X_gg, label_X_YY_ggg, label_X_YY_gtauhtauh, label_QCD_light
\newcommand{\MlstRocJcTwoHurtMinDrop}{2.1}  % experiments/FIGS/data/v2/index.json :: jc2_roc_classes on the class-balanced probes; candidates with at least 1000 test jets :: sha256 c6c4138ee5c86423 :: smallest AUC loss at 17 classes of the hurt, pp
\newcommand{\MlstRocJcTwoHurtMaxDrop}{2.8}  % experiments/FIGS/data/v2/index.json :: jc2_roc_classes on the class-balanced probes; candidates with at least 1000 test jets :: sha256 c6c4138ee5c86423 :: largest AUC loss at 17 classes of the hurt, pp
\newcommand{\MlstRocJcTwoStableMaxChange}{15}  % experiments/FIGS/data/v2/index.json :: jc2_roc_classes on the class-balanced probes; candidates with at least 1000 test jets :: sha256 c6c4138ee5c86423 :: largest change of 1 - AUC at any label set of the stable classes, per cent of the 188-class value
\newcommand{\MlstJcBalNThin}{4}  % experiments/FIGS/data/v2/jc2bal/census.json :: counts.ft_test :: sha256 a4379ee0817254cf :: classes with fewer than 1000 class-balanced test jets (F3 hatching)
\newcommand{\MlstMinBkgPass}{20}  % experiments/FIGS/mlst_v4.py :: JC1_MIN_BKG :: sha256 fa536b4866ffd651 :: fewest background jets behind a drawn ROC point or a quoted rejection
\newcommand{\MlstMassMinJets}{100}  % experiments/EVAL/atlas_resolution.py :: MIN_JETS :: sha256 b82cc219f79f0f14 :: jets below which a mass bin is null
\newcommand{\MlstRocCutBbFull}{0.79}  % experiments/FIGS/data/v2/index.json :: curves_v2/pairs curve bb_vs_cc: largest cut.eps_s over the L188 runs ('not reached' when a run's cut is null or leaves fewer than 2 grid points) :: sha256 c6c4138ee5c86423
\newcommand{\MlstRocCutBbRes}{0.77}  % experiments/FIGS/data/v2/index.json :: curves_v2/pairs curve bb_vs_cc: largest cut.eps_s over the L162 runs ('not reached' when a run's cut is null or leaves fewer than 2 grid points) :: sha256 c6c4138ee5c86423
\newcommand{\MlstRocCutBbMid}{0.79}  % experiments/FIGS/data/v2/index.json :: curves_v2/pairs curve bb_vs_cc: largest cut.eps_s over the R42_Q1 runs ('not reached' when a run's cut is null or leaves fewer than 2 grid points) :: sha256 c6c4138ee5c86423
\newcommand{\MlstRocCutBbCoarse}{0.13}  % experiments/FIGS/data/v2/index.json :: curves_v2/pairs curve bb_vs_cc: largest cut.eps_s over the R16_Q1 runs ('not reached' when a run's cut is null or leaves fewer than 2 grid points) :: sha256 c6c4138ee5c86423
\newcommand{\MlstRocCutBbFzero}{0.03}  % experiments/FIGS/data/v2/index.json :: curves_v2/pairs curve bb_vs_cc: largest cut.eps_s over the FLAV_F0 runs ('not reached' when a run's cut is null or leaves fewer than 2 grid points) :: sha256 c6c4138ee5c86423
\newcommand{\MlstRocCutBbFone}{0.62}  % experiments/FIGS/data/v2/index.json :: curves_v2/pairs curve bb_vs_cc: largest cut.eps_s over the FLAV_F1 runs ('not reached' when a run's cut is null or leaves fewer than 2 grid points) :: sha256 c6c4138ee5c86423
\newcommand{\MlstRocCutBbFoner}{0.51}  % experiments/FIGS/data/v2/index.json :: curves_v2/pairs curve bb_vs_cc: largest cut.eps_s over the FLAV_F1R runs ('not reached' when a run's cut is null or leaves fewer than 2 grid points) :: sha256 c6c4138ee5c86423
\newcommand{\MlstRocNBkgBb}{200{,}000}  % experiments/FIGS/data/v2/index.json :: curves_v2/pairs curve bb_vs_cc: n_background :: sha256 c6c4138ee5c86423
\newcommand{\MlstRocRejMaxBb}{10{,}000}  % experiments/FIGS/data/v2/index.json :: curves_v2/pairs curve bb_vs_cc: n_background / 20 :: sha256 c6c4138ee5c86423
\newcommand{\MlstRocCutBbqqFull}{0.67}  % experiments/FIGS/data/v2/index.json :: curves_v2/pairs curve bbqq_vs_ccqq: largest cut.eps_s over the L188 runs ('not reached' when a run's cut is null or leaves fewer than 2 grid points) :: sha256 c6c4138ee5c86423
\newcommand{\MlstRocCutBbqqRes}{0.66}  % experiments/FIGS/data/v2/index.json :: curves_v2/pairs curve bbqq_vs_ccqq: largest cut.eps_s over the L162 runs ('not reached' when a run's cut is null or leaves fewer than 2 grid points) :: sha256 c6c4138ee5c86423
\newcommand{\MlstRocCutBbqqMid}{0.62}  % experiments/FIGS/data/v2/index.json :: curves_v2/pairs curve bbqq_vs_ccqq: largest cut.eps_s over the R42_Q1 runs ('not reached' when a run's cut is null or leaves fewer than 2 grid points) :: sha256 c6c4138ee5c86423
\newcommand{\MlstRocCutBbqqCoarse}{0.06}  % experiments/FIGS/data/v2/index.json :: curves_v2/pairs curve bbqq_vs_ccqq: largest cut.eps_s over the R16_Q1 runs ('not reached' when a run's cut is null or leaves fewer than 2 grid points) :: sha256 c6c4138ee5c86423
\newcommand{\MlstRocCutBbqqFzero}{0.06}  % experiments/FIGS/data/v2/index.json :: curves_v2/pairs curve bbqq_vs_ccqq: largest cut.eps_s over the FLAV_F0 runs ('not reached' when a run's cut is null or leaves fewer than 2 grid points) :: sha256 c6c4138ee5c86423
\newcommand{\MlstRocCutBbqqFone}{0.56}  % experiments/FIGS/data/v2/index.json :: curves_v2/pairs curve bbqq_vs_ccqq: largest cut.eps_s over the FLAV_F1 runs ('not reached' when a run's cut is null or leaves fewer than 2 grid points) :: sha256 c6c4138ee5c86423
\newcommand{\MlstRocCutBbqqFoner}{0.32}  % experiments/FIGS/data/v2/index.json :: curves_v2/pairs curve bbqq_vs_ccqq: largest cut.eps_s over the FLAV_F1R runs ('not reached' when a run's cut is null or leaves fewer than 2 grid points) :: sha256 c6c4138ee5c86423
\newcommand{\MlstRocNBkgBbqq}{200{,}000}  % experiments/FIGS/data/v2/index.json :: curves_v2/pairs curve bbqq_vs_ccqq: n_background :: sha256 c6c4138ee5c86423
\newcommand{\MlstRocRejMaxBbqq}{10{,}000}  % experiments/FIGS/data/v2/index.json :: curves_v2/pairs curve bbqq_vs_ccqq: n_background / 20 :: sha256 c6c4138ee5c86423
\newcommand{\MlstRocCutBcFull}{0.24}  % experiments/FIGS/data/v2/index.json :: curves_v2/pairs curve bc_vs_bq_cs: largest cut.eps_s over the L188 runs ('not reached' when a run's cut is null or leaves fewer than 2 grid points) :: sha256 c6c4138ee5c86423
\newcommand{\MlstRocCutBcRes}{0.25}  % experiments/FIGS/data/v2/index.json :: curves_v2/pairs curve bc_vs_bq_cs: largest cut.eps_s over the L162 runs ('not reached' when a run's cut is null or leaves fewer than 2 grid points) :: sha256 c6c4138ee5c86423
\newcommand{\MlstRocCutBcMid}{0.14}  % experiments/FIGS/data/v2/index.json :: curves_v2/pairs curve bc_vs_bq_cs: largest cut.eps_s over the R42_Q1 runs ('not reached' when a run's cut is null or leaves fewer than 2 grid points) :: sha256 c6c4138ee5c86423
\newcommand{\MlstRocCutBcCoarse}{0.00}  % experiments/FIGS/data/v2/index.json :: curves_v2/pairs curve bc_vs_bq_cs: largest cut.eps_s over the R16_Q1 runs ('not reached' when a run's cut is null or leaves fewer than 2 grid points) :: sha256 c6c4138ee5c86423
\newcommand{\MlstRocCutBcFzero}{0.00}  % experiments/FIGS/data/v2/index.json :: curves_v2/pairs curve bc_vs_bq_cs: largest cut.eps_s over the FLAV_F0 runs ('not reached' when a run's cut is null or leaves fewer than 2 grid points) :: sha256 c6c4138ee5c86423
\newcommand{\MlstRocCutBcFone}{0.01}  % experiments/FIGS/data/v2/index.json :: curves_v2/pairs curve bc_vs_bq_cs: largest cut.eps_s over the FLAV_F1 runs ('not reached' when a run's cut is null or leaves fewer than 2 grid points) :: sha256 c6c4138ee5c86423
\newcommand{\MlstRocCutBcFoner}{0.01}  % experiments/FIGS/data/v2/index.json :: curves_v2/pairs curve bc_vs_bq_cs: largest cut.eps_s over the FLAV_F1R runs ('not reached' when a run's cut is null or leaves fewer than 2 grid points) :: sha256 c6c4138ee5c86423
\newcommand{\MlstRocNBkgBc}{200{,}000}  % experiments/FIGS/data/v2/index.json :: curves_v2/pairs curve bc_vs_bq_cs: n_background :: sha256 c6c4138ee5c86423
\newcommand{\MlstRocRejMaxBc}{10{,}000}  % experiments/FIGS/data/v2/index.json :: curves_v2/pairs curve bc_vs_bq_cs: n_background / 20 :: sha256 c6c4138ee5c86423
\newcommand{\MlstRocCutEemmFull}{not reached}  % experiments/FIGS/data/v2/index.json :: curves_v2/pairs curve ee_vs_mm: largest cut.eps_s over the L188 runs ('not reached' when a run's cut is null or leaves fewer than 2 grid points) :: sha256 c6c4138ee5c86423
\newcommand{\MlstRocCutEemmRes}{not reached}  % experiments/FIGS/data/v2/index.json :: curves_v2/pairs curve ee_vs_mm: largest cut.eps_s over the L162 runs ('not reached' when a run's cut is null or leaves fewer than 2 grid points) :: sha256 c6c4138ee5c86423
\newcommand{\MlstRocCutEemmMid}{0.00}  % experiments/FIGS/data/v2/index.json :: curves_v2/pairs curve ee_vs_mm: largest cut.eps_s over the R42_Q1 runs ('not reached' when a run's cut is null or leaves fewer than 2 grid points) :: sha256 c6c4138ee5c86423
\newcommand{\MlstRocCutEemmCoarse}{0.00}  % experiments/FIGS/data/v2/index.json :: curves_v2/pairs curve ee_vs_mm: largest cut.eps_s over the R16_Q1 runs ('not reached' when a run's cut is null or leaves fewer than 2 grid points) :: sha256 c6c4138ee5c86423
\newcommand{\MlstRocCutEemmFzero}{not reached}  % experiments/FIGS/data/v2/index.json :: curves_v2/pairs curve ee_vs_mm: largest cut.eps_s over the FLAV_F0 runs ('not reached' when a run's cut is null or leaves fewer than 2 grid points) :: sha256 c6c4138ee5c86423
\newcommand{\MlstRocCutEemmFone}{not reached}  % experiments/FIGS/data/v2/index.json :: curves_v2/pairs curve ee_vs_mm: largest cut.eps_s over the FLAV_F1 runs ('not reached' when a run's cut is null or leaves fewer than 2 grid points) :: sha256 c6c4138ee5c86423
\newcommand{\MlstRocCutEemmFoner}{not reached}  % experiments/FIGS/data/v2/index.json :: curves_v2/pairs curve ee_vs_mm: largest cut.eps_s over the FLAV_F1R runs ('not reached' when a run's cut is null or leaves fewer than 2 grid points) :: sha256 c6c4138ee5c86423
\newcommand{\MlstRocNBkgEemm}{128{,}833}  % experiments/FIGS/data/v2/index.json :: curves_v2/pairs curve ee_vs_mm: n_background :: sha256 c6c4138ee5c86423
\newcommand{\MlstRocRejMaxEemm}{6{,}442}  % experiments/FIGS/data/v2/index.json :: curves_v2/pairs curve ee_vs_mm: n_background / 20 :: sha256 c6c4138ee5c86423
\newcommand{\MlstRocNoCurveEemm}{the \Kfull- and \Kres-class models and the three random partitions with controlled flavor boundaries}  % experiments/FIGS/data/v2/index.json :: curves_v2/pairs curve ee_vs_mm: configurations of F4 with a run whose cut is null or leaves fewer than 2 grid points (F4 draws no curve for them) :: sha256 c6c4138ee5c86423 :: L188, L162, FLAV_F0, FLAV_F1, FLAV_F1R
\newcommand{\MlstRocCutJcOneHbbFull}{0.40}  % experiments/FIGS/data/v2/index.json :: curves_v2/jc1_ft curve label_Hbb: largest cut.eps_s over the 188-class runs ('not reached' when a run's cut is null or leaves fewer than 2 grid points) :: sha256 c6c4138ee5c86423
\newcommand{\MlstRocCutJcOneHccFull}{0.37}  % experiments/FIGS/data/v2/index.json :: curves_v2/jc1_ft curve label_Hcc: largest cut.eps_s over the 188-class runs ('not reached' when a run's cut is null or leaves fewer than 2 grid points) :: sha256 c6c4138ee5c86423
\newcommand{\MlstRocCutJcOneHggFull}{0.04}  % experiments/FIGS/data/v2/index.json :: curves_v2/jc1_ft curve label_Hgg: largest cut.eps_s over the 188-class runs ('not reached' when a run's cut is null or leaves fewer than 2 grid points) :: sha256 c6c4138ee5c86423
\newcommand{\MlstRocCutJcOneHfourqFull}{0.26}  % experiments/FIGS/data/v2/index.json :: curves_v2/jc1_ft curve label_H4q: largest cut.eps_s over the 188-class runs ('not reached' when a run's cut is null or leaves fewer than 2 grid points) :: sha256 c6c4138ee5c86423
\newcommand{\MlstRocCutJcOneHqqlFull}{0.98}  % experiments/FIGS/data/v2/index.json :: curves_v2/jc1_ft curve label_Hqql: largest cut.eps_s over the 188-class runs ('not reached' when a run's cut is null or leaves fewer than 2 grid points) :: sha256 c6c4138ee5c86423
\newcommand{\MlstRocCutJcOneTbqqFull}{0.53}  % experiments/FIGS/data/v2/index.json :: curves_v2/jc1_ft curve label_Tbqq: largest cut.eps_s over the 188-class runs ('not reached' when a run's cut is null or leaves fewer than 2 grid points) :: sha256 c6c4138ee5c86423
\newcommand{\MlstRocCutJcOneTblFull}{not reached}  % experiments/FIGS/data/v2/index.json :: curves_v2/jc1_ft curve label_Tbl: largest cut.eps_s over the 188-class runs ('not reached' when a run's cut is null or leaves fewer than 2 grid points) :: sha256 c6c4138ee5c86423
\newcommand{\MlstRocCutJcOneWqqFull}{0.15}  % experiments/FIGS/data/v2/index.json :: curves_v2/jc1_ft curve label_Wqq: largest cut.eps_s over the 188-class runs ('not reached' when a run's cut is null or leaves fewer than 2 grid points) :: sha256 c6c4138ee5c86423
\newcommand{\MlstRocCutJcOneZqqFull}{0.12}  % experiments/FIGS/data/v2/index.json :: curves_v2/jc1_ft curve label_Zqq: largest cut.eps_s over the 188-class runs ('not reached' when a run's cut is null or leaves fewer than 2 grid points) :: sha256 c6c4138ee5c86423
\newcommand{\MlstFtTopRfiftyCoarseOverFull}{0.932\,\ensuremath{\pm}\,0.018}  % experiments/FIGS/data/v2/index.json :: curves_v2/top_ft curve top: rejection at eps_S = 0.5 of R16_Q1 over L188, run by run :: sha256 c6c4138ee5c86423
\newcommand{\MlstFtTopRfiftyCoarseOverFullPr}{6.8\% less rejection}  % experiments/FIGS/data/v2/index.json :: curves_v2/top_ft curve top: rejection at eps_S = 0.5 of R16_Q1 over L188, run by run :: sha256 c6c4138ee5c86423 :: prose twin: two significant figures, direction in words
\newcommand{\MlstProbeUnconverged}{0}  % experiments/FIGS/data/v2/index.json :: linear probes of pretrained models that did not converge (count) :: sha256 c6c4138ee5c86423
\newcommand{\MlstProbeUnconvergedInit}{0}  % experiments/FIGS/data/v2/index.json :: linear probes of the randomly initialized networks that did not converge (count) :: sha256 c6c4138ee5c86423
\newcommand{\MlstProbeFits}{895}  % experiments/FIGS/data/v2/index.json :: selected linear-probe fits of the reported checkpoints (count) :: sha256 c6c4138ee5c86423
\newcommand{\MlstProbeEdgeC}{16}  % experiments/FIGS/data/v2/index.json :: of these, fits whose selected C is the largest of the grid (count) :: sha256 c6c4138ee5c86423
\newcommand{\MlstMlpUnconverged}{0}  % experiments/FIGS/data/v2/index.json :: nonlinear probes that did not reach their plateau within the epoch cap (count) :: sha256 c6c4138ee5c86423
\newcommand{\MlstMlpCmpJcTwoResOverFull}{1.019\,\ensuremath{\pm}\,0.007}  % experiments/FIGS/data/v2/index.json :: nonlinear probe jc2_188: (1 - AUC) of L162 over L188, run by run :: sha256 c6c4138ee5c86423
\newcommand{\MlstMlpCmpJcTwoResOverFullPr}{1.9\% more errors}  % experiments/FIGS/data/v2/index.json :: nonlinear probe jc2_188: (1 - AUC) of L162 over L188, run by run :: sha256 c6c4138ee5c86423 :: prose twin: two significant figures, direction in words
\newcommand{\MlstMlpCmpJcTwoMidOverFull}{1.170\,\ensuremath{\pm}\,0.006}  % experiments/FIGS/data/v2/index.json :: nonlinear probe jc2_188: (1 - AUC) of R42_Q1 over L188, run by run :: sha256 c6c4138ee5c86423
\newcommand{\MlstMlpCmpJcTwoMidOverFullPr}{17\% more errors}  % experiments/FIGS/data/v2/index.json :: nonlinear probe jc2_188: (1 - AUC) of R42_Q1 over L188, run by run :: sha256 c6c4138ee5c86423 :: prose twin: two significant figures, direction in words
\newcommand{\MlstMlpCmpJcTwoCoarseOverFull}{1.388\,\ensuremath{\pm}\,0.021}  % experiments/FIGS/data/v2/index.json :: nonlinear probe jc2_188: (1 - AUC) of R16_Q1 over L188, run by run :: sha256 c6c4138ee5c86423
\newcommand{\MlstMlpCmpJcTwoCoarseOverFullPr}{39\% more errors}  % experiments/FIGS/data/v2/index.json :: nonlinear probe jc2_188: (1 - AUC) of R16_Q1 over L188, run by run :: sha256 c6c4138ee5c86423 :: prose twin: two significant figures, direction in words
\newcommand{\MlstMlpCmpJcTwoResMassOverRes}{\pending{no run yet}}  % experiments/FIGS/data/v2/index.json :: nonlinear probe jc2_188: (1 - AUC) of L162_MASS over L162, run by run :: sha256 c6c4138ee5c86423
\newcommand{\MlstMlpCmpJcTwoResMassOverResPr}{\pending{no run yet}}  % experiments/FIGS/data/v2/index.json :: nonlinear probe jc2_188: (1 - AUC) of L162_MASS over L162, run by run :: sha256 c6c4138ee5c86423 :: prose twin: two significant figures, direction in words
\newcommand{\MlstMlpCmpJcTwoResMassOverResPrMax}{\pending{no run yet}}  % experiments/FIGS/data/v2/index.json :: nonlinear probe jc2_188: (1 - AUC) of L162_MASS over L162, run by run :: sha256 c6c4138ee5c86423 :: prose twin of the run farther from 1
\newcommand{\MlstMlpCmpJcTwoCoarseMassOverCoarse}{\pending{no run yet}}  % experiments/FIGS/data/v2/index.json :: nonlinear probe jc2_188: (1 - AUC) of R16_Q1_MASS over R16_Q1, run by run :: sha256 c6c4138ee5c86423
\newcommand{\MlstMlpCmpJcTwoCoarseMassOverCoarsePr}{\pending{no run yet}}  % experiments/FIGS/data/v2/index.json :: nonlinear probe jc2_188: (1 - AUC) of R16_Q1_MASS over R16_Q1, run by run :: sha256 c6c4138ee5c86423 :: prose twin: two significant figures, direction in words
\newcommand{\MlstMlpCmpJcTwoCoarseMassOverCoarsePrMax}{\pending{no run yet}}  % experiments/FIGS/data/v2/index.json :: nonlinear probe jc2_188: (1 - AUC) of R16_Q1_MASS over R16_Q1, run by run :: sha256 c6c4138ee5c86423 :: prose twin of the run farther from 1
\newcommand{\MlstMlpCmpJcTwoFzeroOverFull}{1.294 and 1.310}  % experiments/FIGS/data/v2/index.json :: nonlinear probe jc2_188: (1 - AUC) of FLAV_F0 over L188, run by run :: sha256 c6c4138ee5c86423
\newcommand{\MlstMlpCmpJcTwoFzeroOverFullPr}{at least 29\% more errors}  % experiments/FIGS/data/v2/index.json :: nonlinear probe jc2_188: (1 - AUC) of FLAV_F0 over L188, run by run :: sha256 c6c4138ee5c86423 :: prose twin: two significant figures, direction in words
\newcommand{\MlstMlpCmpJcTwoFzeroOverFullPrMax}{up to 31\% more errors}  % experiments/FIGS/data/v2/index.json :: nonlinear probe jc2_188: (1 - AUC) of FLAV_F0 over L188, run by run :: sha256 c6c4138ee5c86423 :: prose twin of the run farther from 1
\newcommand{\MlstMlpCmpJcTwoFoneOverFzero}{0.903 and 0.885}  % experiments/FIGS/data/v2/index.json :: nonlinear probe jc2_188: (1 - AUC) of FLAV_F1 over FLAV_F0, run by run :: sha256 c6c4138ee5c86423
\newcommand{\MlstMlpCmpJcTwoFoneOverFzeroPr}{at least 9.7\% fewer errors}  % experiments/FIGS/data/v2/index.json :: nonlinear probe jc2_188: (1 - AUC) of FLAV_F1 over FLAV_F0, run by run :: sha256 c6c4138ee5c86423 :: prose twin: two significant figures, direction in words
\newcommand{\MlstMlpCmpJcTwoFoneOverFzeroPrMax}{up to 12\% fewer errors}  % experiments/FIGS/data/v2/index.json :: nonlinear probe jc2_188: (1 - AUC) of FLAV_F1 over FLAV_F0, run by run :: sha256 c6c4138ee5c86423 :: prose twin of the run farther from 1
\newcommand{\MlstMlpCmpJcTwoFonerOverFone}{0.9916 and 0.9918}  % experiments/FIGS/data/v2/index.json :: nonlinear probe jc2_188: (1 - AUC) of FLAV_F1R over FLAV_F1, run by run :: sha256 c6c4138ee5c86423
\newcommand{\MlstMlpCmpJcTwoFonerOverFonePr}{equal within 1\%}  % experiments/FIGS/data/v2/index.json :: nonlinear probe jc2_188: (1 - AUC) of FLAV_F1R over FLAV_F1, run by run :: sha256 c6c4138ee5c86423 :: prose twin: two significant figures, direction in words
\newcommand{\MlstMlpCmpJcTwoFonerOverFonePrMax}{equal within 1\%}  % experiments/FIGS/data/v2/index.json :: nonlinear probe jc2_188: (1 - AUC) of FLAV_F1R over FLAV_F1, run by run :: sha256 c6c4138ee5c86423 :: prose twin of the run farther from 1
\newcommand{\MlstMlpCmpJcTwoFonerOverFzero}{0.895 and 0.877}  % experiments/FIGS/data/v2/index.json :: nonlinear probe jc2_188: (1 - AUC) of FLAV_F1R over FLAV_F0, run by run :: sha256 c6c4138ee5c86423
\newcommand{\MlstMlpCmpJcTwoFonerOverFzeroPr}{at least 10\% fewer errors}  % experiments/FIGS/data/v2/index.json :: nonlinear probe jc2_188: (1 - AUC) of FLAV_F1R over FLAV_F0, run by run :: sha256 c6c4138ee5c86423 :: prose twin: two significant figures, direction in words
\newcommand{\MlstMlpCmpJcTwoFonerOverFzeroPrMax}{up to 12\% fewer errors}  % experiments/FIGS/data/v2/index.json :: nonlinear probe jc2_188: (1 - AUC) of FLAV_F1R over FLAV_F0, run by run :: sha256 c6c4138ee5c86423 :: prose twin of the run farther from 1
\newcommand{\MlstMlpCmpJcTwoFzeroOverCoarse}{0.9461 and 0.9456}  % experiments/FIGS/data/v2/index.json :: nonlinear probe jc2_188: (1 - AUC) of FLAV_F0 over R16_Q1, run by run :: sha256 c6c4138ee5c86423
\newcommand{\MlstMlpCmpJcTwoFzeroOverCoarsePr}{at least 5.4\% fewer errors}  % experiments/FIGS/data/v2/index.json :: nonlinear probe jc2_188: (1 - AUC) of FLAV_F0 over R16_Q1, run by run :: sha256 c6c4138ee5c86423 :: prose twin: two significant figures, direction in words
\newcommand{\MlstMlpCmpJcTwoFzeroOverCoarsePrMax}{up to 5.4\% fewer errors}  % experiments/FIGS/data/v2/index.json :: nonlinear probe jc2_188: (1 - AUC) of FLAV_F0 over R16_Q1, run by run :: sha256 c6c4138ee5c86423 :: prose twin of the run farther from 1
\newcommand{\MlstMlpCmpJcTwoRandOverFull}{\pending{no run yet}}  % experiments/FIGS/data/v2/index.json :: nonlinear probe jc2_188: (1 - AUC) of each random partition over the 188-class run of its seed :: sha256 c6c4138ee5c86423
\newcommand{\MlstMlpCmpJcTwoRandOverFullPr}{\pending{no run yet}}  % experiments/FIGS/data/v2/index.json :: nonlinear probe jc2_188: (1 - AUC) of each random partition over the 188-class run of its seed :: sha256 c6c4138ee5c86423 :: prose twin: two significant figures, direction in words
\newcommand{\MlstMlpCmpJcTwoResResOverFull}{1.001\,\ensuremath{\pm}\,0.007}  % experiments/FIGS/data/v2/index.json :: nonlinear probe jc2_188_res: (1 - AUC) of L162 over L188, run by run :: sha256 c6c4138ee5c86423
\newcommand{\MlstMlpCmpJcTwoResResOverFullPr}{equal within 1\%}  % experiments/FIGS/data/v2/index.json :: nonlinear probe jc2_188_res: (1 - AUC) of L162 over L188, run by run :: sha256 c6c4138ee5c86423 :: prose twin: two significant figures, direction in words
\newcommand{\MlstMlpCmpJcTwoResMidOverFull}{1.192\,\ensuremath{\pm}\,0.008}  % experiments/FIGS/data/v2/index.json :: nonlinear probe jc2_188_res: (1 - AUC) of R42_Q1 over L188, run by run :: sha256 c6c4138ee5c86423
\newcommand{\MlstMlpCmpJcTwoResMidOverFullPr}{19\% more errors}  % experiments/FIGS/data/v2/index.json :: nonlinear probe jc2_188_res: (1 - AUC) of R42_Q1 over L188, run by run :: sha256 c6c4138ee5c86423 :: prose twin: two significant figures, direction in words
\newcommand{\MlstMlpCmpJcTwoResCoarseOverFull}{1.408\,\ensuremath{\pm}\,0.026}  % experiments/FIGS/data/v2/index.json :: nonlinear probe jc2_188_res: (1 - AUC) of R16_Q1 over L188, run by run :: sha256 c6c4138ee5c86423
\newcommand{\MlstMlpCmpJcTwoResCoarseOverFullPr}{41\% more errors}  % experiments/FIGS/data/v2/index.json :: nonlinear probe jc2_188_res: (1 - AUC) of R16_Q1 over L188, run by run :: sha256 c6c4138ee5c86423 :: prose twin: two significant figures, direction in words
\newcommand{\MlstMlpCmpJcTwoResResMassOverRes}{\pending{no run yet}}  % experiments/FIGS/data/v2/index.json :: nonlinear probe jc2_188_res: (1 - AUC) of L162_MASS over L162, run by run :: sha256 c6c4138ee5c86423
\newcommand{\MlstMlpCmpJcTwoResResMassOverResPr}{\pending{no run yet}}  % experiments/FIGS/data/v2/index.json :: nonlinear probe jc2_188_res: (1 - AUC) of L162_MASS over L162, run by run :: sha256 c6c4138ee5c86423 :: prose twin: two significant figures, direction in words
\newcommand{\MlstMlpCmpJcTwoResResMassOverResPrMax}{\pending{no run yet}}  % experiments/FIGS/data/v2/index.json :: nonlinear probe jc2_188_res: (1 - AUC) of L162_MASS over L162, run by run :: sha256 c6c4138ee5c86423 :: prose twin of the run farther from 1
\newcommand{\MlstMlpCmpJcTwoResCoarseMassOverCoarse}{\pending{no run yet}}  % experiments/FIGS/data/v2/index.json :: nonlinear probe jc2_188_res: (1 - AUC) of R16_Q1_MASS over R16_Q1, run by run :: sha256 c6c4138ee5c86423
\newcommand{\MlstMlpCmpJcTwoResCoarseMassOverCoarsePr}{\pending{no run yet}}  % experiments/FIGS/data/v2/index.json :: nonlinear probe jc2_188_res: (1 - AUC) of R16_Q1_MASS over R16_Q1, run by run :: sha256 c6c4138ee5c86423 :: prose twin: two significant figures, direction in words
\newcommand{\MlstMlpCmpJcTwoResCoarseMassOverCoarsePrMax}{\pending{no run yet}}  % experiments/FIGS/data/v2/index.json :: nonlinear probe jc2_188_res: (1 - AUC) of R16_Q1_MASS over R16_Q1, run by run :: sha256 c6c4138ee5c86423 :: prose twin of the run farther from 1
\newcommand{\MlstMlpCmpJcTwoResFzeroOverFull}{1.285 and 1.300}  % experiments/FIGS/data/v2/index.json :: nonlinear probe jc2_188_res: (1 - AUC) of FLAV_F0 over L188, run by run :: sha256 c6c4138ee5c86423
\newcommand{\MlstMlpCmpJcTwoResFzeroOverFullPr}{at least 29\% more errors}  % experiments/FIGS/data/v2/index.json :: nonlinear probe jc2_188_res: (1 - AUC) of FLAV_F0 over L188, run by run :: sha256 c6c4138ee5c86423 :: prose twin: two significant figures, direction in words
\newcommand{\MlstMlpCmpJcTwoResFzeroOverFullPrMax}{up to 30\% more errors}  % experiments/FIGS/data/v2/index.json :: nonlinear probe jc2_188_res: (1 - AUC) of FLAV_F0 over L188, run by run :: sha256 c6c4138ee5c86423 :: prose twin of the run farther from 1
\newcommand{\MlstMlpCmpJcTwoResFoneOverFzero}{0.909 and 0.889}  % experiments/FIGS/data/v2/index.json :: nonlinear probe jc2_188_res: (1 - AUC) of FLAV_F1 over FLAV_F0, run by run :: sha256 c6c4138ee5c86423
\newcommand{\MlstMlpCmpJcTwoResFoneOverFzeroPr}{at least 9.1\% fewer errors}  % experiments/FIGS/data/v2/index.json :: nonlinear probe jc2_188_res: (1 - AUC) of FLAV_F1 over FLAV_F0, run by run :: sha256 c6c4138ee5c86423 :: prose twin: two significant figures, direction in words
\newcommand{\MlstMlpCmpJcTwoResFoneOverFzeroPrMax}{up to 11\% fewer errors}  % experiments/FIGS/data/v2/index.json :: nonlinear probe jc2_188_res: (1 - AUC) of FLAV_F1 over FLAV_F0, run by run :: sha256 c6c4138ee5c86423 :: prose twin of the run farther from 1
\newcommand{\MlstMlpCmpJcTwoResFonerOverFone}{0.9793 and 0.9825}  % experiments/FIGS/data/v2/index.json :: nonlinear probe jc2_188_res: (1 - AUC) of FLAV_F1R over FLAV_F1, run by run :: sha256 c6c4138ee5c86423
\newcommand{\MlstMlpCmpJcTwoResFonerOverFonePr}{equal within 3\%}  % experiments/FIGS/data/v2/index.json :: nonlinear probe jc2_188_res: (1 - AUC) of FLAV_F1R over FLAV_F1, run by run :: sha256 c6c4138ee5c86423 :: prose twin: two significant figures, direction in words
\newcommand{\MlstMlpCmpJcTwoResFonerOverFonePrMax}{equal within 3\%}  % experiments/FIGS/data/v2/index.json :: nonlinear probe jc2_188_res: (1 - AUC) of FLAV_F1R over FLAV_F1, run by run :: sha256 c6c4138ee5c86423 :: prose twin of the run farther from 1
\newcommand{\MlstMlpCmpJcTwoResFonerOverFzero}{0.890 and 0.874}  % experiments/FIGS/data/v2/index.json :: nonlinear probe jc2_188_res: (1 - AUC) of FLAV_F1R over FLAV_F0, run by run :: sha256 c6c4138ee5c86423
\newcommand{\MlstMlpCmpJcTwoResFonerOverFzeroPr}{at least 11\% fewer errors}  % experiments/FIGS/data/v2/index.json :: nonlinear probe jc2_188_res: (1 - AUC) of FLAV_F1R over FLAV_F0, run by run :: sha256 c6c4138ee5c86423 :: prose twin: two significant figures, direction in words
\newcommand{\MlstMlpCmpJcTwoResFonerOverFzeroPrMax}{up to 13\% fewer errors}  % experiments/FIGS/data/v2/index.json :: nonlinear probe jc2_188_res: (1 - AUC) of FLAV_F1R over FLAV_F0, run by run :: sha256 c6c4138ee5c86423 :: prose twin of the run farther from 1
\newcommand{\MlstMlpCmpJcTwoResFzeroOverCoarse}{0.9299 and 0.9233}  % experiments/FIGS/data/v2/index.json :: nonlinear probe jc2_188_res: (1 - AUC) of FLAV_F0 over R16_Q1, run by run :: sha256 c6c4138ee5c86423
\newcommand{\MlstMlpCmpJcTwoResFzeroOverCoarsePr}{at least 7.0\% fewer errors}  % experiments/FIGS/data/v2/index.json :: nonlinear probe jc2_188_res: (1 - AUC) of FLAV_F0 over R16_Q1, run by run :: sha256 c6c4138ee5c86423 :: prose twin: two significant figures, direction in words
\newcommand{\MlstMlpCmpJcTwoResFzeroOverCoarsePrMax}{up to 7.7\% fewer errors}  % experiments/FIGS/data/v2/index.json :: nonlinear probe jc2_188_res: (1 - AUC) of FLAV_F0 over R16_Q1, run by run :: sha256 c6c4138ee5c86423 :: prose twin of the run farther from 1
\newcommand{\MlstMlpCmpJcTwoResRandOverFull}{\pending{no run yet}}  % experiments/FIGS/data/v2/index.json :: nonlinear probe jc2_188_res: (1 - AUC) of each random partition over the 188-class run of its seed :: sha256 c6c4138ee5c86423
\newcommand{\MlstMlpCmpJcTwoResRandOverFullPr}{\pending{no run yet}}  % experiments/FIGS/data/v2/index.json :: nonlinear probe jc2_188_res: (1 - AUC) of each random partition over the 188-class run of its seed :: sha256 c6c4138ee5c86423 :: prose twin: two significant figures, direction in words
\newcommand{\MlstMlpCmpJcTwoQcdResOverFull}{1.070\,\ensuremath{\pm}\,0.006}  % experiments/FIGS/data/v2/index.json :: nonlinear probe jc2_188_qcd: (1 - AUC) of L162 over L188, run by run :: sha256 c6c4138ee5c86423
\newcommand{\MlstMlpCmpJcTwoQcdResOverFullPr}{7.0\% more errors}  % experiments/FIGS/data/v2/index.json :: nonlinear probe jc2_188_qcd: (1 - AUC) of L162 over L188, run by run :: sha256 c6c4138ee5c86423 :: prose twin: two significant figures, direction in words
\newcommand{\MlstMlpCmpJcTwoQcdMidOverFull}{1.107\,\ensuremath{\pm}\,0.002}  % experiments/FIGS/data/v2/index.json :: nonlinear probe jc2_188_qcd: (1 - AUC) of R42_Q1 over L188, run by run :: sha256 c6c4138ee5c86423
\newcommand{\MlstMlpCmpJcTwoQcdMidOverFullPr}{11\% more errors}  % experiments/FIGS/data/v2/index.json :: nonlinear probe jc2_188_qcd: (1 - AUC) of R42_Q1 over L188, run by run :: sha256 c6c4138ee5c86423 :: prose twin: two significant figures, direction in words
\newcommand{\MlstMlpCmpJcTwoQcdCoarseOverFull}{1.333\,\ensuremath{\pm}\,0.012}  % experiments/FIGS/data/v2/index.json :: nonlinear probe jc2_188_qcd: (1 - AUC) of R16_Q1 over L188, run by run :: sha256 c6c4138ee5c86423
\newcommand{\MlstMlpCmpJcTwoQcdCoarseOverFullPr}{33\% more errors}  % experiments/FIGS/data/v2/index.json :: nonlinear probe jc2_188_qcd: (1 - AUC) of R16_Q1 over L188, run by run :: sha256 c6c4138ee5c86423 :: prose twin: two significant figures, direction in words
\newcommand{\MlstMlpCmpJcTwoQcdResMassOverRes}{\pending{no run yet}}  % experiments/FIGS/data/v2/index.json :: nonlinear probe jc2_188_qcd: (1 - AUC) of L162_MASS over L162, run by run :: sha256 c6c4138ee5c86423
\newcommand{\MlstMlpCmpJcTwoQcdResMassOverResPr}{\pending{no run yet}}  % experiments/FIGS/data/v2/index.json :: nonlinear probe jc2_188_qcd: (1 - AUC) of L162_MASS over L162, run by run :: sha256 c6c4138ee5c86423 :: prose twin: two significant figures, direction in words
\newcommand{\MlstMlpCmpJcTwoQcdResMassOverResPrMax}{\pending{no run yet}}  % experiments/FIGS/data/v2/index.json :: nonlinear probe jc2_188_qcd: (1 - AUC) of L162_MASS over L162, run by run :: sha256 c6c4138ee5c86423 :: prose twin of the run farther from 1
\newcommand{\MlstMlpCmpJcTwoQcdCoarseMassOverCoarse}{\pending{no run yet}}  % experiments/FIGS/data/v2/index.json :: nonlinear probe jc2_188_qcd: (1 - AUC) of R16_Q1_MASS over R16_Q1, run by run :: sha256 c6c4138ee5c86423
\newcommand{\MlstMlpCmpJcTwoQcdCoarseMassOverCoarsePr}{\pending{no run yet}}  % experiments/FIGS/data/v2/index.json :: nonlinear probe jc2_188_qcd: (1 - AUC) of R16_Q1_MASS over R16_Q1, run by run :: sha256 c6c4138ee5c86423 :: prose twin: two significant figures, direction in words
\newcommand{\MlstMlpCmpJcTwoQcdCoarseMassOverCoarsePrMax}{\pending{no run yet}}  % experiments/FIGS/data/v2/index.json :: nonlinear probe jc2_188_qcd: (1 - AUC) of R16_Q1_MASS over R16_Q1, run by run :: sha256 c6c4138ee5c86423 :: prose twin of the run farther from 1
\newcommand{\MlstMlpCmpJcTwoQcdFzeroOverFull}{1.319 and 1.339}  % experiments/FIGS/data/v2/index.json :: nonlinear probe jc2_188_qcd: (1 - AUC) of FLAV_F0 over L188, run by run :: sha256 c6c4138ee5c86423
\newcommand{\MlstMlpCmpJcTwoQcdFzeroOverFullPr}{at least 32\% more errors}  % experiments/FIGS/data/v2/index.json :: nonlinear probe jc2_188_qcd: (1 - AUC) of FLAV_F0 over L188, run by run :: sha256 c6c4138ee5c86423 :: prose twin: two significant figures, direction in words
\newcommand{\MlstMlpCmpJcTwoQcdFzeroOverFullPrMax}{up to 34\% more errors}  % experiments/FIGS/data/v2/index.json :: nonlinear probe jc2_188_qcd: (1 - AUC) of FLAV_F0 over L188, run by run :: sha256 c6c4138ee5c86423 :: prose twin of the run farther from 1
\newcommand{\MlstMlpCmpJcTwoQcdFoneOverFzero}{0.886 and 0.872}  % experiments/FIGS/data/v2/index.json :: nonlinear probe jc2_188_qcd: (1 - AUC) of FLAV_F1 over FLAV_F0, run by run :: sha256 c6c4138ee5c86423
\newcommand{\MlstMlpCmpJcTwoQcdFoneOverFzeroPr}{at least 11\% fewer errors}  % experiments/FIGS/data/v2/index.json :: nonlinear probe jc2_188_qcd: (1 - AUC) of FLAV_F1 over FLAV_F0, run by run :: sha256 c6c4138ee5c86423 :: prose twin: two significant figures, direction in words
\newcommand{\MlstMlpCmpJcTwoQcdFoneOverFzeroPrMax}{up to 13\% fewer errors}  % experiments/FIGS/data/v2/index.json :: nonlinear probe jc2_188_qcd: (1 - AUC) of FLAV_F1 over FLAV_F0, run by run :: sha256 c6c4138ee5c86423 :: prose twin of the run farther from 1
\newcommand{\MlstMlpCmpJcTwoQcdFonerOverFone}{1.026 and 1.018}  % experiments/FIGS/data/v2/index.json :: nonlinear probe jc2_188_qcd: (1 - AUC) of FLAV_F1R over FLAV_F1, run by run :: sha256 c6c4138ee5c86423
\newcommand{\MlstMlpCmpJcTwoQcdFonerOverFonePr}{at least 1.8\% more errors}  % experiments/FIGS/data/v2/index.json :: nonlinear probe jc2_188_qcd: (1 - AUC) of FLAV_F1R over FLAV_F1, run by run :: sha256 c6c4138ee5c86423 :: prose twin: two significant figures, direction in words
\newcommand{\MlstMlpCmpJcTwoQcdFonerOverFonePrMax}{up to 2.6\% more errors}  % experiments/FIGS/data/v2/index.json :: nonlinear probe jc2_188_qcd: (1 - AUC) of FLAV_F1R over FLAV_F1, run by run :: sha256 c6c4138ee5c86423 :: prose twin of the run farther from 1
\newcommand{\MlstMlpCmpJcTwoQcdFonerOverFzero}{0.909 and 0.887}  % experiments/FIGS/data/v2/index.json :: nonlinear probe jc2_188_qcd: (1 - AUC) of FLAV_F1R over FLAV_F0, run by run :: sha256 c6c4138ee5c86423
\newcommand{\MlstMlpCmpJcTwoQcdFonerOverFzeroPr}{at least 9.1\% fewer errors}  % experiments/FIGS/data/v2/index.json :: nonlinear probe jc2_188_qcd: (1 - AUC) of FLAV_F1R over FLAV_F0, run by run :: sha256 c6c4138ee5c86423 :: prose twin: two significant figures, direction in words
\newcommand{\MlstMlpCmpJcTwoQcdFonerOverFzeroPrMax}{up to 11\% fewer errors}  % experiments/FIGS/data/v2/index.json :: nonlinear probe jc2_188_qcd: (1 - AUC) of FLAV_F1R over FLAV_F0, run by run :: sha256 c6c4138ee5c86423 :: prose twin of the run farther from 1
\newcommand{\MlstMlpCmpJcTwoQcdFzeroOverCoarse}{0.993 and 1.012}  % experiments/FIGS/data/v2/index.json :: nonlinear probe jc2_188_qcd: (1 - AUC) of FLAV_F0 over R16_Q1, run by run :: sha256 c6c4138ee5c86423
\newcommand{\MlstMlpCmpJcTwoQcdFzeroOverCoarsePr}{equal within 2\%}  % experiments/FIGS/data/v2/index.json :: nonlinear probe jc2_188_qcd: (1 - AUC) of FLAV_F0 over R16_Q1, run by run :: sha256 c6c4138ee5c86423 :: prose twin: two significant figures, direction in words
\newcommand{\MlstMlpCmpJcTwoQcdFzeroOverCoarsePrMax}{equal within 2\%}  % experiments/FIGS/data/v2/index.json :: nonlinear probe jc2_188_qcd: (1 - AUC) of FLAV_F0 over R16_Q1, run by run :: sha256 c6c4138ee5c86423 :: prose twin of the run farther from 1
\newcommand{\MlstMlpCmpJcTwoQcdRandOverFull}{\pending{no run yet}}  % experiments/FIGS/data/v2/index.json :: nonlinear probe jc2_188_qcd: (1 - AUC) of each random partition over the 188-class run of its seed :: sha256 c6c4138ee5c86423
\newcommand{\MlstMlpCmpJcTwoQcdRandOverFullPr}{\pending{no run yet}}  % experiments/FIGS/data/v2/index.json :: nonlinear probe jc2_188_qcd: (1 - AUC) of each random partition over the 188-class run of its seed :: sha256 c6c4138ee5c86423 :: prose twin: two significant figures, direction in words
\newcommand{\MlstMlpCmpBcResOverFull}{1.002\,\ensuremath{\pm}\,0.006}  % experiments/FIGS/data/v2/index.json :: nonlinear probe bc_vs_bq_cs: (1 - AUC) of L162 over L188, run by run :: sha256 c6c4138ee5c86423
\newcommand{\MlstMlpCmpBcResOverFullPr}{equal within 1\%}  % experiments/FIGS/data/v2/index.json :: nonlinear probe bc_vs_bq_cs: (1 - AUC) of L162 over L188, run by run :: sha256 c6c4138ee5c86423 :: prose twin: two significant figures, direction in words
\newcommand{\MlstMlpCmpBcMidOverFull}{1.104\,\ensuremath{\pm}\,0.017}  % experiments/FIGS/data/v2/index.json :: nonlinear probe bc_vs_bq_cs: (1 - AUC) of R42_Q1 over L188, run by run :: sha256 c6c4138ee5c86423
\newcommand{\MlstMlpCmpBcMidOverFullPr}{10\% more errors}  % experiments/FIGS/data/v2/index.json :: nonlinear probe bc_vs_bq_cs: (1 - AUC) of R42_Q1 over L188, run by run :: sha256 c6c4138ee5c86423 :: prose twin: two significant figures, direction in words
\newcommand{\MlstMlpCmpBcCoarseOverFull}{1.624\,\ensuremath{\pm}\,0.022}  % experiments/FIGS/data/v2/index.json :: nonlinear probe bc_vs_bq_cs: (1 - AUC) of R16_Q1 over L188, run by run :: sha256 c6c4138ee5c86423
\newcommand{\MlstMlpCmpBcCoarseOverFullPr}{62\% more errors}  % experiments/FIGS/data/v2/index.json :: nonlinear probe bc_vs_bq_cs: (1 - AUC) of R16_Q1 over L188, run by run :: sha256 c6c4138ee5c86423 :: prose twin: two significant figures, direction in words
\newcommand{\MlstMlpCmpBcResMassOverRes}{\pending{no run yet}}  % experiments/FIGS/data/v2/index.json :: nonlinear probe bc_vs_bq_cs: (1 - AUC) of L162_MASS over L162, run by run :: sha256 c6c4138ee5c86423
\newcommand{\MlstMlpCmpBcResMassOverResPr}{\pending{no run yet}}  % experiments/FIGS/data/v2/index.json :: nonlinear probe bc_vs_bq_cs: (1 - AUC) of L162_MASS over L162, run by run :: sha256 c6c4138ee5c86423 :: prose twin: two significant figures, direction in words
\newcommand{\MlstMlpCmpBcResMassOverResPrMax}{\pending{no run yet}}  % experiments/FIGS/data/v2/index.json :: nonlinear probe bc_vs_bq_cs: (1 - AUC) of L162_MASS over L162, run by run :: sha256 c6c4138ee5c86423 :: prose twin of the run farther from 1
\newcommand{\MlstMlpCmpBcCoarseMassOverCoarse}{\pending{no run yet}}  % experiments/FIGS/data/v2/index.json :: nonlinear probe bc_vs_bq_cs: (1 - AUC) of R16_Q1_MASS over R16_Q1, run by run :: sha256 c6c4138ee5c86423
\newcommand{\MlstMlpCmpBcCoarseMassOverCoarsePr}{\pending{no run yet}}  % experiments/FIGS/data/v2/index.json :: nonlinear probe bc_vs_bq_cs: (1 - AUC) of R16_Q1_MASS over R16_Q1, run by run :: sha256 c6c4138ee5c86423 :: prose twin: two significant figures, direction in words
\newcommand{\MlstMlpCmpBcCoarseMassOverCoarsePrMax}{\pending{no run yet}}  % experiments/FIGS/data/v2/index.json :: nonlinear probe bc_vs_bq_cs: (1 - AUC) of R16_Q1_MASS over R16_Q1, run by run :: sha256 c6c4138ee5c86423 :: prose twin of the run farther from 1
\newcommand{\MlstMlpCmpBcFzeroOverFull}{1.633 and 1.679}  % experiments/FIGS/data/v2/index.json :: nonlinear probe bc_vs_bq_cs: (1 - AUC) of FLAV_F0 over L188, run by run :: sha256 c6c4138ee5c86423
\newcommand{\MlstMlpCmpBcFzeroOverFullPr}{at least 63\% more errors}  % experiments/FIGS/data/v2/index.json :: nonlinear probe bc_vs_bq_cs: (1 - AUC) of FLAV_F0 over L188, run by run :: sha256 c6c4138ee5c86423 :: prose twin: two significant figures, direction in words
\newcommand{\MlstMlpCmpBcFzeroOverFullPrMax}{up to 68\% more errors}  % experiments/FIGS/data/v2/index.json :: nonlinear probe bc_vs_bq_cs: (1 - AUC) of FLAV_F0 over L188, run by run :: sha256 c6c4138ee5c86423 :: prose twin of the run farther from 1
\newcommand{\MlstMlpCmpBcFoneOverFzero}{0.848 and 0.810}  % experiments/FIGS/data/v2/index.json :: nonlinear probe bc_vs_bq_cs: (1 - AUC) of FLAV_F1 over FLAV_F0, run by run :: sha256 c6c4138ee5c86423
\newcommand{\MlstMlpCmpBcFoneOverFzeroPr}{at least 15\% fewer errors}  % experiments/FIGS/data/v2/index.json :: nonlinear probe bc_vs_bq_cs: (1 - AUC) of FLAV_F1 over FLAV_F0, run by run :: sha256 c6c4138ee5c86423 :: prose twin: two significant figures, direction in words
\newcommand{\MlstMlpCmpBcFoneOverFzeroPrMax}{up to 19\% fewer errors}  % experiments/FIGS/data/v2/index.json :: nonlinear probe bc_vs_bq_cs: (1 - AUC) of FLAV_F1 over FLAV_F0, run by run :: sha256 c6c4138ee5c86423 :: prose twin of the run farther from 1
\newcommand{\MlstMlpCmpBcFonerOverFone}{0.9812 and 0.9815}  % experiments/FIGS/data/v2/index.json :: nonlinear probe bc_vs_bq_cs: (1 - AUC) of FLAV_F1R over FLAV_F1, run by run :: sha256 c6c4138ee5c86423
\newcommand{\MlstMlpCmpBcFonerOverFonePr}{equal within 2\%}  % experiments/FIGS/data/v2/index.json :: nonlinear probe bc_vs_bq_cs: (1 - AUC) of FLAV_F1R over FLAV_F1, run by run :: sha256 c6c4138ee5c86423 :: prose twin: two significant figures, direction in words
\newcommand{\MlstMlpCmpBcFonerOverFonePrMax}{equal within 2\%}  % experiments/FIGS/data/v2/index.json :: nonlinear probe bc_vs_bq_cs: (1 - AUC) of FLAV_F1R over FLAV_F1, run by run :: sha256 c6c4138ee5c86423 :: prose twin of the run farther from 1
\newcommand{\MlstMlpCmpBcFonerOverFzero}{0.833 and 0.795}  % experiments/FIGS/data/v2/index.json :: nonlinear probe bc_vs_bq_cs: (1 - AUC) of FLAV_F1R over FLAV_F0, run by run :: sha256 c6c4138ee5c86423
\newcommand{\MlstMlpCmpBcFonerOverFzeroPr}{at least 17\% fewer errors}  % experiments/FIGS/data/v2/index.json :: nonlinear probe bc_vs_bq_cs: (1 - AUC) of FLAV_F1R over FLAV_F0, run by run :: sha256 c6c4138ee5c86423 :: prose twin: two significant figures, direction in words
\newcommand{\MlstMlpCmpBcFonerOverFzeroPrMax}{up to 21\% fewer errors}  % experiments/FIGS/data/v2/index.json :: nonlinear probe bc_vs_bq_cs: (1 - AUC) of FLAV_F1R over FLAV_F0, run by run :: sha256 c6c4138ee5c86423 :: prose twin of the run farther from 1
\newcommand{\MlstMlpCmpBcFzeroOverCoarse}{0.994 and 1.030}  % experiments/FIGS/data/v2/index.json :: nonlinear probe bc_vs_bq_cs: (1 - AUC) of FLAV_F0 over R16_Q1, run by run :: sha256 c6c4138ee5c86423
\newcommand{\MlstMlpCmpBcFzeroOverCoarsePr}{equal within 3\%}  % experiments/FIGS/data/v2/index.json :: nonlinear probe bc_vs_bq_cs: (1 - AUC) of FLAV_F0 over R16_Q1, run by run :: sha256 c6c4138ee5c86423 :: prose twin: two significant figures, direction in words
\newcommand{\MlstMlpCmpBcFzeroOverCoarsePrMax}{equal within 3\%}  % experiments/FIGS/data/v2/index.json :: nonlinear probe bc_vs_bq_cs: (1 - AUC) of FLAV_F0 over R16_Q1, run by run :: sha256 c6c4138ee5c86423 :: prose twin of the run farther from 1
\newcommand{\MlstMlpCmpBcRandOverFull}{\pending{no run yet}}  % experiments/FIGS/data/v2/index.json :: nonlinear probe bc_vs_bq_cs: (1 - AUC) of each random partition over the 188-class run of its seed :: sha256 c6c4138ee5c86423
\newcommand{\MlstMlpCmpBcRandOverFullPr}{\pending{no run yet}}  % experiments/FIGS/data/v2/index.json :: nonlinear probe bc_vs_bq_cs: (1 - AUC) of each random partition over the 188-class run of its seed :: sha256 c6c4138ee5c86423 :: prose twin: two significant figures, direction in words
\newcommand{\MlstMlpCmpBbResOverFull}{1.023\,\ensuremath{\pm}\,0.021}  % experiments/FIGS/data/v2/index.json :: nonlinear probe bb_vs_cc: (1 - AUC) of L162 over L188, run by run :: sha256 c6c4138ee5c86423
\newcommand{\MlstMlpCmpBbResOverFullPr}{equal within 3\%}  % experiments/FIGS/data/v2/index.json :: nonlinear probe bb_vs_cc: (1 - AUC) of L162 over L188, run by run :: sha256 c6c4138ee5c86423 :: prose twin: two significant figures, direction in words
\newcommand{\MlstMlpCmpBbMidOverFull}{1.125\,\ensuremath{\pm}\,0.011}  % experiments/FIGS/data/v2/index.json :: nonlinear probe bb_vs_cc: (1 - AUC) of R42_Q1 over L188, run by run :: sha256 c6c4138ee5c86423
\newcommand{\MlstMlpCmpBbMidOverFullPr}{13\% more errors}  % experiments/FIGS/data/v2/index.json :: nonlinear probe bb_vs_cc: (1 - AUC) of R42_Q1 over L188, run by run :: sha256 c6c4138ee5c86423 :: prose twin: two significant figures, direction in words
\newcommand{\MlstMlpCmpBbCoarseOverFull}{3.04\,\ensuremath{\pm}\,0.09}  % experiments/FIGS/data/v2/index.json :: nonlinear probe bb_vs_cc: (1 - AUC) of R16_Q1 over L188, run by run :: sha256 c6c4138ee5c86423
\newcommand{\MlstMlpCmpBbCoarseOverFullPr}{3.0 times the errors}  % experiments/FIGS/data/v2/index.json :: nonlinear probe bb_vs_cc: (1 - AUC) of R16_Q1 over L188, run by run :: sha256 c6c4138ee5c86423 :: prose twin: two significant figures, direction in words
\newcommand{\MlstMlpCmpBbResMassOverRes}{\pending{no run yet}}  % experiments/FIGS/data/v2/index.json :: nonlinear probe bb_vs_cc: (1 - AUC) of L162_MASS over L162, run by run :: sha256 c6c4138ee5c86423
\newcommand{\MlstMlpCmpBbResMassOverResPr}{\pending{no run yet}}  % experiments/FIGS/data/v2/index.json :: nonlinear probe bb_vs_cc: (1 - AUC) of L162_MASS over L162, run by run :: sha256 c6c4138ee5c86423 :: prose twin: two significant figures, direction in words
\newcommand{\MlstMlpCmpBbResMassOverResPrMax}{\pending{no run yet}}  % experiments/FIGS/data/v2/index.json :: nonlinear probe bb_vs_cc: (1 - AUC) of L162_MASS over L162, run by run :: sha256 c6c4138ee5c86423 :: prose twin of the run farther from 1
\newcommand{\MlstMlpCmpBbCoarseMassOverCoarse}{\pending{no run yet}}  % experiments/FIGS/data/v2/index.json :: nonlinear probe bb_vs_cc: (1 - AUC) of R16_Q1_MASS over R16_Q1, run by run :: sha256 c6c4138ee5c86423
\newcommand{\MlstMlpCmpBbCoarseMassOverCoarsePr}{\pending{no run yet}}  % experiments/FIGS/data/v2/index.json :: nonlinear probe bb_vs_cc: (1 - AUC) of R16_Q1_MASS over R16_Q1, run by run :: sha256 c6c4138ee5c86423 :: prose twin: two significant figures, direction in words
\newcommand{\MlstMlpCmpBbCoarseMassOverCoarsePrMax}{\pending{no run yet}}  % experiments/FIGS/data/v2/index.json :: nonlinear probe bb_vs_cc: (1 - AUC) of R16_Q1_MASS over R16_Q1, run by run :: sha256 c6c4138ee5c86423 :: prose twin of the run farther from 1
\newcommand{\MlstMlpCmpBbFzeroOverFull}{3.3 and 4.3}  % experiments/FIGS/data/v2/index.json :: nonlinear probe bb_vs_cc: (1 - AUC) of FLAV_F0 over L188, run by run :: sha256 c6c4138ee5c86423
\newcommand{\MlstMlpCmpBbFzeroOverFullPr}{at least 3.3 times the errors}  % experiments/FIGS/data/v2/index.json :: nonlinear probe bb_vs_cc: (1 - AUC) of FLAV_F0 over L188, run by run :: sha256 c6c4138ee5c86423 :: prose twin: two significant figures, direction in words
\newcommand{\MlstMlpCmpBbFzeroOverFullPrMax}{up to 4.3 times the errors}  % experiments/FIGS/data/v2/index.json :: nonlinear probe bb_vs_cc: (1 - AUC) of FLAV_F0 over L188, run by run :: sha256 c6c4138ee5c86423 :: prose twin of the run farther from 1
\newcommand{\MlstMlpCmpBbFoneOverFzero}{0.52 and 0.38}  % experiments/FIGS/data/v2/index.json :: nonlinear probe bb_vs_cc: (1 - AUC) of FLAV_F1 over FLAV_F0, run by run :: sha256 c6c4138ee5c86423
\newcommand{\MlstMlpCmpBbFoneOverFzeroPr}{at least 48\% fewer errors}  % experiments/FIGS/data/v2/index.json :: nonlinear probe bb_vs_cc: (1 - AUC) of FLAV_F1 over FLAV_F0, run by run :: sha256 c6c4138ee5c86423 :: prose twin: two significant figures, direction in words
\newcommand{\MlstMlpCmpBbFoneOverFzeroPrMax}{down to 0.38 times the errors}  % experiments/FIGS/data/v2/index.json :: nonlinear probe bb_vs_cc: (1 - AUC) of FLAV_F1 over FLAV_F0, run by run :: sha256 c6c4138ee5c86423 :: prose twin of the run farther from 1
\newcommand{\MlstMlpCmpBbFonerOverFone}{1.27 and 1.45}  % experiments/FIGS/data/v2/index.json :: nonlinear probe bb_vs_cc: (1 - AUC) of FLAV_F1R over FLAV_F1, run by run :: sha256 c6c4138ee5c86423
\newcommand{\MlstMlpCmpBbFonerOverFonePr}{at least 27\% more errors}  % experiments/FIGS/data/v2/index.json :: nonlinear probe bb_vs_cc: (1 - AUC) of FLAV_F1R over FLAV_F1, run by run :: sha256 c6c4138ee5c86423 :: prose twin: two significant figures, direction in words
\newcommand{\MlstMlpCmpBbFonerOverFonePrMax}{up to 45\% more errors}  % experiments/FIGS/data/v2/index.json :: nonlinear probe bb_vs_cc: (1 - AUC) of FLAV_F1R over FLAV_F1, run by run :: sha256 c6c4138ee5c86423 :: prose twin of the run farther from 1
\newcommand{\MlstMlpCmpBbFonerOverFzero}{0.66 and 0.55}  % experiments/FIGS/data/v2/index.json :: nonlinear probe bb_vs_cc: (1 - AUC) of FLAV_F1R over FLAV_F0, run by run :: sha256 c6c4138ee5c86423
\newcommand{\MlstMlpCmpBbFonerOverFzeroPr}{at least 34\% fewer errors}  % experiments/FIGS/data/v2/index.json :: nonlinear probe bb_vs_cc: (1 - AUC) of FLAV_F1R over FLAV_F0, run by run :: sha256 c6c4138ee5c86423 :: prose twin: two significant figures, direction in words
\newcommand{\MlstMlpCmpBbFonerOverFzeroPrMax}{up to 45\% fewer errors}  % experiments/FIGS/data/v2/index.json :: nonlinear probe bb_vs_cc: (1 - AUC) of FLAV_F1R over FLAV_F0, run by run :: sha256 c6c4138ee5c86423 :: prose twin of the run farther from 1
\newcommand{\MlstMlpCmpBbFzeroOverCoarse}{1.07 and 1.41}  % experiments/FIGS/data/v2/index.json :: nonlinear probe bb_vs_cc: (1 - AUC) of FLAV_F0 over R16_Q1, run by run :: sha256 c6c4138ee5c86423
\newcommand{\MlstMlpCmpBbFzeroOverCoarsePr}{at least 7.4\% more errors}  % experiments/FIGS/data/v2/index.json :: nonlinear probe bb_vs_cc: (1 - AUC) of FLAV_F0 over R16_Q1, run by run :: sha256 c6c4138ee5c86423 :: prose twin: two significant figures, direction in words
\newcommand{\MlstMlpCmpBbFzeroOverCoarsePrMax}{up to 41\% more errors}  % experiments/FIGS/data/v2/index.json :: nonlinear probe bb_vs_cc: (1 - AUC) of FLAV_F0 over R16_Q1, run by run :: sha256 c6c4138ee5c86423 :: prose twin of the run farther from 1
\newcommand{\MlstMlpCmpBbRandOverFull}{\pending{no run yet}}  % experiments/FIGS/data/v2/index.json :: nonlinear probe bb_vs_cc: (1 - AUC) of each random partition over the 188-class run of its seed :: sha256 c6c4138ee5c86423
\newcommand{\MlstMlpCmpBbRandOverFullPr}{\pending{no run yet}}  % experiments/FIGS/data/v2/index.json :: nonlinear probe bb_vs_cc: (1 - AUC) of each random partition over the 188-class run of its seed :: sha256 c6c4138ee5c86423 :: prose twin: two significant figures, direction in words
\newcommand{\MlstMlpCmpBbqqResOverFull}{1.011\,\ensuremath{\pm}\,0.030}  % experiments/FIGS/data/v2/index.json :: nonlinear probe bbqq_vs_ccqq: (1 - AUC) of L162 over L188, run by run :: sha256 c6c4138ee5c86423
\newcommand{\MlstMlpCmpBbqqResOverFullPr}{equal within 3\%}  % experiments/FIGS/data/v2/index.json :: nonlinear probe bbqq_vs_ccqq: (1 - AUC) of L162 over L188, run by run :: sha256 c6c4138ee5c86423 :: prose twin: two significant figures, direction in words
\newcommand{\MlstMlpCmpBbqqMidOverFull}{1.112\,\ensuremath{\pm}\,0.006}  % experiments/FIGS/data/v2/index.json :: nonlinear probe bbqq_vs_ccqq: (1 - AUC) of R42_Q1 over L188, run by run :: sha256 c6c4138ee5c86423
\newcommand{\MlstMlpCmpBbqqMidOverFullPr}{11\% more errors}  % experiments/FIGS/data/v2/index.json :: nonlinear probe bbqq_vs_ccqq: (1 - AUC) of R42_Q1 over L188, run by run :: sha256 c6c4138ee5c86423 :: prose twin: two significant figures, direction in words
\newcommand{\MlstMlpCmpBbqqCoarseOverFull}{2.93\,\ensuremath{\pm}\,0.17}  % experiments/FIGS/data/v2/index.json :: nonlinear probe bbqq_vs_ccqq: (1 - AUC) of R16_Q1 over L188, run by run :: sha256 c6c4138ee5c86423
\newcommand{\MlstMlpCmpBbqqCoarseOverFullPr}{2.9 times the errors}  % experiments/FIGS/data/v2/index.json :: nonlinear probe bbqq_vs_ccqq: (1 - AUC) of R16_Q1 over L188, run by run :: sha256 c6c4138ee5c86423 :: prose twin: two significant figures, direction in words
\newcommand{\MlstMlpCmpBbqqResMassOverRes}{\pending{no run yet}}  % experiments/FIGS/data/v2/index.json :: nonlinear probe bbqq_vs_ccqq: (1 - AUC) of L162_MASS over L162, run by run :: sha256 c6c4138ee5c86423
\newcommand{\MlstMlpCmpBbqqResMassOverResPr}{\pending{no run yet}}  % experiments/FIGS/data/v2/index.json :: nonlinear probe bbqq_vs_ccqq: (1 - AUC) of L162_MASS over L162, run by run :: sha256 c6c4138ee5c86423 :: prose twin: two significant figures, direction in words
\newcommand{\MlstMlpCmpBbqqResMassOverResPrMax}{\pending{no run yet}}  % experiments/FIGS/data/v2/index.json :: nonlinear probe bbqq_vs_ccqq: (1 - AUC) of L162_MASS over L162, run by run :: sha256 c6c4138ee5c86423 :: prose twin of the run farther from 1
\newcommand{\MlstMlpCmpBbqqCoarseMassOverCoarse}{\pending{no run yet}}  % experiments/FIGS/data/v2/index.json :: nonlinear probe bbqq_vs_ccqq: (1 - AUC) of R16_Q1_MASS over R16_Q1, run by run :: sha256 c6c4138ee5c86423
\newcommand{\MlstMlpCmpBbqqCoarseMassOverCoarsePr}{\pending{no run yet}}  % experiments/FIGS/data/v2/index.json :: nonlinear probe bbqq_vs_ccqq: (1 - AUC) of R16_Q1_MASS over R16_Q1, run by run :: sha256 c6c4138ee5c86423 :: prose twin: two significant figures, direction in words
\newcommand{\MlstMlpCmpBbqqCoarseMassOverCoarsePrMax}{\pending{no run yet}}  % experiments/FIGS/data/v2/index.json :: nonlinear probe bbqq_vs_ccqq: (1 - AUC) of R16_Q1_MASS over R16_Q1, run by run :: sha256 c6c4138ee5c86423 :: prose twin of the run farther from 1
\newcommand{\MlstMlpCmpBbqqFzeroOverFull}{3.20 and 3.78}  % experiments/FIGS/data/v2/index.json :: nonlinear probe bbqq_vs_ccqq: (1 - AUC) of FLAV_F0 over L188, run by run :: sha256 c6c4138ee5c86423
\newcommand{\MlstMlpCmpBbqqFzeroOverFullPr}{at least 3.2 times the errors}  % experiments/FIGS/data/v2/index.json :: nonlinear probe bbqq_vs_ccqq: (1 - AUC) of FLAV_F0 over L188, run by run :: sha256 c6c4138ee5c86423 :: prose twin: two significant figures, direction in words
\newcommand{\MlstMlpCmpBbqqFzeroOverFullPrMax}{up to 3.8 times the errors}  % experiments/FIGS/data/v2/index.json :: nonlinear probe bbqq_vs_ccqq: (1 - AUC) of FLAV_F0 over L188, run by run :: sha256 c6c4138ee5c86423 :: prose twin of the run farther from 1
\newcommand{\MlstMlpCmpBbqqFoneOverFzero}{0.450 and 0.363}  % experiments/FIGS/data/v2/index.json :: nonlinear probe bbqq_vs_ccqq: (1 - AUC) of FLAV_F1 over FLAV_F0, run by run :: sha256 c6c4138ee5c86423
\newcommand{\MlstMlpCmpBbqqFoneOverFzeroPr}{at most 0.45 times the errors}  % experiments/FIGS/data/v2/index.json :: nonlinear probe bbqq_vs_ccqq: (1 - AUC) of FLAV_F1 over FLAV_F0, run by run :: sha256 c6c4138ee5c86423 :: prose twin: two significant figures, direction in words
\newcommand{\MlstMlpCmpBbqqFoneOverFzeroPrMax}{down to 0.36 times the errors}  % experiments/FIGS/data/v2/index.json :: nonlinear probe bbqq_vs_ccqq: (1 - AUC) of FLAV_F1 over FLAV_F0, run by run :: sha256 c6c4138ee5c86423 :: prose twin of the run farther from 1
\newcommand{\MlstMlpCmpBbqqFonerOverFone}{1.311 and 1.343}  % experiments/FIGS/data/v2/index.json :: nonlinear probe bbqq_vs_ccqq: (1 - AUC) of FLAV_F1R over FLAV_F1, run by run :: sha256 c6c4138ee5c86423
\newcommand{\MlstMlpCmpBbqqFonerOverFonePr}{at least 31\% more errors}  % experiments/FIGS/data/v2/index.json :: nonlinear probe bbqq_vs_ccqq: (1 - AUC) of FLAV_F1R over FLAV_F1, run by run :: sha256 c6c4138ee5c86423 :: prose twin: two significant figures, direction in words
\newcommand{\MlstMlpCmpBbqqFonerOverFonePrMax}{up to 34\% more errors}  % experiments/FIGS/data/v2/index.json :: nonlinear probe bbqq_vs_ccqq: (1 - AUC) of FLAV_F1R over FLAV_F1, run by run :: sha256 c6c4138ee5c86423 :: prose twin of the run farther from 1
\newcommand{\MlstMlpCmpBbqqFonerOverFzero}{0.59 and 0.49}  % experiments/FIGS/data/v2/index.json :: nonlinear probe bbqq_vs_ccqq: (1 - AUC) of FLAV_F1R over FLAV_F0, run by run :: sha256 c6c4138ee5c86423
\newcommand{\MlstMlpCmpBbqqFonerOverFzeroPr}{at least 41\% fewer errors}  % experiments/FIGS/data/v2/index.json :: nonlinear probe bbqq_vs_ccqq: (1 - AUC) of FLAV_F1R over FLAV_F0, run by run :: sha256 c6c4138ee5c86423 :: prose twin: two significant figures, direction in words
\newcommand{\MlstMlpCmpBbqqFonerOverFzeroPrMax}{down to 0.49 times the errors}  % experiments/FIGS/data/v2/index.json :: nonlinear probe bbqq_vs_ccqq: (1 - AUC) of FLAV_F1R over FLAV_F0, run by run :: sha256 c6c4138ee5c86423 :: prose twin of the run farther from 1
\newcommand{\MlstMlpCmpBbqqFzeroOverCoarse}{1.170 and 1.236}  % experiments/FIGS/data/v2/index.json :: nonlinear probe bbqq_vs_ccqq: (1 - AUC) of FLAV_F0 over R16_Q1, run by run :: sha256 c6c4138ee5c86423
\newcommand{\MlstMlpCmpBbqqFzeroOverCoarsePr}{at least 17\% more errors}  % experiments/FIGS/data/v2/index.json :: nonlinear probe bbqq_vs_ccqq: (1 - AUC) of FLAV_F0 over R16_Q1, run by run :: sha256 c6c4138ee5c86423 :: prose twin: two significant figures, direction in words
\newcommand{\MlstMlpCmpBbqqFzeroOverCoarsePrMax}{up to 24\% more errors}  % experiments/FIGS/data/v2/index.json :: nonlinear probe bbqq_vs_ccqq: (1 - AUC) of FLAV_F0 over R16_Q1, run by run :: sha256 c6c4138ee5c86423 :: prose twin of the run farther from 1
\newcommand{\MlstMlpCmpBbqqRandOverFull}{\pending{no run yet}}  % experiments/FIGS/data/v2/index.json :: nonlinear probe bbqq_vs_ccqq: (1 - AUC) of each random partition over the 188-class run of its seed :: sha256 c6c4138ee5c86423
\newcommand{\MlstMlpCmpBbqqRandOverFullPr}{\pending{no run yet}}  % experiments/FIGS/data/v2/index.json :: nonlinear probe bbqq_vs_ccqq: (1 - AUC) of each random partition over the 188-class run of its seed :: sha256 c6c4138ee5c86423 :: prose twin: two significant figures, direction in words
\newcommand{\MlstMlpCmpQcdBbResOverFull}{1.085\,\ensuremath{\pm}\,0.030}  % experiments/FIGS/data/v2/index.json :: nonlinear probe qcd_bb_vs_cc: (1 - AUC) of L162 over L188, run by run :: sha256 c6c4138ee5c86423
\newcommand{\MlstMlpCmpQcdBbResOverFullPr}{8.5\% more errors}  % experiments/FIGS/data/v2/index.json :: nonlinear probe qcd_bb_vs_cc: (1 - AUC) of L162 over L188, run by run :: sha256 c6c4138ee5c86423 :: prose twin: two significant figures, direction in words
\newcommand{\MlstMlpCmpQcdBbMidOverFull}{1.121\,\ensuremath{\pm}\,0.017}  % experiments/FIGS/data/v2/index.json :: nonlinear probe qcd_bb_vs_cc: (1 - AUC) of R42_Q1 over L188, run by run :: sha256 c6c4138ee5c86423
\newcommand{\MlstMlpCmpQcdBbMidOverFullPr}{12\% more errors}  % experiments/FIGS/data/v2/index.json :: nonlinear probe qcd_bb_vs_cc: (1 - AUC) of R42_Q1 over L188, run by run :: sha256 c6c4138ee5c86423 :: prose twin: two significant figures, direction in words
\newcommand{\MlstMlpCmpQcdBbCoarseOverFull}{2.31\,\ensuremath{\pm}\,0.04}  % experiments/FIGS/data/v2/index.json :: nonlinear probe qcd_bb_vs_cc: (1 - AUC) of R16_Q1 over L188, run by run :: sha256 c6c4138ee5c86423
\newcommand{\MlstMlpCmpQcdBbCoarseOverFullPr}{2.3 times the errors}  % experiments/FIGS/data/v2/index.json :: nonlinear probe qcd_bb_vs_cc: (1 - AUC) of R16_Q1 over L188, run by run :: sha256 c6c4138ee5c86423 :: prose twin: two significant figures, direction in words
\newcommand{\MlstMlpCmpQcdBbResMassOverRes}{\pending{no run yet}}  % experiments/FIGS/data/v2/index.json :: nonlinear probe qcd_bb_vs_cc: (1 - AUC) of L162_MASS over L162, run by run :: sha256 c6c4138ee5c86423
\newcommand{\MlstMlpCmpQcdBbResMassOverResPr}{\pending{no run yet}}  % experiments/FIGS/data/v2/index.json :: nonlinear probe qcd_bb_vs_cc: (1 - AUC) of L162_MASS over L162, run by run :: sha256 c6c4138ee5c86423 :: prose twin: two significant figures, direction in words
\newcommand{\MlstMlpCmpQcdBbResMassOverResPrMax}{\pending{no run yet}}  % experiments/FIGS/data/v2/index.json :: nonlinear probe qcd_bb_vs_cc: (1 - AUC) of L162_MASS over L162, run by run :: sha256 c6c4138ee5c86423 :: prose twin of the run farther from 1
\newcommand{\MlstMlpCmpQcdBbCoarseMassOverCoarse}{\pending{no run yet}}  % experiments/FIGS/data/v2/index.json :: nonlinear probe qcd_bb_vs_cc: (1 - AUC) of R16_Q1_MASS over R16_Q1, run by run :: sha256 c6c4138ee5c86423
\newcommand{\MlstMlpCmpQcdBbCoarseMassOverCoarsePr}{\pending{no run yet}}  % experiments/FIGS/data/v2/index.json :: nonlinear probe qcd_bb_vs_cc: (1 - AUC) of R16_Q1_MASS over R16_Q1, run by run :: sha256 c6c4138ee5c86423 :: prose twin: two significant figures, direction in words
\newcommand{\MlstMlpCmpQcdBbCoarseMassOverCoarsePrMax}{\pending{no run yet}}  % experiments/FIGS/data/v2/index.json :: nonlinear probe qcd_bb_vs_cc: (1 - AUC) of R16_Q1_MASS over R16_Q1, run by run :: sha256 c6c4138ee5c86423 :: prose twin of the run farther from 1
\newcommand{\MlstMlpCmpQcdBbFzeroOverFull}{2.50 and 2.72}  % experiments/FIGS/data/v2/index.json :: nonlinear probe qcd_bb_vs_cc: (1 - AUC) of FLAV_F0 over L188, run by run :: sha256 c6c4138ee5c86423
\newcommand{\MlstMlpCmpQcdBbFzeroOverFullPr}{at least 2.5 times the errors}  % experiments/FIGS/data/v2/index.json :: nonlinear probe qcd_bb_vs_cc: (1 - AUC) of FLAV_F0 over L188, run by run :: sha256 c6c4138ee5c86423 :: prose twin: two significant figures, direction in words
\newcommand{\MlstMlpCmpQcdBbFzeroOverFullPrMax}{up to 2.7 times the errors}  % experiments/FIGS/data/v2/index.json :: nonlinear probe qcd_bb_vs_cc: (1 - AUC) of FLAV_F0 over L188, run by run :: sha256 c6c4138ee5c86423 :: prose twin of the run farther from 1
\newcommand{\MlstMlpCmpQcdBbFoneOverFzero}{0.573 and 0.505}  % experiments/FIGS/data/v2/index.json :: nonlinear probe qcd_bb_vs_cc: (1 - AUC) of FLAV_F1 over FLAV_F0, run by run :: sha256 c6c4138ee5c86423
\newcommand{\MlstMlpCmpQcdBbFoneOverFzeroPr}{at least 43\% fewer errors}  % experiments/FIGS/data/v2/index.json :: nonlinear probe qcd_bb_vs_cc: (1 - AUC) of FLAV_F1 over FLAV_F0, run by run :: sha256 c6c4138ee5c86423 :: prose twin: two significant figures, direction in words
\newcommand{\MlstMlpCmpQcdBbFoneOverFzeroPrMax}{up to 49\% fewer errors}  % experiments/FIGS/data/v2/index.json :: nonlinear probe qcd_bb_vs_cc: (1 - AUC) of FLAV_F1 over FLAV_F0, run by run :: sha256 c6c4138ee5c86423 :: prose twin of the run farther from 1
\newcommand{\MlstMlpCmpQcdBbFonerOverFone}{1.219 and 1.258}  % experiments/FIGS/data/v2/index.json :: nonlinear probe qcd_bb_vs_cc: (1 - AUC) of FLAV_F1R over FLAV_F1, run by run :: sha256 c6c4138ee5c86423
\newcommand{\MlstMlpCmpQcdBbFonerOverFonePr}{at least 22\% more errors}  % experiments/FIGS/data/v2/index.json :: nonlinear probe qcd_bb_vs_cc: (1 - AUC) of FLAV_F1R over FLAV_F1, run by run :: sha256 c6c4138ee5c86423 :: prose twin: two significant figures, direction in words
\newcommand{\MlstMlpCmpQcdBbFonerOverFonePrMax}{up to 26\% more errors}  % experiments/FIGS/data/v2/index.json :: nonlinear probe qcd_bb_vs_cc: (1 - AUC) of FLAV_F1R over FLAV_F1, run by run :: sha256 c6c4138ee5c86423 :: prose twin of the run farther from 1
\newcommand{\MlstMlpCmpQcdBbFonerOverFzero}{0.699 and 0.635}  % experiments/FIGS/data/v2/index.json :: nonlinear probe qcd_bb_vs_cc: (1 - AUC) of FLAV_F1R over FLAV_F0, run by run :: sha256 c6c4138ee5c86423
\newcommand{\MlstMlpCmpQcdBbFonerOverFzeroPr}{at least 30\% fewer errors}  % experiments/FIGS/data/v2/index.json :: nonlinear probe qcd_bb_vs_cc: (1 - AUC) of FLAV_F1R over FLAV_F0, run by run :: sha256 c6c4138ee5c86423 :: prose twin: two significant figures, direction in words
\newcommand{\MlstMlpCmpQcdBbFonerOverFzeroPrMax}{up to 36\% fewer errors}  % experiments/FIGS/data/v2/index.json :: nonlinear probe qcd_bb_vs_cc: (1 - AUC) of FLAV_F1R over FLAV_F0, run by run :: sha256 c6c4138ee5c86423 :: prose twin of the run farther from 1
\newcommand{\MlstMlpCmpQcdBbFzeroOverCoarse}{1.103 and 1.167}  % experiments/FIGS/data/v2/index.json :: nonlinear probe qcd_bb_vs_cc: (1 - AUC) of FLAV_F0 over R16_Q1, run by run :: sha256 c6c4138ee5c86423
\newcommand{\MlstMlpCmpQcdBbFzeroOverCoarsePr}{at least 10\% more errors}  % experiments/FIGS/data/v2/index.json :: nonlinear probe qcd_bb_vs_cc: (1 - AUC) of FLAV_F0 over R16_Q1, run by run :: sha256 c6c4138ee5c86423 :: prose twin: two significant figures, direction in words
\newcommand{\MlstMlpCmpQcdBbFzeroOverCoarsePrMax}{up to 17\% more errors}  % experiments/FIGS/data/v2/index.json :: nonlinear probe qcd_bb_vs_cc: (1 - AUC) of FLAV_F0 over R16_Q1, run by run :: sha256 c6c4138ee5c86423 :: prose twin of the run farther from 1
\newcommand{\MlstMlpCmpQcdBbRandOverFull}{\pending{no run yet}}  % experiments/FIGS/data/v2/index.json :: nonlinear probe qcd_bb_vs_cc: (1 - AUC) of each random partition over the 188-class run of its seed :: sha256 c6c4138ee5c86423
\newcommand{\MlstMlpCmpQcdBbRandOverFullPr}{\pending{no run yet}}  % experiments/FIGS/data/v2/index.json :: nonlinear probe qcd_bb_vs_cc: (1 - AUC) of each random partition over the 188-class run of its seed :: sha256 c6c4138ee5c86423 :: prose twin: two significant figures, direction in words
\newcommand{\MlstMlpCmpQcdBResOverFull}{1.050\,\ensuremath{\pm}\,0.004}  % experiments/FIGS/data/v2/index.json :: nonlinear probe qcd_b_vs_light: (1 - AUC) of L162 over L188, run by run :: sha256 c6c4138ee5c86423
\newcommand{\MlstMlpCmpQcdBResOverFullPr}{5.0\% more errors}  % experiments/FIGS/data/v2/index.json :: nonlinear probe qcd_b_vs_light: (1 - AUC) of L162 over L188, run by run :: sha256 c6c4138ee5c86423 :: prose twin: two significant figures, direction in words
\newcommand{\MlstMlpCmpQcdBMidOverFull}{1.055\,\ensuremath{\pm}\,0.005}  % experiments/FIGS/data/v2/index.json :: nonlinear probe qcd_b_vs_light: (1 - AUC) of R42_Q1 over L188, run by run :: sha256 c6c4138ee5c86423
\newcommand{\MlstMlpCmpQcdBMidOverFullPr}{5.5\% more errors}  % experiments/FIGS/data/v2/index.json :: nonlinear probe qcd_b_vs_light: (1 - AUC) of R42_Q1 over L188, run by run :: sha256 c6c4138ee5c86423 :: prose twin: two significant figures, direction in words
\newcommand{\MlstMlpCmpQcdBCoarseOverFull}{1.337\,\ensuremath{\pm}\,0.021}  % experiments/FIGS/data/v2/index.json :: nonlinear probe qcd_b_vs_light: (1 - AUC) of R16_Q1 over L188, run by run :: sha256 c6c4138ee5c86423
\newcommand{\MlstMlpCmpQcdBCoarseOverFullPr}{34\% more errors}  % experiments/FIGS/data/v2/index.json :: nonlinear probe qcd_b_vs_light: (1 - AUC) of R16_Q1 over L188, run by run :: sha256 c6c4138ee5c86423 :: prose twin: two significant figures, direction in words
\newcommand{\MlstMlpCmpQcdBResMassOverRes}{\pending{no run yet}}  % experiments/FIGS/data/v2/index.json :: nonlinear probe qcd_b_vs_light: (1 - AUC) of L162_MASS over L162, run by run :: sha256 c6c4138ee5c86423
\newcommand{\MlstMlpCmpQcdBResMassOverResPr}{\pending{no run yet}}  % experiments/FIGS/data/v2/index.json :: nonlinear probe qcd_b_vs_light: (1 - AUC) of L162_MASS over L162, run by run :: sha256 c6c4138ee5c86423 :: prose twin: two significant figures, direction in words
\newcommand{\MlstMlpCmpQcdBResMassOverResPrMax}{\pending{no run yet}}  % experiments/FIGS/data/v2/index.json :: nonlinear probe qcd_b_vs_light: (1 - AUC) of L162_MASS over L162, run by run :: sha256 c6c4138ee5c86423 :: prose twin of the run farther from 1
\newcommand{\MlstMlpCmpQcdBCoarseMassOverCoarse}{\pending{no run yet}}  % experiments/FIGS/data/v2/index.json :: nonlinear probe qcd_b_vs_light: (1 - AUC) of R16_Q1_MASS over R16_Q1, run by run :: sha256 c6c4138ee5c86423
\newcommand{\MlstMlpCmpQcdBCoarseMassOverCoarsePr}{\pending{no run yet}}  % experiments/FIGS/data/v2/index.json :: nonlinear probe qcd_b_vs_light: (1 - AUC) of R16_Q1_MASS over R16_Q1, run by run :: sha256 c6c4138ee5c86423 :: prose twin: two significant figures, direction in words
\newcommand{\MlstMlpCmpQcdBCoarseMassOverCoarsePrMax}{\pending{no run yet}}  % experiments/FIGS/data/v2/index.json :: nonlinear probe qcd_b_vs_light: (1 - AUC) of R16_Q1_MASS over R16_Q1, run by run :: sha256 c6c4138ee5c86423 :: prose twin of the run farther from 1
\newcommand{\MlstMlpCmpQcdBFzeroOverFull}{1.315 and 1.288}  % experiments/FIGS/data/v2/index.json :: nonlinear probe qcd_b_vs_light: (1 - AUC) of FLAV_F0 over L188, run by run :: sha256 c6c4138ee5c86423
\newcommand{\MlstMlpCmpQcdBFzeroOverFullPr}{at least 29\% more errors}  % experiments/FIGS/data/v2/index.json :: nonlinear probe qcd_b_vs_light: (1 - AUC) of FLAV_F0 over L188, run by run :: sha256 c6c4138ee5c86423 :: prose twin: two significant figures, direction in words
\newcommand{\MlstMlpCmpQcdBFzeroOverFullPrMax}{up to 32\% more errors}  % experiments/FIGS/data/v2/index.json :: nonlinear probe qcd_b_vs_light: (1 - AUC) of FLAV_F0 over L188, run by run :: sha256 c6c4138ee5c86423 :: prose twin of the run farther from 1
\newcommand{\MlstMlpCmpQcdBFoneOverFzero}{0.8804 and 0.8849}  % experiments/FIGS/data/v2/index.json :: nonlinear probe qcd_b_vs_light: (1 - AUC) of FLAV_F1 over FLAV_F0, run by run :: sha256 c6c4138ee5c86423
\newcommand{\MlstMlpCmpQcdBFoneOverFzeroPr}{at least 12\% fewer errors}  % experiments/FIGS/data/v2/index.json :: nonlinear probe qcd_b_vs_light: (1 - AUC) of FLAV_F1 over FLAV_F0, run by run :: sha256 c6c4138ee5c86423 :: prose twin: two significant figures, direction in words
\newcommand{\MlstMlpCmpQcdBFoneOverFzeroPrMax}{up to 12\% fewer errors}  % experiments/FIGS/data/v2/index.json :: nonlinear probe qcd_b_vs_light: (1 - AUC) of FLAV_F1 over FLAV_F0, run by run :: sha256 c6c4138ee5c86423 :: prose twin of the run farther from 1
\newcommand{\MlstMlpCmpQcdBFonerOverFone}{1.037 and 1.042}  % experiments/FIGS/data/v2/index.json :: nonlinear probe qcd_b_vs_light: (1 - AUC) of FLAV_F1R over FLAV_F1, run by run :: sha256 c6c4138ee5c86423
\newcommand{\MlstMlpCmpQcdBFonerOverFonePr}{at least 3.7\% more errors}  % experiments/FIGS/data/v2/index.json :: nonlinear probe qcd_b_vs_light: (1 - AUC) of FLAV_F1R over FLAV_F1, run by run :: sha256 c6c4138ee5c86423 :: prose twin: two significant figures, direction in words
\newcommand{\MlstMlpCmpQcdBFonerOverFonePrMax}{up to 4.2\% more errors}  % experiments/FIGS/data/v2/index.json :: nonlinear probe qcd_b_vs_light: (1 - AUC) of FLAV_F1R over FLAV_F1, run by run :: sha256 c6c4138ee5c86423 :: prose twin of the run farther from 1
\newcommand{\MlstMlpCmpQcdBFonerOverFzero}{0.9130 and 0.9217}  % experiments/FIGS/data/v2/index.json :: nonlinear probe qcd_b_vs_light: (1 - AUC) of FLAV_F1R over FLAV_F0, run by run :: sha256 c6c4138ee5c86423
\newcommand{\MlstMlpCmpQcdBFonerOverFzeroPr}{at least 7.8\% fewer errors}  % experiments/FIGS/data/v2/index.json :: nonlinear probe qcd_b_vs_light: (1 - AUC) of FLAV_F1R over FLAV_F0, run by run :: sha256 c6c4138ee5c86423 :: prose twin: two significant figures, direction in words
\newcommand{\MlstMlpCmpQcdBFonerOverFzeroPrMax}{up to 8.7\% fewer errors}  % experiments/FIGS/data/v2/index.json :: nonlinear probe qcd_b_vs_light: (1 - AUC) of FLAV_F1R over FLAV_F0, run by run :: sha256 c6c4138ee5c86423 :: prose twin of the run farther from 1
\newcommand{\MlstMlpCmpQcdBFzeroOverCoarse}{0.968 and 0.978}  % experiments/FIGS/data/v2/index.json :: nonlinear probe qcd_b_vs_light: (1 - AUC) of FLAV_F0 over R16_Q1, run by run :: sha256 c6c4138ee5c86423
\newcommand{\MlstMlpCmpQcdBFzeroOverCoarsePr}{at least 2.2\% fewer errors}  % experiments/FIGS/data/v2/index.json :: nonlinear probe qcd_b_vs_light: (1 - AUC) of FLAV_F0 over R16_Q1, run by run :: sha256 c6c4138ee5c86423 :: prose twin: two significant figures, direction in words
\newcommand{\MlstMlpCmpQcdBFzeroOverCoarsePrMax}{up to 3.2\% fewer errors}  % experiments/FIGS/data/v2/index.json :: nonlinear probe qcd_b_vs_light: (1 - AUC) of FLAV_F0 over R16_Q1, run by run :: sha256 c6c4138ee5c86423 :: prose twin of the run farther from 1
\newcommand{\MlstMlpCmpQcdBRandOverFull}{\pending{no run yet}}  % experiments/FIGS/data/v2/index.json :: nonlinear probe qcd_b_vs_light: (1 - AUC) of each random partition over the 188-class run of its seed :: sha256 c6c4138ee5c86423
\newcommand{\MlstMlpCmpQcdBRandOverFullPr}{\pending{no run yet}}  % experiments/FIGS/data/v2/index.json :: nonlinear probe qcd_b_vs_light: (1 - AUC) of each random partition over the 188-class run of its seed :: sha256 c6c4138ee5c86423 :: prose twin: two significant figures, direction in words
\newcommand{\MlstMlpCmpQcdCResOverFull}{1.039\,\ensuremath{\pm}\,0.005}  % experiments/FIGS/data/v2/index.json :: nonlinear probe qcd_c_vs_light: (1 - AUC) of L162 over L188, run by run :: sha256 c6c4138ee5c86423
\newcommand{\MlstMlpCmpQcdCResOverFullPr}{3.9\% more errors}  % experiments/FIGS/data/v2/index.json :: nonlinear probe qcd_c_vs_light: (1 - AUC) of L162 over L188, run by run :: sha256 c6c4138ee5c86423 :: prose twin: two significant figures, direction in words
\newcommand{\MlstMlpCmpQcdCMidOverFull}{1.055\,\ensuremath{\pm}\,0.006}  % experiments/FIGS/data/v2/index.json :: nonlinear probe qcd_c_vs_light: (1 - AUC) of R42_Q1 over L188, run by run :: sha256 c6c4138ee5c86423
\newcommand{\MlstMlpCmpQcdCMidOverFullPr}{5.5\% more errors}  % experiments/FIGS/data/v2/index.json :: nonlinear probe qcd_c_vs_light: (1 - AUC) of R42_Q1 over L188, run by run :: sha256 c6c4138ee5c86423 :: prose twin: two significant figures, direction in words
\newcommand{\MlstMlpCmpQcdCCoarseOverFull}{1.274\,\ensuremath{\pm}\,0.008}  % experiments/FIGS/data/v2/index.json :: nonlinear probe qcd_c_vs_light: (1 - AUC) of R16_Q1 over L188, run by run :: sha256 c6c4138ee5c86423
\newcommand{\MlstMlpCmpQcdCCoarseOverFullPr}{27\% more errors}  % experiments/FIGS/data/v2/index.json :: nonlinear probe qcd_c_vs_light: (1 - AUC) of R16_Q1 over L188, run by run :: sha256 c6c4138ee5c86423 :: prose twin: two significant figures, direction in words
\newcommand{\MlstMlpCmpQcdCResMassOverRes}{\pending{no run yet}}  % experiments/FIGS/data/v2/index.json :: nonlinear probe qcd_c_vs_light: (1 - AUC) of L162_MASS over L162, run by run :: sha256 c6c4138ee5c86423
\newcommand{\MlstMlpCmpQcdCResMassOverResPr}{\pending{no run yet}}  % experiments/FIGS/data/v2/index.json :: nonlinear probe qcd_c_vs_light: (1 - AUC) of L162_MASS over L162, run by run :: sha256 c6c4138ee5c86423 :: prose twin: two significant figures, direction in words
\newcommand{\MlstMlpCmpQcdCResMassOverResPrMax}{\pending{no run yet}}  % experiments/FIGS/data/v2/index.json :: nonlinear probe qcd_c_vs_light: (1 - AUC) of L162_MASS over L162, run by run :: sha256 c6c4138ee5c86423 :: prose twin of the run farther from 1
\newcommand{\MlstMlpCmpQcdCCoarseMassOverCoarse}{\pending{no run yet}}  % experiments/FIGS/data/v2/index.json :: nonlinear probe qcd_c_vs_light: (1 - AUC) of R16_Q1_MASS over R16_Q1, run by run :: sha256 c6c4138ee5c86423
\newcommand{\MlstMlpCmpQcdCCoarseMassOverCoarsePr}{\pending{no run yet}}  % experiments/FIGS/data/v2/index.json :: nonlinear probe qcd_c_vs_light: (1 - AUC) of R16_Q1_MASS over R16_Q1, run by run :: sha256 c6c4138ee5c86423 :: prose twin: two significant figures, direction in words
\newcommand{\MlstMlpCmpQcdCCoarseMassOverCoarsePrMax}{\pending{no run yet}}  % experiments/FIGS/data/v2/index.json :: nonlinear probe qcd_c_vs_light: (1 - AUC) of R16_Q1_MASS over R16_Q1, run by run :: sha256 c6c4138ee5c86423 :: prose twin of the run farther from 1
\newcommand{\MlstMlpCmpQcdCFzeroOverFull}{1.234 and 1.211}  % experiments/FIGS/data/v2/index.json :: nonlinear probe qcd_c_vs_light: (1 - AUC) of FLAV_F0 over L188, run by run :: sha256 c6c4138ee5c86423
\newcommand{\MlstMlpCmpQcdCFzeroOverFullPr}{at least 21\% more errors}  % experiments/FIGS/data/v2/index.json :: nonlinear probe qcd_c_vs_light: (1 - AUC) of FLAV_F0 over L188, run by run :: sha256 c6c4138ee5c86423 :: prose twin: two significant figures, direction in words
\newcommand{\MlstMlpCmpQcdCFzeroOverFullPrMax}{up to 23\% more errors}  % experiments/FIGS/data/v2/index.json :: nonlinear probe qcd_c_vs_light: (1 - AUC) of FLAV_F0 over L188, run by run :: sha256 c6c4138ee5c86423 :: prose twin of the run farther from 1
\newcommand{\MlstMlpCmpQcdCFoneOverFzero}{0.9395 and 0.9472}  % experiments/FIGS/data/v2/index.json :: nonlinear probe qcd_c_vs_light: (1 - AUC) of FLAV_F1 over FLAV_F0, run by run :: sha256 c6c4138ee5c86423
\newcommand{\MlstMlpCmpQcdCFoneOverFzeroPr}{at least 5.3\% fewer errors}  % experiments/FIGS/data/v2/index.json :: nonlinear probe qcd_c_vs_light: (1 - AUC) of FLAV_F1 over FLAV_F0, run by run :: sha256 c6c4138ee5c86423 :: prose twin: two significant figures, direction in words
\newcommand{\MlstMlpCmpQcdCFoneOverFzeroPrMax}{up to 6.0\% fewer errors}  % experiments/FIGS/data/v2/index.json :: nonlinear probe qcd_c_vs_light: (1 - AUC) of FLAV_F1 over FLAV_F0, run by run :: sha256 c6c4138ee5c86423 :: prose twin of the run farther from 1
\newcommand{\MlstMlpCmpQcdCFonerOverFone}{0.9846 and 0.9767}  % experiments/FIGS/data/v2/index.json :: nonlinear probe qcd_c_vs_light: (1 - AUC) of FLAV_F1R over FLAV_F1, run by run :: sha256 c6c4138ee5c86423
\newcommand{\MlstMlpCmpQcdCFonerOverFonePr}{at least 1.5\% fewer errors}  % experiments/FIGS/data/v2/index.json :: nonlinear probe qcd_c_vs_light: (1 - AUC) of FLAV_F1R over FLAV_F1, run by run :: sha256 c6c4138ee5c86423 :: prose twin: two significant figures, direction in words
\newcommand{\MlstMlpCmpQcdCFonerOverFonePrMax}{up to 2.3\% fewer errors}  % experiments/FIGS/data/v2/index.json :: nonlinear probe qcd_c_vs_light: (1 - AUC) of FLAV_F1R over FLAV_F1, run by run :: sha256 c6c4138ee5c86423 :: prose twin of the run farther from 1
\newcommand{\MlstMlpCmpQcdCFonerOverFzero}{0.9251 and 0.9251}  % experiments/FIGS/data/v2/index.json :: nonlinear probe qcd_c_vs_light: (1 - AUC) of FLAV_F1R over FLAV_F0, run by run :: sha256 c6c4138ee5c86423
\newcommand{\MlstMlpCmpQcdCFonerOverFzeroPr}{at least 7.5\% fewer errors}  % experiments/FIGS/data/v2/index.json :: nonlinear probe qcd_c_vs_light: (1 - AUC) of FLAV_F1R over FLAV_F0, run by run :: sha256 c6c4138ee5c86423 :: prose twin: two significant figures, direction in words
\newcommand{\MlstMlpCmpQcdCFonerOverFzeroPrMax}{up to 7.5\% fewer errors}  % experiments/FIGS/data/v2/index.json :: nonlinear probe qcd_c_vs_light: (1 - AUC) of FLAV_F1R over FLAV_F0, run by run :: sha256 c6c4138ee5c86423 :: prose twin of the run farther from 1
\newcommand{\MlstMlpCmpQcdCFzeroOverCoarse}{0.968 and 0.956}  % experiments/FIGS/data/v2/index.json :: nonlinear probe qcd_c_vs_light: (1 - AUC) of FLAV_F0 over R16_Q1, run by run :: sha256 c6c4138ee5c86423
\newcommand{\MlstMlpCmpQcdCFzeroOverCoarsePr}{at least 3.2\% fewer errors}  % experiments/FIGS/data/v2/index.json :: nonlinear probe qcd_c_vs_light: (1 - AUC) of FLAV_F0 over R16_Q1, run by run :: sha256 c6c4138ee5c86423 :: prose twin: two significant figures, direction in words
\newcommand{\MlstMlpCmpQcdCFzeroOverCoarsePrMax}{up to 4.4\% fewer errors}  % experiments/FIGS/data/v2/index.json :: nonlinear probe qcd_c_vs_light: (1 - AUC) of FLAV_F0 over R16_Q1, run by run :: sha256 c6c4138ee5c86423 :: prose twin of the run farther from 1
\newcommand{\MlstMlpCmpQcdCRandOverFull}{\pending{no run yet}}  % experiments/FIGS/data/v2/index.json :: nonlinear probe qcd_c_vs_light: (1 - AUC) of each random partition over the 188-class run of its seed :: sha256 c6c4138ee5c86423
\newcommand{\MlstMlpCmpQcdCRandOverFullPr}{\pending{no run yet}}  % experiments/FIGS/data/v2/index.json :: nonlinear probe qcd_c_vs_light: (1 - AUC) of each random partition over the 188-class run of its seed :: sha256 c6c4138ee5c86423 :: prose twin: two significant figures, direction in words
\newcommand{\MlstMlpTopRthirtyFullOverInitBn}{1.00\,\ensuremath{\pm}\,0.31}  % experiments/FIGS/data/v2/index.json :: nonlinear probe top, missing inputs zero: R30 of L188 over the BatchNorm-recomputed random network of the same index :: sha256 c6c4138ee5c86423
\newcommand{\MlstMlpTopRthirtyFullOverInitBnPr}{equal within 40\%}  % experiments/FIGS/data/v2/index.json :: nonlinear probe top, missing inputs zero: R30 of L188 over the BatchNorm-recomputed random network of the same index :: sha256 c6c4138ee5c86423 :: prose twin: two significant figures, direction in words
\newcommand{\MlstMlpTopRthirtyResOverInitBn}{0.94\,\ensuremath{\pm}\,0.26}  % experiments/FIGS/data/v2/index.json :: nonlinear probe top, missing inputs zero: R30 of L162 over the BatchNorm-recomputed random network of the same index :: sha256 c6c4138ee5c86423
\newcommand{\MlstMlpTopRthirtyResOverInitBnPr}{equal within 30\%}  % experiments/FIGS/data/v2/index.json :: nonlinear probe top, missing inputs zero: R30 of L162 over the BatchNorm-recomputed random network of the same index :: sha256 c6c4138ee5c86423 :: prose twin: two significant figures, direction in words
\newcommand{\MlstMlpTopRthirtyMidOverInitBn}{0.84\,\ensuremath{\pm}\,0.17}  % experiments/FIGS/data/v2/index.json :: nonlinear probe top, missing inputs zero: R30 of R42_Q1 over the BatchNorm-recomputed random network of the same index :: sha256 c6c4138ee5c86423
\newcommand{\MlstMlpTopRthirtyMidOverInitBnPr}{equal within 20\%}  % experiments/FIGS/data/v2/index.json :: nonlinear probe top, missing inputs zero: R30 of R42_Q1 over the BatchNorm-recomputed random network of the same index :: sha256 c6c4138ee5c86423 :: prose twin: two significant figures, direction in words
\newcommand{\MlstMlpTopRthirtyCoarseOverInitBn}{0.75\,\ensuremath{\pm}\,0.05}  % experiments/FIGS/data/v2/index.json :: nonlinear probe top, missing inputs zero: R30 of R16_Q1 over the BatchNorm-recomputed random network of the same index :: sha256 c6c4138ee5c86423
\newcommand{\MlstMlpTopRthirtyCoarseOverInitBnPr}{25\% less rejection}  % experiments/FIGS/data/v2/index.json :: nonlinear probe top, missing inputs zero: R30 of R16_Q1 over the BatchNorm-recomputed random network of the same index :: sha256 c6c4138ee5c86423 :: prose twin: two significant figures, direction in words
\newcommand{\MlstMlpQgRthirtyFullOverInitBn}{1.083\,\ensuremath{\pm}\,0.026}  % experiments/FIGS/data/v2/index.json :: nonlinear probe qg, missing inputs zero: R30 of L188 over the BatchNorm-recomputed random network of the same index :: sha256 c6c4138ee5c86423
\newcommand{\MlstMlpQgRthirtyFullOverInitBnPr}{8.3\% more rejection}  % experiments/FIGS/data/v2/index.json :: nonlinear probe qg, missing inputs zero: R30 of L188 over the BatchNorm-recomputed random network of the same index :: sha256 c6c4138ee5c86423 :: prose twin: two significant figures, direction in words
\newcommand{\MlstMlpQgRthirtyResOverInitBn}{1.010\,\ensuremath{\pm}\,0.016}  % experiments/FIGS/data/v2/index.json :: nonlinear probe qg, missing inputs zero: R30 of L162 over the BatchNorm-recomputed random network of the same index :: sha256 c6c4138ee5c86423
\newcommand{\MlstMlpQgRthirtyResOverInitBnPr}{equal within 2\%}  % experiments/FIGS/data/v2/index.json :: nonlinear probe qg, missing inputs zero: R30 of L162 over the BatchNorm-recomputed random network of the same index :: sha256 c6c4138ee5c86423 :: prose twin: two significant figures, direction in words
\newcommand{\MlstMlpQgRthirtyMidOverInitBn}{1.06\,\ensuremath{\pm}\,0.05}  % experiments/FIGS/data/v2/index.json :: nonlinear probe qg, missing inputs zero: R30 of R42_Q1 over the BatchNorm-recomputed random network of the same index :: sha256 c6c4138ee5c86423
\newcommand{\MlstMlpQgRthirtyMidOverInitBnPr}{equal within 7\%}  % experiments/FIGS/data/v2/index.json :: nonlinear probe qg, missing inputs zero: R30 of R42_Q1 over the BatchNorm-recomputed random network of the same index :: sha256 c6c4138ee5c86423 :: prose twin: two significant figures, direction in words
\newcommand{\MlstMlpQgRthirtyCoarseOverInitBn}{1.04\,\ensuremath{\pm}\,0.04}  % experiments/FIGS/data/v2/index.json :: nonlinear probe qg, missing inputs zero: R30 of R16_Q1 over the BatchNorm-recomputed random network of the same index :: sha256 c6c4138ee5c86423
\newcommand{\MlstMlpQgRthirtyCoarseOverInitBnPr}{equal within 5\%}  % experiments/FIGS/data/v2/index.json :: nonlinear probe qg, missing inputs zero: R30 of R16_Q1 over the BatchNorm-recomputed random network of the same index :: sha256 c6c4138ee5c86423 :: prose twin: two significant figures, direction in words
\newcommand{\MlstMlpQgHerwigRthirtyFullOverInitBn}{0.971\,\ensuremath{\pm}\,0.034}  % experiments/FIGS/data/v2/index.json :: nonlinear probe qg_herwig, missing inputs zero: R30 of L188 over the BatchNorm-recomputed random network of the same index :: sha256 c6c4138ee5c86423
\newcommand{\MlstMlpQgHerwigRthirtyFullOverInitBnPr}{equal within 4\%}  % experiments/FIGS/data/v2/index.json :: nonlinear probe qg_herwig, missing inputs zero: R30 of L188 over the BatchNorm-recomputed random network of the same index :: sha256 c6c4138ee5c86423 :: prose twin: two significant figures, direction in words
\newcommand{\MlstMlpQgHerwigRthirtyResOverInitBn}{0.94\,\ensuremath{\pm}\,0.05}  % experiments/FIGS/data/v2/index.json :: nonlinear probe qg_herwig, missing inputs zero: R30 of L162 over the BatchNorm-recomputed random network of the same index :: sha256 c6c4138ee5c86423
\newcommand{\MlstMlpQgHerwigRthirtyResOverInitBnPr}{equal within 6\%}  % experiments/FIGS/data/v2/index.json :: nonlinear probe qg_herwig, missing inputs zero: R30 of L162 over the BatchNorm-recomputed random network of the same index :: sha256 c6c4138ee5c86423 :: prose twin: two significant figures, direction in words
\newcommand{\MlstMlpQgHerwigRthirtyMidOverInitBn}{0.945\,\ensuremath{\pm}\,0.030}  % experiments/FIGS/data/v2/index.json :: nonlinear probe qg_herwig, missing inputs zero: R30 of R42_Q1 over the BatchNorm-recomputed random network of the same index :: sha256 c6c4138ee5c86423
\newcommand{\MlstMlpQgHerwigRthirtyMidOverInitBnPr}{equal within 6\%}  % experiments/FIGS/data/v2/index.json :: nonlinear probe qg_herwig, missing inputs zero: R30 of R42_Q1 over the BatchNorm-recomputed random network of the same index :: sha256 c6c4138ee5c86423 :: prose twin: two significant figures, direction in words
\newcommand{\MlstMlpQgHerwigRthirtyCoarseOverInitBn}{0.96\,\ensuremath{\pm}\,0.04}  % experiments/FIGS/data/v2/index.json :: nonlinear probe qg_herwig, missing inputs zero: R30 of R16_Q1 over the BatchNorm-recomputed random network of the same index :: sha256 c6c4138ee5c86423
\newcommand{\MlstMlpQgHerwigRthirtyCoarseOverInitBnPr}{equal within 5\%}  % experiments/FIGS/data/v2/index.json :: nonlinear probe qg_herwig, missing inputs zero: R30 of R16_Q1 over the BatchNorm-recomputed random network of the same index :: sha256 c6c4138ee5c86423 :: prose twin: two significant figures, direction in words
\newcommand{\MlstMlpNatJcTwoResMassOverRes}{1.016 and 1.009}  % experiments/FIGS/data/v2/index.json :: nonlinear probe, natural-composition JetClass-II (fits_v3_mlp2): (1 - AUC) of L162_MASS over L162, run by run :: sha256 c6c4138ee5c86423
\newcommand{\MlstMlpNatJcTwoResMassOverResPr}{equal within 2\%}  % experiments/FIGS/data/v2/index.json :: nonlinear probe, natural-composition JetClass-II (fits_v3_mlp2): (1 - AUC) of L162_MASS over L162, run by run :: sha256 c6c4138ee5c86423 :: prose twin: two significant figures, direction in words
\newcommand{\MlstMlpNatJcTwoResMassOverResPrMax}{equal within 2\%}  % experiments/FIGS/data/v2/index.json :: nonlinear probe, natural-composition JetClass-II (fits_v3_mlp2): (1 - AUC) of L162_MASS over L162, run by run :: sha256 c6c4138ee5c86423 :: prose twin of the run farther from 1
\newcommand{\MlstMlpNatJcTwoCoarseMassOverCoarse}{1.026 and 1.019}  % experiments/FIGS/data/v2/index.json :: nonlinear probe, natural-composition JetClass-II (fits_v3_mlp2): (1 - AUC) of R16_Q1_MASS over R16_Q1, run by run :: sha256 c6c4138ee5c86423
\newcommand{\MlstMlpNatJcTwoCoarseMassOverCoarsePr}{equal within 3\%}  % experiments/FIGS/data/v2/index.json :: nonlinear probe, natural-composition JetClass-II (fits_v3_mlp2): (1 - AUC) of R16_Q1_MASS over R16_Q1, run by run :: sha256 c6c4138ee5c86423 :: prose twin: two significant figures, direction in words
\newcommand{\MlstMlpNatJcTwoCoarseMassOverCoarsePrMax}{equal within 3\%}  % experiments/FIGS/data/v2/index.json :: nonlinear probe, natural-composition JetClass-II (fits_v3_mlp2): (1 - AUC) of R16_Q1_MASS over R16_Q1, run by run :: sha256 c6c4138ee5c86423 :: prose twin of the run farther from 1
\newcommand{\LpSubstitutionTop}{random draws}  % experiments/FIGS/data/v2/index.json :: linprobe/fits_v4_fill: lowest summed validation cross-entropy over every model, largest training set :: sha256 c6c4138ee5c86423 :: zero 8.8428; mean 8.99394; draws 7.00838
\newcommand{\LpSubstitutionQg}{the training mean}  % experiments/FIGS/data/v2/index.json :: linprobe/fits_v4_fill: lowest summed validation cross-entropy over every model, largest training set :: sha256 c6c4138ee5c86423 :: zero 13.6224; mean 13.1964; draws 13.3584
\newcommand{\MlstLpTopRthirtyFullOverInitBn}{4.67\,\ensuremath{\pm}\,0.23}  % experiments/FIGS/data/v2/index.json :: linear probe top, largest set, chosen substitution: R30 of L188 over the BatchNorm-recomputed random network of the same index :: sha256 c6c4138ee5c86423
\newcommand{\MlstLpTopRthirtyFullOverInitBnPr}{4.7 times the rejection}  % experiments/FIGS/data/v2/index.json :: linear probe top, largest set, chosen substitution: R30 of L188 over the BatchNorm-recomputed random network of the same index :: sha256 c6c4138ee5c86423 :: prose twin: two significant figures, direction in words
\newcommand{\MlstLpTopRthirtyFullOverInitBnZero}{1.04\,\ensuremath{\pm}\,0.05}  % experiments/FIGS/data/v2/index.json :: linear probe top, largest set, zero: R30 of L188 over the BatchNorm-recomputed random network of the same index :: sha256 c6c4138ee5c86423
\newcommand{\MlstLpTopRthirtyFullOverInitBnZeroPr}{equal within 5\%}  % experiments/FIGS/data/v2/index.json :: linear probe top, largest set, zero: R30 of L188 over the BatchNorm-recomputed random network of the same index :: sha256 c6c4138ee5c86423 :: prose twin: two significant figures, direction in words
\newcommand{\MlstLpTopRthirtyResOverInitBn}{4.08\,\ensuremath{\pm}\,0.28}  % experiments/FIGS/data/v2/index.json :: linear probe top, largest set, chosen substitution: R30 of L162 over the BatchNorm-recomputed random network of the same index :: sha256 c6c4138ee5c86423
\newcommand{\MlstLpTopRthirtyResOverInitBnPr}{4.1 times the rejection}  % experiments/FIGS/data/v2/index.json :: linear probe top, largest set, chosen substitution: R30 of L162 over the BatchNorm-recomputed random network of the same index :: sha256 c6c4138ee5c86423 :: prose twin: two significant figures, direction in words
\newcommand{\MlstLpTopRthirtyResOverInitBnZero}{1.33\,\ensuremath{\pm}\,0.30}  % experiments/FIGS/data/v2/index.json :: linear probe top, largest set, zero: R30 of L162 over the BatchNorm-recomputed random network of the same index :: sha256 c6c4138ee5c86423
\newcommand{\MlstLpTopRthirtyResOverInitBnZeroPr}{equal within 40\%}  % experiments/FIGS/data/v2/index.json :: linear probe top, largest set, zero: R30 of L162 over the BatchNorm-recomputed random network of the same index :: sha256 c6c4138ee5c86423 :: prose twin: two significant figures, direction in words
\newcommand{\MlstLpTopRthirtyMidOverInitBn}{4.76\,\ensuremath{\pm}\,0.15}  % experiments/FIGS/data/v2/index.json :: linear probe top, largest set, chosen substitution: R30 of R42_Q1 over the BatchNorm-recomputed random network of the same index :: sha256 c6c4138ee5c86423
\newcommand{\MlstLpTopRthirtyMidOverInitBnPr}{4.8 times the rejection}  % experiments/FIGS/data/v2/index.json :: linear probe top, largest set, chosen substitution: R30 of R42_Q1 over the BatchNorm-recomputed random network of the same index :: sha256 c6c4138ee5c86423 :: prose twin: two significant figures, direction in words
\newcommand{\MlstLpTopRthirtyMidOverInitBnZero}{1.01\,\ensuremath{\pm}\,0.28}  % experiments/FIGS/data/v2/index.json :: linear probe top, largest set, zero: R30 of R42_Q1 over the BatchNorm-recomputed random network of the same index :: sha256 c6c4138ee5c86423
\newcommand{\MlstLpTopRthirtyMidOverInitBnZeroPr}{equal within 30\%}  % experiments/FIGS/data/v2/index.json :: linear probe top, largest set, zero: R30 of R42_Q1 over the BatchNorm-recomputed random network of the same index :: sha256 c6c4138ee5c86423 :: prose twin: two significant figures, direction in words
\newcommand{\MlstLpTopRthirtyCoarseOverInitBn}{5.24\,\ensuremath{\pm}\,0.06}  % experiments/FIGS/data/v2/index.json :: linear probe top, largest set, chosen substitution: R30 of R16_Q1 over the BatchNorm-recomputed random network of the same index :: sha256 c6c4138ee5c86423
\newcommand{\MlstLpTopRthirtyCoarseOverInitBnPr}{5.2 times the rejection}  % experiments/FIGS/data/v2/index.json :: linear probe top, largest set, chosen substitution: R30 of R16_Q1 over the BatchNorm-recomputed random network of the same index :: sha256 c6c4138ee5c86423 :: prose twin: two significant figures, direction in words
\newcommand{\MlstLpTopRthirtyCoarseOverInitBnZero}{0.97\,\ensuremath{\pm}\,0.16}  % experiments/FIGS/data/v2/index.json :: linear probe top, largest set, zero: R30 of R16_Q1 over the BatchNorm-recomputed random network of the same index :: sha256 c6c4138ee5c86423
\newcommand{\MlstLpTopRthirtyCoarseOverInitBnZeroPr}{equal within 20\%}  % experiments/FIGS/data/v2/index.json :: linear probe top, largest set, zero: R30 of R16_Q1 over the BatchNorm-recomputed random network of the same index :: sha256 c6c4138ee5c86423 :: prose twin: two significant figures, direction in words
\newcommand{\MlstLpTopRthirtyResMassOverInitBn}{5.71 and 6.25}  % experiments/FIGS/data/v2/index.json :: linear probe top, largest set, chosen substitution: R30 of L162_MASS over the BatchNorm-recomputed random network of the same index :: sha256 c6c4138ee5c86423
\newcommand{\MlstLpTopRthirtyResMassOverInitBnPr}{at least 5.7 times the rejection}  % experiments/FIGS/data/v2/index.json :: linear probe top, largest set, chosen substitution: R30 of L162_MASS over the BatchNorm-recomputed random network of the same index :: sha256 c6c4138ee5c86423 :: prose twin: two significant figures, direction in words
\newcommand{\MlstLpTopRthirtyResMassOverInitBnPrMax}{up to 6.3 times the rejection}  % experiments/FIGS/data/v2/index.json :: linear probe top, largest set, chosen substitution: R30 of L162_MASS over the BatchNorm-recomputed random network of the same index :: sha256 c6c4138ee5c86423 :: prose twin of the run farther from 1
\newcommand{\MlstLpTopRthirtyResMassOverInitBnZero}{1.683 and 1.678}  % experiments/FIGS/data/v2/index.json :: linear probe top, largest set, zero: R30 of L162_MASS over the BatchNorm-recomputed random network of the same index :: sha256 c6c4138ee5c86423
\newcommand{\MlstLpTopRthirtyResMassOverInitBnZeroPr}{at least 68\% more rejection}  % experiments/FIGS/data/v2/index.json :: linear probe top, largest set, zero: R30 of L162_MASS over the BatchNorm-recomputed random network of the same index :: sha256 c6c4138ee5c86423 :: prose twin: two significant figures, direction in words
\newcommand{\MlstLpTopRthirtyResMassOverInitBnZeroPrMax}{up to 68\% more rejection}  % experiments/FIGS/data/v2/index.json :: linear probe top, largest set, zero: R30 of L162_MASS over the BatchNorm-recomputed random network of the same index :: sha256 c6c4138ee5c86423 :: prose twin of the run farther from 1
\newcommand{\MlstLpTopRthirtyCoarseMassOverInitBn}{6.73 and 7.72}  % experiments/FIGS/data/v2/index.json :: linear probe top, largest set, chosen substitution: R30 of R16_Q1_MASS over the BatchNorm-recomputed random network of the same index :: sha256 c6c4138ee5c86423
\newcommand{\MlstLpTopRthirtyCoarseMassOverInitBnPr}{at least 6.7 times the rejection}  % experiments/FIGS/data/v2/index.json :: linear probe top, largest set, chosen substitution: R30 of R16_Q1_MASS over the BatchNorm-recomputed random network of the same index :: sha256 c6c4138ee5c86423 :: prose twin: two significant figures, direction in words
\newcommand{\MlstLpTopRthirtyCoarseMassOverInitBnPrMax}{up to 7.7 times the rejection}  % experiments/FIGS/data/v2/index.json :: linear probe top, largest set, chosen substitution: R30 of R16_Q1_MASS over the BatchNorm-recomputed random network of the same index :: sha256 c6c4138ee5c86423 :: prose twin of the run farther from 1
\newcommand{\MlstLpTopRthirtyCoarseMassOverInitBnZero}{2.3 and 1.2}  % experiments/FIGS/data/v2/index.json :: linear probe top, largest set, zero: R30 of R16_Q1_MASS over the BatchNorm-recomputed random network of the same index :: sha256 c6c4138ee5c86423
\newcommand{\MlstLpTopRthirtyCoarseMassOverInitBnZeroPr}{equal within 200\%}  % experiments/FIGS/data/v2/index.json :: linear probe top, largest set, zero: R30 of R16_Q1_MASS over the BatchNorm-recomputed random network of the same index :: sha256 c6c4138ee5c86423 :: prose twin: two significant figures, direction in words
\newcommand{\MlstLpTopRthirtyCoarseMassOverInitBnZeroPrMax}{equal within 200\%}  % experiments/FIGS/data/v2/index.json :: linear probe top, largest set, zero: R30 of R16_Q1_MASS over the BatchNorm-recomputed random network of the same index :: sha256 c6c4138ee5c86423 :: prose twin of the run farther from 1
\newcommand{\MlstLpTopRthirtyResOverFull}{0.874\,\ensuremath{\pm}\,0.018}  % experiments/FIGS/data/v2/index.json :: linear probe top: R30 of L162 over the 188-class run :: sha256 c6c4138ee5c86423
\newcommand{\MlstLpTopRthirtyResOverFullPr}{13\% less rejection}  % experiments/FIGS/data/v2/index.json :: linear probe top: R30 of L162 over the 188-class run :: sha256 c6c4138ee5c86423 :: prose twin: two significant figures, direction in words
\newcommand{\MlstLpTopRthirtyResOverFullZero}{1.29\,\ensuremath{\pm}\,0.31}  % experiments/FIGS/data/v2/index.json :: linear probe top: R30 of L162 over the 188-class run :: sha256 c6c4138ee5c86423
\newcommand{\MlstLpTopRthirtyResOverFullZeroPr}{equal within 40\%}  % experiments/FIGS/data/v2/index.json :: linear probe top: R30 of L162 over the 188-class run :: sha256 c6c4138ee5c86423 :: prose twin: two significant figures, direction in words
\newcommand{\MlstLpTopRthirtyMidOverFull}{1.02\,\ensuremath{\pm}\,0.08}  % experiments/FIGS/data/v2/index.json :: linear probe top: R30 of R42_Q1 over the 188-class run :: sha256 c6c4138ee5c86423
\newcommand{\MlstLpTopRthirtyMidOverFullPr}{equal within 9\%}  % experiments/FIGS/data/v2/index.json :: linear probe top: R30 of R42_Q1 over the 188-class run :: sha256 c6c4138ee5c86423 :: prose twin: two significant figures, direction in words
\newcommand{\MlstLpTopRthirtyMidOverFullZero}{0.98\,\ensuremath{\pm}\,0.28}  % experiments/FIGS/data/v2/index.json :: linear probe top: R30 of R42_Q1 over the 188-class run :: sha256 c6c4138ee5c86423
\newcommand{\MlstLpTopRthirtyMidOverFullZeroPr}{equal within 30\%}  % experiments/FIGS/data/v2/index.json :: linear probe top: R30 of R42_Q1 over the 188-class run :: sha256 c6c4138ee5c86423 :: prose twin: two significant figures, direction in words
\newcommand{\MlstLpTopRthirtyCoarseOverFull}{1.12\,\ensuremath{\pm}\,0.06}  % experiments/FIGS/data/v2/index.json :: linear probe top: R30 of R16_Q1 over the 188-class run :: sha256 c6c4138ee5c86423
\newcommand{\MlstLpTopRthirtyCoarseOverFullPr}{12\% more rejection}  % experiments/FIGS/data/v2/index.json :: linear probe top: R30 of R16_Q1 over the 188-class run :: sha256 c6c4138ee5c86423 :: prose twin: two significant figures, direction in words
\newcommand{\MlstLpTopRthirtyCoarseOverFullZero}{0.93\,\ensuremath{\pm}\,0.15}  % experiments/FIGS/data/v2/index.json :: linear probe top: R30 of R16_Q1 over the 188-class run :: sha256 c6c4138ee5c86423
\newcommand{\MlstLpTopRthirtyCoarseOverFullZeroPr}{equal within 20\%}  % experiments/FIGS/data/v2/index.json :: linear probe top: R30 of R16_Q1 over the 188-class run :: sha256 c6c4138ee5c86423 :: prose twin: two significant figures, direction in words
\newcommand{\MlstLpTopRthirtyMeanOverZeroFull}{1.25\,\ensuremath{\pm}\,0.27}  % experiments/FIGS/data/v2/index.json :: linear probe top: R30 with mean over R30 with zero, L188 :: sha256 c6c4138ee5c86423
\newcommand{\MlstLpTopRthirtyMeanOverZeroFullPr}{equal within 30\%}  % experiments/FIGS/data/v2/index.json :: linear probe top: R30 with mean over R30 with zero, L188 :: sha256 c6c4138ee5c86423 :: prose twin: two significant figures, direction in words
\newcommand{\MlstLpTopRthirtyDrawOverZeroFull}{3.9\,\ensuremath{\pm}\,0.4}  % experiments/FIGS/data/v2/index.json :: linear probe top: R30 with draws over R30 with zero, L188 :: sha256 c6c4138ee5c86423
\newcommand{\MlstLpTopRthirtyDrawOverZeroFullPr}{3.9 times the rejection}  % experiments/FIGS/data/v2/index.json :: linear probe top: R30 with draws over R30 with zero, L188 :: sha256 c6c4138ee5c86423 :: prose twin: two significant figures, direction in words
\newcommand{\MlstLpTopRthirtyMeanOverZeroInitBn}{1.08\,\ensuremath{\pm}\,0.20}  % experiments/FIGS/data/v2/index.json :: linear probe top: R30 with mean over R30 with zero, INIT_BN :: sha256 c6c4138ee5c86423
\newcommand{\MlstLpTopRthirtyMeanOverZeroInitBnPr}{equal within 30\%}  % experiments/FIGS/data/v2/index.json :: linear probe top: R30 with mean over R30 with zero, INIT_BN :: sha256 c6c4138ee5c86423 :: prose twin: two significant figures, direction in words
\newcommand{\MlstLpTopRthirtyDrawOverZeroInitBn}{0.86\,\ensuremath{\pm}\,0.14}  % experiments/FIGS/data/v2/index.json :: linear probe top: R30 with draws over R30 with zero, INIT_BN :: sha256 c6c4138ee5c86423
\newcommand{\MlstLpTopRthirtyDrawOverZeroInitBnPr}{equal within 20\%}  % experiments/FIGS/data/v2/index.json :: linear probe top: R30 with draws over R30 with zero, INIT_BN :: sha256 c6c4138ee5c86423 :: prose twin: two significant figures, direction in words
\newcommand{\MlstLpTopRthirtyKinOverFull}{3.88}  % experiments/FIGS/data/v2/index.json :: linear probe top: R30 of the kinematic-only model over the mean of the 188-class runs (zero) :: sha256 c6c4138ee5c86423
\newcommand{\MlstLpTopRthirtyKinOverFullPr}{3.9 times the rejection}  % experiments/FIGS/data/v2/index.json :: linear probe top: R30 of the kinematic-only model over the mean of the 188-class runs (zero) :: sha256 c6c4138ee5c86423 :: prose twin: two significant figures, direction in words
\newcommand{\MlstLpTopRthirtyKin}{231}  % experiments/FIGS/data/v2/index.json :: linear probe top, largest set: R30 of the kinematic-only model (one run) :: sha256 c6c4138ee5c86423
\newcommand{\MlstLpTopRthirtyFullZero}{60\,\ensuremath{\pm}\,12}  % experiments/FIGS/data/v2/index.json :: linear probe top, largest set: R30 with zero, L188 :: sha256 c6c4138ee5c86423
\newcommand{\MlstLpTopRthirtyFullMean}{73\,\ensuremath{\pm}\,15}  % experiments/FIGS/data/v2/index.json :: linear probe top, largest set: R30 with mean, L188 :: sha256 c6c4138ee5c86423
\newcommand{\MlstLpTopRthirtyFullDraws}{228\,\ensuremath{\pm}\,25}  % experiments/FIGS/data/v2/index.json :: linear probe top, largest set: R30 with draws, L188 :: sha256 c6c4138ee5c86423
\newcommand{\MlstLpTopRthirtyInitBnZero}{58\,\ensuremath{\pm}\,13}  % experiments/FIGS/data/v2/index.json :: linear probe top, largest set: R30 with zero, INIT_BN :: sha256 c6c4138ee5c86423
\newcommand{\MlstLpTopRthirtyInitBnMean}{60.4\,\ensuremath{\pm}\,2.4}  % experiments/FIGS/data/v2/index.json :: linear probe top, largest set: R30 with mean, INIT_BN :: sha256 c6c4138ee5c86423
\newcommand{\MlstLpTopRthirtyInitBnDraws}{48.6\,\ensuremath{\pm}\,3.0}  % experiments/FIGS/data/v2/index.json :: linear probe top, largest set: R30 with draws, INIT_BN :: sha256 c6c4138ee5c86423
\newcommand{\MlstLpQgRthirtyFullOverInitBn}{1.39\,\ensuremath{\pm}\,0.06}  % experiments/FIGS/data/v2/index.json :: linear probe qg, largest set, chosen substitution: R30 of L188 over the BatchNorm-recomputed random network of the same index :: sha256 c6c4138ee5c86423
\newcommand{\MlstLpQgRthirtyFullOverInitBnPr}{39\% more rejection}  % experiments/FIGS/data/v2/index.json :: linear probe qg, largest set, chosen substitution: R30 of L188 over the BatchNorm-recomputed random network of the same index :: sha256 c6c4138ee5c86423 :: prose twin: two significant figures, direction in words
\newcommand{\MlstLpQgRthirtyFullOverInitBnZero}{1.21\,\ensuremath{\pm}\,0.15}  % experiments/FIGS/data/v2/index.json :: linear probe qg, largest set, zero: R30 of L188 over the BatchNorm-recomputed random network of the same index :: sha256 c6c4138ee5c86423
\newcommand{\MlstLpQgRthirtyFullOverInitBnZeroPr}{equal within 30\%}  % experiments/FIGS/data/v2/index.json :: linear probe qg, largest set, zero: R30 of L188 over the BatchNorm-recomputed random network of the same index :: sha256 c6c4138ee5c86423 :: prose twin: two significant figures, direction in words
\newcommand{\MlstLpQgRthirtyResOverInitBn}{1.360\,\ensuremath{\pm}\,0.025}  % experiments/FIGS/data/v2/index.json :: linear probe qg, largest set, chosen substitution: R30 of L162 over the BatchNorm-recomputed random network of the same index :: sha256 c6c4138ee5c86423
\newcommand{\MlstLpQgRthirtyResOverInitBnPr}{36\% more rejection}  % experiments/FIGS/data/v2/index.json :: linear probe qg, largest set, chosen substitution: R30 of L162 over the BatchNorm-recomputed random network of the same index :: sha256 c6c4138ee5c86423 :: prose twin: two significant figures, direction in words
\newcommand{\MlstLpQgRthirtyResOverInitBnZero}{1.07\,\ensuremath{\pm}\,0.13}  % experiments/FIGS/data/v2/index.json :: linear probe qg, largest set, zero: R30 of L162 over the BatchNorm-recomputed random network of the same index :: sha256 c6c4138ee5c86423
\newcommand{\MlstLpQgRthirtyResOverInitBnZeroPr}{equal within 20\%}  % experiments/FIGS/data/v2/index.json :: linear probe qg, largest set, zero: R30 of L162 over the BatchNorm-recomputed random network of the same index :: sha256 c6c4138ee5c86423 :: prose twin: two significant figures, direction in words
\newcommand{\MlstLpQgRthirtyMidOverInitBn}{1.357\,\ensuremath{\pm}\,0.033}  % experiments/FIGS/data/v2/index.json :: linear probe qg, largest set, chosen substitution: R30 of R42_Q1 over the BatchNorm-recomputed random network of the same index :: sha256 c6c4138ee5c86423
\newcommand{\MlstLpQgRthirtyMidOverInitBnPr}{36\% more rejection}  % experiments/FIGS/data/v2/index.json :: linear probe qg, largest set, chosen substitution: R30 of R42_Q1 over the BatchNorm-recomputed random network of the same index :: sha256 c6c4138ee5c86423 :: prose twin: two significant figures, direction in words
\newcommand{\MlstLpQgRthirtyMidOverInitBnZero}{1.12\,\ensuremath{\pm}\,0.11}  % experiments/FIGS/data/v2/index.json :: linear probe qg, largest set, zero: R30 of R42_Q1 over the BatchNorm-recomputed random network of the same index :: sha256 c6c4138ee5c86423
\newcommand{\MlstLpQgRthirtyMidOverInitBnZeroPr}{equal within 20\%}  % experiments/FIGS/data/v2/index.json :: linear probe qg, largest set, zero: R30 of R42_Q1 over the BatchNorm-recomputed random network of the same index :: sha256 c6c4138ee5c86423 :: prose twin: two significant figures, direction in words
\newcommand{\MlstLpQgRthirtyCoarseOverInitBn}{1.42\,\ensuremath{\pm}\,0.04}  % experiments/FIGS/data/v2/index.json :: linear probe qg, largest set, chosen substitution: R30 of R16_Q1 over the BatchNorm-recomputed random network of the same index :: sha256 c6c4138ee5c86423
\newcommand{\MlstLpQgRthirtyCoarseOverInitBnPr}{42\% more rejection}  % experiments/FIGS/data/v2/index.json :: linear probe qg, largest set, chosen substitution: R30 of R16_Q1 over the BatchNorm-recomputed random network of the same index :: sha256 c6c4138ee5c86423 :: prose twin: two significant figures, direction in words
\newcommand{\MlstLpQgRthirtyCoarseOverInitBnZero}{1.15\,\ensuremath{\pm}\,0.17}  % experiments/FIGS/data/v2/index.json :: linear probe qg, largest set, zero: R30 of R16_Q1 over the BatchNorm-recomputed random network of the same index :: sha256 c6c4138ee5c86423
\newcommand{\MlstLpQgRthirtyCoarseOverInitBnZeroPr}{equal within 20\%}  % experiments/FIGS/data/v2/index.json :: linear probe qg, largest set, zero: R30 of R16_Q1 over the BatchNorm-recomputed random network of the same index :: sha256 c6c4138ee5c86423 :: prose twin: two significant figures, direction in words
\newcommand{\MlstLpQgRthirtyResMassOverInitBn}{1.430 and 1.364}  % experiments/FIGS/data/v2/index.json :: linear probe qg, largest set, chosen substitution: R30 of L162_MASS over the BatchNorm-recomputed random network of the same index :: sha256 c6c4138ee5c86423
\newcommand{\MlstLpQgRthirtyResMassOverInitBnPr}{at least 36\% more rejection}  % experiments/FIGS/data/v2/index.json :: linear probe qg, largest set, chosen substitution: R30 of L162_MASS over the BatchNorm-recomputed random network of the same index :: sha256 c6c4138ee5c86423 :: prose twin: two significant figures, direction in words
\newcommand{\MlstLpQgRthirtyResMassOverInitBnPrMax}{up to 43\% more rejection}  % experiments/FIGS/data/v2/index.json :: linear probe qg, largest set, chosen substitution: R30 of L162_MASS over the BatchNorm-recomputed random network of the same index :: sha256 c6c4138ee5c86423 :: prose twin of the run farther from 1
\newcommand{\MlstLpQgRthirtyResMassOverInitBnZero}{1.213 and 1.175}  % experiments/FIGS/data/v2/index.json :: linear probe qg, largest set, zero: R30 of L162_MASS over the BatchNorm-recomputed random network of the same index :: sha256 c6c4138ee5c86423
\newcommand{\MlstLpQgRthirtyResMassOverInitBnZeroPr}{equal within 30\%}  % experiments/FIGS/data/v2/index.json :: linear probe qg, largest set, zero: R30 of L162_MASS over the BatchNorm-recomputed random network of the same index :: sha256 c6c4138ee5c86423 :: prose twin: two significant figures, direction in words
\newcommand{\MlstLpQgRthirtyResMassOverInitBnZeroPrMax}{equal within 30\%}  % experiments/FIGS/data/v2/index.json :: linear probe qg, largest set, zero: R30 of L162_MASS over the BatchNorm-recomputed random network of the same index :: sha256 c6c4138ee5c86423 :: prose twin of the run farther from 1
\newcommand{\MlstLpQgRthirtyCoarseMassOverInitBn}{1.324 and 1.394}  % experiments/FIGS/data/v2/index.json :: linear probe qg, largest set, chosen substitution: R30 of R16_Q1_MASS over the BatchNorm-recomputed random network of the same index :: sha256 c6c4138ee5c86423
\newcommand{\MlstLpQgRthirtyCoarseMassOverInitBnPr}{at least 32\% more rejection}  % experiments/FIGS/data/v2/index.json :: linear probe qg, largest set, chosen substitution: R30 of R16_Q1_MASS over the BatchNorm-recomputed random network of the same index :: sha256 c6c4138ee5c86423 :: prose twin: two significant figures, direction in words
\newcommand{\MlstLpQgRthirtyCoarseMassOverInitBnPrMax}{up to 39\% more rejection}  % experiments/FIGS/data/v2/index.json :: linear probe qg, largest set, chosen substitution: R30 of R16_Q1_MASS over the BatchNorm-recomputed random network of the same index :: sha256 c6c4138ee5c86423 :: prose twin of the run farther from 1
\newcommand{\MlstLpQgRthirtyCoarseMassOverInitBnZero}{1.224 and 1.261}  % experiments/FIGS/data/v2/index.json :: linear probe qg, largest set, zero: R30 of R16_Q1_MASS over the BatchNorm-recomputed random network of the same index :: sha256 c6c4138ee5c86423
\newcommand{\MlstLpQgRthirtyCoarseMassOverInitBnZeroPr}{equal within 30\%}  % experiments/FIGS/data/v2/index.json :: linear probe qg, largest set, zero: R30 of R16_Q1_MASS over the BatchNorm-recomputed random network of the same index :: sha256 c6c4138ee5c86423 :: prose twin: two significant figures, direction in words
\newcommand{\MlstLpQgRthirtyCoarseMassOverInitBnZeroPrMax}{equal within 30\%}  % experiments/FIGS/data/v2/index.json :: linear probe qg, largest set, zero: R30 of R16_Q1_MASS over the BatchNorm-recomputed random network of the same index :: sha256 c6c4138ee5c86423 :: prose twin of the run farther from 1
\newcommand{\MlstLpQgRthirtyResOverFull}{0.977\,\ensuremath{\pm}\,0.028}  % experiments/FIGS/data/v2/index.json :: linear probe qg: R30 of L162 over the 188-class run :: sha256 c6c4138ee5c86423
\newcommand{\MlstLpQgRthirtyResOverFullPr}{equal within 3\%}  % experiments/FIGS/data/v2/index.json :: linear probe qg: R30 of L162 over the 188-class run :: sha256 c6c4138ee5c86423 :: prose twin: two significant figures, direction in words
\newcommand{\MlstLpQgRthirtyResOverFullZero}{0.89\,\ensuremath{\pm}\,0.20}  % experiments/FIGS/data/v2/index.json :: linear probe qg: R30 of L162 over the 188-class run :: sha256 c6c4138ee5c86423
\newcommand{\MlstLpQgRthirtyResOverFullZeroPr}{equal within 20\%}  % experiments/FIGS/data/v2/index.json :: linear probe qg: R30 of L162 over the 188-class run :: sha256 c6c4138ee5c86423 :: prose twin: two significant figures, direction in words
\newcommand{\MlstLpQgRthirtyMidOverFull}{0.974\,\ensuremath{\pm}\,0.031}  % experiments/FIGS/data/v2/index.json :: linear probe qg: R30 of R42_Q1 over the 188-class run :: sha256 c6c4138ee5c86423
\newcommand{\MlstLpQgRthirtyMidOverFullPr}{equal within 4\%}  % experiments/FIGS/data/v2/index.json :: linear probe qg: R30 of R42_Q1 over the 188-class run :: sha256 c6c4138ee5c86423 :: prose twin: two significant figures, direction in words
\newcommand{\MlstLpQgRthirtyMidOverFullZero}{0.94\,\ensuremath{\pm}\,0.18}  % experiments/FIGS/data/v2/index.json :: linear probe qg: R30 of R42_Q1 over the 188-class run :: sha256 c6c4138ee5c86423
\newcommand{\MlstLpQgRthirtyMidOverFullZeroPr}{equal within 20\%}  % experiments/FIGS/data/v2/index.json :: linear probe qg: R30 of R42_Q1 over the 188-class run :: sha256 c6c4138ee5c86423 :: prose twin: two significant figures, direction in words
\newcommand{\MlstLpQgRthirtyCoarseOverFull}{1.02\,\ensuremath{\pm}\,0.06}  % experiments/FIGS/data/v2/index.json :: linear probe qg: R30 of R16_Q1 over the 188-class run :: sha256 c6c4138ee5c86423
\newcommand{\MlstLpQgRthirtyCoarseOverFullPr}{equal within 6\%}  % experiments/FIGS/data/v2/index.json :: linear probe qg: R30 of R16_Q1 over the 188-class run :: sha256 c6c4138ee5c86423 :: prose twin: two significant figures, direction in words
\newcommand{\MlstLpQgRthirtyCoarseOverFullZero}{0.95\,\ensuremath{\pm}\,0.09}  % experiments/FIGS/data/v2/index.json :: linear probe qg: R30 of R16_Q1 over the 188-class run :: sha256 c6c4138ee5c86423
\newcommand{\MlstLpQgRthirtyCoarseOverFullZeroPr}{equal within 9\%}  % experiments/FIGS/data/v2/index.json :: linear probe qg: R30 of R16_Q1 over the 188-class run :: sha256 c6c4138ee5c86423 :: prose twin: two significant figures, direction in words
\newcommand{\MlstLpQgRthirtyMeanOverZeroFull}{1.12\,\ensuremath{\pm}\,0.17}  % experiments/FIGS/data/v2/index.json :: linear probe qg: R30 with mean over R30 with zero, L188 :: sha256 c6c4138ee5c86423
\newcommand{\MlstLpQgRthirtyMeanOverZeroFullPr}{equal within 20\%}  % experiments/FIGS/data/v2/index.json :: linear probe qg: R30 with mean over R30 with zero, L188 :: sha256 c6c4138ee5c86423 :: prose twin: two significant figures, direction in words
\newcommand{\MlstLpQgRthirtyDrawOverZeroFull}{1.05\,\ensuremath{\pm}\,0.12}  % experiments/FIGS/data/v2/index.json :: linear probe qg: R30 with draws over R30 with zero, L188 :: sha256 c6c4138ee5c86423
\newcommand{\MlstLpQgRthirtyDrawOverZeroFullPr}{equal within 20\%}  % experiments/FIGS/data/v2/index.json :: linear probe qg: R30 with draws over R30 with zero, L188 :: sha256 c6c4138ee5c86423 :: prose twin: two significant figures, direction in words
\newcommand{\MlstLpQgRthirtyMeanOverZeroInitBn}{0.964\,\ensuremath{\pm}\,0.030}  % experiments/FIGS/data/v2/index.json :: linear probe qg: R30 with mean over R30 with zero, INIT_BN :: sha256 c6c4138ee5c86423
\newcommand{\MlstLpQgRthirtyMeanOverZeroInitBnPr}{equal within 4\%}  % experiments/FIGS/data/v2/index.json :: linear probe qg: R30 with mean over R30 with zero, INIT_BN :: sha256 c6c4138ee5c86423 :: prose twin: two significant figures, direction in words
\newcommand{\MlstLpQgRthirtyDrawOverZeroInitBn}{0.97\,\ensuremath{\pm}\,0.04}  % experiments/FIGS/data/v2/index.json :: linear probe qg: R30 with draws over R30 with zero, INIT_BN :: sha256 c6c4138ee5c86423
\newcommand{\MlstLpQgRthirtyDrawOverZeroInitBnPr}{equal within 4\%}  % experiments/FIGS/data/v2/index.json :: linear probe qg: R30 with draws over R30 with zero, INIT_BN :: sha256 c6c4138ee5c86423 :: prose twin: two significant figures, direction in words
\newcommand{\MlstLpQgRthirtyFullZero}{66\,\ensuremath{\pm}\,7}  % experiments/FIGS/data/v2/index.json :: linear probe qg, largest set: R30 with zero, L188 :: sha256 c6c4138ee5c86423
\newcommand{\MlstLpQgRthirtyFullMean}{72.7\,\ensuremath{\pm}\,2.4}  % experiments/FIGS/data/v2/index.json :: linear probe qg, largest set: R30 with mean, L188 :: sha256 c6c4138ee5c86423
\newcommand{\MlstLpQgRthirtyFullDraws}{68.1\,\ensuremath{\pm}\,1.8}  % experiments/FIGS/data/v2/index.json :: linear probe qg, largest set: R30 with draws, L188 :: sha256 c6c4138ee5c86423
\newcommand{\MlstLpQgRthirtyInitBnZero}{54.2\,\ensuremath{\pm}\,2.4}  % experiments/FIGS/data/v2/index.json :: linear probe qg, largest set: R30 with zero, INIT_BN :: sha256 c6c4138ee5c86423
\newcommand{\MlstLpQgRthirtyInitBnMean}{52.2\,\ensuremath{\pm}\,0.8}  % experiments/FIGS/data/v2/index.json :: linear probe qg, largest set: R30 with mean, INIT_BN :: sha256 c6c4138ee5c86423
\newcommand{\MlstLpQgRthirtyInitBnDraws}{52.4\,\ensuremath{\pm}\,0.4}  % experiments/FIGS/data/v2/index.json :: linear probe qg, largest set: R30 with draws, INIT_BN :: sha256 c6c4138ee5c86423
\newcommand{\MlstLpQgHerwigRthirtyFullOverInitBn}{1.14\,\ensuremath{\pm}\,0.10}  % experiments/FIGS/data/v2/index.json :: linear probe qg_herwig, largest set, chosen substitution: R30 of L188 over the BatchNorm-recomputed random network of the same index :: sha256 c6c4138ee5c86423
\newcommand{\MlstLpQgHerwigRthirtyFullOverInitBnPr}{equal within 20\%}  % experiments/FIGS/data/v2/index.json :: linear probe qg_herwig, largest set, chosen substitution: R30 of L188 over the BatchNorm-recomputed random network of the same index :: sha256 c6c4138ee5c86423 :: prose twin: two significant figures, direction in words
\newcommand{\MlstLpQgHerwigRthirtyFullOverInitBnZero}{1.054\,\ensuremath{\pm}\,0.028}  % experiments/FIGS/data/v2/index.json :: linear probe qg_herwig, largest set, zero: R30 of L188 over the BatchNorm-recomputed random network of the same index :: sha256 c6c4138ee5c86423
\newcommand{\MlstLpQgHerwigRthirtyFullOverInitBnZeroPr}{equal within 6\%}  % experiments/FIGS/data/v2/index.json :: linear probe qg_herwig, largest set, zero: R30 of L188 over the BatchNorm-recomputed random network of the same index :: sha256 c6c4138ee5c86423 :: prose twin: two significant figures, direction in words
\newcommand{\MlstLpQgHerwigRthirtyResOverInitBn}{1.17\,\ensuremath{\pm}\,0.04}  % experiments/FIGS/data/v2/index.json :: linear probe qg_herwig, largest set, chosen substitution: R30 of L162 over the BatchNorm-recomputed random network of the same index :: sha256 c6c4138ee5c86423
\newcommand{\MlstLpQgHerwigRthirtyResOverInitBnPr}{17\% more rejection}  % experiments/FIGS/data/v2/index.json :: linear probe qg_herwig, largest set, chosen substitution: R30 of L162 over the BatchNorm-recomputed random network of the same index :: sha256 c6c4138ee5c86423 :: prose twin: two significant figures, direction in words
\newcommand{\MlstLpQgHerwigRthirtyResOverInitBnZero}{0.98\,\ensuremath{\pm}\,0.10}  % experiments/FIGS/data/v2/index.json :: linear probe qg_herwig, largest set, zero: R30 of L162 over the BatchNorm-recomputed random network of the same index :: sha256 c6c4138ee5c86423
\newcommand{\MlstLpQgHerwigRthirtyResOverInitBnZeroPr}{equal within 10\%}  % experiments/FIGS/data/v2/index.json :: linear probe qg_herwig, largest set, zero: R30 of L162 over the BatchNorm-recomputed random network of the same index :: sha256 c6c4138ee5c86423 :: prose twin: two significant figures, direction in words
\newcommand{\MlstLpQgHerwigRthirtyMidOverInitBn}{1.12\,\ensuremath{\pm}\,0.07}  % experiments/FIGS/data/v2/index.json :: linear probe qg_herwig, largest set, chosen substitution: R30 of R42_Q1 over the BatchNorm-recomputed random network of the same index :: sha256 c6c4138ee5c86423
\newcommand{\MlstLpQgHerwigRthirtyMidOverInitBnPr}{equal within 20\%}  % experiments/FIGS/data/v2/index.json :: linear probe qg_herwig, largest set, chosen substitution: R30 of R42_Q1 over the BatchNorm-recomputed random network of the same index :: sha256 c6c4138ee5c86423 :: prose twin: two significant figures, direction in words
\newcommand{\MlstLpQgHerwigRthirtyMidOverInitBnZero}{0.99\,\ensuremath{\pm}\,0.07}  % experiments/FIGS/data/v2/index.json :: linear probe qg_herwig, largest set, zero: R30 of R42_Q1 over the BatchNorm-recomputed random network of the same index :: sha256 c6c4138ee5c86423
\newcommand{\MlstLpQgHerwigRthirtyMidOverInitBnZeroPr}{equal within 7\%}  % experiments/FIGS/data/v2/index.json :: linear probe qg_herwig, largest set, zero: R30 of R42_Q1 over the BatchNorm-recomputed random network of the same index :: sha256 c6c4138ee5c86423 :: prose twin: two significant figures, direction in words
\newcommand{\MlstLpQgHerwigRthirtyCoarseOverInitBn}{1.154\,\ensuremath{\pm}\,0.015}  % experiments/FIGS/data/v2/index.json :: linear probe qg_herwig, largest set, chosen substitution: R30 of R16_Q1 over the BatchNorm-recomputed random network of the same index :: sha256 c6c4138ee5c86423
\newcommand{\MlstLpQgHerwigRthirtyCoarseOverInitBnPr}{15\% more rejection}  % experiments/FIGS/data/v2/index.json :: linear probe qg_herwig, largest set, chosen substitution: R30 of R16_Q1 over the BatchNorm-recomputed random network of the same index :: sha256 c6c4138ee5c86423 :: prose twin: two significant figures, direction in words
\newcommand{\MlstLpQgHerwigRthirtyCoarseOverInitBnZero}{1.07\,\ensuremath{\pm}\,0.05}  % experiments/FIGS/data/v2/index.json :: linear probe qg_herwig, largest set, zero: R30 of R16_Q1 over the BatchNorm-recomputed random network of the same index :: sha256 c6c4138ee5c86423
\newcommand{\MlstLpQgHerwigRthirtyCoarseOverInitBnZeroPr}{equal within 8\%}  % experiments/FIGS/data/v2/index.json :: linear probe qg_herwig, largest set, zero: R30 of R16_Q1 over the BatchNorm-recomputed random network of the same index :: sha256 c6c4138ee5c86423 :: prose twin: two significant figures, direction in words
\newcommand{\MlstLpQgHerwigRthirtyResMassOverInitBn}{1.152 and 1.138}  % experiments/FIGS/data/v2/index.json :: linear probe qg_herwig, largest set, chosen substitution: R30 of L162_MASS over the BatchNorm-recomputed random network of the same index :: sha256 c6c4138ee5c86423
\newcommand{\MlstLpQgHerwigRthirtyResMassOverInitBnPr}{equal within 20\%}  % experiments/FIGS/data/v2/index.json :: linear probe qg_herwig, largest set, chosen substitution: R30 of L162_MASS over the BatchNorm-recomputed random network of the same index :: sha256 c6c4138ee5c86423 :: prose twin: two significant figures, direction in words
\newcommand{\MlstLpQgHerwigRthirtyResMassOverInitBnPrMax}{equal within 20\%}  % experiments/FIGS/data/v2/index.json :: linear probe qg_herwig, largest set, chosen substitution: R30 of L162_MASS over the BatchNorm-recomputed random network of the same index :: sha256 c6c4138ee5c86423 :: prose twin of the run farther from 1
\newcommand{\MlstLpQgHerwigRthirtyResMassOverInitBnZero}{1.13 and 1.03}  % experiments/FIGS/data/v2/index.json :: linear probe qg_herwig, largest set, zero: R30 of L162_MASS over the BatchNorm-recomputed random network of the same index :: sha256 c6c4138ee5c86423
\newcommand{\MlstLpQgHerwigRthirtyResMassOverInitBnZeroPr}{equal within 20\%}  % experiments/FIGS/data/v2/index.json :: linear probe qg_herwig, largest set, zero: R30 of L162_MASS over the BatchNorm-recomputed random network of the same index :: sha256 c6c4138ee5c86423 :: prose twin: two significant figures, direction in words
\newcommand{\MlstLpQgHerwigRthirtyResMassOverInitBnZeroPrMax}{equal within 20\%}  % experiments/FIGS/data/v2/index.json :: linear probe qg_herwig, largest set, zero: R30 of L162_MASS over the BatchNorm-recomputed random network of the same index :: sha256 c6c4138ee5c86423 :: prose twin of the run farther from 1
\newcommand{\MlstLpQgHerwigRthirtyCoarseMassOverInitBn}{1.146 and 1.159}  % experiments/FIGS/data/v2/index.json :: linear probe qg_herwig, largest set, chosen substitution: R30 of R16_Q1_MASS over the BatchNorm-recomputed random network of the same index :: sha256 c6c4138ee5c86423
\newcommand{\MlstLpQgHerwigRthirtyCoarseMassOverInitBnPr}{equal within 20\%}  % experiments/FIGS/data/v2/index.json :: linear probe qg_herwig, largest set, chosen substitution: R30 of R16_Q1_MASS over the BatchNorm-recomputed random network of the same index :: sha256 c6c4138ee5c86423 :: prose twin: two significant figures, direction in words
\newcommand{\MlstLpQgHerwigRthirtyCoarseMassOverInitBnPrMax}{equal within 20\%}  % experiments/FIGS/data/v2/index.json :: linear probe qg_herwig, largest set, chosen substitution: R30 of R16_Q1_MASS over the BatchNorm-recomputed random network of the same index :: sha256 c6c4138ee5c86423 :: prose twin of the run farther from 1
\newcommand{\MlstLpQgHerwigRthirtyCoarseMassOverInitBnZero}{1.070 and 1.100}  % experiments/FIGS/data/v2/index.json :: linear probe qg_herwig, largest set, zero: R30 of R16_Q1_MASS over the BatchNorm-recomputed random network of the same index :: sha256 c6c4138ee5c86423
\newcommand{\MlstLpQgHerwigRthirtyCoarseMassOverInitBnZeroPr}{equal within 20\%}  % experiments/FIGS/data/v2/index.json :: linear probe qg_herwig, largest set, zero: R30 of R16_Q1_MASS over the BatchNorm-recomputed random network of the same index :: sha256 c6c4138ee5c86423 :: prose twin: two significant figures, direction in words
\newcommand{\MlstLpQgHerwigRthirtyCoarseMassOverInitBnZeroPrMax}{equal within 20\%}  % experiments/FIGS/data/v2/index.json :: linear probe qg_herwig, largest set, zero: R30 of R16_Q1_MASS over the BatchNorm-recomputed random network of the same index :: sha256 c6c4138ee5c86423 :: prose twin of the run farther from 1
\newcommand{\MlstLpQgHerwigRthirtyResOverFull}{1.03\,\ensuremath{\pm}\,0.11}  % experiments/FIGS/data/v2/index.json :: linear probe qg_herwig: R30 of L162 over the 188-class run :: sha256 c6c4138ee5c86423
\newcommand{\MlstLpQgHerwigRthirtyResOverFullPr}{equal within 20\%}  % experiments/FIGS/data/v2/index.json :: linear probe qg_herwig: R30 of L162 over the 188-class run :: sha256 c6c4138ee5c86423 :: prose twin: two significant figures, direction in words
\newcommand{\MlstLpQgHerwigRthirtyResOverFullZero}{0.93\,\ensuremath{\pm}\,0.11}  % experiments/FIGS/data/v2/index.json :: linear probe qg_herwig: R30 of L162 over the 188-class run :: sha256 c6c4138ee5c86423
\newcommand{\MlstLpQgHerwigRthirtyResOverFullZeroPr}{equal within 20\%}  % experiments/FIGS/data/v2/index.json :: linear probe qg_herwig: R30 of L162 over the 188-class run :: sha256 c6c4138ee5c86423 :: prose twin: two significant figures, direction in words
\newcommand{\MlstLpQgHerwigRthirtyMidOverFull}{0.98\,\ensuremath{\pm}\,0.06}  % experiments/FIGS/data/v2/index.json :: linear probe qg_herwig: R30 of R42_Q1 over the 188-class run :: sha256 c6c4138ee5c86423
\newcommand{\MlstLpQgHerwigRthirtyMidOverFullPr}{equal within 6\%}  % experiments/FIGS/data/v2/index.json :: linear probe qg_herwig: R30 of R42_Q1 over the 188-class run :: sha256 c6c4138ee5c86423 :: prose twin: two significant figures, direction in words
\newcommand{\MlstLpQgHerwigRthirtyMidOverFullZero}{0.94\,\ensuremath{\pm}\,0.09}  % experiments/FIGS/data/v2/index.json :: linear probe qg_herwig: R30 of R42_Q1 over the 188-class run :: sha256 c6c4138ee5c86423
\newcommand{\MlstLpQgHerwigRthirtyMidOverFullZeroPr}{equal within 9\%}  % experiments/FIGS/data/v2/index.json :: linear probe qg_herwig: R30 of R42_Q1 over the 188-class run :: sha256 c6c4138ee5c86423 :: prose twin: two significant figures, direction in words
\newcommand{\MlstLpQgHerwigRthirtyCoarseOverFull}{1.02\,\ensuremath{\pm}\,0.09}  % experiments/FIGS/data/v2/index.json :: linear probe qg_herwig: R30 of R16_Q1 over the 188-class run :: sha256 c6c4138ee5c86423
\newcommand{\MlstLpQgHerwigRthirtyCoarseOverFullPr}{equal within 10\%}  % experiments/FIGS/data/v2/index.json :: linear probe qg_herwig: R30 of R16_Q1 over the 188-class run :: sha256 c6c4138ee5c86423 :: prose twin: two significant figures, direction in words
\newcommand{\MlstLpQgHerwigRthirtyCoarseOverFullZero}{1.018\,\ensuremath{\pm}\,0.024}  % experiments/FIGS/data/v2/index.json :: linear probe qg_herwig: R30 of R16_Q1 over the 188-class run :: sha256 c6c4138ee5c86423
\newcommand{\MlstLpQgHerwigRthirtyCoarseOverFullZeroPr}{equal within 3\%}  % experiments/FIGS/data/v2/index.json :: linear probe qg_herwig: R30 of R16_Q1 over the 188-class run :: sha256 c6c4138ee5c86423 :: prose twin: two significant figures, direction in words
\newcommand{\MlstLpQgHerwigRthirtyMeanOverZeroFull}{1.07\,\ensuremath{\pm}\,0.11}  % experiments/FIGS/data/v2/index.json :: linear probe qg_herwig: R30 with mean over R30 with zero, L188 :: sha256 c6c4138ee5c86423
\newcommand{\MlstLpQgHerwigRthirtyMeanOverZeroFullPr}{equal within 20\%}  % experiments/FIGS/data/v2/index.json :: linear probe qg_herwig: R30 with mean over R30 with zero, L188 :: sha256 c6c4138ee5c86423 :: prose twin: two significant figures, direction in words
\newcommand{\MlstLpQgHerwigRthirtyDrawOverZeroFull}{1.07\,\ensuremath{\pm}\,0.06}  % experiments/FIGS/data/v2/index.json :: linear probe qg_herwig: R30 with draws over R30 with zero, L188 :: sha256 c6c4138ee5c86423
\newcommand{\MlstLpQgHerwigRthirtyDrawOverZeroFullPr}{equal within 7\%}  % experiments/FIGS/data/v2/index.json :: linear probe qg_herwig: R30 with draws over R30 with zero, L188 :: sha256 c6c4138ee5c86423 :: prose twin: two significant figures, direction in words
\newcommand{\MlstLpQgHerwigRthirtyMeanOverZeroInitBn}{0.988\,\ensuremath{\pm}\,0.019}  % experiments/FIGS/data/v2/index.json :: linear probe qg_herwig: R30 with mean over R30 with zero, INIT_BN :: sha256 c6c4138ee5c86423
\newcommand{\MlstLpQgHerwigRthirtyMeanOverZeroInitBnPr}{equal within 2\%}  % experiments/FIGS/data/v2/index.json :: linear probe qg_herwig: R30 with mean over R30 with zero, INIT_BN :: sha256 c6c4138ee5c86423 :: prose twin: two significant figures, direction in words
\newcommand{\MlstLpQgHerwigRthirtyDrawOverZeroInitBn}{0.983\,\ensuremath{\pm}\,0.032}  % experiments/FIGS/data/v2/index.json :: linear probe qg_herwig: R30 with draws over R30 with zero, INIT_BN :: sha256 c6c4138ee5c86423
\newcommand{\MlstLpQgHerwigRthirtyDrawOverZeroInitBnPr}{equal within 4\%}  % experiments/FIGS/data/v2/index.json :: linear probe qg_herwig: R30 with draws over R30 with zero, INIT_BN :: sha256 c6c4138ee5c86423 :: prose twin: two significant figures, direction in words
\newcommand{\MlstLpQgHerwigRthirtyFullZero}{32.8\,\ensuremath{\pm}\,1.0}  % experiments/FIGS/data/v2/index.json :: linear probe qg_herwig, largest set: R30 with zero, L188 :: sha256 c6c4138ee5c86423
\newcommand{\MlstLpQgHerwigRthirtyFullMean}{35.0\,\ensuremath{\pm}\,2.6}  % experiments/FIGS/data/v2/index.json :: linear probe qg_herwig, largest set: R30 with mean, L188 :: sha256 c6c4138ee5c86423
\newcommand{\MlstLpQgHerwigRthirtyFullDraws}{35.0\,\ensuremath{\pm}\,0.8}  % experiments/FIGS/data/v2/index.json :: linear probe qg_herwig, largest set: R30 with draws, L188 :: sha256 c6c4138ee5c86423
\newcommand{\MlstLpQgHerwigRthirtyInitBnZero}{31.1\,\ensuremath{\pm}\,0.5}  % experiments/FIGS/data/v2/index.json :: linear probe qg_herwig, largest set: R30 with zero, INIT_BN :: sha256 c6c4138ee5c86423
\newcommand{\MlstLpQgHerwigRthirtyInitBnMean}{30.7\,\ensuremath{\pm}\,0.5}  % experiments/FIGS/data/v2/index.json :: linear probe qg_herwig, largest set: R30 with mean, INIT_BN :: sha256 c6c4138ee5c86423
\newcommand{\MlstLpQgHerwigRthirtyInitBnDraws}{30.6\,\ensuremath{\pm}\,0.6}  % experiments/FIGS/data/v2/index.json :: linear probe qg_herwig, largest set: R30 with draws, INIT_BN :: sha256 c6c4138ee5c86423
\newcommand{\MlstLpJcTwoTopZeroAccDropFull}{88.9\,\ensuremath{\pm}\,1.0\%}  % experiments/FIGS/data/v2/index.json :: linear probe JetClass-II, largest set: relative drop in accuracy (%) with the top data's missing slots set to zero, L188 :: sha256 c6c4138ee5c86423
\newcommand{\MlstLpJcTwoTopZeroAccDropFullPr}{89\%}  % experiments/FIGS/data/v2/index.json :: linear probe JetClass-II, largest set: relative drop in accuracy (%) with the top data's missing slots set to zero, L188 :: sha256 c6c4138ee5c86423 :: prose twin, two significant figures
\newcommand{\MlstLpJcTwoTopMeanAccDropFull}{87.9\,\ensuremath{\pm}\,1.1\%}  % experiments/FIGS/data/v2/index.json :: linear probe JetClass-II, largest set: relative drop in accuracy (%) with the top data's missing slots set to mean, L188 :: sha256 c6c4138ee5c86423
\newcommand{\MlstLpJcTwoTopMeanAccDropFullPr}{88\%}  % experiments/FIGS/data/v2/index.json :: linear probe JetClass-II, largest set: relative drop in accuracy (%) with the top data's missing slots set to mean, L188 :: sha256 c6c4138ee5c86423 :: prose twin, two significant figures
\newcommand{\MlstLpJcTwoTopDrawAccDropFull}{85.98\,\ensuremath{\pm}\,0.07\%}  % experiments/FIGS/data/v2/index.json :: linear probe JetClass-II, largest set: relative drop in accuracy (%) with the top data's missing slots set to draws, L188 :: sha256 c6c4138ee5c86423
\newcommand{\MlstLpJcTwoTopDrawAccDropFullPr}{86\%}  % experiments/FIGS/data/v2/index.json :: linear probe JetClass-II, largest set: relative drop in accuracy (%) with the top data's missing slots set to draws, L188 :: sha256 c6c4138ee5c86423 :: prose twin, two significant figures
\newcommand{\MlstLpJcTwoTopZeroAccDropRes}{89.0\,\ensuremath{\pm}\,0.9\%}  % experiments/FIGS/data/v2/index.json :: linear probe JetClass-II, largest set: relative drop in accuracy (%) with the top data's missing slots set to zero, L162 :: sha256 c6c4138ee5c86423
\newcommand{\MlstLpJcTwoTopZeroAccDropResPr}{89\%}  % experiments/FIGS/data/v2/index.json :: linear probe JetClass-II, largest set: relative drop in accuracy (%) with the top data's missing slots set to zero, L162 :: sha256 c6c4138ee5c86423 :: prose twin, two significant figures
\newcommand{\MlstLpJcTwoTopMeanAccDropRes}{87.8\,\ensuremath{\pm}\,0.7\%}  % experiments/FIGS/data/v2/index.json :: linear probe JetClass-II, largest set: relative drop in accuracy (%) with the top data's missing slots set to mean, L162 :: sha256 c6c4138ee5c86423
\newcommand{\MlstLpJcTwoTopMeanAccDropResPr}{88\%}  % experiments/FIGS/data/v2/index.json :: linear probe JetClass-II, largest set: relative drop in accuracy (%) with the top data's missing slots set to mean, L162 :: sha256 c6c4138ee5c86423 :: prose twin, two significant figures
\newcommand{\MlstLpJcTwoTopDrawAccDropRes}{85.89\,\ensuremath{\pm}\,0.11\%}  % experiments/FIGS/data/v2/index.json :: linear probe JetClass-II, largest set: relative drop in accuracy (%) with the top data's missing slots set to draws, L162 :: sha256 c6c4138ee5c86423
\newcommand{\MlstLpJcTwoTopDrawAccDropResPr}{86\%}  % experiments/FIGS/data/v2/index.json :: linear probe JetClass-II, largest set: relative drop in accuracy (%) with the top data's missing slots set to draws, L162 :: sha256 c6c4138ee5c86423 :: prose twin, two significant figures
\newcommand{\MlstLpJcTwoTopZeroAccDropMid}{87.81\,\ensuremath{\pm}\,0.28\%}  % experiments/FIGS/data/v2/index.json :: linear probe JetClass-II, largest set: relative drop in accuracy (%) with the top data's missing slots set to zero, R42_Q1 :: sha256 c6c4138ee5c86423
\newcommand{\MlstLpJcTwoTopZeroAccDropMidPr}{88\%}  % experiments/FIGS/data/v2/index.json :: linear probe JetClass-II, largest set: relative drop in accuracy (%) with the top data's missing slots set to zero, R42_Q1 :: sha256 c6c4138ee5c86423 :: prose twin, two significant figures
\newcommand{\MlstLpJcTwoTopMeanAccDropMid}{86.7\,\ensuremath{\pm}\,0.9\%}  % experiments/FIGS/data/v2/index.json :: linear probe JetClass-II, largest set: relative drop in accuracy (%) with the top data's missing slots set to mean, R42_Q1 :: sha256 c6c4138ee5c86423
\newcommand{\MlstLpJcTwoTopMeanAccDropMidPr}{87\%}  % experiments/FIGS/data/v2/index.json :: linear probe JetClass-II, largest set: relative drop in accuracy (%) with the top data's missing slots set to mean, R42_Q1 :: sha256 c6c4138ee5c86423 :: prose twin, two significant figures
\newcommand{\MlstLpJcTwoTopDrawAccDropMid}{83.17\,\ensuremath{\pm}\,0.29\%}  % experiments/FIGS/data/v2/index.json :: linear probe JetClass-II, largest set: relative drop in accuracy (%) with the top data's missing slots set to draws, R42_Q1 :: sha256 c6c4138ee5c86423
\newcommand{\MlstLpJcTwoTopDrawAccDropMidPr}{83\%}  % experiments/FIGS/data/v2/index.json :: linear probe JetClass-II, largest set: relative drop in accuracy (%) with the top data's missing slots set to draws, R42_Q1 :: sha256 c6c4138ee5c86423 :: prose twin, two significant figures
\newcommand{\MlstLpJcTwoTopZeroAccDropCoarse}{86.7\,\ensuremath{\pm}\,0.7\%}  % experiments/FIGS/data/v2/index.json :: linear probe JetClass-II, largest set: relative drop in accuracy (%) with the top data's missing slots set to zero, R16_Q1 :: sha256 c6c4138ee5c86423
\newcommand{\MlstLpJcTwoTopZeroAccDropCoarsePr}{87\%}  % experiments/FIGS/data/v2/index.json :: linear probe JetClass-II, largest set: relative drop in accuracy (%) with the top data's missing slots set to zero, R16_Q1 :: sha256 c6c4138ee5c86423 :: prose twin, two significant figures
\newcommand{\MlstLpJcTwoTopMeanAccDropCoarse}{86.1\,\ensuremath{\pm}\,0.8\%}  % experiments/FIGS/data/v2/index.json :: linear probe JetClass-II, largest set: relative drop in accuracy (%) with the top data's missing slots set to mean, R16_Q1 :: sha256 c6c4138ee5c86423
\newcommand{\MlstLpJcTwoTopMeanAccDropCoarsePr}{86\%}  % experiments/FIGS/data/v2/index.json :: linear probe JetClass-II, largest set: relative drop in accuracy (%) with the top data's missing slots set to mean, R16_Q1 :: sha256 c6c4138ee5c86423 :: prose twin, two significant figures
\newcommand{\MlstLpJcTwoTopDrawAccDropCoarse}{81.4\,\ensuremath{\pm}\,0.5\%}  % experiments/FIGS/data/v2/index.json :: linear probe JetClass-II, largest set: relative drop in accuracy (%) with the top data's missing slots set to draws, R16_Q1 :: sha256 c6c4138ee5c86423
\newcommand{\MlstLpJcTwoTopDrawAccDropCoarsePr}{81\%}  % experiments/FIGS/data/v2/index.json :: linear probe JetClass-II, largest set: relative drop in accuracy (%) with the top data's missing slots set to draws, R16_Q1 :: sha256 c6c4138ee5c86423 :: prose twin, two significant figures
\newcommand{\MlstLpJcTwoTopZeroAccDropResMass}{89.2, 88.3\%}  % experiments/FIGS/data/v2/index.json :: linear probe JetClass-II, largest set: relative drop in accuracy (%) with the top data's missing slots set to zero, L162_MASS :: sha256 c6c4138ee5c86423
\newcommand{\MlstLpJcTwoTopZeroAccDropResMassPr}{88--89\%}  % experiments/FIGS/data/v2/index.json :: linear probe JetClass-II, largest set: relative drop in accuracy (%) with the top data's missing slots set to zero, L162_MASS :: sha256 c6c4138ee5c86423 :: prose twin, two significant figures
\newcommand{\MlstLpJcTwoTopMeanAccDropResMass}{88.5, 87.8\%}  % experiments/FIGS/data/v2/index.json :: linear probe JetClass-II, largest set: relative drop in accuracy (%) with the top data's missing slots set to mean, L162_MASS :: sha256 c6c4138ee5c86423
\newcommand{\MlstLpJcTwoTopMeanAccDropResMassPr}{88--88\%}  % experiments/FIGS/data/v2/index.json :: linear probe JetClass-II, largest set: relative drop in accuracy (%) with the top data's missing slots set to mean, L162_MASS :: sha256 c6c4138ee5c86423 :: prose twin, two significant figures
\newcommand{\MlstLpJcTwoTopDrawAccDropResMass}{85.4, 85.4\%}  % experiments/FIGS/data/v2/index.json :: linear probe JetClass-II, largest set: relative drop in accuracy (%) with the top data's missing slots set to draws, L162_MASS :: sha256 c6c4138ee5c86423
\newcommand{\MlstLpJcTwoTopDrawAccDropResMassPr}{85--85\%}  % experiments/FIGS/data/v2/index.json :: linear probe JetClass-II, largest set: relative drop in accuracy (%) with the top data's missing slots set to draws, L162_MASS :: sha256 c6c4138ee5c86423 :: prose twin, two significant figures
\newcommand{\MlstLpJcTwoTopZeroAccDropCoarseMass}{84.8, 87.4\%}  % experiments/FIGS/data/v2/index.json :: linear probe JetClass-II, largest set: relative drop in accuracy (%) with the top data's missing slots set to zero, R16_Q1_MASS :: sha256 c6c4138ee5c86423
\newcommand{\MlstLpJcTwoTopZeroAccDropCoarseMassPr}{85--87\%}  % experiments/FIGS/data/v2/index.json :: linear probe JetClass-II, largest set: relative drop in accuracy (%) with the top data's missing slots set to zero, R16_Q1_MASS :: sha256 c6c4138ee5c86423 :: prose twin, two significant figures
\newcommand{\MlstLpJcTwoTopMeanAccDropCoarseMass}{85.9, 86.9\%}  % experiments/FIGS/data/v2/index.json :: linear probe JetClass-II, largest set: relative drop in accuracy (%) with the top data's missing slots set to mean, R16_Q1_MASS :: sha256 c6c4138ee5c86423
\newcommand{\MlstLpJcTwoTopMeanAccDropCoarseMassPr}{86--87\%}  % experiments/FIGS/data/v2/index.json :: linear probe JetClass-II, largest set: relative drop in accuracy (%) with the top data's missing slots set to mean, R16_Q1_MASS :: sha256 c6c4138ee5c86423 :: prose twin, two significant figures
\newcommand{\MlstLpJcTwoTopDrawAccDropCoarseMass}{79.1, 81.7\%}  % experiments/FIGS/data/v2/index.json :: linear probe JetClass-II, largest set: relative drop in accuracy (%) with the top data's missing slots set to draws, R16_Q1_MASS :: sha256 c6c4138ee5c86423
\newcommand{\MlstLpJcTwoTopDrawAccDropCoarseMassPr}{79--82\%}  % experiments/FIGS/data/v2/index.json :: linear probe JetClass-II, largest set: relative drop in accuracy (%) with the top data's missing slots set to draws, R16_Q1_MASS :: sha256 c6c4138ee5c86423 :: prose twin, two significant figures
\newcommand{\MlstLpJcTwoTopZeroAccDropFzero}{88.9, 88.5\%}  % experiments/FIGS/data/v2/index.json :: linear probe JetClass-II, largest set: relative drop in accuracy (%) with the top data's missing slots set to zero, FLAV_F0 :: sha256 c6c4138ee5c86423
\newcommand{\MlstLpJcTwoTopZeroAccDropFzeroPr}{88--89\%}  % experiments/FIGS/data/v2/index.json :: linear probe JetClass-II, largest set: relative drop in accuracy (%) with the top data's missing slots set to zero, FLAV_F0 :: sha256 c6c4138ee5c86423 :: prose twin, two significant figures
\newcommand{\MlstLpJcTwoTopMeanAccDropFzero}{88.3, 87.0\%}  % experiments/FIGS/data/v2/index.json :: linear probe JetClass-II, largest set: relative drop in accuracy (%) with the top data's missing slots set to mean, FLAV_F0 :: sha256 c6c4138ee5c86423
\newcommand{\MlstLpJcTwoTopMeanAccDropFzeroPr}{87--88\%}  % experiments/FIGS/data/v2/index.json :: linear probe JetClass-II, largest set: relative drop in accuracy (%) with the top data's missing slots set to mean, FLAV_F0 :: sha256 c6c4138ee5c86423 :: prose twin, two significant figures
\newcommand{\MlstLpJcTwoTopDrawAccDropFzero}{83.3, 83.9\%}  % experiments/FIGS/data/v2/index.json :: linear probe JetClass-II, largest set: relative drop in accuracy (%) with the top data's missing slots set to draws, FLAV_F0 :: sha256 c6c4138ee5c86423
\newcommand{\MlstLpJcTwoTopDrawAccDropFzeroPr}{83--84\%}  % experiments/FIGS/data/v2/index.json :: linear probe JetClass-II, largest set: relative drop in accuracy (%) with the top data's missing slots set to draws, FLAV_F0 :: sha256 c6c4138ee5c86423 :: prose twin, two significant figures
\newcommand{\MlstLpJcTwoTopZeroAccDropFone}{87.8, 87.6\%}  % experiments/FIGS/data/v2/index.json :: linear probe JetClass-II, largest set: relative drop in accuracy (%) with the top data's missing slots set to zero, FLAV_F1 :: sha256 c6c4138ee5c86423
\newcommand{\MlstLpJcTwoTopZeroAccDropFonePr}{88--88\%}  % experiments/FIGS/data/v2/index.json :: linear probe JetClass-II, largest set: relative drop in accuracy (%) with the top data's missing slots set to zero, FLAV_F1 :: sha256 c6c4138ee5c86423 :: prose twin, two significant figures
\newcommand{\MlstLpJcTwoTopMeanAccDropFone}{88.9, 86.9\%}  % experiments/FIGS/data/v2/index.json :: linear probe JetClass-II, largest set: relative drop in accuracy (%) with the top data's missing slots set to mean, FLAV_F1 :: sha256 c6c4138ee5c86423
\newcommand{\MlstLpJcTwoTopMeanAccDropFonePr}{87--89\%}  % experiments/FIGS/data/v2/index.json :: linear probe JetClass-II, largest set: relative drop in accuracy (%) with the top data's missing slots set to mean, FLAV_F1 :: sha256 c6c4138ee5c86423 :: prose twin, two significant figures
\newcommand{\MlstLpJcTwoTopDrawAccDropFone}{84.4, 85.2\%}  % experiments/FIGS/data/v2/index.json :: linear probe JetClass-II, largest set: relative drop in accuracy (%) with the top data's missing slots set to draws, FLAV_F1 :: sha256 c6c4138ee5c86423
\newcommand{\MlstLpJcTwoTopDrawAccDropFonePr}{84--85\%}  % experiments/FIGS/data/v2/index.json :: linear probe JetClass-II, largest set: relative drop in accuracy (%) with the top data's missing slots set to draws, FLAV_F1 :: sha256 c6c4138ee5c86423 :: prose twin, two significant figures
\newcommand{\MlstLpJcTwoTopZeroAccDropFoner}{89.6, 90.3\%}  % experiments/FIGS/data/v2/index.json :: linear probe JetClass-II, largest set: relative drop in accuracy (%) with the top data's missing slots set to zero, FLAV_F1R :: sha256 c6c4138ee5c86423
\newcommand{\MlstLpJcTwoTopZeroAccDropFonerPr}{90--90\%}  % experiments/FIGS/data/v2/index.json :: linear probe JetClass-II, largest set: relative drop in accuracy (%) with the top data's missing slots set to zero, FLAV_F1R :: sha256 c6c4138ee5c86423 :: prose twin, two significant figures
\newcommand{\MlstLpJcTwoTopMeanAccDropFoner}{89.4, 89.7\%}  % experiments/FIGS/data/v2/index.json :: linear probe JetClass-II, largest set: relative drop in accuracy (%) with the top data's missing slots set to mean, FLAV_F1R :: sha256 c6c4138ee5c86423
\newcommand{\MlstLpJcTwoTopMeanAccDropFonerPr}{89--90\%}  % experiments/FIGS/data/v2/index.json :: linear probe JetClass-II, largest set: relative drop in accuracy (%) with the top data's missing slots set to mean, FLAV_F1R :: sha256 c6c4138ee5c86423 :: prose twin, two significant figures
\newcommand{\MlstLpJcTwoTopDrawAccDropFoner}{84.9, 85.6\%}  % experiments/FIGS/data/v2/index.json :: linear probe JetClass-II, largest set: relative drop in accuracy (%) with the top data's missing slots set to draws, FLAV_F1R :: sha256 c6c4138ee5c86423
\newcommand{\MlstLpJcTwoTopDrawAccDropFonerPr}{85--86\%}  % experiments/FIGS/data/v2/index.json :: linear probe JetClass-II, largest set: relative drop in accuracy (%) with the top data's missing slots set to draws, FLAV_F1R :: sha256 c6c4138ee5c86423 :: prose twin, two significant figures
\newcommand{\MlstLpJcTwoTopZeroAccDropInitBn}{56.4\,\ensuremath{\pm}\,1.4\%}  % experiments/FIGS/data/v2/index.json :: linear probe JetClass-II, largest set: relative drop in accuracy (%) with the top data's missing slots set to zero, INIT_BN :: sha256 c6c4138ee5c86423
\newcommand{\MlstLpJcTwoTopZeroAccDropInitBnPr}{56\%}  % experiments/FIGS/data/v2/index.json :: linear probe JetClass-II, largest set: relative drop in accuracy (%) with the top data's missing slots set to zero, INIT_BN :: sha256 c6c4138ee5c86423 :: prose twin, two significant figures
\newcommand{\MlstLpJcTwoTopMeanAccDropInitBn}{59.8\,\ensuremath{\pm}\,0.8\%}  % experiments/FIGS/data/v2/index.json :: linear probe JetClass-II, largest set: relative drop in accuracy (%) with the top data's missing slots set to mean, INIT_BN :: sha256 c6c4138ee5c86423
\newcommand{\MlstLpJcTwoTopMeanAccDropInitBnPr}{60\%}  % experiments/FIGS/data/v2/index.json :: linear probe JetClass-II, largest set: relative drop in accuracy (%) with the top data's missing slots set to mean, INIT_BN :: sha256 c6c4138ee5c86423 :: prose twin, two significant figures
\newcommand{\MlstLpJcTwoTopDrawAccDropInitBn}{59.8\,\ensuremath{\pm}\,1.2\%}  % experiments/FIGS/data/v2/index.json :: linear probe JetClass-II, largest set: relative drop in accuracy (%) with the top data's missing slots set to draws, INIT_BN :: sha256 c6c4138ee5c86423
\newcommand{\MlstLpJcTwoTopDrawAccDropInitBnPr}{60\%}  % experiments/FIGS/data/v2/index.json :: linear probe JetClass-II, largest set: relative drop in accuracy (%) with the top data's missing slots set to draws, INIT_BN :: sha256 c6c4138ee5c86423 :: prose twin, two significant figures
\newcommand{\MlstLpJcTwoQgZeroAccDropFull}{61.2\,\ensuremath{\pm}\,3.0\%}  % experiments/FIGS/data/v2/index.json :: linear probe JetClass-II, largest set: relative drop in accuracy (%) with the qg data's missing slots set to zero, L188 :: sha256 c6c4138ee5c86423
\newcommand{\MlstLpJcTwoQgZeroAccDropFullPr}{61\%}  % experiments/FIGS/data/v2/index.json :: linear probe JetClass-II, largest set: relative drop in accuracy (%) with the qg data's missing slots set to zero, L188 :: sha256 c6c4138ee5c86423 :: prose twin, two significant figures
\newcommand{\MlstLpJcTwoQgMeanAccDropFull}{48.8\,\ensuremath{\pm}\,1.1\%}  % experiments/FIGS/data/v2/index.json :: linear probe JetClass-II, largest set: relative drop in accuracy (%) with the qg data's missing slots set to mean, L188 :: sha256 c6c4138ee5c86423
\newcommand{\MlstLpJcTwoQgMeanAccDropFullPr}{49\%}  % experiments/FIGS/data/v2/index.json :: linear probe JetClass-II, largest set: relative drop in accuracy (%) with the qg data's missing slots set to mean, L188 :: sha256 c6c4138ee5c86423 :: prose twin, two significant figures
\newcommand{\MlstLpJcTwoQgDrawAccDropFull}{54.30\,\ensuremath{\pm}\,0.21\%}  % experiments/FIGS/data/v2/index.json :: linear probe JetClass-II, largest set: relative drop in accuracy (%) with the qg data's missing slots set to draws, L188 :: sha256 c6c4138ee5c86423
\newcommand{\MlstLpJcTwoQgDrawAccDropFullPr}{54\%}  % experiments/FIGS/data/v2/index.json :: linear probe JetClass-II, largest set: relative drop in accuracy (%) with the qg data's missing slots set to draws, L188 :: sha256 c6c4138ee5c86423 :: prose twin, two significant figures
\newcommand{\MlstLpJcTwoQgZeroAccDropRes}{64\,\ensuremath{\pm}\,5\%}  % experiments/FIGS/data/v2/index.json :: linear probe JetClass-II, largest set: relative drop in accuracy (%) with the qg data's missing slots set to zero, L162 :: sha256 c6c4138ee5c86423
\newcommand{\MlstLpJcTwoQgZeroAccDropResPr}{64\%}  % experiments/FIGS/data/v2/index.json :: linear probe JetClass-II, largest set: relative drop in accuracy (%) with the qg data's missing slots set to zero, L162 :: sha256 c6c4138ee5c86423 :: prose twin, two significant figures
\newcommand{\MlstLpJcTwoQgMeanAccDropRes}{49.5\,\ensuremath{\pm}\,0.5\%}  % experiments/FIGS/data/v2/index.json :: linear probe JetClass-II, largest set: relative drop in accuracy (%) with the qg data's missing slots set to mean, L162 :: sha256 c6c4138ee5c86423
\newcommand{\MlstLpJcTwoQgMeanAccDropResPr}{49\%}  % experiments/FIGS/data/v2/index.json :: linear probe JetClass-II, largest set: relative drop in accuracy (%) with the qg data's missing slots set to mean, L162 :: sha256 c6c4138ee5c86423 :: prose twin, two significant figures
\newcommand{\MlstLpJcTwoQgDrawAccDropRes}{54.6\,\ensuremath{\pm}\,0.6\%}  % experiments/FIGS/data/v2/index.json :: linear probe JetClass-II, largest set: relative drop in accuracy (%) with the qg data's missing slots set to draws, L162 :: sha256 c6c4138ee5c86423
\newcommand{\MlstLpJcTwoQgDrawAccDropResPr}{55\%}  % experiments/FIGS/data/v2/index.json :: linear probe JetClass-II, largest set: relative drop in accuracy (%) with the qg data's missing slots set to draws, L162 :: sha256 c6c4138ee5c86423 :: prose twin, two significant figures
\newcommand{\MlstLpJcTwoQgZeroAccDropMid}{68\,\ensuremath{\pm}\,5\%}  % experiments/FIGS/data/v2/index.json :: linear probe JetClass-II, largest set: relative drop in accuracy (%) with the qg data's missing slots set to zero, R42_Q1 :: sha256 c6c4138ee5c86423
\newcommand{\MlstLpJcTwoQgZeroAccDropMidPr}{68\%}  % experiments/FIGS/data/v2/index.json :: linear probe JetClass-II, largest set: relative drop in accuracy (%) with the qg data's missing slots set to zero, R42_Q1 :: sha256 c6c4138ee5c86423 :: prose twin, two significant figures
\newcommand{\MlstLpJcTwoQgMeanAccDropMid}{47.6\,\ensuremath{\pm}\,1.6\%}  % experiments/FIGS/data/v2/index.json :: linear probe JetClass-II, largest set: relative drop in accuracy (%) with the qg data's missing slots set to mean, R42_Q1 :: sha256 c6c4138ee5c86423
\newcommand{\MlstLpJcTwoQgMeanAccDropMidPr}{48\%}  % experiments/FIGS/data/v2/index.json :: linear probe JetClass-II, largest set: relative drop in accuracy (%) with the qg data's missing slots set to mean, R42_Q1 :: sha256 c6c4138ee5c86423 :: prose twin, two significant figures
\newcommand{\MlstLpJcTwoQgDrawAccDropMid}{53.5\,\ensuremath{\pm}\,1.4\%}  % experiments/FIGS/data/v2/index.json :: linear probe JetClass-II, largest set: relative drop in accuracy (%) with the qg data's missing slots set to draws, R42_Q1 :: sha256 c6c4138ee5c86423
\newcommand{\MlstLpJcTwoQgDrawAccDropMidPr}{54\%}  % experiments/FIGS/data/v2/index.json :: linear probe JetClass-II, largest set: relative drop in accuracy (%) with the qg data's missing slots set to draws, R42_Q1 :: sha256 c6c4138ee5c86423 :: prose twin, two significant figures
\newcommand{\MlstLpJcTwoQgZeroAccDropCoarse}{70.9\,\ensuremath{\pm}\,2.3\%}  % experiments/FIGS/data/v2/index.json :: linear probe JetClass-II, largest set: relative drop in accuracy (%) with the qg data's missing slots set to zero, R16_Q1 :: sha256 c6c4138ee5c86423
\newcommand{\MlstLpJcTwoQgZeroAccDropCoarsePr}{71\%}  % experiments/FIGS/data/v2/index.json :: linear probe JetClass-II, largest set: relative drop in accuracy (%) with the qg data's missing slots set to zero, R16_Q1 :: sha256 c6c4138ee5c86423 :: prose twin, two significant figures
\newcommand{\MlstLpJcTwoQgMeanAccDropCoarse}{50.0\,\ensuremath{\pm}\,1.0\%}  % experiments/FIGS/data/v2/index.json :: linear probe JetClass-II, largest set: relative drop in accuracy (%) with the qg data's missing slots set to mean, R16_Q1 :: sha256 c6c4138ee5c86423
\newcommand{\MlstLpJcTwoQgMeanAccDropCoarsePr}{50\%}  % experiments/FIGS/data/v2/index.json :: linear probe JetClass-II, largest set: relative drop in accuracy (%) with the qg data's missing slots set to mean, R16_Q1 :: sha256 c6c4138ee5c86423 :: prose twin, two significant figures
\newcommand{\MlstLpJcTwoQgDrawAccDropCoarse}{51.7\,\ensuremath{\pm}\,0.4\%}  % experiments/FIGS/data/v2/index.json :: linear probe JetClass-II, largest set: relative drop in accuracy (%) with the qg data's missing slots set to draws, R16_Q1 :: sha256 c6c4138ee5c86423
\newcommand{\MlstLpJcTwoQgDrawAccDropCoarsePr}{52\%}  % experiments/FIGS/data/v2/index.json :: linear probe JetClass-II, largest set: relative drop in accuracy (%) with the qg data's missing slots set to draws, R16_Q1 :: sha256 c6c4138ee5c86423 :: prose twin, two significant figures
\newcommand{\MlstLpJcTwoQgZeroAccDropResMass}{61.8, 67.9\%}  % experiments/FIGS/data/v2/index.json :: linear probe JetClass-II, largest set: relative drop in accuracy (%) with the qg data's missing slots set to zero, L162_MASS :: sha256 c6c4138ee5c86423
\newcommand{\MlstLpJcTwoQgZeroAccDropResMassPr}{62--68\%}  % experiments/FIGS/data/v2/index.json :: linear probe JetClass-II, largest set: relative drop in accuracy (%) with the qg data's missing slots set to zero, L162_MASS :: sha256 c6c4138ee5c86423 :: prose twin, two significant figures
\newcommand{\MlstLpJcTwoQgMeanAccDropResMass}{49.7, 52.2\%}  % experiments/FIGS/data/v2/index.json :: linear probe JetClass-II, largest set: relative drop in accuracy (%) with the qg data's missing slots set to mean, L162_MASS :: sha256 c6c4138ee5c86423
\newcommand{\MlstLpJcTwoQgMeanAccDropResMassPr}{50--52\%}  % experiments/FIGS/data/v2/index.json :: linear probe JetClass-II, largest set: relative drop in accuracy (%) with the qg data's missing slots set to mean, L162_MASS :: sha256 c6c4138ee5c86423 :: prose twin, two significant figures
\newcommand{\MlstLpJcTwoQgDrawAccDropResMass}{54.3, 55.1\%}  % experiments/FIGS/data/v2/index.json :: linear probe JetClass-II, largest set: relative drop in accuracy (%) with the qg data's missing slots set to draws, L162_MASS :: sha256 c6c4138ee5c86423
\newcommand{\MlstLpJcTwoQgDrawAccDropResMassPr}{54--55\%}  % experiments/FIGS/data/v2/index.json :: linear probe JetClass-II, largest set: relative drop in accuracy (%) with the qg data's missing slots set to draws, L162_MASS :: sha256 c6c4138ee5c86423 :: prose twin, two significant figures
\newcommand{\MlstLpJcTwoQgZeroAccDropCoarseMass}{66.7, 65.4\%}  % experiments/FIGS/data/v2/index.json :: linear probe JetClass-II, largest set: relative drop in accuracy (%) with the qg data's missing slots set to zero, R16_Q1_MASS :: sha256 c6c4138ee5c86423
\newcommand{\MlstLpJcTwoQgZeroAccDropCoarseMassPr}{65--67\%}  % experiments/FIGS/data/v2/index.json :: linear probe JetClass-II, largest set: relative drop in accuracy (%) with the qg data's missing slots set to zero, R16_Q1_MASS :: sha256 c6c4138ee5c86423 :: prose twin, two significant figures
\newcommand{\MlstLpJcTwoQgMeanAccDropCoarseMass}{49.0, 49.6\%}  % experiments/FIGS/data/v2/index.json :: linear probe JetClass-II, largest set: relative drop in accuracy (%) with the qg data's missing slots set to mean, R16_Q1_MASS :: sha256 c6c4138ee5c86423
\newcommand{\MlstLpJcTwoQgMeanAccDropCoarseMassPr}{49--50\%}  % experiments/FIGS/data/v2/index.json :: linear probe JetClass-II, largest set: relative drop in accuracy (%) with the qg data's missing slots set to mean, R16_Q1_MASS :: sha256 c6c4138ee5c86423 :: prose twin, two significant figures
\newcommand{\MlstLpJcTwoQgDrawAccDropCoarseMass}{50.8, 51.1\%}  % experiments/FIGS/data/v2/index.json :: linear probe JetClass-II, largest set: relative drop in accuracy (%) with the qg data's missing slots set to draws, R16_Q1_MASS :: sha256 c6c4138ee5c86423
\newcommand{\MlstLpJcTwoQgDrawAccDropCoarseMassPr}{51--51\%}  % experiments/FIGS/data/v2/index.json :: linear probe JetClass-II, largest set: relative drop in accuracy (%) with the qg data's missing slots set to draws, R16_Q1_MASS :: sha256 c6c4138ee5c86423 :: prose twin, two significant figures
\newcommand{\MlstLpJcTwoQgZeroAccDropFzero}{62.7, 62.2\%}  % experiments/FIGS/data/v2/index.json :: linear probe JetClass-II, largest set: relative drop in accuracy (%) with the qg data's missing slots set to zero, FLAV_F0 :: sha256 c6c4138ee5c86423
\newcommand{\MlstLpJcTwoQgZeroAccDropFzeroPr}{62--63\%}  % experiments/FIGS/data/v2/index.json :: linear probe JetClass-II, largest set: relative drop in accuracy (%) with the qg data's missing slots set to zero, FLAV_F0 :: sha256 c6c4138ee5c86423 :: prose twin, two significant figures
\newcommand{\MlstLpJcTwoQgMeanAccDropFzero}{51.9, 51.3\%}  % experiments/FIGS/data/v2/index.json :: linear probe JetClass-II, largest set: relative drop in accuracy (%) with the qg data's missing slots set to mean, FLAV_F0 :: sha256 c6c4138ee5c86423
\newcommand{\MlstLpJcTwoQgMeanAccDropFzeroPr}{51--52\%}  % experiments/FIGS/data/v2/index.json :: linear probe JetClass-II, largest set: relative drop in accuracy (%) with the qg data's missing slots set to mean, FLAV_F0 :: sha256 c6c4138ee5c86423 :: prose twin, two significant figures
\newcommand{\MlstLpJcTwoQgDrawAccDropFzero}{53.9, 53.4\%}  % experiments/FIGS/data/v2/index.json :: linear probe JetClass-II, largest set: relative drop in accuracy (%) with the qg data's missing slots set to draws, FLAV_F0 :: sha256 c6c4138ee5c86423
\newcommand{\MlstLpJcTwoQgDrawAccDropFzeroPr}{53--54\%}  % experiments/FIGS/data/v2/index.json :: linear probe JetClass-II, largest set: relative drop in accuracy (%) with the qg data's missing slots set to draws, FLAV_F0 :: sha256 c6c4138ee5c86423 :: prose twin, two significant figures
\newcommand{\MlstLpJcTwoQgZeroAccDropFone}{66.2, 62.8\%}  % experiments/FIGS/data/v2/index.json :: linear probe JetClass-II, largest set: relative drop in accuracy (%) with the qg data's missing slots set to zero, FLAV_F1 :: sha256 c6c4138ee5c86423
\newcommand{\MlstLpJcTwoQgZeroAccDropFonePr}{63--66\%}  % experiments/FIGS/data/v2/index.json :: linear probe JetClass-II, largest set: relative drop in accuracy (%) with the qg data's missing slots set to zero, FLAV_F1 :: sha256 c6c4138ee5c86423 :: prose twin, two significant figures
\newcommand{\MlstLpJcTwoQgMeanAccDropFone}{52.7, 52.8\%}  % experiments/FIGS/data/v2/index.json :: linear probe JetClass-II, largest set: relative drop in accuracy (%) with the qg data's missing slots set to mean, FLAV_F1 :: sha256 c6c4138ee5c86423
\newcommand{\MlstLpJcTwoQgMeanAccDropFonePr}{53--53\%}  % experiments/FIGS/data/v2/index.json :: linear probe JetClass-II, largest set: relative drop in accuracy (%) with the qg data's missing slots set to mean, FLAV_F1 :: sha256 c6c4138ee5c86423 :: prose twin, two significant figures
\newcommand{\MlstLpJcTwoQgDrawAccDropFone}{55.6, 56.5\%}  % experiments/FIGS/data/v2/index.json :: linear probe JetClass-II, largest set: relative drop in accuracy (%) with the qg data's missing slots set to draws, FLAV_F1 :: sha256 c6c4138ee5c86423
\newcommand{\MlstLpJcTwoQgDrawAccDropFonePr}{56--56\%}  % experiments/FIGS/data/v2/index.json :: linear probe JetClass-II, largest set: relative drop in accuracy (%) with the qg data's missing slots set to draws, FLAV_F1 :: sha256 c6c4138ee5c86423 :: prose twin, two significant figures
\newcommand{\MlstLpJcTwoQgZeroAccDropFoner}{67.2, 69.0\%}  % experiments/FIGS/data/v2/index.json :: linear probe JetClass-II, largest set: relative drop in accuracy (%) with the qg data's missing slots set to zero, FLAV_F1R :: sha256 c6c4138ee5c86423
\newcommand{\MlstLpJcTwoQgZeroAccDropFonerPr}{67--69\%}  % experiments/FIGS/data/v2/index.json :: linear probe JetClass-II, largest set: relative drop in accuracy (%) with the qg data's missing slots set to zero, FLAV_F1R :: sha256 c6c4138ee5c86423 :: prose twin, two significant figures
\newcommand{\MlstLpJcTwoQgMeanAccDropFoner}{52.9, 51.8\%}  % experiments/FIGS/data/v2/index.json :: linear probe JetClass-II, largest set: relative drop in accuracy (%) with the qg data's missing slots set to mean, FLAV_F1R :: sha256 c6c4138ee5c86423
\newcommand{\MlstLpJcTwoQgMeanAccDropFonerPr}{52--53\%}  % experiments/FIGS/data/v2/index.json :: linear probe JetClass-II, largest set: relative drop in accuracy (%) with the qg data's missing slots set to mean, FLAV_F1R :: sha256 c6c4138ee5c86423 :: prose twin, two significant figures
\newcommand{\MlstLpJcTwoQgDrawAccDropFoner}{57.2, 55.4\%}  % experiments/FIGS/data/v2/index.json :: linear probe JetClass-II, largest set: relative drop in accuracy (%) with the qg data's missing slots set to draws, FLAV_F1R :: sha256 c6c4138ee5c86423
\newcommand{\MlstLpJcTwoQgDrawAccDropFonerPr}{55--57\%}  % experiments/FIGS/data/v2/index.json :: linear probe JetClass-II, largest set: relative drop in accuracy (%) with the qg data's missing slots set to draws, FLAV_F1R :: sha256 c6c4138ee5c86423 :: prose twin, two significant figures
\newcommand{\MlstLpJcTwoQgZeroAccDropInitBn}{8.7\,\ensuremath{\pm}\,2.6\%}  % experiments/FIGS/data/v2/index.json :: linear probe JetClass-II, largest set: relative drop in accuracy (%) with the qg data's missing slots set to zero, INIT_BN :: sha256 c6c4138ee5c86423
\newcommand{\MlstLpJcTwoQgZeroAccDropInitBnPr}{8.7\%}  % experiments/FIGS/data/v2/index.json :: linear probe JetClass-II, largest set: relative drop in accuracy (%) with the qg data's missing slots set to zero, INIT_BN :: sha256 c6c4138ee5c86423 :: prose twin, two significant figures
\newcommand{\MlstLpJcTwoQgMeanAccDropInitBn}{10.4\,\ensuremath{\pm}\,1.4\%}  % experiments/FIGS/data/v2/index.json :: linear probe JetClass-II, largest set: relative drop in accuracy (%) with the qg data's missing slots set to mean, INIT_BN :: sha256 c6c4138ee5c86423
\newcommand{\MlstLpJcTwoQgMeanAccDropInitBnPr}{10\%}  % experiments/FIGS/data/v2/index.json :: linear probe JetClass-II, largest set: relative drop in accuracy (%) with the qg data's missing slots set to mean, INIT_BN :: sha256 c6c4138ee5c86423 :: prose twin, two significant figures
\newcommand{\MlstLpJcTwoQgDrawAccDropInitBn}{11.4\,\ensuremath{\pm}\,2.1\%}  % experiments/FIGS/data/v2/index.json :: linear probe JetClass-II, largest set: relative drop in accuracy (%) with the qg data's missing slots set to draws, INIT_BN :: sha256 c6c4138ee5c86423
\newcommand{\MlstLpJcTwoQgDrawAccDropInitBnPr}{11\%}  % experiments/FIGS/data/v2/index.json :: linear probe JetClass-II, largest set: relative drop in accuracy (%) with the qg data's missing slots set to draws, INIT_BN :: sha256 c6c4138ee5c86423 :: prose twin, two significant figures
\newcommand{\MlstMassLambda}{5}  % configs/arms/v2_grid.json :: arms[L162_MASS, R16_Q1_MASS].mass_lambda :: sha256 d8a31b1486992302 :: weight of the mass-regression term in the pretraining loss, both mass models
\newcommand{\MlstMassLossShareResMass}{13.78 and 13.87}  % experiments/FIGS/data/v2/index.json :: share of the mass-regression term in the training loss of L162_MASS, mean over epochs, percent :: sha256 c6c4138ee5c86423
\newcommand{\MlstMassLossShareCoarseMass}{28.65 and 28.67}  % experiments/FIGS/data/v2/index.json :: share of the mass-regression term in the training loss of R16_Q1_MASS, mean over epochs, percent :: sha256 c6c4138ee5c86423
\newcommand{\MlstMassUncorrected}{0.245}  % experiments/FIGS/data/v2/index.json :: head/atlas_res: groomed mass, all jets, ATLAS resolution :: sha256 c6c4138ee5c86423
\newcommand{\MlstMassUncorrectedResponse}{0.929}  % experiments/FIGS/data/v2/index.json :: head/atlas_res: groomed mass, all jets, median response :: sha256 c6c4138ee5c86423
\newcommand{\MlstMassHeadResMass}{0.1471 and 0.1463}  % experiments/FIGS/data/v2/index.json :: regression output of L162_MASS, all jets, ATLAS resolution :: sha256 c6c4138ee5c86423
\newcommand{\MlstMassHeadResMassResponse}{0.995 and 1.003}  % experiments/FIGS/data/v2/index.json :: regression output of L162_MASS, all jets, median response :: sha256 c6c4138ee5c86423
\newcommand{\MlstMassHeadResMassOverUncorr}{0.6009 and 0.5976}  % experiments/FIGS/data/v2/index.json :: resolution of the regression output of L162_MASS over the groomed mass's, all jets :: sha256 c6c4138ee5c86423
\newcommand{\MlstMassHeadResMassOverUncorrPr}{at most 0.60 times the resolution}  % experiments/FIGS/data/v2/index.json :: resolution of the regression output of L162_MASS over the groomed mass's, all jets :: sha256 c6c4138ee5c86423 :: prose twin: two significant figures, direction in words
\newcommand{\MlstMassHeadResMassOverUncorrPrMax}{down to 0.60 times the resolution}  % experiments/FIGS/data/v2/index.json :: resolution of the regression output of L162_MASS over the groomed mass's, all jets :: sha256 c6c4138ee5c86423 :: prose twin of the run farther from 1
\newcommand{\MlstMassHeadResMassOverMlpRes}{0.6645 and 0.6691}  % experiments/FIGS/data/v2/index.json :: resolution of the regression output of L162_MASS over the frozen-feature MLP of L162, run by run, all jets :: sha256 c6c4138ee5c86423
\newcommand{\MlstMassHeadResMassOverMlpResPr}{at most 0.67 times the resolution}  % experiments/FIGS/data/v2/index.json :: resolution of the regression output of L162_MASS over the frozen-feature MLP of L162, run by run, all jets :: sha256 c6c4138ee5c86423 :: prose twin: two significant figures, direction in words
\newcommand{\MlstMassHeadResMassOverMlpResPrMax}{down to 0.66 times the resolution}  % experiments/FIGS/data/v2/index.json :: resolution of the regression output of L162_MASS over the frozen-feature MLP of L162, run by run, all jets :: sha256 c6c4138ee5c86423 :: prose twin of the run farther from 1
\newcommand{\MlstMassHeadCoarseMass}{0.1450 and 0.1461}  % experiments/FIGS/data/v2/index.json :: regression output of R16_Q1_MASS, all jets, ATLAS resolution :: sha256 c6c4138ee5c86423
\newcommand{\MlstMassHeadCoarseMassResponse}{1.005 and 1.005}  % experiments/FIGS/data/v2/index.json :: regression output of R16_Q1_MASS, all jets, median response :: sha256 c6c4138ee5c86423
\newcommand{\MlstMassHeadCoarseMassOverUncorr}{0.5921 and 0.5966}  % experiments/FIGS/data/v2/index.json :: resolution of the regression output of R16_Q1_MASS over the groomed mass's, all jets :: sha256 c6c4138ee5c86423
\newcommand{\MlstMassHeadCoarseMassOverUncorrPr}{at most 0.60 times the resolution}  % experiments/FIGS/data/v2/index.json :: resolution of the regression output of R16_Q1_MASS over the groomed mass's, all jets :: sha256 c6c4138ee5c86423 :: prose twin: two significant figures, direction in words
\newcommand{\MlstMassHeadCoarseMassOverUncorrPrMax}{down to 0.59 times the resolution}  % experiments/FIGS/data/v2/index.json :: resolution of the regression output of R16_Q1_MASS over the groomed mass's, all jets :: sha256 c6c4138ee5c86423 :: prose twin of the run farther from 1
\newcommand{\MlstMassHeadCoarseMassOverMlpCoarse}{0.703 and 0.690}  % experiments/FIGS/data/v2/index.json :: resolution of the regression output of R16_Q1_MASS over the frozen-feature MLP of R16_Q1, run by run, all jets :: sha256 c6c4138ee5c86423
\newcommand{\MlstMassHeadCoarseMassOverMlpCoarsePr}{at most 0.70 times the resolution}  % experiments/FIGS/data/v2/index.json :: resolution of the regression output of R16_Q1_MASS over the frozen-feature MLP of R16_Q1, run by run, all jets :: sha256 c6c4138ee5c86423 :: prose twin: two significant figures, direction in words
\newcommand{\MlstMassHeadCoarseMassOverMlpCoarsePrMax}{down to 0.69 times the resolution}  % experiments/FIGS/data/v2/index.json :: resolution of the regression output of R16_Q1_MASS over the frozen-feature MLP of R16_Q1, run by run, all jets :: sha256 c6c4138ee5c86423 :: prose twin of the run farther from 1
\newcommand{\MlstMassMlpInit}{0.2421\,\ensuremath{\pm}\,0.0018}  % experiments/FIGS/data/v2/index.json :: frozen-feature MLP of INIT, all jets, ATLAS resolution :: sha256 c6c4138ee5c86423
\newcommand{\MlstMassMlpFull}{0.2200\,\ensuremath{\pm}\,0.0014}  % experiments/FIGS/data/v2/index.json :: frozen-feature MLP of L188, all jets, ATLAS resolution :: sha256 c6c4138ee5c86423
\newcommand{\MlstMassMlpRes}{0.2195\,\ensuremath{\pm}\,0.0017}  % experiments/FIGS/data/v2/index.json :: frozen-feature MLP of L162, all jets, ATLAS resolution :: sha256 c6c4138ee5c86423
\newcommand{\MlstMassMlpMid}{0.2158\,\ensuremath{\pm}\,0.0029}  % experiments/FIGS/data/v2/index.json :: frozen-feature MLP of R42_Q1, all jets, ATLAS resolution :: sha256 c6c4138ee5c86423
\newcommand{\MlstMassMlpCoarse}{0.2098\,\ensuremath{\pm}\,0.0031}  % experiments/FIGS/data/v2/index.json :: frozen-feature MLP of R16_Q1, all jets, ATLAS resolution :: sha256 c6c4138ee5c86423
\newcommand{\MlstMassMlpResMass}{0.1591 and 0.1586}  % experiments/FIGS/data/v2/index.json :: frozen-feature MLP of L162_MASS, all jets, ATLAS resolution :: sha256 c6c4138ee5c86423
\newcommand{\MlstMassMlpCoarseMass}{0.1583 and 0.1588}  % experiments/FIGS/data/v2/index.json :: frozen-feature MLP of R16_Q1_MASS, all jets, ATLAS resolution :: sha256 c6c4138ee5c86423
\newcommand{\MlstCmpMlpResOverFull}{0.998\,\ensuremath{\pm}\,0.006}  % experiments/FIGS/data/v2/index.json :: frozen-feature MLP resolution of L162 over L188, run by run, all jets :: sha256 c6c4138ee5c86423
\newcommand{\MlstCmpMlpResOverFullPr}{equal within 1\%}  % experiments/FIGS/data/v2/index.json :: frozen-feature MLP resolution of L162 over L188, run by run, all jets :: sha256 c6c4138ee5c86423 :: prose twin: two significant figures, direction in words
\newcommand{\MlstCmpMlpMidOverFull}{0.981\,\ensuremath{\pm}\,0.017}  % experiments/FIGS/data/v2/index.json :: frozen-feature MLP resolution of R42_Q1 over L188, run by run, all jets :: sha256 c6c4138ee5c86423
\newcommand{\MlstCmpMlpMidOverFullPr}{equal within 2\%}  % experiments/FIGS/data/v2/index.json :: frozen-feature MLP resolution of R42_Q1 over L188, run by run, all jets :: sha256 c6c4138ee5c86423 :: prose twin: two significant figures, direction in words
\newcommand{\MlstCmpMlpCoarseOverFull}{0.954\,\ensuremath{\pm}\,0.018}  % experiments/FIGS/data/v2/index.json :: frozen-feature MLP resolution of R16_Q1 over L188, run by run, all jets :: sha256 c6c4138ee5c86423
\newcommand{\MlstCmpMlpCoarseOverFullPr}{0.95 times the resolution}  % experiments/FIGS/data/v2/index.json :: frozen-feature MLP resolution of R16_Q1 over L188, run by run, all jets :: sha256 c6c4138ee5c86423 :: prose twin: two significant figures, direction in words
\newcommand{\MlstCmpMlpResMassOverRes}{0.7186 and 0.7253}  % experiments/FIGS/data/v2/index.json :: frozen-feature MLP resolution of L162_MASS over L162, run by run, all jets :: sha256 c6c4138ee5c86423
\newcommand{\MlstCmpMlpResMassOverResPr}{at most 0.73 times the resolution}  % experiments/FIGS/data/v2/index.json :: frozen-feature MLP resolution of L162_MASS over L162, run by run, all jets :: sha256 c6c4138ee5c86423 :: prose twin: two significant figures, direction in words
\newcommand{\MlstCmpMlpResMassOverResPrMax}{down to 0.72 times the resolution}  % experiments/FIGS/data/v2/index.json :: frozen-feature MLP resolution of L162_MASS over L162, run by run, all jets :: sha256 c6c4138ee5c86423 :: prose twin of the run farther from 1
\newcommand{\MlstCmpMlpCoarseMassOverCoarse}{0.768 and 0.750}  % experiments/FIGS/data/v2/index.json :: frozen-feature MLP resolution of R16_Q1_MASS over R16_Q1, run by run, all jets :: sha256 c6c4138ee5c86423
\newcommand{\MlstCmpMlpCoarseMassOverCoarsePr}{at most 0.77 times the resolution}  % experiments/FIGS/data/v2/index.json :: frozen-feature MLP resolution of R16_Q1_MASS over R16_Q1, run by run, all jets :: sha256 c6c4138ee5c86423 :: prose twin: two significant figures, direction in words
\newcommand{\MlstCmpMlpCoarseMassOverCoarsePrMax}{down to 0.75 times the resolution}  % experiments/FIGS/data/v2/index.json :: frozen-feature MLP resolution of R16_Q1_MASS over R16_Q1, run by run, all jets :: sha256 c6c4138ee5c86423 :: prose twin of the run farther from 1
\newcommand{\MlstMassUncorrectedResonant}{0.213}  % experiments/FIGS/data/v2/index.json :: head/atlas_res: groomed mass, resonant jets, ATLAS resolution :: sha256 c6c4138ee5c86423
\newcommand{\MlstMassUncorrectedResponseResonant}{0.929}  % experiments/FIGS/data/v2/index.json :: head/atlas_res: groomed mass, resonant jets, median response :: sha256 c6c4138ee5c86423
\newcommand{\MlstMassHeadResMassResonant}{0.1332 and 0.1320}  % experiments/FIGS/data/v2/index.json :: regression output of L162_MASS, resonant jets, ATLAS resolution :: sha256 c6c4138ee5c86423
\newcommand{\MlstMassHeadResMassResponseResonant}{0.996 and 1.004}  % experiments/FIGS/data/v2/index.json :: regression output of L162_MASS, resonant jets, median response :: sha256 c6c4138ee5c86423
\newcommand{\MlstMassHeadResMassOverUncorrResonant}{0.6259 and 0.6205}  % experiments/FIGS/data/v2/index.json :: resolution of the regression output of L162_MASS over the groomed mass's, resonant jets :: sha256 c6c4138ee5c86423
\newcommand{\MlstMassHeadResMassOverUncorrResonantPr}{at most 0.63 times the resolution}  % experiments/FIGS/data/v2/index.json :: resolution of the regression output of L162_MASS over the groomed mass's, resonant jets :: sha256 c6c4138ee5c86423 :: prose twin: two significant figures, direction in words
\newcommand{\MlstMassHeadResMassOverUncorrResonantPrMax}{down to 0.62 times the resolution}  % experiments/FIGS/data/v2/index.json :: resolution of the regression output of L162_MASS over the groomed mass's, resonant jets :: sha256 c6c4138ee5c86423 :: prose twin of the run farther from 1
\newcommand{\MlstMassHeadResMassOverMlpResResonant}{0.6620 and 0.6662}  % experiments/FIGS/data/v2/index.json :: resolution of the regression output of L162_MASS over the frozen-feature MLP of L162, run by run, resonant jets :: sha256 c6c4138ee5c86423
\newcommand{\MlstMassHeadResMassOverMlpResResonantPr}{at most 0.67 times the resolution}  % experiments/FIGS/data/v2/index.json :: resolution of the regression output of L162_MASS over the frozen-feature MLP of L162, run by run, resonant jets :: sha256 c6c4138ee5c86423 :: prose twin: two significant figures, direction in words
\newcommand{\MlstMassHeadResMassOverMlpResResonantPrMax}{down to 0.66 times the resolution}  % experiments/FIGS/data/v2/index.json :: resolution of the regression output of L162_MASS over the frozen-feature MLP of L162, run by run, resonant jets :: sha256 c6c4138ee5c86423 :: prose twin of the run farther from 1
\newcommand{\MlstMassHeadCoarseMassResonant}{0.1312 and 0.1320}  % experiments/FIGS/data/v2/index.json :: regression output of R16_Q1_MASS, resonant jets, ATLAS resolution :: sha256 c6c4138ee5c86423
\newcommand{\MlstMassHeadCoarseMassResponseResonant}{1.006 and 1.006}  % experiments/FIGS/data/v2/index.json :: regression output of R16_Q1_MASS, resonant jets, median response :: sha256 c6c4138ee5c86423
\newcommand{\MlstMassHeadCoarseMassOverUncorrResonant}{0.6165 and 0.6206}  % experiments/FIGS/data/v2/index.json :: resolution of the regression output of R16_Q1_MASS over the groomed mass's, resonant jets :: sha256 c6c4138ee5c86423
\newcommand{\MlstMassHeadCoarseMassOverUncorrResonantPr}{at most 0.62 times the resolution}  % experiments/FIGS/data/v2/index.json :: resolution of the regression output of R16_Q1_MASS over the groomed mass's, resonant jets :: sha256 c6c4138ee5c86423 :: prose twin: two significant figures, direction in words
\newcommand{\MlstMassHeadCoarseMassOverUncorrResonantPrMax}{down to 0.62 times the resolution}  % experiments/FIGS/data/v2/index.json :: resolution of the regression output of R16_Q1_MASS over the groomed mass's, resonant jets :: sha256 c6c4138ee5c86423 :: prose twin of the run farther from 1
\newcommand{\MlstMassHeadCoarseMassOverMlpCoarseResonant}{0.707 and 0.688}  % experiments/FIGS/data/v2/index.json :: resolution of the regression output of R16_Q1_MASS over the frozen-feature MLP of R16_Q1, run by run, resonant jets :: sha256 c6c4138ee5c86423
\newcommand{\MlstMassHeadCoarseMassOverMlpCoarseResonantPr}{at most 0.71 times the resolution}  % experiments/FIGS/data/v2/index.json :: resolution of the regression output of R16_Q1_MASS over the frozen-feature MLP of R16_Q1, run by run, resonant jets :: sha256 c6c4138ee5c86423 :: prose twin: two significant figures, direction in words
\newcommand{\MlstMassHeadCoarseMassOverMlpCoarseResonantPrMax}{down to 0.69 times the resolution}  % experiments/FIGS/data/v2/index.json :: resolution of the regression output of R16_Q1_MASS over the frozen-feature MLP of R16_Q1, run by run, resonant jets :: sha256 c6c4138ee5c86423 :: prose twin of the run farther from 1
\newcommand{\MlstMassMlpInitResonant}{0.2214\,\ensuremath{\pm}\,0.0017}  % experiments/FIGS/data/v2/index.json :: frozen-feature MLP of INIT, resonant jets, ATLAS resolution :: sha256 c6c4138ee5c86423
\newcommand{\MlstMassMlpFullResonant}{0.1990\,\ensuremath{\pm}\,0.0016}  % experiments/FIGS/data/v2/index.json :: frozen-feature MLP of L188, resonant jets, ATLAS resolution :: sha256 c6c4138ee5c86423
\newcommand{\MlstMassMlpResResonant}{0.1989\,\ensuremath{\pm}\,0.0020}  % experiments/FIGS/data/v2/index.json :: frozen-feature MLP of L162, resonant jets, ATLAS resolution :: sha256 c6c4138ee5c86423
\newcommand{\MlstMassMlpMidResonant}{0.1949\,\ensuremath{\pm}\,0.0023}  % experiments/FIGS/data/v2/index.json :: frozen-feature MLP of R42_Q1, resonant jets, ATLAS resolution :: sha256 c6c4138ee5c86423
\newcommand{\MlstMassMlpCoarseResonant}{0.1894\,\ensuremath{\pm}\,0.0035}  % experiments/FIGS/data/v2/index.json :: frozen-feature MLP of R16_Q1, resonant jets, ATLAS resolution :: sha256 c6c4138ee5c86423
\newcommand{\MlstMassMlpResMassResonant}{0.1436 and 0.1430}  % experiments/FIGS/data/v2/index.json :: frozen-feature MLP of L162_MASS, resonant jets, ATLAS resolution :: sha256 c6c4138ee5c86423
\newcommand{\MlstMassMlpCoarseMassResonant}{0.1420 and 0.1431}  % experiments/FIGS/data/v2/index.json :: frozen-feature MLP of R16_Q1_MASS, resonant jets, ATLAS resolution :: sha256 c6c4138ee5c86423
\newcommand{\MlstCmpMlpResOverFullResonant}{1.000\,\ensuremath{\pm}\,0.005}  % experiments/FIGS/data/v2/index.json :: frozen-feature MLP resolution of L162 over L188, run by run, resonant jets :: sha256 c6c4138ee5c86423
\newcommand{\MlstCmpMlpResOverFullResonantPr}{equal within 1\%}  % experiments/FIGS/data/v2/index.json :: frozen-feature MLP resolution of L162 over L188, run by run, resonant jets :: sha256 c6c4138ee5c86423 :: prose twin: two significant figures, direction in words
\newcommand{\MlstCmpMlpMidOverFullResonant}{0.980\,\ensuremath{\pm}\,0.018}  % experiments/FIGS/data/v2/index.json :: frozen-feature MLP resolution of R42_Q1 over L188, run by run, resonant jets :: sha256 c6c4138ee5c86423
\newcommand{\MlstCmpMlpMidOverFullResonantPr}{equal within 3\%}  % experiments/FIGS/data/v2/index.json :: frozen-feature MLP resolution of R42_Q1 over L188, run by run, resonant jets :: sha256 c6c4138ee5c86423 :: prose twin: two significant figures, direction in words
\newcommand{\MlstCmpMlpCoarseOverFullResonant}{0.952\,\ensuremath{\pm}\,0.023}  % experiments/FIGS/data/v2/index.json :: frozen-feature MLP resolution of R16_Q1 over L188, run by run, resonant jets :: sha256 c6c4138ee5c86423
\newcommand{\MlstCmpMlpCoarseOverFullResonantPr}{0.95 times the resolution}  % experiments/FIGS/data/v2/index.json :: frozen-feature MLP resolution of R16_Q1 over L188, run by run, resonant jets :: sha256 c6c4138ee5c86423 :: prose twin: two significant figures, direction in words
\newcommand{\MlstCmpMlpResMassOverResResonant}{0.7137 and 0.7219}  % experiments/FIGS/data/v2/index.json :: frozen-feature MLP resolution of L162_MASS over L162, run by run, resonant jets :: sha256 c6c4138ee5c86423
\newcommand{\MlstCmpMlpResMassOverResResonantPr}{at most 0.72 times the resolution}  % experiments/FIGS/data/v2/index.json :: frozen-feature MLP resolution of L162_MASS over L162, run by run, resonant jets :: sha256 c6c4138ee5c86423 :: prose twin: two significant figures, direction in words
\newcommand{\MlstCmpMlpResMassOverResResonantPrMax}{down to 0.71 times the resolution}  % experiments/FIGS/data/v2/index.json :: frozen-feature MLP resolution of L162_MASS over L162, run by run, resonant jets :: sha256 c6c4138ee5c86423 :: prose twin of the run farther from 1
\newcommand{\MlstCmpMlpCoarseMassOverCoarseResonant}{0.766 and 0.745}  % experiments/FIGS/data/v2/index.json :: frozen-feature MLP resolution of R16_Q1_MASS over R16_Q1, run by run, resonant jets :: sha256 c6c4138ee5c86423
\newcommand{\MlstCmpMlpCoarseMassOverCoarseResonantPr}{at most 0.77 times the resolution}  % experiments/FIGS/data/v2/index.json :: frozen-feature MLP resolution of R16_Q1_MASS over R16_Q1, run by run, resonant jets :: sha256 c6c4138ee5c86423 :: prose twin: two significant figures, direction in words
\newcommand{\MlstCmpMlpCoarseMassOverCoarseResonantPrMax}{down to 0.75 times the resolution}  % experiments/FIGS/data/v2/index.json :: frozen-feature MLP resolution of R16_Q1_MASS over R16_Q1, run by run, resonant jets :: sha256 c6c4138ee5c86423 :: prose twin of the run farther from 1
\newcommand{\MlstMassUncorrectedQcd}{1.52}  % experiments/FIGS/data/v2/index.json :: head/atlas_res: groomed mass, qcd jets, ATLAS resolution :: sha256 c6c4138ee5c86423
\newcommand{\MlstMassUncorrectedResponseQcd}{0.931}  % experiments/FIGS/data/v2/index.json :: head/atlas_res: groomed mass, qcd jets, median response :: sha256 c6c4138ee5c86423
\newcommand{\MlstMassHeadResMassQcd}{0.2263 and 0.2256}  % experiments/FIGS/data/v2/index.json :: regression output of L162_MASS, qcd jets, ATLAS resolution :: sha256 c6c4138ee5c86423
\newcommand{\MlstMassHeadResMassResponseQcd}{0.9897 and 0.9996}  % experiments/FIGS/data/v2/index.json :: regression output of L162_MASS, qcd jets, median response :: sha256 c6c4138ee5c86423
\newcommand{\MlstMassHeadResMassOverUncorrQcd}{0.1485 and 0.1480}  % experiments/FIGS/data/v2/index.json :: resolution of the regression output of L162_MASS over the groomed mass's, qcd jets :: sha256 c6c4138ee5c86423
\newcommand{\MlstMassHeadResMassOverUncorrQcdPr}{at most 0.15 times the resolution}  % experiments/FIGS/data/v2/index.json :: resolution of the regression output of L162_MASS over the groomed mass's, qcd jets :: sha256 c6c4138ee5c86423 :: prose twin: two significant figures, direction in words
\newcommand{\MlstMassHeadResMassOverUncorrQcdPrMax}{down to 0.15 times the resolution}  % experiments/FIGS/data/v2/index.json :: resolution of the regression output of L162_MASS over the groomed mass's, qcd jets :: sha256 c6c4138ee5c86423 :: prose twin of the run farther from 1
\newcommand{\MlstMassHeadResMassOverMlpResQcd}{0.6542 and 0.6632}  % experiments/FIGS/data/v2/index.json :: resolution of the regression output of L162_MASS over the frozen-feature MLP of L162, run by run, qcd jets :: sha256 c6c4138ee5c86423
\newcommand{\MlstMassHeadResMassOverMlpResQcdPr}{at most 0.66 times the resolution}  % experiments/FIGS/data/v2/index.json :: resolution of the regression output of L162_MASS over the frozen-feature MLP of L162, run by run, qcd jets :: sha256 c6c4138ee5c86423 :: prose twin: two significant figures, direction in words
\newcommand{\MlstMassHeadResMassOverMlpResQcdPrMax}{down to 0.65 times the resolution}  % experiments/FIGS/data/v2/index.json :: resolution of the regression output of L162_MASS over the frozen-feature MLP of L162, run by run, qcd jets :: sha256 c6c4138ee5c86423 :: prose twin of the run farther from 1
\newcommand{\MlstMassHeadCoarseMassQcd}{0.2239 and 0.2240}  % experiments/FIGS/data/v2/index.json :: regression output of R16_Q1_MASS, qcd jets, ATLAS resolution :: sha256 c6c4138ee5c86423
\newcommand{\MlstMassHeadCoarseMassResponseQcd}{1.001 and 1.000}  % experiments/FIGS/data/v2/index.json :: regression output of R16_Q1_MASS, qcd jets, median response :: sha256 c6c4138ee5c86423
\newcommand{\MlstMassHeadCoarseMassOverUncorrQcd}{0.1469 and 0.1469}  % experiments/FIGS/data/v2/index.json :: resolution of the regression output of R16_Q1_MASS over the groomed mass's, qcd jets :: sha256 c6c4138ee5c86423
\newcommand{\MlstMassHeadCoarseMassOverUncorrQcdPr}{at most 0.15 times the resolution}  % experiments/FIGS/data/v2/index.json :: resolution of the regression output of R16_Q1_MASS over the groomed mass's, qcd jets :: sha256 c6c4138ee5c86423 :: prose twin: two significant figures, direction in words
\newcommand{\MlstMassHeadCoarseMassOverUncorrQcdPrMax}{down to 0.15 times the resolution}  % experiments/FIGS/data/v2/index.json :: resolution of the regression output of R16_Q1_MASS over the groomed mass's, qcd jets :: sha256 c6c4138ee5c86423 :: prose twin of the run farther from 1
\newcommand{\MlstMassHeadCoarseMassOverMlpCoarseQcd}{0.679 and 0.698}  % experiments/FIGS/data/v2/index.json :: resolution of the regression output of R16_Q1_MASS over the frozen-feature MLP of R16_Q1, run by run, qcd jets :: sha256 c6c4138ee5c86423
\newcommand{\MlstMassHeadCoarseMassOverMlpCoarseQcdPr}{at most 0.70 times the resolution}  % experiments/FIGS/data/v2/index.json :: resolution of the regression output of R16_Q1_MASS over the frozen-feature MLP of R16_Q1, run by run, qcd jets :: sha256 c6c4138ee5c86423 :: prose twin: two significant figures, direction in words
\newcommand{\MlstMassHeadCoarseMassOverMlpCoarseQcdPrMax}{down to 0.68 times the resolution}  % experiments/FIGS/data/v2/index.json :: resolution of the regression output of R16_Q1_MASS over the frozen-feature MLP of R16_Q1, run by run, qcd jets :: sha256 c6c4138ee5c86423 :: prose twin of the run farther from 1
\newcommand{\MlstMassMlpInitQcd}{0.502\,\ensuremath{\pm}\,0.014}  % experiments/FIGS/data/v2/index.json :: frozen-feature MLP of INIT, qcd jets, ATLAS resolution :: sha256 c6c4138ee5c86423
\newcommand{\MlstMassMlpFullQcd}{0.3487\,\ensuremath{\pm}\,0.0029}  % experiments/FIGS/data/v2/index.json :: frozen-feature MLP of L188, qcd jets, ATLAS resolution :: sha256 c6c4138ee5c86423
\newcommand{\MlstMassMlpResQcd}{0.3425\,\ensuremath{\pm}\,0.0031}  % experiments/FIGS/data/v2/index.json :: frozen-feature MLP of L162, qcd jets, ATLAS resolution :: sha256 c6c4138ee5c86423
\newcommand{\MlstMassMlpMidQcd}{0.338\,\ensuremath{\pm}\,0.007}  % experiments/FIGS/data/v2/index.json :: frozen-feature MLP of R42_Q1, qcd jets, ATLAS resolution :: sha256 c6c4138ee5c86423
\newcommand{\MlstMassMlpCoarseQcd}{0.327\,\ensuremath{\pm}\,0.006}  % experiments/FIGS/data/v2/index.json :: frozen-feature MLP of R16_Q1, qcd jets, ATLAS resolution :: sha256 c6c4138ee5c86423
\newcommand{\MlstMassMlpResMassQcd}{0.2473 and 0.2452}  % experiments/FIGS/data/v2/index.json :: frozen-feature MLP of L162_MASS, qcd jets, ATLAS resolution :: sha256 c6c4138ee5c86423
\newcommand{\MlstMassMlpCoarseMassQcd}{0.2465 and 0.2463}  % experiments/FIGS/data/v2/index.json :: frozen-feature MLP of R16_Q1_MASS, qcd jets, ATLAS resolution :: sha256 c6c4138ee5c86423
\newcommand{\MlstCmpMlpResOverFullQcd}{0.982\,\ensuremath{\pm}\,0.017}  % experiments/FIGS/data/v2/index.json :: frozen-feature MLP resolution of L162 over L188, run by run, qcd jets :: sha256 c6c4138ee5c86423
\newcommand{\MlstCmpMlpResOverFullQcdPr}{equal within 2\%}  % experiments/FIGS/data/v2/index.json :: frozen-feature MLP resolution of L162 over L188, run by run, qcd jets :: sha256 c6c4138ee5c86423 :: prose twin: two significant figures, direction in words
\newcommand{\MlstCmpMlpMidOverFullQcd}{0.969\,\ensuremath{\pm}\,0.013}  % experiments/FIGS/data/v2/index.json :: frozen-feature MLP resolution of R42_Q1 over L188, run by run, qcd jets :: sha256 c6c4138ee5c86423
\newcommand{\MlstCmpMlpMidOverFullQcdPr}{0.97 times the resolution}  % experiments/FIGS/data/v2/index.json :: frozen-feature MLP resolution of R42_Q1 over L188, run by run, qcd jets :: sha256 c6c4138ee5c86423 :: prose twin: two significant figures, direction in words
\newcommand{\MlstCmpMlpCoarseOverFullQcd}{0.939\,\ensuremath{\pm}\,0.023}  % experiments/FIGS/data/v2/index.json :: frozen-feature MLP resolution of R16_Q1 over L188, run by run, qcd jets :: sha256 c6c4138ee5c86423
\newcommand{\MlstCmpMlpCoarseOverFullQcdPr}{0.94 times the resolution}  % experiments/FIGS/data/v2/index.json :: frozen-feature MLP resolution of R16_Q1 over L188, run by run, qcd jets :: sha256 c6c4138ee5c86423 :: prose twin: two significant figures, direction in words
\newcommand{\MlstCmpMlpResMassOverResQcd}{0.7147 and 0.7209}  % experiments/FIGS/data/v2/index.json :: frozen-feature MLP resolution of L162_MASS over L162, run by run, qcd jets :: sha256 c6c4138ee5c86423
\newcommand{\MlstCmpMlpResMassOverResQcdPr}{at most 0.72 times the resolution}  % experiments/FIGS/data/v2/index.json :: frozen-feature MLP resolution of L162_MASS over L162, run by run, qcd jets :: sha256 c6c4138ee5c86423 :: prose twin: two significant figures, direction in words
\newcommand{\MlstCmpMlpResMassOverResQcdPrMax}{down to 0.71 times the resolution}  % experiments/FIGS/data/v2/index.json :: frozen-feature MLP resolution of L162_MASS over L162, run by run, qcd jets :: sha256 c6c4138ee5c86423 :: prose twin of the run farther from 1
\newcommand{\MlstCmpMlpCoarseMassOverCoarseQcd}{0.747 and 0.768}  % experiments/FIGS/data/v2/index.json :: frozen-feature MLP resolution of R16_Q1_MASS over R16_Q1, run by run, qcd jets :: sha256 c6c4138ee5c86423
\newcommand{\MlstCmpMlpCoarseMassOverCoarseQcdPr}{at most 0.77 times the resolution}  % experiments/FIGS/data/v2/index.json :: frozen-feature MLP resolution of R16_Q1_MASS over R16_Q1, run by run, qcd jets :: sha256 c6c4138ee5c86423 :: prose twin: two significant figures, direction in words
\newcommand{\MlstCmpMlpCoarseMassOverCoarseQcdPrMax}{down to 0.75 times the resolution}  % experiments/FIGS/data/v2/index.json :: frozen-feature MLP resolution of R16_Q1_MASS over R16_Q1, run by run, qcd jets :: sha256 c6c4138ee5c86423 :: prose twin of the run farther from 1
\newcommand{\MlstMassRespAxesShareQcd}{76\%}  % experiments/FIGS/data/v2/index.json :: head/atlas_res: share of QCD jets whose groomed-mass response lies inside F9a(b)'s x axis :: sha256 c6c4138ee5c86423
\newcommand{\MlstSculptNPassTen}{41{,}028}  % experiments/FIGS/data/v2/index.json :: QCD test jets passing the cut at QCD efficiency 0.1, every run :: sha256 c6c4138ee5c86423
\newcommand{\MlstSculptFloorTen}{0.00944}  % experiments/FIGS/data/v2/index.json :: Jensen-Shannon distance expected for a random selection of as many QCD jets as pass at 0.1 :: sha256 c6c4138ee5c86423
\newcommand{\MlstSculptFullTen}{0.0524\,\ensuremath{\pm}\,0.0034}  % experiments/FIGS/data/v2/index.json :: Jensen-Shannon distance of the QCD m_SD distribution at QCD efficiency 0.1, L188 :: sha256 c6c4138ee5c86423
\newcommand{\MlstSculptFullOverFloorTen}{5.6\,\ensuremath{\pm}\,0.4}  % experiments/FIGS/data/v2/index.json :: distance at 0.1 over the random-selection floor, L188 :: sha256 c6c4138ee5c86423
\newcommand{\MlstSculptFullOverFloorTenPr}{5.6 times the distance}  % experiments/FIGS/data/v2/index.json :: distance at 0.1 over the random-selection floor, L188 :: sha256 c6c4138ee5c86423 :: prose twin: two significant figures, direction in words
\newcommand{\MlstSculptResTen}{0.0532\,\ensuremath{\pm}\,0.0009}  % experiments/FIGS/data/v2/index.json :: Jensen-Shannon distance of the QCD m_SD distribution at QCD efficiency 0.1, L162 :: sha256 c6c4138ee5c86423
\newcommand{\MlstSculptResOverFloorTen}{5.64\,\ensuremath{\pm}\,0.10}  % experiments/FIGS/data/v2/index.json :: distance at 0.1 over the random-selection floor, L162 :: sha256 c6c4138ee5c86423
\newcommand{\MlstSculptResOverFloorTenPr}{5.6 times the distance}  % experiments/FIGS/data/v2/index.json :: distance at 0.1 over the random-selection floor, L162 :: sha256 c6c4138ee5c86423 :: prose twin: two significant figures, direction in words
\newcommand{\MlstSculptMidTen}{0.0560\,\ensuremath{\pm}\,0.0027}  % experiments/FIGS/data/v2/index.json :: Jensen-Shannon distance of the QCD m_SD distribution at QCD efficiency 0.1, R42_Q1 :: sha256 c6c4138ee5c86423
\newcommand{\MlstSculptMidOverFloorTen}{5.94\,\ensuremath{\pm}\,0.28}  % experiments/FIGS/data/v2/index.json :: distance at 0.1 over the random-selection floor, R42_Q1 :: sha256 c6c4138ee5c86423
\newcommand{\MlstSculptMidOverFloorTenPr}{5.9 times the distance}  % experiments/FIGS/data/v2/index.json :: distance at 0.1 over the random-selection floor, R42_Q1 :: sha256 c6c4138ee5c86423 :: prose twin: two significant figures, direction in words
\newcommand{\MlstSculptCoarseTen}{0.0543\,\ensuremath{\pm}\,0.0010}  % experiments/FIGS/data/v2/index.json :: Jensen-Shannon distance of the QCD m_SD distribution at QCD efficiency 0.1, R16_Q1 :: sha256 c6c4138ee5c86423
\newcommand{\MlstSculptCoarseOverFloorTen}{5.76\,\ensuremath{\pm}\,0.10}  % experiments/FIGS/data/v2/index.json :: distance at 0.1 over the random-selection floor, R16_Q1 :: sha256 c6c4138ee5c86423
\newcommand{\MlstSculptCoarseOverFloorTenPr}{5.8 times the distance}  % experiments/FIGS/data/v2/index.json :: distance at 0.1 over the random-selection floor, R16_Q1 :: sha256 c6c4138ee5c86423 :: prose twin: two significant figures, direction in words
\newcommand{\MlstSculptResMassTen}{0.0422 and 0.0464}  % experiments/FIGS/data/v2/index.json :: Jensen-Shannon distance of the QCD m_SD distribution at QCD efficiency 0.1, L162_MASS :: sha256 c6c4138ee5c86423
\newcommand{\MlstSculptResMassOverFloorTen}{4.47 and 4.91}  % experiments/FIGS/data/v2/index.json :: distance at 0.1 over the random-selection floor, L162_MASS :: sha256 c6c4138ee5c86423
\newcommand{\MlstSculptResMassOverFloorTenPr}{at least 4.5 times the distance}  % experiments/FIGS/data/v2/index.json :: distance at 0.1 over the random-selection floor, L162_MASS :: sha256 c6c4138ee5c86423 :: prose twin: two significant figures, direction in words
\newcommand{\MlstSculptResMassOverFloorTenPrMax}{up to 4.9 times the distance}  % experiments/FIGS/data/v2/index.json :: distance at 0.1 over the random-selection floor, L162_MASS :: sha256 c6c4138ee5c86423 :: prose twin of the run farther from 1
\newcommand{\MlstSculptCoarseMassTen}{0.0522 and 0.0423}  % experiments/FIGS/data/v2/index.json :: Jensen-Shannon distance of the QCD m_SD distribution at QCD efficiency 0.1, R16_Q1_MASS :: sha256 c6c4138ee5c86423
\newcommand{\MlstSculptCoarseMassOverFloorTen}{5.5 and 4.5}  % experiments/FIGS/data/v2/index.json :: distance at 0.1 over the random-selection floor, R16_Q1_MASS :: sha256 c6c4138ee5c86423
\newcommand{\MlstSculptCoarseMassOverFloorTenPr}{at least 4.5 times the distance}  % experiments/FIGS/data/v2/index.json :: distance at 0.1 over the random-selection floor, R16_Q1_MASS :: sha256 c6c4138ee5c86423 :: prose twin: two significant figures, direction in words
\newcommand{\MlstSculptCoarseMassOverFloorTenPrMax}{up to 5.5 times the distance}  % experiments/FIGS/data/v2/index.json :: distance at 0.1 over the random-selection floor, R16_Q1_MASS :: sha256 c6c4138ee5c86423 :: prose twin of the run farther from 1
\newcommand{\MlstSculptResMassOverResTen}{0.789 and 0.860}  % experiments/FIGS/data/v2/index.json :: distance at QCD efficiency 0.1: L162_MASS over L162, run by run :: sha256 c6c4138ee5c86423
\newcommand{\MlstSculptResMassOverResTenPr}{at most 0.86 times the distance}  % experiments/FIGS/data/v2/index.json :: distance at QCD efficiency 0.1: L162_MASS over L162, run by run :: sha256 c6c4138ee5c86423 :: prose twin: two significant figures, direction in words
\newcommand{\MlstSculptResMassOverResTenPrMax}{down to 0.79 times the distance}  % experiments/FIGS/data/v2/index.json :: distance at QCD efficiency 0.1: L162_MASS over L162, run by run :: sha256 c6c4138ee5c86423 :: prose twin of the run farther from 1
\newcommand{\MlstSculptCoarseMassOverCoarseTen}{0.94 and 0.79}  % experiments/FIGS/data/v2/index.json :: distance at QCD efficiency 0.1: R16_Q1_MASS over R16_Q1, run by run :: sha256 c6c4138ee5c86423
\newcommand{\MlstSculptCoarseMassOverCoarseTenPr}{equal within 30\%}  % experiments/FIGS/data/v2/index.json :: distance at QCD efficiency 0.1: R16_Q1_MASS over R16_Q1, run by run :: sha256 c6c4138ee5c86423 :: prose twin: two significant figures, direction in words
\newcommand{\MlstSculptCoarseMassOverCoarseTenPrMax}{equal within 30\%}  % experiments/FIGS/data/v2/index.json :: distance at QCD efficiency 0.1: R16_Q1_MASS over R16_Q1, run by run :: sha256 c6c4138ee5c86423 :: prose twin of the run farther from 1
\newcommand{\MlstSculptCoarseOverFullTen}{1.04\,\ensuremath{\pm}\,0.05}  % experiments/FIGS/data/v2/index.json :: distance at QCD efficiency 0.1: R16_Q1 over L188, run by run :: sha256 c6c4138ee5c86423
\newcommand{\MlstSculptCoarseOverFullTenPr}{equal within 6\%}  % experiments/FIGS/data/v2/index.json :: distance at QCD efficiency 0.1: R16_Q1 over L188, run by run :: sha256 c6c4138ee5c86423 :: prose twin: two significant figures, direction in words
\newcommand{\MlstSculptResOverFullTen}{1.02\,\ensuremath{\pm}\,0.06}  % experiments/FIGS/data/v2/index.json :: distance at QCD efficiency 0.1: L162 over L188, run by run :: sha256 c6c4138ee5c86423
\newcommand{\MlstSculptResOverFullTenPr}{equal within 6\%}  % experiments/FIGS/data/v2/index.json :: distance at QCD efficiency 0.1: L162 over L188, run by run :: sha256 c6c4138ee5c86423 :: prose twin: two significant figures, direction in words
\newcommand{\MlstSculptMidOverFullTen}{1.069\,\ensuremath{\pm}\,0.020}  % experiments/FIGS/data/v2/index.json :: distance at QCD efficiency 0.1: R42_Q1 over L188, run by run :: sha256 c6c4138ee5c86423
\newcommand{\MlstSculptMidOverFullTenPr}{1.1 times the distance}  % experiments/FIGS/data/v2/index.json :: distance at QCD efficiency 0.1: R42_Q1 over L188, run by run :: sha256 c6c4138ee5c86423 :: prose twin: two significant figures, direction in words
\newcommand{\MlstSculptCoarseOverMidTen}{0.971\,\ensuremath{\pm}\,0.035}  % experiments/FIGS/data/v2/index.json :: distance at QCD efficiency 0.1: R16_Q1 over R42_Q1, run by run :: sha256 c6c4138ee5c86423
\newcommand{\MlstSculptCoarseOverMidTenPr}{equal within 4\%}  % experiments/FIGS/data/v2/index.json :: distance at QCD efficiency 0.1: R16_Q1 over R42_Q1, run by run :: sha256 c6c4138ee5c86423 :: prose twin: two significant figures, direction in words
\newcommand{\MlstSculptNPassOne}{4{,}103}  % experiments/FIGS/data/v2/index.json :: QCD test jets passing the cut at QCD efficiency 0.01, every run :: sha256 c6c4138ee5c86423
\newcommand{\MlstSculptFloorOne}{0.0313}  % experiments/FIGS/data/v2/index.json :: Jensen-Shannon distance expected for a random selection of as many QCD jets as pass at 0.01 :: sha256 c6c4138ee5c86423
\newcommand{\MlstSculptFullOne}{0.0439\,\ensuremath{\pm}\,0.0024}  % experiments/FIGS/data/v2/index.json :: Jensen-Shannon distance of the QCD m_SD distribution at QCD efficiency 0.01, L188 :: sha256 c6c4138ee5c86423
\newcommand{\MlstSculptFullOverFloorOne}{1.40\,\ensuremath{\pm}\,0.08}  % experiments/FIGS/data/v2/index.json :: distance at 0.01 over the random-selection floor, L188 :: sha256 c6c4138ee5c86423
\newcommand{\MlstSculptFullOverFloorOnePr}{1.4 times the distance}  % experiments/FIGS/data/v2/index.json :: distance at 0.01 over the random-selection floor, L188 :: sha256 c6c4138ee5c86423 :: prose twin: two significant figures, direction in words
\newcommand{\MlstSculptResOne}{0.0429\,\ensuremath{\pm}\,0.0029}  % experiments/FIGS/data/v2/index.json :: Jensen-Shannon distance of the QCD m_SD distribution at QCD efficiency 0.01, L162 :: sha256 c6c4138ee5c86423
\newcommand{\MlstSculptResOverFloorOne}{1.37\,\ensuremath{\pm}\,0.09}  % experiments/FIGS/data/v2/index.json :: distance at 0.01 over the random-selection floor, L162 :: sha256 c6c4138ee5c86423
\newcommand{\MlstSculptResOverFloorOnePr}{1.4 times the distance}  % experiments/FIGS/data/v2/index.json :: distance at 0.01 over the random-selection floor, L162 :: sha256 c6c4138ee5c86423 :: prose twin: two significant figures, direction in words
\newcommand{\MlstSculptMidOne}{0.051\,\ensuremath{\pm}\,0.007}  % experiments/FIGS/data/v2/index.json :: Jensen-Shannon distance of the QCD m_SD distribution at QCD efficiency 0.01, R42_Q1 :: sha256 c6c4138ee5c86423
\newcommand{\MlstSculptMidOverFloorOne}{1.62\,\ensuremath{\pm}\,0.23}  % experiments/FIGS/data/v2/index.json :: distance at 0.01 over the random-selection floor, R42_Q1 :: sha256 c6c4138ee5c86423
\newcommand{\MlstSculptMidOverFloorOnePr}{1.6 times the distance}  % experiments/FIGS/data/v2/index.json :: distance at 0.01 over the random-selection floor, R42_Q1 :: sha256 c6c4138ee5c86423 :: prose twin: two significant figures, direction in words
\newcommand{\MlstSculptCoarseOne}{0.0561\,\ensuremath{\pm}\,0.0017}  % experiments/FIGS/data/v2/index.json :: Jensen-Shannon distance of the QCD m_SD distribution at QCD efficiency 0.01, R16_Q1 :: sha256 c6c4138ee5c86423
\newcommand{\MlstSculptCoarseOverFloorOne}{1.79\,\ensuremath{\pm}\,0.06}  % experiments/FIGS/data/v2/index.json :: distance at 0.01 over the random-selection floor, R16_Q1 :: sha256 c6c4138ee5c86423
\newcommand{\MlstSculptCoarseOverFloorOnePr}{1.8 times the distance}  % experiments/FIGS/data/v2/index.json :: distance at 0.01 over the random-selection floor, R16_Q1 :: sha256 c6c4138ee5c86423 :: prose twin: two significant figures, direction in words
\newcommand{\MlstSculptResMassOne}{0.0424 and 0.0442}  % experiments/FIGS/data/v2/index.json :: Jensen-Shannon distance of the QCD m_SD distribution at QCD efficiency 0.01, L162_MASS :: sha256 c6c4138ee5c86423
\newcommand{\MlstSculptResMassOverFloorOne}{1.356 and 1.411}  % experiments/FIGS/data/v2/index.json :: distance at 0.01 over the random-selection floor, L162_MASS :: sha256 c6c4138ee5c86423
\newcommand{\MlstSculptResMassOverFloorOnePr}{at least 1.4 times the distance}  % experiments/FIGS/data/v2/index.json :: distance at 0.01 over the random-selection floor, L162_MASS :: sha256 c6c4138ee5c86423 :: prose twin: two significant figures, direction in words
\newcommand{\MlstSculptResMassOverFloorOnePrMax}{up to 1.4 times the distance}  % experiments/FIGS/data/v2/index.json :: distance at 0.01 over the random-selection floor, L162_MASS :: sha256 c6c4138ee5c86423 :: prose twin of the run farther from 1
\newcommand{\MlstSculptCoarseMassOne}{0.0405 and 0.0452}  % experiments/FIGS/data/v2/index.json :: Jensen-Shannon distance of the QCD m_SD distribution at QCD efficiency 0.01, R16_Q1_MASS :: sha256 c6c4138ee5c86423
\newcommand{\MlstSculptCoarseMassOverFloorOne}{1.29 and 1.44}  % experiments/FIGS/data/v2/index.json :: distance at 0.01 over the random-selection floor, R16_Q1_MASS :: sha256 c6c4138ee5c86423
\newcommand{\MlstSculptCoarseMassOverFloorOnePr}{at least 1.3 times the distance}  % experiments/FIGS/data/v2/index.json :: distance at 0.01 over the random-selection floor, R16_Q1_MASS :: sha256 c6c4138ee5c86423 :: prose twin: two significant figures, direction in words
\newcommand{\MlstSculptCoarseMassOverFloorOnePrMax}{up to 1.4 times the distance}  % experiments/FIGS/data/v2/index.json :: distance at 0.01 over the random-selection floor, R16_Q1_MASS :: sha256 c6c4138ee5c86423 :: prose twin of the run farther from 1
\newcommand{\MlstSculptResMassOverResOne}{0.962 and 0.983}  % experiments/FIGS/data/v2/index.json :: distance at QCD efficiency 0.01: L162_MASS over L162, run by run :: sha256 c6c4138ee5c86423
\newcommand{\MlstSculptResMassOverResOnePr}{equal within 4\%}  % experiments/FIGS/data/v2/index.json :: distance at QCD efficiency 0.01: L162_MASS over L162, run by run :: sha256 c6c4138ee5c86423 :: prose twin: two significant figures, direction in words
\newcommand{\MlstSculptResMassOverResOnePrMax}{equal within 4\%}  % experiments/FIGS/data/v2/index.json :: distance at QCD efficiency 0.01: L162_MASS over L162, run by run :: sha256 c6c4138ee5c86423 :: prose twin of the run farther from 1
\newcommand{\MlstSculptCoarseMassOverCoarseOne}{0.70 and 0.82}  % experiments/FIGS/data/v2/index.json :: distance at QCD efficiency 0.01: R16_Q1_MASS over R16_Q1, run by run :: sha256 c6c4138ee5c86423
\newcommand{\MlstSculptCoarseMassOverCoarseOnePr}{at most 0.82 times the distance}  % experiments/FIGS/data/v2/index.json :: distance at QCD efficiency 0.01: R16_Q1_MASS over R16_Q1, run by run :: sha256 c6c4138ee5c86423 :: prose twin: two significant figures, direction in words
\newcommand{\MlstSculptCoarseMassOverCoarseOnePrMax}{down to 0.70 times the distance}  % experiments/FIGS/data/v2/index.json :: distance at QCD efficiency 0.01: R16_Q1_MASS over R16_Q1, run by run :: sha256 c6c4138ee5c86423 :: prose twin of the run farther from 1
\newcommand{\MlstSculptCoarseOverFullOne}{1.28\,\ensuremath{\pm}\,0.04}  % experiments/FIGS/data/v2/index.json :: distance at QCD efficiency 0.01: R16_Q1 over L188, run by run :: sha256 c6c4138ee5c86423
\newcommand{\MlstSculptCoarseOverFullOnePr}{1.3 times the distance}  % experiments/FIGS/data/v2/index.json :: distance at QCD efficiency 0.01: R16_Q1 over L188, run by run :: sha256 c6c4138ee5c86423 :: prose twin: two significant figures, direction in words
\newcommand{\MlstSculptResOverFullOne}{0.98\,\ensuremath{\pm}\,0.05}  % experiments/FIGS/data/v2/index.json :: distance at QCD efficiency 0.01: L162 over L188, run by run :: sha256 c6c4138ee5c86423
\newcommand{\MlstSculptResOverFullOnePr}{equal within 5\%}  % experiments/FIGS/data/v2/index.json :: distance at QCD efficiency 0.01: L162 over L188, run by run :: sha256 c6c4138ee5c86423 :: prose twin: two significant figures, direction in words
\newcommand{\MlstSculptMidOverFullOne}{1.15\,\ensuremath{\pm}\,0.12}  % experiments/FIGS/data/v2/index.json :: distance at QCD efficiency 0.01: R42_Q1 over L188, run by run :: sha256 c6c4138ee5c86423
\newcommand{\MlstSculptMidOverFullOnePr}{equal within 20\%}  % experiments/FIGS/data/v2/index.json :: distance at QCD efficiency 0.01: R42_Q1 over L188, run by run :: sha256 c6c4138ee5c86423 :: prose twin: two significant figures, direction in words
\newcommand{\MlstSculptCoarseOverMidOne}{1.13\,\ensuremath{\pm}\,0.16}  % experiments/FIGS/data/v2/index.json :: distance at QCD efficiency 0.01: R16_Q1 over R42_Q1, run by run :: sha256 c6c4138ee5c86423
\newcommand{\MlstSculptCoarseOverMidOnePr}{equal within 20\%}  % experiments/FIGS/data/v2/index.json :: distance at QCD efficiency 0.01: R16_Q1 over R42_Q1, run by run :: sha256 c6c4138ee5c86423 :: prose twin: two significant figures, direction in words
\newcommand{\MlstSculptDiffEffLo}{0.52\%}  % experiments/FIGS/data/v2/index.json :: lowest stored QCD efficiency of the unbroken stretch around the tighter working point at which the 17/188 distance ratio exceeds 1 by twice its SD :: sha256 c6c4138ee5c86423 :: passing grid points: 0.001, 0.005245, 0.007937, 0.012011
\newcommand{\MlstSculptDiffEffHi}{1.2\%}  % experiments/FIGS/data/v2/index.json :: highest efficiency of that stretch :: sha256 c6c4138ee5c86423
\newcommand{\MlstSculptLowestEff}{0.10\%}  % experiments/FIGS/data/v2/index.json :: lowest stored QCD efficiency of the sculpting curves :: sha256 c6c4138ee5c86423
\newcommand{\MlstSculptLowestNPass}{410}  % experiments/FIGS/data/v2/index.json :: QCD test jets passing at the lowest stored efficiency :: sha256 c6c4138ee5c86423
\newcommand{\MlstSculptLowestMaxFloorSd}{1.7}  % experiments/FIGS/data/v2/index.json :: largest distance of a configuration's mean from the random-selection floor at the lowest stored efficiency, in SDs of the floor over random selections :: sha256 c6c4138ee5c86423
\newcommand{\MlstSculptNConfigs}{6}  % experiments/FIGS/data/v2/index.json :: configurations with a sculpting curve (count) :: sha256 c6c4138ee5c86423
\newcommand{\LpJciiTestQcdPercent}{20.5\%}  % experiments/FIGS/data/v2/jc2_test/counts.json :: QCD share of the jets, per cent :: sha256 3574ba1f6f2a59f0
\newcommand{\LpJciiTestTwoProngPercent}{38.7\%}  % experiments/FIGS/data/v2/jc2_test/counts.json :: two-prong share of the jets, per cent :: sha256 3574ba1f6f2a59f0
\newcommand{\LpJciiTrainQcdPercent}{13.5\%}  % experiments/FT/data/jc2_v2_manifest.json :: QCD share of the jets, per cent :: sha256 4bd68655919b14e9
\newcommand{\LpJciiTrainTwoProngPercent}{15.4\%}  % experiments/FT/data/jc2_v2_manifest.json :: two-prong share of the jets, per cent :: sha256 4bd68655919b14e9
\newcommand{\LpJciiTestMinAucJets}{10}  % experiments/FIGS/data/v2/jc2_test/counts.json :: fewest jets of a class among the macro-AUC rows :: sha256 3574ba1f6f2a59f0
\newcommand{\LpJciiTrainMinJets}{11}  % experiments/FT/data/jc2_v2_manifest.json :: fewest jets of a class in the 10^6 training subset :: sha256 4bd68655919b14e9
\newcommand{\LpNTestJci}{2{,}000{,}000}  % experiments/FIGS/data/v2/index.json :: jc1 test jets the probes score :: sha256 c6c4138ee5c86423
\newcommand{\LpNTrainMaxJci}{1{,}000{,}000}  % experiments/FIGS/data/v2/index.json :: jc1 largest training subset :: sha256 c6c4138ee5c86423
\newcommand{\LpNTestTop}{404{,}000}  % experiments/FIGS/data/v2/index.json :: top test jets the probes score :: sha256 c6c4138ee5c86423
\newcommand{\LpNTrainMaxTop}{1{,}200{,}000}  % experiments/FIGS/data/v2/index.json :: top largest training subset :: sha256 c6c4138ee5c86423
\newcommand{\LpNTestQg}{200{,}000}  % experiments/FIGS/data/v2/index.json :: qg test jets the probes score :: sha256 c6c4138ee5c86423
\newcommand{\LpNTrainMaxQg}{1{,}600{,}000}  % experiments/FIGS/data/v2/index.json :: qg largest training subset :: sha256 c6c4138ee5c86423
\newcommand{\LpNTestAucJci}{2{,}000{,}000}  % experiments/FIGS/data/v2/index.json :: JetClass test jets of the macro AUC :: sha256 c6c4138ee5c86423
\newcommand{\LpPairNValMin}{1{,}068}  % experiments/FIGS/data/v2/index.json :: validation jets of the pair probe with the fewest :: sha256 c6c4138ee5c86423
\newcommand{\LpPairNValMax}{20{,}195}  % experiments/FIGS/data/v2/index.json :: validation jets of the pair probe with the most :: sha256 c6c4138ee5c86423
\newcommand{\LpNValJci}{200{,}000}  % experiments/FIGS/data/v2/index.json :: JetClass probe validation jets :: sha256 c6c4138ee5c86423
\newcommand{\SculptEffHigh}{10\%}  % experiments/EVAL/mass_sculpting.py :: WORKING_POINTS :: sha256 e85d125c85dac054 :: QCD efficiency of a sculpting cut
\newcommand{\SculptEffLow}{1\%}  % experiments/EVAL/mass_sculpting.py :: WORKING_POINTS :: sha256 e85d125c85dac054 :: QCD efficiency of a sculpting cut
\newcommand{\LpMlpHidden}{256}  % experiments/EVAL/probe.py :: fit_mlp hidden layer :: sha256 485db0966e660588 :: units of the MLP probe's hidden layer
\newcommand{\LpMlpSeedsWord}{three}  % experiments/EVAL/probe.py :: MLP_SEEDS :: sha256 485db0966e660588 :: seeds of the MLP probe whose scores are averaged
\newcommand{\MassMinTrainPerClass}{50}  % experiments/EVAL/mass_resolution.py :: MIN_TRAIN_PER_CLASS :: sha256 26acb0e315094bd1 :: fewest training jets of a class the mass regression keeps
\newcommand{\MlstFtLr}{\ensuremath{10^{-4}}}  % experiments/FIGS/data/v2/ft_mlst/manifests.json :: cells[best70].lr :: sha256 a689f21e894f48ae
\newcommand{\MlstFtHeadMult}{50}  % experiments/FIGS/data/v2/ft_mlst/manifests.json :: cells[best70].head_lr_mult :: sha256 a689f21e894f48ae
\newcommand{\MlstFtWeightDecay}{0.01}  % experiments/FIGS/data/v2/ft_mlst/manifests.json :: cells[best70].weight_decay :: sha256 a689f21e894f48ae
\newcommand{\MlstFtBatch}{512}  % experiments/FIGS/data/v2/ft_mlst/manifests.json :: cells[best70].batch_size :: sha256 a689f21e894f48ae
\newcommand{\MlstFtBenchEpochs}{20}  % experiments/FIGS/data/v2/ft_mlst/manifests.json :: cells[best70, top/qg].epochs :: sha256 a689f21e894f48ae
\newcommand{\MlstFtEpochsEThree}{50}  % experiments/FIGS/data/v2/ft_mlst/manifests.json :: cells[best70, leg 1/2, n_train 1000].epochs :: sha256 a689f21e894f48ae
\newcommand{\MlstFtEpochsEFour}{50}  % experiments/FIGS/data/v2/ft_mlst/manifests.json :: cells[best70, leg 1/2, n_train 10000].epochs :: sha256 a689f21e894f48ae
\newcommand{\MlstFtEpochsEFive}{30}  % experiments/FIGS/data/v2/ft_mlst/manifests.json :: cells[best70, leg 1/2, n_train 100000].epochs :: sha256 a689f21e894f48ae
\newcommand{\MlstFtEpochsESix}{10}  % experiments/FIGS/data/v2/ft_mlst/manifests.json :: cells[best70, leg 1/2, n_train 1000000].epochs :: sha256 a689f21e894f48ae
\newcommand{\MlstFtHeadDecayShare}{30\%}  % experiments/FT/k8s/job-ft-v2-legs-best70-l188-s3-raunav.yaml :: no --lr-scheduler: weaver 0.4.17 flat+decay, the head's last epochs :: sha256 19f59aa54a385c11
\newcommand{\MlstFtHeadDecayTo}{1\%}  % experiments/FT/k8s/job-ft-v2-legs-best70-l188-s3-raunav.yaml :: no --lr-scheduler: weaver 0.4.17 flat+decay, final factor :: sha256 19f59aa54a385c11
\newcommand{\MlstFtHeadDropout}{0.1}  % experiments/FT/k8s/job-ft-v2-legs-best70-l188-s3-raunav.yaml :: fc_params dropout of the fine-tuning head :: sha256 19f59aa54a385c11
\newcommand{\MlstFtMinEpochJets}{10{,}000}  % experiments/FIGS/data/v2/ft_mlst/manifests.json :: cells[best70].samples_per_epoch (smallest) :: sha256 a689f21e894f48ae
\newcommand{\MlstFtValJcii}{20{,}480}  % experiments/FIGS/data/v2/ft_mlst/manifests.json :: cells[best70, leg 1].val_jets :: sha256 a689f21e894f48ae
\newcommand{\MlstFtValSmall}{20{,}000}  % experiments/FT/k8s/job-ft-v2-legs-best70-l188-s3-raunav.yaml :: JetClass --samples-per-epoch-val; top and q/g at N <= 1e4 :: sha256 19f59aa54a385c11
\newcommand{\MlstFtValLarge}{200{,}000}  % experiments/FT/k8s/job-ft-v2-bench-best70-l188-s1-top-raunav.yaml :: top and q/g --samples-per-epoch-val at N > 1e4 :: sha256 1e38d4f542119ea3
\newcommand{\MlstFtGpuTypes}{V100}  % experiments/FIGS/data/v2/ft_mlst/manifests.json :: cells[best70].gpu_device_name :: sha256 a689f21e894f48ae
\newcommand{\MlstFtCellsDone}{445}  % experiments/FIGS/data/v2/ft_mlst/manifests.json :: DONE paper-run fine-tuning cells (count) :: sha256 a689f21e894f48ae
\newcommand{\MlstFtRefLr}{\ensuremath{10^{-3}}}  % experiments/FIGS/data/v2/ft_mlst/manifests.json :: cells[ref_v100].lr :: sha256 a689f21e894f48ae
\newcommand{\MlstFtRefSeedsWord}{three}  % experiments/FIGS/data/v2/ft_mlst/manifests.json :: cells[ref_v100].ft_seed (count) :: sha256 a689f21e894f48ae
\newcommand{\MlstFtRefGpuTypes}{V100}  % experiments/FIGS/data/v2/ft_mlst/manifests.json :: cells[ref_v100].gpu_device_name :: sha256 a689f21e894f48ae
\newcommand{\MlstFtRefWeightDecayBench}{0.01}  % experiments/FIGS/data/v2/ft_mlst/manifests.json :: cells[ref_v100, top/qg].weight_decay :: sha256 a689f21e894f48ae
\newcommand{\MlstBootResamples}{200}  % experiments/FIGS/data/v2/index.json :: bootstrap resamples of the test jets :: sha256 c6c4138ee5c86423
\newcommand{\MlstBootTestMaxProbe}{0.21}  % experiments/FIGS/data/v2/index.json :: largest test-sample SD of the ln ratio of 1 - AUC on JetClass-II, linear probe, label sets over the 188-class model, per cent :: sha256 c6c4138ee5c86423
\newcommand{\MlstBootTestMaxFt}{0.15}  % experiments/FIGS/data/v2/index.json :: largest test-sample SD of the ln ratio of 1 - AUC on JetClass-II, fine-tuning, label sets over the 188-class model, per cent :: sha256 c6c4138ee5c86423
\newcommand{\MlstBootTestMaxRthirtyProbe}{3.6}  % experiments/FIGS/data/v2/index.json :: largest test-sample SD of the ln R30 ratio on top and quark/gluon, linear probe (counting), label sets over the 188-class model, per cent :: sha256 c6c4138ee5c86423
\newcommand{\MlstBootTestMaxRthirtyFt}{13}  % experiments/FIGS/data/v2/index.json :: largest test-sample SD of the ln R30 ratio on top and quark/gluon, fine-tuning (counting), label sets over the 188-class model, per cent :: sha256 c6c4138ee5c86423
\newcommand{\MlstFtCmpJcTwoResOverFullEFour}{0.989\,\ensuremath{\pm}\,0.011}  % experiments/FIGS/data/v2/index.json :: fine-tuning jc2 N=10000: (1 - macro_auc_ovr) of L162 over L188, run by run :: sha256 c6c4138ee5c86423
\newcommand{\MlstFtCmpJcTwoResOverFullEFourPr}{equal within 2\%}  % experiments/FIGS/data/v2/index.json :: fine-tuning jc2 N=10000: (1 - macro_auc_ovr) of L162 over L188, run by run :: sha256 c6c4138ee5c86423 :: prose twin: two significant figures, direction in words
\newcommand{\MlstFtCmpJcTwoMidOverFullEFour}{1.191\,\ensuremath{\pm}\,0.030}  % experiments/FIGS/data/v2/index.json :: fine-tuning jc2 N=10000: (1 - macro_auc_ovr) of R42_Q1 over L188, run by run :: sha256 c6c4138ee5c86423
\newcommand{\MlstFtCmpJcTwoMidOverFullEFourPr}{19\% more errors}  % experiments/FIGS/data/v2/index.json :: fine-tuning jc2 N=10000: (1 - macro_auc_ovr) of R42_Q1 over L188, run by run :: sha256 c6c4138ee5c86423 :: prose twin: two significant figures, direction in words
\newcommand{\MlstFtCmpJcTwoCoarseOverFullEFour}{1.353\,\ensuremath{\pm}\,0.003}  % experiments/FIGS/data/v2/index.json :: fine-tuning jc2 N=10000: (1 - macro_auc_ovr) of R16_Q1 over L188, run by run :: sha256 c6c4138ee5c86423
\newcommand{\MlstFtCmpJcTwoCoarseOverFullEFourPr}{35\% more errors}  % experiments/FIGS/data/v2/index.json :: fine-tuning jc2 N=10000: (1 - macro_auc_ovr) of R16_Q1 over L188, run by run :: sha256 c6c4138ee5c86423 :: prose twin: two significant figures, direction in words
\newcommand{\MlstFtCmpJcTwoResMassOverResEFour}{1.016 and 1.005}  % experiments/FIGS/data/v2/index.json :: fine-tuning jc2 N=10000: (1 - macro_auc_ovr) of L162_MASS over L162, run by run :: sha256 c6c4138ee5c86423
\newcommand{\MlstFtCmpJcTwoResMassOverResEFourPr}{equal within 2\%}  % experiments/FIGS/data/v2/index.json :: fine-tuning jc2 N=10000: (1 - macro_auc_ovr) of L162_MASS over L162, run by run :: sha256 c6c4138ee5c86423 :: prose twin: two significant figures, direction in words
\newcommand{\MlstFtCmpJcTwoResMassOverResEFourPrMax}{equal within 2\%}  % experiments/FIGS/data/v2/index.json :: fine-tuning jc2 N=10000: (1 - macro_auc_ovr) of L162_MASS over L162, run by run :: sha256 c6c4138ee5c86423 :: prose twin of the run farther from 1
\newcommand{\MlstFtCmpJcTwoCoarseMassOverCoarseEFour}{0.989 and 1.000}  % experiments/FIGS/data/v2/index.json :: fine-tuning jc2 N=10000: (1 - macro_auc_ovr) of R16_Q1_MASS over R16_Q1, run by run :: sha256 c6c4138ee5c86423
\newcommand{\MlstFtCmpJcTwoCoarseMassOverCoarseEFourPr}{equal within 2\%}  % experiments/FIGS/data/v2/index.json :: fine-tuning jc2 N=10000: (1 - macro_auc_ovr) of R16_Q1_MASS over R16_Q1, run by run :: sha256 c6c4138ee5c86423 :: prose twin: two significant figures, direction in words
\newcommand{\MlstFtCmpJcTwoCoarseMassOverCoarseEFourPrMax}{equal within 2\%}  % experiments/FIGS/data/v2/index.json :: fine-tuning jc2 N=10000: (1 - macro_auc_ovr) of R16_Q1_MASS over R16_Q1, run by run :: sha256 c6c4138ee5c86423 :: prose twin of the run farther from 1
\newcommand{\MlstFtCmpJcTwoResOverFullEFive}{1.011\,\ensuremath{\pm}\,0.005}  % experiments/FIGS/data/v2/index.json :: fine-tuning jc2 N=100000: (1 - macro_auc_ovr) of L162 over L188, run by run :: sha256 c6c4138ee5c86423
\newcommand{\MlstFtCmpJcTwoResOverFullEFivePr}{1.1\% more errors}  % experiments/FIGS/data/v2/index.json :: fine-tuning jc2 N=100000: (1 - macro_auc_ovr) of L162 over L188, run by run :: sha256 c6c4138ee5c86423 :: prose twin: two significant figures, direction in words
\newcommand{\MlstFtCmpJcTwoMidOverFullEFive}{1.064\,\ensuremath{\pm}\,0.009}  % experiments/FIGS/data/v2/index.json :: fine-tuning jc2 N=100000: (1 - macro_auc_ovr) of R42_Q1 over L188, run by run :: sha256 c6c4138ee5c86423
\newcommand{\MlstFtCmpJcTwoMidOverFullEFivePr}{6.4\% more errors}  % experiments/FIGS/data/v2/index.json :: fine-tuning jc2 N=100000: (1 - macro_auc_ovr) of R42_Q1 over L188, run by run :: sha256 c6c4138ee5c86423 :: prose twin: two significant figures, direction in words
\newcommand{\MlstFtCmpJcTwoCoarseOverFullEFive}{1.157\,\ensuremath{\pm}\,0.001}  % experiments/FIGS/data/v2/index.json :: fine-tuning jc2 N=100000: (1 - macro_auc_ovr) of R16_Q1 over L188, run by run :: sha256 c6c4138ee5c86423
\newcommand{\MlstFtCmpJcTwoCoarseOverFullEFivePr}{16\% more errors}  % experiments/FIGS/data/v2/index.json :: fine-tuning jc2 N=100000: (1 - macro_auc_ovr) of R16_Q1 over L188, run by run :: sha256 c6c4138ee5c86423 :: prose twin: two significant figures, direction in words
\newcommand{\MlstFtCmpJcTwoResMassOverResEFive}{1.027 and 1.021}  % experiments/FIGS/data/v2/index.json :: fine-tuning jc2 N=100000: (1 - macro_auc_ovr) of L162_MASS over L162, run by run :: sha256 c6c4138ee5c86423
\newcommand{\MlstFtCmpJcTwoResMassOverResEFivePr}{at least 2.1\% more errors}  % experiments/FIGS/data/v2/index.json :: fine-tuning jc2 N=100000: (1 - macro_auc_ovr) of L162_MASS over L162, run by run :: sha256 c6c4138ee5c86423 :: prose twin: two significant figures, direction in words
\newcommand{\MlstFtCmpJcTwoResMassOverResEFivePrMax}{up to 2.7\% more errors}  % experiments/FIGS/data/v2/index.json :: fine-tuning jc2 N=100000: (1 - macro_auc_ovr) of L162_MASS over L162, run by run :: sha256 c6c4138ee5c86423 :: prose twin of the run farther from 1
\newcommand{\MlstFtCmpJcTwoCoarseMassOverCoarseEFive}{1.029 and 1.022}  % experiments/FIGS/data/v2/index.json :: fine-tuning jc2 N=100000: (1 - macro_auc_ovr) of R16_Q1_MASS over R16_Q1, run by run :: sha256 c6c4138ee5c86423
\newcommand{\MlstFtCmpJcTwoCoarseMassOverCoarseEFivePr}{at least 2.2\% more errors}  % experiments/FIGS/data/v2/index.json :: fine-tuning jc2 N=100000: (1 - macro_auc_ovr) of R16_Q1_MASS over R16_Q1, run by run :: sha256 c6c4138ee5c86423 :: prose twin: two significant figures, direction in words
\newcommand{\MlstFtCmpJcTwoCoarseMassOverCoarseEFivePrMax}{up to 2.9\% more errors}  % experiments/FIGS/data/v2/index.json :: fine-tuning jc2 N=100000: (1 - macro_auc_ovr) of R16_Q1_MASS over R16_Q1, run by run :: sha256 c6c4138ee5c86423 :: prose twin of the run farther from 1
\newcommand{\MlstFtCmpJcTwoResOverFullESix}{1.014\,\ensuremath{\pm}\,0.011}  % experiments/FIGS/data/v2/index.json :: fine-tuning jc2 N=1000000: (1 - macro_auc_ovr) of L162 over L188, run by run :: sha256 c6c4138ee5c86423
\newcommand{\MlstFtCmpJcTwoResOverFullESixPr}{equal within 2\%}  % experiments/FIGS/data/v2/index.json :: fine-tuning jc2 N=1000000: (1 - macro_auc_ovr) of L162 over L188, run by run :: sha256 c6c4138ee5c86423 :: prose twin: two significant figures, direction in words
\newcommand{\MlstFtCmpJcTwoMidOverFullESix}{1.043\,\ensuremath{\pm}\,0.002}  % experiments/FIGS/data/v2/index.json :: fine-tuning jc2 N=1000000: (1 - macro_auc_ovr) of R42_Q1 over L188, run by run :: sha256 c6c4138ee5c86423
\newcommand{\MlstFtCmpJcTwoMidOverFullESixPr}{4.3\% more errors}  % experiments/FIGS/data/v2/index.json :: fine-tuning jc2 N=1000000: (1 - macro_auc_ovr) of R42_Q1 over L188, run by run :: sha256 c6c4138ee5c86423 :: prose twin: two significant figures, direction in words
\newcommand{\MlstFtCmpJcTwoCoarseOverFullESix}{1.123\,\ensuremath{\pm}\,0.003}  % experiments/FIGS/data/v2/index.json :: fine-tuning jc2 N=1000000: (1 - macro_auc_ovr) of R16_Q1 over L188, run by run :: sha256 c6c4138ee5c86423
\newcommand{\MlstFtCmpJcTwoCoarseOverFullESixPr}{12\% more errors}  % experiments/FIGS/data/v2/index.json :: fine-tuning jc2 N=1000000: (1 - macro_auc_ovr) of R16_Q1 over L188, run by run :: sha256 c6c4138ee5c86423 :: prose twin: two significant figures, direction in words
\newcommand{\MlstFtCmpJcTwoResMassOverResESix}{1.009 and 1.016}  % experiments/FIGS/data/v2/index.json :: fine-tuning jc2 N=1000000: (1 - macro_auc_ovr) of L162_MASS over L162, run by run :: sha256 c6c4138ee5c86423
\newcommand{\MlstFtCmpJcTwoResMassOverResESixPr}{equal within 2\%}  % experiments/FIGS/data/v2/index.json :: fine-tuning jc2 N=1000000: (1 - macro_auc_ovr) of L162_MASS over L162, run by run :: sha256 c6c4138ee5c86423 :: prose twin: two significant figures, direction in words
\newcommand{\MlstFtCmpJcTwoResMassOverResESixPrMax}{equal within 2\%}  % experiments/FIGS/data/v2/index.json :: fine-tuning jc2 N=1000000: (1 - macro_auc_ovr) of L162_MASS over L162, run by run :: sha256 c6c4138ee5c86423 :: prose twin of the run farther from 1
\newcommand{\MlstFtCmpJcTwoCoarseMassOverCoarseESix}{1.023 and 1.017}  % experiments/FIGS/data/v2/index.json :: fine-tuning jc2 N=1000000: (1 - macro_auc_ovr) of R16_Q1_MASS over R16_Q1, run by run :: sha256 c6c4138ee5c86423
\newcommand{\MlstFtCmpJcTwoCoarseMassOverCoarseESixPr}{at least 1.7\% more errors}  % experiments/FIGS/data/v2/index.json :: fine-tuning jc2 N=1000000: (1 - macro_auc_ovr) of R16_Q1_MASS over R16_Q1, run by run :: sha256 c6c4138ee5c86423 :: prose twin: two significant figures, direction in words
\newcommand{\MlstFtCmpJcTwoCoarseMassOverCoarseESixPrMax}{up to 2.3\% more errors}  % experiments/FIGS/data/v2/index.json :: fine-tuning jc2 N=1000000: (1 - macro_auc_ovr) of R16_Q1_MASS over R16_Q1, run by run :: sha256 c6c4138ee5c86423 :: prose twin of the run farther from 1
\newcommand{\MlstFtCmpJcTwoQcdResOverFullEFour}{1.038\,\ensuremath{\pm}\,0.005}  % experiments/FIGS/data/v2/index.json :: fine-tuning jc2_qcd N=10000: (1 - macro_auc_qcd) of L162 over L188, run by run :: sha256 c6c4138ee5c86423
\newcommand{\MlstFtCmpJcTwoQcdResOverFullEFourPr}{3.8\% more errors}  % experiments/FIGS/data/v2/index.json :: fine-tuning jc2_qcd N=10000: (1 - macro_auc_qcd) of L162 over L188, run by run :: sha256 c6c4138ee5c86423 :: prose twin: two significant figures, direction in words
\newcommand{\MlstFtCmpJcTwoQcdMidOverFullEFour}{1.164\,\ensuremath{\pm}\,0.018}  % experiments/FIGS/data/v2/index.json :: fine-tuning jc2_qcd N=10000: (1 - macro_auc_qcd) of R42_Q1 over L188, run by run :: sha256 c6c4138ee5c86423
\newcommand{\MlstFtCmpJcTwoQcdMidOverFullEFourPr}{16\% more errors}  % experiments/FIGS/data/v2/index.json :: fine-tuning jc2_qcd N=10000: (1 - macro_auc_qcd) of R42_Q1 over L188, run by run :: sha256 c6c4138ee5c86423 :: prose twin: two significant figures, direction in words
\newcommand{\MlstFtCmpJcTwoQcdCoarseOverFullEFour}{1.258\,\ensuremath{\pm}\,0.025}  % experiments/FIGS/data/v2/index.json :: fine-tuning jc2_qcd N=10000: (1 - macro_auc_qcd) of R16_Q1 over L188, run by run :: sha256 c6c4138ee5c86423
\newcommand{\MlstFtCmpJcTwoQcdCoarseOverFullEFourPr}{26\% more errors}  % experiments/FIGS/data/v2/index.json :: fine-tuning jc2_qcd N=10000: (1 - macro_auc_qcd) of R16_Q1 over L188, run by run :: sha256 c6c4138ee5c86423 :: prose twin: two significant figures, direction in words
\newcommand{\MlstFtCmpJcTwoQcdResMassOverResEFour}{1.027 and 0.997}  % experiments/FIGS/data/v2/index.json :: fine-tuning jc2_qcd N=10000: (1 - macro_auc_qcd) of L162_MASS over L162, run by run :: sha256 c6c4138ee5c86423
\newcommand{\MlstFtCmpJcTwoQcdResMassOverResEFourPr}{equal within 3\%}  % experiments/FIGS/data/v2/index.json :: fine-tuning jc2_qcd N=10000: (1 - macro_auc_qcd) of L162_MASS over L162, run by run :: sha256 c6c4138ee5c86423 :: prose twin: two significant figures, direction in words
\newcommand{\MlstFtCmpJcTwoQcdResMassOverResEFourPrMax}{equal within 3\%}  % experiments/FIGS/data/v2/index.json :: fine-tuning jc2_qcd N=10000: (1 - macro_auc_qcd) of L162_MASS over L162, run by run :: sha256 c6c4138ee5c86423 :: prose twin of the run farther from 1
\newcommand{\MlstFtCmpJcTwoQcdCoarseMassOverCoarseEFour}{1.028 and 1.048}  % experiments/FIGS/data/v2/index.json :: fine-tuning jc2_qcd N=10000: (1 - macro_auc_qcd) of R16_Q1_MASS over R16_Q1, run by run :: sha256 c6c4138ee5c86423
\newcommand{\MlstFtCmpJcTwoQcdCoarseMassOverCoarseEFourPr}{equal within 5\%}  % experiments/FIGS/data/v2/index.json :: fine-tuning jc2_qcd N=10000: (1 - macro_auc_qcd) of R16_Q1_MASS over R16_Q1, run by run :: sha256 c6c4138ee5c86423 :: prose twin: two significant figures, direction in words
\newcommand{\MlstFtCmpJcTwoQcdCoarseMassOverCoarseEFourPrMax}{equal within 5\%}  % experiments/FIGS/data/v2/index.json :: fine-tuning jc2_qcd N=10000: (1 - macro_auc_qcd) of R16_Q1_MASS over R16_Q1, run by run :: sha256 c6c4138ee5c86423 :: prose twin of the run farther from 1
\newcommand{\MlstFtCmpJcTwoQcdResOverFullEFive}{1.060\,\ensuremath{\pm}\,0.002}  % experiments/FIGS/data/v2/index.json :: fine-tuning jc2_qcd N=100000: (1 - macro_auc_qcd) of L162 over L188, run by run :: sha256 c6c4138ee5c86423
\newcommand{\MlstFtCmpJcTwoQcdResOverFullEFivePr}{6.0\% more errors}  % experiments/FIGS/data/v2/index.json :: fine-tuning jc2_qcd N=100000: (1 - macro_auc_qcd) of L162 over L188, run by run :: sha256 c6c4138ee5c86423 :: prose twin: two significant figures, direction in words
\newcommand{\MlstFtCmpJcTwoQcdMidOverFullEFive}{1.062\,\ensuremath{\pm}\,0.010}  % experiments/FIGS/data/v2/index.json :: fine-tuning jc2_qcd N=100000: (1 - macro_auc_qcd) of R42_Q1 over L188, run by run :: sha256 c6c4138ee5c86423
\newcommand{\MlstFtCmpJcTwoQcdMidOverFullEFivePr}{6.2\% more errors}  % experiments/FIGS/data/v2/index.json :: fine-tuning jc2_qcd N=100000: (1 - macro_auc_qcd) of R42_Q1 over L188, run by run :: sha256 c6c4138ee5c86423 :: prose twin: two significant figures, direction in words
\newcommand{\MlstFtCmpJcTwoQcdCoarseOverFullEFive}{1.137\,\ensuremath{\pm}\,0.001}  % experiments/FIGS/data/v2/index.json :: fine-tuning jc2_qcd N=100000: (1 - macro_auc_qcd) of R16_Q1 over L188, run by run :: sha256 c6c4138ee5c86423
\newcommand{\MlstFtCmpJcTwoQcdCoarseOverFullEFivePr}{14\% more errors}  % experiments/FIGS/data/v2/index.json :: fine-tuning jc2_qcd N=100000: (1 - macro_auc_qcd) of R16_Q1 over L188, run by run :: sha256 c6c4138ee5c86423 :: prose twin: two significant figures, direction in words
\newcommand{\MlstFtCmpJcTwoQcdResMassOverResEFive}{1.022 and 1.009}  % experiments/FIGS/data/v2/index.json :: fine-tuning jc2_qcd N=100000: (1 - macro_auc_qcd) of L162_MASS over L162, run by run :: sha256 c6c4138ee5c86423
\newcommand{\MlstFtCmpJcTwoQcdResMassOverResEFivePr}{equal within 3\%}  % experiments/FIGS/data/v2/index.json :: fine-tuning jc2_qcd N=100000: (1 - macro_auc_qcd) of L162_MASS over L162, run by run :: sha256 c6c4138ee5c86423 :: prose twin: two significant figures, direction in words
\newcommand{\MlstFtCmpJcTwoQcdResMassOverResEFivePrMax}{equal within 3\%}  % experiments/FIGS/data/v2/index.json :: fine-tuning jc2_qcd N=100000: (1 - macro_auc_qcd) of L162_MASS over L162, run by run :: sha256 c6c4138ee5c86423 :: prose twin of the run farther from 1
\newcommand{\MlstFtCmpJcTwoQcdCoarseMassOverCoarseEFive}{1.028 and 1.027}  % experiments/FIGS/data/v2/index.json :: fine-tuning jc2_qcd N=100000: (1 - macro_auc_qcd) of R16_Q1_MASS over R16_Q1, run by run :: sha256 c6c4138ee5c86423
\newcommand{\MlstFtCmpJcTwoQcdCoarseMassOverCoarseEFivePr}{at least 2.7\% more errors}  % experiments/FIGS/data/v2/index.json :: fine-tuning jc2_qcd N=100000: (1 - macro_auc_qcd) of R16_Q1_MASS over R16_Q1, run by run :: sha256 c6c4138ee5c86423 :: prose twin: two significant figures, direction in words
\newcommand{\MlstFtCmpJcTwoQcdCoarseMassOverCoarseEFivePrMax}{up to 2.8\% more errors}  % experiments/FIGS/data/v2/index.json :: fine-tuning jc2_qcd N=100000: (1 - macro_auc_qcd) of R16_Q1_MASS over R16_Q1, run by run :: sha256 c6c4138ee5c86423 :: prose twin of the run farther from 1
\newcommand{\MlstFtCmpJcTwoQcdResOverFullESix}{1.035\,\ensuremath{\pm}\,0.010}  % experiments/FIGS/data/v2/index.json :: fine-tuning jc2_qcd N=1000000: (1 - macro_auc_qcd) of L162 over L188, run by run :: sha256 c6c4138ee5c86423
\newcommand{\MlstFtCmpJcTwoQcdResOverFullESixPr}{3.5\% more errors}  % experiments/FIGS/data/v2/index.json :: fine-tuning jc2_qcd N=1000000: (1 - macro_auc_qcd) of L162 over L188, run by run :: sha256 c6c4138ee5c86423 :: prose twin: two significant figures, direction in words
\newcommand{\MlstFtCmpJcTwoQcdMidOverFullESix}{1.050\,\ensuremath{\pm}\,0.004}  % experiments/FIGS/data/v2/index.json :: fine-tuning jc2_qcd N=1000000: (1 - macro_auc_qcd) of R42_Q1 over L188, run by run :: sha256 c6c4138ee5c86423
\newcommand{\MlstFtCmpJcTwoQcdMidOverFullESixPr}{5.0\% more errors}  % experiments/FIGS/data/v2/index.json :: fine-tuning jc2_qcd N=1000000: (1 - macro_auc_qcd) of R42_Q1 over L188, run by run :: sha256 c6c4138ee5c86423 :: prose twin: two significant figures, direction in words
\newcommand{\MlstFtCmpJcTwoQcdCoarseOverFullESix}{1.122\,\ensuremath{\pm}\,0.002}  % experiments/FIGS/data/v2/index.json :: fine-tuning jc2_qcd N=1000000: (1 - macro_auc_qcd) of R16_Q1 over L188, run by run :: sha256 c6c4138ee5c86423
\newcommand{\MlstFtCmpJcTwoQcdCoarseOverFullESixPr}{12\% more errors}  % experiments/FIGS/data/v2/index.json :: fine-tuning jc2_qcd N=1000000: (1 - macro_auc_qcd) of R16_Q1 over L188, run by run :: sha256 c6c4138ee5c86423 :: prose twin: two significant figures, direction in words
\newcommand{\MlstFtCmpJcTwoQcdResMassOverResESix}{1.011 and 1.025}  % experiments/FIGS/data/v2/index.json :: fine-tuning jc2_qcd N=1000000: (1 - macro_auc_qcd) of L162_MASS over L162, run by run :: sha256 c6c4138ee5c86423
\newcommand{\MlstFtCmpJcTwoQcdResMassOverResESixPr}{equal within 3\%}  % experiments/FIGS/data/v2/index.json :: fine-tuning jc2_qcd N=1000000: (1 - macro_auc_qcd) of L162_MASS over L162, run by run :: sha256 c6c4138ee5c86423 :: prose twin: two significant figures, direction in words
\newcommand{\MlstFtCmpJcTwoQcdResMassOverResESixPrMax}{equal within 3\%}  % experiments/FIGS/data/v2/index.json :: fine-tuning jc2_qcd N=1000000: (1 - macro_auc_qcd) of L162_MASS over L162, run by run :: sha256 c6c4138ee5c86423 :: prose twin of the run farther from 1
\newcommand{\MlstFtCmpJcTwoQcdCoarseMassOverCoarseESix}{1.015 and 1.015}  % experiments/FIGS/data/v2/index.json :: fine-tuning jc2_qcd N=1000000: (1 - macro_auc_qcd) of R16_Q1_MASS over R16_Q1, run by run :: sha256 c6c4138ee5c86423
\newcommand{\MlstFtCmpJcTwoQcdCoarseMassOverCoarseESixPr}{at least 1.5\% more errors}  % experiments/FIGS/data/v2/index.json :: fine-tuning jc2_qcd N=1000000: (1 - macro_auc_qcd) of R16_Q1_MASS over R16_Q1, run by run :: sha256 c6c4138ee5c86423 :: prose twin: two significant figures, direction in words
\newcommand{\MlstFtCmpJcTwoQcdCoarseMassOverCoarseESixPrMax}{up to 1.5\% more errors}  % experiments/FIGS/data/v2/index.json :: fine-tuning jc2_qcd N=1000000: (1 - macro_auc_qcd) of R16_Q1_MASS over R16_Q1, run by run :: sha256 c6c4138ee5c86423 :: prose twin of the run farther from 1
\newcommand{\MlstFtJcTwoAccFullEFour}{0.3914\,\ensuremath{\pm}\,0.0011}  % experiments/FIGS/data/v2/index.json :: fine-tuning jc2 N=10000: accuracy, L188 :: sha256 c6c4138ee5c86423
\newcommand{\MlstFtJcTwoAccResEFour}{0.3934\,\ensuremath{\pm}\,0.0018}  % experiments/FIGS/data/v2/index.json :: fine-tuning jc2 N=10000: accuracy, L162 :: sha256 c6c4138ee5c86423
\newcommand{\MlstFtJcTwoAccMidEFour}{0.3563\,\ensuremath{\pm}\,0.0025}  % experiments/FIGS/data/v2/index.json :: fine-tuning jc2 N=10000: accuracy, R42_Q1 :: sha256 c6c4138ee5c86423
\newcommand{\MlstFtJcTwoAccCoarseEFour}{0.3219\,\ensuremath{\pm}\,0.0011}  % experiments/FIGS/data/v2/index.json :: fine-tuning jc2 N=10000: accuracy, R16_Q1 :: sha256 c6c4138ee5c86423
\newcommand{\MlstFtJcTwoAccResMassEFour}{0.381 and 0.392}  % experiments/FIGS/data/v2/index.json :: fine-tuning jc2 N=10000: accuracy, L162_MASS :: sha256 c6c4138ee5c86423
\newcommand{\MlstFtJcTwoAccCoarseMassEFour}{0.3156 and 0.3187}  % experiments/FIGS/data/v2/index.json :: fine-tuning jc2 N=10000: accuracy, R16_Q1_MASS :: sha256 c6c4138ee5c86423
\newcommand{\MlstFtJcTwoAccScratchEFour}{0.1492\,\ensuremath{\pm}\,0.0011}  % experiments/FIGS/data/v2/index.json :: fine-tuning jc2 N=10000: accuracy from scratch :: sha256 c6c4138ee5c86423
\newcommand{\MlstFtJcTwoAccFullMinusScratchEFour}{0.2423\,\ensuremath{\pm}\,0.0011}  % experiments/FIGS/data/v2/index.json :: fine-tuning jc2 N=10000: accuracy of each 188-class run minus the from-scratch mean (absolute) :: sha256 c6c4138ee5c86423
\newcommand{\MlstFtJcTwoAccFullEFive}{0.4164\,\ensuremath{\pm}\,0.0015}  % experiments/FIGS/data/v2/index.json :: fine-tuning jc2 N=100000: accuracy, L188 :: sha256 c6c4138ee5c86423
\newcommand{\MlstFtJcTwoAccResEFive}{0.4161\,\ensuremath{\pm}\,0.0015}  % experiments/FIGS/data/v2/index.json :: fine-tuning jc2 N=100000: accuracy, L162 :: sha256 c6c4138ee5c86423
\newcommand{\MlstFtJcTwoAccMidEFive}{0.4007\,\ensuremath{\pm}\,0.0013}  % experiments/FIGS/data/v2/index.json :: fine-tuning jc2 N=100000: accuracy, R42_Q1 :: sha256 c6c4138ee5c86423
\newcommand{\MlstFtJcTwoAccCoarseEFive}{0.3833\,\ensuremath{\pm}\,0.0020}  % experiments/FIGS/data/v2/index.json :: fine-tuning jc2 N=100000: accuracy, R16_Q1 :: sha256 c6c4138ee5c86423
\newcommand{\MlstFtJcTwoAccResMassEFive}{0.4093 and 0.4115}  % experiments/FIGS/data/v2/index.json :: fine-tuning jc2 N=100000: accuracy, L162_MASS :: sha256 c6c4138ee5c86423
\newcommand{\MlstFtJcTwoAccCoarseMassEFive}{0.3758 and 0.3781}  % experiments/FIGS/data/v2/index.json :: fine-tuning jc2 N=100000: accuracy, R16_Q1_MASS :: sha256 c6c4138ee5c86423
\newcommand{\MlstFtJcTwoAccScratchEFive}{0.292\,\ensuremath{\pm}\,0.004}  % experiments/FIGS/data/v2/index.json :: fine-tuning jc2 N=100000: accuracy from scratch :: sha256 c6c4138ee5c86423
\newcommand{\MlstFtJcTwoAccFullMinusScratchEFive}{0.1243\,\ensuremath{\pm}\,0.0015}  % experiments/FIGS/data/v2/index.json :: fine-tuning jc2 N=100000: accuracy of each 188-class run minus the from-scratch mean (absolute) :: sha256 c6c4138ee5c86423
\newcommand{\MlstFtJcTwoAccFullESix}{0.4263\,\ensuremath{\pm}\,0.0005}  % experiments/FIGS/data/v2/index.json :: fine-tuning jc2 N=1000000: accuracy, L188 :: sha256 c6c4138ee5c86423
\newcommand{\MlstFtJcTwoAccResESix}{0.4249\,\ensuremath{\pm}\,0.0021}  % experiments/FIGS/data/v2/index.json :: fine-tuning jc2 N=1000000: accuracy, L162 :: sha256 c6c4138ee5c86423
\newcommand{\MlstFtJcTwoAccMidESix}{0.4166\,\ensuremath{\pm}\,0.0011}  % experiments/FIGS/data/v2/index.json :: fine-tuning jc2 N=1000000: accuracy, R42_Q1 :: sha256 c6c4138ee5c86423
\newcommand{\MlstFtJcTwoAccCoarseESix}{0.4054\,\ensuremath{\pm}\,0.0020}  % experiments/FIGS/data/v2/index.json :: fine-tuning jc2 N=1000000: accuracy, R16_Q1 :: sha256 c6c4138ee5c86423
\newcommand{\MlstFtJcTwoAccResMassESix}{0.4198 and 0.4213}  % experiments/FIGS/data/v2/index.json :: fine-tuning jc2 N=1000000: accuracy, L162_MASS :: sha256 c6c4138ee5c86423
\newcommand{\MlstFtJcTwoAccCoarseMassESix}{0.4012 and 0.4000}  % experiments/FIGS/data/v2/index.json :: fine-tuning jc2 N=1000000: accuracy, R16_Q1_MASS :: sha256 c6c4138ee5c86423
\newcommand{\MlstFtJcTwoAccScratchESix}{0.3677\,\ensuremath{\pm}\,0.0012}  % experiments/FIGS/data/v2/index.json :: fine-tuning jc2 N=1000000: accuracy from scratch :: sha256 c6c4138ee5c86423
\newcommand{\MlstFtJcTwoAccFullMinusScratchESix}{0.0585\,\ensuremath{\pm}\,0.0005}  % experiments/FIGS/data/v2/index.json :: fine-tuning jc2 N=1000000: accuracy of each 188-class run minus the from-scratch mean (absolute) :: sha256 c6c4138ee5c86423
\newcommand{\MlstFtJcTwoAucFull}{0.9797\,\ensuremath{\pm}\,0.0001}  % experiments/FIGS/data/v2/index.json :: fine-tuning class-balanced JetClass-II N=1000000: AUC averaged over the classes, Full :: sha256 c6c4138ee5c86423
\newcommand{\MlstFtJcTwoAucScratch}{0.9734\,\ensuremath{\pm}\,0.0003}  % experiments/FIGS/data/v2/index.json :: fine-tuning class-balanced JetClass-II N=1000000: AUC averaged over the classes, Scratch :: sha256 c6c4138ee5c86423
\newcommand{\MlstFtJcOneAccFullEFour}{0.8117\,\ensuremath{\pm}\,0.0031}  % experiments/FIGS/data/v2/index.json :: fine-tuning jc1 N=10000: accuracy, L188 :: sha256 c6c4138ee5c86423
\newcommand{\MlstFtJcOneAccResEFour}{0.8117\,\ensuremath{\pm}\,0.0013}  % experiments/FIGS/data/v2/index.json :: fine-tuning jc1 N=10000: accuracy, L162 :: sha256 c6c4138ee5c86423
\newcommand{\MlstFtJcOneAccMidEFour}{0.8086\,\ensuremath{\pm}\,0.0031}  % experiments/FIGS/data/v2/index.json :: fine-tuning jc1 N=10000: accuracy, R42_Q1 :: sha256 c6c4138ee5c86423
\newcommand{\MlstFtJcOneAccCoarseEFour}{0.7982\,\ensuremath{\pm}\,0.0023}  % experiments/FIGS/data/v2/index.json :: fine-tuning jc1 N=10000: accuracy, R16_Q1 :: sha256 c6c4138ee5c86423
\newcommand{\MlstFtJcOneAccResMassEFour}{0.8107 and 0.8098}  % experiments/FIGS/data/v2/index.json :: fine-tuning jc1 N=10000: accuracy, L162_MASS :: sha256 c6c4138ee5c86423
\newcommand{\MlstFtJcOneAccCoarseMassEFour}{0.7944 and 0.7990}  % experiments/FIGS/data/v2/index.json :: fine-tuning jc1 N=10000: accuracy, R16_Q1_MASS :: sha256 c6c4138ee5c86423
\newcommand{\MlstFtJcOneAccScratchEFour}{0.677\,\ensuremath{\pm}\,0.004}  % experiments/FIGS/data/v2/index.json :: fine-tuning jc1 N=10000: accuracy from scratch :: sha256 c6c4138ee5c86423
\newcommand{\MlstFtJcOneAccFullMinusScratchEFour}{0.1343\,\ensuremath{\pm}\,0.0031}  % experiments/FIGS/data/v2/index.json :: fine-tuning jc1 N=10000: accuracy of each 188-class run minus the from-scratch mean (absolute) :: sha256 c6c4138ee5c86423
\newcommand{\MlstFtJcOneAccFullEFive}{0.8343\,\ensuremath{\pm}\,0.0035}  % experiments/FIGS/data/v2/index.json :: fine-tuning jc1 N=100000: accuracy, L188 :: sha256 c6c4138ee5c86423
\newcommand{\MlstFtJcOneAccResEFive}{0.8328\,\ensuremath{\pm}\,0.0005}  % experiments/FIGS/data/v2/index.json :: fine-tuning jc1 N=100000: accuracy, L162 :: sha256 c6c4138ee5c86423
\newcommand{\MlstFtJcOneAccMidEFive}{0.8324\,\ensuremath{\pm}\,0.0004}  % experiments/FIGS/data/v2/index.json :: fine-tuning jc1 N=100000: accuracy, R42_Q1 :: sha256 c6c4138ee5c86423
\newcommand{\MlstFtJcOneAccCoarseEFive}{0.8288\,\ensuremath{\pm}\,0.0005}  % experiments/FIGS/data/v2/index.json :: fine-tuning jc1 N=100000: accuracy, R16_Q1 :: sha256 c6c4138ee5c86423
\newcommand{\MlstFtJcOneAccResMassEFive}{0.8326 and 0.8348}  % experiments/FIGS/data/v2/index.json :: fine-tuning jc1 N=100000: accuracy, L162_MASS :: sha256 c6c4138ee5c86423
\newcommand{\MlstFtJcOneAccCoarseMassEFive}{0.8253 and 0.8262}  % experiments/FIGS/data/v2/index.json :: fine-tuning jc1 N=100000: accuracy, R16_Q1_MASS :: sha256 c6c4138ee5c86423
\newcommand{\MlstFtJcOneAccScratchEFive}{0.777\,\ensuremath{\pm}\,0.005}  % experiments/FIGS/data/v2/index.json :: fine-tuning jc1 N=100000: accuracy from scratch :: sha256 c6c4138ee5c86423
\newcommand{\MlstFtJcOneAccFullMinusScratchEFive}{0.0571\,\ensuremath{\pm}\,0.0035}  % experiments/FIGS/data/v2/index.json :: fine-tuning jc1 N=100000: accuracy of each 188-class run minus the from-scratch mean (absolute) :: sha256 c6c4138ee5c86423
\newcommand{\MlstFtJcOneAccFullESix}{0.8468\,\ensuremath{\pm}\,0.0006}  % experiments/FIGS/data/v2/index.json :: fine-tuning jc1 N=1000000: accuracy, L188 :: sha256 c6c4138ee5c86423
\newcommand{\MlstFtJcOneAccResESix}{0.8475\,\ensuremath{\pm}\,0.0009}  % experiments/FIGS/data/v2/index.json :: fine-tuning jc1 N=1000000: accuracy, L162 :: sha256 c6c4138ee5c86423
\newcommand{\MlstFtJcOneAccMidESix}{0.8459\,\ensuremath{\pm}\,0.0002}  % experiments/FIGS/data/v2/index.json :: fine-tuning jc1 N=1000000: accuracy, R42_Q1 :: sha256 c6c4138ee5c86423
\newcommand{\MlstFtJcOneAccCoarseESix}{0.8422\,\ensuremath{\pm}\,0.0005}  % experiments/FIGS/data/v2/index.json :: fine-tuning jc1 N=1000000: accuracy, R16_Q1 :: sha256 c6c4138ee5c86423
\newcommand{\MlstFtJcOneAccResMassESix}{0.8459 and 0.8471}  % experiments/FIGS/data/v2/index.json :: fine-tuning jc1 N=1000000: accuracy, L162_MASS :: sha256 c6c4138ee5c86423
\newcommand{\MlstFtJcOneAccCoarseMassESix}{0.8400 and 0.8398}  % experiments/FIGS/data/v2/index.json :: fine-tuning jc1 N=1000000: accuracy, R16_Q1_MASS :: sha256 c6c4138ee5c86423
\newcommand{\MlstFtJcOneAccScratchESix}{0.8267\,\ensuremath{\pm}\,0.0011}  % experiments/FIGS/data/v2/index.json :: fine-tuning jc1 N=1000000: accuracy from scratch :: sha256 c6c4138ee5c86423
\newcommand{\MlstFtJcOneAccFullMinusScratchESix}{0.0201\,\ensuremath{\pm}\,0.0006}  % experiments/FIGS/data/v2/index.json :: fine-tuning jc1 N=1000000: accuracy of each 188-class run minus the from-scratch mean (absolute) :: sha256 c6c4138ee5c86423
\newcommand{\MlstFtTopRthirtyScratchEThree}{49\,\ensuremath{\pm}\,12}  % experiments/FIGS/data/v2/index.json :: fine-tuning top N=1000: R30 from scratch, every seed :: sha256 c6c4138ee5c86423
\newcommand{\MlstFtTopRthirtyFullEThree}{72\,\ensuremath{\pm}\,19}  % experiments/FIGS/data/v2/index.json :: fine-tuning top N=1000: R30, L188 :: sha256 c6c4138ee5c86423
\newcommand{\MlstFtTopRthirtyFullOverScratchEThree}{1.5\,\ensuremath{\pm}\,0.4}  % experiments/FIGS/data/v2/index.json :: fine-tuning top N=1000: R30 of each L188 run over the from-scratch mean, SD over runs and standard error of that mean in quadrature :: sha256 c6c4138ee5c86423
\newcommand{\MlstFtTopRthirtyFullOverScratchEThreePr}{equal within 50\%}  % experiments/FIGS/data/v2/index.json :: fine-tuning top N=1000: R30 of each L188 run over the from-scratch mean, SD over runs and standard error of that mean in quadrature :: sha256 c6c4138ee5c86423 :: prose twin: two significant figures, direction in words
\newcommand{\MlstFtTopRthirtyResEThree}{135\,\ensuremath{\pm}\,26}  % experiments/FIGS/data/v2/index.json :: fine-tuning top N=1000: R30, L162 :: sha256 c6c4138ee5c86423
\newcommand{\MlstFtTopRthirtyResOverScratchEThree}{2.8\,\ensuremath{\pm}\,0.7}  % experiments/FIGS/data/v2/index.json :: fine-tuning top N=1000: R30 of each L162 run over the from-scratch mean, SD over runs and standard error of that mean in quadrature :: sha256 c6c4138ee5c86423
\newcommand{\MlstFtTopRthirtyResOverScratchEThreePr}{2.8 times the rejection}  % experiments/FIGS/data/v2/index.json :: fine-tuning top N=1000: R30 of each L162 run over the from-scratch mean, SD over runs and standard error of that mean in quadrature :: sha256 c6c4138ee5c86423 :: prose twin: two significant figures, direction in words
\newcommand{\MlstFtTopRthirtyMidEThree}{104\,\ensuremath{\pm}\,20}  % experiments/FIGS/data/v2/index.json :: fine-tuning top N=1000: R30, R42_Q1 :: sha256 c6c4138ee5c86423
\newcommand{\MlstFtTopRthirtyMidOverScratchEThree}{2.1\,\ensuremath{\pm}\,0.5}  % experiments/FIGS/data/v2/index.json :: fine-tuning top N=1000: R30 of each R42_Q1 run over the from-scratch mean, SD over runs and standard error of that mean in quadrature :: sha256 c6c4138ee5c86423
\newcommand{\MlstFtTopRthirtyMidOverScratchEThreePr}{2.1 times the rejection}  % experiments/FIGS/data/v2/index.json :: fine-tuning top N=1000: R30 of each R42_Q1 run over the from-scratch mean, SD over runs and standard error of that mean in quadrature :: sha256 c6c4138ee5c86423 :: prose twin: two significant figures, direction in words
\newcommand{\MlstFtTopRthirtyCoarseEThree}{83\,\ensuremath{\pm}\,25}  % experiments/FIGS/data/v2/index.json :: fine-tuning top N=1000: R30, R16_Q1 :: sha256 c6c4138ee5c86423
\newcommand{\MlstFtTopRthirtyCoarseOverScratchEThree}{1.7\,\ensuremath{\pm}\,0.6}  % experiments/FIGS/data/v2/index.json :: fine-tuning top N=1000: R30 of each R16_Q1 run over the from-scratch mean, SD over runs and standard error of that mean in quadrature :: sha256 c6c4138ee5c86423
\newcommand{\MlstFtTopRthirtyCoarseOverScratchEThreePr}{equal within 80\%}  % experiments/FIGS/data/v2/index.json :: fine-tuning top N=1000: R30 of each R16_Q1 run over the from-scratch mean, SD over runs and standard error of that mean in quadrature :: sha256 c6c4138ee5c86423 :: prose twin: two significant figures, direction in words
\newcommand{\MlstFtTopRthirtyResMassEThree}{105.9 and 112.7}  % experiments/FIGS/data/v2/index.json :: fine-tuning top N=1000: R30, L162_MASS :: sha256 c6c4138ee5c86423
\newcommand{\MlstFtTopRthirtyResMassOverScratchEThree}{2.18 and 2.32}  % experiments/FIGS/data/v2/index.json :: fine-tuning top N=1000: R30 of each L162_MASS run over the from-scratch mean, SD over runs and standard error of that mean in quadrature :: sha256 c6c4138ee5c86423
\newcommand{\MlstFtTopRthirtyResMassOverScratchEThreePr}{at least 2.2 times the rejection}  % experiments/FIGS/data/v2/index.json :: fine-tuning top N=1000: R30 of each L162_MASS run over the from-scratch mean, SD over runs and standard error of that mean in quadrature :: sha256 c6c4138ee5c86423 :: prose twin: two significant figures, direction in words
\newcommand{\MlstFtTopRthirtyResMassOverScratchEThreePrMax}{up to 2.3 times the rejection}  % experiments/FIGS/data/v2/index.json :: fine-tuning top N=1000: R30 of each L162_MASS run over the from-scratch mean, SD over runs and standard error of that mean in quadrature :: sha256 c6c4138ee5c86423 :: prose twin of the run farther from 1
\newcommand{\MlstFtTopRthirtyCoarseMassEThree}{59.1 and 63.9}  % experiments/FIGS/data/v2/index.json :: fine-tuning top N=1000: R30, R16_Q1_MASS :: sha256 c6c4138ee5c86423
\newcommand{\MlstFtTopRthirtyCoarseMassOverScratchEThree}{1.216 and 1.314}  % experiments/FIGS/data/v2/index.json :: fine-tuning top N=1000: R30 of each R16_Q1_MASS run over the from-scratch mean, SD over runs and standard error of that mean in quadrature :: sha256 c6c4138ee5c86423
\newcommand{\MlstFtTopRthirtyCoarseMassOverScratchEThreePr}{at least 22\% more rejection}  % experiments/FIGS/data/v2/index.json :: fine-tuning top N=1000: R30 of each R16_Q1_MASS run over the from-scratch mean, SD over runs and standard error of that mean in quadrature :: sha256 c6c4138ee5c86423 :: prose twin: two significant figures, direction in words
\newcommand{\MlstFtTopRthirtyCoarseMassOverScratchEThreePrMax}{up to 31\% more rejection}  % experiments/FIGS/data/v2/index.json :: fine-tuning top N=1000: R30 of each R16_Q1_MASS run over the from-scratch mean, SD over runs and standard error of that mean in quadrature :: sha256 c6c4138ee5c86423 :: prose twin of the run farther from 1
\newcommand{\MlstFtTopRthirtyResOverFullEThree}{2.0\,\ensuremath{\pm}\,0.8}  % experiments/FIGS/data/v2/index.json :: fine-tuning top N=1000: R30 of L162 over L188, run by run :: sha256 c6c4138ee5c86423
\newcommand{\MlstFtTopRthirtyResOverFullEThreePr}{equal within 100\%}  % experiments/FIGS/data/v2/index.json :: fine-tuning top N=1000: R30 of L162 over L188, run by run :: sha256 c6c4138ee5c86423 :: prose twin: two significant figures, direction in words
\newcommand{\MlstFtTopRthirtyMidOverFullEThree}{1.46\,\ensuremath{\pm}\,0.12}  % experiments/FIGS/data/v2/index.json :: fine-tuning top N=1000: R30 of R42_Q1 over L188, run by run :: sha256 c6c4138ee5c86423
\newcommand{\MlstFtTopRthirtyMidOverFullEThreePr}{46\% more rejection}  % experiments/FIGS/data/v2/index.json :: fine-tuning top N=1000: R30 of R42_Q1 over L188, run by run :: sha256 c6c4138ee5c86423 :: prose twin: two significant figures, direction in words
\newcommand{\MlstFtTopRthirtyCoarseOverFullEThree}{1.3\,\ensuremath{\pm}\,0.7}  % experiments/FIGS/data/v2/index.json :: fine-tuning top N=1000: R30 of R16_Q1 over L188, run by run :: sha256 c6c4138ee5c86423
\newcommand{\MlstFtTopRthirtyCoarseOverFullEThreePr}{equal within 80\%}  % experiments/FIGS/data/v2/index.json :: fine-tuning top N=1000: R30 of R16_Q1 over L188, run by run :: sha256 c6c4138ee5c86423 :: prose twin: two significant figures, direction in words
\newcommand{\MlstFtTopRthirtyResMassOverResEThree}{0.73 and 1.07}  % experiments/FIGS/data/v2/index.json :: fine-tuning top N=1000: R30 of L162_MASS over L162, run by run :: sha256 c6c4138ee5c86423
\newcommand{\MlstFtTopRthirtyResMassOverResEThreePr}{equal within 30\%}  % experiments/FIGS/data/v2/index.json :: fine-tuning top N=1000: R30 of L162_MASS over L162, run by run :: sha256 c6c4138ee5c86423 :: prose twin: two significant figures, direction in words
\newcommand{\MlstFtTopRthirtyResMassOverResEThreePrMax}{equal within 30\%}  % experiments/FIGS/data/v2/index.json :: fine-tuning top N=1000: R30 of L162_MASS over L162, run by run :: sha256 c6c4138ee5c86423 :: prose twin of the run farther from 1
\newcommand{\MlstFtTopRthirtyCoarseMassOverCoarseEThree}{0.56 and 0.75}  % experiments/FIGS/data/v2/index.json :: fine-tuning top N=1000: R30 of R16_Q1_MASS over R16_Q1, run by run :: sha256 c6c4138ee5c86423
\newcommand{\MlstFtTopRthirtyCoarseMassOverCoarseEThreePr}{equal within 50\%}  % experiments/FIGS/data/v2/index.json :: fine-tuning top N=1000: R30 of R16_Q1_MASS over R16_Q1, run by run :: sha256 c6c4138ee5c86423 :: prose twin: two significant figures, direction in words
\newcommand{\MlstFtTopRthirtyCoarseMassOverCoarseEThreePrMax}{equal within 50\%}  % experiments/FIGS/data/v2/index.json :: fine-tuning top N=1000: R30 of R16_Q1_MASS over R16_Q1, run by run :: sha256 c6c4138ee5c86423 :: prose twin of the run farther from 1
\newcommand{\MlstFtTopRthirtyRandOverFullEThree}{1.3\,\ensuremath{\pm}\,0.5}  % experiments/FIGS/data/v2/index.json :: fine-tuning top N=1000: R30 of each random partition over the L188 run of its seed :: sha256 c6c4138ee5c86423
\newcommand{\MlstFtTopRthirtyRandOverFullEThreePr}{equal within 50\%}  % experiments/FIGS/data/v2/index.json :: fine-tuning top N=1000: R30 of each random partition over the L188 run of its seed :: sha256 c6c4138ee5c86423 :: prose twin: two significant figures, direction in words
\newcommand{\MlstFtTopRthirtyRandOverCoarseEThree}{0.87\,\ensuremath{\pm}\,0.12}  % experiments/FIGS/data/v2/index.json :: fine-tuning top N=1000: R30 of each random partition over the R16_Q1 run of its seed :: sha256 c6c4138ee5c86423
\newcommand{\MlstFtTopRthirtyRandOverCoarseEThreePr}{equal within 20\%}  % experiments/FIGS/data/v2/index.json :: fine-tuning top N=1000: R30 of each random partition over the R16_Q1 run of its seed :: sha256 c6c4138ee5c86423 :: prose twin: two significant figures, direction in words
\newcommand{\MlstFtTopRthirtyScratchEFour}{340\,\ensuremath{\pm}\,50}  % experiments/FIGS/data/v2/index.json :: fine-tuning top N=10000: R30 from scratch, every seed :: sha256 c6c4138ee5c86423
\newcommand{\MlstFtTopRthirtyFullEFour}{420\,\ensuremath{\pm}\,9}  % experiments/FIGS/data/v2/index.json :: fine-tuning top N=10000: R30, L188 :: sha256 c6c4138ee5c86423
\newcommand{\MlstFtTopRthirtyFullOverScratchEFour}{1.24\,\ensuremath{\pm}\,0.11}  % experiments/FIGS/data/v2/index.json :: fine-tuning top N=10000: R30 of each L188 run over the from-scratch mean, SD over runs and standard error of that mean in quadrature :: sha256 c6c4138ee5c86423
\newcommand{\MlstFtTopRthirtyFullOverScratchEFourPr}{24\% more rejection}  % experiments/FIGS/data/v2/index.json :: fine-tuning top N=10000: R30 of each L188 run over the from-scratch mean, SD over runs and standard error of that mean in quadrature :: sha256 c6c4138ee5c86423 :: prose twin: two significant figures, direction in words
\newcommand{\MlstFtTopRthirtyResEFour}{430\,\ensuremath{\pm}\,60}  % experiments/FIGS/data/v2/index.json :: fine-tuning top N=10000: R30, L162 :: sha256 c6c4138ee5c86423
\newcommand{\MlstFtTopRthirtyResOverScratchEFour}{1.27\,\ensuremath{\pm}\,0.21}  % experiments/FIGS/data/v2/index.json :: fine-tuning top N=10000: R30 of each L162 run over the from-scratch mean, SD over runs and standard error of that mean in quadrature :: sha256 c6c4138ee5c86423
\newcommand{\MlstFtTopRthirtyResOverScratchEFourPr}{equal within 30\%}  % experiments/FIGS/data/v2/index.json :: fine-tuning top N=10000: R30 of each L162 run over the from-scratch mean, SD over runs and standard error of that mean in quadrature :: sha256 c6c4138ee5c86423 :: prose twin: two significant figures, direction in words
\newcommand{\MlstFtTopRthirtyMidEFour}{455\,\ensuremath{\pm}\,30}  % experiments/FIGS/data/v2/index.json :: fine-tuning top N=10000: R30, R42_Q1 :: sha256 c6c4138ee5c86423
\newcommand{\MlstFtTopRthirtyMidOverScratchEFour}{1.34\,\ensuremath{\pm}\,0.14}  % experiments/FIGS/data/v2/index.json :: fine-tuning top N=10000: R30 of each R42_Q1 run over the from-scratch mean, SD over runs and standard error of that mean in quadrature :: sha256 c6c4138ee5c86423
\newcommand{\MlstFtTopRthirtyMidOverScratchEFourPr}{34\% more rejection}  % experiments/FIGS/data/v2/index.json :: fine-tuning top N=10000: R30 of each R42_Q1 run over the from-scratch mean, SD over runs and standard error of that mean in quadrature :: sha256 c6c4138ee5c86423 :: prose twin: two significant figures, direction in words
\newcommand{\MlstFtTopRthirtyCoarseEFour}{350\,\ensuremath{\pm}\,40}  % experiments/FIGS/data/v2/index.json :: fine-tuning top N=10000: R30, R16_Q1 :: sha256 c6c4138ee5c86423
\newcommand{\MlstFtTopRthirtyCoarseOverScratchEFour}{1.04\,\ensuremath{\pm}\,0.14}  % experiments/FIGS/data/v2/index.json :: fine-tuning top N=10000: R30 of each R16_Q1 run over the from-scratch mean, SD over runs and standard error of that mean in quadrature :: sha256 c6c4138ee5c86423
\newcommand{\MlstFtTopRthirtyCoarseOverScratchEFourPr}{equal within 20\%}  % experiments/FIGS/data/v2/index.json :: fine-tuning top N=10000: R30 of each R16_Q1 run over the from-scratch mean, SD over runs and standard error of that mean in quadrature :: sha256 c6c4138ee5c86423 :: prose twin: two significant figures, direction in words
\newcommand{\MlstFtTopRthirtyResMassEFour}{373.2 and 376.7}  % experiments/FIGS/data/v2/index.json :: fine-tuning top N=10000: R30, L162_MASS :: sha256 c6c4138ee5c86423
\newcommand{\MlstFtTopRthirtyResMassOverScratchEFour}{1.101 and 1.111}  % experiments/FIGS/data/v2/index.json :: fine-tuning top N=10000: R30 of each L162_MASS run over the from-scratch mean, SD over runs and standard error of that mean in quadrature :: sha256 c6c4138ee5c86423
\newcommand{\MlstFtTopRthirtyResMassOverScratchEFourPr}{at least 10\% more rejection}  % experiments/FIGS/data/v2/index.json :: fine-tuning top N=10000: R30 of each L162_MASS run over the from-scratch mean, SD over runs and standard error of that mean in quadrature :: sha256 c6c4138ee5c86423 :: prose twin: two significant figures, direction in words
\newcommand{\MlstFtTopRthirtyResMassOverScratchEFourPrMax}{up to 11\% more rejection}  % experiments/FIGS/data/v2/index.json :: fine-tuning top N=10000: R30 of each L162_MASS run over the from-scratch mean, SD over runs and standard error of that mean in quadrature :: sha256 c6c4138ee5c86423 :: prose twin of the run farther from 1
\newcommand{\MlstFtTopRthirtyCoarseMassEFour}{365.8 and 359.9}  % experiments/FIGS/data/v2/index.json :: fine-tuning top N=10000: R30, R16_Q1_MASS :: sha256 c6c4138ee5c86423
\newcommand{\MlstFtTopRthirtyCoarseMassOverScratchEFour}{1.079 and 1.062}  % experiments/FIGS/data/v2/index.json :: fine-tuning top N=10000: R30 of each R16_Q1_MASS run over the from-scratch mean, SD over runs and standard error of that mean in quadrature :: sha256 c6c4138ee5c86423
\newcommand{\MlstFtTopRthirtyCoarseMassOverScratchEFourPr}{at least 6.2\% more rejection}  % experiments/FIGS/data/v2/index.json :: fine-tuning top N=10000: R30 of each R16_Q1_MASS run over the from-scratch mean, SD over runs and standard error of that mean in quadrature :: sha256 c6c4138ee5c86423 :: prose twin: two significant figures, direction in words
\newcommand{\MlstFtTopRthirtyCoarseMassOverScratchEFourPrMax}{up to 7.9\% more rejection}  % experiments/FIGS/data/v2/index.json :: fine-tuning top N=10000: R30 of each R16_Q1_MASS run over the from-scratch mean, SD over runs and standard error of that mean in quadrature :: sha256 c6c4138ee5c86423 :: prose twin of the run farther from 1
\newcommand{\MlstFtTopRthirtyResOverFullEFour}{1.02\,\ensuremath{\pm}\,0.12}  % experiments/FIGS/data/v2/index.json :: fine-tuning top N=10000: R30 of L162 over L188, run by run :: sha256 c6c4138ee5c86423
\newcommand{\MlstFtTopRthirtyResOverFullEFourPr}{equal within 20\%}  % experiments/FIGS/data/v2/index.json :: fine-tuning top N=10000: R30 of L162 over L188, run by run :: sha256 c6c4138ee5c86423 :: prose twin: two significant figures, direction in words
\newcommand{\MlstFtTopRthirtyMidOverFullEFour}{1.08\,\ensuremath{\pm}\,0.08}  % experiments/FIGS/data/v2/index.json :: fine-tuning top N=10000: R30 of R42_Q1 over L188, run by run :: sha256 c6c4138ee5c86423
\newcommand{\MlstFtTopRthirtyMidOverFullEFourPr}{equal within 9\%}  % experiments/FIGS/data/v2/index.json :: fine-tuning top N=10000: R30 of R42_Q1 over L188, run by run :: sha256 c6c4138ee5c86423 :: prose twin: two significant figures, direction in words
\newcommand{\MlstFtTopRthirtyCoarseOverFullEFour}{0.84\,\ensuremath{\pm}\,0.07}  % experiments/FIGS/data/v2/index.json :: fine-tuning top N=10000: R30 of R16_Q1 over L188, run by run :: sha256 c6c4138ee5c86423
\newcommand{\MlstFtTopRthirtyCoarseOverFullEFourPr}{16\% less rejection}  % experiments/FIGS/data/v2/index.json :: fine-tuning top N=10000: R30 of R16_Q1 over L188, run by run :: sha256 c6c4138ee5c86423 :: prose twin: two significant figures, direction in words
\newcommand{\MlstFtTopRthirtyResMassOverResEFour}{1.00 and 0.76}  % experiments/FIGS/data/v2/index.json :: fine-tuning top N=10000: R30 of L162_MASS over L162, run by run :: sha256 c6c4138ee5c86423
\newcommand{\MlstFtTopRthirtyResMassOverResEFourPr}{equal within 30\%}  % experiments/FIGS/data/v2/index.json :: fine-tuning top N=10000: R30 of L162_MASS over L162, run by run :: sha256 c6c4138ee5c86423 :: prose twin: two significant figures, direction in words
\newcommand{\MlstFtTopRthirtyResMassOverResEFourPrMax}{equal within 30\%}  % experiments/FIGS/data/v2/index.json :: fine-tuning top N=10000: R30 of L162_MASS over L162, run by run :: sha256 c6c4138ee5c86423 :: prose twin of the run farther from 1
\newcommand{\MlstFtTopRthirtyCoarseMassOverCoarseEFour}{1.13 and 0.92}  % experiments/FIGS/data/v2/index.json :: fine-tuning top N=10000: R30 of R16_Q1_MASS over R16_Q1, run by run :: sha256 c6c4138ee5c86423
\newcommand{\MlstFtTopRthirtyCoarseMassOverCoarseEFourPr}{equal within 20\%}  % experiments/FIGS/data/v2/index.json :: fine-tuning top N=10000: R30 of R16_Q1_MASS over R16_Q1, run by run :: sha256 c6c4138ee5c86423 :: prose twin: two significant figures, direction in words
\newcommand{\MlstFtTopRthirtyCoarseMassOverCoarseEFourPrMax}{equal within 20\%}  % experiments/FIGS/data/v2/index.json :: fine-tuning top N=10000: R30 of R16_Q1_MASS over R16_Q1, run by run :: sha256 c6c4138ee5c86423 :: prose twin of the run farther from 1
\newcommand{\MlstFtTopRthirtyRandOverFullEFour}{0.83\,\ensuremath{\pm}\,0.10}  % experiments/FIGS/data/v2/index.json :: fine-tuning top N=10000: R30 of each random partition over the L188 run of its seed :: sha256 c6c4138ee5c86423
\newcommand{\MlstFtTopRthirtyRandOverFullEFourPr}{equal within 20\%}  % experiments/FIGS/data/v2/index.json :: fine-tuning top N=10000: R30 of each random partition over the L188 run of its seed :: sha256 c6c4138ee5c86423 :: prose twin: two significant figures, direction in words
\newcommand{\MlstFtTopRthirtyRandOverCoarseEFour}{0.97\,\ensuremath{\pm}\,0.15}  % experiments/FIGS/data/v2/index.json :: fine-tuning top N=10000: R30 of each random partition over the R16_Q1 run of its seed :: sha256 c6c4138ee5c86423
\newcommand{\MlstFtTopRthirtyRandOverCoarseEFourPr}{equal within 20\%}  % experiments/FIGS/data/v2/index.json :: fine-tuning top N=10000: R30 of each random partition over the R16_Q1 run of its seed :: sha256 c6c4138ee5c86423 :: prose twin: two significant figures, direction in words
\newcommand{\MlstFtTopRthirtyScratchEFive}{560\,\ensuremath{\pm}\,90}  % experiments/FIGS/data/v2/index.json :: fine-tuning top N=100000: R30 from scratch, every seed :: sha256 c6c4138ee5c86423
\newcommand{\MlstFtTopRthirtyFullEFive}{850\,\ensuremath{\pm}\,90}  % experiments/FIGS/data/v2/index.json :: fine-tuning top N=100000: R30, L188 :: sha256 c6c4138ee5c86423
\newcommand{\MlstFtTopRthirtyFullOverScratchEFive}{1.50\,\ensuremath{\pm}\,0.21}  % experiments/FIGS/data/v2/index.json :: fine-tuning top N=100000: R30 of each L188 run over the from-scratch mean, SD over runs and standard error of that mean in quadrature :: sha256 c6c4138ee5c86423
\newcommand{\MlstFtTopRthirtyFullOverScratchEFivePr}{50\% more rejection}  % experiments/FIGS/data/v2/index.json :: fine-tuning top N=100000: R30 of each L188 run over the from-scratch mean, SD over runs and standard error of that mean in quadrature :: sha256 c6c4138ee5c86423 :: prose twin: two significant figures, direction in words
\newcommand{\MlstFtTopRthirtyResEFive}{810\,\ensuremath{\pm}\,60}  % experiments/FIGS/data/v2/index.json :: fine-tuning top N=100000: R30, L162 :: sha256 c6c4138ee5c86423
\newcommand{\MlstFtTopRthirtyResOverScratchEFive}{1.44\,\ensuremath{\pm}\,0.18}  % experiments/FIGS/data/v2/index.json :: fine-tuning top N=100000: R30 of each L162 run over the from-scratch mean, SD over runs and standard error of that mean in quadrature :: sha256 c6c4138ee5c86423
\newcommand{\MlstFtTopRthirtyResOverScratchEFivePr}{44\% more rejection}  % experiments/FIGS/data/v2/index.json :: fine-tuning top N=100000: R30 of each L162 run over the from-scratch mean, SD over runs and standard error of that mean in quadrature :: sha256 c6c4138ee5c86423 :: prose twin: two significant figures, direction in words
\newcommand{\MlstFtTopRthirtyMidEFive}{850\,\ensuremath{\pm}\,80}  % experiments/FIGS/data/v2/index.json :: fine-tuning top N=100000: R30, R42_Q1 :: sha256 c6c4138ee5c86423
\newcommand{\MlstFtTopRthirtyMidOverScratchEFive}{1.51\,\ensuremath{\pm}\,0.20}  % experiments/FIGS/data/v2/index.json :: fine-tuning top N=100000: R30 of each R42_Q1 run over the from-scratch mean, SD over runs and standard error of that mean in quadrature :: sha256 c6c4138ee5c86423
\newcommand{\MlstFtTopRthirtyMidOverScratchEFivePr}{51\% more rejection}  % experiments/FIGS/data/v2/index.json :: fine-tuning top N=100000: R30 of each R42_Q1 run over the from-scratch mean, SD over runs and standard error of that mean in quadrature :: sha256 c6c4138ee5c86423 :: prose twin: two significant figures, direction in words
\newcommand{\MlstFtTopRthirtyCoarseEFive}{830\,\ensuremath{\pm}\,60}  % experiments/FIGS/data/v2/index.json :: fine-tuning top N=100000: R30, R16_Q1 :: sha256 c6c4138ee5c86423
\newcommand{\MlstFtTopRthirtyCoarseOverScratchEFive}{1.47\,\ensuremath{\pm}\,0.17}  % experiments/FIGS/data/v2/index.json :: fine-tuning top N=100000: R30 of each R16_Q1 run over the from-scratch mean, SD over runs and standard error of that mean in quadrature :: sha256 c6c4138ee5c86423
\newcommand{\MlstFtTopRthirtyCoarseOverScratchEFivePr}{47\% more rejection}  % experiments/FIGS/data/v2/index.json :: fine-tuning top N=100000: R30 of each R16_Q1 run over the from-scratch mean, SD over runs and standard error of that mean in quadrature :: sha256 c6c4138ee5c86423 :: prose twin: two significant figures, direction in words
\newcommand{\MlstFtTopRthirtyResMassEFive}{848 and 848}  % experiments/FIGS/data/v2/index.json :: fine-tuning top N=100000: R30, L162_MASS :: sha256 c6c4138ee5c86423
\newcommand{\MlstFtTopRthirtyResMassOverScratchEFive}{1.50 and 1.50}  % experiments/FIGS/data/v2/index.json :: fine-tuning top N=100000: R30 of each L162_MASS run over the from-scratch mean, SD over runs and standard error of that mean in quadrature :: sha256 c6c4138ee5c86423
\newcommand{\MlstFtTopRthirtyResMassOverScratchEFivePr}{at least 50\% more rejection}  % experiments/FIGS/data/v2/index.json :: fine-tuning top N=100000: R30 of each L162_MASS run over the from-scratch mean, SD over runs and standard error of that mean in quadrature :: sha256 c6c4138ee5c86423 :: prose twin: two significant figures, direction in words
\newcommand{\MlstFtTopRthirtyResMassOverScratchEFivePrMax}{up to 50\% more rejection}  % experiments/FIGS/data/v2/index.json :: fine-tuning top N=100000: R30 of each L162_MASS run over the from-scratch mean, SD over runs and standard error of that mean in quadrature :: sha256 c6c4138ee5c86423 :: prose twin of the run farther from 1
\newcommand{\MlstFtTopRthirtyCoarseMassEFive}{856 and 774}  % experiments/FIGS/data/v2/index.json :: fine-tuning top N=100000: R30, R16_Q1_MASS :: sha256 c6c4138ee5c86423
\newcommand{\MlstFtTopRthirtyCoarseMassOverScratchEFive}{1.51 and 1.37}  % experiments/FIGS/data/v2/index.json :: fine-tuning top N=100000: R30 of each R16_Q1_MASS run over the from-scratch mean, SD over runs and standard error of that mean in quadrature :: sha256 c6c4138ee5c86423
\newcommand{\MlstFtTopRthirtyCoarseMassOverScratchEFivePr}{at least 37\% more rejection}  % experiments/FIGS/data/v2/index.json :: fine-tuning top N=100000: R30 of each R16_Q1_MASS run over the from-scratch mean, SD over runs and standard error of that mean in quadrature :: sha256 c6c4138ee5c86423 :: prose twin: two significant figures, direction in words
\newcommand{\MlstFtTopRthirtyCoarseMassOverScratchEFivePrMax}{up to 51\% more rejection}  % experiments/FIGS/data/v2/index.json :: fine-tuning top N=100000: R30 of each R16_Q1_MASS run over the from-scratch mean, SD over runs and standard error of that mean in quadrature :: sha256 c6c4138ee5c86423 :: prose twin of the run farther from 1
\newcommand{\MlstFtTopRthirtyResOverFullEFive}{0.96\,\ensuremath{\pm}\,0.05}  % experiments/FIGS/data/v2/index.json :: fine-tuning top N=100000: R30 of L162 over L188, run by run :: sha256 c6c4138ee5c86423
\newcommand{\MlstFtTopRthirtyResOverFullEFivePr}{equal within 6\%}  % experiments/FIGS/data/v2/index.json :: fine-tuning top N=100000: R30 of L162 over L188, run by run :: sha256 c6c4138ee5c86423 :: prose twin: two significant figures, direction in words
\newcommand{\MlstFtTopRthirtyMidOverFullEFive}{1.02\,\ensuremath{\pm}\,0.18}  % experiments/FIGS/data/v2/index.json :: fine-tuning top N=100000: R30 of R42_Q1 over L188, run by run :: sha256 c6c4138ee5c86423
\newcommand{\MlstFtTopRthirtyMidOverFullEFivePr}{equal within 20\%}  % experiments/FIGS/data/v2/index.json :: fine-tuning top N=100000: R30 of R42_Q1 over L188, run by run :: sha256 c6c4138ee5c86423 :: prose twin: two significant figures, direction in words
\newcommand{\MlstFtTopRthirtyCoarseOverFullEFive}{0.99\,\ensuremath{\pm}\,0.17}  % experiments/FIGS/data/v2/index.json :: fine-tuning top N=100000: R30 of R16_Q1 over L188, run by run :: sha256 c6c4138ee5c86423
\newcommand{\MlstFtTopRthirtyCoarseOverFullEFivePr}{equal within 20\%}  % experiments/FIGS/data/v2/index.json :: fine-tuning top N=100000: R30 of R16_Q1 over L188, run by run :: sha256 c6c4138ee5c86423 :: prose twin: two significant figures, direction in words
\newcommand{\MlstFtTopRthirtyResMassOverResEFive}{1.00 and 1.15}  % experiments/FIGS/data/v2/index.json :: fine-tuning top N=100000: R30 of L162_MASS over L162, run by run :: sha256 c6c4138ee5c86423
\newcommand{\MlstFtTopRthirtyResMassOverResEFivePr}{equal within 20\%}  % experiments/FIGS/data/v2/index.json :: fine-tuning top N=100000: R30 of L162_MASS over L162, run by run :: sha256 c6c4138ee5c86423 :: prose twin: two significant figures, direction in words
\newcommand{\MlstFtTopRthirtyResMassOverResEFivePrMax}{equal within 20\%}  % experiments/FIGS/data/v2/index.json :: fine-tuning top N=100000: R30 of L162_MASS over L162, run by run :: sha256 c6c4138ee5c86423 :: prose twin of the run farther from 1
\newcommand{\MlstFtTopRthirtyCoarseMassOverCoarseEFive}{1.08 and 0.87}  % experiments/FIGS/data/v2/index.json :: fine-tuning top N=100000: R30 of R16_Q1_MASS over R16_Q1, run by run :: sha256 c6c4138ee5c86423
\newcommand{\MlstFtTopRthirtyCoarseMassOverCoarseEFivePr}{equal within 20\%}  % experiments/FIGS/data/v2/index.json :: fine-tuning top N=100000: R30 of R16_Q1_MASS over R16_Q1, run by run :: sha256 c6c4138ee5c86423 :: prose twin: two significant figures, direction in words
\newcommand{\MlstFtTopRthirtyCoarseMassOverCoarseEFivePrMax}{equal within 20\%}  % experiments/FIGS/data/v2/index.json :: fine-tuning top N=100000: R30 of R16_Q1_MASS over R16_Q1, run by run :: sha256 c6c4138ee5c86423 :: prose twin of the run farther from 1
\newcommand{\MlstFtTopRthirtyRandOverFullEFive}{1.01\,\ensuremath{\pm}\,0.22}  % experiments/FIGS/data/v2/index.json :: fine-tuning top N=100000: R30 of each random partition over the L188 run of its seed :: sha256 c6c4138ee5c86423
\newcommand{\MlstFtTopRthirtyRandOverFullEFivePr}{equal within 30\%}  % experiments/FIGS/data/v2/index.json :: fine-tuning top N=100000: R30 of each random partition over the L188 run of its seed :: sha256 c6c4138ee5c86423 :: prose twin: two significant figures, direction in words
\newcommand{\MlstFtTopRthirtyRandOverCoarseEFive}{0.97\,\ensuremath{\pm}\,0.10}  % experiments/FIGS/data/v2/index.json :: fine-tuning top N=100000: R30 of each random partition over the R16_Q1 run of its seed :: sha256 c6c4138ee5c86423
\newcommand{\MlstFtTopRthirtyRandOverCoarseEFivePr}{equal within 10\%}  % experiments/FIGS/data/v2/index.json :: fine-tuning top N=100000: R30 of each random partition over the R16_Q1 run of its seed :: sha256 c6c4138ee5c86423 :: prose twin: two significant figures, direction in words
\newcommand{\MlstFtTopRthirtyScratch}{1{,}412\,\ensuremath{\pm}\,20}  % experiments/FIGS/data/v2/index.json :: fine-tuning top N=1200000: R30 from scratch, every seed :: sha256 c6c4138ee5c86423
\newcommand{\MlstFtTopRthirtyFull}{1{,}800\,\ensuremath{\pm}\,60}  % experiments/FIGS/data/v2/index.json :: fine-tuning top N=1200000: R30, L188 :: sha256 c6c4138ee5c86423
\newcommand{\MlstFtTopRthirtyFullOverScratch}{1.27\,\ensuremath{\pm}\,0.04}  % experiments/FIGS/data/v2/index.json :: fine-tuning top N=1200000: R30 of each L188 run over the from-scratch mean, SD over runs and standard error of that mean in quadrature :: sha256 c6c4138ee5c86423
\newcommand{\MlstFtTopRthirtyFullOverScratchPr}{27\% more rejection}  % experiments/FIGS/data/v2/index.json :: fine-tuning top N=1200000: R30 of each L188 run over the from-scratch mean, SD over runs and standard error of that mean in quadrature :: sha256 c6c4138ee5c86423 :: prose twin: two significant figures, direction in words
\newcommand{\MlstFtTopRthirtyRes}{1{,}790\,\ensuremath{\pm}\,60}  % experiments/FIGS/data/v2/index.json :: fine-tuning top N=1200000: R30, L162 :: sha256 c6c4138ee5c86423
\newcommand{\MlstFtTopRthirtyResOverScratch}{1.27\,\ensuremath{\pm}\,0.04}  % experiments/FIGS/data/v2/index.json :: fine-tuning top N=1200000: R30 of each L162 run over the from-scratch mean, SD over runs and standard error of that mean in quadrature :: sha256 c6c4138ee5c86423
\newcommand{\MlstFtTopRthirtyResOverScratchPr}{27\% more rejection}  % experiments/FIGS/data/v2/index.json :: fine-tuning top N=1200000: R30 of each L162 run over the from-scratch mean, SD over runs and standard error of that mean in quadrature :: sha256 c6c4138ee5c86423 :: prose twin: two significant figures, direction in words
\newcommand{\MlstFtTopRthirtyMid}{1{,}780\,\ensuremath{\pm}\,150}  % experiments/FIGS/data/v2/index.json :: fine-tuning top N=1200000: R30, R42_Q1 :: sha256 c6c4138ee5c86423
\newcommand{\MlstFtTopRthirtyMidOverScratch}{1.26\,\ensuremath{\pm}\,0.11}  % experiments/FIGS/data/v2/index.json :: fine-tuning top N=1200000: R30 of each R42_Q1 run over the from-scratch mean, SD over runs and standard error of that mean in quadrature :: sha256 c6c4138ee5c86423
\newcommand{\MlstFtTopRthirtyMidOverScratchPr}{26\% more rejection}  % experiments/FIGS/data/v2/index.json :: fine-tuning top N=1200000: R30 of each R42_Q1 run over the from-scratch mean, SD over runs and standard error of that mean in quadrature :: sha256 c6c4138ee5c86423 :: prose twin: two significant figures, direction in words
\newcommand{\MlstFtTopRthirtyCoarse}{1{,}740\,\ensuremath{\pm}\,140}  % experiments/FIGS/data/v2/index.json :: fine-tuning top N=1200000: R30, R16_Q1 :: sha256 c6c4138ee5c86423
\newcommand{\MlstFtTopRthirtyCoarseOverScratch}{1.23\,\ensuremath{\pm}\,0.10}  % experiments/FIGS/data/v2/index.json :: fine-tuning top N=1200000: R30 of each R16_Q1 run over the from-scratch mean, SD over runs and standard error of that mean in quadrature :: sha256 c6c4138ee5c86423
\newcommand{\MlstFtTopRthirtyCoarseOverScratchPr}{23\% more rejection}  % experiments/FIGS/data/v2/index.json :: fine-tuning top N=1200000: R30 of each R16_Q1 run over the from-scratch mean, SD over runs and standard error of that mean in quadrature :: sha256 c6c4138ee5c86423 :: prose twin: two significant figures, direction in words
\newcommand{\MlstFtTopRthirtyResMass}{1{,}941 and 1{,}590}  % experiments/FIGS/data/v2/index.json :: fine-tuning top N=1200000: R30, L162_MASS :: sha256 c6c4138ee5c86423
\newcommand{\MlstFtTopRthirtyResMassOverScratch}{1.37 and 1.13}  % experiments/FIGS/data/v2/index.json :: fine-tuning top N=1200000: R30 of each L162_MASS run over the from-scratch mean, SD over runs and standard error of that mean in quadrature :: sha256 c6c4138ee5c86423
\newcommand{\MlstFtTopRthirtyResMassOverScratchPr}{at least 13\% more rejection}  % experiments/FIGS/data/v2/index.json :: fine-tuning top N=1200000: R30 of each L162_MASS run over the from-scratch mean, SD over runs and standard error of that mean in quadrature :: sha256 c6c4138ee5c86423 :: prose twin: two significant figures, direction in words
\newcommand{\MlstFtTopRthirtyResMassOverScratchPrMax}{up to 37\% more rejection}  % experiments/FIGS/data/v2/index.json :: fine-tuning top N=1200000: R30 of each L162_MASS run over the from-scratch mean, SD over runs and standard error of that mean in quadrature :: sha256 c6c4138ee5c86423 :: prose twin of the run farther from 1
\newcommand{\MlstFtTopRthirtyCoarseMass}{1{,}683 and 1{,}628}  % experiments/FIGS/data/v2/index.json :: fine-tuning top N=1200000: R30, R16_Q1_MASS :: sha256 c6c4138ee5c86423
\newcommand{\MlstFtTopRthirtyCoarseMassOverScratch}{1.192 and 1.153}  % experiments/FIGS/data/v2/index.json :: fine-tuning top N=1200000: R30 of each R16_Q1_MASS run over the from-scratch mean, SD over runs and standard error of that mean in quadrature :: sha256 c6c4138ee5c86423
\newcommand{\MlstFtTopRthirtyCoarseMassOverScratchPr}{at least 15\% more rejection}  % experiments/FIGS/data/v2/index.json :: fine-tuning top N=1200000: R30 of each R16_Q1_MASS run over the from-scratch mean, SD over runs and standard error of that mean in quadrature :: sha256 c6c4138ee5c86423 :: prose twin: two significant figures, direction in words
\newcommand{\MlstFtTopRthirtyCoarseMassOverScratchPrMax}{up to 19\% more rejection}  % experiments/FIGS/data/v2/index.json :: fine-tuning top N=1200000: R30 of each R16_Q1_MASS run over the from-scratch mean, SD over runs and standard error of that mean in quadrature :: sha256 c6c4138ee5c86423 :: prose twin of the run farther from 1
\newcommand{\MlstFtTopRthirtyResOverFull}{0.997\,\ensuremath{\pm}\,0.032}  % experiments/FIGS/data/v2/index.json :: fine-tuning top N=1200000: R30 of L162 over L188, run by run :: sha256 c6c4138ee5c86423
\newcommand{\MlstFtTopRthirtyResOverFullPr}{equal within 4\%}  % experiments/FIGS/data/v2/index.json :: fine-tuning top N=1200000: R30 of L162 over L188, run by run :: sha256 c6c4138ee5c86423 :: prose twin: two significant figures, direction in words
\newcommand{\MlstFtTopRthirtyMidOverFull}{0.99\,\ensuremath{\pm}\,0.07}  % experiments/FIGS/data/v2/index.json :: fine-tuning top N=1200000: R30 of R42_Q1 over L188, run by run :: sha256 c6c4138ee5c86423
\newcommand{\MlstFtTopRthirtyMidOverFullPr}{equal within 8\%}  % experiments/FIGS/data/v2/index.json :: fine-tuning top N=1200000: R30 of R42_Q1 over L188, run by run :: sha256 c6c4138ee5c86423 :: prose twin: two significant figures, direction in words
\newcommand{\MlstFtTopRthirtyCoarseOverFull}{0.97\,\ensuremath{\pm}\,0.10}  % experiments/FIGS/data/v2/index.json :: fine-tuning top N=1200000: R30 of R16_Q1 over L188, run by run :: sha256 c6c4138ee5c86423
\newcommand{\MlstFtTopRthirtyCoarseOverFullPr}{equal within 10\%}  % experiments/FIGS/data/v2/index.json :: fine-tuning top N=1200000: R30 of R16_Q1 over L188, run by run :: sha256 c6c4138ee5c86423 :: prose twin: two significant figures, direction in words
\newcommand{\MlstFtTopRthirtyResMassOverRes}{1.12 and 0.89}  % experiments/FIGS/data/v2/index.json :: fine-tuning top N=1200000: R30 of L162_MASS over L162, run by run :: sha256 c6c4138ee5c86423
\newcommand{\MlstFtTopRthirtyResMassOverResPr}{equal within 20\%}  % experiments/FIGS/data/v2/index.json :: fine-tuning top N=1200000: R30 of L162_MASS over L162, run by run :: sha256 c6c4138ee5c86423 :: prose twin: two significant figures, direction in words
\newcommand{\MlstFtTopRthirtyResMassOverResPrMax}{equal within 20\%}  % experiments/FIGS/data/v2/index.json :: fine-tuning top N=1200000: R30 of L162_MASS over L162, run by run :: sha256 c6c4138ee5c86423 :: prose twin of the run farther from 1
\newcommand{\MlstFtTopRthirtyCoarseMassOverCoarse}{0.900 and 0.919}  % experiments/FIGS/data/v2/index.json :: fine-tuning top N=1200000: R30 of R16_Q1_MASS over R16_Q1, run by run :: sha256 c6c4138ee5c86423
\newcommand{\MlstFtTopRthirtyCoarseMassOverCoarsePr}{equal within 10\%}  % experiments/FIGS/data/v2/index.json :: fine-tuning top N=1200000: R30 of R16_Q1_MASS over R16_Q1, run by run :: sha256 c6c4138ee5c86423 :: prose twin: two significant figures, direction in words
\newcommand{\MlstFtTopRthirtyCoarseMassOverCoarsePrMax}{equal within 10\%}  % experiments/FIGS/data/v2/index.json :: fine-tuning top N=1200000: R30 of R16_Q1_MASS over R16_Q1, run by run :: sha256 c6c4138ee5c86423 :: prose twin of the run farther from 1
\newcommand{\MlstFtTopRthirtyRandOverFull}{0.97\,\ensuremath{\pm}\,0.04}  % experiments/FIGS/data/v2/index.json :: fine-tuning top N=1200000: R30 of each random partition over the L188 run of its seed :: sha256 c6c4138ee5c86423
\newcommand{\MlstFtTopRthirtyRandOverFullPr}{equal within 4\%}  % experiments/FIGS/data/v2/index.json :: fine-tuning top N=1200000: R30 of each random partition over the L188 run of its seed :: sha256 c6c4138ee5c86423 :: prose twin: two significant figures, direction in words
\newcommand{\MlstFtTopRthirtyRandOverCoarse}{0.97\,\ensuremath{\pm}\,0.09}  % experiments/FIGS/data/v2/index.json :: fine-tuning top N=1200000: R30 of each random partition over the R16_Q1 run of its seed :: sha256 c6c4138ee5c86423
\newcommand{\MlstFtTopRthirtyRandOverCoarsePr}{equal within 9\%}  % experiments/FIGS/data/v2/index.json :: fine-tuning top N=1200000: R30 of each random partition over the R16_Q1 run of its seed :: sha256 c6c4138ee5c86423 :: prose twin: two significant figures, direction in words
\newcommand{\MlstFtTopRthirtyFullOverPart}{0.650\,\ensuremath{\pm}\,0.020}  % experiments/FIGS/data/v2/index.json :: fine-tuning top N=1200000: R30 of each 188-class run over ParT-f.t.'s published value :: sha256 c6c4138ee5c86423
\newcommand{\MlstFtTopRthirtyFullOverPartPr}{35\% less rejection}  % experiments/FIGS/data/v2/index.json :: fine-tuning top N=1200000: R30 of each 188-class run over ParT-f.t.'s published value :: sha256 c6c4138ee5c86423 :: prose twin: two significant figures, direction in words
\newcommand{\MlstFtTopAucFull}{0.9861\,\ensuremath{\pm}\,0.0001}  % experiments/FIGS/data/v2/index.json :: fine-tuning top N=1200000: AUC of the 188-class runs :: sha256 c6c4138ee5c86423
\newcommand{\MlstFtTopAucScratch}{0.9852\,\ensuremath{\pm}\,0.0001}  % experiments/FIGS/data/v2/index.json :: fine-tuning top N=1200000: AUC from scratch :: sha256 c6c4138ee5c86423
\newcommand{\MlstFtTopRthirtyFullRelSd}{3\%}  % experiments/FIGS/data/v2/index.json :: fine-tuning top, full set: relative SD of R30 over the 188-class runs :: sha256 c6c4138ee5c86423
\newcommand{\MlstFtTopRthirtyKin}{1{,}800}  % experiments/FIGS/data/v2/index.json :: fine-tuning top, full set: R30 of the kinematic-only model (one run) :: sha256 c6c4138ee5c86423
\newcommand{\MlstFtTopRthirtyKinOverFull}{1.00}  % experiments/FIGS/data/v2/index.json :: fine-tuning top, full set: R30 of the kinematic-only model over the mean of the 188-class runs :: sha256 c6c4138ee5c86423
\newcommand{\MlstFtTopRthirtyKinOverFullPr}{0.23\% more rejection}  % experiments/FIGS/data/v2/index.json :: fine-tuning top, full set: R30 of the kinematic-only model over the mean of the 188-class runs :: sha256 c6c4138ee5c86423 :: prose twin: two significant figures, direction in words
\newcommand{\MlstFtTopDataEquivJets}{1{,}200}  % experiments/FIGS/data/v2/index.json :: fine-tuning top: from-scratch training jets whose R30 equals the 188-class mean at 1000 jets (interpolated linearly in R30 against log10 N) :: sha256 c6c4138ee5c86423
\newcommand{\MlstFtTopDataEquivJetsEFour}{23{,}000}  % experiments/FIGS/data/v2/index.json :: fine-tuning top: from-scratch training jets whose R30 equals the 188-class mean at 10000 jets (interpolated linearly in R30 against log10 N) :: sha256 c6c4138ee5c86423
\newcommand{\MlstFtQgRthirtyScratchEThree}{49.7\,\ensuremath{\pm}\,2.1}  % experiments/FIGS/data/v2/index.json :: fine-tuning qg N=1000: R30 from scratch, every seed :: sha256 c6c4138ee5c86423
\newcommand{\MlstFtQgRthirtyFullEThree}{47\,\ensuremath{\pm}\,10}  % experiments/FIGS/data/v2/index.json :: fine-tuning qg N=1000: R30, L188 :: sha256 c6c4138ee5c86423
\newcommand{\MlstFtQgRthirtyFullOverScratchEThree}{0.95\,\ensuremath{\pm}\,0.19}  % experiments/FIGS/data/v2/index.json :: fine-tuning qg N=1000: R30 of each L188 run over the from-scratch mean, SD over runs and standard error of that mean in quadrature :: sha256 c6c4138ee5c86423
\newcommand{\MlstFtQgRthirtyFullOverScratchEThreePr}{equal within 20\%}  % experiments/FIGS/data/v2/index.json :: fine-tuning qg N=1000: R30 of each L188 run over the from-scratch mean, SD over runs and standard error of that mean in quadrature :: sha256 c6c4138ee5c86423 :: prose twin: two significant figures, direction in words
\newcommand{\MlstFtQgRthirtyResEThree}{52.4\,\ensuremath{\pm}\,3.1}  % experiments/FIGS/data/v2/index.json :: fine-tuning qg N=1000: R30, L162 :: sha256 c6c4138ee5c86423
\newcommand{\MlstFtQgRthirtyResOverScratchEThree}{1.06\,\ensuremath{\pm}\,0.07}  % experiments/FIGS/data/v2/index.json :: fine-tuning qg N=1000: R30 of each L162 run over the from-scratch mean, SD over runs and standard error of that mean in quadrature :: sha256 c6c4138ee5c86423
\newcommand{\MlstFtQgRthirtyResOverScratchEThreePr}{equal within 7\%}  % experiments/FIGS/data/v2/index.json :: fine-tuning qg N=1000: R30 of each L162 run over the from-scratch mean, SD over runs and standard error of that mean in quadrature :: sha256 c6c4138ee5c86423 :: prose twin: two significant figures, direction in words
\newcommand{\MlstFtQgRthirtyMidEThree}{54\,\ensuremath{\pm}\,6}  % experiments/FIGS/data/v2/index.json :: fine-tuning qg N=1000: R30, R42_Q1 :: sha256 c6c4138ee5c86423
\newcommand{\MlstFtQgRthirtyMidOverScratchEThree}{1.09\,\ensuremath{\pm}\,0.12}  % experiments/FIGS/data/v2/index.json :: fine-tuning qg N=1000: R30 of each R42_Q1 run over the from-scratch mean, SD over runs and standard error of that mean in quadrature :: sha256 c6c4138ee5c86423
\newcommand{\MlstFtQgRthirtyMidOverScratchEThreePr}{equal within 20\%}  % experiments/FIGS/data/v2/index.json :: fine-tuning qg N=1000: R30 of each R42_Q1 run over the from-scratch mean, SD over runs and standard error of that mean in quadrature :: sha256 c6c4138ee5c86423 :: prose twin: two significant figures, direction in words
\newcommand{\MlstFtQgRthirtyCoarseEThree}{55\,\ensuremath{\pm}\,5}  % experiments/FIGS/data/v2/index.json :: fine-tuning qg N=1000: R30, R16_Q1 :: sha256 c6c4138ee5c86423
\newcommand{\MlstFtQgRthirtyCoarseOverScratchEThree}{1.12\,\ensuremath{\pm}\,0.10}  % experiments/FIGS/data/v2/index.json :: fine-tuning qg N=1000: R30 of each R16_Q1 run over the from-scratch mean, SD over runs and standard error of that mean in quadrature :: sha256 c6c4138ee5c86423
\newcommand{\MlstFtQgRthirtyCoarseOverScratchEThreePr}{equal within 20\%}  % experiments/FIGS/data/v2/index.json :: fine-tuning qg N=1000: R30 of each R16_Q1 run over the from-scratch mean, SD over runs and standard error of that mean in quadrature :: sha256 c6c4138ee5c86423 :: prose twin: two significant figures, direction in words
\newcommand{\MlstFtQgRthirtyResMassEThree}{59.3 and 57.1}  % experiments/FIGS/data/v2/index.json :: fine-tuning qg N=1000: R30, L162_MASS :: sha256 c6c4138ee5c86423
\newcommand{\MlstFtQgRthirtyResMassOverScratchEThree}{1.193 and 1.149}  % experiments/FIGS/data/v2/index.json :: fine-tuning qg N=1000: R30 of each L162_MASS run over the from-scratch mean, SD over runs and standard error of that mean in quadrature :: sha256 c6c4138ee5c86423
\newcommand{\MlstFtQgRthirtyResMassOverScratchEThreePr}{at least 15\% more rejection}  % experiments/FIGS/data/v2/index.json :: fine-tuning qg N=1000: R30 of each L162_MASS run over the from-scratch mean, SD over runs and standard error of that mean in quadrature :: sha256 c6c4138ee5c86423 :: prose twin: two significant figures, direction in words
\newcommand{\MlstFtQgRthirtyResMassOverScratchEThreePrMax}{up to 19\% more rejection}  % experiments/FIGS/data/v2/index.json :: fine-tuning qg N=1000: R30 of each L162_MASS run over the from-scratch mean, SD over runs and standard error of that mean in quadrature :: sha256 c6c4138ee5c86423 :: prose twin of the run farther from 1
\newcommand{\MlstFtQgRthirtyCoarseMassEThree}{54.2 and 57.7}  % experiments/FIGS/data/v2/index.json :: fine-tuning qg N=1000: R30, R16_Q1_MASS :: sha256 c6c4138ee5c86423
\newcommand{\MlstFtQgRthirtyCoarseMassOverScratchEThree}{1.091 and 1.161}  % experiments/FIGS/data/v2/index.json :: fine-tuning qg N=1000: R30 of each R16_Q1_MASS run over the from-scratch mean, SD over runs and standard error of that mean in quadrature :: sha256 c6c4138ee5c86423
\newcommand{\MlstFtQgRthirtyCoarseMassOverScratchEThreePr}{at least 9.1\% more rejection}  % experiments/FIGS/data/v2/index.json :: fine-tuning qg N=1000: R30 of each R16_Q1_MASS run over the from-scratch mean, SD over runs and standard error of that mean in quadrature :: sha256 c6c4138ee5c86423 :: prose twin: two significant figures, direction in words
\newcommand{\MlstFtQgRthirtyCoarseMassOverScratchEThreePrMax}{up to 16\% more rejection}  % experiments/FIGS/data/v2/index.json :: fine-tuning qg N=1000: R30 of each R16_Q1_MASS run over the from-scratch mean, SD over runs and standard error of that mean in quadrature :: sha256 c6c4138ee5c86423 :: prose twin of the run farther from 1
\newcommand{\MlstFtQgRthirtyResOverFullEThree}{1.14\,\ensuremath{\pm}\,0.17}  % experiments/FIGS/data/v2/index.json :: fine-tuning qg N=1000: R30 of L162 over L188, run by run :: sha256 c6c4138ee5c86423
\newcommand{\MlstFtQgRthirtyResOverFullEThreePr}{equal within 20\%}  % experiments/FIGS/data/v2/index.json :: fine-tuning qg N=1000: R30 of L162 over L188, run by run :: sha256 c6c4138ee5c86423 :: prose twin: two significant figures, direction in words
\newcommand{\MlstFtQgRthirtyMidOverFullEThree}{1.17\,\ensuremath{\pm}\,0.14}  % experiments/FIGS/data/v2/index.json :: fine-tuning qg N=1000: R30 of R42_Q1 over L188, run by run :: sha256 c6c4138ee5c86423
\newcommand{\MlstFtQgRthirtyMidOverFullEThreePr}{equal within 20\%}  % experiments/FIGS/data/v2/index.json :: fine-tuning qg N=1000: R30 of R42_Q1 over L188, run by run :: sha256 c6c4138ee5c86423 :: prose twin: two significant figures, direction in words
\newcommand{\MlstFtQgRthirtyCoarseOverFullEThree}{1.21\,\ensuremath{\pm}\,0.27}  % experiments/FIGS/data/v2/index.json :: fine-tuning qg N=1000: R30 of R16_Q1 over L188, run by run :: sha256 c6c4138ee5c86423
\newcommand{\MlstFtQgRthirtyCoarseOverFullEThreePr}{equal within 30\%}  % experiments/FIGS/data/v2/index.json :: fine-tuning qg N=1000: R30 of R16_Q1 over L188, run by run :: sha256 c6c4138ee5c86423 :: prose twin: two significant figures, direction in words
\newcommand{\MlstFtQgRthirtyResMassOverResEThree}{1.108 and 1.042}  % experiments/FIGS/data/v2/index.json :: fine-tuning qg N=1000: R30 of L162_MASS over L162, run by run :: sha256 c6c4138ee5c86423
\newcommand{\MlstFtQgRthirtyResMassOverResEThreePr}{equal within 20\%}  % experiments/FIGS/data/v2/index.json :: fine-tuning qg N=1000: R30 of L162_MASS over L162, run by run :: sha256 c6c4138ee5c86423 :: prose twin: two significant figures, direction in words
\newcommand{\MlstFtQgRthirtyResMassOverResEThreePrMax}{equal within 20\%}  % experiments/FIGS/data/v2/index.json :: fine-tuning qg N=1000: R30 of L162_MASS over L162, run by run :: sha256 c6c4138ee5c86423 :: prose twin of the run farther from 1
\newcommand{\MlstFtQgRthirtyCoarseMassOverCoarseEThree}{0.89 and 1.11}  % experiments/FIGS/data/v2/index.json :: fine-tuning qg N=1000: R30 of R16_Q1_MASS over R16_Q1, run by run :: sha256 c6c4138ee5c86423
\newcommand{\MlstFtQgRthirtyCoarseMassOverCoarseEThreePr}{equal within 20\%}  % experiments/FIGS/data/v2/index.json :: fine-tuning qg N=1000: R30 of R16_Q1_MASS over R16_Q1, run by run :: sha256 c6c4138ee5c86423 :: prose twin: two significant figures, direction in words
\newcommand{\MlstFtQgRthirtyCoarseMassOverCoarseEThreePrMax}{equal within 20\%}  % experiments/FIGS/data/v2/index.json :: fine-tuning qg N=1000: R30 of R16_Q1_MASS over R16_Q1, run by run :: sha256 c6c4138ee5c86423 :: prose twin of the run farther from 1
\newcommand{\MlstFtQgRthirtyRandOverFullEThree}{0.91\,\ensuremath{\pm}\,0.07}  % experiments/FIGS/data/v2/index.json :: fine-tuning qg N=1000: R30 of each random partition over the L188 run of its seed :: sha256 c6c4138ee5c86423
\newcommand{\MlstFtQgRthirtyRandOverFullEThreePr}{equal within 10\%}  % experiments/FIGS/data/v2/index.json :: fine-tuning qg N=1000: R30 of each random partition over the L188 run of its seed :: sha256 c6c4138ee5c86423 :: prose twin: two significant figures, direction in words
\newcommand{\MlstFtQgRthirtyRandOverCoarseEThree}{0.87\,\ensuremath{\pm}\,0.14}  % experiments/FIGS/data/v2/index.json :: fine-tuning qg N=1000: R30 of each random partition over the R16_Q1 run of its seed :: sha256 c6c4138ee5c86423
\newcommand{\MlstFtQgRthirtyRandOverCoarseEThreePr}{equal within 20\%}  % experiments/FIGS/data/v2/index.json :: fine-tuning qg N=1000: R30 of each random partition over the R16_Q1 run of its seed :: sha256 c6c4138ee5c86423 :: prose twin: two significant figures, direction in words
\newcommand{\MlstFtQgRthirtyScratchEFour}{66\,\ensuremath{\pm}\,9}  % experiments/FIGS/data/v2/index.json :: fine-tuning qg N=10000: R30 from scratch, every seed :: sha256 c6c4138ee5c86423
\newcommand{\MlstFtQgRthirtyFullEFour}{75\,\ensuremath{\pm}\,7}  % experiments/FIGS/data/v2/index.json :: fine-tuning qg N=10000: R30, L188 :: sha256 c6c4138ee5c86423
\newcommand{\MlstFtQgRthirtyFullOverScratchEFour}{1.14\,\ensuremath{\pm}\,0.14}  % experiments/FIGS/data/v2/index.json :: fine-tuning qg N=10000: R30 of each L188 run over the from-scratch mean, SD over runs and standard error of that mean in quadrature :: sha256 c6c4138ee5c86423
\newcommand{\MlstFtQgRthirtyFullOverScratchEFourPr}{equal within 20\%}  % experiments/FIGS/data/v2/index.json :: fine-tuning qg N=10000: R30 of each L188 run over the from-scratch mean, SD over runs and standard error of that mean in quadrature :: sha256 c6c4138ee5c86423 :: prose twin: two significant figures, direction in words
\newcommand{\MlstFtQgRthirtyResEFour}{76.4\,\ensuremath{\pm}\,0.6}  % experiments/FIGS/data/v2/index.json :: fine-tuning qg N=10000: R30, L162 :: sha256 c6c4138ee5c86423
\newcommand{\MlstFtQgRthirtyResOverScratchEFour}{1.15\,\ensuremath{\pm}\,0.09}  % experiments/FIGS/data/v2/index.json :: fine-tuning qg N=10000: R30 of each L162 run over the from-scratch mean, SD over runs and standard error of that mean in quadrature :: sha256 c6c4138ee5c86423
\newcommand{\MlstFtQgRthirtyResOverScratchEFourPr}{15\% more rejection}  % experiments/FIGS/data/v2/index.json :: fine-tuning qg N=10000: R30 of each L162 run over the from-scratch mean, SD over runs and standard error of that mean in quadrature :: sha256 c6c4138ee5c86423 :: prose twin: two significant figures, direction in words
\newcommand{\MlstFtQgRthirtyMidEFour}{74\,\ensuremath{\pm}\,5}  % experiments/FIGS/data/v2/index.json :: fine-tuning qg N=10000: R30, R42_Q1 :: sha256 c6c4138ee5c86423
\newcommand{\MlstFtQgRthirtyMidOverScratchEFour}{1.12\,\ensuremath{\pm}\,0.11}  % experiments/FIGS/data/v2/index.json :: fine-tuning qg N=10000: R30 of each R42_Q1 run over the from-scratch mean, SD over runs and standard error of that mean in quadrature :: sha256 c6c4138ee5c86423
\newcommand{\MlstFtQgRthirtyMidOverScratchEFourPr}{equal within 20\%}  % experiments/FIGS/data/v2/index.json :: fine-tuning qg N=10000: R30 of each R42_Q1 run over the from-scratch mean, SD over runs and standard error of that mean in quadrature :: sha256 c6c4138ee5c86423 :: prose twin: two significant figures, direction in words
\newcommand{\MlstFtQgRthirtyCoarseEFour}{80.3\,\ensuremath{\pm}\,1.1}  % experiments/FIGS/data/v2/index.json :: fine-tuning qg N=10000: R30, R16_Q1 :: sha256 c6c4138ee5c86423
\newcommand{\MlstFtQgRthirtyCoarseOverScratchEFour}{1.21\,\ensuremath{\pm}\,0.10}  % experiments/FIGS/data/v2/index.json :: fine-tuning qg N=10000: R30 of each R16_Q1 run over the from-scratch mean, SD over runs and standard error of that mean in quadrature :: sha256 c6c4138ee5c86423
\newcommand{\MlstFtQgRthirtyCoarseOverScratchEFourPr}{21\% more rejection}  % experiments/FIGS/data/v2/index.json :: fine-tuning qg N=10000: R30 of each R16_Q1 run over the from-scratch mean, SD over runs and standard error of that mean in quadrature :: sha256 c6c4138ee5c86423 :: prose twin: two significant figures, direction in words
\newcommand{\MlstFtQgRthirtyResMassEFour}{80.1 and 72.6}  % experiments/FIGS/data/v2/index.json :: fine-tuning qg N=10000: R30, L162_MASS :: sha256 c6c4138ee5c86423
\newcommand{\MlstFtQgRthirtyResMassOverScratchEFour}{1.21 and 1.09}  % experiments/FIGS/data/v2/index.json :: fine-tuning qg N=10000: R30 of each L162_MASS run over the from-scratch mean, SD over runs and standard error of that mean in quadrature :: sha256 c6c4138ee5c86423
\newcommand{\MlstFtQgRthirtyResMassOverScratchEFourPr}{at least 9.3\% more rejection}  % experiments/FIGS/data/v2/index.json :: fine-tuning qg N=10000: R30 of each L162_MASS run over the from-scratch mean, SD over runs and standard error of that mean in quadrature :: sha256 c6c4138ee5c86423 :: prose twin: two significant figures, direction in words
\newcommand{\MlstFtQgRthirtyResMassOverScratchEFourPrMax}{up to 21\% more rejection}  % experiments/FIGS/data/v2/index.json :: fine-tuning qg N=10000: R30 of each L162_MASS run over the from-scratch mean, SD over runs and standard error of that mean in quadrature :: sha256 c6c4138ee5c86423 :: prose twin of the run farther from 1
\newcommand{\MlstFtQgRthirtyCoarseMassEFour}{79.5 and 77.9}  % experiments/FIGS/data/v2/index.json :: fine-tuning qg N=10000: R30, R16_Q1_MASS :: sha256 c6c4138ee5c86423
\newcommand{\MlstFtQgRthirtyCoarseMassOverScratchEFour}{1.198 and 1.173}  % experiments/FIGS/data/v2/index.json :: fine-tuning qg N=10000: R30 of each R16_Q1_MASS run over the from-scratch mean, SD over runs and standard error of that mean in quadrature :: sha256 c6c4138ee5c86423
\newcommand{\MlstFtQgRthirtyCoarseMassOverScratchEFourPr}{at least 17\% more rejection}  % experiments/FIGS/data/v2/index.json :: fine-tuning qg N=10000: R30 of each R16_Q1_MASS run over the from-scratch mean, SD over runs and standard error of that mean in quadrature :: sha256 c6c4138ee5c86423 :: prose twin: two significant figures, direction in words
\newcommand{\MlstFtQgRthirtyCoarseMassOverScratchEFourPrMax}{up to 20\% more rejection}  % experiments/FIGS/data/v2/index.json :: fine-tuning qg N=10000: R30 of each R16_Q1_MASS run over the from-scratch mean, SD over runs and standard error of that mean in quadrature :: sha256 c6c4138ee5c86423 :: prose twin of the run farther from 1
\newcommand{\MlstFtQgRthirtyResOverFullEFour}{1.02\,\ensuremath{\pm}\,0.10}  % experiments/FIGS/data/v2/index.json :: fine-tuning qg N=10000: R30 of L162 over L188, run by run :: sha256 c6c4138ee5c86423
\newcommand{\MlstFtQgRthirtyResOverFullEFourPr}{equal within 10\%}  % experiments/FIGS/data/v2/index.json :: fine-tuning qg N=10000: R30 of L162 over L188, run by run :: sha256 c6c4138ee5c86423 :: prose twin: two significant figures, direction in words
\newcommand{\MlstFtQgRthirtyMidOverFullEFour}{0.99\,\ensuremath{\pm}\,0.14}  % experiments/FIGS/data/v2/index.json :: fine-tuning qg N=10000: R30 of R42_Q1 over L188, run by run :: sha256 c6c4138ee5c86423
\newcommand{\MlstFtQgRthirtyMidOverFullEFourPr}{equal within 20\%}  % experiments/FIGS/data/v2/index.json :: fine-tuning qg N=10000: R30 of R42_Q1 over L188, run by run :: sha256 c6c4138ee5c86423 :: prose twin: two significant figures, direction in words
\newcommand{\MlstFtQgRthirtyCoarseOverFullEFour}{1.07\,\ensuremath{\pm}\,0.13}  % experiments/FIGS/data/v2/index.json :: fine-tuning qg N=10000: R30 of R16_Q1 over L188, run by run :: sha256 c6c4138ee5c86423
\newcommand{\MlstFtQgRthirtyCoarseOverFullEFourPr}{equal within 20\%}  % experiments/FIGS/data/v2/index.json :: fine-tuning qg N=10000: R30 of R16_Q1 over L188, run by run :: sha256 c6c4138ee5c86423 :: prose twin: two significant figures, direction in words
\newcommand{\MlstFtQgRthirtyResMassOverResEFour}{1.06 and 0.95}  % experiments/FIGS/data/v2/index.json :: fine-tuning qg N=10000: R30 of L162_MASS over L162, run by run :: sha256 c6c4138ee5c86423
\newcommand{\MlstFtQgRthirtyResMassOverResEFourPr}{equal within 6\%}  % experiments/FIGS/data/v2/index.json :: fine-tuning qg N=10000: R30 of L162_MASS over L162, run by run :: sha256 c6c4138ee5c86423 :: prose twin: two significant figures, direction in words
\newcommand{\MlstFtQgRthirtyResMassOverResEFourPrMax}{equal within 6\%}  % experiments/FIGS/data/v2/index.json :: fine-tuning qg N=10000: R30 of L162_MASS over L162, run by run :: sha256 c6c4138ee5c86423 :: prose twin of the run farther from 1
\newcommand{\MlstFtQgRthirtyCoarseMassOverCoarseEFour}{0.9769 and 0.9696}  % experiments/FIGS/data/v2/index.json :: fine-tuning qg N=10000: R30 of R16_Q1_MASS over R16_Q1, run by run :: sha256 c6c4138ee5c86423
\newcommand{\MlstFtQgRthirtyCoarseMassOverCoarseEFourPr}{equal within 4\%}  % experiments/FIGS/data/v2/index.json :: fine-tuning qg N=10000: R30 of R16_Q1_MASS over R16_Q1, run by run :: sha256 c6c4138ee5c86423 :: prose twin: two significant figures, direction in words
\newcommand{\MlstFtQgRthirtyCoarseMassOverCoarseEFourPrMax}{equal within 4\%}  % experiments/FIGS/data/v2/index.json :: fine-tuning qg N=10000: R30 of R16_Q1_MASS over R16_Q1, run by run :: sha256 c6c4138ee5c86423 :: prose twin of the run farther from 1
\newcommand{\MlstFtQgRthirtyRandOverFullEFour}{1.03\,\ensuremath{\pm}\,0.18}  % experiments/FIGS/data/v2/index.json :: fine-tuning qg N=10000: R30 of each random partition over the L188 run of its seed :: sha256 c6c4138ee5c86423
\newcommand{\MlstFtQgRthirtyRandOverFullEFourPr}{equal within 20\%}  % experiments/FIGS/data/v2/index.json :: fine-tuning qg N=10000: R30 of each random partition over the L188 run of its seed :: sha256 c6c4138ee5c86423 :: prose twin: two significant figures, direction in words
\newcommand{\MlstFtQgRthirtyRandOverCoarseEFour}{0.93\,\ensuremath{\pm}\,0.12}  % experiments/FIGS/data/v2/index.json :: fine-tuning qg N=10000: R30 of each random partition over the R16_Q1 run of its seed :: sha256 c6c4138ee5c86423
\newcommand{\MlstFtQgRthirtyRandOverCoarseEFourPr}{equal within 20\%}  % experiments/FIGS/data/v2/index.json :: fine-tuning qg N=10000: R30 of each random partition over the R16_Q1 run of its seed :: sha256 c6c4138ee5c86423 :: prose twin: two significant figures, direction in words
\newcommand{\MlstFtQgRthirtyScratchEFive}{88.4\,\ensuremath{\pm}\,0.8}  % experiments/FIGS/data/v2/index.json :: fine-tuning qg N=100000: R30 from scratch, every seed :: sha256 c6c4138ee5c86423
\newcommand{\MlstFtQgRthirtyFullEFive}{99.3\,\ensuremath{\pm}\,1.4}  % experiments/FIGS/data/v2/index.json :: fine-tuning qg N=100000: R30, L188 :: sha256 c6c4138ee5c86423
\newcommand{\MlstFtQgRthirtyFullOverScratchEFive}{1.124\,\ensuremath{\pm}\,0.017}  % experiments/FIGS/data/v2/index.json :: fine-tuning qg N=100000: R30 of each L188 run over the from-scratch mean, SD over runs and standard error of that mean in quadrature :: sha256 c6c4138ee5c86423
\newcommand{\MlstFtQgRthirtyFullOverScratchEFivePr}{12\% more rejection}  % experiments/FIGS/data/v2/index.json :: fine-tuning qg N=100000: R30 of each L188 run over the from-scratch mean, SD over runs and standard error of that mean in quadrature :: sha256 c6c4138ee5c86423 :: prose twin: two significant figures, direction in words
\newcommand{\MlstFtQgRthirtyResEFive}{97.9\,\ensuremath{\pm}\,3.3}  % experiments/FIGS/data/v2/index.json :: fine-tuning qg N=100000: R30, L162 :: sha256 c6c4138ee5c86423
\newcommand{\MlstFtQgRthirtyResOverScratchEFive}{1.11\,\ensuremath{\pm}\,0.04}  % experiments/FIGS/data/v2/index.json :: fine-tuning qg N=100000: R30 of each L162 run over the from-scratch mean, SD over runs and standard error of that mean in quadrature :: sha256 c6c4138ee5c86423
\newcommand{\MlstFtQgRthirtyResOverScratchEFivePr}{11\% more rejection}  % experiments/FIGS/data/v2/index.json :: fine-tuning qg N=100000: R30 of each L162 run over the from-scratch mean, SD over runs and standard error of that mean in quadrature :: sha256 c6c4138ee5c86423 :: prose twin: two significant figures, direction in words
\newcommand{\MlstFtQgRthirtyMidEFive}{96.1\,\ensuremath{\pm}\,1.8}  % experiments/FIGS/data/v2/index.json :: fine-tuning qg N=100000: R30, R42_Q1 :: sha256 c6c4138ee5c86423
\newcommand{\MlstFtQgRthirtyMidOverScratchEFive}{1.088\,\ensuremath{\pm}\,0.021}  % experiments/FIGS/data/v2/index.json :: fine-tuning qg N=100000: R30 of each R42_Q1 run over the from-scratch mean, SD over runs and standard error of that mean in quadrature :: sha256 c6c4138ee5c86423
\newcommand{\MlstFtQgRthirtyMidOverScratchEFivePr}{8.8\% more rejection}  % experiments/FIGS/data/v2/index.json :: fine-tuning qg N=100000: R30 of each R42_Q1 run over the from-scratch mean, SD over runs and standard error of that mean in quadrature :: sha256 c6c4138ee5c86423 :: prose twin: two significant figures, direction in words
\newcommand{\MlstFtQgRthirtyCoarseEFive}{96.4\,\ensuremath{\pm}\,1.0}  % experiments/FIGS/data/v2/index.json :: fine-tuning qg N=100000: R30, R16_Q1 :: sha256 c6c4138ee5c86423
\newcommand{\MlstFtQgRthirtyCoarseOverScratchEFive}{1.091\,\ensuremath{\pm}\,0.013}  % experiments/FIGS/data/v2/index.json :: fine-tuning qg N=100000: R30 of each R16_Q1 run over the from-scratch mean, SD over runs and standard error of that mean in quadrature :: sha256 c6c4138ee5c86423
\newcommand{\MlstFtQgRthirtyCoarseOverScratchEFivePr}{9.1\% more rejection}  % experiments/FIGS/data/v2/index.json :: fine-tuning qg N=100000: R30 of each R16_Q1 run over the from-scratch mean, SD over runs and standard error of that mean in quadrature :: sha256 c6c4138ee5c86423 :: prose twin: two significant figures, direction in words
\newcommand{\MlstFtQgRthirtyResMassEFive}{95.6 and 97.6}  % experiments/FIGS/data/v2/index.json :: fine-tuning qg N=100000: R30, L162_MASS :: sha256 c6c4138ee5c86423
\newcommand{\MlstFtQgRthirtyResMassOverScratchEFive}{1.082 and 1.104}  % experiments/FIGS/data/v2/index.json :: fine-tuning qg N=100000: R30 of each L162_MASS run over the from-scratch mean, SD over runs and standard error of that mean in quadrature :: sha256 c6c4138ee5c86423
\newcommand{\MlstFtQgRthirtyResMassOverScratchEFivePr}{at least 8.2\% more rejection}  % experiments/FIGS/data/v2/index.json :: fine-tuning qg N=100000: R30 of each L162_MASS run over the from-scratch mean, SD over runs and standard error of that mean in quadrature :: sha256 c6c4138ee5c86423 :: prose twin: two significant figures, direction in words
\newcommand{\MlstFtQgRthirtyResMassOverScratchEFivePrMax}{up to 10\% more rejection}  % experiments/FIGS/data/v2/index.json :: fine-tuning qg N=100000: R30 of each L162_MASS run over the from-scratch mean, SD over runs and standard error of that mean in quadrature :: sha256 c6c4138ee5c86423 :: prose twin of the run farther from 1
\newcommand{\MlstFtQgRthirtyCoarseMassEFive}{97.85 and 97.94}  % experiments/FIGS/data/v2/index.json :: fine-tuning qg N=100000: R30, R16_Q1_MASS :: sha256 c6c4138ee5c86423
\newcommand{\MlstFtQgRthirtyCoarseMassOverScratchEFive}{1.107 and 1.108}  % experiments/FIGS/data/v2/index.json :: fine-tuning qg N=100000: R30 of each R16_Q1_MASS run over the from-scratch mean, SD over runs and standard error of that mean in quadrature :: sha256 c6c4138ee5c86423
\newcommand{\MlstFtQgRthirtyCoarseMassOverScratchEFivePr}{at least 11\% more rejection}  % experiments/FIGS/data/v2/index.json :: fine-tuning qg N=100000: R30 of each R16_Q1_MASS run over the from-scratch mean, SD over runs and standard error of that mean in quadrature :: sha256 c6c4138ee5c86423 :: prose twin: two significant figures, direction in words
\newcommand{\MlstFtQgRthirtyCoarseMassOverScratchEFivePrMax}{up to 11\% more rejection}  % experiments/FIGS/data/v2/index.json :: fine-tuning qg N=100000: R30 of each R16_Q1_MASS run over the from-scratch mean, SD over runs and standard error of that mean in quadrature :: sha256 c6c4138ee5c86423 :: prose twin of the run farther from 1
\newcommand{\MlstFtQgRthirtyResOverFullEFive}{0.985\,\ensuremath{\pm}\,0.020}  % experiments/FIGS/data/v2/index.json :: fine-tuning qg N=100000: R30 of L162 over L188, run by run :: sha256 c6c4138ee5c86423
\newcommand{\MlstFtQgRthirtyResOverFullEFivePr}{equal within 2\%}  % experiments/FIGS/data/v2/index.json :: fine-tuning qg N=100000: R30 of L162 over L188, run by run :: sha256 c6c4138ee5c86423 :: prose twin: two significant figures, direction in words
\newcommand{\MlstFtQgRthirtyMidOverFullEFive}{0.969\,\ensuremath{\pm}\,0.030}  % experiments/FIGS/data/v2/index.json :: fine-tuning qg N=100000: R30 of R42_Q1 over L188, run by run :: sha256 c6c4138ee5c86423
\newcommand{\MlstFtQgRthirtyMidOverFullEFivePr}{equal within 4\%}  % experiments/FIGS/data/v2/index.json :: fine-tuning qg N=100000: R30 of R42_Q1 over L188, run by run :: sha256 c6c4138ee5c86423 :: prose twin: two significant figures, direction in words
\newcommand{\MlstFtQgRthirtyCoarseOverFullEFive}{0.972\,\ensuremath{\pm}\,0.020}  % experiments/FIGS/data/v2/index.json :: fine-tuning qg N=100000: R30 of R16_Q1 over L188, run by run :: sha256 c6c4138ee5c86423
\newcommand{\MlstFtQgRthirtyCoarseOverFullEFivePr}{equal within 3\%}  % experiments/FIGS/data/v2/index.json :: fine-tuning qg N=100000: R30 of R16_Q1 over L188, run by run :: sha256 c6c4138ee5c86423 :: prose twin: two significant figures, direction in words
\newcommand{\MlstFtQgRthirtyResMassOverResEFive}{0.942 and 1.027}  % experiments/FIGS/data/v2/index.json :: fine-tuning qg N=100000: R30 of L162_MASS over L162, run by run :: sha256 c6c4138ee5c86423
\newcommand{\MlstFtQgRthirtyResMassOverResEFivePr}{equal within 6\%}  % experiments/FIGS/data/v2/index.json :: fine-tuning qg N=100000: R30 of L162_MASS over L162, run by run :: sha256 c6c4138ee5c86423 :: prose twin: two significant figures, direction in words
\newcommand{\MlstFtQgRthirtyResMassOverResEFivePrMax}{equal within 6\%}  % experiments/FIGS/data/v2/index.json :: fine-tuning qg N=100000: R30 of L162_MASS over L162, run by run :: sha256 c6c4138ee5c86423 :: prose twin of the run farther from 1
\newcommand{\MlstFtQgRthirtyCoarseMassOverCoarseEFive}{1.018 and 1.004}  % experiments/FIGS/data/v2/index.json :: fine-tuning qg N=100000: R30 of R16_Q1_MASS over R16_Q1, run by run :: sha256 c6c4138ee5c86423
\newcommand{\MlstFtQgRthirtyCoarseMassOverCoarseEFivePr}{equal within 2\%}  % experiments/FIGS/data/v2/index.json :: fine-tuning qg N=100000: R30 of R16_Q1_MASS over R16_Q1, run by run :: sha256 c6c4138ee5c86423 :: prose twin: two significant figures, direction in words
\newcommand{\MlstFtQgRthirtyCoarseMassOverCoarseEFivePrMax}{equal within 2\%}  % experiments/FIGS/data/v2/index.json :: fine-tuning qg N=100000: R30 of R16_Q1_MASS over R16_Q1, run by run :: sha256 c6c4138ee5c86423 :: prose twin of the run farther from 1
\newcommand{\MlstFtQgRthirtyRandOverFullEFive}{0.986\,\ensuremath{\pm}\,0.025}  % experiments/FIGS/data/v2/index.json :: fine-tuning qg N=100000: R30 of each random partition over the L188 run of its seed :: sha256 c6c4138ee5c86423
\newcommand{\MlstFtQgRthirtyRandOverFullEFivePr}{equal within 3\%}  % experiments/FIGS/data/v2/index.json :: fine-tuning qg N=100000: R30 of each random partition over the L188 run of its seed :: sha256 c6c4138ee5c86423 :: prose twin: two significant figures, direction in words
\newcommand{\MlstFtQgRthirtyRandOverCoarseEFive}{1.009\,\ensuremath{\pm}\,0.020}  % experiments/FIGS/data/v2/index.json :: fine-tuning qg N=100000: R30 of each random partition over the R16_Q1 run of its seed :: sha256 c6c4138ee5c86423
\newcommand{\MlstFtQgRthirtyRandOverCoarseEFivePr}{equal within 3\%}  % experiments/FIGS/data/v2/index.json :: fine-tuning qg N=100000: R30 of each random partition over the R16_Q1 run of its seed :: sha256 c6c4138ee5c86423 :: prose twin: two significant figures, direction in words
\newcommand{\MlstFtQgRthirtyScratch}{104.0\,\ensuremath{\pm}\,1.0}  % experiments/FIGS/data/v2/index.json :: fine-tuning qg N=1600000: R30 from scratch, every seed :: sha256 c6c4138ee5c86423
\newcommand{\MlstFtQgRthirtyFull}{107.3\,\ensuremath{\pm}\,2.2}  % experiments/FIGS/data/v2/index.json :: fine-tuning qg N=1600000: R30, L188 :: sha256 c6c4138ee5c86423
\newcommand{\MlstFtQgRthirtyFullOverScratch}{1.031\,\ensuremath{\pm}\,0.022}  % experiments/FIGS/data/v2/index.json :: fine-tuning qg N=1600000: R30 of each L188 run over the from-scratch mean, SD over runs and standard error of that mean in quadrature :: sha256 c6c4138ee5c86423
\newcommand{\MlstFtQgRthirtyFullOverScratchPr}{equal within 4\%}  % experiments/FIGS/data/v2/index.json :: fine-tuning qg N=1600000: R30 of each L188 run over the from-scratch mean, SD over runs and standard error of that mean in quadrature :: sha256 c6c4138ee5c86423 :: prose twin: two significant figures, direction in words
\newcommand{\MlstFtQgRthirtyRes}{108.0\,\ensuremath{\pm}\,1.4}  % experiments/FIGS/data/v2/index.json :: fine-tuning qg N=1600000: R30, L162 :: sha256 c6c4138ee5c86423
\newcommand{\MlstFtQgRthirtyResOverScratch}{1.038\,\ensuremath{\pm}\,0.015}  % experiments/FIGS/data/v2/index.json :: fine-tuning qg N=1600000: R30 of each L162 run over the from-scratch mean, SD over runs and standard error of that mean in quadrature :: sha256 c6c4138ee5c86423
\newcommand{\MlstFtQgRthirtyResOverScratchPr}{3.8\% more rejection}  % experiments/FIGS/data/v2/index.json :: fine-tuning qg N=1600000: R30 of each L162 run over the from-scratch mean, SD over runs and standard error of that mean in quadrature :: sha256 c6c4138ee5c86423 :: prose twin: two significant figures, direction in words
\newcommand{\MlstFtQgRthirtyMid}{108.3\,\ensuremath{\pm}\,1.0}  % experiments/FIGS/data/v2/index.json :: fine-tuning qg N=1600000: R30, R42_Q1 :: sha256 c6c4138ee5c86423
\newcommand{\MlstFtQgRthirtyMidOverScratch}{1.041\,\ensuremath{\pm}\,0.012}  % experiments/FIGS/data/v2/index.json :: fine-tuning qg N=1600000: R30 of each R42_Q1 run over the from-scratch mean, SD over runs and standard error of that mean in quadrature :: sha256 c6c4138ee5c86423
\newcommand{\MlstFtQgRthirtyMidOverScratchPr}{4.1\% more rejection}  % experiments/FIGS/data/v2/index.json :: fine-tuning qg N=1600000: R30 of each R42_Q1 run over the from-scratch mean, SD over runs and standard error of that mean in quadrature :: sha256 c6c4138ee5c86423 :: prose twin: two significant figures, direction in words
\newcommand{\MlstFtQgRthirtyCoarse}{106.5\,\ensuremath{\pm}\,0.8}  % experiments/FIGS/data/v2/index.json :: fine-tuning qg N=1600000: R30, R16_Q1 :: sha256 c6c4138ee5c86423
\newcommand{\MlstFtQgRthirtyCoarseOverScratch}{1.024\,\ensuremath{\pm}\,0.010}  % experiments/FIGS/data/v2/index.json :: fine-tuning qg N=1600000: R30 of each R16_Q1 run over the from-scratch mean, SD over runs and standard error of that mean in quadrature :: sha256 c6c4138ee5c86423
\newcommand{\MlstFtQgRthirtyCoarseOverScratchPr}{2.4\% more rejection}  % experiments/FIGS/data/v2/index.json :: fine-tuning qg N=1600000: R30 of each R16_Q1 run over the from-scratch mean, SD over runs and standard error of that mean in quadrature :: sha256 c6c4138ee5c86423 :: prose twin: two significant figures, direction in words
\newcommand{\MlstFtQgRthirtyResMass}{107.0 and 106.6}  % experiments/FIGS/data/v2/index.json :: fine-tuning qg N=1600000: R30, L162_MASS :: sha256 c6c4138ee5c86423
\newcommand{\MlstFtQgRthirtyResMassOverScratch}{1.028 and 1.025}  % experiments/FIGS/data/v2/index.json :: fine-tuning qg N=1600000: R30 of each L162_MASS run over the from-scratch mean, SD over runs and standard error of that mean in quadrature :: sha256 c6c4138ee5c86423
\newcommand{\MlstFtQgRthirtyResMassOverScratchPr}{at least 2.5\% more rejection}  % experiments/FIGS/data/v2/index.json :: fine-tuning qg N=1600000: R30 of each L162_MASS run over the from-scratch mean, SD over runs and standard error of that mean in quadrature :: sha256 c6c4138ee5c86423 :: prose twin: two significant figures, direction in words
\newcommand{\MlstFtQgRthirtyResMassOverScratchPrMax}{up to 2.8\% more rejection}  % experiments/FIGS/data/v2/index.json :: fine-tuning qg N=1600000: R30 of each L162_MASS run over the from-scratch mean, SD over runs and standard error of that mean in quadrature :: sha256 c6c4138ee5c86423 :: prose twin of the run farther from 1
\newcommand{\MlstFtQgRthirtyCoarseMass}{104.9 and 105.6}  % experiments/FIGS/data/v2/index.json :: fine-tuning qg N=1600000: R30, R16_Q1_MASS :: sha256 c6c4138ee5c86423
\newcommand{\MlstFtQgRthirtyCoarseMassOverScratch}{1.009 and 1.015}  % experiments/FIGS/data/v2/index.json :: fine-tuning qg N=1600000: R30 of each R16_Q1_MASS run over the from-scratch mean, SD over runs and standard error of that mean in quadrature :: sha256 c6c4138ee5c86423
\newcommand{\MlstFtQgRthirtyCoarseMassOverScratchPr}{at least 0.87\% more rejection}  % experiments/FIGS/data/v2/index.json :: fine-tuning qg N=1600000: R30 of each R16_Q1_MASS run over the from-scratch mean, SD over runs and standard error of that mean in quadrature :: sha256 c6c4138ee5c86423 :: prose twin: two significant figures, direction in words
\newcommand{\MlstFtQgRthirtyCoarseMassOverScratchPrMax}{up to 1.5\% more rejection}  % experiments/FIGS/data/v2/index.json :: fine-tuning qg N=1600000: R30 of each R16_Q1_MASS run over the from-scratch mean, SD over runs and standard error of that mean in quadrature :: sha256 c6c4138ee5c86423 :: prose twin of the run farther from 1
\newcommand{\MlstFtQgRthirtyResOverFull}{1.007\,\ensuremath{\pm}\,0.026}  % experiments/FIGS/data/v2/index.json :: fine-tuning qg N=1600000: R30 of L162 over L188, run by run :: sha256 c6c4138ee5c86423
\newcommand{\MlstFtQgRthirtyResOverFullPr}{equal within 3\%}  % experiments/FIGS/data/v2/index.json :: fine-tuning qg N=1600000: R30 of L162 over L188, run by run :: sha256 c6c4138ee5c86423 :: prose twin: two significant figures, direction in words
\newcommand{\MlstFtQgRthirtyMidOverFull}{1.010\,\ensuremath{\pm}\,0.029}  % experiments/FIGS/data/v2/index.json :: fine-tuning qg N=1600000: R30 of R42_Q1 over L188, run by run :: sha256 c6c4138ee5c86423
\newcommand{\MlstFtQgRthirtyMidOverFullPr}{equal within 3\%}  % experiments/FIGS/data/v2/index.json :: fine-tuning qg N=1600000: R30 of R42_Q1 over L188, run by run :: sha256 c6c4138ee5c86423 :: prose twin: two significant figures, direction in words
\newcommand{\MlstFtQgRthirtyCoarseOverFull}{0.993\,\ensuremath{\pm}\,0.026}  % experiments/FIGS/data/v2/index.json :: fine-tuning qg N=1600000: R30 of R16_Q1 over L188, run by run :: sha256 c6c4138ee5c86423
\newcommand{\MlstFtQgRthirtyCoarseOverFullPr}{equal within 3\%}  % experiments/FIGS/data/v2/index.json :: fine-tuning qg N=1600000: R30 of R16_Q1 over L188, run by run :: sha256 c6c4138ee5c86423 :: prose twin: two significant figures, direction in words
\newcommand{\MlstFtQgRthirtyResMassOverRes}{0.976 and 0.999}  % experiments/FIGS/data/v2/index.json :: fine-tuning qg N=1600000: R30 of L162_MASS over L162, run by run :: sha256 c6c4138ee5c86423
\newcommand{\MlstFtQgRthirtyResMassOverResPr}{equal within 3\%}  % experiments/FIGS/data/v2/index.json :: fine-tuning qg N=1600000: R30 of L162_MASS over L162, run by run :: sha256 c6c4138ee5c86423 :: prose twin: two significant figures, direction in words
\newcommand{\MlstFtQgRthirtyResMassOverResPrMax}{equal within 3\%}  % experiments/FIGS/data/v2/index.json :: fine-tuning qg N=1600000: R30 of L162_MASS over L162, run by run :: sha256 c6c4138ee5c86423 :: prose twin of the run farther from 1
\newcommand{\MlstFtQgRthirtyCoarseMassOverCoarse}{0.9916 and 0.9926}  % experiments/FIGS/data/v2/index.json :: fine-tuning qg N=1600000: R30 of R16_Q1_MASS over R16_Q1, run by run :: sha256 c6c4138ee5c86423
\newcommand{\MlstFtQgRthirtyCoarseMassOverCoarsePr}{equal within 1\%}  % experiments/FIGS/data/v2/index.json :: fine-tuning qg N=1600000: R30 of R16_Q1_MASS over R16_Q1, run by run :: sha256 c6c4138ee5c86423 :: prose twin: two significant figures, direction in words
\newcommand{\MlstFtQgRthirtyCoarseMassOverCoarsePrMax}{equal within 1\%}  % experiments/FIGS/data/v2/index.json :: fine-tuning qg N=1600000: R30 of R16_Q1_MASS over R16_Q1, run by run :: sha256 c6c4138ee5c86423 :: prose twin of the run farther from 1
\newcommand{\MlstFtQgRthirtyRandOverFull}{0.985\,\ensuremath{\pm}\,0.011}  % experiments/FIGS/data/v2/index.json :: fine-tuning qg N=1600000: R30 of each random partition over the L188 run of its seed :: sha256 c6c4138ee5c86423
\newcommand{\MlstFtQgRthirtyRandOverFullPr}{equal within 2\%}  % experiments/FIGS/data/v2/index.json :: fine-tuning qg N=1600000: R30 of each random partition over the L188 run of its seed :: sha256 c6c4138ee5c86423 :: prose twin: two significant figures, direction in words
\newcommand{\MlstFtQgRthirtyRandOverCoarse}{1.008\,\ensuremath{\pm}\,0.009}  % experiments/FIGS/data/v2/index.json :: fine-tuning qg N=1600000: R30 of each random partition over the R16_Q1 run of its seed :: sha256 c6c4138ee5c86423
\newcommand{\MlstFtQgRthirtyRandOverCoarsePr}{equal within 1\%}  % experiments/FIGS/data/v2/index.json :: fine-tuning qg N=1600000: R30 of each random partition over the R16_Q1 run of its seed :: sha256 c6c4138ee5c86423 :: prose twin: two significant figures, direction in words
\newcommand{\MlstFtQgRthirtyFullOverPart}{0.994\,\ensuremath{\pm}\,0.020}  % experiments/FIGS/data/v2/index.json :: fine-tuning qg N=1600000: R30 of each 188-class run over ParT-f.t.'s published value :: sha256 c6c4138ee5c86423
\newcommand{\MlstFtQgRthirtyFullOverPartPr}{equal within 3\%}  % experiments/FIGS/data/v2/index.json :: fine-tuning qg N=1600000: R30 of each 188-class run over ParT-f.t.'s published value :: sha256 c6c4138ee5c86423 :: prose twin: two significant figures, direction in words
\newcommand{\MlstFtQgAucFull}{0.9148\,\ensuremath{\pm}\,0.0001}  % experiments/FIGS/data/v2/index.json :: fine-tuning qg N=1600000: AUC of the 188-class runs :: sha256 c6c4138ee5c86423
\newcommand{\MlstFtQgAucScratch}{0.9121\,\ensuremath{\pm}\,0.0002}  % experiments/FIGS/data/v2/index.json :: fine-tuning qg N=1600000: AUC from scratch :: sha256 c6c4138ee5c86423
\newcommand{\MlstFtQgDataEquivJets}{\ensuremath{<}1{,}000}  % experiments/FIGS/data/v2/index.json :: fine-tuning qg: from-scratch training jets whose R30 equals the 188-class mean at 1000 jets (interpolated linearly in R30 against log10 N) :: sha256 c6c4138ee5c86423
\newcommand{\MlstFtQgDataEquivJetsEFour}{26{,}000}  % experiments/FIGS/data/v2/index.json :: fine-tuning qg: from-scratch training jets whose R30 equals the 188-class mean at 10000 jets (interpolated linearly in R30 against log10 N) :: sha256 c6c4138ee5c86423
\newcommand{\MlstFtQgHerwigRthirtyScratchEThree}{27.0\,\ensuremath{\pm}\,1.3}  % experiments/FIGS/data/v2/index.json :: fine-tuning qg_herwig N=1000: R30 from scratch, every seed :: sha256 c6c4138ee5c86423
\newcommand{\MlstFtQgHerwigRthirtyFullEThree}{21\,\ensuremath{\pm}\,6}  % experiments/FIGS/data/v2/index.json :: fine-tuning qg_herwig N=1000: R30, L188 :: sha256 c6c4138ee5c86423
\newcommand{\MlstFtQgHerwigRthirtyFullOverScratchEThree}{0.77\,\ensuremath{\pm}\,0.22}  % experiments/FIGS/data/v2/index.json :: fine-tuning qg_herwig N=1000: R30 of each L188 run over the from-scratch mean, SD over runs and standard error of that mean in quadrature :: sha256 c6c4138ee5c86423
\newcommand{\MlstFtQgHerwigRthirtyFullOverScratchEThreePr}{equal within 30\%}  % experiments/FIGS/data/v2/index.json :: fine-tuning qg_herwig N=1000: R30 of each L188 run over the from-scratch mean, SD over runs and standard error of that mean in quadrature :: sha256 c6c4138ee5c86423 :: prose twin: two significant figures, direction in words
\newcommand{\MlstFtQgHerwigRthirtyResEThree}{26.2\,\ensuremath{\pm}\,0.9}  % experiments/FIGS/data/v2/index.json :: fine-tuning qg_herwig N=1000: R30, L162 :: sha256 c6c4138ee5c86423
\newcommand{\MlstFtQgHerwigRthirtyResOverScratchEThree}{0.97\,\ensuremath{\pm}\,0.04}  % experiments/FIGS/data/v2/index.json :: fine-tuning qg_herwig N=1000: R30 of each L162 run over the from-scratch mean, SD over runs and standard error of that mean in quadrature :: sha256 c6c4138ee5c86423
\newcommand{\MlstFtQgHerwigRthirtyResOverScratchEThreePr}{equal within 4\%}  % experiments/FIGS/data/v2/index.json :: fine-tuning qg_herwig N=1000: R30 of each L162 run over the from-scratch mean, SD over runs and standard error of that mean in quadrature :: sha256 c6c4138ee5c86423 :: prose twin: two significant figures, direction in words
\newcommand{\MlstFtQgHerwigRthirtyMidEThree}{26.9\,\ensuremath{\pm}\,2.7}  % experiments/FIGS/data/v2/index.json :: fine-tuning qg_herwig N=1000: R30, R42_Q1 :: sha256 c6c4138ee5c86423
\newcommand{\MlstFtQgHerwigRthirtyMidOverScratchEThree}{1.00\,\ensuremath{\pm}\,0.10}  % experiments/FIGS/data/v2/index.json :: fine-tuning qg_herwig N=1000: R30 of each R42_Q1 run over the from-scratch mean, SD over runs and standard error of that mean in quadrature :: sha256 c6c4138ee5c86423
\newcommand{\MlstFtQgHerwigRthirtyMidOverScratchEThreePr}{equal within 10\%}  % experiments/FIGS/data/v2/index.json :: fine-tuning qg_herwig N=1000: R30 of each R42_Q1 run over the from-scratch mean, SD over runs and standard error of that mean in quadrature :: sha256 c6c4138ee5c86423 :: prose twin: two significant figures, direction in words
\newcommand{\MlstFtQgHerwigRthirtyCoarseEThree}{25.2\,\ensuremath{\pm}\,2.5}  % experiments/FIGS/data/v2/index.json :: fine-tuning qg_herwig N=1000: R30, R16_Q1 :: sha256 c6c4138ee5c86423
\newcommand{\MlstFtQgHerwigRthirtyCoarseOverScratchEThree}{0.94\,\ensuremath{\pm}\,0.10}  % experiments/FIGS/data/v2/index.json :: fine-tuning qg_herwig N=1000: R30 of each R16_Q1 run over the from-scratch mean, SD over runs and standard error of that mean in quadrature :: sha256 c6c4138ee5c86423
\newcommand{\MlstFtQgHerwigRthirtyCoarseOverScratchEThreePr}{equal within 10\%}  % experiments/FIGS/data/v2/index.json :: fine-tuning qg_herwig N=1000: R30 of each R16_Q1 run over the from-scratch mean, SD over runs and standard error of that mean in quadrature :: sha256 c6c4138ee5c86423 :: prose twin: two significant figures, direction in words
\newcommand{\MlstFtQgHerwigRthirtyResMassEThree}{28.0 and 29.0}  % experiments/FIGS/data/v2/index.json :: fine-tuning qg_herwig N=1000: R30, L162_MASS :: sha256 c6c4138ee5c86423
\newcommand{\MlstFtQgHerwigRthirtyResMassOverScratchEThree}{1.038 and 1.076}  % experiments/FIGS/data/v2/index.json :: fine-tuning qg_herwig N=1000: R30 of each L162_MASS run over the from-scratch mean, SD over runs and standard error of that mean in quadrature :: sha256 c6c4138ee5c86423
\newcommand{\MlstFtQgHerwigRthirtyResMassOverScratchEThreePr}{at least 3.8\% more rejection}  % experiments/FIGS/data/v2/index.json :: fine-tuning qg_herwig N=1000: R30 of each L162_MASS run over the from-scratch mean, SD over runs and standard error of that mean in quadrature :: sha256 c6c4138ee5c86423 :: prose twin: two significant figures, direction in words
\newcommand{\MlstFtQgHerwigRthirtyResMassOverScratchEThreePrMax}{up to 7.6\% more rejection}  % experiments/FIGS/data/v2/index.json :: fine-tuning qg_herwig N=1000: R30 of each L162_MASS run over the from-scratch mean, SD over runs and standard error of that mean in quadrature :: sha256 c6c4138ee5c86423 :: prose twin of the run farther from 1
\newcommand{\MlstFtQgHerwigRthirtyCoarseMassEThree}{25.8 and 27.4}  % experiments/FIGS/data/v2/index.json :: fine-tuning qg_herwig N=1000: R30, R16_Q1_MASS :: sha256 c6c4138ee5c86423
\newcommand{\MlstFtQgHerwigRthirtyCoarseMassOverScratchEThree}{0.955 and 1.015}  % experiments/FIGS/data/v2/index.json :: fine-tuning qg_herwig N=1000: R30 of each R16_Q1_MASS run over the from-scratch mean, SD over runs and standard error of that mean in quadrature :: sha256 c6c4138ee5c86423
\newcommand{\MlstFtQgHerwigRthirtyCoarseMassOverScratchEThreePr}{equal within 5\%}  % experiments/FIGS/data/v2/index.json :: fine-tuning qg_herwig N=1000: R30 of each R16_Q1_MASS run over the from-scratch mean, SD over runs and standard error of that mean in quadrature :: sha256 c6c4138ee5c86423 :: prose twin: two significant figures, direction in words
\newcommand{\MlstFtQgHerwigRthirtyCoarseMassOverScratchEThreePrMax}{equal within 5\%}  % experiments/FIGS/data/v2/index.json :: fine-tuning qg_herwig N=1000: R30 of each R16_Q1_MASS run over the from-scratch mean, SD over runs and standard error of that mean in quadrature :: sha256 c6c4138ee5c86423 :: prose twin of the run farther from 1
\newcommand{\MlstFtQgHerwigRthirtyResOverFullEThree}{1.33\,\ensuremath{\pm}\,0.33}  % experiments/FIGS/data/v2/index.json :: fine-tuning qg_herwig N=1000: R30 of L162 over L188, run by run :: sha256 c6c4138ee5c86423
\newcommand{\MlstFtQgHerwigRthirtyResOverFullEThreePr}{equal within 40\%}  % experiments/FIGS/data/v2/index.json :: fine-tuning qg_herwig N=1000: R30 of L162 over L188, run by run :: sha256 c6c4138ee5c86423 :: prose twin: two significant figures, direction in words
\newcommand{\MlstFtQgHerwigRthirtyMidOverFullEThree}{1.36\,\ensuremath{\pm}\,0.31}  % experiments/FIGS/data/v2/index.json :: fine-tuning qg_herwig N=1000: R30 of R42_Q1 over L188, run by run :: sha256 c6c4138ee5c86423
\newcommand{\MlstFtQgHerwigRthirtyMidOverFullEThreePr}{equal within 40\%}  % experiments/FIGS/data/v2/index.json :: fine-tuning qg_herwig N=1000: R30 of R42_Q1 over L188, run by run :: sha256 c6c4138ee5c86423 :: prose twin: two significant figures, direction in words
\newcommand{\MlstFtQgHerwigRthirtyCoarseOverFullEThree}{1.3\,\ensuremath{\pm}\,0.4}  % experiments/FIGS/data/v2/index.json :: fine-tuning qg_herwig N=1000: R30 of R16_Q1 over L188, run by run :: sha256 c6c4138ee5c86423
\newcommand{\MlstFtQgHerwigRthirtyCoarseOverFullEThreePr}{equal within 50\%}  % experiments/FIGS/data/v2/index.json :: fine-tuning qg_herwig N=1000: R30 of R16_Q1 over L188, run by run :: sha256 c6c4138ee5c86423 :: prose twin: two significant figures, direction in words
\newcommand{\MlstFtQgHerwigRthirtyResMassOverResEThree}{1.062 and 1.073}  % experiments/FIGS/data/v2/index.json :: fine-tuning qg_herwig N=1000: R30 of L162_MASS over L162, run by run :: sha256 c6c4138ee5c86423
\newcommand{\MlstFtQgHerwigRthirtyResMassOverResEThreePr}{equal within 8\%}  % experiments/FIGS/data/v2/index.json :: fine-tuning qg_herwig N=1000: R30 of L162_MASS over L162, run by run :: sha256 c6c4138ee5c86423 :: prose twin: two significant figures, direction in words
\newcommand{\MlstFtQgHerwigRthirtyResMassOverResEThreePrMax}{equal within 8\%}  % experiments/FIGS/data/v2/index.json :: fine-tuning qg_herwig N=1000: R30 of L162_MASS over L162, run by run :: sha256 c6c4138ee5c86423 :: prose twin of the run farther from 1
\newcommand{\MlstFtQgHerwigRthirtyCoarseMassOverCoarseEThree}{0.93 and 1.20}  % experiments/FIGS/data/v2/index.json :: fine-tuning qg_herwig N=1000: R30 of R16_Q1_MASS over R16_Q1, run by run :: sha256 c6c4138ee5c86423
\newcommand{\MlstFtQgHerwigRthirtyCoarseMassOverCoarseEThreePr}{equal within 30\%}  % experiments/FIGS/data/v2/index.json :: fine-tuning qg_herwig N=1000: R30 of R16_Q1_MASS over R16_Q1, run by run :: sha256 c6c4138ee5c86423 :: prose twin: two significant figures, direction in words
\newcommand{\MlstFtQgHerwigRthirtyCoarseMassOverCoarseEThreePrMax}{equal within 30\%}  % experiments/FIGS/data/v2/index.json :: fine-tuning qg_herwig N=1000: R30 of R16_Q1_MASS over R16_Q1, run by run :: sha256 c6c4138ee5c86423 :: prose twin of the run farther from 1
\newcommand{\MlstFtQgHerwigRthirtyRandOverFullEThree}{0.99\,\ensuremath{\pm}\,0.20}  % experiments/FIGS/data/v2/index.json :: fine-tuning qg_herwig N=1000: R30 of each random partition over the L188 run of its seed :: sha256 c6c4138ee5c86423
\newcommand{\MlstFtQgHerwigRthirtyRandOverFullEThreePr}{equal within 30\%}  % experiments/FIGS/data/v2/index.json :: fine-tuning qg_herwig N=1000: R30 of each random partition over the L188 run of its seed :: sha256 c6c4138ee5c86423 :: prose twin: two significant figures, direction in words
\newcommand{\MlstFtQgHerwigRthirtyRandOverCoarseEThree}{0.95\,\ensuremath{\pm}\,0.13}  % experiments/FIGS/data/v2/index.json :: fine-tuning qg_herwig N=1000: R30 of each random partition over the R16_Q1 run of its seed :: sha256 c6c4138ee5c86423
\newcommand{\MlstFtQgHerwigRthirtyRandOverCoarseEThreePr}{equal within 20\%}  % experiments/FIGS/data/v2/index.json :: fine-tuning qg_herwig N=1000: R30 of each random partition over the R16_Q1 run of its seed :: sha256 c6c4138ee5c86423 :: prose twin: two significant figures, direction in words
\newcommand{\MlstFtQgHerwigRthirtyScratchEFour}{31\,\ensuremath{\pm}\,7}  % experiments/FIGS/data/v2/index.json :: fine-tuning qg_herwig N=10000: R30 from scratch, every seed :: sha256 c6c4138ee5c86423
\newcommand{\MlstFtQgHerwigRthirtyFullEFour}{27.9\,\ensuremath{\pm}\,2.8}  % experiments/FIGS/data/v2/index.json :: fine-tuning qg_herwig N=10000: R30, L188 :: sha256 c6c4138ee5c86423
\newcommand{\MlstFtQgHerwigRthirtyFullOverScratchEFour}{0.89\,\ensuremath{\pm}\,0.14}  % experiments/FIGS/data/v2/index.json :: fine-tuning qg_herwig N=10000: R30 of each L188 run over the from-scratch mean, SD over runs and standard error of that mean in quadrature :: sha256 c6c4138ee5c86423
\newcommand{\MlstFtQgHerwigRthirtyFullOverScratchEFourPr}{equal within 20\%}  % experiments/FIGS/data/v2/index.json :: fine-tuning qg_herwig N=10000: R30 of each L188 run over the from-scratch mean, SD over runs and standard error of that mean in quadrature :: sha256 c6c4138ee5c86423 :: prose twin: two significant figures, direction in words
\newcommand{\MlstFtQgHerwigRthirtyResEFour}{27.4\,\ensuremath{\pm}\,0.5}  % experiments/FIGS/data/v2/index.json :: fine-tuning qg_herwig N=10000: R30, L162 :: sha256 c6c4138ee5c86423
\newcommand{\MlstFtQgHerwigRthirtyResOverScratchEFour}{0.87\,\ensuremath{\pm}\,0.11}  % experiments/FIGS/data/v2/index.json :: fine-tuning qg_herwig N=10000: R30 of each L162 run over the from-scratch mean, SD over runs and standard error of that mean in quadrature :: sha256 c6c4138ee5c86423
\newcommand{\MlstFtQgHerwigRthirtyResOverScratchEFourPr}{13\% less rejection}  % experiments/FIGS/data/v2/index.json :: fine-tuning qg_herwig N=10000: R30 of each L162 run over the from-scratch mean, SD over runs and standard error of that mean in quadrature :: sha256 c6c4138ee5c86423 :: prose twin: two significant figures, direction in words
\newcommand{\MlstFtQgHerwigRthirtyMidEFour}{28.4\,\ensuremath{\pm}\,2.1}  % experiments/FIGS/data/v2/index.json :: fine-tuning qg_herwig N=10000: R30, R42_Q1 :: sha256 c6c4138ee5c86423
\newcommand{\MlstFtQgHerwigRthirtyMidOverScratchEFour}{0.91\,\ensuremath{\pm}\,0.13}  % experiments/FIGS/data/v2/index.json :: fine-tuning qg_herwig N=10000: R30 of each R42_Q1 run over the from-scratch mean, SD over runs and standard error of that mean in quadrature :: sha256 c6c4138ee5c86423
\newcommand{\MlstFtQgHerwigRthirtyMidOverScratchEFourPr}{equal within 10\%}  % experiments/FIGS/data/v2/index.json :: fine-tuning qg_herwig N=10000: R30 of each R42_Q1 run over the from-scratch mean, SD over runs and standard error of that mean in quadrature :: sha256 c6c4138ee5c86423 :: prose twin: two significant figures, direction in words
\newcommand{\MlstFtQgHerwigRthirtyCoarseEFour}{32.7\,\ensuremath{\pm}\,1.8}  % experiments/FIGS/data/v2/index.json :: fine-tuning qg_herwig N=10000: R30, R16_Q1 :: sha256 c6c4138ee5c86423
\newcommand{\MlstFtQgHerwigRthirtyCoarseOverScratchEFour}{1.05\,\ensuremath{\pm}\,0.14}  % experiments/FIGS/data/v2/index.json :: fine-tuning qg_herwig N=10000: R30 of each R16_Q1 run over the from-scratch mean, SD over runs and standard error of that mean in quadrature :: sha256 c6c4138ee5c86423
\newcommand{\MlstFtQgHerwigRthirtyCoarseOverScratchEFourPr}{equal within 6\%}  % experiments/FIGS/data/v2/index.json :: fine-tuning qg_herwig N=10000: R30 of each R16_Q1 run over the from-scratch mean, SD over runs and standard error of that mean in quadrature :: sha256 c6c4138ee5c86423 :: prose twin: two significant figures, direction in words
\newcommand{\MlstFtQgHerwigRthirtyResMassEFour}{34.3 and 28.5}  % experiments/FIGS/data/v2/index.json :: fine-tuning qg_herwig N=10000: R30, L162_MASS :: sha256 c6c4138ee5c86423
\newcommand{\MlstFtQgHerwigRthirtyResMassOverScratchEFour}{1.09 and 0.91}  % experiments/FIGS/data/v2/index.json :: fine-tuning qg_herwig N=10000: R30 of each L162_MASS run over the from-scratch mean, SD over runs and standard error of that mean in quadrature :: sha256 c6c4138ee5c86423
\newcommand{\MlstFtQgHerwigRthirtyResMassOverScratchEFourPr}{equal within 10\%}  % experiments/FIGS/data/v2/index.json :: fine-tuning qg_herwig N=10000: R30 of each L162_MASS run over the from-scratch mean, SD over runs and standard error of that mean in quadrature :: sha256 c6c4138ee5c86423 :: prose twin: two significant figures, direction in words
\newcommand{\MlstFtQgHerwigRthirtyResMassOverScratchEFourPrMax}{equal within 10\%}  % experiments/FIGS/data/v2/index.json :: fine-tuning qg_herwig N=10000: R30 of each L162_MASS run over the from-scratch mean, SD over runs and standard error of that mean in quadrature :: sha256 c6c4138ee5c86423 :: prose twin of the run farther from 1
\newcommand{\MlstFtQgHerwigRthirtyCoarseMassEFour}{34.4 and 31.2}  % experiments/FIGS/data/v2/index.json :: fine-tuning qg_herwig N=10000: R30, R16_Q1_MASS :: sha256 c6c4138ee5c86423
\newcommand{\MlstFtQgHerwigRthirtyCoarseMassOverScratchEFour}{1.10 and 1.00}  % experiments/FIGS/data/v2/index.json :: fine-tuning qg_herwig N=10000: R30 of each R16_Q1_MASS run over the from-scratch mean, SD over runs and standard error of that mean in quadrature :: sha256 c6c4138ee5c86423
\newcommand{\MlstFtQgHerwigRthirtyCoarseMassOverScratchEFourPr}{equal within 20\%}  % experiments/FIGS/data/v2/index.json :: fine-tuning qg_herwig N=10000: R30 of each R16_Q1_MASS run over the from-scratch mean, SD over runs and standard error of that mean in quadrature :: sha256 c6c4138ee5c86423 :: prose twin: two significant figures, direction in words
\newcommand{\MlstFtQgHerwigRthirtyCoarseMassOverScratchEFourPrMax}{equal within 20\%}  % experiments/FIGS/data/v2/index.json :: fine-tuning qg_herwig N=10000: R30 of each R16_Q1_MASS run over the from-scratch mean, SD over runs and standard error of that mean in quadrature :: sha256 c6c4138ee5c86423 :: prose twin of the run farther from 1
\newcommand{\MlstFtQgHerwigRthirtyResOverFullEFour}{0.99\,\ensuremath{\pm}\,0.12}  % experiments/FIGS/data/v2/index.json :: fine-tuning qg_herwig N=10000: R30 of L162 over L188, run by run :: sha256 c6c4138ee5c86423
\newcommand{\MlstFtQgHerwigRthirtyResOverFullEFourPr}{equal within 20\%}  % experiments/FIGS/data/v2/index.json :: fine-tuning qg_herwig N=10000: R30 of L162 over L188, run by run :: sha256 c6c4138ee5c86423 :: prose twin: two significant figures, direction in words
\newcommand{\MlstFtQgHerwigRthirtyMidOverFullEFour}{1.03\,\ensuremath{\pm}\,0.18}  % experiments/FIGS/data/v2/index.json :: fine-tuning qg_herwig N=10000: R30 of R42_Q1 over L188, run by run :: sha256 c6c4138ee5c86423
\newcommand{\MlstFtQgHerwigRthirtyMidOverFullEFourPr}{equal within 20\%}  % experiments/FIGS/data/v2/index.json :: fine-tuning qg_herwig N=10000: R30 of R42_Q1 over L188, run by run :: sha256 c6c4138ee5c86423 :: prose twin: two significant figures, direction in words
\newcommand{\MlstFtQgHerwigRthirtyCoarseOverFullEFour}{1.18\,\ensuremath{\pm}\,0.08}  % experiments/FIGS/data/v2/index.json :: fine-tuning qg_herwig N=10000: R30 of R16_Q1 over L188, run by run :: sha256 c6c4138ee5c86423
\newcommand{\MlstFtQgHerwigRthirtyCoarseOverFullEFourPr}{18\% more rejection}  % experiments/FIGS/data/v2/index.json :: fine-tuning qg_herwig N=10000: R30 of R16_Q1 over L188, run by run :: sha256 c6c4138ee5c86423 :: prose twin: two significant figures, direction in words
\newcommand{\MlstFtQgHerwigRthirtyResMassOverResEFour}{1.23 and 1.03}  % experiments/FIGS/data/v2/index.json :: fine-tuning qg_herwig N=10000: R30 of L162_MASS over L162, run by run :: sha256 c6c4138ee5c86423
\newcommand{\MlstFtQgHerwigRthirtyResMassOverResEFourPr}{equal within 30\%}  % experiments/FIGS/data/v2/index.json :: fine-tuning qg_herwig N=10000: R30 of L162_MASS over L162, run by run :: sha256 c6c4138ee5c86423 :: prose twin: two significant figures, direction in words
\newcommand{\MlstFtQgHerwigRthirtyResMassOverResEFourPrMax}{equal within 30\%}  % experiments/FIGS/data/v2/index.json :: fine-tuning qg_herwig N=10000: R30 of L162_MASS over L162, run by run :: sha256 c6c4138ee5c86423 :: prose twin of the run farther from 1
\newcommand{\MlstFtQgHerwigRthirtyCoarseMassOverCoarseEFour}{1.12 and 0.91}  % experiments/FIGS/data/v2/index.json :: fine-tuning qg_herwig N=10000: R30 of R16_Q1_MASS over R16_Q1, run by run :: sha256 c6c4138ee5c86423
\newcommand{\MlstFtQgHerwigRthirtyCoarseMassOverCoarseEFourPr}{equal within 20\%}  % experiments/FIGS/data/v2/index.json :: fine-tuning qg_herwig N=10000: R30 of R16_Q1_MASS over R16_Q1, run by run :: sha256 c6c4138ee5c86423 :: prose twin: two significant figures, direction in words
\newcommand{\MlstFtQgHerwigRthirtyCoarseMassOverCoarseEFourPrMax}{equal within 20\%}  % experiments/FIGS/data/v2/index.json :: fine-tuning qg_herwig N=10000: R30 of R16_Q1_MASS over R16_Q1, run by run :: sha256 c6c4138ee5c86423 :: prose twin of the run farther from 1
\newcommand{\MlstFtQgHerwigRthirtyRandOverFullEFour}{1.13\,\ensuremath{\pm}\,0.18}  % experiments/FIGS/data/v2/index.json :: fine-tuning qg_herwig N=10000: R30 of each random partition over the L188 run of its seed :: sha256 c6c4138ee5c86423
\newcommand{\MlstFtQgHerwigRthirtyRandOverFullEFourPr}{equal within 20\%}  % experiments/FIGS/data/v2/index.json :: fine-tuning qg_herwig N=10000: R30 of each random partition over the L188 run of its seed :: sha256 c6c4138ee5c86423 :: prose twin: two significant figures, direction in words
\newcommand{\MlstFtQgHerwigRthirtyRandOverCoarseEFour}{0.93\,\ensuremath{\pm}\,0.14}  % experiments/FIGS/data/v2/index.json :: fine-tuning qg_herwig N=10000: R30 of each random partition over the R16_Q1 run of its seed :: sha256 c6c4138ee5c86423
\newcommand{\MlstFtQgHerwigRthirtyRandOverCoarseEFourPr}{equal within 20\%}  % experiments/FIGS/data/v2/index.json :: fine-tuning qg_herwig N=10000: R30 of each random partition over the R16_Q1 run of its seed :: sha256 c6c4138ee5c86423 :: prose twin: two significant figures, direction in words
\newcommand{\MlstFtQgHerwigRthirtyScratchEFive}{36.7\,\ensuremath{\pm}\,0.9}  % experiments/FIGS/data/v2/index.json :: fine-tuning qg_herwig N=100000: R30 from scratch, every seed :: sha256 c6c4138ee5c86423
\newcommand{\MlstFtQgHerwigRthirtyFullEFive}{34.3\,\ensuremath{\pm}\,1.0}  % experiments/FIGS/data/v2/index.json :: fine-tuning qg_herwig N=100000: R30, L188 :: sha256 c6c4138ee5c86423
\newcommand{\MlstFtQgHerwigRthirtyFullOverScratchEFive}{0.935\,\ensuremath{\pm}\,0.031}  % experiments/FIGS/data/v2/index.json :: fine-tuning qg_herwig N=100000: R30 of each L188 run over the from-scratch mean, SD over runs and standard error of that mean in quadrature :: sha256 c6c4138ee5c86423
\newcommand{\MlstFtQgHerwigRthirtyFullOverScratchEFivePr}{6.5\% less rejection}  % experiments/FIGS/data/v2/index.json :: fine-tuning qg_herwig N=100000: R30 of each L188 run over the from-scratch mean, SD over runs and standard error of that mean in quadrature :: sha256 c6c4138ee5c86423 :: prose twin: two significant figures, direction in words
\newcommand{\MlstFtQgHerwigRthirtyResEFive}{35.5\,\ensuremath{\pm}\,0.5}  % experiments/FIGS/data/v2/index.json :: fine-tuning qg_herwig N=100000: R30, L162 :: sha256 c6c4138ee5c86423
\newcommand{\MlstFtQgHerwigRthirtyResOverScratchEFive}{0.966\,\ensuremath{\pm}\,0.019}  % experiments/FIGS/data/v2/index.json :: fine-tuning qg_herwig N=100000: R30 of each L162 run over the from-scratch mean, SD over runs and standard error of that mean in quadrature :: sha256 c6c4138ee5c86423
\newcommand{\MlstFtQgHerwigRthirtyResOverScratchEFivePr}{3.4\% less rejection}  % experiments/FIGS/data/v2/index.json :: fine-tuning qg_herwig N=100000: R30 of each L162 run over the from-scratch mean, SD over runs and standard error of that mean in quadrature :: sha256 c6c4138ee5c86423 :: prose twin: two significant figures, direction in words
\newcommand{\MlstFtQgHerwigRthirtyMidEFive}{35.5\,\ensuremath{\pm}\,2.1}  % experiments/FIGS/data/v2/index.json :: fine-tuning qg_herwig N=100000: R30, R42_Q1 :: sha256 c6c4138ee5c86423
\newcommand{\MlstFtQgHerwigRthirtyMidOverScratchEFive}{0.97\,\ensuremath{\pm}\,0.06}  % experiments/FIGS/data/v2/index.json :: fine-tuning qg_herwig N=100000: R30 of each R42_Q1 run over the from-scratch mean, SD over runs and standard error of that mean in quadrature :: sha256 c6c4138ee5c86423
\newcommand{\MlstFtQgHerwigRthirtyMidOverScratchEFivePr}{equal within 6\%}  % experiments/FIGS/data/v2/index.json :: fine-tuning qg_herwig N=100000: R30 of each R42_Q1 run over the from-scratch mean, SD over runs and standard error of that mean in quadrature :: sha256 c6c4138ee5c86423 :: prose twin: two significant figures, direction in words
\newcommand{\MlstFtQgHerwigRthirtyCoarseEFive}{34.9\,\ensuremath{\pm}\,1.1}  % experiments/FIGS/data/v2/index.json :: fine-tuning qg_herwig N=100000: R30, R16_Q1 :: sha256 c6c4138ee5c86423
\newcommand{\MlstFtQgHerwigRthirtyCoarseOverScratchEFive}{0.950\,\ensuremath{\pm}\,0.033}  % experiments/FIGS/data/v2/index.json :: fine-tuning qg_herwig N=100000: R30 of each R16_Q1 run over the from-scratch mean, SD over runs and standard error of that mean in quadrature :: sha256 c6c4138ee5c86423
\newcommand{\MlstFtQgHerwigRthirtyCoarseOverScratchEFivePr}{equal within 6\%}  % experiments/FIGS/data/v2/index.json :: fine-tuning qg_herwig N=100000: R30 of each R16_Q1 run over the from-scratch mean, SD over runs and standard error of that mean in quadrature :: sha256 c6c4138ee5c86423 :: prose twin: two significant figures, direction in words
\newcommand{\MlstFtQgHerwigRthirtyResMassEFive}{34.0 and 36.5}  % experiments/FIGS/data/v2/index.json :: fine-tuning qg_herwig N=100000: R30, L162_MASS :: sha256 c6c4138ee5c86423
\newcommand{\MlstFtQgHerwigRthirtyResMassOverScratchEFive}{0.925 and 0.994}  % experiments/FIGS/data/v2/index.json :: fine-tuning qg_herwig N=100000: R30 of each L162_MASS run over the from-scratch mean, SD over runs and standard error of that mean in quadrature :: sha256 c6c4138ee5c86423
\newcommand{\MlstFtQgHerwigRthirtyResMassOverScratchEFivePr}{at least 0.63\% less rejection}  % experiments/FIGS/data/v2/index.json :: fine-tuning qg_herwig N=100000: R30 of each L162_MASS run over the from-scratch mean, SD over runs and standard error of that mean in quadrature :: sha256 c6c4138ee5c86423 :: prose twin: two significant figures, direction in words
\newcommand{\MlstFtQgHerwigRthirtyResMassOverScratchEFivePrMax}{up to 7.5\% less rejection}  % experiments/FIGS/data/v2/index.json :: fine-tuning qg_herwig N=100000: R30 of each L162_MASS run over the from-scratch mean, SD over runs and standard error of that mean in quadrature :: sha256 c6c4138ee5c86423 :: prose twin of the run farther from 1
\newcommand{\MlstFtQgHerwigRthirtyCoarseMassEFive}{35.9 and 39.0}  % experiments/FIGS/data/v2/index.json :: fine-tuning qg_herwig N=100000: R30, R16_Q1_MASS :: sha256 c6c4138ee5c86423
\newcommand{\MlstFtQgHerwigRthirtyCoarseMassOverScratchEFive}{0.979 and 1.064}  % experiments/FIGS/data/v2/index.json :: fine-tuning qg_herwig N=100000: R30 of each R16_Q1_MASS run over the from-scratch mean, SD over runs and standard error of that mean in quadrature :: sha256 c6c4138ee5c86423
\newcommand{\MlstFtQgHerwigRthirtyCoarseMassOverScratchEFivePr}{equal within 7\%}  % experiments/FIGS/data/v2/index.json :: fine-tuning qg_herwig N=100000: R30 of each R16_Q1_MASS run over the from-scratch mean, SD over runs and standard error of that mean in quadrature :: sha256 c6c4138ee5c86423 :: prose twin: two significant figures, direction in words
\newcommand{\MlstFtQgHerwigRthirtyCoarseMassOverScratchEFivePrMax}{equal within 7\%}  % experiments/FIGS/data/v2/index.json :: fine-tuning qg_herwig N=100000: R30 of each R16_Q1_MASS run over the from-scratch mean, SD over runs and standard error of that mean in quadrature :: sha256 c6c4138ee5c86423 :: prose twin of the run farther from 1
\newcommand{\MlstFtQgHerwigRthirtyResOverFullEFive}{1.03\,\ensuremath{\pm}\,0.04}  % experiments/FIGS/data/v2/index.json :: fine-tuning qg_herwig N=100000: R30 of L162 over L188, run by run :: sha256 c6c4138ee5c86423
\newcommand{\MlstFtQgHerwigRthirtyResOverFullEFivePr}{equal within 5\%}  % experiments/FIGS/data/v2/index.json :: fine-tuning qg_herwig N=100000: R30 of L162 over L188, run by run :: sha256 c6c4138ee5c86423 :: prose twin: two significant figures, direction in words
\newcommand{\MlstFtQgHerwigRthirtyMidOverFullEFive}{1.03\,\ensuremath{\pm}\,0.05}  % experiments/FIGS/data/v2/index.json :: fine-tuning qg_herwig N=100000: R30 of R42_Q1 over L188, run by run :: sha256 c6c4138ee5c86423
\newcommand{\MlstFtQgHerwigRthirtyMidOverFullEFivePr}{equal within 6\%}  % experiments/FIGS/data/v2/index.json :: fine-tuning qg_herwig N=100000: R30 of R42_Q1 over L188, run by run :: sha256 c6c4138ee5c86423 :: prose twin: two significant figures, direction in words
\newcommand{\MlstFtQgHerwigRthirtyCoarseOverFullEFive}{1.02\,\ensuremath{\pm}\,0.06}  % experiments/FIGS/data/v2/index.json :: fine-tuning qg_herwig N=100000: R30 of R16_Q1 over L188, run by run :: sha256 c6c4138ee5c86423
\newcommand{\MlstFtQgHerwigRthirtyCoarseOverFullEFivePr}{equal within 7\%}  % experiments/FIGS/data/v2/index.json :: fine-tuning qg_herwig N=100000: R30 of R16_Q1 over L188, run by run :: sha256 c6c4138ee5c86423 :: prose twin: two significant figures, direction in words
\newcommand{\MlstFtQgHerwigRthirtyResMassOverResEFive}{0.971 and 1.028}  % experiments/FIGS/data/v2/index.json :: fine-tuning qg_herwig N=100000: R30 of L162_MASS over L162, run by run :: sha256 c6c4138ee5c86423
\newcommand{\MlstFtQgHerwigRthirtyResMassOverResEFivePr}{equal within 3\%}  % experiments/FIGS/data/v2/index.json :: fine-tuning qg_herwig N=100000: R30 of L162_MASS over L162, run by run :: sha256 c6c4138ee5c86423 :: prose twin: two significant figures, direction in words
\newcommand{\MlstFtQgHerwigRthirtyResMassOverResEFivePrMax}{equal within 3\%}  % experiments/FIGS/data/v2/index.json :: fine-tuning qg_herwig N=100000: R30 of L162_MASS over L162, run by run :: sha256 c6c4138ee5c86423 :: prose twin of the run farther from 1
\newcommand{\MlstFtQgHerwigRthirtyCoarseMassOverCoarseEFive}{1.064 and 1.121}  % experiments/FIGS/data/v2/index.json :: fine-tuning qg_herwig N=100000: R30 of R16_Q1_MASS over R16_Q1, run by run :: sha256 c6c4138ee5c86423
\newcommand{\MlstFtQgHerwigRthirtyCoarseMassOverCoarseEFivePr}{equal within 20\%}  % experiments/FIGS/data/v2/index.json :: fine-tuning qg_herwig N=100000: R30 of R16_Q1_MASS over R16_Q1, run by run :: sha256 c6c4138ee5c86423 :: prose twin: two significant figures, direction in words
\newcommand{\MlstFtQgHerwigRthirtyCoarseMassOverCoarseEFivePrMax}{equal within 20\%}  % experiments/FIGS/data/v2/index.json :: fine-tuning qg_herwig N=100000: R30 of R16_Q1_MASS over R16_Q1, run by run :: sha256 c6c4138ee5c86423 :: prose twin of the run farther from 1
\newcommand{\MlstFtQgHerwigRthirtyRandOverFullEFive}{1.01\,\ensuremath{\pm}\,0.05}  % experiments/FIGS/data/v2/index.json :: fine-tuning qg_herwig N=100000: R30 of each random partition over the L188 run of its seed :: sha256 c6c4138ee5c86423
\newcommand{\MlstFtQgHerwigRthirtyRandOverFullEFivePr}{equal within 5\%}  % experiments/FIGS/data/v2/index.json :: fine-tuning qg_herwig N=100000: R30 of each random partition over the L188 run of its seed :: sha256 c6c4138ee5c86423 :: prose twin: two significant figures, direction in words
\newcommand{\MlstFtQgHerwigRthirtyRandOverCoarseEFive}{1.009\,\ensuremath{\pm}\,0.024}  % experiments/FIGS/data/v2/index.json :: fine-tuning qg_herwig N=100000: R30 of each random partition over the R16_Q1 run of its seed :: sha256 c6c4138ee5c86423
\newcommand{\MlstFtQgHerwigRthirtyRandOverCoarseEFivePr}{equal within 3\%}  % experiments/FIGS/data/v2/index.json :: fine-tuning qg_herwig N=100000: R30 of each random partition over the R16_Q1 run of its seed :: sha256 c6c4138ee5c86423 :: prose twin: two significant figures, direction in words
\newcommand{\MlstFtQgHerwigRthirtyScratch}{35.5\,\ensuremath{\pm}\,1.2}  % experiments/FIGS/data/v2/index.json :: fine-tuning qg_herwig N=1600000: R30 from scratch, every seed :: sha256 c6c4138ee5c86423
\newcommand{\MlstFtQgHerwigRthirtyFull}{35.2\,\ensuremath{\pm}\,2.1}  % experiments/FIGS/data/v2/index.json :: fine-tuning qg_herwig N=1600000: R30, L188 :: sha256 c6c4138ee5c86423
\newcommand{\MlstFtQgHerwigRthirtyFullOverScratch}{0.99\,\ensuremath{\pm}\,0.06}  % experiments/FIGS/data/v2/index.json :: fine-tuning qg_herwig N=1600000: R30 of each L188 run over the from-scratch mean, SD over runs and standard error of that mean in quadrature :: sha256 c6c4138ee5c86423
\newcommand{\MlstFtQgHerwigRthirtyFullOverScratchPr}{equal within 6\%}  % experiments/FIGS/data/v2/index.json :: fine-tuning qg_herwig N=1600000: R30 of each L188 run over the from-scratch mean, SD over runs and standard error of that mean in quadrature :: sha256 c6c4138ee5c86423 :: prose twin: two significant figures, direction in words
\newcommand{\MlstFtQgHerwigRthirtyRes}{34.6\,\ensuremath{\pm}\,2.0}  % experiments/FIGS/data/v2/index.json :: fine-tuning qg_herwig N=1600000: R30, L162 :: sha256 c6c4138ee5c86423
\newcommand{\MlstFtQgHerwigRthirtyResOverScratch}{0.97\,\ensuremath{\pm}\,0.06}  % experiments/FIGS/data/v2/index.json :: fine-tuning qg_herwig N=1600000: R30 of each L162 run over the from-scratch mean, SD over runs and standard error of that mean in quadrature :: sha256 c6c4138ee5c86423
\newcommand{\MlstFtQgHerwigRthirtyResOverScratchPr}{equal within 6\%}  % experiments/FIGS/data/v2/index.json :: fine-tuning qg_herwig N=1600000: R30 of each L162 run over the from-scratch mean, SD over runs and standard error of that mean in quadrature :: sha256 c6c4138ee5c86423 :: prose twin: two significant figures, direction in words
\newcommand{\MlstFtQgHerwigRthirtyMid}{36.1\,\ensuremath{\pm}\,2.7}  % experiments/FIGS/data/v2/index.json :: fine-tuning qg_herwig N=1600000: R30, R42_Q1 :: sha256 c6c4138ee5c86423
\newcommand{\MlstFtQgHerwigRthirtyMidOverScratch}{1.02\,\ensuremath{\pm}\,0.08}  % experiments/FIGS/data/v2/index.json :: fine-tuning qg_herwig N=1600000: R30 of each R42_Q1 run over the from-scratch mean, SD over runs and standard error of that mean in quadrature :: sha256 c6c4138ee5c86423
\newcommand{\MlstFtQgHerwigRthirtyMidOverScratchPr}{equal within 8\%}  % experiments/FIGS/data/v2/index.json :: fine-tuning qg_herwig N=1600000: R30 of each R42_Q1 run over the from-scratch mean, SD over runs and standard error of that mean in quadrature :: sha256 c6c4138ee5c86423 :: prose twin: two significant figures, direction in words
\newcommand{\MlstFtQgHerwigRthirtyCoarse}{36.9\,\ensuremath{\pm}\,0.9}  % experiments/FIGS/data/v2/index.json :: fine-tuning qg_herwig N=1600000: R30, R16_Q1 :: sha256 c6c4138ee5c86423
\newcommand{\MlstFtQgHerwigRthirtyCoarseOverScratch}{1.040\,\ensuremath{\pm}\,0.033}  % experiments/FIGS/data/v2/index.json :: fine-tuning qg_herwig N=1600000: R30 of each R16_Q1 run over the from-scratch mean, SD over runs and standard error of that mean in quadrature :: sha256 c6c4138ee5c86423
\newcommand{\MlstFtQgHerwigRthirtyCoarseOverScratchPr}{equal within 4\%}  % experiments/FIGS/data/v2/index.json :: fine-tuning qg_herwig N=1600000: R30 of each R16_Q1 run over the from-scratch mean, SD over runs and standard error of that mean in quadrature :: sha256 c6c4138ee5c86423 :: prose twin: two significant figures, direction in words
\newcommand{\MlstFtQgHerwigRthirtyResMass}{37.8 and 34.4}  % experiments/FIGS/data/v2/index.json :: fine-tuning qg_herwig N=1600000: R30, L162_MASS :: sha256 c6c4138ee5c86423
\newcommand{\MlstFtQgHerwigRthirtyResMassOverScratch}{1.065 and 0.970}  % experiments/FIGS/data/v2/index.json :: fine-tuning qg_herwig N=1600000: R30 of each L162_MASS run over the from-scratch mean, SD over runs and standard error of that mean in quadrature :: sha256 c6c4138ee5c86423
\newcommand{\MlstFtQgHerwigRthirtyResMassOverScratchPr}{equal within 7\%}  % experiments/FIGS/data/v2/index.json :: fine-tuning qg_herwig N=1600000: R30 of each L162_MASS run over the from-scratch mean, SD over runs and standard error of that mean in quadrature :: sha256 c6c4138ee5c86423 :: prose twin: two significant figures, direction in words
\newcommand{\MlstFtQgHerwigRthirtyResMassOverScratchPrMax}{equal within 7\%}  % experiments/FIGS/data/v2/index.json :: fine-tuning qg_herwig N=1600000: R30 of each L162_MASS run over the from-scratch mean, SD over runs and standard error of that mean in quadrature :: sha256 c6c4138ee5c86423 :: prose twin of the run farther from 1
\newcommand{\MlstFtQgHerwigRthirtyCoarseMass}{37.7 and 32.9}  % experiments/FIGS/data/v2/index.json :: fine-tuning qg_herwig N=1600000: R30, R16_Q1_MASS :: sha256 c6c4138ee5c86423
\newcommand{\MlstFtQgHerwigRthirtyCoarseMassOverScratch}{1.06 and 0.93}  % experiments/FIGS/data/v2/index.json :: fine-tuning qg_herwig N=1600000: R30 of each R16_Q1_MASS run over the from-scratch mean, SD over runs and standard error of that mean in quadrature :: sha256 c6c4138ee5c86423
\newcommand{\MlstFtQgHerwigRthirtyCoarseMassOverScratchPr}{equal within 8\%}  % experiments/FIGS/data/v2/index.json :: fine-tuning qg_herwig N=1600000: R30 of each R16_Q1_MASS run over the from-scratch mean, SD over runs and standard error of that mean in quadrature :: sha256 c6c4138ee5c86423 :: prose twin: two significant figures, direction in words
\newcommand{\MlstFtQgHerwigRthirtyCoarseMassOverScratchPrMax}{equal within 8\%}  % experiments/FIGS/data/v2/index.json :: fine-tuning qg_herwig N=1600000: R30 of each R16_Q1_MASS run over the from-scratch mean, SD over runs and standard error of that mean in quadrature :: sha256 c6c4138ee5c86423 :: prose twin of the run farther from 1
\newcommand{\MlstFtQgHerwigRthirtyResOverFull}{0.9832\,\ensuremath{\pm}\,0.0022}  % experiments/FIGS/data/v2/index.json :: fine-tuning qg_herwig N=1600000: R30 of L162 over L188, run by run :: sha256 c6c4138ee5c86423
\newcommand{\MlstFtQgHerwigRthirtyResOverFullPr}{1.7\% less rejection}  % experiments/FIGS/data/v2/index.json :: fine-tuning qg_herwig N=1600000: R30 of L162 over L188, run by run :: sha256 c6c4138ee5c86423 :: prose twin: two significant figures, direction in words
\newcommand{\MlstFtQgHerwigRthirtyMidOverFull}{1.03\,\ensuremath{\pm}\,0.05}  % experiments/FIGS/data/v2/index.json :: fine-tuning qg_herwig N=1600000: R30 of R42_Q1 over L188, run by run :: sha256 c6c4138ee5c86423
\newcommand{\MlstFtQgHerwigRthirtyMidOverFullPr}{equal within 6\%}  % experiments/FIGS/data/v2/index.json :: fine-tuning qg_herwig N=1600000: R30 of R42_Q1 over L188, run by run :: sha256 c6c4138ee5c86423 :: prose twin: two significant figures, direction in words
\newcommand{\MlstFtQgHerwigRthirtyCoarseOverFull}{1.05\,\ensuremath{\pm}\,0.06}  % experiments/FIGS/data/v2/index.json :: fine-tuning qg_herwig N=1600000: R30 of R16_Q1 over L188, run by run :: sha256 c6c4138ee5c86423
\newcommand{\MlstFtQgHerwigRthirtyCoarseOverFullPr}{equal within 6\%}  % experiments/FIGS/data/v2/index.json :: fine-tuning qg_herwig N=1600000: R30 of R16_Q1 over L188, run by run :: sha256 c6c4138ee5c86423 :: prose twin: two significant figures, direction in words
\newcommand{\MlstFtQgHerwigRthirtyResMassOverRes}{1.14 and 0.93}  % experiments/FIGS/data/v2/index.json :: fine-tuning qg_herwig N=1600000: R30 of L162_MASS over L162, run by run :: sha256 c6c4138ee5c86423
\newcommand{\MlstFtQgHerwigRthirtyResMassOverResPr}{equal within 20\%}  % experiments/FIGS/data/v2/index.json :: fine-tuning qg_herwig N=1600000: R30 of L162_MASS over L162, run by run :: sha256 c6c4138ee5c86423 :: prose twin: two significant figures, direction in words
\newcommand{\MlstFtQgHerwigRthirtyResMassOverResPrMax}{equal within 20\%}  % experiments/FIGS/data/v2/index.json :: fine-tuning qg_herwig N=1600000: R30 of L162_MASS over L162, run by run :: sha256 c6c4138ee5c86423 :: prose twin of the run farther from 1
\newcommand{\MlstFtQgHerwigRthirtyCoarseMassOverCoarse}{1.01 and 0.88}  % experiments/FIGS/data/v2/index.json :: fine-tuning qg_herwig N=1600000: R30 of R16_Q1_MASS over R16_Q1, run by run :: sha256 c6c4138ee5c86423
\newcommand{\MlstFtQgHerwigRthirtyCoarseMassOverCoarsePr}{equal within 20\%}  % experiments/FIGS/data/v2/index.json :: fine-tuning qg_herwig N=1600000: R30 of R16_Q1_MASS over R16_Q1, run by run :: sha256 c6c4138ee5c86423 :: prose twin: two significant figures, direction in words
\newcommand{\MlstFtQgHerwigRthirtyCoarseMassOverCoarsePrMax}{equal within 20\%}  % experiments/FIGS/data/v2/index.json :: fine-tuning qg_herwig N=1600000: R30 of R16_Q1_MASS over R16_Q1, run by run :: sha256 c6c4138ee5c86423 :: prose twin of the run farther from 1
\newcommand{\MlstFtQgHerwigRthirtyRandOverFull}{0.98\,\ensuremath{\pm}\,0.04}  % experiments/FIGS/data/v2/index.json :: fine-tuning qg_herwig N=1600000: R30 of each random partition over the L188 run of its seed :: sha256 c6c4138ee5c86423
\newcommand{\MlstFtQgHerwigRthirtyRandOverFullPr}{equal within 5\%}  % experiments/FIGS/data/v2/index.json :: fine-tuning qg_herwig N=1600000: R30 of each random partition over the L188 run of its seed :: sha256 c6c4138ee5c86423 :: prose twin: two significant figures, direction in words
\newcommand{\MlstFtQgHerwigRthirtyRandOverCoarse}{0.94\,\ensuremath{\pm}\,0.04}  % experiments/FIGS/data/v2/index.json :: fine-tuning qg_herwig N=1600000: R30 of each random partition over the R16_Q1 run of its seed :: sha256 c6c4138ee5c86423
\newcommand{\MlstFtQgHerwigRthirtyRandOverCoarsePr}{equal within 6\%}  % experiments/FIGS/data/v2/index.json :: fine-tuning qg_herwig N=1600000: R30 of each random partition over the R16_Q1 run of its seed :: sha256 c6c4138ee5c86423 :: prose twin: two significant figures, direction in words
\newcommand{\MlstFtQgRthirtyHerwigOverPythiaFull}{0.328\,\ensuremath{\pm}\,0.015}  % experiments/FIGS/data/v2/index.json :: fine-tuning q/g N=1600000: R30 on Herwig over R30 on Pythia, each run of L188 :: sha256 c6c4138ee5c86423
\newcommand{\MlstFtQgRthirtyHerwigOverPythiaFullPr}{0.33 times the rejection}  % experiments/FIGS/data/v2/index.json :: fine-tuning q/g N=1600000: R30 on Herwig over R30 on Pythia, each run of L188 :: sha256 c6c4138ee5c86423 :: prose twin: two significant figures, direction in words
\newcommand{\MlstFtQgRthirtyHerwigOverPythiaRes}{0.320\,\ensuremath{\pm}\,0.023}  % experiments/FIGS/data/v2/index.json :: fine-tuning q/g N=1600000: R30 on Herwig over R30 on Pythia, each run of L162 :: sha256 c6c4138ee5c86423
\newcommand{\MlstFtQgRthirtyHerwigOverPythiaResPr}{0.32 times the rejection}  % experiments/FIGS/data/v2/index.json :: fine-tuning q/g N=1600000: R30 on Herwig over R30 on Pythia, each run of L162 :: sha256 c6c4138ee5c86423 :: prose twin: two significant figures, direction in words
\newcommand{\MlstFtQgRthirtyHerwigOverPythiaMid}{0.334\,\ensuremath{\pm}\,0.023}  % experiments/FIGS/data/v2/index.json :: fine-tuning q/g N=1600000: R30 on Herwig over R30 on Pythia, each run of R42_Q1 :: sha256 c6c4138ee5c86423
\newcommand{\MlstFtQgRthirtyHerwigOverPythiaMidPr}{0.33 times the rejection}  % experiments/FIGS/data/v2/index.json :: fine-tuning q/g N=1600000: R30 on Herwig over R30 on Pythia, each run of R42_Q1 :: sha256 c6c4138ee5c86423 :: prose twin: two significant figures, direction in words
\newcommand{\MlstFtQgRthirtyHerwigOverPythiaCoarse}{0.346\,\ensuremath{\pm}\,0.011}  % experiments/FIGS/data/v2/index.json :: fine-tuning q/g N=1600000: R30 on Herwig over R30 on Pythia, each run of R16_Q1 :: sha256 c6c4138ee5c86423
\newcommand{\MlstFtQgRthirtyHerwigOverPythiaCoarsePr}{0.35 times the rejection}  % experiments/FIGS/data/v2/index.json :: fine-tuning q/g N=1600000: R30 on Herwig over R30 on Pythia, each run of R16_Q1 :: sha256 c6c4138ee5c86423 :: prose twin: two significant figures, direction in words
\newcommand{\MlstFtQgRthirtyHerwigOverPythiaResMass}{0.353 and 0.323}  % experiments/FIGS/data/v2/index.json :: fine-tuning q/g N=1600000: R30 on Herwig over R30 on Pythia, each run of L162_MASS :: sha256 c6c4138ee5c86423
\newcommand{\MlstFtQgRthirtyHerwigOverPythiaResMassPr}{at most 0.35 times the rejection}  % experiments/FIGS/data/v2/index.json :: fine-tuning q/g N=1600000: R30 on Herwig over R30 on Pythia, each run of L162_MASS :: sha256 c6c4138ee5c86423 :: prose twin: two significant figures, direction in words
\newcommand{\MlstFtQgRthirtyHerwigOverPythiaResMassPrMax}{down to 0.32 times the rejection}  % experiments/FIGS/data/v2/index.json :: fine-tuning q/g N=1600000: R30 on Herwig over R30 on Pythia, each run of L162_MASS :: sha256 c6c4138ee5c86423 :: prose twin of the run farther from 1
\newcommand{\MlstFtQgRthirtyHerwigOverPythiaCoarseMass}{0.359 and 0.311}  % experiments/FIGS/data/v2/index.json :: fine-tuning q/g N=1600000: R30 on Herwig over R30 on Pythia, each run of R16_Q1_MASS :: sha256 c6c4138ee5c86423
\newcommand{\MlstFtQgRthirtyHerwigOverPythiaCoarseMassPr}{at most 0.36 times the rejection}  % experiments/FIGS/data/v2/index.json :: fine-tuning q/g N=1600000: R30 on Herwig over R30 on Pythia, each run of R16_Q1_MASS :: sha256 c6c4138ee5c86423 :: prose twin: two significant figures, direction in words
\newcommand{\MlstFtQgRthirtyHerwigOverPythiaCoarseMassPrMax}{down to 0.31 times the rejection}  % experiments/FIGS/data/v2/index.json :: fine-tuning q/g N=1600000: R30 on Herwig over R30 on Pythia, each run of R16_Q1_MASS :: sha256 c6c4138ee5c86423 :: prose twin of the run farther from 1
\newcommand{\MlstFtQgRthirtyHerwigOverPythiaScratch}{0.341\,\ensuremath{\pm}\,0.009}  % experiments/FIGS/data/v2/index.json :: fine-tuning q/g N=1600000: R30 on Herwig over R30 on Pythia, each from-scratch seed :: sha256 c6c4138ee5c86423
\newcommand{\MlstFtQgRthirtyHerwigOverPythiaScratchPr}{0.34 times the rejection}  % experiments/FIGS/data/v2/index.json :: fine-tuning q/g N=1600000: R30 on Herwig over R30 on Pythia, each from-scratch seed :: sha256 c6c4138ee5c86423 :: prose twin: two significant figures, direction in words
\newcommand{\MlstTopCheckFullLong}{\pending{recipe check}}  % experiments/FIGS/data/v2/index.json :: the recipe check is not read out yet :: sha256 c6c4138ee5c86423
\newcommand{\MlstTopCheckScratchLong}{\pending{recipe check}}  % experiments/FIGS/data/v2/index.json :: the recipe check is not read out yet :: sha256 c6c4138ee5c86423
\newcommand{\MlstTopCheckFullShort}{\pending{recipe check}}  % experiments/FIGS/data/v2/index.json :: the recipe check is not read out yet :: sha256 c6c4138ee5c86423
\newcommand{\MlstTopCheckScratchShort}{\pending{recipe check}}  % experiments/FIGS/data/v2/index.json :: the recipe check is not read out yet :: sha256 c6c4138ee5c86423
\newcommand{\MlstLpJcOneGapFracEFour}{56.3\,\ensuremath{\pm}\,1.1\%}  % experiments/FIGS/data/v2/index.json :: jc1 N=10000: (probe - scratch) / (fine-tuned - scratch) accuracy, 188-class runs, per cent :: sha256 c6c4138ee5c86423
\newcommand{\MlstLpJcOneGapFracEFourPr}{56\%}  % experiments/FIGS/data/v2/index.json :: jc1 N=10000: (probe - scratch) / (fine-tuned - scratch) accuracy, 188-class runs, per cent :: sha256 c6c4138ee5c86423 :: prose twin, two significant figures
\newcommand{\MlstLpJcOneGapFracEFive}{\ensuremath{-}15\,\ensuremath{\pm}\,5\%}  % experiments/FIGS/data/v2/index.json :: jc1 N=100000: (probe - scratch) / (fine-tuned - scratch) accuracy, 188-class runs, per cent :: sha256 c6c4138ee5c86423
\newcommand{\MlstLpJcOneGapFracEFivePr}{\ensuremath{-}15\%}  % experiments/FIGS/data/v2/index.json :: jc1 N=100000: (probe - scratch) / (fine-tuned - scratch) accuracy, 188-class runs, per cent :: sha256 c6c4138ee5c86423 :: prose twin, two significant figures
\newcommand{\MlstLpJcTwoGapFracEFour}{96.2\,\ensuremath{\pm}\,0.6\%}  % experiments/FIGS/data/v2/index.json :: jc2_188 N=10000: (probe - scratch) / (fine-tuned - scratch) accuracy, 188-class runs, per cent :: sha256 c6c4138ee5c86423
\newcommand{\MlstLpJcTwoGapFracEFourPr}{96\%}  % experiments/FIGS/data/v2/index.json :: jc2_188 N=10000: (probe - scratch) / (fine-tuned - scratch) accuracy, 188-class runs, per cent :: sha256 c6c4138ee5c86423 :: prose twin, two significant figures
\newcommand{\MlstLpJcTwoGapFracEFive}{102.2\,\ensuremath{\pm}\,1.0\%}  % experiments/FIGS/data/v2/index.json :: jc2_188 N=100000: (probe - scratch) / (fine-tuned - scratch) accuracy, 188-class runs, per cent :: sha256 c6c4138ee5c86423
\newcommand{\MlstLpJcTwoGapFracEFivePr}{100\%}  % experiments/FIGS/data/v2/index.json :: jc2_188 N=100000: (probe - scratch) / (fine-tuned - scratch) accuracy, 188-class runs, per cent :: sha256 c6c4138ee5c86423 :: prose twin, two significant figures
\newcommand{\XbxMx}{\pending{samples not in}}  % experiments/XBS/xbx_report.py :: pending :: sha256 cb6dffe774ff5b1e
\newcommand{\XbxWindowLo}{\pending{samples not in}}  % experiments/XBS/xbx_report.py :: pending :: sha256 cb6dffe774ff5b1e
\newcommand{\XbxWindowHi}{\pending{samples not in}}  % experiments/XBS/xbx_report.py :: pending :: sha256 cb6dffe774ff5b1e
\newcommand{\XbbNTrainPerClass}{\pending{samples not in}}  % experiments/XBS/xbx_report.py :: pending :: sha256 cb6dffe774ff5b1e
\newcommand{\XbbRepetition}{\pending{samples not in}}  % experiments/XBS/xbx_report.py :: pending :: sha256 cb6dffe774ff5b1e
\newcommand{\XbbNEvalSignal}{\pending{samples not in}}  % experiments/XBS/xbx_report.py :: pending :: sha256 cb6dffe774ff5b1e
\newcommand{\XbbNEvalQcd}{\pending{samples not in}}  % experiments/XBS/xbx_report.py :: pending :: sha256 cb6dffe774ff5b1e
\newcommand{\XbsNTrainPerClass}{\pending{samples not in}}  % experiments/XBS/xbx_report.py :: pending :: sha256 cb6dffe774ff5b1e
\newcommand{\XbsRepetition}{\pending{samples not in}}  % experiments/XBS/xbx_report.py :: pending :: sha256 cb6dffe774ff5b1e
\newcommand{\XbsNEvalSignal}{\pending{samples not in}}  % experiments/XBS/xbx_report.py :: pending :: sha256 cb6dffe774ff5b1e
\newcommand{\XbsNEvalQcd}{\pending{samples not in}}  % experiments/XBS/xbx_report.py :: pending :: sha256 cb6dffe774ff5b1e
\newcommand{\MlstPartreproRthirtyPart}{1{,}711}  % experiments/FIGS/data/v2/partrepro/part.json :: r30 :: sha256 0d44b997af9e6448 :: ParT reproduction on top, part: R30 (one run)
\newcommand{\MlstPartreproRthirtyHead}{1{,}577}  % experiments/FIGS/data/v2/partrepro/head.json :: r30 :: sha256 37b82ba0a0fe9299 :: ParT reproduction on top, head: R30 (one run)
\newcommand{\MlstPartreproRthirtyRescale}{1{,}224}  % experiments/FIGS/data/v2/partrepro/rescale.json :: r30 :: sha256 0a9f1401047d07ff :: ParT reproduction on top, rescale: R30 (one run)
\newcommand{\MlstPartreproRthirtyHeadOverPart}{0.922}  % experiments/FIGS/data/v2/partrepro/head.json :: r30 / part r30 :: sha256 37b82ba0a0fe9299
\newcommand{\MlstPartreproRthirtyHeadOverPartPr}{7.8\% less rejection}  % experiments/FIGS/data/v2/partrepro/head.json :: r30 / part r30 :: sha256 37b82ba0a0fe9299 :: prose twin: two significant figures, direction in words
\newcommand{\MlstPartreproRthirtyRescaleOverPart}{0.715}  % experiments/FIGS/data/v2/partrepro/rescale.json :: r30 / part r30 :: sha256 0a9f1401047d07ff
\newcommand{\MlstPartreproRthirtyRescaleOverPartPr}{28\% less rejection}  % experiments/FIGS/data/v2/partrepro/rescale.json :: r30 / part r30 :: sha256 0a9f1401047d07ff :: prose twin: two significant figures, direction in words
\newcommand{\MlstPartreproSliceDropRatio}{\pending{jet-pT slice test not run}}  % experiments/FIGS/data/v2/index.json :: open item O10: R30 in narrow pT slices of partrepro part and rescale :: sha256 c6c4138ee5c86423
\newcommand{\MlstInputBnZeroSigmaMax}{0.96}  % experiments/FIGS/data/v2/head/input_bn_offsets.json :: offsets :: sha256 d6b821509fccecff :: largest |offset| over the five type flags
\newcommand{\MlstFtJcOneNSignal}{nine}  % experiments/FIGS/data/v2/curves/full_budget_long_leg2.json :: JC1_COLS (count) :: sha256 588e935996d1006a :: JetClass signal classes
\newcommand{\MlstFtJcOneRejHbbResOverFull}{0.92\,\ensuremath{\pm}\,0.05}  % experiments/FIGS/data/v2/curves/full_budget_long_leg2.json :: cells.<run>.N1000000.s1.jc1.per_class.label_Hbb: rejection of L162 over L188, run by run :: sha256 588e935996d1006a
\newcommand{\MlstFtJcOneRejHbbResOverFullPr}{equal within 9\%}  % experiments/FIGS/data/v2/curves/full_budget_long_leg2.json :: cells.<run>.N1000000.s1.jc1.per_class.label_Hbb: rejection of L162 over L188, run by run :: sha256 588e935996d1006a :: prose twin: two significant figures, direction in words
\newcommand{\MlstFtJcOneRejHbbMidOverFull}{0.98\,\ensuremath{\pm}\,0.12}  % experiments/FIGS/data/v2/curves/full_budget_long_leg2.json :: cells.<run>.N1000000.s1.jc1.per_class.label_Hbb: rejection of R42_Q1 over L188, run by run :: sha256 588e935996d1006a
\newcommand{\MlstFtJcOneRejHbbMidOverFullPr}{equal within 20\%}  % experiments/FIGS/data/v2/curves/full_budget_long_leg2.json :: cells.<run>.N1000000.s1.jc1.per_class.label_Hbb: rejection of R42_Q1 over L188, run by run :: sha256 588e935996d1006a :: prose twin: two significant figures, direction in words
\newcommand{\MlstFtJcOneRejHbbCoarseOverFull}{1.01\,\ensuremath{\pm}\,0.07}  % experiments/FIGS/data/v2/curves/full_budget_long_leg2.json :: cells.<run>.N1000000.s1.jc1.per_class.label_Hbb: rejection of R16_Q1 over L188, run by run :: sha256 588e935996d1006a
\newcommand{\MlstFtJcOneRejHbbCoarseOverFullPr}{equal within 7\%}  % experiments/FIGS/data/v2/curves/full_budget_long_leg2.json :: cells.<run>.N1000000.s1.jc1.per_class.label_Hbb: rejection of R16_Q1 over L188, run by run :: sha256 588e935996d1006a :: prose twin: two significant figures, direction in words
\newcommand{\MlstFtJcOneRejHbbResMassOverRes}{1.065 and 1.100}  % experiments/FIGS/data/v2/curves/full_budget_long_leg2.json :: cells.<run>.N1000000.s1.jc1.per_class.label_Hbb: rejection of L162_MASS over L162, run by run :: sha256 588e935996d1006a
\newcommand{\MlstFtJcOneRejHbbResMassOverResPr}{equal within 10\%}  % experiments/FIGS/data/v2/curves/full_budget_long_leg2.json :: cells.<run>.N1000000.s1.jc1.per_class.label_Hbb: rejection of L162_MASS over L162, run by run :: sha256 588e935996d1006a :: prose twin: two significant figures, direction in words
\newcommand{\MlstFtJcOneRejHbbResMassOverResPrMax}{equal within 10\%}  % experiments/FIGS/data/v2/curves/full_budget_long_leg2.json :: cells.<run>.N1000000.s1.jc1.per_class.label_Hbb: rejection of L162_MASS over L162, run by run :: sha256 588e935996d1006a :: prose twin of the run farther from 1
\newcommand{\MlstFtJcOneRejHbbCoarseMassOverCoarse}{0.97 and 0.86}  % experiments/FIGS/data/v2/curves/full_budget_long_leg2.json :: cells.<run>.N1000000.s1.jc1.per_class.label_Hbb: rejection of R16_Q1_MASS over R16_Q1, run by run :: sha256 588e935996d1006a
\newcommand{\MlstFtJcOneRejHbbCoarseMassOverCoarsePr}{equal within 20\%}  % experiments/FIGS/data/v2/curves/full_budget_long_leg2.json :: cells.<run>.N1000000.s1.jc1.per_class.label_Hbb: rejection of R16_Q1_MASS over R16_Q1, run by run :: sha256 588e935996d1006a :: prose twin: two significant figures, direction in words
\newcommand{\MlstFtJcOneRejHbbCoarseMassOverCoarsePrMax}{equal within 20\%}  % experiments/FIGS/data/v2/curves/full_budget_long_leg2.json :: cells.<run>.N1000000.s1.jc1.per_class.label_Hbb: rejection of R16_Q1_MASS over R16_Q1, run by run :: sha256 588e935996d1006a :: prose twin of the run farther from 1
\newcommand{\MlstFtJcOneRejHbbRandOverFull}{0.92\,\ensuremath{\pm}\,0.11}  % experiments/FIGS/data/v2/curves/full_budget_long_leg2.json :: per_class.label_Hbb: rejection of each random partition over the 188-class run of its seed :: sha256 588e935996d1006a
\newcommand{\MlstFtJcOneRejHbbRandOverFullPr}{equal within 20\%}  % experiments/FIGS/data/v2/curves/full_budget_long_leg2.json :: per_class.label_Hbb: rejection of each random partition over the 188-class run of its seed :: sha256 588e935996d1006a :: prose twin: two significant figures, direction in words
\newcommand{\MlstFtJcOneRejHbbRandOverFullRange}{0.84--1.1}  % experiments/FIGS/data/v2/curves/full_budget_long_leg2.json :: per_class.label_Hbb: smallest and largest of the partitions' ratios :: sha256 588e935996d1006a
\newcommand{\MlstFtJcOneRejHccResOverFull}{1.07\,\ensuremath{\pm}\,0.12}  % experiments/FIGS/data/v2/curves/full_budget_long_leg2.json :: cells.<run>.N1000000.s1.jc1.per_class.label_Hcc: rejection of L162 over L188, run by run :: sha256 588e935996d1006a
\newcommand{\MlstFtJcOneRejHccResOverFullPr}{equal within 20\%}  % experiments/FIGS/data/v2/curves/full_budget_long_leg2.json :: cells.<run>.N1000000.s1.jc1.per_class.label_Hcc: rejection of L162 over L188, run by run :: sha256 588e935996d1006a :: prose twin: two significant figures, direction in words
\newcommand{\MlstFtJcOneRejHccMidOverFull}{0.99\,\ensuremath{\pm}\,0.10}  % experiments/FIGS/data/v2/curves/full_budget_long_leg2.json :: cells.<run>.N1000000.s1.jc1.per_class.label_Hcc: rejection of R42_Q1 over L188, run by run :: sha256 588e935996d1006a
\newcommand{\MlstFtJcOneRejHccMidOverFullPr}{equal within 20\%}  % experiments/FIGS/data/v2/curves/full_budget_long_leg2.json :: cells.<run>.N1000000.s1.jc1.per_class.label_Hcc: rejection of R42_Q1 over L188, run by run :: sha256 588e935996d1006a :: prose twin: two significant figures, direction in words
\newcommand{\MlstFtJcOneRejHccCoarseOverFull}{0.91\,\ensuremath{\pm}\,0.07}  % experiments/FIGS/data/v2/curves/full_budget_long_leg2.json :: cells.<run>.N1000000.s1.jc1.per_class.label_Hcc: rejection of R16_Q1 over L188, run by run :: sha256 588e935996d1006a
\newcommand{\MlstFtJcOneRejHccCoarseOverFullPr}{equal within 10\%}  % experiments/FIGS/data/v2/curves/full_budget_long_leg2.json :: cells.<run>.N1000000.s1.jc1.per_class.label_Hcc: rejection of R16_Q1 over L188, run by run :: sha256 588e935996d1006a :: prose twin: two significant figures, direction in words
\newcommand{\MlstFtJcOneRejHccResMassOverRes}{0.886 and 0.816}  % experiments/FIGS/data/v2/curves/full_budget_long_leg2.json :: cells.<run>.N1000000.s1.jc1.per_class.label_Hcc: rejection of L162_MASS over L162, run by run :: sha256 588e935996d1006a
\newcommand{\MlstFtJcOneRejHccResMassOverResPr}{equal within 20\%}  % experiments/FIGS/data/v2/curves/full_budget_long_leg2.json :: cells.<run>.N1000000.s1.jc1.per_class.label_Hcc: rejection of L162_MASS over L162, run by run :: sha256 588e935996d1006a :: prose twin: two significant figures, direction in words
\newcommand{\MlstFtJcOneRejHccResMassOverResPrMax}{equal within 20\%}  % experiments/FIGS/data/v2/curves/full_budget_long_leg2.json :: cells.<run>.N1000000.s1.jc1.per_class.label_Hcc: rejection of L162_MASS over L162, run by run :: sha256 588e935996d1006a :: prose twin of the run farther from 1
\newcommand{\MlstFtJcOneRejHccCoarseMassOverCoarse}{0.98 and 0.85}  % experiments/FIGS/data/v2/curves/full_budget_long_leg2.json :: cells.<run>.N1000000.s1.jc1.per_class.label_Hcc: rejection of R16_Q1_MASS over R16_Q1, run by run :: sha256 588e935996d1006a
\newcommand{\MlstFtJcOneRejHccCoarseMassOverCoarsePr}{equal within 20\%}  % experiments/FIGS/data/v2/curves/full_budget_long_leg2.json :: cells.<run>.N1000000.s1.jc1.per_class.label_Hcc: rejection of R16_Q1_MASS over R16_Q1, run by run :: sha256 588e935996d1006a :: prose twin: two significant figures, direction in words
\newcommand{\MlstFtJcOneRejHccCoarseMassOverCoarsePrMax}{equal within 20\%}  % experiments/FIGS/data/v2/curves/full_budget_long_leg2.json :: cells.<run>.N1000000.s1.jc1.per_class.label_Hcc: rejection of R16_Q1_MASS over R16_Q1, run by run :: sha256 588e935996d1006a :: prose twin of the run farther from 1
\newcommand{\MlstFtJcOneRejHccRandOverFull}{1.03\,\ensuremath{\pm}\,0.07}  % experiments/FIGS/data/v2/curves/full_budget_long_leg2.json :: per_class.label_Hcc: rejection of each random partition over the 188-class run of its seed :: sha256 588e935996d1006a
\newcommand{\MlstFtJcOneRejHccRandOverFullPr}{equal within 8\%}  % experiments/FIGS/data/v2/curves/full_budget_long_leg2.json :: per_class.label_Hcc: rejection of each random partition over the 188-class run of its seed :: sha256 588e935996d1006a :: prose twin: two significant figures, direction in words
\newcommand{\MlstFtJcOneRejHccRandOverFullRange}{0.94--1.1}  % experiments/FIGS/data/v2/curves/full_budget_long_leg2.json :: per_class.label_Hcc: smallest and largest of the partitions' ratios :: sha256 588e935996d1006a
\newcommand{\MlstFtJcOneRejHggResOverFull}{0.98\,\ensuremath{\pm}\,0.05}  % experiments/FIGS/data/v2/curves/full_budget_long_leg2.json :: cells.<run>.N1000000.s1.jc1.per_class.label_Hgg: rejection of L162 over L188, run by run :: sha256 588e935996d1006a
\newcommand{\MlstFtJcOneRejHggResOverFullPr}{equal within 6\%}  % experiments/FIGS/data/v2/curves/full_budget_long_leg2.json :: cells.<run>.N1000000.s1.jc1.per_class.label_Hgg: rejection of L162 over L188, run by run :: sha256 588e935996d1006a :: prose twin: two significant figures, direction in words
\newcommand{\MlstFtJcOneRejHggMidOverFull}{1.000\,\ensuremath{\pm}\,0.025}  % experiments/FIGS/data/v2/curves/full_budget_long_leg2.json :: cells.<run>.N1000000.s1.jc1.per_class.label_Hgg: rejection of R42_Q1 over L188, run by run :: sha256 588e935996d1006a
\newcommand{\MlstFtJcOneRejHggMidOverFullPr}{equal within 3\%}  % experiments/FIGS/data/v2/curves/full_budget_long_leg2.json :: cells.<run>.N1000000.s1.jc1.per_class.label_Hgg: rejection of R42_Q1 over L188, run by run :: sha256 588e935996d1006a :: prose twin: two significant figures, direction in words
\newcommand{\MlstFtJcOneRejHggCoarseOverFull}{0.97\,\ensuremath{\pm}\,0.09}  % experiments/FIGS/data/v2/curves/full_budget_long_leg2.json :: cells.<run>.N1000000.s1.jc1.per_class.label_Hgg: rejection of R16_Q1 over L188, run by run :: sha256 588e935996d1006a
\newcommand{\MlstFtJcOneRejHggCoarseOverFullPr}{equal within 10\%}  % experiments/FIGS/data/v2/curves/full_budget_long_leg2.json :: cells.<run>.N1000000.s1.jc1.per_class.label_Hgg: rejection of R16_Q1 over L188, run by run :: sha256 588e935996d1006a :: prose twin: two significant figures, direction in words
\newcommand{\MlstFtJcOneRejHggResMassOverRes}{1.046 and 1.008}  % experiments/FIGS/data/v2/curves/full_budget_long_leg2.json :: cells.<run>.N1000000.s1.jc1.per_class.label_Hgg: rejection of L162_MASS over L162, run by run :: sha256 588e935996d1006a
\newcommand{\MlstFtJcOneRejHggResMassOverResPr}{equal within 5\%}  % experiments/FIGS/data/v2/curves/full_budget_long_leg2.json :: cells.<run>.N1000000.s1.jc1.per_class.label_Hgg: rejection of L162_MASS over L162, run by run :: sha256 588e935996d1006a :: prose twin: two significant figures, direction in words
\newcommand{\MlstFtJcOneRejHggResMassOverResPrMax}{equal within 5\%}  % experiments/FIGS/data/v2/curves/full_budget_long_leg2.json :: cells.<run>.N1000000.s1.jc1.per_class.label_Hgg: rejection of L162_MASS over L162, run by run :: sha256 588e935996d1006a :: prose twin of the run farther from 1
\newcommand{\MlstFtJcOneRejHggCoarseMassOverCoarse}{1.12 and 1.01}  % experiments/FIGS/data/v2/curves/full_budget_long_leg2.json :: cells.<run>.N1000000.s1.jc1.per_class.label_Hgg: rejection of R16_Q1_MASS over R16_Q1, run by run :: sha256 588e935996d1006a
\newcommand{\MlstFtJcOneRejHggCoarseMassOverCoarsePr}{equal within 20\%}  % experiments/FIGS/data/v2/curves/full_budget_long_leg2.json :: cells.<run>.N1000000.s1.jc1.per_class.label_Hgg: rejection of R16_Q1_MASS over R16_Q1, run by run :: sha256 588e935996d1006a :: prose twin: two significant figures, direction in words
\newcommand{\MlstFtJcOneRejHggCoarseMassOverCoarsePrMax}{equal within 20\%}  % experiments/FIGS/data/v2/curves/full_budget_long_leg2.json :: cells.<run>.N1000000.s1.jc1.per_class.label_Hgg: rejection of R16_Q1_MASS over R16_Q1, run by run :: sha256 588e935996d1006a :: prose twin of the run farther from 1
\newcommand{\MlstFtJcOneRejHggRandOverFull}{0.968\,\ensuremath{\pm}\,0.016}  % experiments/FIGS/data/v2/curves/full_budget_long_leg2.json :: per_class.label_Hgg: rejection of each random partition over the 188-class run of its seed :: sha256 588e935996d1006a
\newcommand{\MlstFtJcOneRejHggRandOverFullPr}{equal within 4\%}  % experiments/FIGS/data/v2/curves/full_budget_long_leg2.json :: per_class.label_Hgg: rejection of each random partition over the 188-class run of its seed :: sha256 588e935996d1006a :: prose twin: two significant figures, direction in words
\newcommand{\MlstFtJcOneRejHggRandOverFullRange}{0.94--0.98}  % experiments/FIGS/data/v2/curves/full_budget_long_leg2.json :: per_class.label_Hgg: smallest and largest of the partitions' ratios :: sha256 588e935996d1006a
\newcommand{\MlstFtJcOneRejHfourqResOverFull}{1.05\,\ensuremath{\pm}\,0.06}  % experiments/FIGS/data/v2/curves/full_budget_long_leg2.json :: cells.<run>.N1000000.s1.jc1.per_class.label_H4q: rejection of L162 over L188, run by run :: sha256 588e935996d1006a
\newcommand{\MlstFtJcOneRejHfourqResOverFullPr}{equal within 7\%}  % experiments/FIGS/data/v2/curves/full_budget_long_leg2.json :: cells.<run>.N1000000.s1.jc1.per_class.label_H4q: rejection of L162 over L188, run by run :: sha256 588e935996d1006a :: prose twin: two significant figures, direction in words
\newcommand{\MlstFtJcOneRejHfourqMidOverFull}{0.95\,\ensuremath{\pm}\,0.04}  % experiments/FIGS/data/v2/curves/full_budget_long_leg2.json :: cells.<run>.N1000000.s1.jc1.per_class.label_H4q: rejection of R42_Q1 over L188, run by run :: sha256 588e935996d1006a
\newcommand{\MlstFtJcOneRejHfourqMidOverFullPr}{equal within 6\%}  % experiments/FIGS/data/v2/curves/full_budget_long_leg2.json :: cells.<run>.N1000000.s1.jc1.per_class.label_H4q: rejection of R42_Q1 over L188, run by run :: sha256 588e935996d1006a :: prose twin: two significant figures, direction in words
\newcommand{\MlstFtJcOneRejHfourqCoarseOverFull}{0.962\,\ensuremath{\pm}\,0.025}  % experiments/FIGS/data/v2/curves/full_budget_long_leg2.json :: cells.<run>.N1000000.s1.jc1.per_class.label_H4q: rejection of R16_Q1 over L188, run by run :: sha256 588e935996d1006a
\newcommand{\MlstFtJcOneRejHfourqCoarseOverFullPr}{equal within 4\%}  % experiments/FIGS/data/v2/curves/full_budget_long_leg2.json :: cells.<run>.N1000000.s1.jc1.per_class.label_H4q: rejection of R16_Q1 over L188, run by run :: sha256 588e935996d1006a :: prose twin: two significant figures, direction in words
\newcommand{\MlstFtJcOneRejHfourqResMassOverRes}{0.912 and 0.966}  % experiments/FIGS/data/v2/curves/full_budget_long_leg2.json :: cells.<run>.N1000000.s1.jc1.per_class.label_H4q: rejection of L162_MASS over L162, run by run :: sha256 588e935996d1006a
\newcommand{\MlstFtJcOneRejHfourqResMassOverResPr}{equal within 9\%}  % experiments/FIGS/data/v2/curves/full_budget_long_leg2.json :: cells.<run>.N1000000.s1.jc1.per_class.label_H4q: rejection of L162_MASS over L162, run by run :: sha256 588e935996d1006a :: prose twin: two significant figures, direction in words
\newcommand{\MlstFtJcOneRejHfourqResMassOverResPrMax}{equal within 9\%}  % experiments/FIGS/data/v2/curves/full_budget_long_leg2.json :: cells.<run>.N1000000.s1.jc1.per_class.label_H4q: rejection of L162_MASS over L162, run by run :: sha256 588e935996d1006a :: prose twin of the run farther from 1
\newcommand{\MlstFtJcOneRejHfourqCoarseMassOverCoarse}{0.937 and 0.888}  % experiments/FIGS/data/v2/curves/full_budget_long_leg2.json :: cells.<run>.N1000000.s1.jc1.per_class.label_H4q: rejection of R16_Q1_MASS over R16_Q1, run by run :: sha256 588e935996d1006a
\newcommand{\MlstFtJcOneRejHfourqCoarseMassOverCoarsePr}{equal within 20\%}  % experiments/FIGS/data/v2/curves/full_budget_long_leg2.json :: cells.<run>.N1000000.s1.jc1.per_class.label_H4q: rejection of R16_Q1_MASS over R16_Q1, run by run :: sha256 588e935996d1006a :: prose twin: two significant figures, direction in words
\newcommand{\MlstFtJcOneRejHfourqCoarseMassOverCoarsePrMax}{equal within 20\%}  % experiments/FIGS/data/v2/curves/full_budget_long_leg2.json :: cells.<run>.N1000000.s1.jc1.per_class.label_H4q: rejection of R16_Q1_MASS over R16_Q1, run by run :: sha256 588e935996d1006a :: prose twin of the run farther from 1
\newcommand{\MlstFtJcOneRejHfourqRandOverFull}{0.93\,\ensuremath{\pm}\,0.04}  % experiments/FIGS/data/v2/curves/full_budget_long_leg2.json :: per_class.label_H4q: rejection of each random partition over the 188-class run of its seed :: sha256 588e935996d1006a
\newcommand{\MlstFtJcOneRejHfourqRandOverFullPr}{equal within 8\%}  % experiments/FIGS/data/v2/curves/full_budget_long_leg2.json :: per_class.label_H4q: rejection of each random partition over the 188-class run of its seed :: sha256 588e935996d1006a :: prose twin: two significant figures, direction in words
\newcommand{\MlstFtJcOneRejHfourqRandOverFullRange}{0.86--0.96}  % experiments/FIGS/data/v2/curves/full_budget_long_leg2.json :: per_class.label_H4q: smallest and largest of the partitions' ratios :: sha256 588e935996d1006a
\newcommand{\MlstFtJcOneRejHqqlResOverFull}{1.004\,\ensuremath{\pm}\,0.018}  % experiments/FIGS/data/v2/curves/full_budget_long_leg2.json :: cells.<run>.N1000000.s1.jc1.per_class.label_Hqql: rejection of L162 over L188, run by run :: sha256 588e935996d1006a
\newcommand{\MlstFtJcOneRejHqqlResOverFullPr}{equal within 2\%}  % experiments/FIGS/data/v2/curves/full_budget_long_leg2.json :: cells.<run>.N1000000.s1.jc1.per_class.label_Hqql: rejection of L162 over L188, run by run :: sha256 588e935996d1006a :: prose twin: two significant figures, direction in words
\newcommand{\MlstFtJcOneRejHqqlMidOverFull}{1.04\,\ensuremath{\pm}\,0.18}  % experiments/FIGS/data/v2/curves/full_budget_long_leg2.json :: cells.<run>.N1000000.s1.jc1.per_class.label_Hqql: rejection of R42_Q1 over L188, run by run :: sha256 588e935996d1006a
\newcommand{\MlstFtJcOneRejHqqlMidOverFullPr}{equal within 20\%}  % experiments/FIGS/data/v2/curves/full_budget_long_leg2.json :: cells.<run>.N1000000.s1.jc1.per_class.label_Hqql: rejection of R42_Q1 over L188, run by run :: sha256 588e935996d1006a :: prose twin: two significant figures, direction in words
\newcommand{\MlstFtJcOneRejHqqlCoarseOverFull}{0.93\,\ensuremath{\pm}\,0.16}  % experiments/FIGS/data/v2/curves/full_budget_long_leg2.json :: cells.<run>.N1000000.s1.jc1.per_class.label_Hqql: rejection of R16_Q1 over L188, run by run :: sha256 588e935996d1006a
\newcommand{\MlstFtJcOneRejHqqlCoarseOverFullPr}{equal within 20\%}  % experiments/FIGS/data/v2/curves/full_budget_long_leg2.json :: cells.<run>.N1000000.s1.jc1.per_class.label_Hqql: rejection of R16_Q1 over L188, run by run :: sha256 588e935996d1006a :: prose twin: two significant figures, direction in words
\newcommand{\MlstFtJcOneRejHqqlResMassOverRes}{0.986 and 1.000}  % experiments/FIGS/data/v2/curves/full_budget_long_leg2.json :: cells.<run>.N1000000.s1.jc1.per_class.label_Hqql: rejection of L162_MASS over L162, run by run :: sha256 588e935996d1006a
\newcommand{\MlstFtJcOneRejHqqlResMassOverResPr}{equal within 2\%}  % experiments/FIGS/data/v2/curves/full_budget_long_leg2.json :: cells.<run>.N1000000.s1.jc1.per_class.label_Hqql: rejection of L162_MASS over L162, run by run :: sha256 588e935996d1006a :: prose twin: two significant figures, direction in words
\newcommand{\MlstFtJcOneRejHqqlResMassOverResPrMax}{equal within 2\%}  % experiments/FIGS/data/v2/curves/full_budget_long_leg2.json :: cells.<run>.N1000000.s1.jc1.per_class.label_Hqql: rejection of L162_MASS over L162, run by run :: sha256 588e935996d1006a :: prose twin of the run farther from 1
\newcommand{\MlstFtJcOneRejHqqlCoarseMassOverCoarse}{0.97 and 1.09}  % experiments/FIGS/data/v2/curves/full_budget_long_leg2.json :: cells.<run>.N1000000.s1.jc1.per_class.label_Hqql: rejection of R16_Q1_MASS over R16_Q1, run by run :: sha256 588e935996d1006a
\newcommand{\MlstFtJcOneRejHqqlCoarseMassOverCoarsePr}{equal within 9\%}  % experiments/FIGS/data/v2/curves/full_budget_long_leg2.json :: cells.<run>.N1000000.s1.jc1.per_class.label_Hqql: rejection of R16_Q1_MASS over R16_Q1, run by run :: sha256 588e935996d1006a :: prose twin: two significant figures, direction in words
\newcommand{\MlstFtJcOneRejHqqlCoarseMassOverCoarsePrMax}{equal within 9\%}  % experiments/FIGS/data/v2/curves/full_budget_long_leg2.json :: cells.<run>.N1000000.s1.jc1.per_class.label_Hqql: rejection of R16_Q1_MASS over R16_Q1, run by run :: sha256 588e935996d1006a :: prose twin of the run farther from 1
\newcommand{\MlstFtJcOneRejHqqlRandOverFull}{0.94\,\ensuremath{\pm}\,0.12}  % experiments/FIGS/data/v2/curves/full_budget_long_leg2.json :: per_class.label_Hqql: rejection of each random partition over the 188-class run of its seed :: sha256 588e935996d1006a
\newcommand{\MlstFtJcOneRejHqqlRandOverFullPr}{equal within 20\%}  % experiments/FIGS/data/v2/curves/full_budget_long_leg2.json :: per_class.label_Hqql: rejection of each random partition over the 188-class run of its seed :: sha256 588e935996d1006a :: prose twin: two significant figures, direction in words
\newcommand{\MlstFtJcOneRejHqqlRandOverFullRange}{0.79--1.1}  % experiments/FIGS/data/v2/curves/full_budget_long_leg2.json :: per_class.label_Hqql: smallest and largest of the partitions' ratios :: sha256 588e935996d1006a
\newcommand{\MlstFtJcOneRejTbqqResOverFull}{1.1\,\ensuremath{\pm}\,0.4}  % experiments/FIGS/data/v2/curves/full_budget_long_leg2.json :: cells.<run>.N1000000.s1.jc1.per_class.label_Tbqq: rejection of L162 over L188, run by run :: sha256 588e935996d1006a
\newcommand{\MlstFtJcOneRejTbqqResOverFullPr}{equal within 50\%}  % experiments/FIGS/data/v2/curves/full_budget_long_leg2.json :: cells.<run>.N1000000.s1.jc1.per_class.label_Tbqq: rejection of L162 over L188, run by run :: sha256 588e935996d1006a :: prose twin: two significant figures, direction in words
\newcommand{\MlstFtJcOneRejTbqqMidOverFull}{1.10\,\ensuremath{\pm}\,0.05}  % experiments/FIGS/data/v2/curves/full_budget_long_leg2.json :: cells.<run>.N1000000.s1.jc1.per_class.label_Tbqq: rejection of R42_Q1 over L188, run by run :: sha256 588e935996d1006a
\newcommand{\MlstFtJcOneRejTbqqMidOverFullPr}{equal within 10\%}  % experiments/FIGS/data/v2/curves/full_budget_long_leg2.json :: cells.<run>.N1000000.s1.jc1.per_class.label_Tbqq: rejection of R42_Q1 over L188, run by run :: sha256 588e935996d1006a :: prose twin: two significant figures, direction in words
\newcommand{\MlstFtJcOneRejTbqqCoarseOverFull}{0.85\,\ensuremath{\pm}\,0.32}  % experiments/FIGS/data/v2/curves/full_budget_long_leg2.json :: cells.<run>.N1000000.s1.jc1.per_class.label_Tbqq: rejection of R16_Q1 over L188, run by run :: sha256 588e935996d1006a
\newcommand{\MlstFtJcOneRejTbqqCoarseOverFullPr}{equal within 40\%}  % experiments/FIGS/data/v2/curves/full_budget_long_leg2.json :: cells.<run>.N1000000.s1.jc1.per_class.label_Tbqq: rejection of R16_Q1 over L188, run by run :: sha256 588e935996d1006a :: prose twin: two significant figures, direction in words
\newcommand{\MlstFtJcOneRejTbqqResMassOverRes}{0.89 and 1.21}  % experiments/FIGS/data/v2/curves/full_budget_long_leg2.json :: cells.<run>.N1000000.s1.jc1.per_class.label_Tbqq: rejection of L162_MASS over L162, run by run :: sha256 588e935996d1006a
\newcommand{\MlstFtJcOneRejTbqqResMassOverResPr}{equal within 30\%}  % experiments/FIGS/data/v2/curves/full_budget_long_leg2.json :: cells.<run>.N1000000.s1.jc1.per_class.label_Tbqq: rejection of L162_MASS over L162, run by run :: sha256 588e935996d1006a :: prose twin: two significant figures, direction in words
\newcommand{\MlstFtJcOneRejTbqqResMassOverResPrMax}{equal within 30\%}  % experiments/FIGS/data/v2/curves/full_budget_long_leg2.json :: cells.<run>.N1000000.s1.jc1.per_class.label_Tbqq: rejection of L162_MASS over L162, run by run :: sha256 588e935996d1006a :: prose twin of the run farther from 1
\newcommand{\MlstFtJcOneRejTbqqCoarseMassOverCoarse}{1.00 and 1.79}  % experiments/FIGS/data/v2/curves/full_budget_long_leg2.json :: cells.<run>.N1000000.s1.jc1.per_class.label_Tbqq: rejection of R16_Q1_MASS over R16_Q1, run by run :: sha256 588e935996d1006a
\newcommand{\MlstFtJcOneRejTbqqCoarseMassOverCoarsePr}{equal within 80\%}  % experiments/FIGS/data/v2/curves/full_budget_long_leg2.json :: cells.<run>.N1000000.s1.jc1.per_class.label_Tbqq: rejection of R16_Q1_MASS over R16_Q1, run by run :: sha256 588e935996d1006a :: prose twin: two significant figures, direction in words
\newcommand{\MlstFtJcOneRejTbqqCoarseMassOverCoarsePrMax}{equal within 80\%}  % experiments/FIGS/data/v2/curves/full_budget_long_leg2.json :: cells.<run>.N1000000.s1.jc1.per_class.label_Tbqq: rejection of R16_Q1_MASS over R16_Q1, run by run :: sha256 588e935996d1006a :: prose twin of the run farther from 1
\newcommand{\MlstFtJcOneRejTbqqRandOverFull}{0.79\,\ensuremath{\pm}\,0.15}  % experiments/FIGS/data/v2/curves/full_budget_long_leg2.json :: per_class.label_Tbqq: rejection of each random partition over the 188-class run of its seed :: sha256 588e935996d1006a
\newcommand{\MlstFtJcOneRejTbqqRandOverFullPr}{equal within 30\%}  % experiments/FIGS/data/v2/curves/full_budget_long_leg2.json :: per_class.label_Tbqq: rejection of each random partition over the 188-class run of its seed :: sha256 588e935996d1006a :: prose twin: two significant figures, direction in words
\newcommand{\MlstFtJcOneRejTbqqRandOverFullRange}{0.64--1.0}  % experiments/FIGS/data/v2/curves/full_budget_long_leg2.json :: per_class.label_Tbqq: smallest and largest of the partitions' ratios :: sha256 588e935996d1006a
\newcommand{\MlstFtJcOneRejTblResOverFull}{0.85\,\ensuremath{\pm}\,0.17}  % experiments/FIGS/data/v2/curves/full_budget_long_leg2.json :: cells.<run>.N1000000.s1.jc1.per_class.label_Tbl: rejection of L162 over L188, run by run :: sha256 588e935996d1006a
\newcommand{\MlstFtJcOneRejTblResOverFullPr}{equal within 20\%}  % experiments/FIGS/data/v2/curves/full_budget_long_leg2.json :: cells.<run>.N1000000.s1.jc1.per_class.label_Tbl: rejection of L162 over L188, run by run :: sha256 588e935996d1006a :: prose twin: two significant figures, direction in words
\newcommand{\MlstFtJcOneRejTblMidOverFull}{0.91\,\ensuremath{\pm}\,0.15}  % experiments/FIGS/data/v2/curves/full_budget_long_leg2.json :: cells.<run>.N1000000.s1.jc1.per_class.label_Tbl: rejection of R42_Q1 over L188, run by run :: sha256 588e935996d1006a
\newcommand{\MlstFtJcOneRejTblMidOverFullPr}{equal within 20\%}  % experiments/FIGS/data/v2/curves/full_budget_long_leg2.json :: cells.<run>.N1000000.s1.jc1.per_class.label_Tbl: rejection of R42_Q1 over L188, run by run :: sha256 588e935996d1006a :: prose twin: two significant figures, direction in words
\newcommand{\MlstFtJcOneRejTblCoarseOverFull}{0.94\,\ensuremath{\pm}\,0.28}  % experiments/FIGS/data/v2/curves/full_budget_long_leg2.json :: cells.<run>.N1000000.s1.jc1.per_class.label_Tbl: rejection of R16_Q1 over L188, run by run :: sha256 588e935996d1006a
\newcommand{\MlstFtJcOneRejTblCoarseOverFullPr}{equal within 30\%}  % experiments/FIGS/data/v2/curves/full_budget_long_leg2.json :: cells.<run>.N1000000.s1.jc1.per_class.label_Tbl: rejection of R16_Q1 over L188, run by run :: sha256 588e935996d1006a :: prose twin: two significant figures, direction in words
\newcommand{\MlstFtJcOneRejTblResMassOverRes}{1.04 and 0.93}  % experiments/FIGS/data/v2/curves/full_budget_long_leg2.json :: cells.<run>.N1000000.s1.jc1.per_class.label_Tbl: rejection of L162_MASS over L162, run by run :: sha256 588e935996d1006a
\newcommand{\MlstFtJcOneRejTblResMassOverResPr}{equal within 7\%}  % experiments/FIGS/data/v2/curves/full_budget_long_leg2.json :: cells.<run>.N1000000.s1.jc1.per_class.label_Tbl: rejection of L162_MASS over L162, run by run :: sha256 588e935996d1006a :: prose twin: two significant figures, direction in words
\newcommand{\MlstFtJcOneRejTblResMassOverResPrMax}{equal within 7\%}  % experiments/FIGS/data/v2/curves/full_budget_long_leg2.json :: cells.<run>.N1000000.s1.jc1.per_class.label_Tbl: rejection of L162_MASS over L162, run by run :: sha256 588e935996d1006a :: prose twin of the run farther from 1
\newcommand{\MlstFtJcOneRejTblCoarseMassOverCoarse}{0.84 and 1.22}  % experiments/FIGS/data/v2/curves/full_budget_long_leg2.json :: cells.<run>.N1000000.s1.jc1.per_class.label_Tbl: rejection of R16_Q1_MASS over R16_Q1, run by run :: sha256 588e935996d1006a
\newcommand{\MlstFtJcOneRejTblCoarseMassOverCoarsePr}{equal within 30\%}  % experiments/FIGS/data/v2/curves/full_budget_long_leg2.json :: cells.<run>.N1000000.s1.jc1.per_class.label_Tbl: rejection of R16_Q1_MASS over R16_Q1, run by run :: sha256 588e935996d1006a :: prose twin: two significant figures, direction in words
\newcommand{\MlstFtJcOneRejTblCoarseMassOverCoarsePrMax}{equal within 30\%}  % experiments/FIGS/data/v2/curves/full_budget_long_leg2.json :: cells.<run>.N1000000.s1.jc1.per_class.label_Tbl: rejection of R16_Q1_MASS over R16_Q1, run by run :: sha256 588e935996d1006a :: prose twin of the run farther from 1
\newcommand{\MlstFtJcOneRejTblRandOverFull}{0.88\,\ensuremath{\pm}\,0.26}  % experiments/FIGS/data/v2/curves/full_budget_long_leg2.json :: per_class.label_Tbl: rejection of each random partition over the 188-class run of its seed :: sha256 588e935996d1006a
\newcommand{\MlstFtJcOneRejTblRandOverFullPr}{equal within 30\%}  % experiments/FIGS/data/v2/curves/full_budget_long_leg2.json :: per_class.label_Tbl: rejection of each random partition over the 188-class run of its seed :: sha256 588e935996d1006a :: prose twin: two significant figures, direction in words
\newcommand{\MlstFtJcOneRejTblRandOverFullRange}{0.61--1.2}  % experiments/FIGS/data/v2/curves/full_budget_long_leg2.json :: per_class.label_Tbl: smallest and largest of the partitions' ratios :: sha256 588e935996d1006a
\newcommand{\MlstFtJcOneRejWqqResOverFull}{1.039\,\ensuremath{\pm}\,0.035}  % experiments/FIGS/data/v2/curves/full_budget_long_leg2.json :: cells.<run>.N1000000.s1.jc1.per_class.label_Wqq: rejection of L162 over L188, run by run :: sha256 588e935996d1006a
\newcommand{\MlstFtJcOneRejWqqResOverFullPr}{equal within 4\%}  % experiments/FIGS/data/v2/curves/full_budget_long_leg2.json :: cells.<run>.N1000000.s1.jc1.per_class.label_Wqq: rejection of L162 over L188, run by run :: sha256 588e935996d1006a :: prose twin: two significant figures, direction in words
\newcommand{\MlstFtJcOneRejWqqMidOverFull}{1.005\,\ensuremath{\pm}\,0.032}  % experiments/FIGS/data/v2/curves/full_budget_long_leg2.json :: cells.<run>.N1000000.s1.jc1.per_class.label_Wqq: rejection of R42_Q1 over L188, run by run :: sha256 588e935996d1006a
\newcommand{\MlstFtJcOneRejWqqMidOverFullPr}{equal within 4\%}  % experiments/FIGS/data/v2/curves/full_budget_long_leg2.json :: cells.<run>.N1000000.s1.jc1.per_class.label_Wqq: rejection of R42_Q1 over L188, run by run :: sha256 588e935996d1006a :: prose twin: two significant figures, direction in words
\newcommand{\MlstFtJcOneRejWqqCoarseOverFull}{1.01\,\ensuremath{\pm}\,0.07}  % experiments/FIGS/data/v2/curves/full_budget_long_leg2.json :: cells.<run>.N1000000.s1.jc1.per_class.label_Wqq: rejection of R16_Q1 over L188, run by run :: sha256 588e935996d1006a
\newcommand{\MlstFtJcOneRejWqqCoarseOverFullPr}{equal within 8\%}  % experiments/FIGS/data/v2/curves/full_budget_long_leg2.json :: cells.<run>.N1000000.s1.jc1.per_class.label_Wqq: rejection of R16_Q1 over L188, run by run :: sha256 588e935996d1006a :: prose twin: two significant figures, direction in words
\newcommand{\MlstFtJcOneRejWqqResMassOverRes}{0.996 and 0.964}  % experiments/FIGS/data/v2/curves/full_budget_long_leg2.json :: cells.<run>.N1000000.s1.jc1.per_class.label_Wqq: rejection of L162_MASS over L162, run by run :: sha256 588e935996d1006a
\newcommand{\MlstFtJcOneRejWqqResMassOverResPr}{equal within 4\%}  % experiments/FIGS/data/v2/curves/full_budget_long_leg2.json :: cells.<run>.N1000000.s1.jc1.per_class.label_Wqq: rejection of L162_MASS over L162, run by run :: sha256 588e935996d1006a :: prose twin: two significant figures, direction in words
\newcommand{\MlstFtJcOneRejWqqResMassOverResPrMax}{equal within 4\%}  % experiments/FIGS/data/v2/curves/full_budget_long_leg2.json :: cells.<run>.N1000000.s1.jc1.per_class.label_Wqq: rejection of L162_MASS over L162, run by run :: sha256 588e935996d1006a :: prose twin of the run farther from 1
\newcommand{\MlstFtJcOneRejWqqCoarseMassOverCoarse}{1.025 and 0.926}  % experiments/FIGS/data/v2/curves/full_budget_long_leg2.json :: cells.<run>.N1000000.s1.jc1.per_class.label_Wqq: rejection of R16_Q1_MASS over R16_Q1, run by run :: sha256 588e935996d1006a
\newcommand{\MlstFtJcOneRejWqqCoarseMassOverCoarsePr}{equal within 8\%}  % experiments/FIGS/data/v2/curves/full_budget_long_leg2.json :: cells.<run>.N1000000.s1.jc1.per_class.label_Wqq: rejection of R16_Q1_MASS over R16_Q1, run by run :: sha256 588e935996d1006a :: prose twin: two significant figures, direction in words
\newcommand{\MlstFtJcOneRejWqqCoarseMassOverCoarsePrMax}{equal within 8\%}  % experiments/FIGS/data/v2/curves/full_budget_long_leg2.json :: cells.<run>.N1000000.s1.jc1.per_class.label_Wqq: rejection of R16_Q1_MASS over R16_Q1, run by run :: sha256 588e935996d1006a :: prose twin of the run farther from 1
\newcommand{\MlstFtJcOneRejWqqRandOverFull}{1.02\,\ensuremath{\pm}\,0.07}  % experiments/FIGS/data/v2/curves/full_budget_long_leg2.json :: per_class.label_Wqq: rejection of each random partition over the 188-class run of its seed :: sha256 588e935996d1006a
\newcommand{\MlstFtJcOneRejWqqRandOverFullPr}{equal within 7\%}  % experiments/FIGS/data/v2/curves/full_budget_long_leg2.json :: per_class.label_Wqq: rejection of each random partition over the 188-class run of its seed :: sha256 588e935996d1006a :: prose twin: two significant figures, direction in words
\newcommand{\MlstFtJcOneRejWqqRandOverFullRange}{0.93--1.1}  % experiments/FIGS/data/v2/curves/full_budget_long_leg2.json :: per_class.label_Wqq: smallest and largest of the partitions' ratios :: sha256 588e935996d1006a
\newcommand{\MlstFtJcOneRejZqqResOverFull}{0.987\,\ensuremath{\pm}\,0.019}  % experiments/FIGS/data/v2/curves/full_budget_long_leg2.json :: cells.<run>.N1000000.s1.jc1.per_class.label_Zqq: rejection of L162 over L188, run by run :: sha256 588e935996d1006a
\newcommand{\MlstFtJcOneRejZqqResOverFullPr}{equal within 2\%}  % experiments/FIGS/data/v2/curves/full_budget_long_leg2.json :: cells.<run>.N1000000.s1.jc1.per_class.label_Zqq: rejection of L162 over L188, run by run :: sha256 588e935996d1006a :: prose twin: two significant figures, direction in words
\newcommand{\MlstFtJcOneRejZqqMidOverFull}{0.984\,\ensuremath{\pm}\,0.010}  % experiments/FIGS/data/v2/curves/full_budget_long_leg2.json :: cells.<run>.N1000000.s1.jc1.per_class.label_Zqq: rejection of R42_Q1 over L188, run by run :: sha256 588e935996d1006a
\newcommand{\MlstFtJcOneRejZqqMidOverFullPr}{equal within 2\%}  % experiments/FIGS/data/v2/curves/full_budget_long_leg2.json :: cells.<run>.N1000000.s1.jc1.per_class.label_Zqq: rejection of R42_Q1 over L188, run by run :: sha256 588e935996d1006a :: prose twin: two significant figures, direction in words
\newcommand{\MlstFtJcOneRejZqqCoarseOverFull}{0.91\,\ensuremath{\pm}\,0.05}  % experiments/FIGS/data/v2/curves/full_budget_long_leg2.json :: cells.<run>.N1000000.s1.jc1.per_class.label_Zqq: rejection of R16_Q1 over L188, run by run :: sha256 588e935996d1006a
\newcommand{\MlstFtJcOneRejZqqCoarseOverFullPr}{equal within 9\%}  % experiments/FIGS/data/v2/curves/full_budget_long_leg2.json :: cells.<run>.N1000000.s1.jc1.per_class.label_Zqq: rejection of R16_Q1 over L188, run by run :: sha256 588e935996d1006a :: prose twin: two significant figures, direction in words
\newcommand{\MlstFtJcOneRejZqqResMassOverRes}{0.9803 and 0.9820}  % experiments/FIGS/data/v2/curves/full_budget_long_leg2.json :: cells.<run>.N1000000.s1.jc1.per_class.label_Zqq: rejection of L162_MASS over L162, run by run :: sha256 588e935996d1006a
\newcommand{\MlstFtJcOneRejZqqResMassOverResPr}{equal within 2\%}  % experiments/FIGS/data/v2/curves/full_budget_long_leg2.json :: cells.<run>.N1000000.s1.jc1.per_class.label_Zqq: rejection of L162_MASS over L162, run by run :: sha256 588e935996d1006a :: prose twin: two significant figures, direction in words
\newcommand{\MlstFtJcOneRejZqqResMassOverResPrMax}{equal within 2\%}  % experiments/FIGS/data/v2/curves/full_budget_long_leg2.json :: cells.<run>.N1000000.s1.jc1.per_class.label_Zqq: rejection of L162_MASS over L162, run by run :: sha256 588e935996d1006a :: prose twin of the run farther from 1
\newcommand{\MlstFtJcOneRejZqqCoarseMassOverCoarse}{1.06 and 0.94}  % experiments/FIGS/data/v2/curves/full_budget_long_leg2.json :: cells.<run>.N1000000.s1.jc1.per_class.label_Zqq: rejection of R16_Q1_MASS over R16_Q1, run by run :: sha256 588e935996d1006a
\newcommand{\MlstFtJcOneRejZqqCoarseMassOverCoarsePr}{equal within 7\%}  % experiments/FIGS/data/v2/curves/full_budget_long_leg2.json :: cells.<run>.N1000000.s1.jc1.per_class.label_Zqq: rejection of R16_Q1_MASS over R16_Q1, run by run :: sha256 588e935996d1006a :: prose twin: two significant figures, direction in words
\newcommand{\MlstFtJcOneRejZqqCoarseMassOverCoarsePrMax}{equal within 7\%}  % experiments/FIGS/data/v2/curves/full_budget_long_leg2.json :: cells.<run>.N1000000.s1.jc1.per_class.label_Zqq: rejection of R16_Q1_MASS over R16_Q1, run by run :: sha256 588e935996d1006a :: prose twin of the run farther from 1
\newcommand{\MlstFtJcOneRejZqqRandOverFull}{0.952\,\ensuremath{\pm}\,0.028}  % experiments/FIGS/data/v2/curves/full_budget_long_leg2.json :: per_class.label_Zqq: rejection of each random partition over the 188-class run of its seed :: sha256 588e935996d1006a
\newcommand{\MlstFtJcOneRejZqqRandOverFullPr}{equal within 5\%}  % experiments/FIGS/data/v2/curves/full_budget_long_leg2.json :: per_class.label_Zqq: rejection of each random partition over the 188-class run of its seed :: sha256 588e935996d1006a :: prose twin: two significant figures, direction in words
\newcommand{\MlstFtJcOneRejZqqRandOverFullRange}{0.93--0.99}  % experiments/FIGS/data/v2/curves/full_budget_long_leg2.json :: per_class.label_Zqq: smallest and largest of the partitions' ratios :: sha256 588e935996d1006a
\newcommand{\MlstFtJcOneResNRejDrop}{zero}  % experiments/FIGS/data/v2/curves/full_budget_long_leg2.json :: per_class rejections :: sha256 588e935996d1006a :: classes whose paired rejection ratio to the 188-class model departs from 1 by more than twice its SD and its counting uncertainty, columns with fewer than JC1_MIN_BKG passing jets excluded, ratio below 1 (count)
\newcommand{\MlstFtJcOneResNRejDiff}{zero}  % experiments/FIGS/data/v2/curves/full_budget_long_leg2.json :: per_class rejections :: sha256 588e935996d1006a :: classes whose paired rejection ratio to the 188-class model departs from 1 by more than twice its SD and its counting uncertainty, columns with fewer than JC1_MIN_BKG passing jets excluded (count)
\newcommand{\MlstFtJcOneResRejDropClasses}{none}  % experiments/FIGS/data/v2/curves/full_budget_long_leg2.json :: per_class rejections :: sha256 588e935996d1006a :: the classes of ...NRejDrop, in ParT's column order
\newcommand{\MlstFtJcOneMidNRejDrop}{zero}  % experiments/FIGS/data/v2/curves/full_budget_long_leg2.json :: per_class rejections :: sha256 588e935996d1006a :: classes whose paired rejection ratio to the 188-class model departs from 1 by more than twice its SD and its counting uncertainty, columns with fewer than JC1_MIN_BKG passing jets excluded, ratio below 1 (count)
\newcommand{\MlstFtJcOneMidNRejDiff}{zero}  % experiments/FIGS/data/v2/curves/full_budget_long_leg2.json :: per_class rejections :: sha256 588e935996d1006a :: classes whose paired rejection ratio to the 188-class model departs from 1 by more than twice its SD and its counting uncertainty, columns with fewer than JC1_MIN_BKG passing jets excluded (count)
\newcommand{\MlstFtJcOneMidRejDropClasses}{none}  % experiments/FIGS/data/v2/curves/full_budget_long_leg2.json :: per_class rejections :: sha256 588e935996d1006a :: the classes of ...NRejDrop, in ParT's column order
\newcommand{\MlstFtJcOneCoarseNRejDrop}{zero}  % experiments/FIGS/data/v2/curves/full_budget_long_leg2.json :: per_class rejections :: sha256 588e935996d1006a :: classes whose paired rejection ratio to the 188-class model departs from 1 by more than twice its SD and its counting uncertainty, columns with fewer than JC1_MIN_BKG passing jets excluded, ratio below 1 (count)
\newcommand{\MlstFtJcOneCoarseNRejDiff}{zero}  % experiments/FIGS/data/v2/curves/full_budget_long_leg2.json :: per_class rejections :: sha256 588e935996d1006a :: classes whose paired rejection ratio to the 188-class model departs from 1 by more than twice its SD and its counting uncertainty, columns with fewer than JC1_MIN_BKG passing jets excluded (count)
\newcommand{\MlstFtJcOneCoarseRejDropClasses}{none}  % experiments/FIGS/data/v2/curves/full_budget_long_leg2.json :: per_class rejections :: sha256 588e935996d1006a :: the classes of ...NRejDrop, in ParT's column order
\newcommand{\MlstFtJcOneAccRandMinusCoarse}{+0.002}  % experiments/FIGS/data/v2/curves/full_budget_long_leg2.json :: cells.*.N1000000.s1.jc1.accuracy :: sha256 588e935996d1006a :: mean over the random partitions minus the mean over the 17-class runs
\newcommand{\MlstFtJcOneAccFullBoot}{0.847}  % experiments/FIGS/data/v2/budget/jc1_acc_boot_N1000000.json :: diffs.r16q1-l188.accuracy_ref (mean) :: sha256 7cccb7b0ff900fb8
\newcommand{\MlstFtJcOneAccCoarseBoot}{0.842}  % experiments/FIGS/data/v2/budget/jc1_acc_boot_N1000000.json :: diffs.r16q1-l188.accuracy_config (mean) :: sha256 7cccb7b0ff900fb8
\newcommand{\MlstFtJcOneErrResOverFullBoot}{consistent}  % experiments/FIGS/data/v2/budget/jc1_acc_boot_N1000000.json :: diffs.l162-l188.error_ratio :: sha256 7cccb7b0ff900fb8 :: delta and sigma = hypot(run SD, bootstrap SD) in per cent; 'consistent' when |delta| <= 3 sigma
\newcommand{\MlstFtJcOneErrResOverFullZ}{0.6}  % experiments/FIGS/data/v2/budget/jc1_acc_boot_N1000000.json :: diffs.l162-l188.error_ratio.z :: sha256 7cccb7b0ff900fb8 :: |delta| / sigma
\newcommand{\MlstFtJcOneErrMidOverFullBoot}{consistent}  % experiments/FIGS/data/v2/budget/jc1_acc_boot_N1000000.json :: diffs.r42q1-l188.error_ratio :: sha256 7cccb7b0ff900fb8 :: delta and sigma = hypot(run SD, bootstrap SD) in per cent; 'consistent' when |delta| <= 3 sigma
\newcommand{\MlstFtJcOneErrMidOverFullZ}{1.4}  % experiments/FIGS/data/v2/budget/jc1_acc_boot_N1000000.json :: diffs.r42q1-l188.error_ratio.z :: sha256 7cccb7b0ff900fb8 :: |delta| / sigma
\newcommand{\MlstFtJcOneErrCoarseOverFullBoot}{(3.0\,\ensuremath{\pm}\,0.6)\% more errors}  % experiments/FIGS/data/v2/budget/jc1_acc_boot_N1000000.json :: diffs.r16q1-l188.error_ratio :: sha256 7cccb7b0ff900fb8 :: delta and sigma = hypot(run SD, bootstrap SD) in per cent; 'consistent' when |delta| <= 3 sigma
\newcommand{\MlstFtJcOneErrCoarseOverFullZ}{4.8}  % experiments/FIGS/data/v2/budget/jc1_acc_boot_N1000000.json :: diffs.r16q1-l188.error_ratio.z :: sha256 7cccb7b0ff900fb8 :: |delta| / sigma
\newcommand{\MlstFtJcOneBootReplicas}{200}  % experiments/FIGS/data/v2/budget/jc1_acc_boot_N1000000.json :: replicas :: sha256 7cccb7b0ff900fb8
\newcommand{\MlstFtJcOneBootJets}{2{,}000{,}000}  % experiments/FIGS/data/v2/budget/jc1_acc_boot_N1000000.json :: n_test_jets :: sha256 7cccb7b0ff900fb8
\newcommand{\ComputeTotalGpuHours}{3{,}435}  % experiments/FIGS/data/v2/index.json :: GPU-hours of every reported pretraining run, training and validation passes :: sha256 c6c4138ee5c86423
\newcommand{\MlstComputeTotalGpuHours}{3{,}510}  % experiments/FIGS/data/v2/index.json :: GPU-hours of every pretraining run the paper reports, the kinematic-only model included :: sha256 c6c4138ee5c86423
\newcommand{\ComputeMainGpu}{V100}  % experiments/FIGS/data/v2/index.json :: GPU type with the most hours over the four label sets :: sha256 c6c4138ee5c86423
\newcommand{\ComputeMainGpuShareMin}{56\%}  % experiments/FIGS/data/v2/index.json :: lowest share of a label set's GPU-hours on that type :: sha256 c6c4138ee5c86423
\newcommand{\ComputeMainGpuShareMax}{71\%}  % experiments/FIGS/data/v2/index.json :: highest share of a label set's GPU-hours on that type :: sha256 c6c4138ee5c86423
\newcommand{\ComputeRunsDone}{27}  % experiments/FIGS/data/v2/index.json :: paper runs with a compute record (count) :: sha256 c6c4138ee5c86423
\newcommand{\ComputeGpuTypes}{L4, RTX 3090, RTX 4090 and V100}  % experiments/FIGS/data/v2/index.json :: GPU types any finished paper run trained on :: sha256 c6c4138ee5c86423
\newcommand{\MlstComputeGpuTypes}{L4, RTX 3090, RTX 4090 and V100}  % experiments/FIGS/data/v2/index.json :: GPU types of the pretraining runs the paper reports :: sha256 c6c4138ee5c86423
\newcommand{\MlstFtJcOneErrCoarseOverFullBootEFour}{\pending{the JetClass bootstrap at 1e4 (after the equal-stopping read-out)}}  % experiments/FIGS/data/v2/index.json :: not computed yet :: sha256 c6c4138ee5c86423 :: waits on the JetClass bootstrap at 1e4 (after the equal-stopping read-out)
\newcommand{\MlstLpBbQgZeroAccDropFull}{\pending{pair probes with the quark/gluon inputs zeroed (open item O10, not run)}}  % experiments/FIGS/data/v2/index.json :: not computed yet :: sha256 c6c4138ee5c86423 :: waits on pair probes with the quark/gluon inputs zeroed (open item O10, not run)
\newcommand{\MlstLpBbQgZeroAccDropCoarse}{\pending{pair probes with the quark/gluon inputs zeroed (open item O10, not run)}}  % experiments/FIGS/data/v2/index.json :: not computed yet :: sha256 c6c4138ee5c86423 :: waits on pair probes with the quark/gluon inputs zeroed (open item O10, not run)
\newcommand{\XbsRthirtyFullOverScratch}{\pending{X->bs fits and from-scratch runs}}  % experiments/FIGS/data/v2/index.json :: not computed yet :: sha256 c6c4138ee5c86423 :: waits on X->bs fits and from-scratch runs
\newcommand{\XbsRthirtyFullOverScratchPr}{\pending{X->bs fits and from-scratch runs}}  % experiments/FIGS/data/v2/index.json :: not computed yet :: sha256 c6c4138ee5c86423 :: waits on X->bs fits and from-scratch runs
\newcommand{\XbsRthirtyCoarseOverFull}{\pending{X->bs fits}}  % experiments/FIGS/data/v2/index.json :: not computed yet :: sha256 c6c4138ee5c86423 :: waits on X->bs fits
\newcommand{\XbsRthirtyCoarseOverFullPr}{\pending{X->bs fits}}  % experiments/FIGS/data/v2/index.json :: not computed yet :: sha256 c6c4138ee5c86423 :: waits on X->bs fits
\newcommand{\XbsRthirtyNodeOverTransferFull}{\pending{X->bs fits}}  % experiments/FIGS/data/v2/index.json :: not computed yet :: sha256 c6c4138ee5c86423 :: waits on X->bs fits
\newcommand{\XbsRthirtyNodeOverTransferFullPr}{\pending{X->bs fits}}  % experiments/FIGS/data/v2/index.json :: not computed yet :: sha256 c6c4138ee5c86423 :: waits on X->bs fits
\newcommand{\XbbRthirtyDirectOverTransferFull}{\pending{X->bb fits}}  % experiments/FIGS/data/v2/index.json :: not computed yet :: sha256 c6c4138ee5c86423 :: waits on X->bb fits
\newcommand{\XbbRthirtyDirectOverTransferFullPr}{\pending{X->bb fits}}  % experiments/FIGS/data/v2/index.json :: not computed yet :: sha256 c6c4138ee5c86423 :: waits on X->bb fits
\newcommand{\XbbRthirtyTransferCoarseOverFull}{\pending{X->bb fits}}  % experiments/FIGS/data/v2/index.json :: not computed yet :: sha256 c6c4138ee5c86423 :: waits on X->bb fits
\newcommand{\XbbRthirtyTransferCoarseOverFullPr}{\pending{X->bb fits}}  % experiments/FIGS/data/v2/index.json :: not computed yet :: sha256 c6c4138ee5c86423 :: waits on X->bb fits
             % the counts below are its macros
% Visual language: Sophon's Fig. 1 (arXiv:2405.12972): the document's serif font, light fills,
% rounded boxes, light grey panels behind the stages (PI 2026-11-02). The text is set in the
% document's font at the printed sizes 8.5-10 pt (MLST 8-12 pt), never scaled.
% Text width: 15 cm (MLST full width). The picture is 15.0 cm wide and is drawn at its final
% size, never scaled, so the font sizes hold.
% Optional: \def\FoneDropXbx{} before % Figure 1 (F1 of docs/FIGURE_PLAN.md): the design of the study, in TikZ.
% % Figure 1 (F1 of docs/FIGURE_PLAN.md): the design of the study, in TikZ.
% \input{fig_f1} from main.tex; this file holds one figure environment and nothing else.
%
% WHAT main.tex MUST LOAD (this file loads nothing itself):
%   \usepackage{tikz}
%   \usetikzlibrary{arrows.meta}          % the only TikZ library used (arrow tips)
%   \definecolor{c188}{HTML}{1F77B4}      % the configuration colours of experiments/FIGS/
%   \definecolor{c162}{HTML}{FF7F0E}      % mlst_style.py (matplotlib's tab10 in order, as in
%   \definecolor{c43}{HTML}{2CA02C}       % Sophon's figures); the variants are drawn in their
%   \definecolor{c17}{HTML}{D62728}       % label set's colour with a dotted border, as the
%   \definecolor{crand}{HTML}{17BECF}     % plotted figures draw variants dotted
%   \definecolor{cflav}{HTML}{9467BD}
%   \input{results_generated}             % the counts below are its macros
% Visual language: Sophon's Fig. 1 (arXiv:2405.12972): the document's serif font, light fills,
% rounded boxes, light grey panels behind the stages (PI 2026-11-02). The text is set in the
% document's font at the printed sizes 8.5-10 pt (MLST 8-12 pt), never scaled.
% Text width: 15 cm (MLST full width). The picture is 15.0 cm wide and is drawn at its final
% size, never scaled, so the font sizes hold.
% Optional: \def\FoneDropXbx{} before \input{fig_f1} removes X->bb/X->bs from panel (c)
% (PI 2026-11-02: if the samples are not staged by 12 Nov, the X->bb/X->bs section goes).
%
% Every class box is drawn from the label maps (checked by hand, 2026-11-02):
%   (a) configs/labelmaps/rung_label_maps.v1.csv, columns L188, L162, R42_Q1, R16_Q1. Rows: the ten
%       decays X->QQ/gg in the order bb bc bq | cc cs cq | ss sq qq gg (the 43-class b / c / light
%       groups are then contiguous), then ee mm | tauh-taue tauh-taum, then the 27 QCD classes.
%   (b) column: configs/labelmaps/rand_label_map.v2.csv, RAND2_p1 on the same rows: groups
%       0 0 0 0 0 9 9 9 0 3 | 11 11 | 15 4 | QCD, lettered A..F in order of first appearance.
%       strips: configs/labelmaps/flavour_pair_map.v2.csv, the 22 classes of orbit X_YY_QQQQ.
%       Order with a b quark: bbbb bbcc bbss bbqq qqqb bbcq ccbs ccbq qqbc qqbs bcsq;
%       without: cccc ccss ccqq ssss ssqq qqqq qqqc qqqs ccsq sscq qqcs.
%       FLAV_F0: all 22 in group 9. FLAV_F1: the 11 with a b in group 11, the rest 9.
%       FLAV_F1R: group 11 = bbss bbqq qqqb bbcq ccbq qqbc (with b, positions 2-5, 7, 8 from 0)
%       and ccss ccqq ssss qqqc sscq (without, positions 1-3, 6, 9); filled square = group 11.
% Geometry (cm): row height 0.35 from y = -1.18, gaps of 0.14 between the quark/gluon, lepton
% and QCD blocks, the QCD row 0.70; boxes inset 0.025 top and bottom; columns 1.45 wide with
% left edges 1.75, 3.28, 4.81, 6.34 (hierarchy) and 8.25 (random partition).
\providecommand{\Kfull}{\VocabSizeLoneeighteight}
\providecommand{\Kres}{\VocabSizeLonesixtwo}
\providecommand{\Kmid}{\VocabSizeRfourtwoQone}
\providecommand{\Kcoarse}{\VocabSizeRonesixQone}
\begin{figure}[tbp]
\centering
\ifdim\linewidth<14.99cm\typeout{fig_f1.tex: WARNING: the drawing is 15 cm wide but \string\linewidth\space
  is \the\linewidth; set a 15 cm text width (it is not scaled, to keep its 8.5 pt text)}\fi
\def\FoneText{\rmfamily\fontsize{8.5}{9.5}\selectfont}%
\def\FoneMain{\rmfamily\fontsize{9}{10.5}\selectfont}%
\def\FoneBold{\rmfamily\bfseries\fontsize{9}{10.5}\selectfont}%
\def\FoneTitle{\rmfamily\bfseries\fontsize{10}{12}\selectfont}%
\ifdefined\FoneDropXbx\def\FoneXbx{}\else\def\FoneXbx{\\$X\to b\bar b$, $X\to bs$}\fi
\begin{tikzpicture}[
  x=1cm, y=1cm,
  font=\FoneMain,
  ttl/.style={anchor=base west, font=\FoneTitle, inner sep=0pt},
  cls/.style={draw=black!55, line width=0.45pt, rounded corners=1.2pt},
  hdr/.style={draw=#1!70!black, fill=#1!22, line width=0.6pt, rounded corners=2.5pt, align=center,
              font=\FoneText, inner sep=0pt},
  inbox/.style={align=center, font=\FoneText, inner sep=0pt},
  letter/.style={font=\FoneBold, inner sep=0pt},
  rowlab/.style={anchor=east, inner sep=0pt},
  txt/.style={align=left, font=\FoneText, inner sep=0pt},
  sq/.style={draw=black!65, line width=0.4pt},
  chip/.style={draw=#1!70!black, fill=#1!22, line width=0.6pt, rounded corners=3pt,
               align=center, font=\FoneText, inner sep=0pt, anchor=north west},
  stage/.style={draw=black!45, fill=black!4, line width=0.6pt, rounded corners=3pt, align=left,
                font=\FoneText, inner sep=0.1cm},
  flow/.style={-{Latex[length=1.8mm, width=1.4mm]}, line width=0.7pt, black!75},
]
% ================================================================ (a) the class hierarchy
\node[ttl] at (0,0) {(a) Class hierarchy};
% faint bands on every other row, drawn first so the boxes sit on top
\foreach \rowtop in {-1.53,-2.23,-2.93,-3.63,-4.33,-5.17,-5.87}
  \fill[black!4] (0,\rowtop) rectangle (9.70,\rowtop-0.35);
% column headers
\node[rowlab, align=right, font=\FoneText] at (1.65,-0.68) {JetClass-II\\class};
\node[hdr=c188, minimum width=1.45cm, minimum height=0.80cm, anchor=north west] at (1.75,-0.28) {\Kfull\\classes};
\node[hdr=c162, minimum width=1.45cm, minimum height=0.80cm, anchor=north west] at (3.28,-0.28) {\Kres\\classes};
\node[hdr=c43,  minimum width=1.45cm, minimum height=0.80cm, anchor=north west] at (4.81,-0.28) {\Kmid\\classes};
\node[hdr=c17,  minimum width=1.45cm, minimum height=0.80cm, anchor=north west] at (6.34,-0.28) {\Kcoarse\\classes};
% row labels
\node[rowlab] at (1.65,-1.355) {$X\to b\bar b$};
\node[rowlab] at (1.65,-1.705) {$X\to bc$};
\node[rowlab] at (1.65,-2.055) {$X\to bq$};
\node[rowlab] at (1.65,-2.405) {$X\to c\bar c$};
\node[rowlab] at (1.65,-2.755) {$X\to cs$};
\node[rowlab] at (1.65,-3.105) {$X\to cq$};
\node[rowlab] at (1.65,-3.455) {$X\to s\bar s$};
\node[rowlab] at (1.65,-3.805) {$X\to sq$};
\node[rowlab] at (1.65,-4.155) {$X\to q\bar q$};
\node[rowlab] at (1.65,-4.505) {$X\to gg$};
\node[rowlab] at (1.65,-4.995) {$X\to ee$};
\node[rowlab] at (1.65,-5.345) {$X\to\mu\mu$};
\node[rowlab] at (1.65,-5.695) {$X\to\tau_h\tau_e$};
\node[rowlab] at (1.65,-6.045) {$X\to\tau_h\tau_\mu$};
\node[rowlab] at (1.65,-6.710) {QCD jets};
% 188 and 162 classes: every resonant decay is its own class
\foreach \rowtop in {-1.18,-1.53,-1.88,-2.23,-2.58,-2.93,-3.28,-3.63,-3.98,-4.33,-4.82,-5.17,-5.52,-5.87} {
  \draw[cls, fill=c188!15] (1.75,\rowtop-0.025) rectangle (3.20,\rowtop-0.325);
  \draw[cls, fill=c162!15] (3.28,\rowtop-0.025) rectangle (4.73,\rowtop-0.325);
}
\draw[cls, fill=c188!15] (1.75,-6.385) rectangle (3.20,-7.035);
\node[inbox] at (2.475,-6.71) {\LitJetClassTwoNQcdClasses\\classes};
\draw[cls, fill=c162!15] (3.28,-6.385) rectangle (4.73,-7.035);
\node[inbox] at (4.005,-6.71) {1 class};
% 43 classes: decay topology x {b, else c, else light quark or gluon}; e and mu merged
\draw[cls, fill=c43!15] (4.81,-1.205) rectangle (6.26,-2.205);
\node[inbox] at (5.535,-1.705) {with $b$};
\draw[cls, fill=c43!15] (4.81,-2.255) rectangle (6.26,-3.255);
\node[inbox] at (5.535,-2.755) {with $c$,\\no $b$};
\draw[cls, fill=c43!15] (4.81,-3.305) rectangle (6.26,-4.655);
\node[inbox] at (5.535,-3.980) {$u$, $d$, $s$\\or $g$};
\draw[cls, fill=c43!15] (4.81,-4.845) rectangle (6.26,-5.495);
\node[inbox] at (5.535,-5.170) {$\ell\ell$};
\draw[cls, fill=c43!15] (4.81,-5.545) rectangle (6.26,-6.195);
\node[inbox] at (5.535,-5.870) {$\tau_h\tau_\ell$};
\draw[cls, fill=c43!15] (4.81,-6.385) rectangle (6.26,-7.035);
\node[inbox] at (5.535,-6.710) {1 class};
% 17 classes: number and type of prongs, no flavor
\draw[cls, fill=c17!15] (6.34,-1.205) rectangle (7.79,-4.655);
\node[inbox] at (7.065,-2.930) {two\\quarks\\or gluons};
\draw[cls, fill=c17!15] (6.34,-4.845) rectangle (7.79,-5.495);
\node[inbox] at (7.065,-5.170) {$\ell\ell$};
\draw[cls, fill=c17!15] (6.34,-5.545) rectangle (7.79,-6.195);
\node[inbox] at (7.065,-5.870) {$\tau_h\tau_\ell$};
\draw[cls, fill=c17!15] (6.34,-6.385) rectangle (7.79,-7.035);
\node[inbox] at (7.065,-6.710) {1 class};
% what each merge removes
\node[anchor=north west, font=\FoneBold, inner sep=0pt] at (0,-7.27) {What each merge removes};
\node[txt, anchor=north west, text width=6.04cm] (m1) at (1.75,-7.65)
  {the $b$, $c$ and $s$ content of the QCD jets};
\node[txt, anchor=north west, text width=6.04cm] (m2) at ([yshift=-0.10cm]m1.south west)
  {the number of $b$ and $c$ quarks, light-quark flavor, quark against gluon, and lepton
   flavor ($e$ against $\mu$)};
\node[txt, anchor=north west, text width=6.04cm] (m3) at ([yshift=-0.10cm]m2.south west)
  {quark flavor ($b$, $c$ or light); the number and type of prongs remain};
\node[txt, anchor=north east] at ([xshift=-0.10cm]m1.north west) {\Kfull\,$\to$\,\Kres};
\node[txt, anchor=north east] at ([xshift=-0.10cm]m2.north west) {\Kres\,$\to$\,\Kmid};
\node[txt, anchor=north east] at ([xshift=-0.10cm]m3.north west) {\Kmid\,$\to$\,\Kcoarse};
% ================================================================ (b) the controls
\node[ttl] at (8.25,0) {(b) Controls, \Kcoarse{} classes each};
% the first random partition on the rows of (a); one letter = one class
\node[hdr=crand, minimum width=1.45cm, minimum height=0.80cm, anchor=north west] at (8.25,-0.28) {random\\partition};
\draw[cls, fill=crand!15] (8.25,-1.205) rectangle (9.70,-2.905);
\node[letter] at (8.975,-2.055) {A};
\draw[cls, fill=crand!15] (8.25,-2.955) rectangle (9.70,-3.955);
\node[letter] at (8.975,-3.455) {B};
\draw[cls, fill=crand!15] (8.25,-4.005) rectangle (9.70,-4.305);
\node[letter] at (8.975,-4.155) {A};
\draw[cls, fill=crand!15] (8.25,-4.355) rectangle (9.70,-4.655);
\node[letter] at (8.975,-4.505) {C};
\draw[cls, fill=crand!15] (8.25,-4.845) rectangle (9.70,-5.495);
\node[letter] at (8.975,-5.170) {D};
\draw[cls, fill=crand!15] (8.25,-5.545) rectangle (9.70,-5.845);
\node[letter] at (8.975,-5.695) {E};
\draw[cls, fill=crand!15] (8.25,-5.895) rectangle (9.70,-6.195);
\node[letter] at (8.975,-6.045) {F};
\draw[cls, fill=crand!15] (8.25,-6.385) rectangle (9.70,-7.035);
\node[inbox] at (8.975,-6.710) {1 class};
% how the random partitions are built
\node[hdr=crand, anchor=north west, text width=4.75cm, inner sep=0.1cm, align=left, font=\FoneBold]
  (rh) at (10.05,-0.28) {Random partitions (\RandVtwoNPartitions)};
\node[txt, anchor=north west, text width=4.95cm] (rb) at ([yshift=-0.12cm]rh.south west)
  {\textbullet\ groups of resonant classes drawn at random, plus one class of QCD jets\\[2pt]
   \textbullet\ group shares of training jets within \RandVtwoShareTol{} of the \Kcoarse-class groups\\[2pt]
   \textbullet\ best of \RandVtwoTrials{} draws: the fewest merges shared with the \Kcoarse-class level\\[2pt]
   \textbullet\ left: the first partition; boxes with the same letter are one class};
% how the random partitions with controlled flavor boundaries are built
\node[hdr=cflav, anchor=north west, text width=4.75cm, inner sep=0.1cm, align=left, font=\FoneBold]
  (fh) at ([yshift=-0.22cm]rb.south west) {Random partitions with controlled flavor boundaries};
\node[txt, anchor=north west, text width=4.95cm] (fn) at ([yshift=-0.12cm]fh.south west)
  {whole families of decays that differ only in quark flavor; two of them split one family,
   the \RandFlavSplitOrbitClasses{} decays \RandFlavSplitOrbit:};
\begin{scope}[shift={(fn.south west)}]
  % brackets over the decays with and without a b quark (squares at a 0.2 cm pitch)
  \node[txt, anchor=south] at (1.075,-0.45) {with $b$};
  \node[txt, anchor=south] at (3.425,-0.45) {without $b$};
  \draw[line width=0.5pt] (0,-0.54) -- (0,-0.47) -- (2.15,-0.47) -- (2.15,-0.54);
  \draw[line width=0.5pt] (2.35,-0.54) -- (2.35,-0.47) -- (4.50,-0.47) -- (4.50,-0.54);
  % no flavor boundary: all 22 in one class
  \node[txt, anchor=north west] at (0,-0.62) {no flavor boundary};
  \foreach \i in {0,...,10} {
    \draw[sq, fill=white] (0.2*\i,-1.10) rectangle ++(0.15,0.15);
    \draw[sq, fill=white] (2.35+0.2*\i,-1.10) rectangle ++(0.15,0.15);
  }
  % b-quark boundary: the decays with a b quark moved to another class
  \node[txt, anchor=north west] at (0,-1.22) {$b$-quark boundary};
  \foreach \i in {0,...,10} {
    \draw[sq, fill=cflav!70!black] (0.2*\i,-1.70) rectangle ++(0.15,0.15);
    \draw[sq, fill=white] (2.35+0.2*\i,-1.70) rectangle ++(0.15,0.15);
  }
  % random flavor boundary: the same number of decays, drawn at random
  \node[txt, anchor=north west] at (0,-1.82) {random flavor boundary};
  \foreach \i in {0,1,6,9,10}
    \draw[sq, fill=white] (0.2*\i,-2.30) rectangle ++(0.15,0.15);
  \foreach \i in {2,3,4,5,7,8}
    \draw[sq, fill=cflav!70!black] (0.2*\i,-2.30) rectangle ++(0.15,0.15);
  \foreach \i in {0,4,5,7,8,10}
    \draw[sq, fill=white] (2.35+0.2*\i,-2.30) rectangle ++(0.15,0.15);
  \foreach \i in {1,2,3,6,9}
    \draw[sq, fill=cflav!70!black] (2.35+0.2*\i,-2.30) rectangle ++(0.15,0.15);
  % key
  \draw[sq, fill=cflav!70!black] (0,-2.62) rectangle ++(0.15,0.15);
  \node[txt, anchor=west] (fkey) at (0.25,-2.545) {moved to another class};
\end{scope}
% ================================================================ (c) the pipeline
% (c) starts below the lower of the two columns above, whatever their text heights
\pgfpointanchor{m3}{south}\pgfgetlastxy{\FoneXa}{\FoneYa}%
\pgfpointanchor{fkey}{south}\pgfgetlastxy{\FoneXb}{\FoneYb}%
\pgfmathsetlengthmacro{\FoneYc}{min(\FoneYa,\FoneYb)-0.45cm}%
\coordinate (ctop) at (0pt,\FoneYc);
\begin{scope}[shift={(ctop)}]
  \node[ttl, anchor=north west] at (0,0) {(c) Pretraining and evaluation};
  % what every model of one seed shares
  \node[stage, text width=2.15cm] (seed) at (1.175,-1.96)
    {\textbf{Same for every model of one seed:} the training jets, their order and the
     initial weights, except the output layer};
  % one pretraining per label set: the configuration colours of every later figure
  \node[draw=black!45, fill=black!4, line width=0.6pt, rounded corners=3pt, minimum width=5.15cm,
        minimum height=3.08cm, anchor=north west, inner sep=0pt] (pre) at (2.81,-0.42) {};
  \node[anchor=north west, font=\FoneBold, inner sep=0pt] at (2.96,-0.52) {Pretraining on JetClass-II};
  \node[chip=c188, minimum width=1.13cm, minimum height=0.42cm] at (2.96,-0.95) {\Kfull};
  \node[chip=c162, minimum width=1.13cm, minimum height=0.42cm] at (4.20,-0.95) {\Kres};
  \node[chip=c43,  minimum width=1.13cm, minimum height=0.42cm] at (5.44,-0.95) {\Kmid};
  \node[chip=c17,  minimum width=1.13cm, minimum height=0.42cm] at (6.68,-0.95) {\Kcoarse};
  \node[chip=crand, minimum width=1.70cm, minimum height=0.70cm] at (2.96,-1.47) {random\\partitions};
  \node[chip=cflav, minimum width=3.04cm, minimum height=0.70cm] at (4.77,-1.47) {controlled flavor\\boundaries};
  \node[chip=c162, densely dotted, minimum width=1.13cm, minimum height=0.42cm] at (2.96,-2.32) {\Kres};
  \node[chip=c17,  densely dotted, minimum width=1.13cm, minimum height=0.42cm] at (4.20,-2.32) {\Kcoarse};
  \node[txt, anchor=west] at (5.44,-2.53) {with auxiliary\\mass regression};
  \node[chip=c188, densely dotted, minimum width=1.13cm, minimum height=0.42cm] at (2.96,-2.94) {\Kfull};
  \node[txt, anchor=west] at (4.20,-3.15) {kinematic inputs only};
  % evaluation
  \node[stage, text width=2.4cm] (fz) at (9.68,-0.95) {frozen features:\\linear probe,\\MLP probe};
  \node[stage, text width=2.4cm] (ft) at (9.68,-2.05) {fine-tuning\\of all weights};
  \node[stage, text width=2.4cm, draw=black, fill=white, line width=1.1pt] (sc) at (9.68,-3.15) {trained from\\scratch};
  \node[stage, text width=3.4cm, minimum height=3.2cm] (task) at (13.20,-2.00)
    {\textbf{Downstream tasks}\\JetClass-II: all classes\\JetClass-II: class pairs\\\quad(both class-balanced)\\JetClass\\
     top tagging\\quark/gluon, Pythia\\quark/gluon, Herwig\\jet mass\FoneXbx};
  \draw[flow] (seed.east) -- (seed.east -| pre.west);
  \draw[flow] (fz.west -| pre.east) -- (fz.west);
  \draw[flow] (ft.west -| pre.east) -- (ft.west);
  \draw[flow] (fz.east) -- (fz.east -| task.west);
  \draw[flow] (ft.east) -- (ft.east -| task.west);
  \draw[flow] (sc.east) -- (sc.east -| task.west);
\end{scope}
\end{tikzpicture}
\caption{Design of the study.
(a)~The class hierarchy on the decays of a heavy resonance~$X$ and on the jets from quantum
chromodynamics (QCD); $q$ is a $u$ or $d$ quark, $\ell$ an electron or a muon, and $\tau_h$,
$\tau_e$, $\tau_\mu$ a $\tau$ lepton decaying to hadrons, an electron or a muon. Each box is one
class of the label set above its column, except the QCD box of the \Kfull-class column (Sophon's
label set~\cite{sophon}), which stands for \LitJetClassTwoNQcdClasses{} classes. Each level merges
classes of the level to its left; the lines below name what each merge removes.
(b)~The controls: the first random partition, where boxes with the same letter are one class, and
the random partitions with controlled flavor boundaries on the decays \RandFlavSplitOrbit{} (one
square per decay; filled squares are moved to a second class).
(c)~Models of one seed share the training jets, their order and the initial weights, except the
output layer, the mass regression target, and the first layer and initial weights of the model with
kinematic inputs only (the particle four-momenta). The hierarchy levels have \MlstNRunsWord{} seeds
each, the models with the auxiliary mass regression \MlstMassNRunsWord{}, the random partitions with
controlled flavor boundaries \RandFlavRuns{} and each random partition \RandVtwoRuns{}. Frozen
features, the output of the backbone (the network below the classification head) with its weights held
fixed, feed linear probes (logistic regressions) and multilayer perceptron (MLP) probes; fine-tuning
trains all weights. Colors identify the label sets and controls in every figure; the plots give the three partitions with
controlled flavor boundaries three shades; dotted borders mark the two
variants, which the other figures draw as dotted lines in the color of their label set.}
\label{fig:design}
\end{figure}
 from main.tex; this file holds one figure environment and nothing else.
%
% WHAT main.tex MUST LOAD (this file loads nothing itself):
%   \usepackage{tikz}
%   \usetikzlibrary{arrows.meta}          % the only TikZ library used (arrow tips)
%   \definecolor{c188}{HTML}{1F77B4}      % the configuration colours of experiments/FIGS/
%   \definecolor{c162}{HTML}{FF7F0E}      % mlst_style.py (matplotlib's tab10 in order, as in
%   \definecolor{c43}{HTML}{2CA02C}       % Sophon's figures); the variants are drawn in their
%   \definecolor{c17}{HTML}{D62728}       % label set's colour with a dotted border, as the
%   \definecolor{crand}{HTML}{17BECF}     % plotted figures draw variants dotted
%   \definecolor{cflav}{HTML}{9467BD}
%   % GENERATED by experiments/FIGS/make_tables.py -- do not edit.
% Every line carries: source file :: path inside that file :: first 16 hex of its
% sha256. paper/journal/provenance.json carries the full hash beside the value.
% Regenerate with:  python3 experiments/FIGS/make_tables.py
% Check in CI with: python3 experiments/FIGS/make_tables.py --check
%
% Macro names spell digits as words (162 -> Onesixtwo) because \newcommand takes
% letters only. A number whose source file does not exist yet is NOT emitted --
% see the missing list below, and leave a \pending{} in the text for it.
%
% MISSING INPUTS -- no macro was emitted for these:
%   community benchmarks, bench v3 readout -- experiments/FIGS/data/bench_v3_metrics*/bench_metrics_last.json
%   v2 pretraining rerun: data order and dropout fixed by seed and epoch, identical initial trunk weights, each epoch reading one of five disjoint random parts of the rows of nearly every file, fixed validation sample, checkpoint at the first maximum within the last ten epochs -- experiments/FIGS/data/v2/pretraining/
%   v2 frozen probes on the rerun models, at the first maximum within the last ten epochs and on the model whose weights are the average of the last ten epochs, BatchNorm statistics recomputed, with untrained-network and pooled-embedding reference rows -- experiments/FIGS/data/v2/probe_ladder/
%   v2 paired ratios with run and test-sample errors on the rerun models -- experiments/FIGS/data/v2/paired_errors/
%   v2 the 64- and 30-class vocabularies -- experiments/FIGS/data/v2/levels_64_30/
%   v2 the leave-one-family-out vocabularies -- experiments/FIGS/data/v2/leave_one_family_out/
%   v2 label-tree recovery on the rerun models -- experiments/FIGS/data/v2/label_recovery/
%   v2 the rerun's random partitions and flavour pair -- experiments/FIGS/data/v2/random_partitions/
%   v2 JetClass-II fine-tuning on files held out from pretraining -- experiments/FIGS/data/v2/finetune/
%   v2 fine-tuning from random initialisation and from the self-supervised model -- experiments/FIGS/data/v2/finetune_references/
%   v2 top tagging and quark/gluon tagging (Pythia and Herwig) fine-tuned from every rerun model -- experiments/FIGS/data/v2/benchmarks/
%   v2 the second and third self-supervised pretraining runs -- experiments/FIGS/data/v2/self_supervised/
%   v2 the mass-output models with the loss weight matched in loss share -- experiments/FIGS/data/v2/mass_lambda_matched/
%   v2 anomaly detection on the rerun models, per-run values -- experiments/FIGS/data/v2/anomaly/
%   v2 the open-data fits on the rerun models -- experiments/FIGS/data/v2/real_data/
%   MassLossShareVtwoOnesixtwoMass -- experiments/FIGS/data/v2/mass_lambda_matched/loss_share.json :: shares.162+mass
%   MassLossShareVtwoOnesevenMass -- experiments/FIGS/data/v2/mass_lambda_matched/loss_share.json :: shares.17+mass
%   MassLossShareVtwoOnesevenMassMatched -- experiments/FIGS/data/v2/mass_lambda_matched/loss_share.json :: shares.17+mass_matched
%   DesignVtwoSecondGpu, DesignVtwoSecondGpuRuns -- the grid's job specs put every run index on one product
%   per-arm parameters and MACs -- flops_by_arm.json (see docs/RECORD.md 2.1)
%   fine-tuning read-outs -- experiments/FIGS/data/v2/ft_mlst/: jc2, jc2_qcd
%   top-tagging recipe check -- ft_mlst/topcheck_bench_metrics.json
%   X->bb and X->bs fits -- experiments/FIGS/data/v2/xbx/
%   X->bs from scratch -- experiments/FIGS/data/v2/xbx_scratch/
%   JetClass accuracy bootstrap -- experiments/FIGS/data/v2/budget/jc1_acc_boot_N10000.json
%   MlstFtJcOneErrCoarseOverFullBootEFour -- waits on the JetClass bootstrap at 1e4 (after the equal-stopping read-out)
%   MlstLpBbQgZeroAccDropFull -- waits on pair probes with the quark/gluon inputs zeroed (open item O10, not run)
%   MlstLpBbQgZeroAccDropCoarse -- waits on pair probes with the quark/gluon inputs zeroed (open item O10, not run)
%   XbsRthirtyFullOverScratch -- waits on X->bs fits and from-scratch runs
%   XbsRthirtyFullOverScratchPr -- waits on X->bs fits and from-scratch runs
%   XbsRthirtyCoarseOverFull -- waits on X->bs fits
%   XbsRthirtyCoarseOverFullPr -- waits on X->bs fits
%   XbsRthirtyNodeOverTransferFull -- waits on X->bs fits
%   XbsRthirtyNodeOverTransferFullPr -- waits on X->bs fits
%   XbbRthirtyDirectOverTransferFull -- waits on X->bb fits
%   XbbRthirtyDirectOverTransferFullPr -- waits on X->bb fits
%   XbbRthirtyTransferCoarseOverFull -- waits on X->bb fits
%   XbbRthirtyTransferCoarseOverFullPr -- waits on X->bb fits
%
% SKIPPED TABLES:
%   T4 compute/parameters: no machine-readable source in the repository; docs/RECORD.md 2.1 names /data/results/mtx/flops/flops_by_arm.json, which is not committed

\newcommand{\PendingPretrain}{\pending{pending v2: pretraining rerun: data order and dropout fixed by seed and epoch, identical initial trunk weights, each epoch reading one of five disjoint random parts of the rows of nearly every file, fixed validation sample, checkpoint at the first maximum within the last ten epochs}}  % experiments/FIGS/make_tables.py :: V2_PENDING['Pretrain'] :: sha256 10e6ec08356a4a2b :: a marker, not a number: red in the draft until the input exists
\newcommand{\PendingProbes}{\pending{pending v2: frozen probes on the rerun models, at the first maximum within the last ten epochs and on the model whose weights are the average of the last ten epochs, BatchNorm statistics recomputed, with untrained-network and pooled-embedding reference rows}}  % experiments/FIGS/make_tables.py :: V2_PENDING['Probes'] :: sha256 10e6ec08356a4a2b :: a marker, not a number: red in the draft until the input exists
\newcommand{\PendingPaired}{\pending{pending v2: paired ratios with run and test-sample errors on the rerun models}}  % experiments/FIGS/make_tables.py :: V2_PENDING['Paired'] :: sha256 10e6ec08356a4a2b :: a marker, not a number: red in the draft until the input exists
\newcommand{\PendingLevels}{\pending{pending v2: the 64- and 30-class vocabularies}}  % experiments/FIGS/make_tables.py :: V2_PENDING['Levels'] :: sha256 10e6ec08356a4a2b :: a marker, not a number: red in the draft until the input exists
\newcommand{\PendingLofo}{\pending{pending v2: the leave-one-family-out vocabularies}}  % experiments/FIGS/make_tables.py :: V2_PENDING['Lofo'] :: sha256 10e6ec08356a4a2b :: a marker, not a number: red in the draft until the input exists
\newcommand{\PendingRecovery}{\pending{pending v2: label-tree recovery on the rerun models}}  % experiments/FIGS/make_tables.py :: V2_PENDING['Recovery'] :: sha256 10e6ec08356a4a2b :: a marker, not a number: red in the draft until the input exists
\newcommand{\PendingRandom}{\pending{pending v2: the rerun's random partitions and flavour pair}}  % experiments/FIGS/make_tables.py :: V2_PENDING['Random'] :: sha256 10e6ec08356a4a2b :: a marker, not a number: red in the draft until the input exists
\newcommand{\PendingFtHeldout}{\pending{pending v2: JetClass-II fine-tuning on files held out from pretraining}}  % experiments/FIGS/make_tables.py :: V2_PENDING['FtHeldout'] :: sha256 10e6ec08356a4a2b :: a marker, not a number: red in the draft until the input exists
\newcommand{\PendingFtRefs}{\pending{pending v2: fine-tuning from random initialisation and from the self-supervised model}}  % experiments/FIGS/make_tables.py :: V2_PENDING['FtRefs'] :: sha256 10e6ec08356a4a2b :: a marker, not a number: red in the draft until the input exists
\newcommand{\PendingBench}{\pending{pending v2: top tagging and quark/gluon tagging (Pythia and Herwig) fine-tuned from every rerun model}}  % experiments/FIGS/make_tables.py :: V2_PENDING['Bench'] :: sha256 10e6ec08356a4a2b :: a marker, not a number: red in the draft until the input exists
\newcommand{\PendingSsl}{\pending{pending v2: the second and third self-supervised pretraining runs}}  % experiments/FIGS/make_tables.py :: V2_PENDING['Ssl'] :: sha256 10e6ec08356a4a2b :: a marker, not a number: red in the draft until the input exists
\newcommand{\PendingMassLambda}{\pending{pending v2: the mass-output models with the loss weight matched in loss share}}  % experiments/FIGS/make_tables.py :: V2_PENDING['MassLambda'] :: sha256 10e6ec08356a4a2b :: a marker, not a number: red in the draft until the input exists
\newcommand{\PendingAnomaly}{\pending{pending v2: anomaly detection on the rerun models, per-run values}}  % experiments/FIGS/make_tables.py :: V2_PENDING['Anomaly'] :: sha256 10e6ec08356a4a2b :: a marker, not a number: red in the draft until the input exists
\newcommand{\PendingRealData}{\pending{pending v2: the open-data fits on the rerun models}}  % experiments/FIGS/make_tables.py :: V2_PENDING['RealData'] :: sha256 10e6ec08356a4a2b :: a marker, not a number: red in the draft until the input exists
\newcommand{\PendingConclusion}{\pending{pending v2: written once every input above is in}}  % experiments/FIGS/make_tables.py :: V2_PENDING['Conclusion'] :: sha256 10e6ec08356a4a2b :: a marker, not a number: red in the draft until the input exists
\newcommand{\AojChainDataset}{178}  % experiments/FIGS/data/aoj_checks_v2/selection_chain_v2.json :: n_dataset :: sha256 729d83085aba25f8
\newcommand{\AojChainStaged}{12{,}589{,}467}  % experiments/FIGS/data/aoj_checks_v2/selection_chain_v2.json :: n_staged :: sha256 729d83085aba25f8
\newcommand{\AojChainRhoWindow}{9{,}441{,}209}  % experiments/FIGS/data/aoj_checks_v2/selection_chain_v2.json :: n_rho_window :: sha256 729d83085aba25f8
\newcommand{\AojNJetsFit}{4{,}194{,}846}  % experiments/FIGS/data/aoj_checks_v2/selection_chain_v2.json :: n_fit :: sha256 729d83085aba25f8
\newcommand{\AojChainPtBins}{8}  % experiments/FIGS/data/aoj_checks_v2/selection_chain_v2.json :: n_pt_bins :: sha256 729d83085aba25f8
\newcommand{\AojChainEtaMax}{2.40}  % experiments/FIGS/data/aoj_checks_v2/selection_chain_v2.json :: abs_eta_max :: sha256 729d83085aba25f8
\newcommand{\MassLossShareVtwoOnesixtwoMass}{\pending{pending: realised mass-loss share in the rerun, 162 classes}}  % experiments/FIGS/make_tables.py :: SLOTS[MassLossShareVtwoOnesixtwoMass] :: sha256 10e6ec08356a4a2b :: a marker, not a number: experiments/FIGS/data/v2/mass_lambda_matched/loss_share.json :: shares.162+mass does not exist yet
\newcommand{\MassLossShareVtwoOnesevenMass}{\pending{pending: realised mass-loss share in the rerun, 17 classes}}  % experiments/FIGS/make_tables.py :: SLOTS[MassLossShareVtwoOnesevenMass] :: sha256 10e6ec08356a4a2b :: a marker, not a number: experiments/FIGS/data/v2/mass_lambda_matched/loss_share.json :: shares.17+mass does not exist yet
\newcommand{\MassLossShareVtwoOnesevenMassMatched}{\pending{pending: realised mass-loss share in the rerun, 17 classes, matched weight}}  % experiments/FIGS/make_tables.py :: SLOTS[MassLossShareVtwoOnesevenMassMatched] :: sha256 10e6ec08356a4a2b :: a marker, not a number: experiments/FIGS/data/v2/mass_lambda_matched/loss_share.json :: shares.17+mass_matched does not exist yet
\newcommand{\FtCoverageEThreeNClasses}{160}  % experiments/FT/data/jc2_v2_manifest.json :: class_coverage['train_N1000_s1.parquet'].n_classes_present :: sha256 4bd68655919b14e9
\newcommand{\FtCoverageEFourNClasses}{183}  % experiments/FT/data/jc2_v2_manifest.json :: class_coverage['train_N10000_s1.parquet'].n_classes_present :: sha256 4bd68655919b14e9
\newcommand{\FtCoverageEThreeMedian}{4}  % experiments/FT/data/jc2_v2_manifest.json :: class_coverage['train_N1000_s1.parquet'].median_per_present_class :: sha256 4bd68655919b14e9
\newcommand{\FtCoverageEThreeOneJet}{23}  % experiments/FT/data/jc2_v2_manifest.json :: class_coverage['train_N1000_s1.parquet'].n_classes_with_one_jet :: sha256 4bd68655919b14e9
\newcommand{\FtCoverageVoneEThreeNClasses}{161}  % experiments/FT/data/jc2_v1_coverage.json :: class_coverage['train_N1000_s1.parquet'].n_classes_present :: sha256 82652f57f708dc23
\newcommand{\FtCoverageVoneEThreeMedian}{5}  % experiments/FT/data/jc2_v1_coverage.json :: class_coverage['train_N1000_s1.parquet'].median_per_present_class :: sha256 82652f57f708dc23
\newcommand{\FtCoverageVoneEThreeOneJet}{25}  % experiments/FT/data/jc2_v1_coverage.json :: class_coverage['train_N1000_s1.parquet'].n_classes_with_one_jet :: sha256 82652f57f708dc23
\newcommand{\ProbeNJets}{2{,}000{,}000}  % experiments/FIGS/data/probe_ladder_v2_mlp2/analysis/seed_level_results.json :: provenance.n_jets_total :: sha256 53e1ec81343ac2c1
\newcommand{\ProbeNSeeds}{5}  % experiments/FIGS/data/probe_ladder_v2_mlp2/analysis/seed_level_results.json :: seeds_used (length) :: sha256 53e1ec81343ac2c1
\newcommand{\ProbeNSeedsWord}{five}  % experiments/FIGS/data/probe_ladder_v2_mlp2/analysis/seed_level_results.json :: seeds_used (length) :: sha256 53e1ec81343ac2c1 :: pretraining runs per vocabulary, as a word for the prose
\newcommand{\ProbeRowAlign}{\texttt{a9449fd420c473f6}}  % experiments/FIGS/data/probe_ladder_v2_mlp2/analysis/seed_level_results.json :: provenance.row_alignment_sha256 (first 16) :: sha256 53e1ec81343ac2c1
\newcommand{\ProbeEpsS}{90\%}  % experiments/FIGS/data/probe_ladder_v2_mlp2/analysis/seed_level_results.json :: levels.bvc_4prong.linear[0].headline_eps_s :: sha256 53e1ec81343ac2c1
\newcommand{\ProbeAucBvcFourprongLinearOneeighteight}{0.99320}  % experiments/FIGS/data/probe_ladder_v2_mlp2/analysis/seed_level_results.json :: levels.bvc_4prong.linear[0].mean_auc :: sha256 53e1ec81343ac2c1 :: 5 seeds, 0 saturated at AUC=1
\newcommand{\ProbeAucSdBvcFourprongLinearOneeighteight}{0.00027}  % experiments/FIGS/data/probe_ladder_v2_mlp2/analysis/seed_level_results.json :: table[bvc_4prong,linear,188].auc (SD over seeds, ddof=1) :: sha256 53e1ec81343ac2c1
\newcommand{\ProbeLogOneMinusAucBvcFourprongLinearOneeighteight}{\ensuremath{-}4.9915}  % experiments/FIGS/data/probe_ladder_v2_mlp2/analysis/seed_level_results.json :: levels.bvc_4prong.linear[0].mean :: sha256 53e1ec81343ac2c1 :: natural log of 1-AUC, lower is better
\newcommand{\ProbeLogOneMinusAucSdBvcFourprongLinearOneeighteight}{0.0389}  % experiments/FIGS/data/probe_ladder_v2_mlp2/analysis/seed_level_results.json :: levels.bvc_4prong.linear[0].seed_sd :: sha256 53e1ec81343ac2c1
\newcommand{\ProbeOmaBvcFourprongLinearOneeighteight}{\ensuremath{(6.80 \pm 0.27)\times10^{-3}}}  % experiments/FIGS/data/probe_ladder_v2_mlp2/analysis/seed_level_results.json :: table[bvc_4prong,linear,188].auc :: sha256 53e1ec81343ac2c1 :: 1-AUC, mean +- SD over 5 seeds
\newcommand{\ProbeBkgLeftBvcFourprongLinearOneeighteight}{14.0}  % experiments/FIGS/data/probe_ladder_v2_mlp2/analysis/seed_level_results.json :: levels.bvc_4prong.linear[0].rejection_points['0.90'].mean_n_bkg_pass :: sha256 53e1ec81343ac2c1 :: background jets passing the cut, mean over seeds
\newcommand{\ProbeRejBvcFourprongLinearOneeighteight}{227\,\ensuremath{\pm}\,23}  % experiments/FIGS/data/probe_ladder_v2_mlp2/analysis/seed_level_results.json :: table[bvc_4prong,linear,188].rejection_points['0.90'].rejection :: sha256 53e1ec81343ac2c1 :: at 90% signal efficiency, mean +- SD over seeds; 0 of 5 seeds at the cap
\newcommand{\ProbeAucBvcFourprongLinearOnesixtwo}{0.99311}  % experiments/FIGS/data/probe_ladder_v2_mlp2/analysis/seed_level_results.json :: levels.bvc_4prong.linear[1].mean_auc :: sha256 53e1ec81343ac2c1 :: 5 seeds, 0 saturated at AUC=1
\newcommand{\ProbeAucSdBvcFourprongLinearOnesixtwo}{0.00018}  % experiments/FIGS/data/probe_ladder_v2_mlp2/analysis/seed_level_results.json :: table[bvc_4prong,linear,162].auc (SD over seeds, ddof=1) :: sha256 53e1ec81343ac2c1
\newcommand{\ProbeLogOneMinusAucBvcFourprongLinearOnesixtwo}{\ensuremath{-}4.9777}  % experiments/FIGS/data/probe_ladder_v2_mlp2/analysis/seed_level_results.json :: levels.bvc_4prong.linear[1].mean :: sha256 53e1ec81343ac2c1 :: natural log of 1-AUC, lower is better
\newcommand{\ProbeLogOneMinusAucSdBvcFourprongLinearOnesixtwo}{0.0264}  % experiments/FIGS/data/probe_ladder_v2_mlp2/analysis/seed_level_results.json :: levels.bvc_4prong.linear[1].seed_sd :: sha256 53e1ec81343ac2c1
\newcommand{\ProbeOmaBvcFourprongLinearOnesixtwo}{\ensuremath{(6.89 \pm 0.18)\times10^{-3}}}  % experiments/FIGS/data/probe_ladder_v2_mlp2/analysis/seed_level_results.json :: table[bvc_4prong,linear,162].auc :: sha256 53e1ec81343ac2c1 :: 1-AUC, mean +- SD over 5 seeds
\newcommand{\ProbeBkgLeftBvcFourprongLinearOnesixtwo}{14.4}  % experiments/FIGS/data/probe_ladder_v2_mlp2/analysis/seed_level_results.json :: levels.bvc_4prong.linear[1].rejection_points['0.90'].mean_n_bkg_pass :: sha256 53e1ec81343ac2c1 :: background jets passing the cut, mean over seeds
\newcommand{\ProbeRejBvcFourprongLinearOnesixtwo}{220\,\ensuremath{\pm}\,40}  % experiments/FIGS/data/probe_ladder_v2_mlp2/analysis/seed_level_results.json :: table[bvc_4prong,linear,162].rejection_points['0.90'].rejection :: sha256 53e1ec81343ac2c1 :: at 90% signal efficiency, mean +- SD over seeds; 0 of 5 seeds at the cap
\newcommand{\ProbeAucBvcFourprongLinearFourthree}{0.99250}  % experiments/FIGS/data/probe_ladder_v2_mlp2/analysis/seed_level_results.json :: levels.bvc_4prong.linear[2].mean_auc :: sha256 53e1ec81343ac2c1 :: 5 seeds, 0 saturated at AUC=1
\newcommand{\ProbeAucSdBvcFourprongLinearFourthree}{0.00049}  % experiments/FIGS/data/probe_ladder_v2_mlp2/analysis/seed_level_results.json :: table[bvc_4prong,linear,43].auc (SD over seeds, ddof=1) :: sha256 53e1ec81343ac2c1
\newcommand{\ProbeLogOneMinusAucBvcFourprongLinearFourthree}{\ensuremath{-}4.8946}  % experiments/FIGS/data/probe_ladder_v2_mlp2/analysis/seed_level_results.json :: levels.bvc_4prong.linear[2].mean :: sha256 53e1ec81343ac2c1 :: natural log of 1-AUC, lower is better
\newcommand{\ProbeLogOneMinusAucSdBvcFourprongLinearFourthree}{0.0646}  % experiments/FIGS/data/probe_ladder_v2_mlp2/analysis/seed_level_results.json :: levels.bvc_4prong.linear[2].seed_sd :: sha256 53e1ec81343ac2c1
\newcommand{\ProbeOmaBvcFourprongLinearFourthree}{\ensuremath{(7.5 \pm 0.5)\times10^{-3}}}  % experiments/FIGS/data/probe_ladder_v2_mlp2/analysis/seed_level_results.json :: table[bvc_4prong,linear,43].auc :: sha256 53e1ec81343ac2c1 :: 1-AUC, mean +- SD over 5 seeds
\newcommand{\ProbeBkgLeftBvcFourprongLinearFourthree}{16.6}  % experiments/FIGS/data/probe_ladder_v2_mlp2/analysis/seed_level_results.json :: levels.bvc_4prong.linear[2].rejection_points['0.90'].mean_n_bkg_pass :: sha256 53e1ec81343ac2c1 :: background jets passing the cut, mean over seeds
\newcommand{\ProbeRejBvcFourprongLinearFourthree}{192\,\ensuremath{\pm}\,20}  % experiments/FIGS/data/probe_ladder_v2_mlp2/analysis/seed_level_results.json :: table[bvc_4prong,linear,43].rejection_points['0.90'].rejection :: sha256 53e1ec81343ac2c1 :: at 90% signal efficiency, mean +- SD over seeds; 0 of 5 seeds at the cap
\newcommand{\ProbeAucBvcFourprongLinearOneseven}{0.96936}  % experiments/FIGS/data/probe_ladder_v2_mlp2/analysis/seed_level_results.json :: levels.bvc_4prong.linear[3].mean_auc :: sha256 53e1ec81343ac2c1 :: 5 seeds, 0 saturated at AUC=1
\newcommand{\ProbeAucSdBvcFourprongLinearOneseven}{0.00514}  % experiments/FIGS/data/probe_ladder_v2_mlp2/analysis/seed_level_results.json :: table[bvc_4prong,linear,17].auc (SD over seeds, ddof=1) :: sha256 53e1ec81343ac2c1
\newcommand{\ProbeLogOneMinusAucBvcFourprongLinearOneseven}{\ensuremath{-}3.4967}  % experiments/FIGS/data/probe_ladder_v2_mlp2/analysis/seed_level_results.json :: levels.bvc_4prong.linear[3].mean :: sha256 53e1ec81343ac2c1 :: natural log of 1-AUC, lower is better
\newcommand{\ProbeLogOneMinusAucSdBvcFourprongLinearOneseven}{0.1667}  % experiments/FIGS/data/probe_ladder_v2_mlp2/analysis/seed_level_results.json :: levels.bvc_4prong.linear[3].seed_sd :: sha256 53e1ec81343ac2c1
\newcommand{\ProbeOmaBvcFourprongLinearOneseven}{\ensuremath{(3.1 \pm 0.5)\times10^{-2}}}  % experiments/FIGS/data/probe_ladder_v2_mlp2/analysis/seed_level_results.json :: table[bvc_4prong,linear,17].auc :: sha256 53e1ec81343ac2c1 :: 1-AUC, mean +- SD over 5 seeds
\newcommand{\ProbeBkgLeftBvcFourprongLinearOneseven}{137.4}  % experiments/FIGS/data/probe_ladder_v2_mlp2/analysis/seed_level_results.json :: levels.bvc_4prong.linear[3].rejection_points['0.90'].mean_n_bkg_pass :: sha256 53e1ec81343ac2c1 :: background jets passing the cut, mean over seeds
\newcommand{\ProbeRejBvcFourprongLinearOneseven}{24\,\ensuremath{\pm}\,6}  % experiments/FIGS/data/probe_ladder_v2_mlp2/analysis/seed_level_results.json :: table[bvc_4prong,linear,17].rejection_points['0.90'].rejection :: sha256 53e1ec81343ac2c1 :: at 90% signal efficiency, mean +- SD over seeds; 0 of 5 seeds at the cap
\newcommand{\ProbeOmaRatioBvcFourprongLinearOnesixtwoOverOneeighteight}{1.01}  % experiments/FIGS/data/probe_ladder_v2_mlp2/analysis/seed_level_results.json :: table[bvc_4prong,linear,162/188].auc :: sha256 53e1ec81343ac2c1 :: ratio of the seed means of 1-AUC, coarser over finer
\newcommand{\ProbeOmaRatioBvcFourprongLinearFourthreeOverOneeighteight}{1.10}  % experiments/FIGS/data/probe_ladder_v2_mlp2/analysis/seed_level_results.json :: table[bvc_4prong,linear,43/188].auc :: sha256 53e1ec81343ac2c1 :: ratio of the seed means of 1-AUC, coarser over finer
\newcommand{\ProbeOmaRatioBvcFourprongLinearOnesevenOverOneeighteight}{4.5}  % experiments/FIGS/data/probe_ladder_v2_mlp2/analysis/seed_level_results.json :: table[bvc_4prong,linear,17/188].auc :: sha256 53e1ec81343ac2c1 :: ratio of the seed means of 1-AUC, coarser over finer
\newcommand{\ProbeOmaRatioBvcFourprongLinearFourthreeOverOnesixtwo}{1.09}  % experiments/FIGS/data/probe_ladder_v2_mlp2/analysis/seed_level_results.json :: table[bvc_4prong,linear,43/162].auc :: sha256 53e1ec81343ac2c1 :: ratio of the seed means of 1-AUC, coarser over finer
\newcommand{\ProbeOmaRatioBvcFourprongLinearOnesevenOverOnesixtwo}{4.4}  % experiments/FIGS/data/probe_ladder_v2_mlp2/analysis/seed_level_results.json :: table[bvc_4prong,linear,17/162].auc :: sha256 53e1ec81343ac2c1 :: ratio of the seed means of 1-AUC, coarser over finer
\newcommand{\ProbeOmaRatioBvcFourprongLinearOnesevenOverFourthree}{4.1}  % experiments/FIGS/data/probe_ladder_v2_mlp2/analysis/seed_level_results.json :: table[bvc_4prong,linear,17/43].auc :: sha256 53e1ec81343ac2c1 :: ratio of the seed means of 1-AUC, coarser over finer
\newcommand{\ProbeAucBvcFourprongMlpOneeighteight}{0.99316}  % experiments/FIGS/data/probe_ladder_v2_mlp2/analysis/seed_level_results.json :: levels.bvc_4prong.mlp[0].mean_auc :: sha256 53e1ec81343ac2c1 :: 5 seeds, 0 saturated at AUC=1
\newcommand{\ProbeAucSdBvcFourprongMlpOneeighteight}{0.00020}  % experiments/FIGS/data/probe_ladder_v2_mlp2/analysis/seed_level_results.json :: table[bvc_4prong,mlp,188].auc (SD over seeds, ddof=1) :: sha256 53e1ec81343ac2c1
\newcommand{\ProbeLogOneMinusAucBvcFourprongMlpOneeighteight}{\ensuremath{-}4.9857}  % experiments/FIGS/data/probe_ladder_v2_mlp2/analysis/seed_level_results.json :: levels.bvc_4prong.mlp[0].mean :: sha256 53e1ec81343ac2c1 :: natural log of 1-AUC, lower is better
\newcommand{\ProbeLogOneMinusAucSdBvcFourprongMlpOneeighteight}{0.0290}  % experiments/FIGS/data/probe_ladder_v2_mlp2/analysis/seed_level_results.json :: levels.bvc_4prong.mlp[0].seed_sd :: sha256 53e1ec81343ac2c1
\newcommand{\ProbeOmaBvcFourprongMlpOneeighteight}{\ensuremath{(6.84 \pm 0.20)\times10^{-3}}}  % experiments/FIGS/data/probe_ladder_v2_mlp2/analysis/seed_level_results.json :: table[bvc_4prong,mlp,188].auc :: sha256 53e1ec81343ac2c1 :: 1-AUC, mean +- SD over 5 seeds
\newcommand{\ProbeBkgLeftBvcFourprongMlpOneeighteight}{15.2}  % experiments/FIGS/data/probe_ladder_v2_mlp2/analysis/seed_level_results.json :: levels.bvc_4prong.mlp[0].rejection_points['0.90'].mean_n_bkg_pass :: sha256 53e1ec81343ac2c1 :: background jets passing the cut, mean over seeds
\newcommand{\ProbeRejBvcFourprongMlpOneeighteight}{212\,\ensuremath{\pm}\,34}  % experiments/FIGS/data/probe_ladder_v2_mlp2/analysis/seed_level_results.json :: table[bvc_4prong,mlp,188].rejection_points['0.90'].rejection :: sha256 53e1ec81343ac2c1 :: at 90% signal efficiency, mean +- SD over seeds; 0 of 5 seeds at the cap
\newcommand{\ProbeAucBvcFourprongMlpOnesixtwo}{0.99318}  % experiments/FIGS/data/probe_ladder_v2_mlp2/analysis/seed_level_results.json :: levels.bvc_4prong.mlp[1].mean_auc :: sha256 53e1ec81343ac2c1 :: 5 seeds, 0 saturated at AUC=1
\newcommand{\ProbeAucSdBvcFourprongMlpOnesixtwo}{0.00020}  % experiments/FIGS/data/probe_ladder_v2_mlp2/analysis/seed_level_results.json :: table[bvc_4prong,mlp,162].auc (SD over seeds, ddof=1) :: sha256 53e1ec81343ac2c1
\newcommand{\ProbeLogOneMinusAucBvcFourprongMlpOnesixtwo}{\ensuremath{-}4.9886}  % experiments/FIGS/data/probe_ladder_v2_mlp2/analysis/seed_level_results.json :: levels.bvc_4prong.mlp[1].mean :: sha256 53e1ec81343ac2c1 :: natural log of 1-AUC, lower is better
\newcommand{\ProbeLogOneMinusAucSdBvcFourprongMlpOnesixtwo}{0.0298}  % experiments/FIGS/data/probe_ladder_v2_mlp2/analysis/seed_level_results.json :: levels.bvc_4prong.mlp[1].seed_sd :: sha256 53e1ec81343ac2c1
\newcommand{\ProbeOmaBvcFourprongMlpOnesixtwo}{\ensuremath{(6.82 \pm 0.20)\times10^{-3}}}  % experiments/FIGS/data/probe_ladder_v2_mlp2/analysis/seed_level_results.json :: table[bvc_4prong,mlp,162].auc :: sha256 53e1ec81343ac2c1 :: 1-AUC, mean +- SD over 5 seeds
\newcommand{\ProbeBkgLeftBvcFourprongMlpOnesixtwo}{16.6}  % experiments/FIGS/data/probe_ladder_v2_mlp2/analysis/seed_level_results.json :: levels.bvc_4prong.mlp[1].rejection_points['0.90'].mean_n_bkg_pass :: sha256 53e1ec81343ac2c1 :: background jets passing the cut, mean over seeds
\newcommand{\ProbeRejBvcFourprongMlpOnesixtwo}{200\,\ensuremath{\pm}\,40}  % experiments/FIGS/data/probe_ladder_v2_mlp2/analysis/seed_level_results.json :: table[bvc_4prong,mlp,162].rejection_points['0.90'].rejection :: sha256 53e1ec81343ac2c1 :: at 90% signal efficiency, mean +- SD over seeds; 0 of 5 seeds at the cap
\newcommand{\ProbeAucBvcFourprongMlpFourthree}{0.99232}  % experiments/FIGS/data/probe_ladder_v2_mlp2/analysis/seed_level_results.json :: levels.bvc_4prong.mlp[2].mean_auc :: sha256 53e1ec81343ac2c1 :: 5 seeds, 0 saturated at AUC=1
\newcommand{\ProbeAucSdBvcFourprongMlpFourthree}{0.00058}  % experiments/FIGS/data/probe_ladder_v2_mlp2/analysis/seed_level_results.json :: table[bvc_4prong,mlp,43].auc (SD over seeds, ddof=1) :: sha256 53e1ec81343ac2c1
\newcommand{\ProbeLogOneMinusAucBvcFourprongMlpFourthree}{\ensuremath{-}4.8719}  % experiments/FIGS/data/probe_ladder_v2_mlp2/analysis/seed_level_results.json :: levels.bvc_4prong.mlp[2].mean :: sha256 53e1ec81343ac2c1 :: natural log of 1-AUC, lower is better
\newcommand{\ProbeLogOneMinusAucSdBvcFourprongMlpFourthree}{0.0734}  % experiments/FIGS/data/probe_ladder_v2_mlp2/analysis/seed_level_results.json :: levels.bvc_4prong.mlp[2].seed_sd :: sha256 53e1ec81343ac2c1
\newcommand{\ProbeOmaBvcFourprongMlpFourthree}{\ensuremath{(7.7 \pm 0.6)\times10^{-3}}}  % experiments/FIGS/data/probe_ladder_v2_mlp2/analysis/seed_level_results.json :: table[bvc_4prong,mlp,43].auc :: sha256 53e1ec81343ac2c1 :: 1-AUC, mean +- SD over 5 seeds
\newcommand{\ProbeBkgLeftBvcFourprongMlpFourthree}{17.4}  % experiments/FIGS/data/probe_ladder_v2_mlp2/analysis/seed_level_results.json :: levels.bvc_4prong.mlp[2].rejection_points['0.90'].mean_n_bkg_pass :: sha256 53e1ec81343ac2c1 :: background jets passing the cut, mean over seeds
\newcommand{\ProbeRejBvcFourprongMlpFourthree}{184\,\ensuremath{\pm}\,23}  % experiments/FIGS/data/probe_ladder_v2_mlp2/analysis/seed_level_results.json :: table[bvc_4prong,mlp,43].rejection_points['0.90'].rejection :: sha256 53e1ec81343ac2c1 :: at 90% signal efficiency, mean +- SD over seeds; 0 of 5 seeds at the cap
\newcommand{\ProbeAucBvcFourprongMlpOneseven}{0.97766}  % experiments/FIGS/data/probe_ladder_v2_mlp2/analysis/seed_level_results.json :: levels.bvc_4prong.mlp[3].mean_auc :: sha256 53e1ec81343ac2c1 :: 5 seeds, 0 saturated at AUC=1
\newcommand{\ProbeAucSdBvcFourprongMlpOneseven}{0.00310}  % experiments/FIGS/data/probe_ladder_v2_mlp2/analysis/seed_level_results.json :: table[bvc_4prong,mlp,17].auc (SD over seeds, ddof=1) :: sha256 53e1ec81343ac2c1
\newcommand{\ProbeLogOneMinusAucBvcFourprongMlpOneseven}{\ensuremath{-}3.8087}  % experiments/FIGS/data/probe_ladder_v2_mlp2/analysis/seed_level_results.json :: levels.bvc_4prong.mlp[3].mean :: sha256 53e1ec81343ac2c1 :: natural log of 1-AUC, lower is better
\newcommand{\ProbeLogOneMinusAucSdBvcFourprongMlpOneseven}{0.1366}  % experiments/FIGS/data/probe_ladder_v2_mlp2/analysis/seed_level_results.json :: levels.bvc_4prong.mlp[3].seed_sd :: sha256 53e1ec81343ac2c1
\newcommand{\ProbeOmaBvcFourprongMlpOneseven}{\ensuremath{(2.23 \pm 0.31)\times10^{-2}}}  % experiments/FIGS/data/probe_ladder_v2_mlp2/analysis/seed_level_results.json :: table[bvc_4prong,mlp,17].auc :: sha256 53e1ec81343ac2c1 :: 1-AUC, mean +- SD over 5 seeds
\newcommand{\ProbeBkgLeftBvcFourprongMlpOneseven}{99.0}  % experiments/FIGS/data/probe_ladder_v2_mlp2/analysis/seed_level_results.json :: levels.bvc_4prong.mlp[3].rejection_points['0.90'].mean_n_bkg_pass :: sha256 53e1ec81343ac2c1 :: background jets passing the cut, mean over seeds
\newcommand{\ProbeRejBvcFourprongMlpOneseven}{33\,\ensuremath{\pm}\,5}  % experiments/FIGS/data/probe_ladder_v2_mlp2/analysis/seed_level_results.json :: table[bvc_4prong,mlp,17].rejection_points['0.90'].rejection :: sha256 53e1ec81343ac2c1 :: at 90% signal efficiency, mean +- SD over seeds; 0 of 5 seeds at the cap
\newcommand{\ProbeOmaRatioBvcFourprongMlpOnesixtwoOverOneeighteight}{1.00}  % experiments/FIGS/data/probe_ladder_v2_mlp2/analysis/seed_level_results.json :: table[bvc_4prong,mlp,162/188].auc :: sha256 53e1ec81343ac2c1 :: ratio of the seed means of 1-AUC, coarser over finer
\newcommand{\ProbeOmaRatioBvcFourprongMlpFourthreeOverOneeighteight}{1.12}  % experiments/FIGS/data/probe_ladder_v2_mlp2/analysis/seed_level_results.json :: table[bvc_4prong,mlp,43/188].auc :: sha256 53e1ec81343ac2c1 :: ratio of the seed means of 1-AUC, coarser over finer
\newcommand{\ProbeOmaRatioBvcFourprongMlpOnesevenOverOneeighteight}{3.3}  % experiments/FIGS/data/probe_ladder_v2_mlp2/analysis/seed_level_results.json :: table[bvc_4prong,mlp,17/188].auc :: sha256 53e1ec81343ac2c1 :: ratio of the seed means of 1-AUC, coarser over finer
\newcommand{\ProbeOmaRatioBvcFourprongMlpFourthreeOverOnesixtwo}{1.13}  % experiments/FIGS/data/probe_ladder_v2_mlp2/analysis/seed_level_results.json :: table[bvc_4prong,mlp,43/162].auc :: sha256 53e1ec81343ac2c1 :: ratio of the seed means of 1-AUC, coarser over finer
\newcommand{\ProbeOmaRatioBvcFourprongMlpOnesevenOverOnesixtwo}{3.3}  % experiments/FIGS/data/probe_ladder_v2_mlp2/analysis/seed_level_results.json :: table[bvc_4prong,mlp,17/162].auc :: sha256 53e1ec81343ac2c1 :: ratio of the seed means of 1-AUC, coarser over finer
\newcommand{\ProbeOmaRatioBvcFourprongMlpOnesevenOverFourthree}{2.9}  % experiments/FIGS/data/probe_ladder_v2_mlp2/analysis/seed_level_results.json :: table[bvc_4prong,mlp,17/43].auc :: sha256 53e1ec81343ac2c1 :: ratio of the seed means of 1-AUC, coarser over finer
\newcommand{\ProbeAucBvcQcdLinearOneeighteight}{0.98519}  % experiments/FIGS/data/probe_ladder_v2_mlp2/analysis/seed_level_results.json :: levels.bvc_qcd.linear[0].mean_auc :: sha256 53e1ec81343ac2c1 :: 5 seeds, 0 saturated at AUC=1
\newcommand{\ProbeAucSdBvcQcdLinearOneeighteight}{0.00061}  % experiments/FIGS/data/probe_ladder_v2_mlp2/analysis/seed_level_results.json :: table[bvc_qcd,linear,188].auc (SD over seeds, ddof=1) :: sha256 53e1ec81343ac2c1
\newcommand{\ProbeLogOneMinusAucBvcQcdLinearOneeighteight}{\ensuremath{-}4.2133}  % experiments/FIGS/data/probe_ladder_v2_mlp2/analysis/seed_level_results.json :: levels.bvc_qcd.linear[0].mean :: sha256 53e1ec81343ac2c1 :: natural log of 1-AUC, lower is better
\newcommand{\ProbeLogOneMinusAucSdBvcQcdLinearOneeighteight}{0.0422}  % experiments/FIGS/data/probe_ladder_v2_mlp2/analysis/seed_level_results.json :: levels.bvc_qcd.linear[0].seed_sd :: sha256 53e1ec81343ac2c1
\newcommand{\ProbeOmaBvcQcdLinearOneeighteight}{\ensuremath{(1.48 \pm 0.06)\times10^{-2}}}  % experiments/FIGS/data/probe_ladder_v2_mlp2/analysis/seed_level_results.json :: table[bvc_qcd,linear,188].auc :: sha256 53e1ec81343ac2c1 :: 1-AUC, mean +- SD over 5 seeds
\newcommand{\ProbeBkgLeftBvcQcdLinearOneeighteight}{96.4}  % experiments/FIGS/data/probe_ladder_v2_mlp2/analysis/seed_level_results.json :: levels.bvc_qcd.linear[0].rejection_points['0.90'].mean_n_bkg_pass :: sha256 53e1ec81343ac2c1 :: background jets passing the cut, mean over seeds
\newcommand{\ProbeRejBvcQcdLinearOneeighteight}{56\,\ensuremath{\pm}\,5}  % experiments/FIGS/data/probe_ladder_v2_mlp2/analysis/seed_level_results.json :: table[bvc_qcd,linear,188].rejection_points['0.90'].rejection :: sha256 53e1ec81343ac2c1 :: at 90% signal efficiency, mean +- SD over seeds; 0 of 5 seeds at the cap
\newcommand{\ProbeAucBvcQcdLinearOnesixtwo}{0.98305}  % experiments/FIGS/data/probe_ladder_v2_mlp2/analysis/seed_level_results.json :: levels.bvc_qcd.linear[1].mean_auc :: sha256 53e1ec81343ac2c1 :: 5 seeds, 0 saturated at AUC=1
\newcommand{\ProbeAucSdBvcQcdLinearOnesixtwo}{0.00024}  % experiments/FIGS/data/probe_ladder_v2_mlp2/analysis/seed_level_results.json :: table[bvc_qcd,linear,162].auc (SD over seeds, ddof=1) :: sha256 53e1ec81343ac2c1
\newcommand{\ProbeLogOneMinusAucBvcQcdLinearOnesixtwo}{\ensuremath{-}4.0774}  % experiments/FIGS/data/probe_ladder_v2_mlp2/analysis/seed_level_results.json :: levels.bvc_qcd.linear[1].mean :: sha256 53e1ec81343ac2c1 :: natural log of 1-AUC, lower is better
\newcommand{\ProbeLogOneMinusAucSdBvcQcdLinearOnesixtwo}{0.0142}  % experiments/FIGS/data/probe_ladder_v2_mlp2/analysis/seed_level_results.json :: levels.bvc_qcd.linear[1].seed_sd :: sha256 53e1ec81343ac2c1
\newcommand{\ProbeOmaBvcQcdLinearOnesixtwo}{\ensuremath{(1.695 \pm 0.024)\times10^{-2}}}  % experiments/FIGS/data/probe_ladder_v2_mlp2/analysis/seed_level_results.json :: table[bvc_qcd,linear,162].auc :: sha256 53e1ec81343ac2c1 :: 1-AUC, mean +- SD over 5 seeds
\newcommand{\ProbeBkgLeftBvcQcdLinearOnesixtwo}{121.4}  % experiments/FIGS/data/probe_ladder_v2_mlp2/analysis/seed_level_results.json :: levels.bvc_qcd.linear[1].rejection_points['0.90'].mean_n_bkg_pass :: sha256 53e1ec81343ac2c1 :: background jets passing the cut, mean over seeds
\newcommand{\ProbeRejBvcQcdLinearOnesixtwo}{44.3\,\ensuremath{\pm}\,2.1}  % experiments/FIGS/data/probe_ladder_v2_mlp2/analysis/seed_level_results.json :: table[bvc_qcd,linear,162].rejection_points['0.90'].rejection :: sha256 53e1ec81343ac2c1 :: at 90% signal efficiency, mean +- SD over seeds; 0 of 5 seeds at the cap
\newcommand{\ProbeAucBvcQcdLinearFourthree}{0.98281}  % experiments/FIGS/data/probe_ladder_v2_mlp2/analysis/seed_level_results.json :: levels.bvc_qcd.linear[2].mean_auc :: sha256 53e1ec81343ac2c1 :: 5 seeds, 0 saturated at AUC=1
\newcommand{\ProbeAucSdBvcQcdLinearFourthree}{0.00064}  % experiments/FIGS/data/probe_ladder_v2_mlp2/analysis/seed_level_results.json :: table[bvc_qcd,linear,43].auc (SD over seeds, ddof=1) :: sha256 53e1ec81343ac2c1
\newcommand{\ProbeLogOneMinusAucBvcQcdLinearFourthree}{\ensuremath{-}4.0638}  % experiments/FIGS/data/probe_ladder_v2_mlp2/analysis/seed_level_results.json :: levels.bvc_qcd.linear[2].mean :: sha256 53e1ec81343ac2c1 :: natural log of 1-AUC, lower is better
\newcommand{\ProbeLogOneMinusAucSdBvcQcdLinearFourthree}{0.0370}  % experiments/FIGS/data/probe_ladder_v2_mlp2/analysis/seed_level_results.json :: levels.bvc_qcd.linear[2].seed_sd :: sha256 53e1ec81343ac2c1
\newcommand{\ProbeOmaBvcQcdLinearFourthree}{\ensuremath{(1.72 \pm 0.06)\times10^{-2}}}  % experiments/FIGS/data/probe_ladder_v2_mlp2/analysis/seed_level_results.json :: table[bvc_qcd,linear,43].auc :: sha256 53e1ec81343ac2c1 :: 1-AUC, mean +- SD over 5 seeds
\newcommand{\ProbeBkgLeftBvcQcdLinearFourthree}{126.8}  % experiments/FIGS/data/probe_ladder_v2_mlp2/analysis/seed_level_results.json :: levels.bvc_qcd.linear[2].rejection_points['0.90'].mean_n_bkg_pass :: sha256 53e1ec81343ac2c1 :: background jets passing the cut, mean over seeds
\newcommand{\ProbeRejBvcQcdLinearFourthree}{43\,\ensuremath{\pm}\,4}  % experiments/FIGS/data/probe_ladder_v2_mlp2/analysis/seed_level_results.json :: table[bvc_qcd,linear,43].rejection_points['0.90'].rejection :: sha256 53e1ec81343ac2c1 :: at 90% signal efficiency, mean +- SD over seeds; 0 of 5 seeds at the cap
\newcommand{\ProbeAucBvcQcdLinearOneseven}{0.95504}  % experiments/FIGS/data/probe_ladder_v2_mlp2/analysis/seed_level_results.json :: levels.bvc_qcd.linear[3].mean_auc :: sha256 53e1ec81343ac2c1 :: 5 seeds, 0 saturated at AUC=1
\newcommand{\ProbeAucSdBvcQcdLinearOneseven}{0.00527}  % experiments/FIGS/data/probe_ladder_v2_mlp2/analysis/seed_level_results.json :: table[bvc_qcd,linear,17].auc (SD over seeds, ddof=1) :: sha256 53e1ec81343ac2c1
\newcommand{\ProbeLogOneMinusAucBvcQcdLinearOneseven}{\ensuremath{-}3.1076}  % experiments/FIGS/data/probe_ladder_v2_mlp2/analysis/seed_level_results.json :: levels.bvc_qcd.linear[3].mean :: sha256 53e1ec81343ac2c1 :: natural log of 1-AUC, lower is better
\newcommand{\ProbeLogOneMinusAucSdBvcQcdLinearOneseven}{0.1198}  % experiments/FIGS/data/probe_ladder_v2_mlp2/analysis/seed_level_results.json :: levels.bvc_qcd.linear[3].seed_sd :: sha256 53e1ec81343ac2c1
\newcommand{\ProbeOmaBvcQcdLinearOneseven}{\ensuremath{(4.5 \pm 0.5)\times10^{-2}}}  % experiments/FIGS/data/probe_ladder_v2_mlp2/analysis/seed_level_results.json :: table[bvc_qcd,linear,17].auc :: sha256 53e1ec81343ac2c1 :: 1-AUC, mean +- SD over 5 seeds
\newcommand{\ProbeBkgLeftBvcQcdLinearOneseven}{544.6}  % experiments/FIGS/data/probe_ladder_v2_mlp2/analysis/seed_level_results.json :: levels.bvc_qcd.linear[3].rejection_points['0.90'].mean_n_bkg_pass :: sha256 53e1ec81343ac2c1 :: background jets passing the cut, mean over seeds
\newcommand{\ProbeRejBvcQcdLinearOneseven}{10.1\,\ensuremath{\pm}\,1.9}  % experiments/FIGS/data/probe_ladder_v2_mlp2/analysis/seed_level_results.json :: table[bvc_qcd,linear,17].rejection_points['0.90'].rejection :: sha256 53e1ec81343ac2c1 :: at 90% signal efficiency, mean +- SD over seeds; 0 of 5 seeds at the cap
\newcommand{\ProbeOmaRatioBvcQcdLinearOnesixtwoOverOneeighteight}{1.14}  % experiments/FIGS/data/probe_ladder_v2_mlp2/analysis/seed_level_results.json :: table[bvc_qcd,linear,162/188].auc :: sha256 53e1ec81343ac2c1 :: ratio of the seed means of 1-AUC, coarser over finer
\newcommand{\ProbeOmaRatioBvcQcdLinearFourthreeOverOneeighteight}{1.16}  % experiments/FIGS/data/probe_ladder_v2_mlp2/analysis/seed_level_results.json :: table[bvc_qcd,linear,43/188].auc :: sha256 53e1ec81343ac2c1 :: ratio of the seed means of 1-AUC, coarser over finer
\newcommand{\ProbeOmaRatioBvcQcdLinearOnesevenOverOneeighteight}{3.0}  % experiments/FIGS/data/probe_ladder_v2_mlp2/analysis/seed_level_results.json :: table[bvc_qcd,linear,17/188].auc :: sha256 53e1ec81343ac2c1 :: ratio of the seed means of 1-AUC, coarser over finer
\newcommand{\ProbeOmaRatioBvcQcdLinearFourthreeOverOnesixtwo}{1.01}  % experiments/FIGS/data/probe_ladder_v2_mlp2/analysis/seed_level_results.json :: table[bvc_qcd,linear,43/162].auc :: sha256 53e1ec81343ac2c1 :: ratio of the seed means of 1-AUC, coarser over finer
\newcommand{\ProbeOmaRatioBvcQcdLinearOnesevenOverOnesixtwo}{2.7}  % experiments/FIGS/data/probe_ladder_v2_mlp2/analysis/seed_level_results.json :: table[bvc_qcd,linear,17/162].auc :: sha256 53e1ec81343ac2c1 :: ratio of the seed means of 1-AUC, coarser over finer
\newcommand{\ProbeOmaRatioBvcQcdLinearOnesevenOverFourthree}{2.6}  % experiments/FIGS/data/probe_ladder_v2_mlp2/analysis/seed_level_results.json :: table[bvc_qcd,linear,17/43].auc :: sha256 53e1ec81343ac2c1 :: ratio of the seed means of 1-AUC, coarser over finer
\newcommand{\ProbeAucBvcQcdMlpOneeighteight}{0.98527}  % experiments/FIGS/data/probe_ladder_v2_mlp2/analysis/seed_level_results.json :: levels.bvc_qcd.mlp[0].mean_auc :: sha256 53e1ec81343ac2c1 :: 5 seeds, 0 saturated at AUC=1
\newcommand{\ProbeAucSdBvcQcdMlpOneeighteight}{0.00058}  % experiments/FIGS/data/probe_ladder_v2_mlp2/analysis/seed_level_results.json :: table[bvc_qcd,mlp,188].auc (SD over seeds, ddof=1) :: sha256 53e1ec81343ac2c1
\newcommand{\ProbeLogOneMinusAucBvcQcdMlpOneeighteight}{\ensuremath{-}4.2188}  % experiments/FIGS/data/probe_ladder_v2_mlp2/analysis/seed_level_results.json :: levels.bvc_qcd.mlp[0].mean :: sha256 53e1ec81343ac2c1 :: natural log of 1-AUC, lower is better
\newcommand{\ProbeLogOneMinusAucSdBvcQcdMlpOneeighteight}{0.0400}  % experiments/FIGS/data/probe_ladder_v2_mlp2/analysis/seed_level_results.json :: levels.bvc_qcd.mlp[0].seed_sd :: sha256 53e1ec81343ac2c1
\newcommand{\ProbeOmaBvcQcdMlpOneeighteight}{\ensuremath{(1.47 \pm 0.06)\times10^{-2}}}  % experiments/FIGS/data/probe_ladder_v2_mlp2/analysis/seed_level_results.json :: table[bvc_qcd,mlp,188].auc :: sha256 53e1ec81343ac2c1 :: 1-AUC, mean +- SD over 5 seeds
\newcommand{\ProbeBkgLeftBvcQcdMlpOneeighteight}{94.6}  % experiments/FIGS/data/probe_ladder_v2_mlp2/analysis/seed_level_results.json :: levels.bvc_qcd.mlp[0].rejection_points['0.90'].mean_n_bkg_pass :: sha256 53e1ec81343ac2c1 :: background jets passing the cut, mean over seeds
\newcommand{\ProbeRejBvcQcdMlpOneeighteight}{57\,\ensuremath{\pm}\,6}  % experiments/FIGS/data/probe_ladder_v2_mlp2/analysis/seed_level_results.json :: table[bvc_qcd,mlp,188].rejection_points['0.90'].rejection :: sha256 53e1ec81343ac2c1 :: at 90% signal efficiency, mean +- SD over seeds; 0 of 5 seeds at the cap
\newcommand{\ProbeAucBvcQcdMlpOnesixtwo}{0.98361}  % experiments/FIGS/data/probe_ladder_v2_mlp2/analysis/seed_level_results.json :: levels.bvc_qcd.mlp[1].mean_auc :: sha256 53e1ec81343ac2c1 :: 5 seeds, 0 saturated at AUC=1
\newcommand{\ProbeAucSdBvcQcdMlpOnesixtwo}{0.00039}  % experiments/FIGS/data/probe_ladder_v2_mlp2/analysis/seed_level_results.json :: table[bvc_qcd,mlp,162].auc (SD over seeds, ddof=1) :: sha256 53e1ec81343ac2c1
\newcommand{\ProbeLogOneMinusAucBvcQcdMlpOnesixtwo}{\ensuremath{-}4.1113}  % experiments/FIGS/data/probe_ladder_v2_mlp2/analysis/seed_level_results.json :: levels.bvc_qcd.mlp[1].mean :: sha256 53e1ec81343ac2c1 :: natural log of 1-AUC, lower is better
\newcommand{\ProbeLogOneMinusAucSdBvcQcdMlpOnesixtwo}{0.0241}  % experiments/FIGS/data/probe_ladder_v2_mlp2/analysis/seed_level_results.json :: levels.bvc_qcd.mlp[1].seed_sd :: sha256 53e1ec81343ac2c1
\newcommand{\ProbeOmaBvcQcdMlpOnesixtwo}{\ensuremath{(1.64 \pm 0.04)\times10^{-2}}}  % experiments/FIGS/data/probe_ladder_v2_mlp2/analysis/seed_level_results.json :: table[bvc_qcd,mlp,162].auc :: sha256 53e1ec81343ac2c1 :: 1-AUC, mean +- SD over 5 seeds
\newcommand{\ProbeBkgLeftBvcQcdMlpOnesixtwo}{119.6}  % experiments/FIGS/data/probe_ladder_v2_mlp2/analysis/seed_level_results.json :: levels.bvc_qcd.mlp[1].rejection_points['0.90'].mean_n_bkg_pass :: sha256 53e1ec81343ac2c1 :: background jets passing the cut, mean over seeds
\newcommand{\ProbeRejBvcQcdMlpOnesixtwo}{45.1\,\ensuremath{\pm}\,3.2}  % experiments/FIGS/data/probe_ladder_v2_mlp2/analysis/seed_level_results.json :: table[bvc_qcd,mlp,162].rejection_points['0.90'].rejection :: sha256 53e1ec81343ac2c1 :: at 90% signal efficiency, mean +- SD over seeds; 0 of 5 seeds at the cap
\newcommand{\ProbeAucBvcQcdMlpFourthree}{0.98341}  % experiments/FIGS/data/probe_ladder_v2_mlp2/analysis/seed_level_results.json :: levels.bvc_qcd.mlp[2].mean_auc :: sha256 53e1ec81343ac2c1 :: 5 seeds, 0 saturated at AUC=1
\newcommand{\ProbeAucSdBvcQcdMlpFourthree}{0.00052}  % experiments/FIGS/data/probe_ladder_v2_mlp2/analysis/seed_level_results.json :: table[bvc_qcd,mlp,43].auc (SD over seeds, ddof=1) :: sha256 53e1ec81343ac2c1
\newcommand{\ProbeLogOneMinusAucBvcQcdMlpFourthree}{\ensuremath{-}4.0995}  % experiments/FIGS/data/probe_ladder_v2_mlp2/analysis/seed_level_results.json :: levels.bvc_qcd.mlp[2].mean :: sha256 53e1ec81343ac2c1 :: natural log of 1-AUC, lower is better
\newcommand{\ProbeLogOneMinusAucSdBvcQcdMlpFourthree}{0.0318}  % experiments/FIGS/data/probe_ladder_v2_mlp2/analysis/seed_level_results.json :: levels.bvc_qcd.mlp[2].seed_sd :: sha256 53e1ec81343ac2c1
\newcommand{\ProbeOmaBvcQcdMlpFourthree}{\ensuremath{(1.66 \pm 0.05)\times10^{-2}}}  % experiments/FIGS/data/probe_ladder_v2_mlp2/analysis/seed_level_results.json :: table[bvc_qcd,mlp,43].auc :: sha256 53e1ec81343ac2c1 :: 1-AUC, mean +- SD over 5 seeds
\newcommand{\ProbeBkgLeftBvcQcdMlpFourthree}{119.2}  % experiments/FIGS/data/probe_ladder_v2_mlp2/analysis/seed_level_results.json :: levels.bvc_qcd.mlp[2].rejection_points['0.90'].mean_n_bkg_pass :: sha256 53e1ec81343ac2c1 :: background jets passing the cut, mean over seeds
\newcommand{\ProbeRejBvcQcdMlpFourthree}{45.1\,\ensuremath{\pm}\,1.3}  % experiments/FIGS/data/probe_ladder_v2_mlp2/analysis/seed_level_results.json :: table[bvc_qcd,mlp,43].rejection_points['0.90'].rejection :: sha256 53e1ec81343ac2c1 :: at 90% signal efficiency, mean +- SD over seeds; 0 of 5 seeds at the cap
\newcommand{\ProbeAucBvcQcdMlpOneseven}{0.96226}  % experiments/FIGS/data/probe_ladder_v2_mlp2/analysis/seed_level_results.json :: levels.bvc_qcd.mlp[3].mean_auc :: sha256 53e1ec81343ac2c1 :: 5 seeds, 0 saturated at AUC=1
\newcommand{\ProbeAucSdBvcQcdMlpOneseven}{0.00373}  % experiments/FIGS/data/probe_ladder_v2_mlp2/analysis/seed_level_results.json :: table[bvc_qcd,mlp,17].auc (SD over seeds, ddof=1) :: sha256 53e1ec81343ac2c1
\newcommand{\ProbeLogOneMinusAucBvcQcdMlpOneseven}{\ensuremath{-}3.2811}  % experiments/FIGS/data/probe_ladder_v2_mlp2/analysis/seed_level_results.json :: levels.bvc_qcd.mlp[3].mean :: sha256 53e1ec81343ac2c1 :: natural log of 1-AUC, lower is better
\newcommand{\ProbeLogOneMinusAucSdBvcQcdMlpOneseven}{0.1010}  % experiments/FIGS/data/probe_ladder_v2_mlp2/analysis/seed_level_results.json :: levels.bvc_qcd.mlp[3].seed_sd :: sha256 53e1ec81343ac2c1
\newcommand{\ProbeOmaBvcQcdMlpOneseven}{\ensuremath{(3.8 \pm 0.4)\times10^{-2}}}  % experiments/FIGS/data/probe_ladder_v2_mlp2/analysis/seed_level_results.json :: table[bvc_qcd,mlp,17].auc :: sha256 53e1ec81343ac2c1 :: 1-AUC, mean +- SD over 5 seeds
\newcommand{\ProbeBkgLeftBvcQcdMlpOneseven}{458.8}  % experiments/FIGS/data/probe_ladder_v2_mlp2/analysis/seed_level_results.json :: levels.bvc_qcd.mlp[3].rejection_points['0.90'].mean_n_bkg_pass :: sha256 53e1ec81343ac2c1 :: background jets passing the cut, mean over seeds
\newcommand{\ProbeRejBvcQcdMlpOneseven}{11.8\,\ensuremath{\pm}\,1.5}  % experiments/FIGS/data/probe_ladder_v2_mlp2/analysis/seed_level_results.json :: table[bvc_qcd,mlp,17].rejection_points['0.90'].rejection :: sha256 53e1ec81343ac2c1 :: at 90% signal efficiency, mean +- SD over seeds; 0 of 5 seeds at the cap
\newcommand{\ProbeOmaRatioBvcQcdMlpOnesixtwoOverOneeighteight}{1.11}  % experiments/FIGS/data/probe_ladder_v2_mlp2/analysis/seed_level_results.json :: table[bvc_qcd,mlp,162/188].auc :: sha256 53e1ec81343ac2c1 :: ratio of the seed means of 1-AUC, coarser over finer
\newcommand{\ProbeOmaRatioBvcQcdMlpFourthreeOverOneeighteight}{1.13}  % experiments/FIGS/data/probe_ladder_v2_mlp2/analysis/seed_level_results.json :: table[bvc_qcd,mlp,43/188].auc :: sha256 53e1ec81343ac2c1 :: ratio of the seed means of 1-AUC, coarser over finer
\newcommand{\ProbeOmaRatioBvcQcdMlpOnesevenOverOneeighteight}{2.6}  % experiments/FIGS/data/probe_ladder_v2_mlp2/analysis/seed_level_results.json :: table[bvc_qcd,mlp,17/188].auc :: sha256 53e1ec81343ac2c1 :: ratio of the seed means of 1-AUC, coarser over finer
\newcommand{\ProbeOmaRatioBvcQcdMlpFourthreeOverOnesixtwo}{1.01}  % experiments/FIGS/data/probe_ladder_v2_mlp2/analysis/seed_level_results.json :: table[bvc_qcd,mlp,43/162].auc :: sha256 53e1ec81343ac2c1 :: ratio of the seed means of 1-AUC, coarser over finer
\newcommand{\ProbeOmaRatioBvcQcdMlpOnesevenOverOnesixtwo}{2.3}  % experiments/FIGS/data/probe_ladder_v2_mlp2/analysis/seed_level_results.json :: table[bvc_qcd,mlp,17/162].auc :: sha256 53e1ec81343ac2c1 :: ratio of the seed means of 1-AUC, coarser over finer
\newcommand{\ProbeOmaRatioBvcQcdMlpOnesevenOverFourthree}{2.3}  % experiments/FIGS/data/probe_ladder_v2_mlp2/analysis/seed_level_results.json :: table[bvc_qcd,mlp,17/43].auc :: sha256 53e1ec81343ac2c1 :: ratio of the seed means of 1-AUC, coarser over finer
\newcommand{\ProbeAucBvcResonantLinearOneeighteight}{0.99768}  % experiments/FIGS/data/probe_ladder_v2_mlp2/analysis/seed_level_results.json :: levels.bvc_resonant.linear[0].mean_auc :: sha256 53e1ec81343ac2c1 :: 5 seeds, 0 saturated at AUC=1
\newcommand{\ProbeAucSdBvcResonantLinearOneeighteight}{0.00007}  % experiments/FIGS/data/probe_ladder_v2_mlp2/analysis/seed_level_results.json :: table[bvc_resonant,linear,188].auc (SD over seeds, ddof=1) :: sha256 53e1ec81343ac2c1
\newcommand{\ProbeLogOneMinusAucBvcResonantLinearOneeighteight}{\ensuremath{-}6.0683}  % experiments/FIGS/data/probe_ladder_v2_mlp2/analysis/seed_level_results.json :: levels.bvc_resonant.linear[0].mean :: sha256 53e1ec81343ac2c1 :: natural log of 1-AUC, lower is better
\newcommand{\ProbeLogOneMinusAucSdBvcResonantLinearOneeighteight}{0.0317}  % experiments/FIGS/data/probe_ladder_v2_mlp2/analysis/seed_level_results.json :: levels.bvc_resonant.linear[0].seed_sd :: sha256 53e1ec81343ac2c1
\newcommand{\ProbeOmaBvcResonantLinearOneeighteight}{\ensuremath{(2.32 \pm 0.07)\times10^{-3}}}  % experiments/FIGS/data/probe_ladder_v2_mlp2/analysis/seed_level_results.json :: table[bvc_resonant,linear,188].auc :: sha256 53e1ec81343ac2c1 :: 1-AUC, mean +- SD over 5 seeds
\newcommand{\ProbeBkgLeftBvcResonantLinearOneeighteight}{9.0}  % experiments/FIGS/data/probe_ladder_v2_mlp2/analysis/seed_level_results.json :: levels.bvc_resonant.linear[0].rejection_points['0.90'].mean_n_bkg_pass :: sha256 53e1ec81343ac2c1 :: background jets passing the cut, mean over seeds
\newcommand{\ProbeRejBvcResonantLinearOneeighteight}{1{,}380\,\ensuremath{\pm}\,350}  % experiments/FIGS/data/probe_ladder_v2_mlp2/analysis/seed_level_results.json :: table[bvc_resonant,linear,188].rejection_points['0.90'].rejection :: sha256 53e1ec81343ac2c1 :: at 90% signal efficiency, mean +- SD over seeds; 0 of 5 seeds at the cap
\newcommand{\ProbeAucBvcResonantLinearOnesixtwo}{0.99764}  % experiments/FIGS/data/probe_ladder_v2_mlp2/analysis/seed_level_results.json :: levels.bvc_resonant.linear[1].mean_auc :: sha256 53e1ec81343ac2c1 :: 5 seeds, 0 saturated at AUC=1
\newcommand{\ProbeAucSdBvcResonantLinearOnesixtwo}{0.00006}  % experiments/FIGS/data/probe_ladder_v2_mlp2/analysis/seed_level_results.json :: table[bvc_resonant,linear,162].auc (SD over seeds, ddof=1) :: sha256 53e1ec81343ac2c1
\newcommand{\ProbeLogOneMinusAucBvcResonantLinearOnesixtwo}{\ensuremath{-}6.0493}  % experiments/FIGS/data/probe_ladder_v2_mlp2/analysis/seed_level_results.json :: levels.bvc_resonant.linear[1].mean :: sha256 53e1ec81343ac2c1 :: natural log of 1-AUC, lower is better
\newcommand{\ProbeLogOneMinusAucSdBvcResonantLinearOnesixtwo}{0.0261}  % experiments/FIGS/data/probe_ladder_v2_mlp2/analysis/seed_level_results.json :: levels.bvc_resonant.linear[1].seed_sd :: sha256 53e1ec81343ac2c1
\newcommand{\ProbeOmaBvcResonantLinearOnesixtwo}{\ensuremath{(2.36 \pm 0.06)\times10^{-3}}}  % experiments/FIGS/data/probe_ladder_v2_mlp2/analysis/seed_level_results.json :: table[bvc_resonant,linear,162].auc :: sha256 53e1ec81343ac2c1 :: 1-AUC, mean +- SD over 5 seeds
\newcommand{\ProbeBkgLeftBvcResonantLinearOnesixtwo}{9.4}  % experiments/FIGS/data/probe_ladder_v2_mlp2/analysis/seed_level_results.json :: levels.bvc_resonant.linear[1].rejection_points['0.90'].mean_n_bkg_pass :: sha256 53e1ec81343ac2c1 :: background jets passing the cut, mean over seeds
\newcommand{\ProbeRejBvcResonantLinearOnesixtwo}{1{,}290\,\ensuremath{\pm}\,240}  % experiments/FIGS/data/probe_ladder_v2_mlp2/analysis/seed_level_results.json :: table[bvc_resonant,linear,162].rejection_points['0.90'].rejection :: sha256 53e1ec81343ac2c1 :: at 90% signal efficiency, mean +- SD over seeds; 0 of 5 seeds at the cap
\newcommand{\ProbeAucBvcResonantLinearFourthree}{0.99748}  % experiments/FIGS/data/probe_ladder_v2_mlp2/analysis/seed_level_results.json :: levels.bvc_resonant.linear[2].mean_auc :: sha256 53e1ec81343ac2c1 :: 5 seeds, 0 saturated at AUC=1
\newcommand{\ProbeAucSdBvcResonantLinearFourthree}{0.00007}  % experiments/FIGS/data/probe_ladder_v2_mlp2/analysis/seed_level_results.json :: table[bvc_resonant,linear,43].auc (SD over seeds, ddof=1) :: sha256 53e1ec81343ac2c1
\newcommand{\ProbeLogOneMinusAucBvcResonantLinearFourthree}{\ensuremath{-}5.9834}  % experiments/FIGS/data/probe_ladder_v2_mlp2/analysis/seed_level_results.json :: levels.bvc_resonant.linear[2].mean :: sha256 53e1ec81343ac2c1 :: natural log of 1-AUC, lower is better
\newcommand{\ProbeLogOneMinusAucSdBvcResonantLinearFourthree}{0.0264}  % experiments/FIGS/data/probe_ladder_v2_mlp2/analysis/seed_level_results.json :: levels.bvc_resonant.linear[2].seed_sd :: sha256 53e1ec81343ac2c1
\newcommand{\ProbeOmaBvcResonantLinearFourthree}{\ensuremath{(2.52 \pm 0.07)\times10^{-3}}}  % experiments/FIGS/data/probe_ladder_v2_mlp2/analysis/seed_level_results.json :: table[bvc_resonant,linear,43].auc :: sha256 53e1ec81343ac2c1 :: 1-AUC, mean +- SD over 5 seeds
\newcommand{\ProbeBkgLeftBvcResonantLinearFourthree}{11.0}  % experiments/FIGS/data/probe_ladder_v2_mlp2/analysis/seed_level_results.json :: levels.bvc_resonant.linear[2].rejection_points['0.90'].mean_n_bkg_pass :: sha256 53e1ec81343ac2c1 :: background jets passing the cut, mean over seeds
\newcommand{\ProbeRejBvcResonantLinearFourthree}{1{,}130\,\ensuremath{\pm}\,260}  % experiments/FIGS/data/probe_ladder_v2_mlp2/analysis/seed_level_results.json :: table[bvc_resonant,linear,43].rejection_points['0.90'].rejection :: sha256 53e1ec81343ac2c1 :: at 90% signal efficiency, mean +- SD over seeds; 0 of 5 seeds at the cap
\newcommand{\ProbeAucBvcResonantLinearOneseven}{0.99067}  % experiments/FIGS/data/probe_ladder_v2_mlp2/analysis/seed_level_results.json :: levels.bvc_resonant.linear[3].mean_auc :: sha256 53e1ec81343ac2c1 :: 5 seeds, 0 saturated at AUC=1
\newcommand{\ProbeAucSdBvcResonantLinearOneseven}{0.00147}  % experiments/FIGS/data/probe_ladder_v2_mlp2/analysis/seed_level_results.json :: table[bvc_resonant,linear,17].auc (SD over seeds, ddof=1) :: sha256 53e1ec81343ac2c1
\newcommand{\ProbeLogOneMinusAucBvcResonantLinearOneseven}{\ensuremath{-}4.6848}  % experiments/FIGS/data/probe_ladder_v2_mlp2/analysis/seed_level_results.json :: levels.bvc_resonant.linear[3].mean :: sha256 53e1ec81343ac2c1 :: natural log of 1-AUC, lower is better
\newcommand{\ProbeLogOneMinusAucSdBvcResonantLinearOneseven}{0.1619}  % experiments/FIGS/data/probe_ladder_v2_mlp2/analysis/seed_level_results.json :: levels.bvc_resonant.linear[3].seed_sd :: sha256 53e1ec81343ac2c1
\newcommand{\ProbeOmaBvcResonantLinearOneseven}{\ensuremath{(9.3 \pm 1.5)\times10^{-3}}}  % experiments/FIGS/data/probe_ladder_v2_mlp2/analysis/seed_level_results.json :: table[bvc_resonant,linear,17].auc :: sha256 53e1ec81343ac2c1 :: 1-AUC, mean +- SD over 5 seeds
\newcommand{\ProbeBkgLeftBvcResonantLinearOneseven}{90.8}  % experiments/FIGS/data/probe_ladder_v2_mlp2/analysis/seed_level_results.json :: levels.bvc_resonant.linear[3].rejection_points['0.90'].mean_n_bkg_pass :: sha256 53e1ec81343ac2c1 :: background jets passing the cut, mean over seeds
\newcommand{\ProbeRejBvcResonantLinearOneseven}{134\,\ensuremath{\pm}\,21}  % experiments/FIGS/data/probe_ladder_v2_mlp2/analysis/seed_level_results.json :: table[bvc_resonant,linear,17].rejection_points['0.90'].rejection :: sha256 53e1ec81343ac2c1 :: at 90% signal efficiency, mean +- SD over seeds; 0 of 5 seeds at the cap
\newcommand{\ProbeOmaRatioBvcResonantLinearOnesixtwoOverOneeighteight}{1.02}  % experiments/FIGS/data/probe_ladder_v2_mlp2/analysis/seed_level_results.json :: table[bvc_resonant,linear,162/188].auc :: sha256 53e1ec81343ac2c1 :: ratio of the seed means of 1-AUC, coarser over finer
\newcommand{\ProbeOmaRatioBvcResonantLinearFourthreeOverOneeighteight}{1.09}  % experiments/FIGS/data/probe_ladder_v2_mlp2/analysis/seed_level_results.json :: table[bvc_resonant,linear,43/188].auc :: sha256 53e1ec81343ac2c1 :: ratio of the seed means of 1-AUC, coarser over finer
\newcommand{\ProbeOmaRatioBvcResonantLinearOnesevenOverOneeighteight}{4.0}  % experiments/FIGS/data/probe_ladder_v2_mlp2/analysis/seed_level_results.json :: table[bvc_resonant,linear,17/188].auc :: sha256 53e1ec81343ac2c1 :: ratio of the seed means of 1-AUC, coarser over finer
\newcommand{\ProbeOmaRatioBvcResonantLinearFourthreeOverOnesixtwo}{1.07}  % experiments/FIGS/data/probe_ladder_v2_mlp2/analysis/seed_level_results.json :: table[bvc_resonant,linear,43/162].auc :: sha256 53e1ec81343ac2c1 :: ratio of the seed means of 1-AUC, coarser over finer
\newcommand{\ProbeOmaRatioBvcResonantLinearOnesevenOverOnesixtwo}{4.0}  % experiments/FIGS/data/probe_ladder_v2_mlp2/analysis/seed_level_results.json :: table[bvc_resonant,linear,17/162].auc :: sha256 53e1ec81343ac2c1 :: ratio of the seed means of 1-AUC, coarser over finer
\newcommand{\ProbeOmaRatioBvcResonantLinearOnesevenOverFourthree}{3.7}  % experiments/FIGS/data/probe_ladder_v2_mlp2/analysis/seed_level_results.json :: table[bvc_resonant,linear,17/43].auc :: sha256 53e1ec81343ac2c1 :: ratio of the seed means of 1-AUC, coarser over finer
\newcommand{\ProbeAucBvcResonantMlpOneeighteight}{0.99778}  % experiments/FIGS/data/probe_ladder_v2_mlp2/analysis/seed_level_results.json :: levels.bvc_resonant.mlp[0].mean_auc :: sha256 53e1ec81343ac2c1 :: 5 seeds, 0 saturated at AUC=1
\newcommand{\ProbeAucSdBvcResonantMlpOneeighteight}{0.00009}  % experiments/FIGS/data/probe_ladder_v2_mlp2/analysis/seed_level_results.json :: table[bvc_resonant,mlp,188].auc (SD over seeds, ddof=1) :: sha256 53e1ec81343ac2c1
\newcommand{\ProbeLogOneMinusAucBvcResonantMlpOneeighteight}{\ensuremath{-}6.1102}  % experiments/FIGS/data/probe_ladder_v2_mlp2/analysis/seed_level_results.json :: levels.bvc_resonant.mlp[0].mean :: sha256 53e1ec81343ac2c1 :: natural log of 1-AUC, lower is better
\newcommand{\ProbeLogOneMinusAucSdBvcResonantMlpOneeighteight}{0.0383}  % experiments/FIGS/data/probe_ladder_v2_mlp2/analysis/seed_level_results.json :: levels.bvc_resonant.mlp[0].seed_sd :: sha256 53e1ec81343ac2c1
\newcommand{\ProbeOmaBvcResonantMlpOneeighteight}{\ensuremath{(2.22 \pm 0.09)\times10^{-3}}}  % experiments/FIGS/data/probe_ladder_v2_mlp2/analysis/seed_level_results.json :: table[bvc_resonant,mlp,188].auc :: sha256 53e1ec81343ac2c1 :: 1-AUC, mean +- SD over 5 seeds
\newcommand{\ProbeBkgLeftBvcResonantMlpOneeighteight}{9.8}  % experiments/FIGS/data/probe_ladder_v2_mlp2/analysis/seed_level_results.json :: levels.bvc_resonant.mlp[0].rejection_points['0.90'].mean_n_bkg_pass :: sha256 53e1ec81343ac2c1 :: background jets passing the cut, mean over seeds
\newcommand{\ProbeRejBvcResonantMlpOneeighteight}{1{,}230\,\ensuremath{\pm}\,180}  % experiments/FIGS/data/probe_ladder_v2_mlp2/analysis/seed_level_results.json :: table[bvc_resonant,mlp,188].rejection_points['0.90'].rejection :: sha256 53e1ec81343ac2c1 :: at 90% signal efficiency, mean +- SD over seeds; 0 of 5 seeds at the cap
\newcommand{\ProbeAucBvcResonantMlpOnesixtwo}{0.99768}  % experiments/FIGS/data/probe_ladder_v2_mlp2/analysis/seed_level_results.json :: levels.bvc_resonant.mlp[1].mean_auc :: sha256 53e1ec81343ac2c1 :: 5 seeds, 0 saturated at AUC=1
\newcommand{\ProbeAucSdBvcResonantMlpOnesixtwo}{0.00005}  % experiments/FIGS/data/probe_ladder_v2_mlp2/analysis/seed_level_results.json :: table[bvc_resonant,mlp,162].auc (SD over seeds, ddof=1) :: sha256 53e1ec81343ac2c1
\newcommand{\ProbeLogOneMinusAucBvcResonantMlpOnesixtwo}{\ensuremath{-}6.0667}  % experiments/FIGS/data/probe_ladder_v2_mlp2/analysis/seed_level_results.json :: levels.bvc_resonant.mlp[1].mean :: sha256 53e1ec81343ac2c1 :: natural log of 1-AUC, lower is better
\newcommand{\ProbeLogOneMinusAucSdBvcResonantMlpOnesixtwo}{0.0220}  % experiments/FIGS/data/probe_ladder_v2_mlp2/analysis/seed_level_results.json :: levels.bvc_resonant.mlp[1].seed_sd :: sha256 53e1ec81343ac2c1
\newcommand{\ProbeOmaBvcResonantMlpOnesixtwo}{\ensuremath{(2.32 \pm 0.05)\times10^{-3}}}  % experiments/FIGS/data/probe_ladder_v2_mlp2/analysis/seed_level_results.json :: table[bvc_resonant,mlp,162].auc :: sha256 53e1ec81343ac2c1 :: 1-AUC, mean +- SD over 5 seeds
\newcommand{\ProbeBkgLeftBvcResonantMlpOnesixtwo}{10.4}  % experiments/FIGS/data/probe_ladder_v2_mlp2/analysis/seed_level_results.json :: levels.bvc_resonant.mlp[1].rejection_points['0.90'].mean_n_bkg_pass :: sha256 53e1ec81343ac2c1 :: background jets passing the cut, mean over seeds
\newcommand{\ProbeRejBvcResonantMlpOnesixtwo}{1{,}170\,\ensuremath{\pm}\,220}  % experiments/FIGS/data/probe_ladder_v2_mlp2/analysis/seed_level_results.json :: table[bvc_resonant,mlp,162].rejection_points['0.90'].rejection :: sha256 53e1ec81343ac2c1 :: at 90% signal efficiency, mean +- SD over seeds; 0 of 5 seeds at the cap
\newcommand{\ProbeAucBvcResonantMlpFourthree}{0.99747}  % experiments/FIGS/data/probe_ladder_v2_mlp2/analysis/seed_level_results.json :: levels.bvc_resonant.mlp[2].mean_auc :: sha256 53e1ec81343ac2c1 :: 5 seeds, 0 saturated at AUC=1
\newcommand{\ProbeAucSdBvcResonantMlpFourthree}{0.00004}  % experiments/FIGS/data/probe_ladder_v2_mlp2/analysis/seed_level_results.json :: table[bvc_resonant,mlp,43].auc (SD over seeds, ddof=1) :: sha256 53e1ec81343ac2c1
\newcommand{\ProbeLogOneMinusAucBvcResonantMlpFourthree}{\ensuremath{-}5.9792}  % experiments/FIGS/data/probe_ladder_v2_mlp2/analysis/seed_level_results.json :: levels.bvc_resonant.mlp[2].mean :: sha256 53e1ec81343ac2c1 :: natural log of 1-AUC, lower is better
\newcommand{\ProbeLogOneMinusAucSdBvcResonantMlpFourthree}{0.0164}  % experiments/FIGS/data/probe_ladder_v2_mlp2/analysis/seed_level_results.json :: levels.bvc_resonant.mlp[2].seed_sd :: sha256 53e1ec81343ac2c1
\newcommand{\ProbeOmaBvcResonantMlpFourthree}{\ensuremath{(2.53 \pm 0.04)\times10^{-3}}}  % experiments/FIGS/data/probe_ladder_v2_mlp2/analysis/seed_level_results.json :: table[bvc_resonant,mlp,43].auc :: sha256 53e1ec81343ac2c1 :: 1-AUC, mean +- SD over 5 seeds
\newcommand{\ProbeBkgLeftBvcResonantMlpFourthree}{11.8}  % experiments/FIGS/data/probe_ladder_v2_mlp2/analysis/seed_level_results.json :: levels.bvc_resonant.mlp[2].rejection_points['0.90'].mean_n_bkg_pass :: sha256 53e1ec81343ac2c1 :: background jets passing the cut, mean over seeds
\newcommand{\ProbeRejBvcResonantMlpFourthree}{1{,}030\,\ensuremath{\pm}\,170}  % experiments/FIGS/data/probe_ladder_v2_mlp2/analysis/seed_level_results.json :: table[bvc_resonant,mlp,43].rejection_points['0.90'].rejection :: sha256 53e1ec81343ac2c1 :: at 90% signal efficiency, mean +- SD over seeds; 0 of 5 seeds at the cap
\newcommand{\ProbeAucBvcResonantMlpOneseven}{0.99274}  % experiments/FIGS/data/probe_ladder_v2_mlp2/analysis/seed_level_results.json :: levels.bvc_resonant.mlp[3].mean_auc :: sha256 53e1ec81343ac2c1 :: 5 seeds, 0 saturated at AUC=1
\newcommand{\ProbeAucSdBvcResonantMlpOneseven}{0.00089}  % experiments/FIGS/data/probe_ladder_v2_mlp2/analysis/seed_level_results.json :: table[bvc_resonant,mlp,17].auc (SD over seeds, ddof=1) :: sha256 53e1ec81343ac2c1
\newcommand{\ProbeLogOneMinusAucBvcResonantMlpOneseven}{\ensuremath{-}4.9309}  % experiments/FIGS/data/probe_ladder_v2_mlp2/analysis/seed_level_results.json :: levels.bvc_resonant.mlp[3].mean :: sha256 53e1ec81343ac2c1 :: natural log of 1-AUC, lower is better
\newcommand{\ProbeLogOneMinusAucSdBvcResonantMlpOneseven}{0.1224}  % experiments/FIGS/data/probe_ladder_v2_mlp2/analysis/seed_level_results.json :: levels.bvc_resonant.mlp[3].seed_sd :: sha256 53e1ec81343ac2c1
\newcommand{\ProbeOmaBvcResonantMlpOneseven}{\ensuremath{(7.3 \pm 0.9)\times10^{-3}}}  % experiments/FIGS/data/probe_ladder_v2_mlp2/analysis/seed_level_results.json :: table[bvc_resonant,mlp,17].auc :: sha256 53e1ec81343ac2c1 :: 1-AUC, mean +- SD over 5 seeds
\newcommand{\ProbeBkgLeftBvcResonantMlpOneseven}{62.6}  % experiments/FIGS/data/probe_ladder_v2_mlp2/analysis/seed_level_results.json :: levels.bvc_resonant.mlp[3].rejection_points['0.90'].mean_n_bkg_pass :: sha256 53e1ec81343ac2c1 :: background jets passing the cut, mean over seeds
\newcommand{\ProbeRejBvcResonantMlpOneseven}{192\,\ensuremath{\pm}\,21}  % experiments/FIGS/data/probe_ladder_v2_mlp2/analysis/seed_level_results.json :: table[bvc_resonant,mlp,17].rejection_points['0.90'].rejection :: sha256 53e1ec81343ac2c1 :: at 90% signal efficiency, mean +- SD over seeds; 0 of 5 seeds at the cap
\newcommand{\ProbeOmaRatioBvcResonantMlpOnesixtwoOverOneeighteight}{1.04}  % experiments/FIGS/data/probe_ladder_v2_mlp2/analysis/seed_level_results.json :: table[bvc_resonant,mlp,162/188].auc :: sha256 53e1ec81343ac2c1 :: ratio of the seed means of 1-AUC, coarser over finer
\newcommand{\ProbeOmaRatioBvcResonantMlpFourthreeOverOneeighteight}{1.14}  % experiments/FIGS/data/probe_ladder_v2_mlp2/analysis/seed_level_results.json :: table[bvc_resonant,mlp,43/188].auc :: sha256 53e1ec81343ac2c1 :: ratio of the seed means of 1-AUC, coarser over finer
\newcommand{\ProbeOmaRatioBvcResonantMlpOnesevenOverOneeighteight}{3.3}  % experiments/FIGS/data/probe_ladder_v2_mlp2/analysis/seed_level_results.json :: table[bvc_resonant,mlp,17/188].auc :: sha256 53e1ec81343ac2c1 :: ratio of the seed means of 1-AUC, coarser over finer
\newcommand{\ProbeOmaRatioBvcResonantMlpFourthreeOverOnesixtwo}{1.09}  % experiments/FIGS/data/probe_ladder_v2_mlp2/analysis/seed_level_results.json :: table[bvc_resonant,mlp,43/162].auc :: sha256 53e1ec81343ac2c1 :: ratio of the seed means of 1-AUC, coarser over finer
\newcommand{\ProbeOmaRatioBvcResonantMlpOnesevenOverOnesixtwo}{3.1}  % experiments/FIGS/data/probe_ladder_v2_mlp2/analysis/seed_level_results.json :: table[bvc_resonant,mlp,17/162].auc :: sha256 53e1ec81343ac2c1 :: ratio of the seed means of 1-AUC, coarser over finer
\newcommand{\ProbeOmaRatioBvcResonantMlpOnesevenOverFourthree}{2.9}  % experiments/FIGS/data/probe_ladder_v2_mlp2/analysis/seed_level_results.json :: table[bvc_resonant,mlp,17/43].auc :: sha256 53e1ec81343ac2c1 :: ratio of the seed means of 1-AUC, coarser over finer
\newcommand{\ProbeAucEeVsMmLinearOneeighteight}{\ensuremath{1^{\dagger}}}  % experiments/FIGS/data/probe_ladder_v2_mlp2/analysis/seed_level_results.json :: levels.ee_vs_mm.linear[0].mean_auc :: sha256 53e1ec81343ac2c1 :: 5 seeds, 5 saturated at AUC=1
\newcommand{\ProbeAucSdEeVsMmLinearOneeighteight}{0.00000}  % experiments/FIGS/data/probe_ladder_v2_mlp2/analysis/seed_level_results.json :: table[ee_vs_mm,linear,188].auc (SD over seeds, ddof=1) :: sha256 53e1ec81343ac2c1
\newcommand{\ProbeLogOneMinusAucEeVsMmLinearOneeighteight}{\ensuremath{-}16.9579}  % experiments/FIGS/data/probe_ladder_v2_mlp2/analysis/seed_level_results.json :: levels.ee_vs_mm.linear[0].mean :: sha256 53e1ec81343ac2c1 :: natural log of 1-AUC, lower is better
\newcommand{\ProbeLogOneMinusAucSdEeVsMmLinearOneeighteight}{0.0000}  % experiments/FIGS/data/probe_ladder_v2_mlp2/analysis/seed_level_results.json :: levels.ee_vs_mm.linear[0].seed_sd :: sha256 53e1ec81343ac2c1
\newcommand{\ProbeOmaEeVsMmLinearOneeighteight}{---}  % experiments/FIGS/data/probe_ladder_v2_mlp2/analysis/seed_level_results.json :: table[ee_vs_mm,linear,188].auc :: sha256 53e1ec81343ac2c1 :: not quoted: the AUC reached 1 in a seed, so 1-AUC is only a bound
\newcommand{\ProbeBkgLeftEeVsMmLinearOneeighteight}{0.0}  % experiments/FIGS/data/probe_ladder_v2_mlp2/analysis/seed_level_results.json :: levels.ee_vs_mm.linear[0].rejection_points['0.90'].mean_n_bkg_pass :: sha256 53e1ec81343ac2c1 :: background jets passing the cut, mean over seeds
\newcommand{\ProbeRejEeVsMmLinearOneeighteight}{\ensuremath{>}1{,}571}  % experiments/FIGS/data/probe_ladder_v2_mlp2/analysis/seed_level_results.json :: table[ee_vs_mm,linear,188].rejection_points['0.90'].rejection :: sha256 53e1ec81343ac2c1 :: at 90% signal efficiency, mean +- SD over seeds; 5 of 5 seeds at the cap
\newcommand{\ProbeAucEeVsMmLinearOnesixtwo}{\ensuremath{1^{\dagger}}}  % experiments/FIGS/data/probe_ladder_v2_mlp2/analysis/seed_level_results.json :: levels.ee_vs_mm.linear[1].mean_auc :: sha256 53e1ec81343ac2c1 :: 5 seeds, 5 saturated at AUC=1
\newcommand{\ProbeAucSdEeVsMmLinearOnesixtwo}{0.00000}  % experiments/FIGS/data/probe_ladder_v2_mlp2/analysis/seed_level_results.json :: table[ee_vs_mm,linear,162].auc (SD over seeds, ddof=1) :: sha256 53e1ec81343ac2c1
\newcommand{\ProbeLogOneMinusAucEeVsMmLinearOnesixtwo}{\ensuremath{-}16.9579}  % experiments/FIGS/data/probe_ladder_v2_mlp2/analysis/seed_level_results.json :: levels.ee_vs_mm.linear[1].mean :: sha256 53e1ec81343ac2c1 :: natural log of 1-AUC, lower is better
\newcommand{\ProbeLogOneMinusAucSdEeVsMmLinearOnesixtwo}{0.0000}  % experiments/FIGS/data/probe_ladder_v2_mlp2/analysis/seed_level_results.json :: levels.ee_vs_mm.linear[1].seed_sd :: sha256 53e1ec81343ac2c1
\newcommand{\ProbeOmaEeVsMmLinearOnesixtwo}{---}  % experiments/FIGS/data/probe_ladder_v2_mlp2/analysis/seed_level_results.json :: table[ee_vs_mm,linear,162].auc :: sha256 53e1ec81343ac2c1 :: not quoted: the AUC reached 1 in a seed, so 1-AUC is only a bound
\newcommand{\ProbeBkgLeftEeVsMmLinearOnesixtwo}{0.0}  % experiments/FIGS/data/probe_ladder_v2_mlp2/analysis/seed_level_results.json :: levels.ee_vs_mm.linear[1].rejection_points['0.90'].mean_n_bkg_pass :: sha256 53e1ec81343ac2c1 :: background jets passing the cut, mean over seeds
\newcommand{\ProbeRejEeVsMmLinearOnesixtwo}{\ensuremath{>}1{,}571}  % experiments/FIGS/data/probe_ladder_v2_mlp2/analysis/seed_level_results.json :: table[ee_vs_mm,linear,162].rejection_points['0.90'].rejection :: sha256 53e1ec81343ac2c1 :: at 90% signal efficiency, mean +- SD over seeds; 5 of 5 seeds at the cap
\newcommand{\ProbeAucEeVsMmLinearFourthree}{0.97519}  % experiments/FIGS/data/probe_ladder_v2_mlp2/analysis/seed_level_results.json :: levels.ee_vs_mm.linear[2].mean_auc :: sha256 53e1ec81343ac2c1 :: 5 seeds, 0 saturated at AUC=1
\newcommand{\ProbeAucSdEeVsMmLinearFourthree}{0.02073}  % experiments/FIGS/data/probe_ladder_v2_mlp2/analysis/seed_level_results.json :: table[ee_vs_mm,linear,43].auc (SD over seeds, ddof=1) :: sha256 53e1ec81343ac2c1
\newcommand{\ProbeLogOneMinusAucEeVsMmLinearFourthree}{\ensuremath{-}4.0326}  % experiments/FIGS/data/probe_ladder_v2_mlp2/analysis/seed_level_results.json :: levels.ee_vs_mm.linear[2].mean :: sha256 53e1ec81343ac2c1 :: natural log of 1-AUC, lower is better
\newcommand{\ProbeLogOneMinusAucSdEeVsMmLinearFourthree}{0.9467}  % experiments/FIGS/data/probe_ladder_v2_mlp2/analysis/seed_level_results.json :: levels.ee_vs_mm.linear[2].seed_sd :: sha256 53e1ec81343ac2c1
\newcommand{\ProbeOmaEeVsMmLinearFourthree}{\ensuremath{(2.5 \pm 2.1)\times10^{-2}}}  % experiments/FIGS/data/probe_ladder_v2_mlp2/analysis/seed_level_results.json :: table[ee_vs_mm,linear,43].auc :: sha256 53e1ec81343ac2c1 :: 1-AUC, mean +- SD over 5 seeds
\newcommand{\ProbeBkgLeftEeVsMmLinearFourthree}{155.8}  % experiments/FIGS/data/probe_ladder_v2_mlp2/analysis/seed_level_results.json :: levels.ee_vs_mm.linear[2].rejection_points['0.90'].mean_n_bkg_pass :: sha256 53e1ec81343ac2c1 :: background jets passing the cut, mean over seeds
\newcommand{\ProbeRejEeVsMmLinearFourthree}{80\,\ensuremath{\pm}\,70}  % experiments/FIGS/data/probe_ladder_v2_mlp2/analysis/seed_level_results.json :: table[ee_vs_mm,linear,43].rejection_points['0.90'].rejection :: sha256 53e1ec81343ac2c1 :: at 90% signal efficiency, mean +- SD over seeds; 0 of 5 seeds at the cap
\newcommand{\ProbeAucEeVsMmLinearOneseven}{0.95832}  % experiments/FIGS/data/probe_ladder_v2_mlp2/analysis/seed_level_results.json :: levels.ee_vs_mm.linear[3].mean_auc :: sha256 53e1ec81343ac2c1 :: 5 seeds, 0 saturated at AUC=1
\newcommand{\ProbeAucSdEeVsMmLinearOneseven}{0.00970}  % experiments/FIGS/data/probe_ladder_v2_mlp2/analysis/seed_level_results.json :: table[ee_vs_mm,linear,17].auc (SD over seeds, ddof=1) :: sha256 53e1ec81343ac2c1
\newcommand{\ProbeLogOneMinusAucEeVsMmLinearOneseven}{\ensuremath{-}3.2013}  % experiments/FIGS/data/probe_ladder_v2_mlp2/analysis/seed_level_results.json :: levels.ee_vs_mm.linear[3].mean :: sha256 53e1ec81343ac2c1 :: natural log of 1-AUC, lower is better
\newcommand{\ProbeLogOneMinusAucSdEeVsMmLinearOneseven}{0.2488}  % experiments/FIGS/data/probe_ladder_v2_mlp2/analysis/seed_level_results.json :: levels.ee_vs_mm.linear[3].seed_sd :: sha256 53e1ec81343ac2c1
\newcommand{\ProbeOmaEeVsMmLinearOneseven}{\ensuremath{(4.2 \pm 1.0)\times10^{-2}}}  % experiments/FIGS/data/probe_ladder_v2_mlp2/analysis/seed_level_results.json :: table[ee_vs_mm,linear,17].auc :: sha256 53e1ec81343ac2c1 :: 1-AUC, mean +- SD over 5 seeds
\newcommand{\ProbeBkgLeftEeVsMmLinearOneseven}{289.6}  % experiments/FIGS/data/probe_ladder_v2_mlp2/analysis/seed_level_results.json :: levels.ee_vs_mm.linear[3].rejection_points['0.90'].mean_n_bkg_pass :: sha256 53e1ec81343ac2c1 :: background jets passing the cut, mean over seeds
\newcommand{\ProbeRejEeVsMmLinearOneseven}{17\,\ensuremath{\pm}\,6}  % experiments/FIGS/data/probe_ladder_v2_mlp2/analysis/seed_level_results.json :: table[ee_vs_mm,linear,17].rejection_points['0.90'].rejection :: sha256 53e1ec81343ac2c1 :: at 90% signal efficiency, mean +- SD over seeds; 0 of 5 seeds at the cap
\newcommand{\ProbeOmaRatioEeVsMmLinearOnesevenOverFourthree}{1.68}  % experiments/FIGS/data/probe_ladder_v2_mlp2/analysis/seed_level_results.json :: table[ee_vs_mm,linear,17/43].auc :: sha256 53e1ec81343ac2c1 :: ratio of the seed means of 1-AUC, coarser over finer
\newcommand{\ProbeAucEeVsMmMlpOneeighteight}{\ensuremath{1^{\dagger}}}  % experiments/FIGS/data/probe_ladder_v2_mlp2/analysis/seed_level_results.json :: levels.ee_vs_mm.mlp[0].mean_auc :: sha256 53e1ec81343ac2c1 :: 5 seeds, 5 saturated at AUC=1
\newcommand{\ProbeAucSdEeVsMmMlpOneeighteight}{0.00000}  % experiments/FIGS/data/probe_ladder_v2_mlp2/analysis/seed_level_results.json :: table[ee_vs_mm,mlp,188].auc (SD over seeds, ddof=1) :: sha256 53e1ec81343ac2c1
\newcommand{\ProbeLogOneMinusAucEeVsMmMlpOneeighteight}{\ensuremath{-}16.9579}  % experiments/FIGS/data/probe_ladder_v2_mlp2/analysis/seed_level_results.json :: levels.ee_vs_mm.mlp[0].mean :: sha256 53e1ec81343ac2c1 :: natural log of 1-AUC, lower is better
\newcommand{\ProbeLogOneMinusAucSdEeVsMmMlpOneeighteight}{0.0000}  % experiments/FIGS/data/probe_ladder_v2_mlp2/analysis/seed_level_results.json :: levels.ee_vs_mm.mlp[0].seed_sd :: sha256 53e1ec81343ac2c1
\newcommand{\ProbeOmaEeVsMmMlpOneeighteight}{---}  % experiments/FIGS/data/probe_ladder_v2_mlp2/analysis/seed_level_results.json :: table[ee_vs_mm,mlp,188].auc :: sha256 53e1ec81343ac2c1 :: not quoted: the AUC reached 1 in a seed, so 1-AUC is only a bound
\newcommand{\ProbeBkgLeftEeVsMmMlpOneeighteight}{0.0}  % experiments/FIGS/data/probe_ladder_v2_mlp2/analysis/seed_level_results.json :: levels.ee_vs_mm.mlp[0].rejection_points['0.90'].mean_n_bkg_pass :: sha256 53e1ec81343ac2c1 :: background jets passing the cut, mean over seeds
\newcommand{\ProbeRejEeVsMmMlpOneeighteight}{\ensuremath{>}1{,}571}  % experiments/FIGS/data/probe_ladder_v2_mlp2/analysis/seed_level_results.json :: table[ee_vs_mm,mlp,188].rejection_points['0.90'].rejection :: sha256 53e1ec81343ac2c1 :: at 90% signal efficiency, mean +- SD over seeds; 5 of 5 seeds at the cap
\newcommand{\ProbeAucEeVsMmMlpOnesixtwo}{\ensuremath{1^{\dagger}}}  % experiments/FIGS/data/probe_ladder_v2_mlp2/analysis/seed_level_results.json :: levels.ee_vs_mm.mlp[1].mean_auc :: sha256 53e1ec81343ac2c1 :: 5 seeds, 5 saturated at AUC=1
\newcommand{\ProbeAucSdEeVsMmMlpOnesixtwo}{0.00000}  % experiments/FIGS/data/probe_ladder_v2_mlp2/analysis/seed_level_results.json :: table[ee_vs_mm,mlp,162].auc (SD over seeds, ddof=1) :: sha256 53e1ec81343ac2c1
\newcommand{\ProbeLogOneMinusAucEeVsMmMlpOnesixtwo}{\ensuremath{-}16.9579}  % experiments/FIGS/data/probe_ladder_v2_mlp2/analysis/seed_level_results.json :: levels.ee_vs_mm.mlp[1].mean :: sha256 53e1ec81343ac2c1 :: natural log of 1-AUC, lower is better
\newcommand{\ProbeLogOneMinusAucSdEeVsMmMlpOnesixtwo}{0.0000}  % experiments/FIGS/data/probe_ladder_v2_mlp2/analysis/seed_level_results.json :: levels.ee_vs_mm.mlp[1].seed_sd :: sha256 53e1ec81343ac2c1
\newcommand{\ProbeOmaEeVsMmMlpOnesixtwo}{---}  % experiments/FIGS/data/probe_ladder_v2_mlp2/analysis/seed_level_results.json :: table[ee_vs_mm,mlp,162].auc :: sha256 53e1ec81343ac2c1 :: not quoted: the AUC reached 1 in a seed, so 1-AUC is only a bound
\newcommand{\ProbeBkgLeftEeVsMmMlpOnesixtwo}{0.0}  % experiments/FIGS/data/probe_ladder_v2_mlp2/analysis/seed_level_results.json :: levels.ee_vs_mm.mlp[1].rejection_points['0.90'].mean_n_bkg_pass :: sha256 53e1ec81343ac2c1 :: background jets passing the cut, mean over seeds
\newcommand{\ProbeRejEeVsMmMlpOnesixtwo}{\ensuremath{>}1{,}571}  % experiments/FIGS/data/probe_ladder_v2_mlp2/analysis/seed_level_results.json :: table[ee_vs_mm,mlp,162].rejection_points['0.90'].rejection :: sha256 53e1ec81343ac2c1 :: at 90% signal efficiency, mean +- SD over seeds; 5 of 5 seeds at the cap
\newcommand{\ProbeAucEeVsMmMlpFourthree}{0.99657}  % experiments/FIGS/data/probe_ladder_v2_mlp2/analysis/seed_level_results.json :: levels.ee_vs_mm.mlp[2].mean_auc :: sha256 53e1ec81343ac2c1 :: 5 seeds, 0 saturated at AUC=1
\newcommand{\ProbeAucSdEeVsMmMlpFourthree}{0.00300}  % experiments/FIGS/data/probe_ladder_v2_mlp2/analysis/seed_level_results.json :: table[ee_vs_mm,mlp,43].auc (SD over seeds, ddof=1) :: sha256 53e1ec81343ac2c1
\newcommand{\ProbeLogOneMinusAucEeVsMmMlpFourthree}{\ensuremath{-}6.0371}  % experiments/FIGS/data/probe_ladder_v2_mlp2/analysis/seed_level_results.json :: levels.ee_vs_mm.mlp[2].mean :: sha256 53e1ec81343ac2c1 :: natural log of 1-AUC, lower is better
\newcommand{\ProbeLogOneMinusAucSdEeVsMmMlpFourthree}{0.9962}  % experiments/FIGS/data/probe_ladder_v2_mlp2/analysis/seed_level_results.json :: levels.ee_vs_mm.mlp[2].seed_sd :: sha256 53e1ec81343ac2c1
\newcommand{\ProbeOmaEeVsMmMlpFourthree}{\ensuremath{(3.4 \pm 3.0)\times10^{-3}}}  % experiments/FIGS/data/probe_ladder_v2_mlp2/analysis/seed_level_results.json :: table[ee_vs_mm,mlp,43].auc :: sha256 53e1ec81343ac2c1 :: 1-AUC, mean +- SD over 5 seeds
\newcommand{\ProbeBkgLeftEeVsMmMlpFourthree}{33.4}  % experiments/FIGS/data/probe_ladder_v2_mlp2/analysis/seed_level_results.json :: levels.ee_vs_mm.mlp[2].rejection_points['0.90'].mean_n_bkg_pass :: sha256 53e1ec81343ac2c1 :: background jets passing the cut, mean over seeds
\newcommand{\ProbeRejEeVsMmMlpFourthree}{400\,\ensuremath{\pm}\,500}  % experiments/FIGS/data/probe_ladder_v2_mlp2/analysis/seed_level_results.json :: table[ee_vs_mm,mlp,43].rejection_points['0.90'].rejection :: sha256 53e1ec81343ac2c1 :: at 90% signal efficiency, mean +- SD over seeds; 0 of 5 seeds at the cap
\newcommand{\ProbeAucEeVsMmMlpOneseven}{0.99338}  % experiments/FIGS/data/probe_ladder_v2_mlp2/analysis/seed_level_results.json :: levels.ee_vs_mm.mlp[3].mean_auc :: sha256 53e1ec81343ac2c1 :: 5 seeds, 0 saturated at AUC=1
\newcommand{\ProbeAucSdEeVsMmMlpOneseven}{0.00065}  % experiments/FIGS/data/probe_ladder_v2_mlp2/analysis/seed_level_results.json :: table[ee_vs_mm,mlp,17].auc (SD over seeds, ddof=1) :: sha256 53e1ec81343ac2c1
\newcommand{\ProbeLogOneMinusAucEeVsMmMlpOneseven}{\ensuremath{-}5.0211}  % experiments/FIGS/data/probe_ladder_v2_mlp2/analysis/seed_level_results.json :: levels.ee_vs_mm.mlp[3].mean :: sha256 53e1ec81343ac2c1 :: natural log of 1-AUC, lower is better
\newcommand{\ProbeLogOneMinusAucSdEeVsMmMlpOneseven}{0.0997}  % experiments/FIGS/data/probe_ladder_v2_mlp2/analysis/seed_level_results.json :: levels.ee_vs_mm.mlp[3].seed_sd :: sha256 53e1ec81343ac2c1
\newcommand{\ProbeOmaEeVsMmMlpOneseven}{\ensuremath{(6.6 \pm 0.7)\times10^{-3}}}  % experiments/FIGS/data/probe_ladder_v2_mlp2/analysis/seed_level_results.json :: table[ee_vs_mm,mlp,17].auc :: sha256 53e1ec81343ac2c1 :: 1-AUC, mean +- SD over 5 seeds
\newcommand{\ProbeBkgLeftEeVsMmMlpOneseven}{59.2}  % experiments/FIGS/data/probe_ladder_v2_mlp2/analysis/seed_level_results.json :: levels.ee_vs_mm.mlp[3].rejection_points['0.90'].mean_n_bkg_pass :: sha256 53e1ec81343ac2c1 :: background jets passing the cut, mean over seeds
\newcommand{\ProbeRejEeVsMmMlpOneseven}{81\,\ensuremath{\pm}\,12}  % experiments/FIGS/data/probe_ladder_v2_mlp2/analysis/seed_level_results.json :: table[ee_vs_mm,mlp,17].rejection_points['0.90'].rejection :: sha256 53e1ec81343ac2c1 :: at 90% signal efficiency, mean +- SD over seeds; 0 of 5 seeds at the cap
\newcommand{\ProbeOmaRatioEeVsMmMlpOnesevenOverFourthree}{1.93}  % experiments/FIGS/data/probe_ladder_v2_mlp2/analysis/seed_level_results.json :: table[ee_vs_mm,mlp,17/43].auc :: sha256 53e1ec81343ac2c1 :: ratio of the seed means of 1-AUC, coarser over finer
\newcommand{\ProbeAucRetainedTopologyLinearOneeighteight}{0.99413}  % experiments/FIGS/data/probe_ladder_v2_mlp2/analysis/seed_level_results.json :: levels.retained_topology.linear[0].mean_auc :: sha256 53e1ec81343ac2c1 :: 5 seeds, 0 saturated at AUC=1
\newcommand{\ProbeAucSdRetainedTopologyLinearOneeighteight}{0.00011}  % experiments/FIGS/data/probe_ladder_v2_mlp2/analysis/seed_level_results.json :: table[retained_topology,linear,188].auc (SD over seeds, ddof=1) :: sha256 53e1ec81343ac2c1
\newcommand{\ProbeLogOneMinusAucRetainedTopologyLinearOneeighteight}{\ensuremath{-}5.1388}  % experiments/FIGS/data/probe_ladder_v2_mlp2/analysis/seed_level_results.json :: levels.retained_topology.linear[0].mean :: sha256 53e1ec81343ac2c1 :: natural log of 1-AUC, lower is better
\newcommand{\ProbeLogOneMinusAucSdRetainedTopologyLinearOneeighteight}{0.0187}  % experiments/FIGS/data/probe_ladder_v2_mlp2/analysis/seed_level_results.json :: levels.retained_topology.linear[0].seed_sd :: sha256 53e1ec81343ac2c1
\newcommand{\ProbeOmaRetainedTopologyLinearOneeighteight}{\ensuremath{(5.87 \pm 0.11)\times10^{-3}}}  % experiments/FIGS/data/probe_ladder_v2_mlp2/analysis/seed_level_results.json :: table[retained_topology,linear,188].auc :: sha256 53e1ec81343ac2c1 :: 1-AUC, mean +- SD over 5 seeds
\newcommand{\ProbeBkgLeftRetainedTopologyLinearOneeighteight}{15.4}  % experiments/FIGS/data/probe_ladder_v2_mlp2/analysis/seed_level_results.json :: levels.retained_topology.linear[0].rejection_points['0.90'].mean_n_bkg_pass :: sha256 53e1ec81343ac2c1 :: background jets passing the cut, mean over seeds
\newcommand{\ProbeRejRetainedTopologyLinearOneeighteight}{71\,\ensuremath{\pm}\,7}  % experiments/FIGS/data/probe_ladder_v2_mlp2/analysis/seed_level_results.json :: table[retained_topology,linear,188].rejection_points['0.90'].rejection :: sha256 53e1ec81343ac2c1 :: at 90% signal efficiency, mean +- SD over seeds; 0 of 5 seeds at the cap
\newcommand{\ProbeAucRetainedTopologyLinearOnesixtwo}{0.99413}  % experiments/FIGS/data/probe_ladder_v2_mlp2/analysis/seed_level_results.json :: levels.retained_topology.linear[1].mean_auc :: sha256 53e1ec81343ac2c1 :: 5 seeds, 0 saturated at AUC=1
\newcommand{\ProbeAucSdRetainedTopologyLinearOnesixtwo}{0.00030}  % experiments/FIGS/data/probe_ladder_v2_mlp2/analysis/seed_level_results.json :: table[retained_topology,linear,162].auc (SD over seeds, ddof=1) :: sha256 53e1ec81343ac2c1
\newcommand{\ProbeLogOneMinusAucRetainedTopologyLinearOnesixtwo}{\ensuremath{-}5.1395}  % experiments/FIGS/data/probe_ladder_v2_mlp2/analysis/seed_level_results.json :: levels.retained_topology.linear[1].mean :: sha256 53e1ec81343ac2c1 :: natural log of 1-AUC, lower is better
\newcommand{\ProbeLogOneMinusAucSdRetainedTopologyLinearOnesixtwo}{0.0510}  % experiments/FIGS/data/probe_ladder_v2_mlp2/analysis/seed_level_results.json :: levels.retained_topology.linear[1].seed_sd :: sha256 53e1ec81343ac2c1
\newcommand{\ProbeOmaRetainedTopologyLinearOnesixtwo}{\ensuremath{(5.87 \pm 0.30)\times10^{-3}}}  % experiments/FIGS/data/probe_ladder_v2_mlp2/analysis/seed_level_results.json :: table[retained_topology,linear,162].auc :: sha256 53e1ec81343ac2c1 :: 1-AUC, mean +- SD over 5 seeds
\newcommand{\ProbeBkgLeftRetainedTopologyLinearOnesixtwo}{14.0}  % experiments/FIGS/data/probe_ladder_v2_mlp2/analysis/seed_level_results.json :: levels.retained_topology.linear[1].rejection_points['0.90'].mean_n_bkg_pass :: sha256 53e1ec81343ac2c1 :: background jets passing the cut, mean over seeds
\newcommand{\ProbeRejRetainedTopologyLinearOnesixtwo}{79\,\ensuremath{\pm}\,13}  % experiments/FIGS/data/probe_ladder_v2_mlp2/analysis/seed_level_results.json :: table[retained_topology,linear,162].rejection_points['0.90'].rejection :: sha256 53e1ec81343ac2c1 :: at 90% signal efficiency, mean +- SD over seeds; 0 of 5 seeds at the cap
\newcommand{\ProbeAucRetainedTopologyLinearFourthree}{0.99274}  % experiments/FIGS/data/probe_ladder_v2_mlp2/analysis/seed_level_results.json :: levels.retained_topology.linear[2].mean_auc :: sha256 53e1ec81343ac2c1 :: 5 seeds, 0 saturated at AUC=1
\newcommand{\ProbeAucSdRetainedTopologyLinearFourthree}{0.00012}  % experiments/FIGS/data/probe_ladder_v2_mlp2/analysis/seed_level_results.json :: table[retained_topology,linear,43].auc (SD over seeds, ddof=1) :: sha256 53e1ec81343ac2c1
\newcommand{\ProbeLogOneMinusAucRetainedTopologyLinearFourthree}{\ensuremath{-}4.9253}  % experiments/FIGS/data/probe_ladder_v2_mlp2/analysis/seed_level_results.json :: levels.retained_topology.linear[2].mean :: sha256 53e1ec81343ac2c1 :: natural log of 1-AUC, lower is better
\newcommand{\ProbeLogOneMinusAucSdRetainedTopologyLinearFourthree}{0.0169}  % experiments/FIGS/data/probe_ladder_v2_mlp2/analysis/seed_level_results.json :: levels.retained_topology.linear[2].seed_sd :: sha256 53e1ec81343ac2c1
\newcommand{\ProbeOmaRetainedTopologyLinearFourthree}{\ensuremath{(7.26 \pm 0.12)\times10^{-3}}}  % experiments/FIGS/data/probe_ladder_v2_mlp2/analysis/seed_level_results.json :: table[retained_topology,linear,43].auc :: sha256 53e1ec81343ac2c1 :: 1-AUC, mean +- SD over 5 seeds
\newcommand{\ProbeBkgLeftRetainedTopologyLinearFourthree}{16.4}  % experiments/FIGS/data/probe_ladder_v2_mlp2/analysis/seed_level_results.json :: levels.retained_topology.linear[2].rejection_points['0.90'].mean_n_bkg_pass :: sha256 53e1ec81343ac2c1 :: background jets passing the cut, mean over seeds
\newcommand{\ProbeRejRetainedTopologyLinearFourthree}{66\,\ensuremath{\pm}\,5}  % experiments/FIGS/data/probe_ladder_v2_mlp2/analysis/seed_level_results.json :: table[retained_topology,linear,43].rejection_points['0.90'].rejection :: sha256 53e1ec81343ac2c1 :: at 90% signal efficiency, mean +- SD over seeds; 0 of 5 seeds at the cap
\newcommand{\ProbeAucRetainedTopologyLinearOneseven}{0.99096}  % experiments/FIGS/data/probe_ladder_v2_mlp2/analysis/seed_level_results.json :: levels.retained_topology.linear[3].mean_auc :: sha256 53e1ec81343ac2c1 :: 5 seeds, 0 saturated at AUC=1
\newcommand{\ProbeAucSdRetainedTopologyLinearOneseven}{0.00092}  % experiments/FIGS/data/probe_ladder_v2_mlp2/analysis/seed_level_results.json :: table[retained_topology,linear,17].auc (SD over seeds, ddof=1) :: sha256 53e1ec81343ac2c1
\newcommand{\ProbeLogOneMinusAucRetainedTopologyLinearOneseven}{\ensuremath{-}4.7095}  % experiments/FIGS/data/probe_ladder_v2_mlp2/analysis/seed_level_results.json :: levels.retained_topology.linear[3].mean :: sha256 53e1ec81343ac2c1 :: natural log of 1-AUC, lower is better
\newcommand{\ProbeLogOneMinusAucSdRetainedTopologyLinearOneseven}{0.0979}  % experiments/FIGS/data/probe_ladder_v2_mlp2/analysis/seed_level_results.json :: levels.retained_topology.linear[3].seed_sd :: sha256 53e1ec81343ac2c1
\newcommand{\ProbeOmaRetainedTopologyLinearOneseven}{\ensuremath{(9.0 \pm 0.9)\times10^{-3}}}  % experiments/FIGS/data/probe_ladder_v2_mlp2/analysis/seed_level_results.json :: table[retained_topology,linear,17].auc :: sha256 53e1ec81343ac2c1 :: 1-AUC, mean +- SD over 5 seeds
\newcommand{\ProbeBkgLeftRetainedTopologyLinearOneseven}{21.8}  % experiments/FIGS/data/probe_ladder_v2_mlp2/analysis/seed_level_results.json :: levels.retained_topology.linear[3].rejection_points['0.90'].mean_n_bkg_pass :: sha256 53e1ec81343ac2c1 :: background jets passing the cut, mean over seeds
\newcommand{\ProbeRejRetainedTopologyLinearOneseven}{51\,\ensuremath{\pm}\,11}  % experiments/FIGS/data/probe_ladder_v2_mlp2/analysis/seed_level_results.json :: table[retained_topology,linear,17].rejection_points['0.90'].rejection :: sha256 53e1ec81343ac2c1 :: at 90% signal efficiency, mean +- SD over seeds; 0 of 5 seeds at the cap
\newcommand{\ProbeOmaRatioRetainedTopologyLinearOnesixtwoOverOneeighteight}{1.00}  % experiments/FIGS/data/probe_ladder_v2_mlp2/analysis/seed_level_results.json :: table[retained_topology,linear,162/188].auc :: sha256 53e1ec81343ac2c1 :: ratio of the seed means of 1-AUC, coarser over finer
\newcommand{\ProbeOmaRatioRetainedTopologyLinearFourthreeOverOneeighteight}{1.24}  % experiments/FIGS/data/probe_ladder_v2_mlp2/analysis/seed_level_results.json :: table[retained_topology,linear,43/188].auc :: sha256 53e1ec81343ac2c1 :: ratio of the seed means of 1-AUC, coarser over finer
\newcommand{\ProbeOmaRatioRetainedTopologyLinearOnesevenOverOneeighteight}{1.54}  % experiments/FIGS/data/probe_ladder_v2_mlp2/analysis/seed_level_results.json :: table[retained_topology,linear,17/188].auc :: sha256 53e1ec81343ac2c1 :: ratio of the seed means of 1-AUC, coarser over finer
\newcommand{\ProbeOmaRatioRetainedTopologyLinearFourthreeOverOnesixtwo}{1.24}  % experiments/FIGS/data/probe_ladder_v2_mlp2/analysis/seed_level_results.json :: table[retained_topology,linear,43/162].auc :: sha256 53e1ec81343ac2c1 :: ratio of the seed means of 1-AUC, coarser over finer
\newcommand{\ProbeOmaRatioRetainedTopologyLinearOnesevenOverOnesixtwo}{1.54}  % experiments/FIGS/data/probe_ladder_v2_mlp2/analysis/seed_level_results.json :: table[retained_topology,linear,17/162].auc :: sha256 53e1ec81343ac2c1 :: ratio of the seed means of 1-AUC, coarser over finer
\newcommand{\ProbeOmaRatioRetainedTopologyLinearOnesevenOverFourthree}{1.25}  % experiments/FIGS/data/probe_ladder_v2_mlp2/analysis/seed_level_results.json :: table[retained_topology,linear,17/43].auc :: sha256 53e1ec81343ac2c1 :: ratio of the seed means of 1-AUC, coarser over finer
\newcommand{\ProbeAucRetainedTopologyMlpOneeighteight}{0.99394}  % experiments/FIGS/data/probe_ladder_v2_mlp2/analysis/seed_level_results.json :: levels.retained_topology.mlp[0].mean_auc :: sha256 53e1ec81343ac2c1 :: 5 seeds, 0 saturated at AUC=1
\newcommand{\ProbeAucSdRetainedTopologyMlpOneeighteight}{0.00032}  % experiments/FIGS/data/probe_ladder_v2_mlp2/analysis/seed_level_results.json :: table[retained_topology,mlp,188].auc (SD over seeds, ddof=1) :: sha256 53e1ec81343ac2c1
\newcommand{\ProbeLogOneMinusAucRetainedTopologyMlpOneeighteight}{\ensuremath{-}5.1078}  % experiments/FIGS/data/probe_ladder_v2_mlp2/analysis/seed_level_results.json :: levels.retained_topology.mlp[0].mean :: sha256 53e1ec81343ac2c1 :: natural log of 1-AUC, lower is better
\newcommand{\ProbeLogOneMinusAucSdRetainedTopologyMlpOneeighteight}{0.0522}  % experiments/FIGS/data/probe_ladder_v2_mlp2/analysis/seed_level_results.json :: levels.retained_topology.mlp[0].seed_sd :: sha256 53e1ec81343ac2c1
\newcommand{\ProbeOmaRetainedTopologyMlpOneeighteight}{\ensuremath{(6.06 \pm 0.32)\times10^{-3}}}  % experiments/FIGS/data/probe_ladder_v2_mlp2/analysis/seed_level_results.json :: table[retained_topology,mlp,188].auc :: sha256 53e1ec81343ac2c1 :: 1-AUC, mean +- SD over 5 seeds
\newcommand{\ProbeBkgLeftRetainedTopologyMlpOneeighteight}{15.8}  % experiments/FIGS/data/probe_ladder_v2_mlp2/analysis/seed_level_results.json :: levels.retained_topology.mlp[0].rejection_points['0.90'].mean_n_bkg_pass :: sha256 53e1ec81343ac2c1 :: background jets passing the cut, mean over seeds
\newcommand{\ProbeRejRetainedTopologyMlpOneeighteight}{69\,\ensuremath{\pm}\,9}  % experiments/FIGS/data/probe_ladder_v2_mlp2/analysis/seed_level_results.json :: table[retained_topology,mlp,188].rejection_points['0.90'].rejection :: sha256 53e1ec81343ac2c1 :: at 90% signal efficiency, mean +- SD over seeds; 0 of 5 seeds at the cap
\newcommand{\ProbeAucRetainedTopologyMlpOnesixtwo}{0.99387}  % experiments/FIGS/data/probe_ladder_v2_mlp2/analysis/seed_level_results.json :: levels.retained_topology.mlp[1].mean_auc :: sha256 53e1ec81343ac2c1 :: 5 seeds, 0 saturated at AUC=1
\newcommand{\ProbeAucSdRetainedTopologyMlpOnesixtwo}{0.00048}  % experiments/FIGS/data/probe_ladder_v2_mlp2/analysis/seed_level_results.json :: table[retained_topology,mlp,162].auc (SD over seeds, ddof=1) :: sha256 53e1ec81343ac2c1
\newcommand{\ProbeLogOneMinusAucRetainedTopologyMlpOnesixtwo}{\ensuremath{-}5.0964}  % experiments/FIGS/data/probe_ladder_v2_mlp2/analysis/seed_level_results.json :: levels.retained_topology.mlp[1].mean :: sha256 53e1ec81343ac2c1 :: natural log of 1-AUC, lower is better
\newcommand{\ProbeLogOneMinusAucSdRetainedTopologyMlpOnesixtwo}{0.0773}  % experiments/FIGS/data/probe_ladder_v2_mlp2/analysis/seed_level_results.json :: levels.retained_topology.mlp[1].seed_sd :: sha256 53e1ec81343ac2c1
\newcommand{\ProbeOmaRetainedTopologyMlpOnesixtwo}{\ensuremath{(6.1 \pm 0.5)\times10^{-3}}}  % experiments/FIGS/data/probe_ladder_v2_mlp2/analysis/seed_level_results.json :: table[retained_topology,mlp,162].auc :: sha256 53e1ec81343ac2c1 :: 1-AUC, mean +- SD over 5 seeds
\newcommand{\ProbeBkgLeftRetainedTopologyMlpOnesixtwo}{16.6}  % experiments/FIGS/data/probe_ladder_v2_mlp2/analysis/seed_level_results.json :: levels.retained_topology.mlp[1].rejection_points['0.90'].mean_n_bkg_pass :: sha256 53e1ec81343ac2c1 :: background jets passing the cut, mean over seeds
\newcommand{\ProbeRejRetainedTopologyMlpOnesixtwo}{67\,\ensuremath{\pm}\,14}  % experiments/FIGS/data/probe_ladder_v2_mlp2/analysis/seed_level_results.json :: table[retained_topology,mlp,162].rejection_points['0.90'].rejection :: sha256 53e1ec81343ac2c1 :: at 90% signal efficiency, mean +- SD over seeds; 0 of 5 seeds at the cap
\newcommand{\ProbeAucRetainedTopologyMlpFourthree}{0.99284}  % experiments/FIGS/data/probe_ladder_v2_mlp2/analysis/seed_level_results.json :: levels.retained_topology.mlp[2].mean_auc :: sha256 53e1ec81343ac2c1 :: 5 seeds, 0 saturated at AUC=1
\newcommand{\ProbeAucSdRetainedTopologyMlpFourthree}{0.00040}  % experiments/FIGS/data/probe_ladder_v2_mlp2/analysis/seed_level_results.json :: table[retained_topology,mlp,43].auc (SD over seeds, ddof=1) :: sha256 53e1ec81343ac2c1
\newcommand{\ProbeLogOneMinusAucRetainedTopologyMlpFourthree}{\ensuremath{-}4.9400}  % experiments/FIGS/data/probe_ladder_v2_mlp2/analysis/seed_level_results.json :: levels.retained_topology.mlp[2].mean :: sha256 53e1ec81343ac2c1 :: natural log of 1-AUC, lower is better
\newcommand{\ProbeLogOneMinusAucSdRetainedTopologyMlpFourthree}{0.0570}  % experiments/FIGS/data/probe_ladder_v2_mlp2/analysis/seed_level_results.json :: levels.retained_topology.mlp[2].seed_sd :: sha256 53e1ec81343ac2c1
\newcommand{\ProbeOmaRetainedTopologyMlpFourthree}{\ensuremath{(7.2 \pm 0.4)\times10^{-3}}}  % experiments/FIGS/data/probe_ladder_v2_mlp2/analysis/seed_level_results.json :: table[retained_topology,mlp,43].auc :: sha256 53e1ec81343ac2c1 :: 1-AUC, mean +- SD over 5 seeds
\newcommand{\ProbeBkgLeftRetainedTopologyMlpFourthree}{17.8}  % experiments/FIGS/data/probe_ladder_v2_mlp2/analysis/seed_level_results.json :: levels.retained_topology.mlp[2].rejection_points['0.90'].mean_n_bkg_pass :: sha256 53e1ec81343ac2c1 :: background jets passing the cut, mean over seeds
\newcommand{\ProbeRejRetainedTopologyMlpFourthree}{61\,\ensuremath{\pm}\,7}  % experiments/FIGS/data/probe_ladder_v2_mlp2/analysis/seed_level_results.json :: table[retained_topology,mlp,43].rejection_points['0.90'].rejection :: sha256 53e1ec81343ac2c1 :: at 90% signal efficiency, mean +- SD over seeds; 0 of 5 seeds at the cap
\newcommand{\ProbeAucRetainedTopologyMlpOneseven}{0.99081}  % experiments/FIGS/data/probe_ladder_v2_mlp2/analysis/seed_level_results.json :: levels.retained_topology.mlp[3].mean_auc :: sha256 53e1ec81343ac2c1 :: 5 seeds, 0 saturated at AUC=1
\newcommand{\ProbeAucSdRetainedTopologyMlpOneseven}{0.00062}  % experiments/FIGS/data/probe_ladder_v2_mlp2/analysis/seed_level_results.json :: table[retained_topology,mlp,17].auc (SD over seeds, ddof=1) :: sha256 53e1ec81343ac2c1
\newcommand{\ProbeLogOneMinusAucRetainedTopologyMlpOneseven}{\ensuremath{-}4.6913}  % experiments/FIGS/data/probe_ladder_v2_mlp2/analysis/seed_level_results.json :: levels.retained_topology.mlp[3].mean :: sha256 53e1ec81343ac2c1 :: natural log of 1-AUC, lower is better
\newcommand{\ProbeLogOneMinusAucSdRetainedTopologyMlpOneseven}{0.0658}  % experiments/FIGS/data/probe_ladder_v2_mlp2/analysis/seed_level_results.json :: levels.retained_topology.mlp[3].seed_sd :: sha256 53e1ec81343ac2c1
\newcommand{\ProbeOmaRetainedTopologyMlpOneseven}{\ensuremath{(9.2 \pm 0.6)\times10^{-3}}}  % experiments/FIGS/data/probe_ladder_v2_mlp2/analysis/seed_level_results.json :: table[retained_topology,mlp,17].auc :: sha256 53e1ec81343ac2c1 :: 1-AUC, mean +- SD over 5 seeds
\newcommand{\ProbeBkgLeftRetainedTopologyMlpOneseven}{21.2}  % experiments/FIGS/data/probe_ladder_v2_mlp2/analysis/seed_level_results.json :: levels.retained_topology.mlp[3].rejection_points['0.90'].mean_n_bkg_pass :: sha256 53e1ec81343ac2c1 :: background jets passing the cut, mean over seeds
\newcommand{\ProbeRejRetainedTopologyMlpOneseven}{51\,\ensuremath{\pm}\,4}  % experiments/FIGS/data/probe_ladder_v2_mlp2/analysis/seed_level_results.json :: table[retained_topology,mlp,17].rejection_points['0.90'].rejection :: sha256 53e1ec81343ac2c1 :: at 90% signal efficiency, mean +- SD over seeds; 0 of 5 seeds at the cap
\newcommand{\ProbeOmaRatioRetainedTopologyMlpOnesixtwoOverOneeighteight}{1.01}  % experiments/FIGS/data/probe_ladder_v2_mlp2/analysis/seed_level_results.json :: table[retained_topology,mlp,162/188].auc :: sha256 53e1ec81343ac2c1 :: ratio of the seed means of 1-AUC, coarser over finer
\newcommand{\ProbeOmaRatioRetainedTopologyMlpFourthreeOverOneeighteight}{1.18}  % experiments/FIGS/data/probe_ladder_v2_mlp2/analysis/seed_level_results.json :: table[retained_topology,mlp,43/188].auc :: sha256 53e1ec81343ac2c1 :: ratio of the seed means of 1-AUC, coarser over finer
\newcommand{\ProbeOmaRatioRetainedTopologyMlpOnesevenOverOneeighteight}{1.52}  % experiments/FIGS/data/probe_ladder_v2_mlp2/analysis/seed_level_results.json :: table[retained_topology,mlp,17/188].auc :: sha256 53e1ec81343ac2c1 :: ratio of the seed means of 1-AUC, coarser over finer
\newcommand{\ProbeOmaRatioRetainedTopologyMlpFourthreeOverOnesixtwo}{1.17}  % experiments/FIGS/data/probe_ladder_v2_mlp2/analysis/seed_level_results.json :: table[retained_topology,mlp,43/162].auc :: sha256 53e1ec81343ac2c1 :: ratio of the seed means of 1-AUC, coarser over finer
\newcommand{\ProbeOmaRatioRetainedTopologyMlpOnesevenOverOnesixtwo}{1.50}  % experiments/FIGS/data/probe_ladder_v2_mlp2/analysis/seed_level_results.json :: table[retained_topology,mlp,17/162].auc :: sha256 53e1ec81343ac2c1 :: ratio of the seed means of 1-AUC, coarser over finer
\newcommand{\ProbeOmaRatioRetainedTopologyMlpOnesevenOverFourthree}{1.28}  % experiments/FIGS/data/probe_ladder_v2_mlp2/analysis/seed_level_results.json :: table[retained_topology,mlp,17/43].auc :: sha256 53e1ec81343ac2c1 :: ratio of the seed means of 1-AUC, coarser over finer
\newcommand{\ProbeAucVisibleContentLinearOneeighteight}{0.99917}  % experiments/FIGS/data/probe_ladder_v2_mlp2/analysis/seed_level_results.json :: levels.visible_content.linear[0].mean_auc :: sha256 53e1ec81343ac2c1 :: 5 seeds, 0 saturated at AUC=1
\newcommand{\ProbeAucSdVisibleContentLinearOneeighteight}{0.00003}  % experiments/FIGS/data/probe_ladder_v2_mlp2/analysis/seed_level_results.json :: table[visible_content,linear,188].auc (SD over seeds, ddof=1) :: sha256 53e1ec81343ac2c1
\newcommand{\ProbeLogOneMinusAucVisibleContentLinearOneeighteight}{\ensuremath{-}7.0999}  % experiments/FIGS/data/probe_ladder_v2_mlp2/analysis/seed_level_results.json :: levels.visible_content.linear[0].mean :: sha256 53e1ec81343ac2c1 :: natural log of 1-AUC, lower is better
\newcommand{\ProbeLogOneMinusAucSdVisibleContentLinearOneeighteight}{0.0358}  % experiments/FIGS/data/probe_ladder_v2_mlp2/analysis/seed_level_results.json :: levels.visible_content.linear[0].seed_sd :: sha256 53e1ec81343ac2c1
\newcommand{\ProbeOmaVisibleContentLinearOneeighteight}{\ensuremath{(8.26 \pm 0.30)\times10^{-4}}}  % experiments/FIGS/data/probe_ladder_v2_mlp2/analysis/seed_level_results.json :: table[visible_content,linear,188].auc :: sha256 53e1ec81343ac2c1 :: 1-AUC, mean +- SD over 5 seeds
\newcommand{\ProbeBkgLeftVisibleContentLinearOneeighteight}{0.4}  % experiments/FIGS/data/probe_ladder_v2_mlp2/analysis/seed_level_results.json :: levels.visible_content.linear[0].rejection_points['0.90'].mean_n_bkg_pass :: sha256 53e1ec81343ac2c1 :: background jets passing the cut, mean over seeds
\newcommand{\ProbeRejVisibleContentLinearOneeighteight}{\ensuremath{>}538}  % experiments/FIGS/data/probe_ladder_v2_mlp2/analysis/seed_level_results.json :: table[visible_content,linear,188].rejection_points['0.90'].rejection :: sha256 53e1ec81343ac2c1 :: at 90% signal efficiency, mean +- SD over seeds; 5 of 5 seeds at the cap
\newcommand{\ProbeAucVisibleContentLinearOnesixtwo}{0.99921}  % experiments/FIGS/data/probe_ladder_v2_mlp2/analysis/seed_level_results.json :: levels.visible_content.linear[1].mean_auc :: sha256 53e1ec81343ac2c1 :: 5 seeds, 0 saturated at AUC=1
\newcommand{\ProbeAucSdVisibleContentLinearOnesixtwo}{0.00004}  % experiments/FIGS/data/probe_ladder_v2_mlp2/analysis/seed_level_results.json :: table[visible_content,linear,162].auc (SD over seeds, ddof=1) :: sha256 53e1ec81343ac2c1
\newcommand{\ProbeLogOneMinusAucVisibleContentLinearOnesixtwo}{\ensuremath{-}7.1480}  % experiments/FIGS/data/probe_ladder_v2_mlp2/analysis/seed_level_results.json :: levels.visible_content.linear[1].mean :: sha256 53e1ec81343ac2c1 :: natural log of 1-AUC, lower is better
\newcommand{\ProbeLogOneMinusAucSdVisibleContentLinearOnesixtwo}{0.0449}  % experiments/FIGS/data/probe_ladder_v2_mlp2/analysis/seed_level_results.json :: levels.visible_content.linear[1].seed_sd :: sha256 53e1ec81343ac2c1
\newcommand{\ProbeOmaVisibleContentLinearOnesixtwo}{\ensuremath{(7.87 \pm 0.35)\times10^{-4}}}  % experiments/FIGS/data/probe_ladder_v2_mlp2/analysis/seed_level_results.json :: table[visible_content,linear,162].auc :: sha256 53e1ec81343ac2c1 :: 1-AUC, mean +- SD over 5 seeds
\newcommand{\ProbeBkgLeftVisibleContentLinearOnesixtwo}{0.6}  % experiments/FIGS/data/probe_ladder_v2_mlp2/analysis/seed_level_results.json :: levels.visible_content.linear[1].rejection_points['0.90'].mean_n_bkg_pass :: sha256 53e1ec81343ac2c1 :: background jets passing the cut, mean over seeds
\newcommand{\ProbeRejVisibleContentLinearOnesixtwo}{\ensuremath{>}538}  % experiments/FIGS/data/probe_ladder_v2_mlp2/analysis/seed_level_results.json :: table[visible_content,linear,162].rejection_points['0.90'].rejection :: sha256 53e1ec81343ac2c1 :: at 90% signal efficiency, mean +- SD over seeds; 5 of 5 seeds at the cap
\newcommand{\ProbeAucVisibleContentLinearFourthree}{0.99918}  % experiments/FIGS/data/probe_ladder_v2_mlp2/analysis/seed_level_results.json :: levels.visible_content.linear[2].mean_auc :: sha256 53e1ec81343ac2c1 :: 5 seeds, 0 saturated at AUC=1
\newcommand{\ProbeAucSdVisibleContentLinearFourthree}{0.00009}  % experiments/FIGS/data/probe_ladder_v2_mlp2/analysis/seed_level_results.json :: table[visible_content,linear,43].auc (SD over seeds, ddof=1) :: sha256 53e1ec81343ac2c1
\newcommand{\ProbeLogOneMinusAucVisibleContentLinearFourthree}{\ensuremath{-}7.1124}  % experiments/FIGS/data/probe_ladder_v2_mlp2/analysis/seed_level_results.json :: levels.visible_content.linear[2].mean :: sha256 53e1ec81343ac2c1 :: natural log of 1-AUC, lower is better
\newcommand{\ProbeLogOneMinusAucSdVisibleContentLinearFourthree}{0.1137}  % experiments/FIGS/data/probe_ladder_v2_mlp2/analysis/seed_level_results.json :: levels.visible_content.linear[2].seed_sd :: sha256 53e1ec81343ac2c1
\newcommand{\ProbeOmaVisibleContentLinearFourthree}{\ensuremath{(8.2 \pm 1.0)\times10^{-4}}}  % experiments/FIGS/data/probe_ladder_v2_mlp2/analysis/seed_level_results.json :: table[visible_content,linear,43].auc :: sha256 53e1ec81343ac2c1 :: 1-AUC, mean +- SD over 5 seeds
\newcommand{\ProbeBkgLeftVisibleContentLinearFourthree}{0.4}  % experiments/FIGS/data/probe_ladder_v2_mlp2/analysis/seed_level_results.json :: levels.visible_content.linear[2].rejection_points['0.90'].mean_n_bkg_pass :: sha256 53e1ec81343ac2c1 :: background jets passing the cut, mean over seeds
\newcommand{\ProbeRejVisibleContentLinearFourthree}{\ensuremath{>}538}  % experiments/FIGS/data/probe_ladder_v2_mlp2/analysis/seed_level_results.json :: table[visible_content,linear,43].rejection_points['0.90'].rejection :: sha256 53e1ec81343ac2c1 :: at 90% signal efficiency, mean +- SD over seeds; 5 of 5 seeds at the cap
\newcommand{\ProbeAucVisibleContentLinearOneseven}{0.99768}  % experiments/FIGS/data/probe_ladder_v2_mlp2/analysis/seed_level_results.json :: levels.visible_content.linear[3].mean_auc :: sha256 53e1ec81343ac2c1 :: 5 seeds, 0 saturated at AUC=1
\newcommand{\ProbeAucSdVisibleContentLinearOneseven}{0.00020}  % experiments/FIGS/data/probe_ladder_v2_mlp2/analysis/seed_level_results.json :: table[visible_content,linear,17].auc (SD over seeds, ddof=1) :: sha256 53e1ec81343ac2c1
\newcommand{\ProbeLogOneMinusAucVisibleContentLinearOneseven}{\ensuremath{-}6.0684}  % experiments/FIGS/data/probe_ladder_v2_mlp2/analysis/seed_level_results.json :: levels.visible_content.linear[3].mean :: sha256 53e1ec81343ac2c1 :: natural log of 1-AUC, lower is better
\newcommand{\ProbeLogOneMinusAucSdVisibleContentLinearOneseven}{0.0889}  % experiments/FIGS/data/probe_ladder_v2_mlp2/analysis/seed_level_results.json :: levels.visible_content.linear[3].seed_sd :: sha256 53e1ec81343ac2c1
\newcommand{\ProbeOmaVisibleContentLinearOneseven}{\ensuremath{(2.32 \pm 0.20)\times10^{-3}}}  % experiments/FIGS/data/probe_ladder_v2_mlp2/analysis/seed_level_results.json :: table[visible_content,linear,17].auc :: sha256 53e1ec81343ac2c1 :: 1-AUC, mean +- SD over 5 seeds
\newcommand{\ProbeBkgLeftVisibleContentLinearOneseven}{3.4}  % experiments/FIGS/data/probe_ladder_v2_mlp2/analysis/seed_level_results.json :: levels.visible_content.linear[3].rejection_points['0.90'].mean_n_bkg_pass :: sha256 53e1ec81343ac2c1 :: background jets passing the cut, mean over seeds
\newcommand{\ProbeRejVisibleContentLinearOneseven}{900\,\ensuremath{\pm}\,400}  % experiments/FIGS/data/probe_ladder_v2_mlp2/analysis/seed_level_results.json :: table[visible_content,linear,17].rejection_points['0.90'].rejection :: sha256 53e1ec81343ac2c1 :: at 90% signal efficiency, mean +- SD over seeds; 0 of 5 seeds at the cap
\newcommand{\ProbeOmaRatioVisibleContentLinearOnesixtwoOverOneeighteight}{0.95}  % experiments/FIGS/data/probe_ladder_v2_mlp2/analysis/seed_level_results.json :: table[visible_content,linear,162/188].auc :: sha256 53e1ec81343ac2c1 :: ratio of the seed means of 1-AUC, coarser over finer
\newcommand{\ProbeOmaRatioVisibleContentLinearFourthreeOverOneeighteight}{0.99}  % experiments/FIGS/data/probe_ladder_v2_mlp2/analysis/seed_level_results.json :: table[visible_content,linear,43/188].auc :: sha256 53e1ec81343ac2c1 :: ratio of the seed means of 1-AUC, coarser over finer
\newcommand{\ProbeOmaRatioVisibleContentLinearOnesevenOverOneeighteight}{2.8}  % experiments/FIGS/data/probe_ladder_v2_mlp2/analysis/seed_level_results.json :: table[visible_content,linear,17/188].auc :: sha256 53e1ec81343ac2c1 :: ratio of the seed means of 1-AUC, coarser over finer
\newcommand{\ProbeOmaRatioVisibleContentLinearFourthreeOverOnesixtwo}{1.04}  % experiments/FIGS/data/probe_ladder_v2_mlp2/analysis/seed_level_results.json :: table[visible_content,linear,43/162].auc :: sha256 53e1ec81343ac2c1 :: ratio of the seed means of 1-AUC, coarser over finer
\newcommand{\ProbeOmaRatioVisibleContentLinearOnesevenOverOnesixtwo}{3.0}  % experiments/FIGS/data/probe_ladder_v2_mlp2/analysis/seed_level_results.json :: table[visible_content,linear,17/162].auc :: sha256 53e1ec81343ac2c1 :: ratio of the seed means of 1-AUC, coarser over finer
\newcommand{\ProbeOmaRatioVisibleContentLinearOnesevenOverFourthree}{2.8}  % experiments/FIGS/data/probe_ladder_v2_mlp2/analysis/seed_level_results.json :: table[visible_content,linear,17/43].auc :: sha256 53e1ec81343ac2c1 :: ratio of the seed means of 1-AUC, coarser over finer
\newcommand{\ProbeAucVisibleContentMlpOneeighteight}{0.99911}  % experiments/FIGS/data/probe_ladder_v2_mlp2/analysis/seed_level_results.json :: levels.visible_content.mlp[0].mean_auc :: sha256 53e1ec81343ac2c1 :: 5 seeds, 0 saturated at AUC=1
\newcommand{\ProbeAucSdVisibleContentMlpOneeighteight}{0.00005}  % experiments/FIGS/data/probe_ladder_v2_mlp2/analysis/seed_level_results.json :: table[visible_content,mlp,188].auc (SD over seeds, ddof=1) :: sha256 53e1ec81343ac2c1
\newcommand{\ProbeLogOneMinusAucVisibleContentMlpOneeighteight}{\ensuremath{-}7.0202}  % experiments/FIGS/data/probe_ladder_v2_mlp2/analysis/seed_level_results.json :: levels.visible_content.mlp[0].mean :: sha256 53e1ec81343ac2c1 :: natural log of 1-AUC, lower is better
\newcommand{\ProbeLogOneMinusAucSdVisibleContentMlpOneeighteight}{0.0570}  % experiments/FIGS/data/probe_ladder_v2_mlp2/analysis/seed_level_results.json :: levels.visible_content.mlp[0].seed_sd :: sha256 53e1ec81343ac2c1
\newcommand{\ProbeOmaVisibleContentMlpOneeighteight}{\ensuremath{(8.9 \pm 0.5)\times10^{-4}}}  % experiments/FIGS/data/probe_ladder_v2_mlp2/analysis/seed_level_results.json :: table[visible_content,mlp,188].auc :: sha256 53e1ec81343ac2c1 :: 1-AUC, mean +- SD over 5 seeds
\newcommand{\ProbeBkgLeftVisibleContentMlpOneeighteight}{0.4}  % experiments/FIGS/data/probe_ladder_v2_mlp2/analysis/seed_level_results.json :: levels.visible_content.mlp[0].rejection_points['0.90'].mean_n_bkg_pass :: sha256 53e1ec81343ac2c1 :: background jets passing the cut, mean over seeds
\newcommand{\ProbeRejVisibleContentMlpOneeighteight}{\ensuremath{>}538}  % experiments/FIGS/data/probe_ladder_v2_mlp2/analysis/seed_level_results.json :: table[visible_content,mlp,188].rejection_points['0.90'].rejection :: sha256 53e1ec81343ac2c1 :: at 90% signal efficiency, mean +- SD over seeds; 5 of 5 seeds at the cap
\newcommand{\ProbeAucVisibleContentMlpOnesixtwo}{0.99918}  % experiments/FIGS/data/probe_ladder_v2_mlp2/analysis/seed_level_results.json :: levels.visible_content.mlp[1].mean_auc :: sha256 53e1ec81343ac2c1 :: 5 seeds, 0 saturated at AUC=1
\newcommand{\ProbeAucSdVisibleContentMlpOnesixtwo}{0.00003}  % experiments/FIGS/data/probe_ladder_v2_mlp2/analysis/seed_level_results.json :: table[visible_content,mlp,162].auc (SD over seeds, ddof=1) :: sha256 53e1ec81343ac2c1
\newcommand{\ProbeLogOneMinusAucVisibleContentMlpOnesixtwo}{\ensuremath{-}7.1010}  % experiments/FIGS/data/probe_ladder_v2_mlp2/analysis/seed_level_results.json :: levels.visible_content.mlp[1].mean :: sha256 53e1ec81343ac2c1 :: natural log of 1-AUC, lower is better
\newcommand{\ProbeLogOneMinusAucSdVisibleContentMlpOnesixtwo}{0.0371}  % experiments/FIGS/data/probe_ladder_v2_mlp2/analysis/seed_level_results.json :: levels.visible_content.mlp[1].seed_sd :: sha256 53e1ec81343ac2c1
\newcommand{\ProbeOmaVisibleContentMlpOnesixtwo}{\ensuremath{(8.25 \pm 0.30)\times10^{-4}}}  % experiments/FIGS/data/probe_ladder_v2_mlp2/analysis/seed_level_results.json :: table[visible_content,mlp,162].auc :: sha256 53e1ec81343ac2c1 :: 1-AUC, mean +- SD over 5 seeds
\newcommand{\ProbeBkgLeftVisibleContentMlpOnesixtwo}{0.2}  % experiments/FIGS/data/probe_ladder_v2_mlp2/analysis/seed_level_results.json :: levels.visible_content.mlp[1].rejection_points['0.90'].mean_n_bkg_pass :: sha256 53e1ec81343ac2c1 :: background jets passing the cut, mean over seeds
\newcommand{\ProbeRejVisibleContentMlpOnesixtwo}{\ensuremath{>}538}  % experiments/FIGS/data/probe_ladder_v2_mlp2/analysis/seed_level_results.json :: table[visible_content,mlp,162].rejection_points['0.90'].rejection :: sha256 53e1ec81343ac2c1 :: at 90% signal efficiency, mean +- SD over seeds; 5 of 5 seeds at the cap
\newcommand{\ProbeAucVisibleContentMlpFourthree}{0.99917}  % experiments/FIGS/data/probe_ladder_v2_mlp2/analysis/seed_level_results.json :: levels.visible_content.mlp[2].mean_auc :: sha256 53e1ec81343ac2c1 :: 5 seeds, 0 saturated at AUC=1
\newcommand{\ProbeAucSdVisibleContentMlpFourthree}{0.00006}  % experiments/FIGS/data/probe_ladder_v2_mlp2/analysis/seed_level_results.json :: table[visible_content,mlp,43].auc (SD over seeds, ddof=1) :: sha256 53e1ec81343ac2c1
\newcommand{\ProbeLogOneMinusAucVisibleContentMlpFourthree}{\ensuremath{-}7.0957}  % experiments/FIGS/data/probe_ladder_v2_mlp2/analysis/seed_level_results.json :: levels.visible_content.mlp[2].mean :: sha256 53e1ec81343ac2c1 :: natural log of 1-AUC, lower is better
\newcommand{\ProbeLogOneMinusAucSdVisibleContentMlpFourthree}{0.0773}  % experiments/FIGS/data/probe_ladder_v2_mlp2/analysis/seed_level_results.json :: levels.visible_content.mlp[2].seed_sd :: sha256 53e1ec81343ac2c1
\newcommand{\ProbeOmaVisibleContentMlpFourthree}{\ensuremath{(8.3 \pm 0.6)\times10^{-4}}}  % experiments/FIGS/data/probe_ladder_v2_mlp2/analysis/seed_level_results.json :: table[visible_content,mlp,43].auc :: sha256 53e1ec81343ac2c1 :: 1-AUC, mean +- SD over 5 seeds
\newcommand{\ProbeBkgLeftVisibleContentMlpFourthree}{0.4}  % experiments/FIGS/data/probe_ladder_v2_mlp2/analysis/seed_level_results.json :: levels.visible_content.mlp[2].rejection_points['0.90'].mean_n_bkg_pass :: sha256 53e1ec81343ac2c1 :: background jets passing the cut, mean over seeds
\newcommand{\ProbeRejVisibleContentMlpFourthree}{\ensuremath{>}538}  % experiments/FIGS/data/probe_ladder_v2_mlp2/analysis/seed_level_results.json :: table[visible_content,mlp,43].rejection_points['0.90'].rejection :: sha256 53e1ec81343ac2c1 :: at 90% signal efficiency, mean +- SD over seeds; 5 of 5 seeds at the cap
\newcommand{\ProbeAucVisibleContentMlpOneseven}{0.99793}  % experiments/FIGS/data/probe_ladder_v2_mlp2/analysis/seed_level_results.json :: levels.visible_content.mlp[3].mean_auc :: sha256 53e1ec81343ac2c1 :: 5 seeds, 0 saturated at AUC=1
\newcommand{\ProbeAucSdVisibleContentMlpOneseven}{0.00020}  % experiments/FIGS/data/probe_ladder_v2_mlp2/analysis/seed_level_results.json :: table[visible_content,mlp,17].auc (SD over seeds, ddof=1) :: sha256 53e1ec81343ac2c1
\newcommand{\ProbeLogOneMinusAucVisibleContentMlpOneseven}{\ensuremath{-}6.1839}  % experiments/FIGS/data/probe_ladder_v2_mlp2/analysis/seed_level_results.json :: levels.visible_content.mlp[3].mean :: sha256 53e1ec81343ac2c1 :: natural log of 1-AUC, lower is better
\newcommand{\ProbeLogOneMinusAucSdVisibleContentMlpOneseven}{0.1014}  % experiments/FIGS/data/probe_ladder_v2_mlp2/analysis/seed_level_results.json :: levels.visible_content.mlp[3].seed_sd :: sha256 53e1ec81343ac2c1
\newcommand{\ProbeOmaVisibleContentMlpOneseven}{\ensuremath{(2.07 \pm 0.20)\times10^{-3}}}  % experiments/FIGS/data/probe_ladder_v2_mlp2/analysis/seed_level_results.json :: table[visible_content,mlp,17].auc :: sha256 53e1ec81343ac2c1 :: 1-AUC, mean +- SD over 5 seeds
\newcommand{\ProbeBkgLeftVisibleContentMlpOneseven}{3.4}  % experiments/FIGS/data/probe_ladder_v2_mlp2/analysis/seed_level_results.json :: levels.visible_content.mlp[3].rejection_points['0.90'].mean_n_bkg_pass :: sha256 53e1ec81343ac2c1 :: background jets passing the cut, mean over seeds
\newcommand{\ProbeRejVisibleContentMlpOneseven}{\ensuremath{\geq}638\ensuremath{^{\ast}}}  % experiments/FIGS/data/probe_ladder_v2_mlp2/analysis/seed_level_results.json :: table[visible_content,mlp,17].rejection_points['0.90'].rejection :: sha256 53e1ec81343ac2c1 :: at 90% signal efficiency, mean +- SD over seeds; 1 of 5 seeds at the cap
\newcommand{\ProbeOmaRatioVisibleContentMlpOnesixtwoOverOneeighteight}{0.92}  % experiments/FIGS/data/probe_ladder_v2_mlp2/analysis/seed_level_results.json :: table[visible_content,mlp,162/188].auc :: sha256 53e1ec81343ac2c1 :: ratio of the seed means of 1-AUC, coarser over finer
\newcommand{\ProbeOmaRatioVisibleContentMlpFourthreeOverOneeighteight}{0.93}  % experiments/FIGS/data/probe_ladder_v2_mlp2/analysis/seed_level_results.json :: table[visible_content,mlp,43/188].auc :: sha256 53e1ec81343ac2c1 :: ratio of the seed means of 1-AUC, coarser over finer
\newcommand{\ProbeOmaRatioVisibleContentMlpOnesevenOverOneeighteight}{2.3}  % experiments/FIGS/data/probe_ladder_v2_mlp2/analysis/seed_level_results.json :: table[visible_content,mlp,17/188].auc :: sha256 53e1ec81343ac2c1 :: ratio of the seed means of 1-AUC, coarser over finer
\newcommand{\ProbeOmaRatioVisibleContentMlpFourthreeOverOnesixtwo}{1.01}  % experiments/FIGS/data/probe_ladder_v2_mlp2/analysis/seed_level_results.json :: table[visible_content,mlp,43/162].auc :: sha256 53e1ec81343ac2c1 :: ratio of the seed means of 1-AUC, coarser over finer
\newcommand{\ProbeOmaRatioVisibleContentMlpOnesevenOverOnesixtwo}{2.5}  % experiments/FIGS/data/probe_ladder_v2_mlp2/analysis/seed_level_results.json :: table[visible_content,mlp,17/162].auc :: sha256 53e1ec81343ac2c1 :: ratio of the seed means of 1-AUC, coarser over finer
\newcommand{\ProbeOmaRatioVisibleContentMlpOnesevenOverFourthree}{2.5}  % experiments/FIGS/data/probe_ladder_v2_mlp2/analysis/seed_level_results.json :: table[visible_content,mlp,17/43].auc :: sha256 53e1ec81343ac2c1 :: ratio of the seed means of 1-AUC, coarser over finer
\newcommand{\VocabSizeLoneeighteight}{188}  % configs/labelmaps/rung_label_maps.v1.csv :: column L188 (distinct group count) :: sha256 412f9aa3b58db412
\newcommand{\VocabSizeLonesixtwo}{162}  % configs/labelmaps/rung_label_maps.v1.csv :: column L162 (distinct group count) :: sha256 412f9aa3b58db412
\newcommand{\VocabSizeRsixthreeQone}{64}  % configs/labelmaps/rung_label_maps.v1.csv :: column R63_Q1 (distinct group count) :: sha256 412f9aa3b58db412
\newcommand{\VocabSizeRfourtwoQone}{43}  % configs/labelmaps/rung_label_maps.v1.csv :: column R42_Q1 (distinct group count) :: sha256 412f9aa3b58db412
\newcommand{\VocabSizeRtwonineQone}{30}  % configs/labelmaps/rung_label_maps.v1.csv :: column R29_Q1 (distinct group count) :: sha256 412f9aa3b58db412
\newcommand{\VocabSizeRonesixQone}{17}  % configs/labelmaps/rung_label_maps.v1.csv :: column R16_Q1 (distinct group count) :: sha256 412f9aa3b58db412
\newcommand{\VocabSizeRthreeVIS}{4}  % configs/labelmaps/rung_label_maps.v1.csv :: column R3_VIS (distinct group count) :: sha256 412f9aa3b58db412
\newcommand{\VocabSizeRoneQone}{2}  % configs/labelmaps/rung_label_maps.v1.csv :: column R1_Q1 (distinct group count) :: sha256 412f9aa3b58db412
\newcommand{\ProbeSplitTrain}{60\%}  % experiments/EVAL/probe.py :: SPLIT_FRACTIONS :: sha256 485db0966e660588 :: share of the probe sample; the same jets for every model
\newcommand{\ProbeSplitVal}{20\%}  % experiments/EVAL/probe.py :: SPLIT_FRACTIONS :: sha256 485db0966e660588 :: share of the probe sample; the same jets for every model
\newcommand{\ProbeSplitTest}{20\%}  % experiments/EVAL/probe.py :: SPLIT_FRACTIONS :: sha256 485db0966e660588 :: share of the probe sample; the same jets for every model
\newcommand{\ProbeMlpHidden}{256}  % experiments/EVAL/probe.py :: _fit_mlp :: sha256 485db0966e660588 :: hidden width of the MLP probe
\newcommand{\ProbeMlpInits}{3}  % experiments/EVAL/probe.py :: MLP_SEEDS :: sha256 485db0966e660588 :: MLP initialisations whose scores are averaged per cell
\newcommand{\ProbeMlpLrDrop}{10}  % experiments/EVAL/probe.py :: MLP_SCHEDULE['lr_factor'] :: sha256 485db0966e660588 :: factor the MLP's learning rate drops by on a plateau
\newcommand{\ProbeMlpPlateaus}{three}  % experiments/EVAL/probe.py :: MLP_SCHEDULE: stop_patience / (lr_patience + 1) :: sha256 485db0966e660588 :: patience windows without gain after which the MLP stops
\newcommand{\ProbeCMin}{\ensuremath{10^{-2}}}  % experiments/EVAL/probe.py :: C_GRID :: sha256 485db0966e660588 :: end of the logistic-regression inverse-regularisation grid
\newcommand{\ProbeCMax}{\ensuremath{10^{2}}}  % experiments/EVAL/probe.py :: C_GRID :: sha256 485db0966e660588 :: end of the logistic-regression inverse-regularisation grid
\newcommand{\ProbeNBkgTestBvcFourprong}{3{,}158}  % experiments/FIGS/data/probe_ladder_v2_mlp2/s1.json :: tasks.bvc_4prong.n_background_test :: sha256 f7e1de0e43be89b8 :: test-sample background jets
\newcommand{\ProbeNBkgTestBvcQcd}{5{,}371}  % experiments/FIGS/data/probe_ladder_v2_mlp2/s1.json :: tasks.bvc_qcd.n_background_test :: sha256 f7e1de0e43be89b8 :: test-sample background jets
\newcommand{\ProbeNBkgTestBvcResonant}{11{,}876}  % experiments/FIGS/data/probe_ladder_v2_mlp2/s1.json :: tasks.bvc_resonant.n_background_test :: sha256 f7e1de0e43be89b8 :: test-sample background jets
\newcommand{\ProbeNBkgTestEeVsMm}{4{,}708}  % experiments/FIGS/data/probe_ladder_v2_mlp2/s1.json :: tasks.ee_vs_mm.n_background_test :: sha256 f7e1de0e43be89b8 :: test-sample background jets
\newcommand{\ProbeNBkgTestRetainedTopology}{1{,}078}  % experiments/FIGS/data/probe_ladder_v2_mlp2/s1.json :: tasks.retained_topology.n_background_test :: sha256 f7e1de0e43be89b8 :: test-sample background jets
\newcommand{\ProbeNBkgTestVisibleContent}{2{,}554}  % experiments/FIGS/data/probe_ladder_v2_mlp2/s1.json :: tasks.visible_content.n_background_test :: sha256 f7e1de0e43be89b8 :: test-sample background jets
\newcommand{\UseDiesAtSophonEqfour}{43}  % configs/labelmaps/usecase_survival.v1.json :: sophon_eq4.dies_at :: sha256 122275213c349593 :: first vocabulary size at which it is not constructible
\newcommand{\UseLastAliveSophonEqfour}{64}  % configs/labelmaps/usecase_survival.v1.json :: sophon_eq4.last_rung_alive :: sha256 122275213c349593
\newcommand{\UseDiesAtSophonEqsixAone}{64}  % configs/labelmaps/usecase_survival.v1.json :: sophon_eq6_a1.dies_at :: sha256 122275213c349593 :: first vocabulary size at which it is not constructible
\newcommand{\UseLastAliveSophonEqsixAone}{162}  % configs/labelmaps/usecase_survival.v1.json :: sophon_eq6_a1.last_rung_alive :: sha256 122275213c349593
\newcommand{\UseDiesAtSophonEqsixAthree}{64}  % configs/labelmaps/usecase_survival.v1.json :: sophon_eq6_a3.dies_at :: sha256 122275213c349593 :: first vocabulary size at which it is not constructible
\newcommand{\UseLastAliveSophonEqsixAthree}{162}  % configs/labelmaps/usecase_survival.v1.json :: sophon_eq6_a3.last_rung_alive :: sha256 122275213c349593
\newcommand{\UseDiesAtSophonEqsevenNine}{64}  % configs/labelmaps/usecase_survival.v1.json :: sophon_eq7_9.dies_at :: sha256 122275213c349593 :: first vocabulary size at which it is not constructible
\newcommand{\UseLastAliveSophonEqsevenNine}{162}  % configs/labelmaps/usecase_survival.v1.json :: sophon_eq7_9.last_rung_alive :: sha256 122275213c349593
\newcommand{\UseDiesAtVcbEqone}{64}  % configs/labelmaps/usecase_survival.v1.json :: vcb_eq1.dies_at :: sha256 122275213c349593 :: first vocabulary size at which it is not constructible
\newcommand{\UseLastAliveVcbEqone}{162}  % configs/labelmaps/usecase_survival.v1.json :: vcb_eq1.last_rung_alive :: sha256 122275213c349593
\newcommand{\UseDiesAtGnthreexQcdSubclass}{162}  % configs/labelmaps/usecase_survival.v1.json :: gn3x_qcd_subclass.dies_at :: sha256 122275213c349593 :: first vocabulary size at which it is not constructible
\newcommand{\UseLastAliveGnthreexQcdSubclass}{188}  % configs/labelmaps/usecase_survival.v1.json :: gn3x_qcd_subclass.last_rung_alive :: sha256 122275213c349593
\newcommand{\UseWlikeLastAlive}{162}  % configs/labelmaps/rung_label_maps.v1.csv :: groups at each level vs {label_X_cs, label_X_qq} and that plus the QCD classes :: sha256 412f9aa3b58db412 :: last level at which Sophon's W-like score g_W(2) is constructible
\newcommand{\UseWlikeDiesAt}{64}  % configs/labelmaps/rung_label_maps.v1.csv :: groups at each level vs {label_X_cs, label_X_qq} and that plus the QCD classes :: sha256 412f9aa3b58db412 :: first level at which it is not
\newcommand{\DesignEpochs}{80}  % experiments/MTX/k8s/job-mtx-l188-s1-raunav.yaml :: --num-epochs :: sha256 96fef4b6337aaa27
\newcommand{\DesignSamplesPerEpoch}{10{,}240{,}000}  % experiments/MTX/k8s/job-mtx-l188-s1-raunav.yaml :: --samples-per-epoch :: sha256 96fef4b6337aaa27
\newcommand{\DesignExamplesSeen}{\ensuremath{8.192\times10^{8}}}  % experiments/MTX/k8s/job-mtx-l188-s1-raunav.yaml :: --num-epochs x --samples-per-epoch :: sha256 96fef4b6337aaa27 :: training examples drawn over the run
\newcommand{\DesignPassesSelected}{7.45}  % experiments/FIGS/data/v2/jc2_train/counts.json :: n_selected :: sha256 b0b5a8b63aa808c0 :: --num-epochs x --samples-per-epoch of job-mtx-l188-s1-raunav.yaml / n_selected: passes over the selected training jets
\newcommand{\DesignBatchSize}{512}  % experiments/MTX/k8s/job-mtx-l188-s1-raunav.yaml :: --batch-size :: sha256 96fef4b6337aaa27
\newcommand{\DesignStartLr}{\ensuremath{5\times10^{-4}}}  % experiments/MTX/k8s/job-mtx-l188-s1-raunav.yaml :: --start-lr :: sha256 96fef4b6337aaa27
\newcommand{\DesignHeadHidden}{512}  % experiments/MTX/k8s/job-mtx-l188-s1-raunav.yaml :: -o fc_params (hidden width) :: sha256 96fef4b6337aaa27
\newcommand{\DesignParticleBlocks}{8}  % experiments/E1/ParT_sophon_arch_10c.py :: num_layers=(\d+) :: sha256 5793faa7d55a8b0d
\newcommand{\DesignClassBlocks}{2}  % experiments/E1/ParT_sophon_arch_10c.py :: num_cls_layers=(\d+) :: sha256 5793faa7d55a8b0d
\newcommand{\DesignHeads}{8}  % experiments/E1/ParT_sophon_arch_10c.py :: num_heads=(\d+) :: sha256 5793faa7d55a8b0d
\newcommand{\DesignEmbedDim}{128}  % experiments/E1/ParT_sophon_arch_10c.py :: \bembed_dims=\[\d+, \d+, (\d+)\] :: sha256 5793faa7d55a8b0d
\newcommand{\DesignMpmMaskRate}{40\%}  % experiments/MTX/mpm.py :: DEFAULT_MASK_RATE :: sha256 49f0c8c6666d0211 :: share of each jet's particles hidden from the encoder
\newcommand{\DesignPtMin}{200}  % configs/arms/L188.yaml :: selection :: sha256 d5a188dc80885515 :: GeV
\newcommand{\DesignPtMax}{2{,}500}  % configs/arms/L188.yaml :: selection :: sha256 d5a188dc80885515 :: GeV
\newcommand{\DesignMsdMin}{20}  % configs/arms/L188.yaml :: selection :: sha256 d5a188dc80885515 :: GeV
\newcommand{\DesignMsdMax}{500}  % configs/arms/L188.yaml :: selection :: sha256 d5a188dc80885515 :: GeV
\newcommand{\LitJetClassTwoNJetsMillion}{139}  % experiments/FIGS/data/literature_facts.json :: facts[0].value :: sha256 713dba2a5e6e1dd1 :: arXiv:2405.12972v1 App. A, p. 11
\newcommand{\LitJetClassTwoNResonantClasses}{161}  % experiments/FIGS/data/literature_facts.json :: facts[1].value :: sha256 713dba2a5e6e1dd1 :: arXiv:2405.12972v1 App. A, p. 11
\newcommand{\LitJetClassTwoNQcdClasses}{27}  % experiments/FIGS/data/literature_facts.json :: facts[2].value :: sha256 713dba2a5e6e1dd1 :: arXiv:2405.12972v1 App. A, p. 11
\newcommand{\LitJetClassTwoMaxResonanceMass}{500}  % experiments/FIGS/data/literature_facts.json :: facts[3].value :: sha256 713dba2a5e6e1dd1 :: arXiv:2405.12972v1 Sec. II, p. 2
\newcommand{\LitMadgraphVersion}{2.9.18}  % experiments/FIGS/data/literature_facts.json :: facts[4].value :: sha256 713dba2a5e6e1dd1 :: arXiv:2405.12972v1 App. A, p. 11
\newcommand{\LitPythiaVersion}{8.3}  % experiments/FIGS/data/literature_facts.json :: facts[5].value :: sha256 713dba2a5e6e1dd1 :: arXiv:2405.12972v1 App. A, p. 11
\newcommand{\LitDelphesVersion}{3}  % experiments/FIGS/data/literature_facts.json :: facts[6].value :: sha256 713dba2a5e6e1dd1 :: arXiv:2405.12972v1 App. A, p. 11
\newcommand{\LitJetClassTwoPileupMean}{50}  % experiments/FIGS/data/literature_facts.json :: facts[7].value :: sha256 713dba2a5e6e1dd1 :: arXiv:2405.12972v1 Sec. II, p. 2
\newcommand{\LitJetClassTwoJetRadius}{0.8}  % experiments/FIGS/data/literature_facts.json :: facts[8].value :: sha256 713dba2a5e6e1dd1 :: arXiv:2405.12972v1 Sec. II, p. 2
\newcommand{\LitAtlasUnifiedEffDropPercent}{10}  % experiments/FIGS/data/literature_facts.json :: facts[9].value :: sha256 713dba2a5e6e1dd1 :: ATL-PHYS-PUB-2026-013 Sec. 6.5, p. 22
\newcommand{\LitSophonPaperParticleBlocks}{6}  % experiments/FIGS/data/literature_facts.json :: facts[10].value :: sha256 713dba2a5e6e1dd1 :: arXiv:2405.12972v1 App. B.1, p. 12
\newcommand{\LitSophonSamplesPerEpochMillion}{10}  % experiments/FIGS/data/literature_facts.json :: facts[11].value :: sha256 713dba2a5e6e1dd1 :: arXiv:2405.12972v1 App. B.1, p. 12
\newcommand{\LitSophonSubsetSamplesPerEpochMillion}{2.5}  % experiments/FIGS/data/literature_facts.json :: facts[12].value :: sha256 713dba2a5e6e1dd1 :: arXiv:2405.12972v1 App. B.2, p. 12
\newcommand{\LitGlopartNClasses}{314}  % experiments/FIGS/data/literature_facts.json :: facts[13].value :: sha256 713dba2a5e6e1dd1 :: C. Li, PhD thesis, CERN-THESIS-2024-281 Sec. 11.2.1, p. 188
\newcommand{\LitPartFtEpochs}{20}  % experiments/FIGS/data/literature_facts.json :: facts[14].value :: sha256 713dba2a5e6e1dd1 :: arXiv:2202.03772 Sec. 5.2, p. 8
\newcommand{\LitMehtaNCategories}{eight}  % experiments/FIGS/data/literature_facts.json :: facts[15].value :: sha256 713dba2a5e6e1dd1 :: arXiv:2605.05556v2 Abstract
\newcommand{\LitHongPretrainRuns}{2}  % experiments/FIGS/data/literature_facts.json :: facts[16].value :: sha256 713dba2a5e6e1dd1 :: arXiv:2410.23129 App. A.1.1, p. 18
\newcommand{\LitSophonLookaheadK}{6}  % experiments/FIGS/data/literature_facts.json :: facts[17].value :: sha256 713dba2a5e6e1dd1 :: arXiv:2405.12972v1 App. B.1, p. 12
\newcommand{\LitSophonLookaheadAlpha}{0.5}  % experiments/FIGS/data/literature_facts.json :: facts[18].value :: sha256 713dba2a5e6e1dd1 :: arXiv:2405.12972v1 App. B.1, p. 12
\newcommand{\LitSophonRadamBetaOne}{0.95}  % experiments/FIGS/data/literature_facts.json :: facts[19].value :: sha256 713dba2a5e6e1dd1 :: arXiv:2405.12972v1 App. B.1, p. 12
\newcommand{\LitSophonRadamBetaTwo}{0.999}  % experiments/FIGS/data/literature_facts.json :: facts[20].value :: sha256 713dba2a5e6e1dd1 :: arXiv:2405.12972v1 App. B.1, p. 12
\newcommand{\LitSophonRadamEps}{\ensuremath{10^{-5}}}  % experiments/FIGS/data/literature_facts.json :: facts[21].value :: sha256 713dba2a5e6e1dd1 :: arXiv:2405.12972v1 App. B.1, p. 12
\newcommand{\LitGnTwoAuxRejectionDropPercent}{30}  % experiments/FIGS/data/literature_facts.json :: facts[22].value :: sha256 713dba2a5e6e1dd1 :: arXiv:2505.19689v2 (Nature Commun. 17 (2026) 541) p. 9
\newcommand{\LitCmsPartNClasses}{37}  % experiments/FIGS/data/literature_facts.json :: facts[23].value :: sha256 713dba2a5e6e1dd1 :: arXiv:2604.09809v1 (CMS-JME-25-001) Sec. 5.2, p. 7
\newcommand{\LitCmsPartMassLambda}{0.05}  % experiments/FIGS/data/literature_facts.json :: facts[24].value :: sha256 713dba2a5e6e1dd1 :: arXiv:2604.09809v1 (CMS-JME-25-001) Sec. 5.2, p. 7
\newcommand{\LitGlopartVthreeNClasses}{374}  % experiments/FIGS/data/literature_facts.json :: facts[25].value :: sha256 713dba2a5e6e1dd1 :: CMS-DP-2026-104 (CMS public page) Architecture and outputs
\newcommand{\LitHuhCoarseClasses}{127}  % experiments/FIGS/data/literature_facts.json :: facts[26].value :: sha256 713dba2a5e6e1dd1 :: arXiv:1608.08614v2 Sec. 5.3, p. 5
\newcommand{\LitHuhKnnGapPercent}{15}  % experiments/FIGS/data/literature_facts.json :: facts[27].value :: sha256 713dba2a5e6e1dd1 :: arXiv:1608.08614v2 Sec. 5.3, p. 5
\newcommand{\LitRsthreelRuns}{3}  % experiments/FIGS/data/literature_facts.json :: facts[28].value :: sha256 713dba2a5e6e1dd1 :: arXiv:2403.07066v2 (Phys. Rev. D 111 (2025) 032010) Sec. IV, p. 5
\newcommand{\LitSophonParamsMillion}{2.3}  % experiments/FIGS/data/literature_facts.json :: facts[29].value :: sha256 713dba2a5e6e1dd1 :: arXiv:2405.12972v1 App. B.1, p. 12
\newcommand{\LitGlopartVoneNClasses}{37}  % experiments/FIGS/data/literature_facts.json :: facts[30].value :: sha256 713dba2a5e6e1dd1 :: CMS-DP-2026-104 (CMS public page) Model generations
\newcommand{\LitSophonXMass}{200}  % experiments/FIGS/data/literature_facts.json :: facts[31].value :: sha256 713dba2a5e6e1dd1 :: arXiv:2405.12972v1 Sec. III, p. 3
\newcommand{\LitSophonXLumi}{40}  % experiments/FIGS/data/literature_facts.json :: facts[32].value :: sha256 713dba2a5e6e1dd1 :: arXiv:2405.12972v1 Sec. III, Fig. 2
\newcommand{\LitSophonXWindowLo}{150}  % experiments/FIGS/data/literature_facts.json :: facts[33].value :: sha256 713dba2a5e6e1dd1 :: arXiv:2405.12972v1 Sec. III, Fig. 2
\newcommand{\LitSophonXWindowHi}{230}  % experiments/FIGS/data/literature_facts.json :: facts[34].value :: sha256 713dba2a5e6e1dd1 :: arXiv:2405.12972v1 Sec. III, Fig. 2
\newcommand{\LitSophonAblationNClasses}{42}  % experiments/FIGS/data/literature_facts.json :: facts[35].value :: sha256 713dba2a5e6e1dd1 :: arXiv:2405.12972v1 Sec. III
\newcommand{\LitSophonTransferSamplesMillion}{2.5}  % experiments/FIGS/data/literature_facts.json :: facts[36].value :: sha256 713dba2a5e6e1dd1 :: arXiv:2405.12972v1 App. B.5, p. 13
\newcommand{\LitSophonTransferBatch}{1024}  % experiments/FIGS/data/literature_facts.json :: facts[37].value :: sha256 713dba2a5e6e1dd1 :: arXiv:2405.12972v1 App. B.5, p. 13
\newcommand{\LitSophonScratchEpochs}{50}  % experiments/FIGS/data/literature_facts.json :: facts[38].value :: sha256 713dba2a5e6e1dd1 :: arXiv:2405.12972v1 App. B.3
\newcommand{\LitHuhFineClasses}{1000}  % experiments/FIGS/data/literature_facts.json :: facts[39].value :: sha256 713dba2a5e6e1dd1 :: arXiv:1608.08614v2 Sec. 1, p. 1
\newcommand{\LitHuhPascalDropMap}{2.8}  % experiments/FIGS/data/literature_facts.json :: facts[40].value :: sha256 713dba2a5e6e1dd1 :: arXiv:1608.08614v2 Sec. 1, p. 1
\newcommand{\LitHongSuperclasses}{11}  % experiments/FIGS/data/literature_facts.json :: facts[41].value :: sha256 713dba2a5e6e1dd1 :: arXiv:2410.23129 Sec. 6.2
\newcommand{\LitHongDeclineClasses}{\ensuremath{10^{3}}}  % experiments/FIGS/data/literature_facts.json :: facts[42].value :: sha256 713dba2a5e6e1dd1 :: arXiv:2410.23129 Sec. 6.2
\newcommand{\LitObjectivesFullJets}{\ensuremath{1.2\times10^{6}}}  % experiments/FIGS/data/literature_facts.json :: facts[43].value :: sha256 713dba2a5e6e1dd1 :: arXiv:2606.14870 Sec. IV A
\newcommand{\LitObjectivesFewJets}{\ensuremath{10^{4}}}  % experiments/FIGS/data/literature_facts.json :: facts[44].value :: sha256 713dba2a5e6e1dd1 :: arXiv:2606.14870 Sec. IV A
\newcommand{\LitMpmBenefitJets}{10k}  % experiments/FIGS/data/literature_facts.json :: facts[45].value :: sha256 713dba2a5e6e1dd1 :: arXiv:2401.13537 Sec. V
\newcommand{\LitPartTopRejScratch}{1602}  % experiments/FIGS/data/literature_facts.json :: facts[46].value :: sha256 713dba2a5e6e1dd1 :: arXiv:2202.03772 Table 5
\newcommand{\LitPartTopRejFt}{2766}  % experiments/FIGS/data/literature_facts.json :: facts[47].value :: sha256 713dba2a5e6e1dd1 :: arXiv:2202.03772 Table 5
\newcommand{\LitMahajanLinearGapPercent}{7}  % experiments/FIGS/data/literature_facts.json :: facts[48].value :: sha256 713dba2a5e6e1dd1 :: arXiv:1805.00932 Sec. 3.1.2
\newcommand{\LitBukchinCoarse}{33.83}  % experiments/FIGS/data/literature_facts.json :: facts[49].value :: sha256 713dba2a5e6e1dd1 :: arXiv:2012.03515 Table 2, LIVING-17 all-way
\newcommand{\LitBukchinFine}{58.95}  % experiments/FIGS/data/literature_facts.json :: facts[50].value :: sha256 713dba2a5e6e1dd1 :: arXiv:2012.03515 Table 2, LIVING-17 all-way
\newcommand{\LitKornblithRho}{\ensuremath{-0.93}}  % experiments/FIGS/data/literature_facts.json :: facts[51].value :: sha256 713dba2a5e6e1dd1 :: arXiv:2010.16402 Sec. 3.4
\newcommand{\LitPartTopSigEffPercent}{30}  % experiments/FIGS/data/literature_facts.json :: facts[52].value :: sha256 713dba2a5e6e1dd1 :: arXiv:2202.03772 Table 5 (column Rej30%)
\newcommand{\LitPartTopKinFeatures}{7}  % experiments/FIGS/data/literature_facts.json :: facts[53].value :: sha256 713dba2a5e6e1dd1 :: arXiv:2202.03772 Sec. 5.2
\newcommand{\LitBukchinAncor}{45.14}  % experiments/FIGS/data/literature_facts.json :: facts[54].value :: sha256 713dba2a5e6e1dd1 :: arXiv:2012.03515 Table 2, LIVING-17 all-way
\newcommand{\LitMahajanFtGapPp}{0.1}  % experiments/FIGS/data/literature_facts.json :: facts[55].value :: sha256 713dba2a5e6e1dd1 :: arXiv:1805.00932 Fig. 1, ImageNet-9k target
\newcommand{\LitSophonTransferLr}{5 × 10−4}  % experiments/FIGS/data/literature_facts.json :: facts[56].value :: sha256 713dba2a5e6e1dd1 :: arXiv:2405.12972v1 App. B.5
\newcommand{\LitQgJetRadius}{0.4}  % experiments/FIGS/data/literature_facts.json :: facts[57].value :: sha256 713dba2a5e6e1dd1 :: arXiv:1810.05165 Sec. 4.1
\newcommand{\LitQgPtMin}{500}  % experiments/FIGS/data/literature_facts.json :: facts[58].value :: sha256 713dba2a5e6e1dd1 :: arXiv:1810.05165 Sec. 4.1
\newcommand{\LitQgPtMax}{550}  % experiments/FIGS/data/literature_facts.json :: facts[59].value :: sha256 713dba2a5e6e1dd1 :: arXiv:1810.05165 Sec. 4.1
\newcommand{\LitPartTopAucScratch}{0.9858}  % experiments/FIGS/data/literature_facts.json :: facts[60].value :: sha256 713dba2a5e6e1dd1 :: arXiv:2202.03772 Table 5
\newcommand{\LitPartTopAucFt}{0.9877}  % experiments/FIGS/data/literature_facts.json :: facts[61].value :: sha256 713dba2a5e6e1dd1 :: arXiv:2202.03772 Table 5
\newcommand{\LitPartQgAucScratch}{0.9121}  % experiments/FIGS/data/literature_facts.json :: facts[62].value :: sha256 713dba2a5e6e1dd1 :: arXiv:2202.03772 Table 6 (exp)
\newcommand{\LitPartQgRejScratch}{101.2}  % experiments/FIGS/data/literature_facts.json :: facts[63].value :: sha256 713dba2a5e6e1dd1 :: arXiv:2202.03772 Table 6 (exp)
\newcommand{\LitPartQgAucFt}{0.9151}  % experiments/FIGS/data/literature_facts.json :: facts[64].value :: sha256 713dba2a5e6e1dd1 :: arXiv:2202.03772 Table 6 (exp)
\newcommand{\LitPartQgRejFt}{107.9}  % experiments/FIGS/data/literature_facts.json :: facts[65].value :: sha256 713dba2a5e6e1dd1 :: arXiv:2202.03772 Table 6 (exp)
\newcommand{\LitPartJcTestMillion}{20}  % experiments/FIGS/data/literature_facts.json :: facts[66].value :: sha256 713dba2a5e6e1dd1 :: arXiv:2202.03772 Sec. 3
\newcommand{\LitPartJcTwoMAcc}{0.836}  % experiments/FIGS/data/literature_facts.json :: facts[67].value :: sha256 713dba2a5e6e1dd1 :: arXiv:2202.03772 Table 3
\newcommand{\LitPartSmallSetIters}{100 k}  % experiments/FIGS/data/literature_facts.json :: facts[68].value :: sha256 713dba2a5e6e1dd1 :: arXiv:2202.03772 Sec. 5.1 (training dataset size)
\newcommand{\LitPartSmallSetPercent}{2}  % experiments/FIGS/data/literature_facts.json :: facts[69].value :: sha256 713dba2a5e6e1dd1 :: arXiv:2202.03772 Sec. 5.1 (training dataset size)
\newcommand{\LitSophonEtaMax}{2.5}  % experiments/FIGS/data/literature_facts.json :: facts[70].value :: sha256 713dba2a5e6e1dd1 :: arXiv:2405.12972 App. A
\newcommand{\LitBHadronLifetimePs}{1.5}  % experiments/FIGS/data/literature_facts.json :: facts[71].value :: sha256 713dba2a5e6e1dd1 :: arXiv:1712.07158 Sec. 4.1
\newcommand{\LitCHadronLifetimePs}{1}  % experiments/FIGS/data/literature_facts.json :: facts[72].value :: sha256 713dba2a5e6e1dd1 :: arXiv:1712.07158 Sec. 4.1
\newcommand{\LitOmnilearnPidDropProb}{0.2}  % experiments/FIGS/data/literature_facts.json :: facts[73].value :: sha256 713dba2a5e6e1dd1 :: arXiv:2502.14652 Sec. II.A
\newcommand{\LitMahajanFtNineKLarge}{48.5}  % experiments/FIGS/data/literature_facts.json :: facts[74].value :: sha256 713dba2a5e6e1dd1 :: arXiv:1805.00932 Fig. 1 (IG-1B-17k, IN-9k, full fine-tuning)
\newcommand{\LitMahajanFtNineKSmall}{48.4}  % experiments/FIGS/data/literature_facts.json :: facts[75].value :: sha256 713dba2a5e6e1dd1 :: arXiv:1805.00932 Fig. 1 (IG-940M-1.5k, IN-9k, full fine-tuning)
\newcommand{\LitMahajanFtOneKSmall}{84.2}  % experiments/FIGS/data/literature_facts.json :: facts[76].value :: sha256 713dba2a5e6e1dd1 :: arXiv:1805.00932 Fig. 1 (IG-940M-1.5k, IN-1k, full fine-tuning)
\newcommand{\LitMahajanFtOneKLarge}{83.6}  % experiments/FIGS/data/literature_facts.json :: facts[77].value :: sha256 713dba2a5e6e1dd1 :: arXiv:1805.00932 Fig. 1 (IG-1B-17k, IN-1k, full fine-tuning)
\newcommand{\LitChenGainFifteen}{0.16}  % experiments/FIGS/data/literature_facts.json :: facts[78].value :: sha256 713dba2a5e6e1dd1 :: arXiv:1901.07012 Table VIII (15 coarse classes, CIFAR-100)
\newcommand{\LitChenGainTwenty}{0.01}  % experiments/FIGS/data/literature_facts.json :: facts[79].value :: sha256 713dba2a5e6e1dd1 :: arXiv:1901.07012 Table VIII (20 coarse classes, CIFAR-100)
\newcommand{\LitChenCoarseFifteen}{15}  % experiments/FIGS/data/literature_facts.json :: facts[80].value :: sha256 713dba2a5e6e1dd1 :: arXiv:1901.07012 Table VIII
\newcommand{\LitChenCoarseTwenty}{20}  % experiments/FIGS/data/literature_facts.json :: facts[81].value :: sha256 713dba2a5e6e1dd1 :: arXiv:1901.07012 Table VIII
\newcommand{\LitChenRandomMaxGain}{2.78}  % experiments/FIGS/data/literature_facts.json :: facts[82].value :: sha256 713dba2a5e6e1dd1 :: arXiv:1901.07012 Sec. VI-A, Table VI
\newcommand{\LitGlopartMassLambda}{5}  % experiments/FIGS/data/literature_facts.json :: facts[83].value :: sha256 713dba2a5e6e1dd1 :: CERN-THESIS-2024-281 Eq. (11.7), Sec. 11.2.1
\newcommand{\LitGlopartAddedAucBase}{0.9555}  % experiments/FIGS/data/literature_facts.json :: facts[84].value :: sha256 713dba2a5e6e1dd1 :: CERN-THESIS-2024-281 Sec. 11.4.1, Fig. 11.19 (H->WW 4q vs QCD)
\newcommand{\LitGlopartAddedAucTop}{0.9593}  % experiments/FIGS/data/literature_facts.json :: facts[85].value :: sha256 713dba2a5e6e1dd1 :: CERN-THESIS-2024-281 Sec. 11.4.1, Fig. 11.19 (H->WW 4q vs QCD)
\newcommand{\LitCmsSdZcut}{0.1}  % experiments/FIGS/data/literature_facts.json :: facts[86].value :: sha256 713dba2a5e6e1dd1 :: arXiv:2004.08262 Sec. 6.1
\newcommand{\LitCmsSdBeta}{0}  % experiments/FIGS/data/literature_facts.json :: facts[87].value :: sha256 713dba2a5e6e1dd1 :: arXiv:2004.08262 Sec. 6.1
\newcommand{\LitQgNJetsMillion}{2}  % experiments/FIGS/data/literature_facts.json :: facts[88].value :: sha256 713dba2a5e6e1dd1 :: arXiv:2202.03772 Sec. 5.2
\newcommand{\LitQgTrainMillion}{1.6}  % experiments/FIGS/data/literature_facts.json :: facts[89].value :: sha256 713dba2a5e6e1dd1 :: arXiv:2202.03772 Sec. 5.2
\newcommand{\LitQgValMillion}{0.2}  % experiments/FIGS/data/literature_facts.json :: facts[90].value :: sha256 713dba2a5e6e1dd1 :: arXiv:2202.03772 Sec. 5.2
\newcommand{\LitQgRapidityMax}{1.7}  % experiments/FIGS/data/literature_facts.json :: facts[91].value :: sha256 713dba2a5e6e1dd1 :: Zenodo 10.5281/zenodo.3164691 Zenodo record 3164691, description
\newcommand{\LitHerwigVersion}{7.1.4}  % experiments/FIGS/data/literature_facts.json :: facts[92].value :: sha256 713dba2a5e6e1dd1 :: Zenodo 10.5281/zenodo.3066475 record description
\newcommand{\LitObjectivesFewestJets}{\ensuremath{10^{3}}}  % experiments/FIGS/data/literature_facts.json :: facts[93].value :: sha256 713dba2a5e6e1dd1 :: arXiv:2606.14870 Table XVII (Micro model, Npt = 10^8: Classifier 136 vs MPM 218 at 10^3, 331 vs 539 at 10^4)
\newcommand{\LitPartTopNTrainMillion}{1.2}  % experiments/FIGS/data/literature_facts.json :: facts[94].value :: sha256 713dba2a5e6e1dd1 :: arXiv:2202.03772 Sec. 5.2
\newcommand{\LitLandscapeParticlenetEnsMean}{1412}  % experiments/FIGS/data/literature_facts.json :: facts[95].value :: sha256 713dba2a5e6e1dd1 :: arXiv:1902.09914 Table 1 (ParticleNet, ensemble mean of 1/eps_B at eps_S=0.3)
\newcommand{\LitPartFtHeadReinits}{9}  % experiments/FIGS/data/literature_facts.json :: facts[96].value :: sha256 713dba2a5e6e1dd1 :: arXiv:2202.03772 Sec. 5.2
\newcommand{\LitSophonBlocksCode}{8}  % experiments/FIGS/data/literature_facts.json :: facts[97].value :: sha256 713dba2a5e6e1dd1 :: github.com/jet-universe/sophon@9dd6dd6 networks/example_ParticleTransformer_sophon.py line 67
\newcommand{\LitGrasHerwigVersion}{2.7.1}  % experiments/FIGS/data/literature_facts.json :: facts[98].value :: sha256 713dba2a5e6e1dd1 :: arXiv:1704.03878 Sec. 2 (list of generators)
\newcommand{\DesignMassLambda}{5.0}  % experiments/MTX/k8s/job-mtx-l162_mass-s1-raunav.yaml :: --mass-lambda :: sha256 8882459fe5520795 :: identical in all ten mass-output pretraining specs
\newcommand{\FtRecipeNRuns}{264}  % experiments/FIGS/data/ft_recipes/recipes_w2b_bench_v2_v3.json :: runs (count), reported cells :: sha256 7dc8c791343e1b6f
\newcommand{\FtRecipeLr}{\ensuremath{1\times10^{-4}}}  % experiments/FIGS/data/ft_recipes/recipes_w2b_bench_v2_v3.json :: runs[*].lr, pretrained starts :: sha256 7dc8c791343e1b6f
\newcommand{\FtRecipeHeadMult}{50}  % experiments/FIGS/data/ft_recipes/recipes_w2b_bench_v2_v3.json :: runs[*].head_lr_mult, pretrained starts :: sha256 7dc8c791343e1b6f
\newcommand{\FtRecipeLrHead}{\ensuremath{5\times10^{-3}}}  % experiments/FIGS/data/ft_recipes/recipes_w2b_bench_v2_v3.json :: lr x head_lr_mult :: sha256 7dc8c791343e1b6f
\newcommand{\FtRecipeWeightDecay}{0.01}  % experiments/FIGS/data/ft_recipes/recipes_w2b_bench_v2_v3.json :: runs[*].weight_decay :: sha256 7dc8c791343e1b6f
\newcommand{\FtRecipeBatch}{512}  % experiments/FIGS/data/ft_recipes/recipes_w2b_bench_v2_v3.json :: runs[*].batch_size :: sha256 7dc8c791343e1b6f
\newcommand{\FtRecipeEpochsEThree}{50}  % experiments/FIGS/data/ft_recipes/recipes_w2b_bench_v2_v3.json :: runs[leg 1,2; n_train 1000].epochs :: sha256 7dc8c791343e1b6f
\newcommand{\FtRecipeStepsEThree}{19}  % experiments/FT/k8s/job-ft-legs-w2-raunav.yaml :: job commands: samples per epoch at n_train 1000 // batch size :: sha256 f0962eddb451b7b0 :: optimizer steps per epoch, as weaver floors them
\newcommand{\FtRecipePassesEThree}{10}  % experiments/FT/k8s/job-ft-legs-w2-raunav.yaml :: job commands: samples per epoch at n_train 1000 // n_train :: sha256 f0962eddb451b7b0 :: passes over the training subset per epoch
\newcommand{\FtRecipeEpochsEFour}{50}  % experiments/FIGS/data/ft_recipes/recipes_w2b_bench_v2_v3.json :: runs[leg 1,2; n_train 10000].epochs :: sha256 7dc8c791343e1b6f
\newcommand{\FtRecipeStepsEFour}{19}  % experiments/FT/k8s/job-ft-legs-w2-raunav.yaml :: job commands: samples per epoch at n_train 10000 // batch size :: sha256 f0962eddb451b7b0 :: optimizer steps per epoch, as weaver floors them
\newcommand{\FtRecipePassesEFour}{one}  % experiments/FT/k8s/job-ft-legs-w2-raunav.yaml :: job commands: samples per epoch at n_train 10000 // n_train :: sha256 f0962eddb451b7b0 :: passes over the training subset per epoch
\newcommand{\FtRecipeEpochsEFive}{30}  % experiments/FIGS/data/ft_recipes/recipes_w2b_bench_v2_v3.json :: runs[leg 1,2; n_train 100000].epochs :: sha256 7dc8c791343e1b6f
\newcommand{\FtRecipeStepsEFive}{195}  % experiments/FT/k8s/job-ft-legs-w2-raunav.yaml :: job commands: samples per epoch at n_train 100000 // batch size :: sha256 f0962eddb451b7b0 :: optimizer steps per epoch, as weaver floors them
\newcommand{\FtRecipePassesEFive}{one}  % experiments/FT/k8s/job-ft-legs-w2-raunav.yaml :: job commands: samples per epoch at n_train 100000 // n_train :: sha256 f0962eddb451b7b0 :: passes over the training subset per epoch
\newcommand{\FtRecipeEpochsESix}{10}  % experiments/FIGS/data/ft_recipes/recipes_w2b_bench_v2_v3.json :: runs[leg 1,2; n_train 1000000].epochs :: sha256 7dc8c791343e1b6f
\newcommand{\FtRecipeStepsESix}{1{,}953}  % experiments/FT/k8s/job-ft-legs-w2-raunav.yaml :: job commands: samples per epoch at n_train 1000000 // batch size :: sha256 f0962eddb451b7b0 :: optimizer steps per epoch, as weaver floors them
\newcommand{\FtRecipePassesESix}{one}  % experiments/FT/k8s/job-ft-legs-w2-raunav.yaml :: job commands: samples per epoch at n_train 1000000 // n_train :: sha256 f0962eddb451b7b0 :: passes over the training subset per epoch
\newcommand{\FtRecipeValJets}{19{,}968}  % experiments/FT/k8s/job-ft-legs-w2-raunav.yaml :: --samples-per-epoch-val 20000 // batch x batch :: sha256 f0962eddb451b7b0 :: validation jets per epoch, floored to whole batches
\newcommand{\RecoveryAccOneeighteightByOneeighteight}{0.4123\,\ensuremath{\pm}\,0.0011}  % experiments/FIGS/data/label_recovery_curve_v1err/analysis/label_recovery_curve_summary.json :: table[rung=L188,level=188,probe=linear,n_train=max].accuracy :: sha256 1250d1f00537e0bd :: balanced accuracy, mean +- SD over 5 runs
\newcommand{\RecoveryAccOneeighteightByOnesixtwo}{0.4090\,\ensuremath{\pm}\,0.0012}  % experiments/FIGS/data/label_recovery_curve_v1err/analysis/label_recovery_curve_summary.json :: table[rung=L188,level=162,probe=linear,n_train=max].accuracy :: sha256 1250d1f00537e0bd :: balanced accuracy, mean +- SD over 5 runs
\newcommand{\RecoveryAccOneeighteightByFourthree}{0.356\,\ensuremath{\pm}\,0.008}  % experiments/FIGS/data/label_recovery_curve_v1err/analysis/label_recovery_curve_summary.json :: table[rung=L188,level=43,probe=linear,n_train=max].accuracy :: sha256 1250d1f00537e0bd :: balanced accuracy, mean +- SD over 5 runs
\newcommand{\RecoveryAccOneeighteightByOneseven}{0.306\,\ensuremath{\pm}\,0.008}  % experiments/FIGS/data/label_recovery_curve_v1err/analysis/label_recovery_curve_summary.json :: table[rung=L188,level=17,probe=linear,n_train=max].accuracy :: sha256 1250d1f00537e0bd :: balanced accuracy, mean +- SD over 5 runs
\newcommand{\RecoveryAccOnesixtwoByOneeighteight}{0.4623\,\ensuremath{\pm}\,0.0006}  % experiments/FIGS/data/label_recovery_curve_v1err/analysis/label_recovery_curve_summary.json :: table[rung=L162,level=188,probe=linear,n_train=max].accuracy :: sha256 1250d1f00537e0bd :: balanced accuracy, mean +- SD over 5 runs
\newcommand{\RecoveryAccOnesixtwoByOnesixtwo}{0.4622\,\ensuremath{\pm}\,0.0007}  % experiments/FIGS/data/label_recovery_curve_v1err/analysis/label_recovery_curve_summary.json :: table[rung=L162,level=162,probe=linear,n_train=max].accuracy :: sha256 1250d1f00537e0bd :: balanced accuracy, mean +- SD over 5 runs
\newcommand{\RecoveryAccOnesixtwoByFourthree}{0.402\,\ensuremath{\pm}\,0.009}  % experiments/FIGS/data/label_recovery_curve_v1err/analysis/label_recovery_curve_summary.json :: table[rung=L162,level=43,probe=linear,n_train=max].accuracy :: sha256 1250d1f00537e0bd :: balanced accuracy, mean +- SD over 5 runs
\newcommand{\RecoveryAccOnesixtwoByOneseven}{0.348\,\ensuremath{\pm}\,0.010}  % experiments/FIGS/data/label_recovery_curve_v1err/analysis/label_recovery_curve_summary.json :: table[rung=L162,level=17,probe=linear,n_train=max].accuracy :: sha256 1250d1f00537e0bd :: balanced accuracy, mean +- SD over 5 runs
\newcommand{\RecoveryAccSixfourByOneeighteight}{0.5979\,\ensuremath{\pm}\,0.0008}  % experiments/FIGS/data/label_recovery_curve_v1err/analysis/label_recovery_curve_summary.json :: table[rung=R63_Q1,level=188,probe=linear,n_train=max].accuracy :: sha256 1250d1f00537e0bd :: balanced accuracy, mean +- SD over 5 runs
\newcommand{\RecoveryAccSixfourByOnesixtwo}{0.5973\,\ensuremath{\pm}\,0.0013}  % experiments/FIGS/data/label_recovery_curve_v1err/analysis/label_recovery_curve_summary.json :: table[rung=R63_Q1,level=162,probe=linear,n_train=max].accuracy :: sha256 1250d1f00537e0bd :: balanced accuracy, mean +- SD over 5 runs
\newcommand{\RecoveryAccSixfourByFourthree}{0.5881\,\ensuremath{\pm}\,0.0007}  % experiments/FIGS/data/label_recovery_curve_v1err/analysis/label_recovery_curve_summary.json :: table[rung=R63_Q1,level=43,probe=linear,n_train=max].accuracy :: sha256 1250d1f00537e0bd :: balanced accuracy, mean +- SD over 5 runs
\newcommand{\RecoveryAccSixfourByOneseven}{0.531\,\ensuremath{\pm}\,0.006}  % experiments/FIGS/data/label_recovery_curve_v1err/analysis/label_recovery_curve_summary.json :: table[rung=R63_Q1,level=17,probe=linear,n_train=max].accuracy :: sha256 1250d1f00537e0bd :: balanced accuracy, mean +- SD over 5 runs
\newcommand{\RecoveryAccFourthreeByOneeighteight}{0.6770\,\ensuremath{\pm}\,0.0006}  % experiments/FIGS/data/label_recovery_curve_v1err/analysis/label_recovery_curve_summary.json :: table[rung=R42_Q1,level=188,probe=linear,n_train=max].accuracy :: sha256 1250d1f00537e0bd :: balanced accuracy, mean +- SD over 5 runs
\newcommand{\RecoveryAccFourthreeByOnesixtwo}{0.6766\,\ensuremath{\pm}\,0.0008}  % experiments/FIGS/data/label_recovery_curve_v1err/analysis/label_recovery_curve_summary.json :: table[rung=R42_Q1,level=162,probe=linear,n_train=max].accuracy :: sha256 1250d1f00537e0bd :: balanced accuracy, mean +- SD over 5 runs
\newcommand{\RecoveryAccFourthreeByFourthree}{0.6751\,\ensuremath{\pm}\,0.0011}  % experiments/FIGS/data/label_recovery_curve_v1err/analysis/label_recovery_curve_summary.json :: table[rung=R42_Q1,level=43,probe=linear,n_train=max].accuracy :: sha256 1250d1f00537e0bd :: balanced accuracy, mean +- SD over 5 runs
\newcommand{\RecoveryAccFourthreeByOneseven}{0.633\,\ensuremath{\pm}\,0.006}  % experiments/FIGS/data/label_recovery_curve_v1err/analysis/label_recovery_curve_summary.json :: table[rung=R42_Q1,level=17,probe=linear,n_train=max].accuracy :: sha256 1250d1f00537e0bd :: balanced accuracy, mean +- SD over 5 runs
\newcommand{\RecoveryAccThreezeroByOneeighteight}{0.7001\,\ensuremath{\pm}\,0.0005}  % experiments/FIGS/data/label_recovery_curve_v1err/analysis/label_recovery_curve_summary.json :: table[rung=R29_Q1,level=188,probe=linear,n_train=max].accuracy :: sha256 1250d1f00537e0bd :: balanced accuracy, mean +- SD over 5 runs
\newcommand{\RecoveryAccThreezeroByOnesixtwo}{0.7004\,\ensuremath{\pm}\,0.0007}  % experiments/FIGS/data/label_recovery_curve_v1err/analysis/label_recovery_curve_summary.json :: table[rung=R29_Q1,level=162,probe=linear,n_train=max].accuracy :: sha256 1250d1f00537e0bd :: balanced accuracy, mean +- SD over 5 runs
\newcommand{\RecoveryAccThreezeroByFourthree}{0.6998\,\ensuremath{\pm}\,0.0008}  % experiments/FIGS/data/label_recovery_curve_v1err/analysis/label_recovery_curve_summary.json :: table[rung=R29_Q1,level=43,probe=linear,n_train=max].accuracy :: sha256 1250d1f00537e0bd :: balanced accuracy, mean +- SD over 5 runs
\newcommand{\RecoveryAccThreezeroByOneseven}{0.6778\,\ensuremath{\pm}\,0.0030}  % experiments/FIGS/data/label_recovery_curve_v1err/analysis/label_recovery_curve_summary.json :: table[rung=R29_Q1,level=17,probe=linear,n_train=max].accuracy :: sha256 1250d1f00537e0bd :: balanced accuracy, mean +- SD over 5 runs
\newcommand{\RecoveryAccOnesevenByOneeighteight}{0.7675\,\ensuremath{\pm}\,0.0010}  % experiments/FIGS/data/label_recovery_curve_v1err/analysis/label_recovery_curve_summary.json :: table[rung=R16_Q1,level=188,probe=linear,n_train=max].accuracy :: sha256 1250d1f00537e0bd :: balanced accuracy, mean +- SD over 5 runs
\newcommand{\RecoveryAccOnesevenByOnesixtwo}{0.7681\,\ensuremath{\pm}\,0.0009}  % experiments/FIGS/data/label_recovery_curve_v1err/analysis/label_recovery_curve_summary.json :: table[rung=R16_Q1,level=162,probe=linear,n_train=max].accuracy :: sha256 1250d1f00537e0bd :: balanced accuracy, mean +- SD over 5 runs
\newcommand{\RecoveryAccOnesevenByFourthree}{0.7680\,\ensuremath{\pm}\,0.0010}  % experiments/FIGS/data/label_recovery_curve_v1err/analysis/label_recovery_curve_summary.json :: table[rung=R16_Q1,level=43,probe=linear,n_train=max].accuracy :: sha256 1250d1f00537e0bd :: balanced accuracy, mean +- SD over 5 runs
\newcommand{\RecoveryAccOnesevenByOneseven}{0.7681\,\ensuremath{\pm}\,0.0003}  % experiments/FIGS/data/label_recovery_curve_v1err/analysis/label_recovery_curve_summary.json :: table[rung=R16_Q1,level=17,probe=linear,n_train=max].accuracy :: sha256 1250d1f00537e0bd :: balanced accuracy, mean +- SD over 5 runs
\newcommand{\RecoveryAccFourByOneeighteight}{0.7526\,\ensuremath{\pm}\,0.0002}  % experiments/FIGS/data/label_recovery_curve_v1err/analysis/label_recovery_curve_summary.json :: table[rung=R3_VIS,level=188,probe=linear,n_train=max].accuracy :: sha256 1250d1f00537e0bd :: balanced accuracy, mean +- SD over 5 runs
\newcommand{\RecoveryAccFourByOnesixtwo}{0.7538\,\ensuremath{\pm}\,0.0017}  % experiments/FIGS/data/label_recovery_curve_v1err/analysis/label_recovery_curve_summary.json :: table[rung=R3_VIS,level=162,probe=linear,n_train=max].accuracy :: sha256 1250d1f00537e0bd :: balanced accuracy, mean +- SD over 5 runs
\newcommand{\RecoveryAccFourByFourthree}{0.7557\,\ensuremath{\pm}\,0.0015}  % experiments/FIGS/data/label_recovery_curve_v1err/analysis/label_recovery_curve_summary.json :: table[rung=R3_VIS,level=43,probe=linear,n_train=max].accuracy :: sha256 1250d1f00537e0bd :: balanced accuracy, mean +- SD over 5 runs
\newcommand{\RecoveryAccFourByOneseven}{0.7586\,\ensuremath{\pm}\,0.0008}  % experiments/FIGS/data/label_recovery_curve_v1err/analysis/label_recovery_curve_summary.json :: table[rung=R3_VIS,level=17,probe=linear,n_train=max].accuracy :: sha256 1250d1f00537e0bd :: balanced accuracy, mean +- SD over 5 runs
\newcommand{\RecoveryAccTwoByOneeighteight}{0.8688\,\ensuremath{\pm}\,0.0005}  % experiments/FIGS/data/label_recovery_curve_v1err/analysis/label_recovery_curve_summary.json :: table[rung=R1_Q1,level=188,probe=linear,n_train=max].accuracy :: sha256 1250d1f00537e0bd :: balanced accuracy, mean +- SD over 5 runs
\newcommand{\RecoveryAccTwoByOnesixtwo}{0.8703\,\ensuremath{\pm}\,0.0008}  % experiments/FIGS/data/label_recovery_curve_v1err/analysis/label_recovery_curve_summary.json :: table[rung=R1_Q1,level=162,probe=linear,n_train=max].accuracy :: sha256 1250d1f00537e0bd :: balanced accuracy, mean +- SD over 5 runs
\newcommand{\RecoveryAccTwoByFourthree}{0.8709\,\ensuremath{\pm}\,0.0011}  % experiments/FIGS/data/label_recovery_curve_v1err/analysis/label_recovery_curve_summary.json :: table[rung=R1_Q1,level=43,probe=linear,n_train=max].accuracy :: sha256 1250d1f00537e0bd :: balanced accuracy, mean +- SD over 5 runs
\newcommand{\RecoveryAccTwoByOneseven}{0.8713\,\ensuremath{\pm}\,0.0002}  % experiments/FIGS/data/label_recovery_curve_v1err/analysis/label_recovery_curve_summary.json :: table[rung=R1_Q1,level=17,probe=linear,n_train=max].accuracy :: sha256 1250d1f00537e0bd :: balanced accuracy, mean +- SD over 5 runs
\newcommand{\RecoveryOneeighteightOnesixtwoMaxAbs}{0.0033}  % experiments/FIGS/data/label_recovery_curve_v1err/analysis/label_recovery_curve_summary.json :: table[rung=*,level=188/162,probe=linear,n_train=max].accuracy :: sha256 1250d1f00537e0bd :: max |difference of the run means| over the tree
\newcommand{\RecoveryPairedOnesevenLevelOnesevenMinusOneeighteight}{+0.0006~[\ensuremath{-}0.0004, +0.0015]}  % experiments/FIGS/data/label_recovery_curve_v1err/analysis/label_recovery_curve_summary.json :: table[rung=R16_Q1,level=17 minus 188,probe=linear,n_train=max].accuracy :: sha256 1250d1f00537e0bd :: paired by run, mean difference [95% Student-t interval]
\newcommand{\RecoveryPairedOnesevenLevelOnesevenMinusOnesixtwo}{+0.00004~[\ensuremath{-}0.00123, +0.00131]}  % experiments/FIGS/data/label_recovery_curve_v1err/analysis/label_recovery_curve_summary.json :: table[rung=R16_Q1,level=17 minus 162,probe=linear,n_train=max].accuracy :: sha256 1250d1f00537e0bd :: paired by run, mean difference [95% Student-t interval]
\newcommand{\RecoveryPairedFourLevelOnesevenMinusOneeighteight}{+0.0060~[+0.0051, +0.0069]}  % experiments/FIGS/data/label_recovery_curve_v1err/analysis/label_recovery_curve_summary.json :: table[rung=R3_VIS,level=17 minus 188,probe=linear,n_train=max].accuracy :: sha256 1250d1f00537e0bd :: paired by run, mean difference [95% Student-t interval]
\newcommand{\RecoveryPairedFourLevelOnesevenMinusOnesixtwo}{+0.0048~[+0.0028, +0.0068]}  % experiments/FIGS/data/label_recovery_curve_v1err/analysis/label_recovery_curve_summary.json :: table[rung=R3_VIS,level=17 minus 162,probe=linear,n_train=max].accuracy :: sha256 1250d1f00537e0bd :: paired by run, mean difference [95% Student-t interval]
\newcommand{\RecoveryPairedTwoLevelOnesevenMinusOneeighteight}{+0.0025~[+0.0019, +0.0032]}  % experiments/FIGS/data/label_recovery_curve_v1err/analysis/label_recovery_curve_summary.json :: table[rung=R1_Q1,level=17 minus 188,probe=linear,n_train=max].accuracy :: sha256 1250d1f00537e0bd :: paired by run, mean difference [95% Student-t interval]
\newcommand{\RecoveryPairedTwoLevelOnesevenMinusOnesixtwo}{+0.00101~[+0.00003, +0.00198]}  % experiments/FIGS/data/label_recovery_curve_v1err/analysis/label_recovery_curve_summary.json :: table[rung=R1_Q1,level=17 minus 162,probe=linear,n_train=max].accuracy :: sha256 1250d1f00537e0bd :: paired by run, mean difference [95% Student-t interval]
\newcommand{\RecoveryMlpGainOneseven}{+0.0099~[+0.0078, +0.0120]}  % experiments/FIGS/data/label_recovery_curve_v1err/analysis/label_recovery_curve_summary.json :: mlp_minus_linear_at_largest[rung=L188,level=17].per_run :: sha256 1250d1f00537e0bd :: nonlinear minus linear probe, 188 native labels from the 17-class models
\newcommand{\RecoveryLastGainOneseven}{0.0125\,\ensuremath{\pm}\,0.0021}  % experiments/FIGS/data/label_recovery_curve_v1err/analysis/label_recovery_curve_summary.json :: data_gain_last_step[rung=L188,level=17] :: sha256 1250d1f00537e0bd :: balanced accuracy gained from 443000 to 1400000 training jets
\newcommand{\RecoveryLastGainOneeighteight}{0.0147\,\ensuremath{\pm}\,0.0014}  % experiments/FIGS/data/label_recovery_curve_v1err/analysis/label_recovery_curve_summary.json :: data_gain_last_step[rung=L188,level=188] :: sha256 1250d1f00537e0bd :: balanced accuracy gained from 443000 to 1400000 training jets
\newcommand{\RecoveryNTrain}{1{,}400{,}000}  % experiments/FIGS/data/label_recovery_curve_v1err/analysis/label_recovery_curve_summary.json :: sizes (max) :: sha256 1250d1f00537e0bd :: jets the linear probe is fit on at the end of the curve
\newcommand{\RecoveryNTrainMin}{14{,}000}  % experiments/FIGS/data/label_recovery_curve_v1err/analysis/label_recovery_curve_summary.json :: sizes (min) :: sha256 1250d1f00537e0bd
\newcommand{\RecoveryNTrainPrev}{443{,}000}  % experiments/FIGS/data/label_recovery_curve_v1err/analysis/label_recovery_curve_summary.json :: sizes (second largest) :: sha256 1250d1f00537e0bd
\newcommand{\RecoveryNTest}{400{,}000}  % experiments/FIGS/data/label_recovery_curve_v1err/l188-s1/label_recovery_curve.json :: n_test :: sha256 65651e4c06c7766f :: jets the probes are scored on, recorded in all 20 per-run files
\newcommand{\RecoveryMlpHidden}{512}  % experiments/FIGS/data/label_recovery_curve_v1err/l188-s1/label_recovery_curve.json :: mlp.hidden :: sha256 65651e4c06c7766f :: hidden units of the nonlinear probe
\newcommand{\RandNDraws}{three}  % experiments/FIGS/data/probe_ladder_randcontrol_mlp2/analysis/c4_random_control.json :: table[*].draw :: sha256 c675e609ff4fe3c9 :: random partitions drawn
\newcommand{\RandNDrawsCap}{Three}  % experiments/FIGS/data/probe_ladder_randcontrol_mlp2/analysis/c4_random_control.json :: table[*].draw :: sha256 c675e609ff4fe3c9 :: random partitions drawn, capitalised to open a sentence
\newcommand{\RandOmaBvcFourprongDrawOneLinear}{\ensuremath{1.05\times10^{-2}}}  % experiments/FIGS/data/probe_ladder_randcontrol_mlp2/analysis/c4_random_control.json :: table[0].control_log1m_auc :: sha256 c675e609ff4fe3c9 :: 1-AUC = exp of it, one run
\newcommand{\RandOmaBvcFourprongDrawTwoLinear}{\ensuremath{1.11\times10^{-2}}}  % experiments/FIGS/data/probe_ladder_randcontrol_mlp2/analysis/c4_random_control.json :: table[4].control_log1m_auc :: sha256 c675e609ff4fe3c9 :: 1-AUC = exp of it, one run
\newcommand{\RandOmaBvcFourprongDrawThreeLinear}{\ensuremath{9.98\times10^{-3}}}  % experiments/FIGS/data/probe_ladder_randcontrol_mlp2/analysis/c4_random_control.json :: table[8].control_log1m_auc :: sha256 c675e609ff4fe3c9 :: 1-AUC = exp of it, one run
\newcommand{\RandOmaMeanBvcFourprongLinear}{\ensuremath{(1.05 \pm 0.06)\times10^{-2}}}  % experiments/FIGS/data/probe_ladder_randcontrol_mlp2/analysis/c4_random_control.json :: table[task=bvc_4prong,probe=linear].control_log1m_auc :: sha256 c675e609ff4fe3c9 :: 1-AUC, mean +- SD over 3 draws
\newcommand{\RandOmaBvcFourprongDrawOneMlp}{\ensuremath{1.05\times10^{-2}}}  % experiments/FIGS/data/probe_ladder_randcontrol_mlp2/analysis/c4_random_control.json :: table[1].control_log1m_auc :: sha256 c675e609ff4fe3c9 :: 1-AUC = exp of it, one run
\newcommand{\RandOmaBvcFourprongDrawTwoMlp}{\ensuremath{1.04\times10^{-2}}}  % experiments/FIGS/data/probe_ladder_randcontrol_mlp2/analysis/c4_random_control.json :: table[5].control_log1m_auc :: sha256 c675e609ff4fe3c9 :: 1-AUC = exp of it, one run
\newcommand{\RandOmaBvcFourprongDrawThreeMlp}{\ensuremath{9.28\times10^{-3}}}  % experiments/FIGS/data/probe_ladder_randcontrol_mlp2/analysis/c4_random_control.json :: table[9].control_log1m_auc :: sha256 c675e609ff4fe3c9 :: 1-AUC = exp of it, one run
\newcommand{\RandOmaMeanBvcFourprongMlp}{\ensuremath{(1.01 \pm 0.07)\times10^{-2}}}  % experiments/FIGS/data/probe_ladder_randcontrol_mlp2/analysis/c4_random_control.json :: table[task=bvc_4prong,probe=mlp].control_log1m_auc :: sha256 c675e609ff4fe3c9 :: 1-AUC, mean +- SD over 3 draws
\newcommand{\RandOmaVisibleContentDrawOneLinear}{\ensuremath{1.25\times10^{-3}}}  % experiments/FIGS/data/probe_ladder_randcontrol_mlp2/analysis/c4_random_control.json :: table[2].control_log1m_auc :: sha256 c675e609ff4fe3c9 :: 1-AUC = exp of it, one run
\newcommand{\RandOmaVisibleContentDrawTwoLinear}{\ensuremath{9.89\times10^{-4}}}  % experiments/FIGS/data/probe_ladder_randcontrol_mlp2/analysis/c4_random_control.json :: table[6].control_log1m_auc :: sha256 c675e609ff4fe3c9 :: 1-AUC = exp of it, one run
\newcommand{\RandOmaVisibleContentDrawThreeLinear}{\ensuremath{1.08\times10^{-3}}}  % experiments/FIGS/data/probe_ladder_randcontrol_mlp2/analysis/c4_random_control.json :: table[10].control_log1m_auc :: sha256 c675e609ff4fe3c9 :: 1-AUC = exp of it, one run
\newcommand{\RandOmaMeanVisibleContentLinear}{\ensuremath{(1.11 \pm 0.13)\times10^{-3}}}  % experiments/FIGS/data/probe_ladder_randcontrol_mlp2/analysis/c4_random_control.json :: table[task=visible_content,probe=linear].control_log1m_auc :: sha256 c675e609ff4fe3c9 :: 1-AUC, mean +- SD over 3 draws
\newcommand{\RandOmaVisibleContentDrawOneMlp}{\ensuremath{1.21\times10^{-3}}}  % experiments/FIGS/data/probe_ladder_randcontrol_mlp2/analysis/c4_random_control.json :: table[3].control_log1m_auc :: sha256 c675e609ff4fe3c9 :: 1-AUC = exp of it, one run
\newcommand{\RandOmaVisibleContentDrawTwoMlp}{\ensuremath{1.07\times10^{-3}}}  % experiments/FIGS/data/probe_ladder_randcontrol_mlp2/analysis/c4_random_control.json :: table[7].control_log1m_auc :: sha256 c675e609ff4fe3c9 :: 1-AUC = exp of it, one run
\newcommand{\RandOmaVisibleContentDrawThreeMlp}{\ensuremath{1.13\times10^{-3}}}  % experiments/FIGS/data/probe_ladder_randcontrol_mlp2/analysis/c4_random_control.json :: table[11].control_log1m_auc :: sha256 c675e609ff4fe3c9 :: 1-AUC = exp of it, one run
\newcommand{\RandOmaMeanVisibleContentMlp}{\ensuremath{(1.14 \pm 0.07)\times10^{-3}}}  % experiments/FIGS/data/probe_ladder_randcontrol_mlp2/analysis/c4_random_control.json :: table[task=visible_content,probe=mlp].control_log1m_auc :: sha256 c675e609ff4fe3c9 :: 1-AUC, mean +- SD over 3 draws
\newcommand{\RandCostFactorBvcFourprongLinear}{1.49}  % experiments/FIGS/data/probe_ladder_randcontrol_mlp2/analysis/c4_random_control.json :: grouping_cost_post_hoc.tasks.bvc_4prong.linear.factor_control :: sha256 c675e609ff4fe3c9 :: POST HOC: random grouping vs finer models, in 1-AUC
\newcommand{\RandSemCostFactorBvcFourprongLinear}{4.42}  % experiments/FIGS/data/probe_ladder_randcontrol_mlp2/analysis/c4_random_control.json :: grouping_cost_post_hoc.tasks.bvc_4prong.linear.factor_semantic :: sha256 c675e609ff4fe3c9 :: POST HOC: semantic 17-class vs finer models
\newcommand{\RandCostFactorBvcFourprongMlp}{1.41}  % experiments/FIGS/data/probe_ladder_randcontrol_mlp2/analysis/c4_random_control.json :: grouping_cost_post_hoc.tasks.bvc_4prong.mlp.factor_control :: sha256 c675e609ff4fe3c9 :: POST HOC: random grouping vs finer models, in 1-AUC
\newcommand{\RandSemCostFactorBvcFourprongMlp}{3.26}  % experiments/FIGS/data/probe_ladder_randcontrol_mlp2/analysis/c4_random_control.json :: grouping_cost_post_hoc.tasks.bvc_4prong.mlp.factor_semantic :: sha256 c675e609ff4fe3c9 :: POST HOC: semantic 17-class vs finer models
\newcommand{\RandCostFactorVisibleContentLinear}{1.35}  % experiments/FIGS/data/probe_ladder_randcontrol_mlp2/analysis/c4_random_control.json :: grouping_cost_post_hoc.tasks.visible_content.linear.factor_control :: sha256 c675e609ff4fe3c9 :: POST HOC: random grouping vs finer models, in 1-AUC
\newcommand{\RandSemCostFactorVisibleContentLinear}{2.82}  % experiments/FIGS/data/probe_ladder_randcontrol_mlp2/analysis/c4_random_control.json :: grouping_cost_post_hoc.tasks.visible_content.linear.factor_semantic :: sha256 c675e609ff4fe3c9 :: POST HOC: semantic 17-class vs finer models
\newcommand{\RandCostFactorVisibleContentMlp}{1.33}  % experiments/FIGS/data/probe_ladder_randcontrol_mlp2/analysis/c4_random_control.json :: grouping_cost_post_hoc.tasks.visible_content.mlp.factor_control :: sha256 c675e609ff4fe3c9 :: POST HOC: random grouping vs finer models, in 1-AUC
\newcommand{\RandSemCostFactorVisibleContentMlp}{2.37}  % experiments/FIGS/data/probe_ladder_randcontrol_mlp2/analysis/c4_random_control.json :: grouping_cost_post_hoc.tasks.visible_content.mlp.factor_semantic :: sha256 c675e609ff4fe3c9 :: POST HOC: semantic 17-class vs finer models
\newcommand{\MassResNTest}{398{,}625}  % experiments/FIGS/data/mass_resolution/s1.json :: centering_detail.split[2] :: sha256 88a57bbad270d51c :: test jets sigma_eff is computed on, the same in every per-seed file and arm
\newcommand{\MassResNClasses}{184}  % experiments/FIGS/data/mass_resolution/analysis_holm/s7_mass_resolution.json :: provenance.n_classes_used :: sha256 e03cde45ece9f385 :: native classes with enough training jets to centre
\newcommand{\MassResSigmaEffRidgeOneeighteight}{0.3393\,\ensuremath{\pm}\,0.0019}  % experiments/FIGS/data/mass_resolution/analysis_holm/s7_mass_resolution.json :: table[cell=188,probe=ridge].sigma_eff :: sha256 e03cde45ece9f385 :: mean +- SD over 5 seeds
\newcommand{\MassResSigmaEffRidgeOnesixtwo}{0.3302\,\ensuremath{\pm}\,0.0029}  % experiments/FIGS/data/mass_resolution/analysis_holm/s7_mass_resolution.json :: table[cell=162,probe=ridge].sigma_eff :: sha256 e03cde45ece9f385 :: mean +- SD over 5 seeds
\newcommand{\MassResSigmaEffRidgeFourthree}{0.3228\,\ensuremath{\pm}\,0.0019}  % experiments/FIGS/data/mass_resolution/analysis_holm/s7_mass_resolution.json :: table[cell=43,probe=ridge].sigma_eff :: sha256 e03cde45ece9f385 :: mean +- SD over 5 seeds
\newcommand{\MassResSigmaEffRidgeOneseven}{0.300\,\ensuremath{\pm}\,0.010}  % experiments/FIGS/data/mass_resolution/analysis_holm/s7_mass_resolution.json :: table[cell=17,probe=ridge].sigma_eff :: sha256 e03cde45ece9f385 :: mean +- SD over 5 seeds
\newcommand{\MassResSigmaEffRidgeOnesixtwoMass}{0.2238\,\ensuremath{\pm}\,0.0019}  % experiments/FIGS/data/mass_resolution/analysis_holm/s7_mass_resolution.json :: table[cell=162+mass,probe=ridge].sigma_eff :: sha256 e03cde45ece9f385 :: mean +- SD over 5 seeds
\newcommand{\MassResSigmaEffRidgeOnesevenMass}{0.2170\,\ensuremath{\pm}\,0.0034}  % experiments/FIGS/data/mass_resolution/analysis_holm/s7_mass_resolution.json :: table[cell=17+mass,probe=ridge].sigma_eff :: sha256 e03cde45ece9f385 :: mean +- SD over 5 seeds
\newcommand{\MassResGainRidgeOnesixtwo}{\ensuremath{-}0.1063}  % experiments/FIGS/data/mass_resolution/analysis_holm/s7_mass_resolution.json :: table[cell=162+mass/162,probe=ridge].sigma_eff :: sha256 e03cde45ece9f385 :: sigma_eff with the mass output minus without, difference of the seed means; negative is better
\newcommand{\MassResGainRidgeOneseven}{\ensuremath{-}0.0831}  % experiments/FIGS/data/mass_resolution/analysis_holm/s7_mass_resolution.json :: table[cell=17+mass/17,probe=ridge].sigma_eff :: sha256 e03cde45ece9f385 :: sigma_eff with the mass output minus without, difference of the seed means; negative is better
\newcommand{\MassResDidRidge}{\ensuremath{-}0.0233}  % experiments/FIGS/data/mass_resolution/analysis_holm/s7_mass_resolution.json :: table[cell=162+mass/162/17+mass/17,probe=ridge].sigma_eff :: sha256 e03cde45ece9f385 :: gain at 162 classes minus gain at 17
\newcommand{\MassResSigmaEffMlpOneeighteight}{0.2463\,\ensuremath{\pm}\,0.0022}  % experiments/FIGS/data/mass_resolution/analysis_holm/s7_mass_resolution.json :: table[cell=188,probe=mlp].sigma_eff :: sha256 e03cde45ece9f385 :: mean +- SD over 5 seeds
\newcommand{\MassResSigmaEffMlpOnesixtwo}{0.2438\,\ensuremath{\pm}\,0.0018}  % experiments/FIGS/data/mass_resolution/analysis_holm/s7_mass_resolution.json :: table[cell=162,probe=mlp].sigma_eff :: sha256 e03cde45ece9f385 :: mean +- SD over 5 seeds
\newcommand{\MassResSigmaEffMlpFourthree}{0.2407\,\ensuremath{\pm}\,0.0009}  % experiments/FIGS/data/mass_resolution/analysis_holm/s7_mass_resolution.json :: table[cell=43,probe=mlp].sigma_eff :: sha256 e03cde45ece9f385 :: mean +- SD over 5 seeds
\newcommand{\MassResSigmaEffMlpOneseven}{0.2358\,\ensuremath{\pm}\,0.0025}  % experiments/FIGS/data/mass_resolution/analysis_holm/s7_mass_resolution.json :: table[cell=17,probe=mlp].sigma_eff :: sha256 e03cde45ece9f385 :: mean +- SD over 5 seeds
\newcommand{\MassResSigmaEffMlpOnesixtwoMass}{0.1943\,\ensuremath{\pm}\,0.0001}  % experiments/FIGS/data/mass_resolution/analysis_holm/s7_mass_resolution.json :: table[cell=162+mass,probe=mlp].sigma_eff :: sha256 e03cde45ece9f385 :: mean +- SD over 5 seeds
\newcommand{\MassResSigmaEffMlpOnesevenMass}{0.1949\,\ensuremath{\pm}\,0.0006}  % experiments/FIGS/data/mass_resolution/analysis_holm/s7_mass_resolution.json :: table[cell=17+mass,probe=mlp].sigma_eff :: sha256 e03cde45ece9f385 :: mean +- SD over 5 seeds
\newcommand{\MassResGainMlpOnesixtwo}{\ensuremath{-}0.0496}  % experiments/FIGS/data/mass_resolution/analysis_holm/s7_mass_resolution.json :: table[cell=162+mass/162,probe=mlp].sigma_eff :: sha256 e03cde45ece9f385 :: sigma_eff with the mass output minus without, difference of the seed means; negative is better
\newcommand{\MassResGainMlpOneseven}{\ensuremath{-}0.0409}  % experiments/FIGS/data/mass_resolution/analysis_holm/s7_mass_resolution.json :: table[cell=17+mass/17,probe=mlp].sigma_eff :: sha256 e03cde45ece9f385 :: sigma_eff with the mass output minus without, difference of the seed means; negative is better
\newcommand{\MassResDidMlp}{\ensuremath{-}0.0087}  % experiments/FIGS/data/mass_resolution/analysis_holm/s7_mass_resolution.json :: table[cell=162+mass/162/17+mass/17,probe=mlp].sigma_eff :: sha256 e03cde45ece9f385 :: gain at 162 classes minus gain at 17
\newcommand{\MassResTargetSigmaEff}{0.2986}  % experiments/FIGS/data/mass_resolution/analysis_holm/s7_mass_resolution.json :: table[*].target_sigma_eff :: sha256 e03cde45ece9f385 :: sigma_eff of the class-centred target itself (a model that predicts the class mean)
\newcommand{\MassResRidgeRtwoMin}{20\%}  % experiments/FIGS/data/mass_resolution/s1.json :: arms.{l188,l162,r42q1,r16q1}-s*.ridge.val_r2 :: sha256 88a57bbad270d51c :: validation R^2 of the ridge probe, smallest of 20 models without the mass output, all 5 per-seed files
\newcommand{\MassResRidgeRtwoMax}{30\%}  % experiments/FIGS/data/mass_resolution/s1.json :: arms.{l188,l162,r42q1,r16q1}-s*.ridge.val_r2 :: sha256 88a57bbad270d51c :: validation R^2 of the ridge probe, largest of 20 models without the mass output, all 5 per-seed files
\newcommand{\MassResSigmaEffOnesixtwo}{0.3302}  % experiments/FIGS/data/mass_resolution/analysis_holm/s7_mass_resolution.json :: table[cell=162,probe=ridge].sigma_eff :: sha256 e03cde45ece9f385 :: mean over 5 seeds
\newcommand{\MassResSdOnesixtwo}{0.626}  % experiments/FIGS/data/mass_resolution/analysis_holm/s7_mass_resolution.json :: table[cell=162,probe=ridge].sd :: sha256 e03cde45ece9f385 :: mean over 5 seeds
\newcommand{\MassResTailFractionOnesixtwo}{0.062}  % experiments/FIGS/data/mass_resolution/analysis_holm/s7_mass_resolution.json :: table[cell=162,probe=ridge].tail_fraction :: sha256 e03cde45ece9f385 :: mean over 5 seeds
\newcommand{\MassResSigmaEffOnesixtwoMass}{0.2238}  % experiments/FIGS/data/mass_resolution/analysis_holm/s7_mass_resolution.json :: table[cell=162+mass,probe=ridge].sigma_eff :: sha256 e03cde45ece9f385 :: mean over 5 seeds
\newcommand{\MassResSdOnesixtwoMass}{0.494}  % experiments/FIGS/data/mass_resolution/analysis_holm/s7_mass_resolution.json :: table[cell=162+mass,probe=ridge].sd :: sha256 e03cde45ece9f385 :: mean over 5 seeds
\newcommand{\MassResTailFractionOnesixtwoMass}{0.034}  % experiments/FIGS/data/mass_resolution/analysis_holm/s7_mass_resolution.json :: table[cell=162+mass,probe=ridge].tail_fraction :: sha256 e03cde45ece9f385 :: mean over 5 seeds
\newcommand{\MassResSigmaEffOneseven}{0.3001}  % experiments/FIGS/data/mass_resolution/analysis_holm/s7_mass_resolution.json :: table[cell=17,probe=ridge].sigma_eff :: sha256 e03cde45ece9f385 :: mean over 5 seeds
\newcommand{\MassResSdOneseven}{0.601}  % experiments/FIGS/data/mass_resolution/analysis_holm/s7_mass_resolution.json :: table[cell=17,probe=ridge].sd :: sha256 e03cde45ece9f385 :: mean over 5 seeds
\newcommand{\MassResTailFractionOneseven}{0.056}  % experiments/FIGS/data/mass_resolution/analysis_holm/s7_mass_resolution.json :: table[cell=17,probe=ridge].tail_fraction :: sha256 e03cde45ece9f385 :: mean over 5 seeds
\newcommand{\MassResSigmaEffOnesevenMass}{0.2170}  % experiments/FIGS/data/mass_resolution/analysis_holm/s7_mass_resolution.json :: table[cell=17+mass,probe=ridge].sigma_eff :: sha256 e03cde45ece9f385 :: mean over 5 seeds
\newcommand{\MassResSdOnesevenMass}{0.490}  % experiments/FIGS/data/mass_resolution/analysis_holm/s7_mass_resolution.json :: table[cell=17+mass,probe=ridge].sd :: sha256 e03cde45ece9f385 :: mean over 5 seeds
\newcommand{\MassResTailFractionOnesevenMass}{0.033}  % experiments/FIGS/data/mass_resolution/analysis_holm/s7_mass_resolution.json :: table[cell=17+mass,probe=ridge].tail_fraction :: sha256 e03cde45ece9f385 :: mean over 5 seeds
\newcommand{\MassResSigmaEffOneeighteight}{0.3393}  % experiments/FIGS/data/mass_resolution/analysis_holm/s7_mass_resolution.json :: table[cell=188,probe=ridge].sigma_eff :: sha256 e03cde45ece9f385 :: mean over 5 seeds
\newcommand{\MassResSdOneeighteight}{0.631}  % experiments/FIGS/data/mass_resolution/analysis_holm/s7_mass_resolution.json :: table[cell=188,probe=ridge].sd :: sha256 e03cde45ece9f385 :: mean over 5 seeds
\newcommand{\MassResTailFractionOneeighteight}{0.064}  % experiments/FIGS/data/mass_resolution/analysis_holm/s7_mass_resolution.json :: table[cell=188,probe=ridge].tail_fraction :: sha256 e03cde45ece9f385 :: mean over 5 seeds
\newcommand{\MassResSigmaEffFourthree}{0.3228}  % experiments/FIGS/data/mass_resolution/analysis_holm/s7_mass_resolution.json :: table[cell=43,probe=ridge].sigma_eff :: sha256 e03cde45ece9f385 :: mean over 5 seeds
\newcommand{\MassResSdFourthree}{0.618}  % experiments/FIGS/data/mass_resolution/analysis_holm/s7_mass_resolution.json :: table[cell=43,probe=ridge].sd :: sha256 e03cde45ece9f385 :: mean over 5 seeds
\newcommand{\MassResTailFractionFourthree}{0.060}  % experiments/FIGS/data/mass_resolution/analysis_holm/s7_mass_resolution.json :: table[cell=43,probe=ridge].tail_fraction :: sha256 e03cde45ece9f385 :: mean over 5 seeds
\newcommand{\AojEffPercent}{1\%}  % experiments/FIGS/data/aoj_full_v1/fit_v6/results.json :: eff :: sha256 2416fc21f93fcfa6 :: the score cut's pass fraction outside the mass windows
\newcommand{\AojPtMin}{500}  % experiments/AOJ/peak_fit.py :: PT_RANGE[0] :: sha256 cad06229a74532be :: GeV
\newcommand{\AojPtMax}{2{,}500}  % experiments/AOJ/peak_fit.py :: PT_RANGE[1] :: sha256 cad06229a74532be :: GeV
\newcommand{\AojRhoMin}{\ensuremath{-}5.5}  % experiments/AOJ/peak_fit.py :: RHO_RANGE[0] :: sha256 cad06229a74532be
\newcommand{\AojRhoMax}{\ensuremath{-}2.0}  % experiments/AOJ/peak_fit.py :: RHO_RANGE[1] :: sha256 cad06229a74532be
\newcommand{\AojWindowWLo}{65}  % experiments/AOJ/peak_fit.py :: PEAKS[W].window[0] :: sha256 cad06229a74532be :: GeV
\newcommand{\AojWindowWHi}{105}  % experiments/AOJ/peak_fit.py :: PEAKS[W].window[1] :: sha256 cad06229a74532be :: GeV
\newcommand{\AojWindowTopLo}{140}  % experiments/AOJ/peak_fit.py :: PEAKS[top].window[0] :: sha256 cad06229a74532be :: GeV
\newcommand{\AojWindowTopHi}{220}  % experiments/AOJ/peak_fit.py :: PEAKS[top].window[1] :: sha256 cad06229a74532be :: GeV
\newcommand{\AojFitLo}{105}  % experiments/AOJ/peak_fit.py :: PEAKS[top].fit_range[0] :: sha256 cad06229a74532be :: GeV
\newcommand{\AojFitHi}{300}  % experiments/AOJ/peak_fit.py :: PEAKS[top].fit_range[1] :: sha256 cad06229a74532be :: GeV
\newcommand{\AojMassBin}{5}  % experiments/AOJ/peak_fit.py :: MASS_BIN :: sha256 cad06229a74532be :: GeV
\newcommand{\AojMapOrderRho}{3}  % experiments/AOJ/peak_fit.py :: MAP_ORDER[0] :: sha256 cad06229a74532be :: polynomial order in rho
\newcommand{\AojMapOrderPt}{2}  % experiments/AOJ/peak_fit.py :: MAP_ORDER[1] :: sha256 cad06229a74532be :: polynomial order in ln pT
\newcommand{\AojNRunsWord}{five}  % experiments/FIGS/data/aoj_full_v1/analysis_v6/aoj_top.json :: per_label_set.label_sets.*.signal_yield.n :: sha256 5dfef902a5fe75ba :: pretraining runs per vocabulary
\newcommand{\AojNJetsRhoWindow}{9{,}441{,}209}  % experiments/FIGS/data/aoj_full_v1/fit_v6/results.json :: n_jets :: sha256 2416fc21f93fcfa6 :: staged jets inside the rho window, all files; not the jets in the mass fit
\newcommand{\AojYieldOneeighteight}{2{,}400\,\ensuremath{\pm}\,500}  % experiments/FIGS/data/aoj_full_v1/analysis_v6/aoj_top.json :: per_label_set.label_sets.188.signal_yield :: sha256 5dfef902a5fe75ba :: mean +- SD over 5 runs, shared shape
\newcommand{\AojStatErrOneeighteight}{211}  % experiments/FIGS/data/aoj_full_v1/analysis_v6/aoj_top.json :: per_label_set.label_sets.188.median_stat_err :: sha256 5dfef902a5fe75ba :: median per-fit statistical error
\newcommand{\AojYieldOwnShapeOneeighteight}{2{,}400\,\ensuremath{\pm}\,600}  % experiments/FIGS/data/aoj_full_v1/analysis_v6/aoj_top.json :: per_label_set.label_sets.188.by_shape.own_shape :: sha256 5dfef902a5fe75ba :: every fit with its own floated shape
\newcommand{\AojSpreadOverStatOneeighteight}{2.4}  % experiments/FIGS/data/aoj_full_v1/analysis_v6/aoj_top.json :: per_label_set.label_sets.188.signal_yield.sd / median_stat_err :: sha256 5dfef902a5fe75ba :: spread between runs over the median per-fit statistical error
\newcommand{\AojYieldOnesixtwo}{2{,}400\,\ensuremath{\pm}\,900}  % experiments/FIGS/data/aoj_full_v1/analysis_v6/aoj_top.json :: per_label_set.label_sets.162.signal_yield :: sha256 5dfef902a5fe75ba :: mean +- SD over 5 runs, shared shape
\newcommand{\AojStatErrOnesixtwo}{199}  % experiments/FIGS/data/aoj_full_v1/analysis_v6/aoj_top.json :: per_label_set.label_sets.162.median_stat_err :: sha256 5dfef902a5fe75ba :: median per-fit statistical error
\newcommand{\AojYieldOwnShapeOnesixtwo}{2{,}200\,\ensuremath{\pm}\,800}  % experiments/FIGS/data/aoj_full_v1/analysis_v6/aoj_top.json :: per_label_set.label_sets.162.by_shape.own_shape :: sha256 5dfef902a5fe75ba :: every fit with its own floated shape
\newcommand{\AojSpreadOverStatOnesixtwo}{4.5}  % experiments/FIGS/data/aoj_full_v1/analysis_v6/aoj_top.json :: per_label_set.label_sets.162.signal_yield.sd / median_stat_err :: sha256 5dfef902a5fe75ba :: spread between runs over the median per-fit statistical error
\newcommand{\AojYieldFourthree}{2{,}400\,\ensuremath{\pm}\,800}  % experiments/FIGS/data/aoj_full_v1/analysis_v6/aoj_top.json :: per_label_set.label_sets.43.signal_yield :: sha256 5dfef902a5fe75ba :: mean +- SD over 5 runs, shared shape
\newcommand{\AojStatErrFourthree}{184}  % experiments/FIGS/data/aoj_full_v1/analysis_v6/aoj_top.json :: per_label_set.label_sets.43.median_stat_err :: sha256 5dfef902a5fe75ba :: median per-fit statistical error
\newcommand{\AojYieldOwnShapeFourthree}{2{,}300\,\ensuremath{\pm}\,800}  % experiments/FIGS/data/aoj_full_v1/analysis_v6/aoj_top.json :: per_label_set.label_sets.43.by_shape.own_shape :: sha256 5dfef902a5fe75ba :: every fit with its own floated shape
\newcommand{\AojSpreadOverStatFourthree}{4.2}  % experiments/FIGS/data/aoj_full_v1/analysis_v6/aoj_top.json :: per_label_set.label_sets.43.signal_yield.sd / median_stat_err :: sha256 5dfef902a5fe75ba :: spread between runs over the median per-fit statistical error
\newcommand{\AojYieldOneseven}{2{,}350\,\ensuremath{\pm}\,240}  % experiments/FIGS/data/aoj_full_v1/analysis_v6/aoj_top.json :: per_label_set.label_sets.17.signal_yield :: sha256 5dfef902a5fe75ba :: mean +- SD over 5 runs, shared shape
\newcommand{\AojStatErrOneseven}{201}  % experiments/FIGS/data/aoj_full_v1/analysis_v6/aoj_top.json :: per_label_set.label_sets.17.median_stat_err :: sha256 5dfef902a5fe75ba :: median per-fit statistical error
\newcommand{\AojYieldOwnShapeOneseven}{2{,}480\,\ensuremath{\pm}\,290}  % experiments/FIGS/data/aoj_full_v1/analysis_v6/aoj_top.json :: per_label_set.label_sets.17.by_shape.own_shape :: sha256 5dfef902a5fe75ba :: every fit with its own floated shape
\newcommand{\AojSpreadOverStatOneseven}{1.2}  % experiments/FIGS/data/aoj_full_v1/analysis_v6/aoj_top.json :: per_label_set.label_sets.17.signal_yield.sd / median_stat_err :: sha256 5dfef902a5fe75ba :: spread between runs over the median per-fit statistical error
\newcommand{\AojYieldOnesixtwoMass}{2{,}420\,\ensuremath{\pm}\,240}  % experiments/FIGS/data/aoj_full_v1/analysis_v6/aoj_top.json :: per_label_set.label_sets.162+mass.signal_yield :: sha256 5dfef902a5fe75ba :: mean +- SD over 5 runs, shared shape
\newcommand{\AojStatErrOnesixtwoMass}{179}  % experiments/FIGS/data/aoj_full_v1/analysis_v6/aoj_top.json :: per_label_set.label_sets.162+mass.median_stat_err :: sha256 5dfef902a5fe75ba :: median per-fit statistical error
\newcommand{\AojYieldOwnShapeOnesixtwoMass}{2{,}350\,\ensuremath{\pm}\,320}  % experiments/FIGS/data/aoj_full_v1/analysis_v6/aoj_top.json :: per_label_set.label_sets.162+mass.by_shape.own_shape :: sha256 5dfef902a5fe75ba :: every fit with its own floated shape
\newcommand{\AojSpreadOverStatOnesixtwoMass}{1.3}  % experiments/FIGS/data/aoj_full_v1/analysis_v6/aoj_top.json :: per_label_set.label_sets.162+mass.signal_yield.sd / median_stat_err :: sha256 5dfef902a5fe75ba :: spread between runs over the median per-fit statistical error
\newcommand{\AojYieldOnesevenMass}{2{,}500\,\ensuremath{\pm}\,600}  % experiments/FIGS/data/aoj_full_v1/analysis_v6/aoj_top.json :: per_label_set.label_sets.17+mass.signal_yield :: sha256 5dfef902a5fe75ba :: mean +- SD over 5 runs, shared shape
\newcommand{\AojStatErrOnesevenMass}{200}  % experiments/FIGS/data/aoj_full_v1/analysis_v6/aoj_top.json :: per_label_set.label_sets.17+mass.median_stat_err :: sha256 5dfef902a5fe75ba :: median per-fit statistical error
\newcommand{\AojYieldOwnShapeOnesevenMass}{2{,}900\,\ensuremath{\pm}\,500}  % experiments/FIGS/data/aoj_full_v1/analysis_v6/aoj_top.json :: per_label_set.label_sets.17+mass.by_shape.own_shape :: sha256 5dfef902a5fe75ba :: every fit with its own floated shape
\newcommand{\AojSpreadOverStatOnesevenMass}{2.9}  % experiments/FIGS/data/aoj_full_v1/analysis_v6/aoj_top.json :: per_label_set.label_sets.17+mass.signal_yield.sd / median_stat_err :: sha256 5dfef902a5fe75ba :: spread between runs over the median per-fit statistical error
\newcommand{\AojOwnShapeShiftMin}{\ensuremath{-}7\%}  % experiments/FIGS/data/aoj_full_v1/analysis_v6/aoj_top.json :: per_label_set.label_sets.*.by_shape.own_shape.mean / signal_yield.mean - 1 :: sha256 5dfef902a5fe75ba :: the vocabulary means with each model's own shape, relative to the shared shape
\newcommand{\AojOwnShapeShiftMax}{+13\%}  % experiments/FIGS/data/aoj_full_v1/analysis_v6/aoj_top.json :: per_label_set.label_sets.*.by_shape.own_shape.mean / signal_yield.mean - 1 :: sha256 5dfef902a5fe75ba :: the vocabulary means with each model's own shape, relative to the shared shape
\newcommand{\AojPooledMean}{182.8}  % experiments/FIGS/data/aoj_full_v1/analysis_v6/aoj_top.json :: per_label_set.shapes.pooled :: sha256 5dfef902a5fe75ba :: GeV
\newcommand{\AojPooledWidth}{11.7}  % experiments/FIGS/data/aoj_full_v1/analysis_v6/aoj_top.json :: per_label_set.shapes.pooled :: sha256 5dfef902a5fe75ba :: GeV
\newcommand{\AojPoolSize}{30}  % experiments/FIGS/data/aoj_full_v1/analysis_v6/aoj_top.json :: per_label_set.shapes.pool (count) :: sha256 5dfef902a5fe75ba :: pretrained models the shared shape is fitted over
\newcommand{\AojNModels}{30}  % experiments/FIGS/data/aoj_full_v1/fit_v6/results.json :: models (pretrained) :: sha256 2416fc21f93fcfa6 :: fits
\newcommand{\AojOwnMeanMin}{181.6}  % experiments/FIGS/data/aoj_full_v1/fit_v6/results.json :: models.*.top.own_shape_fit.mean (min) :: sha256 2416fc21f93fcfa6 :: GeV
\newcommand{\AojOwnMeanMax}{184.0}  % experiments/FIGS/data/aoj_full_v1/fit_v6/results.json :: models.*.top.own_shape_fit.mean (max) :: sha256 2416fc21f93fcfa6 :: GeV
\newcommand{\AojOwnWidthMin}{10.0}  % experiments/FIGS/data/aoj_full_v1/fit_v6/results.json :: models.*.top.own_shape_fit.width (min) :: sha256 2416fc21f93fcfa6 :: GeV
\newcommand{\AojOwnWidthMax}{13.8}  % experiments/FIGS/data/aoj_full_v1/fit_v6/results.json :: models.*.top.own_shape_fit.width (max) :: sha256 2416fc21f93fcfa6 :: GeV
\newcommand{\AojShapeStepMean}{0.7}  % experiments/FIGS/data/aoj_full_v1/fit_v6/results.json :: pooled_shape.systematic_step.mean :: sha256 2416fc21f93fcfa6 :: GeV: the SD of the models' own floated values, one shift for every model together
\newcommand{\AojShapeSystMeanMax}{1\%}  % experiments/FIGS/data/aoj_full_v1/fit_v6/results.json :: models.*.top.shape_systematic.pooled_mean_(down|up) / signal_yield - 1 (largest |.|) :: sha256 2416fc21f93fcfa6 :: the largest yield change of any model under that shift
\newcommand{\AojShapeStepWidth}{1.0}  % experiments/FIGS/data/aoj_full_v1/fit_v6/results.json :: pooled_shape.systematic_step.width :: sha256 2416fc21f93fcfa6 :: GeV: the SD of the models' own floated values, one shift for every model together
\newcommand{\AojShapeSystWidthMax}{12\%}  % experiments/FIGS/data/aoj_full_v1/fit_v6/results.json :: models.*.top.shape_systematic.pooled_width_(down|up) / signal_yield - 1 (largest |.|) :: sha256 2416fc21f93fcfa6 :: the largest yield change of any model under that shift
\newcommand{\AojLeakEpsHi}{0.6}  % experiments/FIGS/data/aoj_full_v1/fit_v6/results.json :: models.*.top.leak_systematic.eps_ref (max) :: sha256 2416fc21f93fcfa6 :: the reference's assumed efficiency for data tops, upper variation
\newcommand{\AojLeakEpsLo}{0.4}  % experiments/FIGS/data/aoj_full_v1/fit_v6/results.json :: models.*.top.leak_systematic.eps_ref (min) :: sha256 2416fc21f93fcfa6 :: the same, lower variation
\newcommand{\AojLeakShiftMin}{+97}  % experiments/FIGS/data/aoj_full_v1/fit_v6/results.json :: models.*.top.leak_systematic.shift['0.6'] (min) :: sha256 2416fc21f93fcfa6 :: tops added to a yield
\newcommand{\AojLeakShiftMax}{+271}  % experiments/FIGS/data/aoj_full_v1/fit_v6/results.json :: models.*.top.leak_systematic.shift['0.4'] (max) :: sha256 2416fc21f93fcfa6 :: tops added to a yield
\newcommand{\AojFitPMin}{0.010}  % experiments/FIGS/data/aoj_full_v1/fit_v6/results.json :: models.*.top.asymptotic_p (min) :: sha256 2416fc21f93fcfa6 :: goodness of fit of the working-point fit, pretrained models
\newcommand{\AojFitPMax}{0.974}  % experiments/FIGS/data/aoj_full_v1/fit_v6/results.json :: models.*.top.asymptotic_p (max) :: sha256 2416fc21f93fcfa6
\newcommand{\AojWeakestYield}{885}  % experiments/FIGS/data/aoj_full_v1/fit_v6/results.json :: models.l162-s2.top.signal_yield :: sha256 2416fc21f93fcfa6 :: the smallest yield of any pretrained model
\newcommand{\AojWeakestYieldErr}{211}  % experiments/FIGS/data/aoj_full_v1/fit_v6/results.json :: models.l162-s2.top.signal_yield_err :: sha256 2416fc21f93fcfa6
\newcommand{\AojWeakestFitP}{0.94}  % experiments/FIGS/data/aoj_full_v1/fit_v6/results.json :: models.l162-s2.top.asymptotic_p :: sha256 2416fc21f93fcfa6
\newcommand{\AojSidebandZMin}{\ensuremath{-}0.43}  % experiments/FIGS/data/aoj_full_v1/fit_v6/histograms.npz :: experiments/AOJ/sideband_residuals.py: z (min over the pretrained models) :: sha256 2bd1e3d5b32d4a12 :: passing jets minus the fit, summed from the fit's lower edge to the top window
\newcommand{\AojSidebandZMax}{+0.58}  % experiments/FIGS/data/aoj_full_v1/fit_v6/histograms.npz :: experiments/AOJ/sideband_residuals.py: z (max over the pretrained models) :: sha256 2bd1e3d5b32d4a12
\newcommand{\AojWeakestSidebandZ}{0.00}  % experiments/FIGS/data/aoj_full_v1/fit_v6/histograms.npz :: experiments/AOJ/sideband_residuals.py: z of l162-s2 :: sha256 2bd1e3d5b32d4a12 :: the smallest-yield model
\newcommand{\AojSidebandMaxPull}{2.7}  % experiments/FIGS/data/aoj_full_v1/fit_v6/histograms.npz :: experiments/AOJ/sideband_residuals.py: max_abs_pull (max over the pretrained models) :: sha256 2bd1e3d5b32d4a12 :: largest single-bin pull below the top window, any pretrained model
\newcommand{\AojSharedBinLo}{125}  % experiments/FIGS/data/aoj_full_v1/fit_v6/histograms.npz :: the bin holding the most models' largest pull :: sha256 2bd1e3d5b32d4a12 :: GeV
\newcommand{\AojSharedBinHi}{130}  % experiments/FIGS/data/aoj_full_v1/fit_v6/histograms.npz :: its upper edge :: sha256 2bd1e3d5b32d4a12 :: GeV
\newcommand{\AojSharedBinNModels}{18}  % experiments/FIGS/data/aoj_full_v1/fit_v6/histograms.npz :: models whose largest pull is in that bin :: sha256 2bd1e3d5b32d4a12
\newcommand{\AojSharedBinMeanPull}{+1.4}  % experiments/FIGS/data/aoj_full_v1/fit_v6/histograms.npz :: experiments/AOJ/sideband_residuals.py mean_pulls: that bin :: sha256 2bd1e3d5b32d4a12 :: mean pull over the pretrained models
\newcommand{\AojSharedBinRank}{third}  % experiments/FIGS/data/aoj_full_v1/fit_v6/histograms.npz :: experiments/AOJ/sideband_residuals.py mean_pulls: rank of that bin's |mean pull| :: sha256 2bd1e3d5b32d4a12
\newcommand{\AojFitNBins}{39}  % experiments/FIGS/data/aoj_full_v1/fit_v6/histograms.npz :: mass bins of the fit range :: sha256 2bd1e3d5b32d4a12
\newcommand{\AojPairedYieldOnesixtwoOverOneeighteight}{0.94~[0.45, 1.95]}  % experiments/FIGS/data/aoj_full_v1/analysis_v6/aoj_top.json :: per_label_set.label_sets.{162,188}.signal_yields :: sha256 5dfef902a5fe75ba :: yield ratio paired by run, geometric mean [95% Student-t interval]
\newcommand{\AojPairedYieldFourthreeOverOneeighteight}{0.97~[0.81, 1.16]}  % experiments/FIGS/data/aoj_full_v1/analysis_v6/aoj_top.json :: per_label_set.label_sets.{43,188}.signal_yields :: sha256 5dfef902a5fe75ba :: yield ratio paired by run, geometric mean [95% Student-t interval]
\newcommand{\AojPairedYieldOnesevenOverOneeighteight}{1.01~[0.75, 1.35]}  % experiments/FIGS/data/aoj_full_v1/analysis_v6/aoj_top.json :: per_label_set.label_sets.{17,188}.signal_yields :: sha256 5dfef902a5fe75ba :: yield ratio paired by run, geometric mean [95% Student-t interval]
\newcommand{\AojPairedYieldOnesevenOverOnesixtwo}{1.07~[0.51, 2.24]}  % experiments/FIGS/data/aoj_full_v1/analysis_v6/aoj_top.json :: per_label_set.label_sets.{17,162}.signal_yields :: sha256 5dfef902a5fe75ba :: yield ratio paired by run, geometric mean [95% Student-t interval]
\newcommand{\AojPairedYieldOnesixtwoMassOverOnesixtwo}{1.10~[0.57, 2.14]}  % experiments/FIGS/data/aoj_full_v1/analysis_v6/aoj_top.json :: per_label_set.label_sets.{162+mass,162}.signal_yields :: sha256 5dfef902a5fe75ba :: yield ratio paired by run, geometric mean [95% Student-t interval]
\newcommand{\AojPairedYieldOnesevenMassOverOneseven}{1.05~[0.74, 1.49]}  % experiments/FIGS/data/aoj_full_v1/analysis_v6/aoj_top.json :: per_label_set.label_sets.{17+mass,17}.signal_yields :: sha256 5dfef902a5fe75ba :: yield ratio paired by run, geometric mean [95% Student-t interval]
\newcommand{\AojLeakPairedShiftMax}{2\%}  % experiments/FIGS/data/aoj_full_v1/fit_v6/results.json :: models.*.top.leak_systematic.signal_yield against signal_yield, paired ratios :: sha256 2416fc21f93fcfa6 :: largest change of a paired yield ratio when the fail-region tops are scaled up
\newcommand{\AojRefEffPercent}{1.38\%}  % experiments/FIGS/data/aoj_full_v1/fit_v6/results.json :: reference.top.data_efficiency :: sha256 2416fc21f93fcfa6 :: fraction of all data jets passing its cut
\newcommand{\AojRefFitP}{\ensuremath{2.56\times10^{-11}}}  % experiments/FIGS/data/aoj_full_v1/fit_v6/results.json :: reference.top.asymptotic_p :: sha256 2416fc21f93fcfa6 :: goodness of fit of the reference's own top fit
\newcommand{\AojRefValidationP}{\ensuremath{p \leq 0.005} (0 of 200 toys reach the observed deviance)}  % experiments/FIGS/data/aoj_full_v1/fit_v6/results.json :: reference.top.validation.toy_p, n_toys :: sha256 2416fc21f93fcfa6 :: toy p-value of the background validation, (k + 1) / (n + 1)
\newcommand{\AojRefBandSignalZ}{2.5}  % experiments/FIGS/data/aoj_full_v1/fit_v6/results.json :: reference.top.validation.band_signal_z :: sha256 2416fc21f93fcfa6 :: significance of a peak of the score's own shape in the validation band
\newcommand{\AojModelEffMin}{1.01\%}  % experiments/FIGS/data/aoj_full_v1/fit_v6/results.json :: models.*.top.data_efficiency (min) :: sha256 2416fc21f93fcfa6 :: pretrained models
\newcommand{\AojModelEffMax}{1.05\%}  % experiments/FIGS/data/aoj_full_v1/fit_v6/results.json :: models.*.top.data_efficiency (max) :: sha256 2416fc21f93fcfa6 :: pretrained models
\newcommand{\AojModelsValidationPMin}{0.11}  % experiments/FIGS/data/aoj_full_v1/fit_v6/results.json :: models.*.top.validation.toy_p (min) :: sha256 2416fc21f93fcfa6 :: pretrained models
\newcommand{\AojRefYield}{29{,}143}  % experiments/FIGS/data/aoj_full_v1/fit_v6/results.json :: reference.top.signal_yield :: sha256 2416fc21f93fcfa6
\newcommand{\AojRefYieldErr}{865}  % experiments/FIGS/data/aoj_full_v1/fit_v6/results.json :: reference.top.signal_yield_err :: sha256 2416fc21f93fcfa6
\newcommand{\AojRefMass}{181.0}  % experiments/FIGS/data/aoj_full_v1/fit_v6/results.json :: reference.top.floated_mean :: sha256 2416fc21f93fcfa6 :: GeV
\newcommand{\AojRefWidth}{17.8}  % experiments/FIGS/data/aoj_full_v1/fit_v6/results.json :: reference.top.floated_width :: sha256 2416fc21f93fcfa6 :: GeV
\newcommand{\AojPublicYield}{1{,}813}  % experiments/FIGS/data/aoj_full_v1/fit_v6/results.json :: models.sophon-public.top.signal_yield :: sha256 2416fc21f93fcfa6
\newcommand{\AojPublicYieldErr}{193}  % experiments/FIGS/data/aoj_full_v1/fit_v6/results.json :: models.sophon-public.top.signal_yield_err :: sha256 2416fc21f93fcfa6
\newcommand{\AojPublicMass}{182.8}  % experiments/FIGS/data/aoj_full_v1/fit_v6/results.json :: models.sophon-public.top.floated_mean :: sha256 2416fc21f93fcfa6 :: GeV
\newcommand{\AojPublicWidth}{11.7}  % experiments/FIGS/data/aoj_full_v1/fit_v6/results.json :: models.sophon-public.top.floated_width :: sha256 2416fc21f93fcfa6 :: GeV
\newcommand{\AojInjSizeMin}{1{,}000}  % experiments/FIGS/data/aoj_injection_v1/summary.json :: groups[region=top,mode=ensemble].size (min) :: sha256 d9b66466ae24b584 :: injected tops
\newcommand{\AojInjSizeMax}{4{,}000}  % experiments/FIGS/data/aoj_injection_v1/summary.json :: groups[region=top,mode=ensemble].size (max) :: sha256 d9b66466ae24b584
\newcommand{\AojInjFloatPullSdMax}{2.0}  % experiments/FIGS/data/aoj_injection_v1/summary.json :: groups[region=top,mode=bootstrap,variant=float].pull.sd (max) :: sha256 d9b66466ae24b584 :: pull width with each model's own floated shape
\newcommand{\AojInjEnsembleRecoveryMin}{0.983}  % experiments/FIGS/data/aoj_injection_v1/summary.json :: groups[region=top,mode=ensemble].ratio (min) :: sha256 d9b66466ae24b584 :: the whole procedure, shape re-derived per toy set
\newcommand{\AojInjEnsembleRecoveryMax}{0.998}  % experiments/FIGS/data/aoj_injection_v1/summary.json :: groups[region=top,mode=ensemble].ratio (max) :: sha256 d9b66466ae24b584
\newcommand{\AojInjEnsemblePullSdMin}{1.02}  % experiments/FIGS/data/aoj_injection_v1/summary.json :: groups[region=top,mode=ensemble].pull (min) :: sha256 d9b66466ae24b584 :: the whole procedure, shape re-derived per toy set
\newcommand{\AojInjEnsemblePullSdMax}{1.09}  % experiments/FIGS/data/aoj_injection_v1/summary.json :: groups[region=top,mode=ensemble].pull (max) :: sha256 d9b66466ae24b584
\newcommand{\AojClosureNShards}{10}  % experiments/FIGS/data/aoj_checks_v2/closure_all_shards.json :: n_shards :: sha256 6a6018351d7048d2 :: staged shards of the open data
\newcommand{\AojClosurePtLo}{500}  % experiments/FIGS/data/aoj_checks_v2/closure_all_shards.json :: pt_window[0] :: sha256 6a6018351d7048d2 :: GeV
\newcommand{\AojClosurePtHi}{700}  % experiments/FIGS/data/aoj_checks_v2/closure_all_shards.json :: pt_window[1] :: sha256 6a6018351d7048d2 :: GeV
\newcommand{\AojClosureDzeroIqr}{8.7}  % experiments/FIGS/data/aoj_checks_v2/closure_all_shards.json :: domain_shift.part_d0.iqr_ratio :: sha256 6a6018351d7048d2 :: open data over JetClass-II QCD, pooled over the shards
\newcommand{\AojClosureDzIqr}{10.6}  % experiments/FIGS/data/aoj_checks_v2/closure_all_shards.json :: domain_shift.part_dz.iqr_ratio :: sha256 6a6018351d7048d2 :: open data over JetClass-II QCD, pooled over the shards
\newcommand{\AojClosureNeutral}{0.65}  % experiments/FIGS/data/aoj_checks_v2/closure_all_shards.json :: features.part_isNeutralHadron.ratio.per_shard (mean) :: sha256 6a6018351d7048d2 :: open data over JetClass-II QCD
\newcommand{\AojClosureCharged}{1.22}  % experiments/FIGS/data/aoj_checks_v2/closure_all_shards.json :: features.part_isChargedHadron.ratio.per_shard (mean) :: sha256 6a6018351d7048d2 :: open data over JetClass-II QCD
\newcommand{\AojClosureNPartAoj}{61}  % experiments/FIGS/data/aoj_checks_v2/closure_all_shards.json :: features.n_particles.aoj.per_shard (mean of the shard medians) :: sha256 6a6018351d7048d2 :: median
\newcommand{\AojClosureNPartRef}{48}  % experiments/FIGS/data/aoj_checks_v2/closure_all_shards.json :: features.n_particles.reference.per_shard (mean) :: sha256 6a6018351d7048d2 :: median
\newcommand{\RandMergingDrawsBvcFourprong}{three}  % configs/labelmaps/rand_label_map.v1.csv :: RAND_d* of label_X_YY_bbqq and label_X_YY_ccqq :: sha256 f395a25f52e3cfad :: first-grid partitions that merge the task's pair
\newcommand{\RandSplittingDrawsBvcFourprong}{one and two}  % configs/labelmaps/rand_label_map.v1.csv :: RAND_d* of label_X_YY_bbqq and label_X_YY_ccqq :: sha256 f395a25f52e3cfad :: first-grid partitions that keep the pair apart
\newcommand{\RandMergingDrawsVisibleContent}{one and three}  % configs/labelmaps/rand_label_map.v1.csv :: RAND_d* of label_X_YY_bbqq and label_X_YY_cqtauhv :: sha256 f395a25f52e3cfad :: first-grid partitions that merge the task's pair
\newcommand{\RandSplittingDrawsVisibleContent}{two}  % configs/labelmaps/rand_label_map.v1.csv :: RAND_d* of label_X_YY_bbqq and label_X_YY_cqtauhv :: sha256 f395a25f52e3cfad :: first-grid partitions that keep the pair apart
\newcommand{\RandVtwoNPartitions}{five}  % configs/arms/v2_grid.json :: arms[RAND2_*] (count) :: sha256 d8a31b1486992302 :: random partitions in the rerun
\newcommand{\RandVtwoRuns}{one}  % configs/arms/v2_paper_runs.txt :: runs of RAND2_* :: sha256 b04562b14ea8cbc7 :: runs per partition
\newcommand{\RandFlavRuns}{two}  % configs/arms/v2_paper_runs.txt :: runs of FLAV_* :: sha256 b04562b14ea8cbc7 :: runs per flavour-pair model
\newcommand{\RandVtwoNPairs}{seven}  % configs/labelmaps/rand_v2_selection.json :: pairs (count) :: sha256 98cf52d49c7d2f02 :: probe pairs the balance rule covers
\newcommand{\RandVtwoMergedMin}{two}  % configs/labelmaps/rand_v2_selection.json :: merged_in :: sha256 98cf52d49c7d2f02 :: partitions merging each probe pair, checked on the map
\newcommand{\RandVtwoMergedMax}{three}  % configs/labelmaps/rand_v2_selection.json :: merged_in :: sha256 98cf52d49c7d2f02 :: partitions merging each probe pair, checked on the map
\newcommand{\RandVtwoShareTol}{5\%}  % configs/labelmaps/rand_v2_selection.json :: realised_tolerance :: sha256 98cf52d49c7d2f02 :: largest relative deviation of a realised group share from its target (rule 3)
\newcommand{\RandVtwoTrials}{32}  % scripts/build_rand_control.py :: TRIALS :: sha256 5bc91d77a3f2e81f :: share-matched assignments drawn per random partition; the one sharing the fewest merged pairs with the 17-class label set is kept
\newcommand{\RandVtwoPoolDraws}{6{,}600}  % configs/labelmaps/rand_v2_selection.json :: pool.draws :: sha256 98cf52d49c7d2f02 :: candidate partitions drawn (seeds) before the realised-share rule
\newcommand{\RandVtwoPoolPass}{27}  % configs/labelmaps/rand_v2_selection.json :: rule3_seeds (count) :: sha256 98cf52d49c7d2f02 :: candidate partitions meeting the realised-share rule (rule 3)
\newcommand{\RandFlavSplitOrbit}{\ensuremath{X\to YY\to QQQQ}}  % configs/labelmaps/flavour_pair.v2.json :: split_orbit :: sha256 3b79ce5e44c1e2f4 :: the decays F1 cuts, Q any quark
\newcommand{\RandFlavSplitOrbitClasses}{22}  % configs/labelmaps/flavour_pair_map.v2.csv :: classes whose orbit is split_orbit (count) :: sha256 9aab235099ceac43 :: classes of the split orbit
\newcommand{\RandFlavCutClasses}{11}  % configs/labelmaps/flavour_pair.v2.json :: split_classes_moved (count) :: sha256 3b79ce5e44c1e2f4 :: classes of the split orbit that F1 moves: those with a b quark
\newcommand{\RandFlavOrbitsMoved}{three}  % configs/labelmaps/flavour_pair.v2.json :: orbits_moved_B_to_A (count) :: sha256 3b79ce5e44c1e2f4 :: whole orbits F1 moves the other way
\newcommand{\RandFlavOrbitsMovedList}{\ensuremath{X\to YY\to QQ\mu\mu}, \ensuremath{X\to YY\to QQ\tau_h\tau_h} and \ensuremath{X\to \mu\mu}}  % configs/labelmaps/flavour_pair.v2.json :: orbits_moved_B_to_A :: sha256 3b79ce5e44c1e2f4 :: the orbits F1 moves the other way, Q any quark
\newcommand{\RandFlavOrbitClassesMoved}{9}  % configs/labelmaps/flavour_pair_map.v2.csv :: classes whose orbit is in orbits_moved_B_to_A :: sha256 9aab235099ceac43 :: classes in those orbits
\newcommand{\RandFlavFonerCutWithB}{6}  % configs/labelmaps/flavour_pair_map.v2.csv :: FLAV_F1R != FLAV_F0 within split_orbit, with a b quark (count) :: sha256 9aab235099ceac43 :: classes with a b quark among those the F1r cut moves
\newcommand{\RandFlavBcPairs}{11}  % configs/labelmaps/flavour_pair.v2.json :: b_c_pairs_in_split_orbit (count) :: sha256 3b79ce5e44c1e2f4 :: pairs of the split orbit's classes that differ only by exchanging b and c
\newcommand{\RandFlavFoneBcSplit}{8}  % configs/labelmaps/flavour_pair_map.v2.csv :: b_c_pairs_in_split_orbit split by FLAV_F1 (count) :: sha256 9aab235099ceac43
\newcommand{\RandFlavFonerBcSplit}{none}  % configs/labelmaps/flavour_pair_map.v2.csv :: b_c_pairs_in_split_orbit split by FLAV_F1R (count) :: sha256 9aab235099ceac43 :: PRESPEC A14: none
\newcommand{\RandFlavFonerShareTol}{1\%}  % configs/labelmaps/flavour_pair.v2.json :: F1R.realised_tolerance_vs_F1 :: sha256 3b79ce5e44c1e2f4 :: largest allowed relative difference of a group's realised share from F1's
\newcommand{\RandFlavFonerShareDev}{0.77\%}  % configs/labelmaps/realised_native_shares.v2.json :: share summed per FLAV_F1R and FLAV_F1 group of flavour_pair_map.v2.csv (largest relative difference) :: sha256 5728c12e88fa3683 :: realised shares from the final loader's dry run
\newcommand{\MlstNRunsWord}{three}  % configs/arms/v2_paper_runs.txt :: runs of L188/L162/R42_Q1/R16_Q1 :: sha256 b04562b14ea8cbc7 :: pretraining runs per configuration
\newcommand{\MlstMassNRunsWord}{two}  % configs/arms/v2_paper_runs.txt :: runs of L162_MASS/R16_Q1_MASS :: sha256 b04562b14ea8cbc7 :: pretraining runs per configuration
\newcommand{\LpCMin}{\ensuremath{1\times10^{-3}}}  % experiments/EVAL/linear_probe_v2.py :: C_GRID :: sha256 4c576b44812c18a2 :: smallest inverse regularisation
\newcommand{\LpCMax}{100}  % experiments/EVAL/linear_probe_v2.py :: C_GRID :: sha256 4c576b44812c18a2 :: largest inverse regularisation
\newcommand{\LpNTestJcii}{2{,}000{,}000}  % experiments/EVAL/linear_probe_v2.py :: TEST_FIRST.jc2 :: sha256 4c576b44812c18a2 :: JetClass-II test jets
\newcommand{\LpNTestAucJcii}{500{,}000}  % experiments/EVAL/linear_probe_v2.py :: TEST_FIRST.jc2 / AUC_STRIDE :: sha256 4c576b44812c18a2 :: JetClass-II test jets of the macro AUC
\newcommand{\LpConvergedGrad}{\ensuremath{1\times10^{-5}}}  % experiments/EVAL/linear_probe_v2.py :: CONVERGED_GRAD :: sha256 4c576b44812c18a2 :: largest gradient component a converged probe ends below
\newcommand{\FillNFeatures}{17}  % configs/finetune/TopReference.yaml :: inputs.pf_features.vars (count) :: sha256 501a6ffe5d0a8051 :: per-particle input features of the network
\newcommand{\FillTopNZero}{10}  % configs/finetune/TopReference.yaml :: inputs.pf_features.vars, part_zero_* (count) :: sha256 501a6ffe5d0a8051 :: features the data set lacks, set to zero
\newcommand{\FillTopNPresent}{7}  % configs/finetune/TopReference.yaml :: inputs.pf_features.vars, not part_zero_* (count) :: sha256 501a6ffe5d0a8051 :: features the data set provides
\newcommand{\FillQgNZero}{4}  % configs/finetune/EnergyFlowQG.yaml :: inputs.pf_features.vars, part_zero_* (count) :: sha256 db42d6471e8b865c :: features the data set lacks, set to zero
\newcommand{\FillQgNPresent}{13}  % configs/finetune/EnergyFlowQG.yaml :: inputs.pf_features.vars, not part_zero_* (count) :: sha256 db42d6471e8b865c :: features the data set provides
\newcommand{\DesignVtwoDataFraction}{20\%}  % experiments/MTX/k8s/v2/grid/job-mtx2-l188-s1-raunav.yaml :: --data-fraction :: sha256 d28ce02267a6b87a :: share of every file read per epoch
\newcommand{\DesignVtwoNValFiles}{25}  % experiments/MTX/k8s/v2/grid/job-mtx2-l188-s1-raunav.yaml :: --data-val (files, brace ranges expanded) :: sha256 d28ce02267a6b87a :: validation files of the fixed sample
\newcommand{\DesignWavgFirstEpoch}{70}  % experiments/MTX/pretrain_v2.py :: --num-epochs - WAVG_EPOCHS :: sha256 95da5a6e029855ee :: first epoch of the weight average (epochs count from 0)
\newcommand{\DesignWavgLastEpoch}{79}  % experiments/MTX/k8s/v2/grid/job-mtx2-l188-s1-raunav.yaml :: --num-epochs - 1 :: sha256 d28ce02267a6b87a :: last epoch
\newcommand{\DesignBnJets}{200{,}000}  % experiments/MTX/pretrain_v2.py :: BN_JETS :: sha256 95da5a6e029855ee :: training jets the BatchNorm statistics are recomputed on
\newcommand{\DesignVtwoNWindows}{five}  % experiments/MTX/k8s/v2/grid/job-mtx2-l188-s1-raunav.yaml :: 1 / --data-fraction :: sha256 d28ce02267a6b87a :: disjoint windows of each file's rows, one read per epoch
\newcommand{\DesignLofoNWindows}{three}  % experiments/MTX/k8s/v2/cut/grid/job-mtx2-l188lofo4p-s1-raunav.yaml :: --data-windows :: sha256 44a9573929efb7ac :: windows of the models trained without one family
\newcommand{\DesignVtwoPlanUsed}{97\%}  % experiments/FIGS/data/v2_loader/loader_dryrun/dryrun_s176_seed1.json :: epochs[*].(max_fetch_id + 1) / args.data_split_num, mean :: sha256 03b6586b26ad7a2f :: share of an epoch's planned reading used before its jets are reached, 20 epochs
\newcommand{\DesignLofoPlanUsed}{76\%}  % experiments/FIGS/data/v2_loader/loader_dryrun/dryrun_lofo4p_seed1.json :: epochs[*].(max_fetch_id + 1) / fetches_per_pass, mean :: sha256 2a1315e611d13ff5 :: the same, leaving the family out
\newcommand{\DesignLofoRowFactor}{1.29}  % experiments/FIGS/data/v2_loader/loader_dryrun/dryrun_lofo4p_seed1.json :: (mean plan used / --data-windows) over (mean plan used / k) of the full models' dry run :: sha256 2a1315e611d13ff5 :: how much more often a row of a file is read per epoch when the family is left out
\newcommand{\DesignDistinctJetsCycle}{\ensuremath{3.3\times10^{7}}}  % experiments/FIGS/data/v2_loader/loader_dryrun/dryrun_s176_seed1.json :: sum of epochs[c*k:(c+1)*k].distinct_jets, mean over 4 cycles :: sha256 03b6586b26ad7a2f :: distinct jets the full models train on in one cycle of windows
\newcommand{\DesignLofoFamilyShare}{22.4\%}  % experiments/FIGS/data/v2_loader/loader_dryrun/dryrun_s176_seed1.json :: epochs[*].native_counts over summary.labels_absent_every_epoch of the family-out run :: sha256 03b6586b26ad7a2f :: the left-out family's share of the full models' stream
\newcommand{\DesignLofoNClasses}{27}  % experiments/FIGS/data/v2_loader/loader_dryrun/dryrun_lofo4p_seed1.json :: summary.labels_absent_every_epoch (count) :: sha256 2a1315e611d13ff5 :: native classes of the left-out family
\newcommand{\DesignLofoExposureFactor}{1.29}  % experiments/FIGS/data/v2_loader/loader_dryrun/dryrun_s176_seed1.json :: 1 / (1 - family share) :: sha256 03b6586b26ad7a2f :: how much more often every other class is seen
\newcommand{\DesignLofoQcdShareFull}{14.4\%}  % experiments/FIGS/data/v2_loader/loader_dryrun/dryrun_s176_seed1.json :: summary.mean_qcd_share :: sha256 03b6586b26ad7a2f :: QCD share of the realised stream
\newcommand{\DesignLofoQcdShareLofo}{18.6\%}  % experiments/FIGS/data/v2_loader/loader_dryrun/dryrun_lofo4p_seed1.json :: summary.mean_qcd_share :: sha256 2a1315e611d13ff5 :: QCD share of the realised stream
\newcommand{\DesignLrFlatPercent}{70\%}  % experiments/MTX/pretrain_v2.py :: make_optimizer: 1 - the decay share :: sha256 95da5a6e029855ee :: share of the epochs at the initial rate
\newcommand{\DesignLrDecayFirstEpoch}{56}  % experiments/MTX/pretrain_v2.py :: num_epochs - int(num_epochs * decay share) :: sha256 95da5a6e029855ee :: first epoch at a reduced rate (epochs from 0)
\newcommand{\DesignLrEndPercent}{1\%}  % experiments/MTX/pretrain_v2.py :: make_optimizer: gamma ** n_decay :: sha256 95da5a6e029855ee :: the last epoch's rate over the initial rate
\newcommand{\DesignEndLr}{\ensuremath{5\times10^{-6}}}  % experiments/MTX/k8s/v2/grid/job-mtx2-l188-s1-raunav.yaml :: --start-lr x the final factor :: sha256 d28ce02267a6b87a :: the rate of the last epoch
\newcommand{\DesignWeaverVersion}{0.4.17}  % experiments/MTX/pretrain_v2.py :: the weaver release whose loop it reproduces :: sha256 95da5a6e029855ee
\newcommand{\DesignMeiFactor}{1.1}  % src/stats/mde.py :: MEI_LOG = ln(factor) :: sha256 626468dcc37940bb :: the smallest effect of interest, as a factor
\newcommand{\DesignVtwoGpu}{RTX 3090}  % experiments/MTX/k8s/v2/grid/job-mtx2-flavf0-s1-raunav.yaml :: nodeAffinity nvidia.com/gpu.product :: sha256 960cad1691314980 :: the product of run index 1, every configuration
\newcommand{\DesignVtwoSecondGpu}{\pending{pending: the GPU type of each run index}}  % experiments/FIGS/make_tables.py :: emit_v2_recipe[DesignVtwoSecondGpu] :: sha256 10e6ec08356a4a2b :: a marker, not a number: the grid's job specs put every run index on one product
\newcommand{\DesignVtwoSecondGpuRuns}{\pending{pending: the GPU type of each run index}}  % experiments/FIGS/make_tables.py :: emit_v2_recipe[DesignVtwoSecondGpuRuns] :: sha256 10e6ec08356a4a2b :: a marker, not a number: the grid's job specs put every run index on one product
\newcommand{\MassNRunsWord}{five}  % experiments/FIGS/data/probe_ladder_mass2x2_mlp2/s1.json :: per-seed files (count) :: sha256 906e778d7043fe67 :: pretraining runs of each mass-output configuration
\newcommand{\MassOmaLinearOnesixtwo}{\ensuremath{(2.36 \pm 0.06)\times10^{-3}}}  % experiments/FIGS/data/probe_ladder_v2_mlp2/analysis/seed_level_results.json :: table[bvc_resonant,linear,162].auc :: sha256 53e1ec81343ac2c1 :: 1-AUC, mean +- SD over 5 seeds
\newcommand{\MassOmaLinearOnesixtwoMass}{\ensuremath{(2.62 \pm 0.08)\times10^{-3}}}  % experiments/FIGS/data/probe_ladder_mass2x2_mlp2/s1.json :: tasks.bvc_resonant.arms.l162mass-s*.linear.auc :: sha256 906e778d7043fe67 :: 1-AUC, mean +- SD over the 5 per-seed files
\newcommand{\MassOmaRatioLinearOnesixtwo}{1.11}  % experiments/FIGS/data/probe_ladder_mass2x2_mlp2/s1.json :: tasks.bvc_resonant.arms.l162mass-s*.linear.auc over table[bvc_resonant,linear,162].auc :: sha256 906e778d7043fe67 :: seed mean of 1-AUC with the mass output over without
\newcommand{\MassOmaLinearOneseven}{\ensuremath{(9.3 \pm 1.5)\times10^{-3}}}  % experiments/FIGS/data/probe_ladder_v2_mlp2/analysis/seed_level_results.json :: table[bvc_resonant,linear,17].auc :: sha256 53e1ec81343ac2c1 :: 1-AUC, mean +- SD over 5 seeds
\newcommand{\MassOmaLinearOnesevenMass}{\ensuremath{(3.0 \pm 1.2)\times10^{-2}}}  % experiments/FIGS/data/probe_ladder_mass2x2_mlp2/s1.json :: tasks.bvc_resonant.arms.r16q1mass-s*.linear.auc :: sha256 906e778d7043fe67 :: 1-AUC, mean +- SD over the 5 per-seed files
\newcommand{\MassOmaRatioLinearOneseven}{3.2}  % experiments/FIGS/data/probe_ladder_mass2x2_mlp2/s1.json :: tasks.bvc_resonant.arms.r16q1mass-s*.linear.auc over table[bvc_resonant,linear,17].auc :: sha256 906e778d7043fe67 :: seed mean of 1-AUC with the mass output over without
\newcommand{\MassOmaMlpOnesixtwo}{\ensuremath{(2.32 \pm 0.05)\times10^{-3}}}  % experiments/FIGS/data/probe_ladder_v2_mlp2/analysis/seed_level_results.json :: table[bvc_resonant,mlp,162].auc :: sha256 53e1ec81343ac2c1 :: 1-AUC, mean +- SD over 5 seeds
\newcommand{\MassOmaMlpOnesixtwoMass}{\ensuremath{(2.52 \pm 0.06)\times10^{-3}}}  % experiments/FIGS/data/probe_ladder_mass2x2_mlp2/s1.json :: tasks.bvc_resonant.arms.l162mass-s*.mlp.auc :: sha256 906e778d7043fe67 :: 1-AUC, mean +- SD over the 5 per-seed files
\newcommand{\MassOmaRatioMlpOnesixtwo}{1.09}  % experiments/FIGS/data/probe_ladder_mass2x2_mlp2/s1.json :: tasks.bvc_resonant.arms.l162mass-s*.mlp.auc over table[bvc_resonant,mlp,162].auc :: sha256 906e778d7043fe67 :: seed mean of 1-AUC with the mass output over without
\newcommand{\MassOmaMlpOneseven}{\ensuremath{(7.3 \pm 0.9)\times10^{-3}}}  % experiments/FIGS/data/probe_ladder_v2_mlp2/analysis/seed_level_results.json :: table[bvc_resonant,mlp,17].auc :: sha256 53e1ec81343ac2c1 :: 1-AUC, mean +- SD over 5 seeds
\newcommand{\MassOmaMlpOnesevenMass}{\ensuremath{(1.5 \pm 0.4)\times10^{-2}}}  % experiments/FIGS/data/probe_ladder_mass2x2_mlp2/s1.json :: tasks.bvc_resonant.arms.r16q1mass-s*.mlp.auc :: sha256 906e778d7043fe67 :: 1-AUC, mean +- SD over the 5 per-seed files
\newcommand{\MassOmaRatioMlpOneseven}{2.0}  % experiments/FIGS/data/probe_ladder_mass2x2_mlp2/s1.json :: tasks.bvc_resonant.arms.r16q1mass-s*.mlp.auc over table[bvc_resonant,mlp,17].auc :: sha256 906e778d7043fe67 :: seed mean of 1-AUC with the mass output over without
\newcommand{\FtAucJciiEThreeOneeighteight}{0.902\,\ensuremath{\pm}\,0.004}  % experiments/FIGS/data/w2b_leg1_metrics_v2.json :: cells.{l188-s1,l188-s2,l188-s3,l188-s4,l188-s5}.N1000.s1.macro_auc_ovr :: sha256 b56bf70be9982044 :: mean +- SD over the row's 5 pretrained models
\newcommand{\FtAucJciiEFourOneeighteight}{0.9723\,\ensuremath{\pm}\,0.0004}  % experiments/FIGS/data/w2b_leg1_metrics_v2.json :: cells.{l188-s1,l188-s2,l188-s3,l188-s4,l188-s5}.N10000.s1.macro_auc_ovr :: sha256 b56bf70be9982044 :: mean +- SD over the row's 5 pretrained models
\newcommand{\FtAucJciiEFiveOneeighteight}{0.9793\,\ensuremath{\pm}\,0.0002}  % experiments/FIGS/data/w2b_leg1_metrics_v2.json :: cells.{l188-s1,l188-s2,l188-s3,l188-s4,l188-s5}.N100000.s1.macro_auc_ovr :: sha256 b56bf70be9982044 :: mean +- SD over the row's 5 pretrained models
\newcommand{\FtAucJciiESixOneeighteight}{0.9809\,\ensuremath{\pm}\,\ensuremath{<}0.0001}  % experiments/FIGS/data/w2b_leg1_metrics_v2.json :: cells.{l188-s1,l188-s2,l188-s3,l188-s4,l188-s5}.N1000000.s1.macro_auc_ovr :: sha256 b56bf70be9982044 :: mean +- SD over the row's 5 pretrained models
\newcommand{\FtAucJciiEThreeOnesixtwo}{0.906\,\ensuremath{\pm}\,0.007}  % experiments/FIGS/data/w2b_leg1_metrics_v2.json :: cells.{l162-s1b,l162-s2,l162-s3,l162-s4,l162-s5}.N1000.s1.macro_auc_ovr :: sha256 b56bf70be9982044 :: mean +- SD over the row's 5 pretrained models
\newcommand{\FtAucJciiEFourOnesixtwo}{0.9724\,\ensuremath{\pm}\,0.0005}  % experiments/FIGS/data/w2b_leg1_metrics_v2.json :: cells.{l162-s1b,l162-s2,l162-s3,l162-s4,l162-s5}.N10000.s1.macro_auc_ovr :: sha256 b56bf70be9982044 :: mean +- SD over the row's 5 pretrained models
\newcommand{\FtAucJciiEFiveOnesixtwo}{0.9794\,\ensuremath{\pm}\,0.0002}  % experiments/FIGS/data/w2b_leg1_metrics_v2.json :: cells.{l162-s1b,l162-s2,l162-s3,l162-s4,l162-s5}.N100000.s1.macro_auc_ovr :: sha256 b56bf70be9982044 :: mean +- SD over the row's 5 pretrained models
\newcommand{\FtAucJciiESixOnesixtwo}{0.9808\,\ensuremath{\pm}\,0.0001}  % experiments/FIGS/data/w2b_leg1_metrics_v2.json :: cells.{l162-s1b,l162-s2,l162-s3,l162-s4,l162-s5}.N1000000.s1.macro_auc_ovr :: sha256 b56bf70be9982044 :: mean +- SD over the row's 5 pretrained models
\newcommand{\FtAucJciiEThreeFourthree}{0.918\,\ensuremath{\pm}\,0.005}  % experiments/FIGS/data/w2b_leg1_metrics_v2.json :: cells.{r42q1-s1,r42q1-s2,r42q1-s3,r42q1-s4,r42q1-s5}.N1000.s1.macro_auc_ovr :: sha256 b56bf70be9982044 :: mean +- SD over the row's 5 pretrained models
\newcommand{\FtAucJciiEFourFourthree}{0.9667\,\ensuremath{\pm}\,0.0005}  % experiments/FIGS/data/w2b_leg1_metrics_v2.json :: cells.{r42q1-s1,r42q1-s2,r42q1-s3,r42q1-s4,r42q1-s5}.N10000.s1.macro_auc_ovr :: sha256 b56bf70be9982044 :: mean +- SD over the row's 5 pretrained models
\newcommand{\FtAucJciiEFiveFourthree}{0.9775\,\ensuremath{\pm}\,0.0003}  % experiments/FIGS/data/w2b_leg1_metrics_v2.json :: cells.{r42q1-s1,r42q1-s2,r42q1-s3,r42q1-s4,r42q1-s5}.N100000.s1.macro_auc_ovr :: sha256 b56bf70be9982044 :: mean +- SD over the row's 5 pretrained models
\newcommand{\FtAucJciiESixFourthree}{0.9800\,\ensuremath{\pm}\,0.0001}  % experiments/FIGS/data/w2b_leg1_metrics_v2.json :: cells.{r42q1-s1,r42q1-s2,r42q1-s3,r42q1-s4,r42q1-s5}.N1000000.s1.macro_auc_ovr :: sha256 b56bf70be9982044 :: mean +- SD over the row's 5 pretrained models
\newcommand{\FtAucJciiEThreeOneseven}{0.905\,\ensuremath{\pm}\,0.008}  % experiments/FIGS/data/w2b_leg1_metrics_v2.json :: cells.{r16q1-s1,r16q1-s2,r16q1-s3,r16q1-s4,r16q1-s5}.N1000.s1.macro_auc_ovr :: sha256 b56bf70be9982044 :: mean +- SD over the row's 5 pretrained models
\newcommand{\FtAucJciiEFourOneseven}{0.9635\,\ensuremath{\pm}\,0.0005}  % experiments/FIGS/data/w2b_leg1_metrics_v2.json :: cells.{r16q1-s1,r16q1-s2,r16q1-s3,r16q1-s4,r16q1-s5}.N10000.s1.macro_auc_ovr :: sha256 b56bf70be9982044 :: mean +- SD over the row's 5 pretrained models
\newcommand{\FtAucJciiEFiveOneseven}{0.9748\,\ensuremath{\pm}\,0.0003}  % experiments/FIGS/data/w2b_leg1_metrics_v2.json :: cells.{r16q1-s1,r16q1-s2,r16q1-s3,r16q1-s4,r16q1-s5}.N100000.s1.macro_auc_ovr :: sha256 b56bf70be9982044 :: mean +- SD over the row's 5 pretrained models
\newcommand{\FtAucJciiESixOneseven}{0.9781\,\ensuremath{\pm}\,0.0001}  % experiments/FIGS/data/w2b_leg1_metrics_v2.json :: cells.{r16q1-s1,r16q1-s2,r16q1-s3,r16q1-s4,r16q1-s5}.N1000000.s1.macro_auc_ovr :: sha256 b56bf70be9982044 :: mean +- SD over the row's 5 pretrained models
\newcommand{\FtAucJciiEThreeOnesixtwoMass}{0.921\,\ensuremath{\pm}\,0.005}  % experiments/FIGS/data/w2b_leg1_metrics_v2.json :: cells.{l162mass-s1,l162mass-s2,l162mass-s3,l162mass-s4,l162mass-s5}.N1000.s1.macro_auc_ovr :: sha256 b56bf70be9982044 :: mean +- SD over the row's 5 pretrained models
\newcommand{\FtAucJciiEFourOnesixtwoMass}{0.9724\,\ensuremath{\pm}\,0.0013}  % experiments/FIGS/data/w2b_leg1_metrics_v2.json :: cells.{l162mass-s1,l162mass-s2,l162mass-s3,l162mass-s4,l162mass-s5}.N10000.s1.macro_auc_ovr :: sha256 b56bf70be9982044 :: mean +- SD over the row's 5 pretrained models
\newcommand{\FtAucJciiEFiveOnesixtwoMass}{0.9790\,\ensuremath{\pm}\,0.0001}  % experiments/FIGS/data/w2b_leg1_metrics_v2.json :: cells.{l162mass-s1,l162mass-s2,l162mass-s3,l162mass-s4,l162mass-s5}.N100000.s1.macro_auc_ovr :: sha256 b56bf70be9982044 :: mean +- SD over the row's 5 pretrained models
\newcommand{\FtAucJciiESixOnesixtwoMass}{0.9805\,\ensuremath{\pm}\,0.0001}  % experiments/FIGS/data/w2b_leg1_metrics_v2.json :: cells.{l162mass-s1,l162mass-s2,l162mass-s3,l162mass-s4,l162mass-s5}.N1000000.s1.macro_auc_ovr :: sha256 b56bf70be9982044 :: mean +- SD over the row's 5 pretrained models
\newcommand{\FtAucJciiEThreeOnesevenMass}{0.904\,\ensuremath{\pm}\,0.007}  % experiments/FIGS/data/w2b_leg1_metrics_v2.json :: cells.{r16q1mass-s1,r16q1mass-s2,r16q1mass-s3,r16q1mass-s4,r16q1mass-s5}.N1000.s1.macro_auc_ovr :: sha256 b56bf70be9982044 :: mean +- SD over the row's 5 pretrained models
\newcommand{\FtAucJciiEFourOnesevenMass}{0.9639\,\ensuremath{\pm}\,0.0007}  % experiments/FIGS/data/w2b_leg1_metrics_v2.json :: cells.{r16q1mass-s1,r16q1mass-s2,r16q1mass-s3,r16q1mass-s4,r16q1mass-s5}.N10000.s1.macro_auc_ovr :: sha256 b56bf70be9982044 :: mean +- SD over the row's 5 pretrained models
\newcommand{\FtAucJciiEFiveOnesevenMass}{0.9740\,\ensuremath{\pm}\,0.0004}  % experiments/FIGS/data/w2b_leg1_metrics_v2.json :: cells.{r16q1mass-s1,r16q1mass-s2,r16q1mass-s3,r16q1mass-s4,r16q1mass-s5}.N100000.s1.macro_auc_ovr :: sha256 b56bf70be9982044 :: mean +- SD over the row's 5 pretrained models
\newcommand{\FtAucJciiESixOnesevenMass}{0.9774\,\ensuremath{\pm}\,0.0001}  % experiments/FIGS/data/w2b_leg1_metrics_v2.json :: cells.{r16q1mass-s1,r16q1mass-s2,r16q1mass-s3,r16q1mass-s4,r16q1mass-s5}.N1000000.s1.macro_auc_ovr :: sha256 b56bf70be9982044 :: mean +- SD over the row's 5 pretrained models
\newcommand{\FtAucJciiEThreeRandDrawOne}{0.921}  % experiments/FIGS/data/w2b_leg1_metrics_v2.json :: cells.{rand-d1-s1b}.N1000.s1.macro_auc_ovr :: sha256 b56bf70be9982044 :: one run
\newcommand{\FtAucJciiEFourRandDrawOne}{0.9675}  % experiments/FIGS/data/w2b_leg1_metrics_v2.json :: cells.{rand-d1-s1b}.N10000.s1.macro_auc_ovr :: sha256 b56bf70be9982044 :: one run
\newcommand{\FtAucJciiEFiveRandDrawOne}{0.97617}  % experiments/FIGS/data/w2b_leg1_metrics_v2.json :: cells.{rand-d1-s1b}.N100000.s1.macro_auc_ovr :: sha256 b56bf70be9982044 :: one run
\newcommand{\FtAucJciiESixRandDrawOne}{0.97941}  % experiments/FIGS/data/w2b_leg1_metrics_v2.json :: cells.{rand-d1-s1b}.N1000000.s1.macro_auc_ovr :: sha256 b56bf70be9982044 :: one run
\newcommand{\FtAucJciiEThreeRandDrawTwo}{0.914}  % experiments/FIGS/data/w2b_leg1_metrics_v2.json :: cells.{rand-d2-s2}.N1000.s1.macro_auc_ovr :: sha256 b56bf70be9982044 :: one run
\newcommand{\FtAucJciiEFourRandDrawTwo}{0.9675}  % experiments/FIGS/data/w2b_leg1_metrics_v2.json :: cells.{rand-d2-s2}.N10000.s1.macro_auc_ovr :: sha256 b56bf70be9982044 :: one run
\newcommand{\FtAucJciiEFiveRandDrawTwo}{0.97668}  % experiments/FIGS/data/w2b_leg1_metrics_v2.json :: cells.{rand-d2-s2}.N100000.s1.macro_auc_ovr :: sha256 b56bf70be9982044 :: one run
\newcommand{\FtAucJciiESixRandDrawTwo}{0.97943}  % experiments/FIGS/data/w2b_leg1_metrics_v2.json :: cells.{rand-d2-s2}.N1000000.s1.macro_auc_ovr :: sha256 b56bf70be9982044 :: one run
\newcommand{\FtAucJciiEThreeRandDrawThree}{0.918}  % experiments/FIGS/data/w2b_leg1_metrics_v2.json :: cells.{rand-d3-s3}.N1000.s1.macro_auc_ovr :: sha256 b56bf70be9982044 :: one run
\newcommand{\FtAucJciiEFourRandDrawThree}{0.9686}  % experiments/FIGS/data/w2b_leg1_metrics_v2.json :: cells.{rand-d3-s3}.N10000.s1.macro_auc_ovr :: sha256 b56bf70be9982044 :: one run
\newcommand{\FtAucJciiEFiveRandDrawThree}{0.97605}  % experiments/FIGS/data/w2b_leg1_metrics_v2.json :: cells.{rand-d3-s3}.N100000.s1.macro_auc_ovr :: sha256 b56bf70be9982044 :: one run
\newcommand{\FtAucJciiESixRandDrawThree}{0.97950}  % experiments/FIGS/data/w2b_leg1_metrics_v2.json :: cells.{rand-d3-s3}.N1000000.s1.macro_auc_ovr :: sha256 b56bf70be9982044 :: one run
\newcommand{\FtAccJciiEThreeOneeighteight}{0.445\,\ensuremath{\pm}\,0.004}  % experiments/FIGS/data/w2b_leg1_metrics_v2.json :: cells.{l188-s1,l188-s2,l188-s3,l188-s4,l188-s5}.N1000.s1.accuracy :: sha256 b56bf70be9982044 :: mean +- SD over the row's 5 pretrained models
\newcommand{\FtAccJciiEFourOneeighteight}{0.5199\,\ensuremath{\pm}\,0.0018}  % experiments/FIGS/data/w2b_leg1_metrics_v2.json :: cells.{l188-s1,l188-s2,l188-s3,l188-s4,l188-s5}.N10000.s1.accuracy :: sha256 b56bf70be9982044 :: mean +- SD over the row's 5 pretrained models
\newcommand{\FtAccJciiEFiveOneeighteight}{0.5379\,\ensuremath{\pm}\,0.0010}  % experiments/FIGS/data/w2b_leg1_metrics_v2.json :: cells.{l188-s1,l188-s2,l188-s3,l188-s4,l188-s5}.N100000.s1.accuracy :: sha256 b56bf70be9982044 :: mean +- SD over the row's 5 pretrained models
\newcommand{\FtAccJciiESixOneeighteight}{0.5443\,\ensuremath{\pm}\,0.0003}  % experiments/FIGS/data/w2b_leg1_metrics_v2.json :: cells.{l188-s1,l188-s2,l188-s3,l188-s4,l188-s5}.N1000000.s1.accuracy :: sha256 b56bf70be9982044 :: mean +- SD over the row's 5 pretrained models
\newcommand{\FtAccJciiEThreeOnesixtwo}{0.448\,\ensuremath{\pm}\,0.006}  % experiments/FIGS/data/w2b_leg1_metrics_v2.json :: cells.{l162-s1b,l162-s2,l162-s3,l162-s4,l162-s5}.N1000.s1.accuracy :: sha256 b56bf70be9982044 :: mean +- SD over the row's 5 pretrained models
\newcommand{\FtAccJciiEFourOnesixtwo}{0.5212\,\ensuremath{\pm}\,0.0013}  % experiments/FIGS/data/w2b_leg1_metrics_v2.json :: cells.{l162-s1b,l162-s2,l162-s3,l162-s4,l162-s5}.N10000.s1.accuracy :: sha256 b56bf70be9982044 :: mean +- SD over the row's 5 pretrained models
\newcommand{\FtAccJciiEFiveOnesixtwo}{0.5387\,\ensuremath{\pm}\,0.0008}  % experiments/FIGS/data/w2b_leg1_metrics_v2.json :: cells.{l162-s1b,l162-s2,l162-s3,l162-s4,l162-s5}.N100000.s1.accuracy :: sha256 b56bf70be9982044 :: mean +- SD over the row's 5 pretrained models
\newcommand{\FtAccJciiESixOnesixtwo}{0.5440\,\ensuremath{\pm}\,0.0006}  % experiments/FIGS/data/w2b_leg1_metrics_v2.json :: cells.{l162-s1b,l162-s2,l162-s3,l162-s4,l162-s5}.N1000000.s1.accuracy :: sha256 b56bf70be9982044 :: mean +- SD over the row's 5 pretrained models
\newcommand{\FtAccJciiEThreeFourthree}{0.380\,\ensuremath{\pm}\,0.005}  % experiments/FIGS/data/w2b_leg1_metrics_v2.json :: cells.{r42q1-s1,r42q1-s2,r42q1-s3,r42q1-s4,r42q1-s5}.N1000.s1.accuracy :: sha256 b56bf70be9982044 :: mean +- SD over the row's 5 pretrained models
\newcommand{\FtAccJciiEFourFourthree}{0.481\,\ensuremath{\pm}\,0.004}  % experiments/FIGS/data/w2b_leg1_metrics_v2.json :: cells.{r42q1-s1,r42q1-s2,r42q1-s3,r42q1-s4,r42q1-s5}.N10000.s1.accuracy :: sha256 b56bf70be9982044 :: mean +- SD over the row's 5 pretrained models
\newcommand{\FtAccJciiEFiveFourthree}{0.5219\,\ensuremath{\pm}\,0.0017}  % experiments/FIGS/data/w2b_leg1_metrics_v2.json :: cells.{r42q1-s1,r42q1-s2,r42q1-s3,r42q1-s4,r42q1-s5}.N100000.s1.accuracy :: sha256 b56bf70be9982044 :: mean +- SD over the row's 5 pretrained models
\newcommand{\FtAccJciiESixFourthree}{0.5362\,\ensuremath{\pm}\,0.0007}  % experiments/FIGS/data/w2b_leg1_metrics_v2.json :: cells.{r42q1-s1,r42q1-s2,r42q1-s3,r42q1-s4,r42q1-s5}.N1000000.s1.accuracy :: sha256 b56bf70be9982044 :: mean +- SD over the row's 5 pretrained models
\newcommand{\FtAccJciiEThreeOneseven}{0.332\,\ensuremath{\pm}\,0.004}  % experiments/FIGS/data/w2b_leg1_metrics_v2.json :: cells.{r16q1-s1,r16q1-s2,r16q1-s3,r16q1-s4,r16q1-s5}.N1000.s1.accuracy :: sha256 b56bf70be9982044 :: mean +- SD over the row's 5 pretrained models
\newcommand{\FtAccJciiEFourOneseven}{0.447\,\ensuremath{\pm}\,0.013}  % experiments/FIGS/data/w2b_leg1_metrics_v2.json :: cells.{r16q1-s1,r16q1-s2,r16q1-s3,r16q1-s4,r16q1-s5}.N10000.s1.accuracy :: sha256 b56bf70be9982044 :: mean +- SD over the row's 5 pretrained models
\newcommand{\FtAccJciiEFiveOneseven}{0.5078\,\ensuremath{\pm}\,0.0011}  % experiments/FIGS/data/w2b_leg1_metrics_v2.json :: cells.{r16q1-s1,r16q1-s2,r16q1-s3,r16q1-s4,r16q1-s5}.N100000.s1.accuracy :: sha256 b56bf70be9982044 :: mean +- SD over the row's 5 pretrained models
\newcommand{\FtAccJciiESixOneseven}{0.5227\,\ensuremath{\pm}\,0.0010}  % experiments/FIGS/data/w2b_leg1_metrics_v2.json :: cells.{r16q1-s1,r16q1-s2,r16q1-s3,r16q1-s4,r16q1-s5}.N1000000.s1.accuracy :: sha256 b56bf70be9982044 :: mean +- SD over the row's 5 pretrained models
\newcommand{\FtAccJciiEThreeOnesixtwoMass}{0.451\,\ensuremath{\pm}\,0.007}  % experiments/FIGS/data/w2b_leg1_metrics_v2.json :: cells.{l162mass-s1,l162mass-s2,l162mass-s3,l162mass-s4,l162mass-s5}.N1000.s1.accuracy :: sha256 b56bf70be9982044 :: mean +- SD over the row's 5 pretrained models
\newcommand{\FtAccJciiEFourOnesixtwoMass}{0.515\,\ensuremath{\pm}\,0.004}  % experiments/FIGS/data/w2b_leg1_metrics_v2.json :: cells.{l162mass-s1,l162mass-s2,l162mass-s3,l162mass-s4,l162mass-s5}.N10000.s1.accuracy :: sha256 b56bf70be9982044 :: mean +- SD over the row's 5 pretrained models
\newcommand{\FtAccJciiEFiveOnesixtwoMass}{0.5357\,\ensuremath{\pm}\,0.0007}  % experiments/FIGS/data/w2b_leg1_metrics_v2.json :: cells.{l162mass-s1,l162mass-s2,l162mass-s3,l162mass-s4,l162mass-s5}.N100000.s1.accuracy :: sha256 b56bf70be9982044 :: mean +- SD over the row's 5 pretrained models
\newcommand{\FtAccJciiESixOnesixtwoMass}{0.5409\,\ensuremath{\pm}\,0.0010}  % experiments/FIGS/data/w2b_leg1_metrics_v2.json :: cells.{l162mass-s1,l162mass-s2,l162mass-s3,l162mass-s4,l162mass-s5}.N1000000.s1.accuracy :: sha256 b56bf70be9982044 :: mean +- SD over the row's 5 pretrained models
\newcommand{\FtAccJciiEThreeOnesevenMass}{0.331\,\ensuremath{\pm}\,0.005}  % experiments/FIGS/data/w2b_leg1_metrics_v2.json :: cells.{r16q1mass-s1,r16q1mass-s2,r16q1mass-s3,r16q1mass-s4,r16q1mass-s5}.N1000.s1.accuracy :: sha256 b56bf70be9982044 :: mean +- SD over the row's 5 pretrained models
\newcommand{\FtAccJciiEFourOnesevenMass}{0.4507\,\ensuremath{\pm}\,0.0034}  % experiments/FIGS/data/w2b_leg1_metrics_v2.json :: cells.{r16q1mass-s1,r16q1mass-s2,r16q1mass-s3,r16q1mass-s4,r16q1mass-s5}.N10000.s1.accuracy :: sha256 b56bf70be9982044 :: mean +- SD over the row's 5 pretrained models
\newcommand{\FtAccJciiEFiveOnesevenMass}{0.4993\,\ensuremath{\pm}\,0.0022}  % experiments/FIGS/data/w2b_leg1_metrics_v2.json :: cells.{r16q1mass-s1,r16q1mass-s2,r16q1mass-s3,r16q1mass-s4,r16q1mass-s5}.N100000.s1.accuracy :: sha256 b56bf70be9982044 :: mean +- SD over the row's 5 pretrained models
\newcommand{\FtAccJciiESixOnesevenMass}{0.5166\,\ensuremath{\pm}\,0.0013}  % experiments/FIGS/data/w2b_leg1_metrics_v2.json :: cells.{r16q1mass-s1,r16q1mass-s2,r16q1mass-s3,r16q1mass-s4,r16q1mass-s5}.N1000000.s1.accuracy :: sha256 b56bf70be9982044 :: mean +- SD over the row's 5 pretrained models
\newcommand{\FtAccJciiEThreeRandDrawOne}{0.424}  % experiments/FIGS/data/w2b_leg1_metrics_v2.json :: cells.{rand-d1-s1b}.N1000.s1.accuracy :: sha256 b56bf70be9982044 :: one run
\newcommand{\FtAccJciiEFourRandDrawOne}{0.5070}  % experiments/FIGS/data/w2b_leg1_metrics_v2.json :: cells.{rand-d1-s1b}.N10000.s1.accuracy :: sha256 b56bf70be9982044 :: one run
\newcommand{\FtAccJciiEFiveRandDrawOne}{0.5226}  % experiments/FIGS/data/w2b_leg1_metrics_v2.json :: cells.{rand-d1-s1b}.N100000.s1.accuracy :: sha256 b56bf70be9982044 :: one run
\newcommand{\FtAccJciiESixRandDrawOne}{0.53354}  % experiments/FIGS/data/w2b_leg1_metrics_v2.json :: cells.{rand-d1-s1b}.N1000000.s1.accuracy :: sha256 b56bf70be9982044 :: one run
\newcommand{\FtAccJciiEThreeRandDrawTwo}{0.420}  % experiments/FIGS/data/w2b_leg1_metrics_v2.json :: cells.{rand-d2-s2}.N1000.s1.accuracy :: sha256 b56bf70be9982044 :: one run
\newcommand{\FtAccJciiEFourRandDrawTwo}{0.5074}  % experiments/FIGS/data/w2b_leg1_metrics_v2.json :: cells.{rand-d2-s2}.N10000.s1.accuracy :: sha256 b56bf70be9982044 :: one run
\newcommand{\FtAccJciiEFiveRandDrawTwo}{0.5243}  % experiments/FIGS/data/w2b_leg1_metrics_v2.json :: cells.{rand-d2-s2}.N100000.s1.accuracy :: sha256 b56bf70be9982044 :: one run
\newcommand{\FtAccJciiESixRandDrawTwo}{0.53441}  % experiments/FIGS/data/w2b_leg1_metrics_v2.json :: cells.{rand-d2-s2}.N1000000.s1.accuracy :: sha256 b56bf70be9982044 :: one run
\newcommand{\FtAccJciiEThreeRandDrawThree}{0.423}  % experiments/FIGS/data/w2b_leg1_metrics_v2.json :: cells.{rand-d3-s3}.N1000.s1.accuracy :: sha256 b56bf70be9982044 :: one run
\newcommand{\FtAccJciiEFourRandDrawThree}{0.5114}  % experiments/FIGS/data/w2b_leg1_metrics_v2.json :: cells.{rand-d3-s3}.N10000.s1.accuracy :: sha256 b56bf70be9982044 :: one run
\newcommand{\FtAccJciiEFiveRandDrawThree}{0.5245}  % experiments/FIGS/data/w2b_leg1_metrics_v2.json :: cells.{rand-d3-s3}.N100000.s1.accuracy :: sha256 b56bf70be9982044 :: one run
\newcommand{\FtAccJciiESixRandDrawThree}{0.53473}  % experiments/FIGS/data/w2b_leg1_metrics_v2.json :: cells.{rand-d3-s3}.N1000000.s1.accuracy :: sha256 b56bf70be9982044 :: one run
\newcommand{\FtOmaRatioJciiEThreeOnesixtwoOverOneeighteight}{0.96}  % experiments/FIGS/data/w2b_leg1_metrics_v2.json :: cells.*.N1000.s1.macro_auc_ovr :: sha256 b56bf70be9982044 :: ratio of the seed means of 1 - macro AUC, this row over the reference
\newcommand{\FtOmaRatioJciiEThreeFourthreeOverOneeighteight}{0.84}  % experiments/FIGS/data/w2b_leg1_metrics_v2.json :: cells.*.N1000.s1.macro_auc_ovr :: sha256 b56bf70be9982044 :: ratio of the seed means of 1 - macro AUC, this row over the reference
\newcommand{\FtOmaRatioJciiEThreeOnesevenOverOneeighteight}{0.97}  % experiments/FIGS/data/w2b_leg1_metrics_v2.json :: cells.*.N1000.s1.macro_auc_ovr :: sha256 b56bf70be9982044 :: ratio of the seed means of 1 - macro AUC, this row over the reference
\newcommand{\FtOmaRatioJciiEThreeOnesixtwoMassOverOneeighteight}{0.81}  % experiments/FIGS/data/w2b_leg1_metrics_v2.json :: cells.*.N1000.s1.macro_auc_ovr :: sha256 b56bf70be9982044 :: ratio of the seed means of 1 - macro AUC, this row over the reference
\newcommand{\FtOmaRatioJciiEThreeOnesevenMassOverOneeighteight}{0.98}  % experiments/FIGS/data/w2b_leg1_metrics_v2.json :: cells.*.N1000.s1.macro_auc_ovr :: sha256 b56bf70be9982044 :: ratio of the seed means of 1 - macro AUC, this row over the reference
\newcommand{\FtOmaRatioJciiEThreeRandDrawOneOverOneeighteight}{0.80}  % experiments/FIGS/data/w2b_leg1_metrics_v2.json :: cells.*.N1000.s1.macro_auc_ovr :: sha256 b56bf70be9982044 :: ratio of the seed means of 1 - macro AUC, this row over the reference
\newcommand{\FtOmaRatioJciiEThreeRandDrawTwoOverOneeighteight}{0.87}  % experiments/FIGS/data/w2b_leg1_metrics_v2.json :: cells.*.N1000.s1.macro_auc_ovr :: sha256 b56bf70be9982044 :: ratio of the seed means of 1 - macro AUC, this row over the reference
\newcommand{\FtOmaRatioJciiEThreeRandDrawThreeOverOneeighteight}{0.83}  % experiments/FIGS/data/w2b_leg1_metrics_v2.json :: cells.*.N1000.s1.macro_auc_ovr :: sha256 b56bf70be9982044 :: ratio of the seed means of 1 - macro AUC, this row over the reference
\newcommand{\FtOmaRatioJciiEThreeOneeighteightOverOnesixtwo}{1.05}  % experiments/FIGS/data/w2b_leg1_metrics_v2.json :: cells.*.N1000.s1.macro_auc_ovr :: sha256 b56bf70be9982044 :: ratio of the seed means of 1 - macro AUC, this row over the reference
\newcommand{\FtOmaRatioJciiEThreeFourthreeOverOnesixtwo}{0.88}  % experiments/FIGS/data/w2b_leg1_metrics_v2.json :: cells.*.N1000.s1.macro_auc_ovr :: sha256 b56bf70be9982044 :: ratio of the seed means of 1 - macro AUC, this row over the reference
\newcommand{\FtOmaRatioJciiEThreeOnesevenOverOnesixtwo}{1.02}  % experiments/FIGS/data/w2b_leg1_metrics_v2.json :: cells.*.N1000.s1.macro_auc_ovr :: sha256 b56bf70be9982044 :: ratio of the seed means of 1 - macro AUC, this row over the reference
\newcommand{\FtOmaRatioJciiEThreeOnesixtwoMassOverOnesixtwo}{0.85}  % experiments/FIGS/data/w2b_leg1_metrics_v2.json :: cells.*.N1000.s1.macro_auc_ovr :: sha256 b56bf70be9982044 :: ratio of the seed means of 1 - macro AUC, this row over the reference
\newcommand{\FtOmaRatioJciiEThreeOnesevenMassOverOnesixtwo}{1.02}  % experiments/FIGS/data/w2b_leg1_metrics_v2.json :: cells.*.N1000.s1.macro_auc_ovr :: sha256 b56bf70be9982044 :: ratio of the seed means of 1 - macro AUC, this row over the reference
\newcommand{\FtOmaRatioJciiEThreeRandDrawOneOverOnesixtwo}{0.84}  % experiments/FIGS/data/w2b_leg1_metrics_v2.json :: cells.*.N1000.s1.macro_auc_ovr :: sha256 b56bf70be9982044 :: ratio of the seed means of 1 - macro AUC, this row over the reference
\newcommand{\FtOmaRatioJciiEThreeRandDrawTwoOverOnesixtwo}{0.91}  % experiments/FIGS/data/w2b_leg1_metrics_v2.json :: cells.*.N1000.s1.macro_auc_ovr :: sha256 b56bf70be9982044 :: ratio of the seed means of 1 - macro AUC, this row over the reference
\newcommand{\FtOmaRatioJciiEThreeRandDrawThreeOverOnesixtwo}{0.87}  % experiments/FIGS/data/w2b_leg1_metrics_v2.json :: cells.*.N1000.s1.macro_auc_ovr :: sha256 b56bf70be9982044 :: ratio of the seed means of 1 - macro AUC, this row over the reference
\newcommand{\FtOmaRatioJciiEFourOnesixtwoOverOneeighteight}{0.99}  % experiments/FIGS/data/w2b_leg1_metrics_v2.json :: cells.*.N10000.s1.macro_auc_ovr :: sha256 b56bf70be9982044 :: ratio of the seed means of 1 - macro AUC, this row over the reference
\newcommand{\FtOmaRatioJciiEFourFourthreeOverOneeighteight}{1.20}  % experiments/FIGS/data/w2b_leg1_metrics_v2.json :: cells.*.N10000.s1.macro_auc_ovr :: sha256 b56bf70be9982044 :: ratio of the seed means of 1 - macro AUC, this row over the reference
\newcommand{\FtOmaRatioJciiEFourOnesevenOverOneeighteight}{1.32}  % experiments/FIGS/data/w2b_leg1_metrics_v2.json :: cells.*.N10000.s1.macro_auc_ovr :: sha256 b56bf70be9982044 :: ratio of the seed means of 1 - macro AUC, this row over the reference
\newcommand{\FtOmaRatioJciiEFourOnesixtwoMassOverOneeighteight}{0.99}  % experiments/FIGS/data/w2b_leg1_metrics_v2.json :: cells.*.N10000.s1.macro_auc_ovr :: sha256 b56bf70be9982044 :: ratio of the seed means of 1 - macro AUC, this row over the reference
\newcommand{\FtOmaRatioJciiEFourOnesevenMassOverOneeighteight}{1.30}  % experiments/FIGS/data/w2b_leg1_metrics_v2.json :: cells.*.N10000.s1.macro_auc_ovr :: sha256 b56bf70be9982044 :: ratio of the seed means of 1 - macro AUC, this row over the reference
\newcommand{\FtOmaRatioJciiEFourRandDrawOneOverOneeighteight}{1.17}  % experiments/FIGS/data/w2b_leg1_metrics_v2.json :: cells.*.N10000.s1.macro_auc_ovr :: sha256 b56bf70be9982044 :: ratio of the seed means of 1 - macro AUC, this row over the reference
\newcommand{\FtOmaRatioJciiEFourRandDrawTwoOverOneeighteight}{1.17}  % experiments/FIGS/data/w2b_leg1_metrics_v2.json :: cells.*.N10000.s1.macro_auc_ovr :: sha256 b56bf70be9982044 :: ratio of the seed means of 1 - macro AUC, this row over the reference
\newcommand{\FtOmaRatioJciiEFourRandDrawThreeOverOneeighteight}{1.13}  % experiments/FIGS/data/w2b_leg1_metrics_v2.json :: cells.*.N10000.s1.macro_auc_ovr :: sha256 b56bf70be9982044 :: ratio of the seed means of 1 - macro AUC, this row over the reference
\newcommand{\FtOmaRatioJciiEFourOneeighteightOverOnesixtwo}{1.01}  % experiments/FIGS/data/w2b_leg1_metrics_v2.json :: cells.*.N10000.s1.macro_auc_ovr :: sha256 b56bf70be9982044 :: ratio of the seed means of 1 - macro AUC, this row over the reference
\newcommand{\FtOmaRatioJciiEFourFourthreeOverOnesixtwo}{1.21}  % experiments/FIGS/data/w2b_leg1_metrics_v2.json :: cells.*.N10000.s1.macro_auc_ovr :: sha256 b56bf70be9982044 :: ratio of the seed means of 1 - macro AUC, this row over the reference
\newcommand{\FtOmaRatioJciiEFourOnesevenOverOnesixtwo}{1.32}  % experiments/FIGS/data/w2b_leg1_metrics_v2.json :: cells.*.N10000.s1.macro_auc_ovr :: sha256 b56bf70be9982044 :: ratio of the seed means of 1 - macro AUC, this row over the reference
\newcommand{\FtOmaRatioJciiEFourOnesixtwoMassOverOnesixtwo}{1.00}  % experiments/FIGS/data/w2b_leg1_metrics_v2.json :: cells.*.N10000.s1.macro_auc_ovr :: sha256 b56bf70be9982044 :: ratio of the seed means of 1 - macro AUC, this row over the reference
\newcommand{\FtOmaRatioJciiEFourOnesevenMassOverOnesixtwo}{1.31}  % experiments/FIGS/data/w2b_leg1_metrics_v2.json :: cells.*.N10000.s1.macro_auc_ovr :: sha256 b56bf70be9982044 :: ratio of the seed means of 1 - macro AUC, this row over the reference
\newcommand{\FtOmaRatioJciiEFourRandDrawOneOverOnesixtwo}{1.18}  % experiments/FIGS/data/w2b_leg1_metrics_v2.json :: cells.*.N10000.s1.macro_auc_ovr :: sha256 b56bf70be9982044 :: ratio of the seed means of 1 - macro AUC, this row over the reference
\newcommand{\FtOmaRatioJciiEFourRandDrawTwoOverOnesixtwo}{1.18}  % experiments/FIGS/data/w2b_leg1_metrics_v2.json :: cells.*.N10000.s1.macro_auc_ovr :: sha256 b56bf70be9982044 :: ratio of the seed means of 1 - macro AUC, this row over the reference
\newcommand{\FtOmaRatioJciiEFourRandDrawThreeOverOnesixtwo}{1.14}  % experiments/FIGS/data/w2b_leg1_metrics_v2.json :: cells.*.N10000.s1.macro_auc_ovr :: sha256 b56bf70be9982044 :: ratio of the seed means of 1 - macro AUC, this row over the reference
\newcommand{\FtOmaRatioJciiEFiveOnesixtwoOverOneeighteight}{1.00}  % experiments/FIGS/data/w2b_leg1_metrics_v2.json :: cells.*.N100000.s1.macro_auc_ovr :: sha256 b56bf70be9982044 :: ratio of the seed means of 1 - macro AUC, this row over the reference
\newcommand{\FtOmaRatioJciiEFiveFourthreeOverOneeighteight}{1.09}  % experiments/FIGS/data/w2b_leg1_metrics_v2.json :: cells.*.N100000.s1.macro_auc_ovr :: sha256 b56bf70be9982044 :: ratio of the seed means of 1 - macro AUC, this row over the reference
\newcommand{\FtOmaRatioJciiEFiveOnesevenOverOneeighteight}{1.22}  % experiments/FIGS/data/w2b_leg1_metrics_v2.json :: cells.*.N100000.s1.macro_auc_ovr :: sha256 b56bf70be9982044 :: ratio of the seed means of 1 - macro AUC, this row over the reference
\newcommand{\FtOmaRatioJciiEFiveOnesixtwoMassOverOneeighteight}{1.01}  % experiments/FIGS/data/w2b_leg1_metrics_v2.json :: cells.*.N100000.s1.macro_auc_ovr :: sha256 b56bf70be9982044 :: ratio of the seed means of 1 - macro AUC, this row over the reference
\newcommand{\FtOmaRatioJciiEFiveOnesevenMassOverOneeighteight}{1.26}  % experiments/FIGS/data/w2b_leg1_metrics_v2.json :: cells.*.N100000.s1.macro_auc_ovr :: sha256 b56bf70be9982044 :: ratio of the seed means of 1 - macro AUC, this row over the reference
\newcommand{\FtOmaRatioJciiEFiveRandDrawOneOverOneeighteight}{1.15}  % experiments/FIGS/data/w2b_leg1_metrics_v2.json :: cells.*.N100000.s1.macro_auc_ovr :: sha256 b56bf70be9982044 :: ratio of the seed means of 1 - macro AUC, this row over the reference
\newcommand{\FtOmaRatioJciiEFiveRandDrawTwoOverOneeighteight}{1.12}  % experiments/FIGS/data/w2b_leg1_metrics_v2.json :: cells.*.N100000.s1.macro_auc_ovr :: sha256 b56bf70be9982044 :: ratio of the seed means of 1 - macro AUC, this row over the reference
\newcommand{\FtOmaRatioJciiEFiveRandDrawThreeOverOneeighteight}{1.15}  % experiments/FIGS/data/w2b_leg1_metrics_v2.json :: cells.*.N100000.s1.macro_auc_ovr :: sha256 b56bf70be9982044 :: ratio of the seed means of 1 - macro AUC, this row over the reference
\newcommand{\FtOmaRatioJciiEFiveOneeighteightOverOnesixtwo}{1.00}  % experiments/FIGS/data/w2b_leg1_metrics_v2.json :: cells.*.N100000.s1.macro_auc_ovr :: sha256 b56bf70be9982044 :: ratio of the seed means of 1 - macro AUC, this row over the reference
\newcommand{\FtOmaRatioJciiEFiveFourthreeOverOnesixtwo}{1.09}  % experiments/FIGS/data/w2b_leg1_metrics_v2.json :: cells.*.N100000.s1.macro_auc_ovr :: sha256 b56bf70be9982044 :: ratio of the seed means of 1 - macro AUC, this row over the reference
\newcommand{\FtOmaRatioJciiEFiveOnesevenOverOnesixtwo}{1.22}  % experiments/FIGS/data/w2b_leg1_metrics_v2.json :: cells.*.N100000.s1.macro_auc_ovr :: sha256 b56bf70be9982044 :: ratio of the seed means of 1 - macro AUC, this row over the reference
\newcommand{\FtOmaRatioJciiEFiveOnesixtwoMassOverOnesixtwo}{1.02}  % experiments/FIGS/data/w2b_leg1_metrics_v2.json :: cells.*.N100000.s1.macro_auc_ovr :: sha256 b56bf70be9982044 :: ratio of the seed means of 1 - macro AUC, this row over the reference
\newcommand{\FtOmaRatioJciiEFiveOnesevenMassOverOnesixtwo}{1.26}  % experiments/FIGS/data/w2b_leg1_metrics_v2.json :: cells.*.N100000.s1.macro_auc_ovr :: sha256 b56bf70be9982044 :: ratio of the seed means of 1 - macro AUC, this row over the reference
\newcommand{\FtOmaRatioJciiEFiveRandDrawOneOverOnesixtwo}{1.15}  % experiments/FIGS/data/w2b_leg1_metrics_v2.json :: cells.*.N100000.s1.macro_auc_ovr :: sha256 b56bf70be9982044 :: ratio of the seed means of 1 - macro AUC, this row over the reference
\newcommand{\FtOmaRatioJciiEFiveRandDrawTwoOverOnesixtwo}{1.13}  % experiments/FIGS/data/w2b_leg1_metrics_v2.json :: cells.*.N100000.s1.macro_auc_ovr :: sha256 b56bf70be9982044 :: ratio of the seed means of 1 - macro AUC, this row over the reference
\newcommand{\FtOmaRatioJciiEFiveRandDrawThreeOverOnesixtwo}{1.16}  % experiments/FIGS/data/w2b_leg1_metrics_v2.json :: cells.*.N100000.s1.macro_auc_ovr :: sha256 b56bf70be9982044 :: ratio of the seed means of 1 - macro AUC, this row over the reference
\newcommand{\FtOmaRatioJciiESixOnesixtwoOverOneeighteight}{1.00}  % experiments/FIGS/data/w2b_leg1_metrics_v2.json :: cells.*.N1000000.s1.macro_auc_ovr :: sha256 b56bf70be9982044 :: ratio of the seed means of 1 - macro AUC, this row over the reference
\newcommand{\FtOmaRatioJciiESixFourthreeOverOneeighteight}{1.05}  % experiments/FIGS/data/w2b_leg1_metrics_v2.json :: cells.*.N1000000.s1.macro_auc_ovr :: sha256 b56bf70be9982044 :: ratio of the seed means of 1 - macro AUC, this row over the reference
\newcommand{\FtOmaRatioJciiESixOnesevenOverOneeighteight}{1.14}  % experiments/FIGS/data/w2b_leg1_metrics_v2.json :: cells.*.N1000000.s1.macro_auc_ovr :: sha256 b56bf70be9982044 :: ratio of the seed means of 1 - macro AUC, this row over the reference
\newcommand{\FtOmaRatioJciiESixOnesixtwoMassOverOneeighteight}{1.02}  % experiments/FIGS/data/w2b_leg1_metrics_v2.json :: cells.*.N1000000.s1.macro_auc_ovr :: sha256 b56bf70be9982044 :: ratio of the seed means of 1 - macro AUC, this row over the reference
\newcommand{\FtOmaRatioJciiESixOnesevenMassOverOneeighteight}{1.18}  % experiments/FIGS/data/w2b_leg1_metrics_v2.json :: cells.*.N1000000.s1.macro_auc_ovr :: sha256 b56bf70be9982044 :: ratio of the seed means of 1 - macro AUC, this row over the reference
\newcommand{\FtOmaRatioJciiESixRandDrawOneOverOneeighteight}{1.08}  % experiments/FIGS/data/w2b_leg1_metrics_v2.json :: cells.*.N1000000.s1.macro_auc_ovr :: sha256 b56bf70be9982044 :: ratio of the seed means of 1 - macro AUC, this row over the reference
\newcommand{\FtOmaRatioJciiESixRandDrawTwoOverOneeighteight}{1.07}  % experiments/FIGS/data/w2b_leg1_metrics_v2.json :: cells.*.N1000000.s1.macro_auc_ovr :: sha256 b56bf70be9982044 :: ratio of the seed means of 1 - macro AUC, this row over the reference
\newcommand{\FtOmaRatioJciiESixRandDrawThreeOverOneeighteight}{1.07}  % experiments/FIGS/data/w2b_leg1_metrics_v2.json :: cells.*.N1000000.s1.macro_auc_ovr :: sha256 b56bf70be9982044 :: ratio of the seed means of 1 - macro AUC, this row over the reference
\newcommand{\FtOmaRatioJciiESixOneeighteightOverOnesixtwo}{1.00}  % experiments/FIGS/data/w2b_leg1_metrics_v2.json :: cells.*.N1000000.s1.macro_auc_ovr :: sha256 b56bf70be9982044 :: ratio of the seed means of 1 - macro AUC, this row over the reference
\newcommand{\FtOmaRatioJciiESixFourthreeOverOnesixtwo}{1.04}  % experiments/FIGS/data/w2b_leg1_metrics_v2.json :: cells.*.N1000000.s1.macro_auc_ovr :: sha256 b56bf70be9982044 :: ratio of the seed means of 1 - macro AUC, this row over the reference
\newcommand{\FtOmaRatioJciiESixOnesevenOverOnesixtwo}{1.14}  % experiments/FIGS/data/w2b_leg1_metrics_v2.json :: cells.*.N1000000.s1.macro_auc_ovr :: sha256 b56bf70be9982044 :: ratio of the seed means of 1 - macro AUC, this row over the reference
\newcommand{\FtOmaRatioJciiESixOnesixtwoMassOverOnesixtwo}{1.02}  % experiments/FIGS/data/w2b_leg1_metrics_v2.json :: cells.*.N1000000.s1.macro_auc_ovr :: sha256 b56bf70be9982044 :: ratio of the seed means of 1 - macro AUC, this row over the reference
\newcommand{\FtOmaRatioJciiESixOnesevenMassOverOnesixtwo}{1.18}  % experiments/FIGS/data/w2b_leg1_metrics_v2.json :: cells.*.N1000000.s1.macro_auc_ovr :: sha256 b56bf70be9982044 :: ratio of the seed means of 1 - macro AUC, this row over the reference
\newcommand{\FtOmaRatioJciiESixRandDrawOneOverOnesixtwo}{1.07}  % experiments/FIGS/data/w2b_leg1_metrics_v2.json :: cells.*.N1000000.s1.macro_auc_ovr :: sha256 b56bf70be9982044 :: ratio of the seed means of 1 - macro AUC, this row over the reference
\newcommand{\FtOmaRatioJciiESixRandDrawTwoOverOnesixtwo}{1.07}  % experiments/FIGS/data/w2b_leg1_metrics_v2.json :: cells.*.N1000000.s1.macro_auc_ovr :: sha256 b56bf70be9982044 :: ratio of the seed means of 1 - macro AUC, this row over the reference
\newcommand{\FtOmaRatioJciiESixRandDrawThreeOverOnesixtwo}{1.07}  % experiments/FIGS/data/w2b_leg1_metrics_v2.json :: cells.*.N1000000.s1.macro_auc_ovr :: sha256 b56bf70be9982044 :: ratio of the seed means of 1 - macro AUC, this row over the reference
\newcommand{\PairedFtAccJciiEThreeOnesevenMinusOneeighteight}{\ensuremath{-}0.1132~[\ensuremath{-}0.1214, \ensuremath{-}0.1051]}  % experiments/FIGS/data/w2b_leg1_metrics_v2.json :: cells.{r16q1-s1,r16q1-s2,r16q1-s3,r16q1-s4,r16q1-s5,l188-s1,l188-s2,l188-s3,l188-s4,l188-s5}.N1000.s1.accuracy :: sha256 b56bf70be9982044 :: accuracy, coarse minus fine, paired by run [95% Student-t interval]
\newcommand{\PairedFtAccJciiEFourOnesevenMinusOneeighteight}{\ensuremath{-}0.0730~[\ensuremath{-}0.0902, \ensuremath{-}0.0558]}  % experiments/FIGS/data/w2b_leg1_metrics_v2.json :: cells.{r16q1-s1,r16q1-s2,r16q1-s3,r16q1-s4,r16q1-s5,l188-s1,l188-s2,l188-s3,l188-s4,l188-s5}.N10000.s1.accuracy :: sha256 b56bf70be9982044 :: accuracy, coarse minus fine, paired by run [95% Student-t interval]
\newcommand{\PairedFtAccJciiEFiveOnesevenMinusOneeighteight}{\ensuremath{-}0.0301~[\ensuremath{-}0.0319, \ensuremath{-}0.0284]}  % experiments/FIGS/data/w2b_leg1_metrics_v2.json :: cells.{r16q1-s1,r16q1-s2,r16q1-s3,r16q1-s4,r16q1-s5,l188-s1,l188-s2,l188-s3,l188-s4,l188-s5}.N100000.s1.accuracy :: sha256 b56bf70be9982044 :: accuracy, coarse minus fine, paired by run [95% Student-t interval]
\newcommand{\PairedFtAccJciiESixOnesevenMinusOneeighteight}{\ensuremath{-}0.0216~[\ensuremath{-}0.0232, \ensuremath{-}0.0199]}  % experiments/FIGS/data/w2b_leg1_metrics_v2.json :: cells.{r16q1-s1,r16q1-s2,r16q1-s3,r16q1-s4,r16q1-s5,l188-s1,l188-s2,l188-s3,l188-s4,l188-s5}.N1000000.s1.accuracy :: sha256 b56bf70be9982044 :: accuracy, coarse minus fine, paired by run [95% Student-t interval]
\newcommand{\FtNClassesJcii}{162}  % experiments/FIGS/data/w2b_leg1_metrics_v2.json :: cells.*.*.s1.n_classes_present :: sha256 b56bf70be9982044 :: classes of the fine-tuning task
\newcommand{\FtNTestAucJcii}{125{,}000}  % experiments/FIGS/data/w2b_leg1_metrics_v2.json :: cells.*.*.s1.n_jets_auc :: sha256 b56bf70be9982044 :: test jets the macro AUC is computed on
\newcommand{\FtNTestAccJcii}{500{,}000}  % experiments/FIGS/data/w2b_leg1_metrics_v2.json :: cells.*.*.s1.n_jets :: sha256 b56bf70be9982044 :: test jets the accuracy is computed on
\newcommand{\FtAucJcEThreeOneeighteight}{0.950\,\ensuremath{\pm}\,0.004}  % experiments/FIGS/data/w2b_leg2_metrics_v2.json :: cells.{l188-s1,l188-s2,l188-s3,l188-s4,l188-s5}.N1000.s1.macro_auc_ovr :: sha256 9cd059d3e2f51798 :: mean +- SD over the row's 5 pretrained models
\newcommand{\FtAucJcEFourOneeighteight}{0.9765\,\ensuremath{\pm}\,0.0008}  % experiments/FIGS/data/w2b_leg2_metrics_v2.json :: cells.{l188-s1,l188-s2,l188-s3,l188-s4,l188-s5}.N10000.s1.macro_auc_ovr :: sha256 9cd059d3e2f51798 :: mean +- SD over the row's 5 pretrained models
\newcommand{\FtAucJcEFiveOneeighteight}{0.9833\,\ensuremath{\pm}\,0.0002}  % experiments/FIGS/data/w2b_leg2_metrics_v2.json :: cells.{l188-s1,l188-s2,l188-s3,l188-s4,l188-s5}.N100000.s1.macro_auc_ovr :: sha256 9cd059d3e2f51798 :: mean +- SD over the row's 5 pretrained models
\newcommand{\FtAucJcESixOneeighteight}{0.9858\,\ensuremath{\pm}\,0.0001}  % experiments/FIGS/data/w2b_leg2_metrics_v2.json :: cells.{l188-s1,l188-s2,l188-s3,l188-s4,l188-s5}.N1000000.s1.macro_auc_ovr :: sha256 9cd059d3e2f51798 :: mean +- SD over the row's 5 pretrained models
\newcommand{\FtAucJcEThreeOnesixtwo}{0.9526\,\ensuremath{\pm}\,0.0020}  % experiments/FIGS/data/w2b_leg2_metrics_v2.json :: cells.{l162-s1b,l162-s2,l162-s3,l162-s4,l162-s5}.N1000.s1.macro_auc_ovr :: sha256 9cd059d3e2f51798 :: mean +- SD over the row's 5 pretrained models
\newcommand{\FtAucJcEFourOnesixtwo}{0.9770\,\ensuremath{\pm}\,0.0007}  % experiments/FIGS/data/w2b_leg2_metrics_v2.json :: cells.{l162-s1b,l162-s2,l162-s3,l162-s4,l162-s5}.N10000.s1.macro_auc_ovr :: sha256 9cd059d3e2f51798 :: mean +- SD over the row's 5 pretrained models
\newcommand{\FtAucJcEFiveOnesixtwo}{0.9836\,\ensuremath{\pm}\,0.0001}  % experiments/FIGS/data/w2b_leg2_metrics_v2.json :: cells.{l162-s1b,l162-s2,l162-s3,l162-s4,l162-s5}.N100000.s1.macro_auc_ovr :: sha256 9cd059d3e2f51798 :: mean +- SD over the row's 5 pretrained models
\newcommand{\FtAucJcESixOnesixtwo}{0.9858\,\ensuremath{\pm}\,0.0001}  % experiments/FIGS/data/w2b_leg2_metrics_v2.json :: cells.{l162-s1b,l162-s2,l162-s3,l162-s4,l162-s5}.N1000000.s1.macro_auc_ovr :: sha256 9cd059d3e2f51798 :: mean +- SD over the row's 5 pretrained models
\newcommand{\FtAucJcEThreeFourthree}{0.950\,\ensuremath{\pm}\,0.004}  % experiments/FIGS/data/w2b_leg2_metrics_v2.json :: cells.{r42q1-s1,r42q1-s2,r42q1-s3,r42q1-s4,r42q1-s5}.N1000.s1.macro_auc_ovr :: sha256 9cd059d3e2f51798 :: mean +- SD over the row's 5 pretrained models
\newcommand{\FtAucJcEFourFourthree}{0.9769\,\ensuremath{\pm}\,0.0005}  % experiments/FIGS/data/w2b_leg2_metrics_v2.json :: cells.{r42q1-s1,r42q1-s2,r42q1-s3,r42q1-s4,r42q1-s5}.N10000.s1.macro_auc_ovr :: sha256 9cd059d3e2f51798 :: mean +- SD over the row's 5 pretrained models
\newcommand{\FtAucJcEFiveFourthree}{0.9831\,\ensuremath{\pm}\,0.0002}  % experiments/FIGS/data/w2b_leg2_metrics_v2.json :: cells.{r42q1-s1,r42q1-s2,r42q1-s3,r42q1-s4,r42q1-s5}.N100000.s1.macro_auc_ovr :: sha256 9cd059d3e2f51798 :: mean +- SD over the row's 5 pretrained models
\newcommand{\FtAucJcESixFourthree}{0.9855\,\ensuremath{\pm}\,0.0001}  % experiments/FIGS/data/w2b_leg2_metrics_v2.json :: cells.{r42q1-s1,r42q1-s2,r42q1-s3,r42q1-s4,r42q1-s5}.N1000000.s1.macro_auc_ovr :: sha256 9cd059d3e2f51798 :: mean +- SD over the row's 5 pretrained models
\newcommand{\FtAucJcEThreeOneseven}{0.9470\,\ensuremath{\pm}\,0.0024}  % experiments/FIGS/data/w2b_leg2_metrics_v2.json :: cells.{r16q1-s1,r16q1-s2,r16q1-s3,r16q1-s4,r16q1-s5}.N1000.s1.macro_auc_ovr :: sha256 9cd059d3e2f51798 :: mean +- SD over the row's 5 pretrained models
\newcommand{\FtAucJcEFourOneseven}{0.9749\,\ensuremath{\pm}\,0.0006}  % experiments/FIGS/data/w2b_leg2_metrics_v2.json :: cells.{r16q1-s1,r16q1-s2,r16q1-s3,r16q1-s4,r16q1-s5}.N10000.s1.macro_auc_ovr :: sha256 9cd059d3e2f51798 :: mean +- SD over the row's 5 pretrained models
\newcommand{\FtAucJcEFiveOneseven}{0.9821\,\ensuremath{\pm}\,\ensuremath{<}0.0001}  % experiments/FIGS/data/w2b_leg2_metrics_v2.json :: cells.{r16q1-s1,r16q1-s2,r16q1-s3,r16q1-s4,r16q1-s5}.N100000.s1.macro_auc_ovr :: sha256 9cd059d3e2f51798 :: mean +- SD over the row's 5 pretrained models
\newcommand{\FtAucJcESixOneseven}{0.9846\,\ensuremath{\pm}\,0.0001}  % experiments/FIGS/data/w2b_leg2_metrics_v2.json :: cells.{r16q1-s1,r16q1-s2,r16q1-s3,r16q1-s4,r16q1-s5}.N1000000.s1.macro_auc_ovr :: sha256 9cd059d3e2f51798 :: mean +- SD over the row's 5 pretrained models
\newcommand{\FtAucJcEThreeOnesixtwoMass}{0.9569\,\ensuremath{\pm}\,0.0016}  % experiments/FIGS/data/w2b_leg2_metrics_v2.json :: cells.{l162mass-s1,l162mass-s2,l162mass-s3,l162mass-s4,l162mass-s5}.N1000.s1.macro_auc_ovr :: sha256 9cd059d3e2f51798 :: mean +- SD over the row's 5 pretrained models
\newcommand{\FtAucJcEFourOnesixtwoMass}{0.9768\,\ensuremath{\pm}\,0.0005}  % experiments/FIGS/data/w2b_leg2_metrics_v2.json :: cells.{l162mass-s1,l162mass-s2,l162mass-s3,l162mass-s4,l162mass-s5}.N10000.s1.macro_auc_ovr :: sha256 9cd059d3e2f51798 :: mean +- SD over the row's 5 pretrained models
\newcommand{\FtAucJcEFiveOnesixtwoMass}{0.9834\,\ensuremath{\pm}\,0.0001}  % experiments/FIGS/data/w2b_leg2_metrics_v2.json :: cells.{l162mass-s1,l162mass-s2,l162mass-s3,l162mass-s4,l162mass-s5}.N100000.s1.macro_auc_ovr :: sha256 9cd059d3e2f51798 :: mean +- SD over the row's 5 pretrained models
\newcommand{\FtAucJcESixOnesixtwoMass}{0.9856\,\ensuremath{\pm}\,0.0001}  % experiments/FIGS/data/w2b_leg2_metrics_v2.json :: cells.{l162mass-s1,l162mass-s2,l162mass-s3,l162mass-s4,l162mass-s5}.N1000000.s1.macro_auc_ovr :: sha256 9cd059d3e2f51798 :: mean +- SD over the row's 5 pretrained models
\newcommand{\FtAucJcEThreeOnesevenMass}{0.9460\,\ensuremath{\pm}\,0.0029}  % experiments/FIGS/data/w2b_leg2_metrics_v2.json :: cells.{r16q1mass-s1,r16q1mass-s2,r16q1mass-s3,r16q1mass-s4,r16q1mass-s5}.N1000.s1.macro_auc_ovr :: sha256 9cd059d3e2f51798 :: mean +- SD over the row's 5 pretrained models
\newcommand{\FtAucJcEFourOnesevenMass}{0.9737\,\ensuremath{\pm}\,0.0008}  % experiments/FIGS/data/w2b_leg2_metrics_v2.json :: cells.{r16q1mass-s1,r16q1mass-s2,r16q1mass-s3,r16q1mass-s4,r16q1mass-s5}.N10000.s1.macro_auc_ovr :: sha256 9cd059d3e2f51798 :: mean +- SD over the row's 5 pretrained models
\newcommand{\FtAucJcEFiveOnesevenMass}{0.9817\,\ensuremath{\pm}\,0.0002}  % experiments/FIGS/data/w2b_leg2_metrics_v2.json :: cells.{r16q1mass-s1,r16q1mass-s2,r16q1mass-s3,r16q1mass-s4,r16q1mass-s5}.N100000.s1.macro_auc_ovr :: sha256 9cd059d3e2f51798 :: mean +- SD over the row's 5 pretrained models
\newcommand{\FtAucJcESixOnesevenMass}{0.9842\,\ensuremath{\pm}\,0.0002}  % experiments/FIGS/data/w2b_leg2_metrics_v2.json :: cells.{r16q1mass-s1,r16q1mass-s2,r16q1mass-s3,r16q1mass-s4,r16q1mass-s5}.N1000000.s1.macro_auc_ovr :: sha256 9cd059d3e2f51798 :: mean +- SD over the row's 5 pretrained models
\newcommand{\FtAucJcEThreeRandDrawOne}{0.9477}  % experiments/FIGS/data/w2b_leg2_metrics_v2.json :: cells.{rand-d1-s1b}.N1000.s1.macro_auc_ovr :: sha256 9cd059d3e2f51798 :: one run
\newcommand{\FtAucJcEFourRandDrawOne}{0.9745}  % experiments/FIGS/data/w2b_leg2_metrics_v2.json :: cells.{rand-d1-s1b}.N10000.s1.macro_auc_ovr :: sha256 9cd059d3e2f51798 :: one run
\newcommand{\FtAucJcEFiveRandDrawOne}{0.98265}  % experiments/FIGS/data/w2b_leg2_metrics_v2.json :: cells.{rand-d1-s1b}.N100000.s1.macro_auc_ovr :: sha256 9cd059d3e2f51798 :: one run
\newcommand{\FtAucJcESixRandDrawOne}{0.98513}  % experiments/FIGS/data/w2b_leg2_metrics_v2.json :: cells.{rand-d1-s1b}.N1000000.s1.macro_auc_ovr :: sha256 9cd059d3e2f51798 :: one run
\newcommand{\FtAucJcEThreeRandDrawTwo}{0.9501}  % experiments/FIGS/data/w2b_leg2_metrics_v2.json :: cells.{rand-d2-s2}.N1000.s1.macro_auc_ovr :: sha256 9cd059d3e2f51798 :: one run
\newcommand{\FtAucJcEFourRandDrawTwo}{0.9756}  % experiments/FIGS/data/w2b_leg2_metrics_v2.json :: cells.{rand-d2-s2}.N10000.s1.macro_auc_ovr :: sha256 9cd059d3e2f51798 :: one run
\newcommand{\FtAucJcEFiveRandDrawTwo}{0.98278}  % experiments/FIGS/data/w2b_leg2_metrics_v2.json :: cells.{rand-d2-s2}.N100000.s1.macro_auc_ovr :: sha256 9cd059d3e2f51798 :: one run
\newcommand{\FtAucJcESixRandDrawTwo}{0.98512}  % experiments/FIGS/data/w2b_leg2_metrics_v2.json :: cells.{rand-d2-s2}.N1000000.s1.macro_auc_ovr :: sha256 9cd059d3e2f51798 :: one run
\newcommand{\FtAucJcEThreeRandDrawThree}{0.9470}  % experiments/FIGS/data/w2b_leg2_metrics_v2.json :: cells.{rand-d3-s3}.N1000.s1.macro_auc_ovr :: sha256 9cd059d3e2f51798 :: one run
\newcommand{\FtAucJcEFourRandDrawThree}{0.9743}  % experiments/FIGS/data/w2b_leg2_metrics_v2.json :: cells.{rand-d3-s3}.N10000.s1.macro_auc_ovr :: sha256 9cd059d3e2f51798 :: one run
\newcommand{\FtAucJcEFiveRandDrawThree}{0.98253}  % experiments/FIGS/data/w2b_leg2_metrics_v2.json :: cells.{rand-d3-s3}.N100000.s1.macro_auc_ovr :: sha256 9cd059d3e2f51798 :: one run
\newcommand{\FtAucJcESixRandDrawThree}{0.98506}  % experiments/FIGS/data/w2b_leg2_metrics_v2.json :: cells.{rand-d3-s3}.N1000000.s1.macro_auc_ovr :: sha256 9cd059d3e2f51798 :: one run
\newcommand{\FtAccJcEThreeOneeighteight}{0.7183\,\ensuremath{\pm}\,0.0032}  % experiments/FIGS/data/w2b_leg2_metrics_v2.json :: cells.{l188-s1,l188-s2,l188-s3,l188-s4,l188-s5}.N1000.s1.accuracy :: sha256 9cd059d3e2f51798 :: mean +- SD over the row's 5 pretrained models
\newcommand{\FtAccJcEFourOneeighteight}{0.8039\,\ensuremath{\pm}\,0.0030}  % experiments/FIGS/data/w2b_leg2_metrics_v2.json :: cells.{l188-s1,l188-s2,l188-s3,l188-s4,l188-s5}.N10000.s1.accuracy :: sha256 9cd059d3e2f51798 :: mean +- SD over the row's 5 pretrained models
\newcommand{\FtAccJcEFiveOneeighteight}{0.8351\,\ensuremath{\pm}\,0.0008}  % experiments/FIGS/data/w2b_leg2_metrics_v2.json :: cells.{l188-s1,l188-s2,l188-s3,l188-s4,l188-s5}.N100000.s1.accuracy :: sha256 9cd059d3e2f51798 :: mean +- SD over the row's 5 pretrained models
\newcommand{\FtAccJcESixOneeighteight}{0.8486\,\ensuremath{\pm}\,0.0005}  % experiments/FIGS/data/w2b_leg2_metrics_v2.json :: cells.{l188-s1,l188-s2,l188-s3,l188-s4,l188-s5}.N1000000.s1.accuracy :: sha256 9cd059d3e2f51798 :: mean +- SD over the row's 5 pretrained models
\newcommand{\FtAccJcEThreeOnesixtwo}{0.7268\,\ensuremath{\pm}\,0.0017}  % experiments/FIGS/data/w2b_leg2_metrics_v2.json :: cells.{l162-s1b,l162-s2,l162-s3,l162-s4,l162-s5}.N1000.s1.accuracy :: sha256 9cd059d3e2f51798 :: mean +- SD over the row's 5 pretrained models
\newcommand{\FtAccJcEFourOnesixtwo}{0.8069\,\ensuremath{\pm}\,0.0025}  % experiments/FIGS/data/w2b_leg2_metrics_v2.json :: cells.{l162-s1b,l162-s2,l162-s3,l162-s4,l162-s5}.N10000.s1.accuracy :: sha256 9cd059d3e2f51798 :: mean +- SD over the row's 5 pretrained models
\newcommand{\FtAccJcEFiveOnesixtwo}{0.8363\,\ensuremath{\pm}\,0.0008}  % experiments/FIGS/data/w2b_leg2_metrics_v2.json :: cells.{l162-s1b,l162-s2,l162-s3,l162-s4,l162-s5}.N100000.s1.accuracy :: sha256 9cd059d3e2f51798 :: mean +- SD over the row's 5 pretrained models
\newcommand{\FtAccJcESixOnesixtwo}{0.8485\,\ensuremath{\pm}\,0.0005}  % experiments/FIGS/data/w2b_leg2_metrics_v2.json :: cells.{l162-s1b,l162-s2,l162-s3,l162-s4,l162-s5}.N1000000.s1.accuracy :: sha256 9cd059d3e2f51798 :: mean +- SD over the row's 5 pretrained models
\newcommand{\FtAccJcEThreeFourthree}{0.722\,\ensuremath{\pm}\,0.007}  % experiments/FIGS/data/w2b_leg2_metrics_v2.json :: cells.{r42q1-s1,r42q1-s2,r42q1-s3,r42q1-s4,r42q1-s5}.N1000.s1.accuracy :: sha256 9cd059d3e2f51798 :: mean +- SD over the row's 5 pretrained models
\newcommand{\FtAccJcEFourFourthree}{0.8050\,\ensuremath{\pm}\,0.0026}  % experiments/FIGS/data/w2b_leg2_metrics_v2.json :: cells.{r42q1-s1,r42q1-s2,r42q1-s3,r42q1-s4,r42q1-s5}.N10000.s1.accuracy :: sha256 9cd059d3e2f51798 :: mean +- SD over the row's 5 pretrained models
\newcommand{\FtAccJcEFiveFourthree}{0.8336\,\ensuremath{\pm}\,0.0008}  % experiments/FIGS/data/w2b_leg2_metrics_v2.json :: cells.{r42q1-s1,r42q1-s2,r42q1-s3,r42q1-s4,r42q1-s5}.N100000.s1.accuracy :: sha256 9cd059d3e2f51798 :: mean +- SD over the row's 5 pretrained models
\newcommand{\FtAccJcESixFourthree}{0.8468\,\ensuremath{\pm}\,0.0003}  % experiments/FIGS/data/w2b_leg2_metrics_v2.json :: cells.{r42q1-s1,r42q1-s2,r42q1-s3,r42q1-s4,r42q1-s5}.N1000000.s1.accuracy :: sha256 9cd059d3e2f51798 :: mean +- SD over the row's 5 pretrained models
\newcommand{\FtAccJcEThreeOneseven}{0.714\,\ensuremath{\pm}\,0.005}  % experiments/FIGS/data/w2b_leg2_metrics_v2.json :: cells.{r16q1-s1,r16q1-s2,r16q1-s3,r16q1-s4,r16q1-s5}.N1000.s1.accuracy :: sha256 9cd059d3e2f51798 :: mean +- SD over the row's 5 pretrained models
\newcommand{\FtAccJcEFourOneseven}{0.7964\,\ensuremath{\pm}\,0.0019}  % experiments/FIGS/data/w2b_leg2_metrics_v2.json :: cells.{r16q1-s1,r16q1-s2,r16q1-s3,r16q1-s4,r16q1-s5}.N10000.s1.accuracy :: sha256 9cd059d3e2f51798 :: mean +- SD over the row's 5 pretrained models
\newcommand{\FtAccJcEFiveOneseven}{0.8286\,\ensuremath{\pm}\,0.0006}  % experiments/FIGS/data/w2b_leg2_metrics_v2.json :: cells.{r16q1-s1,r16q1-s2,r16q1-s3,r16q1-s4,r16q1-s5}.N100000.s1.accuracy :: sha256 9cd059d3e2f51798 :: mean +- SD over the row's 5 pretrained models
\newcommand{\FtAccJcESixOneseven}{0.8415\,\ensuremath{\pm}\,0.0007}  % experiments/FIGS/data/w2b_leg2_metrics_v2.json :: cells.{r16q1-s1,r16q1-s2,r16q1-s3,r16q1-s4,r16q1-s5}.N1000000.s1.accuracy :: sha256 9cd059d3e2f51798 :: mean +- SD over the row's 5 pretrained models
\newcommand{\FtAccJcEThreeOnesixtwoMass}{0.728\,\ensuremath{\pm}\,0.004}  % experiments/FIGS/data/w2b_leg2_metrics_v2.json :: cells.{l162mass-s1,l162mass-s2,l162mass-s3,l162mass-s4,l162mass-s5}.N1000.s1.accuracy :: sha256 9cd059d3e2f51798 :: mean +- SD over the row's 5 pretrained models
\newcommand{\FtAccJcEFourOnesixtwoMass}{0.8043\,\ensuremath{\pm}\,0.0020}  % experiments/FIGS/data/w2b_leg2_metrics_v2.json :: cells.{l162mass-s1,l162mass-s2,l162mass-s3,l162mass-s4,l162mass-s5}.N10000.s1.accuracy :: sha256 9cd059d3e2f51798 :: mean +- SD over the row's 5 pretrained models
\newcommand{\FtAccJcEFiveOnesixtwoMass}{0.8349\,\ensuremath{\pm}\,0.0005}  % experiments/FIGS/data/w2b_leg2_metrics_v2.json :: cells.{l162mass-s1,l162mass-s2,l162mass-s3,l162mass-s4,l162mass-s5}.N100000.s1.accuracy :: sha256 9cd059d3e2f51798 :: mean +- SD over the row's 5 pretrained models
\newcommand{\FtAccJcESixOnesixtwoMass}{0.8475\,\ensuremath{\pm}\,0.0009}  % experiments/FIGS/data/w2b_leg2_metrics_v2.json :: cells.{l162mass-s1,l162mass-s2,l162mass-s3,l162mass-s4,l162mass-s5}.N1000000.s1.accuracy :: sha256 9cd059d3e2f51798 :: mean +- SD over the row's 5 pretrained models
\newcommand{\FtAccJcEThreeOnesevenMass}{0.702\,\ensuremath{\pm}\,0.005}  % experiments/FIGS/data/w2b_leg2_metrics_v2.json :: cells.{r16q1mass-s1,r16q1mass-s2,r16q1mass-s3,r16q1mass-s4,r16q1mass-s5}.N1000.s1.accuracy :: sha256 9cd059d3e2f51798 :: mean +- SD over the row's 5 pretrained models
\newcommand{\FtAccJcEFourOnesevenMass}{0.7909\,\ensuremath{\pm}\,0.0031}  % experiments/FIGS/data/w2b_leg2_metrics_v2.json :: cells.{r16q1mass-s1,r16q1mass-s2,r16q1mass-s3,r16q1mass-s4,r16q1mass-s5}.N10000.s1.accuracy :: sha256 9cd059d3e2f51798 :: mean +- SD over the row's 5 pretrained models
\newcommand{\FtAccJcEFiveOnesevenMass}{0.8265\,\ensuremath{\pm}\,0.0007}  % experiments/FIGS/data/w2b_leg2_metrics_v2.json :: cells.{r16q1mass-s1,r16q1mass-s2,r16q1mass-s3,r16q1mass-s4,r16q1mass-s5}.N100000.s1.accuracy :: sha256 9cd059d3e2f51798 :: mean +- SD over the row's 5 pretrained models
\newcommand{\FtAccJcESixOnesevenMass}{0.8390\,\ensuremath{\pm}\,0.0011}  % experiments/FIGS/data/w2b_leg2_metrics_v2.json :: cells.{r16q1mass-s1,r16q1mass-s2,r16q1mass-s3,r16q1mass-s4,r16q1mass-s5}.N1000000.s1.accuracy :: sha256 9cd059d3e2f51798 :: mean +- SD over the row's 5 pretrained models
\newcommand{\FtAccJcEThreeRandDrawOne}{0.7049}  % experiments/FIGS/data/w2b_leg2_metrics_v2.json :: cells.{rand-d1-s1b}.N1000.s1.accuracy :: sha256 9cd059d3e2f51798 :: one run
\newcommand{\FtAccJcEFourRandDrawOne}{0.7974}  % experiments/FIGS/data/w2b_leg2_metrics_v2.json :: cells.{rand-d1-s1b}.N10000.s1.accuracy :: sha256 9cd059d3e2f51798 :: one run
\newcommand{\FtAccJcEFiveRandDrawOne}{0.8326}  % experiments/FIGS/data/w2b_leg2_metrics_v2.json :: cells.{rand-d1-s1b}.N100000.s1.accuracy :: sha256 9cd059d3e2f51798 :: one run
\newcommand{\FtAccJcESixRandDrawOne}{0.84467}  % experiments/FIGS/data/w2b_leg2_metrics_v2.json :: cells.{rand-d1-s1b}.N1000000.s1.accuracy :: sha256 9cd059d3e2f51798 :: one run
\newcommand{\FtAccJcEThreeRandDrawTwo}{0.7135}  % experiments/FIGS/data/w2b_leg2_metrics_v2.json :: cells.{rand-d2-s2}.N1000.s1.accuracy :: sha256 9cd059d3e2f51798 :: one run
\newcommand{\FtAccJcEFourRandDrawTwo}{0.7998}  % experiments/FIGS/data/w2b_leg2_metrics_v2.json :: cells.{rand-d2-s2}.N10000.s1.accuracy :: sha256 9cd059d3e2f51798 :: one run
\newcommand{\FtAccJcEFiveRandDrawTwo}{0.8332}  % experiments/FIGS/data/w2b_leg2_metrics_v2.json :: cells.{rand-d2-s2}.N100000.s1.accuracy :: sha256 9cd059d3e2f51798 :: one run
\newcommand{\FtAccJcESixRandDrawTwo}{0.84449}  % experiments/FIGS/data/w2b_leg2_metrics_v2.json :: cells.{rand-d2-s2}.N1000000.s1.accuracy :: sha256 9cd059d3e2f51798 :: one run
\newcommand{\FtAccJcEThreeRandDrawThree}{0.7090}  % experiments/FIGS/data/w2b_leg2_metrics_v2.json :: cells.{rand-d3-s3}.N1000.s1.accuracy :: sha256 9cd059d3e2f51798 :: one run
\newcommand{\FtAccJcEFourRandDrawThree}{0.7948}  % experiments/FIGS/data/w2b_leg2_metrics_v2.json :: cells.{rand-d3-s3}.N10000.s1.accuracy :: sha256 9cd059d3e2f51798 :: one run
\newcommand{\FtAccJcEFiveRandDrawThree}{0.8310}  % experiments/FIGS/data/w2b_leg2_metrics_v2.json :: cells.{rand-d3-s3}.N100000.s1.accuracy :: sha256 9cd059d3e2f51798 :: one run
\newcommand{\FtAccJcESixRandDrawThree}{0.84409}  % experiments/FIGS/data/w2b_leg2_metrics_v2.json :: cells.{rand-d3-s3}.N1000000.s1.accuracy :: sha256 9cd059d3e2f51798 :: one run
\newcommand{\FtOmaRatioJcEThreeOnesixtwoOverOneeighteight}{0.96}  % experiments/FIGS/data/w2b_leg2_metrics_v2.json :: cells.*.N1000.s1.macro_auc_ovr :: sha256 9cd059d3e2f51798 :: ratio of the seed means of 1 - macro AUC, this row over the reference
\newcommand{\FtOmaRatioJcEThreeFourthreeOverOneeighteight}{1.01}  % experiments/FIGS/data/w2b_leg2_metrics_v2.json :: cells.*.N1000.s1.macro_auc_ovr :: sha256 9cd059d3e2f51798 :: ratio of the seed means of 1 - macro AUC, this row over the reference
\newcommand{\FtOmaRatioJcEThreeOnesevenOverOneeighteight}{1.07}  % experiments/FIGS/data/w2b_leg2_metrics_v2.json :: cells.*.N1000.s1.macro_auc_ovr :: sha256 9cd059d3e2f51798 :: ratio of the seed means of 1 - macro AUC, this row over the reference
\newcommand{\FtOmaRatioJcEThreeOnesixtwoMassOverOneeighteight}{0.87}  % experiments/FIGS/data/w2b_leg2_metrics_v2.json :: cells.*.N1000.s1.macro_auc_ovr :: sha256 9cd059d3e2f51798 :: ratio of the seed means of 1 - macro AUC, this row over the reference
\newcommand{\FtOmaRatioJcEThreeOnesevenMassOverOneeighteight}{1.09}  % experiments/FIGS/data/w2b_leg2_metrics_v2.json :: cells.*.N1000.s1.macro_auc_ovr :: sha256 9cd059d3e2f51798 :: ratio of the seed means of 1 - macro AUC, this row over the reference
\newcommand{\FtOmaRatioJcEThreeRandDrawOneOverOneeighteight}{1.06}  % experiments/FIGS/data/w2b_leg2_metrics_v2.json :: cells.*.N1000.s1.macro_auc_ovr :: sha256 9cd059d3e2f51798 :: ratio of the seed means of 1 - macro AUC, this row over the reference
\newcommand{\FtOmaRatioJcEThreeRandDrawTwoOverOneeighteight}{1.01}  % experiments/FIGS/data/w2b_leg2_metrics_v2.json :: cells.*.N1000.s1.macro_auc_ovr :: sha256 9cd059d3e2f51798 :: ratio of the seed means of 1 - macro AUC, this row over the reference
\newcommand{\FtOmaRatioJcEThreeRandDrawThreeOverOneeighteight}{1.07}  % experiments/FIGS/data/w2b_leg2_metrics_v2.json :: cells.*.N1000.s1.macro_auc_ovr :: sha256 9cd059d3e2f51798 :: ratio of the seed means of 1 - macro AUC, this row over the reference
\newcommand{\FtOmaRatioJcEThreeOneeighteightOverOnesixtwo}{1.05}  % experiments/FIGS/data/w2b_leg2_metrics_v2.json :: cells.*.N1000.s1.macro_auc_ovr :: sha256 9cd059d3e2f51798 :: ratio of the seed means of 1 - macro AUC, this row over the reference
\newcommand{\FtOmaRatioJcEThreeFourthreeOverOnesixtwo}{1.05}  % experiments/FIGS/data/w2b_leg2_metrics_v2.json :: cells.*.N1000.s1.macro_auc_ovr :: sha256 9cd059d3e2f51798 :: ratio of the seed means of 1 - macro AUC, this row over the reference
\newcommand{\FtOmaRatioJcEThreeOnesevenOverOnesixtwo}{1.12}  % experiments/FIGS/data/w2b_leg2_metrics_v2.json :: cells.*.N1000.s1.macro_auc_ovr :: sha256 9cd059d3e2f51798 :: ratio of the seed means of 1 - macro AUC, this row over the reference
\newcommand{\FtOmaRatioJcEThreeOnesixtwoMassOverOnesixtwo}{0.91}  % experiments/FIGS/data/w2b_leg2_metrics_v2.json :: cells.*.N1000.s1.macro_auc_ovr :: sha256 9cd059d3e2f51798 :: ratio of the seed means of 1 - macro AUC, this row over the reference
\newcommand{\FtOmaRatioJcEThreeOnesevenMassOverOnesixtwo}{1.14}  % experiments/FIGS/data/w2b_leg2_metrics_v2.json :: cells.*.N1000.s1.macro_auc_ovr :: sha256 9cd059d3e2f51798 :: ratio of the seed means of 1 - macro AUC, this row over the reference
\newcommand{\FtOmaRatioJcEThreeRandDrawOneOverOnesixtwo}{1.10}  % experiments/FIGS/data/w2b_leg2_metrics_v2.json :: cells.*.N1000.s1.macro_auc_ovr :: sha256 9cd059d3e2f51798 :: ratio of the seed means of 1 - macro AUC, this row over the reference
\newcommand{\FtOmaRatioJcEThreeRandDrawTwoOverOnesixtwo}{1.05}  % experiments/FIGS/data/w2b_leg2_metrics_v2.json :: cells.*.N1000.s1.macro_auc_ovr :: sha256 9cd059d3e2f51798 :: ratio of the seed means of 1 - macro AUC, this row over the reference
\newcommand{\FtOmaRatioJcEThreeRandDrawThreeOverOnesixtwo}{1.12}  % experiments/FIGS/data/w2b_leg2_metrics_v2.json :: cells.*.N1000.s1.macro_auc_ovr :: sha256 9cd059d3e2f51798 :: ratio of the seed means of 1 - macro AUC, this row over the reference
\newcommand{\FtOmaRatioJcEFourOnesixtwoOverOneeighteight}{0.98}  % experiments/FIGS/data/w2b_leg2_metrics_v2.json :: cells.*.N10000.s1.macro_auc_ovr :: sha256 9cd059d3e2f51798 :: ratio of the seed means of 1 - macro AUC, this row over the reference
\newcommand{\FtOmaRatioJcEFourFourthreeOverOneeighteight}{0.98}  % experiments/FIGS/data/w2b_leg2_metrics_v2.json :: cells.*.N10000.s1.macro_auc_ovr :: sha256 9cd059d3e2f51798 :: ratio of the seed means of 1 - macro AUC, this row over the reference
\newcommand{\FtOmaRatioJcEFourOnesevenOverOneeighteight}{1.07}  % experiments/FIGS/data/w2b_leg2_metrics_v2.json :: cells.*.N10000.s1.macro_auc_ovr :: sha256 9cd059d3e2f51798 :: ratio of the seed means of 1 - macro AUC, this row over the reference
\newcommand{\FtOmaRatioJcEFourOnesixtwoMassOverOneeighteight}{0.99}  % experiments/FIGS/data/w2b_leg2_metrics_v2.json :: cells.*.N10000.s1.macro_auc_ovr :: sha256 9cd059d3e2f51798 :: ratio of the seed means of 1 - macro AUC, this row over the reference
\newcommand{\FtOmaRatioJcEFourOnesevenMassOverOneeighteight}{1.12}  % experiments/FIGS/data/w2b_leg2_metrics_v2.json :: cells.*.N10000.s1.macro_auc_ovr :: sha256 9cd059d3e2f51798 :: ratio of the seed means of 1 - macro AUC, this row over the reference
\newcommand{\FtOmaRatioJcEFourRandDrawOneOverOneeighteight}{1.08}  % experiments/FIGS/data/w2b_leg2_metrics_v2.json :: cells.*.N10000.s1.macro_auc_ovr :: sha256 9cd059d3e2f51798 :: ratio of the seed means of 1 - macro AUC, this row over the reference
\newcommand{\FtOmaRatioJcEFourRandDrawTwoOverOneeighteight}{1.04}  % experiments/FIGS/data/w2b_leg2_metrics_v2.json :: cells.*.N10000.s1.macro_auc_ovr :: sha256 9cd059d3e2f51798 :: ratio of the seed means of 1 - macro AUC, this row over the reference
\newcommand{\FtOmaRatioJcEFourRandDrawThreeOverOneeighteight}{1.09}  % experiments/FIGS/data/w2b_leg2_metrics_v2.json :: cells.*.N10000.s1.macro_auc_ovr :: sha256 9cd059d3e2f51798 :: ratio of the seed means of 1 - macro AUC, this row over the reference
\newcommand{\FtOmaRatioJcEFourOneeighteightOverOnesixtwo}{1.02}  % experiments/FIGS/data/w2b_leg2_metrics_v2.json :: cells.*.N10000.s1.macro_auc_ovr :: sha256 9cd059d3e2f51798 :: ratio of the seed means of 1 - macro AUC, this row over the reference
\newcommand{\FtOmaRatioJcEFourFourthreeOverOnesixtwo}{1.01}  % experiments/FIGS/data/w2b_leg2_metrics_v2.json :: cells.*.N10000.s1.macro_auc_ovr :: sha256 9cd059d3e2f51798 :: ratio of the seed means of 1 - macro AUC, this row over the reference
\newcommand{\FtOmaRatioJcEFourOnesevenOverOnesixtwo}{1.09}  % experiments/FIGS/data/w2b_leg2_metrics_v2.json :: cells.*.N10000.s1.macro_auc_ovr :: sha256 9cd059d3e2f51798 :: ratio of the seed means of 1 - macro AUC, this row over the reference
\newcommand{\FtOmaRatioJcEFourOnesixtwoMassOverOnesixtwo}{1.01}  % experiments/FIGS/data/w2b_leg2_metrics_v2.json :: cells.*.N10000.s1.macro_auc_ovr :: sha256 9cd059d3e2f51798 :: ratio of the seed means of 1 - macro AUC, this row over the reference
\newcommand{\FtOmaRatioJcEFourOnesevenMassOverOnesixtwo}{1.14}  % experiments/FIGS/data/w2b_leg2_metrics_v2.json :: cells.*.N10000.s1.macro_auc_ovr :: sha256 9cd059d3e2f51798 :: ratio of the seed means of 1 - macro AUC, this row over the reference
\newcommand{\FtOmaRatioJcEFourRandDrawOneOverOnesixtwo}{1.11}  % experiments/FIGS/data/w2b_leg2_metrics_v2.json :: cells.*.N10000.s1.macro_auc_ovr :: sha256 9cd059d3e2f51798 :: ratio of the seed means of 1 - macro AUC, this row over the reference
\newcommand{\FtOmaRatioJcEFourRandDrawTwoOverOnesixtwo}{1.06}  % experiments/FIGS/data/w2b_leg2_metrics_v2.json :: cells.*.N10000.s1.macro_auc_ovr :: sha256 9cd059d3e2f51798 :: ratio of the seed means of 1 - macro AUC, this row over the reference
\newcommand{\FtOmaRatioJcEFourRandDrawThreeOverOnesixtwo}{1.12}  % experiments/FIGS/data/w2b_leg2_metrics_v2.json :: cells.*.N10000.s1.macro_auc_ovr :: sha256 9cd059d3e2f51798 :: ratio of the seed means of 1 - macro AUC, this row over the reference
\newcommand{\FtOmaRatioJcEFiveOnesixtwoOverOneeighteight}{0.98}  % experiments/FIGS/data/w2b_leg2_metrics_v2.json :: cells.*.N100000.s1.macro_auc_ovr :: sha256 9cd059d3e2f51798 :: ratio of the seed means of 1 - macro AUC, this row over the reference
\newcommand{\FtOmaRatioJcEFiveFourthreeOverOneeighteight}{1.02}  % experiments/FIGS/data/w2b_leg2_metrics_v2.json :: cells.*.N100000.s1.macro_auc_ovr :: sha256 9cd059d3e2f51798 :: ratio of the seed means of 1 - macro AUC, this row over the reference
\newcommand{\FtOmaRatioJcEFiveOnesevenOverOneeighteight}{1.07}  % experiments/FIGS/data/w2b_leg2_metrics_v2.json :: cells.*.N100000.s1.macro_auc_ovr :: sha256 9cd059d3e2f51798 :: ratio of the seed means of 1 - macro AUC, this row over the reference
\newcommand{\FtOmaRatioJcEFiveOnesixtwoMassOverOneeighteight}{1.00}  % experiments/FIGS/data/w2b_leg2_metrics_v2.json :: cells.*.N100000.s1.macro_auc_ovr :: sha256 9cd059d3e2f51798 :: ratio of the seed means of 1 - macro AUC, this row over the reference
\newcommand{\FtOmaRatioJcEFiveOnesevenMassOverOneeighteight}{1.10}  % experiments/FIGS/data/w2b_leg2_metrics_v2.json :: cells.*.N100000.s1.macro_auc_ovr :: sha256 9cd059d3e2f51798 :: ratio of the seed means of 1 - macro AUC, this row over the reference
\newcommand{\FtOmaRatioJcEFiveRandDrawOneOverOneeighteight}{1.04}  % experiments/FIGS/data/w2b_leg2_metrics_v2.json :: cells.*.N100000.s1.macro_auc_ovr :: sha256 9cd059d3e2f51798 :: ratio of the seed means of 1 - macro AUC, this row over the reference
\newcommand{\FtOmaRatioJcEFiveRandDrawTwoOverOneeighteight}{1.03}  % experiments/FIGS/data/w2b_leg2_metrics_v2.json :: cells.*.N100000.s1.macro_auc_ovr :: sha256 9cd059d3e2f51798 :: ratio of the seed means of 1 - macro AUC, this row over the reference
\newcommand{\FtOmaRatioJcEFiveRandDrawThreeOverOneeighteight}{1.05}  % experiments/FIGS/data/w2b_leg2_metrics_v2.json :: cells.*.N100000.s1.macro_auc_ovr :: sha256 9cd059d3e2f51798 :: ratio of the seed means of 1 - macro AUC, this row over the reference
\newcommand{\FtOmaRatioJcEFiveOneeighteightOverOnesixtwo}{1.02}  % experiments/FIGS/data/w2b_leg2_metrics_v2.json :: cells.*.N100000.s1.macro_auc_ovr :: sha256 9cd059d3e2f51798 :: ratio of the seed means of 1 - macro AUC, this row over the reference
\newcommand{\FtOmaRatioJcEFiveFourthreeOverOnesixtwo}{1.03}  % experiments/FIGS/data/w2b_leg2_metrics_v2.json :: cells.*.N100000.s1.macro_auc_ovr :: sha256 9cd059d3e2f51798 :: ratio of the seed means of 1 - macro AUC, this row over the reference
\newcommand{\FtOmaRatioJcEFiveOnesevenOverOnesixtwo}{1.09}  % experiments/FIGS/data/w2b_leg2_metrics_v2.json :: cells.*.N100000.s1.macro_auc_ovr :: sha256 9cd059d3e2f51798 :: ratio of the seed means of 1 - macro AUC, this row over the reference
\newcommand{\FtOmaRatioJcEFiveOnesixtwoMassOverOnesixtwo}{1.02}  % experiments/FIGS/data/w2b_leg2_metrics_v2.json :: cells.*.N100000.s1.macro_auc_ovr :: sha256 9cd059d3e2f51798 :: ratio of the seed means of 1 - macro AUC, this row over the reference
\newcommand{\FtOmaRatioJcEFiveOnesevenMassOverOnesixtwo}{1.11}  % experiments/FIGS/data/w2b_leg2_metrics_v2.json :: cells.*.N100000.s1.macro_auc_ovr :: sha256 9cd059d3e2f51798 :: ratio of the seed means of 1 - macro AUC, this row over the reference
\newcommand{\FtOmaRatioJcEFiveRandDrawOneOverOnesixtwo}{1.06}  % experiments/FIGS/data/w2b_leg2_metrics_v2.json :: cells.*.N100000.s1.macro_auc_ovr :: sha256 9cd059d3e2f51798 :: ratio of the seed means of 1 - macro AUC, this row over the reference
\newcommand{\FtOmaRatioJcEFiveRandDrawTwoOverOnesixtwo}{1.05}  % experiments/FIGS/data/w2b_leg2_metrics_v2.json :: cells.*.N100000.s1.macro_auc_ovr :: sha256 9cd059d3e2f51798 :: ratio of the seed means of 1 - macro AUC, this row over the reference
\newcommand{\FtOmaRatioJcEFiveRandDrawThreeOverOnesixtwo}{1.07}  % experiments/FIGS/data/w2b_leg2_metrics_v2.json :: cells.*.N100000.s1.macro_auc_ovr :: sha256 9cd059d3e2f51798 :: ratio of the seed means of 1 - macro AUC, this row over the reference
\newcommand{\FtOmaRatioJcESixOnesixtwoOverOneeighteight}{1.00}  % experiments/FIGS/data/w2b_leg2_metrics_v2.json :: cells.*.N1000000.s1.macro_auc_ovr :: sha256 9cd059d3e2f51798 :: ratio of the seed means of 1 - macro AUC, this row over the reference
\newcommand{\FtOmaRatioJcESixFourthreeOverOneeighteight}{1.02}  % experiments/FIGS/data/w2b_leg2_metrics_v2.json :: cells.*.N1000000.s1.macro_auc_ovr :: sha256 9cd059d3e2f51798 :: ratio of the seed means of 1 - macro AUC, this row over the reference
\newcommand{\FtOmaRatioJcESixOnesevenOverOneeighteight}{1.08}  % experiments/FIGS/data/w2b_leg2_metrics_v2.json :: cells.*.N1000000.s1.macro_auc_ovr :: sha256 9cd059d3e2f51798 :: ratio of the seed means of 1 - macro AUC, this row over the reference
\newcommand{\FtOmaRatioJcESixOnesixtwoMassOverOneeighteight}{1.01}  % experiments/FIGS/data/w2b_leg2_metrics_v2.json :: cells.*.N1000000.s1.macro_auc_ovr :: sha256 9cd059d3e2f51798 :: ratio of the seed means of 1 - macro AUC, this row over the reference
\newcommand{\FtOmaRatioJcESixOnesevenMassOverOneeighteight}{1.11}  % experiments/FIGS/data/w2b_leg2_metrics_v2.json :: cells.*.N1000000.s1.macro_auc_ovr :: sha256 9cd059d3e2f51798 :: ratio of the seed means of 1 - macro AUC, this row over the reference
\newcommand{\FtOmaRatioJcESixRandDrawOneOverOneeighteight}{1.05}  % experiments/FIGS/data/w2b_leg2_metrics_v2.json :: cells.*.N1000000.s1.macro_auc_ovr :: sha256 9cd059d3e2f51798 :: ratio of the seed means of 1 - macro AUC, this row over the reference
\newcommand{\FtOmaRatioJcESixRandDrawTwoOverOneeighteight}{1.05}  % experiments/FIGS/data/w2b_leg2_metrics_v2.json :: cells.*.N1000000.s1.macro_auc_ovr :: sha256 9cd059d3e2f51798 :: ratio of the seed means of 1 - macro AUC, this row over the reference
\newcommand{\FtOmaRatioJcESixRandDrawThreeOverOneeighteight}{1.05}  % experiments/FIGS/data/w2b_leg2_metrics_v2.json :: cells.*.N1000000.s1.macro_auc_ovr :: sha256 9cd059d3e2f51798 :: ratio of the seed means of 1 - macro AUC, this row over the reference
\newcommand{\FtOmaRatioJcESixOneeighteightOverOnesixtwo}{1.00}  % experiments/FIGS/data/w2b_leg2_metrics_v2.json :: cells.*.N1000000.s1.macro_auc_ovr :: sha256 9cd059d3e2f51798 :: ratio of the seed means of 1 - macro AUC, this row over the reference
\newcommand{\FtOmaRatioJcESixFourthreeOverOnesixtwo}{1.02}  % experiments/FIGS/data/w2b_leg2_metrics_v2.json :: cells.*.N1000000.s1.macro_auc_ovr :: sha256 9cd059d3e2f51798 :: ratio of the seed means of 1 - macro AUC, this row over the reference
\newcommand{\FtOmaRatioJcESixOnesevenOverOnesixtwo}{1.08}  % experiments/FIGS/data/w2b_leg2_metrics_v2.json :: cells.*.N1000000.s1.macro_auc_ovr :: sha256 9cd059d3e2f51798 :: ratio of the seed means of 1 - macro AUC, this row over the reference
\newcommand{\FtOmaRatioJcESixOnesixtwoMassOverOnesixtwo}{1.01}  % experiments/FIGS/data/w2b_leg2_metrics_v2.json :: cells.*.N1000000.s1.macro_auc_ovr :: sha256 9cd059d3e2f51798 :: ratio of the seed means of 1 - macro AUC, this row over the reference
\newcommand{\FtOmaRatioJcESixOnesevenMassOverOnesixtwo}{1.11}  % experiments/FIGS/data/w2b_leg2_metrics_v2.json :: cells.*.N1000000.s1.macro_auc_ovr :: sha256 9cd059d3e2f51798 :: ratio of the seed means of 1 - macro AUC, this row over the reference
\newcommand{\FtOmaRatioJcESixRandDrawOneOverOnesixtwo}{1.05}  % experiments/FIGS/data/w2b_leg2_metrics_v2.json :: cells.*.N1000000.s1.macro_auc_ovr :: sha256 9cd059d3e2f51798 :: ratio of the seed means of 1 - macro AUC, this row over the reference
\newcommand{\FtOmaRatioJcESixRandDrawTwoOverOnesixtwo}{1.05}  % experiments/FIGS/data/w2b_leg2_metrics_v2.json :: cells.*.N1000000.s1.macro_auc_ovr :: sha256 9cd059d3e2f51798 :: ratio of the seed means of 1 - macro AUC, this row over the reference
\newcommand{\FtOmaRatioJcESixRandDrawThreeOverOnesixtwo}{1.05}  % experiments/FIGS/data/w2b_leg2_metrics_v2.json :: cells.*.N1000000.s1.macro_auc_ovr :: sha256 9cd059d3e2f51798 :: ratio of the seed means of 1 - macro AUC, this row over the reference
\newcommand{\PairedFtAccJcEThreeOnesevenMinusOneeighteight}{\ensuremath{-}0.0040~[\ensuremath{-}0.0144, +0.0064]}  % experiments/FIGS/data/w2b_leg2_metrics_v2.json :: cells.{r16q1-s1,r16q1-s2,r16q1-s3,r16q1-s4,r16q1-s5,l188-s1,l188-s2,l188-s3,l188-s4,l188-s5}.N1000.s1.accuracy :: sha256 9cd059d3e2f51798 :: accuracy, coarse minus fine, paired by run [95% Student-t interval]
\newcommand{\PairedFtAccJcEFourOnesevenMinusOneeighteight}{\ensuremath{-}0.0075~[\ensuremath{-}0.0118, \ensuremath{-}0.0033]}  % experiments/FIGS/data/w2b_leg2_metrics_v2.json :: cells.{r16q1-s1,r16q1-s2,r16q1-s3,r16q1-s4,r16q1-s5,l188-s1,l188-s2,l188-s3,l188-s4,l188-s5}.N10000.s1.accuracy :: sha256 9cd059d3e2f51798 :: accuracy, coarse minus fine, paired by run [95% Student-t interval]
\newcommand{\PairedFtAccJcEFiveOnesevenMinusOneeighteight}{\ensuremath{-}0.0066~[\ensuremath{-}0.0083, \ensuremath{-}0.0049]}  % experiments/FIGS/data/w2b_leg2_metrics_v2.json :: cells.{r16q1-s1,r16q1-s2,r16q1-s3,r16q1-s4,r16q1-s5,l188-s1,l188-s2,l188-s3,l188-s4,l188-s5}.N100000.s1.accuracy :: sha256 9cd059d3e2f51798 :: accuracy, coarse minus fine, paired by run [95% Student-t interval]
\newcommand{\PairedFtAccJcESixOnesevenMinusOneeighteight}{\ensuremath{-}0.0071~[\ensuremath{-}0.0080, \ensuremath{-}0.0062]}  % experiments/FIGS/data/w2b_leg2_metrics_v2.json :: cells.{r16q1-s1,r16q1-s2,r16q1-s3,r16q1-s4,r16q1-s5,l188-s1,l188-s2,l188-s3,l188-s4,l188-s5}.N1000000.s1.accuracy :: sha256 9cd059d3e2f51798 :: accuracy, coarse minus fine, paired by run [95% Student-t interval]
\newcommand{\FtNClassesJc}{10}  % experiments/FIGS/data/w2b_leg2_metrics_v2.json :: cells.*.*.s1.n_classes_present :: sha256 9cd059d3e2f51798 :: classes of the fine-tuning task
\newcommand{\FtNTestAucJc}{500{,}000}  % experiments/FIGS/data/w2b_leg2_metrics_v2.json :: cells.*.*.s1.n_jets_auc :: sha256 9cd059d3e2f51798 :: test jets the macro AUC is computed on
\newcommand{\FtNTestAccJc}{2{,}000{,}000}  % experiments/FIGS/data/w2b_leg2_metrics_v2.json :: cells.*.*.s1.n_jets :: sha256 9cd059d3e2f51798 :: test jets the accuracy is computed on
\newcommand{\FtSizeEThree}{\ensuremath{10^{3}}}  % experiments/FIGS/data/w2b_leg1_metrics_v2.json :: cells.*.N1000 :: sha256 b56bf70be9982044 :: fine-tuning training jets
\newcommand{\FtSizeEFour}{\ensuremath{10^{4}}}  % experiments/FIGS/data/w2b_leg1_metrics_v2.json :: cells.*.N10000 :: sha256 b56bf70be9982044 :: fine-tuning training jets
\newcommand{\FtSizeEFive}{\ensuremath{10^{5}}}  % experiments/FIGS/data/w2b_leg1_metrics_v2.json :: cells.*.N100000 :: sha256 b56bf70be9982044 :: fine-tuning training jets
\newcommand{\FtSizeESix}{\ensuremath{10^{6}}}  % experiments/FIGS/data/w2b_leg1_metrics_v2.json :: cells.*.N1000000 :: sha256 b56bf70be9982044 :: fine-tuning training jets
\newcommand{\FtNRunsWord}{five}  % experiments/FIGS/data/w2b_leg1_metrics_v2.json :: cells (pretrained models per vocabulary row) :: sha256 b56bf70be9982044 :: pretraining runs per vocabulary
\newcommand{\PairedProbeBcVsRestLinearOnesevenOverOnesixtwo}{1.64~[1.52, 1.78]}  % experiments/FIGS/data/paired_v1err/probes/ratios.json :: ratios[0].ratio, ratios[0].ci95 :: sha256 14dd51402305d147 :: 17 over 162, bc_vs_rest linear 1-auc: paired geometric mean [95% interval]
\newcommand{\PairedProbeBcVsRestMlpOnesevenOverOnesixtwo}{1.50~[1.39, 1.62]}  % experiments/FIGS/data/paired_v1err/probes/ratios.json :: ratios[3].ratio, ratios[3].ci95 :: sha256 14dd51402305d147 :: 17 over 162, bc_vs_rest mlp 1-auc: paired geometric mean [95% interval]
\newcommand{\PairedProbeBvcFourprongLinearOnesixtwoOverOneeighteight}{1.01~[0.95, 1.08]}  % experiments/FIGS/data/paired_v1err/probes/ratios.json :: ratios[6].ratio, ratios[6].ci95 :: sha256 14dd51402305d147 :: 162 over 188, bvc_4prong linear 1-auc: paired geometric mean [95% interval]
\newcommand{\PairedProbeBvcFourprongLinearFourthreeOverOneeighteight}{1.10~[1.01, 1.20]}  % experiments/FIGS/data/paired_v1err/probes/ratios.json :: ratios[7].ratio, ratios[7].ci95 :: sha256 14dd51402305d147 :: 43 over 188, bvc_4prong linear 1-auc: paired geometric mean [95% interval]
\newcommand{\PairedProbeBvcFourprongLinearOnesevenOverOneeighteight}{4.5~[3.5, 5.7]}  % experiments/FIGS/data/paired_v1err/probes/ratios.json :: ratios[8].ratio, ratios[8].ci95 :: sha256 14dd51402305d147 :: 17 over 188, bvc_4prong linear 1-auc: paired geometric mean [95% interval]
\newcommand{\PairedProbeBvcFourprongLinearFourthreeOverOnesixtwo}{1.09~[0.98, 1.20]}  % experiments/FIGS/data/paired_v1err/probes/ratios.json :: ratios[9].ratio, ratios[9].ci95 :: sha256 14dd51402305d147 :: 43 over 162, bvc_4prong linear 1-auc: paired geometric mean [95% interval]
\newcommand{\PairedProbeBvcFourprongLinearOnesevenOverOnesixtwo}{4.4~[3.4, 5.6]}  % experiments/FIGS/data/paired_v1err/probes/ratios.json :: ratios[10].ratio, ratios[10].ci95 :: sha256 14dd51402305d147 :: 17 over 162, bvc_4prong linear 1-auc: paired geometric mean [95% interval]
\newcommand{\PairedProbeBvcFourprongLinearOnesevenOverFourthree}{4.0~[3.1, 5.4]}  % experiments/FIGS/data/paired_v1err/probes/ratios.json :: ratios[11].ratio, ratios[11].ci95 :: sha256 14dd51402305d147 :: 17 over 43, bvc_4prong linear 1-auc: paired geometric mean [95% interval]
\newcommand{\PairedProbeBvcFourprongLinearOnesixtwoMassOverOnesixtwo}{1.05~[0.97, 1.13]}  % experiments/FIGS/data/paired_v1err/probes/ratios.json :: ratios[12].ratio, ratios[12].ci95 :: sha256 14dd51402305d147 :: 162+mass over 162, bvc_4prong linear 1-auc: paired geometric mean [95% interval]
\newcommand{\PairedProbeBvcFourprongLinearOnesevenMassOverOneseven}{2.1~[1.6, 2.7]}  % experiments/FIGS/data/paired_v1err/probes/ratios.json :: ratios[13].ratio, ratios[13].ci95 :: sha256 14dd51402305d147 :: 17+mass over 17, bvc_4prong linear 1-auc: paired geometric mean [95% interval]
\newcommand{\PairedProbeBvcFourprongLinearRandomOverOneeighteight}{1.54~[1.37, 1.74]}  % experiments/FIGS/data/paired_v1err/probes/ratios.json :: ratios[14].ratio, ratios[14].ci95 :: sha256 14dd51402305d147 :: random over 188, bvc_4prong linear 1-auc: paired geometric mean [95% interval]
\newcommand{\PairedProbeBvcFourprongLinearRandomOverOnesixtwo}{1.55~[1.36, 1.77]}  % experiments/FIGS/data/paired_v1err/probes/ratios.json :: ratios[15].ratio, ratios[15].ci95 :: sha256 14dd51402305d147 :: random over 162, bvc_4prong linear 1-auc: paired geometric mean [95% interval]
\newcommand{\PairedProbeBvcFourprongLinearRandomOverFourthree}{1.39~[1.27, 1.52]}  % experiments/FIGS/data/paired_v1err/probes/ratios.json :: ratios[16].ratio, ratios[16].ci95 :: sha256 14dd51402305d147 :: random over 43, bvc_4prong linear 1-auc: paired geometric mean [95% interval]
\newcommand{\PairedProbeBvcFourprongLinearRandomOverOneseven}{0.34~[0.20, 0.58]}  % experiments/FIGS/data/paired_v1err/probes/ratios.json :: ratios[17].ratio, ratios[17].ci95 :: sha256 14dd51402305d147 :: random over 17, bvc_4prong linear 1-auc: paired geometric mean [95% interval]
\newcommand{\PairedProbeBvcFourprongLinearRandomDrawTwoOverRandomDrawOne}{1.05~[0.94, 1.18]}  % experiments/FIGS/data/paired_v1err/probes/ratios.json :: ratios[18].ratio, ratios[18].ci95 :: sha256 14dd51402305d147 :: random draw 2 over random draw 1, bvc_4prong linear 1-auc: paired geometric mean [95% interval]
\newcommand{\PairedProbeBvcFourprongLinearRandomDrawThreeOverRandomDrawOne}{0.95~[0.85, 1.06]}  % experiments/FIGS/data/paired_v1err/probes/ratios.json :: ratios[19].ratio, ratios[19].ci95 :: sha256 14dd51402305d147 :: random draw 3 over random draw 1, bvc_4prong linear 1-auc: paired geometric mean [95% interval]
\newcommand{\PairedProbeBvcFourprongLinearRandomDrawThreeOverRandomDrawTwo}{0.90~[0.79, 1.02]}  % experiments/FIGS/data/paired_v1err/probes/ratios.json :: ratios[20].ratio, ratios[20].ci95 :: sha256 14dd51402305d147 :: random draw 3 over random draw 2, bvc_4prong linear 1-auc: paired geometric mean [95% interval]
\newcommand{\PairedProbeBvcFourprongMlpOnesixtwoOverOneeighteight}{1.00~[0.93, 1.07]}  % experiments/FIGS/data/paired_v1err/probes/ratios.json :: ratios[66].ratio, ratios[66].ci95 :: sha256 14dd51402305d147 :: 162 over 188, bvc_4prong mlp 1-auc: paired geometric mean [95% interval]
\newcommand{\PairedProbeBvcFourprongMlpFourthreeOverOneeighteight}{1.12~[1.03, 1.22]}  % experiments/FIGS/data/paired_v1err/probes/ratios.json :: ratios[67].ratio, ratios[67].ci95 :: sha256 14dd51402305d147 :: 43 over 188, bvc_4prong mlp 1-auc: paired geometric mean [95% interval]
\newcommand{\PairedProbeBvcFourprongMlpOnesevenOverOneeighteight}{3.2~[2.6, 4.1]}  % experiments/FIGS/data/paired_v1err/probes/ratios.json :: ratios[68].ratio, ratios[68].ci95 :: sha256 14dd51402305d147 :: 17 over 188, bvc_4prong mlp 1-auc: paired geometric mean [95% interval]
\newcommand{\PairedProbeBvcFourprongMlpFourthreeOverOnesixtwo}{1.12~[1.01, 1.26]}  % experiments/FIGS/data/paired_v1err/probes/ratios.json :: ratios[69].ratio, ratios[69].ci95 :: sha256 14dd51402305d147 :: 43 over 162, bvc_4prong mlp 1-auc: paired geometric mean [95% interval]
\newcommand{\PairedProbeBvcFourprongMlpOnesevenOverOnesixtwo}{3.3~[2.6, 4.0]}  % experiments/FIGS/data/paired_v1err/probes/ratios.json :: ratios[70].ratio, ratios[70].ci95 :: sha256 14dd51402305d147 :: 17 over 162, bvc_4prong mlp 1-auc: paired geometric mean [95% interval]
\newcommand{\PairedProbeBvcFourprongMlpOnesevenOverFourthree}{2.9~[2.3, 3.7]}  % experiments/FIGS/data/paired_v1err/probes/ratios.json :: ratios[71].ratio, ratios[71].ci95 :: sha256 14dd51402305d147 :: 17 over 43, bvc_4prong mlp 1-auc: paired geometric mean [95% interval]
\newcommand{\PairedProbeBvcFourprongMlpOnesixtwoMassOverOnesixtwo}{1.07~[0.99, 1.16]}  % experiments/FIGS/data/paired_v1err/probes/ratios.json :: ratios[72].ratio, ratios[72].ci95 :: sha256 14dd51402305d147 :: 162+mass over 162, bvc_4prong mlp 1-auc: paired geometric mean [95% interval]
\newcommand{\PairedProbeBvcFourprongMlpOnesevenMassOverOneseven}{1.97~[1.46, 2.67]}  % experiments/FIGS/data/paired_v1err/probes/ratios.json :: ratios[73].ratio, ratios[73].ci95 :: sha256 14dd51402305d147 :: 17+mass over 17, bvc_4prong mlp 1-auc: paired geometric mean [95% interval]
\newcommand{\PairedProbeBvcFourprongMlpRandomOverOneeighteight}{1.47~[1.30, 1.65]}  % experiments/FIGS/data/paired_v1err/probes/ratios.json :: ratios[74].ratio, ratios[74].ci95 :: sha256 14dd51402305d147 :: random over 188, bvc_4prong mlp 1-auc: paired geometric mean [95% interval]
\newcommand{\PairedProbeBvcFourprongMlpRandomOverOnesixtwo}{1.49~[1.24, 1.79]}  % experiments/FIGS/data/paired_v1err/probes/ratios.json :: ratios[75].ratio, ratios[75].ci95 :: sha256 14dd51402305d147 :: random over 162, bvc_4prong mlp 1-auc: paired geometric mean [95% interval]
\newcommand{\PairedProbeBvcFourprongMlpRandomOverFourthree}{1.28~[1.16, 1.43]}  % experiments/FIGS/data/paired_v1err/probes/ratios.json :: ratios[76].ratio, ratios[76].ci95 :: sha256 14dd51402305d147 :: random over 43, bvc_4prong mlp 1-auc: paired geometric mean [95% interval]
\newcommand{\PairedProbeBvcFourprongMlpRandomOverOneseven}{0.43~[0.28, 0.67]}  % experiments/FIGS/data/paired_v1err/probes/ratios.json :: ratios[77].ratio, ratios[77].ci95 :: sha256 14dd51402305d147 :: random over 17, bvc_4prong mlp 1-auc: paired geometric mean [95% interval]
\newcommand{\PairedProbeBvcFourprongMlpRandomDrawTwoOverRandomDrawOne}{1.00~[0.90, 1.10]}  % experiments/FIGS/data/paired_v1err/probes/ratios.json :: ratios[78].ratio, ratios[78].ci95 :: sha256 14dd51402305d147 :: random draw 2 over random draw 1, bvc_4prong mlp 1-auc: paired geometric mean [95% interval]
\newcommand{\PairedProbeBvcFourprongMlpRandomDrawThreeOverRandomDrawOne}{0.89~[0.79, 0.99]}  % experiments/FIGS/data/paired_v1err/probes/ratios.json :: ratios[79].ratio, ratios[79].ci95 :: sha256 14dd51402305d147 :: random draw 3 over random draw 1, bvc_4prong mlp 1-auc: paired geometric mean [95% interval]
\newcommand{\PairedProbeBvcFourprongMlpRandomDrawThreeOverRandomDrawTwo}{0.89~[0.79, 0.99]}  % experiments/FIGS/data/paired_v1err/probes/ratios.json :: ratios[80].ratio, ratios[80].ci95 :: sha256 14dd51402305d147 :: random draw 3 over random draw 2, bvc_4prong mlp 1-auc: paired geometric mean [95% interval]
\newcommand{\PairedProbeBvcQcdLinearOnesixtwoOverOneeighteight}{1.15~[1.07, 1.23]}  % experiments/FIGS/data/paired_v1err/probes/ratios.json :: ratios[126].ratio, ratios[126].ci95 :: sha256 14dd51402305d147 :: 162 over 188, bvc_qcd linear 1-auc: paired geometric mean [95% interval]
\newcommand{\PairedProbeBvcQcdLinearFourthreeOverOneeighteight}{1.16~[1.08, 1.25]}  % experiments/FIGS/data/paired_v1err/probes/ratios.json :: ratios[127].ratio, ratios[127].ci95 :: sha256 14dd51402305d147 :: 43 over 188, bvc_qcd linear 1-auc: paired geometric mean [95% interval]
\newcommand{\PairedProbeBvcQcdLinearOnesevenOverOneeighteight}{3.0~[2.5, 3.7]}  % experiments/FIGS/data/paired_v1err/probes/ratios.json :: ratios[128].ratio, ratios[128].ci95 :: sha256 14dd51402305d147 :: 17 over 188, bvc_qcd linear 1-auc: paired geometric mean [95% interval]
\newcommand{\PairedProbeBvcQcdLinearFourthreeOverOnesixtwo}{1.01~[0.96, 1.07]}  % experiments/FIGS/data/paired_v1err/probes/ratios.json :: ratios[129].ratio, ratios[129].ci95 :: sha256 14dd51402305d147 :: 43 over 162, bvc_qcd linear 1-auc: paired geometric mean [95% interval]
\newcommand{\PairedProbeBvcQcdLinearOnesevenOverOnesixtwo}{2.6~[2.2, 3.1]}  % experiments/FIGS/data/paired_v1err/probes/ratios.json :: ratios[130].ratio, ratios[130].ci95 :: sha256 14dd51402305d147 :: 17 over 162, bvc_qcd linear 1-auc: paired geometric mean [95% interval]
\newcommand{\PairedProbeBvcQcdLinearOnesevenOverFourthree}{2.6~[2.2, 3.0]}  % experiments/FIGS/data/paired_v1err/probes/ratios.json :: ratios[131].ratio, ratios[131].ci95 :: sha256 14dd51402305d147 :: 17 over 43, bvc_qcd linear 1-auc: paired geometric mean [95% interval]
\newcommand{\PairedProbeBvcQcdLinearOnesixtwoMassOverOnesixtwo}{1.060~[1.002, 1.123]}  % experiments/FIGS/data/paired_v1err/probes/ratios.json :: ratios[132].ratio, ratios[132].ci95 :: sha256 14dd51402305d147 :: 162+mass over 162, bvc_qcd linear 1-auc: paired geometric mean [95% interval]
\newcommand{\PairedProbeBvcQcdLinearOnesevenMassOverOneseven}{1.89~[1.52, 2.34]}  % experiments/FIGS/data/paired_v1err/probes/ratios.json :: ratios[133].ratio, ratios[133].ci95 :: sha256 14dd51402305d147 :: 17+mass over 17, bvc_qcd linear 1-auc: paired geometric mean [95% interval]
\newcommand{\PairedProbeBvcQcdMlpOnesixtwoOverOneeighteight}{1.11~[1.04, 1.19]}  % experiments/FIGS/data/paired_v1err/probes/ratios.json :: ratios[158].ratio, ratios[158].ci95 :: sha256 14dd51402305d147 :: 162 over 188, bvc_qcd mlp 1-auc: paired geometric mean [95% interval]
\newcommand{\PairedProbeBvcQcdMlpFourthreeOverOneeighteight}{1.13~[1.05, 1.21]}  % experiments/FIGS/data/paired_v1err/probes/ratios.json :: ratios[159].ratio, ratios[159].ci95 :: sha256 14dd51402305d147 :: 43 over 188, bvc_qcd mlp 1-auc: paired geometric mean [95% interval]
\newcommand{\PairedProbeBvcQcdMlpOnesevenOverOneeighteight}{2.6~[2.2, 3.0]}  % experiments/FIGS/data/paired_v1err/probes/ratios.json :: ratios[160].ratio, ratios[160].ci95 :: sha256 14dd51402305d147 :: 17 over 188, bvc_qcd mlp 1-auc: paired geometric mean [95% interval]
\newcommand{\PairedProbeBvcQcdMlpFourthreeOverOnesixtwo}{1.01~[0.96, 1.06]}  % experiments/FIGS/data/paired_v1err/probes/ratios.json :: ratios[161].ratio, ratios[161].ci95 :: sha256 14dd51402305d147 :: 43 over 162, bvc_qcd mlp 1-auc: paired geometric mean [95% interval]
\newcommand{\PairedProbeBvcQcdMlpOnesevenOverOnesixtwo}{2.3~[2.0, 2.6]}  % experiments/FIGS/data/paired_v1err/probes/ratios.json :: ratios[162].ratio, ratios[162].ci95 :: sha256 14dd51402305d147 :: 17 over 162, bvc_qcd mlp 1-auc: paired geometric mean [95% interval]
\newcommand{\PairedProbeBvcQcdMlpOnesevenOverFourthree}{2.3~[2.0, 2.6]}  % experiments/FIGS/data/paired_v1err/probes/ratios.json :: ratios[163].ratio, ratios[163].ci95 :: sha256 14dd51402305d147 :: 17 over 43, bvc_qcd mlp 1-auc: paired geometric mean [95% interval]
\newcommand{\PairedProbeBvcQcdMlpOnesixtwoMassOverOnesixtwo}{1.06~[1.01, 1.12]}  % experiments/FIGS/data/paired_v1err/probes/ratios.json :: ratios[164].ratio, ratios[164].ci95 :: sha256 14dd51402305d147 :: 162+mass over 162, bvc_qcd mlp 1-auc: paired geometric mean [95% interval]
\newcommand{\PairedProbeBvcQcdMlpOnesevenMassOverOneseven}{1.64~[1.33, 2.02]}  % experiments/FIGS/data/paired_v1err/probes/ratios.json :: ratios[165].ratio, ratios[165].ci95 :: sha256 14dd51402305d147 :: 17+mass over 17, bvc_qcd mlp 1-auc: paired geometric mean [95% interval]
\newcommand{\PairedProbeBvcResonantLinearOnesixtwoOverOneeighteight}{1.02~[0.96, 1.08]}  % experiments/FIGS/data/paired_v1err/probes/ratios.json :: ratios[190].ratio, ratios[190].ci95 :: sha256 14dd51402305d147 :: 162 over 188, bvc_resonant linear 1-auc: paired geometric mean [95% interval]
\newcommand{\PairedProbeBvcResonantLinearFourthreeOverOneeighteight}{1.09~[1.03, 1.15]}  % experiments/FIGS/data/paired_v1err/probes/ratios.json :: ratios[191].ratio, ratios[191].ci95 :: sha256 14dd51402305d147 :: 43 over 188, bvc_resonant linear 1-auc: paired geometric mean [95% interval]
\newcommand{\PairedProbeBvcResonantLinearOnesevenOverOneeighteight}{4.0~[3.3, 4.9]}  % experiments/FIGS/data/paired_v1err/probes/ratios.json :: ratios[192].ratio, ratios[192].ci95 :: sha256 14dd51402305d147 :: 17 over 188, bvc_resonant linear 1-auc: paired geometric mean [95% interval]
\newcommand{\PairedProbeBvcResonantLinearFourthreeOverOnesixtwo}{1.07~[1.02, 1.12]}  % experiments/FIGS/data/paired_v1err/probes/ratios.json :: ratios[193].ratio, ratios[193].ci95 :: sha256 14dd51402305d147 :: 43 over 162, bvc_resonant linear 1-auc: paired geometric mean [95% interval]
\newcommand{\PairedProbeBvcResonantLinearOnesevenOverOnesixtwo}{3.9~[3.1, 4.9]}  % experiments/FIGS/data/paired_v1err/probes/ratios.json :: ratios[194].ratio, ratios[194].ci95 :: sha256 14dd51402305d147 :: 17 over 162, bvc_resonant linear 1-auc: paired geometric mean [95% interval]
\newcommand{\PairedProbeBvcResonantLinearOnesevenOverFourthree}{3.7~[3.0, 4.5]}  % experiments/FIGS/data/paired_v1err/probes/ratios.json :: ratios[195].ratio, ratios[195].ci95 :: sha256 14dd51402305d147 :: 17 over 43, bvc_resonant linear 1-auc: paired geometric mean [95% interval]
\newcommand{\PairedProbeBvcResonantLinearOnesixtwoMassOverOnesixtwo}{1.11~[1.06, 1.16]}  % experiments/FIGS/data/paired_v1err/probes/ratios.json :: ratios[196].ratio, ratios[196].ci95 :: sha256 14dd51402305d147 :: 162+mass over 162, bvc_resonant linear 1-auc: paired geometric mean [95% interval]
\newcommand{\PairedProbeBvcResonantLinearOnesevenMassOverOneseven}{3.0~[1.9, 4.7]}  % experiments/FIGS/data/paired_v1err/probes/ratios.json :: ratios[197].ratio, ratios[197].ci95 :: sha256 14dd51402305d147 :: 17+mass over 17, bvc_resonant linear 1-auc: paired geometric mean [95% interval]
\newcommand{\PairedProbeBvcResonantMlpOnesixtwoOverOneeighteight}{1.04~[0.98, 1.12]}  % experiments/FIGS/data/paired_v1err/probes/ratios.json :: ratios[222].ratio, ratios[222].ci95 :: sha256 14dd51402305d147 :: 162 over 188, bvc_resonant mlp 1-auc: paired geometric mean [95% interval]
\newcommand{\PairedProbeBvcResonantMlpFourthreeOverOneeighteight}{1.14~[1.08, 1.21]}  % experiments/FIGS/data/paired_v1err/probes/ratios.json :: ratios[223].ratio, ratios[223].ci95 :: sha256 14dd51402305d147 :: 43 over 188, bvc_resonant mlp 1-auc: paired geometric mean [95% interval]
\newcommand{\PairedProbeBvcResonantMlpOnesevenOverOneeighteight}{3.3~[2.8, 3.8]}  % experiments/FIGS/data/paired_v1err/probes/ratios.json :: ratios[224].ratio, ratios[224].ci95 :: sha256 14dd51402305d147 :: 17 over 188, bvc_resonant mlp 1-auc: paired geometric mean [95% interval]
\newcommand{\PairedProbeBvcResonantMlpFourthreeOverOnesixtwo}{1.09~[1.04, 1.14]}  % experiments/FIGS/data/paired_v1err/probes/ratios.json :: ratios[225].ratio, ratios[225].ci95 :: sha256 14dd51402305d147 :: 43 over 162, bvc_resonant mlp 1-auc: paired geometric mean [95% interval]
\newcommand{\PairedProbeBvcResonantMlpOnesevenOverOnesixtwo}{3.1~[2.6, 3.8]}  % experiments/FIGS/data/paired_v1err/probes/ratios.json :: ratios[226].ratio, ratios[226].ci95 :: sha256 14dd51402305d147 :: 17 over 162, bvc_resonant mlp 1-auc: paired geometric mean [95% interval]
\newcommand{\PairedProbeBvcResonantMlpOnesevenOverFourthree}{2.9~[2.4, 3.4]}  % experiments/FIGS/data/paired_v1err/probes/ratios.json :: ratios[227].ratio, ratios[227].ci95 :: sha256 14dd51402305d147 :: 17 over 43, bvc_resonant mlp 1-auc: paired geometric mean [95% interval]
\newcommand{\PairedProbeBvcResonantMlpOnesixtwoMassOverOnesixtwo}{1.09~[1.04, 1.14]}  % experiments/FIGS/data/paired_v1err/probes/ratios.json :: ratios[228].ratio, ratios[228].ci95 :: sha256 14dd51402305d147 :: 162+mass over 162, bvc_resonant mlp 1-auc: paired geometric mean [95% interval]
\newcommand{\PairedProbeBvcResonantMlpOnesevenMassOverOneseven}{1.95~[1.45, 2.61]}  % experiments/FIGS/data/paired_v1err/probes/ratios.json :: ratios[229].ratio, ratios[229].ci95 :: sha256 14dd51402305d147 :: 17+mass over 17, bvc_resonant mlp 1-auc: paired geometric mean [95% interval]
\newcommand{\PairedProbeEeVsMmLinearOnesevenOverFourthree}{2.3~[0.8, 6.5]}  % experiments/FIGS/data/paired_v1err/probes/ratios.json :: ratios[259].ratio, ratios[259].ci95 :: sha256 14dd51402305d147 :: 17 over 43, ee_vs_mm linear 1-auc: paired geometric mean [95% interval]
\newcommand{\PairedProbeEeVsMmLinearOnesevenMassOverOneseven}{0.23~[0.08, 0.64]}  % experiments/FIGS/data/paired_v1err/probes/ratios.json :: ratios[261].ratio, ratios[261].ci95 :: sha256 14dd51402305d147 :: 17+mass over 17, ee_vs_mm linear 1-auc: paired geometric mean [95% interval]
\newcommand{\PairedProbeEeVsMmMlpOnesevenOverFourthree}{2.8~[0.8, 9.0]}  % experiments/FIGS/data/paired_v1err/probes/ratios.json :: ratios[291].ratio, ratios[291].ci95 :: sha256 14dd51402305d147 :: 17 over 43, ee_vs_mm mlp 1-auc: paired geometric mean [95% interval]
\newcommand{\PairedProbeEeVsMmMlpOnesevenMassOverOneseven}{0.172~[0.026, 1.125]}  % experiments/FIGS/data/paired_v1err/probes/ratios.json :: ratios[293].ratio, ratios[293].ci95 :: sha256 14dd51402305d147 :: 17+mass over 17, ee_vs_mm mlp 1-auc: paired geometric mean [95% interval]
\newcommand{\PairedProbeRetainedTopologyLinearOnesixtwoOverOneeighteight}{1.00~[0.92, 1.08]}  % experiments/FIGS/data/paired_v1err/probes/ratios.json :: ratios[318].ratio, ratios[318].ci95 :: sha256 14dd51402305d147 :: 162 over 188, retained_topology linear 1-auc: paired geometric mean [95% interval]
\newcommand{\PairedProbeRetainedTopologyLinearFourthreeOverOneeighteight}{1.24~[1.15, 1.34]}  % experiments/FIGS/data/paired_v1err/probes/ratios.json :: ratios[319].ratio, ratios[319].ci95 :: sha256 14dd51402305d147 :: 43 over 188, retained_topology linear 1-auc: paired geometric mean [95% interval]
\newcommand{\PairedProbeRetainedTopologyLinearOnesevenOverOneeighteight}{1.54~[1.35, 1.75]}  % experiments/FIGS/data/paired_v1err/probes/ratios.json :: ratios[320].ratio, ratios[320].ci95 :: sha256 14dd51402305d147 :: 17 over 188, retained_topology linear 1-auc: paired geometric mean [95% interval]
\newcommand{\PairedProbeRetainedTopologyLinearFourthreeOverOnesixtwo}{1.24~[1.14, 1.34]}  % experiments/FIGS/data/paired_v1err/probes/ratios.json :: ratios[321].ratio, ratios[321].ci95 :: sha256 14dd51402305d147 :: 43 over 162, retained_topology linear 1-auc: paired geometric mean [95% interval]
\newcommand{\PairedProbeRetainedTopologyLinearOnesevenOverOnesixtwo}{1.54~[1.33, 1.78]}  % experiments/FIGS/data/paired_v1err/probes/ratios.json :: ratios[322].ratio, ratios[322].ci95 :: sha256 14dd51402305d147 :: 17 over 162, retained_topology linear 1-auc: paired geometric mean [95% interval]
\newcommand{\PairedProbeRetainedTopologyLinearOnesevenOverFourthree}{1.24~[1.08, 1.42]}  % experiments/FIGS/data/paired_v1err/probes/ratios.json :: ratios[323].ratio, ratios[323].ci95 :: sha256 14dd51402305d147 :: 17 over 43, retained_topology linear 1-auc: paired geometric mean [95% interval]
\newcommand{\PairedProbeRetainedTopologyLinearOnesixtwoMassOverOnesixtwo}{1.06~[0.99, 1.14]}  % experiments/FIGS/data/paired_v1err/probes/ratios.json :: ratios[324].ratio, ratios[324].ci95 :: sha256 14dd51402305d147 :: 162+mass over 162, retained_topology linear 1-auc: paired geometric mean [95% interval]
\newcommand{\PairedProbeRetainedTopologyLinearOnesevenMassOverOneseven}{1.12~[0.97, 1.28]}  % experiments/FIGS/data/paired_v1err/probes/ratios.json :: ratios[325].ratio, ratios[325].ci95 :: sha256 14dd51402305d147 :: 17+mass over 17, retained_topology linear 1-auc: paired geometric mean [95% interval]
\newcommand{\PairedProbeRetainedTopologyMlpOnesixtwoOverOneeighteight}{1.01~[0.87, 1.17]}  % experiments/FIGS/data/paired_v1err/probes/ratios.json :: ratios[350].ratio, ratios[350].ci95 :: sha256 14dd51402305d147 :: 162 over 188, retained_topology mlp 1-auc: paired geometric mean [95% interval]
\newcommand{\PairedProbeRetainedTopologyMlpFourthreeOverOneeighteight}{1.18~[1.03, 1.36]}  % experiments/FIGS/data/paired_v1err/probes/ratios.json :: ratios[351].ratio, ratios[351].ci95 :: sha256 14dd51402305d147 :: 43 over 188, retained_topology mlp 1-auc: paired geometric mean [95% interval]
\newcommand{\PairedProbeRetainedTopologyMlpOnesevenOverOneeighteight}{1.52~[1.34, 1.72]}  % experiments/FIGS/data/paired_v1err/probes/ratios.json :: ratios[352].ratio, ratios[352].ci95 :: sha256 14dd51402305d147 :: 17 over 188, retained_topology mlp 1-auc: paired geometric mean [95% interval]
\newcommand{\PairedProbeRetainedTopologyMlpFourthreeOverOnesixtwo}{1.17~[1.07, 1.28]}  % experiments/FIGS/data/paired_v1err/probes/ratios.json :: ratios[353].ratio, ratios[353].ci95 :: sha256 14dd51402305d147 :: 43 over 162, retained_topology mlp 1-auc: paired geometric mean [95% interval]
\newcommand{\PairedProbeRetainedTopologyMlpOnesevenOverOnesixtwo}{1.50~[1.29, 1.75]}  % experiments/FIGS/data/paired_v1err/probes/ratios.json :: ratios[354].ratio, ratios[354].ci95 :: sha256 14dd51402305d147 :: 17 over 162, retained_topology mlp 1-auc: paired geometric mean [95% interval]
\newcommand{\PairedProbeRetainedTopologyMlpOnesevenOverFourthree}{1.28~[1.12, 1.47]}  % experiments/FIGS/data/paired_v1err/probes/ratios.json :: ratios[355].ratio, ratios[355].ci95 :: sha256 14dd51402305d147 :: 17 over 43, retained_topology mlp 1-auc: paired geometric mean [95% interval]
\newcommand{\PairedProbeRetainedTopologyMlpOnesixtwoMassOverOnesixtwo}{1.04~[0.92, 1.17]}  % experiments/FIGS/data/paired_v1err/probes/ratios.json :: ratios[356].ratio, ratios[356].ci95 :: sha256 14dd51402305d147 :: 162+mass over 162, retained_topology mlp 1-auc: paired geometric mean [95% interval]
\newcommand{\PairedProbeRetainedTopologyMlpOnesevenMassOverOneseven}{1.11~[0.99, 1.23]}  % experiments/FIGS/data/paired_v1err/probes/ratios.json :: ratios[357].ratio, ratios[357].ci95 :: sha256 14dd51402305d147 :: 17+mass over 17, retained_topology mlp 1-auc: paired geometric mean [95% interval]
\newcommand{\PairedProbeVisibleContentLinearOnesixtwoOverOneeighteight}{0.95~[0.88, 1.03]}  % experiments/FIGS/data/paired_v1err/probes/ratios.json :: ratios[382].ratio, ratios[382].ci95 :: sha256 14dd51402305d147 :: 162 over 188, visible_content linear 1-auc: paired geometric mean [95% interval]
\newcommand{\PairedProbeVisibleContentLinearFourthreeOverOneeighteight}{0.99~[0.85, 1.15]}  % experiments/FIGS/data/paired_v1err/probes/ratios.json :: ratios[383].ratio, ratios[383].ci95 :: sha256 14dd51402305d147 :: 43 over 188, visible_content linear 1-auc: paired geometric mean [95% interval]
\newcommand{\PairedProbeVisibleContentLinearOnesevenOverOneeighteight}{2.8~[2.1, 3.7]}  % experiments/FIGS/data/paired_v1err/probes/ratios.json :: ratios[384].ratio, ratios[384].ci95 :: sha256 14dd51402305d147 :: 17 over 188, visible_content linear 1-auc: paired geometric mean [95% interval]
\newcommand{\PairedProbeVisibleContentLinearFourthreeOverOnesixtwo}{1.04~[0.89, 1.20]}  % experiments/FIGS/data/paired_v1err/probes/ratios.json :: ratios[385].ratio, ratios[385].ci95 :: sha256 14dd51402305d147 :: 43 over 162, visible_content linear 1-auc: paired geometric mean [95% interval]
\newcommand{\PairedProbeVisibleContentLinearOnesevenOverOnesixtwo}{2.9~[2.3, 3.8]}  % experiments/FIGS/data/paired_v1err/probes/ratios.json :: ratios[386].ratio, ratios[386].ci95 :: sha256 14dd51402305d147 :: 17 over 162, visible_content linear 1-auc: paired geometric mean [95% interval]
\newcommand{\PairedProbeVisibleContentLinearOnesevenOverFourthree}{2.8~[2.2, 3.7]}  % experiments/FIGS/data/paired_v1err/probes/ratios.json :: ratios[387].ratio, ratios[387].ci95 :: sha256 14dd51402305d147 :: 17 over 43, visible_content linear 1-auc: paired geometric mean [95% interval]
\newcommand{\PairedProbeVisibleContentLinearOnesixtwoMassOverOnesixtwo}{1.10~[0.98, 1.23]}  % experiments/FIGS/data/paired_v1err/probes/ratios.json :: ratios[388].ratio, ratios[388].ci95 :: sha256 14dd51402305d147 :: 162+mass over 162, visible_content linear 1-auc: paired geometric mean [95% interval]
\newcommand{\PairedProbeVisibleContentLinearOnesevenMassOverOneseven}{1.44~[1.26, 1.65]}  % experiments/FIGS/data/paired_v1err/probes/ratios.json :: ratios[389].ratio, ratios[389].ci95 :: sha256 14dd51402305d147 :: 17+mass over 17, visible_content linear 1-auc: paired geometric mean [95% interval]
\newcommand{\PairedProbeVisibleContentLinearRandomOverOneeighteight}{1.31~[1.01, 1.70]}  % experiments/FIGS/data/paired_v1err/probes/ratios.json :: ratios[390].ratio, ratios[390].ci95 :: sha256 14dd51402305d147 :: random over 188, visible_content linear 1-auc: paired geometric mean [95% interval]
\newcommand{\PairedProbeVisibleContentLinearRandomOverOnesixtwo}{1.41~[1.17, 1.70]}  % experiments/FIGS/data/paired_v1err/probes/ratios.json :: ratios[391].ratio, ratios[391].ci95 :: sha256 14dd51402305d147 :: random over 162, visible_content linear 1-auc: paired geometric mean [95% interval]
\newcommand{\PairedProbeVisibleContentLinearRandomOverFourthree}{1.33~[1.12, 1.57]}  % experiments/FIGS/data/paired_v1err/probes/ratios.json :: ratios[392].ratio, ratios[392].ci95 :: sha256 14dd51402305d147 :: random over 43, visible_content linear 1-auc: paired geometric mean [95% interval]
\newcommand{\PairedProbeVisibleContentLinearRandomOverOneseven}{0.48~[0.38, 0.60]}  % experiments/FIGS/data/paired_v1err/probes/ratios.json :: ratios[393].ratio, ratios[393].ci95 :: sha256 14dd51402305d147 :: random over 17, visible_content linear 1-auc: paired geometric mean [95% interval]
\newcommand{\PairedProbeVisibleContentLinearRandomDrawTwoOverRandomDrawOne}{0.79~[0.64, 0.98]}  % experiments/FIGS/data/paired_v1err/probes/ratios.json :: ratios[394].ratio, ratios[394].ci95 :: sha256 14dd51402305d147 :: random draw 2 over random draw 1, visible_content linear 1-auc: paired geometric mean [95% interval]
\newcommand{\PairedProbeVisibleContentLinearRandomDrawThreeOverRandomDrawOne}{0.86~[0.71, 1.05]}  % experiments/FIGS/data/paired_v1err/probes/ratios.json :: ratios[395].ratio, ratios[395].ci95 :: sha256 14dd51402305d147 :: random draw 3 over random draw 1, visible_content linear 1-auc: paired geometric mean [95% interval]
\newcommand{\PairedProbeVisibleContentLinearRandomDrawThreeOverRandomDrawTwo}{1.09~[0.88, 1.35]}  % experiments/FIGS/data/paired_v1err/probes/ratios.json :: ratios[396].ratio, ratios[396].ci95 :: sha256 14dd51402305d147 :: random draw 3 over random draw 2, visible_content linear 1-auc: paired geometric mean [95% interval]
\newcommand{\PairedProbeVisibleContentMlpOnesixtwoOverOneeighteight}{0.92~[0.85, 1.01]}  % experiments/FIGS/data/paired_v1err/probes/ratios.json :: ratios[442].ratio, ratios[442].ci95 :: sha256 14dd51402305d147 :: 162 over 188, visible_content mlp 1-auc: paired geometric mean [95% interval]
\newcommand{\PairedProbeVisibleContentMlpFourthreeOverOneeighteight}{0.93~[0.80, 1.07]}  % experiments/FIGS/data/paired_v1err/probes/ratios.json :: ratios[443].ratio, ratios[443].ci95 :: sha256 14dd51402305d147 :: 43 over 188, visible_content mlp 1-auc: paired geometric mean [95% interval]
\newcommand{\PairedProbeVisibleContentMlpOnesevenOverOneeighteight}{2.3~[1.7, 3.1]}  % experiments/FIGS/data/paired_v1err/probes/ratios.json :: ratios[444].ratio, ratios[444].ci95 :: sha256 14dd51402305d147 :: 17 over 188, visible_content mlp 1-auc: paired geometric mean [95% interval]
\newcommand{\PairedProbeVisibleContentMlpFourthreeOverOnesixtwo}{1.01~[0.89, 1.14]}  % experiments/FIGS/data/paired_v1err/probes/ratios.json :: ratios[445].ratio, ratios[445].ci95 :: sha256 14dd51402305d147 :: 43 over 162, visible_content mlp 1-auc: paired geometric mean [95% interval]
\newcommand{\PairedProbeVisibleContentMlpOnesevenOverOnesixtwo}{2.5~[1.9, 3.2]}  % experiments/FIGS/data/paired_v1err/probes/ratios.json :: ratios[446].ratio, ratios[446].ci95 :: sha256 14dd51402305d147 :: 17 over 162, visible_content mlp 1-auc: paired geometric mean [95% interval]
\newcommand{\PairedProbeVisibleContentMlpOnesevenOverFourthree}{2.5~[2.0, 3.1]}  % experiments/FIGS/data/paired_v1err/probes/ratios.json :: ratios[447].ratio, ratios[447].ci95 :: sha256 14dd51402305d147 :: 17 over 43, visible_content mlp 1-auc: paired geometric mean [95% interval]
\newcommand{\PairedProbeVisibleContentMlpOnesixtwoMassOverOnesixtwo}{1.08~[0.94, 1.24]}  % experiments/FIGS/data/paired_v1err/probes/ratios.json :: ratios[448].ratio, ratios[448].ci95 :: sha256 14dd51402305d147 :: 162+mass over 162, visible_content mlp 1-auc: paired geometric mean [95% interval]
\newcommand{\PairedProbeVisibleContentMlpOnesevenMassOverOneseven}{1.41~[1.19, 1.67]}  % experiments/FIGS/data/paired_v1err/probes/ratios.json :: ratios[449].ratio, ratios[449].ci95 :: sha256 14dd51402305d147 :: 17+mass over 17, visible_content mlp 1-auc: paired geometric mean [95% interval]
\newcommand{\PairedProbeVisibleContentMlpRandomOverOneeighteight}{1.27~[1.01, 1.59]}  % experiments/FIGS/data/paired_v1err/probes/ratios.json :: ratios[450].ratio, ratios[450].ci95 :: sha256 14dd51402305d147 :: random over 188, visible_content mlp 1-auc: paired geometric mean [95% interval]
\newcommand{\PairedProbeVisibleContentMlpRandomOverOnesixtwo}{1.37~[1.12, 1.68]}  % experiments/FIGS/data/paired_v1err/probes/ratios.json :: ratios[451].ratio, ratios[451].ci95 :: sha256 14dd51402305d147 :: random over 162, visible_content mlp 1-auc: paired geometric mean [95% interval]
\newcommand{\PairedProbeVisibleContentMlpRandomOverFourthree}{1.35~[1.13, 1.61]}  % experiments/FIGS/data/paired_v1err/probes/ratios.json :: ratios[452].ratio, ratios[452].ci95 :: sha256 14dd51402305d147 :: random over 43, visible_content mlp 1-auc: paired geometric mean [95% interval]
\newcommand{\PairedProbeVisibleContentMlpRandomOverOneseven}{0.56~[0.45, 0.70]}  % experiments/FIGS/data/paired_v1err/probes/ratios.json :: ratios[453].ratio, ratios[453].ci95 :: sha256 14dd51402305d147 :: random over 17, visible_content mlp 1-auc: paired geometric mean [95% interval]
\newcommand{\PairedProbeVisibleContentMlpRandomDrawTwoOverRandomDrawOne}{0.88~[0.71, 1.09]}  % experiments/FIGS/data/paired_v1err/probes/ratios.json :: ratios[454].ratio, ratios[454].ci95 :: sha256 14dd51402305d147 :: random draw 2 over random draw 1, visible_content mlp 1-auc: paired geometric mean [95% interval]
\newcommand{\PairedProbeVisibleContentMlpRandomDrawThreeOverRandomDrawOne}{0.94~[0.75, 1.18]}  % experiments/FIGS/data/paired_v1err/probes/ratios.json :: ratios[455].ratio, ratios[455].ci95 :: sha256 14dd51402305d147 :: random draw 3 over random draw 1, visible_content mlp 1-auc: paired geometric mean [95% interval]
\newcommand{\PairedProbeVisibleContentMlpRandomDrawThreeOverRandomDrawTwo}{1.06~[0.84, 1.35]}  % experiments/FIGS/data/paired_v1err/probes/ratios.json :: ratios[456].ratio, ratios[456].ci95 :: sha256 14dd51402305d147 :: random draw 3 over random draw 2, visible_content mlp 1-auc: paired geometric mean [95% interval]
\newcommand{\PairedFtJciiEThreeOnesixtwoOverOneeighteight}{0.95~[0.87, 1.05]}  % experiments/FIGS/data/paired_v1err/finetune/ratios.json :: ratios[0].ratio, ratios[0].ci95 :: sha256 3e5cf25e859b68c2 :: 162 over 188, leg1 N1000 1-macro_auc: paired geometric mean [95% interval]
\newcommand{\PairedFtJciiEThreeFourthreeOverOneeighteight}{0.84~[0.77, 0.91]}  % experiments/FIGS/data/paired_v1err/finetune/ratios.json :: ratios[1].ratio, ratios[1].ci95 :: sha256 3e5cf25e859b68c2 :: 43 over 188, leg1 N1000 1-macro_auc: paired geometric mean [95% interval]
\newcommand{\PairedFtJciiEThreeOnesevenOverOneeighteight}{0.97~[0.86, 1.09]}  % experiments/FIGS/data/paired_v1err/finetune/ratios.json :: ratios[2].ratio, ratios[2].ci95 :: sha256 3e5cf25e859b68c2 :: 17 over 188, leg1 N1000 1-macro_auc: paired geometric mean [95% interval]
\newcommand{\PairedFtJciiEThreeFourthreeOverOnesixtwo}{0.88~[0.80, 0.97]}  % experiments/FIGS/data/paired_v1err/finetune/ratios.json :: ratios[3].ratio, ratios[3].ci95 :: sha256 3e5cf25e859b68c2 :: 43 over 162, leg1 N1000 1-macro_auc: paired geometric mean [95% interval]
\newcommand{\PairedFtJciiEThreeOnesevenOverOnesixtwo}{1.02~[0.88, 1.18]}  % experiments/FIGS/data/paired_v1err/finetune/ratios.json :: ratios[4].ratio, ratios[4].ci95 :: sha256 3e5cf25e859b68c2 :: 17 over 162, leg1 N1000 1-macro_auc: paired geometric mean [95% interval]
\newcommand{\PairedFtJciiEThreeOnesevenOverFourthree}{1.16~[1.08, 1.24]}  % experiments/FIGS/data/paired_v1err/finetune/ratios.json :: ratios[5].ratio, ratios[5].ci95 :: sha256 3e5cf25e859b68c2 :: 17 over 43, leg1 N1000 1-macro_auc: paired geometric mean [95% interval]
\newcommand{\PairedFtJciiEThreeOnesixtwoMassOverOnesixtwo}{0.85~[0.77, 0.93]}  % experiments/FIGS/data/paired_v1err/finetune/ratios.json :: ratios[6].ratio, ratios[6].ci95 :: sha256 3e5cf25e859b68c2 :: 162+mass over 162, leg1 N1000 1-macro_auc: paired geometric mean [95% interval]
\newcommand{\PairedFtJciiEThreeOnesevenMassOverOneseven}{1.00~[0.86, 1.17]}  % experiments/FIGS/data/paired_v1err/finetune/ratios.json :: ratios[7].ratio, ratios[7].ci95 :: sha256 3e5cf25e859b68c2 :: 17+mass over 17, leg1 N1000 1-macro_auc: paired geometric mean [95% interval]
\newcommand{\PairedFtJciiEThreeRandomOverOneeighteight}{0.82~[0.73, 0.93]}  % experiments/FIGS/data/paired_v1err/finetune/ratios.json :: ratios[8].ratio, ratios[8].ci95 :: sha256 3e5cf25e859b68c2 :: random over 188, leg1 N1000 1-macro_auc: paired geometric mean [95% interval]
\newcommand{\PairedFtJciiEThreeRandomOverOnesixtwo}{0.90~[0.85, 0.96]}  % experiments/FIGS/data/paired_v1err/finetune/ratios.json :: ratios[9].ratio, ratios[9].ci95 :: sha256 3e5cf25e859b68c2 :: random over 162, leg1 N1000 1-macro_auc: paired geometric mean [95% interval]
\newcommand{\PairedFtJciiEThreeRandomOverFourthree}{0.99~[0.75, 1.30]}  % experiments/FIGS/data/paired_v1err/finetune/ratios.json :: ratios[10].ratio, ratios[10].ci95 :: sha256 3e5cf25e859b68c2 :: random over 43, leg1 N1000 1-macro_auc: paired geometric mean [95% interval]
\newcommand{\PairedFtJciiEThreeRandomOverOneseven}{0.86~[0.58, 1.27]}  % experiments/FIGS/data/paired_v1err/finetune/ratios.json :: ratios[11].ratio, ratios[11].ci95 :: sha256 3e5cf25e859b68c2 :: random over 17, leg1 N1000 1-macro_auc: paired geometric mean [95% interval]
\newcommand{\PairedFtJciiEThreeRandomDrawTwoOverRandomDrawOne}{1.09~[1.07, 1.10]}  % experiments/FIGS/data/paired_v1err/finetune/ratios.json :: ratios[12].ratio, ratios[12].ci95 :: sha256 3e5cf25e859b68c2 :: random draw 2 over random draw 1, leg1 N1000 1-macro_auc: paired geometric mean [95% interval]
\newcommand{\PairedFtJciiEThreeRandomDrawThreeOverRandomDrawOne}{1.04~[1.02, 1.05]}  % experiments/FIGS/data/paired_v1err/finetune/ratios.json :: ratios[13].ratio, ratios[13].ci95 :: sha256 3e5cf25e859b68c2 :: random draw 3 over random draw 1, leg1 N1000 1-macro_auc: paired geometric mean [95% interval]
\newcommand{\PairedFtJciiEThreeRandomDrawThreeOverRandomDrawTwo}{0.95~[0.94, 0.97]}  % experiments/FIGS/data/paired_v1err/finetune/ratios.json :: ratios[14].ratio, ratios[14].ci95 :: sha256 3e5cf25e859b68c2 :: random draw 3 over random draw 2, leg1 N1000 1-macro_auc: paired geometric mean [95% interval]
\newcommand{\PairedFtJciiEFourOnesixtwoOverOneeighteight}{0.99~[0.97, 1.02]}  % experiments/FIGS/data/paired_v1err/finetune/ratios.json :: ratios[15].ratio, ratios[15].ci95 :: sha256 3e5cf25e859b68c2 :: 162 over 188, leg1 N10000 1-macro_auc: paired geometric mean [95% interval]
\newcommand{\PairedFtJciiEFourFourthreeOverOneeighteight}{1.20~[1.17, 1.23]}  % experiments/FIGS/data/paired_v1err/finetune/ratios.json :: ratios[16].ratio, ratios[16].ci95 :: sha256 3e5cf25e859b68c2 :: 43 over 188, leg1 N10000 1-macro_auc: paired geometric mean [95% interval]
\newcommand{\PairedFtJciiEFourOnesevenOverOneeighteight}{1.32~[1.29, 1.34]}  % experiments/FIGS/data/paired_v1err/finetune/ratios.json :: ratios[17].ratio, ratios[17].ci95 :: sha256 3e5cf25e859b68c2 :: 17 over 188, leg1 N10000 1-macro_auc: paired geometric mean [95% interval]
\newcommand{\PairedFtJciiEFourFourthreeOverOnesixtwo}{1.21~[1.19, 1.23]}  % experiments/FIGS/data/paired_v1err/finetune/ratios.json :: ratios[18].ratio, ratios[18].ci95 :: sha256 3e5cf25e859b68c2 :: 43 over 162, leg1 N10000 1-macro_auc: paired geometric mean [95% interval]
\newcommand{\PairedFtJciiEFourOnesevenOverOnesixtwo}{1.32~[1.29, 1.36]}  % experiments/FIGS/data/paired_v1err/finetune/ratios.json :: ratios[19].ratio, ratios[19].ci95 :: sha256 3e5cf25e859b68c2 :: 17 over 162, leg1 N10000 1-macro_auc: paired geometric mean [95% interval]
\newcommand{\PairedFtJciiEFourOnesevenOverFourthree}{1.10~[1.07, 1.13]}  % experiments/FIGS/data/paired_v1err/finetune/ratios.json :: ratios[20].ratio, ratios[20].ci95 :: sha256 3e5cf25e859b68c2 :: 17 over 43, leg1 N10000 1-macro_auc: paired geometric mean [95% interval]
\newcommand{\PairedFtJciiEFourOnesixtwoMassOverOnesixtwo}{1.00~[0.94, 1.06]}  % experiments/FIGS/data/paired_v1err/finetune/ratios.json :: ratios[21].ratio, ratios[21].ci95 :: sha256 3e5cf25e859b68c2 :: 162+mass over 162, leg1 N10000 1-macro_auc: paired geometric mean [95% interval]
\newcommand{\PairedFtJciiEFourOnesevenMassOverOneseven}{0.990~[0.976, 1.004]}  % experiments/FIGS/data/paired_v1err/finetune/ratios.json :: ratios[22].ratio, ratios[22].ci95 :: sha256 3e5cf25e859b68c2 :: 17+mass over 17, leg1 N10000 1-macro_auc: paired geometric mean [95% interval]
\newcommand{\PairedFtJciiEFourRandomOverOneeighteight}{1.15~[1.09, 1.22]}  % experiments/FIGS/data/paired_v1err/finetune/ratios.json :: ratios[23].ratio, ratios[23].ci95 :: sha256 3e5cf25e859b68c2 :: random over 188, leg1 N10000 1-macro_auc: paired geometric mean [95% interval]
\newcommand{\PairedFtJciiEFourRandomOverOnesixtwo}{1.18~[1.13, 1.22]}  % experiments/FIGS/data/paired_v1err/finetune/ratios.json :: ratios[24].ratio, ratios[24].ci95 :: sha256 3e5cf25e859b68c2 :: random over 162, leg1 N10000 1-macro_auc: paired geometric mean [95% interval]
\newcommand{\PairedFtJciiEFourRandomOverFourthree}{0.97~[0.96, 0.99]}  % experiments/FIGS/data/paired_v1err/finetune/ratios.json :: ratios[25].ratio, ratios[25].ci95 :: sha256 3e5cf25e859b68c2 :: random over 43, leg1 N10000 1-macro_auc: paired geometric mean [95% interval]
\newcommand{\PairedFtJciiEFourRandomOverOneseven}{0.88~[0.82, 0.95]}  % experiments/FIGS/data/paired_v1err/finetune/ratios.json :: ratios[26].ratio, ratios[26].ci95 :: sha256 3e5cf25e859b68c2 :: random over 17, leg1 N10000 1-macro_auc: paired geometric mean [95% interval]
\newcommand{\PairedFtJciiEFourRandomDrawTwoOverRandomDrawOne}{1.00~[0.98, 1.02]}  % experiments/FIGS/data/paired_v1err/finetune/ratios.json :: ratios[27].ratio, ratios[27].ci95 :: sha256 3e5cf25e859b68c2 :: random draw 2 over random draw 1, leg1 N10000 1-macro_auc: paired geometric mean [95% interval]
\newcommand{\PairedFtJciiEFourRandomDrawThreeOverRandomDrawOne}{0.97~[0.95, 0.98]}  % experiments/FIGS/data/paired_v1err/finetune/ratios.json :: ratios[28].ratio, ratios[28].ci95 :: sha256 3e5cf25e859b68c2 :: random draw 3 over random draw 1, leg1 N10000 1-macro_auc: paired geometric mean [95% interval]
\newcommand{\PairedFtJciiEFourRandomDrawThreeOverRandomDrawTwo}{0.97~[0.95, 0.98]}  % experiments/FIGS/data/paired_v1err/finetune/ratios.json :: ratios[29].ratio, ratios[29].ci95 :: sha256 3e5cf25e859b68c2 :: random draw 3 over random draw 2, leg1 N10000 1-macro_auc: paired geometric mean [95% interval]
\newcommand{\PairedFtJciiEFiveOnesixtwoOverOneeighteight}{1.00~[0.98, 1.01]}  % experiments/FIGS/data/paired_v1err/finetune/ratios.json :: ratios[30].ratio, ratios[30].ci95 :: sha256 3e5cf25e859b68c2 :: 162 over 188, leg1 N100000 1-macro_auc: paired geometric mean [95% interval]
\newcommand{\PairedFtJciiEFiveFourthreeOverOneeighteight}{1.09~[1.06, 1.11]}  % experiments/FIGS/data/paired_v1err/finetune/ratios.json :: ratios[31].ratio, ratios[31].ci95 :: sha256 3e5cf25e859b68c2 :: 43 over 188, leg1 N100000 1-macro_auc: paired geometric mean [95% interval]
\newcommand{\PairedFtJciiEFiveOnesevenOverOneeighteight}{1.22~[1.20, 1.23]}  % experiments/FIGS/data/paired_v1err/finetune/ratios.json :: ratios[32].ratio, ratios[32].ci95 :: sha256 3e5cf25e859b68c2 :: 17 over 188, leg1 N100000 1-macro_auc: paired geometric mean [95% interval]
\newcommand{\PairedFtJciiEFiveFourthreeOverOnesixtwo}{1.09~[1.07, 1.11]}  % experiments/FIGS/data/paired_v1err/finetune/ratios.json :: ratios[33].ratio, ratios[33].ci95 :: sha256 3e5cf25e859b68c2 :: 43 over 162, leg1 N100000 1-macro_auc: paired geometric mean [95% interval]
\newcommand{\PairedFtJciiEFiveOnesevenOverOnesixtwo}{1.22~[1.20, 1.24]}  % experiments/FIGS/data/paired_v1err/finetune/ratios.json :: ratios[34].ratio, ratios[34].ci95 :: sha256 3e5cf25e859b68c2 :: 17 over 162, leg1 N100000 1-macro_auc: paired geometric mean [95% interval]
\newcommand{\PairedFtJciiEFiveOnesevenOverFourthree}{1.12~[1.10, 1.14]}  % experiments/FIGS/data/paired_v1err/finetune/ratios.json :: ratios[35].ratio, ratios[35].ci95 :: sha256 3e5cf25e859b68c2 :: 17 over 43, leg1 N100000 1-macro_auc: paired geometric mean [95% interval]
\newcommand{\PairedFtJciiEFiveOnesixtwoMassOverOnesixtwo}{1.018~[1.002, 1.035]}  % experiments/FIGS/data/paired_v1err/finetune/ratios.json :: ratios[36].ratio, ratios[36].ci95 :: sha256 3e5cf25e859b68c2 :: 162+mass over 162, leg1 N100000 1-macro_auc: paired geometric mean [95% interval]
\newcommand{\PairedFtJciiEFiveOnesevenMassOverOneseven}{1.03~[1.01, 1.06]}  % experiments/FIGS/data/paired_v1err/finetune/ratios.json :: ratios[37].ratio, ratios[37].ci95 :: sha256 3e5cf25e859b68c2 :: 17+mass over 17, leg1 N100000 1-macro_auc: paired geometric mean [95% interval]
\newcommand{\PairedFtJciiEFiveRandomOverOneeighteight}{1.15~[1.09, 1.20]}  % experiments/FIGS/data/paired_v1err/finetune/ratios.json :: ratios[38].ratio, ratios[38].ci95 :: sha256 3e5cf25e859b68c2 :: random over 188, leg1 N100000 1-macro_auc: paired geometric mean [95% interval]
\newcommand{\PairedFtJciiEFiveRandomOverOnesixtwo}{1.15~[1.09, 1.20]}  % experiments/FIGS/data/paired_v1err/finetune/ratios.json :: ratios[39].ratio, ratios[39].ci95 :: sha256 3e5cf25e859b68c2 :: random over 162, leg1 N100000 1-macro_auc: paired geometric mean [95% interval]
\newcommand{\PairedFtJciiEFiveRandomOverFourthree}{1.05~[1.04, 1.06]}  % experiments/FIGS/data/paired_v1err/finetune/ratios.json :: ratios[40].ratio, ratios[40].ci95 :: sha256 3e5cf25e859b68c2 :: random over 43, leg1 N100000 1-macro_auc: paired geometric mean [95% interval]
\newcommand{\PairedFtJciiEFiveRandomOverOneseven}{0.95~[0.93, 0.96]}  % experiments/FIGS/data/paired_v1err/finetune/ratios.json :: ratios[41].ratio, ratios[41].ci95 :: sha256 3e5cf25e859b68c2 :: random over 17, leg1 N100000 1-macro_auc: paired geometric mean [95% interval]
\newcommand{\PairedFtJciiEFiveRandomDrawTwoOverRandomDrawOne}{0.98~[0.97, 0.99]}  % experiments/FIGS/data/paired_v1err/finetune/ratios.json :: ratios[42].ratio, ratios[42].ci95 :: sha256 3e5cf25e859b68c2 :: random draw 2 over random draw 1, leg1 N100000 1-macro_auc: paired geometric mean [95% interval]
\newcommand{\PairedFtJciiEFiveRandomDrawThreeOverRandomDrawOne}{1.00~[0.99, 1.02]}  % experiments/FIGS/data/paired_v1err/finetune/ratios.json :: ratios[43].ratio, ratios[43].ci95 :: sha256 3e5cf25e859b68c2 :: random draw 3 over random draw 1, leg1 N100000 1-macro_auc: paired geometric mean [95% interval]
\newcommand{\PairedFtJciiEFiveRandomDrawThreeOverRandomDrawTwo}{1.03~[1.01, 1.04]}  % experiments/FIGS/data/paired_v1err/finetune/ratios.json :: ratios[44].ratio, ratios[44].ci95 :: sha256 3e5cf25e859b68c2 :: random draw 3 over random draw 2, leg1 N100000 1-macro_auc: paired geometric mean [95% interval]
\newcommand{\PairedFtJciiESixOnesixtwoOverOneeighteight}{1.002~[0.998, 1.005]}  % experiments/FIGS/data/paired_v1err/finetune/ratios.json :: ratios[45].ratio, ratios[45].ci95 :: sha256 3e5cf25e859b68c2 :: 162 over 188, leg1 N1000000 1-macro_auc: paired geometric mean [95% interval]
\newcommand{\PairedFtJciiESixFourthreeOverOneeighteight}{1.045~[1.037, 1.053]}  % experiments/FIGS/data/paired_v1err/finetune/ratios.json :: ratios[46].ratio, ratios[46].ci95 :: sha256 3e5cf25e859b68c2 :: 43 over 188, leg1 N1000000 1-macro_auc: paired geometric mean [95% interval]
\newcommand{\PairedFtJciiESixOnesevenOverOneeighteight}{1.14~[1.13, 1.16]}  % experiments/FIGS/data/paired_v1err/finetune/ratios.json :: ratios[47].ratio, ratios[47].ci95 :: sha256 3e5cf25e859b68c2 :: 17 over 188, leg1 N1000000 1-macro_auc: paired geometric mean [95% interval]
\newcommand{\PairedFtJciiESixFourthreeOverOnesixtwo}{1.044~[1.036, 1.051]}  % experiments/FIGS/data/paired_v1err/finetune/ratios.json :: ratios[48].ratio, ratios[48].ci95 :: sha256 3e5cf25e859b68c2 :: 43 over 162, leg1 N1000000 1-macro_auc: paired geometric mean [95% interval]
\newcommand{\PairedFtJciiESixOnesevenOverOnesixtwo}{1.14~[1.13, 1.15]}  % experiments/FIGS/data/paired_v1err/finetune/ratios.json :: ratios[49].ratio, ratios[49].ci95 :: sha256 3e5cf25e859b68c2 :: 17 over 162, leg1 N1000000 1-macro_auc: paired geometric mean [95% interval]
\newcommand{\PairedFtJciiESixOnesevenOverFourthree}{1.09~[1.08, 1.11]}  % experiments/FIGS/data/paired_v1err/finetune/ratios.json :: ratios[50].ratio, ratios[50].ci95 :: sha256 3e5cf25e859b68c2 :: 17 over 43, leg1 N1000000 1-macro_auc: paired geometric mean [95% interval]
\newcommand{\PairedFtJciiESixOnesixtwoMassOverOnesixtwo}{1.019~[1.014, 1.023]}  % experiments/FIGS/data/paired_v1err/finetune/ratios.json :: ratios[51].ratio, ratios[51].ci95 :: sha256 3e5cf25e859b68c2 :: 162+mass over 162, leg1 N1000000 1-macro_auc: paired geometric mean [95% interval]
\newcommand{\PairedFtJciiESixOnesevenMassOverOneseven}{1.03~[1.02, 1.04]}  % experiments/FIGS/data/paired_v1err/finetune/ratios.json :: ratios[52].ratio, ratios[52].ci95 :: sha256 3e5cf25e859b68c2 :: 17+mass over 17, leg1 N1000000 1-macro_auc: paired geometric mean [95% interval]
\newcommand{\PairedFtJciiESixRandomOverOneeighteight}{1.073~[1.067, 1.080]}  % experiments/FIGS/data/paired_v1err/finetune/ratios.json :: ratios[53].ratio, ratios[53].ci95 :: sha256 3e5cf25e859b68c2 :: random over 188, leg1 N1000000 1-macro_auc: paired geometric mean [95% interval]
\newcommand{\PairedFtJciiESixRandomOverOnesixtwo}{1.072~[1.065, 1.078]}  % experiments/FIGS/data/paired_v1err/finetune/ratios.json :: ratios[54].ratio, ratios[54].ci95 :: sha256 3e5cf25e859b68c2 :: random over 162, leg1 N1000000 1-macro_auc: paired geometric mean [95% interval]
\newcommand{\PairedFtJciiESixRandomOverFourthree}{1.03~[1.02, 1.04]}  % experiments/FIGS/data/paired_v1err/finetune/ratios.json :: ratios[55].ratio, ratios[55].ci95 :: sha256 3e5cf25e859b68c2 :: random over 43, leg1 N1000000 1-macro_auc: paired geometric mean [95% interval]
\newcommand{\PairedFtJciiESixRandomOverOneseven}{0.94~[0.93, 0.95]}  % experiments/FIGS/data/paired_v1err/finetune/ratios.json :: ratios[56].ratio, ratios[56].ci95 :: sha256 3e5cf25e859b68c2 :: random over 17, leg1 N1000000 1-macro_auc: paired geometric mean [95% interval]
\newcommand{\PairedFtJciiESixRandomDrawTwoOverRandomDrawOne}{1.00~[0.99, 1.01]}  % experiments/FIGS/data/paired_v1err/finetune/ratios.json :: ratios[57].ratio, ratios[57].ci95 :: sha256 3e5cf25e859b68c2 :: random draw 2 over random draw 1, leg1 N1000000 1-macro_auc: paired geometric mean [95% interval]
\newcommand{\PairedFtJciiESixRandomDrawThreeOverRandomDrawOne}{0.996~[0.989, 1.003]}  % experiments/FIGS/data/paired_v1err/finetune/ratios.json :: ratios[58].ratio, ratios[58].ci95 :: sha256 3e5cf25e859b68c2 :: random draw 3 over random draw 1, leg1 N1000000 1-macro_auc: paired geometric mean [95% interval]
\newcommand{\PairedFtJciiESixRandomDrawThreeOverRandomDrawTwo}{0.997~[0.989, 1.004]}  % experiments/FIGS/data/paired_v1err/finetune/ratios.json :: ratios[59].ratio, ratios[59].ci95 :: sha256 3e5cf25e859b68c2 :: random draw 3 over random draw 2, leg1 N1000000 1-macro_auc: paired geometric mean [95% interval]
\newcommand{\PairedFtJcEThreeOnesixtwoOverOneeighteight}{0.96~[0.83, 1.11]}  % experiments/FIGS/data/paired_v1err/finetune/ratios.json :: ratios[60].ratio, ratios[60].ci95 :: sha256 3e5cf25e859b68c2 :: 162 over 188, leg2 N1000 1-macro_auc: paired geometric mean [95% interval]
\newcommand{\PairedFtJcEThreeFourthreeOverOneeighteight}{1.01~[0.84, 1.21]}  % experiments/FIGS/data/paired_v1err/finetune/ratios.json :: ratios[61].ratio, ratios[61].ci95 :: sha256 3e5cf25e859b68c2 :: 43 over 188, leg2 N1000 1-macro_auc: paired geometric mean [95% interval]
\newcommand{\PairedFtJcEThreeOnesevenOverOneeighteight}{1.07~[0.93, 1.23]}  % experiments/FIGS/data/paired_v1err/finetune/ratios.json :: ratios[62].ratio, ratios[62].ci95 :: sha256 3e5cf25e859b68c2 :: 17 over 188, leg2 N1000 1-macro_auc: paired geometric mean [95% interval]
\newcommand{\PairedFtJcEThreeFourthreeOverOnesixtwo}{1.05~[0.96, 1.15]}  % experiments/FIGS/data/paired_v1err/finetune/ratios.json :: ratios[63].ratio, ratios[63].ci95 :: sha256 3e5cf25e859b68c2 :: 43 over 162, leg2 N1000 1-macro_auc: paired geometric mean [95% interval]
\newcommand{\PairedFtJcEThreeOnesevenOverOnesixtwo}{1.12~[1.07, 1.17]}  % experiments/FIGS/data/paired_v1err/finetune/ratios.json :: ratios[64].ratio, ratios[64].ci95 :: sha256 3e5cf25e859b68c2 :: 17 over 162, leg2 N1000 1-macro_auc: paired geometric mean [95% interval]
\newcommand{\PairedFtJcEThreeOnesevenOverFourthree}{1.06~[0.98, 1.15]}  % experiments/FIGS/data/paired_v1err/finetune/ratios.json :: ratios[65].ratio, ratios[65].ci95 :: sha256 3e5cf25e859b68c2 :: 17 over 43, leg2 N1000 1-macro_auc: paired geometric mean [95% interval]
\newcommand{\PairedFtJcEThreeOnesixtwoMassOverOnesixtwo}{0.91~[0.88, 0.94]}  % experiments/FIGS/data/paired_v1err/finetune/ratios.json :: ratios[66].ratio, ratios[66].ci95 :: sha256 3e5cf25e859b68c2 :: 162+mass over 162, leg2 N1000 1-macro_auc: paired geometric mean [95% interval]
\newcommand{\PairedFtJcEThreeOnesevenMassOverOneseven}{1.02~[0.92, 1.13]}  % experiments/FIGS/data/paired_v1err/finetune/ratios.json :: ratios[67].ratio, ratios[67].ci95 :: sha256 3e5cf25e859b68c2 :: 17+mass over 17, leg2 N1000 1-macro_auc: paired geometric mean [95% interval]
\newcommand{\PairedFtJcEThreeRandomOverOneeighteight}{1.01~[0.78, 1.32]}  % experiments/FIGS/data/paired_v1err/finetune/ratios.json :: ratios[68].ratio, ratios[68].ci95 :: sha256 3e5cf25e859b68c2 :: random over 188, leg2 N1000 1-macro_auc: paired geometric mean [95% interval]
\newcommand{\PairedFtJcEThreeRandomOverOnesixtwo}{1.12~[1.04, 1.20]}  % experiments/FIGS/data/paired_v1err/finetune/ratios.json :: ratios[69].ratio, ratios[69].ci95 :: sha256 3e5cf25e859b68c2 :: random over 162, leg2 N1000 1-macro_auc: paired geometric mean [95% interval]
\newcommand{\PairedFtJcEThreeRandomOverFourthree}{1.06~[0.79, 1.41]}  % experiments/FIGS/data/paired_v1err/finetune/ratios.json :: ratios[70].ratio, ratios[70].ci95 :: sha256 3e5cf25e859b68c2 :: random over 43, leg2 N1000 1-macro_auc: paired geometric mean [95% interval]
\newcommand{\PairedFtJcEThreeRandomOverOneseven}{0.98~[0.84, 1.13]}  % experiments/FIGS/data/paired_v1err/finetune/ratios.json :: ratios[71].ratio, ratios[71].ci95 :: sha256 3e5cf25e859b68c2 :: random over 17, leg2 N1000 1-macro_auc: paired geometric mean [95% interval]
\newcommand{\PairedFtJcEThreeRandomDrawTwoOverRandomDrawOne}{0.954~[0.949, 0.958]}  % experiments/FIGS/data/paired_v1err/finetune/ratios.json :: ratios[72].ratio, ratios[72].ci95 :: sha256 3e5cf25e859b68c2 :: random draw 2 over random draw 1, leg2 N1000 1-macro_auc: paired geometric mean [95% interval]
\newcommand{\PairedFtJcEThreeRandomDrawThreeOverRandomDrawOne}{1.013~[1.009, 1.017]}  % experiments/FIGS/data/paired_v1err/finetune/ratios.json :: ratios[73].ratio, ratios[73].ci95 :: sha256 3e5cf25e859b68c2 :: random draw 3 over random draw 1, leg2 N1000 1-macro_auc: paired geometric mean [95% interval]
\newcommand{\PairedFtJcEThreeRandomDrawThreeOverRandomDrawTwo}{1.062~[1.057, 1.067]}  % experiments/FIGS/data/paired_v1err/finetune/ratios.json :: ratios[74].ratio, ratios[74].ci95 :: sha256 3e5cf25e859b68c2 :: random draw 3 over random draw 2, leg2 N1000 1-macro_auc: paired geometric mean [95% interval]
\newcommand{\PairedFtJcEFourOnesixtwoOverOneeighteight}{0.98~[0.93, 1.02]}  % experiments/FIGS/data/paired_v1err/finetune/ratios.json :: ratios[75].ratio, ratios[75].ci95 :: sha256 3e5cf25e859b68c2 :: 162 over 188, leg2 N10000 1-macro_auc: paired geometric mean [95% interval]
\newcommand{\PairedFtJcEFourFourthreeOverOneeighteight}{0.98~[0.93, 1.04]}  % experiments/FIGS/data/paired_v1err/finetune/ratios.json :: ratios[76].ratio, ratios[76].ci95 :: sha256 3e5cf25e859b68c2 :: 43 over 188, leg2 N10000 1-macro_auc: paired geometric mean [95% interval]
\newcommand{\PairedFtJcEFourOnesevenOverOneeighteight}{1.07~[1.02, 1.11]}  % experiments/FIGS/data/paired_v1err/finetune/ratios.json :: ratios[77].ratio, ratios[77].ci95 :: sha256 3e5cf25e859b68c2 :: 17 over 188, leg2 N10000 1-macro_auc: paired geometric mean [95% interval]
\newcommand{\PairedFtJcEFourFourthreeOverOnesixtwo}{1.01~[0.97, 1.04]}  % experiments/FIGS/data/paired_v1err/finetune/ratios.json :: ratios[78].ratio, ratios[78].ci95 :: sha256 3e5cf25e859b68c2 :: 43 over 162, leg2 N10000 1-macro_auc: paired geometric mean [95% interval]
\newcommand{\PairedFtJcEFourOnesevenOverOnesixtwo}{1.09~[1.06, 1.12]}  % experiments/FIGS/data/paired_v1err/finetune/ratios.json :: ratios[79].ratio, ratios[79].ci95 :: sha256 3e5cf25e859b68c2 :: 17 over 162, leg2 N10000 1-macro_auc: paired geometric mean [95% interval]
\newcommand{\PairedFtJcEFourOnesevenOverFourthree}{1.08~[1.07, 1.10]}  % experiments/FIGS/data/paired_v1err/finetune/ratios.json :: ratios[80].ratio, ratios[80].ci95 :: sha256 3e5cf25e859b68c2 :: 17 over 43, leg2 N10000 1-macro_auc: paired geometric mean [95% interval]
\newcommand{\PairedFtJcEFourOnesixtwoMassOverOnesixtwo}{1.01~[0.98, 1.04]}  % experiments/FIGS/data/paired_v1err/finetune/ratios.json :: ratios[81].ratio, ratios[81].ci95 :: sha256 3e5cf25e859b68c2 :: 162+mass over 162, leg2 N10000 1-macro_auc: paired geometric mean [95% interval]
\newcommand{\PairedFtJcEFourOnesevenMassOverOneseven}{1.05~[1.02, 1.08]}  % experiments/FIGS/data/paired_v1err/finetune/ratios.json :: ratios[82].ratio, ratios[82].ci95 :: sha256 3e5cf25e859b68c2 :: 17+mass over 17, leg2 N10000 1-macro_auc: paired geometric mean [95% interval]
\newcommand{\PairedFtJcEFourRandomOverOneeighteight}{1.08~[1.03, 1.14]}  % experiments/FIGS/data/paired_v1err/finetune/ratios.json :: ratios[83].ratio, ratios[83].ci95 :: sha256 3e5cf25e859b68c2 :: random over 188, leg2 N10000 1-macro_auc: paired geometric mean [95% interval]
\newcommand{\PairedFtJcEFourRandomOverOnesixtwo}{1.11~[0.99, 1.23]}  % experiments/FIGS/data/paired_v1err/finetune/ratios.json :: ratios[84].ratio, ratios[84].ci95 :: sha256 3e5cf25e859b68c2 :: random over 162, leg2 N10000 1-macro_auc: paired geometric mean [95% interval]
\newcommand{\PairedFtJcEFourRandomOverFourthree}{1.10~[0.99, 1.21]}  % experiments/FIGS/data/paired_v1err/finetune/ratios.json :: ratios[85].ratio, ratios[85].ci95 :: sha256 3e5cf25e859b68c2 :: random over 43, leg2 N10000 1-macro_auc: paired geometric mean [95% interval]
\newcommand{\PairedFtJcEFourRandomOverOneseven}{1.02~[0.94, 1.10]}  % experiments/FIGS/data/paired_v1err/finetune/ratios.json :: ratios[86].ratio, ratios[86].ci95 :: sha256 3e5cf25e859b68c2 :: random over 17, leg2 N10000 1-macro_auc: paired geometric mean [95% interval]
\newcommand{\PairedFtJcEFourRandomDrawTwoOverRandomDrawOne}{0.960~[0.955, 0.965]}  % experiments/FIGS/data/paired_v1err/finetune/ratios.json :: ratios[87].ratio, ratios[87].ci95 :: sha256 3e5cf25e859b68c2 :: random draw 2 over random draw 1, leg2 N10000 1-macro_auc: paired geometric mean [95% interval]
\newcommand{\PairedFtJcEFourRandomDrawThreeOverRandomDrawOne}{1.008~[1.004, 1.012]}  % experiments/FIGS/data/paired_v1err/finetune/ratios.json :: ratios[88].ratio, ratios[88].ci95 :: sha256 3e5cf25e859b68c2 :: random draw 3 over random draw 1, leg2 N10000 1-macro_auc: paired geometric mean [95% interval]
\newcommand{\PairedFtJcEFourRandomDrawThreeOverRandomDrawTwo}{1.050~[1.045, 1.055]}  % experiments/FIGS/data/paired_v1err/finetune/ratios.json :: ratios[89].ratio, ratios[89].ci95 :: sha256 3e5cf25e859b68c2 :: random draw 3 over random draw 2, leg2 N10000 1-macro_auc: paired geometric mean [95% interval]
\newcommand{\PairedFtJcEFiveOnesixtwoOverOneeighteight}{0.983~[0.968, 0.998]}  % experiments/FIGS/data/paired_v1err/finetune/ratios.json :: ratios[90].ratio, ratios[90].ci95 :: sha256 3e5cf25e859b68c2 :: 162 over 188, leg2 N100000 1-macro_auc: paired geometric mean [95% interval]
\newcommand{\PairedFtJcEFiveFourthreeOverOneeighteight}{1.015~[1.002, 1.029]}  % experiments/FIGS/data/paired_v1err/finetune/ratios.json :: ratios[91].ratio, ratios[91].ci95 :: sha256 3e5cf25e859b68c2 :: 43 over 188, leg2 N100000 1-macro_auc: paired geometric mean [95% interval]
\newcommand{\PairedFtJcEFiveOnesevenOverOneeighteight}{1.07~[1.06, 1.09]}  % experiments/FIGS/data/paired_v1err/finetune/ratios.json :: ratios[92].ratio, ratios[92].ci95 :: sha256 3e5cf25e859b68c2 :: 17 over 188, leg2 N100000 1-macro_auc: paired geometric mean [95% interval]
\newcommand{\PairedFtJcEFiveFourthreeOverOnesixtwo}{1.03~[1.02, 1.05]}  % experiments/FIGS/data/paired_v1err/finetune/ratios.json :: ratios[93].ratio, ratios[93].ci95 :: sha256 3e5cf25e859b68c2 :: 43 over 162, leg2 N100000 1-macro_auc: paired geometric mean [95% interval]
\newcommand{\PairedFtJcEFiveOnesevenOverOnesixtwo}{1.091~[1.088, 1.094]}  % experiments/FIGS/data/paired_v1err/finetune/ratios.json :: ratios[94].ratio, ratios[94].ci95 :: sha256 3e5cf25e859b68c2 :: 17 over 162, leg2 N100000 1-macro_auc: paired geometric mean [95% interval]
\newcommand{\PairedFtJcEFiveOnesevenOverFourthree}{1.06~[1.04, 1.07]}  % experiments/FIGS/data/paired_v1err/finetune/ratios.json :: ratios[95].ratio, ratios[95].ci95 :: sha256 3e5cf25e859b68c2 :: 17 over 43, leg2 N100000 1-macro_auc: paired geometric mean [95% interval]
\newcommand{\PairedFtJcEFiveOnesixtwoMassOverOnesixtwo}{1.016~[1.008, 1.024]}  % experiments/FIGS/data/paired_v1err/finetune/ratios.json :: ratios[96].ratio, ratios[96].ci95 :: sha256 3e5cf25e859b68c2 :: 162+mass over 162, leg2 N100000 1-macro_auc: paired geometric mean [95% interval]
\newcommand{\PairedFtJcEFiveOnesevenMassOverOneseven}{1.02~[1.01, 1.03]}  % experiments/FIGS/data/paired_v1err/finetune/ratios.json :: ratios[97].ratio, ratios[97].ci95 :: sha256 3e5cf25e859b68c2 :: 17+mass over 17, leg2 N100000 1-macro_auc: paired geometric mean [95% interval]
\newcommand{\PairedFtJcEFiveRandomOverOneeighteight}{1.044~[1.039, 1.049]}  % experiments/FIGS/data/paired_v1err/finetune/ratios.json :: ratios[98].ratio, ratios[98].ci95 :: sha256 3e5cf25e859b68c2 :: random over 188, leg2 N100000 1-macro_auc: paired geometric mean [95% interval]
\newcommand{\PairedFtJcEFiveRandomOverOnesixtwo}{1.06~[1.04, 1.09]}  % experiments/FIGS/data/paired_v1err/finetune/ratios.json :: ratios[99].ratio, ratios[99].ci95 :: sha256 3e5cf25e859b68c2 :: random over 162, leg2 N100000 1-macro_auc: paired geometric mean [95% interval]
\newcommand{\PairedFtJcEFiveRandomOverFourthree}{1.03~[1.01, 1.04]}  % experiments/FIGS/data/paired_v1err/finetune/ratios.json :: ratios[100].ratio, ratios[100].ci95 :: sha256 3e5cf25e859b68c2 :: random over 43, leg2 N100000 1-macro_auc: paired geometric mean [95% interval]
\newcommand{\PairedFtJcEFiveRandomOverOneseven}{0.97~[0.95, 0.99]}  % experiments/FIGS/data/paired_v1err/finetune/ratios.json :: ratios[101].ratio, ratios[101].ci95 :: sha256 3e5cf25e859b68c2 :: random over 17, leg2 N100000 1-macro_auc: paired geometric mean [95% interval]
\newcommand{\PairedFtJcEFiveRandomDrawTwoOverRandomDrawOne}{0.992~[0.988, 0.997]}  % experiments/FIGS/data/paired_v1err/finetune/ratios.json :: ratios[102].ratio, ratios[102].ci95 :: sha256 3e5cf25e859b68c2 :: random draw 2 over random draw 1, leg2 N100000 1-macro_auc: paired geometric mean [95% interval]
\newcommand{\PairedFtJcEFiveRandomDrawThreeOverRandomDrawOne}{1.007~[1.003, 1.011]}  % experiments/FIGS/data/paired_v1err/finetune/ratios.json :: ratios[103].ratio, ratios[103].ci95 :: sha256 3e5cf25e859b68c2 :: random draw 3 over random draw 1, leg2 N100000 1-macro_auc: paired geometric mean [95% interval]
\newcommand{\PairedFtJcEFiveRandomDrawThreeOverRandomDrawTwo}{1.015~[1.011, 1.020]}  % experiments/FIGS/data/paired_v1err/finetune/ratios.json :: ratios[104].ratio, ratios[104].ci95 :: sha256 3e5cf25e859b68c2 :: random draw 3 over random draw 2, leg2 N100000 1-macro_auc: paired geometric mean [95% interval]
\newcommand{\PairedFtJcESixOnesixtwoOverOneeighteight}{1.00~[0.99, 1.01]}  % experiments/FIGS/data/paired_v1err/finetune/ratios.json :: ratios[105].ratio, ratios[105].ci95 :: sha256 3e5cf25e859b68c2 :: 162 over 188, leg2 N1000000 1-macro_auc: paired geometric mean [95% interval]
\newcommand{\PairedFtJcESixFourthreeOverOneeighteight}{1.02~[1.01, 1.03]}  % experiments/FIGS/data/paired_v1err/finetune/ratios.json :: ratios[106].ratio, ratios[106].ci95 :: sha256 3e5cf25e859b68c2 :: 43 over 188, leg2 N1000000 1-macro_auc: paired geometric mean [95% interval]
\newcommand{\PairedFtJcESixOnesevenOverOneeighteight}{1.08~[1.07, 1.09]}  % experiments/FIGS/data/paired_v1err/finetune/ratios.json :: ratios[107].ratio, ratios[107].ci95 :: sha256 3e5cf25e859b68c2 :: 17 over 188, leg2 N1000000 1-macro_auc: paired geometric mean [95% interval]
\newcommand{\PairedFtJcESixFourthreeOverOnesixtwo}{1.02~[1.01, 1.03]}  % experiments/FIGS/data/paired_v1err/finetune/ratios.json :: ratios[108].ratio, ratios[108].ci95 :: sha256 3e5cf25e859b68c2 :: 43 over 162, leg2 N1000000 1-macro_auc: paired geometric mean [95% interval]
\newcommand{\PairedFtJcESixOnesevenOverOnesixtwo}{1.08~[1.07, 1.09]}  % experiments/FIGS/data/paired_v1err/finetune/ratios.json :: ratios[109].ratio, ratios[109].ci95 :: sha256 3e5cf25e859b68c2 :: 17 over 162, leg2 N1000000 1-macro_auc: paired geometric mean [95% interval]
\newcommand{\PairedFtJcESixOnesevenOverFourthree}{1.06~[1.05, 1.07]}  % experiments/FIGS/data/paired_v1err/finetune/ratios.json :: ratios[110].ratio, ratios[110].ci95 :: sha256 3e5cf25e859b68c2 :: 17 over 43, leg2 N1000000 1-macro_auc: paired geometric mean [95% interval]
\newcommand{\PairedFtJcESixOnesixtwoMassOverOnesixtwo}{1.012~[1.004, 1.020]}  % experiments/FIGS/data/paired_v1err/finetune/ratios.json :: ratios[111].ratio, ratios[111].ci95 :: sha256 3e5cf25e859b68c2 :: 162+mass over 162, leg2 N1000000 1-macro_auc: paired geometric mean [95% interval]
\newcommand{\PairedFtJcESixOnesevenMassOverOneseven}{1.02~[1.01, 1.04]}  % experiments/FIGS/data/paired_v1err/finetune/ratios.json :: ratios[112].ratio, ratios[112].ci95 :: sha256 3e5cf25e859b68c2 :: 17+mass over 17, leg2 N1000000 1-macro_auc: paired geometric mean [95% interval]
\newcommand{\PairedFtJcESixRandomOverOneeighteight}{1.046~[1.038, 1.053]}  % experiments/FIGS/data/paired_v1err/finetune/ratios.json :: ratios[113].ratio, ratios[113].ci95 :: sha256 3e5cf25e859b68c2 :: random over 188, leg2 N1000000 1-macro_auc: paired geometric mean [95% interval]
\newcommand{\PairedFtJcESixRandomOverOnesixtwo}{1.047~[1.038, 1.055]}  % experiments/FIGS/data/paired_v1err/finetune/ratios.json :: ratios[114].ratio, ratios[114].ci95 :: sha256 3e5cf25e859b68c2 :: random over 162, leg2 N1000000 1-macro_auc: paired geometric mean [95% interval]
\newcommand{\PairedFtJcESixRandomOverFourthree}{1.03~[1.02, 1.04]}  % experiments/FIGS/data/paired_v1err/finetune/ratios.json :: ratios[115].ratio, ratios[115].ci95 :: sha256 3e5cf25e859b68c2 :: random over 43, leg2 N1000000 1-macro_auc: paired geometric mean [95% interval]
\newcommand{\PairedFtJcESixRandomOverOneseven}{0.97~[0.96, 0.98]}  % experiments/FIGS/data/paired_v1err/finetune/ratios.json :: ratios[116].ratio, ratios[116].ci95 :: sha256 3e5cf25e859b68c2 :: random over 17, leg2 N1000000 1-macro_auc: paired geometric mean [95% interval]
\newcommand{\PairedFtJcESixRandomDrawTwoOverRandomDrawOne}{1.001~[0.998, 1.004]}  % experiments/FIGS/data/paired_v1err/finetune/ratios.json :: ratios[117].ratio, ratios[117].ci95 :: sha256 3e5cf25e859b68c2 :: random draw 2 over random draw 1, leg2 N1000000 1-macro_auc: paired geometric mean [95% interval]
\newcommand{\PairedFtJcESixRandomDrawThreeOverRandomDrawOne}{1.005~[1.002, 1.008]}  % experiments/FIGS/data/paired_v1err/finetune/ratios.json :: ratios[118].ratio, ratios[118].ci95 :: sha256 3e5cf25e859b68c2 :: random draw 3 over random draw 1, leg2 N1000000 1-macro_auc: paired geometric mean [95% interval]
\newcommand{\PairedFtJcESixRandomDrawThreeOverRandomDrawTwo}{1.004~[1.001, 1.007]}  % experiments/FIGS/data/paired_v1err/finetune/ratios.json :: ratios[119].ratio, ratios[119].ci95 :: sha256 3e5cf25e859b68c2 :: random draw 3 over random draw 2, leg2 N1000000 1-macro_auc: paired geometric mean [95% interval]
\newcommand{\PairedMassResMlpOnesixtwoOverOneeighteight}{0.990~[0.977, 1.003]}  % experiments/FIGS/data/paired_v1err/mass/ratios.json :: ratios[0].ratio, ratios[0].ci95 :: sha256 d877a7f087260f64 :: 162 over 188, resolution mlp sigma_eff: paired geometric mean [95% interval]
\newcommand{\PairedMassResMlpFourthreeOverOneeighteight}{0.98~[0.96, 0.99]}  % experiments/FIGS/data/paired_v1err/mass/ratios.json :: ratios[1].ratio, ratios[1].ci95 :: sha256 d877a7f087260f64 :: 43 over 188, resolution mlp sigma_eff: paired geometric mean [95% interval]
\newcommand{\PairedMassResMlpOnesevenOverOneeighteight}{0.96~[0.94, 0.98]}  % experiments/FIGS/data/paired_v1err/mass/ratios.json :: ratios[2].ratio, ratios[2].ci95 :: sha256 d877a7f087260f64 :: 17 over 188, resolution mlp sigma_eff: paired geometric mean [95% interval]
\newcommand{\PairedMassResMlpFourthreeOverOnesixtwo}{0.989~[0.978, 0.999]}  % experiments/FIGS/data/paired_v1err/mass/ratios.json :: ratios[3].ratio, ratios[3].ci95 :: sha256 d877a7f087260f64 :: 43 over 162, resolution mlp sigma_eff: paired geometric mean [95% interval]
\newcommand{\PairedMassResMlpOnesevenOverOnesixtwo}{0.97~[0.94, 0.99]}  % experiments/FIGS/data/paired_v1err/mass/ratios.json :: ratios[4].ratio, ratios[4].ci95 :: sha256 d877a7f087260f64 :: 17 over 162, resolution mlp sigma_eff: paired geometric mean [95% interval]
\newcommand{\PairedMassResMlpOnesevenOverFourthree}{0.979~[0.961, 0.997]}  % experiments/FIGS/data/paired_v1err/mass/ratios.json :: ratios[5].ratio, ratios[5].ci95 :: sha256 d877a7f087260f64 :: 17 over 43, resolution mlp sigma_eff: paired geometric mean [95% interval]
\newcommand{\PairedMassResMlpOnesixtwoMassOverOnesixtwo}{0.80~[0.79, 0.81]}  % experiments/FIGS/data/paired_v1err/mass/ratios.json :: ratios[6].ratio, ratios[6].ci95 :: sha256 d877a7f087260f64 :: 162+mass over 162, resolution mlp sigma_eff: paired geometric mean [95% interval]
\newcommand{\PairedMassResMlpOnesevenMassOverOneseven}{0.83~[0.82, 0.84]}  % experiments/FIGS/data/paired_v1err/mass/ratios.json :: ratios[7].ratio, ratios[7].ci95 :: sha256 d877a7f087260f64 :: 17+mass over 17, resolution mlp sigma_eff: paired geometric mean [95% interval]
\newcommand{\PairedMassResRidgeOnesixtwoOverOneeighteight}{0.97~[0.96, 0.98]}  % experiments/FIGS/data/paired_v1err/mass/ratios.json :: ratios[8].ratio, ratios[8].ci95 :: sha256 d877a7f087260f64 :: 162 over 188, resolution ridge sigma_eff: paired geometric mean [95% interval]
\newcommand{\PairedMassResRidgeFourthreeOverOneeighteight}{0.95~[0.94, 0.96]}  % experiments/FIGS/data/paired_v1err/mass/ratios.json :: ratios[9].ratio, ratios[9].ci95 :: sha256 d877a7f087260f64 :: 43 over 188, resolution ridge sigma_eff: paired geometric mean [95% interval]
\newcommand{\PairedMassResRidgeOnesevenOverOneeighteight}{0.88~[0.85, 0.92]}  % experiments/FIGS/data/paired_v1err/mass/ratios.json :: ratios[10].ratio, ratios[10].ci95 :: sha256 d877a7f087260f64 :: 17 over 188, resolution ridge sigma_eff: paired geometric mean [95% interval]
\newcommand{\PairedMassResRidgeFourthreeOverOnesixtwo}{0.98~[0.97, 0.99]}  % experiments/FIGS/data/paired_v1err/mass/ratios.json :: ratios[11].ratio, ratios[11].ci95 :: sha256 d877a7f087260f64 :: 43 over 162, resolution ridge sigma_eff: paired geometric mean [95% interval]
\newcommand{\PairedMassResRidgeOnesevenOverOnesixtwo}{0.91~[0.87, 0.95]}  % experiments/FIGS/data/paired_v1err/mass/ratios.json :: ratios[12].ratio, ratios[12].ci95 :: sha256 d877a7f087260f64 :: 17 over 162, resolution ridge sigma_eff: paired geometric mean [95% interval]
\newcommand{\PairedMassResRidgeOnesevenOverFourthree}{0.93~[0.89, 0.97]}  % experiments/FIGS/data/paired_v1err/mass/ratios.json :: ratios[13].ratio, ratios[13].ci95 :: sha256 d877a7f087260f64 :: 17 over 43, resolution ridge sigma_eff: paired geometric mean [95% interval]
\newcommand{\PairedMassResRidgeOnesixtwoMassOverOnesixtwo}{0.678~[0.672, 0.684]}  % experiments/FIGS/data/paired_v1err/mass/ratios.json :: ratios[14].ratio, ratios[14].ci95 :: sha256 d877a7f087260f64 :: 162+mass over 162, resolution ridge sigma_eff: paired geometric mean [95% interval]
\newcommand{\PairedMassResRidgeOnesevenMassOverOneseven}{0.72~[0.69, 0.76]}  % experiments/FIGS/data/paired_v1err/mass/ratios.json :: ratios[15].ratio, ratios[15].ci95 :: sha256 d877a7f087260f64 :: 17+mass over 17, resolution ridge sigma_eff: paired geometric mean [95% interval]
\newcommand{\VcbNSignal}{1{,}659}  % experiments/FIGS/data/probe_ladder_vcbwindow_mlp2/sall.json :: tasks.bc_vs_rest.n_signal_test :: sha256 5b8b5c7c28602f93 :: X->bc test jets in the window
\newcommand{\VcbNBackground}{19{,}244}  % experiments/FIGS/data/probe_ladder_vcbwindow_mlp2/sall.json :: tasks.bc_vs_rest.n_background_test :: sha256 5b8b5c7c28602f93 :: bq, cs, bqq and QCD test jets in the window
\newcommand{\VcbEpsSixty}{60}  % experiments/FIGS/data/probe_ladder_vcbwindow_mlp2/sall.json :: tasks.bc_vs_rest.eps_s :: sha256 5b8b5c7c28602f93 :: signal efficiency, percent
\newcommand{\VcbOmaLinearOnesixtwo}{\ensuremath{(3.17 \pm 0.05)\times10^{-2}}}  % experiments/FIGS/data/probe_ladder_vcbwindow_mlp2/sall.json :: tasks.bc_vs_rest.arms.l162-s*.linear.log1m_auc :: sha256 5b8b5c7c28602f93 :: 1-AUC = exp of it, mean +- SD over 5 seeds
\newcommand{\VcbLogOmaLinearOnesixtwo}{\ensuremath{-}3.451}  % experiments/FIGS/data/probe_ladder_vcbwindow_mlp2/sall.json :: tasks.bc_vs_rest.arms.l162-s*.linear.log1m_auc :: sha256 5b8b5c7c28602f93 :: mean over seeds
\newcommand{\VcbRejLinearSixtyOnesixtwo}{207\,\ensuremath{\pm}\,19}  % experiments/FIGS/data/probe_ladder_vcbwindow_mlp2/sall.json :: tasks.bc_vs_rest.arms.l162-s*.linear.rejection_at['0.60'].rejection :: sha256 5b8b5c7c28602f93 :: background rejection at this signal efficiency, mean +- SD over seeds
\newcommand{\VcbBkgLeftLinearSixtyOnesixtwo}{93.4}  % experiments/FIGS/data/probe_ladder_vcbwindow_mlp2/sall.json :: tasks.bc_vs_rest.arms.l162-s*.linear.rejection_at['0.60'].n_bkg_pass :: sha256 5b8b5c7c28602f93 :: background jets passing the cut, mean over seeds
\newcommand{\VcbOmaLinearOneseven}{\ensuremath{(5.21 \pm 0.23)\times10^{-2}}}  % experiments/FIGS/data/probe_ladder_vcbwindow_mlp2/sall.json :: tasks.bc_vs_rest.arms.r16q1-s*.linear.log1m_auc :: sha256 5b8b5c7c28602f93 :: 1-AUC = exp of it, mean +- SD over 5 seeds
\newcommand{\VcbLogOmaLinearOneseven}{\ensuremath{-}2.955}  % experiments/FIGS/data/probe_ladder_vcbwindow_mlp2/sall.json :: tasks.bc_vs_rest.arms.r16q1-s*.linear.log1m_auc :: sha256 5b8b5c7c28602f93 :: mean over seeds
\newcommand{\VcbRejLinearSixtyOneseven}{67\,\ensuremath{\pm}\,14}  % experiments/FIGS/data/probe_ladder_vcbwindow_mlp2/sall.json :: tasks.bc_vs_rest.arms.r16q1-s*.linear.rejection_at['0.60'].rejection :: sha256 5b8b5c7c28602f93 :: background rejection at this signal efficiency, mean +- SD over seeds
\newcommand{\VcbBkgLeftLinearSixtyOneseven}{295.0}  % experiments/FIGS/data/probe_ladder_vcbwindow_mlp2/sall.json :: tasks.bc_vs_rest.arms.r16q1-s*.linear.rejection_at['0.60'].n_bkg_pass :: sha256 5b8b5c7c28602f93 :: background jets passing the cut, mean over seeds
\newcommand{\VcbOmaRatioLinear}{1.64}  % experiments/FIGS/data/probe_ladder_vcbwindow_mlp2/sall.json :: tasks.bc_vs_rest.arms.*.linear.log1m_auc :: sha256 5b8b5c7c28602f93 :: seed mean of 1-AUC, 17 classes over 162
\newcommand{\VcbFactorLinear}{1.64}  % experiments/FIGS/data/probe_ladder_vcbwindow_mlp2/sall.json :: tasks.bc_vs_rest.arms.*.linear.log1m_auc :: sha256 5b8b5c7c28602f93 :: exp of the difference of the seed means of log(1-AUC), 17 classes over 162
\newcommand{\VcbOmaMlpOnesixtwo}{\ensuremath{(3.102 \pm 0.023)\times10^{-2}}}  % experiments/FIGS/data/probe_ladder_vcbwindow_mlp2/sall.json :: tasks.bc_vs_rest.arms.l162-s*.mlp.log1m_auc :: sha256 5b8b5c7c28602f93 :: 1-AUC = exp of it, mean +- SD over 5 seeds
\newcommand{\VcbLogOmaMlpOnesixtwo}{\ensuremath{-}3.473}  % experiments/FIGS/data/probe_ladder_vcbwindow_mlp2/sall.json :: tasks.bc_vs_rest.arms.l162-s*.mlp.log1m_auc :: sha256 5b8b5c7c28602f93 :: mean over seeds
\newcommand{\VcbRejMlpSixtyOnesixtwo}{230\,\ensuremath{\pm}\,31}  % experiments/FIGS/data/probe_ladder_vcbwindow_mlp2/sall.json :: tasks.bc_vs_rest.arms.l162-s*.mlp.rejection_at['0.60'].rejection :: sha256 5b8b5c7c28602f93 :: background rejection at this signal efficiency, mean +- SD over seeds
\newcommand{\VcbBkgLeftMlpSixtyOnesixtwo}{84.8}  % experiments/FIGS/data/probe_ladder_vcbwindow_mlp2/sall.json :: tasks.bc_vs_rest.arms.l162-s*.mlp.rejection_at['0.60'].n_bkg_pass :: sha256 5b8b5c7c28602f93 :: background jets passing the cut, mean over seeds
\newcommand{\VcbOmaMlpOneseven}{\ensuremath{(4.66 \pm 0.19)\times10^{-2}}}  % experiments/FIGS/data/probe_ladder_vcbwindow_mlp2/sall.json :: tasks.bc_vs_rest.arms.r16q1-s*.mlp.log1m_auc :: sha256 5b8b5c7c28602f93 :: 1-AUC = exp of it, mean +- SD over 5 seeds
\newcommand{\VcbLogOmaMlpOneseven}{\ensuremath{-}3.066}  % experiments/FIGS/data/probe_ladder_vcbwindow_mlp2/sall.json :: tasks.bc_vs_rest.arms.r16q1-s*.mlp.log1m_auc :: sha256 5b8b5c7c28602f93 :: mean over seeds
\newcommand{\VcbRejMlpSixtyOneseven}{78\,\ensuremath{\pm}\,13}  % experiments/FIGS/data/probe_ladder_vcbwindow_mlp2/sall.json :: tasks.bc_vs_rest.arms.r16q1-s*.mlp.rejection_at['0.60'].rejection :: sha256 5b8b5c7c28602f93 :: background rejection at this signal efficiency, mean +- SD over seeds
\newcommand{\VcbBkgLeftMlpSixtyOneseven}{250.8}  % experiments/FIGS/data/probe_ladder_vcbwindow_mlp2/sall.json :: tasks.bc_vs_rest.arms.r16q1-s*.mlp.rejection_at['0.60'].n_bkg_pass :: sha256 5b8b5c7c28602f93 :: background jets passing the cut, mean over seeds
\newcommand{\VcbOmaRatioMlp}{1.50}  % experiments/FIGS/data/probe_ladder_vcbwindow_mlp2/sall.json :: tasks.bc_vs_rest.arms.*.mlp.log1m_auc :: sha256 5b8b5c7c28602f93 :: seed mean of 1-AUC, 17 classes over 162
\newcommand{\VcbFactorMlp}{1.50}  % experiments/FIGS/data/probe_ladder_vcbwindow_mlp2/sall.json :: tasks.bc_vs_rest.arms.*.mlp.log1m_auc :: sha256 5b8b5c7c28602f93 :: exp of the difference of the seed means of log(1-AUC), 17 classes over 162
\newcommand{\AccLegOneOnesixtwoSonebEFour}{0.32273}  % experiments/FIGS/data/leg1_metrics.json :: summary.l162-s1b.N10000.accuracy_mean :: sha256 3f564e960b2c86a7 :: wave 1, superseded; 3 fine-tuning seeds
\newcommand{\AccSdLegOneOnesixtwoSonebEFour}{0.01894}  % experiments/FIGS/data/leg1_metrics.json :: summary.l162-s1b.N10000.accuracy_sd :: sha256 3f564e960b2c86a7 :: wave 1, superseded
\newcommand{\AccLegOneOnesixtwoSonebEFive}{0.53506}  % experiments/FIGS/data/leg1_metrics.json :: summary.l162-s1b.N100000.accuracy_mean :: sha256 3f564e960b2c86a7 :: wave 1, superseded; 3 fine-tuning seeds
\newcommand{\AccSdLegOneOnesixtwoSonebEFive}{0.00044}  % experiments/FIGS/data/leg1_metrics.json :: summary.l162-s1b.N100000.accuracy_sd :: sha256 3f564e960b2c86a7 :: wave 1, superseded
\newcommand{\AccLegOneOnesixtwoSonebESix}{0.54337}  % experiments/FIGS/data/leg1_metrics.json :: summary.l162-s1b.N1000000.accuracy_mean :: sha256 3f564e960b2c86a7 :: wave 1, superseded; 3 fine-tuning seeds
\newcommand{\AccSdLegOneOnesixtwoSonebESix}{0.00074}  % experiments/FIGS/data/leg1_metrics.json :: summary.l162-s1b.N1000000.accuracy_sd :: sha256 3f564e960b2c86a7 :: wave 1, superseded
\newcommand{\AccLegOneOnesevenStwoEFour}{0.22684}  % experiments/FIGS/data/leg1_metrics.json :: summary.r16q1-s2.N10000.accuracy_mean :: sha256 3f564e960b2c86a7 :: wave 1, superseded; 3 fine-tuning seeds
\newcommand{\AccSdLegOneOnesevenStwoEFour}{0.00805}  % experiments/FIGS/data/leg1_metrics.json :: summary.r16q1-s2.N10000.accuracy_sd :: sha256 3f564e960b2c86a7 :: wave 1, superseded
\newcommand{\AccLegOneOnesevenStwoEFive}{0.49650}  % experiments/FIGS/data/leg1_metrics.json :: summary.r16q1-s2.N100000.accuracy_mean :: sha256 3f564e960b2c86a7 :: wave 1, superseded; 3 fine-tuning seeds
\newcommand{\AccSdLegOneOnesevenStwoEFive}{0.00110}  % experiments/FIGS/data/leg1_metrics.json :: summary.r16q1-s2.N100000.accuracy_sd :: sha256 3f564e960b2c86a7 :: wave 1, superseded
\newcommand{\AccLegOneOnesevenStwoESix}{0.51969}  % experiments/FIGS/data/leg1_metrics.json :: summary.r16q1-s2.N1000000.accuracy_mean :: sha256 3f564e960b2c86a7 :: wave 1, superseded; 3 fine-tuning seeds
\newcommand{\AccSdLegOneOnesevenStwoESix}{0.00072}  % experiments/FIGS/data/leg1_metrics.json :: summary.r16q1-s2.N1000000.accuracy_sd :: sha256 3f564e960b2c86a7 :: wave 1, superseded
\newcommand{\AccLegOneOnesevenSthreeEFour}{0.24203}  % experiments/FIGS/data/leg1_metrics.json :: summary.r16q1-s3.N10000.accuracy_mean :: sha256 3f564e960b2c86a7 :: wave 1, superseded; 3 fine-tuning seeds
\newcommand{\AccSdLegOneOnesevenSthreeEFour}{0.00213}  % experiments/FIGS/data/leg1_metrics.json :: summary.r16q1-s3.N10000.accuracy_sd :: sha256 3f564e960b2c86a7 :: wave 1, superseded
\newcommand{\AccLegOneOnesevenSthreeEFive}{0.49313}  % experiments/FIGS/data/leg1_metrics.json :: summary.r16q1-s3.N100000.accuracy_mean :: sha256 3f564e960b2c86a7 :: wave 1, superseded; 3 fine-tuning seeds
\newcommand{\AccSdLegOneOnesevenSthreeEFive}{0.00294}  % experiments/FIGS/data/leg1_metrics.json :: summary.r16q1-s3.N100000.accuracy_sd :: sha256 3f564e960b2c86a7 :: wave 1, superseded
\newcommand{\AccLegOneOnesevenSthreeESix}{0.51804}  % experiments/FIGS/data/leg1_metrics.json :: summary.r16q1-s3.N1000000.accuracy_mean :: sha256 3f564e960b2c86a7 :: wave 1, superseded; 3 fine-tuning seeds
\newcommand{\AccSdLegOneOnesevenSthreeESix}{0.00049}  % experiments/FIGS/data/leg1_metrics.json :: summary.r16q1-s3.N1000000.accuracy_sd :: sha256 3f564e960b2c86a7 :: wave 1, superseded
\newcommand{\AccLegOneOnesevenSfourEFour}{0.24487}  % experiments/FIGS/data/leg1_metrics.json :: summary.r16q1-s4.N10000.accuracy_mean :: sha256 3f564e960b2c86a7 :: wave 1, superseded; 3 fine-tuning seeds
\newcommand{\AccSdLegOneOnesevenSfourEFour}{0.00645}  % experiments/FIGS/data/leg1_metrics.json :: summary.r16q1-s4.N10000.accuracy_sd :: sha256 3f564e960b2c86a7 :: wave 1, superseded
\newcommand{\AccLegOneOnesevenSfourEFive}{0.49510}  % experiments/FIGS/data/leg1_metrics.json :: summary.r16q1-s4.N100000.accuracy_mean :: sha256 3f564e960b2c86a7 :: wave 1, superseded; 3 fine-tuning seeds
\newcommand{\AccSdLegOneOnesevenSfourEFive}{0.00133}  % experiments/FIGS/data/leg1_metrics.json :: summary.r16q1-s4.N100000.accuracy_sd :: sha256 3f564e960b2c86a7 :: wave 1, superseded
\newcommand{\AccLegOneOnesevenSfourESix}{0.51872}  % experiments/FIGS/data/leg1_metrics.json :: summary.r16q1-s4.N1000000.accuracy_mean :: sha256 3f564e960b2c86a7 :: wave 1, superseded; 3 fine-tuning seeds
\newcommand{\AccSdLegOneOnesevenSfourESix}{0.00049}  % experiments/FIGS/data/leg1_metrics.json :: summary.r16q1-s4.N1000000.accuracy_sd :: sha256 3f564e960b2c86a7 :: wave 1, superseded
\newcommand{\AccLegOneScratchEFour}{0.31369}  % experiments/FIGS/data/leg1_metrics.json :: summary.scratch.N10000.accuracy_mean :: sha256 3f564e960b2c86a7 :: wave 1, superseded; 3 fine-tuning seeds
\newcommand{\AccSdLegOneScratchEFour}{0.00712}  % experiments/FIGS/data/leg1_metrics.json :: summary.scratch.N10000.accuracy_sd :: sha256 3f564e960b2c86a7 :: wave 1, superseded
\newcommand{\AccLegOneScratchEFive}{0.42015}  % experiments/FIGS/data/leg1_metrics.json :: summary.scratch.N100000.accuracy_mean :: sha256 3f564e960b2c86a7 :: wave 1, superseded; 3 fine-tuning seeds
\newcommand{\AccSdLegOneScratchEFive}{0.00267}  % experiments/FIGS/data/leg1_metrics.json :: summary.scratch.N100000.accuracy_sd :: sha256 3f564e960b2c86a7 :: wave 1, superseded
\newcommand{\AccLegOneScratchESix}{0.48621}  % experiments/FIGS/data/leg1_metrics.json :: summary.scratch.N1000000.accuracy_mean :: sha256 3f564e960b2c86a7 :: wave 1, superseded; 3 fine-tuning seeds
\newcommand{\AccSdLegOneScratchESix}{0.00100}  % experiments/FIGS/data/leg1_metrics.json :: summary.scratch.N1000000.accuracy_sd :: sha256 3f564e960b2c86a7 :: wave 1, superseded
\newcommand{\AccLegOneSophonPublicEFour}{0.21613}  % experiments/FIGS/data/leg1_metrics.json :: summary.sophon-public.N10000.accuracy_mean :: sha256 3f564e960b2c86a7 :: wave 1, superseded; 3 fine-tuning seeds
\newcommand{\AccSdLegOneSophonPublicEFour}{0.00734}  % experiments/FIGS/data/leg1_metrics.json :: summary.sophon-public.N10000.accuracy_sd :: sha256 3f564e960b2c86a7 :: wave 1, superseded
\newcommand{\AccLegOneSophonPublicEFive}{0.53187}  % experiments/FIGS/data/leg1_metrics.json :: summary.sophon-public.N100000.accuracy_mean :: sha256 3f564e960b2c86a7 :: wave 1, superseded; 3 fine-tuning seeds
\newcommand{\AccSdLegOneSophonPublicEFive}{0.00231}  % experiments/FIGS/data/leg1_metrics.json :: summary.sophon-public.N100000.accuracy_sd :: sha256 3f564e960b2c86a7 :: wave 1, superseded
\newcommand{\AccLegOneSophonPublicESix}{0.54273}  % experiments/FIGS/data/leg1_metrics.json :: summary.sophon-public.N1000000.accuracy_mean :: sha256 3f564e960b2c86a7 :: wave 1, superseded; 3 fine-tuning seeds
\newcommand{\AccSdLegOneSophonPublicESix}{0.00064}  % experiments/FIGS/data/leg1_metrics.json :: summary.sophon-public.N1000000.accuracy_sd :: sha256 3f564e960b2c86a7 :: wave 1, superseded
\newcommand{\AccLegTwoOnesixtwoSonebEFour}{0.78606}  % experiments/FIGS/data/leg2_metrics.json :: summary.l162-s1b.N10000.accuracy_mean :: sha256 701c58c3240274bb :: wave 1, superseded; 3 fine-tuning seeds
\newcommand{\AccSdLegTwoOnesixtwoSonebEFour}{0.00111}  % experiments/FIGS/data/leg2_metrics.json :: summary.l162-s1b.N10000.accuracy_sd :: sha256 701c58c3240274bb :: wave 1, superseded
\newcommand{\AccLegTwoOnesixtwoSonebEFive}{0.83598}  % experiments/FIGS/data/leg2_metrics.json :: summary.l162-s1b.N100000.accuracy_mean :: sha256 701c58c3240274bb :: wave 1, superseded; 3 fine-tuning seeds
\newcommand{\AccSdLegTwoOnesixtwoSonebEFive}{0.00037}  % experiments/FIGS/data/leg2_metrics.json :: summary.l162-s1b.N100000.accuracy_sd :: sha256 701c58c3240274bb :: wave 1, superseded
\newcommand{\AccLegTwoOnesixtwoSonebESix}{0.84704}  % experiments/FIGS/data/leg2_metrics.json :: summary.l162-s1b.N1000000.accuracy_mean :: sha256 701c58c3240274bb :: wave 1, superseded; 3 fine-tuning seeds
\newcommand{\AccSdLegTwoOnesixtwoSonebESix}{0.00027}  % experiments/FIGS/data/leg2_metrics.json :: summary.l162-s1b.N1000000.accuracy_sd :: sha256 701c58c3240274bb :: wave 1, superseded
\newcommand{\AccLegTwoOnesevenStwoEFour}{0.72156}  % experiments/FIGS/data/leg2_metrics.json :: summary.r16q1-s2.N10000.accuracy_mean :: sha256 701c58c3240274bb :: wave 1, superseded; 3 fine-tuning seeds
\newcommand{\AccSdLegTwoOnesevenStwoEFour}{0.01049}  % experiments/FIGS/data/leg2_metrics.json :: summary.r16q1-s2.N10000.accuracy_sd :: sha256 701c58c3240274bb :: wave 1, superseded
\newcommand{\AccLegTwoOnesevenStwoEFive}{0.82614}  % experiments/FIGS/data/leg2_metrics.json :: summary.r16q1-s2.N100000.accuracy_mean :: sha256 701c58c3240274bb :: wave 1, superseded; 3 fine-tuning seeds
\newcommand{\AccSdLegTwoOnesevenStwoEFive}{0.00058}  % experiments/FIGS/data/leg2_metrics.json :: summary.r16q1-s2.N100000.accuracy_sd :: sha256 701c58c3240274bb :: wave 1, superseded
\newcommand{\AccLegTwoOnesevenStwoESix}{0.83943}  % experiments/FIGS/data/leg2_metrics.json :: summary.r16q1-s2.N1000000.accuracy_mean :: sha256 701c58c3240274bb :: wave 1, superseded; 3 fine-tuning seeds
\newcommand{\AccSdLegTwoOnesevenStwoESix}{0.00052}  % experiments/FIGS/data/leg2_metrics.json :: summary.r16q1-s2.N1000000.accuracy_sd :: sha256 701c58c3240274bb :: wave 1, superseded
\newcommand{\AccLegTwoOnesevenSthreeEFour}{0.72103}  % experiments/FIGS/data/leg2_metrics.json :: summary.r16q1-s3.N10000.accuracy_mean :: sha256 701c58c3240274bb :: wave 1, superseded; 3 fine-tuning seeds
\newcommand{\AccSdLegTwoOnesevenSthreeEFour}{0.01651}  % experiments/FIGS/data/leg2_metrics.json :: summary.r16q1-s3.N10000.accuracy_sd :: sha256 701c58c3240274bb :: wave 1, superseded
\newcommand{\AccLegTwoOnesevenSthreeEFive}{0.82629}  % experiments/FIGS/data/leg2_metrics.json :: summary.r16q1-s3.N100000.accuracy_mean :: sha256 701c58c3240274bb :: wave 1, superseded; 3 fine-tuning seeds
\newcommand{\AccSdLegTwoOnesevenSthreeEFive}{0.00035}  % experiments/FIGS/data/leg2_metrics.json :: summary.r16q1-s3.N100000.accuracy_sd :: sha256 701c58c3240274bb :: wave 1, superseded
\newcommand{\AccLegTwoOnesevenSthreeESix}{0.83944}  % experiments/FIGS/data/leg2_metrics.json :: summary.r16q1-s3.N1000000.accuracy_mean :: sha256 701c58c3240274bb :: wave 1, superseded; 3 fine-tuning seeds
\newcommand{\AccSdLegTwoOnesevenSthreeESix}{0.00019}  % experiments/FIGS/data/leg2_metrics.json :: summary.r16q1-s3.N1000000.accuracy_sd :: sha256 701c58c3240274bb :: wave 1, superseded
\newcommand{\AccLegTwoOnesevenSfourEFour}{0.72953}  % experiments/FIGS/data/leg2_metrics.json :: summary.r16q1-s4.N10000.accuracy_mean :: sha256 701c58c3240274bb :: wave 1, superseded; 3 fine-tuning seeds
\newcommand{\AccSdLegTwoOnesevenSfourEFour}{0.01069}  % experiments/FIGS/data/leg2_metrics.json :: summary.r16q1-s4.N10000.accuracy_sd :: sha256 701c58c3240274bb :: wave 1, superseded
\newcommand{\AccLegTwoOnesevenSfourEFive}{0.82564}  % experiments/FIGS/data/leg2_metrics.json :: summary.r16q1-s4.N100000.accuracy_mean :: sha256 701c58c3240274bb :: wave 1, superseded; 3 fine-tuning seeds
\newcommand{\AccSdLegTwoOnesevenSfourEFive}{0.00041}  % experiments/FIGS/data/leg2_metrics.json :: summary.r16q1-s4.N100000.accuracy_sd :: sha256 701c58c3240274bb :: wave 1, superseded
\newcommand{\AccLegTwoOnesevenSfourESix}{0.83909}  % experiments/FIGS/data/leg2_metrics.json :: summary.r16q1-s4.N1000000.accuracy_mean :: sha256 701c58c3240274bb :: wave 1, superseded; 3 fine-tuning seeds
\newcommand{\AccSdLegTwoOnesevenSfourESix}{0.00037}  % experiments/FIGS/data/leg2_metrics.json :: summary.r16q1-s4.N1000000.accuracy_sd :: sha256 701c58c3240274bb :: wave 1, superseded
\newcommand{\AccLegTwoRefEonearmsSoneRef}{0.81770}  % experiments/FIGS/data/leg2_metrics.json :: summary.ref_e1arms-s1.ref.accuracy_mean :: sha256 701c58c3240274bb :: wave 1, superseded; 1 fine-tuning seeds
\newcommand{\AccSdLegTwoRefEonearmsSoneRef}{---}  % experiments/FIGS/data/leg2_metrics.json :: summary.ref_e1arms-s1.ref.accuracy_sd :: sha256 701c58c3240274bb :: wave 1, superseded
\newcommand{\AccLegTwoRefEonearmsStwoRef}{0.81602}  % experiments/FIGS/data/leg2_metrics.json :: summary.ref_e1arms-s2.ref.accuracy_mean :: sha256 701c58c3240274bb :: wave 1, superseded; 1 fine-tuning seeds
\newcommand{\AccSdLegTwoRefEonearmsStwoRef}{---}  % experiments/FIGS/data/leg2_metrics.json :: summary.ref_e1arms-s2.ref.accuracy_sd :: sha256 701c58c3240274bb :: wave 1, superseded
\newcommand{\AccLegTwoRefEonearmsSthreeRef}{0.81625}  % experiments/FIGS/data/leg2_metrics.json :: summary.ref_e1arms-s3.ref.accuracy_mean :: sha256 701c58c3240274bb :: wave 1, superseded; 1 fine-tuning seeds
\newcommand{\AccSdLegTwoRefEonearmsSthreeRef}{---}  % experiments/FIGS/data/leg2_metrics.json :: summary.ref_e1arms-s3.ref.accuracy_sd :: sha256 701c58c3240274bb :: wave 1, superseded
\newcommand{\AccLegTwoScratchEFour}{0.68230}  % experiments/FIGS/data/leg2_metrics.json :: summary.scratch.N10000.accuracy_mean :: sha256 701c58c3240274bb :: wave 1, superseded; 3 fine-tuning seeds
\newcommand{\AccSdLegTwoScratchEFour}{0.00282}  % experiments/FIGS/data/leg2_metrics.json :: summary.scratch.N10000.accuracy_sd :: sha256 701c58c3240274bb :: wave 1, superseded
\newcommand{\AccLegTwoScratchEFive}{0.77769}  % experiments/FIGS/data/leg2_metrics.json :: summary.scratch.N100000.accuracy_mean :: sha256 701c58c3240274bb :: wave 1, superseded; 3 fine-tuning seeds
\newcommand{\AccSdLegTwoScratchEFive}{0.00208}  % experiments/FIGS/data/leg2_metrics.json :: summary.scratch.N100000.accuracy_sd :: sha256 701c58c3240274bb :: wave 1, superseded
\newcommand{\AccLegTwoScratchESix}{0.82303}  % experiments/FIGS/data/leg2_metrics.json :: summary.scratch.N1000000.accuracy_mean :: sha256 701c58c3240274bb :: wave 1, superseded; 3 fine-tuning seeds
\newcommand{\AccSdLegTwoScratchESix}{0.00134}  % experiments/FIGS/data/leg2_metrics.json :: summary.scratch.N1000000.accuracy_sd :: sha256 701c58c3240274bb :: wave 1, superseded
\newcommand{\AccLegTwoSophonPublicEFour}{0.72739}  % experiments/FIGS/data/leg2_metrics.json :: summary.sophon-public.N10000.accuracy_mean :: sha256 701c58c3240274bb :: wave 1, superseded; 3 fine-tuning seeds
\newcommand{\AccSdLegTwoSophonPublicEFour}{0.00470}  % experiments/FIGS/data/leg2_metrics.json :: summary.sophon-public.N10000.accuracy_sd :: sha256 701c58c3240274bb :: wave 1, superseded
\newcommand{\AccLegTwoSophonPublicEFive}{0.83492}  % experiments/FIGS/data/leg2_metrics.json :: summary.sophon-public.N100000.accuracy_mean :: sha256 701c58c3240274bb :: wave 1, superseded; 3 fine-tuning seeds
\newcommand{\AccSdLegTwoSophonPublicEFive}{0.00049}  % experiments/FIGS/data/leg2_metrics.json :: summary.sophon-public.N100000.accuracy_sd :: sha256 701c58c3240274bb :: wave 1, superseded
\newcommand{\AccLegTwoSophonPublicESix}{0.84632}  % experiments/FIGS/data/leg2_metrics.json :: summary.sophon-public.N1000000.accuracy_mean :: sha256 701c58c3240274bb :: wave 1, superseded; 3 fine-tuning seeds
\newcommand{\AccSdLegTwoSophonPublicESix}{0.00031}  % experiments/FIGS/data/leg2_metrics.json :: summary.sophon-public.N1000000.accuracy_sd :: sha256 701c58c3240274bb :: wave 1, superseded
\newcommand{\FtBudgetSteps}{20{,}000}  % scripts/build_ft_budget_jobs.py :: EPOCHS * STEPS :: sha256 f2ed9faa8e3bd622 :: the largest number of optimizer steps below the full training set
\newcommand{\FtLongSteps}{80{,}000}  % scripts/build_ft_budget_jobs.py :: LONG_EPOCHS * STEPS :: sha256 f2ed9faa8e3bd622 :: the largest number of optimizer steps of the long budget at LONG_N jets (JetClass, JetClass-II)
\newcommand{\DesignNTrainSelected}{109{,}938{,}036}  % experiments/FIGS/data/v2/jc2_train/counts.json :: n_selected :: sha256 b0b5a8b63aa808c0 :: selected jets of the pretrain_train files (experiments/FT/count_jc2_train.py)
\newcommand{\FtBudgetEval}{500}  % scripts/build_ft_budget_jobs.py :: STEPS :: sha256 f2ed9faa8e3bd622 :: steps between validation passes at and above SMALL_N training jets
\newcommand{\FtBudgetEvalSmall}{50}  % scripts/build_ft_budget_jobs.py :: SMALL_STEPS :: sha256 f2ed9faa8e3bd622 :: steps between validation passes below SMALL_N training jets
\newcommand{\FtBudgetEvalSmallBelow}{\ensuremath{10^{5}}}  % scripts/build_ft_budget_jobs.py :: SMALL_N :: sha256 f2ed9faa8e3bd622 :: training jets below which the validation runs every SMALL_STEPS steps
\newcommand{\FtBudgetPatience}{10}  % scripts/build_ft_budget_jobs.py :: PATIENCE :: sha256 f2ed9faa8e3bd622 :: validation passes without improvement before training stops
\newcommand{\FtHeadMult}{50}  % scripts/build_ft_budget_jobs.py :: HEAD_MULT :: sha256 f2ed9faa8e3bd622 :: the head's learning-rate multiplier
\newcommand{\FtHeadCap}{\ensuremath{5\times10^{-3}}}  % scripts/build_ft_budget_jobs.py :: HEAD_CAP :: sha256 f2ed9faa8e3bd622 :: the head's largest learning rate
\newcommand{\FtBudgetBatch}{512}  % scripts/build_ft_budget_jobs.py :: BATCH :: sha256 f2ed9faa8e3bd622
\newcommand{\FtBudgetWeightDecay}{0.01}  % scripts/build_ft_budget_jobs.py :: --optimizer-option weight_decay (every cell) :: sha256 f2ed9faa8e3bd622
\newcommand{\FtSweepPre}{\ensuremath{10^{-4}}, \ensuremath{3\times10^{-4}} and \ensuremath{10^{-3}}}  % scripts/build_ft_budget_jobs.py :: SWEEP_PRE :: sha256 f2ed9faa8e3bd622 :: backbone rates of the sweep
\newcommand{\FtSweepScratch}{\ensuremath{3\times10^{-4}}, \ensuremath{10^{-3}} and \ensuremath{3\times10^{-3}}}  % scripts/build_ft_budget_jobs.py :: SWEEP_SCRATCH :: sha256 f2ed9faa8e3bd622 :: from-scratch rates of the sweep
\newcommand{\FtSweepN}{10{,}000}  % scripts/build_ft_budget_jobs.py :: SWEEP_N :: sha256 f2ed9faa8e3bd622 :: training jets of the sweep
\newcommand{\FtValJets}{200{,}000}  % scripts/build_ft_budget_jobs.py :: VAL_N :: sha256 f2ed9faa8e3bd622 :: the fixed validation set
\newcommand{\FtValJetsBal}{20{,}480}  % scripts/build_ft_budget_jobs.py :: VAL_SIZE[jc2b] :: sha256 f2ed9faa8e3bd622 :: the balanced JetClass-II validation set
\newcommand{\LpVfourCMax}{\ensuremath{10^{4}}}  % experiments/EVAL/probe_bal.py :: C_GRID_V4 :: sha256 a7804216e3179d9d :: largest inverse regularisation of the v4 probes
\newcommand{\JcBalTestPerClass}{5{,}000}  % experiments/FT/stage_jc2_balanced.py :: TEST_PER_CLASS :: sha256 0260341d5ed75fec
\newcommand{\JcBalValPerClass}{109}  % experiments/FT/stage_jc2_balanced.py :: VAL_PER_CLASS :: sha256 0260341d5ed75fec
\newcommand{\JcBalPairTrainPerSide}{\ensuremath{10^{5}}}  % experiments/FT/stage_jc2_balanced.py :: PAIR_TRAIN_PER_SIDE :: sha256 0260341d5ed75fec
\newcommand{\JcBalPairTestPerSide}{\ensuremath{2\times10^{5}}}  % experiments/FT/stage_jc2_balanced.py :: PAIR_TEST_PER_SIDE :: sha256 0260341d5ed75fec
\newcommand{\JcBalNShort}{11}  % experiments/FIGS/data/v2/jc2bal/census.json :: counts.ft_test :: sha256 a4379ee0817254cf :: classes with fewer than 5000 selected test jets
\newcommand{\JcBalRarestTest}{224}  % experiments/FIGS/data/v2/jc2bal/census.json :: counts.ft_test (minimum) :: sha256 a4379ee0817254cf :: label_QCD_bcs
\newcommand{\JcBalShortTrainESix}{11}  % experiments/FIGS/data/v2/jc2bal/census.json :: counts.ft_heldout :: sha256 a4379ee0817254cf :: classes with fewer held-out jets than the 1e6 quota
\newcommand{\JcBalNTest}{903{,}922}  % experiments/FIGS/data/v2/jc2bal/manifest.json :: outputs[test.parquet].rows :: sha256 9c7775be6ff56ac3 :: jets of the balanced test set
\newcommand{\FtLrPre}{\ensuremath{10^{-3}}}  % experiments/FIGS/data/v2/budget/lrsweep_metrics.json :: argmax of the summed validation metric, pretrained :: sha256 7cdc148054ca902c
\newcommand{\FtLrScratchTop}{\ensuremath{10^{-3}}}  % experiments/FIGS/data/v2/budget/lrsweep_metrics.json :: argmax of the validation metric on top, from scratch :: sha256 7cdc148054ca902c :: agrees with LR of the launched job-ft-budget-bench, -jc1, -jc2b, -scratch-full and -long-scratch-full from-scratch cells
\newcommand{\FtLrScratchQg}{\ensuremath{3\times10^{-3}}}  % experiments/FIGS/data/v2/budget/lrsweep_metrics.json :: argmax of the validation metric on quark/gluon, from scratch :: sha256 7cdc148054ca902c :: agrees with LR of the launched job-ft-budget-bench and -scratch-full from-scratch cells
\newcommand{\MlstNParamsFt}{2{,}209{,}270}  % experiments/FIGS/mlst_v4.py :: our_params(2, 17) :: sha256 fa536b4866ffd651 :: trainable parameters of our fine-tuned network, two outputs (top, quark/gluon)
\newcommand{\MlstNParamsFtKin}{2{,}207{,}950}  % experiments/FIGS/mlst_v4.py :: our_params(2, 7) :: sha256 fa536b4866ffd651 :: trainable parameters of our fine-tuned network, two outputs, the 7 kinematic inputs
\newcommand{\MlstNParamsJc}{2{,}213{,}374}  % experiments/FIGS/mlst_v4.py :: our_params(10, 17) :: sha256 fa536b4866ffd651 :: trainable parameters of our fine-tuned network, the ten JetClass classes
\newcommand{\MlstProbeRunsDone}{27}  % experiments/FIGS/data/v2/index.json :: runs with a linear-probe fit (count) :: sha256 c6c4138ee5c86423 :: of the paper's runs
\newcommand{\MlstPaperRuns}{27}  % configs/arms/v2_paper_runs.txt :: runs (count) :: sha256 b04562b14ea8cbc7
\newcommand{\MlstSdRelUnc}{50\%}  % configs/arms/v2_paper_runs.txt :: relative standard error of a standard deviation from the runs of a hierarchy level, 1/sqrt(2(n-1)) :: sha256 b04562b14ea8cbc7
\newcommand{\MlstCmpJcTwoResOverFull}{1.026\,\ensuremath{\pm}\,0.005}  % experiments/FIGS/data/v2/index.json :: linear probe jc2_188 N=1000000: (1 - macro_auc_ovr) of L162 over L188, run by run :: sha256 c6c4138ee5c86423
\newcommand{\MlstCmpJcTwoResOverFullPr}{2.6\% more errors}  % experiments/FIGS/data/v2/index.json :: linear probe jc2_188 N=1000000: (1 - macro_auc_ovr) of L162 over L188, run by run :: sha256 c6c4138ee5c86423 :: prose twin: two significant figures, direction in words
\newcommand{\MlstCmpJcTwoMidOverFull}{1.227\,\ensuremath{\pm}\,0.009}  % experiments/FIGS/data/v2/index.json :: linear probe jc2_188 N=1000000: (1 - macro_auc_ovr) of R42_Q1 over L188, run by run :: sha256 c6c4138ee5c86423
\newcommand{\MlstCmpJcTwoMidOverFullPr}{23\% more errors}  % experiments/FIGS/data/v2/index.json :: linear probe jc2_188 N=1000000: (1 - macro_auc_ovr) of R42_Q1 over L188, run by run :: sha256 c6c4138ee5c86423 :: prose twin: two significant figures, direction in words
\newcommand{\MlstCmpJcTwoCoarseOverFull}{1.557\,\ensuremath{\pm}\,0.019}  % experiments/FIGS/data/v2/index.json :: linear probe jc2_188 N=1000000: (1 - macro_auc_ovr) of R16_Q1 over L188, run by run :: sha256 c6c4138ee5c86423
\newcommand{\MlstCmpJcTwoCoarseOverFullPr}{56\% more errors}  % experiments/FIGS/data/v2/index.json :: linear probe jc2_188 N=1000000: (1 - macro_auc_ovr) of R16_Q1 over L188, run by run :: sha256 c6c4138ee5c86423 :: prose twin: two significant figures, direction in words
\newcommand{\MlstCmpJcTwoResMassOverRes}{1.022 and 1.016}  % experiments/FIGS/data/v2/index.json :: linear probe jc2_188 N=1000000: (1 - macro_auc_ovr) of L162_MASS over L162, run by run :: sha256 c6c4138ee5c86423
\newcommand{\MlstCmpJcTwoResMassOverResPr}{equal within 3\%}  % experiments/FIGS/data/v2/index.json :: linear probe jc2_188 N=1000000: (1 - macro_auc_ovr) of L162_MASS over L162, run by run :: sha256 c6c4138ee5c86423 :: prose twin: two significant figures, direction in words
\newcommand{\MlstCmpJcTwoResMassOverResPrMax}{equal within 3\%}  % experiments/FIGS/data/v2/index.json :: linear probe jc2_188 N=1000000: (1 - macro_auc_ovr) of L162_MASS over L162, run by run :: sha256 c6c4138ee5c86423 :: prose twin of the run farther from 1
\newcommand{\MlstCmpJcTwoCoarseMassOverCoarse}{1.054 and 0.977}  % experiments/FIGS/data/v2/index.json :: linear probe jc2_188 N=1000000: (1 - macro_auc_ovr) of R16_Q1_MASS over R16_Q1, run by run :: sha256 c6c4138ee5c86423
\newcommand{\MlstCmpJcTwoCoarseMassOverCoarsePr}{equal within 6\%}  % experiments/FIGS/data/v2/index.json :: linear probe jc2_188 N=1000000: (1 - macro_auc_ovr) of R16_Q1_MASS over R16_Q1, run by run :: sha256 c6c4138ee5c86423 :: prose twin: two significant figures, direction in words
\newcommand{\MlstCmpJcTwoCoarseMassOverCoarsePrMax}{equal within 6\%}  % experiments/FIGS/data/v2/index.json :: linear probe jc2_188 N=1000000: (1 - macro_auc_ovr) of R16_Q1_MASS over R16_Q1, run by run :: sha256 c6c4138ee5c86423 :: prose twin of the run farther from 1
\newcommand{\MlstCmpJcTwoResMassOverFull}{1.053 and 1.037}  % experiments/FIGS/data/v2/index.json :: linear probe jc2_188 N=1000000: (1 - macro_auc_ovr) of L162_MASS over L188, run by run :: sha256 c6c4138ee5c86423
\newcommand{\MlstCmpJcTwoResMassOverFullPr}{at least 3.7\% more errors}  % experiments/FIGS/data/v2/index.json :: linear probe jc2_188 N=1000000: (1 - macro_auc_ovr) of L162_MASS over L188, run by run :: sha256 c6c4138ee5c86423 :: prose twin: two significant figures, direction in words
\newcommand{\MlstCmpJcTwoResMassOverFullPrMax}{up to 5.3\% more errors}  % experiments/FIGS/data/v2/index.json :: linear probe jc2_188 N=1000000: (1 - macro_auc_ovr) of L162_MASS over L188, run by run :: sha256 c6c4138ee5c86423 :: prose twin of the run farther from 1
\newcommand{\MlstCmpJcTwoCoarseMassOverFull}{1.617 and 1.528}  % experiments/FIGS/data/v2/index.json :: linear probe jc2_188 N=1000000: (1 - macro_auc_ovr) of R16_Q1_MASS over L188, run by run :: sha256 c6c4138ee5c86423
\newcommand{\MlstCmpJcTwoCoarseMassOverFullPr}{at least 53\% more errors}  % experiments/FIGS/data/v2/index.json :: linear probe jc2_188 N=1000000: (1 - macro_auc_ovr) of R16_Q1_MASS over L188, run by run :: sha256 c6c4138ee5c86423 :: prose twin: two significant figures, direction in words
\newcommand{\MlstCmpJcTwoCoarseMassOverFullPrMax}{up to 62\% more errors}  % experiments/FIGS/data/v2/index.json :: linear probe jc2_188 N=1000000: (1 - macro_auc_ovr) of R16_Q1_MASS over L188, run by run :: sha256 c6c4138ee5c86423 :: prose twin of the run farther from 1
\newcommand{\MlstCmpJcTwoFzeroOverFull}{1.424 and 1.462}  % experiments/FIGS/data/v2/index.json :: linear probe jc2_188 N=1000000: (1 - macro_auc_ovr) of FLAV_F0 over L188, run by run :: sha256 c6c4138ee5c86423
\newcommand{\MlstCmpJcTwoFzeroOverFullPr}{at least 42\% more errors}  % experiments/FIGS/data/v2/index.json :: linear probe jc2_188 N=1000000: (1 - macro_auc_ovr) of FLAV_F0 over L188, run by run :: sha256 c6c4138ee5c86423 :: prose twin: two significant figures, direction in words
\newcommand{\MlstCmpJcTwoFzeroOverFullPrMax}{up to 46\% more errors}  % experiments/FIGS/data/v2/index.json :: linear probe jc2_188 N=1000000: (1 - macro_auc_ovr) of FLAV_F0 over L188, run by run :: sha256 c6c4138ee5c86423 :: prose twin of the run farther from 1
\newcommand{\MlstCmpJcTwoFoneOverFull}{1.218 and 1.212}  % experiments/FIGS/data/v2/index.json :: linear probe jc2_188 N=1000000: (1 - macro_auc_ovr) of FLAV_F1 over L188, run by run :: sha256 c6c4138ee5c86423
\newcommand{\MlstCmpJcTwoFoneOverFullPr}{at least 21\% more errors}  % experiments/FIGS/data/v2/index.json :: linear probe jc2_188 N=1000000: (1 - macro_auc_ovr) of FLAV_F1 over L188, run by run :: sha256 c6c4138ee5c86423 :: prose twin: two significant figures, direction in words
\newcommand{\MlstCmpJcTwoFoneOverFullPrMax}{up to 22\% more errors}  % experiments/FIGS/data/v2/index.json :: linear probe jc2_188 N=1000000: (1 - macro_auc_ovr) of FLAV_F1 over L188, run by run :: sha256 c6c4138ee5c86423 :: prose twin of the run farther from 1
\newcommand{\MlstCmpJcTwoFonerOverFull}{1.206 and 1.208}  % experiments/FIGS/data/v2/index.json :: linear probe jc2_188 N=1000000: (1 - macro_auc_ovr) of FLAV_F1R over L188, run by run :: sha256 c6c4138ee5c86423
\newcommand{\MlstCmpJcTwoFonerOverFullPr}{at least 21\% more errors}  % experiments/FIGS/data/v2/index.json :: linear probe jc2_188 N=1000000: (1 - macro_auc_ovr) of FLAV_F1R over L188, run by run :: sha256 c6c4138ee5c86423 :: prose twin: two significant figures, direction in words
\newcommand{\MlstCmpJcTwoFonerOverFullPrMax}{up to 21\% more errors}  % experiments/FIGS/data/v2/index.json :: linear probe jc2_188 N=1000000: (1 - macro_auc_ovr) of FLAV_F1R over L188, run by run :: sha256 c6c4138ee5c86423 :: prose twin of the run farther from 1
\newcommand{\MlstCmpJcTwoFzeroOverCoarse}{0.9275 and 0.9345}  % experiments/FIGS/data/v2/index.json :: linear probe jc2_188 N=1000000: (1 - macro_auc_ovr) of FLAV_F0 over R16_Q1, run by run :: sha256 c6c4138ee5c86423
\newcommand{\MlstCmpJcTwoFzeroOverCoarsePr}{at least 6.5\% fewer errors}  % experiments/FIGS/data/v2/index.json :: linear probe jc2_188 N=1000000: (1 - macro_auc_ovr) of FLAV_F0 over R16_Q1, run by run :: sha256 c6c4138ee5c86423 :: prose twin: two significant figures, direction in words
\newcommand{\MlstCmpJcTwoFzeroOverCoarsePrMax}{up to 7.2\% fewer errors}  % experiments/FIGS/data/v2/index.json :: linear probe jc2_188 N=1000000: (1 - macro_auc_ovr) of FLAV_F0 over R16_Q1, run by run :: sha256 c6c4138ee5c86423 :: prose twin of the run farther from 1
\newcommand{\MlstCmpJcTwoFoneOverCoarse}{0.794 and 0.775}  % experiments/FIGS/data/v2/index.json :: linear probe jc2_188 N=1000000: (1 - macro_auc_ovr) of FLAV_F1 over R16_Q1, run by run :: sha256 c6c4138ee5c86423
\newcommand{\MlstCmpJcTwoFoneOverCoarsePr}{at least 21\% fewer errors}  % experiments/FIGS/data/v2/index.json :: linear probe jc2_188 N=1000000: (1 - macro_auc_ovr) of FLAV_F1 over R16_Q1, run by run :: sha256 c6c4138ee5c86423 :: prose twin: two significant figures, direction in words
\newcommand{\MlstCmpJcTwoFoneOverCoarsePrMax}{up to 23\% fewer errors}  % experiments/FIGS/data/v2/index.json :: linear probe jc2_188 N=1000000: (1 - macro_auc_ovr) of FLAV_F1 over R16_Q1, run by run :: sha256 c6c4138ee5c86423 :: prose twin of the run farther from 1
\newcommand{\MlstCmpJcTwoFonerOverCoarse}{0.786 and 0.772}  % experiments/FIGS/data/v2/index.json :: linear probe jc2_188 N=1000000: (1 - macro_auc_ovr) of FLAV_F1R over R16_Q1, run by run :: sha256 c6c4138ee5c86423
\newcommand{\MlstCmpJcTwoFonerOverCoarsePr}{at least 21\% fewer errors}  % experiments/FIGS/data/v2/index.json :: linear probe jc2_188 N=1000000: (1 - macro_auc_ovr) of FLAV_F1R over R16_Q1, run by run :: sha256 c6c4138ee5c86423 :: prose twin: two significant figures, direction in words
\newcommand{\MlstCmpJcTwoFonerOverCoarsePrMax}{up to 23\% fewer errors}  % experiments/FIGS/data/v2/index.json :: linear probe jc2_188 N=1000000: (1 - macro_auc_ovr) of FLAV_F1R over R16_Q1, run by run :: sha256 c6c4138ee5c86423 :: prose twin of the run farther from 1
\newcommand{\MlstCmpJcTwoFoneOverFzero}{0.856 and 0.829}  % experiments/FIGS/data/v2/index.json :: linear probe jc2_188 N=1000000: (1 - macro_auc_ovr) of FLAV_F1 over FLAV_F0, run by run :: sha256 c6c4138ee5c86423
\newcommand{\MlstCmpJcTwoFoneOverFzeroPr}{at least 14\% fewer errors}  % experiments/FIGS/data/v2/index.json :: linear probe jc2_188 N=1000000: (1 - macro_auc_ovr) of FLAV_F1 over FLAV_F0, run by run :: sha256 c6c4138ee5c86423 :: prose twin: two significant figures, direction in words
\newcommand{\MlstCmpJcTwoFoneOverFzeroPrMax}{up to 17\% fewer errors}  % experiments/FIGS/data/v2/index.json :: linear probe jc2_188 N=1000000: (1 - macro_auc_ovr) of FLAV_F1 over FLAV_F0, run by run :: sha256 c6c4138ee5c86423 :: prose twin of the run farther from 1
\newcommand{\MlstCmpJcTwoFonerOverFzero}{0.847 and 0.826}  % experiments/FIGS/data/v2/index.json :: linear probe jc2_188 N=1000000: (1 - macro_auc_ovr) of FLAV_F1R over FLAV_F0, run by run :: sha256 c6c4138ee5c86423
\newcommand{\MlstCmpJcTwoFonerOverFzeroPr}{at least 15\% fewer errors}  % experiments/FIGS/data/v2/index.json :: linear probe jc2_188 N=1000000: (1 - macro_auc_ovr) of FLAV_F1R over FLAV_F0, run by run :: sha256 c6c4138ee5c86423 :: prose twin: two significant figures, direction in words
\newcommand{\MlstCmpJcTwoFonerOverFzeroPrMax}{up to 17\% fewer errors}  % experiments/FIGS/data/v2/index.json :: linear probe jc2_188 N=1000000: (1 - macro_auc_ovr) of FLAV_F1R over FLAV_F0, run by run :: sha256 c6c4138ee5c86423 :: prose twin of the run farther from 1
\newcommand{\MlstCmpJcTwoFonerOverFone}{0.9899 and 0.9966}  % experiments/FIGS/data/v2/index.json :: linear probe jc2_188 N=1000000: (1 - macro_auc_ovr) of FLAV_F1R over FLAV_F1, run by run :: sha256 c6c4138ee5c86423
\newcommand{\MlstCmpJcTwoFonerOverFonePr}{equal within 2\%}  % experiments/FIGS/data/v2/index.json :: linear probe jc2_188 N=1000000: (1 - macro_auc_ovr) of FLAV_F1R over FLAV_F1, run by run :: sha256 c6c4138ee5c86423 :: prose twin: two significant figures, direction in words
\newcommand{\MlstCmpJcTwoFonerOverFonePrMax}{equal within 2\%}  % experiments/FIGS/data/v2/index.json :: linear probe jc2_188 N=1000000: (1 - macro_auc_ovr) of FLAV_F1R over FLAV_F1, run by run :: sha256 c6c4138ee5c86423 :: prose twin of the run farther from 1
\newcommand{\MlstCmpJcTwoMidOverRes}{1.197\,\ensuremath{\pm}\,0.013}  % experiments/FIGS/data/v2/index.json :: linear probe jc2_188 N=1000000: (1 - macro_auc_ovr) of R42_Q1 over L162, run by run :: sha256 c6c4138ee5c86423
\newcommand{\MlstCmpJcTwoMidOverResPr}{20\% more errors}  % experiments/FIGS/data/v2/index.json :: linear probe jc2_188 N=1000000: (1 - macro_auc_ovr) of R42_Q1 over L162, run by run :: sha256 c6c4138ee5c86423 :: prose twin: two significant figures, direction in words
\newcommand{\MlstCmpJcTwoCoarseOverRes}{1.518\,\ensuremath{\pm}\,0.025}  % experiments/FIGS/data/v2/index.json :: linear probe jc2_188 N=1000000: (1 - macro_auc_ovr) of R16_Q1 over L162, run by run :: sha256 c6c4138ee5c86423
\newcommand{\MlstCmpJcTwoCoarseOverResPr}{52\% more errors}  % experiments/FIGS/data/v2/index.json :: linear probe jc2_188 N=1000000: (1 - macro_auc_ovr) of R16_Q1 over L162, run by run :: sha256 c6c4138ee5c86423 :: prose twin: two significant figures, direction in words
\newcommand{\MlstCmpJcTwoCoarseOverMid}{1.269\,\ensuremath{\pm}\,0.007}  % experiments/FIGS/data/v2/index.json :: linear probe jc2_188 N=1000000: (1 - macro_auc_ovr) of R16_Q1 over R42_Q1, run by run :: sha256 c6c4138ee5c86423
\newcommand{\MlstCmpJcTwoCoarseOverMidPr}{27\% more errors}  % experiments/FIGS/data/v2/index.json :: linear probe jc2_188 N=1000000: (1 - macro_auc_ovr) of R16_Q1 over R42_Q1, run by run :: sha256 c6c4138ee5c86423 :: prose twin: two significant figures, direction in words
\newcommand{\MlstCmpJcTwoRandOverFull}{1.125\,\ensuremath{\pm}\,0.008}  % experiments/FIGS/data/v2/index.json :: linear probe jc2_188 N=1000000: (1 - macro_auc_ovr) of each random partition over the L188 run of its seed, mean +- SD over the partitions :: sha256 c6c4138ee5c86423
\newcommand{\MlstCmpJcTwoRandOverFullPr}{12\% more errors}  % experiments/FIGS/data/v2/index.json :: linear probe jc2_188 N=1000000: (1 - macro_auc_ovr) of each random partition over the L188 run of its seed, mean +- SD over the partitions :: sha256 c6c4138ee5c86423 :: prose twin: two significant figures, direction in words
\newcommand{\MlstCmpJcTwoRandOverMid}{0.918\,\ensuremath{\pm}\,0.007}  % experiments/FIGS/data/v2/index.json :: linear probe jc2_188 N=1000000: (1 - macro_auc_ovr) of each random partition over the R42_Q1 run of its seed, mean +- SD over the partitions :: sha256 c6c4138ee5c86423
\newcommand{\MlstCmpJcTwoRandOverMidPr}{8.2\% fewer errors}  % experiments/FIGS/data/v2/index.json :: linear probe jc2_188 N=1000000: (1 - macro_auc_ovr) of each random partition over the R42_Q1 run of its seed, mean +- SD over the partitions :: sha256 c6c4138ee5c86423 :: prose twin: two significant figures, direction in words
\newcommand{\MlstCmpJcTwoRandOverCoarse}{0.724\,\ensuremath{\pm}\,0.008}  % experiments/FIGS/data/v2/index.json :: linear probe jc2_188 N=1000000: (1 - macro_auc_ovr) of each random partition over the R16_Q1 run of its seed, mean +- SD over the partitions :: sha256 c6c4138ee5c86423
\newcommand{\MlstCmpJcTwoRandOverCoarsePr}{28\% fewer errors}  % experiments/FIGS/data/v2/index.json :: linear probe jc2_188 N=1000000: (1 - macro_auc_ovr) of each random partition over the R16_Q1 run of its seed, mean +- SD over the partitions :: sha256 c6c4138ee5c86423 :: prose twin: two significant figures, direction in words
\newcommand{\MlstCmpJcTwoTrainedResOverFull}{1.026\,\ensuremath{\pm}\,0.005}  % experiments/FIGS/data/v2/index.json :: linear probe jc2_188 N=1000000: (1 - macro_auc_trained) of L162 over L188, run by run :: sha256 c6c4138ee5c86423
\newcommand{\MlstCmpJcTwoTrainedResOverFullPr}{2.6\% more errors}  % experiments/FIGS/data/v2/index.json :: linear probe jc2_188 N=1000000: (1 - macro_auc_trained) of L162 over L188, run by run :: sha256 c6c4138ee5c86423 :: prose twin: two significant figures, direction in words
\newcommand{\MlstCmpJcTwoTrainedMidOverFull}{1.227\,\ensuremath{\pm}\,0.009}  % experiments/FIGS/data/v2/index.json :: linear probe jc2_188 N=1000000: (1 - macro_auc_trained) of R42_Q1 over L188, run by run :: sha256 c6c4138ee5c86423
\newcommand{\MlstCmpJcTwoTrainedMidOverFullPr}{23\% more errors}  % experiments/FIGS/data/v2/index.json :: linear probe jc2_188 N=1000000: (1 - macro_auc_trained) of R42_Q1 over L188, run by run :: sha256 c6c4138ee5c86423 :: prose twin: two significant figures, direction in words
\newcommand{\MlstCmpJcTwoTrainedCoarseOverFull}{1.557\,\ensuremath{\pm}\,0.019}  % experiments/FIGS/data/v2/index.json :: linear probe jc2_188 N=1000000: (1 - macro_auc_trained) of R16_Q1 over L188, run by run :: sha256 c6c4138ee5c86423
\newcommand{\MlstCmpJcTwoTrainedCoarseOverFullPr}{56\% more errors}  % experiments/FIGS/data/v2/index.json :: linear probe jc2_188 N=1000000: (1 - macro_auc_trained) of R16_Q1 over L188, run by run :: sha256 c6c4138ee5c86423 :: prose twin: two significant figures, direction in words
\newcommand{\MlstCmpJcTwoTrainedResMassOverRes}{1.022 and 1.016}  % experiments/FIGS/data/v2/index.json :: linear probe jc2_188 N=1000000: (1 - macro_auc_trained) of L162_MASS over L162, run by run :: sha256 c6c4138ee5c86423
\newcommand{\MlstCmpJcTwoTrainedResMassOverResPr}{equal within 3\%}  % experiments/FIGS/data/v2/index.json :: linear probe jc2_188 N=1000000: (1 - macro_auc_trained) of L162_MASS over L162, run by run :: sha256 c6c4138ee5c86423 :: prose twin: two significant figures, direction in words
\newcommand{\MlstCmpJcTwoTrainedResMassOverResPrMax}{equal within 3\%}  % experiments/FIGS/data/v2/index.json :: linear probe jc2_188 N=1000000: (1 - macro_auc_trained) of L162_MASS over L162, run by run :: sha256 c6c4138ee5c86423 :: prose twin of the run farther from 1
\newcommand{\MlstCmpJcTwoTrainedCoarseMassOverCoarse}{1.054 and 0.977}  % experiments/FIGS/data/v2/index.json :: linear probe jc2_188 N=1000000: (1 - macro_auc_trained) of R16_Q1_MASS over R16_Q1, run by run :: sha256 c6c4138ee5c86423
\newcommand{\MlstCmpJcTwoTrainedCoarseMassOverCoarsePr}{equal within 6\%}  % experiments/FIGS/data/v2/index.json :: linear probe jc2_188 N=1000000: (1 - macro_auc_trained) of R16_Q1_MASS over R16_Q1, run by run :: sha256 c6c4138ee5c86423 :: prose twin: two significant figures, direction in words
\newcommand{\MlstCmpJcTwoTrainedCoarseMassOverCoarsePrMax}{equal within 6\%}  % experiments/FIGS/data/v2/index.json :: linear probe jc2_188 N=1000000: (1 - macro_auc_trained) of R16_Q1_MASS over R16_Q1, run by run :: sha256 c6c4138ee5c86423 :: prose twin of the run farther from 1
\newcommand{\MlstCmpJcTwoTrainedResMassOverFull}{1.053 and 1.037}  % experiments/FIGS/data/v2/index.json :: linear probe jc2_188 N=1000000: (1 - macro_auc_trained) of L162_MASS over L188, run by run :: sha256 c6c4138ee5c86423
\newcommand{\MlstCmpJcTwoTrainedResMassOverFullPr}{at least 3.7\% more errors}  % experiments/FIGS/data/v2/index.json :: linear probe jc2_188 N=1000000: (1 - macro_auc_trained) of L162_MASS over L188, run by run :: sha256 c6c4138ee5c86423 :: prose twin: two significant figures, direction in words
\newcommand{\MlstCmpJcTwoTrainedResMassOverFullPrMax}{up to 5.3\% more errors}  % experiments/FIGS/data/v2/index.json :: linear probe jc2_188 N=1000000: (1 - macro_auc_trained) of L162_MASS over L188, run by run :: sha256 c6c4138ee5c86423 :: prose twin of the run farther from 1
\newcommand{\MlstCmpJcTwoTrainedCoarseMassOverFull}{1.617 and 1.528}  % experiments/FIGS/data/v2/index.json :: linear probe jc2_188 N=1000000: (1 - macro_auc_trained) of R16_Q1_MASS over L188, run by run :: sha256 c6c4138ee5c86423
\newcommand{\MlstCmpJcTwoTrainedCoarseMassOverFullPr}{at least 53\% more errors}  % experiments/FIGS/data/v2/index.json :: linear probe jc2_188 N=1000000: (1 - macro_auc_trained) of R16_Q1_MASS over L188, run by run :: sha256 c6c4138ee5c86423 :: prose twin: two significant figures, direction in words
\newcommand{\MlstCmpJcTwoTrainedCoarseMassOverFullPrMax}{up to 62\% more errors}  % experiments/FIGS/data/v2/index.json :: linear probe jc2_188 N=1000000: (1 - macro_auc_trained) of R16_Q1_MASS over L188, run by run :: sha256 c6c4138ee5c86423 :: prose twin of the run farther from 1
\newcommand{\MlstCmpJcTwoTrainedFzeroOverFull}{1.424 and 1.462}  % experiments/FIGS/data/v2/index.json :: linear probe jc2_188 N=1000000: (1 - macro_auc_trained) of FLAV_F0 over L188, run by run :: sha256 c6c4138ee5c86423
\newcommand{\MlstCmpJcTwoTrainedFzeroOverFullPr}{at least 42\% more errors}  % experiments/FIGS/data/v2/index.json :: linear probe jc2_188 N=1000000: (1 - macro_auc_trained) of FLAV_F0 over L188, run by run :: sha256 c6c4138ee5c86423 :: prose twin: two significant figures, direction in words
\newcommand{\MlstCmpJcTwoTrainedFzeroOverFullPrMax}{up to 46\% more errors}  % experiments/FIGS/data/v2/index.json :: linear probe jc2_188 N=1000000: (1 - macro_auc_trained) of FLAV_F0 over L188, run by run :: sha256 c6c4138ee5c86423 :: prose twin of the run farther from 1
\newcommand{\MlstCmpJcTwoTrainedFoneOverFull}{1.218 and 1.212}  % experiments/FIGS/data/v2/index.json :: linear probe jc2_188 N=1000000: (1 - macro_auc_trained) of FLAV_F1 over L188, run by run :: sha256 c6c4138ee5c86423
\newcommand{\MlstCmpJcTwoTrainedFoneOverFullPr}{at least 21\% more errors}  % experiments/FIGS/data/v2/index.json :: linear probe jc2_188 N=1000000: (1 - macro_auc_trained) of FLAV_F1 over L188, run by run :: sha256 c6c4138ee5c86423 :: prose twin: two significant figures, direction in words
\newcommand{\MlstCmpJcTwoTrainedFoneOverFullPrMax}{up to 22\% more errors}  % experiments/FIGS/data/v2/index.json :: linear probe jc2_188 N=1000000: (1 - macro_auc_trained) of FLAV_F1 over L188, run by run :: sha256 c6c4138ee5c86423 :: prose twin of the run farther from 1
\newcommand{\MlstCmpJcTwoTrainedFonerOverFull}{1.206 and 1.208}  % experiments/FIGS/data/v2/index.json :: linear probe jc2_188 N=1000000: (1 - macro_auc_trained) of FLAV_F1R over L188, run by run :: sha256 c6c4138ee5c86423
\newcommand{\MlstCmpJcTwoTrainedFonerOverFullPr}{at least 21\% more errors}  % experiments/FIGS/data/v2/index.json :: linear probe jc2_188 N=1000000: (1 - macro_auc_trained) of FLAV_F1R over L188, run by run :: sha256 c6c4138ee5c86423 :: prose twin: two significant figures, direction in words
\newcommand{\MlstCmpJcTwoTrainedFonerOverFullPrMax}{up to 21\% more errors}  % experiments/FIGS/data/v2/index.json :: linear probe jc2_188 N=1000000: (1 - macro_auc_trained) of FLAV_F1R over L188, run by run :: sha256 c6c4138ee5c86423 :: prose twin of the run farther from 1
\newcommand{\MlstCmpJcTwoTrainedFzeroOverCoarse}{0.9275 and 0.9345}  % experiments/FIGS/data/v2/index.json :: linear probe jc2_188 N=1000000: (1 - macro_auc_trained) of FLAV_F0 over R16_Q1, run by run :: sha256 c6c4138ee5c86423
\newcommand{\MlstCmpJcTwoTrainedFzeroOverCoarsePr}{at least 6.5\% fewer errors}  % experiments/FIGS/data/v2/index.json :: linear probe jc2_188 N=1000000: (1 - macro_auc_trained) of FLAV_F0 over R16_Q1, run by run :: sha256 c6c4138ee5c86423 :: prose twin: two significant figures, direction in words
\newcommand{\MlstCmpJcTwoTrainedFzeroOverCoarsePrMax}{up to 7.2\% fewer errors}  % experiments/FIGS/data/v2/index.json :: linear probe jc2_188 N=1000000: (1 - macro_auc_trained) of FLAV_F0 over R16_Q1, run by run :: sha256 c6c4138ee5c86423 :: prose twin of the run farther from 1
\newcommand{\MlstCmpJcTwoTrainedFoneOverCoarse}{0.794 and 0.775}  % experiments/FIGS/data/v2/index.json :: linear probe jc2_188 N=1000000: (1 - macro_auc_trained) of FLAV_F1 over R16_Q1, run by run :: sha256 c6c4138ee5c86423
\newcommand{\MlstCmpJcTwoTrainedFoneOverCoarsePr}{at least 21\% fewer errors}  % experiments/FIGS/data/v2/index.json :: linear probe jc2_188 N=1000000: (1 - macro_auc_trained) of FLAV_F1 over R16_Q1, run by run :: sha256 c6c4138ee5c86423 :: prose twin: two significant figures, direction in words
\newcommand{\MlstCmpJcTwoTrainedFoneOverCoarsePrMax}{up to 23\% fewer errors}  % experiments/FIGS/data/v2/index.json :: linear probe jc2_188 N=1000000: (1 - macro_auc_trained) of FLAV_F1 over R16_Q1, run by run :: sha256 c6c4138ee5c86423 :: prose twin of the run farther from 1
\newcommand{\MlstCmpJcTwoTrainedFonerOverCoarse}{0.786 and 0.772}  % experiments/FIGS/data/v2/index.json :: linear probe jc2_188 N=1000000: (1 - macro_auc_trained) of FLAV_F1R over R16_Q1, run by run :: sha256 c6c4138ee5c86423
\newcommand{\MlstCmpJcTwoTrainedFonerOverCoarsePr}{at least 21\% fewer errors}  % experiments/FIGS/data/v2/index.json :: linear probe jc2_188 N=1000000: (1 - macro_auc_trained) of FLAV_F1R over R16_Q1, run by run :: sha256 c6c4138ee5c86423 :: prose twin: two significant figures, direction in words
\newcommand{\MlstCmpJcTwoTrainedFonerOverCoarsePrMax}{up to 23\% fewer errors}  % experiments/FIGS/data/v2/index.json :: linear probe jc2_188 N=1000000: (1 - macro_auc_trained) of FLAV_F1R over R16_Q1, run by run :: sha256 c6c4138ee5c86423 :: prose twin of the run farther from 1
\newcommand{\MlstCmpJcTwoTrainedFoneOverFzero}{0.856 and 0.829}  % experiments/FIGS/data/v2/index.json :: linear probe jc2_188 N=1000000: (1 - macro_auc_trained) of FLAV_F1 over FLAV_F0, run by run :: sha256 c6c4138ee5c86423
\newcommand{\MlstCmpJcTwoTrainedFoneOverFzeroPr}{at least 14\% fewer errors}  % experiments/FIGS/data/v2/index.json :: linear probe jc2_188 N=1000000: (1 - macro_auc_trained) of FLAV_F1 over FLAV_F0, run by run :: sha256 c6c4138ee5c86423 :: prose twin: two significant figures, direction in words
\newcommand{\MlstCmpJcTwoTrainedFoneOverFzeroPrMax}{up to 17\% fewer errors}  % experiments/FIGS/data/v2/index.json :: linear probe jc2_188 N=1000000: (1 - macro_auc_trained) of FLAV_F1 over FLAV_F0, run by run :: sha256 c6c4138ee5c86423 :: prose twin of the run farther from 1
\newcommand{\MlstCmpJcTwoTrainedFonerOverFzero}{0.847 and 0.826}  % experiments/FIGS/data/v2/index.json :: linear probe jc2_188 N=1000000: (1 - macro_auc_trained) of FLAV_F1R over FLAV_F0, run by run :: sha256 c6c4138ee5c86423
\newcommand{\MlstCmpJcTwoTrainedFonerOverFzeroPr}{at least 15\% fewer errors}  % experiments/FIGS/data/v2/index.json :: linear probe jc2_188 N=1000000: (1 - macro_auc_trained) of FLAV_F1R over FLAV_F0, run by run :: sha256 c6c4138ee5c86423 :: prose twin: two significant figures, direction in words
\newcommand{\MlstCmpJcTwoTrainedFonerOverFzeroPrMax}{up to 17\% fewer errors}  % experiments/FIGS/data/v2/index.json :: linear probe jc2_188 N=1000000: (1 - macro_auc_trained) of FLAV_F1R over FLAV_F0, run by run :: sha256 c6c4138ee5c86423 :: prose twin of the run farther from 1
\newcommand{\MlstCmpJcTwoTrainedFonerOverFone}{0.9899 and 0.9966}  % experiments/FIGS/data/v2/index.json :: linear probe jc2_188 N=1000000: (1 - macro_auc_trained) of FLAV_F1R over FLAV_F1, run by run :: sha256 c6c4138ee5c86423
\newcommand{\MlstCmpJcTwoTrainedFonerOverFonePr}{equal within 2\%}  % experiments/FIGS/data/v2/index.json :: linear probe jc2_188 N=1000000: (1 - macro_auc_trained) of FLAV_F1R over FLAV_F1, run by run :: sha256 c6c4138ee5c86423 :: prose twin: two significant figures, direction in words
\newcommand{\MlstCmpJcTwoTrainedFonerOverFonePrMax}{equal within 2\%}  % experiments/FIGS/data/v2/index.json :: linear probe jc2_188 N=1000000: (1 - macro_auc_trained) of FLAV_F1R over FLAV_F1, run by run :: sha256 c6c4138ee5c86423 :: prose twin of the run farther from 1
\newcommand{\MlstCmpJcTwoTrainedMidOverRes}{1.197\,\ensuremath{\pm}\,0.013}  % experiments/FIGS/data/v2/index.json :: linear probe jc2_188 N=1000000: (1 - macro_auc_trained) of R42_Q1 over L162, run by run :: sha256 c6c4138ee5c86423
\newcommand{\MlstCmpJcTwoTrainedMidOverResPr}{20\% more errors}  % experiments/FIGS/data/v2/index.json :: linear probe jc2_188 N=1000000: (1 - macro_auc_trained) of R42_Q1 over L162, run by run :: sha256 c6c4138ee5c86423 :: prose twin: two significant figures, direction in words
\newcommand{\MlstCmpJcTwoTrainedCoarseOverRes}{1.518\,\ensuremath{\pm}\,0.025}  % experiments/FIGS/data/v2/index.json :: linear probe jc2_188 N=1000000: (1 - macro_auc_trained) of R16_Q1 over L162, run by run :: sha256 c6c4138ee5c86423
\newcommand{\MlstCmpJcTwoTrainedCoarseOverResPr}{52\% more errors}  % experiments/FIGS/data/v2/index.json :: linear probe jc2_188 N=1000000: (1 - macro_auc_trained) of R16_Q1 over L162, run by run :: sha256 c6c4138ee5c86423 :: prose twin: two significant figures, direction in words
\newcommand{\MlstCmpJcTwoTrainedCoarseOverMid}{1.269\,\ensuremath{\pm}\,0.007}  % experiments/FIGS/data/v2/index.json :: linear probe jc2_188 N=1000000: (1 - macro_auc_trained) of R16_Q1 over R42_Q1, run by run :: sha256 c6c4138ee5c86423
\newcommand{\MlstCmpJcTwoTrainedCoarseOverMidPr}{27\% more errors}  % experiments/FIGS/data/v2/index.json :: linear probe jc2_188 N=1000000: (1 - macro_auc_trained) of R16_Q1 over R42_Q1, run by run :: sha256 c6c4138ee5c86423 :: prose twin: two significant figures, direction in words
\newcommand{\MlstCmpJcTwoTrainedRandOverFull}{1.125\,\ensuremath{\pm}\,0.008}  % experiments/FIGS/data/v2/index.json :: linear probe jc2_188 N=1000000: (1 - macro_auc_trained) of each random partition over the L188 run of its seed, mean +- SD over the partitions :: sha256 c6c4138ee5c86423
\newcommand{\MlstCmpJcTwoTrainedRandOverFullPr}{12\% more errors}  % experiments/FIGS/data/v2/index.json :: linear probe jc2_188 N=1000000: (1 - macro_auc_trained) of each random partition over the L188 run of its seed, mean +- SD over the partitions :: sha256 c6c4138ee5c86423 :: prose twin: two significant figures, direction in words
\newcommand{\MlstCmpJcTwoTrainedRandOverMid}{0.918\,\ensuremath{\pm}\,0.007}  % experiments/FIGS/data/v2/index.json :: linear probe jc2_188 N=1000000: (1 - macro_auc_trained) of each random partition over the R42_Q1 run of its seed, mean +- SD over the partitions :: sha256 c6c4138ee5c86423
\newcommand{\MlstCmpJcTwoTrainedRandOverMidPr}{8.2\% fewer errors}  % experiments/FIGS/data/v2/index.json :: linear probe jc2_188 N=1000000: (1 - macro_auc_trained) of each random partition over the R42_Q1 run of its seed, mean +- SD over the partitions :: sha256 c6c4138ee5c86423 :: prose twin: two significant figures, direction in words
\newcommand{\MlstCmpJcTwoTrainedRandOverCoarse}{0.724\,\ensuremath{\pm}\,0.008}  % experiments/FIGS/data/v2/index.json :: linear probe jc2_188 N=1000000: (1 - macro_auc_trained) of each random partition over the R16_Q1 run of its seed, mean +- SD over the partitions :: sha256 c6c4138ee5c86423
\newcommand{\MlstCmpJcTwoTrainedRandOverCoarsePr}{28\% fewer errors}  % experiments/FIGS/data/v2/index.json :: linear probe jc2_188 N=1000000: (1 - macro_auc_trained) of each random partition over the R16_Q1 run of its seed, mean +- SD over the partitions :: sha256 c6c4138ee5c86423 :: prose twin: two significant figures, direction in words
\newcommand{\MlstCmpJcTwoResResOverFull}{1.003\,\ensuremath{\pm}\,0.005}  % experiments/FIGS/data/v2/index.json :: linear probe jc2_188 N=1000000: (1 - macro_auc_resonant) of L162 over L188, run by run :: sha256 c6c4138ee5c86423
\newcommand{\MlstCmpJcTwoResResOverFullPr}{equal within 1\%}  % experiments/FIGS/data/v2/index.json :: linear probe jc2_188 N=1000000: (1 - macro_auc_resonant) of L162 over L188, run by run :: sha256 c6c4138ee5c86423 :: prose twin: two significant figures, direction in words
\newcommand{\MlstCmpJcTwoResMidOverFull}{1.262\,\ensuremath{\pm}\,0.012}  % experiments/FIGS/data/v2/index.json :: linear probe jc2_188 N=1000000: (1 - macro_auc_resonant) of R42_Q1 over L188, run by run :: sha256 c6c4138ee5c86423
\newcommand{\MlstCmpJcTwoResMidOverFullPr}{26\% more errors}  % experiments/FIGS/data/v2/index.json :: linear probe jc2_188 N=1000000: (1 - macro_auc_resonant) of R42_Q1 over L188, run by run :: sha256 c6c4138ee5c86423 :: prose twin: two significant figures, direction in words
\newcommand{\MlstCmpJcTwoResCoarseOverFull}{1.590\,\ensuremath{\pm}\,0.032}  % experiments/FIGS/data/v2/index.json :: linear probe jc2_188 N=1000000: (1 - macro_auc_resonant) of R16_Q1 over L188, run by run :: sha256 c6c4138ee5c86423
\newcommand{\MlstCmpJcTwoResCoarseOverFullPr}{59\% more errors}  % experiments/FIGS/data/v2/index.json :: linear probe jc2_188 N=1000000: (1 - macro_auc_resonant) of R16_Q1 over L188, run by run :: sha256 c6c4138ee5c86423 :: prose twin: two significant figures, direction in words
\newcommand{\MlstCmpJcTwoResResMassOverRes}{1.025 and 1.020}  % experiments/FIGS/data/v2/index.json :: linear probe jc2_188 N=1000000: (1 - macro_auc_resonant) of L162_MASS over L162, run by run :: sha256 c6c4138ee5c86423
\newcommand{\MlstCmpJcTwoResResMassOverResPr}{equal within 3\%}  % experiments/FIGS/data/v2/index.json :: linear probe jc2_188 N=1000000: (1 - macro_auc_resonant) of L162_MASS over L162, run by run :: sha256 c6c4138ee5c86423 :: prose twin: two significant figures, direction in words
\newcommand{\MlstCmpJcTwoResResMassOverResPrMax}{equal within 3\%}  % experiments/FIGS/data/v2/index.json :: linear probe jc2_188 N=1000000: (1 - macro_auc_resonant) of L162_MASS over L162, run by run :: sha256 c6c4138ee5c86423 :: prose twin of the run farther from 1
\newcommand{\MlstCmpJcTwoResCoarseMassOverCoarse}{1.08 and 0.96}  % experiments/FIGS/data/v2/index.json :: linear probe jc2_188 N=1000000: (1 - macro_auc_resonant) of R16_Q1_MASS over R16_Q1, run by run :: sha256 c6c4138ee5c86423
\newcommand{\MlstCmpJcTwoResCoarseMassOverCoarsePr}{equal within 8\%}  % experiments/FIGS/data/v2/index.json :: linear probe jc2_188 N=1000000: (1 - macro_auc_resonant) of R16_Q1_MASS over R16_Q1, run by run :: sha256 c6c4138ee5c86423 :: prose twin: two significant figures, direction in words
\newcommand{\MlstCmpJcTwoResCoarseMassOverCoarsePrMax}{equal within 8\%}  % experiments/FIGS/data/v2/index.json :: linear probe jc2_188 N=1000000: (1 - macro_auc_resonant) of R16_Q1_MASS over R16_Q1, run by run :: sha256 c6c4138ee5c86423 :: prose twin of the run farther from 1
\newcommand{\MlstCmpJcTwoResResMassOverFull}{1.033 and 1.018}  % experiments/FIGS/data/v2/index.json :: linear probe jc2_188 N=1000000: (1 - macro_auc_resonant) of L162_MASS over L188, run by run :: sha256 c6c4138ee5c86423
\newcommand{\MlstCmpJcTwoResResMassOverFullPr}{equal within 4\%}  % experiments/FIGS/data/v2/index.json :: linear probe jc2_188 N=1000000: (1 - macro_auc_resonant) of L162_MASS over L188, run by run :: sha256 c6c4138ee5c86423 :: prose twin: two significant figures, direction in words
\newcommand{\MlstCmpJcTwoResResMassOverFullPrMax}{equal within 4\%}  % experiments/FIGS/data/v2/index.json :: linear probe jc2_188 N=1000000: (1 - macro_auc_resonant) of L162_MASS over L188, run by run :: sha256 c6c4138ee5c86423 :: prose twin of the run farther from 1
\newcommand{\MlstCmpJcTwoResCoarseMassOverFull}{1.67 and 1.55}  % experiments/FIGS/data/v2/index.json :: linear probe jc2_188 N=1000000: (1 - macro_auc_resonant) of R16_Q1_MASS over L188, run by run :: sha256 c6c4138ee5c86423
\newcommand{\MlstCmpJcTwoResCoarseMassOverFullPr}{at least 55\% more errors}  % experiments/FIGS/data/v2/index.json :: linear probe jc2_188 N=1000000: (1 - macro_auc_resonant) of R16_Q1_MASS over L188, run by run :: sha256 c6c4138ee5c86423 :: prose twin: two significant figures, direction in words
\newcommand{\MlstCmpJcTwoResCoarseMassOverFullPrMax}{up to 67\% more errors}  % experiments/FIGS/data/v2/index.json :: linear probe jc2_188 N=1000000: (1 - macro_auc_resonant) of R16_Q1_MASS over L188, run by run :: sha256 c6c4138ee5c86423 :: prose twin of the run farther from 1
\newcommand{\MlstCmpJcTwoResFzeroOverFull}{1.408 and 1.449}  % experiments/FIGS/data/v2/index.json :: linear probe jc2_188 N=1000000: (1 - macro_auc_resonant) of FLAV_F0 over L188, run by run :: sha256 c6c4138ee5c86423
\newcommand{\MlstCmpJcTwoResFzeroOverFullPr}{at least 41\% more errors}  % experiments/FIGS/data/v2/index.json :: linear probe jc2_188 N=1000000: (1 - macro_auc_resonant) of FLAV_F0 over L188, run by run :: sha256 c6c4138ee5c86423 :: prose twin: two significant figures, direction in words
\newcommand{\MlstCmpJcTwoResFzeroOverFullPrMax}{up to 45\% more errors}  % experiments/FIGS/data/v2/index.json :: linear probe jc2_188 N=1000000: (1 - macro_auc_resonant) of FLAV_F0 over L188, run by run :: sha256 c6c4138ee5c86423 :: prose twin of the run farther from 1
\newcommand{\MlstCmpJcTwoResFoneOverFull}{1.210 and 1.204}  % experiments/FIGS/data/v2/index.json :: linear probe jc2_188 N=1000000: (1 - macro_auc_resonant) of FLAV_F1 over L188, run by run :: sha256 c6c4138ee5c86423
\newcommand{\MlstCmpJcTwoResFoneOverFullPr}{at least 20\% more errors}  % experiments/FIGS/data/v2/index.json :: linear probe jc2_188 N=1000000: (1 - macro_auc_resonant) of FLAV_F1 over L188, run by run :: sha256 c6c4138ee5c86423 :: prose twin: two significant figures, direction in words
\newcommand{\MlstCmpJcTwoResFoneOverFullPrMax}{up to 21\% more errors}  % experiments/FIGS/data/v2/index.json :: linear probe jc2_188 N=1000000: (1 - macro_auc_resonant) of FLAV_F1 over L188, run by run :: sha256 c6c4138ee5c86423 :: prose twin of the run farther from 1
\newcommand{\MlstCmpJcTwoResFonerOverFull}{1.187 and 1.188}  % experiments/FIGS/data/v2/index.json :: linear probe jc2_188 N=1000000: (1 - macro_auc_resonant) of FLAV_F1R over L188, run by run :: sha256 c6c4138ee5c86423
\newcommand{\MlstCmpJcTwoResFonerOverFullPr}{at least 19\% more errors}  % experiments/FIGS/data/v2/index.json :: linear probe jc2_188 N=1000000: (1 - macro_auc_resonant) of FLAV_F1R over L188, run by run :: sha256 c6c4138ee5c86423 :: prose twin: two significant figures, direction in words
\newcommand{\MlstCmpJcTwoResFonerOverFullPrMax}{up to 19\% more errors}  % experiments/FIGS/data/v2/index.json :: linear probe jc2_188 N=1000000: (1 - macro_auc_resonant) of FLAV_F1R over L188, run by run :: sha256 c6c4138ee5c86423 :: prose twin of the run farther from 1
\newcommand{\MlstCmpJcTwoResFzeroOverCoarse}{0.9069 and 0.9029}  % experiments/FIGS/data/v2/index.json :: linear probe jc2_188 N=1000000: (1 - macro_auc_resonant) of FLAV_F0 over R16_Q1, run by run :: sha256 c6c4138ee5c86423
\newcommand{\MlstCmpJcTwoResFzeroOverCoarsePr}{at least 9.3\% fewer errors}  % experiments/FIGS/data/v2/index.json :: linear probe jc2_188 N=1000000: (1 - macro_auc_resonant) of FLAV_F0 over R16_Q1, run by run :: sha256 c6c4138ee5c86423 :: prose twin: two significant figures, direction in words
\newcommand{\MlstCmpJcTwoResFzeroOverCoarsePrMax}{up to 9.7\% fewer errors}  % experiments/FIGS/data/v2/index.json :: linear probe jc2_188 N=1000000: (1 - macro_auc_resonant) of FLAV_F0 over R16_Q1, run by run :: sha256 c6c4138ee5c86423 :: prose twin of the run farther from 1
\newcommand{\MlstCmpJcTwoResFoneOverCoarse}{0.779 and 0.750}  % experiments/FIGS/data/v2/index.json :: linear probe jc2_188 N=1000000: (1 - macro_auc_resonant) of FLAV_F1 over R16_Q1, run by run :: sha256 c6c4138ee5c86423
\newcommand{\MlstCmpJcTwoResFoneOverCoarsePr}{at least 22\% fewer errors}  % experiments/FIGS/data/v2/index.json :: linear probe jc2_188 N=1000000: (1 - macro_auc_resonant) of FLAV_F1 over R16_Q1, run by run :: sha256 c6c4138ee5c86423 :: prose twin: two significant figures, direction in words
\newcommand{\MlstCmpJcTwoResFoneOverCoarsePrMax}{up to 25\% fewer errors}  % experiments/FIGS/data/v2/index.json :: linear probe jc2_188 N=1000000: (1 - macro_auc_resonant) of FLAV_F1 over R16_Q1, run by run :: sha256 c6c4138ee5c86423 :: prose twin of the run farther from 1
\newcommand{\MlstCmpJcTwoResFonerOverCoarse}{0.764 and 0.740}  % experiments/FIGS/data/v2/index.json :: linear probe jc2_188 N=1000000: (1 - macro_auc_resonant) of FLAV_F1R over R16_Q1, run by run :: sha256 c6c4138ee5c86423
\newcommand{\MlstCmpJcTwoResFonerOverCoarsePr}{at least 24\% fewer errors}  % experiments/FIGS/data/v2/index.json :: linear probe jc2_188 N=1000000: (1 - macro_auc_resonant) of FLAV_F1R over R16_Q1, run by run :: sha256 c6c4138ee5c86423 :: prose twin: two significant figures, direction in words
\newcommand{\MlstCmpJcTwoResFonerOverCoarsePrMax}{up to 26\% fewer errors}  % experiments/FIGS/data/v2/index.json :: linear probe jc2_188 N=1000000: (1 - macro_auc_resonant) of FLAV_F1R over R16_Q1, run by run :: sha256 c6c4138ee5c86423 :: prose twin of the run farther from 1
\newcommand{\MlstCmpJcTwoResFoneOverFzero}{0.859 and 0.831}  % experiments/FIGS/data/v2/index.json :: linear probe jc2_188 N=1000000: (1 - macro_auc_resonant) of FLAV_F1 over FLAV_F0, run by run :: sha256 c6c4138ee5c86423
\newcommand{\MlstCmpJcTwoResFoneOverFzeroPr}{at least 14\% fewer errors}  % experiments/FIGS/data/v2/index.json :: linear probe jc2_188 N=1000000: (1 - macro_auc_resonant) of FLAV_F1 over FLAV_F0, run by run :: sha256 c6c4138ee5c86423 :: prose twin: two significant figures, direction in words
\newcommand{\MlstCmpJcTwoResFoneOverFzeroPrMax}{up to 17\% fewer errors}  % experiments/FIGS/data/v2/index.json :: linear probe jc2_188 N=1000000: (1 - macro_auc_resonant) of FLAV_F1 over FLAV_F0, run by run :: sha256 c6c4138ee5c86423 :: prose twin of the run farther from 1
\newcommand{\MlstCmpJcTwoResFonerOverFzero}{0.843 and 0.820}  % experiments/FIGS/data/v2/index.json :: linear probe jc2_188 N=1000000: (1 - macro_auc_resonant) of FLAV_F1R over FLAV_F0, run by run :: sha256 c6c4138ee5c86423
\newcommand{\MlstCmpJcTwoResFonerOverFzeroPr}{at least 16\% fewer errors}  % experiments/FIGS/data/v2/index.json :: linear probe jc2_188 N=1000000: (1 - macro_auc_resonant) of FLAV_F1R over FLAV_F0, run by run :: sha256 c6c4138ee5c86423 :: prose twin: two significant figures, direction in words
\newcommand{\MlstCmpJcTwoResFonerOverFzeroPrMax}{up to 18\% fewer errors}  % experiments/FIGS/data/v2/index.json :: linear probe jc2_188 N=1000000: (1 - macro_auc_resonant) of FLAV_F1R over FLAV_F0, run by run :: sha256 c6c4138ee5c86423 :: prose twin of the run farther from 1
\newcommand{\MlstCmpJcTwoResFonerOverFone}{0.9806 and 0.9867}  % experiments/FIGS/data/v2/index.json :: linear probe jc2_188 N=1000000: (1 - macro_auc_resonant) of FLAV_F1R over FLAV_F1, run by run :: sha256 c6c4138ee5c86423
\newcommand{\MlstCmpJcTwoResFonerOverFonePr}{equal within 2\%}  % experiments/FIGS/data/v2/index.json :: linear probe jc2_188 N=1000000: (1 - macro_auc_resonant) of FLAV_F1R over FLAV_F1, run by run :: sha256 c6c4138ee5c86423 :: prose twin: two significant figures, direction in words
\newcommand{\MlstCmpJcTwoResFonerOverFonePrMax}{equal within 2\%}  % experiments/FIGS/data/v2/index.json :: linear probe jc2_188 N=1000000: (1 - macro_auc_resonant) of FLAV_F1R over FLAV_F1, run by run :: sha256 c6c4138ee5c86423 :: prose twin of the run farther from 1
\newcommand{\MlstCmpJcTwoResMidOverRes}{1.258\,\ensuremath{\pm}\,0.017}  % experiments/FIGS/data/v2/index.json :: linear probe jc2_188 N=1000000: (1 - macro_auc_resonant) of R42_Q1 over L162, run by run :: sha256 c6c4138ee5c86423
\newcommand{\MlstCmpJcTwoResMidOverResPr}{26\% more errors}  % experiments/FIGS/data/v2/index.json :: linear probe jc2_188 N=1000000: (1 - macro_auc_resonant) of R42_Q1 over L162, run by run :: sha256 c6c4138ee5c86423 :: prose twin: two significant figures, direction in words
\newcommand{\MlstCmpJcTwoResCoarseOverRes}{1.58\,\ensuremath{\pm}\,0.04}  % experiments/FIGS/data/v2/index.json :: linear probe jc2_188 N=1000000: (1 - macro_auc_resonant) of R16_Q1 over L162, run by run :: sha256 c6c4138ee5c86423
\newcommand{\MlstCmpJcTwoResCoarseOverResPr}{58\% more errors}  % experiments/FIGS/data/v2/index.json :: linear probe jc2_188 N=1000000: (1 - macro_auc_resonant) of R16_Q1 over L162, run by run :: sha256 c6c4138ee5c86423 :: prose twin: two significant figures, direction in words
\newcommand{\MlstCmpJcTwoResCoarseOverMid}{1.260\,\ensuremath{\pm}\,0.015}  % experiments/FIGS/data/v2/index.json :: linear probe jc2_188 N=1000000: (1 - macro_auc_resonant) of R16_Q1 over R42_Q1, run by run :: sha256 c6c4138ee5c86423
\newcommand{\MlstCmpJcTwoResCoarseOverMidPr}{26\% more errors}  % experiments/FIGS/data/v2/index.json :: linear probe jc2_188 N=1000000: (1 - macro_auc_resonant) of R16_Q1 over R42_Q1, run by run :: sha256 c6c4138ee5c86423 :: prose twin: two significant figures, direction in words
\newcommand{\MlstCmpJcTwoResRandOverFull}{1.098\,\ensuremath{\pm}\,0.008}  % experiments/FIGS/data/v2/index.json :: linear probe jc2_188 N=1000000: (1 - macro_auc_resonant) of each random partition over the L188 run of its seed, mean +- SD over the partitions :: sha256 c6c4138ee5c86423
\newcommand{\MlstCmpJcTwoResRandOverFullPr}{9.8\% more errors}  % experiments/FIGS/data/v2/index.json :: linear probe jc2_188 N=1000000: (1 - macro_auc_resonant) of each random partition over the L188 run of its seed, mean +- SD over the partitions :: sha256 c6c4138ee5c86423 :: prose twin: two significant figures, direction in words
\newcommand{\MlstCmpJcTwoResRandOverMid}{0.870\,\ensuremath{\pm}\,0.008}  % experiments/FIGS/data/v2/index.json :: linear probe jc2_188 N=1000000: (1 - macro_auc_resonant) of each random partition over the R42_Q1 run of its seed, mean +- SD over the partitions :: sha256 c6c4138ee5c86423
\newcommand{\MlstCmpJcTwoResRandOverMidPr}{13\% fewer errors}  % experiments/FIGS/data/v2/index.json :: linear probe jc2_188 N=1000000: (1 - macro_auc_resonant) of each random partition over the R42_Q1 run of its seed, mean +- SD over the partitions :: sha256 c6c4138ee5c86423 :: prose twin: two significant figures, direction in words
\newcommand{\MlstCmpJcTwoResRandOverCoarse}{0.693\,\ensuremath{\pm}\,0.011}  % experiments/FIGS/data/v2/index.json :: linear probe jc2_188 N=1000000: (1 - macro_auc_resonant) of each random partition over the R16_Q1 run of its seed, mean +- SD over the partitions :: sha256 c6c4138ee5c86423
\newcommand{\MlstCmpJcTwoResRandOverCoarsePr}{31\% fewer errors}  % experiments/FIGS/data/v2/index.json :: linear probe jc2_188 N=1000000: (1 - macro_auc_resonant) of each random partition over the R16_Q1 run of its seed, mean +- SD over the partitions :: sha256 c6c4138ee5c86423 :: prose twin: two significant figures, direction in words
\newcommand{\MlstCmpJcTwoQcdResOverFull}{1.087\,\ensuremath{\pm}\,0.006}  % experiments/FIGS/data/v2/index.json :: linear probe jc2_188 N=1000000: (1 - macro_auc_qcd) of L162 over L188, run by run :: sha256 c6c4138ee5c86423
\newcommand{\MlstCmpJcTwoQcdResOverFullPr}{8.7\% more errors}  % experiments/FIGS/data/v2/index.json :: linear probe jc2_188 N=1000000: (1 - macro_auc_qcd) of L162 over L188, run by run :: sha256 c6c4138ee5c86423 :: prose twin: two significant figures, direction in words
\newcommand{\MlstCmpJcTwoQcdMidOverFull}{1.131\,\ensuremath{\pm}\,0.007}  % experiments/FIGS/data/v2/index.json :: linear probe jc2_188 N=1000000: (1 - macro_auc_qcd) of R42_Q1 over L188, run by run :: sha256 c6c4138ee5c86423
\newcommand{\MlstCmpJcTwoQcdMidOverFullPr}{13\% more errors}  % experiments/FIGS/data/v2/index.json :: linear probe jc2_188 N=1000000: (1 - macro_auc_qcd) of R42_Q1 over L188, run by run :: sha256 c6c4138ee5c86423 :: prose twin: two significant figures, direction in words
\newcommand{\MlstCmpJcTwoQcdCoarseOverFull}{1.466\,\ensuremath{\pm}\,0.017}  % experiments/FIGS/data/v2/index.json :: linear probe jc2_188 N=1000000: (1 - macro_auc_qcd) of R16_Q1 over L188, run by run :: sha256 c6c4138ee5c86423
\newcommand{\MlstCmpJcTwoQcdCoarseOverFullPr}{47\% more errors}  % experiments/FIGS/data/v2/index.json :: linear probe jc2_188 N=1000000: (1 - macro_auc_qcd) of R16_Q1 over L188, run by run :: sha256 c6c4138ee5c86423 :: prose twin: two significant figures, direction in words
\newcommand{\MlstCmpJcTwoQcdResMassOverRes}{1.012 and 1.006}  % experiments/FIGS/data/v2/index.json :: linear probe jc2_188 N=1000000: (1 - macro_auc_qcd) of L162_MASS over L162, run by run :: sha256 c6c4138ee5c86423
\newcommand{\MlstCmpJcTwoQcdResMassOverResPr}{equal within 2\%}  % experiments/FIGS/data/v2/index.json :: linear probe jc2_188 N=1000000: (1 - macro_auc_qcd) of L162_MASS over L162, run by run :: sha256 c6c4138ee5c86423 :: prose twin: two significant figures, direction in words
\newcommand{\MlstCmpJcTwoQcdResMassOverResPrMax}{equal within 2\%}  % experiments/FIGS/data/v2/index.json :: linear probe jc2_188 N=1000000: (1 - macro_auc_qcd) of L162_MASS over L162, run by run :: sha256 c6c4138ee5c86423 :: prose twin of the run farther from 1
\newcommand{\MlstCmpJcTwoQcdCoarseMassOverCoarse}{0.991 and 1.014}  % experiments/FIGS/data/v2/index.json :: linear probe jc2_188 N=1000000: (1 - macro_auc_qcd) of R16_Q1_MASS over R16_Q1, run by run :: sha256 c6c4138ee5c86423
\newcommand{\MlstCmpJcTwoQcdCoarseMassOverCoarsePr}{equal within 2\%}  % experiments/FIGS/data/v2/index.json :: linear probe jc2_188 N=1000000: (1 - macro_auc_qcd) of R16_Q1_MASS over R16_Q1, run by run :: sha256 c6c4138ee5c86423 :: prose twin: two significant figures, direction in words
\newcommand{\MlstCmpJcTwoQcdCoarseMassOverCoarsePrMax}{equal within 2\%}  % experiments/FIGS/data/v2/index.json :: linear probe jc2_188 N=1000000: (1 - macro_auc_qcd) of R16_Q1_MASS over R16_Q1, run by run :: sha256 c6c4138ee5c86423 :: prose twin of the run farther from 1
\newcommand{\MlstCmpJcTwoQcdResMassOverFull}{1.107 and 1.092}  % experiments/FIGS/data/v2/index.json :: linear probe jc2_188 N=1000000: (1 - macro_auc_qcd) of L162_MASS over L188, run by run :: sha256 c6c4138ee5c86423
\newcommand{\MlstCmpJcTwoQcdResMassOverFullPr}{at least 9.2\% more errors}  % experiments/FIGS/data/v2/index.json :: linear probe jc2_188 N=1000000: (1 - macro_auc_qcd) of L162_MASS over L188, run by run :: sha256 c6c4138ee5c86423 :: prose twin: two significant figures, direction in words
\newcommand{\MlstCmpJcTwoQcdResMassOverFullPrMax}{up to 11\% more errors}  % experiments/FIGS/data/v2/index.json :: linear probe jc2_188 N=1000000: (1 - macro_auc_qcd) of L162_MASS over L188, run by run :: sha256 c6c4138ee5c86423 :: prose twin of the run farther from 1
\newcommand{\MlstCmpJcTwoQcdCoarseMassOverFull}{1.472 and 1.474}  % experiments/FIGS/data/v2/index.json :: linear probe jc2_188 N=1000000: (1 - macro_auc_qcd) of R16_Q1_MASS over L188, run by run :: sha256 c6c4138ee5c86423
\newcommand{\MlstCmpJcTwoQcdCoarseMassOverFullPr}{at least 47\% more errors}  % experiments/FIGS/data/v2/index.json :: linear probe jc2_188 N=1000000: (1 - macro_auc_qcd) of R16_Q1_MASS over L188, run by run :: sha256 c6c4138ee5c86423 :: prose twin: two significant figures, direction in words
\newcommand{\MlstCmpJcTwoQcdCoarseMassOverFullPrMax}{up to 47\% more errors}  % experiments/FIGS/data/v2/index.json :: linear probe jc2_188 N=1000000: (1 - macro_auc_qcd) of R16_Q1_MASS over L188, run by run :: sha256 c6c4138ee5c86423 :: prose twin of the run farther from 1
\newcommand{\MlstCmpJcTwoQcdFzeroOverFull}{1.466 and 1.497}  % experiments/FIGS/data/v2/index.json :: linear probe jc2_188 N=1000000: (1 - macro_auc_qcd) of FLAV_F0 over L188, run by run :: sha256 c6c4138ee5c86423
\newcommand{\MlstCmpJcTwoQcdFzeroOverFullPr}{at least 47\% more errors}  % experiments/FIGS/data/v2/index.json :: linear probe jc2_188 N=1000000: (1 - macro_auc_qcd) of FLAV_F0 over L188, run by run :: sha256 c6c4138ee5c86423 :: prose twin: two significant figures, direction in words
\newcommand{\MlstCmpJcTwoQcdFzeroOverFullPrMax}{up to 50\% more errors}  % experiments/FIGS/data/v2/index.json :: linear probe jc2_188 N=1000000: (1 - macro_auc_qcd) of FLAV_F0 over L188, run by run :: sha256 c6c4138ee5c86423 :: prose twin of the run farther from 1
\newcommand{\MlstCmpJcTwoQcdFoneOverFull}{1.241 and 1.235}  % experiments/FIGS/data/v2/index.json :: linear probe jc2_188 N=1000000: (1 - macro_auc_qcd) of FLAV_F1 over L188, run by run :: sha256 c6c4138ee5c86423
\newcommand{\MlstCmpJcTwoQcdFoneOverFullPr}{at least 23\% more errors}  % experiments/FIGS/data/v2/index.json :: linear probe jc2_188 N=1000000: (1 - macro_auc_qcd) of FLAV_F1 over L188, run by run :: sha256 c6c4138ee5c86423 :: prose twin: two significant figures, direction in words
\newcommand{\MlstCmpJcTwoQcdFoneOverFullPrMax}{up to 24\% more errors}  % experiments/FIGS/data/v2/index.json :: linear probe jc2_188 N=1000000: (1 - macro_auc_qcd) of FLAV_F1 over L188, run by run :: sha256 c6c4138ee5c86423 :: prose twin of the run farther from 1
\newcommand{\MlstCmpJcTwoQcdFonerOverFull}{1.260 and 1.263}  % experiments/FIGS/data/v2/index.json :: linear probe jc2_188 N=1000000: (1 - macro_auc_qcd) of FLAV_F1R over L188, run by run :: sha256 c6c4138ee5c86423
\newcommand{\MlstCmpJcTwoQcdFonerOverFullPr}{at least 26\% more errors}  % experiments/FIGS/data/v2/index.json :: linear probe jc2_188 N=1000000: (1 - macro_auc_qcd) of FLAV_F1R over L188, run by run :: sha256 c6c4138ee5c86423 :: prose twin: two significant figures, direction in words
\newcommand{\MlstCmpJcTwoQcdFonerOverFullPrMax}{up to 26\% more errors}  % experiments/FIGS/data/v2/index.json :: linear probe jc2_188 N=1000000: (1 - macro_auc_qcd) of FLAV_F1R over L188, run by run :: sha256 c6c4138ee5c86423 :: prose twin of the run farther from 1
\newcommand{\MlstCmpJcTwoQcdFzeroOverCoarse}{0.987 and 1.031}  % experiments/FIGS/data/v2/index.json :: linear probe jc2_188 N=1000000: (1 - macro_auc_qcd) of FLAV_F0 over R16_Q1, run by run :: sha256 c6c4138ee5c86423
\newcommand{\MlstCmpJcTwoQcdFzeroOverCoarsePr}{equal within 4\%}  % experiments/FIGS/data/v2/index.json :: linear probe jc2_188 N=1000000: (1 - macro_auc_qcd) of FLAV_F0 over R16_Q1, run by run :: sha256 c6c4138ee5c86423 :: prose twin: two significant figures, direction in words
\newcommand{\MlstCmpJcTwoQcdFzeroOverCoarsePrMax}{equal within 4\%}  % experiments/FIGS/data/v2/index.json :: linear probe jc2_188 N=1000000: (1 - macro_auc_qcd) of FLAV_F0 over R16_Q1, run by run :: sha256 c6c4138ee5c86423 :: prose twin of the run farther from 1
\newcommand{\MlstCmpJcTwoQcdFoneOverCoarse}{0.836 and 0.850}  % experiments/FIGS/data/v2/index.json :: linear probe jc2_188 N=1000000: (1 - macro_auc_qcd) of FLAV_F1 over R16_Q1, run by run :: sha256 c6c4138ee5c86423
\newcommand{\MlstCmpJcTwoQcdFoneOverCoarsePr}{at least 15\% fewer errors}  % experiments/FIGS/data/v2/index.json :: linear probe jc2_188 N=1000000: (1 - macro_auc_qcd) of FLAV_F1 over R16_Q1, run by run :: sha256 c6c4138ee5c86423 :: prose twin: two significant figures, direction in words
\newcommand{\MlstCmpJcTwoQcdFoneOverCoarsePrMax}{up to 16\% fewer errors}  % experiments/FIGS/data/v2/index.json :: linear probe jc2_188 N=1000000: (1 - macro_auc_qcd) of FLAV_F1 over R16_Q1, run by run :: sha256 c6c4138ee5c86423 :: prose twin of the run farther from 1
\newcommand{\MlstCmpJcTwoQcdFonerOverCoarse}{0.848 and 0.870}  % experiments/FIGS/data/v2/index.json :: linear probe jc2_188 N=1000000: (1 - macro_auc_qcd) of FLAV_F1R over R16_Q1, run by run :: sha256 c6c4138ee5c86423
\newcommand{\MlstCmpJcTwoQcdFonerOverCoarsePr}{at least 13\% fewer errors}  % experiments/FIGS/data/v2/index.json :: linear probe jc2_188 N=1000000: (1 - macro_auc_qcd) of FLAV_F1R over R16_Q1, run by run :: sha256 c6c4138ee5c86423 :: prose twin: two significant figures, direction in words
\newcommand{\MlstCmpJcTwoQcdFonerOverCoarsePrMax}{up to 15\% fewer errors}  % experiments/FIGS/data/v2/index.json :: linear probe jc2_188 N=1000000: (1 - macro_auc_qcd) of FLAV_F1R over R16_Q1, run by run :: sha256 c6c4138ee5c86423 :: prose twin of the run farther from 1
\newcommand{\MlstCmpJcTwoQcdFoneOverFzero}{0.847 and 0.825}  % experiments/FIGS/data/v2/index.json :: linear probe jc2_188 N=1000000: (1 - macro_auc_qcd) of FLAV_F1 over FLAV_F0, run by run :: sha256 c6c4138ee5c86423
\newcommand{\MlstCmpJcTwoQcdFoneOverFzeroPr}{at least 15\% fewer errors}  % experiments/FIGS/data/v2/index.json :: linear probe jc2_188 N=1000000: (1 - macro_auc_qcd) of FLAV_F1 over FLAV_F0, run by run :: sha256 c6c4138ee5c86423 :: prose twin: two significant figures, direction in words
\newcommand{\MlstCmpJcTwoQcdFoneOverFzeroPrMax}{up to 18\% fewer errors}  % experiments/FIGS/data/v2/index.json :: linear probe jc2_188 N=1000000: (1 - macro_auc_qcd) of FLAV_F1 over FLAV_F0, run by run :: sha256 c6c4138ee5c86423 :: prose twin of the run farther from 1
\newcommand{\MlstCmpJcTwoQcdFonerOverFzero}{0.859 and 0.844}  % experiments/FIGS/data/v2/index.json :: linear probe jc2_188 N=1000000: (1 - macro_auc_qcd) of FLAV_F1R over FLAV_F0, run by run :: sha256 c6c4138ee5c86423
\newcommand{\MlstCmpJcTwoQcdFonerOverFzeroPr}{at least 14\% fewer errors}  % experiments/FIGS/data/v2/index.json :: linear probe jc2_188 N=1000000: (1 - macro_auc_qcd) of FLAV_F1R over FLAV_F0, run by run :: sha256 c6c4138ee5c86423 :: prose twin: two significant figures, direction in words
\newcommand{\MlstCmpJcTwoQcdFonerOverFzeroPrMax}{up to 16\% fewer errors}  % experiments/FIGS/data/v2/index.json :: linear probe jc2_188 N=1000000: (1 - macro_auc_qcd) of FLAV_F1R over FLAV_F0, run by run :: sha256 c6c4138ee5c86423 :: prose twin of the run farther from 1
\newcommand{\MlstCmpJcTwoQcdFonerOverFone}{1.015 and 1.023}  % experiments/FIGS/data/v2/index.json :: linear probe jc2_188 N=1000000: (1 - macro_auc_qcd) of FLAV_F1R over FLAV_F1, run by run :: sha256 c6c4138ee5c86423
\newcommand{\MlstCmpJcTwoQcdFonerOverFonePr}{equal within 3\%}  % experiments/FIGS/data/v2/index.json :: linear probe jc2_188 N=1000000: (1 - macro_auc_qcd) of FLAV_F1R over FLAV_F1, run by run :: sha256 c6c4138ee5c86423 :: prose twin: two significant figures, direction in words
\newcommand{\MlstCmpJcTwoQcdFonerOverFonePrMax}{equal within 3\%}  % experiments/FIGS/data/v2/index.json :: linear probe jc2_188 N=1000000: (1 - macro_auc_qcd) of FLAV_F1R over FLAV_F1, run by run :: sha256 c6c4138ee5c86423 :: prose twin of the run farther from 1
\newcommand{\MlstCmpJcTwoQcdMidOverRes}{1.040\,\ensuremath{\pm}\,0.011}  % experiments/FIGS/data/v2/index.json :: linear probe jc2_188 N=1000000: (1 - macro_auc_qcd) of R42_Q1 over L162, run by run :: sha256 c6c4138ee5c86423
\newcommand{\MlstCmpJcTwoQcdMidOverResPr}{4.0\% more errors}  % experiments/FIGS/data/v2/index.json :: linear probe jc2_188 N=1000000: (1 - macro_auc_qcd) of R42_Q1 over L162, run by run :: sha256 c6c4138ee5c86423 :: prose twin: two significant figures, direction in words
\newcommand{\MlstCmpJcTwoQcdCoarseOverRes}{1.348\,\ensuremath{\pm}\,0.010}  % experiments/FIGS/data/v2/index.json :: linear probe jc2_188 N=1000000: (1 - macro_auc_qcd) of R16_Q1 over L162, run by run :: sha256 c6c4138ee5c86423
\newcommand{\MlstCmpJcTwoQcdCoarseOverResPr}{35\% more errors}  % experiments/FIGS/data/v2/index.json :: linear probe jc2_188 N=1000000: (1 - macro_auc_qcd) of R16_Q1 over L162, run by run :: sha256 c6c4138ee5c86423 :: prose twin: two significant figures, direction in words
\newcommand{\MlstCmpJcTwoQcdCoarseOverMid}{1.296\,\ensuremath{\pm}\,0.016}  % experiments/FIGS/data/v2/index.json :: linear probe jc2_188 N=1000000: (1 - macro_auc_qcd) of R16_Q1 over R42_Q1, run by run :: sha256 c6c4138ee5c86423
\newcommand{\MlstCmpJcTwoQcdCoarseOverMidPr}{30\% more errors}  % experiments/FIGS/data/v2/index.json :: linear probe jc2_188 N=1000000: (1 - macro_auc_qcd) of R16_Q1 over R42_Q1, run by run :: sha256 c6c4138ee5c86423 :: prose twin: two significant figures, direction in words
\newcommand{\MlstCmpJcTwoQcdRandOverFull}{1.198\,\ensuremath{\pm}\,0.011}  % experiments/FIGS/data/v2/index.json :: linear probe jc2_188 N=1000000: (1 - macro_auc_qcd) of each random partition over the L188 run of its seed, mean +- SD over the partitions :: sha256 c6c4138ee5c86423
\newcommand{\MlstCmpJcTwoQcdRandOverFullPr}{20\% more errors}  % experiments/FIGS/data/v2/index.json :: linear probe jc2_188 N=1000000: (1 - macro_auc_qcd) of each random partition over the L188 run of its seed, mean +- SD over the partitions :: sha256 c6c4138ee5c86423 :: prose twin: two significant figures, direction in words
\newcommand{\MlstCmpJcTwoQcdRandOverMid}{1.064\,\ensuremath{\pm}\,0.011}  % experiments/FIGS/data/v2/index.json :: linear probe jc2_188 N=1000000: (1 - macro_auc_qcd) of each random partition over the R42_Q1 run of its seed, mean +- SD over the partitions :: sha256 c6c4138ee5c86423
\newcommand{\MlstCmpJcTwoQcdRandOverMidPr}{6.4\% more errors}  % experiments/FIGS/data/v2/index.json :: linear probe jc2_188 N=1000000: (1 - macro_auc_qcd) of each random partition over the R42_Q1 run of its seed, mean +- SD over the partitions :: sha256 c6c4138ee5c86423 :: prose twin: two significant figures, direction in words
\newcommand{\MlstCmpJcTwoQcdRandOverCoarse}{0.818\,\ensuremath{\pm}\,0.013}  % experiments/FIGS/data/v2/index.json :: linear probe jc2_188 N=1000000: (1 - macro_auc_qcd) of each random partition over the R16_Q1 run of its seed, mean +- SD over the partitions :: sha256 c6c4138ee5c86423
\newcommand{\MlstCmpJcTwoQcdRandOverCoarsePr}{18\% fewer errors}  % experiments/FIGS/data/v2/index.json :: linear probe jc2_188 N=1000000: (1 - macro_auc_qcd) of each random partition over the R16_Q1 run of its seed, mean +- SD over the partitions :: sha256 c6c4138ee5c86423 :: prose twin: two significant figures, direction in words
\newcommand{\MlstCmpBcResOverFull}{0.998\,\ensuremath{\pm}\,0.007}  % experiments/FIGS/data/v2/index.json :: binary probe bc_vs_bq_cs: (1 - AUC) of L162 over L188, run by run :: sha256 c6c4138ee5c86423
\newcommand{\MlstCmpBcResOverFullPr}{equal within 1\%}  % experiments/FIGS/data/v2/index.json :: binary probe bc_vs_bq_cs: (1 - AUC) of L162 over L188, run by run :: sha256 c6c4138ee5c86423 :: prose twin: two significant figures, direction in words
\newcommand{\MlstCmpBcMidOverFull}{1.151\,\ensuremath{\pm}\,0.029}  % experiments/FIGS/data/v2/index.json :: binary probe bc_vs_bq_cs: (1 - AUC) of R42_Q1 over L188, run by run :: sha256 c6c4138ee5c86423
\newcommand{\MlstCmpBcMidOverFullPr}{15\% more errors}  % experiments/FIGS/data/v2/index.json :: binary probe bc_vs_bq_cs: (1 - AUC) of R42_Q1 over L188, run by run :: sha256 c6c4138ee5c86423 :: prose twin: two significant figures, direction in words
\newcommand{\MlstCmpBcCoarseOverFull}{1.95\,\ensuremath{\pm}\,0.06}  % experiments/FIGS/data/v2/index.json :: binary probe bc_vs_bq_cs: (1 - AUC) of R16_Q1 over L188, run by run :: sha256 c6c4138ee5c86423
\newcommand{\MlstCmpBcCoarseOverFullPr}{95\% more errors}  % experiments/FIGS/data/v2/index.json :: binary probe bc_vs_bq_cs: (1 - AUC) of R16_Q1 over L188, run by run :: sha256 c6c4138ee5c86423 :: prose twin: two significant figures, direction in words
\newcommand{\MlstCmpBcFoneOverFzero}{0.759 and 0.698}  % experiments/FIGS/data/v2/index.json :: binary probe bc_vs_bq_cs: (1 - AUC) of FLAV_F1 over FLAV_F0, run by run :: sha256 c6c4138ee5c86423
\newcommand{\MlstCmpBcFoneOverFzeroPr}{at least 24\% fewer errors}  % experiments/FIGS/data/v2/index.json :: binary probe bc_vs_bq_cs: (1 - AUC) of FLAV_F1 over FLAV_F0, run by run :: sha256 c6c4138ee5c86423 :: prose twin: two significant figures, direction in words
\newcommand{\MlstCmpBcFoneOverFzeroPrMax}{up to 30\% fewer errors}  % experiments/FIGS/data/v2/index.json :: binary probe bc_vs_bq_cs: (1 - AUC) of FLAV_F1 over FLAV_F0, run by run :: sha256 c6c4138ee5c86423 :: prose twin of the run farther from 1
\newcommand{\MlstCmpBcFonerOverFone}{1.008 and 0.990}  % experiments/FIGS/data/v2/index.json :: binary probe bc_vs_bq_cs: (1 - AUC) of FLAV_F1R over FLAV_F1, run by run :: sha256 c6c4138ee5c86423
\newcommand{\MlstCmpBcFonerOverFonePr}{equal within 1\%}  % experiments/FIGS/data/v2/index.json :: binary probe bc_vs_bq_cs: (1 - AUC) of FLAV_F1R over FLAV_F1, run by run :: sha256 c6c4138ee5c86423 :: prose twin: two significant figures, direction in words
\newcommand{\MlstCmpBcFonerOverFonePrMax}{equal within 1\%}  % experiments/FIGS/data/v2/index.json :: binary probe bc_vs_bq_cs: (1 - AUC) of FLAV_F1R over FLAV_F1, run by run :: sha256 c6c4138ee5c86423 :: prose twin of the run farther from 1
\newcommand{\MlstCmpBcFonerOverFzero}{0.765 and 0.691}  % experiments/FIGS/data/v2/index.json :: binary probe bc_vs_bq_cs: (1 - AUC) of FLAV_F1R over FLAV_F0, run by run :: sha256 c6c4138ee5c86423
\newcommand{\MlstCmpBcFonerOverFzeroPr}{at least 23\% fewer errors}  % experiments/FIGS/data/v2/index.json :: binary probe bc_vs_bq_cs: (1 - AUC) of FLAV_F1R over FLAV_F0, run by run :: sha256 c6c4138ee5c86423 :: prose twin: two significant figures, direction in words
\newcommand{\MlstCmpBcFonerOverFzeroPrMax}{up to 31\% fewer errors}  % experiments/FIGS/data/v2/index.json :: binary probe bc_vs_bq_cs: (1 - AUC) of FLAV_F1R over FLAV_F0, run by run :: sha256 c6c4138ee5c86423 :: prose twin of the run farther from 1
\newcommand{\MlstCmpBcFzeroOverCoarse}{0.99 and 1.13}  % experiments/FIGS/data/v2/index.json :: binary probe bc_vs_bq_cs: (1 - AUC) of FLAV_F0 over R16_Q1, run by run :: sha256 c6c4138ee5c86423
\newcommand{\MlstCmpBcFzeroOverCoarsePr}{equal within 20\%}  % experiments/FIGS/data/v2/index.json :: binary probe bc_vs_bq_cs: (1 - AUC) of FLAV_F0 over R16_Q1, run by run :: sha256 c6c4138ee5c86423 :: prose twin: two significant figures, direction in words
\newcommand{\MlstCmpBcFzeroOverCoarsePrMax}{equal within 20\%}  % experiments/FIGS/data/v2/index.json :: binary probe bc_vs_bq_cs: (1 - AUC) of FLAV_F0 over R16_Q1, run by run :: sha256 c6c4138ee5c86423 :: prose twin of the run farther from 1
\newcommand{\MlstCmpBcFoneOverCoarse}{0.752 and 0.786}  % experiments/FIGS/data/v2/index.json :: binary probe bc_vs_bq_cs: (1 - AUC) of FLAV_F1 over R16_Q1, run by run :: sha256 c6c4138ee5c86423
\newcommand{\MlstCmpBcFoneOverCoarsePr}{at least 21\% fewer errors}  % experiments/FIGS/data/v2/index.json :: binary probe bc_vs_bq_cs: (1 - AUC) of FLAV_F1 over R16_Q1, run by run :: sha256 c6c4138ee5c86423 :: prose twin: two significant figures, direction in words
\newcommand{\MlstCmpBcFoneOverCoarsePrMax}{up to 25\% fewer errors}  % experiments/FIGS/data/v2/index.json :: binary probe bc_vs_bq_cs: (1 - AUC) of FLAV_F1 over R16_Q1, run by run :: sha256 c6c4138ee5c86423 :: prose twin of the run farther from 1
\newcommand{\MlstCmpBcFonerOverCoarse}{0.758 and 0.778}  % experiments/FIGS/data/v2/index.json :: binary probe bc_vs_bq_cs: (1 - AUC) of FLAV_F1R over R16_Q1, run by run :: sha256 c6c4138ee5c86423
\newcommand{\MlstCmpBcFonerOverCoarsePr}{at least 22\% fewer errors}  % experiments/FIGS/data/v2/index.json :: binary probe bc_vs_bq_cs: (1 - AUC) of FLAV_F1R over R16_Q1, run by run :: sha256 c6c4138ee5c86423 :: prose twin: two significant figures, direction in words
\newcommand{\MlstCmpBcFonerOverCoarsePrMax}{up to 24\% fewer errors}  % experiments/FIGS/data/v2/index.json :: binary probe bc_vs_bq_cs: (1 - AUC) of FLAV_F1R over R16_Q1, run by run :: sha256 c6c4138ee5c86423 :: prose twin of the run farther from 1
\newcommand{\MlstCmpBcFzeroOverFull}{2.00 and 2.19}  % experiments/FIGS/data/v2/index.json :: binary probe bc_vs_bq_cs: (1 - AUC) of FLAV_F0 over L188, run by run :: sha256 c6c4138ee5c86423
\newcommand{\MlstCmpBcFzeroOverFullPr}{at least 2.0 times the errors}  % experiments/FIGS/data/v2/index.json :: binary probe bc_vs_bq_cs: (1 - AUC) of FLAV_F0 over L188, run by run :: sha256 c6c4138ee5c86423 :: prose twin: two significant figures, direction in words
\newcommand{\MlstCmpBcFzeroOverFullPrMax}{up to 2.2 times the errors}  % experiments/FIGS/data/v2/index.json :: binary probe bc_vs_bq_cs: (1 - AUC) of FLAV_F0 over L188, run by run :: sha256 c6c4138ee5c86423 :: prose twin of the run farther from 1
\newcommand{\MlstCmpBcFoneOverFull}{1.518 and 1.530}  % experiments/FIGS/data/v2/index.json :: binary probe bc_vs_bq_cs: (1 - AUC) of FLAV_F1 over L188, run by run :: sha256 c6c4138ee5c86423
\newcommand{\MlstCmpBcFoneOverFullPr}{at least 52\% more errors}  % experiments/FIGS/data/v2/index.json :: binary probe bc_vs_bq_cs: (1 - AUC) of FLAV_F1 over L188, run by run :: sha256 c6c4138ee5c86423 :: prose twin: two significant figures, direction in words
\newcommand{\MlstCmpBcFoneOverFullPrMax}{up to 53\% more errors}  % experiments/FIGS/data/v2/index.json :: binary probe bc_vs_bq_cs: (1 - AUC) of FLAV_F1 over L188, run by run :: sha256 c6c4138ee5c86423 :: prose twin of the run farther from 1
\newcommand{\MlstCmpBcFonerOverFull}{1.530 and 1.515}  % experiments/FIGS/data/v2/index.json :: binary probe bc_vs_bq_cs: (1 - AUC) of FLAV_F1R over L188, run by run :: sha256 c6c4138ee5c86423
\newcommand{\MlstCmpBcFonerOverFullPr}{at least 51\% more errors}  % experiments/FIGS/data/v2/index.json :: binary probe bc_vs_bq_cs: (1 - AUC) of FLAV_F1R over L188, run by run :: sha256 c6c4138ee5c86423 :: prose twin: two significant figures, direction in words
\newcommand{\MlstCmpBcFonerOverFullPrMax}{up to 53\% more errors}  % experiments/FIGS/data/v2/index.json :: binary probe bc_vs_bq_cs: (1 - AUC) of FLAV_F1R over L188, run by run :: sha256 c6c4138ee5c86423 :: prose twin of the run farther from 1
\newcommand{\MlstCmpBcRandOverFull}{1.15\,\ensuremath{\pm}\,0.07}  % experiments/FIGS/data/v2/index.json :: binary probe bc_vs_bq_cs: (1 - AUC) of each random partition over the 188-class run of its seed :: sha256 c6c4138ee5c86423
\newcommand{\MlstCmpBcRandOverFullPr}{15\% more errors}  % experiments/FIGS/data/v2/index.json :: binary probe bc_vs_bq_cs: (1 - AUC) of each random partition over the 188-class run of its seed :: sha256 c6c4138ee5c86423 :: prose twin: two significant figures, direction in words
\newcommand{\MlstPairRandMergedBcN}{2}  % experiments/FIGS/data/v2/index.json :: random partitions that merge bc_vs_bq_cs (count) :: sha256 c6c4138ee5c86423
\newcommand{\MlstPairRandApartBcN}{2}  % experiments/FIGS/data/v2/index.json :: random partitions that keep apart bc_vs_bq_cs (count) :: sha256 c6c4138ee5c86423
\newcommand{\MlstPairRandMergedOverApartBc}{1.12}  % experiments/FIGS/data/v2/index.json :: binary probe bc_vs_bq_cs: mean (1 - AUC) of the partitions that merge it over that of those that keep it apart (one run each) :: sha256 c6c4138ee5c86423
\newcommand{\MlstPairRandMergedOverApartBcPr}{12\% more errors}  % experiments/FIGS/data/v2/index.json :: binary probe bc_vs_bq_cs: mean (1 - AUC) of the partitions that merge it over that of those that keep it apart (one run each) :: sha256 c6c4138ee5c86423 :: prose twin: two significant figures, direction in words
\newcommand{\MlstCmpBbResOverFull}{1.015\,\ensuremath{\pm}\,0.027}  % experiments/FIGS/data/v2/index.json :: binary probe bb_vs_cc: (1 - AUC) of L162 over L188, run by run :: sha256 c6c4138ee5c86423
\newcommand{\MlstCmpBbResOverFullPr}{equal within 3\%}  % experiments/FIGS/data/v2/index.json :: binary probe bb_vs_cc: (1 - AUC) of L162 over L188, run by run :: sha256 c6c4138ee5c86423 :: prose twin: two significant figures, direction in words
\newcommand{\MlstCmpBbMidOverFull}{1.102\,\ensuremath{\pm}\,0.013}  % experiments/FIGS/data/v2/index.json :: binary probe bb_vs_cc: (1 - AUC) of R42_Q1 over L188, run by run :: sha256 c6c4138ee5c86423
\newcommand{\MlstCmpBbMidOverFullPr}{10\% more errors}  % experiments/FIGS/data/v2/index.json :: binary probe bb_vs_cc: (1 - AUC) of R42_Q1 over L188, run by run :: sha256 c6c4138ee5c86423 :: prose twin: two significant figures, direction in words
\newcommand{\MlstCmpBbCoarseOverFull}{4.22\,\ensuremath{\pm}\,0.29}  % experiments/FIGS/data/v2/index.json :: binary probe bb_vs_cc: (1 - AUC) of R16_Q1 over L188, run by run :: sha256 c6c4138ee5c86423
\newcommand{\MlstCmpBbCoarseOverFullPr}{4.2 times the errors}  % experiments/FIGS/data/v2/index.json :: binary probe bb_vs_cc: (1 - AUC) of R16_Q1 over L188, run by run :: sha256 c6c4138ee5c86423 :: prose twin: two significant figures, direction in words
\newcommand{\MlstCmpBbFoneOverFzero}{0.31 and 0.13}  % experiments/FIGS/data/v2/index.json :: binary probe bb_vs_cc: (1 - AUC) of FLAV_F1 over FLAV_F0, run by run :: sha256 c6c4138ee5c86423
\newcommand{\MlstCmpBbFoneOverFzeroPr}{at most 0.31 times the errors}  % experiments/FIGS/data/v2/index.json :: binary probe bb_vs_cc: (1 - AUC) of FLAV_F1 over FLAV_F0, run by run :: sha256 c6c4138ee5c86423 :: prose twin: two significant figures, direction in words
\newcommand{\MlstCmpBbFoneOverFzeroPrMax}{down to 0.13 times the errors}  % experiments/FIGS/data/v2/index.json :: binary probe bb_vs_cc: (1 - AUC) of FLAV_F1 over FLAV_F0, run by run :: sha256 c6c4138ee5c86423 :: prose twin of the run farther from 1
\newcommand{\MlstCmpBbFonerOverFone}{1.40 and 1.86}  % experiments/FIGS/data/v2/index.json :: binary probe bb_vs_cc: (1 - AUC) of FLAV_F1R over FLAV_F1, run by run :: sha256 c6c4138ee5c86423
\newcommand{\MlstCmpBbFonerOverFonePr}{at least 40\% more errors}  % experiments/FIGS/data/v2/index.json :: binary probe bb_vs_cc: (1 - AUC) of FLAV_F1R over FLAV_F1, run by run :: sha256 c6c4138ee5c86423 :: prose twin: two significant figures, direction in words
\newcommand{\MlstCmpBbFonerOverFonePrMax}{up to 86\% more errors}  % experiments/FIGS/data/v2/index.json :: binary probe bb_vs_cc: (1 - AUC) of FLAV_F1R over FLAV_F1, run by run :: sha256 c6c4138ee5c86423 :: prose twin of the run farther from 1
\newcommand{\MlstCmpBbFonerOverFzero}{0.44 and 0.24}  % experiments/FIGS/data/v2/index.json :: binary probe bb_vs_cc: (1 - AUC) of FLAV_F1R over FLAV_F0, run by run :: sha256 c6c4138ee5c86423
\newcommand{\MlstCmpBbFonerOverFzeroPr}{at most 0.44 times the errors}  % experiments/FIGS/data/v2/index.json :: binary probe bb_vs_cc: (1 - AUC) of FLAV_F1R over FLAV_F0, run by run :: sha256 c6c4138ee5c86423 :: prose twin: two significant figures, direction in words
\newcommand{\MlstCmpBbFonerOverFzeroPrMax}{down to 0.24 times the errors}  % experiments/FIGS/data/v2/index.json :: binary probe bb_vs_cc: (1 - AUC) of FLAV_F1R over FLAV_F0, run by run :: sha256 c6c4138ee5c86423 :: prose twin of the run farther from 1
\newcommand{\MlstCmpBbFzeroOverCoarse}{1.3 and 3.2}  % experiments/FIGS/data/v2/index.json :: binary probe bb_vs_cc: (1 - AUC) of FLAV_F0 over R16_Q1, run by run :: sha256 c6c4138ee5c86423
\newcommand{\MlstCmpBbFzeroOverCoarsePr}{at least 30\% more errors}  % experiments/FIGS/data/v2/index.json :: binary probe bb_vs_cc: (1 - AUC) of FLAV_F0 over R16_Q1, run by run :: sha256 c6c4138ee5c86423 :: prose twin: two significant figures, direction in words
\newcommand{\MlstCmpBbFzeroOverCoarsePrMax}{up to 3.2 times the errors}  % experiments/FIGS/data/v2/index.json :: binary probe bb_vs_cc: (1 - AUC) of FLAV_F0 over R16_Q1, run by run :: sha256 c6c4138ee5c86423 :: prose twin of the run farther from 1
\newcommand{\MlstCmpBbFoneOverCoarse}{0.4077 and 0.4126}  % experiments/FIGS/data/v2/index.json :: binary probe bb_vs_cc: (1 - AUC) of FLAV_F1 over R16_Q1, run by run :: sha256 c6c4138ee5c86423
\newcommand{\MlstCmpBbFoneOverCoarsePr}{at most 0.41 times the errors}  % experiments/FIGS/data/v2/index.json :: binary probe bb_vs_cc: (1 - AUC) of FLAV_F1 over R16_Q1, run by run :: sha256 c6c4138ee5c86423 :: prose twin: two significant figures, direction in words
\newcommand{\MlstCmpBbFoneOverCoarsePrMax}{down to 0.41 times the errors}  % experiments/FIGS/data/v2/index.json :: binary probe bb_vs_cc: (1 - AUC) of FLAV_F1 over R16_Q1, run by run :: sha256 c6c4138ee5c86423 :: prose twin of the run farther from 1
\newcommand{\MlstCmpBbFonerOverCoarse}{0.57 and 0.77}  % experiments/FIGS/data/v2/index.json :: binary probe bb_vs_cc: (1 - AUC) of FLAV_F1R over R16_Q1, run by run :: sha256 c6c4138ee5c86423
\newcommand{\MlstCmpBbFonerOverCoarsePr}{at least 23\% fewer errors}  % experiments/FIGS/data/v2/index.json :: binary probe bb_vs_cc: (1 - AUC) of FLAV_F1R over R16_Q1, run by run :: sha256 c6c4138ee5c86423 :: prose twin: two significant figures, direction in words
\newcommand{\MlstCmpBbFonerOverCoarsePrMax}{up to 43\% fewer errors}  % experiments/FIGS/data/v2/index.json :: binary probe bb_vs_cc: (1 - AUC) of FLAV_F1R over R16_Q1, run by run :: sha256 c6c4138ee5c86423 :: prose twin of the run farther from 1
\newcommand{\MlstCmpBbFzeroOverFull}{5.9 and 13.4}  % experiments/FIGS/data/v2/index.json :: binary probe bb_vs_cc: (1 - AUC) of FLAV_F0 over L188, run by run :: sha256 c6c4138ee5c86423
\newcommand{\MlstCmpBbFzeroOverFullPr}{at least 5.9 times the errors}  % experiments/FIGS/data/v2/index.json :: binary probe bb_vs_cc: (1 - AUC) of FLAV_F0 over L188, run by run :: sha256 c6c4138ee5c86423 :: prose twin: two significant figures, direction in words
\newcommand{\MlstCmpBbFzeroOverFullPrMax}{up to 13 times the errors}  % experiments/FIGS/data/v2/index.json :: binary probe bb_vs_cc: (1 - AUC) of FLAV_F0 over L188, run by run :: sha256 c6c4138ee5c86423 :: prose twin of the run farther from 1
\newcommand{\MlstCmpBbFoneOverFull}{1.837 and 1.748}  % experiments/FIGS/data/v2/index.json :: binary probe bb_vs_cc: (1 - AUC) of FLAV_F1 over L188, run by run :: sha256 c6c4138ee5c86423
\newcommand{\MlstCmpBbFoneOverFullPr}{at least 75\% more errors}  % experiments/FIGS/data/v2/index.json :: binary probe bb_vs_cc: (1 - AUC) of FLAV_F1 over L188, run by run :: sha256 c6c4138ee5c86423 :: prose twin: two significant figures, direction in words
\newcommand{\MlstCmpBbFoneOverFullPrMax}{up to 84\% more errors}  % experiments/FIGS/data/v2/index.json :: binary probe bb_vs_cc: (1 - AUC) of FLAV_F1 over L188, run by run :: sha256 c6c4138ee5c86423 :: prose twin of the run farther from 1
\newcommand{\MlstCmpBbFonerOverFull}{2.58 and 3.25}  % experiments/FIGS/data/v2/index.json :: binary probe bb_vs_cc: (1 - AUC) of FLAV_F1R over L188, run by run :: sha256 c6c4138ee5c86423
\newcommand{\MlstCmpBbFonerOverFullPr}{at least 2.6 times the errors}  % experiments/FIGS/data/v2/index.json :: binary probe bb_vs_cc: (1 - AUC) of FLAV_F1R over L188, run by run :: sha256 c6c4138ee5c86423 :: prose twin: two significant figures, direction in words
\newcommand{\MlstCmpBbFonerOverFullPrMax}{up to 3.2 times the errors}  % experiments/FIGS/data/v2/index.json :: binary probe bb_vs_cc: (1 - AUC) of FLAV_F1R over L188, run by run :: sha256 c6c4138ee5c86423 :: prose twin of the run farther from 1
\newcommand{\MlstCmpBbRandOverFull}{1.53\,\ensuremath{\pm}\,0.11}  % experiments/FIGS/data/v2/index.json :: binary probe bb_vs_cc: (1 - AUC) of each random partition over the 188-class run of its seed :: sha256 c6c4138ee5c86423
\newcommand{\MlstCmpBbRandOverFullPr}{53\% more errors}  % experiments/FIGS/data/v2/index.json :: binary probe bb_vs_cc: (1 - AUC) of each random partition over the 188-class run of its seed :: sha256 c6c4138ee5c86423 :: prose twin: two significant figures, direction in words
\newcommand{\MlstPairRandMergedBbN}{3}  % experiments/FIGS/data/v2/index.json :: random partitions that merge bb_vs_cc (count) :: sha256 c6c4138ee5c86423
\newcommand{\MlstPairRandApartBbN}{2}  % experiments/FIGS/data/v2/index.json :: random partitions that keep apart bb_vs_cc (count) :: sha256 c6c4138ee5c86423
\newcommand{\MlstPairRandMergedOverApartBb}{1.07}  % experiments/FIGS/data/v2/index.json :: binary probe bb_vs_cc: mean (1 - AUC) of the partitions that merge it over that of those that keep it apart (one run each) :: sha256 c6c4138ee5c86423
\newcommand{\MlstPairRandMergedOverApartBbPr}{7.3\% more errors}  % experiments/FIGS/data/v2/index.json :: binary probe bb_vs_cc: mean (1 - AUC) of the partitions that merge it over that of those that keep it apart (one run each) :: sha256 c6c4138ee5c86423 :: prose twin: two significant figures, direction in words
\newcommand{\MlstCmpBbqqResOverFull}{1.005\,\ensuremath{\pm}\,0.032}  % experiments/FIGS/data/v2/index.json :: binary probe bbqq_vs_ccqq: (1 - AUC) of L162 over L188, run by run :: sha256 c6c4138ee5c86423
\newcommand{\MlstCmpBbqqResOverFullPr}{equal within 4\%}  % experiments/FIGS/data/v2/index.json :: binary probe bbqq_vs_ccqq: (1 - AUC) of L162 over L188, run by run :: sha256 c6c4138ee5c86423 :: prose twin: two significant figures, direction in words
\newcommand{\MlstCmpBbqqMidOverFull}{1.117\,\ensuremath{\pm}\,0.008}  % experiments/FIGS/data/v2/index.json :: binary probe bbqq_vs_ccqq: (1 - AUC) of R42_Q1 over L188, run by run :: sha256 c6c4138ee5c86423
\newcommand{\MlstCmpBbqqMidOverFullPr}{12\% more errors}  % experiments/FIGS/data/v2/index.json :: binary probe bbqq_vs_ccqq: (1 - AUC) of R42_Q1 over L188, run by run :: sha256 c6c4138ee5c86423 :: prose twin: two significant figures, direction in words
\newcommand{\MlstCmpBbqqCoarseOverFull}{4.7\,\ensuremath{\pm}\,0.5}  % experiments/FIGS/data/v2/index.json :: binary probe bbqq_vs_ccqq: (1 - AUC) of R16_Q1 over L188, run by run :: sha256 c6c4138ee5c86423
\newcommand{\MlstCmpBbqqCoarseOverFullPr}{4.7 times the errors}  % experiments/FIGS/data/v2/index.json :: binary probe bbqq_vs_ccqq: (1 - AUC) of R16_Q1 over L188, run by run :: sha256 c6c4138ee5c86423 :: prose twin: two significant figures, direction in words
\newcommand{\MlstCmpBbqqFoneOverFzero}{0.236 and 0.154}  % experiments/FIGS/data/v2/index.json :: binary probe bbqq_vs_ccqq: (1 - AUC) of FLAV_F1 over FLAV_F0, run by run :: sha256 c6c4138ee5c86423
\newcommand{\MlstCmpBbqqFoneOverFzeroPr}{at most 0.24 times the errors}  % experiments/FIGS/data/v2/index.json :: binary probe bbqq_vs_ccqq: (1 - AUC) of FLAV_F1 over FLAV_F0, run by run :: sha256 c6c4138ee5c86423 :: prose twin: two significant figures, direction in words
\newcommand{\MlstCmpBbqqFoneOverFzeroPrMax}{down to 0.15 times the errors}  % experiments/FIGS/data/v2/index.json :: binary probe bbqq_vs_ccqq: (1 - AUC) of FLAV_F1 over FLAV_F0, run by run :: sha256 c6c4138ee5c86423 :: prose twin of the run farther from 1
\newcommand{\MlstCmpBbqqFonerOverFone}{1.505 and 1.577}  % experiments/FIGS/data/v2/index.json :: binary probe bbqq_vs_ccqq: (1 - AUC) of FLAV_F1R over FLAV_F1, run by run :: sha256 c6c4138ee5c86423
\newcommand{\MlstCmpBbqqFonerOverFonePr}{at least 50\% more errors}  % experiments/FIGS/data/v2/index.json :: binary probe bbqq_vs_ccqq: (1 - AUC) of FLAV_F1R over FLAV_F1, run by run :: sha256 c6c4138ee5c86423 :: prose twin: two significant figures, direction in words
\newcommand{\MlstCmpBbqqFonerOverFonePrMax}{up to 58\% more errors}  % experiments/FIGS/data/v2/index.json :: binary probe bbqq_vs_ccqq: (1 - AUC) of FLAV_F1R over FLAV_F1, run by run :: sha256 c6c4138ee5c86423 :: prose twin of the run farther from 1
\newcommand{\MlstCmpBbqqFonerOverFzero}{0.36 and 0.24}  % experiments/FIGS/data/v2/index.json :: binary probe bbqq_vs_ccqq: (1 - AUC) of FLAV_F1R over FLAV_F0, run by run :: sha256 c6c4138ee5c86423
\newcommand{\MlstCmpBbqqFonerOverFzeroPr}{at most 0.36 times the errors}  % experiments/FIGS/data/v2/index.json :: binary probe bbqq_vs_ccqq: (1 - AUC) of FLAV_F1R over FLAV_F0, run by run :: sha256 c6c4138ee5c86423 :: prose twin: two significant figures, direction in words
\newcommand{\MlstCmpBbqqFonerOverFzeroPrMax}{down to 0.24 times the errors}  % experiments/FIGS/data/v2/index.json :: binary probe bbqq_vs_ccqq: (1 - AUC) of FLAV_F1R over FLAV_F0, run by run :: sha256 c6c4138ee5c86423 :: prose twin of the run farther from 1
\newcommand{\MlstCmpBbqqFzeroOverCoarse}{1.51 and 1.79}  % experiments/FIGS/data/v2/index.json :: binary probe bbqq_vs_ccqq: (1 - AUC) of FLAV_F0 over R16_Q1, run by run :: sha256 c6c4138ee5c86423
\newcommand{\MlstCmpBbqqFzeroOverCoarsePr}{at least 51\% more errors}  % experiments/FIGS/data/v2/index.json :: binary probe bbqq_vs_ccqq: (1 - AUC) of FLAV_F0 over R16_Q1, run by run :: sha256 c6c4138ee5c86423 :: prose twin: two significant figures, direction in words
\newcommand{\MlstCmpBbqqFzeroOverCoarsePrMax}{up to 79\% more errors}  % experiments/FIGS/data/v2/index.json :: binary probe bbqq_vs_ccqq: (1 - AUC) of FLAV_F0 over R16_Q1, run by run :: sha256 c6c4138ee5c86423 :: prose twin of the run farther from 1
\newcommand{\MlstCmpBbqqFoneOverCoarse}{0.356 and 0.276}  % experiments/FIGS/data/v2/index.json :: binary probe bbqq_vs_ccqq: (1 - AUC) of FLAV_F1 over R16_Q1, run by run :: sha256 c6c4138ee5c86423
\newcommand{\MlstCmpBbqqFoneOverCoarsePr}{at most 0.36 times the errors}  % experiments/FIGS/data/v2/index.json :: binary probe bbqq_vs_ccqq: (1 - AUC) of FLAV_F1 over R16_Q1, run by run :: sha256 c6c4138ee5c86423 :: prose twin: two significant figures, direction in words
\newcommand{\MlstCmpBbqqFoneOverCoarsePrMax}{down to 0.28 times the errors}  % experiments/FIGS/data/v2/index.json :: binary probe bbqq_vs_ccqq: (1 - AUC) of FLAV_F1 over R16_Q1, run by run :: sha256 c6c4138ee5c86423 :: prose twin of the run farther from 1
\newcommand{\MlstCmpBbqqFonerOverCoarse}{0.54 and 0.44}  % experiments/FIGS/data/v2/index.json :: binary probe bbqq_vs_ccqq: (1 - AUC) of FLAV_F1R over R16_Q1, run by run :: sha256 c6c4138ee5c86423
\newcommand{\MlstCmpBbqqFonerOverCoarsePr}{at least 46\% fewer errors}  % experiments/FIGS/data/v2/index.json :: binary probe bbqq_vs_ccqq: (1 - AUC) of FLAV_F1R over R16_Q1, run by run :: sha256 c6c4138ee5c86423 :: prose twin: two significant figures, direction in words
\newcommand{\MlstCmpBbqqFonerOverCoarsePrMax}{down to 0.44 times the errors}  % experiments/FIGS/data/v2/index.json :: binary probe bbqq_vs_ccqq: (1 - AUC) of FLAV_F1R over R16_Q1, run by run :: sha256 c6c4138ee5c86423 :: prose twin of the run farther from 1
\newcommand{\MlstCmpBbqqFzeroOverFull}{6.4 and 9.4}  % experiments/FIGS/data/v2/index.json :: binary probe bbqq_vs_ccqq: (1 - AUC) of FLAV_F0 over L188, run by run :: sha256 c6c4138ee5c86423
\newcommand{\MlstCmpBbqqFzeroOverFullPr}{at least 6.4 times the errors}  % experiments/FIGS/data/v2/index.json :: binary probe bbqq_vs_ccqq: (1 - AUC) of FLAV_F0 over L188, run by run :: sha256 c6c4138ee5c86423 :: prose twin: two significant figures, direction in words
\newcommand{\MlstCmpBbqqFzeroOverFullPrMax}{up to 9.4 times the errors}  % experiments/FIGS/data/v2/index.json :: binary probe bbqq_vs_ccqq: (1 - AUC) of FLAV_F0 over L188, run by run :: sha256 c6c4138ee5c86423 :: prose twin of the run farther from 1
\newcommand{\MlstCmpBbqqFoneOverFull}{1.510 and 1.447}  % experiments/FIGS/data/v2/index.json :: binary probe bbqq_vs_ccqq: (1 - AUC) of FLAV_F1 over L188, run by run :: sha256 c6c4138ee5c86423
\newcommand{\MlstCmpBbqqFoneOverFullPr}{at least 45\% more errors}  % experiments/FIGS/data/v2/index.json :: binary probe bbqq_vs_ccqq: (1 - AUC) of FLAV_F1 over L188, run by run :: sha256 c6c4138ee5c86423 :: prose twin: two significant figures, direction in words
\newcommand{\MlstCmpBbqqFoneOverFullPrMax}{up to 51\% more errors}  % experiments/FIGS/data/v2/index.json :: binary probe bbqq_vs_ccqq: (1 - AUC) of FLAV_F1 over L188, run by run :: sha256 c6c4138ee5c86423 :: prose twin of the run farther from 1
\newcommand{\MlstCmpBbqqFonerOverFull}{2.272 and 2.281}  % experiments/FIGS/data/v2/index.json :: binary probe bbqq_vs_ccqq: (1 - AUC) of FLAV_F1R over L188, run by run :: sha256 c6c4138ee5c86423
\newcommand{\MlstCmpBbqqFonerOverFullPr}{at least 2.3 times the errors}  % experiments/FIGS/data/v2/index.json :: binary probe bbqq_vs_ccqq: (1 - AUC) of FLAV_F1R over L188, run by run :: sha256 c6c4138ee5c86423 :: prose twin: two significant figures, direction in words
\newcommand{\MlstCmpBbqqFonerOverFullPrMax}{up to 2.3 times the errors}  % experiments/FIGS/data/v2/index.json :: binary probe bbqq_vs_ccqq: (1 - AUC) of FLAV_F1R over L188, run by run :: sha256 c6c4138ee5c86423 :: prose twin of the run farther from 1
\newcommand{\MlstCmpBbqqRandOverFull}{1.42\,\ensuremath{\pm}\,0.06}  % experiments/FIGS/data/v2/index.json :: binary probe bbqq_vs_ccqq: (1 - AUC) of each random partition over the 188-class run of its seed :: sha256 c6c4138ee5c86423
\newcommand{\MlstCmpBbqqRandOverFullPr}{42\% more errors}  % experiments/FIGS/data/v2/index.json :: binary probe bbqq_vs_ccqq: (1 - AUC) of each random partition over the 188-class run of its seed :: sha256 c6c4138ee5c86423 :: prose twin: two significant figures, direction in words
\newcommand{\MlstPairRandMergedBbqqN}{2}  % experiments/FIGS/data/v2/index.json :: random partitions that merge bbqq_vs_ccqq (count) :: sha256 c6c4138ee5c86423
\newcommand{\MlstPairRandApartBbqqN}{3}  % experiments/FIGS/data/v2/index.json :: random partitions that keep apart bbqq_vs_ccqq (count) :: sha256 c6c4138ee5c86423
\newcommand{\MlstPairRandMergedOverApartBbqq}{0.996}  % experiments/FIGS/data/v2/index.json :: binary probe bbqq_vs_ccqq: mean (1 - AUC) of the partitions that merge it over that of those that keep it apart (one run each) :: sha256 c6c4138ee5c86423
\newcommand{\MlstPairRandMergedOverApartBbqqPr}{0.44\% fewer errors}  % experiments/FIGS/data/v2/index.json :: binary probe bbqq_vs_ccqq: mean (1 - AUC) of the partitions that merge it over that of those that keep it apart (one run each) :: sha256 c6c4138ee5c86423 :: prose twin: two significant figures, direction in words
\newcommand{\MlstCmpEemmResOverFull}{--}  % experiments/FIGS/data/v2/index.json :: binary probe ee_vs_mm: (1 - AUC) of L162 over L188, run by run :: sha256 c6c4138ee5c86423 :: the reference's 1 - AUC is 0
\newcommand{\MlstCmpEemmResOverFullPr}{undefined, since the reference separates the test jets completely}  % experiments/FIGS/data/v2/index.json :: binary probe ee_vs_mm: (1 - AUC) of L162 over L188, run by run :: sha256 c6c4138ee5c86423 :: the reference's 1 - AUC is 0
\newcommand{\MlstCmpEemmMidOverFull}{--}  % experiments/FIGS/data/v2/index.json :: binary probe ee_vs_mm: (1 - AUC) of R42_Q1 over L188, run by run :: sha256 c6c4138ee5c86423 :: the reference's 1 - AUC is 0
\newcommand{\MlstCmpEemmMidOverFullPr}{undefined, since the reference separates the test jets completely}  % experiments/FIGS/data/v2/index.json :: binary probe ee_vs_mm: (1 - AUC) of R42_Q1 over L188, run by run :: sha256 c6c4138ee5c86423 :: the reference's 1 - AUC is 0
\newcommand{\MlstCmpEemmCoarseOverFull}{--}  % experiments/FIGS/data/v2/index.json :: binary probe ee_vs_mm: (1 - AUC) of R16_Q1 over L188, run by run :: sha256 c6c4138ee5c86423 :: the reference's 1 - AUC is 0
\newcommand{\MlstCmpEemmCoarseOverFullPr}{undefined, since the reference separates the test jets completely}  % experiments/FIGS/data/v2/index.json :: binary probe ee_vs_mm: (1 - AUC) of R16_Q1 over L188, run by run :: sha256 c6c4138ee5c86423 :: the reference's 1 - AUC is 0
\newcommand{\MlstCmpEemmFoneOverFzero}{--}  % experiments/FIGS/data/v2/index.json :: binary probe ee_vs_mm: (1 - AUC) of FLAV_F1 over FLAV_F0, run by run :: sha256 c6c4138ee5c86423 :: the reference's 1 - AUC is 0
\newcommand{\MlstCmpEemmFoneOverFzeroPr}{undefined, since the reference separates the test jets completely}  % experiments/FIGS/data/v2/index.json :: binary probe ee_vs_mm: (1 - AUC) of FLAV_F1 over FLAV_F0, run by run :: sha256 c6c4138ee5c86423 :: the reference's 1 - AUC is 0
\newcommand{\MlstCmpEemmFoneOverFzeroPrMax}{undefined, since the reference separates the test jets completely}  % experiments/FIGS/data/v2/index.json :: binary probe ee_vs_mm: (1 - AUC) of FLAV_F1 over FLAV_F0, run by run :: sha256 c6c4138ee5c86423 :: the reference's 1 - AUC is 0
\newcommand{\MlstCmpEemmFonerOverFone}{--}  % experiments/FIGS/data/v2/index.json :: binary probe ee_vs_mm: (1 - AUC) of FLAV_F1R over FLAV_F1, run by run :: sha256 c6c4138ee5c86423 :: the reference's 1 - AUC is 0
\newcommand{\MlstCmpEemmFonerOverFonePr}{undefined, since the reference separates the test jets completely}  % experiments/FIGS/data/v2/index.json :: binary probe ee_vs_mm: (1 - AUC) of FLAV_F1R over FLAV_F1, run by run :: sha256 c6c4138ee5c86423 :: the reference's 1 - AUC is 0
\newcommand{\MlstCmpEemmFonerOverFonePrMax}{undefined, since the reference separates the test jets completely}  % experiments/FIGS/data/v2/index.json :: binary probe ee_vs_mm: (1 - AUC) of FLAV_F1R over FLAV_F1, run by run :: sha256 c6c4138ee5c86423 :: the reference's 1 - AUC is 0
\newcommand{\MlstCmpEemmFonerOverFzero}{--}  % experiments/FIGS/data/v2/index.json :: binary probe ee_vs_mm: (1 - AUC) of FLAV_F1R over FLAV_F0, run by run :: sha256 c6c4138ee5c86423 :: the reference's 1 - AUC is 0
\newcommand{\MlstCmpEemmFonerOverFzeroPr}{undefined, since the reference separates the test jets completely}  % experiments/FIGS/data/v2/index.json :: binary probe ee_vs_mm: (1 - AUC) of FLAV_F1R over FLAV_F0, run by run :: sha256 c6c4138ee5c86423 :: the reference's 1 - AUC is 0
\newcommand{\MlstCmpEemmFonerOverFzeroPrMax}{undefined, since the reference separates the test jets completely}  % experiments/FIGS/data/v2/index.json :: binary probe ee_vs_mm: (1 - AUC) of FLAV_F1R over FLAV_F0, run by run :: sha256 c6c4138ee5c86423 :: the reference's 1 - AUC is 0
\newcommand{\MlstCmpEemmFzeroOverCoarse}{0.0000022 and 0.0000003}  % experiments/FIGS/data/v2/index.json :: binary probe ee_vs_mm: (1 - AUC) of FLAV_F0 over R16_Q1, run by run :: sha256 c6c4138ee5c86423
\newcommand{\MlstCmpEemmFzeroOverCoarsePr}{at most 0.0000022 times the errors}  % experiments/FIGS/data/v2/index.json :: binary probe ee_vs_mm: (1 - AUC) of FLAV_F0 over R16_Q1, run by run :: sha256 c6c4138ee5c86423 :: prose twin: two significant figures, direction in words
\newcommand{\MlstCmpEemmFzeroOverCoarsePrMax}{down to 0.00000034 times the errors}  % experiments/FIGS/data/v2/index.json :: binary probe ee_vs_mm: (1 - AUC) of FLAV_F0 over R16_Q1, run by run :: sha256 c6c4138ee5c86423 :: prose twin of the run farther from 1
\newcommand{\MlstCmpEemmFoneOverCoarse}{0.000000 and 0.000038}  % experiments/FIGS/data/v2/index.json :: binary probe ee_vs_mm: (1 - AUC) of FLAV_F1 over R16_Q1, run by run :: sha256 c6c4138ee5c86423
\newcommand{\MlstCmpEemmFoneOverCoarsePr}{at most 0.000038 times the errors}  % experiments/FIGS/data/v2/index.json :: binary probe ee_vs_mm: (1 - AUC) of FLAV_F1 over R16_Q1, run by run :: sha256 c6c4138ee5c86423 :: prose twin: two significant figures, direction in words
\newcommand{\MlstCmpEemmFoneOverCoarsePrMax}{down to 0.000000042 times the errors}  % experiments/FIGS/data/v2/index.json :: binary probe ee_vs_mm: (1 - AUC) of FLAV_F1 over R16_Q1, run by run :: sha256 c6c4138ee5c86423 :: prose twin of the run farther from 1
\newcommand{\MlstCmpEemmFonerOverCoarse}{0.00000005 and 0.00000017}  % experiments/FIGS/data/v2/index.json :: binary probe ee_vs_mm: (1 - AUC) of FLAV_F1R over R16_Q1, run by run :: sha256 c6c4138ee5c86423
\newcommand{\MlstCmpEemmFonerOverCoarsePr}{at most 0.00000017 times the errors}  % experiments/FIGS/data/v2/index.json :: binary probe ee_vs_mm: (1 - AUC) of FLAV_F1R over R16_Q1, run by run :: sha256 c6c4138ee5c86423 :: prose twin: two significant figures, direction in words
\newcommand{\MlstCmpEemmFonerOverCoarsePrMax}{down to 0.000000055 times the errors}  % experiments/FIGS/data/v2/index.json :: binary probe ee_vs_mm: (1 - AUC) of FLAV_F1R over R16_Q1, run by run :: sha256 c6c4138ee5c86423 :: prose twin of the run farther from 1
\newcommand{\MlstCmpEemmFzeroOverFull}{--}  % experiments/FIGS/data/v2/index.json :: binary probe ee_vs_mm: (1 - AUC) of FLAV_F0 over L188, run by run :: sha256 c6c4138ee5c86423 :: the reference's 1 - AUC is 0
\newcommand{\MlstCmpEemmFzeroOverFullPr}{undefined, since the reference separates the test jets completely}  % experiments/FIGS/data/v2/index.json :: binary probe ee_vs_mm: (1 - AUC) of FLAV_F0 over L188, run by run :: sha256 c6c4138ee5c86423 :: the reference's 1 - AUC is 0
\newcommand{\MlstCmpEemmFzeroOverFullPrMax}{undefined, since the reference separates the test jets completely}  % experiments/FIGS/data/v2/index.json :: binary probe ee_vs_mm: (1 - AUC) of FLAV_F0 over L188, run by run :: sha256 c6c4138ee5c86423 :: the reference's 1 - AUC is 0
\newcommand{\MlstCmpEemmFoneOverFull}{--}  % experiments/FIGS/data/v2/index.json :: binary probe ee_vs_mm: (1 - AUC) of FLAV_F1 over L188, run by run :: sha256 c6c4138ee5c86423 :: the reference's 1 - AUC is 0
\newcommand{\MlstCmpEemmFoneOverFullPr}{undefined, since the reference separates the test jets completely}  % experiments/FIGS/data/v2/index.json :: binary probe ee_vs_mm: (1 - AUC) of FLAV_F1 over L188, run by run :: sha256 c6c4138ee5c86423 :: the reference's 1 - AUC is 0
\newcommand{\MlstCmpEemmFoneOverFullPrMax}{undefined, since the reference separates the test jets completely}  % experiments/FIGS/data/v2/index.json :: binary probe ee_vs_mm: (1 - AUC) of FLAV_F1 over L188, run by run :: sha256 c6c4138ee5c86423 :: the reference's 1 - AUC is 0
\newcommand{\MlstCmpEemmFonerOverFull}{--}  % experiments/FIGS/data/v2/index.json :: binary probe ee_vs_mm: (1 - AUC) of FLAV_F1R over L188, run by run :: sha256 c6c4138ee5c86423 :: the reference's 1 - AUC is 0
\newcommand{\MlstCmpEemmFonerOverFullPr}{undefined, since the reference separates the test jets completely}  % experiments/FIGS/data/v2/index.json :: binary probe ee_vs_mm: (1 - AUC) of FLAV_F1R over L188, run by run :: sha256 c6c4138ee5c86423 :: the reference's 1 - AUC is 0
\newcommand{\MlstCmpEemmFonerOverFullPrMax}{undefined, since the reference separates the test jets completely}  % experiments/FIGS/data/v2/index.json :: binary probe ee_vs_mm: (1 - AUC) of FLAV_F1R over L188, run by run :: sha256 c6c4138ee5c86423 :: the reference's 1 - AUC is 0
\newcommand{\MlstCmpEemmRandOverFull}{--}  % experiments/FIGS/data/v2/index.json :: binary probe ee_vs_mm: (1 - AUC) of each random partition over the 188-class run of its seed :: sha256 c6c4138ee5c86423 :: the reference's 1 - AUC is 0
\newcommand{\MlstCmpEemmRandOverFullPr}{undefined, since the reference separates the test jets completely}  % experiments/FIGS/data/v2/index.json :: binary probe ee_vs_mm: (1 - AUC) of each random partition over the 188-class run of its seed :: sha256 c6c4138ee5c86423 :: the reference's 1 - AUC is 0
\newcommand{\MlstPairRandMergedEemmN}{3}  % experiments/FIGS/data/v2/index.json :: random partitions that merge ee_vs_mm (count) :: sha256 c6c4138ee5c86423
\newcommand{\MlstPairRandApartEemmN}{2}  % experiments/FIGS/data/v2/index.json :: random partitions that keep apart ee_vs_mm (count) :: sha256 c6c4138ee5c86423
\newcommand{\MlstPairRandMergedOverApartEemm}{75.3}  % experiments/FIGS/data/v2/index.json :: binary probe ee_vs_mm: mean (1 - AUC) of the partitions that merge it over that of those that keep it apart (one run each) :: sha256 c6c4138ee5c86423
\newcommand{\MlstPairRandMergedOverApartEemmPr}{75 times the errors}  % experiments/FIGS/data/v2/index.json :: binary probe ee_vs_mm: mean (1 - AUC) of the partitions that merge it over that of those that keep it apart (one run each) :: sha256 c6c4138ee5c86423 :: prose twin: two significant figures, direction in words
\newcommand{\MlstCmpQcdBbResOverFull}{1.15\,\ensuremath{\pm}\,0.04}  % experiments/FIGS/data/v2/index.json :: binary probe qcd_bb_vs_cc: (1 - AUC) of L162 over L188, run by run :: sha256 c6c4138ee5c86423
\newcommand{\MlstCmpQcdBbResOverFullPr}{15\% more errors}  % experiments/FIGS/data/v2/index.json :: binary probe qcd_bb_vs_cc: (1 - AUC) of L162 over L188, run by run :: sha256 c6c4138ee5c86423 :: prose twin: two significant figures, direction in words
\newcommand{\MlstCmpQcdBbMidOverFull}{1.163\,\ensuremath{\pm}\,0.016}  % experiments/FIGS/data/v2/index.json :: binary probe qcd_bb_vs_cc: (1 - AUC) of R42_Q1 over L188, run by run :: sha256 c6c4138ee5c86423
\newcommand{\MlstCmpQcdBbMidOverFullPr}{16\% more errors}  % experiments/FIGS/data/v2/index.json :: binary probe qcd_bb_vs_cc: (1 - AUC) of R42_Q1 over L188, run by run :: sha256 c6c4138ee5c86423 :: prose twin: two significant figures, direction in words
\newcommand{\MlstCmpQcdBbCoarseOverFull}{3.294\,\ensuremath{\pm}\,0.033}  % experiments/FIGS/data/v2/index.json :: binary probe qcd_bb_vs_cc: (1 - AUC) of R16_Q1 over L188, run by run :: sha256 c6c4138ee5c86423
\newcommand{\MlstCmpQcdBbCoarseOverFullPr}{3.3 times the errors}  % experiments/FIGS/data/v2/index.json :: binary probe qcd_bb_vs_cc: (1 - AUC) of R16_Q1 over L188, run by run :: sha256 c6c4138ee5c86423 :: prose twin: two significant figures, direction in words
\newcommand{\MlstCmpQcdBbFoneOverFzero}{0.368 and 0.294}  % experiments/FIGS/data/v2/index.json :: binary probe qcd_bb_vs_cc: (1 - AUC) of FLAV_F1 over FLAV_F0, run by run :: sha256 c6c4138ee5c86423
\newcommand{\MlstCmpQcdBbFoneOverFzeroPr}{at most 0.37 times the errors}  % experiments/FIGS/data/v2/index.json :: binary probe qcd_bb_vs_cc: (1 - AUC) of FLAV_F1 over FLAV_F0, run by run :: sha256 c6c4138ee5c86423 :: prose twin: two significant figures, direction in words
\newcommand{\MlstCmpQcdBbFoneOverFzeroPrMax}{down to 0.29 times the errors}  % experiments/FIGS/data/v2/index.json :: binary probe qcd_bb_vs_cc: (1 - AUC) of FLAV_F1 over FLAV_F0, run by run :: sha256 c6c4138ee5c86423 :: prose twin of the run farther from 1
\newcommand{\MlstCmpQcdBbFonerOverFone}{1.361 and 1.429}  % experiments/FIGS/data/v2/index.json :: binary probe qcd_bb_vs_cc: (1 - AUC) of FLAV_F1R over FLAV_F1, run by run :: sha256 c6c4138ee5c86423
\newcommand{\MlstCmpQcdBbFonerOverFonePr}{at least 36\% more errors}  % experiments/FIGS/data/v2/index.json :: binary probe qcd_bb_vs_cc: (1 - AUC) of FLAV_F1R over FLAV_F1, run by run :: sha256 c6c4138ee5c86423 :: prose twin: two significant figures, direction in words
\newcommand{\MlstCmpQcdBbFonerOverFonePrMax}{up to 43\% more errors}  % experiments/FIGS/data/v2/index.json :: binary probe qcd_bb_vs_cc: (1 - AUC) of FLAV_F1R over FLAV_F1, run by run :: sha256 c6c4138ee5c86423 :: prose twin of the run farther from 1
\newcommand{\MlstCmpQcdBbFonerOverFzero}{0.501 and 0.420}  % experiments/FIGS/data/v2/index.json :: binary probe qcd_bb_vs_cc: (1 - AUC) of FLAV_F1R over FLAV_F0, run by run :: sha256 c6c4138ee5c86423
\newcommand{\MlstCmpQcdBbFonerOverFzeroPr}{at least 50\% fewer errors}  % experiments/FIGS/data/v2/index.json :: binary probe qcd_bb_vs_cc: (1 - AUC) of FLAV_F1R over FLAV_F0, run by run :: sha256 c6c4138ee5c86423 :: prose twin: two significant figures, direction in words
\newcommand{\MlstCmpQcdBbFonerOverFzeroPrMax}{down to 0.42 times the errors}  % experiments/FIGS/data/v2/index.json :: binary probe qcd_bb_vs_cc: (1 - AUC) of FLAV_F1R over FLAV_F0, run by run :: sha256 c6c4138ee5c86423 :: prose twin of the run farther from 1
\newcommand{\MlstCmpQcdBbFzeroOverCoarse}{1.28 and 1.58}  % experiments/FIGS/data/v2/index.json :: binary probe qcd_bb_vs_cc: (1 - AUC) of FLAV_F0 over R16_Q1, run by run :: sha256 c6c4138ee5c86423
\newcommand{\MlstCmpQcdBbFzeroOverCoarsePr}{at least 28\% more errors}  % experiments/FIGS/data/v2/index.json :: binary probe qcd_bb_vs_cc: (1 - AUC) of FLAV_F0 over R16_Q1, run by run :: sha256 c6c4138ee5c86423 :: prose twin: two significant figures, direction in words
\newcommand{\MlstCmpQcdBbFzeroOverCoarsePrMax}{up to 58\% more errors}  % experiments/FIGS/data/v2/index.json :: binary probe qcd_bb_vs_cc: (1 - AUC) of FLAV_F0 over R16_Q1, run by run :: sha256 c6c4138ee5c86423 :: prose twin of the run farther from 1
\newcommand{\MlstCmpQcdBbFoneOverCoarse}{0.4692 and 0.4648}  % experiments/FIGS/data/v2/index.json :: binary probe qcd_bb_vs_cc: (1 - AUC) of FLAV_F1 over R16_Q1, run by run :: sha256 c6c4138ee5c86423
\newcommand{\MlstCmpQcdBbFoneOverCoarsePr}{at most 0.47 times the errors}  % experiments/FIGS/data/v2/index.json :: binary probe qcd_bb_vs_cc: (1 - AUC) of FLAV_F1 over R16_Q1, run by run :: sha256 c6c4138ee5c86423 :: prose twin: two significant figures, direction in words
\newcommand{\MlstCmpQcdBbFoneOverCoarsePrMax}{down to 0.46 times the errors}  % experiments/FIGS/data/v2/index.json :: binary probe qcd_bb_vs_cc: (1 - AUC) of FLAV_F1 over R16_Q1, run by run :: sha256 c6c4138ee5c86423 :: prose twin of the run farther from 1
\newcommand{\MlstCmpQcdBbFonerOverCoarse}{0.639 and 0.664}  % experiments/FIGS/data/v2/index.json :: binary probe qcd_bb_vs_cc: (1 - AUC) of FLAV_F1R over R16_Q1, run by run :: sha256 c6c4138ee5c86423
\newcommand{\MlstCmpQcdBbFonerOverCoarsePr}{at least 34\% fewer errors}  % experiments/FIGS/data/v2/index.json :: binary probe qcd_bb_vs_cc: (1 - AUC) of FLAV_F1R over R16_Q1, run by run :: sha256 c6c4138ee5c86423 :: prose twin: two significant figures, direction in words
\newcommand{\MlstCmpQcdBbFonerOverCoarsePrMax}{up to 36\% fewer errors}  % experiments/FIGS/data/v2/index.json :: binary probe qcd_bb_vs_cc: (1 - AUC) of FLAV_F1R over R16_Q1, run by run :: sha256 c6c4138ee5c86423 :: prose twin of the run farther from 1
\newcommand{\MlstCmpQcdBbFzeroOverFull}{4.24 and 5.21}  % experiments/FIGS/data/v2/index.json :: binary probe qcd_bb_vs_cc: (1 - AUC) of FLAV_F0 over L188, run by run :: sha256 c6c4138ee5c86423
\newcommand{\MlstCmpQcdBbFzeroOverFullPr}{at least 4.2 times the errors}  % experiments/FIGS/data/v2/index.json :: binary probe qcd_bb_vs_cc: (1 - AUC) of FLAV_F0 over L188, run by run :: sha256 c6c4138ee5c86423 :: prose twin: two significant figures, direction in words
\newcommand{\MlstCmpQcdBbFzeroOverFullPrMax}{up to 5.2 times the errors}  % experiments/FIGS/data/v2/index.json :: binary probe qcd_bb_vs_cc: (1 - AUC) of FLAV_F0 over L188, run by run :: sha256 c6c4138ee5c86423 :: prose twin of the run farther from 1
\newcommand{\MlstCmpQcdBbFoneOverFull}{1.561 and 1.532}  % experiments/FIGS/data/v2/index.json :: binary probe qcd_bb_vs_cc: (1 - AUC) of FLAV_F1 over L188, run by run :: sha256 c6c4138ee5c86423
\newcommand{\MlstCmpQcdBbFoneOverFullPr}{at least 53\% more errors}  % experiments/FIGS/data/v2/index.json :: binary probe qcd_bb_vs_cc: (1 - AUC) of FLAV_F1 over L188, run by run :: sha256 c6c4138ee5c86423 :: prose twin: two significant figures, direction in words
\newcommand{\MlstCmpQcdBbFoneOverFullPrMax}{up to 56\% more errors}  % experiments/FIGS/data/v2/index.json :: binary probe qcd_bb_vs_cc: (1 - AUC) of FLAV_F1 over L188, run by run :: sha256 c6c4138ee5c86423 :: prose twin of the run farther from 1
\newcommand{\MlstCmpQcdBbFonerOverFull}{2.124 and 2.189}  % experiments/FIGS/data/v2/index.json :: binary probe qcd_bb_vs_cc: (1 - AUC) of FLAV_F1R over L188, run by run :: sha256 c6c4138ee5c86423
\newcommand{\MlstCmpQcdBbFonerOverFullPr}{at least 2.1 times the errors}  % experiments/FIGS/data/v2/index.json :: binary probe qcd_bb_vs_cc: (1 - AUC) of FLAV_F1R over L188, run by run :: sha256 c6c4138ee5c86423 :: prose twin: two significant figures, direction in words
\newcommand{\MlstCmpQcdBbFonerOverFullPrMax}{up to 2.2 times the errors}  % experiments/FIGS/data/v2/index.json :: binary probe qcd_bb_vs_cc: (1 - AUC) of FLAV_F1R over L188, run by run :: sha256 c6c4138ee5c86423 :: prose twin of the run farther from 1
\newcommand{\MlstCmpQcdBResOverFull}{1.085\,\ensuremath{\pm}\,0.004}  % experiments/FIGS/data/v2/index.json :: binary probe qcd_b_vs_light: (1 - AUC) of L162 over L188, run by run :: sha256 c6c4138ee5c86423
\newcommand{\MlstCmpQcdBResOverFullPr}{8.5\% more errors}  % experiments/FIGS/data/v2/index.json :: binary probe qcd_b_vs_light: (1 - AUC) of L162 over L188, run by run :: sha256 c6c4138ee5c86423 :: prose twin: two significant figures, direction in words
\newcommand{\MlstCmpQcdBMidOverFull}{1.081\,\ensuremath{\pm}\,0.008}  % experiments/FIGS/data/v2/index.json :: binary probe qcd_b_vs_light: (1 - AUC) of R42_Q1 over L188, run by run :: sha256 c6c4138ee5c86423
\newcommand{\MlstCmpQcdBMidOverFullPr}{8.1\% more errors}  % experiments/FIGS/data/v2/index.json :: binary probe qcd_b_vs_light: (1 - AUC) of R42_Q1 over L188, run by run :: sha256 c6c4138ee5c86423 :: prose twin: two significant figures, direction in words
\newcommand{\MlstCmpQcdBCoarseOverFull}{1.54\,\ensuremath{\pm}\,0.07}  % experiments/FIGS/data/v2/index.json :: binary probe qcd_b_vs_light: (1 - AUC) of R16_Q1 over L188, run by run :: sha256 c6c4138ee5c86423
\newcommand{\MlstCmpQcdBCoarseOverFullPr}{54\% more errors}  % experiments/FIGS/data/v2/index.json :: binary probe qcd_b_vs_light: (1 - AUC) of R16_Q1 over L188, run by run :: sha256 c6c4138ee5c86423 :: prose twin: two significant figures, direction in words
\newcommand{\MlstCmpQcdBFoneOverFzero}{0.778 and 0.835}  % experiments/FIGS/data/v2/index.json :: binary probe qcd_b_vs_light: (1 - AUC) of FLAV_F1 over FLAV_F0, run by run :: sha256 c6c4138ee5c86423
\newcommand{\MlstCmpQcdBFoneOverFzeroPr}{at least 17\% fewer errors}  % experiments/FIGS/data/v2/index.json :: binary probe qcd_b_vs_light: (1 - AUC) of FLAV_F1 over FLAV_F0, run by run :: sha256 c6c4138ee5c86423 :: prose twin: two significant figures, direction in words
\newcommand{\MlstCmpQcdBFoneOverFzeroPrMax}{up to 22\% fewer errors}  % experiments/FIGS/data/v2/index.json :: binary probe qcd_b_vs_light: (1 - AUC) of FLAV_F1 over FLAV_F0, run by run :: sha256 c6c4138ee5c86423 :: prose twin of the run farther from 1
\newcommand{\MlstCmpQcdBFonerOverFone}{1.054 and 1.061}  % experiments/FIGS/data/v2/index.json :: binary probe qcd_b_vs_light: (1 - AUC) of FLAV_F1R over FLAV_F1, run by run :: sha256 c6c4138ee5c86423
\newcommand{\MlstCmpQcdBFonerOverFonePr}{equal within 7\%}  % experiments/FIGS/data/v2/index.json :: binary probe qcd_b_vs_light: (1 - AUC) of FLAV_F1R over FLAV_F1, run by run :: sha256 c6c4138ee5c86423 :: prose twin: two significant figures, direction in words
\newcommand{\MlstCmpQcdBFonerOverFonePrMax}{equal within 7\%}  % experiments/FIGS/data/v2/index.json :: binary probe qcd_b_vs_light: (1 - AUC) of FLAV_F1R over FLAV_F1, run by run :: sha256 c6c4138ee5c86423 :: prose twin of the run farther from 1
\newcommand{\MlstCmpQcdBFonerOverFzero}{0.820 and 0.885}  % experiments/FIGS/data/v2/index.json :: binary probe qcd_b_vs_light: (1 - AUC) of FLAV_F1R over FLAV_F0, run by run :: sha256 c6c4138ee5c86423
\newcommand{\MlstCmpQcdBFonerOverFzeroPr}{at least 11\% fewer errors}  % experiments/FIGS/data/v2/index.json :: binary probe qcd_b_vs_light: (1 - AUC) of FLAV_F1R over FLAV_F0, run by run :: sha256 c6c4138ee5c86423 :: prose twin: two significant figures, direction in words
\newcommand{\MlstCmpQcdBFonerOverFzeroPrMax}{up to 18\% fewer errors}  % experiments/FIGS/data/v2/index.json :: binary probe qcd_b_vs_light: (1 - AUC) of FLAV_F1R over FLAV_F0, run by run :: sha256 c6c4138ee5c86423 :: prose twin of the run farther from 1
\newcommand{\MlstCmpQcdBFzeroOverCoarse}{0.962 and 0.981}  % experiments/FIGS/data/v2/index.json :: binary probe qcd_b_vs_light: (1 - AUC) of FLAV_F0 over R16_Q1, run by run :: sha256 c6c4138ee5c86423
\newcommand{\MlstCmpQcdBFzeroOverCoarsePr}{equal within 4\%}  % experiments/FIGS/data/v2/index.json :: binary probe qcd_b_vs_light: (1 - AUC) of FLAV_F0 over R16_Q1, run by run :: sha256 c6c4138ee5c86423 :: prose twin: two significant figures, direction in words
\newcommand{\MlstCmpQcdBFzeroOverCoarsePrMax}{equal within 4\%}  % experiments/FIGS/data/v2/index.json :: binary probe qcd_b_vs_light: (1 - AUC) of FLAV_F0 over R16_Q1, run by run :: sha256 c6c4138ee5c86423 :: prose twin of the run farther from 1
\newcommand{\MlstCmpQcdBFoneOverCoarse}{0.749 and 0.819}  % experiments/FIGS/data/v2/index.json :: binary probe qcd_b_vs_light: (1 - AUC) of FLAV_F1 over R16_Q1, run by run :: sha256 c6c4138ee5c86423
\newcommand{\MlstCmpQcdBFoneOverCoarsePr}{at least 18\% fewer errors}  % experiments/FIGS/data/v2/index.json :: binary probe qcd_b_vs_light: (1 - AUC) of FLAV_F1 over R16_Q1, run by run :: sha256 c6c4138ee5c86423 :: prose twin: two significant figures, direction in words
\newcommand{\MlstCmpQcdBFoneOverCoarsePrMax}{up to 25\% fewer errors}  % experiments/FIGS/data/v2/index.json :: binary probe qcd_b_vs_light: (1 - AUC) of FLAV_F1 over R16_Q1, run by run :: sha256 c6c4138ee5c86423 :: prose twin of the run farther from 1
\newcommand{\MlstCmpQcdBFonerOverCoarse}{0.789 and 0.869}  % experiments/FIGS/data/v2/index.json :: binary probe qcd_b_vs_light: (1 - AUC) of FLAV_F1R over R16_Q1, run by run :: sha256 c6c4138ee5c86423
\newcommand{\MlstCmpQcdBFonerOverCoarsePr}{at least 13\% fewer errors}  % experiments/FIGS/data/v2/index.json :: binary probe qcd_b_vs_light: (1 - AUC) of FLAV_F1R over R16_Q1, run by run :: sha256 c6c4138ee5c86423 :: prose twin: two significant figures, direction in words
\newcommand{\MlstCmpQcdBFonerOverCoarsePrMax}{up to 21\% fewer errors}  % experiments/FIGS/data/v2/index.json :: binary probe qcd_b_vs_light: (1 - AUC) of FLAV_F1R over R16_Q1, run by run :: sha256 c6c4138ee5c86423 :: prose twin of the run farther from 1
\newcommand{\MlstCmpQcdBFzeroOverFull}{1.56 and 1.45}  % experiments/FIGS/data/v2/index.json :: binary probe qcd_b_vs_light: (1 - AUC) of FLAV_F0 over L188, run by run :: sha256 c6c4138ee5c86423
\newcommand{\MlstCmpQcdBFzeroOverFullPr}{at least 45\% more errors}  % experiments/FIGS/data/v2/index.json :: binary probe qcd_b_vs_light: (1 - AUC) of FLAV_F0 over L188, run by run :: sha256 c6c4138ee5c86423 :: prose twin: two significant figures, direction in words
\newcommand{\MlstCmpQcdBFzeroOverFullPrMax}{up to 56\% more errors}  % experiments/FIGS/data/v2/index.json :: binary probe qcd_b_vs_light: (1 - AUC) of FLAV_F0 over L188, run by run :: sha256 c6c4138ee5c86423 :: prose twin of the run farther from 1
\newcommand{\MlstCmpQcdBFoneOverFull}{1.212 and 1.213}  % experiments/FIGS/data/v2/index.json :: binary probe qcd_b_vs_light: (1 - AUC) of FLAV_F1 over L188, run by run :: sha256 c6c4138ee5c86423
\newcommand{\MlstCmpQcdBFoneOverFullPr}{at least 21\% more errors}  % experiments/FIGS/data/v2/index.json :: binary probe qcd_b_vs_light: (1 - AUC) of FLAV_F1 over L188, run by run :: sha256 c6c4138ee5c86423 :: prose twin: two significant figures, direction in words
\newcommand{\MlstCmpQcdBFoneOverFullPrMax}{up to 21\% more errors}  % experiments/FIGS/data/v2/index.json :: binary probe qcd_b_vs_light: (1 - AUC) of FLAV_F1 over L188, run by run :: sha256 c6c4138ee5c86423 :: prose twin of the run farther from 1
\newcommand{\MlstCmpQcdBFonerOverFull}{1.277 and 1.286}  % experiments/FIGS/data/v2/index.json :: binary probe qcd_b_vs_light: (1 - AUC) of FLAV_F1R over L188, run by run :: sha256 c6c4138ee5c86423
\newcommand{\MlstCmpQcdBFonerOverFullPr}{at least 28\% more errors}  % experiments/FIGS/data/v2/index.json :: binary probe qcd_b_vs_light: (1 - AUC) of FLAV_F1R over L188, run by run :: sha256 c6c4138ee5c86423 :: prose twin: two significant figures, direction in words
\newcommand{\MlstCmpQcdBFonerOverFullPrMax}{up to 29\% more errors}  % experiments/FIGS/data/v2/index.json :: binary probe qcd_b_vs_light: (1 - AUC) of FLAV_F1R over L188, run by run :: sha256 c6c4138ee5c86423 :: prose twin of the run farther from 1
\newcommand{\MlstCmpQcdCResOverFull}{1.061\,\ensuremath{\pm}\,0.005}  % experiments/FIGS/data/v2/index.json :: binary probe qcd_c_vs_light: (1 - AUC) of L162 over L188, run by run :: sha256 c6c4138ee5c86423
\newcommand{\MlstCmpQcdCResOverFullPr}{6.1\% more errors}  % experiments/FIGS/data/v2/index.json :: binary probe qcd_c_vs_light: (1 - AUC) of L162 over L188, run by run :: sha256 c6c4138ee5c86423 :: prose twin: two significant figures, direction in words
\newcommand{\MlstCmpQcdCMidOverFull}{1.072\,\ensuremath{\pm}\,0.008}  % experiments/FIGS/data/v2/index.json :: binary probe qcd_c_vs_light: (1 - AUC) of R42_Q1 over L188, run by run :: sha256 c6c4138ee5c86423
\newcommand{\MlstCmpQcdCMidOverFullPr}{7.2\% more errors}  % experiments/FIGS/data/v2/index.json :: binary probe qcd_c_vs_light: (1 - AUC) of R42_Q1 over L188, run by run :: sha256 c6c4138ee5c86423 :: prose twin: two significant figures, direction in words
\newcommand{\MlstCmpQcdCCoarseOverFull}{1.407\,\ensuremath{\pm}\,0.021}  % experiments/FIGS/data/v2/index.json :: binary probe qcd_c_vs_light: (1 - AUC) of R16_Q1 over L188, run by run :: sha256 c6c4138ee5c86423
\newcommand{\MlstCmpQcdCCoarseOverFullPr}{41\% more errors}  % experiments/FIGS/data/v2/index.json :: binary probe qcd_c_vs_light: (1 - AUC) of R16_Q1 over L188, run by run :: sha256 c6c4138ee5c86423 :: prose twin: two significant figures, direction in words
\newcommand{\MlstCmpQcdCFoneOverFzero}{0.915 and 0.944}  % experiments/FIGS/data/v2/index.json :: binary probe qcd_c_vs_light: (1 - AUC) of FLAV_F1 over FLAV_F0, run by run :: sha256 c6c4138ee5c86423
\newcommand{\MlstCmpQcdCFoneOverFzeroPr}{at least 5.6\% fewer errors}  % experiments/FIGS/data/v2/index.json :: binary probe qcd_c_vs_light: (1 - AUC) of FLAV_F1 over FLAV_F0, run by run :: sha256 c6c4138ee5c86423 :: prose twin: two significant figures, direction in words
\newcommand{\MlstCmpQcdCFoneOverFzeroPrMax}{up to 8.5\% fewer errors}  % experiments/FIGS/data/v2/index.json :: binary probe qcd_c_vs_light: (1 - AUC) of FLAV_F1 over FLAV_F0, run by run :: sha256 c6c4138ee5c86423 :: prose twin of the run farther from 1
\newcommand{\MlstCmpQcdCFonerOverFone}{0.9535 and 0.9520}  % experiments/FIGS/data/v2/index.json :: binary probe qcd_c_vs_light: (1 - AUC) of FLAV_F1R over FLAV_F1, run by run :: sha256 c6c4138ee5c86423
\newcommand{\MlstCmpQcdCFonerOverFonePr}{at least 4.6\% fewer errors}  % experiments/FIGS/data/v2/index.json :: binary probe qcd_c_vs_light: (1 - AUC) of FLAV_F1R over FLAV_F1, run by run :: sha256 c6c4138ee5c86423 :: prose twin: two significant figures, direction in words
\newcommand{\MlstCmpQcdCFonerOverFonePrMax}{up to 4.8\% fewer errors}  % experiments/FIGS/data/v2/index.json :: binary probe qcd_c_vs_light: (1 - AUC) of FLAV_F1R over FLAV_F1, run by run :: sha256 c6c4138ee5c86423 :: prose twin of the run farther from 1
\newcommand{\MlstCmpQcdCFonerOverFzero}{0.872 and 0.899}  % experiments/FIGS/data/v2/index.json :: binary probe qcd_c_vs_light: (1 - AUC) of FLAV_F1R over FLAV_F0, run by run :: sha256 c6c4138ee5c86423
\newcommand{\MlstCmpQcdCFonerOverFzeroPr}{at least 10\% fewer errors}  % experiments/FIGS/data/v2/index.json :: binary probe qcd_c_vs_light: (1 - AUC) of FLAV_F1R over FLAV_F0, run by run :: sha256 c6c4138ee5c86423 :: prose twin: two significant figures, direction in words
\newcommand{\MlstCmpQcdCFonerOverFzeroPrMax}{up to 13\% fewer errors}  % experiments/FIGS/data/v2/index.json :: binary probe qcd_c_vs_light: (1 - AUC) of FLAV_F1R over FLAV_F0, run by run :: sha256 c6c4138ee5c86423 :: prose twin of the run farther from 1
\newcommand{\MlstCmpQcdCFzeroOverCoarse}{0.952 and 0.934}  % experiments/FIGS/data/v2/index.json :: binary probe qcd_c_vs_light: (1 - AUC) of FLAV_F0 over R16_Q1, run by run :: sha256 c6c4138ee5c86423
\newcommand{\MlstCmpQcdCFzeroOverCoarsePr}{at least 4.8\% fewer errors}  % experiments/FIGS/data/v2/index.json :: binary probe qcd_c_vs_light: (1 - AUC) of FLAV_F0 over R16_Q1, run by run :: sha256 c6c4138ee5c86423 :: prose twin: two significant figures, direction in words
\newcommand{\MlstCmpQcdCFzeroOverCoarsePrMax}{up to 6.6\% fewer errors}  % experiments/FIGS/data/v2/index.json :: binary probe qcd_c_vs_light: (1 - AUC) of FLAV_F0 over R16_Q1, run by run :: sha256 c6c4138ee5c86423 :: prose twin of the run farther from 1
\newcommand{\MlstCmpQcdCFoneOverCoarse}{0.871 and 0.882}  % experiments/FIGS/data/v2/index.json :: binary probe qcd_c_vs_light: (1 - AUC) of FLAV_F1 over R16_Q1, run by run :: sha256 c6c4138ee5c86423
\newcommand{\MlstCmpQcdCFoneOverCoarsePr}{at least 12\% fewer errors}  % experiments/FIGS/data/v2/index.json :: binary probe qcd_c_vs_light: (1 - AUC) of FLAV_F1 over R16_Q1, run by run :: sha256 c6c4138ee5c86423 :: prose twin: two significant figures, direction in words
\newcommand{\MlstCmpQcdCFoneOverCoarsePrMax}{up to 13\% fewer errors}  % experiments/FIGS/data/v2/index.json :: binary probe qcd_c_vs_light: (1 - AUC) of FLAV_F1 over R16_Q1, run by run :: sha256 c6c4138ee5c86423 :: prose twin of the run farther from 1
\newcommand{\MlstCmpQcdCFonerOverCoarse}{0.8301 and 0.8396}  % experiments/FIGS/data/v2/index.json :: binary probe qcd_c_vs_light: (1 - AUC) of FLAV_F1R over R16_Q1, run by run :: sha256 c6c4138ee5c86423
\newcommand{\MlstCmpQcdCFonerOverCoarsePr}{at least 16\% fewer errors}  % experiments/FIGS/data/v2/index.json :: binary probe qcd_c_vs_light: (1 - AUC) of FLAV_F1R over R16_Q1, run by run :: sha256 c6c4138ee5c86423 :: prose twin: two significant figures, direction in words
\newcommand{\MlstCmpQcdCFonerOverCoarsePrMax}{up to 17\% fewer errors}  % experiments/FIGS/data/v2/index.json :: binary probe qcd_c_vs_light: (1 - AUC) of FLAV_F1R over R16_Q1, run by run :: sha256 c6c4138ee5c86423 :: prose twin of the run farther from 1
\newcommand{\MlstCmpQcdCFzeroOverFull}{1.352 and 1.293}  % experiments/FIGS/data/v2/index.json :: binary probe qcd_c_vs_light: (1 - AUC) of FLAV_F0 over L188, run by run :: sha256 c6c4138ee5c86423
\newcommand{\MlstCmpQcdCFzeroOverFullPr}{at least 29\% more errors}  % experiments/FIGS/data/v2/index.json :: binary probe qcd_c_vs_light: (1 - AUC) of FLAV_F0 over L188, run by run :: sha256 c6c4138ee5c86423 :: prose twin: two significant figures, direction in words
\newcommand{\MlstCmpQcdCFzeroOverFullPrMax}{up to 35\% more errors}  % experiments/FIGS/data/v2/index.json :: binary probe qcd_c_vs_light: (1 - AUC) of FLAV_F0 over L188, run by run :: sha256 c6c4138ee5c86423 :: prose twin of the run farther from 1
\newcommand{\MlstCmpQcdCFoneOverFull}{1.237 and 1.220}  % experiments/FIGS/data/v2/index.json :: binary probe qcd_c_vs_light: (1 - AUC) of FLAV_F1 over L188, run by run :: sha256 c6c4138ee5c86423
\newcommand{\MlstCmpQcdCFoneOverFullPr}{at least 22\% more errors}  % experiments/FIGS/data/v2/index.json :: binary probe qcd_c_vs_light: (1 - AUC) of FLAV_F1 over L188, run by run :: sha256 c6c4138ee5c86423 :: prose twin: two significant figures, direction in words
\newcommand{\MlstCmpQcdCFoneOverFullPrMax}{up to 24\% more errors}  % experiments/FIGS/data/v2/index.json :: binary probe qcd_c_vs_light: (1 - AUC) of FLAV_F1 over L188, run by run :: sha256 c6c4138ee5c86423 :: prose twin of the run farther from 1
\newcommand{\MlstCmpQcdCFonerOverFull}{1.179 and 1.162}  % experiments/FIGS/data/v2/index.json :: binary probe qcd_c_vs_light: (1 - AUC) of FLAV_F1R over L188, run by run :: sha256 c6c4138ee5c86423
\newcommand{\MlstCmpQcdCFonerOverFullPr}{at least 16\% more errors}  % experiments/FIGS/data/v2/index.json :: binary probe qcd_c_vs_light: (1 - AUC) of FLAV_F1R over L188, run by run :: sha256 c6c4138ee5c86423 :: prose twin: two significant figures, direction in words
\newcommand{\MlstCmpQcdCFonerOverFullPrMax}{up to 18\% more errors}  % experiments/FIGS/data/v2/index.json :: binary probe qcd_c_vs_light: (1 - AUC) of FLAV_F1R over L188, run by run :: sha256 c6c4138ee5c86423 :: prose twin of the run farther from 1
\newcommand{\MlstLpEemmAucFull}{1.0000\,\ensuremath{\pm}\,\ensuremath{<}0.0001}  % experiments/FIGS/data/v2/index.json :: linear probe ee_vs_mm (fits_v4_sib): AUC, L188 :: sha256 c6c4138ee5c86423
\newcommand{\MlstLpEemmAucRes}{1.0000\,\ensuremath{\pm}\,\ensuremath{<}0.0001}  % experiments/FIGS/data/v2/index.json :: linear probe ee_vs_mm (fits_v4_sib): AUC, L162 :: sha256 c6c4138ee5c86423
\newcommand{\MlstLpEemmAucMid}{0.9610\,\ensuremath{\pm}\,0.0030}  % experiments/FIGS/data/v2/index.json :: linear probe ee_vs_mm (fits_v4_sib): AUC, R42_Q1 :: sha256 c6c4138ee5c86423
\newcommand{\MlstLpEemmAucCoarse}{0.960\,\ensuremath{\pm}\,0.020}  % experiments/FIGS/data/v2/index.json :: linear probe ee_vs_mm (fits_v4_sib): AUC, R16_Q1 :: sha256 c6c4138ee5c86423
\newcommand{\MlstLpEemmAucFzero}{1.0000 and 1.0000}  % experiments/FIGS/data/v2/index.json :: linear probe ee_vs_mm (fits_v4_sib): AUC, FLAV_F0 :: sha256 c6c4138ee5c86423
\newcommand{\MlstLpEemmAucFone}{1.0000 and 1.0000}  % experiments/FIGS/data/v2/index.json :: linear probe ee_vs_mm (fits_v4_sib): AUC, FLAV_F1 :: sha256 c6c4138ee5c86423
\newcommand{\MlstLpEemmAucFoner}{1.0000 and 1.0000}  % experiments/FIGS/data/v2/index.json :: linear probe ee_vs_mm (fits_v4_sib): AUC, FLAV_F1R :: sha256 c6c4138ee5c86423
\newcommand{\MlstLpEemmAucInit}{0.9995\,\ensuremath{\pm}\,0.0003}  % experiments/FIGS/data/v2/index.json :: linear probe ee_vs_mm (fits_v4_sib): AUC, INIT :: sha256 c6c4138ee5c86423
\newcommand{\MlstLpEemmAucInitBn}{1.0000\,\ensuremath{\pm}\,\ensuremath{<}0.0001}  % experiments/FIGS/data/v2/index.json :: linear probe ee_vs_mm (fits_v4_sib): AUC, INIT_BN :: sha256 c6c4138ee5c86423
\newcommand{\MlstLpEemmAucRandMerged}{1.0000\,\ensuremath{\pm}\,\ensuremath{<}0.0001}  % experiments/FIGS/data/v2/index.json :: linear probe ee_vs_mm (fits_v4_sib): AUC, the random partitions that merge the pair :: sha256 c6c4138ee5c86423
\newcommand{\MlstPairEemmNTrain}{9{,}571}  % experiments/FIGS/data/v2/index.json :: fits_v4_sib ee_vs_mm: n_train :: sha256 c6c4138ee5c86423
\newcommand{\MlstPairEemmNVal}{1{,}068}  % experiments/FIGS/data/v2/index.json :: fits_v4_sib ee_vs_mm: n_val :: sha256 c6c4138ee5c86423
\newcommand{\MlstPairEemmNTestPerSide}{128{,}833--135{,}245}  % experiments/FIGS/data/v2/index.json :: fits_v4_sib ee_vs_mm: n_test_signal and n_test_background :: sha256 c6c4138ee5c86423
\newcommand{\MlstLpJcTwoAucFull}{0.9801\,\ensuremath{\pm}\,\ensuremath{<}0.0001}  % experiments/FIGS/data/v2/index.json :: linear probe JetClass-II N=1000000: AUC averaged over the classes, L188 :: sha256 c6c4138ee5c86423
\newcommand{\MlstLpJcTwoAccFull}{0.4348\,\ensuremath{\pm}\,0.0003}  % experiments/FIGS/data/v2/index.json :: linear probe JetClass-II N=1000000: accuracy, L188 :: sha256 c6c4138ee5c86423
\newcommand{\MlstLpJcTwoAucRes}{0.9796\,\ensuremath{\pm}\,0.0001}  % experiments/FIGS/data/v2/index.json :: linear probe JetClass-II N=1000000: AUC averaged over the classes, L162 :: sha256 c6c4138ee5c86423
\newcommand{\MlstLpJcTwoAccRes}{0.4328\,\ensuremath{\pm}\,0.0008}  % experiments/FIGS/data/v2/index.json :: linear probe JetClass-II N=1000000: accuracy, L162 :: sha256 c6c4138ee5c86423
\newcommand{\MlstLpJcTwoAucMid}{0.9756\,\ensuremath{\pm}\,0.0002}  % experiments/FIGS/data/v2/index.json :: linear probe JetClass-II N=1000000: AUC averaged over the classes, R42_Q1 :: sha256 c6c4138ee5c86423
\newcommand{\MlstLpJcTwoAccMid}{0.361\,\ensuremath{\pm}\,0.005}  % experiments/FIGS/data/v2/index.json :: linear probe JetClass-II N=1000000: accuracy, R42_Q1 :: sha256 c6c4138ee5c86423
\newcommand{\MlstLpJcTwoAucCoarse}{0.9691\,\ensuremath{\pm}\,0.0004}  % experiments/FIGS/data/v2/index.json :: linear probe JetClass-II N=1000000: AUC averaged over the classes, R16_Q1 :: sha256 c6c4138ee5c86423
\newcommand{\MlstLpJcTwoAccCoarse}{0.334\,\ensuremath{\pm}\,0.007}  % experiments/FIGS/data/v2/index.json :: linear probe JetClass-II N=1000000: accuracy, R16_Q1 :: sha256 c6c4138ee5c86423
\newcommand{\MlstLpJcTwoAucResMass}{0.9791 and 0.9794}  % experiments/FIGS/data/v2/index.json :: linear probe JetClass-II N=1000000: AUC averaged over the classes, L162_MASS :: sha256 c6c4138ee5c86423
\newcommand{\MlstLpJcTwoAccResMass}{0.4271 and 0.4305}  % experiments/FIGS/data/v2/index.json :: linear probe JetClass-II N=1000000: accuracy, L162_MASS :: sha256 c6c4138ee5c86423
\newcommand{\MlstLpJcTwoAucCoarseMass}{0.9679 and 0.9696}  % experiments/FIGS/data/v2/index.json :: linear probe JetClass-II N=1000000: AUC averaged over the classes, R16_Q1_MASS :: sha256 c6c4138ee5c86423
\newcommand{\MlstLpJcTwoAccCoarseMass}{0.312 and 0.350}  % experiments/FIGS/data/v2/index.json :: linear probe JetClass-II N=1000000: accuracy, R16_Q1_MASS :: sha256 c6c4138ee5c86423
\newcommand{\MlstLpJcTwoAucFzero}{0.9717 and 0.9710}  % experiments/FIGS/data/v2/index.json :: linear probe JetClass-II N=1000000: AUC averaged over the classes, FLAV_F0 :: sha256 c6c4138ee5c86423
\newcommand{\MlstLpJcTwoAccFzero}{0.3769 and 0.3723}  % experiments/FIGS/data/v2/index.json :: linear probe JetClass-II N=1000000: accuracy, FLAV_F0 :: sha256 c6c4138ee5c86423
\newcommand{\MlstLpJcTwoAucFone}{0.9758 and 0.9759}  % experiments/FIGS/data/v2/index.json :: linear probe JetClass-II N=1000000: AUC averaged over the classes, FLAV_F1 :: sha256 c6c4138ee5c86423
\newcommand{\MlstLpJcTwoAccFone}{0.3945 and 0.3963}  % experiments/FIGS/data/v2/index.json :: linear probe JetClass-II N=1000000: accuracy, FLAV_F1 :: sha256 c6c4138ee5c86423
\newcommand{\MlstLpJcTwoAucFoner}{0.9760 and 0.9760}  % experiments/FIGS/data/v2/index.json :: linear probe JetClass-II N=1000000: AUC averaged over the classes, FLAV_F1R :: sha256 c6c4138ee5c86423
\newcommand{\MlstLpJcTwoAccFoner}{0.4071 and 0.4089}  % experiments/FIGS/data/v2/index.json :: linear probe JetClass-II N=1000000: accuracy, FLAV_F1R :: sha256 c6c4138ee5c86423
\newcommand{\MlstLpJcTwoAucInit}{0.8121\,\ensuremath{\pm}\,0.0019}  % experiments/FIGS/data/v2/index.json :: linear probe JetClass-II N=1000000: AUC averaged over the classes, INIT :: sha256 c6c4138ee5c86423
\newcommand{\MlstLpJcTwoAccInit}{0.0724\,\ensuremath{\pm}\,0.0015}  % experiments/FIGS/data/v2/index.json :: linear probe JetClass-II N=1000000: accuracy, INIT :: sha256 c6c4138ee5c86423
\newcommand{\MlstLpJcTwoAucInitBn}{0.866\,\ensuremath{\pm}\,0.004}  % experiments/FIGS/data/v2/index.json :: linear probe JetClass-II N=1000000: AUC averaged over the classes, INIT_BN :: sha256 c6c4138ee5c86423
\newcommand{\MlstLpJcTwoAccInitBn}{0.1017\,\ensuremath{\pm}\,0.0017}  % experiments/FIGS/data/v2/index.json :: linear probe JetClass-II N=1000000: accuracy, INIT_BN :: sha256 c6c4138ee5c86423
\newcommand{\MlstLpJcTwoAucRand}{0.9777\,\ensuremath{\pm}\,0.0002}  % experiments/FIGS/data/v2/index.json :: linear probe JetClass-II N=1000000: AUC averaged over the classes, the random partitions :: sha256 c6c4138ee5c86423
\newcommand{\MlstLpJcOneAccInitBn}{0.475\,\ensuremath{\pm}\,0.005}  % experiments/FIGS/data/v2/index.json :: linear probe JetClass N=1000000: accuracy, INIT_BN :: sha256 c6c4138ee5c86423
\newcommand{\MlstBnMaxDev}{0.2\%}  % experiments/FIGS/data/v2/index.json :: largest |r(best70_bn) / r(best70) - 1| of a JetClass-II ratio of 1 - AUC to the 188-class run, per cent :: sha256 c6c4138ee5c86423
\newcommand{\MlstBnMaxDevPair}{0.7\%}  % experiments/FIGS/data/v2/index.json :: largest |r(best70_bn) / r(best70) - 1| of a ratio of 1 - AUC on a class pair to the 188-class run, levels of the hierarchy, mass models and controlled partitions, per cent :: sha256 c6c4138ee5c86423 :: pairs with a check: bc_vs_bq_cs, bb_vs_cc, bbqq_vs_ccqq, qcd_bb_vs_cc, qcd_b_vs_light, qcd_c_vs_light; without: ee_vs_mm
\newcommand{\MlstGpuTypeBound}{1.4\%}  % experiments/FIGS/data/v2/index.json :: largest |label-set ratio of 1 - AUC of index 3 / mean of the other indices - 1|, JetClass-II largest set, per cent :: sha256 c6c4138ee5c86423 :: mtx-l188-s3 trained on NVIDIA GeForce RTX 3090; largest for R16_Q1; label sets L162 0.09%, R42_Q1 0.70%, R16_Q1 1.37%
\newcommand{\MlstGpuOtherSeedsDiff}{1.9\%}  % experiments/FIGS/data/v2/index.json :: largest (max - min) / mean of a label-set ratio of 1 - AUC over the indices other than 3, JetClass-II largest set, per cent :: sha256 c6c4138ee5c86423 :: label sets L162 0.90%, R42_Q1 1.14%, R16_Q1 1.93%
\newcommand{\MlstPerClassResNDoubled}{0}  % experiments/FIGS/data/v2/index.json :: JetClass-II classes whose mean ratio of 1 - AUC (L162 over 188) is at least 2 (count) :: sha256 c6c4138ee5c86423
\newcommand{\MlstPerClassResSeedSdMedian}{0.011}  % experiments/FIGS/data/v2/index.json :: median over classes of the SD over runs of the per-class ratio, L162 over 188 :: sha256 c6c4138ee5c86423
\newcommand{\MlstPerClassResTopClass}{QCD, \ensuremath{ccss}}  % experiments/FIGS/data/v2/index.json :: class with the largest mean ratio, L162 :: sha256 c6c4138ee5c86423
\newcommand{\MlstPerClassResTopRatio}{1.129\,\ensuremath{\pm}\,0.010}  % experiments/FIGS/data/v2/index.json :: per-class ratio of 1 - AUC of that class, L162 over 188, run by run :: sha256 c6c4138ee5c86423
\newcommand{\MlstPerClassResTopRatioPr}{13\% more errors}  % experiments/FIGS/data/v2/index.json :: per-class ratio of 1 - AUC of that class, L162 over 188, run by run :: sha256 c6c4138ee5c86423 :: prose twin: two significant figures, direction in words
\newcommand{\MlstPerClassResLeptonN}{0}  % experiments/FIGS/data/v2/index.json :: resonant classes sharing their L162 group with a decay that differs only in lepton flavor (count) :: sha256 c6c4138ee5c86423
\newcommand{\MlstPerClassResLeptonMedian}{--}  % experiments/FIGS/data/v2/index.json :: median over those classes of the mean per-class ratio, L162 over 188 :: sha256 c6c4138ee5c86423
\newcommand{\MlstPerClassResLeptonMedianPr}{\pending{no run yet}}  % experiments/FIGS/data/v2/index.json :: median over those classes of the mean per-class ratio, L162 over 188 :: sha256 c6c4138ee5c86423 :: prose twin: two significant figures, direction in words
\newcommand{\MlstPerClassResLeptonShare}{--}  % experiments/FIGS/data/v2/index.json :: their share of the summed ln ratio over the resonant classes, L162 :: sha256 c6c4138ee5c86423
\newcommand{\MlstPerClassResQuarkN}{0}  % experiments/FIGS/data/v2/index.json :: resonant classes sharing their L162 group with a decay that differs only in quark flavor (count) :: sha256 c6c4138ee5c86423
\newcommand{\MlstPerClassResQuarkMedian}{--}  % experiments/FIGS/data/v2/index.json :: median over those classes of the mean per-class ratio, L162 over 188 :: sha256 c6c4138ee5c86423
\newcommand{\MlstPerClassResQuarkMedianPr}{\pending{no run yet}}  % experiments/FIGS/data/v2/index.json :: median over those classes of the mean per-class ratio, L162 over 188 :: sha256 c6c4138ee5c86423 :: prose twin: two significant figures, direction in words
\newcommand{\MlstPerClassResQuarkShare}{--}  % experiments/FIGS/data/v2/index.json :: their share of the summed ln ratio over the resonant classes, L162 :: sha256 c6c4138ee5c86423
\newcommand{\MlstPerClassResQuarkLeptonN}{0}  % experiments/FIGS/data/v2/index.json :: resonant classes counted in both the quark- and the lepton-flavor set, L162 (count) :: sha256 c6c4138ee5c86423
\newcommand{\MlstPerClassMidNDoubled}{40}  % experiments/FIGS/data/v2/index.json :: JetClass-II classes whose mean ratio of 1 - AUC (R42_Q1 over 188) is at least 2 (count) :: sha256 c6c4138ee5c86423
\newcommand{\MlstPerClassMidSeedSdMedian}{0.023}  % experiments/FIGS/data/v2/index.json :: median over classes of the SD over runs of the per-class ratio, R42_Q1 over 188 :: sha256 c6c4138ee5c86423
\newcommand{\MlstPerClassMidTopClass}{\ensuremath{X\to YY\to bb\mu\mu}}  % experiments/FIGS/data/v2/index.json :: class with the largest mean ratio, R42_Q1 :: sha256 c6c4138ee5c86423
\newcommand{\MlstPerClassMidTopRatio}{4.06\,\ensuremath{\pm}\,0.14}  % experiments/FIGS/data/v2/index.json :: per-class ratio of 1 - AUC of that class, R42_Q1 over 188, run by run :: sha256 c6c4138ee5c86423
\newcommand{\MlstPerClassMidTopRatioPr}{4.1 times the errors}  % experiments/FIGS/data/v2/index.json :: per-class ratio of 1 - AUC of that class, R42_Q1 over 188, run by run :: sha256 c6c4138ee5c86423 :: prose twin: two significant figures, direction in words
\newcommand{\MlstPerClassMidLeptonN}{78}  % experiments/FIGS/data/v2/index.json :: resonant classes sharing their R42_Q1 group with a decay that differs only in lepton flavor (count) :: sha256 c6c4138ee5c86423
\newcommand{\MlstPerClassMidLeptonMedian}{2.02}  % experiments/FIGS/data/v2/index.json :: median over those classes of the mean per-class ratio, R42_Q1 over 188 :: sha256 c6c4138ee5c86423
\newcommand{\MlstPerClassMidLeptonMedianPr}{2.0 times the errors}  % experiments/FIGS/data/v2/index.json :: median over those classes of the mean per-class ratio, R42_Q1 over 188 :: sha256 c6c4138ee5c86423 :: prose twin: two significant figures, direction in words
\newcommand{\MlstPerClassMidLeptonShare}{90\%}  % experiments/FIGS/data/v2/index.json :: their share of the summed ln ratio over the resonant classes, R42_Q1 :: sha256 c6c4138ee5c86423
\newcommand{\MlstPerClassMidQuarkN}{111}  % experiments/FIGS/data/v2/index.json :: resonant classes sharing their R42_Q1 group with a decay that differs only in quark flavor (count) :: sha256 c6c4138ee5c86423
\newcommand{\MlstPerClassMidQuarkMedian}{1.14}  % experiments/FIGS/data/v2/index.json :: median over those classes of the mean per-class ratio, R42_Q1 over 188 :: sha256 c6c4138ee5c86423
\newcommand{\MlstPerClassMidQuarkMedianPr}{14\% more errors}  % experiments/FIGS/data/v2/index.json :: median over those classes of the mean per-class ratio, R42_Q1 over 188 :: sha256 c6c4138ee5c86423 :: prose twin: two significant figures, direction in words
\newcommand{\MlstPerClassMidQuarkShare}{53\%}  % experiments/FIGS/data/v2/index.json :: their share of the summed ln ratio over the resonant classes, R42_Q1 :: sha256 c6c4138ee5c86423
\newcommand{\MlstPerClassMidQuarkLeptonN}{44}  % experiments/FIGS/data/v2/index.json :: resonant classes counted in both the quark- and the lepton-flavor set, R42_Q1 (count) :: sha256 c6c4138ee5c86423
\newcommand{\MlstPerClassCoarseNDoubled}{60}  % experiments/FIGS/data/v2/index.json :: JetClass-II classes whose mean ratio of 1 - AUC (R16_Q1 over 188) is at least 2 (count) :: sha256 c6c4138ee5c86423
\newcommand{\MlstPerClassCoarseSeedSdMedian}{0.055}  % experiments/FIGS/data/v2/index.json :: median over classes of the SD over runs of the per-class ratio, R16_Q1 over 188 :: sha256 c6c4138ee5c86423
\newcommand{\MlstPerClassCoarseTopClass}{\ensuremath{X\to YY\to bbee}}  % experiments/FIGS/data/v2/index.json :: class with the largest mean ratio, R16_Q1 :: sha256 c6c4138ee5c86423
\newcommand{\MlstPerClassCoarseTopRatio}{5.2\,\ensuremath{\pm}\,0.4}  % experiments/FIGS/data/v2/index.json :: per-class ratio of 1 - AUC of that class, R16_Q1 over 188, run by run :: sha256 c6c4138ee5c86423
\newcommand{\MlstPerClassCoarseTopRatioPr}{5.2 times the errors}  % experiments/FIGS/data/v2/index.json :: per-class ratio of 1 - AUC of that class, R16_Q1 over 188, run by run :: sha256 c6c4138ee5c86423 :: prose twin: two significant figures, direction in words
\newcommand{\MlstPerClassCoarseLeptonN}{78}  % experiments/FIGS/data/v2/index.json :: resonant classes sharing their R16_Q1 group with a decay that differs only in lepton flavor (count) :: sha256 c6c4138ee5c86423
\newcommand{\MlstPerClassCoarseLeptonMedian}{2.27}  % experiments/FIGS/data/v2/index.json :: median over those classes of the mean per-class ratio, R16_Q1 over 188 :: sha256 c6c4138ee5c86423
\newcommand{\MlstPerClassCoarseLeptonMedianPr}{2.3 times the errors}  % experiments/FIGS/data/v2/index.json :: median over those classes of the mean per-class ratio, R16_Q1 over 188 :: sha256 c6c4138ee5c86423 :: prose twin: two significant figures, direction in words
\newcommand{\MlstPerClassCoarseLeptonShare}{68\%}  % experiments/FIGS/data/v2/index.json :: their share of the summed ln ratio over the resonant classes, R16_Q1 :: sha256 c6c4138ee5c86423
\newcommand{\MlstPerClassCoarseQuarkN}{141}  % experiments/FIGS/data/v2/index.json :: resonant classes sharing their R16_Q1 group with a decay that differs only in quark flavor (count) :: sha256 c6c4138ee5c86423
\newcommand{\MlstPerClassCoarseQuarkMedian}{1.65}  % experiments/FIGS/data/v2/index.json :: median over those classes of the mean per-class ratio, R16_Q1 over 188 :: sha256 c6c4138ee5c86423
\newcommand{\MlstPerClassCoarseQuarkMedianPr}{65\% more errors}  % experiments/FIGS/data/v2/index.json :: median over those classes of the mean per-class ratio, R16_Q1 over 188 :: sha256 c6c4138ee5c86423 :: prose twin: two significant figures, direction in words
\newcommand{\MlstPerClassCoarseQuarkShare}{87\%}  % experiments/FIGS/data/v2/index.json :: their share of the summed ln ratio over the resonant classes, R16_Q1 :: sha256 c6c4138ee5c86423
\newcommand{\MlstPerClassCoarseQuarkLeptonN}{64}  % experiments/FIGS/data/v2/index.json :: resonant classes counted in both the quark- and the lepton-flavor set, R16_Q1 (count) :: sha256 c6c4138ee5c86423
\newcommand{\MlstPerClassRandFlavMergedOverApart}{1.010\,\ensuremath{\pm}\,0.008}  % experiments/FIGS/data/v2/index.json :: random partitions, linear probe: median per-class ratio of the classes merged with a decay that differs only in quark flavor over that of the classes kept apart from all of them, mean +- SD over the partitions :: sha256 c6c4138ee5c86423
\newcommand{\MlstPerClassRandFlavMergedOverApartPr}{equal within 1\%}  % experiments/FIGS/data/v2/index.json :: random partitions, linear probe: median per-class ratio of the classes merged with a decay that differs only in quark flavor over that of the classes kept apart from all of them, mean +- SD over the partitions :: sha256 c6c4138ee5c86423 :: prose twin: two significant figures, direction in words
\newcommand{\MlstPerClassRandLeptonMergedOverApart}{0.997\,\ensuremath{\pm}\,0.020}  % experiments/FIGS/data/v2/index.json :: random partitions, linear probe: median per-class ratio of the classes merged with a decay that differs only in lepton flavor over that of the classes kept apart from all of them, mean +- SD over the partitions :: sha256 c6c4138ee5c86423
\newcommand{\MlstPerClassRandLeptonMergedOverApartPr}{equal within 3\%}  % experiments/FIGS/data/v2/index.json :: random partitions, linear probe: median per-class ratio of the classes merged with a decay that differs only in lepton flavor over that of the classes kept apart from all of them, mean +- SD over the partitions :: sha256 c6c4138ee5c86423 :: prose twin: two significant figures, direction in words
\newcommand{\MlstMlpNatPerClassRandFlavMergedOverApart}{1.003\,\ensuremath{\pm}\,0.007}  % experiments/FIGS/data/v2/index.json :: random partitions, nonlinear probe, natural composition: median per-class ratio of the classes merged with a decay that differs only in quark flavor over that of the classes kept apart from all of them, mean +- SD over the partitions :: sha256 c6c4138ee5c86423
\newcommand{\MlstMlpNatPerClassRandFlavMergedOverApartPr}{equal within 1\%}  % experiments/FIGS/data/v2/index.json :: random partitions, nonlinear probe, natural composition: median per-class ratio of the classes merged with a decay that differs only in quark flavor over that of the classes kept apart from all of them, mean +- SD over the partitions :: sha256 c6c4138ee5c86423 :: prose twin: two significant figures, direction in words
\newcommand{\MlstMlpNatPerClassRandLeptonMergedOverApart}{1.01\,\ensuremath{\pm}\,0.07}  % experiments/FIGS/data/v2/index.json :: random partitions, nonlinear probe, natural composition: median per-class ratio of the classes merged with a decay that differs only in lepton flavor over that of the classes kept apart from all of them, mean +- SD over the partitions :: sha256 c6c4138ee5c86423
\newcommand{\MlstMlpNatPerClassRandLeptonMergedOverApartPr}{equal within 8\%}  % experiments/FIGS/data/v2/index.json :: random partitions, nonlinear probe, natural composition: median per-class ratio of the classes merged with a decay that differs only in lepton flavor over that of the classes kept apart from all of them, mean +- SD over the partitions :: sha256 c6c4138ee5c86423 :: prose twin: two significant figures, direction in words
\newcommand{\MlstPerClassAucDiffResSeedSdMedian}{0.015}  % experiments/FIGS/data/v2/index.json :: median over classes of the SD over indices of 100 (AUC_c of L162 - AUC_c of 188), pp :: sha256 c6c4138ee5c86423
\newcommand{\MlstPerClassAucDiffResMin}{\ensuremath{-}0.73}  % experiments/FIGS/data/v2/index.json :: most negative mean over indices of 100 (AUC_c of L162 - AUC_c of 188), pp :: sha256 c6c4138ee5c86423 :: label_QCD_bcss
\newcommand{\MlstPerClassAucDiffResMinClass}{QCD, \ensuremath{bcss}}  % experiments/FIGS/data/v2/index.json :: its class, L162 :: sha256 c6c4138ee5c86423
\newcommand{\MlstPerClassAucDiffMidSeedSdMedian}{0.032}  % experiments/FIGS/data/v2/index.json :: median over classes of the SD over indices of 100 (AUC_c of R42_Q1 - AUC_c of 188), pp :: sha256 c6c4138ee5c86423
\newcommand{\MlstPerClassAucDiffMidMin}{\ensuremath{-}1.4}  % experiments/FIGS/data/v2/index.json :: most negative mean over indices of 100 (AUC_c of R42_Q1 - AUC_c of 188), pp :: sha256 c6c4138ee5c86423 :: label_X_YY_cqtauev
\newcommand{\MlstPerClassAucDiffMidMinClass}{\ensuremath{X\to YY\to cq\tau_e\nu}}  % experiments/FIGS/data/v2/index.json :: its class, R42_Q1 :: sha256 c6c4138ee5c86423
\newcommand{\MlstPerClassAucDiffCoarseSeedSdMedian}{0.077}  % experiments/FIGS/data/v2/index.json :: median over classes of the SD over indices of 100 (AUC_c of R16_Q1 - AUC_c of 188), pp :: sha256 c6c4138ee5c86423
\newcommand{\MlstPerClassAucDiffCoarseMin}{\ensuremath{-}3.1}  % experiments/FIGS/data/v2/index.json :: most negative mean over indices of 100 (AUC_c of R16_Q1 - AUC_c of 188), pp :: sha256 c6c4138ee5c86423 :: label_QCD_bcss
\newcommand{\MlstPerClassAucDiffCoarseMinClass}{QCD, \ensuremath{bcss}}  % experiments/FIGS/data/v2/index.json :: its class, R16_Q1 :: sha256 c6c4138ee5c86423
\newcommand{\MlstRocJcTwoHurtList}{\ensuremath{X\to YY\to ccq}; \ensuremath{X\to YY\to ssb}; \ensuremath{X\to YY\to ccss}; QCD, \ensuremath{bcc}; QCD, \ensuremath{bccss}; QCD, \ensuremath{bss}; QCD, \ensuremath{ccs} and QCD, \ensuremath{ccss}}  % experiments/FIGS/data/v2/index.json :: jc2_roc_classes on the class-balanced probes; candidates with at least 1000 test jets :: sha256 c6c4138ee5c86423 :: label_X_YY_ccq, label_X_YY_ssb, label_X_YY_ccss, label_QCD_bcc, label_QCD_bccss, label_QCD_bss, label_QCD_ccs, label_QCD_ccss
\newcommand{\MlstRocJcTwoStableList}{\ensuremath{X\to gg}; \ensuremath{X\to YY\to ggg}; \ensuremath{X\to YY\to g\tau_h\tau_h} and QCD, light}  % experiments/FIGS/data/v2/index.json :: jc2_roc_classes on the class-balanced probes; candidates with at least 1000 test jets :: sha256 c6c4138ee5c86423 :: label_X_gg, label_X_YY_ggg, label_X_YY_gtauhtauh, label_QCD_light
\newcommand{\MlstRocJcTwoHurtMinDrop}{2.1}  % experiments/FIGS/data/v2/index.json :: jc2_roc_classes on the class-balanced probes; candidates with at least 1000 test jets :: sha256 c6c4138ee5c86423 :: smallest AUC loss at 17 classes of the hurt, pp
\newcommand{\MlstRocJcTwoHurtMaxDrop}{2.8}  % experiments/FIGS/data/v2/index.json :: jc2_roc_classes on the class-balanced probes; candidates with at least 1000 test jets :: sha256 c6c4138ee5c86423 :: largest AUC loss at 17 classes of the hurt, pp
\newcommand{\MlstRocJcTwoStableMaxChange}{15}  % experiments/FIGS/data/v2/index.json :: jc2_roc_classes on the class-balanced probes; candidates with at least 1000 test jets :: sha256 c6c4138ee5c86423 :: largest change of 1 - AUC at any label set of the stable classes, per cent of the 188-class value
\newcommand{\MlstJcBalNThin}{4}  % experiments/FIGS/data/v2/jc2bal/census.json :: counts.ft_test :: sha256 a4379ee0817254cf :: classes with fewer than 1000 class-balanced test jets (F3 hatching)
\newcommand{\MlstMinBkgPass}{20}  % experiments/FIGS/mlst_v4.py :: JC1_MIN_BKG :: sha256 fa536b4866ffd651 :: fewest background jets behind a drawn ROC point or a quoted rejection
\newcommand{\MlstMassMinJets}{100}  % experiments/EVAL/atlas_resolution.py :: MIN_JETS :: sha256 b82cc219f79f0f14 :: jets below which a mass bin is null
\newcommand{\MlstRocCutBbFull}{0.79}  % experiments/FIGS/data/v2/index.json :: curves_v2/pairs curve bb_vs_cc: largest cut.eps_s over the L188 runs ('not reached' when a run's cut is null or leaves fewer than 2 grid points) :: sha256 c6c4138ee5c86423
\newcommand{\MlstRocCutBbRes}{0.77}  % experiments/FIGS/data/v2/index.json :: curves_v2/pairs curve bb_vs_cc: largest cut.eps_s over the L162 runs ('not reached' when a run's cut is null or leaves fewer than 2 grid points) :: sha256 c6c4138ee5c86423
\newcommand{\MlstRocCutBbMid}{0.79}  % experiments/FIGS/data/v2/index.json :: curves_v2/pairs curve bb_vs_cc: largest cut.eps_s over the R42_Q1 runs ('not reached' when a run's cut is null or leaves fewer than 2 grid points) :: sha256 c6c4138ee5c86423
\newcommand{\MlstRocCutBbCoarse}{0.13}  % experiments/FIGS/data/v2/index.json :: curves_v2/pairs curve bb_vs_cc: largest cut.eps_s over the R16_Q1 runs ('not reached' when a run's cut is null or leaves fewer than 2 grid points) :: sha256 c6c4138ee5c86423
\newcommand{\MlstRocCutBbFzero}{0.03}  % experiments/FIGS/data/v2/index.json :: curves_v2/pairs curve bb_vs_cc: largest cut.eps_s over the FLAV_F0 runs ('not reached' when a run's cut is null or leaves fewer than 2 grid points) :: sha256 c6c4138ee5c86423
\newcommand{\MlstRocCutBbFone}{0.62}  % experiments/FIGS/data/v2/index.json :: curves_v2/pairs curve bb_vs_cc: largest cut.eps_s over the FLAV_F1 runs ('not reached' when a run's cut is null or leaves fewer than 2 grid points) :: sha256 c6c4138ee5c86423
\newcommand{\MlstRocCutBbFoner}{0.51}  % experiments/FIGS/data/v2/index.json :: curves_v2/pairs curve bb_vs_cc: largest cut.eps_s over the FLAV_F1R runs ('not reached' when a run's cut is null or leaves fewer than 2 grid points) :: sha256 c6c4138ee5c86423
\newcommand{\MlstRocNBkgBb}{200{,}000}  % experiments/FIGS/data/v2/index.json :: curves_v2/pairs curve bb_vs_cc: n_background :: sha256 c6c4138ee5c86423
\newcommand{\MlstRocRejMaxBb}{10{,}000}  % experiments/FIGS/data/v2/index.json :: curves_v2/pairs curve bb_vs_cc: n_background / 20 :: sha256 c6c4138ee5c86423
\newcommand{\MlstRocCutBbqqFull}{0.67}  % experiments/FIGS/data/v2/index.json :: curves_v2/pairs curve bbqq_vs_ccqq: largest cut.eps_s over the L188 runs ('not reached' when a run's cut is null or leaves fewer than 2 grid points) :: sha256 c6c4138ee5c86423
\newcommand{\MlstRocCutBbqqRes}{0.66}  % experiments/FIGS/data/v2/index.json :: curves_v2/pairs curve bbqq_vs_ccqq: largest cut.eps_s over the L162 runs ('not reached' when a run's cut is null or leaves fewer than 2 grid points) :: sha256 c6c4138ee5c86423
\newcommand{\MlstRocCutBbqqMid}{0.62}  % experiments/FIGS/data/v2/index.json :: curves_v2/pairs curve bbqq_vs_ccqq: largest cut.eps_s over the R42_Q1 runs ('not reached' when a run's cut is null or leaves fewer than 2 grid points) :: sha256 c6c4138ee5c86423
\newcommand{\MlstRocCutBbqqCoarse}{0.06}  % experiments/FIGS/data/v2/index.json :: curves_v2/pairs curve bbqq_vs_ccqq: largest cut.eps_s over the R16_Q1 runs ('not reached' when a run's cut is null or leaves fewer than 2 grid points) :: sha256 c6c4138ee5c86423
\newcommand{\MlstRocCutBbqqFzero}{0.06}  % experiments/FIGS/data/v2/index.json :: curves_v2/pairs curve bbqq_vs_ccqq: largest cut.eps_s over the FLAV_F0 runs ('not reached' when a run's cut is null or leaves fewer than 2 grid points) :: sha256 c6c4138ee5c86423
\newcommand{\MlstRocCutBbqqFone}{0.56}  % experiments/FIGS/data/v2/index.json :: curves_v2/pairs curve bbqq_vs_ccqq: largest cut.eps_s over the FLAV_F1 runs ('not reached' when a run's cut is null or leaves fewer than 2 grid points) :: sha256 c6c4138ee5c86423
\newcommand{\MlstRocCutBbqqFoner}{0.32}  % experiments/FIGS/data/v2/index.json :: curves_v2/pairs curve bbqq_vs_ccqq: largest cut.eps_s over the FLAV_F1R runs ('not reached' when a run's cut is null or leaves fewer than 2 grid points) :: sha256 c6c4138ee5c86423
\newcommand{\MlstRocNBkgBbqq}{200{,}000}  % experiments/FIGS/data/v2/index.json :: curves_v2/pairs curve bbqq_vs_ccqq: n_background :: sha256 c6c4138ee5c86423
\newcommand{\MlstRocRejMaxBbqq}{10{,}000}  % experiments/FIGS/data/v2/index.json :: curves_v2/pairs curve bbqq_vs_ccqq: n_background / 20 :: sha256 c6c4138ee5c86423
\newcommand{\MlstRocCutBcFull}{0.24}  % experiments/FIGS/data/v2/index.json :: curves_v2/pairs curve bc_vs_bq_cs: largest cut.eps_s over the L188 runs ('not reached' when a run's cut is null or leaves fewer than 2 grid points) :: sha256 c6c4138ee5c86423
\newcommand{\MlstRocCutBcRes}{0.25}  % experiments/FIGS/data/v2/index.json :: curves_v2/pairs curve bc_vs_bq_cs: largest cut.eps_s over the L162 runs ('not reached' when a run's cut is null or leaves fewer than 2 grid points) :: sha256 c6c4138ee5c86423
\newcommand{\MlstRocCutBcMid}{0.14}  % experiments/FIGS/data/v2/index.json :: curves_v2/pairs curve bc_vs_bq_cs: largest cut.eps_s over the R42_Q1 runs ('not reached' when a run's cut is null or leaves fewer than 2 grid points) :: sha256 c6c4138ee5c86423
\newcommand{\MlstRocCutBcCoarse}{0.00}  % experiments/FIGS/data/v2/index.json :: curves_v2/pairs curve bc_vs_bq_cs: largest cut.eps_s over the R16_Q1 runs ('not reached' when a run's cut is null or leaves fewer than 2 grid points) :: sha256 c6c4138ee5c86423
\newcommand{\MlstRocCutBcFzero}{0.00}  % experiments/FIGS/data/v2/index.json :: curves_v2/pairs curve bc_vs_bq_cs: largest cut.eps_s over the FLAV_F0 runs ('not reached' when a run's cut is null or leaves fewer than 2 grid points) :: sha256 c6c4138ee5c86423
\newcommand{\MlstRocCutBcFone}{0.01}  % experiments/FIGS/data/v2/index.json :: curves_v2/pairs curve bc_vs_bq_cs: largest cut.eps_s over the FLAV_F1 runs ('not reached' when a run's cut is null or leaves fewer than 2 grid points) :: sha256 c6c4138ee5c86423
\newcommand{\MlstRocCutBcFoner}{0.01}  % experiments/FIGS/data/v2/index.json :: curves_v2/pairs curve bc_vs_bq_cs: largest cut.eps_s over the FLAV_F1R runs ('not reached' when a run's cut is null or leaves fewer than 2 grid points) :: sha256 c6c4138ee5c86423
\newcommand{\MlstRocNBkgBc}{200{,}000}  % experiments/FIGS/data/v2/index.json :: curves_v2/pairs curve bc_vs_bq_cs: n_background :: sha256 c6c4138ee5c86423
\newcommand{\MlstRocRejMaxBc}{10{,}000}  % experiments/FIGS/data/v2/index.json :: curves_v2/pairs curve bc_vs_bq_cs: n_background / 20 :: sha256 c6c4138ee5c86423
\newcommand{\MlstRocCutEemmFull}{not reached}  % experiments/FIGS/data/v2/index.json :: curves_v2/pairs curve ee_vs_mm: largest cut.eps_s over the L188 runs ('not reached' when a run's cut is null or leaves fewer than 2 grid points) :: sha256 c6c4138ee5c86423
\newcommand{\MlstRocCutEemmRes}{not reached}  % experiments/FIGS/data/v2/index.json :: curves_v2/pairs curve ee_vs_mm: largest cut.eps_s over the L162 runs ('not reached' when a run's cut is null or leaves fewer than 2 grid points) :: sha256 c6c4138ee5c86423
\newcommand{\MlstRocCutEemmMid}{0.00}  % experiments/FIGS/data/v2/index.json :: curves_v2/pairs curve ee_vs_mm: largest cut.eps_s over the R42_Q1 runs ('not reached' when a run's cut is null or leaves fewer than 2 grid points) :: sha256 c6c4138ee5c86423
\newcommand{\MlstRocCutEemmCoarse}{0.00}  % experiments/FIGS/data/v2/index.json :: curves_v2/pairs curve ee_vs_mm: largest cut.eps_s over the R16_Q1 runs ('not reached' when a run's cut is null or leaves fewer than 2 grid points) :: sha256 c6c4138ee5c86423
\newcommand{\MlstRocCutEemmFzero}{not reached}  % experiments/FIGS/data/v2/index.json :: curves_v2/pairs curve ee_vs_mm: largest cut.eps_s over the FLAV_F0 runs ('not reached' when a run's cut is null or leaves fewer than 2 grid points) :: sha256 c6c4138ee5c86423
\newcommand{\MlstRocCutEemmFone}{not reached}  % experiments/FIGS/data/v2/index.json :: curves_v2/pairs curve ee_vs_mm: largest cut.eps_s over the FLAV_F1 runs ('not reached' when a run's cut is null or leaves fewer than 2 grid points) :: sha256 c6c4138ee5c86423
\newcommand{\MlstRocCutEemmFoner}{not reached}  % experiments/FIGS/data/v2/index.json :: curves_v2/pairs curve ee_vs_mm: largest cut.eps_s over the FLAV_F1R runs ('not reached' when a run's cut is null or leaves fewer than 2 grid points) :: sha256 c6c4138ee5c86423
\newcommand{\MlstRocNBkgEemm}{128{,}833}  % experiments/FIGS/data/v2/index.json :: curves_v2/pairs curve ee_vs_mm: n_background :: sha256 c6c4138ee5c86423
\newcommand{\MlstRocRejMaxEemm}{6{,}442}  % experiments/FIGS/data/v2/index.json :: curves_v2/pairs curve ee_vs_mm: n_background / 20 :: sha256 c6c4138ee5c86423
\newcommand{\MlstRocNoCurveEemm}{the \Kfull- and \Kres-class models and the three random partitions with controlled flavor boundaries}  % experiments/FIGS/data/v2/index.json :: curves_v2/pairs curve ee_vs_mm: configurations of F4 with a run whose cut is null or leaves fewer than 2 grid points (F4 draws no curve for them) :: sha256 c6c4138ee5c86423 :: L188, L162, FLAV_F0, FLAV_F1, FLAV_F1R
\newcommand{\MlstRocCutJcOneHbbFull}{0.40}  % experiments/FIGS/data/v2/index.json :: curves_v2/jc1_ft curve label_Hbb: largest cut.eps_s over the 188-class runs ('not reached' when a run's cut is null or leaves fewer than 2 grid points) :: sha256 c6c4138ee5c86423
\newcommand{\MlstRocCutJcOneHccFull}{0.37}  % experiments/FIGS/data/v2/index.json :: curves_v2/jc1_ft curve label_Hcc: largest cut.eps_s over the 188-class runs ('not reached' when a run's cut is null or leaves fewer than 2 grid points) :: sha256 c6c4138ee5c86423
\newcommand{\MlstRocCutJcOneHggFull}{0.04}  % experiments/FIGS/data/v2/index.json :: curves_v2/jc1_ft curve label_Hgg: largest cut.eps_s over the 188-class runs ('not reached' when a run's cut is null or leaves fewer than 2 grid points) :: sha256 c6c4138ee5c86423
\newcommand{\MlstRocCutJcOneHfourqFull}{0.26}  % experiments/FIGS/data/v2/index.json :: curves_v2/jc1_ft curve label_H4q: largest cut.eps_s over the 188-class runs ('not reached' when a run's cut is null or leaves fewer than 2 grid points) :: sha256 c6c4138ee5c86423
\newcommand{\MlstRocCutJcOneHqqlFull}{0.98}  % experiments/FIGS/data/v2/index.json :: curves_v2/jc1_ft curve label_Hqql: largest cut.eps_s over the 188-class runs ('not reached' when a run's cut is null or leaves fewer than 2 grid points) :: sha256 c6c4138ee5c86423
\newcommand{\MlstRocCutJcOneTbqqFull}{0.53}  % experiments/FIGS/data/v2/index.json :: curves_v2/jc1_ft curve label_Tbqq: largest cut.eps_s over the 188-class runs ('not reached' when a run's cut is null or leaves fewer than 2 grid points) :: sha256 c6c4138ee5c86423
\newcommand{\MlstRocCutJcOneTblFull}{not reached}  % experiments/FIGS/data/v2/index.json :: curves_v2/jc1_ft curve label_Tbl: largest cut.eps_s over the 188-class runs ('not reached' when a run's cut is null or leaves fewer than 2 grid points) :: sha256 c6c4138ee5c86423
\newcommand{\MlstRocCutJcOneWqqFull}{0.15}  % experiments/FIGS/data/v2/index.json :: curves_v2/jc1_ft curve label_Wqq: largest cut.eps_s over the 188-class runs ('not reached' when a run's cut is null or leaves fewer than 2 grid points) :: sha256 c6c4138ee5c86423
\newcommand{\MlstRocCutJcOneZqqFull}{0.12}  % experiments/FIGS/data/v2/index.json :: curves_v2/jc1_ft curve label_Zqq: largest cut.eps_s over the 188-class runs ('not reached' when a run's cut is null or leaves fewer than 2 grid points) :: sha256 c6c4138ee5c86423
\newcommand{\MlstFtTopRfiftyCoarseOverFull}{0.932\,\ensuremath{\pm}\,0.018}  % experiments/FIGS/data/v2/index.json :: curves_v2/top_ft curve top: rejection at eps_S = 0.5 of R16_Q1 over L188, run by run :: sha256 c6c4138ee5c86423
\newcommand{\MlstFtTopRfiftyCoarseOverFullPr}{6.8\% less rejection}  % experiments/FIGS/data/v2/index.json :: curves_v2/top_ft curve top: rejection at eps_S = 0.5 of R16_Q1 over L188, run by run :: sha256 c6c4138ee5c86423 :: prose twin: two significant figures, direction in words
\newcommand{\MlstProbeUnconverged}{0}  % experiments/FIGS/data/v2/index.json :: linear probes of pretrained models that did not converge (count) :: sha256 c6c4138ee5c86423
\newcommand{\MlstProbeUnconvergedInit}{0}  % experiments/FIGS/data/v2/index.json :: linear probes of the randomly initialized networks that did not converge (count) :: sha256 c6c4138ee5c86423
\newcommand{\MlstProbeFits}{895}  % experiments/FIGS/data/v2/index.json :: selected linear-probe fits of the reported checkpoints (count) :: sha256 c6c4138ee5c86423
\newcommand{\MlstProbeEdgeC}{16}  % experiments/FIGS/data/v2/index.json :: of these, fits whose selected C is the largest of the grid (count) :: sha256 c6c4138ee5c86423
\newcommand{\MlstMlpUnconverged}{0}  % experiments/FIGS/data/v2/index.json :: nonlinear probes that did not reach their plateau within the epoch cap (count) :: sha256 c6c4138ee5c86423
\newcommand{\MlstMlpCmpJcTwoResOverFull}{1.019\,\ensuremath{\pm}\,0.007}  % experiments/FIGS/data/v2/index.json :: nonlinear probe jc2_188: (1 - AUC) of L162 over L188, run by run :: sha256 c6c4138ee5c86423
\newcommand{\MlstMlpCmpJcTwoResOverFullPr}{1.9\% more errors}  % experiments/FIGS/data/v2/index.json :: nonlinear probe jc2_188: (1 - AUC) of L162 over L188, run by run :: sha256 c6c4138ee5c86423 :: prose twin: two significant figures, direction in words
\newcommand{\MlstMlpCmpJcTwoMidOverFull}{1.170\,\ensuremath{\pm}\,0.006}  % experiments/FIGS/data/v2/index.json :: nonlinear probe jc2_188: (1 - AUC) of R42_Q1 over L188, run by run :: sha256 c6c4138ee5c86423
\newcommand{\MlstMlpCmpJcTwoMidOverFullPr}{17\% more errors}  % experiments/FIGS/data/v2/index.json :: nonlinear probe jc2_188: (1 - AUC) of R42_Q1 over L188, run by run :: sha256 c6c4138ee5c86423 :: prose twin: two significant figures, direction in words
\newcommand{\MlstMlpCmpJcTwoCoarseOverFull}{1.388\,\ensuremath{\pm}\,0.021}  % experiments/FIGS/data/v2/index.json :: nonlinear probe jc2_188: (1 - AUC) of R16_Q1 over L188, run by run :: sha256 c6c4138ee5c86423
\newcommand{\MlstMlpCmpJcTwoCoarseOverFullPr}{39\% more errors}  % experiments/FIGS/data/v2/index.json :: nonlinear probe jc2_188: (1 - AUC) of R16_Q1 over L188, run by run :: sha256 c6c4138ee5c86423 :: prose twin: two significant figures, direction in words
\newcommand{\MlstMlpCmpJcTwoResMassOverRes}{\pending{no run yet}}  % experiments/FIGS/data/v2/index.json :: nonlinear probe jc2_188: (1 - AUC) of L162_MASS over L162, run by run :: sha256 c6c4138ee5c86423
\newcommand{\MlstMlpCmpJcTwoResMassOverResPr}{\pending{no run yet}}  % experiments/FIGS/data/v2/index.json :: nonlinear probe jc2_188: (1 - AUC) of L162_MASS over L162, run by run :: sha256 c6c4138ee5c86423 :: prose twin: two significant figures, direction in words
\newcommand{\MlstMlpCmpJcTwoResMassOverResPrMax}{\pending{no run yet}}  % experiments/FIGS/data/v2/index.json :: nonlinear probe jc2_188: (1 - AUC) of L162_MASS over L162, run by run :: sha256 c6c4138ee5c86423 :: prose twin of the run farther from 1
\newcommand{\MlstMlpCmpJcTwoCoarseMassOverCoarse}{\pending{no run yet}}  % experiments/FIGS/data/v2/index.json :: nonlinear probe jc2_188: (1 - AUC) of R16_Q1_MASS over R16_Q1, run by run :: sha256 c6c4138ee5c86423
\newcommand{\MlstMlpCmpJcTwoCoarseMassOverCoarsePr}{\pending{no run yet}}  % experiments/FIGS/data/v2/index.json :: nonlinear probe jc2_188: (1 - AUC) of R16_Q1_MASS over R16_Q1, run by run :: sha256 c6c4138ee5c86423 :: prose twin: two significant figures, direction in words
\newcommand{\MlstMlpCmpJcTwoCoarseMassOverCoarsePrMax}{\pending{no run yet}}  % experiments/FIGS/data/v2/index.json :: nonlinear probe jc2_188: (1 - AUC) of R16_Q1_MASS over R16_Q1, run by run :: sha256 c6c4138ee5c86423 :: prose twin of the run farther from 1
\newcommand{\MlstMlpCmpJcTwoFzeroOverFull}{1.294 and 1.310}  % experiments/FIGS/data/v2/index.json :: nonlinear probe jc2_188: (1 - AUC) of FLAV_F0 over L188, run by run :: sha256 c6c4138ee5c86423
\newcommand{\MlstMlpCmpJcTwoFzeroOverFullPr}{at least 29\% more errors}  % experiments/FIGS/data/v2/index.json :: nonlinear probe jc2_188: (1 - AUC) of FLAV_F0 over L188, run by run :: sha256 c6c4138ee5c86423 :: prose twin: two significant figures, direction in words
\newcommand{\MlstMlpCmpJcTwoFzeroOverFullPrMax}{up to 31\% more errors}  % experiments/FIGS/data/v2/index.json :: nonlinear probe jc2_188: (1 - AUC) of FLAV_F0 over L188, run by run :: sha256 c6c4138ee5c86423 :: prose twin of the run farther from 1
\newcommand{\MlstMlpCmpJcTwoFoneOverFzero}{0.903 and 0.885}  % experiments/FIGS/data/v2/index.json :: nonlinear probe jc2_188: (1 - AUC) of FLAV_F1 over FLAV_F0, run by run :: sha256 c6c4138ee5c86423
\newcommand{\MlstMlpCmpJcTwoFoneOverFzeroPr}{at least 9.7\% fewer errors}  % experiments/FIGS/data/v2/index.json :: nonlinear probe jc2_188: (1 - AUC) of FLAV_F1 over FLAV_F0, run by run :: sha256 c6c4138ee5c86423 :: prose twin: two significant figures, direction in words
\newcommand{\MlstMlpCmpJcTwoFoneOverFzeroPrMax}{up to 12\% fewer errors}  % experiments/FIGS/data/v2/index.json :: nonlinear probe jc2_188: (1 - AUC) of FLAV_F1 over FLAV_F0, run by run :: sha256 c6c4138ee5c86423 :: prose twin of the run farther from 1
\newcommand{\MlstMlpCmpJcTwoFonerOverFone}{0.9916 and 0.9918}  % experiments/FIGS/data/v2/index.json :: nonlinear probe jc2_188: (1 - AUC) of FLAV_F1R over FLAV_F1, run by run :: sha256 c6c4138ee5c86423
\newcommand{\MlstMlpCmpJcTwoFonerOverFonePr}{equal within 1\%}  % experiments/FIGS/data/v2/index.json :: nonlinear probe jc2_188: (1 - AUC) of FLAV_F1R over FLAV_F1, run by run :: sha256 c6c4138ee5c86423 :: prose twin: two significant figures, direction in words
\newcommand{\MlstMlpCmpJcTwoFonerOverFonePrMax}{equal within 1\%}  % experiments/FIGS/data/v2/index.json :: nonlinear probe jc2_188: (1 - AUC) of FLAV_F1R over FLAV_F1, run by run :: sha256 c6c4138ee5c86423 :: prose twin of the run farther from 1
\newcommand{\MlstMlpCmpJcTwoFonerOverFzero}{0.895 and 0.877}  % experiments/FIGS/data/v2/index.json :: nonlinear probe jc2_188: (1 - AUC) of FLAV_F1R over FLAV_F0, run by run :: sha256 c6c4138ee5c86423
\newcommand{\MlstMlpCmpJcTwoFonerOverFzeroPr}{at least 10\% fewer errors}  % experiments/FIGS/data/v2/index.json :: nonlinear probe jc2_188: (1 - AUC) of FLAV_F1R over FLAV_F0, run by run :: sha256 c6c4138ee5c86423 :: prose twin: two significant figures, direction in words
\newcommand{\MlstMlpCmpJcTwoFonerOverFzeroPrMax}{up to 12\% fewer errors}  % experiments/FIGS/data/v2/index.json :: nonlinear probe jc2_188: (1 - AUC) of FLAV_F1R over FLAV_F0, run by run :: sha256 c6c4138ee5c86423 :: prose twin of the run farther from 1
\newcommand{\MlstMlpCmpJcTwoFzeroOverCoarse}{0.9461 and 0.9456}  % experiments/FIGS/data/v2/index.json :: nonlinear probe jc2_188: (1 - AUC) of FLAV_F0 over R16_Q1, run by run :: sha256 c6c4138ee5c86423
\newcommand{\MlstMlpCmpJcTwoFzeroOverCoarsePr}{at least 5.4\% fewer errors}  % experiments/FIGS/data/v2/index.json :: nonlinear probe jc2_188: (1 - AUC) of FLAV_F0 over R16_Q1, run by run :: sha256 c6c4138ee5c86423 :: prose twin: two significant figures, direction in words
\newcommand{\MlstMlpCmpJcTwoFzeroOverCoarsePrMax}{up to 5.4\% fewer errors}  % experiments/FIGS/data/v2/index.json :: nonlinear probe jc2_188: (1 - AUC) of FLAV_F0 over R16_Q1, run by run :: sha256 c6c4138ee5c86423 :: prose twin of the run farther from 1
\newcommand{\MlstMlpCmpJcTwoRandOverFull}{\pending{no run yet}}  % experiments/FIGS/data/v2/index.json :: nonlinear probe jc2_188: (1 - AUC) of each random partition over the 188-class run of its seed :: sha256 c6c4138ee5c86423
\newcommand{\MlstMlpCmpJcTwoRandOverFullPr}{\pending{no run yet}}  % experiments/FIGS/data/v2/index.json :: nonlinear probe jc2_188: (1 - AUC) of each random partition over the 188-class run of its seed :: sha256 c6c4138ee5c86423 :: prose twin: two significant figures, direction in words
\newcommand{\MlstMlpCmpJcTwoResResOverFull}{1.001\,\ensuremath{\pm}\,0.007}  % experiments/FIGS/data/v2/index.json :: nonlinear probe jc2_188_res: (1 - AUC) of L162 over L188, run by run :: sha256 c6c4138ee5c86423
\newcommand{\MlstMlpCmpJcTwoResResOverFullPr}{equal within 1\%}  % experiments/FIGS/data/v2/index.json :: nonlinear probe jc2_188_res: (1 - AUC) of L162 over L188, run by run :: sha256 c6c4138ee5c86423 :: prose twin: two significant figures, direction in words
\newcommand{\MlstMlpCmpJcTwoResMidOverFull}{1.192\,\ensuremath{\pm}\,0.008}  % experiments/FIGS/data/v2/index.json :: nonlinear probe jc2_188_res: (1 - AUC) of R42_Q1 over L188, run by run :: sha256 c6c4138ee5c86423
\newcommand{\MlstMlpCmpJcTwoResMidOverFullPr}{19\% more errors}  % experiments/FIGS/data/v2/index.json :: nonlinear probe jc2_188_res: (1 - AUC) of R42_Q1 over L188, run by run :: sha256 c6c4138ee5c86423 :: prose twin: two significant figures, direction in words
\newcommand{\MlstMlpCmpJcTwoResCoarseOverFull}{1.408\,\ensuremath{\pm}\,0.026}  % experiments/FIGS/data/v2/index.json :: nonlinear probe jc2_188_res: (1 - AUC) of R16_Q1 over L188, run by run :: sha256 c6c4138ee5c86423
\newcommand{\MlstMlpCmpJcTwoResCoarseOverFullPr}{41\% more errors}  % experiments/FIGS/data/v2/index.json :: nonlinear probe jc2_188_res: (1 - AUC) of R16_Q1 over L188, run by run :: sha256 c6c4138ee5c86423 :: prose twin: two significant figures, direction in words
\newcommand{\MlstMlpCmpJcTwoResResMassOverRes}{\pending{no run yet}}  % experiments/FIGS/data/v2/index.json :: nonlinear probe jc2_188_res: (1 - AUC) of L162_MASS over L162, run by run :: sha256 c6c4138ee5c86423
\newcommand{\MlstMlpCmpJcTwoResResMassOverResPr}{\pending{no run yet}}  % experiments/FIGS/data/v2/index.json :: nonlinear probe jc2_188_res: (1 - AUC) of L162_MASS over L162, run by run :: sha256 c6c4138ee5c86423 :: prose twin: two significant figures, direction in words
\newcommand{\MlstMlpCmpJcTwoResResMassOverResPrMax}{\pending{no run yet}}  % experiments/FIGS/data/v2/index.json :: nonlinear probe jc2_188_res: (1 - AUC) of L162_MASS over L162, run by run :: sha256 c6c4138ee5c86423 :: prose twin of the run farther from 1
\newcommand{\MlstMlpCmpJcTwoResCoarseMassOverCoarse}{\pending{no run yet}}  % experiments/FIGS/data/v2/index.json :: nonlinear probe jc2_188_res: (1 - AUC) of R16_Q1_MASS over R16_Q1, run by run :: sha256 c6c4138ee5c86423
\newcommand{\MlstMlpCmpJcTwoResCoarseMassOverCoarsePr}{\pending{no run yet}}  % experiments/FIGS/data/v2/index.json :: nonlinear probe jc2_188_res: (1 - AUC) of R16_Q1_MASS over R16_Q1, run by run :: sha256 c6c4138ee5c86423 :: prose twin: two significant figures, direction in words
\newcommand{\MlstMlpCmpJcTwoResCoarseMassOverCoarsePrMax}{\pending{no run yet}}  % experiments/FIGS/data/v2/index.json :: nonlinear probe jc2_188_res: (1 - AUC) of R16_Q1_MASS over R16_Q1, run by run :: sha256 c6c4138ee5c86423 :: prose twin of the run farther from 1
\newcommand{\MlstMlpCmpJcTwoResFzeroOverFull}{1.285 and 1.300}  % experiments/FIGS/data/v2/index.json :: nonlinear probe jc2_188_res: (1 - AUC) of FLAV_F0 over L188, run by run :: sha256 c6c4138ee5c86423
\newcommand{\MlstMlpCmpJcTwoResFzeroOverFullPr}{at least 29\% more errors}  % experiments/FIGS/data/v2/index.json :: nonlinear probe jc2_188_res: (1 - AUC) of FLAV_F0 over L188, run by run :: sha256 c6c4138ee5c86423 :: prose twin: two significant figures, direction in words
\newcommand{\MlstMlpCmpJcTwoResFzeroOverFullPrMax}{up to 30\% more errors}  % experiments/FIGS/data/v2/index.json :: nonlinear probe jc2_188_res: (1 - AUC) of FLAV_F0 over L188, run by run :: sha256 c6c4138ee5c86423 :: prose twin of the run farther from 1
\newcommand{\MlstMlpCmpJcTwoResFoneOverFzero}{0.909 and 0.889}  % experiments/FIGS/data/v2/index.json :: nonlinear probe jc2_188_res: (1 - AUC) of FLAV_F1 over FLAV_F0, run by run :: sha256 c6c4138ee5c86423
\newcommand{\MlstMlpCmpJcTwoResFoneOverFzeroPr}{at least 9.1\% fewer errors}  % experiments/FIGS/data/v2/index.json :: nonlinear probe jc2_188_res: (1 - AUC) of FLAV_F1 over FLAV_F0, run by run :: sha256 c6c4138ee5c86423 :: prose twin: two significant figures, direction in words
\newcommand{\MlstMlpCmpJcTwoResFoneOverFzeroPrMax}{up to 11\% fewer errors}  % experiments/FIGS/data/v2/index.json :: nonlinear probe jc2_188_res: (1 - AUC) of FLAV_F1 over FLAV_F0, run by run :: sha256 c6c4138ee5c86423 :: prose twin of the run farther from 1
\newcommand{\MlstMlpCmpJcTwoResFonerOverFone}{0.9793 and 0.9825}  % experiments/FIGS/data/v2/index.json :: nonlinear probe jc2_188_res: (1 - AUC) of FLAV_F1R over FLAV_F1, run by run :: sha256 c6c4138ee5c86423
\newcommand{\MlstMlpCmpJcTwoResFonerOverFonePr}{equal within 3\%}  % experiments/FIGS/data/v2/index.json :: nonlinear probe jc2_188_res: (1 - AUC) of FLAV_F1R over FLAV_F1, run by run :: sha256 c6c4138ee5c86423 :: prose twin: two significant figures, direction in words
\newcommand{\MlstMlpCmpJcTwoResFonerOverFonePrMax}{equal within 3\%}  % experiments/FIGS/data/v2/index.json :: nonlinear probe jc2_188_res: (1 - AUC) of FLAV_F1R over FLAV_F1, run by run :: sha256 c6c4138ee5c86423 :: prose twin of the run farther from 1
\newcommand{\MlstMlpCmpJcTwoResFonerOverFzero}{0.890 and 0.874}  % experiments/FIGS/data/v2/index.json :: nonlinear probe jc2_188_res: (1 - AUC) of FLAV_F1R over FLAV_F0, run by run :: sha256 c6c4138ee5c86423
\newcommand{\MlstMlpCmpJcTwoResFonerOverFzeroPr}{at least 11\% fewer errors}  % experiments/FIGS/data/v2/index.json :: nonlinear probe jc2_188_res: (1 - AUC) of FLAV_F1R over FLAV_F0, run by run :: sha256 c6c4138ee5c86423 :: prose twin: two significant figures, direction in words
\newcommand{\MlstMlpCmpJcTwoResFonerOverFzeroPrMax}{up to 13\% fewer errors}  % experiments/FIGS/data/v2/index.json :: nonlinear probe jc2_188_res: (1 - AUC) of FLAV_F1R over FLAV_F0, run by run :: sha256 c6c4138ee5c86423 :: prose twin of the run farther from 1
\newcommand{\MlstMlpCmpJcTwoResFzeroOverCoarse}{0.9299 and 0.9233}  % experiments/FIGS/data/v2/index.json :: nonlinear probe jc2_188_res: (1 - AUC) of FLAV_F0 over R16_Q1, run by run :: sha256 c6c4138ee5c86423
\newcommand{\MlstMlpCmpJcTwoResFzeroOverCoarsePr}{at least 7.0\% fewer errors}  % experiments/FIGS/data/v2/index.json :: nonlinear probe jc2_188_res: (1 - AUC) of FLAV_F0 over R16_Q1, run by run :: sha256 c6c4138ee5c86423 :: prose twin: two significant figures, direction in words
\newcommand{\MlstMlpCmpJcTwoResFzeroOverCoarsePrMax}{up to 7.7\% fewer errors}  % experiments/FIGS/data/v2/index.json :: nonlinear probe jc2_188_res: (1 - AUC) of FLAV_F0 over R16_Q1, run by run :: sha256 c6c4138ee5c86423 :: prose twin of the run farther from 1
\newcommand{\MlstMlpCmpJcTwoResRandOverFull}{\pending{no run yet}}  % experiments/FIGS/data/v2/index.json :: nonlinear probe jc2_188_res: (1 - AUC) of each random partition over the 188-class run of its seed :: sha256 c6c4138ee5c86423
\newcommand{\MlstMlpCmpJcTwoResRandOverFullPr}{\pending{no run yet}}  % experiments/FIGS/data/v2/index.json :: nonlinear probe jc2_188_res: (1 - AUC) of each random partition over the 188-class run of its seed :: sha256 c6c4138ee5c86423 :: prose twin: two significant figures, direction in words
\newcommand{\MlstMlpCmpJcTwoQcdResOverFull}{1.070\,\ensuremath{\pm}\,0.006}  % experiments/FIGS/data/v2/index.json :: nonlinear probe jc2_188_qcd: (1 - AUC) of L162 over L188, run by run :: sha256 c6c4138ee5c86423
\newcommand{\MlstMlpCmpJcTwoQcdResOverFullPr}{7.0\% more errors}  % experiments/FIGS/data/v2/index.json :: nonlinear probe jc2_188_qcd: (1 - AUC) of L162 over L188, run by run :: sha256 c6c4138ee5c86423 :: prose twin: two significant figures, direction in words
\newcommand{\MlstMlpCmpJcTwoQcdMidOverFull}{1.107\,\ensuremath{\pm}\,0.002}  % experiments/FIGS/data/v2/index.json :: nonlinear probe jc2_188_qcd: (1 - AUC) of R42_Q1 over L188, run by run :: sha256 c6c4138ee5c86423
\newcommand{\MlstMlpCmpJcTwoQcdMidOverFullPr}{11\% more errors}  % experiments/FIGS/data/v2/index.json :: nonlinear probe jc2_188_qcd: (1 - AUC) of R42_Q1 over L188, run by run :: sha256 c6c4138ee5c86423 :: prose twin: two significant figures, direction in words
\newcommand{\MlstMlpCmpJcTwoQcdCoarseOverFull}{1.333\,\ensuremath{\pm}\,0.012}  % experiments/FIGS/data/v2/index.json :: nonlinear probe jc2_188_qcd: (1 - AUC) of R16_Q1 over L188, run by run :: sha256 c6c4138ee5c86423
\newcommand{\MlstMlpCmpJcTwoQcdCoarseOverFullPr}{33\% more errors}  % experiments/FIGS/data/v2/index.json :: nonlinear probe jc2_188_qcd: (1 - AUC) of R16_Q1 over L188, run by run :: sha256 c6c4138ee5c86423 :: prose twin: two significant figures, direction in words
\newcommand{\MlstMlpCmpJcTwoQcdResMassOverRes}{\pending{no run yet}}  % experiments/FIGS/data/v2/index.json :: nonlinear probe jc2_188_qcd: (1 - AUC) of L162_MASS over L162, run by run :: sha256 c6c4138ee5c86423
\newcommand{\MlstMlpCmpJcTwoQcdResMassOverResPr}{\pending{no run yet}}  % experiments/FIGS/data/v2/index.json :: nonlinear probe jc2_188_qcd: (1 - AUC) of L162_MASS over L162, run by run :: sha256 c6c4138ee5c86423 :: prose twin: two significant figures, direction in words
\newcommand{\MlstMlpCmpJcTwoQcdResMassOverResPrMax}{\pending{no run yet}}  % experiments/FIGS/data/v2/index.json :: nonlinear probe jc2_188_qcd: (1 - AUC) of L162_MASS over L162, run by run :: sha256 c6c4138ee5c86423 :: prose twin of the run farther from 1
\newcommand{\MlstMlpCmpJcTwoQcdCoarseMassOverCoarse}{\pending{no run yet}}  % experiments/FIGS/data/v2/index.json :: nonlinear probe jc2_188_qcd: (1 - AUC) of R16_Q1_MASS over R16_Q1, run by run :: sha256 c6c4138ee5c86423
\newcommand{\MlstMlpCmpJcTwoQcdCoarseMassOverCoarsePr}{\pending{no run yet}}  % experiments/FIGS/data/v2/index.json :: nonlinear probe jc2_188_qcd: (1 - AUC) of R16_Q1_MASS over R16_Q1, run by run :: sha256 c6c4138ee5c86423 :: prose twin: two significant figures, direction in words
\newcommand{\MlstMlpCmpJcTwoQcdCoarseMassOverCoarsePrMax}{\pending{no run yet}}  % experiments/FIGS/data/v2/index.json :: nonlinear probe jc2_188_qcd: (1 - AUC) of R16_Q1_MASS over R16_Q1, run by run :: sha256 c6c4138ee5c86423 :: prose twin of the run farther from 1
\newcommand{\MlstMlpCmpJcTwoQcdFzeroOverFull}{1.319 and 1.339}  % experiments/FIGS/data/v2/index.json :: nonlinear probe jc2_188_qcd: (1 - AUC) of FLAV_F0 over L188, run by run :: sha256 c6c4138ee5c86423
\newcommand{\MlstMlpCmpJcTwoQcdFzeroOverFullPr}{at least 32\% more errors}  % experiments/FIGS/data/v2/index.json :: nonlinear probe jc2_188_qcd: (1 - AUC) of FLAV_F0 over L188, run by run :: sha256 c6c4138ee5c86423 :: prose twin: two significant figures, direction in words
\newcommand{\MlstMlpCmpJcTwoQcdFzeroOverFullPrMax}{up to 34\% more errors}  % experiments/FIGS/data/v2/index.json :: nonlinear probe jc2_188_qcd: (1 - AUC) of FLAV_F0 over L188, run by run :: sha256 c6c4138ee5c86423 :: prose twin of the run farther from 1
\newcommand{\MlstMlpCmpJcTwoQcdFoneOverFzero}{0.886 and 0.872}  % experiments/FIGS/data/v2/index.json :: nonlinear probe jc2_188_qcd: (1 - AUC) of FLAV_F1 over FLAV_F0, run by run :: sha256 c6c4138ee5c86423
\newcommand{\MlstMlpCmpJcTwoQcdFoneOverFzeroPr}{at least 11\% fewer errors}  % experiments/FIGS/data/v2/index.json :: nonlinear probe jc2_188_qcd: (1 - AUC) of FLAV_F1 over FLAV_F0, run by run :: sha256 c6c4138ee5c86423 :: prose twin: two significant figures, direction in words
\newcommand{\MlstMlpCmpJcTwoQcdFoneOverFzeroPrMax}{up to 13\% fewer errors}  % experiments/FIGS/data/v2/index.json :: nonlinear probe jc2_188_qcd: (1 - AUC) of FLAV_F1 over FLAV_F0, run by run :: sha256 c6c4138ee5c86423 :: prose twin of the run farther from 1
\newcommand{\MlstMlpCmpJcTwoQcdFonerOverFone}{1.026 and 1.018}  % experiments/FIGS/data/v2/index.json :: nonlinear probe jc2_188_qcd: (1 - AUC) of FLAV_F1R over FLAV_F1, run by run :: sha256 c6c4138ee5c86423
\newcommand{\MlstMlpCmpJcTwoQcdFonerOverFonePr}{at least 1.8\% more errors}  % experiments/FIGS/data/v2/index.json :: nonlinear probe jc2_188_qcd: (1 - AUC) of FLAV_F1R over FLAV_F1, run by run :: sha256 c6c4138ee5c86423 :: prose twin: two significant figures, direction in words
\newcommand{\MlstMlpCmpJcTwoQcdFonerOverFonePrMax}{up to 2.6\% more errors}  % experiments/FIGS/data/v2/index.json :: nonlinear probe jc2_188_qcd: (1 - AUC) of FLAV_F1R over FLAV_F1, run by run :: sha256 c6c4138ee5c86423 :: prose twin of the run farther from 1
\newcommand{\MlstMlpCmpJcTwoQcdFonerOverFzero}{0.909 and 0.887}  % experiments/FIGS/data/v2/index.json :: nonlinear probe jc2_188_qcd: (1 - AUC) of FLAV_F1R over FLAV_F0, run by run :: sha256 c6c4138ee5c86423
\newcommand{\MlstMlpCmpJcTwoQcdFonerOverFzeroPr}{at least 9.1\% fewer errors}  % experiments/FIGS/data/v2/index.json :: nonlinear probe jc2_188_qcd: (1 - AUC) of FLAV_F1R over FLAV_F0, run by run :: sha256 c6c4138ee5c86423 :: prose twin: two significant figures, direction in words
\newcommand{\MlstMlpCmpJcTwoQcdFonerOverFzeroPrMax}{up to 11\% fewer errors}  % experiments/FIGS/data/v2/index.json :: nonlinear probe jc2_188_qcd: (1 - AUC) of FLAV_F1R over FLAV_F0, run by run :: sha256 c6c4138ee5c86423 :: prose twin of the run farther from 1
\newcommand{\MlstMlpCmpJcTwoQcdFzeroOverCoarse}{0.993 and 1.012}  % experiments/FIGS/data/v2/index.json :: nonlinear probe jc2_188_qcd: (1 - AUC) of FLAV_F0 over R16_Q1, run by run :: sha256 c6c4138ee5c86423
\newcommand{\MlstMlpCmpJcTwoQcdFzeroOverCoarsePr}{equal within 2\%}  % experiments/FIGS/data/v2/index.json :: nonlinear probe jc2_188_qcd: (1 - AUC) of FLAV_F0 over R16_Q1, run by run :: sha256 c6c4138ee5c86423 :: prose twin: two significant figures, direction in words
\newcommand{\MlstMlpCmpJcTwoQcdFzeroOverCoarsePrMax}{equal within 2\%}  % experiments/FIGS/data/v2/index.json :: nonlinear probe jc2_188_qcd: (1 - AUC) of FLAV_F0 over R16_Q1, run by run :: sha256 c6c4138ee5c86423 :: prose twin of the run farther from 1
\newcommand{\MlstMlpCmpJcTwoQcdRandOverFull}{\pending{no run yet}}  % experiments/FIGS/data/v2/index.json :: nonlinear probe jc2_188_qcd: (1 - AUC) of each random partition over the 188-class run of its seed :: sha256 c6c4138ee5c86423
\newcommand{\MlstMlpCmpJcTwoQcdRandOverFullPr}{\pending{no run yet}}  % experiments/FIGS/data/v2/index.json :: nonlinear probe jc2_188_qcd: (1 - AUC) of each random partition over the 188-class run of its seed :: sha256 c6c4138ee5c86423 :: prose twin: two significant figures, direction in words
\newcommand{\MlstMlpCmpBcResOverFull}{1.002\,\ensuremath{\pm}\,0.006}  % experiments/FIGS/data/v2/index.json :: nonlinear probe bc_vs_bq_cs: (1 - AUC) of L162 over L188, run by run :: sha256 c6c4138ee5c86423
\newcommand{\MlstMlpCmpBcResOverFullPr}{equal within 1\%}  % experiments/FIGS/data/v2/index.json :: nonlinear probe bc_vs_bq_cs: (1 - AUC) of L162 over L188, run by run :: sha256 c6c4138ee5c86423 :: prose twin: two significant figures, direction in words
\newcommand{\MlstMlpCmpBcMidOverFull}{1.104\,\ensuremath{\pm}\,0.017}  % experiments/FIGS/data/v2/index.json :: nonlinear probe bc_vs_bq_cs: (1 - AUC) of R42_Q1 over L188, run by run :: sha256 c6c4138ee5c86423
\newcommand{\MlstMlpCmpBcMidOverFullPr}{10\% more errors}  % experiments/FIGS/data/v2/index.json :: nonlinear probe bc_vs_bq_cs: (1 - AUC) of R42_Q1 over L188, run by run :: sha256 c6c4138ee5c86423 :: prose twin: two significant figures, direction in words
\newcommand{\MlstMlpCmpBcCoarseOverFull}{1.624\,\ensuremath{\pm}\,0.022}  % experiments/FIGS/data/v2/index.json :: nonlinear probe bc_vs_bq_cs: (1 - AUC) of R16_Q1 over L188, run by run :: sha256 c6c4138ee5c86423
\newcommand{\MlstMlpCmpBcCoarseOverFullPr}{62\% more errors}  % experiments/FIGS/data/v2/index.json :: nonlinear probe bc_vs_bq_cs: (1 - AUC) of R16_Q1 over L188, run by run :: sha256 c6c4138ee5c86423 :: prose twin: two significant figures, direction in words
\newcommand{\MlstMlpCmpBcResMassOverRes}{\pending{no run yet}}  % experiments/FIGS/data/v2/index.json :: nonlinear probe bc_vs_bq_cs: (1 - AUC) of L162_MASS over L162, run by run :: sha256 c6c4138ee5c86423
\newcommand{\MlstMlpCmpBcResMassOverResPr}{\pending{no run yet}}  % experiments/FIGS/data/v2/index.json :: nonlinear probe bc_vs_bq_cs: (1 - AUC) of L162_MASS over L162, run by run :: sha256 c6c4138ee5c86423 :: prose twin: two significant figures, direction in words
\newcommand{\MlstMlpCmpBcResMassOverResPrMax}{\pending{no run yet}}  % experiments/FIGS/data/v2/index.json :: nonlinear probe bc_vs_bq_cs: (1 - AUC) of L162_MASS over L162, run by run :: sha256 c6c4138ee5c86423 :: prose twin of the run farther from 1
\newcommand{\MlstMlpCmpBcCoarseMassOverCoarse}{\pending{no run yet}}  % experiments/FIGS/data/v2/index.json :: nonlinear probe bc_vs_bq_cs: (1 - AUC) of R16_Q1_MASS over R16_Q1, run by run :: sha256 c6c4138ee5c86423
\newcommand{\MlstMlpCmpBcCoarseMassOverCoarsePr}{\pending{no run yet}}  % experiments/FIGS/data/v2/index.json :: nonlinear probe bc_vs_bq_cs: (1 - AUC) of R16_Q1_MASS over R16_Q1, run by run :: sha256 c6c4138ee5c86423 :: prose twin: two significant figures, direction in words
\newcommand{\MlstMlpCmpBcCoarseMassOverCoarsePrMax}{\pending{no run yet}}  % experiments/FIGS/data/v2/index.json :: nonlinear probe bc_vs_bq_cs: (1 - AUC) of R16_Q1_MASS over R16_Q1, run by run :: sha256 c6c4138ee5c86423 :: prose twin of the run farther from 1
\newcommand{\MlstMlpCmpBcFzeroOverFull}{1.633 and 1.679}  % experiments/FIGS/data/v2/index.json :: nonlinear probe bc_vs_bq_cs: (1 - AUC) of FLAV_F0 over L188, run by run :: sha256 c6c4138ee5c86423
\newcommand{\MlstMlpCmpBcFzeroOverFullPr}{at least 63\% more errors}  % experiments/FIGS/data/v2/index.json :: nonlinear probe bc_vs_bq_cs: (1 - AUC) of FLAV_F0 over L188, run by run :: sha256 c6c4138ee5c86423 :: prose twin: two significant figures, direction in words
\newcommand{\MlstMlpCmpBcFzeroOverFullPrMax}{up to 68\% more errors}  % experiments/FIGS/data/v2/index.json :: nonlinear probe bc_vs_bq_cs: (1 - AUC) of FLAV_F0 over L188, run by run :: sha256 c6c4138ee5c86423 :: prose twin of the run farther from 1
\newcommand{\MlstMlpCmpBcFoneOverFzero}{0.848 and 0.810}  % experiments/FIGS/data/v2/index.json :: nonlinear probe bc_vs_bq_cs: (1 - AUC) of FLAV_F1 over FLAV_F0, run by run :: sha256 c6c4138ee5c86423
\newcommand{\MlstMlpCmpBcFoneOverFzeroPr}{at least 15\% fewer errors}  % experiments/FIGS/data/v2/index.json :: nonlinear probe bc_vs_bq_cs: (1 - AUC) of FLAV_F1 over FLAV_F0, run by run :: sha256 c6c4138ee5c86423 :: prose twin: two significant figures, direction in words
\newcommand{\MlstMlpCmpBcFoneOverFzeroPrMax}{up to 19\% fewer errors}  % experiments/FIGS/data/v2/index.json :: nonlinear probe bc_vs_bq_cs: (1 - AUC) of FLAV_F1 over FLAV_F0, run by run :: sha256 c6c4138ee5c86423 :: prose twin of the run farther from 1
\newcommand{\MlstMlpCmpBcFonerOverFone}{0.9812 and 0.9815}  % experiments/FIGS/data/v2/index.json :: nonlinear probe bc_vs_bq_cs: (1 - AUC) of FLAV_F1R over FLAV_F1, run by run :: sha256 c6c4138ee5c86423
\newcommand{\MlstMlpCmpBcFonerOverFonePr}{equal within 2\%}  % experiments/FIGS/data/v2/index.json :: nonlinear probe bc_vs_bq_cs: (1 - AUC) of FLAV_F1R over FLAV_F1, run by run :: sha256 c6c4138ee5c86423 :: prose twin: two significant figures, direction in words
\newcommand{\MlstMlpCmpBcFonerOverFonePrMax}{equal within 2\%}  % experiments/FIGS/data/v2/index.json :: nonlinear probe bc_vs_bq_cs: (1 - AUC) of FLAV_F1R over FLAV_F1, run by run :: sha256 c6c4138ee5c86423 :: prose twin of the run farther from 1
\newcommand{\MlstMlpCmpBcFonerOverFzero}{0.833 and 0.795}  % experiments/FIGS/data/v2/index.json :: nonlinear probe bc_vs_bq_cs: (1 - AUC) of FLAV_F1R over FLAV_F0, run by run :: sha256 c6c4138ee5c86423
\newcommand{\MlstMlpCmpBcFonerOverFzeroPr}{at least 17\% fewer errors}  % experiments/FIGS/data/v2/index.json :: nonlinear probe bc_vs_bq_cs: (1 - AUC) of FLAV_F1R over FLAV_F0, run by run :: sha256 c6c4138ee5c86423 :: prose twin: two significant figures, direction in words
\newcommand{\MlstMlpCmpBcFonerOverFzeroPrMax}{up to 21\% fewer errors}  % experiments/FIGS/data/v2/index.json :: nonlinear probe bc_vs_bq_cs: (1 - AUC) of FLAV_F1R over FLAV_F0, run by run :: sha256 c6c4138ee5c86423 :: prose twin of the run farther from 1
\newcommand{\MlstMlpCmpBcFzeroOverCoarse}{0.994 and 1.030}  % experiments/FIGS/data/v2/index.json :: nonlinear probe bc_vs_bq_cs: (1 - AUC) of FLAV_F0 over R16_Q1, run by run :: sha256 c6c4138ee5c86423
\newcommand{\MlstMlpCmpBcFzeroOverCoarsePr}{equal within 3\%}  % experiments/FIGS/data/v2/index.json :: nonlinear probe bc_vs_bq_cs: (1 - AUC) of FLAV_F0 over R16_Q1, run by run :: sha256 c6c4138ee5c86423 :: prose twin: two significant figures, direction in words
\newcommand{\MlstMlpCmpBcFzeroOverCoarsePrMax}{equal within 3\%}  % experiments/FIGS/data/v2/index.json :: nonlinear probe bc_vs_bq_cs: (1 - AUC) of FLAV_F0 over R16_Q1, run by run :: sha256 c6c4138ee5c86423 :: prose twin of the run farther from 1
\newcommand{\MlstMlpCmpBcRandOverFull}{\pending{no run yet}}  % experiments/FIGS/data/v2/index.json :: nonlinear probe bc_vs_bq_cs: (1 - AUC) of each random partition over the 188-class run of its seed :: sha256 c6c4138ee5c86423
\newcommand{\MlstMlpCmpBcRandOverFullPr}{\pending{no run yet}}  % experiments/FIGS/data/v2/index.json :: nonlinear probe bc_vs_bq_cs: (1 - AUC) of each random partition over the 188-class run of its seed :: sha256 c6c4138ee5c86423 :: prose twin: two significant figures, direction in words
\newcommand{\MlstMlpCmpBbResOverFull}{1.023\,\ensuremath{\pm}\,0.021}  % experiments/FIGS/data/v2/index.json :: nonlinear probe bb_vs_cc: (1 - AUC) of L162 over L188, run by run :: sha256 c6c4138ee5c86423
\newcommand{\MlstMlpCmpBbResOverFullPr}{equal within 3\%}  % experiments/FIGS/data/v2/index.json :: nonlinear probe bb_vs_cc: (1 - AUC) of L162 over L188, run by run :: sha256 c6c4138ee5c86423 :: prose twin: two significant figures, direction in words
\newcommand{\MlstMlpCmpBbMidOverFull}{1.125\,\ensuremath{\pm}\,0.011}  % experiments/FIGS/data/v2/index.json :: nonlinear probe bb_vs_cc: (1 - AUC) of R42_Q1 over L188, run by run :: sha256 c6c4138ee5c86423
\newcommand{\MlstMlpCmpBbMidOverFullPr}{13\% more errors}  % experiments/FIGS/data/v2/index.json :: nonlinear probe bb_vs_cc: (1 - AUC) of R42_Q1 over L188, run by run :: sha256 c6c4138ee5c86423 :: prose twin: two significant figures, direction in words
\newcommand{\MlstMlpCmpBbCoarseOverFull}{3.04\,\ensuremath{\pm}\,0.09}  % experiments/FIGS/data/v2/index.json :: nonlinear probe bb_vs_cc: (1 - AUC) of R16_Q1 over L188, run by run :: sha256 c6c4138ee5c86423
\newcommand{\MlstMlpCmpBbCoarseOverFullPr}{3.0 times the errors}  % experiments/FIGS/data/v2/index.json :: nonlinear probe bb_vs_cc: (1 - AUC) of R16_Q1 over L188, run by run :: sha256 c6c4138ee5c86423 :: prose twin: two significant figures, direction in words
\newcommand{\MlstMlpCmpBbResMassOverRes}{\pending{no run yet}}  % experiments/FIGS/data/v2/index.json :: nonlinear probe bb_vs_cc: (1 - AUC) of L162_MASS over L162, run by run :: sha256 c6c4138ee5c86423
\newcommand{\MlstMlpCmpBbResMassOverResPr}{\pending{no run yet}}  % experiments/FIGS/data/v2/index.json :: nonlinear probe bb_vs_cc: (1 - AUC) of L162_MASS over L162, run by run :: sha256 c6c4138ee5c86423 :: prose twin: two significant figures, direction in words
\newcommand{\MlstMlpCmpBbResMassOverResPrMax}{\pending{no run yet}}  % experiments/FIGS/data/v2/index.json :: nonlinear probe bb_vs_cc: (1 - AUC) of L162_MASS over L162, run by run :: sha256 c6c4138ee5c86423 :: prose twin of the run farther from 1
\newcommand{\MlstMlpCmpBbCoarseMassOverCoarse}{\pending{no run yet}}  % experiments/FIGS/data/v2/index.json :: nonlinear probe bb_vs_cc: (1 - AUC) of R16_Q1_MASS over R16_Q1, run by run :: sha256 c6c4138ee5c86423
\newcommand{\MlstMlpCmpBbCoarseMassOverCoarsePr}{\pending{no run yet}}  % experiments/FIGS/data/v2/index.json :: nonlinear probe bb_vs_cc: (1 - AUC) of R16_Q1_MASS over R16_Q1, run by run :: sha256 c6c4138ee5c86423 :: prose twin: two significant figures, direction in words
\newcommand{\MlstMlpCmpBbCoarseMassOverCoarsePrMax}{\pending{no run yet}}  % experiments/FIGS/data/v2/index.json :: nonlinear probe bb_vs_cc: (1 - AUC) of R16_Q1_MASS over R16_Q1, run by run :: sha256 c6c4138ee5c86423 :: prose twin of the run farther from 1
\newcommand{\MlstMlpCmpBbFzeroOverFull}{3.3 and 4.3}  % experiments/FIGS/data/v2/index.json :: nonlinear probe bb_vs_cc: (1 - AUC) of FLAV_F0 over L188, run by run :: sha256 c6c4138ee5c86423
\newcommand{\MlstMlpCmpBbFzeroOverFullPr}{at least 3.3 times the errors}  % experiments/FIGS/data/v2/index.json :: nonlinear probe bb_vs_cc: (1 - AUC) of FLAV_F0 over L188, run by run :: sha256 c6c4138ee5c86423 :: prose twin: two significant figures, direction in words
\newcommand{\MlstMlpCmpBbFzeroOverFullPrMax}{up to 4.3 times the errors}  % experiments/FIGS/data/v2/index.json :: nonlinear probe bb_vs_cc: (1 - AUC) of FLAV_F0 over L188, run by run :: sha256 c6c4138ee5c86423 :: prose twin of the run farther from 1
\newcommand{\MlstMlpCmpBbFoneOverFzero}{0.52 and 0.38}  % experiments/FIGS/data/v2/index.json :: nonlinear probe bb_vs_cc: (1 - AUC) of FLAV_F1 over FLAV_F0, run by run :: sha256 c6c4138ee5c86423
\newcommand{\MlstMlpCmpBbFoneOverFzeroPr}{at least 48\% fewer errors}  % experiments/FIGS/data/v2/index.json :: nonlinear probe bb_vs_cc: (1 - AUC) of FLAV_F1 over FLAV_F0, run by run :: sha256 c6c4138ee5c86423 :: prose twin: two significant figures, direction in words
\newcommand{\MlstMlpCmpBbFoneOverFzeroPrMax}{down to 0.38 times the errors}  % experiments/FIGS/data/v2/index.json :: nonlinear probe bb_vs_cc: (1 - AUC) of FLAV_F1 over FLAV_F0, run by run :: sha256 c6c4138ee5c86423 :: prose twin of the run farther from 1
\newcommand{\MlstMlpCmpBbFonerOverFone}{1.27 and 1.45}  % experiments/FIGS/data/v2/index.json :: nonlinear probe bb_vs_cc: (1 - AUC) of FLAV_F1R over FLAV_F1, run by run :: sha256 c6c4138ee5c86423
\newcommand{\MlstMlpCmpBbFonerOverFonePr}{at least 27\% more errors}  % experiments/FIGS/data/v2/index.json :: nonlinear probe bb_vs_cc: (1 - AUC) of FLAV_F1R over FLAV_F1, run by run :: sha256 c6c4138ee5c86423 :: prose twin: two significant figures, direction in words
\newcommand{\MlstMlpCmpBbFonerOverFonePrMax}{up to 45\% more errors}  % experiments/FIGS/data/v2/index.json :: nonlinear probe bb_vs_cc: (1 - AUC) of FLAV_F1R over FLAV_F1, run by run :: sha256 c6c4138ee5c86423 :: prose twin of the run farther from 1
\newcommand{\MlstMlpCmpBbFonerOverFzero}{0.66 and 0.55}  % experiments/FIGS/data/v2/index.json :: nonlinear probe bb_vs_cc: (1 - AUC) of FLAV_F1R over FLAV_F0, run by run :: sha256 c6c4138ee5c86423
\newcommand{\MlstMlpCmpBbFonerOverFzeroPr}{at least 34\% fewer errors}  % experiments/FIGS/data/v2/index.json :: nonlinear probe bb_vs_cc: (1 - AUC) of FLAV_F1R over FLAV_F0, run by run :: sha256 c6c4138ee5c86423 :: prose twin: two significant figures, direction in words
\newcommand{\MlstMlpCmpBbFonerOverFzeroPrMax}{up to 45\% fewer errors}  % experiments/FIGS/data/v2/index.json :: nonlinear probe bb_vs_cc: (1 - AUC) of FLAV_F1R over FLAV_F0, run by run :: sha256 c6c4138ee5c86423 :: prose twin of the run farther from 1
\newcommand{\MlstMlpCmpBbFzeroOverCoarse}{1.07 and 1.41}  % experiments/FIGS/data/v2/index.json :: nonlinear probe bb_vs_cc: (1 - AUC) of FLAV_F0 over R16_Q1, run by run :: sha256 c6c4138ee5c86423
\newcommand{\MlstMlpCmpBbFzeroOverCoarsePr}{at least 7.4\% more errors}  % experiments/FIGS/data/v2/index.json :: nonlinear probe bb_vs_cc: (1 - AUC) of FLAV_F0 over R16_Q1, run by run :: sha256 c6c4138ee5c86423 :: prose twin: two significant figures, direction in words
\newcommand{\MlstMlpCmpBbFzeroOverCoarsePrMax}{up to 41\% more errors}  % experiments/FIGS/data/v2/index.json :: nonlinear probe bb_vs_cc: (1 - AUC) of FLAV_F0 over R16_Q1, run by run :: sha256 c6c4138ee5c86423 :: prose twin of the run farther from 1
\newcommand{\MlstMlpCmpBbRandOverFull}{\pending{no run yet}}  % experiments/FIGS/data/v2/index.json :: nonlinear probe bb_vs_cc: (1 - AUC) of each random partition over the 188-class run of its seed :: sha256 c6c4138ee5c86423
\newcommand{\MlstMlpCmpBbRandOverFullPr}{\pending{no run yet}}  % experiments/FIGS/data/v2/index.json :: nonlinear probe bb_vs_cc: (1 - AUC) of each random partition over the 188-class run of its seed :: sha256 c6c4138ee5c86423 :: prose twin: two significant figures, direction in words
\newcommand{\MlstMlpCmpBbqqResOverFull}{1.011\,\ensuremath{\pm}\,0.030}  % experiments/FIGS/data/v2/index.json :: nonlinear probe bbqq_vs_ccqq: (1 - AUC) of L162 over L188, run by run :: sha256 c6c4138ee5c86423
\newcommand{\MlstMlpCmpBbqqResOverFullPr}{equal within 3\%}  % experiments/FIGS/data/v2/index.json :: nonlinear probe bbqq_vs_ccqq: (1 - AUC) of L162 over L188, run by run :: sha256 c6c4138ee5c86423 :: prose twin: two significant figures, direction in words
\newcommand{\MlstMlpCmpBbqqMidOverFull}{1.112\,\ensuremath{\pm}\,0.006}  % experiments/FIGS/data/v2/index.json :: nonlinear probe bbqq_vs_ccqq: (1 - AUC) of R42_Q1 over L188, run by run :: sha256 c6c4138ee5c86423
\newcommand{\MlstMlpCmpBbqqMidOverFullPr}{11\% more errors}  % experiments/FIGS/data/v2/index.json :: nonlinear probe bbqq_vs_ccqq: (1 - AUC) of R42_Q1 over L188, run by run :: sha256 c6c4138ee5c86423 :: prose twin: two significant figures, direction in words
\newcommand{\MlstMlpCmpBbqqCoarseOverFull}{2.93\,\ensuremath{\pm}\,0.17}  % experiments/FIGS/data/v2/index.json :: nonlinear probe bbqq_vs_ccqq: (1 - AUC) of R16_Q1 over L188, run by run :: sha256 c6c4138ee5c86423
\newcommand{\MlstMlpCmpBbqqCoarseOverFullPr}{2.9 times the errors}  % experiments/FIGS/data/v2/index.json :: nonlinear probe bbqq_vs_ccqq: (1 - AUC) of R16_Q1 over L188, run by run :: sha256 c6c4138ee5c86423 :: prose twin: two significant figures, direction in words
\newcommand{\MlstMlpCmpBbqqResMassOverRes}{\pending{no run yet}}  % experiments/FIGS/data/v2/index.json :: nonlinear probe bbqq_vs_ccqq: (1 - AUC) of L162_MASS over L162, run by run :: sha256 c6c4138ee5c86423
\newcommand{\MlstMlpCmpBbqqResMassOverResPr}{\pending{no run yet}}  % experiments/FIGS/data/v2/index.json :: nonlinear probe bbqq_vs_ccqq: (1 - AUC) of L162_MASS over L162, run by run :: sha256 c6c4138ee5c86423 :: prose twin: two significant figures, direction in words
\newcommand{\MlstMlpCmpBbqqResMassOverResPrMax}{\pending{no run yet}}  % experiments/FIGS/data/v2/index.json :: nonlinear probe bbqq_vs_ccqq: (1 - AUC) of L162_MASS over L162, run by run :: sha256 c6c4138ee5c86423 :: prose twin of the run farther from 1
\newcommand{\MlstMlpCmpBbqqCoarseMassOverCoarse}{\pending{no run yet}}  % experiments/FIGS/data/v2/index.json :: nonlinear probe bbqq_vs_ccqq: (1 - AUC) of R16_Q1_MASS over R16_Q1, run by run :: sha256 c6c4138ee5c86423
\newcommand{\MlstMlpCmpBbqqCoarseMassOverCoarsePr}{\pending{no run yet}}  % experiments/FIGS/data/v2/index.json :: nonlinear probe bbqq_vs_ccqq: (1 - AUC) of R16_Q1_MASS over R16_Q1, run by run :: sha256 c6c4138ee5c86423 :: prose twin: two significant figures, direction in words
\newcommand{\MlstMlpCmpBbqqCoarseMassOverCoarsePrMax}{\pending{no run yet}}  % experiments/FIGS/data/v2/index.json :: nonlinear probe bbqq_vs_ccqq: (1 - AUC) of R16_Q1_MASS over R16_Q1, run by run :: sha256 c6c4138ee5c86423 :: prose twin of the run farther from 1
\newcommand{\MlstMlpCmpBbqqFzeroOverFull}{3.20 and 3.78}  % experiments/FIGS/data/v2/index.json :: nonlinear probe bbqq_vs_ccqq: (1 - AUC) of FLAV_F0 over L188, run by run :: sha256 c6c4138ee5c86423
\newcommand{\MlstMlpCmpBbqqFzeroOverFullPr}{at least 3.2 times the errors}  % experiments/FIGS/data/v2/index.json :: nonlinear probe bbqq_vs_ccqq: (1 - AUC) of FLAV_F0 over L188, run by run :: sha256 c6c4138ee5c86423 :: prose twin: two significant figures, direction in words
\newcommand{\MlstMlpCmpBbqqFzeroOverFullPrMax}{up to 3.8 times the errors}  % experiments/FIGS/data/v2/index.json :: nonlinear probe bbqq_vs_ccqq: (1 - AUC) of FLAV_F0 over L188, run by run :: sha256 c6c4138ee5c86423 :: prose twin of the run farther from 1
\newcommand{\MlstMlpCmpBbqqFoneOverFzero}{0.450 and 0.363}  % experiments/FIGS/data/v2/index.json :: nonlinear probe bbqq_vs_ccqq: (1 - AUC) of FLAV_F1 over FLAV_F0, run by run :: sha256 c6c4138ee5c86423
\newcommand{\MlstMlpCmpBbqqFoneOverFzeroPr}{at most 0.45 times the errors}  % experiments/FIGS/data/v2/index.json :: nonlinear probe bbqq_vs_ccqq: (1 - AUC) of FLAV_F1 over FLAV_F0, run by run :: sha256 c6c4138ee5c86423 :: prose twin: two significant figures, direction in words
\newcommand{\MlstMlpCmpBbqqFoneOverFzeroPrMax}{down to 0.36 times the errors}  % experiments/FIGS/data/v2/index.json :: nonlinear probe bbqq_vs_ccqq: (1 - AUC) of FLAV_F1 over FLAV_F0, run by run :: sha256 c6c4138ee5c86423 :: prose twin of the run farther from 1
\newcommand{\MlstMlpCmpBbqqFonerOverFone}{1.311 and 1.343}  % experiments/FIGS/data/v2/index.json :: nonlinear probe bbqq_vs_ccqq: (1 - AUC) of FLAV_F1R over FLAV_F1, run by run :: sha256 c6c4138ee5c86423
\newcommand{\MlstMlpCmpBbqqFonerOverFonePr}{at least 31\% more errors}  % experiments/FIGS/data/v2/index.json :: nonlinear probe bbqq_vs_ccqq: (1 - AUC) of FLAV_F1R over FLAV_F1, run by run :: sha256 c6c4138ee5c86423 :: prose twin: two significant figures, direction in words
\newcommand{\MlstMlpCmpBbqqFonerOverFonePrMax}{up to 34\% more errors}  % experiments/FIGS/data/v2/index.json :: nonlinear probe bbqq_vs_ccqq: (1 - AUC) of FLAV_F1R over FLAV_F1, run by run :: sha256 c6c4138ee5c86423 :: prose twin of the run farther from 1
\newcommand{\MlstMlpCmpBbqqFonerOverFzero}{0.59 and 0.49}  % experiments/FIGS/data/v2/index.json :: nonlinear probe bbqq_vs_ccqq: (1 - AUC) of FLAV_F1R over FLAV_F0, run by run :: sha256 c6c4138ee5c86423
\newcommand{\MlstMlpCmpBbqqFonerOverFzeroPr}{at least 41\% fewer errors}  % experiments/FIGS/data/v2/index.json :: nonlinear probe bbqq_vs_ccqq: (1 - AUC) of FLAV_F1R over FLAV_F0, run by run :: sha256 c6c4138ee5c86423 :: prose twin: two significant figures, direction in words
\newcommand{\MlstMlpCmpBbqqFonerOverFzeroPrMax}{down to 0.49 times the errors}  % experiments/FIGS/data/v2/index.json :: nonlinear probe bbqq_vs_ccqq: (1 - AUC) of FLAV_F1R over FLAV_F0, run by run :: sha256 c6c4138ee5c86423 :: prose twin of the run farther from 1
\newcommand{\MlstMlpCmpBbqqFzeroOverCoarse}{1.170 and 1.236}  % experiments/FIGS/data/v2/index.json :: nonlinear probe bbqq_vs_ccqq: (1 - AUC) of FLAV_F0 over R16_Q1, run by run :: sha256 c6c4138ee5c86423
\newcommand{\MlstMlpCmpBbqqFzeroOverCoarsePr}{at least 17\% more errors}  % experiments/FIGS/data/v2/index.json :: nonlinear probe bbqq_vs_ccqq: (1 - AUC) of FLAV_F0 over R16_Q1, run by run :: sha256 c6c4138ee5c86423 :: prose twin: two significant figures, direction in words
\newcommand{\MlstMlpCmpBbqqFzeroOverCoarsePrMax}{up to 24\% more errors}  % experiments/FIGS/data/v2/index.json :: nonlinear probe bbqq_vs_ccqq: (1 - AUC) of FLAV_F0 over R16_Q1, run by run :: sha256 c6c4138ee5c86423 :: prose twin of the run farther from 1
\newcommand{\MlstMlpCmpBbqqRandOverFull}{\pending{no run yet}}  % experiments/FIGS/data/v2/index.json :: nonlinear probe bbqq_vs_ccqq: (1 - AUC) of each random partition over the 188-class run of its seed :: sha256 c6c4138ee5c86423
\newcommand{\MlstMlpCmpBbqqRandOverFullPr}{\pending{no run yet}}  % experiments/FIGS/data/v2/index.json :: nonlinear probe bbqq_vs_ccqq: (1 - AUC) of each random partition over the 188-class run of its seed :: sha256 c6c4138ee5c86423 :: prose twin: two significant figures, direction in words
\newcommand{\MlstMlpCmpQcdBbResOverFull}{1.085\,\ensuremath{\pm}\,0.030}  % experiments/FIGS/data/v2/index.json :: nonlinear probe qcd_bb_vs_cc: (1 - AUC) of L162 over L188, run by run :: sha256 c6c4138ee5c86423
\newcommand{\MlstMlpCmpQcdBbResOverFullPr}{8.5\% more errors}  % experiments/FIGS/data/v2/index.json :: nonlinear probe qcd_bb_vs_cc: (1 - AUC) of L162 over L188, run by run :: sha256 c6c4138ee5c86423 :: prose twin: two significant figures, direction in words
\newcommand{\MlstMlpCmpQcdBbMidOverFull}{1.121\,\ensuremath{\pm}\,0.017}  % experiments/FIGS/data/v2/index.json :: nonlinear probe qcd_bb_vs_cc: (1 - AUC) of R42_Q1 over L188, run by run :: sha256 c6c4138ee5c86423
\newcommand{\MlstMlpCmpQcdBbMidOverFullPr}{12\% more errors}  % experiments/FIGS/data/v2/index.json :: nonlinear probe qcd_bb_vs_cc: (1 - AUC) of R42_Q1 over L188, run by run :: sha256 c6c4138ee5c86423 :: prose twin: two significant figures, direction in words
\newcommand{\MlstMlpCmpQcdBbCoarseOverFull}{2.31\,\ensuremath{\pm}\,0.04}  % experiments/FIGS/data/v2/index.json :: nonlinear probe qcd_bb_vs_cc: (1 - AUC) of R16_Q1 over L188, run by run :: sha256 c6c4138ee5c86423
\newcommand{\MlstMlpCmpQcdBbCoarseOverFullPr}{2.3 times the errors}  % experiments/FIGS/data/v2/index.json :: nonlinear probe qcd_bb_vs_cc: (1 - AUC) of R16_Q1 over L188, run by run :: sha256 c6c4138ee5c86423 :: prose twin: two significant figures, direction in words
\newcommand{\MlstMlpCmpQcdBbResMassOverRes}{\pending{no run yet}}  % experiments/FIGS/data/v2/index.json :: nonlinear probe qcd_bb_vs_cc: (1 - AUC) of L162_MASS over L162, run by run :: sha256 c6c4138ee5c86423
\newcommand{\MlstMlpCmpQcdBbResMassOverResPr}{\pending{no run yet}}  % experiments/FIGS/data/v2/index.json :: nonlinear probe qcd_bb_vs_cc: (1 - AUC) of L162_MASS over L162, run by run :: sha256 c6c4138ee5c86423 :: prose twin: two significant figures, direction in words
\newcommand{\MlstMlpCmpQcdBbResMassOverResPrMax}{\pending{no run yet}}  % experiments/FIGS/data/v2/index.json :: nonlinear probe qcd_bb_vs_cc: (1 - AUC) of L162_MASS over L162, run by run :: sha256 c6c4138ee5c86423 :: prose twin of the run farther from 1
\newcommand{\MlstMlpCmpQcdBbCoarseMassOverCoarse}{\pending{no run yet}}  % experiments/FIGS/data/v2/index.json :: nonlinear probe qcd_bb_vs_cc: (1 - AUC) of R16_Q1_MASS over R16_Q1, run by run :: sha256 c6c4138ee5c86423
\newcommand{\MlstMlpCmpQcdBbCoarseMassOverCoarsePr}{\pending{no run yet}}  % experiments/FIGS/data/v2/index.json :: nonlinear probe qcd_bb_vs_cc: (1 - AUC) of R16_Q1_MASS over R16_Q1, run by run :: sha256 c6c4138ee5c86423 :: prose twin: two significant figures, direction in words
\newcommand{\MlstMlpCmpQcdBbCoarseMassOverCoarsePrMax}{\pending{no run yet}}  % experiments/FIGS/data/v2/index.json :: nonlinear probe qcd_bb_vs_cc: (1 - AUC) of R16_Q1_MASS over R16_Q1, run by run :: sha256 c6c4138ee5c86423 :: prose twin of the run farther from 1
\newcommand{\MlstMlpCmpQcdBbFzeroOverFull}{2.50 and 2.72}  % experiments/FIGS/data/v2/index.json :: nonlinear probe qcd_bb_vs_cc: (1 - AUC) of FLAV_F0 over L188, run by run :: sha256 c6c4138ee5c86423
\newcommand{\MlstMlpCmpQcdBbFzeroOverFullPr}{at least 2.5 times the errors}  % experiments/FIGS/data/v2/index.json :: nonlinear probe qcd_bb_vs_cc: (1 - AUC) of FLAV_F0 over L188, run by run :: sha256 c6c4138ee5c86423 :: prose twin: two significant figures, direction in words
\newcommand{\MlstMlpCmpQcdBbFzeroOverFullPrMax}{up to 2.7 times the errors}  % experiments/FIGS/data/v2/index.json :: nonlinear probe qcd_bb_vs_cc: (1 - AUC) of FLAV_F0 over L188, run by run :: sha256 c6c4138ee5c86423 :: prose twin of the run farther from 1
\newcommand{\MlstMlpCmpQcdBbFoneOverFzero}{0.573 and 0.505}  % experiments/FIGS/data/v2/index.json :: nonlinear probe qcd_bb_vs_cc: (1 - AUC) of FLAV_F1 over FLAV_F0, run by run :: sha256 c6c4138ee5c86423
\newcommand{\MlstMlpCmpQcdBbFoneOverFzeroPr}{at least 43\% fewer errors}  % experiments/FIGS/data/v2/index.json :: nonlinear probe qcd_bb_vs_cc: (1 - AUC) of FLAV_F1 over FLAV_F0, run by run :: sha256 c6c4138ee5c86423 :: prose twin: two significant figures, direction in words
\newcommand{\MlstMlpCmpQcdBbFoneOverFzeroPrMax}{up to 49\% fewer errors}  % experiments/FIGS/data/v2/index.json :: nonlinear probe qcd_bb_vs_cc: (1 - AUC) of FLAV_F1 over FLAV_F0, run by run :: sha256 c6c4138ee5c86423 :: prose twin of the run farther from 1
\newcommand{\MlstMlpCmpQcdBbFonerOverFone}{1.219 and 1.258}  % experiments/FIGS/data/v2/index.json :: nonlinear probe qcd_bb_vs_cc: (1 - AUC) of FLAV_F1R over FLAV_F1, run by run :: sha256 c6c4138ee5c86423
\newcommand{\MlstMlpCmpQcdBbFonerOverFonePr}{at least 22\% more errors}  % experiments/FIGS/data/v2/index.json :: nonlinear probe qcd_bb_vs_cc: (1 - AUC) of FLAV_F1R over FLAV_F1, run by run :: sha256 c6c4138ee5c86423 :: prose twin: two significant figures, direction in words
\newcommand{\MlstMlpCmpQcdBbFonerOverFonePrMax}{up to 26\% more errors}  % experiments/FIGS/data/v2/index.json :: nonlinear probe qcd_bb_vs_cc: (1 - AUC) of FLAV_F1R over FLAV_F1, run by run :: sha256 c6c4138ee5c86423 :: prose twin of the run farther from 1
\newcommand{\MlstMlpCmpQcdBbFonerOverFzero}{0.699 and 0.635}  % experiments/FIGS/data/v2/index.json :: nonlinear probe qcd_bb_vs_cc: (1 - AUC) of FLAV_F1R over FLAV_F0, run by run :: sha256 c6c4138ee5c86423
\newcommand{\MlstMlpCmpQcdBbFonerOverFzeroPr}{at least 30\% fewer errors}  % experiments/FIGS/data/v2/index.json :: nonlinear probe qcd_bb_vs_cc: (1 - AUC) of FLAV_F1R over FLAV_F0, run by run :: sha256 c6c4138ee5c86423 :: prose twin: two significant figures, direction in words
\newcommand{\MlstMlpCmpQcdBbFonerOverFzeroPrMax}{up to 36\% fewer errors}  % experiments/FIGS/data/v2/index.json :: nonlinear probe qcd_bb_vs_cc: (1 - AUC) of FLAV_F1R over FLAV_F0, run by run :: sha256 c6c4138ee5c86423 :: prose twin of the run farther from 1
\newcommand{\MlstMlpCmpQcdBbFzeroOverCoarse}{1.103 and 1.167}  % experiments/FIGS/data/v2/index.json :: nonlinear probe qcd_bb_vs_cc: (1 - AUC) of FLAV_F0 over R16_Q1, run by run :: sha256 c6c4138ee5c86423
\newcommand{\MlstMlpCmpQcdBbFzeroOverCoarsePr}{at least 10\% more errors}  % experiments/FIGS/data/v2/index.json :: nonlinear probe qcd_bb_vs_cc: (1 - AUC) of FLAV_F0 over R16_Q1, run by run :: sha256 c6c4138ee5c86423 :: prose twin: two significant figures, direction in words
\newcommand{\MlstMlpCmpQcdBbFzeroOverCoarsePrMax}{up to 17\% more errors}  % experiments/FIGS/data/v2/index.json :: nonlinear probe qcd_bb_vs_cc: (1 - AUC) of FLAV_F0 over R16_Q1, run by run :: sha256 c6c4138ee5c86423 :: prose twin of the run farther from 1
\newcommand{\MlstMlpCmpQcdBbRandOverFull}{\pending{no run yet}}  % experiments/FIGS/data/v2/index.json :: nonlinear probe qcd_bb_vs_cc: (1 - AUC) of each random partition over the 188-class run of its seed :: sha256 c6c4138ee5c86423
\newcommand{\MlstMlpCmpQcdBbRandOverFullPr}{\pending{no run yet}}  % experiments/FIGS/data/v2/index.json :: nonlinear probe qcd_bb_vs_cc: (1 - AUC) of each random partition over the 188-class run of its seed :: sha256 c6c4138ee5c86423 :: prose twin: two significant figures, direction in words
\newcommand{\MlstMlpCmpQcdBResOverFull}{1.050\,\ensuremath{\pm}\,0.004}  % experiments/FIGS/data/v2/index.json :: nonlinear probe qcd_b_vs_light: (1 - AUC) of L162 over L188, run by run :: sha256 c6c4138ee5c86423
\newcommand{\MlstMlpCmpQcdBResOverFullPr}{5.0\% more errors}  % experiments/FIGS/data/v2/index.json :: nonlinear probe qcd_b_vs_light: (1 - AUC) of L162 over L188, run by run :: sha256 c6c4138ee5c86423 :: prose twin: two significant figures, direction in words
\newcommand{\MlstMlpCmpQcdBMidOverFull}{1.055\,\ensuremath{\pm}\,0.005}  % experiments/FIGS/data/v2/index.json :: nonlinear probe qcd_b_vs_light: (1 - AUC) of R42_Q1 over L188, run by run :: sha256 c6c4138ee5c86423
\newcommand{\MlstMlpCmpQcdBMidOverFullPr}{5.5\% more errors}  % experiments/FIGS/data/v2/index.json :: nonlinear probe qcd_b_vs_light: (1 - AUC) of R42_Q1 over L188, run by run :: sha256 c6c4138ee5c86423 :: prose twin: two significant figures, direction in words
\newcommand{\MlstMlpCmpQcdBCoarseOverFull}{1.337\,\ensuremath{\pm}\,0.021}  % experiments/FIGS/data/v2/index.json :: nonlinear probe qcd_b_vs_light: (1 - AUC) of R16_Q1 over L188, run by run :: sha256 c6c4138ee5c86423
\newcommand{\MlstMlpCmpQcdBCoarseOverFullPr}{34\% more errors}  % experiments/FIGS/data/v2/index.json :: nonlinear probe qcd_b_vs_light: (1 - AUC) of R16_Q1 over L188, run by run :: sha256 c6c4138ee5c86423 :: prose twin: two significant figures, direction in words
\newcommand{\MlstMlpCmpQcdBResMassOverRes}{\pending{no run yet}}  % experiments/FIGS/data/v2/index.json :: nonlinear probe qcd_b_vs_light: (1 - AUC) of L162_MASS over L162, run by run :: sha256 c6c4138ee5c86423
\newcommand{\MlstMlpCmpQcdBResMassOverResPr}{\pending{no run yet}}  % experiments/FIGS/data/v2/index.json :: nonlinear probe qcd_b_vs_light: (1 - AUC) of L162_MASS over L162, run by run :: sha256 c6c4138ee5c86423 :: prose twin: two significant figures, direction in words
\newcommand{\MlstMlpCmpQcdBResMassOverResPrMax}{\pending{no run yet}}  % experiments/FIGS/data/v2/index.json :: nonlinear probe qcd_b_vs_light: (1 - AUC) of L162_MASS over L162, run by run :: sha256 c6c4138ee5c86423 :: prose twin of the run farther from 1
\newcommand{\MlstMlpCmpQcdBCoarseMassOverCoarse}{\pending{no run yet}}  % experiments/FIGS/data/v2/index.json :: nonlinear probe qcd_b_vs_light: (1 - AUC) of R16_Q1_MASS over R16_Q1, run by run :: sha256 c6c4138ee5c86423
\newcommand{\MlstMlpCmpQcdBCoarseMassOverCoarsePr}{\pending{no run yet}}  % experiments/FIGS/data/v2/index.json :: nonlinear probe qcd_b_vs_light: (1 - AUC) of R16_Q1_MASS over R16_Q1, run by run :: sha256 c6c4138ee5c86423 :: prose twin: two significant figures, direction in words
\newcommand{\MlstMlpCmpQcdBCoarseMassOverCoarsePrMax}{\pending{no run yet}}  % experiments/FIGS/data/v2/index.json :: nonlinear probe qcd_b_vs_light: (1 - AUC) of R16_Q1_MASS over R16_Q1, run by run :: sha256 c6c4138ee5c86423 :: prose twin of the run farther from 1
\newcommand{\MlstMlpCmpQcdBFzeroOverFull}{1.315 and 1.288}  % experiments/FIGS/data/v2/index.json :: nonlinear probe qcd_b_vs_light: (1 - AUC) of FLAV_F0 over L188, run by run :: sha256 c6c4138ee5c86423
\newcommand{\MlstMlpCmpQcdBFzeroOverFullPr}{at least 29\% more errors}  % experiments/FIGS/data/v2/index.json :: nonlinear probe qcd_b_vs_light: (1 - AUC) of FLAV_F0 over L188, run by run :: sha256 c6c4138ee5c86423 :: prose twin: two significant figures, direction in words
\newcommand{\MlstMlpCmpQcdBFzeroOverFullPrMax}{up to 32\% more errors}  % experiments/FIGS/data/v2/index.json :: nonlinear probe qcd_b_vs_light: (1 - AUC) of FLAV_F0 over L188, run by run :: sha256 c6c4138ee5c86423 :: prose twin of the run farther from 1
\newcommand{\MlstMlpCmpQcdBFoneOverFzero}{0.8804 and 0.8849}  % experiments/FIGS/data/v2/index.json :: nonlinear probe qcd_b_vs_light: (1 - AUC) of FLAV_F1 over FLAV_F0, run by run :: sha256 c6c4138ee5c86423
\newcommand{\MlstMlpCmpQcdBFoneOverFzeroPr}{at least 12\% fewer errors}  % experiments/FIGS/data/v2/index.json :: nonlinear probe qcd_b_vs_light: (1 - AUC) of FLAV_F1 over FLAV_F0, run by run :: sha256 c6c4138ee5c86423 :: prose twin: two significant figures, direction in words
\newcommand{\MlstMlpCmpQcdBFoneOverFzeroPrMax}{up to 12\% fewer errors}  % experiments/FIGS/data/v2/index.json :: nonlinear probe qcd_b_vs_light: (1 - AUC) of FLAV_F1 over FLAV_F0, run by run :: sha256 c6c4138ee5c86423 :: prose twin of the run farther from 1
\newcommand{\MlstMlpCmpQcdBFonerOverFone}{1.037 and 1.042}  % experiments/FIGS/data/v2/index.json :: nonlinear probe qcd_b_vs_light: (1 - AUC) of FLAV_F1R over FLAV_F1, run by run :: sha256 c6c4138ee5c86423
\newcommand{\MlstMlpCmpQcdBFonerOverFonePr}{at least 3.7\% more errors}  % experiments/FIGS/data/v2/index.json :: nonlinear probe qcd_b_vs_light: (1 - AUC) of FLAV_F1R over FLAV_F1, run by run :: sha256 c6c4138ee5c86423 :: prose twin: two significant figures, direction in words
\newcommand{\MlstMlpCmpQcdBFonerOverFonePrMax}{up to 4.2\% more errors}  % experiments/FIGS/data/v2/index.json :: nonlinear probe qcd_b_vs_light: (1 - AUC) of FLAV_F1R over FLAV_F1, run by run :: sha256 c6c4138ee5c86423 :: prose twin of the run farther from 1
\newcommand{\MlstMlpCmpQcdBFonerOverFzero}{0.9130 and 0.9217}  % experiments/FIGS/data/v2/index.json :: nonlinear probe qcd_b_vs_light: (1 - AUC) of FLAV_F1R over FLAV_F0, run by run :: sha256 c6c4138ee5c86423
\newcommand{\MlstMlpCmpQcdBFonerOverFzeroPr}{at least 7.8\% fewer errors}  % experiments/FIGS/data/v2/index.json :: nonlinear probe qcd_b_vs_light: (1 - AUC) of FLAV_F1R over FLAV_F0, run by run :: sha256 c6c4138ee5c86423 :: prose twin: two significant figures, direction in words
\newcommand{\MlstMlpCmpQcdBFonerOverFzeroPrMax}{up to 8.7\% fewer errors}  % experiments/FIGS/data/v2/index.json :: nonlinear probe qcd_b_vs_light: (1 - AUC) of FLAV_F1R over FLAV_F0, run by run :: sha256 c6c4138ee5c86423 :: prose twin of the run farther from 1
\newcommand{\MlstMlpCmpQcdBFzeroOverCoarse}{0.968 and 0.978}  % experiments/FIGS/data/v2/index.json :: nonlinear probe qcd_b_vs_light: (1 - AUC) of FLAV_F0 over R16_Q1, run by run :: sha256 c6c4138ee5c86423
\newcommand{\MlstMlpCmpQcdBFzeroOverCoarsePr}{at least 2.2\% fewer errors}  % experiments/FIGS/data/v2/index.json :: nonlinear probe qcd_b_vs_light: (1 - AUC) of FLAV_F0 over R16_Q1, run by run :: sha256 c6c4138ee5c86423 :: prose twin: two significant figures, direction in words
\newcommand{\MlstMlpCmpQcdBFzeroOverCoarsePrMax}{up to 3.2\% fewer errors}  % experiments/FIGS/data/v2/index.json :: nonlinear probe qcd_b_vs_light: (1 - AUC) of FLAV_F0 over R16_Q1, run by run :: sha256 c6c4138ee5c86423 :: prose twin of the run farther from 1
\newcommand{\MlstMlpCmpQcdBRandOverFull}{\pending{no run yet}}  % experiments/FIGS/data/v2/index.json :: nonlinear probe qcd_b_vs_light: (1 - AUC) of each random partition over the 188-class run of its seed :: sha256 c6c4138ee5c86423
\newcommand{\MlstMlpCmpQcdBRandOverFullPr}{\pending{no run yet}}  % experiments/FIGS/data/v2/index.json :: nonlinear probe qcd_b_vs_light: (1 - AUC) of each random partition over the 188-class run of its seed :: sha256 c6c4138ee5c86423 :: prose twin: two significant figures, direction in words
\newcommand{\MlstMlpCmpQcdCResOverFull}{1.039\,\ensuremath{\pm}\,0.005}  % experiments/FIGS/data/v2/index.json :: nonlinear probe qcd_c_vs_light: (1 - AUC) of L162 over L188, run by run :: sha256 c6c4138ee5c86423
\newcommand{\MlstMlpCmpQcdCResOverFullPr}{3.9\% more errors}  % experiments/FIGS/data/v2/index.json :: nonlinear probe qcd_c_vs_light: (1 - AUC) of L162 over L188, run by run :: sha256 c6c4138ee5c86423 :: prose twin: two significant figures, direction in words
\newcommand{\MlstMlpCmpQcdCMidOverFull}{1.055\,\ensuremath{\pm}\,0.006}  % experiments/FIGS/data/v2/index.json :: nonlinear probe qcd_c_vs_light: (1 - AUC) of R42_Q1 over L188, run by run :: sha256 c6c4138ee5c86423
\newcommand{\MlstMlpCmpQcdCMidOverFullPr}{5.5\% more errors}  % experiments/FIGS/data/v2/index.json :: nonlinear probe qcd_c_vs_light: (1 - AUC) of R42_Q1 over L188, run by run :: sha256 c6c4138ee5c86423 :: prose twin: two significant figures, direction in words
\newcommand{\MlstMlpCmpQcdCCoarseOverFull}{1.274\,\ensuremath{\pm}\,0.008}  % experiments/FIGS/data/v2/index.json :: nonlinear probe qcd_c_vs_light: (1 - AUC) of R16_Q1 over L188, run by run :: sha256 c6c4138ee5c86423
\newcommand{\MlstMlpCmpQcdCCoarseOverFullPr}{27\% more errors}  % experiments/FIGS/data/v2/index.json :: nonlinear probe qcd_c_vs_light: (1 - AUC) of R16_Q1 over L188, run by run :: sha256 c6c4138ee5c86423 :: prose twin: two significant figures, direction in words
\newcommand{\MlstMlpCmpQcdCResMassOverRes}{\pending{no run yet}}  % experiments/FIGS/data/v2/index.json :: nonlinear probe qcd_c_vs_light: (1 - AUC) of L162_MASS over L162, run by run :: sha256 c6c4138ee5c86423
\newcommand{\MlstMlpCmpQcdCResMassOverResPr}{\pending{no run yet}}  % experiments/FIGS/data/v2/index.json :: nonlinear probe qcd_c_vs_light: (1 - AUC) of L162_MASS over L162, run by run :: sha256 c6c4138ee5c86423 :: prose twin: two significant figures, direction in words
\newcommand{\MlstMlpCmpQcdCResMassOverResPrMax}{\pending{no run yet}}  % experiments/FIGS/data/v2/index.json :: nonlinear probe qcd_c_vs_light: (1 - AUC) of L162_MASS over L162, run by run :: sha256 c6c4138ee5c86423 :: prose twin of the run farther from 1
\newcommand{\MlstMlpCmpQcdCCoarseMassOverCoarse}{\pending{no run yet}}  % experiments/FIGS/data/v2/index.json :: nonlinear probe qcd_c_vs_light: (1 - AUC) of R16_Q1_MASS over R16_Q1, run by run :: sha256 c6c4138ee5c86423
\newcommand{\MlstMlpCmpQcdCCoarseMassOverCoarsePr}{\pending{no run yet}}  % experiments/FIGS/data/v2/index.json :: nonlinear probe qcd_c_vs_light: (1 - AUC) of R16_Q1_MASS over R16_Q1, run by run :: sha256 c6c4138ee5c86423 :: prose twin: two significant figures, direction in words
\newcommand{\MlstMlpCmpQcdCCoarseMassOverCoarsePrMax}{\pending{no run yet}}  % experiments/FIGS/data/v2/index.json :: nonlinear probe qcd_c_vs_light: (1 - AUC) of R16_Q1_MASS over R16_Q1, run by run :: sha256 c6c4138ee5c86423 :: prose twin of the run farther from 1
\newcommand{\MlstMlpCmpQcdCFzeroOverFull}{1.234 and 1.211}  % experiments/FIGS/data/v2/index.json :: nonlinear probe qcd_c_vs_light: (1 - AUC) of FLAV_F0 over L188, run by run :: sha256 c6c4138ee5c86423
\newcommand{\MlstMlpCmpQcdCFzeroOverFullPr}{at least 21\% more errors}  % experiments/FIGS/data/v2/index.json :: nonlinear probe qcd_c_vs_light: (1 - AUC) of FLAV_F0 over L188, run by run :: sha256 c6c4138ee5c86423 :: prose twin: two significant figures, direction in words
\newcommand{\MlstMlpCmpQcdCFzeroOverFullPrMax}{up to 23\% more errors}  % experiments/FIGS/data/v2/index.json :: nonlinear probe qcd_c_vs_light: (1 - AUC) of FLAV_F0 over L188, run by run :: sha256 c6c4138ee5c86423 :: prose twin of the run farther from 1
\newcommand{\MlstMlpCmpQcdCFoneOverFzero}{0.9395 and 0.9472}  % experiments/FIGS/data/v2/index.json :: nonlinear probe qcd_c_vs_light: (1 - AUC) of FLAV_F1 over FLAV_F0, run by run :: sha256 c6c4138ee5c86423
\newcommand{\MlstMlpCmpQcdCFoneOverFzeroPr}{at least 5.3\% fewer errors}  % experiments/FIGS/data/v2/index.json :: nonlinear probe qcd_c_vs_light: (1 - AUC) of FLAV_F1 over FLAV_F0, run by run :: sha256 c6c4138ee5c86423 :: prose twin: two significant figures, direction in words
\newcommand{\MlstMlpCmpQcdCFoneOverFzeroPrMax}{up to 6.0\% fewer errors}  % experiments/FIGS/data/v2/index.json :: nonlinear probe qcd_c_vs_light: (1 - AUC) of FLAV_F1 over FLAV_F0, run by run :: sha256 c6c4138ee5c86423 :: prose twin of the run farther from 1
\newcommand{\MlstMlpCmpQcdCFonerOverFone}{0.9846 and 0.9767}  % experiments/FIGS/data/v2/index.json :: nonlinear probe qcd_c_vs_light: (1 - AUC) of FLAV_F1R over FLAV_F1, run by run :: sha256 c6c4138ee5c86423
\newcommand{\MlstMlpCmpQcdCFonerOverFonePr}{at least 1.5\% fewer errors}  % experiments/FIGS/data/v2/index.json :: nonlinear probe qcd_c_vs_light: (1 - AUC) of FLAV_F1R over FLAV_F1, run by run :: sha256 c6c4138ee5c86423 :: prose twin: two significant figures, direction in words
\newcommand{\MlstMlpCmpQcdCFonerOverFonePrMax}{up to 2.3\% fewer errors}  % experiments/FIGS/data/v2/index.json :: nonlinear probe qcd_c_vs_light: (1 - AUC) of FLAV_F1R over FLAV_F1, run by run :: sha256 c6c4138ee5c86423 :: prose twin of the run farther from 1
\newcommand{\MlstMlpCmpQcdCFonerOverFzero}{0.9251 and 0.9251}  % experiments/FIGS/data/v2/index.json :: nonlinear probe qcd_c_vs_light: (1 - AUC) of FLAV_F1R over FLAV_F0, run by run :: sha256 c6c4138ee5c86423
\newcommand{\MlstMlpCmpQcdCFonerOverFzeroPr}{at least 7.5\% fewer errors}  % experiments/FIGS/data/v2/index.json :: nonlinear probe qcd_c_vs_light: (1 - AUC) of FLAV_F1R over FLAV_F0, run by run :: sha256 c6c4138ee5c86423 :: prose twin: two significant figures, direction in words
\newcommand{\MlstMlpCmpQcdCFonerOverFzeroPrMax}{up to 7.5\% fewer errors}  % experiments/FIGS/data/v2/index.json :: nonlinear probe qcd_c_vs_light: (1 - AUC) of FLAV_F1R over FLAV_F0, run by run :: sha256 c6c4138ee5c86423 :: prose twin of the run farther from 1
\newcommand{\MlstMlpCmpQcdCFzeroOverCoarse}{0.968 and 0.956}  % experiments/FIGS/data/v2/index.json :: nonlinear probe qcd_c_vs_light: (1 - AUC) of FLAV_F0 over R16_Q1, run by run :: sha256 c6c4138ee5c86423
\newcommand{\MlstMlpCmpQcdCFzeroOverCoarsePr}{at least 3.2\% fewer errors}  % experiments/FIGS/data/v2/index.json :: nonlinear probe qcd_c_vs_light: (1 - AUC) of FLAV_F0 over R16_Q1, run by run :: sha256 c6c4138ee5c86423 :: prose twin: two significant figures, direction in words
\newcommand{\MlstMlpCmpQcdCFzeroOverCoarsePrMax}{up to 4.4\% fewer errors}  % experiments/FIGS/data/v2/index.json :: nonlinear probe qcd_c_vs_light: (1 - AUC) of FLAV_F0 over R16_Q1, run by run :: sha256 c6c4138ee5c86423 :: prose twin of the run farther from 1
\newcommand{\MlstMlpCmpQcdCRandOverFull}{\pending{no run yet}}  % experiments/FIGS/data/v2/index.json :: nonlinear probe qcd_c_vs_light: (1 - AUC) of each random partition over the 188-class run of its seed :: sha256 c6c4138ee5c86423
\newcommand{\MlstMlpCmpQcdCRandOverFullPr}{\pending{no run yet}}  % experiments/FIGS/data/v2/index.json :: nonlinear probe qcd_c_vs_light: (1 - AUC) of each random partition over the 188-class run of its seed :: sha256 c6c4138ee5c86423 :: prose twin: two significant figures, direction in words
\newcommand{\MlstMlpTopRthirtyFullOverInitBn}{1.00\,\ensuremath{\pm}\,0.31}  % experiments/FIGS/data/v2/index.json :: nonlinear probe top, missing inputs zero: R30 of L188 over the BatchNorm-recomputed random network of the same index :: sha256 c6c4138ee5c86423
\newcommand{\MlstMlpTopRthirtyFullOverInitBnPr}{equal within 40\%}  % experiments/FIGS/data/v2/index.json :: nonlinear probe top, missing inputs zero: R30 of L188 over the BatchNorm-recomputed random network of the same index :: sha256 c6c4138ee5c86423 :: prose twin: two significant figures, direction in words
\newcommand{\MlstMlpTopRthirtyResOverInitBn}{0.94\,\ensuremath{\pm}\,0.26}  % experiments/FIGS/data/v2/index.json :: nonlinear probe top, missing inputs zero: R30 of L162 over the BatchNorm-recomputed random network of the same index :: sha256 c6c4138ee5c86423
\newcommand{\MlstMlpTopRthirtyResOverInitBnPr}{equal within 30\%}  % experiments/FIGS/data/v2/index.json :: nonlinear probe top, missing inputs zero: R30 of L162 over the BatchNorm-recomputed random network of the same index :: sha256 c6c4138ee5c86423 :: prose twin: two significant figures, direction in words
\newcommand{\MlstMlpTopRthirtyMidOverInitBn}{0.84\,\ensuremath{\pm}\,0.17}  % experiments/FIGS/data/v2/index.json :: nonlinear probe top, missing inputs zero: R30 of R42_Q1 over the BatchNorm-recomputed random network of the same index :: sha256 c6c4138ee5c86423
\newcommand{\MlstMlpTopRthirtyMidOverInitBnPr}{equal within 20\%}  % experiments/FIGS/data/v2/index.json :: nonlinear probe top, missing inputs zero: R30 of R42_Q1 over the BatchNorm-recomputed random network of the same index :: sha256 c6c4138ee5c86423 :: prose twin: two significant figures, direction in words
\newcommand{\MlstMlpTopRthirtyCoarseOverInitBn}{0.75\,\ensuremath{\pm}\,0.05}  % experiments/FIGS/data/v2/index.json :: nonlinear probe top, missing inputs zero: R30 of R16_Q1 over the BatchNorm-recomputed random network of the same index :: sha256 c6c4138ee5c86423
\newcommand{\MlstMlpTopRthirtyCoarseOverInitBnPr}{25\% less rejection}  % experiments/FIGS/data/v2/index.json :: nonlinear probe top, missing inputs zero: R30 of R16_Q1 over the BatchNorm-recomputed random network of the same index :: sha256 c6c4138ee5c86423 :: prose twin: two significant figures, direction in words
\newcommand{\MlstMlpQgRthirtyFullOverInitBn}{1.083\,\ensuremath{\pm}\,0.026}  % experiments/FIGS/data/v2/index.json :: nonlinear probe qg, missing inputs zero: R30 of L188 over the BatchNorm-recomputed random network of the same index :: sha256 c6c4138ee5c86423
\newcommand{\MlstMlpQgRthirtyFullOverInitBnPr}{8.3\% more rejection}  % experiments/FIGS/data/v2/index.json :: nonlinear probe qg, missing inputs zero: R30 of L188 over the BatchNorm-recomputed random network of the same index :: sha256 c6c4138ee5c86423 :: prose twin: two significant figures, direction in words
\newcommand{\MlstMlpQgRthirtyResOverInitBn}{1.010\,\ensuremath{\pm}\,0.016}  % experiments/FIGS/data/v2/index.json :: nonlinear probe qg, missing inputs zero: R30 of L162 over the BatchNorm-recomputed random network of the same index :: sha256 c6c4138ee5c86423
\newcommand{\MlstMlpQgRthirtyResOverInitBnPr}{equal within 2\%}  % experiments/FIGS/data/v2/index.json :: nonlinear probe qg, missing inputs zero: R30 of L162 over the BatchNorm-recomputed random network of the same index :: sha256 c6c4138ee5c86423 :: prose twin: two significant figures, direction in words
\newcommand{\MlstMlpQgRthirtyMidOverInitBn}{1.06\,\ensuremath{\pm}\,0.05}  % experiments/FIGS/data/v2/index.json :: nonlinear probe qg, missing inputs zero: R30 of R42_Q1 over the BatchNorm-recomputed random network of the same index :: sha256 c6c4138ee5c86423
\newcommand{\MlstMlpQgRthirtyMidOverInitBnPr}{equal within 7\%}  % experiments/FIGS/data/v2/index.json :: nonlinear probe qg, missing inputs zero: R30 of R42_Q1 over the BatchNorm-recomputed random network of the same index :: sha256 c6c4138ee5c86423 :: prose twin: two significant figures, direction in words
\newcommand{\MlstMlpQgRthirtyCoarseOverInitBn}{1.04\,\ensuremath{\pm}\,0.04}  % experiments/FIGS/data/v2/index.json :: nonlinear probe qg, missing inputs zero: R30 of R16_Q1 over the BatchNorm-recomputed random network of the same index :: sha256 c6c4138ee5c86423
\newcommand{\MlstMlpQgRthirtyCoarseOverInitBnPr}{equal within 5\%}  % experiments/FIGS/data/v2/index.json :: nonlinear probe qg, missing inputs zero: R30 of R16_Q1 over the BatchNorm-recomputed random network of the same index :: sha256 c6c4138ee5c86423 :: prose twin: two significant figures, direction in words
\newcommand{\MlstMlpQgHerwigRthirtyFullOverInitBn}{0.971\,\ensuremath{\pm}\,0.034}  % experiments/FIGS/data/v2/index.json :: nonlinear probe qg_herwig, missing inputs zero: R30 of L188 over the BatchNorm-recomputed random network of the same index :: sha256 c6c4138ee5c86423
\newcommand{\MlstMlpQgHerwigRthirtyFullOverInitBnPr}{equal within 4\%}  % experiments/FIGS/data/v2/index.json :: nonlinear probe qg_herwig, missing inputs zero: R30 of L188 over the BatchNorm-recomputed random network of the same index :: sha256 c6c4138ee5c86423 :: prose twin: two significant figures, direction in words
\newcommand{\MlstMlpQgHerwigRthirtyResOverInitBn}{0.94\,\ensuremath{\pm}\,0.05}  % experiments/FIGS/data/v2/index.json :: nonlinear probe qg_herwig, missing inputs zero: R30 of L162 over the BatchNorm-recomputed random network of the same index :: sha256 c6c4138ee5c86423
\newcommand{\MlstMlpQgHerwigRthirtyResOverInitBnPr}{equal within 6\%}  % experiments/FIGS/data/v2/index.json :: nonlinear probe qg_herwig, missing inputs zero: R30 of L162 over the BatchNorm-recomputed random network of the same index :: sha256 c6c4138ee5c86423 :: prose twin: two significant figures, direction in words
\newcommand{\MlstMlpQgHerwigRthirtyMidOverInitBn}{0.945\,\ensuremath{\pm}\,0.030}  % experiments/FIGS/data/v2/index.json :: nonlinear probe qg_herwig, missing inputs zero: R30 of R42_Q1 over the BatchNorm-recomputed random network of the same index :: sha256 c6c4138ee5c86423
\newcommand{\MlstMlpQgHerwigRthirtyMidOverInitBnPr}{equal within 6\%}  % experiments/FIGS/data/v2/index.json :: nonlinear probe qg_herwig, missing inputs zero: R30 of R42_Q1 over the BatchNorm-recomputed random network of the same index :: sha256 c6c4138ee5c86423 :: prose twin: two significant figures, direction in words
\newcommand{\MlstMlpQgHerwigRthirtyCoarseOverInitBn}{0.96\,\ensuremath{\pm}\,0.04}  % experiments/FIGS/data/v2/index.json :: nonlinear probe qg_herwig, missing inputs zero: R30 of R16_Q1 over the BatchNorm-recomputed random network of the same index :: sha256 c6c4138ee5c86423
\newcommand{\MlstMlpQgHerwigRthirtyCoarseOverInitBnPr}{equal within 5\%}  % experiments/FIGS/data/v2/index.json :: nonlinear probe qg_herwig, missing inputs zero: R30 of R16_Q1 over the BatchNorm-recomputed random network of the same index :: sha256 c6c4138ee5c86423 :: prose twin: two significant figures, direction in words
\newcommand{\MlstMlpNatJcTwoResMassOverRes}{1.016 and 1.009}  % experiments/FIGS/data/v2/index.json :: nonlinear probe, natural-composition JetClass-II (fits_v3_mlp2): (1 - AUC) of L162_MASS over L162, run by run :: sha256 c6c4138ee5c86423
\newcommand{\MlstMlpNatJcTwoResMassOverResPr}{equal within 2\%}  % experiments/FIGS/data/v2/index.json :: nonlinear probe, natural-composition JetClass-II (fits_v3_mlp2): (1 - AUC) of L162_MASS over L162, run by run :: sha256 c6c4138ee5c86423 :: prose twin: two significant figures, direction in words
\newcommand{\MlstMlpNatJcTwoResMassOverResPrMax}{equal within 2\%}  % experiments/FIGS/data/v2/index.json :: nonlinear probe, natural-composition JetClass-II (fits_v3_mlp2): (1 - AUC) of L162_MASS over L162, run by run :: sha256 c6c4138ee5c86423 :: prose twin of the run farther from 1
\newcommand{\MlstMlpNatJcTwoCoarseMassOverCoarse}{1.026 and 1.019}  % experiments/FIGS/data/v2/index.json :: nonlinear probe, natural-composition JetClass-II (fits_v3_mlp2): (1 - AUC) of R16_Q1_MASS over R16_Q1, run by run :: sha256 c6c4138ee5c86423
\newcommand{\MlstMlpNatJcTwoCoarseMassOverCoarsePr}{equal within 3\%}  % experiments/FIGS/data/v2/index.json :: nonlinear probe, natural-composition JetClass-II (fits_v3_mlp2): (1 - AUC) of R16_Q1_MASS over R16_Q1, run by run :: sha256 c6c4138ee5c86423 :: prose twin: two significant figures, direction in words
\newcommand{\MlstMlpNatJcTwoCoarseMassOverCoarsePrMax}{equal within 3\%}  % experiments/FIGS/data/v2/index.json :: nonlinear probe, natural-composition JetClass-II (fits_v3_mlp2): (1 - AUC) of R16_Q1_MASS over R16_Q1, run by run :: sha256 c6c4138ee5c86423 :: prose twin of the run farther from 1
\newcommand{\LpSubstitutionTop}{random draws}  % experiments/FIGS/data/v2/index.json :: linprobe/fits_v4_fill: lowest summed validation cross-entropy over every model, largest training set :: sha256 c6c4138ee5c86423 :: zero 8.8428; mean 8.99394; draws 7.00838
\newcommand{\LpSubstitutionQg}{the training mean}  % experiments/FIGS/data/v2/index.json :: linprobe/fits_v4_fill: lowest summed validation cross-entropy over every model, largest training set :: sha256 c6c4138ee5c86423 :: zero 13.6224; mean 13.1964; draws 13.3584
\newcommand{\MlstLpTopRthirtyFullOverInitBn}{4.67\,\ensuremath{\pm}\,0.23}  % experiments/FIGS/data/v2/index.json :: linear probe top, largest set, chosen substitution: R30 of L188 over the BatchNorm-recomputed random network of the same index :: sha256 c6c4138ee5c86423
\newcommand{\MlstLpTopRthirtyFullOverInitBnPr}{4.7 times the rejection}  % experiments/FIGS/data/v2/index.json :: linear probe top, largest set, chosen substitution: R30 of L188 over the BatchNorm-recomputed random network of the same index :: sha256 c6c4138ee5c86423 :: prose twin: two significant figures, direction in words
\newcommand{\MlstLpTopRthirtyFullOverInitBnZero}{1.04\,\ensuremath{\pm}\,0.05}  % experiments/FIGS/data/v2/index.json :: linear probe top, largest set, zero: R30 of L188 over the BatchNorm-recomputed random network of the same index :: sha256 c6c4138ee5c86423
\newcommand{\MlstLpTopRthirtyFullOverInitBnZeroPr}{equal within 5\%}  % experiments/FIGS/data/v2/index.json :: linear probe top, largest set, zero: R30 of L188 over the BatchNorm-recomputed random network of the same index :: sha256 c6c4138ee5c86423 :: prose twin: two significant figures, direction in words
\newcommand{\MlstLpTopRthirtyResOverInitBn}{4.08\,\ensuremath{\pm}\,0.28}  % experiments/FIGS/data/v2/index.json :: linear probe top, largest set, chosen substitution: R30 of L162 over the BatchNorm-recomputed random network of the same index :: sha256 c6c4138ee5c86423
\newcommand{\MlstLpTopRthirtyResOverInitBnPr}{4.1 times the rejection}  % experiments/FIGS/data/v2/index.json :: linear probe top, largest set, chosen substitution: R30 of L162 over the BatchNorm-recomputed random network of the same index :: sha256 c6c4138ee5c86423 :: prose twin: two significant figures, direction in words
\newcommand{\MlstLpTopRthirtyResOverInitBnZero}{1.33\,\ensuremath{\pm}\,0.30}  % experiments/FIGS/data/v2/index.json :: linear probe top, largest set, zero: R30 of L162 over the BatchNorm-recomputed random network of the same index :: sha256 c6c4138ee5c86423
\newcommand{\MlstLpTopRthirtyResOverInitBnZeroPr}{equal within 40\%}  % experiments/FIGS/data/v2/index.json :: linear probe top, largest set, zero: R30 of L162 over the BatchNorm-recomputed random network of the same index :: sha256 c6c4138ee5c86423 :: prose twin: two significant figures, direction in words
\newcommand{\MlstLpTopRthirtyMidOverInitBn}{4.76\,\ensuremath{\pm}\,0.15}  % experiments/FIGS/data/v2/index.json :: linear probe top, largest set, chosen substitution: R30 of R42_Q1 over the BatchNorm-recomputed random network of the same index :: sha256 c6c4138ee5c86423
\newcommand{\MlstLpTopRthirtyMidOverInitBnPr}{4.8 times the rejection}  % experiments/FIGS/data/v2/index.json :: linear probe top, largest set, chosen substitution: R30 of R42_Q1 over the BatchNorm-recomputed random network of the same index :: sha256 c6c4138ee5c86423 :: prose twin: two significant figures, direction in words
\newcommand{\MlstLpTopRthirtyMidOverInitBnZero}{1.01\,\ensuremath{\pm}\,0.28}  % experiments/FIGS/data/v2/index.json :: linear probe top, largest set, zero: R30 of R42_Q1 over the BatchNorm-recomputed random network of the same index :: sha256 c6c4138ee5c86423
\newcommand{\MlstLpTopRthirtyMidOverInitBnZeroPr}{equal within 30\%}  % experiments/FIGS/data/v2/index.json :: linear probe top, largest set, zero: R30 of R42_Q1 over the BatchNorm-recomputed random network of the same index :: sha256 c6c4138ee5c86423 :: prose twin: two significant figures, direction in words
\newcommand{\MlstLpTopRthirtyCoarseOverInitBn}{5.24\,\ensuremath{\pm}\,0.06}  % experiments/FIGS/data/v2/index.json :: linear probe top, largest set, chosen substitution: R30 of R16_Q1 over the BatchNorm-recomputed random network of the same index :: sha256 c6c4138ee5c86423
\newcommand{\MlstLpTopRthirtyCoarseOverInitBnPr}{5.2 times the rejection}  % experiments/FIGS/data/v2/index.json :: linear probe top, largest set, chosen substitution: R30 of R16_Q1 over the BatchNorm-recomputed random network of the same index :: sha256 c6c4138ee5c86423 :: prose twin: two significant figures, direction in words
\newcommand{\MlstLpTopRthirtyCoarseOverInitBnZero}{0.97\,\ensuremath{\pm}\,0.16}  % experiments/FIGS/data/v2/index.json :: linear probe top, largest set, zero: R30 of R16_Q1 over the BatchNorm-recomputed random network of the same index :: sha256 c6c4138ee5c86423
\newcommand{\MlstLpTopRthirtyCoarseOverInitBnZeroPr}{equal within 20\%}  % experiments/FIGS/data/v2/index.json :: linear probe top, largest set, zero: R30 of R16_Q1 over the BatchNorm-recomputed random network of the same index :: sha256 c6c4138ee5c86423 :: prose twin: two significant figures, direction in words
\newcommand{\MlstLpTopRthirtyResMassOverInitBn}{5.71 and 6.25}  % experiments/FIGS/data/v2/index.json :: linear probe top, largest set, chosen substitution: R30 of L162_MASS over the BatchNorm-recomputed random network of the same index :: sha256 c6c4138ee5c86423
\newcommand{\MlstLpTopRthirtyResMassOverInitBnPr}{at least 5.7 times the rejection}  % experiments/FIGS/data/v2/index.json :: linear probe top, largest set, chosen substitution: R30 of L162_MASS over the BatchNorm-recomputed random network of the same index :: sha256 c6c4138ee5c86423 :: prose twin: two significant figures, direction in words
\newcommand{\MlstLpTopRthirtyResMassOverInitBnPrMax}{up to 6.3 times the rejection}  % experiments/FIGS/data/v2/index.json :: linear probe top, largest set, chosen substitution: R30 of L162_MASS over the BatchNorm-recomputed random network of the same index :: sha256 c6c4138ee5c86423 :: prose twin of the run farther from 1
\newcommand{\MlstLpTopRthirtyResMassOverInitBnZero}{1.683 and 1.678}  % experiments/FIGS/data/v2/index.json :: linear probe top, largest set, zero: R30 of L162_MASS over the BatchNorm-recomputed random network of the same index :: sha256 c6c4138ee5c86423
\newcommand{\MlstLpTopRthirtyResMassOverInitBnZeroPr}{at least 68\% more rejection}  % experiments/FIGS/data/v2/index.json :: linear probe top, largest set, zero: R30 of L162_MASS over the BatchNorm-recomputed random network of the same index :: sha256 c6c4138ee5c86423 :: prose twin: two significant figures, direction in words
\newcommand{\MlstLpTopRthirtyResMassOverInitBnZeroPrMax}{up to 68\% more rejection}  % experiments/FIGS/data/v2/index.json :: linear probe top, largest set, zero: R30 of L162_MASS over the BatchNorm-recomputed random network of the same index :: sha256 c6c4138ee5c86423 :: prose twin of the run farther from 1
\newcommand{\MlstLpTopRthirtyCoarseMassOverInitBn}{6.73 and 7.72}  % experiments/FIGS/data/v2/index.json :: linear probe top, largest set, chosen substitution: R30 of R16_Q1_MASS over the BatchNorm-recomputed random network of the same index :: sha256 c6c4138ee5c86423
\newcommand{\MlstLpTopRthirtyCoarseMassOverInitBnPr}{at least 6.7 times the rejection}  % experiments/FIGS/data/v2/index.json :: linear probe top, largest set, chosen substitution: R30 of R16_Q1_MASS over the BatchNorm-recomputed random network of the same index :: sha256 c6c4138ee5c86423 :: prose twin: two significant figures, direction in words
\newcommand{\MlstLpTopRthirtyCoarseMassOverInitBnPrMax}{up to 7.7 times the rejection}  % experiments/FIGS/data/v2/index.json :: linear probe top, largest set, chosen substitution: R30 of R16_Q1_MASS over the BatchNorm-recomputed random network of the same index :: sha256 c6c4138ee5c86423 :: prose twin of the run farther from 1
\newcommand{\MlstLpTopRthirtyCoarseMassOverInitBnZero}{2.3 and 1.2}  % experiments/FIGS/data/v2/index.json :: linear probe top, largest set, zero: R30 of R16_Q1_MASS over the BatchNorm-recomputed random network of the same index :: sha256 c6c4138ee5c86423
\newcommand{\MlstLpTopRthirtyCoarseMassOverInitBnZeroPr}{equal within 200\%}  % experiments/FIGS/data/v2/index.json :: linear probe top, largest set, zero: R30 of R16_Q1_MASS over the BatchNorm-recomputed random network of the same index :: sha256 c6c4138ee5c86423 :: prose twin: two significant figures, direction in words
\newcommand{\MlstLpTopRthirtyCoarseMassOverInitBnZeroPrMax}{equal within 200\%}  % experiments/FIGS/data/v2/index.json :: linear probe top, largest set, zero: R30 of R16_Q1_MASS over the BatchNorm-recomputed random network of the same index :: sha256 c6c4138ee5c86423 :: prose twin of the run farther from 1
\newcommand{\MlstLpTopRthirtyResOverFull}{0.874\,\ensuremath{\pm}\,0.018}  % experiments/FIGS/data/v2/index.json :: linear probe top: R30 of L162 over the 188-class run :: sha256 c6c4138ee5c86423
\newcommand{\MlstLpTopRthirtyResOverFullPr}{13\% less rejection}  % experiments/FIGS/data/v2/index.json :: linear probe top: R30 of L162 over the 188-class run :: sha256 c6c4138ee5c86423 :: prose twin: two significant figures, direction in words
\newcommand{\MlstLpTopRthirtyResOverFullZero}{1.29\,\ensuremath{\pm}\,0.31}  % experiments/FIGS/data/v2/index.json :: linear probe top: R30 of L162 over the 188-class run :: sha256 c6c4138ee5c86423
\newcommand{\MlstLpTopRthirtyResOverFullZeroPr}{equal within 40\%}  % experiments/FIGS/data/v2/index.json :: linear probe top: R30 of L162 over the 188-class run :: sha256 c6c4138ee5c86423 :: prose twin: two significant figures, direction in words
\newcommand{\MlstLpTopRthirtyMidOverFull}{1.02\,\ensuremath{\pm}\,0.08}  % experiments/FIGS/data/v2/index.json :: linear probe top: R30 of R42_Q1 over the 188-class run :: sha256 c6c4138ee5c86423
\newcommand{\MlstLpTopRthirtyMidOverFullPr}{equal within 9\%}  % experiments/FIGS/data/v2/index.json :: linear probe top: R30 of R42_Q1 over the 188-class run :: sha256 c6c4138ee5c86423 :: prose twin: two significant figures, direction in words
\newcommand{\MlstLpTopRthirtyMidOverFullZero}{0.98\,\ensuremath{\pm}\,0.28}  % experiments/FIGS/data/v2/index.json :: linear probe top: R30 of R42_Q1 over the 188-class run :: sha256 c6c4138ee5c86423
\newcommand{\MlstLpTopRthirtyMidOverFullZeroPr}{equal within 30\%}  % experiments/FIGS/data/v2/index.json :: linear probe top: R30 of R42_Q1 over the 188-class run :: sha256 c6c4138ee5c86423 :: prose twin: two significant figures, direction in words
\newcommand{\MlstLpTopRthirtyCoarseOverFull}{1.12\,\ensuremath{\pm}\,0.06}  % experiments/FIGS/data/v2/index.json :: linear probe top: R30 of R16_Q1 over the 188-class run :: sha256 c6c4138ee5c86423
\newcommand{\MlstLpTopRthirtyCoarseOverFullPr}{12\% more rejection}  % experiments/FIGS/data/v2/index.json :: linear probe top: R30 of R16_Q1 over the 188-class run :: sha256 c6c4138ee5c86423 :: prose twin: two significant figures, direction in words
\newcommand{\MlstLpTopRthirtyCoarseOverFullZero}{0.93\,\ensuremath{\pm}\,0.15}  % experiments/FIGS/data/v2/index.json :: linear probe top: R30 of R16_Q1 over the 188-class run :: sha256 c6c4138ee5c86423
\newcommand{\MlstLpTopRthirtyCoarseOverFullZeroPr}{equal within 20\%}  % experiments/FIGS/data/v2/index.json :: linear probe top: R30 of R16_Q1 over the 188-class run :: sha256 c6c4138ee5c86423 :: prose twin: two significant figures, direction in words
\newcommand{\MlstLpTopRthirtyMeanOverZeroFull}{1.25\,\ensuremath{\pm}\,0.27}  % experiments/FIGS/data/v2/index.json :: linear probe top: R30 with mean over R30 with zero, L188 :: sha256 c6c4138ee5c86423
\newcommand{\MlstLpTopRthirtyMeanOverZeroFullPr}{equal within 30\%}  % experiments/FIGS/data/v2/index.json :: linear probe top: R30 with mean over R30 with zero, L188 :: sha256 c6c4138ee5c86423 :: prose twin: two significant figures, direction in words
\newcommand{\MlstLpTopRthirtyDrawOverZeroFull}{3.9\,\ensuremath{\pm}\,0.4}  % experiments/FIGS/data/v2/index.json :: linear probe top: R30 with draws over R30 with zero, L188 :: sha256 c6c4138ee5c86423
\newcommand{\MlstLpTopRthirtyDrawOverZeroFullPr}{3.9 times the rejection}  % experiments/FIGS/data/v2/index.json :: linear probe top: R30 with draws over R30 with zero, L188 :: sha256 c6c4138ee5c86423 :: prose twin: two significant figures, direction in words
\newcommand{\MlstLpTopRthirtyMeanOverZeroInitBn}{1.08\,\ensuremath{\pm}\,0.20}  % experiments/FIGS/data/v2/index.json :: linear probe top: R30 with mean over R30 with zero, INIT_BN :: sha256 c6c4138ee5c86423
\newcommand{\MlstLpTopRthirtyMeanOverZeroInitBnPr}{equal within 30\%}  % experiments/FIGS/data/v2/index.json :: linear probe top: R30 with mean over R30 with zero, INIT_BN :: sha256 c6c4138ee5c86423 :: prose twin: two significant figures, direction in words
\newcommand{\MlstLpTopRthirtyDrawOverZeroInitBn}{0.86\,\ensuremath{\pm}\,0.14}  % experiments/FIGS/data/v2/index.json :: linear probe top: R30 with draws over R30 with zero, INIT_BN :: sha256 c6c4138ee5c86423
\newcommand{\MlstLpTopRthirtyDrawOverZeroInitBnPr}{equal within 20\%}  % experiments/FIGS/data/v2/index.json :: linear probe top: R30 with draws over R30 with zero, INIT_BN :: sha256 c6c4138ee5c86423 :: prose twin: two significant figures, direction in words
\newcommand{\MlstLpTopRthirtyKinOverFull}{3.88}  % experiments/FIGS/data/v2/index.json :: linear probe top: R30 of the kinematic-only model over the mean of the 188-class runs (zero) :: sha256 c6c4138ee5c86423
\newcommand{\MlstLpTopRthirtyKinOverFullPr}{3.9 times the rejection}  % experiments/FIGS/data/v2/index.json :: linear probe top: R30 of the kinematic-only model over the mean of the 188-class runs (zero) :: sha256 c6c4138ee5c86423 :: prose twin: two significant figures, direction in words
\newcommand{\MlstLpTopRthirtyKin}{231}  % experiments/FIGS/data/v2/index.json :: linear probe top, largest set: R30 of the kinematic-only model (one run) :: sha256 c6c4138ee5c86423
\newcommand{\MlstLpTopRthirtyFullZero}{60\,\ensuremath{\pm}\,12}  % experiments/FIGS/data/v2/index.json :: linear probe top, largest set: R30 with zero, L188 :: sha256 c6c4138ee5c86423
\newcommand{\MlstLpTopRthirtyFullMean}{73\,\ensuremath{\pm}\,15}  % experiments/FIGS/data/v2/index.json :: linear probe top, largest set: R30 with mean, L188 :: sha256 c6c4138ee5c86423
\newcommand{\MlstLpTopRthirtyFullDraws}{228\,\ensuremath{\pm}\,25}  % experiments/FIGS/data/v2/index.json :: linear probe top, largest set: R30 with draws, L188 :: sha256 c6c4138ee5c86423
\newcommand{\MlstLpTopRthirtyInitBnZero}{58\,\ensuremath{\pm}\,13}  % experiments/FIGS/data/v2/index.json :: linear probe top, largest set: R30 with zero, INIT_BN :: sha256 c6c4138ee5c86423
\newcommand{\MlstLpTopRthirtyInitBnMean}{60.4\,\ensuremath{\pm}\,2.4}  % experiments/FIGS/data/v2/index.json :: linear probe top, largest set: R30 with mean, INIT_BN :: sha256 c6c4138ee5c86423
\newcommand{\MlstLpTopRthirtyInitBnDraws}{48.6\,\ensuremath{\pm}\,3.0}  % experiments/FIGS/data/v2/index.json :: linear probe top, largest set: R30 with draws, INIT_BN :: sha256 c6c4138ee5c86423
\newcommand{\MlstLpQgRthirtyFullOverInitBn}{1.39\,\ensuremath{\pm}\,0.06}  % experiments/FIGS/data/v2/index.json :: linear probe qg, largest set, chosen substitution: R30 of L188 over the BatchNorm-recomputed random network of the same index :: sha256 c6c4138ee5c86423
\newcommand{\MlstLpQgRthirtyFullOverInitBnPr}{39\% more rejection}  % experiments/FIGS/data/v2/index.json :: linear probe qg, largest set, chosen substitution: R30 of L188 over the BatchNorm-recomputed random network of the same index :: sha256 c6c4138ee5c86423 :: prose twin: two significant figures, direction in words
\newcommand{\MlstLpQgRthirtyFullOverInitBnZero}{1.21\,\ensuremath{\pm}\,0.15}  % experiments/FIGS/data/v2/index.json :: linear probe qg, largest set, zero: R30 of L188 over the BatchNorm-recomputed random network of the same index :: sha256 c6c4138ee5c86423
\newcommand{\MlstLpQgRthirtyFullOverInitBnZeroPr}{equal within 30\%}  % experiments/FIGS/data/v2/index.json :: linear probe qg, largest set, zero: R30 of L188 over the BatchNorm-recomputed random network of the same index :: sha256 c6c4138ee5c86423 :: prose twin: two significant figures, direction in words
\newcommand{\MlstLpQgRthirtyResOverInitBn}{1.360\,\ensuremath{\pm}\,0.025}  % experiments/FIGS/data/v2/index.json :: linear probe qg, largest set, chosen substitution: R30 of L162 over the BatchNorm-recomputed random network of the same index :: sha256 c6c4138ee5c86423
\newcommand{\MlstLpQgRthirtyResOverInitBnPr}{36\% more rejection}  % experiments/FIGS/data/v2/index.json :: linear probe qg, largest set, chosen substitution: R30 of L162 over the BatchNorm-recomputed random network of the same index :: sha256 c6c4138ee5c86423 :: prose twin: two significant figures, direction in words
\newcommand{\MlstLpQgRthirtyResOverInitBnZero}{1.07\,\ensuremath{\pm}\,0.13}  % experiments/FIGS/data/v2/index.json :: linear probe qg, largest set, zero: R30 of L162 over the BatchNorm-recomputed random network of the same index :: sha256 c6c4138ee5c86423
\newcommand{\MlstLpQgRthirtyResOverInitBnZeroPr}{equal within 20\%}  % experiments/FIGS/data/v2/index.json :: linear probe qg, largest set, zero: R30 of L162 over the BatchNorm-recomputed random network of the same index :: sha256 c6c4138ee5c86423 :: prose twin: two significant figures, direction in words
\newcommand{\MlstLpQgRthirtyMidOverInitBn}{1.357\,\ensuremath{\pm}\,0.033}  % experiments/FIGS/data/v2/index.json :: linear probe qg, largest set, chosen substitution: R30 of R42_Q1 over the BatchNorm-recomputed random network of the same index :: sha256 c6c4138ee5c86423
\newcommand{\MlstLpQgRthirtyMidOverInitBnPr}{36\% more rejection}  % experiments/FIGS/data/v2/index.json :: linear probe qg, largest set, chosen substitution: R30 of R42_Q1 over the BatchNorm-recomputed random network of the same index :: sha256 c6c4138ee5c86423 :: prose twin: two significant figures, direction in words
\newcommand{\MlstLpQgRthirtyMidOverInitBnZero}{1.12\,\ensuremath{\pm}\,0.11}  % experiments/FIGS/data/v2/index.json :: linear probe qg, largest set, zero: R30 of R42_Q1 over the BatchNorm-recomputed random network of the same index :: sha256 c6c4138ee5c86423
\newcommand{\MlstLpQgRthirtyMidOverInitBnZeroPr}{equal within 20\%}  % experiments/FIGS/data/v2/index.json :: linear probe qg, largest set, zero: R30 of R42_Q1 over the BatchNorm-recomputed random network of the same index :: sha256 c6c4138ee5c86423 :: prose twin: two significant figures, direction in words
\newcommand{\MlstLpQgRthirtyCoarseOverInitBn}{1.42\,\ensuremath{\pm}\,0.04}  % experiments/FIGS/data/v2/index.json :: linear probe qg, largest set, chosen substitution: R30 of R16_Q1 over the BatchNorm-recomputed random network of the same index :: sha256 c6c4138ee5c86423
\newcommand{\MlstLpQgRthirtyCoarseOverInitBnPr}{42\% more rejection}  % experiments/FIGS/data/v2/index.json :: linear probe qg, largest set, chosen substitution: R30 of R16_Q1 over the BatchNorm-recomputed random network of the same index :: sha256 c6c4138ee5c86423 :: prose twin: two significant figures, direction in words
\newcommand{\MlstLpQgRthirtyCoarseOverInitBnZero}{1.15\,\ensuremath{\pm}\,0.17}  % experiments/FIGS/data/v2/index.json :: linear probe qg, largest set, zero: R30 of R16_Q1 over the BatchNorm-recomputed random network of the same index :: sha256 c6c4138ee5c86423
\newcommand{\MlstLpQgRthirtyCoarseOverInitBnZeroPr}{equal within 20\%}  % experiments/FIGS/data/v2/index.json :: linear probe qg, largest set, zero: R30 of R16_Q1 over the BatchNorm-recomputed random network of the same index :: sha256 c6c4138ee5c86423 :: prose twin: two significant figures, direction in words
\newcommand{\MlstLpQgRthirtyResMassOverInitBn}{1.430 and 1.364}  % experiments/FIGS/data/v2/index.json :: linear probe qg, largest set, chosen substitution: R30 of L162_MASS over the BatchNorm-recomputed random network of the same index :: sha256 c6c4138ee5c86423
\newcommand{\MlstLpQgRthirtyResMassOverInitBnPr}{at least 36\% more rejection}  % experiments/FIGS/data/v2/index.json :: linear probe qg, largest set, chosen substitution: R30 of L162_MASS over the BatchNorm-recomputed random network of the same index :: sha256 c6c4138ee5c86423 :: prose twin: two significant figures, direction in words
\newcommand{\MlstLpQgRthirtyResMassOverInitBnPrMax}{up to 43\% more rejection}  % experiments/FIGS/data/v2/index.json :: linear probe qg, largest set, chosen substitution: R30 of L162_MASS over the BatchNorm-recomputed random network of the same index :: sha256 c6c4138ee5c86423 :: prose twin of the run farther from 1
\newcommand{\MlstLpQgRthirtyResMassOverInitBnZero}{1.213 and 1.175}  % experiments/FIGS/data/v2/index.json :: linear probe qg, largest set, zero: R30 of L162_MASS over the BatchNorm-recomputed random network of the same index :: sha256 c6c4138ee5c86423
\newcommand{\MlstLpQgRthirtyResMassOverInitBnZeroPr}{equal within 30\%}  % experiments/FIGS/data/v2/index.json :: linear probe qg, largest set, zero: R30 of L162_MASS over the BatchNorm-recomputed random network of the same index :: sha256 c6c4138ee5c86423 :: prose twin: two significant figures, direction in words
\newcommand{\MlstLpQgRthirtyResMassOverInitBnZeroPrMax}{equal within 30\%}  % experiments/FIGS/data/v2/index.json :: linear probe qg, largest set, zero: R30 of L162_MASS over the BatchNorm-recomputed random network of the same index :: sha256 c6c4138ee5c86423 :: prose twin of the run farther from 1
\newcommand{\MlstLpQgRthirtyCoarseMassOverInitBn}{1.324 and 1.394}  % experiments/FIGS/data/v2/index.json :: linear probe qg, largest set, chosen substitution: R30 of R16_Q1_MASS over the BatchNorm-recomputed random network of the same index :: sha256 c6c4138ee5c86423
\newcommand{\MlstLpQgRthirtyCoarseMassOverInitBnPr}{at least 32\% more rejection}  % experiments/FIGS/data/v2/index.json :: linear probe qg, largest set, chosen substitution: R30 of R16_Q1_MASS over the BatchNorm-recomputed random network of the same index :: sha256 c6c4138ee5c86423 :: prose twin: two significant figures, direction in words
\newcommand{\MlstLpQgRthirtyCoarseMassOverInitBnPrMax}{up to 39\% more rejection}  % experiments/FIGS/data/v2/index.json :: linear probe qg, largest set, chosen substitution: R30 of R16_Q1_MASS over the BatchNorm-recomputed random network of the same index :: sha256 c6c4138ee5c86423 :: prose twin of the run farther from 1
\newcommand{\MlstLpQgRthirtyCoarseMassOverInitBnZero}{1.224 and 1.261}  % experiments/FIGS/data/v2/index.json :: linear probe qg, largest set, zero: R30 of R16_Q1_MASS over the BatchNorm-recomputed random network of the same index :: sha256 c6c4138ee5c86423
\newcommand{\MlstLpQgRthirtyCoarseMassOverInitBnZeroPr}{equal within 30\%}  % experiments/FIGS/data/v2/index.json :: linear probe qg, largest set, zero: R30 of R16_Q1_MASS over the BatchNorm-recomputed random network of the same index :: sha256 c6c4138ee5c86423 :: prose twin: two significant figures, direction in words
\newcommand{\MlstLpQgRthirtyCoarseMassOverInitBnZeroPrMax}{equal within 30\%}  % experiments/FIGS/data/v2/index.json :: linear probe qg, largest set, zero: R30 of R16_Q1_MASS over the BatchNorm-recomputed random network of the same index :: sha256 c6c4138ee5c86423 :: prose twin of the run farther from 1
\newcommand{\MlstLpQgRthirtyResOverFull}{0.977\,\ensuremath{\pm}\,0.028}  % experiments/FIGS/data/v2/index.json :: linear probe qg: R30 of L162 over the 188-class run :: sha256 c6c4138ee5c86423
\newcommand{\MlstLpQgRthirtyResOverFullPr}{equal within 3\%}  % experiments/FIGS/data/v2/index.json :: linear probe qg: R30 of L162 over the 188-class run :: sha256 c6c4138ee5c86423 :: prose twin: two significant figures, direction in words
\newcommand{\MlstLpQgRthirtyResOverFullZero}{0.89\,\ensuremath{\pm}\,0.20}  % experiments/FIGS/data/v2/index.json :: linear probe qg: R30 of L162 over the 188-class run :: sha256 c6c4138ee5c86423
\newcommand{\MlstLpQgRthirtyResOverFullZeroPr}{equal within 20\%}  % experiments/FIGS/data/v2/index.json :: linear probe qg: R30 of L162 over the 188-class run :: sha256 c6c4138ee5c86423 :: prose twin: two significant figures, direction in words
\newcommand{\MlstLpQgRthirtyMidOverFull}{0.974\,\ensuremath{\pm}\,0.031}  % experiments/FIGS/data/v2/index.json :: linear probe qg: R30 of R42_Q1 over the 188-class run :: sha256 c6c4138ee5c86423
\newcommand{\MlstLpQgRthirtyMidOverFullPr}{equal within 4\%}  % experiments/FIGS/data/v2/index.json :: linear probe qg: R30 of R42_Q1 over the 188-class run :: sha256 c6c4138ee5c86423 :: prose twin: two significant figures, direction in words
\newcommand{\MlstLpQgRthirtyMidOverFullZero}{0.94\,\ensuremath{\pm}\,0.18}  % experiments/FIGS/data/v2/index.json :: linear probe qg: R30 of R42_Q1 over the 188-class run :: sha256 c6c4138ee5c86423
\newcommand{\MlstLpQgRthirtyMidOverFullZeroPr}{equal within 20\%}  % experiments/FIGS/data/v2/index.json :: linear probe qg: R30 of R42_Q1 over the 188-class run :: sha256 c6c4138ee5c86423 :: prose twin: two significant figures, direction in words
\newcommand{\MlstLpQgRthirtyCoarseOverFull}{1.02\,\ensuremath{\pm}\,0.06}  % experiments/FIGS/data/v2/index.json :: linear probe qg: R30 of R16_Q1 over the 188-class run :: sha256 c6c4138ee5c86423
\newcommand{\MlstLpQgRthirtyCoarseOverFullPr}{equal within 6\%}  % experiments/FIGS/data/v2/index.json :: linear probe qg: R30 of R16_Q1 over the 188-class run :: sha256 c6c4138ee5c86423 :: prose twin: two significant figures, direction in words
\newcommand{\MlstLpQgRthirtyCoarseOverFullZero}{0.95\,\ensuremath{\pm}\,0.09}  % experiments/FIGS/data/v2/index.json :: linear probe qg: R30 of R16_Q1 over the 188-class run :: sha256 c6c4138ee5c86423
\newcommand{\MlstLpQgRthirtyCoarseOverFullZeroPr}{equal within 9\%}  % experiments/FIGS/data/v2/index.json :: linear probe qg: R30 of R16_Q1 over the 188-class run :: sha256 c6c4138ee5c86423 :: prose twin: two significant figures, direction in words
\newcommand{\MlstLpQgRthirtyMeanOverZeroFull}{1.12\,\ensuremath{\pm}\,0.17}  % experiments/FIGS/data/v2/index.json :: linear probe qg: R30 with mean over R30 with zero, L188 :: sha256 c6c4138ee5c86423
\newcommand{\MlstLpQgRthirtyMeanOverZeroFullPr}{equal within 20\%}  % experiments/FIGS/data/v2/index.json :: linear probe qg: R30 with mean over R30 with zero, L188 :: sha256 c6c4138ee5c86423 :: prose twin: two significant figures, direction in words
\newcommand{\MlstLpQgRthirtyDrawOverZeroFull}{1.05\,\ensuremath{\pm}\,0.12}  % experiments/FIGS/data/v2/index.json :: linear probe qg: R30 with draws over R30 with zero, L188 :: sha256 c6c4138ee5c86423
\newcommand{\MlstLpQgRthirtyDrawOverZeroFullPr}{equal within 20\%}  % experiments/FIGS/data/v2/index.json :: linear probe qg: R30 with draws over R30 with zero, L188 :: sha256 c6c4138ee5c86423 :: prose twin: two significant figures, direction in words
\newcommand{\MlstLpQgRthirtyMeanOverZeroInitBn}{0.964\,\ensuremath{\pm}\,0.030}  % experiments/FIGS/data/v2/index.json :: linear probe qg: R30 with mean over R30 with zero, INIT_BN :: sha256 c6c4138ee5c86423
\newcommand{\MlstLpQgRthirtyMeanOverZeroInitBnPr}{equal within 4\%}  % experiments/FIGS/data/v2/index.json :: linear probe qg: R30 with mean over R30 with zero, INIT_BN :: sha256 c6c4138ee5c86423 :: prose twin: two significant figures, direction in words
\newcommand{\MlstLpQgRthirtyDrawOverZeroInitBn}{0.97\,\ensuremath{\pm}\,0.04}  % experiments/FIGS/data/v2/index.json :: linear probe qg: R30 with draws over R30 with zero, INIT_BN :: sha256 c6c4138ee5c86423
\newcommand{\MlstLpQgRthirtyDrawOverZeroInitBnPr}{equal within 4\%}  % experiments/FIGS/data/v2/index.json :: linear probe qg: R30 with draws over R30 with zero, INIT_BN :: sha256 c6c4138ee5c86423 :: prose twin: two significant figures, direction in words
\newcommand{\MlstLpQgRthirtyFullZero}{66\,\ensuremath{\pm}\,7}  % experiments/FIGS/data/v2/index.json :: linear probe qg, largest set: R30 with zero, L188 :: sha256 c6c4138ee5c86423
\newcommand{\MlstLpQgRthirtyFullMean}{72.7\,\ensuremath{\pm}\,2.4}  % experiments/FIGS/data/v2/index.json :: linear probe qg, largest set: R30 with mean, L188 :: sha256 c6c4138ee5c86423
\newcommand{\MlstLpQgRthirtyFullDraws}{68.1\,\ensuremath{\pm}\,1.8}  % experiments/FIGS/data/v2/index.json :: linear probe qg, largest set: R30 with draws, L188 :: sha256 c6c4138ee5c86423
\newcommand{\MlstLpQgRthirtyInitBnZero}{54.2\,\ensuremath{\pm}\,2.4}  % experiments/FIGS/data/v2/index.json :: linear probe qg, largest set: R30 with zero, INIT_BN :: sha256 c6c4138ee5c86423
\newcommand{\MlstLpQgRthirtyInitBnMean}{52.2\,\ensuremath{\pm}\,0.8}  % experiments/FIGS/data/v2/index.json :: linear probe qg, largest set: R30 with mean, INIT_BN :: sha256 c6c4138ee5c86423
\newcommand{\MlstLpQgRthirtyInitBnDraws}{52.4\,\ensuremath{\pm}\,0.4}  % experiments/FIGS/data/v2/index.json :: linear probe qg, largest set: R30 with draws, INIT_BN :: sha256 c6c4138ee5c86423
\newcommand{\MlstLpQgHerwigRthirtyFullOverInitBn}{1.14\,\ensuremath{\pm}\,0.10}  % experiments/FIGS/data/v2/index.json :: linear probe qg_herwig, largest set, chosen substitution: R30 of L188 over the BatchNorm-recomputed random network of the same index :: sha256 c6c4138ee5c86423
\newcommand{\MlstLpQgHerwigRthirtyFullOverInitBnPr}{equal within 20\%}  % experiments/FIGS/data/v2/index.json :: linear probe qg_herwig, largest set, chosen substitution: R30 of L188 over the BatchNorm-recomputed random network of the same index :: sha256 c6c4138ee5c86423 :: prose twin: two significant figures, direction in words
\newcommand{\MlstLpQgHerwigRthirtyFullOverInitBnZero}{1.054\,\ensuremath{\pm}\,0.028}  % experiments/FIGS/data/v2/index.json :: linear probe qg_herwig, largest set, zero: R30 of L188 over the BatchNorm-recomputed random network of the same index :: sha256 c6c4138ee5c86423
\newcommand{\MlstLpQgHerwigRthirtyFullOverInitBnZeroPr}{equal within 6\%}  % experiments/FIGS/data/v2/index.json :: linear probe qg_herwig, largest set, zero: R30 of L188 over the BatchNorm-recomputed random network of the same index :: sha256 c6c4138ee5c86423 :: prose twin: two significant figures, direction in words
\newcommand{\MlstLpQgHerwigRthirtyResOverInitBn}{1.17\,\ensuremath{\pm}\,0.04}  % experiments/FIGS/data/v2/index.json :: linear probe qg_herwig, largest set, chosen substitution: R30 of L162 over the BatchNorm-recomputed random network of the same index :: sha256 c6c4138ee5c86423
\newcommand{\MlstLpQgHerwigRthirtyResOverInitBnPr}{17\% more rejection}  % experiments/FIGS/data/v2/index.json :: linear probe qg_herwig, largest set, chosen substitution: R30 of L162 over the BatchNorm-recomputed random network of the same index :: sha256 c6c4138ee5c86423 :: prose twin: two significant figures, direction in words
\newcommand{\MlstLpQgHerwigRthirtyResOverInitBnZero}{0.98\,\ensuremath{\pm}\,0.10}  % experiments/FIGS/data/v2/index.json :: linear probe qg_herwig, largest set, zero: R30 of L162 over the BatchNorm-recomputed random network of the same index :: sha256 c6c4138ee5c86423
\newcommand{\MlstLpQgHerwigRthirtyResOverInitBnZeroPr}{equal within 10\%}  % experiments/FIGS/data/v2/index.json :: linear probe qg_herwig, largest set, zero: R30 of L162 over the BatchNorm-recomputed random network of the same index :: sha256 c6c4138ee5c86423 :: prose twin: two significant figures, direction in words
\newcommand{\MlstLpQgHerwigRthirtyMidOverInitBn}{1.12\,\ensuremath{\pm}\,0.07}  % experiments/FIGS/data/v2/index.json :: linear probe qg_herwig, largest set, chosen substitution: R30 of R42_Q1 over the BatchNorm-recomputed random network of the same index :: sha256 c6c4138ee5c86423
\newcommand{\MlstLpQgHerwigRthirtyMidOverInitBnPr}{equal within 20\%}  % experiments/FIGS/data/v2/index.json :: linear probe qg_herwig, largest set, chosen substitution: R30 of R42_Q1 over the BatchNorm-recomputed random network of the same index :: sha256 c6c4138ee5c86423 :: prose twin: two significant figures, direction in words
\newcommand{\MlstLpQgHerwigRthirtyMidOverInitBnZero}{0.99\,\ensuremath{\pm}\,0.07}  % experiments/FIGS/data/v2/index.json :: linear probe qg_herwig, largest set, zero: R30 of R42_Q1 over the BatchNorm-recomputed random network of the same index :: sha256 c6c4138ee5c86423
\newcommand{\MlstLpQgHerwigRthirtyMidOverInitBnZeroPr}{equal within 7\%}  % experiments/FIGS/data/v2/index.json :: linear probe qg_herwig, largest set, zero: R30 of R42_Q1 over the BatchNorm-recomputed random network of the same index :: sha256 c6c4138ee5c86423 :: prose twin: two significant figures, direction in words
\newcommand{\MlstLpQgHerwigRthirtyCoarseOverInitBn}{1.154\,\ensuremath{\pm}\,0.015}  % experiments/FIGS/data/v2/index.json :: linear probe qg_herwig, largest set, chosen substitution: R30 of R16_Q1 over the BatchNorm-recomputed random network of the same index :: sha256 c6c4138ee5c86423
\newcommand{\MlstLpQgHerwigRthirtyCoarseOverInitBnPr}{15\% more rejection}  % experiments/FIGS/data/v2/index.json :: linear probe qg_herwig, largest set, chosen substitution: R30 of R16_Q1 over the BatchNorm-recomputed random network of the same index :: sha256 c6c4138ee5c86423 :: prose twin: two significant figures, direction in words
\newcommand{\MlstLpQgHerwigRthirtyCoarseOverInitBnZero}{1.07\,\ensuremath{\pm}\,0.05}  % experiments/FIGS/data/v2/index.json :: linear probe qg_herwig, largest set, zero: R30 of R16_Q1 over the BatchNorm-recomputed random network of the same index :: sha256 c6c4138ee5c86423
\newcommand{\MlstLpQgHerwigRthirtyCoarseOverInitBnZeroPr}{equal within 8\%}  % experiments/FIGS/data/v2/index.json :: linear probe qg_herwig, largest set, zero: R30 of R16_Q1 over the BatchNorm-recomputed random network of the same index :: sha256 c6c4138ee5c86423 :: prose twin: two significant figures, direction in words
\newcommand{\MlstLpQgHerwigRthirtyResMassOverInitBn}{1.152 and 1.138}  % experiments/FIGS/data/v2/index.json :: linear probe qg_herwig, largest set, chosen substitution: R30 of L162_MASS over the BatchNorm-recomputed random network of the same index :: sha256 c6c4138ee5c86423
\newcommand{\MlstLpQgHerwigRthirtyResMassOverInitBnPr}{equal within 20\%}  % experiments/FIGS/data/v2/index.json :: linear probe qg_herwig, largest set, chosen substitution: R30 of L162_MASS over the BatchNorm-recomputed random network of the same index :: sha256 c6c4138ee5c86423 :: prose twin: two significant figures, direction in words
\newcommand{\MlstLpQgHerwigRthirtyResMassOverInitBnPrMax}{equal within 20\%}  % experiments/FIGS/data/v2/index.json :: linear probe qg_herwig, largest set, chosen substitution: R30 of L162_MASS over the BatchNorm-recomputed random network of the same index :: sha256 c6c4138ee5c86423 :: prose twin of the run farther from 1
\newcommand{\MlstLpQgHerwigRthirtyResMassOverInitBnZero}{1.13 and 1.03}  % experiments/FIGS/data/v2/index.json :: linear probe qg_herwig, largest set, zero: R30 of L162_MASS over the BatchNorm-recomputed random network of the same index :: sha256 c6c4138ee5c86423
\newcommand{\MlstLpQgHerwigRthirtyResMassOverInitBnZeroPr}{equal within 20\%}  % experiments/FIGS/data/v2/index.json :: linear probe qg_herwig, largest set, zero: R30 of L162_MASS over the BatchNorm-recomputed random network of the same index :: sha256 c6c4138ee5c86423 :: prose twin: two significant figures, direction in words
\newcommand{\MlstLpQgHerwigRthirtyResMassOverInitBnZeroPrMax}{equal within 20\%}  % experiments/FIGS/data/v2/index.json :: linear probe qg_herwig, largest set, zero: R30 of L162_MASS over the BatchNorm-recomputed random network of the same index :: sha256 c6c4138ee5c86423 :: prose twin of the run farther from 1
\newcommand{\MlstLpQgHerwigRthirtyCoarseMassOverInitBn}{1.146 and 1.159}  % experiments/FIGS/data/v2/index.json :: linear probe qg_herwig, largest set, chosen substitution: R30 of R16_Q1_MASS over the BatchNorm-recomputed random network of the same index :: sha256 c6c4138ee5c86423
\newcommand{\MlstLpQgHerwigRthirtyCoarseMassOverInitBnPr}{equal within 20\%}  % experiments/FIGS/data/v2/index.json :: linear probe qg_herwig, largest set, chosen substitution: R30 of R16_Q1_MASS over the BatchNorm-recomputed random network of the same index :: sha256 c6c4138ee5c86423 :: prose twin: two significant figures, direction in words
\newcommand{\MlstLpQgHerwigRthirtyCoarseMassOverInitBnPrMax}{equal within 20\%}  % experiments/FIGS/data/v2/index.json :: linear probe qg_herwig, largest set, chosen substitution: R30 of R16_Q1_MASS over the BatchNorm-recomputed random network of the same index :: sha256 c6c4138ee5c86423 :: prose twin of the run farther from 1
\newcommand{\MlstLpQgHerwigRthirtyCoarseMassOverInitBnZero}{1.070 and 1.100}  % experiments/FIGS/data/v2/index.json :: linear probe qg_herwig, largest set, zero: R30 of R16_Q1_MASS over the BatchNorm-recomputed random network of the same index :: sha256 c6c4138ee5c86423
\newcommand{\MlstLpQgHerwigRthirtyCoarseMassOverInitBnZeroPr}{equal within 20\%}  % experiments/FIGS/data/v2/index.json :: linear probe qg_herwig, largest set, zero: R30 of R16_Q1_MASS over the BatchNorm-recomputed random network of the same index :: sha256 c6c4138ee5c86423 :: prose twin: two significant figures, direction in words
\newcommand{\MlstLpQgHerwigRthirtyCoarseMassOverInitBnZeroPrMax}{equal within 20\%}  % experiments/FIGS/data/v2/index.json :: linear probe qg_herwig, largest set, zero: R30 of R16_Q1_MASS over the BatchNorm-recomputed random network of the same index :: sha256 c6c4138ee5c86423 :: prose twin of the run farther from 1
\newcommand{\MlstLpQgHerwigRthirtyResOverFull}{1.03\,\ensuremath{\pm}\,0.11}  % experiments/FIGS/data/v2/index.json :: linear probe qg_herwig: R30 of L162 over the 188-class run :: sha256 c6c4138ee5c86423
\newcommand{\MlstLpQgHerwigRthirtyResOverFullPr}{equal within 20\%}  % experiments/FIGS/data/v2/index.json :: linear probe qg_herwig: R30 of L162 over the 188-class run :: sha256 c6c4138ee5c86423 :: prose twin: two significant figures, direction in words
\newcommand{\MlstLpQgHerwigRthirtyResOverFullZero}{0.93\,\ensuremath{\pm}\,0.11}  % experiments/FIGS/data/v2/index.json :: linear probe qg_herwig: R30 of L162 over the 188-class run :: sha256 c6c4138ee5c86423
\newcommand{\MlstLpQgHerwigRthirtyResOverFullZeroPr}{equal within 20\%}  % experiments/FIGS/data/v2/index.json :: linear probe qg_herwig: R30 of L162 over the 188-class run :: sha256 c6c4138ee5c86423 :: prose twin: two significant figures, direction in words
\newcommand{\MlstLpQgHerwigRthirtyMidOverFull}{0.98\,\ensuremath{\pm}\,0.06}  % experiments/FIGS/data/v2/index.json :: linear probe qg_herwig: R30 of R42_Q1 over the 188-class run :: sha256 c6c4138ee5c86423
\newcommand{\MlstLpQgHerwigRthirtyMidOverFullPr}{equal within 6\%}  % experiments/FIGS/data/v2/index.json :: linear probe qg_herwig: R30 of R42_Q1 over the 188-class run :: sha256 c6c4138ee5c86423 :: prose twin: two significant figures, direction in words
\newcommand{\MlstLpQgHerwigRthirtyMidOverFullZero}{0.94\,\ensuremath{\pm}\,0.09}  % experiments/FIGS/data/v2/index.json :: linear probe qg_herwig: R30 of R42_Q1 over the 188-class run :: sha256 c6c4138ee5c86423
\newcommand{\MlstLpQgHerwigRthirtyMidOverFullZeroPr}{equal within 9\%}  % experiments/FIGS/data/v2/index.json :: linear probe qg_herwig: R30 of R42_Q1 over the 188-class run :: sha256 c6c4138ee5c86423 :: prose twin: two significant figures, direction in words
\newcommand{\MlstLpQgHerwigRthirtyCoarseOverFull}{1.02\,\ensuremath{\pm}\,0.09}  % experiments/FIGS/data/v2/index.json :: linear probe qg_herwig: R30 of R16_Q1 over the 188-class run :: sha256 c6c4138ee5c86423
\newcommand{\MlstLpQgHerwigRthirtyCoarseOverFullPr}{equal within 10\%}  % experiments/FIGS/data/v2/index.json :: linear probe qg_herwig: R30 of R16_Q1 over the 188-class run :: sha256 c6c4138ee5c86423 :: prose twin: two significant figures, direction in words
\newcommand{\MlstLpQgHerwigRthirtyCoarseOverFullZero}{1.018\,\ensuremath{\pm}\,0.024}  % experiments/FIGS/data/v2/index.json :: linear probe qg_herwig: R30 of R16_Q1 over the 188-class run :: sha256 c6c4138ee5c86423
\newcommand{\MlstLpQgHerwigRthirtyCoarseOverFullZeroPr}{equal within 3\%}  % experiments/FIGS/data/v2/index.json :: linear probe qg_herwig: R30 of R16_Q1 over the 188-class run :: sha256 c6c4138ee5c86423 :: prose twin: two significant figures, direction in words
\newcommand{\MlstLpQgHerwigRthirtyMeanOverZeroFull}{1.07\,\ensuremath{\pm}\,0.11}  % experiments/FIGS/data/v2/index.json :: linear probe qg_herwig: R30 with mean over R30 with zero, L188 :: sha256 c6c4138ee5c86423
\newcommand{\MlstLpQgHerwigRthirtyMeanOverZeroFullPr}{equal within 20\%}  % experiments/FIGS/data/v2/index.json :: linear probe qg_herwig: R30 with mean over R30 with zero, L188 :: sha256 c6c4138ee5c86423 :: prose twin: two significant figures, direction in words
\newcommand{\MlstLpQgHerwigRthirtyDrawOverZeroFull}{1.07\,\ensuremath{\pm}\,0.06}  % experiments/FIGS/data/v2/index.json :: linear probe qg_herwig: R30 with draws over R30 with zero, L188 :: sha256 c6c4138ee5c86423
\newcommand{\MlstLpQgHerwigRthirtyDrawOverZeroFullPr}{equal within 7\%}  % experiments/FIGS/data/v2/index.json :: linear probe qg_herwig: R30 with draws over R30 with zero, L188 :: sha256 c6c4138ee5c86423 :: prose twin: two significant figures, direction in words
\newcommand{\MlstLpQgHerwigRthirtyMeanOverZeroInitBn}{0.988\,\ensuremath{\pm}\,0.019}  % experiments/FIGS/data/v2/index.json :: linear probe qg_herwig: R30 with mean over R30 with zero, INIT_BN :: sha256 c6c4138ee5c86423
\newcommand{\MlstLpQgHerwigRthirtyMeanOverZeroInitBnPr}{equal within 2\%}  % experiments/FIGS/data/v2/index.json :: linear probe qg_herwig: R30 with mean over R30 with zero, INIT_BN :: sha256 c6c4138ee5c86423 :: prose twin: two significant figures, direction in words
\newcommand{\MlstLpQgHerwigRthirtyDrawOverZeroInitBn}{0.983\,\ensuremath{\pm}\,0.032}  % experiments/FIGS/data/v2/index.json :: linear probe qg_herwig: R30 with draws over R30 with zero, INIT_BN :: sha256 c6c4138ee5c86423
\newcommand{\MlstLpQgHerwigRthirtyDrawOverZeroInitBnPr}{equal within 4\%}  % experiments/FIGS/data/v2/index.json :: linear probe qg_herwig: R30 with draws over R30 with zero, INIT_BN :: sha256 c6c4138ee5c86423 :: prose twin: two significant figures, direction in words
\newcommand{\MlstLpQgHerwigRthirtyFullZero}{32.8\,\ensuremath{\pm}\,1.0}  % experiments/FIGS/data/v2/index.json :: linear probe qg_herwig, largest set: R30 with zero, L188 :: sha256 c6c4138ee5c86423
\newcommand{\MlstLpQgHerwigRthirtyFullMean}{35.0\,\ensuremath{\pm}\,2.6}  % experiments/FIGS/data/v2/index.json :: linear probe qg_herwig, largest set: R30 with mean, L188 :: sha256 c6c4138ee5c86423
\newcommand{\MlstLpQgHerwigRthirtyFullDraws}{35.0\,\ensuremath{\pm}\,0.8}  % experiments/FIGS/data/v2/index.json :: linear probe qg_herwig, largest set: R30 with draws, L188 :: sha256 c6c4138ee5c86423
\newcommand{\MlstLpQgHerwigRthirtyInitBnZero}{31.1\,\ensuremath{\pm}\,0.5}  % experiments/FIGS/data/v2/index.json :: linear probe qg_herwig, largest set: R30 with zero, INIT_BN :: sha256 c6c4138ee5c86423
\newcommand{\MlstLpQgHerwigRthirtyInitBnMean}{30.7\,\ensuremath{\pm}\,0.5}  % experiments/FIGS/data/v2/index.json :: linear probe qg_herwig, largest set: R30 with mean, INIT_BN :: sha256 c6c4138ee5c86423
\newcommand{\MlstLpQgHerwigRthirtyInitBnDraws}{30.6\,\ensuremath{\pm}\,0.6}  % experiments/FIGS/data/v2/index.json :: linear probe qg_herwig, largest set: R30 with draws, INIT_BN :: sha256 c6c4138ee5c86423
\newcommand{\MlstLpJcTwoTopZeroAccDropFull}{88.9\,\ensuremath{\pm}\,1.0\%}  % experiments/FIGS/data/v2/index.json :: linear probe JetClass-II, largest set: relative drop in accuracy (%) with the top data's missing slots set to zero, L188 :: sha256 c6c4138ee5c86423
\newcommand{\MlstLpJcTwoTopZeroAccDropFullPr}{89\%}  % experiments/FIGS/data/v2/index.json :: linear probe JetClass-II, largest set: relative drop in accuracy (%) with the top data's missing slots set to zero, L188 :: sha256 c6c4138ee5c86423 :: prose twin, two significant figures
\newcommand{\MlstLpJcTwoTopMeanAccDropFull}{87.9\,\ensuremath{\pm}\,1.1\%}  % experiments/FIGS/data/v2/index.json :: linear probe JetClass-II, largest set: relative drop in accuracy (%) with the top data's missing slots set to mean, L188 :: sha256 c6c4138ee5c86423
\newcommand{\MlstLpJcTwoTopMeanAccDropFullPr}{88\%}  % experiments/FIGS/data/v2/index.json :: linear probe JetClass-II, largest set: relative drop in accuracy (%) with the top data's missing slots set to mean, L188 :: sha256 c6c4138ee5c86423 :: prose twin, two significant figures
\newcommand{\MlstLpJcTwoTopDrawAccDropFull}{85.98\,\ensuremath{\pm}\,0.07\%}  % experiments/FIGS/data/v2/index.json :: linear probe JetClass-II, largest set: relative drop in accuracy (%) with the top data's missing slots set to draws, L188 :: sha256 c6c4138ee5c86423
\newcommand{\MlstLpJcTwoTopDrawAccDropFullPr}{86\%}  % experiments/FIGS/data/v2/index.json :: linear probe JetClass-II, largest set: relative drop in accuracy (%) with the top data's missing slots set to draws, L188 :: sha256 c6c4138ee5c86423 :: prose twin, two significant figures
\newcommand{\MlstLpJcTwoTopZeroAccDropRes}{89.0\,\ensuremath{\pm}\,0.9\%}  % experiments/FIGS/data/v2/index.json :: linear probe JetClass-II, largest set: relative drop in accuracy (%) with the top data's missing slots set to zero, L162 :: sha256 c6c4138ee5c86423
\newcommand{\MlstLpJcTwoTopZeroAccDropResPr}{89\%}  % experiments/FIGS/data/v2/index.json :: linear probe JetClass-II, largest set: relative drop in accuracy (%) with the top data's missing slots set to zero, L162 :: sha256 c6c4138ee5c86423 :: prose twin, two significant figures
\newcommand{\MlstLpJcTwoTopMeanAccDropRes}{87.8\,\ensuremath{\pm}\,0.7\%}  % experiments/FIGS/data/v2/index.json :: linear probe JetClass-II, largest set: relative drop in accuracy (%) with the top data's missing slots set to mean, L162 :: sha256 c6c4138ee5c86423
\newcommand{\MlstLpJcTwoTopMeanAccDropResPr}{88\%}  % experiments/FIGS/data/v2/index.json :: linear probe JetClass-II, largest set: relative drop in accuracy (%) with the top data's missing slots set to mean, L162 :: sha256 c6c4138ee5c86423 :: prose twin, two significant figures
\newcommand{\MlstLpJcTwoTopDrawAccDropRes}{85.89\,\ensuremath{\pm}\,0.11\%}  % experiments/FIGS/data/v2/index.json :: linear probe JetClass-II, largest set: relative drop in accuracy (%) with the top data's missing slots set to draws, L162 :: sha256 c6c4138ee5c86423
\newcommand{\MlstLpJcTwoTopDrawAccDropResPr}{86\%}  % experiments/FIGS/data/v2/index.json :: linear probe JetClass-II, largest set: relative drop in accuracy (%) with the top data's missing slots set to draws, L162 :: sha256 c6c4138ee5c86423 :: prose twin, two significant figures
\newcommand{\MlstLpJcTwoTopZeroAccDropMid}{87.81\,\ensuremath{\pm}\,0.28\%}  % experiments/FIGS/data/v2/index.json :: linear probe JetClass-II, largest set: relative drop in accuracy (%) with the top data's missing slots set to zero, R42_Q1 :: sha256 c6c4138ee5c86423
\newcommand{\MlstLpJcTwoTopZeroAccDropMidPr}{88\%}  % experiments/FIGS/data/v2/index.json :: linear probe JetClass-II, largest set: relative drop in accuracy (%) with the top data's missing slots set to zero, R42_Q1 :: sha256 c6c4138ee5c86423 :: prose twin, two significant figures
\newcommand{\MlstLpJcTwoTopMeanAccDropMid}{86.7\,\ensuremath{\pm}\,0.9\%}  % experiments/FIGS/data/v2/index.json :: linear probe JetClass-II, largest set: relative drop in accuracy (%) with the top data's missing slots set to mean, R42_Q1 :: sha256 c6c4138ee5c86423
\newcommand{\MlstLpJcTwoTopMeanAccDropMidPr}{87\%}  % experiments/FIGS/data/v2/index.json :: linear probe JetClass-II, largest set: relative drop in accuracy (%) with the top data's missing slots set to mean, R42_Q1 :: sha256 c6c4138ee5c86423 :: prose twin, two significant figures
\newcommand{\MlstLpJcTwoTopDrawAccDropMid}{83.17\,\ensuremath{\pm}\,0.29\%}  % experiments/FIGS/data/v2/index.json :: linear probe JetClass-II, largest set: relative drop in accuracy (%) with the top data's missing slots set to draws, R42_Q1 :: sha256 c6c4138ee5c86423
\newcommand{\MlstLpJcTwoTopDrawAccDropMidPr}{83\%}  % experiments/FIGS/data/v2/index.json :: linear probe JetClass-II, largest set: relative drop in accuracy (%) with the top data's missing slots set to draws, R42_Q1 :: sha256 c6c4138ee5c86423 :: prose twin, two significant figures
\newcommand{\MlstLpJcTwoTopZeroAccDropCoarse}{86.7\,\ensuremath{\pm}\,0.7\%}  % experiments/FIGS/data/v2/index.json :: linear probe JetClass-II, largest set: relative drop in accuracy (%) with the top data's missing slots set to zero, R16_Q1 :: sha256 c6c4138ee5c86423
\newcommand{\MlstLpJcTwoTopZeroAccDropCoarsePr}{87\%}  % experiments/FIGS/data/v2/index.json :: linear probe JetClass-II, largest set: relative drop in accuracy (%) with the top data's missing slots set to zero, R16_Q1 :: sha256 c6c4138ee5c86423 :: prose twin, two significant figures
\newcommand{\MlstLpJcTwoTopMeanAccDropCoarse}{86.1\,\ensuremath{\pm}\,0.8\%}  % experiments/FIGS/data/v2/index.json :: linear probe JetClass-II, largest set: relative drop in accuracy (%) with the top data's missing slots set to mean, R16_Q1 :: sha256 c6c4138ee5c86423
\newcommand{\MlstLpJcTwoTopMeanAccDropCoarsePr}{86\%}  % experiments/FIGS/data/v2/index.json :: linear probe JetClass-II, largest set: relative drop in accuracy (%) with the top data's missing slots set to mean, R16_Q1 :: sha256 c6c4138ee5c86423 :: prose twin, two significant figures
\newcommand{\MlstLpJcTwoTopDrawAccDropCoarse}{81.4\,\ensuremath{\pm}\,0.5\%}  % experiments/FIGS/data/v2/index.json :: linear probe JetClass-II, largest set: relative drop in accuracy (%) with the top data's missing slots set to draws, R16_Q1 :: sha256 c6c4138ee5c86423
\newcommand{\MlstLpJcTwoTopDrawAccDropCoarsePr}{81\%}  % experiments/FIGS/data/v2/index.json :: linear probe JetClass-II, largest set: relative drop in accuracy (%) with the top data's missing slots set to draws, R16_Q1 :: sha256 c6c4138ee5c86423 :: prose twin, two significant figures
\newcommand{\MlstLpJcTwoTopZeroAccDropResMass}{89.2, 88.3\%}  % experiments/FIGS/data/v2/index.json :: linear probe JetClass-II, largest set: relative drop in accuracy (%) with the top data's missing slots set to zero, L162_MASS :: sha256 c6c4138ee5c86423
\newcommand{\MlstLpJcTwoTopZeroAccDropResMassPr}{88--89\%}  % experiments/FIGS/data/v2/index.json :: linear probe JetClass-II, largest set: relative drop in accuracy (%) with the top data's missing slots set to zero, L162_MASS :: sha256 c6c4138ee5c86423 :: prose twin, two significant figures
\newcommand{\MlstLpJcTwoTopMeanAccDropResMass}{88.5, 87.8\%}  % experiments/FIGS/data/v2/index.json :: linear probe JetClass-II, largest set: relative drop in accuracy (%) with the top data's missing slots set to mean, L162_MASS :: sha256 c6c4138ee5c86423
\newcommand{\MlstLpJcTwoTopMeanAccDropResMassPr}{88--88\%}  % experiments/FIGS/data/v2/index.json :: linear probe JetClass-II, largest set: relative drop in accuracy (%) with the top data's missing slots set to mean, L162_MASS :: sha256 c6c4138ee5c86423 :: prose twin, two significant figures
\newcommand{\MlstLpJcTwoTopDrawAccDropResMass}{85.4, 85.4\%}  % experiments/FIGS/data/v2/index.json :: linear probe JetClass-II, largest set: relative drop in accuracy (%) with the top data's missing slots set to draws, L162_MASS :: sha256 c6c4138ee5c86423
\newcommand{\MlstLpJcTwoTopDrawAccDropResMassPr}{85--85\%}  % experiments/FIGS/data/v2/index.json :: linear probe JetClass-II, largest set: relative drop in accuracy (%) with the top data's missing slots set to draws, L162_MASS :: sha256 c6c4138ee5c86423 :: prose twin, two significant figures
\newcommand{\MlstLpJcTwoTopZeroAccDropCoarseMass}{84.8, 87.4\%}  % experiments/FIGS/data/v2/index.json :: linear probe JetClass-II, largest set: relative drop in accuracy (%) with the top data's missing slots set to zero, R16_Q1_MASS :: sha256 c6c4138ee5c86423
\newcommand{\MlstLpJcTwoTopZeroAccDropCoarseMassPr}{85--87\%}  % experiments/FIGS/data/v2/index.json :: linear probe JetClass-II, largest set: relative drop in accuracy (%) with the top data's missing slots set to zero, R16_Q1_MASS :: sha256 c6c4138ee5c86423 :: prose twin, two significant figures
\newcommand{\MlstLpJcTwoTopMeanAccDropCoarseMass}{85.9, 86.9\%}  % experiments/FIGS/data/v2/index.json :: linear probe JetClass-II, largest set: relative drop in accuracy (%) with the top data's missing slots set to mean, R16_Q1_MASS :: sha256 c6c4138ee5c86423
\newcommand{\MlstLpJcTwoTopMeanAccDropCoarseMassPr}{86--87\%}  % experiments/FIGS/data/v2/index.json :: linear probe JetClass-II, largest set: relative drop in accuracy (%) with the top data's missing slots set to mean, R16_Q1_MASS :: sha256 c6c4138ee5c86423 :: prose twin, two significant figures
\newcommand{\MlstLpJcTwoTopDrawAccDropCoarseMass}{79.1, 81.7\%}  % experiments/FIGS/data/v2/index.json :: linear probe JetClass-II, largest set: relative drop in accuracy (%) with the top data's missing slots set to draws, R16_Q1_MASS :: sha256 c6c4138ee5c86423
\newcommand{\MlstLpJcTwoTopDrawAccDropCoarseMassPr}{79--82\%}  % experiments/FIGS/data/v2/index.json :: linear probe JetClass-II, largest set: relative drop in accuracy (%) with the top data's missing slots set to draws, R16_Q1_MASS :: sha256 c6c4138ee5c86423 :: prose twin, two significant figures
\newcommand{\MlstLpJcTwoTopZeroAccDropFzero}{88.9, 88.5\%}  % experiments/FIGS/data/v2/index.json :: linear probe JetClass-II, largest set: relative drop in accuracy (%) with the top data's missing slots set to zero, FLAV_F0 :: sha256 c6c4138ee5c86423
\newcommand{\MlstLpJcTwoTopZeroAccDropFzeroPr}{88--89\%}  % experiments/FIGS/data/v2/index.json :: linear probe JetClass-II, largest set: relative drop in accuracy (%) with the top data's missing slots set to zero, FLAV_F0 :: sha256 c6c4138ee5c86423 :: prose twin, two significant figures
\newcommand{\MlstLpJcTwoTopMeanAccDropFzero}{88.3, 87.0\%}  % experiments/FIGS/data/v2/index.json :: linear probe JetClass-II, largest set: relative drop in accuracy (%) with the top data's missing slots set to mean, FLAV_F0 :: sha256 c6c4138ee5c86423
\newcommand{\MlstLpJcTwoTopMeanAccDropFzeroPr}{87--88\%}  % experiments/FIGS/data/v2/index.json :: linear probe JetClass-II, largest set: relative drop in accuracy (%) with the top data's missing slots set to mean, FLAV_F0 :: sha256 c6c4138ee5c86423 :: prose twin, two significant figures
\newcommand{\MlstLpJcTwoTopDrawAccDropFzero}{83.3, 83.9\%}  % experiments/FIGS/data/v2/index.json :: linear probe JetClass-II, largest set: relative drop in accuracy (%) with the top data's missing slots set to draws, FLAV_F0 :: sha256 c6c4138ee5c86423
\newcommand{\MlstLpJcTwoTopDrawAccDropFzeroPr}{83--84\%}  % experiments/FIGS/data/v2/index.json :: linear probe JetClass-II, largest set: relative drop in accuracy (%) with the top data's missing slots set to draws, FLAV_F0 :: sha256 c6c4138ee5c86423 :: prose twin, two significant figures
\newcommand{\MlstLpJcTwoTopZeroAccDropFone}{87.8, 87.6\%}  % experiments/FIGS/data/v2/index.json :: linear probe JetClass-II, largest set: relative drop in accuracy (%) with the top data's missing slots set to zero, FLAV_F1 :: sha256 c6c4138ee5c86423
\newcommand{\MlstLpJcTwoTopZeroAccDropFonePr}{88--88\%}  % experiments/FIGS/data/v2/index.json :: linear probe JetClass-II, largest set: relative drop in accuracy (%) with the top data's missing slots set to zero, FLAV_F1 :: sha256 c6c4138ee5c86423 :: prose twin, two significant figures
\newcommand{\MlstLpJcTwoTopMeanAccDropFone}{88.9, 86.9\%}  % experiments/FIGS/data/v2/index.json :: linear probe JetClass-II, largest set: relative drop in accuracy (%) with the top data's missing slots set to mean, FLAV_F1 :: sha256 c6c4138ee5c86423
\newcommand{\MlstLpJcTwoTopMeanAccDropFonePr}{87--89\%}  % experiments/FIGS/data/v2/index.json :: linear probe JetClass-II, largest set: relative drop in accuracy (%) with the top data's missing slots set to mean, FLAV_F1 :: sha256 c6c4138ee5c86423 :: prose twin, two significant figures
\newcommand{\MlstLpJcTwoTopDrawAccDropFone}{84.4, 85.2\%}  % experiments/FIGS/data/v2/index.json :: linear probe JetClass-II, largest set: relative drop in accuracy (%) with the top data's missing slots set to draws, FLAV_F1 :: sha256 c6c4138ee5c86423
\newcommand{\MlstLpJcTwoTopDrawAccDropFonePr}{84--85\%}  % experiments/FIGS/data/v2/index.json :: linear probe JetClass-II, largest set: relative drop in accuracy (%) with the top data's missing slots set to draws, FLAV_F1 :: sha256 c6c4138ee5c86423 :: prose twin, two significant figures
\newcommand{\MlstLpJcTwoTopZeroAccDropFoner}{89.6, 90.3\%}  % experiments/FIGS/data/v2/index.json :: linear probe JetClass-II, largest set: relative drop in accuracy (%) with the top data's missing slots set to zero, FLAV_F1R :: sha256 c6c4138ee5c86423
\newcommand{\MlstLpJcTwoTopZeroAccDropFonerPr}{90--90\%}  % experiments/FIGS/data/v2/index.json :: linear probe JetClass-II, largest set: relative drop in accuracy (%) with the top data's missing slots set to zero, FLAV_F1R :: sha256 c6c4138ee5c86423 :: prose twin, two significant figures
\newcommand{\MlstLpJcTwoTopMeanAccDropFoner}{89.4, 89.7\%}  % experiments/FIGS/data/v2/index.json :: linear probe JetClass-II, largest set: relative drop in accuracy (%) with the top data's missing slots set to mean, FLAV_F1R :: sha256 c6c4138ee5c86423
\newcommand{\MlstLpJcTwoTopMeanAccDropFonerPr}{89--90\%}  % experiments/FIGS/data/v2/index.json :: linear probe JetClass-II, largest set: relative drop in accuracy (%) with the top data's missing slots set to mean, FLAV_F1R :: sha256 c6c4138ee5c86423 :: prose twin, two significant figures
\newcommand{\MlstLpJcTwoTopDrawAccDropFoner}{84.9, 85.6\%}  % experiments/FIGS/data/v2/index.json :: linear probe JetClass-II, largest set: relative drop in accuracy (%) with the top data's missing slots set to draws, FLAV_F1R :: sha256 c6c4138ee5c86423
\newcommand{\MlstLpJcTwoTopDrawAccDropFonerPr}{85--86\%}  % experiments/FIGS/data/v2/index.json :: linear probe JetClass-II, largest set: relative drop in accuracy (%) with the top data's missing slots set to draws, FLAV_F1R :: sha256 c6c4138ee5c86423 :: prose twin, two significant figures
\newcommand{\MlstLpJcTwoTopZeroAccDropInitBn}{56.4\,\ensuremath{\pm}\,1.4\%}  % experiments/FIGS/data/v2/index.json :: linear probe JetClass-II, largest set: relative drop in accuracy (%) with the top data's missing slots set to zero, INIT_BN :: sha256 c6c4138ee5c86423
\newcommand{\MlstLpJcTwoTopZeroAccDropInitBnPr}{56\%}  % experiments/FIGS/data/v2/index.json :: linear probe JetClass-II, largest set: relative drop in accuracy (%) with the top data's missing slots set to zero, INIT_BN :: sha256 c6c4138ee5c86423 :: prose twin, two significant figures
\newcommand{\MlstLpJcTwoTopMeanAccDropInitBn}{59.8\,\ensuremath{\pm}\,0.8\%}  % experiments/FIGS/data/v2/index.json :: linear probe JetClass-II, largest set: relative drop in accuracy (%) with the top data's missing slots set to mean, INIT_BN :: sha256 c6c4138ee5c86423
\newcommand{\MlstLpJcTwoTopMeanAccDropInitBnPr}{60\%}  % experiments/FIGS/data/v2/index.json :: linear probe JetClass-II, largest set: relative drop in accuracy (%) with the top data's missing slots set to mean, INIT_BN :: sha256 c6c4138ee5c86423 :: prose twin, two significant figures
\newcommand{\MlstLpJcTwoTopDrawAccDropInitBn}{59.8\,\ensuremath{\pm}\,1.2\%}  % experiments/FIGS/data/v2/index.json :: linear probe JetClass-II, largest set: relative drop in accuracy (%) with the top data's missing slots set to draws, INIT_BN :: sha256 c6c4138ee5c86423
\newcommand{\MlstLpJcTwoTopDrawAccDropInitBnPr}{60\%}  % experiments/FIGS/data/v2/index.json :: linear probe JetClass-II, largest set: relative drop in accuracy (%) with the top data's missing slots set to draws, INIT_BN :: sha256 c6c4138ee5c86423 :: prose twin, two significant figures
\newcommand{\MlstLpJcTwoQgZeroAccDropFull}{61.2\,\ensuremath{\pm}\,3.0\%}  % experiments/FIGS/data/v2/index.json :: linear probe JetClass-II, largest set: relative drop in accuracy (%) with the qg data's missing slots set to zero, L188 :: sha256 c6c4138ee5c86423
\newcommand{\MlstLpJcTwoQgZeroAccDropFullPr}{61\%}  % experiments/FIGS/data/v2/index.json :: linear probe JetClass-II, largest set: relative drop in accuracy (%) with the qg data's missing slots set to zero, L188 :: sha256 c6c4138ee5c86423 :: prose twin, two significant figures
\newcommand{\MlstLpJcTwoQgMeanAccDropFull}{48.8\,\ensuremath{\pm}\,1.1\%}  % experiments/FIGS/data/v2/index.json :: linear probe JetClass-II, largest set: relative drop in accuracy (%) with the qg data's missing slots set to mean, L188 :: sha256 c6c4138ee5c86423
\newcommand{\MlstLpJcTwoQgMeanAccDropFullPr}{49\%}  % experiments/FIGS/data/v2/index.json :: linear probe JetClass-II, largest set: relative drop in accuracy (%) with the qg data's missing slots set to mean, L188 :: sha256 c6c4138ee5c86423 :: prose twin, two significant figures
\newcommand{\MlstLpJcTwoQgDrawAccDropFull}{54.30\,\ensuremath{\pm}\,0.21\%}  % experiments/FIGS/data/v2/index.json :: linear probe JetClass-II, largest set: relative drop in accuracy (%) with the qg data's missing slots set to draws, L188 :: sha256 c6c4138ee5c86423
\newcommand{\MlstLpJcTwoQgDrawAccDropFullPr}{54\%}  % experiments/FIGS/data/v2/index.json :: linear probe JetClass-II, largest set: relative drop in accuracy (%) with the qg data's missing slots set to draws, L188 :: sha256 c6c4138ee5c86423 :: prose twin, two significant figures
\newcommand{\MlstLpJcTwoQgZeroAccDropRes}{64\,\ensuremath{\pm}\,5\%}  % experiments/FIGS/data/v2/index.json :: linear probe JetClass-II, largest set: relative drop in accuracy (%) with the qg data's missing slots set to zero, L162 :: sha256 c6c4138ee5c86423
\newcommand{\MlstLpJcTwoQgZeroAccDropResPr}{64\%}  % experiments/FIGS/data/v2/index.json :: linear probe JetClass-II, largest set: relative drop in accuracy (%) with the qg data's missing slots set to zero, L162 :: sha256 c6c4138ee5c86423 :: prose twin, two significant figures
\newcommand{\MlstLpJcTwoQgMeanAccDropRes}{49.5\,\ensuremath{\pm}\,0.5\%}  % experiments/FIGS/data/v2/index.json :: linear probe JetClass-II, largest set: relative drop in accuracy (%) with the qg data's missing slots set to mean, L162 :: sha256 c6c4138ee5c86423
\newcommand{\MlstLpJcTwoQgMeanAccDropResPr}{49\%}  % experiments/FIGS/data/v2/index.json :: linear probe JetClass-II, largest set: relative drop in accuracy (%) with the qg data's missing slots set to mean, L162 :: sha256 c6c4138ee5c86423 :: prose twin, two significant figures
\newcommand{\MlstLpJcTwoQgDrawAccDropRes}{54.6\,\ensuremath{\pm}\,0.6\%}  % experiments/FIGS/data/v2/index.json :: linear probe JetClass-II, largest set: relative drop in accuracy (%) with the qg data's missing slots set to draws, L162 :: sha256 c6c4138ee5c86423
\newcommand{\MlstLpJcTwoQgDrawAccDropResPr}{55\%}  % experiments/FIGS/data/v2/index.json :: linear probe JetClass-II, largest set: relative drop in accuracy (%) with the qg data's missing slots set to draws, L162 :: sha256 c6c4138ee5c86423 :: prose twin, two significant figures
\newcommand{\MlstLpJcTwoQgZeroAccDropMid}{68\,\ensuremath{\pm}\,5\%}  % experiments/FIGS/data/v2/index.json :: linear probe JetClass-II, largest set: relative drop in accuracy (%) with the qg data's missing slots set to zero, R42_Q1 :: sha256 c6c4138ee5c86423
\newcommand{\MlstLpJcTwoQgZeroAccDropMidPr}{68\%}  % experiments/FIGS/data/v2/index.json :: linear probe JetClass-II, largest set: relative drop in accuracy (%) with the qg data's missing slots set to zero, R42_Q1 :: sha256 c6c4138ee5c86423 :: prose twin, two significant figures
\newcommand{\MlstLpJcTwoQgMeanAccDropMid}{47.6\,\ensuremath{\pm}\,1.6\%}  % experiments/FIGS/data/v2/index.json :: linear probe JetClass-II, largest set: relative drop in accuracy (%) with the qg data's missing slots set to mean, R42_Q1 :: sha256 c6c4138ee5c86423
\newcommand{\MlstLpJcTwoQgMeanAccDropMidPr}{48\%}  % experiments/FIGS/data/v2/index.json :: linear probe JetClass-II, largest set: relative drop in accuracy (%) with the qg data's missing slots set to mean, R42_Q1 :: sha256 c6c4138ee5c86423 :: prose twin, two significant figures
\newcommand{\MlstLpJcTwoQgDrawAccDropMid}{53.5\,\ensuremath{\pm}\,1.4\%}  % experiments/FIGS/data/v2/index.json :: linear probe JetClass-II, largest set: relative drop in accuracy (%) with the qg data's missing slots set to draws, R42_Q1 :: sha256 c6c4138ee5c86423
\newcommand{\MlstLpJcTwoQgDrawAccDropMidPr}{54\%}  % experiments/FIGS/data/v2/index.json :: linear probe JetClass-II, largest set: relative drop in accuracy (%) with the qg data's missing slots set to draws, R42_Q1 :: sha256 c6c4138ee5c86423 :: prose twin, two significant figures
\newcommand{\MlstLpJcTwoQgZeroAccDropCoarse}{70.9\,\ensuremath{\pm}\,2.3\%}  % experiments/FIGS/data/v2/index.json :: linear probe JetClass-II, largest set: relative drop in accuracy (%) with the qg data's missing slots set to zero, R16_Q1 :: sha256 c6c4138ee5c86423
\newcommand{\MlstLpJcTwoQgZeroAccDropCoarsePr}{71\%}  % experiments/FIGS/data/v2/index.json :: linear probe JetClass-II, largest set: relative drop in accuracy (%) with the qg data's missing slots set to zero, R16_Q1 :: sha256 c6c4138ee5c86423 :: prose twin, two significant figures
\newcommand{\MlstLpJcTwoQgMeanAccDropCoarse}{50.0\,\ensuremath{\pm}\,1.0\%}  % experiments/FIGS/data/v2/index.json :: linear probe JetClass-II, largest set: relative drop in accuracy (%) with the qg data's missing slots set to mean, R16_Q1 :: sha256 c6c4138ee5c86423
\newcommand{\MlstLpJcTwoQgMeanAccDropCoarsePr}{50\%}  % experiments/FIGS/data/v2/index.json :: linear probe JetClass-II, largest set: relative drop in accuracy (%) with the qg data's missing slots set to mean, R16_Q1 :: sha256 c6c4138ee5c86423 :: prose twin, two significant figures
\newcommand{\MlstLpJcTwoQgDrawAccDropCoarse}{51.7\,\ensuremath{\pm}\,0.4\%}  % experiments/FIGS/data/v2/index.json :: linear probe JetClass-II, largest set: relative drop in accuracy (%) with the qg data's missing slots set to draws, R16_Q1 :: sha256 c6c4138ee5c86423
\newcommand{\MlstLpJcTwoQgDrawAccDropCoarsePr}{52\%}  % experiments/FIGS/data/v2/index.json :: linear probe JetClass-II, largest set: relative drop in accuracy (%) with the qg data's missing slots set to draws, R16_Q1 :: sha256 c6c4138ee5c86423 :: prose twin, two significant figures
\newcommand{\MlstLpJcTwoQgZeroAccDropResMass}{61.8, 67.9\%}  % experiments/FIGS/data/v2/index.json :: linear probe JetClass-II, largest set: relative drop in accuracy (%) with the qg data's missing slots set to zero, L162_MASS :: sha256 c6c4138ee5c86423
\newcommand{\MlstLpJcTwoQgZeroAccDropResMassPr}{62--68\%}  % experiments/FIGS/data/v2/index.json :: linear probe JetClass-II, largest set: relative drop in accuracy (%) with the qg data's missing slots set to zero, L162_MASS :: sha256 c6c4138ee5c86423 :: prose twin, two significant figures
\newcommand{\MlstLpJcTwoQgMeanAccDropResMass}{49.7, 52.2\%}  % experiments/FIGS/data/v2/index.json :: linear probe JetClass-II, largest set: relative drop in accuracy (%) with the qg data's missing slots set to mean, L162_MASS :: sha256 c6c4138ee5c86423
\newcommand{\MlstLpJcTwoQgMeanAccDropResMassPr}{50--52\%}  % experiments/FIGS/data/v2/index.json :: linear probe JetClass-II, largest set: relative drop in accuracy (%) with the qg data's missing slots set to mean, L162_MASS :: sha256 c6c4138ee5c86423 :: prose twin, two significant figures
\newcommand{\MlstLpJcTwoQgDrawAccDropResMass}{54.3, 55.1\%}  % experiments/FIGS/data/v2/index.json :: linear probe JetClass-II, largest set: relative drop in accuracy (%) with the qg data's missing slots set to draws, L162_MASS :: sha256 c6c4138ee5c86423
\newcommand{\MlstLpJcTwoQgDrawAccDropResMassPr}{54--55\%}  % experiments/FIGS/data/v2/index.json :: linear probe JetClass-II, largest set: relative drop in accuracy (%) with the qg data's missing slots set to draws, L162_MASS :: sha256 c6c4138ee5c86423 :: prose twin, two significant figures
\newcommand{\MlstLpJcTwoQgZeroAccDropCoarseMass}{66.7, 65.4\%}  % experiments/FIGS/data/v2/index.json :: linear probe JetClass-II, largest set: relative drop in accuracy (%) with the qg data's missing slots set to zero, R16_Q1_MASS :: sha256 c6c4138ee5c86423
\newcommand{\MlstLpJcTwoQgZeroAccDropCoarseMassPr}{65--67\%}  % experiments/FIGS/data/v2/index.json :: linear probe JetClass-II, largest set: relative drop in accuracy (%) with the qg data's missing slots set to zero, R16_Q1_MASS :: sha256 c6c4138ee5c86423 :: prose twin, two significant figures
\newcommand{\MlstLpJcTwoQgMeanAccDropCoarseMass}{49.0, 49.6\%}  % experiments/FIGS/data/v2/index.json :: linear probe JetClass-II, largest set: relative drop in accuracy (%) with the qg data's missing slots set to mean, R16_Q1_MASS :: sha256 c6c4138ee5c86423
\newcommand{\MlstLpJcTwoQgMeanAccDropCoarseMassPr}{49--50\%}  % experiments/FIGS/data/v2/index.json :: linear probe JetClass-II, largest set: relative drop in accuracy (%) with the qg data's missing slots set to mean, R16_Q1_MASS :: sha256 c6c4138ee5c86423 :: prose twin, two significant figures
\newcommand{\MlstLpJcTwoQgDrawAccDropCoarseMass}{50.8, 51.1\%}  % experiments/FIGS/data/v2/index.json :: linear probe JetClass-II, largest set: relative drop in accuracy (%) with the qg data's missing slots set to draws, R16_Q1_MASS :: sha256 c6c4138ee5c86423
\newcommand{\MlstLpJcTwoQgDrawAccDropCoarseMassPr}{51--51\%}  % experiments/FIGS/data/v2/index.json :: linear probe JetClass-II, largest set: relative drop in accuracy (%) with the qg data's missing slots set to draws, R16_Q1_MASS :: sha256 c6c4138ee5c86423 :: prose twin, two significant figures
\newcommand{\MlstLpJcTwoQgZeroAccDropFzero}{62.7, 62.2\%}  % experiments/FIGS/data/v2/index.json :: linear probe JetClass-II, largest set: relative drop in accuracy (%) with the qg data's missing slots set to zero, FLAV_F0 :: sha256 c6c4138ee5c86423
\newcommand{\MlstLpJcTwoQgZeroAccDropFzeroPr}{62--63\%}  % experiments/FIGS/data/v2/index.json :: linear probe JetClass-II, largest set: relative drop in accuracy (%) with the qg data's missing slots set to zero, FLAV_F0 :: sha256 c6c4138ee5c86423 :: prose twin, two significant figures
\newcommand{\MlstLpJcTwoQgMeanAccDropFzero}{51.9, 51.3\%}  % experiments/FIGS/data/v2/index.json :: linear probe JetClass-II, largest set: relative drop in accuracy (%) with the qg data's missing slots set to mean, FLAV_F0 :: sha256 c6c4138ee5c86423
\newcommand{\MlstLpJcTwoQgMeanAccDropFzeroPr}{51--52\%}  % experiments/FIGS/data/v2/index.json :: linear probe JetClass-II, largest set: relative drop in accuracy (%) with the qg data's missing slots set to mean, FLAV_F0 :: sha256 c6c4138ee5c86423 :: prose twin, two significant figures
\newcommand{\MlstLpJcTwoQgDrawAccDropFzero}{53.9, 53.4\%}  % experiments/FIGS/data/v2/index.json :: linear probe JetClass-II, largest set: relative drop in accuracy (%) with the qg data's missing slots set to draws, FLAV_F0 :: sha256 c6c4138ee5c86423
\newcommand{\MlstLpJcTwoQgDrawAccDropFzeroPr}{53--54\%}  % experiments/FIGS/data/v2/index.json :: linear probe JetClass-II, largest set: relative drop in accuracy (%) with the qg data's missing slots set to draws, FLAV_F0 :: sha256 c6c4138ee5c86423 :: prose twin, two significant figures
\newcommand{\MlstLpJcTwoQgZeroAccDropFone}{66.2, 62.8\%}  % experiments/FIGS/data/v2/index.json :: linear probe JetClass-II, largest set: relative drop in accuracy (%) with the qg data's missing slots set to zero, FLAV_F1 :: sha256 c6c4138ee5c86423
\newcommand{\MlstLpJcTwoQgZeroAccDropFonePr}{63--66\%}  % experiments/FIGS/data/v2/index.json :: linear probe JetClass-II, largest set: relative drop in accuracy (%) with the qg data's missing slots set to zero, FLAV_F1 :: sha256 c6c4138ee5c86423 :: prose twin, two significant figures
\newcommand{\MlstLpJcTwoQgMeanAccDropFone}{52.7, 52.8\%}  % experiments/FIGS/data/v2/index.json :: linear probe JetClass-II, largest set: relative drop in accuracy (%) with the qg data's missing slots set to mean, FLAV_F1 :: sha256 c6c4138ee5c86423
\newcommand{\MlstLpJcTwoQgMeanAccDropFonePr}{53--53\%}  % experiments/FIGS/data/v2/index.json :: linear probe JetClass-II, largest set: relative drop in accuracy (%) with the qg data's missing slots set to mean, FLAV_F1 :: sha256 c6c4138ee5c86423 :: prose twin, two significant figures
\newcommand{\MlstLpJcTwoQgDrawAccDropFone}{55.6, 56.5\%}  % experiments/FIGS/data/v2/index.json :: linear probe JetClass-II, largest set: relative drop in accuracy (%) with the qg data's missing slots set to draws, FLAV_F1 :: sha256 c6c4138ee5c86423
\newcommand{\MlstLpJcTwoQgDrawAccDropFonePr}{56--56\%}  % experiments/FIGS/data/v2/index.json :: linear probe JetClass-II, largest set: relative drop in accuracy (%) with the qg data's missing slots set to draws, FLAV_F1 :: sha256 c6c4138ee5c86423 :: prose twin, two significant figures
\newcommand{\MlstLpJcTwoQgZeroAccDropFoner}{67.2, 69.0\%}  % experiments/FIGS/data/v2/index.json :: linear probe JetClass-II, largest set: relative drop in accuracy (%) with the qg data's missing slots set to zero, FLAV_F1R :: sha256 c6c4138ee5c86423
\newcommand{\MlstLpJcTwoQgZeroAccDropFonerPr}{67--69\%}  % experiments/FIGS/data/v2/index.json :: linear probe JetClass-II, largest set: relative drop in accuracy (%) with the qg data's missing slots set to zero, FLAV_F1R :: sha256 c6c4138ee5c86423 :: prose twin, two significant figures
\newcommand{\MlstLpJcTwoQgMeanAccDropFoner}{52.9, 51.8\%}  % experiments/FIGS/data/v2/index.json :: linear probe JetClass-II, largest set: relative drop in accuracy (%) with the qg data's missing slots set to mean, FLAV_F1R :: sha256 c6c4138ee5c86423
\newcommand{\MlstLpJcTwoQgMeanAccDropFonerPr}{52--53\%}  % experiments/FIGS/data/v2/index.json :: linear probe JetClass-II, largest set: relative drop in accuracy (%) with the qg data's missing slots set to mean, FLAV_F1R :: sha256 c6c4138ee5c86423 :: prose twin, two significant figures
\newcommand{\MlstLpJcTwoQgDrawAccDropFoner}{57.2, 55.4\%}  % experiments/FIGS/data/v2/index.json :: linear probe JetClass-II, largest set: relative drop in accuracy (%) with the qg data's missing slots set to draws, FLAV_F1R :: sha256 c6c4138ee5c86423
\newcommand{\MlstLpJcTwoQgDrawAccDropFonerPr}{55--57\%}  % experiments/FIGS/data/v2/index.json :: linear probe JetClass-II, largest set: relative drop in accuracy (%) with the qg data's missing slots set to draws, FLAV_F1R :: sha256 c6c4138ee5c86423 :: prose twin, two significant figures
\newcommand{\MlstLpJcTwoQgZeroAccDropInitBn}{8.7\,\ensuremath{\pm}\,2.6\%}  % experiments/FIGS/data/v2/index.json :: linear probe JetClass-II, largest set: relative drop in accuracy (%) with the qg data's missing slots set to zero, INIT_BN :: sha256 c6c4138ee5c86423
\newcommand{\MlstLpJcTwoQgZeroAccDropInitBnPr}{8.7\%}  % experiments/FIGS/data/v2/index.json :: linear probe JetClass-II, largest set: relative drop in accuracy (%) with the qg data's missing slots set to zero, INIT_BN :: sha256 c6c4138ee5c86423 :: prose twin, two significant figures
\newcommand{\MlstLpJcTwoQgMeanAccDropInitBn}{10.4\,\ensuremath{\pm}\,1.4\%}  % experiments/FIGS/data/v2/index.json :: linear probe JetClass-II, largest set: relative drop in accuracy (%) with the qg data's missing slots set to mean, INIT_BN :: sha256 c6c4138ee5c86423
\newcommand{\MlstLpJcTwoQgMeanAccDropInitBnPr}{10\%}  % experiments/FIGS/data/v2/index.json :: linear probe JetClass-II, largest set: relative drop in accuracy (%) with the qg data's missing slots set to mean, INIT_BN :: sha256 c6c4138ee5c86423 :: prose twin, two significant figures
\newcommand{\MlstLpJcTwoQgDrawAccDropInitBn}{11.4\,\ensuremath{\pm}\,2.1\%}  % experiments/FIGS/data/v2/index.json :: linear probe JetClass-II, largest set: relative drop in accuracy (%) with the qg data's missing slots set to draws, INIT_BN :: sha256 c6c4138ee5c86423
\newcommand{\MlstLpJcTwoQgDrawAccDropInitBnPr}{11\%}  % experiments/FIGS/data/v2/index.json :: linear probe JetClass-II, largest set: relative drop in accuracy (%) with the qg data's missing slots set to draws, INIT_BN :: sha256 c6c4138ee5c86423 :: prose twin, two significant figures
\newcommand{\MlstMassLambda}{5}  % configs/arms/v2_grid.json :: arms[L162_MASS, R16_Q1_MASS].mass_lambda :: sha256 d8a31b1486992302 :: weight of the mass-regression term in the pretraining loss, both mass models
\newcommand{\MlstMassLossShareResMass}{13.78 and 13.87}  % experiments/FIGS/data/v2/index.json :: share of the mass-regression term in the training loss of L162_MASS, mean over epochs, percent :: sha256 c6c4138ee5c86423
\newcommand{\MlstMassLossShareCoarseMass}{28.65 and 28.67}  % experiments/FIGS/data/v2/index.json :: share of the mass-regression term in the training loss of R16_Q1_MASS, mean over epochs, percent :: sha256 c6c4138ee5c86423
\newcommand{\MlstMassUncorrected}{0.245}  % experiments/FIGS/data/v2/index.json :: head/atlas_res: groomed mass, all jets, ATLAS resolution :: sha256 c6c4138ee5c86423
\newcommand{\MlstMassUncorrectedResponse}{0.929}  % experiments/FIGS/data/v2/index.json :: head/atlas_res: groomed mass, all jets, median response :: sha256 c6c4138ee5c86423
\newcommand{\MlstMassHeadResMass}{0.1471 and 0.1463}  % experiments/FIGS/data/v2/index.json :: regression output of L162_MASS, all jets, ATLAS resolution :: sha256 c6c4138ee5c86423
\newcommand{\MlstMassHeadResMassResponse}{0.995 and 1.003}  % experiments/FIGS/data/v2/index.json :: regression output of L162_MASS, all jets, median response :: sha256 c6c4138ee5c86423
\newcommand{\MlstMassHeadResMassOverUncorr}{0.6009 and 0.5976}  % experiments/FIGS/data/v2/index.json :: resolution of the regression output of L162_MASS over the groomed mass's, all jets :: sha256 c6c4138ee5c86423
\newcommand{\MlstMassHeadResMassOverUncorrPr}{at most 0.60 times the resolution}  % experiments/FIGS/data/v2/index.json :: resolution of the regression output of L162_MASS over the groomed mass's, all jets :: sha256 c6c4138ee5c86423 :: prose twin: two significant figures, direction in words
\newcommand{\MlstMassHeadResMassOverUncorrPrMax}{down to 0.60 times the resolution}  % experiments/FIGS/data/v2/index.json :: resolution of the regression output of L162_MASS over the groomed mass's, all jets :: sha256 c6c4138ee5c86423 :: prose twin of the run farther from 1
\newcommand{\MlstMassHeadResMassOverMlpRes}{0.6645 and 0.6691}  % experiments/FIGS/data/v2/index.json :: resolution of the regression output of L162_MASS over the frozen-feature MLP of L162, run by run, all jets :: sha256 c6c4138ee5c86423
\newcommand{\MlstMassHeadResMassOverMlpResPr}{at most 0.67 times the resolution}  % experiments/FIGS/data/v2/index.json :: resolution of the regression output of L162_MASS over the frozen-feature MLP of L162, run by run, all jets :: sha256 c6c4138ee5c86423 :: prose twin: two significant figures, direction in words
\newcommand{\MlstMassHeadResMassOverMlpResPrMax}{down to 0.66 times the resolution}  % experiments/FIGS/data/v2/index.json :: resolution of the regression output of L162_MASS over the frozen-feature MLP of L162, run by run, all jets :: sha256 c6c4138ee5c86423 :: prose twin of the run farther from 1
\newcommand{\MlstMassHeadCoarseMass}{0.1450 and 0.1461}  % experiments/FIGS/data/v2/index.json :: regression output of R16_Q1_MASS, all jets, ATLAS resolution :: sha256 c6c4138ee5c86423
\newcommand{\MlstMassHeadCoarseMassResponse}{1.005 and 1.005}  % experiments/FIGS/data/v2/index.json :: regression output of R16_Q1_MASS, all jets, median response :: sha256 c6c4138ee5c86423
\newcommand{\MlstMassHeadCoarseMassOverUncorr}{0.5921 and 0.5966}  % experiments/FIGS/data/v2/index.json :: resolution of the regression output of R16_Q1_MASS over the groomed mass's, all jets :: sha256 c6c4138ee5c86423
\newcommand{\MlstMassHeadCoarseMassOverUncorrPr}{at most 0.60 times the resolution}  % experiments/FIGS/data/v2/index.json :: resolution of the regression output of R16_Q1_MASS over the groomed mass's, all jets :: sha256 c6c4138ee5c86423 :: prose twin: two significant figures, direction in words
\newcommand{\MlstMassHeadCoarseMassOverUncorrPrMax}{down to 0.59 times the resolution}  % experiments/FIGS/data/v2/index.json :: resolution of the regression output of R16_Q1_MASS over the groomed mass's, all jets :: sha256 c6c4138ee5c86423 :: prose twin of the run farther from 1
\newcommand{\MlstMassHeadCoarseMassOverMlpCoarse}{0.703 and 0.690}  % experiments/FIGS/data/v2/index.json :: resolution of the regression output of R16_Q1_MASS over the frozen-feature MLP of R16_Q1, run by run, all jets :: sha256 c6c4138ee5c86423
\newcommand{\MlstMassHeadCoarseMassOverMlpCoarsePr}{at most 0.70 times the resolution}  % experiments/FIGS/data/v2/index.json :: resolution of the regression output of R16_Q1_MASS over the frozen-feature MLP of R16_Q1, run by run, all jets :: sha256 c6c4138ee5c86423 :: prose twin: two significant figures, direction in words
\newcommand{\MlstMassHeadCoarseMassOverMlpCoarsePrMax}{down to 0.69 times the resolution}  % experiments/FIGS/data/v2/index.json :: resolution of the regression output of R16_Q1_MASS over the frozen-feature MLP of R16_Q1, run by run, all jets :: sha256 c6c4138ee5c86423 :: prose twin of the run farther from 1
\newcommand{\MlstMassMlpInit}{0.2421\,\ensuremath{\pm}\,0.0018}  % experiments/FIGS/data/v2/index.json :: frozen-feature MLP of INIT, all jets, ATLAS resolution :: sha256 c6c4138ee5c86423
\newcommand{\MlstMassMlpFull}{0.2200\,\ensuremath{\pm}\,0.0014}  % experiments/FIGS/data/v2/index.json :: frozen-feature MLP of L188, all jets, ATLAS resolution :: sha256 c6c4138ee5c86423
\newcommand{\MlstMassMlpRes}{0.2195\,\ensuremath{\pm}\,0.0017}  % experiments/FIGS/data/v2/index.json :: frozen-feature MLP of L162, all jets, ATLAS resolution :: sha256 c6c4138ee5c86423
\newcommand{\MlstMassMlpMid}{0.2158\,\ensuremath{\pm}\,0.0029}  % experiments/FIGS/data/v2/index.json :: frozen-feature MLP of R42_Q1, all jets, ATLAS resolution :: sha256 c6c4138ee5c86423
\newcommand{\MlstMassMlpCoarse}{0.2098\,\ensuremath{\pm}\,0.0031}  % experiments/FIGS/data/v2/index.json :: frozen-feature MLP of R16_Q1, all jets, ATLAS resolution :: sha256 c6c4138ee5c86423
\newcommand{\MlstMassMlpResMass}{0.1591 and 0.1586}  % experiments/FIGS/data/v2/index.json :: frozen-feature MLP of L162_MASS, all jets, ATLAS resolution :: sha256 c6c4138ee5c86423
\newcommand{\MlstMassMlpCoarseMass}{0.1583 and 0.1588}  % experiments/FIGS/data/v2/index.json :: frozen-feature MLP of R16_Q1_MASS, all jets, ATLAS resolution :: sha256 c6c4138ee5c86423
\newcommand{\MlstCmpMlpResOverFull}{0.998\,\ensuremath{\pm}\,0.006}  % experiments/FIGS/data/v2/index.json :: frozen-feature MLP resolution of L162 over L188, run by run, all jets :: sha256 c6c4138ee5c86423
\newcommand{\MlstCmpMlpResOverFullPr}{equal within 1\%}  % experiments/FIGS/data/v2/index.json :: frozen-feature MLP resolution of L162 over L188, run by run, all jets :: sha256 c6c4138ee5c86423 :: prose twin: two significant figures, direction in words
\newcommand{\MlstCmpMlpMidOverFull}{0.981\,\ensuremath{\pm}\,0.017}  % experiments/FIGS/data/v2/index.json :: frozen-feature MLP resolution of R42_Q1 over L188, run by run, all jets :: sha256 c6c4138ee5c86423
\newcommand{\MlstCmpMlpMidOverFullPr}{equal within 2\%}  % experiments/FIGS/data/v2/index.json :: frozen-feature MLP resolution of R42_Q1 over L188, run by run, all jets :: sha256 c6c4138ee5c86423 :: prose twin: two significant figures, direction in words
\newcommand{\MlstCmpMlpCoarseOverFull}{0.954\,\ensuremath{\pm}\,0.018}  % experiments/FIGS/data/v2/index.json :: frozen-feature MLP resolution of R16_Q1 over L188, run by run, all jets :: sha256 c6c4138ee5c86423
\newcommand{\MlstCmpMlpCoarseOverFullPr}{0.95 times the resolution}  % experiments/FIGS/data/v2/index.json :: frozen-feature MLP resolution of R16_Q1 over L188, run by run, all jets :: sha256 c6c4138ee5c86423 :: prose twin: two significant figures, direction in words
\newcommand{\MlstCmpMlpResMassOverRes}{0.7186 and 0.7253}  % experiments/FIGS/data/v2/index.json :: frozen-feature MLP resolution of L162_MASS over L162, run by run, all jets :: sha256 c6c4138ee5c86423
\newcommand{\MlstCmpMlpResMassOverResPr}{at most 0.73 times the resolution}  % experiments/FIGS/data/v2/index.json :: frozen-feature MLP resolution of L162_MASS over L162, run by run, all jets :: sha256 c6c4138ee5c86423 :: prose twin: two significant figures, direction in words
\newcommand{\MlstCmpMlpResMassOverResPrMax}{down to 0.72 times the resolution}  % experiments/FIGS/data/v2/index.json :: frozen-feature MLP resolution of L162_MASS over L162, run by run, all jets :: sha256 c6c4138ee5c86423 :: prose twin of the run farther from 1
\newcommand{\MlstCmpMlpCoarseMassOverCoarse}{0.768 and 0.750}  % experiments/FIGS/data/v2/index.json :: frozen-feature MLP resolution of R16_Q1_MASS over R16_Q1, run by run, all jets :: sha256 c6c4138ee5c86423
\newcommand{\MlstCmpMlpCoarseMassOverCoarsePr}{at most 0.77 times the resolution}  % experiments/FIGS/data/v2/index.json :: frozen-feature MLP resolution of R16_Q1_MASS over R16_Q1, run by run, all jets :: sha256 c6c4138ee5c86423 :: prose twin: two significant figures, direction in words
\newcommand{\MlstCmpMlpCoarseMassOverCoarsePrMax}{down to 0.75 times the resolution}  % experiments/FIGS/data/v2/index.json :: frozen-feature MLP resolution of R16_Q1_MASS over R16_Q1, run by run, all jets :: sha256 c6c4138ee5c86423 :: prose twin of the run farther from 1
\newcommand{\MlstMassUncorrectedResonant}{0.213}  % experiments/FIGS/data/v2/index.json :: head/atlas_res: groomed mass, resonant jets, ATLAS resolution :: sha256 c6c4138ee5c86423
\newcommand{\MlstMassUncorrectedResponseResonant}{0.929}  % experiments/FIGS/data/v2/index.json :: head/atlas_res: groomed mass, resonant jets, median response :: sha256 c6c4138ee5c86423
\newcommand{\MlstMassHeadResMassResonant}{0.1332 and 0.1320}  % experiments/FIGS/data/v2/index.json :: regression output of L162_MASS, resonant jets, ATLAS resolution :: sha256 c6c4138ee5c86423
\newcommand{\MlstMassHeadResMassResponseResonant}{0.996 and 1.004}  % experiments/FIGS/data/v2/index.json :: regression output of L162_MASS, resonant jets, median response :: sha256 c6c4138ee5c86423
\newcommand{\MlstMassHeadResMassOverUncorrResonant}{0.6259 and 0.6205}  % experiments/FIGS/data/v2/index.json :: resolution of the regression output of L162_MASS over the groomed mass's, resonant jets :: sha256 c6c4138ee5c86423
\newcommand{\MlstMassHeadResMassOverUncorrResonantPr}{at most 0.63 times the resolution}  % experiments/FIGS/data/v2/index.json :: resolution of the regression output of L162_MASS over the groomed mass's, resonant jets :: sha256 c6c4138ee5c86423 :: prose twin: two significant figures, direction in words
\newcommand{\MlstMassHeadResMassOverUncorrResonantPrMax}{down to 0.62 times the resolution}  % experiments/FIGS/data/v2/index.json :: resolution of the regression output of L162_MASS over the groomed mass's, resonant jets :: sha256 c6c4138ee5c86423 :: prose twin of the run farther from 1
\newcommand{\MlstMassHeadResMassOverMlpResResonant}{0.6620 and 0.6662}  % experiments/FIGS/data/v2/index.json :: resolution of the regression output of L162_MASS over the frozen-feature MLP of L162, run by run, resonant jets :: sha256 c6c4138ee5c86423
\newcommand{\MlstMassHeadResMassOverMlpResResonantPr}{at most 0.67 times the resolution}  % experiments/FIGS/data/v2/index.json :: resolution of the regression output of L162_MASS over the frozen-feature MLP of L162, run by run, resonant jets :: sha256 c6c4138ee5c86423 :: prose twin: two significant figures, direction in words
\newcommand{\MlstMassHeadResMassOverMlpResResonantPrMax}{down to 0.66 times the resolution}  % experiments/FIGS/data/v2/index.json :: resolution of the regression output of L162_MASS over the frozen-feature MLP of L162, run by run, resonant jets :: sha256 c6c4138ee5c86423 :: prose twin of the run farther from 1
\newcommand{\MlstMassHeadCoarseMassResonant}{0.1312 and 0.1320}  % experiments/FIGS/data/v2/index.json :: regression output of R16_Q1_MASS, resonant jets, ATLAS resolution :: sha256 c6c4138ee5c86423
\newcommand{\MlstMassHeadCoarseMassResponseResonant}{1.006 and 1.006}  % experiments/FIGS/data/v2/index.json :: regression output of R16_Q1_MASS, resonant jets, median response :: sha256 c6c4138ee5c86423
\newcommand{\MlstMassHeadCoarseMassOverUncorrResonant}{0.6165 and 0.6206}  % experiments/FIGS/data/v2/index.json :: resolution of the regression output of R16_Q1_MASS over the groomed mass's, resonant jets :: sha256 c6c4138ee5c86423
\newcommand{\MlstMassHeadCoarseMassOverUncorrResonantPr}{at most 0.62 times the resolution}  % experiments/FIGS/data/v2/index.json :: resolution of the regression output of R16_Q1_MASS over the groomed mass's, resonant jets :: sha256 c6c4138ee5c86423 :: prose twin: two significant figures, direction in words
\newcommand{\MlstMassHeadCoarseMassOverUncorrResonantPrMax}{down to 0.62 times the resolution}  % experiments/FIGS/data/v2/index.json :: resolution of the regression output of R16_Q1_MASS over the groomed mass's, resonant jets :: sha256 c6c4138ee5c86423 :: prose twin of the run farther from 1
\newcommand{\MlstMassHeadCoarseMassOverMlpCoarseResonant}{0.707 and 0.688}  % experiments/FIGS/data/v2/index.json :: resolution of the regression output of R16_Q1_MASS over the frozen-feature MLP of R16_Q1, run by run, resonant jets :: sha256 c6c4138ee5c86423
\newcommand{\MlstMassHeadCoarseMassOverMlpCoarseResonantPr}{at most 0.71 times the resolution}  % experiments/FIGS/data/v2/index.json :: resolution of the regression output of R16_Q1_MASS over the frozen-feature MLP of R16_Q1, run by run, resonant jets :: sha256 c6c4138ee5c86423 :: prose twin: two significant figures, direction in words
\newcommand{\MlstMassHeadCoarseMassOverMlpCoarseResonantPrMax}{down to 0.69 times the resolution}  % experiments/FIGS/data/v2/index.json :: resolution of the regression output of R16_Q1_MASS over the frozen-feature MLP of R16_Q1, run by run, resonant jets :: sha256 c6c4138ee5c86423 :: prose twin of the run farther from 1
\newcommand{\MlstMassMlpInitResonant}{0.2214\,\ensuremath{\pm}\,0.0017}  % experiments/FIGS/data/v2/index.json :: frozen-feature MLP of INIT, resonant jets, ATLAS resolution :: sha256 c6c4138ee5c86423
\newcommand{\MlstMassMlpFullResonant}{0.1990\,\ensuremath{\pm}\,0.0016}  % experiments/FIGS/data/v2/index.json :: frozen-feature MLP of L188, resonant jets, ATLAS resolution :: sha256 c6c4138ee5c86423
\newcommand{\MlstMassMlpResResonant}{0.1989\,\ensuremath{\pm}\,0.0020}  % experiments/FIGS/data/v2/index.json :: frozen-feature MLP of L162, resonant jets, ATLAS resolution :: sha256 c6c4138ee5c86423
\newcommand{\MlstMassMlpMidResonant}{0.1949\,\ensuremath{\pm}\,0.0023}  % experiments/FIGS/data/v2/index.json :: frozen-feature MLP of R42_Q1, resonant jets, ATLAS resolution :: sha256 c6c4138ee5c86423
\newcommand{\MlstMassMlpCoarseResonant}{0.1894\,\ensuremath{\pm}\,0.0035}  % experiments/FIGS/data/v2/index.json :: frozen-feature MLP of R16_Q1, resonant jets, ATLAS resolution :: sha256 c6c4138ee5c86423
\newcommand{\MlstMassMlpResMassResonant}{0.1436 and 0.1430}  % experiments/FIGS/data/v2/index.json :: frozen-feature MLP of L162_MASS, resonant jets, ATLAS resolution :: sha256 c6c4138ee5c86423
\newcommand{\MlstMassMlpCoarseMassResonant}{0.1420 and 0.1431}  % experiments/FIGS/data/v2/index.json :: frozen-feature MLP of R16_Q1_MASS, resonant jets, ATLAS resolution :: sha256 c6c4138ee5c86423
\newcommand{\MlstCmpMlpResOverFullResonant}{1.000\,\ensuremath{\pm}\,0.005}  % experiments/FIGS/data/v2/index.json :: frozen-feature MLP resolution of L162 over L188, run by run, resonant jets :: sha256 c6c4138ee5c86423
\newcommand{\MlstCmpMlpResOverFullResonantPr}{equal within 1\%}  % experiments/FIGS/data/v2/index.json :: frozen-feature MLP resolution of L162 over L188, run by run, resonant jets :: sha256 c6c4138ee5c86423 :: prose twin: two significant figures, direction in words
\newcommand{\MlstCmpMlpMidOverFullResonant}{0.980\,\ensuremath{\pm}\,0.018}  % experiments/FIGS/data/v2/index.json :: frozen-feature MLP resolution of R42_Q1 over L188, run by run, resonant jets :: sha256 c6c4138ee5c86423
\newcommand{\MlstCmpMlpMidOverFullResonantPr}{equal within 3\%}  % experiments/FIGS/data/v2/index.json :: frozen-feature MLP resolution of R42_Q1 over L188, run by run, resonant jets :: sha256 c6c4138ee5c86423 :: prose twin: two significant figures, direction in words
\newcommand{\MlstCmpMlpCoarseOverFullResonant}{0.952\,\ensuremath{\pm}\,0.023}  % experiments/FIGS/data/v2/index.json :: frozen-feature MLP resolution of R16_Q1 over L188, run by run, resonant jets :: sha256 c6c4138ee5c86423
\newcommand{\MlstCmpMlpCoarseOverFullResonantPr}{0.95 times the resolution}  % experiments/FIGS/data/v2/index.json :: frozen-feature MLP resolution of R16_Q1 over L188, run by run, resonant jets :: sha256 c6c4138ee5c86423 :: prose twin: two significant figures, direction in words
\newcommand{\MlstCmpMlpResMassOverResResonant}{0.7137 and 0.7219}  % experiments/FIGS/data/v2/index.json :: frozen-feature MLP resolution of L162_MASS over L162, run by run, resonant jets :: sha256 c6c4138ee5c86423
\newcommand{\MlstCmpMlpResMassOverResResonantPr}{at most 0.72 times the resolution}  % experiments/FIGS/data/v2/index.json :: frozen-feature MLP resolution of L162_MASS over L162, run by run, resonant jets :: sha256 c6c4138ee5c86423 :: prose twin: two significant figures, direction in words
\newcommand{\MlstCmpMlpResMassOverResResonantPrMax}{down to 0.71 times the resolution}  % experiments/FIGS/data/v2/index.json :: frozen-feature MLP resolution of L162_MASS over L162, run by run, resonant jets :: sha256 c6c4138ee5c86423 :: prose twin of the run farther from 1
\newcommand{\MlstCmpMlpCoarseMassOverCoarseResonant}{0.766 and 0.745}  % experiments/FIGS/data/v2/index.json :: frozen-feature MLP resolution of R16_Q1_MASS over R16_Q1, run by run, resonant jets :: sha256 c6c4138ee5c86423
\newcommand{\MlstCmpMlpCoarseMassOverCoarseResonantPr}{at most 0.77 times the resolution}  % experiments/FIGS/data/v2/index.json :: frozen-feature MLP resolution of R16_Q1_MASS over R16_Q1, run by run, resonant jets :: sha256 c6c4138ee5c86423 :: prose twin: two significant figures, direction in words
\newcommand{\MlstCmpMlpCoarseMassOverCoarseResonantPrMax}{down to 0.75 times the resolution}  % experiments/FIGS/data/v2/index.json :: frozen-feature MLP resolution of R16_Q1_MASS over R16_Q1, run by run, resonant jets :: sha256 c6c4138ee5c86423 :: prose twin of the run farther from 1
\newcommand{\MlstMassUncorrectedQcd}{1.52}  % experiments/FIGS/data/v2/index.json :: head/atlas_res: groomed mass, qcd jets, ATLAS resolution :: sha256 c6c4138ee5c86423
\newcommand{\MlstMassUncorrectedResponseQcd}{0.931}  % experiments/FIGS/data/v2/index.json :: head/atlas_res: groomed mass, qcd jets, median response :: sha256 c6c4138ee5c86423
\newcommand{\MlstMassHeadResMassQcd}{0.2263 and 0.2256}  % experiments/FIGS/data/v2/index.json :: regression output of L162_MASS, qcd jets, ATLAS resolution :: sha256 c6c4138ee5c86423
\newcommand{\MlstMassHeadResMassResponseQcd}{0.9897 and 0.9996}  % experiments/FIGS/data/v2/index.json :: regression output of L162_MASS, qcd jets, median response :: sha256 c6c4138ee5c86423
\newcommand{\MlstMassHeadResMassOverUncorrQcd}{0.1485 and 0.1480}  % experiments/FIGS/data/v2/index.json :: resolution of the regression output of L162_MASS over the groomed mass's, qcd jets :: sha256 c6c4138ee5c86423
\newcommand{\MlstMassHeadResMassOverUncorrQcdPr}{at most 0.15 times the resolution}  % experiments/FIGS/data/v2/index.json :: resolution of the regression output of L162_MASS over the groomed mass's, qcd jets :: sha256 c6c4138ee5c86423 :: prose twin: two significant figures, direction in words
\newcommand{\MlstMassHeadResMassOverUncorrQcdPrMax}{down to 0.15 times the resolution}  % experiments/FIGS/data/v2/index.json :: resolution of the regression output of L162_MASS over the groomed mass's, qcd jets :: sha256 c6c4138ee5c86423 :: prose twin of the run farther from 1
\newcommand{\MlstMassHeadResMassOverMlpResQcd}{0.6542 and 0.6632}  % experiments/FIGS/data/v2/index.json :: resolution of the regression output of L162_MASS over the frozen-feature MLP of L162, run by run, qcd jets :: sha256 c6c4138ee5c86423
\newcommand{\MlstMassHeadResMassOverMlpResQcdPr}{at most 0.66 times the resolution}  % experiments/FIGS/data/v2/index.json :: resolution of the regression output of L162_MASS over the frozen-feature MLP of L162, run by run, qcd jets :: sha256 c6c4138ee5c86423 :: prose twin: two significant figures, direction in words
\newcommand{\MlstMassHeadResMassOverMlpResQcdPrMax}{down to 0.65 times the resolution}  % experiments/FIGS/data/v2/index.json :: resolution of the regression output of L162_MASS over the frozen-feature MLP of L162, run by run, qcd jets :: sha256 c6c4138ee5c86423 :: prose twin of the run farther from 1
\newcommand{\MlstMassHeadCoarseMassQcd}{0.2239 and 0.2240}  % experiments/FIGS/data/v2/index.json :: regression output of R16_Q1_MASS, qcd jets, ATLAS resolution :: sha256 c6c4138ee5c86423
\newcommand{\MlstMassHeadCoarseMassResponseQcd}{1.001 and 1.000}  % experiments/FIGS/data/v2/index.json :: regression output of R16_Q1_MASS, qcd jets, median response :: sha256 c6c4138ee5c86423
\newcommand{\MlstMassHeadCoarseMassOverUncorrQcd}{0.1469 and 0.1469}  % experiments/FIGS/data/v2/index.json :: resolution of the regression output of R16_Q1_MASS over the groomed mass's, qcd jets :: sha256 c6c4138ee5c86423
\newcommand{\MlstMassHeadCoarseMassOverUncorrQcdPr}{at most 0.15 times the resolution}  % experiments/FIGS/data/v2/index.json :: resolution of the regression output of R16_Q1_MASS over the groomed mass's, qcd jets :: sha256 c6c4138ee5c86423 :: prose twin: two significant figures, direction in words
\newcommand{\MlstMassHeadCoarseMassOverUncorrQcdPrMax}{down to 0.15 times the resolution}  % experiments/FIGS/data/v2/index.json :: resolution of the regression output of R16_Q1_MASS over the groomed mass's, qcd jets :: sha256 c6c4138ee5c86423 :: prose twin of the run farther from 1
\newcommand{\MlstMassHeadCoarseMassOverMlpCoarseQcd}{0.679 and 0.698}  % experiments/FIGS/data/v2/index.json :: resolution of the regression output of R16_Q1_MASS over the frozen-feature MLP of R16_Q1, run by run, qcd jets :: sha256 c6c4138ee5c86423
\newcommand{\MlstMassHeadCoarseMassOverMlpCoarseQcdPr}{at most 0.70 times the resolution}  % experiments/FIGS/data/v2/index.json :: resolution of the regression output of R16_Q1_MASS over the frozen-feature MLP of R16_Q1, run by run, qcd jets :: sha256 c6c4138ee5c86423 :: prose twin: two significant figures, direction in words
\newcommand{\MlstMassHeadCoarseMassOverMlpCoarseQcdPrMax}{down to 0.68 times the resolution}  % experiments/FIGS/data/v2/index.json :: resolution of the regression output of R16_Q1_MASS over the frozen-feature MLP of R16_Q1, run by run, qcd jets :: sha256 c6c4138ee5c86423 :: prose twin of the run farther from 1
\newcommand{\MlstMassMlpInitQcd}{0.502\,\ensuremath{\pm}\,0.014}  % experiments/FIGS/data/v2/index.json :: frozen-feature MLP of INIT, qcd jets, ATLAS resolution :: sha256 c6c4138ee5c86423
\newcommand{\MlstMassMlpFullQcd}{0.3487\,\ensuremath{\pm}\,0.0029}  % experiments/FIGS/data/v2/index.json :: frozen-feature MLP of L188, qcd jets, ATLAS resolution :: sha256 c6c4138ee5c86423
\newcommand{\MlstMassMlpResQcd}{0.3425\,\ensuremath{\pm}\,0.0031}  % experiments/FIGS/data/v2/index.json :: frozen-feature MLP of L162, qcd jets, ATLAS resolution :: sha256 c6c4138ee5c86423
\newcommand{\MlstMassMlpMidQcd}{0.338\,\ensuremath{\pm}\,0.007}  % experiments/FIGS/data/v2/index.json :: frozen-feature MLP of R42_Q1, qcd jets, ATLAS resolution :: sha256 c6c4138ee5c86423
\newcommand{\MlstMassMlpCoarseQcd}{0.327\,\ensuremath{\pm}\,0.006}  % experiments/FIGS/data/v2/index.json :: frozen-feature MLP of R16_Q1, qcd jets, ATLAS resolution :: sha256 c6c4138ee5c86423
\newcommand{\MlstMassMlpResMassQcd}{0.2473 and 0.2452}  % experiments/FIGS/data/v2/index.json :: frozen-feature MLP of L162_MASS, qcd jets, ATLAS resolution :: sha256 c6c4138ee5c86423
\newcommand{\MlstMassMlpCoarseMassQcd}{0.2465 and 0.2463}  % experiments/FIGS/data/v2/index.json :: frozen-feature MLP of R16_Q1_MASS, qcd jets, ATLAS resolution :: sha256 c6c4138ee5c86423
\newcommand{\MlstCmpMlpResOverFullQcd}{0.982\,\ensuremath{\pm}\,0.017}  % experiments/FIGS/data/v2/index.json :: frozen-feature MLP resolution of L162 over L188, run by run, qcd jets :: sha256 c6c4138ee5c86423
\newcommand{\MlstCmpMlpResOverFullQcdPr}{equal within 2\%}  % experiments/FIGS/data/v2/index.json :: frozen-feature MLP resolution of L162 over L188, run by run, qcd jets :: sha256 c6c4138ee5c86423 :: prose twin: two significant figures, direction in words
\newcommand{\MlstCmpMlpMidOverFullQcd}{0.969\,\ensuremath{\pm}\,0.013}  % experiments/FIGS/data/v2/index.json :: frozen-feature MLP resolution of R42_Q1 over L188, run by run, qcd jets :: sha256 c6c4138ee5c86423
\newcommand{\MlstCmpMlpMidOverFullQcdPr}{0.97 times the resolution}  % experiments/FIGS/data/v2/index.json :: frozen-feature MLP resolution of R42_Q1 over L188, run by run, qcd jets :: sha256 c6c4138ee5c86423 :: prose twin: two significant figures, direction in words
\newcommand{\MlstCmpMlpCoarseOverFullQcd}{0.939\,\ensuremath{\pm}\,0.023}  % experiments/FIGS/data/v2/index.json :: frozen-feature MLP resolution of R16_Q1 over L188, run by run, qcd jets :: sha256 c6c4138ee5c86423
\newcommand{\MlstCmpMlpCoarseOverFullQcdPr}{0.94 times the resolution}  % experiments/FIGS/data/v2/index.json :: frozen-feature MLP resolution of R16_Q1 over L188, run by run, qcd jets :: sha256 c6c4138ee5c86423 :: prose twin: two significant figures, direction in words
\newcommand{\MlstCmpMlpResMassOverResQcd}{0.7147 and 0.7209}  % experiments/FIGS/data/v2/index.json :: frozen-feature MLP resolution of L162_MASS over L162, run by run, qcd jets :: sha256 c6c4138ee5c86423
\newcommand{\MlstCmpMlpResMassOverResQcdPr}{at most 0.72 times the resolution}  % experiments/FIGS/data/v2/index.json :: frozen-feature MLP resolution of L162_MASS over L162, run by run, qcd jets :: sha256 c6c4138ee5c86423 :: prose twin: two significant figures, direction in words
\newcommand{\MlstCmpMlpResMassOverResQcdPrMax}{down to 0.71 times the resolution}  % experiments/FIGS/data/v2/index.json :: frozen-feature MLP resolution of L162_MASS over L162, run by run, qcd jets :: sha256 c6c4138ee5c86423 :: prose twin of the run farther from 1
\newcommand{\MlstCmpMlpCoarseMassOverCoarseQcd}{0.747 and 0.768}  % experiments/FIGS/data/v2/index.json :: frozen-feature MLP resolution of R16_Q1_MASS over R16_Q1, run by run, qcd jets :: sha256 c6c4138ee5c86423
\newcommand{\MlstCmpMlpCoarseMassOverCoarseQcdPr}{at most 0.77 times the resolution}  % experiments/FIGS/data/v2/index.json :: frozen-feature MLP resolution of R16_Q1_MASS over R16_Q1, run by run, qcd jets :: sha256 c6c4138ee5c86423 :: prose twin: two significant figures, direction in words
\newcommand{\MlstCmpMlpCoarseMassOverCoarseQcdPrMax}{down to 0.75 times the resolution}  % experiments/FIGS/data/v2/index.json :: frozen-feature MLP resolution of R16_Q1_MASS over R16_Q1, run by run, qcd jets :: sha256 c6c4138ee5c86423 :: prose twin of the run farther from 1
\newcommand{\MlstMassRespAxesShareQcd}{76\%}  % experiments/FIGS/data/v2/index.json :: head/atlas_res: share of QCD jets whose groomed-mass response lies inside F9a(b)'s x axis :: sha256 c6c4138ee5c86423
\newcommand{\MlstSculptNPassTen}{41{,}028}  % experiments/FIGS/data/v2/index.json :: QCD test jets passing the cut at QCD efficiency 0.1, every run :: sha256 c6c4138ee5c86423
\newcommand{\MlstSculptFloorTen}{0.00944}  % experiments/FIGS/data/v2/index.json :: Jensen-Shannon distance expected for a random selection of as many QCD jets as pass at 0.1 :: sha256 c6c4138ee5c86423
\newcommand{\MlstSculptFullTen}{0.0524\,\ensuremath{\pm}\,0.0034}  % experiments/FIGS/data/v2/index.json :: Jensen-Shannon distance of the QCD m_SD distribution at QCD efficiency 0.1, L188 :: sha256 c6c4138ee5c86423
\newcommand{\MlstSculptFullOverFloorTen}{5.6\,\ensuremath{\pm}\,0.4}  % experiments/FIGS/data/v2/index.json :: distance at 0.1 over the random-selection floor, L188 :: sha256 c6c4138ee5c86423
\newcommand{\MlstSculptFullOverFloorTenPr}{5.6 times the distance}  % experiments/FIGS/data/v2/index.json :: distance at 0.1 over the random-selection floor, L188 :: sha256 c6c4138ee5c86423 :: prose twin: two significant figures, direction in words
\newcommand{\MlstSculptResTen}{0.0532\,\ensuremath{\pm}\,0.0009}  % experiments/FIGS/data/v2/index.json :: Jensen-Shannon distance of the QCD m_SD distribution at QCD efficiency 0.1, L162 :: sha256 c6c4138ee5c86423
\newcommand{\MlstSculptResOverFloorTen}{5.64\,\ensuremath{\pm}\,0.10}  % experiments/FIGS/data/v2/index.json :: distance at 0.1 over the random-selection floor, L162 :: sha256 c6c4138ee5c86423
\newcommand{\MlstSculptResOverFloorTenPr}{5.6 times the distance}  % experiments/FIGS/data/v2/index.json :: distance at 0.1 over the random-selection floor, L162 :: sha256 c6c4138ee5c86423 :: prose twin: two significant figures, direction in words
\newcommand{\MlstSculptMidTen}{0.0560\,\ensuremath{\pm}\,0.0027}  % experiments/FIGS/data/v2/index.json :: Jensen-Shannon distance of the QCD m_SD distribution at QCD efficiency 0.1, R42_Q1 :: sha256 c6c4138ee5c86423
\newcommand{\MlstSculptMidOverFloorTen}{5.94\,\ensuremath{\pm}\,0.28}  % experiments/FIGS/data/v2/index.json :: distance at 0.1 over the random-selection floor, R42_Q1 :: sha256 c6c4138ee5c86423
\newcommand{\MlstSculptMidOverFloorTenPr}{5.9 times the distance}  % experiments/FIGS/data/v2/index.json :: distance at 0.1 over the random-selection floor, R42_Q1 :: sha256 c6c4138ee5c86423 :: prose twin: two significant figures, direction in words
\newcommand{\MlstSculptCoarseTen}{0.0543\,\ensuremath{\pm}\,0.0010}  % experiments/FIGS/data/v2/index.json :: Jensen-Shannon distance of the QCD m_SD distribution at QCD efficiency 0.1, R16_Q1 :: sha256 c6c4138ee5c86423
\newcommand{\MlstSculptCoarseOverFloorTen}{5.76\,\ensuremath{\pm}\,0.10}  % experiments/FIGS/data/v2/index.json :: distance at 0.1 over the random-selection floor, R16_Q1 :: sha256 c6c4138ee5c86423
\newcommand{\MlstSculptCoarseOverFloorTenPr}{5.8 times the distance}  % experiments/FIGS/data/v2/index.json :: distance at 0.1 over the random-selection floor, R16_Q1 :: sha256 c6c4138ee5c86423 :: prose twin: two significant figures, direction in words
\newcommand{\MlstSculptResMassTen}{0.0422 and 0.0464}  % experiments/FIGS/data/v2/index.json :: Jensen-Shannon distance of the QCD m_SD distribution at QCD efficiency 0.1, L162_MASS :: sha256 c6c4138ee5c86423
\newcommand{\MlstSculptResMassOverFloorTen}{4.47 and 4.91}  % experiments/FIGS/data/v2/index.json :: distance at 0.1 over the random-selection floor, L162_MASS :: sha256 c6c4138ee5c86423
\newcommand{\MlstSculptResMassOverFloorTenPr}{at least 4.5 times the distance}  % experiments/FIGS/data/v2/index.json :: distance at 0.1 over the random-selection floor, L162_MASS :: sha256 c6c4138ee5c86423 :: prose twin: two significant figures, direction in words
\newcommand{\MlstSculptResMassOverFloorTenPrMax}{up to 4.9 times the distance}  % experiments/FIGS/data/v2/index.json :: distance at 0.1 over the random-selection floor, L162_MASS :: sha256 c6c4138ee5c86423 :: prose twin of the run farther from 1
\newcommand{\MlstSculptCoarseMassTen}{0.0522 and 0.0423}  % experiments/FIGS/data/v2/index.json :: Jensen-Shannon distance of the QCD m_SD distribution at QCD efficiency 0.1, R16_Q1_MASS :: sha256 c6c4138ee5c86423
\newcommand{\MlstSculptCoarseMassOverFloorTen}{5.5 and 4.5}  % experiments/FIGS/data/v2/index.json :: distance at 0.1 over the random-selection floor, R16_Q1_MASS :: sha256 c6c4138ee5c86423
\newcommand{\MlstSculptCoarseMassOverFloorTenPr}{at least 4.5 times the distance}  % experiments/FIGS/data/v2/index.json :: distance at 0.1 over the random-selection floor, R16_Q1_MASS :: sha256 c6c4138ee5c86423 :: prose twin: two significant figures, direction in words
\newcommand{\MlstSculptCoarseMassOverFloorTenPrMax}{up to 5.5 times the distance}  % experiments/FIGS/data/v2/index.json :: distance at 0.1 over the random-selection floor, R16_Q1_MASS :: sha256 c6c4138ee5c86423 :: prose twin of the run farther from 1
\newcommand{\MlstSculptResMassOverResTen}{0.789 and 0.860}  % experiments/FIGS/data/v2/index.json :: distance at QCD efficiency 0.1: L162_MASS over L162, run by run :: sha256 c6c4138ee5c86423
\newcommand{\MlstSculptResMassOverResTenPr}{at most 0.86 times the distance}  % experiments/FIGS/data/v2/index.json :: distance at QCD efficiency 0.1: L162_MASS over L162, run by run :: sha256 c6c4138ee5c86423 :: prose twin: two significant figures, direction in words
\newcommand{\MlstSculptResMassOverResTenPrMax}{down to 0.79 times the distance}  % experiments/FIGS/data/v2/index.json :: distance at QCD efficiency 0.1: L162_MASS over L162, run by run :: sha256 c6c4138ee5c86423 :: prose twin of the run farther from 1
\newcommand{\MlstSculptCoarseMassOverCoarseTen}{0.94 and 0.79}  % experiments/FIGS/data/v2/index.json :: distance at QCD efficiency 0.1: R16_Q1_MASS over R16_Q1, run by run :: sha256 c6c4138ee5c86423
\newcommand{\MlstSculptCoarseMassOverCoarseTenPr}{equal within 30\%}  % experiments/FIGS/data/v2/index.json :: distance at QCD efficiency 0.1: R16_Q1_MASS over R16_Q1, run by run :: sha256 c6c4138ee5c86423 :: prose twin: two significant figures, direction in words
\newcommand{\MlstSculptCoarseMassOverCoarseTenPrMax}{equal within 30\%}  % experiments/FIGS/data/v2/index.json :: distance at QCD efficiency 0.1: R16_Q1_MASS over R16_Q1, run by run :: sha256 c6c4138ee5c86423 :: prose twin of the run farther from 1
\newcommand{\MlstSculptCoarseOverFullTen}{1.04\,\ensuremath{\pm}\,0.05}  % experiments/FIGS/data/v2/index.json :: distance at QCD efficiency 0.1: R16_Q1 over L188, run by run :: sha256 c6c4138ee5c86423
\newcommand{\MlstSculptCoarseOverFullTenPr}{equal within 6\%}  % experiments/FIGS/data/v2/index.json :: distance at QCD efficiency 0.1: R16_Q1 over L188, run by run :: sha256 c6c4138ee5c86423 :: prose twin: two significant figures, direction in words
\newcommand{\MlstSculptResOverFullTen}{1.02\,\ensuremath{\pm}\,0.06}  % experiments/FIGS/data/v2/index.json :: distance at QCD efficiency 0.1: L162 over L188, run by run :: sha256 c6c4138ee5c86423
\newcommand{\MlstSculptResOverFullTenPr}{equal within 6\%}  % experiments/FIGS/data/v2/index.json :: distance at QCD efficiency 0.1: L162 over L188, run by run :: sha256 c6c4138ee5c86423 :: prose twin: two significant figures, direction in words
\newcommand{\MlstSculptMidOverFullTen}{1.069\,\ensuremath{\pm}\,0.020}  % experiments/FIGS/data/v2/index.json :: distance at QCD efficiency 0.1: R42_Q1 over L188, run by run :: sha256 c6c4138ee5c86423
\newcommand{\MlstSculptMidOverFullTenPr}{1.1 times the distance}  % experiments/FIGS/data/v2/index.json :: distance at QCD efficiency 0.1: R42_Q1 over L188, run by run :: sha256 c6c4138ee5c86423 :: prose twin: two significant figures, direction in words
\newcommand{\MlstSculptCoarseOverMidTen}{0.971\,\ensuremath{\pm}\,0.035}  % experiments/FIGS/data/v2/index.json :: distance at QCD efficiency 0.1: R16_Q1 over R42_Q1, run by run :: sha256 c6c4138ee5c86423
\newcommand{\MlstSculptCoarseOverMidTenPr}{equal within 4\%}  % experiments/FIGS/data/v2/index.json :: distance at QCD efficiency 0.1: R16_Q1 over R42_Q1, run by run :: sha256 c6c4138ee5c86423 :: prose twin: two significant figures, direction in words
\newcommand{\MlstSculptNPassOne}{4{,}103}  % experiments/FIGS/data/v2/index.json :: QCD test jets passing the cut at QCD efficiency 0.01, every run :: sha256 c6c4138ee5c86423
\newcommand{\MlstSculptFloorOne}{0.0313}  % experiments/FIGS/data/v2/index.json :: Jensen-Shannon distance expected for a random selection of as many QCD jets as pass at 0.01 :: sha256 c6c4138ee5c86423
\newcommand{\MlstSculptFullOne}{0.0439\,\ensuremath{\pm}\,0.0024}  % experiments/FIGS/data/v2/index.json :: Jensen-Shannon distance of the QCD m_SD distribution at QCD efficiency 0.01, L188 :: sha256 c6c4138ee5c86423
\newcommand{\MlstSculptFullOverFloorOne}{1.40\,\ensuremath{\pm}\,0.08}  % experiments/FIGS/data/v2/index.json :: distance at 0.01 over the random-selection floor, L188 :: sha256 c6c4138ee5c86423
\newcommand{\MlstSculptFullOverFloorOnePr}{1.4 times the distance}  % experiments/FIGS/data/v2/index.json :: distance at 0.01 over the random-selection floor, L188 :: sha256 c6c4138ee5c86423 :: prose twin: two significant figures, direction in words
\newcommand{\MlstSculptResOne}{0.0429\,\ensuremath{\pm}\,0.0029}  % experiments/FIGS/data/v2/index.json :: Jensen-Shannon distance of the QCD m_SD distribution at QCD efficiency 0.01, L162 :: sha256 c6c4138ee5c86423
\newcommand{\MlstSculptResOverFloorOne}{1.37\,\ensuremath{\pm}\,0.09}  % experiments/FIGS/data/v2/index.json :: distance at 0.01 over the random-selection floor, L162 :: sha256 c6c4138ee5c86423
\newcommand{\MlstSculptResOverFloorOnePr}{1.4 times the distance}  % experiments/FIGS/data/v2/index.json :: distance at 0.01 over the random-selection floor, L162 :: sha256 c6c4138ee5c86423 :: prose twin: two significant figures, direction in words
\newcommand{\MlstSculptMidOne}{0.051\,\ensuremath{\pm}\,0.007}  % experiments/FIGS/data/v2/index.json :: Jensen-Shannon distance of the QCD m_SD distribution at QCD efficiency 0.01, R42_Q1 :: sha256 c6c4138ee5c86423
\newcommand{\MlstSculptMidOverFloorOne}{1.62\,\ensuremath{\pm}\,0.23}  % experiments/FIGS/data/v2/index.json :: distance at 0.01 over the random-selection floor, R42_Q1 :: sha256 c6c4138ee5c86423
\newcommand{\MlstSculptMidOverFloorOnePr}{1.6 times the distance}  % experiments/FIGS/data/v2/index.json :: distance at 0.01 over the random-selection floor, R42_Q1 :: sha256 c6c4138ee5c86423 :: prose twin: two significant figures, direction in words
\newcommand{\MlstSculptCoarseOne}{0.0561\,\ensuremath{\pm}\,0.0017}  % experiments/FIGS/data/v2/index.json :: Jensen-Shannon distance of the QCD m_SD distribution at QCD efficiency 0.01, R16_Q1 :: sha256 c6c4138ee5c86423
\newcommand{\MlstSculptCoarseOverFloorOne}{1.79\,\ensuremath{\pm}\,0.06}  % experiments/FIGS/data/v2/index.json :: distance at 0.01 over the random-selection floor, R16_Q1 :: sha256 c6c4138ee5c86423
\newcommand{\MlstSculptCoarseOverFloorOnePr}{1.8 times the distance}  % experiments/FIGS/data/v2/index.json :: distance at 0.01 over the random-selection floor, R16_Q1 :: sha256 c6c4138ee5c86423 :: prose twin: two significant figures, direction in words
\newcommand{\MlstSculptResMassOne}{0.0424 and 0.0442}  % experiments/FIGS/data/v2/index.json :: Jensen-Shannon distance of the QCD m_SD distribution at QCD efficiency 0.01, L162_MASS :: sha256 c6c4138ee5c86423
\newcommand{\MlstSculptResMassOverFloorOne}{1.356 and 1.411}  % experiments/FIGS/data/v2/index.json :: distance at 0.01 over the random-selection floor, L162_MASS :: sha256 c6c4138ee5c86423
\newcommand{\MlstSculptResMassOverFloorOnePr}{at least 1.4 times the distance}  % experiments/FIGS/data/v2/index.json :: distance at 0.01 over the random-selection floor, L162_MASS :: sha256 c6c4138ee5c86423 :: prose twin: two significant figures, direction in words
\newcommand{\MlstSculptResMassOverFloorOnePrMax}{up to 1.4 times the distance}  % experiments/FIGS/data/v2/index.json :: distance at 0.01 over the random-selection floor, L162_MASS :: sha256 c6c4138ee5c86423 :: prose twin of the run farther from 1
\newcommand{\MlstSculptCoarseMassOne}{0.0405 and 0.0452}  % experiments/FIGS/data/v2/index.json :: Jensen-Shannon distance of the QCD m_SD distribution at QCD efficiency 0.01, R16_Q1_MASS :: sha256 c6c4138ee5c86423
\newcommand{\MlstSculptCoarseMassOverFloorOne}{1.29 and 1.44}  % experiments/FIGS/data/v2/index.json :: distance at 0.01 over the random-selection floor, R16_Q1_MASS :: sha256 c6c4138ee5c86423
\newcommand{\MlstSculptCoarseMassOverFloorOnePr}{at least 1.3 times the distance}  % experiments/FIGS/data/v2/index.json :: distance at 0.01 over the random-selection floor, R16_Q1_MASS :: sha256 c6c4138ee5c86423 :: prose twin: two significant figures, direction in words
\newcommand{\MlstSculptCoarseMassOverFloorOnePrMax}{up to 1.4 times the distance}  % experiments/FIGS/data/v2/index.json :: distance at 0.01 over the random-selection floor, R16_Q1_MASS :: sha256 c6c4138ee5c86423 :: prose twin of the run farther from 1
\newcommand{\MlstSculptResMassOverResOne}{0.962 and 0.983}  % experiments/FIGS/data/v2/index.json :: distance at QCD efficiency 0.01: L162_MASS over L162, run by run :: sha256 c6c4138ee5c86423
\newcommand{\MlstSculptResMassOverResOnePr}{equal within 4\%}  % experiments/FIGS/data/v2/index.json :: distance at QCD efficiency 0.01: L162_MASS over L162, run by run :: sha256 c6c4138ee5c86423 :: prose twin: two significant figures, direction in words
\newcommand{\MlstSculptResMassOverResOnePrMax}{equal within 4\%}  % experiments/FIGS/data/v2/index.json :: distance at QCD efficiency 0.01: L162_MASS over L162, run by run :: sha256 c6c4138ee5c86423 :: prose twin of the run farther from 1
\newcommand{\MlstSculptCoarseMassOverCoarseOne}{0.70 and 0.82}  % experiments/FIGS/data/v2/index.json :: distance at QCD efficiency 0.01: R16_Q1_MASS over R16_Q1, run by run :: sha256 c6c4138ee5c86423
\newcommand{\MlstSculptCoarseMassOverCoarseOnePr}{at most 0.82 times the distance}  % experiments/FIGS/data/v2/index.json :: distance at QCD efficiency 0.01: R16_Q1_MASS over R16_Q1, run by run :: sha256 c6c4138ee5c86423 :: prose twin: two significant figures, direction in words
\newcommand{\MlstSculptCoarseMassOverCoarseOnePrMax}{down to 0.70 times the distance}  % experiments/FIGS/data/v2/index.json :: distance at QCD efficiency 0.01: R16_Q1_MASS over R16_Q1, run by run :: sha256 c6c4138ee5c86423 :: prose twin of the run farther from 1
\newcommand{\MlstSculptCoarseOverFullOne}{1.28\,\ensuremath{\pm}\,0.04}  % experiments/FIGS/data/v2/index.json :: distance at QCD efficiency 0.01: R16_Q1 over L188, run by run :: sha256 c6c4138ee5c86423
\newcommand{\MlstSculptCoarseOverFullOnePr}{1.3 times the distance}  % experiments/FIGS/data/v2/index.json :: distance at QCD efficiency 0.01: R16_Q1 over L188, run by run :: sha256 c6c4138ee5c86423 :: prose twin: two significant figures, direction in words
\newcommand{\MlstSculptResOverFullOne}{0.98\,\ensuremath{\pm}\,0.05}  % experiments/FIGS/data/v2/index.json :: distance at QCD efficiency 0.01: L162 over L188, run by run :: sha256 c6c4138ee5c86423
\newcommand{\MlstSculptResOverFullOnePr}{equal within 5\%}  % experiments/FIGS/data/v2/index.json :: distance at QCD efficiency 0.01: L162 over L188, run by run :: sha256 c6c4138ee5c86423 :: prose twin: two significant figures, direction in words
\newcommand{\MlstSculptMidOverFullOne}{1.15\,\ensuremath{\pm}\,0.12}  % experiments/FIGS/data/v2/index.json :: distance at QCD efficiency 0.01: R42_Q1 over L188, run by run :: sha256 c6c4138ee5c86423
\newcommand{\MlstSculptMidOverFullOnePr}{equal within 20\%}  % experiments/FIGS/data/v2/index.json :: distance at QCD efficiency 0.01: R42_Q1 over L188, run by run :: sha256 c6c4138ee5c86423 :: prose twin: two significant figures, direction in words
\newcommand{\MlstSculptCoarseOverMidOne}{1.13\,\ensuremath{\pm}\,0.16}  % experiments/FIGS/data/v2/index.json :: distance at QCD efficiency 0.01: R16_Q1 over R42_Q1, run by run :: sha256 c6c4138ee5c86423
\newcommand{\MlstSculptCoarseOverMidOnePr}{equal within 20\%}  % experiments/FIGS/data/v2/index.json :: distance at QCD efficiency 0.01: R16_Q1 over R42_Q1, run by run :: sha256 c6c4138ee5c86423 :: prose twin: two significant figures, direction in words
\newcommand{\MlstSculptDiffEffLo}{0.52\%}  % experiments/FIGS/data/v2/index.json :: lowest stored QCD efficiency of the unbroken stretch around the tighter working point at which the 17/188 distance ratio exceeds 1 by twice its SD :: sha256 c6c4138ee5c86423 :: passing grid points: 0.001, 0.005245, 0.007937, 0.012011
\newcommand{\MlstSculptDiffEffHi}{1.2\%}  % experiments/FIGS/data/v2/index.json :: highest efficiency of that stretch :: sha256 c6c4138ee5c86423
\newcommand{\MlstSculptLowestEff}{0.10\%}  % experiments/FIGS/data/v2/index.json :: lowest stored QCD efficiency of the sculpting curves :: sha256 c6c4138ee5c86423
\newcommand{\MlstSculptLowestNPass}{410}  % experiments/FIGS/data/v2/index.json :: QCD test jets passing at the lowest stored efficiency :: sha256 c6c4138ee5c86423
\newcommand{\MlstSculptLowestMaxFloorSd}{1.7}  % experiments/FIGS/data/v2/index.json :: largest distance of a configuration's mean from the random-selection floor at the lowest stored efficiency, in SDs of the floor over random selections :: sha256 c6c4138ee5c86423
\newcommand{\MlstSculptNConfigs}{6}  % experiments/FIGS/data/v2/index.json :: configurations with a sculpting curve (count) :: sha256 c6c4138ee5c86423
\newcommand{\LpJciiTestQcdPercent}{20.5\%}  % experiments/FIGS/data/v2/jc2_test/counts.json :: QCD share of the jets, per cent :: sha256 3574ba1f6f2a59f0
\newcommand{\LpJciiTestTwoProngPercent}{38.7\%}  % experiments/FIGS/data/v2/jc2_test/counts.json :: two-prong share of the jets, per cent :: sha256 3574ba1f6f2a59f0
\newcommand{\LpJciiTrainQcdPercent}{13.5\%}  % experiments/FT/data/jc2_v2_manifest.json :: QCD share of the jets, per cent :: sha256 4bd68655919b14e9
\newcommand{\LpJciiTrainTwoProngPercent}{15.4\%}  % experiments/FT/data/jc2_v2_manifest.json :: two-prong share of the jets, per cent :: sha256 4bd68655919b14e9
\newcommand{\LpJciiTestMinAucJets}{10}  % experiments/FIGS/data/v2/jc2_test/counts.json :: fewest jets of a class among the macro-AUC rows :: sha256 3574ba1f6f2a59f0
\newcommand{\LpJciiTrainMinJets}{11}  % experiments/FT/data/jc2_v2_manifest.json :: fewest jets of a class in the 10^6 training subset :: sha256 4bd68655919b14e9
\newcommand{\LpNTestJci}{2{,}000{,}000}  % experiments/FIGS/data/v2/index.json :: jc1 test jets the probes score :: sha256 c6c4138ee5c86423
\newcommand{\LpNTrainMaxJci}{1{,}000{,}000}  % experiments/FIGS/data/v2/index.json :: jc1 largest training subset :: sha256 c6c4138ee5c86423
\newcommand{\LpNTestTop}{404{,}000}  % experiments/FIGS/data/v2/index.json :: top test jets the probes score :: sha256 c6c4138ee5c86423
\newcommand{\LpNTrainMaxTop}{1{,}200{,}000}  % experiments/FIGS/data/v2/index.json :: top largest training subset :: sha256 c6c4138ee5c86423
\newcommand{\LpNTestQg}{200{,}000}  % experiments/FIGS/data/v2/index.json :: qg test jets the probes score :: sha256 c6c4138ee5c86423
\newcommand{\LpNTrainMaxQg}{1{,}600{,}000}  % experiments/FIGS/data/v2/index.json :: qg largest training subset :: sha256 c6c4138ee5c86423
\newcommand{\LpNTestAucJci}{2{,}000{,}000}  % experiments/FIGS/data/v2/index.json :: JetClass test jets of the macro AUC :: sha256 c6c4138ee5c86423
\newcommand{\LpPairNValMin}{1{,}068}  % experiments/FIGS/data/v2/index.json :: validation jets of the pair probe with the fewest :: sha256 c6c4138ee5c86423
\newcommand{\LpPairNValMax}{20{,}195}  % experiments/FIGS/data/v2/index.json :: validation jets of the pair probe with the most :: sha256 c6c4138ee5c86423
\newcommand{\LpNValJci}{200{,}000}  % experiments/FIGS/data/v2/index.json :: JetClass probe validation jets :: sha256 c6c4138ee5c86423
\newcommand{\SculptEffHigh}{10\%}  % experiments/EVAL/mass_sculpting.py :: WORKING_POINTS :: sha256 e85d125c85dac054 :: QCD efficiency of a sculpting cut
\newcommand{\SculptEffLow}{1\%}  % experiments/EVAL/mass_sculpting.py :: WORKING_POINTS :: sha256 e85d125c85dac054 :: QCD efficiency of a sculpting cut
\newcommand{\LpMlpHidden}{256}  % experiments/EVAL/probe.py :: fit_mlp hidden layer :: sha256 485db0966e660588 :: units of the MLP probe's hidden layer
\newcommand{\LpMlpSeedsWord}{three}  % experiments/EVAL/probe.py :: MLP_SEEDS :: sha256 485db0966e660588 :: seeds of the MLP probe whose scores are averaged
\newcommand{\MassMinTrainPerClass}{50}  % experiments/EVAL/mass_resolution.py :: MIN_TRAIN_PER_CLASS :: sha256 26acb0e315094bd1 :: fewest training jets of a class the mass regression keeps
\newcommand{\MlstFtLr}{\ensuremath{10^{-4}}}  % experiments/FIGS/data/v2/ft_mlst/manifests.json :: cells[best70].lr :: sha256 a689f21e894f48ae
\newcommand{\MlstFtHeadMult}{50}  % experiments/FIGS/data/v2/ft_mlst/manifests.json :: cells[best70].head_lr_mult :: sha256 a689f21e894f48ae
\newcommand{\MlstFtWeightDecay}{0.01}  % experiments/FIGS/data/v2/ft_mlst/manifests.json :: cells[best70].weight_decay :: sha256 a689f21e894f48ae
\newcommand{\MlstFtBatch}{512}  % experiments/FIGS/data/v2/ft_mlst/manifests.json :: cells[best70].batch_size :: sha256 a689f21e894f48ae
\newcommand{\MlstFtBenchEpochs}{20}  % experiments/FIGS/data/v2/ft_mlst/manifests.json :: cells[best70, top/qg].epochs :: sha256 a689f21e894f48ae
\newcommand{\MlstFtEpochsEThree}{50}  % experiments/FIGS/data/v2/ft_mlst/manifests.json :: cells[best70, leg 1/2, n_train 1000].epochs :: sha256 a689f21e894f48ae
\newcommand{\MlstFtEpochsEFour}{50}  % experiments/FIGS/data/v2/ft_mlst/manifests.json :: cells[best70, leg 1/2, n_train 10000].epochs :: sha256 a689f21e894f48ae
\newcommand{\MlstFtEpochsEFive}{30}  % experiments/FIGS/data/v2/ft_mlst/manifests.json :: cells[best70, leg 1/2, n_train 100000].epochs :: sha256 a689f21e894f48ae
\newcommand{\MlstFtEpochsESix}{10}  % experiments/FIGS/data/v2/ft_mlst/manifests.json :: cells[best70, leg 1/2, n_train 1000000].epochs :: sha256 a689f21e894f48ae
\newcommand{\MlstFtHeadDecayShare}{30\%}  % experiments/FT/k8s/job-ft-v2-legs-best70-l188-s3-raunav.yaml :: no --lr-scheduler: weaver 0.4.17 flat+decay, the head's last epochs :: sha256 19f59aa54a385c11
\newcommand{\MlstFtHeadDecayTo}{1\%}  % experiments/FT/k8s/job-ft-v2-legs-best70-l188-s3-raunav.yaml :: no --lr-scheduler: weaver 0.4.17 flat+decay, final factor :: sha256 19f59aa54a385c11
\newcommand{\MlstFtHeadDropout}{0.1}  % experiments/FT/k8s/job-ft-v2-legs-best70-l188-s3-raunav.yaml :: fc_params dropout of the fine-tuning head :: sha256 19f59aa54a385c11
\newcommand{\MlstFtMinEpochJets}{10{,}000}  % experiments/FIGS/data/v2/ft_mlst/manifests.json :: cells[best70].samples_per_epoch (smallest) :: sha256 a689f21e894f48ae
\newcommand{\MlstFtValJcii}{20{,}480}  % experiments/FIGS/data/v2/ft_mlst/manifests.json :: cells[best70, leg 1].val_jets :: sha256 a689f21e894f48ae
\newcommand{\MlstFtValSmall}{20{,}000}  % experiments/FT/k8s/job-ft-v2-legs-best70-l188-s3-raunav.yaml :: JetClass --samples-per-epoch-val; top and q/g at N <= 1e4 :: sha256 19f59aa54a385c11
\newcommand{\MlstFtValLarge}{200{,}000}  % experiments/FT/k8s/job-ft-v2-bench-best70-l188-s1-top-raunav.yaml :: top and q/g --samples-per-epoch-val at N > 1e4 :: sha256 1e38d4f542119ea3
\newcommand{\MlstFtGpuTypes}{V100}  % experiments/FIGS/data/v2/ft_mlst/manifests.json :: cells[best70].gpu_device_name :: sha256 a689f21e894f48ae
\newcommand{\MlstFtCellsDone}{445}  % experiments/FIGS/data/v2/ft_mlst/manifests.json :: DONE paper-run fine-tuning cells (count) :: sha256 a689f21e894f48ae
\newcommand{\MlstFtRefLr}{\ensuremath{10^{-3}}}  % experiments/FIGS/data/v2/ft_mlst/manifests.json :: cells[ref_v100].lr :: sha256 a689f21e894f48ae
\newcommand{\MlstFtRefSeedsWord}{three}  % experiments/FIGS/data/v2/ft_mlst/manifests.json :: cells[ref_v100].ft_seed (count) :: sha256 a689f21e894f48ae
\newcommand{\MlstFtRefGpuTypes}{V100}  % experiments/FIGS/data/v2/ft_mlst/manifests.json :: cells[ref_v100].gpu_device_name :: sha256 a689f21e894f48ae
\newcommand{\MlstFtRefWeightDecayBench}{0.01}  % experiments/FIGS/data/v2/ft_mlst/manifests.json :: cells[ref_v100, top/qg].weight_decay :: sha256 a689f21e894f48ae
\newcommand{\MlstBootResamples}{200}  % experiments/FIGS/data/v2/index.json :: bootstrap resamples of the test jets :: sha256 c6c4138ee5c86423
\newcommand{\MlstBootTestMaxProbe}{0.21}  % experiments/FIGS/data/v2/index.json :: largest test-sample SD of the ln ratio of 1 - AUC on JetClass-II, linear probe, label sets over the 188-class model, per cent :: sha256 c6c4138ee5c86423
\newcommand{\MlstBootTestMaxFt}{0.15}  % experiments/FIGS/data/v2/index.json :: largest test-sample SD of the ln ratio of 1 - AUC on JetClass-II, fine-tuning, label sets over the 188-class model, per cent :: sha256 c6c4138ee5c86423
\newcommand{\MlstBootTestMaxRthirtyProbe}{3.6}  % experiments/FIGS/data/v2/index.json :: largest test-sample SD of the ln R30 ratio on top and quark/gluon, linear probe (counting), label sets over the 188-class model, per cent :: sha256 c6c4138ee5c86423
\newcommand{\MlstBootTestMaxRthirtyFt}{13}  % experiments/FIGS/data/v2/index.json :: largest test-sample SD of the ln R30 ratio on top and quark/gluon, fine-tuning (counting), label sets over the 188-class model, per cent :: sha256 c6c4138ee5c86423
\newcommand{\MlstFtCmpJcTwoResOverFullEFour}{0.989\,\ensuremath{\pm}\,0.011}  % experiments/FIGS/data/v2/index.json :: fine-tuning jc2 N=10000: (1 - macro_auc_ovr) of L162 over L188, run by run :: sha256 c6c4138ee5c86423
\newcommand{\MlstFtCmpJcTwoResOverFullEFourPr}{equal within 2\%}  % experiments/FIGS/data/v2/index.json :: fine-tuning jc2 N=10000: (1 - macro_auc_ovr) of L162 over L188, run by run :: sha256 c6c4138ee5c86423 :: prose twin: two significant figures, direction in words
\newcommand{\MlstFtCmpJcTwoMidOverFullEFour}{1.191\,\ensuremath{\pm}\,0.030}  % experiments/FIGS/data/v2/index.json :: fine-tuning jc2 N=10000: (1 - macro_auc_ovr) of R42_Q1 over L188, run by run :: sha256 c6c4138ee5c86423
\newcommand{\MlstFtCmpJcTwoMidOverFullEFourPr}{19\% more errors}  % experiments/FIGS/data/v2/index.json :: fine-tuning jc2 N=10000: (1 - macro_auc_ovr) of R42_Q1 over L188, run by run :: sha256 c6c4138ee5c86423 :: prose twin: two significant figures, direction in words
\newcommand{\MlstFtCmpJcTwoCoarseOverFullEFour}{1.353\,\ensuremath{\pm}\,0.003}  % experiments/FIGS/data/v2/index.json :: fine-tuning jc2 N=10000: (1 - macro_auc_ovr) of R16_Q1 over L188, run by run :: sha256 c6c4138ee5c86423
\newcommand{\MlstFtCmpJcTwoCoarseOverFullEFourPr}{35\% more errors}  % experiments/FIGS/data/v2/index.json :: fine-tuning jc2 N=10000: (1 - macro_auc_ovr) of R16_Q1 over L188, run by run :: sha256 c6c4138ee5c86423 :: prose twin: two significant figures, direction in words
\newcommand{\MlstFtCmpJcTwoResMassOverResEFour}{1.016 and 1.005}  % experiments/FIGS/data/v2/index.json :: fine-tuning jc2 N=10000: (1 - macro_auc_ovr) of L162_MASS over L162, run by run :: sha256 c6c4138ee5c86423
\newcommand{\MlstFtCmpJcTwoResMassOverResEFourPr}{equal within 2\%}  % experiments/FIGS/data/v2/index.json :: fine-tuning jc2 N=10000: (1 - macro_auc_ovr) of L162_MASS over L162, run by run :: sha256 c6c4138ee5c86423 :: prose twin: two significant figures, direction in words
\newcommand{\MlstFtCmpJcTwoResMassOverResEFourPrMax}{equal within 2\%}  % experiments/FIGS/data/v2/index.json :: fine-tuning jc2 N=10000: (1 - macro_auc_ovr) of L162_MASS over L162, run by run :: sha256 c6c4138ee5c86423 :: prose twin of the run farther from 1
\newcommand{\MlstFtCmpJcTwoCoarseMassOverCoarseEFour}{0.989 and 1.000}  % experiments/FIGS/data/v2/index.json :: fine-tuning jc2 N=10000: (1 - macro_auc_ovr) of R16_Q1_MASS over R16_Q1, run by run :: sha256 c6c4138ee5c86423
\newcommand{\MlstFtCmpJcTwoCoarseMassOverCoarseEFourPr}{equal within 2\%}  % experiments/FIGS/data/v2/index.json :: fine-tuning jc2 N=10000: (1 - macro_auc_ovr) of R16_Q1_MASS over R16_Q1, run by run :: sha256 c6c4138ee5c86423 :: prose twin: two significant figures, direction in words
\newcommand{\MlstFtCmpJcTwoCoarseMassOverCoarseEFourPrMax}{equal within 2\%}  % experiments/FIGS/data/v2/index.json :: fine-tuning jc2 N=10000: (1 - macro_auc_ovr) of R16_Q1_MASS over R16_Q1, run by run :: sha256 c6c4138ee5c86423 :: prose twin of the run farther from 1
\newcommand{\MlstFtCmpJcTwoResOverFullEFive}{1.011\,\ensuremath{\pm}\,0.005}  % experiments/FIGS/data/v2/index.json :: fine-tuning jc2 N=100000: (1 - macro_auc_ovr) of L162 over L188, run by run :: sha256 c6c4138ee5c86423
\newcommand{\MlstFtCmpJcTwoResOverFullEFivePr}{1.1\% more errors}  % experiments/FIGS/data/v2/index.json :: fine-tuning jc2 N=100000: (1 - macro_auc_ovr) of L162 over L188, run by run :: sha256 c6c4138ee5c86423 :: prose twin: two significant figures, direction in words
\newcommand{\MlstFtCmpJcTwoMidOverFullEFive}{1.064\,\ensuremath{\pm}\,0.009}  % experiments/FIGS/data/v2/index.json :: fine-tuning jc2 N=100000: (1 - macro_auc_ovr) of R42_Q1 over L188, run by run :: sha256 c6c4138ee5c86423
\newcommand{\MlstFtCmpJcTwoMidOverFullEFivePr}{6.4\% more errors}  % experiments/FIGS/data/v2/index.json :: fine-tuning jc2 N=100000: (1 - macro_auc_ovr) of R42_Q1 over L188, run by run :: sha256 c6c4138ee5c86423 :: prose twin: two significant figures, direction in words
\newcommand{\MlstFtCmpJcTwoCoarseOverFullEFive}{1.157\,\ensuremath{\pm}\,0.001}  % experiments/FIGS/data/v2/index.json :: fine-tuning jc2 N=100000: (1 - macro_auc_ovr) of R16_Q1 over L188, run by run :: sha256 c6c4138ee5c86423
\newcommand{\MlstFtCmpJcTwoCoarseOverFullEFivePr}{16\% more errors}  % experiments/FIGS/data/v2/index.json :: fine-tuning jc2 N=100000: (1 - macro_auc_ovr) of R16_Q1 over L188, run by run :: sha256 c6c4138ee5c86423 :: prose twin: two significant figures, direction in words
\newcommand{\MlstFtCmpJcTwoResMassOverResEFive}{1.027 and 1.021}  % experiments/FIGS/data/v2/index.json :: fine-tuning jc2 N=100000: (1 - macro_auc_ovr) of L162_MASS over L162, run by run :: sha256 c6c4138ee5c86423
\newcommand{\MlstFtCmpJcTwoResMassOverResEFivePr}{at least 2.1\% more errors}  % experiments/FIGS/data/v2/index.json :: fine-tuning jc2 N=100000: (1 - macro_auc_ovr) of L162_MASS over L162, run by run :: sha256 c6c4138ee5c86423 :: prose twin: two significant figures, direction in words
\newcommand{\MlstFtCmpJcTwoResMassOverResEFivePrMax}{up to 2.7\% more errors}  % experiments/FIGS/data/v2/index.json :: fine-tuning jc2 N=100000: (1 - macro_auc_ovr) of L162_MASS over L162, run by run :: sha256 c6c4138ee5c86423 :: prose twin of the run farther from 1
\newcommand{\MlstFtCmpJcTwoCoarseMassOverCoarseEFive}{1.029 and 1.022}  % experiments/FIGS/data/v2/index.json :: fine-tuning jc2 N=100000: (1 - macro_auc_ovr) of R16_Q1_MASS over R16_Q1, run by run :: sha256 c6c4138ee5c86423
\newcommand{\MlstFtCmpJcTwoCoarseMassOverCoarseEFivePr}{at least 2.2\% more errors}  % experiments/FIGS/data/v2/index.json :: fine-tuning jc2 N=100000: (1 - macro_auc_ovr) of R16_Q1_MASS over R16_Q1, run by run :: sha256 c6c4138ee5c86423 :: prose twin: two significant figures, direction in words
\newcommand{\MlstFtCmpJcTwoCoarseMassOverCoarseEFivePrMax}{up to 2.9\% more errors}  % experiments/FIGS/data/v2/index.json :: fine-tuning jc2 N=100000: (1 - macro_auc_ovr) of R16_Q1_MASS over R16_Q1, run by run :: sha256 c6c4138ee5c86423 :: prose twin of the run farther from 1
\newcommand{\MlstFtCmpJcTwoResOverFullESix}{1.014\,\ensuremath{\pm}\,0.011}  % experiments/FIGS/data/v2/index.json :: fine-tuning jc2 N=1000000: (1 - macro_auc_ovr) of L162 over L188, run by run :: sha256 c6c4138ee5c86423
\newcommand{\MlstFtCmpJcTwoResOverFullESixPr}{equal within 2\%}  % experiments/FIGS/data/v2/index.json :: fine-tuning jc2 N=1000000: (1 - macro_auc_ovr) of L162 over L188, run by run :: sha256 c6c4138ee5c86423 :: prose twin: two significant figures, direction in words
\newcommand{\MlstFtCmpJcTwoMidOverFullESix}{1.043\,\ensuremath{\pm}\,0.002}  % experiments/FIGS/data/v2/index.json :: fine-tuning jc2 N=1000000: (1 - macro_auc_ovr) of R42_Q1 over L188, run by run :: sha256 c6c4138ee5c86423
\newcommand{\MlstFtCmpJcTwoMidOverFullESixPr}{4.3\% more errors}  % experiments/FIGS/data/v2/index.json :: fine-tuning jc2 N=1000000: (1 - macro_auc_ovr) of R42_Q1 over L188, run by run :: sha256 c6c4138ee5c86423 :: prose twin: two significant figures, direction in words
\newcommand{\MlstFtCmpJcTwoCoarseOverFullESix}{1.123\,\ensuremath{\pm}\,0.003}  % experiments/FIGS/data/v2/index.json :: fine-tuning jc2 N=1000000: (1 - macro_auc_ovr) of R16_Q1 over L188, run by run :: sha256 c6c4138ee5c86423
\newcommand{\MlstFtCmpJcTwoCoarseOverFullESixPr}{12\% more errors}  % experiments/FIGS/data/v2/index.json :: fine-tuning jc2 N=1000000: (1 - macro_auc_ovr) of R16_Q1 over L188, run by run :: sha256 c6c4138ee5c86423 :: prose twin: two significant figures, direction in words
\newcommand{\MlstFtCmpJcTwoResMassOverResESix}{1.009 and 1.016}  % experiments/FIGS/data/v2/index.json :: fine-tuning jc2 N=1000000: (1 - macro_auc_ovr) of L162_MASS over L162, run by run :: sha256 c6c4138ee5c86423
\newcommand{\MlstFtCmpJcTwoResMassOverResESixPr}{equal within 2\%}  % experiments/FIGS/data/v2/index.json :: fine-tuning jc2 N=1000000: (1 - macro_auc_ovr) of L162_MASS over L162, run by run :: sha256 c6c4138ee5c86423 :: prose twin: two significant figures, direction in words
\newcommand{\MlstFtCmpJcTwoResMassOverResESixPrMax}{equal within 2\%}  % experiments/FIGS/data/v2/index.json :: fine-tuning jc2 N=1000000: (1 - macro_auc_ovr) of L162_MASS over L162, run by run :: sha256 c6c4138ee5c86423 :: prose twin of the run farther from 1
\newcommand{\MlstFtCmpJcTwoCoarseMassOverCoarseESix}{1.023 and 1.017}  % experiments/FIGS/data/v2/index.json :: fine-tuning jc2 N=1000000: (1 - macro_auc_ovr) of R16_Q1_MASS over R16_Q1, run by run :: sha256 c6c4138ee5c86423
\newcommand{\MlstFtCmpJcTwoCoarseMassOverCoarseESixPr}{at least 1.7\% more errors}  % experiments/FIGS/data/v2/index.json :: fine-tuning jc2 N=1000000: (1 - macro_auc_ovr) of R16_Q1_MASS over R16_Q1, run by run :: sha256 c6c4138ee5c86423 :: prose twin: two significant figures, direction in words
\newcommand{\MlstFtCmpJcTwoCoarseMassOverCoarseESixPrMax}{up to 2.3\% more errors}  % experiments/FIGS/data/v2/index.json :: fine-tuning jc2 N=1000000: (1 - macro_auc_ovr) of R16_Q1_MASS over R16_Q1, run by run :: sha256 c6c4138ee5c86423 :: prose twin of the run farther from 1
\newcommand{\MlstFtCmpJcTwoQcdResOverFullEFour}{1.038\,\ensuremath{\pm}\,0.005}  % experiments/FIGS/data/v2/index.json :: fine-tuning jc2_qcd N=10000: (1 - macro_auc_qcd) of L162 over L188, run by run :: sha256 c6c4138ee5c86423
\newcommand{\MlstFtCmpJcTwoQcdResOverFullEFourPr}{3.8\% more errors}  % experiments/FIGS/data/v2/index.json :: fine-tuning jc2_qcd N=10000: (1 - macro_auc_qcd) of L162 over L188, run by run :: sha256 c6c4138ee5c86423 :: prose twin: two significant figures, direction in words
\newcommand{\MlstFtCmpJcTwoQcdMidOverFullEFour}{1.164\,\ensuremath{\pm}\,0.018}  % experiments/FIGS/data/v2/index.json :: fine-tuning jc2_qcd N=10000: (1 - macro_auc_qcd) of R42_Q1 over L188, run by run :: sha256 c6c4138ee5c86423
\newcommand{\MlstFtCmpJcTwoQcdMidOverFullEFourPr}{16\% more errors}  % experiments/FIGS/data/v2/index.json :: fine-tuning jc2_qcd N=10000: (1 - macro_auc_qcd) of R42_Q1 over L188, run by run :: sha256 c6c4138ee5c86423 :: prose twin: two significant figures, direction in words
\newcommand{\MlstFtCmpJcTwoQcdCoarseOverFullEFour}{1.258\,\ensuremath{\pm}\,0.025}  % experiments/FIGS/data/v2/index.json :: fine-tuning jc2_qcd N=10000: (1 - macro_auc_qcd) of R16_Q1 over L188, run by run :: sha256 c6c4138ee5c86423
\newcommand{\MlstFtCmpJcTwoQcdCoarseOverFullEFourPr}{26\% more errors}  % experiments/FIGS/data/v2/index.json :: fine-tuning jc2_qcd N=10000: (1 - macro_auc_qcd) of R16_Q1 over L188, run by run :: sha256 c6c4138ee5c86423 :: prose twin: two significant figures, direction in words
\newcommand{\MlstFtCmpJcTwoQcdResMassOverResEFour}{1.027 and 0.997}  % experiments/FIGS/data/v2/index.json :: fine-tuning jc2_qcd N=10000: (1 - macro_auc_qcd) of L162_MASS over L162, run by run :: sha256 c6c4138ee5c86423
\newcommand{\MlstFtCmpJcTwoQcdResMassOverResEFourPr}{equal within 3\%}  % experiments/FIGS/data/v2/index.json :: fine-tuning jc2_qcd N=10000: (1 - macro_auc_qcd) of L162_MASS over L162, run by run :: sha256 c6c4138ee5c86423 :: prose twin: two significant figures, direction in words
\newcommand{\MlstFtCmpJcTwoQcdResMassOverResEFourPrMax}{equal within 3\%}  % experiments/FIGS/data/v2/index.json :: fine-tuning jc2_qcd N=10000: (1 - macro_auc_qcd) of L162_MASS over L162, run by run :: sha256 c6c4138ee5c86423 :: prose twin of the run farther from 1
\newcommand{\MlstFtCmpJcTwoQcdCoarseMassOverCoarseEFour}{1.028 and 1.048}  % experiments/FIGS/data/v2/index.json :: fine-tuning jc2_qcd N=10000: (1 - macro_auc_qcd) of R16_Q1_MASS over R16_Q1, run by run :: sha256 c6c4138ee5c86423
\newcommand{\MlstFtCmpJcTwoQcdCoarseMassOverCoarseEFourPr}{equal within 5\%}  % experiments/FIGS/data/v2/index.json :: fine-tuning jc2_qcd N=10000: (1 - macro_auc_qcd) of R16_Q1_MASS over R16_Q1, run by run :: sha256 c6c4138ee5c86423 :: prose twin: two significant figures, direction in words
\newcommand{\MlstFtCmpJcTwoQcdCoarseMassOverCoarseEFourPrMax}{equal within 5\%}  % experiments/FIGS/data/v2/index.json :: fine-tuning jc2_qcd N=10000: (1 - macro_auc_qcd) of R16_Q1_MASS over R16_Q1, run by run :: sha256 c6c4138ee5c86423 :: prose twin of the run farther from 1
\newcommand{\MlstFtCmpJcTwoQcdResOverFullEFive}{1.060\,\ensuremath{\pm}\,0.002}  % experiments/FIGS/data/v2/index.json :: fine-tuning jc2_qcd N=100000: (1 - macro_auc_qcd) of L162 over L188, run by run :: sha256 c6c4138ee5c86423
\newcommand{\MlstFtCmpJcTwoQcdResOverFullEFivePr}{6.0\% more errors}  % experiments/FIGS/data/v2/index.json :: fine-tuning jc2_qcd N=100000: (1 - macro_auc_qcd) of L162 over L188, run by run :: sha256 c6c4138ee5c86423 :: prose twin: two significant figures, direction in words
\newcommand{\MlstFtCmpJcTwoQcdMidOverFullEFive}{1.062\,\ensuremath{\pm}\,0.010}  % experiments/FIGS/data/v2/index.json :: fine-tuning jc2_qcd N=100000: (1 - macro_auc_qcd) of R42_Q1 over L188, run by run :: sha256 c6c4138ee5c86423
\newcommand{\MlstFtCmpJcTwoQcdMidOverFullEFivePr}{6.2\% more errors}  % experiments/FIGS/data/v2/index.json :: fine-tuning jc2_qcd N=100000: (1 - macro_auc_qcd) of R42_Q1 over L188, run by run :: sha256 c6c4138ee5c86423 :: prose twin: two significant figures, direction in words
\newcommand{\MlstFtCmpJcTwoQcdCoarseOverFullEFive}{1.137\,\ensuremath{\pm}\,0.001}  % experiments/FIGS/data/v2/index.json :: fine-tuning jc2_qcd N=100000: (1 - macro_auc_qcd) of R16_Q1 over L188, run by run :: sha256 c6c4138ee5c86423
\newcommand{\MlstFtCmpJcTwoQcdCoarseOverFullEFivePr}{14\% more errors}  % experiments/FIGS/data/v2/index.json :: fine-tuning jc2_qcd N=100000: (1 - macro_auc_qcd) of R16_Q1 over L188, run by run :: sha256 c6c4138ee5c86423 :: prose twin: two significant figures, direction in words
\newcommand{\MlstFtCmpJcTwoQcdResMassOverResEFive}{1.022 and 1.009}  % experiments/FIGS/data/v2/index.json :: fine-tuning jc2_qcd N=100000: (1 - macro_auc_qcd) of L162_MASS over L162, run by run :: sha256 c6c4138ee5c86423
\newcommand{\MlstFtCmpJcTwoQcdResMassOverResEFivePr}{equal within 3\%}  % experiments/FIGS/data/v2/index.json :: fine-tuning jc2_qcd N=100000: (1 - macro_auc_qcd) of L162_MASS over L162, run by run :: sha256 c6c4138ee5c86423 :: prose twin: two significant figures, direction in words
\newcommand{\MlstFtCmpJcTwoQcdResMassOverResEFivePrMax}{equal within 3\%}  % experiments/FIGS/data/v2/index.json :: fine-tuning jc2_qcd N=100000: (1 - macro_auc_qcd) of L162_MASS over L162, run by run :: sha256 c6c4138ee5c86423 :: prose twin of the run farther from 1
\newcommand{\MlstFtCmpJcTwoQcdCoarseMassOverCoarseEFive}{1.028 and 1.027}  % experiments/FIGS/data/v2/index.json :: fine-tuning jc2_qcd N=100000: (1 - macro_auc_qcd) of R16_Q1_MASS over R16_Q1, run by run :: sha256 c6c4138ee5c86423
\newcommand{\MlstFtCmpJcTwoQcdCoarseMassOverCoarseEFivePr}{at least 2.7\% more errors}  % experiments/FIGS/data/v2/index.json :: fine-tuning jc2_qcd N=100000: (1 - macro_auc_qcd) of R16_Q1_MASS over R16_Q1, run by run :: sha256 c6c4138ee5c86423 :: prose twin: two significant figures, direction in words
\newcommand{\MlstFtCmpJcTwoQcdCoarseMassOverCoarseEFivePrMax}{up to 2.8\% more errors}  % experiments/FIGS/data/v2/index.json :: fine-tuning jc2_qcd N=100000: (1 - macro_auc_qcd) of R16_Q1_MASS over R16_Q1, run by run :: sha256 c6c4138ee5c86423 :: prose twin of the run farther from 1
\newcommand{\MlstFtCmpJcTwoQcdResOverFullESix}{1.035\,\ensuremath{\pm}\,0.010}  % experiments/FIGS/data/v2/index.json :: fine-tuning jc2_qcd N=1000000: (1 - macro_auc_qcd) of L162 over L188, run by run :: sha256 c6c4138ee5c86423
\newcommand{\MlstFtCmpJcTwoQcdResOverFullESixPr}{3.5\% more errors}  % experiments/FIGS/data/v2/index.json :: fine-tuning jc2_qcd N=1000000: (1 - macro_auc_qcd) of L162 over L188, run by run :: sha256 c6c4138ee5c86423 :: prose twin: two significant figures, direction in words
\newcommand{\MlstFtCmpJcTwoQcdMidOverFullESix}{1.050\,\ensuremath{\pm}\,0.004}  % experiments/FIGS/data/v2/index.json :: fine-tuning jc2_qcd N=1000000: (1 - macro_auc_qcd) of R42_Q1 over L188, run by run :: sha256 c6c4138ee5c86423
\newcommand{\MlstFtCmpJcTwoQcdMidOverFullESixPr}{5.0\% more errors}  % experiments/FIGS/data/v2/index.json :: fine-tuning jc2_qcd N=1000000: (1 - macro_auc_qcd) of R42_Q1 over L188, run by run :: sha256 c6c4138ee5c86423 :: prose twin: two significant figures, direction in words
\newcommand{\MlstFtCmpJcTwoQcdCoarseOverFullESix}{1.122\,\ensuremath{\pm}\,0.002}  % experiments/FIGS/data/v2/index.json :: fine-tuning jc2_qcd N=1000000: (1 - macro_auc_qcd) of R16_Q1 over L188, run by run :: sha256 c6c4138ee5c86423
\newcommand{\MlstFtCmpJcTwoQcdCoarseOverFullESixPr}{12\% more errors}  % experiments/FIGS/data/v2/index.json :: fine-tuning jc2_qcd N=1000000: (1 - macro_auc_qcd) of R16_Q1 over L188, run by run :: sha256 c6c4138ee5c86423 :: prose twin: two significant figures, direction in words
\newcommand{\MlstFtCmpJcTwoQcdResMassOverResESix}{1.011 and 1.025}  % experiments/FIGS/data/v2/index.json :: fine-tuning jc2_qcd N=1000000: (1 - macro_auc_qcd) of L162_MASS over L162, run by run :: sha256 c6c4138ee5c86423
\newcommand{\MlstFtCmpJcTwoQcdResMassOverResESixPr}{equal within 3\%}  % experiments/FIGS/data/v2/index.json :: fine-tuning jc2_qcd N=1000000: (1 - macro_auc_qcd) of L162_MASS over L162, run by run :: sha256 c6c4138ee5c86423 :: prose twin: two significant figures, direction in words
\newcommand{\MlstFtCmpJcTwoQcdResMassOverResESixPrMax}{equal within 3\%}  % experiments/FIGS/data/v2/index.json :: fine-tuning jc2_qcd N=1000000: (1 - macro_auc_qcd) of L162_MASS over L162, run by run :: sha256 c6c4138ee5c86423 :: prose twin of the run farther from 1
\newcommand{\MlstFtCmpJcTwoQcdCoarseMassOverCoarseESix}{1.015 and 1.015}  % experiments/FIGS/data/v2/index.json :: fine-tuning jc2_qcd N=1000000: (1 - macro_auc_qcd) of R16_Q1_MASS over R16_Q1, run by run :: sha256 c6c4138ee5c86423
\newcommand{\MlstFtCmpJcTwoQcdCoarseMassOverCoarseESixPr}{at least 1.5\% more errors}  % experiments/FIGS/data/v2/index.json :: fine-tuning jc2_qcd N=1000000: (1 - macro_auc_qcd) of R16_Q1_MASS over R16_Q1, run by run :: sha256 c6c4138ee5c86423 :: prose twin: two significant figures, direction in words
\newcommand{\MlstFtCmpJcTwoQcdCoarseMassOverCoarseESixPrMax}{up to 1.5\% more errors}  % experiments/FIGS/data/v2/index.json :: fine-tuning jc2_qcd N=1000000: (1 - macro_auc_qcd) of R16_Q1_MASS over R16_Q1, run by run :: sha256 c6c4138ee5c86423 :: prose twin of the run farther from 1
\newcommand{\MlstFtJcTwoAccFullEFour}{0.3914\,\ensuremath{\pm}\,0.0011}  % experiments/FIGS/data/v2/index.json :: fine-tuning jc2 N=10000: accuracy, L188 :: sha256 c6c4138ee5c86423
\newcommand{\MlstFtJcTwoAccResEFour}{0.3934\,\ensuremath{\pm}\,0.0018}  % experiments/FIGS/data/v2/index.json :: fine-tuning jc2 N=10000: accuracy, L162 :: sha256 c6c4138ee5c86423
\newcommand{\MlstFtJcTwoAccMidEFour}{0.3563\,\ensuremath{\pm}\,0.0025}  % experiments/FIGS/data/v2/index.json :: fine-tuning jc2 N=10000: accuracy, R42_Q1 :: sha256 c6c4138ee5c86423
\newcommand{\MlstFtJcTwoAccCoarseEFour}{0.3219\,\ensuremath{\pm}\,0.0011}  % experiments/FIGS/data/v2/index.json :: fine-tuning jc2 N=10000: accuracy, R16_Q1 :: sha256 c6c4138ee5c86423
\newcommand{\MlstFtJcTwoAccResMassEFour}{0.381 and 0.392}  % experiments/FIGS/data/v2/index.json :: fine-tuning jc2 N=10000: accuracy, L162_MASS :: sha256 c6c4138ee5c86423
\newcommand{\MlstFtJcTwoAccCoarseMassEFour}{0.3156 and 0.3187}  % experiments/FIGS/data/v2/index.json :: fine-tuning jc2 N=10000: accuracy, R16_Q1_MASS :: sha256 c6c4138ee5c86423
\newcommand{\MlstFtJcTwoAccScratchEFour}{0.1492\,\ensuremath{\pm}\,0.0011}  % experiments/FIGS/data/v2/index.json :: fine-tuning jc2 N=10000: accuracy from scratch :: sha256 c6c4138ee5c86423
\newcommand{\MlstFtJcTwoAccFullMinusScratchEFour}{0.2423\,\ensuremath{\pm}\,0.0011}  % experiments/FIGS/data/v2/index.json :: fine-tuning jc2 N=10000: accuracy of each 188-class run minus the from-scratch mean (absolute) :: sha256 c6c4138ee5c86423
\newcommand{\MlstFtJcTwoAccFullEFive}{0.4164\,\ensuremath{\pm}\,0.0015}  % experiments/FIGS/data/v2/index.json :: fine-tuning jc2 N=100000: accuracy, L188 :: sha256 c6c4138ee5c86423
\newcommand{\MlstFtJcTwoAccResEFive}{0.4161\,\ensuremath{\pm}\,0.0015}  % experiments/FIGS/data/v2/index.json :: fine-tuning jc2 N=100000: accuracy, L162 :: sha256 c6c4138ee5c86423
\newcommand{\MlstFtJcTwoAccMidEFive}{0.4007\,\ensuremath{\pm}\,0.0013}  % experiments/FIGS/data/v2/index.json :: fine-tuning jc2 N=100000: accuracy, R42_Q1 :: sha256 c6c4138ee5c86423
\newcommand{\MlstFtJcTwoAccCoarseEFive}{0.3833\,\ensuremath{\pm}\,0.0020}  % experiments/FIGS/data/v2/index.json :: fine-tuning jc2 N=100000: accuracy, R16_Q1 :: sha256 c6c4138ee5c86423
\newcommand{\MlstFtJcTwoAccResMassEFive}{0.4093 and 0.4115}  % experiments/FIGS/data/v2/index.json :: fine-tuning jc2 N=100000: accuracy, L162_MASS :: sha256 c6c4138ee5c86423
\newcommand{\MlstFtJcTwoAccCoarseMassEFive}{0.3758 and 0.3781}  % experiments/FIGS/data/v2/index.json :: fine-tuning jc2 N=100000: accuracy, R16_Q1_MASS :: sha256 c6c4138ee5c86423
\newcommand{\MlstFtJcTwoAccScratchEFive}{0.292\,\ensuremath{\pm}\,0.004}  % experiments/FIGS/data/v2/index.json :: fine-tuning jc2 N=100000: accuracy from scratch :: sha256 c6c4138ee5c86423
\newcommand{\MlstFtJcTwoAccFullMinusScratchEFive}{0.1243\,\ensuremath{\pm}\,0.0015}  % experiments/FIGS/data/v2/index.json :: fine-tuning jc2 N=100000: accuracy of each 188-class run minus the from-scratch mean (absolute) :: sha256 c6c4138ee5c86423
\newcommand{\MlstFtJcTwoAccFullESix}{0.4263\,\ensuremath{\pm}\,0.0005}  % experiments/FIGS/data/v2/index.json :: fine-tuning jc2 N=1000000: accuracy, L188 :: sha256 c6c4138ee5c86423
\newcommand{\MlstFtJcTwoAccResESix}{0.4249\,\ensuremath{\pm}\,0.0021}  % experiments/FIGS/data/v2/index.json :: fine-tuning jc2 N=1000000: accuracy, L162 :: sha256 c6c4138ee5c86423
\newcommand{\MlstFtJcTwoAccMidESix}{0.4166\,\ensuremath{\pm}\,0.0011}  % experiments/FIGS/data/v2/index.json :: fine-tuning jc2 N=1000000: accuracy, R42_Q1 :: sha256 c6c4138ee5c86423
\newcommand{\MlstFtJcTwoAccCoarseESix}{0.4054\,\ensuremath{\pm}\,0.0020}  % experiments/FIGS/data/v2/index.json :: fine-tuning jc2 N=1000000: accuracy, R16_Q1 :: sha256 c6c4138ee5c86423
\newcommand{\MlstFtJcTwoAccResMassESix}{0.4198 and 0.4213}  % experiments/FIGS/data/v2/index.json :: fine-tuning jc2 N=1000000: accuracy, L162_MASS :: sha256 c6c4138ee5c86423
\newcommand{\MlstFtJcTwoAccCoarseMassESix}{0.4012 and 0.4000}  % experiments/FIGS/data/v2/index.json :: fine-tuning jc2 N=1000000: accuracy, R16_Q1_MASS :: sha256 c6c4138ee5c86423
\newcommand{\MlstFtJcTwoAccScratchESix}{0.3677\,\ensuremath{\pm}\,0.0012}  % experiments/FIGS/data/v2/index.json :: fine-tuning jc2 N=1000000: accuracy from scratch :: sha256 c6c4138ee5c86423
\newcommand{\MlstFtJcTwoAccFullMinusScratchESix}{0.0585\,\ensuremath{\pm}\,0.0005}  % experiments/FIGS/data/v2/index.json :: fine-tuning jc2 N=1000000: accuracy of each 188-class run minus the from-scratch mean (absolute) :: sha256 c6c4138ee5c86423
\newcommand{\MlstFtJcTwoAucFull}{0.9797\,\ensuremath{\pm}\,0.0001}  % experiments/FIGS/data/v2/index.json :: fine-tuning class-balanced JetClass-II N=1000000: AUC averaged over the classes, Full :: sha256 c6c4138ee5c86423
\newcommand{\MlstFtJcTwoAucScratch}{0.9734\,\ensuremath{\pm}\,0.0003}  % experiments/FIGS/data/v2/index.json :: fine-tuning class-balanced JetClass-II N=1000000: AUC averaged over the classes, Scratch :: sha256 c6c4138ee5c86423
\newcommand{\MlstFtJcOneAccFullEFour}{0.8117\,\ensuremath{\pm}\,0.0031}  % experiments/FIGS/data/v2/index.json :: fine-tuning jc1 N=10000: accuracy, L188 :: sha256 c6c4138ee5c86423
\newcommand{\MlstFtJcOneAccResEFour}{0.8117\,\ensuremath{\pm}\,0.0013}  % experiments/FIGS/data/v2/index.json :: fine-tuning jc1 N=10000: accuracy, L162 :: sha256 c6c4138ee5c86423
\newcommand{\MlstFtJcOneAccMidEFour}{0.8086\,\ensuremath{\pm}\,0.0031}  % experiments/FIGS/data/v2/index.json :: fine-tuning jc1 N=10000: accuracy, R42_Q1 :: sha256 c6c4138ee5c86423
\newcommand{\MlstFtJcOneAccCoarseEFour}{0.7982\,\ensuremath{\pm}\,0.0023}  % experiments/FIGS/data/v2/index.json :: fine-tuning jc1 N=10000: accuracy, R16_Q1 :: sha256 c6c4138ee5c86423
\newcommand{\MlstFtJcOneAccResMassEFour}{0.8107 and 0.8098}  % experiments/FIGS/data/v2/index.json :: fine-tuning jc1 N=10000: accuracy, L162_MASS :: sha256 c6c4138ee5c86423
\newcommand{\MlstFtJcOneAccCoarseMassEFour}{0.7944 and 0.7990}  % experiments/FIGS/data/v2/index.json :: fine-tuning jc1 N=10000: accuracy, R16_Q1_MASS :: sha256 c6c4138ee5c86423
\newcommand{\MlstFtJcOneAccScratchEFour}{0.677\,\ensuremath{\pm}\,0.004}  % experiments/FIGS/data/v2/index.json :: fine-tuning jc1 N=10000: accuracy from scratch :: sha256 c6c4138ee5c86423
\newcommand{\MlstFtJcOneAccFullMinusScratchEFour}{0.1343\,\ensuremath{\pm}\,0.0031}  % experiments/FIGS/data/v2/index.json :: fine-tuning jc1 N=10000: accuracy of each 188-class run minus the from-scratch mean (absolute) :: sha256 c6c4138ee5c86423
\newcommand{\MlstFtJcOneAccFullEFive}{0.8343\,\ensuremath{\pm}\,0.0035}  % experiments/FIGS/data/v2/index.json :: fine-tuning jc1 N=100000: accuracy, L188 :: sha256 c6c4138ee5c86423
\newcommand{\MlstFtJcOneAccResEFive}{0.8328\,\ensuremath{\pm}\,0.0005}  % experiments/FIGS/data/v2/index.json :: fine-tuning jc1 N=100000: accuracy, L162 :: sha256 c6c4138ee5c86423
\newcommand{\MlstFtJcOneAccMidEFive}{0.8324\,\ensuremath{\pm}\,0.0004}  % experiments/FIGS/data/v2/index.json :: fine-tuning jc1 N=100000: accuracy, R42_Q1 :: sha256 c6c4138ee5c86423
\newcommand{\MlstFtJcOneAccCoarseEFive}{0.8288\,\ensuremath{\pm}\,0.0005}  % experiments/FIGS/data/v2/index.json :: fine-tuning jc1 N=100000: accuracy, R16_Q1 :: sha256 c6c4138ee5c86423
\newcommand{\MlstFtJcOneAccResMassEFive}{0.8326 and 0.8348}  % experiments/FIGS/data/v2/index.json :: fine-tuning jc1 N=100000: accuracy, L162_MASS :: sha256 c6c4138ee5c86423
\newcommand{\MlstFtJcOneAccCoarseMassEFive}{0.8253 and 0.8262}  % experiments/FIGS/data/v2/index.json :: fine-tuning jc1 N=100000: accuracy, R16_Q1_MASS :: sha256 c6c4138ee5c86423
\newcommand{\MlstFtJcOneAccScratchEFive}{0.777\,\ensuremath{\pm}\,0.005}  % experiments/FIGS/data/v2/index.json :: fine-tuning jc1 N=100000: accuracy from scratch :: sha256 c6c4138ee5c86423
\newcommand{\MlstFtJcOneAccFullMinusScratchEFive}{0.0571\,\ensuremath{\pm}\,0.0035}  % experiments/FIGS/data/v2/index.json :: fine-tuning jc1 N=100000: accuracy of each 188-class run minus the from-scratch mean (absolute) :: sha256 c6c4138ee5c86423
\newcommand{\MlstFtJcOneAccFullESix}{0.8468\,\ensuremath{\pm}\,0.0006}  % experiments/FIGS/data/v2/index.json :: fine-tuning jc1 N=1000000: accuracy, L188 :: sha256 c6c4138ee5c86423
\newcommand{\MlstFtJcOneAccResESix}{0.8475\,\ensuremath{\pm}\,0.0009}  % experiments/FIGS/data/v2/index.json :: fine-tuning jc1 N=1000000: accuracy, L162 :: sha256 c6c4138ee5c86423
\newcommand{\MlstFtJcOneAccMidESix}{0.8459\,\ensuremath{\pm}\,0.0002}  % experiments/FIGS/data/v2/index.json :: fine-tuning jc1 N=1000000: accuracy, R42_Q1 :: sha256 c6c4138ee5c86423
\newcommand{\MlstFtJcOneAccCoarseESix}{0.8422\,\ensuremath{\pm}\,0.0005}  % experiments/FIGS/data/v2/index.json :: fine-tuning jc1 N=1000000: accuracy, R16_Q1 :: sha256 c6c4138ee5c86423
\newcommand{\MlstFtJcOneAccResMassESix}{0.8459 and 0.8471}  % experiments/FIGS/data/v2/index.json :: fine-tuning jc1 N=1000000: accuracy, L162_MASS :: sha256 c6c4138ee5c86423
\newcommand{\MlstFtJcOneAccCoarseMassESix}{0.8400 and 0.8398}  % experiments/FIGS/data/v2/index.json :: fine-tuning jc1 N=1000000: accuracy, R16_Q1_MASS :: sha256 c6c4138ee5c86423
\newcommand{\MlstFtJcOneAccScratchESix}{0.8267\,\ensuremath{\pm}\,0.0011}  % experiments/FIGS/data/v2/index.json :: fine-tuning jc1 N=1000000: accuracy from scratch :: sha256 c6c4138ee5c86423
\newcommand{\MlstFtJcOneAccFullMinusScratchESix}{0.0201\,\ensuremath{\pm}\,0.0006}  % experiments/FIGS/data/v2/index.json :: fine-tuning jc1 N=1000000: accuracy of each 188-class run minus the from-scratch mean (absolute) :: sha256 c6c4138ee5c86423
\newcommand{\MlstFtTopRthirtyScratchEThree}{49\,\ensuremath{\pm}\,12}  % experiments/FIGS/data/v2/index.json :: fine-tuning top N=1000: R30 from scratch, every seed :: sha256 c6c4138ee5c86423
\newcommand{\MlstFtTopRthirtyFullEThree}{72\,\ensuremath{\pm}\,19}  % experiments/FIGS/data/v2/index.json :: fine-tuning top N=1000: R30, L188 :: sha256 c6c4138ee5c86423
\newcommand{\MlstFtTopRthirtyFullOverScratchEThree}{1.5\,\ensuremath{\pm}\,0.4}  % experiments/FIGS/data/v2/index.json :: fine-tuning top N=1000: R30 of each L188 run over the from-scratch mean, SD over runs and standard error of that mean in quadrature :: sha256 c6c4138ee5c86423
\newcommand{\MlstFtTopRthirtyFullOverScratchEThreePr}{equal within 50\%}  % experiments/FIGS/data/v2/index.json :: fine-tuning top N=1000: R30 of each L188 run over the from-scratch mean, SD over runs and standard error of that mean in quadrature :: sha256 c6c4138ee5c86423 :: prose twin: two significant figures, direction in words
\newcommand{\MlstFtTopRthirtyResEThree}{135\,\ensuremath{\pm}\,26}  % experiments/FIGS/data/v2/index.json :: fine-tuning top N=1000: R30, L162 :: sha256 c6c4138ee5c86423
\newcommand{\MlstFtTopRthirtyResOverScratchEThree}{2.8\,\ensuremath{\pm}\,0.7}  % experiments/FIGS/data/v2/index.json :: fine-tuning top N=1000: R30 of each L162 run over the from-scratch mean, SD over runs and standard error of that mean in quadrature :: sha256 c6c4138ee5c86423
\newcommand{\MlstFtTopRthirtyResOverScratchEThreePr}{2.8 times the rejection}  % experiments/FIGS/data/v2/index.json :: fine-tuning top N=1000: R30 of each L162 run over the from-scratch mean, SD over runs and standard error of that mean in quadrature :: sha256 c6c4138ee5c86423 :: prose twin: two significant figures, direction in words
\newcommand{\MlstFtTopRthirtyMidEThree}{104\,\ensuremath{\pm}\,20}  % experiments/FIGS/data/v2/index.json :: fine-tuning top N=1000: R30, R42_Q1 :: sha256 c6c4138ee5c86423
\newcommand{\MlstFtTopRthirtyMidOverScratchEThree}{2.1\,\ensuremath{\pm}\,0.5}  % experiments/FIGS/data/v2/index.json :: fine-tuning top N=1000: R30 of each R42_Q1 run over the from-scratch mean, SD over runs and standard error of that mean in quadrature :: sha256 c6c4138ee5c86423
\newcommand{\MlstFtTopRthirtyMidOverScratchEThreePr}{2.1 times the rejection}  % experiments/FIGS/data/v2/index.json :: fine-tuning top N=1000: R30 of each R42_Q1 run over the from-scratch mean, SD over runs and standard error of that mean in quadrature :: sha256 c6c4138ee5c86423 :: prose twin: two significant figures, direction in words
\newcommand{\MlstFtTopRthirtyCoarseEThree}{83\,\ensuremath{\pm}\,25}  % experiments/FIGS/data/v2/index.json :: fine-tuning top N=1000: R30, R16_Q1 :: sha256 c6c4138ee5c86423
\newcommand{\MlstFtTopRthirtyCoarseOverScratchEThree}{1.7\,\ensuremath{\pm}\,0.6}  % experiments/FIGS/data/v2/index.json :: fine-tuning top N=1000: R30 of each R16_Q1 run over the from-scratch mean, SD over runs and standard error of that mean in quadrature :: sha256 c6c4138ee5c86423
\newcommand{\MlstFtTopRthirtyCoarseOverScratchEThreePr}{equal within 80\%}  % experiments/FIGS/data/v2/index.json :: fine-tuning top N=1000: R30 of each R16_Q1 run over the from-scratch mean, SD over runs and standard error of that mean in quadrature :: sha256 c6c4138ee5c86423 :: prose twin: two significant figures, direction in words
\newcommand{\MlstFtTopRthirtyResMassEThree}{105.9 and 112.7}  % experiments/FIGS/data/v2/index.json :: fine-tuning top N=1000: R30, L162_MASS :: sha256 c6c4138ee5c86423
\newcommand{\MlstFtTopRthirtyResMassOverScratchEThree}{2.18 and 2.32}  % experiments/FIGS/data/v2/index.json :: fine-tuning top N=1000: R30 of each L162_MASS run over the from-scratch mean, SD over runs and standard error of that mean in quadrature :: sha256 c6c4138ee5c86423
\newcommand{\MlstFtTopRthirtyResMassOverScratchEThreePr}{at least 2.2 times the rejection}  % experiments/FIGS/data/v2/index.json :: fine-tuning top N=1000: R30 of each L162_MASS run over the from-scratch mean, SD over runs and standard error of that mean in quadrature :: sha256 c6c4138ee5c86423 :: prose twin: two significant figures, direction in words
\newcommand{\MlstFtTopRthirtyResMassOverScratchEThreePrMax}{up to 2.3 times the rejection}  % experiments/FIGS/data/v2/index.json :: fine-tuning top N=1000: R30 of each L162_MASS run over the from-scratch mean, SD over runs and standard error of that mean in quadrature :: sha256 c6c4138ee5c86423 :: prose twin of the run farther from 1
\newcommand{\MlstFtTopRthirtyCoarseMassEThree}{59.1 and 63.9}  % experiments/FIGS/data/v2/index.json :: fine-tuning top N=1000: R30, R16_Q1_MASS :: sha256 c6c4138ee5c86423
\newcommand{\MlstFtTopRthirtyCoarseMassOverScratchEThree}{1.216 and 1.314}  % experiments/FIGS/data/v2/index.json :: fine-tuning top N=1000: R30 of each R16_Q1_MASS run over the from-scratch mean, SD over runs and standard error of that mean in quadrature :: sha256 c6c4138ee5c86423
\newcommand{\MlstFtTopRthirtyCoarseMassOverScratchEThreePr}{at least 22\% more rejection}  % experiments/FIGS/data/v2/index.json :: fine-tuning top N=1000: R30 of each R16_Q1_MASS run over the from-scratch mean, SD over runs and standard error of that mean in quadrature :: sha256 c6c4138ee5c86423 :: prose twin: two significant figures, direction in words
\newcommand{\MlstFtTopRthirtyCoarseMassOverScratchEThreePrMax}{up to 31\% more rejection}  % experiments/FIGS/data/v2/index.json :: fine-tuning top N=1000: R30 of each R16_Q1_MASS run over the from-scratch mean, SD over runs and standard error of that mean in quadrature :: sha256 c6c4138ee5c86423 :: prose twin of the run farther from 1
\newcommand{\MlstFtTopRthirtyResOverFullEThree}{2.0\,\ensuremath{\pm}\,0.8}  % experiments/FIGS/data/v2/index.json :: fine-tuning top N=1000: R30 of L162 over L188, run by run :: sha256 c6c4138ee5c86423
\newcommand{\MlstFtTopRthirtyResOverFullEThreePr}{equal within 100\%}  % experiments/FIGS/data/v2/index.json :: fine-tuning top N=1000: R30 of L162 over L188, run by run :: sha256 c6c4138ee5c86423 :: prose twin: two significant figures, direction in words
\newcommand{\MlstFtTopRthirtyMidOverFullEThree}{1.46\,\ensuremath{\pm}\,0.12}  % experiments/FIGS/data/v2/index.json :: fine-tuning top N=1000: R30 of R42_Q1 over L188, run by run :: sha256 c6c4138ee5c86423
\newcommand{\MlstFtTopRthirtyMidOverFullEThreePr}{46\% more rejection}  % experiments/FIGS/data/v2/index.json :: fine-tuning top N=1000: R30 of R42_Q1 over L188, run by run :: sha256 c6c4138ee5c86423 :: prose twin: two significant figures, direction in words
\newcommand{\MlstFtTopRthirtyCoarseOverFullEThree}{1.3\,\ensuremath{\pm}\,0.7}  % experiments/FIGS/data/v2/index.json :: fine-tuning top N=1000: R30 of R16_Q1 over L188, run by run :: sha256 c6c4138ee5c86423
\newcommand{\MlstFtTopRthirtyCoarseOverFullEThreePr}{equal within 80\%}  % experiments/FIGS/data/v2/index.json :: fine-tuning top N=1000: R30 of R16_Q1 over L188, run by run :: sha256 c6c4138ee5c86423 :: prose twin: two significant figures, direction in words
\newcommand{\MlstFtTopRthirtyResMassOverResEThree}{0.73 and 1.07}  % experiments/FIGS/data/v2/index.json :: fine-tuning top N=1000: R30 of L162_MASS over L162, run by run :: sha256 c6c4138ee5c86423
\newcommand{\MlstFtTopRthirtyResMassOverResEThreePr}{equal within 30\%}  % experiments/FIGS/data/v2/index.json :: fine-tuning top N=1000: R30 of L162_MASS over L162, run by run :: sha256 c6c4138ee5c86423 :: prose twin: two significant figures, direction in words
\newcommand{\MlstFtTopRthirtyResMassOverResEThreePrMax}{equal within 30\%}  % experiments/FIGS/data/v2/index.json :: fine-tuning top N=1000: R30 of L162_MASS over L162, run by run :: sha256 c6c4138ee5c86423 :: prose twin of the run farther from 1
\newcommand{\MlstFtTopRthirtyCoarseMassOverCoarseEThree}{0.56 and 0.75}  % experiments/FIGS/data/v2/index.json :: fine-tuning top N=1000: R30 of R16_Q1_MASS over R16_Q1, run by run :: sha256 c6c4138ee5c86423
\newcommand{\MlstFtTopRthirtyCoarseMassOverCoarseEThreePr}{equal within 50\%}  % experiments/FIGS/data/v2/index.json :: fine-tuning top N=1000: R30 of R16_Q1_MASS over R16_Q1, run by run :: sha256 c6c4138ee5c86423 :: prose twin: two significant figures, direction in words
\newcommand{\MlstFtTopRthirtyCoarseMassOverCoarseEThreePrMax}{equal within 50\%}  % experiments/FIGS/data/v2/index.json :: fine-tuning top N=1000: R30 of R16_Q1_MASS over R16_Q1, run by run :: sha256 c6c4138ee5c86423 :: prose twin of the run farther from 1
\newcommand{\MlstFtTopRthirtyRandOverFullEThree}{1.3\,\ensuremath{\pm}\,0.5}  % experiments/FIGS/data/v2/index.json :: fine-tuning top N=1000: R30 of each random partition over the L188 run of its seed :: sha256 c6c4138ee5c86423
\newcommand{\MlstFtTopRthirtyRandOverFullEThreePr}{equal within 50\%}  % experiments/FIGS/data/v2/index.json :: fine-tuning top N=1000: R30 of each random partition over the L188 run of its seed :: sha256 c6c4138ee5c86423 :: prose twin: two significant figures, direction in words
\newcommand{\MlstFtTopRthirtyRandOverCoarseEThree}{0.87\,\ensuremath{\pm}\,0.12}  % experiments/FIGS/data/v2/index.json :: fine-tuning top N=1000: R30 of each random partition over the R16_Q1 run of its seed :: sha256 c6c4138ee5c86423
\newcommand{\MlstFtTopRthirtyRandOverCoarseEThreePr}{equal within 20\%}  % experiments/FIGS/data/v2/index.json :: fine-tuning top N=1000: R30 of each random partition over the R16_Q1 run of its seed :: sha256 c6c4138ee5c86423 :: prose twin: two significant figures, direction in words
\newcommand{\MlstFtTopRthirtyScratchEFour}{340\,\ensuremath{\pm}\,50}  % experiments/FIGS/data/v2/index.json :: fine-tuning top N=10000: R30 from scratch, every seed :: sha256 c6c4138ee5c86423
\newcommand{\MlstFtTopRthirtyFullEFour}{420\,\ensuremath{\pm}\,9}  % experiments/FIGS/data/v2/index.json :: fine-tuning top N=10000: R30, L188 :: sha256 c6c4138ee5c86423
\newcommand{\MlstFtTopRthirtyFullOverScratchEFour}{1.24\,\ensuremath{\pm}\,0.11}  % experiments/FIGS/data/v2/index.json :: fine-tuning top N=10000: R30 of each L188 run over the from-scratch mean, SD over runs and standard error of that mean in quadrature :: sha256 c6c4138ee5c86423
\newcommand{\MlstFtTopRthirtyFullOverScratchEFourPr}{24\% more rejection}  % experiments/FIGS/data/v2/index.json :: fine-tuning top N=10000: R30 of each L188 run over the from-scratch mean, SD over runs and standard error of that mean in quadrature :: sha256 c6c4138ee5c86423 :: prose twin: two significant figures, direction in words
\newcommand{\MlstFtTopRthirtyResEFour}{430\,\ensuremath{\pm}\,60}  % experiments/FIGS/data/v2/index.json :: fine-tuning top N=10000: R30, L162 :: sha256 c6c4138ee5c86423
\newcommand{\MlstFtTopRthirtyResOverScratchEFour}{1.27\,\ensuremath{\pm}\,0.21}  % experiments/FIGS/data/v2/index.json :: fine-tuning top N=10000: R30 of each L162 run over the from-scratch mean, SD over runs and standard error of that mean in quadrature :: sha256 c6c4138ee5c86423
\newcommand{\MlstFtTopRthirtyResOverScratchEFourPr}{equal within 30\%}  % experiments/FIGS/data/v2/index.json :: fine-tuning top N=10000: R30 of each L162 run over the from-scratch mean, SD over runs and standard error of that mean in quadrature :: sha256 c6c4138ee5c86423 :: prose twin: two significant figures, direction in words
\newcommand{\MlstFtTopRthirtyMidEFour}{455\,\ensuremath{\pm}\,30}  % experiments/FIGS/data/v2/index.json :: fine-tuning top N=10000: R30, R42_Q1 :: sha256 c6c4138ee5c86423
\newcommand{\MlstFtTopRthirtyMidOverScratchEFour}{1.34\,\ensuremath{\pm}\,0.14}  % experiments/FIGS/data/v2/index.json :: fine-tuning top N=10000: R30 of each R42_Q1 run over the from-scratch mean, SD over runs and standard error of that mean in quadrature :: sha256 c6c4138ee5c86423
\newcommand{\MlstFtTopRthirtyMidOverScratchEFourPr}{34\% more rejection}  % experiments/FIGS/data/v2/index.json :: fine-tuning top N=10000: R30 of each R42_Q1 run over the from-scratch mean, SD over runs and standard error of that mean in quadrature :: sha256 c6c4138ee5c86423 :: prose twin: two significant figures, direction in words
\newcommand{\MlstFtTopRthirtyCoarseEFour}{350\,\ensuremath{\pm}\,40}  % experiments/FIGS/data/v2/index.json :: fine-tuning top N=10000: R30, R16_Q1 :: sha256 c6c4138ee5c86423
\newcommand{\MlstFtTopRthirtyCoarseOverScratchEFour}{1.04\,\ensuremath{\pm}\,0.14}  % experiments/FIGS/data/v2/index.json :: fine-tuning top N=10000: R30 of each R16_Q1 run over the from-scratch mean, SD over runs and standard error of that mean in quadrature :: sha256 c6c4138ee5c86423
\newcommand{\MlstFtTopRthirtyCoarseOverScratchEFourPr}{equal within 20\%}  % experiments/FIGS/data/v2/index.json :: fine-tuning top N=10000: R30 of each R16_Q1 run over the from-scratch mean, SD over runs and standard error of that mean in quadrature :: sha256 c6c4138ee5c86423 :: prose twin: two significant figures, direction in words
\newcommand{\MlstFtTopRthirtyResMassEFour}{373.2 and 376.7}  % experiments/FIGS/data/v2/index.json :: fine-tuning top N=10000: R30, L162_MASS :: sha256 c6c4138ee5c86423
\newcommand{\MlstFtTopRthirtyResMassOverScratchEFour}{1.101 and 1.111}  % experiments/FIGS/data/v2/index.json :: fine-tuning top N=10000: R30 of each L162_MASS run over the from-scratch mean, SD over runs and standard error of that mean in quadrature :: sha256 c6c4138ee5c86423
\newcommand{\MlstFtTopRthirtyResMassOverScratchEFourPr}{at least 10\% more rejection}  % experiments/FIGS/data/v2/index.json :: fine-tuning top N=10000: R30 of each L162_MASS run over the from-scratch mean, SD over runs and standard error of that mean in quadrature :: sha256 c6c4138ee5c86423 :: prose twin: two significant figures, direction in words
\newcommand{\MlstFtTopRthirtyResMassOverScratchEFourPrMax}{up to 11\% more rejection}  % experiments/FIGS/data/v2/index.json :: fine-tuning top N=10000: R30 of each L162_MASS run over the from-scratch mean, SD over runs and standard error of that mean in quadrature :: sha256 c6c4138ee5c86423 :: prose twin of the run farther from 1
\newcommand{\MlstFtTopRthirtyCoarseMassEFour}{365.8 and 359.9}  % experiments/FIGS/data/v2/index.json :: fine-tuning top N=10000: R30, R16_Q1_MASS :: sha256 c6c4138ee5c86423
\newcommand{\MlstFtTopRthirtyCoarseMassOverScratchEFour}{1.079 and 1.062}  % experiments/FIGS/data/v2/index.json :: fine-tuning top N=10000: R30 of each R16_Q1_MASS run over the from-scratch mean, SD over runs and standard error of that mean in quadrature :: sha256 c6c4138ee5c86423
\newcommand{\MlstFtTopRthirtyCoarseMassOverScratchEFourPr}{at least 6.2\% more rejection}  % experiments/FIGS/data/v2/index.json :: fine-tuning top N=10000: R30 of each R16_Q1_MASS run over the from-scratch mean, SD over runs and standard error of that mean in quadrature :: sha256 c6c4138ee5c86423 :: prose twin: two significant figures, direction in words
\newcommand{\MlstFtTopRthirtyCoarseMassOverScratchEFourPrMax}{up to 7.9\% more rejection}  % experiments/FIGS/data/v2/index.json :: fine-tuning top N=10000: R30 of each R16_Q1_MASS run over the from-scratch mean, SD over runs and standard error of that mean in quadrature :: sha256 c6c4138ee5c86423 :: prose twin of the run farther from 1
\newcommand{\MlstFtTopRthirtyResOverFullEFour}{1.02\,\ensuremath{\pm}\,0.12}  % experiments/FIGS/data/v2/index.json :: fine-tuning top N=10000: R30 of L162 over L188, run by run :: sha256 c6c4138ee5c86423
\newcommand{\MlstFtTopRthirtyResOverFullEFourPr}{equal within 20\%}  % experiments/FIGS/data/v2/index.json :: fine-tuning top N=10000: R30 of L162 over L188, run by run :: sha256 c6c4138ee5c86423 :: prose twin: two significant figures, direction in words
\newcommand{\MlstFtTopRthirtyMidOverFullEFour}{1.08\,\ensuremath{\pm}\,0.08}  % experiments/FIGS/data/v2/index.json :: fine-tuning top N=10000: R30 of R42_Q1 over L188, run by run :: sha256 c6c4138ee5c86423
\newcommand{\MlstFtTopRthirtyMidOverFullEFourPr}{equal within 9\%}  % experiments/FIGS/data/v2/index.json :: fine-tuning top N=10000: R30 of R42_Q1 over L188, run by run :: sha256 c6c4138ee5c86423 :: prose twin: two significant figures, direction in words
\newcommand{\MlstFtTopRthirtyCoarseOverFullEFour}{0.84\,\ensuremath{\pm}\,0.07}  % experiments/FIGS/data/v2/index.json :: fine-tuning top N=10000: R30 of R16_Q1 over L188, run by run :: sha256 c6c4138ee5c86423
\newcommand{\MlstFtTopRthirtyCoarseOverFullEFourPr}{16\% less rejection}  % experiments/FIGS/data/v2/index.json :: fine-tuning top N=10000: R30 of R16_Q1 over L188, run by run :: sha256 c6c4138ee5c86423 :: prose twin: two significant figures, direction in words
\newcommand{\MlstFtTopRthirtyResMassOverResEFour}{1.00 and 0.76}  % experiments/FIGS/data/v2/index.json :: fine-tuning top N=10000: R30 of L162_MASS over L162, run by run :: sha256 c6c4138ee5c86423
\newcommand{\MlstFtTopRthirtyResMassOverResEFourPr}{equal within 30\%}  % experiments/FIGS/data/v2/index.json :: fine-tuning top N=10000: R30 of L162_MASS over L162, run by run :: sha256 c6c4138ee5c86423 :: prose twin: two significant figures, direction in words
\newcommand{\MlstFtTopRthirtyResMassOverResEFourPrMax}{equal within 30\%}  % experiments/FIGS/data/v2/index.json :: fine-tuning top N=10000: R30 of L162_MASS over L162, run by run :: sha256 c6c4138ee5c86423 :: prose twin of the run farther from 1
\newcommand{\MlstFtTopRthirtyCoarseMassOverCoarseEFour}{1.13 and 0.92}  % experiments/FIGS/data/v2/index.json :: fine-tuning top N=10000: R30 of R16_Q1_MASS over R16_Q1, run by run :: sha256 c6c4138ee5c86423
\newcommand{\MlstFtTopRthirtyCoarseMassOverCoarseEFourPr}{equal within 20\%}  % experiments/FIGS/data/v2/index.json :: fine-tuning top N=10000: R30 of R16_Q1_MASS over R16_Q1, run by run :: sha256 c6c4138ee5c86423 :: prose twin: two significant figures, direction in words
\newcommand{\MlstFtTopRthirtyCoarseMassOverCoarseEFourPrMax}{equal within 20\%}  % experiments/FIGS/data/v2/index.json :: fine-tuning top N=10000: R30 of R16_Q1_MASS over R16_Q1, run by run :: sha256 c6c4138ee5c86423 :: prose twin of the run farther from 1
\newcommand{\MlstFtTopRthirtyRandOverFullEFour}{0.83\,\ensuremath{\pm}\,0.10}  % experiments/FIGS/data/v2/index.json :: fine-tuning top N=10000: R30 of each random partition over the L188 run of its seed :: sha256 c6c4138ee5c86423
\newcommand{\MlstFtTopRthirtyRandOverFullEFourPr}{equal within 20\%}  % experiments/FIGS/data/v2/index.json :: fine-tuning top N=10000: R30 of each random partition over the L188 run of its seed :: sha256 c6c4138ee5c86423 :: prose twin: two significant figures, direction in words
\newcommand{\MlstFtTopRthirtyRandOverCoarseEFour}{0.97\,\ensuremath{\pm}\,0.15}  % experiments/FIGS/data/v2/index.json :: fine-tuning top N=10000: R30 of each random partition over the R16_Q1 run of its seed :: sha256 c6c4138ee5c86423
\newcommand{\MlstFtTopRthirtyRandOverCoarseEFourPr}{equal within 20\%}  % experiments/FIGS/data/v2/index.json :: fine-tuning top N=10000: R30 of each random partition over the R16_Q1 run of its seed :: sha256 c6c4138ee5c86423 :: prose twin: two significant figures, direction in words
\newcommand{\MlstFtTopRthirtyScratchEFive}{560\,\ensuremath{\pm}\,90}  % experiments/FIGS/data/v2/index.json :: fine-tuning top N=100000: R30 from scratch, every seed :: sha256 c6c4138ee5c86423
\newcommand{\MlstFtTopRthirtyFullEFive}{850\,\ensuremath{\pm}\,90}  % experiments/FIGS/data/v2/index.json :: fine-tuning top N=100000: R30, L188 :: sha256 c6c4138ee5c86423
\newcommand{\MlstFtTopRthirtyFullOverScratchEFive}{1.50\,\ensuremath{\pm}\,0.21}  % experiments/FIGS/data/v2/index.json :: fine-tuning top N=100000: R30 of each L188 run over the from-scratch mean, SD over runs and standard error of that mean in quadrature :: sha256 c6c4138ee5c86423
\newcommand{\MlstFtTopRthirtyFullOverScratchEFivePr}{50\% more rejection}  % experiments/FIGS/data/v2/index.json :: fine-tuning top N=100000: R30 of each L188 run over the from-scratch mean, SD over runs and standard error of that mean in quadrature :: sha256 c6c4138ee5c86423 :: prose twin: two significant figures, direction in words
\newcommand{\MlstFtTopRthirtyResEFive}{810\,\ensuremath{\pm}\,60}  % experiments/FIGS/data/v2/index.json :: fine-tuning top N=100000: R30, L162 :: sha256 c6c4138ee5c86423
\newcommand{\MlstFtTopRthirtyResOverScratchEFive}{1.44\,\ensuremath{\pm}\,0.18}  % experiments/FIGS/data/v2/index.json :: fine-tuning top N=100000: R30 of each L162 run over the from-scratch mean, SD over runs and standard error of that mean in quadrature :: sha256 c6c4138ee5c86423
\newcommand{\MlstFtTopRthirtyResOverScratchEFivePr}{44\% more rejection}  % experiments/FIGS/data/v2/index.json :: fine-tuning top N=100000: R30 of each L162 run over the from-scratch mean, SD over runs and standard error of that mean in quadrature :: sha256 c6c4138ee5c86423 :: prose twin: two significant figures, direction in words
\newcommand{\MlstFtTopRthirtyMidEFive}{850\,\ensuremath{\pm}\,80}  % experiments/FIGS/data/v2/index.json :: fine-tuning top N=100000: R30, R42_Q1 :: sha256 c6c4138ee5c86423
\newcommand{\MlstFtTopRthirtyMidOverScratchEFive}{1.51\,\ensuremath{\pm}\,0.20}  % experiments/FIGS/data/v2/index.json :: fine-tuning top N=100000: R30 of each R42_Q1 run over the from-scratch mean, SD over runs and standard error of that mean in quadrature :: sha256 c6c4138ee5c86423
\newcommand{\MlstFtTopRthirtyMidOverScratchEFivePr}{51\% more rejection}  % experiments/FIGS/data/v2/index.json :: fine-tuning top N=100000: R30 of each R42_Q1 run over the from-scratch mean, SD over runs and standard error of that mean in quadrature :: sha256 c6c4138ee5c86423 :: prose twin: two significant figures, direction in words
\newcommand{\MlstFtTopRthirtyCoarseEFive}{830\,\ensuremath{\pm}\,60}  % experiments/FIGS/data/v2/index.json :: fine-tuning top N=100000: R30, R16_Q1 :: sha256 c6c4138ee5c86423
\newcommand{\MlstFtTopRthirtyCoarseOverScratchEFive}{1.47\,\ensuremath{\pm}\,0.17}  % experiments/FIGS/data/v2/index.json :: fine-tuning top N=100000: R30 of each R16_Q1 run over the from-scratch mean, SD over runs and standard error of that mean in quadrature :: sha256 c6c4138ee5c86423
\newcommand{\MlstFtTopRthirtyCoarseOverScratchEFivePr}{47\% more rejection}  % experiments/FIGS/data/v2/index.json :: fine-tuning top N=100000: R30 of each R16_Q1 run over the from-scratch mean, SD over runs and standard error of that mean in quadrature :: sha256 c6c4138ee5c86423 :: prose twin: two significant figures, direction in words
\newcommand{\MlstFtTopRthirtyResMassEFive}{848 and 848}  % experiments/FIGS/data/v2/index.json :: fine-tuning top N=100000: R30, L162_MASS :: sha256 c6c4138ee5c86423
\newcommand{\MlstFtTopRthirtyResMassOverScratchEFive}{1.50 and 1.50}  % experiments/FIGS/data/v2/index.json :: fine-tuning top N=100000: R30 of each L162_MASS run over the from-scratch mean, SD over runs and standard error of that mean in quadrature :: sha256 c6c4138ee5c86423
\newcommand{\MlstFtTopRthirtyResMassOverScratchEFivePr}{at least 50\% more rejection}  % experiments/FIGS/data/v2/index.json :: fine-tuning top N=100000: R30 of each L162_MASS run over the from-scratch mean, SD over runs and standard error of that mean in quadrature :: sha256 c6c4138ee5c86423 :: prose twin: two significant figures, direction in words
\newcommand{\MlstFtTopRthirtyResMassOverScratchEFivePrMax}{up to 50\% more rejection}  % experiments/FIGS/data/v2/index.json :: fine-tuning top N=100000: R30 of each L162_MASS run over the from-scratch mean, SD over runs and standard error of that mean in quadrature :: sha256 c6c4138ee5c86423 :: prose twin of the run farther from 1
\newcommand{\MlstFtTopRthirtyCoarseMassEFive}{856 and 774}  % experiments/FIGS/data/v2/index.json :: fine-tuning top N=100000: R30, R16_Q1_MASS :: sha256 c6c4138ee5c86423
\newcommand{\MlstFtTopRthirtyCoarseMassOverScratchEFive}{1.51 and 1.37}  % experiments/FIGS/data/v2/index.json :: fine-tuning top N=100000: R30 of each R16_Q1_MASS run over the from-scratch mean, SD over runs and standard error of that mean in quadrature :: sha256 c6c4138ee5c86423
\newcommand{\MlstFtTopRthirtyCoarseMassOverScratchEFivePr}{at least 37\% more rejection}  % experiments/FIGS/data/v2/index.json :: fine-tuning top N=100000: R30 of each R16_Q1_MASS run over the from-scratch mean, SD over runs and standard error of that mean in quadrature :: sha256 c6c4138ee5c86423 :: prose twin: two significant figures, direction in words
\newcommand{\MlstFtTopRthirtyCoarseMassOverScratchEFivePrMax}{up to 51\% more rejection}  % experiments/FIGS/data/v2/index.json :: fine-tuning top N=100000: R30 of each R16_Q1_MASS run over the from-scratch mean, SD over runs and standard error of that mean in quadrature :: sha256 c6c4138ee5c86423 :: prose twin of the run farther from 1
\newcommand{\MlstFtTopRthirtyResOverFullEFive}{0.96\,\ensuremath{\pm}\,0.05}  % experiments/FIGS/data/v2/index.json :: fine-tuning top N=100000: R30 of L162 over L188, run by run :: sha256 c6c4138ee5c86423
\newcommand{\MlstFtTopRthirtyResOverFullEFivePr}{equal within 6\%}  % experiments/FIGS/data/v2/index.json :: fine-tuning top N=100000: R30 of L162 over L188, run by run :: sha256 c6c4138ee5c86423 :: prose twin: two significant figures, direction in words
\newcommand{\MlstFtTopRthirtyMidOverFullEFive}{1.02\,\ensuremath{\pm}\,0.18}  % experiments/FIGS/data/v2/index.json :: fine-tuning top N=100000: R30 of R42_Q1 over L188, run by run :: sha256 c6c4138ee5c86423
\newcommand{\MlstFtTopRthirtyMidOverFullEFivePr}{equal within 20\%}  % experiments/FIGS/data/v2/index.json :: fine-tuning top N=100000: R30 of R42_Q1 over L188, run by run :: sha256 c6c4138ee5c86423 :: prose twin: two significant figures, direction in words
\newcommand{\MlstFtTopRthirtyCoarseOverFullEFive}{0.99\,\ensuremath{\pm}\,0.17}  % experiments/FIGS/data/v2/index.json :: fine-tuning top N=100000: R30 of R16_Q1 over L188, run by run :: sha256 c6c4138ee5c86423
\newcommand{\MlstFtTopRthirtyCoarseOverFullEFivePr}{equal within 20\%}  % experiments/FIGS/data/v2/index.json :: fine-tuning top N=100000: R30 of R16_Q1 over L188, run by run :: sha256 c6c4138ee5c86423 :: prose twin: two significant figures, direction in words
\newcommand{\MlstFtTopRthirtyResMassOverResEFive}{1.00 and 1.15}  % experiments/FIGS/data/v2/index.json :: fine-tuning top N=100000: R30 of L162_MASS over L162, run by run :: sha256 c6c4138ee5c86423
\newcommand{\MlstFtTopRthirtyResMassOverResEFivePr}{equal within 20\%}  % experiments/FIGS/data/v2/index.json :: fine-tuning top N=100000: R30 of L162_MASS over L162, run by run :: sha256 c6c4138ee5c86423 :: prose twin: two significant figures, direction in words
\newcommand{\MlstFtTopRthirtyResMassOverResEFivePrMax}{equal within 20\%}  % experiments/FIGS/data/v2/index.json :: fine-tuning top N=100000: R30 of L162_MASS over L162, run by run :: sha256 c6c4138ee5c86423 :: prose twin of the run farther from 1
\newcommand{\MlstFtTopRthirtyCoarseMassOverCoarseEFive}{1.08 and 0.87}  % experiments/FIGS/data/v2/index.json :: fine-tuning top N=100000: R30 of R16_Q1_MASS over R16_Q1, run by run :: sha256 c6c4138ee5c86423
\newcommand{\MlstFtTopRthirtyCoarseMassOverCoarseEFivePr}{equal within 20\%}  % experiments/FIGS/data/v2/index.json :: fine-tuning top N=100000: R30 of R16_Q1_MASS over R16_Q1, run by run :: sha256 c6c4138ee5c86423 :: prose twin: two significant figures, direction in words
\newcommand{\MlstFtTopRthirtyCoarseMassOverCoarseEFivePrMax}{equal within 20\%}  % experiments/FIGS/data/v2/index.json :: fine-tuning top N=100000: R30 of R16_Q1_MASS over R16_Q1, run by run :: sha256 c6c4138ee5c86423 :: prose twin of the run farther from 1
\newcommand{\MlstFtTopRthirtyRandOverFullEFive}{1.01\,\ensuremath{\pm}\,0.22}  % experiments/FIGS/data/v2/index.json :: fine-tuning top N=100000: R30 of each random partition over the L188 run of its seed :: sha256 c6c4138ee5c86423
\newcommand{\MlstFtTopRthirtyRandOverFullEFivePr}{equal within 30\%}  % experiments/FIGS/data/v2/index.json :: fine-tuning top N=100000: R30 of each random partition over the L188 run of its seed :: sha256 c6c4138ee5c86423 :: prose twin: two significant figures, direction in words
\newcommand{\MlstFtTopRthirtyRandOverCoarseEFive}{0.97\,\ensuremath{\pm}\,0.10}  % experiments/FIGS/data/v2/index.json :: fine-tuning top N=100000: R30 of each random partition over the R16_Q1 run of its seed :: sha256 c6c4138ee5c86423
\newcommand{\MlstFtTopRthirtyRandOverCoarseEFivePr}{equal within 10\%}  % experiments/FIGS/data/v2/index.json :: fine-tuning top N=100000: R30 of each random partition over the R16_Q1 run of its seed :: sha256 c6c4138ee5c86423 :: prose twin: two significant figures, direction in words
\newcommand{\MlstFtTopRthirtyScratch}{1{,}412\,\ensuremath{\pm}\,20}  % experiments/FIGS/data/v2/index.json :: fine-tuning top N=1200000: R30 from scratch, every seed :: sha256 c6c4138ee5c86423
\newcommand{\MlstFtTopRthirtyFull}{1{,}800\,\ensuremath{\pm}\,60}  % experiments/FIGS/data/v2/index.json :: fine-tuning top N=1200000: R30, L188 :: sha256 c6c4138ee5c86423
\newcommand{\MlstFtTopRthirtyFullOverScratch}{1.27\,\ensuremath{\pm}\,0.04}  % experiments/FIGS/data/v2/index.json :: fine-tuning top N=1200000: R30 of each L188 run over the from-scratch mean, SD over runs and standard error of that mean in quadrature :: sha256 c6c4138ee5c86423
\newcommand{\MlstFtTopRthirtyFullOverScratchPr}{27\% more rejection}  % experiments/FIGS/data/v2/index.json :: fine-tuning top N=1200000: R30 of each L188 run over the from-scratch mean, SD over runs and standard error of that mean in quadrature :: sha256 c6c4138ee5c86423 :: prose twin: two significant figures, direction in words
\newcommand{\MlstFtTopRthirtyRes}{1{,}790\,\ensuremath{\pm}\,60}  % experiments/FIGS/data/v2/index.json :: fine-tuning top N=1200000: R30, L162 :: sha256 c6c4138ee5c86423
\newcommand{\MlstFtTopRthirtyResOverScratch}{1.27\,\ensuremath{\pm}\,0.04}  % experiments/FIGS/data/v2/index.json :: fine-tuning top N=1200000: R30 of each L162 run over the from-scratch mean, SD over runs and standard error of that mean in quadrature :: sha256 c6c4138ee5c86423
\newcommand{\MlstFtTopRthirtyResOverScratchPr}{27\% more rejection}  % experiments/FIGS/data/v2/index.json :: fine-tuning top N=1200000: R30 of each L162 run over the from-scratch mean, SD over runs and standard error of that mean in quadrature :: sha256 c6c4138ee5c86423 :: prose twin: two significant figures, direction in words
\newcommand{\MlstFtTopRthirtyMid}{1{,}780\,\ensuremath{\pm}\,150}  % experiments/FIGS/data/v2/index.json :: fine-tuning top N=1200000: R30, R42_Q1 :: sha256 c6c4138ee5c86423
\newcommand{\MlstFtTopRthirtyMidOverScratch}{1.26\,\ensuremath{\pm}\,0.11}  % experiments/FIGS/data/v2/index.json :: fine-tuning top N=1200000: R30 of each R42_Q1 run over the from-scratch mean, SD over runs and standard error of that mean in quadrature :: sha256 c6c4138ee5c86423
\newcommand{\MlstFtTopRthirtyMidOverScratchPr}{26\% more rejection}  % experiments/FIGS/data/v2/index.json :: fine-tuning top N=1200000: R30 of each R42_Q1 run over the from-scratch mean, SD over runs and standard error of that mean in quadrature :: sha256 c6c4138ee5c86423 :: prose twin: two significant figures, direction in words
\newcommand{\MlstFtTopRthirtyCoarse}{1{,}740\,\ensuremath{\pm}\,140}  % experiments/FIGS/data/v2/index.json :: fine-tuning top N=1200000: R30, R16_Q1 :: sha256 c6c4138ee5c86423
\newcommand{\MlstFtTopRthirtyCoarseOverScratch}{1.23\,\ensuremath{\pm}\,0.10}  % experiments/FIGS/data/v2/index.json :: fine-tuning top N=1200000: R30 of each R16_Q1 run over the from-scratch mean, SD over runs and standard error of that mean in quadrature :: sha256 c6c4138ee5c86423
\newcommand{\MlstFtTopRthirtyCoarseOverScratchPr}{23\% more rejection}  % experiments/FIGS/data/v2/index.json :: fine-tuning top N=1200000: R30 of each R16_Q1 run over the from-scratch mean, SD over runs and standard error of that mean in quadrature :: sha256 c6c4138ee5c86423 :: prose twin: two significant figures, direction in words
\newcommand{\MlstFtTopRthirtyResMass}{1{,}941 and 1{,}590}  % experiments/FIGS/data/v2/index.json :: fine-tuning top N=1200000: R30, L162_MASS :: sha256 c6c4138ee5c86423
\newcommand{\MlstFtTopRthirtyResMassOverScratch}{1.37 and 1.13}  % experiments/FIGS/data/v2/index.json :: fine-tuning top N=1200000: R30 of each L162_MASS run over the from-scratch mean, SD over runs and standard error of that mean in quadrature :: sha256 c6c4138ee5c86423
\newcommand{\MlstFtTopRthirtyResMassOverScratchPr}{at least 13\% more rejection}  % experiments/FIGS/data/v2/index.json :: fine-tuning top N=1200000: R30 of each L162_MASS run over the from-scratch mean, SD over runs and standard error of that mean in quadrature :: sha256 c6c4138ee5c86423 :: prose twin: two significant figures, direction in words
\newcommand{\MlstFtTopRthirtyResMassOverScratchPrMax}{up to 37\% more rejection}  % experiments/FIGS/data/v2/index.json :: fine-tuning top N=1200000: R30 of each L162_MASS run over the from-scratch mean, SD over runs and standard error of that mean in quadrature :: sha256 c6c4138ee5c86423 :: prose twin of the run farther from 1
\newcommand{\MlstFtTopRthirtyCoarseMass}{1{,}683 and 1{,}628}  % experiments/FIGS/data/v2/index.json :: fine-tuning top N=1200000: R30, R16_Q1_MASS :: sha256 c6c4138ee5c86423
\newcommand{\MlstFtTopRthirtyCoarseMassOverScratch}{1.192 and 1.153}  % experiments/FIGS/data/v2/index.json :: fine-tuning top N=1200000: R30 of each R16_Q1_MASS run over the from-scratch mean, SD over runs and standard error of that mean in quadrature :: sha256 c6c4138ee5c86423
\newcommand{\MlstFtTopRthirtyCoarseMassOverScratchPr}{at least 15\% more rejection}  % experiments/FIGS/data/v2/index.json :: fine-tuning top N=1200000: R30 of each R16_Q1_MASS run over the from-scratch mean, SD over runs and standard error of that mean in quadrature :: sha256 c6c4138ee5c86423 :: prose twin: two significant figures, direction in words
\newcommand{\MlstFtTopRthirtyCoarseMassOverScratchPrMax}{up to 19\% more rejection}  % experiments/FIGS/data/v2/index.json :: fine-tuning top N=1200000: R30 of each R16_Q1_MASS run over the from-scratch mean, SD over runs and standard error of that mean in quadrature :: sha256 c6c4138ee5c86423 :: prose twin of the run farther from 1
\newcommand{\MlstFtTopRthirtyResOverFull}{0.997\,\ensuremath{\pm}\,0.032}  % experiments/FIGS/data/v2/index.json :: fine-tuning top N=1200000: R30 of L162 over L188, run by run :: sha256 c6c4138ee5c86423
\newcommand{\MlstFtTopRthirtyResOverFullPr}{equal within 4\%}  % experiments/FIGS/data/v2/index.json :: fine-tuning top N=1200000: R30 of L162 over L188, run by run :: sha256 c6c4138ee5c86423 :: prose twin: two significant figures, direction in words
\newcommand{\MlstFtTopRthirtyMidOverFull}{0.99\,\ensuremath{\pm}\,0.07}  % experiments/FIGS/data/v2/index.json :: fine-tuning top N=1200000: R30 of R42_Q1 over L188, run by run :: sha256 c6c4138ee5c86423
\newcommand{\MlstFtTopRthirtyMidOverFullPr}{equal within 8\%}  % experiments/FIGS/data/v2/index.json :: fine-tuning top N=1200000: R30 of R42_Q1 over L188, run by run :: sha256 c6c4138ee5c86423 :: prose twin: two significant figures, direction in words
\newcommand{\MlstFtTopRthirtyCoarseOverFull}{0.97\,\ensuremath{\pm}\,0.10}  % experiments/FIGS/data/v2/index.json :: fine-tuning top N=1200000: R30 of R16_Q1 over L188, run by run :: sha256 c6c4138ee5c86423
\newcommand{\MlstFtTopRthirtyCoarseOverFullPr}{equal within 10\%}  % experiments/FIGS/data/v2/index.json :: fine-tuning top N=1200000: R30 of R16_Q1 over L188, run by run :: sha256 c6c4138ee5c86423 :: prose twin: two significant figures, direction in words
\newcommand{\MlstFtTopRthirtyResMassOverRes}{1.12 and 0.89}  % experiments/FIGS/data/v2/index.json :: fine-tuning top N=1200000: R30 of L162_MASS over L162, run by run :: sha256 c6c4138ee5c86423
\newcommand{\MlstFtTopRthirtyResMassOverResPr}{equal within 20\%}  % experiments/FIGS/data/v2/index.json :: fine-tuning top N=1200000: R30 of L162_MASS over L162, run by run :: sha256 c6c4138ee5c86423 :: prose twin: two significant figures, direction in words
\newcommand{\MlstFtTopRthirtyResMassOverResPrMax}{equal within 20\%}  % experiments/FIGS/data/v2/index.json :: fine-tuning top N=1200000: R30 of L162_MASS over L162, run by run :: sha256 c6c4138ee5c86423 :: prose twin of the run farther from 1
\newcommand{\MlstFtTopRthirtyCoarseMassOverCoarse}{0.900 and 0.919}  % experiments/FIGS/data/v2/index.json :: fine-tuning top N=1200000: R30 of R16_Q1_MASS over R16_Q1, run by run :: sha256 c6c4138ee5c86423
\newcommand{\MlstFtTopRthirtyCoarseMassOverCoarsePr}{equal within 10\%}  % experiments/FIGS/data/v2/index.json :: fine-tuning top N=1200000: R30 of R16_Q1_MASS over R16_Q1, run by run :: sha256 c6c4138ee5c86423 :: prose twin: two significant figures, direction in words
\newcommand{\MlstFtTopRthirtyCoarseMassOverCoarsePrMax}{equal within 10\%}  % experiments/FIGS/data/v2/index.json :: fine-tuning top N=1200000: R30 of R16_Q1_MASS over R16_Q1, run by run :: sha256 c6c4138ee5c86423 :: prose twin of the run farther from 1
\newcommand{\MlstFtTopRthirtyRandOverFull}{0.97\,\ensuremath{\pm}\,0.04}  % experiments/FIGS/data/v2/index.json :: fine-tuning top N=1200000: R30 of each random partition over the L188 run of its seed :: sha256 c6c4138ee5c86423
\newcommand{\MlstFtTopRthirtyRandOverFullPr}{equal within 4\%}  % experiments/FIGS/data/v2/index.json :: fine-tuning top N=1200000: R30 of each random partition over the L188 run of its seed :: sha256 c6c4138ee5c86423 :: prose twin: two significant figures, direction in words
\newcommand{\MlstFtTopRthirtyRandOverCoarse}{0.97\,\ensuremath{\pm}\,0.09}  % experiments/FIGS/data/v2/index.json :: fine-tuning top N=1200000: R30 of each random partition over the R16_Q1 run of its seed :: sha256 c6c4138ee5c86423
\newcommand{\MlstFtTopRthirtyRandOverCoarsePr}{equal within 9\%}  % experiments/FIGS/data/v2/index.json :: fine-tuning top N=1200000: R30 of each random partition over the R16_Q1 run of its seed :: sha256 c6c4138ee5c86423 :: prose twin: two significant figures, direction in words
\newcommand{\MlstFtTopRthirtyFullOverPart}{0.650\,\ensuremath{\pm}\,0.020}  % experiments/FIGS/data/v2/index.json :: fine-tuning top N=1200000: R30 of each 188-class run over ParT-f.t.'s published value :: sha256 c6c4138ee5c86423
\newcommand{\MlstFtTopRthirtyFullOverPartPr}{35\% less rejection}  % experiments/FIGS/data/v2/index.json :: fine-tuning top N=1200000: R30 of each 188-class run over ParT-f.t.'s published value :: sha256 c6c4138ee5c86423 :: prose twin: two significant figures, direction in words
\newcommand{\MlstFtTopAucFull}{0.9861\,\ensuremath{\pm}\,0.0001}  % experiments/FIGS/data/v2/index.json :: fine-tuning top N=1200000: AUC of the 188-class runs :: sha256 c6c4138ee5c86423
\newcommand{\MlstFtTopAucScratch}{0.9852\,\ensuremath{\pm}\,0.0001}  % experiments/FIGS/data/v2/index.json :: fine-tuning top N=1200000: AUC from scratch :: sha256 c6c4138ee5c86423
\newcommand{\MlstFtTopRthirtyFullRelSd}{3\%}  % experiments/FIGS/data/v2/index.json :: fine-tuning top, full set: relative SD of R30 over the 188-class runs :: sha256 c6c4138ee5c86423
\newcommand{\MlstFtTopRthirtyKin}{1{,}800}  % experiments/FIGS/data/v2/index.json :: fine-tuning top, full set: R30 of the kinematic-only model (one run) :: sha256 c6c4138ee5c86423
\newcommand{\MlstFtTopRthirtyKinOverFull}{1.00}  % experiments/FIGS/data/v2/index.json :: fine-tuning top, full set: R30 of the kinematic-only model over the mean of the 188-class runs :: sha256 c6c4138ee5c86423
\newcommand{\MlstFtTopRthirtyKinOverFullPr}{0.23\% more rejection}  % experiments/FIGS/data/v2/index.json :: fine-tuning top, full set: R30 of the kinematic-only model over the mean of the 188-class runs :: sha256 c6c4138ee5c86423 :: prose twin: two significant figures, direction in words
\newcommand{\MlstFtTopDataEquivJets}{1{,}200}  % experiments/FIGS/data/v2/index.json :: fine-tuning top: from-scratch training jets whose R30 equals the 188-class mean at 1000 jets (interpolated linearly in R30 against log10 N) :: sha256 c6c4138ee5c86423
\newcommand{\MlstFtTopDataEquivJetsEFour}{23{,}000}  % experiments/FIGS/data/v2/index.json :: fine-tuning top: from-scratch training jets whose R30 equals the 188-class mean at 10000 jets (interpolated linearly in R30 against log10 N) :: sha256 c6c4138ee5c86423
\newcommand{\MlstFtQgRthirtyScratchEThree}{49.7\,\ensuremath{\pm}\,2.1}  % experiments/FIGS/data/v2/index.json :: fine-tuning qg N=1000: R30 from scratch, every seed :: sha256 c6c4138ee5c86423
\newcommand{\MlstFtQgRthirtyFullEThree}{47\,\ensuremath{\pm}\,10}  % experiments/FIGS/data/v2/index.json :: fine-tuning qg N=1000: R30, L188 :: sha256 c6c4138ee5c86423
\newcommand{\MlstFtQgRthirtyFullOverScratchEThree}{0.95\,\ensuremath{\pm}\,0.19}  % experiments/FIGS/data/v2/index.json :: fine-tuning qg N=1000: R30 of each L188 run over the from-scratch mean, SD over runs and standard error of that mean in quadrature :: sha256 c6c4138ee5c86423
\newcommand{\MlstFtQgRthirtyFullOverScratchEThreePr}{equal within 20\%}  % experiments/FIGS/data/v2/index.json :: fine-tuning qg N=1000: R30 of each L188 run over the from-scratch mean, SD over runs and standard error of that mean in quadrature :: sha256 c6c4138ee5c86423 :: prose twin: two significant figures, direction in words
\newcommand{\MlstFtQgRthirtyResEThree}{52.4\,\ensuremath{\pm}\,3.1}  % experiments/FIGS/data/v2/index.json :: fine-tuning qg N=1000: R30, L162 :: sha256 c6c4138ee5c86423
\newcommand{\MlstFtQgRthirtyResOverScratchEThree}{1.06\,\ensuremath{\pm}\,0.07}  % experiments/FIGS/data/v2/index.json :: fine-tuning qg N=1000: R30 of each L162 run over the from-scratch mean, SD over runs and standard error of that mean in quadrature :: sha256 c6c4138ee5c86423
\newcommand{\MlstFtQgRthirtyResOverScratchEThreePr}{equal within 7\%}  % experiments/FIGS/data/v2/index.json :: fine-tuning qg N=1000: R30 of each L162 run over the from-scratch mean, SD over runs and standard error of that mean in quadrature :: sha256 c6c4138ee5c86423 :: prose twin: two significant figures, direction in words
\newcommand{\MlstFtQgRthirtyMidEThree}{54\,\ensuremath{\pm}\,6}  % experiments/FIGS/data/v2/index.json :: fine-tuning qg N=1000: R30, R42_Q1 :: sha256 c6c4138ee5c86423
\newcommand{\MlstFtQgRthirtyMidOverScratchEThree}{1.09\,\ensuremath{\pm}\,0.12}  % experiments/FIGS/data/v2/index.json :: fine-tuning qg N=1000: R30 of each R42_Q1 run over the from-scratch mean, SD over runs and standard error of that mean in quadrature :: sha256 c6c4138ee5c86423
\newcommand{\MlstFtQgRthirtyMidOverScratchEThreePr}{equal within 20\%}  % experiments/FIGS/data/v2/index.json :: fine-tuning qg N=1000: R30 of each R42_Q1 run over the from-scratch mean, SD over runs and standard error of that mean in quadrature :: sha256 c6c4138ee5c86423 :: prose twin: two significant figures, direction in words
\newcommand{\MlstFtQgRthirtyCoarseEThree}{55\,\ensuremath{\pm}\,5}  % experiments/FIGS/data/v2/index.json :: fine-tuning qg N=1000: R30, R16_Q1 :: sha256 c6c4138ee5c86423
\newcommand{\MlstFtQgRthirtyCoarseOverScratchEThree}{1.12\,\ensuremath{\pm}\,0.10}  % experiments/FIGS/data/v2/index.json :: fine-tuning qg N=1000: R30 of each R16_Q1 run over the from-scratch mean, SD over runs and standard error of that mean in quadrature :: sha256 c6c4138ee5c86423
\newcommand{\MlstFtQgRthirtyCoarseOverScratchEThreePr}{equal within 20\%}  % experiments/FIGS/data/v2/index.json :: fine-tuning qg N=1000: R30 of each R16_Q1 run over the from-scratch mean, SD over runs and standard error of that mean in quadrature :: sha256 c6c4138ee5c86423 :: prose twin: two significant figures, direction in words
\newcommand{\MlstFtQgRthirtyResMassEThree}{59.3 and 57.1}  % experiments/FIGS/data/v2/index.json :: fine-tuning qg N=1000: R30, L162_MASS :: sha256 c6c4138ee5c86423
\newcommand{\MlstFtQgRthirtyResMassOverScratchEThree}{1.193 and 1.149}  % experiments/FIGS/data/v2/index.json :: fine-tuning qg N=1000: R30 of each L162_MASS run over the from-scratch mean, SD over runs and standard error of that mean in quadrature :: sha256 c6c4138ee5c86423
\newcommand{\MlstFtQgRthirtyResMassOverScratchEThreePr}{at least 15\% more rejection}  % experiments/FIGS/data/v2/index.json :: fine-tuning qg N=1000: R30 of each L162_MASS run over the from-scratch mean, SD over runs and standard error of that mean in quadrature :: sha256 c6c4138ee5c86423 :: prose twin: two significant figures, direction in words
\newcommand{\MlstFtQgRthirtyResMassOverScratchEThreePrMax}{up to 19\% more rejection}  % experiments/FIGS/data/v2/index.json :: fine-tuning qg N=1000: R30 of each L162_MASS run over the from-scratch mean, SD over runs and standard error of that mean in quadrature :: sha256 c6c4138ee5c86423 :: prose twin of the run farther from 1
\newcommand{\MlstFtQgRthirtyCoarseMassEThree}{54.2 and 57.7}  % experiments/FIGS/data/v2/index.json :: fine-tuning qg N=1000: R30, R16_Q1_MASS :: sha256 c6c4138ee5c86423
\newcommand{\MlstFtQgRthirtyCoarseMassOverScratchEThree}{1.091 and 1.161}  % experiments/FIGS/data/v2/index.json :: fine-tuning qg N=1000: R30 of each R16_Q1_MASS run over the from-scratch mean, SD over runs and standard error of that mean in quadrature :: sha256 c6c4138ee5c86423
\newcommand{\MlstFtQgRthirtyCoarseMassOverScratchEThreePr}{at least 9.1\% more rejection}  % experiments/FIGS/data/v2/index.json :: fine-tuning qg N=1000: R30 of each R16_Q1_MASS run over the from-scratch mean, SD over runs and standard error of that mean in quadrature :: sha256 c6c4138ee5c86423 :: prose twin: two significant figures, direction in words
\newcommand{\MlstFtQgRthirtyCoarseMassOverScratchEThreePrMax}{up to 16\% more rejection}  % experiments/FIGS/data/v2/index.json :: fine-tuning qg N=1000: R30 of each R16_Q1_MASS run over the from-scratch mean, SD over runs and standard error of that mean in quadrature :: sha256 c6c4138ee5c86423 :: prose twin of the run farther from 1
\newcommand{\MlstFtQgRthirtyResOverFullEThree}{1.14\,\ensuremath{\pm}\,0.17}  % experiments/FIGS/data/v2/index.json :: fine-tuning qg N=1000: R30 of L162 over L188, run by run :: sha256 c6c4138ee5c86423
\newcommand{\MlstFtQgRthirtyResOverFullEThreePr}{equal within 20\%}  % experiments/FIGS/data/v2/index.json :: fine-tuning qg N=1000: R30 of L162 over L188, run by run :: sha256 c6c4138ee5c86423 :: prose twin: two significant figures, direction in words
\newcommand{\MlstFtQgRthirtyMidOverFullEThree}{1.17\,\ensuremath{\pm}\,0.14}  % experiments/FIGS/data/v2/index.json :: fine-tuning qg N=1000: R30 of R42_Q1 over L188, run by run :: sha256 c6c4138ee5c86423
\newcommand{\MlstFtQgRthirtyMidOverFullEThreePr}{equal within 20\%}  % experiments/FIGS/data/v2/index.json :: fine-tuning qg N=1000: R30 of R42_Q1 over L188, run by run :: sha256 c6c4138ee5c86423 :: prose twin: two significant figures, direction in words
\newcommand{\MlstFtQgRthirtyCoarseOverFullEThree}{1.21\,\ensuremath{\pm}\,0.27}  % experiments/FIGS/data/v2/index.json :: fine-tuning qg N=1000: R30 of R16_Q1 over L188, run by run :: sha256 c6c4138ee5c86423
\newcommand{\MlstFtQgRthirtyCoarseOverFullEThreePr}{equal within 30\%}  % experiments/FIGS/data/v2/index.json :: fine-tuning qg N=1000: R30 of R16_Q1 over L188, run by run :: sha256 c6c4138ee5c86423 :: prose twin: two significant figures, direction in words
\newcommand{\MlstFtQgRthirtyResMassOverResEThree}{1.108 and 1.042}  % experiments/FIGS/data/v2/index.json :: fine-tuning qg N=1000: R30 of L162_MASS over L162, run by run :: sha256 c6c4138ee5c86423
\newcommand{\MlstFtQgRthirtyResMassOverResEThreePr}{equal within 20\%}  % experiments/FIGS/data/v2/index.json :: fine-tuning qg N=1000: R30 of L162_MASS over L162, run by run :: sha256 c6c4138ee5c86423 :: prose twin: two significant figures, direction in words
\newcommand{\MlstFtQgRthirtyResMassOverResEThreePrMax}{equal within 20\%}  % experiments/FIGS/data/v2/index.json :: fine-tuning qg N=1000: R30 of L162_MASS over L162, run by run :: sha256 c6c4138ee5c86423 :: prose twin of the run farther from 1
\newcommand{\MlstFtQgRthirtyCoarseMassOverCoarseEThree}{0.89 and 1.11}  % experiments/FIGS/data/v2/index.json :: fine-tuning qg N=1000: R30 of R16_Q1_MASS over R16_Q1, run by run :: sha256 c6c4138ee5c86423
\newcommand{\MlstFtQgRthirtyCoarseMassOverCoarseEThreePr}{equal within 20\%}  % experiments/FIGS/data/v2/index.json :: fine-tuning qg N=1000: R30 of R16_Q1_MASS over R16_Q1, run by run :: sha256 c6c4138ee5c86423 :: prose twin: two significant figures, direction in words
\newcommand{\MlstFtQgRthirtyCoarseMassOverCoarseEThreePrMax}{equal within 20\%}  % experiments/FIGS/data/v2/index.json :: fine-tuning qg N=1000: R30 of R16_Q1_MASS over R16_Q1, run by run :: sha256 c6c4138ee5c86423 :: prose twin of the run farther from 1
\newcommand{\MlstFtQgRthirtyRandOverFullEThree}{0.91\,\ensuremath{\pm}\,0.07}  % experiments/FIGS/data/v2/index.json :: fine-tuning qg N=1000: R30 of each random partition over the L188 run of its seed :: sha256 c6c4138ee5c86423
\newcommand{\MlstFtQgRthirtyRandOverFullEThreePr}{equal within 10\%}  % experiments/FIGS/data/v2/index.json :: fine-tuning qg N=1000: R30 of each random partition over the L188 run of its seed :: sha256 c6c4138ee5c86423 :: prose twin: two significant figures, direction in words
\newcommand{\MlstFtQgRthirtyRandOverCoarseEThree}{0.87\,\ensuremath{\pm}\,0.14}  % experiments/FIGS/data/v2/index.json :: fine-tuning qg N=1000: R30 of each random partition over the R16_Q1 run of its seed :: sha256 c6c4138ee5c86423
\newcommand{\MlstFtQgRthirtyRandOverCoarseEThreePr}{equal within 20\%}  % experiments/FIGS/data/v2/index.json :: fine-tuning qg N=1000: R30 of each random partition over the R16_Q1 run of its seed :: sha256 c6c4138ee5c86423 :: prose twin: two significant figures, direction in words
\newcommand{\MlstFtQgRthirtyScratchEFour}{66\,\ensuremath{\pm}\,9}  % experiments/FIGS/data/v2/index.json :: fine-tuning qg N=10000: R30 from scratch, every seed :: sha256 c6c4138ee5c86423
\newcommand{\MlstFtQgRthirtyFullEFour}{75\,\ensuremath{\pm}\,7}  % experiments/FIGS/data/v2/index.json :: fine-tuning qg N=10000: R30, L188 :: sha256 c6c4138ee5c86423
\newcommand{\MlstFtQgRthirtyFullOverScratchEFour}{1.14\,\ensuremath{\pm}\,0.14}  % experiments/FIGS/data/v2/index.json :: fine-tuning qg N=10000: R30 of each L188 run over the from-scratch mean, SD over runs and standard error of that mean in quadrature :: sha256 c6c4138ee5c86423
\newcommand{\MlstFtQgRthirtyFullOverScratchEFourPr}{equal within 20\%}  % experiments/FIGS/data/v2/index.json :: fine-tuning qg N=10000: R30 of each L188 run over the from-scratch mean, SD over runs and standard error of that mean in quadrature :: sha256 c6c4138ee5c86423 :: prose twin: two significant figures, direction in words
\newcommand{\MlstFtQgRthirtyResEFour}{76.4\,\ensuremath{\pm}\,0.6}  % experiments/FIGS/data/v2/index.json :: fine-tuning qg N=10000: R30, L162 :: sha256 c6c4138ee5c86423
\newcommand{\MlstFtQgRthirtyResOverScratchEFour}{1.15\,\ensuremath{\pm}\,0.09}  % experiments/FIGS/data/v2/index.json :: fine-tuning qg N=10000: R30 of each L162 run over the from-scratch mean, SD over runs and standard error of that mean in quadrature :: sha256 c6c4138ee5c86423
\newcommand{\MlstFtQgRthirtyResOverScratchEFourPr}{15\% more rejection}  % experiments/FIGS/data/v2/index.json :: fine-tuning qg N=10000: R30 of each L162 run over the from-scratch mean, SD over runs and standard error of that mean in quadrature :: sha256 c6c4138ee5c86423 :: prose twin: two significant figures, direction in words
\newcommand{\MlstFtQgRthirtyMidEFour}{74\,\ensuremath{\pm}\,5}  % experiments/FIGS/data/v2/index.json :: fine-tuning qg N=10000: R30, R42_Q1 :: sha256 c6c4138ee5c86423
\newcommand{\MlstFtQgRthirtyMidOverScratchEFour}{1.12\,\ensuremath{\pm}\,0.11}  % experiments/FIGS/data/v2/index.json :: fine-tuning qg N=10000: R30 of each R42_Q1 run over the from-scratch mean, SD over runs and standard error of that mean in quadrature :: sha256 c6c4138ee5c86423
\newcommand{\MlstFtQgRthirtyMidOverScratchEFourPr}{equal within 20\%}  % experiments/FIGS/data/v2/index.json :: fine-tuning qg N=10000: R30 of each R42_Q1 run over the from-scratch mean, SD over runs and standard error of that mean in quadrature :: sha256 c6c4138ee5c86423 :: prose twin: two significant figures, direction in words
\newcommand{\MlstFtQgRthirtyCoarseEFour}{80.3\,\ensuremath{\pm}\,1.1}  % experiments/FIGS/data/v2/index.json :: fine-tuning qg N=10000: R30, R16_Q1 :: sha256 c6c4138ee5c86423
\newcommand{\MlstFtQgRthirtyCoarseOverScratchEFour}{1.21\,\ensuremath{\pm}\,0.10}  % experiments/FIGS/data/v2/index.json :: fine-tuning qg N=10000: R30 of each R16_Q1 run over the from-scratch mean, SD over runs and standard error of that mean in quadrature :: sha256 c6c4138ee5c86423
\newcommand{\MlstFtQgRthirtyCoarseOverScratchEFourPr}{21\% more rejection}  % experiments/FIGS/data/v2/index.json :: fine-tuning qg N=10000: R30 of each R16_Q1 run over the from-scratch mean, SD over runs and standard error of that mean in quadrature :: sha256 c6c4138ee5c86423 :: prose twin: two significant figures, direction in words
\newcommand{\MlstFtQgRthirtyResMassEFour}{80.1 and 72.6}  % experiments/FIGS/data/v2/index.json :: fine-tuning qg N=10000: R30, L162_MASS :: sha256 c6c4138ee5c86423
\newcommand{\MlstFtQgRthirtyResMassOverScratchEFour}{1.21 and 1.09}  % experiments/FIGS/data/v2/index.json :: fine-tuning qg N=10000: R30 of each L162_MASS run over the from-scratch mean, SD over runs and standard error of that mean in quadrature :: sha256 c6c4138ee5c86423
\newcommand{\MlstFtQgRthirtyResMassOverScratchEFourPr}{at least 9.3\% more rejection}  % experiments/FIGS/data/v2/index.json :: fine-tuning qg N=10000: R30 of each L162_MASS run over the from-scratch mean, SD over runs and standard error of that mean in quadrature :: sha256 c6c4138ee5c86423 :: prose twin: two significant figures, direction in words
\newcommand{\MlstFtQgRthirtyResMassOverScratchEFourPrMax}{up to 21\% more rejection}  % experiments/FIGS/data/v2/index.json :: fine-tuning qg N=10000: R30 of each L162_MASS run over the from-scratch mean, SD over runs and standard error of that mean in quadrature :: sha256 c6c4138ee5c86423 :: prose twin of the run farther from 1
\newcommand{\MlstFtQgRthirtyCoarseMassEFour}{79.5 and 77.9}  % experiments/FIGS/data/v2/index.json :: fine-tuning qg N=10000: R30, R16_Q1_MASS :: sha256 c6c4138ee5c86423
\newcommand{\MlstFtQgRthirtyCoarseMassOverScratchEFour}{1.198 and 1.173}  % experiments/FIGS/data/v2/index.json :: fine-tuning qg N=10000: R30 of each R16_Q1_MASS run over the from-scratch mean, SD over runs and standard error of that mean in quadrature :: sha256 c6c4138ee5c86423
\newcommand{\MlstFtQgRthirtyCoarseMassOverScratchEFourPr}{at least 17\% more rejection}  % experiments/FIGS/data/v2/index.json :: fine-tuning qg N=10000: R30 of each R16_Q1_MASS run over the from-scratch mean, SD over runs and standard error of that mean in quadrature :: sha256 c6c4138ee5c86423 :: prose twin: two significant figures, direction in words
\newcommand{\MlstFtQgRthirtyCoarseMassOverScratchEFourPrMax}{up to 20\% more rejection}  % experiments/FIGS/data/v2/index.json :: fine-tuning qg N=10000: R30 of each R16_Q1_MASS run over the from-scratch mean, SD over runs and standard error of that mean in quadrature :: sha256 c6c4138ee5c86423 :: prose twin of the run farther from 1
\newcommand{\MlstFtQgRthirtyResOverFullEFour}{1.02\,\ensuremath{\pm}\,0.10}  % experiments/FIGS/data/v2/index.json :: fine-tuning qg N=10000: R30 of L162 over L188, run by run :: sha256 c6c4138ee5c86423
\newcommand{\MlstFtQgRthirtyResOverFullEFourPr}{equal within 10\%}  % experiments/FIGS/data/v2/index.json :: fine-tuning qg N=10000: R30 of L162 over L188, run by run :: sha256 c6c4138ee5c86423 :: prose twin: two significant figures, direction in words
\newcommand{\MlstFtQgRthirtyMidOverFullEFour}{0.99\,\ensuremath{\pm}\,0.14}  % experiments/FIGS/data/v2/index.json :: fine-tuning qg N=10000: R30 of R42_Q1 over L188, run by run :: sha256 c6c4138ee5c86423
\newcommand{\MlstFtQgRthirtyMidOverFullEFourPr}{equal within 20\%}  % experiments/FIGS/data/v2/index.json :: fine-tuning qg N=10000: R30 of R42_Q1 over L188, run by run :: sha256 c6c4138ee5c86423 :: prose twin: two significant figures, direction in words
\newcommand{\MlstFtQgRthirtyCoarseOverFullEFour}{1.07\,\ensuremath{\pm}\,0.13}  % experiments/FIGS/data/v2/index.json :: fine-tuning qg N=10000: R30 of R16_Q1 over L188, run by run :: sha256 c6c4138ee5c86423
\newcommand{\MlstFtQgRthirtyCoarseOverFullEFourPr}{equal within 20\%}  % experiments/FIGS/data/v2/index.json :: fine-tuning qg N=10000: R30 of R16_Q1 over L188, run by run :: sha256 c6c4138ee5c86423 :: prose twin: two significant figures, direction in words
\newcommand{\MlstFtQgRthirtyResMassOverResEFour}{1.06 and 0.95}  % experiments/FIGS/data/v2/index.json :: fine-tuning qg N=10000: R30 of L162_MASS over L162, run by run :: sha256 c6c4138ee5c86423
\newcommand{\MlstFtQgRthirtyResMassOverResEFourPr}{equal within 6\%}  % experiments/FIGS/data/v2/index.json :: fine-tuning qg N=10000: R30 of L162_MASS over L162, run by run :: sha256 c6c4138ee5c86423 :: prose twin: two significant figures, direction in words
\newcommand{\MlstFtQgRthirtyResMassOverResEFourPrMax}{equal within 6\%}  % experiments/FIGS/data/v2/index.json :: fine-tuning qg N=10000: R30 of L162_MASS over L162, run by run :: sha256 c6c4138ee5c86423 :: prose twin of the run farther from 1
\newcommand{\MlstFtQgRthirtyCoarseMassOverCoarseEFour}{0.9769 and 0.9696}  % experiments/FIGS/data/v2/index.json :: fine-tuning qg N=10000: R30 of R16_Q1_MASS over R16_Q1, run by run :: sha256 c6c4138ee5c86423
\newcommand{\MlstFtQgRthirtyCoarseMassOverCoarseEFourPr}{equal within 4\%}  % experiments/FIGS/data/v2/index.json :: fine-tuning qg N=10000: R30 of R16_Q1_MASS over R16_Q1, run by run :: sha256 c6c4138ee5c86423 :: prose twin: two significant figures, direction in words
\newcommand{\MlstFtQgRthirtyCoarseMassOverCoarseEFourPrMax}{equal within 4\%}  % experiments/FIGS/data/v2/index.json :: fine-tuning qg N=10000: R30 of R16_Q1_MASS over R16_Q1, run by run :: sha256 c6c4138ee5c86423 :: prose twin of the run farther from 1
\newcommand{\MlstFtQgRthirtyRandOverFullEFour}{1.03\,\ensuremath{\pm}\,0.18}  % experiments/FIGS/data/v2/index.json :: fine-tuning qg N=10000: R30 of each random partition over the L188 run of its seed :: sha256 c6c4138ee5c86423
\newcommand{\MlstFtQgRthirtyRandOverFullEFourPr}{equal within 20\%}  % experiments/FIGS/data/v2/index.json :: fine-tuning qg N=10000: R30 of each random partition over the L188 run of its seed :: sha256 c6c4138ee5c86423 :: prose twin: two significant figures, direction in words
\newcommand{\MlstFtQgRthirtyRandOverCoarseEFour}{0.93\,\ensuremath{\pm}\,0.12}  % experiments/FIGS/data/v2/index.json :: fine-tuning qg N=10000: R30 of each random partition over the R16_Q1 run of its seed :: sha256 c6c4138ee5c86423
\newcommand{\MlstFtQgRthirtyRandOverCoarseEFourPr}{equal within 20\%}  % experiments/FIGS/data/v2/index.json :: fine-tuning qg N=10000: R30 of each random partition over the R16_Q1 run of its seed :: sha256 c6c4138ee5c86423 :: prose twin: two significant figures, direction in words
\newcommand{\MlstFtQgRthirtyScratchEFive}{88.4\,\ensuremath{\pm}\,0.8}  % experiments/FIGS/data/v2/index.json :: fine-tuning qg N=100000: R30 from scratch, every seed :: sha256 c6c4138ee5c86423
\newcommand{\MlstFtQgRthirtyFullEFive}{99.3\,\ensuremath{\pm}\,1.4}  % experiments/FIGS/data/v2/index.json :: fine-tuning qg N=100000: R30, L188 :: sha256 c6c4138ee5c86423
\newcommand{\MlstFtQgRthirtyFullOverScratchEFive}{1.124\,\ensuremath{\pm}\,0.017}  % experiments/FIGS/data/v2/index.json :: fine-tuning qg N=100000: R30 of each L188 run over the from-scratch mean, SD over runs and standard error of that mean in quadrature :: sha256 c6c4138ee5c86423
\newcommand{\MlstFtQgRthirtyFullOverScratchEFivePr}{12\% more rejection}  % experiments/FIGS/data/v2/index.json :: fine-tuning qg N=100000: R30 of each L188 run over the from-scratch mean, SD over runs and standard error of that mean in quadrature :: sha256 c6c4138ee5c86423 :: prose twin: two significant figures, direction in words
\newcommand{\MlstFtQgRthirtyResEFive}{97.9\,\ensuremath{\pm}\,3.3}  % experiments/FIGS/data/v2/index.json :: fine-tuning qg N=100000: R30, L162 :: sha256 c6c4138ee5c86423
\newcommand{\MlstFtQgRthirtyResOverScratchEFive}{1.11\,\ensuremath{\pm}\,0.04}  % experiments/FIGS/data/v2/index.json :: fine-tuning qg N=100000: R30 of each L162 run over the from-scratch mean, SD over runs and standard error of that mean in quadrature :: sha256 c6c4138ee5c86423
\newcommand{\MlstFtQgRthirtyResOverScratchEFivePr}{11\% more rejection}  % experiments/FIGS/data/v2/index.json :: fine-tuning qg N=100000: R30 of each L162 run over the from-scratch mean, SD over runs and standard error of that mean in quadrature :: sha256 c6c4138ee5c86423 :: prose twin: two significant figures, direction in words
\newcommand{\MlstFtQgRthirtyMidEFive}{96.1\,\ensuremath{\pm}\,1.8}  % experiments/FIGS/data/v2/index.json :: fine-tuning qg N=100000: R30, R42_Q1 :: sha256 c6c4138ee5c86423
\newcommand{\MlstFtQgRthirtyMidOverScratchEFive}{1.088\,\ensuremath{\pm}\,0.021}  % experiments/FIGS/data/v2/index.json :: fine-tuning qg N=100000: R30 of each R42_Q1 run over the from-scratch mean, SD over runs and standard error of that mean in quadrature :: sha256 c6c4138ee5c86423
\newcommand{\MlstFtQgRthirtyMidOverScratchEFivePr}{8.8\% more rejection}  % experiments/FIGS/data/v2/index.json :: fine-tuning qg N=100000: R30 of each R42_Q1 run over the from-scratch mean, SD over runs and standard error of that mean in quadrature :: sha256 c6c4138ee5c86423 :: prose twin: two significant figures, direction in words
\newcommand{\MlstFtQgRthirtyCoarseEFive}{96.4\,\ensuremath{\pm}\,1.0}  % experiments/FIGS/data/v2/index.json :: fine-tuning qg N=100000: R30, R16_Q1 :: sha256 c6c4138ee5c86423
\newcommand{\MlstFtQgRthirtyCoarseOverScratchEFive}{1.091\,\ensuremath{\pm}\,0.013}  % experiments/FIGS/data/v2/index.json :: fine-tuning qg N=100000: R30 of each R16_Q1 run over the from-scratch mean, SD over runs and standard error of that mean in quadrature :: sha256 c6c4138ee5c86423
\newcommand{\MlstFtQgRthirtyCoarseOverScratchEFivePr}{9.1\% more rejection}  % experiments/FIGS/data/v2/index.json :: fine-tuning qg N=100000: R30 of each R16_Q1 run over the from-scratch mean, SD over runs and standard error of that mean in quadrature :: sha256 c6c4138ee5c86423 :: prose twin: two significant figures, direction in words
\newcommand{\MlstFtQgRthirtyResMassEFive}{95.6 and 97.6}  % experiments/FIGS/data/v2/index.json :: fine-tuning qg N=100000: R30, L162_MASS :: sha256 c6c4138ee5c86423
\newcommand{\MlstFtQgRthirtyResMassOverScratchEFive}{1.082 and 1.104}  % experiments/FIGS/data/v2/index.json :: fine-tuning qg N=100000: R30 of each L162_MASS run over the from-scratch mean, SD over runs and standard error of that mean in quadrature :: sha256 c6c4138ee5c86423
\newcommand{\MlstFtQgRthirtyResMassOverScratchEFivePr}{at least 8.2\% more rejection}  % experiments/FIGS/data/v2/index.json :: fine-tuning qg N=100000: R30 of each L162_MASS run over the from-scratch mean, SD over runs and standard error of that mean in quadrature :: sha256 c6c4138ee5c86423 :: prose twin: two significant figures, direction in words
\newcommand{\MlstFtQgRthirtyResMassOverScratchEFivePrMax}{up to 10\% more rejection}  % experiments/FIGS/data/v2/index.json :: fine-tuning qg N=100000: R30 of each L162_MASS run over the from-scratch mean, SD over runs and standard error of that mean in quadrature :: sha256 c6c4138ee5c86423 :: prose twin of the run farther from 1
\newcommand{\MlstFtQgRthirtyCoarseMassEFive}{97.85 and 97.94}  % experiments/FIGS/data/v2/index.json :: fine-tuning qg N=100000: R30, R16_Q1_MASS :: sha256 c6c4138ee5c86423
\newcommand{\MlstFtQgRthirtyCoarseMassOverScratchEFive}{1.107 and 1.108}  % experiments/FIGS/data/v2/index.json :: fine-tuning qg N=100000: R30 of each R16_Q1_MASS run over the from-scratch mean, SD over runs and standard error of that mean in quadrature :: sha256 c6c4138ee5c86423
\newcommand{\MlstFtQgRthirtyCoarseMassOverScratchEFivePr}{at least 11\% more rejection}  % experiments/FIGS/data/v2/index.json :: fine-tuning qg N=100000: R30 of each R16_Q1_MASS run over the from-scratch mean, SD over runs and standard error of that mean in quadrature :: sha256 c6c4138ee5c86423 :: prose twin: two significant figures, direction in words
\newcommand{\MlstFtQgRthirtyCoarseMassOverScratchEFivePrMax}{up to 11\% more rejection}  % experiments/FIGS/data/v2/index.json :: fine-tuning qg N=100000: R30 of each R16_Q1_MASS run over the from-scratch mean, SD over runs and standard error of that mean in quadrature :: sha256 c6c4138ee5c86423 :: prose twin of the run farther from 1
\newcommand{\MlstFtQgRthirtyResOverFullEFive}{0.985\,\ensuremath{\pm}\,0.020}  % experiments/FIGS/data/v2/index.json :: fine-tuning qg N=100000: R30 of L162 over L188, run by run :: sha256 c6c4138ee5c86423
\newcommand{\MlstFtQgRthirtyResOverFullEFivePr}{equal within 2\%}  % experiments/FIGS/data/v2/index.json :: fine-tuning qg N=100000: R30 of L162 over L188, run by run :: sha256 c6c4138ee5c86423 :: prose twin: two significant figures, direction in words
\newcommand{\MlstFtQgRthirtyMidOverFullEFive}{0.969\,\ensuremath{\pm}\,0.030}  % experiments/FIGS/data/v2/index.json :: fine-tuning qg N=100000: R30 of R42_Q1 over L188, run by run :: sha256 c6c4138ee5c86423
\newcommand{\MlstFtQgRthirtyMidOverFullEFivePr}{equal within 4\%}  % experiments/FIGS/data/v2/index.json :: fine-tuning qg N=100000: R30 of R42_Q1 over L188, run by run :: sha256 c6c4138ee5c86423 :: prose twin: two significant figures, direction in words
\newcommand{\MlstFtQgRthirtyCoarseOverFullEFive}{0.972\,\ensuremath{\pm}\,0.020}  % experiments/FIGS/data/v2/index.json :: fine-tuning qg N=100000: R30 of R16_Q1 over L188, run by run :: sha256 c6c4138ee5c86423
\newcommand{\MlstFtQgRthirtyCoarseOverFullEFivePr}{equal within 3\%}  % experiments/FIGS/data/v2/index.json :: fine-tuning qg N=100000: R30 of R16_Q1 over L188, run by run :: sha256 c6c4138ee5c86423 :: prose twin: two significant figures, direction in words
\newcommand{\MlstFtQgRthirtyResMassOverResEFive}{0.942 and 1.027}  % experiments/FIGS/data/v2/index.json :: fine-tuning qg N=100000: R30 of L162_MASS over L162, run by run :: sha256 c6c4138ee5c86423
\newcommand{\MlstFtQgRthirtyResMassOverResEFivePr}{equal within 6\%}  % experiments/FIGS/data/v2/index.json :: fine-tuning qg N=100000: R30 of L162_MASS over L162, run by run :: sha256 c6c4138ee5c86423 :: prose twin: two significant figures, direction in words
\newcommand{\MlstFtQgRthirtyResMassOverResEFivePrMax}{equal within 6\%}  % experiments/FIGS/data/v2/index.json :: fine-tuning qg N=100000: R30 of L162_MASS over L162, run by run :: sha256 c6c4138ee5c86423 :: prose twin of the run farther from 1
\newcommand{\MlstFtQgRthirtyCoarseMassOverCoarseEFive}{1.018 and 1.004}  % experiments/FIGS/data/v2/index.json :: fine-tuning qg N=100000: R30 of R16_Q1_MASS over R16_Q1, run by run :: sha256 c6c4138ee5c86423
\newcommand{\MlstFtQgRthirtyCoarseMassOverCoarseEFivePr}{equal within 2\%}  % experiments/FIGS/data/v2/index.json :: fine-tuning qg N=100000: R30 of R16_Q1_MASS over R16_Q1, run by run :: sha256 c6c4138ee5c86423 :: prose twin: two significant figures, direction in words
\newcommand{\MlstFtQgRthirtyCoarseMassOverCoarseEFivePrMax}{equal within 2\%}  % experiments/FIGS/data/v2/index.json :: fine-tuning qg N=100000: R30 of R16_Q1_MASS over R16_Q1, run by run :: sha256 c6c4138ee5c86423 :: prose twin of the run farther from 1
\newcommand{\MlstFtQgRthirtyRandOverFullEFive}{0.986\,\ensuremath{\pm}\,0.025}  % experiments/FIGS/data/v2/index.json :: fine-tuning qg N=100000: R30 of each random partition over the L188 run of its seed :: sha256 c6c4138ee5c86423
\newcommand{\MlstFtQgRthirtyRandOverFullEFivePr}{equal within 3\%}  % experiments/FIGS/data/v2/index.json :: fine-tuning qg N=100000: R30 of each random partition over the L188 run of its seed :: sha256 c6c4138ee5c86423 :: prose twin: two significant figures, direction in words
\newcommand{\MlstFtQgRthirtyRandOverCoarseEFive}{1.009\,\ensuremath{\pm}\,0.020}  % experiments/FIGS/data/v2/index.json :: fine-tuning qg N=100000: R30 of each random partition over the R16_Q1 run of its seed :: sha256 c6c4138ee5c86423
\newcommand{\MlstFtQgRthirtyRandOverCoarseEFivePr}{equal within 3\%}  % experiments/FIGS/data/v2/index.json :: fine-tuning qg N=100000: R30 of each random partition over the R16_Q1 run of its seed :: sha256 c6c4138ee5c86423 :: prose twin: two significant figures, direction in words
\newcommand{\MlstFtQgRthirtyScratch}{104.0\,\ensuremath{\pm}\,1.0}  % experiments/FIGS/data/v2/index.json :: fine-tuning qg N=1600000: R30 from scratch, every seed :: sha256 c6c4138ee5c86423
\newcommand{\MlstFtQgRthirtyFull}{107.3\,\ensuremath{\pm}\,2.2}  % experiments/FIGS/data/v2/index.json :: fine-tuning qg N=1600000: R30, L188 :: sha256 c6c4138ee5c86423
\newcommand{\MlstFtQgRthirtyFullOverScratch}{1.031\,\ensuremath{\pm}\,0.022}  % experiments/FIGS/data/v2/index.json :: fine-tuning qg N=1600000: R30 of each L188 run over the from-scratch mean, SD over runs and standard error of that mean in quadrature :: sha256 c6c4138ee5c86423
\newcommand{\MlstFtQgRthirtyFullOverScratchPr}{equal within 4\%}  % experiments/FIGS/data/v2/index.json :: fine-tuning qg N=1600000: R30 of each L188 run over the from-scratch mean, SD over runs and standard error of that mean in quadrature :: sha256 c6c4138ee5c86423 :: prose twin: two significant figures, direction in words
\newcommand{\MlstFtQgRthirtyRes}{108.0\,\ensuremath{\pm}\,1.4}  % experiments/FIGS/data/v2/index.json :: fine-tuning qg N=1600000: R30, L162 :: sha256 c6c4138ee5c86423
\newcommand{\MlstFtQgRthirtyResOverScratch}{1.038\,\ensuremath{\pm}\,0.015}  % experiments/FIGS/data/v2/index.json :: fine-tuning qg N=1600000: R30 of each L162 run over the from-scratch mean, SD over runs and standard error of that mean in quadrature :: sha256 c6c4138ee5c86423
\newcommand{\MlstFtQgRthirtyResOverScratchPr}{3.8\% more rejection}  % experiments/FIGS/data/v2/index.json :: fine-tuning qg N=1600000: R30 of each L162 run over the from-scratch mean, SD over runs and standard error of that mean in quadrature :: sha256 c6c4138ee5c86423 :: prose twin: two significant figures, direction in words
\newcommand{\MlstFtQgRthirtyMid}{108.3\,\ensuremath{\pm}\,1.0}  % experiments/FIGS/data/v2/index.json :: fine-tuning qg N=1600000: R30, R42_Q1 :: sha256 c6c4138ee5c86423
\newcommand{\MlstFtQgRthirtyMidOverScratch}{1.041\,\ensuremath{\pm}\,0.012}  % experiments/FIGS/data/v2/index.json :: fine-tuning qg N=1600000: R30 of each R42_Q1 run over the from-scratch mean, SD over runs and standard error of that mean in quadrature :: sha256 c6c4138ee5c86423
\newcommand{\MlstFtQgRthirtyMidOverScratchPr}{4.1\% more rejection}  % experiments/FIGS/data/v2/index.json :: fine-tuning qg N=1600000: R30 of each R42_Q1 run over the from-scratch mean, SD over runs and standard error of that mean in quadrature :: sha256 c6c4138ee5c86423 :: prose twin: two significant figures, direction in words
\newcommand{\MlstFtQgRthirtyCoarse}{106.5\,\ensuremath{\pm}\,0.8}  % experiments/FIGS/data/v2/index.json :: fine-tuning qg N=1600000: R30, R16_Q1 :: sha256 c6c4138ee5c86423
\newcommand{\MlstFtQgRthirtyCoarseOverScratch}{1.024\,\ensuremath{\pm}\,0.010}  % experiments/FIGS/data/v2/index.json :: fine-tuning qg N=1600000: R30 of each R16_Q1 run over the from-scratch mean, SD over runs and standard error of that mean in quadrature :: sha256 c6c4138ee5c86423
\newcommand{\MlstFtQgRthirtyCoarseOverScratchPr}{2.4\% more rejection}  % experiments/FIGS/data/v2/index.json :: fine-tuning qg N=1600000: R30 of each R16_Q1 run over the from-scratch mean, SD over runs and standard error of that mean in quadrature :: sha256 c6c4138ee5c86423 :: prose twin: two significant figures, direction in words
\newcommand{\MlstFtQgRthirtyResMass}{107.0 and 106.6}  % experiments/FIGS/data/v2/index.json :: fine-tuning qg N=1600000: R30, L162_MASS :: sha256 c6c4138ee5c86423
\newcommand{\MlstFtQgRthirtyResMassOverScratch}{1.028 and 1.025}  % experiments/FIGS/data/v2/index.json :: fine-tuning qg N=1600000: R30 of each L162_MASS run over the from-scratch mean, SD over runs and standard error of that mean in quadrature :: sha256 c6c4138ee5c86423
\newcommand{\MlstFtQgRthirtyResMassOverScratchPr}{at least 2.5\% more rejection}  % experiments/FIGS/data/v2/index.json :: fine-tuning qg N=1600000: R30 of each L162_MASS run over the from-scratch mean, SD over runs and standard error of that mean in quadrature :: sha256 c6c4138ee5c86423 :: prose twin: two significant figures, direction in words
\newcommand{\MlstFtQgRthirtyResMassOverScratchPrMax}{up to 2.8\% more rejection}  % experiments/FIGS/data/v2/index.json :: fine-tuning qg N=1600000: R30 of each L162_MASS run over the from-scratch mean, SD over runs and standard error of that mean in quadrature :: sha256 c6c4138ee5c86423 :: prose twin of the run farther from 1
\newcommand{\MlstFtQgRthirtyCoarseMass}{104.9 and 105.6}  % experiments/FIGS/data/v2/index.json :: fine-tuning qg N=1600000: R30, R16_Q1_MASS :: sha256 c6c4138ee5c86423
\newcommand{\MlstFtQgRthirtyCoarseMassOverScratch}{1.009 and 1.015}  % experiments/FIGS/data/v2/index.json :: fine-tuning qg N=1600000: R30 of each R16_Q1_MASS run over the from-scratch mean, SD over runs and standard error of that mean in quadrature :: sha256 c6c4138ee5c86423
\newcommand{\MlstFtQgRthirtyCoarseMassOverScratchPr}{at least 0.87\% more rejection}  % experiments/FIGS/data/v2/index.json :: fine-tuning qg N=1600000: R30 of each R16_Q1_MASS run over the from-scratch mean, SD over runs and standard error of that mean in quadrature :: sha256 c6c4138ee5c86423 :: prose twin: two significant figures, direction in words
\newcommand{\MlstFtQgRthirtyCoarseMassOverScratchPrMax}{up to 1.5\% more rejection}  % experiments/FIGS/data/v2/index.json :: fine-tuning qg N=1600000: R30 of each R16_Q1_MASS run over the from-scratch mean, SD over runs and standard error of that mean in quadrature :: sha256 c6c4138ee5c86423 :: prose twin of the run farther from 1
\newcommand{\MlstFtQgRthirtyResOverFull}{1.007\,\ensuremath{\pm}\,0.026}  % experiments/FIGS/data/v2/index.json :: fine-tuning qg N=1600000: R30 of L162 over L188, run by run :: sha256 c6c4138ee5c86423
\newcommand{\MlstFtQgRthirtyResOverFullPr}{equal within 3\%}  % experiments/FIGS/data/v2/index.json :: fine-tuning qg N=1600000: R30 of L162 over L188, run by run :: sha256 c6c4138ee5c86423 :: prose twin: two significant figures, direction in words
\newcommand{\MlstFtQgRthirtyMidOverFull}{1.010\,\ensuremath{\pm}\,0.029}  % experiments/FIGS/data/v2/index.json :: fine-tuning qg N=1600000: R30 of R42_Q1 over L188, run by run :: sha256 c6c4138ee5c86423
\newcommand{\MlstFtQgRthirtyMidOverFullPr}{equal within 3\%}  % experiments/FIGS/data/v2/index.json :: fine-tuning qg N=1600000: R30 of R42_Q1 over L188, run by run :: sha256 c6c4138ee5c86423 :: prose twin: two significant figures, direction in words
\newcommand{\MlstFtQgRthirtyCoarseOverFull}{0.993\,\ensuremath{\pm}\,0.026}  % experiments/FIGS/data/v2/index.json :: fine-tuning qg N=1600000: R30 of R16_Q1 over L188, run by run :: sha256 c6c4138ee5c86423
\newcommand{\MlstFtQgRthirtyCoarseOverFullPr}{equal within 3\%}  % experiments/FIGS/data/v2/index.json :: fine-tuning qg N=1600000: R30 of R16_Q1 over L188, run by run :: sha256 c6c4138ee5c86423 :: prose twin: two significant figures, direction in words
\newcommand{\MlstFtQgRthirtyResMassOverRes}{0.976 and 0.999}  % experiments/FIGS/data/v2/index.json :: fine-tuning qg N=1600000: R30 of L162_MASS over L162, run by run :: sha256 c6c4138ee5c86423
\newcommand{\MlstFtQgRthirtyResMassOverResPr}{equal within 3\%}  % experiments/FIGS/data/v2/index.json :: fine-tuning qg N=1600000: R30 of L162_MASS over L162, run by run :: sha256 c6c4138ee5c86423 :: prose twin: two significant figures, direction in words
\newcommand{\MlstFtQgRthirtyResMassOverResPrMax}{equal within 3\%}  % experiments/FIGS/data/v2/index.json :: fine-tuning qg N=1600000: R30 of L162_MASS over L162, run by run :: sha256 c6c4138ee5c86423 :: prose twin of the run farther from 1
\newcommand{\MlstFtQgRthirtyCoarseMassOverCoarse}{0.9916 and 0.9926}  % experiments/FIGS/data/v2/index.json :: fine-tuning qg N=1600000: R30 of R16_Q1_MASS over R16_Q1, run by run :: sha256 c6c4138ee5c86423
\newcommand{\MlstFtQgRthirtyCoarseMassOverCoarsePr}{equal within 1\%}  % experiments/FIGS/data/v2/index.json :: fine-tuning qg N=1600000: R30 of R16_Q1_MASS over R16_Q1, run by run :: sha256 c6c4138ee5c86423 :: prose twin: two significant figures, direction in words
\newcommand{\MlstFtQgRthirtyCoarseMassOverCoarsePrMax}{equal within 1\%}  % experiments/FIGS/data/v2/index.json :: fine-tuning qg N=1600000: R30 of R16_Q1_MASS over R16_Q1, run by run :: sha256 c6c4138ee5c86423 :: prose twin of the run farther from 1
\newcommand{\MlstFtQgRthirtyRandOverFull}{0.985\,\ensuremath{\pm}\,0.011}  % experiments/FIGS/data/v2/index.json :: fine-tuning qg N=1600000: R30 of each random partition over the L188 run of its seed :: sha256 c6c4138ee5c86423
\newcommand{\MlstFtQgRthirtyRandOverFullPr}{equal within 2\%}  % experiments/FIGS/data/v2/index.json :: fine-tuning qg N=1600000: R30 of each random partition over the L188 run of its seed :: sha256 c6c4138ee5c86423 :: prose twin: two significant figures, direction in words
\newcommand{\MlstFtQgRthirtyRandOverCoarse}{1.008\,\ensuremath{\pm}\,0.009}  % experiments/FIGS/data/v2/index.json :: fine-tuning qg N=1600000: R30 of each random partition over the R16_Q1 run of its seed :: sha256 c6c4138ee5c86423
\newcommand{\MlstFtQgRthirtyRandOverCoarsePr}{equal within 1\%}  % experiments/FIGS/data/v2/index.json :: fine-tuning qg N=1600000: R30 of each random partition over the R16_Q1 run of its seed :: sha256 c6c4138ee5c86423 :: prose twin: two significant figures, direction in words
\newcommand{\MlstFtQgRthirtyFullOverPart}{0.994\,\ensuremath{\pm}\,0.020}  % experiments/FIGS/data/v2/index.json :: fine-tuning qg N=1600000: R30 of each 188-class run over ParT-f.t.'s published value :: sha256 c6c4138ee5c86423
\newcommand{\MlstFtQgRthirtyFullOverPartPr}{equal within 3\%}  % experiments/FIGS/data/v2/index.json :: fine-tuning qg N=1600000: R30 of each 188-class run over ParT-f.t.'s published value :: sha256 c6c4138ee5c86423 :: prose twin: two significant figures, direction in words
\newcommand{\MlstFtQgAucFull}{0.9148\,\ensuremath{\pm}\,0.0001}  % experiments/FIGS/data/v2/index.json :: fine-tuning qg N=1600000: AUC of the 188-class runs :: sha256 c6c4138ee5c86423
\newcommand{\MlstFtQgAucScratch}{0.9121\,\ensuremath{\pm}\,0.0002}  % experiments/FIGS/data/v2/index.json :: fine-tuning qg N=1600000: AUC from scratch :: sha256 c6c4138ee5c86423
\newcommand{\MlstFtQgDataEquivJets}{\ensuremath{<}1{,}000}  % experiments/FIGS/data/v2/index.json :: fine-tuning qg: from-scratch training jets whose R30 equals the 188-class mean at 1000 jets (interpolated linearly in R30 against log10 N) :: sha256 c6c4138ee5c86423
\newcommand{\MlstFtQgDataEquivJetsEFour}{26{,}000}  % experiments/FIGS/data/v2/index.json :: fine-tuning qg: from-scratch training jets whose R30 equals the 188-class mean at 10000 jets (interpolated linearly in R30 against log10 N) :: sha256 c6c4138ee5c86423
\newcommand{\MlstFtQgHerwigRthirtyScratchEThree}{27.0\,\ensuremath{\pm}\,1.3}  % experiments/FIGS/data/v2/index.json :: fine-tuning qg_herwig N=1000: R30 from scratch, every seed :: sha256 c6c4138ee5c86423
\newcommand{\MlstFtQgHerwigRthirtyFullEThree}{21\,\ensuremath{\pm}\,6}  % experiments/FIGS/data/v2/index.json :: fine-tuning qg_herwig N=1000: R30, L188 :: sha256 c6c4138ee5c86423
\newcommand{\MlstFtQgHerwigRthirtyFullOverScratchEThree}{0.77\,\ensuremath{\pm}\,0.22}  % experiments/FIGS/data/v2/index.json :: fine-tuning qg_herwig N=1000: R30 of each L188 run over the from-scratch mean, SD over runs and standard error of that mean in quadrature :: sha256 c6c4138ee5c86423
\newcommand{\MlstFtQgHerwigRthirtyFullOverScratchEThreePr}{equal within 30\%}  % experiments/FIGS/data/v2/index.json :: fine-tuning qg_herwig N=1000: R30 of each L188 run over the from-scratch mean, SD over runs and standard error of that mean in quadrature :: sha256 c6c4138ee5c86423 :: prose twin: two significant figures, direction in words
\newcommand{\MlstFtQgHerwigRthirtyResEThree}{26.2\,\ensuremath{\pm}\,0.9}  % experiments/FIGS/data/v2/index.json :: fine-tuning qg_herwig N=1000: R30, L162 :: sha256 c6c4138ee5c86423
\newcommand{\MlstFtQgHerwigRthirtyResOverScratchEThree}{0.97\,\ensuremath{\pm}\,0.04}  % experiments/FIGS/data/v2/index.json :: fine-tuning qg_herwig N=1000: R30 of each L162 run over the from-scratch mean, SD over runs and standard error of that mean in quadrature :: sha256 c6c4138ee5c86423
\newcommand{\MlstFtQgHerwigRthirtyResOverScratchEThreePr}{equal within 4\%}  % experiments/FIGS/data/v2/index.json :: fine-tuning qg_herwig N=1000: R30 of each L162 run over the from-scratch mean, SD over runs and standard error of that mean in quadrature :: sha256 c6c4138ee5c86423 :: prose twin: two significant figures, direction in words
\newcommand{\MlstFtQgHerwigRthirtyMidEThree}{26.9\,\ensuremath{\pm}\,2.7}  % experiments/FIGS/data/v2/index.json :: fine-tuning qg_herwig N=1000: R30, R42_Q1 :: sha256 c6c4138ee5c86423
\newcommand{\MlstFtQgHerwigRthirtyMidOverScratchEThree}{1.00\,\ensuremath{\pm}\,0.10}  % experiments/FIGS/data/v2/index.json :: fine-tuning qg_herwig N=1000: R30 of each R42_Q1 run over the from-scratch mean, SD over runs and standard error of that mean in quadrature :: sha256 c6c4138ee5c86423
\newcommand{\MlstFtQgHerwigRthirtyMidOverScratchEThreePr}{equal within 10\%}  % experiments/FIGS/data/v2/index.json :: fine-tuning qg_herwig N=1000: R30 of each R42_Q1 run over the from-scratch mean, SD over runs and standard error of that mean in quadrature :: sha256 c6c4138ee5c86423 :: prose twin: two significant figures, direction in words
\newcommand{\MlstFtQgHerwigRthirtyCoarseEThree}{25.2\,\ensuremath{\pm}\,2.5}  % experiments/FIGS/data/v2/index.json :: fine-tuning qg_herwig N=1000: R30, R16_Q1 :: sha256 c6c4138ee5c86423
\newcommand{\MlstFtQgHerwigRthirtyCoarseOverScratchEThree}{0.94\,\ensuremath{\pm}\,0.10}  % experiments/FIGS/data/v2/index.json :: fine-tuning qg_herwig N=1000: R30 of each R16_Q1 run over the from-scratch mean, SD over runs and standard error of that mean in quadrature :: sha256 c6c4138ee5c86423
\newcommand{\MlstFtQgHerwigRthirtyCoarseOverScratchEThreePr}{equal within 10\%}  % experiments/FIGS/data/v2/index.json :: fine-tuning qg_herwig N=1000: R30 of each R16_Q1 run over the from-scratch mean, SD over runs and standard error of that mean in quadrature :: sha256 c6c4138ee5c86423 :: prose twin: two significant figures, direction in words
\newcommand{\MlstFtQgHerwigRthirtyResMassEThree}{28.0 and 29.0}  % experiments/FIGS/data/v2/index.json :: fine-tuning qg_herwig N=1000: R30, L162_MASS :: sha256 c6c4138ee5c86423
\newcommand{\MlstFtQgHerwigRthirtyResMassOverScratchEThree}{1.038 and 1.076}  % experiments/FIGS/data/v2/index.json :: fine-tuning qg_herwig N=1000: R30 of each L162_MASS run over the from-scratch mean, SD over runs and standard error of that mean in quadrature :: sha256 c6c4138ee5c86423
\newcommand{\MlstFtQgHerwigRthirtyResMassOverScratchEThreePr}{at least 3.8\% more rejection}  % experiments/FIGS/data/v2/index.json :: fine-tuning qg_herwig N=1000: R30 of each L162_MASS run over the from-scratch mean, SD over runs and standard error of that mean in quadrature :: sha256 c6c4138ee5c86423 :: prose twin: two significant figures, direction in words
\newcommand{\MlstFtQgHerwigRthirtyResMassOverScratchEThreePrMax}{up to 7.6\% more rejection}  % experiments/FIGS/data/v2/index.json :: fine-tuning qg_herwig N=1000: R30 of each L162_MASS run over the from-scratch mean, SD over runs and standard error of that mean in quadrature :: sha256 c6c4138ee5c86423 :: prose twin of the run farther from 1
\newcommand{\MlstFtQgHerwigRthirtyCoarseMassEThree}{25.8 and 27.4}  % experiments/FIGS/data/v2/index.json :: fine-tuning qg_herwig N=1000: R30, R16_Q1_MASS :: sha256 c6c4138ee5c86423
\newcommand{\MlstFtQgHerwigRthirtyCoarseMassOverScratchEThree}{0.955 and 1.015}  % experiments/FIGS/data/v2/index.json :: fine-tuning qg_herwig N=1000: R30 of each R16_Q1_MASS run over the from-scratch mean, SD over runs and standard error of that mean in quadrature :: sha256 c6c4138ee5c86423
\newcommand{\MlstFtQgHerwigRthirtyCoarseMassOverScratchEThreePr}{equal within 5\%}  % experiments/FIGS/data/v2/index.json :: fine-tuning qg_herwig N=1000: R30 of each R16_Q1_MASS run over the from-scratch mean, SD over runs and standard error of that mean in quadrature :: sha256 c6c4138ee5c86423 :: prose twin: two significant figures, direction in words
\newcommand{\MlstFtQgHerwigRthirtyCoarseMassOverScratchEThreePrMax}{equal within 5\%}  % experiments/FIGS/data/v2/index.json :: fine-tuning qg_herwig N=1000: R30 of each R16_Q1_MASS run over the from-scratch mean, SD over runs and standard error of that mean in quadrature :: sha256 c6c4138ee5c86423 :: prose twin of the run farther from 1
\newcommand{\MlstFtQgHerwigRthirtyResOverFullEThree}{1.33\,\ensuremath{\pm}\,0.33}  % experiments/FIGS/data/v2/index.json :: fine-tuning qg_herwig N=1000: R30 of L162 over L188, run by run :: sha256 c6c4138ee5c86423
\newcommand{\MlstFtQgHerwigRthirtyResOverFullEThreePr}{equal within 40\%}  % experiments/FIGS/data/v2/index.json :: fine-tuning qg_herwig N=1000: R30 of L162 over L188, run by run :: sha256 c6c4138ee5c86423 :: prose twin: two significant figures, direction in words
\newcommand{\MlstFtQgHerwigRthirtyMidOverFullEThree}{1.36\,\ensuremath{\pm}\,0.31}  % experiments/FIGS/data/v2/index.json :: fine-tuning qg_herwig N=1000: R30 of R42_Q1 over L188, run by run :: sha256 c6c4138ee5c86423
\newcommand{\MlstFtQgHerwigRthirtyMidOverFullEThreePr}{equal within 40\%}  % experiments/FIGS/data/v2/index.json :: fine-tuning qg_herwig N=1000: R30 of R42_Q1 over L188, run by run :: sha256 c6c4138ee5c86423 :: prose twin: two significant figures, direction in words
\newcommand{\MlstFtQgHerwigRthirtyCoarseOverFullEThree}{1.3\,\ensuremath{\pm}\,0.4}  % experiments/FIGS/data/v2/index.json :: fine-tuning qg_herwig N=1000: R30 of R16_Q1 over L188, run by run :: sha256 c6c4138ee5c86423
\newcommand{\MlstFtQgHerwigRthirtyCoarseOverFullEThreePr}{equal within 50\%}  % experiments/FIGS/data/v2/index.json :: fine-tuning qg_herwig N=1000: R30 of R16_Q1 over L188, run by run :: sha256 c6c4138ee5c86423 :: prose twin: two significant figures, direction in words
\newcommand{\MlstFtQgHerwigRthirtyResMassOverResEThree}{1.062 and 1.073}  % experiments/FIGS/data/v2/index.json :: fine-tuning qg_herwig N=1000: R30 of L162_MASS over L162, run by run :: sha256 c6c4138ee5c86423
\newcommand{\MlstFtQgHerwigRthirtyResMassOverResEThreePr}{equal within 8\%}  % experiments/FIGS/data/v2/index.json :: fine-tuning qg_herwig N=1000: R30 of L162_MASS over L162, run by run :: sha256 c6c4138ee5c86423 :: prose twin: two significant figures, direction in words
\newcommand{\MlstFtQgHerwigRthirtyResMassOverResEThreePrMax}{equal within 8\%}  % experiments/FIGS/data/v2/index.json :: fine-tuning qg_herwig N=1000: R30 of L162_MASS over L162, run by run :: sha256 c6c4138ee5c86423 :: prose twin of the run farther from 1
\newcommand{\MlstFtQgHerwigRthirtyCoarseMassOverCoarseEThree}{0.93 and 1.20}  % experiments/FIGS/data/v2/index.json :: fine-tuning qg_herwig N=1000: R30 of R16_Q1_MASS over R16_Q1, run by run :: sha256 c6c4138ee5c86423
\newcommand{\MlstFtQgHerwigRthirtyCoarseMassOverCoarseEThreePr}{equal within 30\%}  % experiments/FIGS/data/v2/index.json :: fine-tuning qg_herwig N=1000: R30 of R16_Q1_MASS over R16_Q1, run by run :: sha256 c6c4138ee5c86423 :: prose twin: two significant figures, direction in words
\newcommand{\MlstFtQgHerwigRthirtyCoarseMassOverCoarseEThreePrMax}{equal within 30\%}  % experiments/FIGS/data/v2/index.json :: fine-tuning qg_herwig N=1000: R30 of R16_Q1_MASS over R16_Q1, run by run :: sha256 c6c4138ee5c86423 :: prose twin of the run farther from 1
\newcommand{\MlstFtQgHerwigRthirtyRandOverFullEThree}{0.99\,\ensuremath{\pm}\,0.20}  % experiments/FIGS/data/v2/index.json :: fine-tuning qg_herwig N=1000: R30 of each random partition over the L188 run of its seed :: sha256 c6c4138ee5c86423
\newcommand{\MlstFtQgHerwigRthirtyRandOverFullEThreePr}{equal within 30\%}  % experiments/FIGS/data/v2/index.json :: fine-tuning qg_herwig N=1000: R30 of each random partition over the L188 run of its seed :: sha256 c6c4138ee5c86423 :: prose twin: two significant figures, direction in words
\newcommand{\MlstFtQgHerwigRthirtyRandOverCoarseEThree}{0.95\,\ensuremath{\pm}\,0.13}  % experiments/FIGS/data/v2/index.json :: fine-tuning qg_herwig N=1000: R30 of each random partition over the R16_Q1 run of its seed :: sha256 c6c4138ee5c86423
\newcommand{\MlstFtQgHerwigRthirtyRandOverCoarseEThreePr}{equal within 20\%}  % experiments/FIGS/data/v2/index.json :: fine-tuning qg_herwig N=1000: R30 of each random partition over the R16_Q1 run of its seed :: sha256 c6c4138ee5c86423 :: prose twin: two significant figures, direction in words
\newcommand{\MlstFtQgHerwigRthirtyScratchEFour}{31\,\ensuremath{\pm}\,7}  % experiments/FIGS/data/v2/index.json :: fine-tuning qg_herwig N=10000: R30 from scratch, every seed :: sha256 c6c4138ee5c86423
\newcommand{\MlstFtQgHerwigRthirtyFullEFour}{27.9\,\ensuremath{\pm}\,2.8}  % experiments/FIGS/data/v2/index.json :: fine-tuning qg_herwig N=10000: R30, L188 :: sha256 c6c4138ee5c86423
\newcommand{\MlstFtQgHerwigRthirtyFullOverScratchEFour}{0.89\,\ensuremath{\pm}\,0.14}  % experiments/FIGS/data/v2/index.json :: fine-tuning qg_herwig N=10000: R30 of each L188 run over the from-scratch mean, SD over runs and standard error of that mean in quadrature :: sha256 c6c4138ee5c86423
\newcommand{\MlstFtQgHerwigRthirtyFullOverScratchEFourPr}{equal within 20\%}  % experiments/FIGS/data/v2/index.json :: fine-tuning qg_herwig N=10000: R30 of each L188 run over the from-scratch mean, SD over runs and standard error of that mean in quadrature :: sha256 c6c4138ee5c86423 :: prose twin: two significant figures, direction in words
\newcommand{\MlstFtQgHerwigRthirtyResEFour}{27.4\,\ensuremath{\pm}\,0.5}  % experiments/FIGS/data/v2/index.json :: fine-tuning qg_herwig N=10000: R30, L162 :: sha256 c6c4138ee5c86423
\newcommand{\MlstFtQgHerwigRthirtyResOverScratchEFour}{0.87\,\ensuremath{\pm}\,0.11}  % experiments/FIGS/data/v2/index.json :: fine-tuning qg_herwig N=10000: R30 of each L162 run over the from-scratch mean, SD over runs and standard error of that mean in quadrature :: sha256 c6c4138ee5c86423
\newcommand{\MlstFtQgHerwigRthirtyResOverScratchEFourPr}{13\% less rejection}  % experiments/FIGS/data/v2/index.json :: fine-tuning qg_herwig N=10000: R30 of each L162 run over the from-scratch mean, SD over runs and standard error of that mean in quadrature :: sha256 c6c4138ee5c86423 :: prose twin: two significant figures, direction in words
\newcommand{\MlstFtQgHerwigRthirtyMidEFour}{28.4\,\ensuremath{\pm}\,2.1}  % experiments/FIGS/data/v2/index.json :: fine-tuning qg_herwig N=10000: R30, R42_Q1 :: sha256 c6c4138ee5c86423
\newcommand{\MlstFtQgHerwigRthirtyMidOverScratchEFour}{0.91\,\ensuremath{\pm}\,0.13}  % experiments/FIGS/data/v2/index.json :: fine-tuning qg_herwig N=10000: R30 of each R42_Q1 run over the from-scratch mean, SD over runs and standard error of that mean in quadrature :: sha256 c6c4138ee5c86423
\newcommand{\MlstFtQgHerwigRthirtyMidOverScratchEFourPr}{equal within 10\%}  % experiments/FIGS/data/v2/index.json :: fine-tuning qg_herwig N=10000: R30 of each R42_Q1 run over the from-scratch mean, SD over runs and standard error of that mean in quadrature :: sha256 c6c4138ee5c86423 :: prose twin: two significant figures, direction in words
\newcommand{\MlstFtQgHerwigRthirtyCoarseEFour}{32.7\,\ensuremath{\pm}\,1.8}  % experiments/FIGS/data/v2/index.json :: fine-tuning qg_herwig N=10000: R30, R16_Q1 :: sha256 c6c4138ee5c86423
\newcommand{\MlstFtQgHerwigRthirtyCoarseOverScratchEFour}{1.05\,\ensuremath{\pm}\,0.14}  % experiments/FIGS/data/v2/index.json :: fine-tuning qg_herwig N=10000: R30 of each R16_Q1 run over the from-scratch mean, SD over runs and standard error of that mean in quadrature :: sha256 c6c4138ee5c86423
\newcommand{\MlstFtQgHerwigRthirtyCoarseOverScratchEFourPr}{equal within 6\%}  % experiments/FIGS/data/v2/index.json :: fine-tuning qg_herwig N=10000: R30 of each R16_Q1 run over the from-scratch mean, SD over runs and standard error of that mean in quadrature :: sha256 c6c4138ee5c86423 :: prose twin: two significant figures, direction in words
\newcommand{\MlstFtQgHerwigRthirtyResMassEFour}{34.3 and 28.5}  % experiments/FIGS/data/v2/index.json :: fine-tuning qg_herwig N=10000: R30, L162_MASS :: sha256 c6c4138ee5c86423
\newcommand{\MlstFtQgHerwigRthirtyResMassOverScratchEFour}{1.09 and 0.91}  % experiments/FIGS/data/v2/index.json :: fine-tuning qg_herwig N=10000: R30 of each L162_MASS run over the from-scratch mean, SD over runs and standard error of that mean in quadrature :: sha256 c6c4138ee5c86423
\newcommand{\MlstFtQgHerwigRthirtyResMassOverScratchEFourPr}{equal within 10\%}  % experiments/FIGS/data/v2/index.json :: fine-tuning qg_herwig N=10000: R30 of each L162_MASS run over the from-scratch mean, SD over runs and standard error of that mean in quadrature :: sha256 c6c4138ee5c86423 :: prose twin: two significant figures, direction in words
\newcommand{\MlstFtQgHerwigRthirtyResMassOverScratchEFourPrMax}{equal within 10\%}  % experiments/FIGS/data/v2/index.json :: fine-tuning qg_herwig N=10000: R30 of each L162_MASS run over the from-scratch mean, SD over runs and standard error of that mean in quadrature :: sha256 c6c4138ee5c86423 :: prose twin of the run farther from 1
\newcommand{\MlstFtQgHerwigRthirtyCoarseMassEFour}{34.4 and 31.2}  % experiments/FIGS/data/v2/index.json :: fine-tuning qg_herwig N=10000: R30, R16_Q1_MASS :: sha256 c6c4138ee5c86423
\newcommand{\MlstFtQgHerwigRthirtyCoarseMassOverScratchEFour}{1.10 and 1.00}  % experiments/FIGS/data/v2/index.json :: fine-tuning qg_herwig N=10000: R30 of each R16_Q1_MASS run over the from-scratch mean, SD over runs and standard error of that mean in quadrature :: sha256 c6c4138ee5c86423
\newcommand{\MlstFtQgHerwigRthirtyCoarseMassOverScratchEFourPr}{equal within 20\%}  % experiments/FIGS/data/v2/index.json :: fine-tuning qg_herwig N=10000: R30 of each R16_Q1_MASS run over the from-scratch mean, SD over runs and standard error of that mean in quadrature :: sha256 c6c4138ee5c86423 :: prose twin: two significant figures, direction in words
\newcommand{\MlstFtQgHerwigRthirtyCoarseMassOverScratchEFourPrMax}{equal within 20\%}  % experiments/FIGS/data/v2/index.json :: fine-tuning qg_herwig N=10000: R30 of each R16_Q1_MASS run over the from-scratch mean, SD over runs and standard error of that mean in quadrature :: sha256 c6c4138ee5c86423 :: prose twin of the run farther from 1
\newcommand{\MlstFtQgHerwigRthirtyResOverFullEFour}{0.99\,\ensuremath{\pm}\,0.12}  % experiments/FIGS/data/v2/index.json :: fine-tuning qg_herwig N=10000: R30 of L162 over L188, run by run :: sha256 c6c4138ee5c86423
\newcommand{\MlstFtQgHerwigRthirtyResOverFullEFourPr}{equal within 20\%}  % experiments/FIGS/data/v2/index.json :: fine-tuning qg_herwig N=10000: R30 of L162 over L188, run by run :: sha256 c6c4138ee5c86423 :: prose twin: two significant figures, direction in words
\newcommand{\MlstFtQgHerwigRthirtyMidOverFullEFour}{1.03\,\ensuremath{\pm}\,0.18}  % experiments/FIGS/data/v2/index.json :: fine-tuning qg_herwig N=10000: R30 of R42_Q1 over L188, run by run :: sha256 c6c4138ee5c86423
\newcommand{\MlstFtQgHerwigRthirtyMidOverFullEFourPr}{equal within 20\%}  % experiments/FIGS/data/v2/index.json :: fine-tuning qg_herwig N=10000: R30 of R42_Q1 over L188, run by run :: sha256 c6c4138ee5c86423 :: prose twin: two significant figures, direction in words
\newcommand{\MlstFtQgHerwigRthirtyCoarseOverFullEFour}{1.18\,\ensuremath{\pm}\,0.08}  % experiments/FIGS/data/v2/index.json :: fine-tuning qg_herwig N=10000: R30 of R16_Q1 over L188, run by run :: sha256 c6c4138ee5c86423
\newcommand{\MlstFtQgHerwigRthirtyCoarseOverFullEFourPr}{18\% more rejection}  % experiments/FIGS/data/v2/index.json :: fine-tuning qg_herwig N=10000: R30 of R16_Q1 over L188, run by run :: sha256 c6c4138ee5c86423 :: prose twin: two significant figures, direction in words
\newcommand{\MlstFtQgHerwigRthirtyResMassOverResEFour}{1.23 and 1.03}  % experiments/FIGS/data/v2/index.json :: fine-tuning qg_herwig N=10000: R30 of L162_MASS over L162, run by run :: sha256 c6c4138ee5c86423
\newcommand{\MlstFtQgHerwigRthirtyResMassOverResEFourPr}{equal within 30\%}  % experiments/FIGS/data/v2/index.json :: fine-tuning qg_herwig N=10000: R30 of L162_MASS over L162, run by run :: sha256 c6c4138ee5c86423 :: prose twin: two significant figures, direction in words
\newcommand{\MlstFtQgHerwigRthirtyResMassOverResEFourPrMax}{equal within 30\%}  % experiments/FIGS/data/v2/index.json :: fine-tuning qg_herwig N=10000: R30 of L162_MASS over L162, run by run :: sha256 c6c4138ee5c86423 :: prose twin of the run farther from 1
\newcommand{\MlstFtQgHerwigRthirtyCoarseMassOverCoarseEFour}{1.12 and 0.91}  % experiments/FIGS/data/v2/index.json :: fine-tuning qg_herwig N=10000: R30 of R16_Q1_MASS over R16_Q1, run by run :: sha256 c6c4138ee5c86423
\newcommand{\MlstFtQgHerwigRthirtyCoarseMassOverCoarseEFourPr}{equal within 20\%}  % experiments/FIGS/data/v2/index.json :: fine-tuning qg_herwig N=10000: R30 of R16_Q1_MASS over R16_Q1, run by run :: sha256 c6c4138ee5c86423 :: prose twin: two significant figures, direction in words
\newcommand{\MlstFtQgHerwigRthirtyCoarseMassOverCoarseEFourPrMax}{equal within 20\%}  % experiments/FIGS/data/v2/index.json :: fine-tuning qg_herwig N=10000: R30 of R16_Q1_MASS over R16_Q1, run by run :: sha256 c6c4138ee5c86423 :: prose twin of the run farther from 1
\newcommand{\MlstFtQgHerwigRthirtyRandOverFullEFour}{1.13\,\ensuremath{\pm}\,0.18}  % experiments/FIGS/data/v2/index.json :: fine-tuning qg_herwig N=10000: R30 of each random partition over the L188 run of its seed :: sha256 c6c4138ee5c86423
\newcommand{\MlstFtQgHerwigRthirtyRandOverFullEFourPr}{equal within 20\%}  % experiments/FIGS/data/v2/index.json :: fine-tuning qg_herwig N=10000: R30 of each random partition over the L188 run of its seed :: sha256 c6c4138ee5c86423 :: prose twin: two significant figures, direction in words
\newcommand{\MlstFtQgHerwigRthirtyRandOverCoarseEFour}{0.93\,\ensuremath{\pm}\,0.14}  % experiments/FIGS/data/v2/index.json :: fine-tuning qg_herwig N=10000: R30 of each random partition over the R16_Q1 run of its seed :: sha256 c6c4138ee5c86423
\newcommand{\MlstFtQgHerwigRthirtyRandOverCoarseEFourPr}{equal within 20\%}  % experiments/FIGS/data/v2/index.json :: fine-tuning qg_herwig N=10000: R30 of each random partition over the R16_Q1 run of its seed :: sha256 c6c4138ee5c86423 :: prose twin: two significant figures, direction in words
\newcommand{\MlstFtQgHerwigRthirtyScratchEFive}{36.7\,\ensuremath{\pm}\,0.9}  % experiments/FIGS/data/v2/index.json :: fine-tuning qg_herwig N=100000: R30 from scratch, every seed :: sha256 c6c4138ee5c86423
\newcommand{\MlstFtQgHerwigRthirtyFullEFive}{34.3\,\ensuremath{\pm}\,1.0}  % experiments/FIGS/data/v2/index.json :: fine-tuning qg_herwig N=100000: R30, L188 :: sha256 c6c4138ee5c86423
\newcommand{\MlstFtQgHerwigRthirtyFullOverScratchEFive}{0.935\,\ensuremath{\pm}\,0.031}  % experiments/FIGS/data/v2/index.json :: fine-tuning qg_herwig N=100000: R30 of each L188 run over the from-scratch mean, SD over runs and standard error of that mean in quadrature :: sha256 c6c4138ee5c86423
\newcommand{\MlstFtQgHerwigRthirtyFullOverScratchEFivePr}{6.5\% less rejection}  % experiments/FIGS/data/v2/index.json :: fine-tuning qg_herwig N=100000: R30 of each L188 run over the from-scratch mean, SD over runs and standard error of that mean in quadrature :: sha256 c6c4138ee5c86423 :: prose twin: two significant figures, direction in words
\newcommand{\MlstFtQgHerwigRthirtyResEFive}{35.5\,\ensuremath{\pm}\,0.5}  % experiments/FIGS/data/v2/index.json :: fine-tuning qg_herwig N=100000: R30, L162 :: sha256 c6c4138ee5c86423
\newcommand{\MlstFtQgHerwigRthirtyResOverScratchEFive}{0.966\,\ensuremath{\pm}\,0.019}  % experiments/FIGS/data/v2/index.json :: fine-tuning qg_herwig N=100000: R30 of each L162 run over the from-scratch mean, SD over runs and standard error of that mean in quadrature :: sha256 c6c4138ee5c86423
\newcommand{\MlstFtQgHerwigRthirtyResOverScratchEFivePr}{3.4\% less rejection}  % experiments/FIGS/data/v2/index.json :: fine-tuning qg_herwig N=100000: R30 of each L162 run over the from-scratch mean, SD over runs and standard error of that mean in quadrature :: sha256 c6c4138ee5c86423 :: prose twin: two significant figures, direction in words
\newcommand{\MlstFtQgHerwigRthirtyMidEFive}{35.5\,\ensuremath{\pm}\,2.1}  % experiments/FIGS/data/v2/index.json :: fine-tuning qg_herwig N=100000: R30, R42_Q1 :: sha256 c6c4138ee5c86423
\newcommand{\MlstFtQgHerwigRthirtyMidOverScratchEFive}{0.97\,\ensuremath{\pm}\,0.06}  % experiments/FIGS/data/v2/index.json :: fine-tuning qg_herwig N=100000: R30 of each R42_Q1 run over the from-scratch mean, SD over runs and standard error of that mean in quadrature :: sha256 c6c4138ee5c86423
\newcommand{\MlstFtQgHerwigRthirtyMidOverScratchEFivePr}{equal within 6\%}  % experiments/FIGS/data/v2/index.json :: fine-tuning qg_herwig N=100000: R30 of each R42_Q1 run over the from-scratch mean, SD over runs and standard error of that mean in quadrature :: sha256 c6c4138ee5c86423 :: prose twin: two significant figures, direction in words
\newcommand{\MlstFtQgHerwigRthirtyCoarseEFive}{34.9\,\ensuremath{\pm}\,1.1}  % experiments/FIGS/data/v2/index.json :: fine-tuning qg_herwig N=100000: R30, R16_Q1 :: sha256 c6c4138ee5c86423
\newcommand{\MlstFtQgHerwigRthirtyCoarseOverScratchEFive}{0.950\,\ensuremath{\pm}\,0.033}  % experiments/FIGS/data/v2/index.json :: fine-tuning qg_herwig N=100000: R30 of each R16_Q1 run over the from-scratch mean, SD over runs and standard error of that mean in quadrature :: sha256 c6c4138ee5c86423
\newcommand{\MlstFtQgHerwigRthirtyCoarseOverScratchEFivePr}{equal within 6\%}  % experiments/FIGS/data/v2/index.json :: fine-tuning qg_herwig N=100000: R30 of each R16_Q1 run over the from-scratch mean, SD over runs and standard error of that mean in quadrature :: sha256 c6c4138ee5c86423 :: prose twin: two significant figures, direction in words
\newcommand{\MlstFtQgHerwigRthirtyResMassEFive}{34.0 and 36.5}  % experiments/FIGS/data/v2/index.json :: fine-tuning qg_herwig N=100000: R30, L162_MASS :: sha256 c6c4138ee5c86423
\newcommand{\MlstFtQgHerwigRthirtyResMassOverScratchEFive}{0.925 and 0.994}  % experiments/FIGS/data/v2/index.json :: fine-tuning qg_herwig N=100000: R30 of each L162_MASS run over the from-scratch mean, SD over runs and standard error of that mean in quadrature :: sha256 c6c4138ee5c86423
\newcommand{\MlstFtQgHerwigRthirtyResMassOverScratchEFivePr}{at least 0.63\% less rejection}  % experiments/FIGS/data/v2/index.json :: fine-tuning qg_herwig N=100000: R30 of each L162_MASS run over the from-scratch mean, SD over runs and standard error of that mean in quadrature :: sha256 c6c4138ee5c86423 :: prose twin: two significant figures, direction in words
\newcommand{\MlstFtQgHerwigRthirtyResMassOverScratchEFivePrMax}{up to 7.5\% less rejection}  % experiments/FIGS/data/v2/index.json :: fine-tuning qg_herwig N=100000: R30 of each L162_MASS run over the from-scratch mean, SD over runs and standard error of that mean in quadrature :: sha256 c6c4138ee5c86423 :: prose twin of the run farther from 1
\newcommand{\MlstFtQgHerwigRthirtyCoarseMassEFive}{35.9 and 39.0}  % experiments/FIGS/data/v2/index.json :: fine-tuning qg_herwig N=100000: R30, R16_Q1_MASS :: sha256 c6c4138ee5c86423
\newcommand{\MlstFtQgHerwigRthirtyCoarseMassOverScratchEFive}{0.979 and 1.064}  % experiments/FIGS/data/v2/index.json :: fine-tuning qg_herwig N=100000: R30 of each R16_Q1_MASS run over the from-scratch mean, SD over runs and standard error of that mean in quadrature :: sha256 c6c4138ee5c86423
\newcommand{\MlstFtQgHerwigRthirtyCoarseMassOverScratchEFivePr}{equal within 7\%}  % experiments/FIGS/data/v2/index.json :: fine-tuning qg_herwig N=100000: R30 of each R16_Q1_MASS run over the from-scratch mean, SD over runs and standard error of that mean in quadrature :: sha256 c6c4138ee5c86423 :: prose twin: two significant figures, direction in words
\newcommand{\MlstFtQgHerwigRthirtyCoarseMassOverScratchEFivePrMax}{equal within 7\%}  % experiments/FIGS/data/v2/index.json :: fine-tuning qg_herwig N=100000: R30 of each R16_Q1_MASS run over the from-scratch mean, SD over runs and standard error of that mean in quadrature :: sha256 c6c4138ee5c86423 :: prose twin of the run farther from 1
\newcommand{\MlstFtQgHerwigRthirtyResOverFullEFive}{1.03\,\ensuremath{\pm}\,0.04}  % experiments/FIGS/data/v2/index.json :: fine-tuning qg_herwig N=100000: R30 of L162 over L188, run by run :: sha256 c6c4138ee5c86423
\newcommand{\MlstFtQgHerwigRthirtyResOverFullEFivePr}{equal within 5\%}  % experiments/FIGS/data/v2/index.json :: fine-tuning qg_herwig N=100000: R30 of L162 over L188, run by run :: sha256 c6c4138ee5c86423 :: prose twin: two significant figures, direction in words
\newcommand{\MlstFtQgHerwigRthirtyMidOverFullEFive}{1.03\,\ensuremath{\pm}\,0.05}  % experiments/FIGS/data/v2/index.json :: fine-tuning qg_herwig N=100000: R30 of R42_Q1 over L188, run by run :: sha256 c6c4138ee5c86423
\newcommand{\MlstFtQgHerwigRthirtyMidOverFullEFivePr}{equal within 6\%}  % experiments/FIGS/data/v2/index.json :: fine-tuning qg_herwig N=100000: R30 of R42_Q1 over L188, run by run :: sha256 c6c4138ee5c86423 :: prose twin: two significant figures, direction in words
\newcommand{\MlstFtQgHerwigRthirtyCoarseOverFullEFive}{1.02\,\ensuremath{\pm}\,0.06}  % experiments/FIGS/data/v2/index.json :: fine-tuning qg_herwig N=100000: R30 of R16_Q1 over L188, run by run :: sha256 c6c4138ee5c86423
\newcommand{\MlstFtQgHerwigRthirtyCoarseOverFullEFivePr}{equal within 7\%}  % experiments/FIGS/data/v2/index.json :: fine-tuning qg_herwig N=100000: R30 of R16_Q1 over L188, run by run :: sha256 c6c4138ee5c86423 :: prose twin: two significant figures, direction in words
\newcommand{\MlstFtQgHerwigRthirtyResMassOverResEFive}{0.971 and 1.028}  % experiments/FIGS/data/v2/index.json :: fine-tuning qg_herwig N=100000: R30 of L162_MASS over L162, run by run :: sha256 c6c4138ee5c86423
\newcommand{\MlstFtQgHerwigRthirtyResMassOverResEFivePr}{equal within 3\%}  % experiments/FIGS/data/v2/index.json :: fine-tuning qg_herwig N=100000: R30 of L162_MASS over L162, run by run :: sha256 c6c4138ee5c86423 :: prose twin: two significant figures, direction in words
\newcommand{\MlstFtQgHerwigRthirtyResMassOverResEFivePrMax}{equal within 3\%}  % experiments/FIGS/data/v2/index.json :: fine-tuning qg_herwig N=100000: R30 of L162_MASS over L162, run by run :: sha256 c6c4138ee5c86423 :: prose twin of the run farther from 1
\newcommand{\MlstFtQgHerwigRthirtyCoarseMassOverCoarseEFive}{1.064 and 1.121}  % experiments/FIGS/data/v2/index.json :: fine-tuning qg_herwig N=100000: R30 of R16_Q1_MASS over R16_Q1, run by run :: sha256 c6c4138ee5c86423
\newcommand{\MlstFtQgHerwigRthirtyCoarseMassOverCoarseEFivePr}{equal within 20\%}  % experiments/FIGS/data/v2/index.json :: fine-tuning qg_herwig N=100000: R30 of R16_Q1_MASS over R16_Q1, run by run :: sha256 c6c4138ee5c86423 :: prose twin: two significant figures, direction in words
\newcommand{\MlstFtQgHerwigRthirtyCoarseMassOverCoarseEFivePrMax}{equal within 20\%}  % experiments/FIGS/data/v2/index.json :: fine-tuning qg_herwig N=100000: R30 of R16_Q1_MASS over R16_Q1, run by run :: sha256 c6c4138ee5c86423 :: prose twin of the run farther from 1
\newcommand{\MlstFtQgHerwigRthirtyRandOverFullEFive}{1.01\,\ensuremath{\pm}\,0.05}  % experiments/FIGS/data/v2/index.json :: fine-tuning qg_herwig N=100000: R30 of each random partition over the L188 run of its seed :: sha256 c6c4138ee5c86423
\newcommand{\MlstFtQgHerwigRthirtyRandOverFullEFivePr}{equal within 5\%}  % experiments/FIGS/data/v2/index.json :: fine-tuning qg_herwig N=100000: R30 of each random partition over the L188 run of its seed :: sha256 c6c4138ee5c86423 :: prose twin: two significant figures, direction in words
\newcommand{\MlstFtQgHerwigRthirtyRandOverCoarseEFive}{1.009\,\ensuremath{\pm}\,0.024}  % experiments/FIGS/data/v2/index.json :: fine-tuning qg_herwig N=100000: R30 of each random partition over the R16_Q1 run of its seed :: sha256 c6c4138ee5c86423
\newcommand{\MlstFtQgHerwigRthirtyRandOverCoarseEFivePr}{equal within 3\%}  % experiments/FIGS/data/v2/index.json :: fine-tuning qg_herwig N=100000: R30 of each random partition over the R16_Q1 run of its seed :: sha256 c6c4138ee5c86423 :: prose twin: two significant figures, direction in words
\newcommand{\MlstFtQgHerwigRthirtyScratch}{35.5\,\ensuremath{\pm}\,1.2}  % experiments/FIGS/data/v2/index.json :: fine-tuning qg_herwig N=1600000: R30 from scratch, every seed :: sha256 c6c4138ee5c86423
\newcommand{\MlstFtQgHerwigRthirtyFull}{35.2\,\ensuremath{\pm}\,2.1}  % experiments/FIGS/data/v2/index.json :: fine-tuning qg_herwig N=1600000: R30, L188 :: sha256 c6c4138ee5c86423
\newcommand{\MlstFtQgHerwigRthirtyFullOverScratch}{0.99\,\ensuremath{\pm}\,0.06}  % experiments/FIGS/data/v2/index.json :: fine-tuning qg_herwig N=1600000: R30 of each L188 run over the from-scratch mean, SD over runs and standard error of that mean in quadrature :: sha256 c6c4138ee5c86423
\newcommand{\MlstFtQgHerwigRthirtyFullOverScratchPr}{equal within 6\%}  % experiments/FIGS/data/v2/index.json :: fine-tuning qg_herwig N=1600000: R30 of each L188 run over the from-scratch mean, SD over runs and standard error of that mean in quadrature :: sha256 c6c4138ee5c86423 :: prose twin: two significant figures, direction in words
\newcommand{\MlstFtQgHerwigRthirtyRes}{34.6\,\ensuremath{\pm}\,2.0}  % experiments/FIGS/data/v2/index.json :: fine-tuning qg_herwig N=1600000: R30, L162 :: sha256 c6c4138ee5c86423
\newcommand{\MlstFtQgHerwigRthirtyResOverScratch}{0.97\,\ensuremath{\pm}\,0.06}  % experiments/FIGS/data/v2/index.json :: fine-tuning qg_herwig N=1600000: R30 of each L162 run over the from-scratch mean, SD over runs and standard error of that mean in quadrature :: sha256 c6c4138ee5c86423
\newcommand{\MlstFtQgHerwigRthirtyResOverScratchPr}{equal within 6\%}  % experiments/FIGS/data/v2/index.json :: fine-tuning qg_herwig N=1600000: R30 of each L162 run over the from-scratch mean, SD over runs and standard error of that mean in quadrature :: sha256 c6c4138ee5c86423 :: prose twin: two significant figures, direction in words
\newcommand{\MlstFtQgHerwigRthirtyMid}{36.1\,\ensuremath{\pm}\,2.7}  % experiments/FIGS/data/v2/index.json :: fine-tuning qg_herwig N=1600000: R30, R42_Q1 :: sha256 c6c4138ee5c86423
\newcommand{\MlstFtQgHerwigRthirtyMidOverScratch}{1.02\,\ensuremath{\pm}\,0.08}  % experiments/FIGS/data/v2/index.json :: fine-tuning qg_herwig N=1600000: R30 of each R42_Q1 run over the from-scratch mean, SD over runs and standard error of that mean in quadrature :: sha256 c6c4138ee5c86423
\newcommand{\MlstFtQgHerwigRthirtyMidOverScratchPr}{equal within 8\%}  % experiments/FIGS/data/v2/index.json :: fine-tuning qg_herwig N=1600000: R30 of each R42_Q1 run over the from-scratch mean, SD over runs and standard error of that mean in quadrature :: sha256 c6c4138ee5c86423 :: prose twin: two significant figures, direction in words
\newcommand{\MlstFtQgHerwigRthirtyCoarse}{36.9\,\ensuremath{\pm}\,0.9}  % experiments/FIGS/data/v2/index.json :: fine-tuning qg_herwig N=1600000: R30, R16_Q1 :: sha256 c6c4138ee5c86423
\newcommand{\MlstFtQgHerwigRthirtyCoarseOverScratch}{1.040\,\ensuremath{\pm}\,0.033}  % experiments/FIGS/data/v2/index.json :: fine-tuning qg_herwig N=1600000: R30 of each R16_Q1 run over the from-scratch mean, SD over runs and standard error of that mean in quadrature :: sha256 c6c4138ee5c86423
\newcommand{\MlstFtQgHerwigRthirtyCoarseOverScratchPr}{equal within 4\%}  % experiments/FIGS/data/v2/index.json :: fine-tuning qg_herwig N=1600000: R30 of each R16_Q1 run over the from-scratch mean, SD over runs and standard error of that mean in quadrature :: sha256 c6c4138ee5c86423 :: prose twin: two significant figures, direction in words
\newcommand{\MlstFtQgHerwigRthirtyResMass}{37.8 and 34.4}  % experiments/FIGS/data/v2/index.json :: fine-tuning qg_herwig N=1600000: R30, L162_MASS :: sha256 c6c4138ee5c86423
\newcommand{\MlstFtQgHerwigRthirtyResMassOverScratch}{1.065 and 0.970}  % experiments/FIGS/data/v2/index.json :: fine-tuning qg_herwig N=1600000: R30 of each L162_MASS run over the from-scratch mean, SD over runs and standard error of that mean in quadrature :: sha256 c6c4138ee5c86423
\newcommand{\MlstFtQgHerwigRthirtyResMassOverScratchPr}{equal within 7\%}  % experiments/FIGS/data/v2/index.json :: fine-tuning qg_herwig N=1600000: R30 of each L162_MASS run over the from-scratch mean, SD over runs and standard error of that mean in quadrature :: sha256 c6c4138ee5c86423 :: prose twin: two significant figures, direction in words
\newcommand{\MlstFtQgHerwigRthirtyResMassOverScratchPrMax}{equal within 7\%}  % experiments/FIGS/data/v2/index.json :: fine-tuning qg_herwig N=1600000: R30 of each L162_MASS run over the from-scratch mean, SD over runs and standard error of that mean in quadrature :: sha256 c6c4138ee5c86423 :: prose twin of the run farther from 1
\newcommand{\MlstFtQgHerwigRthirtyCoarseMass}{37.7 and 32.9}  % experiments/FIGS/data/v2/index.json :: fine-tuning qg_herwig N=1600000: R30, R16_Q1_MASS :: sha256 c6c4138ee5c86423
\newcommand{\MlstFtQgHerwigRthirtyCoarseMassOverScratch}{1.06 and 0.93}  % experiments/FIGS/data/v2/index.json :: fine-tuning qg_herwig N=1600000: R30 of each R16_Q1_MASS run over the from-scratch mean, SD over runs and standard error of that mean in quadrature :: sha256 c6c4138ee5c86423
\newcommand{\MlstFtQgHerwigRthirtyCoarseMassOverScratchPr}{equal within 8\%}  % experiments/FIGS/data/v2/index.json :: fine-tuning qg_herwig N=1600000: R30 of each R16_Q1_MASS run over the from-scratch mean, SD over runs and standard error of that mean in quadrature :: sha256 c6c4138ee5c86423 :: prose twin: two significant figures, direction in words
\newcommand{\MlstFtQgHerwigRthirtyCoarseMassOverScratchPrMax}{equal within 8\%}  % experiments/FIGS/data/v2/index.json :: fine-tuning qg_herwig N=1600000: R30 of each R16_Q1_MASS run over the from-scratch mean, SD over runs and standard error of that mean in quadrature :: sha256 c6c4138ee5c86423 :: prose twin of the run farther from 1
\newcommand{\MlstFtQgHerwigRthirtyResOverFull}{0.9832\,\ensuremath{\pm}\,0.0022}  % experiments/FIGS/data/v2/index.json :: fine-tuning qg_herwig N=1600000: R30 of L162 over L188, run by run :: sha256 c6c4138ee5c86423
\newcommand{\MlstFtQgHerwigRthirtyResOverFullPr}{1.7\% less rejection}  % experiments/FIGS/data/v2/index.json :: fine-tuning qg_herwig N=1600000: R30 of L162 over L188, run by run :: sha256 c6c4138ee5c86423 :: prose twin: two significant figures, direction in words
\newcommand{\MlstFtQgHerwigRthirtyMidOverFull}{1.03\,\ensuremath{\pm}\,0.05}  % experiments/FIGS/data/v2/index.json :: fine-tuning qg_herwig N=1600000: R30 of R42_Q1 over L188, run by run :: sha256 c6c4138ee5c86423
\newcommand{\MlstFtQgHerwigRthirtyMidOverFullPr}{equal within 6\%}  % experiments/FIGS/data/v2/index.json :: fine-tuning qg_herwig N=1600000: R30 of R42_Q1 over L188, run by run :: sha256 c6c4138ee5c86423 :: prose twin: two significant figures, direction in words
\newcommand{\MlstFtQgHerwigRthirtyCoarseOverFull}{1.05\,\ensuremath{\pm}\,0.06}  % experiments/FIGS/data/v2/index.json :: fine-tuning qg_herwig N=1600000: R30 of R16_Q1 over L188, run by run :: sha256 c6c4138ee5c86423
\newcommand{\MlstFtQgHerwigRthirtyCoarseOverFullPr}{equal within 6\%}  % experiments/FIGS/data/v2/index.json :: fine-tuning qg_herwig N=1600000: R30 of R16_Q1 over L188, run by run :: sha256 c6c4138ee5c86423 :: prose twin: two significant figures, direction in words
\newcommand{\MlstFtQgHerwigRthirtyResMassOverRes}{1.14 and 0.93}  % experiments/FIGS/data/v2/index.json :: fine-tuning qg_herwig N=1600000: R30 of L162_MASS over L162, run by run :: sha256 c6c4138ee5c86423
\newcommand{\MlstFtQgHerwigRthirtyResMassOverResPr}{equal within 20\%}  % experiments/FIGS/data/v2/index.json :: fine-tuning qg_herwig N=1600000: R30 of L162_MASS over L162, run by run :: sha256 c6c4138ee5c86423 :: prose twin: two significant figures, direction in words
\newcommand{\MlstFtQgHerwigRthirtyResMassOverResPrMax}{equal within 20\%}  % experiments/FIGS/data/v2/index.json :: fine-tuning qg_herwig N=1600000: R30 of L162_MASS over L162, run by run :: sha256 c6c4138ee5c86423 :: prose twin of the run farther from 1
\newcommand{\MlstFtQgHerwigRthirtyCoarseMassOverCoarse}{1.01 and 0.88}  % experiments/FIGS/data/v2/index.json :: fine-tuning qg_herwig N=1600000: R30 of R16_Q1_MASS over R16_Q1, run by run :: sha256 c6c4138ee5c86423
\newcommand{\MlstFtQgHerwigRthirtyCoarseMassOverCoarsePr}{equal within 20\%}  % experiments/FIGS/data/v2/index.json :: fine-tuning qg_herwig N=1600000: R30 of R16_Q1_MASS over R16_Q1, run by run :: sha256 c6c4138ee5c86423 :: prose twin: two significant figures, direction in words
\newcommand{\MlstFtQgHerwigRthirtyCoarseMassOverCoarsePrMax}{equal within 20\%}  % experiments/FIGS/data/v2/index.json :: fine-tuning qg_herwig N=1600000: R30 of R16_Q1_MASS over R16_Q1, run by run :: sha256 c6c4138ee5c86423 :: prose twin of the run farther from 1
\newcommand{\MlstFtQgHerwigRthirtyRandOverFull}{0.98\,\ensuremath{\pm}\,0.04}  % experiments/FIGS/data/v2/index.json :: fine-tuning qg_herwig N=1600000: R30 of each random partition over the L188 run of its seed :: sha256 c6c4138ee5c86423
\newcommand{\MlstFtQgHerwigRthirtyRandOverFullPr}{equal within 5\%}  % experiments/FIGS/data/v2/index.json :: fine-tuning qg_herwig N=1600000: R30 of each random partition over the L188 run of its seed :: sha256 c6c4138ee5c86423 :: prose twin: two significant figures, direction in words
\newcommand{\MlstFtQgHerwigRthirtyRandOverCoarse}{0.94\,\ensuremath{\pm}\,0.04}  % experiments/FIGS/data/v2/index.json :: fine-tuning qg_herwig N=1600000: R30 of each random partition over the R16_Q1 run of its seed :: sha256 c6c4138ee5c86423
\newcommand{\MlstFtQgHerwigRthirtyRandOverCoarsePr}{equal within 6\%}  % experiments/FIGS/data/v2/index.json :: fine-tuning qg_herwig N=1600000: R30 of each random partition over the R16_Q1 run of its seed :: sha256 c6c4138ee5c86423 :: prose twin: two significant figures, direction in words
\newcommand{\MlstFtQgRthirtyHerwigOverPythiaFull}{0.328\,\ensuremath{\pm}\,0.015}  % experiments/FIGS/data/v2/index.json :: fine-tuning q/g N=1600000: R30 on Herwig over R30 on Pythia, each run of L188 :: sha256 c6c4138ee5c86423
\newcommand{\MlstFtQgRthirtyHerwigOverPythiaFullPr}{0.33 times the rejection}  % experiments/FIGS/data/v2/index.json :: fine-tuning q/g N=1600000: R30 on Herwig over R30 on Pythia, each run of L188 :: sha256 c6c4138ee5c86423 :: prose twin: two significant figures, direction in words
\newcommand{\MlstFtQgRthirtyHerwigOverPythiaRes}{0.320\,\ensuremath{\pm}\,0.023}  % experiments/FIGS/data/v2/index.json :: fine-tuning q/g N=1600000: R30 on Herwig over R30 on Pythia, each run of L162 :: sha256 c6c4138ee5c86423
\newcommand{\MlstFtQgRthirtyHerwigOverPythiaResPr}{0.32 times the rejection}  % experiments/FIGS/data/v2/index.json :: fine-tuning q/g N=1600000: R30 on Herwig over R30 on Pythia, each run of L162 :: sha256 c6c4138ee5c86423 :: prose twin: two significant figures, direction in words
\newcommand{\MlstFtQgRthirtyHerwigOverPythiaMid}{0.334\,\ensuremath{\pm}\,0.023}  % experiments/FIGS/data/v2/index.json :: fine-tuning q/g N=1600000: R30 on Herwig over R30 on Pythia, each run of R42_Q1 :: sha256 c6c4138ee5c86423
\newcommand{\MlstFtQgRthirtyHerwigOverPythiaMidPr}{0.33 times the rejection}  % experiments/FIGS/data/v2/index.json :: fine-tuning q/g N=1600000: R30 on Herwig over R30 on Pythia, each run of R42_Q1 :: sha256 c6c4138ee5c86423 :: prose twin: two significant figures, direction in words
\newcommand{\MlstFtQgRthirtyHerwigOverPythiaCoarse}{0.346\,\ensuremath{\pm}\,0.011}  % experiments/FIGS/data/v2/index.json :: fine-tuning q/g N=1600000: R30 on Herwig over R30 on Pythia, each run of R16_Q1 :: sha256 c6c4138ee5c86423
\newcommand{\MlstFtQgRthirtyHerwigOverPythiaCoarsePr}{0.35 times the rejection}  % experiments/FIGS/data/v2/index.json :: fine-tuning q/g N=1600000: R30 on Herwig over R30 on Pythia, each run of R16_Q1 :: sha256 c6c4138ee5c86423 :: prose twin: two significant figures, direction in words
\newcommand{\MlstFtQgRthirtyHerwigOverPythiaResMass}{0.353 and 0.323}  % experiments/FIGS/data/v2/index.json :: fine-tuning q/g N=1600000: R30 on Herwig over R30 on Pythia, each run of L162_MASS :: sha256 c6c4138ee5c86423
\newcommand{\MlstFtQgRthirtyHerwigOverPythiaResMassPr}{at most 0.35 times the rejection}  % experiments/FIGS/data/v2/index.json :: fine-tuning q/g N=1600000: R30 on Herwig over R30 on Pythia, each run of L162_MASS :: sha256 c6c4138ee5c86423 :: prose twin: two significant figures, direction in words
\newcommand{\MlstFtQgRthirtyHerwigOverPythiaResMassPrMax}{down to 0.32 times the rejection}  % experiments/FIGS/data/v2/index.json :: fine-tuning q/g N=1600000: R30 on Herwig over R30 on Pythia, each run of L162_MASS :: sha256 c6c4138ee5c86423 :: prose twin of the run farther from 1
\newcommand{\MlstFtQgRthirtyHerwigOverPythiaCoarseMass}{0.359 and 0.311}  % experiments/FIGS/data/v2/index.json :: fine-tuning q/g N=1600000: R30 on Herwig over R30 on Pythia, each run of R16_Q1_MASS :: sha256 c6c4138ee5c86423
\newcommand{\MlstFtQgRthirtyHerwigOverPythiaCoarseMassPr}{at most 0.36 times the rejection}  % experiments/FIGS/data/v2/index.json :: fine-tuning q/g N=1600000: R30 on Herwig over R30 on Pythia, each run of R16_Q1_MASS :: sha256 c6c4138ee5c86423 :: prose twin: two significant figures, direction in words
\newcommand{\MlstFtQgRthirtyHerwigOverPythiaCoarseMassPrMax}{down to 0.31 times the rejection}  % experiments/FIGS/data/v2/index.json :: fine-tuning q/g N=1600000: R30 on Herwig over R30 on Pythia, each run of R16_Q1_MASS :: sha256 c6c4138ee5c86423 :: prose twin of the run farther from 1
\newcommand{\MlstFtQgRthirtyHerwigOverPythiaScratch}{0.341\,\ensuremath{\pm}\,0.009}  % experiments/FIGS/data/v2/index.json :: fine-tuning q/g N=1600000: R30 on Herwig over R30 on Pythia, each from-scratch seed :: sha256 c6c4138ee5c86423
\newcommand{\MlstFtQgRthirtyHerwigOverPythiaScratchPr}{0.34 times the rejection}  % experiments/FIGS/data/v2/index.json :: fine-tuning q/g N=1600000: R30 on Herwig over R30 on Pythia, each from-scratch seed :: sha256 c6c4138ee5c86423 :: prose twin: two significant figures, direction in words
\newcommand{\MlstTopCheckFullLong}{\pending{recipe check}}  % experiments/FIGS/data/v2/index.json :: the recipe check is not read out yet :: sha256 c6c4138ee5c86423
\newcommand{\MlstTopCheckScratchLong}{\pending{recipe check}}  % experiments/FIGS/data/v2/index.json :: the recipe check is not read out yet :: sha256 c6c4138ee5c86423
\newcommand{\MlstTopCheckFullShort}{\pending{recipe check}}  % experiments/FIGS/data/v2/index.json :: the recipe check is not read out yet :: sha256 c6c4138ee5c86423
\newcommand{\MlstTopCheckScratchShort}{\pending{recipe check}}  % experiments/FIGS/data/v2/index.json :: the recipe check is not read out yet :: sha256 c6c4138ee5c86423
\newcommand{\MlstLpJcOneGapFracEFour}{56.3\,\ensuremath{\pm}\,1.1\%}  % experiments/FIGS/data/v2/index.json :: jc1 N=10000: (probe - scratch) / (fine-tuned - scratch) accuracy, 188-class runs, per cent :: sha256 c6c4138ee5c86423
\newcommand{\MlstLpJcOneGapFracEFourPr}{56\%}  % experiments/FIGS/data/v2/index.json :: jc1 N=10000: (probe - scratch) / (fine-tuned - scratch) accuracy, 188-class runs, per cent :: sha256 c6c4138ee5c86423 :: prose twin, two significant figures
\newcommand{\MlstLpJcOneGapFracEFive}{\ensuremath{-}15\,\ensuremath{\pm}\,5\%}  % experiments/FIGS/data/v2/index.json :: jc1 N=100000: (probe - scratch) / (fine-tuned - scratch) accuracy, 188-class runs, per cent :: sha256 c6c4138ee5c86423
\newcommand{\MlstLpJcOneGapFracEFivePr}{\ensuremath{-}15\%}  % experiments/FIGS/data/v2/index.json :: jc1 N=100000: (probe - scratch) / (fine-tuned - scratch) accuracy, 188-class runs, per cent :: sha256 c6c4138ee5c86423 :: prose twin, two significant figures
\newcommand{\MlstLpJcTwoGapFracEFour}{96.2\,\ensuremath{\pm}\,0.6\%}  % experiments/FIGS/data/v2/index.json :: jc2_188 N=10000: (probe - scratch) / (fine-tuned - scratch) accuracy, 188-class runs, per cent :: sha256 c6c4138ee5c86423
\newcommand{\MlstLpJcTwoGapFracEFourPr}{96\%}  % experiments/FIGS/data/v2/index.json :: jc2_188 N=10000: (probe - scratch) / (fine-tuned - scratch) accuracy, 188-class runs, per cent :: sha256 c6c4138ee5c86423 :: prose twin, two significant figures
\newcommand{\MlstLpJcTwoGapFracEFive}{102.2\,\ensuremath{\pm}\,1.0\%}  % experiments/FIGS/data/v2/index.json :: jc2_188 N=100000: (probe - scratch) / (fine-tuned - scratch) accuracy, 188-class runs, per cent :: sha256 c6c4138ee5c86423
\newcommand{\MlstLpJcTwoGapFracEFivePr}{100\%}  % experiments/FIGS/data/v2/index.json :: jc2_188 N=100000: (probe - scratch) / (fine-tuned - scratch) accuracy, 188-class runs, per cent :: sha256 c6c4138ee5c86423 :: prose twin, two significant figures
\newcommand{\XbxMx}{\pending{samples not in}}  % experiments/XBS/xbx_report.py :: pending :: sha256 cb6dffe774ff5b1e
\newcommand{\XbxWindowLo}{\pending{samples not in}}  % experiments/XBS/xbx_report.py :: pending :: sha256 cb6dffe774ff5b1e
\newcommand{\XbxWindowHi}{\pending{samples not in}}  % experiments/XBS/xbx_report.py :: pending :: sha256 cb6dffe774ff5b1e
\newcommand{\XbbNTrainPerClass}{\pending{samples not in}}  % experiments/XBS/xbx_report.py :: pending :: sha256 cb6dffe774ff5b1e
\newcommand{\XbbRepetition}{\pending{samples not in}}  % experiments/XBS/xbx_report.py :: pending :: sha256 cb6dffe774ff5b1e
\newcommand{\XbbNEvalSignal}{\pending{samples not in}}  % experiments/XBS/xbx_report.py :: pending :: sha256 cb6dffe774ff5b1e
\newcommand{\XbbNEvalQcd}{\pending{samples not in}}  % experiments/XBS/xbx_report.py :: pending :: sha256 cb6dffe774ff5b1e
\newcommand{\XbsNTrainPerClass}{\pending{samples not in}}  % experiments/XBS/xbx_report.py :: pending :: sha256 cb6dffe774ff5b1e
\newcommand{\XbsRepetition}{\pending{samples not in}}  % experiments/XBS/xbx_report.py :: pending :: sha256 cb6dffe774ff5b1e
\newcommand{\XbsNEvalSignal}{\pending{samples not in}}  % experiments/XBS/xbx_report.py :: pending :: sha256 cb6dffe774ff5b1e
\newcommand{\XbsNEvalQcd}{\pending{samples not in}}  % experiments/XBS/xbx_report.py :: pending :: sha256 cb6dffe774ff5b1e
\newcommand{\MlstPartreproRthirtyPart}{1{,}711}  % experiments/FIGS/data/v2/partrepro/part.json :: r30 :: sha256 0d44b997af9e6448 :: ParT reproduction on top, part: R30 (one run)
\newcommand{\MlstPartreproRthirtyHead}{1{,}577}  % experiments/FIGS/data/v2/partrepro/head.json :: r30 :: sha256 37b82ba0a0fe9299 :: ParT reproduction on top, head: R30 (one run)
\newcommand{\MlstPartreproRthirtyRescale}{1{,}224}  % experiments/FIGS/data/v2/partrepro/rescale.json :: r30 :: sha256 0a9f1401047d07ff :: ParT reproduction on top, rescale: R30 (one run)
\newcommand{\MlstPartreproRthirtyHeadOverPart}{0.922}  % experiments/FIGS/data/v2/partrepro/head.json :: r30 / part r30 :: sha256 37b82ba0a0fe9299
\newcommand{\MlstPartreproRthirtyHeadOverPartPr}{7.8\% less rejection}  % experiments/FIGS/data/v2/partrepro/head.json :: r30 / part r30 :: sha256 37b82ba0a0fe9299 :: prose twin: two significant figures, direction in words
\newcommand{\MlstPartreproRthirtyRescaleOverPart}{0.715}  % experiments/FIGS/data/v2/partrepro/rescale.json :: r30 / part r30 :: sha256 0a9f1401047d07ff
\newcommand{\MlstPartreproRthirtyRescaleOverPartPr}{28\% less rejection}  % experiments/FIGS/data/v2/partrepro/rescale.json :: r30 / part r30 :: sha256 0a9f1401047d07ff :: prose twin: two significant figures, direction in words
\newcommand{\MlstPartreproSliceDropRatio}{\pending{jet-pT slice test not run}}  % experiments/FIGS/data/v2/index.json :: open item O10: R30 in narrow pT slices of partrepro part and rescale :: sha256 c6c4138ee5c86423
\newcommand{\MlstInputBnZeroSigmaMax}{0.96}  % experiments/FIGS/data/v2/head/input_bn_offsets.json :: offsets :: sha256 d6b821509fccecff :: largest |offset| over the five type flags
\newcommand{\MlstFtJcOneNSignal}{nine}  % experiments/FIGS/data/v2/curves/full_budget_long_leg2.json :: JC1_COLS (count) :: sha256 588e935996d1006a :: JetClass signal classes
\newcommand{\MlstFtJcOneRejHbbResOverFull}{0.92\,\ensuremath{\pm}\,0.05}  % experiments/FIGS/data/v2/curves/full_budget_long_leg2.json :: cells.<run>.N1000000.s1.jc1.per_class.label_Hbb: rejection of L162 over L188, run by run :: sha256 588e935996d1006a
\newcommand{\MlstFtJcOneRejHbbResOverFullPr}{equal within 9\%}  % experiments/FIGS/data/v2/curves/full_budget_long_leg2.json :: cells.<run>.N1000000.s1.jc1.per_class.label_Hbb: rejection of L162 over L188, run by run :: sha256 588e935996d1006a :: prose twin: two significant figures, direction in words
\newcommand{\MlstFtJcOneRejHbbMidOverFull}{0.98\,\ensuremath{\pm}\,0.12}  % experiments/FIGS/data/v2/curves/full_budget_long_leg2.json :: cells.<run>.N1000000.s1.jc1.per_class.label_Hbb: rejection of R42_Q1 over L188, run by run :: sha256 588e935996d1006a
\newcommand{\MlstFtJcOneRejHbbMidOverFullPr}{equal within 20\%}  % experiments/FIGS/data/v2/curves/full_budget_long_leg2.json :: cells.<run>.N1000000.s1.jc1.per_class.label_Hbb: rejection of R42_Q1 over L188, run by run :: sha256 588e935996d1006a :: prose twin: two significant figures, direction in words
\newcommand{\MlstFtJcOneRejHbbCoarseOverFull}{1.01\,\ensuremath{\pm}\,0.07}  % experiments/FIGS/data/v2/curves/full_budget_long_leg2.json :: cells.<run>.N1000000.s1.jc1.per_class.label_Hbb: rejection of R16_Q1 over L188, run by run :: sha256 588e935996d1006a
\newcommand{\MlstFtJcOneRejHbbCoarseOverFullPr}{equal within 7\%}  % experiments/FIGS/data/v2/curves/full_budget_long_leg2.json :: cells.<run>.N1000000.s1.jc1.per_class.label_Hbb: rejection of R16_Q1 over L188, run by run :: sha256 588e935996d1006a :: prose twin: two significant figures, direction in words
\newcommand{\MlstFtJcOneRejHbbResMassOverRes}{1.065 and 1.100}  % experiments/FIGS/data/v2/curves/full_budget_long_leg2.json :: cells.<run>.N1000000.s1.jc1.per_class.label_Hbb: rejection of L162_MASS over L162, run by run :: sha256 588e935996d1006a
\newcommand{\MlstFtJcOneRejHbbResMassOverResPr}{equal within 10\%}  % experiments/FIGS/data/v2/curves/full_budget_long_leg2.json :: cells.<run>.N1000000.s1.jc1.per_class.label_Hbb: rejection of L162_MASS over L162, run by run :: sha256 588e935996d1006a :: prose twin: two significant figures, direction in words
\newcommand{\MlstFtJcOneRejHbbResMassOverResPrMax}{equal within 10\%}  % experiments/FIGS/data/v2/curves/full_budget_long_leg2.json :: cells.<run>.N1000000.s1.jc1.per_class.label_Hbb: rejection of L162_MASS over L162, run by run :: sha256 588e935996d1006a :: prose twin of the run farther from 1
\newcommand{\MlstFtJcOneRejHbbCoarseMassOverCoarse}{0.97 and 0.86}  % experiments/FIGS/data/v2/curves/full_budget_long_leg2.json :: cells.<run>.N1000000.s1.jc1.per_class.label_Hbb: rejection of R16_Q1_MASS over R16_Q1, run by run :: sha256 588e935996d1006a
\newcommand{\MlstFtJcOneRejHbbCoarseMassOverCoarsePr}{equal within 20\%}  % experiments/FIGS/data/v2/curves/full_budget_long_leg2.json :: cells.<run>.N1000000.s1.jc1.per_class.label_Hbb: rejection of R16_Q1_MASS over R16_Q1, run by run :: sha256 588e935996d1006a :: prose twin: two significant figures, direction in words
\newcommand{\MlstFtJcOneRejHbbCoarseMassOverCoarsePrMax}{equal within 20\%}  % experiments/FIGS/data/v2/curves/full_budget_long_leg2.json :: cells.<run>.N1000000.s1.jc1.per_class.label_Hbb: rejection of R16_Q1_MASS over R16_Q1, run by run :: sha256 588e935996d1006a :: prose twin of the run farther from 1
\newcommand{\MlstFtJcOneRejHbbRandOverFull}{0.92\,\ensuremath{\pm}\,0.11}  % experiments/FIGS/data/v2/curves/full_budget_long_leg2.json :: per_class.label_Hbb: rejection of each random partition over the 188-class run of its seed :: sha256 588e935996d1006a
\newcommand{\MlstFtJcOneRejHbbRandOverFullPr}{equal within 20\%}  % experiments/FIGS/data/v2/curves/full_budget_long_leg2.json :: per_class.label_Hbb: rejection of each random partition over the 188-class run of its seed :: sha256 588e935996d1006a :: prose twin: two significant figures, direction in words
\newcommand{\MlstFtJcOneRejHbbRandOverFullRange}{0.84--1.1}  % experiments/FIGS/data/v2/curves/full_budget_long_leg2.json :: per_class.label_Hbb: smallest and largest of the partitions' ratios :: sha256 588e935996d1006a
\newcommand{\MlstFtJcOneRejHccResOverFull}{1.07\,\ensuremath{\pm}\,0.12}  % experiments/FIGS/data/v2/curves/full_budget_long_leg2.json :: cells.<run>.N1000000.s1.jc1.per_class.label_Hcc: rejection of L162 over L188, run by run :: sha256 588e935996d1006a
\newcommand{\MlstFtJcOneRejHccResOverFullPr}{equal within 20\%}  % experiments/FIGS/data/v2/curves/full_budget_long_leg2.json :: cells.<run>.N1000000.s1.jc1.per_class.label_Hcc: rejection of L162 over L188, run by run :: sha256 588e935996d1006a :: prose twin: two significant figures, direction in words
\newcommand{\MlstFtJcOneRejHccMidOverFull}{0.99\,\ensuremath{\pm}\,0.10}  % experiments/FIGS/data/v2/curves/full_budget_long_leg2.json :: cells.<run>.N1000000.s1.jc1.per_class.label_Hcc: rejection of R42_Q1 over L188, run by run :: sha256 588e935996d1006a
\newcommand{\MlstFtJcOneRejHccMidOverFullPr}{equal within 20\%}  % experiments/FIGS/data/v2/curves/full_budget_long_leg2.json :: cells.<run>.N1000000.s1.jc1.per_class.label_Hcc: rejection of R42_Q1 over L188, run by run :: sha256 588e935996d1006a :: prose twin: two significant figures, direction in words
\newcommand{\MlstFtJcOneRejHccCoarseOverFull}{0.91\,\ensuremath{\pm}\,0.07}  % experiments/FIGS/data/v2/curves/full_budget_long_leg2.json :: cells.<run>.N1000000.s1.jc1.per_class.label_Hcc: rejection of R16_Q1 over L188, run by run :: sha256 588e935996d1006a
\newcommand{\MlstFtJcOneRejHccCoarseOverFullPr}{equal within 10\%}  % experiments/FIGS/data/v2/curves/full_budget_long_leg2.json :: cells.<run>.N1000000.s1.jc1.per_class.label_Hcc: rejection of R16_Q1 over L188, run by run :: sha256 588e935996d1006a :: prose twin: two significant figures, direction in words
\newcommand{\MlstFtJcOneRejHccResMassOverRes}{0.886 and 0.816}  % experiments/FIGS/data/v2/curves/full_budget_long_leg2.json :: cells.<run>.N1000000.s1.jc1.per_class.label_Hcc: rejection of L162_MASS over L162, run by run :: sha256 588e935996d1006a
\newcommand{\MlstFtJcOneRejHccResMassOverResPr}{equal within 20\%}  % experiments/FIGS/data/v2/curves/full_budget_long_leg2.json :: cells.<run>.N1000000.s1.jc1.per_class.label_Hcc: rejection of L162_MASS over L162, run by run :: sha256 588e935996d1006a :: prose twin: two significant figures, direction in words
\newcommand{\MlstFtJcOneRejHccResMassOverResPrMax}{equal within 20\%}  % experiments/FIGS/data/v2/curves/full_budget_long_leg2.json :: cells.<run>.N1000000.s1.jc1.per_class.label_Hcc: rejection of L162_MASS over L162, run by run :: sha256 588e935996d1006a :: prose twin of the run farther from 1
\newcommand{\MlstFtJcOneRejHccCoarseMassOverCoarse}{0.98 and 0.85}  % experiments/FIGS/data/v2/curves/full_budget_long_leg2.json :: cells.<run>.N1000000.s1.jc1.per_class.label_Hcc: rejection of R16_Q1_MASS over R16_Q1, run by run :: sha256 588e935996d1006a
\newcommand{\MlstFtJcOneRejHccCoarseMassOverCoarsePr}{equal within 20\%}  % experiments/FIGS/data/v2/curves/full_budget_long_leg2.json :: cells.<run>.N1000000.s1.jc1.per_class.label_Hcc: rejection of R16_Q1_MASS over R16_Q1, run by run :: sha256 588e935996d1006a :: prose twin: two significant figures, direction in words
\newcommand{\MlstFtJcOneRejHccCoarseMassOverCoarsePrMax}{equal within 20\%}  % experiments/FIGS/data/v2/curves/full_budget_long_leg2.json :: cells.<run>.N1000000.s1.jc1.per_class.label_Hcc: rejection of R16_Q1_MASS over R16_Q1, run by run :: sha256 588e935996d1006a :: prose twin of the run farther from 1
\newcommand{\MlstFtJcOneRejHccRandOverFull}{1.03\,\ensuremath{\pm}\,0.07}  % experiments/FIGS/data/v2/curves/full_budget_long_leg2.json :: per_class.label_Hcc: rejection of each random partition over the 188-class run of its seed :: sha256 588e935996d1006a
\newcommand{\MlstFtJcOneRejHccRandOverFullPr}{equal within 8\%}  % experiments/FIGS/data/v2/curves/full_budget_long_leg2.json :: per_class.label_Hcc: rejection of each random partition over the 188-class run of its seed :: sha256 588e935996d1006a :: prose twin: two significant figures, direction in words
\newcommand{\MlstFtJcOneRejHccRandOverFullRange}{0.94--1.1}  % experiments/FIGS/data/v2/curves/full_budget_long_leg2.json :: per_class.label_Hcc: smallest and largest of the partitions' ratios :: sha256 588e935996d1006a
\newcommand{\MlstFtJcOneRejHggResOverFull}{0.98\,\ensuremath{\pm}\,0.05}  % experiments/FIGS/data/v2/curves/full_budget_long_leg2.json :: cells.<run>.N1000000.s1.jc1.per_class.label_Hgg: rejection of L162 over L188, run by run :: sha256 588e935996d1006a
\newcommand{\MlstFtJcOneRejHggResOverFullPr}{equal within 6\%}  % experiments/FIGS/data/v2/curves/full_budget_long_leg2.json :: cells.<run>.N1000000.s1.jc1.per_class.label_Hgg: rejection of L162 over L188, run by run :: sha256 588e935996d1006a :: prose twin: two significant figures, direction in words
\newcommand{\MlstFtJcOneRejHggMidOverFull}{1.000\,\ensuremath{\pm}\,0.025}  % experiments/FIGS/data/v2/curves/full_budget_long_leg2.json :: cells.<run>.N1000000.s1.jc1.per_class.label_Hgg: rejection of R42_Q1 over L188, run by run :: sha256 588e935996d1006a
\newcommand{\MlstFtJcOneRejHggMidOverFullPr}{equal within 3\%}  % experiments/FIGS/data/v2/curves/full_budget_long_leg2.json :: cells.<run>.N1000000.s1.jc1.per_class.label_Hgg: rejection of R42_Q1 over L188, run by run :: sha256 588e935996d1006a :: prose twin: two significant figures, direction in words
\newcommand{\MlstFtJcOneRejHggCoarseOverFull}{0.97\,\ensuremath{\pm}\,0.09}  % experiments/FIGS/data/v2/curves/full_budget_long_leg2.json :: cells.<run>.N1000000.s1.jc1.per_class.label_Hgg: rejection of R16_Q1 over L188, run by run :: sha256 588e935996d1006a
\newcommand{\MlstFtJcOneRejHggCoarseOverFullPr}{equal within 10\%}  % experiments/FIGS/data/v2/curves/full_budget_long_leg2.json :: cells.<run>.N1000000.s1.jc1.per_class.label_Hgg: rejection of R16_Q1 over L188, run by run :: sha256 588e935996d1006a :: prose twin: two significant figures, direction in words
\newcommand{\MlstFtJcOneRejHggResMassOverRes}{1.046 and 1.008}  % experiments/FIGS/data/v2/curves/full_budget_long_leg2.json :: cells.<run>.N1000000.s1.jc1.per_class.label_Hgg: rejection of L162_MASS over L162, run by run :: sha256 588e935996d1006a
\newcommand{\MlstFtJcOneRejHggResMassOverResPr}{equal within 5\%}  % experiments/FIGS/data/v2/curves/full_budget_long_leg2.json :: cells.<run>.N1000000.s1.jc1.per_class.label_Hgg: rejection of L162_MASS over L162, run by run :: sha256 588e935996d1006a :: prose twin: two significant figures, direction in words
\newcommand{\MlstFtJcOneRejHggResMassOverResPrMax}{equal within 5\%}  % experiments/FIGS/data/v2/curves/full_budget_long_leg2.json :: cells.<run>.N1000000.s1.jc1.per_class.label_Hgg: rejection of L162_MASS over L162, run by run :: sha256 588e935996d1006a :: prose twin of the run farther from 1
\newcommand{\MlstFtJcOneRejHggCoarseMassOverCoarse}{1.12 and 1.01}  % experiments/FIGS/data/v2/curves/full_budget_long_leg2.json :: cells.<run>.N1000000.s1.jc1.per_class.label_Hgg: rejection of R16_Q1_MASS over R16_Q1, run by run :: sha256 588e935996d1006a
\newcommand{\MlstFtJcOneRejHggCoarseMassOverCoarsePr}{equal within 20\%}  % experiments/FIGS/data/v2/curves/full_budget_long_leg2.json :: cells.<run>.N1000000.s1.jc1.per_class.label_Hgg: rejection of R16_Q1_MASS over R16_Q1, run by run :: sha256 588e935996d1006a :: prose twin: two significant figures, direction in words
\newcommand{\MlstFtJcOneRejHggCoarseMassOverCoarsePrMax}{equal within 20\%}  % experiments/FIGS/data/v2/curves/full_budget_long_leg2.json :: cells.<run>.N1000000.s1.jc1.per_class.label_Hgg: rejection of R16_Q1_MASS over R16_Q1, run by run :: sha256 588e935996d1006a :: prose twin of the run farther from 1
\newcommand{\MlstFtJcOneRejHggRandOverFull}{0.968\,\ensuremath{\pm}\,0.016}  % experiments/FIGS/data/v2/curves/full_budget_long_leg2.json :: per_class.label_Hgg: rejection of each random partition over the 188-class run of its seed :: sha256 588e935996d1006a
\newcommand{\MlstFtJcOneRejHggRandOverFullPr}{equal within 4\%}  % experiments/FIGS/data/v2/curves/full_budget_long_leg2.json :: per_class.label_Hgg: rejection of each random partition over the 188-class run of its seed :: sha256 588e935996d1006a :: prose twin: two significant figures, direction in words
\newcommand{\MlstFtJcOneRejHggRandOverFullRange}{0.94--0.98}  % experiments/FIGS/data/v2/curves/full_budget_long_leg2.json :: per_class.label_Hgg: smallest and largest of the partitions' ratios :: sha256 588e935996d1006a
\newcommand{\MlstFtJcOneRejHfourqResOverFull}{1.05\,\ensuremath{\pm}\,0.06}  % experiments/FIGS/data/v2/curves/full_budget_long_leg2.json :: cells.<run>.N1000000.s1.jc1.per_class.label_H4q: rejection of L162 over L188, run by run :: sha256 588e935996d1006a
\newcommand{\MlstFtJcOneRejHfourqResOverFullPr}{equal within 7\%}  % experiments/FIGS/data/v2/curves/full_budget_long_leg2.json :: cells.<run>.N1000000.s1.jc1.per_class.label_H4q: rejection of L162 over L188, run by run :: sha256 588e935996d1006a :: prose twin: two significant figures, direction in words
\newcommand{\MlstFtJcOneRejHfourqMidOverFull}{0.95\,\ensuremath{\pm}\,0.04}  % experiments/FIGS/data/v2/curves/full_budget_long_leg2.json :: cells.<run>.N1000000.s1.jc1.per_class.label_H4q: rejection of R42_Q1 over L188, run by run :: sha256 588e935996d1006a
\newcommand{\MlstFtJcOneRejHfourqMidOverFullPr}{equal within 6\%}  % experiments/FIGS/data/v2/curves/full_budget_long_leg2.json :: cells.<run>.N1000000.s1.jc1.per_class.label_H4q: rejection of R42_Q1 over L188, run by run :: sha256 588e935996d1006a :: prose twin: two significant figures, direction in words
\newcommand{\MlstFtJcOneRejHfourqCoarseOverFull}{0.962\,\ensuremath{\pm}\,0.025}  % experiments/FIGS/data/v2/curves/full_budget_long_leg2.json :: cells.<run>.N1000000.s1.jc1.per_class.label_H4q: rejection of R16_Q1 over L188, run by run :: sha256 588e935996d1006a
\newcommand{\MlstFtJcOneRejHfourqCoarseOverFullPr}{equal within 4\%}  % experiments/FIGS/data/v2/curves/full_budget_long_leg2.json :: cells.<run>.N1000000.s1.jc1.per_class.label_H4q: rejection of R16_Q1 over L188, run by run :: sha256 588e935996d1006a :: prose twin: two significant figures, direction in words
\newcommand{\MlstFtJcOneRejHfourqResMassOverRes}{0.912 and 0.966}  % experiments/FIGS/data/v2/curves/full_budget_long_leg2.json :: cells.<run>.N1000000.s1.jc1.per_class.label_H4q: rejection of L162_MASS over L162, run by run :: sha256 588e935996d1006a
\newcommand{\MlstFtJcOneRejHfourqResMassOverResPr}{equal within 9\%}  % experiments/FIGS/data/v2/curves/full_budget_long_leg2.json :: cells.<run>.N1000000.s1.jc1.per_class.label_H4q: rejection of L162_MASS over L162, run by run :: sha256 588e935996d1006a :: prose twin: two significant figures, direction in words
\newcommand{\MlstFtJcOneRejHfourqResMassOverResPrMax}{equal within 9\%}  % experiments/FIGS/data/v2/curves/full_budget_long_leg2.json :: cells.<run>.N1000000.s1.jc1.per_class.label_H4q: rejection of L162_MASS over L162, run by run :: sha256 588e935996d1006a :: prose twin of the run farther from 1
\newcommand{\MlstFtJcOneRejHfourqCoarseMassOverCoarse}{0.937 and 0.888}  % experiments/FIGS/data/v2/curves/full_budget_long_leg2.json :: cells.<run>.N1000000.s1.jc1.per_class.label_H4q: rejection of R16_Q1_MASS over R16_Q1, run by run :: sha256 588e935996d1006a
\newcommand{\MlstFtJcOneRejHfourqCoarseMassOverCoarsePr}{equal within 20\%}  % experiments/FIGS/data/v2/curves/full_budget_long_leg2.json :: cells.<run>.N1000000.s1.jc1.per_class.label_H4q: rejection of R16_Q1_MASS over R16_Q1, run by run :: sha256 588e935996d1006a :: prose twin: two significant figures, direction in words
\newcommand{\MlstFtJcOneRejHfourqCoarseMassOverCoarsePrMax}{equal within 20\%}  % experiments/FIGS/data/v2/curves/full_budget_long_leg2.json :: cells.<run>.N1000000.s1.jc1.per_class.label_H4q: rejection of R16_Q1_MASS over R16_Q1, run by run :: sha256 588e935996d1006a :: prose twin of the run farther from 1
\newcommand{\MlstFtJcOneRejHfourqRandOverFull}{0.93\,\ensuremath{\pm}\,0.04}  % experiments/FIGS/data/v2/curves/full_budget_long_leg2.json :: per_class.label_H4q: rejection of each random partition over the 188-class run of its seed :: sha256 588e935996d1006a
\newcommand{\MlstFtJcOneRejHfourqRandOverFullPr}{equal within 8\%}  % experiments/FIGS/data/v2/curves/full_budget_long_leg2.json :: per_class.label_H4q: rejection of each random partition over the 188-class run of its seed :: sha256 588e935996d1006a :: prose twin: two significant figures, direction in words
\newcommand{\MlstFtJcOneRejHfourqRandOverFullRange}{0.86--0.96}  % experiments/FIGS/data/v2/curves/full_budget_long_leg2.json :: per_class.label_H4q: smallest and largest of the partitions' ratios :: sha256 588e935996d1006a
\newcommand{\MlstFtJcOneRejHqqlResOverFull}{1.004\,\ensuremath{\pm}\,0.018}  % experiments/FIGS/data/v2/curves/full_budget_long_leg2.json :: cells.<run>.N1000000.s1.jc1.per_class.label_Hqql: rejection of L162 over L188, run by run :: sha256 588e935996d1006a
\newcommand{\MlstFtJcOneRejHqqlResOverFullPr}{equal within 2\%}  % experiments/FIGS/data/v2/curves/full_budget_long_leg2.json :: cells.<run>.N1000000.s1.jc1.per_class.label_Hqql: rejection of L162 over L188, run by run :: sha256 588e935996d1006a :: prose twin: two significant figures, direction in words
\newcommand{\MlstFtJcOneRejHqqlMidOverFull}{1.04\,\ensuremath{\pm}\,0.18}  % experiments/FIGS/data/v2/curves/full_budget_long_leg2.json :: cells.<run>.N1000000.s1.jc1.per_class.label_Hqql: rejection of R42_Q1 over L188, run by run :: sha256 588e935996d1006a
\newcommand{\MlstFtJcOneRejHqqlMidOverFullPr}{equal within 20\%}  % experiments/FIGS/data/v2/curves/full_budget_long_leg2.json :: cells.<run>.N1000000.s1.jc1.per_class.label_Hqql: rejection of R42_Q1 over L188, run by run :: sha256 588e935996d1006a :: prose twin: two significant figures, direction in words
\newcommand{\MlstFtJcOneRejHqqlCoarseOverFull}{0.93\,\ensuremath{\pm}\,0.16}  % experiments/FIGS/data/v2/curves/full_budget_long_leg2.json :: cells.<run>.N1000000.s1.jc1.per_class.label_Hqql: rejection of R16_Q1 over L188, run by run :: sha256 588e935996d1006a
\newcommand{\MlstFtJcOneRejHqqlCoarseOverFullPr}{equal within 20\%}  % experiments/FIGS/data/v2/curves/full_budget_long_leg2.json :: cells.<run>.N1000000.s1.jc1.per_class.label_Hqql: rejection of R16_Q1 over L188, run by run :: sha256 588e935996d1006a :: prose twin: two significant figures, direction in words
\newcommand{\MlstFtJcOneRejHqqlResMassOverRes}{0.986 and 1.000}  % experiments/FIGS/data/v2/curves/full_budget_long_leg2.json :: cells.<run>.N1000000.s1.jc1.per_class.label_Hqql: rejection of L162_MASS over L162, run by run :: sha256 588e935996d1006a
\newcommand{\MlstFtJcOneRejHqqlResMassOverResPr}{equal within 2\%}  % experiments/FIGS/data/v2/curves/full_budget_long_leg2.json :: cells.<run>.N1000000.s1.jc1.per_class.label_Hqql: rejection of L162_MASS over L162, run by run :: sha256 588e935996d1006a :: prose twin: two significant figures, direction in words
\newcommand{\MlstFtJcOneRejHqqlResMassOverResPrMax}{equal within 2\%}  % experiments/FIGS/data/v2/curves/full_budget_long_leg2.json :: cells.<run>.N1000000.s1.jc1.per_class.label_Hqql: rejection of L162_MASS over L162, run by run :: sha256 588e935996d1006a :: prose twin of the run farther from 1
\newcommand{\MlstFtJcOneRejHqqlCoarseMassOverCoarse}{0.97 and 1.09}  % experiments/FIGS/data/v2/curves/full_budget_long_leg2.json :: cells.<run>.N1000000.s1.jc1.per_class.label_Hqql: rejection of R16_Q1_MASS over R16_Q1, run by run :: sha256 588e935996d1006a
\newcommand{\MlstFtJcOneRejHqqlCoarseMassOverCoarsePr}{equal within 9\%}  % experiments/FIGS/data/v2/curves/full_budget_long_leg2.json :: cells.<run>.N1000000.s1.jc1.per_class.label_Hqql: rejection of R16_Q1_MASS over R16_Q1, run by run :: sha256 588e935996d1006a :: prose twin: two significant figures, direction in words
\newcommand{\MlstFtJcOneRejHqqlCoarseMassOverCoarsePrMax}{equal within 9\%}  % experiments/FIGS/data/v2/curves/full_budget_long_leg2.json :: cells.<run>.N1000000.s1.jc1.per_class.label_Hqql: rejection of R16_Q1_MASS over R16_Q1, run by run :: sha256 588e935996d1006a :: prose twin of the run farther from 1
\newcommand{\MlstFtJcOneRejHqqlRandOverFull}{0.94\,\ensuremath{\pm}\,0.12}  % experiments/FIGS/data/v2/curves/full_budget_long_leg2.json :: per_class.label_Hqql: rejection of each random partition over the 188-class run of its seed :: sha256 588e935996d1006a
\newcommand{\MlstFtJcOneRejHqqlRandOverFullPr}{equal within 20\%}  % experiments/FIGS/data/v2/curves/full_budget_long_leg2.json :: per_class.label_Hqql: rejection of each random partition over the 188-class run of its seed :: sha256 588e935996d1006a :: prose twin: two significant figures, direction in words
\newcommand{\MlstFtJcOneRejHqqlRandOverFullRange}{0.79--1.1}  % experiments/FIGS/data/v2/curves/full_budget_long_leg2.json :: per_class.label_Hqql: smallest and largest of the partitions' ratios :: sha256 588e935996d1006a
\newcommand{\MlstFtJcOneRejTbqqResOverFull}{1.1\,\ensuremath{\pm}\,0.4}  % experiments/FIGS/data/v2/curves/full_budget_long_leg2.json :: cells.<run>.N1000000.s1.jc1.per_class.label_Tbqq: rejection of L162 over L188, run by run :: sha256 588e935996d1006a
\newcommand{\MlstFtJcOneRejTbqqResOverFullPr}{equal within 50\%}  % experiments/FIGS/data/v2/curves/full_budget_long_leg2.json :: cells.<run>.N1000000.s1.jc1.per_class.label_Tbqq: rejection of L162 over L188, run by run :: sha256 588e935996d1006a :: prose twin: two significant figures, direction in words
\newcommand{\MlstFtJcOneRejTbqqMidOverFull}{1.10\,\ensuremath{\pm}\,0.05}  % experiments/FIGS/data/v2/curves/full_budget_long_leg2.json :: cells.<run>.N1000000.s1.jc1.per_class.label_Tbqq: rejection of R42_Q1 over L188, run by run :: sha256 588e935996d1006a
\newcommand{\MlstFtJcOneRejTbqqMidOverFullPr}{equal within 10\%}  % experiments/FIGS/data/v2/curves/full_budget_long_leg2.json :: cells.<run>.N1000000.s1.jc1.per_class.label_Tbqq: rejection of R42_Q1 over L188, run by run :: sha256 588e935996d1006a :: prose twin: two significant figures, direction in words
\newcommand{\MlstFtJcOneRejTbqqCoarseOverFull}{0.85\,\ensuremath{\pm}\,0.32}  % experiments/FIGS/data/v2/curves/full_budget_long_leg2.json :: cells.<run>.N1000000.s1.jc1.per_class.label_Tbqq: rejection of R16_Q1 over L188, run by run :: sha256 588e935996d1006a
\newcommand{\MlstFtJcOneRejTbqqCoarseOverFullPr}{equal within 40\%}  % experiments/FIGS/data/v2/curves/full_budget_long_leg2.json :: cells.<run>.N1000000.s1.jc1.per_class.label_Tbqq: rejection of R16_Q1 over L188, run by run :: sha256 588e935996d1006a :: prose twin: two significant figures, direction in words
\newcommand{\MlstFtJcOneRejTbqqResMassOverRes}{0.89 and 1.21}  % experiments/FIGS/data/v2/curves/full_budget_long_leg2.json :: cells.<run>.N1000000.s1.jc1.per_class.label_Tbqq: rejection of L162_MASS over L162, run by run :: sha256 588e935996d1006a
\newcommand{\MlstFtJcOneRejTbqqResMassOverResPr}{equal within 30\%}  % experiments/FIGS/data/v2/curves/full_budget_long_leg2.json :: cells.<run>.N1000000.s1.jc1.per_class.label_Tbqq: rejection of L162_MASS over L162, run by run :: sha256 588e935996d1006a :: prose twin: two significant figures, direction in words
\newcommand{\MlstFtJcOneRejTbqqResMassOverResPrMax}{equal within 30\%}  % experiments/FIGS/data/v2/curves/full_budget_long_leg2.json :: cells.<run>.N1000000.s1.jc1.per_class.label_Tbqq: rejection of L162_MASS over L162, run by run :: sha256 588e935996d1006a :: prose twin of the run farther from 1
\newcommand{\MlstFtJcOneRejTbqqCoarseMassOverCoarse}{1.00 and 1.79}  % experiments/FIGS/data/v2/curves/full_budget_long_leg2.json :: cells.<run>.N1000000.s1.jc1.per_class.label_Tbqq: rejection of R16_Q1_MASS over R16_Q1, run by run :: sha256 588e935996d1006a
\newcommand{\MlstFtJcOneRejTbqqCoarseMassOverCoarsePr}{equal within 80\%}  % experiments/FIGS/data/v2/curves/full_budget_long_leg2.json :: cells.<run>.N1000000.s1.jc1.per_class.label_Tbqq: rejection of R16_Q1_MASS over R16_Q1, run by run :: sha256 588e935996d1006a :: prose twin: two significant figures, direction in words
\newcommand{\MlstFtJcOneRejTbqqCoarseMassOverCoarsePrMax}{equal within 80\%}  % experiments/FIGS/data/v2/curves/full_budget_long_leg2.json :: cells.<run>.N1000000.s1.jc1.per_class.label_Tbqq: rejection of R16_Q1_MASS over R16_Q1, run by run :: sha256 588e935996d1006a :: prose twin of the run farther from 1
\newcommand{\MlstFtJcOneRejTbqqRandOverFull}{0.79\,\ensuremath{\pm}\,0.15}  % experiments/FIGS/data/v2/curves/full_budget_long_leg2.json :: per_class.label_Tbqq: rejection of each random partition over the 188-class run of its seed :: sha256 588e935996d1006a
\newcommand{\MlstFtJcOneRejTbqqRandOverFullPr}{equal within 30\%}  % experiments/FIGS/data/v2/curves/full_budget_long_leg2.json :: per_class.label_Tbqq: rejection of each random partition over the 188-class run of its seed :: sha256 588e935996d1006a :: prose twin: two significant figures, direction in words
\newcommand{\MlstFtJcOneRejTbqqRandOverFullRange}{0.64--1.0}  % experiments/FIGS/data/v2/curves/full_budget_long_leg2.json :: per_class.label_Tbqq: smallest and largest of the partitions' ratios :: sha256 588e935996d1006a
\newcommand{\MlstFtJcOneRejTblResOverFull}{0.85\,\ensuremath{\pm}\,0.17}  % experiments/FIGS/data/v2/curves/full_budget_long_leg2.json :: cells.<run>.N1000000.s1.jc1.per_class.label_Tbl: rejection of L162 over L188, run by run :: sha256 588e935996d1006a
\newcommand{\MlstFtJcOneRejTblResOverFullPr}{equal within 20\%}  % experiments/FIGS/data/v2/curves/full_budget_long_leg2.json :: cells.<run>.N1000000.s1.jc1.per_class.label_Tbl: rejection of L162 over L188, run by run :: sha256 588e935996d1006a :: prose twin: two significant figures, direction in words
\newcommand{\MlstFtJcOneRejTblMidOverFull}{0.91\,\ensuremath{\pm}\,0.15}  % experiments/FIGS/data/v2/curves/full_budget_long_leg2.json :: cells.<run>.N1000000.s1.jc1.per_class.label_Tbl: rejection of R42_Q1 over L188, run by run :: sha256 588e935996d1006a
\newcommand{\MlstFtJcOneRejTblMidOverFullPr}{equal within 20\%}  % experiments/FIGS/data/v2/curves/full_budget_long_leg2.json :: cells.<run>.N1000000.s1.jc1.per_class.label_Tbl: rejection of R42_Q1 over L188, run by run :: sha256 588e935996d1006a :: prose twin: two significant figures, direction in words
\newcommand{\MlstFtJcOneRejTblCoarseOverFull}{0.94\,\ensuremath{\pm}\,0.28}  % experiments/FIGS/data/v2/curves/full_budget_long_leg2.json :: cells.<run>.N1000000.s1.jc1.per_class.label_Tbl: rejection of R16_Q1 over L188, run by run :: sha256 588e935996d1006a
\newcommand{\MlstFtJcOneRejTblCoarseOverFullPr}{equal within 30\%}  % experiments/FIGS/data/v2/curves/full_budget_long_leg2.json :: cells.<run>.N1000000.s1.jc1.per_class.label_Tbl: rejection of R16_Q1 over L188, run by run :: sha256 588e935996d1006a :: prose twin: two significant figures, direction in words
\newcommand{\MlstFtJcOneRejTblResMassOverRes}{1.04 and 0.93}  % experiments/FIGS/data/v2/curves/full_budget_long_leg2.json :: cells.<run>.N1000000.s1.jc1.per_class.label_Tbl: rejection of L162_MASS over L162, run by run :: sha256 588e935996d1006a
\newcommand{\MlstFtJcOneRejTblResMassOverResPr}{equal within 7\%}  % experiments/FIGS/data/v2/curves/full_budget_long_leg2.json :: cells.<run>.N1000000.s1.jc1.per_class.label_Tbl: rejection of L162_MASS over L162, run by run :: sha256 588e935996d1006a :: prose twin: two significant figures, direction in words
\newcommand{\MlstFtJcOneRejTblResMassOverResPrMax}{equal within 7\%}  % experiments/FIGS/data/v2/curves/full_budget_long_leg2.json :: cells.<run>.N1000000.s1.jc1.per_class.label_Tbl: rejection of L162_MASS over L162, run by run :: sha256 588e935996d1006a :: prose twin of the run farther from 1
\newcommand{\MlstFtJcOneRejTblCoarseMassOverCoarse}{0.84 and 1.22}  % experiments/FIGS/data/v2/curves/full_budget_long_leg2.json :: cells.<run>.N1000000.s1.jc1.per_class.label_Tbl: rejection of R16_Q1_MASS over R16_Q1, run by run :: sha256 588e935996d1006a
\newcommand{\MlstFtJcOneRejTblCoarseMassOverCoarsePr}{equal within 30\%}  % experiments/FIGS/data/v2/curves/full_budget_long_leg2.json :: cells.<run>.N1000000.s1.jc1.per_class.label_Tbl: rejection of R16_Q1_MASS over R16_Q1, run by run :: sha256 588e935996d1006a :: prose twin: two significant figures, direction in words
\newcommand{\MlstFtJcOneRejTblCoarseMassOverCoarsePrMax}{equal within 30\%}  % experiments/FIGS/data/v2/curves/full_budget_long_leg2.json :: cells.<run>.N1000000.s1.jc1.per_class.label_Tbl: rejection of R16_Q1_MASS over R16_Q1, run by run :: sha256 588e935996d1006a :: prose twin of the run farther from 1
\newcommand{\MlstFtJcOneRejTblRandOverFull}{0.88\,\ensuremath{\pm}\,0.26}  % experiments/FIGS/data/v2/curves/full_budget_long_leg2.json :: per_class.label_Tbl: rejection of each random partition over the 188-class run of its seed :: sha256 588e935996d1006a
\newcommand{\MlstFtJcOneRejTblRandOverFullPr}{equal within 30\%}  % experiments/FIGS/data/v2/curves/full_budget_long_leg2.json :: per_class.label_Tbl: rejection of each random partition over the 188-class run of its seed :: sha256 588e935996d1006a :: prose twin: two significant figures, direction in words
\newcommand{\MlstFtJcOneRejTblRandOverFullRange}{0.61--1.2}  % experiments/FIGS/data/v2/curves/full_budget_long_leg2.json :: per_class.label_Tbl: smallest and largest of the partitions' ratios :: sha256 588e935996d1006a
\newcommand{\MlstFtJcOneRejWqqResOverFull}{1.039\,\ensuremath{\pm}\,0.035}  % experiments/FIGS/data/v2/curves/full_budget_long_leg2.json :: cells.<run>.N1000000.s1.jc1.per_class.label_Wqq: rejection of L162 over L188, run by run :: sha256 588e935996d1006a
\newcommand{\MlstFtJcOneRejWqqResOverFullPr}{equal within 4\%}  % experiments/FIGS/data/v2/curves/full_budget_long_leg2.json :: cells.<run>.N1000000.s1.jc1.per_class.label_Wqq: rejection of L162 over L188, run by run :: sha256 588e935996d1006a :: prose twin: two significant figures, direction in words
\newcommand{\MlstFtJcOneRejWqqMidOverFull}{1.005\,\ensuremath{\pm}\,0.032}  % experiments/FIGS/data/v2/curves/full_budget_long_leg2.json :: cells.<run>.N1000000.s1.jc1.per_class.label_Wqq: rejection of R42_Q1 over L188, run by run :: sha256 588e935996d1006a
\newcommand{\MlstFtJcOneRejWqqMidOverFullPr}{equal within 4\%}  % experiments/FIGS/data/v2/curves/full_budget_long_leg2.json :: cells.<run>.N1000000.s1.jc1.per_class.label_Wqq: rejection of R42_Q1 over L188, run by run :: sha256 588e935996d1006a :: prose twin: two significant figures, direction in words
\newcommand{\MlstFtJcOneRejWqqCoarseOverFull}{1.01\,\ensuremath{\pm}\,0.07}  % experiments/FIGS/data/v2/curves/full_budget_long_leg2.json :: cells.<run>.N1000000.s1.jc1.per_class.label_Wqq: rejection of R16_Q1 over L188, run by run :: sha256 588e935996d1006a
\newcommand{\MlstFtJcOneRejWqqCoarseOverFullPr}{equal within 8\%}  % experiments/FIGS/data/v2/curves/full_budget_long_leg2.json :: cells.<run>.N1000000.s1.jc1.per_class.label_Wqq: rejection of R16_Q1 over L188, run by run :: sha256 588e935996d1006a :: prose twin: two significant figures, direction in words
\newcommand{\MlstFtJcOneRejWqqResMassOverRes}{0.996 and 0.964}  % experiments/FIGS/data/v2/curves/full_budget_long_leg2.json :: cells.<run>.N1000000.s1.jc1.per_class.label_Wqq: rejection of L162_MASS over L162, run by run :: sha256 588e935996d1006a
\newcommand{\MlstFtJcOneRejWqqResMassOverResPr}{equal within 4\%}  % experiments/FIGS/data/v2/curves/full_budget_long_leg2.json :: cells.<run>.N1000000.s1.jc1.per_class.label_Wqq: rejection of L162_MASS over L162, run by run :: sha256 588e935996d1006a :: prose twin: two significant figures, direction in words
\newcommand{\MlstFtJcOneRejWqqResMassOverResPrMax}{equal within 4\%}  % experiments/FIGS/data/v2/curves/full_budget_long_leg2.json :: cells.<run>.N1000000.s1.jc1.per_class.label_Wqq: rejection of L162_MASS over L162, run by run :: sha256 588e935996d1006a :: prose twin of the run farther from 1
\newcommand{\MlstFtJcOneRejWqqCoarseMassOverCoarse}{1.025 and 0.926}  % experiments/FIGS/data/v2/curves/full_budget_long_leg2.json :: cells.<run>.N1000000.s1.jc1.per_class.label_Wqq: rejection of R16_Q1_MASS over R16_Q1, run by run :: sha256 588e935996d1006a
\newcommand{\MlstFtJcOneRejWqqCoarseMassOverCoarsePr}{equal within 8\%}  % experiments/FIGS/data/v2/curves/full_budget_long_leg2.json :: cells.<run>.N1000000.s1.jc1.per_class.label_Wqq: rejection of R16_Q1_MASS over R16_Q1, run by run :: sha256 588e935996d1006a :: prose twin: two significant figures, direction in words
\newcommand{\MlstFtJcOneRejWqqCoarseMassOverCoarsePrMax}{equal within 8\%}  % experiments/FIGS/data/v2/curves/full_budget_long_leg2.json :: cells.<run>.N1000000.s1.jc1.per_class.label_Wqq: rejection of R16_Q1_MASS over R16_Q1, run by run :: sha256 588e935996d1006a :: prose twin of the run farther from 1
\newcommand{\MlstFtJcOneRejWqqRandOverFull}{1.02\,\ensuremath{\pm}\,0.07}  % experiments/FIGS/data/v2/curves/full_budget_long_leg2.json :: per_class.label_Wqq: rejection of each random partition over the 188-class run of its seed :: sha256 588e935996d1006a
\newcommand{\MlstFtJcOneRejWqqRandOverFullPr}{equal within 7\%}  % experiments/FIGS/data/v2/curves/full_budget_long_leg2.json :: per_class.label_Wqq: rejection of each random partition over the 188-class run of its seed :: sha256 588e935996d1006a :: prose twin: two significant figures, direction in words
\newcommand{\MlstFtJcOneRejWqqRandOverFullRange}{0.93--1.1}  % experiments/FIGS/data/v2/curves/full_budget_long_leg2.json :: per_class.label_Wqq: smallest and largest of the partitions' ratios :: sha256 588e935996d1006a
\newcommand{\MlstFtJcOneRejZqqResOverFull}{0.987\,\ensuremath{\pm}\,0.019}  % experiments/FIGS/data/v2/curves/full_budget_long_leg2.json :: cells.<run>.N1000000.s1.jc1.per_class.label_Zqq: rejection of L162 over L188, run by run :: sha256 588e935996d1006a
\newcommand{\MlstFtJcOneRejZqqResOverFullPr}{equal within 2\%}  % experiments/FIGS/data/v2/curves/full_budget_long_leg2.json :: cells.<run>.N1000000.s1.jc1.per_class.label_Zqq: rejection of L162 over L188, run by run :: sha256 588e935996d1006a :: prose twin: two significant figures, direction in words
\newcommand{\MlstFtJcOneRejZqqMidOverFull}{0.984\,\ensuremath{\pm}\,0.010}  % experiments/FIGS/data/v2/curves/full_budget_long_leg2.json :: cells.<run>.N1000000.s1.jc1.per_class.label_Zqq: rejection of R42_Q1 over L188, run by run :: sha256 588e935996d1006a
\newcommand{\MlstFtJcOneRejZqqMidOverFullPr}{equal within 2\%}  % experiments/FIGS/data/v2/curves/full_budget_long_leg2.json :: cells.<run>.N1000000.s1.jc1.per_class.label_Zqq: rejection of R42_Q1 over L188, run by run :: sha256 588e935996d1006a :: prose twin: two significant figures, direction in words
\newcommand{\MlstFtJcOneRejZqqCoarseOverFull}{0.91\,\ensuremath{\pm}\,0.05}  % experiments/FIGS/data/v2/curves/full_budget_long_leg2.json :: cells.<run>.N1000000.s1.jc1.per_class.label_Zqq: rejection of R16_Q1 over L188, run by run :: sha256 588e935996d1006a
\newcommand{\MlstFtJcOneRejZqqCoarseOverFullPr}{equal within 9\%}  % experiments/FIGS/data/v2/curves/full_budget_long_leg2.json :: cells.<run>.N1000000.s1.jc1.per_class.label_Zqq: rejection of R16_Q1 over L188, run by run :: sha256 588e935996d1006a :: prose twin: two significant figures, direction in words
\newcommand{\MlstFtJcOneRejZqqResMassOverRes}{0.9803 and 0.9820}  % experiments/FIGS/data/v2/curves/full_budget_long_leg2.json :: cells.<run>.N1000000.s1.jc1.per_class.label_Zqq: rejection of L162_MASS over L162, run by run :: sha256 588e935996d1006a
\newcommand{\MlstFtJcOneRejZqqResMassOverResPr}{equal within 2\%}  % experiments/FIGS/data/v2/curves/full_budget_long_leg2.json :: cells.<run>.N1000000.s1.jc1.per_class.label_Zqq: rejection of L162_MASS over L162, run by run :: sha256 588e935996d1006a :: prose twin: two significant figures, direction in words
\newcommand{\MlstFtJcOneRejZqqResMassOverResPrMax}{equal within 2\%}  % experiments/FIGS/data/v2/curves/full_budget_long_leg2.json :: cells.<run>.N1000000.s1.jc1.per_class.label_Zqq: rejection of L162_MASS over L162, run by run :: sha256 588e935996d1006a :: prose twin of the run farther from 1
\newcommand{\MlstFtJcOneRejZqqCoarseMassOverCoarse}{1.06 and 0.94}  % experiments/FIGS/data/v2/curves/full_budget_long_leg2.json :: cells.<run>.N1000000.s1.jc1.per_class.label_Zqq: rejection of R16_Q1_MASS over R16_Q1, run by run :: sha256 588e935996d1006a
\newcommand{\MlstFtJcOneRejZqqCoarseMassOverCoarsePr}{equal within 7\%}  % experiments/FIGS/data/v2/curves/full_budget_long_leg2.json :: cells.<run>.N1000000.s1.jc1.per_class.label_Zqq: rejection of R16_Q1_MASS over R16_Q1, run by run :: sha256 588e935996d1006a :: prose twin: two significant figures, direction in words
\newcommand{\MlstFtJcOneRejZqqCoarseMassOverCoarsePrMax}{equal within 7\%}  % experiments/FIGS/data/v2/curves/full_budget_long_leg2.json :: cells.<run>.N1000000.s1.jc1.per_class.label_Zqq: rejection of R16_Q1_MASS over R16_Q1, run by run :: sha256 588e935996d1006a :: prose twin of the run farther from 1
\newcommand{\MlstFtJcOneRejZqqRandOverFull}{0.952\,\ensuremath{\pm}\,0.028}  % experiments/FIGS/data/v2/curves/full_budget_long_leg2.json :: per_class.label_Zqq: rejection of each random partition over the 188-class run of its seed :: sha256 588e935996d1006a
\newcommand{\MlstFtJcOneRejZqqRandOverFullPr}{equal within 5\%}  % experiments/FIGS/data/v2/curves/full_budget_long_leg2.json :: per_class.label_Zqq: rejection of each random partition over the 188-class run of its seed :: sha256 588e935996d1006a :: prose twin: two significant figures, direction in words
\newcommand{\MlstFtJcOneRejZqqRandOverFullRange}{0.93--0.99}  % experiments/FIGS/data/v2/curves/full_budget_long_leg2.json :: per_class.label_Zqq: smallest and largest of the partitions' ratios :: sha256 588e935996d1006a
\newcommand{\MlstFtJcOneResNRejDrop}{zero}  % experiments/FIGS/data/v2/curves/full_budget_long_leg2.json :: per_class rejections :: sha256 588e935996d1006a :: classes whose paired rejection ratio to the 188-class model departs from 1 by more than twice its SD and its counting uncertainty, columns with fewer than JC1_MIN_BKG passing jets excluded, ratio below 1 (count)
\newcommand{\MlstFtJcOneResNRejDiff}{zero}  % experiments/FIGS/data/v2/curves/full_budget_long_leg2.json :: per_class rejections :: sha256 588e935996d1006a :: classes whose paired rejection ratio to the 188-class model departs from 1 by more than twice its SD and its counting uncertainty, columns with fewer than JC1_MIN_BKG passing jets excluded (count)
\newcommand{\MlstFtJcOneResRejDropClasses}{none}  % experiments/FIGS/data/v2/curves/full_budget_long_leg2.json :: per_class rejections :: sha256 588e935996d1006a :: the classes of ...NRejDrop, in ParT's column order
\newcommand{\MlstFtJcOneMidNRejDrop}{zero}  % experiments/FIGS/data/v2/curves/full_budget_long_leg2.json :: per_class rejections :: sha256 588e935996d1006a :: classes whose paired rejection ratio to the 188-class model departs from 1 by more than twice its SD and its counting uncertainty, columns with fewer than JC1_MIN_BKG passing jets excluded, ratio below 1 (count)
\newcommand{\MlstFtJcOneMidNRejDiff}{zero}  % experiments/FIGS/data/v2/curves/full_budget_long_leg2.json :: per_class rejections :: sha256 588e935996d1006a :: classes whose paired rejection ratio to the 188-class model departs from 1 by more than twice its SD and its counting uncertainty, columns with fewer than JC1_MIN_BKG passing jets excluded (count)
\newcommand{\MlstFtJcOneMidRejDropClasses}{none}  % experiments/FIGS/data/v2/curves/full_budget_long_leg2.json :: per_class rejections :: sha256 588e935996d1006a :: the classes of ...NRejDrop, in ParT's column order
\newcommand{\MlstFtJcOneCoarseNRejDrop}{zero}  % experiments/FIGS/data/v2/curves/full_budget_long_leg2.json :: per_class rejections :: sha256 588e935996d1006a :: classes whose paired rejection ratio to the 188-class model departs from 1 by more than twice its SD and its counting uncertainty, columns with fewer than JC1_MIN_BKG passing jets excluded, ratio below 1 (count)
\newcommand{\MlstFtJcOneCoarseNRejDiff}{zero}  % experiments/FIGS/data/v2/curves/full_budget_long_leg2.json :: per_class rejections :: sha256 588e935996d1006a :: classes whose paired rejection ratio to the 188-class model departs from 1 by more than twice its SD and its counting uncertainty, columns with fewer than JC1_MIN_BKG passing jets excluded (count)
\newcommand{\MlstFtJcOneCoarseRejDropClasses}{none}  % experiments/FIGS/data/v2/curves/full_budget_long_leg2.json :: per_class rejections :: sha256 588e935996d1006a :: the classes of ...NRejDrop, in ParT's column order
\newcommand{\MlstFtJcOneAccRandMinusCoarse}{+0.002}  % experiments/FIGS/data/v2/curves/full_budget_long_leg2.json :: cells.*.N1000000.s1.jc1.accuracy :: sha256 588e935996d1006a :: mean over the random partitions minus the mean over the 17-class runs
\newcommand{\MlstFtJcOneAccFullBoot}{0.847}  % experiments/FIGS/data/v2/budget/jc1_acc_boot_N1000000.json :: diffs.r16q1-l188.accuracy_ref (mean) :: sha256 7cccb7b0ff900fb8
\newcommand{\MlstFtJcOneAccCoarseBoot}{0.842}  % experiments/FIGS/data/v2/budget/jc1_acc_boot_N1000000.json :: diffs.r16q1-l188.accuracy_config (mean) :: sha256 7cccb7b0ff900fb8
\newcommand{\MlstFtJcOneErrResOverFullBoot}{consistent}  % experiments/FIGS/data/v2/budget/jc1_acc_boot_N1000000.json :: diffs.l162-l188.error_ratio :: sha256 7cccb7b0ff900fb8 :: delta and sigma = hypot(run SD, bootstrap SD) in per cent; 'consistent' when |delta| <= 3 sigma
\newcommand{\MlstFtJcOneErrResOverFullZ}{0.6}  % experiments/FIGS/data/v2/budget/jc1_acc_boot_N1000000.json :: diffs.l162-l188.error_ratio.z :: sha256 7cccb7b0ff900fb8 :: |delta| / sigma
\newcommand{\MlstFtJcOneErrMidOverFullBoot}{consistent}  % experiments/FIGS/data/v2/budget/jc1_acc_boot_N1000000.json :: diffs.r42q1-l188.error_ratio :: sha256 7cccb7b0ff900fb8 :: delta and sigma = hypot(run SD, bootstrap SD) in per cent; 'consistent' when |delta| <= 3 sigma
\newcommand{\MlstFtJcOneErrMidOverFullZ}{1.4}  % experiments/FIGS/data/v2/budget/jc1_acc_boot_N1000000.json :: diffs.r42q1-l188.error_ratio.z :: sha256 7cccb7b0ff900fb8 :: |delta| / sigma
\newcommand{\MlstFtJcOneErrCoarseOverFullBoot}{(3.0\,\ensuremath{\pm}\,0.6)\% more errors}  % experiments/FIGS/data/v2/budget/jc1_acc_boot_N1000000.json :: diffs.r16q1-l188.error_ratio :: sha256 7cccb7b0ff900fb8 :: delta and sigma = hypot(run SD, bootstrap SD) in per cent; 'consistent' when |delta| <= 3 sigma
\newcommand{\MlstFtJcOneErrCoarseOverFullZ}{4.8}  % experiments/FIGS/data/v2/budget/jc1_acc_boot_N1000000.json :: diffs.r16q1-l188.error_ratio.z :: sha256 7cccb7b0ff900fb8 :: |delta| / sigma
\newcommand{\MlstFtJcOneBootReplicas}{200}  % experiments/FIGS/data/v2/budget/jc1_acc_boot_N1000000.json :: replicas :: sha256 7cccb7b0ff900fb8
\newcommand{\MlstFtJcOneBootJets}{2{,}000{,}000}  % experiments/FIGS/data/v2/budget/jc1_acc_boot_N1000000.json :: n_test_jets :: sha256 7cccb7b0ff900fb8
\newcommand{\ComputeTotalGpuHours}{3{,}435}  % experiments/FIGS/data/v2/index.json :: GPU-hours of every reported pretraining run, training and validation passes :: sha256 c6c4138ee5c86423
\newcommand{\MlstComputeTotalGpuHours}{3{,}510}  % experiments/FIGS/data/v2/index.json :: GPU-hours of every pretraining run the paper reports, the kinematic-only model included :: sha256 c6c4138ee5c86423
\newcommand{\ComputeMainGpu}{V100}  % experiments/FIGS/data/v2/index.json :: GPU type with the most hours over the four label sets :: sha256 c6c4138ee5c86423
\newcommand{\ComputeMainGpuShareMin}{56\%}  % experiments/FIGS/data/v2/index.json :: lowest share of a label set's GPU-hours on that type :: sha256 c6c4138ee5c86423
\newcommand{\ComputeMainGpuShareMax}{71\%}  % experiments/FIGS/data/v2/index.json :: highest share of a label set's GPU-hours on that type :: sha256 c6c4138ee5c86423
\newcommand{\ComputeRunsDone}{27}  % experiments/FIGS/data/v2/index.json :: paper runs with a compute record (count) :: sha256 c6c4138ee5c86423
\newcommand{\ComputeGpuTypes}{L4, RTX 3090, RTX 4090 and V100}  % experiments/FIGS/data/v2/index.json :: GPU types any finished paper run trained on :: sha256 c6c4138ee5c86423
\newcommand{\MlstComputeGpuTypes}{L4, RTX 3090, RTX 4090 and V100}  % experiments/FIGS/data/v2/index.json :: GPU types of the pretraining runs the paper reports :: sha256 c6c4138ee5c86423
\newcommand{\MlstFtJcOneErrCoarseOverFullBootEFour}{\pending{the JetClass bootstrap at 1e4 (after the equal-stopping read-out)}}  % experiments/FIGS/data/v2/index.json :: not computed yet :: sha256 c6c4138ee5c86423 :: waits on the JetClass bootstrap at 1e4 (after the equal-stopping read-out)
\newcommand{\MlstLpBbQgZeroAccDropFull}{\pending{pair probes with the quark/gluon inputs zeroed (open item O10, not run)}}  % experiments/FIGS/data/v2/index.json :: not computed yet :: sha256 c6c4138ee5c86423 :: waits on pair probes with the quark/gluon inputs zeroed (open item O10, not run)
\newcommand{\MlstLpBbQgZeroAccDropCoarse}{\pending{pair probes with the quark/gluon inputs zeroed (open item O10, not run)}}  % experiments/FIGS/data/v2/index.json :: not computed yet :: sha256 c6c4138ee5c86423 :: waits on pair probes with the quark/gluon inputs zeroed (open item O10, not run)
\newcommand{\XbsRthirtyFullOverScratch}{\pending{X->bs fits and from-scratch runs}}  % experiments/FIGS/data/v2/index.json :: not computed yet :: sha256 c6c4138ee5c86423 :: waits on X->bs fits and from-scratch runs
\newcommand{\XbsRthirtyFullOverScratchPr}{\pending{X->bs fits and from-scratch runs}}  % experiments/FIGS/data/v2/index.json :: not computed yet :: sha256 c6c4138ee5c86423 :: waits on X->bs fits and from-scratch runs
\newcommand{\XbsRthirtyCoarseOverFull}{\pending{X->bs fits}}  % experiments/FIGS/data/v2/index.json :: not computed yet :: sha256 c6c4138ee5c86423 :: waits on X->bs fits
\newcommand{\XbsRthirtyCoarseOverFullPr}{\pending{X->bs fits}}  % experiments/FIGS/data/v2/index.json :: not computed yet :: sha256 c6c4138ee5c86423 :: waits on X->bs fits
\newcommand{\XbsRthirtyNodeOverTransferFull}{\pending{X->bs fits}}  % experiments/FIGS/data/v2/index.json :: not computed yet :: sha256 c6c4138ee5c86423 :: waits on X->bs fits
\newcommand{\XbsRthirtyNodeOverTransferFullPr}{\pending{X->bs fits}}  % experiments/FIGS/data/v2/index.json :: not computed yet :: sha256 c6c4138ee5c86423 :: waits on X->bs fits
\newcommand{\XbbRthirtyDirectOverTransferFull}{\pending{X->bb fits}}  % experiments/FIGS/data/v2/index.json :: not computed yet :: sha256 c6c4138ee5c86423 :: waits on X->bb fits
\newcommand{\XbbRthirtyDirectOverTransferFullPr}{\pending{X->bb fits}}  % experiments/FIGS/data/v2/index.json :: not computed yet :: sha256 c6c4138ee5c86423 :: waits on X->bb fits
\newcommand{\XbbRthirtyTransferCoarseOverFull}{\pending{X->bb fits}}  % experiments/FIGS/data/v2/index.json :: not computed yet :: sha256 c6c4138ee5c86423 :: waits on X->bb fits
\newcommand{\XbbRthirtyTransferCoarseOverFullPr}{\pending{X->bb fits}}  % experiments/FIGS/data/v2/index.json :: not computed yet :: sha256 c6c4138ee5c86423 :: waits on X->bb fits
             % the counts below are its macros
% Visual language: Sophon's Fig. 1 (arXiv:2405.12972): the document's serif font, light fills,
% rounded boxes, light grey panels behind the stages (PI 2026-11-02). The text is set in the
% document's font at the printed sizes 8.5-10 pt (MLST 8-12 pt), never scaled.
% Text width: 15 cm (MLST full width). The picture is 15.0 cm wide and is drawn at its final
% size, never scaled, so the font sizes hold.
% Optional: \def\FoneDropXbx{} before % Figure 1 (F1 of docs/FIGURE_PLAN.md): the design of the study, in TikZ.
% \input{fig_f1} from main.tex; this file holds one figure environment and nothing else.
%
% WHAT main.tex MUST LOAD (this file loads nothing itself):
%   \usepackage{tikz}
%   \usetikzlibrary{arrows.meta}          % the only TikZ library used (arrow tips)
%   \definecolor{c188}{HTML}{1F77B4}      % the configuration colours of experiments/FIGS/
%   \definecolor{c162}{HTML}{FF7F0E}      % mlst_style.py (matplotlib's tab10 in order, as in
%   \definecolor{c43}{HTML}{2CA02C}       % Sophon's figures); the variants are drawn in their
%   \definecolor{c17}{HTML}{D62728}       % label set's colour with a dotted border, as the
%   \definecolor{crand}{HTML}{17BECF}     % plotted figures draw variants dotted
%   \definecolor{cflav}{HTML}{9467BD}
%   \input{results_generated}             % the counts below are its macros
% Visual language: Sophon's Fig. 1 (arXiv:2405.12972): the document's serif font, light fills,
% rounded boxes, light grey panels behind the stages (PI 2026-11-02). The text is set in the
% document's font at the printed sizes 8.5-10 pt (MLST 8-12 pt), never scaled.
% Text width: 15 cm (MLST full width). The picture is 15.0 cm wide and is drawn at its final
% size, never scaled, so the font sizes hold.
% Optional: \def\FoneDropXbx{} before \input{fig_f1} removes X->bb/X->bs from panel (c)
% (PI 2026-11-02: if the samples are not staged by 12 Nov, the X->bb/X->bs section goes).
%
% Every class box is drawn from the label maps (checked by hand, 2026-11-02):
%   (a) configs/labelmaps/rung_label_maps.v1.csv, columns L188, L162, R42_Q1, R16_Q1. Rows: the ten
%       decays X->QQ/gg in the order bb bc bq | cc cs cq | ss sq qq gg (the 43-class b / c / light
%       groups are then contiguous), then ee mm | tauh-taue tauh-taum, then the 27 QCD classes.
%   (b) column: configs/labelmaps/rand_label_map.v2.csv, RAND2_p1 on the same rows: groups
%       0 0 0 0 0 9 9 9 0 3 | 11 11 | 15 4 | QCD, lettered A..F in order of first appearance.
%       strips: configs/labelmaps/flavour_pair_map.v2.csv, the 22 classes of orbit X_YY_QQQQ.
%       Order with a b quark: bbbb bbcc bbss bbqq qqqb bbcq ccbs ccbq qqbc qqbs bcsq;
%       without: cccc ccss ccqq ssss ssqq qqqq qqqc qqqs ccsq sscq qqcs.
%       FLAV_F0: all 22 in group 9. FLAV_F1: the 11 with a b in group 11, the rest 9.
%       FLAV_F1R: group 11 = bbss bbqq qqqb bbcq ccbq qqbc (with b, positions 2-5, 7, 8 from 0)
%       and ccss ccqq ssss qqqc sscq (without, positions 1-3, 6, 9); filled square = group 11.
% Geometry (cm): row height 0.35 from y = -1.18, gaps of 0.14 between the quark/gluon, lepton
% and QCD blocks, the QCD row 0.70; boxes inset 0.025 top and bottom; columns 1.45 wide with
% left edges 1.75, 3.28, 4.81, 6.34 (hierarchy) and 8.25 (random partition).
\providecommand{\Kfull}{\VocabSizeLoneeighteight}
\providecommand{\Kres}{\VocabSizeLonesixtwo}
\providecommand{\Kmid}{\VocabSizeRfourtwoQone}
\providecommand{\Kcoarse}{\VocabSizeRonesixQone}
\begin{figure}[tbp]
\centering
\ifdim\linewidth<14.99cm\typeout{fig_f1.tex: WARNING: the drawing is 15 cm wide but \string\linewidth\space
  is \the\linewidth; set a 15 cm text width (it is not scaled, to keep its 8.5 pt text)}\fi
\def\FoneText{\rmfamily\fontsize{8.5}{9.5}\selectfont}%
\def\FoneMain{\rmfamily\fontsize{9}{10.5}\selectfont}%
\def\FoneBold{\rmfamily\bfseries\fontsize{9}{10.5}\selectfont}%
\def\FoneTitle{\rmfamily\bfseries\fontsize{10}{12}\selectfont}%
\ifdefined\FoneDropXbx\def\FoneXbx{}\else\def\FoneXbx{\\$X\to b\bar b$, $X\to bs$}\fi
\begin{tikzpicture}[
  x=1cm, y=1cm,
  font=\FoneMain,
  ttl/.style={anchor=base west, font=\FoneTitle, inner sep=0pt},
  cls/.style={draw=black!55, line width=0.45pt, rounded corners=1.2pt},
  hdr/.style={draw=#1!70!black, fill=#1!22, line width=0.6pt, rounded corners=2.5pt, align=center,
              font=\FoneText, inner sep=0pt},
  inbox/.style={align=center, font=\FoneText, inner sep=0pt},
  letter/.style={font=\FoneBold, inner sep=0pt},
  rowlab/.style={anchor=east, inner sep=0pt},
  txt/.style={align=left, font=\FoneText, inner sep=0pt},
  sq/.style={draw=black!65, line width=0.4pt},
  chip/.style={draw=#1!70!black, fill=#1!22, line width=0.6pt, rounded corners=3pt,
               align=center, font=\FoneText, inner sep=0pt, anchor=north west},
  stage/.style={draw=black!45, fill=black!4, line width=0.6pt, rounded corners=3pt, align=left,
                font=\FoneText, inner sep=0.1cm},
  flow/.style={-{Latex[length=1.8mm, width=1.4mm]}, line width=0.7pt, black!75},
]
% ================================================================ (a) the class hierarchy
\node[ttl] at (0,0) {(a) Class hierarchy};
% faint bands on every other row, drawn first so the boxes sit on top
\foreach \rowtop in {-1.53,-2.23,-2.93,-3.63,-4.33,-5.17,-5.87}
  \fill[black!4] (0,\rowtop) rectangle (9.70,\rowtop-0.35);
% column headers
\node[rowlab, align=right, font=\FoneText] at (1.65,-0.68) {JetClass-II\\class};
\node[hdr=c188, minimum width=1.45cm, minimum height=0.80cm, anchor=north west] at (1.75,-0.28) {\Kfull\\classes};
\node[hdr=c162, minimum width=1.45cm, minimum height=0.80cm, anchor=north west] at (3.28,-0.28) {\Kres\\classes};
\node[hdr=c43,  minimum width=1.45cm, minimum height=0.80cm, anchor=north west] at (4.81,-0.28) {\Kmid\\classes};
\node[hdr=c17,  minimum width=1.45cm, minimum height=0.80cm, anchor=north west] at (6.34,-0.28) {\Kcoarse\\classes};
% row labels
\node[rowlab] at (1.65,-1.355) {$X\to b\bar b$};
\node[rowlab] at (1.65,-1.705) {$X\to bc$};
\node[rowlab] at (1.65,-2.055) {$X\to bq$};
\node[rowlab] at (1.65,-2.405) {$X\to c\bar c$};
\node[rowlab] at (1.65,-2.755) {$X\to cs$};
\node[rowlab] at (1.65,-3.105) {$X\to cq$};
\node[rowlab] at (1.65,-3.455) {$X\to s\bar s$};
\node[rowlab] at (1.65,-3.805) {$X\to sq$};
\node[rowlab] at (1.65,-4.155) {$X\to q\bar q$};
\node[rowlab] at (1.65,-4.505) {$X\to gg$};
\node[rowlab] at (1.65,-4.995) {$X\to ee$};
\node[rowlab] at (1.65,-5.345) {$X\to\mu\mu$};
\node[rowlab] at (1.65,-5.695) {$X\to\tau_h\tau_e$};
\node[rowlab] at (1.65,-6.045) {$X\to\tau_h\tau_\mu$};
\node[rowlab] at (1.65,-6.710) {QCD jets};
% 188 and 162 classes: every resonant decay is its own class
\foreach \rowtop in {-1.18,-1.53,-1.88,-2.23,-2.58,-2.93,-3.28,-3.63,-3.98,-4.33,-4.82,-5.17,-5.52,-5.87} {
  \draw[cls, fill=c188!15] (1.75,\rowtop-0.025) rectangle (3.20,\rowtop-0.325);
  \draw[cls, fill=c162!15] (3.28,\rowtop-0.025) rectangle (4.73,\rowtop-0.325);
}
\draw[cls, fill=c188!15] (1.75,-6.385) rectangle (3.20,-7.035);
\node[inbox] at (2.475,-6.71) {\LitJetClassTwoNQcdClasses\\classes};
\draw[cls, fill=c162!15] (3.28,-6.385) rectangle (4.73,-7.035);
\node[inbox] at (4.005,-6.71) {1 class};
% 43 classes: decay topology x {b, else c, else light quark or gluon}; e and mu merged
\draw[cls, fill=c43!15] (4.81,-1.205) rectangle (6.26,-2.205);
\node[inbox] at (5.535,-1.705) {with $b$};
\draw[cls, fill=c43!15] (4.81,-2.255) rectangle (6.26,-3.255);
\node[inbox] at (5.535,-2.755) {with $c$,\\no $b$};
\draw[cls, fill=c43!15] (4.81,-3.305) rectangle (6.26,-4.655);
\node[inbox] at (5.535,-3.980) {$u$, $d$, $s$\\or $g$};
\draw[cls, fill=c43!15] (4.81,-4.845) rectangle (6.26,-5.495);
\node[inbox] at (5.535,-5.170) {$\ell\ell$};
\draw[cls, fill=c43!15] (4.81,-5.545) rectangle (6.26,-6.195);
\node[inbox] at (5.535,-5.870) {$\tau_h\tau_\ell$};
\draw[cls, fill=c43!15] (4.81,-6.385) rectangle (6.26,-7.035);
\node[inbox] at (5.535,-6.710) {1 class};
% 17 classes: number and type of prongs, no flavor
\draw[cls, fill=c17!15] (6.34,-1.205) rectangle (7.79,-4.655);
\node[inbox] at (7.065,-2.930) {two\\quarks\\or gluons};
\draw[cls, fill=c17!15] (6.34,-4.845) rectangle (7.79,-5.495);
\node[inbox] at (7.065,-5.170) {$\ell\ell$};
\draw[cls, fill=c17!15] (6.34,-5.545) rectangle (7.79,-6.195);
\node[inbox] at (7.065,-5.870) {$\tau_h\tau_\ell$};
\draw[cls, fill=c17!15] (6.34,-6.385) rectangle (7.79,-7.035);
\node[inbox] at (7.065,-6.710) {1 class};
% what each merge removes
\node[anchor=north west, font=\FoneBold, inner sep=0pt] at (0,-7.27) {What each merge removes};
\node[txt, anchor=north west, text width=6.04cm] (m1) at (1.75,-7.65)
  {the $b$, $c$ and $s$ content of the QCD jets};
\node[txt, anchor=north west, text width=6.04cm] (m2) at ([yshift=-0.10cm]m1.south west)
  {the number of $b$ and $c$ quarks, light-quark flavor, quark against gluon, and lepton
   flavor ($e$ against $\mu$)};
\node[txt, anchor=north west, text width=6.04cm] (m3) at ([yshift=-0.10cm]m2.south west)
  {quark flavor ($b$, $c$ or light); the number and type of prongs remain};
\node[txt, anchor=north east] at ([xshift=-0.10cm]m1.north west) {\Kfull\,$\to$\,\Kres};
\node[txt, anchor=north east] at ([xshift=-0.10cm]m2.north west) {\Kres\,$\to$\,\Kmid};
\node[txt, anchor=north east] at ([xshift=-0.10cm]m3.north west) {\Kmid\,$\to$\,\Kcoarse};
% ================================================================ (b) the controls
\node[ttl] at (8.25,0) {(b) Controls, \Kcoarse{} classes each};
% the first random partition on the rows of (a); one letter = one class
\node[hdr=crand, minimum width=1.45cm, minimum height=0.80cm, anchor=north west] at (8.25,-0.28) {random\\partition};
\draw[cls, fill=crand!15] (8.25,-1.205) rectangle (9.70,-2.905);
\node[letter] at (8.975,-2.055) {A};
\draw[cls, fill=crand!15] (8.25,-2.955) rectangle (9.70,-3.955);
\node[letter] at (8.975,-3.455) {B};
\draw[cls, fill=crand!15] (8.25,-4.005) rectangle (9.70,-4.305);
\node[letter] at (8.975,-4.155) {A};
\draw[cls, fill=crand!15] (8.25,-4.355) rectangle (9.70,-4.655);
\node[letter] at (8.975,-4.505) {C};
\draw[cls, fill=crand!15] (8.25,-4.845) rectangle (9.70,-5.495);
\node[letter] at (8.975,-5.170) {D};
\draw[cls, fill=crand!15] (8.25,-5.545) rectangle (9.70,-5.845);
\node[letter] at (8.975,-5.695) {E};
\draw[cls, fill=crand!15] (8.25,-5.895) rectangle (9.70,-6.195);
\node[letter] at (8.975,-6.045) {F};
\draw[cls, fill=crand!15] (8.25,-6.385) rectangle (9.70,-7.035);
\node[inbox] at (8.975,-6.710) {1 class};
% how the random partitions are built
\node[hdr=crand, anchor=north west, text width=4.75cm, inner sep=0.1cm, align=left, font=\FoneBold]
  (rh) at (10.05,-0.28) {Random partitions (\RandVtwoNPartitions)};
\node[txt, anchor=north west, text width=4.95cm] (rb) at ([yshift=-0.12cm]rh.south west)
  {\textbullet\ groups of resonant classes drawn at random, plus one class of QCD jets\\[2pt]
   \textbullet\ group shares of training jets within \RandVtwoShareTol{} of the \Kcoarse-class groups\\[2pt]
   \textbullet\ best of \RandVtwoTrials{} draws: the fewest merges shared with the \Kcoarse-class level\\[2pt]
   \textbullet\ left: the first partition; boxes with the same letter are one class};
% how the random partitions with controlled flavor boundaries are built
\node[hdr=cflav, anchor=north west, text width=4.75cm, inner sep=0.1cm, align=left, font=\FoneBold]
  (fh) at ([yshift=-0.22cm]rb.south west) {Random partitions with controlled flavor boundaries};
\node[txt, anchor=north west, text width=4.95cm] (fn) at ([yshift=-0.12cm]fh.south west)
  {whole families of decays that differ only in quark flavor; two of them split one family,
   the \RandFlavSplitOrbitClasses{} decays \RandFlavSplitOrbit:};
\begin{scope}[shift={(fn.south west)}]
  % brackets over the decays with and without a b quark (squares at a 0.2 cm pitch)
  \node[txt, anchor=south] at (1.075,-0.45) {with $b$};
  \node[txt, anchor=south] at (3.425,-0.45) {without $b$};
  \draw[line width=0.5pt] (0,-0.54) -- (0,-0.47) -- (2.15,-0.47) -- (2.15,-0.54);
  \draw[line width=0.5pt] (2.35,-0.54) -- (2.35,-0.47) -- (4.50,-0.47) -- (4.50,-0.54);
  % no flavor boundary: all 22 in one class
  \node[txt, anchor=north west] at (0,-0.62) {no flavor boundary};
  \foreach \i in {0,...,10} {
    \draw[sq, fill=white] (0.2*\i,-1.10) rectangle ++(0.15,0.15);
    \draw[sq, fill=white] (2.35+0.2*\i,-1.10) rectangle ++(0.15,0.15);
  }
  % b-quark boundary: the decays with a b quark moved to another class
  \node[txt, anchor=north west] at (0,-1.22) {$b$-quark boundary};
  \foreach \i in {0,...,10} {
    \draw[sq, fill=cflav!70!black] (0.2*\i,-1.70) rectangle ++(0.15,0.15);
    \draw[sq, fill=white] (2.35+0.2*\i,-1.70) rectangle ++(0.15,0.15);
  }
  % random flavor boundary: the same number of decays, drawn at random
  \node[txt, anchor=north west] at (0,-1.82) {random flavor boundary};
  \foreach \i in {0,1,6,9,10}
    \draw[sq, fill=white] (0.2*\i,-2.30) rectangle ++(0.15,0.15);
  \foreach \i in {2,3,4,5,7,8}
    \draw[sq, fill=cflav!70!black] (0.2*\i,-2.30) rectangle ++(0.15,0.15);
  \foreach \i in {0,4,5,7,8,10}
    \draw[sq, fill=white] (2.35+0.2*\i,-2.30) rectangle ++(0.15,0.15);
  \foreach \i in {1,2,3,6,9}
    \draw[sq, fill=cflav!70!black] (2.35+0.2*\i,-2.30) rectangle ++(0.15,0.15);
  % key
  \draw[sq, fill=cflav!70!black] (0,-2.62) rectangle ++(0.15,0.15);
  \node[txt, anchor=west] (fkey) at (0.25,-2.545) {moved to another class};
\end{scope}
% ================================================================ (c) the pipeline
% (c) starts below the lower of the two columns above, whatever their text heights
\pgfpointanchor{m3}{south}\pgfgetlastxy{\FoneXa}{\FoneYa}%
\pgfpointanchor{fkey}{south}\pgfgetlastxy{\FoneXb}{\FoneYb}%
\pgfmathsetlengthmacro{\FoneYc}{min(\FoneYa,\FoneYb)-0.45cm}%
\coordinate (ctop) at (0pt,\FoneYc);
\begin{scope}[shift={(ctop)}]
  \node[ttl, anchor=north west] at (0,0) {(c) Pretraining and evaluation};
  % what every model of one seed shares
  \node[stage, text width=2.15cm] (seed) at (1.175,-1.96)
    {\textbf{Same for every model of one seed:} the training jets, their order and the
     initial weights, except the output layer};
  % one pretraining per label set: the configuration colours of every later figure
  \node[draw=black!45, fill=black!4, line width=0.6pt, rounded corners=3pt, minimum width=5.15cm,
        minimum height=3.08cm, anchor=north west, inner sep=0pt] (pre) at (2.81,-0.42) {};
  \node[anchor=north west, font=\FoneBold, inner sep=0pt] at (2.96,-0.52) {Pretraining on JetClass-II};
  \node[chip=c188, minimum width=1.13cm, minimum height=0.42cm] at (2.96,-0.95) {\Kfull};
  \node[chip=c162, minimum width=1.13cm, minimum height=0.42cm] at (4.20,-0.95) {\Kres};
  \node[chip=c43,  minimum width=1.13cm, minimum height=0.42cm] at (5.44,-0.95) {\Kmid};
  \node[chip=c17,  minimum width=1.13cm, minimum height=0.42cm] at (6.68,-0.95) {\Kcoarse};
  \node[chip=crand, minimum width=1.70cm, minimum height=0.70cm] at (2.96,-1.47) {random\\partitions};
  \node[chip=cflav, minimum width=3.04cm, minimum height=0.70cm] at (4.77,-1.47) {controlled flavor\\boundaries};
  \node[chip=c162, densely dotted, minimum width=1.13cm, minimum height=0.42cm] at (2.96,-2.32) {\Kres};
  \node[chip=c17,  densely dotted, minimum width=1.13cm, minimum height=0.42cm] at (4.20,-2.32) {\Kcoarse};
  \node[txt, anchor=west] at (5.44,-2.53) {with auxiliary\\mass regression};
  \node[chip=c188, densely dotted, minimum width=1.13cm, minimum height=0.42cm] at (2.96,-2.94) {\Kfull};
  \node[txt, anchor=west] at (4.20,-3.15) {kinematic inputs only};
  % evaluation
  \node[stage, text width=2.4cm] (fz) at (9.68,-0.95) {frozen features:\\linear probe,\\MLP probe};
  \node[stage, text width=2.4cm] (ft) at (9.68,-2.05) {fine-tuning\\of all weights};
  \node[stage, text width=2.4cm, draw=black, fill=white, line width=1.1pt] (sc) at (9.68,-3.15) {trained from\\scratch};
  \node[stage, text width=3.4cm, minimum height=3.2cm] (task) at (13.20,-2.00)
    {\textbf{Downstream tasks}\\JetClass-II: all classes\\JetClass-II: class pairs\\\quad(both class-balanced)\\JetClass\\
     top tagging\\quark/gluon, Pythia\\quark/gluon, Herwig\\jet mass\FoneXbx};
  \draw[flow] (seed.east) -- (seed.east -| pre.west);
  \draw[flow] (fz.west -| pre.east) -- (fz.west);
  \draw[flow] (ft.west -| pre.east) -- (ft.west);
  \draw[flow] (fz.east) -- (fz.east -| task.west);
  \draw[flow] (ft.east) -- (ft.east -| task.west);
  \draw[flow] (sc.east) -- (sc.east -| task.west);
\end{scope}
\end{tikzpicture}
\caption{Design of the study.
(a)~The class hierarchy on the decays of a heavy resonance~$X$ and on the jets from quantum
chromodynamics (QCD); $q$ is a $u$ or $d$ quark, $\ell$ an electron or a muon, and $\tau_h$,
$\tau_e$, $\tau_\mu$ a $\tau$ lepton decaying to hadrons, an electron or a muon. Each box is one
class of the label set above its column, except the QCD box of the \Kfull-class column (Sophon's
label set~\cite{sophon}), which stands for \LitJetClassTwoNQcdClasses{} classes. Each level merges
classes of the level to its left; the lines below name what each merge removes.
(b)~The controls: the first random partition, where boxes with the same letter are one class, and
the random partitions with controlled flavor boundaries on the decays \RandFlavSplitOrbit{} (one
square per decay; filled squares are moved to a second class).
(c)~Models of one seed share the training jets, their order and the initial weights, except the
output layer, the mass regression target, and the first layer and initial weights of the model with
kinematic inputs only (the particle four-momenta). The hierarchy levels have \MlstNRunsWord{} seeds
each, the models with the auxiliary mass regression \MlstMassNRunsWord{}, the random partitions with
controlled flavor boundaries \RandFlavRuns{} and each random partition \RandVtwoRuns{}. Frozen
features, the output of the backbone (the network below the classification head) with its weights held
fixed, feed linear probes (logistic regressions) and multilayer perceptron (MLP) probes; fine-tuning
trains all weights. Colors identify the label sets and controls in every figure; the plots give the three partitions with
controlled flavor boundaries three shades; dotted borders mark the two
variants, which the other figures draw as dotted lines in the color of their label set.}
\label{fig:design}
\end{figure}
 removes X->bb/X->bs from panel (c)
% (PI 2026-11-02: if the samples are not staged by 12 Nov, the X->bb/X->bs section goes).
%
% Every class box is drawn from the label maps (checked by hand, 2026-11-02):
%   (a) configs/labelmaps/rung_label_maps.v1.csv, columns L188, L162, R42_Q1, R16_Q1. Rows: the ten
%       decays X->QQ/gg in the order bb bc bq | cc cs cq | ss sq qq gg (the 43-class b / c / light
%       groups are then contiguous), then ee mm | tauh-taue tauh-taum, then the 27 QCD classes.
%   (b) column: configs/labelmaps/rand_label_map.v2.csv, RAND2_p1 on the same rows: groups
%       0 0 0 0 0 9 9 9 0 3 | 11 11 | 15 4 | QCD, lettered A..F in order of first appearance.
%       strips: configs/labelmaps/flavour_pair_map.v2.csv, the 22 classes of orbit X_YY_QQQQ.
%       Order with a b quark: bbbb bbcc bbss bbqq qqqb bbcq ccbs ccbq qqbc qqbs bcsq;
%       without: cccc ccss ccqq ssss ssqq qqqq qqqc qqqs ccsq sscq qqcs.
%       FLAV_F0: all 22 in group 9. FLAV_F1: the 11 with a b in group 11, the rest 9.
%       FLAV_F1R: group 11 = bbss bbqq qqqb bbcq ccbq qqbc (with b, positions 2-5, 7, 8 from 0)
%       and ccss ccqq ssss qqqc sscq (without, positions 1-3, 6, 9); filled square = group 11.
% Geometry (cm): row height 0.35 from y = -1.18, gaps of 0.14 between the quark/gluon, lepton
% and QCD blocks, the QCD row 0.70; boxes inset 0.025 top and bottom; columns 1.45 wide with
% left edges 1.75, 3.28, 4.81, 6.34 (hierarchy) and 8.25 (random partition).
\providecommand{\Kfull}{\VocabSizeLoneeighteight}
\providecommand{\Kres}{\VocabSizeLonesixtwo}
\providecommand{\Kmid}{\VocabSizeRfourtwoQone}
\providecommand{\Kcoarse}{\VocabSizeRonesixQone}
\begin{figure}[tbp]
\centering
\ifdim\linewidth<14.99cm\typeout{fig_f1.tex: WARNING: the drawing is 15 cm wide but \string\linewidth\space
  is \the\linewidth; set a 15 cm text width (it is not scaled, to keep its 8.5 pt text)}\fi
\def\FoneText{\rmfamily\fontsize{8.5}{9.5}\selectfont}%
\def\FoneMain{\rmfamily\fontsize{9}{10.5}\selectfont}%
\def\FoneBold{\rmfamily\bfseries\fontsize{9}{10.5}\selectfont}%
\def\FoneTitle{\rmfamily\bfseries\fontsize{10}{12}\selectfont}%
\ifdefined\FoneDropXbx\def\FoneXbx{}\else\def\FoneXbx{\\$X\to b\bar b$, $X\to bs$}\fi
\begin{tikzpicture}[
  x=1cm, y=1cm,
  font=\FoneMain,
  ttl/.style={anchor=base west, font=\FoneTitle, inner sep=0pt},
  cls/.style={draw=black!55, line width=0.45pt, rounded corners=1.2pt},
  hdr/.style={draw=#1!70!black, fill=#1!22, line width=0.6pt, rounded corners=2.5pt, align=center,
              font=\FoneText, inner sep=0pt},
  inbox/.style={align=center, font=\FoneText, inner sep=0pt},
  letter/.style={font=\FoneBold, inner sep=0pt},
  rowlab/.style={anchor=east, inner sep=0pt},
  txt/.style={align=left, font=\FoneText, inner sep=0pt},
  sq/.style={draw=black!65, line width=0.4pt},
  chip/.style={draw=#1!70!black, fill=#1!22, line width=0.6pt, rounded corners=3pt,
               align=center, font=\FoneText, inner sep=0pt, anchor=north west},
  stage/.style={draw=black!45, fill=black!4, line width=0.6pt, rounded corners=3pt, align=left,
                font=\FoneText, inner sep=0.1cm},
  flow/.style={-{Latex[length=1.8mm, width=1.4mm]}, line width=0.7pt, black!75},
]
% ================================================================ (a) the class hierarchy
\node[ttl] at (0,0) {(a) Class hierarchy};
% faint bands on every other row, drawn first so the boxes sit on top
\foreach \rowtop in {-1.53,-2.23,-2.93,-3.63,-4.33,-5.17,-5.87}
  \fill[black!4] (0,\rowtop) rectangle (9.70,\rowtop-0.35);
% column headers
\node[rowlab, align=right, font=\FoneText] at (1.65,-0.68) {JetClass-II\\class};
\node[hdr=c188, minimum width=1.45cm, minimum height=0.80cm, anchor=north west] at (1.75,-0.28) {\Kfull\\classes};
\node[hdr=c162, minimum width=1.45cm, minimum height=0.80cm, anchor=north west] at (3.28,-0.28) {\Kres\\classes};
\node[hdr=c43,  minimum width=1.45cm, minimum height=0.80cm, anchor=north west] at (4.81,-0.28) {\Kmid\\classes};
\node[hdr=c17,  minimum width=1.45cm, minimum height=0.80cm, anchor=north west] at (6.34,-0.28) {\Kcoarse\\classes};
% row labels
\node[rowlab] at (1.65,-1.355) {$X\to b\bar b$};
\node[rowlab] at (1.65,-1.705) {$X\to bc$};
\node[rowlab] at (1.65,-2.055) {$X\to bq$};
\node[rowlab] at (1.65,-2.405) {$X\to c\bar c$};
\node[rowlab] at (1.65,-2.755) {$X\to cs$};
\node[rowlab] at (1.65,-3.105) {$X\to cq$};
\node[rowlab] at (1.65,-3.455) {$X\to s\bar s$};
\node[rowlab] at (1.65,-3.805) {$X\to sq$};
\node[rowlab] at (1.65,-4.155) {$X\to q\bar q$};
\node[rowlab] at (1.65,-4.505) {$X\to gg$};
\node[rowlab] at (1.65,-4.995) {$X\to ee$};
\node[rowlab] at (1.65,-5.345) {$X\to\mu\mu$};
\node[rowlab] at (1.65,-5.695) {$X\to\tau_h\tau_e$};
\node[rowlab] at (1.65,-6.045) {$X\to\tau_h\tau_\mu$};
\node[rowlab] at (1.65,-6.710) {QCD jets};
% 188 and 162 classes: every resonant decay is its own class
\foreach \rowtop in {-1.18,-1.53,-1.88,-2.23,-2.58,-2.93,-3.28,-3.63,-3.98,-4.33,-4.82,-5.17,-5.52,-5.87} {
  \draw[cls, fill=c188!15] (1.75,\rowtop-0.025) rectangle (3.20,\rowtop-0.325);
  \draw[cls, fill=c162!15] (3.28,\rowtop-0.025) rectangle (4.73,\rowtop-0.325);
}
\draw[cls, fill=c188!15] (1.75,-6.385) rectangle (3.20,-7.035);
\node[inbox] at (2.475,-6.71) {\LitJetClassTwoNQcdClasses\\classes};
\draw[cls, fill=c162!15] (3.28,-6.385) rectangle (4.73,-7.035);
\node[inbox] at (4.005,-6.71) {1 class};
% 43 classes: decay topology x {b, else c, else light quark or gluon}; e and mu merged
\draw[cls, fill=c43!15] (4.81,-1.205) rectangle (6.26,-2.205);
\node[inbox] at (5.535,-1.705) {with $b$};
\draw[cls, fill=c43!15] (4.81,-2.255) rectangle (6.26,-3.255);
\node[inbox] at (5.535,-2.755) {with $c$,\\no $b$};
\draw[cls, fill=c43!15] (4.81,-3.305) rectangle (6.26,-4.655);
\node[inbox] at (5.535,-3.980) {$u$, $d$, $s$\\or $g$};
\draw[cls, fill=c43!15] (4.81,-4.845) rectangle (6.26,-5.495);
\node[inbox] at (5.535,-5.170) {$\ell\ell$};
\draw[cls, fill=c43!15] (4.81,-5.545) rectangle (6.26,-6.195);
\node[inbox] at (5.535,-5.870) {$\tau_h\tau_\ell$};
\draw[cls, fill=c43!15] (4.81,-6.385) rectangle (6.26,-7.035);
\node[inbox] at (5.535,-6.710) {1 class};
% 17 classes: number and type of prongs, no flavor
\draw[cls, fill=c17!15] (6.34,-1.205) rectangle (7.79,-4.655);
\node[inbox] at (7.065,-2.930) {two\\quarks\\or gluons};
\draw[cls, fill=c17!15] (6.34,-4.845) rectangle (7.79,-5.495);
\node[inbox] at (7.065,-5.170) {$\ell\ell$};
\draw[cls, fill=c17!15] (6.34,-5.545) rectangle (7.79,-6.195);
\node[inbox] at (7.065,-5.870) {$\tau_h\tau_\ell$};
\draw[cls, fill=c17!15] (6.34,-6.385) rectangle (7.79,-7.035);
\node[inbox] at (7.065,-6.710) {1 class};
% what each merge removes
\node[anchor=north west, font=\FoneBold, inner sep=0pt] at (0,-7.27) {What each merge removes};
\node[txt, anchor=north west, text width=6.04cm] (m1) at (1.75,-7.65)
  {the $b$, $c$ and $s$ content of the QCD jets};
\node[txt, anchor=north west, text width=6.04cm] (m2) at ([yshift=-0.10cm]m1.south west)
  {the number of $b$ and $c$ quarks, light-quark flavor, quark against gluon, and lepton
   flavor ($e$ against $\mu$)};
\node[txt, anchor=north west, text width=6.04cm] (m3) at ([yshift=-0.10cm]m2.south west)
  {quark flavor ($b$, $c$ or light); the number and type of prongs remain};
\node[txt, anchor=north east] at ([xshift=-0.10cm]m1.north west) {\Kfull\,$\to$\,\Kres};
\node[txt, anchor=north east] at ([xshift=-0.10cm]m2.north west) {\Kres\,$\to$\,\Kmid};
\node[txt, anchor=north east] at ([xshift=-0.10cm]m3.north west) {\Kmid\,$\to$\,\Kcoarse};
% ================================================================ (b) the controls
\node[ttl] at (8.25,0) {(b) Controls, \Kcoarse{} classes each};
% the first random partition on the rows of (a); one letter = one class
\node[hdr=crand, minimum width=1.45cm, minimum height=0.80cm, anchor=north west] at (8.25,-0.28) {random\\partition};
\draw[cls, fill=crand!15] (8.25,-1.205) rectangle (9.70,-2.905);
\node[letter] at (8.975,-2.055) {A};
\draw[cls, fill=crand!15] (8.25,-2.955) rectangle (9.70,-3.955);
\node[letter] at (8.975,-3.455) {B};
\draw[cls, fill=crand!15] (8.25,-4.005) rectangle (9.70,-4.305);
\node[letter] at (8.975,-4.155) {A};
\draw[cls, fill=crand!15] (8.25,-4.355) rectangle (9.70,-4.655);
\node[letter] at (8.975,-4.505) {C};
\draw[cls, fill=crand!15] (8.25,-4.845) rectangle (9.70,-5.495);
\node[letter] at (8.975,-5.170) {D};
\draw[cls, fill=crand!15] (8.25,-5.545) rectangle (9.70,-5.845);
\node[letter] at (8.975,-5.695) {E};
\draw[cls, fill=crand!15] (8.25,-5.895) rectangle (9.70,-6.195);
\node[letter] at (8.975,-6.045) {F};
\draw[cls, fill=crand!15] (8.25,-6.385) rectangle (9.70,-7.035);
\node[inbox] at (8.975,-6.710) {1 class};
% how the random partitions are built
\node[hdr=crand, anchor=north west, text width=4.75cm, inner sep=0.1cm, align=left, font=\FoneBold]
  (rh) at (10.05,-0.28) {Random partitions (\RandVtwoNPartitions)};
\node[txt, anchor=north west, text width=4.95cm] (rb) at ([yshift=-0.12cm]rh.south west)
  {\textbullet\ groups of resonant classes drawn at random, plus one class of QCD jets\\[2pt]
   \textbullet\ group shares of training jets within \RandVtwoShareTol{} of the \Kcoarse-class groups\\[2pt]
   \textbullet\ best of \RandVtwoTrials{} draws: the fewest merges shared with the \Kcoarse-class level\\[2pt]
   \textbullet\ left: the first partition; boxes with the same letter are one class};
% how the random partitions with controlled flavor boundaries are built
\node[hdr=cflav, anchor=north west, text width=4.75cm, inner sep=0.1cm, align=left, font=\FoneBold]
  (fh) at ([yshift=-0.22cm]rb.south west) {Random partitions with controlled flavor boundaries};
\node[txt, anchor=north west, text width=4.95cm] (fn) at ([yshift=-0.12cm]fh.south west)
  {whole families of decays that differ only in quark flavor; two of them split one family,
   the \RandFlavSplitOrbitClasses{} decays \RandFlavSplitOrbit:};
\begin{scope}[shift={(fn.south west)}]
  % brackets over the decays with and without a b quark (squares at a 0.2 cm pitch)
  \node[txt, anchor=south] at (1.075,-0.45) {with $b$};
  \node[txt, anchor=south] at (3.425,-0.45) {without $b$};
  \draw[line width=0.5pt] (0,-0.54) -- (0,-0.47) -- (2.15,-0.47) -- (2.15,-0.54);
  \draw[line width=0.5pt] (2.35,-0.54) -- (2.35,-0.47) -- (4.50,-0.47) -- (4.50,-0.54);
  % no flavor boundary: all 22 in one class
  \node[txt, anchor=north west] at (0,-0.62) {no flavor boundary};
  \foreach \i in {0,...,10} {
    \draw[sq, fill=white] (0.2*\i,-1.10) rectangle ++(0.15,0.15);
    \draw[sq, fill=white] (2.35+0.2*\i,-1.10) rectangle ++(0.15,0.15);
  }
  % b-quark boundary: the decays with a b quark moved to another class
  \node[txt, anchor=north west] at (0,-1.22) {$b$-quark boundary};
  \foreach \i in {0,...,10} {
    \draw[sq, fill=cflav!70!black] (0.2*\i,-1.70) rectangle ++(0.15,0.15);
    \draw[sq, fill=white] (2.35+0.2*\i,-1.70) rectangle ++(0.15,0.15);
  }
  % random flavor boundary: the same number of decays, drawn at random
  \node[txt, anchor=north west] at (0,-1.82) {random flavor boundary};
  \foreach \i in {0,1,6,9,10}
    \draw[sq, fill=white] (0.2*\i,-2.30) rectangle ++(0.15,0.15);
  \foreach \i in {2,3,4,5,7,8}
    \draw[sq, fill=cflav!70!black] (0.2*\i,-2.30) rectangle ++(0.15,0.15);
  \foreach \i in {0,4,5,7,8,10}
    \draw[sq, fill=white] (2.35+0.2*\i,-2.30) rectangle ++(0.15,0.15);
  \foreach \i in {1,2,3,6,9}
    \draw[sq, fill=cflav!70!black] (2.35+0.2*\i,-2.30) rectangle ++(0.15,0.15);
  % key
  \draw[sq, fill=cflav!70!black] (0,-2.62) rectangle ++(0.15,0.15);
  \node[txt, anchor=west] (fkey) at (0.25,-2.545) {moved to another class};
\end{scope}
% ================================================================ (c) the pipeline
% (c) starts below the lower of the two columns above, whatever their text heights
\pgfpointanchor{m3}{south}\pgfgetlastxy{\FoneXa}{\FoneYa}%
\pgfpointanchor{fkey}{south}\pgfgetlastxy{\FoneXb}{\FoneYb}%
\pgfmathsetlengthmacro{\FoneYc}{min(\FoneYa,\FoneYb)-0.45cm}%
\coordinate (ctop) at (0pt,\FoneYc);
\begin{scope}[shift={(ctop)}]
  \node[ttl, anchor=north west] at (0,0) {(c) Pretraining and evaluation};
  % what every model of one seed shares
  \node[stage, text width=2.15cm] (seed) at (1.175,-1.96)
    {\textbf{Same for every model of one seed:} the training jets, their order and the
     initial weights, except the output layer};
  % one pretraining per label set: the configuration colours of every later figure
  \node[draw=black!45, fill=black!4, line width=0.6pt, rounded corners=3pt, minimum width=5.15cm,
        minimum height=3.08cm, anchor=north west, inner sep=0pt] (pre) at (2.81,-0.42) {};
  \node[anchor=north west, font=\FoneBold, inner sep=0pt] at (2.96,-0.52) {Pretraining on JetClass-II};
  \node[chip=c188, minimum width=1.13cm, minimum height=0.42cm] at (2.96,-0.95) {\Kfull};
  \node[chip=c162, minimum width=1.13cm, minimum height=0.42cm] at (4.20,-0.95) {\Kres};
  \node[chip=c43,  minimum width=1.13cm, minimum height=0.42cm] at (5.44,-0.95) {\Kmid};
  \node[chip=c17,  minimum width=1.13cm, minimum height=0.42cm] at (6.68,-0.95) {\Kcoarse};
  \node[chip=crand, minimum width=1.70cm, minimum height=0.70cm] at (2.96,-1.47) {random\\partitions};
  \node[chip=cflav, minimum width=3.04cm, minimum height=0.70cm] at (4.77,-1.47) {controlled flavor\\boundaries};
  \node[chip=c162, densely dotted, minimum width=1.13cm, minimum height=0.42cm] at (2.96,-2.32) {\Kres};
  \node[chip=c17,  densely dotted, minimum width=1.13cm, minimum height=0.42cm] at (4.20,-2.32) {\Kcoarse};
  \node[txt, anchor=west] at (5.44,-2.53) {with auxiliary\\mass regression};
  \node[chip=c188, densely dotted, minimum width=1.13cm, minimum height=0.42cm] at (2.96,-2.94) {\Kfull};
  \node[txt, anchor=west] at (4.20,-3.15) {kinematic inputs only};
  % evaluation
  \node[stage, text width=2.4cm] (fz) at (9.68,-0.95) {frozen features:\\linear probe,\\MLP probe};
  \node[stage, text width=2.4cm] (ft) at (9.68,-2.05) {fine-tuning\\of all weights};
  \node[stage, text width=2.4cm, draw=black, fill=white, line width=1.1pt] (sc) at (9.68,-3.15) {trained from\\scratch};
  \node[stage, text width=3.4cm, minimum height=3.2cm] (task) at (13.20,-2.00)
    {\textbf{Downstream tasks}\\JetClass-II: all classes\\JetClass-II: class pairs\\\quad(both class-balanced)\\JetClass\\
     top tagging\\quark/gluon, Pythia\\quark/gluon, Herwig\\jet mass\FoneXbx};
  \draw[flow] (seed.east) -- (seed.east -| pre.west);
  \draw[flow] (fz.west -| pre.east) -- (fz.west);
  \draw[flow] (ft.west -| pre.east) -- (ft.west);
  \draw[flow] (fz.east) -- (fz.east -| task.west);
  \draw[flow] (ft.east) -- (ft.east -| task.west);
  \draw[flow] (sc.east) -- (sc.east -| task.west);
\end{scope}
\end{tikzpicture}
\caption{Design of the study.
(a)~The class hierarchy on the decays of a heavy resonance~$X$ and on the jets from quantum
chromodynamics (QCD); $q$ is a $u$ or $d$ quark, $\ell$ an electron or a muon, and $\tau_h$,
$\tau_e$, $\tau_\mu$ a $\tau$ lepton decaying to hadrons, an electron or a muon. Each box is one
class of the label set above its column, except the QCD box of the \Kfull-class column (Sophon's
label set~\cite{sophon}), which stands for \LitJetClassTwoNQcdClasses{} classes. Each level merges
classes of the level to its left; the lines below name what each merge removes.
(b)~The controls: the first random partition, where boxes with the same letter are one class, and
the random partitions with controlled flavor boundaries on the decays \RandFlavSplitOrbit{} (one
square per decay; filled squares are moved to a second class).
(c)~Models of one seed share the training jets, their order and the initial weights, except the
output layer, the mass regression target, and the first layer and initial weights of the model with
kinematic inputs only (the particle four-momenta). The hierarchy levels have \MlstNRunsWord{} seeds
each, the models with the auxiliary mass regression \MlstMassNRunsWord{}, the random partitions with
controlled flavor boundaries \RandFlavRuns{} and each random partition \RandVtwoRuns{}. Frozen
features, the output of the backbone (the network below the classification head) with its weights held
fixed, feed linear probes (logistic regressions) and multilayer perceptron (MLP) probes; fine-tuning
trains all weights. Colors identify the label sets and controls in every figure; the plots give the three partitions with
controlled flavor boundaries three shades; dotted borders mark the two
variants, which the other figures draw as dotted lines in the color of their label set.}
\label{fig:design}
\end{figure}
 removes X->bb/X->bs from panel (c)
% (PI 2026-11-02: if the samples are not staged by 12 Nov, the X->bb/X->bs section goes).
%
% Every class box is drawn from the label maps (checked by hand, 2026-11-02):
%   (a) configs/labelmaps/rung_label_maps.v1.csv, columns L188, L162, R42_Q1, R16_Q1. Rows: the ten
%       decays X->QQ/gg in the order bb bc bq | cc cs cq | ss sq qq gg (the 43-class b / c / light
%       groups are then contiguous), then ee mm | tauh-taue tauh-taum, then the 27 QCD classes.
%   (b) column: configs/labelmaps/rand_label_map.v2.csv, RAND2_p1 on the same rows: groups
%       0 0 0 0 0 9 9 9 0 3 | 11 11 | 15 4 | QCD, lettered A..F in order of first appearance.
%       strips: configs/labelmaps/flavour_pair_map.v2.csv, the 22 classes of orbit X_YY_QQQQ.
%       Order with a b quark: bbbb bbcc bbss bbqq qqqb bbcq ccbs ccbq qqbc qqbs bcsq;
%       without: cccc ccss ccqq ssss ssqq qqqq qqqc qqqs ccsq sscq qqcs.
%       FLAV_F0: all 22 in group 9. FLAV_F1: the 11 with a b in group 11, the rest 9.
%       FLAV_F1R: group 11 = bbss bbqq qqqb bbcq ccbq qqbc (with b, positions 2-5, 7, 8 from 0)
%       and ccss ccqq ssss qqqc sscq (without, positions 1-3, 6, 9); filled square = group 11.
% Geometry (cm): row height 0.35 from y = -1.18, gaps of 0.14 between the quark/gluon, lepton
% and QCD blocks, the QCD row 0.70; boxes inset 0.025 top and bottom; columns 1.45 wide with
% left edges 1.75, 3.28, 4.81, 6.34 (hierarchy) and 8.25 (random partition).
\providecommand{\Kfull}{\VocabSizeLoneeighteight}
\providecommand{\Kres}{\VocabSizeLonesixtwo}
\providecommand{\Kmid}{\VocabSizeRfourtwoQone}
\providecommand{\Kcoarse}{\VocabSizeRonesixQone}
\begin{figure}[tbp]
\centering
\ifdim\linewidth<14.99cm\typeout{fig_f1.tex: WARNING: the drawing is 15 cm wide but \string\linewidth\space
  is \the\linewidth; set a 15 cm text width (it is not scaled, to keep its 8.5 pt text)}\fi
\def\FoneText{\rmfamily\fontsize{8.5}{9.5}\selectfont}%
\def\FoneMain{\rmfamily\fontsize{9}{10.5}\selectfont}%
\def\FoneBold{\rmfamily\bfseries\fontsize{9}{10.5}\selectfont}%
\def\FoneTitle{\rmfamily\bfseries\fontsize{10}{12}\selectfont}%
\ifdefined\FoneDropXbx\def\FoneXbx{}\else\def\FoneXbx{\\$X\to b\bar b$, $X\to bs$}\fi
\begin{tikzpicture}[
  x=1cm, y=1cm,
  font=\FoneMain,
  ttl/.style={anchor=base west, font=\FoneTitle, inner sep=0pt},
  cls/.style={draw=black!55, line width=0.45pt, rounded corners=1.2pt},
  hdr/.style={draw=#1!70!black, fill=#1!22, line width=0.6pt, rounded corners=2.5pt, align=center,
              font=\FoneText, inner sep=0pt},
  inbox/.style={align=center, font=\FoneText, inner sep=0pt},
  letter/.style={font=\FoneBold, inner sep=0pt},
  rowlab/.style={anchor=east, inner sep=0pt},
  txt/.style={align=left, font=\FoneText, inner sep=0pt},
  sq/.style={draw=black!65, line width=0.4pt},
  chip/.style={draw=#1!70!black, fill=#1!22, line width=0.6pt, rounded corners=3pt,
               align=center, font=\FoneText, inner sep=0pt, anchor=north west},
  stage/.style={draw=black!45, fill=black!4, line width=0.6pt, rounded corners=3pt, align=left,
                font=\FoneText, inner sep=0.1cm},
  flow/.style={-{Latex[length=1.8mm, width=1.4mm]}, line width=0.7pt, black!75},
]
% ================================================================ (a) the class hierarchy
\node[ttl] at (0,0) {(a) Class hierarchy};
% faint bands on every other row, drawn first so the boxes sit on top
\foreach \rowtop in {-1.53,-2.23,-2.93,-3.63,-4.33,-5.17,-5.87}
  \fill[black!4] (0,\rowtop) rectangle (9.70,\rowtop-0.35);
% column headers
\node[rowlab, align=right, font=\FoneText] at (1.65,-0.68) {JetClass-II\\class};
\node[hdr=c188, minimum width=1.45cm, minimum height=0.80cm, anchor=north west] at (1.75,-0.28) {\Kfull\\classes};
\node[hdr=c162, minimum width=1.45cm, minimum height=0.80cm, anchor=north west] at (3.28,-0.28) {\Kres\\classes};
\node[hdr=c43,  minimum width=1.45cm, minimum height=0.80cm, anchor=north west] at (4.81,-0.28) {\Kmid\\classes};
\node[hdr=c17,  minimum width=1.45cm, minimum height=0.80cm, anchor=north west] at (6.34,-0.28) {\Kcoarse\\classes};
% row labels
\node[rowlab] at (1.65,-1.355) {$X\to b\bar b$};
\node[rowlab] at (1.65,-1.705) {$X\to bc$};
\node[rowlab] at (1.65,-2.055) {$X\to bq$};
\node[rowlab] at (1.65,-2.405) {$X\to c\bar c$};
\node[rowlab] at (1.65,-2.755) {$X\to cs$};
\node[rowlab] at (1.65,-3.105) {$X\to cq$};
\node[rowlab] at (1.65,-3.455) {$X\to s\bar s$};
\node[rowlab] at (1.65,-3.805) {$X\to sq$};
\node[rowlab] at (1.65,-4.155) {$X\to q\bar q$};
\node[rowlab] at (1.65,-4.505) {$X\to gg$};
\node[rowlab] at (1.65,-4.995) {$X\to ee$};
\node[rowlab] at (1.65,-5.345) {$X\to\mu\mu$};
\node[rowlab] at (1.65,-5.695) {$X\to\tau_h\tau_e$};
\node[rowlab] at (1.65,-6.045) {$X\to\tau_h\tau_\mu$};
\node[rowlab] at (1.65,-6.710) {QCD jets};
% 188 and 162 classes: every resonant decay is its own class
\foreach \rowtop in {-1.18,-1.53,-1.88,-2.23,-2.58,-2.93,-3.28,-3.63,-3.98,-4.33,-4.82,-5.17,-5.52,-5.87} {
  \draw[cls, fill=c188!15] (1.75,\rowtop-0.025) rectangle (3.20,\rowtop-0.325);
  \draw[cls, fill=c162!15] (3.28,\rowtop-0.025) rectangle (4.73,\rowtop-0.325);
}
\draw[cls, fill=c188!15] (1.75,-6.385) rectangle (3.20,-7.035);
\node[inbox] at (2.475,-6.71) {\LitJetClassTwoNQcdClasses\\classes};
\draw[cls, fill=c162!15] (3.28,-6.385) rectangle (4.73,-7.035);
\node[inbox] at (4.005,-6.71) {1 class};
% 43 classes: decay topology x {b, else c, else light quark or gluon}; e and mu merged
\draw[cls, fill=c43!15] (4.81,-1.205) rectangle (6.26,-2.205);
\node[inbox] at (5.535,-1.705) {with $b$};
\draw[cls, fill=c43!15] (4.81,-2.255) rectangle (6.26,-3.255);
\node[inbox] at (5.535,-2.755) {with $c$,\\no $b$};
\draw[cls, fill=c43!15] (4.81,-3.305) rectangle (6.26,-4.655);
\node[inbox] at (5.535,-3.980) {$u$, $d$, $s$\\or $g$};
\draw[cls, fill=c43!15] (4.81,-4.845) rectangle (6.26,-5.495);
\node[inbox] at (5.535,-5.170) {$\ell\ell$};
\draw[cls, fill=c43!15] (4.81,-5.545) rectangle (6.26,-6.195);
\node[inbox] at (5.535,-5.870) {$\tau_h\tau_\ell$};
\draw[cls, fill=c43!15] (4.81,-6.385) rectangle (6.26,-7.035);
\node[inbox] at (5.535,-6.710) {1 class};
% 17 classes: number and type of prongs, no flavor
\draw[cls, fill=c17!15] (6.34,-1.205) rectangle (7.79,-4.655);
\node[inbox] at (7.065,-2.930) {two\\quarks\\or gluons};
\draw[cls, fill=c17!15] (6.34,-4.845) rectangle (7.79,-5.495);
\node[inbox] at (7.065,-5.170) {$\ell\ell$};
\draw[cls, fill=c17!15] (6.34,-5.545) rectangle (7.79,-6.195);
\node[inbox] at (7.065,-5.870) {$\tau_h\tau_\ell$};
\draw[cls, fill=c17!15] (6.34,-6.385) rectangle (7.79,-7.035);
\node[inbox] at (7.065,-6.710) {1 class};
% what each merge removes
\node[anchor=north west, font=\FoneBold, inner sep=0pt] at (0,-7.27) {What each merge removes};
\node[txt, anchor=north west, text width=6.04cm] (m1) at (1.75,-7.65)
  {the $b$, $c$ and $s$ content of the QCD jets};
\node[txt, anchor=north west, text width=6.04cm] (m2) at ([yshift=-0.10cm]m1.south west)
  {the number of $b$ and $c$ quarks, light-quark flavor, quark against gluon, and lepton
   flavor ($e$ against $\mu$)};
\node[txt, anchor=north west, text width=6.04cm] (m3) at ([yshift=-0.10cm]m2.south west)
  {quark flavor ($b$, $c$ or light); the number and type of prongs remain};
\node[txt, anchor=north east] at ([xshift=-0.10cm]m1.north west) {\Kfull\,$\to$\,\Kres};
\node[txt, anchor=north east] at ([xshift=-0.10cm]m2.north west) {\Kres\,$\to$\,\Kmid};
\node[txt, anchor=north east] at ([xshift=-0.10cm]m3.north west) {\Kmid\,$\to$\,\Kcoarse};
% ================================================================ (b) the controls
\node[ttl] at (8.25,0) {(b) Controls, \Kcoarse{} classes each};
% the first random partition on the rows of (a); one letter = one class
\node[hdr=crand, minimum width=1.45cm, minimum height=0.80cm, anchor=north west] at (8.25,-0.28) {random\\partition};
\draw[cls, fill=crand!15] (8.25,-1.205) rectangle (9.70,-2.905);
\node[letter] at (8.975,-2.055) {A};
\draw[cls, fill=crand!15] (8.25,-2.955) rectangle (9.70,-3.955);
\node[letter] at (8.975,-3.455) {B};
\draw[cls, fill=crand!15] (8.25,-4.005) rectangle (9.70,-4.305);
\node[letter] at (8.975,-4.155) {A};
\draw[cls, fill=crand!15] (8.25,-4.355) rectangle (9.70,-4.655);
\node[letter] at (8.975,-4.505) {C};
\draw[cls, fill=crand!15] (8.25,-4.845) rectangle (9.70,-5.495);
\node[letter] at (8.975,-5.170) {D};
\draw[cls, fill=crand!15] (8.25,-5.545) rectangle (9.70,-5.845);
\node[letter] at (8.975,-5.695) {E};
\draw[cls, fill=crand!15] (8.25,-5.895) rectangle (9.70,-6.195);
\node[letter] at (8.975,-6.045) {F};
\draw[cls, fill=crand!15] (8.25,-6.385) rectangle (9.70,-7.035);
\node[inbox] at (8.975,-6.710) {1 class};
% how the random partitions are built
\node[hdr=crand, anchor=north west, text width=4.75cm, inner sep=0.1cm, align=left, font=\FoneBold]
  (rh) at (10.05,-0.28) {Random partitions (\RandVtwoNPartitions)};
\node[txt, anchor=north west, text width=4.95cm] (rb) at ([yshift=-0.12cm]rh.south west)
  {\textbullet\ groups of resonant classes drawn at random, plus one class of QCD jets\\[2pt]
   \textbullet\ group shares of training jets within \RandVtwoShareTol{} of the \Kcoarse-class groups\\[2pt]
   \textbullet\ best of \RandVtwoTrials{} draws: the fewest merges shared with the \Kcoarse-class level\\[2pt]
   \textbullet\ left: the first partition; boxes with the same letter are one class};
% how the random partitions with controlled flavor boundaries are built
\node[hdr=cflav, anchor=north west, text width=4.75cm, inner sep=0.1cm, align=left, font=\FoneBold]
  (fh) at ([yshift=-0.22cm]rb.south west) {Random partitions with controlled flavor boundaries};
\node[txt, anchor=north west, text width=4.95cm] (fn) at ([yshift=-0.12cm]fh.south west)
  {whole families of decays that differ only in quark flavor; two of them split one family,
   the \RandFlavSplitOrbitClasses{} decays \RandFlavSplitOrbit:};
\begin{scope}[shift={(fn.south west)}]
  % brackets over the decays with and without a b quark (squares at a 0.2 cm pitch)
  \node[txt, anchor=south] at (1.075,-0.45) {with $b$};
  \node[txt, anchor=south] at (3.425,-0.45) {without $b$};
  \draw[line width=0.5pt] (0,-0.54) -- (0,-0.47) -- (2.15,-0.47) -- (2.15,-0.54);
  \draw[line width=0.5pt] (2.35,-0.54) -- (2.35,-0.47) -- (4.50,-0.47) -- (4.50,-0.54);
  % no flavor boundary: all 22 in one class
  \node[txt, anchor=north west] at (0,-0.62) {no flavor boundary};
  \foreach \i in {0,...,10} {
    \draw[sq, fill=white] (0.2*\i,-1.10) rectangle ++(0.15,0.15);
    \draw[sq, fill=white] (2.35+0.2*\i,-1.10) rectangle ++(0.15,0.15);
  }
  % b-quark boundary: the decays with a b quark moved to another class
  \node[txt, anchor=north west] at (0,-1.22) {$b$-quark boundary};
  \foreach \i in {0,...,10} {
    \draw[sq, fill=cflav!70!black] (0.2*\i,-1.70) rectangle ++(0.15,0.15);
    \draw[sq, fill=white] (2.35+0.2*\i,-1.70) rectangle ++(0.15,0.15);
  }
  % random flavor boundary: the same number of decays, drawn at random
  \node[txt, anchor=north west] at (0,-1.82) {random flavor boundary};
  \foreach \i in {0,1,6,9,10}
    \draw[sq, fill=white] (0.2*\i,-2.30) rectangle ++(0.15,0.15);
  \foreach \i in {2,3,4,5,7,8}
    \draw[sq, fill=cflav!70!black] (0.2*\i,-2.30) rectangle ++(0.15,0.15);
  \foreach \i in {0,4,5,7,8,10}
    \draw[sq, fill=white] (2.35+0.2*\i,-2.30) rectangle ++(0.15,0.15);
  \foreach \i in {1,2,3,6,9}
    \draw[sq, fill=cflav!70!black] (2.35+0.2*\i,-2.30) rectangle ++(0.15,0.15);
  % key
  \draw[sq, fill=cflav!70!black] (0,-2.62) rectangle ++(0.15,0.15);
  \node[txt, anchor=west] (fkey) at (0.25,-2.545) {moved to another class};
\end{scope}
% ================================================================ (c) the pipeline
% (c) starts below the lower of the two columns above, whatever their text heights
\pgfpointanchor{m3}{south}\pgfgetlastxy{\FoneXa}{\FoneYa}%
\pgfpointanchor{fkey}{south}\pgfgetlastxy{\FoneXb}{\FoneYb}%
\pgfmathsetlengthmacro{\FoneYc}{min(\FoneYa,\FoneYb)-0.45cm}%
\coordinate (ctop) at (0pt,\FoneYc);
\begin{scope}[shift={(ctop)}]
  \node[ttl, anchor=north west] at (0,0) {(c) Pretraining and evaluation};
  % what every model of one seed shares
  \node[stage, text width=2.15cm] (seed) at (1.175,-1.96)
    {\textbf{Same for every model of one seed:} the training jets, their order and the
     initial weights, except the output layer};
  % one pretraining per label set: the configuration colours of every later figure
  \node[draw=black!45, fill=black!4, line width=0.6pt, rounded corners=3pt, minimum width=5.15cm,
        minimum height=3.08cm, anchor=north west, inner sep=0pt] (pre) at (2.81,-0.42) {};
  \node[anchor=north west, font=\FoneBold, inner sep=0pt] at (2.96,-0.52) {Pretraining on JetClass-II};
  \node[chip=c188, minimum width=1.13cm, minimum height=0.42cm] at (2.96,-0.95) {\Kfull};
  \node[chip=c162, minimum width=1.13cm, minimum height=0.42cm] at (4.20,-0.95) {\Kres};
  \node[chip=c43,  minimum width=1.13cm, minimum height=0.42cm] at (5.44,-0.95) {\Kmid};
  \node[chip=c17,  minimum width=1.13cm, minimum height=0.42cm] at (6.68,-0.95) {\Kcoarse};
  \node[chip=crand, minimum width=1.70cm, minimum height=0.70cm] at (2.96,-1.47) {random\\partitions};
  \node[chip=cflav, minimum width=3.04cm, minimum height=0.70cm] at (4.77,-1.47) {controlled flavor\\boundaries};
  \node[chip=c162, densely dotted, minimum width=1.13cm, minimum height=0.42cm] at (2.96,-2.32) {\Kres};
  \node[chip=c17,  densely dotted, minimum width=1.13cm, minimum height=0.42cm] at (4.20,-2.32) {\Kcoarse};
  \node[txt, anchor=west] at (5.44,-2.53) {with auxiliary\\mass regression};
  \node[chip=c188, densely dotted, minimum width=1.13cm, minimum height=0.42cm] at (2.96,-2.94) {\Kfull};
  \node[txt, anchor=west] at (4.20,-3.15) {kinematic inputs only};
  % evaluation
  \node[stage, text width=2.4cm] (fz) at (9.68,-0.95) {frozen features:\\linear probe,\\MLP probe};
  \node[stage, text width=2.4cm] (ft) at (9.68,-2.05) {fine-tuning\\of all weights};
  \node[stage, text width=2.4cm, draw=black, fill=white, line width=1.1pt] (sc) at (9.68,-3.15) {trained from\\scratch};
  \node[stage, text width=3.4cm, minimum height=3.2cm] (task) at (13.20,-2.00)
    {\textbf{Downstream tasks}\\JetClass-II: all classes\\JetClass-II: class pairs\\\quad(both class-balanced)\\JetClass\\
     top tagging\\quark/gluon, Pythia\\quark/gluon, Herwig\\jet mass\FoneXbx};
  \draw[flow] (seed.east) -- (seed.east -| pre.west);
  \draw[flow] (fz.west -| pre.east) -- (fz.west);
  \draw[flow] (ft.west -| pre.east) -- (ft.west);
  \draw[flow] (fz.east) -- (fz.east -| task.west);
  \draw[flow] (ft.east) -- (ft.east -| task.west);
  \draw[flow] (sc.east) -- (sc.east -| task.west);
\end{scope}
\end{tikzpicture}
\caption{Design of the study.
(a)~The class hierarchy on the decays of a heavy resonance~$X$ and on the jets from quantum
chromodynamics (QCD); $q$ is a $u$ or $d$ quark, $\ell$ an electron or a muon, and $\tau_h$,
$\tau_e$, $\tau_\mu$ a $\tau$ lepton decaying to hadrons, an electron or a muon. Each box is one
class of the label set above its column, except the QCD box of the \Kfull-class column (Sophon's
label set~\cite{sophon}), which stands for \LitJetClassTwoNQcdClasses{} classes. Each level merges
classes of the level to its left; the lines below name what each merge removes.
(b)~The controls: the first random partition, where boxes with the same letter are one class, and
the random partitions with controlled flavor boundaries on the decays \RandFlavSplitOrbit{} (one
square per decay; filled squares are moved to a second class).
(c)~Models of one seed share the training jets, their order and the initial weights, except the
output layer, the mass regression target, and the first layer and initial weights of the model with
kinematic inputs only (the particle four-momenta). The hierarchy levels have \MlstNRunsWord{} seeds
each, the models with the auxiliary mass regression \MlstMassNRunsWord{}, the random partitions with
controlled flavor boundaries \RandFlavRuns{} and each random partition \RandVtwoRuns{}. Frozen
features, the output of the backbone (the network below the classification head) with its weights held
fixed, feed linear probes (logistic regressions) and multilayer perceptron (MLP) probes; fine-tuning
trains all weights. Colors identify the label sets and controls in every figure; the plots give the three partitions with
controlled flavor boundaries three shades; dotted borders mark the two
variants, which the other figures draw as dotted lines in the color of their label set.}
\label{fig:design}
\end{figure}
 removes X->bb/X->bs from panel (c)
% (PI 2026-11-02: if the samples are not staged by 12 Nov, the X->bb/X->bs section goes).
%
% Every class box is drawn from the label maps (checked by hand, 2026-11-02):
%   (a) configs/labelmaps/rung_label_maps.v1.csv, columns L188, L162, R42_Q1, R16_Q1. Rows: the ten
%       decays X->QQ/gg in the order bb bc bq | cc cs cq | ss sq qq gg (the 43-class b / c / light
%       groups are then contiguous), then ee mm | tauh-taue tauh-taum, then the 27 QCD classes.
%   (b) column: configs/labelmaps/rand_label_map.v2.csv, RAND2_p1 on the same rows: groups
%       0 0 0 0 0 9 9 9 0 3 | 11 11 | 15 4 | QCD, lettered A..F in order of first appearance.
%       strips: configs/labelmaps/flavour_pair_map.v2.csv, the 22 classes of orbit X_YY_QQQQ.
%       Order with a b quark: bbbb bbcc bbss bbqq qqqb bbcq ccbs ccbq qqbc qqbs bcsq;
%       without: cccc ccss ccqq ssss ssqq qqqq qqqc qqqs ccsq sscq qqcs.
%       FLAV_F0: all 22 in group 9. FLAV_F1: the 11 with a b in group 11, the rest 9.
%       FLAV_F1R: group 11 = bbss bbqq qqqb bbcq ccbq qqbc (with b, positions 2-5, 7, 8 from 0)
%       and ccss ccqq ssss qqqc sscq (without, positions 1-3, 6, 9); filled square = group 11.
% Geometry (cm): row height 0.35 from y = -1.18, gaps of 0.14 between the quark/gluon, lepton
% and QCD blocks, the QCD row 0.70; boxes inset 0.025 top and bottom; columns 1.45 wide with
% left edges 1.75, 3.28, 4.81, 6.34 (hierarchy) and 8.25 (random partition).
\providecommand{\Kfull}{\VocabSizeLoneeighteight}
\providecommand{\Kres}{\VocabSizeLonesixtwo}
\providecommand{\Kmid}{\VocabSizeRfourtwoQone}
\providecommand{\Kcoarse}{\VocabSizeRonesixQone}
\begin{figure}[tbp]
\centering
\ifdim\linewidth<14.99cm\typeout{fig_f1.tex: WARNING: the drawing is 15 cm wide but \string\linewidth\space
  is \the\linewidth; set a 15 cm text width (it is not scaled, to keep its 8.5 pt text)}\fi
\def\FoneText{\rmfamily\fontsize{8.5}{9.5}\selectfont}%
\def\FoneMain{\rmfamily\fontsize{9}{10.5}\selectfont}%
\def\FoneBold{\rmfamily\bfseries\fontsize{9}{10.5}\selectfont}%
\def\FoneTitle{\rmfamily\bfseries\fontsize{10}{12}\selectfont}%
\ifdefined\FoneDropXbx\def\FoneXbx{}\else\def\FoneXbx{\\$X\to b\bar b$, $X\to bs$}\fi
\begin{tikzpicture}[
  x=1cm, y=1cm,
  font=\FoneMain,
  ttl/.style={anchor=base west, font=\FoneTitle, inner sep=0pt},
  cls/.style={draw=black!55, line width=0.45pt, rounded corners=1.2pt},
  hdr/.style={draw=#1!70!black, fill=#1!22, line width=0.6pt, rounded corners=2.5pt, align=center,
              font=\FoneText, inner sep=0pt},
  inbox/.style={align=center, font=\FoneText, inner sep=0pt},
  letter/.style={font=\FoneBold, inner sep=0pt},
  rowlab/.style={anchor=east, inner sep=0pt},
  txt/.style={align=left, font=\FoneText, inner sep=0pt},
  sq/.style={draw=black!65, line width=0.4pt},
  chip/.style={draw=#1!70!black, fill=#1!22, line width=0.6pt, rounded corners=3pt,
               align=center, font=\FoneText, inner sep=0pt, anchor=north west},
  stage/.style={draw=black!45, fill=black!4, line width=0.6pt, rounded corners=3pt, align=left,
                font=\FoneText, inner sep=0.1cm},
  flow/.style={-{Latex[length=1.8mm, width=1.4mm]}, line width=0.7pt, black!75},
]
% ================================================================ (a) the class hierarchy
\node[ttl] at (0,0) {(a) Class hierarchy};
% faint bands on every other row, drawn first so the boxes sit on top
\foreach \rowtop in {-1.53,-2.23,-2.93,-3.63,-4.33,-5.17,-5.87}
  \fill[black!4] (0,\rowtop) rectangle (9.70,\rowtop-0.35);
% column headers
\node[rowlab, align=right, font=\FoneText] at (1.65,-0.68) {JetClass-II\\class};
\node[hdr=c188, minimum width=1.45cm, minimum height=0.80cm, anchor=north west] at (1.75,-0.28) {\Kfull\\classes};
\node[hdr=c162, minimum width=1.45cm, minimum height=0.80cm, anchor=north west] at (3.28,-0.28) {\Kres\\classes};
\node[hdr=c43,  minimum width=1.45cm, minimum height=0.80cm, anchor=north west] at (4.81,-0.28) {\Kmid\\classes};
\node[hdr=c17,  minimum width=1.45cm, minimum height=0.80cm, anchor=north west] at (6.34,-0.28) {\Kcoarse\\classes};
% row labels
\node[rowlab] at (1.65,-1.355) {$X\to b\bar b$};
\node[rowlab] at (1.65,-1.705) {$X\to bc$};
\node[rowlab] at (1.65,-2.055) {$X\to bq$};
\node[rowlab] at (1.65,-2.405) {$X\to c\bar c$};
\node[rowlab] at (1.65,-2.755) {$X\to cs$};
\node[rowlab] at (1.65,-3.105) {$X\to cq$};
\node[rowlab] at (1.65,-3.455) {$X\to s\bar s$};
\node[rowlab] at (1.65,-3.805) {$X\to sq$};
\node[rowlab] at (1.65,-4.155) {$X\to q\bar q$};
\node[rowlab] at (1.65,-4.505) {$X\to gg$};
\node[rowlab] at (1.65,-4.995) {$X\to ee$};
\node[rowlab] at (1.65,-5.345) {$X\to\mu\mu$};
\node[rowlab] at (1.65,-5.695) {$X\to\tau_h\tau_e$};
\node[rowlab] at (1.65,-6.045) {$X\to\tau_h\tau_\mu$};
\node[rowlab] at (1.65,-6.710) {QCD jets};
% 188 and 162 classes: every resonant decay is its own class
\foreach \rowtop in {-1.18,-1.53,-1.88,-2.23,-2.58,-2.93,-3.28,-3.63,-3.98,-4.33,-4.82,-5.17,-5.52,-5.87} {
  \draw[cls, fill=c188!15] (1.75,\rowtop-0.025) rectangle (3.20,\rowtop-0.325);
  \draw[cls, fill=c162!15] (3.28,\rowtop-0.025) rectangle (4.73,\rowtop-0.325);
}
\draw[cls, fill=c188!15] (1.75,-6.385) rectangle (3.20,-7.035);
\node[inbox] at (2.475,-6.71) {\LitJetClassTwoNQcdClasses\\classes};
\draw[cls, fill=c162!15] (3.28,-6.385) rectangle (4.73,-7.035);
\node[inbox] at (4.005,-6.71) {1 class};
% 43 classes: decay topology x {b, else c, else light quark or gluon}; e and mu merged
\draw[cls, fill=c43!15] (4.81,-1.205) rectangle (6.26,-2.205);
\node[inbox] at (5.535,-1.705) {with $b$};
\draw[cls, fill=c43!15] (4.81,-2.255) rectangle (6.26,-3.255);
\node[inbox] at (5.535,-2.755) {with $c$,\\no $b$};
\draw[cls, fill=c43!15] (4.81,-3.305) rectangle (6.26,-4.655);
\node[inbox] at (5.535,-3.980) {$u$, $d$, $s$\\or $g$};
\draw[cls, fill=c43!15] (4.81,-4.845) rectangle (6.26,-5.495);
\node[inbox] at (5.535,-5.170) {$\ell\ell$};
\draw[cls, fill=c43!15] (4.81,-5.545) rectangle (6.26,-6.195);
\node[inbox] at (5.535,-5.870) {$\tau_h\tau_\ell$};
\draw[cls, fill=c43!15] (4.81,-6.385) rectangle (6.26,-7.035);
\node[inbox] at (5.535,-6.710) {1 class};
% 17 classes: number and type of prongs, no flavor
\draw[cls, fill=c17!15] (6.34,-1.205) rectangle (7.79,-4.655);
\node[inbox] at (7.065,-2.930) {two\\quarks\\or gluons};
\draw[cls, fill=c17!15] (6.34,-4.845) rectangle (7.79,-5.495);
\node[inbox] at (7.065,-5.170) {$\ell\ell$};
\draw[cls, fill=c17!15] (6.34,-5.545) rectangle (7.79,-6.195);
\node[inbox] at (7.065,-5.870) {$\tau_h\tau_\ell$};
\draw[cls, fill=c17!15] (6.34,-6.385) rectangle (7.79,-7.035);
\node[inbox] at (7.065,-6.710) {1 class};
% what each merge removes
\node[anchor=north west, font=\FoneBold, inner sep=0pt] at (0,-7.27) {What each merge removes};
\node[txt, anchor=north west, text width=6.04cm] (m1) at (1.75,-7.65)
  {the $b$, $c$ and $s$ content of the QCD jets};
\node[txt, anchor=north west, text width=6.04cm] (m2) at ([yshift=-0.10cm]m1.south west)
  {the number of $b$ and $c$ quarks, light-quark flavor, quark against gluon, and lepton
   flavor ($e$ against $\mu$)};
\node[txt, anchor=north west, text width=6.04cm] (m3) at ([yshift=-0.10cm]m2.south west)
  {quark flavor ($b$, $c$ or light); the number and type of prongs remain};
\node[txt, anchor=north east] at ([xshift=-0.10cm]m1.north west) {\Kfull\,$\to$\,\Kres};
\node[txt, anchor=north east] at ([xshift=-0.10cm]m2.north west) {\Kres\,$\to$\,\Kmid};
\node[txt, anchor=north east] at ([xshift=-0.10cm]m3.north west) {\Kmid\,$\to$\,\Kcoarse};
% ================================================================ (b) the controls
\node[ttl] at (8.25,0) {(b) Controls, \Kcoarse{} classes each};
% the first random partition on the rows of (a); one letter = one class
\node[hdr=crand, minimum width=1.45cm, minimum height=0.80cm, anchor=north west] at (8.25,-0.28) {random\\partition};
\draw[cls, fill=crand!15] (8.25,-1.205) rectangle (9.70,-2.905);
\node[letter] at (8.975,-2.055) {A};
\draw[cls, fill=crand!15] (8.25,-2.955) rectangle (9.70,-3.955);
\node[letter] at (8.975,-3.455) {B};
\draw[cls, fill=crand!15] (8.25,-4.005) rectangle (9.70,-4.305);
\node[letter] at (8.975,-4.155) {A};
\draw[cls, fill=crand!15] (8.25,-4.355) rectangle (9.70,-4.655);
\node[letter] at (8.975,-4.505) {C};
\draw[cls, fill=crand!15] (8.25,-4.845) rectangle (9.70,-5.495);
\node[letter] at (8.975,-5.170) {D};
\draw[cls, fill=crand!15] (8.25,-5.545) rectangle (9.70,-5.845);
\node[letter] at (8.975,-5.695) {E};
\draw[cls, fill=crand!15] (8.25,-5.895) rectangle (9.70,-6.195);
\node[letter] at (8.975,-6.045) {F};
\draw[cls, fill=crand!15] (8.25,-6.385) rectangle (9.70,-7.035);
\node[inbox] at (8.975,-6.710) {1 class};
% how the random partitions are built
\node[hdr=crand, anchor=north west, text width=4.75cm, inner sep=0.1cm, align=left, font=\FoneBold]
  (rh) at (10.05,-0.28) {Random partitions (\RandVtwoNPartitions)};
\node[txt, anchor=north west, text width=4.95cm] (rb) at ([yshift=-0.12cm]rh.south west)
  {\textbullet\ groups of resonant classes drawn at random, plus one class of QCD jets\\[2pt]
   \textbullet\ group shares of training jets within \RandVtwoShareTol{} of the \Kcoarse-class groups\\[2pt]
   \textbullet\ best of \RandVtwoTrials{} draws: the fewest merges shared with the \Kcoarse-class level\\[2pt]
   \textbullet\ left: the first partition; boxes with the same letter are one class};
% how the random partitions with controlled flavor boundaries are built
\node[hdr=cflav, anchor=north west, text width=4.75cm, inner sep=0.1cm, align=left, font=\FoneBold]
  (fh) at ([yshift=-0.22cm]rb.south west) {Random partitions with controlled flavor boundaries};
\node[txt, anchor=north west, text width=4.95cm] (fn) at ([yshift=-0.12cm]fh.south west)
  {whole families of decays that differ only in quark flavor; two of them split one family,
   the \RandFlavSplitOrbitClasses{} decays \RandFlavSplitOrbit:};
\begin{scope}[shift={(fn.south west)}]
  % brackets over the decays with and without a b quark (squares at a 0.2 cm pitch)
  \node[txt, anchor=south] at (1.075,-0.45) {with $b$};
  \node[txt, anchor=south] at (3.425,-0.45) {without $b$};
  \draw[line width=0.5pt] (0,-0.54) -- (0,-0.47) -- (2.15,-0.47) -- (2.15,-0.54);
  \draw[line width=0.5pt] (2.35,-0.54) -- (2.35,-0.47) -- (4.50,-0.47) -- (4.50,-0.54);
  % no flavor boundary: all 22 in one class
  \node[txt, anchor=north west] at (0,-0.62) {no flavor boundary};
  \foreach \i in {0,...,10} {
    \draw[sq, fill=white] (0.2*\i,-1.10) rectangle ++(0.15,0.15);
    \draw[sq, fill=white] (2.35+0.2*\i,-1.10) rectangle ++(0.15,0.15);
  }
  % b-quark boundary: the decays with a b quark moved to another class
  \node[txt, anchor=north west] at (0,-1.22) {$b$-quark boundary};
  \foreach \i in {0,...,10} {
    \draw[sq, fill=cflav!70!black] (0.2*\i,-1.70) rectangle ++(0.15,0.15);
    \draw[sq, fill=white] (2.35+0.2*\i,-1.70) rectangle ++(0.15,0.15);
  }
  % random flavor boundary: the same number of decays, drawn at random
  \node[txt, anchor=north west] at (0,-1.82) {random flavor boundary};
  \foreach \i in {0,1,6,9,10}
    \draw[sq, fill=white] (0.2*\i,-2.30) rectangle ++(0.15,0.15);
  \foreach \i in {2,3,4,5,7,8}
    \draw[sq, fill=cflav!70!black] (0.2*\i,-2.30) rectangle ++(0.15,0.15);
  \foreach \i in {0,4,5,7,8,10}
    \draw[sq, fill=white] (2.35+0.2*\i,-2.30) rectangle ++(0.15,0.15);
  \foreach \i in {1,2,3,6,9}
    \draw[sq, fill=cflav!70!black] (2.35+0.2*\i,-2.30) rectangle ++(0.15,0.15);
  % key
  \draw[sq, fill=cflav!70!black] (0,-2.62) rectangle ++(0.15,0.15);
  \node[txt, anchor=west] (fkey) at (0.25,-2.545) {moved to another class};
\end{scope}
% ================================================================ (c) the pipeline
% (c) starts below the lower of the two columns above, whatever their text heights
\pgfpointanchor{m3}{south}\pgfgetlastxy{\FoneXa}{\FoneYa}%
\pgfpointanchor{fkey}{south}\pgfgetlastxy{\FoneXb}{\FoneYb}%
\pgfmathsetlengthmacro{\FoneYc}{min(\FoneYa,\FoneYb)-0.45cm}%
\coordinate (ctop) at (0pt,\FoneYc);
\begin{scope}[shift={(ctop)}]
  \node[ttl, anchor=north west] at (0,0) {(c) Pretraining and evaluation};
  % what every model of one seed shares
  \node[stage, text width=2.15cm] (seed) at (1.175,-1.96)
    {\textbf{Same for every model of one seed:} the training jets, their order and the
     initial weights, except the output layer};
  % one pretraining per label set: the configuration colours of every later figure
  \node[draw=black!45, fill=black!4, line width=0.6pt, rounded corners=3pt, minimum width=5.15cm,
        minimum height=3.08cm, anchor=north west, inner sep=0pt] (pre) at (2.81,-0.42) {};
  \node[anchor=north west, font=\FoneBold, inner sep=0pt] at (2.96,-0.52) {Pretraining on JetClass-II};
  \node[chip=c188, minimum width=1.13cm, minimum height=0.42cm] at (2.96,-0.95) {\Kfull};
  \node[chip=c162, minimum width=1.13cm, minimum height=0.42cm] at (4.20,-0.95) {\Kres};
  \node[chip=c43,  minimum width=1.13cm, minimum height=0.42cm] at (5.44,-0.95) {\Kmid};
  \node[chip=c17,  minimum width=1.13cm, minimum height=0.42cm] at (6.68,-0.95) {\Kcoarse};
  \node[chip=crand, minimum width=1.70cm, minimum height=0.70cm] at (2.96,-1.47) {random\\partitions};
  \node[chip=cflav, minimum width=3.04cm, minimum height=0.70cm] at (4.77,-1.47) {controlled flavor\\boundaries};
  \node[chip=c162, densely dotted, minimum width=1.13cm, minimum height=0.42cm] at (2.96,-2.32) {\Kres};
  \node[chip=c17,  densely dotted, minimum width=1.13cm, minimum height=0.42cm] at (4.20,-2.32) {\Kcoarse};
  \node[txt, anchor=west] at (5.44,-2.53) {with auxiliary\\mass regression};
  \node[chip=c188, densely dotted, minimum width=1.13cm, minimum height=0.42cm] at (2.96,-2.94) {\Kfull};
  \node[txt, anchor=west] at (4.20,-3.15) {kinematic inputs only};
  % evaluation
  \node[stage, text width=2.4cm] (fz) at (9.68,-0.95) {frozen features:\\linear probe,\\MLP probe};
  \node[stage, text width=2.4cm] (ft) at (9.68,-2.05) {fine-tuning\\of all weights};
  \node[stage, text width=2.4cm, draw=black, fill=white, line width=1.1pt] (sc) at (9.68,-3.15) {trained from\\scratch};
  \node[stage, text width=3.4cm, minimum height=3.2cm] (task) at (13.20,-2.00)
    {\textbf{Downstream tasks}\\JetClass-II: all classes\\JetClass-II: class pairs\\\quad(both class-balanced)\\JetClass\\
     top tagging\\quark/gluon, Pythia\\quark/gluon, Herwig\\jet mass\FoneXbx};
  \draw[flow] (seed.east) -- (seed.east -| pre.west);
  \draw[flow] (fz.west -| pre.east) -- (fz.west);
  \draw[flow] (ft.west -| pre.east) -- (ft.west);
  \draw[flow] (fz.east) -- (fz.east -| task.west);
  \draw[flow] (ft.east) -- (ft.east -| task.west);
  \draw[flow] (sc.east) -- (sc.east -| task.west);
\end{scope}
\end{tikzpicture}
\caption{Design of the study.
(a)~The class hierarchy on the decays of a heavy resonance~$X$ and on the jets from quantum
chromodynamics (QCD); $q$ is a $u$ or $d$ quark, $\ell$ an electron or a muon, and $\tau_h$,
$\tau_e$, $\tau_\mu$ a $\tau$ lepton decaying to hadrons, an electron or a muon. Each box is one
class of the label set above its column, except the QCD box of the \Kfull-class column (Sophon's
label set~\cite{sophon}), which stands for \LitJetClassTwoNQcdClasses{} classes. Each level merges
classes of the level to its left; the lines below name what each merge removes.
(b)~The controls: the first random partition, where boxes with the same letter are one class, and
the random partitions with controlled flavor boundaries on the decays \RandFlavSplitOrbit{} (one
square per decay; filled squares are moved to a second class).
(c)~Models of one seed share the training jets, their order and the initial weights, except the
output layer, the mass regression target, and the first layer and initial weights of the model with
kinematic inputs only (the particle four-momenta). The hierarchy levels have \MlstNRunsWord{} seeds
each, the models with the auxiliary mass regression \MlstMassNRunsWord{}, the random partitions with
controlled flavor boundaries \RandFlavRuns{} and each random partition \RandVtwoRuns{}. Frozen
features, the output of the backbone (the network below the classification head) with its weights held
fixed, feed linear probes (logistic regressions) and multilayer perceptron (MLP) probes; fine-tuning
trains all weights. Colors identify the label sets and controls in every figure; the plots give the three partitions with
controlled flavor boundaries three shades; dotted borders mark the two
variants, which the other figures draw as dotted lines in the color of their label set.}
\label{fig:design}
\end{figure}
